\documentclass{book}
%\usepackage[letterpaper, portrait, margin=1cm]{geometry}
%\usepackage[letterpaper, bindingoffset=0.2in, left=1in,right=1in,top=.5in,bottom=.5in,footskip=.25in,marginparwidth=5em]{geometry}
\usepackage[letterpaper, left=1in,right=1in,top=.5in,bottom=.5in,footskip=.25in,marginparwidth=1cm]{geometry}
% ---------------------
% mini-table-of-content
% ---------------------
\usepackage{minitoc}
\setcounter{minitocdepth}{1}
\setlength{\mtcindent}{24pt}
\setcounter{secnumdepth}{-2}
%\renewcommand{\mtcfont}{\small\rm}
%\renewcommand{\mtcSfont}{\small\bf}
%\usepackage{setspace}
%\usepackage{tocloft}
%\setlength\cftparskip{-1.2pt}
%\setlength\cftbeforesecskip{1.3pt}
%\setlength\cftaftertoctitleskip{2pt}
%\renewcommand{\cftsecafterpnum}{\hspace*{02.0em}}
%\renewcommand{\cftsubsecafterpnum}{\hspace*{02.0em}}

% ---------------------------
% Chinese Characters Packages
% ---------------------------
\usepackage{fontspec} 
\usepackage{xeCJK}
\setmainfont{Times}
\setCJKmainfont{BiauKai}
\newfontfamily\sblgoodhebrew{SBL BibLit}[Script=Hebrew,Contextuals=Alternate]
\newfontfamily\sblgoodgreek{SBL BibLit}[Script=Greek,Contextuals=Alternate]

\usepackage{ifpdf,cite,algorithmic,url,tikz}
\usepackage[cmex10]{amsmath}

% ---------------------------
% Hebrew Characters Packages
% ---------------------------
\usepackage{polyglossia}
\setmainfont{Times New Roman}

% -------
% General
% -------
\usepackage{multicol}
\usepackage{multirow}
\usepackage{color,colortbl}
\usepackage{xparse}
\usepackage{pbox}
\usepackage{stackengine}
\usepackage{titlesec}% http://ctan.org/pkg/titlesec
\usepackage{tabularx}
\usepackage{xltabular}
\usepackage{titlesec}
\usepackage{makecell}
\newcommand{\sectionbreak}{\clearpage}

\author{
  Editor, Michael Chan\\
  \texttt{michaelchan\_wahyan@yahoo.com.hk}
}
\usepackage{tocloft}

\usepackage{hyperref}
\hypersetup{
    colorlinks=true, % set true if you want colored links
    linktoc   =all , % set to all if you want both sections and subsections linked
    linkcolor =blue, % choose some color if you want links to stand out
}

% ----------
% Afterword
% ----------
\usepackage{marginnote}
\usepackage{sectsty}
\usepackage{ragged2e}
\usepackage{lineno}
\usepackage{xcolor}
\usepackage{paracol}

\begin{document}

\clearpage
%% temporary titles
% command to provide stretchy vertical space in proportion
\newcommand\nbvspace[1][3]{\vspace*{\stretch{#1}}}
% allow some slack to avoid under/overfull boxes
\newcommand\nbstretchyspace{\spaceskip0.5em plus 0.25em minus 0.25em}
% To improve spacing on titlepages
\newcommand{\nbtitlestretch}{\spaceskip0.6em}
\pagestyle{empty}
\begin{center}
\bfseries
\nbvspace[1]
\Huge
{%\nbtitlestretch
\Large
\textbf{華人教會網絡 粵語講道逐字稿 第19卷 \\
       church.com.hk
       }}

\nbvspace[1]

{\large
Editor: Michael\\
\texttt{michaelchan\_wahyan@yahoo.com.hk}
}

\nbvspace[1]

{\large
Revision: \texttt{v1.3}\\
Last Update: \today
}


\vfill
\begin{tikzpicture}
    %remove comment for OT cover%\node (0,0) [opacity=0.03]{\includegraphics[width=15cm]{../bible_out/ot_frontcover.png}} ;
    %remove comment for NT cover%\node (0,0) [opacity=0.03]{\includegraphics[width=15cm]{../bible_out/christ_on_cross.png}} ;
    %remove comment for Bible cover%\node (0,0) [xshift=0.8cm, yshift=+2cm, opacity=0.03]{\includegraphics[width=10cm]{./christ_on_cross.png}} ;
    %remove comment for Bible cover%\node (0,0) [              yshift=-2cm, opacity=0.03]{\includegraphics[width=14cm]{./ot_frontcover.png}} ;
\end{tikzpicture}
\vfill

\end{center}

\newpage

\setcounter{tocdepth}{0}
\dominitoc
\begin{multicols}{3}
\addtocontents{toc}{\protect\hypertarget{toc}{}}
\tableofcontents
\end{multicols}

\large
%\twocolumn

% the color definition syntax is as follow:
% \definecolor{name}{system}{definition}
% example: a mono-channel color can be defined as
%          \definecolor{Gray}{gray}{0.9}
% example: an rgb-3-channel color can be defined as
%          \definecolor{LightCyan}{rgb}{0.88,1,1}
%          \definecolor{pink}{rgb}{0.68,0,0.68}

\definecolor{CUV1LightRed}{rgb}{1,0.75,0.75}     % for CUV1
\definecolor{LZZVLightGray}{rgb}{0.9,0.9,0.9}    % for LZZ
\definecolor{KJVVLightGreen}{rgb}{0.75,1,0.85}   % for KJV
\definecolor{CUV2LightYellow}{rgb}{1,1,0.75}     % for CUV2
\definecolor{CNVVLightBrown}{rgb}{1,0.85,0.7}    % for CNV
\definecolor{NRSVLightBlue}{rgb}{0.75,1,1}       % for NRSV
\definecolor{WENLLightPurple}{rgb}{0.95,0.85,0.9}% for WENL
\definecolor{TCV19PaleGreen}{rgb}{0.85,1,0.95}   % for TCV19
\definecolor{MSGVLightWhite}{rgb}{0.98,0.98,0.98}% for MSGV
\definecolor{NETSLightRed}{rgb}{1,0.75,0.75}     % for NETS
\definecolor{JPS1917LightYellow}{rgb}{1,1,0.75}  % for JPS1917
\definecolor{SBLGNTPaleRed}{rgb}{1,0.85,0.80}    % for SBLGNT

\section{目錄\small{(順時)}}
\label{sec:index_chronic}
{ \scriptsize


\begin{xltabular}{\textwidth}{|p{0.15\textwidth} p{0.15\textwidth} p{0.45\textwidth}|p{0.07\textwidth} p{0.1\textwidth}|}
\hline
傳道書 12:1-8 & 官志達牧師 & \hyperref[sec:031601]{不要等} & 2026-02-01 & \href{http://www.church.com.hk/acms/content.asp?site=cdc\&op=show\&type=product\&code=031601\&layout=sermon}{\texttt{031601}} \\
士師記 17:1-18:6 & 潘建堂牧師 & \hyperref[sec:031602]{迷失世代 何去何從?} & 2025-07-13 & \href{http://www.church.com.hk/acms/content.asp?site=cdc\&op=show\&type=product\&code=031602\&layout=sermon}{\texttt{031602}} \\
哥林多後書 8:1-10:18 & 劉少康牧師 & \hyperref[sec:031603]{哥林多前後書 (5) 學習奉獻的功課} & 2026-02-04 & \href{http://www.church.com.hk/acms/content.asp?site=cdc\&op=show\&type=product\&code=031603\&layout=sermon}{\texttt{031603}} \\
歌羅西書 1:3-8 & 黃紹權牧師 & \hyperref[sec:031604]{明白真正神委任的牧者心腸} & 2026-02-08 & \href{http://www.church.com.hk/acms/content.asp?site=cdc\&op=show\&type=product\&code=031604\&layout=sermon}{\texttt{031604}} \\
羅馬書 12:1-2 & 吳磐德牧師 & \hyperref[sec:031605]{建立健康的使命教會 (4) 活出敬拜人生} & 2026-02-08 & \href{http://www.church.com.hk/acms/content.asp?site=cdc\&op=show\&type=product\&code=031605\&layout=sermon}{\texttt{031605}} \\
使徒行傳 9:32-43 & 禤凱榮牧師 & \hyperref[sec:031606]{耶穌戰勝死亡} & 2026-02-08 & \href{http://www.church.com.hk/acms/content.asp?site=cdc\&op=show\&type=product\&code=031606\&layout=sermon}{\texttt{031606}} \\
  & 黃克勤牧師博士 & \hyperref[sec:031607]{聖經人聖經事:聖靈與火} & 2026-01-28 & \href{http://www.church.com.hk/acms/content.asp?site=cdc\&op=show\&type=product\&code=031607\&layout=sermon}{\texttt{031607}} \\
哥林多前書 2:1-5 & 林玉英牧師 & \hyperref[sec:031608]{非靠口才,乃靠上帝的大能} & 2026-02-08 & \href{http://www.church.com.hk/acms/content.asp?site=cdc\&op=show\&type=product\&code=031608\&layout=sermon}{\texttt{031608}} \\
約翰福音 13:1-5 & 陳榮超博士 & \hyperref[sec:031609]{耶穌洗門徒的腳:魔鬼在細節裡} & 2026-02-08 & \href{http://www.church.com.hk/acms/content.asp?site=cdc\&op=show\&type=product\&code=031609\&layout=sermon}{\texttt{031609}} \\
馬太福音 6:9-15 & 石偉漢牧師 & \hyperref[sec:031610]{求主教導我們禱告} & 2026-02-08 & \href{http://www.church.com.hk/acms/content.asp?site=cdc\&op=show\&type=product\&code=031610\&layout=sermon}{\texttt{031610}} \\
約翰一書 1:1-2:2 & 鄧榮仲傳道 & \hyperref[sec:031611]{如何活出合神心意的團契生活} & 2026-02-08 & \href{http://www.church.com.hk/acms/content.asp?site=cdc\&op=show\&type=product\&code=031611\&layout=sermon}{\texttt{031611}} \\
路加福音 10:25-37 & 鄺偉志牧師 & \hyperref[sec:031612]{邊gwo 係我ge 鄰舍?} & 2026-02-08 & \href{http://www.church.com.hk/acms/content.asp?site=cdc\&op=show\&type=product\&code=031612\&layout=sermon}{\texttt{031612}} \\
民數記 6:22-27 & 許永勝傳道 & \hyperref[sec:031613]{祝福人的人有福了!(粵/普)} & 2026-02-08 & \href{http://www.church.com.hk/acms/content.asp?site=cdc\&op=show\&type=product\&code=031613\&layout=sermon}{\texttt{031613}} \\
以弗所書 2:1-10 & 劉燕飛教師 & \hyperref[sec:031614]{活著 InG} & 2026-02-08 & \href{http://www.church.com.hk/acms/content.asp?site=cdc\&op=show\&type=product\&code=031614\&layout=sermon}{\texttt{031614}} \\
馬太福音 3:17 & 黃克勤牧師博士 & \hyperref[sec:031615]{聖經人聖經事:這是我的愛子} & 2026-01-29 & \href{http://www.church.com.hk/acms/content.asp?site=cdc\&op=show\&type=product\&code=031615\&layout=sermon}{\texttt{031615}} \\
箴言 2:1-22 & 周依娜牧師 & \hyperref[sec:031616]{留心聽智慧,專心求聰明 (普)} & 2025-07-20 & \href{http://www.church.com.hk/acms/content.asp?site=cdc\&op=show\&type=product\&code=031616\&layout=sermon}{\texttt{031616}} \\
出埃及記 33:12-23 & 張志怡傳道 & \hyperref[sec:031617]{渴求榮耀的禱告} & 2026-02-08 & \href{http://www.church.com.hk/acms/content.asp?site=cdc\&op=show\&type=product\&code=031617\&layout=sermon}{\texttt{031617}} \\
使徒行傳 9:1-6 & 曾興才牧師 & \hyperref[sec:031618]{大馬士革奇遇記 (普)} & 2026-02-08 & \href{http://www.church.com.hk/acms/content.asp?site=cdc\&op=show\&type=product\&code=031618\&layout=sermon}{\texttt{031618}} \\
以弗所書 6:10-12 & 招世超牧師 & \hyperref[sec:031619]{鞏固屬靈爭戰的大後方} & 2026-02-08 & \href{http://www.church.com.hk/acms/content.asp?site=cdc\&op=show\&type=product\&code=031619\&layout=sermon}{\texttt{031619}} \\
士師記 4:1-24 & 楊志溢傳道 & \hyperref[sec:031620]{信實的上帝} & 2026-02-08 & \href{http://www.church.com.hk/acms/content.asp?site=cdc\&op=show\&type=product\&code=031620\&layout=sermon}{\texttt{031620}} \\
詩篇 1:1-3 & 李慧珊傳道 & \hyperref[sec:031621]{珍貴的「日」與「夜」} & 2026-02-08 & \href{http://www.church.com.hk/acms/content.asp?site=cdc\&op=show\&type=product\&code=031621\&layout=sermon}{\texttt{031621}} \\
路加福音 1:13 & 黃克勤牧師博士 & \hyperref[sec:031622]{聖經人聖經事:祈禱已被聽見} & 2026-01-30 & \href{http://www.church.com.hk/acms/content.asp?site=cdc\&op=show\&type=product\&code=031622\&layout=sermon}{\texttt{031622}} \\
箴言 3:1-12 & 蕭壽華牧師 & \hyperref[sec:031623]{不輕看管教,不推辭施與} & 2025-07-27 & \href{http://www.church.com.hk/acms/content.asp?site=cdc\&op=show\&type=product\&code=031623\&layout=sermon}{\texttt{031623}} \\
約翰福音 5:31-47 & 林鍵牧師 & \hyperref[sec:031624]{榮耀君王:使你得救} & 2026-02-07 & \href{http://www.church.com.hk/acms/content.asp?site=cdc\&op=show\&type=product\&code=031624\&layout=sermon}{\texttt{031624}} \\
約伯記 1:1-2:13 & 俞大鳴牧師 & \hyperref[sec:031625]{全心全意} & 2026-02-08 & \href{http://www.church.com.hk/acms/content.asp?site=cdc\&op=show\&type=product\&code=031625\&layout=sermon}{\texttt{031625}} \\
加拉太書 4:28-31 & 葛國光牧師 & \hyperref[sec:031626]{作屬靈的兒女 (普)} & 2026-02-08 & \href{http://www.church.com.hk/acms/content.asp?site=cdc\&op=show\&type=product\&code=031626\&layout=sermon}{\texttt{031626}} \\
撒母耳記上 17:1-58 & 李思聰牧師 & \hyperref[sec:031627]{你也可以做巨人} & 2026-02-08 & \href{http://www.church.com.hk/acms/content.asp?site=cdc\&op=show\&type=product\&code=031627\&layout=sermon}{\texttt{031627}} \\
以賽亞書 14:12-17 & 李應新牧師 & \hyperref[sec:031628]{審判早晨之子} & 2026-02-08 & \href{http://www.church.com.hk/acms/content.asp?site=cdc\&op=show\&type=product\&code=031628\&layout=sermon}{\texttt{031628}} \\
路加福音 1:17 & 黃克勤牧師博士 & \hyperref[sec:031629]{聖經人聖經事:在主面前為大} & 2026-01-31 & \href{http://www.church.com.hk/acms/content.asp?site=cdc\&op=show\&type=product\&code=031629\&layout=sermon}{\texttt{031629}} \\
民數記 4:1-19 & 林誠信牧師 & \hyperref[sec:031630]{各司其職,有序同行} & 2026-02-08 & \href{http://www.church.com.hk/acms/content.asp?site=cdc\&op=show\&type=product\&code=031630\&layout=sermon}{\texttt{031630}} \\
箴言 21:31 & 梁永善牧師 & \hyperref[sec:031631]{馬年話馬} & 2026-02-11 & \href{http://www.church.com.hk/acms/content.asp?site=cdc\&op=show\&type=product\&code=031631\&layout=sermon}{\texttt{031631}} \\
羅馬書 12:2 & 鄭兆奇牧師 & \hyperref[sec:031632]{如何察驗神的旨意} & 2026-02-08 & \href{http://www.church.com.hk/acms/content.asp?site=cdc\&op=show\&type=product\&code=031632\&layout=sermon}{\texttt{031632}} \\
撒母耳記上 3:1-21 & 田俐力傳道 & \hyperref[sec:031633]{孩子是神給慕道園的禮物 (普/粵)} & 2026-02-08 & \href{http://www.church.com.hk/acms/content.asp?site=cdc\&op=show\&type=product\&code=031633\&layout=sermon}{\texttt{031633}} \\
羅馬書 13:8-14 & 袁惠鈞牧師博士 & \hyperref[sec:031634]{得救比初信時更近了} & 2026-02-08 & \href{http://www.church.com.hk/acms/content.asp?site=cdc\&op=show\&type=product\&code=031634\&layout=sermon}{\texttt{031634}} \\
路加福音 1:20 & 黃克勤牧師博士 & \hyperref[sec:031635]{聖經人聖經事:這話必然應驗} & 2026-02-01 & \href{http://www.church.com.hk/acms/content.asp?site=cdc\&op=show\&type=product\&code=031635\&layout=sermon}{\texttt{031635}} \\
約書亞記 4:1-24 & 魏文俊傳道 & \hyperref[sec:031636]{帶著過去往前走} & 2026-02-08 & \href{http://www.church.com.hk/acms/content.asp?site=cdc\&op=show\&type=product\&code=031636\&layout=sermon}{\texttt{031636}} \\
申命記 1:19-40 & 羅皓文傳道 & \hyperref[sec:031637]{過去與現在的對話 (1) YA 好聲音} & 2026-02-07 & \href{http://www.church.com.hk/acms/content.asp?site=cdc\&op=show\&type=product\&code=031637\&layout=sermon}{\texttt{031637}} \\
創世記 6:5-22 & 王啟恩傳道 & \hyperref[sec:031638]{多代同行•與神同行系列:信心的同行} & 2026-02-08 & \href{http://www.church.com.hk/acms/content.asp?site=cdc\&op=show\&type=product\&code=031638\&layout=sermon}{\texttt{031638}} \\
哥林多後書 1:15-24 & 陳雅莉傳道 & \hyperref[sec:031639]{第二次的恩典} & 2026-02-08 & \href{http://www.church.com.hk/acms/content.asp?site=cdc\&op=show\&type=product\&code=031639\&layout=sermon}{\texttt{031639}} \\
馬太福音 12:22-28 & 黃健羒牧師 & \hyperref[sec:031640]{神國臨在} & 2026-02-08 & \href{http://www.church.com.hk/acms/content.asp?site=cdc\&op=show\&type=product\&code=031640\&layout=sermon}{\texttt{031640}} \\
約書亞記 14:6-12 & 梁明財牧師 & \hyperref[sec:031641]{活出使命的跟從者——迦勒 I (普)} & 2026-02-08 & \href{http://www.church.com.hk/acms/content.asp?site=cdc\&op=show\&type=product\&code=031641\&layout=sermon}{\texttt{031641}} \\
路加福音 1:66 & 黃克勤牧師博士 & \hyperref[sec:031642]{聖經人聖經事:有主與他同在} & 2026-02-02 & \href{http://www.church.com.hk/acms/content.asp?site=cdc\&op=show\&type=product\&code=031642\&layout=sermon}{\texttt{031642}} \\
路加福音 16:1-13 & 陳傳華牧師博士 & \hyperref[sec:031643]{今世錢財,永恆眼光} & 2026-02-08 & \href{http://www.church.com.hk/acms/content.asp?site=cdc\&op=show\&type=product\&code=031643\&layout=sermon}{\texttt{031643}} \\
使徒行傳 9:32-43 & 禤凱榮牧師 & \hyperref[sec:031644]{福音超越界限} & 2026-02-15 & \href{http://www.church.com.hk/acms/content.asp?site=cdc\&op=show\&type=product\&code=031644\&layout=sermon}{\texttt{031644}} \\
羅馬書 1:16 & 趙崇健牧師 & \hyperref[sec:031645]{今天的福音} & 2026-02-15 & \href{http://www.church.com.hk/acms/content.asp?site=cdc\&op=show\&type=product\&code=031645\&layout=sermon}{\texttt{031645}} \\
希伯來書 10:19-39 & 余大峰傳道 & \hyperref[sec:031646]{你的信仰是否愈信愈後退?} & 2026-02-15 & \href{http://www.church.com.hk/acms/content.asp?site=cdc\&op=show\&type=product\&code=031646\&layout=sermon}{\texttt{031646}} \\
以賽亞書 29:13-14 & 羅婉明傳道 & \hyperref[sec:031647]{重得盼望} & 2026-02-15 & \href{http://www.church.com.hk/acms/content.asp?site=cdc\&op=show\&type=product\&code=031647\&layout=sermon}{\texttt{031647}} \\
出埃及記 24:12-18 & 蘇善美宣教師 & \hyperref[sec:031648]{再登山上} & 2026-02-15 & \href{http://www.church.com.hk/acms/content.asp?site=cdc\&op=show\&type=product\&code=031648\&layout=sermon}{\texttt{031648}} \\
創世記 11:3-4 & 黃克勤牧師博士 & \hyperref[sec:031649]{聖經人聖經事:製造品牌} & 2026-02-03 & \href{http://www.church.com.hk/acms/content.asp?site=cdc\&op=show\&type=product\&code=031649\&layout=sermon}{\texttt{031649}} \\
馬太福音 5:38-48 & 鄧瑞強牧師博士 & \hyperref[sec:031650]{不切實際的耶穌} & 2026-02-15 & \href{http://www.church.com.hk/acms/content.asp?site=cdc\&op=show\&type=product\&code=031650\&layout=sermon}{\texttt{031650}} \\
馬太福音 6:19-24 & 石偉漢牧師 & \hyperref[sec:031651]{見錢眼開?} & 2026-02-15 & \href{http://www.church.com.hk/acms/content.asp?site=cdc\&op=show\&type=product\&code=031651\&layout=sermon}{\texttt{031651}} \\
歌羅西書 1:9-14 & 黃紹權牧師 & \hyperref[sec:031652]{如何代禱?如何接受代禱?} & 2026-02-15 & \href{http://www.church.com.hk/acms/content.asp?site=cdc\&op=show\&type=product\&code=031652\&layout=sermon}{\texttt{031652}} \\
腓立比書 3:13-14 & 劉彤牧師 & \hyperref[sec:031653]{奔向標竿人生 (普)} & 2026-02-15 & \href{http://www.church.com.hk/acms/content.asp?site=cdc\&op=show\&type=product\&code=031653\&layout=sermon}{\texttt{031653}} \\
創世記 11:9 & 黃克勤牧師博士 & \hyperref[sec:031654]{聖經人聖經事:分散在全地} & 2026-02-05 & \href{http://www.church.com.hk/acms/content.asp?site=cdc\&op=show\&type=product\&code=031654\&layout=sermon}{\texttt{031654}} \\
約書亞記 5:1-15 & 劉銳鋒牧師 & \hyperref[sec:031655]{在神聖使命中更新與順服} & 2026-02-15 & \href{http://www.church.com.hk/acms/content.asp?site=cdc\&op=show\&type=product\&code=031655\&layout=sermon}{\texttt{031655}} \\
雅歌 8:6-7 & 賴若瀚牧師 & \hyperref[sec:031656]{真愛律} & 2026-02-15 & \href{http://www.church.com.hk/acms/content.asp?site=cdc\&op=show\&type=product\&code=031656\&layout=sermon}{\texttt{031656}} \\
啟示錄 19:11-16 & 陳崇基牧師 & \hyperref[sec:031657]{白馬之王} & 2026-02-15 & \href{http://www.church.com.hk/acms/content.asp?site=cdc\&op=show\&type=product\&code=031657\&layout=sermon}{\texttt{031657}} \\
羅馬書 15:14-33 & 袁惠鈞牧師博士 & \hyperref[sec:031658]{以行動活出因信稱義} & 2026-02-15 & \href{http://www.church.com.hk/acms/content.asp?site=cdc\&op=show\&type=product\&code=031658\&layout=sermon}{\texttt{031658}} \\
哈巴谷書 3:2 & 羅泰然牧師 & \hyperref[sec:031659]{復興:憐憫與回轉的盼望} & 2026-02-15 & \href{http://www.church.com.hk/acms/content.asp?site=cdc\&op=show\&type=product\&code=031659\&layout=sermon}{\texttt{031659}} \\
創世記 11:31 & 黃克勤牧師博士 & \hyperref[sec:031660]{聖經人聖經事:要往迦南} & 2026-02-06 & \href{http://www.church.com.hk/acms/content.asp?site=cdc\&op=show\&type=product\&code=031660\&layout=sermon}{\texttt{031660}} \\
哥林多前書 12:1-14:40 & 林羅麗斯牧師 & \hyperref[sec:031661]{成全的愛} & 2026-02-15 & \href{http://www.church.com.hk/acms/content.asp?site=cdc\&op=show\&type=product\&code=031661\&layout=sermon}{\texttt{031661}} \\
哥林多前書 13:1-8 & 冼嘉琪傳道 & \hyperref[sec:031662]{給屬靈暴發戶的妙道} & 2026-02-15 & \href{http://www.church.com.hk/acms/content.asp?site=cdc\&op=show\&type=product\&code=031662\&layout=sermon}{\texttt{031662}} \\
詩篇 8:1-9 & 張磊牧師 & \hyperref[sec:031663]{星空與塵土 (普)} & 2026-02-15 & \href{http://www.church.com.hk/acms/content.asp?site=cdc\&op=show\&type=product\&code=031663\&layout=sermon}{\texttt{031663}} \\
撒母耳記上 18:1-23:29 & 李思聰牧師 & \hyperref[sec:031664]{做個良朋益友} & 2026-02-15 & \href{http://www.church.com.hk/acms/content.asp?site=cdc\&op=show\&type=product\&code=031664\&layout=sermon}{\texttt{031664}} \\
列王記上 1:15-31 & 廖武洪牧師 & \hyperref[sec:031665]{應許消失了?} & 2026-02-15 & \href{http://www.church.com.hk/acms/content.asp?site=cdc\&op=show\&type=product\&code=031665\&layout=sermon}{\texttt{031665}} \\
創世記 12:4 & 黃克勤牧師博士 & \hyperref[sec:031666]{聖經人聖經事:離開本地本族父家} & 2026-02-07 & \href{http://www.church.com.hk/acms/content.asp?site=cdc\&op=show\&type=product\&code=031666\&layout=sermon}{\texttt{031666}} \\
約翰福音 6:1-15 & 林鍵牧師 & \hyperref[sec:031667]{榮耀君王:要他作王} & 2026-02-14 & \href{http://www.church.com.hk/acms/content.asp?site=cdc\&op=show\&type=product\&code=031667\&layout=sermon}{\texttt{031667}} \\
哥林多後書 4:7 & 湯傳哲牧師 & \hyperref[sec:031668]{大能的勇士 (普/粵)} & 2026-02-20 & \href{http://www.church.com.hk/acms/content.asp?site=cdc\&op=show\&type=product\&code=031668\&layout=sermon}{\texttt{031668}} \\
提摩太前書 3:1-7 & 張景恒牧師 & \hyperref[sec:031669]{長老的崇高呼召,資格和品格 (普/粵)} & 2026-02-15 & \href{http://www.church.com.hk/acms/content.asp?site=cdc\&op=show\&type=product\&code=031669\&layout=sermon}{\texttt{031669}} \\
詩篇 62:1-12 & 陳永就牧師 & \hyperref[sec:031670]{實在靠得住} & 2026-02-15 & \href{http://www.church.com.hk/acms/content.asp?site=cdc\&op=show\&type=product\&code=031670\&layout=sermon}{\texttt{031670}} \\
創世記 14:20 & 黃克勤牧師博士 & \hyperref[sec:031671]{聖經人聖經事:麥基洗德} & 2026-02-08 & \href{http://www.church.com.hk/acms/content.asp?site=cdc\&op=show\&type=product\&code=031671\&layout=sermon}{\texttt{031671}} \\
路加福音 16:19-31 & 梁成裕牧師 & \hyperref[sec:031672]{今生決定永生} & 2026-02-15 & \href{http://www.church.com.hk/acms/content.asp?site=cdc\&op=show\&type=product\&code=031672\&layout=sermon}{\texttt{031672}} \\
約伯記 39:19-25 & 曹子光牧師 & \hyperref[sec:031673]{無限馬力} & 2026-02-15 & \href{http://www.church.com.hk/acms/content.asp?site=cdc\&op=show\&type=product\&code=031673\&layout=sermon}{\texttt{031673}} \\
申命記 2:1-3:29 & 何淑斌傳道 & \hyperref[sec:031674]{過去和現在的對話 (2) 聽歷史,講歷史,寫歷史} & 2026-02-14 & \href{http://www.church.com.hk/acms/content.asp?site=cdc\&op=show\&type=product\&code=031674\&layout=sermon}{\texttt{031674}} \\
箴言 4:23 & 潘梁其德傳道 & \hyperref[sec:031675]{祝福 •由心而發} & 2026-02-15 & \href{http://www.church.com.hk/acms/content.asp?site=cdc\&op=show\&type=product\&code=031675\&layout=sermon}{\texttt{031675}} \\
創世記 22:1-19 & 黃浩君傳道 & \hyperref[sec:031676]{通過「考驗」成為敬畏神的人} & 2026-02-14 & \href{http://www.church.com.hk/acms/content.asp?site=cdc\&op=show\&type=product\&code=031676\&layout=sermon}{\texttt{031676}} \\
羅馬書 15:22-33 & 秦民德傳道 & \hyperref[sec:031677]{「無我」的事奉} & 2026-02-15 & \href{http://www.church.com.hk/acms/content.asp?site=cdc\&op=show\&type=product\&code=031677\&layout=sermon}{\texttt{031677}} \\
創世記 17:1 & 黃克勤牧師博士 & \hyperref[sec:031678]{聖經人聖經事:作完全的人} & 2026-02-09 & \href{http://www.church.com.hk/acms/content.asp?site=cdc\&op=show\&type=product\&code=031678\&layout=sermon}{\texttt{031678}} \\
使徒行傳 11:19-30 & 駱瑞麟傳道 & \hyperref[sec:031679]{落在異鄉,不為異客} & 2026-02-22 & \href{http://www.church.com.hk/acms/content.asp?site=cdc\&op=show\&type=product\&code=031679\&layout=sermon}{\texttt{031679}} \\
創世記 32:22-32 & 林穎豪傳道 & \hyperref[sec:031680]{雅各之爭三部曲:與神之爭} & 2026-02-22 & \href{http://www.church.com.hk/acms/content.asp?site=cdc\&op=show\&type=product\&code=031680\&layout=sermon}{\texttt{031680}} \\
約書亞記 1:1-18 & 吳銘淇傳道 & \hyperref[sec:031681]{求主給我看見異象} & 2026-02-22 & \href{http://www.church.com.hk/acms/content.asp?site=cdc\&op=show\&type=product\&code=031681\&layout=sermon}{\texttt{031681}} \\
羅馬書 5:12-17 & 林玉英牧師 & \hyperref[sec:031682]{不對稱的救贖} & 2026-02-21 & \href{http://www.church.com.hk/acms/content.asp?site=cdc\&op=show\&type=product\&code=031682\&layout=sermon}{\texttt{031682}} \\
馬太福音 4:1-11 & 高國雄牧師 & \hyperref[sec:031683]{如何對付那試探者——魔鬼撒旦} & 2026-02-22 & \href{http://www.church.com.hk/acms/content.asp?site=cdc\&op=show\&type=product\&code=031683\&layout=sermon}{\texttt{031683}} \\
創世記 17:16 & 黃克勤牧師博士 & \hyperref[sec:031684]{聖經人聖經事:從她得一個兒子} & 2026-02-11 & \href{http://www.church.com.hk/acms/content.asp?site=cdc\&op=show\&type=product\&code=031684\&layout=sermon}{\texttt{031684}} \\
歌羅西書 1:15-20 & 黃紹權牧師 & \hyperref[sec:031685]{耶穌是永恆超級屬靈接合劑} & 2026-02-22 & \href{http://www.church.com.hk/acms/content.asp?site=cdc\&op=show\&type=product\&code=031685\&layout=sermon}{\texttt{031685}} \\
歷代志下 20:1-3 & 卓爾君牧師 & \hyperref[sec:031686]{絕境中的神操作 (華)} & 2026-02-22 & \href{http://www.church.com.hk/acms/content.asp?site=cdc\&op=show\&type=product\&code=031686\&layout=sermon}{\texttt{031686}} \\
馬可福音 16:15 & 張瑪莉女士 & \hyperref[sec:031687]{張瑪莉突破生命困境分享} & 2026-02-22 & \href{http://www.church.com.hk/acms/content.asp?site=cdc\&op=show\&type=product\&code=031687\&layout=sermon}{\texttt{031687}} \\
約書亞記 6:1-27 & 張智聰博士 & \hyperref[sec:031688]{同心打美好的仗} & 2026-02-22 & \href{http://www.church.com.hk/acms/content.asp?site=cdc\&op=show\&type=product\&code=031688\&layout=sermon}{\texttt{031688}} \\
創世記 19:36 & 黃克勤牧師博士 & \hyperref[sec:031689]{聖經人聖經事:人的方法} & 2026-02-12 & \href{http://www.church.com.hk/acms/content.asp?site=cdc\&op=show\&type=product\&code=031689\&layout=sermon}{\texttt{031689}} \\
詩篇 104:24-35 & 江海薇傳道 & \hyperref[sec:031690]{總是記得創造主} & 2026-02-22 & \href{http://www.church.com.hk/acms/content.asp?site=cdc\&op=show\&type=product\&code=031690\&layout=sermon}{\texttt{031690}} \\
哥林多後書 11:1-13:14 & 劉少康牧師 & \hyperref[sec:031691]{哥林多前後書 (6) 保羅以患難為榮} & 2026-02-11 & \href{http://www.church.com.hk/acms/content.asp?site=cdc\&op=show\&type=product\&code=031691\&layout=sermon}{\texttt{031691}} \\
以弗所書 4:14-16 & 鍾韵悠牧師 & \hyperref[sec:031692]{建立健康的使命教會 (5) 以真理和恩典建立關係} & 2026-02-22 & \href{http://www.church.com.hk/acms/content.asp?site=cdc\&op=show\&type=product\&code=031692\&layout=sermon}{\texttt{031692}} \\
以斯帖記 8:3-9:16 & 黃國棟醫生 & \hyperref[sec:031693]{神的救贖 人的責任} & 2026-02-22 & \href{http://www.church.com.hk/acms/content.asp?site=cdc\&op=show\&type=product\&code=031693\&layout=sermon}{\texttt{031693}} \\
馬太福音 5:27-30 & 張健庭牧師博士 & \hyperref[sec:031694]{重塑生命系列 (4) 色慾與貞潔} & 2026-02-22 & \href{http://www.church.com.hk/acms/content.asp?site=cdc\&op=show\&type=product\&code=031694\&layout=sermon}{\texttt{031694}} \\
創世記 21:13 & 黃克勤牧師博士 & \hyperref[sec:031695]{聖經人聖經事:因他是你所生的} & 2026-02-13 & \href{http://www.church.com.hk/acms/content.asp?site=cdc\&op=show\&type=product\&code=031695\&layout=sermon}{\texttt{031695}} \\
傳道書 1:2-7 & 周朗先生 & \hyperref[sec:031696]{它,他/她,祂} & 2026-02-22 & \href{http://www.church.com.hk/acms/content.asp?site=cdc\&op=show\&type=product\&code=031696\&layout=sermon}{\texttt{031696}} \\
馬太福音 4:1-17 & 林豪恩先生 & \hyperref[sec:031697]{三項試探,一個中心} & 2026-02-22 & \href{http://www.church.com.hk/acms/content.asp?site=cdc\&op=show\&type=product\&code=031697\&layout=sermon}{\texttt{031697}} \\
馬可福音 9:14-29 & 梁永善牧師 & \hyperref[sec:031698]{變臉後傳 I 心弱力微} & 2026-02-25 & \href{http://www.church.com.hk/acms/content.asp?site=cdc\&op=show\&type=product\&code=031698\&layout=sermon}{\texttt{031698}} \\
加拉太書 5:1 & 葛國光牧師 & \hyperref[sec:031699]{真自由 (華)} & 2026-02-22 & \href{http://www.church.com.hk/acms/content.asp?site=cdc\&op=show\&type=product\&code=031699\&layout=sermon}{\texttt{031699}} \\
撒迦利亞書 8:1-23 & 黃馮潔霞師母 & \hyperref[sec:031700]{有神同在,恩福滿滿} & 2026-02-22 & \href{http://www.church.com.hk/acms/content.asp?site=cdc\&op=show\&type=product\&code=031700\&layout=sermon}{\texttt{031700}} \\
創世記 22:12 & 黃克勤牧師博士 & \hyperref[sec:031701]{聖經人聖經事:我知道} & 2026-02-14 & \href{http://www.church.com.hk/acms/content.asp?site=cdc\&op=show\&type=product\&code=031701\&layout=sermon}{\texttt{031701}} \\
約翰福音 6:13-25 & 林鍵牧師 & \hyperref[sec:031702]{榮耀君王:不要害怕} & 2026-02-21 & \href{http://www.church.com.hk/acms/content.asp?site=cdc\&op=show\&type=product\&code=031702\&layout=sermon}{\texttt{031702}} \\
約伯記 39:19-25 & 吳宗文牧師 & \hyperref[sec:031703]{一馬當先,誰與爭鋒!} & 2026-02-22 & \href{http://www.church.com.hk/acms/content.asp?site=cdc\&op=show\&type=product\&code=031703\&layout=sermon}{\texttt{031703}} \\
馬太福音 25:1-13 & 劉少康牧師 & \hyperref[sec:031704]{我的燈需要油} & 2026-02-22 & \href{http://www.church.com.hk/acms/content.asp?site=cdc\&op=show\&type=product\&code=031704\&layout=sermon}{\texttt{031704}} \\
希伯來書 3:7-4:11 & 歐世達牧師 & \hyperref[sec:031705]{創造之始的約:工作關係 (英/粵)} & 2026-02-22 & \href{http://www.church.com.hk/acms/content.asp?site=cdc\&op=show\&type=product\&code=031705\&layout=sermon}{\texttt{031705}} \\
約翰福音 3:5-15 & 駱柱成牧師 & \hyperref[sec:031706]{耶穌與尼哥底母談重生} & 2026-02-22 & \href{http://www.church.com.hk/acms/content.asp?site=cdc\&op=show\&type=product\&code=031706\&layout=sermon}{\texttt{031706}} \\
創世記 23:13 & 黃克勤牧師博士 & \hyperref[sec:031707]{聖經人聖經事:人情世故} & 2026-02-15 & \href{http://www.church.com.hk/acms/content.asp?site=cdc\&op=show\&type=product\&code=031707\&layout=sermon}{\texttt{031707}} \\
撒母耳記上 17:21-26 & 許永勝傳道 & \hyperref[sec:031708]{成就不可能的任務 (粵/華)} & 2026-02-22 & \href{http://www.church.com.hk/acms/content.asp?site=cdc\&op=show\&type=product\&code=031708\&layout=sermon}{\texttt{031708}} \\
列王記下 2:1-18 & 溫艾衛傳道 & \hyperref[sec:031709]{雙倍恩膏,薪火傳承} & 2026-02-22 & \href{http://www.church.com.hk/acms/content.asp?site=cdc\&op=show\&type=product\&code=031709\&layout=sermon}{\texttt{031709}} \\
哈該書 2:1-9 & 蕭麗琴傳道 & \hyperref[sec:031710]{榮耀的殿} & 2026-02-22 & \href{http://www.church.com.hk/acms/content.asp?site=cdc\&op=show\&type=product\&code=031710\&layout=sermon}{\texttt{031710}} \\
馬太福音 5:1-6 & 陳雅莉傳道 & \hyperref[sec:031711]{有福了!屬神的人 (1)} & 2026-02-22 & \href{http://www.church.com.hk/acms/content.asp?site=cdc\&op=show\&type=product\&code=031711\&layout=sermon}{\texttt{031711}} \\
加拉太書 1:1-3:29 & 劉少康牧師 & \hyperref[sec:031712]{加拉太書 (1) 因信稱義的真理} & 2026-02-25 & \href{http://www.church.com.hk/acms/content.asp?site=cdc\&op=show\&type=product\&code=031712\&layout=sermon}{\texttt{031712}} \\
創世記 24:7 & 黃克勤牧師博士 & \hyperref[sec:031713]{聖經人聖經事:娶一個妻子} & 2026-02-17 & \href{http://www.church.com.hk/acms/content.asp?site=cdc\&op=show\&type=product\&code=031713\&layout=sermon}{\texttt{031713}} \\
申命記 4:1-14 & 李觀南傳道 & \hyperref[sec:031714]{過去和現在的對話 (3) 聽誰說} & 2026-02-21 & \href{http://www.church.com.hk/acms/content.asp?site=cdc\&op=show\&type=product\&code=031714\&layout=sermon}{\texttt{031714}} \\
創世記 12:1-4 & 陳黎韻怡傳道 & \hyperref[sec:031715]{多代同行•與神同行系列:當人生遇上呼召時} & 2026-02-22 & \href{http://www.church.com.hk/acms/content.asp?site=cdc\&op=show\&type=product\&code=031715\&layout=sermon}{\texttt{031715}} \\
詩篇 121:1-8 & 藍鷺平姊妹 & \hyperref[sec:031716]{繼續前行 (華/粵)} & 2026-02-22 & \href{http://www.church.com.hk/acms/content.asp?site=cdc\&op=show\&type=product\&code=031716\&layout=sermon}{\texttt{031716}} \\
馬太福音 5:7-12 & 陳雅莉傳道 & \hyperref[sec:031717]{有福了!屬神的人 (2)} & 2026-03-01 & \href{http://www.church.com.hk/acms/content.asp?site=cdc\&op=show\&type=product\&code=031717\&layout=sermon}{\texttt{031717}} \\
歌羅西書 1:21-23 & 黃紹權牧師 & \hyperref[sec:031718]{主耶穌已經開始了,我們又怎樣呢?} & 2026-03-01 & \href{http://www.church.com.hk/acms/content.asp?site=cdc\&op=show\&type=product\&code=031718\&layout=sermon}{\texttt{031718}} \\
以弗所書 4:1-3 & 鍾梓昭傳道 & \hyperref[sec:031719]{建立健康的使命教會 (6) 竭力同心同行} & 2026-03-01 & \href{http://www.church.com.hk/acms/content.asp?site=cdc\&op=show\&type=product\&code=031719\&layout=sermon}{\texttt{031719}} \\
創世記 12:1-4 & 林家棋宣教師 & \hyperref[sec:031720]{Rebbot 404 ID} & 2026-03-01 & \href{http://www.church.com.hk/acms/content.asp?site=cdc\&op=show\&type=product\&code=031720\&layout=sermon}{\texttt{031720}} \\
箴言 4:10-27 & 周敬敏牧師 & \hyperref[sec:031721]{保守你心勝過保守一切} & 2025-08-03 & \href{http://www.church.com.hk/acms/content.asp?site=cdc\&op=show\&type=product\&code=031721\&layout=sermon}{\texttt{031721}} \\
約翰福音 3:1-17 & 高國雄牧師 & \hyperref[sec:031722]{上帝愛世界} & 2026-03-01 & \href{http://www.church.com.hk/acms/content.asp?site=cdc\&op=show\&type=product\&code=031722\&layout=sermon}{\texttt{031722}} \\
創世記 24:27 & 黃克勤牧師博士 & \hyperref[sec:031723]{聖經人聖經事:耶和華引領我} & 2026-03-18 & \href{http://www.church.com.hk/acms/content.asp?site=cdc\&op=show\&type=product\&code=031723\&layout=sermon}{\texttt{031723}} \\
使徒行傳 12:1-24 & 楊芹牧師 & \hyperref[sec:031724]{神掌權 (華)} & 2026-03-01 & \href{http://www.church.com.hk/acms/content.asp?site=cdc\&op=show\&type=product\&code=031724\&layout=sermon}{\texttt{031724}} \\
啟示錄 1:1-2 & 趙崇健牧師 & \hyperref[sec:031725]{序:看見啟示錄} & 2026-03-01 & \href{http://www.church.com.hk/acms/content.asp?site=cdc\&op=show\&type=product\&code=031725\&layout=sermon}{\texttt{031725}} \\
約翰福音 4:46-54 & 謝友德牧師 & \hyperref[sec:031726]{信心成長進程} & 2026-03-01 & \href{http://www.church.com.hk/acms/content.asp?site=cdc\&op=show\&type=product\&code=031726\&layout=sermon}{\texttt{031726}} \\
創世記 12:1-4 & 王家輝牧師 & \hyperref[sec:031727]{冒險的祝福} & 2026-03-01 & \href{http://www.church.com.hk/acms/content.asp?site=cdc\&op=show\&type=product\&code=031727\&layout=sermon}{\texttt{031727}} \\
創世記 24:56 & 黃克勤牧師博士 & \hyperref[sec:031728]{聖經人聖經事:不要耽誤我} & 2026-02-19 & \href{http://www.church.com.hk/acms/content.asp?site=cdc\&op=show\&type=product\&code=031728\&layout=sermon}{\texttt{031728}} \\
利未記 1:1-2 & 羅泰然牧師 & \hyperref[sec:031729]{從祭壇到心靈奉獻} & 2026-03-01 & \href{http://www.church.com.hk/acms/content.asp?site=cdc\&op=show\&type=product\&code=031729\&layout=sermon}{\texttt{031729}} \\
馬太福音 6:25-34 & 石偉漢牧師 & \hyperref[sec:031730]{憂慮?何足掛齒?} & 2026-03-01 & \href{http://www.church.com.hk/acms/content.asp?site=cdc\&op=show\&type=product\&code=031730\&layout=sermon}{\texttt{031730}} \\
雅各書 5:13-16 & 劉彤牧師 & \hyperref[sec:031731]{如何在生活中支取神的能力 (華)} & 2026-03-01 & \href{http://www.church.com.hk/acms/content.asp?site=cdc\&op=show\&type=product\&code=031731\&layout=sermon}{\texttt{031731}} \\
詩篇 42:1-43:5 & 李業基牧師 & \hyperref[sec:031732]{茅塞頓開} & 2026-03-01 & \href{http://www.church.com.hk/acms/content.asp?site=cdc\&op=show\&type=product\&code=031732\&layout=sermon}{\texttt{031732}} \\
詩篇 22:24 & 司徒永富博士 & \hyperref[sec:031733]{被看見的力量} & 2026-03-01 & \href{http://www.church.com.hk/acms/content.asp?site=cdc\&op=show\&type=product\&code=031733\&layout=sermon}{\texttt{031733}} \\
創世記 24:67 & 黃克勤牧師博士 & \hyperref[sec:031734]{聖經人聖經事:得了安慰} & 2026-02-20 & \href{http://www.church.com.hk/acms/content.asp?site=cdc\&op=show\&type=product\&code=031734\&layout=sermon}{\texttt{031734}} \\
箴言 5:15-23 & 袁兆添傳道 & \hyperref[sec:031735]{快掙脫纏繞 當速速醒來} & 2025-08-10 & \href{http://www.church.com.hk/acms/content.asp?site=cdc\&op=show\&type=product\&code=031735\&layout=sermon}{\texttt{031735}} \\
路加福音 17:11-19 & 梁成裕牧師 & \hyperref[sec:031736]{從醫治到拯救} & 2026-03-01 & \href{http://www.church.com.hk/acms/content.asp?site=cdc\&op=show\&type=product\&code=031736\&layout=sermon}{\texttt{031736}} \\
箴言 14:26-35 & 梁永善牧師 & \hyperref[sec:031737]{生命泉源} & 2026-03-01 & \href{http://www.church.com.hk/acms/content.asp?site=cdc\&op=show\&type=product\&code=031737\&layout=sermon}{\texttt{031737}} \\
詩篇 150:1-6 & 陳崇基牧師 & \hyperref[sec:031738]{全人的敬拜} & 2026-03-01 & \href{http://www.church.com.hk/acms/content.asp?site=cdc\&op=show\&type=product\&code=031738\&layout=sermon}{\texttt{031738}} \\
以弗所書 1:9-10 & 葛國光牧師 & \hyperref[sec:031739]{以「合一」迎主來 (華)} & 2026-03-01 & \href{http://www.church.com.hk/acms/content.asp?site=cdc\&op=show\&type=product\&code=031739\&layout=sermon}{\texttt{031739}} \\
撒母耳記上 18:1-21:15 & 李思聰牧師 & \hyperref[sec:031740]{甚麼是我最終的依靠} & 2026-03-01 & \href{http://www.church.com.hk/acms/content.asp?site=cdc\&op=show\&type=product\&code=031740\&layout=sermon}{\texttt{031740}} \\
創世記 25:16 & 黃克勤牧師博士 & \hyperref[sec:031741]{聖經人聖經事:十二族長} & 2026-02-21 & \href{http://www.church.com.hk/acms/content.asp?site=cdc\&op=show\&type=product\&code=031741\&layout=sermon}{\texttt{031741}} \\
民數記 5:1-31 & 李志輝傳道 & \hyperref[sec:031742]{神聖群體的健康學} & 2026-02-28 & \href{http://www.church.com.hk/acms/content.asp?site=cdc\&op=show\&type=product\&code=031742\&layout=sermon}{\texttt{031742}} \\
民數記 5:1-31 & 何慧芬牧師 & \hyperref[sec:031743]{神住在我們中間:聖潔的要求和實踐} & 2026-03-01 & \href{http://www.church.com.hk/acms/content.asp?site=cdc\&op=show\&type=product\&code=031743\&layout=sermon}{\texttt{031743}} \\
馬太福音 17:24-27 & 梁永善牧師 & \hyperref[sec:031744]{變臉後傳 II  掘銀交稅} & 2026-03-04 & \href{http://www.church.com.hk/acms/content.asp?site=cdc\&op=show\&type=product\&code=031744\&layout=sermon}{\texttt{031744}} \\
約翰福音 15:16 & Lydia Yeung姊妹 & \hyperref[sec:031745]{呼召,回應,行動——帶來的祝福} & 2026-03-01 & \href{http://www.church.com.hk/acms/content.asp?site=cdc\&op=show\&type=product\&code=031745\&layout=sermon}{\texttt{031745}} \\
羅馬書 16:1-27 & 袁惠鈞牧師博士 & \hyperref[sec:031746]{堅固人心的福音奧秘} & 2026-03-01 & \href{http://www.church.com.hk/acms/content.asp?site=cdc\&op=show\&type=product\&code=031746\&layout=sermon}{\texttt{031746}} \\
以賽亞書 42:1-9 & 周弘毅牧師 & \hyperref[sec:031747]{看哪!我的僕人} & 2026-02-14 & \href{http://www.church.com.hk/acms/content.asp?site=cdc\&op=show\&type=product\&code=031747\&layout=sermon}{\texttt{031747}} \\
創世記 25:21 & 黃克勤牧師博士 & \hyperref[sec:031748]{聖經人聖經事:祈求耶和華} & 2026-03-01 & \href{http://www.church.com.hk/acms/content.asp?site=cdc\&op=show\&type=product\&code=031748\&layout=sermon}{\texttt{031748}} \\
申命記 4:15-30 & 葉向榮傳道 & \hyperref[sec:031749]{過去和現在的對話 (4) 入鄉除俗} & 2026-02-28 & \href{http://www.church.com.hk/acms/content.asp?site=cdc\&op=show\&type=product\&code=031749\&layout=sermon}{\texttt{031749}} \\
列王記下 2:1-14 & 唐郭麗婷牧師 & \hyperref[sec:031750]{多代同行•與神同行系列:主,我要你更多:以利沙的加倍之靈} & 2026-03-01 & \href{http://www.church.com.hk/acms/content.asp?site=cdc\&op=show\&type=product\&code=031750\&layout=sermon}{\texttt{031750}} \\
雅各書 1:19-21 & 張健庭牧師博士 & \hyperref[sec:031751]{重塑生命系列 (5) 忿怒與溫柔} & 2026-03-01 & \href{http://www.church.com.hk/acms/content.asp?site=cdc\&op=show\&type=product\&code=031751\&layout=sermon}{\texttt{031751}} \\
馬太福音 5:13-16 & 伍宇衛長老 & \hyperref[sec:031752]{鹽和光} & 2026-03-01 & \href{http://www.church.com.hk/acms/content.asp?site=cdc\&op=show\&type=product\&code=031752\&layout=sermon}{\texttt{031752}} \\
以弗所書 6:1-4 & 葉慧敏傳道 & \hyperref[sec:031753]{在基督裡的親子關係} & 2026-03-01 & \href{http://www.church.com.hk/acms/content.asp?site=cdc\&op=show\&type=product\&code=031753\&layout=sermon}{\texttt{031753}} \\
以賽亞書 44:1-23 & 周弘毅牧師 & \hyperref[sec:031754]{告別虛空} & 2026-02-28 & \href{http://www.church.com.hk/acms/content.asp?site=cdc\&op=show\&type=product\&code=031754\&layout=sermon}{\texttt{031754}} \\
出埃及記 35:30-36:1 & 杜崇基牧師 & \hyperref[sec:031755]{比撒列的自白} & 2026-03-01 & \href{http://www.church.com.hk/acms/content.asp?site=cdc\&op=show\&type=product\&code=031755\&layout=sermon}{\texttt{031755}} \\
創世記 25:24 & 黃克勤牧師博士 & \hyperref[sec:031756]{聖經人聖經事:果然是雙子} & 2026-02-23 & \href{http://www.church.com.hk/acms/content.asp?site=cdc\&op=show\&type=product\&code=031756\&layout=sermon}{\texttt{031756}} \\
羅馬書 16:1-16 & 嚴增森弟兄 & \hyperref[sec:031757]{我們這一家•坐言起行} & 2026-03-01 & \href{http://www.church.com.hk/acms/content.asp?site=cdc\&op=show\&type=product\&code=031757\&layout=sermon}{\texttt{031757}} \\
路加福音 15:4-7 & 姚永光傳道 & \hyperref[sec:031758]{一比九十九} & 2025-08-17 & \href{http://www.church.com.hk/acms/content.asp?site=cdc\&op=show\&type=product\&code=031758\&layout=sermon}{\texttt{031758}} \\
使徒行傳 13:1-12 & 禤凱榮牧師 & \hyperref[sec:031759]{黑暗敵擋光明} & 2026-03-08 & \href{http://www.church.com.hk/acms/content.asp?site=cdc\&op=show\&type=product\&code=031759\&layout=sermon}{\texttt{031759}} \\
羅馬書 5:18-21 & 梁錫強醫生 & \hyperref[sec:031760]{失敗中的英雄} & 2026-03-08 & \href{http://www.church.com.hk/acms/content.asp?site=cdc\&op=show\&type=product\&code=031760\&layout=sermon}{\texttt{031760}} \\
加拉太書 4:1-6:18 & 劉少康牧師 & \hyperref[sec:031761]{加拉太書 (2) 基督使我得自由} & 2026-03-04 & \href{http://www.church.com.hk/acms/content.asp?site=cdc\&op=show\&type=product\&code=031761\&layout=sermon}{\texttt{031761}} \\
歌羅西書 1:24-29 & 黃紹權牧師 & \hyperref[sec:031762]{事奉主的喜痛交雜} & 2026-03-08 & \href{http://www.church.com.hk/acms/content.asp?site=cdc\&op=show\&type=product\&code=031762\&layout=sermon}{\texttt{031762}} \\
創世記 26:35 & 黃克勤牧師博士 & \hyperref[sec:031763]{聖經人聖經事:心裡愁煩} & 2026-02-25 & \href{http://www.church.com.hk/acms/content.asp?site=cdc\&op=show\&type=product\&code=031763\&layout=sermon}{\texttt{031763}} \\
馬太福音 9:14-17 & 王正傑牧師 & \hyperref[sec:031764]{變!定唔變?} & 2026-03-01 & \href{http://www.church.com.hk/acms/content.asp?site=cdc\&op=show\&type=product\&code=031764\&layout=sermon}{\texttt{031764}} \\
約翰福音 4:5-5:42 & 李均熊牧師博士 & \hyperref[sec:031765]{活水之道} & 2026-03-08 & \href{http://www.church.com.hk/acms/content.asp?site=cdc\&op=show\&type=product\&code=031765\&layout=sermon}{\texttt{031765}} \\
以弗所書 4:25-32 & 陳燕華牧師 & \hyperref[sec:031766]{建立健康的使命教會 (7) 只說造就人的話} & 2026-03-08 & \href{http://www.church.com.hk/acms/content.asp?site=cdc\&op=show\&type=product\&code=031766\&layout=sermon}{\texttt{031766}} \\
詩篇 100:1-5 & 林廣潤牧師 & \hyperref[sec:031767]{朝聖路上——歡欣敬拜} & 2026-03-08 & \href{http://www.church.com.hk/acms/content.asp?site=cdc\&op=show\&type=product\&code=031767\&layout=sermon}{\texttt{031767}} \\
約書亞記 2:1-24 & 吳銘淇傳道 & \hyperref[sec:031768]{抓著應許} & 2026-03-08 & \href{http://www.church.com.hk/acms/content.asp?site=cdc\&op=show\&type=product\&code=031768\&layout=sermon}{\texttt{031768}} \\
約翰福音 17:18 & 羅錫堅牧師 & \hyperref[sec:031769]{上帝為何不差派天使作宣教士?} & 2026-03-08 & \href{http://www.church.com.hk/acms/content.asp?site=cdc\&op=show\&type=product\&code=031769\&layout=sermon}{\texttt{031769}} \\
創世記 27:2 & 黃克勤牧師博士 & \hyperref[sec:031770]{聖經人聖經事:不知哪天死} & 2026-02-26 & \href{http://www.church.com.hk/acms/content.asp?site=cdc\&op=show\&type=product\&code=031770\&layout=sermon}{\texttt{031770}} \\
創世記 1:1-3:24 & 林羅麗斯牧師 & \hyperref[sec:031771]{女人的使命} & 2026-03-08 & \href{http://www.church.com.hk/acms/content.asp?site=cdc\&op=show\&type=product\&code=031771\&layout=sermon}{\texttt{031771}} \\
歌羅西書 1:15-23 & 俞大鳴牧師 & \hyperref[sec:031772]{在基督裡} & 2026-03-08 & \href{http://www.church.com.hk/acms/content.asp?site=cdc\&op=show\&type=product\&code=031772\&layout=sermon}{\texttt{031772}} \\
以弗所書 6:13 & 王亞辰牧師 & \hyperref[sec:031773]{成就了一切,還能站立得住 (華)} & 2026-03-08 & \href{http://www.church.com.hk/acms/content.asp?site=cdc\&op=show\&type=product\&code=031773\&layout=sermon}{\texttt{031773}} \\
提摩太後書 2:3-4 & 招世超牧師 & \hyperref[sec:031774]{突破知易難行的瓶頸} & 2026-03-08 & \href{http://www.church.com.hk/acms/content.asp?site=cdc\&op=show\&type=product\&code=031774\&layout=sermon}{\texttt{031774}} \\
出埃及記 17:1-7 & 林玉英牧師 & \hyperref[sec:031775]{誰為妳受擊打?} & 2026-03-08 & \href{http://www.church.com.hk/acms/content.asp?site=cdc\&op=show\&type=product\&code=031775\&layout=sermon}{\texttt{031775}} \\
創世記 27:13 & 黃克勤牧師博士 & \hyperref[sec:031776]{聖經人聖經事:我兒,咒詛歸我} & 2026-02-27 & \href{http://www.church.com.hk/acms/content.asp?site=cdc\&op=show\&type=product\&code=031776\&layout=sermon}{\texttt{031776}} \\
尼希米記 1:1-4 & 顏沛恩牧師 & \hyperref[sec:031777]{生命無牆,好快收檔 (1)} & 2026-03-08 & \href{http://www.church.com.hk/acms/content.asp?site=cdc\&op=show\&type=product\&code=031777\&layout=sermon}{\texttt{031777}} \\
路加福音 17:20-37 & 陳傳華牧師博士 & \hyperref[sec:031778]{隨時預備,警醒迎接神的國} & 2026-03-08 & \href{http://www.church.com.hk/acms/content.asp?site=cdc\&op=show\&type=product\&code=031778\&layout=sermon}{\texttt{031778}} \\
出埃及記 3:1-22 & 江芷敏傳道 & \hyperref[sec:031779]{呼召的關係} & 2026-03-08 & \href{http://www.church.com.hk/acms/content.asp?site=cdc\&op=show\&type=product\&code=031779\&layout=sermon}{\texttt{031779}} \\
約翰福音 1:1-6 & 潘智翔牧師 & \hyperref[sec:031780]{在恩典裡面對新的挑戰} & 2026-03-08 & \href{http://www.church.com.hk/acms/content.asp?site=cdc\&op=show\&type=product\&code=031780\&layout=sermon}{\texttt{031780}} \\
腓立比書 3:18-19 & 張健庭牧師博士 & \hyperref[sec:031781]{重塑生命系列 (6) 貪食與節制——當慾望失去界線} & 2026-03-08 & \href{http://www.church.com.hk/acms/content.asp?site=cdc\&op=show\&type=product\&code=031781\&layout=sermon}{\texttt{031781}} \\
創世記 27:19 & 黃克勤牧師博士 & \hyperref[sec:031782]{聖經人聖經事:我是你的長子以掃} & 2026-02-28 & \href{http://www.church.com.hk/acms/content.asp?site=cdc\&op=show\&type=product\&code=031782\&layout=sermon}{\texttt{031782}} \\
約翰福音 6:22-40 & 林鍵牧師 & \hyperref[sec:031783]{榮耀君王:永生天糧} & 2026-03-07 & \href{http://www.church.com.hk/acms/content.asp?site=cdc\&op=show\&type=product\&code=031783\&layout=sermon}{\texttt{031783}} \\
箴言 6:20-35 & 莊明偉牧師 & \hyperref[sec:031784]{管教的責備是生命的道} & 2025-08-24 & \href{http://www.church.com.hk/acms/content.asp?site=cdc\&op=show\&type=product\&code=031784\&layout=sermon}{\texttt{031784}} \\
撒母耳記上 18:1-20:42 & 李觀南傳道 & \hyperref[sec:031785]{多代同行 (2) PeerWalk (1) Godly Friends, Godly Wisdom} & 2026-03-07 & \href{http://www.church.com.hk/acms/content.asp?site=cdc\&op=show\&type=product\&code=031785\&layout=sermon}{\texttt{031785}} \\
馬可福音 5:21-34 & 莫樹堅牧師 & \hyperref[sec:031786]{耶穌,麻煩快點} & 2026-03-08 & \href{http://www.church.com.hk/acms/content.asp?site=cdc\&op=show\&type=product\&code=031786\&layout=sermon}{\texttt{031786}} \\
加拉太書 5:1 & 葛國光牧師 & \hyperref[sec:031787]{得自由之路 (華)} & 2026-03-08 & \href{http://www.church.com.hk/acms/content.asp?site=cdc\&op=show\&type=product\&code=031787\&layout=sermon}{\texttt{031787}} \\
使徒行傳 9:1-21 & 李黎秀兒師母 & \hyperref[sec:031788]{不一樣的人生:新方向} & 2026-03-08 & \href{http://www.church.com.hk/acms/content.asp?site=cdc\&op=show\&type=product\&code=031788\&layout=sermon}{\texttt{031788}} \\
約翰福音 3:16-21 & 駱柱成牧師 & \hyperref[sec:031789]{神愛世人} & 2026-03-08 & \href{http://www.church.com.hk/acms/content.asp?site=cdc\&op=show\&type=product\&code=031789\&layout=sermon}{\texttt{031789}} \\
創世記 29:9 & 黃克勤牧師博士 & \hyperref[sec:031790]{聖經人聖經事:正說話的時候} & 2026-03-06 & \href{http://www.church.com.hk/acms/content.asp?site=cdc\&op=show\&type=product\&code=031790\&layout=sermon}{\texttt{031790}} \\
馬太福音 18:21-35 & 曹偉彤牧師博士 & \hyperref[sec:031791]{未被寬恕?} & 2026-03-01 & \href{http://www.church.com.hk/acms/content.asp?site=cdc\&op=show\&type=product\&code=031791\&layout=sermon}{\texttt{031791}} \\
民數記 6:1-27 & 吳周婉貞傳道 & \hyperref[sec:031792]{上乘之選} & 2026-03-07 & \href{http://www.church.com.hk/acms/content.asp?site=cdc\&op=show\&type=product\&code=031792\&layout=sermon}{\texttt{031792}} \\
民數記 6:1-21 & 孫靜敏牧師 & \hyperref[sec:031793]{分別為聖:把自己交給神} & 2026-03-08 & \href{http://www.church.com.hk/acms/content.asp?site=cdc\&op=show\&type=product\&code=031793\&layout=sermon}{\texttt{031793}} \\
以弗所書 5:25-32 & 馬天嵐牧師 & \hyperref[sec:031794]{榮耀的基督和祂最愛的教會 (華/粵)} & 2026-03-08 & \href{http://www.church.com.hk/acms/content.asp?site=cdc\&op=show\&type=product\&code=031794\&layout=sermon}{\texttt{031794}} \\
約書亞記 7:1-26 & 李美儀傳道 & \hyperref[sec:031795]{暗中的罪,眾人的代價} & 2026-03-08 & \href{http://www.church.com.hk/acms/content.asp?site=cdc\&op=show\&type=product\&code=031795\&layout=sermon}{\texttt{031795}} \\
箴言 8:12-31 & 蕭壽華牧師 & \hyperref[sec:031796]{懇切尋求智慧的必尋見} & 2025-08-31 & \href{http://www.church.com.hk/acms/content.asp?site=cdc\&op=show\&type=product\&code=031796\&layout=sermon}{\texttt{031796}} \\
約翰福音 13:34-35 & 林誠信牧師 & \hyperref[sec:031797]{彼此相愛的標準與實踐} & 2026-03-08 & \href{http://www.church.com.hk/acms/content.asp?site=cdc\&op=show\&type=product\&code=031797\&layout=sermon}{\texttt{031797}} \\
創世記 29:20 & 黃克勤牧師博士 & \hyperref[sec:031798]{聖經人聖經事:七年如同幾天} & 2026-03-07 & \href{http://www.church.com.hk/acms/content.asp?site=cdc\&op=show\&type=product\&code=031798\&layout=sermon}{\texttt{031798}} \\
馬可福音 6:34-44 & 劉卓聰宣教士 & \hyperref[sec:031799]{上帝的邀請} & 2026-03-15 & \href{http://www.church.com.hk/acms/content.asp?site=cdc\&op=show\&type=product\&code=031799\&layout=sermon}{\texttt{031799}} \\
約書亞記 3:1-6 & 劉少康牧師 & \hyperref[sec:031800]{你有沒有期待看見神的作為?} & 2010-03-01 & \href{http://www.church.com.hk/acms/content.asp?site=cdc\&op=show\&type=product\&code=031800\&layout=sermon}{\texttt{031800}} \\
歌羅西書 2:1-5 & 黃紹權牧師 & \hyperref[sec:031801]{屬靈上做醜人} & 2026-03-15 & \href{http://www.church.com.hk/acms/content.asp?site=cdc\&op=show\&type=product\&code=031801\&layout=sermon}{\texttt{031801}} \\
使徒行傳 15:1-21 & 禤凱榮牧師 & \hyperref[sec:031802]{當紛爭再起時} & 2026-03-15 & \href{http://www.church.com.hk/acms/content.asp?site=cdc\&op=show\&type=product\&code=031802\&layout=sermon}{\texttt{031802}} \\
路加福音 10:38-42 & 鄺偉志牧師 & \hyperref[sec:031803]{為馬大殷勤服侍平反?} & 2026-03-15 & \href{http://www.church.com.hk/acms/content.asp?site=cdc\&op=show\&type=product\&code=031803\&layout=sermon}{\texttt{031803}} \\
創世記 12:1-9 & 任卓賢教授 & \hyperref[sec:031804]{召命人生} & 2026-03-15 & \href{http://www.church.com.hk/acms/content.asp?site=cdc\&op=show\&type=product\&code=031804\&layout=sermon}{\texttt{031804}} \\
創世記 37:12-36 & 黃克勤牧師博士 & \hyperref[sec:031805]{聖經人聖經事:上帝的手在帶領} & 2026-03-12 & \href{http://www.church.com.hk/acms/content.asp?site=cdc\&op=show\&type=product\&code=031805\&layout=sermon}{\texttt{031805}} \\
約翰福音 9:1-41 & 葉菁華教授 & \hyperref[sec:031806]{瞎眼今得看見} & 2026-03-15 & \href{http://www.church.com.hk/acms/content.asp?site=cdc\&op=show\&type=product\&code=031806\&layout=sermon}{\texttt{031806}} \\
列王記下 4:8-37 & 曹子光牧師 & \hyperref[sec:031807]{信心系列:信心奇女子} & 2026-03-08 & \href{http://www.church.com.hk/acms/content.asp?site=cdc\&op=show\&type=product\&code=031807\&layout=sermon}{\texttt{031807}} \\
約翰福音 17:20-23 & 江山偉牧師 & \hyperref[sec:031808]{建立健康的使命教會 (8) 有質素的相交} & 2026-03-15 & \href{http://www.church.com.hk/acms/content.asp?site=cdc\&op=show\&type=product\&code=031808\&layout=sermon}{\texttt{031808}} \\
約書亞記 3:7-17 & 劉少康牧師 & \hyperref[sec:031809]{神蹟的出現} & 2010-03-02 & \href{http://www.church.com.hk/acms/content.asp?site=cdc\&op=show\&type=product\&code=031809\&layout=sermon}{\texttt{031809}} \\
馬太福音 7:7-12 & 石偉漢牧師 & \hyperref[sec:031810]{哪裡會有好東西?} & 2026-03-15 & \href{http://www.church.com.hk/acms/content.asp?site=cdc\&op=show\&type=product\&code=031810\&layout=sermon}{\texttt{031810}} \\
創世記 39:21 & 黃克勤牧師博士 & \hyperref[sec:031811]{聖經人聖經事:耶和華在監房與約瑟同在} & 2026-03-13 & \href{http://www.church.com.hk/acms/content.asp?site=cdc\&op=show\&type=product\&code=031811\&layout=sermon}{\texttt{031811}} \\
歌羅西書 3:5 & 林天行牧師 & \hyperref[sec:031812]{貪婪,後果} & 2026-03-15 & \href{http://www.church.com.hk/acms/content.asp?site=cdc\&op=show\&type=product\&code=031812\&layout=sermon}{\texttt{031812}} \\
約書亞記 9:1-27 & 陳沛偉牧師 & \hyperref[sec:031813]{汲取教訓} & 2026-03-15 & \href{http://www.church.com.hk/acms/content.asp?site=cdc\&op=show\&type=product\&code=031813\&layout=sermon}{\texttt{031813}} \\
腓立比書 2:12-16 & 李適清博士 & \hyperref[sec:031814]{成就主的美意} & 2026-03-15 & \href{http://www.church.com.hk/acms/content.asp?site=cdc\&op=show\&type=product\&code=031814\&layout=sermon}{\texttt{031814}} \\
馬太福音 26:36-39 & 曾興才牧師 & \hyperref[sec:031815]{客西馬尼園的爭戰 (華)} & 2026-03-15 & \href{http://www.church.com.hk/acms/content.asp?site=cdc\&op=show\&type=product\&code=031815\&layout=sermon}{\texttt{031815}} \\
馬可福音 14:3-9 & 梁惠玲傳道 & \hyperref[sec:031816]{為主打破玉瓶的女人} & 2026-03-15 & \href{http://www.church.com.hk/acms/content.asp?site=cdc\&op=show\&type=product\&code=031816\&layout=sermon}{\texttt{031816}} \\
約翰福音 9:1-41 & 周曉沂傳道 & \hyperref[sec:031817]{能看見的瞎子} & 2026-03-15 & \href{http://www.church.com.hk/acms/content.asp?site=cdc\&op=show\&type=product\&code=031817\&layout=sermon}{\texttt{031817}} \\
腓立比書 1:3-11 & 王永揚牧師 & \hyperref[sec:031818]{竭力成為傳福音的團契} & 2026-03-15 & \href{http://www.church.com.hk/acms/content.asp?site=cdc\&op=show\&type=product\&code=031818\&layout=sermon}{\texttt{031818}} \\
創世記 37:12-36 & 黃克勤牧師博士 & \hyperref[sec:031819]{聖經人聖經事:上主的手在帶領} & 2026-03-14 & \href{http://www.church.com.hk/acms/content.asp?site=cdc\&op=show\&type=product\&code=031819\&layout=sermon}{\texttt{031819}} \\
約翰福音 6:35-59 & 林鍵牧師 & \hyperref[sec:031820]{榮耀君王:永不飢渴} & 2026-03-14 & \href{http://www.church.com.hk/acms/content.asp?site=cdc\&op=show\&type=product\&code=031820\&layout=sermon}{\texttt{031820}} \\
使徒行傳 1:4-8 & 李適清博士 & \hyperref[sec:031821]{生活就是使命} & 2026-01-25 & \href{http://www.church.com.hk/acms/content.asp?site=cdc\&op=show\&type=product\&code=031821\&layout=sermon}{\texttt{031821}} \\
以弗所書 1:1-23 & 劉少康牧師 & \hyperref[sec:031822]{以弗所書 (1) 導論,求啟示的禱告} & 2026-03-18 & \href{http://www.church.com.hk/acms/content.asp?site=cdc\&op=show\&type=product\&code=031822\&layout=sermon}{\texttt{031822}} \\
約翰福音 5:1-18 & 鍾智愷傳道 & \hyperref[sec:031823]{屬靈的力量} & 2026-03-15 & \href{http://www.church.com.hk/acms/content.asp?site=cdc\&op=show\&type=product\&code=031823\&layout=sermon}{\texttt{031823}} \\
哥林多後書 3:17 & 葛國光牧師 & \hyperref[sec:031824]{聖靈使你自由 (華)} & 2026-03-15 & \href{http://www.church.com.hk/acms/content.asp?site=cdc\&op=show\&type=product\&code=031824\&layout=sermon}{\texttt{031824}} \\
撒母耳記上 22:1-2 & 李思聰牧師 & \hyperref[sec:031825]{我們需要避難所} & 2026-03-15 & \href{http://www.church.com.hk/acms/content.asp?site=cdc\&op=show\&type=product\&code=031825\&layout=sermon}{\texttt{031825}} \\
傳道書 4:9-10 & 吳燕清傳道 & \hyperref[sec:031826]{多代同行 (2) PeerWalk (2) 誰是你的同行者?} & 2026-03-14 & \href{http://www.church.com.hk/acms/content.asp?site=cdc\&op=show\&type=product\&code=031826\&layout=sermon}{\texttt{031826}} \\
使徒行傳 16:6-15 & 關鄧婥寧傳道 & \hyperref[sec:031827]{甚麼是你的人生異象?} & 2026-03-15 & \href{http://www.church.com.hk/acms/content.asp?site=cdc\&op=show\&type=product\&code=031827\&layout=sermon}{\texttt{031827}} \\
雅各書 1:19-22 & 冼嘉琪傳道 & \hyperref[sec:031828]{有入有出的執道人 (粵/華)} & 2026-03-15 & \href{http://www.church.com.hk/acms/content.asp?site=cdc\&op=show\&type=product\&code=031828\&layout=sermon}{\texttt{031828}} \\
路加福音 15:1-2 & 梁誠德牧師 & \hyperref[sec:031829]{寶貴,保證,賬簿} & 2026-03-15 & \href{http://www.church.com.hk/acms/content.asp?site=cdc\&op=show\&type=product\&code=031829\&layout=sermon}{\texttt{031829}} \\
創世記 40:8 & 黃克勤牧師博士 & \hyperref[sec:031830]{聖經人聖經事:夢出於神} & 2026-03-15 & \href{http://www.church.com.hk/acms/content.asp?site=cdc\&op=show\&type=product\&code=031830\&layout=sermon}{\texttt{031830}} \\
詩篇 127:1-5 & 陳德蓉傳道 & \hyperref[sec:031831]{建造,看守,發展} & 2026-03-15 & \href{http://www.church.com.hk/acms/content.asp?site=cdc\&op=show\&type=product\&code=031831\&layout=sermon}{\texttt{031831}} \\
民數記 7:1-89 & 林誠信牧師 & \hyperref[sec:031832]{合神心意的奉獻} & 2026-03-15 & \href{http://www.church.com.hk/acms/content.asp?site=cdc\&op=show\&type=product\&code=031832\&layout=sermon}{\texttt{031832}} \\
約翰福音 2:1-12 & 許開明牧師 & \hyperref[sec:031833]{慈母的難題:“喜酒飲盡了!”} & 2026-03-15 & \href{http://www.church.com.hk/acms/content.asp?site=cdc\&op=show\&type=product\&code=031833\&layout=sermon}{\texttt{031833}} \\
以弗所書 4:1-7 & 林穎豪傳道 & \hyperref[sec:031834]{基督榮美的身體——歸屬,服侍,成長 (英/粵)} & 2026-03-15 & \href{http://www.church.com.hk/acms/content.asp?site=cdc\&op=show\&type=product\&code=031834\&layout=sermon}{\texttt{031834}} \\
以弗所書 2:11-22 & 陳義生傳道 & \hyperref[sec:031835]{福音的本質:在憂傷中帶來的盼望} & 2026-03-15 & \href{http://www.church.com.hk/acms/content.asp?site=cdc\&op=show\&type=product\&code=031835\&layout=sermon}{\texttt{031835}} \\
哥林多前書 3:1-15 & 王正傑牧師 & \hyperref[sec:031836]{你要食飯嗎?} & 2026-03-15 & \href{http://www.church.com.hk/acms/content.asp?site=cdc\&op=show\&type=product\&code=031836\&layout=sermon}{\texttt{031836}} \\
創世記 39:6-18 & 黃克勤牧師博士 & \hyperref[sec:031837]{聖經人聖經事:約瑟一時失慎} & 2026-03-16 & \href{http://www.church.com.hk/acms/content.asp?site=cdc\&op=show\&type=product\&code=031837\&layout=sermon}{\texttt{031837}} \\
雅各書 2:14-15 & 陳雅莉傳道 & \hyperref[sec:031838]{愛的施予} & 2026-03-22 & \href{http://www.church.com.hk/acms/content.asp?site=cdc\&op=show\&type=product\&code=031838\&layout=sermon}{\texttt{031838}} \\
使徒行傳 15:36-41 & 駱瑞麟傳道 & \hyperref[sec:031839]{路不同,道相同} & 2026-03-22 & \href{http://www.church.com.hk/acms/content.asp?site=cdc\&op=show\&type=product\&code=031839\&layout=sermon}{\texttt{031839}} \\
以賽亞書 60:1-18 & 陳穎康牧師 & \hyperref[sec:031840]{先嚐天國的味道} & 2026-03-22 & \href{http://www.church.com.hk/acms/content.asp?site=cdc\&op=show\&type=product\&code=031840\&layout=sermon}{\texttt{031840}} \\
以西結書 37:1-14 & 林家棋宣教師 & \hyperref[sec:031841]{我現在是「活」的嗎?} & 2026-03-21 & \href{http://www.church.com.hk/acms/content.asp?site=cdc\&op=show\&type=product\&code=031841\&layout=sermon}{\texttt{031841}} \\
創世記 40:23 & 黃克勤牧師博士 & \hyperref[sec:031842]{聖經人聖經事:人會忘記你} & 2026-03-18 & \href{http://www.church.com.hk/acms/content.asp?site=cdc\&op=show\&type=product\&code=031842\&layout=sermon}{\texttt{031842}} \\
詩篇 130:1-8 & 郭志榮牧師 & \hyperref[sec:031843]{從深處的覺醒} & 2026-03-22 & \href{http://www.church.com.hk/acms/content.asp?site=cdc\&op=show\&type=product\&code=031843\&layout=sermon}{\texttt{031843}} \\
詩篇 1:1-3 & 林耀楠牧師 & \hyperref[sec:031844]{建立健康的使命教會 (9) 愛慕遵從神話語} & 2026-03-22 & \href{http://www.church.com.hk/acms/content.asp?site=cdc\&op=show\&type=product\&code=031844\&layout=sermon}{\texttt{031844}} \\
創世記 1:1-3 & 卓爾君牧師 & \hyperref[sec:031845]{恢復生命的榮美 (華)} & 2026-03-22 & \href{http://www.church.com.hk/acms/content.asp?site=cdc\&op=show\&type=product\&code=031845\&layout=sermon}{\texttt{031845}} \\
約翰福音 5:17-21 & 招世超牧師 & \hyperref[sec:031846]{你為何仍然活著?} & 2026-03-22 & \href{http://www.church.com.hk/acms/content.asp?site=cdc\&op=show\&type=product\&code=031846\&layout=sermon}{\texttt{031846}} \\
以西結書 37:1-14 & 趙承恩宣教師 & \hyperref[sec:031847]{復活的生命} & 2026-03-22 & \href{http://www.church.com.hk/acms/content.asp?site=cdc\&op=show\&type=product\&code=031847\&layout=sermon}{\texttt{031847}} \\
創世記 41:16 & 黃克勤牧師博士 & \hyperref[sec:031848]{聖經人聖經事:這不在乎我} & 2026-03-19 & \href{http://www.church.com.hk/acms/content.asp?site=cdc\&op=show\&type=product\&code=031848\&layout=sermon}{\texttt{031848}} \\
撒迦利亞書 4:1-14 & 黃梅玲牧師 & \hyperref[sec:031849]{創始成終的三一神} & 2026-03-22 & \href{http://www.church.com.hk/acms/content.asp?site=cdc\&op=show\&type=product\&code=031849\&layout=sermon}{\texttt{031849}} \\
以弗所書 1:13-19 & 俞大鳴牧師 & \hyperref[sec:031850]{受了印記} & 2026-03-22 & \href{http://www.church.com.hk/acms/content.asp?site=cdc\&op=show\&type=product\&code=031850\&layout=sermon}{\texttt{031850}} \\
約翰福音 2:1-12 & 梁紫健傳道 & \hyperref[sec:031851]{行奇事的神} & 2026-03-22 & \href{http://www.church.com.hk/acms/content.asp?site=cdc\&op=show\&type=product\&code=031851\&layout=sermon}{\texttt{031851}} \\
雅各書 2:1-26 & 司徒贊傳道 & \hyperref[sec:031852]{在信裡沒有TACO} & 2026-03-22 & \href{http://www.church.com.hk/acms/content.asp?site=cdc\&op=show\&type=product\&code=031852\&layout=sermon}{\texttt{031852}} \\
以賽亞書 37:1-4 & 李應新牧師 & \hyperref[sec:031853]{危難當前} & 2026-03-22 & \href{http://www.church.com.hk/acms/content.asp?site=cdc\&op=show\&type=product\&code=031853\&layout=sermon}{\texttt{031853}} \\
創世記 41:32 & 黃克勤牧師博士 & \hyperref[sec:031854]{聖經人聖經事:神命定的,必速速成就} & 2026-03-20 & \href{http://www.church.com.hk/acms/content.asp?site=cdc\&op=show\&type=product\&code=031854\&layout=sermon}{\texttt{031854}} \\
約翰福音 6:59-71 & 林鍵牧師 & \hyperref[sec:031855]{榮耀君王:永生之道} & 2026-03-21 & \href{http://www.church.com.hk/acms/content.asp?site=cdc\&op=show\&type=product\&code=031855\&layout=sermon}{\texttt{031855}} \\
以弗所書 2:1-4:32 & 劉少康牧師 & \hyperref[sec:031856]{以弗所書 (2) 教會——基督的身體} & 2026-03-25 & \href{http://www.church.com.hk/acms/content.asp?site=cdc\&op=show\&type=product\&code=031856\&layout=sermon}{\texttt{031856}} \\
詩篇 34:1-7 & 高志翔牧師 & \hyperref[sec:031857]{當我害怕時 (華)} & 2026-03-22 & \href{http://www.church.com.hk/acms/content.asp?site=cdc\&op=show\&type=product\&code=031857\&layout=sermon}{\texttt{031857}} \\
撒母耳記上 23:1-29 & 李思聰牧師 & \hyperref[sec:031858]{面對困境有出路} & 2026-03-22 & \href{http://www.church.com.hk/acms/content.asp?site=cdc\&op=show\&type=product\&code=031858\&layout=sermon}{\texttt{031858}} \\
約書亞記 14:6-12 & 梁明財牧師 & \hyperref[sec:031859]{活出使命的跟從者——迦勒 II (華)} & 2026-03-22 & \href{http://www.church.com.hk/acms/content.asp?site=cdc\&op=show\&type=product\&code=031859\&layout=sermon}{\texttt{031859}} \\
約書亞記 10:1-12:24 & 方兆雄牧師 & \hyperref[sec:031860]{有神帥領的爭戰} & 2026-03-22 & \href{http://www.church.com.hk/acms/content.asp?site=cdc\&op=show\&type=product\&code=031860\&layout=sermon}{\texttt{031860}} \\
創世記 41:38 & 黃克勤牧師博士 & \hyperref[sec:031861]{聖經人聖經事:神的靈在他裡頭} & 2026-03-21 & \href{http://www.church.com.hk/acms/content.asp?site=cdc\&op=show\&type=product\&code=031861\&layout=sermon}{\texttt{031861}} \\
羅馬書 16:17-27 & 黃浩君傳道 & \hyperref[sec:031862]{在福音中被堅固與合一} & 2026-03-22 & \href{http://www.church.com.hk/acms/content.asp?site=cdc\&op=show\&type=product\&code=031862\&layout=sermon}{\texttt{031862}} \\
哥林多前書 13:4-7 & 張健庭牧師博士 & \hyperref[sec:031863]{重塑生命系列 (7) 妒忌與愛} & 2026-03-22 & \href{http://www.church.com.hk/acms/content.asp?site=cdc\&op=show\&type=product\&code=031863\&layout=sermon}{\texttt{031863}} \\
路加福音 5:33-35 & 黃健羒牧師 & \hyperref[sec:031864]{為主禁食} & 2026-03-22 & \href{http://www.church.com.hk/acms/content.asp?site=cdc\&op=show\&type=product\&code=031864\&layout=sermon}{\texttt{031864}} \\
約翰福音 20:19-29 & 麥慕思牧師 & \hyperref[sec:031865]{多代同行 (2) PeerWalk (3) 你們的……與我何干?} & 2026-03-21 & \href{http://www.church.com.hk/acms/content.asp?site=cdc\&op=show\&type=product\&code=031865\&layout=sermon}{\texttt{031865}} \\
詩篇 1:1-6 & 高銘謙牧師博士 & \hyperref[sec:031866]{畫夜默想之福} & 2026-03-22 & \href{http://www.church.com.hk/acms/content.asp?site=cdc\&op=show\&type=product\&code=031866\&layout=sermon}{\texttt{031866}} \\
創世記 45:5 & 黃克勤牧師博士 & \hyperref[sec:031867]{聖經人聖經事:是上帝差我先來} & 2026-03-22 & \href{http://www.church.com.hk/acms/content.asp?site=cdc\&op=show\&type=product\&code=031867\&layout=sermon}{\texttt{031867}} \\
撒母耳記上 8:4-10 & 袁惠鈞牧師博士 & \hyperref[sec:031868]{以色列興衰史:合一王朝系列 (1)} & 2026-03-15 & \href{http://www.church.com.hk/acms/content.asp?site=cdc\&op=show\&type=product\&code=031868\&layout=sermon}{\texttt{031868}} \\
約翰福音 3:22-30 & 駱柱成牧師 & \hyperref[sec:031869]{事奉的榜樣} & 2026-03-22 & \href{http://www.church.com.hk/acms/content.asp?site=cdc\&op=show\&type=product\&code=031869\&layout=sermon}{\texttt{031869}} \\
箴言 18:24 & 曹子光牧師 & \hyperref[sec:031870]{摯愛的知己好友} & 2026-03-22 & \href{http://www.church.com.hk/acms/content.asp?site=cdc\&op=show\&type=product\&code=031870\&layout=sermon}{\texttt{031870}} \\
路加福音 18:9-14 & 梁成裕牧師 & \hyperref[sec:031871]{法利賽人與稅利禱告的真偽:生命的反思} & 2026-03-22 & \href{http://www.church.com.hk/acms/content.asp?site=cdc\&op=show\&type=product\&code=031871\&layout=sermon}{\texttt{031871}} \\
創世記 49:9 & 黃克勤牧師博士 & \hyperref[sec:031872]{聖經人聖經事:猶大是個小獅子} & 2026-03-23 & \href{http://www.church.com.hk/acms/content.asp?site=cdc\&op=show\&type=product\&code=031872\&layout=sermon}{\texttt{031872}} \\
詩篇 8:1-9 & 梁燕城博士 & \hyperref[sec:031873]{從美學沉思宇宙與人生} & 2026-03-22 & \href{http://www.church.com.hk/acms/content.asp?site=cdc\&op=show\&type=product\&code=031873\&layout=sermon}{\texttt{031873}} \\
哥林多前書 1:26-31 & 梁錫強醫生 & \hyperref[sec:031874]{小人物,大使命} & 2026-03-29 & \href{http://www.church.com.hk/acms/content.asp?site=cdc\&op=show\&type=product\&code=031874\&layout=sermon}{\texttt{031874}} \\
使徒行傳 16:6-15 & 楊芹牧師 & \hyperref[sec:031875]{神的美意 (華)} & 2026-03-29 & \href{http://www.church.com.hk/acms/content.asp?site=cdc\&op=show\&type=product\&code=031875\&layout=sermon}{\texttt{031875}} \\
啟示錄 1:7-8 & 譚綺敏傳道 & \hyperref[sec:031876]{耶穌基督的啟示} & 2026-03-29 & \href{http://www.church.com.hk/acms/content.asp?site=cdc\&op=show\&type=product\&code=031876\&layout=sermon}{\texttt{031876}} \\
創世記 50:19 & 黃克勤牧師博士 & \hyperref[sec:031877]{聖經人聖經事:我豈能代替神呢?} & 2026-03-24 & \href{http://www.church.com.hk/acms/content.asp?site=cdc\&op=show\&type=product\&code=031877\&layout=sermon}{\texttt{031877}} \\
腓立比書 2:5-11 & 蘇善美宣教師 & \hyperref[sec:031878]{以基督的心為心} & 2026-03-29 & \href{http://www.church.com.hk/acms/content.asp?site=cdc\&op=show\&type=product\&code=031878\&layout=sermon}{\texttt{031878}} \\
馬太福音 21:1-11 & 林豪恩先生 & \hyperref[sec:031879]{路不能亂走,車不能亂坐} & 2026-03-29 & \href{http://www.church.com.hk/acms/content.asp?site=cdc\&op=show\&type=product\&code=031879\&layout=sermon}{\texttt{031879}} \\
馬太福音 27:45-50 & 石偉漢牧師 & \hyperref[sec:031880]{我的上帝!我的上帝!為甚麼離棄我?} & 2026-03-29 & \href{http://www.church.com.hk/acms/content.asp?site=cdc\&op=show\&type=product\&code=031880\&layout=sermon}{\texttt{031880}} \\
歌羅西書 2:6-15 & 黃紹權牧師 & \hyperref[sec:031881]{超脫人生的塑造基礎 (1)} & 2026-03-29 & \href{http://www.church.com.hk/acms/content.asp?site=cdc\&op=show\&type=product\&code=031881\&layout=sermon}{\texttt{031881}} \\
羅馬書 8:28-30 & 吳磐德牧師 & \hyperref[sec:031882]{建立健康的使命教會 (10) 塑造活像基督的生命} & 2026-03-29 & \href{http://www.church.com.hk/acms/content.asp?site=cdc\&op=show\&type=product\&code=031882\&layout=sermon}{\texttt{031882}} \\
創世記 13:6-18 & 江惠明牧師 & \hyperref[sec:031883]{世上最「正」的眼鏡} & 2026-03-29 & \href{http://www.church.com.hk/acms/content.asp?site=cdc\&op=show\&type=product\&code=031883\&layout=sermon}{\texttt{031883}} \\
出埃及記 2:10 & 黃克勤牧師博士 & \hyperref[sec:031884]{聖經人聖經事:從水裡拉出來} & 2026-03-25 & \href{http://www.church.com.hk/acms/content.asp?site=cdc\&op=show\&type=product\&code=031884\&layout=sermon}{\texttt{031884}} \\
約書亞記 18:1-10 & 劉銳鋒牧師 & \hyperref[sec:031885]{你們耽延不去得,要到幾時呢?} & 2026-03-29 & \href{http://www.church.com.hk/acms/content.asp?site=cdc\&op=show\&type=product\&code=031885\&layout=sermon}{\texttt{031885}} \\
以弗所書 1:15-19 & 劉彤牧師 & \hyperref[sec:031886]{成為轉化世代的器皿 (華)} & 2026-03-29 & \href{http://www.church.com.hk/acms/content.asp?site=cdc\&op=show\&type=product\&code=031886\&layout=sermon}{\texttt{031886}} \\
馬可福音 10:42-45 & 招世超牧師 & \hyperref[sec:031887]{耶穌改變世界的方法} & 2026-03-29 & \href{http://www.church.com.hk/acms/content.asp?site=cdc\&op=show\&type=product\&code=031887\&layout=sermon}{\texttt{031887}} \\
傳道書 3:16-4:3 & 譚文倩莊傳道 & \hyperref[sec:031888]{公義缺席} & 2026-03-29 & \href{http://www.church.com.hk/acms/content.asp?site=cdc\&op=show\&type=product\&code=031888\&layout=sermon}{\texttt{031888}} \\
出埃及記 3:7 & 黃克勤牧師博士 & \hyperref[sec:031889]{聖經人聖經事:我原知道他們的痛苦} & 2026-03-28 & \href{http://www.church.com.hk/acms/content.asp?site=cdc\&op=show\&type=product\&code=031889\&layout=sermon}{\texttt{031889}} \\
約翰福音 11:17-32 & 鄭子禧傳道 & \hyperref[sec:031890]{俯伏主耶穌腳前} & 2026-03-29 & \href{http://www.church.com.hk/acms/content.asp?site=cdc\&op=show\&type=product\&code=031890\&layout=sermon}{\texttt{031890}} \\
腓立比書 3:17-4:7 & 盧偉心弟兄 & \hyperref[sec:031891]{腓書回顧 (3) 天民} & 2026-03-28 & \href{http://www.church.com.hk/acms/content.asp?site=cdc\&op=show\&type=product\&code=031891\&layout=sermon}{\texttt{031891}} \\
羅馬書 14:1-23 & 郭建恩傳道 & \hyperref[sec:031892]{論斷,還是接納} & 2026-03-29 & \href{http://www.church.com.hk/acms/content.asp?site=cdc\&op=show\&type=product\&code=031892\&layout=sermon}{\texttt{031892}} \\
馬太福音 21:1-9 & 葛國光牧師 & \hyperref[sec:031893]{謙和的君王 (華)} & 2026-03-29 & \href{http://www.church.com.hk/acms/content.asp?site=cdc\&op=show\&type=product\&code=031893\&layout=sermon}{\texttt{031893}} \\
馬太福音 21:1-11 & 黃馮潔霞師母 & \hyperref[sec:031894]{看哪,你的王來到你這裡!} & 2026-03-29 & \href{http://www.church.com.hk/acms/content.asp?site=cdc\&op=show\&type=product\&code=031894\&layout=sermon}{\texttt{031894}} \\
出埃及記 3:12 & 黃克勤牧師博士 & \hyperref[sec:031895]{聖經人聖經事:我必與你同在} & 2026-03-29 & \href{http://www.church.com.hk/acms/content.asp?site=cdc\&op=show\&type=product\&code=031895\&layout=sermon}{\texttt{031895}} \\
箴言 27:17 & 沈炳華傳道 & \hyperref[sec:031896]{多代同行 (2) PeerWalk IV Iron Sharpens Iron} & 2026-03-28 & \href{http://www.church.com.hk/acms/content.asp?site=cdc\&op=show\&type=product\&code=031896\&layout=sermon}{\texttt{031896}} \\
申命記 6:6-9 & 何羅乃萱師母 & \hyperref[sec:031897]{跨代傳承} & 2026-03-29 & \href{http://www.church.com.hk/acms/content.asp?site=cdc\&op=show\&type=product\&code=031897\&layout=sermon}{\texttt{031897}} \\
雅各書 5:7-12 & 秦潤祥牧師 & \hyperref[sec:031898]{百忍成聖} & 2026-03-29 & \href{http://www.church.com.hk/acms/content.asp?site=cdc\&op=show\&type=product\&code=031898\&layout=sermon}{\texttt{031898}} \\
馬太福音 24:15-24 & 袁惠鈞牧師博士 & \hyperref[sec:031899]{末世三個伊斯蘭教的徵兆與真假救主} & 2026-03-29 & \href{http://www.church.com.hk/acms/content.asp?site=cdc\&op=show\&type=product\&code=031899\&layout=sermon}{\texttt{031899}} \\
箴言 6:6-11 & 張健庭牧師博士 & \hyperref[sec:031900]{重塑生命系列 (8) 懶惰與安息} & 2026-03-29 & \href{http://www.church.com.hk/acms/content.asp?site=cdc\&op=show\&type=product\&code=031900\&layout=sermon}{\texttt{031900}} \\
出埃及記 3:15 & 黃克勤牧師博士 & \hyperref[sec:031901]{聖經人聖經事:耶和華是我的名} & 2026-03-30 & \href{http://www.church.com.hk/acms/content.asp?site=cdc\&op=show\&type=product\&code=031901\&layout=sermon}{\texttt{031901}} \\
民數記 8:1-26 & 甄盧鳳雯傳道 & \hyperref[sec:031902]{分別為聖,轉眼為耳} & 2026-03-28 & \href{http://www.church.com.hk/acms/content.asp?site=cdc\&op=show\&type=product\&code=031902\&layout=sermon}{\texttt{031902}} \\
民數記 8:1-26 & 吳安強牧師 & \hyperref[sec:031903]{分別為聖的服侍} & 2026-03-29 & \href{http://www.church.com.hk/acms/content.asp?site=cdc\&op=show\&type=product\&code=031903\&layout=sermon}{\texttt{031903}} \\
詩篇 95:1-11 & 歐世達牧師 & \hyperref[sec:031904]{洪水中的約——安睡在浮動的動物園 (英/粵)} & 2026-03-29 & \href{http://www.church.com.hk/acms/content.asp?site=cdc\&op=show\&type=product\&code=031904\&layout=sermon}{\texttt{031904}} \\
約翰福音 12:12-33 & 周曉暉牧師 & \hyperref[sec:031905]{當你高唱迎接榮耀君王} & 2026-03-29 & \href{http://www.church.com.hk/acms/content.asp?site=cdc\&op=show\&type=product\&code=031905\&layout=sermon}{\texttt{031905}} \\
出埃及記 4:2 & 黃克勤牧師博士 & \hyperref[sec:031906]{聖經人聖經事:你手裡是甚麼} & 2026-03-31 & \href{http://www.church.com.hk/acms/content.asp?site=cdc\&op=show\&type=product\&code=031906\&layout=sermon}{\texttt{031906}} \\
以弗所書 5:1-6:24 & 劉少康牧師 & \hyperref[sec:031907]{以弗所書 (3) 仇敵與全副軍裝} & 2026-04-01 & \href{http://www.church.com.hk/acms/content.asp?site=cdc\&op=show\&type=product\&code=031907\&layout=sermon}{\texttt{031907}} \\
以賽亞書 50:4-9 & 林佩華姊妹 & \hyperref[sec:031908]{誰能定我有罪呢?} & 2026-03-29 & \href{http://www.church.com.hk/acms/content.asp?site=cdc\&op=show\&type=product\&code=031908\&layout=sermon}{\texttt{031908}} \\
創世記 24:50-67 & 翁曉怡傳道 & \hyperref[sec:031909]{活於應許的故事裡} & 2026-03-28 & \href{http://www.church.com.hk/acms/content.asp?site=cdc\&op=show\&type=product\&code=031909\&layout=sermon}{\texttt{031909}} \\
馬可福音 15:42-16:8 & 禤凱榮牧師 & \hyperref[sec:031910]{從恐懼中走出來} & 2026-04-05 & \href{http://www.church.com.hk/acms/content.asp?site=cdc\&op=show\&type=product\&code=031910\&layout=sermon}{\texttt{031910}} \\
路加福音 24:36-53 & 譚綺敏傳道 & \hyperref[sec:031911]{復活主帶來翻轉} & 2026-04-05 & \href{http://www.church.com.hk/acms/content.asp?site=cdc\&op=show\&type=product\&code=031911\&layout=sermon}{\texttt{031911}} \\
以賽亞書 52:13-53:12 & 高銘謙牧師博士 & \hyperref[sec:031912]{擔當我們憂患的主} & 2026-04-05 & \href{http://www.church.com.hk/acms/content.asp?site=cdc\&op=show\&type=product\&code=031912\&layout=sermon}{\texttt{031912}} \\
出埃及記 4:11 & 黃克勤牧師博士 & \hyperref[sec:031913]{聖經人聖經事:誰造人的口呢?} & 2026-04-01 & \href{http://www.church.com.hk/acms/content.asp?site=cdc\&op=show\&type=product\&code=031913\&layout=sermon}{\texttt{031913}} \\
使徒行傳 10:34-43 & 林玉英牧師 & \hyperref[sec:031914]{跨越邊界的愛} & 2026-04-05 & \href{http://www.church.com.hk/acms/content.asp?site=cdc\&op=show\&type=product\&code=031914\&layout=sermon}{\texttt{031914}} \\
哥林多前書 15:1-11 & 王正傑牧師 & \hyperref[sec:031915]{基督復活的大能} & 2026-04-05 & \href{http://www.church.com.hk/acms/content.asp?site=cdc\&op=show\&type=product\&code=031915\&layout=sermon}{\texttt{031915}} \\
哥林多前書 15:50-58 & 黃紹權牧師 & \hyperref[sec:031916]{叫人全然改變的復活主} & 2026-04-05 & \href{http://www.church.com.hk/acms/content.asp?site=cdc\&op=show\&type=product\&code=031916\&layout=sermon}{\texttt{031916}} \\
路加福音 24:13-35 & 李適清博士 & \hyperref[sec:031917]{遇見復活的主 (粵/英)} & 2026-04-05 & \href{http://www.church.com.hk/acms/content.asp?site=cdc\&op=show\&type=product\&code=031917\&layout=sermon}{\texttt{031917}} \\
哥林多前書 15:1-8 & 鄭兆奇牧師 & \hyperref[sec:031918]{復活的主} & 2026-04-05 & \href{http://www.church.com.hk/acms/content.asp?site=cdc\&op=show\&type=product\&code=031918\&layout=sermon}{\texttt{031918}} \\
約翰福音 19:38-42 & 林家棋宣教師 & \hyperref[sec:031919]{明唔明,唔明就黎明} & 2026-03-28 & \href{http://www.church.com.hk/acms/content.asp?site=cdc\&op=show\&type=product\&code=031919\&layout=sermon}{\texttt{031919}} \\
出埃及記 8:19 & 黃克勤牧師博士 & \hyperref[sec:031920]{聖經人聖經事:這是神的手段} & 2026-04-02 & \href{http://www.church.com.hk/acms/content.asp?site=cdc\&op=show\&type=product\&code=031920\&layout=sermon}{\texttt{031920}} \\
馬太福音 28:1-10 & 陳沛偉牧師 & \hyperref[sec:031921]{空墳墓的訊息} & 2026-04-05 & \href{http://www.church.com.hk/acms/content.asp?site=cdc\&op=show\&type=product\&code=031921\&layout=sermon}{\texttt{031921}} \\
馬太福音 28:5-6 & 劉彤牧師 & \hyperref[sec:031922]{主復活帶來七大恩典 (華)} & 2026-04-05 & \href{http://www.church.com.hk/acms/content.asp?site=cdc\&op=show\&type=product\&code=031922\&layout=sermon}{\texttt{031922}} \\
馬太福音 16:25 & 李業基牧師 & \hyperref[sec:031923]{天國戰略} & 2026-04-08 & \href{http://www.church.com.hk/acms/content.asp?site=cdc\&op=show\&type=product\&code=031923\&layout=sermon}{\texttt{031923}} \\
約翰福音 21:1-19 & 俞大鳴牧師 & \hyperref[sec:031924]{全人得痊癒} & 2026-04-05 & \href{http://www.church.com.hk/acms/content.asp?site=cdc\&op=show\&type=product\&code=031924\&layout=sermon}{\texttt{031924}} \\
出埃及記 10:21 & 黃克勤牧師博士 & \hyperref[sec:031925]{聖經人聖經事:摩西向天伸杖} & 2026-04-03 & \href{http://www.church.com.hk/acms/content.asp?site=cdc\&op=show\&type=product\&code=031925\&layout=sermon}{\texttt{031925}} \\
以賽亞書 53:4-5 & 李健明牧師 & \hyperref[sec:031926]{苦難中的出路} & 2026-04-05 & \href{http://www.church.com.hk/acms/content.asp?site=cdc\&op=show\&type=product\&code=031926\&layout=sermon}{\texttt{031926}} \\
以賽亞書 52:13-53:6 & 李應新牧師 & \hyperref[sec:031927]{受苦義僕} & 2026-04-05 & \href{http://www.church.com.hk/acms/content.asp?site=cdc\&op=show\&type=product\&code=031927\&layout=sermon}{\texttt{031927}} \\
羅馬書 1:1-4 & 葛國光牧師 & \hyperref[sec:031928]{基督復活的大能 (華)} & 2026-04-05 & \href{http://www.church.com.hk/acms/content.asp?site=cdc\&op=show\&type=product\&code=031928\&layout=sermon}{\texttt{031928}} \\
彼得前書 1:3-9 & 李思聰牧師 & \hyperref[sec:031929]{基督復活的大能} & 2026-04-05 & \href{http://www.church.com.hk/acms/content.asp?site=cdc\&op=show\&type=product\&code=031929\&layout=sermon}{\texttt{031929}} \\
哥林多前書 8:6 & 袁惠鈞牧師博士 & \hyperref[sec:031930]{不信三位一體的神能得救嗎?} & 2026-04-05 & \href{http://www.church.com.hk/acms/content.asp?site=cdc\&op=show\&type=product\&code=031930\&layout=sermon}{\texttt{031930}} \\
出埃及記 12:30 & 黃克勤牧師博士 & \hyperref[sec:031931]{聖經人聖經事:無一家不死一個人} & 2026-04-04 & \href{http://www.church.com.hk/acms/content.asp?site=cdc\&op=show\&type=product\&code=031931\&layout=sermon}{\texttt{031931}} \\
使徒行傳 12:1-25 & 林鍵牧師 & \hyperref[sec:031932]{越獄風雲 (1)} & 2026-04-04 & \href{http://www.church.com.hk/acms/content.asp?site=cdc\&op=show\&type=product\&code=031932\&layout=sermon}{\texttt{031932}} \\
約翰福音 20:1-18 & 高國雄牧師 & \hyperref[sec:031933]{基督今復活} & 2026-04-05 & \href{http://www.church.com.hk/acms/content.asp?site=cdc\&op=show\&type=product\&code=031933\&layout=sermon}{\texttt{031933}} \\
  & 張健庭牧師博士 & \hyperref[sec:031934]{當黑暗籠罩世界:神在十字架上的回應 (粵/華)} & 2026-04-03 & \href{http://www.church.com.hk/acms/content.asp?site=cdc\&op=show\&type=product\&code=031934\&layout=sermon}{\texttt{031934}} \\
哥林多前書 15:1-4 & 姚臣泰長老 & \hyperref[sec:031935]{復活和基督徒的盼望 (華/粵)} & 2026-04-05 & \href{http://www.church.com.hk/acms/content.asp?site=cdc\&op=show\&type=product\&code=031935\&layout=sermon}{\texttt{031935}} \\
腓立比書 3:10-11 & 杜崇基牧師 & \hyperref[sec:031936]{基督復活的大能} & 2026-04-05 & \href{http://www.church.com.hk/acms/content.asp?site=cdc\&op=show\&type=product\&code=031936\&layout=sermon}{\texttt{031936}} \\
出埃及記 4:23 & 黃克勤牧師博士 & \hyperref[sec:031937]{聖經人聖經事:為何要殺長子} & 2026-04-05 & \href{http://www.church.com.hk/acms/content.asp?site=cdc\&op=show\&type=product\&code=031937\&layout=sermon}{\texttt{031937}} \\
帖撒羅尼迦前書 5:16-19 & 謝又生牧師 & \hyperref[sec:031938]{復活:恢復感恩的心} & 2026-04-05 & \href{http://www.church.com.hk/acms/content.asp?site=cdc\&op=show\&type=product\&code=031938\&layout=sermon}{\texttt{031938}} \\
約翰福音 20:19-23 & 葉慧敏傳道 & \hyperref[sec:031939]{門關了,主在中間} & 2026-04-05 & \href{http://www.church.com.hk/acms/content.asp?site=cdc\&op=show\&type=product\&code=031939\&layout=sermon}{\texttt{031939}} \\
腓立比書 1:1-4:23 & 劉少康牧師 & \hyperref[sec:031940]{腓立比書:喜樂的秘訣 (1) 全書概論} & 2026-04-08 & \href{http://www.church.com.hk/acms/content.asp?site=cdc\&op=show\&type=product\&code=031940\&layout=sermon}{\texttt{031940}} \\
約翰福音 20:24-29 & 甘雪芬傳道 & \hyperref[sec:031941]{復活主榮耀的記號} & 2026-04-05 & \href{http://www.church.com.hk/acms/content.asp?site=cdc\&op=show\&type=product\&code=031941\&layout=sermon}{\texttt{031941}} \\
約翰福音 21:15-19 & 李炎強牧師 & \hyperref[sec:031942]{愛得更深} & 2026-04-04 & \href{http://www.church.com.hk/acms/content.asp?site=cdc\&op=show\&type=product\&code=031942\&layout=sermon}{\texttt{031942}} \\
出埃及記 12:42 & 黃克勤牧師博士 & \hyperref[sec:031943]{聖經人聖經事:這夜是耶和華的夜} & 2026-04-06 & \href{http://www.church.com.hk/acms/content.asp?site=cdc\&op=show\&type=product\&code=031943\&layout=sermon}{\texttt{031943}} \\
詩篇 90:1-17 & 廖諾文牧師 & \hyperref[sec:031944]{原來我們一直錯判了人生} & 2026-04-05 & \href{http://www.church.com.hk/acms/content.asp?site=cdc\&op=show\&type=product\&code=031944\&layout=sermon}{\texttt{031944}} \\
詩篇 16:1-2 & 呂君望牧師 & \hyperref[sec:031945]{投靠神的人有福了} & 2026-04-05 & \href{http://www.church.com.hk/acms/content.asp?site=cdc\&op=show\&type=product\&code=031945\&layout=sermon}{\texttt{031945}} \\
列王記上 8:37-45 & 潘建堂牧師 & \hyperref[sec:031946]{禱告讓主名被高舉} & 2025-09-07 & \href{http://www.church.com.hk/acms/content.asp?site=cdc\&op=show\&type=product\&code=031946\&layout=sermon}{\texttt{031946}} \\
歌羅西書 2:6-15 & 黃紹權牧師 & \hyperref[sec:031947]{超脫人生的塑造基礎 (2)} & 2026-04-12 & \href{http://www.church.com.hk/acms/content.asp?site=cdc\&op=show\&type=product\&code=031947\&layout=sermon}{\texttt{031947}} \\
約翰福音 21:1-14 & 禤凱榮牧師 & \hyperref[sec:031948]{主為你做早餐} & 2026-04-12 & \href{http://www.church.com.hk/acms/content.asp?site=cdc\&op=show\&type=product\&code=031948\&layout=sermon}{\texttt{031948}} \\
出埃及記 12:11 & 黃克勤牧師博士 & \hyperref[sec:031949]{聖經人聖經事:這是耶和華的逾越節} & 2026-04-07 & \href{http://www.church.com.hk/acms/content.asp?site=cdc\&op=show\&type=product\&code=031949\&layout=sermon}{\texttt{031949}} \\
使徒行傳 2:22-32 & 趙承恩宣教師 & \hyperref[sec:031950]{我們都是這事的見證人} & 2026-04-12 & \href{http://www.church.com.hk/acms/content.asp?site=cdc\&op=show\&type=product\&code=031950\&layout=sermon}{\texttt{031950}} \\
約翰福音 20:19-31 & 李均熊牧師博士 & \hyperref[sec:031951]{我信身體復活} & 2026-04-12 & \href{http://www.church.com.hk/acms/content.asp?site=cdc\&op=show\&type=product\&code=031951\&layout=sermon}{\texttt{031951}} \\
馬可福音 16:1-8 & 黃順成牧師 & \hyperref[sec:031952]{石頭滾開的奇蹟} & 2026-04-12 & \href{http://www.church.com.hk/acms/content.asp?site=cdc\&op=show\&type=product\&code=031952\&layout=sermon}{\texttt{031952}} \\
路加福音 24:13-32 & 石偉漢牧師 & \hyperref[sec:031953]{愁容不再,火熱常在} & 2026-04-12 & \href{http://www.church.com.hk/acms/content.asp?site=cdc\&op=show\&type=product\&code=031953\&layout=sermon}{\texttt{031953}} \\
馬可福音 8:31-38 & 許永勝傳道 & \hyperref[sec:031954]{我要與基督同釘十字架 (粵/華)} & 2026-04-05 & \href{http://www.church.com.hk/acms/content.asp?site=cdc\&op=show\&type=product\&code=031954\&layout=sermon}{\texttt{031954}} \\
詩篇 25:1-22 & 江山偉牧師 & \hyperref[sec:031955]{建立健康的使命教會 (11) 辨識並遵從神的引導} & 2026-04-12 & \href{http://www.church.com.hk/acms/content.asp?site=cdc\&op=show\&type=product\&code=031955\&layout=sermon}{\texttt{031955}} \\
啟示錄 2:1-7 & 林穎豪傳道 & \hyperref[sec:031956]{以弗所:信仰純正卻愛心冷淡的教會} & 2026-04-12 & \href{http://www.church.com.hk/acms/content.asp?site=cdc\&op=show\&type=product\&code=031956\&layout=sermon}{\texttt{031956}} \\
出埃及記 13:8 & 黃克勤牧師博士 & \hyperref[sec:031957]{聖經人聖經事:要告訴你的兒子} & 2026-04-08 & \href{http://www.church.com.hk/acms/content.asp?site=cdc\&op=show\&type=product\&code=031957\&layout=sermon}{\texttt{031957}} \\
撒母耳記上 1:1-11 & 李玥坤傳道 & \hyperref[sec:031958]{一心•倚靠} & 2026-04-12 & \href{http://www.church.com.hk/acms/content.asp?site=cdc\&op=show\&type=product\&code=031958\&layout=sermon}{\texttt{031958}} \\
歌羅西書 4:1-6 & 梁永善牧師 & \hyperref[sec:031959]{從歌羅西書看成熟生命表現} & 2026-04-12 & \href{http://www.church.com.hk/acms/content.asp?site=cdc\&op=show\&type=product\&code=031959\&layout=sermon}{\texttt{031959}} \\
路加福音 11:1-13 & 鄺偉志牧師 & \hyperref[sec:031960]{禱告蒙應允的秘訣?} & 2026-04-12 & \href{http://www.church.com.hk/acms/content.asp?site=cdc\&op=show\&type=product\&code=031960\&layout=sermon}{\texttt{031960}} \\
創世記 22:15-19 & 陳崇基牧師 & \hyperref[sec:031961]{海沙} & 2026-04-12 & \href{http://www.church.com.hk/acms/content.asp?site=cdc\&op=show\&type=product\&code=031961\&layout=sermon}{\texttt{031961}} \\
出埃及記 13:18 & 黃克勤牧師博士 & \hyperref[sec:031962]{聖經人聖經事:神領百姓繞道} & 2026-04-08 & \href{http://www.church.com.hk/acms/content.asp?site=cdc\&op=show\&type=product\&code=031962\&layout=sermon}{\texttt{031962}} \\
馬可福音 6:31-34 & 周弘毅牧師 & \hyperref[sec:031963]{無處不在的基督} & 2026-04-11 & \href{http://www.church.com.hk/acms/content.asp?site=cdc\&op=show\&type=product\&code=031963\&layout=sermon}{\texttt{031963}} \\
以西結書 36:26-27 & 林羅麗斯牧師 & \hyperref[sec:031964]{全人更新} & 2026-04-12 & \href{http://www.church.com.hk/acms/content.asp?site=cdc\&op=show\&type=product\&code=031964\&layout=sermon}{\texttt{031964}} \\
路加福音 14:1-24 & 俞大鳴牧師 & \hyperref[sec:031965]{筵席上的教導} & 2026-04-12 & \href{http://www.church.com.hk/acms/content.asp?site=cdc\&op=show\&type=product\&code=031965\&layout=sermon}{\texttt{031965}} \\
士師記 1:1-2:5 & 梁永善牧師 & \hyperref[sec:031966]{士師記析讀 (1) 鳥瞰士師記} & 2026-04-15 & \href{http://www.church.com.hk/acms/content.asp?site=cdc\&op=show\&type=product\&code=031966\&layout=sermon}{\texttt{031966}} \\
路加福音 24:13-21 & 曾興才牧師 & \hyperref[sec:031967]{以馬忤斯路上的驚喜 (華)} & 2026-04-12 & \href{http://www.church.com.hk/acms/content.asp?site=cdc\&op=show\&type=product\&code=031967\&layout=sermon}{\texttt{031967}} \\
約翰福音 21:1-11 & 吳天凱牧師 & \hyperref[sec:031968]{開往水深之處} & 2026-04-12 & \href{http://www.church.com.hk/acms/content.asp?site=cdc\&op=show\&type=product\&code=031968\&layout=sermon}{\texttt{031968}} \\
出埃及記 14:4 & 黃克勤牧師博士 & \hyperref[sec:031969]{聖經人聖經事:我在法老和他全軍得榮耀} & 2026-04-10 & \href{http://www.church.com.hk/acms/content.asp?site=cdc\&op=show\&type=product\&code=031969\&layout=sermon}{\texttt{031969}} \\
申命記 4:44-5:33 & 沈炳華傳道 & \hyperref[sec:031970]{上帝也講KPI?(1) 重述十誡} & 2026-04-11 & \href{http://www.church.com.hk/acms/content.asp?site=cdc\&op=show\&type=product\&code=031970\&layout=sermon}{\texttt{031970}} \\
約翰福音 17:21-23 & 唐郭麗婷牧師 & \hyperref[sec:031971]{與同輩同行的起點:合而為一} & 2026-04-12 & \href{http://www.church.com.hk/acms/content.asp?site=cdc\&op=show\&type=product\&code=031971\&layout=sermon}{\texttt{031971}} \\
出埃及記 19:1-6 & 丘放河牧師博士 & \hyperref[sec:031972]{信徒皆具雙職的身份} & 2026-04-12 & \href{http://www.church.com.hk/acms/content.asp?site=cdc\&op=show\&type=product\&code=031972\&layout=sermon}{\texttt{031972}} \\
約書亞記 24:20-24 & 丘放河牧師博士 & \hyperref[sec:031973]{在應許地活出「雙職」的使命} & 2026-04-12 & \href{http://www.church.com.hk/acms/content.asp?site=cdc\&op=show\&type=product\&code=031973\&layout=sermon}{\texttt{031973}} \\
詩篇 51:1-19 & 張健庭牧師博士 & \hyperref[sec:031974]{內裡的更新} & 2026-04-12 & \href{http://www.church.com.hk/acms/content.asp?site=cdc\&op=show\&type=product\&code=031974\&layout=sermon}{\texttt{031974}} \\
腓立比書 1:3-11 & 陳明正牧師 & \hyperref[sec:031975]{同心合意,興旺福音 (華)} & 2026-04-12 & \href{http://www.church.com.hk/acms/content.asp?site=cdc\&op=show\&type=product\&code=031975\&layout=sermon}{\texttt{031975}} \\
撒母耳記上 24:1-19 & 李思聰牧師 & \hyperref[sec:031976]{以善勝惡} & 2026-04-12 & \href{http://www.church.com.hk/acms/content.asp?site=cdc\&op=show\&type=product\&code=031976\&layout=sermon}{\texttt{031976}} \\
出埃及記 14:13 & 黃克勤牧師博士 & \hyperref[sec:031977]{聖經人聖經事:看耶和華施行的救恩} & 2026-04-11 & \href{http://www.church.com.hk/acms/content.asp?site=cdc\&op=show\&type=product\&code=031977\&layout=sermon}{\texttt{031977}} \\
民數記 9:1-23 & 鄭玉兒傳道 & \hyperref[sec:031978]{生命更新•新旅程} & 2026-04-11 & \href{http://www.church.com.hk/acms/content.asp?site=cdc\&op=show\&type=product\&code=031978\&layout=sermon}{\texttt{031978}} \\
民數記 9:1-23 & 董家驊牧師 & \hyperref[sec:031979]{走一條沒走過的新路!(華)} & 2026-04-12 & \href{http://www.church.com.hk/acms/content.asp?site=cdc\&op=show\&type=product\&code=031979\&layout=sermon}{\texttt{031979}} \\
以賽亞書 66:1-6 & 黃茂林傳道 & \hyperref[sec:031980]{蒙主恩待的敬拜者} & 2026-04-12 & \href{http://www.church.com.hk/acms/content.asp?site=cdc\&op=show\&type=product\&code=031980\&layout=sermon}{\texttt{031980}} \\
列王記上 1:1-53 & 袁惠鈞牧師博士 & \hyperref[sec:031981]{王位誰屬:人看智慧,神看應許} & 2026-04-12 & \href{http://www.church.com.hk/acms/content.asp?site=cdc\&op=show\&type=product\&code=031981\&layout=sermon}{\texttt{031981}} \\
列王記下 6:8-23 & 曹子光牧師 & \hyperref[sec:031982]{請開我雙眼} & 2026-04-12 & \href{http://www.church.com.hk/acms/content.asp?site=cdc\&op=show\&type=product\&code=031982\&layout=sermon}{\texttt{031982}} \\
使徒行傳 11:27-12:25 & 林鍵牧師 & \hyperref[sec:031983]{越獄風雲 (2)} & 2026-04-11 & \href{http://www.church.com.hk/acms/content.asp?site=cdc\&op=show\&type=product\&code=031983\&layout=sermon}{\texttt{031983}} \\
出埃及記 14:18 & 黃克勤牧師博士 & \hyperref[sec:031984]{聖經人聖經事:知道我是耶和華了} & 2026-04-12 & \href{http://www.church.com.hk/acms/content.asp?site=cdc\&op=show\&type=product\&code=031984\&layout=sermon}{\texttt{031984}} \\
尼希米記 4:1-11 & 顏沛恩牧師 & \hyperref[sec:031985]{爭戰到底•祈禱包底} & 2026-04-12 & \href{http://www.church.com.hk/acms/content.asp?site=cdc\&op=show\&type=product\&code=031985\&layout=sermon}{\texttt{031985}} \\
約翰福音 6:1-15 & 鄭偉誠牧師 & \hyperref[sec:031986]{你生命可以參與神蹟} & 2026-04-19 & \href{http://www.church.com.hk/acms/content.asp?site=cdc\&op=show\&type=product\&code=031986\&layout=sermon}{\texttt{031986}} \\
約翰福音 21:15-25 & 禤凱榮牧師 & \hyperref[sec:031987]{當失敗遇上復活主} & 2026-04-19 & \href{http://www.church.com.hk/acms/content.asp?site=cdc\&op=show\&type=product\&code=031987\&layout=sermon}{\texttt{031987}} \\
啟示錄 2:8-12 & 趙崇健牧師 & \hyperref[sec:031988]{士每拿:貧窮卻又富足的教會} & 2026-04-19 & \href{http://www.church.com.hk/acms/content.asp?site=cdc\&op=show\&type=product\&code=031988\&layout=sermon}{\texttt{031988}} \\
出埃及記 15:3 & 黃克勤牧師博士 & \hyperref[sec:031989]{聖經人聖經事:我的神,他名是耶和華} & 2026-04-13 & \href{http://www.church.com.hk/acms/content.asp?site=cdc\&op=show\&type=product\&code=031989\&layout=sermon}{\texttt{031989}} \\
歷代志上 29:10-19 & 莊明偉牧師 & \hyperref[sec:031990]{一切屬你樂意獻上} & 2025-09-21 & \href{http://www.church.com.hk/acms/content.asp?site=cdc\&op=show\&type=product\&code=031990\&layout=sermon}{\texttt{031990}} \\
馬太福音 7:1-12 & 王正傑牧師 & \hyperref[sec:031991]{最好的東西} & 2026-04-19 & \href{http://www.church.com.hk/acms/content.asp?site=cdc\&op=show\&type=product\&code=031991\&layout=sermon}{\texttt{031991}} \\
馬可福音 8:29-9:1 & 鄧瑞強牧師博士 & \hyperref[sec:031992]{只有一位受苦的上帝能救我們} & 2026-04-19 & \href{http://www.church.com.hk/acms/content.asp?site=cdc\&op=show\&type=product\&code=031992\&layout=sermon}{\texttt{031992}} \\
彼得後書 1:5-8 & 吳磐德牧師 & \hyperref[sec:031993]{建立健康的使命教會 (12) 勤成長的靈命} & 2026-04-19 & \href{http://www.church.com.hk/acms/content.asp?site=cdc\&op=show\&type=product\&code=031993\&layout=sermon}{\texttt{031993}} \\
使徒行傳 10:9-16 & 何俊文牧師 & \hyperref[sec:031994]{求主藉異象激動我} & 2026-04-19 & \href{http://www.church.com.hk/acms/content.asp?site=cdc\&op=show\&type=product\&code=031994\&layout=sermon}{\texttt{031994}} \\
馬太福音 7:13-27 & 石偉漢牧師 & \hyperref[sec:031995]{二揀一} & 2026-04-19 & \href{http://www.church.com.hk/acms/content.asp?site=cdc\&op=show\&type=product\&code=031995\&layout=sermon}{\texttt{031995}} \\
出埃及記 14:2 & 黃克勤牧師博士 & \hyperref[sec:031996]{聖經人聖經事:吩咐以色列人回轉} & 2026-04-14 & \href{http://www.church.com.hk/acms/content.asp?site=cdc\&op=show\&type=product\&code=031996\&layout=sermon}{\texttt{031996}} \\
歌羅西書 2:1-8 & 羅錫堅牧師 & \hyperref[sec:031997]{牧人的掙扎} & 2026-04-19 & \href{http://www.church.com.hk/acms/content.asp?site=cdc\&op=show\&type=product\&code=031997\&layout=sermon}{\texttt{031997}} \\
列王記上 19:9-16 & 文麗芳傳道 & \hyperref[sec:031998]{只剩下我一個人} & 2026-04-19 & \href{http://www.church.com.hk/acms/content.asp?site=cdc\&op=show\&type=product\&code=031998\&layout=sermon}{\texttt{031998}} \\
箴言 9:10 & 蕭壽華牧師 & \hyperref[sec:031999]{認識至聖者便是聰明} & 2025-09-28 & \href{http://www.church.com.hk/acms/content.asp?site=cdc\&op=show\&type=product\&code=031999\&layout=sermon}{\texttt{031999}} \\
詩篇 139:1-24 & 張健庭牧師博士 & \hyperref[sec:032015]{認識神無所不在的生命轉化} & 2026-04-19 & \href{http://www.church.com.hk/acms/content.asp?site=cdc\&op=show\&type=product\&code=032015\&layout=sermon}{\texttt{032015}} \\
使徒行傳 27:1-28:1 & 林鍵牧師 & \hyperref[sec:032016]{征服世界的復活救主} & 2026-04-18 & \href{http://www.church.com.hk/acms/content.asp?site=cdc\&op=show\&type=product\&code=032016\&layout=sermon}{\texttt{032016}} \\
出埃及記 16:8 & 黃克勤牧師博士 & \hyperref[sec:032017]{聖經人聖經事:你們的怨言是向耶和華發的} & 2026-04-17 & \href{http://www.church.com.hk/acms/content.asp?site=cdc\&op=show\&type=product\&code=032017\&layout=sermon}{\texttt{032017}} \\
哥林多前書 1:26-29 & 林可豐傳道 & \hyperref[sec:032018]{神國的精英主義} & 2026-04-19 & \href{http://www.church.com.hk/acms/content.asp?site=cdc\&op=show\&type=product\&code=032018\&layout=sermon}{\texttt{032018}} \\
民數記 10:1-36 & 陳傳華牧師博士 & \hyperref[sec:032019]{起行與安營:聆聽神的信號,跟從祂的節奏} & 2026-04-19 & \href{http://www.church.com.hk/acms/content.asp?site=cdc\&op=show\&type=product\&code=032019\&layout=sermon}{\texttt{032019}} \\
約翰福音 21:20-23 & 黃健羒牧師 & \hyperref[sec:032020]{你跟從我吧} & 2026-04-19 & \href{http://www.church.com.hk/acms/content.asp?site=cdc\&op=show\&type=product\&code=032020\&layout=sermon}{\texttt{032020}} \\
哥林多前書 4:1-13 & 黎炳鍵弟兄 & \hyperref[sec:032021]{使徒的真面目} & 2026-04-19 & \href{http://www.church.com.hk/acms/content.asp?site=cdc\&op=show\&type=product\&code=032021\&layout=sermon}{\texttt{032021}} \\
出埃及記 16:35 & 黃克勤牧師博士 & \hyperref[sec:032022]{聖經人聖經事:以色人吃嗎哪四十年} & 2026-04-18 & \href{http://www.church.com.hk/acms/content.asp?site=cdc\&op=show\&type=product\&code=032022\&layout=sermon}{\texttt{032022}} \\
出埃及記 2:23-25 & 馮永和博士 & \hyperref[sec:032023]{改變現狀的禱告——嘆息哀求} & 2026-04-18 & \href{http://www.church.com.hk/acms/content.asp?site=cdc\&op=show\&type=product\&code=032023\&layout=sermon}{\texttt{032023}} \\
撒母耳記上 15:17-23 & 曾敬賢傳道 & \hyperref[sec:032024]{始終一僕人} & 2026-04-19 & \href{http://www.church.com.hk/acms/content.asp?site=cdc\&op=show\&type=product\&code=032024\&layout=sermon}{\texttt{032024}} \\
路加福音 18:15-30 & 梁成裕牧師 & \hyperref[sec:032025]{神國的真諦:單純的信靠和全心的跟隨} & 2026-04-19 & \href{http://www.church.com.hk/acms/content.asp?site=cdc\&op=show\&type=product\&code=032025\&layout=sermon}{\texttt{032025}} \\
創世記 25:19-34 & 梁智程傳道 & \hyperref[sec:032026]{長子或紅湯} & 2026-04-18 & \href{http://www.church.com.hk/acms/content.asp?site=cdc\&op=show\&type=product\&code=032026\&layout=sermon}{\texttt{032026}} \\
使徒行傳 1:8 & 李文康傳道 & \hyperref[sec:032027]{與你同行,直到地極} & 2026-04-19 & \href{http://www.church.com.hk/acms/content.asp?site=cdc\&op=show\&type=product\&code=032027\&layout=sermon}{\texttt{032027}} \\
歌羅西書 1:1-4:18 & 劉少康牧師 & \hyperref[sec:032028]{歌羅西書全書概論:超越的基督} & 2026-04-15 & \href{http://www.church.com.hk/acms/content.asp?site=cdc\&op=show\&type=product\&code=032028\&layout=sermon}{\texttt{032028}} \\
羅馬書 15:1-6 & 陳雅莉傳道 & \hyperref[sec:032029]{不求自己的喜悅} & 2026-04-26 & \href{http://www.church.com.hk/acms/content.asp?site=cdc\&op=show\&type=product\&code=032029\&layout=sermon}{\texttt{032029}} \\
約翰福音 4:23 & 王亞辰牧師 & \hyperref[sec:032030]{父神在尋找:真正的敬拜者 (華)} & 2026-04-26 & \href{http://www.church.com.hk/acms/content.asp?site=cdc\&op=show\&type=product\&code=032030\&layout=sermon}{\texttt{032030}} \\
馬太福音 18:1-20 & 駱瑞麟傳道 & \hyperref[sec:032031]{一個都不能少} & 2026-04-26 & \href{http://www.church.com.hk/acms/content.asp?site=cdc\&op=show\&type=product\&code=032031\&layout=sermon}{\texttt{032031}} \\
出埃及記 17:6 & 黃克勤牧師博士 & \hyperref[sec:032032]{聖經人聖經事:我必站在你面前} & 2026-04-19 & \href{http://www.church.com.hk/acms/content.asp?site=cdc\&op=show\&type=product\&code=032032\&layout=sermon}{\texttt{032032}} \\
以西結書 47:1-12 & 郭張玉清姊妹 & \hyperref[sec:032033]{活水江河} & 2026-04-26 & \href{http://www.church.com.hk/acms/content.asp?site=cdc\&op=show\&type=product\&code=032033\&layout=sermon}{\texttt{032033}} \\
約翰福音 10:1-10 & 林津牧師 & \hyperref[sec:032034]{得更豐盛的「生命」} & 2026-04-26 & \href{http://www.church.com.hk/acms/content.asp?site=cdc\&op=show\&type=product\&code=032034\&layout=sermon}{\texttt{032034}} \\
詩篇 136:1-26 & 張健庭牧師博士 & \hyperref[sec:032035]{因祂的慈愛永遠長存} & 2026-04-26 & \href{http://www.church.com.hk/acms/content.asp?site=cdc\&op=show\&type=product\&code=032035\&layout=sermon}{\texttt{032035}} \\
以弗所書 2:1-10 & 歐陽萬廷弟兄 & \hyperref[sec:032036]{天父之恩,聖徒之行} & 2026-04-26 & \href{http://www.church.com.hk/acms/content.asp?site=cdc\&op=show\&type=product\&code=032036\&layout=sermon}{\texttt{032036}} \\
帖撒羅尼迦前書 1:1-5:28 & 劉少康牧師 & \hyperref[sec:032037]{帖撒羅尼迦書 (1) 主快再來} & 2026-04-15 & \href{http://www.church.com.hk/acms/content.asp?site=cdc\&op=show\&type=product\&code=032037\&layout=sermon}{\texttt{032037}} \\
出埃及記 17:5 & 黃克勤牧師博士 & \hyperref[sec:032038]{聖經人聖經事:耶和華尼西} & 2026-04-20 & \href{http://www.church.com.hk/acms/content.asp?site=cdc\&op=show\&type=product\&code=032038\&layout=sermon}{\texttt{032038}} \\
馬太福音 17:1-9 & 謝鵬輝傳道 & \hyperref[sec:032039]{起來!下山去吧!} & 2026-04-26 & \href{http://www.church.com.hk/acms/content.asp?site=cdc\&op=show\&type=product\&code=032039\&layout=sermon}{\texttt{032039}} \\
路加福音 18:31-43 & 梁成裕牧師 & \hyperref[sec:032040]{看見和跟隨耶穌的道路} & 2026-04-26 & \href{http://www.church.com.hk/acms/content.asp?site=cdc\&op=show\&type=product\&code=032040\&layout=sermon}{\texttt{032040}} \\
出埃及記 17:8-16 & 梁永善牧師 & \hyperref[sec:032041]{托手zoung} & 2026-04-26 & \href{http://www.church.com.hk/acms/content.asp?site=cdc\&op=show\&type=product\&code=032041\&layout=sermon}{\texttt{032041}} \\
羅馬書 8:18-30 & 張路加牧師 & \hyperref[sec:032042]{從廢墟到盼望 (華)} & 2026-04-26 & \href{http://www.church.com.hk/acms/content.asp?site=cdc\&op=show\&type=product\&code=032042\&layout=sermon}{\texttt{032042}} \\
約伯記 19:23-29 & 黃馮潔霞師母 & \hyperref[sec:032043]{我的救贖主活著} & 2026-04-26 & \href{http://www.church.com.hk/acms/content.asp?site=cdc\&op=show\&type=product\&code=032043\&layout=sermon}{\texttt{032043}} \\
出埃及記 20:2 & 黃克勤牧師博士 & \hyperref[sec:032044]{聖經人聖經事:我曾領你出來} & 2026-04-21 & \href{http://www.church.com.hk/acms/content.asp?site=cdc\&op=show\&type=product\&code=032044\&layout=sermon}{\texttt{032044}} \\
申命記 6:1-25 & 符氣松牧師 & \hyperref[sec:032045]{上帝也講KPI?(2) 點解要做?} & 2026-04-25 & \href{http://www.church.com.hk/acms/content.asp?site=cdc\&op=show\&type=product\&code=032045\&layout=sermon}{\texttt{032045}} \\
馬太福音 18:21-22 & 唐志濠牧師 & \hyperref[sec:032046]{得饒人處且饒人} & 2026-04-26 & \href{http://www.church.com.hk/acms/content.asp?site=cdc\&op=show\&type=product\&code=032046\&layout=sermon}{\texttt{032046}} \\
約翰福音 4:43-54 & 譚偉明宣教士 & \hyperref[sec:032047]{你是耶穌的粉絲嗎?} & 2026-04-26 & \href{http://www.church.com.hk/acms/content.asp?site=cdc\&op=show\&type=product\&code=032047\&layout=sermon}{\texttt{032047}} \\
約書亞記 24:14-15 & 招世超牧師 & \hyperref[sec:032048]{家庭文化2.0} & 2026-04-26 & \href{http://www.church.com.hk/acms/content.asp?site=cdc\&op=show\&type=product\&code=032048\&layout=sermon}{\texttt{032048}} \\
約翰福音 21:15-17 & 秦潤祥牧師 & \hyperref[sec:032049]{陀飛輪:從睇化,到睇透,再到激活} & 2026-04-26 & \href{http://www.church.com.hk/acms/content.asp?site=cdc\&op=show\&type=product\&code=032049\&layout=sermon}{\texttt{032049}} \\
出埃及記 32:13 & 黃克勤牧師博士 & \hyperref[sec:032050]{聖經人聖經事:你曾指著自己起誓} & 2026-04-22 & \href{http://www.church.com.hk/acms/content.asp?site=cdc\&op=show\&type=product\&code=032050\&layout=sermon}{\texttt{032050}} \\
約書亞記 21:1-45 & 方兆雄牧師 & \hyperref[sec:032051]{應許} & 2026-04-26 & \href{http://www.church.com.hk/acms/content.asp?site=cdc\&op=show\&type=product\&code=032051\&layout=sermon}{\texttt{032051}} \\
馬太福音 5:3-10 & 朱卓慈牧師 & \hyperref[sec:032052]{真正有福的人} & 2026-04-25 & \href{http://www.church.com.hk/acms/content.asp?site=cdc\&op=show\&type=product\&code=032052\&layout=sermon}{\texttt{032052}} \\
約翰福音 13:1 & 余德淳博士 & \hyperref[sec:032053]{愛到底的意思} & 2026-04-26 & \href{http://www.church.com.hk/acms/content.asp?site=cdc\&op=show\&type=product\&code=032053\&layout=sermon}{\texttt{032053}} \\
路加福音 15:1-32 & 俞大鳴牧師 & \hyperref[sec:032054]{理應歡喜} & 2026-04-26 & \href{http://www.church.com.hk/acms/content.asp?site=cdc\&op=show\&type=product\&code=032054\&layout=sermon}{\texttt{032054}} \\
歌羅西書 2:16-23 & 黃紹權牧師 & \hyperref[sec:032055]{教會尋求社會認同的危機} & 2026-04-26 & \href{http://www.church.com.hk/acms/content.asp?site=cdc\&op=show\&type=product\&code=032055\&layout=sermon}{\texttt{032055}} \\
出埃及記 23:12 & 黃克勤牧師博士 & \hyperref[sec:032056]{聖經人聖經事:第七日要安息} & 2026-04-23 & \href{http://www.church.com.hk/acms/content.asp?site=cdc\&op=show\&type=product\&code=032056\&layout=sermon}{\texttt{032056}} \\
創世記 5:21-29 & 黃若娟傳道 & \hyperref[sec:032057]{與神同行} & 2026-04-26 & \href{http://www.church.com.hk/acms/content.asp?site=cdc\&op=show\&type=product\&code=032057\&layout=sermon}{\texttt{032057}} \\
腓立比書 2:1-5 & 林耀楠牧師 & \hyperref[sec:032058]{建立健康的使命教會 (13) 同心事奉的秘訣} & 2026-04-26 & \href{http://www.church.com.hk/acms/content.asp?site=cdc\&op=show\&type=product\&code=032058\&layout=sermon}{\texttt{032058}} \\
彼得前書 2:2-10 & 李均熊牧師博士 & \hyperref[sec:032071]{活石} & 2026-05-03 & \href{http://www.church.com.hk/acms/content.asp?site=cdc\&op=show\&type=product\&code=032071\&layout=sermon}{\texttt{032071}} \\
約翰福音 4:1-42 & 石偉漢牧師 & \hyperref[sec:032072]{天涯何處無芳草} & 2026-05-03 & \href{http://www.church.com.hk/acms/content.asp?site=cdc\&op=show\&type=product\&code=032072\&layout=sermon}{\texttt{032072}} \\
彼得前書 4:7-11 & 江山偉牧師 & \hyperref[sec:032073]{建立健康的使命教會 (14) 作恩賜的管家} & 2026-05-03 & \href{http://www.church.com.hk/acms/content.asp?site=cdc\&op=show\&type=product\&code=032073\&layout=sermon}{\texttt{032073}} \\
帖撒羅尼迦後書 1:1-3:18 & 劉少康牧師 & \hyperref[sec:032074]{帖撒羅尼迦書 (2) 主快再來} & 2026-04-28 & \href{http://www.church.com.hk/acms/content.asp?site=cdc\&op=show\&type=product\&code=032074\&layout=sermon}{\texttt{032074}} \\
出埃及記 25:10 & 黃克勤牧師博士 & \hyperref[sec:032075]{聖經人聖經事:用皂莢木做約櫃} & 2026-04-25 & \href{http://www.church.com.hk/acms/content.asp?site=cdc\&op=show\&type=product\&code=032075\&layout=sermon}{\texttt{032075}} \\
腓立比書 3:12-16 & 莫潤輝牧師 & \hyperref[sec:032076]{努力!繼續進擊} & 2026-05-03 & \href{http://www.church.com.hk/acms/content.asp?site=cdc\&op=show\&type=product\&code=032076\&layout=sermon}{\texttt{032076}} \\
創世記 35:1-15 & 羅泰然牧師 & \hyperref[sec:032077]{把偶像埋在橡樹下} & 2026-05-03 & \href{http://www.church.com.hk/acms/content.asp?site=cdc\&op=show\&type=product\&code=032077\&layout=sermon}{\texttt{032077}} \\
馬太福音 28:1-7 & 趙崇健牧師 & \hyperref[sec:032078]{When Hope Returns (英/粵)} & 2026-05-03 & \href{http://www.church.com.hk/acms/content.asp?site=cdc\&op=show\&type=product\&code=032078\&layout=sermon}{\texttt{032078}} \\
哥林多前書 12:27 & 劉彤牧師 & \hyperref[sec:032079]{同心合意,建造教會 (華/英)} & 2026-05-03 & \href{http://www.church.com.hk/acms/content.asp?site=cdc\&op=show\&type=product\&code=032079\&layout=sermon}{\texttt{032079}} \\
以弗所書 5:21-33 & 陳志輝牧師 & \hyperref[sec:032080]{共同的使命:婚姻} & 2026-05-03 & \href{http://www.church.com.hk/acms/content.asp?site=cdc\&op=show\&type=product\&code=032080\&layout=sermon}{\texttt{032080}} \\
以賽亞書 49:1-6 & 李應新牧師 & \hyperref[sec:032081]{外邦人的光} & 2026-05-03 & \href{http://www.church.com.hk/acms/content.asp?site=cdc\&op=show\&type=product\&code=032081\&layout=sermon}{\texttt{032081}} \\
出埃及記 32:4 & 黃克勤牧師博士 & \hyperref[sec:032082]{聖經人聖經事:阿倫的妥協} & 2026-04-30 & \href{http://www.church.com.hk/acms/content.asp?site=cdc\&op=show\&type=product\&code=032082\&layout=sermon}{\texttt{032082}} \\
約翰福音 7:1-13 & 林鍵牧師 & \hyperref[sec:032083]{榮耀君王:兄弟情深} & 2026-05-02 & \href{http://www.church.com.hk/acms/content.asp?site=cdc\&op=show\&type=product\&code=032083\&layout=sermon}{\texttt{032083}} \\
腓立比書 4:4 & 陳淳牧師 & \hyperref[sec:032084]{喜樂,喜樂:生命的顛覆 (粵/華)} & 2026-05-03 & \href{http://www.church.com.hk/acms/content.asp?site=cdc\&op=show\&type=product\&code=032084\&layout=sermon}{\texttt{032084}} \\
士師記 1:1-4 & 梁永善牧師 & \hyperref[sec:032085]{士師記析讀 (4) 互相扶助 迎挑戰} & 2026-05-06 & \href{http://www.church.com.hk/acms/content.asp?site=cdc\&op=show\&type=product\&code=032085\&layout=sermon}{\texttt{032085}} \\
馬太福音 6:25-33 & 張健庭牧師博士 & \hyperref[sec:032086]{先求祂的國} & 2026-05-03 & \href{http://www.church.com.hk/acms/content.asp?site=cdc\&op=show\&type=product\&code=032086\&layout=sermon}{\texttt{032086}} \\
民數記 20:12 & 黃克勤牧師博士 & \hyperref[sec:032087]{聖經人聖經事:摩西發脾氣出事} & 2026-05-02 & \href{http://www.church.com.hk/acms/content.asp?site=cdc\&op=show\&type=product\&code=032087\&layout=sermon}{\texttt{032087}} \\
以弗所書 5:15-20 & 杜崇基牧師 & \hyperref[sec:032088]{要明白主的旨意:七要三不要} & 2026-05-03 & \href{http://www.church.com.hk/acms/content.asp?site=cdc\&op=show\&type=product\&code=032088\&layout=sermon}{\texttt{032088}} \\
羅馬書 1:1-15 & 黃志剛牧師 & \hyperref[sec:032089]{持定福音使命的人生} & 2025-10-04 & \href{http://www.church.com.hk/acms/content.asp?site=cdc\&op=show\&type=product\&code=032089\&layout=sermon}{\texttt{032089}} \\
約翰福音 4:1-18 & 駱柱成牧師 & \hyperref[sec:032090]{主耶穌傳福音的智慧} & 2026-05-03 & \href{http://www.church.com.hk/acms/content.asp?site=cdc\&op=show\&type=product\&code=032090\&layout=sermon}{\texttt{032090}} \\
馬可福音 9:30-37 & 羅杰才牧師 & \hyperref[sec:032091]{耶穌如何回應門徒的爭論} & 2026-04-26 & \href{http://www.church.com.hk/acms/content.asp?site=cdc\&op=show\&type=product\&code=032091\&layout=sermon}{\texttt{032091}} \\
路加福音 16:1-13 & 俞大鳴牧師 & \hyperref[sec:032092]{門徒的錢財觀} & 2026-05-03 & \href{http://www.church.com.hk/acms/content.asp?site=cdc\&op=show\&type=product\&code=032092\&layout=sermon}{\texttt{032092}} \\
約翰福音 16:16-24 & 謝喬遠弟兄 & \hyperref[sec:032093]{憂愁變喜樂} & 2026-05-03 & \href{http://www.church.com.hk/acms/content.asp?site=cdc\&op=show\&type=product\&code=032093\&layout=sermon}{\texttt{032093}} \\
撒母耳記上 1:1-2:36 & 梁廷益牧師 & \hyperref[sec:032094]{一位母親的奉獻} & 2026-05-03 & \href{http://www.church.com.hk/acms/content.asp?site=cdc\&op=show\&type=product\&code=032094\&layout=sermon}{\texttt{032094}} \\
約書亞記 2:11 & 黃克勤牧師博士 & \hyperref[sec:032095]{勇敢機智信靠神的妓女喇合} & 2026-05-07 & \href{http://www.church.com.hk/acms/content.asp?site=cdc\&op=show\&type=product\&code=032095\&layout=sermon}{\texttt{032095}} \\
申命記 7:1-26 & 李觀南傳道 & \hyperref[sec:032096]{上帝也講KPI?(3) 超出預計的愛} & 2026-05-02 & \href{http://www.church.com.hk/acms/content.asp?site=cdc\&op=show\&type=product\&code=032096\&layout=sermon}{\texttt{032096}} \\
士師記 3:12-30 & 謝又生牧師 & \hyperref[sec:032097]{士師記系列 (2) 委身那意想不到的神} & 2026-05-03 & \href{http://www.church.com.hk/acms/content.asp?site=cdc\&op=show\&type=product\&code=032097\&layout=sermon}{\texttt{032097}} \\
創世記 12:1-50:26 & 陳崇基牧師 & \hyperref[sec:032098]{碎片} & 2026-05-03 & \href{http://www.church.com.hk/acms/content.asp?site=cdc\&op=show\&type=product\&code=032098\&layout=sermon}{\texttt{032098}} \\
詩篇 36:1-9 & 張磊牧師 & \hyperref[sec:032099]{人的邪惡,神的美善 (華)} & 2026-05-03 & \href{http://www.church.com.hk/acms/content.asp?site=cdc\&op=show\&type=product\&code=032099\&layout=sermon}{\texttt{032099}} \\
撒母耳記上 25:1-44 & 李思聰牧師 & \hyperref[sec:032100]{作個好男人,作個好女人——大衛與亞比該 (1)} & 2026-05-03 & \href{http://www.church.com.hk/acms/content.asp?site=cdc\&op=show\&type=product\&code=032100\&layout=sermon}{\texttt{032100}} \\
約書亞記 5:2-12 & 黃克勤牧師博士 & \hyperref[sec:032101]{聖經人聖經事:入門三件事} & 2026-05-08 & \href{http://www.church.com.hk/acms/content.asp?site=cdc\&op=show\&type=product\&code=032101\&layout=sermon}{\texttt{032101}} \\
提摩太前書 1:1-6:21 & 劉少康牧師 & \hyperref[sec:032102]{牧者手冊:提摩太前後書 (1)} & 2026-05-06 & \href{http://www.church.com.hk/acms/content.asp?site=cdc\&op=show\&type=product\&code=032102\&layout=sermon}{\texttt{032102}} \\
哥林多前書 9:1-6 & 陳嘉信牧師 & \hyperref[sec:032103]{意想不到的自由} & 2026-05-03 & \href{http://www.church.com.hk/acms/content.asp?site=cdc\&op=show\&type=product\&code=032103\&layout=sermon}{\texttt{032103}} \\
創世記 25:24-28 & 譚綺敏傳道 & \hyperref[sec:032104]{利百加:愛與掌控的母親} & 2026-05-10 & \href{http://www.church.com.hk/acms/content.asp?site=cdc\&op=show\&type=product\&code=032104\&layout=sermon}{\texttt{032104}} \\
約翰福音 7:11-36 & 林鍵牧師 & \hyperref[sec:032105]{榮耀君王:終極問答} & 2026-05-09 & \href{http://www.church.com.hk/acms/content.asp?site=cdc\&op=show\&type=product\&code=032105\&layout=sermon}{\texttt{032105}} \\
以賽亞書 49:1-6 & 周弘毅牧師 & \hyperref[sec:032107]{未竟全功?} & 2026-05-09 & \href{http://www.church.com.hk/acms/content.asp?site=cdc\&op=show\&type=product\&code=032107\&layout=sermon}{\texttt{032107}} \\
以弗所書 2:1-10 & 禤凱榮牧師 & \hyperref[sec:032108]{超越七大奇蹟的奇蹟} & 2026-05-10 & \href{http://www.church.com.hk/acms/content.asp?site=cdc\&op=show\&type=product\&code=032108\&layout=sermon}{\texttt{032108}} \\
彼得前書 3:13-22 & 林豪恩先生 & \hyperref[sec:032109]{耶穌不在的時候} & 2026-05-10 & \href{http://www.church.com.hk/acms/content.asp?site=cdc\&op=show\&type=product\&code=032109\&layout=sermon}{\texttt{032109}} \\
約書亞記 5:14 & 黃克勤牧師博士 & \hyperref[sec:032110]{聖經人聖經事:我來是要作元帥} & 2026-05-09 & \href{http://www.church.com.hk/acms/content.asp?site=cdc\&op=show\&type=product\&code=032110\&layout=sermon}{\texttt{032110}} \\
箴言 22:6 & 林羅麗斯牧師 & \hyperref[sec:032111]{為母則剛} & 2026-05-10 & \href{http://www.church.com.hk/acms/content.asp?site=cdc\&op=show\&type=product\&code=032111\&layout=sermon}{\texttt{032111}} \\
使徒行傳 2:42-47 & 鄺偉志牧師 & \hyperref[sec:032112]{教會增長的迷思 The Growing Church} & 2026-05-10 & \href{http://www.church.com.hk/acms/content.asp?site=cdc\&op=show\&type=product\&code=032112\&layout=sermon}{\texttt{032112}} \\
腓立比書 4:4-7 & 張健庭牧師博士 & \hyperref[sec:032113]{在主裡得滿足} & 2026-05-10 & \href{http://www.church.com.hk/acms/content.asp?site=cdc\&op=show\&type=product\&code=032113\&layout=sermon}{\texttt{032113}} \\
以弗所書 6:1-4 & 王正傑牧師 & \hyperref[sec:032114]{當孝敬父母} & 2026-05-10 & \href{http://www.church.com.hk/acms/content.asp?site=cdc\&op=show\&type=product\&code=032114\&layout=sermon}{\texttt{032114}} \\
歌羅西書 3:5-11 & 黃紹權牧師 & \hyperref[sec:032115]{我們怎樣彰顯神的榮耀?} & 2026-05-10 & \href{http://www.church.com.hk/acms/content.asp?site=cdc\&op=show\&type=product\&code=032115\&layout=sermon}{\texttt{032115}} \\
箴言 22:6 & 劉彤牧師 & \hyperref[sec:032116]{如何把兒女帶到神面前} & 2026-05-10 & \href{http://www.church.com.hk/acms/content.asp?site=cdc\&op=show\&type=product\&code=032116\&layout=sermon}{\texttt{032116}} \\
以弗所書 5:21 & 招世超牧師 & \hyperref[sec:032117]{與子偕老的見證} & 2026-05-10 & \href{http://www.church.com.hk/acms/content.asp?site=cdc\&op=show\&type=product\&code=032117\&layout=sermon}{\texttt{032117}} \\
撒母耳記上 25:1-44 & 李思聰牧師 & \hyperref[sec:032118]{作個好男人,作個好女人—大衛與亞比該 (下)} & 2026-05-10 & \href{http://www.church.com.hk/acms/content.asp?site=cdc\&op=show\&type=product\&code=032118\&layout=sermon}{\texttt{032118}} \\
箴言 31:1-31 & 袁惠鈞牧師博士 & \hyperref[sec:032119]{智慧與恩慈:生命的織布者} & 2026-05-10 & \href{http://www.church.com.hk/acms/content.asp?site=cdc\&op=show\&type=product\&code=032119\&layout=sermon}{\texttt{032119}} \\
瑪拉基書 1:1-5 & 石偉漢牧師 & \hyperref[sec:032120]{誰曾愛我?} & 2026-05-10 & \href{http://www.church.com.hk/acms/content.asp?site=cdc\&op=show\&type=product\&code=032120\&layout=sermon}{\texttt{032120}} \\
約拿書 1:1-4:11 & 陳雅莉傳道 & \hyperref[sec:032121]{出走的先知} & 2026-05-10 & \href{http://www.church.com.hk/acms/content.asp?site=cdc\&op=show\&type=product\&code=032121\&layout=sermon}{\texttt{032121}} \\
歌羅西書 3:12-21 & 林瑞貞傳道 & \hyperref[sec:032122]{最好的禮物} & 2026-05-10 & \href{http://www.church.com.hk/acms/content.asp?site=cdc\&op=show\&type=product\&code=032122\&layout=sermon}{\texttt{032122}} \\
瑪拉基書 3:7-18 & 梁永善牧師 & \hyperref[sec:032123]{失感情——十一奉獻的探討} & 2026-05-10 & \href{http://www.church.com.hk/acms/content.asp?site=cdc\&op=show\&type=product\&code=032123\&layout=sermon}{\texttt{032123}} \\
約書亞記 9:27 & 黃克勤牧師博士 & \hyperref[sec:032124]{聖經人聖經事:自己吃虧也不改誓} & 2026-05-14 & \href{http://www.church.com.hk/acms/content.asp?site=cdc\&op=show\&type=product\&code=032124\&layout=sermon}{\texttt{032124}} \\
路加福音 15:11-32 & 冼嘉琪傳道 & \hyperref[sec:032125]{當層出不窮的愛,遇上深不見底的愛} & 2026-05-10 & \href{http://www.church.com.hk/acms/content.asp?site=cdc\&op=show\&type=product\&code=032125\&layout=sermon}{\texttt{032125}} \\
以弗所書 1:15-23 & 姚臣泰長老 & \hyperref[sec:032126]{保羅為教會禱告 (華/粵)} & 2026-05-10 & \href{http://www.church.com.hk/acms/content.asp?site=cdc\&op=show\&type=product\&code=032126\&layout=sermon}{\texttt{032126}} \\
以賽亞書 49:14-18 & 馮潔儀傳道 & \hyperref[sec:032127]{必不忘記} & 2026-05-10 & \href{http://www.church.com.hk/acms/content.asp?site=cdc\&op=show\&type=product\&code=032127\&layout=sermon}{\texttt{032127}} \\
列王記下 4:8-37 & 溫艾衛傳道 & \hyperref[sec:032128]{婦人的愛與信心} & 2026-05-10 & \href{http://www.church.com.hk/acms/content.asp?site=cdc\&op=show\&type=product\&code=032128\&layout=sermon}{\texttt{032128}} \\
路得記 1:6-22 & 毛譚美虹傳道 & \hyperref[sec:032129]{在神的季節中選擇} & 2026-05-09 & \href{http://www.church.com.hk/acms/content.asp?site=cdc\&op=show\&type=product\&code=032129\&layout=sermon}{\texttt{032129}} \\
約翰福音 19:25-27 & 李頌華傳道 & \hyperref[sec:032130]{看,你的母親} & 2026-05-10 & \href{http://www.church.com.hk/acms/content.asp?site=cdc\&op=show\&type=product\&code=032130\&layout=sermon}{\texttt{032130}} \\
撒母耳記上 1:1-28 & 梁永善牧師 & \hyperref[sec:032131]{偉人母親的質素} & 2026-05-10 & \href{http://www.church.com.hk/acms/content.asp?site=cdc\&op=show\&type=product\&code=032131\&layout=sermon}{\texttt{032131}} \\
路加福音 1:46-55 & 梁錫強醫生 & \hyperref[sec:032132]{生命的逆轉} & 2026-05-17 & \href{http://www.church.com.hk/acms/content.asp?site=cdc\&op=show\&type=product\&code=032132\&layout=sermon}{\texttt{032132}} \\
歌羅西書 3:12-17 & 黃紹權牧師 & \hyperref[sec:032133]{信徒,你們在攪甚麼?} & 2026-05-17 & \href{http://www.church.com.hk/acms/content.asp?site=cdc\&op=show\&type=product\&code=032133\&layout=sermon}{\texttt{032133}} \\
約書亞記 11:23 & 黃克勤牧師博士 & \hyperref[sec:032134]{聖經人聖經事:靠主奪了那全地} & 2026-05-16 & \href{http://www.church.com.hk/acms/content.asp?site=cdc\&op=show\&type=product\&code=032134\&layout=sermon}{\texttt{032134}} \\
使徒行傳 1:6-14 & 趙承恩宣教師 & \hyperref[sec:032135]{天上天下 唯主獨尊} & 2026-05-17 & \href{http://www.church.com.hk/acms/content.asp?site=cdc\&op=show\&type=product\&code=032135\&layout=sermon}{\texttt{032135}} \\
路加福音 7:11-17 & 鄧瑞強牧師博士 & \hyperref[sec:032136]{夕陽西下 斷腸人在天涯} & 2026-05-17 & \href{http://www.church.com.hk/acms/content.asp?site=cdc\&op=show\&type=product\&code=032136\&layout=sermon}{\texttt{032136}} \\
瑪拉基書 1:6-14 & 石偉漢牧師 & \hyperref[sec:032137]{「A貨」的敬拜} & 2026-05-17 & \href{http://www.church.com.hk/acms/content.asp?site=cdc\&op=show\&type=product\&code=032137\&layout=sermon}{\texttt{032137}} \\
提摩太後書 1:1-4:22 & 劉少康牧師 & \hyperref[sec:032138]{牧者手冊:提摩太前後書 (2)} & 2026-05-13 & \href{http://www.church.com.hk/acms/content.asp?site=cdc\&op=show\&type=product\&code=032138\&layout=sermon}{\texttt{032138}} \\
約翰福音 13:3-10 & 蔡允瑤傳道 & \hyperref[sec:032139]{愛的醒覺} & 2026-05-17 & \href{http://www.church.com.hk/acms/content.asp?site=cdc\&op=show\&type=product\&code=032139\&layout=sermon}{\texttt{032139}} \\
歌羅西書 3:1-4 & 羅錫堅牧師 & \hyperref[sec:032140]{基督徒的現在與未來} & 2026-05-17 & \href{http://www.church.com.hk/acms/content.asp?site=cdc\&op=show\&type=product\&code=032140\&layout=sermon}{\texttt{032140}} \\
詩篇 127:1-5 & 李美儀傳道 & \hyperref[sec:032141]{建造的人不枉然勞力} & 2026-05-17 & \href{http://www.church.com.hk/acms/content.asp?site=cdc\&op=show\&type=product\&code=032141\&layout=sermon}{\texttt{032141}} \\
約書亞記 4:8 & 黃克勤牧師博士 & \hyperref[sec:032142]{聖經人聖經事:迦勤 專心跟從耶和華} & 2026-05-19 & \href{http://www.church.com.hk/acms/content.asp?site=cdc\&op=show\&type=product\&code=032142\&layout=sermon}{\texttt{032142}} \\
羅馬書 1:16-32 & 蕭壽華牧師 & \hyperref[sec:032143]{以信為始 以信為終} & 2025-10-11 & \href{http://www.church.com.hk/acms/content.asp?site=cdc\&op=show\&type=product\&code=032143\&layout=sermon}{\texttt{032143}} \\
以弗所書 2:11-22 & 禤凱榮牧師 & \hyperref[sec:032144]{近在咫尺,遠在天涯} & 2026-05-17 & \href{http://www.church.com.hk/acms/content.asp?site=cdc\&op=show\&type=product\&code=032144\&layout=sermon}{\texttt{032144}} \\
啟示錄 2:12-17 & 林穎豪傳道 & \hyperref[sec:032145]{別迦摩:勇敢卻走偏的教會} & 2026-05-17 & \href{http://www.church.com.hk/acms/content.asp?site=cdc\&op=show\&type=product\&code=032145\&layout=sermon}{\texttt{032145}} \\
使徒行傳 8:9-25 & 吳銘淇傳道 & \hyperref[sec:032146]{無法複製的聖靈工作} & 2026-05-17 & \href{http://www.church.com.hk/acms/content.asp?site=cdc\&op=show\&type=product\&code=032146\&layout=sermon}{\texttt{032146}} \\
約書亞記 12:7 & 黃克勤牧師博士 & \hyperref[sec:032147]{聖經人聖經事:從北到南的勝利,因應許而來} & 2026-05-17 & \href{http://www.church.com.hk/acms/content.asp?site=cdc\&op=show\&type=product\&code=032147\&layout=sermon}{\texttt{032147}} \\
利未記 6:8-13 & 周曉沂傳道 & \hyperref[sec:032148]{常常燒著的火} & 2026-05-17 & \href{http://www.church.com.hk/acms/content.asp?site=cdc\&op=show\&type=product\&code=032148\&layout=sermon}{\texttt{032148}} \\
撒迦利亞書 5:1-4 & 黃梅玲牧師 & \hyperref[sec:032149]{審判全地的神} & 2026-05-17 & \href{http://www.church.com.hk/acms/content.asp?site=cdc\&op=show\&type=product\&code=032149\&layout=sermon}{\texttt{032149}} \\
路加福音 2:41-52 & 曾興才牧師 & \hyperref[sec:032150]{成為一個充滿恩寵的人 (華)} & 2026-05-17 & \href{http://www.church.com.hk/acms/content.asp?site=cdc\&op=show\&type=product\&code=032150\&layout=sermon}{\texttt{032150}} \\
創世記 18:17-19 & 李業基牧師 & \hyperref[sec:032151]{何日再歸家?} & 2026-05-17 & \href{http://www.church.com.hk/acms/content.asp?site=cdc\&op=show\&type=product\&code=032151\&layout=sermon}{\texttt{032151}} \\
出埃及記 20:1-26 & 陳崇基牧師 & \hyperref[sec:032152]{石板} & 2026-05-17 & \href{http://www.church.com.hk/acms/content.asp?site=cdc\&op=show\&type=product\&code=032152\&layout=sermon}{\texttt{032152}} \\
約書亞記 17:14 & 黃克勤牧師博士 & \hyperref[sec:032153]{聖經人聖經事:不可僅有一鬮之地} & 2026-05-21 & \href{http://www.church.com.hk/acms/content.asp?site=cdc\&op=show\&type=product\&code=032153\&layout=sermon}{\texttt{032153}} \\
但以理書 6:10 & 容永祺先生 & \hyperref[sec:032154]{與素常一樣} & 2026-05-17 & \href{http://www.church.com.hk/acms/content.asp?site=cdc\&op=show\&type=product\&code=032154\&layout=sermon}{\texttt{032154}} \\
約翰福音 7:25-52 & 林鍵牧師 & \hyperref[sec:032155]{榮耀君王:活水靈風} & 2026-05-16 & \href{http://www.church.com.hk/acms/content.asp?site=cdc\&op=show\&type=product\&code=032155\&layout=sermon}{\texttt{032155}} \\
路加福音 16:19-31 & 俞大鳴牧師 & \hyperref[sec:032156]{錢財無效時} & 2026-05-17 & \href{http://www.church.com.hk/acms/content.asp?site=cdc\&op=show\&type=product\&code=032156\&layout=sermon}{\texttt{032156}} \\
使徒行傳 4:1-22 & 司徒贊傳道 & \hyperref[sec:032157]{聖靈本色} & 2026-05-17 & \href{http://www.church.com.hk/acms/content.asp?site=cdc\&op=show\&type=product\&code=032157\&layout=sermon}{\texttt{032157}} \\
列王記上 4:1-5:18 & 袁惠鈞牧師博士 & \hyperref[sec:032158]{繁榮中的伏筆,國度中的危機} & 2026-05-17 & \href{http://www.church.com.hk/acms/content.asp?site=cdc\&op=show\&type=product\&code=032158\&layout=sermon}{\texttt{032158}} \\
約書亞記 18:3 & 黃克勤牧師博士 & \hyperref[sec:032159]{聖經人聖經事:你們耽延不去得} & 2026-05-22 & \href{http://www.church.com.hk/acms/content.asp?site=cdc\&op=show\&type=product\&code=032159\&layout=sermon}{\texttt{032159}} \\
使徒行傳 15:36-41 & 李應新牧師 & \hyperref[sec:032160]{起了爭端,彼此分開} & 2026-05-17 & \href{http://www.church.com.hk/acms/content.asp?site=cdc\&op=show\&type=product\&code=032160\&layout=sermon}{\texttt{032160}} \\
約書亞記 1:1-9 & 馮永和博士 & \hyperref[sec:032161]{給所有母親的寄語:「剛強壯膽」} & 2026-05-10 & \href{http://www.church.com.hk/acms/content.asp?site=cdc\&op=show\&type=product\&code=032161\&layout=sermon}{\texttt{032161}} \\
詩篇 62:1-12 & 陳永就牧師 & \hyperref[sec:032162]{實在靠得住} & 2026-05-17 & \href{http://www.church.com.hk/acms/content.asp?site=cdc\&op=show\&type=product\&code=032162\&layout=sermon}{\texttt{032162}} \\
約翰福音 4:46-54 & 曹子光牧師 & \hyperref[sec:032163]{信心的建立} & 2026-05-10 & \href{http://www.church.com.hk/acms/content.asp?site=cdc\&op=show\&type=product\&code=032163\&layout=sermon}{\texttt{032163}} \\
士師記 3:8-31 & 梁永善牧師 & \hyperref[sec:032164]{士師記析讀 (5) 三師匯} & 2026-05-13 & \href{http://www.church.com.hk/acms/content.asp?site=cdc\&op=show\&type=product\&code=032164\&layout=sermon}{\texttt{032164}} \\
哥林多後書 5:18-20 & 葛國光牧師 & \hyperref[sec:032165]{作基督的大使 (華)} & 2026-05-17 & \href{http://www.church.com.hk/acms/content.asp?site=cdc\&op=show\&type=product\&code=032165\&layout=sermon}{\texttt{032165}} \\
路加福音 10:1-12 & 慎勇牧師 & \hyperref[sec:032166]{莊稼已熟,工人在哪裡?} & 2026-05-17 & \href{http://www.church.com.hk/acms/content.asp?site=cdc\&op=show\&type=product\&code=032166\&layout=sermon}{\texttt{032166}} \\
羅馬書 1:16-23 & 梁誠德牧師 & \hyperref[sec:032167]{懺悔錄} & 2026-05-17 & \href{http://www.church.com.hk/acms/content.asp?site=cdc\&op=show\&type=product\&code=032167\&layout=sermon}{\texttt{032167}} \\
使徒行傳 1:6-11 & 黃健羒牧師 & \hyperref[sec:032168]{耶穌升天} & 2026-05-17 & \href{http://www.church.com.hk/acms/content.asp?site=cdc\&op=show\&type=product\&code=032168\&layout=sermon}{\texttt{032168}} \\
以弗所書 6:10-24 & 葉慧敏傳道 & \hyperref[sec:032169]{在基督裡作剛強的人} & 2026-05-17 & \href{http://www.church.com.hk/acms/content.asp?site=cdc\&op=show\&type=product\&code=032169\&layout=sermon}{\texttt{032169}} \\
創世記 26:1-11 & 甘雪芬傳道 & \hyperref[sec:032170]{應許的傳承——神在危機中的保守} & 2026-05-16 & \href{http://www.church.com.hk/acms/content.asp?site=cdc\&op=show\&type=product\&code=032170\&layout=sermon}{\texttt{032170}} \\
使徒行傳 1:1-26 & 羅翠雲傳道 & \hyperref[sec:032171]{使命接力:靠聖靈而行的群體} & 2026-05-17 & \href{http://www.church.com.hk/acms/content.asp?site=cdc\&op=show\&type=product\&code=032171\&layout=sermon}{\texttt{032171}} \\
約書亞記 20:2 & 黃克勤牧師博士 & \hyperref[sec:032172]{聖經人聖經事:設立逃城} & 2026-05-22 & \href{http://www.church.com.hk/acms/content.asp?site=cdc\&op=show\&type=product\&code=032172\&layout=sermon}{\texttt{032172}} \\
以弗所書 3:1-21 & 駱瑞麟傳道 & \hyperref[sec:032173]{在受限中活出無限} & 2026-05-24 & \href{http://www.church.com.hk/acms/content.asp?site=cdc\&op=show\&type=product\&code=032173\&layout=sermon}{\texttt{032173}} \\
啟示錄 2:18-29 & 林穎豪傳道 & \hyperref[sec:032174]{推雅推喇:成長卻容忍錯誤的教會} & 2026-05-24 & \href{http://www.church.com.hk/acms/content.asp?site=cdc\&op=show\&type=product\&code=032174\&layout=sermon}{\texttt{032174}} \\
士師記 4:1-24 & 梁永善牧師 & \hyperref[sec:032175]{士師記析讀 (6) 謙招損,滿受益} & 2026-05-20 & \href{http://www.church.com.hk/acms/content.asp?site=cdc\&op=show\&type=product\&code=032175\&layout=sermon}{\texttt{032175}} \\
馬太福音 7:13-29 & 王正傑牧師 & \hyperref[sec:032176]{兩道門,你揀jor 邊道門?} & 2026-05-24 & \href{http://www.church.com.hk/acms/content.asp?site=cdc\&op=show\&type=product\&code=032176\&layout=sermon}{\texttt{032176}} \\
約翰福音 20:19-23 & 林玉英牧師 & \hyperref[sec:032177]{從囚禁到釋放} & 2026-05-24 & \href{http://www.church.com.hk/acms/content.asp?site=cdc\&op=show\&type=product\&code=032177\&layout=sermon}{\texttt{032177}} \\
使徒行傳 2:1-21 & 高國雄牧師 & \hyperref[sec:032178]{我們今天一同得這能力吧!} & 2026-05-24 & \href{http://www.church.com.hk/acms/content.asp?site=cdc\&op=show\&type=product\&code=032178\&layout=sermon}{\texttt{032178}} \\
約書亞記 1:8 & 王偉德傳道 & \hyperref[sec:032179]{一切從神的話開始} & 2026-05-24 & \href{http://www.church.com.hk/acms/content.asp?site=cdc\&op=show\&type=product\&code=032179\&layout=sermon}{\texttt{032179}} \\
哥林多前書 13:4-8 & 羅杰才牧師 & \hyperref[sec:032180]{愛的四重奏} & 2026-05-10 & \href{http://www.church.com.hk/acms/content.asp?site=cdc\&op=show\&type=product\&code=032180\&layout=sermon}{\texttt{032180}} \\
歌羅西書 3:1-4 & 黃紹權牧師 & \hyperref[sec:032181]{神榮耀的倒影——神的子民} & 2026-05-03 & \href{http://www.church.com.hk/acms/content.asp?site=cdc\&op=show\&type=product\&code=032181\&layout=sermon}{\texttt{032181}} \\
約書亞記 21:2 & 黃克勤牧師博士 & \hyperref[sec:032182]{聖經人聖經事:利未的城邑} & 2026-05-23 & \href{http://www.church.com.hk/acms/content.asp?site=cdc\&op=show\&type=product\&code=032182\&layout=sermon}{\texttt{032182}} \\
創世記 4:3-9 & 梁明財牧師 & \hyperref[sec:032183]{活出召命——生命的覺醒與承擔 (華)} & 2026-05-24 & \href{http://www.church.com.hk/acms/content.asp?site=cdc\&op=show\&type=product\&code=032183\&layout=sermon}{\texttt{032183}} \\
腓立比書 4:7 & 陳淳牧師 & \hyperref[sec:032184]{平安,平安:生命的守護} & 2026-05-24 & \href{http://www.church.com.hk/acms/content.asp?site=cdc\&op=show\&type=product\&code=032184\&layout=sermon}{\texttt{032184}} \\
使徒行傳 2:1-4 & 劉彤牧師 & \hyperref[sec:032185]{歡迎聖靈的風吹來 (華)} & 2026-05-24 & \href{http://www.church.com.hk/acms/content.asp?site=cdc\&op=show\&type=product\&code=032185\&layout=sermon}{\texttt{032185}} \\
創世記 2:7 & 招世超牧師 & \hyperref[sec:032186]{人是被創造,而非被製造} & 2026-05-24 & \href{http://www.church.com.hk/acms/content.asp?site=cdc\&op=show\&type=product\&code=032186\&layout=sermon}{\texttt{032186}} \\
歌羅西書 3:18-4:1 & 黃紹權牧師 & \hyperref[sec:032187]{群策群力所建造的基督群體} & 2026-05-24 & \href{http://www.church.com.hk/acms/content.asp?site=cdc\&op=show\&type=product\&code=032187\&layout=sermon}{\texttt{032187}} \\
約書亞記 22:34 & 黃克勤牧師博士 & \hyperref[sec:032188]{聖經人聖經事:證壇——證明耶和華是上帝} & 2026-05-25 & \href{http://www.church.com.hk/acms/content.asp?site=cdc\&op=show\&type=product\&code=032188\&layout=sermon}{\texttt{032188}} \\
歌羅西書 4:7-14 & 羅錫堅牧師 & \hyperref[sec:032189]{使徒保羅的監獄合照} & 2026-05-24 & \href{http://www.church.com.hk/acms/content.asp?site=cdc\&op=show\&type=product\&code=032189\&layout=sermon}{\texttt{032189}} \\
羅馬書 3:21-31 & 潘建堂牧師 & \hyperref[sec:032190]{最好的福氣} & 2025-10-26 & \href{http://www.church.com.hk/acms/content.asp?site=cdc\&op=show\&type=product\&code=032190\&layout=sermon}{\texttt{032190}} \\
以賽亞書 61:1-4 & 羅泰然牧師 & \hyperref[sec:032191]{破碎世代的城市宣教} & 2026-05-24 & \href{http://www.church.com.hk/acms/content.asp?site=cdc\&op=show\&type=product\&code=032191\&layout=sermon}{\texttt{032191}} \\
以弗所書 2:10 & 吳磐德牧師 & \hyperref[sec:032192]{建立健康的使命教會 (15) 得發揮的恩賜} & 2026-05-24 & \href{http://www.church.com.hk/acms/content.asp?site=cdc\&op=show\&type=product\&code=032192\&layout=sermon}{\texttt{032192}} \\
使徒行傳 11:1-30 & 余大峰傳道 & \hyperref[sec:032193]{同感一靈,突破規限的合一} & 2026-05-24 & \href{http://www.church.com.hk/acms/content.asp?site=cdc\&op=show\&type=product\&code=032193\&layout=sermon}{\texttt{032193}} \\
傳道書 4:4-5:7 & 譚文倩莊傳道 & \hyperref[sec:032194]{關係的失衡} & 2026-05-24 & \href{http://www.church.com.hk/acms/content.asp?site=cdc\&op=show\&type=product\&code=032194\&layout=sermon}{\texttt{032194}} \\
約書亞記 23:11 & 黃克勤牧師博士 & \hyperref[sec:032195]{聖經人聖經事:約書亞贈別 (1) 愛耶和華} & 2026-05-26 & \href{http://www.church.com.hk/acms/content.asp?site=cdc\&op=show\&type=product\&code=032195\&layout=sermon}{\texttt{032195}} \\
以賽亞書 56:1-8 & 文麗芳傳道 & \hyperref[sec:032196]{安息有時} & 2026-05-24 & \href{http://www.church.com.hk/acms/content.asp?site=cdc\&op=show\&type=product\&code=032196\&layout=sermon}{\texttt{032196}} \\
腓立比書 1:27 & 盧偉心弟兄 & \hyperref[sec:032197]{腓書回顧 (4) 相稱} & 2026-05-23 & \href{http://www.church.com.hk/acms/content.asp?site=cdc\&op=show\&type=product\&code=032197\&layout=sermon}{\texttt{032197}} \\
提多書 1:1-3:15 & 劉少康牧師 & \hyperref[sec:032198]{提多書:行為手冊} & 2026-05-20 & \href{http://www.church.com.hk/acms/content.asp?site=cdc\&op=show\&type=product\&code=032198\&layout=sermon}{\texttt{032198}} \\
士師記 5:7 & 梁永善牧師 & \hyperref[sec:032199]{士師記析讀 (7) 女權當道} & 2026-05-27 & \href{http://www.church.com.hk/acms/content.asp?site=cdc\&op=show\&type=product\&code=032199\&layout=sermon}{\texttt{032199}} \\
使徒行傳 2:1-47 & 鄭兆奇牧師 & \hyperref[sec:032200]{五旬節:教會的誕生} & 2026-05-24 & \href{http://www.church.com.hk/acms/content.asp?site=cdc\&op=show\&type=product\&code=032200\&layout=sermon}{\texttt{032200}} \\
加拉太書 5:2-6 & 葛國光牧師 & \hyperref[sec:032201]{靠律法 VS 靠信心 (華)} & 2026-05-24 & \href{http://www.church.com.hk/acms/content.asp?site=cdc\&op=show\&type=product\&code=032201\&layout=sermon}{\texttt{032201}} \\
羅馬書 15:1-6 & 張健庭牧師博士 & \hyperref[sec:032202]{彼此擔當} & 2026-05-24 & \href{http://www.church.com.hk/acms/content.asp?site=cdc\&op=show\&type=product\&code=032202\&layout=sermon}{\texttt{032202}} \\
民數記 11:1-35 & 林誠信牧師 & \hyperref[sec:032203]{百姓之怨言和摩西之不滿} & 2026-05-24 & \href{http://www.church.com.hk/acms/content.asp?site=cdc\&op=show\&type=product\&code=032203\&layout=sermon}{\texttt{032203}} \\
約書亞記 23:16 & 黃克勤牧師博士 & \hyperref[sec:032204]{聖經人聖經事:約書亞贈別 (2) 勿事奉別神} & 2026-05-27 & \href{http://www.church.com.hk/acms/content.asp?site=cdc\&op=show\&type=product\&code=032204\&layout=sermon}{\texttt{032204}} \\
歌羅西書 1:1-14 & 鄭卓雄牧師 & \hyperref[sec:032205]{在基督裡} & 2026-05-24 & \href{http://www.church.com.hk/acms/content.asp?site=cdc\&op=show\&type=product\&code=032205\&layout=sermon}{\texttt{032205}} \\
約翰福音 4:19-30 & 駱柱成牧師 & \hyperref[sec:032206]{真正的敬拜} & 2026-05-24 & \href{http://www.church.com.hk/acms/content.asp?site=cdc\&op=show\&type=product\&code=032206\&layout=sermon}{\texttt{032206}} \\
腓利門書 1:1-25 & 劉少康牧師 & \hyperref[sec:032207]{腓利門書:接納你的弟兄} & 2026-05-27 & \href{http://www.church.com.hk/acms/content.asp?site=cdc\&op=show\&type=product\&code=032207\&layout=sermon}{\texttt{032207}} \\
哥林多前書 12:12-13:13 & 王啟恩牧師 & \hyperref[sec:032208]{「將勤補拙」的肢體相交} & 2026-05-24 & \href{http://www.church.com.hk/acms/content.asp?site=cdc\&op=show\&type=product\&code=032208\&layout=sermon}{\texttt{032208}} \\
約書亞記 24:15 & 黃克勤牧師博士 & \hyperref[sec:032209]{聖經人聖經事:約書亞贈別 (3) 我家事奉神} & 2026-05-28 & \href{http://www.church.com.hk/acms/content.asp?site=cdc\&op=show\&type=product\&code=032209\&layout=sermon}{\texttt{032209}} \\
路加福音 19:1-10 & 陳傳華牧師博士 & \hyperref[sec:032210]{撒該:尋回真我} & 2026-05-24 & \href{http://www.church.com.hk/acms/content.asp?site=cdc\&op=show\&type=product\&code=032210\&layout=sermon}{\texttt{032210}} \\
以弗所書 4:1-16 & 禤凱榮牧師 & \hyperref[sec:032211]{沒有輪心的車輪走不動} & 2026-05-31 & \href{http://www.church.com.hk/acms/content.asp?site=cdc\&op=show\&type=product\&code=032211\&layout=sermon}{\texttt{032211}} \\
哥林多前書 8:1-13 & 馮志成牧師 & \hyperref[sec:032212]{因基督的愛而造就人} & 2026-05-31 & \href{http://www.church.com.hk/acms/content.asp?site=cdc\&op=show\&type=product\&code=032212\&layout=sermon}{\texttt{032212}} \\
使徒行傳 20:13-38 & 謝友德牧師 & \hyperref[sec:032213]{基督裡合而為一} & 2026-05-31 & \href{http://www.church.com.hk/acms/content.asp?site=cdc\&op=show\&type=product\&code=032213\&layout=sermon}{\texttt{032213}} \\
約書亞記 24:24 & 黃克勤牧師博士 & \hyperref[sec:032214]{聖經人聖經事:我們必事奉耶和華} & 2026-06-29 & \href{http://www.church.com.hk/acms/content.asp?site=cdc\&op=show\&type=product\&code=032214\&layout=sermon}{\texttt{032214}} \\
馬太福音 28:16-20 & 林家棋宣教師 & \hyperref[sec:032215]{大使命的謎思} & 2026-05-31 & \href{http://www.church.com.hk/acms/content.asp?site=cdc\&op=show\&type=product\&code=032215\&layout=sermon}{\texttt{032215}} \\
約拿書 1:1-17 & 張淑嫻牧師 & \hyperref[sec:032216]{神要找的人} & 2026-05-31 & \href{http://www.church.com.hk/acms/content.asp?site=cdc\&op=show\&type=product\&code=032216\&layout=sermon}{\texttt{032216}} \\
腓立比書 1:12-30 & 王永揚牧師 & \hyperref[sec:032221]{一切為福音!} & 2026-05-31 & \href{http://www.church.com.hk/acms/content.asp?site=cdc\&op=show\&type=product\&code=032221\&layout=sermon}{\texttt{032221}} \\
約翰福音 7:37-39 & 林鍵牧師 & \hyperref[sec:032222]{榮耀君王:世界的光} & 2026-05-30 & \href{http://www.church.com.hk/acms/content.asp?site=cdc\&op=show\&type=product\&code=032222\&layout=sermon}{\texttt{032222}} \\
詩篇 29:1-11 & 卓爾君牧師 & \hyperref[sec:032223]{風暴之上的平安:化解焦慮的三個行動 (華)} & 2026-05-31 & \href{http://www.church.com.hk/acms/content.asp?site=cdc\&op=show\&type=product\&code=032223\&layout=sermon}{\texttt{032223}} \\
瑪拉基書 2:1-9 & 石偉漢牧師 & \hyperref[sec:032224]{與神有約} & 2026-05-31 & \href{http://www.church.com.hk/acms/content.asp?site=cdc\&op=show\&type=product\&code=032224\&layout=sermon}{\texttt{032224}} \\
哈巴谷書 3:16-19 & 陳雅莉傳道 & \hyperref[sec:032225]{義人必因信得生} & 2026-05-31 & \href{http://www.church.com.hk/acms/content.asp?site=cdc\&op=show\&type=product\&code=032225\&layout=sermon}{\texttt{032225}} \\
約書亞記 23:1-16 & 方兆雄牧師 & \hyperref[sec:032226]{傳承敬虔與恩典} & 2026-05-31 & \href{http://www.church.com.hk/acms/content.asp?site=cdc\&op=show\&type=product\&code=032226\&layout=sermon}{\texttt{032226}} \\
士師記 2:12 & 黃克勤牧師博士 & \hyperref[sec:032227]{聖經人聖經事:惹耶和華發怒} & 2026-06-01 & \href{http://www.church.com.hk/acms/content.asp?site=cdc\&op=show\&type=product\&code=032227\&layout=sermon}{\texttt{032227}} \\
路加福音 19:11-27 & 梁成裕牧師 & \hyperref[sec:032228]{在王離去與再來之間:作忠心的僕人} & 2026-05-31 & \href{http://www.church.com.hk/acms/content.asp?site=cdc\&op=show\&type=product\&code=032228\&layout=sermon}{\texttt{032228}} \\
士師記 6:1-40 & 梁永善牧師 & \hyperref[sec:032229]{士師記析讀 (8) 要有憑據} & 2026-06-03 & \href{http://www.church.com.hk/acms/content.asp?site=cdc\&op=show\&type=product\&code=032229\&layout=sermon}{\texttt{032229}} \\
詩篇 61:1-8 & 高志翔牧師 & \hyperref[sec:032230]{領我到更高的磐石 (華)} & 2026-05-31 & \href{http://www.church.com.hk/acms/content.asp?site=cdc\&op=show\&type=product\&code=032230\&layout=sermon}{\texttt{032230}} \\
詩篇 73:1-28 & 張健庭牧師博士 & \hyperref[sec:032231]{在聖所中明白人生} & 2026-05-31 & \href{http://www.church.com.hk/acms/content.asp?site=cdc\&op=show\&type=product\&code=032231\&layout=sermon}{\texttt{032231}} \\
士師記 2:19 & 黃克勤牧師博士 & \hyperref[sec:032232]{聖經人聖經事:士師記的惡性循環} & 2026-06-02 & \href{http://www.church.com.hk/acms/content.asp?site=cdc\&op=show\&type=product\&code=032232\&layout=sermon}{\texttt{032232}} \\
民數記 12:1-16 & 何嘉敏傳道 & \hyperref[sec:032233]{從比較到復和的兄姊弟情} & 2026-05-30 & \href{http://www.church.com.hk/acms/content.asp?site=cdc\&op=show\&type=product\&code=032233\&layout=sermon}{\texttt{032233}} \\
民數記 12:1-16 & 吳安強牧師 & \hyperref[sec:032234]{嫉妒與毀謗的可惡} & 2026-05-31 & \href{http://www.church.com.hk/acms/content.asp?site=cdc\&op=show\&type=product\&code=032234\&layout=sermon}{\texttt{032234}} \\
以弗所書 2:1-10 & 姚臣泰長老 & \hyperref[sec:032235]{神極大的救恩 (華/粵)} & 2026-05-31 & \href{http://www.church.com.hk/acms/content.asp?site=cdc\&op=show\&type=product\&code=032235\&layout=sermon}{\texttt{032235}} \\
約伯記 22:21 & 曹子光牧師 & \hyperref[sec:032236]{你要認識神} & 2026-05-31 & \href{http://www.church.com.hk/acms/content.asp?site=cdc\&op=show\&type=product\&code=032236\&layout=sermon}{\texttt{032236}} \\
士師記 3:1 & 黃克勤牧師博士 & \hyperref[sec:032237]{聖經人聖經事:留下這幾族 試驗以色列} & 2026-06-03 & \href{http://www.church.com.hk/acms/content.asp?site=cdc\&op=show\&type=product\&code=032237\&layout=sermon}{\texttt{032237}} \\
以賽亞書 65:17-25 & 羅婉明傳道 & \hyperref[sec:032238]{新天新地} & 2026-05-31 & \href{http://www.church.com.hk/acms/content.asp?site=cdc\&op=show\&type=product\&code=032238\&layout=sermon}{\texttt{032238}} \\
申命記 8:1-10 & 符氣松牧師 & \hyperref[sec:032239]{感恩 MODE (1) 記念思想} & 2026-05-30 & \href{http://www.church.com.hk/acms/content.asp?site=cdc\&op=show\&type=product\&code=032239\&layout=sermon}{\texttt{032239}} \\
創世記 4:1-16 & 何志滌牧師 & \hyperref[sec:032240]{勝過嫉妒} & 2026-05-31 & \href{http://www.church.com.hk/acms/content.asp?site=cdc\&op=show\&type=product\&code=032240\&layout=sermon}{\texttt{032240}} \\
路加福音 17:1-10 & 俞大鳴牧師 & \hyperref[sec:032241]{門徒要順服} & 2026-05-24 & \href{http://www.church.com.hk/acms/content.asp?site=cdc\&op=show\&type=product\&code=032241\&layout=sermon}{\texttt{032241}} \\
路加福音 17:11-19 & 俞大鳴牧師 & \hyperref[sec:032242]{感恩} & 2026-05-31 & \href{http://www.church.com.hk/acms/content.asp?site=cdc\&op=show\&type=product\&code=032242\&layout=sermon}{\texttt{032242}} \\
創世記 26:12-33 & 秦民德傳道 & \hyperref[sec:032243]{行在「恩約」中} & 2026-05-30 & \href{http://www.church.com.hk/acms/content.asp?site=cdc\&op=show\&type=product\&code=032243\&layout=sermon}{\texttt{032243}} \\
使徒行傳 2:14-36 & 黃浩君傳道 & \hyperref[sec:032244]{仰賴基督與聖靈作見證} & 2026-05-31 & \href{http://www.church.com.hk/acms/content.asp?site=cdc\&op=show\&type=product\&code=032244\&layout=sermon}{\texttt{032244}} \\
士師記 3:9 & 黃克勤牧師博士 & \hyperref[sec:032245]{聖經人聖經事:第一位士師 俄陀聶} & 2026-06-04 & \href{http://www.church.com.hk/acms/content.asp?site=cdc\&op=show\&type=product\&code=032245\&layout=sermon}{\texttt{032245}} \\
羅馬書 4:20-25 & 莊明偉牧師 & \hyperref[sec:032246]{因信而剛強} & 2025-11-02 & \href{http://www.church.com.hk/acms/content.asp?site=cdc\&op=show\&type=product\&code=032246\&layout=sermon}{\texttt{032246}} \\
創世記 12:1-9 & 李均熊牧師博士 & \hyperref[sec:032247]{放心……你的信治好你了} & 2026-06-07 & \href{http://www.church.com.hk/acms/content.asp?site=cdc\&op=show\&type=product\&code=032247\&layout=sermon}{\texttt{032247}} \\
以弗所書 4:17-32 & 禤凱榮牧師 & \hyperref[sec:032248]{皇帝賜的新衣} & 2026-06-07 & \href{http://www.church.com.hk/acms/content.asp?site=cdc\&op=show\&type=product\&code=032248\&layout=sermon}{\texttt{032248}} \\
歌羅西書 4:1-18 & 黃紹權牧師 & \hyperref[sec:032249]{建構在屬靈支援中的基督社群} & 2026-06-07 & \href{http://www.church.com.hk/acms/content.asp?site=cdc\&op=show\&type=product\&code=032249\&layout=sermon}{\texttt{032249}} \\
啟示錄 3:1-6 & 趙崇健牧師 & \hyperref[sec:032250]{撒狄:好像活著卻是死的教會} & 2026-06-07 & \href{http://www.church.com.hk/acms/content.asp?site=cdc\&op=show\&type=product\&code=032250\&layout=sermon}{\texttt{032250}} \\
何西阿書 5:15-6:6 & 趙承恩宣教師 & \hyperref[sec:032251]{真誠的侮罪} & 2026-06-07 & \href{http://www.church.com.hk/acms/content.asp?site=cdc\&op=show\&type=product\&code=032251\&layout=sermon}{\texttt{032251}} \\
士師記 3:15 & 黃克勤牧師博士 & \hyperref[sec:032252]{聖經人聖經事:第二位士師 以芴} & 2026-06-05 & \href{http://www.church.com.hk/acms/content.asp?site=cdc\&op=show\&type=product\&code=032252\&layout=sermon}{\texttt{032252}} \\
馬可福音 3:1-6 & 羅泰然牧師 & \hyperref[sec:032253]{忘我的愛} & 2026-06-07 & \href{http://www.church.com.hk/acms/content.asp?site=cdc\&op=show\&type=product\&code=032253\&layout=sermon}{\texttt{032253}} \\
哥林多後書 5:14-19 & 林耀楠牧師 & \hyperref[sec:032254]{建立健康的使命教會 (17) 立定心志作基督的使者} & 2026-06-07 & \href{http://www.church.com.hk/acms/content.asp?site=cdc\&op=show\&type=product\&code=032254\&layout=sermon}{\texttt{032254}} \\
以弗所書 1:3-14 & 謝友德牧師 & \hyperref[sec:032255]{無限靈福} & 2026-06-07 & \href{http://www.church.com.hk/acms/content.asp?site=cdc\&op=show\&type=product\&code=032255\&layout=sermon}{\texttt{032255}} \\
出埃及記 25:1-40:38 & 陳崇基牧師 & \hyperref[sec:032256]{會幕} & 2026-06-07 & \href{http://www.church.com.hk/acms/content.asp?site=cdc\&op=show\&type=product\&code=032256\&layout=sermon}{\texttt{032256}} \\
尼希米記 13:1-31 & 王正傑牧師 & \hyperref[sec:032257]{更新} & 2026-06-07 & \href{http://www.church.com.hk/acms/content.asp?site=cdc\&op=show\&type=product\&code=032257\&layout=sermon}{\texttt{032257}} \\
士師記 4:4 & 黃克勤牧師博士 & \hyperref[sec:032258]{聖經人聖經事:第三位士師  底波拉 巴拉} & 2026-06-06 & \href{http://www.church.com.hk/acms/content.asp?site=cdc\&op=show\&type=product\&code=032258\&layout=sermon}{\texttt{032258}} \\
詩篇 126:1-6 & 羅錫堅牧師 & \hyperref[sec:032259]{耶和華為我們行了大事} & 2026-06-09 & \href{http://www.church.com.hk/acms/content.asp?site=cdc\&op=show\&type=product\&code=032259\&layout=sermon}{\texttt{032259}} \\
瑪拉基書 2:10-16 & 石偉漢牧師 & \hyperref[sec:032260]{向行事詭詐說:不!} & 2026-06-07 & \href{http://www.church.com.hk/acms/content.asp?site=cdc\&op=show\&type=product\&code=032260\&layout=sermon}{\texttt{032260}} \\
哥林多後書 2:14 & 劉彤牧師 & \hyperref[sec:032261]{你是個得勝者 (華)} & 2026-06-07 & \href{http://www.church.com.hk/acms/content.asp?site=cdc\&op=show\&type=product\&code=032261\&layout=sermon}{\texttt{032261}} \\
約書亞記 2:1-24 & 李業基牧師 & \hyperref[sec:032262]{恩典的記號} & 2026-06-07 & \href{http://www.church.com.hk/acms/content.asp?site=cdc\&op=show\&type=product\&code=032262\&layout=sermon}{\texttt{032262}} \\
士師記 6:1 & 黃克勤牧師博士 & \hyperref[sec:032263]{聖經人聖經事:米甸人壓制 以色列人} & 2026-06-07 & \href{http://www.church.com.hk/acms/content.asp?site=cdc\&op=show\&type=product\&code=032263\&layout=sermon}{\texttt{032263}} \\
創世記 4:10-16 & 梁明財牧師 & \hyperref[sec:032264]{活出召命——生命的覺醒與承擔 II (華)} & 2026-06-07 & \href{http://www.church.com.hk/acms/content.asp?site=cdc\&op=show\&type=product\&code=032264\&layout=sermon}{\texttt{032264}} \\
使徒行傳 18:1-3 & 劉銳鋒牧師 & \hyperref[sec:032265]{宣教的神國俠侶} & 2026-06-07 & \href{http://www.church.com.hk/acms/content.asp?site=cdc\&op=show\&type=product\&code=032265\&layout=sermon}{\texttt{032265}} \\
羅馬書 5:17-21 & 蕭壽華牧師 & \hyperref[sec:032266]{靠著耶穌得勝} & 2025-11-09 & \href{http://www.church.com.hk/acms/content.asp?site=cdc\&op=show\&type=product\&code=032266\&layout=sermon}{\texttt{032266}} \\
馬可福音 1:1-45 & 秦潤祥牧師 & \hyperref[sec:032267]{不再無間道?} & 2026-06-07 & \href{http://www.church.com.hk/acms/content.asp?site=cdc\&op=show\&type=product\&code=032267\&layout=sermon}{\texttt{032267}} \\
羅馬書 8:18-27 & 葛國光牧師 & \hyperref[sec:032268]{苦難世界的盼望 (華)} & 2026-06-07 & \href{http://www.church.com.hk/acms/content.asp?site=cdc\&op=show\&type=product\&code=032268\&layout=sermon}{\texttt{032268}} \\
約翰福音 4:31-42 & 駱柱成牧師 & \hyperref[sec:032269]{遵行神旨,作成主工} & 2026-06-07 & \href{http://www.church.com.hk/acms/content.asp?site=cdc\&op=show\&type=product\&code=032269\&layout=sermon}{\texttt{032269}} \\
士師記 6:12 & 黃克勤牧師博士 & \hyperref[sec:032270]{聖經人聖經事:耶和華的使者 呼召基甸} & 2026-06-08 & \href{http://www.church.com.hk/acms/content.asp?site=cdc\&op=show\&type=product\&code=032270\&layout=sermon}{\texttt{032270}} \\
約翰福音 8:21-29 & 林鍵牧師 & \hyperref[sec:032271]{榮耀君王:死在罪中} & 2026-06-06 & \href{http://www.church.com.hk/acms/content.asp?site=cdc\&op=show\&type=product\&code=032271\&layout=sermon}{\texttt{032271}} \\
民數記 13:1-14:45 & 林誠信牧師 & \hyperref[sec:032272]{信不信神應許的結果} & 2026-06-07 & \href{http://www.church.com.hk/acms/content.asp?site=cdc\&op=show\&type=product\&code=032272\&layout=sermon}{\texttt{032272}} \\
希伯來書 1:1-3 & 張健庭牧師博士 & \hyperref[sec:032273]{希伯來書系列 (1) 進入聖經中的奇異世界} & 2026-06-07 & \href{http://www.church.com.hk/acms/content.asp?site=cdc\&op=show\&type=product\&code=032273\&layout=sermon}{\texttt{032273}} \\
士師記 7:1-25 & 梁永善牧師 & \hyperref[sec:032274]{士師記析讀 (9) 以少勝多} & 2026-06-10 & \href{http://www.church.com.hk/acms/content.asp?site=cdc\&op=show\&type=product\&code=032274\&layout=sermon}{\texttt{032274}} \\
列王記上 9:3-9 & 袁惠鈞牧師博士 & \hyperref[sec:032275]{華麗殿宇背後:神看重的不是建築,而是內心} & 2026-06-07 & \href{http://www.church.com.hk/acms/content.asp?site=cdc\&op=show\&type=product\&code=032275\&layout=sermon}{\texttt{032275}} \\
尼希米記 1:1-4 & 顏沛恩牧師 & \hyperref[sec:032276]{Keep Burning} & 2026-06-07 & \href{http://www.church.com.hk/acms/content.asp?site=cdc\&op=show\&type=product\&code=032276\&layout=sermon}{\texttt{032276}} \\
士師記 6:17 & 黃克勤牧師博士 & \hyperref[sec:032277]{聖經人聖經事:基甸第一次求證據} & 2026-06-12 & \href{http://www.church.com.hk/acms/content.asp?site=cdc\&op=show\&type=product\&code=032277\&layout=sermon}{\texttt{032277}} \\
希伯來書 9:22-28 & 謝又生牧師 & \hyperref[sec:032278]{向過去說再見,向未來說嗨} & 2026-06-07 & \href{http://www.church.com.hk/acms/content.asp?site=cdc\&op=show\&type=product\&code=032278\&layout=sermon}{\texttt{032278}} \\
馬太福音 3:1-3 & 王梯青長老 & \hyperref[sec:032279]{天國的福音} & 2026-06-07 & \href{http://www.church.com.hk/acms/content.asp?site=cdc\&op=show\&type=product\&code=032279\&layout=sermon}{\texttt{032279}} \\
詩篇 78:1-8 & 葉慧敏傳道 & \hyperref[sec:032280]{信仰的接力} & 2026-06-07 & \href{http://www.church.com.hk/acms/content.asp?site=cdc\&op=show\&type=product\&code=032280\&layout=sermon}{\texttt{032280}} \\
出埃及記 20:1-5 & 招世超牧師 & \hyperref[sec:032281]{當A.I.成為偶像} & 2026-06-07 & \href{http://www.church.com.hk/acms/content.asp?site=cdc\&op=show\&type=product\&code=032281\&layout=sermon}{\texttt{032281}} \\
士師記 6:32 & 黃克勤牧師博士 & \hyperref[sec:032282]{聖經人聖經事:基甸第一次勝仗 耶路•巴力} & 2026-06-10 & \href{http://www.church.com.hk/acms/content.asp?site=cdc\&op=show\&type=product\&code=032282\&layout=sermon}{\texttt{032282}} \\
馬太福音 9:35-10:15 & 陳傳華牧師博士 & \hyperref[sec:032283]{門徒的天國使命} & 2026-06-07 & \href{http://www.church.com.hk/acms/content.asp?site=cdc\&op=show\&type=product\&code=032283\&layout=sermon}{\texttt{032283}} \\
馬太福音 25:14-30 & 李文康傳道 & \hyperref[sec:032284]{善用今生,拓展神國} & 2026-06-06 & \href{http://www.church.com.hk/acms/content.asp?site=cdc\&op=show\&type=product\&code=032284\&layout=sermon}{\texttt{032284}} \\
使徒行傳 2:37-47 & 翁曉怡傳道 & \hyperref[sec:032285]{屬乎耶穌基督的群體} & 2026-06-07 & \href{http://www.church.com.hk/acms/content.asp?site=cdc\&op=show\&type=product\&code=032285\&layout=sermon}{\texttt{032285}} \\
哥林多後書 10:1-5 & 梁錫強醫生 & \hyperref[sec:032286]{勝利不在乎刀劍} & 2026-06-14 & \href{http://www.church.com.hk/acms/content.asp?site=cdc\&op=show\&type=product\&code=032286\&layout=sermon}{\texttt{032286}} \\
約翰福音 13:31-38 & 黃紹權牧師 & \hyperref[sec:032287]{主的榮光往哪裡去了?} & 2026-06-14 & \href{http://www.church.com.hk/acms/content.asp?site=cdc\&op=show\&type=product\&code=032287\&layout=sermon}{\texttt{032287}} \\
帖撒羅尼迦前書 1:1-10 & 梁康民牧師 & \hyperref[sec:032288]{培育生命,建立教會:感恩有你!(粵/英)} & 2026-06-14 & \href{http://www.church.com.hk/acms/content.asp?site=cdc\&op=show\&type=product\&code=032288\&layout=sermon}{\texttt{032288}} \\
士師記 6:39 & 黃克勤牧師博士 & \hyperref[sec:032289]{聖經人聖經事:基甸用第二,第三次求證} & 2026-06-11 & \href{http://www.church.com.hk/acms/content.asp?site=cdc\&op=show\&type=product\&code=032289\&layout=sermon}{\texttt{032289}} \\
羅馬書 5:1-8 & 林玉英牧師 & \hyperref[sec:032290]{看透我,依然愛我} & 2026-06-14 & \href{http://www.church.com.hk/acms/content.asp?site=cdc\&op=show\&type=product\&code=032290\&layout=sermon}{\texttt{032290}} \\
羅馬書 7:13-25 & 周曉暉牧師 & \hyperref[sec:032291]{別害怕,祂知你難受!} & 2025-11-15 & \href{http://www.church.com.hk/acms/content.asp?site=cdc\&op=show\&type=product\&code=032291\&layout=sermon}{\texttt{032291}} \\
創世記 18:1-15 & 高國雄牧師 & \hyperref[sec:032292]{在主前蒙恩} & 2026-06-14 & \href{http://www.church.com.hk/acms/content.asp?site=cdc\&op=show\&type=product\&code=032292\&layout=sermon}{\texttt{032292}} \\
瑪拉基書 2:17-3:5 & 石偉漢牧師 & \hyperref[sec:032293]{你都有今日} & 2026-06-14 & \href{http://www.church.com.hk/acms/content.asp?site=cdc\&op=show\&type=product\&code=032293\&layout=sermon}{\texttt{032293}} \\
以弗所書 5:1-20 & 孫雷雷牧師 & \hyperref[sec:032294]{活出你是誰 (華)} & 2026-06-14 & \href{http://www.church.com.hk/acms/content.asp?site=cdc\&op=show\&type=product\&code=032294\&layout=sermon}{\texttt{032294}} \\
啟示錄 3:7-13 & 林穎豪傳道 & \hyperref[sec:032295]{非拉鐵非教會:軟弱卻又強壯的教會} & 2026-06-14 & \href{http://www.church.com.hk/acms/content.asp?site=cdc\&op=show\&type=product\&code=032295\&layout=sermon}{\texttt{032295}} \\
士師記 7:2 & 黃克勤牧師博士 & \hyperref[sec:032296]{聖經人聖經事:免得以色列人向我自誇} & 2026-06-12 & \href{http://www.church.com.hk/acms/content.asp?site=cdc\&op=show\&type=product\&code=032296\&layout=sermon}{\texttt{032296}} \\
約翰福音 8:30-36 & 林鍵牧師 & \hyperref[sec:032297]{榮耀君王:上帝兒女} & 2026-06-13 & \href{http://www.church.com.hk/acms/content.asp?site=cdc\&op=show\&type=product\&code=032297\&layout=sermon}{\texttt{032297}} \\
歷代志上 22:1-13 & 林于謙宣教士 & \hyperref[sec:032298]{上帝的居所:從聖殿到子民} & 2026-06-14 & \href{http://www.church.com.hk/acms/content.asp?site=cdc\&op=show\&type=product\&code=032298\&layout=sermon}{\texttt{032298}} \\
以弗所書 2:11-22 & 鄺偉志牧師 & \hyperref[sec:032299]{合一和睦ge 代價唔簡單} & 2026-06-14 & \href{http://www.church.com.hk/acms/content.asp?site=cdc\&op=show\&type=product\&code=032299\&layout=sermon}{\texttt{032299}} \\
希伯來書 10:19-20 & 曾興才牧師 & \hyperref[sec:032300]{每一次你禱告 (華)} & 2026-06-14 & \href{http://www.church.com.hk/acms/content.asp?site=cdc\&op=show\&type=product\&code=032300\&layout=sermon}{\texttt{032300}} \\
約翰福音 4:1-26 & 梁紫健傳道 & \hyperref[sec:032301]{放下水桶} & 2026-06-14 & \href{http://www.church.com.hk/acms/content.asp?site=cdc\&op=show\&type=product\&code=032301\&layout=sermon}{\texttt{032301}} \\
士師記 7:7 & 黃克勤牧師博士 & \hyperref[sec:032302]{聖經人聖經事:用手捧著舔水的三百人} & 2026-06-13 & \href{http://www.church.com.hk/acms/content.asp?site=cdc\&op=show\&type=product\&code=032302\&layout=sermon}{\texttt{032302}} \\
馬可福音 4:35-41 & 倫國平牧師 & \hyperref[sec:032303]{真正的平安  (見證:張瑪莉姊妹)} & 2026-06-14 & \href{http://www.church.com.hk/acms/content.asp?site=cdc\&op=show\&type=product\&code=032303\&layout=sermon}{\texttt{032303}} \\
俄巴底亞書 1:1-9 & 鄭兆奇牧師 & \hyperref[sec:032304]{驕傲的以東} & 2026-06-14 & \href{http://www.church.com.hk/acms/content.asp?site=cdc\&op=show\&type=product\&code=032304\&layout=sermon}{\texttt{032304}} \\
加拉太書 5:7-12 & 葛國光牧師 & \hyperref[sec:032305]{直奔天路 (華)} & 2026-06-14 & \href{http://www.church.com.hk/acms/content.asp?site=cdc\&op=show\&type=product\&code=032305\&layout=sermon}{\texttt{032305}} \\
撒母耳記上 26:1-25 & 李思聰牧師 & \hyperref[sec:032306]{兩種人生,兩種死亡} & 2026-06-14 & \href{http://www.church.com.hk/acms/content.asp?site=cdc\&op=show\&type=product\&code=032306\&layout=sermon}{\texttt{032306}} \\
士師記 7:13 & 黃克勤牧師博士 & \hyperref[sec:032307]{聖經人聖經事:一個夢:一個大麥餅} & 2026-06-14 & \href{http://www.church.com.hk/acms/content.asp?site=cdc\&op=show\&type=product\&code=032307\&layout=sermon}{\texttt{032307}} \\
以弗所書 6:1-4 & 杜崇基牧師 & \hyperref[sec:032308]{把孩子帶到主面前} & 2026-06-14 & \href{http://www.church.com.hk/acms/content.asp?site=cdc\&op=show\&type=product\&code=032308\&layout=sermon}{\texttt{032308}} \\
申命記 9:1-29 & 王啟恩牧師 & \hyperref[sec:032309]{感恩 MODE (2) 恩典中的感恩} & 2026-06-13 & \href{http://www.church.com.hk/acms/content.asp?site=cdc\&op=show\&type=product\&code=032309\&layout=sermon}{\texttt{032309}} \\
腓立比書 1:3-6 & 唐郭麗婷牧師 & \hyperref[sec:032310]{多代同行•與同輩同行:追求共同目標} & 2026-06-14 & \href{http://www.church.com.hk/acms/content.asp?site=cdc\&op=show\&type=product\&code=032310\&layout=sermon}{\texttt{032310}} \\
列王記上 8:1-29 & 袁惠鈞牧師博士 & \hyperref[sec:032311]{耶和華的榮耀充滿聖殿:神同在的高峰,人心順服的考驗} & 2026-06-14 & \href{http://www.church.com.hk/acms/content.asp?site=cdc\&op=show\&type=product\&code=032311\&layout=sermon}{\texttt{032311}} \\
希伯來書 5:7-10 & 戴永富博士 & \hyperref[sec:032312]{在基督裡效法基督 (華/粵)} & 2026-06-14 & \href{http://www.church.com.hk/acms/content.asp?site=cdc\&op=show\&type=product\&code=032312\&layout=sermon}{\texttt{032312}} \\
約翰福音 1:29-36 & 李應新牧師 & \hyperref[sec:032313]{神的羔羊,以靈施洗} & 2026-06-14 & \href{http://www.church.com.hk/acms/content.asp?site=cdc\&op=show\&type=product\&code=032313\&layout=sermon}{\texttt{032313}} \\
士師記 7:18 & 黃克勤牧師博士 & \hyperref[sec:032314]{聖經人聖經事:耶和華和基甸的刀} & 2026-06-15 & \href{http://www.church.com.hk/acms/content.asp?site=cdc\&op=show\&type=product\&code=032314\&layout=sermon}{\texttt{032314}} \\
羅馬書 8:1-17 & 周依娜牧師 & \hyperref[sec:032315]{惟獨在基督裡} & 2025-11-23 & \href{http://www.church.com.hk/acms/content.asp?site=cdc\&op=show\&type=product\&code=032315\&layout=sermon}{\texttt{032315}} \\
馬可福音 5:21-43 & 李成傑傳道 & \hyperref[sec:032316]{兩極連結於基督,絕望中的信心轉向} & 2026-06-14 & \href{http://www.church.com.hk/acms/content.asp?site=cdc\&op=show\&type=product\&code=032316\&layout=sermon}{\texttt{032316}} \\
以弗所書 4:1-6 & 歐陽萬廷弟兄 & \hyperref[sec:032317]{在基督裡合一} & 2026-06-14 & \href{http://www.church.com.hk/acms/content.asp?site=cdc\&op=show\&type=product\&code=032317\&layout=sermon}{\texttt{032317}} \\
箴言 25:11 & 曹子光牧師 & \hyperref[sec:032318]{金蘋果} & 2026-06-14 & \href{http://www.church.com.hk/acms/content.asp?site=cdc\&op=show\&type=product\&code=032318\&layout=sermon}{\texttt{032318}} \\
路加福音 24:13-35 & 李適清博士 & \hyperref[sec:032319]{原來主一直同在} & 2026-05-03 & \href{http://www.church.com.hk/acms/content.asp?site=cdc\&op=show\&type=product\&code=032319\&layout=sermon}{\texttt{032319}} \\
羅馬書 8:22-39 & 周曉暉牧師 & \hyperref[sec:032320]{靠著愛我們的那一位} & 2025-11-30 & \href{http://www.church.com.hk/acms/content.asp?site=cdc\&op=show\&type=product\&code=032320\&layout=sermon}{\texttt{032320}} \\
士師記 8:1 & 黃克勤牧師博士 & \hyperref[sec:032321]{聖經人聖經事:為甚麼沒有招我們同去} & 2026-06-16 & \href{http://www.church.com.hk/acms/content.asp?site=cdc\&op=show\&type=product\&code=032321\&layout=sermon}{\texttt{032321}} \\
以弗所書 5:21-6:9 & 禤凱榮牧師 & \hyperref[sec:032322]{飛翔需要兩隻翅膀} & 2026-06-21 & \href{http://www.church.com.hk/acms/content.asp?site=cdc\&op=show\&type=product\&code=032322\&layout=sermon}{\texttt{032322}} \\
俄巴底亞書 1:1-21 & 黃紹權牧師 & \hyperref[sec:032323]{不要沾沾自喜?} & 2026-06-21 & \href{http://www.church.com.hk/acms/content.asp?site=cdc\&op=show\&type=product\&code=032323\&layout=sermon}{\texttt{032323}} \\
啟示錄 3:14-22 & 譚綺敏傳道 & \hyperref[sec:032324]{老底嘉:富足卻又貧窮的教會} & 2026-06-21 & \href{http://www.church.com.hk/acms/content.asp?site=cdc\&op=show\&type=product\&code=032324\&layout=sermon}{\texttt{032324}} \\
創世記 21:8-20 & 蘇善美宣教師 & \hyperref[sec:032325]{天父與人同渡時艱} & 2026-06-21 & \href{http://www.church.com.hk/acms/content.asp?site=cdc\&op=show\&type=product\&code=032325\&layout=sermon}{\texttt{032325}} \\
士師記 8:11 & 黃克勤牧師博士 & \hyperref[sec:032326]{聖經人聖經事:因為他們坦然無懼!} & 2026-06-17 & \href{http://www.church.com.hk/acms/content.asp?site=cdc\&op=show\&type=product\&code=032326\&layout=sermon}{\texttt{032326}} \\
約翰福音 10:1-18 & 鄧瑞強牧師博士 & \hyperref[sec:032327]{從一扇門,到心深處} & 2026-06-21 & \href{http://www.church.com.hk/acms/content.asp?site=cdc\&op=show\&type=product\&code=032327\&layout=sermon}{\texttt{032327}} \\
歌羅西書 3:18-25 & 嚴建平博士 & \hyperref[sec:032328]{不要惹兒女的氣 (英/粵)} & 2026-06-23 & \href{http://www.church.com.hk/acms/content.asp?site=cdc\&op=show\&type=product\&code=032328\&layout=sermon}{\texttt{032328}} \\
瑪拉基書 3:6-12 & 石偉漢牧師 & \hyperref[sec:032329]{福浸眼眉的驚喜} & 2026-06-14 & \href{http://www.church.com.hk/acms/content.asp?site=cdc\&op=show\&type=product\&code=032329\&layout=sermon}{\texttt{032329}} \\
以弗所書 4:1-16 & 吳銘淇傳道 & \hyperref[sec:032330]{合一的關鍵} & 2026-06-21 & \href{http://www.church.com.hk/acms/content.asp?site=cdc\&op=show\&type=product\&code=032330\&layout=sermon}{\texttt{032330}} \\
士師記 8:21 & 黃克勤牧師博士 & \hyperref[sec:032331]{聖經人聖經事:人如何,力量也如何} & 2026-06-18 & \href{http://www.church.com.hk/acms/content.asp?site=cdc\&op=show\&type=product\&code=032331\&layout=sermon}{\texttt{032331}} \\
路加福音 19:1-10 & 麥偉祺牧師 & \hyperref[sec:032332]{從大起大落到大喜大樂的人生} & 2026-06-21 & \href{http://www.church.com.hk/acms/content.asp?site=cdc\&op=show\&type=product\&code=032332\&layout=sermon}{\texttt{032332}} \\
馬太福音 11:20-30 & 魏文俊傳道 & \hyperref[sec:032333]{一同找到安息的智慧} & 2026-06-14 & \href{http://www.church.com.hk/acms/content.asp?site=cdc\&op=show\&type=product\&code=032333\&layout=sermon}{\texttt{032333}} \\
申命記 6:6-7 & 劉彤牧師 & \hyperref[sec:032334]{信仰的傳承 (華)} & 2026-06-21 & \href{http://www.church.com.hk/acms/content.asp?site=cdc\&op=show\&type=product\&code=032334\&layout=sermon}{\texttt{032334}} \\
傳道書 5:8-6:9 & 譚文倩莊傳道 & \hyperref[sec:032335]{錢值幾錢?} & 2026-06-21 & \href{http://www.church.com.hk/acms/content.asp?site=cdc\&op=show\&type=product\&code=032335\&layout=sermon}{\texttt{032335}} \\
馬可福音 2:1-12 & 梁廷益牧師 & \hyperref[sec:032336]{爸爸的心 (粵/華)} & 2026-06-21 & \href{http://www.church.com.hk/acms/content.asp?site=cdc\&op=show\&type=product\&code=032336\&layout=sermon}{\texttt{032336}} \\
士師記 8:23 & 黃克勤牧師博士 & \hyperref[sec:032337]{聖經人聖經事:唯有耶和華管理你們} & 2026-06-19 & \href{http://www.church.com.hk/acms/content.asp?site=cdc\&op=show\&type=product\&code=032337\&layout=sermon}{\texttt{032337}} \\
以弗所書 6:10-24 & 呂君望牧師 & \hyperref[sec:032338]{門徒當自強} & 2026-06-21 & \href{http://www.church.com.hk/acms/content.asp?site=cdc\&op=show\&type=product\&code=032338\&layout=sermon}{\texttt{032338}} \\
出埃及記 32:1-10 & 羅泰然牧師 & \hyperref[sec:032339]{偷天換日} & 2026-06-21 & \href{http://www.church.com.hk/acms/content.asp?site=cdc\&op=show\&type=product\&code=032339\&layout=sermon}{\texttt{032339}} \\
歌羅西書 1:23 & 招世超牧師 & \hyperref[sec:032340]{A.I.機器人襲地球} & 2026-06-21 & \href{http://www.church.com.hk/acms/content.asp?site=cdc\&op=show\&type=product\&code=032340\&layout=sermon}{\texttt{032340}} \\
羅馬書 4:17 & 葛國光牧師 & \hyperref[sec:032341]{模範的父親 (華)} & 2026-06-21 & \href{http://www.church.com.hk/acms/content.asp?site=cdc\&op=show\&type=product\&code=032341\&layout=sermon}{\texttt{032341}} \\
列王記上 3:3 & 李思聰牧師 & \hyperref[sec:032342]{從大衛家庭看父親對兒女的影響} & 2026-06-21 & \href{http://www.church.com.hk/acms/content.asp?site=cdc\&op=show\&type=product\&code=032342\&layout=sermon}{\texttt{032342}} \\
士師記 8:33 & 黃克勤牧師博士 & \hyperref[sec:032343]{聖經人聖經事:忘恩負義的以色列人} & 2026-06-20 & \href{http://www.church.com.hk/acms/content.asp?site=cdc\&op=show\&type=product\&code=032343\&layout=sermon}{\texttt{032343}} \\
路加福音 15:11-12 & 李慧珊傳道 & \hyperref[sec:032344]{一位父親的心} & 2026-06-21 & \href{http://www.church.com.hk/acms/content.asp?site=cdc\&op=show\&type=product\&code=032344\&layout=sermon}{\texttt{032344}} \\
路加福音 11:9-14 & 王安洋牧師 & \hyperref[sec:032345]{懂愛敢愛的爸爸} & 2026-06-21 & \href{http://www.church.com.hk/acms/content.asp?site=cdc\&op=show\&type=product\&code=032345\&layout=sermon}{\texttt{032345}} \\
士師記 8:1-35 & 梁永善牧師 & \hyperref[sec:032346]{士師記析讀 (10) 失敗之父} & 2026-06-24 & \href{http://www.church.com.hk/acms/content.asp?site=cdc\&op=show\&type=product\&code=032346\&layout=sermon}{\texttt{032346}} \\
列王記上 8:1-66 & 袁惠鈞牧師博士 & \hyperref[sec:032347]{萬國禱告的殿:悔改與復興的真正道路} & 2026-06-21 & \href{http://www.church.com.hk/acms/content.asp?site=cdc\&op=show\&type=product\&code=032347\&layout=sermon}{\texttt{032347}} \\
路加福音 15:11-32 & 蕭海波牧師博士 & \hyperref[sec:032348]{慈愛天父心:從被愛到愛人的呼召} & 2026-06-21 & \href{http://www.church.com.hk/acms/content.asp?site=cdc\&op=show\&type=product\&code=032348\&layout=sermon}{\texttt{032348}} \\
士師記 9:5 & 黃克勤牧師博士 & \hyperref[sec:032349]{聖經人聖經事:亞比米勒用計殺兄弟作王} & 2026-06-21 & \href{http://www.church.com.hk/acms/content.asp?site=cdc\&op=show\&type=product\&code=032349\&layout=sermon}{\texttt{032349}} \\
申命記 10:12-22 & 謝又生牧師 & \hyperref[sec:032350]{感恩 MODE ON (3) 喂,係時候「落場」喇!} & 2026-06-20 & \href{http://www.church.com.hk/acms/content.asp?site=cdc\&op=show\&type=product\&code=032350\&layout=sermon}{\texttt{032350}} \\
路加福音 12:13-34 & 李觀南傳道 & \hyperref[sec:032351]{跳出填倉人生} & 2026-06-21 & \href{http://www.church.com.hk/acms/content.asp?site=cdc\&op=show\&type=product\&code=032351\&layout=sermon}{\texttt{032351}} \\
路加福音 17:20-37 & 俞大鳴牧師 & \hyperref[sec:032352]{天國的來臨} & 2026-06-07 & \href{http://www.church.com.hk/acms/content.asp?site=cdc\&op=show\&type=product\&code=032352\&layout=sermon}{\texttt{032352}} \\
尼希米記 1:1-7 & 劉振惠弟兄 & \hyperref[sec:032353]{尼希米:以神的異象引領,以神的使命前行} & 2026-06-21 & \href{http://www.church.com.hk/acms/content.asp?site=cdc\&op=show\&type=product\&code=032353\&layout=sermon}{\texttt{032353}} \\
約翰福音 4:43-54 & 許開明牧師 & \hyperref[sec:032354]{慈父憑信見神蹟,得人間三寶} & 2026-06-21 & \href{http://www.church.com.hk/acms/content.asp?site=cdc\&op=show\&type=product\&code=032354\&layout=sermon}{\texttt{032354}} \\
士師記 11:30-31 & 黃克勤牧師博士 & \hyperref[sec:032355]{聖經人聖經事:耶弗他以他女兒許願} & 2026-06-27 & \href{http://www.church.com.hk/acms/content.asp?site=cdc\&op=show\&type=product\&code=032355\&layout=sermon}{\texttt{032355}} \\
約書亞記 24:14-28 & 梁成裕牧師 & \hyperref[sec:032356]{至於我和我家,我們必定事奉耶和華} & 2026-06-21 & \href{http://www.church.com.hk/acms/content.asp?site=cdc\&op=show\&type=product\&code=032356\&layout=sermon}{\texttt{032356}} \\
約翰福音 1:37-51 & 羅婉明傳道 & \hyperref[sec:032357]{蒙主呼召,更大的事} & 2026-06-21 & \href{http://www.church.com.hk/acms/content.asp?site=cdc\&op=show\&type=product\&code=032357\&layout=sermon}{\texttt{032357}} \\
馬太福音 11:25-30 & 陳雅莉傳道 & \hyperref[sec:032431]{主裡得享安息} & 2026-07-12 & \href{http://www.church.com.hk/acms/content.asp?site=cdc\&op=show\&type=product\&code=032431\&layout=sermon}{\texttt{032431}} \\
俄巴底亞書 1:10 & 黃紹權牧師 & \hyperref[sec:032432]{自作孽,滿招損} & 2026-07-12 & \href{http://www.church.com.hk/acms/content.asp?site=cdc\&op=show\&type=product\&code=032432\&layout=sermon}{\texttt{032432}} \\
申命記 2:2-25 & 禤凱榮牧師 & \hyperref[sec:032433]{人生的冤枉路} & 2026-07-12 & \href{http://www.church.com.hk/acms/content.asp?site=cdc\&op=show\&type=product\&code=032433\&layout=sermon}{\texttt{032433}} \\
士師記 16:19 & 黃克勤牧師博士 & \hyperref[sec:032434]{聖經人聖經事:耶和華已經離開了} & 2026-07-09 & \href{http://www.church.com.hk/acms/content.asp?site=cdc\&op=show\&type=product\&code=032434\&layout=sermon}{\texttt{032434}} \\
羅馬書 8:1-11 & 林玉英牧師 & \hyperref[sec:032435]{有靠山,活得帥!} & 2026-07-12 & \href{http://www.church.com.hk/acms/content.asp?site=cdc\&op=show\&type=product\&code=032435\&layout=sermon}{\texttt{032435}} \\
創世記 25:19-34 & 李均熊牧師博士 & \hyperref[sec:032436]{我都就死,這於我有甚麼益處?} & 2026-07-12 & \href{http://www.church.com.hk/acms/content.asp?site=cdc\&op=show\&type=product\&code=032436\&layout=sermon}{\texttt{032436}} \\
詩篇 119:49-56 & 陸超明牧師博士 & \hyperref[sec:032437]{靈命成長筆記 (2) 回歸永活的聖言} & 2026-07-12 & \href{http://www.church.com.hk/acms/content.asp?site=cdc\&op=show\&type=product\&code=032437\&layout=sermon}{\texttt{032437}} \\
啟示錄 6:1-8 & 譚綺敏傳道 & \hyperref[sec:032438]{審判的七印——逼近的災難} & 2026-07-12 & \href{http://www.church.com.hk/acms/content.asp?site=cdc\&op=show\&type=product\&code=032438\&layout=sermon}{\texttt{032438}} \\
列王記上 11:1-43 & 袁惠鈞牧師博士 & \hyperref[sec:032439]{愚昧的智者:心偏離神一寸,人生偏離神一里} & 2026-07-12 & \href{http://www.church.com.hk/acms/content.asp?site=cdc\&op=show\&type=product\&code=032439\&layout=sermon}{\texttt{032439}} \\
士師記 16:26 & 黃克勤牧師博士 & \hyperref[sec:032440]{聖經人聖經事:我要靠一靠} & 2026-07-10 & \href{http://www.church.com.hk/acms/content.asp?site=cdc\&op=show\&type=product\&code=032440\&layout=sermon}{\texttt{032440}} \\
以弗所書 2:14 & 陳淳牧師 & \hyperref[sec:032441]{拆毀與更新} & 2026-07-12 & \href{http://www.church.com.hk/acms/content.asp?site=cdc\&op=show\&type=product\&code=032441\&layout=sermon}{\texttt{032441}} \\
腓立比書 4:1-9 & 陸幸泉牧師 & \hyperref[sec:032442]{靠主與同心} & 2026-07-12 & \href{http://www.church.com.hk/acms/content.asp?site=cdc\&op=show\&type=product\&code=032442\&layout=sermon}{\texttt{032442}} \\
瑪拉基書 4:4-6 & 石偉漢牧師 & \hyperref[sec:032443]{主角出場,君王降臨} & 2026-07-12 & \href{http://www.church.com.hk/acms/content.asp?site=cdc\&op=show\&type=product\&code=032443\&layout=sermon}{\texttt{032443}} \\
路加福音 2:1-7 & 羅杰才牧師 & \hyperref[sec:032444]{耶穌的出生與一生} & 2025-12-25 & \href{http://www.church.com.hk/acms/content.asp?site=cdc\&op=show\&type=product\&code=032444\&layout=sermon}{\texttt{032444}} \\
腓立比書 2:1-11 & 余大峰傳道 & \hyperref[sec:032445]{與基督同一心思心志的合一表現} & 2026-07-12 & \href{http://www.church.com.hk/acms/content.asp?site=cdc\&op=show\&type=product\&code=032445\&layout=sermon}{\texttt{032445}} \\
士師記 16:30 & 黃克勤牧師博士 & \hyperref[sec:032446]{聖經人聖經事:死時殺的人比活著時更多} & 2026-07-11 & \href{http://www.church.com.hk/acms/content.asp?site=cdc\&op=show\&type=product\&code=032446\&layout=sermon}{\texttt{032446}} \\
約翰一書 1:1-4 & 劉燕飛教師 & \hyperref[sec:032447]{生命之道,團契生活} & 2026-07-12 & \href{http://www.church.com.hk/acms/content.asp?site=cdc\&op=show\&type=product\&code=032447\&layout=sermon}{\texttt{032447}} \\
創世記 4:1-7 & 林羅麗斯牧師 & \hyperref[sec:032448]{討主喜悅} & 2026-07-12 & \href{http://www.church.com.hk/acms/content.asp?site=cdc\&op=show\&type=product\&code=032448\&layout=sermon}{\texttt{032448}} \\
哥林多前書 10:23-33 & 梁永善牧師 & \hyperref[sec:032449]{得自由,失自由?} & 2026-07-12 & \href{http://www.church.com.hk/acms/content.asp?site=cdc\&op=show\&type=product\&code=032449\&layout=sermon}{\texttt{032449}} \\
以西結書 47:9 & 劉彤牧師 & \hyperref[sec:032450]{得著腳掌所踏之地 (華)} & 2026-07-12 & \href{http://www.church.com.hk/acms/content.asp?site=cdc\&op=show\&type=product\&code=032450\&layout=sermon}{\texttt{032450}} \\
哥林多前書 11:2-16 & 陳崇基牧師 & \hyperref[sec:032451]{再思女性的事奉 (2)} & 2026-07-12 & \href{http://www.church.com.hk/acms/content.asp?site=cdc\&op=show\&type=product\&code=032451\&layout=sermon}{\texttt{032451}} \\
士師記 17:6 & 黃克勤牧師博士 & \hyperref[sec:032452]{聖經人聖經事:沒有王,各人任意而行 (1)} & 2026-07-12 & \href{http://www.church.com.hk/acms/content.asp?site=cdc\&op=show\&type=product\&code=032452\&layout=sermon}{\texttt{032452}} \\
羅馬書 10:4-21 & 周曉暉牧師 & \hyperref[sec:032453]{宣認耶穌是主} & 2025-12-14 & \href{http://www.church.com.hk/acms/content.asp?site=cdc\&op=show\&type=product\&code=032453\&layout=sermon}{\texttt{032453}} \\
約翰福音 8:48-59 & 林鍵牧師 & \hyperref[sec:032454]{榮耀君王:我是救主} & 2026-07-11 & \href{http://www.church.com.hk/acms/content.asp?site=cdc\&op=show\&type=product\&code=032454\&layout=sermon}{\texttt{032454}} \\
加拉太書 5:13-15 & 葛國光牧師 & \hyperref[sec:032455]{自由的真諦 (華)} & 2026-07-12 & \href{http://www.church.com.hk/acms/content.asp?site=cdc\&op=show\&type=product\&code=032455\&layout=sermon}{\texttt{032455}} \\
撒母耳記下 2:1-32 & 李思聰牧師 & \hyperref[sec:032456]{毋讓玉帛變干戈} & 2026-07-12 & \href{http://www.church.com.hk/acms/content.asp?site=cdc\&op=show\&type=product\&code=032456\&layout=sermon}{\texttt{032456}} \\
馬可福音 3:13-19 & 冼嘉琪牧師 & \hyperref[sec:032457]{耶穌為何呼召你} & 2026-07-12 & \href{http://www.church.com.hk/acms/content.asp?site=cdc\&op=show\&type=product\&code=032457\&layout=sermon}{\texttt{032457}} \\
士師記 18:9 & 黃克勤牧師博士 & \hyperref[sec:032458]{聖經人聖經事:沒有王,各人任意而行 (2) 但人以神之名謀劃殺人奪地} & 2026-07-13 & \href{http://www.church.com.hk/acms/content.asp?site=cdc\&op=show\&type=product\&code=032458\&layout=sermon}{\texttt{032458}} \\
馬可福音 1:1-8 & 梁柏堅弟兄 & \hyperref[sec:032459]{為福音,好好預備自己} & 2026-07-11 & \href{http://www.church.com.hk/acms/content.asp?site=cdc\&op=show\&type=product\&code=032459\&layout=sermon}{\texttt{032459}} \\
啟示錄 2:1-7 & 劉沙崙牧師 & \hyperref[sec:032460]{教會如何在黑暗中作光} & 2026-07-12 & \href{http://www.church.com.hk/acms/content.asp?site=cdc\&op=show\&type=product\&code=032460\&layout=sermon}{\texttt{032460}} \\
希伯來書 2:5-18 & 張健庭牧師博士 & \hyperref[sec:032461]{希伯來書系列 (3) 救贖的元帥} & 2026-07-12 & \href{http://www.church.com.hk/acms/content.asp?site=cdc\&op=show\&type=product\&code=032461\&layout=sermon}{\texttt{032461}} \\
箴言 12:27 & 曹子光牧師 & \hyperref[sec:032462]{懶到出汁} & 2026-07-12 & \href{http://www.church.com.hk/acms/content.asp?site=cdc\&op=show\&type=product\&code=032462\&layout=sermon}{\texttt{032462}} \\
利未記 23:39-43 & 陳崇基牧師 & \hyperref[sec:032463]{住棚} & 2026-06-28 & \href{http://www.church.com.hk/acms/content.asp?site=cdc\&op=show\&type=product\&code=032463\&layout=sermon}{\texttt{032463}} \\
士師記 18:17 & 黃克勤牧師博士 & \hyperref[sec:032464]{聖經人聖經事:沒有王,各人任意而行 (3) 六百但人搶劫神像,挾持祭司} & 2026-07-14 & \href{http://www.church.com.hk/acms/content.asp?site=cdc\&op=show\&type=product\&code=032464\&layout=sermon}{\texttt{032464}} \\
創世記 30:25-31:16 & 林天行牧師 & \hyperref[sec:032465]{參與神的作為} & 2026-07-12 & \href{http://www.church.com.hk/acms/content.asp?site=cdc\&op=show\&type=product\&code=032465\&layout=sermon}{\texttt{032465}} \\
路加福音 21:25-38 & 陳傳華牧師博士 & \hyperref[sec:032466]{AI,UAP與末世:僥倖或挺身昂首} & 2026-07-12 & \href{http://www.church.com.hk/acms/content.asp?site=cdc\&op=show\&type=product\&code=032466\&layout=sermon}{\texttt{032466}} \\
民數記 15:1-41 & 林誠信牧師 & \hyperref[sec:032467]{審判,恩典,敬拜} & 2026-07-12 & \href{http://www.church.com.hk/acms/content.asp?site=cdc\&op=show\&type=product\&code=032467\&layout=sermon}{\texttt{032467}} \\
路加福音 19:11-27 & 俞大鳴牧師 & \hyperref[sec:032468]{天國人的使命} & 2026-07-05 & \href{http://www.church.com.hk/acms/content.asp?site=cdc\&op=show\&type=product\&code=032468\&layout=sermon}{\texttt{032468}} \\
申命記 2:26-3:11 & 禤凱榮牧師 & \hyperref[sec:032469]{戰勝巨人三步曲} & 2026-07-19 & \href{http://www.church.com.hk/acms/content.asp?site=cdc\&op=show\&type=product\&code=032469\&layout=sermon}{\texttt{032469}} \\
羅馬書 1:20-23 & 梁錫強醫生 & \hyperref[sec:032470]{誰是我「王」} & 2026-07-19 & \href{http://www.church.com.hk/acms/content.asp?site=cdc\&op=show\&type=product\&code=032470\&layout=sermon}{\texttt{032470}} \\
士師記 18:31 & 黃克勤牧師博士 & \hyperref[sec:032471]{聖經人聖經事:沒有王,各人任意而行 (4) 六百但人屠城,奪地,私設神堂} & 2026-07-15 & \href{http://www.church.com.hk/acms/content.asp?site=cdc\&op=show\&type=product\&code=032471\&layout=sermon}{\texttt{032471}} \\
以賽亞書 44:6-8 & 林家棋宣教師 & \hyperref[sec:032472]{真實與虛無} & 2026-07-19 & \href{http://www.church.com.hk/acms/content.asp?site=cdc\&op=show\&type=product\&code=032472\&layout=sermon}{\texttt{032472}} \\
那鴻書 1:2-3 & 蕭壽華牧師 & \hyperref[sec:032473]{榮耀的釋放} & 2025-12-21 & \href{http://www.church.com.hk/acms/content.asp?site=cdc\&op=show\&type=product\&code=032473\&layout=sermon}{\texttt{032473}} \\
使徒行傳 16:6-10 & 鄧瑞強牧師博士 & \hyperref[sec:032474]{直搗羅馬} & 2026-07-19 & \href{http://www.church.com.hk/acms/content.asp?site=cdc\&op=show\&type=product\&code=032474\&layout=sermon}{\texttt{032474}} \\
\end{xltabular}
}
\newpage



\section{傳道書 12:1-8}
\label{sec:031601}
\textbf{不要等}~~~~官志達牧師
\newline
\newline
連結: \href{http://www.church.com.hk/acms/content.asp?site=cdc\&op=show\&type=product\&code=031601\&layout=sermon}{\texttt{031601}} ~~~~ 語音日期: 2026-02-01
\newline
\newline
\hyperref[sec:032474]{\small{< < < PREV SERMON < < <}}
~
\hyperref[sec:index_chronic]{\small{[返順時目]}}
~
\hyperref[sec:031602]{\small{> > > NEXT SERMON > > >}}
\newline
\newline
傳道書 12:1-8
\newline
\begin{longtable}{cl}
\hline
\hline
章節 & 經文 (和合本修訂版)\\
\hline
12:1 & \begin{tabularx}{0.7\textwidth}{X} 你趁著年輕、衰老的日子尚未來到，就是你所說，我毫無喜悅的那些歲月來臨之前，當記念造你的主。 \end{tabularx} \\ \\ \relax
12:2 & \begin{tabularx}{0.7\textwidth}{X} 不要等到太陽、光明、月亮、星宿變為黑暗，雨後雲又返回； \end{tabularx} \\ \\ \relax
12:3 & \begin{tabularx}{0.7\textwidth}{X} 看守房屋的發顫，強壯的屈身，推磨的婦女因人少而停工，從窗戶往外看的眼光變為昏暗； \end{tabularx} \\ \\ \relax
12:4 & \begin{tabularx}{0.7\textwidth}{X} 街門關閉，推磨的聲音微小，鳥一叫，就驚醒，唱歌女子的聲音也都微弱； \end{tabularx} \\ \\ \relax
12:5 & \begin{tabularx}{0.7\textwidth}{X} 人怕高處，路上有驚慌；杏樹開花，蚱蜢成為重擔，慾望不再挑起；因為人歸他永遠的家，弔喪的在街上往來。 \end{tabularx} \\ \\ \relax
12:6 & \begin{tabularx}{0.7\textwidth}{X} 不要等到銀鏈折斷，金罐破裂，瓶子在泉旁損壞，水輪在井口斷裂， \end{tabularx} \\ \\ \relax
12:7 & \begin{tabularx}{0.7\textwidth}{X} 塵土仍歸於地，像原來一樣，氣息仍歸於賜氣息的神。 \end{tabularx} \\ \\ \relax
12:8 & \begin{tabularx}{0.7\textwidth}{X} 傳道者說：「虛空的虛空，全是虛空。」 \end{tabularx} \\ \\
[1ex]
\hline
\hline
\end{longtable}
$^{1}$今天大家看的題目叫做「不要等」.
我相信在香港生活的人.
這些題目其實不用說的了.
大家都知道.
中國人所說的.
「執書行頭」是怎樣的呢?.
慘過敗家.
什麼都要搶,抓,等.
不要過去等花了任何時間.
沒什麼機會了.
但是這個「不要等」這個情況.
通常用在我們自己生命當中.
跟自己有關,有好處的情況之下.
我們才會很急切地追趕.
趕著做完一件事.
不過對於在屬靈的事情上.
或者對神的事情.
我們往往有一個「拖延症」.
神的事情.
遲一點做吧.
不用那麼急.
慢慢做.
不需要那麼緊張.
或者我們在神的事情裡面.
時刻未到.
我們先禱告等候神.
用一個比較拖延的態度去回應神.
祂自己對我們心裡面的呼召和感動.
回看《聖經》裡面的時候.
又不是.
《聖經》有些比較舊的詩歌.
譬如唐從明那首《雅歌》.
第一句是什麼呢?.
「你吸引我,我就慢跑跟隨你」.
是「快跑跟隨你」.
在《雅歌》第一章第四節.
它形容到.
當一件神那麼吸引祂的時候.
祂真的能夠撲過去.
我們曾幾何時對神那種愛慕之情.

$^{41}$仰慕之心有那麼激動呢?.
如果大家是一些粵劇歌迷的話.
你就明白我說什麼了.
哪個大腦官唱完了.
我們走到後台裡面去看看他.
然後跟他握握手.
或者那些粉絲.
握握手.
很開心見到他.
究竟對神的態度.
有沒有那種渴慕,追求,熱切的心呢?.
我們是慢跑還是超慢跑去跟隨神呢?.
中學階段或者小學階段.
都聽過一首詩.
「勸君莫惜金柳衣」.
「勸君惜取少年時」.
大家明白了.
鼓勵年輕人不要想著穿漂亮的衣服.
不需要穿名牌.
要趁著年輕的時候.
要好好地讀書.
現在的小朋友都不用我說了.
他小時候.
父母兩歲已經叫他讀書了.
很慘.
但是這首詩提醒.
要把握機會.
不要等到花開很盛放的時候.
應該剪下來做插花的.
今天好像沒有插花.
可能花凋謝了.
「莫待無花空」.
不要等花凋謝的時候才剪下來.
要把握時間,把握機會.
鼓勵年輕人去把握這個時機.
剪轉剪轉這首歌唱了幾十年了.
我唱了差不多有50年.
我已經拿了路由卡了.
要唱另一首歌了.
不要唱這首歌那麼珍惜機會了.

$^{81}$時間無多剩了.
「老人之歌」.
不知道大家聽過沒有.
「一晚通宵唔捱得」.
「兩腳無力唔行得」.
「三粒藥丸唔少得」.
「四次夜尿又唔忍得」.
「五更醒jor 唔訓得」.
「六種骨痛唔識得」.
「七樣凍品唔受得」.
「八卦小器唔激得」.
「九屎垃圾唔捨得」.
「十分鐘後唔記得」.
記得嗎?.
如果不記得就馬上拍一張照.
我想給大家最主要的.
如果我是年輕人.
家裡有長輩.
就想一下他們的情況.
我們要關心老人家,長者.
他們有這樣的情況.
如果我們是長者的時候.
這些更加是會心微笑.
不要緊,大家是有份的.
理不孤單的.
總有一件事在你身上發生.
我們.
先不想太興奮.
在這樣的情況之下.
怎麼辦呢?.
怎麼看我們餘下的人生.
或者怎麼看我們有的生命呢?.
看回《聖經》的一些人物.
雅各很出名.
但他的生命裡充滿了很多軟弱.
騙爸爸騙媽媽.
被人騙,被兒子騙.
臨老的時候.
他的孫子被人說不見了.
最晚年的時候.

$^{121}$大饑荒.
縱然他有錢也好.
大饑荒要走去埃及擲米擲糧.
見到法老王的時候.
他就說.
「我就是世紀居日子一百三十年」.
「我一生的年月」.
「又短又苦」.
是個有錢人.
他也說「又短又苦」.
可能是向法老炸死了.
「比不上我祖先在世紀居日子」.
很快的.
「又短又苦」.
摩西他自己的這首詩裡面.
也有這樣的說法.
「我們其中所驚誇的」.
「可以牙刷的是什麼呢?」.
「奴婦受煩轉眼即逝」.
「如便飛而去」.
一下子就過去了.
我記得我第一次在這裡分享的時候.
應該是二十多年前了.
不知道大家當中有沒有記得.
印象我當時很年輕的一個小智.
也很年輕的.
轉過了二十多年的時候.
真的很快就過去了.
未來不知道有沒有機會再過二十年.
如果是這樣的話.
我怎樣去把握我今天的日子呢?.
人生是不是這麼苦這麼短呢?.
我想到最近在香港那場大火裡面.
原來人生的終局.
不是一些老人家的事情.
可能很年輕的人.
隨時一刻都可能灰飛煙滅.
到當人生裡面最後一刻的時候.
我們有沒有一些遺憾呢?.
有時候那些危機管理就覺得.

$^{161}$如果我今天只剩下二十四小時的生命.
我有沒有一些遺憾的事情.
我還沒做.
我要趁著這二十四個小時儘快做了.
我們有沒有這樣的驚覺呢?.
當保羅自己看他人生的時候.
他就不同了.
他說我離世的時候到了.
在他離世之前兩三年就寫了這本提摩泰的後書.
給他的徒弟.
他說美好的仗我已經打過了.
當跑的路我已經跑盡了.
我該信的道已經守住了.
他毫無後悔的勇往直前的面對主面.
這是不是我們期盼的心態去過我們的日子呢?.
我們不知道有多少個時日.
不過此時此刻我們怎樣看待我們的生命呢?.
我就引用了今天所讀的經文.
在《傳道書》第十二章的經文.
彼此分享和鼓勵.
這段經文已經看過了.
有些明白有些不明白.
他說你趁著年輕衰老的日子尚未來到.
當紀念你的主.
這段經文提醒一些年輕人.
或者未曾衰老的人.
今天我們怎樣看待我們的生命呢?.
你覺不覺得你年輕呢?.
年輕不代表你是一個年紀小.
年輕是指在心境上.
例如現在的青少年.
老氣橫秋.
老職之外.
人生好像沒有什麼盼望.
捱到比長者還要慘.
年輕你怎樣給自己定義呢?.
你是不是一個年輕的人呢?.
不是在講心智上.
不是在講年紀上.
究竟我們怎樣看待我們今天的生命呢?.

$^{201}$無論你從一個年紀或者從一個心境看的時候.
無論什麼情況都好.
當紀念做你的主.
想想我們是由上帝創造的.
我們有一天要回到上帝的時候.
好像有一個少年要走到老年的時候.
這段階段我們可以做些什麼呢?.
施人華提醒我們要紀念做我們的主.
在這段經文裡面.
很實際講到究竟.
真是講年紀老大的情況.
看看和我們剛才老人之歌有什麼相似的地方.
他說.
衰敗的只是什麼呢?.
看守房屋的發展.
就是古代.
如果你想著幾千年前的時候.
拿著支棍去守住房屋的發展代表什麼呢?.
手開始有震動.
用現代的醫學去可能有什麼病呢?.
柏金遜.
拿著支棍都不可以很鎮靜.
看守房屋的發展.
強壯的鬱身.
以前是挺起腰板走路的.
就是趾高氣揚的.
老了.
為什麼彎下身呢?.
是什麼呢?.
現代醫學的其中解釋.
骨質疏鬆.
你看到老人家是彎起身去走路的.
就是不能再挺起腰去走路了.
身體已經彎曲了.
推磨的婦女因人少而停工.
大家乘巴士.
等巴士.
經常看到巴士廣告.
如果想去深圳購物.
其中一個廣告是去深圳的.

$^{241}$是什麼呢?.
按摩.
沒看過這個廣告.
整容.
不是整容是美顏.
是一種整牙.
搞到很多牙醫牙痛的聲音.
整隻牙幾萬元到幾千元就搞定了.
推磨牙是磨東西的.
年紀老大就開始脫牙.
大假牙.
中牙.
牙少就開始停工.
第四個情況.
從窗戶往外看的眼光昏暗.
超過四十歲就開始明白多些.
未到四十歲就感受少些.
開始老化.
六十歲有什麼呢?.
白內障.
再差些是青光眼.
或者黃斑病變.
大家很有經驗.
看東西會蒙起眼睛.
不過不要緊.
弄完白內障.
有些三度變焦的鏡頭.
放在眼睛都可以將視力改變.
看東西開始蒙蒙查查.
不是單是身體蒙查.
可能心靈上也蒙蒙查查.
對世界事物已經不認識.
看東西不清不楚.
這是另一個討論.
門口關閉推磨的聲音微小.
跟第三個有關係.
嘴巴不想合起來.
吃東西的份量少了.
不要請老人去吃自助餐.
虧本.

$^{281}$他吃不到太多東西.
已經沒有胃口.
吃東西已經有限制.
胃已經不能再消化很多.
門口關閉的口.
推磨的聲音微小.
不會吃太多東西.
尿一叫就驚醒.
大家住在城市比較少.
新界五六點就聽到鳥叫.
然後就起床.
醒了就睡不著.
這是講長者的情況.
或是衰敗日子的情況.
唱歌女子的聲音也都微弱.
姐妹年輕時唱頌士班很大聲.
年紀大了想唱都唱不出.
聲帶已經沉了.
不應該大聲唱.
罵人都小聲一點.
人怕高處路上有驚慌.
是甚麼.
所以有些長者很怕出街.
為甚麼.
怕高處有驚慌.
逛街隨時被人撞.
年輕時被人撞你撞我撞你無所謂.
老人家被人撞會怎樣.
跌倒跌倒會怎樣.
骨折會怎樣.
行行得不行得.
年輕人不行得無所謂.
但老人家跌倒不行得.
睡床要睡幾個月.
如果再處理不好睡床.
長期睡床有甚麼困難.
壓力瘡.
如果身體再差有肺炎.
跟住就再見.
所以老人家很怕跌.

$^{321}$所以年輕人記住拖住老人家出街.
行樹開花.
白色頭殼.
頭髮斑白.
馬渣成為重擔.
拿些東西.
細細隻很難拿.
拿不到東西.
慾望不再挑起.
沒心機沒衝動去做任何事.
在講衰敗的日子.
跟住下來.
因為差不多有結論.
一切歸於原位.
人歸他永遠的家.
吊喪的在街上往來.
我們很明白.
做年紀大時.
經常如果在紅磡出入.
去完萬國就去世界.
去完世界就去萬國.
吊喪的在街上往來.
是講什麼.
出席很多喪禮.
教會也有些喪禮.
去到他告訴大家.
有些這樣的情況出現.
人去行到人生的結局.
終結了.
銀鏈截斷.
金冠破裂.
全體在旁邊損壞.
一切的東西都歸於原處了.
水輪在井口斷裂.
塵土亦不歸於地.
過完了.
少年老年最終就是.
去到人生最後的終局.
在這樣的情況.
詩人提醒我們.

$^{361}$未到這樣的情況的時候.
去紀念做我們的主.
我們從哪裡來的時候.
也從哪裡回去.
在這個人生過程裡面的時候.
我們要紀念神.
在我們生命當中.
做過一切一切的事情.
當我們今天還沒去到.
人生最後一刻的時候.
我們記住神什麼呢.
神有記住我們什麼呢.
紀念主的意念是一個很重要的.
對信徒來說很大的提醒.
這裡說了.
紀念我們出於塵土.
而歸於塵土.
神在聖經裡面很多次說.
神要記住羅亞的約.
說到人類犯罪.
但神要用洪水滅了世界.
亦都用了一個彩虹和人立約.
見到彩虹之後.
不會再用洪水毀滅這個世界.
人是有軟弱.
不過神有恩典提醒我們.
神的恩典是不會因為我們犯罪而去抹殺.
軟弱我們迴轉之後.
神的約仍然有效.
亞伯拉罕以撒亞國的約.
亦都是在說這件事.
以色列人祖先.
他做了很多的事情.
好的壞的神都知道.
唯有我們怎樣去看我們的神呢.
所以在詩篇有一個回應.
無論如何不要忘記神的恩惠.
在人的困難當中.
我們不會忘記神的恩惠.
但當人的豐富的時候.

$^{401}$環境安定的時候.
我們有沒有忘記神的恩惠呢.
會不會忘記了神對於我們生命裡面的提醒.
或者忘記了神託付我們生命的任務要完成呢.
才是說不要忘記神的恩惠.
神的使命.
詩篇137篇差不多是一個.
講到以色列人被擄巴比倫回歸之後.
他們看見神在他們身上的作為的時候.
提醒他們.
今天的日子.
怎樣去過他敬畏神的生活.
剛才說了.
這個是私人對於生命的體會.
我也嘗試投入一下.
回顧我過去二十多年的時候.
究竟紀念著神什麼事情呢.
神又紀念著我什麼事情呢.
這是我過往二十五年裡.
在一些事奉裡的日子的經歷.
有三個問題.
我在這二十多年裡.
在海外事奉裡.
有幾個問題經常營繞著我.
今天我們怎樣看呢.
嘗試和大家一起分享.
第一個問題是.
我經常說我們是關愛使者.
或者在普遍教會.
猜拳的宣教.
往往有一個題目.
令到我有些迷惑.
覺得宣教士很偉大.
是不是.
是喔.
有些宣教士.
如果說回基督教的歷史來說.
那些人從英國去到中國.
在中國很艱難的學習.
大家去澳門看馬利遜的墓.

$^{441}$就有這樣的記載.
做了很多偉大的事情.
他一個人很偉大.
做了很多偉大的事情.
回到英華寺殿建了一個學校.
影響很深.
但是否所有的宣教士都這麼偉大呢.
不是喔.
但很多時候我們將這些人物.
捧得很高.
做宣教工作很偉大.
如果有這樣的想法.
其實難了我們很多弟兄姊妹.
參與宣教的工作.
這些偉大的事情.
當然留給偉大的人去做.
牧師傳道你做對了.
你這麼偉大.
我小小的信徒.
我不能參與這些宣教.
參與這些猜拳工作.
以致我們整間教會裡面.
有小撮人有心.
而不是整體一起去做這個工作.
因為這些這麼偉大的工作.
留給偉大的人去做.
我做很簡單的卑微工作就算了.
有些弟兄姊妹願意去突破.
參與一些宣教體驗.
一年兩年.
但因為這個光環的緣故.
當他生命裡面在工場裡面.
遇到一些挑戰困難.
我怎麼可能從一個宣教工場退下來.
你會覺得很軟弱.
很多的問題.
人們怎麼看我.
於是他就有一些情況.
去完這個工場一段的短時間.
不適應.

$^{481}$再去第二個工場又不適應.
去第三個工場又不適應.
結果我就稱為長期的短宣.
又好像變成一個長期宣教.
但其實這個短暫短暫短暫加起來.
其實就不是一個長期.
我們做猜拳工作.
也見到很多弟兄姊妹.
在這個工場裡面很大的掙扎.
對當地的教會.
對當地的文化不認識.
往往用了自己一套既有的觀念.
放在當地當中.
我第三點就說到.
對文化的不認識不掌握.
以致碰了很多的壁.
這個是我們自己在經歷裡面的時候.
講猜拳講宣教的時候.
一定要和大家弟兄姊妹講.
我們不可以輕視在工場裡面的文化差異.
用我自己一個簡單在這幾點裡面的體會.
其實每一個基督徒的跟隨者.
每一個耶穌基督的跟隨者.
就是我們.
是用道成肉身的一個生活態度.
去認識我們的文化.
認識當地的文化.
以致我們能夠將我們的信仰.
透過文化的貫通.
而將我們的信息.
還一個位置.
以致能夠完成基督託付的任務.
任務的意思就是說.
我們是從神領受一個工作體會.
希望我們能夠.
藉著我們完成一些託付.
這件事我們覺得將宣教變成很偉大.
其實我看如果是一個簡單的定義來看.
一個任務一個宣教的mission.
mission應該翻譯做任務.

$^{521}$如果是一個微信的人會這樣翻譯.
就是說神給我們生命裡面.
有一個任務要做的.
你做完就行了.
因為宣教這個字.
或者這個mission這個字.
基本上是講猜拳.
當講猜拳的時候就是在講.
你帶著一個任務.
去到另一個地方去完成這個任務.
基本上就是猜拳.
猜險這個聖經的意思.
而宣教或者猜拳.
根本不是聖經出現過的.
沒有聖經出現過的.
聖經只是在強調一個猜險.
上帝猜險我們完成這個任務.
的一個這樣的動作.
每一個人都在神面前有一個任務.
我們有任務 你有任務.
你的任務可能是一個.
在你的家庭當中.
可能在你的職場當中.
可能在你的事奉環境當中.
在你的生活當中.
神有一樣東西.
有一樣東西給你做的.
很可能你的專業.
很可能你的專業再加上.
你自己一些的體會.
神的呼召.
你就在不同的地方裡面.
完成你的任務.
你的任務是什麼.
其實不是很偉大的.
除了我們在宗教裡面.
講這個宣教的事.
推到這麼高之外.
其實在我們身邊裡面有很多人.
他未信耶穌也好.

$^{561}$他都有很清楚的任務.
有時看見一些運動員.
他怎樣要堅持.
他要做成這件事.
有時看見一些藝術工作者.
他沒有政府的資助.
他願意知道這件事.
是好的時候.
他一年兩年三年四年十年的.
投放他的生命.
在他覺得需要完成的事情當中.
其實不要把宗教.
性化在一個超然的行動.
上帝給我們每一個人.
在生命當中有一樣東西要完成.
你清不清楚.
這就是神託付的差遣.
如果進入教會的背景來說.
我們說福音.
我們說傳道.
我在這幾年當中有這樣的體會.
究竟傳福音還是傳教會文化給當地人呢.
兩者之間有什麼分別呢.
很多時候福音.
教會文化就是我們現有的教會模式.
有祈禱會 有崇拜 有主日學 有探訪.
一切基於一個環境中有這樣的活動.
於是我們很容易說.
我們傳福音.
將這一套東西.
抄到另一個地區.
譬如說.
傳福音在香港會學什麼呢.
三福四律.
生命橋.
於是將這些東西.
告訴我們的工場.
但發覺.
不通.
因為我們.

$^{601}$有一個文化的差異.
我們只是將我們在教會裡面的.
信仰文化的模式.
抄到最初.
舉個例子.
譬如說.
我們在香港的生活裡面.
一個星期七天.
很穩定.
但對於無論在英國或在澳門.
那些人的生活模式.
跟香港是很不同的.
我們要因應他們的生活模式.
習慣而有一個手術調動.
簡單舉個例子.
譬如說我們去到英國裡面.
他們去做餐館的.
即是那些人士.
要晚上九點多十點才放工.
在十一點才放工.
我們要做什麼呢.
就要去到晚上十二點多.
用茶經班.
但這些情況已經少了.
我們要轉化.
我們希望能夠幫助到聖經.
於是我們就用茶經班.
但發覺.
這樣的模式不行.
因為他們聚集的時間夠困難.
他們住的地方很遠.
他們回家要很艱難.
我們要用一個方法.
我們目的是想他們明白聖經.
我們目的不是想他們回到茶經班.
於是就要改變策略.
改變策略.
他們不能來到這個教會.
或是這個集合地方嗎.
我們就去.

$^{641}$目的是想他們帶到神的面前.
他們要認識聖經.
而不是認識茶經班.
我講類似的教會文化.
我們有時會有一個迷思.
把教會文化.
變成一個福音.
告訴他們.
教會.
或是福音.
令人幸福和成功.
沒錯.
福音的其中一個結果.
是令人生活豐盛.
但在經驗裡.
講到.
當我們講福音的時候.
他信了耶穌之後.
反而有些情況.
沒有我們想像中理想.
可能信了耶穌之後.
生活更加艱難更加困苦.
他有沒有信錯耶穌呢.
其實最大的關鍵.
我們很容易把.
信耶穌簡單化.
基本上是漏了一個字.
是信靠.
耶穌過日子.
這個是我們.
在不同工場裡面的體驗.
我們目的是想.
那個人能夠和他.
幫助到這個人.
能夠和神建立這個關係.
而不是讓他帶來一個.
教會文化.
幫助他在無論何時何地.
困苦.
貧窮都要信靠他.

$^{681}$好像在.
分盟當中.
一生一世無論任何的.
境況都要信靠他.
以他為滿足.
這個和我們第二點.
有關係.
因為很多時候我們信了耶穌之後就想著.
耶穌會令到我.
凡事順利.
一帆風順.
有困難的時候我呼求神.
多謝神.
我做我自己的事.
福音是在說.
我們和神之間的那種關係.
我們要迴轉和神的.
生命結亂.
第三個是一樣的.
猜傳福音和教會.
是一脈相承.
當自己看見.
在海外的教會.
狀態的時候.
變成了一些同鄉會.
或者屬靈俱樂部.
在英國的情況更加明顯.
很多香港人.
無論是信主還是不信主.
都會回到教會.
所以在教會當中.
充滿了很多人回教會有不同的目的.
在教會當中.
會強調「我」很重要.
我們用英文「me church」.
我回教會是為了「我」.
我要好的音樂,要好的教導.
我要好的關心,我要好的栽培.
但教會是屬於神.
這些概念上.

$^{721}$我們有時會歪了一點.
我們重新再看看.
這是我自己關心神的事情.
神的教會.
給我的定義是.
教會是一群基督跟隨者.
去實踐信仰的群體.
無論逆境還是順境.
我們一起去面對.
這個挑戰.
在不同的時區.
不同的時間.
去見證耶穌.
見證耶穌的意思是.
耶穌有多少的東西.
告訴其他人.
未信的人知道.
原來耶穌是一位又真又活的神.
而不是說一些.
我們額外的宗教活動.
那種禮儀.
說完這幾方面.
我嘗試一個總結.
神究竟關心什麼.
或者我們關不關心神的事情.
兩節經文在我自己過往.
裡面提醒又提醒.
究竟我們現在.
自己在做什麼.
他說.
以塞翁書第49章第6節.
他說:你作我的僕人.
即是無論是解釋上.
是教會也好.
是神的以色列人也好.
或者是神的一個人也好.
使亞國眾支派復興.
使以色列眾蒙保存的人歸回.
當時是在說.
以色列人面對.

$^{761}$被擄之後要回歸的情況.
我們要令到.
以色列回到加拿大.
回到神所應許的地方.
但現在還不行.
那群以色列人很關心在政治上.
經濟上去關心這件事.
不過神說.
然而此事尚小.
不是不重要.
重要,不過有一件事情對比.
這些事情更加重要.
是什麼呢?我還要鮮你.
即是這個以色列人或者教會.
或者個人也好.
有一個任務交給他的僕人.
作萬邦之光.
即是說.
在這個世界上.
有什麼地方是黑暗的.
你就要去到.
你見到哪個地方的時候.
你作為神的僕人.
或者神的兒女.
你就要將你的光照去黑暗的地方.
哪裡是黑暗?.
這裡很光.
黑暗不是單是.
物質上那種黑暗.
而是在心靈上.
在物體上,道德上.
社會上.
是不見得光的一些情況.
你就要去.
神就掛著這些黑暗的地方.
很想黑暗的地方能夠有光去到.
神不是掛著.
要令到以色列人.
要再次偉大.
這是重要,但是對於.

$^{801}$黑暗幫人有光.
這件事情神所更加關心.
這件事情.
先不用太重要.
使你施行我的救恩.
直到地極.
舊約裡面.
神的心意是很想.
每個人,每個群體.
能夠將世界.
將任何的地方.
有光去到.
這只是舊約.
新約是否講這件事呢?.
一樣,很熟悉的.
門徒見到耶穌.
哇,復活升天了.
哇,這麼有能力的神.
厲害.
復活了,還沒升天.
耶穌這麼厲害,我們什麼時候復國.
這是第四行傳第一章.
第八節之前的對話.
耶穌說,不是.
聖靈降臨在你們身上.
你們只要得到這個能力.
耶路撒冷,猶太傳遞.
撒瑪利亞,直到地極.
作我的見證.
就將我們經歷耶穌.
那種的能力.
經歷耶穌的生命.
讓外邦人知道.
因為外邦人在一個黑暗的地方.
在一個黑暗.
不見得光的地方.
就要照著他們.
這是神的心意.
神的期盼.
講完這些.

$^{841}$我自己對神的.
紀念神的作為的時候.
接著講.
在傳道書結束的時候.
他說,這些事你都要聽見了.
聽完了.
結論是什麼呢.
敬畏神,謹守.
祂的誡命.
這是人當盡的本分.
神對期望.
我們作為一個人.
有個期望.
本分就是.
懂得敬畏神.
進行祂的誡命.
意思是什麼呢.
我們要以神為我們懼怕的對象.
因為人所做的事情.
連一切隱藏的事情.
無論是善是惡.
神都必審問.
是不是這麼嚴重呢.
是這麼嚴重的.
神會審問我們在世界上.
由少年人到老年人.
人生這麼多階段.
你做過好事壞事.
見得光不見得光.
神都要審我們.
很嚴重.
怎麼辦呢.
今天.
這段經文.
可能很快會過去.
但今天離開這個會堂之後.
我們仍然好像.
私人所說的.
你可能不懼怕神.
但我們可能懼怕另一件事.

$^{881}$今天.
身在香港這個地方.
你有沒有恐懼.
看看.
時間無多.
不知道有多少時間.
不是嚇大家.
真的不知道有多少時間.
好像.
說到.
打仗的聲音周圍起的時候.
不知道有多少時間.
美元可能要貶值.
拿著美金怎麼辦.
不知道怎麼辦.
然後整個世界變幻.
不確定性.
有變化性.
不確定性 複雜性 模糊性.
脆弱性 焦慮性.
不可理解的情況.
整個世界.
你不知怎麼辦.
星期五股票跌了五百多點.
星期一不知道升多少點.
都不知道買哪隻股票好.
沒有把握.
可以有把握.
怕不怕.
這個狀態是我們今天身處的.
一個相對.
好像一個.
經濟不錯的社會裡.
會有擔心.
我們身為在香港的一份子.
都難免受住.
這樣的情況影響.
時間.
無多了 世界變幻.
都當在教會的情況.

$^{921}$會更加擔心.
整個世界變得那麼多.
我們教會怎樣追得上.
這個變化呢.
會不會怕被淘汰呢.
這些我們在教會服事的人.
都會擔心.
怎麼辦呢.
外面頂展會有那麼多難處.
外面有福音工作那麼多.
教會內部.
內憂外患都來了.
不知道教會的長輩.
有什麼感覺.
真的 怎麼辦.
經濟又不肯定.
很多事情我們都害怕.
你怕.
究竟你怕這些環境.
還是像聖經所說.
選擇去怕上帝.
神必審問.
關心神關心的事.
不好人事.
神的心意.
是想人心歸回.
人心歸回.
包括那些在黑暗當中.
需要歸回給神.
我們這班信徒.
都要人心歸回給上帝.
我們的恐懼.
不是再恐懼這些東西.
是恐懼上帝.
回到上帝的身邊.
如果我們的人心.
生命不歸回給神的時候.
我們怎可能去鼓勵.
其他人歸回給神呢.
這是我們的矛盾.

$^{961}$時候到了.
審判.
要從神的家開始.
大家不用怕.
不用那麼直證.
真的.
這是時代你願意.
還是不願意.
聖經很清楚告訴你.
我們要面對這樣的終局.
不用怕.
有人在大廳.
約翰對那出來.
要受洗的人說.
毒蛇種類誰指示你.
逃避那將來要的憤怒呢.
這句說話講了二千多年.
你們要結出果子.
來悔改心傷症.
施洗約翰很簡單.
當你們當中.
那些毒蛇種類害怕的時候.
不要擔心.
生命結出果子就可以了.
不要憑著說你是心裡.
亞伯拉罕是我的祖宗.
不要因為你的宗教背景.
你的歷史.
跟亞伯拉罕有關係.
不是這樣的.
我告訴你們神能從.
這石頭中吸亞伯拉罕.
興起子孫來.
現在虎子已經放在樹根上.
凡不結好果子的.
樹就.
撼下來 凋在火裡.
很嚴厲的.
很嚴厲的警告.
不要因為你背景是怎樣.

$^{1001}$不要因為你出身是怎樣.
如果我們生命沒有結出.
這個跟悔改心傷症的.
行為的話.
表現出你跟神的.
那種關係的話.
劈下來了.
斧頭已經在.
樹根上.
不結好果子的樹.
就撼下來.
凋在火裡.
那群人聽完之後怎麼辦?.
怎麼辦?.
眾人問他這樣.
我該做什麼呢?.
眾人問他.
也有衰來來要受洗.
對他說老師我該做什麼呢?.
也有士兵問他說.
我該做什麼呢?.
在那個場景裡面.
作者用了三個.
總結了.
可以做什麼呢?.
我們怎樣能夠將我們的.
生命能夠跟我們的.
信仰結連呢?.
最基本.
還原基本步.
重新再看看我們的生命.
施生說.
有兩件以上的.
分給沒有的.
有食物就該怎樣做.
是什麼意思呢?.
當你有的時候.
見到人沒有的話.
幫他分享.
當有很多食物的時候.

$^{1041}$見到人沒有東西吃的時候.
將你的食物分享.
尤其是在我們.
今天的世界裡面.
相對來說是豐富的.
但在我們周遭裡面.
還有很多人被我們缺乏的.
我們有沒有賭注連吻的心.
去將我們有的跟我們分享呢?.
生命結束這個果節之後.
那些人就會好奇.
為什麼這群跟隨耶穌的人.
這麼願意有的.
幫我們分享呢?.
你不用說耶穌是救主.
你不用說這件事.
他自然會問你.
為什麼你的生命與眾不同.
你這麼願意跟別人分享呢?.
會這樣的.
但是要小心.
有些非信徒.
比我們做這些更加做得厲害.
如果你認識佛教徒的話.
他更加願意捐錢.
如果在台灣.
捐錢救災是第一時間.
去災場救那些人.
我們不要覺得.
他做得好像很偉大.
這些其實對於那些.
未信的人來說只是小事.
但我們覺得基督徒.
跟別人分享已經是很偉大的事情.
其實不是很簡單.
第二個場景.
有稅利.
在工作裡面的人.
你做一個稅利.
收稅的.

$^{1081}$除了規定數目.
不要多收.
是什麼意思呢?.
做一個忠心的工人.
不要偷懶.
不要正正常常.
做一個基督徒的工人.
忠心完成你的工作.
那些人就會看到.
你信耶穌的人.
跟隨耶穌的人.
真的很正直.
其他人不正直的時候.
就會看到分別.
你把光照在黑暗當中的時候.
就覺得你的生命與眾不同.
你為什麼可以有.
這種生命力.
可以做到這樣的事情呢?.
這是第二個情況.
我們要說一個公平.
一個忠心的生活.
不要疏疏忽忽.
不要輕輕忽忽.
完成你所作的事情.
因為這些工作是神所教給你的.
第三個場景.
士兵可以做什麼呢?.
不要勒索任何人.
不可以敲詐人.
在工作中.
一些比較.
做領袖的.
有權力的人.
不可以以權謀私.
不可以用你的職位.
去勒索或欺詐任何人.
那些人會看到.
原來基督徒的老闆.
做人這麼正直.

$^{1121}$這麼有誠信.
他為什麼可以這樣做呢?.
做老闆.
做老闆可以很貪心.
做老闆可以有很多手段.
但是這個人.
這麼正直的時候.
他覺得這個基督徒老闆.
生命與眾不同.
這樣的生命見證.
這樣的生活態度.
就會吸引一些人.
為什麼你的生命.
可以這麼平安喜樂.
你告訴我.
你的秘訣是什麼.
當我們生命有了.
學習之後.
在我們的場景裡面.
人家見到我的生命.
很自然的.
去問我們信仰的緣由.
這個就是.
有人問你心中盼望的緣由.
你就要常作準備.
要溫柔敬畏的心回答他們的問題.
是一個.
主動的被動.
人家問我.
我們信仰的緣由.
而不是說.
我派一張單張.
他不想要的就硬要告訴他.
那個過後是相反的.
當我們在.
宣教工場裡面.
聽見一些見證的時候.
在說什麼呢.
有些在中東的地方.
他們要不可以.

$^{1161}$因為簽證問題.
現在宣教世界因為簽證.
很嚴重的問題.
很難去用宗教簽證.
入境.
譬如剛才說中東國家的時候.
我們怎麼可以去到.
這些地區裡面.
去見證耶穌呢.
那個同事就跟我說.
我們需要一些專業人士.
他在生命當中.
在職場當中.
他做得是有見證的.
他就可以憑著他的專業.
去到這些中東的國家.
做工作.
透過他工作裡面的見證.
在當中生活的見證.
能夠.
講出信仰的源由.
去吸引人去問他信仰.
這樣去.
作成耶穌的見證.
現在我們.
如果憑著一個專職的.
去到那些地方去講聖經的時候.
人家未必.
去給你入境.
但如果我是一個.
某個行業上面是專業的.
他請到我的話.
我就可以這樣進去.
在當中見證我的生命.
有人這樣說.
我們的生命.
如果在職場裡面講的時候.
那些人看見我的生命.
一天.
做十個小時的話.

$^{1201}$你就在當中見證十個小時.
耶穌的事情.
如果你在職場裡面.
一個星期.
五六天.
還是這麼弱的時候.
你說第七,第六天,第一,第七天.
你回來教會啦 教會很可愛的.
他就會講一句話.
你講什麼.
你平時信耶穌都信成這樣.
你不要想著你回教會就變得好人.
那些人不會聽我們.
那種的見證.
除非我們在生活當中.
日常的見證裡面.
已經緊緊結連耶穌的話.
我們所做的一切一切.
去見證到耶穌的作為的話.
我們生命有一個屬靈果節之後.
人們見得到.
很願意去吸引那些人.
去問我們的信仰.
這個就是基本上對猜拳.
對使命.
對教會的一個反省.
盼望我們每個等姐妹.
由最基本步開始.
生命裡面.
一個禮拜168小時.
除了這兩三個小時在教會.
見到等姐妹之外.
你見得最多的是你的同事.
你的朋友你的家人.
就在他們的生命當中.
活現出這個屬靈的過節.
去吸引他們明白.
有一天神就說.
你的生命都挺好.
或者都值得我.

$^{1241}$願意去發更加多的.
猜拳的工作.
在這裡.
接著再去到海外.
或者去做其他.
如果我們說要做海外猜拳.
連最基本.
生命見證都做不好的時候.
怎麼可能.
出到去說你是一個.
好的傳道者呢.
盼望我們每個等姐妹.
由我們今天開始.
多多少少.
不會遲.
不要等.
生命邁向神的時候.
奇妙的事情就會發生.
神的猜拳就在我身上.
去彰顯了出來.
\newpage



\section{士師記 17:1-18:6}
\label{sec:031602}
\textbf{迷失世代 何去何從?}~~~~潘建堂牧師
\newline
\newline
連結: \href{http://www.church.com.hk/acms/content.asp?site=cdc\&op=show\&type=product\&code=031602\&layout=sermon}{\texttt{031602}} ~~~~ 語音日期: 2025-07-13
\newline
\newline
\hyperref[sec:031601]{\small{< < < PREV SERMON < < <}}
~
\hyperref[sec:index_chronic]{\small{[返順時目]}}
~
\hyperref[sec:031603]{\small{> > > NEXT SERMON > > >}}
\newline
\newline
士師記 17:1-18:6
\newline
\begin{longtable}{cl}
\hline
\hline
章節 & 經文 (和合本修訂版)\\
\hline
17:1 & \begin{tabularx}{0.7\textwidth}{X} 以法蓮山區有一個人，名叫米迦。 \end{tabularx} \\ \\ \relax
17:2 & \begin{tabularx}{0.7\textwidth}{X} 他對母親說：「你的一千一百塊銀子被人拿走了，為此你發咒起誓，也說給我聽。看哪，銀子在我這裡，是我拿的。」他母親說：「願我兒蒙耶和華賜福！」 \end{tabularx} \\ \\ \relax
17:3 & \begin{tabularx}{0.7\textwidth}{X} 米迦把這一千一百塊銀子還他母親。他母親說：「我把這銀子分別為聖，親手獻給耶和華，為我兒子造一尊雕刻的像，以及一尊鑄成的像。現在我把銀子交給你。」 \end{tabularx} \\ \\ \relax
17:4 & \begin{tabularx}{0.7\textwidth}{X} 米迦把銀子還他母親，他母親把二百塊銀子交給銀匠，去造一尊雕刻的像，以及一尊鑄成的像，安置在米迦的房子裡。 \end{tabularx} \\ \\ \relax
17:5 & \begin{tabularx}{0.7\textwidth}{X} 米迦這個人有了神堂，又造了以弗得和家中的神像，派他的一個兒子作祭司。 \end{tabularx} \\ \\ \relax
17:6 & \begin{tabularx}{0.7\textwidth}{X} 那時，以色列中沒有王，各人照自己眼中看為對的去做。 \end{tabularx} \\ \\ \relax
17:7 & \begin{tabularx}{0.7\textwidth}{X} 猶大的伯利恆有一個年輕人，是猶大族的人。他是利未人，寄居在那裡。 \end{tabularx} \\ \\ \relax
17:8 & \begin{tabularx}{0.7\textwidth}{X} 這人離開猶大的伯利恆城，要找一個可住的地方。他來到以法蓮山區米迦的家，還要往前走。 \end{tabularx} \\ \\ \relax
17:9 & \begin{tabularx}{0.7\textwidth}{X} 米迦對他說：「你從哪裡來？」他說：「我從猶大的伯利恆來。我是利未人，要找一個可住的地方。」 \end{tabularx} \\ \\ \relax
17:10 & \begin{tabularx}{0.7\textwidth}{X} 米迦說：「你就住在我這裡吧！我以你為父為祭司，每年給你十塊銀子和一套衣服，以及生活所需的食物。」利未人就來了。 \end{tabularx} \\ \\ \relax
17:11 & \begin{tabularx}{0.7\textwidth}{X} 利未人願意和這人同住；他待這年輕人如自己的兒子一樣。 \end{tabularx} \\ \\ \relax
17:12 & \begin{tabularx}{0.7\textwidth}{X} 米迦授這年輕的利未人祭司的職任，他就住在米迦的家裡。 \end{tabularx} \\ \\ \relax
17:13 & \begin{tabularx}{0.7\textwidth}{X} 米迦說：「現在我知道耶和華必恩待我，因為我有利未人作我的祭司。」 \end{tabularx} \\ \\ \relax
18:1 & \begin{tabularx}{0.7\textwidth}{X} 那時，以色列中沒有王。但支派的人還在覓地居住，因為直到那日，他們還沒有在以色列支派中抽籤得地為業。 \end{tabularx} \\ \\ \relax
18:2 & \begin{tabularx}{0.7\textwidth}{X} 但人從瑣拉和以實陶派本族中的五個勇士，去窺探偵察那地，對他們說：「你們去偵察那地。」他們來到以法蓮山區米迦的家中，就在那裡住宿。 \end{tabularx} \\ \\ \relax
18:3 & \begin{tabularx}{0.7\textwidth}{X} 他們臨近米迦的家，聽出那年輕的利未人的口音，就繞到那裡，對他說：「誰領你到這裡來？你在這裡做甚麼？你在這裡得了甚麼？」 \end{tabularx} \\ \\ \relax
18:4 & \begin{tabularx}{0.7\textwidth}{X} 他對他們說：「米迦如此如此待我，他雇用我，我就作了他的祭司。」 \end{tabularx} \\ \\ \relax
18:5 & \begin{tabularx}{0.7\textwidth}{X} 他們對他說：「請你求問神，使我們知道所走的道路是否通達。」 \end{tabularx} \\ \\ \relax
18:6 & \begin{tabularx}{0.7\textwidth}{X} 祭司對他們說：「你們平平安安去吧，你們所行的道路是在耶和華面前的。」 \end{tabularx} \\ \\
[1ex]
\hline
\hline
\end{longtable}
$^{1}$《錄音》或《錄影》的版權是屬於基督教宣導會不覺堂所有.
我們的屬靈傳統之一是以真理教導為本.
歡迎你將此《錄音》或《錄影》的連結傳給親友.
懇請注意未經授權請勿轉錄或剪輯本《錄音》或《錄影》.
甚願好像《使徒行傳》二十章三十二節下所說.
「這道能建立你們,叫你們和一切誠誠的人同得基業」.
又讓我們知道如何去跟從主你的心意.
為我們這個世代去受望.
我們從心祈禱是奉耶穌基督的名字.
阿們.
今天是社關主日.
我們會透過時事記十七,十八章的訊息去思想.
我們應該或可以怎樣去到.
為我們所關心的這個世代去受望,祈禱.
講到時事記的時候.
大家通常都會記得很多稱為士師的人物.
那些人物多數都是一些勇士.
很有力量.
打仗很厲害.
神給他們有這種力量.
他們透過神給他們的大能大力.
去戰勝那些侵害以色列人的仇敵.
但是當以色列人離開上帝的時候.
他們就被周圍的敵人攻打.
這是神給他們的懲罰.
當他們悔改的時候.
神就興起士師去拯救他們.
但是當以色列人得救之後又犯罪.
犯罪之後又悔改.
反反覆覆經過很多次的犯罪.
很多次的悔改.
很多個的事先.
這些士師的人物.
我們比較熟悉的可能是.
基奮,底波拉,貪孫,耶不他,摩托列.
其實還有十幾個.
所以他們反覆了很多次.
但是今天我們所讀的英文.
剛才你聽.
沒有.

$^{41}$沒有聽見士師.
任何士師的名字.
只會聽到.
這一段是非常混亂的日子.
這是第十七章.
第十七章到尾大概有兩三件事.
特別是講述當時的混亂情況.
按照士師記的編排.
這個編排.
學者普遍認為.
是講述了那麼多士師之後的幾個附錄.
這是第一個附錄.
這個附錄的事件.
在什麼時間發生是什麼.
是很難去確定.
因為整個段落裡面沒有透露一個很清楚的背景.
甚至可能就是.
其實士師記的背景就是.
摩西死後有約書亞去接任.
但是約書亞死後就沒有人接任.
可能就是根本這段時間裡面.
是還沒有士師出現都不一定的.
學者形容那段士師的時間.
就是十二支派的聯盟.
十二支派都有在這裡.
入了艱難.
但是他們十二支派是自己管理自己.
他們不單止沒有可能叫做聯邦政府.
沒有的.
實際上他們在宗教上和在政治上.
都是非常的混亂.
記得在幾個月之前.
我在晚上.
就是和太太坐地鐵回家.
都開始比較晚了.
由金鐘.
經過金鐘.
然後回去港島線.
去到魚則魚涌站.
又近了家裡.

$^{81}$在魚則魚涌站.
多些人上車.
有一個婆婆上了車之後.
就直向我那邊走.
坐在我旁邊.
就問我.
先生.
請問彩虹站在哪裡下車呢.
如果你熟悉港島線熟悉地鐵.
彩虹在哪裡呢.
在觀塘線.
綠色線上面那裡.
我當時身處的是.
我們附近的魚則魚涌站.
往柴灣方向.
當她問完這個問題的時候.
車門已經關了.
很明顯.
我告訴她.
不是這輛車.
我和太太在對望.
大家都明白.
一個老人家在魚則魚涌站.
一個這麼多選擇.
路線這麼複雜.
行人路線都很複雜的地方.
她坐錯車.
怎樣三言兩語就教到她.
去那個地方呢.
還開車去太古勒.
我們就跟她說.
我們帶她坐對那輛車.
如果要坐對那輛車.
起碼要在中間.
幫她去到油塘.
然後她可能都要再問人.
想著想著.
我們就問婆婆.
一路走的時候就說.
有沒有家人去接她.

$^{121}$她說不用了.
我習慣自己回家.
我們都是鼓勵她.
你有沒有電話.
打電話給她.
聯絡之後.
我們覺得都不太足以因.
都是送婆婆去到彩虹站.
不用家人到處找.
去到彩虹站.
我們就把婆婆交給家人.
也為家人打氣.
我們表示明白.
她應該很艱辛.
她說真的不是第一次.
原來婆婆喜歡自己一個人去參加活動.
所以我們都稍為知道.
就是說周圍全港都有很多社區活動.
我鼓勵她.
她住的區應該都有很多社區活動.
那裡還有教會.
教會都有很多好的活動.
包括有長者活動.
這個過程帶給我自己一個提醒.
就是香港人都愈來愈長壽.
我都可能會活長一點.
不知道將來我會不會都是.
在進到地鐵之後去問人.
先生請問炮台山站怎樣去.
我想去銅牌.
老人家年紀大.
走迷了路是一個可以理解.
但是我們又知道.
是很需要關心和照顧的事.
不過走迷路是不是只是老人家的事呢.
可能是整個社會的事.
我們今天所讀的《時之記》.
就是整個社會都走迷了路.
特別是在信仰上面.
他們走迷了路.

$^{161}$事實上人類的問題總是一代接一代.
解決了一個問題.
又有另一個問題去產生.
每一代的人都需要面對自己.
眼前環境的局面.
《時之記》第十七章.
我們剛才所讀的一個人.
是二百年山區的人.
這個家庭我們數到有三代.
他有媽媽 他有兒子.
三代同堂.
米家這個人他偷了錢之後.
就聽到母親救助偷錢的人救了他.
他就害怕了.
於是就拿回來.
而母親就說錢是會分別為性.
她說要親手獻給他.
大概有時候我們都會.
遇到困難的時間.
不見了東西.
有時候都會跟上帝說.
如果我找得到的話.
上帝我就這樣這樣.
如果看經文的脈絡.
大概都是.
不見了錢就許願.
好像舊約那樣.
我們可能就說.
不見了錢就會還給他.
那好了 現在找到了.
所以很多小學者在這裡都會問.
明明說近一千一百的銀子.
現在後面那裡說的.
說二百交給銀長.
那九百去了哪裡呢.
這是一個謎.
但是這二百銀子交出去.
其實最終又是製造了一些偶像.
而米家有了錢.
他身上有錢.

$^{201}$有銀 有神像.
製造了祭司穿的袍.
祭司穿的袍.
那個外袍就叫做二不得.
有了這些服裝.
當然穿了.
就派他的兒子.
做了他家裡的祭司.
就實行他家裡的祭司的DIY.
Do it yourself.
自己搞定.
當然神聖的角色是私有化.
讓他進一步.
他後來遇到一個寄居.
在那裡的一個利尾人.
這個利尾人.
可以幫他忙.
他就寵金禮聘.
做他家裡的祭司.
大概他覺得自己的兒子.
做祭司.
家裡的祭司都不夠功力.
所以既然祭司本身.
如果在以色列人當中.
就是出自利尾之派.
當然不是所有利尾之派.
都是做祭司.
要是亞倫的後裔.
起碼利尾之派.
是辦理會幕的事情.
即是說是政治的事.
都算是一個.
有正式宗教事務的角色.
所以有一個利尾人.
在他家裡做祭司的話.
他在他家裡.
即是社區的環境當中.
可以說是一個宗教的富戶.
資源特別多.
這個情況其實真的看到.

$^{241}$整件事是錢作怪.
錢令到米家和利尾人.
都落入信仰的網絡.
在這個家庭裡.
錢的迷惑.
是由他的母親.
以及可能已至他的下一代.
都是這樣去思考.
還有就是.
米家會相信一些好像魔法的東西.
他聽到母親咒詛就會很害怕.
會覺得這些咒詛是有效的.
他還有就是覺得家裡.
做祭司.
特別有利尾人做祭司.
會有更多福氣.
是他說的.
所以他說.
我知道耶和華必因待我.
因為我有這個資源.
這些就好像相當於我們今天.
民間宗教的心態.
它的實物.
好像在代表一些神聖的福氣.
所以如果搬回來.
好像如果對應我們自己的資源來說.
如果我們真的以為.
家裡多掛幾句金句.
然後又多掛幾個十字架.
這樣就會更屬靈.
這樣就會更多福氣.
當然掛多一點.
如果真的提醒多一點.
那當然是有作用的.
不過就是宗教的用品.
本身沒有神奇力量.
我們也不會拿著十字架.
好像影片那樣驅魔.
它是幫助我們學習真理的工具.
十字架是提醒我們.

$^{281}$我們不單止不可以敬拜別神.
也不可以用任何地上的東西.
去代表上帝.
今天我們在崇拜裡面也不會.
我們的禮儀.
我們的講章.
音樂等等.
我們所做的一切.
不是直接代表上帝.
但是那些東西是有用的.
也不是直接代表上帝的榮耀.
只能夠說見證上帝的榮耀.
我們只是在見證上帝的真實.
上帝的慈愛.
所以我們只是透過不同的方法.
去學習真理.
透過學習真理去親近神.
節子記接著去到記載.
十八章裡面.
這個米家以為這樣就會得到福氣.
最終他所設立的東西.
就在混亂當中被人搶走了.
因為彈支派的人.
他們還沒有分地.
所以他們要尋找自己的地方.
例如11個支派.
他們甚至跨越地區.
他們去到200年的山區.
他們經過了米家的家.
他們想到一個對象.
就是拉益的那個地方.
那裡挺好的.
如果住那裡.
他們預備去攻打.
攻打之前還忘記.
這裡還有些好東西可以拿.
於是就搶走了米家的祭司袍和神像.
結果他們用他們的勢力.
去到米家毫無反抗之力.
這個家庭的信仰觀念.

$^{321}$受到金錢和私慾影響.
一代接一代.
最終蒙受損失.
而我們去看以色列人.
他們的歷史.
由約書記到節子記.
他們還沒有完成.
很多東西都還沒有完成.
還有很多問題還沒有解決.
留給下一代的人去面對.
今天我們每個時代都是這樣.
我們面對一個世代.
去到另一個世代的時候.
除了我們身體的DNA.
我們去問究竟我們整個社會.
存了些什麼.
留下了些什麼.
給下一個世代呢.
我們做成年人的.
我們都是曾經年輕過.
都是曾經是一個年輕人.
但是經過我們成長.
成為成年人.
有些時候可能都忘記了.
在年輕的時候的狀況和需要.
我自己想會不會.
我都可能是為了超越自己的限制.
所以不斷去打倒以前的我.
然後我可以去到更加成功呢.
如果打倒以前的我.
不就是打倒年輕的我嗎.
會不會又不知不覺地變成.
是在壓抑新一代他們的成長和貢獻呢.
另一個極端就是.
都有時候會聽到他說.
可能是過於保護.
過於保護年輕一代.
其實年輕的一代.
他們在面對他們的環境的時候.
有些自然的過程.

$^{361}$他們需要面對一些真實的挑戰.
這些挑戰有些地方是對他們有幫助.
令他們更加強壯.
但過於保護就失去了這些面對挑戰的能力.
讓他們日後真的面對困難的時候.
真的不知如何處理.
當然在社會上也有不少的人.
有很多的家庭.
很多的家長.
是真正明白青少年人的需要.
用適切的方式和他們去溝通和相處.
幫助他們去成長.
我們很方便.
現在在網上我們都很容易去認識一下.
年輕人他們的需要.
我用不同的顏色.
有生理方面他們的成長.
他們心理方面的需要.
他們在社交層面的需要.
他們在整個大社會裡面的需要.
如果用不同的顏色.
每一項在細分的時候.
在生理上他們面對快速的成長.
有很多的變化.
他們的健康狀況也都很需要去關心.
他們越來越掛著自己的自我的形象.
自我的意識.
他們的情緒授著身體和他遭遇的東西.
有很多很容易的變化.
他們面對很多的壓力.
需要適切的成長.
他們面對家人的關係.
面對他周圍學校等等的環境.
不同的人.
同學 老師 不同的朋友.
這些都是他們所面對.
他們面對整個社會的價值觀.
給他們的方向.
但有時候可能是壓力.
他們很需要社會的支持.

$^{401}$不單止是因為他們面對的人的支持.
是整個社會環境的安排.
能不能給他們支持呢.
這些都是我們可以去到.
去到明白他們更多的地方.
當然我們做父母的都很留心.
很直接就是我們家庭.
我們可以怎樣去到對他們有好的影響呢.
不過我們今天寫這本書.
我們就可以再想多一點.
我們社會的環境會怎樣影響他們呢.
剛才的經文都提到環境對於人的影響.
經文我們剛才所讀的有兩次提及.
那時以色列中沒有王.
其實由第十七章至到第二章.
我講的幾個附錄裡面.
總共是出現了四次.
而第一次和最後那次.
即是史記最後那一節.
都補上一句說話.
就是各人照自己眼中看為對的事.
去到做.
很多學者認為.
史記最後重覆去到說.
那時以色列中沒有王.
這句說話可能在後期.
當大衛王朝興盛的時間.
加上一句評語.
可能表示以色列人.
他們在史記那段很混亂的時間.
正正就是需要一個王的出現.
好像大衛王朝那樣.
不過我們知道.
以色列人他們興盛的時間.
主要都是大衛王和所羅門的時代.
之後很快就沒落和分裂.
所以實際上王的出現.
在整個以色列人當中.
救到的時間其實不算很多.
其實我們看回紙之記.

$^{441}$緊接紙之記的撒慕爾記.
就已經告訴我們.
以色列人真正缺乏的是甚麼呢.
很重要.
少了甚麼.
耶和華的言語.
不常有異象.
但是很少跟他們說話.
當以色列人他們很需要王.
求納王的時間.
然後就跟撒慕爾說.
你只管聽從百姓.
向你說的一切話.
因為他們不是厭棄你.
而是厭棄我.
不要我作他們的王.
之後又警戒他們說.
你們若不斷作惡.
你們和你們的王.
必一同滅亡.
所以以色列人的問題.
不是在於他們有沒有人.
去繼承約書亞.
有沒有王.
而是因為他們沒有以耶和華為他們的王.
所以實際上紙之記是告訴我們.
每個世代最需要的是上帝的說話.
沒有上帝的說話.
那個時代就要陷入混亂當中.
當說了那麼多王.
我們現在有沒有王.
我們現今仍然有的.
叫做君主.
都有一些是君主的制度.
網上告訴我們大概有二十多個.
百度百科就告訴我們有二十六個.
我們熟悉的.
有英國 西班牙 荷蘭 歐洲等等.
亞洲我們熟悉的可能是.
沙特 阿拉伯 日本 泰國 馬來西亞.

$^{481}$非洲都有摩洛哥.
這些叫做王 君主.
只是扮演一些儀式性的角色.
負責簽署一下文件.
重要的日子出現一下.
通常是這樣的.
不過無論是不是用君主制度.
或者其他的制度.
任何一種領導的模式.
人的制度都是有它的限制.
也都有它的困難.
任何的制度.
它能不能夠帶給人民真正的幸福呢.
聖經給我們看的就是.
好在於它能不能夠跟隨上帝的創造.
世界的盡力創造.
都是最終由它去管治.
就是君王究竟能不能夠.
像耶和華那樣做一個好的牧人呢.
在古代的時代.
君王都被比喻為一個好的牧人.
去到牧羊羊.
好像詩篇二十三篇那樣.
就是那些如果得到照顧的時間.
那些羊是可以睡在那裡吃草那麼安樂的.
是不擔心有任何的威脅的.
很安樂的.
究竟我們有沒有跟隨上帝的心願呢.
回望我們的歷史.
很漫長的歷史.
有人就說.
這個世界的歷史是由很多的戰爭不斷發生.
間中有和平.
有些人說在和平之中偶然有戰爭.
相同.
我們相信不同的地區.
不同的時代.
不同的人都有不同的感覺.
不過其實回望我們的歷史.
是沒有人能夠真正絕對去掌管整個的局勢.

$^{521}$大局是在上帝的手裡.
所以我們很需要.
我們向這位上帝去祈禱.
去為每一個世代去守望.
去求平安.
正如在《箴言書》.
我們今季讀《箴言書》.
裡面都有說.
王的心在耶和華手中.
像河水.
他能使他隨意運轉.
君王有什麼想法.
絕對在上帝的手中.
人一切所行的.
在自己眼中看為正.
唯有耶和華衡量人心.
剛才說到.
各人按自己的意思.
看為好的事去做.
這裡也說.
聖經從來都不否認.
王和其他人的本質是一樣.
所以他們都是做他們看為好的事.
所以最終的掌權.
是在上帝的手裡.
而保羅也在新約裡面提醒我們.
字製了一點.
所以我勸你.
首先要為人人祈禱.
然後就是.
祈求禱告.
為君王和一切在位的也要如此.
使我們能夠敬虔端正地過平靜寧靜的生活.
這裡好像分開了次序.
為人人和為君王.
不過在新漢語曆本裡面.
是由一氣康成去說的.
所以萬人為君王.
都是這樣一致的說.
是這樣祈禱.

$^{561}$意思是什麼呢.
其實是相關的.
就是我們要為所有的人.
和那些影響所有人.
就是那些領袖.
去祈禱.
因為他們很影響.
所有人他們生活的環境.
我講到影響的時間.
我們不要只想那些領袖.
現在的影響力不單止是那些管治的系統.
在現代的社會裡面.
影響我們的生活.
影響我們周圍的人.
特別可能影響多一些青少年人的.
是那些社會文化和媒體.
我們知道社交媒體是很方便的.
我們差不多個個都用.
不見了電話好像不知道怎麼辦.
對青少年人更加重要.
和他們很密切.
社交媒體對於青少年人的影響.
是真的有兩面的.
有積極的一面.
也有消極的一面.
積極的一面就是社交媒體.
可以幫助青少年人建立社交的聯繫.
擴大社交的圈子.
增強歸屬感和自信心.
也可以獲取資訊.
我這些資訊就是從這裡來的.
也可以學習.
不懂什麼在Google打一句.
就有很多幫助我們.
AI是不是.
但是社交媒體或者網上的資源.
也都有消極方面的影響.
我們也不可以去到不良的.
不良的使用的情況.
會可以導致網絡欺凌.

$^{601}$過度的比較.
憂慮.
抑鬱.
對青少年人甚至可能令他們對面不足.
注意力分散等等很多的問題.
很嚴重的.
甚至可以影響大陸的發育和認知的功能.
這些都是需要我們同心的留心去關注.
其實香港的基督徒群體.
在很多年前已經開始去關注.
傳播媒體對於整個社會的影響.
不專家的電子媒體.
也有參與一些傳媒監察的行動.
當時我也是和電子媒體一起有參與的情況.
在過程當中我們都會想.
有些機構會幫助.
可以在媒體或者公佈裡面.
表達一些意見.
對社會的關注.
對於傳媒的一些情況.
在過程裡面會發覺.
家長這個角色在社會裡面.
在家庭的身份也有一定的影響力.
當然我們今天不是每個人都是家長.
最基本我們都是在社會裡面的一個公民.
都是一份子.
我們都可以想一下我們可以怎樣做.
很重要最少我們作為基督徒.
作為社會一份子.
是可以為這個社會.
這個世代去祈禱.
在聖經裡面也是這樣去鼓勵我們.
聖經就告訴我們對以前的人說.
耶利米茲.
即使他們都被擄去到其他地方.
那個地方是外邦的地方.
沒有聖殿.
這樣生活和以前很不同.
祂說.
「使你們被擄到的那城.

$^{641}$你們要為那城求平安.
為那城向耶和華祈求.
因為那城得平安.
你們也隨著得平安」.
這個就是我們每一個基督徒都可以做.
所以最後我可以和大家分享.
其實我們一齊思想可以有幾個祈禱的方向.
第一.
我們可以為身處這個世代去祈禱.
因為每一個世代的大局都有它的困難.
每一個年代的人都有他所面對的挑戰.
第二就是我們特別為年輕人的需要去祈禱.
他們極其需要我們的關注.
我們成年人都曾經是年輕人.
所以我們只要回想一下就會明白.
多一點他們的需要.
我們要為他們本身的需要去祈禱.
特別幫助他們.
求主幫助他們在新心靈裡都有好的成長.
我們看到有不同的範圍.
他們需要被關注.
求神也幫助他們有智慧去保護自己.
忍受這個方法可能未必真是很有效.
有些人可能會有智慧幫到他們.
求神幫助他們都在困難的裡面.
找到合適的人去幫助.
我們也為青年人他們所面對的整個環境去祈禱.
求神幫助他們脫離那些會影響他們的風險.
無論是間接或者直接傷害他們的力量或者資訊.
求主都是難助.
也求主幫助我們.
我們每一個人每一個基督徒.
善用我們的家庭角色.
或者我們社群的角色.
可以去改善周圍的環境.
就在這段時間.
大家按照剛才所聽的感動.
去為我們的社會.
特別為年青人去禱告.
然後我帶領大家總結.

$^{681}$我們同心祈禱.
愛愛天父我們.
仰望你的恩典.
我們深深知道.
你是因著愛創造我們.
你也因著愛拯救我們.
我們所面對的整個世界.
特別近幾年的裡面.
都有我們常常聽到的戰爭.
各樣不同的混亂.
求主在這些地方.
去彰顯出來的公義和慈愛.
我們所身處的香港的社會.
求主去對領.
我們都有面對不同的困難.
求主去幫助.
讓我們的政府.
讓我們不同的社會人士.
我們都能夠同心的努力.
為香港的發展去付出力量.
求主幫助我們.
讓我們關心不同年代的人.
無論是長者.
成年人.
青年人.
兒童.
求主幫助.
社會能夠對不同的人士.
都有好的照顧.
我們今天.
我們特別去關心年青人.
求主眷顧他們.
幫助他們.
他們的需要.
都是能夠得到好的安排.
他們的心靈得到關心.
求主幫助我們有智慧.
怎樣去跟他們相處.
成為他們的幫助.
成為他們成長正面的力量.

$^{721}$求主幫助我們在社會當中.
善用我們的角色.
去改善我們周圍社會的環境.
以致他們能夠有更美好的成長.
脫離一切的凶險.
求主順應我們的禱告.
奉耶穌基督的名字.
阿們.
\newpage



\section{哥林多後書 8:1-10:18}
\label{sec:031603}
\textbf{哥林多前後書 (5) 學習奉獻的功課}~~~~劉少康牧師
\newline
\newline
連結: \href{http://www.church.com.hk/acms/content.asp?site=cdc\&op=show\&type=product\&code=031603\&layout=sermon}{\texttt{031603}} ~~~~ 語音日期: 2026-02-04
\newline
\newline
\hyperref[sec:031602]{\small{< < < PREV SERMON < < <}}
~
\hyperref[sec:index_chronic]{\small{[返順時目]}}
~
\hyperref[sec:031604]{\small{> > > NEXT SERMON > > >}}
\newline
\newline
哥林多後書 8:1-10:18
\newline
\begin{longtable}{cl}
\hline
\hline
章節 & 經文 (和合本修訂版)\\
\hline
8:1 & \begin{tabularx}{0.7\textwidth}{X} 弟兄們，我們要把神賜給馬其頓眾教會的恩惠告訴你們： \end{tabularx} \\ \\ \relax
8:2 & \begin{tabularx}{0.7\textwidth}{X} 他們在患難中受大考驗的時候，仍然滿有喜樂，在極度貧窮中還格外顯出他們樂捐的慷慨。 \end{tabularx} \\ \\ \relax
8:3 & \begin{tabularx}{0.7\textwidth}{X} 我可以證明，他們是按著能力，而且超過了能力來捐助，主動 \end{tabularx} \\ \\ \relax
8:4 & \begin{tabularx}{0.7\textwidth}{X} 再三懇求我們，准他們在這供給聖徒的善事上有份； \end{tabularx} \\ \\ \relax
8:5 & \begin{tabularx}{0.7\textwidth}{X} 並且他們所做的，不但照我們所期望的，更照神的旨意先把自己獻給主，又給了我們。 \end{tabularx} \\ \\ \relax
8:6 & \begin{tabularx}{0.7\textwidth}{X} 因此，我們勸提多，既然在你們中間開始這慈善的事，就當把它辦成。 \end{tabularx} \\ \\ \relax
8:7 & \begin{tabularx}{0.7\textwidth}{X} 既然你們在信心、口才、知識、萬分的熱忱，以及我們對你們的愛心上，都勝人一等，那麼，當在這慈善的事上也要勝人一等。 \end{tabularx} \\ \\ \relax
8:8 & \begin{tabularx}{0.7\textwidth}{X} 我說這話，並不是命令你們，而是藉著別人的熱忱來考驗你們愛心的真誠。 \end{tabularx} \\ \\ \relax
8:9 & \begin{tabularx}{0.7\textwidth}{X} 你們知道我們主耶穌基督的恩典：他本是富足，卻為你們成了貧窮，好使你們因他的貧窮而成為富足。 \end{tabularx} \\ \\ \relax
8:10 & \begin{tabularx}{0.7\textwidth}{X} 我在這事上把我的意見告訴你們，是對你們有益，因為你們開始辦這事，而且起此心意已經有一年了。 \end{tabularx} \\ \\ \relax
8:11 & \begin{tabularx}{0.7\textwidth}{X} 如今就當辦成這事，既然有願做的心，也當照你們所有的去辦成。 \end{tabularx} \\ \\ \relax
8:12 & \begin{tabularx}{0.7\textwidth}{X} 因為人只要有願做的心，必照他所有的蒙悅納，並不是照他所沒有的。 \end{tabularx} \\ \\ \relax
8:13 & \begin{tabularx}{0.7\textwidth}{X} 我不是要別人輕鬆，你們受累，而是要均勻： \end{tabularx} \\ \\ \relax
8:14 & \begin{tabularx}{0.7\textwidth}{X} 就是要你們現在的富餘補他們的不足，使他們的富餘將來也可以補你們的不足，這就均勻了。 \end{tabularx} \\ \\ \relax
8:15 & \begin{tabularx}{0.7\textwidth}{X} 如經上所記：多收的沒有餘，少收的也沒有缺。 \end{tabularx} \\ \\ \relax
8:16 & \begin{tabularx}{0.7\textwidth}{X} 感謝神，把我對你們的熱忱同樣放在提多心裡。 \end{tabularx} \\ \\ \relax
8:17 & \begin{tabularx}{0.7\textwidth}{X} 他固然聽了我的勸告，但自己更加熱心，自願往你們那裡去。 \end{tabularx} \\ \\ \relax
8:18 & \begin{tabularx}{0.7\textwidth}{X} 我們還差遣一位弟兄和他同去，這人在傳福音的事上得了眾教會的稱讚； \end{tabularx} \\ \\ \relax
8:19 & \begin{tabularx}{0.7\textwidth}{X} 不但這樣，他也被眾教會選派跟我們同行，把所交託我們的這捐款送到了，為的是榮耀主，也表明我們的好意。 \end{tabularx} \\ \\ \relax
8:20 & \begin{tabularx}{0.7\textwidth}{X} 我們這樣做，免得有人因我們收的捐款多而挑剔我們。 \end{tabularx} \\ \\ \relax
8:21 & \begin{tabularx}{0.7\textwidth}{X} 我們留心做好事，不但在主面前，就是在人面前也是這樣。 \end{tabularx} \\ \\ \relax
8:22 & \begin{tabularx}{0.7\textwidth}{X} 我們又差遣一位弟兄同去。這人的熱忱，我們在許多事上屢次考驗過，現在他因為深深信任你們，就更加熱心了。 \end{tabularx} \\ \\ \relax
8:23 & \begin{tabularx}{0.7\textwidth}{X} 至於提多，他是我的夥伴，為服事你們作我的同工。至於那兩位弟兄，他們是眾教會的使者，是基督的榮耀。 \end{tabularx} \\ \\ \relax
8:24 & \begin{tabularx}{0.7\textwidth}{X} 所以，你們務要在眾教會面前向他們顯明你們的愛心和我所誇獎你們的憑據。 \end{tabularx} \\ \\ \relax
9:1 & \begin{tabularx}{0.7\textwidth}{X} 關於供給聖徒的事，我本來不必寫信給你們； \end{tabularx} \\ \\ \relax
9:2 & \begin{tabularx}{0.7\textwidth}{X} 因為我知道你們的好意，常對馬其頓人誇獎你們，說亞該亞人預備好已經有一年了。你們的熱心感動了許多人。 \end{tabularx} \\ \\ \relax
9:3 & \begin{tabularx}{0.7\textwidth}{X} 但我差遣那幾位弟兄去，要使你們照我的話預備妥當，免得我們在這事上誇獎你們的話落了空。 \end{tabularx} \\ \\ \relax
9:4 & \begin{tabularx}{0.7\textwidth}{X} 萬一有馬其頓人與我同去，見你們沒有預備好，就使我們所確信的反成了羞愧；你們的羞愧更不用說了。 \end{tabularx} \\ \\ \relax
9:5 & \begin{tabularx}{0.7\textwidth}{X} 因此，我想必須鼓勵那幾位弟兄先到你們那裡去，把從前所應許的捐款預備妥當，好顯出你們所捐的是出於樂意，不是出於勉強。 \end{tabularx} \\ \\ \relax
9:6 & \begin{tabularx}{0.7\textwidth}{X} 還有一點：「少種的少收；多種的多收。」 \end{tabularx} \\ \\ \relax
9:7 & \begin{tabularx}{0.7\textwidth}{X} 各人要隨心所願，不要為難，不要勉強，因為神愛樂捐的人。 \end{tabularx} \\ \\ \relax
9:8 & \begin{tabularx}{0.7\textwidth}{X} 神能將各樣的恩惠多多加給你們，使你們凡事常常充足，能多做各樣善事。 \end{tabularx} \\ \\ \relax
9:9 & \begin{tabularx}{0.7\textwidth}{X} 如經上所記：「他施捨，賙濟貧窮；他的義行存到永遠。」 \end{tabularx} \\ \\ \relax
9:10 & \begin{tabularx}{0.7\textwidth}{X} 那賜種子給撒種的，賜糧食給人吃的，必多多加給你們種地的種子，又增添你們仁義的果子。 \end{tabularx} \\ \\ \relax
9:11 & \begin{tabularx}{0.7\textwidth}{X} 你們必凡事富足，能多多施捨，使人藉著我們而生感謝神的心。 \end{tabularx} \\ \\ \relax
9:12 & \begin{tabularx}{0.7\textwidth}{X} 因為辦這供給的事，不但補聖徒的缺乏，而且使許多人對神充滿更多的感謝。 \end{tabularx} \\ \\ \relax
9:13 & \begin{tabularx}{0.7\textwidth}{X} 他們從這供給的事上得了憑據，知道你們宣認基督，順服他的福音，慷慨捐助給他們和眾人，把榮耀歸給神。 \end{tabularx} \\ \\ \relax
9:14 & \begin{tabularx}{0.7\textwidth}{X} 他們也因神極大的恩賜顯在你們身上而切切想念你們，為你們祈禱。 \end{tabularx} \\ \\ \relax
9:15 & \begin{tabularx}{0.7\textwidth}{X} 感謝神，因他有說不盡的恩賜！ \end{tabularx} \\ \\ \relax
10:1 & \begin{tabularx}{0.7\textwidth}{X} 我－保羅與你們見面的時候是溫和的，不在你們那裡的時候向你們是勇敢的，如今親自藉著基督的溫柔和慈祥勸你們。 \end{tabularx} \\ \\ \relax
10:2 & \begin{tabularx}{0.7\textwidth}{X} 有人認為我們是憑著血氣行事，我認為必須敢於對付這等人；我但求在那裡的時候，不必這樣勇敢。 \end{tabularx} \\ \\ \relax
10:3 & \begin{tabularx}{0.7\textwidth}{X} 我們雖然在血氣中行事，卻不憑著血氣爭戰。 \end{tabularx} \\ \\ \relax
10:4 & \begin{tabularx}{0.7\textwidth}{X} 因為我們爭戰的兵器本不是屬血氣的，而是憑著神的能力，能夠攻破堅固的營壘。我們攻破各樣的計謀， \end{tabularx} \\ \\ \relax
10:5 & \begin{tabularx}{0.7\textwidth}{X} 和各樣攔阻人認識神的高壘，又奪回人心來順服基督。 \end{tabularx} \\ \\ \relax
10:6 & \begin{tabularx}{0.7\textwidth}{X} 我已經預備好了，等你們完全順服的時候來懲罰所有不順服的人。 \end{tabularx} \\ \\ \relax
10:7 & \begin{tabularx}{0.7\textwidth}{X} 你們只看事情的外表。倘若有人自信是屬基督的，他要再想想，他屬基督，我們也屬基督。 \end{tabularx} \\ \\ \relax
10:8 & \begin{tabularx}{0.7\textwidth}{X} 主賜給我們權柄，是要造就你們，並不是要拆毀你們；我就是為這權柄稍微誇口也不覺得慚愧。 \end{tabularx} \\ \\ \relax
10:9 & \begin{tabularx}{0.7\textwidth}{X} 我說這話，免得你們以為我寫信是要恐嚇你們。 \end{tabularx} \\ \\ \relax
10:10 & \begin{tabularx}{0.7\textwidth}{X} 因為有人說：「他信上的語氣既嚴厲又強硬，他本人卻軟弱無能，言語粗俗。」 \end{tabularx} \\ \\ \relax
10:11 & \begin{tabularx}{0.7\textwidth}{X} 這等人當明白，我們不在那裡時信上怎麼說，見面時也必怎麼做。 \end{tabularx} \\ \\ \relax
10:12 & \begin{tabularx}{0.7\textwidth}{X} 因為我們不敢將自己和某些自我推薦的人並列相比；他們用自己度量自己，用自己比較自己，是不明智的。 \end{tabularx} \\ \\ \relax
10:13 & \begin{tabularx}{0.7\textwidth}{X} 我們不願意過分誇口，但是我們只在神劃定的界限內誇口。這界限甚至擴展到你們那裡。 \end{tabularx} \\ \\ \relax
10:14 & \begin{tabularx}{0.7\textwidth}{X} 我們擴展到你們那裡時並沒有越過了自己的界限，其實我們是首先到你們那裡傳基督福音的。 \end{tabularx} \\ \\ \relax
10:15 & \begin{tabularx}{0.7\textwidth}{X} 我們不靠別人所勞碌的過分誇口；我們只希望你們信心增長的時候，所劃定給我們的範圍也能夠因著你們更加擴展， \end{tabularx} \\ \\ \relax
10:16 & \begin{tabularx}{0.7\textwidth}{X} 使福音得以傳到你們以外的地方，而不在別人的範圍之內，以別人所成就的事誇口。 \end{tabularx} \\ \\ \relax
10:17 & \begin{tabularx}{0.7\textwidth}{X} 但「要誇耀的，該誇耀主」。 \end{tabularx} \\ \\ \relax
10:18 & \begin{tabularx}{0.7\textwidth}{X} 因為蒙悅納的，不是自我稱許的，而是主所稱許的。 \end{tabularx} \\ \\
[1ex]
\hline
\hline
\end{longtable}
$^{1}$各位聽眾平安,我們繼續黎到言讀哥林多後書.
今日已經黎到第五講.
就係由第八章一直到第十章.
保羅好主要ge 就係話.
我地需要學習奉獻ge 功課.
咁有邊個榜樣呢.
馬其頓教會.
就係我地ge 榜樣啦.
我地學佢咁樣.
原來馬其頓教會好犀利呀.
呢間教會保羅讚到不得了.
佢話佢地係極窮之間.
仲顯出.
佢地樂捐ge 厚恩利.
一間教會自己好窮.
唔夠錢用架喇.
但係係極窮ge 裡面.
佢地仍然.
抽出一部分ge 奉獻出黎去幫其他ge 教會.
真係好難得呀.
我做jor 教會ge 牧師幾十年.
真係知道有時教會ge 財政.
會有.
不足ge 時候有D 時候頂住咪好樂意奉獻.
但唔知點解有D 月份.
可能要交稅gwo D 月份.
咁D 奉獻就會減少架喇.
見到教會D 收入唔夠喎.
咁有時自己都會擔心架.
咁點算呢出糧夠唔夠呢.
應用各樣日常ge 開支夠唔夠呢.
但係我地ge 神.
係一位信實ge 神.
佢必定會公認架.
我做jor 幾十年牧師.
我好知道我地ge 神.
係好實在ge .
佢從來唔誤事ge .
佢係最適當ge 時候.
就會將.

$^{41}$奉獻.
帶到黎比我地.
有時見到年底.
好似今年真係會出現赤字.
但係好奇妙.
到jor 12月.
D 奉獻又慢慢慢慢入黎.
於是又再足夠.
我做咁多年牧師.
未試過教會欠缺ge .
我地ge 神真係好信實好信實.
所以保羅做jor 呢間教會.
佢地咁樂於奉獻.
有幾節ge 聖經.
特別第八章第二節gwo 度話.
係患難中受大思念ge 時候.
佢地.
係仍然有滿足.
ge 快樂.
係大患難.
受思念ge 時候.
佢仍然有滿足.
ge 快樂.
係極窮之間.
仲格外顯出佢地.
樂.
捐ge .
厚恩.
呢個恩典真係好厚呀.
樂於奉獻.
原來我地奉獻第一個態度.
必須要樂於.
你笑住咁黎捐架.
我寫jor 張check十萬蚊架.
但我笑住黎奉獻.
因為呢個係你樂意架.
捐得樂意.
係神.
所喜悅架.
我地ge 奉獻必須要樂意.

$^{81}$唔係話呢個係我十一奉獻.
我一定要將薪水ge 十分之一拎出黎好似好唔捨得咁樣.
唔係.
讓我地學習.
奉獻.
必須要樂意.
係極窮ge 裡面.
我地仍然顯出.
樂捐ge .
厚恩黎.
仲有到第九章.
保留jor 誇讚.
亞該亞信徒gwo 個ge 熱心.
所以第八第九章保留都係讚.
有D 教會可以做我地ge 榜樣呀.
亞該亞教會有乜榜樣比我地學呢.
就係佢地ge 熱心.
係第七節下面gwo 度話.
因為捐得樂意ge 人係神所喜悅ge .
神喜歡gwo D .
你樂意去奉獻ge 人.
所以各位睇完節目我鼓勵你.
做一個樂捐ge 信徒.
每次奉獻.
你會樂意捐ge .
而且呢你每個主日.
有時我地好多教會全奉獻袋.
你係帶定黎.
唔係到時先打開個銀包.
抄D 碎紙啦係咁已經dump 落去啦.
唔係.
係你預備好ge .
所以樂意ge 奉獻.
係神所喜悅ge .
然後係第八節gwo 度話.
神能夠將各樣ge 恩惠.
係多多ge 加比我地.
使我地凡事都充足.
然後呢能夠行.
各樣ge 善事.

$^{121}$神係叫我地自己先充足先.
自己好夠啦.
然後我地好自然.
就會去關心gwo D 有需要ge .
所以我地就會行善.
做一個樂捐ge 人.
然後係第15節gwo 度就話.
神ge 恩賜神ge 恩典.
係我地講說不盡.
如果叫我數下神有乜ye 恩典.
我真係可以數好耐.
當你回顧過去你人生ge 歷史.
你所經歷ge ye .
真係神ge 恩典.
樣樣都要數.
我地事事都要數算主恩.
主ge 恩典.
係數算不盡.
所以哥羅莫多後書第八.
第九章.
都係保羅讚一些ge 教會.
讓一些個榜樣.
能夠激勵哥林多教會同樣.
做咁樣ge 事.
然後去到第十章.
就教佢地一個屬靈ge 功課.
唔好憑.
血氣.
黎到去爭戰.
有時我地人呢.
好鍾意憑自己ge 血氣.
憑血氣即係話.
我要做返比你睇.
我要贏過你.
我可以做得更加好,呢D 就係血氣.
因為你想有返個面.
咁所以要勉強自己去做.
唔係.
我地ge 爭戰.
唔係憑血氣.

$^{161}$第四節gwo 度話我地爭戰ge 兵器.
並唔係屬血氣.
你係係神面前.
有能力.
可以攻破.
gwo D 堅固ge 邢女.
有好多ge 邢女係好堅固ge .
靠住上帝ge 能力.
我地先至能夠攻破佢.
所以唔係憑自己一時之爭,我要出返啖氣.
我要出啖氣,gwo D 血氣.
然後係第五節gwo 度話.
神能夠將各樣ge 計謀.
各樣難阻人認識神ge gwo D 自誇ge 事.
一概.
攻破jor .
有時人好驕傲,諗住我要誇下口,你睇下我ge 成績幾犀利.
我比個成績表你睇下,我做過咩,我做過咩,我做過咩.
神唔係要我地睇呢D .
神要我地仰望佢.
倚靠佢.
我地ge 能力係從我地至高ge 神而黎.
所以我地仰望ge 係天上面ge 父.
佢係我地ge 幫助.
佢係我地能力ge 源頭.
我地知取ge 係從天上黎.
屬靈ge 能力,屬靈ge 資源.
又將人所有ge 心意.
奪回呀.
使佢歸於耶穌基督.
所以我地唔係憑.
自己ge 血氣.
去爭D 乜ye .
我地係倚靠神.
然後係最後第17節.
保羅好似做jor 一個conclusion.
一個總結咁.
誇口ge .
應當指住主.
黎到誇口.

$^{201}$如果你想沙塵呀.
嘩,你睇下我幾威.
你指住主.
誇口.
我之所以有今日.
如果我過去有D 乜ye 做得好ge .
唔係我有咩叻.
係天父.
幫助我.
主耶穌係我裡面.
主耶穌加我力量.
主耶穌賜我口才.
主耶穌叫我能夠講道.
所以我乜都冇.
如果唔係主ge 緣故.
我乜都做唔到.
所以我地如果要誇口.
就指住主.
黎到誇口.
主.
係我ge 榜樣.
主.
係我ge 能力.
主.
與我同在.
主.
祝福我ge 工作.
所以讓我地上上而隱藏.
係基督裡面.
基督.
被顯偉大.
我.
必須明.
祂必興旺.
讓我學習施浸約翰.
那個謙卑ge 態度.
佢指住耶穌.
佢呢個係.
神ge .
羔羊.

$^{241}$除去世人罪孽.
所以施浸約翰佢話.
祂必興旺.
我必遂明.
讓我地誇口ge .
指指住主.
誇口.
求主幫助我地.
讓我地學習一D ge 教會.
成為我地一個好好ge 榜樣.
又讓我地.
上上冇憑血氣去爭戰.
單單仰望主.
倚靠主.
主成為我地ge 幫助.
主成為我地ge 能力.
願天父.
賜福比各位.
\newpage



\section{歌羅西書 1:3-8}
\label{sec:031604}
\textbf{明白真正神委任的牧者心腸}~~~~黃紹權牧師
\newline
\newline
連結: \href{http://www.church.com.hk/acms/content.asp?site=cdc\&op=show\&type=product\&code=031604\&layout=sermon}{\texttt{031604}} ~~~~ 語音日期: 2026-02-08
\newline
\newline
\hyperref[sec:031603]{\small{< < < PREV SERMON < < <}}
~
\hyperref[sec:index_chronic]{\small{[返順時目]}}
~
\hyperref[sec:031605]{\small{> > > NEXT SERMON > > >}}
\newline
\newline
歌羅西書 1:3-8
\newline
\begin{longtable}{cl}
\hline
\hline
章節 & 經文 (和合本修訂版)\\
\hline
1:3 & \begin{tabularx}{0.7\textwidth}{X} 我們為你們禱告的時候，常常感謝我們主耶穌基督的父神， \end{tabularx} \\ \\ \relax
1:4 & \begin{tabularx}{0.7\textwidth}{X} 因為聽見你們對基督耶穌的信心，並對眾聖徒有的愛心。 \end{tabularx} \\ \\ \relax
1:5 & \begin{tabularx}{0.7\textwidth}{X} 這都是因著那給你們存在天上的盼望，它就是你們從前所聽見真理的道，就是福音； \end{tabularx} \\ \\ \relax
1:6 & \begin{tabularx}{0.7\textwidth}{X} 這福音傳到你們那裡，也傳到普天下，並且繼續增長，不斷結果，正如自從你們聽見福音，真正知道神恩惠的日子起，在你們中間也是這樣。 \end{tabularx} \\ \\ \relax
1:7 & \begin{tabularx}{0.7\textwidth}{X} 這福音是你們從我們所親愛、一同作僕人的以巴弗學到的。他為我們作了基督的忠心僕役， \end{tabularx} \\ \\ \relax
1:8 & \begin{tabularx}{0.7\textwidth}{X} 也把聖靈賜給你們的愛告訴我們。 \end{tabularx} \\ \\
[1ex]
\hline
\hline
\end{longtable}
$^{1}$徒長千輝,讓我們在聆聽主的說話,意思就是同心低頭禱告..
父神多謝你,讓我們再次回來,在你的說話的根前,一同去仰望在你的裡面..
求你的靈與我們同在,帶領我們能夠在今天,在你的啟示當中,讓我們能夠更準確,拿捏..
上帝你怎樣去牧養我們,要讓我們怎樣去辨明,真正在上帝你裡面的教導,.
以及我們應當帶著一份甚麼的心智,在你的當中去事奉..
求你的帶領幫助,讓我們能夠一同在聆聽主的說話裡得以成長,.
好使我們能夠更加,更尊敬上上帝你..
多謝你給我們這段時間,無論善上善下的每一個,都讓我們的身心靈在聖靈引導之下,.
能夠去看到怎樣走這條屬靈的人生路,禱告奉靠主耶穌基督,得勝名自求,阿們..
在我牧養這麼多年的裡面,有不少不是以馬來華人浸信會的會友,都曾經和我聯絡過..
我想這個可能是上帝給我在事奉當中的其中一個主要的環節..
因為很多這些弟兄姊妹,雖然他們不是以馬來的會眾,.
但是透過網絡的事供,他們能夠和我聯絡得到,或者認識得到,然後主動來聯絡..
他們無獨有偶都分享到一件事,就是說他們在靈性的上面,曾經經歷過不同程度上的傷害..
他們對我很多時候的發問都問到一個問題,.
就是我怎樣才找到一個真正能夠是上帝所委任的牧者來牧養,或者帶領我們一同事奉呢?.
我覺得這條問題可以說是有點奇怪,.
因為上帝給我們不是因為我們要去跟某一個人來事奉,.
但是事實上他們所留意到的是一個真真實實的情況..
很多時候在教會當中,我們見到的牧者,大大小小不同的堂會的牧者,.
有出名的,有不出名的,有些搞到自己很出名的,有些不出名的,.
都利用很多各樣的方法去吸引人..
究竟怎樣才辨認到一個傳道人,甚至簡單到一個基督徒,.
是不是真的有上帝的呼召,或者有上帝的委任呢?.
這個題目我相信不單止很多這些弟兄姊妹有的問題,.
或者可能甚至是你的問題,連我自己在事奉的過程當中,.
都不斷地去問我自己,上帝你的呼召我能不能夠準確去拿捏..
因此,去辨別一個真正的傳道人,.
或者一個真正的信徒,其實的確是需要一些原則性的教導來讓他們明白,.
辨認清楚,好使他們所聽的道理,能夠知道怎樣去接受..
有我在這裡加插一點點,學習真理是不一定由一個很好的傳道人身上才學得到的..
你說牧師這句話何以見得呢?.
在聖經裡面,甚至我們知道撒旦都講聖經的,.
一個成熟的基督徒,或者有上帝呼召的基督徒,.
其實他們是能夠辨別到錯謀的教導的,.
即是說是撒旦傳異道,異樣的道,不是上帝的道,.
但是用聖經的說話來包裝也好,信徒是聞到那一股味道的..
我記得在神學院的時候,特別我去學到所謂經文批判學,.
一個被譽為很新派的神學,識經技巧的一個環節..
我問到我的教授,我的教授他八十幾歲了,.

$^{41}$我說你用了八十幾年,因為他從出世沒多久,.
所以他父母已經將他放在一間修道院裡面,去學做一個傳道人..
他在修道院裡面長大,我問他,.
你這八十幾年裡面,你怎樣學會辨別你做的識經是合上帝的心意?.
他沒有說過答案給我聽,他只是拿起手指一指著鼻子,然後他吸了一口氣,.
他說我一聞就聞到了,其實當他這樣說的時候,.
就讓我想起我有個同學的珠寶鑒證專家,.
或者可能對珠寶,金器有研究的人都會知道,.
其實你怎樣才辨別到真的珠寶,真的金銀呢?.
其實當你長期對著真的時候,你一摸下去,一看下去就知道哪些是真,哪些是假..
所以今天在這段聖經裡面,保羅特爾在透過歌羅西書裡面,.
首先不單止將自己介紹給他知道,並且同時間告訴哥羅西的教會知道,.
你應該怎樣去辨別真正有神所委任的牧者,.
而從這個委任開始的辨明,你就可以知道應該怎樣去接收他們所說的話..
或許他未必真是神的僕人,你辨別得到,他是一個假的傳道人,.
但是你也能夠從他辨別他講假的道理的過程當中,.
能夠更準確知道上帝真正的啟示是什麼..
為什麼我覺得這段聖經是很值得我們去思想,.
因為尤其是在這近二十多年,我不敢講其他世界國家的地方,.
我只講在多倫多,我認識在多倫多現在有很多很多很多的聲稱自己是自由傳道宣教的牧者,.
流傳在宗教會當中,你不要問我他究竟因為什麼原因流傳在,.
隱藏在不同的教會裡面,但是當他們傳流在你和我身邊的時候,.
我們應該要學會一件事,能夠辨別到究竟他是傳真正的道理,.
還是只不過是混在其中想吃教飯的呢?.
更可怕的可能是今天很流行的一個說法,是不是一種滲透呢?.
這些等等的事情,讓我們去好好辨別,從而準確去領受神的教導..
首先讓我們去明白保羅他給了一個很重要的因素去辨別一個真偽的牧者,.
就是一個人,他縱使他稱為自己是一個傳道人,.
他究竟他有一種什麼的心腸是在他的事奉當中呢?.
所以今天和大家一起講道的題目,也是叫做明白真正神委任的牧者的心腸,.
透過認識他的動機,來辨別他的真偽..
好,讓我們回到歌羅西書第一章第三節第八節,我們一起再讀一次,好不好?.
好,找到之後我們一起去讀了..
好,我們一起讀..
因輕見你們在基督耶穌裡的信心,並向眾星圖的愛心,.
是為那給你們存在天上的盼望..
這盼望就是你們從前在福音真理的道上所聽見的..
這福音傳到你們那裡,也傳到普天下,.
並且結果,增長如同在你們中間..
自從你們聽見福音,真知道神恩惠的日子一樣,.

$^{81}$正如你們從我們所親愛,一同作僕人的已不可學的,.
他為我們作了基督中心的執事,也把你們因聖靈所傳的愛心,.
告訴了我們..
今天我會分兩個部分和大家一起去看的..
不過先看回上文第一召第二節,.
有不少的弟兄姊妹說,那些都是書信的慣常開頭,.
牧師又說成過中,其實想說什麼呢?.
其實我想再重覆說多一次,.
歌羅西書雖然是很短一卷的書信,.
但是他卻是有一位已經年事已高,.
但仍然關心著和他沒有直接關係的哥羅西教會的保羅,.
來為他們提供最適切的牧養..
因此保羅在整個開端裡面,.
我們會留意到與其他保羅書信自然不同的地方,.
就是保羅是沒有任何的陳述,是希望有什麼謀求在教會裡面,.
反而他的目的是要幫助哥羅西教會的信徒,.
在日後的發展裡面能夠修正他們那條道路,.
使他們在屬靈的道路上知道怎樣去走,.
專心仰望的是上帝,而不是人..
所以保羅這份心思,.
讓我們看見好像昔日但以理一樣的心態,.
對於猶大國的子民的未來,.
即使自己沒有事做得到,.
但是他仍然在末後的地方,.
能夠將他作出最大的一個關懷,.
去使到將來的猶大國的子民知道怎樣走..
而第一章的第一節,我們看到保羅稱自己是神旨以下委任的仕途,.
因此哥羅西教會怎樣才能夠真正辨別到神委任的仕途呢?.
不是因為那個人他講得多漂亮,.
或者他口才多好,.
就去斷定這個一定是有恩賜的神委任的牧者來到我們當中..
保羅沒有用這些說話,.
反而我們在第一章第一節和第二節裡面,.
我們充分看到保羅是毫無任何自己的動機,.
有什麼索求而去寫這封書信..
在這個大前提之下,進入到第三至第八節的經文,.
第一部分我們去看的,.
就是從第三節至到第六節的上半節,.
有一班充滿信望愛齊全的哥羅西教會的聖徒..
請聽清楚這個標題,.

$^{121}$我們看回開始第三節,.
我們主耶穌基督的父,.
常常為你們禱告..
所以一開始他把焦點放在上帝的身上,.
為什麼我這麼強調這一點呢?.
曾經有些信徒問過我,.
牧師,你平時在教會辦公室做什麼?.
你有沒有這條問題?.
你想問我在教會辦公室做什麼?.
是不是在睡覺呢?.
還是在上網看Netflix?.
還是上網看一些無無畏畏的東西?.
是不是做這些事情呢?.
保羅在這裡提供了一個答案,.
一個牧者,.
第一件必然在教會裡經常做的事情,.
就是為你們禱告..
有些信徒覺得禱告而已,.
不用做了,.
坐著閉上眼睛,.
說說話,.
是不是這麼簡單的事情呢?.
或許很多信徒的祈禱是這樣,.
或許甚至有些牧者的祈禱真的是這樣,.
但真正我們留心去看看,.
保羅在這裡講到一個牧者,.
是為信徒怎樣禱告呢?.
主耶穌基督的父..
其實弟兄姊妹你祈禱,.
有時你的禱告生活是怎樣,.
其實我不多不少都會知道的,.
當你一開口祈禱我就會知道的,.
有很多信徒祈禱,.
全部的禱告內容都是一些屬靈的慣用詞,.
你無限偉大,.
上帝你是無所不在,.
無所不能,.
我依靠你,.
我委身在你裡面,.
你知不知道當你這樣講這些禱告,.

$^{161}$發出這類型的禱告,.
其實我都知道你都是平日的祈禱生活,.
是沒有留意你自己的屬靈狀態,.
聽清楚這句話,.
必然是清楚知道哥羅西教會,.
我們一會兒會詳細去看,.
不是空洞地講幾句話,.
我作為一個牧者上上,.
雖然我們教會沒有辦公室,.
我是work from home,.
但是每次每一日,.
我的生活都是為了弟兄姊妹,.
你的靈性狀態去禱告,.
要理解你的狀態是如何,.
而去預備我接下來其他的事情,.
包括講道,.
歌羅西書不是偶然,.
講完舊約隨便找一卷新約來講,.
不是這樣,.
而是歌羅西書的講道是承接,.
因為內容承接,.
但以理書裡面,.
那位牧者好像但以一樣的神的僕人,.
怎樣去關懷他的信徒,.
關懷他的子民,.
所以當我們即使作為傳道人,.
平日在教會裡面,.
辦公室也好,.
我work from home也好,.
我們的禱告其實都是為了弟兄姊妹,.
你的個別特別的需要去代禱,.
我覺得有一位弟兄,.
我曾經很欣賞他的代禱,.
他走來問我牧師,.
那個弟兄或者那個姊妹,.
曾經怎樣怎樣,.
他近況如何,.
你有沒有去update,.
我想為他代禱,.
弟兄姊妹,.

$^{201}$因為他關心弟兄姊妹的靈性狀況,.
所以他知道一點的時候去代禱,.
而當他代禱的時候,.
他想更知道更準確,.
為他必定的需要而去代禱,.
所以他主動去找多些資料,.
來為他更作仔細的禱告,.
這個就是保羅的禱告的第一個重點,.
他要明白到教會的狀態需要,.
來去按需要禱告,.
並且調整他的教導,.
或者怎樣去牧養教會,.
就好像一個媽媽一樣,.
我不知道大家現在怎樣去做媽媽爸爸,.
我們在加拿大通常的做法就是,.
買菜煮飯,.
就看看超級市場減價是在賣什麼,.
但除此之外,.
我們往往做父母的買菜煮飯,.
其實是在想子女或者家人,.
他們有什麼健康需要,.
譬如最簡單,.
如果我見到我自己的子女,.
出滿了暗瘡的話,.
我就很自然去想,.
去超級市場不是買看今天特價是什麼,.
而是去看看有沒有貴妃展,.
貴貴的貴妃展,.
買回去煲花旗參,.
希望幫他們清清熱,.
做牧者的心腸就好像做父母一樣,.
人們經常說我不懂做父母,.
因為你沒有關心過你的子女,.
你當然不會做,.
但當你肯關心他的時候,.
你必然會知道他的需要是什麼,.
由他小時候,他剛出生,.
他哭你就會去看看他的紙尿片是不是大了,.
是不是肚子餓了,.
你會按照他的需要去做餵養,.

$^{241}$同樣做牧者一樣,.
為他們禱告,.
因為保羅要知道的是他們真正的屬靈需要是什麼,.
什麼屬靈需要,.
第四節至第六節上半節就詳細記載了給我們知道,.
他說:因我聽見你們在基督裡的信心,.
並向眾聖徒的愛心,.
是為那給你們存在天上的盼望,.
這盼望就是你們從前在福音真理的道上所聽見的,.
這福音傳到你們那裡,.
也傳到普天之下,.
並且結果子增長,.
如同在你們中間..
這句聖經說什麼呢?.
這句聖經其實是保羅原來收到消息,.
聽到哥羅西教會他們那個狀況,.
保羅從哪一個人聽到他們的狀況呢?.
如果我們偷偷看,.
就可以知道原來在第七節那裡介紹到有一個人物,.
叫做爾巴弗,.
原來保羅有一位童工叫做爾巴弗,.
爾巴弗向他報告,.
哥羅西教會的眾聖徒的狀態如何,.
這些不是叫小報告,.
這些叫做對教會的關懷理解,.
而保羅去問及一位童工,.
為什麼要問爾巴弗這位童工,.
因為爾巴弗就是哥羅西本地人,.
他問一個本地的童工,.
老遠來到保羅的身邊,.
他信了主,.
知道他們教會的狀況,.
保羅見到哥羅西教會他們的情況就非常之喜悅,.
喜悅在什麼呢?.
因為看到他們在基督耶穌裡有信心,.
並且向眾聖徒有愛心,.
而且是為了他們在傳在天上的那份盼望,.
而繼續在這份盼望上繼續努力,.
我們去分析一下,.
或者探究一下,.

$^{281}$這幾樣東西是什麼,.
第一,在基督耶穌裡的信心,.
這是表明保羅清楚知道哥羅西教會的信徒,.
他們沒有離開基督耶穌這個核心拯救者的救恩來源,.
他們所得到的結果是完全信任在基督耶穌裡面,.
或者用另一個翻譯的方法,.
就是你們的信心在基督裡,.
耶穌基督給你們的信心不是任何其他地方的信心,.
而是有耶穌存在其中的,.
簡單來說,他們有的信心是有耶穌同處的信心,.
或者說得更直白一點,.
就是跟耶穌的信是相同的,.
耶穌的信是什麼呢?.
就是對上帝父神的完全信賴,.
為什麼我們知道呢?.
因為在之前第三節不是說了,.
耶穌和父神之間的關係是父與子之間的關係,.
兒子不會懷疑父親,.
而且兒子還要為了父親所期望的事情,.
去願意犧牲父上自己的生命,.
碰巧今天我爸爸生日,我做兒子的都會為爸爸慶祝生日,.
雖然不在身邊,人到都錢到,.
為他付出一些我能夠付出的東西,.
希望父親喜悅的事,.
生活在西方社會內那裡,buy a lunch,buy a big dinner,.
請他吃頓飯都可以,.
耶穌基督對父神的心就是這樣,.
為了父神拯救人,.
恢復人離開罪,.
這個計劃,這個心願,.
他願意將自己的生命放上,.
完全無保留地去為父神付出,.
而哥羅西教會你們就是有一個這樣的心..
第二,並向眾星圖的愛心,.
請留意,這句說話講到愛心,我們就將它騰空講到無限大,.
愛心呀,我們要用愛去接待人家,.
對,但是請留意,向眾星圖的愛心是什麼呢?.
在這裡其實就講到在教會裡面的事奉的問題了,.
哥羅西教會的信徒,.
就是為了幫助在教會其他的人,.

$^{321}$來讓他們更加的明白上帝的心思意念,來參與事奉,.
這個所謂愛心,其實就是為他們最大的益處的考量,.
不只是討好,很多信徒誤會,.
愛心就是說什麼都不要追究了,有錯就不要講了,愛心呀,.
聖經其實愛心這個說法,.
是包括指督責的功效的,.
你想那個人好,你才會去提出他的問題,.
所以弟兄姊妹,.
哥羅西教會原來他不單止在耶穌基督身上,.
有同向父神的心一樣,.
他向著同樣自己教會裡面實在的每一個會眾,.
他們都是以愛心相待,希望大家能夠同樣在上帝裡面,.
得以成長成熟,之後會講的,.
但是這個行動,保羅就看見,.
原來哥羅西教會真是做到這麼美好的見證,.
我們繼續看,.
為什麼哥羅西教會他們會這麼著力去做呢?.
原來他們看見他們已經在天上有一份盼望,.
存留在當中,他們的目標不是只是在地上討好人,.
或者教會裡面,大家你好我好,.
而是希望大家向著天上的盼望而走過去,.
而這個盼望是從福音的真理他們早在之前聽過而來的,.
而這個福音就是耶穌基督所提造出來,.
因此,哥羅西教會從保羅的觀察裡面所聽聞到的就是,.
原來他們是很努力地帶著所有整個教會的群體,.
來朝向上帝修和的關係去進發,.
所以第六節上半節這樣說,.
「這福音傳到你們那裡也傳到普天之下,.
並且結果只增長與同在你們中間.」.
這裡所說到的,原來你們是繼續希望這個福音能夠在你們裡面不斷膨脹地成長起來,.
有兩個很重要的字眼,.
「結果」和「增長」,.
在希臘文裡面這兩個字是福音的特性,.
福音的特性是什麼呢?.
是要讓每一個得到福音的人,.
他們能夠在福音的栽培之下去結出更多似上帝的生命的果子,.
並且能夠multiply不斷繁衍..
這兩個字眼其實就等同於先前哥羅西教會他們的眾聖徒的愛心的行動,.
他們去傳這個福音,.
在教會裡栽培弟兄姊妹更加成熟,.

$^{361}$無非都是想大家更似這個福音的本質,.
他自己會不斷繼續成長,.
不是死的,.
是會成長的,.
而且他會不斷的multiply,.
會繼續的去生長更多,.
他自己不單自己會長大,.
而且他會更加延伸去其他的地方..
所以從這個經文裡面,.
我們看見保羅見到哥羅西教會有這麼美好的成長,.
這麼似福音的本質,.
亦都履行上帝好似耶穌履行父神心意一樣,.
來到不斷成長,.
請留意,.
他沒有讚許這班哥羅西教會的信徒,.
你們有沒有留意?.
有沒有任何一句說話去讚賞哥羅西教會,.
你們厲害啊!.
弟兄姊妹,.
這種情況正正是違背了今天很多教育的觀念,.
甚至教會很多牧者的所謂牧會觀,.
我們在OTS讀書的時候,.
安省神學院,.
安省神學院原本很好,後來改名為天頂,.
開始有變質了,.
其中一位牧者教牧會學的,.
在多倫多環教會當中非常出名的一位牧者,.
那位我不告訴你了,.
你不需要知道,.
他公然在神學院教牧會的時候,.
他說成功的牧會你們想不想,.
全班的同學都說很想,.
他說成功的牧養就是要使教會的人開開心心就行了,.
他說了之後我覺得非常之嘩然,.
但很多同學很同意,.
我們回來看看保羅,.
保羅和哥羅西教會都不相識,.
沒有直接關係,.
但是保羅雖然沒有直接關係,.
但是他現在寫一封信給哥羅西教會,.

$^{401}$目的何在?.
搭關係?.
如果他是搭關係的,.
他應該在信裡面寫,.
哥羅西教會你們真是頂呱呱,.
在宗教會之間,.
最嘻哈那一間,沒有哪一間比你們對耶穌基督的福音更理解,.
你們又有愛心,又彼此接待,.
你們又彼此怎樣關懷,.
我現在怎樣,我落難,.
遲些他會說的,保羅,.
但是他有沒有在這個時候用這種方法去拉關係爭你們?.
我還沒說完,.
剛才那個著名牧者,.
他死了,你們不用找他,.
他死了,你們不需要找他,.
不過那些門生眾多很恐怖,.
他說除了要會友開開心心,.
你要有空多讚他們,.
有個同學,.
情不超那些,好像我這些,.
就去問他,.
點讚呀 點讚呀 你就要好好利用你的「港台」.
「港台」? 「港台」就是宣講神的說話.
他說你記住呀 我們的華人教會講道就叫「Free Point Sermon」.
什麼都Free Point.
我就沒有的 經文講什麼我就照講的.
分開三步 三步 四步 四步 五步 五步 兩步 兩步 一步 一步.
他說不是的 三步.
為什麼三步?原來有原因的 你知不知道?.
你們將來出去聽講道你看看.
我們Free Point Sermon 三點式 不是那些榮譽.
三點式的講道.
你要首先用頭兩點讚許教會的信徒.
我說很難的 糟了 暴露了我的身份問這條問題.
我說很難的.
他說你沒有辦法 你都要找些東西來讚他.
然後到第三點.
縱使是罵都好.
你都要將那個罵的程度降到最低.

$^{441}$用字就要很小心.
不要譴責 不要直斥其非.
你們要做什麼?.
要有進步 要有改善 要有更成聖的機會.
這些字眼 你在教會經常聽牧者這樣講.
你聽完以為很順耳都好 其實他在下毒.
保羅在這裡沒有讚過哥羅西教會.
縱使做得這麼好都沒有讚過.
曾經有位姐妹在網上寫了一篇文章.
然後發給我.
她說黃牧師我在你身上學到什麼什麼.
我跟她說 我只是做我的職責.
但我都很多謝你.
多謝了神就多謝了我.
我只是做我的職任 上帝交給我的功夫的部分.
為什麼要這樣做?.
因為不要透過上帝給你的恩賜.
去成為偷上帝榮耀的機會.
我重複再講多一次.
千萬不要用上帝給我們的恩賜.
去偷上帝榮耀的機會.
你縱使識經多好 多厲害都好.
什麼人在你身上都學到上帝的話都好.
那個感恩是上帝應該配得的.
上帝給我的恩賜 我在這個職任上承擔多一點.
所以為什麼聖經常常這樣說.
多給誰 向誰多取.
就是這個原因.
千萬不要用上帝給我們的恩賜.
去偷上帝的榮耀.
哥羅西教會做得好 我都見到.
保羅說我聽到了.
但是怎麼樣?.
所以弟兄姊妹 你們有沒有留意.
在教會的事奉.
我至少和你們做事奉的時候我會說多謝.
因為你不是做給我.
你是做給上帝.
你搬一張椅子也好 去接一個程序表也好.
接觸一個新來賓也好.

$^{481}$全部是你做給上帝.
所以教會的敬拜是什麼?.
我不是說一定弟兄姊妹要實體崇拜.
但是如果你是來教會實體崇拜的時候.
你能夠做到的事 是向上帝去做.
所以教會的崇拜叫做什麼?.
英文叫做Worship Service.
你是來服侍上帝.
你不是來買一個服務.
保羅在這個時候.
是哥羅西教會的靈性與福音.
有一致性的質素和果效.
因為你們同樣的結果只會增長.
所以請大家留意.
在這裡到第六節最後有一個字眼.
一會兒我們會再重複再說.
你們有一樣東西與上帝的福音是一樣的.
這個字眼.
「一樣」這個字眼.
其實在接下來的經文.
用了中文 中文就真的很難搞.
中文聖經有時會用「一樣」.
有時會用「一同」.
「如同」.
其實是說給我們知道.
保羅欣賞到.
哥羅西教會與福音完全配搭一模一樣.
能夠達成父神帶領人的心願.
能夠讓人在天國裡有盼望.
所以保羅在聖經裡.
是希望能夠與哥羅西教會繼續成長.
流露出一種.
大家一起專心仰望在上帝裡面.
上帝所喜悅的事.
你們增長.
我們又增長.
一起在上帝裡面增長.
這個就是福音的果效.
這個也是為什麼教會今年的主題.
在哪裡呢?.

$^{521}$歌羅西書第三章第四節說.
「與基督同顯他的榮耀」.
以禱告理解教會的需要.
我們看看保羅.
從第六節的下半節.
一直到第八節.
他怎樣將福音的本質.
也與他作為一個牧者的心腸.
將它連結起來.
我讀給大家聽.
我們從第六節讀完.
我們繼續讀下半節.
第六節這樣說.
他為我們作了基督中心的執事.
也把你們恩聖靈所傳的愛心.
告訴了你們.
保羅在這裡.
是要將福音真正落實到.
每一個信徒的身上.
所以當我們回顧看回.
第一章的第四節裡面.
保羅說你們都做到.
向眾聖徒彼此發出愛心.
而我們看回第一章第八節最後.
保羅也說你們是恩聖靈所傳的愛心.
告訴了我們.
怎樣的情況.
保羅一方面看見.
哥羅西教會的愛心的事奉.
但同時間也都恩著他們的愛心.
讓我們去印證到.
有聖靈存在在我們當中.
保羅基本上可以說是完全體會到.
上帝怎樣在哥羅西教會裡面去做事.
那我怎樣去做事呢.
所以第七節.
第六節說了原來這個福音的本質.
有結果和增長的功效.
如同在你們中間.
已經在你們裡面了.

$^{561}$看不看到.
在你們裡面了.
現在我真的從樣子看見.
我真知道神的恩惠的日子是一樣.
什麼一樣.
你們那個生命的增長.
就是落實在你們平日的每一天.
上帝和你們同坐的日子.
所以正如又是同一個字眼.
這裡已經同樣用同一個希臘文.
叫Kaphos.
你們的生命有增長.
有結果.
福音的日子一路發展下去.
又繼續和你們一樣.
上帝和你們同坐.
而又正如你們從我們所親愛.
一同僕人的以巴弗所身上學到的.
原來連保羅的同宮.
都同樣在裡面和你們一模一樣去事奉.
你看不看到這個連貫性.
保羅.
哥羅西教會.
以及他的同宮.
以巴弗.
全部人都在增長中.
他成為了保羅的中心的執事.
活在你們的當中.
換言之.
保羅所看見的整個哥羅西教會.
以及他自己的關係.
透過以巴弗的連結.
都大家同樣都在這個福音裡面.
一起去努力.
不過在這裡有一點值得我們去留意的就是.
保羅沒有因為自己不是直接建立哥羅西教會.
而親疏有別.
保羅他.
即使哥羅西教會.
不是由他自己直接報道去傳道都好.

$^{601}$只要是他的同宮能夠去不斷的延展.
去不同的地方建立教會.
因為.
因為以巴弗其實就是哥羅西教會建立其中一個人士.
這裡說了.
他在你們當中作為中心的執事.
他在你們當中去做這份工作.
保羅只要見到他所帶領所信了主的以巴弗.
能夠擴展到上帝的福音工作.
去到哥羅西教會.
保羅同樣寫這一封的歌羅西書.
按照他們的需要來去模樣.
換言之.
保羅的事奉從來都沒有一種叫做親自部署.
親自指揮.
也都不需要親自執行.
只需要福音的工作.
在有福音生命的人身上.
能夠擴展的話.
保羅都會將他的模樣延伸到那個地方.
我再重複講多一次.
保羅的事奉官不是因為.
你是不是我帶領信主.
然後你才得到我的模樣.
只需要我傳道給任何一個人.
那個人繼續將這個福音擴展下去.
好像福音本身自己的本質會不斷擴展.
會不斷結果的.
我保羅雖然隔了以巴弗都好.
跟哥羅西教會你們不是直接認識.
我仍然將最好的模樣帶到你們當中.
什麼人才是樣樣東西都要親自部署.
親自指揮.
親自去執行.
這種人叫做control freak.
控制狂.
武神生命的人就是這樣.
我回目了這麼多年.
曾經有一個老牧者這樣跟我說.
Ellen 你很難目的.

$^{641}$以馬來尼.
我說為什麼.
他說因為以馬來尼環振信會.
不是你親手模樣帶他們信主的信徒的會友.
我覺得很奇怪.
他說這番話.
他說如果你教會的會友.
絕大部分都是你親身帶他信主.
親身為他施浸.
他就會更加信服你.
你就容易目會.
我聽完他這番話之後.
我覺得這位牧者可以丟掉.
他還用了一個很人性的描述.
你的會友不是你的親生子女.
你怎樣目樣他.
我跟著我想頂他一句.
你既然這麼講人性的話.
我講一句人的說話給你聽.
你知不知道.
養良不及生良親.
不過我沒有這樣串他.
因為給兩分薄面這位老牧者.
生良.
怎樣.
生良.
我不是這樣講嗎.
所以你知道我今天什麼狀態.
生良不及養良.
我跟著問.
這樣是不是真的不用目.
他說你要不就趕走你現在那群.
自己帶過一群信主.
你就目到.
不過我看你都很難的了.
弟兄姊妹.
保羅有沒有這樣的心態.
我又不認識他.
他又不認識我.
他認識的誰.

$^{681}$現在我們知道了.
他認識的是以巴弗.
因為以巴弗是保羅帶他信了主.
然後他回到他自己的哥羅西城.
裡面去傳道.
建立教會.
作你們的執事不是我說的.
但是他願意在你們當中.
去使你們有這麼美好的見證.
而且我希望藉著這封書信.
來去牧養你們.
因為我仔細聽到.
所有你們的境況是怎麼樣.
弟兄姊妹.
我在開頭的時候不是講過.
有很多弟兄姊妹.
和我有接觸的.
其實我.
我真的不認識他們.
但是上帝用了.
我能夠能力做到的事.
譬如做broadcast.
最近我才和一個弟兄聊天.
全球我可以這麼說.
華人教會做網絡的廣播的.
以馬萊利是其中一間.
不是那些大教會.
我很記得1996年的時候.
我回來的時候.
來到加拿大.
我已經計劃去這麼做.
透過認識上帝的說話.
上帝用了這麼多的媒介媒體.
你又看看以馬萊利環浸信會的網站有沒有宣傳.
沒有的.
我又沒有賣廣告.
我們教會現在連報紙都不賣.
新來移民的人都不用在報紙找到.
以馬萊利環浸信會在哪裡.
因為沒有賣廣告.

$^{721}$但是來得以馬萊利環浸信會參與的.
每一個弟兄姊妹.
無論先上先下的.
我都希望你能夠在上帝的道理上得到造就.
當你在靈性上遇見挑戰.
我作為一個牧者.
不是因為你是不是我們的會友.
不是因為你有沒有實體崇拜.
不是因為你有沒有奉獻.
而是我真的希望的是.
我在主的裡面.
上帝給我的恩賜能夠服侍到每一個弟兄姊妹.
所需要的靈性需要.
我不是服侍你.
我是服侍你的靈性需要.
因為這個是上帝交給我的功夫.
保羅母親疏有別.
離去牧羊.
保羅相反的是在這裡讓我看見.
也把你們因聖靈所傳的愛心告訴了我們.
激勵了保羅回頭.
所以真正的牧者.
他的心腸是什麼呢.
只要信徒得到成長.
他不管你究竟是哪一個教會.
哪一個宗派.
只需要靈命成長就是他的目的.
我知道有些弟兄姊妹聽完.
《馬來華正宗》我講到《釋經》之後.
他會對於他自己的教會更加掙扎.
因為他越來越發覺.
他自己的教會越來越不像《聖經》裡面所講的教會.
但是讓我在這裡講多一句.
上帝用他的說話打開了你的心眼之後.
不是叫你去改革你的教會.
而是讓你在所處的教會.
在所處的環境裡面.
活出一個更加像上帝有成長.
會multiply 會結果子的基督徒.
讓人更加紮實地認順上主.

$^{761}$這個福音和這個福音本質是相接連.
當你願意這樣做的時候.
有一次我接觸到一批.
我完全從來不認識的信徒.
後來當我對他們多了點認識之後.
我就發現到有一個問題.
原來在他們當中有很多弟兄姊妹是.
叫做信徒或者甚至受了洗.
但是他們沒有信主.
他們沒有經過認罪悔改.
看不到自己的罪是什麼.
因為教會為了當時要組織一個教堂登記.
所以他們就成為了教會的創堂會友.
後來我去到這間教會認識這班弟兄姊妹.
多了認識他們之後.
我覺得這是一個很大的問題.
問題是怎樣可以幫到他們.
回歸正當的福音.
好像這裡所說.
他們可以真的結果子.
結的是上帝喜悅的果子.
是可以領受天上的盼望的果子.
並且可以multiply.
於是乎我就跟他說.
我留意到你們教會裡面有一個情況.
有一個人物很有影響力.
但是他的信是很有問題.
先停了他的事奉.
慢慢再處理好他的靈性.
當我提出了這一點之後.
原來他們教會全部的長輩.
全部教會弟兄姊妹知道這件事.
老早留意到這個人.
其實在教會裡面做很多非法的事情.
但是因為他在教會的影響力.
他的影響力有很多種方法.
在這裡不詳細說了.
但是其中一樣東西.
他能夠幫到很多人做翻譯.
所以很多人要靠他.

$^{801}$不敢出聲.
全教會全部人都知道他的問題.
沒有一個人出聲.
為什麼一間教會裡面的人.
不會一起處理這個罪的問題呢.
哥羅西教會就不會這樣.
他會全部一起努力的.
因為他們怎麼樣.
他們會一起祈求.
可以在天上面得到.
儲備到這個盼望.
當我這樣說之後.
教會說好.
黃牧師你叫停了他的事奉.
就停了他的事奉.
我說停了之後再慢慢處理這個弟兄.
我和這個弟兄處理.
他完全不認為他自己有問題.
他認為他自己所作的事情.
完全沒有問題.
就算是違法都好.
他無所謂的.
違法不是違聖經的法律.
違社會的法律.
不要說道德層面.
道德更加混亂.
我說如果你不悔改的話.
我沒有辦法.
在教會讓你繼續事奉.
全教會的人好像很高興.
好像有人終於處理這個罪的問題.
但誰不知.
不到一年.
這間教會因為這個弟兄的影響力.
於是乎再次給他在教會.
最後成為教會的主幹事奉者.
鄧小慕你知不知道這件事.
在我心裡面我掙扎了多久.
掙扎到現在.
一班很好的弟兄姊妹.

$^{841}$有真信的弟兄姊妹.
為什麼不能夠面對他們的罪.
哥羅西教會有沒有這樣的問題.
已經很好了.
哥羅西教會.
哥羅西有他自己另外一隻的問題.
所以為什麼保羅要繼續寫下去.
但牧者的心是什麼心.
就是希望你們能夠.
好像福音的本質一樣.
成長起來.
multiply起來.
不斷擴展.
而不是死咕咕地動在家裡.
或者動在你的工作環境裡面.
毫無影響力.
我說的影響力不是去搞什麼改革.
而是別人看不看到.
你是一個有神生命的人.
一個真正的牧者.
有神的委任.
他一定有一個特徵.
就是要讓他的教會的弟兄姊妹.
他所認識的.
能夠朝向上帝的方向去改變.
改得越來越像上帝.
他所生出的影響力.
也是越來越像基督耶穌.
那個福音的本質.
讓人回到舊恩裡面.
期望在舊恩裡面有份.
更重要是期望和上帝可以走在一起.
弟兄姊妹.
你和我的崇靈生命又如何呢?.
明白一位牧者的心腸是重要.
但是更重要的是.
你能不能夠在明白過後.
你自己也有這份心腸.
繼續擴展下去.
使你的生命能夠成長.

$^{881}$讓我們一起和心低頭禱告.
看到保羅如何.
由哥羅西教會的見證.
看到福音的本質.
看到以巴弗他的童宮的事奉.
以至保羅自己的模樣.
維繫在一個很重要的訊息.
正如福音的本質.
可以成長起來.
可以生出果子.
multiply.
繼續生養眾多.
這個本來就是你在創世的時候的原意.
很可惜.
人走了樣.
被罪所牽連.
如今你透過主耶穌基督.
讓我們可以重歸於這個生命.
可以為此獻上感恩.
求你讓我們同樣可以.
好像福音一樣.
繼續在你所擺我們的地方.
幫助更多的人.
認識這個福音的本身.
那個本意.
活得像你.
求你的帶領.
禱告奉教.
主耶穌基督.
得勝名致救.
阿們.
\newpage



\section{羅馬書 12:1-2}
\label{sec:031605}
\textbf{建立健康的使命教會 (4) 活出敬拜人生}~~~~吳磐德牧師
\newline
\newline
連結: \href{http://www.church.com.hk/acms/content.asp?site=cdc\&op=show\&type=product\&code=031605\&layout=sermon}{\texttt{031605}} ~~~~ 語音日期: 2026-02-08
\newline
\newline
\hyperref[sec:031604]{\small{< < < PREV SERMON < < <}}
~
\hyperref[sec:index_chronic]{\small{[返順時目]}}
~
\hyperref[sec:031606]{\small{> > > NEXT SERMON > > >}}
\newline
\newline
羅馬書 12:1-2
\newline
\begin{longtable}{cl}
\hline
\hline
章節 & 經文 (和合本修訂版)\\
\hline
12:1 & \begin{tabularx}{0.7\textwidth}{X} 所以，弟兄們，我以神的慈悲勸你們，將身體獻上當作活祭，是聖潔的，是神所喜悅的，你們如此事奉乃是理所當然的。 \end{tabularx} \\ \\ \relax
12:2 & \begin{tabularx}{0.7\textwidth}{X} 不要效法這個世界，只要心意更新而變化，叫你們察驗何為神的善良、純全、可喜悅的旨意。 \end{tabularx} \\ \\
[1ex]
\hline
\hline
\end{longtable}
$^{1}$各位大型節目, 早晨!.
我們在教會里面, 過去幾個禮拜,.
我們一直在學習, 如何建立一個健康的市民教會..
當大家看圖畫的時候,.
很想你去留心, 那些一點點是在移動的..
就好像我們的細胞, 就好像我們教會里面的五個健康元素,.
是不斷增長的, 成長的, 活動的, 活潑的..
其實教會, 自從耶穌基督創立教會以來,.
二千年來, 都是這五方面, 五大元素,.
去建立一個健康的教會..
但是究竟我們有多認識健康的教會, 究竟我們如何活出呢?.
過去這幾個星期, 我們一起去《塵聖經》裡面看,.
如何活出敬拜的人生..
今天是第四講,.
我會將活出敬拜的人生, 一個很重要的信息,.
特別是在《羅馬書》, 如何教導我們, 和大家分享..
但是我很想大家繼續為自己祈禱,.
在未來這半年, 究竟我們在我們的生活當中,.
這些字我們全都知道, 活出敬拜的人生,.
我們如何有質素的相交, 在小組, 不僅是普通的社交..
我們都知道的, 但是究竟我們如何深入活出呢?.
如何在《聖經》裡面看到, 究竟我們如何回應呢?.
勤成長的靈明, 我們都知道要做了,.
但是究竟我們如何做呢?.
《聖經》裡面有什麼要求呢?.
或者是得發揮的恩賜,.
我們如何去服侍他, 用什麼心態呢?.
我們如何去服侍神呢?.
我們在洛傳揚的恩典, 作為整個教會, 一個使命的教會,.
必須要將福音傳開, 但是我們如何去回應神呢?.
我想和大家說, 這五方面,.
其實我們不多不少有聽過, 不多不少有認識,.
其實在我自己的事奉生涯裡面,.
很多時候, 由依仗到實踐, 我自己的經歷,.
有時真的經過十年的時間, 聽同一段經文,.
同一個活出敬拜人生,.
其實當我預備今天的講道的時候,.
神將羅馬書十二章一至二節給我, 是一個更新的體會,.
我自己覺得, 好像重新去看這段經文,.
對我一個很大的提醒,.

$^{41}$也將這個提醒, 好想和大家分享,.
所以願意我們一起, 在未來這半年,.
我們是很謙卑的, 在神的面前,.
重新認識這五方面, 也掌握,.
也明白神的心意, 然後一起活出,.
互相守望, 在我們的教會當中,.
以至我們能夠成為一個健康的善明教會,.
我們一起同心低頭祈禱,.
當天父來到我們面前的時候,.
就好像詩歌剛纔所提醒我們的,.
主耶穌, 禰是深愛我們的,.
我們真的要謹記,.
主有時我們都是知道禰愛我們,.
我們去明白禰愛我們,.
但壞在我們未必每天每時每刻,.
感受主禰的愛是多深,.
證明我們求禰幫助我們,.
感受到主禰的擁抱,.
禰對我們每一個人的愛,.
使我們懂得怎樣去愛禰,.
也都是回應禰,.
證明我們求禰挪開我們每一個人的舍心,.
賜給我們肉心,.
讓我們的心靈是柔軟柔和謙卑的,.
容得下禰的話語,.
又願意實踐在我們的生命當中,.
集體一起去見證禰的榮美,.
祈禱奉祭祖基督的名求..
我又考考大家,.
聽了三個禮拜敬拜人生的道理,.
究竟你來到今天崇拜,.
你覺得你是觀眾還是演員?.
觀眾的舉手好不好?.
有什麼都挺好,.
演員的舉手好不好?.
都不舉,.
又不是觀眾,又不是演員,.
你覺得自己不是演員是什麼意思?.
因為我都沒有演員,.
我就這樣坐著,.

$^{81}$但是你知不知道崇拜神是唯一的觀眾,.
他看過來,見到你的樣子和你的心就覺得慘了,.
為什麼說慘呢?.
有時被神見到的心好像很赤裸裸的,.
是很不舒服的,.
因為有些東西想要遮住,.
好像亞當夏娃那樣,.
但是原來神是唯一的觀眾嘛,.
我們經常說,.
什麼意思呢?.
神是觀眾,你是什麼?.
你也是觀眾?.
那不就看來看去?.
敬拜就變成我看上你,你看上我,.
敬拜不是這樣的,.
敬拜是真是,.
你就是唯一的,.
你所有人都是演員,.
我都是演員,.
我都是在神的面前過去幾個星期,.
領受祂的信息以至到我正在敬拜祂,.
我來到這裡是和你分享我敬拜的成果,.
我們一起去領受,.
所以請大家認定你是演員的角色,.
如果你不認定這個角色,.
難怪給個通告你,.
說11點半你就11點7來,.
或者給個通告你,.
你就不現場在網上,.
當然有原因在網上,.
我不是雜備網上的,.
只不過你的心是怎樣的呢?.
究竟你作為一個演員好不好呢?.
還是你來到現場纔看那些劇本?.
你來到現場纔看劇本?.
還是你已經預備好?.
你過去一個星期有沒有在靈脩裡面,.
在你的生活裡面,.
你預備好的,.
你的劇情已經預備好,.

$^{121}$你來到這裡,.
正華唱的詩歌,.
你能夠投入的什麼原因?.
因為你經歷了,.
否則呢,.
心心謹記著,.
你心心愛我,.
耶穌心心的愛你,.
你怎樣感受到心心呢?.
我知道這些字每個人都寫在這裡,.
我看到這些字,.
不過我的心是否真的這麼感受到呢?.
就因為看你過去一個星期,.
你怎樣活在神的面前,.
以至今天我們集體的時候,.
一起去敬拜祂,.
我們就是一個見證,.
我們是一個回應,.
所以我又想問你,.
又考考你,.
在崇拜當中,.
你來到的時候,.
你覺得詩歌,.
廣道或者其他的東西,.
就是讓你吸取屬靈電力,.
充足電,.
插著,.
禮拜一開始就可以用到禮拜六,.
差不多沒電了,.
回來就好像Tesla一樣,.
不是,.
就不是任何電動車,.
就插著,.
就充電,.
是不是這樣呢?.
這也不是不好的,.
我不是說不好的,.
起碼你也想充電先,.
不是來睡覺,.
但是,.

$^{161}$當你不覺得自己是演員,.
當你來到只是充電的時候難怪,.
哪個走音你覺得他走音,.
哪個唱得拍子不對,.
你覺得他只是看到拍子不對,.
哪個講道令你睡不著,.
你又會覺得,.
哎呀,不如下個禮拜再來,.
就是,.
你的心態沒有調整的時候,.
你就會在崇拜裡面,.
反映著你那個敬拜的心態,.
你怎麼敬拜神呢?.
你有沒有預備好來敬拜祂呢?.
而你在崇拜裡面又怎麼去,.
崇拜這個聚會里面怎麼去回應祂,.
在崇拜以外的24小時,.
和每一天的24小時又怎麼去敬拜祂呢?.
你有沒有想過原來聖經裡面邀請你的,.
是一個,.
每一個崇拜你來到的時候,.
神正在邀請你,.
將過去一個禮拜從神領受的,.
是被神更新的生命,.
向神獻誠,.
和做付出代價的獻祭,.
你有沒有這樣預備去崇拜呢?.
我第一次說崇拜是這樣的,.
不只是唱歌聽道,散會,.
但崇拜其實是神發出一個邀請,.
其實不只是在這裡發生的事,.
而是你一生命發生的事,.
來到這裡你是怎麼去回應神的呢?.
聖經怎麼說呢?.
今天我們讀的經文,羅馬書十二章,.
羅馬書是一個很重要的書卷,.
特別是講到我們基督徒的生命,.
應該在神的面前怎麼活出的,.
特別去到十二節,.
差不多是敬拜的核心,.

$^{201}$十二章第一節怎麼說呢?.
「所以弟兄們,我以神的慈悲勸你們,.
將身體獻上當作活祭,.
是聖潔的,是神所喜悅的,.
你們如此事奉乃是理所當然的.」.
有多少人會背這段經文呢?.
又不舉手,.
應該很多人會,真的很多人會的,.
那麼有多少人會活出這段經文呢?.
又難一點了,.
我也覺得,.
究竟是不是這麼容易活出呢?.
哇,很大的挑戰,很高的標準,.
其實信耶穌是不是隻是說我們應罪悔改,.
就是說我願意承認罪人,.
耶穌是我的救主,就願不呢?.
其實信耶穌就是神發出一個邀請,.
就是給了你,.
這個邀請是給我們每一個人的,.
不是給某些人的,.
所以《聖卷》的羅馬書,.
如果從中文的詩篇書,.
第一章至第十一章是整大段,.
然後到十二章一節,.
這裡是一個轉捩點,.
由十二章一路到第十六章,.
就是如何活出,.
所以一章到十一章是說究竟我們經歷了什麼,.
信了耶穌之後,十二章一節開始說,.
之後要發生的事,.
以至到第十六章,.
簡單來說,.
羅馬書第一章至十一章,.
它最重要的信息就是福音的核心,.
你信什麼呢?.
我想請問,你是信什麼呢?.
其實這個也是保羅想向每一個基督徒問的問題,.
保羅是很清晰地跟我們說,.
我們本來是失喪的罪人,.
你和我都是失喪的,.

$^{241}$你和我都是遠離神的,.
你和我都是同神沒有關係的,.
但卻是因著恩典稱為義,.
所以我們說羅馬書是一個恩訊稱義的一個很重要的基礎的福音書卷,.
然後他又說,.
我們是被罪捆綁的,.
有時我們信耶穌仍然被罪纏擾,.
但什麼釋放我們?.
我們犯罪的時候,什麼能夠減輕我們,.
除去我們的罪?.
是主耶穌基督,.
於是在基督裡我們可以得到真正的自由,.
我們是不配的,.
我們來到神的面前,.
誰可以說我配得這個福音,.
配得耶穌為我死,.
配得神給我那麼多的恩惠,.
沒有一個人可以說,.
我們是不配的,.
我自己也是不配的,.
但我們是什麼?.
我們是一群蒙神揀選,.
被祂收納,接納成為祂的義女,.
耶穌說我和你是祂的朋友,.
是這樣的關係,.
是在祂裡我們得到保守,.
沒有任何東西能夠使我們與神的愛隔絕,.
所以纔是祂的愛對我們有多深,.
羅馬書十一章,.
在整卷書有十六章,.
用了三分二的篇幅跟我們說,.
我們原先是怎樣的,.
然後神做了些什麼,.
以至我們得到這個恩惠和救恩,.
於是我們就清楚明白,.
救恩不是我們賺取回來的,.
不是說我們多做些好事,.
我們要怎麼做纔有神的恩典,.
神的恩典已經給過我們,.
祂的憐憫已經給過我們,.

$^{281}$是神百百次給我們的憐憫,.
這是何等大的恩典,.
神的愛是何等的深,.
於是在第一節裡才說「所以」,.
因著剛纔所說的所有東西纔是「所以」,.
所以當你和我信了耶穌,.
得著你的福音,得著恩惠的時候,.
保羅就告訴我們我們要做什麼,.
我們自然就會做這些事,.
就是保羅說「弟兄們,我以神的慈悲勸你們」,.
慈悲是什麼?.
神的慈悲就是神對我們的愛,.
祂對我們深切的憐憫,.
祂對我們犯的罪,.
對我們的違抗,.
對我們對祂的不敬,.
神憐憫我們,.
感動我們,赦免我們,.
以至我們能夠回應祂,.
祂是信實的,.
只要我們回到祂的面前,祂就赦免我們的罪,.
祂是私恩的,.
我們每個人都經歷過,.
在我們一生裡面,神給我們很多不同的恩典,.
在不同的階段裡面,.
祂是守約的,.
祂說過的話算數的,.
祂這個慈愛,.
所以保羅說「弟兄們」,.
祂是虎口婆心地勸大家,.
勸你勸我,.
所以當我再一次看這段經文的時候,.
當我默想這段經文的時候,.
心裡面都有很大的打動,.
就是保羅和我說,.
如果今天保羅和你說,.
「弟兄,姊妹,你信什麼?」.
「你今天是信什麼?」.
如果這一切都是真的,.
如果福音是真的,.

$^{321}$你的生命是如何回應這個私恩的神呢?.
纔會說「所以」.
其實敬拜就是一個.
神為我們做了這麼大的恩典之後的回應,.
當我們回應祂,.
我們有一個感恩的心,.
我們就充滿著感恩的地步不得不回應,.
是一個自然的反應來的,.
是一個自然的反應,.
於是就呼籲信徒,.
敬拜簡單的定義就是向神的回應,.
一個自然的回應,.
但我們是如何回應祂的呢?.
神這麼大的恩典,這麼大的愛,.
於是保羅就說,.
「將身體獻上當作活祭,.
是聖潔的,是神所喜悅的,.
你們如此事奉乃是理所當然的.」.
哇!是整個身體獻上給祂做活祭,.
「如此事奉」是指「獻上活祭」,.
英文譯本其實不是用「事奉」,.
是用「Spiritual Act of Worship」,.
換句話來說,.
我們如此敬拜,.
如此去崇拜神,.
是理所當然,是全然的獻上,.
是當我們回應神的時候,.
每一次的回應就是敬拜,.
我們是如何回應神的呢?.
來到崇拜當中我們是如何回應祂的呢?.
我們平時每天的生活是如何回應祂的呢?.
於是,如果我們看不同的英文譯本,.
特別是Message譯本我很喜歡,.
什麼是身體獻上當作活祭呢?.
是什麼呢?.
英文譯本,Eugene Peterson 譯得很好,.
他說,求神幫助我們,.
「Take your everyday ordinary life」,.
什麼是獻上活祭,每天敬拜神,回應神呢?.
就是將你平常每天的生命,.

$^{361}$包括你的睡眠,.
不要以為睡眠不好,睡眠很神聖的,.
因為你不休息怎麼做得好呢?.
喫也很重要的,.
不喫飽怎麼做事呢?.
然後going to work,.
然後纔去做事,.
walking around life,.
你在生命裡面做的任何事,.
然後將這些每天的事,.
換句話說就是,.
你整個生命的全部,.
你的時間,金錢,資源各種,.
place before God as offering,.
就是擺在神的面前,.
「神啊,我給了禰,我奉獻給禰,.
這個是屬於禰的,就交給祂.」.
這個就是敬拜,.
其實我們基督徒一生要學習的是什麼呢?.
就是如何在我們每天的生活裡面,.
去敬拜神,回應祂,.
就是在我們每一個生活的範疇去回應祂,.
求神幫助我們,.
讓我們能夠真正活出,.
一個生命全然獻誠的生命,.
其實神是邀請我們每一個信徒,.
既然祂已經為我們成就了救恩,.
既然福音是真實的,.
既然耶穌的救贖,.
耶穌的愛對你是真實的,.
你是會這樣去回應,.
於是你就會去獻誠,.
你的獻誠不是另外一件事,.
就是奉獻袋傳到來的時候,.
奉獻金錢,不是這樣的,.
其實這裡所說的,.
其實保羅想邀請我們的,.
是整個的生命獻誠給神,.
是先獻給祂,.
然後你就會去安排你的生活如何獻誠,.

$^{401}$身體是什麼呢?.
身體,在希臘語文這個字,.
其實是說我們生命的全部,.
除了剛纔Eugene Peterson的演繹,.
就是生命的全部,.
有些解經家就說,.
身體是包括你的時間,精力,習慣,.
你的身體各方面的功能,.
工作,休息,.
所以包括你的工作態度,.
你如何使用金錢,.
你如何安排時間,.
身體和生活的節奏等等,.
換句話說,一句話就是你成就人,.
但是我想給大家一個邀請,.
就是羅馬蘇保羅給我們的邀請,.
你會如何將身體,.
就是剛纔所說的整全的生命,.
逐步逐步分,.
或者全部一起,.
如果你能夠的話,.
就是去想想,.
我在每一個崗位,.
每一個事情上,.
我是如何回應神的呢?.
因為每一個的回應就是敬拜嘛,.
那你如何去敬拜神,.
就是獻賞的呢?.
獻賞是一個行動來的,.
是要做一些決定的,.
在舊約聖經裡面,.
祭壇就是將你最好的牛,羊,.
獻祭給耶和華神,.
然後就燒了給祂,.
那個是獻賞,.
當你要決定獻祭的時候,.
首先你要決定哪一隻牛好一點,.
看到這只是Prime Rib,.
喫了它好一點,.
獻給神,.

$^{441}$譁!浪費的,很大浪費的,.
但是你要決定了,.
Prime Rib也要獻給神,.
先這樣的時候,.
你其實是一個決定,.
你一天不決定,.
其實一天不會發生的,.
所以其實這裡說的是,.
身體就是我們全個人每一天要做的什麼事呢?.
就是我們是獻賞,.
就是每一天都要做那個決定,.
我們就是那個祭物,.
我們的生活就是那個祭壇,.
我們在這個生活當中,.
我們是怎樣每一天獻誠給神的呢?.
我們是怎樣去燃燒給祂的呢?.
活祭是很特別的,.
活祭這個字是保羅自宰的,.
你發不發覺?.
因為祭是什麼?.
在舊約聖經裡面,.
獻祭就是死得,.
是不是?.
你試一下祭羊上去,.
走下來你就把祭羊抓上去,.
因為你有辦法抓羊底,.
或者先宰死它,.
等它死了沒得走,.
才放上祭壇,.
所有都是預備的,.
要選擇好的東西,.
放上去,.
舊約裡面那些祭祭完之後是怎樣的?.
還有沒有?是燒的嘛,.
死了,沒有了有沒有?.
沒有了,只是燒了一些灰,.
聞到很香,.
想起來很餓,.
就是燒了一隻這麼漂亮的牛,.
但是,.

$^{481}$如果你是活祭保羅的意思就是,.
舊約的祭在神面前死了,.
就獻給祂,.
將最好的獻給祂,.
新約的祭是什麼?.
是活祭,是活生生的祭,.
你就放在祭壇上,.
你和我放在祭壇上的時候,.
神學家就說,.
這個真是很大祭,.
很麻煩,為什麼呢?.
因為活祭就不可以殺了它來建祭,.
那個就死祭了,.
你不祭都要祭的,.
死了那隻牛,.
但是如果你是活祭,.
最糟糕的是人,.
神學家就說,.
人有自由意志,.
你上了祭壇,.
你在生活當中,你又走下來,.
你又逃避,.
你又不想獻,.
就是這樣,.
所以呢,.
我們其實保羅是提醒大家,.
你是活祭,.
是每天,.
你的生活就是祭壇,.
你就是祭物,.
你是每天要做的決定,.
今天早上起來問神,.
神啊早晨,我今天有些什麼可以回應你?.
我獻些什麼給你?.
我今天去教書,.
我怎麼對我的學生?.
怎麼用我回應你,.
去教導我的學生成為你的見證?.
在祭壇裡面也是一樣,.
我們生活當中也如是,.

$^{521}$所以,.
這裡保羅就提醒我們說,.
我們的身體就是我們全部,.
我們的獻上是要做決定,.
而我們是要做活祭,.
這個活祭是你活生生每天都要問自己的,.
每天在進行的一個事情,.
當你這樣做的時候,.
這個活祭就是敬拜,.
因為你們如此事奉,.
你如此獻祭,.
就是你的敬拜,.
你每天都在敬拜,.
來到這裡是我們一個集體的見證,.
是理所當然的,.
因為是經過你思考的,.
經過你邏輯推斷的,.
經過你感受的,.
然後你就覺得理所當然的,.
它應分得的,.
是神所喜悅的,.
所以當我們生命全然的獻誠的時候,.
敬拜的核心就是獻上活祭,.
敬拜是做什麼的?.
所以你什麼時候做敬拜?.
是不是每天?.
每一個決定就是一個敬拜,.
因為每天都是一個敬拜,.
你每件事情如何為神去回應這件事,.
而回應神來回應這件事呢?.
那個就是一個獻祭,.
所以是你全部的時間,精神,習慣,所有價值觀,.
所有都是以神為中心,.
以至你可以做一個抉擇,.
而每一個抉擇是回應神的時候,.
這個這樣的抉擇是理所當然的,.
是相等對等地回應神的,.
我們每一個人都不同的狀況,.
沒得和別人比較是吧?.
你在神的面前,你如何回應?.

$^{561}$就是你和神之間,.
但關鍵就是你如何去回應神,.
因為每一次的回應就是你的敬拜,.
這樣纔會保羅繼續到十二章的第二節,.
就說不要效法這個世界,.
只要心意更新而變化,.
叫你們察驗何謂神的善良,純純可悅的旨意,.
為什麼第一節叫我們獻祭,.
第二節就叫我們察驗神的旨意?.
有什麼關係呢?.
關鍵就是當你每天要做決定,.
而將你最好的和增活的每一個範疇,.
要獻呈給神的時候,.
你是需要知道世界的價值觀是怎樣,.
因為效法這個世界的意思就是,.
這個世界是不需要你批准的,.
自然而然就在冒犯你,.
世界的成功觀是怎樣的?.
生產力,做得有效率,身份,成就,.
你的價值觀和別人比較,.
或者對真理有些感受,.
或者抗拒一些人對你面前的真理,.
這些都是世界的價值觀,.
在影響我們和神的關係,.
保羅說你需要抗拒,.
你不要接受,.
關鍵在哪裡呢?.
你要抗拒不被世界蒙造,.
你需要有一個辨識能力,.
知道哪些是世界的價值觀,.
哪些是神所喜悅的,.
如果你沒有辨識能力就沒辦法知道,.
那你要獻什麼給他,你先懂得分辨,.
然後他說要深以更新而變化的,.
就是不斷地由內而外去作出改變,.
為什麼呢?.
因為當我們信了耶穌,.
明白了祂的恩典和愛,.
我們做了決定接受耶穌基督做我們的救主,.
就是所以之後就要獻上活祭,.

$^{601}$就是基督徒一個很自然的反應的時候,.
你的心裡面自然地要跟隨神,信服神,.
於是你就由內而外願意一個信服,聆聽,受教的心,.
神就可以不斷地在你的生命裡面去產生果效,.
以致你能夠回應祂,.
每一天的靈脩你怎麼回應神,.
每一天的決定你怎麼去做,.
以致你能夠榮耀祂,討祂的喜悅,.
記得帕薩艾倫上個星期說,.
榮耀神是什麼?.
就是讓別人驚訝神的作為,.
稱讚神的作為的任何事情,.
你做得對,你做得喜悅神的事,.
別人都會覺得你很值得稱讚,.
就將榮耀歸給神,.
就是這樣的,.
所以不要效法這個世界,.
最關鍵的就是你其實要察驗神的旨意,.
所以Empathize Bible,英文,日本說得好,.
就是不要效法這個世界,.
any longer with its superficial values and customs,.
世界的價值觀是很表面化的,.
它的習慣是你要很警醒的,.
然後他說focusing on godly values and ethical values,.
要效法的是什麼呢?.
就是集中焦點看著神的價值觀和道德倫理,.
聖經裡面的標準,.
然後按著這些分辨哪些是世界,哪些是神的心意的時候,.
英文,日本說in his plan and purpose for you,.
究竟我們檢查,察驗神的心意是什麼?.
是你生命的目的和你生命的計劃,.
察驗神的旨意就是神究竟在這些事情上,.
在你每天的生活當中,.
你每一個決定的現象都是引領你去看見神在你生命當中的旨意和計劃,.
於是你正在經歷神和你同行,.
你就很開心地和神同行,.
你正在經歷改變,你就能夠榮耀神,見證神,.
這是一個互動的關係,.
當《羅馬書》第十二章一至二節,.
這兩節經文,第二節特別,.

$^{641}$就描繪了一個新生命的新生活,.
就是每天的察驗,每天的分辨,每天的回應,.
察驗了這個旨意之後,知道他的計劃,你就回應,.
你就跟著他去做,.
於是如果我們這個新生活,我想邀請大家的就是,.
在你作為基督徒的生命當中,.
其實你每天要做的事情就是起身,祈禱,讀經,靈脩,.
為什麼呢?因為知道什麼是神的旨意,.
知道祂的旨意之後,就拒絕世界的價值觀,.
世界怎麼看你,而是神怎麼看你,.
於是你就會備神,榮耀神,已經在你心裡面了,.
就不斷對你的生命產生變化,.
你又察驗了它產生變化的時候,你就跟隨神去做一些決定,.
很多時候我們經常說,神你的旨意是什麼?我們怎麼跟隨你?.
我們應該找哪份工作?我們應該買哪間房子?.
我們應該買不買?諸如此類,每一個決定的時候,.
其實那些東西,如果當你的生命獻給神,.
你已經是拒絕了世界的價值觀,用神的價值觀去看每一件事,.
你所做的每一個決定,就在神的心意當中,.
那個就是敬拜,那個就是敬拜,.
所以有一個很出名的靈脩神學家,叫做A.W. Torser,《圖書》,.
他說: "Worship is no longer worship when it reflects the culture around us more than the Christ within us.".
在崇拜當中,如果你來到,你的心是很世俗的,去看崇拜的經歷,.
又或者你在工作上,你在職場當中,你在家庭裡面,.
你教子女都是用世界的價值觀的時候,那個就不是敬拜,.
但是當你教你的子女,用聖經的話語,.
其實所有這些就成為了你敬拜的生命,這個新生命裡面,新生活,.
你是每天都可以活出敬拜的人生,土神起源的人生,以神為中心的人生,.
因為當我們看這兩節經文一起看的時候,.
敬拜就是這兩節經文的核心,.
"如此是奉乃"是理所當然的,.
而敬拜就是獻祭,.
當我們看舊約的獻祭,新約的獻祭的時候,.
你就會知道保羅在這裡發出什麼邀請,你是怎麼獻祭的,.
舊約會莫的獻祭是怎樣的?.
拿著生畜,牛和羊,就死在祭壇上,.
放在童祭壇上,一次搞定,燒完沒有了,就祭司去獻祭,.
祭生是被焚燒,祭物是被神燒,.
但是羅馬書保羅特別用"活祭"去形容我們對神的獻祭,.
和應該理所當然的回應,就是信徒自己就是祭物,.

$^{681}$活的祭是生的,.
賢在整個生命放在神的裡面,每一天的決定是天天活出,.
生命就被神塑造,因為我們辨識神的心意,.
信徒就是祭司,也是活祭,祭物,.
你怎麼去回應神呢?.
當我們這樣理解,我們每天要掌管我們的生命,.
知道神的心意,然後作出回應的時候,.
其實你就會明白,.
為什麼當我們教會講"活出敬拜人生"的時候,.
英文是用"good stewardship",.
因為當你已經在獻情給神的生命的時候,.
你要獻的祭是需要有一箇中心的管家的角度去看你的生命的每一部分,.
以致你知道原來所有的東西都是神給過我們的,.
是神給我們的恩惠,在畢歷前書裡四章十字提到,.
人人要照自己所得的恩賜彼此服侍,.
作神各種恩賜的好管家,.
我很想鼓勵弟兄姊妹,.
你不要看信仰是compartment,即分格式的生命,.
今天星期天回崇拜,.
昨天星期六開小組,星期日上午去主日學,.
平時星期一上班,星期二晚上有祈禱會,.
事奉我為哪個崗位做了什麼事奉,.
其實不是每一個單獨的事情,.
而是所有的這些都歸納在一個的獻情,.
就是你的生命,.
一個的獻情就是你的生命,.
所以特別在畢歷前書四章十字當中說到,.
人人要照自己,就是你,就是我,.
所得的恩賜就是你的生命,.
你的財富,你的資源,你的時間,你的能力,.
你事奉的恩賜各種,我們要彼此服侍,.
但我們怎麼知道我有什麼?.
我們知道如何服侍,.
就是需要一個忠於神,已經賜給你的管家,.
你留意這裡又說神所賜,.
即不是我們自己擁有,.
是神所賜的,是給過你的,.
然後你拿著的時候,.
神就說這些是恩賜來的,.
就給過你,你就做各種恩賜,.

$^{721}$就是你生命各種資源的好管家,.
如果你不做一個好的管家,.
你沒辦法去敬拜真正的敬拜神,.
但如果你要真正的敬拜神,.
你就每天要做一個好的管家,.
因為好的管家就是管好你的生命,.
去回應神,具體的回應神,.
於是你那個行為本身就是一個敬拜,.
所以保羅羅馬書在第十二章一節裡說,.
將身體獻上當作活祭,.
這個活祭就是你每天的獻上,.
那我怎麼知道獻上什麼呢?.
什麼時候獻上?在哪裡做?怎麼做呢?.
就是管家的職份,.
安排你的時間,安排你的資源,.
留心社會,留心教會,.
留心家庭和你工作上的需要,.
以致你能夠回應,用神的角度去回應,.
那個就是敬拜,.
然而我們也回應保羅所邀請的,.
就是當我們用羅馬書十二章一節,.
這個敬拜就是獻上的回應,.
自然的回應,感恩的回應的時候,.
我們就會很清楚,.
我們是管家的職份,.
就是管理好我們已經全然獻上給神的東西,.
不是說我什麼時候要捐錢給神,.
什麼時候要給什麼時間給神,不是的,.
而是我主權在祂那裡,.
我先將自己獻給神,.
然後就管理好神所託付給我們的資源,.
然後就按著神的心意有效去運用,.
這個就是一個管家的職份,.
就活出神喜悅的生命,.
究竟敬拜人生的實踐是怎麼樣呢?.
保羅說了很多大道理,.
我們都知道,.
將身體獻上,就是整個人的生命,.
但實際上我們生活的範疇是怎樣呢?.
我今天不可能說完各種可能性,.

$^{761}$不過我想說兩方面,.
一方面就是我們怎樣傳人的塔敬拜,.
當我們實踐的時候,.
來到敬拜的時候,.
譬如崇拜這個地方,.
你怎樣實踐這個回應神的生命呢?.
如果你看江穆斯在一月十八號那一堂道,.
他有說到這段經文,.
我不會重複講,.
我只不過很簡單地想提醒大家看紅色的字,.
或者藍色的字,.
你會發覺,.
我們要在聖所,也就是在穹倉,.
在聖殿,又在世界,.
在自然的環境,.
就是任何地方你都可以敬拜神,.
那你所有的地方都要敬拜神,.
因為祂值得敬拜,.
然後用什麼方法呢?.
國聲,鼓瑟,彈琴,跳舞,擊鼓,.
詩言的樂器,笙的聲音,大筆,大響的筆,.
高聲的筆,意思是什麼呢?.
就是什麼都對,.
總之,神感動你,你回應祂的時候,.
就盡情回應的意思,.
我們教會又沒有笙,又沒有不夠鼓,.
又沒有筆,那又怎樣呢?.
不關事,.
但如果我們有,那又怎樣呢?.
那就好了,.
所以不是說有什麼,沒有什麼,.
而是你帶著一個心,.
就是我全人頭的敬拜,.
神給我們什麼,.
我們就用博百般恩賜的好管家,.
去管理好神給我們教會的資源,.
我們一起去敬拜,全人都要投入,.
我們做好的演員是怎樣呢?.
準時來崇拜,.
我們用一個獻情回應的心的時候,.

$^{801}$我們聽家事分享,.
就是報告,不關我事,.
不是哦,.
當你的心來到全然的獻上,.
我就是回應神的時候,.
究竟這個家,神的家帶著我們來這裡,.
我有什麼可以回應神的呢?.
有什麼可以參與的呢?.
我要做什麼決定呢?.
都要做決定的,.
那些是關你事的,.
所以是整個每一個環節,.
你的禱告,.
我們的家事,.
我們的建堂,每一件事其實,.
當我們全人頭投入敬拜的時候,.
是你的回應,.
你怎樣去在神的面前按住神給你的感動回應,.
都是一個敬拜,.
並且你可以很自由的敬拜,.
由心而發的,.
其實《詩篇》第150篇最精彩的地方就是,.
好像大衛,當時張國光老師也有用大衛的例子,.
大衛是不顧自己的,.
因為神纔是唯一的觀眾,.
他是演員,.
他演給神看的,.
手舞足蹈,連他的妻子都藐視他,.
他做皇帝這樣跳舞,.
但是他跳舞是給神看,.
不是給妻子看,不是給人看,.
所以我很鼓勵大兄姐妹,.
以後敬拜你可以拍手,可以舉手,可以起身,.
你不要說得我起身,這樣很怪,不是的,.
如果那個是神給你的回應,你就這樣回應,.
你不是給人看的,.
給神看,神看在你的心裡,.
你的心是心如止水,.
那沒辦法,那你就心如止水,.
但如果你不心如止水,你就是激盪,.

$^{841}$你就回應激盪,.
應該是自由的,.
你反映了你的內心的敬拜,.
第二方面大家都會猜到我一定會講的,.
就是奉獻,.
一定會講的,.
但我也想大家明白,.
從羅馬書那裡看奉獻,.
有一個很有趣的地方,.
就是奉獻從來在新約聖經,.
十一奉獻只是一個標準,.
你會留意在新約聖經裡面,.
從頭到尾沒有十一奉獻這個標準,.
沒有的,.
你會說牧師你今天講這個真好,.
我最喜歡聽這個,.
因為十一奉獻不是一個要求,.
不是一個命令,.
是一個建議,.
不過我們先聽完後面,.
因為這個建議,.
就是在耶穌成就了律法之後,.
以前舊約是一個律法,.
十一奉獻是律法,.
你將投生的最好的獻給神是律法,.
是神的,.
這個原則來到新約裡面,.
仍然存在的就是所有東西都是神的,.
所以在羅馬詩十二章一節裡面,.
才說把身體獻上就是全部的獻上,.
你說那我怎麼辦?.
買屋怎麼辦?.
生活怎麼辦?.
喫飯怎麼辦?.
他就說要測驗神的心意,.
神也想你住得好一點,.
神也想你開心,.
就是神不會要你,.
神一定有你和祂之間祂給你的指引,.
所以經常在新約聖經裡面的教導就是,.

$^{881}$不要勉強是甘心,.
十分一是一個指引,.
所以你可以十分二,.
就不是十分一,.
為什麼沒有人有反應?.
其實真是,.
高露潔是十分九的,.
大家都聽過高露潔,.
Google一下就可以了,.
KOLGATE,.
十分九,為什麼?.
因為祂就全部獻上,.
祂就拿十個百分比來用,.
因為祂覺得神已經給祂一百個百分比了,.
祂就給祂九十個百分比,.
只是自己用了十個百分比,.
但是每個人的情況都不同,.
所以不需要勉強,.
新約的教導就是,.
不要勉強,不要作難,.
但是一切已經是神的,.
用羅馬書的原則就是你怎麼去辨識,.
你怎麼去辨別呢?.
於是,.
巴基頓教會的理財之道就成為我和你今天很重要的提醒,.
這裡怎麼說的?.
它說,弟兄們,我們要把神賜給馬其頓宗教會的恩惠告訴你,.
你留意,.
保羅說,.
我把這個捐錢的事,.
其實是一個恩惠,.
就是神給了一些恩典給這群人,.
所以是神所給的,.
現在我和你分享他們這個恩典的故事,.
但是這個恩典的故事聽起來一點都不恩典的,.
你有沒有留意?.
為什麼?.
他們在患難中受大考驗的時候仍然滿有喜樂,.
你覺得你考驗的時候滿有喜樂,.
你覺得是恩典還是受苦?.

$^{921}$但是如果你有一個正確的,.
以神為中心的生命的觀,.
那個也可以喜樂的,.
因為神為中心的,.
祂有祂的心意的,.
你願意接受的,.
然後祂說,.
怎樣纔算是恩惠?.
祂還說,.
在極度貧窮之中還格外顯出他們樂捐的慷慨,.
他們的狀況是怎樣的呢?.
極度貧窮,.
那何來恩惠呢?.
他們是否勉強呢?.
不是,.
保羅說,.
我可以證明,.
他們是按照他們的能力,.
而且超過他們的能力來捐助,.
主動再三地懇求我們,.
準他們在這個供給聖徒的善事上有份,.
什麼人,.
當我再看這段經文的時候,.
什麼人會在自己都貧窮,.
極受考驗,.
患難的時候,.
去再三要求保羅,.
表求對方,.
你拜託了,.
你讓我幫你吧,.
你也應該說,.
拜託了,.
你幫一下我吧,.
幫我先解決那間屋的事吧,.
先解決我的工作吧,.
先幫我吧,.
以我為中心,.
以神為中心是怎樣的呢?.
他們看見神在他們生命的恩典的時候,.
他是按照他們的能力,.

$^{961}$這裡已經講得很清楚了,.
不是說超越他們的能力,.
是按照他們的能力雖然很貧窮,.
但是仍然看到神給了他們付足,.
於是第五節是關鍵的,.
就是他們所做的不是照我們所期望的,.
我們期望的是他們不用的,.
誰知道他們也做了,.
這是恩典的故事,.
因為神在他們的恩典感動他們去做,.
不是被逼的,.
所以神從來不是要我們奉獻,.
不是要我們事奉,.
是被逼的,.
是拿杆拿鏡的,.
不是的,.
而是祂給了你一切,.
你也願意給祂一切,.
然後你辨識究竟什麼是你的那一份,.
而關鍵是先後,.
留意最後這一段我們一起讀,.
《更照神的旨意》那裡一起讀,.
更照神的旨意,先把自己獻給主,又給了我們..
臨留意到馬其頓教會的弟兄姊妹,.
他們實踐保羅那個將身體獻上當作活祭的先後次序,.
首先是將你自己獻上,.
如果你不將自己獻上,.
你信耶穌沒有將自己獻上,.
沒有回應神,.
所有都是以自己為中心的時候,.
不是以神為中心,.
你的主權仍然在自己那裡,.
不是在神那裡的時候,.
一切都是被逼的,.
一切都要做給人看的,.
人做你就做,.
但是如果主權在神那裡,.
如果你是一個自然理所當然的回應,.
你就會先將自己獻給主,.
於是你有一個樂淵的心,.

$^{1001}$去將十分一,十分二,十分九,.
好像高老傑一樣,就是神按著你和他的關係,.
你和他的關係,他和你的關係一起去回應他,.
而那個回應就是一個敬拜,.
以下我很想由Kenneth王扎庭弟兄分享一下他的生命見證,.
我看看他在哪裡,.
他很簡單的一個生命見證,.
我也想大家有些真實的生活上如何去見拜神的分享,.
各位,剛纔多謝牧師,多謝主,.
因為給了我一個機會和大家做一個見證的分享,.
這樣的,我是Kenneth,.
我回流了十年,.
我回流回來和我爸爸一起做裝修接工程的生意,.
去年,正正一年,非常愚蠢,.
不小心弄斷了一隻腳,.
沒辦法了,要躺在床上,.
因為做這個行業,手停就接近口停,.
那時候也很擔心,因為擔心前路,擔心家裡,擔心經濟,.
擔心別人擔心我,非常擔心,.
我也是依然嬉皮笑臉,.
唯有躺在床上,每天都是去祈禱,看書,做我的事,.
很奇妙的是,當過了,又是大約一年,.
剛剛那段時間,忽然間有一份兼職,可以在家裡做的,.
做一些我以前本行的事,做一個問卷調查,.
因為那時候的事宜值是一個省選,.
所以有一個月的時間讓我去做,.
挺好的,因為它是由早做到晚的,.
很神奇的,.
因為如果大家有印象的話,.
就是省選到差不多的時候,就變成了聯邦大選了,.
再做多一個月,.
我足足做了兩個月的時候,可以撐住家裡的經濟,.
去到差不多四月的時候,其實我已經可以下床,可以動了,.
也開始接到一些生意,接到一些工程來做,.
神就真的很好,因為祂知道我休息了這麼久,.
是時候要活動了,.
所以就幫我排了很多的工作給我,.
一直排到年底,有些還是挺大的,.
而更加喜出望外的就是,.
在這段日子裡面,神也安排了很多小型的工程給我的,.

$^{1041}$就是可能做一兩天,.
放一些時間出來,星期六日假期都去開一下工,.
那還可以的,我覺得幫到人又有時間,.
又是一個收入去做,.
有一天就是在其中做一個工程的時候,在游泳的時候,.
游泳的時候就會有很多東西想,.
忽然間就想起之前有一次的靈脩經文,.
那個經文就是在約翰福音裡面說,.
我將這些事告訴你們,是要叫你們在我裡面有平安,.
那一刻就有這個感動,.
就覺得,也都時值那時候教會開始批了批文,.
要開始可以去建堂,.
又聽到很多這些鼓勵的說話,.
所以在那一刻就覺得,.
那個工程不如這樣吧,.
將它全數都捐給教會去建堂,.
其實不是很多,只是800塊錢而已,.
捐完之後發覺,go ahead,不要想那麼多了,.
捐完之後,神就可能知道你的工程都可以幫到我建教會,.
所以又多給我一些小小的工作去做,.
尤其是我記得在Bread Friday那時候,.
Bread Friday之後就有很多做一天兩天的工作,.
所以我會將其中一張保安官抄油煙機,.
那天的收入,又捐了給建堂那裡,.
我覺得神在這方面去看顧就真的非常之好,.
因為我想我這麼多年,今年是最忙的一刻,.
我覺得我都受了個馬,.
結果我做到12月底的時候,.
終於都可以休息了,.
因為冬天又是比較淡的季節,.
我是應該去到過完年纔會正式有公開的,.
但是又很奇怪,一月不是很大雪,今年又冷又大雪,.
又不知道為什麼有人找我,可不可以幫他剷雪,.
我前後鏟了四五間屋,.
神奇的是,鏟完雪沒多久之後,又給我找了份兼職,.
每天只做幾個小時早上而已,.
我就覺得這個看顧就非常之大了,.
所以我都立意會將每花五天裡面,.
我都會將其中一天捐給建堂這件事上,.
我覺得神給我最大的功勞是信心,.

$^{1081}$還有祂是我們的以便以謝的神,.
因為祂由始至終都是看顧著我們,.
希望這方面可以鼓勵到大家,.
也可以將我這些小小的分享的見證去和大家分享,.
多謝你們!.
原來當我們全日投入去自由的敬拜神,.
我們就知道神就是一切事情的起點,.
也需要邀請我和你去過土神起源的生活,.
我們的奉獻其實不是一件額外的事情,.
而是我們獻上生命的一個部分,.
而在這個部分裡面,.
當我們做我們金錢的管家,身體的管家,.
人生方向的管家的時候,.
我們就在敬拜神,.
更加神邀請我和你在奉獻的事情上,.
好像剛纔Kenneth的見證和當保羅寫給哥林多教會,.
講到馬其頓教會的時候特別提到的,.
就是他們將主權交給神,.
一切從神而來,.
又將生命獻在祭壇上面,.
然後互相提醒,.
我們不是擁有我們自己的,.
一切從神而來,.
我們只是神的管家,.
神不是因為缺乏要求我們去給祂,.
而是因為我們本身就是屬於神,.
一切從神而來的,.
我們回去給祂,.
這是生命的操練,.
這是獻上的起點,.
也是操練我們信心的功課,.
如何去跟隨神,.
如何去回應神,.
就在這些行動上,.
我們就是在活出一個敬拜的人生,.
我們一起低頭,.
我們思想一下,.
我想大家閉上眼睛,.
閉上眼睛,.
你思想一下,.

$^{1121}$敬拜不僅是星期日,.
而是敬拜是一個生活的方式,.
你向神禱告,.
因為敬拜是向神恩典作出的回應,.
我們的生命是以神為中心,.
我們就做好神所託付給我們資源的一個好的管家,.
每天的求神憐憫,指引,.
開放我們的生命去看,.
去檢視自己的生命,.
我們有沒有做一個好的管家,.
將我們的生命全然的獻上,.
去過討祂喜悅的生活,.
你向神回應,.
你去禱告,.
你為自己去禱告.
稱天父求您的靈,.
幫助我們每個人,.
在每天做合理心意的決定,.
以至我們所獻上的,.
就是合理的心意,.
被主您引靈,.
我們每一個決定是合乎您的帶領,.
永要您的靈,.
祈禱感恩,.
奉著宿命求,.
阿們.
\newpage



\section{使徒行傳 9:32-43}
\label{sec:031606}
\textbf{耶穌戰勝死亡}~~~~禤凱榮牧師
\newline
\newline
連結: \href{http://www.church.com.hk/acms/content.asp?site=cdc\&op=show\&type=product\&code=031606\&layout=sermon}{\texttt{031606}} ~~~~ 語音日期: 2026-02-08
\newline
\newline
\hyperref[sec:031605]{\small{< < < PREV SERMON < < <}}
~
\hyperref[sec:index_chronic]{\small{[返順時目]}}
~
\hyperref[sec:031607]{\small{> > > NEXT SERMON > > >}}
\newline
\newline
使徒行傳 9:32-43
\newline
\begin{longtable}{cl}
\hline
\hline
章節 & 經文 (和合本修訂版)\\
\hline
9:32 & \begin{tabularx}{0.7\textwidth}{X} 彼得在眾信徒中到處奔波的時候，也到了住在呂大的聖徒那裡。 \end{tabularx} \\ \\ \relax
9:33 & \begin{tabularx}{0.7\textwidth}{X} 他在那裡遇見一個人，名叫以尼雅，得了癱瘓，在褥子上躺了八年。 \end{tabularx} \\ \\ \relax
9:34 & \begin{tabularx}{0.7\textwidth}{X} 彼得對他說：「以尼雅，耶穌基督醫好你了，起來！整理你的褥子吧。」他立刻就起來了。 \end{tabularx} \\ \\ \relax
9:35 & \begin{tabularx}{0.7\textwidth}{X} 凡住呂大和沙崙的人都看見了他，就歸向主。 \end{tabularx} \\ \\ \relax
9:36 & \begin{tabularx}{0.7\textwidth}{X} 在約帕有一個女門徒，名叫大比大，翻出來的意思是多加；她廣行善事，多施賙濟。 \end{tabularx} \\ \\ \relax
9:37 & \begin{tabularx}{0.7\textwidth}{X} 當時，她患病死了，有人把她清洗後，停在樓上。 \end{tabularx} \\ \\ \relax
9:38 & \begin{tabularx}{0.7\textwidth}{X} 呂大原與約帕相近；門徒聽見彼得在那裡，就派兩個人去見他，央求他說：「請快到我們那裡去，不要耽延。」 \end{tabularx} \\ \\ \relax
9:39 & \begin{tabularx}{0.7\textwidth}{X} 彼得就起身和他們同去。他到了，就有人領他上樓。眾寡婦都站在彼得旁邊哭，拿多加與她們同在時所做的內衣外衣給他看。 \end{tabularx} \\ \\ \relax
9:40 & \begin{tabularx}{0.7\textwidth}{X} 彼得叫她們都出去，然後跪下禱告，轉身對著屍體說：「大比大，起來！」她就睜開眼睛，看見彼得，就坐了起來。 \end{tabularx} \\ \\ \relax
9:41 & \begin{tabularx}{0.7\textwidth}{X} 彼得伸手扶她起來，叫那些聖徒和寡婦都進來，把多加活活地交給他們。 \end{tabularx} \\ \\ \relax
9:42 & \begin{tabularx}{0.7\textwidth}{X} 這事傳遍了約帕，就有許多人信了主。 \end{tabularx} \\ \\ \relax
9:43 & \begin{tabularx}{0.7\textwidth}{X} 此後，彼得在約帕一個皮革匠西門的家裡住了好些日子。 \end{tabularx} \\ \\
[1ex]
\hline
\hline
\end{longtable}
$^{1}$有人說生老病死是生命的一部分.
我相信沒有人會反對.
因為這是事實.
沒有一個人能夠逃避疾病,衰老和死亡.
不過如果我跟你說.
我說其實生老病死這些所謂生命的事實.
並不是神原初創造的心意.
曾經有一個大家都應該知道,聽過的.
一個英國很出名的作者叫做C.S.Lewis.
他寫了一本書叫做Miracles.
中文譯做奇蹟.
在這本書裡面他講到.
其實在神原初的創造裡面.
人是不會生病,不會癱,不會死的.
有這些情況的原因.
就是因為進入了這個世界的一個結果.
相反來說其實主耶穌和使徒所行的神蹟.
讓我們反而一睹神原初創造.
是應該是怎樣的.
所以C.S.Lewis就認為.
基督徒我們不單止需要接受神蹟這個事實.
還要為了神蹟而歡喜,受到鼓舞.
因為神蹟是神親自參與在他的創造當中的一個證據.
在今天的經文裡面我們看到兩個神蹟.
一個是講到癱子被醫治.
另外一個神蹟是講到一個死了的人得以復活.
兩個事件都可以說是違反了自然定律.
亦都由於這個原因.
這兩個違反了自然定律而發生的事件.
其實就是一個很具體的明證.
講出神在這個世界裡面他如何的介入.
亦都讓我們看到耶穌基督在十字架裡面所成就的.
究竟是甚麼一回事.
讓我們看一看 去學習一下.
在這個故事裡面的一些事情.
在第一個神蹟.
是使徒彼得遇到一個人叫做以利雅.
在一個屢大的地方.
這個人聖經告訴我們.
他病患了八年 八年裡面癱了.

$^{41}$只能夠在一張床裡面生活.
聖經就簡單地介紹完這個人物之後.
聖經沒有給很多其他詳細的資料.
沒有講到究竟在這個醫治的事情之前發生了甚麼事.
沒有講到的.
所以我們不知道究竟這個以利雅是不是基督徒.
亦都不知道這個以利雅是不是他請求彼得醫治他.
還是彼得走過寄仰去醫治他.
聖經沒有講到.
我們就會覺得有點奇怪.
因為這些資料.
如果有了這些其他的資料的時候.
我們豈不是對這個故事更加具體知道多些嗎.
是不是.
如果是的時候.
究竟為甚麼聖經沒有告訴我們這些資料呢.
為甚麼呢.
我相信這個原因絕對不是作者路加忘記了.
大義疏忽.
不是這個原因.
我相信這個原因就是因為很簡單.
因為這些不是聖經需要我們知道的資料.
我們需要明白一件事.
聖經不是一本故事書.
聖經亦都不是一個新聞報道.
它沒有必要將一件事情發生的所有資料.
詳細的記錄以至我們知道整個故事的具體資料.
不需要的.
我們基督徒我們相信是甚麼.
我們相信聖經是神的話語.
意思是甚麼.
意思是聖經是記下了神要我們知道的事情.
所以雖然這個故事只有四節經文.
不過這四節經文絕對包含了神需要我們知道的東西.
大家同不同意.
並不是我們想知道的東西.
究竟這四節經文這個故事要講的是甚麼訊息呢.
很簡單.
第一個訊息.
如果我們大家看.

$^{81}$經文就是說.
彼得周遊四方.
不錯 周遊四方.
誰不想去是holiday.
會否那時候旅遊團便宜價.
所以彼得去旅遊.
很明顯.
彼得所說的周遊四方.
並不是旅遊觀光.
亦不是為了增廣見聞.
究竟是甚麼事.
聖經說.
如果我們看前一點.
在《使徒行傳》第八章說.
從這日起.
耶路撒冷的教會大遭逼迫.
這日是甚麼意思.
這日就是當日.
大家記得吧.
史提反被石頭打死.
有人用石頭打死了史提反之後.
那日教會就開始遭受到大逼迫.
所以當時在耶路撒冷的信徒就走了.
四散.
聖經說他們散到猶太和撒瑪利亞不同的地方.
當然.
這些信徒離開了耶路撒冷.
不論去到哪裡.
他們都是基督徒.
所以他們就帶著基督的信仰.
去到其他地方.
過了一段時間.
我們來到《使徒行傳》第九章.
大家可以想像.
信徒離開了耶路撒冷.
即是說他們得不到教會領袖的牧養.
大致上是這樣的原因.
彼得作為教會的領袖.
就要去到那些地方.
我們都牧養四散了的信徒.

$^{121}$這個發生在.
呂大的神跡.
如果大家看地圖的話.
就在耶路撒冷的西北邊.
大概是25英里左右.
如果我們看地圖的話.
我們不難發現.
這個呂大就在耶路撒冷和沿海位置的城市.
包括我們等一下第二個神跡發生的地方約柏.
在這兩個地方中間.
即是說.
他是在很多主要的城市的必經之路.
如果你要去耶路撒冷或者沿海地帶.
都需要去到這個地方.
即是說這個地方是當時的貿易.
旅遊或者文化的交匯地.
聖經就說彼得去到這個呂大的地方.
探望當地的聖徒.
即是告訴我們.
當時分散了的基督徒.
已經在呂大這個地方.
其實已經建立了一個信仰的群體.
接著彼得就見到以利雅的一個人.
聖經沒有說到究竟是.
彼得主動的去醫好他.
還是他要求彼得去醫好.
因為聖經一開始的時候就已經說了.
彼得對他說.
以利雅.
耶穌基督醫好你了.
起來收拾你的玉子.
這裡就很簡單.
講.
這個舉動以利雅就得到醫治了.
我們需要認真地看一看.
因為聖經只是說了這個醫治的舉動.
即是說這個舉動.
簡單的這個說話就包含了.
很清楚.
神要我們知道的事情.

$^{161}$現在我們認真地看一看.
這麼簡單的一句.
叫他起來的命令.
單看這句命令的時候.
大家有沒有留意.
有沒有提醒你們.
另外一個耶穌.
很類似的一個醫治的經驗.
有沒有提醒大家一個類似的醫治.
如果我們看這個馬太福音的第九章.
是有的.
曾經有一個人.
就抬著一個毯子.
去到耶穌的面前.
就要求耶穌醫治他.
耶穌見到這群人這麼有信心.
耶穌就做了一個他們都沒有要求到的事情.
他們就是求醫治.
但耶穌就說.
你的罪.
赦了.
當時在場的其他宗教領袖聽到之後.
就很不開心.
很不開心就舔你一口水.
為什麼你可以赦罪.
你這樣說的時候.
你是犯了這個潛枉的罪.
你是褻瀆了神.
那當然了.
作為耶穌基督是神的兒子的時候.
耶穌當然有絕對的權力.
來赦罪和醫治任何的疾病.
所以你看到耶穌不需要開藥單.
也不需要做手術.
耶穌一句簡單的說話.
你起來,這個癱子就得到醫治.
好了.
我想大家去想一想.
耶穌這個叫癱子起來.
和彼得那個醫治的命令有沒有類似.

$^{201}$很類似的.
特別他們所說的都是同一樣的.
起來,拿起你的玉紙.
為了一個不同的.
大家有沒有留意到有一個不同的地方.
不同的地方就是.
彼得在這裡加了一句.
耶穌基督,醫好你了.
很清楚,講得清楚一件事是什麼.
彼得知道很清楚.
他能夠醫治病人的能力.
並不是他自己的能力.
是來自耶穌基督的.
所以彼得他沒有說.
我現在醫好你了.
馬上起來.
彼得沒有這樣說的.
他很清楚這個是神的能力.
所以他沒有將神的作為據為己有.
他將一切的榮耀歸給天父上帝.
這個醫治的舉動告訴我們很重要的事情就是.
神其實是通過在地上的信徒.
來將耶穌能夠醫病的權柄與能力.
繼續施行在地上.
只要我們真心相信耶穌基督的時候.
我們帶著信心.
有神在面前給過我們的恩賜的時候.
今天我們都可以醫病.
可以趕鬼.
以致人可以通過信徒能夠行的神蹟.
知道耶穌基督仍然活在世上.
今天的耶穌基督仍然將他的權柄.
施行在地上.
通過他的信徒.
好了,行神蹟,醫病,坦子起來.
結果是怎樣呢?.
結果就很清楚了.
這個困住坦子八年在床上的病馬上被醫好.
這個以尼雅馬上可以起身了.
留意接著發生什麼事.

$^{241}$接著聖經就說.
住在呂大和隔離沙崙的居民.
見到能夠起身行走的以尼雅就歸復主.
其實新約聖經或者使徒行傳都是一個模式.
模式就是說被醫好了之後.
人就因此而信主.
所以我們從這個角度看的話.
神蹟其實不是對自然秩序的干擾.
反而是對神原初創造的秩序恢復過來.
神蹟見證了神的權能.
亦都讓信徒有一個很具體的證據.
我們深信當時候來的時候.
神有絕對的權柄與能力.
可以恢復原初創造的秩序.
我們再看第二個神蹟.
經文在第二個神蹟的記載.
就是一個復活的故事.
發生在一個沿海的城鎮叫做約柏.
是很近旅家的.
這個地方有一個女信徒.
很虔誠的女信徒.
不過因病而離世.
她離開世界之後.
人們知道當時使徒彼得就在隔壁那個鎮.
不如我們就馬上叫人邀請彼得過來.
不要遲疑.
聖經就說彼得去到那裡.
然後就見到很多人哀哭.
彼得就叫其他群眾離開.
讓他可以單獨跟死者.
直接跟他講.
大比大 起來.
就這麼簡單.
這個離世的女信徒.
就張開眼睛可以起來.
看完這個簡單的故事.
這個簡單的叫死人復活的神蹟.
有沒有讓大家想起另外一個聖經裡面的故事.
又是很類似的.
如果我們看馬可芬的第五章.

$^{281}$當時有一個管會堂的女兒.
病得很厲害.
人們就去找耶穌.
希望耶穌去那裡治好病重管會堂的女兒.
不過很可惜.
當耶穌去到這個地方的時候.
這個女兒已經離世了.
不過同樣的.
當耶穌去到的時候.
也都是叫人離開房間.
然後就直接跟那個女兒講起來.
兩個的故事.
如果我們拍在一起看的時候.
有很多類似的地方.
都有很多一幫哀哭的群眾.
都是去到的時候叫群眾離開.
都是發出命令叫死者起來.
然後復活就馬上發生.
特別如果我們看那個命令的時候.
就更加的清晰.
當彼得叫那個死者起來的時候.
他怎麼叫.
大彼大起來.
大彼大就是那個死者的名字.
但是到耶穌那個例子的時候.
都是很類似的.
叫大利大.
古米.
古米就是起來的意思.
連聲音發音都很類似.
是不是巧合呢.
是不是這麼巧合呢.
還是彼得刻意模仿耶穌呢.
無論情況是怎麼樣.
我們由於太多的類同.
我們很難將這兩個的事件.
完全當成一個獨立的個案來看.
一定有一些的關聯.
不過雖然是這樣.
在眾多的類同當中.

$^{321}$仍然有一個很清楚的不同的地方.
有什麼不同的地方.
禱告.
有沒有留意.
彼得在叫這個死者起來之前.
他有跪下禱告的.
不過耶穌基督就沒有.
耶穌基督直接叫那個龜女起來.
為什麼.
我相信很清楚了.
耶穌基督他自己就是生命的傳言.
他對生與死有絕對的權柄.
所以通過彼得所行的這個復活的神蹟.
聖經其實告訴我們.
在耶穌基督升天之後.
耶穌基督的權柄仍然通過使徒施行在地上.
我想大家要理解.
這個能夠行神蹟的證明.
是對早期教會的建立很重要.
大家記不記得保羅曾經在二弗所書那裡.
講到在基督裡面的合一.
保羅曾經這樣講.
當我們不理我們什麼背景什麼的人.
只要我們信了耶穌之後.
就是神家裡的人.
接著講一個很重要的事情.
他說我們每一個信徒.
都是被建造在使徒和先知的根基上面.
耶穌基督就是整個根基的防國石.
我們整個信仰的群體.
就成為主的聖殿.
大家有沒有想過.
為什麼今天的教會一定要建立在使徒和先知上面呢.
我相信對於第一代的信徒來講.
特別是猶太裔的信徒.
很明顯.
一直以來的信仰體制都是建基於什麼.
神是通過舊約的先知來跟別人講說話的.
來做他要做的事情.
所以很明顯.

$^{361}$整個信仰的根基.
就是擺在先知的傳統上面.
不過使徒呢.
為什麼又要建造在使徒上面呢.
大家有沒有這樣想過.
當日很有可能.
當彼得帶著其他的使徒.
一起建立教會的時候.
很可能有人挑戰他.
彼得.
這個人我不是很久之前才跟他一起出海打魚.
我打魚比他厲害.
我打魚還比他多.
為什麼今天要我聽他講道理.
有沒有想過.
馬太.
馬太那天才在路邊打稅.
他對著我還要打超過我要給的稅.
我為此還要跟他吵個架.
為什麼今天我要聽他講道理.
大家有沒有試過在公司工作呢.
在同一個部門裡面.
本來是同事.
突然之間某一個跟你同級的同事.
被升了做你的經理.
我試過很多次.
引發很多的問題.
以前是我跟他同事.
我做事比他厲害.
為什麼今天我要聽他講道理.
很多時候有這樣的事情發生.
所以大家看不看到問題所在.
如果神介入的時候.
這12個所謂的平民百姓.
沒有錢,沒有權柄,沒有能力.
很難在早期的教會裡面.
來建立他們擁有的權柄.
但是神介入.
大家記不記得一個很重要的.
可以說是一個轉捩點的事件.

$^{401}$阿拿尼亞和薩菲拉.
記不記得這兩個人物是什麼人物.
他們欺控聖靈.
由於欺控聖靈.
所以在彼得的面前相繼死亡.
聖經就說.
這件事情之後.
全教會和聽見這事的人都甚懼怕.
主藉著使徒的手.
在民間行了很多的神蹟騎士.
人們看到這些神蹟騎士的時候.
沒有一個敢靠得太近這幫使徒.
但卻尊重他們.
也由於這樣信而歸主的人.
越發增添.
連男帶女很多.
所以告訴我們什麼.
早期的基督徒.
他們相信跟隨,接受使徒的領導.
並不是因為這幫領袖.
說那張嘴巴會過油.
甜言蜜語.
籠絡民心.
又或者他們像江湖術士那樣.
耍花招來讓人跟隨他們.
不是這樣.
而是很清楚通過他們所說的.
所行的神蹟.
證明給所有人看是神揀選他們.
賜予他們權柄與能力.
能夠醫病,趕鬼.
甚至叫死人復活.
奉主的名來建立教會.
耶穌基督就是教會的根基.
而耶穌基督本身.
也是通過三天後的空的墳墓.
證明給整個世界看.
通過神蹟.
告訴整個世界看.
祂已經戰勝死亡與罪惡.

$^{441}$能夠從死亡裡面復活.
從整體圖畫來說.
神蹟無論你是醫治,趕鬼.
還是叫死人復活.
從來都不是一個目的.
為什麼這樣說.
你猜彼得醫治完以尼亞之後.
以尼亞以後還會不會生病呢.
大概會吧.
你猜彼得叫大比大.
從死裡復活之後.
大比大是不是以後不再死呢.
也不會.
也就是說.
其實神蹟本身從來都不是一個目的.
神蹟從來都是一個路標.
指向.
讓人看得到.
神是誰.
讓人看得到.
神在哪裡.
讓人清楚知道.
耶穌基督今天仍然活著.
祂仍然掌管著上天下地的權柄.
仍然通過祂的信徒.
在地上今天在工作.
這就是神蹟.
這兩個神蹟之後聖經告訴我們.
很多人因此而歸服主.
信了主.
讓神得著榮耀.
這就是神蹟最重要的果效.
我相信除了這個果效之外.
聖經有另外一個.
在這段經文裡面.
它也在改變,模造,更新彼得.
在神蹟之後.
聖經說彼得在約柏這個地方.
和一個俏皮匠西門的家女住了幾天.
很簡單,在某某人家那裡住了幾天.

$^{481}$沒什麼特別的資料.
不過如果我們認真翻查歷史.
認真思考主宰這句經文.
其實絕不簡單.
我們要明白這句簡單的經文.
背後的力量.
我們首先需要知道.
什麼是俏皮匠.
他的工作是什麼.
俏皮匠簡單來說.
就是做皮革.
怎樣拿皮革.
就是在動物的身體上.
把皮革剝下來.
要做處理.
把肉弄掉.
處理它.
然後用皮革做各樣的皮革用品.
這就是當日的俏皮匠要做的工作.
那就是說什麼.
我們想像得到.
這個工作又臭又骯髒.
不是一個很乾淨.
是寫字樓的工作.
特別對於當時的猶太人來說.
大家都知道.
猶太人是不會碰死屍的.
為什麼.
任何的死屍碰了之後.
或者碰了一些不潔淨的東西之後.
他們在禮儀上是成為不潔的.
是很煩的.
所以猶太人基本上是不會接近這些俏皮匠的.
更不要說住在他們那裡.
但是現在聖經告訴我們.
這個比猶太人更猶太人的彼得.
今天竟然住在這個俏皮匠的家裡.
多日.
究竟發生什麼事.
我相信.

$^{521}$如果我們思考.
我相信在彼得跟隨耶穌那幾年的裡面.
他聽過很多次.
也親眼看過很多次.
一個重要的訊息.
什麼訊息.
這個訊息就是說.
人心裡面的信服.
人心裡面愛神的心.
遠遠重要過你守外在的禮儀.
我隨便舉聖經幾個的例子.
如果我們看這個馬太方的二十二章.
耶穌講得很清楚.
你們要盡心盡性盡義.
愛主你的神.
因為這個誡命中的第一.
而且是最大的旗幟就是愛人如己.
我想大家去想一想.
耶穌所指出的其實就是.
律法上所要求人的真正的精髓.
其實就在這裡.
愛神.
愛人.
很簡單.
如果你做足所有律法的要求.
你根本沒有愛神的心的時候.
聖經又講得很清楚.
沒有意思的.
如果我們看這個荷西亞書的時候.
那裡更加直接將這個道理講出來.
神說我喜愛連率.
不喜愛祭祀.
我喜愛認識神.
勝於凡祭.
留意.
這一句經文不是告訴我們.
祭祀或者律法是不重要的.
而是告訴我們.
當我們履行這些律法條文的時候.
如果我們內心沒有那種愛的時候.

$^{561}$一切外在的遵守規條都是沒有意思的.
宗教的禮儀.
遠遠不能夠替代內在的敬虔.
如果我們想做完所有的事情的時候.
我們就在神面前100分的時候.
你搞錯了.
外在的禮儀.
或者我們遵從宗教的條文.
是為了表達我們內心對神的敬虔.
所以耶穌在這裡命令.
他的信徒愛人如己的時候.
如果我們看回彼得住在這個俏皮城家裡多天.
就是在實踐這個愛人的道理.
彼得在這裡打破了社會的隔閡.
他展現了一種神對所有人的愛.
無論這班人的背景如何.
他們的職業如何.
不再有這個隔閡.
這個對彼得來說是一個很重大的改變.
因為我們都知道.
一直看經文看聖經的時候.
知道彼得是一個相當傳統的猶太人.
所以如果他願意在這個俏皮城家裡住了這麼多天的時候.
很明顯是一個清楚的證據.
證明這個彼得已經克服了律法主義上的顧慮.
他們都讓他看見原來這個西門.
雖然他的工作是俏皮帳.
但是神仍然在他生命裡面有數不盡的恩典.
神不斷地去讓彼得一步一步地改變.
漸漸地由以前很堅持遵守的猶太的律法主義.
慢慢可以抽離進入神奇妙的恩典.
得著真正的自由.
這個釋放.
這個逐漸從那些沒用.
缺乏愛的禮儀律法主義中.
慢慢地被釋放出來.
得著真正的自由.
這就是神在我們面前的恩典.
從整個故事可以看出.
神即讓彼得有這樣的更新.

$^{601}$所以弟兄姊妹.
我不知道大家有沒有在你們的信仰歷程裡面.
經歷過一些神即般的改變呢.
如果有的話.
求神繼續通過不同的事情.
無造我們,更新我們,讓我們成為神一個手中寶貴的器皿.
如果我們從來沒有經歷過這些更新的經歷的時候.
讓我們通過今天的彼得的例子.
讓我們懂得放手.
讓我們懂得將生命交給神.
讓神一個偉大的生命無造者.
無造我們的生命,讓我們成為神一個合用的器皿.
成為我們在這個世界上其他人的一個祝福.
讓我們同低頭祈禱.
讓這個故事來不斷地提醒我們.
在我們今天走在人生的路上的時候.
我們需要你與我們的同在.
讓我們可以得著神你而來的寶貴的救恩的時候.
我們更加地不斷地更新.
不斷地在我們人生路上.
讓你成為我們的神跡.
來將我們的榮耀高舉.
多謝你給我們一切祈禱,奉主名求.
\newpage



\section{}
\label{sec:031607}
\textbf{聖經人聖經事:聖靈與火}~~~~黃克勤牧師博士
\newline
\newline
連結: \href{http://www.church.com.hk/acms/content.asp?site=cdc\&op=show\&type=product\&code=031607\&layout=sermon}{\texttt{031607}} ~~~~ 語音日期: 2026-01-28
\newline
\newline
\hyperref[sec:031606]{\small{< < < PREV SERMON < < <}}
~
\hyperref[sec:index_chronic]{\small{[返順時目]}}
~
\hyperref[sec:031608]{\small{> > > NEXT SERMON > > >}}
\newline
\newline
$^{1}$各位電影節目大家好.
又一齊讀聖經了.
生命說書的網絡.
跟大家一起看聖經人.
聖經的故事.
最重要是認識背後的上主.
今次的題目是聖靈與火.
仍然是接續著.
講施洗日漢的訊息.
就快講完了.
施洗日漢的訊息.
聖靈與火.
《馬太福音》三章十節開始.
現在虎子已經放在樹根上.
凡不見好果子的樹就撼下來.
凋在火裡.
我是用水給你們施洗.
叫你們悔改的.
這個水叫你們悔改的.
但那在我以後來的.
能力比我更大.
我就是給他提鞋也不配.
他要用聖靈與火.
給你們施洗.
這個就是今天要講的.
大家知道.
紅色字黑色字.
都是經文來的.
和合本經文.
這個光了.
螢光了就是我要講的重點.
分類分細的點.
留意第一點.
現在虎子已經放在樹根上.
虎子即是斧頭.
斧頭已經放在樹根上.
斬不是斬樹枝.
不是斬樹幹.
是放在樹根.
即是一斬整棵就沒有了.

$^{41}$連根拔起就沒有了.
樹就死了.
會斬哪些樹呢.
是不結好果子.
即是有果子.
但果子不行.
這樣都糟糕.
不結好果子的樹就撞下來.
凋在火裡.
現在是講時間到了.
各位猶太人.
各位以色列人.
各位神的子民.
大家留意.
現在的斧頭已經放在樹根上.
即是準備好.
要動手斬.
現在時間到了.
現在已經是那個時間.
這個斧子會對付什麼呢.
是不結好果子的樹.
不結好果子的樹都斬.
沒有果子的就斬.
因此要小心凋在火裡.
這個凋在火裡更加告訴我們.
其實這三章第十節是說.
審判就快到.
如果將第十節改了講.
就是告訴各位以色列人留意.
上主的審判馬上就到.
已經準備好.
審判馬上就到.
現在馬上就到.
第十節的意思就是.
斧頭預備好.
審判馬上到.
第十一節.
我用水給你們施洗.
這個水是用來象徵悔改.
用水施洗是象徵悔改.

$^{81}$你一定要悔改.
其實上面的好果子是什麼意思呢.
就是認認真真地結出好的悔改果子.
如果你不好好的認真地悔改.
作為果子修成.
結果斧頭就會劈下來.
現在已經到了.
所以你一定要好好悔改.
結出悔改的好果子.
你可以用水.
我和你施洗做一個.
你結出悔改好果子的開始與立止.
這個就是施洗向做.
不過他補充.
其實你真的要做.
因為接下來那個用斧頭劈的那個.
在我以後來.
我都算有少少能力.
他的能力比我更大.
你有什麼能力呢.
就是我有上帝的能力.
有聖靈感動.
一叫大家都來悔改.
這個真是很大的能力.
不過當然不是施洗有限的能力.
是聖靈的能力.
感動很多人去悔改.
他比我更大.
這個能力是指能力與權力比我更大.
我和他一比.
很多人來跟著施洗有限.
做老師老師拉比拉比.
其實現在有一個比我厲害很多.
比我大很多.
我和他一比.
我提鞋也不配.
我都講過很多次了.
僕人與奴隸有很多等級.
最低等的就是客人來到.
家人回來.

$^{121}$門口有幾個石缸.
舀水給他洗腳.
但是洗腳的都高級的.
最低級的僕人奴隸是做什麼的.
幫他解鞋帶.
解完鞋帶.
他就洗腳.
洗完腳之後他有屋.
幫他拿著鞋進去.
是和解鞋帶提鞋.
是最低級的僕人做的.
如果我和他比.
我連提鞋帶都未買到身.
我真的和他差很遠.
我只用水.
他用火.
他是用聖靈與火給你們施洗.
第一點就是.
預備好主的審判馬上就到.
所以第二點.
你要結出好的悔改果子.
不要想著這段聖經.
是和以前那些耶穌時代.
施洗有限時代.
那些猶太人說的.
不是.
今天同樣和你和我說的.
要悔改.
好好地結出悔改的果子.
否則就慘了.
好了.
審判就會到了.
我用水.
哇.
這一位.
跟著我來到比我大很多.
這一位基督尼徹.
他用什麼施洗的.
他一洗.
他用聖靈與火的.

$^{161}$哇.
厲害很多.
聖靈的意思是.
與聖靈和上主的聯合.
一舉一動.
與主聯合之後.
舉一動有新生的樣式.
是聖靈在我們裡面.
他來到.
他會帶給我們聖靈的能力.
聖靈在我們當中.
而且火.
火也代表審判.
這個火.
真的好像一把斧頭.
水就是施洗後帶來的.
耶穌帶來的是火.
和聖靈.
所以是不同的.
大家的洗禮是不同的.
羅馬書第六章.
就說到今天的洗禮是怎樣的.
豈不知我們洗禮是.
與主耶穌同死嗎.
然後呢.
同死同埋葬同復活.
一舉一動有新生的樣式.
是這樣的意思.
誰做的呢.
是聖靈將我們重生.
尼哥底姆問耶穌.
你真的很厲害.
你很有能力.
如果你的神蹟不是由神而來.
是無人能行.
耶穌就說.
我實在告訴你.
那個問題很小.
你必須要重生.
人這麼老怎麼重生.

$^{201}$鑽進貓肚生過嗎.
不是.
是聖靈風吹著意思吹.
就是聖靈.
好像風一樣.
將我們這個人改變了.
叫我們與主同死同埋葬同復活.
而且火是表示審判與煉證.
這個很重要.
這是聖靈會做的事.
留意師兄教我們什麼.
現在他手裡拿著波機.
有些人讀「播」.
波機要揚靜他的牆.
把墨紙收在瘡裡.
把腔用不滅的火燒盡了.
這句要特別解釋.
雖然不是我想講的三個重點之一.
三點講道之一.
拿著波機是什麼意思呢.
其實是他的審判.
怎麼審判呢.
就好像拿著一個波機.
這是什麼呢.
其實就像燒機一樣.
要揚靜他的牆.
怎麼揚靜呢.
拋起.
其實他們以前把墨紙.
在下面根的部分.
最貼地的部分割了.
一束束就打墨紙.
墨紙就掉了.
這些幹就沒用了.
就燒了.
墨紙一粒一粒在那裡.
他們就在地上磨.
用牛馬爐.
牛馬爐.
拉著.

$^{241}$兩個人坐在地上.
石頭在木板底.
硬硬的.
硬到那些墨.
墨的肉和殼分離了.
在地上研磨分離了.
那怎麼分呢.
那些殼就不要了.
就拿著一個波機.
拋起.
那些墨紙的肉重就掉了.
有風才做的.
一拋起風就把坑吹散.
被風吹散.
拋起就摺回那些肉.
拋起摺回那些肉.
他們要用這樣來揚靜.
整個和牆的墨紙.
研磨脫殼之後.
不會一粒一粒.
拋起風吹走坑.
然後把墨紙收好.
放在貨倉裡.
坑就不要了.
用火燒盡.
用來做燃料.
網上找到一個圖.
大家看看.
通常要站高一點.
風就大一點.
波機就把墨放進去.
有墨紙的肉和殼.
已經硬到脫了.
拋起.
或者倒下來的時候.
風已經把坑和墨的殼吹走.
這些就留來做燃料.
墨紙就可以堆在一起.
其實有點像.
我們中國人.

$^{281}$拋起這些.
谷糠和米.
這樣的帽子.
應該是越南人的.
我們中國人是用這種的.
拋拋拋.
揚靜他的牆.
不是洗刷.
整個和牆這麼多的谷糠.
沒有一粒會漏的.
都是他曾經拋過的.
表示這個審判.
是會來到的.
根據.
馬太福音三章十字.
一齊讀.
現在虎子已經放在樹根上.
凡不見好果子的樹.
就撼下來凋在火裡.
虎子已經放在樹根上.
在今天.
主都有他的審判.
對每一個人.
我們中國人都說.
時候未到.
若言不報.
時候未到.
其實神都忍耐.
不過他一定會審判.
不是因為中國人說了神就會審判.
而是中國人領會到.
如果你做的壞事.
你做的惡事.
未有報應.
是時候未到.
一定會報.
出來玩.
而他們會歡喜.
用虎子放在樹根上.
你記得.

$^{321}$這個也可以隱喻.
在今天來說.
主耶穌快要回來.
他回來的時候.
騎著白馬.
做一個審判得勝者的角色.
要記得.
虎子已經放在樹根上.
信不信?會不會呢?.
見到很多惡人惡事.
他很需要毀滅他.
要審判他.
我自己的惡很小.
可以容忍.
我自己容忍很多.
對別人不容忍.
但其實在上帝來說.
為何還未審判呢?.
我一審判你也沒有.
因此快些悔改.
虎子已經放在樹根上.
我們一起禱告.
天父上帝我們在你面前謙卑祈禱.
有時候我們會覺得.
以前那些猶太人很差.
法利賽人,薩魯哥人很差.
猶大賣主那些很差.
現在我們看這個世界上.
很多不公義,不敬虔.
很多人害別人很差.
很多人貪心很差.
無親情不解怨.
好說殘言很差.
不過其實我們自己.
多多少少都同樣是這樣.
不過我們裝扮得很好.
修得很好.
沒有人知道.
或者我們很多惡毒藏在心裡.
主要求你更新我們.

$^{361}$給你的靈性靈與火.
來煉證我們.
更新我們.
給我們更多合理的心意.
我們知道你的審判.
很快會來到.
而且馬上就來到.
給我們信靠你的審判.
給我們看到別人不好的時候.
我們默然交託.
求主幫助.
改變我們這個人.
給我們更多的悔改.
更多的更新.
誰聽祈求奉靠耶穌基督的名.
阿們.
\newpage



\section{哥林多前書 2:1-5}
\label{sec:031608}
\textbf{非靠口才,乃靠上帝的大能}~~~~林玉英牧師
\newline
\newline
連結: \href{http://www.church.com.hk/acms/content.asp?site=cdc\&op=show\&type=product\&code=031608\&layout=sermon}{\texttt{031608}} ~~~~ 語音日期: 2026-02-08
\newline
\newline
\hyperref[sec:031607]{\small{< < < PREV SERMON < < <}}
~
\hyperref[sec:index_chronic]{\small{[返順時目]}}
~
\hyperref[sec:031609]{\small{> > > NEXT SERMON > > >}}
\newline
\newline
哥林多前書 2:1-5
\newline
\begin{longtable}{cl}
\hline
\hline
章節 & 經文 (和合本修訂版)\\
\hline
2:1 & \begin{tabularx}{0.7\textwidth}{X} 弟兄們，從前我到你們那裡去，並沒有用高言大智對你們宣講神的奧祕。 \end{tabularx} \\ \\ \relax
2:2 & \begin{tabularx}{0.7\textwidth}{X} 因為我曾定了主意，在你們中間不知道別的，只知道耶穌基督並他釘十字架。 \end{tabularx} \\ \\ \relax
2:3 & \begin{tabularx}{0.7\textwidth}{X} 我在你們那裡時，又軟弱，又懼怕，又戰戰兢兢。 \end{tabularx} \\ \\ \relax
2:4 & \begin{tabularx}{0.7\textwidth}{X} 我說的話、講的道不是用委婉智慧的言語，而是以聖靈的大能來證明， \end{tabularx} \\ \\ \relax
2:5 & \begin{tabularx}{0.7\textwidth}{X} 為要使你們的信不靠著人的智慧，而是靠著神的大能。 \end{tabularx} \\ \\
[1ex]
\hline
\hline
\end{longtable}
$^{1}$今天要和大家分享的經文是.
《哥林多前書2章1至5至10至14節》.
其實經文應該去到16節.
不過因為我給錯了同工.
所以大家可以打開《新約傳書》第205頁.
那裡可以看完《哥林多前書》的15至16節.
首先讓大家一齊去禱告.
讓我們懷著恭敬的心 戰兢的心來到您的面前.
我們願意透過您聖靈在當中的開示.
讓我們開竅.
能夠明白您懷疑中的奧秘.
更加能夠明白.
上帝您對我們每一個獨特的人所說的話.
願聽的和說的都同夢您的月立.
祈禱奉主耶穌基督的名求.
阿們.
我相信大家都認識一個影視明星黃祖藍.
認識的意思是你認識她 她不認識我們.
我們認識她的.
我在想 如果你是她.
假設你都很喜歡港姐冠軍李亞男.
現在是她的太太.
當初你會不會向李亞男展開追求呢?.
會不會呢?.
哦 有人回答我.
可能我看不到他的眼睛.
可能很多人都不會.
可能你心裡有句說話跟你說.
我又沒有黃祖藍那麼有口才.
那麼有名氣 那麼多身家.
當然不會這樣做.
最怕的是什麼?.
最怕那些人跟我說.
你有什麼資格追求李亞男.
矮又矮的一個人.
又沒有名氣 又沒有身家.
你憑什麼去追港姐冠軍呢?.
喂 沒一塊鏡.
都拿一盆水照照自己的樣子.
很多時候我們心裡有這個心魔.

$^{41}$來搗亂.
告訴我們這裡不行 那裡不行.
你不合資格的 你不合資格的.
於是我們原本想去試一下.
都不敢.
包括我們很想我們身邊的親朋戚友.
好像我們一樣領受上帝給我們的福.
我們想去信耶穌.
不過就是不敢開口 是不是?.
今天我們想跟大家透過這段經文.
來看看保羅怎麼樣來鼓勵我們.
可能我們的心魔不斷告訴我們.
不要啦 人家拒絕你 那怎麼辦?.
人家問到你口硬硬 那怎麼回答?.
我們以為全方位好像百萬元卓的銷售員.
Top Sales.
這樣的人才有資格去傳福音.
不單止他要包裝 甚至他要很有口才.
或者他是一個成功的人士.
人家才會聽我們講 人家才會信.
不過保羅告訴我們.
沒錯 他都面對一班很崇拜名嘴.
很喜歡說話的人.
最好是神學家.
他都面對一班這樣的人.
不過保羅在第二章一字一開就說.
我不是這樣做的.
保羅去哥林多之前.
他經過希臘的雅典.
《使律人主》十七章告訴我們.
當地人很喜歡來探討一些新奇的事物.
彼此去討論 甚至激烈的辯論.
保羅都曾經和他們分享過關於鬼神和耶穌復活的事.
不過當地的人嘲笑他.
不太願意聽他講.
於是他來到哥林多城的時候.
他開始和人傳福音的時候.
他做了三個決定.
就在今次經文裡頭三節.
第一 第二 第三節.

$^{81}$他不用高賢大智.
他只講耶穌基督和他的十字架.
他雖然在軟弱和恐懼之中.
他仍然繼續去服侍.
保羅不是不會說話的.
大家都知道他受過很高深的學習訓練.
從小時候就受到很嚴格的思辨訓練.
不過他只是刻意放棄當時大家都崇尚的修辭技巧.
和一些高深的學問.
哥林多城鄰近希臘.
大家都知道希臘有三位很出名的哲學家.
蘇格拉底 柏拉圖和亞歷史多德.
我們後世很多的文學 哲學 藝術.
其實都是從他們那裡繼續發展出來.
當時哥林多城的人和希臘雅典城的人都差不多.
都很崇尚思辨邏輯 形意上的哲學思考.
很喜歡辯論 認為真理是越辯越明的.
他們更加想一些會用很華麗的修辭技巧去演講的人.
例如第一句開頭用個A字.
第二句開頭用個B字 如此類推.
令人讚嘆.
就好像當你聽到我們有人用詩詞歌賦.
吟詩作對來跟你分享的時候.
大家都會覺得那種優美 那種優雅實在令人欣賞.
又例如 大家都認識黃之華 許冠文這些才子.
他們博學 懂很多東西.
又幽默 又有心動 引人發笑.
能夠引起大眾的共鳴.
我們可能都會覺得全分辯局好像他們那樣.
那些人就會容易相信.
我們自己不要預我們了.
我們沒有這個資格去全分辯.
不過 保羅深信全分辯.
不是要用很高深的宗教術語才能夠令人相信.
又或者一定要咬文嚼字去令人驚歎.
保羅告訴我們 其實我們要講真話.
講出上帝怎樣在我們的生命裡祝福我們.
我們要跟別人分享你為什麼會相信耶穌.
又或者要跟別人分享在你這麼多生命的旅程裡.
上帝怎樣跟你同行.

$^{121}$原來最動人的往往是我們最真誠和真摯的經歷.
我們的生命故事.
保羅在傳福音的時候是講什麼的.
第二節他說 他納定主義.
不講深奧的哲學 只是講福音的核心信息.
就是耶穌基督和他釘十字架.
十字架這個很兇殘的凝聚 怎樣去拯救世人.
他怎樣被上帝的愛所感動.
上帝的愛怎樣去改變了他的生命.
以及他對事物的看法 我們大家都知道.
保羅看萬事如同糞土.
這個不是只有一個「講」字.
你見到他生命裡面都是這樣.
當然我們同行傳福音的時候.
都很容易被誘惑走入另一個狀況.
就是一些無謂的爭論裡面.
例如別人會質疑我們的信仰.
難道耶穌給你飯吃嗎.
難道耶穌給你工作嗎.
又或者挑戰我們的上帝.
上帝是怎樣做恐龍的.
聖經好像沒有說.
信任誰都無所謂.
每個宗教都是導人向善的.
好像這個地球的歷史很久很久很久.
亞當究竟是智人還是元人.
你說上帝是愛.
為何會有這麼多天災人禍.
如果我們走入了這個辯論的模式.
就會進入這個爭拗當中.
無論誰贏了.
都已經偏離了福音的核心.
就是有關於耶穌對世人的愛.
所以保羅立定心意.
去指光耶穌基督.
以及祂的十字架.
要相信的人不在乎人的智慧.
只是在乎神的大能.
保羅承認他初到哥林多城去傳福音的時候.
又軟弱 又懼怕 又戰戰兢兢.

$^{161}$相比名嘴和大哲學家.
保羅算得了甚麼呢.
在希臘雅典的時候.
正如我剛才所說.
不是很多人願意去聽保羅所說的話.
甚至嘲笑他.
這個江湖術士在說甚麼.
來到哥林多不知道會不會遭遇同樣的遭遇呢.
你看看 連大使徒都害怕.
所以我們害怕是很正常的.
不過他的選擇值得我們借鏡.
雖然在軟弱 恐懼和戰戰兢兢.
他仍然繼續去服侍.
繼續向人去傳講耶穌基督.
以及祂的十字架.
因為在恐懼裡面.
他更加全心去依靠上帝.
他一定要聽他講福音的人.
相信的人不是在乎他的智慧.
而是在乎上帝的大能.
保羅知道我們不用做到成功人士.
事事完美才去傳福音.
聖經說上帝最喜歡用軟弱的人.
因為軟弱的人才會懂得去依賴上帝的大能.
而不是依賴他自己個人的魅力或能力.
而那些相信耶穌的人.
也不是因為傳福音那個人.
他個人的魅力.
又或者他有很華麗的措辭表演.
來相信上帝.
其實上帝是會透過我們的不完美.
去感化 去感動其他的人.
保羅定義是要相信的人.
不是在乎人的智慧.
只是在乎上帝的大能.
我們或許會擔心.
如果對方我講完後他不相信怎麼辦?.
又或者我講完後他也不明白.
那為什麼要去傳福音?.
對於耶穌基督的福音.

$^{201}$為什麼有些人聽完也不明白呢?.
就是不肯相信.
保羅在這段經文裡解釋了當中的原理.
聖經說 上帝的事只有上帝的靈才知道.
當上帝願意將上帝的事.
包括神的隱密.
奧秘的事 去啟示給人知道.
那個人才明白.
我們可以用翻譯的app來作為類比.
聖靈願意翻譯屬靈的事給一個人知道.
那個人才會開竅.
才會明白有關屬靈的事.
有關上帝的事.
我相信大部分人都會在電話裡安裝一個翻譯的程式.
翻譯的app.
所以你去到世界各地旅行的時候.
你不需要擔心.
啊 那些文字我不懂.
餐牌是什麼來的?.
當然現在會有一個圖在這裡.
你可以開這個翻譯的app.
對著那幅圖畫.
對著那些字.
拍一下 哦 原來是這個意思.
先看看是什麼事.
哦 原來是不可以進入的.
這樣就很清楚很明白了.
又或者你對著一個跟你的母語不同的人.
你又不知道他在說什麼.
他又不知道你在說什麼.
你開了一個翻譯的app.
說了下去.
然後他就會翻譯對方的語言給他知道.
多麼方便 是不是?.
沒錯.
不需要再像以前那樣雞同鴨講.
大家明白了.
全方位就好像對著一個.
沒有裝翻譯app的人去說話.
他不知道你在說什麼.

$^{241}$因為他裡面沒有聖靈.
而聖靈也未曾開導他.
讓他明白關於神的奧秘.
和神隱密的事.
所以聖經說.
他一定覺得我們說的話很愚蠢.
啊 這些你都相信的?.
所以我們不用為對方不相信去負責任.
我們的責任是什麼?.
說.
真誠的分享在我們生命裡的故事.
信你怎樣在我們生命裡來彰顯.
怎樣去祝福我們.
為什麼我們會信耶穌.
以及在我們生命裡.
我們信耶穌之後有什麼不同.
又或者可能你從小就已經上了教會.
好像是伊洛那樣.
啊 沒有前後分別的.
但是相信我們在成長的過程裡.
你會看到上帝怎樣去祝福你.
怎樣跟你同行.
怎樣在你最軟弱的時候.
祂加力量給你.
你可以跟祂分享.
你跟上帝相遇.
上帝怎樣跟你同行的事蹟.
聖靈會隨著祂自己的意思.
向那些未信的人去作出翻譯.
去開導他們和啟示他們.
所以在傳福音之前.
我鼓勵大家.
在傳福音之前.
在傳福音的時間.
祈禱.
讓你練好你的口才.
更加重要.
第十四節保羅說.
屬血氣的人.
即是未信的人.

$^{281}$不接受上帝的靈的事.
他反倒以為如絕.
並且他不能夠了解.
因為這些事唯有屬靈的人才能夠領悟.
我們所傳講的唯有被聖靈開導的人.
才會明白.
所以我們要為未信的親友.
我們想跟他傳福音的朋友去禱告.
求聖靈開導他們.
在他們裡面安裝一個翻譯的應用程式.
以至他們能夠明白關於屬靈的事.
所以我們一定要成功帶人信耶穌這個包袱去放下.
我們的責任只是跟人分享.
聖靈開導過我們的事.
至於對方心裡面那粒種子.
何時發芽.
何時開花結果.
聖靈何時跟他翻譯.
這是上帝的工作.
是聖靈的責任.
我們只是跟上帝同工.
我們有幸成為上帝的左右手.
去祝福其他人.
我們要記住.
我們不是救贖主耶穌基督.
上星期我們上了蘇善美宣教師的課程.
《報道法》和《關係報道》.
蘇宣教師講了一個很好的比喻.
他說每一個人的胃口都不同.
有些人吃一個餅就飽了.
有些人吃十個餅.
或者二十個餅才飽.
我們不知道對方要吃多少個餅.
又或者他之前吃了多少個餅.
他才能夠明白上帝的奧秘.
才願意去相信耶穌.
我們不知道的.
我們的責任就是去分餅.
去傳講.
至於對方何時才飽.

$^{321}$是聖靈的責任.
所以如果你把餅帶過去的時候.
向他傳講福音.
他相信了.
我們當然很開心.
但我們也不需要驕傲.
因為在前面.
很多人做了很多事.
我們有幸是收過那個.
又或者當你把餅帶過去的時候.
他不相信.
他聽過很多次都不相信.
你也不需要氣餒.
因為你不知道他要吃多少個餅才飽.
又或者聖靈在哪個時候.
才能讓他完全明白福音.
其實我們和人分享福音.
不是要和他講一大套理論.
而是邀請他轉用一個全新的眼光.
去看這個世界.
經文第十節說.
只有上帝藉著聖靈.
將這些事向我們顯明.
十五節.
他說屬靈的人能夠看透萬事.
第十六節.
誰曾知道主的心.
誰會教導他.
至於我們.
我們有基督的心.
因為我們有基督的心.
我們有了耶穌基督的思維.
和祂的眼光.
所以我們看世界就不同了.
前幾天我們有機會和一些家長開祈禱會.
當中有些信了.
有些還沒有信.
大家都分享著大家所關心的事情.
有些說他們很關心.
他們的孩子的學業和他們的成長.

$^{361}$我很能夠明白家長的擔憂.
我們的孩子讀到書冠很高興.
但如果你有機會去看坊間的YouTube資訊.
你可能會很害怕.
坊間的資訊告訴你讀書沒有用.
讀大學也沒有用.
現在美國的公司都開始不聘請大學生.
再過兩年.
有80\%的工作將會被AI機械人所取締.
當然這些資訊.
其實沒有人知道是真是假.
但已經令到看的人心裡很膽怯.
又有另一個訊息.
不知道大家知不知道.
我們深水Po 有一間吃小菜的店.
我們吃得津津有味.
價錢又不貴.
這一位廚房裡有幾個人在炒菜.
誰知無線告訴我們.
是有三部AI炒菜機.
全自動的.
不需要人手就完成了.
接著又有無人駕駛.
連白領藍領.
都有機會被AI所取締.
另一邊有些資訊.
有些理財顧問告訴你.
生仔.
你準備好二百萬至八百萬了嗎.
這麼多嗎.
看你窮養還是貴養.
窮養也要二百萬.
如果你想退休.
你有沒有為自己預備了二百五十萬.
就要四百萬.
這個二百五十萬至四百萬.
是維持你基本的伙食.
交通.
而且你還需要為自己.
預備一個居所.

$^{401}$是不用租不用供的.
當你聽完這些訊息的時候.
你就會明白.
為什麼香港人這麼大壓力.
這麼多憂慮.
我們的信仰怎麼幫我們呢.
在重重的壓力和憂慮底下.
我們怎麼找到出路呢.
基督徒和微信的朋友都一樣.
都會遇到困難和苦難.
我們怎麼在這個大大小小的困難和苦難裡.
經歷上讓我們看到出路.
看到盼望.
看到平安呢.
在另一方面.
或者我們每個人都有機會.
跟別人有爭執有頂撞.
在這些爭執和頂撞裡.
我們不是彼此去報仇.
割席.
互不理睬.
而是願意有基督的心.
藉著聖靈的幫助.
我們去饒恕.
和帶來真正的悔改和復活.
在你身上這些真實的故事.
正正展現了你有基督的心.
有不一樣的眼光去看這個世界.
和靠著聖靈作出不同的行動.
讓你和未信者.
雖然現在在相同的處境底下.
在難關在焦慮的裡面.
你仍然擁有一份真實的平安和盼望.
相信我.
你真實的生命故事.
比長篇大論的講道.
更加吸引和感動人.
哥林多前世第二章告訴我們.
要使你們的信不靠著人的智慧.
而是靠著上帝的大能.

$^{441}$福音的大能不是建立在人的智慧上.
而是建立在上帝的大能上.
所以鼓勵你.
也邀請你.
在這個星期找一個.
你一直很害怕和他講福音.
但你又很想和他講福音的人.
去分享你生命的故事.
邀請他來《量髮Cafe》.
順便和他說.
其實我也很害怕.
我不是很會說話.
不過我真的想和你分享.
耶穌怎樣幫過我.
和他分享之前.
你先祈禱.
靠神使用你.
更求聖靈在對方的心裡面.
安裝一個翻譯的應用程式.
讓他明白上帝的奧秘.
當你說完之後.
禱告將一切結果交給上帝.
好嗎?.
讓我們一起努力.
去興旺福音.
讓我們同心低頭祈禱.
親愛的天父多謝你.
因為是聖靈的啟示.
我們才能夠明白屬靈的事.
明白基督的愛.
以致我們可以稱呼你為「阿爸父」.
在充滿壓力和誘惑的城市.
有新的眼光去看周圍的事.
帶領幫助我們.
我們懇求你去使用我們.
成為你福音的使者.
可以向身邊未信的親友.
鄰舍去傳講.
你在我們身上的作為.
你怎樣幫過我們.

$^{481}$天父啊.
在未來一個星期.
我們願意作出行動.
不過我們很害怕.
很軟弱.
多謝你提醒我們.
當我們全心去依賴你的時候.
就能彰顯你的能力.
因為人的信不是靠著人的智慧.
而是靠著神的大能.
所以懇求上帝你使用我們.
成為你恩典的渠道.
我們這樣禱告直至耶穌基督的名而求.
阿們.
\newpage



\section{約翰福音 13:1-5}
\label{sec:031609}
\textbf{耶穌洗門徒的腳:魔鬼在細節裡}~~~~陳榮超博士
\newline
\newline
連結: \href{http://www.church.com.hk/acms/content.asp?site=cdc\&op=show\&type=product\&code=031609\&layout=sermon}{\texttt{031609}} ~~~~ 語音日期: 2026-02-08
\newline
\newline
\hyperref[sec:031608]{\small{< < < PREV SERMON < < <}}
~
\hyperref[sec:index_chronic]{\small{[返順時目]}}
~
\hyperref[sec:031610]{\small{> > > NEXT SERMON > > >}}
\newline
\newline
約翰福音 13:1-5
\newline
\begin{longtable}{cl}
\hline
\hline
章節 & 經文 (和合本修訂版)\\
\hline
13:1 & \begin{tabularx}{0.7\textwidth}{X} 逾越節以前，耶穌知道自己離世歸父的時候到了。他一向愛世間屬自己的人，就愛他們到底。 \end{tabularx} \\ \\ \relax
13:2 & \begin{tabularx}{0.7\textwidth}{X} 晚餐的時候，魔鬼已把出賣耶穌的意思放在加略人西門的兒子猶大心裡。 \end{tabularx} \\ \\ \relax
13:3 & \begin{tabularx}{0.7\textwidth}{X} 耶穌知道父已把萬有交在他手裡，且知道自己是從神出來的，又要回到神那裡去， \end{tabularx} \\ \\ \relax
13:4 & \begin{tabularx}{0.7\textwidth}{X} 就離席站起來，脫了衣服，拿一條手巾束腰， \end{tabularx} \\ \\ \relax
13:5 & \begin{tabularx}{0.7\textwidth}{X} 隨後把水倒在盆裡，開始洗門徒的腳，並用束腰的手巾擦乾。 \end{tabularx} \\ \\
[1ex]
\hline
\hline
\end{longtable}
$^{1}$多謝高牧師的安排.
我感到十分榮耀.
因為能夠分享崇基禮拜堂的港台.
在我心目中的感受.
遠超過各位弟兄姊妹朋友的感受.
你看見我只不過是一個港元.
但我站在這個港台的時候.
真的是戰戰兢兢.
這個港台不是一個簡單的港台.
因為崇基禮拜堂,崇基學院.
在香港的歷史當中.
是佔一個很重要的角色.
今年如果大家有留意崇基學院的新聞的話.
你就知道今年是崇基學院慶祝75週年的慶典.
過去的75年比中文大學的歷史還要久.
75年不是一個短暫的時間.
所以崇基學院大我一年.
所以我和崇基學院都是親兄親弟.
我們在地上一同成長.
所以今年聽到崇基學院的75週年.
就是我這個哥哥的生日.
我這個做弟弟的真的感到很榮幸.
所以那個嘉年華會我也衝下去.
1月底的時候也去參加去感受一下.
所以崇基學院的成立.
這個禮拜堂的存在.
在香港的基督教歷史上.
是一個很重要的傳統.
因為傳統傳下傳下就可能不是很統.
因為很多時間的改變.
政治,經濟.
我們活在一個傳統當中.
我們有時久不久要重新去想一想.
這個75年裡面.
崇基學院帶給我們每一個在香港生活的人.
我們成長當中的經歷.
所以我覺得傳統是我們需要有時.
我們要再去想一想.
我們再去看看.
我們更新.

$^{41}$我們重溫一下.
以致我們對於傳統本身的意義.
是更加的深刻.
所以在這個傳統當中.
我就揀選了一段經文.
剛才列出來的.
就是一個很傳統式的一段經文.
就是耶穌和門徒洗腳.
在座的弟兄姊妹們.
我相信對這段經文.
或者對這個故事是不陌生的.
甚至乎就算不是基督徒.
都聽過在耶穌洗腳這個事情上.
所以今天早上我和各位弟兄姊妹.
我們從基督教很傳統的故事當中.
我們再去想一想.
這個故事.
這個傳統.
在我們今天信徒的生命當中.
我們應該怎樣去回應呢?.
這個榜樣.
我們經常覺得耶穌洗腳這個榜樣.
那個意義是什麼呢?.
我可以說.
今天這個重溫.
可能會帶給我們一些信徒.
對這段傳統會有一些驚訝.
或者有些過去沒怎麼想過的方向.
因為傳統能夠更新.
就是對我們生命一個很大的鼓勵.
我們待會和大家讀這段經文之前.
有幾個細節我想和你們做一些比較.
以致這個經文能夠更加凸顯它的特色出來.
這個洗腳的地點.
就是最後晚餐的房間.
這一點我們首先要把握.
但是這個洗腳的活動.
在最後晚餐的房間.
在另外那三本福音書裡面.
是沒有提及過的.

$^{81}$在那三本福音書裡面所提及過的.
我們認為很重要的東西.
在約翰福音裡面是沒有提及的.
譬如我們最後晚餐.
我們大家都熟悉.
或者我們每一個月都會去守這個聖餐.
最重要的兩個聖品就是這個餅和這個杯.
所以在三福音裡面.
就是很清楚地去畫出來.
約翰福音就沒有怎麼提及這件事.
我們覺得最後晚餐.
那件事大家都說了.
其實約翰福音裡面.
是有把這些物品.
有不同的方式顯露出來.
只不過沒有三福音那麼直接跟我們講.
這個第一個比較.
引起我們對這個經文的一些興趣.
所以我們覺得有意思.
在同一個房間,同一個空間裡面.
竟然有一個活動.
我們可能在領受最後晚餐的時候.
我們沒有怎麼去留意.
我希望今天我們能夠帶出多一點.
令到我們下一次再去領受這個餅和杯的時候.
我們心裡面還多一點的思考.
所以餅和杯的祝福我們知道.
耶穌洗腳我們都知道.
只不過我們過去沒有將這兩樣東西平衡起來.
一起去思考.
將它們放在一起.
同一個空間的時候.
所以約翰在約翰福音裡面.
就給我們一個特寫.
將這個洗腳的事情更加顯露出來.
藉此我們看看經文.
禹節以前.
這個就是最後晚餐之前的日子.
耶穌知道自己離世歸附的時候.
到了祂既然愛世間屬自己的人.

$^{121}$就愛他們到底.
吃晚飯的時候.
最後晚餐.
魔鬼已將賣耶穌的意思放在西門的兒子.
加略人猶大的心裡.
耶穌知道父已將萬有交在他手裡.
且知道自己是從神出來的.
又要歸到神那裡去.
就離席站起來脫掉衣服拿起一條手巾束腰.
隨後把水倒在盆裡.
就洗門徒的腳.
並用自己手足的手巾去擦乾.
所以在這段經文當中.
最突顯出來的.
讓我們看到一個動作.
就是耶穌去洗門徒的腳.
這個動作為何發生呢?.
是不是無緣無故發生呢?.
或者祂想表達一些什麼意思呢?.
後期教會將最後晚餐上.
兩個動作成為一個聖品.
這個杯,這個餅是聖品.
和這個洗腳禮.
或者叫做作足禮.
都是成為歷史上一個聖品.
或者一種聖禮來處理的.
如我剛才說過.
如果我們作比較的話.
三福音就將餅和杯.
成為最後晚餐的中心點.
而漢福音就將耶穌洗腳這件事.
成為最後晚餐上最顯著的一件事.
令我們去思考.
所以各位聽眾姐妹.
我和大家將經文處境作一個比較的話.
我希望引起你心裡面.
可能對這兩個沒有想過.
或者想過也不一定.
我們可以一同去參考一下.
我們一同探討一下.

$^{161}$這兩個動作.
都是最後晚餐上最重要的一個環節.
因為祂帶出來的訊息.
從我們基督教的救恩來說是很震撼的.
耶穌以祂自己的身體.
用餅代表祂的肉體.
用杯代表祂裡面的血液.
血淋淋的赤裸裸的告訴我們.
救贖就是要這樣做.
但漢福音就揀選另外一個角度.
它都是和救贖有關的.
只不過它的角度不一樣.
它就說在救贖當中.
各位群體姐妹,各位朋友,各位門徒.
我們更加要留心.
我為什麼要和你們洗腳.
以致我們不會太驚訝發生什麼事.
我很驕傲地和各位弟兄姊妹.
擺一個打卡點給我們看看.
我們有張崇基的打卡點.
如果你沒去過這裡打卡的話.
待會完結之後你自己走上去.
就在我們神學院.
我們崇基學院神學院.
我們神學樓前面有銅像的雕塑.
就是我們前兩任院長盧龍光牧師.
特別他跟我說.
陳牧師我告訴你.
這個銅像很高的.
我想有六,七呎.
你站過去根本就高過你.
在南京特別打造運入香港.
為慶祝我們新聖堂的成立.
那時候是擺設的.
這是我們崇基中文大學的打卡點.
所以你沒去過的話.
待會完結之後一定要去那裡.
站在那裡.
或者坐在彼得的腳前面.
耶穌幫你洗腳.

$^{201}$所以我都很驕傲.
我作為老師和朋友分享.
你在哪裡教書.
大學,崇基學院,神學院 去看看就知道.
有一個很美的銅像雕塑在那裡.
我們院長的意思.
當然是勸導我們神學生.
你來讀神學.
要學什麼呢.
在洗腳的傳統信息.
或者一般我們認同的信息上面.
或者甚至在世界的文學上面.
對於耶穌的洗腳.
都有一般的認同.
謙卑,愛和服侍.
犧牲,屬靈的潔淨.
服侍的榜樣,彼此相愛彼此服侍.
我數了一堆出來.
Google一下都全部出來了.
沒什麼特別深奧的.
世界認同耶穌洗腳.
是一個謙卑,愛和服侍的榜樣.
約翰提出這一點更加理所當然.
因為他覺得自己是愛的門徒.
所以他將耶穌的愛.
透過洗腳.
更加可以表達自己個人的意思.
在這裡我想停一停.
想一想.
是不是就是這麼簡單呢.
是不是就是為了傳達一個謙卑的信息.
而這麼大件事將杯和餅都不顯露出來.
而是將洗腳顯露出來呢.
我覺得在這點上我們應該停一停.
想一想.
是不是就是單單為了謙卑這麼簡單呢.
其實世界上不同的宗教.
不同的哲學.
都去推崇謙卑.
其實我們有什麼了不起呢.

$^{241}$我們又說這件事.
除了謙卑我們互相傷害之外.
我覺得如果單單從一個.
我們要學習做好人.
我們要學習做一個好基督徒.
單單從這個角度我們去看這段經文的話.
各位電影節目的朋友.
我覺得我們低估了其中的屬靈的價值是什麼.
耶穌為門徒洗腳的意義.
我相信不只是教導門徒他們謙卑.
而是為什麼要謙卑.
為何要謙卑.
因為我們做事.
沒理由無緣無故找這些來做的.
是啊,很修養嘛.
這裡樣樣東西都可以修養的.
何必這麼大件事呢.
所以在這個轉彎點上.
各位電影節目的朋友.
我想和大家再深入去看看這段經文的記載.
在這個記載當中.
特別從約翰福音的揀選.
去表達到耶穌的救贖.
在最後的晚餐上面.
祂洗腳來做這個環節.
意義在哪一處.
像剛才說過.
禹節的晚飯和最後晚飯的連接.
餅和杯.
我們都在猶太人的歷史上.
可以找到源頭.
因為最後的晚餐是由禹節.
我們在舊約可以看到很多蛛絲馬跡.
甚至基督教我們都會嘗試.
我們有些教會會.
我們自己組織一個禹節的晚飯.
去經歷一下當年埃及人.
他們要去出埃及的時候.
或者以後他們要去紀念他們的神的救贖.
我們都可以找到很多蛛絲馬跡.

$^{281}$每一樣的食物 每一樣的餅和杯.
有四個杯 哪個杯打哪個杯.
我們可以找到很多的細節.
但是在最後晚餐的禹節的過程.
那個傳統當中.
我們是找不到洗腳這件事出來的.
所以這個又令我們再想一想.
耶穌的這個行動.
是和一般的禹節我們所預測的.
似乎有一點點好像不是很對勁.
當然猶太人他們洗手洗腳.
這個是經常的習慣.
但是在晚餐的中間還沒吃完的時候.
突然間加插這個動作.
這件事值得我們去想一想.
有什麼事情發生呢.
因為沒有一件事情的發生.
是無緣無故發生的.
我相信你和我都相信這件事.
我們沒有無緣無故去做一些事情.
一般我們要揀選一些重要的事情來做的時候.
沒有啊 因為無緣無故.
如果我們要去做出來的話.
因為一定有些事情我們是要針對的.
這個就是我這個信息.
其中下一半就說耶穌為門徒洗腳.
但是魔鬼就在這個細節當中.
真是魔鬼在這個細節當中.
我們把經文分開來.
我們再看看它裡面有什麼意思.
《禹節》我們有一段經文.
我們回到經文.
我們把經文看到中間有些是灰暗了的字.
所以你們不需要讀.
我們把經文連串起來.
我們找到意思.
在《禹節》之前耶穌知道自己離世.
到第三節耶穌知道.
父將萬有交在祂的手裡.
然後祂說知道.

$^{321}$這段經文我畫了幾圈出來.
你們可以看到不是太難.
從救贖當中.
耶穌知道祂離世.
祂知道救恩快要成就.
祂知道祂上十字架.
知道知道知道.
是很清楚的.
就是說這個救贖.
如果我換一個角度去說另外一個信息.
總之這三個知道.
都可以說出一篇信息.
這個是很直線的去看這段經文的意思.
但我想換一個角度.
把灰暗了的字顯現出來.
我們下一張看看.
灰暗了的字你不要看了.
我把深色的字.
因為歌句的說話很有意思.
因為歌句的說話是加插進去的.
就是說約翰在記載這件事情的時候.
在記載下來之後.
突然間他覺得.
有一樣東西他不知道放在前還是放在後.
另外開一張.
他在這個轉彎點裡面.
我們叫做加插.
加插不是說以後才加上去.
只不過在過程當中.
他發覺到有一樣東西.
都不知道放在哪一處.
他很尷尬的時候.
他放在這個中間裡面.
這個曲線去傳承這段經文.
魔鬼以將賣耶穌的意思.
放在西門的兒子嘉良人猶大的心裡.
這個是其中一個的門徒.
這個曲線的加插句.
令我們對整段的經文.
浮現出一個不同的思念.

$^{361}$或者是有一種的幻想.
或者是有一種的意境.
在整個漫畫裡面.
突然間浮現了一樣東西.
我們不知道將它放在哪一個位置.
但是不說它又不行.
我忘記了第一節有沒有事.
我們會有些想表達的心意.
你真的不知道放在哪一處.
我們可能將它放了出來.
但是它帶來的影響力.
可能超過你和我想.
所以這個叫插入語的句子.
它不是去打斷這個主題.
它只不過在這個主題當中.
它要去更加引人注目.
我們肉眼未必看到的事情.
而是肉眼之外.
讓我們看到另外一個層面.
在約翰福音書的表現.
就是在一個黑暗的背後.
我們要看出一個真實的現象.
所以這個不是一個普通的節日晚餐.
如果門徒跟隨耶穌三十年.
我們猜測聖經沒有怎麼說.
這個也不是他們首次如此的晚餐.
不過這次.
在最後的晚餐當中.
這個表達是以前沒有的.
因為這個加插語.
是帶著一種怎樣的感受.
殺機重重的感受.
殺意重重.
隱藏在背後的一種殺機.
我把這段經文再拉闊一點.
我們再去看一看.
由第二章推到很前.
第二節我們用二十七節來收尾.
將這兩節夾在一起.
讓我們看整個行動的大環境.

$^{401}$因為第二節就說.
吃晚飯的時候.
魔鬼已將賣耶穌的意思.
放在西門的兒子嘉掠人猶大的心裡.
誰出現了?.
魔鬼出現在這段經文當中.
但因為我們對洗腳太興奮.
我們可能沒有留意到.
原來在這個環境當中.
在黑暗的背後.
魔鬼已經在這個御節的晚飯當中存在.
跳到二十七節.
吃完飯.
猶大吃了以後.
撒旦就進入了他的心.
整段御節的晚飯包括洗腳的過程.
就包含在這兩節當中.
這兩節的經文指出了一件事.
背後的黑暗環境不是人般的.
還有兩個很特別的字出現.
第二節是魔鬼.
第二十七節是撒旦.
陳牧師,這不是一件事嗎?.
是一件事.
不過.
用魔鬼來開始.
用撒旦來形容.
就是從一個很隱約的警方.
去到一個赤裸裸的.
意思是開始的時候.
魔鬼都比較客氣.
躲起來,遮遮掩掩.
戴上hood.
到撒旦的時候.
兩隻角,插著兩隻叉.
已經很赤裸裸.
完全沒有任何禁忌.
我就是撒旦.
因為撒旦是一個名字.
魔鬼是一種勢力.

$^{441}$單從這兩個不同的字的挑選.
我們對於殺機重重的這段經文.
心裡面寒一寒.
用一個肉眼看不到的背景.
一個黑暗的背景.
門徒吃得很開心.
他不知道發生什麼事.
但他很開心.
還說起上個星期.
我們的主耶穌騎著雷子.
威威風風的.
走進耶路撒冷的城門.
沿途有很多人拿著棕櫚樹.
多威風.
門徒一邊吃羊腿.
一邊傾談.
喝一杯酒.
很開心.
約翰看到的情況.
耶穌看到的情況.
是不一樣.
當時他還想起.
可能猶大還很怒氣.
瑪利亞有沒有搞錯.
打破了這麼貴的珍娜達香膏.
過去一個多星期有很多活動.
在回大.
在傾談.
興高采烈的時候.
但耶穌看到你們吃歸吃.
你們說歸說.
但我看到後面有勢力在.
我們不可以用.
這是一個屬靈爭戰的時候.
因為魔鬼的勢力.
從耶穌的眼光看.
近在微節.
在殺氣騰騰之下.
所以約翰將杯和餅放下.
其他三福音講得比較清楚.

$^{481}$約翰在這裡將杯和餅放下.
他有一件更重要的事情.
要去處理.
經文引入魔鬼.
煽動侵蝕的勢力.
魔鬼的存在就是一個貪腐.
將謊言講一千次.
我們就當以為是對的.
我們歪曲整個歷史的過程.
惡魔.
邪惡.
魔鬼和撒旦的現實的存在.
已經是赤裸裸的.
現在在十二個門徒當中.
他就出現.
所以約翰福音這個說話.
不只是單單.
我們基督徒做個好人.
我們要學習謙卑.
我們要服事眾人.
各位聽眾.
這是一個殺機重重.
這是一個屬靈爭戰的時候.
那種魔鬼的氣味.
那種氣氛.
到後來我們知道.
到猶大.
他失守.
然後失守.
在這樣的過程當中.
耶穌洗腳這個動作.
顯著出來的意義.
是重大的.
撒旦的謊言.
在曠野裡面.
耶穌領受過.
但是對於耶穌來說.
他應付撒旦的攻勢.
是綽綽有餘.
所以三次如何用騙局來騙他.

$^{521}$石頭變麵包.
擁抱權勢.
擁有所有榮耀.
耶穌說你不就說了.
我有神的話語.
耶穌在曠野被試探.
這個故事傳統我們都知道.
我們覺得耶穌說輕而淡寫就搞定.
但是這次撒旦要針對的.
是他的門徒.
這是耶穌所擔心的.
你來攻擊我.
Make my day.
耶穌說.
你來攻擊我.
來呀試試.
但是這次看到他攻擊他的門徒.
他看到門徒的軟弱.
他看到門徒很容易受到誘惑.
和受到欺騙.
所以耶穌面對一個這樣的情況來說.
是一個很緊張的時刻.
是一個非尋常的挑戰.
所以面對一個非尋常的挑戰.
要採取一個非尋常的方法.
才可以對付到.
這就是耶穌開始有行動的時候.
離席脫衣拿著手巾.
然後將水倒到盆裡洗耶穌的腳.
如果你將第四節經文.
一口氣急速地讀出來的話.
離席脫衣束腰洗腳.
這個行動各位弟兄姊妹.
比西部牛仔電影拔槍那一下動作.
還要迅速.
你看到一個這麼緊迫的急在眉睫.
撒旦宮的攻勢要臨到門口的時候.
耶穌立即站起來.
就用屬靈的增箭去接招.
因為祂要保護的是他的門徒.

$^{561}$好像母親去保護她的小雞一樣.
這四個動作是一氣呵成.
算是一個作戰的狀態.
對於十二個門徒來說.
他們都是呆了的.
從來未試過洗腳.
不但未試過洗腳.
不但未見過他們的主耶穌.
那個動作,那個迅速,那個快速.
還有那種肯定性.
十二個門徒看著耶穌做那幾個行動.
站起來.
脫下衣服.
束上腰.
倒盤水.
然後洗腳.
一連串的行動.
又沒有解釋.
又沒有前又沒有後.
門徒坐在那裡.
個個都是口頂目外.
所以我們有彼得的說話.
不可以搞我.
所以我們可以看到.
門徒在耶穌的工作上.
我們有彼得哥哥很戲劇化地拒絕耶穌洗腿.
當中我們有些教導.
但整個情況不是很和善,和諧,慢節拍.
是一個急不及待,不能再等待的時刻.
我們的謙卑不是為了謙卑而謙卑.
如果我們只是為謙卑而謙卑.
我們抵蠱了耶穌洗腿的重要.
所以在這個動作上我們隱藏著.
和這個對話有兩個很重要的意思.
魔鬼的計劃蠢蠢欲動.
但同一時刻,國度權並榮耀要降臨的時候.
各位姐妹摸到這個時刻,這個時間點.
再重覆一次.
魔鬼的計劃已經蠢蠢欲動.
預備要破開天國的門.

$^{601}$要擄掠所有神的子民.
但同一時刻,國度權並榮耀要降臨的時候.
雙邊的訊息就在這一刻.
耶穌站起來為門徒洗腿顯現出來.
表面上耶穌洗腿是和門徒設立一個謙卑服事的榜樣.
但實質上耶穌洗腿針對的對象.
是虎視眈眈的魔鬼和撒旦.
在洗腿上,耶穌拍台和撒旦說.
「我和你畫一條紅線」.
「你搞我的門徒,你是沒辦法可以搞他們的」.
「因為我今天和你畫一條紅線」.
「這條紅線就是我和他們洗腿」.
在打破撒旦出現的謊言.
他知道門徒是軟弱.
撒旦的滲透是無處不在.
從門徒的口述當中,我們都可以摸索到.
門徒都是蠢蠢欲動地進入魔鬼的教堂.
有些人說「我想坐在耶穌的右邊」.
「我要坐在耶穌的左邊」.
包括當時彼得說「你不可以洗我的腿」.
我們今天回顧一下,我們查經.
彼得那些傻子,如果你是彼得,你當時坐在那裡.
而是你爆出這樣的說話.
因為你要怎樣去爆這個說話出來呢?.
耶穌說「我也可以負責的」.
「你不用洗我的腿」.
「我也可以做大佬的」.
即是在他們心目中的當中.
他們對於洗腿,根本摸不清楚耶穌在這個傳統上.
所以耶穌就在這裡告訴他們.
「我不洗你的腿,你就與我無分」.
然後他們又不停地說.
「主啊,你連我的頭也洗了」.
耶穌說,然後一個很精確的金句就出來.
「凡洗過躁的人,只會把他洗的」.
這些我不多說了.
但在整個過程中,我想帶出一個論點.
讓我們去思考,耶穌整個過程不是為了什麼.
祂就是為了應付撒旦的攻擊.
整個過程是要突破魔鬼的詭計.

$^{641}$曾經說過,若有人做頭的,要做後的.
若有人要做領導的,先要做僕人.
將訊息混入在一個屬靈的爭戰當中.
領袖不是上一課MBA這麼簡單.
不是上一課教會增長學這麼簡單.
而是整個我們不留心的話.
撒旦魔鬼就已經在我們門口.
所以我們要突破這個詭計.
以致我們學習另外一種屬靈的增長.
在最後的晚餐當中,一共有12位門徒.
如果你數這個,你只能數到11個.
為什麼我特別要數12個呢?.
因為這幅畫很有意思.
他將猶大那個第12個.
如果你一眼看下去,不數還可.
你只能數到11個.
因為第12個已經在門口裡.
向著黑暗將耶穌出賣的計劃已經向前走了.
所以他已經在門口當中.
所以這一幅油畫.
更加帶給我們一個感受.
這個黑暗更急在眉睫.
當猶大出去了之後.
耶穌就跟其他11個門徒說.
如今人子得榮耀,記得嗎?雙重.
魔鬼在作弊的時候.
同一時候國度降臨.
神在人子的身上又得你榮耀.
神要因自己榮耀人子.
並且要快快地榮耀他.
然後一句我們眾人皆知的說話.
我要賜給你們一條新命令.
乃是叫你們彼此相愛.
我怎樣愛你們你們也要怎樣相愛.
你們若有彼此相愛的心.
眾人就認出你們是我的門徒了.
這就是紅線.
這就是從撒旦畫的一條紅線.
跟門徒說你們要去做這件事.
你們彼此相愛.

$^{681}$因為彼此相愛這條紅線就會清楚了.
就是別人都知道你是我的門徒.
因為回應洗腳謙卑的服事.
彼此相愛祈求合一.
但不是單單是為了合一而合一.
而是因為背後有一個更重要的意義和方向.
去引導我們去幫助我們.
在結束的時候我用保羅的一句說話.
在《爾弗所書》六章提醒我們.
要穿戴神所賜的全副軍裝.
就能抵擋魔鬼的詭計.
因我們並不是屬血氣的征戰.
乃是與那些執政的掌權的管轄的.
這個幽暗世界的以及天空屬靈氣的惡魔爭霸.
保羅很清楚.
他打仗的戰場不是在教堂裡面.
不是在信徒信徒之間的爭執.
不是在教會教會之間的破裂.
他不是在信徒信徒之間.
牧師牧師之間的不和.
他打仗的戰場.
是因為外面黑暗的世界.
他的勢力就伏在我們門口裡面.
但耶穌畫了一條紅線.
這個撒旦邪靈的魔鬼是進不來我們教會的.
他只不過伏在門口裡面.
但我們要懂得怎樣去管理他.
我們要懂得怎樣去控制他.
我們怎樣以這個合一.
我們怎樣互相傷害.
我們要互相洗腳.
不是單單我們覺得和諧和和善.
而是因為如果不是這樣的話.
我們就會中了這個詭計的策略.
所以各位弟兄姊妹.
這個洗腳.
不需要在一些毛巾與布的行動上面.
他帶來的意義是更深的層次.
就是讓我們更加明白到.
我們今天所看到的處境.

$^{721}$因為這個黑暗的勢力.
是用一個侵蝕,用一個滲透,用一個謊言.
在我們當中運行.
但當我們合一的時候.
我們是應付黑暗的來臨.
而神的榮耀會降臨在我們當中.
我們一同低頭禱告.
但你不是單單因為我們要團結.
而叫我們去做這個動作.
你是告訴我們因為你所看到的.
是一個危機的時刻.
撒旦和他的流落是虎視眈眈.
但當我們今天我們的肢體.
我們能夠聯合在一起.
我們這個合一互相相愛的時候.
我們就是一個德性的群體.
因為主是一個德性的主.
我們為此感恩.
我們奉主名禱告.
阿們.
\newpage



\section{馬太福音 6:9-15}
\label{sec:031610}
\textbf{求主教導我們禱告}~~~~石偉漢牧師
\newline
\newline
連結: \href{http://www.church.com.hk/acms/content.asp?site=cdc\&op=show\&type=product\&code=031610\&layout=sermon}{\texttt{031610}} ~~~~ 語音日期: 2026-02-08
\newline
\newline
\hyperref[sec:031609]{\small{< < < PREV SERMON < < <}}
~
\hyperref[sec:index_chronic]{\small{[返順時目]}}
~
\hyperref[sec:031611]{\small{> > > NEXT SERMON > > >}}
\newline
\newline
馬太福音 6:9-15
\newline
\begin{longtable}{cl}
\hline
\hline
章節 & 經文 (和合本修訂版)\\
\hline
6:9 & \begin{tabularx}{0.7\textwidth}{X} 「所以，你們要這樣禱告：『我們在天上的父：願人都尊你的名為聖。 \end{tabularx} \\ \\ \relax
6:10 & \begin{tabularx}{0.7\textwidth}{X} 願你的國降臨；願你的旨意行在地上，如同行在天上。 \end{tabularx} \\ \\ \relax
6:11 & \begin{tabularx}{0.7\textwidth}{X} 我們日用的飲食，今日賜給我們。 \end{tabularx} \\ \\ \relax
6:12 & \begin{tabularx}{0.7\textwidth}{X} 免我們的債，如同我們免了人的債。 \end{tabularx} \\ \\ \relax
6:13 & \begin{tabularx}{0.7\textwidth}{X} 不叫我們陷入試探；救我們脫離那惡者。因為國度、權柄、榮耀，全是你的，直到永遠。阿們！』 \end{tabularx} \\ \\ \relax
6:14 & \begin{tabularx}{0.7\textwidth}{X} 你們若饒恕人的過犯，你們的天父也必饒恕你們； \end{tabularx} \\ \\ \relax
6:15 & \begin{tabularx}{0.7\textwidth}{X} 你們若不饒恕人，你們的天父也必不饒恕你們的過犯。」 \end{tabularx} \\ \\
[1ex]
\hline
\hline
\end{longtable}
$^{1}$弟兄姊妹主內平安.
我們今天一齊來到崇拜當中.
今天和大家分享的題目是.
求主教導我們禱告.
都是問一個很基本的問題.
大家懂不懂祈禱?.
只有Winbo領導.
一會兒一起教Winbo.
我們說祈禱也不是第一次說.
我們偶爾也會說祈禱的事情.
所以很有趣.
我們也會問一個問題.
牧師我還需要學習祈禱嗎?.
我回教會這麼久了.
我也信了耶穌這麼久了.
祈禱的事情天天都做.
還需要教嗎?.
我今天希望不會嚇到你.
今天我講的經文可能和你平時的理解不一樣.
扣好安全帶.
我們準備經歷重新思考聖經的過程.
我們不懂祈禱就先去祈禱.
我們一起祈禱.
讓我們學習如何去祈禱.
求主讓我們認清祈禱的功課.
祈禱和你的關係.
以致我們更加捉緊神你自己.
幫助我們由不認識到認識.
幫助我們由認識到更明白.
以致我們更加能夠應用我們的禱告.
我們將我們的禱告交在神你的恩手裡.
祈禱是奉耶穌的名字而求.
阿們.
我們之前一直在講馬達福音的登山部分段落.
我們從馬達福音第五章一直開始.
已經講完第五章已經講到第六章.
如果你記得我們之前的幾次.
大約來說我們講到耶穌論到施捨.
論到禱告和論到禁食.
我們說在當時猶太人來說.

$^{41}$三個很重要的屬靈操練.
以致他們能夠更加執著操練.
更加學懂如何建立和神的關係.
所以這也是我們之前看到的段落.
由六章第一節開始.
我們要謹慎.
我們要小心.
我們要學習如何去敬虔度日.
這裡說不要將我們的敬虔.
著重於表現在人的面前.
因為這樣你得了人的賞賜.
卻得不到天父的賞賜.
所以我們透過六章一至四節.
五至八節和十六節十八節.
說了這三個重要操練的課題.
好像有點快.
要再減慢一點.
所以這三個實踐個人敬虔生活的操練.
當我們說到這三個段落的時候.
你會否發現.
其實我還有一段未說.
中間有一段漏了未說.
其實不是漏了.
是因為設定是這樣.
因為上星期的行道團計.
集中說禁食的問題.
所以我特意跳了上一次.
說了禁食.
其實有一個段落我們未說的.
不知道你能否想到.
應該是六章九至十五節.
這個段落不知道你是否很熟悉.
或者你回教會很久的日子.
你都很熟悉.
所以我現在挑戰你一個問題.
你能否背出主套文.
你不需要背和諧本修訂版.
因為可能有些字眼不同了.
你可以背和諧本.
我試試大家能否背出主套文.

$^{81}$好 開始.
大家背了一段聖經.
你是不是剛才背了這一段聖經呢?.
剛才大家背了這一段聖經.
所以我今天要帶你挑戰的是.
這段不是主套文.
所以你們剛才都背錯了.
我告訴你哪些是主套文呢?.
這些就是主套文.
就是馬達福音第26章39節.
耶穌往前走在赫西瑪利園的時候.
祂就俯伏在地.
祂就是這樣祈禱.
耶穌就是這樣祈禱的.
我父啊 如果可能求你使者杯離開我.
然而不是照我所願的.
而是照你所願的.
這才是耶穌的祈禱.
這在聖經裡.
是耶穌自己和神在赫西瑪利園的時候的祈禱.
我父啊 如果可以的話.
將這個杯離開我.
但是不要照我的意思.
照你的意思.
是耶穌自己親自的祈禱.
這就是主的套文.
然後耶穌有第二次.
再一次遠離門徒去祈禱.
我父啊 這個杯如果不能夠離開我.
必須要我喝的話.
就願你的旨意成全.
這就是耶穌的祈禱.
還有一段耶穌的祈禱.
在十字架上面的時候.
馬達福音第27章.
大約下午三點鐘的時候.
耶穌就大聲高呼.
以利拉瑪撒巴國大利 我的神 我的神 為什麼離棄我.
所以耶穌在十字架上的時候.
那種痛苦的掙扎.

$^{121}$要成就救恩.
卻要承擔世人的罪的時候.
在十字架上 耶穌在呼喚著什麼呢?.
我的神 我的神 為什麼離棄我.
是這樣向神祈禱的.
這些就是耶穌的祈禱.
這些就叫做主禱文.
聽得懂嗎?.
這些就是主禱文.
所以你們剛才背的那段.
我就說不是主禱文.
為什麼呢?.
因為你看馬達福音前面怎麼說呢?.
馬達福音在上到登山寶訓的時候.
一開始的時候.
五章一至二節就說.
耶穌看見這一群人.
就上了山坐下後.
門徒來到他的庚前.
耶穌就開口教導他們說.
整個登山寶訓第五至第七章.
是耶穌教導門徒的東西.
而去到第六章 第九節.
當耶穌講完祈禱的問題之後.
耶穌說.
所以你們要這樣禱告.
耶穌不是說.
所以我就這樣祈禱.
你聽得懂嗎?.
耶穌不是說.
所以我就會這樣祈禱.
不是 耶穌說.
所以你們要跟這個形式.
或者這個公式去祈禱.
所以這是什麼呢?.
耶穌叫門徒的祈禱方法.
不是耶穌的祈禱.
你聽得懂嗎?.
我們聖經裡面有另外一段.
同樣講到這一段.

$^{161}$聖經就在路加福音.
一個對觀福音裡面.
對應的段落.
就在路加福音第十一章.
他說什麼呢?.
耶穌在一個地方禱告.
禱告完之後.
有門徒跟他說.
原來有門徒問主.
求你教導我們禱告.
像約翰教導他的門徒一樣.
於是耶穌對他們說.
你們禱告的時候.
就要這樣去說.
所以根據路加福音.
基於門徒其實很想去學習祈禱的功課.
所以祂就求耶穌教他們.
好像約翰教他的門徒去禱告一樣.
因為司祭約翰會教導他的跟隨者.
怎樣去學習祈禱.
去親近神.
因為司祭約翰都是一些.
有時會去曠野安靜退休.
親近神的一個人.
所以你看到司祭約翰的服飾.
都是一些很曠野的服飾.
現在路加福音告訴我們.
門徒正正想學習一個禱告功課.
就是耶穌身上學習禱告的功課.
所以你會看到.
是門徒求耶穌.
要求耶穌教他們怎樣祈禱.
所以耶穌就教他們怎樣祈禱.
所以我們平時說主禱文.
其實我們是很籠統的.
我們說主禱文是很籠統的.
因為我們說.
如果你熟悉中文的話.
你一說主禱文.
就會理解是主自己的禱文.

$^{201}$但剛才我們看到主的禱文不是這個.
但這個其實我們就偷工減料.
其實這個正式的名字不是主禱文.
是主所教導的禱文.
這樣就正確了.
這個是主所教導的禱文.
所以主所教導的禱文.
就不一定是耶穌要用的禱文.
聽到分別嗎?.
就是耶穌自己祈禱的時候.
和耶穌教門徒要祈禱的時候.
會有不一樣.
待會我們去到後面.
我們就會見到.
很明顯你們背開的主禱文.
這個段落是耶穌不適合用的.
我們會見到的.
待會再解釋.
如果不明白我講漏了.
你們再提出疑問吧.
所以希望今天先打破大家傳統理解主禱文的觀念.
希望大家能夠接受得到.
我這裡說這個不是主禱文.
是因為這個講法不清楚.
應該是主所教導的禱文.
所以我們今天的題目是.
求主教導我們禱告.
我們求主教導我們怎樣去祈禱的功課.
所以今天我們讀的段落是來自.
回應前面第六章第五至八節.
正正好因為門徒.
或者很多人對於祈禱的功課.
其實都不能夠很清楚.
或者有些錯誤的理解.
這句在我們之前看到第五至八節已經說了.
包括了.
耶穌提出有些人.
祂說你們禱告的時候.
不可以仗納假冒為善的人.
耶穌正正好指出有些人的祈禱.

$^{241}$就是假冒為善的人.
我們說假冒為善就是一些戴著假面具.
戴著一些棄子的面具.
好像演戲一樣的面具去祈禱的人.
他們錯了.
耶穌說不對.
耶穌也說到第七第八節說.
你們都不可以作為外邦人.
外邦人的祈禱方式都錯了.
因為他們重覆重覆和重覆.
以為重覆得越多.
神就越聽.
耶穌就說不對了.
所以耶穌也說到一個正確的禱告方式.
就是第六節.
祂說你們禱告的時候.
要進入內室.
關上門.
向那在隱蔽中的父禱告.
父在隱蔽中察看必賞賜你.
所以耶穌說了一個正確的方法.
心態的禱告.
要進入內室的禱告.
所以當耶穌說完幾個禱告正確的方法.
或者不正確的方法之後.
直到第九節開始.
耶穌就說.
所以你們要這樣去祈禱.
就是說耶穌說完一個理論.
現在教一個實踐方式.
甚至教一個叫標準禱文的方式.
因為在猶太人的環境裡.
他們都會常常有一些標準的禱文.
就是有一些方式.
讓猶太人在祈禱的時候.
不會祈禱錯了方向.
而有一些標準禱文.
所以就算今天有些教會.
他們都會有一些標準的禱文.
邀請會眾一起去中讀.

$^{281}$可能用啟認的方式.
就是主席讀會眾讀.
接著主席讀會眾讀的方式.
都會有的.
所以事實上當時很多人都想學習.
怎樣有一些好的方式.
一個format.
一個形式去幫他們.
去能夠明白怎樣去祈禱.
所以我們今天就在這個狀況.
我們明白門徒想知道一種方式.
耶穌也教導他們一個formula.
或者一個形式.
怎樣去做一個好一點的祈禱.
就教導了在六章九至十三節的段落.
當我們今天要去看六章九至十三節的時候.
我們好一點的方法.
我們怎樣能夠讀得好呢?.
我們就要對應路加福音.
因為路加福音裡面都有一個主所教導的圖文.
而我們透過對比這兩個圖文.
我們可以大約了解到.
當時耶穌的教導準確性.
會更加準確.
大家明白我的意思吧?.
接著我又要打破你一直讀這段聖經的說法.
我們看看對比.
對比來說.
我們發現有馬太福音的版本.
就是六章九至十三節.
和路加福音二至四節.
我們去到耶穌所教的祈禱的內容.
我們會發現分開了四個部分.
或者四個元素.
第一個就是「父」.
所以在馬太福音就有「我們在天上的父」.
而在路加福音也有「父」.
可以對比到大家都有的部分.
接著到馬太福音第二個部分.
就是講到「你」.

$^{321}$在路加福音也有這個對應.
所以我們經常會背誦.
「願人都尊你的名為聖」.
「願你的國降臨」.
「願你的旨意行在地上如同行在天上」.
所以路加福音也是在講「你的名」「你的國」.
第三個比較.
第三個比較就會發現.
祈願前面的部分.
耶穌會帶到「我們」.
將祂由「你」轉為「我們」.
「我們」日用的飲食.
今日賜給我們.
免我們的債.
如同我們免了人的債.
不叫我們陷入試探.
救我們脫離凶惡.
對應到路加福音.
有「我們」日用的飲食.
赦免我們的罪.
不叫我們陷入試探.
大家看到這個對應.
這個大家清楚了嗎?.
接著到最後.
馬太福音有.
「因為國度權並榮耀」.
「全是你的」.
「直到永遠」.
「阿們」.
而路加福音沒有.
所以我們先處理.
為什麼會這樣.
為什麼我們在馬太福音裡.
有這個狀況.
「因為國度權並榮耀」.
「全是你的」.
「直到永遠」.
「阿們」.
而路加福音沒有.
我的解釋就是這樣.

$^{361}$你可以打開自己的聖經.
也有這個解釋.
因為古抄本裡.
在馬太福音裡.
沒有「因為國度權並榮耀」.
「全是你的」.
「直到永遠」.
「阿們」.
就是說.
古抄本是什麼叫古抄本呢?.
就是說.
馬太福音開始寫成之後.
一直最古舊的抄本.
一直看一些較舊的抄寫版本.
原來都沒有這一句.
你自己打開你的聖經.
都有這個備註.
當沒有這一句的時候.
就是說.
如果古抄本馬太福音.
最原本最舊的版本.
是找不到這一句.
而路加福音.
它本身也沒有這一句.
那就是說什麼呢?.
很大機會.
本來就是沒有這一句.
馬太福音和路加福音.
都沒有這一句.
它就是沒有這一句.
那我們怎樣去思考這個問題呢?.
我們搞清楚.
耶穌教的部分是哪一部分.
在今天.
在這裡告訴大家.
在耶路撒冷.
有很多不同的景點.
去參觀耶路撒冷.
或者猶太的地方.
很多不同的景點.

$^{401}$去看關於耶穌的生平和事蹟.
其中有一個.
在耶路撒冷的橄欖山上.
有一個叫主禱文教堂.
主禱文教堂.
在這個主禱文教堂裡.
就是說.
當時耶穌在登山的部分.
大約就在這個位置.
所以這個教堂.
就成為出名.
因為耶穌在登山的部分.
就在這個位置進行教導.
而在這個教堂裡.
其中一個特色就是.
他們這個教堂裡.
會在牆上.
有很多不同翻譯版本的主禱文.
主禱文就是.
耶穌所教導的禱文.
可以嗎?.
在整個教堂裡.
不同的牆壁裡.
就有這些不同的.
耶穌教導的禱文版本.
就是《馬達福音》.
六章九至十三節的部分.
其中這一塊.
就是用希臘文所寫的主禱文.
就是我們耶穌所教導的禱文.
在這個教堂裡.
找到這個版本.
原來.
這個希臘文的主禱文版本.
原來都是沒有了.
因為國度權並榮耀.
傳是來的.
直到永遠.
這裡你回到希臘文.
去到最後.

$^{441}$其實都是去到.
救我們脫離的那個惡者.
這就是耶路撒冷的版本.
所以我們要問一個問題.
為什麼古的抄本.
《馬達福音》和.
在《路加福音》裡.
沒有了這一句說話.
其實這一句說話是怎樣的呢?.
有少少原因.
因為剛才說到.
耶穌所教導的祈禱.
你覺得重不重要?.
重要嗎?.
是不是每個人都懂得背?.
事實上.
原來耶穌教導的禱文.
不只是今天我們懂得背.
原來在當時耶穌的年代.
其實很多人都會背誦這個經文.
甚至很多場合都會用主禱文.
你知不知道.
今天什麼場合都會用主禱文呢?.
我們會背這個.
我現在說主禱文.
就是9至13節.
9至13節.
今天什麼情況.
你會無端端背這段聖經.
這段聖經.
平時是沒有背的.
包括了什麼呢?.
有些教會沒有牧師的時候.
就會背這個.
有些教會沒有牧師.
沒有主任牧師.
臨最後祝福的時候.
那個主席覺得.
那個主席不適合這樣去祝福.
他們就會和會眾一起背.

$^{481}$9至13節.
這是其中一種用法.
今天另外一些情況.
用法是什麼呢?.
譬如有時候.
喪禮都會背這個禱文.
有時候婚禮都會背這個禱文.
甚至某些場合.
一些陪寧聚會.
有時候都會以這個禱文.
作為一個完結.
所以這個禱文.
其實很多時候都會用.
原來到了耶穌死了.
升天之後的時間.
原來早期的教會都是這樣.
他們覺得耶穌教這個祈禱文.
非常之好.
所以他們都會希望.
那些人要背.
和常常用這個禱文.
甚至乎在當時.
一些基督徒聚會的場合.
他們都成為了.
他們必定會背的一個禱文.
必定會背的一個禱文.
我們試試看一個文獻.
一個叫做十二使徒為婚.
大約公元70年左右.
或者可能再後一點都不一定.
就是說當時.
耶穌年代之後.
幾十年裡面.
開始有些基督徒.
聚集了一些使徒的教婚.
組成了一個叫做十二使徒為婚.
就是說十二使徒交托著.
要去持守信仰的一些要點.
其中第八個章節.
第八個部分.

$^{521}$就是說有關禁食和祈禱.
當你看看這裡.
有關禁食和祈禱的時候.
看8.1.
8.1就說.
他們其中一個的.
十二使徒為婚裡面的教導是什麼呢.
不要和偽善的人一起禁食.
你記得我們剛剛說禁食和祈禱.
他說不要和偽善的人一起禁食.
因為原來猶太人的禁食.
星期幾就是逢星期一和星期四禁食.
你記得我之前上次說過.
一個禮拜有兩次禁食.
你有沒有印象.
我上一次說到這個問題.
他們就逢星期一和星期四禁食.
因為他們認定.
這些是假冒為善的人的禁食.
所以你就不要學他了.
不要學他的話.
你就選擇星期三和星期五做.
行不行.
就是說你認定你不和猶太人那些假冒為善的人一樣.
你就是跟隨耶穌的基督徒.
你如果要做禁食操練.
就選擇星期三和星期五去禁食.
其實就是這樣的.
當時的一些早期的基督徒的教導.
而其中8.2.
就說不要像那些偽善人的祈禱的時候.
你發現原來他們就會在.
有一個公共場合一起祈禱的地方裡.
就要背這一段.
這段是說什麼呢.
我們在天上的父.
願你的名為聖.
你的國度降臨.
你的旨意成就在地上如同在天上.
給我們每日的飲食.

$^{561}$並且免我們的債.
如同我們免了別人的債.
不要引導我們進入私貪.
救我們脫離邪惡.
因為權柄榮耀都歸你.
直到永遠.
你會發現原來早期的.
他們要在崇拜聚會裡面.
就會加了這一句上去.
加了這一句的理由.
就好像我們今天.
如果耶穌教你祈禱.
去救我們脫離兇惡.
你覺得祈完禱了嗎.
你一定是覺得還沒有祈完禱.
因為你最低限度都要說.
奉主耶穌基督的名祈禱.
原來也是他們的土文裡面.
耶穌教到這裡就結束了.
他們覺得還不夠.
所以要有一個正式的方式去說.
好像奉主耶穌的名祈禱.
要有一句叫完結的句子.
而這句句子.
就是他們用到的剛才的段落.
這個段落可以稱為一個Dosology.
就是一個榮耀的頌讚.
即是說祈完這個祈禱之後.
就要用一個榮耀上帝的方式.
作為結束.
所以由十二使徒門訓.
開始到打後的基督教教會.
他們都用了這種方式.
在他們的崇拜聚會裡面.
去一路加了這種方式去數算.
要將榮耀頌讚歸給三一上帝的方式.
而其實當你看其他的書信.
你都會見到有這種方式.
包括羅馬書.
羅馬書寫完一堆因訊清義的道理之後.

$^{601}$去到第十一章的結束.
開始說實際門徒的應用的時候.
他就說因為萬有都事本於他.
倚靠他 歸於他.
願榮耀歸給他 直到永遠.
他都會用這個.
作為他一至十一章的段落結束的時候.
要獻上一個頌讚給神.
在我們之前讀到的猶大書.
他同樣在一章二十五節說.
我們的救主獨一的上帝.
藉著我們主耶穌基督得享榮耀.
威嚴 能力 權柄.
從萬古以前到現今.
直到永永遠遠 阿們.
都會獻上一個頌讚.
所以根據一些研究.
其實去到大約是公元五世紀.
即是四百多年之後.
原來那些聖經抄本.
關於馬太福音的抄本.
才開始有我們剛才的最後一句.
所以其實是說幾百年之後.
才會有這個抄本.
會有這些經文的段落.
所以我們今天的思考.
耶穌教的討論裡面.
最後這個我們就不說了.
明白我的意思嗎?.
你可以自己上網去研究.
既然我們說了這個部分.
我們大概知道耶穌教的討論.
就是前面的頭三個部分.
我們稍後會思考這三個部分.
在思考的時候.
我們再去了解多一些背景.
讓我們讀主禱文會清晰些.
在剛才說到的主禱文教堂裡.
原來還有兩個翻譯本的聖經.
而這個聖經是特別些的.

$^{641}$在這個主禱文教堂裡.
他們有過百個國家語言的翻譯本.
都是寫在磚牆上.
但有兩個版本是刻在石板上.
特別將這兩個版本的耶穌教的禱文.
刻在石板上.
包括那兩種禱文呢?.
就是希伯來文.
希伯來文的段落.
另外一個就是亞蘭文.
為什麼要幫助我們讀這段聖經呢?.
其中一個原因就是.
猶太人講什麼文呢?.
講希伯來文.
我們一直讀舊約聖經都是希伯來文.
新約聖經用什麼文呢?.
就是希臘文.
那我們就要問.
那亞蘭文是什麼呢?.
亞蘭文是什麼呢?.
或者如果你再看一看畫面裡面的字.
其實希伯來文和亞蘭文都很接近.
分別是什麼呢?.
就是去到耶穌的年代.
你可以想像一下.
希伯來文就是宗教和聖經裡面的.
叫做Official Language.
就是正式的語言.
去到舊約聖經裡面.
正式一用的話.
就用希伯來文.
就是主要的語言.
這個大家容易明白.
希臘文就是當時耶穌的年代的時候.
整個巴勒斯坦地方等等.
因為他們有很多商業往來.
以至到能夠和世界各地去溝通.
去聊天.
其中一個重要用的語言就是希臘文.
就是希臘文.

$^{681}$當然有時也有拉丁文的存在.
就是希臘文.
這裡說到原來日常生活的猶太人.
講什麼文呢?.
原來就是講亞蘭文.
如果在生活裡面聊天.
不是一個正式的記錄.
就講亞蘭文.
就好像我們聊天.
就直接用廣東話.
如果正式寫文字.
可能用文言文.
又或者用書信語.
我們寫法是不一樣的.
但我們講口語就會流暢些.
表達到我們的情感多些.
所以這個都會幫助我們明白.
耶穌去教門徒.
教門徒祈禱的時候.
他當時極大機會用亞蘭文.
是他講的語言.
作為他教的語言.
對話裡面講話.
就用亞蘭文.
但當在聖經裡面正式的文字描述.
耶穌的生平.
耶穌的歷史.
耶穌的出生.
就會用希臘文.
主要一開始的時候就用希臘文.
即是說.
當我們今天見到的新約聖經.
最初的時候.
其實當時就需要將耶穌的亞蘭文講的東西.
特別是祈禱文講的東西.
就要轉成希臘文.
才去寫出來.
你會明白有這樣的背景在裡面.
幫助我們處理一些問題.
好,大約都可以.

$^{721}$因為今天可能比較學術性.
但其實就幫我們更加去明白.
那段聖經.
好,所以我們接著很快.
講完三點我們就完結.
不是講到三點,是講三點.
第一點就是.
既然是這樣.
我們就要明白.
耶穌第一個部分是講什麼.
講什麼.
剛才我講.
講一個字.
父.
就是父.
所以我們就看父.
第一個我們思考的地方就是.
耶穌在一開始.
耶穌教我們的祈禱.
第一個要教我們的是什麼呢.
就是教我們父.
所以在這裡.
馬達福音的描述就是.
我們在天上的父.
路加福音說.
父.
是一講父親的角色和形象.
所以你看回剛才我們去到.
之前看到第六章四節.
我們做詩寫的時候.
最重要給誰看.
給父看.
給天上的爸爸看.
禱告的時候給誰看呢.
給父親.
給天上的爸爸看.
我們禁食的時候.
是給別人看我們怎樣面容扭曲嗎.
不是.
我們要給天上的爸爸看.

$^{761}$所以在這裡.
耶穌原來一直在講.
一個很重要的重點.
整個的操練裡.
你要清楚你的對象是誰.
原來今天我們.
其中禱告學習的功課.
我們要搞清楚.
我們祈禱的對象是誰.
是誰才是我們祈禱的對象呢.
祂說是我們的父.
那個父是怎樣呢.
我們試試看.
在《馬可福音》裡就這樣去譯道.
祂說.
《馬可福音》.
耶穌就這樣教過祈禱.
祂說.
阿爸 父啊.
在你凡事都能.
求你將這杯撤去.
然而不是照我所願的.
而是照你所願的.
原來在《馬可福音》裡面.
就已經譯成.
為何譯成阿爸 父呢.
這個就不是我們中文的阿爸.
因為我們廣東話就說.
阿爸 阿媽.
但是這個其實就是.
在亞蘭文裡面.
阿爸就是說爸爸的意思.
所以你看到耶穌的叫天父.
祂自己叫.
祂可以叫得很直接.
就是阿爸 阿爸天父.
耶穌是這樣祈禱的.
耶穌是直接叫天上的父親.
叫做阿爸 父.
這個就是祂直接的父親.

$^{801}$這個描述有什麼重要呢.
當我們理解.
耶穌教導的圖文.
如果都是這樣.
叫我們天上的父.
叫我們的爸爸.
叫我們的阿爸天父的時候.
有什麼重要呢.
因為你何時會叫阿爸天父呢.
或者你何時會叫你爸爸阿爸呢.
你今天會怎樣叫你爸爸.
你會叫他老豆 阿爹.
其實這個阿媽這個字.
就是一些小朋友.
啞啞學語的時候.
第一最快學到的就是這些發音.
所以你知道有些IQ題說什麼.
什麼字是爸爸媽媽.
會說到哥哥姐姐.
這些做不到的.
有沒有玩過這個遊戲.
哥哥姐姐嘴巴不會貼著.
爸爸媽媽嘴巴會貼著.
BB最小的時候.
怎樣會叫你呢.
一定是叫爸爸和媽媽.
而爸爸媽媽往往就是爭.
究竟是叫爸爸還是媽媽.
所以這個字的發音.
是一個幼嫩的小朋友.
第一次叫到他的父親或者母親的時候.
就是這種方式.
這種這麼親切.
當時他不懂什麼叫制度.
不懂讀書.
不懂什麼叫考試.
不懂什麼叫玩.
不懂什麼叫賺錢.
的時候.
一個小朋友就直接叫爸爸.

$^{841}$你要知道你和天父的關係都是這樣的.
其實不是在於你現在懂多少.
你有多少事奉.
你是什麼教會有多高權力.
或者你在外面社會有多少金錢能力.
原來耶穌告訴我們一個最基本的.
就是我們來到神的面前.
重回我們好像一個小朋友一樣.
就是叫爸爸.
這個才重要.
所以有些人經常說.
我不懂得祈禱.
不懂得祈禱.
鄧小姐你搞錯了.
你不懂得祈禱.
可能是你說話不流利.
如果你是一個BB.
你只是對著你爸爸.
那個爸爸已經開心到飛起了.
你做過父母你明白.
天父都是希望我們有這樣的狀況.
而這個問題正正好和當時猶太人.
他們那些宮土的習慣是不一樣的.
為什麼不一樣呢.
因為猶太人在會堂聚集的時候.
他們的祈禱已經變得很正式.
很公式化了.
所以到了他們的宮土猶太人的會堂聚會.
他們會喊主神傳遞的王.
他們只會用一個很敬畏的態度.
不敢亂說話的態度.
確實有時候我們對我們的神.
可能是要恭敬.
但是耶穌更加說.
如果這是一個私密的禱告.
因為你知道.
第六章說的是個人屬靈生命的操練.
不是在公開的場合.
而是你和神的親密關係.
就不應該用一個很公式化.

$^{881}$很正式的方法.
去說你這位天上的父親.
所以第一個我們要知道的.
我們要思考的是.
我們有沒有想過我們和神的關係.
就好像子女對爸爸一樣.
單單告訴他爸爸.
不需要計較你認識多少人.
就是去到我們的父親面前.
他就是這樣愛我們的那位爸爸.
當我們想想.
我們開始要搞定我們的焦點的話.
就是這位的關係.
這個關係就是父的和子的關係.
第二個我們要思考的就是.
你了.
第二個很重要的元素就是.
用了三個的祈願.
三個的願.
去表達了耶穌教我們祈禱的方向.
這個祈禱的願望是什麼呢?.
「願人尊你的名為聖」.
「願你的國降臨」.
「願你的子兒行在地上如同行在天上」.
這裡說到第一個「你」是什麼呢?.
「願你的名為聖」是屬於聖的.
是分別出來的.
因為神的名字是重要的.
我們作為神的子女.
我們要知道.
我們期待著神的名字.
要被人看為聖的.
是出埃及的說法.
「不可妄稱耶和華你上帝的名」.
「因為妄稱耶和華的名」.
「耶和華必不以他為無罪」.
任何人輕看了上帝的名字.
神會定他為有罪.
而且神的名字是滿有能力的.
在詩篇裡說.

$^{921}$「有人靠車有人靠馬」.
「但我們要看到耶和華」.
「我們上帝的名」.
還有你記不記得.
有首兒歌是什麼呢?.
「靠著耶穌聖名」.
懂得唱嗎?.
「靠著耶穌聖名」.
「我們便能得勝」.
「靠著耶穌聖名」.
「我們便能得勝」.
「誰能掃盡新的作為」.
「誰能掃盡他的愛」.
「靠著耶穌寶貝聖名」.
「我們便能得勝」.
所以神的名字是為我們帶來勝利.
因為神的名字是包括了他的權能.
他的權柄.
所以我們去到神的面前.
我們祈禱.
我們要想想.
我們先將我們的焦點.
不是想想自己有多好的生活.
而是想想我們能否將我們的祈願.
放在神身上作為優先.
所以我們用了三個「你」.
去將我們的焦點.
我們的眼光看向天上.
所以神的名字是我們要留意的.
還有什麼呢?.
神的國呢?.
馬達福音耶穌說.
「天國近了,你們應該悔改」.
我們也要像我們的天父一樣.
將我們的眼目.
想想神的國如何臨近地上呢?.
但是你看回門徒.
當在《使徒行傳》第一章的時候.
耶穌已經要升天了.
但門徒問的問題是什麼呢?.

$^{961}$他說:「主啊!你就要在這時候復興以色列國嗎?」.
門徒正在問.
關心地上的以色列國.
是否能夠復國.
但是耶穌所關心的是.
天國要臨近.
所以耶穌也在告訴我們.
要挑戰我們的眼光.
仍然在一個焦點.
就是「你」.
就是天父「你」.
你的名字是怎樣呢?.
你的國是怎樣呢?.
將我們的眼光看向這個景地.
第三個.
就是你的旨意.
你的旨意「行在地上」.
「行在地上」是什麼意思呢?.
就是「被實現」.
我們願意神的旨意能夠在地上實現.
而不是我們的心思在地上能夠被實現.
因為你要知道.
當神的名被彰顯.
神的國臨近.
神的旨意能夠在地上實現.
這是屬於上帝的土地.
上帝必然看顧.
這是屬於祂的土地.
所以今天我們第二個思考的問題.
就是.
在我們的生活中.
我們有沒有具體的行動.
顯明了我們確實.
願人都尊你的名為性.
願天父的國降臨.
願天父的旨意實現呢?.
這是很重要的.
我們要知道.
我們如何能夠將我們的眼光.
在我們祈禱的時候.

$^{1001}$先去尋求神.
當我們願意先去尋求神.
我們知道其他的事情.
上帝會看顧.
上帝會顧惜.
好.
今天很快了.
第三個部分.
這個主禱文教我們的部分.
就是「我們」.
當我們能夠搞清楚.
我們的主題.
我們的方向.
是神的時候.
當我們這樣祈禱.
轉來的是什麼呢?.
就是天父會回應我們的需要.
天父會知道我們身心靈的需要.
因為我們是祂的子女.
當我們能夠認定.
我們是神的子女的話.
我們就知道.
接著這三個問題.
可以因為神.
能夠得到解決.
這裡舉了三個情況.
可以讓我們去想一想.
第一個就是.
我們常常都會面對的困難.
就是我們的飲食.
我們常常擔心.
有沒有得吃.
吃得夠不夠.
這是我們所擔心的.
這裡經文裡面.
飲食這個字.
其實就是在說「糧」.
那個「糧食」.
那個字的意思.
這裡說.

$^{1041}$我們日用的飲食.
日用這個字.
其實是什麼意思呢?.
這裡說.
有兩個可能的翻譯.
第一個就是「明日」.
另一個翻譯就是「今日」.
就是「今日」.
無論你是「明日」.
還是「今日」也好.
重點就是「今日次我」.
你想像一下那個畫面.
如果我們理解這裡.
就是一個「明日之糧」.
就是日用的飲食的意思.
「明日之糧」.
就是明天我需要的食物.
「今日次我」.
代表什麼呢?.
代表當你還在擔心.
你自己有沒有生活的供應.
擔心你有沒有得吃.
擔心你的子女有沒有得吃的時候.
我們可能處於一個焦慮的狀況.
很多擔憂的時候.
「明日之糧」.
今天就已經賜給你了.
即是說.
神已經供應了.
而這一段的段落.
其實都正正好.
回應在摩西的時候的嗎哪.
這裡是同一個畫面.
嗎哪的特色是什麼呢?.
每天降份量.
降了之後怎樣呢?.
拿多一點呢?.
又不行,會變壞的.
你說拿少一點呢?.
它又夠你吃的.

$^{1081}$是不是?.
到安息日之前那天.
又會降多一點的.
平時放一天就變壞了.
但之後有安息日.
安息日就什麼工都不可以做的時候.
那天降的那些.
就兩天都不壞的.
那個時間剛剛好.
剛剛好是什麼呢?.
你每一天.
如果你要想為那天擔憂的時候.
神就在那裡已經供應了給你.
你不用想那麼多了.
那天就給了你.
他要為我們釋除的是什麼呢?.
我們有很多的憂慮.
之後我們會在第六章.
再講一些憂慮的問題.
原來是這樣的.
神對我們的看顧.
有沒有認定了.
祂是你天上的父親.
祂的旨意成就.
難道沒有好處給你嗎?.
難道不看顧你嗎?.
不會的.
所以在今天我們地上的需要.
我們的糧食裡.
神都會看顧我們.
就好像昔日神看顧摩西.
他們一群以色列人.
記住.
神的看顧不是一天.
吃了多少天.
嗎哪.
真的吃到你都悶.
但是又夠你吃.
又吃得你充充足足.
四十年的時間.

$^{1121}$當然準確一點不是四十年的.
不過整體大約是四十年的時間.
四十年裡面上帝的供應.
那你信了耶穌多少年了?.
有沒有人多過四十年信耶穌?.
有啊.
那你四十年有沒有供應?.
有啊.
不捨得還在這裡.
四十年都有啊.
那你們這些三十幾年.
二十幾年.
十幾年那些.
不用擔心.
有個見證告訴你.
四十幾年上帝都有供應.
所以我們今天的飲食.
你知道我們怎樣去祈禱.
天父啊.
你會供應.
告訴神.
我知道你會供應.
第二個是什麼呢?.
免我們的債.
如同我們免了人的債.
這裡有個有趣的地方.
這裡的翻譯是什麼呢?.
免我們的債.
這個債字其實是什麼呢?.
因為原來用希臘文.
希臘文這個債字.
就是講和錢有關的.
如果真的是希臘文.
這個債就是錢債.
所以如果是這樣理解.
我們就知道.
可能是在講耶穌.
今天會免我們的錢債.
如同我們免了人的錢債一樣.
用希臘文這樣去理解.

$^{1161}$不過正如剛才所說.
耶穌可能當時說的時候.
是用亞蘭文的.
原來亞蘭文的債這個字.
又不只是講錢債.
如果你明白耶穌當時.
如果用亞蘭文說的話.
這個債字原來可以解作錢債.
也可以解作罪債.
即是說耶穌當時可能在說.
當我們認定了這位天父的時候.
神會免去我們的錢債.
就如同我們免了人的錢債.
但也可能在說.
神會免去我們的罪債.
就如同我們免去別人.
得罪你的罪債一樣.
原來聖經教我們什麼呢?.
我們要免去別人的債.
因為我們都得到神的赦免.
事實上在舊約裡面.
已經在教這件事情.
你知道在安息日裡面.
他們就會得到安息.
有些的債要赦免.
當有些人欠你的債.
經歷安息日的話.
即是安息七年的時候.
可能欠你的債.
你要一筆勾銷.
你要去嘗試免人的債.
你會經歷上帝免你的債.
甚至乎到欺凌的時候.
本來你的土地賣了給別人.
欺凌五十年的時候.
你的土地可以還給你.
免去欠你的土地.
把土地還給你.
所以你知道.
原來當我們搞清楚.

$^{1201}$我們和神的關係的時候.
想想我們.
今天你欠別人多還是別人欠你多呢?.
那你搞錯了.
我們欠神最多.
別人欠你的.
永遠都不夠我們欠神的.
對不對?.
我們欠神的一定是最多.
如果神願意赦免我們的債.
為什麼我們不赦免別人的債呢?.
你那些區區細細的數目.
真是零頭.
零頭即是跌了都不記得.
但是你欠上帝的是何等大.
因為你欠上帝的.
可能由你出生到現在.
幾十年欠上帝.
而且是不斷加積.
你欠上帝的罪債.
但是上帝都願意寬恕你的時候.
為什麼你不寬恕別人呢?.
所以耶穌帶給我們的角度.
就是要搞清楚.
我們都要學習神給我們的狀況.
就是神願意寬恕我們的罪債.
就好像我們寬恕別人的罪債一樣.
當然這不是代表我們.
赦免別人的罪債.
是神一定要怎樣寬恕我們.
因為我們的恩典是白白的.
是何等的恩典.
但是這正是提醒我們.
在我們地上的生活裡.
就是要回應上帝這種.
免債的教導.
最後就是.
我們的試探.
不要叫我們陷入一個試探裡.
要脫離惡者.

$^{1241}$脫離凶惡.
這裡說.
究竟我們會不會讀了.
就是神試探我們呢?.
我們知道在雅各書裡.
神是不會試探人的.
只是我們會自己.
讓自己陷入試探裡.
所以我們知道這一句.
這句正正提醒我們.
原文亞蘭文翻譯的意思.
去到亞蘭文原文的翻譯.
就是神說.
叫我們去到一種情況.
不需要面對試探.
就是這句意思.
不叫我們遇見試探.
不是說神做一個試探給你.
叫你神這次放手.
不要做一個試探給你.
不是這樣理解.
而是說原來.
求神讓我們.
遇上一種情況.
就是我們不需要面對試探.
那個試探不是來自上帝.
那個試探可能來自我們自己.
但是我們可以求神.
讓這種情況不要讓我們遇見.
不要讓這種情況出現.
因為神是沒有能力的.
所以我們要回到上帝面前.
我們要求神幫助我們.
脫離這些凶惡的試探.
即是說神會幫助我們.
你相不相信神會幫助我們呢?.
如果你相信.
你就可以這樣去祈禱.
所以這裡整個我們的部分.
正正好神正在說.

$^{1281}$耶穌正在說什麼呢?.
我們生活肉體的需要.
你相不相信神會幫助你?.
我們生命裡罪的問題.
你相不相信上帝會赦免你?.
我們生命裡每天.
面對很多凶險的狀況.
一些我們所謂的試探.
或者一些凶惡的事情.
我們相不相信上帝會保守你?.
所以耶穌教導這個圖文.
去到後段的話.
教導我們什麼呢?.
要挑戰我們.
真的要相信那位上帝.
我們真心去信這位神.
我們能夠跨越這種狀況.
所以這是一種信心的表現.
耶穌教導我們的.
不是要一個公式.
耶穌在圖文裡教導我們的.
就是要有我們的焦點.
首先要搞清楚我們和神的關係.
是父子和父女的關係.
第二個.
我們的眼光.
應該先向神.
先望向天.
先以神為優先.
這就是優先次序.
第三.
教導我們去學習.
真的放上我們的信心.
相信上帝在生活實踐.
在生活,靈性,聚齋,困難裡.
上帝都會幫助我們.
耶穌很想我們堅心去依靠祂.
這就是耶穌給我們的.
祈禱教導.
祈禱教導到這裡差不多了.

$^{1321}$這是A部分.
今天沒有B部分.
放心,沒有B部分.
不過你發現還有一段沒說嗎?.
這段是我的總結.
當耶穌教導完這個圖文之後.
你會發現他在第十四,十五節.
說了兩句話.
「你們若不要恕人的過犯.
你們的天父也必饒恕你們的過犯」.
「你們若不要恕人的過犯.
你們的天父也必不要恕你們的過犯」.
這是什麼意思呢?.
即是說我們已經祈禱了.
我們已經祈禱了.
我們覺得我們的祈禱是公式.
我們覺得祈禱只是循例去祈禱.
我們覺得祈禱只是將我們的要求.
寫一張清單交給神.
我們也搞錯了.
耶穌告訴我們我們的祈禱是什麼呢?.
就是我們祈禱的時候.
是進入和神的一種關係和狀況裡.
而這種關係和狀況裡.
你是要被轉化的.
聽得明白嗎?.
一些肉放進攪碎機裡.
那些肉如果出來.
仍然是那塊肉.
那攪碎機行不行?.
是不行的.
你把一些東西放進微波爐裡.
開著微波爐.
拿出來.
仍然是冷的.
行不行?.
不行的.
我們要轉化.
微波爐就是將能量轉化.
攪拌機就是將物質轉化.

$^{1361}$而我們去到神面前.
祈禱就是讓我們和神的屬靈關係.
能夠轉化.
所以為什麼最後有十四十五節呢?.
就是一個轉化的挑戰.
不是祈願就算了.
你有沒有做呢?.
你有沒有讓你生命.
被這種禱告去轉化.
你能夠去實踐在你的生命裡呢?.
因為在中間我們說到.
我們免人的債.
如同我們免了別人的債.
我們得到神赦免我們的罪債.
就好像我們應該免去別人欠你的罪債.
有人得罪了你.
如果你祈禱.
你是應該期望你會轉變的.
你會期望你的心態會轉變.
你的態度會轉變.
以致能夠像十四十五節那樣說.
你要麼就選擇饒恕別人的過犯.
因為你都很期待天父饒恕你.
要麼你就繼續堅持不饒恕別人.
但你要記住.
你欠天父更多.
如果天父都不饒恕你.
你就很糟糕了.
所以最後十四十五節對我們生命帶來一個挑戰.
就是我們願不願意.
我們的祈禱不是一個公式的祈禱.
而是一個轉化生命的祈禱.
當你認真向神祈禱.
你會發現你的心會轉變.
你的心不會轉變.
因為你對神的祈禱是搞錯了方向.
聽得明白嗎?.
當你帶著一個硬心.
帶著一個上帝滿足你的心.
上帝不滿足你.

$^{1401}$他就不聽你的祈禱.
你搞錯了.
而是我們帶著一個謙卑在神面前禱告的心.
求神去轉化我們.
就像我們的禱告一樣.
成為輕香.
所以當我們求神免去你的債的時候.
你也要免去別人的債.
最後一個總結就是一個小故事.
一個真實的故事.
就是我在香港做一些青少年工作.
其中有一次有一個女孩子.
是一個中學生.
她來找我和另一個傳道同工.
她不斷投訴她的同學.
她投訴那個同學怎樣對她不好.
怎樣壞.
甚至乎一邊說一邊投訴的時候.
她不斷說髒話罵那個人.
我們也鼓勵她.
因為她說她信耶穌來找我們.
做基督徒是這樣的.
但是我們也知道和她談了一段時間.
希望能夠幫她梳理一下.
但其實她心裡也很冤屈.
我們也覺得要和她談一談.
接著我們只做一件事.
我們邀請那個年輕人.
那個女孩子去祈禱.
我說這個人真的得罪了你.
不如你為她祈禱吧.
你為她祈禱我們才結束祈禱.
我們邀請那個小女孩.
為這個得罪她的同學祈禱.
她一直為這個祈禱的時候.
她本來在口裡也有很多髒話.
不斷地罵.
但當她真的下頭去祈禱的時候.
她知道她和天父祈禱說話.
她不敢再在神面前說髒話.

$^{1441}$她轉過來真的為這個得罪她的同學去祈禱.
一直聽她祈禱的時候.
我聽到神改變她的心.
她真的關心這個同學的狀況.
求神幫助她.
當這個女孩子願意為這個同學.
求神幫助這個同學的時候.
其實她已經經歷了神的幫助.
她不是在咒詛,埋怨裡.
所以那次我很感恩.
看到原來她的祈禱真的遇上神.
那種遇上神不是問題解決了.
而是她的心境在禱告裡能夠轉化.
第一次每次求主幫助我們.
特別有很多人得罪你的時候.
你為這些人祈禱.
你要經歷那種轉化.
因為我們祈禱就是和神建立關係.
如果你覺得她憐憫不夠.
你就為她求憐憫.
不斷數算自己的抱怨.
因為你欠神的往往比更多.
如果神都原諒你.
你應該要被神原諒.
你也要學習原諒別人.
求主幫助我們怎樣去祈禱.
你要盼望今天讓我們更加明白.
主禱文的意思其實就是主教導我們祈禱的文章.
還有補充一句.
我忘記補充.
為什麼耶穌不可以用這個禱文.
補充完這個就完了.
為什麼耶穌不可以用這個禱文.
為什麼.
聽完了.
因為耶穌不可能免他的罪.
因為耶穌是無罪的.
免我們的債.
免我們的罪債.
如同我們免別人的罪債.

$^{1481}$耶穌怎需要免他的罪債呢.
所以耶穌這個禱文一定不適合他.
搞清楚了.
這個不是耶穌需要的禱文.
而是我們基督徒需要學習的禱文.
求主幫助我們.
我們一起去祈禱.
讓我們重新調整和你的關係.
特別讓我們更加認清我們的焦點.
排好我們的優先次序.
更加去信賴上帝的.
供應,憐憫,赦罪.
使我們更加有力.
在被你轉化的過程裡.
我們去饒恕別人.
我們去幫助別人.
我們去為別人生命裡.
去扶持他們渡過他們的凶險.
求主你就是這樣帶領我們.
聽我們祈禱奉耶穌的名字而求.
阿們.
\newpage



\section{約翰一書 1:1-2:2}
\label{sec:031611}
\textbf{如何活出合神心意的團契生活}~~~~鄧榮仲傳道
\newline
\newline
連結: \href{http://www.church.com.hk/acms/content.asp?site=cdc\&op=show\&type=product\&code=031611\&layout=sermon}{\texttt{031611}} ~~~~ 語音日期: 2026-02-08
\newline
\newline
\hyperref[sec:031610]{\small{< < < PREV SERMON < < <}}
~
\hyperref[sec:index_chronic]{\small{[返順時目]}}
~
\hyperref[sec:031612]{\small{> > > NEXT SERMON > > >}}
\newline
\newline
約翰一書 1:1-2:2
\newline
\begin{longtable}{cl}
\hline
\hline
章節 & 經文 (和合本修訂版)\\
\hline
1:1 & \begin{tabularx}{0.7\textwidth}{X} 論到從起初原有的生命之道，就是我們所聽見、所看見、親眼看過、親手摸過的－ \end{tabularx} \\ \\ \relax
1:2 & \begin{tabularx}{0.7\textwidth}{X} 這生命已經顯現出來，我們看見了，現在又作見證，把原與父同在，並且向我們顯現過的那永遠的生命傳揚給你們－ \end{tabularx} \\ \\ \relax
1:3 & \begin{tabularx}{0.7\textwidth}{X} 我們把所看見、所聽見的傳揚給你們，為要使你們也與我們有團契，而我們的團契是與父和他兒子耶穌基督所共有的。 \end{tabularx} \\ \\ \relax
1:4 & \begin{tabularx}{0.7\textwidth}{X} 我們把這些事寫給你們，使我們的喜樂得以滿足。 \end{tabularx} \\ \\ \relax
1:5 & \begin{tabularx}{0.7\textwidth}{X} 神就是光，在他毫無黑暗；這是我們從主所聽見，又報給你們的信息。 \end{tabularx} \\ \\ \relax
1:6 & \begin{tabularx}{0.7\textwidth}{X} 我們若說，我們與神有團契，卻仍在黑暗裡行走，就是說謊話，不實行真理了。 \end{tabularx} \\ \\ \relax
1:7 & \begin{tabularx}{0.7\textwidth}{X} 我們若在光明中行走，如同神在光明中，就彼此有團契，他兒子耶穌的血就洗淨我們一切的罪。 \end{tabularx} \\ \\ \relax
1:8 & \begin{tabularx}{0.7\textwidth}{X} 我們若說自己沒有罪，就是欺騙自己，真理就不在我們裡面了。 \end{tabularx} \\ \\ \relax
1:9 & \begin{tabularx}{0.7\textwidth}{X} 我們若認自己的罪，神是信實的，是公義的，必要赦免我們的罪，洗淨我們一切的不義。 \end{tabularx} \\ \\ \relax
1:10 & \begin{tabularx}{0.7\textwidth}{X} 我們若說自己沒有犯過罪，就是把神當作說謊的，他的道就不在我們裡面了。 \end{tabularx} \\ \\ \relax
2:1 & \begin{tabularx}{0.7\textwidth}{X} 我的孩子們哪，我把這些話寫給你們，是要你們不犯罪。若有人犯罪，在父那裡我們有一位中保，就是那義者耶穌基督。 \end{tabularx} \\ \\ \relax
2:2 & \begin{tabularx}{0.7\textwidth}{X} 他為我們的罪作了贖罪祭，不單是為我們的罪，也是為普天下人的罪。 \end{tabularx} \\ \\
[1ex]
\hline
\hline
\end{longtable}
$^{1}$各位大家好 早安.
多謝林全圖的邀請來到貴堂和大家分享訊息.
讓我們一起去低頭 先有個祈禱.
因為是你的恩典 讓我們能聚集一齊去敬拜你.
當我們來到你的根前的時候 聆聽你的說話.
求你的聖靈開我們的眼睛 開我們的耳朵.
以致我們能夠明白你的話語.
當我們去跟從你 做你的門徒.
神父我們求你 首先教導我們和你建立一個好的關係.
以致我們能夠和弟兄姊妹.
能夠活出一個合神心意的團體生活.
神父我們求你 教導我們更多地認識你.
因為當我們認識你的時候.
我們才能夠和弟兄姊妹有那份團契.
主啊 我們知道我們很多時候都有很多的軟弱.
我們未必有愛心.
所以上帝我們求你 教導我們回到你的根前.
去懂得認罪悔改 求你的饒恕.
主啊 你是滿有恩典.
因為你應許過 當我們回到你的根前的時候.
你會赦免我們的罪.
你會讓聖靈去轉化我們的生命.
我們的禱告是奉乃基督的命求 阿們.
我們大部份都是說廣東話的.
所以我們是澳洲的移民.
當我們來到澳洲的時候.
我們會選擇怎樣的去找教會呢?.
我們會不會去到教會那裡參加的時候.
我們會看一看教會的團體生活是怎樣的.
如果那些人不理會你.
你回到教會 坐到崇拜.
接著來到第一次 沒人和你溝通.
沒人理會你的時候.
你可能會想 這間教會也未必很歡迎我.
我還是再找另一間吧.
會不會這樣呢?.
當我們去想團契是一個怎樣的處境的時候.
我們會不會覺得應該按照我們的年齡層.
我們有共同的興趣.
以致我們劃分我們的團契.

$^{41}$以致我們能夠志趣相投的時候.
我們能夠更加做到我們的團契生活.
當我們回到教會久一點的時候.
我們會說一個這樣的說法.
就是我們要彼此相交.
其實相交這個字是從聖經而來的.
是聖經的用語.
剛才我們讀這段經文 有讀過.
簡單地說就是相交等同團契生活.
我之後會再解釋.
但是我們想想團契生活應該是怎樣的一回事呢?.
有一個今天的說法.
團契或小組是提供一個靈明成長的環境.
讓我們能夠彼此關顧,服侍,激發愛心.
經歷神的真實.
我想整體上大家對這個定義都會同意.
但是團契是否就是這麼簡單.
說我們跟弟兄姊妹之間的關係呢?.
我們今天看這段經文的時候.
我們可以再看多一點.
當然我們說我們作為一個基督徒的時候.
大家可能初信的時候.
接受栽培可能都見過這幅差不多的圖畫.
就是說我們的生命的中心是基督.
當我們有基督在我們生命的中心的時候.
我們以一個直向的關係.
就是藉著祈禱和讀神的話語.
我們能夠認識上帝.
但是我們更加有一個橫向的關係.
就是我們對這個世界未信的朋友作見證.
見證耶穌基督.
另一方面 我們另一個橫向關係.
就是我們教會弟兄姊妹肢體之間的相交.
大家都很同意.
這是我們整個基督徒生命的重要性.
我們回到今天我們讀這段經文的內容的時候.
《約翰一書》.
大家對《約翰一書》有沒有什麼印象?.
大家有沒有讀過《約翰一書》最少一次的?.
最少一次 OK.

$^{81}$其他沒讀過 是吧? OK.
不要緊 我們就嘗試看一看.
今天這段經文《約翰一書》的第一章.
其實這本書不長 只有五章.
鼓勵大家細讀這本書.
而我們去看這段經文的起點的時候.
大家可能有個發現.
這一章的三到七節當中.
其實連續有四次「相交」這個字出現.
很特別的.
而我們看《新約聖經》的內容的時候.
「相交」這個字在另一本書.
在哥林多後書出現得很多.
不過都是比較分散一點.
但是如果集中地「相交」這個字.
在三到七節當中出現了四次.
就唯一在《約翰一書》出現.
所以我們可以很肯定地說.
《約翰一書》其中一個很重要的主題是「相交」.
那什麼叫「相交」呢?.
我們剛才很簡單地說「相交」就是團契.
而我們再看清楚這個字的意思的時候.
在我們中文的不同翻譯本裡面.
有不同的說法.
《栗鎮忠》說「團契」.
《新約本》說「彼此相通」.
《新漢語》說「契合相交」.
《環球聖經》說「親密的契交」.
不同的形容詞都好像說一個很親密的關係.
但是基本上在大部分的英語翻譯.
Kane,James,NIV,ESV,NASB,NetBible.
其實都是翻譯同一個字.
Fellowship 即是「團契」.
我們剛才說了團契的定義.
小組團的定義.
似乎是說我們和弟兄姊妹之間的關係.
但是我們再看清楚.
「團契」在新約不同的經文裡面.
其實是一個怎樣的理解.
我們看看哥林多前書一章.

$^{121}$這段經文說「神是信實的你們.
原是他所召,好與他而止.
我們主耶穌基督一同得分」.
「得分」這個字就是「相交」,「團契」.
其實這是一個分享的意念.
第二,如果我們從《仕途行傳》二章四十二節說.
「都恆心遵守仕途的教訓彼此交接」.
「交接」這個字也是「相交」.
他們去抹餅祈禱.
這裡似乎和我們一般的意念更加清晰.
就是說弟兄姊妹之間彼此聯絡的處境.
第三個「相交」這個意念有點特別.
在哥林多後書的九章十三節.
這裡說「他們從者供給事上得了憑據.
知道你們承認基督.
順服他的福音,多多地捐錢」.
其實這個字是「相交」.
「給他們和眾人使將榮耀歸於神」.
所以這裡所說的「相交」或「團契」.
是一個捐獻的意思.
所以中文的翻譯是這樣做的.
其實這個「相交」或「團契」.
在新約的理念裡有不同的層次.
回到今天我們說這段經文的時候.
我想有三個主題想和大家分享.
我們如何活出一個合神心意的團契生活.
第一點:我們首先要和基督的生命連結.
第二:我們需要認識神.
第三:我們要向神認罪.
第一點:一章一到四節的大主題.
是要連結基督的生命.
因為一章三節這樣說.
「我們將所看見所聽見的傳給你們.
使你們與我們相交.
乃是與父並他兒子耶穌基督相交的」.
如果我們仔細看這一節的經文.
其實是在說這件事.
我們弟兄姊妹之間.
能夠有一個好的團契生活.
背後有一個原因.

$^{161}$我們教會不是一個社區團體.
你來參加一些社康活動.
大家來社交一下,玩一下.
不是這麼簡單.
我們背後有一個很重要的原因.
是因為首先我們自己.
是與父和耶穌基督.
首先有一個團契的生活.
因為這樣的緣故.
這是背後的一個原因.
我們能夠與弟兄姊妹之間有相交的生活.
我們如何能夠與父和耶穌基督有相交的生活?.
用剛才我介紹的三個理念分享的理念.
首先有一個分享的理念.
這段經文一章二節說.
「這生命已經顯現出來.
我們也看見過.
現在又作見證.
願與父同在.
並顯現與我們那永遠的生命全給你們」.
這段經文告訴我們.
基督首先願意分享祂的生命給我們.
所以我們能夠與父和耶穌基督.
有一個團契.
另一方面.
我們如何與基督有彼此聯絡的關係?.
回到剛才的十字圖畫.
基督在中心的時候.
我們藉著禱告.
將我們的心思,意念,心底的處境.
向神,向基督說出來.
但另一方面.
我們藉著讀上帝的話語.
我們能夠明白神的心意.
我們能夠明白基督的心意.
藉著這個靈修的過程.
我們與主耶穌基督有彼此聯絡的關係.
至於捐獻的部分是怎樣呢?.
我們看看在約翰福音裡.
我們很熟悉的那段經文.

$^{201}$約翰福音三章十六節.
是我們信仰一個很核心的內容.
這裡說:.
「神愛世人,甚至將祂的獨生子賜給他們.
叫一切相信的不致滅亡,反得永生」.
父上帝愛我們.
將祂的獨生子捐獻出來送給我們.
以致當我們願意相信的時候.
我們就能夠承受祂賜給我們的永生.
所以這裡上帝將祂的獨生子賜給我們.
都是表達了一個捐獻的言言.
所以在這樣的情況之下.
我們就是能夠與父上帝和耶穌基督.
有那份團契相交的生活.
我想和大家分享一個我自己的故事.
我在很多年前,三十多年前.
來澳洲讀醫科.
然後我回到香港做醫生.
在政府醫院工作.
頭十幾年我和我的上司.
這個上司是改裝了的.
你不認識他.
完全不是這樣的樣子.
頭十幾年我和他的關係都非常好.
我基本上是他的得力助手.
基本上是他的左右手.
他很多時候要出去開會.
醫院只剩下我一個.
基本上我已經打了色.
他沒有空做.
但十幾年之後有些事情發生了.
令到他對我的改觀一百八十度.
他甚至去到一個位置.
我有一次早上上班開電郵.
收到一個三千字的文章.
晚上三點鐘送出來.
數過有什麼不是.
他對我的不滿很嚴峻.
有一次有另一個病房護士告訴我.
我的上司下班後留下.

$^{241}$過去兩個星期病房.
不停地看病歷.
抄下很多東西.
我之後知道他做了什麼.
兩個星期後他約我見.
然後數我的罪狀.
二十種還是三十種.
總之我在病房處理病人不理想.
逐項逐項向我表達.
這裡做得不好 這裡做得不對.
在這個過程當中很痛苦.
直至有六年的時間.
直至他退休為止.
這個情況才能夠完結.
但在這個過程當中.
我可以做的就是.
學習和耶穌基督和上帝有傳契.
我要將我的難處向上帝祈禱.
直接讀他的話語.
去嘗試明白上帝的心意.
不是馬上有答案.
在這個過程當中.
有件事發生了.
我上師另外做了一件事.
就是對我的憎恨去到一個地步.
他去到一次門診的處境.
因為我在政府醫院工作.
門診的模式大家可能以往經驗過.
就是你來看專科門診的時候.
你不能選醫生.
到時隨機看哪個醫生的狀態.
所以有一次的狀態就是這樣.
我的上師看了一個病人.
上一次是我看的.
然後他就游說病人去投訴我.
他的護士覺得很奇怪.
為什麼會游說病人去投訴另一個醫生.
她覺得很不明白發生什麼事.
她就將這件事跟門診的護士主管分享.
然後護士主管跟我說這件事.

$^{281}$不知道大家知不知道.
其實在香港我們有一個醫生守則.
這是醫務委員會發出的守則.
醫生的守則就是說.
你不可以去游說病人投訴另一個醫生.
這是違反醫生守則的元素.
當然我也知道.
但是他做了這件事.
之後的做法就是.
我將這件事跟我們的部門主任分享.
部門主任聽完也覺得不對勁.
所以要求我的上師向我道歉.
但是其實我有權去醫務委員會投訴他.
這是違反醫生守則.
這是不可以做的事.
他做了.
但是上帝的話語在我的心裡.
其實我猶豫了一段時間.
但是有一句話.
上帝的話語在我心裡.
不斷地浮現.
就是「身冤在我,我必報應」.
我有權做這件事.
但是上帝的話提醒我.
上帝會對付他.
不用我出手去對付他.
當然最終的結論就是.
我沒有去醫務委員會投訴他.
因為上帝的話告訴我.
「身冤在我,我必報應」.
上帝會知道他的惡的狀態.
所以我們去反省一下.
我們與耶穌基督的關係是怎樣的.
當我們說要跟弟兄姊妹之間.
有一個好的團結生活的時候.
我們首先要看看.
我們自己與耶穌基督的團契是怎樣的.
我們有沒有定時.
藉著禱告,藉著讀神的話語.
與耶穌基督建立團契生活.

$^{321}$第二點我想和大家分享.
認識神.
因為.
認識神.
在這段經文的第一章.
頭兩節的經文說到.
使徒約翰的經驗.
論到生命之道.
就是從起初以有.
我們親眼見過,看過,親手觸摸過的.
這是使徒約翰自己.
親身跟隨耶穌的過程中.
他親身經歷過的經驗.
所以他與耶穌基督有一個很親密的經驗.
而在這個過程當中.
約翰接著說.
一個很大的宣告就是.
第一章五節這裡說.
神是光.
剛才我們崇拜之前.
Nelson的分享.
牧師的祈禱裡也說了.
神是世界的光.
不僅如此.
約翰的宣稱就是.
神是光.
在他裡面毫無黑暗.
這個宣告對約翰來說.
有什麼特別呢?.
他作為猶太人的他.
他在《猶約聖經》裡.
如何理解神是光呢?.
或者我們可以從.
《出埃及記》十三章這段經文裡.
去看到以色列民族.
他們如何經驗神是光.
因為當他們在曠野漂流的時候.
他們已經出到埃及的時候.
上帝如何帶領他們呢?.
《出埃及記》十三章二十一到二十二節說.

$^{361}$日間耀和華在雲柱中領他們的路.
在夜間裡火柱中的光照他們.
使他們日夜都可以行走.
日間雲柱 夜間有火柱.
總不離開百姓的面前.
所以他們在被救贖之後.
去經歷上帝的光.
引領他們 與他們同在.
而約翰自己在《約翰福音》裡.
如何表達耶穌基督神是光呢?.
在《約翰福音》第一章裡.
講到施洗約翰的來臨.
而施洗約翰來到這個世界的目的.
就是要作見證.
而且這裡講到他要為光作見證.
叫眾人因為可以信他.
而他自己本身不是那個光.
要為光作見證.
這個光是真光 要照耀世上的人.
所以繼續在《約翰福音》第八章裡.
講到耶穌基督自己宣告.
祂是世界的光.
跟隨我的就不在黑暗裡.
必要得著生命的光.
所以這裡很明顯告訴我們.
神是光.
這個光如何去照耀他們.
在以色列的歷史裡被救贖之後.
經驗上帝與他們同在 引領他們.
所以耶穌基督也同樣宣告.
祂是光 照耀我們.
令我們不會再在黑暗裡行走.
這裡我想和大家分享一個.
美國華裔的見證.
不知道大家有沒有聽過袁佑軒的故事.
他信主之後去攻讀神學.
曾經做過慕迪神學院的教授.
但在他年幼九歲的時候.
在他家裡發現一些色情雜誌.
開始感受到在這些色情雜誌的刺激下.

$^{401}$對性慾有一個很強的需求.
當他去到高中的時候.
因為美國的教育制度裡.
他們需要有軍訓.
他進入了海軍陸戰隊的軍訓過程裡.
開始有一個轉變.
就是他覺得有一個同性的傾向.
而到了他二十歲的時候.
他向父母宣佈他是一個同性戀者.
當時他的父母不算耶穌.
不過對他的感受有一個很強烈的不滿的反應.
無論當時他的父母的婚姻處境都不是很理想.
他的父母差不多到了一個位置需要離婚.
而他媽媽在那個過程當中很沮喪.
一方面為了他的婚姻沮喪.
另一方面為了他的兒子是同性戀者有另一個沮喪.
他去到一個位置覺得他人生沒有希望.
所以當時袁有軒已經去了另一個省.
而他媽媽在當時有一個想法.
就是去找他的兒子見他一次面.
然後去自殺.
但在這個過程當中.
他買票坐火車去另一個省探望他的兒子.
在過程當中.
他在車廂裡拿到一個福音小冊子.
給他一個心思的轉向.
他下車後嘗試去找當地的教會去理解.
福音是什麼.
最終他沒有自殺.
他信了耶穌.
因為他信了耶穌之後.
他回去和他的丈夫復和.
以致兩夫婦都能夠相信耶穌.
但袁有軒在這個時間保持他同性戀的狀態.
而在這個時間他正在攻讀博士學位.
但大部分其他時間.
白天去讀書.
但晚上他就花時間在同性俱樂部.
他的人生繼續轉向下坡.
以致他接觸了一些朋友去賣毒品.

$^{441}$而在這個惡劣的環境當中.
他繼續販賣毒品.
竟然他都會做到一個狀態.
賣毒品給他的教授和其他同事.
有一天聯邦政府的警察抓了他坐牢.
在這個沮喪的時候.
他當時已經和他的父母沒有什麼接觸.
但他去到這個惡劣的環境.
他想到他還有父母.
他打電話去找他的父母.
他以為父母知道他去到這麼潦倒的狀態.
一定會大罵他一頓.
但他很特別地發覺.
原來父母在和他溝通的過程中.
是關心他的處境.
他開始發覺他人生中為何去到這麼嚴峻的狀態.
而在他坐牢的過程中.
他找到一本聖經開始去閱讀.
他藉著閱讀聖經的過程中尋找到福音.
但在這個過程中他得到一個很震驚的消息.
就是他患上了愛滋病.
但無論如何在這個過程中.
他開始接受福音,接受醫療的療程.
聖靈令他覺得他要脫離同性捆綁的狀態.
這不是即時發生的事.
不過他找到信仰之後.
他選擇了耶穌而不是同性戀.
在這個過程中他找到了真光,找到了耶穌.
因為耶穌賜他永生.
耶穌帶他走出黑暗的處境.
所以他藉著禱告,讀聖經.
他能夠和耶穌基督,和父上有那份團契生活.
他認識到耶穌是真光.
帶他走出黑暗的狀態.
第三點我想和大家分享的是「向神認罪」.
在一章六節到二章二節這段經文.
比較長,不過如果大家細心去讀這段經文.
這裡講述了三個假象和三個真相的狀態.
大家看不清楚,我會慢慢再和大家解釋.
不過大家仔細去看的時候.

$^{481}$在第六節,第八節和第十節.
這裡講述了假象的狀態.
約翰用了「我們若說」的方法去表達出來.
而真相的狀態是怎樣呢?.
就是在第七節,第九節,二章一,二節裡.
作者用「若」帶出真相的處境.
第一個對比假象和真相是怎樣的狀態呢?.
第六節說「我們若說,是與神相交,乃仍在黑暗裡行」.
就是在說方話「不行真理了」.
而真相是「我們若在光明中行,如同神在光明中彼此相交」.
「祂的兒子的血,洗淨我們一切的罪」.
這裡的假象是什麼呢?.
如果我們說我們與神有緊密的團契.
但是在黑暗裡行,是什麼意思呢?.
這裡所說的「行」跟著再說「不行真理」.
第一個「行」和第二個「行」有點不同.
第一個「行」的意思是持續性在黑暗裡.
剛才我們說了若翰早一點的描述.
神就是光,在祂裡面沒有黑暗.
所以這裡說什麼叫做「行」在黑暗裡呢?.
就是告訴我們我們與上帝沒有那份關係.
而且這個狀態是持續性的狀態.
所以我們就不在真理裡.
第七節所說的是「我們行在光明中」.
即是我們按照剛才一章五節所說.
神就是光的狀態.
我們真心真意與耶穌基督,與父上帝連結的時候.
我們就是「行在光明之中」.
表明了我們與上帝有那份的團契.
在這個狀態之下 耶穌基督的血就洗淨了我們一切一切.
若翰有沒有說關於我們與弟兄姊妹相交的假象例子呢?.
《約翰一書》後的第二章九到十一節說.
「人若說自己在光明中 卻恨他的弟兄.
他到如今還是在黑暗裡.
愛弟兄的就是在光明中.
在他並沒有跌倒的緣由.
唯獨恨弟兄仍在黑暗裡 且在黑暗裡行.
也不知道往那裡 因為黑暗叫他黑暗」.
所以這裡說的是我們與弟兄姊妹之間的假象.
當我們去恨弟兄的時候.

$^{521}$我們就在黑暗當中.
我們不在光明當中.
這就是我們在弟兄姊妹之間那份恨弟兄的罪.
所以我們想想.
在我們團契弟兄姊妹相交的過程當中.
我們有沒有那份恨的元素出現呢?.
或者在這一刻大家未必馬上想到.
我們待會再和大家分享一些例子.
我們去第二個假象和真相.
第八節和第九節都這樣說.
「我們若說自己無罪 便是自欺.
真理不在我們裡面」.
真相就是我們約定自己的罪.
神是信實,是公義的.
祂要赦免我們的罪.
洗淨我們一切的不義.
第八節說我們說自己無罪.
這裡所說的罪是單數.
想表達什麼意思呢?.
其實罪是我們人性裡面的狀態.
經文想表達的是.
我們天生就有罪.
這是聖經告訴我們的.
當我們遲早犯罪的時候.
這個罪性就存留在我們人類當中.
世世代代污染了我們上帝所創造的世界.
我們的罪性在我們的心底裡.
所以我們需要耶穌基督的救贖.
第九節說如果我們認自己的罪.
這個罪的表達是一個複數.
即是當我們知道自己有罪性.
但其實我們很多時候.
很容易在我們人生的不同時間裡.
不同處境之下.
我們都有不同的罪惡的狀態出現.
但我們若願意在神的面前承認的時候.
神宣告祂是信實和公義的.
當我們願意向著神承認我們的罪的時候.
祂就必赦免我們的罪.
這件事很特別.

$^{561}$我們再看清楚認罪是什麼一回事.
認罪這個字.
在新約聖經的表達.
從兩個字出現.
一個是相同.
一個是說話.
其實是說什麼意思呢?.
認罪就是我們願意說出真理.
我們願意說出認同上帝說話的內容.
就是認罪.
所以這裡的真理就是說.
神是信實 神是公義.
我們承認上帝是一個這樣的上帝的時候.
我們知道我們有罪.
我們願意回到祂的面前.
這就是一個認罪的行動.
所以這裡所說的.
在團契生活裡面.
我們要承認什麼罪呢?.
約翰繼續給我們一些例子.
剛才我們說了.
恆弟兄這個表達在約翰一書裡重複地說了.
在第四章裡再說一次.
都是恆弟兄.
所以這裡所說的.
恆弟兄是一個什麼樣的處境呢?.
我想和大家分享一個例子.
我之前在另外一間實習教會的時候.
認識一位姊妹.
她也是移民.
她媽媽在香港住.
可能我們在座很多都是.
同一個處境 我自己也是.
但是她媽媽的身體狀態都是很不理想的.
差不多她每年過去的十幾年.
每年都會回香港最少一次去探望母親.
而在她團契的過程當中.
她在澳洲和另一個團契的領袖.
分享過她媽媽的身體狀態.
一段好幾年的時間.

$^{601}$直到這個姊妹的媽媽.
在身體狀態不理想的情況下已經去世.
她也將這個訊息和這個領袖分享.
過了一年半載之後.
這位姊妹有機會回香港旅遊探親.
回來之後.
這位領袖就去關心她.
關心這位姊妹.
這位領袖就這樣說.
「文豪的姊妹 你媽媽的身體怎麼樣?」.
這位姊妹心裡很難受.
「其實我媽媽已經去世一年多兩年了」.
「我之前跟你分享過」.
「在那個時間你也關心過我」.
「你很明白我心情好不好」.
「你知道的」.
「為什麼一年多兩年後」.
「你完全忘記了我媽媽的一生」.
接著這位姊妹對這位領袖很不滿.
心裡有份怨恨.
就是這裡所說的恨弟兄.
這個領袖做得不理想.
忘記了這位姊妹的一些過往歷史.
但這位姊妹心裡積聚了這份怨氣.
都一段時間.
但她在這個過程當中.
她覺得不行了.
我不可能繼續持續這個怨氣在我心裡.
所以她藉著禱告之後.
她找到機會再跟這個團契的領袖分享.
對她的怨氣,對她的不滿.
這位領袖在過程當中都發覺有點怪.
為什麼這位姊妹跟她疏遠了這麼多.
不明白.
但藉著這個大家有和解的機會.
她向這位姊妹道歉.
因為她的疏忽,她的遺忘.
令到這位姊妹很心痛.
在我們當中會不會也有類似的事件發生呢?.
其實我自己也試過.

$^{641}$遺漏了一些東西.
突然間說錯人名.
以前也記得會叫這個人名.
但突然間叫錯了.
這些事件可能也會發生.
當然可能未去到剛才的例子這麼嚴重的處境.
第三個假象和真相.
第十節說:我們說自己沒有犯過罪.
就是神說謊.
祂的道不在我們心中.
而二章一二節的真相是什麼呢?.
就是約翰告訴我們.
如果有人犯罪.
在我們腹中有一個鐘堡.
就是耶穌基督.
祂為我們的罪作了挽回制.
不單只是為我們的罪.
也是為了普天下人的罪.
所以約翰繼續說.
我們信徒相交之中的罪.
是一個怎樣的處境呢?.
約翰一書裡繼續有一個很大的主題.
就是彼此相愛.
在第三章裡說了一次.
而三章十七到十八節說了一個例子.
凡有世上財物的看見.
他弟兄有需要腳硬著心腸.
那怎會有神的愛存在他們裡面呢?.
孩子們.
我們相愛不要用言語和舌頭.
要用行為和真理.
這裡說了一個例子.
我們看見弟兄姊妹有需要.
但我們不理會他們.
這就是我們相交中的罪.
可能這件物質性的事.
我們未必會那麼容易犯.
或者我們也會容易一點.
知道弟兄姊妹在處境上的需要.
物質上的需要.

$^{681}$我們去做到彼此相愛.
而這裡我想和大家分享另一個例子.
是在我們教會發生一個沒有彼此相愛的例子.
幾個月前.
大家都知道香港一個很嚴重的事故.
黃福苑的大火.
在那個時間裡.
其實我們教會有一個很特別的弟兄.
他的背景和我們香港人有少許不同.
他對於大陸政府有一個很特別的看法.
他很容易想到國內政府會做很多事令我們受損.
在這個事件發生的時候.
有弟兄姊妹在我們的教會群組裡表達這些訊息.
慘痛的感受.
這位弟兄就說了一句話.
他說這個事件其實是大陸政府找人做手腳做出來的.
要害死香港人.
當然這個說法在我們理解裡面.
機會率應該不是很高.
但因為他的背景有些很特別的想法.
他有這樣的表達.
這個弟兄其實他和教會很多弟兄姊妹都不太熟絡.
很多弟兄姊妹都不太認識他.
但無論如何他在這個處境之下.
有不知多少個弟兄姊妹.
甚至有一些以往離開了我們教會的教會領袖.
已經不存在了.
所以一些以往的教會領袖都圍攻他.
說他不是.
當然這個離開了教會的以往的教會領袖.
對這個弟兄更加不認識.
我作為傳道.
我明白他為什麼這樣表達.
因為我和他都是一個細胞小組.
我都花了很長的時間去了解他.
嘗試明白他在想什麼.
因為他以往都表達過類似的狀態.
他有這樣的想法.
我明白.
當然我不是說他對.

$^{721}$但問題是當我們不明白另一個人.
他有些很特別的表述的時候.
我們用一個圍攻他的方法.
不去搞清楚他為什麼這樣說話.
而去說他不是的時候.
其實對他都是有傷害的.
我們很多時候會不會在團契生活裡面.
都有類似的事情發生.
其實我們和接觸的弟兄姊妹.
因為他們的認識不足夠.
以致他們有特別的表述的時候.
我們有很大的差異.
以致我們有些不滿的表述.
都傷害了他們.
所以今天我想和大家分享.
怎樣我們可以活出.
合神心意的團契生活.
首先我們要和基督的生命連結.
我們需要去認識神.
我們要知道我們有罪性.
我們很容易犯罪.
我們需要向神認罪.
讓我們一起低頭禱告.
因為你讓我們知道.
我們需要和弟兄姊妹有好的團契生活的時候.
我們首先需要和你的生命連結.
我們需要去認識你.
我們更加需要向你認罪.
因為當我們願意在你面前認罪悔改的時候.
你應許你會赦免我們.
因為你是那位滿有豐盛恩典的上帝.
禱告放在基督的面前.
阿們.
\newpage



\section{路加福音 10:25-37}
\label{sec:031612}
\textbf{邊gwo 係我ge 鄰舍?}~~~~鄺偉志牧師
\newline
\newline
連結: \href{http://www.church.com.hk/acms/content.asp?site=cdc\&op=show\&type=product\&code=031612\&layout=sermon}{\texttt{031612}} ~~~~ 語音日期: 2026-02-08
\newline
\newline
\hyperref[sec:031611]{\small{< < < PREV SERMON < < <}}
~
\hyperref[sec:index_chronic]{\small{[返順時目]}}
~
\hyperref[sec:031613]{\small{> > > NEXT SERMON > > >}}
\newline
\newline
路加福音 10:25-37
\newline
\begin{longtable}{cl}
\hline
\hline
章節 & 經文 (和合本修訂版)\\
\hline
10:25 & \begin{tabularx}{0.7\textwidth}{X} 有一個律法師起來試探耶穌，說：「老師！我該做甚麼才可以承受永生？」 \end{tabularx} \\ \\ \relax
10:26 & \begin{tabularx}{0.7\textwidth}{X} 耶穌對他說：「律法上寫的是甚麼？你是怎樣念的呢？」 \end{tabularx} \\ \\ \relax
10:27 & \begin{tabularx}{0.7\textwidth}{X} 他回答說：「你要盡心、盡性、盡力、盡意愛主—你的神，又要愛鄰如己。」 \end{tabularx} \\ \\ \relax
10:28 & \begin{tabularx}{0.7\textwidth}{X} 耶穌對他說：「你回答得正確，你這樣做就會得永生。」 \end{tabularx} \\ \\ \relax
10:29 & \begin{tabularx}{0.7\textwidth}{X} 那人要證明自己有理，就對耶穌說：「誰是我的鄰舍呢？」 \end{tabularx} \\ \\ \relax
10:30 & \begin{tabularx}{0.7\textwidth}{X} 耶穌回答：「有一個人從耶路撒冷下耶利哥去，落在強盜手中。他們剝去他的衣裳，把他打個半死，丟下他走了。 \end{tabularx} \\ \\ \relax
10:31 & \begin{tabularx}{0.7\textwidth}{X} 偶然有一個祭司從那條路下來，看見他就從另一邊過去了。 \end{tabularx} \\ \\ \relax
10:32 & \begin{tabularx}{0.7\textwidth}{X} 又有一個利未人來到那裡，看見他，也照樣從另一邊過去了。 \end{tabularx} \\ \\ \relax
10:33 & \begin{tabularx}{0.7\textwidth}{X} 可是，有一個撒瑪利亞人路過那裡，看見他就動了慈心， \end{tabularx} \\ \\ \relax
10:34 & \begin{tabularx}{0.7\textwidth}{X} 上前用油和酒倒在他的傷處，包裹好了，扶他騎上自己的牲口，帶他到旅店裡去，照應他。 \end{tabularx} \\ \\ \relax
10:35 & \begin{tabularx}{0.7\textwidth}{X} 第二天，他拿出兩個銀幣來，交給店主，說：『請你照應他，額外的費用，我回來時會還你。』 \end{tabularx} \\ \\ \relax
10:36 & \begin{tabularx}{0.7\textwidth}{X} 你想，這三個人哪一個是落在強盜手中那人的鄰舍呢？」 \end{tabularx} \\ \\ \relax
10:37 & \begin{tabularx}{0.7\textwidth}{X} 他說：「是憐憫他的。」耶穌對他說：「你去，照樣做吧！」 \end{tabularx} \\ \\
[1ex]
\hline
\hline
\end{longtable}
$^{1}$我們一起祈禱.
天父,我們將以下時間交給主.
讓我們真的可以領受經文的內容.
雖然時間不是很長.
但是讓我們可以重新思想這段如此熟悉的經文.
讓我們有重新的亮光鼓勵.
以致我們有愛淪捨.
我們有傳福音.
有愛主的行動.
祈禱,奉耶穌基督的名字.
阿們.
誰是我的淪捨.
好撒瑪利亞人的故事.
很多弟兄姊妹都會聽過.
以及如果你翻譯用主日學剎說過的.
所以這次也用一個方法.
跟大家去思想.
也配合傳道主日的內容.
時間到了,不是吧?.
聽到嘟嘟嘟的聲音.
害怕.
OK,因為我只剩下15分鐘左右.
其實我們在方性查經剛才提過.
有一個課程.
跟這些微信的朋友去思考聖經.
我們就用一些耶穌的事蹟的經文.
還有一個方法.
我今天也會用類似的方法.
就是因為經文很熟悉.
聖經這樣寫.
其實實質是怎樣呢?.
很多耶穌傳.
耶穌的事蹟.
拍了視頻.
我們就會看那些視頻.
然後又對一下經文.
然後一起發現有什麼相同與不同.
就令我們更加深刻.
聖經原來是想表達這些意思的.
在我準備這個經文的時候.

$^{41}$好撒瑪利亞人的故事這麼熟悉.
找來找去都是卡通片而已.
沒理由跟大家放卡通片去對照.
但是我發現有一首穩如回應詩歌.
這首詩歌叫做好撒瑪利亞人.
他們是音樂家.
他們出了他們創作的一首詩歌.
我也問他可不可以給我用.
他說OK的.
所以這段經文.
剛才主席已經讀出了.
大家應該也會認識這個故事.
你說不太熟悉不要緊.
我們這樣對照一下.
那首好撒瑪利亞人的詩歌.
是很少的詩歌.
是用聖經的故事作為題材.
整個故事放在詩歌裡面.
很少的.
我們可以看一下創作的歌詞.
他要表達的東西.
跟本身的經文.
有些對照的時候.
我們就會學習一下.
整首歌.
待會我們回應詩歌也會唱.
你可以學一下.
也蠻難唱的.
四分多鐘.
但是我就用他第一段整個故事.
兩分鐘四十五秒.
大家請聽聽看歌詞.
這個很可愛的.
有些圖片.
讓你回顧整個好撒瑪利亞人的故事.
我現在播出來了.
留意一下.
對照一下.
你拿出聖經或者經文來看就可以了.
從前有一個撒瑪利亞人一天走到山邊.

$^{81}$他沒有名氣也沒有實名.
根本就沒有特別.
但他擁有一顆溫柔善良的心.
一個人們會喜歡上帝也喜悅撒瑪利亞的心.
他偶爾看見山邊有個人躺在地上待久.
沒有遺傳也沒有氣息.
一個人半死掙扎但沒有人求他.
也沒有人幫他.
只有這一個撒瑪利亞人他懂了慈心.
好撒瑪利亞人他獻出溫柔關懷愛心.
去幫助一個受了傷的陌生人.
好撒瑪利亞人他借出生命的果子就像耶穌的教導.
你要盡心盡力也要盡力盡力愛著主.
因為上帝鮮愛我們.
你要盡心盡力還要盡力盡力愛人如今.
這樣才配得永遠的生命.
究竟有什麼特別的地方呢?.
我們對一對.
不問大家問題了.
你會看到整段經文分兩段.
25-29節是一個背景.
這也是重要的.
為什麼耶穌會說撒瑪利亞人的一個例子一個故事.
30-37節就是故事本身.
那首詩歌我細心聽了很多次.
黃色的地方就是他提到的地方.
沒有提到的地方就是白色的地方.
其實他在詩歌裡面描述那件事.
但其實聖經也有描述.
聖經說為什麼那人在山邊.
其實他沒有說山邊.
大概也是因為耶路撒冷下耶利哥.
在山邊他遇到強盜.
很慘很不幸的際遇.
所以就打他半死.
詩歌有提到半死.
還有他很慘的處境.
撒瑪利亞人在詩歌裡面有很多東西加進去.
第一就是這個撒瑪利亞人是好的.
其實很多人解釋聖經.

$^{121}$整段都沒有說撒瑪利亞人是好.
但為什麼我們說撒瑪利亞人是好呢?.
因為普通正常人看這個故事都知道.
那個撒瑪利亞人是好人.
所以加一個好字進去.
他又說撒瑪利亞的心.
我想這個不是聖經的東西.
不過很有意思.
有沒有一個撒瑪利亞的心.
等下跟大家探討一下什麼叫撒瑪利亞的心.
他在詩歌裡面特別強調.
這個人沒有名氣.
只是普通人.
但他得到神和人的喜愛.
他強調了這一點.
他還加了一件事.
就是.
他有聖靈的果子.
溫柔.
還有神才愛我們.
所以我們要愛人愛神.
我覺得是非常好的.
有時候看詩歌的創作.
都看到他幫我們透過音樂.
還有思考.
這個故事重新有一個方向.
你會看到.
詩歌都很好.
我覺得他不是單單說這個故事.
他還說了背景.
所以他有提到第27節.
大家看到這裡.
他說要盡心盡性盡義.
愛主你的神又要愛倫如己.
這個很重要.
因為.
為什麼撒瑪利亞人.
這個故事.
是表達到.
誰是你的倫社.

$^{161}$其實.
因為是在探討這個問題.
你說你愛神.
那個人挑戰耶穌.
走來挑戰耶穌.
我做什麼才可以.
承受永生.
這首歌就說.
我們有一個這樣的生命表演.
才可以配得永遠的生命.
我覺得這個詩歌寫得很好.
這樣才可以配得.
其實這個.
律法師.
挑戰耶穌.
我承受永生.
我應該做什麼.
耶穌就回答他.
你應該怎樣才可以.
做到愛人如己.
因為.
律法師很熟悉聖經.
一回答就能回答.
耶穌一問他.
他就說.
我要盡心盡性盡義愛神.
我們在教會.
做基督徒應該怎樣.
愛神愛人.
跟你真的能不能做到.
是兩回事.
當然耶穌說你答對了.
一小問題.
都要講一講.
你不是說信耶穌就能永生嗎.
為什麼你做這些就能永生.
說真的你一個人愛神愛人.
你就已經是信神的人.
你不信神你怎樣愛神.
你愛神愛人就包含了信神.

$^{201}$但問題是.
怎樣去得到永生.
其實你不是說得到一件事.
而是說.
你得到這個生命.
你活出這個生命.
然後那個人就不服氣.
誰是我的靈寫.
就是這樣.
講出這個撒瑪利亞人的故事.
你看到.
這個撒瑪利亞人的故事.
其實.
情節不用怎樣描述.
他將那個人.
幫助他.
這個受傷的人.
半死的人.
他講到這個撒瑪利亞人.
不止幫他包紮.
之後一日.
還要是包底.
不夠錢的我幫你包.
你看到他.
做到足.
只不過是見到一個人.
他落在一個患難的處境.
然後他有一個同情心.
來幫助他.
當然這個詩歌.
特別是講這個人怎樣好.
但沒有講的其實很重要.
有兩個人.
pass through.
其實有兩個人已經見到.
他還要強調這兩個人是什麼人.
就是律法師那一類.
有祭司,有利未人,宗教人士.
應該那些人.
你是一個宗教人士.

$^{241}$你那麼會律法.
你見到那些人就馬上會幫他.
耶穌講故事.
是很厲害的.
其實這個故事是針對那個人.
你看看這個故事.
其實這個故事肯定沒有發生過.
不過他特意用這個故事.
去針對.
你見不見到兩個人.
都是宗教人士.
應該他是做得好的.
但他們都會見死不救.
但反過來.
猶太人.
就是撒瑪利亞人是混血.
簡單來說就是世仇.
有一個世仇.
但那個人竟然有一個同情心.
去幫助那個人.
是他的鄰居.
所以其實撒瑪利亞人的故事.
如果看到這裡.
其實什麼故事都可以.
為什麼呢?.
因為我們在想.
如果我們跟耶穌有這樣的對話.
當我們問.
誰是我的鄰居的時候.
耶穌可能會和我們講不同的故事.
你想想你最討厭誰?.
你想想你最討厭哪一類人?.
譬如你說.
大家不要猜了.
好像都幾好人.
總之你最討厭那個人.
就是那個撒瑪利亞人.
你明白嗎?.
那你自己就是那兩個人.
那耶穌就去編故事.

$^{281}$如果我聽了.
就是耶穌就講故事給我聽.
有一個傳道人走過.
又不理會.
不是我們的執事.
然後有一個人.
就是一個在外面的黑衣.
他走過會理會.
你明白嗎?.
套入你最不喜歡的人裡面.
套入你自己在故事裡面.
所以你說什麼叫撒瑪利亞的心呢?.
其實就.
可以說沒有這件事.
撒瑪利亞的心就是.
總之一個人.
他見到一個人.
他有一個同情心.
他是想來幫他.
這個就是了.
是很認真去幫他.
這段故事.
講到我們中國人.
很熟悉的儒家.
萬子.
大家很多都讀過.
不是拋書包.
但這個人就有不忍人之心.
或者惻隱之心.
我們是認識的中國人.
惻隱之心.
意思你自己慢慢看.
為什麼?.
他說如果你沒有惻隱之心.
就不是人.
現在他是在講.
是好的.
這就證明本性是好的.
你看到一個小孩.
跌在井裡你都會救.

$^{321}$你不用想.
不是因為要討好他的父母.
不是因為要認識更多朋友.
是.
什麼叫惻隱呢?.
惻就是悲傷.
忍就是痛苦.
你看到人落於患難悲傷痛苦.
我們就有一種同情心.
憐憫的心.
如果你講萬子講的.
人人應該是有的.
沒有就不是人.
還有惻隱之心.
人之端也.
是一個開端.
一個好.
惻隱之心就是一個好的開始.
所以其實我想到.
薩馬利亞人的故事.
其實這個薩馬利亞人.
只不過是他.
見到一個有需要的人.
他動了一個惻隱之心.
接著他就做事.
其實每個人都有的.
每個人都應該這樣做.
但是有一個問題.
你看到.
這裡我不講了.
我想討論一下.
這個插曲.
因為他特別強調.
這兩個.
是一個宗教人士.
應該是好人.
但是他竟然因為一些原因.
不去幫助這個人.
你明白嗎?.
究竟為什麼會這樣呢?.

$^{361}$連良知也壓下去.
有同情心也壓下去.
而且是他的同胞.
我們會想.
就算我們是基督徒.
說我們有淑女的生命.
使我們埋沒了良知和憐憫的行動.
這兩個人就被我們這樣想.
你說薩馬利亞人很好.
但是我們會不會見到.
需要人的事.
需要幫助的事.
我們時而不見.
或者有很多藉口.
但是當一個小孩掉進井裡.
你也會去救他.
我趕時間去上班.
是不是這樣呢?.
但是這個世界真的有這樣的人.
但我相信不是基督徒.
我們要想一個問題.
會不會有時候我們.
有一些想法.
可能是你的身份.
或者一些事.
在你面前.
或者有傳福音的需要.
我們也會埋沒了.
傳福音的需要也有關係.
你信不信我們信主上天堂.
你信不信我們.
一些未信的朋友.
將來會在地獄受苦.
信不信.
如果是信的話.
我們看到身邊很多人未信主.
他會下地獄的.
我們沒有那種心.
你不要下地獄.
我不有空.

$^{401}$其實傳福音是絕對.
和這個有關.
我們知道這個真理的時候.
我們是會回應的.
當然我不是說我們五時三刻.
就是這樣做.
但我們一定有這種.
差人之心.
差人之心的位置是什麼.
你見到人將會受苦.
或者不要說天堂地獄那麼遠.
但是你在我們身邊.
神可能給我們看見.
很多我們需要關心的人.
我們有沒有關心.
而那些可能是一些有需要福音的人.
我們有沒有動一個心.
然後去做.
有一個例子我曾經說過.
我的經歷就是那時候在拍拖.
我太太.
現在太太.
那時候她住竹苑.
我送女朋友回家.
在過程中.
在那個.
屋村裡面.
有一個人她就伏在那裡.
在那個平台.
屋村有很多平台.
有一個人伏在那裡.
手腳縮起來.
她的身材很小.
所以我初初經過.
看不到她的樣子.
我以為是小朋友.
可能十多歲.
經過了.
和太太經過了.
不是太太.

$^{441}$和她一起經過.
送她回家.
但她的心不安.
如果那是小朋友怎麼辦.
在那裡伏在那裡.
早上.
長話短說.
後來我回去.
看看那個人是誰.
找他出來.
原來是一個醉酒的人.
後來慢慢發現.
是一個中年男士.
原來他是吸毒的.
我初初.
幫助他.
問他為何這樣.
原來他喝啤酒醉醺醺.
一直在那裡.
我送我女朋友回家.
回來後發現她仍然在那裡.
原來她一直都在那裡.
一個這樣的人.
到幫助她.
了解了之後.
知道她吸毒.
如果要幫她戒毒怎麼辦.
當時我在短訊中心.
受訓.
有一個同學.
他是國內人.
都是一些戒毒人士.
和我爹一起歡呼.
我問他如何幫助她.
他第一樣跟我說.
我很深刻,一輩子都記得.
他說你幫她.
是要上身的.
不是鬼上身.
是要上身的.

$^{481}$你要幫這些人,是要上身的.
什麼意思.
你要預備,你要付出.
可能你要幫助她.
帶她戒毒.
有很多.
這些不止一次戒毒.
可能有很多問題.
我真的想清楚什麼是上身.
然後他說好.
繼續跟進.
打電話給她的.
姐姐跟她談.
姐姐在家裡完全不理會這個人.
沒有希望.
後來就說服她.
幫她戒毒.
我自己回想.
我不知道有沒有.
這樣的事情在你身邊發生.
你有沒有.
都是那個問題.
你有沒有面對神.
給你看到一個需要.
你禁了你的良知.
禁了聖靈的感動.
以致我們沒有做.
或者我們看到一些.
很明顯的真理.
很多人如果沒有信主.
是一個很不妥的處境.
將來現在都會有很多困難.
如果神.
是給我們看到的.
我們願不願意去回應.
其實從良知開始.
這個撒瑪利亞心.
這個憐憫的心.
是很自然會有的.
何況我們有神的生命.

$^{521}$所以我就問.
大家的問題.
我們有沒有撒瑪利亞心呢.
只要這兩個人.
我們是不是呢.
這個故事我不說了.
因為沒時間了.
但是說一句.
這個奧斯丁.
剛剛星期一.
他是少年英雄.
他游泳也過不了.
但在那個情況.
我是想說什麼呢.
Love will find a way.
他不做也不行.
他就不斷祈禱.
不斷.
見證很好.
接著要受浸.
要答應耶穌.
變成一個超人.
但其實之後是神保守他.
但如果我們因為愛的時候.
我要上身.
我要面對那個需要.
我要接觸他.
你自然會有條路.
撒瑪利亞人就自然會做事.
God will make a way.
神也會給我們一條路.
我們怎樣去服事.
我們一起祈禱吧.
天父我們謝謝你.
讓我們一起思想.
好撒瑪利亞人的故事.
主我們求你.
將你的經文.
你的故事.
你的訊息放在我們心裡.

$^{561}$鼓勵我們.
在我們的事奉.
在我們的愛人.
愛神的誡命裡面.
在我們傳福音的課題裡面.
亦都願意.
來回應我們的主.
多謝主聽我們禱告.
奉耶穌基督的名字.
\newpage



\section{民數記 6:22-27}
\label{sec:031613}
\textbf{祝福人的人有福了!(粵/普)}~~~~許永勝傳道
\newline
\newline
連結: \href{http://www.church.com.hk/acms/content.asp?site=cdc\&op=show\&type=product\&code=031613\&layout=sermon}{\texttt{031613}} ~~~~ 語音日期: 2026-02-08
\newline
\newline
\hyperref[sec:031612]{\small{< < < PREV SERMON < < <}}
~
\hyperref[sec:index_chronic]{\small{[返順時目]}}
~
\hyperref[sec:031614]{\small{> > > NEXT SERMON > > >}}
\newline
\newline
民數記 6:22-27
\newline
\begin{longtable}{cl}
\hline
\hline
章節 & 經文 (和合本修訂版)\\
\hline
6:22 & \begin{tabularx}{0.7\textwidth}{X} 耶和華吩咐摩西說： \end{tabularx} \\ \\ \relax
6:23 & \begin{tabularx}{0.7\textwidth}{X} 「你要吩咐亞倫和他兒子說：你們要這樣為以色列人祝福，對他們說： \end{tabularx} \\ \\ \relax
6:24 & \begin{tabularx}{0.7\textwidth}{X} 『願耶和華賜福給你，保護你。 \end{tabularx} \\ \\ \relax
6:25 & \begin{tabularx}{0.7\textwidth}{X} 願耶和華使他的臉光照你，賜恩給你。 \end{tabularx} \\ \\ \relax
6:26 & \begin{tabularx}{0.7\textwidth}{X} 願耶和華向你仰臉，賜你平安。』 \end{tabularx} \\ \\ \relax
6:27 & \begin{tabularx}{0.7\textwidth}{X} 「他們要如此奉我的名為以色列人祝福；我也要賜福給他們。」 \end{tabularx} \\ \\
[1ex]
\hline
\hline
\end{longtable}
$^{1}$弟兄姊妹平安.
今天是二月一日.
剛踏入二月我就已經感受到農曆新年快到了.
除了中港澳台灣以及全球華人的社區過農曆年之外.
韓國同樣也會慶祝農曆新年.
不過兩個民族在過農曆年的傳統上.
還是有些相同也有些差異.
相同的是中國和韓國慶祝農曆年的核心.
都是以家庭團聚為核心.
不過韓國比起中國在以家庭為核心上會更進一步.
他們在初一的時候穿著傳統的服裝.
小朋友會向家中的家長父母行跪拜大禮.
在過年的飲食文化上世界各地不同的地方都有不同的習慣.
至於過年派利是或者送禮物.
這方面兩個民族都會做的.
但有些分別的就是中國派利是的時候.
或者送禮的時候喜歡在包裝上以紅色為主.
象徵喜慶或者吉祥的意思.
韓國會以白色信封為包利是.
在韓國卻是以白色信封來裝利是.
是不是有點奇怪呢.
他們送禮的時候就不會講究包裝和顏色.
最後有些相同和不相同的分別.
就是在過年的時候向其他人發出新年的祝福.
在韓國無論是除夕後的新年或者農曆的新年.
韓國人都只會講「新年快樂」.
韓國人就會講這句話.
意思就是說新年快樂.
但全世界的華人社會就不一樣了.
因為在新年的時候在祝福方面是多得不得了.
可能在下面聽到那幾位韓國的弟兄姊妹就不認同了.
認為韓國除了新年快樂這句話.
其實還有其他祝福的說話.
但我好像沒怎麼聽過.
我不是說他們完全沒有其他的祝福話.
但我很肯定我們華人是很多的.
是不是很多呢?其實是看不完的.
除了一些很正常的祝福語之外.
我們香港在新年的祝福上面是比中國其他地方更加創新更加有趣的.
例如說要什麼有什麼.

$^{41}$「荷包醬醬」翻譯得很好.
「越撈越掂」「越來越順利」「色擊五肥」.
「無論怎麼吃都不會長胖」「怎麼睡都行」.
「無論在什麼樣的情況下你都可以睡得香」.
「保住工作」「我們可以保存我們的工作」等等.
其實還有很多祝福語是你想不到的都有.
所以我覺得香港人在過年期間是最喜歡祝福別人的人.
弟兄姊妹,作為基督徒我們是否都會追隨過年的傳統呢?.
在大年初一的時候我們會和身邊的人發出不同的祝福.
在大年初一我們向身邊的人發出祝福.
但是我們祝福別人的時候我們和傳統上的祝福有什麼分別呢?.
當我們講出一些祝福語的時候我們的心態是怎麼樣的呢?.
今天我想透過經文來分享基督徒在祝福其他人的時候三個很重要的心態.
第一個心態就是神先是一切祝福的源頭.
其實神是祝福的源頭我相信沒有人會不知道的.
連拜偶像的人都知道這個道理.
有一句話是「拜得神多自有神庇佑」.
向神祈求祝福是人的本性和慾望.
既然大家都知道的事為什麼我還要和大家說在祝福人的時候.
第一個心態就是神先是一切祝福的源頭呢?.
難道我們不是一直都是這樣想的嗎?.
那大家知道為什麼我們老人家會說不要說一些不吉利的話嗎?.
尤其是在過年或者過節的時候更加不可以.
那這種講法背後隱藏的思想是什麼呢?.
是不是講什麼就會發生什麼事的意思呢?.
那講吉利的話就會有好事發生呢?.
那講不吉利的話就會有壞事發生呢?.
是不是這樣的思想呢?.
其實有些人都堅持在節日講祝福和吉利的話.
背後的思想和邏輯其實和念咒語的概念是一樣的.
全世界無論在哪一個地方從古時開始那些人就有這種思想.
在全世界從很古老的時候開始大家都會有這樣的思想.
相信祝福還有咒詛的話都帶著能力.
就好像念咒語一樣.
講吉利的祝福語就會為對方帶來好運.
相反講一些咒詛的話就會為對方帶來惡運的結果.
但是問題來了.
既然所有人都知道神是一切祝福的源頭.
那神在人與人之間彼此祝福的事上有什麼角色和行動呢?.
是不是人用祝福或者咒詛的話語.

$^{81}$就可以命令神去祝福或者咒詛別人呢?.
又或者神所委派的代表可以代表神為人帶來祝福或者咒詛呢?.
還是神以自己的主權去決定賜福給誰或者去咒詛誰呢?.
在聖經的舊約裡面只有神所指派的授權的代表才有資格去祝福別人.
例如做族長的亞伯拉罕,以撒和雅各等等.
還有祭司因為他們是為百姓去獻贖罪祭和執行各種的祭祀的.
他們是人與神之間的中間人.
還有一些是君王和先知他們都是神所揀選和授權去祝福別人.
還有就是做父母的也都被神授權去祝福自己的子女.
我們來看看今天的經文.
第二十二到二十三節.
經文說亞倫和他的兒子們被神授權去祝福百姓.
為什麼呢?.
因為亞倫和他的兒子是耶和華所揀選的祭司.
關於神揀選亞倫和他的後裔做祭司的事我想大家都應該知道的.
考考大家.
耶和華是否只是揀選了亞倫一家和他的後裔做祭司呢?.
亞倫一家所屬的利未人是不是祭司呢?.
認為利未人是祭司的請舉手.
利未人.
後面有人舉手.
利未人是不是祭司呢?.
有些人點頭但沒有舉手.
果然是有人認為是.
但事實上並不是.
並不是所有的利未人都是祭司.
因為在出埃及記二十八章一節.
神只是揀選了亞倫和他的後代成為祭司.
利美支派只是神揀選作為神的長子.
因為神在救以色列人出埃及的時候.
擊殺了埃及所有的長子.
因為神救以色列人出埃及的時候.
擊殺了埃及所有的長子.
而留下以色列的長子.
所以神便以利未人代替傳以色列人的長子.
當作活祭獻給自己.
以長子的身份分別為聖來服侍神.
亞倫家族雖是立位人.
但他們是特別被神揀選成為祭司.
而利未人成為了祭司的幫手.

$^{121}$但是他們沒有限制於代表神祝福別人的權利.
簡單來說就是幫忙幫幫台台.
我們看到神授權代表來祝福.
很符合我們祝福別人還有被人祝福的思想.
被人祝福的人是如何想的呢?.
很多人都覺得在這個世界上有很多人有特別的權柄.
能夠向世人施予祝福.
沒有信仰的人,每年都會見一些風水師,法師等等.
或者尋求一些符咒,希望給自己帶來一些祝福.
有信仰的人,我們知道祝福是來自神.
他們有能力代表神祝福別人.
譬如佛教的師傅.
道教的道長.
天主教的神父.
還有我們基督教的牧師,傳道人.
人們總覺得從他們身上就可以得到救贖,醫治,各種各樣的祝福.
彷彿神的祝福是透過這些少數的人身上傳遞給世人.
以我的家庭為例.
每次當我們許家聚會的時候.
我爸爸就會叫我為家人祝福祈禱.
當然作為傳道的我很樂意為家人祝福祈禱.
但其實我爸爸和任何信主的家人的祈禱.
神都不是一樣照聽的嗎?.
因為神才是一切祝福的源頭.
今天我所選的經文22-27節.
短短6節經文裡,神的名字耶和華出現了4次.
這代表什麼呢?.
還有經文中包含了6個祝福的動詞,全部都是由神主動.
這說明了人在祝福的事實上是完全被動的.
在聖經裡面凡是重覆的都是代表很重要很重要.
說明了一些很重要的真理.
就是神才是一切祝福的源頭.
人是不能夠真正將祝福賜給人的.
唯獨神才可以.
不過神把祂才是祝福源頭的訊息.
交託了給那些代表祂的人.
由他們去告訴其他人知道.
在我們中文的和合本聖經22-27節出現了3組的祝福.
包括24,25,26節.
每一組的祝福的開頭都是有個願耶和華.

$^{161}$其實如果你看看原文和英文的譯本是沒有願字這個字.
而是好像24節那樣用了耶和華賜福.
在這節的經文賜福的動詞上.
它是一個未完成式的事態.
一般是理解為將來式或者必然發生的事.
所以也可以翻譯為耶和華必賜福或者耶和華將賜福.
有時候聖經的翻譯不是說錯不錯的問題.
是用了某些字之後令到讀的人會產生一些誤會或者錯覺.
就好像願這個字.
它本身也可以指向未來.
但是閱讀的人就會感覺好像透過宣佈祝福的人.
是藉著他們的祈願促使神賜福給人.
講白一點就好像我作為傳道為某人去祝福祈禱.
他們就一定可以得到神的祝福.
真的這樣子嗎?.
如果真的是這樣子的話.
祝福的能力和榮耀都是集中在宣佈祝福的人身上.
而不是真正賜福給人的神.
弟兄姊妹.
神才是一切祝福的源頭.
這不是一個老吊牙的認知.
首先你作為被祝福的人.
或者盼望得到祝福的人.
你要知道宣佈祝福和為你祝福的只不過是血有血有肉的人.
真正的福氣是來自神.
所以人是不應該盼望得福的焦點放在人的身上.
或者某某人身上.
而是放在賜福的上帝身上.
另外在新約裡面.
能夠代表神去祝福的不再是一些少數的人.
不再只是集中在少數人身上.
我們一起來誦讀彼得前書二章九節.
所有信主的人你們都被神授權.
可以為別人帶來祝福的訊息.
所以盼望祝福的人.
你只需要找一個真正信耶穌的人.
他們的祈禱和任何牧師或傳道的祈禱一樣有效.
第二如果你是被揀選代表神去祝福別人的人.
或者為人祈禱祝福的人.
或者你盼望為身邊的家人和朋友去送上神的祝福的人.

$^{201}$你就要知道你只是神的工具.
是神的僕人和代表.
真正的能力和祝福都是源自神.
作為神的代表.
我們的責任是告訴別人知道一切的祝福都是從神而來的.
我們要將榮耀和讚美歸給神.
祝福人的目標不是單單宣佈祝福這麼簡單.
而是好像傳福音一樣.
要令對方去認識那個賜福的神.
第二個心態就是要傳揚真正重要的福.
上個禮拜崇拜結束之後.
大家有沒有拿到那個過年的灰春.
這些灰春是我們親愛的姐妹Anita親手寫的.
是不是寫得很漂亮呢.
多謝她的付出.
我想問一下是不是全部的灰春都只有兩組祝福.
因為我得到的是平安喜樂和天恩浩瀚.
Anita說是天恩浩瀚.
Anita所寫的兩組祝福.
剛好很像今天經文中神吩咐亞倫為百姓所做的祝福.
Anita所寫的兩組祝福與核心.
剛好與今天經文的三組祝福的核心是一樣的.
這兩個祝福的核心為平安與恩惠.
Anita寫了平安喜樂.
當然最中心點是平安.
另外寫了天恩浩瀚.
最核心的就是神給我們的恩惠.
Anita你寫的時候是這樣想的嗎.
你是不是真的想將主的平安和恩惠祝福帶給我們.
還是想將神的平安和恩惠帶給別人呢.
我猜都是這樣.
我們一起看看經文的三組祝福的內容是甚麼.
我們看看經文中三組祝福的內容是甚麼.
還有它的意義是甚麼.
首先二十四節.
願耶和華賜福給你保護你.
首先二十四節這樣說.
願耶和華賜福給你保護你.
提到神的保護大家會想到甚麼呢.
當我們提到神的保護大家會想到甚麼呢.

$^{241}$我想到小時候有一種集體的遊戲很受歡迎.
我想到小時候有一些遊戲很受歡迎.
很多時候幼稚園老師很喜歡跟小朋友玩.
很多時候幼稚園的老師很喜歡跟小朋友一起玩.
就是麻鷹捉雞仔.
就是麻鷹捉小雞.
剛才有人已經馬上叫出來了.
剛才已經有人知道.
這種遊戲的玩法重點在於保護.
這個遊戲的重點在於保護.
扮演母雞的人就必須保護小雞免受麻鷹的傷害.
扮演母雞的人要保護他的小雞免受那些麻鷹的傷害.
我們就好像遊戲裡面的小雞一樣.
在危險的時候想抓住一些可靠的人去得到保護.
那我們就像遊戲中的小雞一樣.
當遇到危險的時候想找到一些的保護.
我們的神就是我們的天賦.
祂就是我們隨時的依靠.
我們的神是我們的天賦.
祂是我們隨時的依靠.
至於當時走在曠野的那些以色列百姓.
他們日間就要面對曠野上面的酷熱,口渴.
那對於當時遊走在曠野的以色列百姓.
他們日間面對在曠野的路上的酷熱和口渴.
晚上又要面對可能突發而來的那些猛獸的襲擊.
所以他們急需要一個可以依靠.
而且能夠保護他們.
或者在困境裡面保護的人只有耶和華.
他們只需要一位可以依靠.
而且作為他們的依靠還有保護的耶和華.
他們所盼望的不是摩西,亞倫或者其他的領袖.
他們所盼望的並不是摩西,亞倫還有其他的領袖.
同樣在我們的人生裡面我們需要一位可以依靠和保護我們的.
祂就是我們的神.
25節.
在25節.
願耶和華使祂的臉光照你,賜恩給你.
賜恩就是mercy憐憫的意思.
那與前半句臉光照有什麼關係呢?.
我們大家當中結了婚的弟兄姊妹.

$^{281}$當你做錯事得罪你另一半的時候呢?.
我們在座的弟兄姊妹,結婚的弟兄姊妹.
當你做錯事得罪另一半的時候呢?.
那你的太太或者丈夫會有什麼反應呢?.
你的太太或者丈夫會有什麼反應呢?.
會不會生氣到一個地步呢?.
地步呢就跟你說你走吧,我不想見到你了.
會不會很生氣的說你走吧,我不想見到你.
你不走啊,那我就走了.
如果你不走的話,那我就走了.
那我看結了婚久一點的夫婦.
那你們應該經歷過這個情況.
我看到很多結婚很久的夫婦呢.
可能會有這種情況發生.
因為新婚很甜蜜嘛.
因為新婚的時候比較甜蜜.
那我呢,大家猜一猜,我有沒有經歷過啊?.
那大家猜一猜,我有沒有經歷過呢?.
那激怒太太當然是不會想見到我的了.
那也會把太太弄生氣了,她會不理我.
那我們人面對神的時候也是一樣的.
當我們面對神的時候也會這樣子.
因為我們經常都會犯錯激怒神.
因為我們會常常犯錯激怒神.
那聖經舊約呢,幾乎全部都是記載百姓犯各種各樣的罪去激怒神.
那聖經舊約呢,幾乎全部都是記錄百姓犯各種的罪行.
而惹怒神.
所以這裡耶和,給百姓的祝福是會賜恩憐憫他們.
所以這裡神給百姓的祝福呢,是會賜恩典憐憫他們.
因為神的面是向著我們.
因為神的臉向著我們.
就好像會發光一直這樣照亮著我們.
就好像發光一樣來照亮我們.
令我們不至於跌入這個黑暗裡面.
令我們不至於跌入黑暗裡面.
申命記三十一章十七節那裡說.
如果神氣不再見我們的話.
申命記三十一章十七節這樣說.
如果神生氣而不再見我們的話.
那很多災禍都會臨到我們身上.

$^{321}$那會很多的災禍會臨到我們身上.
所以能夠得到神的憐憫確實是一個極大的祝福.
能夠得到神的憐憫確實是一個極大的祝福.
二十六節是什麼意思呢.
願耶和華向你養臉賜你平安.
那大家說一下.
現今我們世人每天最關注最關注的是什麼呢.
那大家可以說一說.
現今世代的人每一天最關注最關注的事情是什麼呢.
那我就說是手機.
雖然有的人去工作的時候.
在公司電腦的時間很長.
雖然有的人去工作的時候在公司電腦的時間很長.
但你未必是甘心樂意和喜歡的.
你未必是甘心樂意和喜歡的.
但看手機肯定是自願和喜歡的.
那神向我們養臉指的就是神常常去注視著我們.
因為祂喜歡我們.
那這種被神長時間去關注.
令我們可以享受到那種平安的喜樂.
這種被神長時間的關注.
可以使我們享受平安的喜樂.
平安這個字對猶太人來說.
是神所賜祝福中最好的.
那平安這個字對猶太人來說.
一直都是神所賜的.
祝福中是最好的.
最好的.
平安的含義已經包含了人一生所需要的東西.
因為平安的含義已經包括了人一生所需要的.
健康,財富,安穩,和平,滿足等等.
包括健康,財富,安穩,和平,滿足等等.
猶太人每天都會用平安向人打招呼.
猶太人每天都用平安跟別人打招呼.
所以這個習慣和傳統的根據.
就是今天二十六節的經文.
這個習慣和傳統的根據.
就是今天二十六節的經文.
神長時間不斷的關注.
不是僅僅代表一兩次.

$^{361}$神長時間不間斷的關注.
意思是不僅僅代表一次.
也不是代表一兩件祝福的事.
也不只是代表一兩次祝福的話.
神向我們仰念是一種終極的福氣.
就是人與神擁有和諧美好的關係.
弟兄姊妹.
人想得到一切的好處.
全部都在神裡面.
人想得到一切的好處.
盡都在神的裡面.
開始講道的時候給大家看了很多很多祝福語.
開始講道的時候.
我讓大家看了很多的祝福語.
這些都是人想要得到的祝福.
但是怎樣的祝福才值得我們去追求呢?.
那怎樣的祝福才是得到我們的追求呢?.
如果我們在過年的時候.
想為我們最愛最關心的人去祝福的話.
如果在過年的時候.
我們想為自己最愛和最關心的人來祝福的話.
那怎樣的祝福才是真正而且重要的福氣呢?.
怎樣的祝福才是很重要最大的福氣呢?.
我記得五年前我和太太還沒有回韓國生活的時候.
我記得在五年前我和太太還沒有回到韓國生活的時候.
有十九年左右的時間.
我岳父岳母都會每年農曆年期間來香港看我們.
有十九年的時間.
我岳父岳母都會在農曆年期間來香港探望我們.
在我們結婚的時候.
我岳父岳母還不是基督徒.
在我們剛結婚的時候.
我的岳父岳母還沒有信耶穌.
所以我每次他們來香港過年的時候.
我都會定下一個目標.
所以當他們每次來香港過年的時候.
我都會定下一個目標.
我下了決心見到他們的時候.
要將最寶貴的福氣帶給他們.
我下定決心每次見到他們的時候.

$^{401}$我要把最寶貴的福音介紹給他們.
所以連續十六年.
我每一次見到他們我都會跟他們傳福音.
所以連續十六年.
我每一次見到他們的時候我都會跟他們傳福音.
因為所有的祝福都比不上他們能夠得到耶穌救贖的恩惠.
因為所有的祝福都比不上他們得到耶穌的救贖的機會.
有耶穌他們可以跟神建立關係.
有耶穌的話他們可以跟神建立關係.
也可以在神那裡得到各樣的祝福和平安.
也可以在神裡面得到各樣的祝福和平安.
最後一個心態也是今天的主題.
也是今天信息的總結.
就是傳遞祝福的人也會得到祝福.
大家放心因為這個是總結我就不會花很長時間.
大家放心因為是個總結我不會花很長的時間.
我們一起看今天的經文最後一節.
我們一起看今天的經文最後一節27節.
作為神所授權去祝福百姓的亞倫和他的眾子.
他們去祝福人的最重要任務是什麼呢.
他們去祝福人的最後的任務是什麼呢.
祝福的目的又是什麼呢.
是不是說亞倫他們只要將耶和華所說的祝福.
一字不漏地轉告給百姓就可以了.
是不是說亞倫他們只要將耶和華所說的祝福.
一字不漏地轉告給百姓就可以了.
或者是亞倫要奉耶和華的名.
令到百姓不單止聽到祝福的內容.
或者是亞倫要奉耶和華的名.
使百姓不但聽到祝福的內容.
更加要明白祝福的意義.
明白百姓所需要所渴慕的一切福氣都是從神而來的.
明白祝福的意義.
明白百姓所需要所渴慕得到的一切福氣都是從神而來的.
令百姓可以認識神和神建立關係.
可以使百姓可以認識神與神建立關係.
大家覺得哪一種更像是神給亞倫和他一家的使命呢.
前面還是後面呢.
還有亞倫和他的眾子去宣告神的祝福的時候.
是不是為了自己能夠在百姓面前得到尊榮呢.

$^{441}$是不是為了自己在百姓面前得到尊榮呢.
還有在神那裡得到祝福呢.
還有在神的面前得到祝福呢.
還是他們希望能夠完成神交給他們的使命.
還是他們希望能夠完成神交給他們的使命呢.
完成使命是不是已經是神所賜給我們最大的祝福呢.
完成使命是不是已經是神所賜的極大的祝福呢.
最後我想引用哥林多前書九章十六至十七節來鼓勵大家.
最後我想引用哥林多前書九章十六到十七節來鼓勵大家.
\newpage



\section{以弗所書 2:1-10}
\label{sec:031614}
\textbf{活著 InG}~~~~劉燕飛教師
\newline
\newline
連結: \href{http://www.church.com.hk/acms/content.asp?site=cdc\&op=show\&type=product\&code=031614\&layout=sermon}{\texttt{031614}} ~~~~ 語音日期: 2026-02-08
\newline
\newline
\hyperref[sec:031613]{\small{< < < PREV SERMON < < <}}
~
\hyperref[sec:index_chronic]{\small{[返順時目]}}
~
\hyperref[sec:031615]{\small{> > > NEXT SERMON > > >}}
\newline
\newline
以弗所書 2:1-10
\newline
\begin{longtable}{cl}
\hline
\hline
章節 & 經文 (和合本修訂版)\\
\hline
2:1 & \begin{tabularx}{0.7\textwidth}{X} 從前，你們因著自己的過犯罪惡而死了。 \end{tabularx} \\ \\ \relax
2:2 & \begin{tabularx}{0.7\textwidth}{X} 那時，你們在過犯罪惡中生活，隨從今世的風俗，順服空中掌權者的領袖，就是現今在悖逆的人心中運行的邪靈。 \end{tabularx} \\ \\ \relax
2:3 & \begin{tabularx}{0.7\textwidth}{X} 我們從前也都生活在他們當中，放縱肉體的私慾，隨著肉體和心中的意念去做，和別人一樣，生來就是該受懲罰的人。 \end{tabularx} \\ \\ \relax
2:4 & \begin{tabularx}{0.7\textwidth}{X} 然而，神有豐富的憐憫，因著他愛我們的大愛， \end{tabularx} \\ \\ \relax
2:5 & \begin{tabularx}{0.7\textwidth}{X} 竟在我們因過犯而死了的時候，使我們與基督一同活過來—可見你們得救是本乎恩— \end{tabularx} \\ \\ \relax
2:6 & \begin{tabularx}{0.7\textwidth}{X} 他又使我們在基督耶穌裡與他一同復活，一同坐在天上， \end{tabularx} \\ \\ \relax
2:7 & \begin{tabularx}{0.7\textwidth}{X} 為要把他極豐富的恩典，就是他在基督耶穌裡向我們所施的恩慈，顯明給後來的世代。 \end{tabularx} \\ \\ \relax
2:8 & \begin{tabularx}{0.7\textwidth}{X} 你們得救是本乎恩，也因著信；這並不是出於自己，而是神所賜的； \end{tabularx} \\ \\ \relax
2:9 & \begin{tabularx}{0.7\textwidth}{X} 也不是出於行為，免得有人自誇。 \end{tabularx} \\ \\ \relax
2:10 & \begin{tabularx}{0.7\textwidth}{X} 我們是他所造之物，在基督耶穌裡創造的，為要使我們行善，就是神早已預備好要我們做的。 \end{tabularx} \\ \\
[1ex]
\hline
\hline
\end{longtable}
$^{1}$在英文裡面有一個ING.
原來有很多用法.
我們查一查.
原來有這麼多種用法.
不過對我們來說.
最熟悉,最常用.
有什麼字會加ING呢?.
就是「現在進行式」.
加了ING就代表那件事正在發生.
正在進行.
例如.
「We are listening」.
就是「我們在聽的」.
「He is praying」.
「他正在祈禱」.
又或者我們現在打訊息.
譬如你在WhatsApp.
或者其他地方打訊息的時候.
已經很隨意了.
亦都可以有人會將ING.
加進去中文後面.
就代表那件事正在進行.
譬如.
「將要去面試了」.
「將要有重要時刻了」.
「很緊張了」.
「緊張ING」.
「緊張型」.
或者有人催你的時候.
「你走得了嗎?」.
「走緊了」.
「走」就加ING.
代表的意思就是.
「進行中就是這樣了」的時刻.
今天我預備的時候.
都在試玩一下.
就將題目加了ING.
就叫「活著ING」.
其實想一想.
「活著」這件事.

$^{41}$用在我們身上.
理論上都是「現在進行式」的.
是不是?.
我見到你.
我聽到你說話.
生勾勾凍在這裡.
喂 還不是活的?.
但是今天我們讀的.
《已忽所書》的經文.
作者保羅又真的對著一群.
生勾勾的人跟他們說.
「活」的相反.
就是「死」.
咦?.
說他們是曾經死了.
現在又活.
其實在說什麼呢?.
如果根據這段經文的說法.
究竟你和我今天是活的.
還是死的呢?.
基督徒怕不怕死的?.
怕 還是不怕?.
有人輕輕地說.
我會說其實基督徒都怕死的.
你問我 我都怕死的.
因為我都怕死會帶來的辛苦和分離.
基督徒都怕死的.
不過基督徒就不怕說「死」.
甚至我們說我們要好好地去說.
我們要認真地去想.
因為當我們知道「死」是什麼一回事的時候.
我們才懂得怎樣去活.
《已忽所書》是一本講教會的書.
所以是很適合我們一起去讀的.
讓我們知道在教會裡面聚集的是什麼人.
在這個世界.
不同的民族.
不同的語言.
都有教會的.
究竟這麼多間教會.

$^{81}$我們共同的基礎和經驗是什麼呢?.
教會在神的榮耀計劃裡面的角色和地位是什麼呢?.
我們作為教會的一份子.
我們日常生活行事為人.
又應該是怎樣的呢?.
《已忽所書》只有短短六章.
但是它的信息非常豐富.
上個月Andy教師就和我們打頭陣.
開始講第一章.
記不記得他的題目是什麼呢?.
啊!.
你回想起那套戲.
Andy教師的題目叫「神尋記」.
神去找人的一個記錄.
一個傳記.
Andy教師很用心去寫.
去想他的題目.
幫助你記得.
在那一章的這一段裡面.
保羅就去頌讚.
首先去頌讚三一神.
聖父,聖子,聖靈.
然後上個禮拜呂牧師.
又以第一章去為修信的人祈禱的那一段.
鼓勵我們要齊心努力.
認識基督.
接著今天就進入第二章.
可以說是保羅這封寫給教會的信的正文.
而今天讀的二章一至十節.
更加可以說是在聖經裡面.
講救恩講得最清晰的經文之一.
其中的第八,第九節.
更加經常是被選用.
作為保羅的本符因.
也因著信的救恩論的基礎.
在我的PowerPoint裡面.
我就將這一段經文二不所需二章一至十節.
分為三頁.
每一頁開頭都有兩個字.
讓我們留意一下.

$^{121}$這裡順著經文排.
這兩個字分別是「從前」,「然而」和「為要」.
「從前」是怎樣的呢?.
「然而」中間有一個轉捩點出現了.
「為要」之後變成怎樣呢?.
我們首先順著經文去讀一次.
二不所需二章一至十節.
「從前你們因著自己的過犯罪惡而死了.
那次你們再過犯罪惡終生了.
循從今世的宗族.
順服空中掌權者的領袖.
就似現在活躍的人心中運行的邪靈.
我們從前也都生活在他門當中.
放縱肉體的私慾.
隨著肉體和心中的意念去做.
和別人一樣.
生來就是該受罰的人.
然而神有豐富的憐憫.
因著他愛我們的大愛.
竟在我們因過犯而死了的時候.
使我們與基督一同活過來.
可見你們得救是本福恩.
他又使我們在基督耶穌裡.
與他一同復活.
一同坐在天上.
為要把他極豐富的恩典.
就使他在基督耶穌裡.
向我們所施的恩賜.
顯微及後來的世代.
你們得救是本福恩.
也因著信.
這並不是出於自己.
而是神所賜的.
也不是出於行為.
免得有人自誇.
我們是他所做之物.
在基督耶穌裡創造的.
為要使我們行事.
就是神早已預備好要我們做的.
從前.

$^{161}$如果我用一個英文字 就是「I」字.
「I」字是英文的「I」.
代表我的意思.
從前 我是怎樣的.
如果我們將剛才讀過第一至三節的經文.
稍為重做一下.
意思就是這樣.
從前 或者那個時候.
就是大家未信耶穌之前.
你們 就是在說已不所書的讀者.
就是那些不是猶太人的外邦信徒.
「外邦」是指猶太人以外的人.
的確在舊約聖經裡.
猶太人是有神獨特的揀選.
但是到了新約.
耶穌基督的救恩就不分種族了.
不再分猶太人還是外邦人.
總之就是給所有民族的人.
然後保羅說「我們」.
這裡就是包括保羅在內的猶太信徒.
保羅說「我們」「你們」.
未信耶穌之前其實都是一樣的.
沒有因為我們的出身而有什麼特權.
「你們」「我們」一樣都是怎樣呢?.
就是因著自己的過犯罪惡而死.
保羅形容那時候人的狀態.
就好像一個死的狀態一樣.
是被自己所犯的罪害死.
具體來說.
世界既然是神所創造的.
如果人拒絕神.
選擇一個與神割斷關係的生活.
那會是一個怎樣的生活呢?.
那時候行事為人的準則又會是什麼呢?.
經文接著就說.
這種生活主要由三方面的力量.
去主導人怎樣做決定.
第一就是隨從今世的風俗.
現今世界的風俗習慣.
它的價值觀,世界觀其實不是全部都是壞的.

$^{201}$都會有好的,壞的或者是一些相對的.
但是就是與神無關的.
是一個與神隔絕的社會價值系統.
即使有多少的進步,學問,經驗,財富.
也沒有辦法將人引導到世界的創造主那裡.
這些風俗,社會的風氣.
就好像空氣一樣包圍著我們.
在英文 The Message 版本的翻譯也很有趣.
意思就是說.
當人吸進去的肺那些.
全部都是那些污濁的不順.
我們呼出來的很自然就是.
活逆和不順從.
第二就是.
信服空中掌權者的領袖.
就是現今活逆的人心中運行的邪靈.
這個是邪惡的勢力,即是魔鬼撒旦.
聖經很清楚地告訴我們.
是有靈界敵黨神的邪惡勢力的.
祂在人之上.
但是祂不是與神同等的.
祂本質是與神對抗的.
也是統治這個罪惡世界的王.
祂專門就是要破壞神與人.
人與人的關係.
不過祂有多厲害都好.
其實祂仍然都在神的治理之下.
祂的行為同樣是受到神的認可或者控制.
其實魔鬼本身就已經好像一些咆哮的獅子.
聖經是這樣形容的.
祂就好像一隻很餓的獅子.
周圍行趣事.
就找那些祂想吃的人就吃了它.
既然這麼危險.
如果有人主動去接觸這些邪惡的靈.
其實就更加危險.
甚至是一個自殺式的行為.
第三方面就是放縱肉體的私慾.
隨著肉體和心中的意念去做.
各種的私慾.

$^{241}$就包括你想得到的各種肉體的情慾.
包括我們對食物,性,物質,權力等等的慾望.
當人很想去滿足這些慾望的時候.
就引發了很多的貪婪,爭奪,欺騙,放縱等等.
生來就是該受罰的人.
這樣的生活是糟蹋了神原本甚好的創造.
必然就是會自食其果.
死在這些罪惡過犯之中.
而且人的罪不單止是個人的事.
這個人很壞.
但是每個人都是這樣的時候.
人和人之間就會因此去互相碰撞,爭奪,彼此傷害.
如果沒有其他的因素出現.
懲罰,死亡就已經是必然的結局.
沒有希望了.
從前保羅就是在說我們猶太基督徒.
你們非猶太裔的基督徒.
以致今天的我們在未信耶穌之前.
就是活在這三方面的法則底下.
世界的風俗,魔鬼的搞擾.
自己的慾望.
這個也在解釋為什麼我們看新聞的時候.
無論是本地新聞,國際新聞.
有時候都會看到越來越灰.
或者另一種狀態就是.
變得越來越冷酷,冷酷無情.
當每天都是打仗的消息.
再看,跟昨天一樣.
都沒什麼感覺了.
人不是不知道這些問題.
或者不是不知道這些問題所引發的後果是多麼的壞.
但是人裡面的自私,貪婪,仇恨,放縱,承辱等等.
並沒有因為科技發達了.
財富增加了.
或者知識,學問的增長而消失.
反而更加好像耶穌要預言末世的時候.
這個世界的現象.
耶穌其實在二千年前就已經說.
到末世的時候會怎樣.
在《馬太福音》那裡說.

$^{281}$「民要攻打民,國要攻打國.
多處必有饑荒地震,這都是災難的起頭」.
那時候會有很多人跌倒.
也都會彼此陷害,彼此憎恨.
「且有好些假先知起來迷惑許多人.
因為不法的事增多.
許多人的愛心漸漸冷淡了」.
是不是都好像現在的寫照.
現在就算你在路邊見到有個人有需要.
你都要想清楚一點.
因為都不一定是真的.
這樣我們會不會變得越來越冷漠.
面對這個世界.
我們已經沒有那份心想去做呢.
這樣的情況如果沒有其他的因素出現.
懲罰,死亡就是必然的結局.
沒有希望了.
「然而」.
「NG」.
用「N」和「G」去代表神.
當「NG」拍到「I」的時候.
就是「然而」.
「神出手了」.
聖經裡面接著就說.
「然而」.
「神是怎樣的」.
「神有豐盛的憐憫」.
「祂有愛我們的大愛」.
「在我們因為過犯而死的時候」.
「使我們和基督一起活或者復活過來」.
「一起坐在天上面」.
「可見你們得救是本乎因也因著信」.
「這個並不是出於自己或者行為」.
「而是神所賜的」.
「免得有人自誇」.
「我們是祂所做之物」.
「在基督耶穌裡創造的」.
沒錯.
一開始神創造這個世界.
創造了人.

$^{321}$我們當然是祂做的.
但是這個人壞了和神斷絕了.
所以在基督耶穌裡面.
祂再一次給我們新的生命.
祂給我們有重新活著的機會.
是因著神的憐憫和愛.
我看到一個故事是這樣說的.
首先它說在英國有一個牧師.
有一天他在逛街.
逛街的時候他見到有一個男孩子.
手上拿著一隻鳥籠.
向著牧師一直走過來.
籠裡面有三隻小鳥.
好像很冷又很害怕地在顫顫顫.
牧師見到就叫著那個男孩子.
「啊!小孩你拿著什麼呀?」.
那個男孩子就說.
「哦!沒什麼」.
「我剛剛捉了幾隻老鳥而已」.
「那你打算怎樣對牠們呀?」.
「哦!我會拿牠們回家玩」.
「我會拔牠們的毛」.
「讓牠們互相去打對方」.
「這樣就有好戲看了」.
「哇!但是你這樣玩」.
「你遲早都會玩厭的」.
這個男孩子就說.
「哦!不要緊」.
「我家還有一隻貓」.
「牠最喜歡捉小鳥了」.
「我會將牠們放在一起」.
「就看看這隻貓怎樣捉牠了」.
這個牧師聽完之後就安靜了一會兒.
然後他就問.
「啊!我想買了你這幾隻小鳥」.
「你想要多少錢啊?」.
這個男孩子就說.
「咦!先生」.
「你不會真的想買這幾隻小鳥吧?」.
「牠們又不會唱歌」.

$^{361}$「牠們又不是長得特別好看」.
「如果你有錢」.
「你去街市那裡選吧」.
「大把小鳥漂亮很多」.
「可以讓你買啦」.
牧師就說.
「嗯!但是我真的很想買了牠們啊」.
「多少錢啊?」.
這個男孩子看著這個人.
覺得很傻.
然後隨口就說了.
「十英鎊吧」.
這個牧師就在他的袋子裡.
拿了十英鎊給這個男孩子.
拿了那個鳥籠回來.
這個男孩子就走了.
於是這個牧師就很小心翼翼地.
去到街頭那裡.
那裡有一棵樹.
有一個草地.
他就將那個鳥籠慢慢打開.
然後拍一拍牠.
去讓這幾隻小鳥飛出來.
重獲自由.
在這件事之後.
牧師就在一次的講台上.
講了一個故事.
他就說.
「嗯!有一天撒旦就和耶穌聊天了」.
「撒旦剛剛從伊甸園出來」.
他名字很難唸.
「行災樂禍自吹自擂」.
這樣就說.
「我剛剛捉了下面那個」.
「充滿人的世界」.
「我整了一個陷阱,一個圈套出來」.
「他們根本招架不住」.
「被我一網打盡」.
耶穌就問.
「咦!那你打算怎樣?」.

$^{401}$「哦!我要拿他們來玩」.
「叫他們喝酒,吸毒,淫亂」.
「我要教他們怎樣結婚又離婚」.
「彼此仇恨,虐待」.
「我要教他們去發明槍彈,武器」.
「核武器,彼此殘殺」.
「我要令到他們為了這個世界」.
「忙到妥妥齡,休息不到」.
「就算停下來」.
「也只是在看電視,看電腦,看手機」.
「我總是想到一些新的東西」.
「可以折磨他們的」.
耶穌就問他.
「嗯!那你之後打算怎樣處置他們?」.
撒旦就說.
「我要殺了他們」.
「讓他們在地獄裡永遠受折磨」.
嗯!耶穌又停了一會兒.
他就跟撒旦說.
「我想買了他們,開個價」.
撒旦就說.
「你不會真的想買他們吧?」.
「他們一無是處,背信棄義」.
「他們不懂得感恩的」.
「他們會否認你,向你吐口水」.
「咒詛你,最後他們還會殺了你」.
「你不會想買這些人的」.
耶穌就說.
「嗯!你要多少?開個價吧」.
撒旦對著耶穌說.
「好啊!要你的眼淚,你的血和你的生命」.
耶穌就說.
「好啊!成交吧」.
然後他就負上了自己生命的代價.
就是這樣.
耶穌負上了生命的代價.
去換取人可以重新做人的一個機會.
我們只有能夠明白這個恩典是多麼的重要.
神去為我們去負上了多大的代價.
我們才會明白.

$^{441}$這個重新做人的機會是有多寶貴.
耶穌負上這些代價.
為著人可以重新去活著.
祂期望之後我們的生活是怎樣的呢?.
當「I」加上「NG」.
「ING」就變成了一個現在進行式.
就是活著.
是怎樣的呢?.
接下來第三部分的經文.
就是講「為著」兩件事.
第一.
要將祂極豐富在耶穌基督裡面.
向我們所施的恩典,恩慈.
顯明給後來的世代.
第二.
就是為要使我們行善.
就是神早已預備好要我們做的.
能夠按著神的心意去行善.
這是一個與世界不一樣的法則.
這是連魔鬼都沒有我們乎的生命.
這是人可以在神面前謙卑下來.
去克服自我的一個狀態.
人就在我們身上能夠看得到.
是神極豐富的恩典.
想和大家分享一下.
最近我有一個經歷或體會.
可能有部分人知道.
我十二月的時候放了一個月的進修假.
其實那個月已經是我整個進修的尾聲.
過去三年我在公餘時間完成了.
道豐山基督教叢林一個關於靈修的課程.
有很多學習,很多恩典有機會再分享.
但是一月的時候就發生了一些事.
那些事令我有些困擾,有些掙扎.
不知道怎樣才可以好好地去拆解去應對這些事.
我就想到 嗯 神啊.
當我以為我已經完成了一個靈修的課程.
我交完功課的時候.
你竟然馬上給我一個新的功課.
是啊 上課,看書,寫文.

$^{481}$這些功課都不算最難做.
生命的功課才是最難.
啊 這個功課我要怎樣做啊.
上星期一我又去了道豐山.
是因為他們為校友安排了一個定期的保養.
就是每個月有一次的退休.
當日他邀請了方濟會的修士來分享.
之後修士分享完之後.
就每人送了一個禮物給我們.
就是這個了 我帶了來.
我打開之後.
它其實是一張明信片.
然後還有四個磁石的書籤.
裡面所印的東西其實是同一樣東西.
我看完這個就覺得眼前一亮.
我覺得是神要教我怎樣去做這份生命的功課.
它上面印著的是什麼呢?.
就是聖方濟 即是聖Francis的圖文.
當我說的時候可能你馬上就會想到.
哦 很熟悉啊這個.
它說是求神使他成為一個締造和平的人.
裡面是怎樣說的呢?.
它說 在有仇恨的地方讓我去播種仁愛.
在有殘害的地方讓我播種寬恕.
在有猜疑的地方讓我播種信任.
在絕望的地方讓我播種希望.
在有黑暗的地方讓我播種光明.
是的 信了耶穌.
我們就沒有馬上變成聖人或者解決所有問題.
無論是我現在眼前面對的這件事.
或者是其他的事.
其實我們今天每一個人.
我們生命裡面仍然都有不同的挑戰難處.
可能是來自世界的人和事.
可能是來自魔鬼的搞妖.
或者是我們裡面的自我仍然很大很實淨.
但恩典就是從此有耶穌基督和我們一起活著.
祂會指示我們那條正確的路.
祂會教我們怎樣去做生命的功課.
讓我們超越自己本身有的.

$^{521}$讓我們多過世界能夠給到我們的.
讓我們甚至有勝過魔鬼的力量.
神也會保守去保護那些願意謙卑下來.
信服他心意而行的人.
當神的兒女們能夠按著神的心意去生活.
去行善的時候.
就好像經文所說.
這些人能夠讓這個世代看到.
神的恩慈,憐憫是多麼豐富.
神的法則和這個世界不同.
神才可以帶來仁愛,寬恕,信任,希望和光明.
活著ING.
這個就是福音.
從前就是我們未信耶穌之前.
原來我們是這樣生活的.
我們是這樣去做決定,做選擇,這樣做人的.
然而耶穌用祂的寵駕.
給了我們一個重新做人的機會.
現在在基督裡面有新的指引,恩典.
在我們的生命裡面.
這個世界都看得出.
這個就是我們的得救見證.
今天選了一首回應詩叫做.
「活著為要敬拜你」.
我記得是一次團契的生日會.
有人選唱了這首歌.
我覺得挺好的.
想了一下.
對呀.
生日是代表我們活著的.
但基督徒的活著.
不只是肉身的生命.
不只是單憑時間過了.
就可以加的數字.
這裡講的活著是指我們屬靈生命.
是本福恩.
亦都恩著信.
而得以和耶穌基督一起去活著.
是一個我們的眼能夠看到神.
能夠從祂的眼光去看這個世界.

$^{561}$我們的耳朵能夠聽到神的話.
並且去進行.
這樣我們就能夠愛祂.
敬拜祂.
就為我們死.
神的愛就在此向我們顯明.
是出於你極豐富的恩典.
我們信心的回應.
你就給了我們重新做人的機會.
今天我們仍然要面對世界的風俗.
魔鬼的搞擾.
自我的慾望.
這些張力.
這些掙扎.
有時候我們彼此爭奪.
彼此傷害.
求你寬恕我們得罪你.
得罪人的地方.
亦都求你教導我們.
如何按著你的旨意去行善.
在世上活出一個愛你.
敬拜你的生命.
我們祈禱奉主耶穌名求.
阿們.
\newpage



\section{馬太福音 3:17}
\label{sec:031615}
\textbf{聖經人聖經事:這是我的愛子}~~~~黃克勤牧師博士
\newline
\newline
連結: \href{http://www.church.com.hk/acms/content.asp?site=cdc\&op=show\&type=product\&code=031615\&layout=sermon}{\texttt{031615}} ~~~~ 語音日期: 2026-01-29
\newline
\newline
\hyperref[sec:031614]{\small{< < < PREV SERMON < < <}}
~
\hyperref[sec:index_chronic]{\small{[返順時目]}}
~
\hyperref[sec:031616]{\small{> > > NEXT SERMON > > >}}
\newline
\newline
馬太福音 3:17
\newline
\begin{longtable}{cl}
\hline
\hline
章節 & 經文 (和合本修訂版)\\
\hline
3:17 & \begin{tabularx}{0.7\textwidth}{X} 這時，天上有聲音說：「這是我的愛子，我所喜愛的。」 \end{tabularx} \\ \\
[1ex]
\hline
\hline
\end{longtable}
$^{1}$大家好,又來讀聖經,新命樹書架網絡為您提供的聖經人聖經時.
今日講的是謝氏我的愛子,都是施洗約翰講的.
閱讀得多些關於施洗約翰,是越覺得這個人好厲害.
他亦都好有福,因為他能夠為主耶穌做先鋒.
見到主耶穌好多事情.
如果能夠見到主耶穌洗禮,由天上有聲音講,有聖靈降下禮.
主耶穌在水裡面上來,嘩,多厲害.
而且還是我為他施洗的,嘩,真是絕世的福氣.
施洗約翰就看到這些.
先看看思想,謝氏我的愛子.
馬太三章十三節,當下在那個時間,耶穌從加利利留到約旦河.
加利利在當日猶太恆省北部,加利利海那邊.
加利利北部.
他來到南邊近死海耶路撒冷,來到約旦河.
見約翰,然後要受他的洗.
很明顯,這一次主耶穌到南邊找施洗約翰幫他洗禮.
是特意的,從加利利來到約旦河.
來約旦河,不然他會去猶太省.
因為這個約旦河洗禮的地方,我都講過了.
如果你有機會去以色列.
傳說中,主耶穌洗禮的地方,嘩,曠野.
施洗約翰,曠野有人聲喊著說,預備他的路,收拾他的道.
曠野有人聲喊著說.
即是曠野來的,都說我們由耶路撒冷坐旅遊巴.
高速公路由耶路撒冷出發.
去到約旦河,施洗約翰洗禮,耶穌洗禮的地方.
都近三個小時.
嘩,你說走了很遠,他還由北部加利利走下來,沒兩天不行.
我想要睡兩晚在露天.
嘩,露宿,這樣真的不容易.
耶穌是特別下來找施洗約翰洗禮的.
所以他一來都是去約旦河,而不是去耶路撒冷.
因為多數是耶路撒冷聖殿,他是去約旦河.
是找約翰,然後你幫我洗禮,我想你幫我洗禮.
施洗約翰眼睛都大了,不是吧.
哎呀,不好意思,對不起.
又不好,不好,不行,為什麼.
我幫你提鞋也不配,你幫我洗禮是對的.
我怎麼敢幫你洗,我當受你的洗.
你返島上我借你來嗎,如果你在那裡洗禮,我馬上去找你.

$^{41}$你來我這裡,我不配,我不配,我怎麼可以幫你洗禮.
我當受你的洗是對的.
耶穌回答說,你暫且許我,你先讓我洗,你幫我洗.
因為我們理當這樣盡諸般的義.
我們,理當.
有一個字要留意.
他是說我們,因為我們理當盡諸般的義.
我們,其實他除了,這個諸般的義先解釋.
諸般的義是應該做的事情,好的事情.
其實義可以做諸般的禮.
諸般的禮,禮就是儀式.
其實用水施洗是一個儀式,代表我要悔改.
作為一個人,我是應該盡諸般的禮.
所以舉個例子,如果今天主耶穌在世上見到人.
祂會不會叫早晨呢,還是這樣.
等你叫我不是的,祂都會盡諸般的禮.
早晨,晚安,謝謝.
祂會盡諸般的禮.
走了,請請,多謝.
祂會做諸般的禮.
所以你洗禮,我作為一個人,我都應該來的.
而且盡是做到最好,可以做的都做了.
這個我們呢,是表示我是一個人.
而且我是和大家一樣的.
所以叫做我們.
我並沒有高高在上,我堂堂上帝來這個世界.
你們走在一邊,我幫你洗.
我只不過是來到,我暫時是一個人.
而且我和大家一起,所以我們.
我們還要合理,你當既然我是一個人.
我就應該做完所有要做的禮儀.
很好,主耶穌真的很謙卑.
施洗約翰這樣說,於是約翰就許了他.
既然如此,你不要怪我,我就勉強.
我勉為其難,其實你幫我洗是對的.
好了,耶穌受了洗,隨即從水裡上來.
當然是浸了,浸完之後就起來.
不會這樣浸著的.
浸完就起來.
哇,天忽然為他開了.

$^{81}$天為他開了是怎樣呢?.
我都未見過天開了.
不過我都識有人見過.
是怎樣的呢?.
我都無法描述,他也沒有描述.
是一個以前是宣教士的媽媽.
她在醫院善中靈陽院最後那幾天.
我去探望她.
她怕不怕呢?.
她媽媽說,這裡身體不太好,怕不怕呢?.
不怕,為甚麼呢?.
她說我見到很多天使,很多馬.
他們在那裡等我,預備來接我.
她見到,她見到天開了靈界的事情.
哇,一點都不怕.
究竟是怎樣的呢?.
她沒有說得清楚,不過我能夠明白.
就是在這個世界望天或望牆.
她見到一些在現場看到的東西.
背後好像開了一個洞.
天隨即為他開了.
他就看見神的靈彷彿甲子降下.
落在他身上.
他就看見神的靈彷彿甲子降下落在他身上.
他看見這個他應該是指耶穌.
主耶穌在洗完禮起床.
天隨即為他開了.
可能他起床之後,舉目望天.
天隨即為他開了.
有聖靈降臨在他身上.
彷彿甲子.
是不是甲子呢?.
每個畫圖畫都是畫甲子.
應該是似甲子.
說不出是甚麼.
如果是甲子就畫甲子.
所以是似甲子.
所以是彷彿甲子.
不要畫得太真實.
是似甲子落在主耶穌的身上.

$^{121}$他看見的.
另外,天上有聲音.
「這是我的愛子,我所喜悅的」.
這個是我所愛.
他來到這個世界.
是不是神不再愛他.
所以將他貶離凡間.
好像人家這樣說.
不是,是我的愛子.
我所喜悅.
這個固然是三一上帝同時的表現.
很重要.
我的愛子.
我所愛的.
是我愛的.
不是我不愛他.
所以給世人.
是我愛的.
所以我給他.
給人.
這是重要的.
這是施洗後很有福.
他看見,所以說出來.
今天跟大家分享到這裡.
金句就是.
從天上有聲音說.
「這是我的愛子,我所喜悅的」.
你知道上帝為我們做了這些.
他是父神.
猜批他的兒子.
來為我們工作.
而聖靈是在主耶穌的身上.
當主耶穌要離開的時候.
「我們怎麼辦?你走」.
「我不偏下你做孤兒」.
「所以我去,如是與你們有益」.
「我若不去聖靈,保衛師就不去你那裡」.
所以你見到.
在三一上帝.
在救贖上有差遣.

$^{161}$好像有高低之分.
不過其實他本身是三一上帝.
本質上是三一.
是一個神.
但是在差遣上.
他有差遣的次序.
今天跟大家分享到這裡.
「這是我的愛子」.
我們一起祈禱.
讚美上主.
你竟然差遣你的愛子來到這個世界.
在洗禮的時候.
你再次表達.
是天開了地表達.
你是多麼愛你的愛子.
但竟然你的愛子.
都肯將你的愛子為我們寫了.
多謝我們的主.
你就是這樣愛我們.
多謝天父預備的救恩.
多謝主耶穌的成就.
在這個十誡上多謝聖靈.
每一個時刻與我們同在.
承傳救恩.
讚美你.
願你的名被彰顯.
被高舉.
願你得著讚美.
祈禱奉耶穌名求.
阿們.
\newpage



\section{箴言 2:1-22}
\label{sec:031616}
\textbf{留心聽智慧,專心求聰明 (普)}~~~~周依娜牧師
\newline
\newline
連結: \href{http://www.church.com.hk/acms/content.asp?site=cdc\&op=show\&type=product\&code=031616\&layout=sermon}{\texttt{031616}} ~~~~ 語音日期: 2025-07-20
\newline
\newline
\hyperref[sec:031615]{\small{< < < PREV SERMON < < <}}
~
\hyperref[sec:index_chronic]{\small{[返順時目]}}
~
\hyperref[sec:031617]{\small{> > > NEXT SERMON > > >}}
\newline
\newline
箴言 2:1-22
\newline
\begin{longtable}{cl}
\hline
\hline
章節 & 經文 (和合本修訂版)\\
\hline
2:1 & \begin{tabularx}{0.7\textwidth}{X} 我兒啊，你若領受我的言語，珍藏我的命令， \end{tabularx} \\ \\ \relax
2:2 & \begin{tabularx}{0.7\textwidth}{X} 留心聽智慧，專心求聰明； \end{tabularx} \\ \\ \relax
2:3 & \begin{tabularx}{0.7\textwidth}{X} 你若呼求明理，揚聲求聰明， \end{tabularx} \\ \\ \relax
2:4 & \begin{tabularx}{0.7\textwidth}{X} 尋找她，如尋找銀子，搜尋她，如搜尋寶藏， \end{tabularx} \\ \\ \relax
2:5 & \begin{tabularx}{0.7\textwidth}{X} 你就懂得敬畏耶和華，得以認識神。 \end{tabularx} \\ \\ \relax
2:6 & \begin{tabularx}{0.7\textwidth}{X} 因為耶和華賞賜智慧，知識和聰明都由他口而出。 \end{tabularx} \\ \\ \relax
2:7 & \begin{tabularx}{0.7\textwidth}{X} 他為正直人珍藏健全的知識，給行為純正的人作盾牌， \end{tabularx} \\ \\ \relax
2:8 & \begin{tabularx}{0.7\textwidth}{X} 為要保護公正的路，庇護虔誠人的道。 \end{tabularx} \\ \\ \relax
2:9 & \begin{tabularx}{0.7\textwidth}{X} 那時，你就明白公義、公平、正直，和一切完善的道路。 \end{tabularx} \\ \\ \relax
2:10 & \begin{tabularx}{0.7\textwidth}{X} 因為智慧要進入你的心，知識要使你內心歡愉。 \end{tabularx} \\ \\ \relax
2:11 & \begin{tabularx}{0.7\textwidth}{X} 智謀要庇護你，聰明必保護你， \end{tabularx} \\ \\ \relax
2:12 & \begin{tabularx}{0.7\textwidth}{X} 救你脫離惡人的道，脫離言談乖謬的人。 \end{tabularx} \\ \\ \relax
2:13 & \begin{tabularx}{0.7\textwidth}{X} 他們離棄正直的路，行走黑暗的道， \end{tabularx} \\ \\ \relax
2:14 & \begin{tabularx}{0.7\textwidth}{X} 喜歡作惡，喜愛惡人的錯謬。 \end{tabularx} \\ \\ \relax
2:15 & \begin{tabularx}{0.7\textwidth}{X} 他們的路歪曲，他們偏離中道。 \end{tabularx} \\ \\ \relax
2:16 & \begin{tabularx}{0.7\textwidth}{X} 智慧要救你遠離陌生女子，遠離那油嘴滑舌的外邦女子。 \end{tabularx} \\ \\ \relax
2:17 & \begin{tabularx}{0.7\textwidth}{X} 她離棄年輕時的配偶，忘了自己神聖的盟約。 \end{tabularx} \\ \\ \relax
2:18 & \begin{tabularx}{0.7\textwidth}{X} 她的家陷入死亡，她的路偏向陰魂。 \end{tabularx} \\ \\ \relax
2:19 & \begin{tabularx}{0.7\textwidth}{X} 凡到她那裡去的，不得回轉，也得不到生命的路。 \end{tabularx} \\ \\ \relax
2:20 & \begin{tabularx}{0.7\textwidth}{X} 智慧使你行善人的道，守義人的路。 \end{tabularx} \\ \\ \relax
2:21 & \begin{tabularx}{0.7\textwidth}{X} 正直人必在地上居住，完全人必在其上存留； \end{tabularx} \\ \\ \relax
2:22 & \begin{tabularx}{0.7\textwidth}{X} 惟惡人要從地上剪除，奸詐人要被拔出。 \end{tabularx} \\ \\
[1ex]
\hline
\hline
\end{longtable}
$^{1}$此錄音或錄影的版權屬基督教宣道會,北教堂所有.
我們的屬靈傳統之一,是以真理教導為本.
歡迎你將此錄音或錄影的連結傳給親友.
懇請注意,未經授權,請勿轉錄或剪輯本錄音或錄影.
聖願就像十讀行傳二十章三十二節所說.
這道能夠傳遞給我們所有的信徒.
這道能建立你們,叫你們和一切成聖的人同德及義.
弟兄姊妹主日平安.
在這颱風的主日,冤主賜給你們有平安.
在我們領受今天的道儀前,我們有一個禱告.
我們一起來禱告.
滿有智慧的神,我們在今天.
來親近你,我們打開聖經.
尋求你的智慧,願主你今天引導我們.
來留心聽,專心求,使我們的智慧.
在你的保守之下,每一天都有增長.
以致你可以保守我們,走一條人生正直的路.
禱告奉耶穌基督的名求,阿們.
弟兄姊妹,這個主日我們繼續的來.
看看真言的一些教導.
真言是來自父母親對愛兒的一個教導.
是一種家庭的教育,我不知道在你的家庭裡面.
父母親對你有過怎麼樣的教導呢?我自己的父母對我的教導.
就是人不要太短視,要有遠見,要有目標.
為將來做好預備,當他們信主之後,他們對我的.
教導是,我們家得來的一切,我們所擁有的一切.
其實都是上帝所賜的,以前我們會以為.
這一些都是靠我們自己的力量來賺取得來的.
但是現在我們知道這一切都是神賞賜.
神保守的,父母給我的教導讓我受用一生.
他們因為愛我,給我這些的教導,希望我有一個幸福的未來.
而我發現原來這些的教導慢慢慢慢的在塑造.
我的品格,猶太人跟中國人一樣非常.
重視這個家庭教育的,除了是塑造.
他們的品格以外,還有最大的目的呢,他們要來學習.
妥拉神的教會,是要來敬畏神.
在中世紀的猶太的社會呢,小孩子開始學十字的時候.
他們會舉行一個隆重的儀式的,流傳到今天.
他們也會有這樣的一個儀式,你在PPT裡面就看到.
那些的小孩跟大人會戴上象徵很敬虔的.

$^{41}$一個長方形的披肩叫Tsitsit,由父親帶著小孩走向老師.
然後老師會引領男孩坐在他的大腿上,展示一塊石板給他看.
寫著一些上帝的妥拉,一些的律法,然後上面還會寫著一行字.
就是希伯來文的「但願妥拉成為你的終生執志」.
那老師宣讀每一個字,小孩就跟著念.
然後石板上就會沾滿一些的蜂蜜,那小孩就去舔它.
就代表耶和華的律法就好像蜜糖般一樣的甘甜.
那到今天猶太人的母親都會把蜂蜜滴在書本上.
讓孩子來舔一下的,目的就是讓他從小就知道.
書是甜的,神的律法也是甜的.
就建立起他們對書,對上帝律法的一個喜愛之情.
這一篇的《箴言二章》的猶太人的父親.
他的家庭教育也是教導孩子要正確地來對待智慧的態度.
就是讓神的智慧塑造他成為一個正直,純良的人.
好走在一條正道之上.
我們今天也來聽聽這位蜀靈父親有關智慧的教導.
相信也會對我們受用一生.
首先我想帶大家看看《箴言二章》很美麗的一個結構.
它是模仿希伯來文22個字母寫成的詩歌.
就解釋了為什麼會有22節.
它的字母未必每一個連著來.
所以它是叫仿字母詩.
這篇的仿字母詩它是一口氣的講完的.
就是不斷氣的一口氣的講完.
所以就出現了一個很有趣的一個結構.
你看就是如果你一到四節對待智慧有一個正確的態度.
你就會明白五到九節有兩件事.
那你就會獲得十到十九節兩個好處.
最後就有一個結論.
二十到二十二節智慧人行在生命之道上.
從結構我們看見《箴言二章》.
它有一個很巧妙的巧合.
就是數字二.
《箴言二章》.
然後會明白兩件事兩個好處.
然後有二十二節.
我覺得這個是很巧妙的一個安排.
在北方的方言裡面.
有一句就是你真二.
你真二就是你表達這個人有點傻氣和冒失.

$^{81}$希望我們對待智慧不會是二.
而是讓這個數字二.
說明要你記得今天的一個信息.
做一個智慧人.
《箴言二章》的主題就是讓智慧塑造你成為一個正直.
純良的人.
智慧就會保護你行在正道之上.
今天我會跟大家來分享三點.
第一點就是珍藏智慧.
竭力尋求.
珍藏智慧.
竭力尋求.
智者的父親他要愛兒得智慧的第一步.
就是對待智慧的正確的態度.
從第一到第四節你就會看到有很多活潑的動詞.
請受珍藏.
留心聽專心求.
呼求揚聲求.
然後第四節就會出現有四次有關尋找的詞語.
尋找尋找搜尋搜尋.
猶太人的父親他們是擔起教導孩子特別是男孩.
妥拉的一個責任.
所以我的言語.
我的命令.
法源呢都是來自律法的誡命.
然後父親說我兒啊你若領受receive.
我的言語.
然後第六節.
耶和華賞賜.
give.
這裡都說明了智慧有一個源頭和有一個範圍的.
那這個智慧呢.
只有一個源頭.
就是從生而來.
而真智慧呢也只有一個範圍.
你看因為耶和華賞賜智慧.
知識和聰明都由他口而出.
這個就是智慧的範圍.
你看wisdom.
knowledge and understanding.

$^{121}$這三個智慧知識聰明呢.
唯獨從耶和華而來.
就是wisdom knowledge和understanding其實都是有關於.
神對神的一個認識.
對神的律法的那個認識.
我們只是一個領受者.
receiver我們只是一個領受者.
領受者的意思就是我們是.
去聽神的教誨的.
那最近呢香港呢有一位才子他去世了.
香港四大才子都離開了.
他們的作品都留存下來.
其中他們的小說.
一些的作品.
甚至人生觀都讓人的津津樂道特別是年輕人.
很愛去看他們的言論.
那我就想如果世上的智慧都在日常生活裡面.
去影響著我們做人的道理.
那何況神的作品.
聖經.
不是更值得我們去尋求嗎.
弟兄姊妹我不知道你現在面臨人生哪一些的不安.
挫折.
或者是難題.
你習慣去尋求的是什麼人.
或者是什麼渠道.
你尋求的是哪一些的來源的智慧.
是不是好像聖經說世上的小學.
言談.
甚至是虛談.
還是你會尋求神的一個智慧和幫助呢.
你有沒有經歷過.
當我們自己的智慧也無法可施的時候.
你的感受.
和你的行動.
你會怎麼做呢.
真言教導我們要去神那裡尋求智慧.
我很喜歡雅各書一章五節.
你們中間若有缺少智慧的應當求那厚次於眾人.
也.

$^{161}$不斥責人的神.
主就必賜給他.
真智慧的範圍.
只有在神那裡.
有了智慧的範圍.
還有一步.
就看你.
怎麼去對待智慧了.
這個父親說.
你若.
如果你領受我的言語珍藏我的命令.
就會.
懂得上帝的智慧嗎.
所以若其實是一個選擇.
在真言裡面經常有跟生命計一樣有兩條路的一個選擇.
是走一條生命之道.
之路還是.
死亡之路.
你要去選.
所以在這裡.
父親要兒子去選.
你要不要是選擇認真的來對待神的話.
還是隨便來對待神的話.
真言它的格式是上下兩句領受珍藏.
留心聽專心求.
其實這個結構是有一個更進一步更進深的.
你領受了.
然後不是.
不是單單領受了.
而是要去珍藏.
在這裡父親要愛兒不只是很被動被動的領受.
的一些的資訊然後左耳進右耳就出或者是當耳邊風了.
如果不經過珍藏這個階段的話呢.
你就會把神的話給忘掉了.
我們要珍藏來.
承繼神的話.
來珍藏神的話.
原來珍藏它這個意思珍藏.
它原來的意思是high.
treasure.

$^{201}$就是隱藏的珍寶.
而如果你去看第四節的時候呢.
這個寶藏.
也是hidden treasure的意思就是隱藏的珍寶.
神的話神的智慧.
是向願意去尋求他智慧的人.
來慢慢的來打開.
隱藏.
珍藏.
就是.
神的話神的智慧.
就是我們一生的珍寶.
猶太人.
他們大約從五歲的開始呢.
他們就會被送到.
一些的學校裡面.
去聽老師的一些的訓誨.
還有老師的一些的教導然後兩個人之間呢.
還會有一些的對話這樣來學習的.
那之前猶太人的會堂呢.
會堂一般是給成年的男子.
去學習.
律法的地方.
但是這些的會堂呢還是會附設一些的學校.
讓五歲開始的這些的小朋友呢.
來學習上帝的律法.
那特別是摩西五經或者是申命記的.
那老師會透過口傳.
誦讀.
和背誦.
讓那個小朋友呢.
來記住上帝的律法.
所以在猶太的舊約裡面不論是在家庭或者是在這些學校.
裡面呢都會非常強調.
聽和記憶的.
經文的誦讀呢是非常重要的一部分哈我不知道.
弟兄姊妹你是怎麼存記.
或者是珍藏.
神的珍寶就是神的話呢.
我曾經分享過一段的靈修的一些的技巧.

$^{241}$那其中三章的PPT呢可以跟大家分享.
當你讀聖經的時候讀神的話的時候第一步呢.
你可以速讀.
很快的把你的靈修經文.
或者是崇拜的經文.
很快的來讀一次.
然後呢讀完之後呢你停留一下.
然後反覆去回顧一下這個內容是什麼意思呢.
就讓經文呢去刻畫在你的心.
稍留片刻以後呢你再去讀經文.
再讀.
再讀的時候呢就不是速讀是慢讀.
以祈禱的心去讀.
讓神的話慢慢進一步來滋潤你的心田.
然後呢.
就存想經文.
你嘗試.
存記一下.
有一些的經文.
是很深刻的.
你就最好呢.
就是你一些主要的一些的經節關鍵的一些的.
經節.
讓你呢可以在一天裡面呢就反覆的來背誦反覆的來思想.
晝夜思想嘛.
當你發現其中有一句的經文有一句的經文你特別感到是神對你說.
話你特別有感動的.
這個時候呢你就可以有一句的禱告.
你說主啊願你的話進入我的心.
我需要你.
來觸摸我的生命.
我的心.
在聆聽.
然後呢你就可以把這句的經文呢.
給背起來.
我受洗的時候呢其實收到一盒很樸素的一個經句卡.
它是白紙黑字的.
每天一句.
我那時候呢.
每天吃一張.

$^{281}$就是吃聖經.
吃神的話.
我就把那一天的經文我就背起來.
後來我就慢慢發現.
原來我這些已經背好的經文.
在我不同的場景.
在我有需要有危險有試探有軟弱的時候呢.
聖靈就會讓這些的經文很快的浮現在我的腦海.
成為我的幫助.
我也想起探望一對的夫婦.
在晚餐的時候呢.
當我想做飯前禱告的時候.
他們就說等一等.
我說.
為什麼要等呢我們現在不是要飯前禱告了嗎.
原來他們的家庭也有一個家規有一個家庭的教育.
他們在飯前呢先會在經句卡中他們預備了一盒經句卡在飯桌上.
然後呢每人抽一張的經句卡.
丈夫讀.
妻子讀然後小孩讀.
我發現這個是很有意思的.
他們在吃飯以前也讓神的話成為他們的食物.
珍藏在心裡.
以後每當我去退休的時候呢.
我也會帶一個小袋子.
在我們退休以後.
挖掘神的話語之後.
就讓弟兄姊妹呢抽一張的經句卡.
有時候弟兄姊妹呢會跟我分享說.
哎呀那句的經句呢剛剛正好是神.
給我去處理當前難題的一個智慧呢.
那今天其實我們的手機.
也會有每天經句這個APP的.
所以我也鼓勵你.
可以善用這些的每天經句每日一句.
這些的APP呢鼓勵你去珍藏神的話.
我們再細心的來看一到四節的時候呢你會看到.
這個智慧的那個行動升級了.
不單單只是珍藏而已.
我們經常說聽道而行道嘛.

$^{321}$心動還有行動.
所以現在這個父親.
他所描繪的是很多關於.
身體語言是動詞的.
你看留心聽就是豎起你的耳朵.
專心求.
心其實是猶太人整個人的意思.
靜心靜意靜力的愛神這個心就是你整個人的投入.
還有口呼求明理.
揚聲求.
然後你的腳了.
就去尋找尋找搜尋搜尋.
都是一種呢全人全心竭力去尋找智慧的態度.
第二張的智慧呢跟第一張.
是不同形式出現.
第一張呢智慧呢在街外.
向人公開的來呼喚.
但是在這裡呢是以搜尋隱藏的珍寶的這個姿態.
來出現的.
你看這個本來的意思是hide.
就是隱藏的.
然後第四節是尋找就是seek.
那我們都玩過一個遊戲就是hide and seek.
就是捉迷藏.
我覺得呢好像是一種.
神給人的一種的樂趣.
他說.
尋找影子.
尋找智慧好像尋找影子.
那你看到右邊那幅圖呢就是淘金嘛.
大家知不知道在前段時間有澳洲的淘金熱.
就有很多的淘金團.
就是有導遊帶著他們去淘金的.
有一些的參加者呢他就分享說.
其實要付出很多的代價.
要走好多的路.
要尋找好久.
才可以找到一粒金.
當他們找到一粒金的時候呢.
他們就非常的興奮.

$^{361}$他們說非常享受那種尋找的樂趣.
就算很多時候都好像是要花很多的時間很多的代價.
但是他們覺得都是很有價值的.
那尋找神的智慧.
神的話其實也是一樣.
我們需要花很多的代價.
很多的時間.
我們就可以.
神的智慧被我們去尋見的.
我不知道你.
去最大的搜尋樂趣是什麼呢.
是一些的app購物的一些的平台.
還是.
搜尋什麼是你的最大的樂趣呢.
我最近搜尋的一個經歷呢就是陪小朋友呢去抽一些的盲盒.
原來大人呢都很愛抽盲盒.
每個人帶呢都是會掛一支的.
我也尋得一支香薰盒.
因為我自覺自己付不出大量的金錢.
去抽新頭號.
香薰.
無論抽中哪一支起碼實用一點嘛.
那我覺得抽盲盒的設計.
吸引人的地方呢就是把你最心愛的公仔.
隱藏起來.
激發你一定要尋到寶的決心.
你就會享受過程當中的期待感.
還有打開.
你那一下的哇那種的驚喜感.
我覺得神的心思也好像盲盒一樣.
智慧.
好像.
像你.
暫時的隱藏起來.
其實就是吸引你.
去尋寶.
而神他是怎麼對待竭力去尋找他智慧的人呢.
我們一起來讀出第七節.
他為正直人珍藏健全的知識.
給行為純正的人做盾牌.

$^{401}$你看見嗎神真幽默.
當我們珍藏神的話的時候.
神也怎麼樣.
他也.
為我們珍藏.
就是一份的禮物.
這份的禮物是什麼呢.
就是盾牌.
這個盾牌.
就是讓智慧.
成為.
你的一個盾牌.
所以第二點就是智慧.
盾牌.
護衛.
正義.
神為所有竭力尋找他智慧的人.
預備了一份珍貴的禮物.
盾牌.
這個盾牌其實主要是用在軍事的場地裡面.
來抵抗一些的仇敵的.
所以你看七到十一節用的字眼都是軍事性的.
保護.
庇護.
庇護保護.
God protect.
Protect God.
你看神.
專門是為那些什麼人.
會用智慧來保護他們呢.
是正直人.
純正的人.
虔誠人.
換句話說.
智慧就會向那些愚昧的人.
繼續的隱藏.
因為他們根本沒有去搜尋.
上帝的智慧.
所以是一直隱藏起來.
但是為那些.

$^{441}$積極.
去尋找.
去搜尋的那些.
虔誠人,正直人和純正的人呢.
上帝就會用他的智慧來保護他們.
而最重要的就是上帝的這種的護衛和保衛呢.
是有一種的目的的.
就是讓這些虔誠人,正直人,純良的人.
去懂得.
明白.
也就是understand的意思.
是.
去明白,敬畏耶和華.
認識神.
認識神是一種跟神.
有真實交往的一種的經歷.
然後你也會明白公義,公平和正直.
和一切完善的道路.
這裡的認識.
代表是指認識神的屬性.
你看公義,公平,正直.
就是神的屬性.
很多時候我們用世俗的標準來定義什麼叫公義,什麼叫公平.
和正直.
但是因為.
每個人的觀點角度都不一樣.
以至於在政治議題上,社會議題上呢.
都會有不同的標準.
或者有不同的一種的.
界限.
但是神的公義,公平和正直.
就是最終絕對的一個真理.
因為智慧,知識和聰明都是從他而來的.
所以他的公義,公平,正直才是最終的標準.
而公義,公平,正直這個屬性.
其實也是上帝.
要來塑造,要讓神的兒女.
來塑造出同樣的這種三方面的屬性.
三方面的品格.
如果用一句話來概括公義,公平,正直.

$^{481}$這方面的品格.
其實就是上帝.
希望我們.
活出神眼中正確的生活方式.
神眼中正確的生活方式.
是按照他的律法來處世待人.
所以正確的生活方式.
公義,公平,正直.
就是按照他的律法來處世待人.
這樣你就可以走上一條完善的道路.
完善的道路.
就是good path.
good path.
完善就是好的.
好的意思就是和神心意的一個道路.
怎麼樣才是和神心意的道路.
怎麼樣才是危險的道路呢.
在最後一段12到19節.
就特別提出兩個例子.
這兩個例子就回應.
7到12節.
兩種人的邪道.
智慧要救正直,純良的人去脫離.
這兩種的邪道.
所以最後一點就是智慧訴咎.
守護界限.
兩個陷阱,第一個陷阱是惡人惡道.
惡人惡道就是涵蓋了社會上面.
有很多的一些的惡人,惡人的意思就是陷在罪惡裡面的人.
他們不明白上帝的智慧.
他們也不聽上帝的話語.
一個不聽上帝話語的人.
他的言行.
就會出現了歪曲,你看脫離.
言談.
乖妙的人.
要脫離這種的人.
他們行走黑暗的道.
你看言談.
行走.

$^{521}$言談行走,在真言兩章.
都是在說我們的全人嗎?我們的眼睛,我們的口,我們的心,對嗎?.
在這裡他們是言行不一,言行歪曲的人.
這種的人的道路你又避免.
你要逃離,你要脫離的.
第二個陷阱呢,就是外女的迷惑.
外女.
這裡的經文是陌生的女子.
陌生的女子.
其實呢如果你看.
他這個經文裡面呢.
他是一種.
要試探你的.
一種要去.
勾引你進入他的一個陷阱當中的一種的淫亂的一種的誘惑.
這裡做父親的教誨呢.
是要愛兒要脫離.
這些的陷阱.
你看這個救字呢.
這個救並不是普通的一個解救.
而是特別指在一個極端危險.
和邪惡的環境當中呢.
非常快速的把你拉出來.
救出來的.
而這個外女的這個危險.
在今天在這個道德感越來越低的社會當中呢.
也會經常的出現的.
特別在現在手機非常的方便.
有更多的淫亂誘惑.
在不同的平台要拉你進入不同的試探裡面.
我們可以在電視裡面.
看到呢有很多的男子不論是年輕人或者是老年人呢.
都會被女子呢來色誘來騙財的.
他們都很容易被情慾所誘惑而犯罪的.
在最近的一個新聞裡面我們就看到一對偷情的男女.
他們去看演唱會.
但是卻被演唱會的大屏幕呢.
就公開了他們的關係.
然後他們的偷情的關係呢.
就展示在世人的底下.

$^{561}$所以智慧的救.
他強調的是一種的遠離.
遠離再遠離.
在16節你看這個外女.
外女就是外邦的女子陌生的女子.
你看他的用的一種的試探呢.
是油嘴滑舌.
油嘴滑舌其實是一種言語的一種的挑逗.
目的是要刺激人感官的一種的情慾.
經文說他忘了盟約離棄配偶.
看出.
這個女人.
她是一個離婚的一個女人.
然後她是一個不負責任的.
對關係是沒有忠誠偽身.
只是追隨心中慾望的一些的人.
面對這些的外女.
智慧呼籲我們要趕快的逃離.
趕快的速離.
趕快逃趕快逃.
這個就是最有效的一個方法.
當你遇見視覺.
當你遇見言語.
一些的挑逗誘惑的時候.
最好的方法第一個方法就是.
就是速逃.
然後智慧就會來救你.
速救所以速救速逃.
大家要記得.
第十節他說因為要智慧進入你的心知識要使你內心歡愉.
內心歡愉呢.
他的本來的意思是吸引和愉悅.
這個父親讓他的愛兒要知道.
真真正正吸引你的.
不是這些的外女.
的一些的誘惑.
而是智慧.
智慧進入你的心.
你會感受到真正的喜樂.
真正的平安.

$^{601}$真正的愉悅和安息.
但是如果你讓所有一些的情慾.
淫亂的一些的事態進入你心的時候.
你最後的結果就會跟著他進入這個死亡之路.
十一節這個智慧以智謀和聰明來表達.
智謀就是讓你能夠明辨是非.
洞悉這些試探背後的一個動機.
就是要欺騙你.
不論是你的財物.
不論是你的感情.
不論是你的生命.
他全部都拿走.
你要洞悉這個動機.
如果你有這種的智謀.
這個盾牌.
拿起這個盾牌.
他會保護你.
聰明就是你明白真理.
明白上帝的知識.
關於上帝的真理.
這樣如果你有智謀和聰明的時候.
這個盾牌就發揮效力.
他一定會保護你的.
家庭教育.
是非常重要的.
在我們北宣教會裡面經常會有一些夫婦的講座.
親子的講座.
我鼓勵大家來參加.
在有一次的.
保家運動裡面呢.
我最記得那個講員.
他講了一句話.
他說弟兄姊妹.
特別是弟兄.
我們不要太信任自己的情感.
人是軟弱和經不起重複的誘惑.
所以唯一而有效的方法.
就是遠離.
遠離.
我們都記得約瑟.

$^{641}$他經歷波提法的妻子.
一而再再而三的一個勸說.
要跟他發生親密關係.
約瑟他的舉動是怎麼樣.
也是訴討.
保全了他在上帝面前的正直純良和虔誠.
所以.
在這裡.
我們要記得.
在教會裡面.
在家庭裡面.
我們要.
去.
行善人的道.
但是這個守.
也是很重要的.
守住他.
守住.
我們家庭的堡壘.
用上帝的話上帝的智慧.
來堅守.
自己的一個界限.
我自己在教會裡面也做很多的一些的.
教會的一些的紀律.
讓我很難過的是.
有一些的.
試探.
也是關於.
婚外情婚外性的一些的情況.
是最後呢.
也是勝不過的.
我繼續的為他們來禱告希望他們能夠.
來遵從智慧的一個呼求智慧的一個庇護.
以至於他們.
回到家庭來.
回到妻子的身邊.
回到孩子的身邊有一個完整的家.
當我們知道.
主耶穌基督是我們一家之主的時候我們稱他為主的時候.
我們就要保護我們的家庭.

$^{681}$守住.
我們的一個界限.
所以這個守是很重要你行了.
但是還是要守住.
在真言我們非常.
熟悉的就是.
那個經文就是你要保守你心勝過保守一切因為一生的果孝是由心.
發出.
但是呢在下面那一段呢我們經常就會忘記了我們一起來讀.
你要除掉邪鄙的口.
棄絕乖妙的嘴.
你的眼目要向前正看.
你的眼睛當向前直觀.
要修平你腳下的路.
堅定你一切的道.
你看.
這個也是有關於口.
眼.
腳.
就是我們在第一點所講的.
整個人全心.
全力.
的要來守住.
自己.
的一個純良正直的生命和品格.
這樣呢我們就會有一個真言說的一個結論.
這個結論就是.
在21 22節說.
正直人必在地上居住.
完全人必在其上存留.
惟惡人要從地上剪除.
奸詐人要被拔出.
今天.
你也要來選擇一條道路是一個智慧庇護你的一個生命之道.
還是你要走在.
各樣.
誘惑陷阱.
誘惑你走在死亡彎曲之路呢.
你要選擇.
那我們來看今天這三點很簡單的一個解說.

$^{721}$就是第一.
珍藏智慧.
竭力尋求.
智慧盾牌.
護衛正義.
智慧訴求.
守護界限.
讓我們有一個禱告.
智慧的神.
求主.
你來.
教導我們.
你的話.
是隱藏的珍寶.
是讓我們一生來珍藏.
在我們的心中.
主啊我們實在需要你.
觸摸我們的生命.
因為我們.
需要你在我們的生命當中.
主啊求主.
你來保守我們.
成為一個.
有.
公義公正.
正直.
品格的人.
讓智慧來保守我們一生.
一直走在一條正直的路上.
求主來護衛我們.
也幫助我們守住家庭的界限.
以至於我們.
可以.
來保護我們的家庭.
我們尊主.
你為我們一家之主.
使我們對孩子都有一個.
最好的.
循著你的真理的一個家庭教育.
求主.

$^{761}$來幫助我們保守我們.
禱告奉耶穌基督的名求.
阿們.
\newpage



\section{出埃及記 33:12-23}
\label{sec:031617}
\textbf{渴求榮耀的禱告}~~~~張志怡傳道
\newline
\newline
連結: \href{http://www.church.com.hk/acms/content.asp?site=cdc\&op=show\&type=product\&code=031617\&layout=sermon}{\texttt{031617}} ~~~~ 語音日期: 2026-02-08
\newline
\newline
\hyperref[sec:031616]{\small{< < < PREV SERMON < < <}}
~
\hyperref[sec:index_chronic]{\small{[返順時目]}}
~
\hyperref[sec:031618]{\small{> > > NEXT SERMON > > >}}
\newline
\newline
出埃及記 33:12-23
\newline
\begin{longtable}{cl}
\hline
\hline
章節 & 經文 (和合本修訂版)\\
\hline
33:12 & \begin{tabularx}{0.7\textwidth}{X} 摩西對耶和華說：「看，你曾對我說：『將這百姓領上去』；卻沒有讓我知道你要差派誰與我同去。你還說：『我按你的名認識你，你也在我眼前蒙了恩。』 \end{tabularx} \\ \\ \relax
33:13 & \begin{tabularx}{0.7\textwidth}{X} 我如今若在你眼前蒙恩，求你將你的道指示我，使我可以認識你，並在你眼前蒙恩。求你顧念這國是你的子民。」 \end{tabularx} \\ \\ \relax
33:14 & \begin{tabularx}{0.7\textwidth}{X} 耶和華說：「我必親自去，讓你安心。」 \end{tabularx} \\ \\ \relax
33:15 & \begin{tabularx}{0.7\textwidth}{X} 摩西說：「你若不親自去，就不要把我們從這裡領上去。 \end{tabularx} \\ \\ \relax
33:16 & \begin{tabularx}{0.7\textwidth}{X} 現在，人如何得知我和你的百姓在你眼前蒙恩呢？豈不是因為你與我們同去，使我和你的百姓與地面上的萬民有分別嗎？」 \end{tabularx} \\ \\ \relax
33:17 & \begin{tabularx}{0.7\textwidth}{X} 耶和華對摩西說：「你所說的這件事，我也會去做，因為你在我眼前蒙了恩，並且我按你的名認識你。」 \end{tabularx} \\ \\ \relax
33:18 & \begin{tabularx}{0.7\textwidth}{X} 摩西說：「求你顯出你的榮耀給我看。」 \end{tabularx} \\ \\ \relax
33:19 & \begin{tabularx}{0.7\textwidth}{X} 耶和華說：「我要顯示我一切的美善，在你面前經過，並要在你面前宣告耶和華的名。我要恩待誰就恩待誰，要憐憫誰就憐憫誰。」 \end{tabularx} \\ \\ \relax
33:20 & \begin{tabularx}{0.7\textwidth}{X} 他又說：「只是你不能看見我的面，因為沒有人看見我還可以存活。」 \end{tabularx} \\ \\ \relax
33:21 & \begin{tabularx}{0.7\textwidth}{X} 耶和華說：「看哪，靠近我這裡有個地方，你可以站在這磐石上。 \end{tabularx} \\ \\ \relax
33:22 & \begin{tabularx}{0.7\textwidth}{X} 當我的榮耀經過的時候，我必將你放在磐石縫裡，用我的手掌遮掩你，等我過去， \end{tabularx} \\ \\ \relax
33:23 & \begin{tabularx}{0.7\textwidth}{X} 然後我要將我的手掌收回，你就可以看見我的背，卻看不到我的面。」 \end{tabularx} \\ \\
[1ex]
\hline
\hline
\end{longtable}
$^{1}$靚姊妹平安.
多謝主席剛才說我們讀出這段經文.
一間我講到或者在PowerPoint裡面.
呈現的經文都是黃修版的聖經.
請大家留意.
踏入今天是2月8號.
也是教會四十天運動剛剛完成了一個星期.
不知道在座有多少個靚姊妹完成了第一個星期的.
可不可以舉一下手.
大家很厲害.
希望有更多人可以繼續.
我們還有三十多天.
大家可以繼續來為教會.
為不同的需要來祈禱.
今天我講到不是想hard sell這個祈禱運動.
不過很想和大家一起.
來反思一下我們祈禱的生活.
除了我們說祈禱祈多久之外.
其實有兩樣東西都很想大家去想一下.
我們祈禱的內容究竟是什麼呢?.
並且其實我們真正祈禱的時候.
我們求的是什麼呢?.
2023年4月在美國有一個基督教的組織.
叫做天際線SkyKnight.
它向1732個美國成年的基督徒.
做了一個信仰的調查.
在裡面發現了一些有趣.
也是值得我們關注的結果.
第一點都是時間的關係.
調查發現美國成年的基督徒平均每天.
他們祈禱的時間是18分鐘.
我們祈禱運動想大家祈禱多久呢?.
15分鐘這個比較長.
平均他們都18分鐘.
所以我相信在座弟兄姊妹會祈15分鐘.
應該都沒有什麼大問題的.
我們繼續努力.
除了祈禱時間的長短.
第二個就是.
美國人他們多數祈禱是為誰祈禱呢?.

$^{41}$都不是很驚訝的結果.
85\%的人是為家人祈禱.
其次就是為朋友和有需要的人祈禱.
不過我沒想過配偶排到第四位.
52\%的人會為配偶和重要的人物祈禱.
聖經說教我們要為在上位的人祈禱.
為教會的領袖祈禱.
有多少\%的人會為這些祈禱?.
23\%和22\%.
不是很多.
不知道這個會不會反映你的祈禱生活的狀況呢?.
第三.
美國人他們祈禱的原因.
他們為了什麼來祈禱呢?.
有七成的人都是為了與神有一個結連來祈禱.
45\%是為了減少焦慮.
還有四成的人是為了解決一些個人的問題.
其實文卷後面有寫很小的字.
還有其他什麼原因.
都是圍繞著自己的.
譬如想成為更好的自己.
想創傷得到醫治.
自己更有自信等等.
這一切大部分都是環繞著自己和即時的需要.
其實為自己的需要祈禱.
這件事是沒有錯的.
因為在聖經裡面彼得前書五章七節都告訴我們.
我們要將一切的憂慮.
捨給神.
因為祂顧念我們.
但問題是.
如果我們的祈禱只是停留在這裡.
只是為了自己.
我們的家人,朋友.
或者我們今生的需要滿足來祈禱.
我們會不會錯過了一些更加寶貴的東西呢?.
主耶穌在《馬太福音》六章三十二至三十三節說.
「這都是愛邦人所求的」.
「這都是什麼穿什麼的」.
「都是愛邦人所求的」.

$^{81}$「你們需要的這一切東西」.
「你們的天父都是知道的」.
「你們要先求上帝的國和上帝的義」.
「這些東西都要加給你們」.
耶穌在這裡不是說.
我們不需要為自己的需要祈禱.
不過祂告訴我們.
我們要先為上帝的國和上帝的義.
即是上帝的自己來祈禱.
然後祂就會將我們一切所需的.
都會加給我們.
今天我們看的經文.
《初一及記》三十三章十二至二十三節.
摩西的祈禱.
正正是屬於這種先求神的祈禱的向導.
因為在這個祈禱裡面.
摩西沒有為自己的需要.
或者祂的滿足來祈禱.
祂只是為了兩件事來祈禱.
第一件事就是為了整個以色列民族來代求.
第二件事就是祂為了祂對神榮耀的學務來祈禱.
這段經文其實是在《初一及記》.
三十二至三十四章.
這個大段落裡面的其中一個段落.
其實在三十二至三十四章.
摩西是一共為了以色列民祈禱了四次.
為什麼呢?背景發生了什麼事呢?.
就是以色列民.
因為摩西上了山.
領受法阪.
四十晝夜不見了.
他們很害怕.
不知道摩西去了哪裡.
於是就要阿倫幫他們造一隻金牛毒出來.
所以他們拜金牛毒.
他們叛逆神.
導致神發烈怒.
要滅絕他們.
摩西情詞迫切地為著民族.
祂的同胞切切的禱告.

$^{121}$在三十二章.
祂首先發出了兩次為神的指紋來祈禱.
求神饒恕和保存他們.
因著摩西情詞迫切的禱告.
神真的寬恕了以色列民.
並且祂說祂願意差遣使者.
在以色列民的前面.
帶他們繼續進入應許地.
也會為他們趕走當中的敵人.
不過祂表明了一樣東西.
在三十三章三節.
祂說:我不如你們上去.
上帝不去.
祂只是派使者去.
領子妹.
不知道你聽到這裡.
你覺得怎樣呢.
這個結果.
你收不收貨.
覺得是否Happy Ending.
沒事了 犯罪被赦免.
又可以繼續進入迦南地.
還可以把敵人趕走.
耶和華是否真的和我們一起.
真的那麼重要.
是否必須的.
可能有些人會覺得.
OK 收貨.
有祝福 有應許.
那就行了.
但摩西不是.
摩西很清楚知道.
沒有神同在的成功是沒有意義的.
所以他繼續站在這個破口上.
第三次向神發出迫切的禱告.
也是我們今天看的這段經文.
在他的禱告中.
流露出他對神有三樣的渴慕.
第一樣就是不膚淺的認識.
第二樣是不一樣的同在.

$^{161}$第三樣是不知識的渴求.
我們看看第一點.
不膚淺的認識.
摩西的祈禱.
其實流露出他對神的追求.
不只是膚淺地想知道神的計劃.
而是他很想深刻地認識上帝自己.
首先摩西在禱告中.
他知道神說他不和他們同去的時候.
摩西如何祈禱呢.
十二節告訴我們.
「看!你曾對我說.
將這百姓領上去.
卻沒有讓我知道你要猜派誰與我同去」.
摩西在這裡不是抱怨.
為何不去.
其實摩西是一個很有智慧的舉動.
因為他回到他和上帝關係的起始點.
他抓住神最初給他的應許.
作為他祈禱的群隊.
大家還記得上帝如何呼召摩西嗎.
我們回到三章十一至十二節.
當時神找摩西帶領以色列人出埃及時.
摩西很膽怯地說.
「我是甚麼人竟能去見法老.
把以色列人從埃及領出來」.
神如何回應他.
他不是給他一些策略或資源.
而是上帝給他一個關係性的應許.
他說「我必與你同在.
這就是我猜派你去給你的群隊」.
所以當金牛毒事件發生後.
神的同在彷彿要馬上收回的時候.
摩西馬上回到重要的基礎.
彷彿他在提醒上帝.
也在提醒他自己.
你是信實的神.
你說過 你應許過.
你會與我同在.
不要忘記了.

$^{201}$除了摩西提神有這個憑據之外.
他再進一步這樣說.
「你還說 我按你的名認識你.
你在我眼前蒙了恩」.
「在你眼前蒙了恩」.
這個意思不是指我做的事情.
都可以很順利很成功.
而是在說一個關係.
在這個關係裡面.
是指這個人在神面前蒙恩的這個人.
在神面前擁有一個蒙月立.
被愛並且被認識的一個很親密的地位.
摩西深知道他被揀選作為以色列的領袖.
就是基於他有這個蒙恩的關係.
所以他大膽地用這個關係.
向神發出一個祈禱.
十三節這樣說.
「我如今若在你眼前蒙恩.
求你將你的道指示我.
使我可以認識你」.
我們要留意摩西的說話.
摩西不是說讓我認識你的道.
就是讓我知道你有什麼計劃就這樣完結.
而是他的說話是說.
「將你的道指示我」.
目的在哪裡呢?.
就是要讓我可以認識你.
對於摩西來說.
知道認識神的道.
其實只是一個方法.
一個way 一個tool.
但是認識神自己.
才是他終極的目標.
過去摩西見過紅海分開.
見過天降嗎哪.
這些經歷讓他知道上帝的大能.
來到這一刻在這個危機裡面.
他渴求的是更深一層.
他能夠去認識上帝信實的屬性和心儀.
親愛的弟子姐妹.

$^{241}$我們思想一下.
我們代入一下摩西.
如果你是摩西.
你在這個處境裡面.
你會怎樣祈禱呢?.
面對領袖失能.
百姓叛逆.
神在發憤的時候.
你祈禱第一樣會祈什麼?.
當然是解決方法.
最好上帝消滅敵人.
擺平麻煩.
給我方概.
怎樣懲治以色列等等.
我們是很實際的.
但是我們看到.
摩西在這裡的祈禱.
他聚焦在最深的需要.
他知道最重要的是什麼?.
是要求神恢復關係性的同在.
並且讓他更深刻地去認識神.
他只要這兩件事.
他不需要解決的方法.
為什麼?.
因為他清楚知道.
只要神仍然和他們一起.
只要他能夠知道上帝的心意.
其他所有的事情.
都能夠解決.
都能夠找到出路.
耶利密書33章3節這樣說.
「你求供我,我就應運你.
並將你所不知道.
又大又難的事指示你」.
這節經文不是我們知道隱密的事.
你祈禱就知道秘密.
不是.
這節經文其實是允許我們.
在祈禱當中.
能夠更認識那位啟示隱密事的上帝.

$^{281}$心願我們的祈禱.
能夠超越求事情解決的層面.
可以進入求更深認識上帝的親密關係.
因為這才是禱告的核心和禱告的目標.
神如何回應摩西的祈禱呢?.
14節神回答他.
「我必親自去,讓你安心」.
這是一個很寶貴的應許.
不過摩西是否接受呢?.
我們看到摩西沒有因為這樣而滿足.
為什麼呢?.
中文字我們未必看得到.
但摩西知道.
因為神這裡說「讓你安心」的「你」是單數.
只是讓摩西一個可以安心.
因為神的同在只是針對摩西一個人.
但摩西是否想這樣呢?.
其實神不是第一次.
32章已經說過要滅掉所有以色列人.
剩下你的後裔成為大國.
但摩西不為所動.
因為他心裡很深愛他的民族.
所以當神這裡說「讓你安心」的時候.
摩西覺得不夠.
他渴求的是一種屬於以色列民.
整個民族的同在.
所以他馬上把祈禱改了.
由「我」擴展成為「我們」.
目標很明確地跟神說.
「你若不親自去,就不要把我們從這裡領上來.
現在人如何得知我和你的百姓在你的眼前蒙恩呢?.
豈不是因為你與我們同去.
使我和你的百姓在地上.
與地上的萬民有分別呢?」.
摩西不再用「我」的祈禱.
是「我們」或「我」和「百姓」.
摩西這裡很清楚.
跟神的祈禱裡.
他堅持想要的是什麼呢?.
是想要神親身,親密.

$^{321}$和整個群體的同在.
對於他來說.
神的同在是以色列人身份的核心.
因為辨別以色列民和萬民的分別.
不是在於軍事上的強悍.
或土地上的多寡.
而是在於一句說話的實踐.
就是「神與我們同在」.
如果沒有了這句說話的實踐.
就算去到應許之地.
也會失去那個屬靈的意義.
似乎摩西的祈禱抓住了神的心.
因此神再次來到應許之地.
十七字裡說.
「你所說的這件事我也會去做.
因為你在我眼前無了因.
並且我按你的名認識你」.
在這段對話裡.
我們看到摩西期望的同在.
其實有兩個特質.
第一個特質就是群體性的同在.
雖然摩西作為以色列民的領袖.
但是他沒有特權獨享主的同在.
相反他最重要的職責.
其實就是要將民族.
他的百姓帶進神的同在當中.
弟兄姊妹其實我們的祈禱也是.
今天我們的禱告.
不單單只是祈禱.
主啊你祝福我.
而是我們的祈禱.
更加應該禱告.
主啊願你的同在.
臨在我的家人.
我的朋友.
我的教會.
我的社會.
我的城市.
真正的代求.
是我們渴望其他人.

$^{361}$同樣可以經歷.
我們所經歷的上帝.
這是一個群體性的同在.
第二種的特質.
就是一個面對面的同在.
在經文裡所指的同在.
其實在希伯來文原文的意思.
是說face 面.
這個不單止是.
我們共同在一個空間裡的同在.
就算在一個禮堂裡.
大家也可以做不同的事情.
但真正的同在.
是你的臉不會看著對方的臉.
是有一種親密.
是專注的相交在當中.
這個讓我想起.
現代生活.
一個很寫實的圖畫.
就是我們經常說的.
同桌吃飯是怎樣的.
各自修行.
同桌吃飯.
都坐在同一桌.
但這個各自修行.
大家覺得是不是同在呢.
不是的.
今天我們眼睛和心思都放在.
手機和其他不同的東西上.
身在.
但心卻是缺席.
這個不是真正面對面的同在.
我記得我女兒小時候.
幼稚園初小的時候.
她一直都不會因為要買玩具.
買零食.
任何東西.
來對柳仔哭.
不過有一個情景.
她經常都會哭.

$^{401}$這個情景就是.
當我晚上回家吃飯.
因為以前我事奉的教會.
和我家是很近的.
所以我會晚上回去吃飯.
陪陪女兒.
然後回教會開會.
或者做其他事情.
每一次我吃完飯.
說要走.
你知道我女兒會做什麼嗎.
第一步不是哭.
她會先衝到大門.
用她很小的身軀.
擋住大門是很誇張的.
她會張開雙手.
然後抬高雙腿.
用她最大的空間.
蓋住大門.
讓你不能出去.
有時候不是只有我一個出門.
如果只有我一個出門會好一點.
有時候我老公也會和我一起回教會.
這樣就很難了.
幸好有姐姐.
幫忙哄她.
或者清開她.
但很多時候就算我們成功出門.
在走廊都會聽到我女兒爆哭.
因為她很渴望.
她的父母.
可以繼續停留.
和她同在.
我女兒這份單純的渴望.
有時候我都苦心自問.
不知道我對神的渴望.
有沒有像我女兒那樣.
好像我沒有一切.
所有東西都不要緊.
只想你留在這裡.

$^{441}$和我有面對面的相處和接觸.
有不少的研究調查.
指出有父母高質量的陪伴的孩子.
更有自信和更快樂.
不知道大家有沒有想過.
屬靈的生命會不會有一樣的效果.
當我們花更多的時間.
和我們的天父面對面相處溝通.
我們享受在他的同在的時候.
我想我們的生命.
都會經歷很大的喜樂.
滿足和平安.
因為這就是在主裡面.
神給我們的禮物.
主耶穌也稱讚瑪利亞.
說她選擇了那上好的福分.
就是瑪利亞願意放下所有的事務.
坐在主的腳前.
享受那份面對面的同在.
但願我們都不被次等的事物纏繞我們.
天天我們都選擇上好的福分.
讓他不一樣的同在.
能夠塑造我們不一樣的生命.
第三點,不知識的渴求.
其實剛才十七節已經說了神的回覆.
就是祂願意跟著百姓和摩西一起前行.
彷彿這又好像是一個快樂的結局.
祈禱應該可以完結.
為何還有十八,十九,二十三節呢?.
因為摩西並不滿足.
在他和神深入的對話之後.
他的心被更大更深的渴望充滿.
他不再只是滿足於見到神的禮物.
不只是想知道神的計劃和同在.
他更渴望見到送這份禮物給他的那位上帝.
他不單要認識神的作為.
他更渴望認識神的本體.
這是一個靈魂對他創造主最深的渴望和渴望.
因此摩西向神發出了一個很震撼的祈禱.
「求你顯出你的榮耀給我看」.

$^{481}$摩西想見到天父神聖的本體.
神的回應同樣令人感到很震撼.
因為神沒有責備摩西的貪心.
反而他是閱立他這一份的學務.
十九節他說「我要顯示我一切的微善.
在你的面前經過.
並要在你面前宣告耶和華的名」.
神閱立.
不過他也鄭重地告訴摩西.
二十節說「只是你不能看見我的面.
因為沒有人看見我還可以傳喚」.
神的榮耀很偉大.
以致墮落有罪的有限的人.
沒有辦法直視他全部的光芒.
而仍然可以傳喚.
不過很有趣.
在神的裡面有恩典.
神為摩西預備了一條出路.
一個方法.
二十一節說「看,靠近我這裡有一個地方.
你可以站在這盆石的上面.
當我的榮耀經過的時候.
我必將你放在盆石縫裡.
用我的手掌遮掩你.
等我過去.
然後我要將我的手掌收回.
你就可以見我的背.
卻看不到我的面」.
這是一個很有神學意義的圖畫.
神沒有讓摩西看他的面.
因為他的榮光太強.
但他說要讓他看見他的背.
這個背影其實代表著甚麼呢?.
有三個意思.
第一個是恩典的遮蓋.
有限的人不能夠直視上帝的光芒.
但神願意親自保護軟弱的我們.
讓我們能夠在他的榮耀的餘暉中存活.
所以這個背影是恩典的遮蓋.
第二,他的背影也是代表著啟示的漸進.

$^{521}$神樂意向尋求他的人啟示他自己.
但也要按照我們所能承受的程度.
一步一步地讓我們認識.
所以在舊約就是背夢.
當新約耶穌來到的時候.
我們就可以從耶穌的身上看見上帝的榮耀.
第三,這個背影也是代表著榮耀的蹤跡.
很多時候我們在苦難中.
在困境中.
我們未必會見到神的同在.
他的恩典榮耀.
不過當我們過了苦難之後.
回望整個歷程的時候.
我們就會很清晰地看見.
神恩典的作為和榮耀的蹤跡.
摩西發出這個禱告之後.
他不是見到很震撼的宇宙奇觀.
他只是見到神的背影.
但這個背影就是代表著.
神過去一直的同在.
他對摩西的護臂.
和他在生命裡恩典的痕跡.
頂姐不知道當我這樣說的時候.
你會不會又想起.
在你的生命裡曾幾何時.
你都見過上帝的背影呢?.
我預備講章就想起.
我也見過一次.
在2023年前兩年.
在復活節假期裡.
我見到神的背影.
不知道大家有沒有印象.
2023年是COVID差不多完結的日子.
三年封關的日子終於完結.
通常香港人第一件事是去旅行.
我也不例外.
在復活節假期.
我和家人.
我老公女兒和老爺奶奶.
我們五個人去了韓國自由行.

$^{561}$在六日五夜的行程裡.
我奶奶在第六日凌晨.
她突然過世.
其實她在第五日早上已經不舒服.
她開始發燒和周身很痛.
所以她留在酒店休息.
到了晚上她就越來越辛苦.
於是我們就送她去酒店.
附近我們剛好有醫院和急症室.
差不多凌晨12點.
到了凌晨3點半.
她心臟病發搶救無效.
其實搭進這個急症室開始.
第一個面對最大的困難就是溝通.
因為我當然不懂韓文.
她也不懂英文.
任何一個醫護人員都不懂.
我們就一人拿著一部手機用翻譯軟件翻譯.
其實一句說話就要問幾次.
才能準確地得到資訊.
我們在奶奶過世之後.
第一件最關心的事.
當然是如何將遺體運回香港.
以及如何取得死亡證.
因為奶奶過世是在禮拜六的凌晨.
在韓國禮拜六和禮拜日是公眾假期.
醫院和大使館全部都沒有職員來辦公.
我們由凌晨3點多到6點.
一直和職員溝通.
結果是零.
什麼都得不到.
沒有資訊.
職員不斷叫我們先回去休息.
因為已經天亮了.
她真的無法解答我們.
於是我們就回酒店休息.
聯絡一些家人.
發出了一些代禱的請求.
當時我還沒有消化到這個消息.
看著老公和我的老爺.

$^{601}$其實我都不太知道如何祈禱.
我只在心裡呼求神.
神你的同在要臨在.
帶領我們.
其實都不知道如何做.
因為又不認識人.
所有的規矩都不懂.
但我一直在發呆的時候.
突然間收到一個信息.
我有一個前教會的夫婦.
他們已經移民到澳洲.
她知道我的消息.
她發了一個信息給我.
告訴我她認識一個同事是韓國人.
她懂得韓文和英文.
我問她今天有沒有空.
可不可以來附近幫忙.
我覺得都無所謂.
因為我真的不用.
於是她就聯絡我.
很感恩她說剛好這個朋友.
真的在我婆婆的國生醫院附近.
正在做一些事情.
其實首爾很大.
她也可以不在首爾.
但剛好就在附近.
於是一個小時後.
我們就約在醫院等.
她幫我們聊了很久.
幾個小時後.
突然間她就拿了一張死亡證給我.
婆婆的死亡證.
說了星期六不工作.
但不知道為什麼.
她就可以辦理到這張死亡證.
並且她幫我們聯絡了一些.
可以運遺體的代理人.
幫我們拿了報價和手續.
她說得很清楚和寫給我們.
去思考怎樣去做.

$^{641}$最後她整個下午就在陪我們.
當她要離開的時候.
我望著她的背影.
心裡有一句說話就響起.
「這是神派來的使者」.
在當地我們什麼人都不認識.
什麼都不知道.
但神很奇妙地.
就在我很無助的時候.
這個人出現.
就好像上帝出現的一樣.
幫我們整個家庭處理了很多事.
這句說話一出.
望著這個背影.
神的使者.
我接待了神的使者.
我的眼淚終於都流出來.
整個早上都沒有流過.
因為心裡突然間很安慰.
雖然很痛.
但見到主的使者.
就好像見到主一樣.
很踏實.
也很有力.
詩篇145篇18節.
同樣地應許耶和華臨近凡求告他的人.
臨近所有誠心求告他的人.
是真的.
神回應我們禱告的方法.
不一定像神跡般扭轉劣勢或困局.
相反.
好像詩人所說的一樣.
是他親自臨在我們困境裡.
用他的美善.
就像摩西所見到的美善.
用憐憫,恩典,慈愛,誠實.
去遮蓋我們.
以致當我們經歷過這些困難之後.
我們回望就會發現.
他恩典的足跡.

$^{681}$也能夠更深去認識.
他的屬性.
但願我們的祈禱生活.
我們都可以效法摩西.
這個最終極的學務.
就是更深去認識.
神的榮耀屬性.
總結.
摩西的祈禱.
其實帶來了極其榮耀的結果.
因為他修復了破裂的關係.
扭轉了神的審判.
他也恢復了.
神和以色列人寶貴的同行.
神親自去答應和他們入迦南地.
並且這個祈禱最終.
改變了摩西他自己.
如果你繼續看經文.
我們今天看33章末.
如果你繼續看下去34章.
就是摩西發出的第四個為以色列人的祈禱.
他上山去領受新的法板.
然後他下山之後.
人們就從他的念想.
反射了神的榮光出來.
讓所有人看見他.
都知道他真實地和神一起同在.
弟兄姊妹今天其實我們在座的每一個.
我們比摩西更加有福.
摩西只是看到神的背影.
我剛才說.
我們今天藉著耶穌基督可以怎樣.
看見神的榮耀.
哥林多後書四章六節這樣說.
那吩咐光從黑暗裡照出來的神.
已經照在我們心裡.
使我們知道神榮耀的光.
顯在耶穌基督的念想.
所以弟兄姊妹們.
願我們都能夠成為渴求神榮耀的禱告者.

$^{721}$無論我們在什麼境況.
我們都追求對神有更深入的認識.
不只是要解決的方案.
並且我們都追求與神面對面的同在.
以及我們對主帶著一份更深的渴望.
祈求我們最終可以像摩西一樣.
被耶穌的榮耀觸摸和改變.
我們的念想會發出從主而來的光彩.
可以照耀這個黑暗的世代.
並且可以將更多更多的人.
帶領他們進入上帝同在的當中.
和我們同心一切來祈禱.
真是天賦我們到你面前.
我們去多謝你.
多謝你讓我們有禱告的權柄.
並且多謝你.
雖然你知道我們很有限.
我們有罪.
但你卻鍥而不捨.
你都願意展現你自己讓我們認識.
好讓我們的生命和禱告更貼近你.
多謝你今天讓我們一起學習摩西的祈禱.
主願你來熬造我們每一個.
讓我們都能成為像摩西一樣的禱告者.
願你把那份渴望的靈放在我們裡面.
讓我們不滿足於只是認識你的作為.
而是我們渴望更深的去認識你自己.
我們也求你擴張我們的境界.
讓我們不單單為我們自己的需要去祈禱.
主教我們如何為別人.
為社會,國家.
甚至乎是上帝你的國度去禱告.
以致我們從這些禱告裡面.
也與你更深的結連.
神父我們更願你與我們同在.
這句應許都實現在我們每一個生命裡面.
成為我們生命最真實的記號.
天父願你繼續讓你今天的說話.
在我們的心裡面繼續去熬造我們.
繼續去發酵.

$^{761}$幫助我們成為一個你更喜悅的禱告者.
求你就這樣祝福我們每一個.
多謝主替我們祈禱.
奉主耶穌基督得勝的名字祈求.
阿們.
\newpage



\section{使徒行傳 9:1-6}
\label{sec:031618}
\textbf{大馬士革奇遇記 (普)}~~~~曾興才牧師
\newline
\newline
連結: \href{http://www.church.com.hk/acms/content.asp?site=cdc\&op=show\&type=product\&code=031618\&layout=sermon}{\texttt{031618}} ~~~~ 語音日期: 2026-02-08
\newline
\newline
\hyperref[sec:031617]{\small{< < < PREV SERMON < < <}}
~
\hyperref[sec:index_chronic]{\small{[返順時目]}}
~
\hyperref[sec:031619]{\small{> > > NEXT SERMON > > >}}
\newline
\newline
使徒行傳 9:1-6
\newline
\begin{longtable}{cl}
\hline
\hline
章節 & 經文 (和合本修訂版)\\
\hline
9:1 & \begin{tabularx}{0.7\textwidth}{X} 掃羅不斷用威嚇兇悍的口氣向主的門徒說話。他去見大祭司， \end{tabularx} \\ \\ \relax
9:2 & \begin{tabularx}{0.7\textwidth}{X} 要求發信給大馬士革的各會堂，若是找著信奉這道的人，無論男女，都准他捆綁帶到耶路撒冷。 \end{tabularx} \\ \\ \relax
9:3 & \begin{tabularx}{0.7\textwidth}{X} 掃羅在途中，將到大馬士革的時候，忽然有一道光從天上下來，四面照射著他， \end{tabularx} \\ \\ \relax
9:4 & \begin{tabularx}{0.7\textwidth}{X} 他就仆倒在地，聽見有聲音對他說：「掃羅！掃羅！你為甚麼迫害我？」 \end{tabularx} \\ \\ \relax
9:5 & \begin{tabularx}{0.7\textwidth}{X} 他說：「主啊！你是誰？」主說：「我就是你所迫害的耶穌。 \end{tabularx} \\ \\ \relax
9:6 & \begin{tabularx}{0.7\textwidth}{X} 起來！進城去，你應該做的事，必有人告訴你。」 \end{tabularx} \\ \\
[1ex]
\hline
\hline
\end{longtable}
$^{1}$基督教的信仰是改變生命的信仰,.
因為基督教的信仰有那永遠不改變的真理,.
真理可以由生命去經歷的,.
真理也可以憑信心來領受..
所以耶穌基督在約翰福音20-29節對多疑的多瑪說,.
你因為看見了我纔信,.
那沒有看見我就信的人有福了..
我們每個人都有生命中不同的經歷,.
有些人的經歷非常地平淡,.
有些人的經歷卻是非常地精彩刺激..
有一天老張的太太對他說,.
我好無聊啊..
她要她丈夫老張幫她找刺激的事情去做..
然後老張想了一下說,好啊,.
那我們去喫肉丸倉怎麼樣?.
太太很高興地說,好啦..
結果他們就到了一間餐館,.
找到了靠近門口的座位坐下來,.
點了菜以後就開始喫飯..
當太太去洗手間的時候,.
老張就偷偷地趕緊把賬結了,付了錢..
太太回來以後他們就喫飯,.
喫得差不多以後,.
老張就悄悄地問老婆說,準備好了嗎?.
他的太太非常激動地點了頭,.
老張就拉著她,頭不轉地就衝出去,.
直跑,頭都不回來,回頭再看..
一路上太太非常地開心..
在我的老闆看在眼裡,.
就心裡想,每天如果碰上這麼浪漫的夫妻就好了,.
因為剛剛他的太太也付了錢結了賬,.
兩個人頭也不轉地跑了出去,.
怎麼追也追不到..
她的旁邊的人說,不要想喫霸王餐..
今天不是要大家喫霸王餐,.
是要跟大家分享,.
一個人當他遇到耶穌基督以後,.
他的生命就完全被翻轉,被改變..
這個人到底是誰呢?我們一起來看,.
《使徒行章》第九章第一節到第六節,.

$^{41}$我們非常熟悉的..
《使徒行章》第九章第一到第六節,這裡說,掃羅仍然向諸的門徒口吐威嚇兇殺的話,去見大祭司,.
求文書給大馬社或者大馬士革各會堂,.
若是找著信奉這道的人,.
無論男女都準他捆綁帶到耶路撒冷..
第三節說,掃羅行路將到大馬社,.
忽然天上發光,四面照著他,.
他就撲倒在地,.
聽見有聲音對他說,掃羅,掃羅,你為什麼逼迫我?.
他說,主啊,你是誰?主說,我就是你所逼迫的耶穌..
起來,進城去,你所當做的事,必有人告訴你..
掃羅這一段在大馬士革路上的經歷,我稱他為大馬士革的奇遇記..
為什麼說是奇遇記呢?.
他的奇遇記遠遠比1719年出版的小說羅賓生奇遇記更加精彩萬分..
因為掃羅在大馬士革路上的奇遇改變了他的一生,.
不但改變了他的一生,也因此影響了整個基督教界,.
影響了世世代代..
我們先來做幾個重要的觀察..
首先,這裡說掃羅這個人,掃羅的名字是希伯來文的名字,叫掃羅壽..
可能是他父母親想到他們當時以色列第一個國王,記得嗎?.
掃羅王,也同樣是跟掃羅一樣,汴亞敏支派的人,.
因此父母親叫他叫掃羅..
後來,其實他的父親是羅馬的公民,在羅馬時代的生活,.
因此給他另外一個名字叫polus,polus就是保羅..
保羅的意思是什麼呢?保羅的意思就是說微小卑微的..
換句話說,在羅馬世代服侍的保羅,.
後來就一致再用他羅馬的名字,拉丁文的名字,.
就是保羅這個名字,不再用他的名字掃羅..
這裡面告訴我們說,在史羅西納第十三章第九節,.
掃羅也名叫保羅..
第二,這個掃羅到底是什麼樣一個人呢?.
他是猶太人的一個拉比,非常懂得舊約律法,.
而且他在非常出名的法利賽人的老師加瑪略的底下成為他的大弟子..
可以說他是法利賽中的法利賽人,.
非常敬虔遵守猶太教律法的一個人..
因此當他聽說有人說耶穌基督這位被釘死在十字架的囚犯,.
耶穌基督是誰的兒子的時候,他非常生氣,.
就來逼迫所有相信這個道的人..
所以我們可以說掃羅他的個性是非常的幹練,.
也是非常的積極分子的一種,也是非常掌控型的一個人..

$^{81}$因此他到處去抓拿信奉耶穌基督的人,要把他架在監獄裡面,.
就像當年這些法利賽人來逼迫耶穌基督一樣..
然後我們也看到,第三方面,他到大馬士革去,.
拿了文書,從大祭司那裡拿了文書,要到大馬士革去,.
把這些信徒抓到耶路撒冷來受刑罰..
為什麼要到大馬士革去呢?.
大馬士革其實是很多基督徒因為要逃離在耶路撒冷的逼迫,.
他們就紛紛地逃到大馬士革或者大馬士這個地方去..
大馬士革其實離開耶路撒冷一共有兩百二十公里,.
如果用走路來算的話,可能需要七到十天的時間..
如果是騎馬或者騎駱駝的話,可能需要五到七天的時間..
所以在遙遠的路途當中,保羅就帶著這個文件要到大馬士革去,.
想要抓拿這些信奉耶穌基督的人回到耶路撒冷受刑罰..
沒想到就在途中的時候,就有一個大光在那裡,.
所以忽然間他快到大馬士革,忽然間有大光照下來,.
不是太陽的光,乃是耶和華榮耀的榮光照在他身上..
不能照得前面,他四圍都被這個榮光所照住..
你想想看,為什麼耶穌基督復活的主要向這位.
逼迫他門徒的掃羅顯現呢?.
我相信只有一個原因,因為耶穌要使用充滿熱情的人..
看到沒有,耶穌要使用充滿熱情的人..
首先他來尋找生命充滿熱情的人,.
掃羅可以說是生命充滿熱情的一個人,.
他在尋找生命有熱情的人..
請問你對你的生命有熱情嗎?.
但是我們生命熱情的對象,目標往往不一樣..
有些人非常熱情於美食,燒飯,.
記得要將你們的食譜記上來,我們有生命和傳承這個食譜..
有些人的熱情在於畫畫,在於寫作,有些人做於園丁..
那有些弟兄們他們的熱情是什麼?是運動..
今天下午我們會有美式足球的超級盃,.
就是有西雅圖的海鷹隊對英格蘭的愛國者要有一個決賽..
所以請你現在專心聽,球賽在下午,聚會完了以後纔開始的,要專心聽..
所以有不同的熱愛..
請問你熱愛的是什麼?.
保羅他熱愛的是宗教,因為他有這個宗教款..
他非常熱愛猶太教,因此他到處逼迫那些相信耶穌從死裡復活的人,.
相信耶穌是神的兒子,是彌賽的那些人..
所以耶穌基督在尋找的時候去尋找這位充滿熱情的掃羅..
那你看看耶穌如何對待掃羅的熱情呢?.

$^{121}$我們來看看第八節,第九節,十六項第九章..
這裡說什麼呢?.
第九章,我們來看看..
第八節說,掃羅從地上起來,睜開眼睛,竟不能看見什麼?.
有人拉他的手,領他進了大馬寺,三日不能看見,也不喫也不喝..
掃羅是一個領袖型的一個人,.
現在竟然要被人家拉著他的手,才能夠進到大馬寺閣..
對他來說,真是一個學習鮮卑的功課..
三天他看不見東西,三天他在那裡進食禱告,不喫不喝..
請問你有三天什麼都看不見的經歷嗎?.
你可想象這是多麼難受,多麼難過的日子..
我雖然沒有,但是很多年前,我右邊因為眼睛受傷,.
所以我大概五天的時間只能夠用左眼,那時已經很辛苦了..
你可以想象,三天什麼都看不見..
作為這種激進派的一個人,這種非常乾裂型的一個人,.
掌控型的一個人,什麼都不能做,三天的時間..
我相信神就使用這三天的時間,.
讓掃羅好好地反覆思想他過去所做的一切事..
同時也要好好地檢討他過去的生命,.
好好地反省到底他做了什麼事情..
我相信經過這三天的反省,三天的這樣的思想,三天的這樣的檢討,.
保羅領悟了三個非常重要的道理..
首先第一,就是不要把熱情放在錯誤的目標上..
他旁邊說不要把熱情放在錯誤的目標上..
只能不要把熱情放在錯誤的目標上..
他過去真的充滿了熱情啊,充滿宗教的熱情,.
猶太教的熱情,可以說這個宗教狂..
但是現在後來發現到原來他所逼迫的耶穌基督的信徒,.
耶穌基督真的是從死裡復活了..
他知道在大馬色路上他正遇見了這位從死裡復活的耶穌基督..
他才領悟他過去把他的熱情放在錯誤的事情上,.
錯誤的目標上..
好幾年以後,當他繼續宣講耶穌基督從死裡復活的福音以後,.
遭受了很多的逼迫,在耶路撒冷的時候,.
耶路撒冷的百姓就起來攻擊他,逼迫他..
因此保羅就在死路上第二十二章,.
再一次向這些百姓說他在大馬色格的奇遇..
我們來看看第二十二章第四節和第五節..
第二十二章第四節,保羅說我曾逼迫信奉這道的人,.
直到死地,無論男女都鎖納下監..

$^{161}$這時大祭祀的眾長老都可以給我做見證的,.
我又領了他們答於弟兄的書信,.
往大馬色去,要把在那裡奉這道的人鎖納,.
帶到耶路撒冷受刑..
保羅這裡重新在眾百姓面前做這個見證,.
我是要來殺害這些信徒的..
記得嗎?在死路信章第七章,.
第一位為主殉道的是誰?斯蒂凡..
當斯蒂凡被獅頭打死的時候,.
他眼睛往天看,看見耶穌基督站在神的寶座的旁邊,.
他說主啊,請你不要再把罪歸到這些人身上,.
後來就斷了氣..
那時候聖經告訴我們說,第七章六十節說,.
掃羅也非常喜悅斯蒂凡的被殺..
所以你可以想象當時他心中有這種憎恨這些信奉耶穌基督的人,.
連斯蒂凡死的時候他一點憐憫心都沒有,.
反而喜悅他如此被殺..
所以他說其實我就是這樣一個人,.
過去把熱情放在錯誤的目標上,.
因此我拿了文書,.
不但在耶路撒冷要迫害信奉耶穌基督的人,.
我要到大馬士革去,大馬社去,要把他們抓來,帶到耶路撒冷受刑罰..
他明白了,在大馬社的路上,他終於遇見了這位復活的主,.
他知道原來他是一個罪惡極重的人,.
他本來是該死的人,如今卻沒有,因此被神擊打而死..
這個官造作他死,那個榮王造作他死,.
他可以當成被神擊打而死,卻沒有死..
所以他的確明白了,他是一個罪惡極重的人..
這在《羅馬書》裡面他講得非常清楚,.
第三章第十節說,沒有一人,連一個都沒有,沒有一個尋求神的..
第三章二十三節說,四人都犯了罪,虧欠神的榮耀..
第六章二十三節保羅說什麼呢?.
他說,因為罪的功效乃是什麼?.
死,唯有神的恩賜在我們的主基督耶穌裡,乃是永生..
他忽然間明白過來了,.
我本來是一個該死的罪人,往死亡的道路去,.
但是耶穌基督的恩典臨到我的時候,.
讓我有永生的生命..
親愛的弟兄姊妹,你的熱情是什麼?.
你生命的熱情的目標是什麼?.

$^{201}$你過去所追求的是什麼?.
不久前有一個弟兄跟我說,他年輕的時候,.
他心中只有三個有,有錢,有老婆,有孩子..
現在年紀大了,心中有三個無,.
無病,無痛,無傷高..
請問你生命的熱情是什麼?.
你追求的是什麼?.
千萬不要把熱情放在錯誤的目標,錯誤的對象上..
第二保羅領悟了什麼?領悟了耶穌會保護信靠他的人..
當大患照他的時候,他就馬上掉下來,撲倒在地的時候,.
有聲音說,掃羅,掃羅,你為什麼逼迫我?.
我們再看看第九章第五節,.
他說什麼?主啊,你是誰?.
主說,我就是你所逼迫的耶穌..
很有趣,他馬上認出來說,.
就在那個大光,不是太陽的光,是神的榮光照下來,.
馬上知道這位乃是神,是主的意識..
在新哥森林用這個Lord,.
就講了舊的神的意識,你是神啊,神啊,你是誰?.
他不知道這位他所逼迫的是誰,.
主說,我就是你所逼迫的耶穌..
我想當時保羅領悟的感覺很納悶,.
我明明是在逼迫那些信奉耶穌基督的人,.
怎麼會是逼迫你呢?.
其實你們很多地方告訴我們說,.
耶穌說當人們逼迫你的時候就等於逼迫我,.
當人們憎恨你們的時候就等於憎恨了我..
為什麼呢?.
在馬太福音第十六章十八節,耶穌講得很清楚,.
那彼得說你是基督神的兒子以後呢,.
他說我還告訴你你是彼得,.
我要把什麼?我的教會建造在這盤石上,.
陰間的權民不能勝過他,我的教會..
耶穌說是我的教會,我要把我的教會建造在這個盤石上,.
神的兒子耶穌基督的盤石上..
這個教會不是任何牧者的教會,不是領牧師的教會,.
不是我們任何一個人的教會,是耶穌基督的教會..
所以當你在逼迫教會的時候,當你逼迫教會的信徒的時候,.
就等於你在逼迫耶穌了..
保羅忽然間領會過來,.

$^{241}$神在保護他的子民,.
神在保護那些跟隨他的人,攔阻他到大馬士去抓拿更多的信徒,.
回到耶路撒冷受刑法..
這深深地影響了他以後的福祉,.
他以後的福祉遭受多麼多的困難,多麼多的逼迫,多麼多的迫害..
他知道耶穌會保護他,.
刀劍的危險,死亡的危險,.
各樣的危險他都遇到,他仍然如此堅定不移地來信號耶穌基督..
為什麼?.
他有這樣的堅持是因為他在大馬士革的路上遇見了耶穌,.
從死裡復活的耶穌..
第三方面,他所領悟的到底是什麼呢?.
耶穌使用充滿熱情的人..
他過去是充滿熱情的,但是熱情在錯誤的目標,錯誤的對象上,.
他如今他明白了,原來神在找那個充滿熱情的人,向他叫的人,.
過去對他的宗教非常熱忱,.
如今他認識而且遇見了這位在大馬士革的路上向他顯現了耶穌基督以後,.
他的熱情的對象擺對了目標,擺對了對象..
所以他非常清楚知道,神會使用充滿熱情的人..
當他看不見三天三夜不喫不喝的時候,.
神的靈透過異象感動另外一個充滿熱情的人,名叫亞拿尼亞..
亞拿尼亞可以說是那種謹慎型的人,.
一聽到說在異象當中主要他去什麼,.
你要去一個叫直接的找到一個大叔的掃羅,要為他按手禱告..
如果你是亞拿尼亞,你聽了會怎麼樣?.
亞拿尼亞馬上說,主啊,難道你不知道嗎?.
他在一路上人是殺害信徒的,.
如今來到大馬,是要來撞,我們回去要受刑罰的,.
我怎麼能夠去為他禱告呢?.
同樣的問題,我們會有同樣的疑問..
這個人明明是殺害主跟隨你的人,.
怎麼會要我去為他禱告,不是去送死嗎?.
然後主怎麼對他說呢?.
我們來看看第九章十五節..
主對亞拿尼亞說,你只管去,.
他是我所揀選的器皿,要在外邦人和君王,.
並以色列人面前宣揚我的名,.
我也要指示他為我的名必須受許多的苦難..
主說你只管去,因為有三件事情,.
首先我要揀選他,接下來要做什麼?.

$^{281}$要在猶太人當中,還有在外邦人當中宣揚我的名,.
第三,他要為我的名受許多的苦..
如果你去亞拿尼亞,你敢去嗎?.
成為掃羅生命中的第一個貴人,.
幫助他認識,接受耶穌基督..
接下來我們看看第九章第十七節..
亞拿尼亞就去了,進入那家,把手按在掃羅身上,.
按的時候是戰戰兢兢,有點害怕,到底這個人是真是假..
他說什麼?兄弟掃羅,哇,這位本來是他們的仇敵,.
殺害他們,竟然說兄弟,你是我的兄弟掃羅,.
在你來的路上,向你顯現的主,就是耶穌,.
打發我來,叫你能看見,又被聖靈充滿..
十八節是什麼呢?.
掃羅的眼睛上好像有靈立刻掉下來,.
他就能看見,於是起來做什麼?.
受了洗,喫過飯,就見狀了掃羅和大馬士革的門徒,.
同住了些日子,二十節非常寶貴的..
就在各會堂裡宣講耶穌,說他是什麼?.
神的兒子..
亞拿尼亞做了三件非常重要的事情,.
首先為他按手禱告,領受一致的恩高..
所以掃羅那時候眼睛就打開了,.
三天不能看見的眼睛現在打開,.
有靈出來的眼睛掉了下來,他馬上能夠看得見..
第二件事是什麼?.
按手在他身上,讓他領受什麼?.
領受聖靈,他才能夠為神做美好的見證..
第三呢?.
並且為他受洗,使他歸入耶穌基督的名下..
哇!你可以想象保羅生命的轉捩點就是在大馬士革的奇遇..
之後他也沒有閒做..
他在大馬士革和弟兄們住在一段時間裡面,.
他就開始在各會堂裡面做什麼?.
來傳講耶穌,說他是什麼?.
神的兒子,這個非常重要,說他是神的兒子..
過去他完全不相信耶穌是神的兒子..
如今他在大馬士革上遇見了這個父的耶穌,.
他完全明白了,大膽地來宣告他就是神的兒子..
而且大馬士革有二三十間的所謂的會堂,.
你可以想象他那年多麼的忙碌,.

$^{321}$在那裡跟人來傳講耶穌是神的兒子..
親愛的弟兄姊妹,你有熱情嗎?.
怎麼能夠讓你的熱情繼續能夠像耶穌基督一樣,.
像保羅一樣來燃燒?.
怎麼能夠幫助你的熱情能夠繼續被恢復,.
能夠繼續在燃燒呢?.
你要來效法,節力效法,充滿耶穌熱情的人..
要節力效法,充滿耶穌熱情的人..
向你左右看看,有哪一些人,.
他們是對耶穌有熱情的?.
然後找去問他們,你為什麼對耶穌這麼有熱情啊?.
如果你找不到,看不到的話呢?.
那看看我們的劉同牧師吧..
他是我牧者當中最有熱情的人..
每個月到世界各地去奔跑,.
然後傳講神的福音,.
禮拜天回到教會里面還是充滿熱情,.
充滿激情的傳講神的話來感動我們,.
來激勵我們..
找誰去問他,.
他說:林博士你怎麼會有這麼多的熱情啊?.
有熱情的人..
使徒保羅也就是這樣一個充滿熱情的人..
為什麼呢?.
因為他永遠不能忘記他在大馬士革的起義..
剛纔我那邊說,不要忘記你的大馬士革..
什麼是你的大馬士革?.
就是你重生得救的經歷啊..
是你遇見主,是你接受耶穌基督成為你個人救主,.
是你領受神給你的託付,給你的護照的那個日子..
大馬士革的經歷是非常非常的重要..
保羅他一生都沒有忘記這段經歷..
不但他在第九章告訴我們他經歷了耶穌向他顯現,.
後來幾十年以後,.
當他在耶路撒冷繼續地傳講耶穌基督重生的真理的時候,.
當時的百姓就非常的氣憤,就來逼迫他,要來抓拿他..
暫時的保羅他就站起來,再次來述說他在大馬士革的經歷..
我們來看看使徒保羅第十六章第三到第十五節..
當時是像耶路撒冷的百姓,現在是像雅及帕王..
當他在雅及帕王面前要申訴,.

$^{361}$別人要逼迫他,要來控告他的時候,.
他說:王啊,我在路上享武的時候,看見從天發光,.
比日頭還亮,四面照著我,並與我同行的人,.
我們都撲倒在地,我聽見有聲音,用西班牙話向我說:.
掃羅,掃羅,為什麼逼迫我?你用腳踢刺是難的..
我說:主啊,你是誰?主說:我就是你所逼迫的耶穌..
首先第一次他像耶路撒冷的百姓,.
再次見證他在大馬士革路上遇見的耶穌基督,.
這是第二次他在雅及帕王面前申訴他自己的案件的時候,.
他說:我曾經在大馬士革路上遇見了這位復活的耶穌基督..
他重新重述了一次在大馬士革路上的經歷:.
掃羅,掃羅,你為什麼逼迫我?.
主啊,你是誰?我就是你所逼迫的耶穌..
所以我們可以講,像保羅在大馬士革路上的奇遇,.
可以說深深地影響了他整個生命,整個福祉..
所以他在後來第26章,我們看看第19節,.
他對雅及帕王說:我固執沒有違背那從天上來的意向,.
太寶貴了,我沒有違背那從天上來當時所領受的意向,.
所看到的意向,所領受的使命,.
要在外邦人當中,以色列當中傳講耶穌基督使領復活的信息..
大馬士革的經歷,深深地改變了使徒保羅的一生..
請問你記得你的大馬士革嗎?.
你記得你是如何在神面前認罪悔改,.
接受耶穌基督成為你個人的救主嗎?.
你記得你的大馬士革的經歷嗎?.
是如何透過意象或者異夢,把神的夢想託付放在你的心中?.
我自己大馬士革的經歷,就說到1980年2月28號,.
在英國的一個學生福音裡面,在兩百多人面前,.
我在那裡痛哭流淚,為什麼呢?.
覺得自己是一個罪惡深重的罪人..
當時牧師在講道說,你的罪會追上你,.
如果你沒有好好對付你的罪,你將不會有一個美好的生活,豐盛的生命..
當時我聽到這個罪,我就馬上在那裡就痛哭流淚,.
主就好像放電影一樣,就要在我前面放一幕一幕的,.
我過去我所犯的這些錯誤,這些罪過..
偷拿媽媽的錢去看電影,然後跟爸爸媽媽說謊,.
禮拜天沒有去教會,參試班,反而跑到戲院裡面看電影,.
回來告訴媽媽說,我有在試班裡面,你沒看到而已..
這一幕一幕在我面前,我一面看一面痛哭流淚..
在兩百多人面前,我就覺知相敬,耶穌就成為我個人的救主..

$^{401}$從此以後,我的生命就不再一樣..
那你說,我們的生命難道都要像保羅,像你一樣,.
成為傳道人,成為使徒嗎?不是..
如果神有呼召你,帶領你,讓你能夠全身奉祂的話,.
什麼呼召?來跟從我..
我們每個人所領受的使命,要做全時間的,不是傳道人,.
全時間的基督徒,不是主日的基督徒而已,不是part time Christian..
告訴媽媽說,你是全時間的基督徒..
是24小時,一個禮拜七天的基督徒,不是禮拜天,.
做完禮拜以後,走出這個門以後,就是你自己的一個人了..
因為耶穌的生命對我們每個人都是一樣的,來跟從我..
讓我們在每天的生活裡面,不論是在我們的家庭,.
在我們的工作場所,在我們所接觸的地方,在我們去度假的地方,.
我們都要記得,我們是跟隨耶穌基督的人..
在神的國度裡面,只有full time Christian,沒有part time Christian..
你的工作可以是part time的,.
但是基督徒的身份永遠是full time,全時間的..
求神幫助我們..
第二,我們一定要定義追求認識耶穌..
幹嘛要定義追求認識耶穌?.
一定要定義追求認識耶穌..
如果你不下定決心認識耶穌,那你怎麼跟從耶穌呢?.
耶穌要來跟從我,因此你要跟從耶穌,.
必須要好好的來認識耶穌基督..
今天下午的球賽,.
這個西雅圖的海英隊有一個鐵粉,他們叫他超級粉絲,.
他的名字叫Mama Blue,辣媽媽..
他從1976年,當西雅圖這個球隊開始進入到美國世足球隊的競賽的時候,.
他就已經是計票的支持者,沒有缺席過..
每一次他看球賽的時候,都會穿著非常靚麗的這些藍色的服裝,.
還有非常誇張的羽毛在他身上,看看這張相片,就是Mama Blue..
從1976年到他過世的那一年,2023年,.
始終一直都是西雅圖海英隊的鐵粉,.
他們都知道他是超級粉絲..
他認識所有的球員,所有的教練,他都清楚知道,.
因為他定義認識他所熱愛的球隊..
今天不是叫你成為西雅圖的鐵粉,.
要你成為耶穌基督的超級粉絲,Amen?.
要成為耶穌基督的超級粉絲,.
不必穿著藍色的衣服,戴上羽毛,你要花時間來認識耶穌基督..

$^{441}$你我都要成為耶穌基督的超級粉絲,要定義認識耶穌基督,.
你們的熱情纔會不斷地燃燒,不斷地被恢復..
如果你是禮拜天在這裡坐坐,.
禮拜天出去以後就跟教會,跟耶穌基督完全沒有關係的話,.
那你怎麼成為耶穌基督的粉絲呢?.
怎麼知道耶穌基督說什麼,做什麼,要你做什麼,.
他的心意在你身上,在你公司,在你的家庭,.
婚姻,在你個人的工作上,有什麼計劃呢?.
求神幫助我們,那我們真的是願意成為耶穌基督的鐵粉..
因此我鼓勵你,.
如果你沒有辦法每年把聖經從頭到尾看一遍的話,.
每年至少把釋福音書,馬太馬克,路加福音,.
還有約翰福音一條看,至少看一遍..
你把耶穌基督說了什麼話,.
你把耶穌基督做了什麼事情,你要跟從他嗎?.
你要成為他的門徒嗎?.
你不知道耶穌說什麼,你怎麼跟從啊?.
你不知道耶穌基督做什麼,你能夠為他做什麼呢?.
所以求神幫助我們,成為耶穌的超級粉絲,.
花時間定義追求認識他,.
像保羅追求認識耶穌基督一樣..
三十年以後,當保羅寫信給菲利比的教會,.
目前我們小組材料在查的經文,.
在監獄當中,他怎麼告訴我們他如何認識耶穌呢?.
第三章我們來看,菲利比書第三章,.
我們從第七節到第八節,.
他說:只是我先前與我有益的,.
我現在因基督都當作有損的,.
不但如此我也將萬事當作有損的,.
因我以認識我主基督耶穌為至寶,.
我為他已經丟棄萬事,.
看作憤途,我要得著基督..
短短的兩節聖經裡面,基督出現至少三次,.
你可以想象保羅他的心是如何迫切地要認識基督嗎?.
過去我為有益的,當作是有損的,.
講起來很容易對不對?.
你可知道保羅過去當作他有益的有多少東西啊?.
猶太人的拉比啊,高高在上,.
法利賽人中的法利賽人啊,.
名,分,地位,崇高,他都有了..

$^{481}$過去有我有益的這些是因為認識耶穌基督,.
把他當作什麼?當作憤途啊..
什麼叫憤途啊?米田共..
不是垃圾而已啊,憤途啊,.
我們想都不敢想的東西啊,.
碰都不敢碰的東西..
我為了基督的緣故,.
把過去這些名分,地位,權利,.
都完全放下,當作是憤途,.
當作是有損的,為了得著耶穌基督..
你可以想象保羅血義裡面所留的,.
都是充滿基督和耶穌..
你讀保羅十三封書信裡面,.
最中心的他的思想,.
他最中心的神學觀念,.
只有環繞在一個人身上,.
就是基督耶穌身上..
菲利比書特別明顯,.
幾乎每一章,每一節,.
你都看到他講到基督耶穌,.
基督耶穌..
這次我們二月份背誦的經文,.
第三章第十節說什麼呢?.
使我認識基督,.
曉得他復活的大能,.
並且和他一同受苦,.
效法他的死,.
或者我也得從此復活,.
你們可以繼續背下去..
這裡看到保羅使我認識基督,.
曉得他復活的大能,.
因為他的確在大馬色路上,.
遇見了復活的耶穌,.
沒有想到他所逼迫的就是復活的耶穌,.
曉得他復活的大能,.
並且和他一同受苦,.
效法他的死..
你可以想象保羅認識耶穌到什麼地步嗎?.
不但只是認識耶穌從死裡復活,.
他也知道耶穌基督為他,.

$^{521}$為我們世人受了多少的苦,.
曉得和他一同受苦,.
效法他的死..
他就想到耶穌基督,.
這位神的兒子,.
從天上來到地上,.
住在我們人當中,.
取了奴僕的形象,.
有人的樣子,.
自己卑微,.
全心聖服,.
以至於死且死在死之加上..
為了什麼?.
為了完成天父救你我這些罪惡極重的罪人..
保羅他明白,.
所以他要笑他,.
曉得和他一同受苦,.
效法他的死..
換句話說,.
保羅說我如今也是因為基督的緣故,.
哪怕是再多的痛苦,.
中國的苦難,.
哪怕是再多的逼迫,.
我也願意與耶穌基督一同來受苦,.
甚至為了傳揚耶穌基督的福音,.
甚至是死,.
我也在所不辭..
這點神靈意識就是這樣,.
是我認識基督,.
曉得他復活的大能,.
並且曉得和他一同受苦,.
效法他的死..
這是這樣一個充滿熱情的一個人..
在這地方姐妹,.
神也在尋找那些對他充滿熱情的人..
因為神要使用的人,.
不是一個糊里糊塗的基督徒,.
而是對他充滿熱情的人..
因為像保羅一樣,.
付上任何的代價,.

$^{561}$都需要得著耶穌基督的人..
因為充滿耶穌基督熱情的人,.
才能夠影響世代..
而且再講一遍來,.
充滿耶穌熱情必能影響世代..
神非常的公平,.
我們每一個人,.
都只有一次的生命可以活..
你可以選擇為你自己而活,.
因為我們都知道,.
就是虛空的虛空..
只有一次的生命你可以活,.
但是你可以選擇像保羅一樣,.
我是為耶穌基督而活..
你的生命肯定會有意義,.
肯定會影響你的家人,.
影響你作為所接觸的人,.
影響你的工作,.
影響甚至世代,.
像子魯保羅一樣..
我們看看子魯保羅的相片,.
這個最傳統,最被公認是子魯保羅的肖像..
保羅為主傳道了三十幾年,.
結果大概62歲到65歲,.
因為沒有人知道他正確的時期,.
什麼時候為主殉道,.
大概62到65年之間,.
他是為主殉道,.
在羅馬尼魯的統治底下被砍頭的..
子魯保羅他一生因為那大馬士革的奇遇,.
改變了他整個生命..
他一生所盼望所切磨的是什麼呢?.
我非常喜歡這兩節聖經,.
菲律賓第一章二十節二十一節,.
我們一起來讀來..
他說:照我所切磨所盼望的,.
沒有一事叫我羞愧,.
只要凡事放膽,.
無論是生是死,.
總叫基督在我身上照常顯大,.

$^{601}$活著就是基督,.
我死了就有益處..
保羅他說:照我所切磨所盼望的,.
沒有一事叫我羞愧,.
換句話說只有基督在我身上照常顯大,.
什麼意思啊?.
不論我做什麼事情,.
不論是在患難逼迫當中,.
傳染耶穌基督的當中,.
只要人在我身上能夠看見耶穌基督被放大了..
所以保羅的名字是什麼?.
是微小,.
是卑微,.
只要別人在我身上能夠看見耶穌基督的樣子,.
耶穌在我身上照常顯大,.
無論是生是死,.
讓耶穌基督在我身上能夠彰顯出來..
這是保羅一生所切磨的,一生所盼望的,.
你可以想象保羅的熱情到什麼地步嗎?.
只是你說他是顛倒式,.
他是攪亂世界的人,.
羅馬人說他..
再給點多情書,.
他說我是那位癲狂的人,.
不是為別的,.
而是為耶穌基督的緣故,.
我癲狂了,.
I'm out of my mind,.
for him,.
Jesus Christ..
所以他說我活了,.
我活著就是基督,.
我死了就有益處..
每次你看到這句話的時候,.
你有什麼感想?.
你願意跟保羅,.
你敢跟保羅一樣說嗎?.
我活著就是基督,.
我死了就有益處..
因為他既然清楚明白,.

$^{641}$他的生命是基督耶穌基督換來的,.
他本來是一個該死的罪人,.
而且還逼迫耶穌基督他的門徒,.
他理當是走向滅亡,.
死亡的道路,.
如今在大馬斯路上,.
他重生了,.
有機會認識耶穌基督,.
所以他現在我的生命是被撿回來的,.
是耶穌基督買贖回來的,.
因此我活著就是基督,.
沒有別的原因我活著,.
就是因為耶穌基督的緣故,.
因此我死了就有益處,.
我死了就有得著,.
什麼得著?.
得著耶穌基督和我永遠在一起,.
他說這是好得無比的..
向各位弟兄姊妹,.
甚至充滿熱情的同工和弟兄姊妹,.
我們的生命的確影響了我們的城市,.
也影響了至少36個國家,.
390多間分團的弟兄姊妹和他們的社區,.
因為你我心中都有那一股熱情為耶穌,.
怕我們能夠繼續讓這個熱情在我們裡面燃燒,.
讓這個對耶穌熱情的火不要熄滅,.
一直到我們可以跟保羅一樣說,.
我活著,.
就是基督,.
我死了就有益處..
阿們..
我們一起禱告..
親愛的老伯父,.
主要是透過保羅的生命,.
我們再一次看見你是何等地看重你的兒女們,.
我們再一次在你面前說主啊,.
我們需要恢復對你的熱情,.
我們不要忘記我們的大馬詩歌,.
我們再一次能夠緊緊地來跟隨你,.
跟隨你到底,.

$^{681}$再次將我們的生命獻給你,.
背起十字架,.
為天跟隨你..
謝謝耶穌聽我們的禱告,.
奉耶穌基督寶貴的名就..
阿們..
\newpage



\section{以弗所書 6:10-12}
\label{sec:031619}
\textbf{鞏固屬靈爭戰的大後方}~~~~招世超牧師
\newline
\newline
連結: \href{http://www.church.com.hk/acms/content.asp?site=cdc\&op=show\&type=product\&code=031619\&layout=sermon}{\texttt{031619}} ~~~~ 語音日期: 2026-02-08
\newline
\newline
\hyperref[sec:031618]{\small{< < < PREV SERMON < < <}}
~
\hyperref[sec:index_chronic]{\small{[返順時目]}}
~
\hyperref[sec:031620]{\small{> > > NEXT SERMON > > >}}
\newline
\newline
以弗所書 6:10-12
\newline
\begin{longtable}{cl}
\hline
\hline
章節 & 經文 (和合本修訂版)\\
\hline
6:10 & \begin{tabularx}{0.7\textwidth}{X} 最後，你們要靠著主，依賴他的大能大力作剛強的人。 \end{tabularx} \\ \\ \relax
6:11 & \begin{tabularx}{0.7\textwidth}{X} 要穿戴神所賜的全副軍裝，好抵擋魔鬼的詭計。 \end{tabularx} \\ \\ \relax
6:12 & \begin{tabularx}{0.7\textwidth}{X} 因為我們的爭戰並不是對抗有血有肉的人，而是對抗那些執政的、掌權的、管轄這幽暗世界的，以及天空靈界的惡魔。 \end{tabularx} \\ \\
[1ex]
\hline
\hline
\end{longtable}
$^{1}$OK,今天跟大家講的題目是什麼呢?.
講古術靈爭智的大後方.
給個圖片給大家看看.
先看看圖片.
不知你見到這個魔鬼的醜惡面孔的時候.
你有什麼反應呢?.
見到這張照片比較好.
如果魔鬼真的在我們眼前出現的時候.
不知你又會有什麼不同的回應.
可能我們會很緊張,握緊拳頭.
要不就跟他拼了,要不就掉頭就走.
如果魔鬼在這個時候說.
你幫我做些偷騙的事.
或者說這些你當然不會聽.
不如,你跟我離開上帝吧.
會不會聽呢?.
如果魔鬼對你這樣說,你當然不會聽.
聽我講都傻了.
但聖經說,如果魔鬼扮作光明的天使.
下一張是光明的天使.
聖經說得很清楚.
豈不知魔鬼要幫我做光明的天使.
他跟你們說,下星期過年了.
過年,每個人都以為農曆新年對我們中國人來說很重要.
回家過年吧.
但老闆說,不行,現在生活艱難.
可不可以加班.
我給你雙倍加班薪金.
老闆跟你這樣說.
那你接不接呢?.
多兩倍薪金又可以討你老闆的高興.
而且陪家人過年,那就晚點了.
不一定要在年初一.
那你會不會接呢?.
如果是你的話,老闆跟你這樣說.
你會不會接老闆給你的offer呢?.
有沒有人點頭呢?.
我猜你有吧,心裡面點吧.
陪家人是可以的.
不一定要那個時間,找別的時間都可以.

$^{41}$但後來發覺,不單單是過農曆新年.
所有的長週末,春假,寒假.
你都要加班.
很多時候,我們落在這樣的情況下.
你就會發覺,沒錯,賺到一些錢.
但跟家人的時間就越來越少.
所以自然是什麼呢? 陪家人的時間越來越瘋狂.
再過一段時間.
你的家人,朋友就會對你定性.
定性什麼呢?.
You are a workaholic.
你是一個工作狂.
而對我們來說.
我們連陪家人,照顧他們的時間都沒有.
雖然我們相信耶穌.
但我們哪有時間去傳福音.
哪有時間去供奉上帝呢?.
所以我們發覺有些人.
基督徒,我們暮色這麼久.
有時候很熱心的.
但到了某一個階段.
就停下來了.
我們以為我們只是在傳福音和供奉神的戰場上.
暫時休戰.
誰不知魔鬼原來跟以色列學了一個很好的功課.
1967年中東發生了什麼事呢?.
老人家當然會知道.
六日戰爭,你沒聽過六日戰爭嗎?.
1967年的時間,埃及為首.
帶著約旦,敘利亞和阿拉伯聯軍四國圍攻.
突襲以色列.
希望一舉將以色列殲滅.
以後就不要再以色列過了.
誰不知.
以色列的特種部隊.
偷偷繞到敘利亞的大後方.
切斷了他們的後門.
結果短短六日.
以色列一個很不小的國家.
打敗了四國的聯軍.

$^{81}$大家記不記得這個故事.
魔鬼很厲害.
如果你相信神,你也相信有魔鬼的.
魔鬼很厲害.
他知道神要藉著教會.
去完成拯救世界的使命.
魔鬼知道如果他攻擊基督教的教義.
這個教義不對,那個教義不好.
或者阻撓教會去完成使命.
或者攔阻基督徒去傳福音去為信眾見證.
行不行? 正面交鋒.
不是不行,魔鬼很厲害.
他知道效果不大.
而且很費力氣.
但是如果攻擊教會的後方.
攻擊教會的大後方.
如果他成功的話.
魔鬼在屬靈的戰爭上.
他就可以稍微佔上風.
那麼我們教會的大後方是什麼?.
大家知不知道?.
我們正在看耶福所書.
耶福所書第五章尾到第六章.
講了一大堆什麼?.
家人的關係.
你應該怎樣做功.
然後才到第六章.
所以其實魔鬼的策略就是.
要攻擊我們的大後方.
就是我們的家.
和我們的職場.
因為魔鬼知道如果我們家出了問題.
我們教會的事奉會不會受影響?.
當然會受影響.
家裡有個大病人.
你只照顧他,哪有時間服侍?.
所以不單受影響.
甚至我們的事奉全方位會被癱瘓.
現在我們都說工作越來越艱難.
當你工作這麼大的壓力.

$^{121}$淋到我們每個人身上的時候.
會不會令你的家人的關係惡化?.
因為你自己偷不到氣.
回家可能是一個讓你生氣的地方.
抒發你壓力的一個機會.
所以可能不自覺地.
跟家人的關係會越來越惡化.
進一步影響我們的健康.
然後教會的事奉會就放在一旁.
所以大家有沒有留意到.
原來我們的家和我們的職場.
就成為了被魔鬼重點攻擊的目標.
魔鬼知道只要拆毀我們的家庭.
破壞我們的教會這個大後方.
他就會大大削弱了我們.
完成天國事奉的能力.
而且都會分散神的力量.
為什麼呢?.
神王要照顧那些受傷的百姓.
要很努力去幫助重建那些被破碎的家庭.
所以到第六章.
保羅很鄭重地跟我們說.
第十四章說什麼呢?.
「我還有無聊的話」.
保羅很鄭重地說.
「我還有一些話沒說完」.
「要說什麼呢?」.
「你們要靠著主」.
「依賴他的大能大力做剛強的人」.
「要穿戴神所賜的全副軍裝」.
「就能抵擋魔鬼的詭計」.
「因為我們並不是與屬血氣的爭戰」.
「而是與執政的,掌權的,管轄著一岸世界的」.
「以及天空屬靈氣的惡魔爭戰」.
所以保羅說到最後的時候.
他就提醒我們.
「你們要小心了」.
「因為你要知道我們的敵人是誰」.
「我們的敵人不是我們家人」.
「我們的敵人不是我們的同事」.

$^{161}$「不是我們的上司」.
「我們的敵人是什麼呢?」.
「是一個屬靈的惡魔」.
原來中文版的翻譯是什麼呢?.
「因為我們並不是對抗有血有肉的人」.
「而是與天界的邪靈」.
「就是黑暗世代的執政者,掌權者」.
「以及宇宙間邪惡的勢力作戰」.
所以弟兄姊妹.
現在我們坐在這裡似乎平安無事.
但是你們知不知道.
其實我們現在是在一場屬靈的戰爭當中.
我們說得平安無事.
No!.
弟兄姊妹.
我們正在打仗.
而我們打仗的敵人是什麼呢?.
是靈界的惡魔.
這個屬靈的惡魔很厲害.
他懂得玩弄手段.
他能夠操縱黑暗的世界.
聖經形容魔鬼大有能力.
很有力量.
在教會聖經裡面.
有一個很出名的人物.
我們幾個圖應該記得.
這個人叫約伯.
聖經如何形容約伯呢?.
上帝如何為約伯作見證呢?.
約伯是一個什麼人呢?.
大家看到聖經怎麼說嗎?.
很簡單的.
幾句話而已.
約伯是完全正直.
敬畏神.
遠離惡事.
這是聖經的記載.
和上帝為約伯作的見證.
所以魔鬼受不了.
你這麼厲害啊.

$^{201}$我開想方法讓你不相信神.
於是和神討價還價.
我們看到魔鬼嘗試用豬般的手段.
想摧毀約伯對神的信心.
要逼約伯放下他的信仰.
不要再相信神.
你看聖經的時候.
第一章講得很清楚.
這個約伯是很成功的牧場主人.
有超過一萬頭的生畜.
羊七千,駱駝三百.
很厲害啊.
三千,加起來超過一萬.
瀑被成群.
是很有錢的人.
東方而去為最大.
但是魔鬼在一天之內.
用天災人禍.
策動強盜和周邊的仇敵.
外族人搶走約伯所有的家產.
殺死了他很多的僕人婢女.
不但如此.
用強風.
吹倒了他的子女正在吃午餐.
七個兒子,三個女兒.
當場死亡.
停了嗎?還沒停.
然後讓約伯全身長滿了毒瘡.
令他生不如死.
夠慘了嗎?.
比中國的梁天來還要慘嗎?.
還沒有.
最慘的大家知道是什麼嗎?.
最令他難過的是約伯的老婆.
就說她是老公.
你去死了她吧.
你還相信這個上帝?.
這個上帝還值得你去相信嗎?.
你去死吧,不要在這裡獻世.
慘不慘?.

$^{241}$完沒?.
還沒完.
約伯有三個好朋友.
他這麼慘,遇到這麼多的災難.
我們去探望他,安慰他.
但是見到約伯的時候.
約伯,你知不知道你這麼慘.
一定是犯罪得罪神.
這是報應.
你能不能想像約伯.
家裡.
工作.
和妻子.
和朋友的關係.
嚴重受創.
魔鬼厲害嗎?.
魔鬼相當厲害.
我不知道我們在座的弟兄姊妹.
有沒有遇到類似的事情.
我希望沒有.
但是保羅就提醒我們.
弟兄姊妹,你要知道.
我們現在在屬靈的爭戰當中.
我們不要道以輕心.
就好像現在我們在美國.
特別我們在美國.
OK嗎?.
是啊,有一點緊張.
經濟政治不穩定.
我們覺得OK的.
但是原來現在美國.
只是和全世界在打超限戰.
Am I correct?.
你同意嗎?.
而我們的敵人.
教會的敵人.
不是政治.
不是經濟.
這些都不是我們的敵人.
聖經說.

$^{281}$天空屬靈是惡魔.
邪靈.
他沒有能力.
可以用諸般的計謀手段.
去拆毀神的工作.
消除人的信心.
驅使人要離開上帝.
所以保羅提醒我們.
弟兄姊妹,要警醒.
我們必須要警醒.
做好隨時作戰的準備.
魔鬼在我們想不到的時候.
隨時會來攻擊我們.
正面攻擊我們.
我們會防備.
但是不用的,繞個圈.
家裡人有病.
不用大病.
小病都可以.
工作被老闆罵了.
怎麼樣?.
我們就很depressed.
求神幫助我們,弟兄姊妹.
今天有很多基督徒.
我們不是被魔鬼打到靠近.
我們不懂得去求神幫助.
我們沒有做好準備.
去跟魔鬼做一個屬靈的征戰.
所以保羅在同一段經文裡.
講了兩次.
要穿戴,要拿起.
什麼?.
神所賜的全副軍裝.
我們要整裝待發.
隨時做好準備.
魔鬼什麼時候來攻擊我們.
我們不知道.
所以要準備好.
跟魔鬼再征戰.
聖經裡說.

$^{321}$要穿起神所賜的全副軍裝.
保在磨爛的日子.
能夠抵擋魔鬼.
抵擋仇敵.
並且成就了一切.
還能站立不住.
所以我們要站穩.
如果我們在磨爛的日子.
希望能夠抵擋魔鬼.
是很慘,是很艱難.
很多的挑戰.
不知道怎麼去應付.
但是如果我們想能夠抵擋魔鬼.
就要將神所賜給我們的全副軍裝.
穿在我們的身上.
不要將神賜給我們的裝備.
放在一旁.
不去使用,不去處理.
整個裝備包括什麼?.
第一,用真理度大致速要.
我有皮帶.
不過不是打仗.
所以我今天穿得很斯文.
如果打仗的時候.
腰帶很重要.
以前的配件是不是放在這裡.
兵手一邊左一邊右.
水,糧食,各樣的東西.
如果沒有腰帶.
你要打仗的時候.
就不能夠手到拿來.
是不是?.
如果有看武學片的話.
你更加知道.
有沒有看黃飛鴻.
打仗的時候還有什麼?.
穿上長衣.
摺在腰上.
為什麼?.
扎著腰.

$^{361}$靈活一點.
你才可以打仗.
所以腰帶很重要.
我們用什麼做腰帶?.
真理.
用真理做腰帶去促腰.
真理就是代表.
神透過聖經.
透過傳道人.
給我們的教導.
聖經的真理.
用神的話.
用神的真理.
去對抗虛謊.
去對抗那些似是而非的言論.
很多時候.
聖經說的是說謊的始祖.
用謊言來欺騙我們.
攻擊我們.
大家有沒有聽過.
你的同事跟你們說.
你上班.
你的職場.
你的同事說.
這是上面的決定.
不關我的事.
我執行的.
有沒有聽過?.
所以他怎麼做都無所謂.
頭要這樣做.
我是執行的.
只要能夠交差.
用什麼方法都不要緊.
最重要是要搞定它.
有沒有聽過.
上班的時候.
不要想太多.
不用計劃的.
做了再說.
逼到你了.

$^{401}$哪有時間想.
做了再說.
聖經怎麼說?.
聖經怎麼教我們?.
你們作僕人的.
不要只著眼前的事奉.
像是做曲人看的.
總要全心誠實.
更加要敬畏神.
聽上去似乎很有道理.
但其實是推卸責任.
不做計劃.
為了應付工作.
而去做工作.
這個不是神給我們的教導.
大家有沒有聽過.
我不知道你兒子多大.
我兒子這麼小.
當然會鬧彆扭.
長大一點就不會了.
有沒有聽過.
就讓他鬧彆扭.
很多父母又覺得.
現在世界艱難.
千萬不要輸在起跑時.
這個大家已經聽很久了.
所以兒子放學之後.
又要學一樣學那樣.
搞到他們沒有童真.
沒有童年那種喜樂.
為了補償他們.
What can you do?.
你可以做什麼?.
他喜歡買什麼給他.
小事幹而已.
錢買了他開心就行了.
作為一個補償.
我希望我們的家長不會這樣做.
大聖經怎麼說.
要教養孩童.

$^{441}$使得行當行的路.
就是道路也不偏離.
如果我們等到兒子長大了才教.
那時候你想教他還聽不聽你教.
可能太遲了.
所以主佑弟兄姊妹.
求聖靈指教我們.
讓我們能夠按照聖經的教導.
所啟示的真理.
去評估我們在公司.
在家裡.
待人接物的態度.
我們的言語我們的行動.
是否能夠建基於聖經所啟示的真理上.
用公義當作護心鏡正好.
無論我們在職場也好.
在家裡也好.
或者在朋友圈當中.
我們是人嘛.
我們會不會做錯事.
有沒有試過我們當中沒有人做錯事呢.
I think we all do.
我們都有做錯事的時候.
我們對人對事都會有偏見.
你說我沒有.
我很公正.
我很fair.
自己騙自己.
人有軟弱.
這個load up about it.
還有這個世界真的有很多不公不義的事.
你有沒有遇過.
我遇過很多.
大家遇過很多.
這個是現實.
那我們對自己對人.
我們應該怎樣面對這樣的刺激.
例如我開車的時候.
師傅坐在我旁邊.
有時忽然間會大叫.

$^{481}$慢!慢!慢!.
即是叫我慢.
我第一個反應.
前面的車很煞車.
是前面那個人的問題.
你不要見到煞車燈.
你就這麼緊張.
你緊張叫我又害怕.
我離開很遠的.
師傅怎樣說.
師傅說我害怕才叫.
因為你的車速比前面的車速快.
你讓我感覺要撞上去.
我就叫.
無論如何.
師傅說害怕就是我錯.
無論真假也好.
他說害怕就是我錯.
不過有時.
我安靜想一想.
其實師傅是一個很fair的人.
他是一個很公正的人.
他說你的車速比前面的車速快.
雖然你覺得你可以煞停.
但你去到人家車尾才煞停.
我很害怕.
為何你要賴人.
其實問題不是前面那個人開車.
忽然間煞車.
是你的開車習慣不好.
師傅這樣說我.
你明白我的意思嗎.
所以經常要我大叫.
慢!慢!慢!.
但對我來說.
我第一個反應是.
Not me!.
不關我的事.
對我來說不公平.
是你太緊張.

$^{521}$我可以煞停你的車.
不過你太緊張.
是前面那個人忽然間煞車.
關我甚麼事.
人是不是就是這樣.
莫斯都會這樣.
原來我們本能的反應是甚麼.
否認自己犯錯.
這是人本能的反應.
而我們剛才也說過.
在這個世界裡.
真的有很多不公不義的事.
在我們當中.
當不公不義淋到我們身上時.
我們會覺得很無奈.
那你無奈有甚麼反應.
第一為了自保.
不出聲.
或者你會出聲.
但你出聲是偏頗和你的關係密切的人.
人家怎樣我都不理.
總之我和這個人friend.
我就站在他那邊.
不聽之諱.
我們這個世界.
不公不義的事.
我說過很多次.
只會越來越多.
所以為何保羅提醒我們.
要用甚麼.
公義當作護心鏡遮胸.
用公義當作護心鏡遮胸.
無論對家人.
對朋友.
對同事.
要保持你的心地純潔.
良善.
不要掩著良心去做人.
為了保護自己去說謊.
甚至破壞人的名聲.

$^{561}$我們做錯事.
每個人都會的.
分別是甚麼.
我們有沒有勇氣去承認我們做錯事.
做成眼中看為正的事.
今年一月有很多大事發生.
我不是在說戰爭.
在基督教圈子裡面.
有一個很出名的人.
他叫Philip Yancey.
可能大家不認識他.
有沒有聽過.
Today's Christianity 雜誌.
金玉基督教雜誌的主編.
世界知名的作家.
他今年一月的時候.
公開承認自己有.
八年的婚外情.
不是有人告發他.
所以他沒有人發不成敬.
沒有人告發他.
他是自己公開的告訴別人.
I feel sorry.
在過去八年我又婚外情.
這個是出於聖靈的責備.
他不想繼續做一個兩面人.
他在別人面前的時候.
萬人尊敬的基督徒領袖.
但是他知道原來他裡面是這麼污穢.
Philip Yancey 當他公告這件事的時候.
他沒有為自己的失敗去辯解.
他沒有形容自己的行為無傷大雅.
這個世界很多人都有婚外情.
他是將自己的失敗視為罪惡.
他知道他得罪神.
他知道他背叛了妻子.
他知道他自欺欺人.
他承認得罪了神.
傷害了自己的家人.
以及傷害了別人的家庭.

$^{601}$還有他承認.
他破壞了全世界基督徒對他的信任.
在世人的眼中.
Philip Yancey 是另外一個基督教失敗者的例證.
但是我嘗試從屬靈的角度去看這件事.
我認為Philip Yancey 打贏了一場屬靈的仗.
我不知道你同不同意.
有人曾經這樣說.
我忘記作者是誰.
邪惡的核心指標不是罪行本身.
而是拒絕面對自己的罪惡.
邪惡的核心缺陷不是罪惡的本身.
而是在於拒絕承認犯罪.
求天父上帝憐憫我們.
如果我們犯錯 得罪上帝 得罪別人.
求神賜給我們勇氣能夠走在光明當中.
做神捍衛正的事 捍衛好的事.
做一個行公義好憐憫的人.
第三個裝備是什麼呢.
用平安的福音當作走路的鞋.
剛才也說過我們生活在一個不公不義的社會.
再加上我們這個世界 今天不知道明天的事.
信息萬變.
所以大家都承認 越來越多人沒有安全感.
包括彼得道在內.
於是我們不自覺地在我們周圍築起很多圍牆.
導致人與人之間的隔膜和衝突越來越多.
結果怎麼樣呢.
即使在家裡也一樣.
我們的家很多時候就不再是一個安樂窩.
而我們工作的地方就變成一個高壓窩.
大家有沒有這樣的感覺.
大家有沒有發覺我們能夠信任的人越來越少.
而我們的敵人或者我們不喜歡的人越來越多.
有沒有這樣的體會.
但是保羅怎麼說.
保羅說要用平安的福音當作走路的鞋穿在腳上.
意思就是我們要時時做好準備.
將平安和和睦帶到每一個我們去到的地方.
回來教會也好 去小組也好 在家裡也好.

$^{641}$在公司和朋友一起吃飯的時候.
將平安和和睦帶到他們當中.
保羅說若是能行 總要盡力與眾人和睦.
這是保羅給我們的提醒和教導.
在教會要這樣做 在家裡要這樣做.
在職場和朋友圈都要這樣做.
因為我們知道 一言可以輕磅 一言可以瘋磅.
到底我們跟別人說的話使人和睦.
還是使人加深矛盾和掙扎.
聖經的教導是很清楚的.
污穢的言語一句不可錯後.
只要凡事隨時說造就人的好話.
讓聽見的人得益處.
最後一頂帽子我們知道人言可畏.
是中國人的說話 我們都明白.
我們都知道 細細個就聽.
我相信我們在座的人不會說粗口.
不會用污穢怪語和別人說話.
但求聖靈提醒我們 光照我們.
我們會不會在和別人聊天的時候.
不經意地散播是非.
去說一些令人灰心 氣餒 引人反感的說話.
我記得我小時候很自卑.
因為家人不受很多教育.
小朋友做事會不會做錯事.
會不會學慢一點.
會的 所以我小時候怎麼教你都不懂.
不知道你爸爸有沒有跟你妹妹說過.
怎麼說你都不明白.
我媽媽最喜歡說 牛皮燈籠加七.
什麼意思 怎麼點都不明白.
牛皮燈籠還要有些黑漆在上面.
裡面點光沒用 出不來.
現在不是說是非 是說是非.
我岳父經常說我岳母豬這麼蠢.
我聽到也不自在 在我們面前這樣說.
有時候我們不經意地習慣了.
說的原來不是造就人的好話.
是令人灰心 失望 氣餒的說話.
讓保羅的提醒再次在我們心中.

$^{681}$要隨時說造就人的好話.
隨時是什麼 什麼事都是.
無論你遇到什麼事 隨哪件事都好.
無論是什麼事都好 無論在什麼處境都好.
都要說一些對人有益 有鼓勵 有幫助的說話.
叫聽見的人蒙福.
如果說了一句話 令人憎恨對方.
令人不喜歡教會.
令人覺得神很不公平.
那些不是造就人的好話.
我們都不應該說.
求神憐憫我們.
當我們在家裡 在公司.
遇到惡劣的環境.
或者我們自己被惡人批評中傷打壓.
令我們很難受的時候.
保羅要我們做些什麼呢.
我們可以怎麼做呢.
全副軍裝 還沒完的.
要拿著信德作盾牌.
拿著信德作盾牌.
就可以熄滅魔鬼的火箭.
什麼叫信德 信德這個字好像說Virtual.
其實你看英文四經就是Faith.
要用信心 作為一個盾牌.
去抵擋魔鬼的火箭.
而且盾牌用木 用藤造.
射進去之後 火就會熄滅.
要用信心去抵擋和熄滅.
魔鬼利用家人 朋友 同事.
向我們發起的攻擊.
就像我剛才跟大家說過.
還記得約伯.
約伯在身體 心靈極其痛苦的時候.
我剛才說他最痛苦的是什麼呢.
他老婆叫他死.
叫他不要再相信上帝了.
不要獻世了.
受折磨了.
他的好朋友說.

$^{721}$你這樣是因為你犯罪的緣故.
所以這就成了報應.
其實約伯的老婆和他的朋友.
是在約伯的傷口上殺驗.
約伯有沒有離開神.
有沒有埋怨上帝.
沒有.
半句埋怨的話都沒有說.
也沒有離開上帝.
反而去責他老婆是一個無知的婦人.
上帝給我的.
他要拿回是很正常的事.
你怎麼說這種話呢.
我覺得約伯對神的信心.
真的值得神誇讚他.
什麼能夠讓約伯面對這麼多困難痛苦.
他能夠勝過.
因為他對神的信心.
我們要求神一件事.
我們要認識我們的神.
是一個怎樣的神.
首先我們要求神.
幫助我們洞悉魔鬼的詭計.
用什麼來攻擊我們呢.
用什麼火箭來對付我們呢.
然後我們知道.
雖然這些攻擊這麼猛烈.
但是要信得過上帝.
為什麼要信得過上帝.
上帝讓這些東西留在我身上.
我還要信上帝.
這就是約伯的老婆的反應.
今天我想說一件事.
上帝認識約伯.
所以他能夠為約伯做見證.
這個人完全正直.
敬畏神.
遠離惡事.
上帝為約伯做見證.
他知道約伯是一個怎樣的基督徒.

$^{761}$他知道約伯的能耐.
所以大家回去看約伯的記載.
每次當魔鬼和上帝說.
如果你做這些東西.
他們就會離開你.
如果你做這些東西去對付約伯的話.
約伯就不信你了.
每一次當魔鬼和神說的時候.
神其實有條紅線.
你回去看聖經.
好啊,任由你怎樣做.
不過.
不可以過這條線.
你怎樣做都沒問題.
不過,你不要過這條線.
Sorry.
大家明白我說什麼嗎.
因為神知道約伯的能耐.
他能夠頂得住.
去到哪裡.
原來在這個時候.
魔鬼似乎是要驅使約伯離開上帝.
不信上帝.
原來從神的角度來說.
我知道約伯是個怎樣的人.
他能夠抵受你什麼.
這個是試驗約伯信心的時候.
你聽到我的意思嗎.
弟兄姊妹.
可能你遇過一些困難痛苦.
為什麼你還可以坐在這裡.
你沒病過嗎.
你沒遇到困難過嗎.
為什麼你還要信上帝.
因為看到上帝的保守和看顧.
你還可以頂得住.
原來那些是神給我們的磨練.
弟兄姊妹.
約伯知道.
上帝知道約伯的能耐.

$^{801}$上帝也知道你的能耐.
祂明白你的感受.
祂知道你心中的恐懼.
弟兄姊妹.
我們要信一個神.
這個神不單單知道我們的能耐.
神愛我們.
神要我們成為一個更加壯碩的人.
所以弟兄姊妹.
我們在遇到困難挑戰攻擊的時候.
聖經經常說.
要將一切的憂慮交給神.
豈不知這可能是神試驗我們信心的時候.
還有戴上九印的頭盔.
頭盔是用來做什麼的.
我們當中有人踩單車.
今早看到謙謙踩單車回來.
有沒有戴頭盔.
要戴.
為什麼要戴頭盔.
保護頭.
讓我們的頭不會在遇到撞擊的時候.
受到太大的傷害.
這是代表什麼.
是代表我們的思想和行動.
如果我們的頭受傷害.
沒有得到保護的時候.
不單單我們的思想和行動會受傷害.
甚至可能會癱瘓.
從信仰的角度來看.
如果讓我們的心思意念受損.
就會影響對神,對教會的看法.
甚至會癱瘓我們屬靈的生命.
諸葛弟兄姊妹.
大家知不知道今天基督徒所面對最大的危機是什麼.
你知不知道.
基督徒今天面對最大的危機.
不是叛教離開信仰.
而是在不知不覺當中.
用當代的潮流思想文化.

$^{841}$重新定義我們的信仰.
所以所羅門如何提醒我們.
你要保守你的心.
勝過保守一切.
或者翻譯為.
你要切切保守你的心.
因為人一生的果效.
是由心發出.
相對主義認為沒有絕對真理.
你覺得是真的.
對我來說不一定是真的.
於是真理是被私有化.
真理只是個人的感覺.
問題是什麼.
當沒有了絕對的真理.
沒有了絕對的對與錯.
基督教所說的真理.
就變得毫無意義.
今天說個人主義.
強調個人的感受.
我開心不開心.
我快樂不快樂.
比我能不能上天堂來得重要.
有沒有遇到這樣的人.
於是很多人來教會.
他們來教會的.
是希望自己的心靈得安慰.
讓自己能夠舒服一點.
教會變成什麼.
教會變成一個心靈的加油站.
我心靈受傷了,我有癒藥.
來教會,牧師說些安慰的說話行不行.
我們多過看教會是一個門徒訓練中心.
是一個宣教基地.
有沒有這樣的情況.
消費主義成功神學.
現在讓我們不自覺地用財富.
讓人的成就去衡量一個人有沒有價值.
我們用工作的效率.
事業那種成功.

$^{881}$去取代免費的舊恩.
信仰於是被工具化.
信仰和物質拉上了關係.
而現代的懷疑論.
我曾經和大家提過.
解構文化,大家聽過.
我們受到這個影響.
鼓勵人對權威對信仰發出批判.
於是當我們對權威.
對聖經的真理.
用一個批判的眼光去看的時候.
就削弱了我們對神的話語.
甚至對上帝那種敬畏的心.
所以我們不難找到在教會當中.
很多基督徒都不敬畏神.
如果敬畏神怎麼會做這樣的事呢.
很多時候基督徒.
有時候我們自己都是其中一份子.
不自覺地對神沒那麼敬畏.
敬畏父母親多過敬畏上帝.
會不會呢.
還有,懷疑論和解構文化.
也減低了人追求與神建立關係的意欲.
所以更多人站在寬容門外看.
張嘴鼎姐妹剛才和大家說的.
都是最常見的例子.
我和大家分享的可能有點深.
不過其實是很普通的事.
我們經常隨口說的.
我提出來希望能夠幫助大家提高警覺.
保守我們的心思意念.
好像保羅哥要保守我們的心.
聖歌要保守一切.
我們不要被社會的潮流思想不斷衝擊我們.
讓我們與神的關係受虧損.
污染了我們的屬靈生命.
我們讀卷宗完了嗎.
九點半,十點半,答九是嗎.
OK,多講五分鐘.
大家知不知道你們平時每天.

$^{921}$用多少時間看電視,看YouTube.
看社交網站給你的訊息.
有沒有計算過用多少時間.
我做了一個研究,當然問AI.
平均在美國每個大人.
用四到五個小時看電視節目.
上油管,去社交媒體.
TikTok,Facebook那些.
看看有沒有四五個小時.
我想說什麼.
原來我們每天用很多時間.
最低限度四五個小時.
被世界的思潮在洗腦.
那麼你的思想.
每天有多少時間被神洗腦.
不是洗腦,是洗乾淨.
你用多少時間去祈禱.
去靈修默想,去看神的說話.
我們每天除了上班之外.
就是上互聯網.
不知不覺被影響我們都不知道.
反而我們覺得是寶貝的.
是真理的東西.
我們不是用很多時間去看.
所以我們要靠信念.
可不可以改變習慣.
多些跟神親近,多些看聖經,多些祈禱.
多過上網,多過看電視呢.
那著聖靈的寶劍就是神的道.
神的道必須是用悟性去體會.
道這個字不只是說聖經.
道這個字除了聖經之外.
還包括了先知,仕徒,聖經.
神給人的意象,牧師的訊息.
透過不同的環境.
神所創造的一切,將神的真理向我們啟示.
所以道是一個很深奧的字眼.
不只是聖經這麼簡單.
而聖經說,神的道就是聖靈的寶劍.
寶劍是一個攻擊性的武器.

$^{961}$除了可以防衛之外.
神的寶劍可以割斷謊言.
刺穿魔鬼的詭計.
可以戰勝引誘和使徒.
所以耶穌說,等真理的聖靈來了.
祂要使我們了解真理.
真理就能夠釋放了.
真理可以使我們得以自由.
所以最後提醒每一個人.
要時時在家裡和家人.
不單單只是讀聖經祈禱.
大家要去jam一下idea.
好像Prince Tim經常jam一下歌.
應該怎麼唱,這首歌是什麼意思.
怎麼表達它的意義,對我們來說有什麼教導.
很多時候我們看聖經不用腦子.
寬恕目識這樣說.
如果你能夠用多點時間去思想默想的時候.
你會發覺,神的話語真的很寶貴.
讓神的真理成為我們的寶劍.
懂得去分辨是非,對付邪惡.
最後,靠著聖靈隨時禱告祈求.
隨時,隨地,anytime.
第一節我想告訴大家一件事.
是不是秘密呢.
很多時候我們會祈禱.
什麼時候會祈禱呢.
借飯前會祈禱,離教會會祈禱.
有需要的時候會祈禱.
但今天我想告訴大家一個秘密.
這個秘密是什麼呢.
神不介意我們去到祂面前亂說話.
因為很多時候我們不知道自己在想什麼.
神真的很喜歡和我們接觸和親近.
你不用想好再去祈禱.
不需要的.
神是無所不在,anytime,隨時隨地.
想到什麼,想到什麼.
可能還沒想得很清楚.
你跟神越說越清楚.

$^{1001}$我再說,神不介意你跟祂發瘋.
神最喜歡你跟祂發瘋.
多說說你的感覺,思想,看法.
可能還沒很清楚.
你跟神多聊天,多祈禱.
你就會明白神怎麼看這件事.
隨時要多方祈禱代求.
諸多兄弟姐妹.
我不知道魔鬼現在是不是在攻擊你的家.
家裡人有病,剛才說過了.
有癌症,怪病,不知什麼來的.
我的小朋友找不到工作.
我女兒整天躲在房裡不肯出來,不肯吃東西.
不知道是抑鬱還是厭食症.
你工作,每天工作12個小時.
壓力很大,老闆還是不滿意.
我不知道,I don't know.
我不知道現在魔鬼是不是正在攻擊你們的大後房.
如果我們落在類似的情況,應該怎麼做.
如果沒有就更加好.
但是沒有,在這個時候,正是我們要做裝備的時候.
不要到魔鬼來攻擊我們的時候.
我們都亂七八糟,不知道該怎麼回應.
我們一起讀兩封聖經.
耶福所說第六章,第十字,第十一字.
一起讀好不好?.
我們要靠著主,以賴他的大能大力,作剛強的人.
要穿戴神所賜的全副軍裝,就能抵擋魔鬼的詭計.
然後我們看第六章十三節.
所以要拿起神所賜的全副軍裝.
好在磨爛的日子,抵擋仇敵,並且成就了一切,還能站立得住.
當我們的後房,我們的家能夠被保護好.
利用這個時間,如果沒有太多挑戰的時候.
如何能夠鞏固你們家庭之間的關係.
我們的家就能夠成為一個強大有力的見證.
我們的家會遇到問題的,但因為上帝和我們同在的緣故.
我們能夠保護我們的家人.
我們更加可以引領那些未認識救主的人,到神的面前尋求庇護.
求神幫助我們,如果我們能夠在我們的職場站立得穩.
做工現在很難,很多的挑戰,很多的壓力,很多的問題.

$^{1041}$但是如果我們能夠持守信仰,我們就能夠成為黑暗當中的明燈.
照亮人的心田,讓你的同事,讓你的上師渴慕去追求認識神.
你知不知道,你能夠這樣做.
靠著聖靈的能力,基督復活的大能,我們能夠打贏這場宿命的爭戰.
我們一起祈禱.
天保上帝,我們來到你的面前,其實我們都很震驚.
因為我們不經過原來我們正在打一場宿命的仗.
所以我們很鬆散,以致一有問題,我們就跌倒了.
我們就逃避了,我們就失敗了.
靠你饒恕我們的軟弱,今日祈求你令提醒我們要穿著神的全副軍裝.
要隨時做好裝備,因為魔鬼會把握每一個機會去攻擊我們.
它會拆毀我們的家,會在職場裡面讓我們感到難堪.
但是當我們有神與我們同在,當我們穿著全副的軍裝.
我們就能夠抵擋魔鬼一切的火箭,並且能夠以聖靈的寶劍去對付魔鬼.
拆穿它的計謀,求天賦你讓我們成為一個強大的屬靈隊伍.
為主爭戰,並且能夠得勝,奉耶穌基督的聖名.
\newpage



\section{士師記 4:1-24}
\label{sec:031620}
\textbf{信實的上帝}~~~~楊志溢傳道
\newline
\newline
連結: \href{http://www.church.com.hk/acms/content.asp?site=cdc\&op=show\&type=product\&code=031620\&layout=sermon}{\texttt{031620}} ~~~~ 語音日期: 2026-02-08
\newline
\newline
\hyperref[sec:031619]{\small{< < < PREV SERMON < < <}}
~
\hyperref[sec:index_chronic]{\small{[返順時目]}}
~
\hyperref[sec:031621]{\small{> > > NEXT SERMON > > >}}
\newline
\newline
士師記 4:1-24
\newline
\begin{longtable}{cl}
\hline
\hline
章節 & 經文 (和合本修訂版)\\
\hline
4:1 & \begin{tabularx}{0.7\textwidth}{X} 以笏死後，以色列人又行耶和華眼中看為惡的事。 \end{tabularx} \\ \\ \relax
4:2 & \begin{tabularx}{0.7\textwidth}{X} 耶和華把他們交給在夏瑣作王的迦南王耶賓手中；他的將軍是西西拉，住在夏羅設‧哈歌印。 \end{tabularx} \\ \\ \relax
4:3 & \begin{tabularx}{0.7\textwidth}{X} 以色列人呼求耶和華，因為耶賓王有鐵的戰車九百輛，並且殘酷欺壓以色列人二十年。 \end{tabularx} \\ \\ \relax
4:4 & \begin{tabularx}{0.7\textwidth}{X} 有一位女先知底波拉，是拉比多的妻子，當時作以色列的士師。 \end{tabularx} \\ \\ \relax
4:5 & \begin{tabularx}{0.7\textwidth}{X} 她住在以法蓮山區拉瑪和伯特利的中間，在底波拉的棕樹下。以色列人都上到她那裡去聽審判。 \end{tabularx} \\ \\ \relax
4:6 & \begin{tabularx}{0.7\textwidth}{X} 她派人從拿弗他利的基低斯把亞比挪菴的兒子巴拉召來，對他說：「耶和華－以色列的神吩咐你：『你要率領一萬拿弗他利人和西布倫人上他泊山去。 \end{tabularx} \\ \\ \relax
4:7 & \begin{tabularx}{0.7\textwidth}{X} 我必使耶賓的將軍西西拉率領他的戰車和全軍往基順河，到你那裡去，我必把他交在你手中。』」 \end{tabularx} \\ \\ \relax
4:8 & \begin{tabularx}{0.7\textwidth}{X} 巴拉對她說：「你若同我去，我就去；你若不同我去，我就不去。」 \end{tabularx} \\ \\ \relax
4:9 & \begin{tabularx}{0.7\textwidth}{X} 底波拉說：「我一定會與你同去，然而你在所行的路上必得不著榮耀，因為耶和華要把西西拉交給一個婦人的手裡。」於是底波拉起來，與巴拉一同往基低斯去了。 \end{tabularx} \\ \\ \relax
4:10 & \begin{tabularx}{0.7\textwidth}{X} 巴拉召集西布倫人和拿弗他利人到基低斯，跟他上去的有一萬人。底波拉也同他上去。 \end{tabularx} \\ \\ \relax
4:11 & \begin{tabularx}{0.7\textwidth}{X} 摩西岳父何巴的後裔，基尼人希百離開了基尼族，到靠近基低斯的撒拿音橡樹旁支搭帳棚。 \end{tabularx} \\ \\ \relax
4:12 & \begin{tabularx}{0.7\textwidth}{X} 有人告訴西西拉：「亞比挪菴的兒子巴拉已經上了他泊山。」 \end{tabularx} \\ \\ \relax
4:13 & \begin{tabularx}{0.7\textwidth}{X} 西西拉就召集所有的鐵戰車九百輛和隨從的全軍，從夏羅設‧哈歌印出來，到了基順河。 \end{tabularx} \\ \\ \relax
4:14 & \begin{tabularx}{0.7\textwidth}{X} 底波拉對巴拉說：「起來，今日就是耶和華把西西拉交在你手中的日子。耶和華豈不在你前面行嗎?」於是巴拉下了他泊山，跟隨他的有一萬人。 \end{tabularx} \\ \\ \relax
4:15 & \begin{tabularx}{0.7\textwidth}{X} 耶和華使西西拉和他一切的戰車，以及全軍潰亂，在巴拉面前倒在刀下。西西拉下了車，徒步逃跑。 \end{tabularx} \\ \\ \relax
4:16 & \begin{tabularx}{0.7\textwidth}{X} 巴拉追趕戰車、軍隊，直到夏羅設‧哈歌印。西西拉的全軍都倒在刀下，一個也沒有留下。 \end{tabularx} \\ \\ \relax
4:17 & \begin{tabularx}{0.7\textwidth}{X} 只有西西拉徒步逃跑到基尼人希百之妻雅億的帳棚，因為夏瑣王耶賓與基尼人的希百家和平共處。 \end{tabularx} \\ \\ \relax
4:18 & \begin{tabularx}{0.7\textwidth}{X} 雅億出來迎接西西拉，對他說：「請我主進來，進到我這裡來，不要怕。」西西拉就進了她的帳棚，雅億用被子將他蓋住。 \end{tabularx} \\ \\ \relax
4:19 & \begin{tabularx}{0.7\textwidth}{X} 西西拉對雅億說：「我渴了，求你給我一點水喝。」雅億就打開裝奶的皮袋，給他喝，再把他蓋住。 \end{tabularx} \\ \\ \relax
4:20 & \begin{tabularx}{0.7\textwidth}{X} 西西拉對雅億說：「請你站在帳棚門口，若有人來問你說：『有人在這裡嗎？』你就說：『沒有。』」 \end{tabularx} \\ \\ \relax
4:21 & \begin{tabularx}{0.7\textwidth}{X} 西西拉疲乏沉睡了。希百的妻雅億取了帳棚的橛子，手拿著錘子，靜悄悄地到他那裡，將橛子從他的太陽穴釘進去，直釘到地裡。西西拉就死了。 \end{tabularx} \\ \\ \relax
4:22 & \begin{tabularx}{0.7\textwidth}{X} 看哪，巴拉追趕西西拉，雅億出來迎接他，對他說：「來，我給你看你要找的人。」他就進入帳棚，看哪，西西拉已經倒在地上死了，橛子還在他的太陽穴中。 \end{tabularx} \\ \\ \relax
4:23 & \begin{tabularx}{0.7\textwidth}{X} 那日，神在以色列人面前制伏了迦南王耶賓。 \end{tabularx} \\ \\ \relax
4:24 & \begin{tabularx}{0.7\textwidth}{X} 從此，以色列人的手對迦南王耶賓越來越強硬，直到將迦南王耶賓剪除。 \end{tabularx} \\ \\
[1ex]
\hline
\hline
\end{longtable}
$^{1}$給大家看一張圖先.
在歷史上有一個很出名的故事.
就是這位仁慶跟教會的爭論.
大家不知道認不認得他是誰呢?.
有沒有認得?.
什麼?.
是對的 嘉莉略啊.
哇 聰明啊 阿拉奧神真聰明.
他就是嘉莉略一個很出名的數學家 天文學家.
當年嘉莉略用望遠鏡去觀察星空.
他發現太陽不是圍著地球轉.
相反是地球圍著太陽轉.
這個知識對我們來說現在是一個很正常.
小學也讀過 是不是?.
但是當時很多神學家堅信聖經的自然變異事.
聖經有幾句是說太陽又升起又落下.
是來自詩篇十九章和傳道書一章.
所以一定是太陽移動不是地球移動.
結果他們就因為自己對於信仰的傳息那種堅持.
就寧願否認我們看到的事實.
當年這一班自以為為了真理把關的宗教領袖.
就在這個羅馬的審訊廳裡面.
一次又一次盤問嘉莉略.
你是否承認自己錯?.
你是否承認聖經是不可以給你新的理論去挑戰?.
最後這個年老的科學家嘉莉略被迫跪在地上.
他讀他的懺悔文.
公開收回自己的言論是太陽是圍著地球轉的.
不是地球圍著太陽轉.
只因他見到的事實和教廷一直堅持的傳息不一樣.
幾百年之後.
科學證明了嘉莉略說的是真的.
教會就要承認當時的錯誤.
向死了幾百年的嘉莉略道歉.
這個故事提醒我們.
有時我們以為我們一直堅持的信念.
就等於上帝的真理.
可能我們出於虔誠,出於經驗,出於傳統.
覺得自己沒有錯.
但如果我們不肯學習,不肯聆聽.

$^{41}$我們就可能會被真理推翻.
成為了當年教廷的人.
所以今天我們要問自己.
我們堅持的東西是真理本身.
還是我們覺得是好似真理的東西.
兩樣東西是很不同的.
是真理本身還是我覺得那個是真理.
你分得到那兩樣東西是甚麼嗎.
如果大家準備好.
我們就會進入我們今天的講道.
聽這篇講道我們會知道三樣東西.
當時的猶太人為甚麼會拒絕耶穌.
就是十誡為甚麼成為了他們最大信耶穌的阻礙.
而這樣會不會同樣地成為了我們的阻礙呢.
有準備好我們就一齊來.
上一次我們講了約翰福音的頭頭那幾節.
今天我們再繼續講約翰福音的一章.
下面中間的經文.
一章十四至十八節.
道成了肉身住在我們中間.
充充滿滿地有恩典有真理.
我們也見過他的榮光.
正是負獨一而至的榮光.
約翰為他作見證喊著說.
這就是我曾說.
那在我以後來的先於我.
因為在我以前他已經存在.
從他的豐富裡.
我們都領受了恩典.
而且欣賞家恩.
律法是藉著摩西班堡的恩典和真理.
卻是由耶穌基督來的.
從來沒有人見過上帝.
只有在父女.
父懷裡獨一而至.
將他顯明出來.
在這一段的上文.
大家如果看整個約翰福音的第一章.
上文其實會介紹了另外一個約翰.
是施濟約翰.

$^{81}$說他是為真光而作見證.
為耶穌鋪平道路.
接著就到這十五字.
我們正在看的這段經文.
說了一句大家都不太明白的東西.
在我後來的人先於我.
因為在我以前他已經存在.
這個世間你怎樣理解.
在我之前,也在我之後,在我之後,也在我之前呢?.
第十十字為何要提到摩西班堡的律法.
而恩典和真理也是由耶穌來的呢?.
從來都沒有人見過上帝.
那麼耶穌又怎樣將上帝顯明出來呢?.
不知道大家看約翰福音.
會不會通常都有一個印象.
就是很長氣.
那些東西說完又說.
說完又說.
說了一些好像你明白的東西.
但又不是很明白.
大家這些每個字都明白吧?.
擺在一起好像很奇怪.
大家如果有這樣的感覺是非常正常的.
因為約翰福音所針對的思潮.
是接近在二千年前.
我們已經和他隔了很多代了.
很有代溝.
我們現在和下一代隔三年都覺得有代溝.
何況我們隔了二千多年.
我們已經隔了很多代了.
有些東西我們不明白.
所以今天我們要弄清楚他在回應甚麼.
我們要說一下他在回應甚麼.
其實他在回應三個很不同的思潮.
第一個就是諾斯底主義.
就是大家台右這邊.
他在回應的是很強烈的二元論.
就是至高的真神屬於靈界.
而物質的世界是屬於次一級的.
是由邪惡的創造主所做的.

$^{121}$這個思潮就會覺得如果上帝道成肉身.
成為人的話.
就會令上帝承認自己就是邪惡的創造主.
所以約翰福音對這個思潮有很多反駁.
第二 約翰福音還會針對異教世界觀.
特別是希臘的神明.
如果大家有聽過希臘的神話.
很多不同的神明.
太陽神 阿典娜.
時不時他們會落到人間去扮人去欺騙人.
以致那些神明可以滿足到自己的私慾.
這個和約翰福音所說的道成肉身.
本質上是完全不同的.
耶穌就是那個創造萬有的道.
他來到這個世界上要彰顯神的真理.
所以約翰福音一開始.
他很強調真理和道路.
而第三在中間.
你看我今天最想說的就是十誡.
約翰很想回應猶太人對十誡的詮釋.
特別是第一和第二誡.
大家背不背得到第一誡和第二誡.
除愛以愛 保佑別臣.
那第二誡呢?.
不可為自己做偶像.
我們這裡要花一點時間去了解一下.
究竟猶太傳統怎樣去理解十誡頭兩誡.
對他們來說很重要.
但是他們有什麼詮釋的錯誤.
要約翰這麼用力 這麼長氣去反駁他們呢?.
其實在猶太傳統裡面.
上帝是怎樣的?.
是無形無像的.
是無處不在的.
如果有人說這個人就是上帝.
他們就會覺得你肯定是吹水.
這個人一定是人造的偶像.
就好像當年他們在西奈山下面做了一隻金牛一樣.
他說這隻金牛就是帶你們出埃及的上帝.
於是猶太人受到逼迫流散在不同的地方.

$^{161}$受到希臘羅馬人的統治的時候.
他們需要向羅馬人告訴.
其實我們猶太人都是很有文化底蘊的.
很有智慧的.
他們就作了一本書.
叫做《變失拉自焚》.
裡面就很避重就輕地形容.
上帝就是好像一個人.
他不敢說上帝就是一個人.
他說上帝就好像一個人.
好像一個智慧在人的當中.
所以他們就說.
這個上帝依然都是無形無像.
他無處不在.
就好像希臘裡面的Logos智慧一樣.
和人一起.
雖然在人的當中.
但依然都是無形無像.
就好像大家現在用的AI一樣.
叫Siri說你幫我想一下這個決定對不對.
他突然就出聲說這樣這樣做才對.
這個是一個很有智慧的導人.
他這本書就是這樣寫的.
我現在明白多一點.
為什麼猶太人拒絕耶穌道成肉身.
我們一起看看這本書好不好.
《變失拉自焚》.
大家好像會看這本書的.
因為這本是次經來的.
他這樣說的.
大家可能覺得從太初之先.
在永恆的時候已經創造了我.
這一句是很有約翰福音的味道.
如果你掩蓋了《變失拉自焚》的話.
大家可能會覺得這個是約翰福音.
這裡的我其實指的是智慧.
使我的障目安寧.
安置的意思是這個智慧就住在當年.
以色列人出埃及的時候的障目那裡.
跟他們一起移動.

$^{201}$跟他們一起生活.
他們就想強調上帝就好像智慧.
就好像一個人那樣.
去到在人類當中居住.
跟人一起.
但是因為猶太人對於十誡有很特別的理解.
第一誡是什麼.
不可有別臣.
第二誡是不可為自己做偶像.
所以他們就會覺得.
你做有型的事就等於偶像.
他們不敢將智慧.
不敢將上帝形容為一個真實的人.
你只能說好像.
他就好像一個人.
好像一個Siri那樣.
經常跟著你.
以人法.
上帝只能好像一個人那樣.
無形無味地住在人類當中.
跟以色列人一起.
這裡說的是上帝的一部份.
上帝的份子.
上帝的一部份都不能夠成為.
道成肉身.
所以這樣他們才能夠.
繼續他們對十誡的詮釋.
所以去到這裡.
大家就會開始明白.
約翰福音這裡說的道成肉身.
還提到摩西的律法.
摩西的律法就是其實在.
釋拉山拿的法板下來的十誡.
他是在說服聽眾.
這個理解才是對的.
他來到自己的地方.
自己的人不接納他.
如果你看十一節.
你會見到.
他來到自己的地方.

$^{241}$自己的人不接納他.
自己人是甚麼人.
猶太人沒錯.
因為耶穌是猶太人.
自己人不接納他.
是因為他們對於.
摩西的律法的詮釋.
十誡的詮釋.
和約翰福音有一個很根本的分別.
在猶太傳統底下所詮釋的上帝.
是絕對絕對不會成為肉身的.
他絕對不可能用一個人的身份.
出現在他身邊.
所以當約翰說真理和恩典.
出自耶穌彌賽亞的時候.
其實就是在拆解.
猶太人對於十誡的詮釋.
不可有別臣.
不可做偶像.
猶太人世世代代.
這樣傳下來的傳統.
說上帝不可能用人出現的.
這個傳統.
如果有這樣的人.
說自己是上帝.
他一定是一個偶像.
是一個別臣.
這個傳統傳了多少年.
傳了幾千年.
很深刻很牢固地刻在人的生命裡面.
這份信念.
你可以說是執著.
可想而知.
很難去改變.
就算猶太人受到迫害的時候.
他們都沒有將上帝形容成希臘的神明.
以人的身份進入世界.
他極其量說.
他很像.
他很像一個人.

$^{281}$這份信念和執著.
很堅毅.
情況就好像.
大家有沒有見過食物金字塔.
大家可能有見過.
這個新的食物金字塔.
和我們長久以來的教育.
和我們的文化.
和我們家庭的習慣.
都很不同.
告訴我們.
我們以前學習的五穀類是最健康的.
人體所需的東西.
今早大家吃的米線.
都應該是米線多一點.
對我們健康是最大的益處.
最近美國衛生部發表了一個顛覆我們想像的食物金字塔.
和我們以前很堅信的東西完全不同.
他們總結發現原來五穀類不是我們人體最需要的營養.
反而是什麼.
反而是脂肪和蛋白質.
脂肪和蛋白質.
將我們傳統的飲食習慣完全顛倒.
令我們所相信的健康飲食完全不同.
無論我們屬於那一邊的飲食.
你做168還是KETO.
我們自己都要重新去審視對於健康飲食的習慣.
你究竟這樣吃是否健康.
你剛剛吃的米線究竟是否健康.
我們會再一次去審視究竟我們吃的東西是否健康.
現在才跟他說不要吃那麼多不要省吃.
我們會怎樣做呢.
如果將這樣的處境代入猶太人.
在十戒的處境就是他一直相信我們之前信的食物金字塔才是真的.
你現在跟我說倒轉的才是真的.
他們就會覺得頭突然會爆炸.
我小小一代都是這樣信的.
我小小一代都是這樣說上帝是無形無味.
住在人的當中.
你現在跟我說上帝可以倒成肉身.

$^{321}$可以沒有違反十戒.
你可以想像到對於他們的信念的挑戰是多麼的巨大.
多麼的劇烈.
就算我們對著這個新的食物金字塔.
我們都會覺得很衝擊.
我們不斷去找他的科學證據.
質疑他們的立論.
質疑他們的總結.
甚至我們會經過很多輪的思考和驗證.
我們開始慢慢去想究竟這些東西是否真的.
或者有些人就是我不理你.
我吃我的飯.
你吃你的肉.
我就不理你.
我自己的信念是怎樣.
我不需要受到挑戰.
我不需要改變.
同樣地對於猶太人來說.
這個衝擊的過程.
很需要空間和時間.
所以約翰福音很長氣地說.
他說完又說 提完又提.
目的是怎樣.
讓聽眾知道.
有些我們一直以來堅持的信念.
做的傳統.
可能做了很多年.
未必是對的.
上帝是會改變我們對世界的看法.
對我們自己.
對事.
對我們身邊的人的看法.
這個道成肉身的上帝.
就一直在等待著人.
去改變.
好像有些人未信主之前.
他覺得基督教.
都是怎樣.
導人向善.
和其他宗教和佛教一樣.

$^{361}$都是導人向善.
所以大家都是做好事.
然後勸大家做好人.
堅守一些道德的底線.
我怎樣衡量善.
多數都是用自己和別人比較.
自己不算差.
譬如我在M記.
我見到隔壁桌.
吃完東西就不收拾.
我就不同了.
我自己收拾自己的東西.
我多有原則和多有愛心.
這個就叫善.
這個是未信主之前的一個比較.
因為你要和別人比較.
才知道誰壞誰好.
但是被神改變之後.
你現在知道自己是罪人.
發現所謂善的原則.
不是用人的標準和和別人比較.
是上帝完全的公義.
聖潔的標準去看.
本身你會覺得自己都幾好.
但是慢慢會承認自己是一個罪人.
內心有很多自我中心.
計較 驕傲.
只是平時包裝得好一點.
明白罪不只是做錯事這麼簡單.
是離開神自我中心的一個狀態.
以前我覺得對人好.
不傷害人就行.
但信主之後我就發現.
上帝改變我們內心的看法.
原來你表面對人好.
內心都可以很黑暗.
內心都可以很遠離神.
不知道大家有沒有看最新的《尋秦記》.
我應該不會劇透.
我講一些大家知道的東西.

$^{401}$在《尋秦記》裡面.
最後一集已經講了.
《尋秦記》的作者.
將秦王認證的身世重新詮釋.
他將他描述為.
已經死了 在戰場上已經被人打死.
然後是項少龍找了趙國的王子.
趙盤去頂替他.
為了秦王認證之後焚燒坑儒的歷史.
去做解釋.
他不希望有人知道項少龍的存在.
而知道他是一個假冒的秦王.
大概故事線是這樣.
沒有劇透.
如果想看之後就去看電影.
我的看法當然是無從稽考.
是作出來的.
但是都給我們一個新的角度去審視.
秦始皇所做的一切.
究竟他做焚燒坑儒有什麼呢.
上帝作為我們人生的主宰.
我們的作者.
他為我們人生都鋪設了很多很精彩.
很美妙的故事.
他希望我們從他的角度去審視.
我們對人生很多很大的事件的一些看法.
就好像當年猶太人一樣.
他覺得十誡就講了.
不可做偶像不可.
有別神的.
怎麼會有一個神可以道成肉身.
這樣是一個偶像的行為.
但是在我們生命出現的耶穌.
他就成為了一個很大的解法.
就是他怎樣.
當初他上帝的話語.
他由起初已經存在.
他來到這個世界.
在一開始.
在十誡未出現之前他已經在這裡.

$^{441}$他成為了一個話語.
進入了這個世界裡面.
所以當年猶太人對於十誡的詮釋有很大的偏差.
而耶穌就要利用.
約翰福音這麼長的經文.
去重新提及摩西的律法.
神的道路和真理和恩典.
去令他們去審視自己對十誡.
對道成肉身耶穌的看法.
但是呢.
現實是很殘酷的.
人總是想回到自己很熟悉的詮釋和理解.
對自己對別人對事件有一個很安全的位置.
我不會受到挑戰.
其實我信主一段日子之後.
我就被蒙召去做傳道人.
我去讀神學.
當時我就秉承著一個理念.
人活著為了什麼.
敬拜上帝.
在我的詮釋裡面.
敬拜是動詞.
要做些事.
怎麼做都不是.
你不做事.
坐著就叫做敬拜.
我就很拼命地做.
在教會和神學院.
我什麼服事都做.
教會,師他,組長,洗廁所,抹桌椅.
在圖書館執書.
一個年教五十二個禮拜主日課.
情侶輔導,教導輔導.
總之我覺得我懂的東西我都會做.
我把所有東西拿出來服事.
因為我很堅信我自己.
人活著為了什麼.
敬拜上帝.
敬拜一定做些事.
後來有一個神學院老師提醒我.

$^{481}$他問我.
你做這麼多事.
為了什麼.
我聽到之後很愕然.
我心想.
做這麼多事當然是為了上帝.
還用問.
我當然沒有回答他.
他再問.
你做這麼多事.
真正為了什麼.
那一刻我覺得自己好像被冒犯.
我做了這麼多事.
當然是為了上帝.
我竟然換來別人的質疑.
還要有神學院老師這樣跟我說.
後來我回想了一陣子.
我才發現.
這是我的保護機制.
原來我做了這麼多事.
一直是我的生存機制.
一直是我的成長.
告訴我.
我的動機是.
我要用我的表現去換取別人的接納.
我內心有一個很深深處的渴望.
和一個恐懼.
就是很怕上帝因為我做得不夠好.
不夠多.
而不接納我.
就算信了主.
也不代表我們從小到大的信念.
會自動被改變.
當我很努力事奉.
以為我做這麼多事是為了上帝的時候.
後來我才發現.
原來我的自我價值.
很容易綁在我的成績.
綁在我的工作表現.
別人怎樣看我.

$^{521}$我身邊的人有沒有接納我.
原來我做這麼多事.
其中一部份的動機.
是在換取別人對我的接納.
看好我.
於是結果是怎樣呢.
你失敗的時候就會很容易自卑.
很自責.
做得成功的時候就很容易會驕傲.
或者甚至會很害怕.
會失去別人對你一個好的印象.
當我深入被神改變之後.
還要在神學院才開始.
被改變之後.
我才發現我裡面很多堅持的信念.
的動機.
都不是真理.
只是我成長上的一些.
保護機制.
或者我一些生存的狀態.
以前我都很明白.
我是按神的形象被造成.
所以我不是靠我多成功.
多有價值.
上帝才會愛我.
我一定是很明白的.
但我裡面的潛意識告訴我.
我從小到大的生存模式.
建立了一套很牢不可破的信念.
就好像有些人說江山易改.
本性難移.
但在基督裡面又不是那麼絕望.
上帝是給人有機會改變.
他用了很多年去拆解我對於.
所謂敬拜要做很多事的想法.
其實人活著本身.
和上帝建立關係就是敬拜.
人本身的存在.
今天你可以呼吸.
你和他有關係.

$^{561}$那個就已經是敬拜.
也很直白地問我這個問題.
令我當頭棒喝.
重新審視我自己的動機.
給上帝改變.
所以來個總結.
我們應該可以一齊回答.
第一和第二個問題.
就是當時的猶太人為何拒絕耶穌.
十誡為何成為他們最大信耶穌的阻礙.
就是因為世代的猶太傳統.
對於十誡有一個很堅持的理解.
就是不可有別神不可做偶像.
就等於神不能夠成為肉身成為人.
這個到現在他們都在堅持.
如果你去以色列短訓和傳福音的話.
很多猶太拉比都這樣說.
神不能夠成為人.
這個其中一個他們拒絕耶穌的原因.
而第三個.
這樣會不會同樣地成為我們的阻礙.
我猜這個要我們自己才能夠回答.
我鼓勵大家.
第一就是主動檢視自己的堅持.
究竟堅持什麼呢.
我鼓勵大家在上帝面前安靜地問自己.
你現在堅持的是什麼.
是真理 是傳統 是習慣.
還是自己的面子.
我建議大家寫下一兩樣.
自己很不想被人挑戰的觀點.
你帶入祈禱裡面.
你看看上帝怎樣回應你.
也可以在每天讀經裡面特別留意一下.
有沒有聖經對你一些很新的衝擊.
有時候如果上帝透過聖經.
透過祈禱給你一些新的衝擊的時候.
我邀請你誠實地去面對.
勇敢地去看看究竟這件事.
是不是真的要審時.

$^{601}$第二就是鼓勵大家聆聽不同的聲音.
鼓勵大家刻意跟立場不同.
背景不同的朋友聊天.
不是為了辯論.
是為了明白對方 也明白自己.
最近Bella有幾個workshop幫助我們認識自己.
認識我們身邊的人.
對我來說最大的得著就是.
我能夠聽到別人眼中的我是怎樣.
我帶給別人的感覺是怎樣.
我帶給別人的動力是什麼.
也讓我聽到不少令我成長的聲音.
有時候那些成長的聲音對我來說.
本身一些信念是一種挑戰.
但同一時間也令我重新審視.
我一直以來的一些信念.
究竟那些東西是否符合聖經 是否真理.
有時候我們跟朋友傾信仰 傳福音.
我們可以學習一下.
聽一下別人關心什麼.
我們有時候會覺得.
將一套信仰的模式套在別人身上.
他們就會接受.
但我們回想一下 想一下.
當你未信主的時候.
你也不是一下子直接接受上帝所給我們的東西.
心態上的改變 生命上的改變.
是要時間和要關係.
如果你將心比己.
你看人家想一想.
看一看這個世界.
怎樣理解這個世界.
怎樣看身邊的人.
我們才可以找到我們身邊.
跟人有關係的接點.
大家一起祈禱.
天父.
多謝你讓我們有這個時間.
在你的面前.
就像那個食物金字塔一樣.

$^{641}$我們有時候也可能要重新審視一下.
我們人生很多的信念.
很多我們堅持著做的事情.
道理是否正確.
讓我們每一個人.
無論是初信.
還是我們已經跟從很久的弟兄姊妹.
讓我們每一個人都能夠看見我們裡面.
內心深處的東西.
也讓我們互相幫助.
去幫助弟兄姊妹可以成長.
在裡面更加真實.
在裡面更加看見主你的真理.
我這樣禱告.
奉靠耶穌得勝免疾.
阿們.
祝福大家.
\newpage



\section{詩篇 1:1-3}
\label{sec:031621}
\textbf{珍貴的「日」與「夜」}~~~~李慧珊傳道
\newline
\newline
連結: \href{http://www.church.com.hk/acms/content.asp?site=cdc\&op=show\&type=product\&code=031621\&layout=sermon}{\texttt{031621}} ~~~~ 語音日期: 2026-02-08
\newline
\newline
\hyperref[sec:031620]{\small{< < < PREV SERMON < < <}}
~
\hyperref[sec:index_chronic]{\small{[返順時目]}}
~
\hyperref[sec:031622]{\small{> > > NEXT SERMON > > >}}
\newline
\newline
詩篇 1:1-3
\newline
\begin{longtable}{cl}
\hline
\hline
章節 & 經文 (和合本修訂版)\\
\hline
1:1 & \begin{tabularx}{0.7\textwidth}{X} 不從惡人的計謀，不站罪人的道路，不坐傲慢人的座位， \end{tabularx} \\ \\ \relax
1:2 & \begin{tabularx}{0.7\textwidth}{X} 惟喜愛耶和華的律法，晝夜思想他的律法；這人便為有福！ \end{tabularx} \\ \\ \relax
1:3 & \begin{tabularx}{0.7\textwidth}{X} 他要像一棵樹栽在溪水旁，按時候結果子，葉子也不枯乾。凡他所做的盡都順利。 \end{tabularx} \\ \\
[1ex]
\hline
\hline
\end{longtable}
$^{1}$好 鼎姐妹平安.
剛才是不是很開心呢 見到小聰.
最後呢 我已經看了第二次了.
每一次見到最後都很好笑的.
還有都很感恩 我真的見過小聰.
我早兩個星期就放了大假.
去了新加坡旅行.
然後就見到小聰和小聰太太.
我們一起都吃了幾頓飯.
在這個旅程裡面.
我發現一樣東西是非常棒的.
看看是什麼.
有沒有人用過這個功能.
這個就是Google地圖裡面.
一個實時景觀的功能.
就是去到哪裡都不用怕會迷路的功能.
開了它就有箭咀叫你向前走.
或者轉彎.
然後你走到馬路邊.
它會提醒你去看路.
當然你要走的時候都要看路.
不要完全相信科技.
會不會煽惱問題影響了呢.
我覺得科技都是帶來很多方便.
去指引我們的道路.
令到你一個陌生的地方.
都可以很放心的不怕走錯路.
在今天選擇的經文裡.
看到一樣東西.
是很寶貴的.
就是引路方面.
那個珍貴的.
就是神和我們一起.
是日與夜都和我們一起.
很熟悉的經文.
我們先祈禱.
我們一起先有禱告.
親愛的主耶穌.
多謝你.
讓我們回到你的家裡.

$^{41}$一同去敬拜 去讚美.
更加很寶貴的.
你設下你的華語給我們.
一本聖經.
雖然好像很厚.
但我們看不到實體的厚度.
但裡面真的很豐富.
神理子將你的心意顯明.
更加讓我們可以透過聖經.
去認識自己 去認識神.
主耶穌 今天是一個很寶貴的時刻.
求聖靈你親自引導我們個人的心.
你要進度去明白 去領受.
不單是知道 明白 去領受.
這一份寶貴的信息.
無論是宣講的 去聆聽的.
都求真理的聖靈來帶領.
叫我們同盟的去領受祈福.
奉主耶穌基督的聖名自祈求.
阿們.
剛才說引路.
聖經裡面有一個很特別的神蹟.
很熟悉的 先看看.
就是日與夜的引路.
以色列人從前在埃及受苦.
神透過摩西帶領他們的子民.
走出紅海進入曠野的路.
曠野是甚麼地方呢.
是沙漠的地方.
大家知不知道沙漠.
大家知不知道沙漠是怎樣的.
大家有沒有想過天氣是怎樣的.
白天是怎樣的.
很熱.
蒸到沒有水蒸 蒸到人間蒸發.
高度可以超過50度的氣溫.
晚上原來是很冷的.
原來日間的溫度和晚上的溫度.
可以相差最少30度.
甚至可以降到零度.

$^{81}$很想像一群以色列人在走曠野之路的時候是怎樣的.
天氣已經不行了.
但竟然有雲柱 日間雲柱 夜間火柱來引導.
當你由天氣入手的時候.
你會看到雲柱是怎樣的.
看到雲在上空很自由地飄下飄下.
不會靜止的.
即是有什麼.
有風.
很棒 整條雲柱伴著一群以色列人.
會不會有冷氣呢.
大家有沒有這樣想過.
可能真的有一個特別的人的需要.
神都去引導的時候.
裡面都帶來一個涼風.
一個遮陰太陽.
雲柱有多麼的覆蓋 沒人見過.
我覺得這幅圖很棒.
它有一個空間來展現頂頭.
雲柱會不會可以遮陰呢.
然後有火柱.
我們圍爐取暖.
生火使人溫暖.
夜間有火柱.
火的熱力會不會帶來晚上的溫暖呢.
很棒.
雖然在曠野.
雖然在沙漠之地.
不單止身體上的溫度可以承受.
更加神的供應是無斷過.
40年他們的鞋都沒有穿爛.
有沒有人家裡的衣服穿了40年.
你的鞋磨了40年都不爛.
這些全都是奇蹟.
神的引導有神在.
什麼環境呢 是不是問題呢.
環境是一個問題.
自己經歷神的恩典.
問題都可以得到解決.
很棒吧.

$^{121}$很寶貴 很棒.
日與夜可以帶來.
怎樣可以經歷到神這樣愛我們.
神願意愛我們.
經歷祂的保護 同行 帶領.
今日的經文揭示一樣很寶貴的東西.
就是日與夜.
首先我們要選擇對的位置.
我們看看經文.
不從惡人的計謀.
不站罪人的道路.
不坐洩萬人的座位.
我們很清楚什麼是惡人 什麼是罪人.
我們信了很多年.
今天有沒有新朋友在場.
全部都是我們的頂屎妹.
我們要知道什麼是惡人 什麼是罪人.
塞萬人才知道是什麼.
嘴妙妙 態度很傲慢.
很喜歡說驕傲的話 詆毀人.
輕佻 浮躁這類.
就是洩萬人.
關乎我們的生命.
我們有沒有選擇對的路.
就是走在神的道路裡.
我們怎樣去知道神的道路呢.
當然杜牧師有講了兩堂.
聖靈在我們的心裡.
當然會讓我們知罪 引導我們.
更加寶貴的就是引導我們明白真理.
真理從哪裡來.
就是聖經裡面來.
在聖經裡面來.
我們怎樣去領受神的話語.
這個寶貴的禮物 聖經呢.
聖靈是真理的聖靈.
耶穌在約翰福音裡面.
十四章至第十六章.
很清楚的說明.
真理的聖靈就是保衛師.

$^{161}$保衛師就是真理的聖靈.
祂要常與我們同在.
然後告訴我們.
真理的聖靈是來見證耶穌的.
然後告訴我們 真理的聖靈要引導我們明白一切的真理.
所以聖靈在我們生命裡面是很寶貴的.
更加讓我們去明白聖經.
我想問 信耶穌的和不信耶穌的.
看聖經會不會有分別.
一定有分別.
你可以去研究聖經.
甚至乎在外國的大學.
甚至乎我在一間神學院裡面.
會收國內的人.
甚至乎有不信主 讀文學的人.
去研究聖經.
和我一班試過上課.
去微信耶穌.
和我一起收高林多前書.
收分的.
我想不信主都這麼有心讀聖經.
為了他的學業.
原來信主和不信主都可以讀聖經.
剛才我們都說有分別.
有什麼分別呢.
和我們有沒有關係.
聖經所說的 關乎我們生命的事.
和我們有沒有關係.
這個關係 是我們和神的關係.
神接著聖經的真理.
讓我們明白 知道祂 認識祂.
祂是我們的主.
祂是怎樣帶領我們.
祂是怎樣愛我們.
所以聖經是很寶貴.
引導我們將我們的生命.
和我們相信的神連結在一起.
這就是我們讀聖經.
和不信人讀聖經的不同.
而我們去讀的時候.

$^{201}$第二節就到了.
要喜愛.
喜愛耶和華的律法.
怎樣去喜愛呢.
就是晝夜去思想.
接著說 者人便為有福.
這是結果.
有福是一種結果.
而我們喜不喜愛呢.
是我們和神的關係.
有福是怎樣呢 到了第三節.
好像一棵樹再是溪水旁.
按時候結果指.
葉子也不枯乾.
結果是凡他所做的盡到順利.
而這件經文很有趣.
和約書亞記第一章一至八節.
都會有關係.
我們看一看約書亞記.
他說摩西帶領以色列人.
當摩西離世.
他傳承了領袖任務給約書亞.
當時的約書亞是怎樣呢.
震騰騰的.
因為那麼大群以色列人.
超過二百萬人 年男大女.
怎樣去帶領呢.
害怕的 未有經驗的領袖.
一來帶那麼大群人.
整群人都是你帶的.
所以神親自去和約書亞說話.
他跟她說.
這個律法是不可以離開你的口.
總要怎樣呢.
晝夜思想.
好使他謹守遵行書上所寫的一切.
如此 接著寫甚麼.
你的道路就可以亨通.
接著凡事順利.
是否很平衡.

$^{241}$和第一篇第一至三節.
是很平衡的一段經文.
我們看到原來去承擔神的使命.
去跟隨神的路的時候.
要怎樣呢.
跟隨神的話.
更加要晝夜思想神的話.
這是很寶貴的提醒.
我們看看思想.
思想英文是meditate.
希伯來文是客家.
有幾個解釋.
因為希伯來文的字很豐富.
有兩個解釋在這段經文裡很貼切.
第一個是很輕聲的說話matter.
好像呼吸裡很小聲.
很輕聲的說話.
第二是低沉的吼叫聲.
這是甚麼呢.
怎樣低沉的吼叫聲呢.
這個字可以看見以賽亞書.
看看以賽亞書.
第三是一章四字.
這個吼叫聲就是客家.
就是獅子跟少莊.
獅子互食咆哮.
甚麼是互食咆哮呢.
獅子是要獵食的.
有食物.
其實還有其他獅子.
有一份是他的.
就大聲地吼叫.
叫其他獅子吃完再離開.
或者這個地盤是我的.
很享受那個食物.
他是有意圖.
獅子去找食物.
很辛苦去找到一隻野獸.
然後他就要吃他的食物.
在大自然世界.

$^{281}$聲音是怎樣呢.
我們看看這幅圖.
好像很可愛.
一個舞獅和一個小獅子.
這個本身在YouTube是一個影片.
我截了圖出來.
他都很斯文.
不是太大聲.
很小聲.
可能有距離拍攝的時候.
然後我再找到另一條片.
比較血淋淋的.
但我又不想說出來.
但我錄了少少的咆哮聲.
先聽一聽.
對了.
在吃東西的聲音.
這把聲音是有很多鬚的.
很多毛的.
紅絲的聲音.
聲音會比較宏偉一點.
神的創造很奇妙.
就是這種浸淫在裡面.
很投入的.
來進食.
那種喜愛.
飢餓.
來進食.
原來思想.
Meditate.
剛才說的.
思想其中一部分.
要進食.
好像獅子一樣肚餓.
要獵食.
獵食之後.
吃到咆哮聲.
是這樣的狀況.
那種喜愛.
是怎樣去理解呢.

$^{321}$有時候會有一種狀況.
可能是一個環境裡.
逼到你很倚靠神的時候.
才逼出這種喜愛出來.
去投靠主.
我們是不是很多時候都要這樣.
人就是這麼軟弱.
神就是這麼愛我們.
藉著很多環境來呼喚我們.
你來吧.
你肯靠我了.
可能是我們的身體.
可能是我們的工作.
可能是我們的家庭.
可能是我們的關係.
不斷要我們的心去投靠他.
如果我們覺得認同的時候.
即是說.
很多時候.
我們忘記了他.
我們沒有去願意.
好像這麼有企圖.
這麼有意識 這麼喜愛.
去到主的面前.
甚或我們好像.
之前有一些靈修短片.
來到一個故事的形式.
來呈現.
可能我們就是.
做功課.
或者做一件事.
不是很理解原來.
主是很想我們去享受祂的愛.
很想去保護我們.
很想我們在祂裡面.
就好像在沙漠的曠野裡.
都要經歷冷氣.
經歷火爐的溫暖.
祂告訴我們.
祂是我們的主.

$^{361}$我們的神.
我們的好處.
不在祂以外.
這是主對我們的心意.
我們怎樣喜愛呢?.
我會想起一些片段.
青少年的時候.
甚至那種青春期.
和求偶期的時段.
大家知道我想說什麼.
那時很喜歡有些人.
誰喜歡誰.
然後全群人都知道.
誰喜歡誰.
而被喜歡的人是不知道的.
然後有一些情竇初開的男生.
都會比較害羞.
後來被朋友逼迫.
人們就推一推.
然後勇猛出去追女生.
其實那個心.
怎樣去化成行動.
你有多麼喜愛.
多麼有意圖.
多麼想的時候.
你就會走出來.
這一份心.
這一份想望.
都求主給我們恩典.
能夠去經歷它.
另一段片段給大家看看.
他在做什麼呢?.
反戳.
吃下去了.
再來一次.
不斷在他的身體裡.
去重複這個消化的過程.
剛才說的思想.
其中一樣.
就是那種輕聲.

$^{401}$很輕聲.
即是說.
我們讀聖經去默想的時候.
原來是會有聲音的.
剛才兩個字都有聲音.
我們是將神的話輕聲地讀出.
來消化.
來咀嚼.
這樣很有意思.
你讀的時候.
同時耳朵在做什麼呢?.
在聽.
在聽神的話.
我記得蕭二傳道說.
我們聽神的話.
是要行出來.
聽而行道.
我們去咀嚼.
去讀.
是輕聲.
接著去聽到神的話.
接著不斷地去消化.
去默想.
去反覆的思想.
是晝夜的思想.
在聖經裡.
想到但義理.
但義理當時面對的環境.
是很艱難的.
因為他是一個年輕人.
一個年輕的才俊.
青英的分子.
被擄去巴比倫.
而去到巴比倫的時候.
他吃的東西可能有制止別臣.
他說我不吃的.
我要吃青菜就行了.
但是神給了他.
雖然營養看起來好像很少.
但他的樣子是盡美的.

$^{441}$而且智慧也充足.
而當他要面對政權.
跟他說不可以去拜別臣.
只是單單去拜皇帝的時候.
他不妥協.
更加公然有人知道.
因為他打開大窗戶.
一日三次去祈禱.
然後被人告狀.
說他叛國.
然後就被人扔到獅子坑.
為什麼他不怕被扔到獅子坑呢?.
之前試過被人扔到火窯.
死不了.
有神去救他.
就是這一份信心.
這份不是說我有信心.
他真的經歷了跟神緊密的關係.
他定時一日三次去祈禱.
我們一日吃三餐.
他定時三次去祈禱.
我不是想說.
我們要一日多少次去祈禱.
我不是要像伊斯蘭教一樣.
穆斯林他們一日五次去祈禱.
是一種心.
好像我們之前有一條片.
傑傳我在這裡有份說.
我們要關乎我們的心.
不是時鐘上的指針.
那次我們都聽到他笑.
不是關乎時鐘的指針.
是關乎我們的心.
這份心去帶動我們.
晝夜去到主的面前.
Day and night.
是一種我們的操練.
也是靈修的操練.
上次跟大家說.
不知道有沒有人真的回去.

$^{481}$試過感恩試項.
每一日寫三個感恩試項.
有沒有人正在做?.
星期五有人跟我說他正在做.
他老公.
不過不是他.
是我老公.
他老公那天星期五不在.
很有趣.
跟我說很開心有人做了.
我繼續寫三次感恩.
用這本簿.
這本簿有沒有似曾相識?.
見過類似.
這是我自己貼的前面.
一本裡面是白色的.
前面有我寫的東西.
是什麼來的?.
是靈修的札記簿.
有沒有人在寫靈修札記簿?.
有.
在哪裡?.
很大聲.
前面.
很特別.
寫靈修札記和不寫靈修札記的靈修.
有沒有分別?.
寫的人覺得.
很有分別.
除了有寫和沒有寫之外.
你寫的時候不知為何.
有多些東西可以想到.
很具體.
甚至有神的說話.
透過你寫的時候.
突然間多了些東西.
有沒有人試過?.
試一下好不好?.
給大家一個挑戰.
我今年用這兩本簿.

$^{521}$這本是在哪裡買的?.
MUJI買的.
大約六元一本.
這本是朋友送的.
他很老實跟我說.
他是淘寶.
其實不是我們用什麼簿.
是我們的心預備了.
神在我們的生命裡面對我們說話.
待會我想和大家分享.
我的靈修.
我們看看這幅圖.
這個靈修.
有沒有看到這個像?.
這個像叫Enations of Loyola.
是伊納爵羅耀拉.
他在1491年出生.
是西班牙人.
不多說他的故事.
看到這幅像.
向前看.
其實他是耶穌會的創辦人.
更加有一個伊納爵式的靈修.
叫Enation Spirituality.
就是他的靈修.
我們說靈修.
好像是內在的操練.
修煉內在生命的東西.
這個是真的.
另外靈修的果實是向外的.
指向我們門徒的生命.
我們如何服侍向外.
傳福音並且服侍倫社.
所以我很喜歡這個像.
和裡面的精神.
在我的神學學習裡面.
我第一間的神學院.
就是中國宣道神學院.
裡面他們叫中國宣道.
當然是中國宣教.

$^{561}$裡面會結合這套靈修操練.
所以我學靈修.
是在20年前.
2005年開始學這個靈修.
但是當時的我.
是怎樣的呢 靈修.
上靈修堂.
老師叫你安靜的時候.
會怎樣呢.
黑眼睛睡.
很自然就會睡著.
因為很不習慣.
我很喜歡熱情的敬拜.
你也看到我很多時候.
講道之後的回應詩歌.
都比較要澎湃一點.
因為我都是喜歡這一類的.
我是喜歡快的詩歌.
我是一個這樣的人.
要去這套靈修.
是要你安靜.
都有安靜.
是要安靜的.
甚至乎不鼓吹唱詩歌.
鼓吹你讀聖經.
用經文親自對你說話.
更有不同的靈修內容.
可能是大自然.
良心的審查.
即是一種回顧.
生命的回顧.
各方面.
我覺得很奇妙.
神帶領我去學這個靈修.
但是我真的學得不好.
讀了四年神學.
學得不好.
甚至乎我之後.
沒有再用這套靈修操練.
雖然我第一間事奉教會.

$^{601}$我是沒有教過靈修.
但來到修浸.
杜牧師識貨了.
他知道我讀中宣.
但不知道我靈修是這樣的.
他說你讀中宣.
你教我靈修吧.
杜牧師就叫我教靈修.
我答應了.
為何你連靈修都教.
經過一個流淚谷的日子.
進來修浸前一個流淚谷的日子.
神竟然帶我去靈修操練.
是這一套靈修.
我再去信義中神學院.
讀靈修.
讀的時候又遇上.
以前中國宣導神學院的老師.
去了信義中教.
這個老師讀完神學的博士學位.
在美國再回來教.
老師教真的和以前不同.
學生和以前的我都不同.
以前是拿B減.
再讀的時候拿A減.
信心都回來了.
但不是這樣.
是進度到達.
人生的生命到達.
去經驗.
就是這套靈修.
我進到神的說話裡.
將我從流淚谷裡拉上來.
人生命快要枯乾.
但竟然有水支援我的生命.
再次來到去事奉神.
所以很感恩可以去靈修.
更加可以和弟兄姊妹一起分享靈修.
在過去教了三季靈修進心之旅的時候.
在課堂裡.

$^{641}$很感恩有弟兄姊妹在職場經驗到神和祂說話.
最近一次.
我的姊妹.
我都有點害怕.
她哭了.
安靜到哭了.
但她進到一個場景裡.
經文的場景裡.
耶穌敲語.
她在默想的時候哭了.
我很擔心她會否了解到.
她生命裡隱藏了需要去醫治的東西.
後來再分享.
她說見到耶穌的樣子.
耶穌的樣子.
那個愛令她很動容和感動.
這一幕我也覺得很驚訝.
在短短幾分鐘安靜的時候.
有這麼大的神的恩典在她生命裡.
也有一次在第二班裡.
有一個不會畫畫的婦女.
她畫了火柴人的公仔.
畫了兩個女孩.
太陽那些.
好像幼稚園學生畫的畫.
我們之後要交換畫.
抽號碼.
有個弟兄拿了一幅畫.
令我感動.
因為他有兩個女兒.
那幅圖畫也是神祝福她.
在女兒的關係上.
如何經歷神恩典的部分.
我覺得勝靈也在這裡.
原來靈修去到神的面前.
聖靈真的在這裡.
我們真的可以經歷到神的同在.
很感恩神教我要去親近祂.
並且要將人帶到祂的面前.
我很感恩神叫我去不同的地方.

$^{681}$我又願意去學習.
我也喜歡學習.
就學了一個禱讀.
大家有禱讀過吧.
禱讀是一段經文讀完一兩次.
然後用經文來祈禱.
大概是這樣.
或者延伸我們的禱告.
我再去不同的屬靈群體.
去學習禱讀的時候.
台灣那邊.
新店行道會.
張茂松牧師教的禱讀是.
道道 道道 道道 道道.
這是他的口訣.
道道 道道 道道.
他要先讀.
讀完再禱讀.
禱讀完再讀.
是互換的.
不要只是讀完一次就去禱讀.
不要漏了經文.
經文是不斷的.
頗像歷代的靈修操練.
延伸下來的更新版本.
有一段經文.
看回第一篇詩篇.
我示範一次給你看.
應該是怎樣.
道道 道道 道道.
我用一個經文去示範.
沒有稿的.
純粹很自然的.
你也可以做到.
為喜愛耶和華的律法.
晝夜思想.
這人便為有福.
為喜愛耶和華的律法.
晝夜思想.
這人便為有福.

$^{721}$為喜愛耶和華的律法.
晝夜思想.
這人便為有福.
Amen.
大家做到嗎?.
三次.
再讀的時候.
道道 道道.
為喜愛耶和華的律法.
為喜愛.
喜愛.
喜愛耶和華的律法.
是呀主.
我要喜愛你的律法.
你幫助我.
去喜愛你的律法.
主呀 我常常忘記你.
常常忘記你.
但是你說.
要我喜愛你.
喜愛你的律法.
晝夜思想.
主呀 幫助我.
晝夜思想.
Day and night.
晝夜思想.
日與夜.
去經歷你.
經歷你的律法的美好.
在我生命裡.
主呀 我要有福氣.
因為你定義.
我領受從你而來的祝福.
Amen.
這個是隨意的禱告.
剛才我沒有說.
是你的心.
你的心有多麼想.
你的口便會出來.
我們常常說不懂得說.

$^{761}$不懂得祈禱.
就是我們沒有心.
我們的心想的時候.
口便會出來.
好像平時我們說.
口控制不了.
是因為我們的心.
控制不了.
當我們的心去到主的面前.
我們的口便會懂得祈禱.
祈禱的恩賜.
我覺得.
是每一個都可以的.
每一個有心.
不要看自己.
心去到神的面前.
我就學慕主.
我就要去神的面前.
你的口便會出來.
就會說到.
神要你說的話.
甚至是祝福的說話.
不單是祝福你祈禱的對象.
自己的生命.
都可以這樣去領受祝福.
我們的靈修.
是和神建立關係.
更加是.
個人化.
神的話語是要祝福.
個人化我們的學想.
去經歷神的大能.
去更新他的工作.
有弟兄姊妹跟我說.
這兩年.
見到你進步了.
有信心了.
我沒有理會.
這些東西.
其實我一直都是.

$^{801}$主在你調教我的生命.
用了.
去年.
完成了十個月.
跟了一個靈修導師.
他不是教我.
是我的生命和他講.
去揭示.
更加和我同行.
剛才說伊臘靈修.
有一套神操.
是三十天的神操.
但我不是靜修.
困了三十天的那種去做.
是信徒.
在生活上的經驗.
靈修導師同行才做到.
在這個靈修旅程裡.
有很多破碎.
繼續破碎.
繼續去調教自己.
跟神真實的對話.
我不是說.
以前跟神學院老師.
靈修睡覺.
不是,現在閉上眼.
安靜下陪伴的時候.
神真的會提示.
見到圖畫.
他怎樣去愛我.
怎樣在經文裡.
進入到經文場景裡.
去和耶穌相遇.
這是靈修.
靈修當下.
與耶穌相遇.
靈修當下.
領受神話語.
靈修當下.
生命被更新.

$^{841}$這是我繼續追求的靈修.
在耶穌進入到.
一個村莊裡.
他見到一個女人.
就是瑪大.
瑪大見到耶穌的時候.
就接耶穌去自己家裡.
去一個妹妹叫瑪利亞.
瑪利亞當時.
坐在耶穌的前面.
腳前.
去聽耶穌講道.
瑪大要服侍耶穌.
貴賓來到.
去張羅.
去打點.
弄東西吃.
他很忙.
他就和妹妹講.
瑪利亞你為何坐在這裡不幫忙.
是不是這樣.
耶穌回答瑪大.
和瑪大講.
他講什麼呢.
耶穌回答說.
瑪大瑪大.
是思慮煩擾.
但是不可少的.
只有一件.
瑪利亞已經選擇.
上好的福分.
是不能奪去的.
我們很熟悉.
我們會不會漏了去領受呢.
就是耶穌說.
但是不可少的.
只有一件.
你這麼多事去忙.
最重要的那件.
你沒有做.

$^{881}$瑪利亞做了.
去到耶穌的腳前.
去聽他講道.
這個是最重要的.
但是沒有做.
做了其他你覺得很重要.
但是耶穌說.
不是對你的生命最重要.
很大的當頭棒樂.
對瑪大來說.
同樣.
對我們的生命來說.
會不會都是很重要的提醒.
我們很多事去忙.
很多事去做.
只有一件.
去到耶穌的腳前.
我最近.
不是最近.
這兩個月都是很多事做.
排隊.
做完一樣就一樣.
真的沒有停過.
我們教會很多會開.
開會要時間.
和人交接.
要時間.
很不開心的日子.
也要時間.
我區裡.
走了幾個老人家.
也要時間.
做事有感情.
也會愛上.
去醫院奔波要時間.
那裡很少探.
要去探探人的時間.
講道.
要預備講章.
很難的.

$^{921}$有時時間是預不到.
我今次的講章.
沒什麼時間.
也有時間去做.
不斷.
操練日和夜.
我今年要操練日和夜.
早上一定.
幾分鐘也好一定要.
經文多少也好.
一定要.
讀一節經文.
也要去.
領受神的說話.
很寶貴.
Night.
一定要今早的經文.
再去refresh.
再寫感恩試項.
很難的.
我試過第二天才寫感恩試項.
但每天都在寫.
我們在這裡不斷有很多事.
很多事.
我一樣.
我們一起去面對.
去依靠主.
去經歷神.
去幫助我們.
和我們同行的那種關係.
再看這幅圖.
就是.
我們要好像一棵樹.
栽在溪水旁.
這個果子.
是我加上去的.
用canva.
我沒有數錯.
九粒果子.
我們按時候.

$^{961}$結果指.
葉子也不枯乾.
凡他所作的.
盡都順利.
我們不知道這棵樹.
是否那麼堅壯.
甚至有些橫枝被砍斷.
在《約伯記》有段經文很棒.
不知道大家有沒有留意.
十四章七至九節.
那裡說到.
樹若被砍下.
還有子望.
去發芽.
亂枝生長會不息.
即是已經砍斷樹.
根雖然衰老.
在地裡.
幹都死在土中.
及至.
有水氣.
滋潤這棵.
割下來的樹.
還要發芽.
並且長出枝條.
會像.
新栽的樹一樣.
我搜尋了一下.
我相信.
我對這些植物的東西.
大自然的植物都不太認識.
所以再搜尋一下.
印證到.
大家搜尋一下.
有AI的模式.
馬上出現.
原來棵樹砍了之後.
還有根.
這個經文說.
雖然衰老.

$^{1001}$但還有根在土.
但有些水滋潤的時候.
會再生長.
會再發芽.
會再長出枝條.
更加可以.
重新去做一棵樹.
可以再結果子.
即是說我們的生命.
有多難呢.
還有根.
根在哪裡?.
我們還是連結於.
耶穌基督的時候.
我們只欠水.
水從哪裡來?.
神的話語.
就敲罐下來.
來滋潤我們這棵樹根.
以致我們.
不斷可以再生長.
不知道我們生命裡.
現在的狀況是怎樣.
你覺得.
這棵樹已經被刺到.
不同的環枝.
好像不是一棵正常的樹.
在生命的樹裡.
或者.
在你的工作裡.
覺得很.
沒有盼望.
甚至覺得.
隨時會失業.
家庭的關係.
好像沒有什麼喜悅.
可能你的婚姻.
都會覺得沒有滋潤.
沒有愛情滋潤.
都好像已經失落了.

$^{1041}$不敢和別人說.
我和我家人的關係.
我和我老婆和我丈夫的關係.
不敢和別人說.
甚至你的呼召.
你的召命.
都停滯不前.
沒有亮光.
我們需要.
主的金能.
祂的話語再鼓勵我們.
提升我們.
我們一起有個禱告.
靈姐妹不知道今天的.
經文.
信息裡面.
有沒有提醒你.
你這棵生命樹的狀況.
是怎樣.
主要知道.
這一刻的你.
不是等於永遠的你.
我們的生命.
不可以被滋潤.
\newpage



\section{路加福音 1:13}
\label{sec:031622}
\textbf{聖經人聖經事:祈禱已被聽見}~~~~黃克勤牧師博士
\newline
\newline
連結: \href{http://www.church.com.hk/acms/content.asp?site=cdc\&op=show\&type=product\&code=031622\&layout=sermon}{\texttt{031622}} ~~~~ 語音日期: 2026-01-30
\newline
\newline
\hyperref[sec:031621]{\small{< < < PREV SERMON < < <}}
~
\hyperref[sec:index_chronic]{\small{[返順時目]}}
~
\hyperref[sec:031623]{\small{> > > NEXT SERMON > > >}}
\newline
\newline
路加福音 1:13
\newline
\begin{longtable}{cl}
\hline
\hline
章節 & 經文 (和合本修訂版)\\
\hline
1:13 & \begin{tabularx}{0.7\textwidth}{X} 天使對他說：「撒迦利亞，不要害怕，因為你的祈禱已經被聽見了。你的妻子伊利莎白要給你生一個兒子，你要給他起名叫約翰。 \end{tabularx} \\ \\
[1ex]
\hline
\hline
\end{longtable}
$^{1}$弟兄姊妹又讀聖經了.
「祈禱以備聽見」.
生命說書的網絡提供的聖經人聖經故事.
「祈禱以備聽見」.
是什麼來的?.
仍然是施洗約翰的故事.
施洗約翰的故事是非常豐富的.
路加福音的記載.
我們先看約翰福音.
然後看馬太福音.
現在看路加福音.
路加福音一章五節.
「當猶太王希律的時候.
阿比亞班裡有一個祭司.
名叫撒迦利亞.
他妻子是阿倫的後人.
名叫以利沙白」.
先解釋一下.
再提一次.
綠色字是注釋.
黑色字紅色字.
就是和合本原來的文字.
紅色字是想爭取你的眼球.
妻子是阿倫的後人.
名叫以利沙白.
「他們二人在神面前都是二人來的.
遵行主的一切誡命禮儀.
沒有可指責的.
不過他們沒有孩子.
因為以利沙白畢生育兩個人年紀又老邁了.
撒迦利亞按班次在神面前.
供祭司的積罰」.
撒迦利亞是一個祭司.
是阿比亞班的班次.
裡面其中一位是祭司.
「祭司那麼多位.
於是他們按照祭司的規矩祭籤.
抽籤.
看誰可以入主的殿.
笑向祈禱」.

$^{41}$入主的殿笑向祈禱是這樣的.
當猶太王希律的時候.
希律是甚麼意思呢?.
希律是英勇的.
有很多班祭司.
其中一班是阿比亞班.
阿比亞的意思是.
「我父親是耶和華」.
裡面有一個祭司.
這個祭司叫撒迦利亞.
撒迦利亞是被耶和華紀念的意思.
是一個好名字.
撒迦利亞是被耶和華紀念的意思.
他的妻子是亞倫的後人.
亞倫是甚麼意思呢?.
亞倫是帶來光的人.
撒迦利亞的丈夫是祭司.
太太以利沙白也是亞倫的後人.
所以兩個都是亞倫後人才可以做到祭司.
他的妻子也是亞倫後人.
叫做以利沙白.
以利沙白就是我們現在很多人用的英文名.
以利沙白是神的誓言的意思.
神發誓.
神指著自己起誓.
是一個很嚴肅的名字.
不過是一個很好的名字.
是一個很高的名字.
所以如果你有個女孩要改名.
以利沙白也是一個好名字.
不差的.
亞比亞的祭司裡面有一個叫撒迦利亞的祭司.
兩個人都是義人.
甚麼是義人呢?.
聖經說了很多次.
聖經裡面沒有好人壞人.
是比較多寫義人和惡人.
他們兩個是好人.
這個義人特別是指甚麼呢?.
就是他們是遵行一切誡命禮儀.

$^{81}$而且無可指責,你做不足.
做得好足.
即是說做得好足就要盡心盡力.
出力,不單有心還要出力.
譬如返崇拜無可指責.
風雨不改.
有主日崇拜必定祖道.
風雨不改病都動動下之棍黎.
這真是很好,無可指責.
十務封求,有開崇拜我都黎.
嘩,好厲害.
無可指責,大家明白意思嗎?.
即是說這個人做一切誡命律,跟著律法做.
沒有可指責之處.
真是神面前是一個義人.
神都祝福他,只不過他沒有孩子.
當日是一個好肉孫.
是覺得上帝不祝福他的表現.
沒有孩子,因為以利沙白不能生育.
你怎知是以利沙白問題還是撒迦利亞問題?.
如果今日醫學上會這樣問.
不過我們不問.
因為聖經是說以利沙白不生育.
兩個人又年紀老邁,都年紀大了.
更加生不了.
撒迦利亞按班次在神面前供祭司祭佛.
按班次進去,祭司會在香壇面前燒香.
把香料放進去禱告,香爐也要禱告.
還要整好燈台,燈台像我的手這麼高.
有七盞燈,要加油,要剪好燈芯.
很多事情要做.
所以要在神面前供祭司祭佛.
而撒迦利亞剛剛好按照規矩祭籤.
抽籤到他可以進去在聖殿燒香.
這個就是香壇.
壇不是香爐.
香爐是圓形的香壇.
香壇高度大約三尺高.
高度剛好就是我們跪下祈禱的高度.
象徵著跪下祈禱.

$^{121}$這是寓意的解釋.
壇有四個角,畫得很好.
一定要有這些,要把它抬走.
是鋪金的.
上面放一些特製香料.
某個祭司家族做的香料.
主要是乳香.
黃金乳香末藥,乳香是祭司用的.
放進去燒,燒煙出來象徵禱告.
這件衣服是大祭司.
剛剛在萬寺前燒香.
當時燒香的時候,眾百姓在外面都很好.
祭司進去燒香,他們在外面祈禱.
很多時候在聖殿外.
不時有很多人聚集.
如果去耶路撒冷,一定去聖殿.
很想在聖殿前燒香,因為不能進去聖所.
所以在外面燒香,燒香的時候就禱告.
當燒香的時候.
他就去祭司即任.
眾百姓在外面祈禱,真是一個很好的環境.
有主的使者站在香壇的右邊.
香壇,萬寺,至聖所在裡面.
藥鬼在裡面,萬寺,半腰高的香壇在裡面.
主的使者站在香壇的右邊,即是外面的位置.
向他顯現.
撒迦利亞燒了香,旁邊多了一個人,有一個天使,他很害怕.
不要說天使,是人的害怕.
撒迦利亞燒了香,旁邊多了一個人.
撒迦利亞看見天使.
通常他們見到天使都很害怕.
第一,怕時間到了,天使帶我走了.
又或者他出現會不會審判我,上帝叫他來.
會不會有天使出現,會不會帶訊息給我.
比較少人這樣想,因為我很小的一個人物.
怎麼會審判一個天使跟我說話呢?.
所以有天使出現,他們多數是吉是空,不知道,很害怕.
驚慌害怕,天使一定這樣說.
不要害怕,還叫他的名字.
因為你的祈禱已經被聽見了.

$^{161}$第三點,你的祈禱已經被聽見了.
大家留意,他是一個義人.
這個義人按公職進去燒香.
燒香的時候,見到主的使者站在旁邊.
要說的是你的祈禱已經被聽見了.
這句話真的很安慰.
就是神說你的祈禱我聽見了.
你一定要相信這件事.
就是你的祈禱上帝是聽見的.
不是就這樣算的.
以前我小時候有一首英文詩歌.
不記得那個音是什麼了.
不過它就說.
Nice are never busy.
打電話給上帝.
哎呀,現在線路很繁忙.
請遲些再打個來吧.
以前聽這些就很煩躁.
現在線路繁忙,請遲些再打來.
不知道什麼時候打.
現在那些按一二三四.
什麼按一,什麼按二,什麼按三.
現在線路很繁忙,請等候.
上帝的線路沒有繁忙.
他是聽見的.
你的祈禱已經聽見了.
上帝特別祝福你的妻子以利沙白.
要給你生一個兒子.
上帝的名字也改了.
要叫做約翰.
這個就是施洗約翰了.
施洗約翰的父親是撒迦利亞.
母親是以利沙白.
以利沙白是馬利亞主耶穌母親的表姐.
因此施洗約翰是主耶穌的表哥.
就這樣.
你的祈禱已經被聽見了.
很安慰,已經被聽見了.
所以以利沙白要給你生一個兒子.
你叫他做約翰,名字也改了.

$^{201}$今天要跟大家講的就是這次聖經.
這句.
撒迦利亞,不要害怕.
因為你的祈禱已經被聽見了.
你的妻子以利沙白要給你生一個兒子.
叫他做約翰.
特別今天的金句只是紅色的.
撒迦利亞.
叫他的名字,或者叫你的名字.
陳大文,或者王克勤,叫你的名字.
不要害怕.
不要害怕,為什麼?.
哇,我的神啊,你對我講說話.
跟我說我又害怕.
其實很有趣,究竟他自己想怎樣?.
其實不要害怕.
你的祈禱已經被聽見了.
記得,你的禱告是有功效的.
是有聽到的.
他還沒回應我.
沒回應我的時候.
是時間,他的時間還沒到.
所以他沒回應我.
很有趣.
剛剛跟一個同事聊天.
那幾個小孩經常問.
要做什麼?.
他們大人會說.
一會兒吧,一會兒吧.
為什麼經常要一會兒呢?.
為什麼要一會兒呢?.
為什麼現在不行呢?.
我們心急.
一想要就是現在.
但其實神已經聽我們的祈禱.
祂有祂的時間.
祂就會向你顯現.
站在你身邊告訴你,不用怕.
你的禱告已經被應運了.
聽到了.

$^{241}$不過,另一件事,補充一下.
就是.
應運你禱告的時間與方法.
不是我們想的一樣.
祂不跟著我們的計劃去做.
祂有祂的計劃.
很想要一個兒子.
可能神沒有給.
給什麼?.
我給一個養子給你,可以嗎?.
或者我給你另一樣東西.
那樣可能更加好的.
你是不知道的.
很多時候都是這樣.
要記得你的祈禱已經被聽見了.
今日的金句.
你的祈禱已經被聽見了.
不要害怕.
我們同心禱告.
多謝我們的主.
因為你樂意誰聽人祈禱.
不會覺得我們煩著.
你的善惡又不會繁忙.
你無所不知.
你全部聽到世界上的事.
你是知道萬人的心.
主是你知道的.
我們就是坦誠到你面前.
獻給禱告.
將自己所需要的告訴你.
將我們的重擔憂慮.
射給耶和華.
因為你是愛我們到底的.
你是會背負我們的.
你是天天背負我們重擔的主.
願你被榮耀.
願你被高舉.
願你被信靠.
我們信靠等候你.
因為你已經聽我們的祈禱.

$^{281}$多謝主.
祈禱奉耶穌基督的名求.
阿們.
\newpage



\section{箴言 3:1-12}
\label{sec:031623}
\textbf{不輕看管教,不推辭施與}~~~~蕭壽華牧師
\newline
\newline
連結: \href{http://www.church.com.hk/acms/content.asp?site=cdc\&op=show\&type=product\&code=031623\&layout=sermon}{\texttt{031623}} ~~~~ 語音日期: 2025-07-27
\newline
\newline
\hyperref[sec:031622]{\small{< < < PREV SERMON < < <}}
~
\hyperref[sec:index_chronic]{\small{[返順時目]}}
~
\hyperref[sec:031624]{\small{> > > NEXT SERMON > > >}}
\newline
\newline
箴言 3:1-12
\newline
\begin{longtable}{cl}
\hline
\hline
章節 & 經文 (和合本修訂版)\\
\hline
3:1 & \begin{tabularx}{0.7\textwidth}{X} 我兒啊，不要忘記我的教誨，你的心要謹守我的命令， \end{tabularx} \\ \\ \relax
3:2 & \begin{tabularx}{0.7\textwidth}{X} 因為它們必加給你長久的日子，生命的年數與平安。 \end{tabularx} \\ \\ \relax
3:3 & \begin{tabularx}{0.7\textwidth}{X} 不可使慈愛和誠信離開你，要繫在你頸項上，刻在你心版上。 \end{tabularx} \\ \\ \relax
3:4 & \begin{tabularx}{0.7\textwidth}{X} 這樣，你必在神和世人眼前蒙恩惠，有美好的見識。 \end{tabularx} \\ \\ \relax
3:5 & \begin{tabularx}{0.7\textwidth}{X} 你要專心仰賴耶和華，不可倚靠自己的聰明， \end{tabularx} \\ \\ \relax
3:6 & \begin{tabularx}{0.7\textwidth}{X} 在你一切所行的路上都要認定他，他必使你的道路平直。 \end{tabularx} \\ \\ \relax
3:7 & \begin{tabularx}{0.7\textwidth}{X} 不要自以為有智慧；要敬畏耶和華，遠離惡事。 \end{tabularx} \\ \\ \relax
3:8 & \begin{tabularx}{0.7\textwidth}{X} 這便醫治你的肉體，滋潤你的百骨。 \end{tabularx} \\ \\ \relax
3:9 & \begin{tabularx}{0.7\textwidth}{X} 你要以財物和一切初熟的土產尊崇耶和華， \end{tabularx} \\ \\ \relax
3:10 & \begin{tabularx}{0.7\textwidth}{X} 這樣，你的倉庫必充滿有餘，你的酒池有新酒盈溢。 \end{tabularx} \\ \\ \relax
3:11 & \begin{tabularx}{0.7\textwidth}{X} 我兒啊，不可輕看耶和華的管教，也不可厭煩他的責備， \end{tabularx} \\ \\ \relax
3:12 & \begin{tabularx}{0.7\textwidth}{X} 因為耶和華所愛的，他必責備，正如父親責備所喜愛的兒子。 \end{tabularx} \\ \\
[1ex]
\hline
\hline
\end{longtable}
$^{1}$《錄音》或《錄影》的版權 是屬於基督教宣導會北角堂所有.
我們的屬靈傳統之一 是以真理教導為本.
歡迎你將此《錄音》或《錄影》的連結 傳給親友.
懇請注意 未經授權 請勿轉錄或剪輯本《錄音》或《錄影》.
甚願好像《仕途行傳》二十章三十二節下所說.
「這道能建立你們 叫你們和一切誠誠的人 同得基業」.
各位弟兄姊妹平安.
我想問一下 有哪位弟兄姊妹上星期日在這個禮堂聚會的?.
上星期日我是答應了去一間教會講道的.
但是星期六就被他們呼召 去到他們教會.
我去到的時候是一個空的禮堂.
在空的禮堂是六部崇拜的講道.
預備明天在家中放給弟兄姊妹一起參與崇拜.
今天回來和弟兄姊妹一同在教會裡面敬拜.
更加感受到是恩典.
每一次能夠和你的弟兄姊妹一起敬拜上帝不是必然的.
所以我想請你和你身邊的弟兄姊妹一起來到這裡講道.
我請大家和我一起祈禱 我們同心 仰望著.
我們的主 我們的上帝.
你的寶座 納定在天上.
你的權柄 統管萬有.
主在你的權柄 帶著慈愛 帶著憐憫.
當我們在人生經歷不同的事情的時候.
無論是信還是逆.
我們心知道都在你的慈愛裡面.
因此我們到你面前來我們敬拜你.
承認你是我們生命的主宰.
我們願意信從你 信靠你.
我們很願意你自己透過你聖經的說話.
不斷勉勵我們的心 加添我們的力量.
叫我們人生的路上堅定的跟隨你.
天父啊 求你賜寬一段時間.
你遮蓋人的不足.
讓你自己的恩賢淋到我們眾人的心中.
我們仰望祈禱 奉主耶穌 成名祈求.
阿們.
這個地方是一個室外的園亭.
裡面有一個很特別的迷宮.
這個迷宮是用一些很矮的樹叢.
一路造成的.

$^{41}$小朋友和孫兒見到這個情況很開心.
馬上就衝進去.
要到處轉動 到處走.
我看到他們在樹叢下面一路走.
前面有些障礙.
但是發覺右邊有一條直路.
馬上衝去右邊.
但是我知道那條路其實是不通的.
因為我比較高.
我在那個迷宮上已經看了所有通道在哪裡.
我就大聲叫 不要走右邊 走左邊.
但是似乎他們不太肯聽.
走著走著的時候 他們發覺沒有出路.
但是他們想到一些方法.
在樹叢之間有些空位.
因為他就鑽進樹叢裡面.
從另一邊穿出來.
我在旁邊看到的時候覺得很開心.
也體會到我們真的有很短小的眼光.
看我們人生的路.
我們很需要上帝給我們智慧.
能夠看到人生的路應該怎麼走.
我們這段時間讀到《箴言》.
第一章第一第二節一開始就說.
所羅門為什麼會寫《箴言》呢?.
他說這個目的是叫人曉得智慧.
是能夠領受智慧.
領受神的說話 以致可以領受智慧.
《箴言》說到智慧的時候.
不單單說一些知識學問.
不單說一些分析事情的能力.
是更加說一種心靈的洞悉.
看透了事情的根本.
以致你能夠知道哪些是重要 哪些是次要.
而作出適當的抉擇.
《箴言》說到神要將一種屬天的智慧.
賜給敬畏祂的人.
人生的路因此可以走在祂的裡面.
走在光明的路上.
所以你記得上個星期.

$^{81}$當我們讀到第二章的時候.
《箴言》就提到.
智謀要庇護你 聰明必保護你 救你脫離惡人的道.
人生的路你會遇到很多惡人的道.
你不察覺你走了在上面.
你自己就受害.
但請告訴我們知道.
智慧要保護我們 救我們脫離惡人的道.
今天讀到聖經三章就提到.
在你一切所行的路上都要認定祂.
祂必使你的道路平直.
你人生的路會平直.
會能夠走在光明正直的路上.
因為你認定祂.
各位電影姐妹 我們今天思想短經文.
在你的心裡面你可以默默地祈禱.
求神賜給你生命的智慧.
以致你可以在人生的路上.
不會走很多冤枉路 或者很多錯路.
而你一路能夠行在正直平安的路上.
我們看聖經.
在聖經剛才我們讀的第三章裡面.
一開始的時候.
你會發覺聖經裡面提醒我們一件事情.
三章第一第二節.
我現在不要忘記我的教誨.
你的心要謹守我的命令.
聖經很清楚地講述.
智慧從哪裡來.
智慧是從神的說話而來.
因此 真言裡面不斷提醒我們.
要牢記神的說話.
要將神的話藏在心裡.
我將你的話藏在心裡.
就不至於跌倒.
就不至於走錯路.
當你讀到這段經文的時候.
或者你心裡會想.
有些時候我都記不住.
年長一點的更加體會.

$^{121}$剛剛今早讀完.
就問你們今早讀了什麼.
忘記了.
但你放心 你繼續讀.
你繼續重溫.
你繼續不斷去認識.
你願意領受神的心.
領受神話語的心.
神是知道的.
你只管讀.
有些時候可能是忘記.
但你放心.
當面對一些處境.
你會發覺神會幫你的記憶.
想起那些說話.
這是我的經歷.
很多時候有些特別的處境.
突然間有神的說話.
在那些時刻浮現出來.
是一些我都忘記了的說話.
但那個時候出現了.
神透過他的說話.
給我們提醒.
也引導我們做一些決定.
你只管放心.
繼續將神的話放在你的心裡.
正如聖經所說.
謹守我的命令.
意味著.
你的內心是願意持守著神的話.
不會淡忘.
不會輕視.
然後聖經應許.
因為他們必給你長久的日子.
你生命的年數和平安.
是神應許我們的.
然後再去第三第四節的時候.
聖經進一步說.
不要讓恃愛和誠信離開你.
要好好地緊握著它.

$^{161}$繫在你的頸項上.
刻在你的心板上.
如果你將神的話放在心中.
你會發覺在你的生活行為上.
你自然有一個方向.
就是我不會讓恃愛和誠信離開我.
尊重神說話的人.
在他的生活行為.
會常常謹記.
怎樣讓我一切所做的事情.
是出於愛.
出於誠信.
聖經很多時候用這兩個字來講述.
他對我們立約的一種關係.
當他和人立約的時候.
他就說他以恃愛和誠信.
與我們立約.
聖經亦作信實.
神和我們立約的時候表示.
他不會以恃愛和誠信.
讓這些離開我們.
因此你和我跟隨神的人.
我們都當按神怎樣待我們的方式.
去待人對神.
因此聖經說.
你要小心.
讓恃愛和誠信不會離開你.
緊緊的在你的生命裡.
而聖經跟著應許.
只要你必在神和世人眼前蒙恩惟.
神喜悅你.
人都會喜悅你.
因為你領受神給你美好的見識.
第三第四節之後.
很快就來到一段.
我們會很熟悉的經文.
大家知道是今個月的金句.
請告訴我們.
你要專心仰賴耶和華.
不可依靠自己的聰明.

$^{201}$你要專心仰賴耶和華.
這個仰賴.
不單止是一種.
頭腦上的一種心意.
覺得我都好.
我都想來信靠.
依靠上帝.
請提到是要專心.
專心含有一種全心全意的意思.
所以你看到英文的翻譯.
Trust in the Lord with all your heart.
你全心全意來信靠祂.
意味著就算你付上代價.
意味著你可能會因此有些損失.
你仍然去信靠祂.
專心仰賴耶和華.
是一種對不信.
對很多不同的試探的解約.
在頂智慧我們個人的經歷裡面.
你都會體會.
有些時候我們生活遇到一些挑戰.
有些突然的難處.
在那個時候可能會懷疑神.
如果你是愛人的.
為什麼你會讓我經歷這些事呢?.
有時會懷疑神的信實.
懷疑祂的說話是否那麼可靠.
因此那些時候會出現一些試探.
什麼試探?.
轉移依靠自己的聰明.
當然神給我們聰明.
我們用我們聰明行事不是犯罪.
但是何時你依靠你的聰明.
或者說你信任自己的聰明.
過於神自己的說話.
在那個時候你發覺內心出現了驕傲.
而你發覺驕傲是很多罪.
很多惡行的根源.
當你內心自以為是的時候.
你就會走自己的路.

$^{241}$利器神的路.
跌入很多罪惡的圈套裡面.
聖經告訴我們.
專心仰賴耶和華.
過程當中留心.
不要依靠自己的聰明.
我記得一件事情.
就是好幾年之前.
有一個家庭聚會.
有一群電影節目在弟兄的家裡.
一起有聚會.
分享一首詩歌.
這位弟兄就帶我來唱一首詩歌.
這首詩歌我很深刻.
可能很多電影節目都會唱.
這首詩歌叫做.
我心靈得安寧.
It is well with my soul.
我特別打了一些歌詞出來.
字粒小一點希望你們看到.
這首詩歌是一位特別的作者.
這位作者叫做.
Heratio Spetford.
斯伯福弟兄.
這首詩歌其實很多年了.
150年多以前所作的詩歌.
但是這麼多年來.
歷世歷代.
千千萬萬的弟兄姊妹.
因為這首詩歌得著鼓勵.
得著激勵.
這首詩歌講到一位作者斯伯福律師.
他本身是一位很敬虔的弟兄.
在很多年之前.
在美國當地.
如果大家有些印象的話.
在大概1870年代的時候.
芝加哥發生一個很嚴重的大火.
整個城市裡面.
差不多有三分之一的建築物被燒毀.

$^{281}$在那段時間.
斯伯福弟兄.
他自己本身有很多不同的物業.
他都是投資很多物業.
但是那段時間絕大多數被燒毀了.
在那段時間.
他的兒子亦因為患病而喪生.
一段很低沉的時間.
但是他仍然堅定信靠神.
因此他自己不斷努力.
重建一些房屋.
接待了差不多十萬的人.
在那個時刻找到地方居住.
過了一兩年之後.
他和他太太想休息.
因此一家.
當時他仍然有四個女兒.
坐一輛郵輪去歐洲的地方.
因為工作太重.
所以他說你們先去吧.
太太和四個女兒一起坐郵輪.
但郵輪在大海中間.
大西洋的中間.
突然被一艘鐵甲船橫刀撞過去.
這艘船就沉沒了.
二百二十四個乘客喪生.
剩下幾十個.
這件事發生在鐵達尼號沉船之前.
在那段時間裡.
太太獲救.
但四個女兒全部喪生.
我不知道如果你是一位弟兄.
你會有什麼感覺.
很多人問.
為何神會這樣待你.
一波未平一波又起.
為何神那麼殘忍.
一個兒子死了.
然後你四個女兒.
都是幼年的女兒.

$^{321}$全部奪走了他們的命.
這位弟兄很快就坐船去歐洲.
想來安慰太太.
船上發現有很多遭難者的家人.
去到沉船的地方.
他們特別停留.
懷念他們的家人.
師彼福弟兄在那時候.
他心裡有一個很大的感動.
告訴自己.
It is well with my soul.
這位上帝仍然是可信可靠.
因此他寫了這首詩歌.
這首詩歌我簡單看.
有時享平安如江河平又穩.
有些時候很感恩神給我們平穩的生活.
但有時想不到遇到悲傷.
好像波浪滾滾.
在大船上在海洋中間.
他們感受到波浪的那種空肚.
然後不論何環境.
我要蒙主引領.
我心靈得安寧.
一種信心帶來的.
一種意想不到的安寧.
然後這首詩歌說.
撒旦雖來侵.
眾試煉雖來臨.
我有主保證在我心.
基督已清楚我警方無人助.
就為我流寶血救贖我.
在那個時候.
這位弟兄想到的.
是他自己的舊恩.
正如聖經所說.
神既不愛惜他的兒子.
為我們捨了.
但豈不將萬物和他自己.
都賜給我們嗎?.
當你小心去想的時候.

$^{361}$你會明白.
這位上帝.
連他自己最愛的兒子.
都賜給我們.
讓他為我們承受罪的代價.
豈不是將一切美善的.
都賜給我們?.
甚麼能夠幫助施比福弟兄.
能夠面對這個困境.
是因為他知道自己無拯救.
因此他知道上帝的慈愛是堅定的.
不會轉變的.
因此在這些環境當中.
仍然專心信靠上帝.
因此他說.
It is well with my soul.
我很記得.
當我想起這首詩歌的時候.
我都想起這位弟兄.
剛才我說.
帶領我們家中.
來一起唱詩歌的弟兄.
這位弟兄.
講述他自己個人的困難.
我也不會想到一首詩歌.
一百五十年之後.
會安慰另一個時空的弟兄.
因為曾經有一位施比福弟兄.
堅定的信靠神.
也成為了無數人的鼓勵和祝福.
各位聽眾.
聖經勉勵我們.
你要專心仰賴耶和華.
但願神這樣來鼓勵你.
幫助你.
然後聖經再講.
在你一切所行的路上.
要認定他.
無論你走著什麼路程.
你記得要認定他.

$^{401}$聖經所說的認定有兩個含義.
第一.
就是你承認.
神的慈愛的恩手.
掌管著你人生的每一步.
無論發生什麼事情.
都在神的慈愛的恩手當中掌管.
你未必馬上明白.
但肯定是出於神的慈愛.
第二.
你認定他的時候.
英文的翻譯認定是acknowledge.
是很清楚的確認神的本性.
什麼本性?.
聖經清楚說.
神喜愛公義.
恆護罪惡.
因此你就按照他的本性.
或者說按照他的眼光.
去看你自己所經歷的事情.
按照他的本性.
按照他的心意.
作出人生不同的抉擇.
我要不要轉一份工?.
我要不要跟這位弟兄說這些話?.
很多時候.
有很多不同的決定的時候.
你回想.
耶穌如果是你.
在你的處境.
他會怎樣做呢?.
你記得那本書嗎?.
那本書告訴我們.
WWJD.
What would Jesus do?.
什麼時候.
在不同的處境的時候.
你回想.
他是上帝.
他是掌管我生命的主.

$^{441}$他會怎樣做呢?.
然後按照他會做的去做.
這個是認定他.
在你一切所行的事上.
認定他.
聖經就應許.
他必使你的道路平靜.
這裡所說的平靜.
不等於你沒有困難.
沒有風沒有浪.
很安全.
沒錯,你會有困難.
但你會知道這條路是平靜的路.
因為是神指示你走的路.
回說剛才的迷宮.
告訴他.
右邊有條直路.
走那麼遠.
旁邊左邊更快.
好像很辛苦走一條長路.
但其實那條路.
是帶你去出路.
重要的是.
有智慧.
有天上的智慧引導著你.
以致你能夠走在光明的路.
不會跌倒的路上.
你記得勝經就應許.
他必使你的道路平靜.
你放心跟隨主走上去.
我們再看.
當我們一路再讀的時候.
聖經裡面有很多對我們的勉勵.
但接著去到十一,十二節的時候.
聖經突然講到「管教」.
很奇怪.
沒錯,我們也體會.
有些時候我們心裡很願意.
遵從神的說話.
跟隨祂的旨意走.

$^{481}$但有時候我們也會察覺.
自己有些時候也忍不住.
做一些決定.
走了偏差的路.
你留意.
當你走偏差路的時候.
你會遇到「管教」.
突然有些想不到的患難.
突然想不到有很大的挫折出現.
當然我們面對「管教」.
不一定是因為我們的罪.
有些時候我們面對困難,患難.
是因為神要磨練我們.
但何時我們有一些不能夠轉變.
或者長期不肯放棄的惡行的時候.
你留心.
有些時候神會讓「管教」臨到你.
當「管教」臨到你的時候.
正如聖經所說.
有些時候我們會感到厭煩.
為什麼你會讓我遇到這樣的事情?.
但聖經提醒.
不要輕看.
不要輕看「管教」背後有神的意思.
不要厭煩祂的責備.
因為.
聖經接著解釋.
為什麼呢?大家一起讀.
「因為耶和華所愛的他必責備.
正如父親責備所喜愛的兒子」.
是因為神是愛我們的.
祂是出於對我們的愛.
是出於對我們的幫助.
是我們所需要的.
各位電影節目.
如果你心裡面.
有時候會為了所遭遇的一些困難.
覺得很厭煩的時候.
你安靜下來.
你想一想.

$^{521}$神有什麼心意.
或者只不過是神給你的磨練.
又或者可能是你體會到.
自己有一些惡行.
很需要你去轉會.
在過程當中.
小心反思神要你糾正什麼.
神要你怎樣去轉會.
然後你順著祂給你的引導.
走回正路上.
那些管教就發揮了美善的作用.
你一定記得.
我們以前也常常提及.
路易斯先生.
C.S. Lewis.
一位著名的基督徒作家.
他說到一件事情.
他說有時候你面對苦難.
你記得.
有些神正在跟你說話.
有一段話.
他這樣說.
他說神在我們歡樂的當中.
向我們輕輕說話.
告訴我們祂愛我們.
告訴我們.
你放心走你人生的路.
但祂卻在痛苦中向我們呼喊.
痛苦是祂的擴音器.
以此喚醒這耳聾的世界.
這個世界是在享樂.
在自我中心.
在一種強權得勢的情況下.
聽不到神的聲音.
神容許苦難來到.
讓人的心受到提醒.
在困難當中反思人生的時候.
就明白神的旨意.
因此他說痛苦是給予行惡者.
修正自己的唯一機會.

$^{561}$各位弟兄姊妹.
你可以問很多弟兄姊妹.
你身邊認識的.
很多都會告訴你.
在人生某些很低落的階段中.
反而是自己最親近神的時候.
因為有些時候心靈被提醒.
更知道去選擇正直的路.
來跟隨著.
我們再讀下去.
讀下去的時候.
去到三章二十七,二十八節的時候.
你發覺裡面有一個很重要的訊息帶出來.
就是神的同在和祂的親密.
比一切我們所擁有的更加重要.
有時候我們會以為得到那樣東西.
得到那個機會.
那個物質重要.
但其實你人生一路走的時候.
你發覺神的同在更重要.
聖經裡面提到.
你手若有行善的力量.
不可推遲要施於那應得的人.
聖經特別提應得的人.
我若手頭方便.
不可對鄰舍說.
去吧明天再來我必給你.
有時候我們都很願意施予給別人.
去輔助去幫助他一把.
但有時候我們都會有一種感覺.
當考慮可否給對方.
或者是我的時間.
我的關心或者我的金錢.
有時候裡面有一種.
好像不太捨得的感覺.
給了他.
我自己有點損失.
有時候內心有一種掙扎.
就是我付出會帶給我損失.
所以我還是不要了.

$^{601}$或者晚點再說.
但清秀講述到一種智慧.
這種智慧面對著一些機會的時候.
有行善的力量就不推遲.
意味著他把握著現在仍然有的資源.
我還有健康.
我還可以走得走得.
我可以幫到他.
如果他現在有那些資源.
他不會去拖延.
不會告訴鄰舍.
明天回來再說.
明天回來不知道會不會再給.
因為他在那個時候.
他會把握著遇到的機會.
現在就去施予.
一個人當他內心有一種明白.
明白神的掌權的時候.
他就會不推遲不拖延.
因為知道在施予當中有神的應許.
在真言書裡面有很多關於.
神應許施予的人的恩典.
不如我們多看一兩節的經文.
在聖經的十一章真言二十四十五節.
聖經說有施予的錢財損失.
我讀錯聖經了.
你施予你拿出去應該損失了一點.
但聖經不是這樣說.
有施者的錢財增添.
論識過度的反自窮乏.
慷慨待人的必然豐裕.
知潤人的念自己也得知潤.
聖經十九章十七節繼續說.
憐憫貧寒人的就是借給貧寒人.
不是,是借給誰?.
借給耶和華.
你有沒有想過.
當你去幫助別人的時候.
你是在幫神.
而神告訴你知道.

$^{641}$我必不會不報償.
這句英文說話說.
God will never be your adapter.
祂不會虧欠你的.
你按祂的旨意所作的善行.
神一定紀念.
祂會報償.
而且會歸還給你.
甚至是百倍的歸還.
當你明白這些事情的根本.
領受一份智慧的時候.
你就會很歡喜地去施予.
不覺得是缺失.
而是知道是神給你蒙恩的基本.
很多時候我們讀到聖經.
提到關於付出,施予的時候.
我們都會讀到聖經的.
哥林多後書第八章.
講到馬其頓教會落捐的後因.
我很記得當我讀到第八章的時候.
常常都會想起那句話.
就是馬其頓教會.
是非常貧窮的.
大家知道聖經提到他們是極窮的狀況.
而且遭遇著患難.
在那個時候聖經就說.
他們向保羅再三地請求.
不是勉強,無可奈何地地這樣說.
而是再三地請求.
請保羅准他們在供給聖徒的恩情上有份.
他是請求他們有機會付出.
各位弟兄姊妹.
當然我們不需要勉強自己.
不需要有一種假的信心.
說沒有問題.
神一定會給我.
但當你內心平靜地去祈禱.
知道神有應許.
知道神給你一份信心.
你就把握機會.

$^{681}$甚至好像馬其頓教會那樣.
再三請求給我一個機會.
讓我可以付出.
為什麼?.
因為你的智慧會帶給你在光明的路上.
在讀的時候.
有些時候當你一直把握著機會行善.
但行善不等於你不需要去小心謹慎.
免得你變成了惡人.
聖經接著31-32節說.
不可嫉妒殘暴的人.
不可選擇他的任何道路.
有些時候我們很想得到別人所得到的.
或者是他的財富.
或者是他的成就.
或者他所得到別人對他的擁戴.
所以我們很想得到別人所得到的.
有時候不察覺就出現了內心的嫉妒.
聖經說到內心出現嫉妒的時候.
就會有一種試探.
就會用對方所用的方法.
甚至是殘暴的方法.
去得到他已經得到的東西.
聖經提醒我們.
不要嫉妒惡人.
就算他榮華富貴.
成就非凡.
你記得不要嫉妒他.
更加不要選擇他的道路.
因為走偏方向的人.
是以和所爭護的.
正直人為他所親愛.
你會發覺.
沒錯.
某些惡人得到地上某些成就.
可能是被很多人讚賞.
歡呼擁戴.
但你記得.
這些情況是神所厭惡的.
當你想起舊約聖經.

$^{721}$當你想起約瑟的時候.
約瑟面對著主母對他很多的試探.
大家知道故事.
在過程當中其實他會明白.
如果他能夠順勢遷就下一個主人.
其實他會榮華富貴等待著他.
但他說一句話.
他拒絕一位主母.
他說:我豈能作這大惡得罪神呢?.
在他心裡很清楚知道.
我做這件事可能得到眼前的利益.
但我會叫神厭惡我.
我會作這件事得罪那位恩待我的上帝.
他一定問:.
何時你的心裡渴望神的喜悅.
渴望神的同在.
你的內心會得到一種很大的動力去拒絕邪惡.
那不是一種道德的規條.
你不准這樣不准這樣.
而是心裡知道.
你渴望神的愛.
因為祂的愛是你真正的滿足.
是你真實的幫助.
我再說一件事.
我記得有一次我在北美.
一間教會的校令會講道.
在講道期間.
在聚會的時候我遇到一對夫婦.
我看到他們容貌很寬容.
很喜樂的神態.
我就和他們聊幾句.
一邊聊的時候.
這位弟兄就告訴我他的故事.
簡單來說他們是曾經做生意的.
在大概四五年之前.
他和另一位朋友一起拍檔做生意.
生意很好很有成績.
但後來發覺他被一位拍檔欺騙.
拿了他公司很多的資源.
當他發現的時候.

$^{761}$和他理論對方說是你欺騙我.
接著一連串的官司就出現.
是對方告上法庭.
控告他欺騙自己.
這位弟兄很憤怒.
你欺騙我竟然還去官府告我.
裡面有很多很多情緒.
但那段時間他心裡領受神秘的提醒.
你繼續信靠我.
你保持良善.
這位弟兄忍耐著.
但那時候突然發現太太有癌症.
很難過.
想不到有這樣的情況.
因此那段時間他不斷要照顧太太.
太太要經歷很多種不同的治療.
因此那段時間他也不知道如何捱過.
足足四年官司結束.
他要破產.
但太太在過程中一步步康復.
康復之後.
這位弟兄心裡有一種信心.
不要緊 我從頭再來.
因此自己重新開始建立自己的公司.
當我和他見面的時候.
他的公司已經重新穩定地發展.
聽完他說之後我問他一個問題.
弟兄 這些日子不容易.
你在當中有什麼要學習.
他說我有兩件事情.
第一件事情.
令我很深刻地明白.
人生很多路上.
什麼是最重要.
我的金錢 我的健康.
就算是公義都會失去.
我受到很不公義的對待.
但我發覺因為有神的同在.
因為祂親近我.
我可以有力量經過這些幽谷.

$^{801}$我仍然可以站得住.
因此我明白人生最重要的是有神的同在.
然後他說第二件事.
第二就是我很清楚地明白.
人生的路上什麼是我的倚靠.
一切都會失去.
但主耶穌永遠永遠是我的倚靠.
是我的幫助.
成為我人生常常的提醒.
各位弟兄姊妹 無論你遇到什麼處境.
你知道神的親愛 神的同在.
是你最寶貴的.
因此你全心全意去尊心養賴祂.
在你一切所行的路上.
要認定祂 按照祂的本性 祂的旨意.
來走你的路.
祂必指引你的路.
祂必叫你的道路平穩安靜直.
走在無恩的路上.
求主將這份恩典賜給我們每一位弟兄姊妹.
我們一起祈禱.
我們要注意我們的神.
你永遠永遠是我們的倚靠.
主在你的慈愛永不改變.
有時候我們看不清楚.
因為前途為什麼這麼暗淡.
為什麼這麼多障礙.
但是主我們知道.
在其中仍然藏著你的恩惠.
主有求你幫助我們信靠你.
倚靠你 保持對你那份敬畏的心.
以至我們繼續走的時候.
我們能夠走上出路.
也能夠進入你給我們靜止的平靜的路上.
主啊 幫助我們.
奉靠耶穌聖名祈禱.
\newpage



\section{約翰福音 5:31-47}
\label{sec:031624}
\textbf{榮耀君王:使你得救}~~~~林鍵牧師
\newline
\newline
連結: \href{http://www.church.com.hk/acms/content.asp?site=cdc\&op=show\&type=product\&code=031624\&layout=sermon}{\texttt{031624}} ~~~~ 語音日期: 2026-02-07
\newline
\newline
\hyperref[sec:031623]{\small{< < < PREV SERMON < < <}}
~
\hyperref[sec:index_chronic]{\small{[返順時目]}}
~
\hyperref[sec:031625]{\small{> > > NEXT SERMON > > >}}
\newline
\newline
約翰福音 5:31-47
\newline
\begin{longtable}{cl}
\hline
\hline
章節 & 經文 (和合本修訂版)\\
\hline
5:31 & \begin{tabularx}{0.7\textwidth}{X} 「我若為自己作見證，我的見證就不真。 \end{tabularx} \\ \\ \relax
5:32 & \begin{tabularx}{0.7\textwidth}{X} 另有一位為我作見證，我也知道他為我作的見證是真的。 \end{tabularx} \\ \\ \relax
5:33 & \begin{tabularx}{0.7\textwidth}{X} 你們曾差人到約翰那裡，他為真理作過見證。 \end{tabularx} \\ \\ \relax
5:34 & \begin{tabularx}{0.7\textwidth}{X} 其實，我所受的見證不是從人來的；然而，我說這些話是為了使你們得救。 \end{tabularx} \\ \\ \relax
5:35 & \begin{tabularx}{0.7\textwidth}{X} 約翰是點亮的明燈，你們情願因他的光歡欣一時。 \end{tabularx} \\ \\ \relax
5:36 & \begin{tabularx}{0.7\textwidth}{X} 但我有比約翰更大的見證：父交給我去完成的工作，就是我正在做的，為我作證是父差遣了我。 \end{tabularx} \\ \\ \relax
5:37 & \begin{tabularx}{0.7\textwidth}{X} 那差我來的父也為我作了見證。你們從來沒有聽見他的聲音，也沒有看見他的形像。 \end{tabularx} \\ \\ \relax
5:38 & \begin{tabularx}{0.7\textwidth}{X} 你們並沒有他的道存在心裡，因為你們不信他所差來的那一位。 \end{tabularx} \\ \\ \relax
5:39 & \begin{tabularx}{0.7\textwidth}{X} 你們查考聖經，因你們以為其中有永生；而這經正是為我作見證的。 \end{tabularx} \\ \\ \relax
5:40 & \begin{tabularx}{0.7\textwidth}{X} 然而，你們不肯到我這裡來得生命。 \end{tabularx} \\ \\ \relax
5:41 & \begin{tabularx}{0.7\textwidth}{X} 「我不接受從人來的榮耀， \end{tabularx} \\ \\ \relax
5:42 & \begin{tabularx}{0.7\textwidth}{X} 但我知道，你們沒有愛神的心。 \end{tabularx} \\ \\ \relax
5:43 & \begin{tabularx}{0.7\textwidth}{X} 我奉我父的名來了，你們並不接納我；若有別人奉自己的名來，你們倒會接納他。 \end{tabularx} \\ \\ \relax
5:44 & \begin{tabularx}{0.7\textwidth}{X} 你們互相受榮耀，卻不尋求從獨一神來的榮耀，怎能信我呢？ \end{tabularx} \\ \\ \relax
5:45 & \begin{tabularx}{0.7\textwidth}{X} 不要以為我會在父面前告你們；有一位告你們的，就是你們所仰望的摩西。 \end{tabularx} \\ \\ \relax
5:46 & \begin{tabularx}{0.7\textwidth}{X} 如果你們信摩西，也會信我，因為他寫過關於我的事。 \end{tabularx} \\ \\ \relax
5:47 & \begin{tabularx}{0.7\textwidth}{X} 你們若不信他的書，怎能信我的話呢？」 \end{tabularx} \\ \\
[1ex]
\hline
\hline
\end{longtable}
$^{1}$歡迎大家來到我們的周六崇拜.
三個星期前我們離開了約翰福音一至四章.
開始進入另外一個新的部分.
約翰福音的第五章.
故事開始的時候.
耶穌來到耶路撒冷.
在靠近陽門的畢氏大祠旁邊.
醫治了一個騎腿了38年的人.
然而這個不是一個神蹟.
不是一個普通的神蹟.
而是一個十字架救贖的記號.
故事當中有上帝的羔羊.
耶穌已經有一個花名叫上帝的羔羊.
第一章已經出來了.
上帝的羔羊來到聖殿旁邊的陽門.
要經過陽門然後在池水裡面被潔淨.
是為了預備祂要去耶路撒冷的聖殿那裡.
作為一個贖罪祭.
你看這幅圖畫就知道.
祂不是在說一個神蹟.
是在說一個將來的記號.
救贖的記號.
耶穌除了是上帝的羔羊以外.
祂也是生命的活水.
約翰福音第四章.
你已經知道祂是生命的活水.
現在生命的活水浮現在池中.
或者池的旁邊.
去等待要救贖要潔淨人群當中.
那些罪人的其中一位.
接著耶穌在上千人.
如果你記得我說那裡的時候.
這個池是很大的.
大到可能是有上千人在旁邊.
耶穌只選了一個人去救.
這個人是38號.
他是騎腿了38年的人.
耶穌選這個人不是因為他獨有信心.
其實這個人是沒有信心的.
如果你記得我那裡.

$^{41}$也不是因為他特別可憐.
因為可能有一個病了40年都可以.
而是因為38年.
38年就是舊日聖經裡面.
以色列人出了埃及之後.
因為他們叛逆不信在曠野流浪.
多少年?.
進不了上帝的安息.
多少年?.
38年.
整本聖經裡面38.
只出現過兩次的38年.
一個就是在申命記裡面.
38年告訴我們.
以色列人沒辦法進入上帝的安息.
在曠野裡面死了是38年.
另外一個就是這裡38年.
耶穌在做一個天國的記號.
耶穌靈是要呼召要救贖要復活.
要帶領靈女叛逆癱瘓.
罪孽深重的以色列人.
好像舊日聖經裡面那群以色列人.
要引領他們進入上帝所應許的永恆安息.
所以耶穌必須要行這個神蹟.
是在什麼?.
是在安息日那裡行.
因為要帶這個人.
醫治他進入安息.
當年如果你記得以色列人癱瘓.
沒辦法進入上帝那個應許之地.
迦南安息.
是因為他們迷戀埃及的偶像.
於是在這個故事裡面.
耶穌最後和他們說了一句話.
他說不要再犯罪.
這個不是說普通的犯罪.
是不要再回頭貪戀埃及的偶像.
你要相反要愛慕天守上帝的安息.
要奔跑朝向上帝永恆的天國.
這個就是神蹟或者天國的記號.

$^{81}$講到尾 講真的.
施行神蹟很容易.
傳揚福音就很困難.
這個不單單是我們困難.
耶穌傳福音都是成功率極低.
這裡耶穌做完神蹟之後.
要開始講福音.
他要解釋神蹟的時候.
開口夾著舌頭說了兩句.
已經被人shut down.
耶穌都是很困難.
耶穌的神蹟和解說.
說了一句.
不單沒有被猶太人領袖所接納.
相反引來他們很憤怒的迫害.
他們說想要殺他.
你記得故事裡面.
和引起他們很不公的審判.
當地立刻.
由第十九節開始.
耶穌就是被人公審.
已經在聖殿裡面.
來到這裡我們可能會有這樣的想法.
儘管耶穌.
聖子耶穌有憐憫.
有恩慈不輕易發怒.
有豐盛的慈愛和信神.
全部耶和華在舊約聖經裡面.
耶和華的personality就是這樣.
耶和華的character就是這樣.
耶和華的本質就是這樣.
耶穌也是這樣.
儘管耶穌有憐憫.
有恩慈不輕易發怒.
忍耐都有一個極限.
於是來到這裡.
佛都有火.
於是由讓福音第五章.
第十九節開始一直到四十七節.
連續二十九節都是耶穌的獨白.

$^{121}$耶穌一輪召喚.
我們的感覺上好像是.
上帝兒子對猶太領袖的一個重錘反擊.
就是你講的又像你.
我來傳你傳福音.
你現在我講兩句.
你shut me down.
我現在一錘打過來.
這個是我們的感覺上.
但是如果你真的這樣想.
你是大錯特錯.
正如我兩個禮拜之前.
不斷的強調.
在約翰福音裡面有一個金句中的金句.
上帝愛.
我們將世人變成三八.
三八代表意識的人.
上帝當然愛世人.
但是世人當中.
也都包括了意識的人.
上帝愛三八.
甚至將他的獨生子賜給.
他們叫一切信他的不至滅亡.
反得永生.
因為上帝差他的兒子到三八來.
不是要定三八的罪.
而是要三八因怕而得救.
我們必須要這樣去理解.
耶穌的獨白.
當你仔細去研讀.
我們會發現.
耶穌每段的回應.
其實都是為了三八.
為了意識的人的好處.
上次兩個禮拜之前.
我們看耶穌獨白的頭十二節.
那十二節剛剛好分為三段.
那三段是什麼.
三段的起頭都是我實實在在地告訴你.
如果你記憶力好.

$^{161}$你會知道在《讓福音》第三章.
以色列人的神學家尼哥德姆來找耶穌.
耶穌都是用了三個.
我實實在在地告訴你和他回應的.
這裡基本上是一個重覆.
表面上如果你記得《讓福音》第三章.
好像耶穌在咆哮尼哥德姆.
痛罵尼哥德姆.
以色列人的老師什麼都不懂.
其實如果你回想起那時候我講道的時候.
耶穌在這裡循循善誘尼哥德姆.
在那個故事裡面.
在他罵他的過程裡面.
其實給了他很多舊約的反思.
舊約的應許.
引領他去到十字架的面前.
因為耶穌最後說.
正如摩西舉蛇.
人子會被舉棄.
你不斷跟著我.
你看到我被舉的時候.
你就會知道那個是重生的機會.
最後尼哥德姆就在十字架面前重生.
所以我想你記住這一切.
你發現耶穌仍然想將這個福音.
傳給這班叛逆的猶太領袖.
整段三個.
我實在告訴你的.
中心是第五章第24節.
你看耶穌想他們做什麼.
究竟那個獨特的中心.
是我實在告訴你.
那聽我話的有信猜我來.
那為的就有永生.
信我的有永生.
不信的被定罪.
不自被定罪.
而是已經出死入生.
耶穌想我們得著生命.
耶穌想我們得著永生.

$^{201}$耶穌的獨白不是要放棄以色列人.
不是要詛咒以色列人.
耶穌的獨白是要懇求他們.
聽祂信祂.
因為聽祂的話.
信靠祂是上帝差來的.
就有永生.
這個人不自被定罪.
已經從死裡面出生了.
以色列人需要聽信耶穌.
我們也都需要聽信耶穌.
為什麼.
因為原來這裡在座的每一個人.
無論現在再生或者很快要死.
在這個世界裡面.
每一個人無論生與死.
都要面對一個末日終極的審判.
楊福音第五章第28節.
如果你還記得上次.
我們提到這一節聖經.
你們不要對這事感到驚訝.
耶穌說的.
因為時候將到.
凡在墳墓裡的.
都要聽見他的聲音.
說的是耶穌自己的呼喊.
come out.
於是全部的人在墳墓裡都要出來.
然後有一個極大的審判.
行善的復活得生命.
行惡的復活被定罪.
這個就是所謂的末日的復活.
所以你下次遇到去喪禮的時候.
不用這麼傷心.
因為還有得見的.
你不想見都會見到的.
有一個末日的復活最後的審判.
死人活人全部一起受審.
你火葬了都可以將你拿出來.
海葬了都可以出來.

$^{241}$你不想出來都要出來.
信主的會出來.
不信主的都會出來.
最後審判決定.
誰進永生誰進永死.
怎樣審判呢.
耶穌很驚人的跟我們說.
行善的復活得生命.
作惡的復活被定罪.
這是一個翻譯的問題.
什麼意思呢.
這裡所謂的行善.
英文是do good.
原文是good.
但是good在聖經裡面的意思不是good.
good在聖經裡面的意思是God.
因為在創世記的第一章.
當上帝創造這個世界的時候.
每一天做完之後他就說good good good good.
什麼是good.
只有上帝是good.
有一次有人來找耶穌.
跟他說良善的夫子good teacher.
耶穌說你為什麼叫我良善.
除了上帝以外沒有人是良善的.
所以什麼叫good work.
good work的意思不是.
現在普遍宗教的人.
道德宗教說做好事.
good work的意思是God's work.
這個世界有兩種人.
一種是信靠上帝的人.
信靠上帝的人不再為自己活.
而開始為上帝而活地.
do God's work.
they do God's will.
願你的國降臨.
願你的旨意行在地上.
如同行在天上.
你不再做自己的旨意.

$^{281}$你開始進入上帝的國裡面.
做上帝的旨意.
你成為一個屬上帝的人.
所以這個對我們來說是很大的挑戰.
我們以為來教會就學那幾招.
因順稱義.
到時人家問你.
天使問你.
怎樣才知道要進入天堂的密碼.
因順稱義.
好ye .
來.
跟後面那些因順稱義.
總之因順稱義記住.
到時就可以了.
你以為這麼容易嗎.
完全不是.
羅馬書第二章保羅給我們一樣的教導.
保羅說有一天我們要照著每個人的行為.
不是靠信心.
信心要有行為才是真信心.
所以最後的審判.
是要靠各人的行為報應.
各人.
凡行心行善.
不是一日行善.
是行心行善.
Do God's work.
尋求榮耀.
這個榮耀當然不是地上的榮耀.
你說我很努力讀書的.
我很努力工作.
我很努力賺錢的.
牧師.
對不起.
你錯了.
那個是地上的榮耀.
凡行心尋求天上的榮耀.
永恆的榮耀.
C.S. Lewis 說.

$^{321}$All that is not eternal is eternally out of date.
凡是不屬於永恆的東西.
從永恆角度已經過時了.
我們要尋求的是一個不會過時的榮耀.
上帝的榮耀.
尋求榮耀.
尊貴.
不能扭歪的東西.
就有永生去報償他們.
就有榮耀.
永恆的榮耀.
尊貴平安.
家級一切行善.
當然.
反過來說.
我沒有打出來.
凡不信從真理的.
凡追求地上的榮耀的.
復活之後必被定罪.
這是一個何等.
令人驚訝的東西.
你都要回去想一想.
你不要告訴我.
你的教會多久了.
你信了多少年.
你受洗多少年.
你做過什麼執事長老.
沒有用的.
你要開始想.
究竟我們的人生.
我究竟是在做.
Good work.
God's work.
還是其實我都是.
星期日來教會.
星期六來教會.
隨便而已.
其實我的人生.
和我旁邊的人沒什麼分別.
我們追求的榮耀.

$^{361}$一樣的東西.
我都是來拜一下耶穌.
希望耶穌保守我這件事.
如果是這樣的話.
我想你有很大問題.
你去到最後審判的時候.
你不可以側側膊走到那邊.
一樣被人推回來.
你是就是.
不是就不是.
來到這裡.
你開始意識到.
我們每個人.
都在上帝的審判之下.
上帝是我們終極的審判官.
我們就是被告席上面的犯人.
You are in the dark.
你在被告席上面.
的位置就是你的位置.
唯有耶穌基督可以拯救我們.
唯有我們跟隨耶穌基督.
去行天路.
追求永恆的榮耀.
我們才有可能.
去到最後的時候.
得著永恆的榮耀.
然而這個世界.
治癒的罪人很多.
治癒的罪人有一個特徵.
就是他們很黑眼.
眼睛看不見.
他看不到自己的罪孽.
他都看不到上帝的聖潔.
有句說話說.
死到臨頭都不知死.
我們還以為.
我們在掌管自己的生命.
我們也都以為.
自己可以審判上帝.
很多時候我們.

$^{401}$有些所謂的慕道友來到教會的時候.
都問得我口啞啞.
無論你怎樣跟他說.
那個福音怎樣.
他都覺得不對.
我覺得不對.
我覺得上帝這樣不對.
你的感覺是什麼.
好像你可以審判上帝.
你沒有一種謙卑在上帝的面前.
你覺得是我決定基督教的信仰.
究竟是真與假.
耶穌是真與假.
對與錯.
我覺得上帝這樣我不同意.
你可以.
但我想你謙卑下來想一想.
究竟是你正在審判上帝.
還是上帝正在審判你.
這個就是當時猶太領袖.
面對著耶穌天國見證的時候.
表現出來的那種傲慢和無知.
但我們去到今天的那段經文.
你發覺今天的那段經文一開播的時候.
耶穌正在被人告.
耶穌是受審.
於是在受審台上.
耶穌說.
我若為自己作見證.
我的見證是不真的.
我自己作見證好像不太好.
這個故事是這樣發生的.
醫治神蹟之後.
猶太的領袖不斷找耶穌.
於是找了很久.
終於找到耶穌出來了.
於是他們再不等.
就地開庭要審判耶穌.
告他什麼呢.
告他第一.

$^{441}$犯了安息人.
當然是他們的規矩.
第二.
告他稱上帝為他的天父.
將自己與上帝捍衛同等死罪一條.
就是要將耶穌置於死地.
但在他們開庭審耶穌的時候.
從來沒有想過.
耶穌的權命和他的大能.
遠遠凌駕於他們之上.
很多時候我們覺得自己是很應該的.
說上帝.
我覺得你這樣不對.
本身已經有些問題.
耶穌你這樣說有些問題等等.
我們從來沒有想過.
其實不是我們在審判上帝.
是上帝在審判我們.
但到了這裡.
耶穌很有趣.
耶穌順水推舟.
將計就計.
你說要審我.
我又好像扮一個很無奈的樣子.
我有什麼辦法呢.
只可以被你受審.
接著耶穌就說了一大堆東西.
因為他假裝他被受審.
究竟在搞什麼呢.
耶穌之所以這樣做.
是因為他想通過他自辯時候的見證.
就是你被人告的時候.
假裝是公義.
讓你自辯一下.
於是耶穌才有機會撞獨白.
耶穌是拿著自辯的機會.
通過那個見證.
是向以色列的領袖傳講福音.
就是說他要走.
你怎麼做他呢.

$^{481}$但他不走在這裡做什麼呢.
將那段福音留在這裡給你.
希望你有一天.
可以因為他留下的那段福音的錄音.
可以得救.
我想你這樣去明白.
我若為自己作見證.
耶穌說我的見證是不真的.
好像很謙卑.
接著下面那句.
他說我另為有一位人幫我做見證.
我也知道他為我作的見證是真的.
另外一個是誰呢.
其實是在說他的天父上帝.
但是很有趣.
如果我們繼續看33至35節.
說首席的見證人是施體日漢.
但其實耶穌這裡最終說的.
另外一個人是第37節的天父.
他說那差我來的父.
為我作了見證.
這個重錘出擊.
上帝一駁鎚.
有什麼辦法呢.
但是我想你想一下.
我們應該怎樣去理解.
耶穌的思想.
耶穌提出的見證.
其實最終都是出於天父上帝.
因為是天父差他兒子來.
因為天父就是上帝.
上帝一駁鎚來.
你說知道了.
其他人走了.
走到一邊.
但是因為天父是零.
天父沒有形象.
天父不可以出面.
於是你會發現.
耶穌接下來給了一系列.

$^{521}$是天父為聖子在地上留下的痕跡.
那些叫見證.
所以耶穌不單單跟你說天父.
天父是出了名的.
但是你都不明白.
你看不到天父對我的印證.
於是耶穌給了一系列的見證.
是天父的見證.
不過天父通過地上的見證.
去表明的東西.
一系列的見證.
包括有施洗約翰的見證.
包括有父與子一同作供的見證.
你說到耶穌的神蹟騎士.
說到耶穌天國的記號.
說到耶穌很多天國的教導.
那些全部都是他與天父一起做的.
第三他會說救約聖經的見證.
救約聖經也是天父留在這個世界上.
是為著聖子耶穌所作的見證.
最後特別說到有神人摩西.
摩西的見證.
一系列的見證.
這一系列的見證其實滿足了.
救約裡面需要兩三個見證人的要求.
但更重要的是.
為什麼要說一系列的見證呢.
就是希望你有機會通過這些見證.
可以聽得明白那一副音.
而信耶穌.
於是耶穌一口氣.
將所有祂要做的見證.
和盤托出.
這裡說的是十七節的讀白.
由第五章的第三十一節.
去到第五章的四十七節.
耶穌和盤托出.
將祂剛才我說的那些見證.
全部拿出來了.
剛才告訴我們讓福音第五章.

$^{561}$第三十一至四十七節.
是一個ABCDCPA.
前後呼應的架構.
走不開的.
來來去去都是這樣.
但我想你留意.
我想你留意整段聖經.
是從一個正面流向一個負面.
最後形成一個終極的警告.
開始的時候.
另外有一位為我做見證.
結束的時候.
另外有一位再付錢.
控告你們.
由正面開始是很正面的.
不過他知道他不信的.
所以最後變成了負面的.
AA跟著BB.
你曾經有去過約翰那裡.
因為約翰閃亮的光明.
歡欣了一時.
跟著改變主意.
為什麼因為下面那個B.
因為你根本不愛上帝.
因為你不投上帝的榮耀.
所以只是斬一斬.
好像你有心去問一下.
施洗約翰關於我的事.
跟著就沒有了.
跟著CC儘管.
我做了很多的神蹟奇事.
我有很多的教導.
有聖公作見證.
儘管有聖經為我作見證.
等一下我們有機會去講.
你會比較明白一點.
然而最後D.
你們對我所有見證的回應是什麼.
你們不信.
不信我是上帝所差來的.

$^{601}$哪有為.
他們不信是因為.
他們以為他自己可以審判上帝.
讓我想起C.S. Lewis有一本書.
叫做God in the dark.
在被告直上的上帝.
這本書其實不是C.S. Lewis寫的.
是他死了之後.
那些人將他不同的寫作放在一起.
跟著要想一個名字.
那個名字就是那些不信的人.
對上帝的投訴.
或者對我們基督教信仰的一些投訴.
這本書叫做God in the dark.
被告直上的上帝.
將上帝放在被告者那裡.
他明明是法官.
但他將上帝放在被告直上那裡.
現在耶穌就是處於這樣的狀況.
但我想再次強調.
耶穌這段精彩的獨白.
不是為了要定以色列人的罪.
耶穌之後所講的一切.
是通過他自辯時候的見證.
要向猶太領袖宣講天國的福音.
最後的目的是甚麼.
使他們得救.
讓福音第五章第三十三節.
你們曾猜人到約翰那裡.
約翰為真理作過見證.
跟著第三十四節.
其實我所受的見證.
不是從人來的.
而是我說這些話.
是為了使你們得救.
使你們得救.
耶穌已經指名道姓.
我跟著跟你講的一切.
是使你們得救.
首先出場為耶穌作見證的是施洗約翰.

$^{641}$施洗約翰是耶穌的老表.
帶領耶穌半年左右.
更重要的是.
他是以色列人四百年來.
第一個真先知.
他也是上帝揀選.
應許君王的先鋒.
無論你怎麼看.
施洗約翰也是一個很重量級的見證人.
然而耶穌跟著那句.
出人意表.
其實我所受的見證.
不是從人來的.
意思是甚麼.
其實我是不需要施洗約翰的見證.
其實我不需要地上任何的見證.
因為其實我有天賦的見證.
不過我以下會跟你報告一輪.
我地上天賦為我預備的見證.
目的是要引導你們.
要幫助你們.
使你們得救.
跟著我講.
我根本不需要跟你講這些.
我不需要跟你講施洗約翰.
我不需要跟你講聖經.
我不需要跟你講摩西.
不過我跟你講.
是因為我想你通過施洗約翰的見證.
通過聖經的見證.
通過摩西的見證.
通過我的神蹟騎士的見證.
明白福音.
使你們得救.
我們繼續看施洗約翰.
耶穌這樣說.
約翰是明亮的明燈.
你們因他的光.
寬恩一時.
你記得約翰福音第一章.

$^{681}$耶穌是那個真光.
接著施洗約翰出來.
這人是為真光所見證的.
因為施洗約翰.
信靠這個真光.
因為他見證這個真光.
他就成為一個明亮的明燈.
但施洗約翰不是那個真光.
他只是見證那個真光.
當時如果你記得.
施洗約翰聚集了很多.
等待著彌賽亞尼的以色列人.
很多人.
由不同的地方來到.
要施洗約翰為他們受洗.
而當時也都引起了.
猶太領袖的注意.
於是他們派了法利賽人去問.
施洗約翰你是不是那個彌賽.
然而這裡說.
你們因他的光.
歡欣一時.
只歡欣了一時.
接著就沒氣了.
為甚麼呢.
有兩個可能.
第一個可能是.
起初他們去求問約翰.
他問約翰你是不是那個彌賽亞.
約翰說不是不是不是.
那你和彌賽亞有沒有關係.
那約翰就說了一堆東西.
他是那個先知甚麼鬼東西.
於是約翰開始指向耶穌.
其中一個.
歡欣一時的.
定義可能是.
他們起初對施洗約翰很有信心.
但他一指向耶穌.
他們不是很喜歡耶穌.

$^{721}$於是.
只是歡欣一時.
但另外一個可能.
可能更有可能的可能.
就是起初他們對施洗約翰.
很大著期望.
但很快大家知道施洗約翰.
令人失望.
為甚麼.
坐牢.
一坐牢就沒有氣了.
其實坐牢是不是.
一定代表施洗約翰是.
漏東西.
不是的.
我們說某某牧師去坐牢.
那個人又去了坐牢.
看我為甚麼坐牢.
你都未必一定說我漏的.
你可能說這樣才是真的.
但為甚麼他們會覺得.
施洗約翰一入監就漏呢.
這個和他們所追求的角度是有關係的.
無論如何.
他們到了這個時候.
大概施洗約翰已經在坐牢.
已經等著死.
於是他們只要是歡欣了.
一時.
接著耶穌繼續說.
他說不要緊.
無論你拜不拜施洗約翰都不要緊.
因為我有更大的見證.
OK.
我先給你一個株仔出場.
你不拜他不要緊.
第三個目的.
但我有比施洗約翰更大的見證.
嘩 厲害.
我負交給我去完成的工作.

$^{761}$我正在做的工作.
為我作見證.
見證我是負責差來的.
他剛剛做了甚麼工作.
就是行了神蹟.
有個三十八年動不了的人.
他可以將他救出來.
他不是跟你說神蹟.
他說你要明白神蹟背後的記號.
耶穌說.
就算你不看重施洗約翰的見證.
我有比施洗約翰更堅的東西.
就是我現在所做的所有的東西.
所有的東西不單止是剛剛做了神蹟.
是之前約翰福音一至四章裡面.
所有耶穌做過的神蹟奇事.
所有的天國記號.
其實這班人一點都不親.
整天有人盯著耶穌的.
也都看到耶穌怎樣去呼喚那些人悔改回轉.
怎樣去傳揚天國的福音.
所謂天國的福音.
就是講到君王的來臨.
耶穌不斷在宣講你們要悔改.
因為天國近了.
耶穌將天國的福音.
祂自己就是上帝那個君王的信息.
是不斷傳揚的.
祂說你聽聽我講了些什麼.
你看看我做了些什麼.
耶穌說我所有出於我手的東西.
所有出於我口的聖供.
全部都在為我作聖.
只要你清醒一下.
其實你想一下我講過的那些東西.
你想一下我做過的那些東西.
你就知道我是從天父那裡被差遣下來的.
甚至可以這樣說.
耶穌可以加多一句.
祂說我知道你一時之間很難消化這些東西.

$^{801}$我行了這麼多的事.
而不是神蹟.
而是記號.
我想你舊樣子不是那麼熟.
不要緊的.
你不熟是來找我的.
你們那個神學家.
你覺得沒得晚上來找我.
我一會兒照指點他.
其實你們不明白可以來找我的.
為什麼現在要將我受審.
為什麼不謙卑一點.
為什麼要帶著你的錫心硬著.
你們的頸項.
這些全部都是舊約聖經的description.
Stony heart.
Straight neck.
為什麼要掩耳不聽我的說話.
耶穌其實是鼓勵他們醒來.
你見到我所做的事.
你不明白不要緊.
不明白就來找我.
耶穌繼續.
他說那差我來的父.
為我作了見證.
我們問耶穌的父.
怎樣為兒子作見證.
最直接的答案就是.
在耶穌在約旦河受洗的時候.
約翰告訴我們.
在耶穌未受洗之前.
那差我來用水施洗的.
跟我說我不知道是報夢還是什麼.
約翰約翰.
明天你見到有聖靈降下來.
它停留在誰的身上.
那個人將來會用聖靈施洗.
那個就是微塞亞.
這個是約翰說出來的.
當然那個啟示是出於上帝.

$^{841}$是上帝證明的.
或者你也記得當耶穌.
在婦女福音裡面.
馬可馬太路加裡面.
說耶穌受洗的時候.
天上有一把聲音.
這是我的兒子.
我所喜悅的.
天父為耶穌作了一個.
很清晰的見證.
這個是比較擇意一點的.
我們說天父什麼時候出過面.
就是這裡.
但是我們都可以將他廣義的明白.
廣義的明白什麼呢.
聖旨在地上所有的見證.
都是出於天父的.
施洗約翰出於天父.
因為施洗約翰是誰派過來的.
上帝派過來的.
施洗約翰的啟示是哪裡來的.
上帝告訴你的.
跟著後來我們說聖經.
耶穌的神蹟是哪裡來的.
全部都是上帝和他一起做的.
所有東西都是出於天父.
我想你保留這個重點.
等一下我們回來.
你會看到整幅圖畫.
我會幫你整頓一下.
我們繼續看.
第五章第三十七節.
那差我來的婦.
也為我作了見證.
你們從來沒有聽見過婦的聲音.
你們也沒有看見過婦的形象.
你們並沒有婦的道存在你們的心裡面.
你們不信.
因為你們不信他所差來的.
那二位那位是誰.

$^{881}$耶穌.
這兩節聖經其實是引領我們.
是反向閱讀的.
你應該是反過來讀就對了.
反過來怎麼讀呢.
因為你們不信耶穌.
不信耶穌是上帝所差來的那位.
所以你們沒有上帝的道存在你們的心裡面.
所以你們聽不到上帝的聲音.
所以你們看不到上帝的形象.
你知道耶穌在說什麼嗎.
他說你不知道我是誰嗎.
你就是不知道.
你沒有看約翰福音第一章.
你看了約翰福音第一章就知道.
我就是上帝的道.
起初有道,道與上帝同在,早就是上帝.
所以如果你信了我.
你就有上帝的道在你心裡面.
你不信我,那你當然沒有上帝的道.
第二.
你記不記得約翰福音第一章第十八字怎麼說.
從來沒有人見過上帝.
只有在富懷裡獨一的兒子將上帝表明出來.
所以如果你信我的話.
其實你已經聽到上帝的聲音.
其實你信我的時候.
你已經看到上帝的形象.
因為道成了肉身,住在我們當中.
道在我們當中.
你見到上帝的聲音.
你聽到上帝的聲音.
你見到上帝的形象.
但是因為你不信我.
你什麼都沒有.
這個我想當時門徒.
十二門徒都不知道他在說什麼.
但是他回頭看,他開始明白了.
信耶穌什麼都有.
不信耶穌什麼都沒有.

$^{921}$我們繼續看第三十九章.
你們查考聖經.
猶太領袖查考聖經.
因為你們以為聖經中有永生.
而這聖經正是為我作見證的.
你們不肯到我這裡來.
但是你們卻不肯到我這裡來得到生命.
他說你們很努力地在這裡查考聖經.
你們查考聖經的原因是因為聖經裡面有永生.
然而你們不懂得看聖經.
你們不知道這本聖經是為我作見證.
於是你們最後的結局是.
你們不願意到我這裡來得到永生.
你可以這樣想.
他們為了永生去查考聖經.
這個是好事來的.
聖經明明說所有的東西指向耶穌.
但是他們就是差一點點.
看不到聖經是關於耶穌.
於是他們最後沒有辦法來到耶穌那裡.
問耶穌拿永生.
可能有人會這樣說.
他會覺得太可惜了.
猶太領袖查考聖經.
我們教會領袖都不是很查考聖經.
然後猶太領袖尋找永生.
我們教會的弟兄姊妹都不是很尋找永生.
尋找地上財寶多一點.
所以你可以說猶太領袖.
剛剛好一大半都很可惜.
沒有人教他們怎麼讀經.
沒有人教他們正確的讀經方法.
沒有人教他們查經的法則.
於是他們最後功虧一簣.
找不到耶穌.
如果你可以時光倒流.
或者穿越時空.
找林牧師去教他們怎麼讀經.
他馬上就找到耶穌了.
原來事情不是這麼簡單.

$^{961}$如果你繼續讀下去.
你會發現一件很重要的事情.
猶太人所尋求的永生.
和上帝要給人的永生.
是很不一樣的.
所以要重新定義一下.
什麼叫永生呢.
永生就是在上帝永恆天國裡面的生命.
你要這樣想.
你就不可以說永生就是永遠活著.
不是的.
所謂永生就是在上帝永恆天國裡面的生命.
如果我成為天國的一份子.
我的人生將會是怎樣.
所以以色列的領袖.
他們想要的永生.
和耶穌最後要給我們的永生.
是兩回事.
就好像這個世界.
可能有兩個同名同姓的人.
兩個都叫林健.
一個很有錢一個很窮.
你見到兩個名字.
你以為兩個名字.
都是李嘉誠.
原來在街上碰到兩個人.
完全是兩個人.
所以大家都要永生.
但是大家想要的永生.
是完全不一樣的永生.
這也是為什麼他們沒辦法相信耶穌的理由.
因為耶穌根本就不是他們的理想型.
你明白嗎.
不是他們的type(類型).
你問一下人家.
你喜歡哪個type(類型).
耶穌不是他們的type(類型).
他們要的永生.
是發達的永生.
在上帝天國裡面不發達用來做什麼.

$^{1001}$耶穌要賦予的永生.
是將來發達.
今生了倒.
有沒有想過如果你今天去街上.
隨便去街上報道.
有沒有興趣想要永生.
你不定義永生是什麼.
你就是永生.
當然想.
有誰會說我不想永生.
但是如果你解釋給他們聽.
永生就是在上帝永恆天國裡面的生命.
接著他們問你永生的人生是怎樣.
好像保羅彼得約翰.
死了.
保羅彼得約翰的人生就叫永生.
那你要不要.
當然不要.
傻的嗎.
你說做了我有永生.
我會發財.
我有永生好像保羅彼得約翰.
這樣要為主受苦.
背十字架.
不好搞.
你明白嗎.
這就是問題所在.
他們要的永生和耶穌要的永生.
是兩回事.
這不是解經的問題.
大家開始明白了.
完全不是解經的問題.
你找今天弟兄姊妹的聖經學者.
回去解經.
都解不了.
因為這是心的問題.
心結的問題.
我們還有一個見證.
兩節聖經.
在眾多聖經裡面.

$^{1041}$耶穌說摩西尤其重要.
第四十六節.
如果你們信摩西.
又會信我.
因為他寫過關於我的事.
你們若不信摩西的書.
怎能信我的話呢.
摩西這個名字在楊寶欽一共出現過十三次.
十三次不是很多.
但是你要計算.
還有許許多多摩西五經裡引述出來的經文和圖畫.
比如說耶穌.
下次我們回來.
五名耳語.
我是天國的猛男.
誰來的.
摩西來的.
所以你如果留心讀楊寶欽.
你會發現摩西是楊寶欽裡的耶穌.
一個很重要的見證人.
不是單單十三次.
而是你去到哪裡都能看到摩西.
我要你記住摩西.
因為摩西最後使人得救.
就是靠摩西的.
但是現在我想為你做一個總結.
你開始看到耶穌的獨白裡面.
有多項的見證.
On the top.
在最上面的.
其實所有的見證都是.
聖父的見證.
下面的東西都是和聖父有關的.
在聖父見證的底下.
所有見證出於聖父.
我們有第一約翰的見證.
第二有耶穌所做聖公的見證.
第三有聖經的見證.
第四最後有摩西的見證.
那個希臘文.

$^{1081}$Material.
就是見證的意思.
見證在這段聖經.
不斷出來Material Material Material.
天父在這裡Material.
約翰在這裡Material.
聖公耶穌的工作在這裡Material.
聖經在這裡Material.
摩西是沒有Material.
我加上去的.
所以框住了.
但是idea是一樣.
我們問.
耶穌為什麼要列舉這麼多項的見證.
是因為A.
祂為自己自辯.
希望可以洗脫罪名.
怕死.
還是B.
祂列舉這麼多的東西.
是因為祂想和猶太的領袖.
傳講福音.
使他們得救.
我想告訴你答案是B.
耶穌這麼長氣.
真的很老土.
祂不是怕死.
祂是怕他們死.
不是怕自己死.
是要將這些留給我們.
就算他們那群人不信.
我們今天有得看.
我們可以靠著這些見證.
去認識上帝使我們得救.
待會回來再說怎樣使我們得救.
但是我們回到這裡.
我們發覺猶太領袖其實當時還沒有準備好.
他們耶穌就算和盤托出.
說了這些.
其實他們是沒有辦法去聆聽.

$^{1121}$考量耶穌列舉出來的所有見證.
為什麼.
有得解的.
因為他們滿腦子當時.
都是世人的榮耀.
我帶你回去.
讓福音你就開始明白了.
耶穌將他們內心的底牌.
挖了出來.
告訴你.
我現在列舉給你聽.
想你得救.
但是你暫時聽不進耳.
希望以後會聽得進耳.
你聽不進耳的原因是因為.
你滿腦子都是世人的榮耀.
第五章第四十一節.
我不接受從人來的榮耀.
答案其實是反過來.
因為你接受.
我和你不同的地方就是這樣.
你追求的是世人的榮耀.
我不接受從人來的榮耀.
但我知道你們沒有愛上帝的心.
一個人為什麼會沒有愛上帝的心.
有理由的.
他沒有愛上帝的心的緣故.
是因為他心裡面不斷追逐渴求.
從人來的榮耀.
耶穌將他們的底牌爆了出來.
《聖經》也是約翰遜的一書第二章第十五節.
不要愛世界和世上的東西.
因為若有人愛世界.
愛父的心就不在他裡面了.
你的腦子是不可以兩用的.
你的腦子裝了世界.
裝不了上帝.
你不要自以為聰明.
我可以平衡的.
我可以兩邊都投資的.

$^{1161}$我可以平衡 portfolio.
你傻的.
耶穌說不可以就是不可以.
而耶穌自己也給我們很大的警告.
一個人不能事奉兩個主.
《瑪利福音》第六章二十四節.
他不是恨者過就算愛那個.
不是重者過輕那個.
你們不能又事奉上帝又事奉馬門.
以色列的領袖.
這個不是隨便在街上找人出來.
這些是敬拜上帝的人.
這些是帶領群眾敬拜上帝的人.
他們沒有辦法愛上帝.
耶穌告訴你原因是什麼.
因為他們心裡面充滿著世界的榮耀.
第四十三節.
我奉我父的名來了.
你們不接納我.
若有別人奉自己的名來.
你們都會接納他.
你不接納我的原因.
是因為我要給你的榮耀.
是你想要的榮耀.
大家明白嗎.
為什麼他們沒有辦法接納耶穌.
因為耶穌是不是有榮耀要給他.
是有的.
耶穌有很多榮耀要給你.
要和你分享演義.
耶穌要給你的是上帝的榮耀.
天上的榮耀.
永恆的榮耀.
十字架的榮耀.
不是世上人的榮耀.
所以你不要.
耶穌說.
若有別人奉自己的名來.
你們都會接納他.
若有人.

$^{1201}$有別人.
指引你們去追求地上的榮耀.
你們就會接納他.
這是耶穌的意思.
不是隨便找一個人來.
如果有另外一個人.
他來帶領你們去.
找銀.
耶穌也是帶領你們去找銀.
找真銀.
有人帶你去找假銀.
地上的榮耀.
你馬上捉到他.
接納他.
其實這句說話預表了耶穌的將來.
因為有一天.
猶太的領袖要選擇.
between巴拉巴和釘死耶穌.
你記不記得.
他們選擇了釋放巴拉巴.
耶穌怎樣.
crucify him.
巴拉巴是一個謀殺人犯.
為什麼殺人.
叛亂.
巴拉巴是一個失敗者.
但他失敗得來有道理.
失敗得來是他們的type.
追求地上榮耀的人.
所以他們選擇了巴拉巴.
不選擇耶穌.
第五章44節.
你們互相接受彼此的榮耀.
卻不尋求從獨一上帝來的榮耀.
你怎能信我呢.
如果你習慣了去追求地上的榮耀.
你怎可能信耶穌呢.
如果有一個人.
他習以為常的.
每天看見的就是眼前這個世界.

$^{1241}$他不思想人生終極永恆的事.
好像一個動物般生活.
每天都是吃喝玩樂.
但你不會去想人生終極的問題.
人從哪裡來.
人要去哪裡.
人生的意義是什麼.
如果你不想這些終極永恆的事.
這個人其實是很難信耶穌的.
因為耶穌來要指引.
要引導我們去思想的是天上的事.
是永恆的事.
又或者我這樣說.
如果有一個人.
他恆常的去敬拜黃大仙.
他就是追求金錢名利.
你要這個人突然間信著.
基本上是不可能的.
為什麼.
因為你腦裡已經充斥著那些.
他不思想人生終極的問題.
去到這裡有人會說.
林牧師你的教會不行.
我的教會可以.
不知道有多少人去拜黃大仙.
然後信了.
為什麼.
因為金融才俊的見證.
你知道發生什麼事.
你把聖經的耶穌改了.
你試一下找保羅彼得約翰來.
你教會做見證.
你看看有沒有人信主.
有沒有拜黃大仙的人.
突然間因為保羅彼得約翰的見證.
突然間砰.
痛哭流淚.
跟他們一起奔跑天路.
當然沒有.
但是你可以把耶穌改了.

$^{1281}$你可以把福音改了.
你可以當保羅彼得約翰.
什麼都不是.
只是兩個小人物.
已經過時了.
二千年前有沒有說.
今天的時間不一樣了.
保羅彼得約翰是另外一個世界的人.
金融才俊出來.
你把整個福音改了.
然後你說這個人可以信耶穌.
究竟你信了的是什麼的耶穌.
耶穌自己說.
如果你們整天在想.
在交換.
都是地上的榮耀.
你根本沒可能信.
要一個人.
真正的信耶穌.
這個人最起碼.
要開始.
在他現在的生命裡.
開始有一點搖動.
這個人要開始開竅.
可以看到人性的扭曲.
可以看到人生的短暫.
可以看到今生的虛空.
你要這樣.
才有可能讓基督的福音.
在這個人的心裡開始播種.
發芽.
成長.
如果不是.
連殺種的機會.
都沒有.
大家真的要想一下這件事.
耶穌總結.
警告他們.
第45節.
不要以為我會在父面前告你們.

$^{1321}$在父面前已經有.
另外一位.
會告你們.
就是你們.
所仰望的.
摩西.
你這是什麼意思.
他其實是在提醒他們.
在指責他們.
他說你們現在很崇拜摩西.
到末日的時候.
摩西會來告你們.
摩西會成為你們的檢控官.
為什麼?因為摩西在說的是我.
你不相信.
原來當時的以色列人.
是在等待摩西的來臨.
申命記第18章第17節.
耶穌說曾經說過.
在以後.
在你們的定聽中.
會興起另外一個先知像摩西.
就是因為那一句話.
以色列人.
幾百年幾千年都在等摩西的再現.
他們在等摩西2.0.
在等摩西的增強版.
尤其是在如意節的時候.
特別希望摩西來.
再做一個如意節.
如意節拿來做什麼.
殺羊.
不是如意節.
是出埃及.
是擊打埃及.
擊殺埃及.
於是他們在等待摩西的再現.
他們要摩西施行神蹟.
他要摩西摧毀埃及.
他要一個這樣的先知.

$^{1361}$他要一個這樣的基督.
他要一個這樣的君王.
從某個角度來說.
他們恨不得耶穌.
就是那個摩西.
他恨不得他看到耶穌.
禁馬炮.
他霸不得耶穌.
就是用他的神力.
為他們打江山.
聖羅馬.
德榮耀.
懸疑.
耶穌不肯.
死都不肯.
就是因為他死都不肯.
所以就送他上西天.
你明白嗎.
耶穌需什麼.
不是需無能.
就是他有能力.
他不願意.
降伏於人的榮耀底下.
於是他們對耶穌徹底的失望.
因為耶穌死都不肯.
他們對耶穌恨之入骨.
於是他們送耶穌上十字架.
但我想你知道另外一樣東西.
對耶穌失望的何止是猶太的領袖.
耶穌自己十二門徒.
都對耶穌極端的失望.
我想帶你去看這個.
大家很熟悉的一段經文.
《路加福音》第24章.
講到有兩個對耶穌非常失望的門徒.
我們不知道他們的名字.
但我們知道他們對耶穌很失望.
你記得這個故事嗎.
以馬五師路上有兩個黑眼心死的門徒.
他們聽到耶穌.

$^{1401}$他見到耶穌.
但他認不出耶穌.
耶穌令他們看不到.
但也反映了他們的熟悉狀態.
他們黑眼心死.
於是他們向耶穌訴苦.
為什麼這麼慘.
兩位兄台.
他說因為他是耶穌.
因為他說話行事都帶有能力.
這個是誰 摩西.
舊聖經裡面有一個人.
說話行事都帶有能力.
那個叫摩西.
這麼像摩西2.0增強版.
我們數來盼望.
這個盼望和之前的盼望.
是同一個希臘文.
他們很盼望這個摩西.
可以救贖以色列人.
怎樣救贖 打贏羅馬.
《馬太福音》告訴他們.
他的名字叫耶穌.
因為他要將上帝的子民.
從他們的罪裡拯救出來.
他們看不到罪.
他們看到羅馬.
我現在最實際.
我生活要過好一點.
我錢要多一點.
我的命要長一點.
我的健康要改善.
你幫我搞定我就信你.
神經病.
他的名字叫耶穌.
因為他要將上帝的子民.
從罪惡裡拯救出來.
他們看不到.
於是他們對耶穌極其失望.
我們對他有很大的盼望.

$^{1441}$但是他死了.
現在耶穌復活了.
我想你看看耶穌怎樣對他們.
怎樣去救他們.
耶穌第一句說話.
這個很重要.
耶穌還沒有開始查經.
大家說怎樣救他們.
要查經.
還沒有 還沒有查經.
就這樣一開始查經是沒有用的.
耶穌第一個動作.
對他們說.
無知的人.
你們的心信得太遲了.
無知的人.
你們貪愛世界.
不愛上帝.
這個就是耶穌跟他們說的.
跟那些以色列的領袖說的一樣.
為什麼上帝你沒有愛上帝的心呢.
因為你貪愛這個世界.
耶穌一拍下去.
這個東西很重要.
你不幫他.
他不會醒的.
所以這句說話很重要.
他說無知的人.
你們就是貪愛世界.
不愛上帝.
於是你看不清楚.
聽不明白摩西和先知.
如何指向我.
又如何顯明.
基督必須受苦進入榮耀.
如果你要的是錢.
要的是名.
要的是利.
你怎樣查聖經.
都查不到耶穌.

$^{1481}$查不到聖經新約裡面.
那個十字架上的耶穌.
你要twist and turn.
將耶穌變形.
你才可以找到你想要的理想型.
耶穌做了一樣東西.
打醒你.
因為你的眼睛迷糊.
很多眼屎.
擦一擦.
這樣你才有機會.
醒出來.
你醒了我們才可以去查經.
令我想起.
有一套經典的美國電影.
1957年.
法律司法劇.
《12 Angry Men》.
十二怒汗.
這套電影真棒.
在美國電影史上.
其中數一數二的法律戲劇.
非常有名.
有名到一個地步.
台灣,日本,中國大陸,香港.
都有話劇團是重演的.
還有一套.
你可以回去找找.
韓國翻拍的.
《Jura 8》第八號的陪審員.
故事講什麼呢.
故事圍繞著有十二人的陪審團.
他們有個案在手上.
有個貧民區裡有個十八歲的男孩.
被兇謀殺他爸爸.
聲明攸關.
那個法官跟他們說.
你們一定要十二個人達成一同的協議.
而協議的中心思想就是.
有沒有reasonable doubt.

$^{1521}$有沒有合理的懷疑.
如果有合理的懷疑就not guilty.
not guilty不代表他無罪.
是代表他謀殺無罪.
所以他們要決定.
究竟他謀殺有罪還是謀殺無罪.
故事開始的時候.
案件好像鐵證如山.
他們自己十二個人在討論的時候.
第一次投票.
十一個說有罪謀殺.
另一個不是說無罪.
另一個說我不是很肯定.
另一個說生命攸關.
是否我們應該再想想.
再從頭再看一看.
才決定.
我不是說他無罪.
但是我不肯定.
於是十二個人因為一個人.
他是一個architect(建築師).
他想細心的去再看看那些東西.
重新的反思.
然後發現疑點重重.
他們發現在過程裡面.
原來每個人都有他很明顯的偏見.
因為太明顯了看不到.
他們發現每一個人都有他錯誤的假設.
於是他們談了幾個小時.
兩三個小時.
整套戲很短.
結果發現這個謀殺罪是不成立的.
結果十二個人.
not guilty.
由十一個guilty一個I don't know.
去到十二個人同意.
not guilty.
我們有reasonable doubt.
我為什麼要說這個故事給你聽.
因為我想告訴你.

$^{1561}$耶穌一樣.
你想想這班人想定耶穌死罪.
說真的耶穌做了什麼值得他死的.
如果你問他們有沒有reasonable doubt.
耶穌做的事是不是值得他死.
其實只要他們將自己的有色眼鏡.
退去一下.
他是不是起碼覺得耶穌罪不該死.
他們要釘死耶穌說guilty.
是因為他們戴上有色眼鏡.
我也想跟你說.
很多人第一次接觸基督教的時候.
很多時候都有一種有色的眼鏡.
他跟我談的時候.
可能他覺得自己很有道理.
他會給我三五個.
他覺得很有道理.
為什麼基督教思想是不成立的原因.
從來沒有想過他所問的問題.
他所提出的東西.
說真的幾千年.
你以為沒有其他人提出過嗎.
而這個人提出的時候.
根本看不到自己.
心裡面是充滿非常明顯的偏見.
和錯誤的假設.
其實我不需要你將你有色眼鏡脫下.
我只需要你移動一下.
只要你移動一下.
你就開始明白.
基督教流傳了這麼多年.
有它的道理.
罪是明顯的東西.
審判是極有可能的事情.
十字架是一個很獨特的救贖方法.
我想說什麼呢.
你要信耶穌.
第一件事你要移動一下你自己的有色眼鏡.
然後你開始明白.
原來你想釘死基督信仰.

$^{1601}$不是這麼容易的.
於是你開始有一個開口.
有一個可能性.
於是耶穌就可以開始和你查經.
在你沒有移動到你的眼睛之前.
你是不能查經的.
你和你怎樣查.
我看不到.
因為我不同意.
我不同意當然不同意.
你戴著這麼深的有色眼鏡.
願意耶穌打他一下.
打他一下.
愚蠢的人.
然後耶穌就開始和你們查經.
怎樣查經.
從舊約聖經裡面指向耶穌.
於是他們眼睛迷糊.
心裡從眼睛迷糊變成眼睛明亮.
查經之後.
從心遲鈍slow apart.
變成做心火熱burning hot.
然後他們開始的時候.
是走向爾瑪吾斯.
走錯路.
你明白嗎.
然後最後他們u-turn.
他們再不去爾瑪吾斯.
或者去到爾瑪吾斯.
然後扭tie 回去耶路撒冷.
代表整個故事的大反轉.
一個人的迴轉.
因為他們看見耶穌的真實.
在路加福音第24章第33節那裡.
And they rose.
他們起來.
起來的意思其實是原文可以是復活.
他們復活了.
因為耶穌的查經.
爾瑪吾斯路上的查經.

$^{1641}$但是我想和你再提一下.
因為耶穌打了他們一下.
所有的東西從耶穌當頭.
棒喝那一下開始.
醒了醒.
於是他們擦乾淨自己的眼屎.
於是可以開始一起的查經.
於是從摩西和眾先知起.
凡經上所指著耶穌的東西.
都向他們解釋.
這次想得通了.
是喔 make sense.
為什麼上個禮拜不make sense呢.
因為你的頭還沒被人打.
你仍然帶著世界的眼睛去這件事.
但你只要脫一脫眼睛.
你開始明白.
但是我們回到耶穌.
有更多的見證.
我們又鬥了很久才回來.
所有的見證都是為了他們得救.
所有的見證有一個共同點.
所有的見證都是為了要創造一個.
爾瑪吾斯路上的經驗.
你有沒有發現.
施啟約翰講了很多舊約的東西.
人家問他你是哪位.
他介紹他不是說我介紹.
我是某某城市出來那裡讀書的.
他是引用舊約.
講他自己的身份的.
你有沒有發現.
施啟約翰的見證是充滿著舊約.
你有沒有發現耶穌基督的聖經.
是充滿著舊約的影子.
耶穌不是隨便找一些神蹟.
不如上天飛一飛回來.
他不是的.
他做的所有事都是舊約裡面.
預言了的東西.

$^{1681}$你有沒有發現整本聖經的東西.
都是指向耶穌.
你有沒有發現摩西所講的東西.
都是指向耶穌.
時間有限.
但是我想和你快快去看一下.
約翰的見證.
我是用水施洗的.
但有一位在我們當中.
你們不認識的.
他在我耳後來的.
我給他綁鞋帶都不配.
他要用聖靈受洗.
這是哪裡來的.
舊約聖經出來的.
接著施啟約翰講了一個很重要的見證.
The Lamb of God will take away the sin of the world.
耶穌是誰.
上帝的高人.
要除去世人的罪.
這是哪裡來的.
舊約聖經出來的.
施啟約翰的見證.
然後是耶穌的見證.
耶穌的見證真是多到我都數之不盡.
你記得我們剛剛讀完.
約翰福音一至四章.
接著我們再看第五章.
耶穌的另外一件神蹟.
所有的東西其實.
我很快說.
你看到的好像是.
迦南的婚宴水變酒.
其實是十字架的記號.
在第二章裡面.
耶穌為上帝的聖殿召集.
他要告訴他們.
他是大衛的後人.
他是猶太人的君王.
去到第四章.

$^{1721}$耶穌是生命的活泉.
那個活泉從聖殿裡面流出活水.
在第四章裡面亦都說.
他是上帝慈愛信實的彰顯.
他是聖靈帶我們進入的永恆敬拜.
你說要敬拜的.
要是the spirit and the truth.
聖靈.
聖靈從哪裡來.
以西結書47章來.
and the truth.
那個不是真理.
是十字架上上帝承傳救贖的信實.
唉 說不了那麼多.
跟著我們見到第四章最後一個.
是帶神兒子活了.
其實是一個復活的記號.
所有的東西.
都是帶我們回到舊約聖經.
看看耶穌怎樣去承傳舊約聖經.
耶穌第五章來到耶路撒冷的陽門.
在那裡醫治了一個38年騎腿的人.
又是一個天國的記號.
所有的東西.
告訴我們.
不是神蹟乃是記號.
不是地上乃是天上.
不是今生乃是永恆.
這些全部都是我以前講到裡面的精彩片段.
你一定要想起那些.
所有的東西全部有連貫.
跟著中間告訴我們.
我們要信耶穌.
我們要思想天上的事情.
我們要思想上帝的事情.
我們要思想永恆的事情.
我們要信耶穌.
真信耶穌不是假信.
因為有一個人很強調.
見到神蹟信耶穌的是假信.

$^{1761}$要聽到上帝的話語信耶穌.
才是真信耶穌.
耶穌是誰.
耶穌是世界的救主.
在耶穌的面前.
我們只有一個回應.
我們降服在祂面前.
祂必興旺.
我們必衰微.
耶穌是獨一的救主.
耶穌有律法和先知為祂作見證.
回到第一章.
律法書上所寫的.
和眾先知所寫的那一位.
我們預想了.
整本聖經舊的聖經指向耶穌.
祂就是約瑟的兒子耶穌.
祂要為我們做什麼.
摩西給我們一個重要的提示.
或者一個重要的預言.
摩西在曠野怎樣舉蛇.
人子也必須被舉起.
使一切信他的人道德重生.
為什麼耶穌最後的見證是摩西.
因為祂想帶我們回到這個圖畫.
因為祂一早已經和以色列的神學家.
尼葛德姆說了祂最後的命運.
祂說你看到我被舉起來.
你仰望我.
你就可以得著永生.
耶穌說得出做得到.
於是在約翰福音最後的場景.
耶穌被釘死在十字架上.
兵士擲炸機立旁.
有血和水流出來.
這個血和水重生了兩個人.
一個是約瑟一個是尼葛德姆.
這兩個人從前在黑暗中跟隨耶穌.
不想被人知道.
在那一天他們來找耶穌.

$^{1801}$拿身體.
找彼拉多拿耶穌的身體.
他們離開了黑暗.
追隨了真光.
他們奉上他們所有帶了75年薪水.
買的沉香活藥.
他們拋棄世界.
單求天王.
他們被捲入了上帝重生的漩渦裡面.
願我們都是一樣.
領受耶穌的見證.
聆聽祂的說話使我們得救.
耶穌將所有的見證留下給我們.
叫我們跟著聖經.
可以在祂裡面得到重生.
這是祈禱.
天父我們很明顯感覺到自己的軟弱.
就是你的聖經是何等的豐富.
我們的腦袋是何等的狹窄.
太多數據了.
太多不同的經文.
尤其是我們去到第五章的時候.
我們發現原來第五章.
去到這個獨白的結尾.
所有的東西和前面一至四章.
一至五章的東西都有關係.
我求主你幫我們看見耶穌的榮耀.
看見祂那種在敵人面前.
堅持獨白的去見證自己.
不是因為祂要脫罪.
而是因為祂要拯救我們.
耶穌既然將這些見證放在我們面前.
也求主你讓我們有一種追求的心智.
好好地讀這本一千一百八十九章的聖經.
我們真的開始明白為什麼這本書這麼厚.
我們也開始明白.
要我們珍惜這本書的每一章每一節.
因為原來所有的東西都是指向耶穌.
而當我們有以馬五思的經驗的時候.
我們的眼睛就可以明亮.

$^{1841}$我們的心裡就可以火熱.
我們可以好像保羅彼得約翰一樣.
得著那個天國的生命.
幫我們在耶穌.
在聖經裡面尋找你的永生.
我們這樣祈禱你奉主耶穌基督的名求.
\newpage



\section{約伯記 1:1-2:13}
\label{sec:031625}
\textbf{全心全意}~~~~俞大鳴牧師
\newline
\newline
連結: \href{http://www.church.com.hk/acms/content.asp?site=cdc\&op=show\&type=product\&code=031625\&layout=sermon}{\texttt{031625}} ~~~~ 語音日期: 2026-02-08
\newline
\newline
\hyperref[sec:031624]{\small{< < < PREV SERMON < < <}}
~
\hyperref[sec:index_chronic]{\small{[返順時目]}}
~
\hyperref[sec:031626]{\small{> > > NEXT SERMON > > >}}
\newline
\newline
約伯記 1:1-2:13
\newline
\begin{longtable}{cl}
\hline
\hline
章節 & 經文 (和合本修訂版)\\
\hline
1:1 & \begin{tabularx}{0.7\textwidth}{X} 烏斯地有一個人名叫約伯。這人完全、正直、敬畏神、遠離惡事。 \end{tabularx} \\ \\ \relax
1:2 & \begin{tabularx}{0.7\textwidth}{X} 他生了七個兒子，三個女兒。 \end{tabularx} \\ \\ \relax
1:3 & \begin{tabularx}{0.7\textwidth}{X} 他的家產有七千隻羊，三千匹駱駝，五百對牛，五百匹母驢，並有許多僕婢。這人在東方人中為至大。 \end{tabularx} \\ \\ \relax
1:4 & \begin{tabularx}{0.7\textwidth}{X} 他的兒子按著日子各在自己家裡擺設宴席，派人去請他們的三個姊妹來，與他們一同吃喝。 \end{tabularx} \\ \\ \relax
1:5 & \begin{tabularx}{0.7\textwidth}{X} 宴席的日子過了，約伯派人去叫他們自潔。他清早起來，按著他們眾人的數目獻燔祭，因為他說：「恐怕我的兒子犯了罪，心中背棄神。」約伯常常這樣行。 \end{tabularx} \\ \\ \relax
1:6 & \begin{tabularx}{0.7\textwidth}{X} 有一天，神的眾使者來侍立在耶和華面前，撒但也來在其中。 \end{tabularx} \\ \\ \relax
1:7 & \begin{tabularx}{0.7\textwidth}{X} 耶和華對撒但說：「你從哪裡來？」撒但回答耶和華說：「我從地上走來走去，在那裡往返。」 \end{tabularx} \\ \\ \relax
1:8 & \begin{tabularx}{0.7\textwidth}{X} 耶和華對撒但說：「你曾用心察看我的僕人約伯沒有？地上再沒有人像他那樣完全、正直、敬畏神、遠離惡事。」 \end{tabularx} \\ \\ \relax
1:9 & \begin{tabularx}{0.7\textwidth}{X} 撒但回答耶和華說：「約伯敬畏神，豈是無故呢？ \end{tabularx} \\ \\ \relax
1:10 & \begin{tabularx}{0.7\textwidth}{X} 你豈不是四面圈上籬笆圍護他和他的家，以及他一切所有的嗎？他手所做的都蒙你賜福，他的家產也在地上增多。 \end{tabularx} \\ \\ \relax
1:11 & \begin{tabularx}{0.7\textwidth}{X} 但你若伸手毀他一切所有的，他必當面背棄你。」 \end{tabularx} \\ \\ \relax
1:12 & \begin{tabularx}{0.7\textwidth}{X} 耶和華對撒但說：「看哪，凡他所有的都在你手中；只是不可伸手加害於他。」於是撒但從耶和華面前退出去。 \end{tabularx} \\ \\ \relax
1:13 & \begin{tabularx}{0.7\textwidth}{X} 有一天，約伯的兒女正在他們長兄的家裡吃飯喝酒， \end{tabularx} \\ \\ \relax
1:14 & \begin{tabularx}{0.7\textwidth}{X} 有報信的來見約伯，說：「牛正耕地，母驢在旁邊吃草， \end{tabularx} \\ \\ \relax
1:15 & \begin{tabularx}{0.7\textwidth}{X} 示巴人忽然闖來，把牲畜擄去，並用刀殺了僕人；惟有我一人逃脫，來報信給你。」 \end{tabularx} \\ \\ \relax
1:16 & \begin{tabularx}{0.7\textwidth}{X} 他還說話的時候，又有人來說：「神從天上降下火來，把羊群和僕人都吞滅了；惟有我一人逃脫，來報信給你。」 \end{tabularx} \\ \\ \relax
1:17 & \begin{tabularx}{0.7\textwidth}{X} 他還說話的時候，又有人來說：「迦勒底人分成三隊忽然闖來，把駱駝擄去，並用刀殺了僕人；惟有我一人逃脫，來報信給你。」 \end{tabularx} \\ \\ \relax
1:18 & \begin{tabularx}{0.7\textwidth}{X} 他還說話的時候，又有人來說：「你的兒女正在他們長兄的家裡吃飯喝酒， \end{tabularx} \\ \\ \relax
1:19 & \begin{tabularx}{0.7\textwidth}{X} 看哪，有狂風從曠野颳來，襲擊房屋的四角，房屋倒塌在年輕人身上，他們就都死了；惟有我一人逃脫，來報信給你。」 \end{tabularx} \\ \\ \relax
1:20 & \begin{tabularx}{0.7\textwidth}{X} 約伯就起來，撕裂外袍，剃了頭，俯伏在地敬拜， \end{tabularx} \\ \\ \relax
1:21 & \begin{tabularx}{0.7\textwidth}{X} 說：「我赤身出於母胎，也必赤身歸回；賞賜的是耶和華，收取的也是耶和華。耶和華的名是應當稱頌的。」 \end{tabularx} \\ \\ \relax
1:22 & \begin{tabularx}{0.7\textwidth}{X} 在這一切的事上，約伯並沒有犯罪，也不以神為狂妄。 \end{tabularx} \\ \\ \relax
2:1 & \begin{tabularx}{0.7\textwidth}{X} 又有一天，神的眾使者來侍立在耶和華面前，撒但也來在其中。 \end{tabularx} \\ \\ \relax
2:2 & \begin{tabularx}{0.7\textwidth}{X} 耶和華問撒但說：「你從哪裡來？」撒但回答說：「我從地上走來走去，在那裡往返。」 \end{tabularx} \\ \\ \relax
2:3 & \begin{tabularx}{0.7\textwidth}{X} 耶和華對撒但說：「你曾用心察看我的僕人約伯沒有？地上再沒有人像他那樣完全、正直、敬畏神、遠離惡事。你雖激起我攻擊他，無故吞滅他，他仍然持守他的純正。」 \end{tabularx} \\ \\ \relax
2:4 & \begin{tabularx}{0.7\textwidth}{X} 撒但回答耶和華說：「人以皮代皮，情願捨去一切所有的，來保全性命。 \end{tabularx} \\ \\ \relax
2:5 & \begin{tabularx}{0.7\textwidth}{X} 但你若伸手傷他的骨頭和他的肉，他必當面背棄你。」 \end{tabularx} \\ \\ \relax
2:6 & \begin{tabularx}{0.7\textwidth}{X} 耶和華對撒但說：「看哪，他在你手中，只要留下他的性命。」 \end{tabularx} \\ \\ \relax
2:7 & \begin{tabularx}{0.7\textwidth}{X} 於是撒但從耶和華面前退出去，擊打約伯，使他從腳掌到頭頂長毒瘡。 \end{tabularx} \\ \\ \relax
2:8 & \begin{tabularx}{0.7\textwidth}{X} 約伯就坐在灰燼中，拿瓦片刮身體。 \end{tabularx} \\ \\ \relax
2:9 & \begin{tabularx}{0.7\textwidth}{X} 他的妻子對他說：「你仍然持守你的純正嗎？你背棄神，死了吧！」 \end{tabularx} \\ \\ \relax
2:10 & \begin{tabularx}{0.7\textwidth}{X} 約伯卻對她說：「你說話，正如愚頑的婦人。唉！難道我們從神手裡得福，不也受禍嗎？」在這一切的事上，約伯並沒有以口犯罪。 \end{tabularx} \\ \\ \relax
2:11 & \begin{tabularx}{0.7\textwidth}{X} 約伯的三個朋友，提幔人以利法、書亞人比勒達、拿瑪人瑣法，聽說這一切的災禍臨到他身上，各人就從自己的地方相約同來，為他悲傷，安慰他。 \end{tabularx} \\ \\ \relax
2:12 & \begin{tabularx}{0.7\textwidth}{X} 他們遠遠地舉目觀看，認不出他來，就放聲大哭。各人撕裂外袍，向空中撒塵土，落在自己的頭上。 \end{tabularx} \\ \\ \relax
2:13 & \begin{tabularx}{0.7\textwidth}{X} 他們同他七天七夜坐在地上，一句話也不對他說，因為他們見到了極大的痛苦。 \end{tabularx} \\ \\
[1ex]
\hline
\hline
\end{longtable}
$^{1}$我地一齊祈禱下.
多謝你恩典讓我地起住一個清晨.
可以黎到呢度聚集敬拜你.
學習你ge 話語.
求聖靈幫助我地.
堅固我地ge 心.
使我地可以全心全意跟隨主.
奉耶穌基督的名祈禱.
阿們.
猶太人有一段ge 祈禱文呢.
係gwo D 虔誠ge 猶太教徒.
每日都會背三次ge .
就係呢係《新命紀》六章四至九節.
係四至五節呢.
係一開始gwo 度就話.
以色列你要聽.
耶和華我們的上帝.
是獨一的主.
你要盡心盡性盡力.
愛耶和華你的上帝.
當時耶穌被一位律法師問到佢.
解明宗最大ge 係邊條呢.
或者話最重要係邊一條呢.
主耶穌亦都回答話.
你要盡心盡性盡力.
盡意愛主你的上帝.
猶要愛倫如己.
咁簡單黎講呢.
就係要人全心全意咁去愛上帝.
呢個係上帝對佢子民ge 期望.
今年我地教ge 主題就係「全心全意」.
我地對好多事物都可能會全情投入.
包括我地ge 婚姻,家庭,工作或者事業.
有時我地不自覺之間.
就會忽略jor 對上帝ge 全心全意.
係主耶穌好多教導當中.
其實都提到全心全意ge 重要.
例如係《馬太福音》六章二十四節.
佢講過.
「一個人不能事奉兩個主」.

$^{41}$「那是恨者過愛那過,就是重者過輕那過」.
「你們不能又服侍上帝,又服侍馬門」.
咁即係話要全心全意去服侍上帝啦.
係《路加方》九章六十二節.
耶穌又咁樣對一個人講過.
「守扶著泥,向後看的人,不會進上帝的閣」.
呢度就用jor 一個農夫耕地ge 比喻.
扶著泥,就係農夫點解要扶著泥呢?.
希望整鬆gwo D 泥土.
等D 種子可以容易發芽生長.
我聽過一位牧師講過.
猶太人耕種同華人就有D 唔同.
華人未落種子就耕好條地先.
整鬆晒先.
猶太人就調轉黎殺jor 種先.
殺完種佢先去耕鬆gwo D 泥地.
其實無論你做邊樣先都好啦.
我相信耕地ge 目的.
都係整鬆D 泥土同埋.
佢點解要耕地呢?.
就係希望整到D 田地齊整一D .
方便日後你去林水又好收割又好.
都比較容易做.
因此手扶著泥ge 農夫.
係應該前面有隻牛ge .
個牛拖住個泥.
個農夫就係後面幫手扶住佢.
等佢直D 咁樣.
如果一個農夫一邊就扶住泥.
一邊就想向後面睇.
其實幾難做ge .
如果我地當中有人推個剪草機就知.
推剪草機梗係向前推.
你一邊推一邊望後面點做呢?.
好艱難ge ,好唔容易.
《舊約聖經》仲記載jor 一個故事.
就係有關羅德夫人ge 故事.
天使就要救羅德.
因為羅德都算係義人.
天使就唔想將義人同埋.

$^{81}$所多瑪俄摩拉gwo 兩個罪惡之神ge 人一齊毀滅.
所以就准許羅德帶埋佢D 家人逃亡.
快D 走,走去一個小山gwo 度.
但係吩咐佢地你地走ge 盡快.
走去gwo 個山唔好拎返轉頭睇.
結果羅德ge 大大.
佢就唔知忘記jor 定係好奇心太強.
佢就拎返轉頭睇.
一路走叫佢走快D .
佢拎返轉頭睇.
結果變jor 一根鹽柱.
今日我好想同大家思考聖經中一位人物.
睇睇佢係點樣全心全意信靠上帝.
呢個人就係約伯.
好多人讀約伯記.
睇到約伯受jor 好多苦難.
同埋佢三個朋友.
都係同佢討論苦難ge 原因.
就以為約伯記ge 主題係討論苦難.
事實上約伯記ge 重點.
唔係討論苦難ge 源頭或者苦難.
如果你真係想搵苦難ge 源頭.
你讀完約伯記都搵唔到答案.
因為冇呢個答案.
上帝後來向約伯顯現ge 時候.
都冇回答約伯ge 問題.
冇回答過點解會有苦難臨到約伯.
但由於上帝肯同約伯傾下計.
肯同佢顯現.
約伯就好開心.
今日我地會睇約伯記頭兩章.
從上帝同撒旦ge 對話裡面.
去了解約伯記ge 重點討論ge 問題.
我地先睇第一章.
約伯記一章一節就介紹呢個主角.
烏斯地有一個人名叫約伯.
這人完全正值敬畏上帝遠離惡善.
經文一開始就描述約伯.
係一個敬畏上帝ge 人.
由於佢信上帝 敬畏上帝.

$^{121}$係行為上佢係完全正值.
並且遠離惡善.
呢篇經文冇形容約伯係一位異人.
好多時候有D 人總係認為.
唔信主耶穌ge 人就一定落地獄.
而忽略jor 保羅係羅馬書二章一至十六節ge 講論.
gwo 度先講到上帝係一位公義ge 上帝.
係一位公義ge 審判者.
對猶太人就以律法作為審判猶太人行為ge 標準.
而對所有非猶太人就以人ge 良心.
作為標準去審判佢地ge 行為.
呢個就顯出上帝ge 公義.
唔係用一個標準ge 佢唔同人有唔同ge 標準.
因此所有係主耶穌未出生之前ge 非猶太人.
其實都係以佢ge 良心同埋行為.
係接受上帝ge 審判.
其實係舊約裡面都記載jor 幾位ge 非猶太人.
無論上帝有冇稱佢地為異人都好.
我相信佢地都會進入天界.
其中有爾諾就最明顯.
聖經只記載佢與上帝同行三百年.
被上帝接jor 返天界.
跟住落去仲有羅亞.
由於上帝搵到呢一位係罪惡世代裡面ge 好人.
上帝就決定拯救佢一家人.
使佢不至於被洪水毀滅.
仲有另一位係亞伯拉罕遇見ge 撒冷王默基西德.
佢係至高上帝ge 祭司.
呢幾位都係主耶穌成為肉身之前ge 非猶太人.
我相信佢地都會進入天界.
所以假如我地傳福音時被人問到孔夫姐咪可以進入天界呢.
我地可以話有可能ge .
但係佢同樣要接受上主公義ge 審判.
先至知道最後ge 結果.
千萬唔好講無信耶穌ge 人死緊.
又非猶太人去地獄.
孔子都唔掂.
呢個講法實際上係唔合乎聖經ge 教導.
我剛才已經講jor .
聖經本身都已經記載jor 幾位非猶太人ge 義人.

$^{161}$有D 人時時就引聖經.
哎呀 冇有義人 連一個都冇有.
我再提醒大家.
呢個出自詩篇十四章.
唔係第一節 係第三節.
但十四章開始時點講呢.
佢話愚頑人心裡算冇上帝.
他們都敗壞行了可憎惡可憎惡的事.
冇一個人行善.
佢話先將人類分jor 先.
分jor 做兩批.
一批佢叫做愚頑人.
係呢D 愚頑人心裡冇上帝ge 人當中就冇義人.
連一個又冇有.
真ge  冇錯.
但可能仲有另外一批.
唔係愚頑gwo D 點算呀.
係咪會有義人呢.
未必.
返番黎講約伯.
約伯係咪猶太人呢.
大概唔係.
因為一章一節講明佢係烏斯地ge 人.
烏斯地係邊度呢 其實我地唔知道.
應該唔係迦南地呀.
應該唔係迦南地.
聖經形容約伯係完全正直同埋遠離惡事ge 人.
我相信其實佢都會入得天界.
係一章五節提到.
佢係以一個家長ge 身份為兒女獻祭.
專家認為佢應該係同亞伯拉罕時代ge 人.
甚至更早D 都唔定.
因為亞伯拉罕都係自己.
祝壇獻祭.
冇搵祭司做.
係第一章第二節就介紹約伯ge 家庭狀況.
話佢有七個仔三個女.
七加三就係十.
呢個都係一個代表完全ge 數字.
第三章就講到非常富裕.

$^{201}$佢有七千隻羊 三千匹駱駝 五百堆牛 五百匹母鹿.
七千加三千就係一萬.
其實呢D 都係十ge 倍數.
五百加五百又係一樣啦.
都係一千又係十ge 倍數.
深水清ge 弟兄姊妹可能會留意到呢度話佢.
有五百.
咩牛呀?.
五百堆牛喎.
五百堆喎唔係五百隻喎.
一堆呢就兩隻ge .
咁換句話講呢佢其實應該係有一千隻牛.
但作者故意呢就唔話佢有一千隻牛.
話佢五百堆牛.
就同下面gwo 個五百呢係黎到對稱.
五百匹母鹿.
咁事實上係唔係真係咁巧合.
七千三千五百加五百呢.
其實都唔一定ge .
你arm arm 數完又有隻牛 生gwo 隻牛仔點算 點計數呢.
但係作者用一個齊頭數.
去表達佢係唔止佢兒女.
係完全ge 數字.
佢養ge 生畜佢擁有ge 生畜.
都係完全ge 數字.
有D 中國人講十全十美.
中國人用十全十美唔會話九全九美.
都係一個表達一個完全ge 意思.
一個敬畏上帝ge 人.
上帝都係祝福佢.
使佢十全十美.
第十節就講gwo D 兒女真係好.
老豆賺錢仔 享福佢地時時開party.
飲飲食食.
第五節就話約伯為佢地獻祭俾上帝.
原因係怕佢D 兒女.
係無意中犯罪得罪jor 上帝.
咁問題黎o下.
約伯敬畏上帝ge 動機係乜ye 呢.
係咪為jor 上帝ge 祝福呢.

$^{241}$事實上一個人即使係為jor 祝福而敬畏上帝.
都唔係錯ge 事.
邊個人唔想要祝福ge 呢.
我地讀亞伯拉罕ge 故事ge 時候.
其實都係睇到上帝答應祝福佢.
使佢多子多孫.
仲要俾塊地圖佢.
使佢後裔成為大國.
咁亞伯拉罕先至肯離鄉別井跟jor 上帝ge 指示去迦南地.
係第六至十二節就記載jor 天上ge 一幕.
呢一幕.
係唔係上主啟示作者俾佢知道呢.
我地無論作者係點樣知道發生係天上ge 事.
我地要記得故事中ge 主角約伯.
係完全唔知道有天上呢一幕.
不過讀約伯記ge 讀者就會知道.
我地睇一睇呢一幕ge 內容.
第六節我地睇見天使在上帝面前.
最後一句就話.
撒旦也來在其中.
撒旦呢個名ge 意思就係敵對者.
從呢個名我地就知道呢位ge 天使.
係專門同上帝或者同人類作對.
但上帝又冇毀滅佢.
不過到時到後開大會佢都要出席.
唔可以唔黎.
諗起《針言》十六章四字話.
「耶和華,做萬物,各適其用,就似惡人也為和幻的日子所做」.
連撒旦都係有目的.
留佢係度有目的.
但敵對就還敵對.
佢無論點樣敵對都好.
佢仍然係要受上帝ge 管轄.
包括要聽從上帝ge 命令.
係呢段ge 記載中.
我地留意到.
唔係撒旦黎到面前.
控訴約伯.
話俾上帝約伯呢個衰仔.
做jor 好多惡事.

$^{281}$唔係咁講.
反而係上帝首先問撒旦.
你係邊度黎架?.
撒旦話我冇咩做ge je .
我咪係地上誇下囉.
上帝就問佢.
既然你係地上跑黎跑去.
有冇注意到約伯呀?.
佢就指出.
地上再冇人像他那樣.
完全正直敬畏上帝遠離惡性.
上帝真係非常欣賞約伯.
係撒旦ge 面前都一樣係咁樣.
黎到特別表揚約伯.
上帝對約伯ge 表揚呢.
係呢度係同第一章一節完全一樣.
撒旦有冇反駁上帝呢?.
冇喎.
佢冇反對上帝對約伯ge 表揚.
不過佢係挑戰上帝.
佢話如果百敬畏ge 上帝唔係好正ge .
就唔係好純正.
即係約伯唔係全心全意敬畏上帝.
佢係為jor 自己ge 利益.
黎到敬畏上帝.
亦都係間接係話上帝你都蒙查查.
上帝你ge 眼光差jor D .
你未能夠洞察人心呀.
約伯唔係真正全心敬畏你架.
係十一節佢更加挑戰上帝話.
但你若伸手毀他一切所有的.
他別當面背棄你.
上帝ge 反應係點呢?.
十二節上帝話.
凡他所有的都在你手中.
只是不可伸手加害於他.
上帝接受撒旦ge 挑戰.
容許撒旦去奪取約伯所擁有ge 一切.
但係就唔准傷害佢本人.
所以呢本書討論ge 中心.

$^{321}$主要都唔係討論苦難ge 源頭.
而係討論真正ge 信仰.
一個人如果單單係為jor 自己利益.
而敬畏上帝,算唔算係敬畏上帝呢?.
如果用現代D ge 基督教術語黎講.
一個人話信耶穌.
只係為jor lor 張天堂ge 入場券.
算唔算係真係信耶穌呢?.
係約伯記一章十三節到耶二節.
我地冇讀呢段經文.
呢度就記載jor 撒旦得到jor 上帝ge 批准.
佢就真係毫不手軟出手攻擊約伯.
佢都唔使慢慢黎.
記載ge 時間一日之內.
約伯就失去jor 佢所有ge 財物同埋兒女.
約伯反映係記係一章二十一節.
佢話我赤身出於無態.
也必赤身歸回.
賞賜的是耶和華.
收取的也是耶和華.
耶和華的名是應當稱頌的.
好明顯約伯冇係由於遭遇到苦難.
而離開上帝或者埋怨上帝.
反而仍然係度稱頌上帝.
二十二節就係作者注究.
唔係上帝講唔係約伯講.
作者注究話.
作者一看見上帝約伯沒有犯罪.
沒有歸咎於上帝.
呢個係橫性ge 翻譯.
但係呢一場ge 上帝同撒旦ge 博弈當中.
我地知道撒旦輸jor .
撒旦挑戰上帝.
你lor 晒佢D 財物.
佢就當面背棄你.
但約伯冇咁做.
冇背棄上帝.
仲稱頌上帝.
撒旦就輸.
但佢唔服氣.

$^{361}$輸還輸.
佢未服氣.
所以就有第二章.
天庭大會ge 第二幕.
我地睇睇.
第二章頭兩節同先前ge 天庭大會ge 記載一樣.
我地睇第三節.
首先上帝仍然係讚賞約伯.
話佢係完全正直.
敬畏上帝.
遠離惡事.
呢次已經係第三次.
上帝讚賞約伯ge 為人.
然之後三章ge 下半節上帝繼續講.
佢話呢.
你需激起我攻擊他.
無故吞滅他.
他仍然持守他的純正.
係第一章其實我地睇到攻擊約伯ge .
係撒旦.
但呢度上帝話係我攻擊他.
雖然唔係上帝親自出手去攻擊約伯.
而係撒旦攻擊佢.
但由於係上帝批准.
都等於係上帝攻擊約伯.
不過上帝亦都指出約伯仍然係持守佢ge 純正.
咩純正?.
佢信心呀.
佢信心好純正.
約伯敬畏上帝ge 心係純正.
唔係由於上帝賜比佢有兒有女.
有豐富ge 物質.
而係真係全心全意.
由於上帝係上帝.
而敬畏佢.
但撒旦唔服氣.
佢又再挑戰上帝.
但你若伸手傷他的骨頭和他的肉.
他必當面助救你.
呢個亦都係環聖ge 翻譯.

$^{401}$比較容易明白je .
上帝對約伯真係充滿信心.
佢又批准撒旦去攻擊約伯.
但係又比個規範佢.
二章六節約華對撒旦說.
看下他在你手中.
你只要留下他的性命.
你想整佢都得.
唔可以殺jor 佢.
一定要保存佢性命.
二章七節話比佢聽.
撒旦就馬上去行動.
擊打約伯.
一D 都唔手軟呀.
約伯ge 身體.
由腳掌到頭都掌滿毒瘡.
係咩瘡我都唔知道.
但係呢個瘡有個後果.
就令到約伯全身都好癢.
我諗起上次.
gwo 位癌腫科醫生話比我聽.
我可以同你用標靶藥.
不過呢.
標靶藥都會有副作用.
你可能你D 皮膚會好癢.
好似約伯咁.
因為我呢位癌腫醫生.
都係基督徒.
我會行到好似約伯咁.
但係結果都係用jor 隻標靶藥.
不過最近檢查ge 時候.
就顯示呢隻標靶藥都失效.
醫生就要.
要我用另外一D 強力D ge 化療藥.
求主憐憫.
請大家都為我祈禱.
二章九節.
約伯ge 太太出場.
一直以黎.
由第一章開始.

$^{441}$我地冇見到約伯ge 太太.
一見到佢呢.
就俾我地作為讀者ge 印象.
都係好負面.
佢好似係.
撒旦ge 爪牙咁樣.
佢走黎勸呀約伯.
你背起上帝死了吧.
即係係佢心目中諗住.
你上帝都唔保守你.
arm arm 先發生jor 咁多悲慘ge 事.
子女又死哂.
錢財又冇哂.
依家你仲咁樣.
本身都係瘡.
可能gwo D 瘡仲有D 臭味.
佢就非常嫌棄約伯.
勸佢不如你走就死jor 算數.
第十節.
約伯就鬧佢個妻子.
愚頑的婦人.
仲話呢.
難道我地從上帝手裡得福也不也受禍嗎?.
呢句話表示約伯.
認為佢所受ge 禍都係出於上帝.
並且表明佢ge 神學思想裡面.
人生ge 福與禍都係出於上帝.
我地應該接受上帝所賜ge 一切.
唔係只係為求得福.
一遇到禍患.
就連上帝都唔要.
作者最後一句補充咁話.
在這一切的事上.
約伯並沒有異口犯罪.
約伯真係一位全心全意敬畏主ge 人.
唔係為jor 得福.
唔係為jor 得財利.
先至跟隨呢位上帝.
約伯記要討論ge 重心.
就係點樣ge 信仰.

$^{481}$先至係真ge 信仰.
如果外在環境對我地不利.
我地會唔會否定上帝.
懷疑上帝ge 良善同信實呢?.
係以色列人ge 歷史當中.
其實我地時時都睇到佢地.
每當環境不順ge 時候.
佢地都懷疑上帝.
係《說書記》六章十三節.
基顛對向佢顯現ge 天使講.
主呀請容許我說.
時華若與我地同在,我地怎麼會遭遇這一切事呢?.
我唔得列祖告訴我地.
耶和華領我地從埃及上流.
他話奇妙ge 作為在哪裡呢?.
現在耶和華卻丟棄了我地.
把我地交在米顛人ge 手掌中.
呢個好早期ge 事.
到猶太人末期D ge 事又點呢?.
末期D ge 時候係《耶利米說》四十四章十五至十八節.
咁講,咁樣記載.
我那些知道自己妻子向別臣燒香的男人.
與暫在那裡的一大群婦女.
就視著埃及地巴特羅所有ge 百姓.
回答《耶利米說》.
論到李鳳耶和之名向我地所說的話.
我地必不聽從.
我地定要照我地口中所說ge 一切話去做.
向天后燒香.
獻曉酒祭.
按照我地與我地ge 祖先君王官長在猶大城鎮.
和耶路撒冷街市訴償所做都一樣.
因為我地得意吃飽.
因為那時我地得意吃飽.
享福樂並未遇見災禍.
自從我地停止向天后燒香.
獻曉酒祭.
我們都缺乏這一切.
猶因刀劍饑荒滅絕.
這些人都有一些錯誤的神學思想.

$^{521}$好像與上帝交往是一種交易.
人向上帝燒香獻祭物.
上帝就應該供應我這一切.
保護我讓我享福.
如果不是的話,要這個上帝做什麼?.
轉去拜其他神明.
今年我們教會主題是「全心全意」.
我鼓勵大家自我反省.
我們的信仰是否一種功利主義交易式的信仰.
應該好像約伯一樣.
明白到賞賜的是耶和華,收取的也是耶和華.
難道我們從上帝手裡得福也不受禍嗎?.
無論我們遇到什麼樣的境況.
我們都會忠心,全心跟隨上帝.
我們一起祈禱.
更為我們的聖名,寫名.
拯救我們出黑暗,進入奇妙光明的主.
求你幫助我們認識你更多,愛你更多.
不只是為了得福,而是為了回報你的恩典.
以我們的生命事奉你,是理所當然的.
奉耶穌基督的名祈禱,阿們.
\newpage



\section{加拉太書 4:28-31}
\label{sec:031626}
\textbf{作屬靈的兒女 (普)}~~~~葛國光牧師
\newline
\newline
連結: \href{http://www.church.com.hk/acms/content.asp?site=cdc\&op=show\&type=product\&code=031626\&layout=sermon}{\texttt{031626}} ~~~~ 語音日期: 2026-02-08
\newline
\newline
\hyperref[sec:031625]{\small{< < < PREV SERMON < < <}}
~
\hyperref[sec:index_chronic]{\small{[返順時目]}}
~
\hyperref[sec:031627]{\small{> > > NEXT SERMON > > >}}
\newline
\newline
加拉太書 4:28-31
\newline
\begin{longtable}{cl}
\hline
\hline
章節 & 經文 (和合本修訂版)\\
\hline
4:28 & \begin{tabularx}{0.7\textwidth}{X} 弟兄們，你們是憑著應許作兒女的，如同以撒一樣。 \end{tabularx} \\ \\ \relax
4:29 & \begin{tabularx}{0.7\textwidth}{X} 當時，那按著肉體生的迫害了那按著聖靈生的，現在也是這樣。 \end{tabularx} \\ \\ \relax
4:30 & \begin{tabularx}{0.7\textwidth}{X} 然而經上是怎麼說的呢？是說：「把使女和她兒子趕出去！因為使女的兒子絕不能與自由婦人的兒子一同承受產業。」 \end{tabularx} \\ \\ \relax
4:31 & \begin{tabularx}{0.7\textwidth}{X} 弟兄們，這樣看來，我們不是使女的兒女，而是自由婦人的兒女了。 \end{tabularx} \\ \\
[1ex]
\hline
\hline
\end{longtable}
$^{1}$今天我們要跟大家回到我們去年的要講的加拉太書..
哇,加拉太書只有六章,但是你看神的話這麼豐富,字字珠璣..
加拉太書,我們今天要講到加拉太書最後的這一段,.
加拉太書第四章二十八到三十一節..
我們一起來讀這段聖經,好不好?.
好,我們起來,預備,請..
弟兄們,我們是憑著應許做兒女,如同以下一樣,.
當時按著血氣生的,逼迫了按聖靈生的,現在也是這樣..
然而經上是怎麼說的呢?.
是說把使女和她兒子趕出去,.
因為使女的兒子不可以自主婦人的兒子一同承受產業..
弟兄們,這樣看來,我們不是使女的兒女,乃是自主婦人的兒女了..
好,我想在這一段聖經裡面我們稍微瞭解一下,.
保羅在前面裡面講什麼呢?.
他一直在強調的就是針對當時候一些猶太教的基督徒,.
或者是說他們其實還是猶太教,.
那麼他們把猶太教的觀念裡面,就是律法的觀念帶進到基督教裡面來..
所以他們說,你們單單信耶穌不夠,你們還要怎麼樣呢?.
你們還要受割離,你們還要做這個做那個,.
你們要行律法上面的這些事情..
所以保羅非常嚴厲的跟這些加勒賽蘇的弟兄們講說,.
不可以,因為他們所傳的是別的福音,看到沒有?.
那就是唄,所以保羅一直在強調一件事情..
我還記得我們曾經講過一篇信息,講到說我們是誰的兒女啊?.
神的兒女,對不對?.
那麼你怎麼樣子才能夠做神的兒女呢?.
保羅在這裡一開始就說,我們是憑著什麼做兒女?.
硬去做兒女..
然後他接著就講了一個例子,.
他說如同以撒一樣,所以你要對舊約稍微有點概念..
我們來看看,在這個地方講到兩種的兒女,.
第一種的就是所謂的按著血氣生的,有沒有?.
剛纔我們讀了那句話說,.
他說按著血氣生的,逼迫了按著聖靈生的,.
所以第一種叫做按著血氣生的,.
第二種就是叫做按著聖靈生的..
那麼聖經裡面讓我們看見,.
這個按著血氣生的,逼迫了按著聖靈生的呢,.
就可以從聖經裡面我們可以看見,是不是?.
那按著血氣生的是誰?.

$^{41}$以史瑪利.按著聖靈生的是什麼?是以撒..
以史瑪利逼迫以撒..
直到今天,你會發現中東的問題都要推溯到這個聖經裡面來..
當然這個問題我已經講過幾次,.
這個問題是很複雜的一個問題,.
可是不管你喜歡不喜歡,不管你瞭解不瞭解,.
你會發現整個的歷史裡面,.
猶太人跟中東阿拉伯人的歷史的恩恩怨怨,.
沒完沒了..
這個恩恩怨怨會一直到什麼時候?.
到主耶穌基督再來的時候..
如果你仔細去觀察,你去了解,.
你看起書裡面,看聖經裡面,.
你會更清楚的明白..
這就是鼓勵同志們好好讀聖經..
有時間我再跟大家繼續分享..
那麼我們來看看,.
在這裡講到這兩個人呢,.
我們首先來看看,.
保羅在前面裡面我們已經講了很幾次,.
一個是以史瑪利,一個是以撒,對不對?.
以史瑪利是按著血氣生的..
這個血氣就是說,.
他不是按照上帝的應許而生的,.
是按照人意而生的,.
按照天然而生的..
那麼以撒呢,.
他是按照聖靈而生的,.
也就是按照神的應許生的..
上帝跟亞伯拉罕說,.
你要生一個兒子,.
這個兒子要繼承你的產業,.
這個兒子是誰?.
但是亞伯拉罕連撒他也自己發現,.
自己已經過了生育年齡了,.
沒有辦法再生了,.
他生不出來了,怎麼辦呢?.
他自己想了個點子,.
找誰生呢?.
他就找夏甲去生..

$^{81}$夏甲是誰?.
十里,十里是什麼?.
是切..
切,切都還不算..
好,那我們來看看,.
亞伯拉罕怎麼會這樣子,.
糊塗到這種地步,是不是?.
你會發現,.
人到了一個地步,.
很多時候看到絕望了,.
他也沒信心了..
所以,.
以斯瑪利是按血契生的,.
以撒是按聖靈生的,.
然後呢?.
這裡告訴我們啥?.
這個以斯瑪利,.
就是在預表著,.
他是猶太教的..
這是亞伯拉罕的子孫,.
所謂的肉身的子孫..
這個意思就是說,.
保羅用這個例子來預表說,.
今天猶太教徒,.
他還是活在這種律法的觀念裡面,.
活在這種血契的裡面..
那麼可是呢?.
以撒的後裔呢?.
預表著今天我們新約的基督徒,.
我們不再是按照律法了,.
我們是屬靈的,.
因著聖靈的緣故,.
所以我們做了亞伯拉罕屬靈的後裔,.
懂這意思嗎?.
然後呢?.
他們是憑血契生的,.
憑肉體生的,.
以斯瑪利..
那麼,.
以撒呢?.

$^{121}$是預表著基督..
這是一個對照,.
這個對照讓我們看見很清楚的.
以斯瑪利跟以撒的分別..
那我現在要跟他講的,.
剛纔保羅在這個地方強調的說,.
按血契生的,.
按聖靈生的..
所以他這邊講了幾種人啊?.
幾種人?.
一個是按血契生的,.
另外一個是按聖靈生的,.
有兩種人,.
對不對?.
那請問你,.
你是按血契生的,.
還是按聖靈生的?.
你問一下你旁邊..
你是按血契生的,.
還是按聖靈生的?.
那我們首先來看看,.
什麼叫做屬血契的?.
哥林多前書第二章第十四節,.
這裡說,.
然而屬血契的人不領會.
神聖靈的事,.
反倒以為愚拙..
所以這一個屬血契的人,.
指的那些不相信神的人,.
你跟那個不信主的人,.
可能是你的家人,.
可能是你的朋友,.
你跟他講屬靈的事情,.
他聽得懂嗎?.
禮拜天要聚會,.
太太還不信主,.
先生會怎麼說?.
不信主的家人會怎麼說?.
禮拜天你跑教會去,.
你迷了,.

$^{161}$你迷信了,.
你跟他講說,.
他聽得懂嗎?.
聽不懂,.
因為屬血契的人,.
不能夠理會神聖靈的事情..
好了,第二種人,.
按著聖靈生的..
但是按著聖靈生的,.
這裡是指的什麼呢?.
你會發現按著聖靈生的,.
保羅在這裡講到有兩種,.
第一個是什麼?.
叫做屬肉體的..
這個屬肉體的人,.
他們也相信神,.
但是他們在基督裡面是什麼?.
這是陰海,.
這是屬肉體的..
在另外一方面,.
講屬靈是什麼呢?.
這裡說什麼?.
他說屬靈的人能夠看透萬事,.
有沒有?.
丟子妹,.
不是因為你PhD,.
不是因為你有很好的豐富知識,.
屬靈的人,.
可能一個在中國的老太婆,.
一個老姐妹,.
她從來沒有讀過書,.
大志都不是一個,.
但是她卻是屬靈的,.
她能夠參透萬事..
你懂我的意思嗎?.
所以保羅在這裡告訴我們說,.
我們要做屬肉體的呢,.
還是做屬靈的人呢?.
我們來看看,.
這有一個表,.

$^{201}$左邊最左邊是屬血氣的,.
中間是屬肉體的,.
右邊是屬靈的..
這個表呢,.
我想屬血氣的跟屬靈的比較沒問題,.
中間這個屬肉體,.
我稍微說明一下,.
因為他這邊講的可能有點misleading,.
他這裡說是不信靠神的基督徒,.
我想這句話,.
因為這個不是我製作的表,.
所以我不知道怎麼去刪它,.
如果你們懂的話可以教我,.
怎麼把它換成對的..
應該這樣講,.
我們嚴格來講,.
我們就分兩邊,.
一個就是信神,.
一個不信神..
第一個區分,.
我們先把這兩個分出來..
接著第二個區分,.
如果一個人信了耶穌之後呢,.
你可以看見在保羅整個的書信裡面他提到,.
一個人他如果信了耶穌,.
可是他是不是就變成是一個成熟的,.
長大成人的基督徒呢?.
不是,.
是不是?.
所以保羅才告訴我們說要喫乾糧,.
對不對?.
不要像小孩子喫一樣,.
要能夠喫乾糧,.
不要只是喫奶的嬰孩,.
是不是?.
跟你旁邊說你不要做喫奶的嬰孩..
每個禮拜天來這邊等牧師拿個大奶瓶,.
然後你就喫了一滴,.
就靠這一滴來維持我這一個禮拜的生活..
千萬不要做喫奶的嬰孩..

$^{241}$所以保羅在這邊其實在講到一個人信了耶穌之後,.
我們裡面就有一個屬靈的生命,.
這個生命它需要不斷地成長,.
但是有的人他不成長,.
是不是?.
所以就還是做嬰孩..
有的人持續地成長,.
參加門訓,.
不單單只是知識的追求,.
他也去操練實際的追求,.
他的生命就越來越像耶穌基督,.
這就成為屬靈的人..
講到這地方,.
我就想到最近可能你也注意到,.
美國的司法部把有一個叫做艾普斯坦,.
大家都知道嗎?.
把他的文件,.
聽說有幾百萬,.
你可以看看,.
我也想這個人真是厲害,.
他幹這些事情,.
他還都把它記錄起來..
Anyway,.
這個公佈裡面,.
連我女兒最近都在問說,.
Daddy, 你的看法如何?.
就是你能不能想象,.
這些人,.
他們不認識神,.
所做的事情,.
簡直是連禽獸都不如,.
是不是?.
但是這些公佈出來的人當中,.
有沒有基督徒?.
也有啊!.
誰是基督徒?.
川普是基督徒嗎?.
這是一個問號,OK?.
還有,你說還有誰?.
還有誰?.

$^{281}$比爾蓋茨,比爾蓋茨是基督徒嗎?.
我都幫他Google了一下,.
AI告訴我,.
還有一個人叫做克林頓,有沒有?.
克林頓,.
如果你去Google一下,.
AI會告訴你說,.
克林頓是,他是一個基督徒..
大家基督徒眼睛馬上亮起來了,.
他是一個基督徒,.
他是一個長期的美蘭進信會的基督徒,.
他在擔任總統期間經常參加教堂禮拜,.
他從小在進信會長大,.
十歲就接受洗禮,.
在格里汗佈道大會當中,.
對他的生命產生了重大影響..
他的太太,Hillary,也是一個基督徒,.
是虔誠的衛理工會的信徒,.
這是AI說的..
那我請問丁五姐妹,.
但是這樣子一個人,.
克林頓發生什麼事情?.
在他做總統的事情,.
有沒有?.
跟他的實習生髮生了問題,是不是?.
那再加上這個東西出來,.
所以很多人就對他說他是一個基督徒嗎?.
打上了一個很大的問號,是不是?.
請問奧巴馬是不是基督徒?.
不是啊,為什麼?.
你怎麼說他不是基督徒?.
你說他回教徒嗎?.
這是誤會,.
AI說你誤會他了,.
AI說他是一個基督徒..
所以我們今天要把基督徒這個定義搞清楚一點,.
等一下我會跟大家分析..
在剛纔講的這段信息裡面,.
保羅說,現在也是這樣,.
按著血氣生的,逼迫的按著聖靈生的,.

$^{321}$現在也是這樣,.
然而經上是怎麼說的呢?.
是說,把使禮和他兒子趕出去,.
這段信息裡面是在講什麼呢?.
講到創世記,.
在創世記裡面提到這一段故事,.
當以赦他到了一個斷奶的日子,.
請問你們這兒有小孩的,.
什麼時候斷奶?.
猶太人的斷奶的日子就是指的三歲..
太晚了?.
好,不管了,那是猶太人..
這個孩子漸漸長大,.
聖經在創世記二十一章第八節第十節就說了,.
孩子漸長就斷了奶,.
在這裡講說,.
亞伯拉罕就色擺豐盛的宴席,.
當時候撒拉看見埃及人下架給亞伯拉罕所生的兒子西孝,.
就對亞伯拉罕說,.
你把這使禮和他兒子趕出去,.
因為這使禮的兒子不可與我的兒子以撒一同承受產業,.
所以這是從創世記二十一章這個地方來的..
這個上下文裡面就是說,.
當以撒三歲他斷奶的時候,.
那以實瑪利幾歲呢?.
以實瑪利比以撒大十四歲,.
所以請問以實瑪利幾歲?.
加三歲,十七歲..
所以以實瑪利當時候看見爸爸亞伯拉罕在幫這個兒子,.
他的兄弟在擺設宴席的時候聲音上說,.
下架的兒子以實瑪利西孝,西孝是什麼?.
雞孝,取孝他..
所以撒拉就說什麼?.
把這個使禮和他兒子趕出去!.
好狠啊!.
但是你知道嗎?.
在屬靈上面的意義啊,.
他意思就是說啊,.
因為以實瑪利他不是按著應許生的,.
他是屬肉體,他是屬血氣生的,.

$^{361}$所以把他趕出去的意思就是什麼呢?.
就是把那一切的血氣肉體趕出去..
西孝就說什麼?.
在我們的裡面有沒有血氣?.
有沒有肉體?.
有..
這個血氣肉體,我們來看一下,.
我們來讀這段聖經好不好?.
我們讀右邊這個,迦拉太書的,.
這個是迦拉太書的第五章,.
我們來讀這一段,.
情慾的事好不好,我們一起來讀,.
情慾的事都是顯而易見的,.
嫉妒,醉酒,荒諺等等,.
我從前告訴過你們,.
現在又告訴你們,.
行這樣事的人,必不能承受神的果..
這裡提到說,.
行這樣事的人,必不能夠承受上帝的果,.
很嚴重,對不對?.
好,那請告訴丁毅子妹們,.
你有沒有淫亂?.
沒有..
沒有?.
OK..
從來沒有?.
你有沒有誤會?.
有沒有放蕩?.
有沒有拜偶像?.
新邪書?.
有沒有仇恨?.
有沒有?.
有沒有紛爭?.
有沒有嫉恨?.
有沒有憤怒?.
有沒有自私?.
有沒有紛爭?.
有沒有結黨?.
有沒有嫉妒?.
有沒有?.

$^{401}$你可能在那邊選,.
刪掉,刪掉,刪掉,.
但是有沒有這當中你所有都沒有的?.
有沒有?.
請舉手..
沒有..
會不會嫉妒?.
會不會?.
我不嫉妒..
那些被我嫉妒的人都死了..
會不會嫉妒?.
會啊,誰會不會?.
爸爸喜歡老大,你是老二,.
還有一個最小的,.
他就是喜歡老大,老腰又很疼,.
那就是老二怎麼樣?.
不上不下?.
會不會嫉妒?.
親愛的同事們,.
如果照這樣子的標準,.
你我都死了,.
你我都沒有資格承受神的國,.
對不對?.
但是好消息,.
來,我告訴你好消息,.
這就是你需要好好讀聖經,.
在原文裡面,.
這裡說的行這些事的人,.
這邊把它翻成做,.
不管是行或就是do,.
這個字原文是現在分詞,.
什麼意思?.
表明說是一個慣性的,.
持續的去做這樣子的事情的人..
那慣性的持續的會去做這樣子的人,.
其實他根本就等於是沒有重生,.
懂這意思嗎?.
所以親愛的同事們,.
我們會有嫉妒,.
我們會有紛爭,.

$^{441}$可是這些在我們整個追求成聖的過程當中,.
聖經裡面一再地告訴我們,.
要把這些怎麼樣棄絕?.
保祿告訴我們說,.
這些事情等等等等等等,.
要棄絕它,.
要拒絕它,.
一切的苦毒論斷都要棄絕它,.
所以你看消極的一方面來講,.
聖經告訴我們說,.
是的,我們還會紛爭,.
我們還有嫉妒,.
我們還有這些仇恨,.
我們還會有這些的等等,.
不討神喜悅的這些的誤會,.
但是我們要做一件什麼事情?.
我們要棄絕,.
我們要拒絕它,.
我們要把它趕出去,.
跟你旁邊說,.
把一切的紛爭嫉妒趕出去,.
把一切的爭吵從你的家趕出去,.
從你的心中趕出去,.
從那個驕傲,.
把你心中的恣意趕出去,.
如果你這一點沒有看見,.
那你還沒有真正的認識神,.
一個真正認識神的人,.
他會看見這些是不討神喜悅的,.
如果一個人他對這些事情他沒感覺,.
甚至他覺得我就是這樣子,.
那反過來就是他的生命是有問題的,.
所以保羅在這裡說.
要把使禮和他兒子趕出去,.
應用在我們今天的生活裡面,.
在我們的生命當中,.
意思就是說,.
在我們的生命當中,.
我們常常要做一件事情,.
就是拒絕,.

$^{481}$就是要把這一些不合神心意的誤會敗壞,.
這些思想意念趕出去,.
把這些不合神心意的,.
不一定是這些看起來是罪惡的,.
包括憂慮煩惱,.
這些也是不合神心意的,.
懂我意思嗎?.
包括不合神心意的,.
不是神旨意的這些錯誤的世界上的觀念,.
我們要把它趕出去,.
前面一個信息我曾經提到說,.
比如說我們會以功成名就,.
會以我們長相,.
會以我們的外表來去看一個人的價值,.
會不會?.
特別在今天整個社會里面,.
你看看每一個網紅都怎麼樣,.
光光鮮鮮的,.
你看到這些的時候,.
他全都是讓你感覺,.
哇,他真是美啊,.
他真是怎麼樣子啊,.
所以這個就變成是我們基督徒的價值觀,.
但是親愛的同志們,.
聖經上告訴我們說,.
要把這一切錯誤的東西趕出去,.
不要再用外表,.
不要再用學業的成就,.
不要再用一個人他的身份地位,.
去來判斷這個人的價值,.
也不要用這些來判斷我們自己,.
那麼第二個,.
我們積極的一方面,.
我們要做屬靈的兒女,.
耶穌告訴我們說,.
實在告訴你,.
人若不是從水和森林生的,.
就不能夠進神的國,.
從肉生生的就是肉生,.
從靈生的就是靈,.

$^{521}$看到沒有?.
那麼中國人有一句俗話說,.
叫做龍生龍,.
鳳生鳳,.
老鼠生的兒子,.
會打洞,.
其實這句話是從古書出來的,.
它的原文不完全是這樣子,.
但是意思差不多,.
丟子輩,.
請問老鼠生的兒子會打洞?.
是不是?.
這什麼意思?.
對,.
是什麼樣子就什麼樣子,.
你是什麼生的你就什麼樣,.
這就是從血氣生的,.
是不是?.
所以耶穌說,.
從肉生生的就是肉生,.
從靈生的就是靈,.
因此我們需要重生,.
從哪裡生呢?.
從聖靈生,.
丟子輩,.
這是第一步,.
血氣的跟聖靈生的,.
分別就在這個地方,.
一個重生得救的基督徒,.
他就是從聖靈生的,.
也就是憑應許做上帝的兒女,.
有聖靈在我們心中做印記,.
證明我們是屬乎神的,.
丟子輩,.
請問,.
你怎麼知道你到底重生了沒有?.
你怎麼知道你重生了沒有?.
我要問大家一件事情,.
你知道你自己重生了沒有?.
好像都是半隻手,.

$^{561}$是一半去的,.
高一點,.
來,.
好,.
請放下,.
剛纔我們談到的這些名人,.
是不是?.
克林頓,.
他說他是基督徒,.
奧巴馬,.
他也說他是基督徒,.
我們再推更早一點,.
還有誰?.
克林頓,.
克林頓在前面是誰?.
布希,.
布希是不是基督徒?.
肯定是,.
是不是?.
好,.
那請問川普是不是?.
他自己說他是不是?.
他說他是,.
根據統計,.
美國在2024年統計大概有63\%或者64\%的美國人自稱他是基督徒,.
這個比例已經下降,.
在2007年的時候,.
這個比例是78\%,.
好了,.
那麼另外有一個統計,.
他問得更直接了,.
他說你認為你自己是重生的基督徒嗎?.
你看那個藍色的有多少人?.
有29\%的人,.
他自認為自己是重生的基督徒,.
因此,.
丁永子妹,.
這讓我們看清一件事情,.
是不是一個基督徒,.
一個重生的基督徒,.

$^{601}$在這裡是關鍵了,.
是不是?.
這很重要了,.
我們看到這些事情的時候,.
我們就會說這個人也是基督徒,.
這個人也是基督徒,.
那好像基督徒沒什麼了不起,.
基督徒一樣犯罪,.
而且一直做這些甚至連世界上的人都不齒的事情,.
是不是?.
所以到底一個人他怎麼樣子是基督徒呢?.
好,.
所以我們要看見,.
怎麼樣子才能夠成為重生的基督徒?.
第一個,.
約翰福音第三章,.
尼哥底姆問耶穌的時候,.
耶穌剛纔講了一句話很清楚,.
他說,.
從靈生的就是靈,.
從肉生生的就是肉生,.
對不對?.
人若不是從水和聖靈生的,.
就不能夠進神的國,.
所以那個地方講到聖靈的工作,.
耶穌說,.
風隨著自己的意識吹,.
所以聖靈的工作也是如此,.
請問聖靈叫一個人重生,.
你看見嗎?.
你看得見嗎?.
看不見,.
但是你會知道,.
這是聖靈的工作,.
當年山東大復興的時候,.
上帝使用一個僕人叫做孟教師,.
他常常去到中國各地帶領聚會的時候,.
他每次聚會都要問一個問題,.
你重生了嗎?.
因為當時候的教會充滿了許許多多,.

$^{641}$搞不清楚自己信仰的狀況,.
很多這些沒有重生的人,.
一樣在教會里面做執事,.
做長老,.
甚至做牧師,.
所以在那個情況之下,.
孟教師常常在講到要問這句話,.
那丟什麼你知道嗎?.
有一次有一個牧師,.
他怕這個孟教師問他,.
他一看他聚會結束,.
他怕他要問你重生了沒有,.
所以他就趕快趁著在結束禱告的時候,.
他就怎麼樣,.
一個箭步就溜出去了,.
結果他怕從門口去會被人家看到,.
他就從窗戶跑出去,.
那時候的窗戶是有點像是落地窗,.
有點高度,.
但是可以一跨就過去,.
他把窗戶一開,.
他就跨出去了,.
沒想到當他一跨出去,.
孟教師一個箭步就在他面前出現,.
就問他說你重生了嗎?.
兄弟們,.
你重生了嗎?.
淪落不重生就不能夠進神的國,.
聖靈工作,.
我沒有時間講了,.
聖靈工作的時候藉著孟教師,.
許多的人他以前一直覺得自己已經是得救的基督徒,.
他覺得他自己的生命,.
他覺得他自己的行為都很好,.
他也不覺得自己有什麼罪,.
但是當聖靈光照的時候,.
他就悔改在神的面前,.
認罪,.
承認自己的罪,.
聚會常常就是如此,.

$^{681}$聖靈大大的工作,.
很多的人多年來經歷到那種被從罪裡面得釋放的那種喜樂,.
親愛的同志們,.
這是聖靈所需要的工作,.
這是聖靈所做的工作,.
另外一個,.
從約翰福音三章十六節,.
我們都很懂的,.
都會背在這些聖經裡面,.
告訴我們說什麼呢?.
神愛世人,.
聖旨將他的獨生子賜給他們,.
叫一切信他的,.
不致滅亡,.
反得..
所以什麼樣子的人不致滅亡,.
反得永生?.
信他的人..
親愛的同志們,.
信好,.
問題就在這裡了,.
我們今天常常在講信好,.
所以美國的基督徒也是如此,.
為什麼自稱叫基督徒?.
因為我們所謂的信是什麼呢?.
就是他願意禱告,.
接受耶穌基督,.
就是信了..
但是我告訴你,.
我們中國人很會這一套,.
你要我做什麼,.
我就做什麼,.
你懂我意思嗎?.
甚至你去傳福音,.
你會發現,.
很多時候這些中國人,.
他覺得你太煩了,.
但是要把你打發走,.
也不太容易,.
他怎麼辦呢?.

$^{721}$我信,我信,.
然後你說,.
好,那我代你禱告,.
好,禱告,.
這一樣話,.
我給你畫個葫蘆了,.
禱告完之後怎麼樣?.
他們纔不信你這個教,.
很多人是這樣子的,.
包括來到教會里面的,.
有沒有這樣子的人?.
有..
所以為什麼在受洗的時候,.
我常常提醒,.
你如果沒有真正的重生得救,.
你一百遍都沒有用..
好,關鍵在這個信了,.
信它的不致滅亡,.
信是什麼?.
在原文裡面這個信,.
它後面加了一個叫做信路,.
就是believe into,.
這個代表的是一個信是什麼?.
不是頭腦知識,.
這個代表的是一個真正的偽身的信,.
懂這意思嗎?.
第二個,.
原文裡面這個信的實態也是現在分詞,.
什麼意思?.
所以顯然我們過去一個人說,.
在那個moment,.
他說我信耶穌,.
我們就認定他已經信耶穌得救了,.
這樣子的認定法是不準確的,.
懂我意思嗎?.
這個信代表的是一個持續性的,.
不是隻是那一個moment,.
我就做了一個絕世禱告而已,.
所以你怎麼知道他重生得救沒有呢?.
需要時間,.

$^{761}$從信,.
以至於信,.
我們的信仰,.
這也提醒我們,.
我們不要把我們的信仰,.
停留在只是信耶穌受洗的那個階段,.
我們這一生都需要持續from faith to faith,.
從信,.
以至於信,.
一直不斷地在信仰的信心上面,.
不斷地不斷地追求,.
使得我們對神的信心越來越增加,.
最後,.
我們講到做屬靈的兒女,.
怎麼做呢?.
重生得救是第一步,.
接著是什麼呢?.
大衛在這邊的一個禱告,.
在詩篇27篇,.
我們一起來讀,.
有一件事,.
我曾求耶和華,.
我仍要尋求,.
就是一生一世在耶和華的殿中,.
來瞻仰祂的柔美,.
在祂的殿裡求問,.
大衛說我一生一世,.
我有一個禱告,.
我有一個渴慕,.
我就是要在神的殿中,.
神的殿,.
在就業的時候,.
神的殿代表是神所在的地方,.
所以呢,.
今天神在哪裡?.
神在哪裡?.
在我們心裡,.
是不是?.
所以當我們說,.
我們來清靜神的時候,.

$^{801}$早上六點鐘,.
是不是要到禮拜堂來?.
是不是?.
是嗎?.
不是..
你在你家裡面的書房,.
可不可以清清楚?.
可以,.
是不是?.
總之,.
這裡大衛的禱告說,.
我一生一世要住在耶和華的殿中,.
瞻仰祂的柔美,.
請注意這兩件事情,.
第一個,.
瞻仰祂的柔美,.
第二個是什麼?.
在祂的殿裡求問,.
求問就是禱告,.
大衛凡事求問神,.
在他一生當中,.
他一直是這樣子的一個態度,.
他求問神,.
當非利式人來攻擊基拉的時候,.
大衛連續的兩次求問神說,.
可以去攻打嗎?.
我們一般打了就是了,.
但是大衛在打仗的時候,.
他問說,.
主啊,.
可以打嗎?.
可以去打嗎?.
連續兩次,.
後來在喜格拉,.
亞馬利人擄走了大衛的家人跟財物,.
大衛在焦急的當中,.
他也求問神,.
他說,.
我追趕敵軍,.
追得上追不上的?.

$^{841}$哎呀,.
你會覺得大衛,.
come on,.
你這是一個軍人,.
你是打仗已經很有經驗的,.
你還問說追得上追不上?.
還問上帝嗎?.
就像我們今天很多的知識分子,.
我們在教會里面信了耶穌,.
我們也是這樣子,.
哎呀,.
這個common sense,.
還要求問神嗎?.
五年計劃,.
十年計劃擺出來,.
大家break storming,.
一下子我們就出來一個五年計劃了,.
幹嘛還要求問神呢?.
所以今天教會有沒有屬血氣的基督徒,.
屬肉體的基督徒?.
一樣有,.
凡事不按照神的心意來做事的,.
就是屬肉體,.
大衛凡事尋求神,.
所以他一生一世渴望,.
要能夠在神的殿中,.
那個在新約的意思就是什麼?.
渴望住在主的裡面,.
聽懂了沒?.
那是我們門心班每一個聽懂了沒?.
即使你LT畢業了之後,.
你一生要追求的目標就是什麼呀?.
一生一世住在主的裡面,.
每一天來慶敬他,.
不是成島,.
這個門訓的時候起來成島,.
因為要交報告,.
那門訓結束了呢?.
哎呀,.
沒人催我交報告了,.

$^{881}$你知道嗎?.
我有一個侄女在臺灣,.
她去過了幾個教會,.
樹林光景一直都不穩定,.
前一段時間,.
我介紹她去了一個我的學弟的教會,.
她去了之後,.
她跟我講,.
她說,.
我們那個小組長,.
天天的來,.
要不然就是text,.
她一天不知道十幾個text,.
或者是line,.
她就是來問我,.
她說,.
怎麼樣,.
她也不是逼你,.
她也不是push你,.
她就是很關心你,.
哎呀,.
怎麼樣,.
我的這個就叫Lacey,.
你最近這個離休的筆記寫了嗎?.
哎呀,.
沒關係,.
加油,.
我知道你很忙,.
然後過了一下子又來了,.
這叫做愛的捆綁,.
所以她跟我分享的時候,.
我真是為她高興,.
她到了一個好教會,.
Amen?.
盼望我們的小組長都這樣子,.
盼望我們教會弟兄姊妹們,.
我們都是這樣子帶門訊的,.
你懂意思嗎?.
見了她說,.
哎呀,.

$^{921}$最近怎麼樣,.
離休有沒有啊?.
哎呀,.
最近起來沒有啊?.
弟兄姊妹,.
我告訴你,.
我們在這裡,.
就是這樣子提醒彼此的追求,.
因為一生一世住在她的店中,.
是何等的美好?.
Amen?.
\newpage



\section{撒母耳記上 17:1-58}
\label{sec:031627}
\textbf{你也可以做巨人}~~~~李思聰牧師
\newline
\newline
連結: \href{http://www.church.com.hk/acms/content.asp?site=cdc\&op=show\&type=product\&code=031627\&layout=sermon}{\texttt{031627}} ~~~~ 語音日期: 2026-02-08
\newline
\newline
\hyperref[sec:031626]{\small{< < < PREV SERMON < < <}}
~
\hyperref[sec:index_chronic]{\small{[返順時目]}}
~
\hyperref[sec:031628]{\small{> > > NEXT SERMON > > >}}
\newline
\newline
撒母耳記上 17:1-58
\newline
\begin{longtable}{cl}
\hline
\hline
章節 & 經文 (和合本修訂版)\\
\hline
17:1 & \begin{tabularx}{0.7\textwidth}{X} 非利士人召集他們的軍隊來爭戰。他們聚集在猶大的梭哥，在梭哥和亞西加中間的以弗‧大憫安營。 \end{tabularx} \\ \\ \relax
17:2 & \begin{tabularx}{0.7\textwidth}{X} 掃羅和以色列人也聚集，在以拉谷安營，擺陣迎戰，要與非利士人打仗。 \end{tabularx} \\ \\ \relax
17:3 & \begin{tabularx}{0.7\textwidth}{X} 非利士人站在這邊的山上，以色列人站在那邊的山上，當中有谷。 \end{tabularx} \\ \\ \relax
17:4 & \begin{tabularx}{0.7\textwidth}{X} 從非利士營中出來一個挑戰的人，名叫歌利亞，是迦特人，身高六肘一虎口。 \end{tabularx} \\ \\ \relax
17:5 & \begin{tabularx}{0.7\textwidth}{X} 他頭戴銅盔，身穿鎧甲，甲重五千舍客勒銅。 \end{tabularx} \\ \\ \relax
17:6 & \begin{tabularx}{0.7\textwidth}{X} 他腿上有銅護膝，兩肩之中背負銅矛。 \end{tabularx} \\ \\ \relax
17:7 & \begin{tabularx}{0.7\textwidth}{X} 他的槍桿粗如織布機的軸，槍頭的鐵重六百舍客勒。有一個拿盾牌的人走在他前面。 \end{tabularx} \\ \\ \relax
17:8 & \begin{tabularx}{0.7\textwidth}{X} 歌利亞站著，對以色列的軍隊喊叫，對他們說：「你們出來擺陣作戰是為了甚麼呢？我不是非利士人嗎？你們不是掃羅的僕人嗎？你們選一個人出來，叫他下來到我這裡吧。 \end{tabularx} \\ \\ \relax
17:9 & \begin{tabularx}{0.7\textwidth}{X} 他若能與我決鬥，把我殺死，我們就作你們的奴隸；我若勝了他，把他殺死，你們就作我們的奴隸，服事我們。」 \end{tabularx} \\ \\ \relax
17:10 & \begin{tabularx}{0.7\textwidth}{X} 那非利士人又說：「我今日向以色列的軍隊罵陣。你們叫一個人出來，跟我決鬥吧。」 \end{tabularx} \\ \\ \relax
17:11 & \begin{tabularx}{0.7\textwidth}{X} 掃羅和以色列眾人聽見非利士人這些話就驚惶，非常害怕。 \end{tabularx} \\ \\ \relax
17:12 & \begin{tabularx}{0.7\textwidth}{X} 大衛是猶大伯利恆的以法他人耶西的兒子，耶西有八個兒子。在掃羅的時候，這人年老，在眾人中受敬重。 \end{tabularx} \\ \\ \relax
17:13 & \begin{tabularx}{0.7\textwidth}{X} 耶西最大的三個兒子跟隨掃羅出征。出征的三個兒子名字是：長子以利押，次子亞比拿達，三子沙瑪。 \end{tabularx} \\ \\ \relax
17:14 & \begin{tabularx}{0.7\textwidth}{X} 大衛是最小的，最大的三個兒子跟隨掃羅。 \end{tabularx} \\ \\ \relax
17:15 & \begin{tabularx}{0.7\textwidth}{X} 大衛有時離開掃羅，回伯利恆為他父親放羊。 \end{tabularx} \\ \\ \relax
17:16 & \begin{tabularx}{0.7\textwidth}{X} 那非利士人早晚都出來站著，共四十日。 \end{tabularx} \\ \\ \relax
17:17 & \begin{tabularx}{0.7\textwidth}{X} 耶西對他兒子大衛說：「你拿一伊法烘了的穗子和十個餅，跑到營裡去，交給你的哥哥， \end{tabularx} \\ \\ \relax
17:18 & \begin{tabularx}{0.7\textwidth}{X} 再拿這十塊奶餅，送給他們的千夫長，並要問你哥哥好，向他們要個憑據回來。」 \end{tabularx} \\ \\ \relax
17:19 & \begin{tabularx}{0.7\textwidth}{X} 掃羅和大衛的三個哥哥，以及以色列眾人，都在以拉谷與非利士人打仗。 \end{tabularx} \\ \\ \relax
17:20 & \begin{tabularx}{0.7\textwidth}{X} 大衛早晨起來，把羊交託一個看守的人，照耶西所吩咐的帶著食物去了。到了軍營，軍隊剛出到戰場，吶喊叫陣。 \end{tabularx} \\ \\ \relax
17:21 & \begin{tabularx}{0.7\textwidth}{X} 以色列人和非利士人都擺列陣勢，彼此相對。 \end{tabularx} \\ \\ \relax
17:22 & \begin{tabularx}{0.7\textwidth}{X} 大衛把東西留在看守物件的人手中，跑到戰場，問他哥哥好。 \end{tabularx} \\ \\ \relax
17:23 & \begin{tabularx}{0.7\textwidth}{X} 他與他們說話的時候，看哪，那挑戰的人，就是迦特的非利士人歌利亞，從非利士隊伍中上來，說了同樣的話，大衛聽見了。 \end{tabularx} \\ \\ \relax
17:24 & \begin{tabularx}{0.7\textwidth}{X} 以色列眾人看見那人就非常害怕，從他面前逃跑。 \end{tabularx} \\ \\ \relax
17:25 & \begin{tabularx}{0.7\textwidth}{X} 以色列人說：「這上來的人你看見了嗎？他上來是要向以色列人罵陣。若有人能殺他，王必賞賜他大財，將自己的女兒嫁給他，並在以色列人中免除他父家納糧服役。」 \end{tabularx} \\ \\ \relax
17:26 & \begin{tabularx}{0.7\textwidth}{X} 大衛對站在旁邊的人說：「若有人殺這非利士人，除掉以色列人的羞辱，他會怎樣呢？這未受割禮的非利士人是誰，竟敢向永生神的軍隊罵陣！」 \end{tabularx} \\ \\ \relax
17:27 & \begin{tabularx}{0.7\textwidth}{X} 百姓照同樣的話對他說：「若有人殺了那人，必這樣待他。」 \end{tabularx} \\ \\ \relax
17:28 & \begin{tabularx}{0.7\textwidth}{X} 大衛的長兄以利押聽見大衛與他們所說的話，就向他發怒，說：「你下來做甚麼呢？在曠野的那幾隻羊，你交託誰了呢？我知道你的驕傲和你心裡的惡意，你下來只是為了看戰爭！」 \end{tabularx} \\ \\ \relax
17:29 & \begin{tabularx}{0.7\textwidth}{X} 大衛說：「我現在做了甚麼呢？只是問一句話也不可以嗎？」 \end{tabularx} \\ \\ \relax
17:30 & \begin{tabularx}{0.7\textwidth}{X} 大衛離開他轉向別人，問了同樣的事，百姓也照先前的話回答他。 \end{tabularx} \\ \\ \relax
17:31 & \begin{tabularx}{0.7\textwidth}{X} 有人聽見大衛所說的話，就在掃羅面前報告；掃羅就派人叫他來。 \end{tabularx} \\ \\ \relax
17:32 & \begin{tabularx}{0.7\textwidth}{X} 大衛對掃羅說：「人不必因那非利士人灰心。你的僕人要去與他決鬥。」 \end{tabularx} \\ \\ \relax
17:33 & \begin{tabularx}{0.7\textwidth}{X} 掃羅對大衛說：「你不能去與那非利士人決鬥，因為你年紀太輕，他從小就是戰士。」 \end{tabularx} \\ \\ \relax
17:34 & \begin{tabularx}{0.7\textwidth}{X} 大衛對掃羅說：「你僕人為父親放羊，有時獅子來了，有時熊來了，從群中抓走一隻羔羊。 \end{tabularx} \\ \\ \relax
17:35 & \begin{tabularx}{0.7\textwidth}{X} 我就追趕牠，擊打牠，把羔羊從牠口中救出來。牠起來攻擊我，我就揪牠的鬍子，打死牠。 \end{tabularx} \\ \\ \relax
17:36 & \begin{tabularx}{0.7\textwidth}{X} 你僕人曾打死獅子和熊，這未受割禮的非利士人必像獅子和熊一樣，因為他向永生神的軍隊罵陣。」 \end{tabularx} \\ \\ \relax
17:37 & \begin{tabularx}{0.7\textwidth}{X} 大衛又說：「耶和華救我脫離獅子和熊的爪，他必救我脫離這非利士人的手。」掃羅對大衛說：「你去吧！耶和華必與你同在。」 \end{tabularx} \\ \\ \relax
17:38 & \begin{tabularx}{0.7\textwidth}{X} 掃羅把自己的戰衣給大衛穿上，將銅盔戴在他頭上，又給他穿上鎧甲。 \end{tabularx} \\ \\ \relax
17:39 & \begin{tabularx}{0.7\textwidth}{X} 大衛佩刀在戰衣上，試著走走看。因大衛沒有試過，就對掃羅說：「我穿戴這些不能走路，因為我沒有試過。」於是他脫下身上的這些軍裝。 \end{tabularx} \\ \\ \relax
17:40 & \begin{tabularx}{0.7\textwidth}{X} 他手中拿杖，又在溪中挑選了五塊光滑的石子，放在袋裡，就是牧人帶的囊裡，手裡拿著甩石的機弦，迎向那非利士人。 \end{tabularx} \\ \\ \relax
17:41 & \begin{tabularx}{0.7\textwidth}{X} 那非利士人漸漸走近大衛，拿盾牌的人在他前面。 \end{tabularx} \\ \\ \relax
17:42 & \begin{tabularx}{0.7\textwidth}{X} 非利士人觀看，見了大衛，就藐視他，因為他年輕，面色紅潤，容貌俊美。 \end{tabularx} \\ \\ \relax
17:43 & \begin{tabularx}{0.7\textwidth}{X} 非利士人對大衛說：「你拿著杖到我這裡來，我豈是狗嗎？」非利士人就指著自己的神明詛咒大衛。 \end{tabularx} \\ \\ \relax
17:44 & \begin{tabularx}{0.7\textwidth}{X} 非利士人又對大衛說：「來吧！我要把你的肉給空中的飛鳥和田野的走獸。」 \end{tabularx} \\ \\ \relax
17:45 & \begin{tabularx}{0.7\textwidth}{X} 大衛對非利士人說：「你來攻擊我，是靠著刀槍和銅矛，但我來攻擊你，是靠著萬軍之耶和華的名，就是你所辱罵、帶領以色列軍隊的神。 \end{tabularx} \\ \\ \relax
17:46 & \begin{tabularx}{0.7\textwidth}{X} 今日耶和華必將你交在我手裡。我必殺你，砍下你的頭，今日我要把非利士軍兵的屍體給空中的飛鳥和地上的野獸，使全地的人都知道以色列中有神， \end{tabularx} \\ \\ \relax
17:47 & \begin{tabularx}{0.7\textwidth}{X} 又使這裡的全會眾知道，耶和華使人得勝，不是用刀用槍，因為戰爭全在乎耶和華。他必將你們交在我們手裡。」 \end{tabularx} \\ \\ \relax
17:48 & \begin{tabularx}{0.7\textwidth}{X} 那非利士人起來，迎向大衛，走近前來。大衛急忙往戰場，迎向非利士人跑去。 \end{tabularx} \\ \\ \relax
17:49 & \begin{tabularx}{0.7\textwidth}{X} 大衛伸手入囊中，從裡面掏出一塊石子來，用機弦甩去，擊中非利士人的前額，石子進入額內，他就仆倒，面伏於地。 \end{tabularx} \\ \\ \relax
17:50 & \begin{tabularx}{0.7\textwidth}{X} 這樣，大衛用機弦和石子勝了那非利士人，擊中了他，把他殺死；大衛手中沒有刀。 \end{tabularx} \\ \\ \relax
17:51 & \begin{tabularx}{0.7\textwidth}{X} 大衛跑去，站在那非利士人身旁，把他的刀從鞘中拔出來，殺死他，用刀割下他的頭。非利士眾人看見他們的勇士死了，就都逃跑。 \end{tabularx} \\ \\ \relax
17:52 & \begin{tabularx}{0.7\textwidth}{X} 以色列人和猶大人就起來吶喊，追趕非利士人，直到該和以革倫的城門。被殺的非利士人倒在路上，從沙拉音直到迦特和以革倫。 \end{tabularx} \\ \\ \relax
17:53 & \begin{tabularx}{0.7\textwidth}{X} 以色列人追趕非利士人回來，搶奪了他們的軍營。 \end{tabularx} \\ \\ \relax
17:54 & \begin{tabularx}{0.7\textwidth}{X} 大衛拿著那非利士人的頭帶到耶路撒冷，卻把那非利士人的軍裝放在自己的帳棚裡。 \end{tabularx} \\ \\ \relax
17:55 & \begin{tabularx}{0.7\textwidth}{X} 掃羅看見大衛去迎戰非利士人，就問押尼珥元帥說：「押尼珥，那年輕人是誰的兒子？」押尼珥說：「王啊，我在你面前起誓，我不知道。」 \end{tabularx} \\ \\ \relax
17:56 & \begin{tabularx}{0.7\textwidth}{X} 王說：「你可以問問那孩子是誰的兒子。」 \end{tabularx} \\ \\ \relax
17:57 & \begin{tabularx}{0.7\textwidth}{X} 大衛打死那非利士人回來，押尼珥領他到掃羅面前，大衛手中拿著非利士人的頭。 \end{tabularx} \\ \\ \relax
17:58 & \begin{tabularx}{0.7\textwidth}{X} 掃羅問他說：「年輕人，你是誰的兒子？」大衛說：「我是你僕人伯利恆人耶西的兒子。」 \end{tabularx} \\ \\
[1ex]
\hline
\hline
\end{longtable}
$^{1}$聖經問答:誰殺了哥利亞?.
上面那句最後的意思是:When you have not been to church for a while.
所以沒有回來教會很久了.
誰殺了哥利亞?.
答了COVID也有的.
不是David,是COVID.
所以你記得,沒有回來教會一段時間真的生疏了.
讓我們一同繼續彼此勉勵.
說到大衛的生平,我想跳過他,不說也不行.
當然要說,大衛打哥利亞.
有沒有人沒聽過這個故事?.
有沒有誰聽過這個故事?.
也不是全部,所以溫一下.
上星期跟大家說到.
讓神發揮神給你的潛能.
所以大衛真的發揮神給他的潛能.
未必一定要打仗那種潛能.
今天我們要看的題目是.
你都可以做巨人.
你想做誰?.
左邊那個還是右邊那個?.
那你都可以做巨人,究竟誰是巨人呢?.
當然我們不會做哥利亞.
但是我們可以做巨人.
為什麼呢?.
我們先溫習一下這個故事.
我們溫習一下.
用啟認方式讀,好嗎?.
我讀一節,請大家讀一節.
有些字比較深,我也要查一下字典怎麼讀.
所以讀錯不要緊.
不過我們可以學一下.
第一節請聽我讀.
非利士人焦聚他們的軍旅,要來爭戰.
我先讀.
聚集在屬猶大的疏哥.
安營在疏哥和亞西加中間的爾忽大盟.
請讀.
非利士人站在這邊山上.
以色列人站在那邊山上.

$^{41}$當中有谷.
所以我們知道.
讓我控制.
我們知道在兩邊的山上.
大家對壘.
下面的谷,大家明白谷就是下面.
所以在山上派人下來打.
大家可以想像這幅圖畫.
第四節.
綠藻零一富厚是多高?.
一藻即是一尺半.
綠藻即是多少?.
九尺.
九尺零一富厚.
一富厚大約半尺.
即是九尺半.
九尺五寸就這麼準.
九尺六寸即是九尺半.
九尺半有多高?.
NBA那些球星多數有多高?.
六尺多.
最高的那些七尺.
最多七尺一.
有沒有八尺?.
快有了.
即是人工培育那些?.
你有沒有試過走近NBA的球員?.
你發覺他真是很可怕.
如果NBA的球員走近他.
就等於我們走近NBA的球員.
如果我們走近他.
九尺半.
第五節.
頭戴銅盔身穿海甲.
甲重五千尺下立.
五千尺下立即是大約一百五十磅.
一百五十磅不是他體重.
是那些銅盔.
海甲的重量.
即是甲重五千.

$^{81}$即是單是那件盔甲都重一百五十磅.
即是我這麼重.
即是我抱著他.
他就拿著我.
連盔甲都這麼重.
第六節請你讀.
這個讀作「擊」字.
刀槍劍擊.
第七節.
槍罕粗如織布的機軸.
鐵槍頭重六百尺下立.
六百尺下立即是大約十八磅.
即是頭的槍頭十八磅.
一個啞鈴兩個啞鈴.
有一個拿盾牌的人在他面前走.
還要有一個盾牌的人在他面前走.
裝備真的很厲害.
如果今天來說是重型裝備.
第八節請你讀.
他若能與我戰鬥.
將我殺死.
我們就作你們的福人.
我若勝了他.
將我殺死.
他若能與我戰鬥.
將我殺死.
我們就作你們的福人.
他若能與我戰鬥.
將我殺死.
我們就作你們的福人.
他若能與我戰鬥.
將我殺死.
我們就作你們的福人.
我若勝了他.
將他殺死.
你們就作我們的福人.
服侍我們.
所以非利士人和以色列人打仗.
他們先派代表.
贏了就不用再打了.

$^{121}$免得傷亡那麼多.
所以派一個人出來打.
當然要有最厲害的那個.
派最厲害的那個.
非利士人最厲害的那個就是哥利亞.
當然是他.
而且有最好的裝備.
以色列人就沒有人敢出去.
四十天.
如果大家回去再讀一下聖經.
四十天和他等.
等多久都沒有人敢出來.
四十天就是多久?.
一個多月.
天天叫你出來.
都沒有人出來.
所以等了很久.
結果這個大衛.
十幾歲.
剛好去到.
軍營探望哥利亞.
知道這件事.
自告奮勇.
就說我代表以色列人出戰.
這個真是.
出戰打誰?.
打巨人哥利亞.
結果怎樣呢?.
我們跳了一些經文.
一會兒回來再讀.
四十八節一起讀.
非利士人起身.
大衛前來,大衛急忙迎著非利士人.
往戰場跑去.
這個非利士人是指哥利亞.
不是所有非利士人.
四十九節.
大衛用手將囊中.
掏出一塊石子來.
用機緣甩去.

$^{161}$打中非利士人的額.
石子進入額內.
他就撲倒面伏於地.
這樣大衛用機緣甩石.
勝了那非利士人.
打死他.
大衛手中.
卻沒有刀.
連刀都沒有.
相比哥利亞.
真是厲害.
又有刀又有擊又有盾牌.
又有槍.
頭都送十八磅.
穿上盔甲.
前面還有人擋住.
這樣竟然輸給大衛.
然後大衛就用.
哥利亞的刀.
拔掉他的刀.
就殺了他.
將非利士人剷除.
哥利亞雖然是肉身的巨人.
但卻不堪一擊.
相反少年英雄大衛.
他是真正的巨人.
後來受萬人敬仰.
所以我們不需要像哥利亞才是巨人.
可以像誰?.
像大衛.
能夠成為巨人.
為什麼呢?.
第一,我們看到.
因為巨人不是憑肉身衡量的.
是否一陣子大就贏了?.
現在的NBA球星.
不是一定是最好的.
不是的.
有些最厲害的.
好像史蒂夫·克里.

$^{201}$很矮小.
比起他們.
所以.
事實上有很多東西都是.
好像巴西柔術.
有沒有誰學過巴西柔術?.
即是柔道.
所以這是一次真正的比賽.
這個重多少磅呢?.
490磅.
你看他的樣子就知道了.
490磅.
做到我三個以上.
OK.
高六尺八寸.
結果輸在Royce Garcia.
這個只是重.
190磅.
差了300磅的重量.
結果誰贏呢?.
誰贏呢?.
Garcia贏了就告訴你了.
不會給你看.
又好像這個.
這個叫做.
Chicago Bears.
的Perry.
他們參加.
吃熱狗.
後面看到.
吃熱狗大賽.
這麼大隻.
想著一定可以吃很多.
贏定了.
結果誰贏呢?.
輸給這個.
這個Perry.
355磅.
這個日本的.
叫做Takaru.

$^{241}$他是五尺七寸高.
30磅.
腰圍只有30.
有沒有誰腰圍30?.
結果他吃了.
44隻半熱狗.
12分鐘之內.
那個Perry.
吃了幾隻都不夠10隻.
距離太大.
所以大勝.
所以我們看到.
很多東西未必一定.
大隻就贏.
哥利亞是巨人.
他有最好的裝備.
其實他不需要的.
他這麼大隻.
需要這麼多裝備.
揹一百多二百磅的東西.
跟對面的人打.
按住他都贏.
照計.
用不著穿這麼多東西.
為什麼呢?.
證明他的口很狂奧.
口氣很大.
誰敢出來跟我打.
咒罵上帝.
永生神的上帝.
口氣真的很大.
但是內心.
很多人都是這樣.
外面那些人.
好像很厲害.
很多錢,很多權勢.
很多人擁戴.
但裡面.
卻是非常的.
有一種膽怯和自卑.

$^{281}$不知道是不是因為有些特別高大的人.
經常被人取笑.
打仗的人.
好像很威風.
其實看他就知道.
他要裝很多外面的東西.
可見他裡面.
自己的價值不是很高.
相反大衛人家給他盔甲.
他說不用.
我不合用.
所以他那種信心.
不是在乎他的裝備.
不是他的外表.
這就是分別.
所以我們看到.
很多人安全感建立在賺很多錢.
擁有很多物質的東西.
有很多名成利就的東西.
來叫他自己.
感覺安全.
或者覺得.
有很多人擁戴他.
感覺安全.
但他裡面.
其實沒有安全感.
所以靠外面很多東西.
來裝飾.
遮住.
來告訴自己.
我厲害.
這就是那種.
不同的.
分別.
所以一個人沒有上帝.
沒有神.
他的信心根基建立在他自己.
建立在他其他.
外在物質的東西.
其實是很虛的.

$^{321}$但他裡面有上帝.
他建立的信心.
和自己的價值觀.
和自己的.
所擁有的.
不是那些外面的東西.
他裡面的信心.
是鞏固的.
所以我們看到.
在哥林多前書這樣說.
我們一起讀一讀這幾個聖經.
請大家同聲讀出.
上帝卻揀選了.
世上如絕的.
叫有智慧的羞愧.
又揀選了.
世上軟弱的.
叫那強壯的羞愧.
上帝也揀選了.
世上卑賤的.
被人厭惡的.
以及那無有的.
為要棄廢棄那有的.
此一切有血氣的.
在上帝面前.
一個也不能自誇.
但你們得在.
基督耶穌女.
是本乎上帝.
上帝.
又使他成為.
我們的智慧.
公義 聖潔.
救贖.
趙寶樂講出一個很重要的道理.
他說這個世界上的人.
神揀選什麼人.
他說他.
反而揀選那些如絕的.
揀選那些.

$^{361}$好像軟弱的.
叫強壯的羞愧.
強壯只是建立在他自己.
和他所擁有的東西.
但這些如絕的.
軟弱的.
他建立的是在基督耶穌裡面.
所以在基督耶穌裡面.
是本乎上帝.
就是今天我們能夠承受.
在基督裡面.
上帝給我們.
藉著基督成就的智慧.
公義 聖潔和救贖.
這些是世界.
不可以給我們的.
所以當我們能夠擁有這些的時候.
基督徒屬神的子民.
我們能夠在基督裡面.
有這種真正.
不是這個世界給我們.
那種滿足.
而是上帝為我們成就.
在基督裡面救贖的恩典.
智慧 公義 聖潔 救贖.
這些就是.
我們所擁有最寶貴的東西.
所以你們.
我可不可以做巨人.
可不可以.
這個就是勝過了.
其他那些強壯的.
那些.
所謂那些有的東西.
那些一切.
血氣的東西.
因為這些一切都不可以.
真正地在神的面前.
自誇.
所以原來我們都可以.

$^{401}$做巨人 因為我們在基督裡面.
蒙神的簡選.
這個就是.
那個意思.
那究竟怎樣做巨人呢?.
我們在基督裡面蒙簡選.
怎樣發揮.
那個潛能.
怎樣能夠真正.
做一個在神面前.
他看我們為重要的.
巨人.
我們看看是甚麼衡量呢?.
就是在屬靈的品格上.
在品格成熟上做巨人.
所以你和我.
可不可以做巨人呢?.
不是健碩.
不是擁有多少金錢.
不是擁有多少名利.
是在品格成熟上.
做巨人.
所以貼上來.
這個撒母耳記上十七章.
剛才我們讀到第九節.
和四十幾節 中間有些經文.
我們現在補讀一下.
十七節我們開始讀.
我們再用啟認方式讀.
這次請會眾先讀.
請你們讀.
和十個餅.
我讀 請我讀.
到我讀.
再拿這十塊奶餅.
送給他們的千夫長.
且問你哥哥們好.
向他們要一封信來.
(大衛和以拉谷的戰士).
所以蘇羅.

$^{441}$他就立了.
大衛三個哥哥.
其實就是頭三個.
那三個最大的.
成為以色列的戰士.
他們正在.
以拉谷.
剛才我讀過 和比利時人打仗.
所以耶西.
耶西是誰?.
阿達和阿爸.
你現在幫我拿一些.
食物去.
那些餅.
有些麥餅 有些奶餅.
然後就.
拿給哥哥 和順便.
將這些.
給一些.
其他以色列的軍官.
所以就給那些千夫長.
千夫長是高級軍官.
所以呢 然後呢.
叫你哥哥寫一封信.
報平安.
他爸爸很有禮貌.
很照顧兒子.
不過就叫大衛去做信差.
所以就叫大衛去.
就是這個神安排的機會.
大衛就可以.
去到戰場 本來沒有他的份.
他就因為.
聽爸爸的話.
去尊敬哥哥.
又去尊敬那些.
好像爸爸一樣.
去尊敬那些以色列的軍官.
然後家人.
又互相可以通訊.

$^{481}$所以就去了.
二十節.
請聽我讀.
大衛早晨起來.
將羊交託一個看守的人.
其實可以交給別人看守.
上次我們說.
不是一定要他去看守.
不過他爸爸叫他去看守就去看.
這次不是了 爸爸叫他去做別的事.
他就交給別人.
照著他父親所吩咐的話.
留意這句.
大衛帶著食物去.
到了.
支寵營.
軍兵剛出到戰場.
立喊要戰.
請讀.
二十二節大衛把他帶來的食物.
留在看守物件的人手下.
跑到戰場.
問他哥哥門安.
二十八節.
真是無辜.
無端端被他鬧一場.
你有沒有試過無端端被人鬧.
有嗎.
真慘.
一片好心.
爸爸叫我帶些食物給你們.
結果被你鬧一場.
鬧得很激烈.
向他發怒.
你來這裡幹什麼.
你抗議羊交給誰.
為什麼不看守.
還強行給他罪名.
說他驕傲.
還有心裡的惡意.

$^{521}$心裡的惡意.
你們來幹什麼.
想害我們.
想他最壞的事.
真慘.
你們來.
只不過看熱鬧.
用一些很嚴厲.
很惡毒的說話.
這是你們的弟弟.
拿食物給你們.
可見為什麼.
神不揀選這些哥哥.
大衛說我做了什麼.
我來喜沒緣故嗎.
三十字請讀.
大衛說.
我來喜沒緣故嗎.
三十字請讀.
大衛說.
我來喜沒緣故嗎.
三十字請讀.
所以在這整個記載裡.
我們見到很重要.
大衛他能夠.
有一種大將之風.
就是這種大將之風.
那些小人.
那些.
專門.
藐視你.
專門惡意.
不要理他那麼多.
所以.
他就有大將之風.
這個就代表他那種.
屬靈的成熟度.
一種生命的品格的內涵.
這裡就讓我們看見.
作為巨人.

$^{561}$屬靈的品格的成熟度.
來衡量.
是不是一個巨人.
所以.
大衛雖然年少.
但是他比起所有.
其他以色列的軍人.
比起他大個好幾個.
他有七八個兄弟.
頭三個就是距離一段.
這些哥哥都成熟.
因為.
他更加比他們.
認識上帝.
還有他有一種.
屬靈的品格.
大衛就是因為他們在年少的時候.
在抗議的時候已經和神.
有一個建立.
寫詩篇還稱讚神.
靠著神的力量對付那些.
猛獸.
靠著神的力量去面對.
那些人被他藐視.
用神的眼光看他自己.
而且發揮神給他的潛能.
他看重神的事.
他以神的事為念.
他看重神的心.
他以神的心為心.
他看重神的榮耀.
所以他很著緊.
誰來藐視我的上帝.
這就是他屬靈的成熟.
所以他屬靈的成熟不只是.
我和神的關係.
他影響他和人的關係.
這個才重要.
我和神很親近.
但他和人相處有很多問題.

$^{601}$不是真正.
和神的關係.
影響到他的人格.
只不過是一種理念.
影響人格就影響他的品行.
影響他和人的關係.
所以他對爸爸很尊敬.
很聽話.
這些就是反映.
他不是說我不用理你.
不是的.
因為這樣.
他更加對人很尊敬.
很有禮貌.
其實哥哥也很有愛.
雖然哥哥對他很差.
這就是屬靈的成熟.
就是他的品格.
所以在聖經裡形容他.
處事成熟.
說話俠義.
這麼小就說話很俠義.
這就是他.
他知道他的集中力.
不是和哥哥爭辯.
浪費時間.
你怎麼看.
他的集中力就是.
看誰在面對國家的危機.
面對勝敗的國運.
你看他的成熟.
他以大事為念.
這就是他的成熟.
所以他的屬靈品格.
不是建基在別人對他好還是不好.
當然我們都需要反省.
到底是什麼問題.
但如果不是自己的問題.
就很有信心.
知道我不是因為.

$^{641}$活在他的影子下.
我是活在上帝的面前.
討他的喜悅.
所以他的自我形象是很正確的.
他的品格成熟就在這裡.
聖經裡面也是這樣提醒我們.
在提摩太前書第四章.
我們一起讀.
不可叫人小看你年輕.
你信了主耶穌多久.
你年紀多大.
原來不是看我們年紀多大.
不是看我們信了主耶穌.
做了基督徒多久.
保羅和提摩太說.
不要叫人小看你年輕.
著重點就是.
你能夠做別人的榜樣.
你信了耶穌很久.
但對於那些初信的基督徒.
你沒有榜樣給他們.
即是不夠他們成熟.
那麼.
那個榜樣在哪裡看呢.
究竟聖經裡面著重多少呢.
著重什麼呢.
就是著重這些.
這些就是指標.
看我們是年少還是巨人.
所以我們也可以做巨人.
真的.
如果我們的言語.
我們的行為.
我們的愛心.
我們的信心.
我們的清潔.
在這些上面.
可以成熟成長的話.
這個就是巨人.
在屬靈的成熟度上面做巨人.

$^{681}$不是肉身的巨人.
在品格.
屬靈的成熟度上做巨人.
在屬靈的成熟度上做巨人.
頂尖的我們也可以追求.
做巨人.
哪幾樣話.
言語.
行為.
愛心.
信心.
清潔.
第三.
我們看到在恆心.
操練上做巨人.
一步十寒.
是一步登天的,所以逐一來,大衛也是這樣的,大衛不是一下子就打高利亞,他背後有很多東西,我們看看,他自己說,十七章三十四節,我們再來讀,這次我先讀,大衛對蘇羅說,你僕人為父親放羊,有時來鳥獅子,有時來鳥雄,從群中含一隻羊羔去.三十五節,請讀..
大衛對蘇羅說,你僕人為父親放羊,有時來鳥獅子,有時來鳥雄,從群中含一隻羊羔去.三十五節,請讀..
你僕人曾打死獅子和雄,這未受割禮的非利士人,向永生上帝的軍隊罵陣,也必將獅子和雄一半..
這代表他背後的操練,他背後經常靠著神救羊,這是他的愛心,他的責任感,既然託付給我,我就做好它,豈有此理,我父親的羊你來欺負..
所以他現在一樣,豈有此理,我上帝的軍隊你來欺負.這就是他平時的操練.他平時操練習慣了對於屬父親的很緊張,對羊很有愛心,對於獅子和雄他靠著神來擊退..
所以他打割禮,召羊,用回這個樣子.所以他自己說的,這個未受割禮的非利士人,竟然也必將獅子和雄一半.他的眼光就在這裡..
三十七節,到誰讀?到你們讀,請讀..
他手中拿著一個杖,又在溪中挑選了五塊光滑的石子,放在袋裡,就是牧人帶的狼狸,手中拿著脫石的機緣,.
就去迎拿非利士人,四十九節,請讀..
所以大家看到,大衛是有訓練的.他訓練,他平時知道第一粒不行就拿第二粒,就是他平時練的..
然後他面對那些獅子和雄,他也是用心靠著神去面對.他擁有一種裡面正確的動機,要救羊,然後保護爸爸的物業..
所以依靠上帝去做成負責任,很有責任感去完成,他不是跑掉的..
所以他看到這種是他的恆心的操練,過往這麼多年在抗疫中不斷練習這種心性,練習這些技巧,練習這些應對危機的時候的一些策略..
這個就是恆心操練..
很多時候我們做一點點又倒退很多,明知自己有些東西需要改,改一點點又不行..
不要灰心,我們恆心繼續操練,繼續操練,繼續操練..
你知道自己有些性格需要改,你就嘗試做第一步..
這個很好的,這個叫做什麼?.
你走上山,上到山頂,success或者上到山頂,it's a journey of baby steps,.
即是不會一步就登山,就要走很多步,逐一步,逐一步,逐一步..
所以親愛的兄弟姐妹,求神幫助我們,我們新年開始,讓我們有些東西需要改..
我們今年裡面,下個禮拜就馬上改了,多數都不會的..
但是你可以想一些辦法,改一些,再改一些..
我每次來崇拜都遲到的,遲到十分鐘,下個禮拜準時,可能未必..

$^{721}$遲五分鐘,有進步了,是不是?.
接著又再早一點,又再早一點,結果怎麼樣?.
早到..
Baby steps,你知道有些東西要改的,.
聖靈感動了你很久了,人們都看見你了..
所以你開始去改,但是你知道,不要緊,我們每個人都很包容的..
一下子做不到,不要緊,baby steps..
不只是一步登山,一步登天..
這是什麼人?.
登天的,這個人叫做Buzz Aldrin,.
Buzz Aldrin是歷史裡面第二個登上月球的人,一步登天..
是不是一步?不是,他告訴我們他的名言,Bravery comes along as a gradual accumulation of discipline..
你知不知道他練了多久,才能做到太空人?.
勇氣是日積月累,恆心操練而來的..
平時的操練真的重要,.
怎麼會一發即中呢?當然一定有上帝的幫助,.
但是是他平時操練而來,有些成果..
所以,小的事忠心,負責任,大事託付給你,那就放心..
所以我們見到他的信心,他的勇氣,他的技能,.
都是日積月累,不斷的..
所以有些時候,有些事看不到的,你有沒有功德心?.
有些事看不到的,人看不到..
你能夠有功德心,做了一些事,譬如去完洗手間,清理了它,沒人看到..
這是鍛鍊我們自己,有功德心..
你自己病,所以要考慮一下會不會傳染人,這是功德心..
所以,這些是日積月累,你內心的心性操練,.
以致慢慢你的人格健全..
這些都是小的事操練,小的事操練..
可不可以做巨人呢?就看我們願不願意恆心操練..
第四,我們看見在勇氣上,勇氣可嘉上做巨人..
第十一節,我們一起讀..
「蘇羅和以色列眾人聽見非利士的這聲話,就驚慌,極其害怕.」.
二十四節..
「以色列眾人看見那人就逃跑,極其害怕.」.
三十二節..
「大衛對蘇羅說:人都不必因那非利士人膽怯,你的僕人要去與那非利士人戰鬥.」.
先是蘇羅和以色列眾人,聖經形容他:就驚慌,極其害怕..
蘇羅也驚慌,其他所有人都驚慌..
以色列眾人看見那人就逃跑,極其害怕..
大衛對蘇羅說:人都不必因那非利士人膽怯,你的僕人要去與那非利士人戰鬥..

$^{761}$你看見這就是大衛的勇氣,勇氣可嘉..
的確,勇氣和懼怕可能是相反的,而勇氣就能勝過懼怕..
他來到向以色列人罵戰,以色列人都驚慌,士氣低落,皇帝都極其懼怕..
四十日僵持,沒有一個人敢出去,可見他們真的非常害怕..
今日有沒有害怕的事?弟兄姊妹,有沒有我們很害怕的事?.
你說我很害怕老鼠,蟑螂?.
最近我們教會開始捉老鼠,大家記得有時要關門..
我們辦公室有一次有老鼠咬爛我們的東西,在我們辦公室裡面..
所以大家就放東西在那裡捉,竟然真的捉到一隻,所有同工都很害怕,不敢進去看..
我拍拍胸口進去看一看,其實很小隻..
你們害怕老鼠嗎?.
說真的,我都害怕..
不過那些都是一時之間,但有些事經常都會有懼怕..
因為我們人生裡面是一個變幻的時代,不懂就已經是懼怕,不知又是懼怕..
或者有時很多事令我們憂慮,身體,工作,家庭都有不如意的事,.
甚至有些事覺得力不能勝,有沒有感覺力不能勝?.
這些對我們都好像哥利亞一樣,其實我們都會面對哥利亞,因為那些好像很大..
其實可能不是那麼大,不過我們覺得他好像很大,向我們迎戰..
我們那麼渺小微不足道,有時都會害怕,真是人之常情,甚至逃跑..
感覺到沒有勇氣去面對,這個就是我們人的軟弱..
沒有那麼重要,上帝說神揀選軟弱的人..
所以我們今天可以有福的就是在神的裡面,我們能夠依靠祂..
大衛其實他的勇氣都是慢慢操練出來的,面對獅子和熊,開玩笑的..
但是他慢慢操練,慢慢操練,他的勇氣慢慢加增..
他的勇氣的來源,相信他一次再經歷,兩次再經歷,就是上帝是我的乙求,是我的幫助,是我力量的來源..
所以他很多詩篇都是這樣說,他是我的盤石,是我的山寨,是我的避難所..
親愛的兄弟姐妹,今天我們都是面對巨人,不知道你的巨人是什麼?.
是你的經濟,是你的工作,是你的健康,是你的家人,是你的憂慮,是其他一些的重擔..
我們看他真的很大很大,我們都面對著我們心裡面的懼怕..
所以在一位退休的基督徒的心理學家,他的名字叫Fred Grief,.
他有幾樣東西,其實很簡單,不過很有意思..
他說當我們面對懼怕的時候,因為他是心理學家,.
面對不安,面對痛苦,他說有幾樣兵器很重要,不是心理學的方法..
他說出幾樣東西都是基督徒來的,他說很重要是什麼?.
第一,seek Jesus,即是尋求耶穌..
怎麼會不怕呢?他做心理學家,每個人來他那裡都是怕的,我們自己都怕..
有時候我發覺那些心理學家自己都是因為他很怕,所以他做心理學家..
對不起,不是每個人都是,不過有些是..
他那方面其實是最脆弱的,他特別需要加強..
所以他說第一,就是面對這些懼怕的時候,seek Jesus..

$^{801}$如果你還不是基督徒,你面對人生的懼怕,.
他的勇氣真的不是在我們自己那裡,不是拍拍胸口..
拍拍胸口很多以色列人都應該是,怎麼會等了四十天,全軍都沒有一個..
大衛的勇氣是因為他尋求神..
第二,他說迫切祈禱,pray earnestly..
所以好像主耶穌的門徒都經歷那些大風大浪,.
在船上他說,夫子,我門喪命了,救我..
所以這就是我們的祈禱..
所以讓我們都向神祈禱,夫子,救我..
這就是他的忠告..
我們不可以靠我們自己的意志力,或者單憑我們克服壓力的方法..
當然那些都不是不好的,比如放鬆自己,做運動,這些都是很好的..
但是不要忘記,pray earnestly..
第三,他說讚美上帝..
好像保羅和西拉他們面對苦難,面對極重的困憂,生死未卜的時候,.
他們來讚美上帝,praise the Lord..
所以不是用一些東西麻醉自己,很多人吸毒,喝酒,麻醉自己..
這些都是錯誤的方法..
認同上帝的安排,praise the Lord,就是說無論如何我讚美你,.
無論如何我知道我的生命在乎你,我將我前面信任你..
所以保羅和西拉他們就經歷神奇妙的拯救,.
因為他先在未得這些之前,就來稱頌上帝..
當然這位基督徒的心理學家,他說勤讀聖經,read his word..
他說很多時候他自己經歷到很多很多的安慰和力量,.
都是從神的話語那裡來的..
有時候上帝跟我們說一句話,那個力量很大..
一捉住那句話,那個力量很大,那個勇氣就從那裡來..
所以讓我們不是水果鴨背下,而是我們要讀到上帝的話,.
捉住那句話,read his word,靠著他的應許..
接著他也有很好的建議,他說服事別人,serve others..
他說因為很多時候,其實那些事我們自己看得很大,.
但是當你去發覺,你去幫助其他人的時候,.
哎呀,其他人更慘,那你就覺得自己沒那麼慘..
其他人的需要你能夠幫助他們,.
你覺得自己有一種力量還可以給神去用去幫助人..
那個懼怕忽然間就少了很多,.
因為你發覺原來我還可以走,還可以做,還可以幫,還可以給神用,.
所以就覺得還可以,勇氣就從那裡來..
最後他說「委身於神下」,就是commit fully to God..
所以seek Jesus, pray earnestly, praise the Lord, read his word, serve others, commit fully to God..

$^{841}$「耶和華是我的亮光,是我的拯救,我還怕誰呢?」.
這就是大衛的《雌篇二十七篇》..
「耶和華是我性命的保障,我還懼誰呢?」.
《雌篇二十七篇》第一節..
所以這就是他的心態..
他知道耶和華可以將一切託付給他,交託給他..
勇氣可加,在哪裡加?在上帝那裡加..
勇氣可加,我們都要去學習..
可不可以做巨人?求神給我們這些勇氣..
勇氣就代表這種力量..
所以最後一件事我們看到,最重要就是在對神的信靠上做巨人..
我們再來讀這六節經文,一起讀吧..
大衛問站在旁邊的人說:.
「有人殺這非利士人,除掉以色列人的羞恥,怎樣代他呢?」.
「這未受割禮的非利士人是誰呢?」.
「竟敢向永生上帝的軍隊罵陣嗎?」.
「你僕人曾打死獅子和雄.」.
「這未受割禮的非利士人向永生上帝的軍隊罵陣.」.
「也必將獅子和雄一半.」.
「大衛又說:耶和華救我脫離獅子和雄的罩.」.
「也必救我脫離這非利士人的手.」.
「數落對大衛說:是一句話,耶和華願你同在.」.
我讀這段經文最重要的就是這些highlight..
大衛不斷提什麼呢?.
就是提永生神,永生上帝..
耶和華可以救我,祂都可以救我..
讀到這麼多的時候,連數落最後都說:.
「耶華願你同去.」.
這個就是impression..
impression就是大衛不斷提上帝..
最後那三節讀完它,請讀45節..
「大衛對非利士人說:你來攻擊我,是靠刀槍和銅鏡.」.
「我來攻擊你是靠著萬軍之耶和華的命.」.
「就是你所怒罵帶領以色列軍隊的上帝.」.
「今天,耶和華必將你交在我手裡.」.
「我必殺你,斬你的頭.」.
「又將非利士軍兵的屍首,及空中的飛鳥,地上的野獸吃.」.
「使普天下的人都知道,以色列中有上帝.」.
「又使者眾人知道,耶和華使人得勝.」.
「不是用刀用槍,因為征戰的勝敗存在乎.」.

$^{881}$「耶和華,他必將你們交在我們手裡.」.
弟子們看到嗎?.
我引用這些聖經給大家看..
從聖經而出,你就看見..
分析這些聖經,你就看到大衛的心..
他的心充滿著上帝..
連他出去跟哥利亞打仗的時候,.
他跟哥利亞說,不是說他有什麼了不起,.
不是說他自己有什麼了不起,.
他說:「上帝在我的身上,與我同在.」.
你竟然,你知不知道,你跟誰在打?.
不是跟我在打,你跟上帝在打..
事實上,這個征戰也是..
今天我們基督徒去面對這一切,.
背後是掌管天空屬靈氣的惡魔征戰..
我們征戰不是在乎血氣的征戰,.
而是在乎掌管空中屬靈氣的惡魔征戰..
其實當時以色列跟非理性人打仗,.
背後是上帝跟假神打仗..
所以大衛不是憑血氣征戰,.
他是憑著上帝征戰..
很明顯,這裡很清楚..
叫普天下人都知道,.
以色列中有上帝..
所以當我們做完一切的時候,.
其他人都看見,不是我,是因為誰?.
是因為上帝,Amen?.
所以他說,他自己知道,他做見證,.
我不是用刀,不是用槍殺你的,.
是我要殺你的,斬你的頭..
不過我不是用刀,不是用槍,.
因為征戰的勝敗,全在乎耶和華..
祂必將你們交在我們手裡,.
是祂做主動,祂做主權,祂偉大..
所以什麼是巨人?.
原來這裡給我們看見,Great Faith in a Great God..
今天信得過誰大一點?.
上帝大一點..
所以不只是我們做巨人,.
是我們能夠在這位巨大的上帝那裡有份..

$^{921}$在這位巨大的上帝那裡,.
我們能夠將自己交託給他,.
而且能夠投求他..
因為縱使你九尺五寸那麼高,.
九尺六寸那麼高,.
有一個比你更高的..
所以信心就是從那裡來的..
我們看到這個就是今天給我們看見的大衛,.
他不是盲目的樂觀,.
他是有好的預備,平日的操練..
最重要的就是他有對神的信心..
他發出的勇氣充滿著自信,.
是因為他信得過有一個更大的上帝..
聖經說:在我們裡面的比那在世界的更大..
所以今天我們知道我們在主的裡邊.
能夠做真正的巨人..
巴不得我們每一個都是,.
都能夠做真正的巨人..
在屬靈的品格上做巨人,.
在平日的操練上做巨人,.
在我們的勇氣,可嘉上做巨人..
更加的就是對神的信心和認識上..
最後這段經文結束,.
我們一起看後書四章七和十六節,請讀..
「我們有這寶貝放在瓦器裡,.
要顯明這莫大的能力是出於上帝,.
不是出於我們..
所以我們不喪膽,.
外體雖然毀壞,.
內心卻一天生死一天.」.
保羅很明白道理,我們只不過是瓦器,.
瓦器根本不可以和哥利亞打,.
不過他說要顯明這個莫大的能力是出於上帝,.
不是出於我們..
Amen?.
然而我們都是這樣,叫我們不喪膽,.
因為我們有主在我們裡邊,.
叫我們這個瓦器,本來是沒有價值的,.
本來是一碰就爛的,.
本來是很懼怕的,.

$^{961}$但我們可以有神給我們的勇氣,力量和信心,.
因為我們的神是一位巨大,偉大的上帝..
我們低頭地祈禱,請大人同心禱告..
我們慈愛的天父,我們再次投靠仰望你..
愛我們的主耶穌,.
你已經是勝過死亡和陰間一切,.
全世復活的主..
所以我們今天再次宣告你的德性,.
宣告你是我們的王,是我們的主..
主要是因為我們時常謹靠你,.
時常在聖靈的幫助下,.
去抓緊你的話,得著你的應許,.
加添我們的信心,.
叫我們能夠充滿勇氣,.
繼續面對我們所面對的一切..
主叫我們能夠在你的眼光當中,.
不在世人的眼光當中做巨人,.
而在你的眼光當中,.
我們真正成為一個像大衛一樣,.
勝過這個肉身巨人的真正的巨人..
多謝您聽我們的祈禱,.
補償我們每一位主內親愛的弟兄姊妹,.
奉耶穌基督的名..
\newpage



\section{以賽亞書 14:12-17}
\label{sec:031628}
\textbf{審判早晨之子}~~~~李應新牧師
\newline
\newline
連結: \href{http://www.church.com.hk/acms/content.asp?site=cdc\&op=show\&type=product\&code=031628\&layout=sermon}{\texttt{031628}} ~~~~ 語音日期: 2026-02-08
\newline
\newline
\hyperref[sec:031627]{\small{< < < PREV SERMON < < <}}
~
\hyperref[sec:index_chronic]{\small{[返順時目]}}
~
\hyperref[sec:031629]{\small{> > > NEXT SERMON > > >}}
\newline
\newline
以賽亞書 14:12-17
\newline
\begin{longtable}{cl}
\hline
\hline
章節 & 經文 (和合本修訂版)\\
\hline
14:12 & \begin{tabularx}{0.7\textwidth}{X} 「明亮之星，早晨之子啊，你竟然從天墜落！你這攻敗列國的，竟然被砍倒在地上！ \end{tabularx} \\ \\ \relax
14:13 & \begin{tabularx}{0.7\textwidth}{X} 你心裡曾說：『我要升到天上，我要高舉我的寶座在神的眾星之上，我要坐在會眾聚集的山上，在極北的地方。 \end{tabularx} \\ \\ \relax
14:14 & \begin{tabularx}{0.7\textwidth}{X} 我要升到高雲之上，我要與至高者同等。』 \end{tabularx} \\ \\ \relax
14:15 & \begin{tabularx}{0.7\textwidth}{X} 然而，你必墜落陰間，到地府極深之處。 \end{tabularx} \\ \\ \relax
14:16 & \begin{tabularx}{0.7\textwidth}{X} 凡看見你的都要定睛望你，留意看你，說：『就是這個人嗎？他使大地顫抖，使列國震動， \end{tabularx} \\ \\ \relax
14:17 & \begin{tabularx}{0.7\textwidth}{X} 使世界如同荒野，使城鎮傾覆；是他，不釋放被擄的人歸家。』 \end{tabularx} \\ \\
[1ex]
\hline
\hline
\end{longtable}
$^{1}$「沒勢災難不要害怕」.
因為主幫我們隱藏.
我們要看天上在雲裡與主相遇.
今天我就和你講一段經文.
我們回到剛才那段經文.
第十四章.
這裡講到明亮之星 早晨之子.
這裡以蔡霞先知在講什麼呢.
原來這些是當時的偶像.
那時候迦南地的星星.
全都是覺得他們是天上偶像.
中國人有時候也覺得是什麼星星.
什麼偶像.
那不是很可怕嗎.
為什麼以蔡霞先知會講偶像的東西出來呢.
今天牧師要解釋給你聽.
為什麼會這樣做.
有三重含義.
但先講第一重.
以蔡霞先知引用迦南地的傳說.
是有原因的.
或者容易明白.
不知道你有沒有聽過.
有一個人叫做拿渣.
拿渣不為海龍王.
我們自洗口看.
我知道你是誰.
那時候劉德華還是很年輕.
還有梁朝偉.
他們還是很年輕.
這是拿渣故事.
但昨天我再唱的時候.
原來有些年輕人跟我說.
在國內成長的小朋友看了三次.
原來這套拿渣是國內很流行的.
是卡通片.
看了三次.
拿渣是什麼.
我命由我不由天.
我講拿渣的時候不代表我相信拿渣.

$^{41}$整個東西是什麼.
不畏強權.
挑戰權威.
衝破權勢.
或者有權有勢的人.
挑戰海龍王.
那種.
香港也有一個議員.
號稱叫做拿渣.
在將軍澳區.
聽說據聞.
他小時候做同聲做過拿渣.
但他整個生平都很拿渣.
去選立法局.
六次都選不到.
證明他沒有欽點.
六次都中選.
是不是很拿渣.
整件事都很拿渣.
不過這次被他選到.
這個是後話.
我描述他.
我講了這麼多拿渣.
其實我都不是很信拿渣.
我們拜拿渣.
不會當他是偶像.
這個symbol.
整個meaning其實很深入人心.
我一講大家就明白.
以賽亞先知同樣.
明亮之星早晨之子是偶像.
但那時候有些人聽不明白.
不是說他信他.
而是整個故事的package.
第一重的含義.
是巴比倫王.
因為有類似的地方.
他現在描述的不是偶像.
而是巴比倫王.
不過像他.

$^{81}$原來他從天上墜落.
再撞倒在地上.
因為他心裡沉睡.
他很驕傲.
我要升到天上.
我要高舉我的寶座.
在神眾星以上.
剛才講過那些星星粒粒都是神.
粒粒都代表神.
他要高過其他神.
我要在聚會的山上.
在北方的極處.
很多神會聚會.
在山上.
在北方的極處.
很多星的位置.
原來那些神就在那裡.
他要高舉他的寶座.
比其他神還要厲害.
這是當年的.
但是他那麼驕傲.
我要升到高雲之上.
我要與至上者同等.
他特意用這個翻譯.
我們翻譯的聖經.
特意不用至高者.
為什麼呢?.
免得你搞錯.
整個包裝都是講迦南偶像.
至上者.
即是較大粒的神.
我們平時翻譯聖經.
全地至高.
是我主的耶和華.
這不是在講耶和華.
那時候有些至上的神明.
太陽神.
日神.
大粒很多.
太陽當然大粒過那些星星.

$^{121}$將他下的十五節.
將他墜落在地上.
將他打倒.
必墜落在陰間.
在坑中極深的之處.
當時的故事是這樣說.
明亮之星.
早晨之子.
掉落坑中.
大概是這個故事的包裝.
第一層的應用是什麼呢?.
第一層的應用是在講.
巴比倫王.
因為巴比倫王那時候很大權勢.
打到仗都贏.
而且他非常凶殘.
他怎樣描述他呢?.
「凡看見你都定睛看你.
留意看你說.
使大地顫抖.
使列國震動」.
即是說他經過的地方.
人們都害怕.
聞到風聲都害怕.
因為他來到時很殘暴.
很厲害.
打到仗都贏.
世界如同荒野.
軍隊經過寸草不生.
如同荒野.
而且使城邑傾覆.
即是說會拆了你的城牆.
搞爛你的東西.
搶走你的東西.
使城邑傾覆.
而且不釋放被擄的人歸.
即是說.
當時他們的制度是怎樣呢?.
會將皇帝,貴族.
或者祭司捉回他們的家.

$^{161}$你都聽過.
被擄到巴比倫.
就是巴比倫王.
他們就做這些事情.
雖然他只是這樣.
他都好像早晨之子.
明亮之星早晨之子.
將會墜落.
後面我沒有讀下去.
十八字打後一大段是說他將會怎樣死.
你不要以為他那麼厲害.
將會被人打敗.
將會死無全屍.
將會墜落陰間.
將會被上帝審判.
這是第一重的含義.
提醒我們有權有勢的人.
你看到他好像很駕細.
很得戚.
不用看他什麼時候.
就算那些很有權勢.
權力的統治者.
我們換什麼.
搞到百五歲.
我猜你一百歲打後都要坐輪椅.
就算他全換了.
都不是很威風.
我寧願快點回天府.
你搞到我百五歲幹什麼.
又要坐輪椅.
又要別人扶.
又要別人餵我吃東西.
我不願意.
你就算神的審判.
有時我們會見到.
有時我們未必見到.
但上帝一定會審判惡人.
巴比倫王不用威風很久.
你將會受到審判.
因為你很驕傲.

$^{201}$好像明亮之星早晨之子.
將會受到神的審判.
這是第一重的意義.
這是在以賽亞的時候.
最原本想表達的事情.
第二層的意義.
我們繼續看下去.
下面還有一些經文.
盧家福音.
十章十八字.
耶穌對他們說.
我曾看見撒旦從天上墜落.
像閃電一樣.
對稱.
先解釋一件事.
這是傳統教會第二層.
對這段經文的解釋.
為什麼會有魔鬼撒旦的存在呢?.
神這麼公義.
這麼憐憫.
怎麼會做出這麼壞的撒旦呢?.
神沒有做撒旦.
我們將這兩段經文去理解.
教會傳統去解釋.
魔鬼撒旦是怎麼來的?.
魔鬼撒旦是光明天使墜落.
因為他要高舉自己.
好像與神同等.
他很驕傲.
於是抗拒上帝.
上帝沒有做魔鬼撒旦.
只做了光明天使.
但他墜落就變成魔鬼撒旦.
這是第二層的解釋.
在這裡我還想和大家分享.
我一定說過.
不過說多次都可以.
其實我們見到.
魔鬼撒旦.
我們平時說要趕鬼.

$^{241}$魔鬼撒旦只有一種.
不是我們死了之後.
死了之後變魔鬼.
我們死了.
就不是變魔鬼.
我們死了就回天父身邊.
或者去了陰間.
但不會到處走.
做了魔鬼嚇嚇你.
全部魔鬼都是撒旦扮的.
用不同的形式來迷惑你.
使你不敬拜真神上帝.
魔鬼撒旦扮鬼.
這麼複雜.
魔鬼撒旦要令你離開真神.
要令你感覺到魔鬼撒旦.
是魔鬼撒旦.
才能夠迷惑你.
由頭說到尾.
要迷惑你.
迷惑到你才能迷惑你.
要令你迷惑.
簡單來說.
你以為魔鬼撒旦是這樣的樣子.
他就會給你看這個樣子.
否則不認得他是魔鬼.
於是迷惑不到你.
一看就知道是魔鬼.
什麼意思呢.
鬼佬見到的鬼.
鬼佬的鬼樣.
中國人見到的是中國人的鬼樣.
即是說魔鬼在中國人的地方.
是不會穿西裝.
不會有殭屍的樣子.
因為你見到有西裝.
又反過來.
如果那些鬼去到外國的地方.
他不會穿青裝.
他就說不到你.

$^{281}$青裝是什麼.
他不知道這是鬼.
所以他就不會用這個樣子.
其實是有性別歧視的.
因為從文化角度.
鬼佬的鬼通常是男鬼.
中國的鬼一般是女鬼.
是真的.
因為你們的鬼應該是這樣.
所以他們是用這個樣子來表達.
你見到是這個鬼.
所以我要怕他.
所以他們因為你文化.
能夠用到你迷惑到你的方式.
來迷惑.
當然這是最低級的迷惑方法.
有些更高級的.
不過今天沒有時間再講下去.
不能夠再講下去.
但你大概知道魔鬼殺人.
原來是這樣來的.
這是第二層的含義.
還有第三層的含義.
聖經提摩太前書三章第六.
寫到初入教的不可作監督.
恐怕他自高自大.
就落在魔鬼所受的刑罰裡.
這是什麼意思呢.
這裡就說.
剛剛信主的人.
不要給他這麼高的位置.
是不是自高自大.
等於魔鬼撒旦墜落.
就是他們自高自大要高舉自己.
同等.
所以受到同等的刑罰.
這是第三層的解釋.
但我想再延伸一下.
其實除了初入教.
會有自高自大的跌倒.

$^{321}$或者迷惑.
其實我們信了主很久.
同樣會有.
所以兩樣我們都提防小心.
免得我們落在魔鬼的刑罰裡.
什麼意思呢.
我們信了主很久.
信了主最久是誰呢.
彼得.
第一批信徒.
當然是最久的.
彼得耶穌說你不認我.
眾人都會跌倒.
但我總不會.
那天晚上就三次不認住.
因為驕傲.
自高自大.
覺得自己很厲害.
人人都跌倒.
所以提防小心.
我們呢.
不只是初信主會.
我們信了主很久都會.
我知道先父對我很憐憫.
所以給我一件經歷.
我深刻印象.
我都可能說過.
但都要再說一次.
對我提醒我.
很多年前我在馬鞍山做牧師的時候.
我們教會是教幼稚園的.
我們幼稚園有四百多個學生.
所以我們很多學生.
每年都會見我幾次.
例如畢業禮 開學禮.
其實我主要的角色和功用.
是巴斯江年.
就是用來拍照.
所以.
每個人都有和我拍過照.

$^{361}$所以我經常都出鏡.
不是 Mickey Mouse.
是巴斯江年 有型一點.
我不喜歡做老鼠.
anyway.
所以孩子認得我.
但我認不到孩子.
我經常覺得除了我自己的孩子.
其他孩子的樣子都差不多.
不知道你們有沒有這個印象.
其他孩子的樣子都差不多.
我真的沒有辦法認到四百多個孩子.
有一次.
我經常在馬鞍山服侍的時候.
我經常吃早餐.
就會去大家樂吃早餐.
那間在新港城的大家樂.
是很出名的.
難吃和排隊.
很難吃又要排隊.
氣死人.
大家樂.
有一天我去漁場.
吃早餐的時候排隊.
我坐下來.
因為剛才說要排隊很多人.
所以沒辦法有自己的位置.
一定要搭桌.
所以我習慣了搭桌.
那天我坐下來.
早餐放下的時候.
就坐在我對面.
我搭桌的是一個小孩子.
和他可能是阿婆或阿媽.
我坐下來的時候.
那個小孩子就看著我.
那個樣子.
我現在歷歷在目.
想跟他說.
天父啊.

$^{401}$那天.
我立刻定神.
小朋友.
我不是天父啊.
我是.
牧師黎家宅.
這句說話.
我說出口.
藏在我心裡.
十年前.
我還記得.
我是一個牧師黎家宅.
做了幾十年人.
總有些人生的閱歷.
有些不同的位份.
可能是主任.
校長.
經理.
董事.
導師.
蟑螂.
我們都是很軟弱的.
千萬不要驕傲.
每逢驕傲就會跌倒.
小心提防.
我們是什麼來的.
蟑螂.
我不是天父也不是主耶穌.
提醒我們自己小心提防.
自高自大免得我們落在魔鬼的刑罰裡.
我還想跟他說.
神是很愛我們的.
所以經常被我們提醒.
這個星期我常常默想到.
上帝用很多方式來提醒我.
上個星期一.
是我們牧師全套放假的時候.
星期一牧師全套放假有什麼要做的呢.
我約了人吃飯.
我就去吃飯.

$^{441}$吃完飯還有些時間.
我就跟別人去.
不是跟別人去我自己去.
我自己去行山.
最近打卡.
我也沒有留意到.
因為有人跌倒.
破邊洲.
我以前也不太擔心牧師.
破邊洲其實很安全.
漁農處弄得很漂亮.
那裡有圍欄.
圍欄寫著千萬不要橫過.
但人們很喜歡橫過圍欄.
爬出去也是要出去打卡.
這樣就跌倒了.
不關人事.
寫得很清楚.
千萬不要離開.
以前我也經常行山.
但我從來沒有去過破邊洲.
這條路是新建的.
很棒又跌倒.
我先去打卡.
所以我那天放假.
吃完飯大約三點.
我就出發去西貢破邊洲.
東北其實挺好行的.
不是很難.
我去那裡行山.
經過一條新路.
第一次去.
我一開始以為星期一沒有人行.
西貢星期一應該很清靜.
很清幽.
原來不是.
很多人的.
很多人在我前面行走.
這個不是問題.
但我見到前面有一個人.

$^{481}$不知為何他們行山.
特別是女孩子.
會穿高跟鞋.
我真的不明白.
行山怎會穿高跟鞋.
你穿高跟鞋對我有沒有影響.
他們會咯咯咯咯咯.
阻礙我行山.
因為我三點多才出發.
天黑了.
阻礙我行山.
所以我想.
既然是這樣.
我要開渦輪爬你頭.
不爬頭不游自開.
這些山路的東西.
是不能爬頭的.
因為我一開渦輪爬頭.
就馬上PK在他面前.
在他咯咯咯咯的面前.
他還要好好的問我.
先生你沒事嗎.
沒事.
終於爬了.
沒事一點點.
沒事的.
所以不要.
自高自大.
如果是這樣會怎樣.
會跌倒.
所以小朋友提醒我們.
除了physically.
不要自高自大之外.
我們心理上都不可以自高自大.
都是會跌倒的.
同一天發生的事情.
我行完之後都五點多.
黃昏日落的時候.
平日是沒有小巴.
只有的士去.

$^{521}$那個的士站.
有一大堆人在排隊.
都是講普通話排隊.
我估計只有我一個香港人.
我在排隊.
那些人問我有沒有車的.
有車的.
因為我離遠看到.
有些車有些閃閃燈.
除了的士之外哪有人來.
所以我看到.
離遠有些閃閃燈.
那些是的士來的.
所以我不怕有車的.
接著他又問我一個問題.
先生是打標的嗎.
他問我聽得懂嗎.
是打標的嗎.
我很有文化.
香港人自信.
這裡是香港.
是打標的.
很自信的告訴他.
接著的士來的時候.
一百塊一百塊.
大雪站大雪站.
百湖百湖.
百湖連堂連堂.
百湖.
個個都是.
那些人看著我.
我年輕人太年輕了.
我真的不知道.
原來個個都是泥mank.
這些民族自信.
我自己看高了自己.
看高了身邊的人.
我想互相提醒弟兄姊妹.
我又想和你分享一下.
小心提防.

$^{561}$自高自大.
但有時我們教會又不經意.
有時很多都會追求卓越.
有時會標榜要叻的.
譬如我不經意.
有時都說我們約書亞團.
誰的女兒.
是狀元.
我們都會標榜一下.
忍不住.
我們又要講又要考得叻.
又要不驕傲.
有點矛盾.
誰的女兒厲害.
狀元.
現在去了哪裡讀書.
又說不要驕傲.
即是怎樣.
以下這一段的內容沒有聖經根據.
我自己觀察我的經驗和大家分享.
小心提防.
我們心態上的改變.
其實你有沒有參加過.
不是參加.
在電視上看過.
頒獎禮最佳男主角.
最佳女主角.
通常會怎樣.
通常會多謝.
多謝導演.
多謝編劇.
多謝燈光.
多謝茶水部蓮姐.
是不是這樣.
真心.
其實我們有少少的成就.
其實不是.
不只是我的功勞.
不知你同不同意.
擺很多人在我們身邊.

$^{601}$建立扶持.
使我們能夠成為今時今日的我.
不知你看不看得到上帝的恩典.
我們都為我們身邊的人感恩.
所以給幾秒你想想.
身邊我有很多人幫我.
所以我們太古城堂.
都管理得不錯.
其實是因為有.
杜生 羅姑娘 文姑娘.
還有很多弟兄姊妹幫我.
我們才能夠有今時今日.
不是說很完善.
如果不是他們.
靠我死定了.
我們身邊給很多人.
我都要存著感恩的心.
觀察我們身邊的弟兄姊妹和朋友.
不過說開.
這個考得很狀元.
其實我都有位.
狀元的團友.
我一起成長的.
這位弟兄.
是很謙卑的.
起碼他不高調.
我不知道他原來是9A狀元.
我當年會考的時候只考得9科.
後來可以考10科.
當年是9科年代.
他是一個9A狀元.
不過我不知道.
20年後才知道這件事.
我知道他厲害.
但不知道他這麼厲害.
有什麼專訪.
訪問20年前的9A狀元.
我看到他很面熟.
他就是我的團友.
於是我才想起.

$^{641}$原來他當年是9A狀元.
這個就是.
我當年聽他分享的故事.
我希望可以放在我們心中.
我當年聽他說.
他有四個朋友.
是同學.
他有四個同學.
平時他們做什麼活動.
他們平時做的活動就是.
我昨晚計算的數字很難.
我計算了半小時.
我都有條數字.
我計算了九個字.
看你半小時行不行.
他們有四個朋友經常做這些事.
其他人見到就會走開.
只有他們四個這樣做.
但他們四個.
互相傳授.
我後來問他們.
你們幾個的結果是怎樣.
一個是9A.
一個是8A.
一個是7A.
一個是6A.
感不感受到互相提攜.
互相抵禮同行.
一起成長.
給不同的同伴.
和我們一起同行.
不是我們不好.
我們都有些努力.
但是神給我們恩典.
我們可以和其他弟兄姊妹一起成長.
結伴同行.
在說到這裡的時候.
我還想和你分享一件事.
不知道你有沒有這個感覺.
或者感受.

$^{681}$有時有些人不驕傲.
但聽到又覺得有點硬.
有什麼硬.
他好像很謙卑.
但又不覺得他不是謙卑.
他經常說那些.
其實松柏團不少.
牧師我不識字.
我又老又窮.
又不識字.
你不要找我.
不知道有沒有見過.
不少.
好像很威風.
我又老.
你年輕人.
你又識字.
我知道你識字.
但你又不要驕傲.
年輕人.
不要覺得自己識別人不識.
驕傲就會跌倒.
有時我自己的感受.
聽著牧師解釋.
我覺得你都要悔改.
要改變.
為什麼呢.
因為我很多時候都會和這些弟兄姊妹.
提醒他們.
你經常用這些標準來衡量自己的時候.
其實你會忍不住衡量你身邊的人.
你怎麼可能會分得開呢.
如果比你更窮.
更老.
更不識字.
那個人就更差.
你會不會這樣想.
肯定你說謊.
你不停地想自己都會這樣想別人.
所以.

$^{721}$你窮,老,不識字.
那就窮,老,不識字.
上帝都可以使用的.
你都可以做一個.
又忠心又良善的僕人.
個個都做得到.
不是叫你驕傲.
是因為我們個個都可以.
完成上帝給我們的託付使命.
做一個忠心良善的僕人.
你和我都一定做得到.
最後.
牧師還想講一個故事.
給大家聽.
最近發生.
有位長者.
離世.
98歲.
據聞.
98歲.
她不是我的教友.
這個婆婆98歲.
她的孫子是回我們教會的.
所以我都順手去探望她.
我都算試探望她很多.
因為.
其實都是順手.
我們太古城堂.
總有幾個東區醫院.
順手就順手去見.
但我每次.
都只能夠見到一個.
兒子,女兒.
她臨離世之前.
通常醫院就會打電話.
叫你家人來.
那個就是記號.
差不多不行.
他們家人都叫.
其實我不是她教會的.

$^{761}$婆婆不是我的教友.
叫我去不去.
當然要去.
她都是天父的小陽.
她回到天父身邊之前.
見到我回到天父身邊.
牧師都幾勤力.
我有少少.
日後官面.
那天我去到.
給我一個非常強的經歷.
因為我真的一次過.
見到她所有的子女.
我從來沒有.
見到這麼多人.
我感受非常強烈.
她的子女.
個個都很聰明.
其實有數得計.
一個九十八歲的婆婆.
她的子女都幾年輕.
起碼都七歲.
座底都七歲.
所以個個都退休.
但我聽他們這樣說.
這個不知道在醫院.
做什麼治療師.
這個在海關做什麼主任.
這個又不知道什麼處.
又做什麼.
所以個個都很聰明.
還要想想.
七十歲的那些人.
那些聰明的都很聰明.
那年代讀到書.
所以個個都很聰明.
但很特別.
我見到最多.
這個我見到最多.
有一個女兒是不聰明的.

$^{801}$什麼叫不聰明呢.
我不懂形容.
總之跟她說兩句就知道.
她是不聰明的.
一講就知道.
不聰明.
但又不是要被胡公搶.
總之是不聰明.
但那天我為了這位婆婆.
感恩見了她所有的子女.
我為了婆婆感恩.
天父賜給她.
這位女兒.
是最貼身的.
貼身地服侍媽媽.
她媽媽有好幾年的時間.
在老人院.
但她媽媽從不孤單.
從沒有欠缺.
是因為有這位女兒.
我見到.
因為我每次去探望她.
差不多都見到這位女兒.
我見到她跟她媽媽聊天.
說話.
我知道神很賜福給她.
伯母.
反過來.
那些聰明的子女.
我以前沒有見過他們在一起.
我見了之後.
強烈感覺.
多多說話.
麻麻煩煩.
每件事都拿來爭辯.
那一年.
我心想.
幸好我不是牧師.
我和他們開會就死定了.
感受到嗎?.

$^{841}$上帝使用的是.
忠心.
良善.
神越南.
我們每個人都做得到.
你同意嗎?.
\newpage



\section{路加福音 1:17}
\label{sec:031629}
\textbf{聖經人聖經事:在主面前為大}~~~~黃克勤牧師博士
\newline
\newline
連結: \href{http://www.church.com.hk/acms/content.asp?site=cdc\&op=show\&type=product\&code=031629\&layout=sermon}{\texttt{031629}} ~~~~ 語音日期: 2026-01-31
\newline
\newline
\hyperref[sec:031628]{\small{< < < PREV SERMON < < <}}
~
\hyperref[sec:index_chronic]{\small{[返順時目]}}
~
\hyperref[sec:031630]{\small{> > > NEXT SERMON > > >}}
\newline
\newline
路加福音 1:17
\newline
\begin{longtable}{cl}
\hline
\hline
章節 & 經文 (和合本修訂版)\\
\hline
1:17 & \begin{tabularx}{0.7\textwidth}{X} 他將有以利亞的精神和能力，走在主的前面，叫父親的心轉向兒女，叫悖逆的人轉向義人的智慧，又為主預備迎接他的百姓。」 \end{tabularx} \\ \\
[1ex]
\hline
\hline
\end{longtable}
$^{1}$各位大家好.
又讀聖經了.
《生命述書》的網絡.
聖經人聖經士又和大家見面.
今次仍然在講.
「思洗若寒,在主面前為大」.
「在主面前為大」.
這麼重要的?.
在主面前很大是什麼意思呢?.
留意一下.
《若露加福音》第一章十四十五字.
「你必歡喜快樂」.
「有許多人因祂出世也必喜樂」.
「祂在主面前將要為大」.
「探酒,籠酒都不渴」.
「從無福裡就被聖靈充滿了」.
這是在講什麼呢?.
是在講上次昨日.
不知道大家記不記得.
就是撒加里亞在聖所裡面.
祭司來的.
就在那裡公祭,敬拜,禱告.
剛剛抽籤就輪到他了.
這也是主的安排.
就是在祭簽的時候.
就輪到他是亞比亞班的一個祭司.
還要是每一班祭司都有一個時段要公祭.
每一班裡面有好幾位祭司.
那誰進去呢?.
不是全部進去的.
因為時間都很長的.
輪到他.
進去之後.
工作的那段時間就不出來了.
那這邊做什麼呢?.
很睏就睡一會兒.
那不行的.
要做那個燈台.
要看那些陳設餅.
就在那裡燒香祈禱.

$^{41}$要做這些事情的.
那上主就被他抽到進去.
進去的時候他就在工作.
突然間看到香壇旁邊有天使.
哎呀,大喜一驚.
而天使就跟他說不用怕.
然後他就說.
你會生一個兒子的.
生一個兒子?.
他就想.
我已經很老了.
怎麼生呢?.
當你兒子出生的時候.
他的名字就叫做約翰.
名字也說了.
先是說.
嘉伯列.
他理必歡喜快樂.
有很多人因為約翰的出生.
就很喜樂.
那誰會很喜樂呢?.
就是等候主再來的人.
等候不是再來.
我們等候主再來.
等候彌賽亞尼的人.
心是朝向主的人.
就見到約翰施修翰將來會悔改.
重新跟隨主.
很喜樂.
這個是會發生的.
那些人就很喜樂.
這個人約翰施修翰出現之後.
他來這個世界出生之後.
他在主面前張耀為大.
而且他是淡酒濃酒都不渴的.
這些字就表示了特別淡酒濃酒都不渴.
猶太人是喝酒的.
如果我們基督徒不喝酒.
誰說的?聖經說的.
你真的沒有看過聖經.

$^{81}$猶太人或以色列人.
個個男女都喝酒.
而且如果他們有能力.
他們都會在自己的後山.
做一個酒駕來押酒池.
用來釀酒喝.
因為水有些味道.
不是很乾淨.
個個都有自己的酒.
淡酒濃酒都不渴.
這個人一定是舊約聖經裡面的拿西爾人.
是分別歸主為聖的人.
他們除了不剪頭髮.
和清酒濃酒都不渴.
不能夠接近死屍.
歸主為聖.
自己分別為聖的人.
這個原因使他由出世做拿西爾人.
所以他沒有跟隨主耶穌做門徒.
第一點和螢光的第一點.
在主面前將要為大的意思.
不是大小.
高低.
而是被重視.
在主面前他是會被重視的.
這個人我很看重他.
這個意思.
如果主看重你.
上帝看重你.
覺得你不是一個重要的人物.
這件事就很好了.
很大的恩典.
很大的祝福.
他在主面前是要為大.
他是一個拿西爾人.
而且他在母腹裡被聖靈充滿.
這個特別了.
第二點要說的是.
屍死有恨有第二個特別大的恩寵.
是什麼呢?.

$^{121}$就是被聖靈充滿.
被聖靈充滿.
我們也可以是.
被聖靈充滿的意思.
不是忽然之間.
很厲害.
一指你就死了.
這樣躺下.
一看就看出你的內心是怎樣.
你偷了別人東西.
充滿靈力.
病的起心.
不是那些.
被聖靈充滿.
是聖靈在你心裡做一個主導者.
他就好像司機駕馭著我們這個人.
這輛車.
他做我們的司機.
想去哪裡左右他帶領.
就是聖靈充滿.
就是做我們一生人的主.
被聖靈充滿.
聖靈是帶領到我們.
很多人是聖靈不帶領的.
為什麼?.
因為他沒有被他帶領.
把他困住.
聖靈是什麼時候?.
四離幫我.
感動什麼?.
平時是自己做主.
被聖靈充滿.
而且我們也可以被聖靈充滿.
但在施洗行來說.
他是在沒有福女就被聖靈充滿.
所以這是上帝特別的恩典給他的.
在這裡看到.
他有兩個特別大的恩典.
第一個在主面前.
他是被重視.

$^{161}$第二.
沒有福女就被聖靈充滿.
有時候你會看到.
為什麼這樣?.
為什麼特別給他這麼多?.
沒有給我.
因為上主要用施洗約翰.
地道為尼采亞自己的兒子.
預備他的道.
修直他的路.
要在曠野裡生活.
將人呼喚出來.
他吃蝗蟲野蜜.
穿著駱駝毛的皮.
還有.
聞說他時時禱告.
他的膝蓋好像駱駝的膝蓋.
給你行不行?.
給我行不行?.
不要啊.
福就要.
辛苦不要.
我們很多時候是這樣.
要知道不是這樣的.
不容易的.
這個做施洗約翰.
真的不容易.
有時候我們都會說.
為什麼神選擇以色列人做選民.
不選我們中國?.
你看回以色列人的歷史.
和上帝對他的選民的要求.
和他要選民去見證他.
而沒有成功.
你看選民,以色列人.
在他的歷史裡遭遇多大的迫迫.
多大的鞭打.
多大的摧殘.
為什麼?.
因為選民.

$^{201}$這個特別神要在全世界面前.
要做一個樣板.
千萬不要搞.
幸好沒有選我們中國做選民.
中國沒有被選做選民.
都慘成這樣.
如果真的被選了.
真的慘.
因為要你做一個模範生.
怎知道這個模範生.
是失敗.
於是就虎子已經在樹根上要斬下來.
所以你看到真的不容易.
看這些當然好.
他在神面前要重視.
在武福就被聖靈充滿.
不過你知道他一出生.
他就要做一個拿西爾人.
全年歸給主.
清酒濃酒都不渴.
這樣不行.
可以有這些恩典.
但又可以自己為所欲為.
當然不行.
有時候不是那麼容易.
快點看下去.
他要使許多以色列人迴轉歸於主.
他們的神.
他是有以利亞的心智能力.
第三就是他很厲害.
他除了被主重視.
除了在武福就被聖靈充滿.
第三就是他有以利亞的心智和能力.
很強.
以利亞是以色列人最尊敬的先知.
原因就是.
他是這麼多個先知裡面.
是最有能力.
最強行.
最多神蹟.

$^{241}$偉大到不得了.
單挑巴力先知.
450人.
多厲害.
降火下來.
這麼厲害.
一禱告三年零八個月不下雨.
再禱告又下大雨.
以利亞很厲害.
他能夠叫撒勒法寡婦.
吃剩的麵都還有.
很厲害.
這麼厲害.
他有這樣的心智能力.
是神叫他這樣做.
他要做的工作就是能夠扭轉.
使以色列人迴轉歸於主.
他們的神.
而且他能夠叫為父的轉向兒女.
叫搏逆的人轉從二人的智慧.
又為上主預備合用百姓.
他做什麼呢?.
他能夠叫父的心轉向兒女.
這個父的心是指父神.
父神已經很不喜歡他的百姓了.
以色列人.
他來到.
行在主面前.
叫父神的心轉向兒女.
再次疼愛兒女.
而兒女就變成了搏逆.
但是他們轉從二人的路.
走回正路.
會這樣.
為主去預備合用的百姓.
這個就是以利亞做的工作.
因此你看到.
真的不簡單.
這個就是撕洗了汗.
所以在今天來說.

$^{281}$看馬拉基書.
馬拉基書就說.
當主的日子要來.
彌賽亞要來的時候.
才有以利亞會來.
所以今天的禹節晚餐.
以色列來說.
他們整個禹節晚餐的禮儀.
如果你參加過主持.
都會知道.
他們有一張椅子是沒有人坐的.
是給以利亞坐的.
臨聚會完的時候.
禹節晚餐完的時候.
就叫小朋友去門口看看.
以利亞來了沒有.
接著他回來就說.
沒有,沒有見到以利亞.
好了好了,關門回來.
我們繼續明年一起來守禹節.
等以利亞來.
我們一起再在耶路撒冷守禹節好不好.
好啊,明年見.
他們想著.
以利亞出現.
彌賽亞會來了.
以利亞重現.
其實呢.
這個以利亞就是.
我們等到了.
我們看到了.
就是撕洗了汗.
為主預備合容的百姓.
不過可惜當時的人不知道.
撕洗了汗是有這樣的遺憤的.
跟大家看看今天的金句.
我都選了一節.
選這節有一點點夾硬的原因就是.
沒有什麼金句特別對.
不過這節都重要的.

$^{321}$就是路加福音一章十七字.
他必有以利亞的心自和能力.
行在主面前.
叫遺父的心轉向兒女.
叫背逆的人.
轉從二人的智慧.
又為主預備合容的百姓.
我在這裡.
想用這節金句鼓勵大家的就是.
我希望我們.
你,我.
都會想.
我要為主立志.
會不會我都跟隨經文.
我都做一個以利亞.
嘩,這麼大膽.
做以利亞我沒有這樣的能力.
所有的能力都從上主而來.
不用擔心.
那以利亞的能力是以利亞自己的嗎?.
不是的,都是神給他的.
那施洗約翰的能力是誰給的?.
都是上主給的.
他是會行在主的面前.
就是說看著主的心意.
而跑步做的.
做什麼呢?.
就是這句字.
轉向,轉從.
我希望我們都可以做施洗約翰.
可以做以利亞.
是得著上帝的能力.
是時時想的.
怎樣可以令別人轉向上帝呢?.
在自己的活躍的路上.
轉從二人的智慧.
走回正確的路呢?.
這些活躍的人覺得自己很厲害.
這樣我就賺到錢.
這樣我就害到他.

$^{361}$這樣我就陷害到他.
這樣我就嚇到他.
嘩,這樣他死了都不知道.
很多這些活躍的人是這樣的.
但是他覺得自己很聰明,很厲害.
但是其實二人有二人的智慧.
是上主會保守他的.
夜半敲門也不驚.
都沒有做過虧心事.
其實更有智慧的.
是叫人轉向.
我要學習求主給我能力.
給我施洗約翰.
是時時想,時時禱告,時時做.
叫人轉向上主.
叫活躍的人轉從二人的智慧.
二人更有智慧.
而且為人信主預備好.
這樣就好了.
因此.
希望我們都有以利亞的心智能力.
行在主的心意面前.
叫為父上帝轉向兒女.
人又轉向上主.
行二人的路.
這樣就好了.
願主祝福我們.
有以利亞施洗約翰的心智和能力.
不要怕想大些.
想闊些,高些的事.
倚靠主的能力.
是我們能夠做更大的事.
我們一起祈禱.
天護上帝幫助我們.
我們真的很需要你.
我們吃飯,睡覺,喝水.
上班,下班又一天.
然後又一天.
然後又一年.
又一年.

$^{401}$時間過得很快.
然後已經逐漸覺得自己.
在世上的日子已經過了很多.
好像沒有做到些什麼.
主又幫助我們.
這次聖經提醒我們.
讓我們有一個更大的心智.
要在主面前被重視.
就是因為我們要把自己.
靠著主.
要成為一個叫人轉向你的人.
讓我們時時為主預備百姓轉向你.
叫我們有這個心智與能力.
我們等候主.
讓我們學習做以利亞.
學習做施洗約翰.
有大志的.
要將人帶到你面前.
呼叫人悔改.
垂聽祈禱.
奉靠耶穌基督的名字.
阿們.
\newpage



\section{民數記 4:1-19}
\label{sec:031630}
\textbf{各司其職,有序同行}~~~~林誠信牧師
\newline
\newline
連結: \href{http://www.church.com.hk/acms/content.asp?site=cdc\&op=show\&type=product\&code=031630\&layout=sermon}{\texttt{031630}} ~~~~ 語音日期: 2026-02-08
\newline
\newline
\hyperref[sec:031629]{\small{< < < PREV SERMON < < <}}
~
\hyperref[sec:index_chronic]{\small{[返順時目]}}
~
\hyperref[sec:031631]{\small{> > > NEXT SERMON > > >}}
\newline
\newline
民數記 4:1-19
\newline
\begin{longtable}{cl}
\hline
\hline
章節 & 經文 (和合本修訂版)\\
\hline
4:1 & \begin{tabularx}{0.7\textwidth}{X} 耶和華吩咐摩西和亞倫說： \end{tabularx} \\ \\ \relax
4:2 & \begin{tabularx}{0.7\textwidth}{X} 「你要照宗族、父家計算利未人中哥轄子孫的人口， \end{tabularx} \\ \\ \relax
4:3 & \begin{tabularx}{0.7\textwidth}{X} 就是從三十歲到五十歲，凡前來任職，在會幕裡事奉的人。 \end{tabularx} \\ \\ \relax
4:4 & \begin{tabularx}{0.7\textwidth}{X} 這是哥轄子孫在會幕裡有關至聖之物的職責。 \end{tabularx} \\ \\ \relax
4:5 & \begin{tabularx}{0.7\textwidth}{X} 「拔營的時候，亞倫和他兒子要進去，把遮掩的幔子取下，用它來遮蓋法櫃， \end{tabularx} \\ \\ \relax
4:6 & \begin{tabularx}{0.7\textwidth}{X} 又用精美皮料蓋在上面，鋪上純藍色的布，再把槓穿上。 \end{tabularx} \\ \\ \relax
4:7 & \begin{tabularx}{0.7\textwidth}{X} 他們要用藍色的布鋪在供餅的桌上，將盤、碟，以及澆酒祭的杯和壺擺在上面；經常供的餅也要留在桌上。 \end{tabularx} \\ \\ \relax
4:8 & \begin{tabularx}{0.7\textwidth}{X} 他們要在這些東西的上面鋪上朱紅色的布，把精美皮料蓋在上面，再把槓穿上。 \end{tabularx} \\ \\ \relax
4:9 & \begin{tabularx}{0.7\textwidth}{X} 他們要用藍色的布遮蓋供職用的燈臺、燈臺上的燈盞、燈剪、燈盤，以及所有盛油的器皿； \end{tabularx} \\ \\ \relax
4:10 & \begin{tabularx}{0.7\textwidth}{X} 又要用精美皮料把燈臺和燈臺的一切器具包好，放在抬架上。 \end{tabularx} \\ \\ \relax
4:11 & \begin{tabularx}{0.7\textwidth}{X} 他們要用藍色的布鋪在金壇上，用精美皮料蓋在上面，再把槓穿上。 \end{tabularx} \\ \\ \relax
4:12 & \begin{tabularx}{0.7\textwidth}{X} 要用藍色的布把聖所供職用的一切器具包好，再用精美皮料蓋在上面，放在抬架上。 \end{tabularx} \\ \\ \relax
4:13 & \begin{tabularx}{0.7\textwidth}{X} 他們要清理祭壇上的灰，用紫色的布鋪在壇上； \end{tabularx} \\ \\ \relax
4:14 & \begin{tabularx}{0.7\textwidth}{X} 又要把供職用的一切器具，就是祭壇一切的器具，火盆、肉叉、鏟子和盤子，都擺在壇上，鋪上精美皮料，再把槓穿上。 \end{tabularx} \\ \\ \relax
4:15 & \begin{tabularx}{0.7\textwidth}{X} 「拔營的時候，亞倫和他兒子把聖所和聖所一切的器具蓋好之後，哥轄的子孫才好來抬，免得他們摸聖物而死；這是哥轄子孫在會幕裡所當抬的。 \end{tabularx} \\ \\ \relax
4:16 & \begin{tabularx}{0.7\textwidth}{X} 「祭司亞倫的兒子以利亞撒所要照管的是點燈的油和香料，以及常獻的素祭和膏油。他要照管整個帳幕和其中所有的，以及聖所和聖所的器具。」 \end{tabularx} \\ \\ \relax
4:17 & \begin{tabularx}{0.7\textwidth}{X} 耶和華吩咐摩西和亞倫說： \end{tabularx} \\ \\ \relax
4:18 & \begin{tabularx}{0.7\textwidth}{X} 「你們不可使哥轄人宗族的這一支從利未人中剪除。 \end{tabularx} \\ \\ \relax
4:19 & \begin{tabularx}{0.7\textwidth}{X} 他們挨近至聖之物的時候，要向他們這樣做，使他們存活，不致死亡；亞倫和他的兒子要進去，分派各人當做的，當抬的。 \end{tabularx} \\ \\
[1ex]
\hline
\hline
\end{longtable}
$^{1}$讓我們一起來看看今天要講的經文.
其實經文那裡是有一個標籤的.
那是四章一到四十九節.
是整章的.
所以我昨天晚上聽到主席這樣說.
我就嚇了一跳.
我說咦 那我是不是可以早點講完呢.
那其實時間也是一樣.
我就要在很短的時間裡面.
是將整章文授記第四章要跟大家分享的.
其實文授記第三第四章是很有關係的.
但教大家一個秘訣.
如果我們要將這兩章聖經連在一起.
我們可以去看得通透的話.
兩章聖經都是在講利未蒙神揀選.
在十二個支派揀選利未的支派.
利未有三個兒子.
而這三個兒子是被神揀選出來.
去代替所有以色列的家族每一個的長子.
被分派就是要守護回望.
他們來到是蒙神選照.
用兩個字 英文字.
你可以試一下去綜合第三章講的東西就是Who and Where.
他們是什麼人呢.
他們是分別為姓.
專責回望的服侍的一班人.
這個就是他們的身份.
他們也有一個編制.
他們在回望的南邊西邊和北邊是紮營的.
這三班的利未家族的人.
他們有很重要的職守.
是會防止任何包括以色列其他支派的人.
如果要擅闖進去回望裡面.
是你可以將他處死的.
就是這麼嚴重的.
因為舊約裡面讓我們明白.
神的聖潔是人不可以冒犯.
他要讓人看到是有界線的.
上帝怎樣來顯明他自己被摩西來抗的時候.
摩西很想見神的面.

$^{41}$上帝說你只能夠見我背脊.
如果你面對面見我.
你就死定了.
因為我們人是有罪的.
我們是需要在神的面前.
是蒙他自己的一個寶血潔淨.
我們今天新約的人就很不同了.
我們是靠著耶穌基督.
我們可以進到他這個私人寶藏面前.
第四章所說的.
就是利美這三個兒子.
他們在這個會幕的服侍.
和那個分工實際的情況.
實際的情況.
你用兩個英文字就很容易記了.
What and How.
他們做些什麼.
和怎樣去做呢.
怎樣去做去執行呢.
這個就讓我們很容易明白了.
我們今天看第四章的時候.
神吩咐是數點利未人當中的壯年人.
這個和第三章是有不同的.
第三章就叫由一個月大的嬰兒.
數他們整個人口.
當時利美整個家族在第三章是二萬二千人.
但是他們這章裡面.
說到他們數點三十到五十歲.
是只有八千多人.
我們後文會交代了.
他們按家族來分工負責會幕的拆卸.
搬運和看守的事.
好讓會幕在曠野移動的時候.
是有效率和有秩序.
一切都是按照神的吩咐.
你看到他是很仔細.
這三組的人.
他們要做些什麼.
他們分工很清楚的.
上帝說我們要起行了.

$^{81}$我們不是always station在這裡.
因為曠野他們必須要向著.
這個應許之地進發.
不過你知道神的子民是很不長進的.
他們在近四十年裡面.
他們在曠野裡面.
就不斷有很多埋怨.
他們甚至很想回到埃及.
是因為他們那種信心.
只是看眼前看環境看物質.
那種缺乏.
他們忘記那個呼召.
他們離開埃及遺老之家的.
這一位上帝.
是何等的聖潔公義真實.
而且他是豐富傳遞的供應者.
我們今天同樣.
在我們人生裡面.
無論我們今天是身處什麼樣的處境.
你是在流淚谷.
或者你在曠野裡面.
或者你覺得是不是神丟棄你.
其實上帝永遠都沒有丟棄他的子民.
永遠都是我們忘記上帝.
祂無時無刻其實是看顧.
我們人生無論去到哪個階段.
你是年輕人也好.
中年人也好.
或者你現在去到.
是長者了.
或者下游了.
你都要記住.
神沒有丟棄你.
我們今天我們的健康.
我們每天的飲食.
這一切都在神的掌管下.
你們這三個兒子.
他們是代替亞倫加和惠仲.
去辦理這個會望所有的事.
我們在第三章上一課有跟大家說過.

$^{121}$會望的東邊就是摩西亞倫加紮營.
南邊就是利美的第二個兒子.
就是哥核族.
左邊就是西邊.
會望的西邊就是利美的長子.
格遜族.
而會望的北邊就是利美最小的兒子.
就是米拉利族.
我們看看這個國師奇蹟.
其實分三個段落.
由第一節到三十三節.
你看到第一個段落是最長的篇幅.
是與哥核族有關的.
而第二就是二十一到二十八節.
和格遜族這個長子有關的.
而第三和米拉利族有關係.
二十九到三十三節.
而三十四到四十五節.
就是說他們怎樣回應上帝的吩咐.
他們去數點他們各族的人數.
這一章的聖經的脈絡讓大家稍為了解.
但我就會詳細解釋一下.
頭一個大段落.
這三組的人.
由哥核族,格遜族和米拉利族.
究竟他們相似的地方.
和不同的地方是什麼呢?.
我們先來看.
他們有類似的地方.
首先是神吩咐摩西和亞倫.
去數點你每個宗族的.
這個是由三十到五十歲.
讓他們當中最壯年.
最精壯的人出來去事奉神.
第二就是神吩咐每個宗族.
所負責與回眸有關係的工作.
大家是各司其職的.
不會擾亂的.
也不會重複的.
是很清楚的.

$^{161}$下一個平面圖讓大家容易了解.
在這個圖跡裡面.
你看到中間這個位置.
就是聖所和致聖所.
這個在聖所前面.
這個就是祭壇.
這個回眸.
你看見這三組人.
他們的位置和他們所要負責的工作是什麼呢?.
歌客族.
這個在下面南邊.
他負責要搬運的.
全部是致聖之物.
致聖之物是在說.
在聖所裡面.
這些所有器具.
包括陳灼臺.
還有金燈臺.
以及香爐.
進入致聖所.
就是這個藥櫃.
這些全部是由他包辦.
他要負責搬運的.
外面在這個院子裡面.
就是這個祭壇.
祭壇也是他要負責的.
這個是很清楚的.
好了 革信族.
他們負責什麼呢?.
他們就是負責這個院子的.
是關於舉辦所有的那些報匹的東西.
報匹.
包括就是這個院子的門簾.
這裡.
這個聖所.
會幕的門簾.
還有萬紙.
致聖所和聖所之間.
這個是有個萬紙的.
是遮蓋的.

$^{201}$連普通祭司都不可以植入這個致聖所的.
每年只有大祭司.
一年一次在贖罪日.
他是代表以色列整個的民族.
是向神來贖罪.
我們在這裡有一個了解.
還有這個院子的錐子.
這些全部是他們負責的.
好了.
米拉利族就負責那些工程的部件.
米拉利族.
你看到這裡整個會幕.
他是有那些的板,山,柱,座.
全部是最重的東西.
最重型的東西.
就是由這個錐子負責.
他們這個族裡頭.
所有的這些被揀選出來.
就要搬這些的.
這些既是重.
但也都很重要.
因為你拆了之後.
你砌不回就大件事了.
你明不明白.
因為每一個部件.
他們是全部要有專人去負責.
分門別類.
大家移動到某個新的地方.
一來到下山的時候.
你那支釘.
進去多深.
你的角度怎樣釘好那個錐子.
全部都是工程的力學.
所以是專業人士.
是需要來做的.
當你明白了這個分佈的時候.
我們看這三個段落.
也都有第三點有類似的地方.
就是他每個宗族.
都是有supervisor的.

$^{241}$就是哥核族是祭司以利亞薩來管轄的.
加遜和米拉利族就是以他馬.
我想請問大家記不記得.
以利亞薩和以他馬.
是亞倫的第幾個兒子.
第三和第四個兒子.
為什麼不是長子和次子呢.
是的,第三章你們有聽到.
第三章就是拿達阿比胡.
獻梵火.
他們沒有照上帝的吩咐亂來.
祭壇的火.
你是要用祭壇的火種來點燃.
但是他是找別處的地方.
點燃祭壇上這個火.
所以那裡是說一個unholy fire.
英文很多譯本都有翻譯的.
我們再看這三個段落.
也都有一些不同的地方是什麼呢.
首先你看到剛才我們都說.
哥核宗族的描述是最長的篇幅.
這個是看到他比格遜和米拉尼族.
在神的心意裡面更加重要.
是神揀選的.
我們照問一個問題.
為什麼是哥核族呢.
為什麼哥核族是排第幾.
是利尾三個兒子裡面.
排第二.
不是第一.
第一那個是誰.
第一那個是格遜.
我們會問為什麼上帝這麼有趣.
三個兒子為什麼找次子.
去負責聖所,致聖所的那些致聖之物呢.
這個是上帝的主權.
上帝命定的.
其實這個都不奇怪.
如果你看回舊約聖經.
我們有很多的例證.

$^{281}$那個孿生仔最出名的是以素和雅各.
明明應該以素是承繼所謂長子的名份.
但竟然以素因為紅豆湯.
他是出賣了長子的名份.
而雅各是被神揀選的.
你要接觸亞伯拉罕,以撒.
然後不是以素.
然後是雅各.
這個是上帝祂有祂自己主權的揀選.
你不需要跟別人比較.
今天你在家族裡面也是.
你不需要跟你的哥哥姐姐在那裡比較.
如果神揀選你.
你就去順命而行就行了.
我在很多的.
這麼多年牧會裡面.
我都看到有不同的家庭.
有些是神揀了一個長子出來.
做傳道人的.
但有些不是.
有些揀選可能是第二個兒子出來.
有些是揀了最小的女兒出來.
或者最小的兒子出來.
我就是其中一個例子.
我的家族裡面.
是最小的女兒,最小的兒子.
竟然被神揀選.
我看到神的揀選真的很奇妙.
但這件事是提醒我們.
你不需要跟別人比較.
你要忠於上帝給你的身份,位份.
還有你有沒有忠心.
按照神給你的恩賜.
你來擺上.
你一生無論你是傳道人與否.
你是神的兒女.
你就要寶貴你的名份.
因為神的兒女同樣.
是在神的家裡面.
我們是要各司其職的.

$^{321}$我等一下就會跟大家再講下去.
我們怎樣應用在我們「廣福家」裡面.
我們「廣福家」一點都不小.
我們有很多不同方面的需要.
我們再來看第三個分別在哪裡.
提到那個辦事只有哥核族.
他是用這個希伯來文.
叫做「默勒卡」這個字.
在第三節.
去表明這個事是需要特別技能的.
去負責「致聖之物」.
在第三節裡面我們看到.
這些三十到五十歲的人.
是來到會幕裡面事奉的人.
他特別在這裡說.
是哥核子孫.
是要負責有關「致聖之物」這個職責.
這一種專職的事.
是用「默勒卡」這個字.
其他那兩個族.
他們是用這個「阿伯達」.
「阿伯達」這個字是說勞動.
勞動只是說他們搬台搬那些東西.
「阿伯達」.
那你認識了兩個希伯來文.
不過你不要亂拋書包.
人家問你第三個你就不認識了.
我們起碼在這裡.
為了要明白.
原來上帝透過那個文字的敘述.
是有分別的.
這個分別我們再來要看到.
這個描述是要凸顯哥核族的地位.
是高於他的哥哥和他的弟弟.
那個格遜和米拉尼的宗族.
在「抗夜路上」按神的吩咐起行的時候.
是先由哥核族搬運.
他首先是將聖所,致聖所那些「致聖之物」搬來.
我們先來看第一段.
第五到第六節.

$^{361}$就是「約季」.
我們一起讀一下.
「拔營的時候阿倫和他兒子要進去.
擺遮掩的慢紙取下.
用它來遮蓋法季.
又用精美皮料蓋在上面.
鋪上純藍色的布.
再擺供村上」.
這裡是說聖所,致聖所.
這裡面所說的這些.
這裡只是說.
其實是致聖所.
第一樣東西.
是「約季」那些東西.
說了三件事.
是祭司做的.
不是其他人做.
因為要進入聖所,致聖所之後.
祭司才能夠進入.
但是這裡面說的這三樣東西.
「遮掩的慢紙」.
「覆蓋法季」.
但他接著說.
這個是第一層.
第二層是精美的皮料.
這個是和修本的翻譯.
我們和合本的翻譯.
都是在說「海狗皮」.
我參考其他的中文譯本.
都是在說「海狗皮」.
參考希伯來文的原文.
都是指向「海狗皮」.
我不是很明白.
為什麼和修本修訂之後.
竟然.
他可能覺得這樣.
是不是比較容易.
一些精美的皮料.
確實「海狗皮」是很精美.
它是能夠厚,實.

$^{401}$是耐磨的.
就是你要在曠野.
要搬運的時候.
不要有磨損.
這個是有意思的.
最後第三層.
就是純藍色的布.
這個是代表著.
這個是屬天的.
這個是上帝的寶座.
這三樣東西.
接著由第七節到第十四節.
就是其他聖所.
去到院子的祭壇.
因為時間有關係.
我用一個圖表.
解釋給大家知道.
好不好.
這樣你們就可以明白了.
這個圖表.
是有這些.
怎樣去包那個次序.
第一塊布.
第二塊布.
第三塊布.
是什麼的性質.
陳設祥.
這裡面.
他是說.
這個藍色的毯子.
珠紅色的毯子.
然後「海狗皮」.
用這三樣東西.
來包裹著.
燈台也是.
香壇也是.
到外面的祭壇也是.
你會發覺到最後.
他們全部都是.
「海狗皮」在最外層.

$^{441}$來包紮.
這些聖物.
我們看了這個圖表之後.
你再看下去.
第十五節.
這個是重新覆蓋上列.
每一件物件.
是必須由祭司親自執行.
然後才由歌俠族去搬運.
但是歌俠族不可以直接摸聖物.
否則他們會死亡.
這個很嚴重的一件事.
是規矩來的.
所以你是負責致勝之物.
你就要打穩十二分精神.
你不可以說.
我可不可以.
那個藥櫃.
我很好奇,我想看看藥櫃.
他包紮完之後.
你去抬的時候揭開.
揭開,我看一看.
你馬上死亡.
等於獻梵火一樣.
大家明白嗎.
所以後來.
你看到他們舊約聖經也是.
那些祭司.
就算抬藥櫃的時候.
所謂失前提那件事.
跌下來的時候.
他走去參附.
那件事有另外的解釋.
因為根本他們都不應該找那些.
去牛車去搬.
這個是完全兩件事.
不過起碼我們要知道.
神是有法規.
有祂的這些準則.
你要聽話.

$^{481}$你不要好奇.
也都不可以自作主張.
不可以隨隨便便.
在這一切的事情上.
好.
我們看一看哥核族要負責抬這些致聖之物.
在祭司以利亞撒的督導下.
他要負責搬那些東西.
就包括我剛才說的所有這些.
會墨的萬紙.
罩盤,海狗皮,蓋會墨的門簾.
丸子的錐子,門簾和繩子.
圍繞遮蓋會墨有關的所有這些物件.
包裹完之後.
他才可以搬.
全程我們都看到.
革信族負責另外這些.
他們也都在以他馬的祭司的督導下.
這個牽涉到他們布匹那些東西.
遮蓋.
表示神的同在的遮蓋和界線.
但同樣要很忠心和細心去管理.
最近我和師班他們的師班長.
幾個師班我們一起開會.
我也都邀請了我們教會.
一個是時裝設計的一位姐妹.
她在香港做了一個很出名的品牌.
她本身這一方面是神給她很優秀的恩賜.
我請她做我們的一個顧問.
給我們一些意見.
做什麼呢.
就是為了做一件事.
我們想做新的師班舖.
你現在見無論是喜樂師班.
或者我們廣福師班.
他們現在姐妹通常都是.
黑褲.
黑裙.
白色的衣服.
加一條有顏色的絲帶.

$^{521}$如果是指令團就做一個紫色的絲帶.
但是我們發覺姐妹就好一點.
容易一點.
但是男士呢.
今天沒有男士給大家看.
你看廣福師班的時候走出來的男士.
是叫他們穿深色的西裝.
黑色的.
但是黑色原來有深黑色.
有淺黑色.
有不知道什麼不同種類的黑色.
有些是深藍色.
加上白色的恤衫.
領帶就更麻煩了.
一領帶就說紅色的領帶.
誰知道紅色的領帶有棗紅色.
有森紅色.
有鮮紅色.
就好像不是很整齊.
所以我們最近就希望能夠.
做一套新的師班舖.
都是莊重.
希望能夠代表.
他們就好像在昔日的.
那個聖殿裡面那種的服侍.
他們的歌頌讚美神.
代表會眾呈獻禮物給神.
也都帶動會眾一起唱歌讚頌神.
在整個的過程裡面.
這些是很需要.
是大家有這一種.
有莊重.
但是我們又想莊重得來.
和現代的時代.
又有些融合.
所以都有很多意見.
所以請大家為我們祈禱.
我們求主給我們有聰明智慧.
可以做一個新的.
廣福師班的這個師班舖.

$^{561}$會出現.
好了我們再看這個細仔負責.
原來是最重要的.
就是那些工程的結構.
這些組件.
包括帳幕的板.
山柱子.
大樓的座.
和丸子四周的柱子.
蝎子.
繩子和一切器具.
他們是按明.
指定各人所抬之物.
按明有意思.
因為他們到期時拿了那塊東西.
他知道他是砌在第幾塊板那裡.
有所有次序的.
不會亂的.
還有他們就是負責.
是搬這些最重.
但是都是涉及工程的部件.
這些是需要體力.
但是亦都需要是精準的分工.
事奉有時亦都是需要有秩序和周全的安排.
大家記不記得我們在.
前年 現在說回來.
在二四年十月的時候.
我們因為教會沒有那麼大的地方.
容納我們四堂崇拜.
加上我們的國語.
加上英語.
我們是說接近二千人的會眾.
所以我們是租借了恩福堂.
用它的大禮堂.
去舉行我們四十週年感恩崇拜.
但之前我們是做了很多很多的會議.
和在籌組的過程.
我們要做rehearsal的.
即是說是當日星期日.
我們這個四十週年的感恩崇拜之前.

$^{601}$我們有兩次去到恩福堂.
我們要請每一個上台.
有份要負責帶領無論是詩歌敬拜.
或者是負責做主席.
負責講訊息.
每一個人.
每一個傳道同工.
或者執事.
我們有份的就要去.
有個預演.
我們要試麥克風.
試燈光.
試powerpoint.
在整件事裡面.
是一點都不簡單.
你要將整個教會.
說的是二千人的教會.
要搬去另一個地方.
絕對不容易.
我們在十二月.
在海洋公園MARRIOT酒店.
舉行兩場四十週年感恩福音的午宴和晚宴.
午宴都是吃晚宴的那些東西.
因為我們不想被人說我們口齒薄鼻.
所以有些傳道同工.
其實兩餐都要吃.
你不吃就浪費了位置.
所以那次我們其實.
大家好感恩.
林曉樂牧師跟我們講了訊息.
我們都有很多親友.
在那兩場的聚會裡面.
他們得聞福音.
他們一起跟我們在這個很歡喜快樂的日子.
我們去經歷.
就是數算神的恩典.
我們去頌讚神的作為.
我們是扶老攜幼.
就是想我們整個家族.
我們的兒女和我們的兒孫.

$^{641}$都一起去參與.
一起去認識.
去明白上帝奇妙的作為.
利未三子.
他們也都順命去數點人數.
在最後這個段落.
三十四到四十五節.
這些數字我是照抄的.
在聖經裡面.
你可以根據這些經文看到.
從三十歲到五十歲.
你有任職的利未人.
總人數是八千五百八十人.
佔這個總人口.
在三章三十九節的二萬二千人.
如果你心算快.
或者你有計數機.
你一按下去.
這兩個數字正確是三十九個百分比.
三十九個百分比.
三十九個百分比不是聖經講的.
不過我是根據聖經計算出來的.
三十九個百分比給我們一些啟示.
一般教會.
我在北美有住過.
我在北美亦都服侍過.
不同的地方包括東南亞.
香港.
國內都有.
一般很多的教會.
就算是很大型的教會.
有時我們發現一個問題.
會出現是甚麼呢.
那個事奉人數的比例.
是沒有三十九個百分比.
沒有三十九.
如果你看一些教會增長的書.
一般來說.
其實很多教會都停留在十個百分比.
十個百分比的事奉人員.

$^{681}$去服侍九十個百分比的會眾.
這個是不健康的.
在我們前行的日子.
講福家.
我們要繼續起行.
我們要繼續是跨代同行的話.
我們要達至這個理想的時候.
我們必須要鼓勵更多弟兄姊妹出來.
所以今年教會去推動這個門女.
是很重要.
是commit而去參與.
我們在未來這兩年.
是會積極重點培訓他們.
他們每天去靈修去祈禱.
這些是基本功.
這些基本功是很重要的.
沒有了這些基本功的時候.
我們忙於服侍.
我們只會很快就burn out.
我們很快會殘.
而且我們很快就會發怨言.
起爭論.
為什麼你沒有一個好的靈命.
你跟神的關係不密切.
你怎麼能夠跟人一起去相處和同工呢?.
這個是很重要.
我今早起來靈修.
我自己看的經文.
就是《大使命》.
我看到最感動我的地方.
是主耶穌說到.
你們凡我所吩咐你們的.
你們都要遵行.
我就常與你們同在.
直到世界的末了.
我自己深深感到.
今天我們這個世代很混亂.
但我們不需要懼怕.
我們因為最重要.
是有主耶穌活在我們的心裡面.

$^{721}$你照著祂的藍圖和心意.
照著祂所教導.
你去教導弟兄姊妹.
讓我們都能夠成為主所喜悅的門徒的話.
他們由信主到栽培到洗禮.
到認識主.
在聖經裡面所有這些真理的教導之後.
你就會發現一件事.
你從來教會就一直不會缺人.
你可以一直來建立更多更多.
不單止是有傳道人.
我們更需要有很多弟兄姊妹.
一起來服侍.
所以我們嘗試從利未人的服侍裡面.
學習幾件事.
我們作為一個提醒.
第一 他們並非隨意自選.
不是他們喜歡或者不喜歡.
他們很清楚他們服侍的崗位.
是神揀選他們.
有神呼召.
而且第二 是按恩賜.
和家族被指派的各自的位置.
哥核族他們是負責致聖之物.
而米拉尼族他們要負責搬很重的東西.
他們不會抱怨的.
他們不會埋怨的.
他們知道這是他們的位分.
這是神安排他們的崗位.
這種按恩賜和信服神一起去服侍.
是很美的一個圖畫.
而哥核族雖然是賜旨.
但是被神揀選去抬這個致聖之物.
不過你不要以為舒服.
這個服侍要打穩十二精神.
因為稍一不慎.
你是被擊殺的.
因為你是與神的器皿是最親近的.
親近神我們要帶著敬畏和信服.
我們今天看「國師奇蹟」有什麼反思呢?.

$^{761}$利未人他們這些全部是有規格.
他們去看守會望和搬運這些器皿的時候.
包括帳幕和祭壇所有這些物件.
我們今天在崇拜裡面.
你看到不只是台上服侍的人.
其實有大量的對公.
除了敬拜隊之外.
我們中堂,五堂有領師的弟兄或者姐妹.
我們有講道的牧師傳道.
我們有做主席的.
有傳道同工來幫大家祈禱的.
有長執來治理的.
有迎新大使在外面.
歡迎我們每一堂的新朋友.
介紹教會給他們認識.
我們有場地的堂務.
有清潔,有物業管理.
有前台的職員去接待.
去接受別人的查詢.
還有很多的行政人員.
裡面有堂務,有文書幹事.
還有一個專業的團隊.
包括會計,包括人事.
包括做美術設計的.
包括還有IT的.
負責IT的.
這些全部他真的要不認識就不認識.
你不可以隨便找一個人說.
你有空不如幫我弄張海報出來.
你不知道做了些什麼出來.
我們今天是需要大家彼此各按各職.
我們還有教導,牧養.
這些除了崇拜,講壇之外.
我們有團契,講道,教導.
我們除了有教牧同工之外.
還有主學老師,團契組長.
輔導這個也是極其需要.
我們要與這麼多弟兄姊妹同行的時候.
我們必須教牧同工或者團契老師.
他可以在第一站裡面.

$^{801}$作為一個輔導.
我們可以了解這個弟兄姊妹.
她是不是真的有一些情緒病.
她是不是有抑鬱病.
這些我們是有過一些訓練.
我們起碼初步可以有些判斷.
但如果嚴重的話你是需要轉介.
去看專業的輔導.
這些我們也是會陪伴他們去面對.
最後我們看有序同行這個反思.
我們廣福家其實做了分靈牧養很多很多年.
但是我們今天說分靈牧養有它的好處.
你叫牧養年青人那些學生.
你到牧養中年人.
他們在就業裡面去面對那些壓力是不同的.
那個處境是不同的.
我們怎樣能夠和他們有效去同行呢.
這個我們鼓勵大家是跨代同行.
有跨代同行的好處就是.
年資深的一些是夫婦.
可以給年幼一些的夫婦.
給他們一些借鏡.
幫助他們看他們生活上一些的需要.
他們的孩子們.
他們怎樣來和他們同行.
但是最重要.
無論你今天和XYZ世代同行.
最重要也都是最好的.
就是要與主同行.
與主同行.
你沒有與主同行的時候.
你沒有力量.
你也沒有智慧.
那你和什麼人去同行.
你也會同歸於盡.
我告訴你們.
但是我們需要很好的.
我們今天是在這個雖然混亂的世代.
但是你必須要做一個與主同行的一個神喜悅的門徒.
我們又思想另外一件事.

$^{841}$就是你有沒有這個世代裡的父母心.
因為我們很常都會聽到掛在嘴唇邊.
哇 醫者父母心啊.
這個醫生有醫德啊.
他看病很有心機啊.
但同樣牧者我們需要有父母心.
保羅說我們是用母親那種溫柔去育養弟兄姊妹.
但同樣要有父親那種威嚴去督澤或者去教導弟兄姊妹.
去明白上帝的心意而行.
所以恩威並施.
這個都是需要.
這個需要我可以總結來說.
我們今天這個世代的父母心是什麼呢.
最重要的父母心是什麼呢.
我可以告訴大家.
我從聖經看的.
一定離開不了這件事.
這件事你就是必須要帶領你自己的兒女和你的兒孫.
要認識神的話.
要認識耶和華是哪一個.
讓他不要失落.
不要失落.
你看一看在約書亞記.
約書亞林中之前.
他說至於我和我家.
我們必定要事奉耶和華.
他要去得地為業.
要進入迦南地.
他要征戰.
他很清楚這件事.
但他在世的時候.
以色列人是事奉耶和華的.
但他死了之後.
他留下了一些在世的長老.
長老還在的時候.
以色列人事奉耶和華.
但鏡頭一轉.
沒多久在士師記.
他們進入迦南地.
你看看長老還在世的時候.

$^{881}$百姓都事奉耶和華.
但第十節.
士師記二章.
他說另一個世代興起.
但他們不認識耶和華.
也不知道耶和華為以色列所做的事.
以色列行耶和華眼中看為惡的事.
去事奉朱巴力.
很恐怖.
今天我們必須謹記.
我們要和每一個年齡層的弟兄姊妹.
與主同行.
讓我們認識神的話.
我最近看一些資料統計.
在香港大學.
在去年九月發佈一個研究.
二十五歲到五十九歲.
香港的男士自殺率急增.
由二零二一年的十六點九個百分比.
增加到二零二四年.
是二十三點四個百分比.
差不多是四分之一.
四個有一個.
這些壯年的男士.
二十五歲到五十九歲.
尋求自殺.
原因是什麼.
壓力是財務最主要的原因.
欠債或者供樓.
很多問題.
我們都知道.
由整個COVID到二十四年.
其實我們整個香港的經濟都是差的.
這個不只是香港.
有很多其他地方都一樣.
但是我們祈禱求神帶領.
進入二六年.
我們新的一年裡面.
我們會看到.
似乎有很多數據都說.

$^{921}$這些經濟有些起色.
回春.
我們求主讓我們不要灰心.
無論你今天在什麼處境.
你知道神沒有丟棄我們.
耶和華是我們的牧者.
是我.
如果我真的以耶和華是我的牧者.
這個「是」是現在式的.
不是過去的.
不是將來.
是說「是」.
上帝永遠都是所「是」.
祂今天在我們任何處境.
祂都與我們同行.
與我們同在.
問題是.
是不是你願意.
我們一起同心禱告.
天父上帝求你幫助我們港福家.
每一個家庭.
都能夠愛主更深.
我們都看見這個異象.
我們必須跨代同行.
我們要帶領我們的兒女們.
兒孫們.
都不要離開你的話語.
讓他們能夠自由.
都能夠認識你.
一生不偏離你的路.
求主讓我們在這個混亂的世代.
因著你的說話.
我們知道你的心意.
我們行在其中.
我們就能夠真正去經歷.
神你自己的保守.
和給予的平安.
我們也都將我們每一個.
神的兒女交在你手中.
求主讓我們成為.

$^{961}$這個福音基督的使者.
能夠將你的救恩帶進去.
我們的朋友圈裡面.
讓每一個未知未信.
包括我們的家人.
朋友.
同事.
同學.
都能夠認識你.
我們同心合意的祈禱.
奉主耶穌基督聖命給我們.
阿們.
\newpage



\section{箴言 21:31}
\label{sec:031631}
\textbf{馬年話馬}~~~~梁永善牧師
\newline
\newline
連結: \href{http://www.church.com.hk/acms/content.asp?site=cdc\&op=show\&type=product\&code=031631\&layout=sermon}{\texttt{031631}} ~~~~ 語音日期: 2026-02-11
\newline
\newline
\hyperref[sec:031630]{\small{< < < PREV SERMON < < <}}
~
\hyperref[sec:index_chronic]{\small{[返順時目]}}
~
\hyperref[sec:031632]{\small{> > > NEXT SERMON > > >}}
\newline
\newline
箴言 21:31
\newline
\begin{longtable}{cl}
\hline
\hline
章節 & 經文 (和合本修訂版)\\
\hline
21:31 & \begin{tabularx}{0.7\textwidth}{X} 馬是為打仗之日預備的；得勝卻在於耶和華。 \end{tabularx} \\ \\
[1ex]
\hline
\hline
\end{longtable}
$^{1}$從稅守到年終.
神都來帶我們去眷顧,幫助我們.
今天會和大家一起來分享神的話語.
我們會在聖經裡說關於「馬」這件事.
「馬年」有很多關於馬的教導在聖經裡.
第一是提醒我們.
千萬不要只靠馬.
是甚麼意思呢?.
我們先來讀一段聖經.
這段聖經讓我們帶大家看到.
究竟是甚麼一回事.
我們請二汶替我們讀出來.
「馬是為打仗之日預備的.
德勝乃在乎耶和華」.
還有飾演二十篇的第七節.
「有人靠車,有人靠馬.
但我們要提到耶和華我們神的名」.
今天我們不得不承認.
在我們的生活上.
我們會遇到一些不同的問題.
一些困難.
甚至有些重要的抉擇.
甚至有時會陷入一個危機之中.
當我們面對這些困難危機時.
我們會怎樣做呢?.
我們很快就會去想.
會不會找些人來幫我們呢?.
找人幫忙不是錯的.
你有病去看醫生.
你買賣樓宇要找律師.
當你來到英國的時候.
可能你很多地方都不熟悉.
或者法例都不熟悉.
你要靠一些人給你一些提醒和指示.
這是人之常情.
有句成語叫「以馬可代」.
意思是我們倚靠那匹馬.
然後給我們有力量.
因為馬代表了力量.
現在的汽車.

$^{41}$有幾匹馬力.
馬代表力量很大.
古時候有馬車是力量的象徵.
我們遲些會看《士師記》.
這幾天我已經開始看《士師記》的《釋經書》.
特別看到當日在迦南地.
有些外敵有馬,鐵車.
有馬車可以雄霸當時的天下.
而所羅門國勢最盛的時候.
就是用馬來做代表.
我們看《列王記》上第十章二十六節.
「所羅門聚集戰車,馬兵.
有戰車一千四百輛.
馬兵一萬二千名.
安置在屯車的城邑.
和耶路撒冷就是王那裡」.
我們去以色列的時候.
都說我們祈禱了.
四月五日我們去以色列.
那時候我們一定會參觀一個地方.
叫做米吉多.
米吉多是甚麼呢?.
最重要的是當日所羅門.
在這個地方養了很多馬.
有很多馬騎的遺跡.
如果你參觀米吉多的時候.
還有一些佈置.
用馬兵作為標誌.
證明那個地方是養馬的地方.
所羅門國勢最強盛的時候.
有多少戰車和馬呢?.
但如果我們只是將信心,前途.
放在一些人身上.
而將這個焦點放在這些人身上的時候.
而輕忽了神的時候.
就很容易出錯了.
我舉一個例子.
大衛被索羅追殺.
他想投靠撒布爾.
但撒布爾已經無能為力了.

$^{81}$大衛無處可逃.
竟然做甚麼?.
做一些完全不合理的事.
他逃亡到非利士.
不要忘記.
他當日是殺死非利士人的將領.
特別是哥利亞.
他投靠非利士的加特王阿吉.
阿吉又肯接納他.
但最後大衛被困在困境裡.
後來非利士幾個領袖.
預備帶兵去殺希伯來人.
即是大衛的同胞.
而大衛投靠了加特王阿吉.
他也要去打這場仗.
如果他打這場仗的話.
他殺了自己的同胞.
他怎可以來大衛做自己的人民的王呢?.
但神有憐憫.
於是大衛就幫助了他.
不用打這場仗.
大家可以一路上我們的《周三茶經》.
以前我們會讀大衛的.
有一個茶經叫《易戰》.
即是半翼的翼翼戰.
解釋神如何憐憫.
猶大王受著以色列和亞蘭聯合攻擊.
猶大國的王怎麼辦呢?.
阿哈斯向亞述來求助.
結果不會有免費餐食.
要將聖殿和國庫裡的金銀珠寶.
送給亞述王才可以解燃眉之急.
我們看看《列王記》下十六章第七,第八兩節.
你想想阿哈斯要求王.
不是自己所相信天下主宰的王.
是亞述王.
結果要將皇宮和府庫裡的金銀財寶.
都送出去.
最後末代皇帝荷西亞.
首先降服亞述.

$^{121}$但左搖右擺.
後來看到埃及強大的時候.
於是他就偏向埃及.
不再朝貢給亞述王.
結果被亞述王捉了.
甚至有一個很慘的下場.
第四節我們會看到.
我們看看《列王記》下十七章第三,第四兩節.
亞述王撒曼以色上來攻擊荷西亞.
荷西亞就服侍他.
給他縱貢.
荷西亞背叛.
差人去見埃及王所.
不照往年所行的.
與亞述王進貢.
亞述王知道了.
就把他鎖禁.
困囚在監獄.
所以你看到他左搖右擺.
他只是投靠有勢力的人.
而不是倚靠神.
結果就失敗了.
所以你看到.
其實我們不是要找人幫忙.
但我們的焦點是放在神的裡面.
其實先知一直都提醒他們.
千萬不要去投靠有勢力的人.
當時埃及曾經是很強的.
先知二賽說不要投靠他們.
有時候甚麼都不要做.
等候神的工作.
我們基督徒很熟悉的一句說話.
甚至我們抽了這句說話來鼓勵人.
「你們得救在乎歸回安息.
你們得力在乎平靜安穩」.
但我們抽了這句.
接著還有幾個字我們沒有用到.
我們看看這一段聖經.
二賽一書31章第一和第三節.
「和哉!那些下埃及求幫助的.

$^{161}$是因障礙麻匹倚靠甚多的車輛.
並倚靠強壯的馬兵.
卻不仰望以色列的聖者.
也不求問耶和華.
埃及人不過是人並不是神.
他們的馬不過是血肉並不是靈.
耶和華一伸手那幫助人的必半跌.
那受幫助的也必跌倒都一同滅亡」.
這裡很清楚說到.
你去埃及求幫助.
他們勢力很強.
你倚賴他們的馬匹.
那種力量很大.
甚至他們有很多戰車.
但其實沒有用的.
你不去倚賴你所相信的神.
你只會招來禍害.
埃及人不過是人.
也不是神 他們的馬不過是血肉.
所以很糟糕.
他們怎樣強?.
神一伸手甚麼都沒有.
當年出埃及都是這樣.
你不要忘記這是31章.
第30章就是我們很熟悉的.
好 我們一起來看.
30章15,16兩節.
主耶和華以色列的聖者.
曾如此說:.
「你們得救在乎歸回安息.
你們得力在乎平靜安穩.
你們竟自不肯.
你們卻說不然.
我們要騎馬奔走.
所以你們必然奔走.
又說我們要騎飛快的牲口.
所以追趕你們的也必飛快」.
我們經常去送一些經文鼓勵人.
抽取我們認為喜歡的.
看到這裡說到.

$^{201}$「你們得救在乎歸回安息.
你們得力在乎平靜安穩」.
但下面那句是甚麼?.
「你們竟自不肯」.
所謂恩賜你甚麼都不用做.
不如你等一等.
看看我怎樣去工作.
但有時我們困境來到.
我們很徬徨.
我們缺乏這種等候.
甚至我們會怎樣呢?.
我們要去投靠那些有能力的人.
我們要騎馬奔走.
但我說沒有用的.
接著就是我們剛才所讀.
第31章裡面.
埃及王只是人.
馬匹都是血肉之軀.
是沒有用的.
所以求神來幫助我們.
我再說.
當我們在困難.
當我們有需要.
我們不是不需要人幫助.
甚至這一年開始.
多些人會申請「五加一」.
他們要看很多不同人的網站.
教怎樣去申請.
我們不是不看這些.
我們不是不準備.
但我們最大的準備是甚麼?.
是要仰望我們的神.
你記住神都可能透過不同的時段.
透過不同的人.
給我們一些幫助.
但我們要搞清楚.
究竟是神安排這些人.
還是我們去自找呢?.
所以記住.
我們絕對不要將信心.

$^{241}$只放在某些人身上.
我們要依靠的是我們的神.
在這一點裡的總結.
就是詩篇33篇16到17節.
我們一起來看好不好?.
君王不能因兵多得勝.
勇士不能因力大得久.
靠馬得救是枉然的.
馬也不能因力大救人.
記住.
今天我們不只是靠世界的力量.
是靠著我們的神.
好了.
我們提醒.
絕對不要只是靠馬.
另外一句成語大家很熟悉.
就是「伯樂識馬」.
我們常說伯樂很厲害.
懂得選馬.
但其實我們有少許錯誤.
其實不是伯樂.
其實這裡是出自戰國冊.
楚冊裡面的第四章.
它說到「世有伯樂然後才有千里馬」.
春秋後期.
即時期裡面有一個人.
記住.
這個人不是叫伯樂.
這個人叫甚麼呢?.
叫孫陽.
姓孫.
孫公冊的孫.
即兒孫的孫.
太陽的陽.
他有鑒別千里馬的能力.
每逢給他鑒賞.
然後去讚賞.
認定這些馬.
都升價百倍.
而且這些馬真的有實力.

$^{281}$他看這些馬的質素是很重要.
有一天他經過一個市集.
發現有人在賣馬.
這些馬是很死瘦的.
但是原來賣了三天.
都沒有人去摳這隻馬.
孫陽看到這隻馬是非常好的.
原來有一句說話是怎樣說的呢?.
「上馬失之瘦,上士失之貧」.
你認為這隻馬很瘦而沒有用.
你只是看輕了它.
結果這匹馬被孫陽看中了.
實際上真是一隻千里馬.
後來有些人用了他們的神話故事.
在天上掌管馬的神叫做「伯樂」.
「伯樂上馬」.
「伯樂」其實是指孫陽.
後來唐朝的韓語在雜說中說.
西有伯樂然後有千里馬.
千里馬常有.
但是伯樂不常有.
是什麼意思呢?.
很多時候都有很多千里馬.
但是沒有伯樂.
沒有人懂得發掘它的才幹.
結果就很糟糕.
所以伯樂就懂得怎樣來上馬.
不要忘記原來我們的神是我們的伯樂.
神看重我們.
神怎樣看我們呢?.
在《雅歌書》第二章第二節.
「我的佳偶在女子中.
寶像百合花在荊棘內」.
神看我們的比喻是甚麼呢?.
荊棘是沒有用的.
是刺手的.
你不會用荊棘來做裝飾.
但是在荊棘裡竟然有百合花.
是多麼的美麗.
意思是神看我們美麗的一面.

$^{321}$不單看我們不好的一面.
於是神在某個程度是獨具慧眼.
看到我們的美好的質素.
於是來帶我們去選擇.
所以很重要.
所以今天我們要怎樣呢?.
我們要來帶我們去看.
今天往往很多人.
去選擇一個人的時候.
看他的外貌,看他的學識.
看他的社會地位.
但是神不是這樣看的.
《新福音經》第十六章第七節.
我很喜歡這段經文.
我們一起來看.
「耶和華不像人看人.
人是看外貌.
耶和華是看內心」.
稍後會再解釋這段經文的背景.
如果你有查《悉讀大衛》.
我第一次講《大衛》的時候.
題目是「金玉其內」.
不是其外.
是裡面那種質素.
如果你留意到神揀選一些人.
如果真的用人的角度來看.
這些人其實是不值得去選擇的.
雅各是以色列之父.
但是不要忘記雅各是什麼?.
一個很狡猾的人.
只是來釋放一個機會.
經常說謊.
所以一個這麼狡猾的人.
大神竟然讓他成為什麼呢?.
成為以色列的國父.
還有我們稍後在《事實記者》要看.
耶伐他.
耶伐他是什麼人?.
原來是妓女的兒子.
野種.

$^{361}$沒有社會地位.
都不知道生父是誰.
人們歧視他.
但是神揀選他成為事司.
去拯救以色列人.
還有就是我們很熟悉的大衛.
剛才我們讀了一段經文.
原來撒母耳這個一國最高的宗教領袖.
去耶西的家.
即是大衛的父親的家.
他想揀選一個接班人.
接續掃羅.
耶西馬上叫他兒子出來.
大兒子叫做以利亞.
一出來高大衛猛.
很厲害.
想著一定是他.
誰不知神就說.
不是這個.
就是剛才我們所讀的經文.
耶華不像人看人.
人看外貌神看內心.
其他的兒子一個一個出現在撒母耳眼前.
但是神都說不是.
結果撒母耳問耶西的丈母子.
啊!還有一個在看羊.
你看看.
如果大衛在父親心目中是很重要.
撒母耳來的時候.
大衛還在放羊而沒有叫他回來.
所以大衛其實在父親眼中.
並不是一個值得他愛惜的孩子.
我們看看新約的彼得.
彼得說話很衝動.
甚至三次不認主.
但是彼得卻成為教會的基石.
保羅.
我們知道他是一個很偉大的仕途.
但是不要忘記保羅.
他未信主之前.

$^{401}$殺了很多基督徒.
一個殺了那麼多基督徒的人.
竟然可以被神彌大人選為外邦人.
這是我們想像不到的.
可能保羅提攜一個年輕人叫提摩泰.
為什麼會選提摩泰呢?.
對保羅沒有好處.
因為提摩泰是一個混血兒.
不要忘記當時猶太人很不滿意那些混血兒.
特別是有混了猶太的血種.
他媽媽是希伯來人.
但爸爸是希臘人.
提摩泰早年喪父.
在媽媽和外婆的養育下.
很沒有男兒氣概.
很沒有信心.
但卻被神選擇.
所以今天不要忘記.
原來神選擇我們.
絕對不只是看外表.
正如聖經所說.
哥林多前書第1章26至29節和第31節.
「弟兄們!可見你們蒙召的按著肉體有智慧的不多.
有能力的不多.
有尊貴的也不多.
神卻揀選了世上愚絕的.
要有智慧的羞愧.
又揀選了世上軟弱的.
要那強壯的羞愧.
神也揀選了世上卑賤的.
被人厭惡.
以及那無有的.
為要廢掉那有的.
使一切有血氣的.
在神面前一個也不能自誇.
如經上所記.
誇厚的當只著主誇厚」(宗15:1-2).
原來那時的基督徒被人看上來.
沒用了.
不行了.

$^{441}$才去倚靠神.
很多人覺得我們是愚絕.
但其實我們是聰明的.
如果你有聽我上兩集說的.
「一帶一路」.
你帶那個扣安全帶.
你會很生氣的.
那個鬧劇.
我就說到.
你又明白信耶穌實在不容易.
因為你要認錯.
到今天運輸署署長陳美寶.
都不肯去承認自己.
即是看漏眼.
即是有錯處.
所以你留意到.
原來我們有勇氣.
然後先者會信主.
但在世人看上來.
我們基督徒是無能軟弱的.
但神祂是重視我們.
我記得我曾經去過伊朗.
不知道今年年底能否去得成.
伊朗有世界上最大的寶石.
在一間伊朗銀行裡我們可以去參觀.
我看到一幅畫給我很深感受.
是一塊鑽石.
切開了一半.
裡面是紫色的.
即是很漂亮很通透的寶石.
但原來包著的是一塊石頭.
原來是有些人真的獨具慧眼.
才留意到這塊石頭裡面.
是有寶藏在當中.
是要敲磨的.
今日有時別人看我們.
都覺得我們是一塊爛石.
以我自己為例.
很真實地說.
當年我讀神學的時候.

$^{481}$每個人都看不好.
團友都說:不是吧?.
你讀神學?你真是假的?.
他們在質疑我.
就算我讀神學.
有些老師,有些同學.
都不會覺得.
我是怎樣做到一個好的牧者.
當然我不是說現在是一個好牧者.
但至少我今年踏入.
已經是第五十一年的事房.
即是神都憐憫我.
當然不是靠我自己.
有很多人掉眼鏡.
有些被視為屬靈的同學.
但做了幾年就已經退下來.
所以今日我們不要只看我們的外表.
還有我們是被人來賞識.
或者被神來賞識.
於是我們可以為主來做事.
但同樣我們都要成為伯樂.
去找一些接班人.
要揀選一些人.
這是很重要的.
我們怎樣來帶他成為伯樂呢?.
就是我們會不會來帶他.
不要這麼快去跟那個人定論.
這個人是不行的.
千萬不要.
有時我們要給多些機會人.
正如保羅殺了很多基督徒.
但在大馬適路上.
他被主呼召.
他信了主.
他很想去見仕途.
但仕途又比這個反應更對.
仕途說不想見他.
這個殺人兇手.
他來見我.
他說信了.

$^{521}$誰知道是真是假.
會不會做無間道.
到時來抓我們.
沒有人肯去見保羅.
如果保羅多次被拒絕.
他會灰心失望.
他會再成為一個仕途嗎?.
不會.
但有一個人很寶貴.
他看不到之前他跟保羅有什麼交往.
但他很單純地去信保羅.
於是他就成為保羅的伯樂.
好,我們一起看《賜旺傳》第九章.
二十六,二十七兩節.
「蘇羅到了耶路撒冷.
想與門徒結交.
他們卻都怕他.
不信他是門徒.
唯有巴拿巴接待他」.
看不看到?.
仕途不肯相信保羅.
不肯見他.
但聖經說唯有巴拿巴接待他.
巴拿巴相信他.
後來巴拿巴帶著保羅去見仕途.
因為巴拿巴在當時的教會非常有生命力.
有好的見證.
所以聖經裡很清楚說.
是人們信巴拿巴而不是信保羅.
由巴拿巴去引薦保羅.
所以巴拿巴某個程度是保羅的伯樂.
所以今天我們都需要成為伯樂.
我們需要去找一些接班人,使徒.
我們見到人有潛質.
給多些機會他去服侍.
訓練他,提醒他.
教會才會有好的發展.
絕對不是將注意力集中在一,兩人身上.
但今天我們見到社會不是看這些.
是會看你的社會地位.

$^{561}$看你最重要是否忠誠.
是否愛國.
於是你便可以被揀選.
但聰明的人可能被否定.
所以今天神揀選我們的時候.
絕對不是只看我們的外表.
而是看我們最重要是.
裡面的質素是否真的愛我們的神.
好了 釋圖魯瑪.
釋圖魯瑪的背景是甚麼呢?.
我去講道就不會講這些了.
大家都聽過管仲.
和齊桓公一起去北伐.
誰知春去冬來.
他們入錯了一個竹林.
沒有甚麼熟悉的嚮導帶他們去.
在回家的路上他們迷途.
管仲說魯瑪的智慧是可以用的.
於是他就說我們不要騎馬.
我們下馬.
讓那隻馬自己行.
因為他是魯瑪.
他一定會認識那條路途.
我們反而跟著那隻馬去行.
結果真的成功了.
於是他就說到.
我們要承認我們有不足的地方.
這就叫做釋圖魯瑪.
即是說.
應用在今天的宗教.
有經驗的人要帶著一些新人.
剛才我們說我們都要做伯樂.
我們就要做釋圖魯瑪.
好 我們留意到.
聖經都是這樣教導我們.
《提摩太后書》第二章第二節.
「你在許多見證人面前聽見我所教訓的.
也要交託那忠心能教導別人的人」.
這是保羅寫信給提摩太太.
保羅提她提摩太太.

$^{601}$但接著他就教她提摩太太.
你都要找一些接班人的.
接班人是甚麼呢?.
是忠心於神.
而且能夠教導別人的.
所以領袖應該是薪火相傳.
缺乏自信的領袖.
很驚見後起之秀.
受人歡迎.
強過我.
於是就會打壓.
迫他走.
這是很可悲的.
甚至中國的一些所謂武術的前輩.
收徒弟.
他們很流行收起一,兩招不教.
因為很害怕.
教會徒弟沒有師父.
但我們作為牧者.
我們要廣闊的胸襟.
要提拔有潛質的人.
我們不用怕別人搶我們的肛.
你繼續努力.
神自然會給你肯定.
所以保羅全力輔助提摩太太.
給她訓練,鼓勵,教導,栽培.
並且給她有事奉的機會.
使她從事奉中學習.
這是我們應該效法的.
請你記住經驗是要慢慢累積的.
經驗這個字其實是很痛苦的.
即是你經過很多碰碰跌跌,失敗.
流過不少血和汗.
甚至淚水.
但這些經驗.
你要好好地教導你的後輩.
作為指引.
使他們不會再撞這樣的板.
這在神的角度是很重要的.
但你們也好一點.

$^{641}$現在到了我這個年齡.
我可以去跟從的前輩.
已經是越來越少了.
即是今日就算教會事奉的人.
當你退下了你的執事會.
或者領導階層.
你休息是需要的.
但你也要將這些經驗.
跟後輩分享,賦益他們.
這樣才是好事.
不單是事奉的經驗.
甚至有時我們會失敗.
當我們失敗.
等神來挽回我們.
我們再能夠重新站起來時.
我們需要怎樣做呢?.
就是將這些曾經失敗的經驗.
去指教後輩.
免得他們犯錯.
好,我們來看看《盧九福音》第22章第32節.
「但我已經為你祈求.
叫你不至於失去信心.
你回頭以後要兼顧你的弟兄」.
這裡很清楚地說.
彼得臣永台說.
你會三次不認我.
但你會重新站起來.
當你重新站起來時.
你就要好好地兼顧你的弟兄.
所以今天我們如果曾經軟弱失敗.
但是有神的接納.
我們重新站起來時.
我們就需要幫助別人.
正如大衛犯了姦淫的罪.
你有沒有留意到?.
先知拿丹指責他.
他認罪後.
他寫了懺悔詩.
這個懺悔詩是全國的.
來教導他們.

$^{681}$甚至我要指教你們.
不要再犯這個錯.
將自己的失敗提醒給別人.
所以我們要悉途勞碌.
就是勸導別人.
同樣地不一定是犯錯.
你經歷一些患難.
神安慰我們.
於是我們經歷這些安慰.
就可以安慰其他人.
我們看看哥林多後書第一章第四到第六節.
「我們在一切患難中.
他就安慰我們.
叫我們能用神所賜的安慰.
去安慰那遭過各樣患難的人.
我們既多受基督的輔助.
就靠基督多得安慰.
我們受患難.
少為叫你們得安慰.
得拯救.
我們得安慰.
也少為叫你們得安慰.
這安慰能叫你們忍受.
我們所受的那樣的苦楚」.
即是我們經歷這個苦楚.
就可以安慰同樣有需要的人.
很簡單.
黃福縣大火我們教會都有十多個家庭.
我所認識的就已經毀於這個火裡面.
不幸地我大埔的女傳道.
和上水的男傳道他們是夫婦的.
他們都是住在黃福縣.
原來他們剛剛借裝修完.
裝修完之後還去台灣旅行.
回到台灣之後他們就沒有了家園.
我也有去關心安慰.
後來我也知道教會的老人家.
都失去了他們的家.
我聯絡了那位女傳道.
我問她可不可以去關心這些老人家.

$^{721}$她現在肯回教會了.
我說你經歷了這個苦楚.
希望可以幫到他們.
我要接受的.
她就說不行了牧師.
我現在自己都還未行.
所以我不可以去跟她談.
我只可以祈禱求神幫助她.
有更強的承載力.
希望她復元之後可以幫助那些人.
因為大家有相同的經歷.
這個時候分享會更加貼身.
所以求神幫助我們.
記不記得我們剛才說.
我們是要被神揀選.
這是神體寵我們.
很感恩.
但我們都要做伯樂.
我們懂得去揀選人.
我們是老馬識途.
我們有經驗.
我們要指教後輩.
但同樣後輩都需要接受教導.
這一點比我們更重要.
就是有一個人.
不詳細說了.
當年他因為吸毒.
被捕後來不坐牢就去了戶外團.
就順得住了.
後來他回考不成功.
在大一讀設計.
誰知道他讀了設計之後.
竟然英國有大學覺得他的作品不錯.
讓他讀碩士.
他說他英文不行.
那時候不流行.
現在國內很流行找人替槍考試.
但他很努力地去補習.
有弟兄姊妹去教他.
終於可以來到英國.

$^{761}$甚至原來每一年都會選一些.
海外傑出的留學生去獲獎.
他回到香港之後.
他做什麼呢?.
他開設計公司.
我和他做過很多合作.
他講見證我講道.
甚至現在的設計.
入了見歷史世界大全的最大壁畫.
特別是幫那些戒毒吸毒者.
下面有一個我們冥恩堂都有請他.
叫張增森.
當年他是吸毒的.
後來戒了毒.
我和他一起在監獄裡做工作.
是四十多年前的事了.
後來他讀神學的時候.
他快要畢業.
我記得我在旺角一間茶餐廳.
跟他約了一會兒.
我很誠意邀請他去我冥恩堂事奉.
但後來他說.
「不要了牧師,我多謝你」.
他想去戒毒的地方.
因為他有這個經歷.
可以幫到一些人.
所以很寶貴.
當然我們去教人.
但我剛才已經講了.
我們要「管眾隨馬」.
是什麼意思呢?.
就是我們要有謙虛的心.
去接受別人給我們的教導.
有時候很擔心年輕人.
他們覺得自己很行.
但我們就要放下自己.
因為我剛才講的那件事.
就是「管眾」說.
我們不要再騎那隻馬了.
我們下來跟著那隻馬去走.

$^{801}$這樣就會找到出路了.
所以求神來幫助我們.
不要忘記當日羅波安.
所羅門的兒子登基.
那些人就問.
可不可以讓他們休息呢?.
那些工人要求.
我們做得很辛苦.
他問那些老人家.
老人家說「好啊,你讓他們休息吧」.
那些人就會覺得你愛他們.
就會來帶他們去跟隨你.
但他竟然不聽老人家講.
就聽同輩的人講.
「當然不行啦,你現在做王.
他就要這麼多要求.
這樣就行了.
再繼續下去.
他就要更多要求」.
結果他就來罵這班人.
指責他們.
甚至誇大他們說.
「我的母子比我爸爸的腰更大」.
「我爸爸找哪邊來?」.
「找傑氏哪邊來就哪邊來」.
結果這班人就知道.
這個王不值得跟隨.
結果國家就分裂了.
所以今天年輕人都需要怎樣呢?.
是要有謙卑的心.
要有管眾除馬.
要有聽人的忠告.
好了.
一個我自己很渴望的就是馬格果斯.
就是指到一些人很英勇.
戰死沙場.
其後他的遺體.
就被死去的馬的皮包著.
然後送回家鄉裡.
所以這個就叫做馬格果斯.

$^{841}$是表示一個人忠心至死.
我自己期望一天我離開世界的時候.
在我的追思會裡.
或者安息禮.
都引用這一段經文.
好 我們一齊帶大家去看一看.
臺灣太后書第四章第六到第八節.
我現在被囂電我離世的時候到了.
那美好的仗我已經打過了.
當跑的路我已經跑盡了.
所信的道我已經守住了.
從此以後有公義的公民為我存留.
就是按照公義審判的主.
到了那日要賜給我的.
不但賜給我.
也賜給凡愛無他顯現的人.
保羅他很清楚的宣告.
提摩太后書是他的絕筆信.
後來沒多久他就離開世界.
但他可以傲然的說.
那美好的仗我已經打過了.
當跑的路我已經跑盡了.
所信的道我已經持守了.
所以今天很寶貴.
我自己很深的去體會.
我希望我到離世那一天.
我仍然能夠為主擺上.
有一句成語叫做「老馬為驅」.
意思是甚麼呢?.
「練老,但仍然有壯志」.
剛才我所說的.
是越來越少我可以跟隨的長者.
在我自己的教會裡.
滕錦輝牧師是我教會的牧者.
很溫文爾雅.
我今天中午還在聽.
滕牧師的講道.
想不到還有滕牧師.
不知是誰可以放上網.
是幾天前放上網的.

$^{881}$講希伯來書.
他有十講.
我今天聽了第一講.
去年我自己一個遺憾.
就是我所尊敬的曾立華牧師.
也離開了世界.
曾牧師在我十三歲回教會的時候.
那時候他懂得讀神學.
他是神學生.
來我們團體裡面做實習工作.
是我少年團體的導師.
後來我有幸繼續跟隨他.
在神學院裡面他也教我.
也是我一位尊敬的長者.
但是他亦在去年離開了世界.
所以我自己也有很深的感受.
可以跟隨的人是越來越少.
但是正如保羅他自己這樣說.
他說「我活著就是基督.
我死了就是有益處」.
那時候他是願意為主有墳進.
所以死亡並不可怕.
我有兩個人是很影響我的.
一位是我屬靈的父親.
優育靈牧師.
剛才我在網上找到優牧師的生平.
是以前基督教周報裡面的一個刊物.
另一個就是楊牧公牧師.
楊牧公牧師就是我會透過他的生命.
他的對聖經的執著和解經的重要.
而怎樣應用在生活裡面.
我是效法他的.
楊牧師特別在他病了之後.
他有很大的改變.
最近我在網上看到很寶貴的.
就是70歲的孫保齡牧師寫了一大段東西.
說基督教裡面有些牧者是令你搖頭嘆息的.
有些是令你讚嘆.
有些是令你嘆息.
他特別提到楊牧公牧師.

$^{921}$大家可以上網試試找來看.
楊牧師有癌症之後很特別.
有人付錢去私家醫院診治.
但他寧願去威爾斯尼達去排隊.
他接觸了很多平民.
這一次癌症是令他改變了寫作的方向.
是很隱瞞的.
即是說隨訪哲學,讀書謝華儀.
再少寫了一些聖經書.
所以對我來說是一個損失.
但他有更多人會得到幫助.
他當年生病的時候.
他在時代論壇每個星期都寫一篇病歷經歷.
後來緝錄成書.
大家如果有機會可以看這本書.
叫《再生情緣》.
他講到甚麼是死亡呢?.
他說死亡永遠都是引進一道新時代的門檻.
因此不需要為了這件事而害怕.
我想用最後的時刻和最後的經驗作見證.
不用怕.
更加犯不著因為害怕死亡.
一生做他的奴隸.
要挺起胸膛做人.
也要敢於愛人和愛神.
這樣你的生命就保證豐盛.
他說死亡不是問題的所在.
生命才是.
因為死亡只是過渡.
生命卻是延續.
把握現在改變將來是一種負責任的生活.
掌握著生命並且將自己最好的擺上.
這個生命就是不可能一事無成.
是否很寶貴?.
所以我都很感恩.
我要向他學習.
所以我自己一直以來.
我不會去尼德雅懶惰.
正如今晚我已經說了.
我兩個孫女生日.

$^{961}$今晚結束後我會馬上去尼德雅慶祝生日.
我都是重視事奉.
我希望自己每一天都用得好.
所以熟悉我都知道.
我是甚麼呢?.
有一句名句叫「抵死有餘」.
抵死不是錯的.
是讚賞.
就是用得好的時候.
就能夠怎樣呢?.
就會值得了.
我曾經和太太說.
如果我死了就不如放在墳墓裡.
因為太太不肯.
我很感恩.
一位會友上去北京的時候.
他去讀書就幫我找人寫了這四個字.
送了給我.
有個牌匾我帶了來英國的.
就是「抵死有餘」.
所以求神來幫助我們.
希望我們活得好.
事實上我見過有些人很糟糕.
聖經不斷教導我們.
但我們都要留意一件事.
我們絕對不可以輕率.
不可以馬放南山.
馬放南山是甚麼意思呢?.
就是敷衍地來做.
讓他自生自滅.
這就是很糟糕的.
我見過有些人剛出道的時候很熱心.
但慢慢地就越做越懶.
傳道人,牧師都有這些犯錯.
所以我們也很深體會到.
「路遙知馬力,日久見人心」.
更加遑論馬格果斯.
所以如果輕率地去做的時候.
這是教會的悲哀.
所以求神幫助我們.

$^{1001}$讓我們真的能夠有所堅持.
在馬年的時候.
也會跟大家祝福.
祝福教會和你我能夠學習到.
在我們生命裡面記住.
絕對不會有偶然的事發生.
可能有些不是很開心的事.
但我們會失業.
或者我們來到英國找不到工作.
但你怎知道這不是塞翁失馬.
焉知非福呢?.
特別你看一下.
弱勢去坐牢被出賣.
但原來最終可能是一個祝福.
求神幫助我們.
特別欣賞我們物質好像失去色彩的風候.
但是我們應該學到.
就是「馬死落地行」.
是甚麼意思呢?.
就是接受這個現實.
可能我過不了豐富的生活.
我就過一個破車劣馬的生活.
即是我們較為簡樸的生活.
正如保羅怎樣教導我們呢?.
《肥下比書》第四章第十二節.
「我知道怎樣處悲戰 也知道怎樣處豐富.
或飽足或飢餓 或有餘或缺乏.
隨時隨災 我得了被決」.
這裡很清楚讓我們看到.
今天我們知道怎樣處豐富.
知道怎樣處悲戰.
不要忘記保羅說這段話的時候.
他是在坐牢的.
他之前是非常有地位.
但他信了主之後.
中了罪 就來到失敗了.
他要坐牢.
在人看來是失敗.
但在神的眼中他是寶貴的.
所以都求神來到幫助我們.

$^{1041}$在這些處境當中.
有時候我們不得不承認.
就是社會有很多東西引誘我們.
我們面對社會這麼多的誘惑性.
我們會心懸意馬的.
不知為何這樣說.
結果可能我們被世俗扯走.
所以在聖靈的提醒下.
我們要懸崖勒馬.
千萬不要來到犯錯.
因為如果我們來到犯錯的時候.
我們就會成為害群之馬.
其實害群之馬.
我們會說這匹馬是害死人.
但原來真正的害群之馬.
原本不是這個意思.
就是有一個王叫迷路.
見到一個年輕人在放馬.
在問路.
人就不太理會他.
但後來聊天的時候.
他很有智慧.
他就問他.
他說到我怎樣可以治理好國家呢?.
那個人就跟他說.
很簡單.
你養馬和你治理國家一樣.
我怎樣來養馬呢?.
如果有一些東西.
令到那匹馬吃了之後.
這匹馬其實是一隻駱駝.
是我送給我的孫女.
去埃及的時候.
是藥彈.
他說到如果那些東西.
令到那匹馬吃了之後.
會有病的.
我就會將這些東西除掉.
免得害了我的馬.
他說到如果有哪些東西.

$^{1081}$在影響你國家的發展.
你就將它除掉.
後來就演變成為害群之馬.
就好像一匹馬害了其他人.
所以我們都要很小心.
但聖經也提醒我們甚麼呢?.
我們要很小心.
如果我們不如危勒馬的時候.
我們犯錯就很大件事了.
拜託第五章二十九到三十節.
若是你的右眼叫你跌倒.
就挖出來丟掉.
寧可失去八體中的一體.
不叫全身丟在地獄裡.
若是右手叫你跌倒.
就撼下來丟掉.
寧可失去八體中的一體.
不叫全身下入地獄.
所以今天我們要很小心.
免得我們有犯錯.
因為最後環顧周圍.
有很多未信的人.
福音的和尚很大.
所以我們大家要快馬加鞭.
和撒旦競取.
因為撒旦很勤力的.
我們要馬不停蹄地去傳福音.
要為神立下這個汗馬功勞.
但記住.
今天神把我們在群體裡面.
我們不是單槍匹馬去面對這些事.
不是自己請英雄.
教會是一個群體.
教會往往也會招兵買馬.
如果我們願意投入教會事奉的時候.
我們就能夠兵強馬壯.
然後勵兵抹馬.
為神的國度而戰.
最重要是怎樣呢?.
我們要馬首是瞻.

$^{1121}$我們以神為我們的標準.
我們跟隨我們的神.
倚靠祂,仰望祂.
這樣我們就會馬到功成.
最後.
祝福教會.
居水馬龍.
這樣的時候.
即是說人越來越多.
有更多人歸向主.
所以求神幫助我們.
我們一起來祈禱.
願意主你的恩典.
來與我們同在.
你從歲首到年終都眷顧我們.
農曆新年快要到.
一個新的開始.
我們特別在這個星期.
特別是在英國的弟兄姊妹.
他們有休假期.
都求主你幫助到做父母的.
可以與孩子們有好的交往.
亦都教導他們.
中國不單單是一個文化.
更加是聖地牙利大學.
去從聖經教導他們.
你幫助我們.
讓我們知道神你揀選我們.
不是我們有什麼好處.
但你看到我們裡面.
還有一些美麗的質素.
你將這個質素來擴大.
所以你揀選我們.
也叫我們能夠聖地牙利大學.
去成為一個賞識別人.
更加去啟迪別人的人.
願意以神你自己幫助我們.
更加叫我們能夠抵抗.
殺人的誘惑.
叫我們回心圓意.

$^{1161}$幫助我們忠心地服侍你.
以致有一天我們面對你的時候.
我們可以無愧於見到你.
求主你自己親自眷顧.
聽下在面前的禱告交託.
乃奉主耶穌基督得勝的名禱告.
阿們.
\newpage



\section{羅馬書 12:2}
\label{sec:031632}
\textbf{如何察驗神的旨意}~~~~鄭兆奇牧師
\newline
\newline
連結: \href{http://www.church.com.hk/acms/content.asp?site=cdc\&op=show\&type=product\&code=031632\&layout=sermon}{\texttt{031632}} ~~~~ 語音日期: 2026-02-08
\newline
\newline
\hyperref[sec:031631]{\small{< < < PREV SERMON < < <}}
~
\hyperref[sec:index_chronic]{\small{[返順時目]}}
~
\hyperref[sec:031633]{\small{> > > NEXT SERMON > > >}}
\newline
\newline
羅馬書 12:2
\newline
\begin{longtable}{cl}
\hline
\hline
章節 & 經文 (和合本修訂版)\\
\hline
12:2 & \begin{tabularx}{0.7\textwidth}{X} 不要效法這個世界，只要心意更新而變化，叫你們察驗何為神的善良、純全、可喜悅的旨意。 \end{tabularx} \\ \\
[1ex]
\hline
\hline
\end{longtable}
$^{1}$今日來到是崇拜上帝.
大家的心都需要歸向上帝.
剛才詩歌都提醒我們.
無論何時上帝都會陪伴著我們.
而且在我們生命當中我們要學習.
有一個很重要的功課就是學習.
怎樣將我們奉獻自己.
聖經使徒保羅都提醒我們.
我們要將身體獻上當作活祭.
將身體獻上當作活祭之前.
是有一個原因的.
因為在諾曼書裡一直跟我們說.
因為我們因為上帝的恩典.
憑著我們的相信.
我們能夠稱義在上帝面前.
以致我們將自己的身體當作活祭.
當然將身體當作活祭.
是一個很有咎若一種奉獻的色彩.
今日你將身體當作活祭的時候.
你會怎樣將身體當作活祭呢?.
當我們領受了很多真理的時候.
我們走出來.
這就是一個最好的奉獻.
而這個奉獻,付諸行動.
很多時候我們就像聖經所說.
是甘心樂意的.
因為我們真的滿有上帝的恩典.
我們在這裡再一次用一個奉獻的心.
我們用這個奉獻的心去到上帝面前.
在這一刻我們將我們的身心奉獻給祂.
我們一起在這裡祈禱,默想.
上帝在這一刻的時候.
當你與我們同在的時候.
父啊,叫我們的心敬畏你.
父啊,我們在這當中.
弟兄姊妹,如果你和上帝的關係有疏離.
或者你很忙碌,忘記的時候.
今日上帝告訴你.
祂仍然是看顧你.
弟兄姊妹,如果你這一刻.

$^{41}$你是願意將你自己的身心靈奉獻給上帝的時候.
你可以輕輕的跟上帝說.
我願意.
七愛神,我們在你面前.
我們真是很不配.
但是你何等的恩典.
你自己的拯救再一次每個星期去提醒我們.
無論我們走多遠,離你有多遠.
但是上帝你都在呼喚我們.
教導我們再一次將身心去獻上給你.
因為上帝你悅立.
大主到上帝你悅立的是一個聖潔的器皿.
求主你給我們今日受潔心清.
來到你面前的時候.
我們將我們一切的憂慮.
將我們一切的重擔交給你.
我們甚至要在上帝你面前學習去認罪.
更加承認我們有什麼過犯.
以致你釋放我們.
我們又再活過來.
再一次奉獻給你.
多謝你今日給我們集體的敬拜.
讓我們在集體的敬拜當中.
遇見耶穌基督.
祈禱奉主名求.
阿們.
下一個題目叫做如何測驗神的旨意.
測驗神的旨意.
都是我們基督徒.
相信上帝的人.
一個一生的功課.
無論你是信主的第一天.
直到如今.
都是要去學習這件事情.
因為不是說上帝的旨意千變萬化.
而是因為你人生當中會遇見很多不同的事.
不同的階段.
年幼的時候相信主的.
任務在在當中.
那時候很年輕.

$^{81}$那時候也是被人欺騙.
不是 怎麼會被人欺騙.
相信耶穌.
年長才相信耶穌.
有沒有說這裡.
年長才相信.
也有.
這個不是被人欺騙.
是自己欺騙自己.
都是不會的.
其實上帝的旨意.
是沒有變的.
但我們人生當中.
很多時候變幻的時候.
我們需要測驗上帝的旨意.
在羅馬書十二章第二節.
剛才我讀的是第一章第一節.
今天第二章第二節.
他怎麼說.
我們一起讀一節這段聖經.
我們可以.
你要我按嗎.
我們一起讀這節聖經.
我們再提醒一下自己.
一二三.
不要效法這個世界.
只要心意更新而變化.
叫你們測驗何謂神的善良.
純全可喜悅的旨意.
這一句聖經是一句很美麗的聖經.
這一句聖經的美麗.
不是因為它的容詞的緣故.
而是因為它當中.
它教導我們跟隨上帝的人.
有一個心意更新而變化的心態.
目的是什麼呢.
目的就是說.
當我們恩順清儀之後.
當我們把身體獻上當作活祭之後.
其實其中一個功課.

$^{121}$就是怎樣測驗.
神純全善良可喜悅的旨意.
當一個思想來到你心裡的時候.
其實你怎樣知道這個思想是從哪裡來的呢.
到底是從上帝那裡來.
還是從魔鬼那裡來.
還是甚至從自己那裡來.
你怎樣知道那個想法那個意念.
是從哪裡來.
是不是神的旨意呢.
你們會不會認識一些人.
很喜歡自言自語的.
有沒有.
不知道在座當中有沒有自言自語的人.
誰是你最好的朋友.
很多人都說.
我是我自己的最好朋友.
為什麼呢.
因為我經常跟自己說話.
我不知道誰是你最好的朋友.
會不會你是你自己的最好朋友.
但到底這個旨意來的時候.
是你自己還是上帝呢.
其實如果我們不去測驗.
不去分辨神的話.
其實是一個很危險的事情.
後果很嚴重甚至會致命.
為什麼呢.
因為聖經這樣說.
有一條路人以為正路.
但最後卻是死亡之路.
很多人以為這條路是正的.
但其實最後不是正.
而反而是死亡之路.
你知不知道.
現在其實很多人不是跟自己說話.
現在很多人跟什麼人說話.
是跟AI.
為什麼呢.
因為你都知道.

$^{161}$很多人跟AI說話.
尤其是那些獨居的.
不單止是老人.
年輕人也會有.
經常跟AI說話.
早前我看到一個新聞.
有一個年輕人跟AI說話.
他有一個念頭.
是什麼念頭呢.
輕生.
即是他想自殺.
他跟AI說.
到底怎樣可以最舒服地自殺.
AI就真的教他.
他覺得AI教他這個很適合的時候.
他真的做了這件事.
所以人以為是一條正路.
但其實最後是死亡之路.
其實是很危險的.
早前大家都知道.
師母都回到香港.
我自己一個人都是獨居.
雖然不是老人.
我真的不知道找誰說話.
女兒在酒店.
她經常都沒有check in.
我又不經常在接待處等她.
自己一個人的時候.
兩個多三個星期.
都不知道在家跟誰說話.
會自然自如.
有第五姐妹.
教會真的有很多小組.
所以有一天星期五.
早上又看到我.
晚上又看到我.
為什麼呢.
不是因為我無處不在.
是因為我不想自己跟自己說話.
她回來了.

$^{201}$我終於可以跟AI說話.
我的AI就叫Angela Intelligent.
測驗這件事情.
其實你看聖經.
在保羅書信裡.
在羅馬書用過一次.
測驗上帝的旨意.
但他提醒信徒.
在哥林多前書的時候.
他用過這個詞.
用過兩處的地方.
這兩處的地方都提醒我們.
其實我們都需要去測驗.
測驗什麼呢.
是測驗自己.
我們一起讀這次聖經.
因此無論什麼人.
若用不合適的態度吃主的餅.
喝主的杯.
就是得罪主的身體.
主的血尿.
所以人應當醒察自己.
然後才吃這餅.
喝這杯.
這句聖經大家都很熟悉.
我們每次手煮餐的時候.
牧師都鼓勵大家.
勸勉大家要怎樣.
醒察自己.
「醒察」這個詞.
其實就是剛才的「測驗」.
原來保羅提醒我們.
不單止我們要測驗神的旨意.
在測驗神的旨意之前.
其實我們是需要去醒察自己.
他也要講到.
什麼人要醒察自己.
特別.
當然每個人都需要.
他特別講到.

$^{241}$領受主餐之前.
即是領受上帝恩典之前.
你更加需要去測驗.
醒察自己.
醒察自己的心思意念.
原來越領受得多上帝恩典的人.
這個人就應該是越多去.
醒察自己.
在座的你.
你有沒有領受上帝恩典?.
有.
多不多?.
多.
但你相反的時候.
你想想.
你對自己的醒察.
有沒有這麼多呢?.
這是我們值得深思的問題.
使徒保羅這裡再一次提醒我們.
我們越領受上帝恩典多的人.
其實我們應該越去.
醒察自己.
測驗上帝的心意.
除了我們要測驗上帝的旨意.
醒察我們自己.
原來他另一次的用詞.
用同一個字.
他在哥林多前書的第三章.
第三章的十二到十三節.
這一次聖經又怎樣提醒我們呢?.
他用這個字的時候.
我們一起讀.
若有人用金銀寶石.
草木和街.
在這根基上建造.
各人的工程必然顯露.
因為那日子要將他表明出來.
有否發現.
有否要試驗各人的工程怎樣?.
這句經文的意思.

$^{281}$就是不單止我們要測驗.
我們自己.
我們測驗上帝的旨意.
原來這裡這個字.
都提到上帝都會去試驗.
試驗其實就是這個字.
和那個測驗.
和那個醒察是一個字.
就是說上帝會試驗我們所做的一切.
是的.
原來我們所做的一切.
上帝是會測驗.
到底我們所做的是草木和街.
一燒就沒有呢?.
還是金銀寶石?.
大家想要金銀寶石還是草木和街?.
金價很貴.
搶金潮.
所以大家都要用金去建造很多東西.
去讀這段經文的時候.
我記得我年輕的時候.
去看這段聖經的時候.
都很害怕.
其實我做那麼多東西.
會不會是草木和街?.
去到上帝面前.
一燒就沒有了.
到底我怎樣可以做金銀寶石的工作.
工程.
獻給上帝呢?.
所以那時候很多擔心.
不知道最後會怎樣.
其實你再看那三節聖經.
保羅所用這三個字的時候.
其實原來大家不用害怕.
為什麼不用害怕呢?.
答案是.
原來只要你多領受上帝恩典的時候.
你多去醒察自己.
你多去測驗神的心意的時候.

$^{321}$你是不會再怕上帝去試驗你的.
為什麼你會害怕上帝會試驗你呢?.
因為我們平時只是領受了恩典.
而我們沒有去醒察和測驗神的旨意.
以致當時上帝試驗我們的時候.
我們好像很害怕.
這裡聖經保羅再提醒我們.
你不需要怕上帝試驗你的.
只要你在你的生命當中.
常常測驗神的旨意.
醒察自己去領受上帝的恩典.
這個就是聖經給我們的真理.
事實上在我們的生活中.
是有很大挑戰的.
為什麼呢?.
因為上一秒很多時候我們就有上帝的旨意.
下一秒我們很多的旨意.
就好像從魔鬼那裡來.
其實聖經也有一個這樣的例子.
如果你記得在《馬太福音》第十六章.
當耶穌基督和十二門徒一起的時候.
耶穌問彼得一個問題.
祂問他什麼問題?.
祂問你知不知道我是誰?.
耶穌要肯定他們.
你要知不知道我是誰.
所以其實今天你相信神.
耶穌常常要問你.
你知不知道耶穌是誰?.
你知不知道基督是誰?.
因為你說你是基督徒.
你不知道基督是誰是不行的.
彼得在十二門徒當中彈出來.
第一個說.
他說什麼?.
他說主啊!.
你就是基督.
你就是那位救世主.
你就是Messiah.
你就是我們期待的那位救主.

$^{361}$你是永生神的兒子.
然後怎樣?.
就好像穆斯丹問你們問題.
你們的聖經問題一答就答中.
馬上就百分百.
耶穌怎樣跟他們說?.
耶穌說你說得對.
你這個意念,思維.
並不是屬於血氣能夠指示你的.
是在天上的父給你啟示.
是在天上的父給他啟示.
當有人問你基督教最大的特色是什麼的時候.
當然我們三一神.
聖父,聖子,聖靈.
我們都要告訴別人.
但其實你要告訴別人.
我們基督教最大的特色.
就是一個啟示的宗教.
什麼叫啟示的宗教?.
啟示的宗教就是說.
其實我們不是靠著人的力量去尋找上帝.
而是上帝親自來到人間.
來到我們生命當中.
去啟示給我們看.
若非上帝啟示給你.
你不能受感動.
當你不能受感動的時候.
你又怎會獻上呢?.
你又怎會去測驗神的心意呢?.
所以一切我們都知道是一個啟示的緣故.
是上帝啟示給我們.
換句話來說.
彼得剛才的想法是從哪裡來的?.
很明顯是從天父,上帝那裡來的.
但如果你繼續讀下去的時候.
你會記得發生什麼事.
下一秒的時候.
當耶穌告訴彼得.
你覺得是正確的.
100分給你.

$^{401}$但是我現在即將要上耶路撒冷.
我要被釘死.
我要被捕.
受害等等.
彼得怎麼說呢?.
彼得那時候說什麼呢?.
彼得說.
主啊.
你千萬不可以這樣做.
因為這件事不可以發生在你身上.
當他說完之後耶穌轉過頭來.
馬上跟彼得說了一句話.
祂說什麼呢?.
祂說.
撒旦.
推我後去.
推到我後面.
彼得前一秒.
剛剛才在神那裡領受了一個精彩的答案.
耶穌就是基督.
但是轉念下一秒.
祂竟然叫耶穌基督不要成就救恩.
耶穌就跟他說撒旦推去我的後面.
所以我們有很多想法.
有時候是從神來的.
但是有時候我們不能夠否認.
很多都是從魔鬼那裡來的.
很快我們有時候就會想歪了.
我記得有一個弟兄.
他很關心其他人.
另一個弟兄問他最近身體怎麼樣.
他說最近身體不太好.
血糖很高.
經常都很累.
弟兄說你很好.
你從神而來.
你很關心和我說.
知道我血糖很高.
血壓又高.
好像沒有精神.

$^{441}$你關懷我.
但是想不到下一秒這個弟兄就說.
你A1C多少?.
他說六.
他說六不是很高.
我認識一個十多歲的.
十多歲?還生存?.
是的,還生存.
我帶了一些朱古力雪糕給你.
你試一下吧,沒所謂的.
然後那個弟兄跟他說.
撒旦推我後面.
很多事情,很難的.
這一刻我很關心他.
另一刻可能有些意念不知從何而來.
所以在你的生命當中.
其實你想想.
每一天你都可能要說撒旦推我後面.
我不知道你有沒有說過這句話.
很少吧.
撒旦推我後面,讓我吃掉它.
我不知道.
到底這樣的時候我們怎樣去測驗上帝的旨意呢?.
其實沒有什麼所謂的方程式.
但是有四個大概比較向導.
每逢你有一些難做的決定.
你有一些意念.
你都不知道是不是神來的.
還是自己來的.
還是隔壁的人經常叫我做的.
我不知道是不是真的.
其實這四個問題.
你問一問.
是真的可以幫助到你.
我不敢說一定.
這個就是上帝的旨意.
問完之後.
這四個都是是非題.
四個都是yes.
我覺得這個都走不掉.

$^{481}$都應該大概是.
這四個問題是什麼問題呢?.
第一個問題就是.
到底我這個意念,我這個想法,我這個決定.
是不是符合聖經呢?.
這個第一,很重要.
大家都同意嗎?.
第二就是這件事.
這個想法.
我要做這件事.
或者我這樣的意念.
這個我喜歡的.
它到底有沒有讓我更像基督.
這個又更深了.
不是一個客觀的事情.
當然剛才那個比較客觀.
符合聖經.
這個有沒有讓你更像基督.
更加有life.
更加活.
第三個問題就是.
教會的弟兄姊妹.
即是成熟的那些.
其實有沒有都認可的呢?.
有時候我們很多時候都是當行獨斷.
我覺得不行.
我們有時候其實上帝要我們和一群熟齡的同行者.
這一群同行者其實上帝會透過他們跟我們說話.
所謂教會其實就是這些人.
教會是不是都認可這個決定呢?.
最後都是一個比較個人性一點的.
即是你會感覺到很平安.
如果這四件事你在這個決定底下.
你都發覺全部都yes的時候.
都應該不會是撒旦推我後面吧.
即是這樣說.
這些可以是一些很大的事.
人生很多的決定.
但是有些可以很細微的時候.
其實都可以用得著的.

$^{521}$所以我覺得我們可以一起practice多一點.
當我們說參應上帝的旨意的時候.
其實你每一天嘗試都去問這幾個問題.
或者你說一件事.
不要每件都問.
到底吃叉燒還是吃燒鴨呢?.
這些不要問.
你第一條已經含合性經你都解決不了.
你問牧師都不知道的.
你要問人家的醫生才知道.
問一下你老婆才知道.
等等.
這些可以操練一下的.
每一天嘗試問一件事.
問一下這四條問題.
對你整個素人成長是很健康很有幫助的.
我們每個touch base一下.
因為時間關係未必說得完.
但我們盡量吧.
第一個.
是否符合性經.
我們知道其實神的旨意.
是永遠都不會和自己的話語自相矛盾的.
我們剛才說怎樣測驗神的旨意.
測驗神的旨意就是他的話語.
他的話語和他的旨意是不會自相矛盾的.
吉勒神不會說做一件事明天又要改的.
不是朝令夕改的.
你知道很多人都是情緒化的.
比較情緒化的.
我不知道你是否一個情緒化的人.
是被情緒所控制的去做決定.
但上帝是不會的.
上帝是很consistent的.
祂說愛你的時候.
祂不會因為你理起祂.
祂不愛你的.
不會因為你犯罪.
祂不聽話.
祂不愛你的.

$^{561}$祂會的.
但只是用不同形式的意思.
所以祂不會和聖經原則有違背的事情.
如果你想到一件事.
第一個就發覺和聖經原則是相違背的.
其實你都不用想了.
其他都不用再問了.
就已經不是了.
很多時候.
譬如一些很簡單的例子.
例如謊話,欺騙.
這些很多時候都不合乎聖經的.
但很多時候你的意念都會有的.
為甚麼呢?.
因為欺騙的益處是甚麼呢?.
很多時候它帶來的就是不用煩.
說了一個小小的謊話.
又不是很大.
那就不用煩那麼多了.
我不知道你有沒有這樣的經歷.
一定有的.
我們從小就成長一個謊言欺騙的世代.
很多時候都是這樣成長的.
你想一下.
很多時候你的同學都會欺騙你.
我也會欺騙我的同學.
從小就玩那個遊戲.
那個撲克牌的遊戲.
那個撲克牌的遊戲叫甚麼呢?.
講謊話.
有沒有玩過?.
我出一張牌.
是綠色的.
你信不信?.
玩很多這些是欺騙.
到底是不是呢?.
猜一下呢?.
就是這樣去成長的.
我小時候我媽媽.
我想你都會的.

$^{601}$媽媽問你去哪裡?.
你去哪裡呢?.
你去玩.
然後你當然告訴你媽媽.
她在圖書館.
我記得我媽媽問我.
全世界的謊言.
我當然說全都做了.
不要煩我.
我媽媽問我.
考試考得怎樣?.
我明明是肥佬的.
我告訴她.
我合格.
不要煩.
但是當然了.
這樣講下去的時候.
最後招致的.
是一個很麻煩.
是不是?.
被人打到.
因為這些謊言累積的時候.
最後是爆煲.
這些是經常見到的.
最近都聽到.
很多KK園那些.
是不是?.
那個緬甸炸毀了整個KK園.
這個詐騙集團.
當我看的時候.
他在這方面有很多思想.
我不知道他這次炸.
是不是真炸還是假炸.
可能這個炸都是一個詐騙.
讓中國大陸知道.
我們緬甸做的那些事.
不要再煩我們了.
其實他現在就是說.
不知道在隔壁.
其實已經再萌生一支.

$^{641}$詐騙集團.
我不知道真假.
又說那些很多時候.
炸完之後.
很多主腦逃離了.
緬甸KK園那些地方.
去了哪裡?.
去了緬甸巴提瓦.
我心想.
你逃難都去巴提瓦?.
我們這些人.
大英前輩去旅行.
又去巴提瓦?.
真的是假的.
不止是真是假.
我經常思想.
為什麼可以.
這麼大一個集團.
捉到千多人.
很多人.
那些電話詐騙.
大家其實都經常收到.
誰來找你.
我不認識他.
我不明白.
為什麼可以這麼多.
我再想想.
原來.
詐騙這件事.
他最後有一個逃難了出來.
他問他.
其實裡面情況是怎樣.
他說其實裡面.
不是真的這麼多賣豬仔的.
是很多.
是自願去的.
為什麼呢?.
我心想.
其實為什麼可以.
去訓練一個人.

$^{681}$去騙其他人呢?.
人是可以被訓練到.
去騙其他人的嗎?.
是可以的.
原來我們每個人.
有一個本性.
去欺騙.
去講謊話.
去....
無論什麼原因都好.
而我們其實可以被訓練.
去做一個詐騙的人.
這件事是很恐怖的時候.
但其實你再想想.
又不是.
我們經常都說.
學好很多年.
學壞的又如何.
不用三天.
三分鐘就可以學壞.
其實我們人去學壞.
是很容易的.
所以其實我們滿腦子很多旨意.
很多思想.
都是不合神心意的.
無論是詐騙.
無論是欺騙.
或者是其他的罪惡.
很多思想一來的時候.
如果我們問一問.
其實是不是合乎聖經呢?.
其實就已經幫助很多.
幫助很多.
當然譬如聖經也有說這句話.
我們一起讀.
「說方的嘴唇,為耶和華所憎恨,幸是誠實的,為神所喜悅.」.
這些是上帝的說話.
有時候這些念頭.
我們就是因為沒有去測驗.
我們就希望.

$^{721}$說方比較容易.
永遠都好像不用有什麼責任.
但其實不是.
很大責任.
所以這是很普通的例子.
我們第一個問到底是不是符合聖經呢?.
好,第二個.
第二個我們就問.
到底這個是不是讓我們更加像耶穌基督.
耶穌基督是我們的主.
我們叫基督徒.
我們基督徒的時候.
我們跟隨耶穌.
所以耶穌是我們每一個人.
應該思想上,意見上,感受上的一個標準.
你有沒有試過想過.
為什麼上帝仍然容許我們活在這個世界呢?.
為什麼我們不相信主之後.
就立刻帶我們上天堂呢?.
我們講了很多次這個原因.
你一定知道為什麼.
我們就說.
因為上帝對我們每一個人的生命.
是仍然有目的的.
仍然有使命的.
所以祂不帶我們走.
要我們完成了這些使命才走.
當然其實這是其中一個目的.
不單單是這樣.
更加不單單上帝想我們成為一個好人.
很多不信上帝的人說.
我們好人不就行了嗎?.
其實上帝不單止想我們成為一個好人.
祂想我們更加像耶穌基督.
活得更加像耶穌基督.
有這方面的品格.
在以下我們的生命當中.
在我們還未見到天父之前.
我們要變得更加像耶穌基督.
所以當你有一個意念.

$^{761}$有一個決定要出來的時候.
你要問一問.
這件事會不會讓我更加像耶穌基督呢?.
大家都同意現在這個世界很混亂的.
同不同意嗎?.
有人說現在世界發生的事情很混亂.
為什麼呢?.
關係破裂.
家庭破碎.
國家破碎.
很多個破碎在這裡.
這幅圖大家都知道是什麼事.
知道什麼叫Ice嗎?.
知道吧?那塊冰.
那些是什麼冰呢?.
我上個星期和youth講.
我問他們什麼是Ice.
他們說是從雪櫃拿出來的Ice.
我不是想講那些.
Ice是Immigration and Customs Enforcement.
大家都知道這件事.
不是那塊冰.
不過他這個示威挺好的.
他怎麼寫呢?.
No room for ice in the melting pot.
很有趣.
我們美國是melting pot.
把Ice放在這裡.
走了.
有時候你會發覺很奇怪.
為什麼這個政府可以容許Ice去槍殺無辜的市民呢?.
是不是這樣的呢?.
為什麼我們國家可以去逮捕另一個國家的總統.
什麼原因都好.
我經常想.
如果特朗普被伊朗抓了.
我們會怎樣呢?.
你想想這件事會怎樣.
是不是OK呢?.
我覺得不是OK的.

$^{801}$所以為什麼我們這樣抓人是OK的呢?.
為什麼我們可以明明那個地方不是我們的.
為什麼我們的政府可以說.
這個地方如果我們擁有是更好的呢?.
格陵蘭明明是丹麥的.
為什麼可以這樣說呢?.
為什麼我們的國家可以羞辱我們最親密的盟友.
在我們旁邊.
鄰居來的.
加拿大.
為什麼可以說他可以成為我們第51個依賴呢?.
可不可以呢?.
我覺得世界很混亂.
這個世界是沒有完美的政府.
沒有完美的家庭.
沒有完美的教會.
我們都不完美.
所以當這些事發生的時候.
如果你是沒有耶穌基督的品格.
你不像耶穌的時候.
你會怎樣呢?.
你看到這些事.
你會不安.
你會焦慮.
你會很憤怒.
甚至可能有些人會很生氣.
你會很煩惱.
但是如果你是像耶穌基督的時候.
我不是說你不會生氣.
你不會沮喪.
你不會冷漠等等.
但是你因為耶穌基督的緣故.
你就可能會更多平安.
你更加嘗試怎樣在這個世代混亂當中.
真的成為耶穌基督徒.
跟隨耶穌的基督徒.
成為這個世界的炎石光.
否則很多時候.
有幾句說話我們都會經常說.
我不能夠就這樣算了.

$^{841}$他這樣對我.
我怎樣可以就這樣算呢?.
我只是討回公道.
他傷我先,先療者賤.
我反擊而已.
如果我不反擊.
我這口氣吞不下.
我不知道這幾句說話.
你平常會說哪一句多.
四句都說完了.
不用選擇.
就算說哪一句都好.
其實這些都是受傷後很委屈的想法.
這些意念,憤怒.
是可以理解的.
但是我們知道反擊的憤怒.
是反映不了耶穌基督的養成.
耶穌會怎樣做呢?.
我們經常問耶穌會怎樣做.
有人說耶穌很軟弱.
聖經怎樣說呢?.
它都教導我們紅綠燈.
過馬路要小心.
要慢慢的動腦.
快快的聽.
快快的.
慢慢的說.
這個不說了.
說回經文.
有人說耶穌很軟弱.
為什麼?.
因為經文說.
祂避罵不還口.
受害的時候不算威嚇的話.
其實這個經文的著重點在哪裡呢?.
它的著重點在後面.
這個經文不是說.
耶穌基督避罵不還口.
受害的時候不算威嚇的話.
這個經文說.

$^{881}$耶穌不是不見公義.
我們一起讀一句.
祂只將自己交託給那位.
按公義審判人的神.
當你知道哪一個是審判者的時候.
哪一個真正的大法官的時候.
耶穌完全知道什麼叫不公不義.
因為祂就是被不公不義的對待.
但是祂拒絕讓這種仇恨.
成為祂餘下人生的主人.
耶穌沒有說對方是對的.
沒有說傷害祂的人是對的.
但是祂就將傷害祂的人.
交託給滿有審判權.
歸還給上帝.
當然上帝叫我們更像耶穌的時候.
一切謙卑溫柔忍耐.
是我們一生的功課.
其實你知不知道.
你認識了耶穌之後.
在今天這個混亂的世代.
你仍然可以活出聖靈的果子.
聖靈的果子不只是在天下太平.
大家歡樂的氣氛的時候才有.
在此時刻我們仍然可以有.
因為上帝更加關注的是你的品格.
有沒有更像耶穌基督.
你的個性是否真是耶穌基督.
接著我們會問.
這個意念到底教會弟兄姊妹.
成熟的基督徒會否都認可呢?.
我們經常透過耶穌基督的身體.
神會透過耶穌基督的身體.
跟我們說話.
即我們身邊成熟的基督徒.
他不單可以幫我們確認.
這個想法,這個意念.
有時候他可以溫柔地提醒我們.
這個很重要.
因為上帝從來都不想我們獨自一個生活.

$^{921}$當然有時候一個人生活是沒有辦法的.
但是我們感受上帝的血.
讓我們有一個崇靈的群體.
我明白有時候我們不願意去邊界別人.
不要煩其他人.
不要問他們.
你們有什麼感覺嗎?.
其實不要怕煩.
其實我覺得有時候你問了一句.
你會有很多出人意表的回應.
幫你想通很多事情.
而當你不去問的時候.
其實很多時候是什麼原因呢?.
因為很多時候我們有些驕傲.
老實說.
很多時候我們驕傲.
不願意被接受別人的幫助.
有時候可能你心知肚明 知道有些事情你自己想都不太對.
所以都不想去問別人.
都可能會是這樣.
我自己都試過.
我記得我讀神學的時候.
我經常都有一個意念.
那個意念就是神要我退學.
神要我退學.
我不知道.
我讀書是這樣的.
我從小到大讀神學都是這樣的.
經常都去退學.
但當然讀神學的時候.
神有句說話告訴我.
這句說話是什麼呢?.
這句說話是神在二弗所書的第三章第十節.
祂說.
神要藉著教會.
使天上執政的,掌權的.
現在都要得知到神百般的智慧.
原來上上神真的透過教會去啟示祂的智慧.
教會我再說.
是一班屬靈的弟兄姊妹.

$^{961}$我將這個掙扎帶到當時有些我認識的屬靈弟兄姊妹.
去幫我分辨我是否真的要退學.
到底他們最後怎麼說呢?.
有一個很好的,他不問我什麼原因.
他說我會繼續為你祈禱.
你一有空有這樣的意念.
打給我和我祈禱.
他沒有問我為什麼要有這樣的意念.
我很欣賞他.
另一個他問我為什麼.
為什麼你想退學?.
因為我讀得很辛苦.
因為我是讀電腦的.
所謂理科.
轉了文科要寫很多東西.
很辛苦,要看很多書.
我不是說讀電腦不用看書,不用工作.
大家讀電腦的人都知道.
打機的.
怎樣也好.
我說很辛苦.
他的回應很直接.
他說對了,你這麼辛苦,這是神的旨意.
你不用退學.
讀神學不辛苦,你怎麼讀呢?.
所以如果大家被神呼召你進修.
你想過一定是辛苦的.
如果你想著很單調.
不用交,老師也給你A.
不用想了,這是妄想症.
怎樣也好.
透過和很多弟兄姊妹聊的時候.
我發覺這個負擔是繼續在我身上.
當然上次有一個挑戰我.
他說你先試試退學.
你退一個學期.
如果你退了學期之後.
你的心仍然很混亂.
仍然覺得神真的要你去做.
你就再去報.

$^{1001}$很多很多的建議.
到底教會弟兄姊妹對這件事是怎樣認可呢?.
我覺得是很重要.
其實我們開始這個活因教會的時候.
我覺得這個也是我其中一個問的問題.
當我有這樣的意念的時候.
當然我不只是和幾個同工一起去印證.
到底是不是.
我們教會當中其實有一個CALL TEAM.
這個CALL TEAM.
我們由開始.
最初期我們一起去尋求上帝的心意的時候.
我們看到其實上帝是慢慢一步一步.
將我們的意念更加去convert.
更加見到同一個異象.
即是見到同一樣東西.
最初不是的.
慢慢一路溝通一路去認可的時候.
大家原來真的見到的時候.
這個我覺得是一個很重要的問題.
我們需要問.
每當你有一個意念.
最後因為時間緣故不好意思.
很快最後就是說.
到底這個是不是有上帝的平安呢?.
在裡面.
剛才那個問題我們都會發覺.
其實上帝不會是一個不安的源頭.
如果祂叫你做一些事.
祂不會令你覺得很害怕.
也會的.
祂不會讓你很困難.
因為其實在困難的時間.
上帝仍然可以給我們平安.
聖經有句說話.
祂說.
一齊讀.
神所賜出人以外的平安.
必在耶穌基督裡.
保守你們的心懷意念.

$^{1041}$這句話每讀的時候.
我都覺得很安慰.
即是神會保守我們的心懷意念.
以致你有平安.
因為我們的心懷意念出了問題.
以致沒有平安.
當我剛才說.
我和一些call group和一些同工.
我們去尋求神的旨意.
去建立會長教會的時候.
我們同心吻合我們見到.
但其實我自己心裡面.
仍然有很多的憂慮.
很多的不安.
這個時候我就問.
為什麼我仍然有不安的情緒?.
為什麼我仍然有這麼多的掙扎?.
我和上帝說.
當然這句聖經.
其實之前還有一句.
這句聖經之前那句是說什麼?.
這句聖經之前我們一起讀.
應當音無掛慮.
只要凡事直著禱告.
將你們所要的.
告訴神.
神所賜出人以外的平安.
必在耶穌基督裡.
保守你心懷念.
剛才原來那句聖經之前.
就是說應當音無掛慮.
原來我發覺我一路去.
做神的工作的時候.
我經常都有一種不安的感覺.
而我怎樣能夠得到上帝的平安呢?.
我不是為了得到上帝的平安而去做這些.
這句聖經告訴我們.
我要將一切的憂慮.
直著禱告和祈求.
交給上帝.

$^{1081}$就是剛才我們那個call group.
我們由開始會忍教會到現在.
我們每一個星期.
其實都一起祈禱.
在過去這一段時間.
其實我有很多事情.
有好幾件事情.
我自己都有很多不安的.
但都是透過那個祈禱會.
我們一班人禱告的時候.
上帝又再保守我的心懷念.
原來有些事情我想多了.
有些事情其實我不需要那麼擔憂.
而那個平安就再回來了.
最近當我們祈禱的時候.
我們會忍教會.
神帶領我們去到一個很重要的里程碑.
這個里程碑就是我們.
拿到這個501(c)(3)的status.
這個status我不知道大家明白什麼意思.
這個不單止是一個法律上的身份.
是在美國社會裡面.
是一個被認可,被信任.
是可以運作成為一個教會.
的一個基礎的status.
我過去都有很多不安的.
為了這件事.
但在我禱告的時候.
上帝又再給我很多的平安.
今天我們有這個status.
我經常都說沒有IG沒有LG.
不知道為什麼有很多G.
因為我們其實一直都有IG去扶植我們.
去幫助我們.
做個預告吧.
3月8日PG就會來講道.
那天我們就要把他放上台.
他代表IG.
他們很成功.
一年不夠就植了一堂出來.

$^{1121}$怎樣也好.
我們將一切榮耀歸給上帝.
大家在上帝的平安裡面.
是會有不安的很多時候.
做上帝的工作都會.
但當我們一直去尋求.
去測驗上帝的旨意的時候.
剛才那四個問題.
記不記得第一個是什麼.
含脈扶聖經.
我們有沒有更加像耶穌基督.
第三.
教會弟兄姊妹會不會都認可.
第四.
你們是不是都感覺到有這份平安.
我們一起祈禱吧.
在你的帶領當中.
仍然去經歷你每一個的恩典.
求你繼續幫助我們.
用你的話語來.
測驗我們每一個的意念.
塑造我們更加像耶穌基督一樣的品格.
無論是遇到困難的時候.
是不容易的時候.
我們都能夠有這份謙卑柔和的心.
當我們不確定的時候.
主人我們亦都願意去請教.
教會一些屬靈的同伴.
我們求主你繼續讓你的平安充滿我們.
讓我們跟隨你的路上.
是不動搖.
無論是春夏秋冬.
禱告奉主名求.
阿姆!.
\newpage



\section{撒母耳記上 3:1-21}
\label{sec:031633}
\textbf{孩子是神給慕道園的禮物 (普/粵)}~~~~田俐力傳道
\newline
\newline
連結: \href{http://www.church.com.hk/acms/content.asp?site=cdc\&op=show\&type=product\&code=031633\&layout=sermon}{\texttt{031633}} ~~~~ 語音日期: 2026-02-08
\newline
\newline
\hyperref[sec:031632]{\small{< < < PREV SERMON < < <}}
~
\hyperref[sec:index_chronic]{\small{[返順時目]}}
~
\hyperref[sec:031634]{\small{> > > NEXT SERMON > > >}}
\newline
\newline
撒母耳記上 3:1-21
\newline
\begin{longtable}{cl}
\hline
\hline
章節 & 經文 (和合本修訂版)\\
\hline
3:1 & \begin{tabularx}{0.7\textwidth}{X} 那孩子撒母耳在以利面前事奉耶和華。在那些日子，耶和華的言語稀少，不常有異象。 \end{tabularx} \\ \\ \relax
3:2 & \begin{tabularx}{0.7\textwidth}{X} 那時，以利在自己的地方睡覺；他眼目開始昏花，不能看見。 \end{tabularx} \\ \\ \relax
3:3 & \begin{tabularx}{0.7\textwidth}{X} 神的燈還沒有熄滅，撒母耳睡在耶和華的殿內，神的約櫃就在那裡。 \end{tabularx} \\ \\ \relax
3:4 & \begin{tabularx}{0.7\textwidth}{X} 耶和華呼喚撒母耳，撒母耳說：「我在這裡！」 \end{tabularx} \\ \\ \relax
3:5 & \begin{tabularx}{0.7\textwidth}{X} 他跑到以利那裡，說：「你叫我嗎？我在這裡。」以利說：「我沒有叫你，回去睡吧。」他就回去睡了。 \end{tabularx} \\ \\ \relax
3:6 & \begin{tabularx}{0.7\textwidth}{X} 耶和華又呼喚撒母耳。撒母耳起來，到以利那裡，說：「你叫我嗎？我在這裡。」以利說：「我兒，我沒有叫你，回去睡吧。」 \end{tabularx} \\ \\ \relax
3:7 & \begin{tabularx}{0.7\textwidth}{X} 那時撒母耳還未認識耶和華，耶和華的話也未曾向他啟示。 \end{tabularx} \\ \\ \relax
3:8 & \begin{tabularx}{0.7\textwidth}{X} 耶和華第三次再呼喚撒母耳。撒母耳起來，到以利那裡，說：「你叫我嗎？我在這裡。」以利才明白是耶和華呼喚這小孩。 \end{tabularx} \\ \\ \relax
3:9 & \begin{tabularx}{0.7\textwidth}{X} 以利對撒母耳說：「你回去睡吧。他若再叫你，你就說：『耶和華啊，請說，僕人敬聽！』」撒母耳就回去，仍睡在原處。 \end{tabularx} \\ \\ \relax
3:10 & \begin{tabularx}{0.7\textwidth}{X} 耶和華來站著，像前幾次呼喚：「撒母耳！撒母耳！」撒母耳說：「請說，僕人敬聽！」 \end{tabularx} \\ \\ \relax
3:11 & \begin{tabularx}{0.7\textwidth}{X} 耶和華對撒母耳說：「看哪，我在以色列中必行一件事，凡聽見的人都必雙耳齊鳴。 \end{tabularx} \\ \\ \relax
3:12 & \begin{tabularx}{0.7\textwidth}{X} 我指著以利家所說的話，到了時候，必從頭到尾應驗在以利身上。 \end{tabularx} \\ \\ \relax
3:13 & \begin{tabularx}{0.7\textwidth}{X} 我曾告訴他，我必永遠懲罰他的家，因為他知道自己的兒子作惡，褻瀆神，卻不禁止他們。 \end{tabularx} \\ \\ \relax
3:14 & \begin{tabularx}{0.7\textwidth}{X} 所以我向以利家起誓：『以利家的罪孽，就是獻祭物和供物，也永不得贖。』」 \end{tabularx} \\ \\ \relax
3:15 & \begin{tabularx}{0.7\textwidth}{X} 撒母耳睡到天亮，就開了耶和華殿的門。撒母耳害怕，不敢將異象告訴以利。 \end{tabularx} \\ \\ \relax
3:16 & \begin{tabularx}{0.7\textwidth}{X} 以利呼喚撒母耳說：「我兒撒母耳！」撒母耳說：「我在這裡！」 \end{tabularx} \\ \\ \relax
3:17 & \begin{tabularx}{0.7\textwidth}{X} 以利說：「他對你說了甚麼話，你不要向我隱瞞。你若將他對你所說的話向我隱瞞一句，願神重重懲罰你。」 \end{tabularx} \\ \\ \relax
3:18 & \begin{tabularx}{0.7\textwidth}{X} 撒母耳就把一切話都告訴以利，並沒有隱瞞。以利說：「他是耶和華，願他照他看為好的去做。」 \end{tabularx} \\ \\ \relax
3:19 & \begin{tabularx}{0.7\textwidth}{X} 撒母耳長大了，耶和華與他同在，使他所說的話一句都不落空。 \end{tabularx} \\ \\ \relax
3:20 & \begin{tabularx}{0.7\textwidth}{X} 從但到別是巴，所有的以色列人都知道耶和華立撒母耳為先知。 \end{tabularx} \\ \\ \relax
3:21 & \begin{tabularx}{0.7\textwidth}{X} 耶和華又在示羅顯現，因為耶和華在示羅藉他的話向撒母耳啟示他自己。 \end{tabularx} \\ \\
[1ex]
\hline
\hline
\end{longtable}
$^{1}$各位弟兄姊妹平安.
各位弟兄姊妹們平安.
好 我們一起來低頭禱告.
我們一起來低頭禱告.
恩賜給我們慕道園華人精心會這個大家庭.
讓我們弟兄姊妹可以在一起彼此相愛.
彼此扶持 彼此建立.
讓我們有機會能夠帶領他們認識你 親近你.
今天主啊就求你的聖靈.
風風滿滿地充滿我們整個會堂.
充滿我們每一位弟兄姊妹的心.
也求你藉著使女的口來傳講你的信息.
主啊 也求你用你自己的話語.
來觸摸每一位弟兄姊妹的心.
來啟示你自己的旨意.
以上禱告是奉主耶穌基督聖名祈求.
阿們.
好 今天我想跟大家分享一個非常重要.
但是又常常被人忽視的一個話題.
今天我想跟大家分享一個非常重要.
但是時時被忽略的話題 就是小朋友.
當我們每次思想到孩子的時候.
我們可能會想到孩子是下一代.
我們有時思想孩子的時候.
我們會以為他們是我們下一代.
我們會覺得他們是教會的未來.
我們或者會覺得他們是教會的未來.
是還需要成長的信徒.
但是聖經當中不是這麼說的.
我們來看看詩篇127篇第三節.
我們一起來看看詩篇第127篇第三節.
孩子是耶和華所賜的產業.
是神的賞賜 是神親自託付的禮物.
孩子是耶和華所賜的產業.
是神的賞賜 是神親自託付的禮物.
所以在神的計劃當中.
孩子並不是未來才重要.
所以在神的計劃當中.
孩子不是未來才重要的.
孩子是現在就寶貴.

$^{41}$現在就被神使用.
現在就值得全教會的投入.
孩子現在就是寶貴的.
現在就被神所使用.
現在就值得全教會投入的.
好了 今天我想跟大家分享三個部分.
第一是從聖經看孩子的價值.
孩子在神眼中何其珍貴.
今天跟大家想分享的是三部分.
第一是從聖經看到孩子的價值.
孩子在神眼中是何其的珍貴.
第二個部分是從慕道員的孩子.
看見神的作為.
孩子正在成為祝福.
第二是從慕道員的孩子中.
看見神的作為.
孩子正在被神所祝福.
第三部分是全教會的呼召.
如何一起建立進前的下一代.
第三部分是全教會的呼召.
怎樣一起建立敬虔的下一代.
好 我們先來看看聖經當中孩子的價值.
我們先來看一下聖經當中孩子的價值.
當我們翻開聖經的時候.
我們會發現神對孩子的心意清晰又強烈.
當我們翻開聖經的時候.
我們會發現神對孩子的心意清晰又強烈.
孩子不是次要的.
孩子是神要祝福的.
孩子是神要使用的.
孩子是神同在的對象.
今天我就想跟大家分享.
聖經當中四個比較突出的孩子的故事.
今天我想跟大家分享.
聖經當中四個比較突出的孩子的故事.
每一個都是神在小小生命當中.
大大作為的見證.
每一個都是在這些孩子的小小生命之中.
有大大作為的見證.
好 我們第一個來看一下撒母耳.

$^{81}$一個孩子的傾聽改變一個國家的命運.
好 我們第一個看的就是撒母耳.
一個孩子的傾聽來改變一個國家.
根據經上記載.
撒母耳剛剛斷奶不久.
就被送到聖殿裡面去服侍神.
聖經記載撒母耳剛剛斷奶不久的時候.
就被送到聖殿去服侍主.
大家記不記得當時的以色列民族是什麼樣的境地.
大家記不記得當時以色列的民族是什麼樣的境地.
整個以色列民族屬靈黑暗.
祭司以利年老.
兒子悖逆.
百姓也遠離神.
但是神就在這個時候選擇了一個孩子.
他既不是君王.
也不是先知.
更不是勇士.
但是只是一個非常單純的能夠傾聽神聲音的孩子.
當神第一次呼召他的時候.
他甚至分不清楚是誰在說話.
他甚至分不清楚究竟是誰在跟他說話.
但是祭司以利告訴他如何回答之後.
但是當以利祭司教他如何回答之後.
他那句單純的回應卻震動了整個世代.
撒母耳說:請說,僕人敬聽.
親愛的弟兄姊妹.
一個孩子順服.
足以改變一個民族的屬靈方向.
親愛的弟兄姊妹.
一個小孩子的順服.
足以改變一個民族的屬靈方向.
孩子的心是神最容易澆灌.
最容易觸摸.
也是最容易感動的土壤.
孩子的心是神最容易澆灌.
最容易觸摸.
也是最容易感動的土壤.
好,我們再來看第二個孩子約西亞.
一個八歲登基.

$^{121}$但是帶來了復興的孩子王.
我們來看看第二個小孩子約西亞.
是八歲登基帶來了復興的王子.
他的父親祖父都行惡.
百姓拜偶像尼契神.
但是聖經當中說.
約西亞他行耶和華眼中看為正的事.
約西亞行耶和華眼中看為正的事.
那究竟約西亞行了什麼正事呢.
大家記不記得他行了什麼正事.
大家記不記得他行了什麼正事呢.
他十幾歲的時候.
約西亞就開始潔淨聖殿.
十幾歲的時候約西亞就開始潔淨聖殿.
二十幾歲的時候就帶領全國人民讀律法.
二十幾歲的時候他帶領全國人民讀律法.
他拆毀偶像.
修繕神的殿.
與神立約.
恢復愉悅節.
帶來全國的屬靈復興.
所以神的復興不是從成人才開始的.
所以神的復興不是從成人開始的.
而是從一個願意尋求神的少年開始.
所以這裡告訴我們.
孩子並不是太小而沒辦法影響世界.
相反孩子常常是神復興計劃當中.
最柔軟最敞開的器皿.
相反的孩子永遠是神復興計劃中最柔軟最敞開的器皿.
第三個我想分享的是五餅二魚的孩子.
小小的奉獻卻被神大大使用.
第三個我想分享的是五餅二魚的孩子.
他所奉獻的小小的奉獻是大大被神所使用的.
當我們讀到新約聖經的時候.
五餅二魚的故事大家應該都非常熟悉了吧.
當我們讀到新約聖經的時候.
五餅二魚的故事大家應該都非常熟悉了吧.
哪怕是非基督徒可能都聽過這個神蹟了.
甚至有些非基督徒都聽過這個神蹟了.
當我們很多的人都關注在五個餅兩條魚.

$^{161}$卻餵飽了五千多人的神奇的時候.
當我們的人只關注五餅二魚餵飽了五千人的神蹟的時候.
又有多少個人思考過.
這五個餅兩條魚究竟從何而來呢?.
有多少人會思考這五餅二魚究竟從何而來呢?.
當時有五千個男人,還不包括婦女和孩子.
當門徒只看到缺乏,困難,不可能供應的時候.
當門徒看到缺乏,困難,不能夠供應的時候.
有一個孩子願意拿出他全部的午餐.
對於人來說可能那只是一頓簡單的午餐.
對於普通人來說可能是一個簡單的午餐.
但是對於神來說卻能餵飽幾千人.
甚至還能剩下十二籃.
但是對於神來說就能夠餵飽五千人.
甚至可以剩下十二籃的.
所以弟兄姊妹,神不是看我們有多少.
所以弟兄姊妹們,神不是看我們有多少.
能做多少.
我們能夠做多少.
懂多少.
我們究竟明白多少.
神看的是.
我們是否願意把自己所有的都獻給祂.
神看重的是我們願不願意將我們所有的都獻給祂.
所以孩子的心往往比大人更單純.
更願意相信在神凡事都能.
好,弟兄姊妹,當我們在聊起孩子的時候.
不要忘記我們的主耶穌基督曾經也是.
好,弟兄姊妹們,當我們講起孩子的時候.
我們不要忘記我們的耶穌基督都曾經是.
是一個嬰孩.
是一個孩子.
是一位可慕神話語的少年.
是一個嬰孩.
是一個孩子.
也是學慕神話語的少年.
神本可以讓救主以成年人的形象降臨.
神可以讓救主是用一個成年人的形象降臨的.
但是沒有.
但是他沒有這樣做.

$^{201}$他選擇讓耶穌從一個嬰孩開始.
他選擇耶穌從一個嬰孩開始.
在母親的懷抱裡成長.
在家人的陪伴下學習.
在家人的陪伴之下學習.
在聖殿裡尋求神的話語.
我們的主耶穌從嬰兒開始.
就被託付給了一個家庭.
我們的主耶穌基督從嬰孩開始就被託付給一個家庭.
路加福音說.
他生了投胎的兒子.
被布包起來放在馬槽裡.
路加福音說他生了投胎的兒子.
被布包起來放在馬槽裡.
弟兄姊妹那是救贖的開始.
卻不是成年的耶穌.
而是一個極其微弱的小嬰孩.
是一個極其微弱的小嬰孩.
救主的降臨不是以天君萬馬的榮耀.
救主的降臨不是天兵萬馬的榮耀.
而是以最需要被呵護.
最需要被照顧的方式來臨.
而是以最需要被呵護.
最需要被照顧的方式來臨.
為什麼呢?.
大家有沒有想過為什麼?.
因為神要告訴我們.
孩子的生命是神用來榮耀他的方式.
一個嬰孩的誕生.
就是神計劃的開始.
就是救恩歷史的起點.
主耶穌也經歷了孩子的成長.
主耶穌基督也經歷了孩子的成長.
聖經告訴我們.
孩子漸漸長大,強健起來.
充滿智慧,又有神的恩在他身上.
孩子漸漸長大,強健起來.
充滿智慧,又有神的恩在他身上.
想像一下.
耶穌也曾經一一壓壓地學說話.

$^{241}$大家想像一下.
耶穌基督曾經是以一壓壓地學說話.
跌跌絆絆地學走路.
學習如何與家人相處.
和他屬世的父親約瑟學做木匠活.
學習和他地上的父親約瑟做木匠生活.
在聖殿裡學習禱告.
在聖念中學習禱告.
誦讀經文.
學習順服.
耶穌並不是跳過了童年.
耶穌基督沒有跳過了童年.
而是親自進入孩童的生命.
而是親自進入了孩童的生命.
成為所有孩子的榜樣.
也成為我們父母和教會的提醒.
如果神的救贖計劃選擇以一個孩子開始.
如果神的救贖計劃選擇用一個孩子來開始的話.
但我們又怎麼能輕看孩子呢?.
我們又怎麼能輕看孩子呢?.
主耶穌還為我們樹立了一個榜樣.
讓我們能夠知道孩子也能夠可悟神的話語.
主耶穌基督也為我們樹立了一個榜樣.
就是孩子也能夠愛慕神的話語.
當耶穌十二歲的時候.
祂在聖殿裡面聆聽,發問,思想.
讓眾人都驚訝祂的智慧.
讓所有的人都驚訝祂的智慧.
祂的父母找了祂三天,非常的焦急.
祂的爸爸媽媽找了祂三天,非常的焦急.
當找到祂的時候,主耶穌說了什麼?.
當找到祂的時候,主耶穌跟他們怎麼說?.
耶穌說:為什麼找我呢?.
豈不知我當以我父的事為念嗎?.
耶穌說:為什麼找我呢?.
你豈不知道我當以我父的事為念嗎?.
弟兄姊妹,主耶穌在最青春,最活躍,最跳皮,最容易分心的年紀裡.
祂沒有迷失.
弟兄姊妹們,主耶穌在最青春,最活躍,最跳皮,最容易分心的時候.
祂自己沒有迷失.

$^{281}$祂選擇渴慕神的家.
祂選擇渴慕神的話語.
祂選擇渴慕神的心意.
所以這裡提醒我們.
孩子不僅能認識神.
他們還可以非常地渴慕神.
他們可以理解屬靈的事.
愛神的話語.
聽見神的聲音.
回應神的呼召.
他們可以在教會裡成長.
他們也可以在教會裡成長.
也可以在教會裡經歷呼召.
而且主耶穌也向我們顯明了.
祂愛孩子.
而且主耶穌也向我們顯明了.
祂愛孩子.
我們也應該愛孩子.
親愛的弟兄姊妹.
你們是否明白.
親愛的弟兄姊妹.
你們是否明白.
當我們擁抱孩子.
我們就擁抱了主耶穌.
來過的生命的階段.
我們是擁抱著主耶穌.
來過的生命的階段.
當我們陪伴孩子成長.
我們就在參與神對他兒子.
當初最初的信念.
就好像神當初對他自己兒子的.
當初的心意.
當我們投資孩子的屬靈生命.
當我們投資孩子的屬靈生命的時候.
我們就在回應主耶穌的那句.
「讓孩子到我這裡來」.
「照顧孩子」.
「陪伴孩子」.
「教導孩子」.
「愛孩子」.

$^{321}$「照顧孩子」.
「陪伴孩子」.
「教導孩子」.
「愛孩子」.
「這不是兒童事工而已」.
「這是耶穌基督的心」.
「因為祂自己曾經就是一個孩子」.
「祂也明確地告訴我們」.
「祂喜悅孩子」.
「祂很喜悅孩子」.
「在天國裡正是孩子們」.
「在天國裡面正是孩子們」.
「當我們回顧聖經當中」.
「一個個鮮活孩子的例子的時候」.
「當我們回顧聖經」.
「一個個鮮活孩子的小孩故事的時候」.
「神讓我們看見」.
「孩子能聽見神」.
「孩子能影響國家」.
「孩子能奉獻並經歷神蹟」.
「孩子能成為天國的榜樣」.
「所以孩子不是教會的未來」.
「孩子是教會的現在」.
「所以孩子不是教會的未來」.
「孩子是教會的現在」.
「弟兄姊妹」.
「我想再帶大家看看」.
「慕道元孩子身上的神蹟」.
「弟兄姊妹」.
「我想給大家看看」.
「慕道元孩子身上的神蹟」.
「聖經的故事並沒有結束」.
「聖經的故事沒有結束」.
「神仍然在孩子們身上做工」.
「神仍然在聖恩仔的身上做工」.
「今天在我們慕道元華人敬信會」.
「在我們自己孩子身上」.
「今天在我們慕道元環贈信會」.
「在我們自己的聖恩仔的身上」.
「神仍然行神蹟」.

$^{361}$「神仍然在行神蹟」.
「作為兒童傳道」.
「我想跟大家分享一下」.
「關於我們自己孩子真實的故事」.
「作為一個兒童的傳道」.
「我想跟大家分享」.
「我們的聖恩仔真實的故事」.
「首先我想說的是」.
「孩子的禱告常常比大人更真誠」.
「首先我要說」.
「聖恩仔的禱告」.
「往往比大人更加真誠」.
「我想分享三個禱告小勇士」.
「我想跟大家分享」.
「三個禱告的小勇士」.
「第一個是真誠的禱告者Aiden」.
「大家看看這張照片」.
「就是我們的小Aiden」.
「你看看這張照片」.
「這個是我們的小Aiden」.
「他就是我們國語堂」.
「全弟兄和Wendy姊妹的兒子」.
「他是我們國語堂」.
「全弟兄和Wendy姊妹的兒子」.
「在洛杉磯山火的時候」.
「他們家中連續停電」.
「也斷了網路」.
「他們家中連續停電」.
「也斷了網路」.
「那個時候小Aiden」.
「第一次開始向神禱告」.
「這個時候小Aiden」.
「第一次向神禱告」.
「他說神啊」.
「求你讓我家來電吧」.
「恢復網路吧」.
「他說神啊」.
「求你讓我家有電」.
「和恢復網路」.
「你看他的禱告」.

$^{401}$「沒有華麗的詞藻」.
「沒有華麗的詞操」.
「也沒有神學術語」.
「只有單純的需要」.
「只是單單的需要」.
「但是耶穌就是最喜愛」.
「這種單純信靠的心」.
「但是耶穌基督最喜愛」.
「這種單純信靠的心」.
「所以神就神奇地」.
「應允了孩子的禱告」.
「所以神就很神奇地」.
「應允了他的禱告」.
「他家中立刻來電」.
「也恢復了網路」.
「他家中立刻來電」.
「和恢復了網路」.
「重要的是」.
「最重要的是」.
「從那以後」.
「從那天開始」.
「Aiden每天都開始禱告」.
「而且他還和爸爸一起」.
「帶領家裡的飯前禱告」.
「而且和他爸爸一起」.
「帶領家裡的飯前禱告」.
「睡覺之前也會禱告」.
「睡覺之前他自己會禱告」.
「而Aiden也會運用屬靈的恩賜」.
「用自己的畫筆」.
「畫出聖經的故事」.
「Aiden也能夠運用」.
「屬靈的恩賜」.
「用畫畫畫出聖經的故事」.
「這就是Aiden畫的諾亞方舟」.
「這就是他畫的諾亞方針」.
「非常多的細節」.
「整個聖經故事裡面的小細節」.
「全都在裡面」.
「很細節」.

$^{441}$「將聖經裡面每一個細節」.
「都畫在他的畫中間」.
「這幅畫就掛在我們兒童堂」.
「這幅畫就掛在我們兒童的堂裡面」.
「如果弟兄姊妹感興趣」.
「歡迎去參觀」.
「如果你有興趣的」.
「歡迎大家去參觀」.
「好,而且孩子的屬靈生命的改變」.
「也影響了慰信主的媽媽」.
「還有這個小孩屬靈的生命改變」.
「也影響了他微信主的媽媽」.
「我們的Wendy姊妹」.
「在去年的感恩節受洗歸主」.
「我們的Wendy姊妹」.
「在去年的感恩節受浸歸主」.
「Amen」.
「好,第二個我想介紹的」.
「是一個充滿愛的禱告者」.
「星星」.
「第二個我想介紹給大家」.
「充滿愛的禱告者」.
「星星」.
「星星一家」.
「他是英文堂Jason弟兄的兒子」.
「星星是我們英文堂Jason弟兄的兒子」.
「他們一家是從前年的野餐遊園會」.
「之後來到教會」.
「他們一家是從前年的野餐遊園會」.
「之後來到教會崇拜」.
「但是孩子的媽媽非常忙碌」.
「但是這個小孩的媽媽工作非常忙碌」.
「沒法參加崇拜」.
「但是當孩子們學會如何禱告之後」.
「他們開始為媽媽禱告」.
「當這個小孩學會了如何禱告之後」.
「他開始為他的媽媽禱告」.
「這讓一個佛教徒的媽媽深受感動」.
「另一個佛教徒的媽媽深受感動」.
「也開始求告我們的神」.

$^{481}$「這一家人也非常受祝福」.
「爸爸Jason也是在去年的感恩節受洗歸主」.
「爸爸Jason上一年也是感恩節受洗歸主」.
「Amen」.
「這裡我有一個小片段,小video」.
「想跟大家分享是星星為他媽媽的禱告」.
「這裡我有一個小片段想跟大家分享」.
「就是星星為他媽媽的禱告」.
「親愛的天父」.
「我今天在這裡和您說話是因為」.
「我想為我一個家庭成員」.
「我的媽媽學校去禱告」.
「因為她每天早上很早去照顧工作」.
「很晚回來」.
「我求您讓她的身體健康」.
「讓她在過去和回來的路上一路順風」.
「我祈禱基督的名字」.
「Amen」.
「愛您」.
「非常地感恩」.
「我們也為這位媽媽來禱告」.
「希望她能夠早日也來到神的面前」.
「我也為星星的媽媽禱告」.
「希望她早日能夠來到神的面前」.
「好,第三個我想分享的是一對追求和睦的禱告者」.
「Eason和Jackson」.
「第三個我想和大家分享是追求和睦的禱告者」.
「Eason和Jackson」.
「Eason和Jackson是我們歡歡姊妹的兒子」.
「Eason和Jackson是我們Angel姊妹的兒子」.
「他們還有一個姊姊Kitty」.
「在我禱告生命當中」.
「最讓我感動的一次禱告」.
「在我禱告生命當中」.
「最讓我能夠感動的一次禱告就是」.
「就是和Connie,Esther」.
「還有Eason,Jackson的一次禱告」.
「就是和Connie,Esther,導師和Eason,Jackson的一次禱告」.
「那一次Eason,Jackson用他們稚嫩的小眼睛」.
「看著我們每一個老師」.

$^{521}$「那一次Eason和Jackson用他們自戀的小眼睛」.
「看著我們這些導師」.
「他說老師我有什麼可以為你們禱告的嗎?」.
「他說老師我有什麼可以為你們代禱啊?」.
「但是當我問到Eason,Jackson」.
「有什麼老師可以為你們禱告的嗎?」.
「但是當我問到Eason,Jackson」.
「有什麼老師可以為你們禱告的呢?」.
「孩子們說」.
「細蚊仔就說」.
「我希望媽媽能夠身體健康」.
「我希望我媽媽能夠身體健康」.
「我希望我們的關係能夠更好」.
「我希望我們的關係會更加好」.
「我希望我能夠更進步」.
「我希望我能夠更加進步」.
「我的英語能夠更好」.
「我的英文會更加好」.
「大家看看孩子的禱告」.
「你看看細蚊仔的禱告」.
「他們知道生命當中最重要的是什麼」.
「不是金錢」.
「不是物質」.
「也不是地位」.
「是關係」.
「所以我相信這樣合乎神心意的禱告」.
「是一定會被應允的」.
「所以我相信這些合乎神心意的禱告」.
「一定會蒙神所應允的」.
「感恩的是前段時間歡歡告訴我」.
「前一段時間呢」.
「方方Angel姊妹就告訴我」.
「說Jackson被老師表揚了」.
「Jackson被老師表揚了」.
「有進步」.
「大家看看」.
「神是不是就在當中做工了」.
「大家看看神是不是在當中做工呢」.
「是的」.
「所以我們要相信」.

$^{561}$「神喜悅孩子們的禱告」.
「所以我們要相信」.
「神喜悅孩子們的禱告」.
「孩子們的禱告是真誠的」.
「有力量的」.
「孩子們的禱告是真誠的」.
「是有力量的」.
「另外一個我想分享的是」.
「孩子是勇敢的傳福音者」.
「另外想和大家分享的是」.
「孩子是勇敢的傳福音者」.
「我們一直很擔心」.
「在美國的公立教育體系裡面」.
「不可以談論信仰」.
「我們一直都很擔心」.
「在美國的公立教育體系裡面」.
「不可以分享信仰」.
「但是當我和Emily」.
「Hailey」.
「Mia」.
「這幾個大一點的孩子交談的時候」.
「但是當我和Emily」.
「Hailey」.
「Mia」.
「這幾個大一點的小女生交談的時候」.
「她們落落大方的告訴我沒問題」.
「她們很大方的跟我說沒問題的」.
「而且年齡小的孩子」.
「還會在學校跟朋友談論神」.
「而且還有一些年齡小的信仰者」.
「會跟同學在學校討論神」.
「他們還會帶朋友來教會」.
「他會帶朋友來我們教會」.
「Jackson就帶來了Bolin」.
「我的兒子Joe呢」.
「也會經常帶好朋友來參加主日學」.
「我兒子Joe呢」.
「也都時時會帶他的好朋友來」.
「參加我們的主日學」.
「我們也要持續地」.

$^{601}$「為這些孩子的父母禱告」.
「我們也都要持續地」.
「為這些小朋友的父母禱告」.
「願他們有一天」.
「能夠親近神認識神」.
「讓他們有一天都能夠親近神」.
「跟隨神」.
「弟兄姊妹不要忘記」.
「弟兄姊妹們不要忘記」.
「孩子們才是公立教育體系裡」.
「福音的核心傳播者」.
「小朋友才是公立教育體系裡面」.
「最福音的最熱心的傳播者」.
「孩子的信心往往能改變」.
「大人的世界」.
「小朋友的信心往往都能夠」.
「改變大人的世界」.
「我們再來看看孩子們的敬拜」.
「我們再看一看小朋友的敬拜」.
「孩子投入的敬拜」.
「能激勵整間教會」.
「小朋友投入的敬拜」.
「能夠激勵整個教會」.
「你們有沒有看過孩子們」.
「舉起雙手來敬拜」.
「你們有沒有看過小朋友」.
「舉起雙手來敬拜」.
「特別是小班的小朋友」.
「尤其是小班的小朋友」.
「他們沒有顧慮」.
「沒有羞澀」.
「不會害羞」.
「他們完全單純地把心獻給神」.
「我們的 Annie 姊妹的寶貝 Isabel」.
「我們 Annie 姊妹的寶貝 Isabel」.
「一開始被 Annie 抱在懷裡面崇拜」.
「一開始的時候是 Annie 抱在懷中來敬拜的」.
「如今可以站在台上手舞足蹈」.
「全情投入」.
「如今可以站在台上手舞足蹈」.

$^{641}$「全情投入」.
「還有我們國語堂陸弟兄的女兒小恩怡」.
「她自己可以唱很多首敬拜歌曲」.
「她自己可以唱很多敬拜歌曲」.
「每次聽到她甜美的聲音」.
「每一次聽到她甜美的聲音」.
「看到她帶著小酒窩的笑」.
「看到她帶著小酒窩的笑容」.
「我都能感受得到神的光和愛」.
「當然我們也有不知道什麼是敬拜的孩子」.
「當然我們也有不知道什麼是敬拜的小朋友」.
「但是當我們的老師教導和帶領之後」.
「我們的孩子們都知道敬拜的重要性」.
「我們的小朋友都知道敬拜的重要性」.
「也都能夠認真的來敬拜神」.
「所以孩子們敬拜的活力與熱情」.
「也激勵著我們更投入地將自己全然獻上」.
「所以孩子們敬拜的活力和能夠激勵我們」.
「更加投入將自己傳言的獻上」.
「而且我們的孩子們現在也在成為未來的領袖」.
「而且我們的孩子們也在成為未來的領袖」.
「Kalista時常會幫助照顧小一點的弟弟妹妹」.
「Joshua,Joe,星星是背經文的高手」.
「Joshua,Joe,星星是背誦金句的高手」.
「Aidan,Jackson,Ethan也會在家裡和父母一起讀經禱告」.
「還有我們的Hayden經常會挑戰我們的老師」.
「問出非常深奧的問題」.
「還有Hayden時時都會更正我們老師問一些很深奧的問題」.
「而且大家應該也開始發現我們的孩子也進入服事了」.
「現在大家也都發現我們的小孩子已經開始進入服事了」.
「有的孩子是Usher」.
「有的是做招待」.
「有些孩子是Grader」.
「有的是做歡迎的」.
「這些都不是偶然喔」.
「這些是神在建立祂的教會」.
「這些是神在培養下一代的屬靈領袖」.
「所以弟兄姊妹們」.
「我們慕道園的孩子本身就是神蹟」.
「我們慕道園的小孩子本身就是神蹟」.

$^{681}$「他們不是被需要解決的問題」.
「他們是神賜給教會的未來和祝福」.
「這個是神賜給我們教會的未來和祝福」.
「好,接下來我們再看一看」.
「一起建立敬虔的下一代」.
「我們一起來看」.
「一起建立敬虔的下一代」.
「親愛的慕道園家人們」.
「親愛的慕道園家園們」.
「我想代表兒童事工」.
「也代表神的心」.
「我想代表兒童事工」.
「也代表神的心」.
「向你們發出一個重要的邀請」.
「向大家發出一個重要的邀請」.
「讓我們全教會一起參與建立下一代」.
「兒童事工並不是幾個主日學老師的事情」.
「兒童事工不是幾個主日學老師的事情」.
「也都不是幾個家長的需要」.
「兒童事工不是禮拜堂旁邊的幾個小房間或者一棟樓」.
「兒童事工是教會的使命」.
「是教會的未來」.
「是教會的心跳」.
「那有人會問了」.
「那為什麼全教會都需要投入呢?」.
「那你們或者會發問」.
「為什麼全個教會都要投入呢?」.
「有三個原因我想分享一下」.
「第一,因為孩子需要的是屬靈的家庭」.
「而不是一堂課」.
「第一,孩子需要的是屬靈的家庭」.
「而不是一堂課」.
「孩子每個星期日過來不是來上知識課」.
「他們來尋找什麼呢?」.
「他們在尋找什麼呢?」.
「他們來尋求安全感,歸屬感,愛,鼓勵,榜樣」.
「和屬靈的支持」.
「他們在尋找什麼呢?」.
「安全感,歸屬感,愛,鼓勵,榜樣」.
「和屬靈的支持」.

$^{721}$「孩子不會因為某段聖經故事而信主」.
「他們是因為感受到」.
「但是他們感受得到」.
「有人愛我」.
「有人聽我」.
「有人陪我」.
「有人知我」.
「有人懂我」.
「而認識耶穌」.
「而認識到耶穌基督的」.
「第二個原因」.
「因為孩子在觀察每一個大人」.
「你以為孩子沒有注意你嗎?」.
「他們的小眼睛時刻都在看著你們」.
「他們的小眼睛時時都在看著大家」.
「他們在看大人如何禱告」.
「大人如何敬拜」.
「大人如何愛神」.
「如何愛他人」.
「大人如何愛他人」.
「大人如何面對困難」.
「大人如何活出神的道」.
「孩子不是在聽你說什麼」.
「孩子是在看你們怎麼活」.
「孩子從成人的身上學習什麼叫基督徒」.
「第三個原因是沒有哪一個人」.
「太忙,太老,太沒用」.
「第三個原因是沒有哪一個人」.
「太忙,太老,太忙,太沒用」.
「我們來看看聖經」.
「摩西八十歲受呼召」.
「摩西八十歲才受呼召」.
「撒拉九十歲生以撒」.
「大衛是個牧羊童」.
「以利沙是農夫」.
「以利沙是個農夫」.
「而使徒彼得是個漁夫」.
「保羅曾經是迪基督的罪人」.
「神使用他們了嗎?」.
「神有沒有使用他們呀?」.

$^{761}$「當然」.
「因為神使用願意的人」.
「因為神會使用願意的人」.
「而不是完美的人」.
「而不是一個完美的人」.
「你不需要會教學」.
「你不需要懂得怎樣教他們」.
「只要願意愛」.
「願意陪伴」.
「願意支持」.
「好,可能大家又會問了」.
「那我應該怎麼參與呢?」.
「或者大家會問我怎樣參與呢?」.
「我又不會教」.
「那我能做什麼呢?」.
「我又不懂得教」.
「那我能夠做些什麼呢?」.
「但兒童事工有很多可以做的喲」.
「但是兒童事工有很多可以做的喲」.
「你可以禱告」.
「你可以為某一個孩子」.
「或者為某一個班級禱告」.
「你可以為某一個小朋友」.
「或者為某一個班級禱告」.
「禱告是屬靈戰場上最強的力量」.
「你還可以為孩子們講故事」.
「你可以為孩子們講故事」.
「大家知不知道」.
「我們的孩子有多麼喜歡聽故事?」.
「大家知不知道」.
「我們的孩子是多麼喜歡聽故事?」.
「每個星期我們的Uncle Joe」.
「來給孩子講故事的時候」.
「每個星期我們的Uncle Joe」.
「來跟孩子們講故事的時候」.
「孩子們瞪著眼睛認真地聽著」.
「你們或許不曾想過」.
「從你們口中所出的某一個話語」.
「會影響一個孩子」.
「甚至改變他的生命」.

$^{801}$「會影響一個孩子」.
「甚至改變他的生命」.
「你們還可以做助教」.
「做陪伴者」.
「或者你們可以做助教或者陪伴者」.
「不需要教學」.
「不需要教書的」.
「只需要維持秩序」.
「幫我們維持秩序」.
「陪孩子一起做手工」.
「陪小朋友一起做手工」.
「一起敬拜」.
「你的一個微笑」.
「你一個微笑」.
「一聲問候」.
「一個蹲下來對孩子說歡迎你來」.
「也可以讓孩子們感受到被愛」.
「你們還可以幫忙做點心佈置環境」.
「我知道有很多弟兄姊妹心靈手巧」.
「有些人可以做美味的點心」.
「有些人可以做精緻的小手工」.
「有些人可以做很精緻的小手工」.
「你們所有的恩賜都可以被神所使用」.
「你知道嗎?」.
「有些孩子就是因為我們教會美味的食物」.
「而每個星期準時出現」.
「你知道有些小朋友因為教會的美味食物」.
「每個星期會準時來到我們教會嗎?」.
「好,你們還可以成為屬靈的榜樣」.
「你們還可以成為屬靈的榜樣」.
「孩子們看著你們」.
「你們的生命就是最好的教材」.
「好,最後你們可能會問了」.
「如果我們現在投入」.
「未來將會如何呢?」.
「你或者會問」.
「如果我們現在能夠投入的時候」.
「教會未來會怎樣呢?」.
「如果我們今天和孩子們站在一起」.
「十年,二十年」.

$^{841}$「你們會看見」.
「你會看到」.
「他們會成為小組長」.
「他們會成為敬拜者」.
「他們會成為宣教士」.
「他們會成為牧者」.
「他們會成為職場榮耀神的領袖」.
「他們會成為成熟的基督徒父母」.
「他們會成為教會復興的力量」.
「所以現在對孩子的投資」.
「就是牧道園華人敬信會未來的樣子」.
「所以現在我們對小孩子的投資」.
「就是牧道園環浸信會未來的樣子」.
「親愛的弟兄姊妹」.
「親愛的弟兄姊妹們」.
「當我們服侍孩子」.
「當我們服侍小孩子」.
「當我們歡迎孩子」.
「我們不是在做一件次要的事情」.
「我們是在參與神最初的計劃」.
「我們是在跟隨主耶穌成長的腳步」.
「我們是跟隨著耶穌成長的腳步」.
「孩子是耶穌最在乎的人之一」.
「是神眼中寶貴的銅人」.
「如果教會愛孩子」.
「這就是神所喜悅的教會」.
「好,我們最後一起來低頭禱告」.
「親愛的天父,孩子」.
「恩賜給我們這樣美妙的清晨」.
「主啊,就請您來大大地使用」.
「我們慕道園華人敬信會」.
「使用我們每一個人」.
「使用我們每一個人的恩賜」.
「來帶領孩子認識您,親近您」.
「讓每一個孩子來到教會」.
「都能夠感受得到您屬天的愛」.
「都能夠感受得到安全感」.
「讓每一個孩子來到教會」.
「都能夠經歷您的祝福」.
「主啊,就請您親自地」.

$^{881}$「帶領我們每一個弟兄姊妹」.
「來服侍您所珍愛的孩子」.
「讓孩子們在我們弟兄姊妹的生命當中」.
「看到主耶穌基督」.
「以上禱告是奉主耶穌基督的聖名祈求」.
「Amen」.
\newpage



\section{羅馬書 13:8-14}
\label{sec:031634}
\textbf{得救比初信時更近了}~~~~袁惠鈞牧師博士
\newline
\newline
連結: \href{http://www.church.com.hk/acms/content.asp?site=cdc\&op=show\&type=product\&code=031634\&layout=sermon}{\texttt{031634}} ~~~~ 語音日期: 2026-02-08
\newline
\newline
\hyperref[sec:031633]{\small{< < < PREV SERMON < < <}}
~
\hyperref[sec:index_chronic]{\small{[返順時目]}}
~
\hyperref[sec:031635]{\small{> > > NEXT SERMON > > >}}
\newline
\newline
羅馬書 13:8-14
\newline
\begin{longtable}{cl}
\hline
\hline
章節 & 經文 (和合本修訂版)\\
\hline
13:8 & \begin{tabularx}{0.7\textwidth}{X} 你們除了彼此相愛，對任何人都不可虧欠甚麼，因為那愛人的就成全了律法。 \end{tabularx} \\ \\ \relax
13:9 & \begin{tabularx}{0.7\textwidth}{X} 那不可姦淫，不可殺人，不可偷盜，不可貪婪，或別的誡命，都包括在「愛鄰如己」這一句話之內了。 \end{tabularx} \\ \\ \relax
13:10 & \begin{tabularx}{0.7\textwidth}{X} 愛是不對鄰人作惡，所以愛就成全了律法。 \end{tabularx} \\ \\ \relax
13:11 & \begin{tabularx}{0.7\textwidth}{X} 還有，你們要知道，現在正是該從睡夢中醒來的時候了；因為我們得救，現在比初信的時候更近了。 \end{tabularx} \\ \\ \relax
13:12 & \begin{tabularx}{0.7\textwidth}{X} 黑夜已深，白晝將近。所以我們該除去暗昧的行為，帶上光明的兵器。 \end{tabularx} \\ \\ \relax
13:13 & \begin{tabularx}{0.7\textwidth}{X} 行事為人要端正，好像在白晝行走。不可荒宴醉酒；不可好色淫蕩；不可紛爭嫉妒。 \end{tabularx} \\ \\ \relax
13:14 & \begin{tabularx}{0.7\textwidth}{X} 總要披戴主耶穌基督，不要只顧滿足肉體，去放縱私慾。 \end{tabularx} \\ \\
[1ex]
\hline
\hline
\end{longtable}
$^{1}$羅馬書第35講,羅馬書第13章8至14節.
今日又要同大家提醒一下.
我講好多聖經的系列.
都是要幫大家一齊讀經.
如果你一個星期聽我五篇道的話.
你可以在兩年半讀完整部聖經.
而且你對聖經會有一個非常深入的認識.
因為我每一部聖經的系列.
第一講都是一個介紹.
是會講很多的背景.
給你有很多的資料.
能夠認識這一部書.
而第二個好處.
如果你跟著我們一起讀經.
聽我講道的話.
你會有一個一致解經的方法.
不會有矛盾.
因為我是用聖經解釋聖經.
是用字面解釋聖經.
這個是真正讀經的方法.
以致所有的經文都會非常一致.
完全沒有矛盾.
所以你想讀經讀到完全沒有矛盾.
又有深入認識的話.
希望你能夠跟我一起讀聖經.
今日我們要講的題目.
就是得救比初訊時更勤了.
為什麼保羅突然間會講這樣的訊息.
原因很簡單.
因為上一次我們才講過.
保羅教我們要服從政府.
無條件的來服從.
當時的軍王羅馬王是誰.
是尼祿一個非常殘暴的軍王.
而且他是北基督徒.
所以當保羅說要服從政府的時候.
他也是給很多羅馬的信徒.
有很大的鼓勵.
叫他們不要放棄.
你們得救比初訊時更勤了.

$^{41}$這個也是給我們現在的基督徒.
我不知道可能你在聽我這篇道.
你是在受很大的逼迫.
有很大的困難.
但是不用怕.
你得救比初訊時更勤了.
你很快就會達到你的目的地.
所以千萬不要放棄.
這個訊息也是給你們的.
我上一次也講過.
在羅馬書的第十三章一至七節.
有一個很精彩的首尾結構.
講到我們要服從政府.
現在十三章的八節至十四節.
都是有一個很精彩的首尾結構.
我們今天就要用這個首尾結構來解經.
而當我們是用首尾結構.
形成一個大綱來講這些經文.
就不容易解錯經文了.
這個是作者原本有的容易.
所以他才用這樣的首尾結構來寫出經文.
你說今天的首尾結構是什麼?.
一共有四層.
第一層就是十三章的八節和十三章的十四節.
十三章八節說我們要愛人完全律法.
十三章十四節說我們要披戴基督.
然後有很多事我們不會去做.
我們不會去犯罪.
為什麼十三章八節是講愛人.
十三章十四節是講披戴基督?.
披戴基督和愛人有什麼關係?.
等一下我們會來看.
第二層這個首尾結構.
就是十三章九節和十三章十三節.
都是講要守律法,不要去犯罪.
第三層十三章十一節說我們該趁早睡醒.
你們快點醒吧.
十三章十二節說八就將近天光.
這是什麼時候?.
在第四層整個首尾結構最中心的地方.

$^{81}$就是羅馬書十三章十一節.
說我們得救的日子越來越近.
十三章十一節尾說.
因為我們得救 現今比初信的時候更近了.
這個就是我們今天的主題.
就是保羅要鼓勵當時的基督徒.
這個也是給我們每一個在困難當中的基督徒的鼓勵.
所以我們今天大綱一共有四點.
第一點就是講我們要愛人.
我們要披上基督.
在羅馬書十三章八節和十三章十四節.
第二點就是我們要遵守律法.
十三章九節和十三節.
第三點就是時候近了.
羅馬書十三章十一節頭和十二節.
不要放棄 時候近了.
什麼近了?.
得救比初信的時候更近了.
羅馬書十三章十一節.
所以我們現在先來看第一層 第一點.
就是我們要愛人 我們要披上基督.
你說愛和披上耶穌基督有什麼關係?.
羅馬書十三章八節說愛人就完全了律法.
在聖經裡還告訴我們誰成全了律法.
耶穌基督 《馬太福音》五章十七節.
「莫想我來要廢掉律法.
我先知我來不是要廢掉乃是要成全」.
所以無論是愛或者是耶穌基督.
都是成全了律法.
這就是他們相似的地方.
另外羅馬書十三章十四節頭說.
總要披戴主耶穌基督.
英文就是「Put on Christ」.
在歌羅西書的三章十四節.
英文說要「Put on love」.
有沒有看到?.
「Put on Christ」和「Put on love」.
披上愛和披上基督.
雖然和合本歌羅西書的三章十四節.
說不定你看不到「穿上愛」這些字面.

$^{121}$但是原文的希臘文的確有這樣的意念.
所以我們要愛人.
我們要披上基督.
這就是羅馬書十三章八節.
「愛人的就完全了律法」.
十三章十四節.
「總要披戴主耶穌基督.
不要為肉體安排去放縱私慾」.
有沒有看到?.
不要為肉體安排.
不可以放縱私慾.
你不可以犯罪.
這個就是負面的要我們遵守神的律法.
正面是十三章八節.
「我們愛人就完全了律法」.
「成全了律法」.
所以這兩節的經文都是說一樣的東西.
就是我們今天手尾結構的第一層.
其實如果你注意的話.
耶穌基督在約翰福音的十五章.
說到葡萄樹和枝子的關係.
都有很相似的經文.
其實那裡也有一個手尾結構.
在約翰福音的十五章七節.
就是手尾結構的上半部.
耶穌基督說「你們若常在我裡面」.
我們怎麼能常在耶穌基督裡面?.
當我們披戴耶穌基督.
我們就時時在耶穌基督的裡面.
我們披戴著祂.
這個手尾結構的下半部分.
十五章十節尾就「常在我愛裡」.
有沒有看到?.
上面是說披戴基督.
下面是說愛.
耶穌基督的愛.
我們在祂的愛裡面.
又是披戴耶穌基督和愛.
這個手尾結構最中心的地方.
十五章的七節.

$^{161}$「我的話也常在你們裡面」.
剛才說我們愛人就完全了律法.
我們披戴耶穌基督.
我們就不犯罪.
這裡說什麼?.
如果我們常在耶穌基督裡面.
我們常在耶穌基督的愛裡面.
神的話語就在我們的裡面.
然後十五章的十節頭.
我們就會遵守耶穌基督的命令.
所以這裡所說的.
其實和羅馬書我們今天的經文.
羅馬書的十三章是很相似的.
所以有兩個部分.
一個部分就是我們要愛人.
另外一個部分我們要披上耶穌基督.
所以我們先說愛的部分.
羅馬書的十三章八節說.
「凡事都不可虧欠人.
唯有彼此相愛.
要常以為虧欠.
因為愛人的就完全了律法」.
有沒有看到這裡很簡單的一節經文.
我們看到有三點.
第一點就是不可虧欠人.
第二點是要彼此相愛.
第三點我們要完全律法.
所以我們先看第一點.
不可虧欠人.
羅馬書的十三章八節頭說.
「凡事都不可虧欠人」.
你說為什麼保羅無端端來到這裡.
就說不要虧欠人?.
我們可以看回前面的經文.
十三章的六節和七節是說什麼?.
保羅說我們要服從政府.
其中一樣服從政府的就是要納稅.
你不要虧欠政府.
你要納稅.
但是今天的經文也可以解釋.

$^{201}$我們不要虧欠任何人的債.
的確在歷史當中.
有很多很好的基督徒.
他們從來不會去借錢.
好像斯布珍.
好像內地會的戴德森.
克森·貝勒.
他們都不會借錢.
原因他們說.
就是羅馬書的十三章八節說.
「凡事都不可虧欠人」.
但是這一節的經文.
是不是真的這樣解釋呢?.
如果我們注意到其實在舊約.
聖經是有教我們關於借錢的律例典章.
例如在出埃及記的二十二章.
和申命記的二十三章.
經文說不可取利.
你借錢給人.
你不要去收貴利.
你不要做大耳窿.
另外申命記的十五章.
告訴我們.
有人借錢我們不可以忍著心.
我們要借給他們.
所以在摩西五經裡面.
有教我們關於借錢的律例典章.
在詩篇的三十七篇二十六節.
「他終日恩待人,借給人.
他的後裔也蒙福」.
蒙福的是甚麼人呢?.
借錢給人的人.
然後在述詩篇的三十七篇二十一節說.
「惡人借貸而不償還」.
蒙福的人就時時借錢給人.
惡人又怎樣呢?.
借了錢不還.
所以第一點.
我們一定要不要虧欠人.
如果欠人錢就要還錢.

$^{241}$千萬不要像劉備借荊州.
一借無回頭.
這裡說不要虧欠人的債.
我們剛才說要愛人.
第二點就是要彼此相愛.
經文說唯有彼此相愛.
要常以為虧欠.
羅馬書十三章八節尾.
剛才說不要欠人的債.
如果你要欠人的債.
你欠甚麼?.
欠人愛的債.
你好像還不清別人愛的債一樣.
你拼命去愛人.
好像還不清一樣.
這個就是保羅教我們.
那你說愛是甚麼?.
怎樣才能表達出愛?.
你問愛是否一種感覺?.
是否一種感動?.
你跟我說我愛你.
我又跟你說我愛你.
大家說得雞皮疙瘩.
是否這樣就是愛?.
當然不是.
我們看看聖經告訴我們愛是甚麼.
首先第一點.
愛是出於真理.
你說真是?.
是呀.
保羅在《敘哥林多後書》第六章.
他跟哥林多教會的人說.
我為了你們經歷了很多苦難.
你們鞭打,監禁,擾亂.
我勤勞,警醒,不食.
吃都沒得吃,為了你們.
在《敘哥林多後書》第六章六節.
保羅說.
我對你們有無謂的愛心.
你說這個愛心是怎樣?.

$^{281}$接著第六章七節保羅說.
我傳真實的道理給你們.
原來傳真實的道理給人.
就是有無謂的愛心.
原來愛是出自真理.
這是保羅告訴我們.
在《說耳乏》所書的四章十五節.
保羅說.
唯用愛心說誠實話.
凡事長進.
憐於元首基督.
有沒有看到.
我們要憑著愛心來講實話.
我們要憑著愛心來講真理.
因為愛是出於真理.
還有第二點.
愛是事後聖徒.
在希伯來書的六章十節.
經文說.
你們為他的命所顯的愛心.
就是先前事後聖徒.
如今還是事後.
有沒有看到事後出現了兩次.
愛是什麼?.
愛是事後聖徒.
格拉提書的五章十三節.
保羅也說.
弟兄們.
你們蒙召事要得自由.
只是不可將你們的自由.
當作放縱情慾的機會.
總要用愛心互相服侍.
有沒有看到.
希伯來書說.
有愛心就是事後聖徒.
格拉提書五章怎麼說.
我們要用愛心來服侍人.
愛是一種行動.
愛不是單單一種感覺.
事後或者服侍這個字.

$^{321}$是一個動詞.
不是一個名詞.
不是說我愛你我愛你.
而是真真正正有行為.
所以我們看到愛是什麼.
第一點.
愛是出於真理.
是教導人有關神的知識.
這是愛.
第二點.
愛也是事後聖徒.
事後聖徒.
還有第三點.
愛是活出美好的見證.
基爾弗所書的四章二節.
經文說.
凡事謙虛溫柔忍耐.
用愛心互相寬容.
這個就是有美好的見證.
在彼得前書的四章八節.
最要緊的是彼此切實相愛.
因為愛能遮掩許多的罪.
愛是什麼.
是包容別人的過犯.
愛是什麼.
是急於揭發別人的短處.
愛是什麼.
是隱藏別人的不是.
這裡我要解釋.
不是說愛我們就怎樣.
是包容別人的罪.
別人犯了罪你當作看不到.
不是的.
如果別人犯罪.
你是要跟他說清楚.
你要他面對自己的罪.
但是這裡是告訴我們.
當我們愛別人的時候.
我們不會去揭露別人的短處.
我們會隱藏別人無心的不是.

$^{361}$在書哥林多前書的十三章四到六節.
保羅告訴我們.
愛是恆久忍耐.
又有恩慈.
愛是不疾度.
愛是不自誇.
不張狂.
不作害羞的事.
不求自己的益處.
不輕易發怒.
不計算人的惡.
不喜歡不義.
只喜歡真理.
有沒有看到最後保羅說什麼.
愛是不喜歡不義.
只喜歡真理.
我們剛才第一點說愛是出於真理.
我們現在又回到我們第一點.
所以愛不是一種單單感情上的感覺.
愛是對人的責任.
去事奉人.
愛是虧欠了別人的愛.
去還債一樣.
還清債這樣來還.
這個才是愛.
那你說這種這樣的愛.
我哪裡有能力這樣去愛人.
我怎麼能夠做得到.
在羅馬書的五章五節尾.
因為所賜給我們的聖靈.
將神的愛攪拌在我們的心裡.
當我們每一個人重新得救的時候.
就有聖靈在我們的身上.
聖靈是會幫助我們.
我們成為一個新人是有新的能力.
現在我們有愛的債務.
我們要去愛人.
這個新的能力就能夠幫助我們.
去還債去愛人.
在我們得救之前.

$^{401}$我們沒有辦法可以做得到.
但是當我們得救之後.
靠著聖靈我們是能夠做得到.
我們剛才說愛人是三點.
第一不要虧欠人.
第二要彼此相愛.
我們也解釋了.
愛是出於真理.
愛是事奉聖徒.
愛是有好的見證.
現在我們說第三點關於愛.
就是愛承傳律法.
羅馬書的十三章八節尾.
因為愛人的就完全了律法.
在聖經裡有很多關於愛的命令.
保羅在菲勒比書的一章九節.
說要憑著愛心來做事.
在哥林多前書的十四章一節.
說我們要追求愛.
在提摩太前書的二章十五節.
說我們要常存愛心.
而保羅在鐵撒羅尼加前書的四章九節尾.
說你們自己以蒙了神的教訓.
叫你們彼此相愛.
有沒有看到?.
這裡我們得了神的教訓.
叫我們彼此相愛.
剛才說愛是什麼?.
是出於真理.
這裡又讓我們看到.
我們有了神的教訓.
才能彼此相愛.
還有使徒約翰.
他是被稱為愛的使徒.
所以在他的寫作裡.
都是充滿了愛的教導.
在約翰福音十三章三十五節.
耶穌基督說.
「你們若有彼此相愛的心.
眾人因此就認出你們是我的門徒了」.

$^{441}$我們能夠被人認出.
是否只用口說?.
不是.
認出什麼?.
是我們的行為.
是我們的行動.
彼此時候.
剛才說愛就是出於真理.
愛就是彼此時候.
約翰一書的二章十節.
約翰說.
「愛弟兄的就是住在光明中」.
一個真正得救的人在光明裡.
如果我們真的得救的話.
我們會愛我們的弟兄.
約翰一書的四章二十一節.
很相似的經文.
「愛神的也當愛弟兄」.
我們愛弟兄.
就是反映出我們愛神的心.
所以當我們愛神的時候.
我們也會愛我們的弟兄.
連彼得在彼得前書的一章二十二節也說.
「要愛弟兄沒有虛假.
從心裡彼此切實相愛」.
聖經裡甚至教我們要愛我們的敵人.
《馬太福音》第五章和《羅馬書》第十二章.
愛我們的敵人真的不容易.
所以我們要揀選去愛.
我們的敵人會怎樣?.
會欺負我們.
會侮辱我們.
會攻擊我們.
令我們很憤怒.
當我們憤怒的時候.
還能不能去愛人呢?.
很艱難.
所以我們真的要選擇去愛人.
這就是保羅在《羅馬書》十三章所教我們.
我們現在已經講完愛的部分.

$^{481}$剛才我們說有兩個部分.
愛和披戴基督.
《歌羅西書》三章十四節.
說我們要披上愛.
或者是「Put on love」.
我們剛才說.
現在《羅馬書》十三章十四節.
教我們要披上基督.
就是《羅馬書》十三節.
「總要披戴主耶穌基督」.
我們每個人生出來要披戴什麼?.
有人說要披戴尿片.
這樣就很現實了.
我們每個人生出來都要披戴尿片.
但是我們屬靈的方面.
我們每個人生出來就是披戴罪.
我們都是與神隔絕的.
所以如果我們要與神和好的話.
不是靠我們自己的行為.
是靠披上耶穌基督的義.
記不記得在《羅馬書》裡我們說.
怎樣能夠與神和好?.
就是我們接受耶穌基督.
然後披戴耶穌基督的義.
神就會接受我們與我們和好.
因為神看見我們.
祂只看見祂愛子的義.
而看不到我們的罪.
在哥林多後書的五章二十一節.
經文說.
「神使那無罪的替我們成為罪.
好叫我們在他裡面成為神的義」.
那無罪的替我們成為罪.
以致我們能夠披上耶穌基督的義.
記不記得在《馬太福音》二十二章.
有一個國王為自己的兒子擺設婚宴.
他叫他的僕人到處去大街小巷.
去請人來參加這個婚宴.
後來在婚宴裡.
發覺有一個人沒有穿禮服.

$^{521}$大件事.
王就叫他的僕人捆起這個人的手腳.
把他雕在外邊的黑暗裡.
在那裡必要哀哭切齒了.
我說過很多次.
當有人被扔到外面黑暗的地方.
哀哭切齒就是扔到地獄裡.
你說這個人來參加婚宴.
很明顯這裡有個比喻.
是參加羔羊的婚宴.
這個應該是一個得救的人.
如果不是為什麼會參加婚宴呢?.
這個禮服是代表什麼呢?.
是代表耶穌基督的義.
這個人沒有披上耶穌基督的義.
這個人是想靠自己的行為.
來參加這個君王的婚宴.
沒有可能的.
我們每個人都要披上耶穌基督的義.
很明顯這個不是一個真的得救的人.
所以被扔到外邊哀哭切齒.
在聖經裡其實有很多披上耶穌基督的象徵.
在羅馬書第六章.
當我們受洗歸入耶穌基督的死.
和祂一起同釘十字架.
一同被埋葬.
經文說.
如果我們和基督同死.
相信我們也必和祂同活.
我們向罪也.
當看自己是死的.
向神在基督耶穌裡.
卻當看自己是活的.
我們要當我們的舊人已經死了.
我們要脫去我們的舊人.
但是我們也需要穿上我們的新人.
有新的樣式.
我們需要披上耶穌基督.
世人叫我們做基督徒.
什麼意思?.

$^{561}$我們有耶穌基督的名字在我們身上.
就像我們披上耶穌基督一樣.
古時的拉比時說.
當猶太人敬拜神的時候.
他們就是穿上神的榮耀.
同樣當我們接受耶穌基督.
我們披上祂的義.
就是披上祂的榮耀.
披上祂的尊貴.
在聖經裡.
有很多披戴耶穌基督.
穿上新人的經文.
其實早在舊約已經有這樣的經文.
你說真的?.
以賽亞書61章第十節.
經文說:耶和華以拯救為衣給我穿上.
以公義為袍給我披上.
有沒有看到?.
這表示出將來耶穌基督將會為我們釘十字架.
祂的義將會遮蓋我們.
在新約也有加拉太書3章27節.
你們受洗歸入基督的都要披戴基督了.
另外有兩節的經文.
我們在這篇道裡會時時回來看這兩節的經文.
一節就是以弗所書的4章22節.
說我們要脫去舊人.
然後24節說要穿上新人.
另外一節是歌羅西書的3章9節.
說因為你們已經脫去舊人和舊人的行為.
歌羅西書的3章10節說.
所以我們要穿上新人.
我們會不斷回來看這兩節的經文.
因為和我們今天所說的息息相關.
所以聖經告訴我們.
我們要披上愛.
我們要披上耶穌基督.
剛才說愛有三點.
記不記得?.
愛是什麼?.
愛是出於真理.

$^{601}$第二點是愛是事後聖徒.
第三點是愛要有好的見證.
同樣披上基督.
都是有這三點.
所以聖經真的非常一致.
因為是由一位聖靈來寫.
所以是很一致.
我們披上基督.
都是出於真理.
為什麼?.
剛才就是我們所說的.
以弗所書的4章和歌羅西書的3章.
以弗所書的4章21節.
如果你們聽過他的道.
領了他的教.
學了他的真理.
他是誰?.
是耶穌基督.
你領了他的教.
你學了他的真理.
就會怎樣?.
22節.
就要脫去你們從前行為上的舊人.
然後24節要穿上新人.
有沒有看到?.
披上耶穌基督和愛一樣.
都是出於真理.
歌羅西書都是一樣.
3章10節說.
穿上了新人.
然後3章16節說.
當用各樣智慧把基督的道理.
豐豐富富藏在心裡.
所以披上基督.
就是出於真理.
這是第一點.
還有第二點.
披上基督是要事奉聖徒.
同樣都是這兩節的經文.
歌羅西書的3章10節.

$^{641}$和已忽所書的4章24.
3章10節.
歌羅西書說.
穿上了新人.
然後3章15節說.
你們為此夢照歸為一體.
喂!歸為一體?.
這一體是什麼?.
是教會的肢體.
是耶穌基督的身體.
是我們眾聖徒.
穿上新人是為了什麼?.
是要侍候聖徒.
然後在歌羅西書的3章11至12節.
說出如何侍候聖徒.
要全憐憫.
恩慈.
謙虛.
溫柔.
忍耐的心.
來彼此侍候.
在已忽所書的4章24節.
也說要穿上新人.
然後25節.
我們對待肢體.
告訴我們應該要怎樣做.
要使肢體得益.
要造就教會.
這就是侍候聖徒.
說了兩點.
披上基督是出於真理.
披上基督是要侍候聖徒.
第三點.
披上基督是要作好見證.
歌羅西書的3章10節.
又是這節經文.
要穿上了新人.
然後在3章11至25節.
保羅說出有很多不同的關係.
有什麼關係?.

$^{681}$夫妻的關係.
父母和子女的關係.
主人和僕人的關係.
為什麼?.
因為我們要有美好的見證.
在已忽所書的4章24節.
也說要穿上新人.
然後在4章25至32節.
保羅具體地告訴我們.
應該要怎樣做.
以致我們能夠活出美好的見證.
所以愛人和披上基督.
都是出於真理.
第一點.
第二點.
都是要來侍候聖徒.
第三點.
都是要我們有美好的見證.
當我們披上耶穌基督.
我們就會越來越像耶穌基督.
在彼得後書的1章3節.
說我們就得與神的聖情有份.
哇!神的聖情我們有份?.
是什麼?.
是神的義.
我們披上了耶穌基督的義.
但是要注意.
這只不過是地位上.
不是說我們屬靈的情況.
已經去到像耶穌基督一樣的聖潔.
這只不過是地位.
而保羅在以弗所書的4章1節說.
我們蒙召的人要與蒙召的恩相稱.
就是要與我們的地位相稱.
在羅馬書的13章14節.
保羅說.
總要披戴主耶穌基督.
注意「披戴」這個字的文法是現代式.
不是過去式.
很多人你去問他.

$^{721}$你是否基督徒?.
他會怎樣回答?.
我是基督徒.
我曾經某年某月某日決志的.
也可能他說.
我在某年某月某日曾經受洗.
你也會說曾經.
這些都是過去的事.
怎樣證明我們真是一個基督徒?.
我們是越來越像耶穌基督.
這是現在的事.
不是過去的事.
你不要只顧著你的過去.
你要問自己.
你現在是否越來越像耶穌基督?.
這才證明你是一個真正的基督徒.
保羅在加拉太書的四章十九節說.
我為你們好像受生產之苦.
直等到基督成形在你們的心裡.
基督成形在你們的心裡.
保羅說我教導你們.
我訓練你們.
直到基督成形在你們的心裡.
這就是神聖的路上所會發生的事.
我們會越來越像耶穌基督.
保羅也說在哥林多後書的三章十八節.
說我們這些接受主的人.
我們面上的面紗已經揭下來.
我們能夠明白神的話語.
得以看見主的榮光.
我們會變成主的形象.
榮上加榮.
如同從主的靈變成的.
什麼意思?.
榮上加榮就是由一個榮耀的地步.
去到另外一個榮耀的地步.
這就是越來越像耶穌基督.
我們每一個真正接受耶穌基督的人.
我們都應該是這樣.
我們說完第一點.

$^{761}$就是我們要愛人.
我們要披戴耶穌基督.
第二點.
保羅說我們要遵守律法.
羅馬書十三章九節和十三節.
剛才已經說過.
我們還要怎樣遵守?.
要好像欠了人債一樣.
還不清去愛人.
你說愛人和守律法又有什麼關係?.
關係就大了.
首先我們要愛神.
然後我們愛人的心.
就是反映出我們愛神的心.
如果我們真的愛神的時候.
我們就不會犯罪,得罪人.
這就是保羅告訴我們的.
其實保羅在這裡.
在十節當中.
保羅只說了四節.
然後他補充了.
我們要愛人如己.
其實全部神的律法.
都是包含在我們要愛神,愛人.
保羅在這裡說四節的律法.
都是人與人之間的律法.
第七節的不可姦淫.
第六節的不可殺人.
第八節的不可偷渡.
和第十節的不可貪婪別人的一切.
其實十節當中.
頭四節是和神的關係.
然後後面的六節是和人的關係.
保羅只說漏了第五節當孝敬父母.
和第九節不可作假見證.
而在這裡所說的.
四條與人相處的律例的次序.
好像和我們相熟的聖經有些不同.
保羅先說第七節不可姦淫.
然後才說第六節不可殺人.

$^{801}$其實這個次序和在七十字譯本中.
的申命記五章十七至二十一節是一樣.
七十字譯本就是向舊約的希伯來文.
翻譯成希臘文的第一部聖經.
其實很多新約的作者.
都很喜歡引用七十字譯本.
這裡告訴我們.
當我們真的愛神的時候.
我們就會愛人.
當我們愛人的時候.
我們就不會犯姦淫.
當我們犯姦淫是什麼意思?.
就是因為我們不夠愛神.
不夠愛人.
我們的慾望太強.
但我們的愛心太小.
有很多人都會說.
我沒有犯姦淫.
是呀,你可能沒有犯姦淫外在的罪.
但你不要太高興.
耶穌基督在數馬太福音第五章說.
「你們聽見有話說」.
你們聽見有人說不可殺人.
但當你們憎恨你們的弟兄的時候.
就是等於殺人.
耶穌基督也說.
「你們聽見有話說不可姦淫」.
但如果你們很淫慰地看女人的話.
你們就等於和她們犯了姦淫.
所以我們見到守律法.
不只是用外在的行為來守.
而是用我們的心來守.
甚至可以說是用我們的愛來守.
剛才說過.
當我們愛神愛人的時候.
我們就不會犯罪.
有一件事很重要.
很多人不去犯罪.
因為他們怕神.
怕神的懲罰.

$^{841}$如果你怕神的懲罰.
怕神的審判.
如果你這樣來守神的律法.
你就受不了.
我們要用愛來守神的律法.
這也是原因.
為何當神那些十戒比以色列人.
無論在《出埃及記》的20章.
和《申命記》的第5章.
敘述的時候.
神都說.
「我是耶和華你們的神.
曾將你們從埃及地為奴之家領出來」.
有沒有看到神在賦予十戒比以色列人的時候.
祂先說.
「我就是那位打救你們的神.
將你們從為奴之家領出來.
你們要愛我.
你們要回報我.
你們要回應我」.
所以我們應該要盡心盡意盡性來愛神.
很多時候我們都會盡心盡意盡性.
來愛我們的家人.
來照顧我們的家人.
就是因為我們愛他們.
如果你是請了一個傭人回來.
他會不會盡心盡意盡性來照顧你呢?.
不會的.
你看到他的時候.
他就會做得很好.
照顧人.
但如果你看不到的時候.
他就不會去做.
不過且過.
以色列人守神的律法.
就好像一個被請回來的傭人一樣.
不過且過.
所以他們不能夠真正守到神的律法.
聖經不是說.
你要盡心盡性盡意盡力來懼怕神.

$^{881}$你看清楚.
聖經是說什麼?.
你要盡心盡性盡意盡力來愛神.
當他們愛神的時候.
才能夠守得到神的律法.
耶穌基督也說.
「這百姓用嘴唇尊敬我.
心卻遠離我」.
我剛才也說.
如果你是因為懼怕去守神的律法.
可能是會有一些效果的.
但你沒有辦法可以做得好.
因為你一定要愛神.
你才能夠守得到神的律法.
所以十誡裡頭那四誡是同神的關係.
當我們是愛神的時候.
我們就不會有其他的神.
這個就是第一誡.
當我們是愛神.
我們就不會為自己雕刻偶像.
有很多其他的偶像.
例如錢財.
很多物質的東西.
這些都可以成為我們的偶像.
當我們是愛神的時候.
我們不會忘清祂的名字.
這個就是第三誡.
當我們是愛神的時候.
我們也會記得將自己分別為性.
其實原本的十誡.
第四誡是「當守安息日」.
「當守安息日」是唯一在十誡當中.
在新約再沒有提到.
那你說為什麼?.
是不是不用再遵守安息日?.
是的.
因為節旗和食物都是神給以色列人.
要將他們分別為性.
和他們身邊的外邦人分開.
以致他們不會去學習外邦人來拜偶像.

$^{921}$現在因為耶穌基督為我們十字架所成就.
祂已經拆毀了.
以色列人和外邦人中間的牆.
所以以色列人不再需要用節旗禮儀食物.
來將自己和外邦人分開.
保羅在羅馬書的十二章第一節說過.
我們要愛神將自己獻上當作活祭.
所以在這裡他沒有重複.
當我們愛神的話.
當我們守前面的四誡.
之後我們也會愛人如己.
就是利未記的十九章十八節.
早在舊約我們已經看到.
當我們愛神的話.
我們也會愛人如己.
這就是十誡後面的六誡和人的關係.
當我們愛人的時候.
我們會愛我們的父母.
孝敬我們的父母.
當我們愛人的時候.
我們就不會謀殺.
不會姦淫.
不會偷渡.
不會作假見證.
不會貪婪別人一切的事.
聖經有很多很多的律例典章.
但是所有的律例典章都可以總結為.
要愛神愛人.
這就是在馬太福音的二十二章.
耶穌基督告訴我們.
而在今天的經文我們也看到.
當我們愛人的時候.
亦都等於披戴耶穌基督.
完全是一樣的道理.
我們說完了我們要完全律法.
我們要守法.
這是正面的.
現在保羅是用負面來告訴我們.
一樣的東西.
我們要守律法不要犯罪.

$^{961}$但是從負面來看.
羅馬書十三章十四節.
不要為肉體安排去放縱私慾.
安排這個字就是提前去思考.
有預謀的.
為了滿足我們自己的慾望.
就去計劃犯罪.
保羅說你們不要這樣.
一個人犯罪之前.
首先他是思想上來犯罪.
然後就引起他的慾望.
然後在他的行為表現出來.
這樣就犯罪.
在詩篇的三十六篇第四節.
他在床上圖謀罪孽.
定義行不善的道.
不爭惡惡事.
在床上一睡醒就已經在籌備.
策劃怎樣去犯罪.
有時我們犯的罪是一時的衝動.
就好像大衛和拔士巴犯的姦淫.
但是我們千萬不要去策劃.
不要去籌備罪.
在《箴言》的二十四章第八節.
設計作惡的必稱為姦淫.
我們要披大既獨.
我們千萬不要去策劃.
不要去籌備犯罪.
在希臘文「方宴」這個字.
「Commos」這個字.
在加拉太書的五章二十一節.
和彼得前書的四章三節都有用詞.
就是一種狂歡的派對.
就是任何形式的淫亂.
就是用到這個字.
「Commos」這個字.
最初這個字是用在哪裡?.
是士兵打完仗海船回來.
慶祝他們的勝利.
有狂歡來慶祝.

$^{1001}$也可能是當時參加了運動會.
好像奧林匹克的運動會.
得到勝利又是來慶祝.
慢慢這個字就變成了負面.
是那些狂歡醉酒淫亂的聚會.
也代表俗世間所做的淫亂的事.
還有就是「不可好色邪道」.
「好色」這個字就是床.
一張床的床.
原文的希臘文就是「kaide」這個字.
在新約聖經一共出現四次.
在希伯來書的十三章四節.
婚姻人人都當尊重.
床也不可污穢.
因為苟合行人的人神別審判.
還有「不可尊敬嫉妒」.
「尊敬」這個字是Aries這個字.
是一個人非常自我.
他吃要多打要贏.
他不肯屈居第二.
人家拿到第一他就會嫉妒.
所以這裡說「不可好色邪道」.
「不可尊敬嫉妒」.
其實只不過是幾個例子.
最終保羅教我們千萬不要犯罪.
如果我們愛神愛人的話.
我們就不會犯罪.
一個人犯罪的問題在哪裡?.
我講過羅馬書第七章的時候.
我講過.
因為是我們的肉體.
我們的肉體未被救贖.
是人性和罪性所在的地方.
我們仍然是有犯罪的傾向.
我們的身體是包括了我們的心思.
我們的意念和我們的動機.
一個真正不救的人.
我們是有新的本性.
我們是一個新人.
但是很可惜.

$^{1041}$我們仍然住在我們這個罪身裡面.
仍然未救贖的這個罪身.
我們的靈被救贖.
但我們的身體未被救贖.
有一天當耶穌基督再次回來.
我們與耶穌基督在半空相遇.
我們這個取死的身體.
就會搖身一變.
變成榮耀的身體.
就好像羅馬書的八章二十三節.
和哥林多前書的十五章告訴我們.
但是在這個之前.
我們應該要怎麼做?.
就好像哥林多前書的九章二十七節.
保羅說:我是攻克己心.
叫身服我.
恐怕我全福音級別人.
自己反被棄絕了.
保羅說:在耶穌基督回來之前.
我要攻克己心.
我要叫我的身體來服我.
如果不是的話.
我雖然叫人家是得救.
反而我自己可能不是一個得救的人.
我講了兩點.
第一點.
我們要愛人.
要披上耶穌基督.
第二點.
我們要遵守神的律法.
第三點.
就是告訴我們原因.
為什麼.
因為時候將近了.
羅馬書十三章十一節頭和十二節.
十三節頭說.
「在這你們曉得.
現今就是該趁早睡醒的時候」.
十二節說.
「黑夜已深.

$^{1081}$白晝將近.
我們就當脫去暗媒的行為.
帶上光明的兵器」.
這裡保羅給我們一個軍事上的命令.
要穿上基督.
非常迫切性.
保羅給了我們三個動詞.
第一是要醒來.
第二是要脫去.
第三就是要穿上.
時間越來越近了.
我們立即要採取行動.
事後無多了.
是很緊張的時候.
你不要自我滿足.
不要還在這裡休閒.
貪訓.
要醒來.
天亮了.
這就是保羅告訴我們.
很多人說.
用不著這麼緊張.
讓我再享受一下吧.
希伯來書十章二十四至二十五節說.
「又要彼此相顧.
激發外心.
勉勵行善.
你們不可停止聚會.
好像那些停止慣了的人.
都要彼此勸勉.
既知道那日子臨近.
就更當如此」.
希伯來書十章三十七節.
「因為還有一點點的時候.
那要來的就來.
並不遲」.
有很多人說.
牧師.
教會說耶穌基督會來.
已經說了二千年了.

$^{1121}$耶穌基督還沒來.
用不著這麼緊張吧.
可能祂永遠都不回來呢.
不如我們隨波逐流吧.
賺錢享樂一下吧.
顧著我們的事業.
顧著我們有舒適的生活.
顧著我們的財物.
最重要的是有房子有車.
有物質上的東西.
可能還有很多時間.
不用這麼緊張.
耶穌基督怎麼說呢.
祂說「洪水以前的日子.
人照常吃喝嫁娶.
直到羅亞進方舟的那日.
不知不覺洪水來了.
把他們全都沖去.
人子降臨也要這樣」.
《馬太福音》的二十四章.
三十八至三十九節.
我們要警醒.
要好像十個童女當中的五個童女.
不要等新郎來到.
發覺我們燈裡沒有油.
在等待耶穌基督的時候.
我們要潔淨自己.
我們要用神的話語來潔淨自己.
最後一點.
就是「得救近勤了」.
羅馬書的十三章十一節.
就是羅馬書的十一章尾說.
「因為我們得救.
現在比初信的時候更勤了」.
喂,你說,牧師.
你現在在說什麼.
你不是告訴我們.
我們全部都已經得救了嗎.
接受耶穌基督的人.
是的,但是得救有三個階段.

$^{1161}$第一個階段就是清二.
是過去式的.
不是因為我們的行為.
是因為我們接受耶穌基督.
被神稱我們為二人.
不是因為我們的行為.
就好像以弗所書的二章八節.
「我們得救是本乎恩恩著信.
不是靠行為免得由人自誇」.
所以第一個步驟就是過去式的.
是清二.
第二個步驟是成聖.
是現在的.
剛才說清二不是因為我們的行為.
但是成聖.
我們的行為就表現出.
我們真是一個清二的人.
是一個真正得救的人.
所以保羅說.
在庶於肥納比書的二章十二節.
「作成你們得救的功夫」.
是有功夫的.
我們的功夫就表現出來.
我們是一個真正得救的人.
這個也是我剛才所說的.
披上耶穌基督.
和成聖越來越像耶穌基督.
還有第三點.
是得榮耀.
這個是將來.
還未到的.
我們現在就在等.
我剛才說我們有取死的身體.
我們在等耶穌基督回來.
以致我們有榮耀的身體.
這個就是我們今天的主題.
羅馬書十三章十一節尾.
「我們得救然今比初信的時候更勤了」.
距離那個日子.
越來越近了.

$^{1201}$阿摩斯書的五章十八節.
「賞望那耶和華日子來到的有禍了.
你們為何賞望耶和華的日子.
那日黑暗沒有光明」.
喂!牧師你又說我們要盼望.
我們得救.
現在為何阿摩斯書裡面.
阿摩斯才說.
盼望這個日子的人有和.
喂!當時有很多以色列人.
他們希望主的日子快些來臨.
但他們根本沒有準備好自己.
他們以為主的日子來臨.
就會帶給他們榮耀.
為甚麼?進入千禧年國.
進入大衛的王朝.
再一次興起大衛的王朝.
他們以為是帶給他們榮耀.
但他們沒有準備好.
所以當他們沒有準備好的時候.
這個日子來臨.
就是審判他們的日子.
所以我們都要很小心.
在結束的時候.
我們不妨看看.
《使徒行傳》的一章六節.
門徒問耶穌基督.
主啊!你復興以色列國.
就在這時候嗎?.
喂!有很多人又說沒有千禧年國.
門徒也問.
復興以色列國是何時?.
就是大衛的王朝.
實實在在地上的國.
耶穌基督如何回答?.
祂說:我回來的日子.
不是你們可以知道的.
但聖靈降臨在你們的身上.
你們就必得著能力.
並要在耶路撒冷.

$^{1241}$猶大全地和撒瑪利亞.
直到地極作我的見證.
《使徒行傳》的一章八節.
說完了之後.
就有一朵雲彩.
接了耶穌基督去.
門徒再也見不到祂了.
於是他們怎樣?.
個個都向天來望.
有兩個天使走過來說.
喂!你們張開嘴巴.
開個洞望著天做甚麼?.
或者離開你們被接升天的耶穌.
你們見祂怎樣往天上去.
祂還要怎樣回來?.
等待耶穌基督再次回來.
就成為我們基督徒的動力.
來傳福音.
將福音由耶路撒冷傳到猶大全地.
撒瑪利亞直到地極.
我們有沒有現在在傳福音?.
有人說耶穌基督說回來已經二千年了.
為甚麼我們要相信祂會回來?.
因為耶穌基督第一次的回來.
聖經裡面有很多的寓言.
全部都已經實驗了.
例如有關耶穌基督的出生.
祂的生平,祂的死,祂的復活.
全部都實驗了.
所以我們能夠肯定.
祂第二次回來的寓言一樣是這麼準確.
我可以告訴你.
每一個寓言有關耶穌基督第一次的來臨.
就有八個寓言關於祂第二次的來臨.
所以如果祂第一次的來臨都是這麼準確的話.
第二次的來臨就會更加準確.
就會更加確實.
所以我們千萬不要輕心.
千萬不要放棄.
我們要等待因為我們得救的日子.

$^{1281}$比我們初信的時候更緊了.
願神祝福祂自己的話語.
讓我們低頭禱告.
讓我們看到保羅怎樣來鼓勵信徒.
是,他教我們要服從政府.
但是他也知道當我們服從政府的時候.
會有很多困難,艱難,苦難臨到我們的身上.
所以他給信徒有鼓勵.
在這段時間我們等待的時候.
我們要愛神愛人.
我們要披上耶穌基督.
我們要遵守你的律例典章.
我們要做那美好的見證.
以致當你回來的時候.
我們能夠被你稱為一個忠心良善的僕人.
將一切榮耀頌讚歸給你.
我們今天禱告.
奉主耶穌基督的名義而求.
阿們.
\newpage



\section{路加福音 1:20}
\label{sec:031635}
\textbf{聖經人聖經事:這話必然應驗}~~~~黃克勤牧師博士
\newline
\newline
連結: \href{http://www.church.com.hk/acms/content.asp?site=cdc\&op=show\&type=product\&code=031635\&layout=sermon}{\texttt{031635}} ~~~~ 語音日期: 2026-02-01
\newline
\newline
\hyperref[sec:031634]{\small{< < < PREV SERMON < < <}}
~
\hyperref[sec:index_chronic]{\small{[返順時目]}}
~
\hyperref[sec:031636]{\small{> > > NEXT SERMON > > >}}
\newline
\newline
路加福音 1:20
\newline
\begin{longtable}{cl}
\hline
\hline
章節 & 經文 (和合本修訂版)\\
\hline
1:20 & \begin{tabularx}{0.7\textwidth}{X} 到了時候，這些話必然應驗；只因你不信我的話，你會成為啞巴，不能說話，直到這些事實現的日子。」 \end{tabularx} \\ \\
[1ex]
\hline
\hline
\end{longtable}
$^{1}$各位電影姐妹,大家平安.
又到時間一齊讀聖經了.
生命樹書教網絡.
和大家一起看聖經人,聖經事.
聖經裡面有些故事.
記得可以看到所有在背後的工作.
讓我們認識主更加多.
今次的主題是.
「這話必然應驗」.
仍然是讀跟屍體有關的東西.
我自己預備的時候.
都沒有發覺,以前沒有留意.
原來司禧約翰是這麼多的蝙蝠.
在聖經裡面.
司禧約翰原來是一個很偉大的人物.
無論靈性上.
他自己的位身.
上主對他的體重.
原來都是非常好的.
是一個很難得,很少有這樣的人.
為主做先鋒.
但是自己很謙卑.
今天看看經文.
在路加一章十八節.
撒迦利亞就是司禧約翰的父親.
撒迦利亞這位祭司.
他在聖所裡面公職.
剛剛他是亞比亞班的其中一位祭司.
很多位的.
大家抽籤就輪到.
可能是上主的安排.
抽到他.
於是他就進了聖所裡面.
在那裡祈禱.
忽然之間.
在聖所裡面禱告.
在香壇放雨香禱告的時候.
忽然間.
見到有天使站在香壇旁邊.
嚇了他一跳.

$^{41}$就跟他說.
你的祈禱蒙英韻.
你會生一個兒子.
撒迦利亞就在這裡.
十八節對天使說.
撒迦利亞其實是一個很老實的人.
他很正直.
他就跟天使說.
我憑著什麼可知道這事呢?.
我會生一個兒子?.
我怎麼知道?憑什麼知道?.
第一,我已經老了.
我妻子二十年多老邁了.
原來.
他這麼直言.
上帝說你會生一個兒子.
叫天使跟他說.
他說.
會不會呀?.
我老了.
我太太老了.
哪裡會得?.
我憑什麼可以知道這事呢?.
其實他的意思就是.
哪裡有可能呀?.
我老了,我太太也老了.
因此你見到下一句.
是有點生氣的.
天使回答說.
你知不知道我是誰?.
我是站在神面前的加伯列.
即是說.
天使裡面是很崇高的一位.
站在上主旁邊.
他有什麼都叫我說的.
是上帝叫我來的.
我不是自己編的.
加伯列非常厲害.
你看《大耶利書》.
就是大耶利誠心向上主禱告.

$^{81}$我看《耶利米書》.
就知道我們回歸祖國.
是差不多時間了.
禱告神.
你指示我們關於以色列民會怎樣?.
然後加伯列晚上就來了.
加伯列就說.
我是在神面前的加伯列.
神叫我來,我就急忙飛來告訴你.
很厲害的.
神有些很重要的話.
是會跟加伯列說.
然後叫他去傳話的.
加伯列也傳訊息給瑪利亞.
說你會生一個兒子.
他就稱為至高者的兒子.
又是加伯列.
所以我是站在神面前的加伯列.
你知不知道我是誰?.
我是奉差而來對你說話.
將這個好消息傳給你.
你竟然說這種話.
竟然不相信.
二十字.
到了時候,這話必然應言.
不過,只因你不信.
你會成為啞巴.
不能說話,直到這是成就的日子.
百姓就在外面等撒迦利亞出來.
為什麼在聖殿內這麼久也不能出來?.
什麼事?沒什麼吧?.
你看到剛才這一節.
如果天使加伯列生氣.
也代表上帝生氣.
我是站在神面前的加伯列.
奉差而來對你說話.
將這個好消息傳給你.
你竟然不相信.
所以要罰你.
你要成為啞巴.

$^{121}$不能說話,直到這是成就.
但是「你會生個兒子」這件事.
必然會應驗.
外面的人等了很久還沒出來.
終於出來了.
沒東西買了這麼久,是不是不舒服?.
不能和他說話.
說不出話來.
他們就知道差在電裡見了異象.
因為他一直向他們打手勢.
警醒了啞巴.
他出來的時候.
沒事吧?祭司?撒迦利亞?.
撒迦利亞說不出話來,被打手勢.
說不出話來,說不出話來.
他們就知道難怪這麼久了.
他在裡面一定有異象見到了.
成為啞巴.
他公職的日子滿了,就回家了.
我跟大家分享一下,有三點跟大家看的.
第一,我憑著什麼可知道寫事呢?.
這一句是什麼意思呢?.
就說過了,怎麼會成功?.
很多時候,怎麼會成功?.
為什麼會這樣問?.
因為照人來看不行.
我太太老了,怎麼可以生育呢?.
生育了的,已經生了.
怎麼可以成功?.
很多時候,我們無論事奉,信主,屬靈的經歷.
很多時候都遭遇這個問題.
說出來當然容易.
但是遇到這個事情的時候.
明明是人去想不行的.
我衡量是不可能的.
所以就會問.
憑什麼可以知道寫事呢?.
好東西.
我是這麼想的.
我就說出來,他很正直地說.

$^{161}$我怎麼知道可以成功呢?.
其實是不可能的.
所以第一點是想說.
是不可能的.
不可能的.
但是你要知道.
在神面前第二點就是.
必定應驗是因為上帝說的.
不過很多人聽到必應驗.
他的心都不服的.
都認為是上帝這麼說而已.
哪裡會成功呢?.
不信.
你不信.
不信於是成為啞巴.
很有趣.
不信成為啞巴.
當我們要為主做見證.
要為感要為感.
講聖經.
但有時如果你本身不信.
其實你就像一個啞巴.
你都會變成一個啞巴.
第一,你慢慢不敢說.
那些信息不敢說.
很多時候傳道人都避了.
不會講《創世記》第一章.
為什麼呢?.
因為怕說出來會被人挑戰.
一天就做一樣東西.
這麼厲害?.
但是進化論你怎麼看?.
你怎麼解釋這些事情?.
解釋不了.
不敢說.
他很少說.
為什麼?.
因為他怕被人問.
他的不信會令他成為啞巴.
不敢說《創世記》第一章.

$^{201}$又或者是.
他不想奉獻兩個養生的小錢.
所以他一定不會說.
寡婦獻上兩個養生小錢那段聖經.
因為他不能說.
因為他不信那段經文.
不信他就會成為啞巴.
很有趣.
所以.
亞伯列是罰他.
你會成為啞巴.
是真的.
因為你的不信你成為啞巴.
大家要小心.
上帝的話語一定應驗.
這些日子以後.
他的妻子以利沙白懷了孕.
隱藏了五個月.
主在眷顧我的日子.
這樣看待我.
將我在人間的羞恥除掉.
隱藏了五個月的意思.
不是說二十一懷孕.
快點躲起來五個月.
不讓別人見.
懷孕之後.
這個隱藏是指他.
很小心地休養生息.
不去周圍活動.
過一個更平淡的生活.
如果有機會.
我們去到耶路撒冷.
有時間的話.
我們會去一個山.
旅遊班去到然後走上山.
走到很累.
也挺高的.
到半山有一間修院.
那間修院相傳是.
施洗約翰的家.

$^{241}$就是撒迦利亞和以利沙白的家.
那裡有教堂.
有些古蹟.
走得很辛苦上去.
我想平時也沒有什麼人會上去.
他就在那裡.
回到他住的地方.
不出來了.
幾個月休息.
我推測是很小心養胎.
因為主眷顧我.
除掉我在人間的羞恥.
終於給了我一個兒子.
把我在人間的羞恥除掉.
意思是他給了我一個兒子.
除掉我人間的羞恥.
所以我要好好地去保存.
要小心.
這樣的意思.
好了.
今日的金句.
路加福音一章二十節.
多了時候者說必應矣.
不信你成為啞巴.
直到這件事成就的日子.
再留意多說一次.
我們要對於聖經有信心.
對於神的話語有信心.
當你不信.
你就不能說那些經文.
是很真的.
有些傳道人很少說復活.
為什麼?.
因為說復活又怕被人問.
說創世記日一章又怕被人問.
說很多事他都怕被人問.
或者他自己都不太相信那段經文.
他就會在那些經文上成為啞巴.
當你自己的信仰都不清楚的時候.
你又會成為啞巴.

$^{281}$因為你信得不夠深.
不過你要知道.
我們信不信不要緊.
上帝的話語必然應驗是真的.
啟示錄說出來哪有人信的.
不過當你信的時候.
當你研究內心看的時候.
嘩!你發覺.
用不著這麼準確.
用不著這樣應驗.
真是現在正在發生.
你見到是很震驚的.
在神面前震驚.
而且直到這事成就的日子.
那日子會來到.
有一天會來到.
因此我們信任上帝.
聖經裡面見到不少事情.
是照著上帝的說話.
一步一步一步一步應驗.
今天我們看主耶穌回來的日子很近了.
你要有一個信心.
就是主會回來.
祂還沒回來.
差不多我都會去見祂了.
事情會成就的.
祂是為我們開了壽身誘惑的路.
去到父神面前.
而且我們跟祂在一起.
這是一個很大的信心保證.
好了,我今天跟大家分享到這裡.
就是神的話必然應驗.
就是今天的主題.
上帝說得出就做得到.
而且祂有能力做得到.
願主給我們更大的信心.
我們一起禱告吧.
天父上帝.
我們跟你真的沒得比.
你滿有能力.

$^{321}$無所不能.
不過我們可能是.
從創世的時候開始.
亞當夏娃辭職的時候開始.
就被蛇引誘.
我們經常都在想.
神喜事真說.
是不是真的?.
少少超乎了我們的思考力.
我們都很被人民主義去影響.
過了我們的理性.
我們就不能夠相信了.
怎麼能夠有這件事呢?.
憑什麼知道這件事呢?.
於是我們會向主發出挑戰.
不過你夠高.
我們有時候不敢挑戰.
就藏在心裡面.
但是你自己的使者加巴列就說.
你不信,你就會成為阿爸.
有時候自己反省一下.
很多事情我們自己都信得不夠.
我就不敢說了,成為阿爸.
但是.
雖然我們軟弱.
不過你是剛強的.
你是強壯的.
以致你說的話是必然應言.
讓我們有這個信心.
讓我們有這個醒覺.
更多信靠你的話語.
我們祈禱是奉靠耶穌基督的名字.
阿們.
\newpage



\section{約書亞記 4:1-24}
\label{sec:031636}
\textbf{帶著過去往前走}~~~~魏文俊傳道
\newline
\newline
連結: \href{http://www.church.com.hk/acms/content.asp?site=cdc\&op=show\&type=product\&code=031636\&layout=sermon}{\texttt{031636}} ~~~~ 語音日期: 2026-02-08
\newline
\newline
\hyperref[sec:031635]{\small{< < < PREV SERMON < < <}}
~
\hyperref[sec:index_chronic]{\small{[返順時目]}}
~
\hyperref[sec:031637]{\small{> > > NEXT SERMON > > >}}
\newline
\newline
約書亞記 4:1-24
\newline
\begin{longtable}{cl}
\hline
\hline
章節 & 經文 (和合本修訂版)\\
\hline
4:1 & \begin{tabularx}{0.7\textwidth}{X} 當全民都過了約旦河，耶和華對約書亞說： \end{tabularx} \\ \\ \relax
4:2 & \begin{tabularx}{0.7\textwidth}{X} 「你要從百姓中選出十二個人，每支派一人， \end{tabularx} \\ \\ \relax
4:3 & \begin{tabularx}{0.7\textwidth}{X} 吩咐他們說：『你們從這裡，從約旦河中祭司的腳穩穩站立的地方，取十二塊石頭，一起帶過去，放在你們今夜住宿的地方。』」 \end{tabularx} \\ \\ \relax
4:4 & \begin{tabularx}{0.7\textwidth}{X} 於是約書亞召集了他從以色列人中所選的十二個人，每支派一人。 \end{tabularx} \\ \\ \relax
4:5 & \begin{tabularx}{0.7\textwidth}{X} 約書亞對他們說：「你們要過去，到約旦河中，耶和華－你們神的約櫃前面，按以色列人支派的數目，每人各取一塊石頭扛在肩上。 \end{tabularx} \\ \\ \relax
4:6 & \begin{tabularx}{0.7\textwidth}{X} 這些石頭在你們中間將成為記號。日後，你們的子孫問你們說：『這些石頭對你們有甚麼意思呢？』 \end{tabularx} \\ \\ \relax
4:7 & \begin{tabularx}{0.7\textwidth}{X} 你們就對他們說：『這是因為約旦河的水在耶和華的約櫃前中斷；約櫃過約旦河的時候，約旦河的水就中斷了。這些石頭要作以色列人永遠的紀念。』」 \end{tabularx} \\ \\ \relax
4:8 & \begin{tabularx}{0.7\textwidth}{X} 以色列人就照約書亞所吩咐的做了。他們按以色列人支派的數目，從約旦河中取了十二塊石頭，正如耶和華所吩咐約書亞的。他們把石頭帶過去，到他們所住宿的地方，就放在那裡。 \end{tabularx} \\ \\ \relax
4:9 & \begin{tabularx}{0.7\textwidth}{X} 約書亞另外把十二塊石頭立在約旦河的中間，在抬約櫃祭司的腳站立的地方；直到今日，石頭還在那裡。 \end{tabularx} \\ \\ \relax
4:10 & \begin{tabularx}{0.7\textwidth}{X} 抬約櫃的祭司站在約旦河的中間，直到耶和華命令約書亞告訴百姓的一切事辦完為止，正如摩西所吩咐約書亞的一切話。於是，百姓急速過了河。 \end{tabularx} \\ \\ \relax
4:11 & \begin{tabularx}{0.7\textwidth}{X} 全體百姓都過了河之後，耶和華的約櫃和祭司才過去，到百姓的前面。 \end{tabularx} \\ \\ \relax
4:12 & \begin{tabularx}{0.7\textwidth}{X} 呂便人、迦得人、瑪拿西半支派的人都照摩西所吩咐他們的，帶著兵器在以色列人的前面過去。 \end{tabularx} \\ \\ \relax
4:13 & \begin{tabularx}{0.7\textwidth}{X} 約有四萬帶兵器的軍隊在耶和華面前過去，到耶利哥的平原，準備上陣。 \end{tabularx} \\ \\ \relax
4:14 & \begin{tabularx}{0.7\textwidth}{X} 在那日，耶和華使約書亞在以色列眾人眼前被尊為大。在他一生的年日中，百姓敬服他，像從前敬服摩西一樣。 \end{tabularx} \\ \\ \relax
4:15 & \begin{tabularx}{0.7\textwidth}{X} 耶和華對約書亞說： \end{tabularx} \\ \\ \relax
4:16 & \begin{tabularx}{0.7\textwidth}{X} 「你吩咐抬法櫃的祭司從約旦河上來。」 \end{tabularx} \\ \\ \relax
4:17 & \begin{tabularx}{0.7\textwidth}{X} 約書亞就吩咐祭司說：「你們從約旦河上來。」 \end{tabularx} \\ \\ \relax
4:18 & \begin{tabularx}{0.7\textwidth}{X} 抬耶和華約櫃的祭司從約旦河中上來，腳掌一落乾地，約旦河的水就流回原處，仍舊漲滿兩岸。 \end{tabularx} \\ \\ \relax
4:19 & \begin{tabularx}{0.7\textwidth}{X} 正月初十，百姓從約旦河上來，就在耶利哥東邊的吉甲安營。 \end{tabularx} \\ \\ \relax
4:20 & \begin{tabularx}{0.7\textwidth}{X} 約書亞把他們從約旦河取來的那十二塊石頭立在吉甲， \end{tabularx} \\ \\ \relax
4:21 & \begin{tabularx}{0.7\textwidth}{X} 對以色列人說：「日後，你們的子孫問他們的父親說：『這些石頭是甚麼意思呢？』 \end{tabularx} \\ \\ \relax
4:22 & \begin{tabularx}{0.7\textwidth}{X} 你們就讓你們的子孫知道，說：『以色列人曾走乾地過這約旦河。』 \end{tabularx} \\ \\ \relax
4:23 & \begin{tabularx}{0.7\textwidth}{X} 因為耶和華－你們的神在你們前面使約旦河的水乾了，直到你們過來，就如耶和華－你們的神從前在我們前面使紅海乾了，直到我們過來一樣， \end{tabularx} \\ \\ \relax
4:24 & \begin{tabularx}{0.7\textwidth}{X} 要使地上萬民都知道，耶和華的手大有能力，也要使你們天天敬畏耶和華－你們的神。」 \end{tabularx} \\ \\
[1ex]
\hline
\hline
\end{longtable}
$^{1}$各位電影姐妹平安.
今日的經文是接著上星期.
方牧師所說的以色列人過河的經文.
所以是第四章一至二十四節.
這段經文是神透過一群人.
為整個以色列的民族設立恩典的記號.
在這段經文裡的領受.
盼望在分享當中.
和大家領受同一個祝福.
就是我們不單單個人.
或者是我們成群的群體.
我們領受了神的豐盛恩典的時候.
仍然記得那些恩典的記號.
我們勇敢地走向未知的前路.
經歷更多神的恩典.
為的是讓社區 讓世界得著 祝福 造就.
特別是承振在建堂的進程裡.
讓我們在一陣經文的思考裡.
去紀念神過去八十多年的帶領.
我們去學習倚靠神迎向未知的日子.
好讓神能夠用我們更加多.
使社區得著 祝福 造就.
整段經文我會分兩個部分和大家一起看.
大家讀經 如果打開聖經可以留意.
特別是第一部分是第一節到第九節.
接著會跳到第十九節.
應該是二十四節才對.
這部分是說.
為什麼神要叫以色列人.
把這些石頭放在河的裡面和河的另一邊.
這些石頭如何成為以色列人恩典的記號.
第二部分就是中間.
第十節到十九節和二十四節.
第二部分是說.
我們看這些恩典的記號.
如何成為以色列人的鼓勵和提醒.
他們成為他們行前的動力.
也嘗試在當中.
看看對我們有什麼提醒和鼓勵.
我會用兩個英文字母.

$^{41}$用一個英文字母P.
兩個英文字送給大家.
去總結這兩個部分.
第一個是P是Past.
是一個名詞 過去.
過去代表神過去的恩典.
神去吩咐以色列人.
去放這些石頭做記號.
其實是想他們整個民族.
包括他們的子孫.
都記得神是曾經施行大能拯救的神.
首先我想邀請大家讀出這部分經文.
看到我們一起讀吧.
準備起.
於是約書亞和他們.
以色列人中有預備的那十二個人.
每支派一人都召了來.
對他們說.
你們下約旦河中.
過到耶和華你們神的約櫃前頭.
按著以色列人十二支派的數目.
每人取一塊石頭放在肩上.
這些石頭在你們中間可以作為證據.
日後你們的子孫問你們說.
這些石頭是什麼意思.
你們就對他們說.
這是因為約旦河的水.
在耶和華的約櫃前斷絕.
越過約旦河的時候.
約旦河的水就斷絕了.
這些石頭要作以色列人永遠的紀念.
以色列人就照約書亞所吩咐的.
按著以色列人支派的數目.
從約旦河中取了十二塊石頭.
都遵耶和華所吩咐約書亞的行了.
他們把石頭帶過去到他們所住宿的地方.
就放在那裡.
約書亞另把十二塊石頭.
納在約旦河中.
在台約櫃的祭司腳站納的地方.

$^{81}$直到今日那石頭還在那裡.
接著是十九至二十四節.
正月初十日.
百姓從約旦河裡上來.
就在吉甲在耶利哥的東邊安營.
他們從約旦河中取來的那十二塊石頭.
約書亞就納在吉甲.
對以色列人說.
日後你們的子孫問他們的父親說.
這些石頭是什麼意思.
你們就告訴他們說.
以色列人曾走過這約定.
因為耶和華你們的神.
在你們前面使約旦河的水乾了.
等著你們過來.
就如耶和華你們的神.
從前在我們前面使紅海乾了.
等著我們過來一樣.
要使地上萬民都知道.
耶和華的手大有能力.
也要使你們永遠敬畏耶和華你們的神.
謝謝大家.
在很長的經文裡面.
大家剛剛一路讀.
不知道有沒有發現很多不同重複的地方呢.
其實在經文裡面不斷重複一點.
應該是很明顯的.
就是在裡面是說很多次.
在十二支派裡面要選一個人.
那個人要每人拿一塊石頭.
即是拿十二塊石頭.
第二至第五節已經重複了幾次了.
那十二個人每個人代表一個支派.
要按著十二支派的數目拿一塊石頭.
接著第八節又說要按著以色列人支派的數目.
又要拿十二塊石頭.
十九至二十節又再重複說.
在他們約旦河裡面拿那十二塊石頭.
約書亞就立在吉甲裡面.
重重複複地去到說.

$^{121}$其實還包括約書亞.
去到將那塊石頭放在祭司旗的地方.
為什麼要說這麼多.
誰拿多少塊呢.
其實我想我們不難去想像的.
就好像我們成長班裡面.
如果我們要每個人派一個人做代表.
拿一塊石頭放在紅地毯這個通道裡面.
不知道夠不夠擺呢.
如果大家想像一下.
應該都夠擺的.
那就塞滿條路.
大家就走兩邊出去.
其實就是在說以色列人.
每個支派的代表.
就成為那個支派的代表.
而這十二個人就是整個民族的代表.
包括約書亞都是代表整個民族.
所以那十二塊石頭.
不是為了約書亞一個人去立的記號.
不是為了某一個支派去立的記號.
而是為整個民族立的一個記號.
那個記號是很清楚的.
第六第七節去到說.
那些記號這些石頭是做什麼的呢.
第六節開頭第一句就說.
這些石頭在你們中間可以成為證據.
接著第七節最後那一句就說.
這些石頭是要作以色列人永遠的紀念.
然後二十一至二十二節又補充.
紀念什麼呢.
就是紀念他們日後子孫.
當看到這些石頭的時候.
就可以告訴他們.
以色列人以前是曾經走過乾地.
過這個約旦河.
約書亞說要成為那個證據.
是證明神曾經真的讓約旦河的河水斷絕.
真的成為乾地.
另外一個有重複的.

$^{161}$除了十二支派這個代表之外.
裡面有一個不是很起眼的重複.
可能大家剛才讀完都會覺得很普通.
其實這麼多節經文裡面就是重複.
你們 我們 他們.
所以我都不打出來了.
因為大家一定讀很多.
你們 我們 他們.
其實就是想告訴以色列人和我們知道.
那個恩典的記號不是為誰.
是為整個民族去設立.
所以要知道恩典的記號在這裡.
以色列人就要相信神是多麼的信實.
不單止走過約旦河.
更加要相信他之後會帶領他.
可能去到這裡你就會想.
我都知道 看這段經文就是講神的信實.
我就回去數算一下恩典就知道神是信實.
但是我們又想一下.
其實我們面對生活.
我們有沒有對神缺乏信心的時候呢?.
有沒有面對一些難關.
然後我們覺得好像神都不知道能不能幫助我們呢?.
我想大家都會想到有吧.
其實沒有人會說我完全對神很有信心.
我每一次都能夠很有信心.
一定有些信心軟弱的時候.
什麼事令我們原本充滿信心的自己.
對神很有信心.
慢慢會覺得可能神都幫不了.
我自己想可能是現實那種殘酷.
是將我們本身滿滿的信心.
攻擊到體無完膚.
舉個例子.
可能是上班的時候.
我們都上班為工作祈禱.
希望可以準時收工.
或者我們順利一點.
但是祈禱了很久.
不單止沒有準時收工.

$^{201}$OT由七點開始到八點,九點,十點.
甚至有些弟兄姊妹告訴我.
其實他通宵都在家裡工作.
健康可能病患傷風感冒會康復.
但是我們都經歷過有病患的肢體.
可能你為他祈禱.
他未必能夠康復.
未必像我們祈禱那樣說.
神真的有醫治他.
家裡呢?.
可能我們都為家裡的事情去祈禱.
希望家人和睦.
子女可以孝順自己.
但是可能事情又不是這樣的.
我都為他祈禱.
回教會很久了.
我聽很多家長跟我說.
但是他很不長進.
最近很頑皮.
他不回教會.
種種的事情.
我想我們都經歷過.
弟兄姊妹.
我不是想很負面地說.
神不行了.
或者不是想諷刺大家.
我們就是因為沒有信心.
所以你生活得不好.
接著我就說一個我很有信心的例子.
叫大家對神有信心.
我並不是這樣.
我是想指出.
上述我有說的情況.
其實我們活了這麼多年.
我們做了這麼多年人.
其實我們都經歷過.
那些不理想不如意的地方.
那些人際關係的以如我詐.
那些情況我們一一都經歷過.
所以為什麼我們對神沒有信心.

$^{241}$其實很可能是我們經歷了這麼多.
我們做了這麼多年人.
我們不敢再相信神那種信實.
都不行的.
不如我自己想方法.
健康不如多看幾個醫生.
多問幾個醫生的意見.
可能那件事就可以了.
我工作上想想辦法.
怎樣跟我上司溝通.
叫下屬做些事情.
子女學上什麼課.
我們就聽話一些.
會去教會.
可能我能夠幫助到他.
我們用很多不同的方法.
唯一只有一個目的.
就是我們不再相信神那種信實.
我們覺得自己心裡的焦慮.
可以有方法可以掃除.
我們那種不安.
是我們有方法可以解決到.
我想背後不單單是不相信神的信實.
而背後是我們真的覺得.
那個恩典是過去了.
不單單只是恩典過去.
連神都過去了.
恩典當然是過去了.
好像以色列人那樣.
他納石為記.
那塊石頭就放在那裡.
他不會動的.
但最大問題是.
如果我們連背後的神.
我都覺得他是過去了.
以後他都不會幫助我的了.
他不會和我同在的了.
他沒有能力.
去令我承載今天的困難.
這個才是最大.

$^{281}$最大我們不相信神信實的原因.
弟兄姊妹.
你有沒有這樣的情況.
不敢向神說.
我怎樣祈禱.
我怎樣相信你的信實.
我想唯有直到一件事發生.
我們有一天再次去體驗.
再次認清.
耶穌基督是那個在十字架上.
寫生救贖我們那位主.
我們再次發現.
世界上是有一位愛人到底的上帝.
他一直同在.
我們先將剛才我說的那一堆.
現實很殘酷的那一堆例子.
我們先不會再視它為永恆.
健康.
工作.
我們的子女.
可能我們的事奉.
甚至是我們的前途.
我們都不會再視作永恆.
是很重要.
不過不是永恆.
我們深信即使多艱難.
神都同在.
他信實不變.
所以我們生活才能夠有力.
在聖經裡面.
為何要看23,24節呢.
就是補充解釋.
為何要放石頭在這裡的意思.
在23,24節裡面.
因為耶和華你們的神.
在你們前面.
令到約旦河的水乾了.
等著你們過來.
就如耶和華你們的神.
從前在我們前面.

$^{321}$使紅海乾了.
等著我們過來.
一樣.
這段經文去到勾起.
以色列人他們的記憶.
記得原來上一代.
或者自己去到過紅海.
是神去到幫助他們.
以前為何能夠分開紅海.
是否因為以色列人自己會變魔術.
是因為他們倚靠神.
紅海才分開.
他們才能夠過去.
現在呢.
約旦河為何能夠成為乾地.
上星期荒牧師有講到.
其實當時是水漲的時間.
這樣的情況是不可能發生的.
以色列人是因為倚靠神.
所以神將約旦河的水分開.
讓約旦河成為乾地.
讓他們走過去到吉格.
將來呢.
我們其實看過聖經第五章.
第六章第七章.
你就會發現.
當然是神有能力.
不過你代入當時過約旦河的那群以色列人.
他知不知道是這樣的.
他知不知道圍住耶利哥城的七天.
然後城就倒塌.
他不知道的.
其實他知道迦南地的人很難打.
打硬仗.
但是神告訴他們.
將來仍然神都和他們同在.
無論前面的仗有多凶險.
神都會和他們一起同行.
神是信實不變.
所以我想第一個學習是.

$^{361}$我們要知道神是過去,現在,將來都在和我們同在的神.
祂過去和以色列人同在.
現在祂也會和我們同在.
因為祂是那位守施慈愛的神.
從來從來都沒有改變過.
過去的恩典提醒我們.
神是一直都在.
主耶穌基督如何介入大家的生命裡.
如何讓你對神原本充滿滿滿的信心.
我們回想那些恩典的記號.
當我們每一次想的時候.
其實就讓我們記得.
神是信實不變.
現實的殘酷不會激動我們.
所以為什麼學習數算主因是重要.
因為我們常常都記得神是信實不變.
想到這裡其實我想了很久.
用什麼的例子.
最後我覺得用這個例子比較合適.
我大學畢業就做小學老師.
教中文.
我記得剛剛第一年教書的時候.
上班的日子.
我通常上班都很早上學.
因為我是要很早坐下.
吃早餐 改補.
要定神 知道自己一天要做什麼.
有多少課 我要做多少.
我才會安心一點.
我很早已經在改補.
有一天我記得我在改補的時候.
同事陸續回來.
差不多時間.
我同級同科的幾個同事.
正在討論一些東西.
沒有叫我當然不討論.
我就繼續改補 不關我的事.
但當我聽的時候就覺得好像有點問題.
聽起來的氣氛有點不對勁.
有點緊張.

$^{401}$後來聽著聽著好像有點火藥味.
接著我就認真一點聽.
原來是因為某個課題可能教得不好.
想某個同事教得好一點.
去幫他 大家明白幫是什麼意思.
我看到這麼重火藥味.
我就做了一件事.
你猜我做了什麼事.
你代入一下我的情況.
你會做些什麼.
有沒有人會立即起來說公道話.
有沒有人會繼續改補.
不要叫我嘛.
有沒有.
是 是 是.
可能你正在想其他東西.
其實都可以不拒收.
大家自己知道都可以.
我做了第三個選擇.
老師的桌子是這樣的.
有個高一點的 我坐在那裡.
我在改補.
我縮低了頭.
然後繼續改補.
我怕其他人見到我.
叫我起來.
但是有趣的地方就是.
我這邊有個同事.
他坐在我一條巷的旁邊.
同事這樣看我.
然後他就叫我.
老師叫代號.
我是魏文俊.
所以這個姓氏都沒什麼人用.
就叫阿魏.
他轉頭看過來就跟我說.
喂 阿魏.
你起來一起聊吧.
我已經被人點名了.
當然不能走了.

$^{441}$所以就起來聊吧.
故事的.
不是故事 是真的.
這件事的最後結果.
其實我已經忘記了.
總之我無論如何都要參與討論.
但是我想起.
其實幾年的.
做教師的生涯.
裡面.
神常常都有這樣的情況.
那個同事是一個很好的基督徒.
他的樣子那個畫面.
我時常都回想.
起來吧 一起聊吧.
我覺得那一刻.
我回想好像神捉住我一樣.
你不要躲起來.
有個位置.
在改補 假裝在改補.
你改補都不夠.
還要縮起來 不想被人看到.
我自己回想.
是神去藉著那些.
情況 這些恩典.
去提醒我.
去看到別人的好處.
要勇於面對自己.
那種軟弱 學習.
主耶穌基督的樣式.
的確是很不容易.
很怕衝突 但是.
在學校裡面那幾年.
我經歷了很多衝突.
而神都讓我慢慢去學習.
怎樣去面對和處理.
為什麼覺得這個例子很合適呢.
就是在教會裡面 服事八年多.
我自己經歷社運疫情.
和移民潮.

$^{481}$其實牧羊的事情.
是有很多轉變.
特別大家都經歷疫情的時候.
你坐不住在那裡崇拜.
要妄想.
但妄想崇拜都不是突然間.
生出來的 是有一段時間.
大家不知道怎麼辦.
大家會迷不迷惘呢.
我也很迷惘.
我沒有試過這樣做傳道人.
我那時候就形容.
是一個虛擬的傳道.
虛擬就是網上的傳道.
虛擬的傳道是怎樣做的呢.
見不到弟兄姊妹 又要牧羊.
又要和別人建立關係.
很難做.
很迷惘 有時候.
都和神說.
不如算了吧.
當我很想說算了.
那一刻.
那個畫面又再出來了.
那個同時的畫面又再出來了.
神又再捉我出來.
其實為弟兄姊妹的好處.
是有很多事情可以去做.
所以今天在教會.
能夠有些網上的事情.
去做.
自己都覺得能夠接受.
都是過往的恩典.
讓我相信神的信神.
弟兄姊妹 你有沒有這樣.
的恩典記號.
在你的生命裡.
主耶穌基督有沒有試過.
這樣介入你的生命.
捉住你.

$^{521}$讓你記得神真的信實.
即使現在.
很難 很困難.
很殘酷 我的現實.
但是我都記得神是信實.
我想就是.
這些恩典記號.
讓我們不斷去相信主耶穌基督.
祂那份信實和那份真.
不是馬上消除困難.
就好像疫情已經過了幾年.
原來先走過了.
現在回望.
我們都是.
這是今天第一個.
學習的功課.
我們要想起那些恩典的記號.
目的是提醒我們.
神是色在.
金在 淨在.
來臨的神.
祂一直都沒有變.
祂是很真實.
第二個字.
都是P字.
第一個字是名詞.
第二個字是動詞.
要動一下.
不過沒有T音.
是前進的動力.
是行前的意思.
原來神給我們的恩典記號.
是要推動我們整個群體.
不只是我們自己.
是整個群體.
要學習信服.
去更新生命.
特別是面前未知.
很難的路.
我們怎樣靠著相信神.

$^{561}$走下去呢?.
我們一起讀10-19節.
耶和華的約季和祭司就在百姓面前過去.
路邊人 加特人 馬來西半支派的人.
都照以摩西所吩咐他們的.
帶著兵器在以色列人前頭過去.
約有四萬人都準備打仗.
在耶和華面前過去.
到耶利哥的平原等候上陣.
當那日.
耶和華使約書亞在以色列眾人眼前尊大.
在他平生的日子.
百姓敬畏他.
像從前敬畏摩西一樣.
耶和華曉遇約書亞說.
你吩咐臺法季的祭司從約旦河裡上來.
約書亞就吩咐祭司說.
你們從約旦河裡上來.
臺耶和華約季的祭司從約旦河裡上來.
腳掌乾落寒地.
約旦河的水就流到原處.
仍夠漲過兩岸.
正月初十日.
百姓從約旦河上來.
就在吉甲.
在耶利哥的東邊安營.
在剛才第六第七節.
約書亞說要將石頭作為證據.
又要成為以色列人永遠的紀念.
不知道大家有沒有一些紀念的東西.
放在家裡提醒自己呢.
我就想起十幾年前受浸.
有一張教會的受浸納字卡.
大家有沒有想起那張卡是什麼顏色.
上面有少少紫色.
下面白色.
然後有五行剔納字.
我以前就將納字卡放在我工作的地方.
玻璃的下面.
今天就放在我辦公室玻璃下面的地方.

$^{601}$提醒自己原來當日對神都很有信心.
因為那五行我想都沒想.
就覺得很容易就剔了.
大家有沒有剔了呢.
我試過見到考問信德的時候.
有弟兄姊妹說.
這兩行我暫時不剔了.
我放在這裡都想提醒自己.
對神曾經有很大的信心.
你有沒有什麼紀念品.
提醒自己那些時刻或日子呢.
或者情侶有些紀念日.
或者結婚的人有戒指.
或者受浸有些弟兄姊妹送十字架給你.
放在桌上你就會記得.
我那天受浸了.
我什麼時候受浸了.
不過原來希伯來人那種紀念.
不是只是提一提有哪一天.
不是農曆新年初一就初一.
就初一沒事做.
原來希伯來人那種紀念.
是說他們要有行動.
那種紀念不是就這樣站在那裡.
而是有行動.
要去反思.
要為了那件事去做一些事情.
意思是什麼呢.
就是納石遺跡之後.
那些石頭其實是有個意思.
叫他們要做一些事情.
如果我們今天納石遺跡.
只是放在這裡做一個紀念.
是一個普通的紀念品的話.
可能我們想像.
以色列人今天就會走過去.
可以了 放了石頭.
就過去拍照.
打卡.
點個讚.

$^{641}$然後就走了.
沒事了 沒事發生.
但原來神叫以色列人納石遺跡.
是要以色列人成為一個守約信實的民族.
也同時讓世界知道神是獨一的真神.
去見證神自己的偉大.
有時候我們很容易將神.
局限在一個很小的空間裡.
比如什麼呢.
就好像這個港島powerpoint.
祈禱希望快點出個powerpoint給大家看.
就知道是什麼了.
其實有時候局限在這麼小的空間裡.
其實原來神是很大.
不單單只是在我們家庭工作.
或者我們個人的興趣上.
不是那個powerpoint出了就沒事.
而是原來在過程裡.
神都要我們去想.
怎樣彰顯祂自己的榮耀.
再講得落地一點.
就算在過程裡可能很有阻滯也好.
我們都要見證著神是信實.
我們仍然是生活得有力.
那些事本身都很好.
比如出powerpoint都很好.
給大家看那個字是什麼.
有這些字.
比如給大家看盒子.
有個圖像給大家看.
不過是不是這些就是終極呢.
就是我們要關注永恆的地方呢.
我想不是.
神要讓以色列人納石為基.
不是想他們圍爐取暖.
不是要打卡騙個like就完了.
不是集體回憶這麼簡單.
而是要這個民族成為一個信心的民族.
讓他們對神有足夠的信心.
向世界做見證.

$^{681}$所以為什麼第十至十一節會這樣說.
在第四行的尾說.
於是百姓就急速過去了.
眾百姓盡都過了河.
約桂和祭司就在百姓面前過去.
意思是說.
其實他們沒有在打卡騙個like.
也沒有慶祝.
如果我們想像的話.
他們只是站著不走.
他們就切個蛋糕.
很開心就完了.
但他們就是急速過去.
是過去約旦河的另外一岸.
除了他們急速過去之外.
祭司也是一直站著.
直到神的吩咐完成.
他才走.
他雙腳離開乾地的時候.
約旦河的水就回到原處.
除了百姓和祭司之外.
兩個半支派的人其實都一樣.
他們都沒有就這樣完結成為一個終點.
他們都好像承諾了摩西和神一樣.
他們就過去約旦河的另外一岸.
拿著兵器準備打仗上陣.
所以我們看到這幅圖畫是這樣的.
納什維記之後.
以色列人整個群體都是跟從神的說話.
馬上採取行動.
而不是停留在我們騙個like.
打卡去拍神的恩典當中.
所以那種過去不是.
我們不要想像.
我們今天回去拍張照.
很有神的恩典就完了.
而是以色列人知道那種過去.
不只是歷史.
也不是終點.
那種過去是幫助他們.

$^{721}$去相信神是信實不變.
而他們能夠依靠神繼續前行.
進入迦南和各族去打仗.
去到承受產業的起點.
所以弟兄姊妹.
當你下次經歷神的恩典的時候.
除了當然po過社交媒體.
打過卡.
都要告訴大家.
你經歷了神的恩典之後.
其實你都要多做一件事.
你想想這種信實不變的恩典.
這個上帝.
其實要我們多做些什麼呢?.
以色列人就知道.
他要依靠神去打仗.
所以往後攻打耶利哥城.
當約書亞這樣吩咐.
他們就這樣去做.
正正是一個典範.
弟兄姊妹.
你現在在你生命的關口.
有沒有什麼.
是要你急速過去.
你有沒有在河邊仍然在猶豫呢?.
有沒有站在那裡.
神是不是信實呢?.
或者站在那裡打完卡.
很有恩典啊.
我今天工作很順利啊.
我今天做了一件事.
很感恩啊.
就完了呢?.
有沒有再想一下.
其實神在這件事上.
想我們怎樣為祂去見證祂自己呢?.
我認識一所小教會.
這個教會多年來.
都要去不同地方租借地方聚會.
可能我們城鎮弟兄姊妹.

$^{761}$聚會在這裡.
就不太明白小教會的難處.
小教會租借地方有很多原因.
包括他們沒有那麼多錢去買物業.
他們租借的地方.
可能不續租.
他們就要搬.
甚至他們租了一個地方.
其實不是很合適.
不過那時候沒什麼辦法.
就先租住.
其實他們知道不合適聚會和崇拜.
所以他們又想搬.
近來這個教會又搬了.
感恩的是他們在另外一區.
有一所學校借了一個地方給他們崇拜和團契.
而奇妙的地方就在於.
這個教會本身是沒有兒童侍工的.
只有零人聚會的.
但是當他們搬到新的地方.
就是那個學校聚會的時候.
竟然有七至八個小朋友來到這裡.
還要是恆常聚會.
不是來一來就算了.
讓這個教會能夠開展兒童侍工.
去接觸其他區內的居民.
去關心社區.
弟兄姊妹聽到這些見證.
是不是很想城鎮都立刻能夠領受這個恩典.
最好我們明天就決定搬.
然後突然間有很多侍工能夠發展.
很多人來到.
但是弟兄姊妹.
我們知道可見堂進程的底下.
其實並不是這樣.
幾十年來都是走兩步 退兩步.
但是我們想一想.
城鎮其實有沒有恩典的記號.
有沒有想到的.
我聽過會回來城鎮聚會的人.

$^{801}$或者弟兄姊妹.
他們會有幾個不同的原因.
比較籠統的就是.
禮堂都很漂亮.
覺得很莊嚴.
都覺得回來很感受到神.
或者音樂很好聽.
講道也不錯.
或者弟兄姊妹很熱心查經.
因為我們一季有八次查經.
而其他的教會可能是一季一次.
半年一次.
甚至全年都可能沒有一兩次.
這些都可能是其中城鎮很恩典的記號.
不過我想到其中一樣就是成長班.
成長班上年是二十五週年.
今年是二十六週年.
可能對於我們來說.
運作很正常.
我們都知道是怎樣.
有祈禱 唱班歌 背班份.
但是回想二十多年前.
雖然我還沒有回教會.
但是我看文字和一些屬靈長輩分享.
其實二十多年前的城鎮.
這個成長班是一個很全新的模式.
它將主日學和團契二合一結合.
當時的轉變.
我們今天未必能想像到.
但是你可以想像到這麼龐大的轉變.
有什麼會一起發生呢?.
就是大家要有不同的適應.
當中醞釀的時候有很多討論.
可能有人不同意.
可能有人同意.
總之就有很多事情發生.
不過二十多年後的今天.
大家回望一下.
其實成長班這個模式.
是祝福了很多人.

$^{841}$祝福了很多弟兄姊妹.
包括我自己回來.
其實成長班都差不多十週年了.
神樣教會走到今天.
不是很肯定地說.
不是為我們自己安殊.
或者有平安就足夠.
而是很想藉著我們.
向社區和世界去見證主.
以前以色列人去過約旦河.
是一條未知的路.
入迦南地都不知道怎樣打仗.
你以為打仗很好玩嗎?.
不如站在那裡好過吧.
但是神就是想和以色列人去見證.
這個世界有一個獨一的真神.
所以有一節.
剛才是沒有讀到的.
在五章一節裡面說.
在納石遺跡之後.
其他地方發生什麼事呢?.
看到我們一起讀這段經文.
《約書亞記》五章一節.
約旦河西阿摩利人的諸王.
和靠海迦南人的諸王.
聽見耶和華在以色列人前面.
使約旦河的水乾了.
等到我們過去.
他們的心因以色列人的緣故.
就消化了.
不再有膽氣.
弟兄姊妹.
今天如果和身邊的弟兄姊妹.
或者其他未信的人說.
說起城鎮.
其他人的反應是怎樣的呢?.
或者我們和弟兄姊妹之間.
說起城鎮建堂.
和其他教會認識的弟兄姊妹.
說起城鎮建堂.

$^{881}$大家的反應是怎樣的呢?.
是不是好像經文這樣說.
其他人的心就消化.
不再有膽氣.
見到神是獨一的真神呢?.
在建堂計劃這一年多進行中.
如果我告訴你建堂計劃只有一個意見.
大家都很同心.
大家都不相信.
因為你有參加建堂交流會.
有參加祈禱會.
你會見到其實有不同意見.
ABCDE不同的方案大家都會有.
我想說的是.
在商討和等候的過程中.
我們有沒有讓城鎮建堂這個進程.
讓我們自己城鎮的群體.
或者其他人看到神呢?.
如果我們和弟兄姊妹說起建堂.
是否看到神的信實在裡面呢?.
我的意思不是說.
看到城鎮建堂有信實.
就是知道決定了就去走.
然後有很多供應.
而是當我們商討的時候.
我們是否看到在弟兄姊妹中間.
有耶穌基督的樣式在當中.
有一份榜樣在裡面.
所以我自己覺得在我們現況裡面.
有兩件事可以做.
第一件就是.
就算我們意見真的很不同.
有很多不同方案.
我們可以選擇用信心和愛心去做回應.
意思是什麼呢?.
是我們不要採取一個非此即彼的態度.
別人不聽我的意見.
我就反抬離場.
別人不信納不採納我的意見.
我就選擇抗議.

$^{921}$我不支持.
不和你一起走.
當神呼召我們成為一個群體.
回到城鎮當中.
即使不同意見.
我們相信神的信實.
背後的意思是.
可能我的意見和大家主流不同.
但如果確定了.
大家的決定是這樣走的.
我都願意同行.
和一起去跟進.
我就去看看神是怎樣信實不變.
怎樣利用這個建堂的進程.
在過程裡面去令其他人得著.
祝福和造就.
我想這個才是我們應該擺上的信心.
因為我們其實這麼多個會眾.
未必個個都參與在建堂的進程.
給予很多意見.
但我想我們自己的心態.
怎樣參與在裡面.
是很重要的.
這是第一件事.
第二件事就是撇開建堂.
其實我們還有很大一段時間.
就算見不見也好.
這段時間裡面.
我們有很多不同的人加入教會.
如果你打開崇拜的周刊.
你會發現同工.
或者我們身邊的弟兄姊妹.
其實有很多不是城鎮長大的弟兄姊妹.
加入成為同工.
有些不是在這個教會長大的弟兄姊妹.
在其他教會過來聚會.
然後落地生根.
我不是想說我們辨認誰是城鎮的出身.
然後敵我很分明.
而是我們真是同一個群體.

$^{961}$上帝呼召我們來到這個地方成為一個群體.
我們無論是年長年幼.
在屬靈上年長年幼.
或者加入這個教會年紀多少.
我們都要記得一件事.
是神對我們這個群體有未完成的託付.
祂是託付我們將祂的福音傳開.
所以我想其中一件可以做的事.
就是在這麼漫長的過程裡.
我們回想自己.
我們的生命.
我們的學校.
我們的職場和家庭裡.
有沒有見證到神那種恩典的記號.
而那種恩典的記號.
可不可以讓我們身邊未信的人去見到呢?.
我們可不可以讓未信的人去看到神的同在.
和聽到神的聲音呢?.
如果我們真是能夠有這樣的生活見證.
我相信在這個過程裡.
無論最後建堂的過程是去或留.
或者有第三第四第五個方向都好.
神是想使用我們讓我們這個群體去祝福其他人.
我們不要錯了焦點.
將那個建堂放在我們第一個很重要的位置.
能否做到就等於我們教會的成敗.
我們教會的成敗應該取決於我們對社區的見證.
我們自己群體裡面的見證是怎樣.
等於我們都失見證的話.
即使我們怎樣建造一個好的堂.
或者我們在這裡繼續修葺我們原有的堂祠.
這都是一個失敗的故事.
弟兄姊妹請不要錯焦點.
請起來去做一個好的見證.
不單單是你自己.
是我們整個群體都應該有一個好的見證.
去見證主耶穌基督的信實.
所以今天的講題其實是充斥著一個經文裡面的.
那個過去那兩個字大家剛才讀.
不知道能不能讀了很多次.

$^{1001}$帶著過去往前走.
其實我很想很俏口地寫.
帶著過去走過去.
因為經文是很真實地告訴我們.
以色列人就是帶著神恩典的記號.
走過約旦河.
我們其實都應該要帶著神恩典.
過去給我們的那種記號.
我們踏出那個信心的一步.
好讓我們能夠.
因著神能夠使用我們.
因著我們能夠讓其他人認識神這位大能的主.
弟兄姊妹那種信實我再次都講.
不是我們就相信建堂是可以.
還是不建堂是可以.
而是那種信心是.
我們在這麼漫長的進程裡面.
我們會不會相信神在過程裡面都使用我們呢.
而我們群體是否又能夠活出一個好的見證呢.
特別當人提到城鎮或城鎮建堂的時候.
神的名是否得著榮耀呢.
在這一刻我邀請大家在座位裡面.
去想一件神給你的恩典那種記號.
在禱告裡面請你向神獻上感恩的禱告.
去求神使用你的生命.
讓人認識耶穌基督.
以後一段默討的時間.
我們同心禱告.
我們死而復活的主耶穌.
因為你十字架的大愛.
我們得著那份屬天的盼望.
現實的殘酷是不能夠奪去我們這份盼望.
現實的殘酷是不能夠擊倒我們.
我們能夠靠著主你在這個世代裡面生活有力.
剛才主耶穌和我們一起回想.
你在我們生命中恩典的記號.
我很肯定你的恩典是信實.
是豐盛.
是永遠不變的.
這些恩典記號在大家心裡面一定不止一個.

$^{1041}$求你藉著這些記號.
讓弟兄姊妹更加確信.
你自己是信實不變.
我們又求主使用我們.
無論我們在哪一個成長班.
什麼地方.
什麼警方.
求你讓我們堅心去信賴你.
相信你那種信實不變.
以致我們能夠每天靠著主你活出真道.
使我們身邊的人得著祝福和造就.
多謝主你聽我們同心合意禱告.
奉主耶穌基督得勝之名祈求.
阿們.
\newpage



\section{申命記 1:19-40}
\label{sec:031637}
\textbf{過去與現在的對話 (1) YA 好聲音}~~~~羅皓文傳道
\newline
\newline
連結: \href{http://www.church.com.hk/acms/content.asp?site=cdc\&op=show\&type=product\&code=031637\&layout=sermon}{\texttt{031637}} ~~~~ 語音日期: 2026-02-07
\newline
\newline
\hyperref[sec:031636]{\small{< < < PREV SERMON < < <}}
~
\hyperref[sec:index_chronic]{\small{[返順時目]}}
~
\hyperref[sec:031638]{\small{> > > NEXT SERMON > > >}}
\newline
\newline
申命記 1:19-40
\newline
\begin{longtable}{cl}
\hline
\hline
章節 & 經文 (和合本修訂版)\\
\hline
1:19 & \begin{tabularx}{0.7\textwidth}{X} 「我們照著耶和華－我們神所吩咐的，從何烈山起行，經過你們所看見那一切大而可怕的曠野，往亞摩利人的山區去，到了加低斯‧巴尼亞。 \end{tabularx} \\ \\ \relax
1:20 & \begin{tabularx}{0.7\textwidth}{X} 我對你們說：『你們已經到了耶和華－我們神所賜給我們的亞摩利人之山區。 \end{tabularx} \\ \\ \relax
1:21 & \begin{tabularx}{0.7\textwidth}{X} 看，耶和華－你的神已將那地擺在你面前，你要照耶和華－你列祖的神所說的，上去得那地為業。不要懼怕，也不要驚惶。』 \end{tabularx} \\ \\ \relax
1:22 & \begin{tabularx}{0.7\textwidth}{X} 你們都來到我這裡，說：『讓我們先派人去，為我們窺探那地，把我們上去該走的路線和該進的城鎮回報我們。』 \end{tabularx} \\ \\ \relax
1:23 & \begin{tabularx}{0.7\textwidth}{X} 這話我看為美，就從你們中間選取十二個人，每支派一人。 \end{tabularx} \\ \\ \relax
1:24 & \begin{tabularx}{0.7\textwidth}{X} 於是他們起身上山區去，到以實各谷，窺探那地。 \end{tabularx} \\ \\ \relax
1:25 & \begin{tabularx}{0.7\textwidth}{X} 他們的手帶著那地的一些果子，下到我們這裡，回報我們說：『耶和華－我們的神所賜給我們的是美地。』 \end{tabularx} \\ \\ \relax
1:26 & \begin{tabularx}{0.7\textwidth}{X} 「你們卻不肯上去，竟違背了耶和華---你們神的指示， \end{tabularx} \\ \\ \relax
1:27 & \begin{tabularx}{0.7\textwidth}{X} 在帳棚內發怨言說：『耶和華因為恨我們，所以將我們從埃及地領出來，要把我們交在亞摩利人的手中，除滅我們。 \end{tabularx} \\ \\ \relax
1:28 & \begin{tabularx}{0.7\textwidth}{X} 我們上哪裡去呢？我們的弟兄使我們膽戰心驚，說那裡的百姓比我們又大又高 ，那裡的城鎮又大，城牆又堅固，如天一樣高，並且我們在那裡看見亞衲族人。』 \end{tabularx} \\ \\ \relax
1:29 & \begin{tabularx}{0.7\textwidth}{X} 我就對你們說：『不要驚惶，也不要怕他們。 \end{tabularx} \\ \\ \relax
1:30 & \begin{tabularx}{0.7\textwidth}{X} 在你們前面行的耶和華－你們的神必為你們爭戰，正如他在埃及，在你們眼前為你們所做的一樣； \end{tabularx} \\ \\ \relax
1:31 & \begin{tabularx}{0.7\textwidth}{X} 並且你們在曠野所行的一切路上，也看見了耶和華---你們的神背著你們，如同人背自己的兒子一樣，直到你們來到這地方。』 \end{tabularx} \\ \\ \relax
1:32 & \begin{tabularx}{0.7\textwidth}{X} 你們在這事上卻不信耶和華---你們的神。 \end{tabularx} \\ \\ \relax
1:33 & \begin{tabularx}{0.7\textwidth}{X} 他一路行在你們前面，為你們尋找安營的地方；他夜間在火中，日間在雲中，指示你們當走的路。」 \end{tabularx} \\ \\ \relax
1:34 & \begin{tabularx}{0.7\textwidth}{X} 「耶和華聽見你們的怨言，就發怒，起誓說： \end{tabularx} \\ \\ \relax
1:35 & \begin{tabularx}{0.7\textwidth}{X} 『這邪惡世代的人，一個也不得看見我起誓要賜給你們列祖的美地； \end{tabularx} \\ \\ \relax
1:36 & \begin{tabularx}{0.7\textwidth}{X} 惟有耶孚尼的兒子迦勒必得看見，並且我要將他所踏過的地賜給他和他的子孫，因為他專心跟從我。』 \end{tabularx} \\ \\ \relax
1:37 & \begin{tabularx}{0.7\textwidth}{X} 耶和華也因你們的緣故向我發怒，說：『你也不得進入那地。 \end{tabularx} \\ \\ \relax
1:38 & \begin{tabularx}{0.7\textwidth}{X} 那侍候你，嫩的兒子約書亞必得進入那地。你要勉勵他，因為他要使以色列承受那地為業。 \end{tabularx} \\ \\ \relax
1:39 & \begin{tabularx}{0.7\textwidth}{X} 你們的孩子，你們說要成為擄物的，就是今日尚不知善惡的兒女，必進入那地。我要將那地賜給他們，他們必得為業。 \end{tabularx} \\ \\ \relax
1:40 & \begin{tabularx}{0.7\textwidth}{X} 至於你們，要轉回，從紅海的路往曠野去。』 \end{tabularx} \\ \\
[1ex]
\hline
\hline
\end{longtable}
$^{1}$今天會和大家分享.
摩西和以色列人.
其實他們.
在出埃及的時候.
他們已經選擇了一條.
我相信大家都同意是對的道路.
因為耶穌.
應該是耶和華上帝.
叫他們去選擇這條道路.
這條道路的目標.
就是要擺到.
進入加勒美地當中.
不過.
當以色列人.
一出埃及很久.
就出事了.
出什麼事了?.
開始有怨言.
他就嫌三嫌四.
嫌什麼?.
嫌這裡沒有水喝.
嫌這裡很熱,嫌這裡很辛苦.
甚至說到.
其實為什麼.
要找我們出來呢?.
我在埃及不是很好嗎?.
比這裡好,這裡真的很辛苦.
辛苦.
不知道你有沒有經歷過.
例如.
有沒有人喜歡行山?.
這一代你們不行山吧?.
未老到喜歡.
會行.
我年輕的時候.
參加制服團隊.
我都會行山.
都會帶一些年輕人去行.
我就.
對於我來說.

$^{41}$行山是讓我發現上帝的創造.
讓我能夠.
將我自己很繁忙的.
生活裡面,我的老思緒很繁忙.
能夠回到.
大自然神創造裡面.
我有個安靜去賞神.
但是我帶的.
年輕人就不是這樣想.
他們會問.
什麼時候可以回家?.
行山是一個很好的鍛鍊.
因為你上到山的中間.
如果他們說要回家,那怎麼辦?.
不是不能回家.
是要繼續行,你就能回家.
你不能說回頭.
沒問題,不過走到中間.
你回頭還是要走那麼多路.
所以其實是一個很好的鍛鍊.
不過我想表達的是那個年輕人的心態.
就是那種.
很辛苦.
為什麼你要帶我出來?.
我在家,我現在對著電腦.
碳冷氣.
對著.
現在的社交媒體.
我們可以很開心.
為什麼要我動呢?.
神在這個裡面.
要以色列人.
經歷這種.
鍛鍊.
更重要的是什麼呢?.
在這個第一次叫他們要走出去.
加拿美地.
要走去加拿美地的時候.
在這個抗疫的過程裡面.
他們抱怨.

$^{81}$他們就向.
他們的領袖摩西.
去抱怨.
為什麼你會帶我們出來.
不過在這裡.
其實看得出.
這個表面上.
就是向領袖去抱怨.
你們覺得.
為什麼會叫你們.
現在又要做那麼多.
服事.
要那麼多的時候.
你們會抱怨.
抱怨誰呢?.
你們的領袖應該是.
藍傳道.
他們.
心裡面都會有些抱怨的話.
第一個.
應該會指責的是他.
但在這裡.
想反映什麼呢?.
其實表面上你是向領導去.
抱怨.
但其實你心裡面.
在做什麼呢?.
這班以色列人.
我們叫第一代.
第一代出埃及的以色列人.
他們的抱怨背後.
其實是對神的不信.
為什麼這樣說呢?.
是因為.
其實在這個過程裡面.
明明在出埃及的過程裡面.
在未出埃及之前已經做了很多神蹟奇事.
被傳遞的人看到.
埃及人又知道.
有十災.

$^{121}$表明是因為.
當時一班以色列人覺得.
埃及的生活.
很困難.
不如神拯救我們.
帶我們出去.
其實是這樣的.
但出來的時候.
又反過來.
究竟是什麼事呢?.
其中一件事就是.
他們出來.
好像轉換環境.
但他們沒有完全相信神.
是保護他們.
帶領他們.
帶領到底的人.
所以就出現了這件事.
說回今天這個主題.
提一提.
可能一開始的時候也有展示.
今年教會.
全年的主題就是.
多待同行.
而今天.
透過我收到的.
安排.
我們也會根據大堂的.
四個主題.
第一個主題就是與神同行.
我也很想透過今天.
按照這個主題.
跟大家提一提.
藉著剛才我所說.
以色列人.
第一代的以色列人.
用他們這四十年在曠野.
走這些流浪的冤枉路.
去借鑒.
去看一看.

$^{161}$我們往後的時間.
透過申命記在這個位置.
正正就是神.
當第一代.
走過了這四十年的.
以色列人.
到差不多的時候.
而神再一次.
藉著摩西.
向這一班新一代的.
以色列人.
去說一些聲音.
而這些聲音.
我定為是一個好的聲音.
所以今天我也定了.
我的主題是.
給大家看看.
不知道大家有沒有看過.
Y.A號聲音.
看得出.
其實我有些心思在裡面.
看得明白紅色那張椅子是甚麼來的嗎.
有人在笑.
知不知道是甚麼來的.
有沒有想像過類似的東西.
是會轉的.
一按鍵就會轉.
我要你啊.
I want you.
或者一些音樂比賽.
聽到一個好聲音就會轉身.
看著它.
希望你先記住這個標誌.
接下來就看看.
究竟我說甚麼.
我們一起看看究竟在申命記裡面.
和.
耶和華上帝.
想和這一班新一代.
甚麼叫新一代呢.

$^{201}$我要先說清楚.
其實走了四十年.
就是上一代.
他們因為不信.
而白白走了.
四十年的曠野.
那個是上一代.
而走了四十年.
新一代已經出來了.
正正摩西就在這個位置.
上帝要跟他們說.
上帝要透過摩西跟他們說.
你們要準備了.
進入迦南地.
而正正這個背景.
就是在申命記裡面的一個開始.
不如我們就透過.
今天就說一說.
申命記裡面的一些背景簡介.
在申命記裡面.
其實本身的名字.
申命記.
是中文的翻譯.
如果我們看回原來.
七十事譯本.
大家都應該有聽過.
這個版本.
在希臘文裡的.
七十事譯本.
本身是叫第二律法.
再說一次 叫第二律法.
意思是什麼呢.
第二次去.
申述律法.
是什麼意思呢.
有第二次就是有第一次.
第一次是什麼意思呢.
所以我現在在螢光幕上.
那幅圖.
給大家看.

$^{241}$這幅圖是時間表.
出埃及到曠野.
行走了四十年.
的關鍵時間表.
將一些關鍵的時刻都告訴你.
大概讓你知道一出埃及.
第一年發生什麼事.
就是一開頭.
乳節的夜晚.
然後以色列人.
執實行莊.
出埃及 過紅海.
然後進入了.
西賴的半島.
這是第一年.
大家投入得到.
跟著行程就行.
然後進入到第一至第二年.
之間.
整群以色列人.
去到西賴山.
上帝在那裡.
要求這群以色列人.
在那裡紮營.
然後在那裡.
頒佈十誡.
和律法.
那個就為之.
第一次.
頒佈律法.
然後.
隨著第二年就開始.
進入到.
迦底斯巴尼亞.
的地方去流浪.
你見不見到.
這個流浪.
這裡說是因為.
他們不信.
所以他們就被罰.

$^{281}$在曠野那裡流浪.
四十年.
OK.
是四十年.
有多少個四十年.
我想不是很多.
四十年.
起碼都有兩個.
跟著.
第三個沒有.
當時的人就長命.
摩西都很大年紀.
才能夠做領袖.
但這裡想說.
那四十年其實.
可以很長.
可以是他們一代人的命.
所以這裡都是值得我們留意.
然後就去到.
四十年過後.
他們就進入到一個叫.
米利巴水的事件.
其實這個地方就正正.
發生了他們缺水.
他們向摩西去投訴.
然後摩西就恩著.
他都很不滿.
所以他就敲打了石頭.
沒有按上帝說.
上帝本身是叫他要.
用口去.
說水就會出.
但摩西卻.
按自己的己意敲打了石頭.
以致摩西都很慘.
他都不能進入.
迦南地.
這是一個簡介.
我相信大家對聖經有一定的認識.
這裡我就不再詳細解釋.

$^{321}$不過想給你一個時間表.
你知道我今天即將要說的.
就是最後那個點.
就是去到.
那個第四十年.
他要摩西再一次.
被神呼召.
向新一代的以色列人.
要跟他們說準備好了.
大家準備好了.
我們要.
再聽多一次神說.
他的十誡.
並且說得更深入.
為了什麼?.
就是進入神為你們預備的迦南美地.
我們不如就一起來看看.
聖經.
其實今天.
我主要想和大家.
說的聖經其實在第一章裡面.
而最主要.
第一章其實.
我都要先解釋一下.
第一章其實主要.
在整個申命記裡面.
其實是.
摩西當時候.
預備要和那群以色列人.
開始說.
大家新一代的.
我要和你們說.
他們的歷史.
衰得很厲害.
因為他們不是很好.
所以他們.
足足走了四十年.
你們都看到的.
或者你們都跟著的.
現在差不多時候了.

$^{361}$上帝就要和你們說.
再一次起來.
要去.
明白上帝是怎樣的上帝.
以致要進入迦南美地.
這個第一章其實就是.
你可以這樣說.
可以理解摩西將那個歷史.
再重溫一次.
所以在今天裡面.
我不會將整個第一章和大家一起讀一次.
我只是會集中.
我想將重點擺在申命記的一章.
二十到二十二節那裡.
不如都邀請大家.
一起來動動你的口.
一起來讀一次.
預備 123.
我對你們說.
你們已經到了耶和華.
神所賜給我們的墓.
看啊.
耶和華.
你的神已將那地.
擺在你面前.
你要照.
耶和華你列祖的神.
所說的上去.
得那地為業.
不要懼怕.
也不要驚慌.
你們都就跟我來說.
我們要先打發人去.
為我們窺探那地.
將我們上去.
該走何道.
必進何城.
都回報我們.
在這裡你會見到.
其實.

$^{401}$已經在.
這個過程裡面.
你見到他所有描述.
其實已經去到.
阿摩利人的山地.
接著.
其實摩西就說.
看啊.
耶和華你的神.
已將那地擺在你面前.
你記住.
21節.
這一句我想你弄清楚.
在《申命記》裡.
這一節其實是在複述.
第一次.
其實.
他在第一代.
的以色列人.
已經說過這句說話.
而他這次.
他再告訴他們.
之前.
阿南就跟我說過.
跟你們說過.
叫你們要.
努力崇拜.
努力事奉.
現在我講這一番說話.
其實是在複述.
同樣他也是在複述.
不過他不是在複述阿南說的話.
他是在複述.
耶和華在第一次.
其實透過摩西已經講了一次.
這一次再講的時候.
為什麼我要集中在這一句說話.
去給大家知道呢?.
因為當時第一代的以色列人.
他聽過了.

$^{441}$不過我剛才說他們發生什麼事.
他們因為他們的心.
其實都未完全相信.
耶和華.
以致到他們.
我粗俗一點說他們衰了.
他們衰了.
同意嗎?.
而這次所以.
前車可鑑.
這一班新一代的以色列人.
要再一次重申.
你們要記住.
耶和華.
你的神已經將那地.
擺在你的面前.
接著.
你要照耶和華你列祖的神所說.
請幫我一起讀紅色那三個字.
可不可以讀多一次.
這個字很重要.
讀多一次.
重要的字要讀三次.
上去得.
我特別很想.
focus說一下什麼叫上去得.
你們有沒有經歷過.
什麼叫上去就得.
即是你去那裡.
你就得了.
有沒有經歷過類似的事.
通常我這樣問很難.
因為我都要用時間想想什麼叫上去就得.
我就想到了.
我就想我剛剛信耶穌.
沒多久.
就誤導信耶穌.
我和你們差不多.
即將步入.
職場.

$^{481}$中間之前這段時間.
我就在教會裡面.
參加了報道隊.
接受裝備.
在過程裡面.
要求我.
要學怎樣去街頭報道.
在座有沒有人.
試過街頭報道.
有的.
請放下手.
有沒有人試過接受了裝備.
去街頭報道.
沒有的.
就馬上去街頭報道.
這些就真的對了.
這些就叫做.
你上去就得了.
你明不明白意思.
走出去街報道就得了.
這裡想表達的意思誇張到這樣.
當然.
我都是人來的.
我想講講我的經歷.
當時候.
我學過報道訓練.
有本書叫.
四律 如何與神和好.
有一條橋.
過去不知道你們有沒有聽過.
我當時就學了這東西.
當時的牧者就叫.
得了.
你們出去吧.
兩個兩個.
街頭報道.
你見到人就跟他說就得了.
我心想.
真的嗎?.
真的有人停下來嗎?.

$^{521}$第二個問題是.
他願意聽我說耶穌嗎?.
第三個問題是.
他聽我說耶穌.
他願意馬上接受耶穌嗎?.
真的嗎?.
但是牧者就說.
是的.
你上去就得了.
我就真的.
坦白講我當時很火熱.
牧者告訴我去.
我就順服.
結果一出去.
帶著戰經的心情.
接著.
停下來.
你好.
當然一定是試過吃檸檬.
一停下來.
不好意思我來自哪一間教會.
可不可以跟你聊聊天.
真的有人停下來.
有人停下來.
接著就開始說.
請問你有沒有.
上過教會.
沒有.
今天其實我出來.
想跟你分享一下耶穌.
其實就是這樣.
講完一輪.
其實你會不會想.
今天這一刻接受耶穌.
好啊.
真的有嗎?.
真的我不是開玩笑.
我真的有這樣的經歷.
真的試過在這個經歷裡.
讓我明白了原來.

$^{561}$上去.
就真的可以.
你真的肯為耶穌.
為神去.
就真的做到一些.
難成的事.
在這個裡面.
我透過這個經歷.
很想告訴你.
其實當時神很簡單.
就是向著第一代的.
跟他們說.
那句經文可不可以再給他們看一次.
就是.
其實前面那句.
他說.
耶穌說我們預備的地.
已經放在前面.
你知道做的動作是什麼嗎?.
一起讀一次.
上去就可以了.
但是以你.
認識以色列人.
他們當時.
他們的回應是什麼?.
知道嗎?.
他們的回應就是.
喂.
我們看下一句.
我們一起讀一次.
預備 1 2 3.
你們都.
在這裡.
你會見到.
他們很聰明.
他們不是說不去.
不過他們去之前先看清楚.
怎樣看清楚?.
他們更聰明的是.
他們不是親自和耶和華說的.

$^{601}$他們是去找他們的領袖.
你們就好像你們那樣.
找阿南說.
阿南,我們不太想跟旅行.
不如.
阿南,你幫我們.
派一些人去探一探.
看看好不好玩.
好玩就去.
誰知道回來.
12個有多少人說不好玩?.
有10個說不好玩.
只有2個說好玩.
通常在這個世界裡.
少數服從多數.
結果就.
因為這樣衰了.
在這個過程裡.
其實這裡反映了什麼?.
其實他們已經有一個不信.
為什麼我剛才那麼堅持.
我說第一代的以色列人.
他們出現了那個不信.
他們背後.
表面上.
雖然因為那些抱怨.
那些埋怨.
摩西說.
為什麼你這樣?怨他.
但其實在這個第一步裡.
這個回應已經反映了.
他們根本就已經不是完全相信.
神已經為他們.
所有的一切.
在這個裡面.
我們.
有什麼功課要學習呢?.
我們.
可以看看.
一件事就是.

$^{641}$其實.
始終我們都是人.
很明白一件事.
要做到信.
其實並不容易.
在聖經裡面有一個經文.
我們經常.
我們經常都會有一個.
提醒.
我們一起讀一次聖經.
預備 1 2 3.
希伯來書的十一章一節.
其實這個我相信大家都有機會聽過.
就是關於什麼叫信.
信就是.
簡單來說.
很簡單說.
信原來.
是你未見到.
你已經確認.
它會成真.
你會相信.
它會出現.
而這個就是.
我們要對上帝有這份信心.
上帝呼召我們.
或者叫我們要做某些事情.
縱然那件事還未發生.
我們就要相信.
這個才是信心.
如果.
上帝叫我們去做.
而真的出現了.
那就是事實.
那就不需要有信心.
分不分得到.
其實對於以色列人來說.
最難學的.
對我們來說都是很難學的.
就是信心.

$^{681}$我講一下.
我自己的經歷.
其實在我剛才說的中午.
信耶穌.
是真的.
在我中午信耶穌的時候.
我心裡已經很確認.
我一生要跟從耶穌.
所以.
在這個成長的道路裡.
我不斷地.
去等待.
上帝何時呼召我.
去服侍你.
好了.
直到幾年前.
幾年前.
有一位牧者.
就跟我說.
Amon 不如.
你.
進入教會全職服侍.
你猜.
我怎樣回應.
我已經預備好.
打從我信耶穌的時候.
我已經決定我一生要跟隨主.
你猜我怎樣回應.
很噁心.
幾年前.
答案是.
不是我這樣想的.
原來.
我一開始好像很有熱心.
但其實我一直在預備.
我自己經常說.
我學這件事都是預備將來服侍神.
我做這件事.
就算我在職場上遇到任何挫折.
我都是為神而經歷.

$^{721}$但幾年前.
剛才說的這位牧者.
他邀請我.
那一刻.
我反而回應.
我覺得我還沒有預備好.
原來真實的情況.
是什麼事呢.
原來當時.
我已經結了婚.
我也有兩個小朋友.
我要去.
照顧家庭.
我要去賺錢養家.
神啊.
其實如果.
很簡單.
神如果能夠給我三年時間.
知不知道三年做什麼.
這三年時間.
我會想方法.
你教我也好或我去學.
學一下怎樣做到財務自由.
你們有沒有人聽過財務自由.
即是.
找一些方法.
令到錢放進去.
錢會自己賺錢.
那些錢.
賺錢出來.
可以照顧到家庭.
一切的需要.
我就安心了.
全心地服侍你.
上帝.
你覺得這件事.
是否現實.
我告訴你,是很現實的.
很現實的意思是這個心態.
不知你們.

$^{761}$每一位神對你們有感動.
或叫你去服侍神的時候.
你會不會有很多.
你的掙扎.
我現在的工作.
我已經工作了很長時間.
要我回來.
投入這麼多的服侍.
或者要我回去崇拜.
會不會.
你給我一份工作.
舒舒服服的.
早就.
早上九點.
晚上五點.
早點下班.
預備心情.
睡一晚.
星期六回來崇拜.
正常的.
我也會有這樣的心情.
在這個過程裡面.
我發覺一件事.
我真的用了這幾年的時間.
當牧者跟我說.
呼籲我.
然後我拒絕了.
要用幾年時間去做的時候.
我得出的結果是什麼呢.
我結果是.
越做.
越想賺錢,是越窮的.
是越不夠用的.
最慘是什麼呢.
最慘就是.
最近這幾天.
我發現了一個真相.
就是.
我同工.
介紹了我去做推拿.

$^{801}$那麼.
推拿的師傅就跟我說.
因為這個推拿.
是有醫治性的.
他就會說.
羅弟兄.
你全身都很不妥.
你全身都塞了.
你那些.
心肝脾肺腎.
全部都.
很不妥.
你真的要搞了,你再不搞.
什麼三高都會有.
那一刻.
我聽完之後就馬上想起這件事.
我就醒覺.
其實在過往的日子.
我心裡面是很想去.
全身的伺服神.
不過.
我卻.
發現.
當我用人的方法.
去完成神給我的任務.
我是做不到的.
是越做越壞.
連身體的健康.
都沒有了.
所以在這裡.
其實想帶出什麼呢.
想帶出.
其實那個.
神的呼召.
其實很簡單.
祂呼召你.
你有信心,你就去回應.
其實天父爸爸.
是很好的.
在我講的見證裡面.

$^{841}$到最後,到今天.
他出現了什麼事情呢.
就是在中間裡面.
他彷彿就像聖經裡面.
所講的那位浪子的爸爸.
他一直都在等待著.
我的.
那個迴轉.
那個將心.
歸向給上帝.
我的.
狀態.
而更重要的是.
他讓我要明白.
其實我的身份是什麼.
就是天父.
兒女.
天父寶貴兒女的身份.
以致讓我.
明白這個身份的時候.
我就活回我這個身份.
應該要有的生活裡面.
在過程裡面.
過往這一年.
應該說是這一年前.
他就開始不斷用一些.
不同的方法.
去發出一些聲音.
他令到我真的一路一路.
歸向他,迴轉向.
神那一邊.
在這個迴轉裡面,令到我最終.
他的感動是大到.
我就好像那張椅子.
我又要展示一下PowerPoint給大家看.
像這張椅子.
我真的.
要.
急不及待.
要很大聲的.

$^{881}$跟上帝說.
I want you.
跟上帝說,我需要你.
這個經歷.
這個處境.
因為之前的生活.
太磨難.
遭遇太多的低谷.
太多的挫折.
令到我.
經歷到像浪子一樣.
我才想起.
原來愛我們的天父.
爸爸已經在等著我.
我只要迴轉.
歸向他的時候.
其實他在等待著我.
所以一年前.
當真的再有.
牧者又再和我提及.
Amon.
你會不會考慮.
全職的.
去教會去服侍上帝.
在那裡.
我就二話不說.
我心裡就立刻說.
上帝,我來了.
我願意.
不過,還有一些事.
就是.
在我現在的狀態裡.
我不是一個人.
我是一家,我有太太.
太太也有來支持我.
還有.
我一對子女.
我要和上帝.
談一談.
上帝,如果.

$^{921}$我願意去全職服侍你.
不過,如果.
我一家人.
他們不是一起.
受感動.
或者同意支持我.
去全職服侍的話.
我很難作這個決定.
因為我們是一家.
一起.
很感恩.
有聽我的聲音.
或者聽我的禱告.
上帝.
親自的.
讓我太太.
在不同的經歷裡.
她感受到神.
是支持.
讓她支持這件事.
支持我去正式的.
去作為一個傳道同工.
而.
太太也同樣地.
被神感動.
我相信她這一刻也很努力.
在聖經學院裡.
繼續讀神學.
繼續像我一樣.
預備為神服侍.
還有.
我的兩個子女.
我也問他們.
爸爸有這個打算.
你怎樣看.
他們兩個.
基本上回應都很感恩.
令我很感動.
他們說.
既然上帝讓你去做.

$^{961}$你當然要去做.
你沒理由做其他事.
既然連.
下一代.
我的子女.
他們也有這個.
決心的時候.
這麼大的肯定的時候.
我沒辦法.
我一定要去做.
所以在這裡.
我很感恩.
讓我的浪子.
再一次回到天父爸爸當中.
我也想.
多謝她在這個過程裡.
她的支持.
因為她真的很重要.
作為一家.
一個群體.
我們這個家裡.
是很需要.
不是一個人.
是一個群體.
我今天也很想.
和大家說.
天父其實真的很好.
他對我們每一個在在你們每一位.
其實他也有一個呼召.
就好像今天.
經文裡面所說.
神早就告訴你們.
各位.
神已經將.
祂為你們預備的最好的.
放在前面.
只要你們願意去.
就可以了.
我今天很希望.
透過我自己的見證.

$^{1001}$透過我的分享去鼓勵大家.
神真的被我現在.
事奉了.
在教會裡全職服侍.
已經有一年多.
我可以和大家說.
當一個人.
當我自己明白了.
自己是一個屬天.
寶貴兒女的身份.
然後我將我過往.
神給我經歷到的一切.
我的恩賜和能力去發揮.
而走在.
神為我預備的命定裡面的時候.
那份喜樂.
是沒有.
其他東西可以代替的.
我也很想.
去鼓勵大家.
各位.
其實你們都是我的家人.
在同福這個大家庭裡.
我很想鼓勵大家.
一起.
好像今天這段經文.
鼓勵我們.
神其實已經看重我們.
你們每一個都是神寶貴兒女.
我們是一家人.
我們是一個群體.
今天神呼召不會只是呼召Amon.
是會呼召你們.
在座每一位.
你不一定像我這個角色.
要站在台上去講訊息.
或者做傳道.
但是你一定有.
屬天的身份和屬天的命定.
我鼓勵你.

$^{1041}$你就去尋找.
因為神很願意.
你去尋找的時候.
祂必定讓你尋得見.
不如.
這一刻我很想有個祈禱.
很想.
為大家有個祈禱.
七天父 七天耶穌 七天聖靈.
我們.
去呼求你.
我們呼求你.
聖靈.
親自來到我們當中.
你幫助我們.
在我們的生活裡.
我們有不同的光景.
上帝我們.
肯定的是我們只有一個神.
就是耶穌基督你.
願你繼續幫助我們.
就如今天所講的.
我們很願意.
清晰知道我們的身份.
就是你的寶貴兒女.
我們願意.
在我們的生命裡.
再一次的向你迴轉.
歸向你.
在這個過程裡.
可能有很多的艱難.
但我們深信.
你是.
必定帶領我們進入你的美地.
所以藉此.
我們一定會繼續.
在你的裡面堅持下去.
多謝你 聽我們禱告.
奉耶穌基督的聖名.
Amen.

\newpage



\section{創世記 6:5-22}
\label{sec:031638}
\textbf{多代同行•與神同行系列:信心的同行}~~~~王啟恩傳道
\newline
\newline
連結: \href{http://www.church.com.hk/acms/content.asp?site=cdc\&op=show\&type=product\&code=031638\&layout=sermon}{\texttt{031638}} ~~~~ 語音日期: 2026-02-08
\newline
\newline
\hyperref[sec:031637]{\small{< < < PREV SERMON < < <}}
~
\hyperref[sec:index_chronic]{\small{[返順時目]}}
~
\hyperref[sec:031639]{\small{> > > NEXT SERMON > > >}}
\newline
\newline
創世記 6:5-22
\newline
\begin{longtable}{cl}
\hline
\hline
章節 & 經文 (和合本修訂版)\\
\hline
6:5 & \begin{tabularx}{0.7\textwidth}{X} 耶和華見人在地上罪大惡極，終日心裡所想的盡都是惡事， \end{tabularx} \\ \\ \relax
6:6 & \begin{tabularx}{0.7\textwidth}{X} 耶和華就因造人在地上感到遺憾，心中憂傷。 \end{tabularx} \\ \\ \relax
6:7 & \begin{tabularx}{0.7\textwidth}{X} 耶和華說：「我要把所造的人和走獸，爬行動物，以及天空的飛鳥，都從地面上除滅，因為我造了他們感到遺憾。」 \end{tabularx} \\ \\ \relax
6:8 & \begin{tabularx}{0.7\textwidth}{X} 只有挪亞在耶和華眼前蒙恩。 \end{tabularx} \\ \\ \relax
6:9 & \begin{tabularx}{0.7\textwidth}{X} 這是挪亞的後代。挪亞是個義人，在他的世代中是個完全人。挪亞與神同行。 \end{tabularx} \\ \\ \relax
6:10 & \begin{tabularx}{0.7\textwidth}{X} 挪亞生了三個兒子，就是閃、含和雅弗。 \end{tabularx} \\ \\ \relax
6:11 & \begin{tabularx}{0.7\textwidth}{X} 這地在神面前敗壞了，地上充滿了暴力。 \end{tabularx} \\ \\ \relax
6:12 & \begin{tabularx}{0.7\textwidth}{X} 神觀看這地，看哪，它敗壞了，因為凡血肉之軀在地上的行為都敗壞了。 \end{tabularx} \\ \\ \relax
6:13 & \begin{tabularx}{0.7\textwidth}{X} 神對挪亞說：「在我面前，凡血肉之軀的結局已經臨到，因著他們，地上充滿了暴力。看哪，我要把他們和這地一起毀滅。 \end{tabularx} \\ \\ \relax
6:14 & \begin{tabularx}{0.7\textwidth}{X} 你要為自己用歌斐木造一艘方舟，並在方舟內造房間，內外都要抹上瀝青。 \end{tabularx} \\ \\ \relax
6:15 & \begin{tabularx}{0.7\textwidth}{X} 方舟的造法是這樣：要長三百肘，寬五十肘，高三十肘。 \end{tabularx} \\ \\ \relax
6:16 & \begin{tabularx}{0.7\textwidth}{X} 方舟上面要造天窗，向上一肘。方舟的門要開在旁邊。方舟要分上、中、下三層。 \end{tabularx} \\ \\ \relax
6:17 & \begin{tabularx}{0.7\textwidth}{X} 看哪，我要使洪水氾濫在地上，毀滅天下凡有生命氣息的血肉之軀，地上的一切都要滅亡。 \end{tabularx} \\ \\ \relax
6:18 & \begin{tabularx}{0.7\textwidth}{X} 但我要與你立約；你同你的兒子、妻子和媳婦都要進入方舟。 \end{tabularx} \\ \\ \relax
6:19 & \begin{tabularx}{0.7\textwidth}{X} 凡有血肉的動物，每樣一對，一公一母，你要帶進方舟，好跟你一起保全生命。 \end{tabularx} \\ \\ \relax
6:20 & \begin{tabularx}{0.7\textwidth}{X} 飛鳥各從其類，牲畜各從其類，地上的爬行動物各從其類，每樣一對，都要到你那裡，好保全生命。 \end{tabularx} \\ \\ \relax
6:21 & \begin{tabularx}{0.7\textwidth}{X} 你要拿各樣可吃的食物，儲存在你那裡，作你和牠們的糧食。」 \end{tabularx} \\ \\ \relax
6:22 & \begin{tabularx}{0.7\textwidth}{X} 挪亞就去做了；凡神吩咐他的，他都照樣去做。 \end{tabularx} \\ \\
[1ex]
\hline
\hline
\end{longtable}
$^{1}$相信很多人都會做健身.
為了什麼呢?為了健康.
為了追求健美的身形.
我又不是 大家看我應該不是.
維持體能和減壓等等.
為了保持身形好的緣故.
你會很自律地去鍛鍊.
還有你可能會請了一些健身教練.
當你這樣做的時候.
因為你相信他.
還有你付了錢.
所以你一定會聽話.
照足他的吩咐去做.
照足他的指示.
因為你知道你這樣做的時候.
必定對你有很好的益處.
所以你是信得過你的健身教練.
但是對神又怎麼樣呢?.
如果我問你信神嗎?.
大家會怎樣回應?.
你信神嗎?.
有些很肯定 有些不知所措.
不知怎樣回答.
可能是心中說了.
我信 我信.
但是如果我再追問.
你是不是真的信得過神呢?.
這樣就去到另一個層次了.
就好像健身教練一樣.
你會不會跟足神的指示去做呢?.
面對逆境的時候.
你會不會仍然是信靠耶穌基督呢?.
如果你說不是的話.
又或者做一點不做一點.
這樣就不是真正信得過神了.
正如《雅各書》二章十七節這樣說.
我們一起去讀.
一二三.
信心若沒有行為就是死的.
一節大家很熟悉的經文.

$^{41}$大家都知道.
神的話對我們是有益處的.
提摩特侯書也說得很清楚.
但是為什麼有時候又不願意走出來呢?.
很想大家問一下這個問題.
我都自己去問自己這個問題.
好像聖靈所結的果子是什麼?.
忍耐 喜樂 和平 忍耐.
恩慈 良善 信實 溫柔 節制.
這是說一個的果子是九個的向度.
不是說我一向性格忍耐不了.
我又沒有溫柔 我又沒有節制.
是可以不去走的.
而是一個果子九個的向度.
都是同樣的.
我們都是要聽神的指示去做.
問題就是為什麼我們有時候做不到呢?.
原因是什麼呢?.
有可能就是我根本都不是很熟悉.
迦拉提書五章二十二節.
聖靈所結的果子是什麼.
今天你知道了.
你會不會放在心裡呢?.
走出來吧.
你自己問自己.
可能還有其他的原因.
希望今天我的分享訊息可以啟發到大家.
羅亞身處是什麼的世代呢?.
首先說羅亞.
羅亞的意思是休息的意思.
羅亞的爸爸叫做拉麥.
他為羅亞起名的時候就這樣說.
在創世記五章二十九節.
我們一起讀吧.
一二三.
這個兒子必為我們的操作和守中的奴婦.
安慰我們.
而且操作奴婦是因為耶和華咒助地.
意思就是說羅亞必定會使得我們免除痛苦.
就是從奴婦操作當中得到歇息.

$^{81}$休息的意思就是歇息.
並且從中得到安慰.
這裡說了為什麼人會奴婦呢?.
就是因為耶和華咒助地.
咒助地的意思是什麼呢?.
咒助地會長出一些荊棘.
農作物是很難生長的.
就算生長了也會很容易變壞.
而人類是要付上很大的勞力去耕種.
但是收種也是不多的意思.
事源是什麼呢?.
事源就是在創世記三章十七節.
阿當犯罪.
吃了禁果.
然後神就這樣跟他說.
地必為你的緣故受咒助.
你必終生奴婦才能從地裡得吃的.
這個咒助.
這個奴婦.
去到羅亞時代仍然是影響到的.
羅亞使得人得安慰.
得到什麼安慰呢?.
很大可能就是知道羅亞在洪水之後.
可以保存到人類和其他的生物.
成為一個新的里程碑的代表.
那時候的世代是什麼世代呢?.
是一個罪惡滔天的世代.
我們看看經文.
經文講到.
耶和華見人在地上罪惡很大.
終日所思想的盡都是惡.
還有在六章十一節裡補充了.
人類是敗壞的.
是充滿強暴的.
終日的意思是什麼呢?.
大家都意會到了.
不是一次過的.
是恆久的.
經常地的時時刻刻的.
而思想這裡講的是一些意念.

$^{121}$是籌算的意思.
整句說話的意思就是說.
當時每一個人時時刻刻都是犯罪.
隨己而行.
甚至所籌劃的籌算的都是惡.
不只是做的都是惡.
籌算的都是惡.
罪惡氾濫.
人人都是沉溺在罪的當中.
弟兄姊妹.
和今天的世代有沒有分別呢?.
大家認為.
我覺得不遑多讓.
只不過做出來的一些是很顯然的.
這裡是說.
是人都在做的.
大家都知道.
是明目張膽的.
但是今天的世代可能是暗暗的.
有個例子給大家.
最近電視每天都會講這個人.
四個字就是愛潑斯坦.
每天都聽的.
大家都知道.
這是美國一個所謂的淫媒.
做一些很恐怖的事情.
令人髮指的事情.
牽連的範圍很廣.
有國家元首.
有一些政治的領袖.
甚至是商界的領袖.
國家元首,商界領袖,政治領袖.
其實應該要做些什麼的?.
去建設國家要為國家人民謀福祉.
但是他們卻是沉溺在罪的當中.
他們以為暗暗的就可以做到.
但是爆了出來.
揚了出來.
弟兄姊妹.
今天我們有沒有暗暗的呢?.

$^{161}$人家聽不到.
只有我自己知道.
弟兄可能有弟兄的暗暗的.
姊妹有姊妹的暗暗的.
弟兄姊妹.
神不但看到人的外在行為.
祂還看到人內在的心思意念.
我們逃不過神的眼目.
所以我們千萬不要暗暗的.
以為神不知鬼不覺.
但是神是知道的.
一個值得每天去反省的問題.
我想去鼓勵大家的.
就是你平時中日思想的是什麼?.
是你自己看為好還是神看為好?.
我重複一次.
你平時中日所思想的是什麼?.
是自己看為好還是神看為好?.
很大的分別.
神的感受是什麼呢?.
耶和說後悔做人在地上.
心中憂傷.
大家可能會很驚訝.
神會後悔?.
後悔的意思不是做錯事那種後悔.
也不是早知今日何必當初那種後悔.
而是神感到遺憾.
感到嘆息.
甚至憂傷.
感到很悲傷.
很仇苦那種.
於是神的審判就來到了.
我們看看經文怎麼說.
有兩節的經文.
耶和話說我要將所做的人和走獸.
並昆蟲.
以及空中的飛鳥都從地上去除滅.
再強調多一次.
因為我做他們後悔了.
另外一節經文就是我要把他們和地一併毀滅.

$^{201}$神定要用洪水去毀滅世界.
剛才我問過一個問題 大家可能都在想這個問題.
為什麼神那麼殘忍呢?.
弟兄姊妹.
我相信神是萬分不願意用洪水去懲罰人.
他這樣做.
其實心裡都很惋惜很憂傷.
神就好像父母一樣 在座也有不少父母.
如果你的孩子或小朋友犯錯.
受刑罰的時候 你的心會不會很憂傷呢?.
會的.
你不忍心看到你的子女受苦 但你又不能不這樣做.
神也是一樣 神不能不這樣做 因為當時的人罪大到.
很驚人.
甚至去到一個.
不可以收拾的地步.
而在六章十三節.
聖經說.
去到盡頭.
去形容這一件事 所以神.
不能夠不去除滅.
人類和地上天空一切的活物 為了要剷除.
罪惡 不讓罪惡去蔓延.
在這個時候.
我們知道當世界將要被神審判的時候.
當要被毀滅的時候.
神做事的法則.
一向都是這樣的.
就是.
他同時之間.
會.
施行極度的審判.
以及會有極度的拯救.
神不會不將人類不利 因為神是愛.
原因就是神是愛.
神是滿有恩典憐憫.
好像耶利米蘇二十九章十一節所說.
神所懷的意念是什麼意念?.
平安的意念.
不是幹災禍的意念.

$^{241}$使得你們.
末後.
有指望.
神將盼望種在我們的心 如果我們沒有盼望.
我們很難.
去維持下去 很難在這個世間上.
生存的那種動力.
所以在這個時候.
出現了一個與眾不同的人物 大家都知道 就是羅亞.
聖經怎麼說他呢?.
羅亞.
一起讀 一二三.
「惟有羅亞在耶和華眼前蒙恩」 恩是什麼意思呢?.
恩的意思是恩寵.
因為.
甚至是好處.
為什麼.
只是羅亞.
在神的眼前蒙恩呢?.
之後那節就給了答案 我們一起讀.
一二三.
「羅亞是個異人.
在當時的世代是個完全人.
羅亞與神同行.
他的意思是什麼呢? 我們來看看.
異人的意思就是一個很正當.
好的人.
是做對的.
恰當的事情.
那你對什麼為之恰當呢?.
就是由神去設立 神說恰當就是恰當 神說對就是對.
完全的意思是完全無可指責 在當時的世代.
羅亞是相對.
完全的.
因為羅亞是.
與神同行.
一般的解釋就是.
持續地.
很親密地.
與神相交.

$^{281}$希伯來文的.
希伯來原文的意思.
又有少少不同.
意思就是說.
日日一起.
即是一起 與神來來去去.
與神散步.
即是意思.
很緊緊地.
很自然地.
去跟隨著神.
來來往往 即是神去到哪裡 你就走去哪裡.
神去哪裡 你就走去哪裡.
有一種散步的感覺.
在任何環境底下.
都可以與神去享受.
各位姐妹.
看到畫面嗎? 即是想像的圖畫.
可以很輕散.
可以很享受.
在當中.
好像詩篇23篇 清楚地溪水旁.
大家不是不知道.
大家應該都知道.
而是在歷史告訴我們.
人不是不知道.
只不過他們.
卻是.
不肯.
人竟自不肯.
聖經都提過很多很多次.
這麼好的東西 為甚麼人.
竟自不肯呢? 今天我都想問大家這個問題.
你肯不肯呢?.
還是你竟自不肯呢?.
希伯來書十一章七節.
總括了與神同行的深層意義.
看看這一節的經文.
我們一起讀 1 2 3.
諾亞因著信.

$^{321}$既夢神指示他未見的事.
動了敬畏的心.
預備了一隻方舟.
使他全家得救.
因此就定了那世代的罪.
自己也承受了那從信而來的疑.
希伯來書十一章七節.
意思就是說.
他.
這裡說的是.
這一節經文是對與神同行.
下了一個定義.
大家留意這三個字.
諾亞因著信.
欣著信與神同行.
縱然他未見過一生未見過這麼大的洪水.
但他相信大洪水將會來臨.
於是他就懷著敬畏的心去依靠神.
順命去做那隻方舟.
結果.
他在神的眼前去蒙恩.
全家得救人類和動物的存在.
可以得到延續.
同時之間.
也都開創一個.
參生的時代.
經文的重點就是這三個字.
欣著信.
而有的相應的態度和行為.
表示他是真實的相信.
而不是表面上信.
耶和華是神信耶穌基督是神.
這個就是我開頭所說.
我們是不是信得過神.
信得過神的意義就是這樣.
而伴著的就是.
伴隨著的.
就是與神同行是緊扣著一切的.
在創世的第六章和第七章我們可以看到.
羅亞欣著信與神同行.

$^{361}$有三個特點.
第一個特點.
剛才就是這個我所分享的.
第一個特點.
不疑惑的信靠.
首先在第六章.
第七章.
記錄了神與羅亞講話.
當中羅亞是這樣的.
他沒有駁嘴的.
我們經常所說.
沒有一言九鼎的.
我們看看聖經是怎樣說的.
首先.
神就對羅亞說在創世記六章十三節.
接著說什麼呢.
接著.
就向羅亞宣告.
要將人類與地一起毀滅.
然後吩咐羅亞.
你要做一隻方舟.
當中的規格是怎樣.
並且還與羅亞立約.
去一家百口以及各種其類的.
雀鳥,飛鳥,新竹,昆蟲等等.
都要一弓一武.
入方舟.
好保存性命.
和吩咐羅亞.
要預備食糧等等.
如是者羅亞就去做方舟.
接著到創世記七章一節.
這裡寫出來了.
耶和華又對羅亞說.
跟他講什麼呢.
神就說七天之後就要降雨在世上四十個晝夜.
就將各種活物從地上去除滅.
吩咐羅亞一家要入方舟.
同時除了之前我所說的動物之外.
還要加了一個要求是什麼呢.

$^{401}$是潔淨的畜類和空中的飛鳥.
要帶七弓七武.
之前我說是一弓一武.
現在好像多了一點.
潔淨的畜類和空中的飛鳥要帶七弓七武.
以至他們可以去流腫.
在當初.
羅亞可以有大條道理去質問神.
甚至去埋怨神.
這麼大的方舟要怎麼做.
這艘船要什麼時候造起.
是不是真的.
為什麼會有這麼大的洪水呢.
沒可能的.
很多時候我們會做的動作就是.
或者會問的問題.
給我一些印證.
也好.
到了後期.
七章一節之後.
神還說七天就要發生了.
時間會不會很倉促.
弟兄姊妹.
當事情不是你預期或者不是你掌控的時候.
你會怎麼反應呢.
而當教會的事情.
不合乎你的想法.
你又會怎麼反應呢.
你會不會埋怨啊.
會不會質問啊.
不是拉拉鏈 而是開拉鏈.
當然態度會不會是持著質問.
埋怨的態度.
應該這麼做 應該這麼做嘛.
大家看我呢 都是溫文爾雅的人 是不是.
好像不是.
都是斯斯文文 是不是 也很容易相處 認識我的人都會知道.
而且很容易合拍.
還有我有一些 弟兄姊妹也說 哇 Larry.
你的情緒智商真的很高.

$^{441}$不見你發脾氣的.
未曾見過你發脾氣 最低限度.
但是我告訴大家.
不要被我騙到.
我也有爆人的時候 不過不是經常.
很少有的 但也有的.
講回一個事件.
就是一年多前.
我負責.
教會負責一個項目.
這個項目開會 好像有七八位弟兄姊妹一起開會.
當中有個同工跟我說.
我覺得她是一個很挑釁的.
很埋怨的.
很挑戰性的.
一些語氣跟我說.
喂 你是做大那個.
時間那麼緊速 你要做決定了.
好問的時間 你一定要這樣這樣.
她好像在教我做事.
我當時是.
負責這個項目.
的IC.
有面子的.
於是 我真的爆了她.
真的受不了.
我爆了她.
我就說 你的態度可不可以好一點 你不要用這個脅迫.
用一個挑戰 一個埋怨的.
口吻跟我說話行不行.
接著.
在會議裡面的同工.
嘔了嘴.
閉了嘴 大家呆了.
這個不是Larry 不是我.
Larry平時不是這樣的.
當時的氣氛很僵硬.
於是就要開會了.
大家又不可以若無其事 心中還是有些東西.
我看到西式不對.

$^{481}$喂 先停一停 我們去洗手間.
先停一停 冷靜一下 安靜一下.
那時候 神.
聖靈就感動我.
跟我說.
你要跟這個同工說對不起.
我聽到之後.
當然就馬上去做了.
跟同工說.
不好意思啊.
我火爆了.
我真的跟她說sorry 不好意思我火爆了 我自己覺得.
可能埋怨的是我自己 可能態度不好的是我自己.
只不過我感覺到她態度不好.
她很緊張.
接著她說了一句話.
我們時常都會發生這些事情.
不要介懷 不要介意 我不會記在心裡面.
我們是對事.
接著呢.
不是對人.
如是者.
我們經常說一句話.
不打不相識 我就認識了這位同工 他又認識了我.
之後那個項目就很順利地進行.
問題就是.
我都會有一些缺點 我都會有一些.
死穴的時間 原來那些就是我的火爆點.
那些就是我的導火線.
我會失控.
那時候我會.
不是走在神的心意裡面.
我是信不過神.
在那些時刻當中.
但是這樣.
會發生的.
在座可能都會發生 但問題就是.
我們要是竭力保守.
聖靈所賜 合而為一的心.
我們要去.

$^{521}$和好.
我們要和好 要去.
求神在當中去幫助.
所以羅亞憑信心去聆聽 他沒有埋怨 他沒有一言九鼎.
就照著去辦.
充分體驗到他對神的敬畏和馴服.
在馴服神的話語之先 大家說最重要的是什麼.
我要明白神的心意 怎樣明白神的心意.
要聽到神的話.
羅亞怎樣分辨是神的話 不是自己的心意.
也都不是世界的雜音呢.
相信羅亞必定.
他是.
有聆聽神的聲音.
那種的操練.
弟兄姊妹我們都是一樣.
需要有持久穩定.
和神有約會那種的關係.
在安靜不受騷擾的環境裡面去親近神.
聆聽神並且去讀神的話語.
還記不記得羅亞的意思是什麼.
休息 休息的意思不是躺平 只是睡覺.
是安靜在神的裡面去親近神.
使得我們可以.
重新.
得力.
弟兄姊妹.
第二個.
是我們要活出神的價值觀.
世界的政治 工作 甚至家庭的環境都是講這一件事的.
選邊站 這裡講到.
可能你要選媽媽那邊 還是選爸爸那邊.
選爺爺那邊 或者選媽媽那邊 我們要懂得選邊站.
否則的話 很多事情就可能沒你的份了.
否則的話 可能你就會惹來.
很多的麻煩等等.
羅亞因著信.
他選擇在上帝那邊.
他持守上帝所設的法則.
他選擇一件事.

$^{561}$就是這件事 他是反潮流的 弟兄姊妹.
他不跟隨當時世界的價值觀 當時是什麼.
吃喝 玩樂 行淫 強暴.
隨意去縱慾.
那種風氣.
堅持神的標準 神的法則.
行在神的心意裡面.
最低限度有兩個值得我們學習的.
在羅亞的身上.
我們是可以看到的.
第一個 就是他持守一夫一妻的制度.
在《創世記》七章.
十三節 這裡也說了.
羅亞一家八口.
四對夫婦是入方舟的.
大家都有數得計了 一家八口 四對夫婦 即是一夫一妻.
所以很清楚說到羅亞.
一家是一夫一妻那種關係.
第二個.
他敢於面對很多人的冷嘲熱諷.
當人人在作惡.
那時候羅亞不但沒有跟隨.
默默地去做他的方舟.
神要求他所做的方舟.
難免他會被當世的人說的嘲諷.
那種排斥 但是羅亞頂住壓力.
毫不搖動.
很專心地去做他的方舟.
各位姐妹 現今的世代社會講什麼的.
多元主義 大家認不認同.
自由主義.
我想怎樣就怎樣.
神所設立的一夫一妻的制度.
已經被扭曲了.
人種育 同性婚姻都可以合法化.
連大麻都可以合法化.
還有一些以基督教立國的國家.
竟然可以容許去基督化.
怎麼搞呢.
小學或中學不可以在學校讀聖經.

$^{601}$各位姐妹.
神嘆息.
神憂傷.
不是因為神沒有能力去改變.
而是因為他賦予人類那種自由意志.
可以去做選擇.
他不想我們做現在很流行的.
機械人 AI robot.
不想我們這樣.
否則我們做人是沒什麼意義 沒什麼意思.
他原本希望人可以選擇.
去進行他的命令.
以致到人類可以一生的蒙福.
享受神的祝福.
同行來來去去去散步.
那種享受當中.
但是人類卻是一再選擇去搏逆神.
不聽祂的吩咐.
不走在祂的真理裡面.
當.
人脫離了上帝的法則.
表面上是不是好像很自由.
自在無憂無慮 暗中地又可以做這些事.
結局.
就好像羅亞的世代一樣.
失控.
混亂.
死亡.
自己要承擔後果.
好比.
一隻紙妖 大家都放過紙妖.
如果紙妖不想被約束.
不想被放紙妖的人.
去擺佈.
會怎樣呢 一旦那條線.
斷了的時候.
必然會失控.
必然它的作用就沒有了 同時之間是會毀滅的.
我相信大家都 我分享過我曾經移民過加拿大.
我移民加拿大之前要做一個抉擇.

$^{641}$我那時候是住公屋的.
住公屋 我是那個戶主 用我的名字.
那個單位.
我有一個抉擇.
我要做的就是.
需不需要將那間屋的單位還給政府呢.
那時候我的家人.
就會跟我說 阿欣 你有沒有想過 我叫阿欣 王啟欣 我個人叫阿欣.
阿欣 你有沒有想過你將來會不會回來.
會不會有一間公屋 給公屋政府 你有沒有想過.
我那時候的意念就是一個.
我就想過.
我那時候已經是基督徒了 雖然我不是很投入.
那時候教會的生活.
但是.
我心裡跟自己說.
我要做一個正確的事情.
做一個正確的事情就是我要將公屋交還給政府.
同時之間會想到 很多人在排隊.
我這樣做是應該的.
然後我的親戚朋友跟我說.
唉 你真是傻子 你怎麼這麼傻.
你回來後真的沒有房子住了.
同時之間那個單位是非常好的 貝山面水.
看到很美的海景.
公屋單位都很小的.
OK.
於是去了加拿大 大概四年多.
回來 積蓄不多.
錢不多 盤穿不多.
但是那時候的租金都很貴.
真的想到 又要租房子 又怎麼樣呢.
我的親朋戚友又跟我說.
是不是呢 Larry 阿欣.
都說叫你不要這樣 你是不是傻子呢.
當很多人跟我這樣說的時候.
這個好像是潮流.
你可以暗暗地去做這件事情.
但是.
我卻是.

$^{681}$怎麼說呢.
我只覺得應該要反潮流.
我不應該跟著這個世界.
我應該要做正確的事情.
我那時候就是這麼簡單.
要做正確的事情.
OK.
第三點.
第三點是說到.
甘願冒付出代價的風險.
冒付出代價的風險.
羅亞恩爵信 我要記住這三個字.
恩 爵 信.
義無反顧 堅持不屈地去遵行神的吩咐.
看聖經怎麼說.
我們一起讀.
羅亞就這樣行.
凡神所吩咐的.
他都照樣行了.
不單止一節經文.
其餘有.
另外三節經文 都是說到同樣的.
一件事情 在過程當中.
相信羅亞會有.
不少 是大大小小的困難.
例如.
龐大的物料的需要.
做這麼大的船.
儲糧在哪裡找 還要種.
怎麼去運輸呢.
做船的專業方法是怎麼樣呢.
什麼時候需要員工 都會有壓力.
人手問題.
長時間建造那種.
就是做的.
放奏的那種持久力.
還要安排各類的動物.
入荒州妥善的.
幫我們.
存放.

$^{721}$還有.
這個我相信才是難頂.
長期的陷入一個孤單的.
境況裡面 大家都知道.
除了他之外.
他一家八口之外 其他的人都是怎麼樣.
長期孤單.
但是羅亞怎麼樣 羅亞你看看.
羅亞他是照足神的規格.
他也照樣行 神吩咐他也照樣行 怎麼做怎麼做.
他照樣行 他沒有靠自己的聰明.
丁姐妹.
我們有沒有靠我們的聰明呢.
我們在工作裡面有沒有靠我們的聰明呢.
我們在我們的.
教會裡面有沒有靠我們的聰明呢 有沒有把我們工作或者在世間的價值觀帶入教會呢.
例如我們有沒有出位呢.
博出位.
有沒有競爭呢 有沒有山頭主義 或者有沒有.
我要拿這個資源 拿這個資源 我要認第一 認叻.
這些可能會發生.
但是有一個我相信.
大家都不可以.
忽略的 或者大家都是其中的 包括我自己 是什麼呢.
剛才羅亞的意思是什麼.
休息.
聖經在詩篇第46篇說到.
你們要休息.
要知道我是神.
你們要休息 要知道我是神 我們做神的事.
我們工作其實都是等同是事奉神的 不單止在教會.
都是等同我們去事奉 做神的事.
我們不單止是做.
我們還要做得對.
當我們神說.
你們要休息 要知道我是神 我們有多好可以做到呢.
包括我自己.
兩天前我去看醫生.
我時刻都要看著我的血壓.
有高低.

$^{761}$這次又高了一點 當然不是很高 紮戒那種.
他跟我說.
Larry 你平時睡覺睡得怎樣.
最近這兩個禮拜都不太好 有時候都會驚醒.
他問你幾點睡覺.
我說最早都十二點多 有時候一整點.
他就苦苦的跟我說.
Larry.
你睡眠要睡得好.
知道你有很多壓力 有很多事情需要照顧.
但是你需要.
休息.
所以他跟我說.
你睡覺之前的那幾個小時.
不要看WhatsApp 別看其他東西.
安靜 他沒有叫我.
這樣說安靜 總之叫我要休息 我知道這件事情.
我是要.
雖然是繁忙.
但是神的吩咐是什麼.
你們要休息 要知道我是神.
我相信我們很多弟兄姊妹都應該做不到這件事情.
今天我想分享給大家 大家彼此激勵.
我都要好好的聽這句話.
另外就是不沮喪而半途而廢 他很孤單很寂寞.
做船做了很久 相信都超過幾十年.
要看聖經 都有一些顯示 大概是這樣.
他沒有沮喪.
沒有半途而廢 弟兄姊妹.
一樣 我們是否會沮喪 跟他不合.
我都是.
退了 退了 退了.
昨天說到之前 崇拜之前 有代禱者為我禱告.
我很注重這個時刻 因為需要很多禱告 很多力量去支持.
他跟我說 有位姊妹 代禱者.
Larry 見到你很累.
見到你樣子很累 你需要休息 又是休息.
又是跟我說.
接著.
他說 你不要沮喪.

$^{801}$你不要沮喪 要耕耘賜力.
我不可以沮喪.
都是跟我說.
另外有兩個大問題 相信.
羅亞是會面對的 第一個 就是他怎樣叫家人入方舟呢.
今時今日.
可能叫子女跟父母住都很困難.
或者你是子女的.
你都可能不太願意跟你父母.
你的祖父母住.
還有.
有時你都說 老婆.
叫老婆做這件事 老婆不一定會聽你的.
或者你叫老公做這件事.
不一定會聽 現在叫我們入方舟.
都看似容易 但都不容易.
還有.
在彼得厚書二章五節說到 他要傳義道給當代世人.
容易嗎.
那時惡人當道 又敗壞又強暴 你跟他講福音 會否被人殺 斬了他.
好像早期的宣教士一樣.
會否有人聽呢.
如果沒有人悔改的話 羅亞要忍受被拒絕和沒有人相信的那種挫敗感.
面對以上種種的難題.
羅亞不是不知道要冒.
負重大 很大的代價 那種風險.
但他照做.
沒有退縮 為什麼 因為他知道神要他做的事情.
神一定會給他能力 一定會給他資源 同時之間神會包底.
羅亞要做的就是.
憑信心.
就是做.
但還要用神那種方法去做 不要用自己的方法去做.
就是這樣 做正確的事情是有風險的.
我不知道大家.
同意嗎 在工作界裡面 如果有同事說要買六合彩 你會買嗎.
我不知道 如果你不買的話 同事可能會說你不合群.
還有一些正確的事情.
曾經有一個情況.
我曾經有一份工作.

$^{841}$是做物流的主管.
當時有一批貨 有個貨櫃被扣關.
被扣關 因為有些文件不對 報關的文件錯了.
但那批貨很急趕 要送到客人的手裡.
於是我的老闆就跟我說.
Larry 你知道怎麼做的吧 這件事.
大家會怎麼意會 你知道怎麼做吧.
我意會 當然啦 每個人都是這樣做的.
就是桌底交易.
我是意會的.
但我問他一個問題.
就是怎麼做.
你說我知道怎麼做.
他沒有說話 很大力地關了門.
bang 一聲 全寫字樓都聽到.
我是第一次被我老闆這樣霸門的.
所以有風險的 信耶穌.
這是我們的原則.
但信得來 我們卻是有恩典 有憐憫 有祝福在裡面.
當我們做正確的事情的時候.
就會在神的眼前去蒙恩.
當然這件事我沒有做到.
桌底那些東西.
而那批貨不是很久就出來了.
只不過我老闆很心急 想要送到客人的手裡.
而客人也沒有很大的投訴.
在這裡有人讚美我們的主.
最後一個總結給大家.
神對我們的心意是什麼呢.
是我們在日常的生活裡面.
在我們的工作裡面.
與神同行.
意思就是說凡神所吩咐的.
我們就是憑信心.
並且存著敬畏 馴服的心去進行.
這樣才是我一開始所說的.
或經常所說的 這樣才是信得過神的表現.
而不是留於信有神.
我信有耶穌.
但是所做的一切.

$^{881}$與耶穌亂不上關係.
當我們願意這樣做的時候.
我們都可以像螺一樣.
能夠在耶和華所咒詛的地上.
必使我們從工作和手裡面的勞苦得到安慰.
自己在神的眼前蒙恩.
得到祝福.
別人也可以因著你得到祝福.
有兩句說話我想和大家分享.
第一個是巴斯卡.
如果讀科學的人都會知道.
這是一個壓力的單位.
所以這個人物都非常出名.
他是法國的神學家,哲學家和數學家.
他說了一句說話.
每個人的心裡都有一個無限的空洞.
唯有無限的上帝才能夠填滿.
另外一個就是奧古斯丁.
奧古斯丁大家都應該聽過.
四至五世紀一個神學家和哲學家.
他這樣向上帝禱告.
主啊,你為自己的緣故創造了我們.
只要找到你.
我們那種焦躁不安的心.
才能夠在你裡面得享平靜和安息.
弟兄姊妹.
大家可以咀嚼這兩句說話.
上帝的創造的原意是人類要對上帝有渴望.
還有上帝對人類都有渴望.
不是其他的事物可以取代的.
當我們憑信心與神同行的時候.
我們的渴望,神的渴望才能得到滿足.
我們一起去回應上帝.
因為你安排了很多很多的路.
就是我們可以與你同行.
你亦都很渴望,渴求與我們同行.
所以幫助我們.
在我們的生命裡面.
我們願意與你同行.
願意去時刻思想.

$^{921}$你在我們的生命裡面.
是在甚麼位置.
是否最重的位置呢.
是最想蒙你喜悅的位置.
所以神,將我們每一個人都交上.
求你幫助我們.
叫我們因著你,我們要堅持.
因著你,我們不被世間事物所影響.
而是願意的.
去到你的私人寶座面前.
每天去投靠你.
不用自己的方法,不用自己的心意去行.
而是只能行在你的心裡面.
去蒙你的喜悅.
我們今日向你納至.
我們今日在我們思想裡面.
神,你與我們去講說話.
讓我們一一記在心裡面.
成為我們向前行的動力.
\newpage



\section{哥林多後書 1:15-24}
\label{sec:031639}
\textbf{第二次的恩典}~~~~陳雅莉傳道
\newline
\newline
連結: \href{http://www.church.com.hk/acms/content.asp?site=cdc\&op=show\&type=product\&code=031639\&layout=sermon}{\texttt{031639}} ~~~~ 語音日期: 2026-02-08
\newline
\newline
\hyperref[sec:031638]{\small{< < < PREV SERMON < < <}}
~
\hyperref[sec:index_chronic]{\small{[返順時目]}}
~
\hyperref[sec:031640]{\small{> > > NEXT SERMON > > >}}
\newline
\newline
哥林多後書 1:15-24
\newline
\begin{longtable}{cl}
\hline
\hline
章節 & 經文 (和合本修訂版)\\
\hline
1:15 & \begin{tabularx}{0.7\textwidth}{X} 既然我這樣深信，早就有意先到你們那裡去，讓你們得加倍的益處。 \end{tabularx} \\ \\ \relax
1:16 & \begin{tabularx}{0.7\textwidth}{X} 我要路過你們那裡往馬其頓去，再從馬其頓回到你們那裡，讓你們給我送行往猶太去。 \end{tabularx} \\ \\ \relax
1:17 & \begin{tabularx}{0.7\textwidth}{X} 我有此意，難道是反覆不定嗎？難道我的意願是從私慾起的，以致我忽是忽非嗎？ \end{tabularx} \\ \\ \relax
1:18 & \begin{tabularx}{0.7\textwidth}{X} 我指著信實的神說，我們向你們所傳的道並非又是又非的。 \end{tabularx} \\ \\ \relax
1:19 & \begin{tabularx}{0.7\textwidth}{X} 因為，我、西拉和提摩太在你們中間傳神的兒子耶穌基督，從沒有「又是又非」的；在他只有一個「是」。 \end{tabularx} \\ \\ \relax
1:20 & \begin{tabularx}{0.7\textwidth}{X} 神的應許，不論有多少，在基督都是「是」的。所以，我們藉著他說「阿們」，使神因我們得榮耀。 \end{tabularx} \\ \\ \relax
1:21 & \begin{tabularx}{0.7\textwidth}{X} 那在基督裡堅固我們和你們，並且膏抹我們的，就是神。 \end{tabularx} \\ \\ \relax
1:22 & \begin{tabularx}{0.7\textwidth}{X} 他在我們身上蓋了印，並賜聖靈在我們心裡作憑據。 \end{tabularx} \\ \\ \relax
1:23 & \begin{tabularx}{0.7\textwidth}{X} 我指著我的性命求告神作證，我沒有再往哥林多去是為了要寬容你們。 \end{tabularx} \\ \\ \relax
1:24 & \begin{tabularx}{0.7\textwidth}{X} 我們並不是要控制你們的信心，而是要作你們的同工，讓你們得快樂，因為你們在信仰上已經站得穩了。 \end{tabularx} \\ \\
[1ex]
\hline
\hline
\end{longtable}
$^{1}$請師妹平安,感恩.
讓我們一起回到加神的殿裡面去敬拜神.
上星期梁銓導帶我們看了《哥林多後書》.
不知道大家有沒有記得呢?.
他跟我們看了《約書亞記》和《哥林多後書》.
上星期梁銓導帶我們看《哥林多後書》.
一段很經典的經文.
其實我們說《哥林多後書》是有很多很經典的經文的.
其中上星期看了的就是.
保羅怎樣在軟弱上面變得剛強.
是神怎樣用祂的恩典在軟弱的人身上.
將恩典變得完全.
其實當我們認真去看《哥林多後書》的話.
我們會發現其實這個信息幾乎是貫穿整本書.
整卷書.
而從我們今天的經文我們就看得到.
這封書信在開始的時候.
剛才主席讀的時候.
我不知道你會不會覺得很複雜.
很多又是,又非,又忽是忽非,又是的是是的.
好像很複雜.
其實剛才我們讀的經文你可以看到.
似乎保羅一開始就很努力的為自己去澄清一些東西.
去澄清一些東西,去辯護一些東西.
從一張17,18字我們看到.
他說有這個意思.
由此難道我是反覆不定嗎?.
難道我的意願是從私欲起的嗎?.
難道我是忽是忽非的嗎?.
你似乎看到他想解釋一些東西.
想澄清一些東西.
甚至有一種意味是.
我在為自己辯護中.
他甚至指著信實的神說.
我向你們所傳的道不是有是有非的.
他似乎是在辯護中.
奇怪了.
為什麼保羅要無緣無故去辯護呢?.
要去澄清呢?.
如果當你去想像.

$^{41}$你什麼時候會去跟別人辯護或者澄清的?.
就是當你遇到一些誤會的時候.
或者可能有人指控你的時候.
一些無理的控告的時候.
可能你就會為自己去辯護.
保羅也一樣.
當時他被人攻擊.
他遇到一些惡意的批評和指控.
來到這裡就更奇怪了.
哥林多教會是保羅親手建立的.
是保羅親手所建立.
他不是很愛這個教會嗎?.
他不是處處都為這個教會著想嗎?.
從哥林多前書.
我們還看到保羅說要去探望他們.
神差派一個這麼好的僕人.
去他們中間教導他們.
幫助他們 建立他們 餵養他們.
應該感恩才是.
為什麼他們竟然會去攻擊保羅呢?.
原來也有原因的.
在哥林多前後書我們對比的時候.
就發現原來保羅曾經在哥林多前書.
有講過他一些行程.
是他要先去馬其頓.
然後就會去哥林多教會探望他們.
還說要和他們一起過冬.
令到哥林多教會一直在等待他們來.
又很期待他們來.
但是來到後書我們就發現.
保羅不是去了馬其頓再去哥林多.
而是他先路過了他們再去馬其頓.
然後他又會再從馬其頓回到他們那裡.
就是這個行程上的改變.
一個保羅計劃上的改變.
直到他被人攻擊.
有人惡意對他批評和指控.
說他是一個相當不可靠的人.
他經常都反覆不定.
又是又非.

$^{81}$又說來又說不來.
不來又來.
就是這樣.
他們就攻擊他的.
甚至攻擊他不單只是改變行程這麼簡單.
說他是從私慾起義.
說他是自己有私慾.
要起義這樣去做.
一個愛教會.
愛哥林多教會的使徒.
他改變他的行程.
但就成為了被攻擊,被誣告.
甚至一個屬靈的人.
按著屬靈的事,按著屬靈去行事.
竟然被誣告成他是按私慾.
按著他的肉體去行事.
其實保羅是不是真的出爾反爾.
行事這麼反覆不定.
讓人無所適從呢?.
其實從哥林多後書我們一路看下去.
今天我們看的是第一章.
但是整本書一路看下去.
你會發現不是.
保羅不是真的忽是忽非的人.
他也不是一個出爾反爾的人.
保羅原來改變行程是有原因的.
因為其實他已經發現.
哥林多教會對他有猜疑,有猜忌.
因為他之前對他們有過嚴厲的責備.
這在哥林多後書二章有記載.
保羅曾經對他們有過很嚴厲的責備.
又托提多帶了一封嚴厲責備他們的信.
所以保羅就一直在等待消息.
看看他們會不會回轉.
看看他們有沒有機會回頭悔改.
所以保羅延後他的行程.
改變他的計劃.
保羅說其實他早有意去那裡了.
去你們那裡了.
不過其實他是讓他們得加倍的益處.

$^{121}$他改變行程是為了讓他們得加倍的益處.
而加倍的益處在原來的文字.
直譯為第二次的恩典.
保羅是想給哥林多教會第二次恩典.
其實他去建立教會.
去造就他們,教導他們,輔助他們.
已經是神給哥林多教會莫大的恩典.
但當他們有走迷的時候.
有迷失的時候.
有走差的時候.
保羅責備他們,提醒他們.
保羅不再去他們中間停一停.
等一等.
改變了行程.
為的是要給他們機會回轉.
給他們時間改變.
保羅的等待.
是為了要給他們第二次的恩典.
甚至保羅說.
我指著我的性命求告神作證.
我沒有再來哥林多.
不是忽是忽非.
不是反反覆覆.
而是為了要寬容你們.
他是有意思不去.
他是有原因.
他是為了愛教會給他們空間挽回他們.
等候他們改過.
是為了寬容他們不再去.
因為保羅一去到.
見到他們未改,未回轉.
他就會再責備他們.
在哥林多後書.
保羅甚至覺得自己曾經對他們的責備太重.
恐怕他們受不了.
所以保羅願意這樣去.
給他們有第二次的恩典.
為了要寬容他們.
他才改變行程.
改變他們的計劃.

$^{161}$但我們發現.
這種寬容.
這種願意再給他們第二次恩典.
給他們再多一次機會的好意.
竟然被曲解了.
竟然相反成為他們攻擊保羅.
責備保羅.
質疑保羅.
在當中在信徒當中混亂.
傳謠言去重傷保羅.
可能我們來到這裡也會問一問.
其實改變行程有什麼大不了呢?.
保羅也有他的自由去邊的.
改變行程有什麼大不了呢?.
難道保羅不可以改變的嗎?.
這個攻擊似乎不太合理.
不過我們會發現.
其實當哥林多教會去攻擊保羅的時候.
相信已經不是因著他改變行程那麼簡單.
而是當中早早就已經有人.
不滿他的責備.
不滿他的教導.
不喜歡聽責備的話.
還要是嚴厲的責備.
這個改變行程.
直到他們攻擊保羅,誣告保羅.
只不過是一個導火線.
只不過是一個藉口.
給他們一個機會.
我還記得.
可能就是這些.
保羅曾經責備過他們.
以致他們有不愉快的經歷.
相信已經是積怨了可能一段時間.
所以保羅這次.
這個改變行程就成為了一個把柄.
我們甚至會想.
即使是改變行程.
可能我們會覺得他誠信有問題.
他說話出爾反爾.

$^{201}$但他們甚至攻擊保羅.
起意願是從私慾起的.
這個謠傳一直在信徒中間傳開.
牽連的人可以相信是很多.
至少保羅要寫信去澄清.
他們指控保羅是按肉體行事的人.
其實這個對於一個屬靈的領袖來說.
不是一個簡單的指控.
是非常嚴重的指控.
是相當嚴重的.
一個屬靈的領袖,一個事奉神的人.
竟然被指控為.
他是按著私慾,肉體去行事.
我深信保羅真的經歷了.
面對哥林多教的時候.
他難免真的會想.
我給你們的良心.
你們當狗吠.
我們宗教會這樣說.
我的良心你當狗吠.
我的好意竟然被你曲解.
好心無好報.
確實在我們跟隨神的路上.
我們事奉神的人的路上.
難免都會有這樣的經歷.
難免會有被誤會的時候.
被中傷的時候.
可能沒有保羅那麼嚴重被攻擊.
但難免都有被人說閒話的時候.
被人不滿的時候.
實在是 特別是為主的緣故.
你盡心去事奉.
盡心盡力去愛主.
盡心盡力去傳福音,作見證.
都難免會遇到逼迫.
難免遭遇到誣告.
如果你是保羅.
你的反應會是怎樣呢?.
可能你會發脾氣.
很生氣.

$^{241}$我要罵人.
或者可能你會很傷心,很難過.
受夠了.
很厲害嗎?我從此不事奉了.
事奉那麼艱難.
不做了 踢爆不生氣.
或者會很生氣.
裡面儲起了一些苦毒.
儲起了一些冤仇.
對著你弟兄姊妹.
可能從此就成為陌路人.
不想見到對方.
見到就眼淵.
會不會是這樣呢?.
確實無心無好報.
這樣被重傷.
確實是很傷害.
不能否認這些感受是很真實的.
如果你是保羅.
你還願不願意愛這些弟兄姊妹呢?.
如果你是保羅.
你還會不會回哥林多教會?.
你還會不會去探望他們?.
你還會不會寫信給他們?.
從聖經裡我們看到.
保羅沒有放棄.
他仍然寫哥林多後書給他們.
我記得從前在我初信的時候.
我曾經聽過一位屬靈的長者跟我分享.
當時我初信.
回到教會.
有屬靈的長者跟我聊天的時候.
看到教會有很多的問題.
很多的不足.
很多的缺陷.
於是我就問他.
當時很天真.
初信無狼我就問他.
教會這麼多問題.
你這麼不滿.

$^{281}$又有這麼多不足.
又這麼缺乏.
為什麼你不轉教會?.
你走不就行了?.
為什麼我會這樣說呢?.
因為我的母會.
旁邊就有一間教會.
旁邊是旁邊.
真的走出門口走三步.
就有另一間教會.
我們在同一層樓有兩間教會.
同一個位置我們共用.
然後進入兩個門口.
是可以轉教會的.
走不就行了?.
誰知道這位屬靈的長者.
很有趣他跟我分享.
他說教會是他屬靈的家.
他說家裡有事.
家裡有事.
他不會只是懂得問為什麼.
他不會只是問為什麼.
而是他會問我當作什麼.
他跟我分享.
他說弟兄姊妹之間有嚴規.
我不是只是在這裡吃花生看電影.
我有沒有圍著弟兄姊妹之間的嚴規.
好好去禱告神.
甚至乎有人和我有嚴規.
我有沒有去尋求神.
我看到教會裡面有傷害.
有不滿的時候.
我是不是只是在這裡議論紛紛.
在背後說三道四.
還是我願意在當中做和平之子.
願意在當中去禱告守望.
他又跟我分享.
當看到教會奉獻不足的時候.
我是不是只是坐在這裡.
什麼都不做.

$^{321}$而是應該好好地.
為他教會需要供應禱告.
好好求神.
讓我們有更甘心樂意奉獻的心禱告.
他又說.
當他看到大家不勤奮讀經.
沒有人想追求的時候.
我是不是只是坐在這裡問為什麼.
還是我會努力鼓勵大家一起讀聖經.
招聚弟兄姊妹一起來讀.
陪伴那些讀不明白的人一起讀.
我有沒有將一切的問題.
一切我看到的事.
放在神的面前.
好好禱告.
求主幫助我們.
他跟我說.
他不離開.
他不離開教會.
不是因為他沒有選擇.
大把選擇.
也不是因為教會什麼問題都沒有.
很好.
所以要繼續留下.
而是.
他說這個是神給他屬靈的家.
除非有一天.
神要帶他離開這個家.
如果神帶他走.
他就是順服神的帶領.
去別的地方.
但如果不是.
他一天還是這個.
他說他仍然會每一天都問.
主啊.
我當作什麼.
是.
確實.
如果我們問為什麼.
真是很論盡.

$^{361}$但如果我們願意問上帝.
主啊 我看到這一切.
那你怎麼使用我.
我當作什麼.
我聽過一位師母說.
墨西珍師母她曾經說過.
她說為什麼會叫我們消極.
退喪 憤恨 退縮.
我當作什麼.
會叫我們清醒 積極 果敢 前進.
我相信保羅也是.
他遇到哥林多教會的攻擊時候.
我相信他不是只是坐著問.
為什麼.
為什麼他們攻擊我.
我相信保羅不是這樣問.
我深信保羅是問.
主啊 我當作什麼.
於是他當作什麼呢.
他就寫信給教會澄清.
去嘗試澄清 嘗試為自己辯護.
解釋給他們知道.
他不是一個出爾反爾的人.
他不是一個忽是忽非的人.
保羅不單止為自己辯護.
因為當時的人已經去到一個地步.
不單止是誣告保羅.
甚至乎誣告保羅所傳的福音.
甚至乎誣告保羅所講的道.
保羅要鄭重地跟他們聲明.
保羅說我主自信實的神去講的.
我向你們所傳的道不是有是有非的.
耶穌基督不是有是有非的.
沒有 沒有有是有非.
神的應許無論多少.
在基督裡都是時的.
他要發出一個鄭重的聲明.
就是我在你們中間所傳的.
沒有有是有非 沒有忽是忽非.
在耶穌基督裡都是時的.

$^{401}$保羅再一次強調.
不單止他傳的道.
他的人是這樣 他傳的道是這樣.
他更加用哥林多教會作為見證.
哥林多教會是因他所傳的道而去相信.
而建立成為教會.
在基督裡 神在基督裡堅固他們.
這個教會不只是保羅自己一個人.
是整個教會有弟兄姊妹在裡面.
他們都是被聖靈 被耶穌基督所堅固的.
被聖靈蓋上了印.
被聖靈在弟兄姊妹心裡作憑證的.
不是保羅忽是忽非.
更加不是他所說的道忽是忽非.
弟兄姊妹你們都相信了.
為什麼會被一些誣衊.
一些重傷 一些攻擊就混亂了你們呢?.
保羅的意思是這樣.
保羅作為基督的僕人.
他怎麼會用忽是忽非的態度.
有是有非 出而反而的態度.
來去傳揚福音呢?.
甚至保羅最後再澄清.
我指著我的聖名求告神作證.
我沒有再來哥林多教會.
不是因為我忽是忽非.
是為了要寬容你們 是為了要寬容.
我們可能會覺得這一番澄清的說話.
保羅憑著良心說.
憑著神的見證說 憑著聖靈說.
憑著他在三位一體 真神的面前作證說.
但是有時真是被攻擊 被誣告.
會有理說不清的我們發現.
真的會有這樣的情況.
所以當我們在教會裡面的時候.
我們真的發現.
我們很需要謹慎.
好好謹慎我們的言語.
好好謹慎我們的舌頭.
有時候當我們說出一些誣告的說話.

$^{441}$一些中傷的說話.
甚至傳來傳去 將一些傳言傳出來的時候.
我們是不是真的能夠確實知道.
其實所有事情背後發生的是什麼呢?.
難道我們真的覺得我們無所不知?.
其實有沒有發現 我們見到的都只不過是很片面.
我們所聽到的其實都是.
不知道那件事的百分之幾多.
可能百分之一都沒有.
但我們是不是就按著我們所看到的.
所聽到的 就隨便去論斷呢?.
我們豈能不代替神 作審判官呢?.
在教會裡面 在弟兄姊妹之間的交流裡面.
不論是公開的 私下的.
我們都要謹慎.
不要輕易說中傷人的說話.
正所謂「謠言止於智者」.
如果有謠言傳到你那裡的時候.
應該就要在你那裡去止息.
在教會裡面.
更加是應該止息在那些凡敬畏神的人身上.
不只是智者那麼簡單.
當我們說我們是敬畏神的人.
我們是屬神的兒女.
我們是怎樣去敬畏我們這位神的呢?.
箴言說:多言多語難免有過.
節制嘴唇是有智慧的.
在弟兄姊妹之間發生一些小事.
我們很容易就議論紛紛.
不過當然 在教會裡面.
我們又不是不說話的 是不是?.
都很難吧?我回來叫你不說話.
剛才叫大家回來會怎樣?.
彼此問候,關心 是不是?.
但我想我們懂得怎樣分.
甚麼是關心,問候.
和甚麼是諸事八卦.
我們懂得分甚麼是真心.
是想關心弟兄姊妹.
還是想多傳幾句說三道四.

$^{481}$求主幫助我們.
我曾經看過一個很有趣的統計.
他說一個人平均去搜集一些資料.
去統計一個人平均一天說多少分鐘閒話.
統計發現一個人平均每天說52分鐘閒話.
你說哪有啊 是不是?.
我沒有啊.
平均數來的 我真的希望我們教會裡面不是.
是平均數來的.
這些閒話裡面當然都包括了一些有益處造就人的說話.
不過亦都包括了更多破壞性的,傷害性的.
我們真的要謹慎.
在《異忽所書》說.
一句壞話都不可以出口.
一句都不可以.
只要隨著需要說造就人的好話.
不是不可以說話.
是說好話 造就人的好話.
讓聽見的人得益處.
《雅各書》亦都提醒我們.
舌頭是我們肢體裡面很小的 是不是?.
很小 我們的人這麼大 舌頭很小.
但是它卻能說很大的話.
不是說謊話 是很大的話.
真的 我們可以很小的口說出大到不得了的話.
一個小小的火就能夠點著一個森林.
我相信這個我們不難明白.
我們用口唇進讚我們的神 是不是?.
我們怎能夠同一個口出詛咒的言語.
誣告別人,重傷別人的言語呢?.
是不應該的.
《雅各書》直說這是不應該的.
確實是 人與人之間需要用很長的時間去建立.
我們有時會發現.
我們認識我們弟兄姊妹同一個人建立關係.
是在說以年計.
甚至乎用上幾年的時間.
我們都未必覺得真的與對方建立很深的關係.
或者很認識對方.
但是一句話就可以拆毀你所建立的一切.

$^{521}$求主幫助我們真的勒緊我們的舌頭.
有時候我們都很喜歡用我們自己的標準.
我們的看法 很難免的.
有我們的想法,看法去評論別人,去評斷.
求主幫助我們吧.
我們多看別人所付出的,所用心的.
而不是只看那個結果.
我們多看我們弟兄姊妹怎樣盡心盡力.
她背後的動機,留心她的需要.
多於只是看到她做出來的果效.
她能否完成那件事.
在教會裡面.
用言語傷害別人的事實在太多了.
我們真的要求神憐憫我們.
求神幫助我們,求神醫治我們.
禁止我們的嘴唇說不應該說的話.
傷害我們的弟兄姊妹.
如果倘若我們真的有受傷.
在我們中間.
你覺得你真的曾經受傷.
你真的覺得不容易面對.
甚至你真的被別人的說話傷害過.
我們都要學習.
不要只是將眼目放在那些傷害你的人身上.
我深信如果保羅被攻擊之後.
他只是將眼目放在那些誣告他的人.
罵他的人,說他的人身上.
他就不用走下去了.
但是他懂得轉眼仰望耶穌.
他懂得將他雖然受傷.
但他轉眼回去看神.
看耶穌基督.
他再一次說.
那個是基督裡面兼顧我們和你們的.
並且是高沒我們的.
是神.
他回去看神.
他不是看人怎樣攻擊他.
傷害他,誣告他.
他不只是憤恨地罵人.

$^{561}$他是告訴他們回來基督那裡.
高沒我們的是神.
我們有聖靈.
我回去看聖靈怎樣幫助我們.
回去看我們是一個在神面前怎樣的人.
我們是有神高沒的.
我們是有神兼顧的.
我們是有聖靈在我們心裡面為印記的.
是的,真的.
當你認定在神裡面那份保護.
認定神裡面的幫助.
認定你是屬神的人.
即使有再多重傷你的說話.
你感覺到受傷害了.
感受是真實的.
但不叫你失去盼望.
不叫你放棄繼續事奉神.
不叫你沒有動力.
繼續在教會裡面去付出愛.
而且你會發現當這些傷心難過.
讓你受傷害的事一次又一次發生.
而你又一次又一次回去神那裡.
一次又一次在神那裡得保護得安慰得提醒.
你會發現你的屬靈生命成長了.
你會發現你長大了.
你裡面的空間能夠勝在這一切的空間大了.
闊了.
是透過一次又一次回去尋求神.
透過一次又一次回去求神幫助你.
你得著神的幫助.
真實地經歷神跟你同在走過艱難的日子.
是這樣令你壯大.
是這樣令你的屬靈生命得以成長.
或者在人的看來.
這種我們會叫做抗疫力.
抗疫力是不會在順境建立出來的.
抗疫力是在逆境建立出來的.
而且抗疫力不是天生的.
它是可以操練的.
它是可以被培養的.

$^{601}$它是可以提升的.
但怎樣能夠培養提升建立呢?.
是透過面對逆境.
你會發現你自己一次比一次反彈力越來越快.
雖然受傷害,跌落低谷.
很快就能夠反彈上來.
這是經歷了一次又一次神與你同在.
你將你的信心投放在上帝那裡.
而祂一次又一次用聖靈的大能大力幫助你.
以致你發現終有一天.
你可以好像保羅一樣.
很快就反彈過來.
即使有委屈.
能夠仍然靠主站立.
保羅即使面對哥林多教會中的攻擊.
他不是說你們很厲害嗎?.
我從此不來了.
我不教你們了.
你們好好反省一下吧.
不是.
保羅竟然說.
我還要做你們的同工.
保羅說我不是要控制你們.
不是.
而是我要來做你們的同工.
保羅在受傷之後.
他不是選擇放棄.
選擇撇棄他們.
覺得他們沒救了.
保羅在受傷之後.
仍然選擇.
我要來做你們的同工.
保羅辨明我不是來管你們的.
是上帝有話要對你們說.
是上帝要教導你們.
而且他很相信.
其實大部份哥林多的信徒.
都是被迷惑的.
被那些小撮的人誣告攻擊.
傳來傳去.

$^{641}$這樣就影響了信徒.
所以保羅說.
其實你們在信仰上.
是已經站立得穩的.
不要忘記了.
你們本來就已經站立得穩了.
我要你再繼續做你們的同工.
使你們得喜樂.
保羅仍然要他們得快樂.
他們使保羅受傷傷心.
但保羅說我要來和你們一起.
使你們得快樂.
這個是保羅.
保羅能夠超越一切的傷害.
不以惡報惡.
他以愛還愛.
他選擇繼續去教導他們.
繼續和他們做同工.
繼續以基督的愛去愛這間教會.
這個是保羅愛哥林多教會.
他看出他們當中的軟弱.
明白他們有不足的地方.
保羅願意再一次和他們一起.
我要總結了.
從今天的經文裡面我們看到.
即使忠心盡心愛主的僕人.
仍然都會受到傷害.
甚至乎是無告攻擊.
如果你在事奉的旅程裡面.
受過傷害.
不要覺得稀奇.
這是使你更壯大.
我們又看到原來在教會裡面.
我們都不是原理.
我們難免都有軟弱的時候.
我們會出口傷害別人的時候.
我們會在我們自以為.
自以為自己很厲害的時候.
去攻擊別人.
求主幫助我們.

$^{681}$面對無理.
不合理的.
我們懂得依靠神.
仰望神.
繼續奔走.
我們傷害過人的.
不是不知者不醉.
今天你知.
神提醒你.
求主幫助你.
納緊我們的舌頭.
謹慎我們的言語.
倘若有傷害過的.
願主教導我們.
幫助我們.
能夠再去一次彼此認罪.
彼此饒恕.
從昔日的哥林多教會.
到今天的普世教會.
我們會發現.
根本沒有一間教會是完美的.
這個也正正反映了.
就是因為教會不完美.
就是因為我們是一群不完美的人.
我們才更加需要上帝的恩典.
我們更加需要上帝的恩典.
在我們軟弱的人上面.
顯得剛強.
最大的信心.
不是因為我們是一群很了不起的門徒.
最大的信心是因為.
我們有主耶穌基督.
教會之所以能夠成為明亮的燈台.
不是因為教會很了不起.
有什麼了不起.
是因為教會是神的教會.
有神在我們中間.
我們也不是什麼很了不起的人.
我們能夠了不起.
是因為有主耶穌基督在我們的生命裡.

$^{721}$有人的地方就自然有誤會.
有人的地方就自然有不明白的時候.
會有中傷的時候.
但這個絕對今天不是成為我們的藉口.
就是因為我軟弱.
所以我可以隨便攻擊人.
因為我軟弱.
我又控制不了的時候.
我就可以隨便出口傷人.
不是.
不是這個意思.
更加不是.
叫我們弟兄姊妹之間.
就是說我們要懦弱一點.
受傷的時候我們不要出聲.
死忍.
不是這個意思.
兩者都不是.
不要走極端.
保羅教導我們.
有過錯的時候.
就去承認.
有誤會的時候.
就嘗試去澄清.
用理性明辨.
在主裡靠主的安慰.
有被攻擊的時候.
求神幫助.
求神憐憫.
同時今天的經文也讓我們再一次去反思.
反思我們自己.
我也反思一下.
我一直預備這段經文的時候.
我也深深地反思.
如果我們曾經在言語上真的有過失.
真的有得罪人的地方.
讓我們學習在主裡面認罪.
甚至主也許可讓我們彼此認罪.
彼此饒恕.
求主幫助我們.

$^{761}$我們真的深深相信.
我們真的可以在教會裡面.
給予人第二次恩典.
並且深深地承認.
其實我們每一個人都很需要對方的第二次恩典.
求主幫助我們.
我們一起禱告.
是被人愛的天賦.
是你的愛挽回我們.
確實是你的愛先拯救我們.
你知道我們能夠在你面前站立.
我們能夠得稱為義.
不是因著我們真的有何等的義.
而是當我們披戴主人的蘇基督在我們身上.
當聖靈進入我們心裡.
當我們成為一個新造的人的時候.
我們在你面前能夠無所懼怕.
得稱為義.
主也求主你幫助教會.
幫助我們.
能夠細細地聽.
慢慢地說.
實在是在你裡面.
人的憤怒不能夠成就神的義.
求主你幫助我們.
在教會裡重新建立我們.
復興我們.
醫治我們.
改變我們.
你知道我們懂得.
以你的真理.
聆聽你的真理.
遵行你的真理.
繼續前行.
讓我們能夠靠著你的恩.
再一次以你的愛.
去愛我們身邊.
我們覺得不可愛的人.
願主你幫助我們.
去深深地承認.

$^{801}$我們願意給予對方第二次恩典.
而我們也是何等需要.
對方所給予的第二次恩典.
願主你幫助教會.
幫助我們.
阿們.
\newpage



\section{馬太福音 12:22-28}
\label{sec:031640}
\textbf{神國臨在}~~~~黃健羒牧師
\newline
\newline
連結: \href{http://www.church.com.hk/acms/content.asp?site=cdc\&op=show\&type=product\&code=031640\&layout=sermon}{\texttt{031640}} ~~~~ 語音日期: 2026-02-08
\newline
\newline
\hyperref[sec:031639]{\small{< < < PREV SERMON < < <}}
~
\hyperref[sec:index_chronic]{\small{[返順時目]}}
~
\hyperref[sec:031641]{\small{> > > NEXT SERMON > > >}}
\newline
\newline
馬太福音 12:22-28
\newline
\begin{longtable}{cl}
\hline
\hline
章節 & 經文 (和合本修訂版)\\
\hline
12:22 & \begin{tabularx}{0.7\textwidth}{X} 當時，有人把一個被鬼附，又盲又啞的人帶到耶穌那裡，耶穌醫治他，那啞巴就能說話，又能看見。 \end{tabularx} \\ \\ \relax
12:23 & \begin{tabularx}{0.7\textwidth}{X} 眾人都驚奇，說：「這不是大衛之子嗎？」 \end{tabularx} \\ \\ \relax
12:24 & \begin{tabularx}{0.7\textwidth}{X} 但法利賽人聽見，就說：「這個人趕鬼，無非是靠著鬼王別西卜罷了。」 \end{tabularx} \\ \\ \relax
12:25 & \begin{tabularx}{0.7\textwidth}{X} 耶穌知道他們的心思，就對他們說：「一國自相紛爭，必定荒蕪；一城一家自相紛爭，必立不住。 \end{tabularx} \\ \\ \relax
12:26 & \begin{tabularx}{0.7\textwidth}{X} 若撒但趕出撒但，就是自相紛爭，他的國怎能立得住呢？ \end{tabularx} \\ \\ \relax
12:27 & \begin{tabularx}{0.7\textwidth}{X} 我若靠著別西卜趕鬼，你們的子弟趕鬼又靠著誰呢？這樣，他們要作你們的判官。 \end{tabularx} \\ \\ \relax
12:28 & \begin{tabularx}{0.7\textwidth}{X} 我若靠著神的靈趕鬼，那麼，神的國就已臨到你們了。 \end{tabularx} \\ \\
[1ex]
\hline
\hline
\end{longtable}
$^{1}$先解釋一下這個講題.
雖然馬太福音比較多用天國這個字.
但同樣或同一段經文在其他福音書.
馬可或路加會用神的國.
我認為天國就是神的國.
可以交替使用.
因為馬太福音是寫給猶太人看的.
有時他會用神的國.
但很多時候都是用天國.
是為了尊敬.
把神這個字用天代替.
有點像中文的避諱.
如果皇帝的名字有這個字.
如果你想改名就不要用這個字.
有時用同音但用不同的字.
或者有些長輩用這個字.
我們就避一避.
有少少這個意思.
所以馬太福音很多時候是用天國.
其實天國就是神的國.
還有大概十天就是過農曆年.
馬年.
新年期間由小到大父母要教導.
要講吉利的話.
有些吉利話跟我們信仰沒有衝突.
我覺得無所謂可以用.
譬如龍馬精神.
笑口常開.
跟信仰沒有衝突無所謂可以講.
不過有些就要避免.
例如橫財就守.
橫財不是很正當.
那些就不要講了.
或者是鴻運當頭.
因為我們不相信運氣.
一切都在神的手中.
沒有偶然發生的運氣.
所以那些我們可以少些或者避免講.
跟信仰就沒有問題.
新年期間有些顧忌.

$^{41}$譬如新年的時候不要講「死」字.
不要講「輸」.
有些人賭錢輸就不要了.
甚至不要講「空」.
或者是你要租一間屋.
不叫「空屋出租」.
叫「吉屋出租」.
香港才有這些東西.
或者「吉樓交吉」.
不是交「空」.
其實怕什麼?.
我們說避開這些不要講.
新年期間不要講這些.
是怕什麼呢?.
很多人其實心裡是怕得罪別人.
這個我個人是會的.
或者是連累別人.
或者是別人真的遇到一些不順利的事.
就怪你.
就怕這些.
而有可能其實是怕鬼.
怕一些靈界的事情.
因為一些民間的習俗等等.
我希望今天講完之後.
心裡面比較安心一點.
不要怕.
過年期間不要害怕.
當然講話我們小心一點.
不要得罪別人.
或者不要令別人心不安樂.
但我們自己起碼心裡面.
是沒有一些這樣的顧忌.
從經文裡面我們看到有幾方面.
第一就是耶穌醫治這個被鬼附的人.
看見他受攻擊.
就引起很多毛病.
耶穌當年除了傳講神的道.
勸人認罪悔改.
歸向真神.
也同時醫病趕鬼.

$^{81}$剛才我們所讀的經文.
就是記載他和門徒在加里尼那一帶的工作.
首先我們要想一想我們怎樣看待一個人.
我們這些在二十一世紀的人.
通常我們看一個人是一個組合.
主要就是有生理,有心理和有靈性那方面.
英文有時是生理和心理.
就是Psychosomatic Being.
就是生理和心理是一個結合.
我們還有第三方面就是靈性那方面.
以前我們在醫院實習過做醫院的牧師.
我們就說對病人的治療.
除了生理和心理還有靈性那方面.
一起結合對一個病人很快就好一點.
生理的毛病我們當然知道很容易.
就是腎臟,呼吸,癌症.
有問題我們就找醫生.
甚至找專科的醫生做一些檢查.
心理方面有沒有毛病.
有些是抑鬱或者是精神分裂.
我們找心理醫生.
如果靈性有毛病.
或者連續兩三個月.
每晚都發惡夢.
可能你們未知道或者未聽過有這樣的報告.
連續兩三個月每晚都發惡夢.
肯定是睡得不好,白天就沒有精神.
不是生理的病.
是靈性方面接觸到一些邪靈的攻擊.
那就找牧師,或者找一些專門的.
牧師都有很多種.
我就沒有這種趕鬼的恩賜.
有些牧師有,你就找他談談.
其實人是一個整體.
而且這三方面是相互影響.
譬如你身體不舒服的時候.
真的不想祈禱,沒有精神.
或者靈性有毛病的時候.
這樣也會影響你的心理.
有些人是這兩方面.

$^{121}$這幾方面是相互影響的.
身體不好的時候會影響你的心理.
也會影響你的靈性.
有這樣的影響.
所以剛才我們所讀的經文.
這個人看不見,是盲眼的.
也不懂得說話,是因為被鬼附.
絕對有可能,是一個整體.
被邪靈或者魔鬼攻擊的時候.
身體就會生病.
不是所有的病都是因為邪靈攻擊.
有些病是被病毒攻擊.
譬如COVID病毒攻擊.
或者有些細菌攻擊.
總之被攻擊的時候.
我們身體就會有毛病.
所以被邪靈攻擊的時候.
就一定會出毛病.
有些觀眾可能會問.
究竟邪靈是什麼?.
這個圖我在一些團體聚會也分享過.
大概我們看這個世界是怎樣的.
剛才大家唱詩也唱到.
神是最大的.
祂是創造一切.
祂是高高在上.
無論發生什麼事.
神依然是最高的.
上面這裡我們叫靈界.
下面就是物界.
物界當然包括動物,植物,物質.
人類唯一神的創造.
就是在靈界和物界之間.
兩樣都會牽涉到.
所以我們說雞,狗沒有靈魂.
雖然有時說All dogs go to heaven.
狗的靈魂去天堂.
但我們不相信狗有靈魂.
是動物.
所以吃牛肉,豬肉無所謂.

$^{161}$但人就不同了.
人是有靈魂的.
神創造這些東西.
是不是神創造好的天使.
創造壞的天使呢?.
我們說撒旦其實是一個墮落的天使掌.
其實神創造天使的時候.
所有都是好的.
人都一樣.
神不是說做好人,做壞人.
所有人都是好的.
神創造的時候是好的.
但在靈界上的天使.
我們不知道有多少.
但起碼有米加勒,戈培奧.
有幾個天使掌的名字我們知道.
但可能會是少數.
我們不太知道.
其中一個就是墮落的天使掌.
我們就叫做撒旦.
名字叫撒旦.
或者在這段經文裡叫做別西卜.
其實都是指同一個.
這個天使掌想說權力超過神.
他想做神.
但神就限制了他的權力.
所以有一些其他的天使.
本來是好的.
後來就跟著他離開神.
那些我們叫做邪靈.
所以有很多的天使在這裡.
很多的邪靈.
天使不是由人變的.
有些人以為是好人.
將來變成天使.
不是.
或者你做壞事的時候.
將來變成邪靈.
不是.
天使云天使.

$^{201}$人歸人.
所以聖經特別是歌羅西書裡面.
一切所看見的,不看得見的.
都是神所創造.
所以不看得見的就是靈界上的存在.
天使.
這裡有一個比較好的了解.
我們就知道這個世界是這樣.
至於離開了世界的人是怎樣.
那些靈魂是怎樣.
這個就打做問號.
我們不太清楚.
聖經有提到一點.
但就不是很全面.
只不過我們今天相信人離開世界之後.
就不會影響仍然在世界上的人.
所以離開了世界的祖宗.
他們的靈.
不會說是保佑後人.
或者降禍後人.
即是影響人的是邪靈.
可能他們假扮祖先.
或者假扮某一個人.
那麼現身就影響到人.
所以今天我們對於死了的人.
過世的人的靈魂.
我們就不相信他們對我們有什麼特別的影響.
紀念他們是應該的.
記得他們無所謂.
但就不是以為他們有什麼活動影響我們.
這裡就是鬼和.
我們叫鬼其實就是邪靈.
所以在魔鬼的角度裡面.
我們統稱就是撒旦.
有時也包括所有的邪靈.
現今的所謂正邪鬥爭.
其實是好的天使和壞的天使.
或者邪靈之間的鬥爭.
並不是魔鬼和神的鬥爭.
因為神高高在上.

$^{241}$魔鬼一定會輸的.
即是不能打.
所以是一個天使和邪靈之間的攻擊.
我們知道魔鬼如果不能打的時候.
他怎樣傷害神呢?.
就是以神最愛的對象.
傷害他的對象.
神在所有創造裡面.
最愛的是什麼?.
人.
就好像兩夫妻吵架.
明知道打不過你.
就扔爛你的東西.
你最愛的紫砂茶壺就扔那個.
以前有一對朋友.
兩夫妻吵架.
婚姻不好.
有時太太見到先生拍的照片.
有美女在那裡.
兩人吵架.
到最後他們住在二樓.
開窗看下面沒有人.
把相機扔了.
幾百元一個.
鏡頭也幾百元.
扔了.
撒旦是沒有能力和神作鬥.
所以他就要傷害人.
所以撒旦所有的工作.
邪靈所有的工作都是要攻擊人.
因為人是神最愛的對象.
藉著傷害人來傷害神.
這是魔鬼最終的目的.
撒旦和他權下掌控的那些邪靈.
附在一個人身上.
可以嗎?可以.
可以附在一個人身上.
所以剛才我們看見聖經上說.
邪靈附在人身上.
或者附在人身上做壞事.

$^{281}$所以今天我們看到.
親生子女把父母都殺了.
或者在戰爭當中殘忍,折磨人.
那些都是邪靈在壞人身上.
藉著他們做這些事.
邪靈是有能力的.
新約聖經也說.
有些假先知出來.
他也懂得行騎士.
神蹟騎士.
不只是神或者天使可以做的.
邪靈也可以做的.
但為什麼神給他這樣的權力.
我們就不太清楚了.
不過這是暫時的.
所以所有的占卜,算命,星座.
或者香港時的問米,碟仙.
我們年輕時去玩碟仙.
以為沒有害.
不是,那些背後都有邪靈.
邪靈的工作就在那些地方.
神是厭惡這些事的.
所以我們要離開.
遠離,不要接觸這些.
不要算命.
如果我們去看廟.
參觀一下,拍張照片也無所謂.
不過不要裝香.
不要給香油錢.
那類的事不要接觸.
所以神是和邪靈的事.
這兩個世界是避開的.
為什麼邪靈會這樣呢?.
就是要引起一種害怕.
一種恐懼.
如果你怕了我的時候.
叫你捐錢,你就捐錢.
叫你做壞事,你就做壞事.
因為你怕了我.
所以在這樣的情況下.

$^{321}$我們就知道整個靈界會發生這樣的事.
但我們不怕.
因為以弗所書第六章.
剛才Vincent也稍為提到.
我們的征戰並不是對抗有血有肉的人.
而是對抗那些執政的,掌權的.
管轄這個幽暗世界的.
以及天空屬靈的惡魔.
我們的對象是邪靈或魔鬼的世界.
對抗魔鬼.
但是不用怕.
為什麼?.
因為以弗所書第六章也說.
神給我們有這些武器.
真理,公義,信心.
當然聖經的時代是羅馬的時代.
我們現在要現代化.
你拿著聖靈所賜的自動步槍.
可以打得到的.
打完之後怎麼樣?.
壯烈犧牲嗎?.
不是.
聖經說我們贏了之後還可以站立得住.
還可以站著.
這是以弗所書給我們的一個應許.
你拿著神的武器去打屬靈的戰爭的時候.
你會打得贏的.
你不用怕的.
還可以抵抗仇敵完成之後.
還能夠站立得住.
所以凡是有一些攻擊來到我們的身上.
不用怕.
藉著神賜給我們的武器.
有些是攻擊的,有些是防衛的.
可以打得贏的.
第二方面就是.
由主的醫治之下.
我們可以恢復完整.
經文說當耶穌.
將這個附在人身的邪靈趕走之後.

$^{361}$病人就懂得說話了.
又可以看得見.
耶穌醫治了他之後.
他就恢復一個完整的生命.
能夠好像正常人一樣生活.
在那個時代.
為什麼耶穌的功績是包括醫病趕鬼呢?.
是因為一般人將疾病看為神的一種刑罰.
譬如《宇航福音》第九章也記載.
耶穌遇見一個一出生就盲眼的人.
門徒就問了一個問題.
這個人一出生就看不見是他自己犯罪.
還是他的父母犯罪牽連他呢?.
因為罪引起他看不見盲眼呢?.
門徒問這個問題.
但是耶穌回答.
不是他犯罪.
也不是他父母犯罪.
而是要在他身上顯出神的作為.
於是拜他自豪.
所以耶穌的說話已經否定了.
人有病的時候是由於犯了罪.
犯罪當然可能會引起一些毛病.
例如吸煙吸得多.
可能肺部有問題.
或者賭錢輸了有問題.
當然我們今天看一些比較廣泛的.
神喜悅的我們都叫做罪.
犯了罪可能會引起一些毛病.
但是有病不一定是因為犯罪.
這是耶穌給我們的一個解釋.
無論如何耶穌醫治了一個人之後.
就叫人能夠有一個完整的生命.
但是問題是什麼叫做完整的生命?.
很多人在新年的期間.
有時都說五福臨門.
五福是什麼呢?.
就是長壽,富貴,康寧,好德,善終.
新年可能不是說善終.
總之五福就是.

$^{401}$通常我們說有子女有孫.
四代五代同堂很有福氣.
又有錢又有健康可以去花錢.
這樣多好啊.
五福臨門.
通常我們說這個就是幸福的人生.
這樣多好啊.
聖經裡面有一個人都是很幸福的.
這個人是完全正直,敬畏神,遠離惡事.
好人一個.
他生了七個兒子,三個女兒.
有兒子有女兒這個好字.
好幾個好字.
他的家產有七千個羊.
有三千個駱駝.
有五百對牛.
還有五百匹母雷.
不知道為什麼聖經只記載母雷.
應該和他供的.
不計算的.
還有很多僕人婢女.
有錢有財產.
又有子女.
多好啊.
有僕人.
當然大家一定知道他的名字.
就是駱駝.
他以為行得正站得正.
樂善好施敬畏神.
神就賜福給他.
他就有幸福.
後來怎麼樣?.
看過《約伯記》的人知道.
一日之間.
子女全部死了.
同一天.
所有羊,駱駝都被人搶了.
那些僕人被人殺了.
剩下幾個.
很苦啊.

$^{441}$痛苦啊.
再過一會兒.
連身體健康都沒有了.
全身生瘡.
站著也不是.
躺著也不是.
坐著也不是.
都不知道怎麼回事.
全身都是這樣.
受苦.
為什麼《約伯記》的苦難當中.
他要問一個問題.
為什麼好人會受苦?.
我們時常說.
善有善報.
惡有惡報.
做好人一定有好報.
為什麼好人都會受苦呢?.
這是《約伯記》所提出的一個主要問題.
是他要探索的一個問題.
他那三個朋友來探望他的時候.
跟他的想法差不多.
如果你是好人.
做好事.
神就一定賜福給你.
如果你是壞人.
你又做壞事的話.
神就一定會應罰你.
給你有苦難.
所以《約伯記》的三個朋友.
就有這個想法.
你現在有苦難.
當然你犯了罪.
當然你有些不妥.
他們以為神的作為就是這樣.
如果你是做好人.
就一定會有福氣.
如果你是壞人.
就一定會有苦難.
《約伯記》很辛苦.

$^{481}$想來想去都不知道.
哪裡做錯了.
哪裡得罪了神.
為什麼會這樣.
他很希望問神.
甚至說我們質問神.
為什麼你這樣對我.
我對你這麼好.
我又做了這麼多好事.
為什麼你對我不公平.
但是他找不到神.
我相信這是他最痛苦的地方.
神沒有出現.
他連問都沒有得問.
他不明白很多事情.
後來神出現了.
《約伯記》看了很久.
神出現了.
神出現的時候.
並不是回答他這個問題.
還問他一大堆問題.
他一句都回答不出來.
但是他在那個時候.
好像比較能夠經歷那個苦難.
沒有那麼痛苦.
因為神出現了.
起碼出現了.
他回答不了他的問題.
但是他謙卑下來.
承認自己不懂.
這也是約伯第四個朋友的態度.
來到安慰約伯的時候.
說我都不知道.
神的作為我都不懂.
這是謙卑在神面前的態度.
是面對苦難的一個方式.
雖然我們看不到神.
或者經歷不到神.
但是神依然是看顧萬有.
管理萬有.

$^{521}$這是我們基督徒應該有的態度.
雖然我現在感覺不到神同在.
神當然是同在.
但是我們感覺不到.
我們看不見神.
聽不到祂的話.
但是我們相信.
神依然是掌管萬有.
神依然是做王.
是掌管萬有.
這樣就可以一起經歷那個苦難.
所以約伯記所探討的.
我怎樣去渡過苦難呢?.
他指出一個道理.
神同在我們就可以渡過苦難.
人很多時候我們以為.
我們知道那個苦難的原因.
我們就可以解決了.
所以很多人說我病了.
很多人想安慰一下.
當然是問.
你為什麼有這個病?.
為什麼這麼痛?.
為什麼會跌倒?.
知道原因之後.
原來是這樣.
好像就可以面對那個痛苦.
不是的.
約伯記所指出的道理就是.
神同在我們可以經歷.
神同在一起的時候.
我們就是一個完整的生命.
就是一個豐盛的生命.
就是一個幸福的生命.
可惜到今天.
有些基督徒還是好像約伯當初的想法.
好像他那三個朋友的想法.
只要我是樂善好施.
只要我做好人神.
就一定會賜福給我.

$^{561}$其實他們最缺少的是什麼?.
同神的互動.
同神同在同行那一種的經歷.
所以尼加書這節經文大家都懂得讀.
神要我們做什麼?.
就是行公義好憐憫.
全謙卑的心與神同行.
有時候我們讀過就忘記了.
其實重點是謙卑.
謙卑是什麼?.
承認我不懂.
承認我不知道.
我不能夠完全認識神.
不能夠完全知道.
神的旨意就是這樣.
不是這樣.
我們謙卑.
而且要與神同行.
日日與神同行.
有一首聖詩.
有些年長的好像我這樣.
會唱過「他自己 Himself」.
歌詞第一字是.
「前所要是祝福.
今所要是主.
前所憑是感覺.
今靠主言語.
前所要是恩賜.
今要賜恩主.
前我尋求醫治.
今要主自己」.
副歌是「永遠高舉耶穌.
讚美主不厭.
一切在耶穌裡.
為他示我一切」.
有些人懂得唱歌.
我不像長老一樣.
懂得唱歌.
這個「Himself」.
是宣道會的創辦人.

$^{601}$A.B. Simpson.
Simpson他在.
十九世紀二十世紀初寫的.
他當時是長老會的牧師.
已經做了二十多年.
但他體會到自己以前的想法.
有少許不同.
其實他當時的體會是.
他真正需要的是主自己.
是神自己.
以前可能很追求神的祝福.
神的恩賜,醫治,能力.
原來忽略了最重要的是神自己.
所以我盼望大家能夠分得清楚.
過新年的時候彼此祝賀.
都希望得到福.
世俗人當然不理會.
我只想得到五福.
六福,七福都可以.
什麼神,什麼神仙都無所謂.
去什麼廟拜都無所謂.
因為我所得到的只是福氣.
誰給我都可以.
我們信耶穌的人不是這樣.
因為我們不只是說那些福氣就好.
而是得到神自己.
得到耶穌自己.
那才是真正與神同行的福氣.
那才是豐盛的生命.
那才是真正完整的生命.
所以耶穌基督的福音.
什麼時候都強調.
我們不單止從最終得到釋放.
並且要天天跟隨祂做祂的門徒.
與祂的互動.
這才是豐盛的生命.
所以你自己有沒有豐盛的生命.
你就要看看自己有沒有與神一起同行.
一起天天有互動.
第三方面就是信靠神.

$^{641}$我們必定會得到保護.
回到今天的經文.
當耶穌把鬼父治好之後.
有些法利賽人說.
他靠著的能力就是靠鬼王.
即是撒旦或者別西卜.
靠著他的能力趕鬼.
即是大鬼趕小鬼.
耶穌說:可能嗎?.
在撒旦的國度裡面.
如果有權力鬥爭的時候.
這個國怎樣能夠存留呢?.
這個道理我們今天都看得很清楚.
無論古今中外.
如果一個政權內部有內鬥.
權力鬥爭的時候.
遲早會倒台.
不會太久.
很多獨裁者的倒台不是真的有人反他.
而是他懷疑有人反他.
那些最親密的戰友.
一起玩大的.
他會不會害我?.
獨裁者一定會慢慢懷疑到這樣.
所以連那些最親的.
後來都整了他.
到最後怎樣?.
個個都信不過.
剩下他自己一個.
遲早會倒台.
所以這些一個國.
或者一個政權內部.
如果彼此不和的時候.
就會得到繼續.
所以耶穌會繼續說.
如果不是靠鬼王的能力來趕鬼.
而是靠聖靈的話.
這個就是神的國.
好像我們剛才所說的.
神掌權萬有.

$^{681}$祂是最高的權威.
祂的權力在哪裡彰顯?.
神的國就在那裡.
如果神的能力在.
C4C這個教會裡彰顯.
神的國就在這個教會裡.
如果神的能力在你的家裡彰顯.
神的國就在你的家裡.
如果神的能力在你的生命裡彰顯.
神的國就在你的生命裡.
神的能力.
是否在你的生命裡?.
如果你信服在聖靈的引領之下.
神的能力一定會在你的生命裡彰顯出來.
問題是你相不相信.
聖靈就在你的裡面呢?.
有人以為我相信.
甚麼都不需要做.
只是相信,很簡單.
神就自動保護我.
所以信耶穌好.
要洗禮也好.
信洗禮受洗.
如果你只是認為.
口裡這樣說.
神就保護你.
你甚麼都不做的時候.
我怕你真的遇到撒旦攻擊的時候.
你自己都懷疑.
我是否真的相信了?.
神是否真的會保護我?.
有些信心比較軟弱的人.
在這些時候.
就求神給他一個證據.
或者給他一個神跡.
給他看到甚麼.
其實這種信心不是真正的信心.
因為你仍然是要摸得到.
看得到才相信.
這個信心我們應該加強一點.

$^{721}$我怎樣知道聖靈是住在我裡面呢?.
比較妥當一點.
我覺得信心應該建立在神的話語上.
有時信心軟弱的時候.
你揭聖經或者問人.
聖經在哪裡說?.
有神的話語在當中.
你就會知道.
聖靈是住在你裡面.
我介紹了三類的經文.
第一類是講到神的應許.
在約翰福音第十四章.
約翰一書第四章.
和羅馬書第八章都有提到.
我要求父.
耶穌說的.
父就賜給你們聖靈.
使祂永遠與你們同在.
聖靈就是真理的靈.
要引導你們進入真理.
羅馬書也提到.
保羅說.
神差派的聖靈住在我們裡面.
所以神的應許.
聖靈在你裡面.
第二類的經文就是我要自我反省.
自我省測.
約翰一書第四章裡面提到.
神就是愛.
住在愛裡面的.
就住在神的裡面.
神也在你裡面.
所以如果我們有一個愛人愛神的心.
就證明神住在我們裡面.
藉著聖靈住在我們裡面.
如果你發現你的愛心不夠.
你重溫一下.
耶穌怎樣饒恕你.
怎樣接納你.
愛你.

$^{761}$你就會愛他.
你就知道.
聖靈就在你裡面.
第三就是他人的印證.
加泰書第五章二疑節.
聖靈的果子.
很多人會背.
一起背.
忍耐喜樂和平.
忍耐恩慈良善.
信實溫柔節制.
九方面.
聖靈的果子.
我說是他人的印證.
不是說你自己說.
我有信心.
人家說你怎樣.
以前我們教會.
時不時都會有些.
在中國國內來探親的.
他們的子女在美國工作.
探親有時會幫他們看孫子.
子女在這裡信主.
他們來探親的.
還是在黨裡面.
這些幹部 高幹.
當然是不信.
無神論.
後來信了.
相信耶穌.
為什麼?.
因為看到他的兒子.
看到他的媳婦.
很大的改變.
有聖靈的果子在裡面.
以前兒子和媳婦經常吵架.
現在好了.
吵少了.
或者沒有怎樣吵架.
特別是.

$^{801}$記得有一次做婆婆.
受洗禮時的見證.
稱讚媳婦.
不容易的.
教婆稱讚媳婦.
是因為她.
我的想法改變了.
不是說她懂得很多.
而是看見她兒子媳婦的改變.
這是人家的見證.
如果你在身上有聖靈的果子.
聖靈就在你裡面.
如果我們確定.
是神藉著聖靈住在我們裡面.
我們就有神的能力.
能夠抗拒魔鬼的攻擊.
抗拒邪靈的攻擊.
所以當魔鬼攻擊我們時.
我們就不怕.
好像雅各書所說.
我們拒絕他.
奉耶穌的名.
請他離開.
魔鬼就離開我們.
遠離我們.
就不怕.
附在我們的身上.
信靠神就有平安.
信靠神有勇氣.
不會害怕邪靈.
這是神能力的彰顯.
在你的生命裡面.
所以耶穌當年.
祂除了傳道.
也是醫治.
趕鬼醫病.
祂的門徒.
也是繼續做這樣的工作.
他們的能力.
不是來自魔鬼.

$^{841}$而是來自聖靈.
神藉著聖靈.
耶穌的靈.
能力彰顯.
就表明神國的臨在.
整個國就出現了.
清朝有些兵的衣服.
寫著「勇」字.
有些前有些後.
其實「勇」是所謂「勇刑」的兵.
「勇刑」其實是文兵.
不是正規的.
正規的叫「兵」.
「勇刑」叫「勇」.
那些兵就是八旗.
清朝的八旗後代.
才做得正規軍.
文兵有個「勇」字.
如果你相信耶穌.
確定聖靈在你裡面.
就有膽量.
你不會怕.
「勇」.
不是說魔鬼放過你.
或者不攻擊你.
是會的.
不過你有足夠的抵抗能力.
神給你有足夠的武器.
來保衛自己.
如果有這樣的信心.
不單是過年過節.
無論什麼時候.
我們都是百無禁忌.
其實另一件事.
我們什麼都不怕.
不過如果是為了尊敬.
尊重別人.
過年的時候.
多說吉利的話.
無所謂.

$^{881}$都是應該的.
我們就在神的面前.
才能夠有這樣的平安.
請大家低頭.
我們同心祈禱.
天父上帝.
你是高於自高之大的神.
除了你以外.
沒有其他的神.
除了你以外.
沒有真正的主.
我們甚至是.
跟著聖經所記載.
在耶穌基督裡面.
我們的地位甚至超乎天使.
可以跟耶穌一樣.
在天上.
能夠有這樣的信心.
今天我們能夠勝過.
赤靈的攻擊.
勝過撒旦的引誘.
都是靠聖靈的力量.
每一個信主的人心中.
你住下來.
永遠不會離開.
求主繼續幫助我們.
有這樣的信心.
知道聖靈和我們同在.
這個國在我們當中.
臨在.
讓我們都能夠有.
心裡面的平安.
並且有這樣的能力.
到我們的生命.
就是一個完整的生命.
阿們.
\newpage



\section{約書亞記 14:6-12}
\label{sec:031641}
\textbf{活出使命的跟從者——迦勒 I (普)}~~~~梁明財牧師
\newline
\newline
連結: \href{http://www.church.com.hk/acms/content.asp?site=cdc\&op=show\&type=product\&code=031641\&layout=sermon}{\texttt{031641}} ~~~~ 語音日期: 2026-02-08
\newline
\newline
\hyperref[sec:031640]{\small{< < < PREV SERMON < < <}}
~
\hyperref[sec:index_chronic]{\small{[返順時目]}}
~
\hyperref[sec:031642]{\small{> > > NEXT SERMON > > >}}
\newline
\newline
約書亞記 14:6-12
\newline
\begin{longtable}{cl}
\hline
\hline
章節 & 經文 (和合本修訂版)\\
\hline
14:6 & \begin{tabularx}{0.7\textwidth}{X} 猶大人來到吉甲，約書亞那裡，基尼洗族耶孚尼的兒子迦勒對約書亞說：「耶和華在加低斯‧巴尼亞指著我和你對神人摩西所說的話，你都知道。 \end{tabularx} \\ \\ \relax
14:7 & \begin{tabularx}{0.7\textwidth}{X} 耶和華的僕人摩西從加低斯‧巴尼亞差派我窺探這地的時候，我剛四十歲。我把心裡的話向他報告。 \end{tabularx} \\ \\ \relax
14:8 & \begin{tabularx}{0.7\textwidth}{X} 雖然同我上去的眾弟兄使百姓膽戰心驚，我仍然專心跟從耶和華－我的神。 \end{tabularx} \\ \\ \relax
14:9 & \begin{tabularx}{0.7\textwidth}{X} 那日，摩西起誓說：『你腳所踏之地必要歸你和你的子孫永遠為業，因為你專心跟從耶和華－我的神。』 \end{tabularx} \\ \\ \relax
14:10 & \begin{tabularx}{0.7\textwidth}{X} 現在，看哪，耶和華照他所說的使我活了這四十五年。當以色列人在曠野飄流的時候，耶和華曾對摩西說了這話。現在，看哪，我已經八十五歲了。 \end{tabularx} \\ \\ \relax
14:11 & \begin{tabularx}{0.7\textwidth}{X} 現今我還很健壯，像摩西差派我去的那天一樣；無論是戰爭，是出入，我現在的力量和那時的力量一樣。 \end{tabularx} \\ \\ \relax
14:12 & \begin{tabularx}{0.7\textwidth}{X} 請你將耶和華那日所說的這山區給我。那日你也曾聽說，這裡有亞衲族人，以及寬大堅固的城，或許耶和華會照他所說的與我同在，我就把他們趕出去。」 \end{tabularx} \\ \\
[1ex]
\hline
\hline
\end{longtable}
$^{1}$主持人:各位網友好!今天我跟各位可以開始一個新的講題了,大家不用再聽那個十二,十三篇了..
主持人:今天我大概,如果未來我還會在《普通話崇拜》跟各位講到的時候,我會按照教會的一個主題..
主持人:今天這個二零二六年,你知道教會的一個主要的主題是什麼嗎?這邊有一個可以看到的..
主持人:我們這幾年都是以敬畏上主作為一個總題,一個最重要的..
主持人:我們有三方面希望弟兄姊妹能夠一起在當中,我們來依靠神來做到的.這是紮根在神的話裡面..
主持人:去年是鼓勵我們多傳揚福音,今年就是活出我們的使命..
主持人:所以,我希望能夠在崇拜當中跟各位按照教會的重點主題,今年是敬畏上主,活出使命..
主持人:如果你有機會拿到周刊或者程序表的時候,就會很清楚地給我們這個使命是什麼..
主持人:所以,敬畏上主,活出使命,我大概希望能夠從一些聖經人物的身上來幫助我們去活出神給我們的使命..
主持人:教會在我們大禮堂的崇拜裡面,我們主要是用到粵書雅技的經文來幫助我們去看怎麼樣可以活出神給我們的使命..
主持人:所以我也參考一下,在廣東話崇拜裡面他們講得會比較細節一點..
主持人:他們從粵書雅技的第一章開始,我大概會講完粵書雅技..
主持人:我去從當中把一位聖經人物拿出來跟各位來分享,怎麼樣可以幫助我們活出神給我們的使命..
主持人:我就用粵書雅的一個好朋友,你知道是誰嗎?好,都看到了..
主持人:最佳的就是粵書雅的好朋友,我們就說他是死黨的最佳拍檔,是跟粵書雅密切關係的..
主持人:如果你很想多聽一下粵書雅的話,你就可以到廣東話的崇拜裡面聽到粵書雅..
主持人:我跟各位今天來看迦勒這位聖經人物的時候,其實也會講到一點點粵書雅,但主角是迦勒..
主持人:我跟各位來看迦勒他怎麼樣回應神的召命,然後活出神給他的使命.我們來看一段的經文,在《粵書雅記》十四章的六到十二節..
主持人:我還是邀請各位弟兄姊妹們至少來讀一遍.如果你可以的話,你大概可以回家把這段經文多看幾遍.我今天大概只會講到這段經文的大概一部分,我三月份的講到會把其餘下去,所以今天講不完的,放心,你不要以為我不會再講..
主持人:三月份我會講第一節和下一節,在三月份的時候會講.我們先來讀一遍,讓我們多認識一下迦勒到底是一個怎麼樣的人.看到這段經文的話,我們一起來讀一遍..
主持人:那時尤大人來到吉甲見約書亞,有基里洗足耶夫尼的兒子迦勒對約書亞說:「耶和華在迦底斯巴黎雅,指著我與你對神人摩西所說的話,你都知道了..
耶和華的僕人摩西從迦底斯巴黎雅打發我揮炭這地,那時我正四十歲,我按著心意回報他,然而同我上去的眾弟兄使百姓的心騷化,但我專心跟從耶和華我的神.」.
巴爾魯摩西啟示說:「你早說他之地定要歸你和你的子孫永遠為也,因為你專心跟從耶和華我的神.」.
自從耶和華對摩西說這話的時候,耶和華照他所應許的使我全活著四十五年,期間以色列人在曠野行走,看啊,我已經我八十五歲了,還細長壯,像摩西打發我去的那天一樣,無論是征戰出入,我的力量那時如何,現在還是如何..
求你將耶和華那日應許我的這三弟給我,讓你有雅納族人,並歡大兼顧的臣,也臣聽見了,或者耶和華照他所應許的與我同在,我就把他們趕出去.」.
主持人:通常現在我們看電影也好,看什麼書本故事也好,我們通常這個故事做完以後,我們通常會再拍一套電影,我們把它叫做什麼電影的前傳,這個當然是賺錢,因為那套電影可能是很紅,很賣座,所以他就再拍一個前傳,讓我們知道到底之前發生了什麼事情..
這個《月書雅記》我大概有一個前傳,它主要是在《明述記》裡面,讓我們多明白或者說多知道那個背景是怎麼樣的..
所以我跟各位看一下《月書雅記》的一個前傳,在《明述記》有幾段經文.這幾段經文的最重要是什麼呢?.
剛才我們在《月書雅記》十四章裡面,我們講到聖經怎麼樣形容迦勒..
這裡大概就是它是怎麼樣形容迦勒的,因為他是一個專心跟從耶和華的人.但在他的前傳裡面,其實這個已經在自己的模式也好,耶和華說也好,他就這樣子讓我們去認識,明白或者說知道迦勒到底是一個怎麼樣的人..
這個前傳在《明述記》裡面說,耶和華小玉摩西說,你打犯人去窺探我所賜給以色列人的迦南地..
這個大概是在當中一個,我不曉得弟兄姊妹們對以色列人的歷史有多熟悉..
這個其實記載一個很特別的故事,告訴我們他們還沒有到迦南地理之前,他們離開了埃及,然後他們先派十二個探子去看一下神所給以色列人的應許地到底是怎麼樣的..
所以他們就檢選了十二個,都是他們很有代表性的,是他們做首領的.如果你看到這一段經文,就打發他們去探子,去看看那個地方是怎麼樣的..
然後他們十二個探子代表十二個支派出去,當中就有迦勒,也有月舒雅.這個講到任的兒子和希雅,最後那個其實就是月舒雅,當時的翻譯有點不一樣而已..
然後十二個探子回來以後,你們知道發生什麼事嗎?當中有十個探子說那個地方,其實他們都說那個地方很好,那個物產很豐富..
就像神應許給他們一個,我們通常叫應許地是叫做流乃與密之地,就是很多的豐富的產物,就是你進到裡面去,你就可以不用擔心,憂愁你的生活是怎麼樣..
當他們十個人都很認同,但有十個跟兩個對比,有十個人說他們不能夠勝過那些在當地的居民,他們說那些人是很高大,很厲害的人,就是他們是很會打仗,很會打架的人..
但這十個人就說我們不要去了,另外兩個就是月舒雅跟扎勒,他們說什麼?他們說可以去,這是神應許給他們的地,他們就可以去了..
所以,當他們這樣子十對二,你說誰會贏?要投票的時候,我們都知道,現代人如果我們喜歡要用民主方式去投票的話,十對二誰會贏?如果信你的話,你會跟從十個人還是兩個人的意見?.

$^{41}$我們可能大部分人都跟從十個人的意見.但這個就是他們對神那種不相信,他們也對神應許給他們那個保證,他們沒有相信神可以帶領他們,幫助他們進到那個應許地..
所以,當這十個人說不要去的時候,神就跟摩西來說話了,他就在明述這十四章裡面,他說這些人他們沒有接受神給他們的應許,神就要發誓啟示,應許他們去那個地方..
只有兩個說可以去的,扎勒跟約書仰可以去,其他人就會在以後他們要在四十年在曠野的時間,他們最後都不能進到那個應許地..
所以這裡也形容扎勒跟約書仰他們都可以進去那個應許地,他們就是一個為二兩格的人,所以他們是很密切的關係..
其實在明述這三十二章裡面,他們講到,再一次的形容是,扎勒跟約書仰都是專心跟從耶和華神的人,這就是兩個人,讓我們去多認識在他們的背景裡面,他們是怎麼樣的一個狀況..
在前傳裡面,或者說跟約書亞記所記載的,我把它有點形容是這兩個人.但其實這兩個人都是一樣.約書亞大概是代表他是領受了神的召命..
然後他在約書亞記第五章裡面,我們看到耶和華神出現跟約書亞的對話.當神出現的時候,這個人在約書亞預備去攻打耶利哥城的時候,他在跟約書亞的對話..
約書亞看到這個人的出現,他這邊形容他是一個人,但其實我們都相信他是耶和華神.所以,約書亞問他,你是幫助我的,還是幫助我的敵人呢?他回答說,不是的.我來是要做什麼?耶和華軍隊的元帥..
所以,約書亞就伏伏殺敗他.就是說,我主有什麼話,吩咐仆人.耶和華軍隊的元帥吩咐約書亞說,把你的腳上的鞋脫下來,因為你所站的地方是聖的.約書亞就照著行..
所以這裡,我想跟各位看,約書亞是很直接的,以前是神透過摩斯來吩咐他們以色列人要做什麼.這一次,讓我們看到,約書亞在神面前,他就領受神給他的一個,以前從摩西給他的一個召命..
但這個神要來的,不是成為約書亞的助力,不是攔阻他,但也不是幫助他.這個人來到,不是必須要幫助那個約書亞,幫助以色列人進迦南帝,他是怎麼樣?他是主力..
意思是怎麼樣?他是要領導,要帶領.所以,約書亞只是能夠聽從這位主耶和華神的話,這就是讓我們看到,我們的神是不單止我們要敬畏祂,敬畏祂是因為祂是神聖的上主,萬有都是屬於祂的..
所以,祂來在以色列人當中,在約書亞的領導裡面,祂才是真正的領導,祂是帶領以色列人成為耶和華軍隊的元帥.所以,約書亞可以代表著,祂是親自領受神給祂的召命..
我想跟各位用另一個角度看,扎勒就是一個很會回應召命的人.當然,你可以說,扎勒和約書亞都領受了神的召命,但他們也是去回應神的召命..
但我相信,在扎勒的身上,從今天我們要看的經文裡面,他是去把回應召命從他的一生人當中來實踐出來,所以他就能夠成為一個活出神給他召命使命的人..
其實,在整個約書亞和扎勒他們在帶領以色列人離開埃及進到迦南地的時候,他們都是一同領受神給他們的召命,他們也一同去回應,但他們有不同的角色,他們有不同的工作要負責..
約書亞是做帶領的,他很明顯就是神透過摩西給他的一個很直接的跟他說,他就是摩西的接班人,摩西的棒子就交給約書亞..
然後,扎勒不是那個接班人,扎勒也要順服跟從約書亞的帶領..
但扎勒有他的角色,就是他是跑到前線去打仗的那一個,甚至說他是打一場最難打的仗..
如果我們看那個約書亞第十四章,剛才所讀的經文後半部,他是講到他要去打一個在迦南地叫雅納族的人,那個雅納族的人是最高大,最威猛的,他在當中是打最難打的仗..
雖然他不是領導,但他是在回應神的召命,活出神給他的使命的時候,他反而是能夠更實踐這個召命去承擔的最難打的一場仗..
所以我就把這個扎勒跟約書亞讓各位弟兄姊妹可以多認識一下,他們都領受神的召命,都回應神的召命.但回應神的召命當中,我相信扎勒就是一個活出神使命的一個跟從者..
就像聖經裡面形容他是一個專心跟從耶和華的人,他的跟從就是跟從神去打一場最難打的仗.他就是把這個回應召命,雖然很艱難,但他能夠活出神給他的使命..
所以我想跟各位來看,這個活出使命的跟從者扎勒的時候,聖經怎麼樣形容他?聖經形容他是一個專心跟從耶和華的人.這在約書亞跟那個明述記,以及在申命記當中都有這樣子形容扎勒..
另外在聖經裡面,這個形容詞除了形容扎勒,而且大概也形容了約書亞之外,他另外一個形容的聖經人物就是大衛.他對大衛也是這樣子的形容,他是一個專心順服神的人..
待會我再跟各位來看這個專心到底是什麼意思.這個專心,我讓各位來看另外一個聖經的譯本,這個叫《呂正宗譯本》的聖經.他把那個專心跟從耶和華翻譯作"滿心滿懷跟從永恆主"..
他把那個專心跟從耶和華翻譯作"滿心滿懷跟從永恆主".其實在原來的原文聖經裡面所講的,他沒有一個"心"這個字,他只是用一個"滿",英文叫"full",所以他用另外一個英文的翻譯叫"fill up after Jehovah",就是說他是一個跟從耶和華的滿..
如果我要讓各位來多認識就是說,那個重點就是在一個"滿",這個"fullness",或者說"fill up after Jehovah",就是一個人他充滿了什麼在他的生命裡面.在他這個人裡面,他到底充滿了什麼東西在他的生命裡面?.
所以這個是神對迦勒的一個肯定.我想先跟各位來看,在活出使命當中,迦勒他怎麼樣活出神的使命呢?這是神對迦勒的一個肯定..
所以從迦勒神對他的一個他自己在神面前所領受,也是神姐姐摩西所講的.這個是迦勒他自己講他是專心跟從耶和華的神..
然後在剛才我們所讀的經文裡面也講到,耶和華神是讓摩西跟迦勒他能夠明白認識迦勒是一個怎麼樣的人.這個滿心跟從,或者說神對迦勒的肯定是什麼呢?.
這裡講到耶和華教他所應許的,使我存活這四十五年.然後他說我的力量那時如何,現在還是如何.這裡最重要,我想各位來看,是耶和華神教他所應許的,使我存活這四十五年..
能夠充滿在迦勒這個人生命裡面的,不是他專心,也不是他真的是很有能力,很有力量.其實這個充滿在迦勒的生命裡面,是神的應許..
無論在中文聖經的翻譯裡面去說他專心跟從耶和華神,或者說他是一個怎麼樣,有一個對神那種信靠,這個都是他的滿..
他的滿是什麼呢?他的充滿是什麼呢?是因為他有神的應許在他的生命當中,放在他的生命當中.或者說他緊握著的是什麼呢?他緊握著的不是他擁有多少的軍兵,軍馬,他有多少的能力..
我相信在剛才提到那個渣男帝的人,他們都是比以色列人還高大的.我相信迦勒也不是一個很強壯的人,也不是很高大,跟那些渣男人可以差不多..
但我相信迦勒能夠知道他自己擁有的,充滿的,就是神給他的應許.他擁有的就是從神給他的那一份的能力..
所以這些在這段經文裡面提到,迦勒的見證是他很明白神對他的肯定.迦勒他分享,從四十歲到八十五歲的時候,他對神視中如一,他沒有離開神,他沒有改變..
從他被差派到迦勒去做探子的時候,他已經相信充滿擁有神的應許,他很確信神的應許一定會實現..
所以在四十歲到八十五歲人生的每一個不同的階段裡面,其實神都知道.神是透過經文來讓我們看到,在迦勒的身上,無論一個人的人生,從他出生開始到走完人生的路,.

$^{81}$我們相信無論任何的年紀,我們人生任何的階段,神都知道我們,認識我們.他認識我們這個人是一個怎麼樣的人,就像他對迦勒的認識一樣..
所以從這段經文,雖然是迦勒他跟約書亞的分享,其實也是分享到神對迦勒的肯定.迦勒很明白他不是做以色列人的領導,他不是做以色列人的元帥,他不會跟約書亞來比較..
有時候我們會覺得,在一個群體或者說在你的公司裡面,你的角色可能是,為什麼有時候你會覺得不是我當領導,還是別人當領導呢?.
好像你跟人比較的時候,你就發現好像為什麼不是我呢?但這裡讓我們看到,迦勒他對自己的肯定,相信神對他的肯定是,他不需要跟約書亞比較,他就是那個要跑到前線去打仗的人..
他甚至是要打那個最艱難的仗,最難認負的敵人,他就去上去,要求那個約書亞給那個地給他.他就負責去把那個最強壯的迦勒人,雅納人來打下來..
所以他知道,約書亞他是領導,迦勒明白約書亞是領導,他就是那個角色,就是他要去執行神給他的應許,得到神給他的應許..
所以神對迦勒的肯定,也是迦勒他自己的確認跟確信..
另外,我想跟各位來看迦勒他擁有的滿,他是一個怎麼樣的滿心滿懷,滿意跟滿足了..
這個滿大概的意思,我相信就是跟我們通常用的心滿意足,或者說我們已經很足夠了,那個滿的意思..
但到底是什麼充滿我們這個人生命?什麼是充滿了迦勒這個人的生命呢?他整個人..
剛才我們提到那個專心跟從,耶和華神另外一個形容是在大衛王的身上..
我想跟各位來先看看,大衛王他的滿是什麼呢?.
《列王記》上的第十一節,他用所羅門跟大衛這個比較,所羅門追求的,他的心去追隨的是什麼,跟大衛追隨的,追求的是什麼來比較..
這裡告訴我們說,所羅門在年老的時候,他的妃嫔誘惑他的心去隨從別神,他沒有效法他父親大衛,誠誠實實地順服耶和華他的神..
因為所羅門隨從斯頓人的女神亞西塔魯和阿們可真的神米勒宮,所羅門行耶和華眼中看為惡的事,不效法他父親大衛專心順從耶和華..
大衛專心順從耶和華,就是跟那個扎勒所形容的專心跟從神,耶和華神一樣,在原來希伯來文是一樣的..
所以這裡我們看到,所羅門追求跟大衛的很完全的不同,所羅門追求可能是一種滿足是什麼,跟從他的妃嫔去拜他們的偶像..
但大衛是專心順從,這是他的滿是以順服耶和華神,作為他的最大的滿足,兩個不同的心,兩個不同的跟從,一個是跟從他的妃嫔去拜偶像,一個是願意順服這位耶和華神..
所以這是大衛跟所羅門不同.大衛的那種滿是以他願意充滿對神的那種順服..
主持人:新年的時候,通常我們以前就有說新年的勵志,有什麼在新年的勵志?我大概不會問這個了,我會問大家有什麼新年的願望呢?.
大家有沒有想過你今年新年有什麼願望呢?大家已經沒有什麼願望了,大家覺得願望都不太實際,都不要再想什麼東西了,過年就算了.我們說每年都是年年難過年年過,有新年就好了..
我們中國人在過春節的時候,我們大概都喜歡貼一些春聯,我們香港人叫"飛春",就是貼在家裡面,門口的春聯..
我不曉得您會不會有這樣子的習慣,或者說您都會貼一些,買一些,還是自己寫.我不曉得您今年會貼什麼,今年春節是馬年,有誰會貼什麼馬道功臣,還是香港人以前古老的時候有一些狗馬亨通,那些喜歡獨馬的人..
然後是什麼一馬當先,馬上成功,還是很多這些不同的飛春,春聯的貼出來..
我不曉得您有沒有看過,這個"福"字通常很多人都會有些把它倒過來,有它的意思.有人見過有一個叫"長門"的春聯嗎?年輕的應該沒有看過,也沒有知道這個"長門"是什麼意思..
我們貼在米缸裡面的,我家以前真的是古老的時候有一個很大的米缸,很厲害的大米缸.通常過年前就要把它換一下,就把一個新的"長門"貼在裡面..
因為我們中國人都是吃飯的,米缸常常很擔心沒有米,希望未來的一年裡面的米缸裡面的米長長滿.這就是我們都希望的一個"滿",我不曉得您在您的人生當中,或者說在春節的時候,新的一年開始,您會有什麼樣的願望,會不會這個"長門"在您的未來是一個..
但到底說,如果真的我的願望是一個"滿",是一個"長門"的話,那到底這個"長門"是承載一些什麼樣的"門"呢?.
我不曉得您有沒有想過,如果在您未來一年的生活裡面很滿,但這個"滿"是承載一些什麼東西,充滿了一些什麼東西呢?.
更重要的是,在您的生命裡面,您的生活裡面,您的"滿載"是什麼呢?其實從迦勒的身上,他跟隨耶和華的"滿"是什麼呢?.
其實這個"滿"就是充滿神的應許,跟神跟人之間的約,這是在迦勒的生命當中.如果我們回去看剛才,您還有印象的話,在十四章的第十節裡面,就是講到神的應許..
迦勒能夠來到約書亞面前,跟他要求給他那個地,作為迦勒的子孫要的地,地業的話,其實不是迦勒他自己要求,是神透過某些應許,迦勒要給他的,是神的應許..
所以,迦勒他記得的,不是因為他硬要把一個好的地方拿,就是要約書亞給他,他是以神的應許作為他的"滿",他記得是神應許給他的..
其實這個"滿",我想提醒弟兄姐們,不一定是多.在神的應許裡面,不是說那個"滿"就是給你很多..
我們有時候會想到,在剛才說的那個應許地是很多的流來與蜜的地方,但如果大家有機會去讀一下以色列的地理環境的話,他們會告訴你,以色列的地理環境是不是真的是所有地方都是流來與蜜的地方呢?.
你就會發現其實不是.其實以色列的地方也真的不太好,如果你到過以色列的話.所以,"滿"不是多.當然,神應許給他們的地方真的是最好的,但不是我們想像的很多,很豐富的流來與蜜的地方..
所以,"滿"不是關乎多,跟多不一定要掛勾,相等於多."滿"是告訴我們,我們充滿什麼在我們的生命裡面."加勒"充滿的是神的應許,是神跟人之間的約..
如果你回到剛才我們讀到的第八,第九節裡面,加勒講到,他跟神的關係是講到"我的神",他引用摩西的說話也是說"我的神",是表達他跟神之間那種立約的關係..
在約書亞記的第五章裡面,更提到一個很特別的事情,就是當時候神吩咐約書亞要給以色列人進行第二次的割離..
他們離開土耳其的時候,實行了第一次的割離.當他們進到迦蘭地以後,在四十年後也就沒有進行過割離了.他們第二次的割離就在約書亞記的第五章..

$^{121}$其實,這個割離對以色列人跟神之間的表達是一個立約的關係,是重申他們跟神之間的約的關係是沒有改變的.神仍然是接納這群的以色列人,雖然他們的祖先不好,他們不相信神的應許,所以他們四十年都死在曠野裡面..
這就是神的應許跟神仍積約充滿在扎勒的生命裡面.扎勒那種根從英合法的門,不是因為他很有能力,不是因為他很會打仗,是因為他有神的應許,有神跟他的約的關係在他的裡面..
這就是扎勒的門.大衛他那種順從神的門,也是一個讓我們可以多一點在我們的生命裡面學習的一個機會..
雖然在剛才立王祭上所羅門跟大衛的比照裡面,我們大衛覺得所羅門是一個很差勁的人,他的心根從偶像,根從他的妃嫔,根從那些不好的人,大衛則是一個專心順服,充滿順服對神的順服..
但是我們也知道,在大衛的生平當中,大衛其實也有他軟弱,有他跌倒,有他犯罪,得罪神的時候.所以他有一些時間他在充滿,他的門在不一定是順服,他的門在可能是一些很黑暗,很骯髒的東西..
所以我想跟各位從大衛的身上,也可以學習一些我們要被提醒的功課,到底我們是一個怎樣滿的人.我們從大衛的一個詩篇裡面,我想跟各位來看一下..
詩篇五十一篇,如果大家回去的時候可以再看這個詩篇.我把幾節的經文拿出來,不是整個詩篇.我們一起來讀一遍.你看看,在大衛的生平裡面,他曾經充滿了些什麼東西?.
他後來寫這個詩篇的時候,他也充滿了一些另外什麼東西?他也期望,祈求,將來可以充滿什麼東西?我們來看這個詩篇五十一篇的幾節的經文.看到了我們一起來讀一遍..
"神啊!求你按你的慈愛而廉恤我,按你豐盛的慈悲塗抹我的過犯.求你將我的罪念洗除盡盡,並解除我的罪.神啊!求你為我造清潔的心,使我裡面重新有正直的靈..
你本不喜愛罪惡,若喜愛,我就獻上.凡罪,你也不喜悅.神所要的罪,就是憂傷的靈.神啊!憂傷痛悔的心,你必不輕看.".
大家從這個大衛的詩篇裡面,你可以看到,大衛曾經充滿過什麼東西?在寫這個詩篇之前,大衛充滿什麼?充滿了很多的過犯,充滿了很多的罪念..
他的門是什麼?是罪的門.這是大衛在寫這個詩篇之前,他的反省,他的認罪,他知道自己的錯,他改過..
然後他,這是門,他懇求神把他除去,把他拿掉.然後他寫這個詩篇的時候,他祈求神在那個現在的時候,他求神給他什麼充滿?.
求神的慈悲,憐思的門.當然,他也同時會經歷到憂傷,痛悔的門.這就是大衛的門..
將來,我相信,他去祈求神給他有清潔,正直的門,讓他的生命重新的開始.他去帶著滿滿的去跟從耶和華所得到的門,不再是那些罪的門,也不再是憂傷,痛悔..
他很想在他的未來,他去跟從耶和華神的時候,是有清潔,正直的門,是一種真正的滿足,真正的放心,真正的足夠,是跟他在神面前的重新建立那種關係..
是因為他有神的應許,神按著他的慈悲,按著他的慈愛,來讓他重新再來.所以,我盼望這也是我們在生命當中,如果我們真的有不好的門在我們的生命裡面,我們求神把他拿走,重新有一個正確的門,好的門在我們的生命裡面..
我想跟各位來看施彬赤三篇,詩人他所經歷到的.這個我在下次講到的時候,最後也會再講..
不過,今天我們講那個福的話,在過春節的時候我們要講福的話,這裡提醒我們什麼是他的福分,對這個詩人來說..
縱然他真的是……可能我們每一個人都會經歷到的了.有神的同在,神拖著我們的手,帶領我們,他用他的訓練引導我們,除了神以外,我們在天上地下都沒有什麼是值得我們去愛慕,去追求的了..
雖然我們的肉體跟我們的心腸會衰殘,可能你還年輕,你沒有感覺衰殘太多,但是當你工作忙碌,壓力太大的時候,你都覺得你的身體和心靈都被摧殘,不是衰殘,是摧殘..
但這裡提醒我們說,神是我們的力量,也是迦勒的力量如何,他的日子也如何.神更是詩人的福分,直到永遠..
所以今天我們的門,我不知道你的長門,你會期望的是什麼.今天我們的福氣,我也不知道你會盼望哪一種福氣.我盼望的是,神就是你跟我的福氣.我們就是能夠過一個生活是"跟從耶和華神的門",不是追求其他的東西..
像迦勒一樣,他有神的應許,他有神跟人之間的約,這個約就是成為迦勒生命當中的充滿,最滿足的,不是他得到任何的東西,不是他擁有很大的,強壯的軍兵,軍馬,不是他拿到神給他的地業,作為他的應許地的時候,他的滿足..
我相信迦勒他給我們的一個看見就是,他是跟從耶和華神的那種滿足,就是他有神的應許.他很確信神跟他之間的約是永遠不會改變的.他就帶著這樣子的應許,守著這個約來進到迦南地,進到神給他的地方..
所以我盼望在新春,在新年開始的時候,你的福氣真的是以神為你的福分,你的場門真的是把神的應許,把神人的約在你的生命裡面,你就充充滿滿地帶著這個福氣去跟從我們的神.我們一起來禱告..
求你幫助我們,讓我們從迦勒的身上,他怎麼樣活出神給他的使命,進到神你應許給他的地方的時候,他就是帶著對神這份應許的滿足.他的生命裡面充滿神給他的這份應許,他也確信神跟他所立的約的時候,他有從神而來的能力,他就去打最難打的仗..
但我相信他就是這樣子,因為跟從耶和華的門,他就得到神給他的應許.願神在我們新的一年開始的時候,也讓我們真的在生命裡面,我們有這上迦勒這份的滿足在我們的生命裡面.願主把你的應許,願主你自己成為我們生命裡面的最大的福氣.謝謝主聽我們的禱告,奉耶穌救的民主,Amen..
\newpage



\section{路加福音 1:66}
\label{sec:031642}
\textbf{聖經人聖經事:有主與他同在}~~~~黃克勤牧師博士
\newline
\newline
連結: \href{http://www.church.com.hk/acms/content.asp?site=cdc\&op=show\&type=product\&code=031642\&layout=sermon}{\texttt{031642}} ~~~~ 語音日期: 2026-02-02
\newline
\newline
\hyperref[sec:031641]{\small{< < < PREV SERMON < < <}}
~
\hyperref[sec:index_chronic]{\small{[返順時目]}}
~
\hyperref[sec:031643]{\small{> > > NEXT SERMON > > >}}
\newline
\newline
路加福音 1:66
\newline
\begin{longtable}{cl}
\hline
\hline
章節 & 經文 (和合本修訂版)\\
\hline
1:66 & \begin{tabularx}{0.7\textwidth}{X} 凡聽見的人都把這事放在心裡，他們說：「這個孩子將來會怎麼樣呢？」因為有主的手與他同在。 \end{tabularx} \\ \\
[1ex]
\hline
\hline
\end{longtable}
$^{1}$大家午安.
時間過得好快.
眨眼又一日.
聖經人聖經士在生命說書的網絡向大家提供.
有主與他同在是今天的主題.
有主與他同在仍然是講施洗約翰有關的經文.
路加呼音一章57節開始.
伊利沙伯即是施洗約翰的媽媽.
產其道了生了一個兒子.
倫理親族聽見主向他大施憐憫.
和他一同歡樂.
真是好.
沒有人說你真是甚麼.
大家很誠心來恭賀他.
主真是眷顧你.
大施憐憫.
和他一同歡樂.
我看這段聖經和他一同歡樂.
我都覺得幾好.
有時見到有些人有感恩的事.
自己不開心.
就打打彈彈.
當然是甚麼.
其實不緊要.
與喜樂的人一同喜樂.
他好ye 就喜樂.
哪個兒子又拿到甚麼冠軍.
好ye 呀,好感恩.
一同喜樂.
就無謂拿到冠軍.
當然是媽媽逼到他.
不要說這些.
我們和他一同喜樂.
讚美他的兒子好ye .
他教得好.
都學到些東西.
一同喜樂.
第八日就要割禮.
大家知道可能我都要講講.
割禮就是今天的現代.

$^{41}$男生割包皮的手術.
所有的猶太人.
他們在第八日就要行割禮.
這個割禮是上帝吩咐必定要做.
是甚麼呢.
就好像是一個印記.
屬於我的子民百姓的男性.
會在第八日就要行割禮.
有時會認出這是猶太人.
所以當父大講一下.
當在二次大戰希特拉德國人要捉猶太人的時候.
他們會由他有沒有受割禮來檢查.
就知道他是否猶太人.
當他是猶太人.
第八日就要行割禮.
因為他有嬰兒.
傳聞是不痛的.
嬰兒不會吱吱叫.
行割禮是好處的清潔一點.
並且照他父親的名字叫他撒加里亞.
由於撒加里亞這麼大年紀才有個兒子.
不如就勸他改回撒加里亞第二.
就叫兒子做撒加里亞.
好不好呀.
他媽媽就說.
如果本來也好.
但是不可以.
不過這次要叫他做約翰.
親屬就說.
喂!約翰?.
我們成族人都沒有一個人叫約翰這些名字.
我們都是撒加里亞什麼什麼.
到時候那些名字約翰沒有叫過.
為什麼要用這麼特別的名字呢?.
他們就說不如問一下父親.
於是就打手勢問父親.
大家不妨試玩一下.
今晚想像一下.
你怎樣用手勢問一個不懂說話的人.
你兒子叫什麼名字.

$^{81}$這些手勢不知怎樣做.
就向他打手勢問這個孩子叫什麼名字.
其實哪用打手勢.
撒加里亞是說不到.
他沒有說聽不到.
你問他撒加里亞叫什麼名字.
就不用打手勢.
向他打手勢.
所以有時候太勇敢.
其實他聽到的.
不過說不到.
於是他就以一塊寫字的板.
寫上他的名字是約翰.
很稀奇.
你真的說他叫約翰?.
那好吧.
挺有趣的.
撒加里亞的口立時開了.
舌頭也舒展了.
說出話來稱頌實.
周圍居住的人都懼怕.
很害怕這件事傳遍了那裡的山地.
凡聽見的人都把這件事放在心裡.
這個孩子出生得這麼特別.
將來會怎樣呢?.
很明顯是有主和他同在.
其實我覺得撒加里亞和施洗漢的親屬也不錯.
沒有說太奇怪,撞到什麼.
沒有說不好的事.
全部都是這樣.
這個孩子很正面.
有主在他身上.
將來會怎樣呢?.
沒有說看看他將來怎樣.
很正面,挺好的.
我想和大家看一些經文.
第一.
不可叫撒加里亞.
他媽媽說不可以.
一定要叫約翰.

$^{121}$為什麼呢?.
因為他堅持著.
我這個孩子有在.
是加伯列尼特別和我丈夫說的.
起了名字.
所以我不想第二個名字叫約翰.
你會知道其實.
最容易做是什麼呢?.
叫撒加里亞約翰.
約翰撒加里亞.
很多時候都會這樣.
他很厲害.
或者改什麼呢?.
巴撒加里亞.
Son of 撒加里亞.
Ben.
兒子,很多時候都有巴的姓.
其實譯做巴字,Ben就是兒子.
其實他可以叫約翰撒加里亞.
或者撒加里亞約翰.
這樣就滿足到很多人.
不過不可以,我只是跟隨主叫約翰.
這一點值得我們學習.
第二,他爸爸寫字.
他的名字叫約翰.
第一點和第二點都是.
無論媽媽和爸爸都是指明.
以前我們真的沒用.
不相信.
終於啞了.
現在我一定見到主的能力.
我沒別的選擇,一定要叫他約翰.
這種堅定的信靠.
其實是很好的.
雖然這個名字我的家族都沒有人叫.
見到愛納時口就書張了.
開始稱頌神.
周圍居住的人都懼怕.
懼怕,懼怕這個字是什麼呢?.
Phobos.

$^{161}$Phobos這個懼怕真是.
驚慌,驚慌警覺尊重.
其實驚慌為主.
大家都知道有個英文字叫Phobia.
Phobia就是驚恐症.
驚恐症,不是驚恐症.
Phobia這個字是恐懼症.
恐懼症的意思就是.
見到一件事就害怕.
譬如有些人有恐懼症,有Phobia是什麼呢?.
見到雞毛.
嘩,不行,見到蜥蜴.
見到蟑螂,見到什麼呢?.
他整個人快要暈軟了,說不出話,走也走不了.
那些叫恐懼症.
Phobos這個字其實是從這裡來的.
Phobia,恐懼.
即是說很害怕.
嘩,又不說話那麼久.
忽然間說得出來,他們驚恐什麼呢?.
不是驚恐撒迦利亞.
是驚恐上帝的能力真是大.
上帝的能力真是厲害.
真是和我們同在.
所以他們將這件事放在心裡.
不看這個小朋友將來會怎樣.
又主動地同在.
好了,今天的金句.
若路加一張落路紙.
凡聽見的人將這事放在心裡.
鄧小妹,我很想你學到一件事.
很多事你靜心去觀察.
然後將它放在心裡.
特別將什麼放在心裡呢?.
其實我已經將很多事放在心裡.
記起很多不開心的事.
我心裡太多事了.
其實意思是.
你將一些很重要的事.
最好如果你的心空得那麼小.

$^{201}$你的腦海那麼小.
不要放那麼多垃圾東西.
將上帝的作為放在你的腦裡.
時時將這件神的事放在心裡.
看看將來神怎樣成就才行.
神開始了這件事.
看看將來怎樣成就才行.
這件事在教會裡開始.
大家都贊成開始.
不要說什麼時候不行.
又停了,人又走了不做了.
又有頭沒尾陣了.
將這件事放在神的手裡.
自己心裡記住.
看看神怎樣成就這件事.
當你有這樣的信心.
你是越會看到這個成就.
要記得這件事的成就.
是因為有主與他同在.
記得嗎?.
神的話必定應驗.
事情一定成就.
是因為有主與他同在.
這件事放在心裡.
多些將神的事放在心裡.
我們祈禱吧.
主我們要向你承認.
很多時候我們將很多事.
真的有些不太重要.
放在心裡.
哪個人對我做過什麼不好的事.
我記住.
旁邊有什麼危險困難.
我記住.
自己的遭遇怎樣.
記住.
教會裡發生了多少不好的事.
記住.
反而沒有將一些好的事.
沒有將神的作為記住.

$^{241}$我們看到很多事情是你的手在工作.
很多人不會信主.
不過現在又信了.
教會裡很多事應該是不發生的.
不過又發生了.
這個教會很難生存.
不過又維持到今天.
所以其實這些全部都是你的作為.
主我們看到你的同在.
看到你的話是必定應驗的.
因此我們將這些事放在心裡.
讓我們更多將聖經裡我們看到的話語.
看到在舊約新約裡.
你成就過的事.
我們記在心裡.
以致我們的信心是一路增加.
我們一路跟隨你.
是更加放心.
更加有信心.
求主教導我們.
將你的事放在心裡.
隨聽祈禱.
奉耶穌基督的名.
阿們.
\newpage



\section{路加福音 16:1-13}
\label{sec:031643}
\textbf{今世錢財,永恆眼光}~~~~陳傳華牧師博士
\newline
\newline
連結: \href{http://www.church.com.hk/acms/content.asp?site=cdc\&op=show\&type=product\&code=031643\&layout=sermon}{\texttt{031643}} ~~~~ 語音日期: 2026-02-08
\newline
\newline
\hyperref[sec:031642]{\small{< < < PREV SERMON < < <}}
~
\hyperref[sec:index_chronic]{\small{[返順時目]}}
~
\hyperref[sec:031644]{\small{> > > NEXT SERMON > > >}}
\newline
\newline
路加福音 16:1-13
\newline
\begin{longtable}{cl}
\hline
\hline
章節 & 經文 (和合本修訂版)\\
\hline
16:1 & \begin{tabularx}{0.7\textwidth}{X} 耶穌又對門徒說：「某財主有一個管家，有人向主人告管家浪費他的財物。 \end{tabularx} \\ \\ \relax
16:2 & \begin{tabularx}{0.7\textwidth}{X} 主人叫他來，對他說：『我聽到了，你做的是甚麼事？把你所經管的交代清楚，你不能再作我的管家了。』 \end{tabularx} \\ \\ \relax
16:3 & \begin{tabularx}{0.7\textwidth}{X} 那管家心裡說：『主人辭我，不用我再作管家，我將來做甚麼呢？鋤地嘛，沒有力氣；討飯嘛，怕羞。 \end{tabularx} \\ \\ \relax
16:4 & \begin{tabularx}{0.7\textwidth}{X} 我知道怎麼做，好叫人們在我不作管家之後，接我到他們家裡去。』 \end{tabularx} \\ \\ \relax
16:5 & \begin{tabularx}{0.7\textwidth}{X} 於是他把欠他主人債的，一個一個地叫了來，問頭一個說：『你欠我主人多少？』 \end{tabularx} \\ \\ \relax
16:6 & \begin{tabularx}{0.7\textwidth}{X} 他說：『一百簍油。』管家對他說：『拿你的賬，快坐下，寫五十。』 \end{tabularx} \\ \\ \relax
16:7 & \begin{tabularx}{0.7\textwidth}{X} 他問另一個說：『你欠多少？』他說：『一百石麥子。』管家對他說：『拿你的賬，寫八十。』 \end{tabularx} \\ \\ \relax
16:8 & \begin{tabularx}{0.7\textwidth}{X} 主人就誇獎這不義的管家做事精明，因為今世之子應付自己的世代比光明之子更加精明。 \end{tabularx} \\ \\ \relax
16:9 & \begin{tabularx}{0.7\textwidth}{X} 我又告訴你們，要藉著那不義的錢財結交朋友，到了錢財無用的時候，他們可以接你們到永遠的住處去。 \end{tabularx} \\ \\ \relax
16:10 & \begin{tabularx}{0.7\textwidth}{X} 人在最小的事上忠心，在大事上也忠心；在最小的事上不義，在大事上也不義。 \end{tabularx} \\ \\ \relax
16:11 & \begin{tabularx}{0.7\textwidth}{X} 若是你們在不義的錢財上不忠心，誰還把那真實的錢財託付你們呢？ \end{tabularx} \\ \\ \relax
16:12 & \begin{tabularx}{0.7\textwidth}{X} 如果你們在別人的東西上不忠心，誰還把你們自己的東西給你們呢？ \end{tabularx} \\ \\ \relax
16:13 & \begin{tabularx}{0.7\textwidth}{X} 一個僕人不能服侍兩個主；他不是恨這個愛那個，就是重這個輕那個。你們不能又服侍神，又服侍瑪門。」 \end{tabularx} \\ \\
[1ex]
\hline
\hline
\end{longtable}
$^{1}$今天的講題可以說是聖經裡面四福音.
甚至全本正經是最難解的經文之一.
甚至有學者認為這段經文是最難解的經文.
這段經文說什麼呢?.
在題目裡面就看到了.
今世前才永恆眼光.
今 永 時間性.
我大概有這個簡介.
物質性前才.
眼光是心裡的選擇.
這些難解經文以前當然要先破冰.
有些謎語難解的.
讓你有些感覺.
看看這些難解困惑的東西.
當然破冰應該是挺好玩的.
挺好笑的.
先問一下大家.
有什麼東西越洗越髒的?.
銀紙.
銀紙 還沒想到.
如果有髒的可以給我就行了.
也可能是的.
還有呢?.
水.
水 厲害 滿有智慧.
好了.
有什麼書人人要看.
但是很多人從來都不看的?.
《聖經》.
很傷心.
又好像很有道理.
雖然我的答案不是.
也適合下次加起來.
還有呢?.
我想問一下大家.
剛剛買完新的電話.
有沒有人看了說明書才用的?.
如果你看了說明書.
才用你的新電話的.
舉一下手.

$^{41}$嘩!Gabriel.
很準確.
真是奇人.
四堂崇拜.
我想有八百多人.
暫時只有兩個人舉手.
一個是你.
一個是我老婆.
厲害.
說明書.
《聖經》我們要多看.
《聖經》要多看.
稍後再問下一堂崇拜.
再問愛倫姐.
是不是?.
OK 好.
什麼門是不需要留鎖匙的?.
什麼?.
阿們.
阿們?嘩!.
你的思維我真的跟不上.
真是厲害.
梁牧師就說如意門.
真是厲害.
阿們.
我們是想正常一點.
龍門.
自從阿阿們.
厲害.
加上去.
這個.
人類歷史中最長的車.
是什麼車?.
不夠長.
塞車.
嘩!誰想的?.
塞車你都知道.
真的.
塞車你知不知道.
塞過十幾二十公里.

$^{81}$你知不知道.
真是.
什麼人你都想到.
稍後在門口等我.
好.
什麼車是沒有輪子的?.
E車仍然夠適合.
都是新任務.
還有呢?.
風車.
這個真是跟我差不多.
還有自然浮列車.
好 最後這個了.
為什麼9+5=2.
12+4=4.
上一堂的人都這樣說.
為什麼?.
因為計錯了.
為什麼?.
3.
2.
1.
終於陰到你了.
為什麼?.
因為.
是時鐘.
我想說.
有些很困惑的問題.
是要有一個處境才能解答.
行不行?.
車沒有車都是處境.
或者月洗月滿就是一個處境.
今天呢.
這段難解的經文.
我們怎麼嘗試解釋.
也是一個處境.
這段經文呢.
是關於一個不義管家的比喻.
行不行?.
我們慢慢試一下.

$^{121}$為什麼這麼難解釋.
我們試一下能不能解釋.
也不知道.
我們嘗試努力地去寫一寫.
看看大家這樣看進去.
又會經歷些什麼.
相信是十分寶貴的教導.
要不然不會用這樣的進度.
耶穌所說的比喻.
是最令人困惑的一個.
主人為什麼會稱讚一個作弊的管家.
行不行?.
就是一個KK圓的人.
例如.
為什麼呢?.
那麼.
是不是耶穌贊成不誠實呢?.
贊成蒙騙呢?.
當然不是.
我們先看看是說什麼.
先看看經文的背景.
在盧家邦第十六章一到二.
一起讀.
一二三.
耶穌又對門徒說:反你所經管的,因你不能再作我的管家.
背景很簡單.
有人被炒魷魚.
因為.
他老闆.
財主.
請了一個管家.
他收到消息.
這個管家浪費.
浪費.
我為什麼說會和十五章有關呢?.
十五章是一個浪子.
耗盡.
浪費.
不應該做的事已經做了.
叫了你來.

$^{161}$他沒有解釋是什麼.
這件事的困惑就是很多東西留白了.
我們要看處境才能明白.
並且.
他告訴你.
你不能再.
作我管家.
交代明白.
這個背景.
炒了.
首先我們今天.
要知道我們要學幾點.
三點.
一個是關於時空有關.
一個是關於物質有關.
一個是關於你的選擇.
先看第一個.
還有第二個.
接著這件事發生了.
這個管家.
他想什麼做什麼呢?.
我們讀一讀第三到第七節.
預備起.
那管家心裡說.
主人持我.
不容我再作管家.
我尚能做什麼?.
做事呢?.
武力呢?.
道德呢?.
下手呢?.
我只做什麼?.
好幾年來.
尚可不作管家之徒.
自我道德無能.
於是把欠他主人債的一個一個地叫來.
你欠我主人多少?.
他說.
管家說.
一百樓有.

$^{201}$管家說.
拿你的帳.
寫了五十.
又問一個說.
你欠我多少?.
管家說.
一百四十.
管家說.
拿你的帳.
寫了八十.
對了.
他很清楚他的處境.
他快沒得.
什麼?.
做事.
你和我一樣地道.
沒得做事.
於是.
他又很清楚自己.
他說.
沒得做事.
怎樣?.
坐地.
武力.
有沒有人現在還可以懂得耕田.
可以賺到錢的.
玩玩有.
家中種一些辣椒.
辣椒可能.
香葉.
薄荷葉.
要找刀吃.
真是難.
黑錢討飯又怎樣?.
怕醜.
他很明白.
所以他想到一個辦法.
這個辦法就是.
我相信他的.
這個主人.

$^{241}$財主.
是借錢給人的.
所以他代他的主人.
他的財主.
去借錢收錢.
他就說.
我借了錢.
其中一個.
是一百口油.
其實直譯.
如果說對換價.
是八百加侖橄欖油.
等於多少?.
三四年的工資.
他說借給很多人.
其中一個三四年工資.
你想一想多少?.
你先想一想你多少錢工資.
然後告訴你旁邊的人.
你多少錢工資.
是不用的.
他真的說.
你想一想.
可能都二百萬.
有些多些.
有些少些.
幾十萬.
三四年工資.
可能賣樓.
不知道多少.
一斬.
斬多少?.
斬一半.
你開心嗎?.
如果你借錢買樓.
欠五四百萬.
突然間明天接了一張單.
你現在還二百萬就可以了.
你開心嗎?.
好像一般般.

$^{281}$如果你覺得不開心.
你給我吧.
當然開心.
另外一個.
一百石墨紙是什麼呢?.
是一百英畝的墨紙.
即是十年工資這麼大.
又斬.
斬了兩年.
不得了.
這個問題很有趣.
他說這個管家.
慷他人之概.
即是幫他老闆倒霉.
很奇怪.
如果你是受惠者.
你開心嗎?.
開心.
不過如果你是老闆怎麼辦?.
哭吧.
亂刀斬他.
什麼事?.
你都明白.
你一定不會感到開心.
不過這個事件奇怪的是.
耶穌繼續講解.
祂的主人.
即是祂的老闆的回應.
就是很特別.
祂這樣回答第八節第九節.
一二三.
主人就誇獎者不義的管家.
做事聰明.
之子更加聰明.
我又告訴你們.
要借助那不義的賤才結交朋友.
到要賤才做好事的時候.
他們可以自便到永存的帳幕裡去.
嘩.
耶穌讚他.

$^{321}$明天上班會做.
把老闆的錢都給客人.
然後讚賞客人的誇獎.
叫你過去做.
那就糟糕了.
絕對不是這個意思.
監也有.
祂不是講這件事.
我會今天.
剛才講有三個角度.
第一個角度.
祂突然間這個不義的管家.
不義的.
祂突然間知道.
有些東西突然間沒用了.
無用的時候.
你看到黃色是不義.
不義的錢財.
錢財沒用.
綠色的是誇獎聰明.
聰明結交朋友.
永存的帳幕.
祂突然間發現一件事.
可以沒了的.
原來有一種永存的東西.
或者是超越了現有的東西.
祂遺忘了.
第一件事我們一定要學的.
原來有些東西.
在祂這麼聰明的東西.
才發現到.
現有的東西.
不是等於.
你可以永遠保留.
第一件事.
我先解釋一下.
一般的聖經學者.
有三個可能性.
可能你想到第四個也不一定.
覺得這個處境是這麼發生的.

$^{361}$什麼叫三個可能性.
第一.
這個管家突然減少主人的債.
當時.
就是說.
一百變八十.
那個動作是不義的.
這是第一個可能性.
第二個可能性就是.
他之前放貸的時候.
他收多了.
就是說本來借一百元給人.
是應該收十元利息的.
誰知道他就收四十.
那也是不義.
現在減了.
其實常常都有的.
以前六七十年代警察也是.
收錢的.
到處收.
又或者.
在當時羅馬也是.
收稅的人.
政府要收稅一元.
但是收稅的稅利.
你自己可以收兩元.
自己一元也可以.
但是不公義的.
所以第二個可能性就是.
他當時放貸的時候.
已經是九出.
十三龜.
你真聰明.
不對的.
所以他開始.
第三個可能性.
就是當時他放貸的時候.
主人已經收多了.
所以他斬了主人那個.
對我來說.

$^{401}$我催向相信是第二個可能.
因為他是不義.
所以他現在要扭轉.
他不義的管理.
如果這樣的話.
就是說他本來可能.
剛才說的那些.
一百樓友.
或者是一百石物資.
本來可能不是收這麼多的.
斬了.
如果假設.
這個想法是.
他開始知道.
他搞錯了.
怎樣運用神本來給他的資源.
第二點我會多講一些.
調動資源.
第一點我先講一下.
又是beyond.
是超越了你自己可以考慮的時間.
突然間去知道.
所以耶穌說的不是要學管家的繼續不義.
或者以往的不義.
是突然間有種迫切感.
和前瞻性.
明白錢財用途.
很短的.
沒有用的時候.
並且可以因為物質世界沒有用的錢財.
竟然可以轉為.
永存的賬目.
這裡要聽清楚.
他不是說你去到永存的賬目.
看看你會怎樣.
是你沒得去的.
就是說如果你搞不清楚.
上帝給你的金錢管理好.
是沒得去的.
明白了嗎.

$^{441}$如果你明白你開始慢慢get到.
我很記得.
我有時候搞錯了.
計錯數.
長遠等等.
在COVID的時候我們知道常常有什麼.
限聚令.
二人限聚.
對我來說.
其實也不是很大影響.
因為二人和老婆拍拖.
很開心的.
四人限聚也不是很影響.
因為和子女一起去走.
所以我們會有行山.
我記得有一次我太太.
就安排行山的路程.
我們不是行山銀行.
我們真的告了假去行山.
從哪裡呢.
從荃灣.
我老婆設計.
走回火炭行不行.
走回我家.
不知道多遠.
保守估計也四五個小時.
那就走吧.
我們拿了水.
我慢慢走.
走到快要一公里到.
因為我們計算到走幾公里.
我就喝了水.
喝了水還有一個公里.
怎麼越走越向上.
越來越風景繼續漂亮.
我就想.
情況不對.
於是就問那些行山人士.
很專業的那些.
樣子好像經常走的那些.

$^{481}$雖然兩個限聚令都要問一問.
火炭在哪裡.
然後那些專家就說.
你去火炭.
你完全走錯方向.
你至少還有一半路程.
我結果走了八個小時.
我差點渴死.
三十幾度.
我老婆昨天知道我講這個例子.
她找回一張照片.
我真的看.
然後當時跟我拍一張照片.
我的樣子是怎樣呢.
我真的差點死了.
水又沒得喝.
沒得喝三個小時.
然後我記得一下下去.
沙田那個位置.
有間餐廳.
第一件事你知道我做什麼嗎.
馬上喝東西和吃西多士.
是人生最好吃的.
很恐怖.
然後我就看著老婆.
為什麼看錯了.
然後我老婆就說.
其實是這樣.
你記不記得一開始走的時候.
我拿個電話出來.
你揹著我.
因為我很怕大家一起出去.
你又看電話.
不是一起享受風景.
和享受美麗的老公.
然後我就揹著她.
我老婆就把電話放回去.
然後她就說.
你知不知道我為什麼拿電話出來.
我說不知道.

$^{521}$她當時在找路.
在看Google Map.
然後指著我.
是你害我.
是啊.
我害了她.
我算錯了.
我不知道原來路程可以很遠.
這個畢業管家他搞錯了.
他說他可以掌控地上的東西.
他不知道還有很多路要走.
但是他可能走不了.
所以這一次的困難.
是幫他醒過來.
非常重要.
永存的將來.
不只是今生.
如果我們今天不去管理自己.
恐怕去不了永存的將來.
還有未來的.
你知不知道.
我們怎樣去準備好.
所以就接到第二點.
第一點我叫永恆短暫.
其實有時候我們真的不知道.
很多事情突然就會發生.
畢業的管家.
共德偉自己開始將來打算.
發現他現有的模式.
是行不通的.
Doesn't work.
所以他要用現有的資源開始鋪路.
所以這樣被稱為聰明.
當然開始支援不同的工作.
我們會在想.
第二.
我們繼續學.
就是著重於開始調動資源.
如果這樣的話.
我們怎樣調動今天的物質世界.

$^{561}$好叫做影響那種永恆.
和屬靈世界.
這個非常奇怪.
好了我再看第八第九節說什麼.
1 2 3 主人就誇張者不義的管家.
做事聰明.
告訴你們.
要借著那不義的善財結交朋友.
他誇張和聰明結交.
這裡說開始我們看第二個部分.
就是原來他開始動得調動.
我們看到他在之前怎樣調動.
所以我們在學的是什麼呢.
我們開始心態要改變.
是從我的錢.
我的時間.
我的資源.
我的財幹.
我的恩賜.
改為是神託盤我管理的資源.
管理的財幹.
管理的恩賜.
管理的時間.
這個叫做很大的Planetary Shift.
而這個Planetary Shift.
令到我們可以知道.
上帝供應給我們.
用在地上之外.
我們還可以進入到.
永存的帳幕.
好了.
就這件事.
我也有些分享.
我在迎賓臺也分享過.
我記得我中午信主.
中六大學.
我一信主的時候.
我決定一定要比服事主好.
無論在教會.
在職場.

$^{601}$或者在讀書的時候.
傳福音.
團契.
不同的地方去服事.
而同時我記得在大學的時候.
一直學習真理.
我也決定一件事.
就是將來出來工作的時候.
我知道福音是需要金錢去到的.
宣教士.
扶貧.
所以我說一定會盡忠去到奉獻.
不過當我進入職場.
我發現這個世界太多東西玩.
我要去旅行.
我要儲錢結婚.
我還要買東西.
開始我就奉獻一些.
又不奉獻一些.
頭一兩年開始有掙扎.
我明明也很決定.
為什麼開始有掙扎.
去到第一,二年.
我有一次,有兩次的感受.
祂開始跟我說話.
我記得.
不知道第一,二年.
出來工作九三九四左右.
我們有不同的團契.
不同的小組.
坐下來.
有一次臨過節就圍圈分享.
就有人拿了一些金句卡.
就每人抽一張.
金句卡每人抽一張.
通常金句卡是說什麼呢.
耶和華是我的牧者.
很溫暖的.
耶和華是我的避難所.
又或者將你的憂慮交給主.

$^{641}$大家一抽就說.
我抽到這一張.
因為上帝要將我的憂慮交給.
我一抽.
就是當立十分之一的風雲.
什麼金句啊.
誰放這些東西在這裡啊.
最慘就是要逐個分享.
多尷尬啊.
又不能說謊.
我是真的有些掙扎.
那次之後過了兩三個月.
我不知道是聖誕節還是農年.
我掙扎的.
也好一點.
也是掙扎.
又第二個聚會.
另一群人.
又有人抽金句卡.
這次我想一定是另一些.
有些人也是這樣抽.
耶和華是衛王.
又什麼很好.
又是供應者.
我一抽.
當立十分之一.
不是吧.
我還沒分享.
我跟主說.
明白了.
一定會.
以後都會.
明白了.
不單止.
一向都是.
之後就.
還有可能騙我們.
或者有需要.
他家裡有病.
或者宣教等等.

$^{681}$明白了.
馬上明白.
我覺得很好.
上帝提醒我.
我要集中是什麼.
因為金錢從來都不是忠誠.
一會兒最後怎麼說.
很吸引.
好誘惑.
多麼堅定.
都會被他吸走.
我知道.
因為牧會了二十年.
我做了職場十幾年.
就是牧會.
接著我又看到.
想起很多人.
加薪.
或者升職.
但是沒有反應.
記住.
這段經文不是說.
有需要貧乏的人.
貧乏一天.
是我們會供應.
這裡是說.
那個管家是有資源的.
所以.
這一次的經文.
主要集中有資源.
沒有資源的.
就是供給他.
因為受惠者.
都有提.
好了.
有些就多謝主.
都沒有了.
不過我看到很多.
愛主的弟兄姊妹.
或者門徒等等.

$^{721}$他知道.
臣賜予他.
他不單是自己受惠.
又可以給家人.
以外.
他要在主的國裡面.
去分配.
不單止他們奉獻.
不單止他有宣教士.
又支持.
打風的.
不斷追.
陳牧師有什麼需要.
菲律賓我們有短村隊.
緬甸有什麼需要.
還有我們在街坊的需要.
很多都說.
有些人生了.
有什麼可以幫他.
有些突然間街坊病了.
抓不到工.
有什麼可以幫他.
有些不單止這樣.
奉獻十三.
十四.
因為給他足夠有餘.
有些感恩.
順著祝福.
又買多一層樓.
他不是很貴.
租出去賺到盡.
其實這個經文是說到一件事.
叫賺到盡.
一百.
賺到盡.
斬.
有些人就將他.
50\%很便宜.
就租給貧乏人.
租給那些很壞的人.

$^{761}$這種是調動神之源.
在神國有份.
又不一定是職場的.
有些機構.
我不知道大家有沒有聽過.
雲彩機構.
就是1月來的分享.
將緬甸雲南那些老蛋的小朋友.
幾百計一千計.
甚至有些去了KK園救出來的.
都是這個機構幫他的.
做了很多年.
他分享異象.
很多人就要奉獻幫他.
才可以完成.
運用資源.
資源不是私有化的.
在主國裡面一起去運用.
我聽了很多.
不單止這樣.
有些已經可以退休.
他說不行.
我要為主多做幾年.
繼續叫到神加油.
光明之子和今世之子是什麼人呢?.
今世之子當然是屬世的人.
不是因為屬世人就沒有智慧.
有的.
他有種敏銳度.
迫切性.
光明之子當然是屬神的人.
不過他有些東西不足.
這裡我說的很重要.
就是今世之子.
在世的好處就是.
他有策略的.
他想到東西的.
聰明的.
他懂得鋪路的.
光明之子要學習.

$^{801}$不應該為永恆.
起不應該更為永恆的歸宿.
積極調動資源.
很多時候光明之子.
神會供應.
神會安排.
神會供應會安排.
不過直著神的指紋.
我們要計劃的.
非常重要的.
光明之子常常缺乏這種調動的資源.
所以更加要學習.
積極迫切的智慧.
可以有策略為永恆將來投資.
所以你看到我們教會.
扶貧探訪是我們的主旨.
是不會不做的.
是一定要這樣走的.
不單止是教會.
教會是什麼?.
就是你.
教會就是你.
你跟別人說教會就是你.
所以你去探訪.
你去做事.
很重要.
而我們更加知道教會和職場的弟兄姊妹.
其實是攜手同行.
裡面也有天國用我們的地方.
有幾件事我們要注意的.
第一.
我們職場的不是等於我們不是在做神的工作.
我們肯定是呼召我們在職場不同的領域.
並且在教會的職份以外.
我們天天都是神眼中有價值.
服侍主的人.
第二我們在職場要明白我們的確有那種壓力.
和有一種試探.
那種吸引.
那種倫理的誘惑.

$^{841}$那種金錢很容易跌倒.
我們要小心.
這個經文提醒我們.
這個人在職場.
他就要賺到盡.
不義地賺到盡.
我們不可以這樣.
然後我們更加要知道.
其實職場也是一個宣教的地方.
我們要好好做這個職場.
叫到人認識主.
這個不義的管家.
他將他不義賺到盡的東西減去.
令到人受惠.
我們要很小心.
所以因此這樣的話.
我們要跟職場的領事館攜手合作.
聰明調動.
就好像昨天我們在其中一個例子.
不單止在職場.
在教會也一樣這樣合作.
我們昨天開執事會的時候.
我們很感恩.
分堂.
三年前的職堂.
當時四十多人去.
今天已經有三倍的人了.
有人信主有人受洗.
然後昨天就去探訪了.
榮福縣的很多居民.
其實我們知道.
榮福縣需要我們過去.
報道他立刻去籌款.
很短時間有七八萬.
每一元每一毫都給了他們.
我們不單止這樣去探訪.
是很積極.
一直去到探訪關心.
這個很真實.
而我們昨天知道.

$^{881}$我們開會的時候.
分堂說.
他的租約滿了.
要再談了.
最好談租約的當然不是牧者.
我們對商場已經有些轉折了.
那怎麼辦.
當然是執事.
那些在職場打滾的智慧執事.
你去搞定他.
很重要這種合作.
互相配搭.
在教會服事.
在外面服事.
非常重要.
然後.
我還想在第三點前說一件事.
剛才我講到.
這個畢業不單止是他開始走回那些東西.
他將過去扭曲了的東西都要調配.
所以我們如果要很運用金錢.
我們有一樣東西.
就是我們自己要管理好自己的財富.
很多大型節目都是簽卡過度.
很多大型節目都是亂消費.
這個要很小心.
其中.
我有時聽到很多的頂住.
哎呀.
今日屋企又爆了.
有人來追債.
某個家人又不知做了甚麼.
真的很害怕.
網上的笑話關於這件事都幾好笑.
那裡有個叫陳生陳太.
實述.
巧合而已.
兩個work from home 家裡工作.
突然間有個訪客按門鈴.
陳太開門.

$^{921}$見到訪客很陌生.
不知道是誰.
問你找誰.
找陳生.
他都知道.
立刻不要說陳生在家.
你找他甚麼事.
訪客說關於一筆貸款.
陳太立刻很驚訝.
立刻說陳生.
他今早出門了.
然後訪客說.
真是可惜了.
我本來借了陳生幾十萬.
今日打算還給他.
接著陳太立刻說.
等等.
他今早出門.
下午回來.
有時很害怕.
按門鈴究竟是甚麼事.
收甚麼信.
作為牧者.
我聽到真的太多了.
先好好管理你的財富.
你才可以好好運用上帝給你.
上帝才可以給你更多東西管.
其中一個見證很真實.
我認識一位姊妹.
她說到十多年前.
她工作很厲害.
工作很好.
高薪厚職.
不過高薪厚職不等於.
你很運用金錢.
她不斷買東西.
更加借很多錢.
更加買很多東西.
越買越多.
結果她欠了很多錢.

$^{961}$她有一次在零收退收裡面.
陳太立說你這樣不行.
你要和你老公說.
他老公都不知道.
她說說.
他會和我離婚的.
結果她和老公坦白.
一說之後.
她老公很驚訝.
但是很支持她.
一起去面對.
結果我們怎樣還呢.
她賣了房子還.
賣完房子都不夠還.
還欠幾十萬.
你說多困難.
結果我們很努力還完之後.
同時我們都很努力上班.
一向都很努力上班.
去到今天十多年之後.
我為她感恩.
她好好管理好財富.
她今天已經有一個老闆.
在台灣,香港,日本都有生意.
並且她很忠心的奉獻.
不單止這樣.
再奉獻之外.
還撥十分之一出來.
專門去照顧那些婚姻狀況.
婚姻輔導等等有需要的人.
並且不單止這樣.
她要決定好.
石木姐過年過節都會買東西給她吃.
我昨天就收到去買.
日本餅一盒.
我老婆說又有人送餅.
肥仔你小心點.
我想說甚麼呢.
她有這種空間.
是因為她好好管理.

$^{1001}$而令到她有一種自由.
一方面經歷上帝公認.
並且她可以去祝福不同的群體.
所以最後就關於這件事.
我們要知道.
我們怎樣去走呢.
我們要懂得選擇.
這段經文很重要.
耶穌將這段經文推到高峰.
講完這麼多.
那你怎樣呢.
今天永恆.
物質.
那你怎樣呢.
你怎樣呢.
我們讀一讀最後四節經文.
第十節到十三節.
一二三.
人在最小的時常中心.
在大事上也中心.
在最小的時常中心.
在大事上也中心.
人在最小的時常中心.
在大事上也中心.
在最小的時常中心.
恍若你們在別人的東西上不中心.
誰說我們在別人的東西上不中心呢.
一般人不得事無能助處.
誰說我們在別人的東西上不中心呢.
你們不得事無能助處.
這段很有趣.
她跟著講一件事.
她說你要懂得選擇.
並且你要細心地知道.
神會用細小的測試你.
中不中心.
我一看就很記得.
我剛剛大學畢業的時候.
我開始閃了.
但感恩啊神是.

$^{1041}$並且第十三節.
一個僕人不農.
事奉兩個主.
不是活者一個.
愛人一個.
重者一個 輕人一個.
所以不可以事奉神.
又不可以事奉猛鬼.
我們背一背.
就變成了雙層樓的巨無霸.
忘記了裡面講什麼.
她講什麼呢 其實這個講很重要.
她說馬門.
好像神一樣吸引.
它並排 看到嗎.
事奉的原文是Doulos.
有奴隸.
它可以吸引你成為奴隸.
今天很多人金城奴隸.
而吸引到你被綁架了.
並不自由.
所以她說你要選擇.
選擇什麼呢.
第三點.
就是你要實現你的真實的忠心.
小事上要通過測試.
跟這個世界智慧不同.
這個世界經常說.
做大事.
大事先行.
神說不是.
用小事先行.
好叫到你可以慢慢學習.
好叫到你可以.
一直學習的時候.
你累積著這種.
上帝的原則.
上帝的法則.
上帝的裝備.
所以忠心的東西.

$^{1081}$在小事上開始.
終極的抉擇.
不能又事奉神.
又事奉馬門.
錢從來都不是忠誠.
一忠誠就很危險.
一是擺在主裡.
主求你幫助.
好好的運用.
如果再不是.
你就私有化.
成為很危險的事.
這種對資源的理解.
和世俗文化價值觀背道而馳.
文化貫穿我們的思想.
我就是我的.
我喜歡.
隨心所欲.
耶穌說出.
這個是神託管.
並且會觀察我們怎樣使用.
而我們要記得.
上帝沒有不記得我們.
祂會供應足夠.
但是我們要很小心.
供應足夠的同時.
私有化了.
就將這個上帝.
所透過我們祝福的時候.
沒錯.
稍後的經文我們再看看.
不行不行.
我要解釋完剛才的經文.
我漏了一部分.
你看看第十一,十二節.
有個很有趣的東西.
他說.
如果你小事上不義.
大事上不義.
如果你不義.

$^{1121}$錢財上不忠心.
聽著.
誰會給你真實的錢財來管.
即是說神有些真實的東西在給你.
職業真正.
灌輸今天和永恆.
如果小測試你不忠心.
沒有了.
行不行.
本來祂準備了東西給你.
沒有了.
你再看下去更奇怪.
第十二節.
他說.
倘若你在別人的東西上不忠心.
不義管家別人.
誰把你自己的東西給你.
你聽清楚很奇怪.
原來測試著你怎樣對別人.
他才會知道.
他給不給你應該是你的東西.
你自己的東西.
不過你不忠心.
Amen.
記住.
神的東西好像流動的水喉一樣.
不斷流轉.
你越堵塞.
越不被使用.
所以這樣的緣故我們更加明白.
所以這篇道很多人聽了兩次.
他正在消化.
我覺得很精彩耶穌.
我們要堅守永恆的價值.
近來我們聽到.
有些組員突然間大病.
街坊說.
有些人沒那麼大儲備.
手停口停.
為了祈禱平平安安地去吧.

$^{1161}$不是的.
一定要籌款給他.
一起渡過難關.
額國書正正是評擊這種.
平平安安地去吧.
願你吃得飽槍得亂.
No.
不是這樣的.
很多地方我們可以好好運用這些資源.
好叫整個群體被建立.
並且如果我們調轉了.
追逐了金錢多些好些我的.
忽略上帝.
輕視家庭.
踐踏他人.
利用他人.
工作也是.
現在很多詐騙也是.
剛剛都有來.
有個稅局的WhatsApp電郵.
他說有退稅.
起發了.
Peter馬上提醒我們.
不要揭進去.
整個樣子是稅局的網站.
行不行.
你進去放你銀行的資料.
你明天的銀行就會轉了給他.
小心.
到處都是.
不是.
我們是神先行.
道德.
我們要關懷先行.
錢就不是那回事.
那回事不可以再控制我們.
這一種新的法則.
最近我和另一個信徒聊天.
他很忠心.
又釋放.

$^{1201}$又做了這裡那裡.
他什麼都做.
然後我就知道.
他又很積極建行.
很積極外面的需要.
奉獻很積極.
並且.
每個月以幾萬來奉獻.
不是多少個問題聽清楚.
是忠心的問題.
我問他為什麼這麼忠心.
他說我開始畢業的時候.
我都用幾百塊開始.
人工幾千.
神一直祝福他.
他說不要.
記住不要.
神越祝福就越害怕.
留下這些錢.
不是的.
神有你的份.
也用在神的家.
這一種新的思維.
這個世界不是這樣教的.
越多.
越在自己的袋子裡.
我想表達.
這一種新的思維.
有一個小的例子.
我想問問大家.
如果上去太空旅行.
上去太空旅行.
你會坐什麼交通工具.
才能上去太空.
坐什麼.
巴士.
可以嗎.
坐船可以嗎.
坐飛機可以嗎.
坐什麼.

$^{1241}$太空船.
火箭.
上到去.
如果你在太空艙的時候.
你作為一個太空人去旅行.
你很想離開太空艙.
去太空漫遊.
你要穿什麼.
穿棉襖.
不是吧.
羽絨.
頸巾可以嗎.
穿什麼.
太空衣.
因為那裡沒有空氣.
你要穿了才有氧氣.
新的裝備.
一個新的環境.
舊的思維.
舊的衣著不行的.
最近太空衣有一個新的設計.
叫Take Me Home.
什麼是帶我回家.
原來很多太空漫遊去做研究的時候.
因為可能太黑可能太光.
在太空是很難尋找方向的.
常常在做的時候.
找不到去哪裡.
雖然有條繩.
不過玩來玩去.
結果科學家.
發明了一個叫Take Me Home.
帶我回家的按鈕.
當你找不到一按.
它會怎樣呢.
帶你回太空艙.
今天耶穌給了一個Take Me Home的按鈕.
帶你回去永恆.
很重要.
不要忘記那個按鈕.

$^{1281}$也不要忘記它代表什麼.
我們也不要忘記.
原來上帝在小事上測試我們是很重要的.
你不要覺得我剛開始很少控制.
我覺得不是的.
覺得我的運功很枯燥.
不是的.
收入也很少.
不是的.
家庭的責任也很少.
不是的.
每一樣東西.
上帝都是測試.
我們對主的忠心.
兄弟姐妹.
上帝給什麼你.
給什麼你去管理.
總結和呼召.
今天我們有三點.
發揮贖天.
我們請勁拜隊的聰明.
為永恆而投資.
向世界之子學習.
懂得調動世上之元.
展現現實的忠心.
在小事上要通過考驗.
而中間我們要懂得怎樣運用.
其中有人問.
那我可以怎樣呢.
我簡單有兩個運用給你.
第一個.
我們開始要好好檢視我們的賬目清單.
每個月看看你的時間怎樣去分配.
你的銀行收入支出有沒有浪費.
你的消費紀錄又是怎樣.
你的財幹有沒有好好運用.
好叫你土共中民主啊.
我所經營這一切.
有哪些是浪費了.
有沒有將神的資源.

$^{1321}$浪費在短暫的事物上.
哪些可以有更好的帶來.
到永存的賬目.
永恆的投資.
你記住.
這段講不是我們給的人.
是受惠者帶我們去永存.
可能你要問我們救恩.
不是耶穌寫十字架嗎.
當然是啦.
那個最基本最重要.
耶穌寫十字架對我們的拯救赦免.
但是我們就要悔改.
要跟從.
要有恆敬.
才能完成這個救恩.
是不是.
而這個金錢管理是其中一個必然的行動.
如果我們經歷主的救恩帶來必然的行動.
就是我們生命改變.
悔改我們過去行出主的心意.
是不是.
明白嗎.
這樣的話第二我們就開始了.
當然我們的金錢管理神會照顧我們家人子女父母.
這個必然.
而永恆的投資.
那個當然和永恆有關.
當然我們就想.
咦這樣我們怎樣好好運用上帝給我們的資源.
奉獻到要奉獻.
要到位.
並且這個不是用淨才奉獻.
我當日也是上帝教得很清楚.
希望你不要收兩張金句卡.
一會兒門口抽一張金句卡.
慈惠.
怎樣去不同的機構宣教士.
或者資源時間運用.
真正事奉的是誰.

$^{1361}$好讓我們知道我們正在事奉的是神.
並不是馬門.
在《馬大福音》他這樣說.
最後一次經文我們一起讀.
好嗎?一二三.
好,你這又良善又忠心的僕人.
我不知道你們想不想學習成為又良善又忠心的僕人.
這是一個末世的經文.
二十五章.
他說如果我們忠心神會給我們更多的事委派.
並且在主人的快樂中有份.
我們一起祈禱.
天王多謝你.
今天你藉著這段經文向我們說話.
讓我們學習到原來上帝可以用物質世界影響屬靈世界.
今天又影響永恆.
用我們生命的選擇去作為我們神為神僕人的代表.
所以今天很想為你禱告.
主啊,求你教導我,加力給我.
讓我的財幹金錢時間資源運用合你心意.
如果你們其中一個很想這樣禱告.
這樣立志.
求主所用.
這一刻邀請你將手放在你的心.
我為你禱告.
天王我們多謝你.
我自己也是不斷禱告,讓上帝教導我.
不要又偏離主的帶領.
我們多謝你.
每一個訂閱節目都很真誠.
今天在你面前學習運用.
我知道所有資源本身都是上帝.
也都豐富充足有餘.
求主教導我們.
好好學習.
讓我們學習可以更加管理主給我們.
許多許多的事.
神我們多謝你.
你何等的真實.
教導我們.

$^{1401}$給我們有智慧.
有決心.
讓聖靈引導.
我們多謝主.
聽我們禱告.
奉耶穌名下.
非常讚的加了一個.
\newpage



\section{使徒行傳 9:32-43}
\label{sec:031644}
\textbf{福音超越界限}~~~~禤凱榮牧師
\newline
\newline
連結: \href{http://www.church.com.hk/acms/content.asp?site=cdc\&op=show\&type=product\&code=031644\&layout=sermon}{\texttt{031644}} ~~~~ 語音日期: 2026-02-15
\newline
\newline
\hyperref[sec:031643]{\small{< < < PREV SERMON < < <}}
~
\hyperref[sec:index_chronic]{\small{[返順時目]}}
~
\hyperref[sec:031645]{\small{> > > NEXT SERMON > > >}}
\newline
\newline
使徒行傳 9:32-43
\newline
\begin{longtable}{cl}
\hline
\hline
章節 & 經文 (和合本修訂版)\\
\hline
9:32 & \begin{tabularx}{0.7\textwidth}{X} 彼得在眾信徒中到處奔波的時候，也到了住在呂大的聖徒那裡。 \end{tabularx} \\ \\ \relax
9:33 & \begin{tabularx}{0.7\textwidth}{X} 他在那裡遇見一個人，名叫以尼雅，得了癱瘓，在褥子上躺了八年。 \end{tabularx} \\ \\ \relax
9:34 & \begin{tabularx}{0.7\textwidth}{X} 彼得對他說：「以尼雅，耶穌基督醫好你了，起來！整理你的褥子吧。」他立刻就起來了。 \end{tabularx} \\ \\ \relax
9:35 & \begin{tabularx}{0.7\textwidth}{X} 凡住呂大和沙崙的人都看見了他，就歸向主。 \end{tabularx} \\ \\ \relax
9:36 & \begin{tabularx}{0.7\textwidth}{X} 在約帕有一個女門徒，名叫大比大，翻出來的意思是多加；她廣行善事，多施賙濟。 \end{tabularx} \\ \\ \relax
9:37 & \begin{tabularx}{0.7\textwidth}{X} 當時，她患病死了，有人把她清洗後，停在樓上。 \end{tabularx} \\ \\ \relax
9:38 & \begin{tabularx}{0.7\textwidth}{X} 呂大原與約帕相近；門徒聽見彼得在那裡，就派兩個人去見他，央求他說：「請快到我們那裡去，不要耽延。」 \end{tabularx} \\ \\ \relax
9:39 & \begin{tabularx}{0.7\textwidth}{X} 彼得就起身和他們同去。他到了，就有人領他上樓。眾寡婦都站在彼得旁邊哭，拿多加與她們同在時所做的內衣外衣給他看。 \end{tabularx} \\ \\ \relax
9:40 & \begin{tabularx}{0.7\textwidth}{X} 彼得叫她們都出去，然後跪下禱告，轉身對著屍體說：「大比大，起來！」她就睜開眼睛，看見彼得，就坐了起來。 \end{tabularx} \\ \\ \relax
9:41 & \begin{tabularx}{0.7\textwidth}{X} 彼得伸手扶她起來，叫那些聖徒和寡婦都進來，把多加活活地交給他們。 \end{tabularx} \\ \\ \relax
9:42 & \begin{tabularx}{0.7\textwidth}{X} 這事傳遍了約帕，就有許多人信了主。 \end{tabularx} \\ \\ \relax
9:43 & \begin{tabularx}{0.7\textwidth}{X} 此後，彼得在約帕一個皮革匠西門的家裡住了好些日子。 \end{tabularx} \\ \\
[1ex]
\hline
\hline
\end{longtable}
$^{1}$如果二三十年之前大家想在家看電影.
你會做些甚麼.
當然是去影帶舖租帶看.
影帶舖很久都沒有聽過這個名詞.
如果大家在二三十年之前在澳洲居住的話.
大家記不記得最出名的影帶舖是甚麼.
不是說唐人那些.
一是CIVIC 一是BLOCKBUSTER.
這個是最出名的 差不多每一個區都有一間BLOCKBUSTER.
不過現在一間都找不到 為甚麼.
因為時代的巨輪底下.
影帶已經被另一個科技所取代了.
當這個電影開始在網上的串流媒體裡面.
成為一個潮流一個趨勢的時候.
當年的BLOCKBUSTER這個影帶舖.
它拒絕去改變它的商業模式.
拒絕不改變.
而當時另一間剛剛冒起頭的小公司.
叫NETFLIX.
現在大家可能都在用NETFLIX看電影.
這間小公司擁抱當時.
剛剛起步的一個科技的新趨勢.
最後在2010年的時候.
這個BLOCKBUSTER就宣佈破產.
而NETFLIX慢慢就成為了網上串流平台看電影的一個龍頭大哥.
可以這樣說 從這個事件我們發覺.
當我們在一個日新月異的商業世界裡面.
如果我們不不斷求變的話.
恐怕我們很快就會被時代所淘汰.
其實這個現象可以說.
在屬靈的世界某個程度上都是一樣的.
當然我們不是說改變我們聖經的真理.
來遷就迎合這個世界.
我們在說的是很多時候我們需要改變和突破一些.
我們因為自己的成長背景或者文化的差異.
而自己都不自覺地畫下了一些不正確的界線.
這些界線其實並不是神所設立的.
在今天的經文裡面.
聖經在告訴我們 神用了一個異象來挑戰當時的彼得.
來放下他一些的成見.

$^{41}$一些的偏見而所建立出來的牆.
這些牆對當時來講.
就將猶太人和外邦人分開了.
當彼得有改變之後.
整個福音的門就向外邦人的世界躺開.
這是一個很特別的時刻.
讓我們一看彼得在當中面對這個異象而改變的一個故事.
讓我們一同學習當中一些很重要的屬靈功課.
故事一開始的時候.
聖經就提到一個人物.
這個人物就叫做哥利劉.
他是一個該撒尼亞的人.
大家要留意.
他是一個外邦人.
他並不是猶太人.
聖經還說他是一個意大利型的白夫長.
就是說他是管轄著一百個羅馬人的.
從各方面來看.
根據這個描述.
你可以說哥利劉在以色列人或者猶太人的圈子裡面.
的社群裡面.
絕對是一個局外人.
不過如果我們再看聖經怎樣來形容描述這個人的時候.
我們會發覺很大很強烈的對比在當中.
聖經說哥利劉這個人是一個虔誠人.
他整個家族都是敬拜神的.
敬畏神的.
還多多周濟百姓.
常常禱告神.
如果我們這樣看的時候.
他所祈禱.
向神祈禱.
根據大部分的解經家和這個描述.
解經家都相信.
這個哥利劉當時所敬拜所禱告的神.
就是以色列那個神.
為什麼我這麼強調這個事實呢.
因為如果我們看回歷史的時候.
當時的羅馬帝國.
他們沒有一個單一的神.

$^{81}$是叫整個帝國的子民都來敬拜的.
因為當時有很多很多的神.
甚至當時的羅馬皇帝.
甚至要求他們的子民.
你要拜我.
你當我是神那樣拜.
當時大家知道.
羅馬帝國都有很多不同的民族.
有很多不同的文化.
他們各自的民族.
各自的文化都有他們自己的神.
當然了.
以色列的神就是其中一個.
我們要問一個很有趣的問題.
既然羅馬帝國當時有這麼多的神.
我們的問題就是說.
為什麼哥利劉.
偏偏選中.
我們的以色列獨一的真神呢.
我相信這個絕對不是.
點子冰冰點出來隨便.
信一個就剛剛被他找到這個真神.
我相信不是這麼簡單的.
我相信哥利劉是經過深思熟慮.
經過探求.
去做一個決定.
這個以色列的神.
就是一個我願意.
毀身.
跟隨.
相信的神.
或者我這樣問大家.
如果在這麼多的神在你面前.
你要選一個神來相信.
來跟隨的時候.
你會怎麼選的.
這麼多的神你怎麼選.
我不知道大家.
如果給我的話.
我當然選一個.

$^{121}$有大能力.
最厲害,最棒,最有能力的神.
我記得.
大家都知道我沒多久之前.
我去巴黎.
在巴黎的時候.
師母就看一個秀.
我說那些秀我不喜歡看.
我寧願在外面逛一逛.
逛的時候就看到.
很多的牆壁介紹印度教.
大家知道.
雖然印尼是伊斯蘭教.
但是巴厘島就是印度教.
Hindu.
我在牆壁裡面看.
他們怎麼介紹他們的Hindu神.
我發覺一件事.
原來Hindu有很多神.
而這些神不是每一個都有同樣無限的能力.
有些是厲害一點.
有些是弱一點.
沒那麼有能力.
不同的神是有不同的事物.
對我來說我覺得很奇怪.
既然能力不是很高的時候.
我為什麼要相信他呢.
是不是.
你瞬間要相信一個最有能力的那個.
除了能力之外.
我覺得另外一個我一定需要的就是.
你的神一定要很可靠的.
我可以靠得到他的.
不知道大家知不知道.
希臘有很多神.
希臘有一些神.
就是每逢冬天的時候.
他就.
可能是放冬假.
冬眠我不知道.

$^{161}$但是他不在這裡.
你在冬天的時候找他找不到的.
所以有一些希臘的神.
要到春天快要來的時候.
那些信徒那些信眾.
就會在他們的神廟那裡唱歌.
跳舞.
為什麼這麼興奮呢.
就是因為要做很多的聲量.
獻祭.
為什麼要做這麼多聲量呢.
叫醒那些睡了覺或者走開了的神.
你回來了回來了.
以致他們的五穀可以得到豐收.
沒睡覺就回來了.
大家覺得很奇怪吧.
如果我相信一個神.
我要他的時候.
他放了假.
他走去冬眠的時候.
這個神就值得我相信的.
除了能力.
可不可靠之外.
我覺得另外一個很重要的就是.
這個神要有憐憫的.
是可愛的神.
不知道大家有沒有留意.
你去一些神廟看的時候.
有一些的神是凶神惡煞的.
有一些的神是有動物的形象的.
老實講.
如果我相信這些神的時候.
我怕我晚上不用看恐怖片.
我可能每晚都會做噩夢.
我就想.
為什麼我要相信這些神.
而另外我覺得.
更加重要的就是.
我要相信一個是可信才行的.
合乎我自己過到我的理性那一關才行.

$^{201}$我記得.
當我還是很小的時候.
遠遠我還沒有認識基督教.
我的母親就跟我說.
你要葬香給祖先.
因為他們會保佑你讀書好.
健康.
將來找美麗的老婆.
保佑你的.
要葬香.
我說很棒.
死了的親人原來變了神.
祂會保佑我們.
這樣就容易了.
我只是葬一支香.
接著.
我的母親就跟我說.
某個時候.
你要放一些雞,水果.
放一些不知什麼東西.
來供奉祂.
接著要燒一些東西給祂.
燒一些錢.
燒一些車,小妹.
什麼都燒給祂.
我問為什麼.
因為祂要吃東西.
祂要錢花.
我問我母親.
母親.
究竟是祂們在照顧我們.
還是我們在照顧祂.
我感覺好像是我們在照顧祂多一點.
對不起.
如果我要選一個神的話.
要過到我的理性.
我相信大概哥尼留.
他找一個神.
其實.
和我很類似.

$^{241}$他找一個值得信的.
值得毀身.
值得去跟隨的神.
我相信哥尼留在神面前的信心.
是真實的.
為什麼.
因為他在行為上表現出一個真實信徒.
應該有一個行為.
因為聖經說的.
他周濟窮人敬畏神.
常常禱告.
而我相信他的禱告.
也都達到神的面前.
為什麼.
因為聖經說的.
聖經說.
有一天他見到天使.
天使就跟他說.
你的禱告和你的周濟.
達到神面前.
已蒙紀念了.
然後他還叫哥尼留派人去約柏這個地方.
叫一個叫彼得的人來到你家裡.
大家記得約柏這個地方.
上個星期我都說過.
約柏這個地方.
就是彼得在那裡將一個死了的女信徒.
叫她復活.
然後彼得就留在那個地方.
住在一個皺皮匠的家裡.
而當他仍然留在約柏的地方的時候.
神的天使就去到哥尼留那裡.
叫哥尼留派人邀請這個在約柏的彼得.
去到他家裡.
雖然天使沒有說接著會發生什麼事.
不過如果我們看.
天使和彼得所說的話.
我們都知道裡面大概會是怎麼樣.
天使就跟他說.
你的禱告已經達到神的面前.

$^{281}$已蒙紀念了.
很明顯.
就是說神聽了彼得的祈禱.
也都會應允他的祈禱.
我們問究竟這個哥尼留和神祈禱的是什麼呢.
如果我們看遠一點的話.
我們看後一些的44和45節.
彼得去到哥尼留那裡.
跟他們講福音.
講到還沒講完.
聖經就說聖靈降在當時一切聽到的人身上.
當聖靈降臨的時候.
就是說這班人得著救恩.
好了.
既然聖靈降臨.
甚至在彼得還沒講完之前.
這樣很明顯.
神差派彼得去哥尼留這個地方.
有什麼目的.
目的就是.
將這個哥尼留或當時的外邦人.
他們知道這個神.
但是彼得就幫助他們.
將這個對神的認知.
轉化成為在基督裡面的信心.
也都是將他們本來是一個尋道者.
將他成為一個耶穌基督的跟隨者.
記不記得主耶穌曾經這樣講過.
主耶穌曾經說.
凡祈求的就得著.
尋找的就尋見.
叩門的就給他們開門.
所以在座的朋友們.
如果你仍然在尋找.
哪一個神值得我信.
值得我跟隨的話.
和你今天聽到這個信息的話.
很有可能.
神正在聽你的禱告.
神正在應運你的禱告.

$^{321}$神聽到你敲門的聲音.
現在為你開門.
等著你來歸回神的家裡面.
二千年以前.
神聽哥尼留的禱告.
今天同樣都會聽我和你的禱告.
好了.
或者我們將鏡頭轉一轉去看看彼得.
彼得.
這裡說他有一天.
祈禱的時候.
見到一個異象.
那個異象就像一塊布一樣.
裡面包了很多不同類型的.
走獸,昆蟲等等.
從天降下來.
然後有聲音跟他說.
彼得彼得.
你起來.
再聊天.
對很多人來說.
正呀有東西沒吃.
不過.
對於當時的彼得來說.
他一直很嚴守聖經裡面.
舊約的食物的規例.
所以要叫他吃這些野味的話.
對當時的彼得來說.
是難以想像.
難以接受的.
而聖經怎麼說.
彼得就說.
主呀.
這是不可的.
我從來都不吃這些東西.
你現在叫我吃.
搞不定搞不定.
然後那些異象,走獸就全部回到上去.
同樣的異象.
同樣的對答.

$^{361}$發生了三次.
三次.
我最喜歡問為什麼.
為什麼三次.
為什麼發生三次.
有些人說.
中國人說.
重要的事講三次.
牧師這麼簡單.
重要的事當然要講三次.
那就發生三次.
如果真是重要的事講三次的話.
那麼每一個聖經裡面所講的異象.
都應該發生三次.
因為神給的異象全部都重要.
你敢說神給的異象不重要.
但不是.
聖經裡面.
沒錯曾經有一些是重複過.
而也有一些是用不同的方式.
來傳遞同一個意思.
但是你說.
完全一模一樣的異象發生三次的話.
是這裡而已.
為什麼三次呢.
我相信原因很簡單.
就是因為這個彼得.
彼得這個人相當的頑固.
很固執.
不容易.
不喜歡.
妥協的一個人.
如果大家記得的話.
這個發生三次的情況.
其實都不是第一次發生在彼得身上.
有什麼例子又是三次發生在彼得身上.
如果大家記得的話.
就是在大祭司的院子.
當時有人抓了耶穌去.
彼得就不敢過去.

$^{401}$遠遠跟隨.
就在大祭司的院子裡面被人認出.
你不是和耶穌一黨的嗎.
耶穌說不是啊,我不認識他.
三次發生了.
三次.
另外一次呢.
另外一次就是耶穌在提比利亞海邊.
當耶穌在復活之後向門徒顯現.
耶穌就三次跟彼得說.
約翰的兒子西門,你愛我嗎.
三次.
彼得說,哇,你為什麼問我三次.
神啊,你什麼都知道,你知道我愛你.
耶穌問了三次.
看到嗎.
彼得似乎每天都要三次.
也是三次.
我想問問大家.
如果今天神有一些訊息跟你說.
或者有些事情需要改變.
或者有些事情你需要去做.
你猜神要跟你說多少次.
你才會收到.
彼得就要三次.
才聽得到這個訊息.
什麼訊息啊.
神所潔淨的,你不可當作活物.
我想留意大家聽一聽這句話.
神所潔淨的,你不可當作活物.
其實這句話不是.
單是在說食物.
或者不是在說食物.
是在說其他任何我們對某些事物的迷.
不斷地執著它.
當作真理地執著它.
對彼得來說.
他的世界其實很簡單.
分開潔淨與不潔淨.
樣樣都是這樣看的.

$^{441}$無論是食物,人,機構,組織.
總之就是你潔淨就是潔淨.
你不是潔淨就不是潔淨.
這個界線成為了彼得一個.
他的人就是這樣.
他的生命的身份.
他對神的信徒.
全部都是兩樣東西.
一是潔淨,一是不潔淨.
對彼得來說.
也許是食物的規例.
但對於我們來說.
未必是食物.
可能是什麼.
可能是政見的不同.
可能是文化的偏見.
或者是事工的模式.
我是要這樣做這個事工的.
又或者是個人對好與壞的判斷.
這些都可以成為我們.
潔淨與不潔淨的區分.
我們是無意之中建立了很多的牆.
這些牆是上帝從來沒有設立的牆.
是我們自己做出來的.
而這些牆對我們來說是好東西.
但是卻是可以成為一個牆.
拒人於牆外的.
好了.
耶穌來到這個世界.
在十字架上其實付了很大的代價.
為什麼要付這個代價.
就是要拆毀這些牆.
這個意象.
彼得所看到的意象.
並不是在說食物.
而是在說人.
神說.
我所潔淨的.
你不可當作俗物.
好了.

$^{481}$這個現象完了之後.
彼得還在想.
這個意象究竟是什麼意思.
神其實就是一個宣稱.
以前我們覺得不潔淨的東西.
其實現在已經是潔淨了.
為什麼會有這樣的改變.
全因耶穌基督在十字架上所成就的救恩.
神已經因為耶穌所成就的救恩.
將人一直以為的潔淨與不潔淨之間.
那條分界線.
中間所建立的牆.
神將它完全拆毀.
我們再看.
雖然彼得.
正經說形在差異之間.
究竟這個意象是什麼意思.
雖然同一個意象剛才說過發生過三次.
但彼得仍然在拗頭.
究竟為什麼.
仍然在考慮和困惑的時候.
哥尼流所派的人已經去到他那裡.
聖靈就跟他說.
有三個人來找你.
你跟著他們去吧.
不要在那裡想.
因為是我猜他們來的.
雖然彼得未完全明白意象的意思.
但很明顯.
神這樣說的時候.
這三個來找他的人.
很明顯跟這個意象是有關係的.
縱然他不明白.
但他仍然選擇聽從聖靈的指示.
這個故事告訴我們什麼呢.
很明顯.
當時意象對彼得的果效已經在發揮.
彼得已經開始思考這個看不到的界限.
一直困了他這麼多年.
讓他定義什麼叫潔淨.

$^{521}$什麼叫不潔淨.
什麼叫猶太人.
什麼叫外邦人.
這種議員的觀念.
重新在這個意象裡面挑戰彼得.
彼得也開始重新思考.
我們再看下去.
接著彼得就下去見那些人.
留意這些哥尼流派的人.
不是猶太人.
是外邦人.
如果我們對彼得有認識的話.
根據他以前一直以來的老外的性情來說.
我代入彼得.
我猜他應該會這樣說.
對不起.
我一直都不跟外邦人打交道.
更加不要叫我去他家.
跟他聊天.
跟他吃東西.
不要搞我.
以前的彼得是這樣回應的.
但是在三次同樣發生重複的意象之後.
再聽天使怎麼跟哥尼流的對話之後.
彼得也不是蠢不是傻子.
他將所有的事情聯繫在一起的時候.
他馬上知道神想他做些什麼.
所以彼得就邀請這三個外邦人.
跟他進去住了一宿.
留意我.
叫他們一起過夜.
不是喝杯茶.
你們留下地址.
我自己去就可以了.
謝謝.
拜拜.
不是這樣.
是叫這三個外邦人進來跟他一起共宿一宵.
對於一個極度猶太人傳統的彼得來說.
這個是很難想像的.

$^{561}$不過彼得做到.
而說到外邦人.
我們不得不提.
使徒保羅一句已忽所思的說話.
保羅這樣說的.
當你們這班外邦人仍然未認識耶穌基督之前.
在神所應許的諸約上面.
你們都是局外人.
不過現在你們在耶穌基督裡面.
信了耶穌.
靠著他的血.
已經得與神親近了.
跟我們猶太人可以合而為一.
以前被視為潔淨的猶太人與不潔淨的外邦人.
現在因為耶穌基督的救恩的緣故.
可以在基督裡合而為一.
看不看到這個世界有很多牆.
不知道大家同不同意.
民與民之間.
國與國之間.
種族之間.
社會地位.
財富.
名譽.
權力.
每一個生活上面.
生命上面的範疇.
都有一道牆.
但是主耶穌說.
無論你是要什麼牆.
在基督裡.
都成為一了.
這就是耶穌基督在十字架上.
成就救恩的果後.
你說這個故事.
有沒有告訴我們.
究竟如何拆毀這些牆呢?.
我相信有三個大家可以思考.
第一.
是神主動來作出這個改變.

$^{601}$神去跟哥利樓說話.
神給彼得有意象.
神差派聖靈來引導他們.
如果我們有些牆要拆毀的時候.
兄弟姐妹.
不能夠靠我們自己的能力.
來拆毀這些牆.
是要靠神的介入.
讓一些我們不應該存在的界限.
或者牆可以得到拆毀.
所以弟弟姐妹不要靠自己.
來嘗試一些重大的改變.
我們必須來神面前.
將事情交託給神祈禱.
第二.
這個故事也告訴我們.
拆毀這些牆.
是需要我們挑戰我們固有的觀念.
我們的一些偏見.
偏執.
比喻在這裡.
他很要面對他一生以來.
對潔淨與不潔淨的觀念.
所以對我們來說.
如果我們有一些錯誤的觀念的時候.
我們都需要有一個心智.
來重新檢視我們一些人生態度.
或者人生一些假設.
簡單來說.
弟弟姐妹.
不要以為自己一定是對的.
很多人都是這樣.
我當然是對的.
這個態度是完全將我們封鎖在我們自己固有的框框裡面.
很難突破出來的.
所以我們來到神面前.
我們必須謙卑.
讓神有足夠的空間來去.
模造我們.
改變我們.

$^{641}$更新我們.
第三.
這個故事也告訴我們.
如果我們真的要拆毀我們一些固有的框框的時候.
我們必須要踏出我們的舒適領域.
你想一想彼得.
彼得以前整個人生都是跟猶太人一起生活的.
交往的.
現在他踏出他的舒適圈.
跟一些他不認識.
他感到生疏.
以前感到不潔淨的社群.
來相交.
很大的一個改變.
很大的一個挑戰.
很大的一個踏出他的舒適區.
但在這裡我們看到神是挑戰彼得.
以致彼得能夠踏出他固有的光景.
讓神帶到他屬靈的高處.
在高處裡面.
我們才能夠看得遠.
在高處裡面.
我們才能夠看得到神所賜給我們真實的平安和喜樂.
稍後我請我們低頭祈禱.
讓這個故事來提醒我們.
耶穌是藉著他的死和復活.
已經拆毀了世上所有不應該有.
甚至是帶來人與人之間分裂的牆.
這個故事也讓我們看到.
無論在什麼地方.
如果我們將傳統變成了你的真理的時候.
神求你使我們成為彼得.
願意作出改變.
願意去到你所差遣我們的地方.
願意接受你的教導.
打破自己錯誤的迷執.
以致你可以將我們帶到屬靈的高處.
看得遠.
看得準.
看得清.

$^{681}$被你所使用.
榮耀你的名.
多謝您給予我們一切祈禱 奉耶穌名字而祈禱.
阿們.
\newpage



\section{羅馬書 1:16}
\label{sec:031645}
\textbf{今天的福音}~~~~趙崇健牧師
\newline
\newline
連結: \href{http://www.church.com.hk/acms/content.asp?site=cdc\&op=show\&type=product\&code=031645\&layout=sermon}{\texttt{031645}} ~~~~ 語音日期: 2026-02-15
\newline
\newline
\hyperref[sec:031644]{\small{< < < PREV SERMON < < <}}
~
\hyperref[sec:index_chronic]{\small{[返順時目]}}
~
\hyperref[sec:031646]{\small{> > > NEXT SERMON > > >}}
\newline
\newline
羅馬書 1:16
\newline
\begin{longtable}{cl}
\hline
\hline
章節 & 經文 (和合本修訂版)\\
\hline
1:16 & \begin{tabularx}{0.7\textwidth}{X} 我不以福音為恥；這福音本是神的大能，要救一切相信的，先是猶太人，後是希臘人。 \end{tabularx} \\ \\
[1ex]
\hline
\hline
\end{longtable}
$^{1}$今天的訊息主題就是今天的福音.
現在已經踏入了二月已經有一半了.
如果沒有搞錯的話.
農曆新年就好像快到了.
我都預先想祝大家新年快樂 身體健康.
今年是馬年.
我祝大家農馬精神.
馬到功成.
我聽到有個基督徒跟我說.
他說基督徒今年可以祝大家馬上得福.
馬上得福.
即是你馬上就行了 馬上得福.
我都很記得這句說話.
今日的訊息開始之前.
我就想這樣說.
農曆新年是講福氣的.
就想大家都得福.
但這篇訊息又並非單單是安慰的訊息.
反而是一個提醒我們的訊息.
我就想問大家一個問題.
你有沒有曾經試過.
有意無意之間做了一些出醜的事情.
甚至你覺得真的很醜怪的事情.
你想一想有還是沒有.
喝了多多少少.
你可能會覺得很奇怪.
為什麼牧師今天一來就問這樣的問題.
如果你可能會覺得.
哎呀 我說沒有.
你又不相信 怎麼會沒有.
但我說有 但我又不好意思說.
你都說醜事了 我怎麼會想說呢.
這就是我們的心境.
傳統華人有句說話.
家醜什麼.
不出外傳.
很多事情是怎樣呢.
自己知道就算了.
不要傳出去.
正所謂 關上門打兒子.

$^{41}$有人問起 最近你兒子怎麼樣.
最近我兒子特別乖.
打完當然特別乖.
這個傳統思想我就不知道有沒有過時.
但是今天的世代.
很多事情都不同了.
我曾經在網上見過一句說話.
這句說話叫做.
只要你不怕醜.
怕醜的就是別人.
你明白就明白了.
也有聽過別人說.
現今世界在網上.
甚至你做人無恥就是最什麼.
有人說無敵.
對呀 網上真的說了這些說話.
我不是贊同.
但是現在網上新的一代.
就是覺得無恥是一種能力.
我記得我小時候.
有一次我都會人有三急的.
四處找廁所.
很不容易終於找到.
一推開門.
一看進去.
為什麼沒有尿兜.
看真一點.
為什麼粉紅色的.
小孩子推開的那些.
是什麼顏色的.
男孩子應該是粉藍色的.
接著就可能見到一位女士.
推開廁格的門出來.
小孩子都害怕的.
我就馬上走出去.
如果這些醜事.
今天發生.
你看到我推開後面的女廁.
進去裡面的話.
現在就很大件事了.

$^{81}$我覺得出醜.
醜事.
多多少少人人都會有.
如果沒有人知道.
那就當沒事了.
同意嗎.
但是你想一想.
如果你的醜事.
羞恥的事情.
被人知道了.
不止.
還要你都知道人家知道.
人家都知道你知道他們知道的時候.
那你又會怎樣呢.
今天的經文.
保羅正正就是面對一個這樣的處境.
保羅他知道的.
不少人覺得他很醜怪.
很羞恥.
因為在這一班人的眼中.
保羅所信過的福音.
本身是一件很羞恥的事情.
訊息主題是今天的福音.
我之所以要著重今天.
是因為今天的經文.
其實整段保羅都是用上一個現在式去講的.
保羅他不是著重我們覺得過去福音有多重要.
而是在今天的當中.
今天我們的生命.
今天我們的呼吸.
福音意義是什麼.
我們怎樣看福音.
我們應該怎樣去著重這個福音.
讓我們低頭祈禱.
是因為你的拯救.
是因為有人將福音傳給我們.
我們今天才能夠相信.
我們才有新的生命.
我們願主你幫助引領我們.
叫我們明白我們的福音是從別人所傳而來.

$^{121}$也叫我們明白我們自己將福音傳給別人.
對他們生命的重要.
於是今天我們有所欠.
我們有所做得不足.
願主你教導.
於是今天我們膽怯軟弱.
為了福音我們不願站起來.
求主你堅強.
於是今天我們看輕了福音的寶貴.
求主你赦免原諒.
更重要.
願你的聖靈幫助我們.
看見福音在我們生命當中何等的重要.
主耶穌你是何等重要.
福音,看重福音在我們生命當中.
所表現出來的那種能力.
那種大能.
我主你幫助.
祈禱奉耶穌基督的名叩.
阿們.
經文很短.
保羅第一句就說.
「我不以福音為恥」.
保羅其實是在引用一個我們必須要明白的概念.
榮譽和恥辱.
大家可能在一些《釋經》書.
在《培靈會》都會聽過.
第一世紀是一個榮辱很分明的世界.
不過我想著重的一點.
是其實我們今天的世界同樣如此.
當時的世界.
什麼是榮譽.
有一套很清楚的理解.
簡單來說.
一個人.
他的出身.
他的身份.
本身就是一個榮譽.
出身於顯赫的家族.
是榮譽.

$^{161}$一個人如果有能力.
他又有影響力.
無論你說是武力也好.
你說是知識也好.
當中都是一種榮譽.
因為可以貢獻國家.
貢獻社會.
自然人就有光榮.
更重要.
榮譽在當時的社會.
不是自己說了我有.
就當你是有.
這樣就算了.
是要由別人說你有榮譽的.
舉個例子.
羅馬社會其實有選舉的.
當時的地區選舉裡面.
候選人是不會為自己拉票.
好像現在的政府那樣告訴你.
其實我有多配得去服侍大家.
因為當時的理解是如果你這樣說.
你就沒有榮譽了.
榮譽是要由候選人身邊.
去說回你知不知道.
這個候選人.
他曾經這樣.
如果他做我們的候選人.
那我們就最好.
所以大家要投票給他.
榮譽是要由別人的口.
去說出來.
今天的社會其實都是一樣的.
我們對榮耀都有一套理解的.
譬如每年都有澳洲人.
會有人因為對社會的貢獻.
而去獲得一個叫做.
Australian of the Year的獎項.
很大的榮耀來的.
又或者一個人今天.
你做什麼工作.

$^{201}$你擁有多少財富.
都可以是一種榮耀.
不過這一切的榮耀.
就跟保羅其實都沒有什麼關係.
因為按照當時的標準.
保羅是不會被歸為榮譽那一邊.
相反他是會被歸為恥辱.
羞恥那一邊.
當時的人.
是不會好像現在社會的人.
或者基督徒那樣.
哇 保羅啊.
你為神擺上那麼多.
你真是很愛神的.
你真是有榮譽的.
不會的.
相反.
保羅正正就是因為福音.
他成為一個羞恥的人.
在當時的社會.
對羞恥都有一套理解的.
沒有受過教育.
是羞恥來的.
保羅受過高等教育.
卻不做一個有身份的人應該做的事.
更大的羞恥.
對羅馬這個國家.
對民族忠誠.
如果你不忠誠的話.
你就是羞恥.
是不是.
保羅是羅馬公民來的.
他不是在說.
凱撒是主啊.
他說什麼呢.
耶穌是主啊.
你不忠於羅馬帝國.
但是他享有特權.
更大的羞恥.
當時最大的羞恥是什麼呢.

$^{241}$就是一個人被人全剝了衣服.
只掛著一塊布.
烈日當空.
在十字架下面.
曬著釘著.
慢慢地死.
這個目的是遊街示眾.
保羅傳耶穌基督的十字架.
在當時是極大的羞恥.
極大的荒謬.
所以保羅才會說.
在這裡他這樣說.
「我不以福音為恥」.
為什麼呢.
因為他很清楚.
很多人當他是羞恥.
當他所信的福音為恥.
所以這句話.
其實是保羅宣告.
你們.
你們這群不信的人.
當為羞恥的事.
我現在告訴你.
我們當做是光榮.
這句話是保羅的抗辯.
他在拒絕當時社會對榮譽的定義.
他在拒絕當時普羅大眾.
所有人都追求並且接受的生活方式.
對保羅來說.
福音是榮譽.
按著福音的方式去生活.
就算要承擔社會的羞恥.
都是榮譽.
這就是保羅的想法.
弟兄姊妹.
我相信人到了某些階段.
自然都會想安穩下來.
想少一點衝突.
保羅寫這句話的時候.
他不年輕了.

$^{281}$可能他也會想.
保羅這句話.
即使你這樣說出來.
都會覺得不容易去聽.
今天對你來說.
你覺得什麼是榮譽?.
你覺得什麼是羞恥?.
今天我們追求的是哪一種榮譽?.
我很記得我最初回教會的時候.
我聽過一個牧師.
分享當時教會裡面.
一個弟兄的見證.
我到現在都記得很清楚.
那個弟兄當時是四五十歲.
他做會計師.
他一直在一個大型的會計師樓裡面工作.
大學畢業.
做了二十多年會計.
事業都挺好 挺順利.
職位都做得不錯.
有一次.
他負責幫一間公司去核數.
上頭就暗示他.
會不會明示暗示.
這次要你幫忙.
出一點點狡猾.
這個弟兄就知道這樣是不老實的.
他就回去和牧師談.
牧師應該怎麼做.
大家都明白你們都有工作.
社會就是這樣.
你肯幫忙.
大家就是什麼?.
自己人.
我們坐在一條船.
我們同舟共濟.
將來有什麼事一起幫忙.
如果你不肯幫忙.
甚至你踢爆我.
那就不是自己人了.

$^{321}$不是敵人的.
不過肯定不是朋友.
那你怎麼辦?.
後果是怎樣?.
心中有數.
普遍的信徒可能會以為是這樣.
我們一般來說.
我們都有一種孩童的天真.
神會為他開路的.
我們祈禱.
神求你讓他避過這個難關.
沒事的.
又或者即使這個弟兄拒絕.
神都會保守他一個好結果.
他不會離職的.
他老闆不知道什麼原因被解僱了.
沒事的.
好像迪士尼樂園一樣.
永遠都發發落落地生活下去.
我們想這樣.
你想不想?.
想.
這個見證的結局就是他沒有妥協.
到了最後.
很多壓力了.
他受不了.
他自己辭了一份工作.
我記得牧師有個主人說這句信息.
當時我很年輕.
雖然我覺得我現在都很年輕.
當時的我.
我覺得這個見證我很深印象.
剛剛信主嘛.
是不是?.
我心裡面想.
福音不是好消息嗎?.
如果福音是這麼苦.
福音應該改錯名.
福音應該叫做苦音.
都是f字頭.

$^{361}$不過音不同.
今天的世界.
其實同樣.
你和我活在一套榮辱的系統裡面.
什麼是成功.
什麼是失敗.
有什麼值得羨慕.
社會有一套答案.
老實說.
我覺得我們今天都很幸福.
起碼在這個澳洲社會.
我們的社會.
不會.
因為.
我們是基督徒.
就極端地排斥你.
給你面色看.
甚至覺得.
我們這些人是羞恥.
某程度上.
我們承受的壓力.
少了很多.
這樣就幸福了.
少一點壓力.
就福音了.
是不是?.
但正正因為這樣.
有時我們付出了另外一個代價.
弟兄姊妹.
我想我們忽略了.
就是我們慢慢慢慢再分不出.
世界那種的福.
和福音那種的福.
有沒有不同呢?.
我們或許以為.
福音的福.
就必定帶來世界的福.
福音的榮譽.
就必定會換來.
世界的榮譽.

$^{401}$最好了.
如果你不換都換回七八成給我.
我就要講清楚.
我不是否定.
福音能夠為大家.
帶來世上的福氣.
我恨不得為大家祈禱.
我們一同在主的裡面.
我們身體健康.
我們有主給我們的福氣.
但今天的經文只是提醒了我們.
對保羅來說.
福音.
原來.
是可以被這個世界.
視為羞恥的事.
我們當然渴望神賜福給我們.
但如果有一天.
福音.
不再等同世界的成功.
世界的榮譽.
我們還願意去持這個信仰嗎?.
我們會否為福音.
說一句抗辯的話呢?.
保羅這一句.
「我不以福音為恥」.
其實對我們來說.
是一個很需要的提醒.
其實我們要預備我們的心態.
有時候.
或者為福音.
成為別人眼中如住的人.
成為別人眼中的笑柄.
其實是值得的.
是榮譽來的.
因為好像保羅所說.
其實我們不以福音為恥.
我們願意承擔這個福音.
好嗎? 弟兄姊妹.
接著保羅這樣說.

$^{441}$「這福音本是神的大能.
要救一切相信的.
先是猶太人.
後是希利尼人」.
這裡說到大能.
當時大能.
就是軍事大能.
知識大能.
財富的大能.
保羅這裡所說的大能.
就是主耶穌基督十字架受死復活的大能.
而這個十字架正正就是當時的人認為.
最沒有大能的事.
保羅說他是大能.
別人說是羞恥的事.
他說這個是榮譽.
這裡我覺得放在今天.
彷彿那個衝擊力就是.
其實做賊是這個世界上最好的事一樣.
當時的人聽下去.
他們覺得是這樣.
做賊那個是最好.
但保羅說不是吧.
這個羞恥極荒唐所謂的大能.
保羅說要拯救世上一切的能.
一切的人.
無論是猶太人還是希利尼人.
希利尼人不是民族.
意思是說希臘話的人.
即是說外語的人.
猶太人很簡單.
這個世界只有兩種人.
一種是猶太人.
另一種就是非猶太人.
混在一起就可以了.
不要緊的.
所以猶太人和希利尼人混在一起.
即是所有的人.
這個大能保羅說.
要拯救世上所有的人.

$^{481}$拯救Delivered.
拯救世界.
或者你可以翻譯成解放世界.
是當時羅馬這個國家.
靠力量去宣揚的事.
他們一直都相信.
只要羅馬所到之地.
那個地方就有永遠的福樂.
但是保羅卻說福音是十字架.
他會反而去拯救所有人.
你可能會懷疑.
你還懷疑嗎?.
我今天問你.
福音是不是真的拯救所有人?.
即是你最憎恨的那個.
福音也拯救.
我們會不會懷疑.
保羅說拯救所有相信的人.
其實極度荒唐.
並且你要留意這節經文.
不是說我們應該怎樣傳福音.
而是說福音的本質是什麼.
我聽過一篇訊息這樣說.
神其實很有大能.
何必要我們傳福音呢?.
很容易的.
如果神要人相信.
神只需要把這個世界.
所有的人一分為二.
基督徒和非基督徒.
基督徒是不會生病的.
我們連肚子痛都不會.
非基督徒他們會生病的.
很容易會肚瀉.
統計每一年一出.
所有的統計都分兩類.
基督徒病的比率是多少?.
零.
非基督徒病的比率是多少?.
普通人這樣.

$^{521}$只要你這麼簡單.
一分一分神有沒有能力有.
我們不用傳了.
教會外面今天排隊排到圖書館那裡.
但神就是不是這樣.
今天我們知道傳福音不容易.
我們想想為什麼呢?.
原因在哪裡呢?.
經文裡面提出另一個概念.
因為神的義正在這福音上顯明出來.
除了榮譽和恥辱.
保羅突然間跳出一個很重要的概念.
就是義與不義這個概念.
經文在提醒我們.
一個顛覆的現實.
小心聽.
現今世代以為羞恥的事.
在神眼中可以是義.
而這個世界上以為是榮譽的事.
在神眼中可以是不義.
大家或許就應該問.
有沒有一些例子.
世上的榮耀卻是神眼中的不義.
有的 商業世界很容易想得到.
很明顯.
如果一間公司的文化是肥上瘦下.
長期剝削.
把人當成一顆螺絲.
用完就即棄.
公司很成功的.
很有榮譽的.
股價不斷上升.
但是就是不義.
華人有句說話你懂得很有趣.
這句說話叫做發財甚麼.
立品 發財立品.
這句說話很有味道.
要深入一點.
不是立品發財.
是先發財 後立品.

$^{561}$那我就問.
我們發財之前.
是否我們可以一起沒有品.
是否我們為了發財.
就可以繼續沒有品.
發了財 才立品.
是不是呢.
剛才我問.
為什麼今天傳福音不容易.
你想想.
這個長期剝削別人.
肥上瘦下的上司.
跟他的同事說.
其實神真的很愛你.
今天OT我求神加力給你.
福音要顯出神的義.
這種顯明.
其實不只是我們內心.
更加是一種廣義的覆蓋.
家庭 朋友 商業文化 政府.
甚至乎我們祈禱也有.
今天是祈禱事項.
國家方向.
我們可不可以在個人上.
不把人當成一粒螺絲.
我們基督徒相信.
永生是信主那一刻就開始.
同樣的道理.
福音所顯明的那種神的義.
其實也應該在世上多少顯明出來.
拉丁姐妹在這個世界.
能夠顯明福音.
顯明神的義.
有什麼人呢?.
沒有其他人了.
不會是佛教徒.
不會是回教徒.
不會是其他信仰的人.
只有我們基督徒.
唯有是基督徒.

$^{601}$雖然我們都是罪人.
我們有不義,我們有軟弱.
但如果我們不願意福音為恥.
我們願意馴服福音多少少.
或者我們可以學習一下.
調整一下我們的生活.
在各種的地方.
顯出神的義多少少.
我看得出我們當中不少人.
都是作父母.
我們自己也是.
我自己也是.
作父母總會想兒女身體健康.
生活順利.
有一條安穩的路去走.
希望他們讀好書.
有好的工作.
照顧好自己.
照顧好家庭.
開枝散葉.
生一隊足球隊更好.
這些事情都是很正常的.
但當然下一代走哪條路.
有時又輪不到自己父母出聲.
我兒子今年都十四歲了.
慢慢都輪不到我出聲了.
有一次在香港.
我和神學院裡面的一班傳道人.
有團契.
不知為何我們聊著聊著.
談了一條問題.
我們說大家想不想.
大家的兒女都做傳道人呢?.
我當時沒有馬上回答.
我心裡就想.
如果神父照顧他們.
我會祝福的.
但我都不想催逼.
又不想引導.
我又不想試探他們.

$^{641}$更加沒有辦法去操控.
我個人覺得.
事奉創天造地的永生神.
是世上最美好的工作.
比大家好.
但要留意這是我的個人想法.
大家未必一定要同意.
但也不需要反對.
有些事情聽了我們就算了.
當時有一位六十歲的傳道人.
他這樣說.
我不希望自己的兒女.
像我這樣做傳道人.
因為好像社會的收入,地位,身份.
好像沒有那麼吸引.
你也知道現在美國.
傳道人是其中一個最不受信任的職位.
他寧願他的兒女去做信徒.
有份穩定的工作.
好好的.
首先.
我想我們千萬不要.
將這段說話聽成負面.
這段說話的出發點.
真的可以純粹是一個父母.
對兒女的關心.
或者擔心的矛盾心情.
但是弟兄姊妹.
我想問一下.
如果我們的兒女.
深孝一些很重視的人.
要事奉神.
好像Cindy那樣.
去校園做宣教士.
我們心裡有沒有什麼顧慮呢?.
我們會不會因為收入,地位.
我們反對呢?.
或者現在個人上.
信仰來說.
有人傷害過我們.

$^{681}$我們會不會像以前信主那樣.
我們會生氣.
我們想報復呢?.
再舉一個實際的例子.
很簡單.
每天很多事我們都遇到.
在雪梨我們經常開車.
你看到前面那條小路.
有一個人等了很久.
他的樣子已經很可憐.
他的車有一個牌.
黃色上面有個英文字母.
L.
他等了很久.
沒有人讓他.
你會讓他.
還是你覺得.
你當然出不了.
這條路我正在走.
你慢慢等吧.
明天到你應該出道了.
我們會不會這樣還是我們讓他呢?.
這個世界鼓吹的是什麼呢?.
不停的消費.
買多一點.
未必有用.
買.
開心.
不斷的剝削.
為什麼剝削呢?.
自己有多一點.
別人有少一點.
他會搞定自己.
我有多一點.
我要搞定我自己.
不停的去購買.
消費.
不停的浪費.
用盡各種的事.
去塞滿 塞爆我們的生命.

$^{721}$就是現代的生活方式.
但是其實.
簡樸生活.
本身是基督徒的美德.
我們今天願意為福音.
為了神的義.
我們願意調整生活到什麼地步.
這一句神的義.
在這福音上要顯明出來.
正如保羅提醒我們.
福音不是內在.
是要調整我們的生活.
去顯明出來.
我們的生命.
是因為有福音而不同.
馬太兄姊妹.
剛才這麼多的分享.
我們還是在說.
經文上的反省.
不過問題就來了.
每一次聽完反省.
我們會問.
實際來說.
在生活上.
我們應該怎樣去實踐呢?.
回顧保羅.
他出身其實不差.
他受過高等教育.
老師是加瑪列名師.
他家境不錯.
家境做一些生意.
做跡象棚.
所以他也有一門手藝.
就是跡象棚.
他本身是一個有榮譽.
有地位的人.
今天的經文其實是.
保羅的自我獨白.
他在回顧自己的出生經歷.
反思那種從被看為榮譽.

$^{761}$去到承擔羞恥的轉變.
並且他為福音去抗辯.
我想帶大家去看保羅三句說話.
就是他這種轉變是怎樣發生的呢?.
主耶穌在他生命當中做過些什麼.
令他變成這樣的人.
保羅怎樣去理解耶穌.
令他願意這樣放下.
這樣去轉變.
三段經文.
畫面沒有的.
大家要聽我說.
第一是虛己.
《希臘比書》二章六節.
保羅這樣說.
「他本有神的形象反倒虛己.
取了奴僕的形象.
成為人的樣式」.
我聽過很熟悉.
虛己我想說是一種心態的調整.
耶穌回顧他.
本身他很高的.
他卻是將自己放下.
這就是耶穌怎樣看待自己.
你和我想一想.
耶穌是神來到世上.
如果耶穌一直覺得自己比別人高一等.
他樂高過頭望別人.
你覺得身邊的人.
知不知道耶穌看低自己?.
他們會否知道.
耶穌你當然覺得自己高.
我們這些是低的.
如果是這樣的話.
耶穌和身邊的人.
會有多大的距離感.
但耶穌不是這樣.
耶穌就是願意放下自己.
先去想別人.
將焦點放在別人那裡去虛己.

$^{801}$不是要我們否定自己.
只不過是要我們避免過分去計算.
自己應得些什麼.
或者自己賠得些什麼.
這是第二件事情.
之前我也想邀請大家.
我想邀請你思想一下.
你今天有沒有什麼執著?.
你正在做的事.
你覺得你需要學習虛己.
放下你自己的事.
想一想吧.
第二是「施予」.
《使徒行傳》二十章三十五節.
保羅不認識耶穌.
但他聽過耶穌在別人口中所說的說話.
保羅說要紀念主耶穌基督的話.
施比受.
接著是什麼?.
更為有福.
施予就是我們對別人的一種調整.
按照世界計算的方式.
施予就等於我所擁有的減少了.
就等於我失去了.
但我想說.
施予不是付出.
施予不只是慈善.
施予是透過讓別人得到更多.
讓自己都進入更大的祝福之中.
所以施予不是少了.
而是施那個和受那個都多了.
如果不是這樣理解.
這一節經文不是很明白.
每次施予都是損失.
不是.
每次施予我們自己都多了.
今天我就邀請大家去思想一個.
你覺得可以施予的人.
有什麼人你可以支持.
去思想一樣你可以施予的事.

$^{841}$不用很大.
只要你的能力能夠做到就行了.
你想一想.
最後獻予.
Hand over.
已不鎖書五章二節.
基督愛我們.
為我們捨了自己.
當作傾向的供物和祭物獻予神.
耶穌就將自己獻給了神.
這是福音.
我就想你思想.
你在生活上面.
有什麼事情有什麼習慣.
你可以調整一下.
甚至你寫它.
有時候獻予不是給多一點.
或者你做多一點.
獻上也可以自己找少一點.
捉少一點.
放低一點.
英文有現在式有過去式.
今天的經文提醒我們.
福音給我們這種生命.
是不會是過去式.
永遠都是今天的現在式.
二人必因順得生.
是現在式.
這個新生命是現在式.
影響我們的生活.
塑造今天的我們.
影響我們的家庭和社會.
保羅到底為什麼不以福音為恥呢?.
因為他真的覺得福音很寶貴.
我就問我自己.
或者大家也可以問.
我覺得神的愛價值多少?.
有個客觀答案.
十字架和救恩是很寶貴的.
主耶穌基督在十字架所流的血和血心.

$^{881}$是最貴重的事.
客觀的答案.
主觀呢?.
有沒有一些主觀的經歷.
個人的.
告訴我們.
我們自己覺得.
主耶穌基督的愛價值多少?.
有沒有?.
我想到一個答案給我自己提醒.
或者當我們願意為神的愛擺上多少.
就正好引證.
神的愛對我們來說價值多少.
我重複這句話.
或者我們願意為神的愛擺上多少.
就正好引證.
神的愛對我們來說價值多少.
今天的福音不是要我們改變世界.
但是我相信最起碼.
是要求我們.
不以這個世界的榮譽去定義我們自己.
也提醒我們.
不是單單去追求世界的榮譽.
讓我們低頭祈禱.
我們每次都記得主你的愛何等寶貴.
但是父神要示我們尋求一個答案.
我求主你的聖靈光照我們.
我們看得出.
不單單只是客觀的證據去說出十字架.
而是一種刻在心上.
刻在我們的心上.
記在我們的腦海當中的一個價值.
好叫我們思想主你的救恩的時候.
我們不是停留在頭腦當中.
更加是去到我們的生命.
我們的情感當中.
好能夠叫我們被主你的聖靈所催促.
去奉獻更多給主耶穌基督.
不是單單是我們所擁有.
更加是我們自己.

$^{921}$父神今天要示有很多事情.
我們不知道如何去照顧.
我們擺在主你的面前.
要示今天我們需要放手的事.
原諒的事,原諒的人.
求主你教導.
願你的聖靈在我們生命當中.
將不可能的事轉為可能.
哪怕是我們在主你的當中.
要學習放下我們自己.
在我們來說何等困難.
何等不可能.
但是主你說過.
在神凡事都能.
給我們有新的生命.
願以我們每一天起來.
我們有新的氣息.
我們更加思想得到.
其實每一天都是主你在我們生命當中.
顯出大能的一天.
我們祈禱奉耶穌你的名求.
阿們.
\newpage



\section{希伯來書 10:19-39}
\label{sec:031646}
\textbf{你的信仰是否愈信愈後退?}~~~~余大峰傳道
\newline
\newline
連結: \href{http://www.church.com.hk/acms/content.asp?site=cdc\&op=show\&type=product\&code=031646\&layout=sermon}{\texttt{031646}} ~~~~ 語音日期: 2026-02-15
\newline
\newline
\hyperref[sec:031645]{\small{< < < PREV SERMON < < <}}
~
\hyperref[sec:index_chronic]{\small{[返順時目]}}
~
\hyperref[sec:031647]{\small{> > > NEXT SERMON > > >}}
\newline
\newline
希伯來書 10:19-39
\newline
\begin{longtable}{cl}
\hline
\hline
章節 & 經文 (和合本修訂版)\\
\hline
10:19 & \begin{tabularx}{0.7\textwidth}{X} 所以，弟兄們，既然我們靠著耶穌的血得以坦然進入至聖所， \end{tabularx} \\ \\ \relax
10:20 & \begin{tabularx}{0.7\textwidth}{X} 是藉著他給我們開了一條又新又活的路，從幔子經過，這幔子就是他的身體。 \end{tabularx} \\ \\ \relax
10:21 & \begin{tabularx}{0.7\textwidth}{X} 既然我們有一位偉大祭司治理神的家， \end{tabularx} \\ \\ \relax
10:22 & \begin{tabularx}{0.7\textwidth}{X} 那麼，我們該用誠心和充足的信心，同已蒙潔淨、無虧的良心，和清水洗淨了的身體來親近神。 \end{tabularx} \\ \\ \relax
10:23 & \begin{tabularx}{0.7\textwidth}{X} 我們要堅守所宣認的指望，毫不動搖，因為應許我們的那位是信實的。 \end{tabularx} \\ \\ \relax
10:24 & \begin{tabularx}{0.7\textwidth}{X} 我們要彼此相顧，激發愛心，勉勵行善； \end{tabularx} \\ \\ \relax
10:25 & \begin{tabularx}{0.7\textwidth}{X} 不可停止聚會，好像那些停止慣了的人，倒要彼此勸勉，既然知道那日子臨近，就更當如此。 \end{tabularx} \\ \\ \relax
10:26 & \begin{tabularx}{0.7\textwidth}{X} 如果我們領受真理的知識以後仍故意犯罪，就不再有贖罪的祭物， \end{tabularx} \\ \\ \relax
10:27 & \begin{tabularx}{0.7\textwidth}{X} 惟有戰戰兢兢等候審判和那將吞滅眾敵人的烈火了。 \end{tabularx} \\ \\ \relax
10:28 & \begin{tabularx}{0.7\textwidth}{X} 任何人干犯摩西的律法，憑兩個或三個證人，尚且必須處死，不得寬赦， \end{tabularx} \\ \\ \relax
10:29 & \begin{tabularx}{0.7\textwidth}{X} 更何況踐踏神兒子的人，他們將那使他成聖之約的血當作不潔淨，又褻慢施恩的聖靈的人，你們想，他不該受更嚴厲的懲罰嗎？ \end{tabularx} \\ \\ \relax
10:30 & \begin{tabularx}{0.7\textwidth}{X} 因為我們知道誰說：「伸冤在我，我必報應。」又說：「主要審判他的百姓。」 \end{tabularx} \\ \\ \relax
10:31 & \begin{tabularx}{0.7\textwidth}{X} 落在永生神的手裡真是可怕呀！ \end{tabularx} \\ \\ \relax
10:32 & \begin{tabularx}{0.7\textwidth}{X} 你們要追念往日；你們蒙了光照以後，忍受了許多痛苦的掙扎： \end{tabularx} \\ \\ \relax
10:33 & \begin{tabularx}{0.7\textwidth}{X} 一面在眾人面前公然被毀謗，遭患難；一面陪伴那些受這樣苦難的人。 \end{tabularx} \\ \\ \relax
10:34 & \begin{tabularx}{0.7\textwidth}{X} 你們同情那些遭監禁的人，也欣然忍受你們的家業被人搶去，因為你們知道自己有更美好更長存的家業。 \end{tabularx} \\ \\ \relax
10:35 & \begin{tabularx}{0.7\textwidth}{X} 所以，不可丟棄你們無懼的心，存這樣的心必得大賞賜。 \end{tabularx} \\ \\ \relax
10:36 & \begin{tabularx}{0.7\textwidth}{X} 你們必須忍耐，使你們行完了神的旨意，可以獲得所應許的。 \end{tabularx} \\ \\ \relax
10:37 & \begin{tabularx}{0.7\textwidth}{X} 因為「還有一點點時候，那要來的就來，必不遲延。 \end{tabularx} \\ \\ \relax
10:38 & \begin{tabularx}{0.7\textwidth}{X} 只是我的義人必因信得生；他若退縮，我心就不喜歡他。」 \end{tabularx} \\ \\ \relax
10:39 & \begin{tabularx}{0.7\textwidth}{X} 我們卻不是退縮以致沉淪的那等人，而是有信心以致得生命的人。 \end{tabularx} \\ \\
[1ex]
\hline
\hline
\end{longtable}
$^{1}$弟兄姊妹 主佬平安.
你的信仰會不會越信越後退.
這個主題會不會覺得有些詫異.
有沒有想過大家.
最近大家可能都.
如果有在返九一內的話.
或者有看一些屬靈書籍的話.
可能對Philip Yancey是不陌生的.
但若然你沒有看過.
就容許我介紹Philip Yancey是一個什麼人.
他21歲就已經是全世界的事奉.
在大學當中他已經是寫稿的.
在這些campus的live magazine開始寫稿.
他有50多年的信主和事奉的生涯.
出了25本書 全球賣1500本.
去visit過60多個國家.
有40個語言的譯本.
作者當中可能一本書沒有寫過.
基督教的書.
他寫什麼書特別是怎樣去鼓勵弟兄姊妹.
面對痛苦 面對艱難的時候.
怎能夠去持守信仰.
但是他76歲那一年.
他竟然自己的去透露.
原來他是有第三者.
和一個有夫之婦.
有這個姦情八年.
隱瞞了他家人.
他這樣去透露之後.
他在主面前認錯 在人面前認錯.
他說以後都不再寫作了.
以後都不再check post了.
他的書曾幾何時真是感動很多人.
我家裡也有.
很多在座的都看過.
不難去想.
為什麼他信主這麼多年.
事奉了這麼多.
也這麼受人敬重.
這麼多這麼深的屬靈的經歷.

$^{41}$他寫的每句話是他自己的個人經歷.
不是抄襲別人 不是說其他.
是說他自己的經歷 自己的領受.
但是都跌倒了.
會不會問你和我.
會不會有這樣的可能性.
為什麼會越信越退後.
我想一個好的形容詞就是.
離婚的人.
你結婚那一天.
那剎那做完宣誓之後.
接著那一秒你應該不打算離婚了.
之後的機會少一點.
但是當你的婚姻越走越久的時候.
可能用很多的挑戰.
越來越有機會想離婚.
或者就算不離婚也好.
相對都無疑.
感情都越來越冷淡.
沒有了實體.
有外表 飯都照煮.
但是已經沒有了實體.
今天這個訊息.
特別是我會用新譯本.
因為剛才多謝主席用了和本.
所以兩個譯本大家看了是非常好的.
特別是這兩句話.
想大家深入思考一下.
因為我們太多神學思想.
在我們腦海裡我們去記住.
但是這些思想對我們沒有幫助.
是不是對我們生活有幫助.
我的義人必因信得生.
現在在說的是什麼人.
是我的義人必因信.
因信菁兒.
既然我已經信了.
我就得道永生.
這句大家同意的.
如果他後退.

$^{81}$他是說誰.
未得救的人.
未信的人.
還是信了的義人.
是因信而稱義的義人.
因信而稱義的義人.
如果他.
所以根本整個希伯來書.
不是和未信的人.
是和信了.
特別是說和信了的猶太背景的基督徒.
信了他們已經.
他也是因信.
不是因為他做舊約那些手行為.
憲制那些.
而是因為耶穌基督的保權.
他相信他.
信了.
但是.
他受到猶太教.
又受到羅馬的大逼迫.
一會兒會見到怎樣逼迫.
我剛才已經說了怎樣逼迫.
他的財物被人搶.
被人抓去坐牢.
在大逼迫之後.
之後他後退.
後退是什麼意思呢.
他本來是在這裡的.
現在後退.
本來和太太的感情是很好的.
牽著手的.
昨天情人節牽著手的.
到後來已經吵架了.
大家睡覺都是這樣分開睡.
分房睡.
那種的情況.
他.
即是在說誰.
在座如果是信耶穌的人.

$^{121}$就是在說你和我.
我是因信稱義的位置.
就是在說你.
就是在說我.
我的心.
如果他後退.
我的心就不起穴.
神的心就不起穴.
但是我們不是那些後退.
以致滅亡.
後退.
原來是不是後退又立刻滅亡呢.
如果後退立刻滅亡.
那就要小心了.
可能我們隨時後退一步.
就已經滅亡了.
立刻死光了.
但是後退.
以致滅亡.
換句話說可能後退後退後退.
接著又會滅亡.
好像我上次的希伯來書.
我講了三堂.
今天是最後一堂.
我希望希伯來書這麼複雜的書.
這麼長的書.
叫做在普通書信最長的書.
我嘗試用三堂講那個精髓.
上一次是講什麼呢.
就是說我們可不可以信呢.
信耶穌.
我們沒有成長.
不追求成長可不可以.
如果忘記了就回去聽一下.
不進就不退.
我和太太的關係.
我不追求成長了.
我不追求愛她了.
都愛了二十年了還要愛她.
那你就已經後退了.

$^{161}$你講這句話這樣想就已經後退.
但是我們那些不是後退以致滅亡的人.
反過來.
我們是有信心.
以致保存生命的人.
在和盤就說.
後退以致滅亡.
他用沉淪.
保存生命.
他用靈魂得救.
我靈魂不是得救了嗎.
我信耶穌的時候已經得救了.
但是如果我後退.
我就會滅亡.
我如果有信心不後退.
我就得救.
那就是說.
後退的相反詞是什麼.
相反詞是什麼.
要有信心.
滅亡的相反詞.
就是得生命.
換句話.
基督徒.
你一是後退.
你一是保持信心.
繼續向前走.
你後退可能會致滅亡.
你有信心.
就會保住你的生命.
保住你得救的生命.
希望萊斯的信息.
當中你說不是.
我覺得一次得救永遠得救.
我講了很多次.
一次得救永遠得救這個真理.
是從神的角度.
祂保住你的.
祂當然救到你.
但問題是你不是神.

$^{201}$你怎麼知道你得不得救.
聖經說.
你知道.
你持守到底.
你就知道.
你不持守到底.
你自己都知道不是.
不要騙自己.
是不是.
就是講持守到底.
你的婚姻.
你的愛情.
你愛到底.
你就是持守你的愛情.
你不到底.
你就是沒有持守愛情.
這麼簡單的道理.
你不用講什麼神學.
所以同時間是對的.
一次得救永遠得救.
神保住你是對的.
但問題是你不是神.
你不知道.
但神說你要知道的.
就是你要持守.
所以希望萊斯有五個.
是對原來基督徒是可以.
真心信耶穌的基督徒的人.
是可以後退.
以致滅亡的.
有五個警告.
但之前我講過了.
我不再重複講了.
第一個警告是講到.
怎麼形容呢.
就是說你要聽.
真的要再聽.
每次聽到都要聽.
你不要在這裡睡覺.
因為你真的要聽神的道理.

$^{241}$聖經的教導.
你要聽.
免得你垂留失去.
如果你這樣垂留的時候.
這麼大的救恩.
你怎麼逃脫呢.
講過了.
之後第三章他再講.
他說你不要硬著心.
你要堅持到底.
否則你會跌倒的.
最終你會膠著的.
接著他再講第六章.
之前已經講過了.
就是什麼.
就是說.
因為那些曾經蒙了光照的人.
今天也講過了.
剛才那個經文也講過.
蒙了光照.
常過主恩的那些.
跟聖靈都有份的.
也就是常過神美善惡道的人.
如果他偏離正道.
偏離正道.
即是他本身是不是正道.
偏離正道.
他本身是不是正道.
他在正道.
現在偏離了.
退後.
即是他本身是前.
現在是後.
有個前後.
有個正現在是退.
不是跟本身是斜的那些人.
是跟本身正的那些人.
現在是退了.
換句話說.
我是講信徒.

$^{281}$是很清楚不過.
你這樣再做下去的時候.
如果你這樣經歷過這些事情.
你這樣繼續偏離正道的時候.
偏離不是一次過偏離.
倒退不是一次過倒退.
是繼續不斷地偏離.
繼續不斷地倒退的時候.
你就不可能再有悔改的機會.
你就好像將耶穌基督再釘十字架.
再公然侮辱他.
而他給的圖畫是什麼呢.
就好像這個地方.
這個泥土.
很多雨供應他.
常過這些恩怨.
你應該獎.
獎果實.
你沒有獎.
反而是上這種荊棘.
結果是什麼呢.
是會燒你的.
這個是一個好的比喻.
這三個都講了上次.
這次是講第三段希伯來書.
就是第十章到第十三章.
再繼續警告.
如果我們領受了真理知識以後.
還是故意持續犯罪.
就沒有留下贖罪的制.
沒有留下贖罪的制.
整個希伯來書.
就是講耶穌基督一次的制.
就已經是完成了.
他已經完成了救贖的工作.
潔淨了我們.
那為什麼不留下制呢.
是不是有些衝突.
他完成了.
過去現在將來.

$^{321}$都已經卸下他的罪了.
那為什麼沒有留下制給我們呢.
他不就是我們的制.
他正正就想講.
他做了他的事.
你沒有做你的事.
你的事.
就是你要領受這個恩典.
他這個一次過的制的恩典.
你要去領受才行.
你放棄了他.
你輾踏了他.
你當他沒有了.
你忽略了他.
那還有哪個制可以給你呢.
贖罪制給你.
是這樣的意思.
到第十二章.
他再加警告.
不要棄絕那位說話的.
因為從前那些人棄絕那位.
在地上境界就先知.
都不能夠逃罪了.
在舊約.
到今天你棄絕.
從天而來.
道成肉身那位耶穌基督的教導.
你就更加不可以逃絕.
所以很明顯不過這五個警告.
五個警告是什麼呢.
就是基督徒千萬不可以後退.
接著你的問題就會來了.
剛才有個經文.
如果我不後退.
我怕時間不夠.
因為經常有很多人說.
你時間不夠.
你又說得很快.
所以最精果的說不到.
現在先講這個答案.

$^{361}$講完這麼久.
有五個警告.
不要後退了.
我怎樣可以不後退.
兩個答案.
不要丟棄坦然無懼的心.
坦然無懼.
你看我是出自新漢語.
如果你看喜伯來斯新漢語.
你發覺全聯英文.
多次出現坦然無懼這個字.
即是我們要帶著一個坦然無懼的心.
心態要坦然無懼.
第一樣東西.
第二樣東西要忍.
記住這兩樣東西.
坦然無懼和忍.
坦然無懼的心.
第二樣東西要忍.
一會兒你會看到是甚麼意思.
再講剛才那裡.
經文講得很清楚.
因為我們知道.
誰說過.
身冤在我.
我必報應.
主必定審判他自己的子民.
落在永無一人手裡.
真可怕.
是講神審判他自己的子民.
是多麼慘.
是一個報應.
是講信徒如果你偏離正道.
你後退.
你最終滅亡.
是多麼慘.
後退.
在講基督徒在哪方面.
這裡所講的後退.
是哪方面的後退呢.

$^{401}$大家想不想到.
好像有婚姻一樣.
婚姻的後退.
你可能說婚姻後退很明顯.
就有第三者.
這個有明顯的後退.
剛才還有甚麼後退.
你心裡面已經沒有對方.
你不尊重他.
你不體恤他.
你不聆聽他.
你不扶持他.
我那天看到電影.
我覺得非常之好.
他講一句說話.
I like you because.
but I love you despite.
我因為這些愛你.
但是哪怕那些.
我仍然愛你.
前者是我愛你.
我喜歡你的優點.
後者是哪怕你的缺點.
我都仍然愛你.
對著別人的缺點是難受的.
是在傷害你的.
但我仍然愛你.
我不會離棄你.
我不會倒退.
我不會偏離.
因為這是一個盟誓.
原來人和神的關係都一樣.
而在經文來說.
倒退是講甚麼.
信任哪方面呢.
倒退不是講你有沒有教會在事奉.
你在教會事奉就叫做行前.
沒有事奉就叫做倒退.
你事奉就叫做正路.
不事奉就叫做偏行正路.

$^{441}$不是.
在聖經來說.
大家都知道倒退是甚麼意思.
倒退指哪方面呢.
其實只有一方面.
倒退就是我們經常祝福的時候.
你會講這句說話.
其實祝福這句說話很有意思.
是提醒我們.
其實我們基督徒的生命有甚麼目的.
目的是甚麼呢.
就是願賜平安的神.
在我們那位貧著永血.
我時間已經不夠了.
在一切善事上要完備你們.
做甚麼.
有目的的救恩.
有目的的願賜給你們.
好像你們是遵行祂的旨意.
行祂所喜悅的事.
行祂的旨意就是行祂所喜悅的事.
兩件同一樣的事.
所以甚麼是行神的旨意.
甚麼是神喜悅的事呢.
其實羅馬十二章一節講得很清楚.
大家都會唸.
所以我不講了.
就是基本上將身體獻上當作活祭.
換句說話.
基督徒上了耶穌之後.
整個希伯來書就講耶穌就是那位祭司.
基西德的祭司.
他獻的那個祭都是最偉大的祭.
因為一次過獻完之後.
你以為整本書是講他.
當然是講耶穌基督.
透過他得到他的幫助.
但是再想清楚.
咦.
原來我們上耶穌之後.

$^{481}$我們都要獻自己當活祭的.
我們會成為祭品的.
因為耶穌自己成為祭品.
祂是我們的兄長所謂.
我們都是神的兒子.
我們要跟著祂.
祂走過的路我們要跟著祂.
我們軟弱祂都經歷過.
但是祂沒有跌倒.
祂走過的路我們都要跟著祂.
祂在營外面受苦.
我們都要在營外面受苦.
希伯來書講了這件事.
換句話說.
希伯來書到最後的時候.
十字十三章.
他用了祭來形容我們的生活.
你當自己的生活是一個祭.
你的心境是一個祭.
而他就說在五方面.
你可以作行神的旨意.
做神喜悅的事.
在這五方面.
讓大家在五方面要持守.
大家去看看.
我包括自己去看一下.
我當這個準備查經的時候自己去看一下.
在信心,在愛心,在聖潔,在盼望.
和在坦然無懼.
去到主面前求祂的幫助.
在這五方面我們有沒有持守.
你可能說我信心持守.
愛心持守.
聖潔持守.
盼望持守.
但是其實我是靠自己的.
我沒有靠耶穌的.
我靠自己就已經做到了.
我這麼行.
我用不著靠主.

$^{521}$我信祂的實立要靠主.
得永生要靠祂主.
但是其他的東西我不用靠祂.
我靠自己的了.
那你就沒有持守到.
你就已經後退了.
因為祂已經不是你生命的主.
祂救了你.
但是你生命不需要靠祂.
好這五點.
我沒有還有一講.
很快只剩下十分鐘時間.
第一點是堅持信仰的到底.
當其時猶太人是被逼迫.
是被外面的人逼迫.
所以他是承受很多代價.
但是經文想講的.
就是你不要因為信仰的代價逼迫.
被人的排擠.
或者你受的苦.
或者你覺得已經沒有力量再去承受了.
我走回頭路.
去回猶太教.
經文就是這樣講.
你應該要堅持你所宣認的盼望.
要有完備的信.
真誠的心.
你要繼續以你的頌讚.
就是你宣揚耶穌基督.
在你生命的作為.
繼續以這件事為祭品.
是祭.
可以用什麼祭來形容.
就是你的稱讚是一個祭.
頂多今天唱詩.
麻煩唱大聲一點.
因為唱詩就是一個祭.
你唱一大聲就是一個祭.
因為我走音.
越走音越唱.

$^{561}$因為神喜歡你走音的那種唱.
是不喜歡你因為沒有走音而不唱.
是從心而發的神喜歡.
在微信的朋友去講你生命的時候.
哪怕他們說信耶穌的人是傻的.
又上教會又不升職.
又不會踩著別人上去.
你打算沒得升嗎?.
是啊沒得升.
我已經足夠了.
除了在信心堅持這件事之外.
在當時猶太教與我們有什麼關係.
我們不是猶太教.
我們現在不會回去就路猶太教.
到今天的應用是什麼呢?.
今天應用當時猶太教最大的問題是什麼呢?.
在耶穌的看法.
是你只是獻祭物和禮物.
你的心是沒有.
舊約的律法你表面守了.
其實你沒有守你內心的律法.
但是新約就將律法印在你的心.
與今天我們所講什麼呢?.
今天所有的宗教行為.
沒有心的宗教行為.
就是猶太教.
再講一次.
今天在教會的事奉.
沒有心的事奉.
就是猶太教.
今天在教會的教導.
沒有心的教導.
愛心,信心的教導.
你就是犯錯人.
與我們大家聚在一起去學神學.
去學主一學.
所有的門訓.
只是學技巧.
你只是學外面那層.
外面那層的信仰.

$^{601}$就是猶太教.
你就是走在回頭路.
今天在教會.
你有沒有認識這樣的人.
今天你問自己.
是不是這樣的情況.
我們只是問你有沒有事奉.
沒有問你事奉的心是怎樣.
你為什麼事奉.
你繼續以這個心去事奉嗎?.
你在事奉當中.
有沒有承受傷害,承受痛苦?.
有.
有沒有被人排斥?.
有.
但是你的心有沒有堅守?.
而你這個堅守.
是更加要除了信心之外.
你的愛心要堅守.
你越事奉.
你的愛心可能會越冷淡.
因為你的愛心出去了.
原來出去了.
是被別人糟蹋的.
你的真誠出去了.
是被人到處捧著捧著的.
所以你為了要保護這個.
你自己要自保了.
你的愛心就開始冷淡了.
我開始不要沾在群體上了.
不要說自己的事了.
不要分享這麼多了.
我掙扎不掙扎.
不要跟別人說這麼多.
我那時候到處傳的.
我越來越不是帶著真誠的心態說話了.
但是反過來.
你要堅持愛心的話.
你就不要這樣做.
不要以冷淡,離群,自我保護.

$^{641}$去回應信仰.
這樣是一個亂水的青蛙.
你遲了信仰就死了.
你這樣信仰越信越久之後.
你越來越後退的時候.
你好像越來越感受不到神的愛.
因為你越付出少.
你越感受少.
你去到團契小組.
你只要保護自己的時候.
你越跟別人分享少了.
你是經歷少了.
所以為什麼我對著猶太教背景的人.
叫他們不要離教.
為什麼講完這件事之後.
會講一個人與人之間的關係呢.
因為當時他們面對猶太教的衝擊.
迫不及待的時候.
他其中一個做法就是.
我不要那麼積極回教會.
甚至回少一點.
別人就不知道我信耶穌.
我回少一點聚會.
別人就看不到我.
別人不會抓我財產.
不會抓我坐牢.
我就越來越少回聚會就可以了.
今天我有網上聚會.
我越來越多網上聚會回少一點.
實體聚會回少一點.
我就很少會受人傷害了.
也不用被全黨叫我事奉了.
煩死了.
抓著我經常抓著我事奉.
我在網上敬拜不知有多好.
你就倒退了.
因為你不需要去愛人.
你躲在房間裡.
閉門造軀.
你當然不用洗碗.

$^{681}$你不用洗碗.
當然跟你太太和丈夫沒什麼爭執.
但是在爭執當中才是你的事.
試一下你是否彼此關心.
激發愛心.
勉勵行善.
今天我們教會搞很多很多聚會.
每年的行事曆滿滿的.
沒有一個會問.
究竟每一個聚會.
我們是否去到彼此關心.
激發愛心.
勉勵行善.
或許你會想.
這些都是傳道人的責任.
那你就不要獻這個祭了.
留下這個活祭給傳道人去獻.
因為你不需要愛.
全部都是由傳道人去愛.
組長去愛.
那你就沒有獻活祭了.
你的愛心就已經退縮了.
你擺明是退縮了.
總要保持弟兄的愛.
不要忘了用愛心去接待人.
要紀念.
行善,眷恕.
這些都是祭.
剛才是信仰的祭.
現在是愛心,行善的祭.
要堅守聖潔.
不要只講信,忘,愛.
講到一口都是神學的理論.
你有沒有過聖潔的生活.
你有沒有妒忌.
有沒有戀權.
有沒有淫慾.
有沒有小看人.
各方面這些是追求聖潔.
我們上一次問自己.

$^{721}$原來我這麼壞的自己.
這麼壞的自己.
這麼小人的.
這麼小家的.
有沒有這樣去問過自己.
我真的不聖潔.
我真的和我的.
我真的不相稱.
持守著聖潔.
不要自私,不要私慾.
不要覺得出公數就狂洗.
沒有所謂的.
出公數的.
有得去享樂就去享樂.
我們來教會是事奉.
不要在私慾和淫亂去低頭妥協.
但很多人都是會很容易.
在工作,在家庭,在事奉.
都可能會私慾和淫亂.
低頭妥協是很容易的.
難一點的.
是真的以神為滿足.
我覺得守聖潔.
是今天教會是難講的.
Very hard to swallow.
很難去講你的罪.
或者我的罪.
是很難講.
因為很難聽.
我們大家都是要政治正確的.
所以要疼愛著大家.
所以我不會講你的私慾.
你也不會講我的私慾.
或者你講的說話很毒.
毒舌大狀的毒.
我都不會說你不要這麼毒.
我最多和你避一避.
和你保持距離.
結果大家就不是互相去鼓勵.
這個就引申.

$^{761}$原來我們自己私生活.
就見到你是不是聖潔.
應該尊重婚姻.
不是你有沒有第三者這麼簡單.
是你有沒有尊重你太太.
有沒有真的好好尊重你的婚姻.
因為神是一定會審判.
也不要貪戀外財.
或者貪戀任何的東西.
因為神講過.
祂是不會離棄你的.
祂不離棄你就已經足夠了.
你還貪甚麼.
你貪那些.
貪完也好.
神離棄你就更糟.
我們要竭力追求這兩樣東西.
人的和睦,聖潔.
兩者是互相.
聖潔的人.
他自然會追求和睦.
和睦的人.
是因為他從上而來的那種聖潔.
所以我們要小心.
免得失去了神的恩典.
原來我們不追求這東西的時候.
我們會失去神的恩典.
我們越不追求.
越少.
很奇怪.
你不追求和你太太有好關係.
你發覺關係愈差.
恩典愈少.
你有沒有想過.
我發覺如果我不疼愛太太的時候.
我覺得好像.
這個婚姻沒有甚麼恩典.
但是你很有心追求神.
你幫我吧.
做一個好的丈夫.

$^{801}$你覺得恩典有難度.
你在教會,在小組.
你真的關心別人的時候.
當你關心別人的時候.
你發覺雖然有少少受傷害.
但是你覺得你經歷了神的恩典.
在你關心的那一剎那.
在你得力的那一剎那.
你經歷了神的恩典.
真的很美,很奇怪.
非常之好.
但是你不要說.
不要緊.
過了之後我再後悔.
時間不到你決定.
以素之後失去了.
賣了名份之後.
他想要承受這個祝福.
對不起.
我賣錯了.
可以還給我嗎.
太遲了.
所以問題就是你後退.
後退到何時.
是成為你的滅亡呢.
這個就問到.
究竟菲利普耶穌犯了八年.
他後退了八年.
是不是他已經滅亡了.
你覺得菲利普耶穌有沒有滅亡.
我覺得他沒有.
因為經文是鼓勵我們.
你只要不繼續這樣做.
you never know.
你只要迴轉的時候.
那一剎那你得回生命.
我相信菲利普耶穌.
和大衛一樣.
他只要迴轉.
認罪.

$^{841}$神是給他的.
要持守.
很奇怪.
在最後那三章.
怎麼會說這段經文呢.
為什麼要持守.
這個是剛才那幾個持守之外.
他加了另一個持守.
是什麼呢.
要以持守信仰的領袖作榜樣.
不要跟隨那些背道者.
有樣學樣.
教會裡面有很多背道者.
剛才我講了.
很多人事奉.
但是沒有了心.
他事奉他沒有了心.
我聽過一個.
我不可以提名.
一個教神學的人.
他說.
我不需要有朋友.
我只需要有結果.
你教什麼神學啊.
你不需要有朋友.
那有弟兄嗎.
你不是當我朋友嗎.
你不是當我弟兄嗎.
你只要結果.
你什麼神學啊.
教會只要人數.
我只要人數.
我不需要有弟兄才能成長.
什麼教會啊.
舉例.
持守信仰的領袖.
很奇怪.
在希伯來書.
他講了兩班人.
講舊的一班人.

$^{881}$他硬著心就跌倒了.
但是同時間.
在十一章的時候.
他講到一班舊的信心偉人.
他說什麼.
他說他還沒得到那些東西.
但是他是憑著信.
這些人我們要作榜樣.
所以很奇怪.
在希伯來書最後的時候.
他叫我們要以領袖作榜樣.
但是以什麼領袖呢.
不是每逢領袖都作榜樣.
千萬不要我講一句廢話.
廢到不得了.
我是牧師.
我要馴服他.
余全土我要馴服你.
廢到不得了.
因為聖經不是這樣教的.
就這樣廢.
他怎麼說.
你要紀念那些領導過你們的.
有尾注.
把神的道傳給你們的人.
你們要什麼.
觀察他們一生的成果.
要效法他們的信心.
前後夾住.
下面夾住.
就是你要聽從那些領導你們的人.
要馴服他們.
因為他們是為你們的靈魂警醒的.
他們的心是願意凡事進行正道.
這些傳道人.
他的心是為你的靈魂警醒.
而他自己也是進行神道.
那些傳道人.
你就要看他做一個榜樣.
好像舊約那些信心英雄.

$^{921}$要看著一個榜樣.
不要硬著心.
看著一個榜樣.
臨死的時候.
應該他還沒有得到.
但是他仍然咬著牙根去忍.
你要看他們.
但是其他的傳道人.
牧者.
若然只是講數字.
只是講外表.
只是叫你馴服不關心你的.
你千萬學都不要學.
我不同意.
你說我不同意.
我最喜歡學這些.
就是中間夾住了.
你們千萬不要受這些各樣的.
什麼?.
太小的字.
各樣怪異的教訓勾引他.
這些就是怪異教訓.
所以.
終結.
兩個字.
要忍.
要有個心.
經文怎麼說?.
你要有坦然無懼的心.
什麼叫坦然無懼呢?.
其實原文是講一個.
你是不怕講話的.
你很有信心講話的.
背景是什麼呢?.
有兩個背景.
一些奴隸在主人面前是不敢講話的.
一些打仗的人去到戰場.
他不敢走.
因為怕會死.
但是我們已經是得自由的人.

$^{961}$我們現在在天家這個致聖所.
耶穌基督的寶血.
已經將致聖所裡面的膜.
聖所和不聖所的膜.
他的身體的寶血已經拆開了.
所以意思是什麼呢?.
我們現在已經直上天庭.
我們已經去到天堂.
我們現在已經可以去到天堂.
可以和神一起聊天.
既然你可以和神聊天的時候.
你就可以在神面前什麼都講.
我很痛苦.
我有這個試探.
我被人逼迫求主你幫我.
而坐在神旁邊的是誰?.
耶穌基督.
祂是管神的家.
祂說我會幫你.
因為我經歷過這件事.
你經歷過的我經歷過.
我幫你.
只要你望著我.
你思想我.
我會幫你.
但是我們有沒有這樣做?.
我們有沒有這樣做?.
我們跌倒.
一次又一次又跌倒.
一次又一次跌倒.
但是當我們轉眼又望著耶穌基督的時候.
我們就能夠去得勝.
所以最終.
很多人發覺很多字眼.
你叫什麼?.
原來是要忍.
忍受忍受.
耶穌祂都說忍受忍受.
有時你不是被逼迫的.
是神對你的管教.

$^{1001}$是神對你的試煉.
你要怎樣?.
你都要忍受神的管教.
這樣叫忍.
所以回到家裡.
你對著你太太.
要怎樣?.
要忍.
但是忍是什麼?.
那種好像你跑步的時候.
跑長跑.
你知道再跑多沒多久.
我就去終點.
終點就可以和你的兄弟去吃早餐.
我就忍忍忍忍.
好像我上去洗手間.
很急的時候.
我就祈禱.
我知道我一上去洗手間的時候.
會很釋放.
會很舒服.
忍忍忍忍.
因為我有個盼望.
我有個很實在的盼望.
我經歷過有個盼望.
整本希伯來書.
就是講你的盼望.
所以你要持守你的盼望.
你的盼望就幫助你今天的忍.
而我忍.
在婚姻上忍的是什麼?.
那個盼望就是.
我能夠我和太太的愛.
是經歷過任何的試探都好.
到最後我都是拖著手.
進老人院.
或者一手拖著手上天家.
而這個我覺得是我婚姻愛情.
最好的結晶品.
而現在所經歷的事.

$^{1041}$小小的事是值得忍.
所謂完結的時候.
怎樣忍?.
是靠著主.
捉著這個不會失去的應許.
因為這個應許.
是神祂自己用祂的誓言.
祂的信實兩樣東西.
是和我們去作的應許.
而到最後什麼是盼望呢?.
盼望就是將來會發生的事.
你見不到的.
你未見的.
你未見的事.
你當作已經發生了.
好像一個證據.
我知道是肯定會發生.
好像前陣子.
有一個小孩一家人.
走到海裡.
不是走到海裡.
被海浪沖沖沖沖沖沖到很遠.
然後慘了.
媽媽帶著三個小孩.
慘了越沖越遠.
那怎麼辦?.
三個小孩.
結果大哥.
大哥十三歲.
大哥.
你游幾公里回去吧.
救我們吧.
因為媽媽走了.
那幾個怎麼辦?.
唯有你了.
但游幾公里很難游.
同時你是逆水流.
都沒有辦法.
幾個小時兒子都未來.
慘了兒子都不堪設想.

$^{1081}$兒子游得很辛苦游到去.
打完電話都收集了.
心想.
媽媽和弟弟都沒有了.
但當他報了警.
那些人說.
飛機來到.
當媽媽見到飛機的時候.
上了飛機沒有?.
未上.
但心裡面已經知道什麼?.
我兒子有沒有事?.
沒事了.
我們有沒有事?.
沒事了.
因為我的信心就是.
見到這件事.
我就知道我已經得救了.
雖然未去到醫院.
但見到.
我就已經知道得救了.
當是今天見到耶穌基督.
祂已經進入了節聖水的時候.
我們都知道.
我們已經得救的時候.
我們就能夠忍受.
今天所有的苦痛.
所以到最終.
你問.
菲利普恩師.
最終是怎樣呢?.
主有恩典憐憫.
讓他能夠回轉悔改.
在座當中的兄姊妹.
可能你曾經倒退.
現在可能倒退中.
可能在愛心.
在信心.
在盼望.
在聖義上倒退了.

$^{1121}$但只要你望著耶穌基督.
有回轉.
我們仍然是能夠得救得勝.
有很多人在祈禱.
所以教父千呼千喚主.
我們去到你面前.
主因為希伯來書這個訊息.
實際上是hard to swallow.
我們很難去面對自己.
原來是慢慢慢慢倒退的處境.
主在我們群體當中.
有很多的雲彩見證人.
在聖經也是.
去鼓勵我們.
和我們同行.
叫我們不要氣餒.
不要放棄.
我們一起去跑天國路.
將身體現象當作活祭.
今天祂的詩.
正正就是你給我們的訊息.
基督加害奉照面而求.
阿們.
\newpage



\section{以賽亞書 29:13-14}
\label{sec:031647}
\textbf{重得盼望}~~~~羅婉明傳道
\newline
\newline
連結: \href{http://www.church.com.hk/acms/content.asp?site=cdc\&op=show\&type=product\&code=031647\&layout=sermon}{\texttt{031647}} ~~~~ 語音日期: 2026-02-15
\newline
\newline
\hyperref[sec:031646]{\small{< < < PREV SERMON < < <}}
~
\hyperref[sec:index_chronic]{\small{[返順時目]}}
~
\hyperref[sec:031648]{\small{> > > NEXT SERMON > > >}}
\newline
\newline
以賽亞書 29:13-14
\newline
\begin{longtable}{cl}
\hline
\hline
章節 & 經文 (和合本修訂版)\\
\hline
29:13 & \begin{tabularx}{0.7\textwidth}{X} 主說：「因這百姓以口親近我，用嘴唇尊敬我，心卻遠離我；他們敬畏我，不過是領受前人的命令。 \end{tabularx} \\ \\ \relax
29:14 & \begin{tabularx}{0.7\textwidth}{X} 所以，看哪，我要在這百姓中行奇妙的事，就是奇妙又奇妙的事。他們智慧人的智慧必然消滅，聰明人的聰明必然消失。」 \end{tabularx} \\ \\
[1ex]
\hline
\hline
\end{longtable}
$^{1}$鼎鼎鼓鼓 早安.
想問一下.
大家有沒有玩過密室逃脫?.
(笑).
有些有 有些沒有.
我記得有一年.
我一家四口.
和我弟弟一家五口.
好好當當去旅行.
我們去西安玩.
其中一個節目就是去玩密室逃脫.
(笑).
因為我們覺得.
在內地玩會闊落很多.
而且更刺激.
很特別的.
其實玩密室逃脫.
你會見到不同的人.
是不同的反應.
(笑).
而且你進去.
就算你選了主題.
其實你猜不到裡面玩什麼.
都是一個挺刺激的.
看到人進入一個新的環境.
那些反應.
和你內心的心態.
例如我自己.
一進去玩的時候.
我都會覺得有點不安.
你會想 怎樣玩?.
然後就想.
快點玩完吧.
(笑).
我不知道你會不會這樣.
去到一些新的環境.
快點玩完出來.
恢復你之前熟悉的環境.
但玩著玩著又會投入.
但當然有不同人的反應.

$^{41}$其中.
我又爆我兒子.
玩完沒有?.
快點玩完吧.
他不是很投入去玩.
但有些人會很投入去找線索.
我想.
說這個.
我想說.
人面對新的環境.
或者面對一些他意想不到的環境.
或者一些不同的環境.
是一個很多不同的反應.
但總是.
去到差不多的時候.
你會覺得 真是好了 玩完了.
會有一種盼望在你裡面.
我們今天就說.
在《異塞俠書》裡面的崇德盼望.
這個崇德盼望.
既然是崇德盼望.
到底當日他們是面對一個什麼的情景呢?.
好.
在29章 剛才沒有讀.
29章第一節.
就出了這個.
阿利伊勒這四個字.
其實我相信這四個中文字.
大家都認識的.
黏在一起就不知道是什麼.
其實這個是希伯來文.
那裡翻譯出來的.
其實是為什麼呢?.
就是在說祭壇的奴.
祭壇 你會知道是什麼.
在聖殿裡面就一定有祭壇.
因為那個祭壇要獻祭.
那個奴就是.
獻梵祭的那個奴.
當你說.

$^{81}$祭壇 說聖殿.
你很快就會想到.
神同在.
所以這裡.
再說下去.
你再推就是在說.
祂在對什麼人說.
神的子民.
代表耶路撒冷.
在第一節到第八節.
就出現了六次.
通常這樣重複.
你和我.
都會這麼聰明.
就會知道 應該是很著重.
就是說 喂 和神的子民在說.
和耶路撒冷在說.
今天就是說 和你和我在說.
我們都是屬於神的人.
屬於神的子民.
和他們在說 說什麼呢?.
這一段經文.
很快你就會知道.
對象是和我們說話的.
和神的子民說話的.
不是和外邦人說話的.
有對象了.
到底當日在經歷什麼呢?.
你看這幅地圖.
在這幅地圖裡面.
你第一眼看到.
最大的那一塊.
是亞述.
然後呢.
我就說.
找不找到耶路撒冷呢?.
大家找一找.
找到的 大家.
不過呢 很小.
很小.

$^{121}$耶路撒冷屬於猶大.
南國 北國 南國.
猶大 看到這裡淺色了嗎?.
淺色了.
很淺色的.
你就會發現耶路撒冷.
跟著亞述.
你會看到什麼呢?.
這個大到不得了.
這個小到不得了 當日發生什麼事呢?.
亞述打.
耶路撒冷.
亞述打猶大 圍攻耶路撒冷.
只是看這幅地圖.
你就會看到.
怎樣打.
是不是?.
我們說的是什麼呢?.
大石砸死蟹.
或者大象的腳一踏下來.
其實怎樣打?.
沒得打的.
所以你會想像到.
這段經文是跟神的子民說.
當日他們就是被圍困的.
亞述正在打他們.
你和我一看.
我們這樣看就知道.
其實他們都知道.
沒得打.
知道結果嗎? 知道.
一定輸的.
出路是什麼? 都不知道可以怎樣.
你會想像到.
其實這段經文說的對象是誰.
和他們正在經歷什麼事.
面對一個這樣的環境.
到底裡面的人是怎樣呢?.
我們就看下去.
這段經文.

$^{161}$由九節到十六節.
對三類人說話.
有三類人.
我們逐個說.
第一類是昏蒙的人.
我們看下去.
第九節說.
你們艷樂.
昏迷.
你們醉了不是因為有酒.
你們東倒西歪.
不是因為龍酒.
其實背後都在說.
他們都在艷樂.
其實是沉迷在裡面.
再加上.
其實他們很迷失.
蒙蔽著 沉迷著.
再看下去.
第十節我們一起讀.
因為耶和華將沉睡的靈.
囂冠你們.
封閉你們的眼.
蒙蓋你們的頭.
你們的眼就是先知.
你們的頭就是先見.
等一會我們看東西.
靠眼睛看.
想事情用頭去想.
用腦去想.
但這裡是說.
我們如何得到耶和華的話呢.
就是透過先知.
先知和先見都是一類人.
就是耶和華的代言人.
因為先知書常常說.
耶和華的話臨到誰.
臨到他之後.
他就將耶和華的話.
跟神的百姓說.

$^{201}$但這裡說的是.
蒙蓋頭.
封閉眼睛.
什麼都看不到.
什麼都聽不到.
即是說.
連可以聽到神的啟示.
的先知什麼都聽不到.
都沒有神的啟示.
沒有東西在裡面.
然後這裡說.
將沉睡的靈囂冠你們.
其實是說.
神正在容許.
沉睡的靈囂冠你們.
即是說.
大家都聽不到.
以致.
大家都不知道.
發生什麼事.
再加上.
還有一個.
你發現.
再諷刺一點.
諷刺是什麼呢.
他說.
所有的密士.
你們看封住的書卷.
我交給你讀.
識字的人.
他就說讀不到.
封住了.
不識字的又說讀不到.
讀不到.
即是說識字又讀不到.
不識字又讀不到.
即是怎樣.
拒絕去聆聽.
和接受.
對神的話語完全沒有興趣.

$^{241}$沒有興趣.
即是說.
腦又封閉.
眼又盲.
然後腦閉塞.
剛才說過昏迷.
如果醫學上昏迷.
你會好多.
醫護.
醫生就會馬上去搶救你.
即是馬上去搶救.
大家不要介意.
過新年說這些.
但問題是如果在醫學上昏迷.
我們要馬上去搶救.
但弟兄姊妹.
如果在屬靈上我們昏迷了.
你會馬上去搶救嗎.
還是你都不知道自己昏迷了.
這裡都是很諷刺.
你都問問自己.
或者你都問問你自己.
你和我都一定識字.
你識字 但會不會看到.
就已經睡著了.
都看不到下去.
看不到神的話語.
這個真是上帝提醒我們.
我們會不會對神的話語.
沒有興趣.
或者我們的心都塞實了.
腦又閉塞了.
求主去幫助我們.
我記得初返教會的時候.
就會聽一些見證.
當年都有一些長者.
那些婆婆.
他們是不識字的.
但我相信今天我們在座的很多都不會.
但他們有些是真的不識字.

$^{281}$從小沒有讀書.
但我們就會聽.
牧師再分享他們.
很特別 他說那些婆婆回來後.
他們會很大聲去讀土文.
唱詩歌.
然後他們又會嘗試一起去看聖經.
去到一段時間.
他們會說.
他們現在整本聖經都看得到.
這是一種很大的.
覺得上帝很奇妙.
所以求神幫助我們.
今天你和我都識字的時候.
我們不要再說.
看不明白 看不到.
求神幫助我們.
對神的話語有一份渴慕.
以致不會.
今天我們失卻了.
對上帝追求的那份認識.
不再昏迷.
求主憐憫我們.
不要沉睡.
不要封閉.
不要蓋住了.
第二類人.
第二類就是虛偽的人.
剛才大家都有讀.
他說那些百姓親近我.
然後用嘴唇尊敬我.
他的心是遠離了我.
這個親近.
其實是一個獻祭.
類似術語.
他們把祭獻給神.
就是親近.
但上帝就說.
如果你只是用嘴唇親近我.
你的心遠離我的話.

$^{321}$你根本就不是親近我.
很快我們就想.
剛才.
或者當你回到教會的時候.
你的人在.
你的心在不在.
這都是上帝.
向我們發出的提醒.
親近.
什麼是親近.
可能我的嘴唇和心很近.
但會不會真實的.
裡面是隔了十萬八千里.
又或者.
你回來的時候.
應酬.
無論如何星期日都要回來.
星期六都要回來.
我就想起.
小朋友回教會都會這樣.
要跟爸爸回去.
沒有選擇.
然後.
爸爸媽媽就會告訴你.
你要跟我回去.
不然不能獨留兒童在家.
大家都會回來.
但可能在他們裡面.
我應酬爸爸媽媽.
但今天我想問.
你和我回來的時候.
會不會我們.
其實都在敷衍神.
而上帝很想.
真的敬畏祂.
領受祂的吩咐.
我想這也是對我的提醒.
可能有時候我們都會說.
祈禱求神高沒講完的口.
我都會提醒自己.

$^{361}$不要只把口說.
我整個生命都在.
我講的東西是我真的相信.
我講的東西都是我自己走出來的.
這也是我給自己.
一個很大的提醒.
不然就好像保羅說.
我全福音但自己凡都被棄絕.
所以.
我們真的在上帝裡面.
我們求主你幫助我.
我不要有敬虔的外表.
而是我內心都是這樣敬畏你.
專注為大.
我想這是上帝很想.
我們走去經歷的.
可能有時候我們都會.
很多疑問在我們裡面.
我都鼓勵其實不要緊.
有疑問就問神.
神可以完全接納我們.
接收我們的疑惑.
我們的困難.
有時候容許我這樣說.
質問都可以.
因為神真的很喜悅.
我們那種真實在祂的面前.
當我們真實在祂的面前的時候.
祂會回答我們.
我自己都經歷很多.
很多事情真的不明白.
去問的時候.
祂未必即時回答.
有時候祂會慢慢一步步回答我.
所以我覺得神真的很好.
好在主我不要.
一個敬虔的外表.
而是我真的很想真實的我來到你面前.
因為你接納我.
所以這裡說不要虛偽.

$^{401}$不要表裡不一.
真實來到神的面前.
因為神接納真實的你.
好.
第三樣.
就是自以為是.
自以為是的人.
我們讀15節.
15節說什麼呢.
和在那些向耶和華.
深藏謀略的.
又在暗中行事.
說誰看見我們呢.
誰知道我們呢.
這裡說.
深藏謀略.
暗中行事.
說什麼呢.
見到剛才的政治環境.
亞述在打耶路撒冷.
這裡的暗中行事.
就是在說.
希西加王.
暗中去找埃及軍.
去支援.
給你地圖再看一次.
你就會明白.
耶路撒冷.
亞述.
打下來.
如果你是耶路撒冷王.
最直接最快.
找誰幫忙.
埃及.
大家都是耶路撒冷王.
找埃及幫忙.
和他結盟.
很正路.
你即時一定會這樣想.
有時我們.

$^{441}$面對困難的時候.
我面對困難.
面對困境.
即時出現.
馬上找人.
這裡就.
所以呢.
這裡說.
因為當日的亞述君.
這裡說.
好像飛腔.
很多.
所以是即時的反應.
但亭姐妹.
亦都是這個即時.
很合理很正常.
這種想法.
其實都是在問我們.
今天問我們自己.
到底這種很正常.
很合理.
是這樣的.
是不是真的上帝氣血.
是真的很顛覆.
有時信仰也是.
有些事情是.
真的想不通.
然後呢.
再加上第十六節.
為什麼會罵他們.
第十六節.
你們將事情顛倒了.
你們呢.
看yiu 匠好像泥一樣.
這裡大家會.
馬上想到耶律美書說什麼.
神是誰.
yiu 匠.
我們是泥土.
我們是祂手中的器皿.

$^{481}$但這裡說反轉了.
下面說.
被製作的物起可論製作的物說.
祂沒有製作我.
或是被創造的物論造物的說.
祂沒有聰明.
這裡很繞口.
不明白他說什麼.
其實就是說.
我們顛倒了整個位置.
神成為了.
我們手中的泥.
神成為了被造的那個.
所以祂沒有我那麼聰明.
因為我即時就想到方法.
我即時就知道怎樣做.
請姐妹想想.
有時我們.
和神之間的關係是怎樣.
會不會好像去到餐廳.
拿起餐牌.
我呢.
就是點餐的人.
我想要.
牛扒.
六成熟.
煎到剛剛好.
碟子要很軟.
很嫩.
碟子要熱過.
才可以放牛扒.
接著我要什麼.
會不會其實.
在我們的生命裡.
在我們的信仰裡.
我們和神的關係變了.
神彷彿在服務我們.
我想要什麼.
這就是在說的.
求主幫助我們.

$^{521}$猶大.
自以為是 自作主張.
很驕傲.
以至祂蔑視了神.
以神為儒媒.
這就是自以為是.
聽起來挺沉重的.
但也要再問.
弟兄姊妹.
當你看完這幾節的經文之後.
你覺得.
你是哪一類人呢.
我都問我自己.
可能我告訴自己.
三類都是.
上帝你幫我.
其實.
今天去說這段經文.
或者當你一直看耳塞.
才發覺很多的責備.
不要緊.
當責備過後.
是上帝將祂的安慰.
將祂的恩典放在我們的生命裡.
你是哪一類.
神想我們.
再一次檢視我們的生命.
主你幫助我們.
就是.
以色列人.
耶路撒冷的人.
他們面對一個這麼大的危機.
然後.
他們正在經歷一件事.
就是他們想都想不到.
十四字說.
神說他要做奇妙的事.
什麼叫做奇妙的事呢.
耶路撒冷的人想都想不到.
就是神會用愛幫人.

$^{561}$去傾福他們.
去懲罰他們.
神想都想不到.
以致他們所有智慧人的智慧.
都想不到.
聰明人的聰明都想不到.
但同時間.
當我們以為神正在傾福以色列.
但同時上帝也在傾福亞述.
在一夜之間.
在主前701年.
有一晚.
上帝讓整個亞述軍潰敗.
死傷無數.
然後大軍潰敗就退了.
在傾克間.
神的拯救就臨到.
這就是我們的上帝.
縱使有艱難.
縱使我們有得罪他的地方.
神會懲罰我們.
但同時在警告的同時.
其實你會知道.
上帝的恩典和憐憫.
都同樣存在.
他的安慰和盼望.
都放在我們的生命裡.
這就是我們的上帝.
靠著幫助我們.
醒察的時候.
我們再一次轉眼去仰望耶穌.
轉眼去仰望我們上帝的時候.
他的盼望和安慰就臨到.
剛才我們有讀.
22至24節.
那其實是一個安慰.
盼望和出路.
因為上帝要救贖.
救贖他的子民.
所以.

$^{601}$雅各不再羞愧.
不再辯色.
羞愧是什麼.
詩篇說.
倚靠耶和華的人必不自羞愧.
必不自羞愧就是.
很多事情想不通.
為什麼會這樣.
搞不定.
我真的無法向任何人說.
我信耶穌是這樣的.
但上帝說.
你倚靠神的人必不自羞愧.
不自辯色,不會再懼怕.
23節.
我們一起讀紅色字.
尊我的名為聖.
尊雅各的聖者為聖.
敬畏以色列的神.
上帝要的是.
再一次去.
甦醒和帶領他的民.
子民就是.
尊耶華為聖.
剛才說那種虛偽.
那種心口不一.
但神今天幫助我們.
重新將我們的生命降服在神裡面.
我們知道是祂帶領我們.
當我們面對困難.
困境的時候.
我不是即時去找專家.
上網.
當然.
這不是不做.
而是第一下你覺得很正常.
很合理的反應.
其實神是要我們再一次回到祂裡面.
以致我們知道祂怎樣帶領我們.
24節.

$^{641}$紅色的字.
一起讀.
迷糊的必得明白.
發怨言的必受憤悔.
昏迷的人醒了.
腦閉塞的人再一次明白.
再一次在神的話語裡面.
去甦醒.
重新得著上帝的話語.
有時也是這樣.
有時你和我都會不明白.
生命裡面走過不同的路.
其實是很難明白.
但也是在那些難明.
艱難的裡面.
很多電影節目的分享就是可能患病.
無論自己患病或家人患病.
你會發覺.
在患病裡面.
你感受到好心是神同在.
神安慰.
神保護.
還有神猜很多不同的人.
將那種愛連結在一起.
是很奇妙.
是一個paradox.
有時是不明白的.
但這就是神給我們.
一個很大的應許.
全心信靠.
清心倚靠.
你不會再懼怕.
你敬畏神 尊主為大.
最後.
我用一段經文.
送給大家.
阿倫的祝福.
大家都會熟悉的.
願耶和華賜福給你.
保護你.

$^{681}$願耶和華使他的念光照你.
賜恩給你.
願耶和華向你養念.
賜你平安.
這段經文很熟悉.
我們有一次.
在一個聚會.
當天是一個猶太人.
拉比講.
他信了耶穌.
他就講這段經文.
其中一句.
令我印象很深刻.
當天他就問.
大家.
有沒有想過為什麼耶和華.
要向你養念.
待會我們試一下.
用養念的動作來看.
這樣吧.
他在想.
他說完之後.
我真的在想.
為什麼耶和華要向我養念.
耶和華這麼偉大.
我自己想像.
這麼高大 這麼偉大.
為什麼要向我養念.
他問了幾次.
我都回答不了.
真的.
我們讀這段經文.
(笑).
他就給了我一幅圖畫.
我就明白了.
真的很感動.
我真的那一刻.
看完之後覺得很感動.
感動就是.
原來養念就是.

$^{721}$神抱起我.
當神抱起我的時候.
你會看到他用他的笑臉.
看著我.
他會很開心地看著我.
他以我為榮.
以他很深的愛去愛我.
我覺得很感動.
今天我每次想起這幅圖畫.
我就覺得很感動.
他抱起我.
任何艱難的裡邊.
困難的裡邊.
上帝抱起我.
無論我是怎樣的人.
上帝都會抱起我.
用他的笑臉看著我.
愛我.
耶和華用他永遠的愛.
愛我.
所以今天第一天節目.
在神裡面得著這份盼望.
我們醒察我們自己.
然後在上帝裡面.
我們知道神開我為寶貴.
神和我同在.
我們一起祈禱.
主耶穌.
你心知我們.
你明白我們.
你認識我們.
你也看見我們.
主耶穌.
你幫助我們.
真實地面對.
我們有不合心意的地方.
剛才那三類人.
主耶穌.
你憐憫我們.
幫助我們.

$^{761}$讓我們願意在你面前.
呼求你的幫助.
以致我們在裡面.
重新找緊.
這份愛.
這份盼望.
這份真是在上帝裡面.
主啊我重新尊你為聖.
主啊我的心裡面.
不再迷糊.
而是我認識你 知道你.
我不自羞愧.
也不再害怕.
主啊幫助我們.
求親愛的聖靈.
賜下你的力量.
賜下你的恩典.
賜下你的亮光在我們生命當中.
以致我們回應你的時候.
我們就是這樣說.
主啊我尊你為大.
我跟從你.
我單單的信靠你.
耶穌多謝你 讚美你.
我們這樣禱告 奉耶穌基督名求.
阿們.
\newpage



\section{出埃及記 24:12-18}
\label{sec:031648}
\textbf{再登山上}~~~~蘇善美宣教師
\newline
\newline
連結: \href{http://www.church.com.hk/acms/content.asp?site=cdc\&op=show\&type=product\&code=031648\&layout=sermon}{\texttt{031648}} ~~~~ 語音日期: 2026-02-15
\newline
\newline
\hyperref[sec:031647]{\small{< < < PREV SERMON < < <}}
~
\hyperref[sec:index_chronic]{\small{[返順時目]}}
~
\hyperref[sec:031649]{\small{> > > NEXT SERMON > > >}}
\newline
\newline
出埃及記 24:12-18
\newline
\begin{longtable}{cl}
\hline
\hline
章節 & 經文 (和合本修訂版)\\
\hline
24:12 & \begin{tabularx}{0.7\textwidth}{X} 耶和華對摩西說：「你上山到我這裡來，就在那裡，我要將石版，就是我所寫的律法和誡命賜給你，使你可以教導他們。」 \end{tabularx} \\ \\ \relax
24:13 & \begin{tabularx}{0.7\textwidth}{X} 摩西和他的助手約書亞站起來；摩西上了神的山。 \end{tabularx} \\ \\ \relax
24:14 & \begin{tabularx}{0.7\textwidth}{X} 摩西對長老們說：「你們在這裡等我們，直到我們再回到你們這裡。看哪，亞倫和戶珥與你們同在。誰有訴訟，可以去找他們。」 \end{tabularx} \\ \\ \relax
24:15 & \begin{tabularx}{0.7\textwidth}{X} 摩西上山，有雲彩把山遮蓋。 \end{tabularx} \\ \\ \relax
24:16 & \begin{tabularx}{0.7\textwidth}{X} 耶和華的榮耀駐在西奈山，雲彩遮蓋了山六天，第七天他從雲中呼叫摩西。 \end{tabularx} \\ \\ \relax
24:17 & \begin{tabularx}{0.7\textwidth}{X} 耶和華的榮耀在山頂上，在以色列人眼前，形狀如吞噬的火。 \end{tabularx} \\ \\ \relax
24:18 & \begin{tabularx}{0.7\textwidth}{X} 摩西進入雲中，登上了山。摩西在山上四十晝夜。 \end{tabularx} \\ \\
[1ex]
\hline
\hline
\end{longtable}
$^{1}$梁美芬議員.
\newpage



\section{創世記 11:3-4}
\label{sec:031649}
\textbf{聖經人聖經事:製造品牌}~~~~黃克勤牧師博士
\newline
\newline
連結: \href{http://www.church.com.hk/acms/content.asp?site=cdc\&op=show\&type=product\&code=031649\&layout=sermon}{\texttt{031649}} ~~~~ 語音日期: 2026-02-03
\newline
\newline
\hyperref[sec:031648]{\small{< < < PREV SERMON < < <}}
~
\hyperref[sec:index_chronic]{\small{[返順時目]}}
~
\hyperref[sec:031650]{\small{> > > NEXT SERMON > > >}}
\newline
\newline
創世記 11:3-4
\newline
\begin{longtable}{cl}
\hline
\hline
章節 & 經文 (和合本修訂版)\\
\hline
11:3 & \begin{tabularx}{0.7\textwidth}{X} 他們彼此商量說：「來，讓我們來做磚，把磚燒透了。」他們就拿磚當石頭，又拿柏油當泥漿。 \end{tabularx} \\ \\ \relax
11:4 & \begin{tabularx}{0.7\textwidth}{X} 他們說：「來，讓我們建造一座城和一座塔，塔頂通天。我們要為自己立名，免得我們分散在全地面上。」 \end{tabularx} \\ \\
[1ex]
\hline
\hline
\end{longtable}
$^{1}$各位電影姐妹大家好.
一起讀聖經.
《生命樹宣教網絡》.
呈獻給大家的《聖經人聖經時》.
讓我們一起在聖經人和事後.
看到我們這位受弟.
今天的題目是製造品牌.
製造品牌這四個字在聖經裡面是沒有的.
不過以我讀商業的背景.
我就將這個名字放進去就不是聖經裡面的事.
不過他們是很早.
原來聖經已經有人在製造品牌了.
今天和大家看看《創世記十一章》第一節.
看看這幫人怎樣製造品牌.
那時天下人的口音言語都是一樣的.
為什麼一樣呢?.
因為他們是羅亞後裔.
「閃含瓦忽三兄弟」.
他們說的語言是一樣的.
其實天下人.
天下人其實羅亞之前的人說的話.
當然由於地域有少少不同.
但羅亞方舟洪水之後.
只有一家人.
他們一直繁衍出來的人.
都是口音語言都是一家人.
他們逐漸到他們往東遷移的時候.
往東遷移.
他們往東遷移.
是否全部羅亞後裔「閃含瓦忽」.
這麼多族人.
如果你看第十章.
裡面有很多族人.
很多都繁衍了.
是否全部都一起向東遷移呢?.
可能是,可能不是.
不過應該是有一個相當的部分.
是向東遷移.
所以才用到他們往東邊遷移的時候.
一起走去到一塊地.

$^{41}$原來是想去找一個可以長久居住的地方.
找些好地方去.
如果是當時.
當然找有河流的地方.
有水嘛.
於是他們一直往東去找.
聽聞東邊有河.
去到一個地方叫做士那.
往東遷移的時候去到一個地方叫做士那.
士那叫Shinar.
Shinar是什麼意思呢?.
士那綠色字是解釋.
Shinar是兩條河流的土地.
Shinar的意思就是兩條河流的土地.
去到士那地遇見一片平原.
就住在那裡.
所以這塊地在哪裡呢?.
Shinar在哪裡呢?.
很簡單.
兩條河流的土地.
即是兩條河流之間的土地.
兩條河流之間的土地.
這兩條河流是什麼呢?.
極有可能是幼法拉底河.
底格里斯河.
如果大家看最早期伊甸園的時候.
伊甸園有四條河.
其中有底格里斯河.
和這兩條河的流域.
幼法拉底河.
四條河中有兩條.
在那裡有兩條就不知道在哪裡.
可能都是他們交界的地方.
所以其實伊甸園不是一個很小的地方.
即是兩條河兩河流域.
兩河流域.
美索不達米亞.
Mesopotamia.
這個字是兩河流域的意思.
很有可能在兩河流域遇見一個平原.

$^{81}$當然兩條河中間.
是一個平原.
於是他們住在那裡.
是合理的.
整件事是合理的.
於是他們彼此商量.
說奈巴我們要做磚.
把磚燒透.
於是他們拿磚當石頭.
又拿石漆做灰泥.
他們說奈巴我們要建一座城.
和一座塔.
塔頂通天.
為要傳揚我們的名.
免得我們分散在全地.
第三第四節有很多字.
留意一下.
他們有兩個奈巴.
奈巴我們做磚.
奈巴我們做城和塔.
兩個奈巴.
喂你呀.
即是甚麼呢.
集中.
大家集中.
他們彼此商量.
我們一起做磚.
進步了.
懂得做磚.
代替他們以前用石頭來建屋.
石頭應該怎樣.
用小塊的石頭.
如果不夠小就打碎.
用大塊的石頭打碎.
例如小石頭.
小石頭這樣一塊.
然後放在這裡.
要砌一幢石的牆.
再小的石塊放在那裡.
讓牆不會倒塌.

$^{121}$然後用東西黏住石頭.
用甚麼呢.
用石漆當灰泥.
可能他們發現某些地面上.
有一些石油產品之類.
黏黏的.
就好像今天的那些叫做.
石漆.
就好像今天的那些叫做.
蠟青.
蠟青這樣弄在那裡.
於是又防水又黏住石頭之間.
又有小石塞住大石.
那就做好了.
一幢牆.
是比較原始.
我們現在新式了.
我們發現了原來用泥.
燒它做磚.
燒到透水.
原來它好像石一樣硬.
而且在做泥的時候.
可以照我們的形狀大小.
四四方方長方可以疊起.
就當做石頭疊起它.
不用以前那樣凹凸凸的石.
又要找小石塞住.
我們是用石漆當灰泥.
以前那些石就用一些灰,泥,土.
疊住它就起了一幢灰泥的牆.
現在我們厲害很多了.
用石漆.
用蠟青做的.
或者是用石油產品做的.
是會防水兼黏住的.
進步了.
我們有技術了.
第二次內壩.
內壩我們要建造一座城.
和一座塔.

$^{161}$塔頂通天.
這裡講城和塔.
建造城做什麼呢?.
我們要住.
一起住就行了.
為什麼要建城呢?.
城就只有城牆和城門.
於是裡面的人住在一起.
大家互通.
是一個社區.
另外是對抗外面有些人.
對抗外敵.
外敵有什麼人呢?.
都是我們以前羅亞下來的兄弟.
但很有可能這麼大族人之後.
大家會沒那麼認識.
多了不同的搶奪或者仇口.
或者有牙齒印等等.
怕別人來攻擊.
來吧.
我們建造一座城.
再加上建一座塔.
而這裡出事了.
塔頂通天.
塔頂通天這個故事.
令到我們很多人都說.
其實他們建一座塔.
塔頂通天.
去到天那麼高.
要來做什麼呢?.
為要傳揚我們的名.
免得我們分散在地上.
很多人選這段經文.
說當日的人就頑皮了.
有少少能力.
人多就想著要做一座塔.
與天比高.
去到和上帝那麼高.
塔頂通天.
我可以成為上帝那麼厲害.

$^{201}$我個人看這段經文.
看這個上下文.
我就不認為是這樣的.
其實只是想建一座高一點的塔.
就可以望得很遠.
另一方面.
表示我們的威望.
表示我們的能力.
表示我們的國.
不是國勢.
表示我們這班人的勢力.
表示我們多厲害.
塔頂通天.
而且有座城有座塔.
如果要買東西去哪裡.
城裡很豐富.
很多東西賣.
有市場我們去.
其他民族就會來這個城.
裡面是集中.
正如你今天買東西.
你不會隨街逛.
為什麼?.
去商場.
因為那裡是一個商城.
基本上你去到就買到.
任何大部份你想要買的東西.
打擊其他很多小店.
一起去一個城.
塔頂是一個標誌.
是一個標誌.
Tower.
標誌.
好了.
今天我們拍攝一個大城市.
香港.
幾座高樓大廈.
上海最高.
台北.
這是吉隆坡.

$^{241}$他們喜歡拍攝這些高樓.
代表一個地標.
很厲害.
建這麼高是否想通天與神比高呢?.
其實不是的.
只不過是我想建高一點.
讓人家有個地標.
和有座城一起.
大家可以來我們這裡做買賣.
在這裡居住.
我們這群人的勢力越來越強.
在一起.
好了.
其實目標講得很清楚.
免得我們分散在全地上.
因此.
其實他們是想集中.
第一點.
遷移.
我們去找一個好地方.
對的.
遷移就遇到一塊好地.
然後他們就說.
來吧.
我們建一些東西.
要來傳揚我們的名.
傳揚這個字.
Asha.
Asha呢.
Asha這個字是製造.
為要製造我們的名.
製造我們的名.
即是製作品牌.
就是我要建一座塔.
建一座城.
要來Asha.
要來製造我的名.
讓人知道.
嘩.
那裡很厲害.

$^{281}$於是.
免得我們這群族人.
我們會分散.
越住越遠.
你在這裡耕田.
我又要去.
於是就.
這麼小.
我去別的地方.
不要.
住在一起.
開始過城市生活.
多數以往的地方都是.
城市聚在一起.
很多工商業.
遠些就是田地.
所以他們出去耕田.
回來這裡買.
免得我們分散.
所以這個免得我們分散.
即是說我們聚集.
仍然堅固地聚集.
第一就是遷移一個好地方.
然後建一些東西.
製造一個品牌.
傳揚製造我們的名.
免得我們分開.
我們想聚集.
基本上.
我不認為塔頂通天.
是要來像上帝那樣通天.
當然是和上帝爭鬥大.
和上帝一樣.
很厲害.
其實我想當時的人.
未必有這樣的想法.
其實他只是想製造一個.
我們的品牌.
如果想遷移一個好地方.
製造一個品牌.

$^{321}$然後聚集一些人.
基本上不能夠算為錯的.
都是一件好事.
這個可以說是都市化的.
在聖經裡面最早出現.
比今天很多人說都市化.
他還都市化得早.
不是差的.
等一下我想和大家解釋一下.
這段聖經.
想聚集在一起.
聚集在一起有沒有錯?.
沒有錯的.
想有個品牌有沒有錯?.
沒有錯的.
去找一個好的地方.
都沒有錯的.
去找一個好的生活.
可以有更加有名望.
是更加能夠大家聚在一起.
於是有一個協作.
造成更大力量.
是沒有錯的.
當時的人都沒有想錯.
而且是善用科技.
開始做磚.
而不是用石砌牆.
所以看到這裡.
其實是好事.
不過很多人解釋經文.
由於他有塔頂通天.
就說他是差的.
因為這座塔叫巴比塔.
於是我們就說他錯了.
我只是讀到這裡.
我不覺得他們錯的.
是好事.
不過如果在自己的發展.
發展,發展裡面.
而忘記了上帝.

$^{361}$於是就會開始.
認為自己的品牌夠靈.
我對手夠好,夠巧妙.
我嘴巴能洗.
我的思想厲害.
於是你會開始依靠自己.
我懂得遷移,我懂得找地.
我懂得製造品牌.
我懂得聚集人氣.
嘩!厲害.
我最厲害了.
如果這樣就錯在這裡.
似乎當日的人走錯了路.
就不在塔頂通天.
通天的意思是高.
高.
所以這座大廈叫摩天大廈.
你以為建造一個人.
不知道摸不到天嗎?.
不過叫摩天大廈是形容.
很高,好像摸到天.
但建造摩天大廈.
我們沒有人認為它與上帝比高.
因為我們知道.
其實我相信他們不是這個意思.
我只是派他們聚集聚集人.
聚集起來覺得自己有勢力.
覺得自己厲害.
腦子想事對手厲害.
這樣開始他們就不要上帝了.
這是最危險的地方.
今天跟大家說這個故事.
我是想嘗試.
特別是很多人認為.
塔頂通天這個字是對上帝的冒犯.
我認為未必是.
當時只不過是想建造高一點的塔.
那麼有多少層?.
有沒有一百層?.
如果以當時的技術應該沒有.

$^{401}$很有可能他們覺得.
五層十層已經很了不起.
已經通天了.
叫摩天塔頂.
不過其實自己也知道摸不到天.
只是做一個牌子.
很容易看到地標而已.
其實他們也很厲害.
不過有時候是厲害才會忘記上帝.
這是最危險的地方.
所以如果你有個很厲害的兒子.
如果你自己很厲害.
那你要小心.
因為很厲害的.
很多時候會靠自己的手.
自己的腦袋.
而忘記了上帝的祝福.
如果沒有神.
你哪有那麼厲害?.
最怕這件事.
所以最重要.
你記得.
要紀念做你的主.
這是傳道書說的主題.
有機會跟大家讀傳道書.
要紀念做你的主.
我們一起祈禱.
因為你創造我們人.
也將你自己的形象讓識給我們.
所以我們也會有一個創作力.
我們也會有很多巧妙的手.
有很多計算.
知道怎樣做一個城.
怎樣建一個塔更加好.
用石怎麼建都那麼高.
一塊塊石砌起來.
然後放些泥灰.
怎麼砌都那麼高.
於是他們開始做磚.
會砌出堅固的城.

$^{441}$然後又可以建高一點的塔.
在那裡他們要製造一個品牌.
這個品牌成為一個標誌.
吸引人在一起.
在那裡可以發生更大的互動.
發展更大的能力.
其實是對的.
也是上帝叫我們.
在這裡做這些事.
如果是好的.
是上帝喜悅的.
不過只不過怕.
這樣聚集起來.
覺得自己能力很高.
品牌很強.
會忘記你.
幫助我們無論自己厲害到多高.
會厲害到多厲害.
富有到多多.
求你幫助我們.
不會忘記你.
知道你是我們創造的主.
紀念你創造我們的大能.
紀念你給我們智慧的大能.
多謝主.
願你得了榮耀.
不是我們自己得了榮耀.
將榮耀歸給你.
祈禱奉靠耶穌基督名求.
阿們.
\newpage



\section{馬太福音 5:38-48}
\label{sec:031650}
\textbf{不切實際的耶穌}~~~~鄧瑞強牧師博士
\newline
\newline
連結: \href{http://www.church.com.hk/acms/content.asp?site=cdc\&op=show\&type=product\&code=031650\&layout=sermon}{\texttt{031650}} ~~~~ 語音日期: 2026-02-15
\newline
\newline
\hyperref[sec:031649]{\small{< < < PREV SERMON < < <}}
~
\hyperref[sec:index_chronic]{\small{[返順時目]}}
~
\hyperref[sec:031651]{\small{> > > NEXT SERMON > > >}}
\newline
\newline
馬太福音 5:38-48
\newline
\begin{longtable}{cl}
\hline
\hline
章節 & 經文 (和合本修訂版)\\
\hline
5:38 & \begin{tabularx}{0.7\textwidth}{X} 「你們聽過有話說：『以眼還眼，以牙還牙。』 \end{tabularx} \\ \\ \relax
5:39 & \begin{tabularx}{0.7\textwidth}{X} 但是我告訴你們：不要與惡人作對。有人打你的右臉，連另一邊也轉過去由他打。 \end{tabularx} \\ \\ \relax
5:40 & \begin{tabularx}{0.7\textwidth}{X} 有人想要告你，要拿你的裡衣，連外衣也由他拿去。 \end{tabularx} \\ \\ \relax
5:41 & \begin{tabularx}{0.7\textwidth}{X} 有人強迫你走一里路，你就跟他走二里。 \end{tabularx} \\ \\ \relax
5:42 & \begin{tabularx}{0.7\textwidth}{X} 有求你的，就給他；有向你借貸的，不可推辭。」 \end{tabularx} \\ \\ \relax
5:43 & \begin{tabularx}{0.7\textwidth}{X} 「你們聽過有話說：『要愛你的鄰舍，恨你的仇敵。』 \end{tabularx} \\ \\ \relax
5:44 & \begin{tabularx}{0.7\textwidth}{X} 但是我告訴你們：要愛你們的仇敵，為那迫害你們的禱告。 \end{tabularx} \\ \\ \relax
5:45 & \begin{tabularx}{0.7\textwidth}{X} 這樣，你們就可以作天父的兒女了。因為他叫太陽照好人，也照壞人；降雨給義人，也給不義的人。 \end{tabularx} \\ \\ \relax
5:46 & \begin{tabularx}{0.7\textwidth}{X} 你們若只愛那愛你們的人，有甚麼賞賜呢？就是稅吏不也是這樣做嗎？ \end{tabularx} \\ \\ \relax
5:47 & \begin{tabularx}{0.7\textwidth}{X} 你們若只請你弟兄的安，有甚麼比別人強呢？就是外邦人不也是這樣做嗎？ \end{tabularx} \\ \\ \relax
5:48 & \begin{tabularx}{0.7\textwidth}{X} 所以，你們要完全，如同你們的天父是完全的。」 \end{tabularx} \\ \\
[1ex]
\hline
\hline
\end{longtable}
$^{1}$各位弟兄姊妹早晨.
元神的恩惠,常與你們同在.
舊約講完了.
我的講道終於來到新約.
天賦若許可的話.
我未來二十七次的講道.
都會一次講一卷新約的書卷.
直到講完新約全部二十七卷書卷為止.
我其實有一個很卑微的渴望.
是渴望弟兄姊妹能夠對整本聖經.
有一個整全的理解.
聖經其實不只是幫我們信耶穌得永生.
聖經更加承托著我們整個人的人生.
讓我們無論面對什麼光景都看到神的真實和恩典.
就以舊約為例.
三十九卷的書卷其實涵蓋了人生的方方面面.
無論在什麼光景,神都在.
在遠古荒涼的時代.
那個不知道人生意義為何的亞伯拉罕.
他聽到神的呼喚讓他明白自己的天命.
當那個亞伯拉罕仍然只有一個很小的家庭的時候.
神在,讓他開啟一個民族.
當這個民族在異鄉,在埃及做奴隸的時候.
神在,讓他們得到釋放.
並且在迦南地建立家園.
當他們在迦南地群龍無首的時候.
神在,為他們揀選王者,建立君主的制度.
當君王腐敗的時候.
神在,差派先知去提點那些君王.
當整個社會禮崩樂壞的時候.
神在,興起列國來審判他們.
當國破家亡的時候.
神在,讓他們有盼望.
當神的子民從被擄的地方返回故鄉.
面對頹垣敗瓦而失望的時候.
神在,教他們建立生命.
以生命為殿.
無論是喜或悲,無論是得或失.
聖經都教我們看到神在.
祂以無限的人詞支持著我們.

$^{41}$聖經的不同經卷.
其實就對應著不同的人生處境.
讓我們每一個都看到.
無論我們面對什麼人生大問題.
神都在,真切地指引我們.
舊約的所有經卷.
透露神就是那位與我們同在的神.
祂不會放棄我們.
這就是神與我們之間的約的連繫.
就因為神人之間的緊密的連繫.
最後神來到人間成為人.
馬太福音一開始就表明.
耶穌降生的意義就是「以馬來利」.
意思就是神與我們同在.
三一神透過耶穌.
啟示這個神人連繫的世界的真相.
來到新約.
就是來到神啟示的一個新的階段.
透過耶穌.
我們更加清楚地看到.
這個世界的真正面目.
今天我們一來看新約的第一卷書.
馬太福音.
我們會跟著耶穌的腳步.
我們將會看到一個奇異的.
全新的世界.
在馬太福音裡面.
耶穌有多重的角色.
祂是神的兒子.
祂是天國的王.
祂是來到這個世界的救主.
祂還有一個角色.
祂是一個新的摩西.
祂教導我們如何做天國的子民.
馬太福音.
其中一段最著名的段落.
就是五到第七章的登山寶訓.
耶穌好像摩西那樣.
上到山.
然後向祂的子民頒佈天國的誡命.

$^{81}$今天我們看登山寶訓的其中一段.
馬太福音第五章第三十八節.
至到第四十八節.
經文都印製在我們的程序表.
教背後的部分.
請弟兄姊妹打開經文.
我們都會跟著經文來講.
馬太福音第五章第三十八節.
耶穌這樣說.
你們聽見有話說.
以眼還眼.
以牙還牙.
只是我告訴你們.
不要與惡人作對.
有人打你的右臉.
連左臉也轉過來由他打.
有人想要告你.
要拿你的女衣.
連外衣也由他拿去.
有人強迫你走一里路.
你就跟他走兩里.
有求你的就給他.
有向你借貸的不可推辭.
你們聽見有話說.
要愛你們的鄰舍.
恨你們的仇敵.
只是我告訴你們.
要愛你們的仇敵.
為那逼迫你們的祈禱.
這樣就可以作你們天父的兒女.
因為他叫日頭照好人.
也照歹人.
降與給異人.
也給不異的人.
你們若單愛那愛你們的.
有什麼賞賜呢.
就是歲利不也是這樣行嗎.
你們若單請你弟兄的安.
給人有什麼賞賜呢.
就是愛邦人不也是這樣行嗎.

$^{121}$所以你們要完全.
像你們的天父.
完全一樣.
主耶穌一開授就說.
你們聽見有話說.
這樣這樣這樣.
只是我告訴你們.
應該這樣這樣這樣.
這裡其實已經做了一個.
很大的對比.
這個對比大到.
後者要否定前者.
主耶穌來到這個世界.
不是要肯定.
世界的人.
習以為常的做法.
而是去否定這種做法.
展示另外一種做法.
一種我們看起來.
不太合理的做法.
耶穌來到這個世界.
不是教我們.
如何在這個世界的人.
建構的這個世界得到成功.
而是教我們想像.
另外一個世界.
一個神聖的世界.
登山寶訓的誡命.
在識經上.
在倫理上.
在神學上.
都引起極大的爭論.
究竟這些誡命.
登山寶訓的誡命.
的性質到底是什麼呢.
有人打你的右臉.
你就連左臉也轉過來.
由他打.
有人要奪你的女衣.
你就連外衣.

$^{161}$都給他.
如果你真的照字面去做的話.
有人認為.
這樣只會令到這個世界亂七八糟.
這種教導.
看起來是不能制止惡人作惡的.
反而令到那些惡人.
變本加厲.
如果是這樣的話.
這個不義的世界.
如何管理呢.
於是不同的人提出很多不同的想法.
有人認為.
耶穌的教導.
是一種很明智的策略.
你對敵人和善.
就會令到那些敵人變成你的朋友.
你為你的敵人多走一步.
會激發他的同理心.
害你炒魷魚的那位同事.
你送禮物給他.
他會變得良善.
回過頭來.
會介紹你一份更好的工作給你做.
如果你真的這樣想的話.
你就是天真得可憐.
從現實的策略來說.
登山寶訓並不是好的策略.
有人認為.
耶穌講的是一種正向的思考.
教我們積極地面對逆境.
我們都知道正向思考.
其實不質疑世人建構的世界.
也不質疑世人追求的目標.
耶穌顯然不是這樣.
耶穌顯然要和現實世界的一切.
做一種分別.
做一個對比.
他根本不認同現實的一切.
如果他只不過是一個正向的人.

$^{201}$教人在現實的世界獲得成功.
他根本不會被人釘在十字架上.
反而會成為社會中.
個個都拍掌而認同的成功典範.
又有人認為.
由於登山寶訓的教導並不現實.
故此不能作為一個信仰的社群.
長遠運作的指引.
故此.
這些說法應該是一種特殊時期的做法.
當時有一個觀點.
認為主耶穌很快會回來.
就因為主耶穌很快會回來.
所以信徒採取了一種極端的做法.
弟兄姊妹.
如果你肯定主耶穌下星期會回來.
我相信.
你是應該願意變賣你一切.
現在所有的東西.
分給窮人.
一個信徒若認為主耶穌很快會回來.
有人拿他一件女衣.
他連外衣都給他.
其實一點都不出奇.
但如果你真的這樣想.
登山寶訓便會變得.
與我們悠長的歷史生活無關.
只不過是短暫時期的激烈做法.
但如果主耶穌不是這麼快回來.
一個這麼悠長的歷史.
這些登山寶訓便不適切.
有人認為.
由於登山寶訓難以實踐.
故此它只是一種遙不可及的理想.
如果是理想.
做不到便情有可原.
每個人都有很多理想.
哪一個能夠完全實現呢?.
由於人人現實上都做不到.
於是有人認為.

$^{241}$登山寶訓的真意.
其實是凸顯我們做不到.
多麼需要恩典.
這是強調恩典.
若果真是這樣.
登山寶訓不僅不是信徒努力實踐的倫理.
更加有些反效果.
就是誘導信徒不要去實踐.
因為那些信徒可能會這樣想.
硬要去實踐這些做不到的行為.
這樣便靠行為得救.
遠離了恩典.
於是有些人假恩典之名.
讓信徒可以安心放下登山寶訓的教導.
放下了.
這樣每一個人便可以依然固我.
不同學者千方百計為登山寶訓的教導.
找不同的理由.
因為這些教導看起來不切實際.
助紂為虐.
難以實行.
大家想想.
若果有人打你的右臉.
你連左臉都轉過來由他打.
他會不會收手?.
若果有人要搶你的女衣.
你連外衣都由他奪走.
那你還剩下什麼?.
愛你的仇敵.
為逼迫你的人禱告.
這樣會不會縱容了那些壞人作惡呢?.
面對這些教導.
不同的學者找不同的理由.
但大家想清楚一點.
這些理由其實都沒有什麼理由.
問題出在哪裡呢?.
為什麼這些教導看起來不切實際.
主耶穌又語重心長地教導我們呢?.
問題出在哪裡呢?.
問題出在.

$^{281}$我們將這些教導.
從主耶穌所建構的天國中抽出來.
然後將這一套東西.
本來是在天國中使用的.
放回我們習慣的世俗領域.
結果就會造成格格不入.
大家想像一下.
有一個外星人來到這個地球.
這個外星人和我們地球人.
有一個很不同的地方.
就是他不用吃東西.
他可以像植物一樣進行光合作用.
這樣已經獲得充分的養分.
這個外星人就告訴我們地球人.
不要為了爭奪生存的資源而開戰.
是不需要這樣做的.
你只需要到戶外走一走.
曬一下太陽吸一口炭氣.
這樣就已經有足夠的生存養分.
我們聽到這個外星人這樣說.
我們會覺得他所說的話.
完全不切實際.
瘋的.
但是又想像一下.
如果這個外星人.
真的可以在我們的地球.
建構一個平衡的宇宙.
如果你願意的話.
我們地球人是可以進入他建構的世界.
一進入他建構的世界.
我們的生命馬上會轉化.
就真的能夠自己進行光合作用.
如果你試一下進入.
你明白了.
你就會發覺原來他說的話是真的.
其實他說的話只是一個事實.
不過那個事實是在他的世界裡.
但那個世界和我們這個世界是一個平衡的世界.
他說的話原來不是無法實踐的道理.
他說的話只是呈現另外一個世界的真相.

$^{321}$問題出現了.
其實大部份人都不願意.
或者沒有進入他的平衡宇宙.
然後我們只不過是堅固地站在我們現在的習慣的世界.
然後我們在這個立場去評價那個人說的任何話.
你就會發覺那個外星人說的任何話.
都不現實,都不符合實際,假的.
問題正正就是出在這裡.
耶穌當然不是外星人.
祂比外星人強得多.
祂是神.
祂見到的世界,祂見到的真相.
比我們習慣的真相差異大得多.
對比更加的強烈.
耶穌來到這個世界.
以祂自己獨特的生命.
呈現一個又奇異又神聖的世界.
耶穌以祂自己的一生.
去呈現祂眼中見到的真相.
真相是怎樣的呢?.
耶穌是怎樣活的呢?.
祂的真相就是,在祂的世界是很真的.
就是祂完全相信天父.
祂相信天父比相信人類的各種想法更加好.
如果天父的旨意是叫祂背起十字架.
祂就會很開心很積極地全力以赴.
並不認為千方百計放開十字架.
去增加自己的利益更加划算.
祂不會這樣想的.
因為天父的想法比我們世界人習慣的想法更加好.
人生追求的目標.
不應該是由世界的世俗來設定.
而是由天父來設定.
祂覺得更加合理.
祂見的世界就是這樣.
祂見到天父是大地和歷史的主宰.
祂相信天父的人愛.
祂不同意人類用各種殘酷的手段.
會得到更高的文明進步的想法.
祂相信主宰歷史的是一個人愛的天父.

$^{361}$而不是人類那種暴力的法則.
當有人要捉耶穌的時候.
耶穌吩咐彼得.
你放下武器.
耶穌看到人生真正的保護.
只是來自天父.
人心真正的平安.
是不能夠用以暴易暴的方法來保證的.
當耶穌上到十字架上的時候.
耶穌第一句說話就說.
天父啊,赦免他們每一個.
祂的眼睛真的看到.
人間的那種痛苦和暴力.
並不是終極的真相.
並不需要用痛苦和暴力.
來主宰我們的理性和情感.
真正的終極的真相.
是人詞的天父.
耶穌完全信任他.
信任天父的人詞.
而不信任世間的暴力.
如果你看到耶穌所看到的真相.
你會明白.
登山寶訓的教導.
並不是匪夷所思,超乎常理.
登山寶訓的教導.
放在耶穌所呈現的真相裡.
是完全合理的.
但當我們抽離了.
耶穌所說的終極的真相.
我們把這些道理放在我們.
習慣的世俗所建構的世界裡.
我們這樣合起來.
顯得不切實際和格格不入.
如果我們追求的東西.
是用最有效的策略.
帶來最大的果效的話.
有人強迫你走一里路.
而你要跟他走夠兩里.
這只能夠是下下之策.

$^{401}$如果你執著個人的成敗.
是大過信服天父的旨意的話.
有人打你的右臉.
然後你連左臉都被他打了.
這根本是背離了正義.
如果你認為誰大誰惡誰正確.
就是人間的真相.
不認為人池的天父.
在主管歷史.
不認為人池的天父.
才是終極的真相的話.
愛你們的仇敵.
為逼迫你們的禱告.
只能夠是荒謬至極的歪理.
登山寶訓.
其實他不單單是講一種很古怪的倫理.
他其實是呈現一個很奇異的世界.
天父透過耶穌的一生.
呈示這個世界的真相.
在這個真實.
但我們覺得很古怪的世界裡.
你要放下世俗.
一切去跟從主.
其實這才是一個最好的人生.
在這個真實的世界裡.
施比受更為有福.
在這個真實的世界裡.
寬恕比寄愁好.
吾還手比還手好.
十字架比利劍好.
這才是人生的真相.
是的.
耶穌就是這樣活的.
祂就是活這個真相給我們看.
上帝千方百計來到人間.
就是向我們啟示這個真相.
耶穌是神.
祂就是這樣活的.
祂順從天父.
而天父也是這樣活的.

$^{441}$因為天父叫日頭照好人.
也照大人.
降雨給義人.
也給不義的人.
天父也是這樣做.
真正的信耶穌.
其實就是跟著耶穌走.
在這個跟從裡.
我們慢慢離開世俗的軌跡.
進入神的軌跡.
我們的世界.
我們見到的世界.
英文叫vision.
就會隨著這個軌跡來開闊.
慢慢你見到的世界.
你的視域開闊的時候.
你就慢慢見到耶穌所看到的世界.
有一天.
如果你真的跟著這個腳蹤而行.
有一天.
你會看到.
原來耶穌來.
祂真的開了一個新的世界.
呈現了一個新的世界.
活現了一個新的世界.
耶穌不斷呼召我們.
跟從祂.
祂期望我們.
在我們的生活裡.
跟從祂的緣故.
都能夠活出這個世界.
一個新的世界出來.
耶穌以祂的一生.
呈現了一個新的世界.
也都期望我們.
用我們的人生.
活現這個世界.
登山寶訓.
不是拿來作為一個策略.
去在一個舊的世界裡成功.

$^{481}$是不可以的.
其實是活現一個全新的世界.
今日講道的經文.
最後的一節是這樣的.
《華大福音》五章四十八節.
所以你們要完全.
像你們的天父.
完全一樣.
這裡講的完全.
當然不是指完美無瑕.
而是指一種屬乎天父的印記.
一個記號.
天父是完全的.
而屬乎他的人.
就會打下一個屬乎他的印記.
有了這個印記的人.
他就應該信服天父的旨意.
以天父的人愛為目標.
這些人有了天父印記的人.
他就應該放下世界的種種暴力.
平安自在地.
在人間創造和平.
創造復和.
這樣的人.
他用一種很獨特的人生.
活現一個全新的世界.
這個就是天國的真相.
天國的真實.
耶穌在登山寶訓裡面.
曾經教過我們祈禱.
在登山寶訓裡面.
教我們用主禱文來祈禱.
主禱文其實是登山寶訓的一個段落.
主禱文是這樣說的.
天父啊,願你的國降臨.
願你的旨意行在地上.
如同行在天上.
當我們看到耶穌生活.
背後的那個神聖的世界.
而我們又願意跟著耶穌來走的時候.

$^{521}$其實我們就是將天父的旨意.
行在地上.
而當我們這樣做.
這樣實踐的時候.
其實天國的真相.
就會呈現在人間.
每次的主日崇拜.
我們都念主禱文.
念主禱文每一次都提醒我們登山寶訓.
登山寶訓的時候.
告訴我們.
其實我們正在活在一個平衡宇宙.
另外一個世界.
這個世界行事的法則.
和這個世俗.
我們習慣的世界行事的法則.
剛剛相反.
弟兄姊妹.
希望你們能夠完全.
像你們的天父.
完全一樣.
願神祝福你們.
願榮耀歸給聖父.
聖子.
聖靈的名.
阿們.
\newpage



\section{馬太福音 6:19-24}
\label{sec:031651}
\textbf{見錢眼開?}~~~~石偉漢牧師
\newline
\newline
連結: \href{http://www.church.com.hk/acms/content.asp?site=cdc\&op=show\&type=product\&code=031651\&layout=sermon}{\texttt{031651}} ~~~~ 語音日期: 2026-02-15
\newline
\newline
\hyperref[sec:031650]{\small{< < < PREV SERMON < < <}}
~
\hyperref[sec:index_chronic]{\small{[返順時目]}}
~
\hyperref[sec:031652]{\small{> > > NEXT SERMON > > >}}
\newline
\newline
馬太福音 6:19-24
\newline
\begin{longtable}{cl}
\hline
\hline
章節 & 經文 (和合本修訂版)\\
\hline
6:19 & \begin{tabularx}{0.7\textwidth}{X} 「不要為自己在地上積蓄財寶；地上有蟲子咬，能鏽壞，也有賊挖洞來偷。 \end{tabularx} \\ \\ \relax
6:20 & \begin{tabularx}{0.7\textwidth}{X} 要在天上積蓄財寶；天上沒有蟲子咬，不會鏽壞，也沒有賊挖洞來偷。 \end{tabularx} \\ \\ \relax
6:21 & \begin{tabularx}{0.7\textwidth}{X} 因為你的財寶在哪裡，你的心也在哪裡。」 \end{tabularx} \\ \\ \relax
6:22 & \begin{tabularx}{0.7\textwidth}{X} 「眼睛是身體的燈。你的眼睛若明亮，全身就光明； \end{tabularx} \\ \\ \relax
6:23 & \begin{tabularx}{0.7\textwidth}{X} 你的眼睛若昏花，全身就黑暗。你裡面的光若黑暗了，那黑暗是何等大呢！」 \end{tabularx} \\ \\ \relax
6:24 & \begin{tabularx}{0.7\textwidth}{X} 「一個人不能服侍兩個主；他不是恨這個愛那個，就是重這個輕那個。你們不能又服侍神，又服侍瑪門。」 \end{tabularx} \\ \\
[1ex]
\hline
\hline
\end{longtable}
$^{1}$我們一起祈禱.
天父多謝你讓我們藉著你的話去屠宰我們.
特別藉著耶穌在山上的補糞.
讓我們明白上帝的旨意.
讓我們知道如何在地上活得更合理心意.
求主你幫助我們.
讓我們明白今天所思考的聖經.
祈禱奉耶穌名字而求.
阿們.
今天跟大家說的題目叫見淺開眼.
不知道大家知不知道.
見淺開眼這個字都有典故.
不是亂說的.
有典故.
原來根據一些書.
《金瓶梅》這本書在第八十回裡都說.
漢官聽說院中唱的.
以賣肖為活的.
將脂粉作生涯.
稀究盈身.
見淺開眼自然之利.
在妓院裡面的人.
以他們的技能求生.
他們會稀究盈身.
原來稀究盈身在這裡都有出現.
見淺開眼原來是自然之利.
不知道大家有沒有覺得.
見淺開眼是一個很自然的道理.
也是一個很正常的行為.
不過今天我們會.
一直看下去.
我告訴你見淺開眼是不合理的.
待會我們有機會再到後面去看一看.
當然我們先說這個題目.
我們當然大家都拜過年.
馬年馬年.
我們2026年農曆新年.
我們下個星期就是了.
剛才有不同人的公開說過馬.
我用什麼呢.

$^{41}$我用的馬字都和聖經有關.
舊約的.
大家猜到我用哪一句祝大家馬年嗎.
馬太福音.
馬太福音都對了.
馬太福音都有新約裡面有二馬來離.
二馬來離就有馬.
這是新約的.
今天我選的祝福的說話就是來自舊約的.
這個大家可能很少聽的.
這個舊約.
不過大家儘管學一下.
叫與馬同跑.
有沒有人聽過與馬同跑這個概念.
真的沒有嗎.
與馬同跑是來自耶利米書.
十二章五節的.
裡面說你與步行的人同跑.
尚且覺得累.
怎能與馬賽跑呢.
你在安全之地.
尚且會跌倒.
在一旦河邊的叢林.
要怎麼辦呢.
不過這個我用了和修本.
或者和學本都用了與馬賽跑的.
有書用了.
有些寫作的書用了與馬同跑.
其實是出自耶利米書.
這個的意思是什麼呢.
就是說我們如果日常連一些瑣碎事.
好像步行一樣.
都常常覺得很累.
很辛苦.
沒力氣.
那你怎樣能夠有力量.
與馬一起賽跑呢.
這裡都是一個勉勵.
我們不要停留在我們現在的階段.
已經覺得很疲累很累.

$^{81}$我們應該追求卓越.
突破安息圈.
我們要在那個辛苦的日子裡面.
我們要堅持下去.
盼望大家在今年裡面都努力去堅持.
不過牧師和你就與馬同跑.
祝你.
但是你在農曆新年的日子.
可能你遇到其他的人.
就不一定是這樣祝賀你的.
特別我想起我小時候.
經常聽到有一個很出名的歌星.
在新年的日子.
通常都會播這些歌星的歌.
這個歌星就是歌神.
知道是誰嗎.
大家不同的年代.
這是張學友.
這是許冠傑.
我就是Johnny的年代.
許冠傑.
許冠傑他的歌就很多好聽的歌.
鐵塔靈魂那些.
大家當然一定不會忘記.
不過農曆新年大家通常有兩首歌.
大家一定記得他會唱的.
第一首是什麼.
財神到.
因為大家一定記得.
財神到 財神到.
好心就得好報.
財神話 財神話.
搵錢醫正勞.
財神到 財神到.
好走快兩步.
得到他看起你.
你又前途.
這首歌小時候經常聽.
新年就聽到聾了.
所以原來新年.

$^{121}$有些人都會恭祝你財神到.
希望你家將財字有時候都反轉.
希望他財神到.
有些人真的期望錢多點來.
所以有另外一首.
Sam Huy 許冠傑的歌是什麼呢.
歌名叫天才白痴錢錢錢.
聽過沒有.
聽過了.
會唱.
最重要他的前奏是.
錢錢錢錢.
一路整個前奏都是說錢錢錢.
前奏就是這樣.
貧賤咬頸多.
生仔都諗過.
阿女嗌肚餓.
屋租都要拖.
有什麼會令你心花怒放呢.
諗得多.
終於癲癲喪喪.
有了他當然是大跛腳.
雙腳都沒傷過.
無得擋無得擋.
就是和錢有關.
所以Sam Huy這兩首歌.
很多時候大家一想到錢.
就想起他這兩首很經典的歌曲.
不過今天牧師講到.
就不會講Sam Huy的歌的內容.
或者那個重點.
不過今天都確實和錢有關.
事實上因為.
你知道耶穌做了很多的教導.
但耶穌其實他的教導.
不只是我們上次講到這麼簡單.
只是說施捨.
說禱告.
說禁食.
在馬達福音.

$^{161}$第六章十九節.
原來耶穌都開始和門徒講錢.
講錢銀的事.
所以今天我們看到.
馬達福音第六章十九節.
我們同樣看到耶穌.
怎樣和我們講錢.
耶穌的講錢.
和我們聽到的歌裡.
Sam Huy講的錢.
的想法又是否一樣呢.
求主幫助我們.
在這個段落我們讀到.
Alex帶我們讀的段落.
十九節到二十四節.
其實這個段落很簡單.
結構可以分成三個比較對比的概念.
由十九節到二十一節.
就是第一個小段落.
二十二節到二十三節.
就是第二個小段落.
第三個小段落就是第二十四節.
耶穌就用了這三個對比.
去將我們帶到錢財的討論.
所以今天我和大家去思考的.
就用了三個法則.
三個耶穌教導我們.
怎樣去看錢財的法則.
幫助我們從耶穌的眼裡.
怎樣去教導他的門徒.
他所愛的徒弟.
第一個讀到的經文.
就是在第十九節到二十一節.
十九二十一節裡面說.
不要為自己在地上積儲財寶.
在地上就有蟲子咬.
能受壞也有賊挖洞來偷.
要在天上積儲財寶.
天上沒有蟲子咬.
不會受壞也沒有賊挖洞來偷.

$^{201}$因為你的財寶在那裡.
你的心也在那裡.
當耶穌講財寶處理的教導的時候.
其實在路加福音也有.
路加福音記載了在不同的部分.
不是一次過連續三個記錄.
但路加福音確實這三個法則.
其實在其他的段落也有提及.
我們今天可以看到.
在路加福音十二章三十三至三十四節.
就有這個同樣類似接近的描述.
說到怎樣去面對第一個法則.
就是儲蓄法則.
儲蓄法則.
在這個儲蓄法則裡面.
我們會看到.
其實耶穌在說一個很重要的對比.
一個比較.
比較就是有地上積儲財寶.
和有天上積儲財寶.
特別第十九節一開始的時候.
說到在地上有什麼問題呢.
當我們去積儲一些財富的時候.
財寶的時候.
他說到有三樣的事情.
三樣事情發生.
而在原文裡面.
其實他的翻譯應該再準確一點.
不要為自己在地上積儲財寶.
地上那裡有蟲子咬能受壞.
那裡也有滑動來偷.
即是說那裡有蟲子咬能受壞.
那裡有滑動來偷.
就是這樣的排列.
有一個概念.
說了兩次的那裡.
一個那裡是說到.
在地上有蟲子咬和能受壞.
我們試試去看那個比較.
第一個說到有蟲子咬.

$^{241}$為什麼會有蟲子咬呢.
因為原來你要知道.
地上的財寶有什麼財寶呢.
你地上有什麼財寶可以擁有呢.
其中有一樣的財寶.
就是有些財物會藏起來.
有些財物藏起來.
對你來說是貴重.
不過他不會每次都拿出來用.
這些是什麼呢.
原來就是一些衣服.
或者一些毛皮.
或者有時候書籍.
都會有這樣的情況.
原來有錢人會怎樣呢.
買很多衣服.
古代也是.
現在也是.
當你有錢.
你買很多不同的衣服.
事實上在古代來說.
一些漂亮的衣服.
其實他們用了很多.
不同的珍貴的材料.
如果你讀聖經.
你會看到賣紙布的女人.
有些顏料紫色是貴重一點.
另外有些布料的質量.
也是比較好的質量.
所以一些有錢的人.
就會買了一些很貴的衣服.
但又未必會經常穿出來.
相對窮人.
衣服就少一點.
如果你有記得.
聖經裡面說.
有些窮人.
沒有錢到一個地步.
要當成外衣.
那件保暖的衣服.

$^{281}$都要當成.
裡面只剩下一件單薄的內衣.
所以一些窮人.
很簡單只有兩件衣服.
一件冷的時候.
也可以保暖.
一件很簡單的內衣.
就是這樣.
所以你會看到.
這裡說的是.
一些有錢的人.
他們就有貴價的衣服.
而且數量也很多.
象徵著他是很富有.
因為你也知道.
貴價衣服.
都在不同場合.
不同時候.
才會穿出來.
所以不穿的時候.
你會收在籠底.
知道什麼叫籠嗎?.
籠底.
其實是希望保護它.
希望儲存好.
但通常這些籠底.
不是今天這樣.
今天可以長開冷氣.
但在古代來說.
在耶穌的年代.
卻又不是這樣.
擺在籠底.
不是有一個.
很好的控制濕度.
溫度的一個地方.
所以他們這些地方.
有時候都比較陰濕.
不是說人.
是那些地方比較陰濕.
比較陰暗.

$^{321}$所以那些地方.
反而適合有一些蟲出現.
學者就說.
這些蟲是什麼呢?.
就是現在我在畫面上.
看到的那隻蟲.
這種叫伊蛾.
叫伊蛾的蟲.
你可以看到.
在左邊.
你可以看到.
它有很噁心的樣子.
有很多毛.
顏色是黃藹色.
不同的顏色.
重要的是.
右邊那裡.
原來這些蟲的特性.
是什麼呢?.
專愛吃一些衣物和書本.
你不會吃衣物.
但原來蟲很喜歡.
吃這些衣物和書本.
而且很喜歡潮濕的地方.
一些暗的角落.
在那裡.
慢慢地生長.
又會有蟲.
又會有成蟲.
一路一路長大.
以致在那裡.
慢慢地破壞了你的衣服.
所以原來耶穌說的.
很實際地說.
在他們的環境裡.
原來他們以為.
有很多漂亮的衣服.
收藏起來.
保住它.
有一天可以拿出來.

$^{361}$二手衣服賣.
或者有些不同的價值.
或者很漂亮的場合穿.
原來這些衣服.
都會被蟲咬了.
你以為放在籠子裡保護.
原來又不可以.
另外一個情況是什麼呢?.
就是能受壞.
能受壞是什麼情況呢?.
因為你會知道.
金屬有時候都會容易.
隨著一些環境的變化.
而出現一些受壞.
你看到這幅圖裡.
有三種不同的金屬狀況.
一個就是侵蝕.
中間那個就是氧化.
第三個就是生鏽.
一些鐵.
原來都是這樣的.
在耶穌的年代裡.
有一些的金屬.
因為你知道.
在古代來說.
金屬都是一些貴價的東西.
他們的金屬.
也不只是只有黃金.
黃金,銀,銅,鐵.
這些都是他們用到的金屬.
而且如果你記得.
在舊約的年代裡.
有一些金屬.
其實可以造成一些武器.
成為在戰場上很好用的東西.
所以他們都會用很多的金屬.
成為他們一些重要的物品.
但是這裡說.
原來這些金屬.
原來放得久.

$^{401}$放得久都會隨著.
一些的時代.
一些的氧化.
或者一些侵蝕.
等等的原因.
都會令到生鏽.
或者這麼說.
我們的翻譯叫做什麼呢?.
叫做光澤變了.
原文裡面說.
那個光澤.
由一個可能光的地方.
的狀況.
成為一個暗淡的狀況.
所以你發現.
有時候我們在一些電視報道.
一些考古的物品裡面.
就好像最幼的幅圖那樣.
有時候本來很漂亮的.
但是隨著時間.
都會變成一些綠綠的色.
好像有些髮毛.
有些顏色是啞的一些狀況.
有什麼問題呢?.
你想一下.
就是出現了.
你這些金屬的價值.
就會跌了.
你可以想像一下.
剛剛純的物料.
一定是漂亮一點.
值錢一點的.
當你發霉發爛.
接著生鏽的.
賣一條生鏽的鐵給你.
你都不願意給這麼高的價錢.
所以原來地上裡面.
就有這些狀況.
衣服會被蟲咬.
你的金屬會被受壞.

$^{441}$這裡就是說.
第一類型的一些物品.
有什麼共通點呢?.
因為我們說這裡.
經文說.
哪裡有蟲子能受壞.
這種就一定叫.
天然大自然引發的.
一些財物損失.
是透過大自然裡面的動物.
或者環境氣候等等.
所帶來的損失.
原來有時候.
我們都會面對一些.
還有一些其他損失是什麼呢?.
來自人禍.
來自人所帶來的損失.
所以到後面就說.
哪裡也有賊挖洞來偷.
耶穌就說.
一個情況就是.
有一個狀況就是.
有些賊會來偷你的東西.
當然來偷你的東西.
就不一定.
不會告訴你.
我現在來偷了.
所以你先收起來.
我一會兒來找.
就不會這樣.
趁你不知不覺的時候.
就去偷取你本身擁有的物品.
令到你本來擁有的財物.
就帶來損失.
原來耶穌就說.
一個狀況就是.
無論大自然.
或者是一個人禍的原因.
都會令到我們損失了一些物資.
我不知道你們覺不覺得.

$^{481}$今天我們都會這樣.
好端端正常工作賺錢.
不知道為什麼.
什麼時候.
某時某刻.
就會有一些事發生.
令到你帶來損失.
好像有很多財富.
轉頭又沒有了.
原來會這樣的.
我記得我初出來的時候.
一來到.
轉頭就壞了冷氣機.
正常的.
第一個月正常.
第二個月就壞了.
有時候突然發生一些事.
可能是一些物質.
可能來自大自然.
甚至有時候.
無端端又不知如何.
好端端正常放一些錢.
在一些你認為安全的地方.
原來轉頭都會沒有了.
最近在中國有些新聞.
有一間公司叫傑我銳.
傑我銳有沒有讀錯.
傑我銳.
這間公司是做黃金買賣的.
你想想黃金應該是最安全.
最穩陣的.
黃金是什麼.
通常黃金.
香港的標準叫999金.
金的成份去到99.99\%.
是很足金的.
所以你想想黃金是最穩陣的.
加上近來金價是怎樣.
升到很厲害.
人們爭著買金.

$^{521}$理論上你買一些黃金.
應該想著你發達了.
越來越多錢.
但這間公司怎樣.
爆煲.
你想著很安全買一些黃金.
或者投資在這間公司裡面.
怎知道原來這間公司.
都做了一些期貨買賣的公司.
結果因為黃金升得太厲害.
人們想套現拿回錢.
一個人套現拿錢.
公司當然有足夠的錢退給他.
但如果有十個人.
有一百個人.
個個都說黃金好價錢套現.
公司突然沒有現金付.
代表資金鏈斷了.
資金鏈斷了.
這裡不能付款.
那裡不能付款.
人們就怕你不能付款.
就追你錢.
所以原來這單新聞.
好端端投資在這裡.
原來有十五萬人.
沒有了一百三十三億人民幣.
真是不少.
今天可能你沒有買黃金.
不過你要知道.
原來我們有時候生活裡面.
總是有一些東西.
會令到我們地上財寶.
會受到虧損.
被人偷.
不知不覺少了.
耶穌說原來我們地上的財寶.
就是這樣的.
所以耶穌帶來另一個對比.
在第二十節說.

$^{561}$既然是這樣.
你地上的投資.
會有這麼高風險的狀況.
不如你做一個更好的投資.
耶穌說更好的投資是什麼呢.
就是天上積蓄你的財寶.
積蓄你的財寶在天上.
有什麼特點呢.
第二十節說.
在那裡原文都是說.
天上那裡沒有蟲子咬不會受壞.
那裡也沒有賊挖洞來偷.
原來耶穌說在天上.
當你儲一些財寶的時候.
很有趣.
原來沒有天然災害.
沒有天然災害會發生.
也不會有賊進來.
知不知道為什麼.
你有沒有想過.
我也在想為什麼不會有賊呢.
因為賊上不到天堂.
所以他上不到天堂來偷.
也是合理的.
所以耶穌說.
原來你投資在天上的財寶.
是萬無一失的.
耶穌就告訴我們.
原來有一種方法.
讓我們處理我們的財寶更好.
就是讓我們積賺財寶在天上.
那我們就要問.
如何能積賺財寶在天上.
而不是積賺財寶在地上呢.
事實上我們今天能夠出去工作.
出去生活賺錢等等.
我們也只能夠賺到地上的財寶.
我們憑什麼可以賺到天上的財寶呢.
又或者這樣說.
現在這個段落裡面說.

$^{601}$如何將你地上的財寶.
轉化成為天上的財寶.
聽得明白嗎.
所以這個段落不是說.
叫你傳福音的問題.
聽得明白我的意思嗎.
因為傳福音.
你不需要用地上的財寶也能做到.
當然有些地方.
你要用地上的財寶傳福音.
但如果你純粹傳福音.
你不需要用地上的財寶.
你也可以積賺天上的獎賞.
但這裡特別說的是.
你如何用地上的財寶.
轉化成為天上的財寶呢.
原來是有投資工具的.
剛才我們買黃金也有投資工具.
那如何做天上的投資工具呢.
原來路加福音.
剛才我們對應的段落.
就說了一個方法.
那個方法是什麼呢.
就在路加福音12章33節的時候說.
你們要變賣財產去周祭人.
為自己預備永不壞的前郎.
和用不盡的財寶在天上.
習不能更 從不能住的地方.
你會發現字眼和我們剛才讀的相類似.
不過這裡原來.
耶穌說原來我們能夠.
藉著我們變賣財產去周祭人.
就是利用我們現在地上的財富.
去轉化成為一種積聚天上財富的方法.
就是周祭人.
所以你又記得.
我們在6章第三至四節的時候.
耶穌教我們三個門徒的操練是什麼呢.
其中一個的屬靈操練.
就是學習施捨的功課.

$^{641}$這裡說.
你施捨的時候.
不要讓左手知道右手所做的.
好像你隱蔽地施捨.
當你去隱蔽地施捨的時候.
原來你就會得到賞賜.
得到負在暗中的賞賜.
也就是說你在天上的戶口就多了.
你天上的戶口就會多了.
因為天父就會幫你加了你的信用.
所以原來這裡就說一個狀況.
就是原來耶穌教導我們.
處理我們地上的錢財的時候.
當然我們錢財可以用來自己.
去用一些生活上的需要.
但是耶穌更加提醒我們.
我們要怎樣去將我們那個.
叫做投資的Portfolio.
就是投資組合整理一下.
我們今天很多人將我們的投資組合.
就是用100\%都是在地上裡面.
讓自己有更豐富.
有更好的享受.
更好的生活.
但是我們會不會將投資組合轉變一下.
讓我們真的能夠將我們.
天上那個儲存的財寶.
比重越來越增多呢.
所以這個第一個耶穌教我們的法則.
我們就是要儲存.
儲蓄.
我們的儲蓄裡面.
不只是儲蓄地上的.
而是能夠儲蓄天上的財寶.
有一個很出名的人.
我不知道你們認不認識.
應該之前早幾個星期.
我都有提及他的名字.
他的名字叫做約翰·韋斯利.
John Wesley.

$^{681}$這個人物.
一個很出名的基督徒.
是一個教會的領袖.
他其中說了.
有一次他有一篇的講道.
叫做The Use of Money.
他說發表了一篇講道的訊息.
就是怎樣去用你的錢.
他當中提出了三個.
很重要的原則.
就是他成為他應用.
怎樣處理地上錢財的原則.
這個原則就是這樣.
第一個是Earn all you can.
第二個是Save all you can.
第三個是Give all you can.
賺錢是為了令到有能力去付出.
我們儲錢是為了我們有空間去付出.
奉獻就是為了讓我們能夠完成管家的使命.
試下簡單的說一說.
他用了儲蓄的法則.
或者說怎樣去面對錢財的法則.
他說道.
為何要Earn all you can.
怎樣去賺你能賺到的.
即是說原來基督徒是可以賺錢的.
今天有時我們又去到另一個極端.
就是基督徒不應該賺錢.
不應該貪好世界的錢財.
所以就應該不用那麼勤力.
我們又搞錯了.
這裡說Earn all you can.
你可以盡力去賺錢.
盡力去賺錢不是因為你有很多很多用不盡的財富.
而是幫助我們透過賺錢.
讓我們有能力去應付我們自己生活的需要.
而不需要去依靠或者過度依賴別人.
聽得明白嗎.
你可以去工作賺錢.
令你有能力去維持自己的生活.

$^{721}$而不需要過度依賴別人.
因為今時今日有些人就慣了.
靠政府的救濟.
不知道這裡是不是這樣.
我知道我們香港也有些人是這樣的.
就是生多點.
生多點小朋友有什麼好處.
政府的資助就大.
如果你已經在拿取綜援.
我們叫做社會政府資助的話.
你生得越多資助越多.
你的子女讀書就不用給多點錢.
原來那些人的想法就是.
他能夠藉著這樣的方法.
不用工作就繼續拿政府多點的資源去工作.
我在香港也見過這樣的情況.
但這裡耶穌說.
John Wesley說這個原則是什麼呢.
就是你賺多點就是因為.
正正因為我足夠養自己的生活.
我就不用靠政府.
又不用靠親戚.
又不用靠什麼方法.
起碼我自己可以應付自己生活的責任.
以致到我不需要.
為了我們的金錢的煩惱.
而令到自己的身體和心智受到損害.
你會知道當你的生活不夠錢用的時候.
你經常都很煩的.
經常都不開心的.
你對你的身體和心靈都受到損害.
甚至乎有時候會令到我們.
進入一個引誘和試探裡面.
不如用一些方法賺快錢.
所以有些人用一些方法去犯法.
就是因為他們不夠了.
想拿多點就拿多點.
他們就用這種方法.
但如果我們真的用一個正常合理的方法.
去賺足夠的錢.

$^{761}$John Wesley說是正確的.
這個就是他提出的.
如何面對我們的生活金錢的方法.
當你賺很多了.
賺很多的時候.
第二個重點是Save all you can.
儲蓄是令到你有空間.
代表什麼呢.
當我們賺到很多的金錢的時候.
就並不等於你揮霍.
聽得明白嗎.
賺很多錢不代表等於你揮霍.
因為很多人賺得越來越多錢.
開始就怎樣呢.
當然是花錢.
不花那麼多就不滿足.
但John Wesley說.
當你賺夠的時候.
你仍然保持著一個簡約的生活.
Simplicity.
這個簡約的生活是什麼呢.
就是什麼是你基本需要的.
你就照買照用.
你需要的東西你仍然可以去使用.
但因為這些錢其實都是上帝給你的.
你就不要浪費了上帝給你的錢.
是為了滿足自己一些奢華的慾望.
一些揮霍的慾望而去買.
原來可以這樣的.
你可以賺最多.
賺了最多之後.
你可以用你應該正常要用的.
接著剩下的就是儲下來.
很簡單.
你想想我們有時候家裡找一找.
我們家裡的手袋有多少個.
可能我們家裡有十個手袋.
但當你看到第十一個手袋在街上的時候.
你覺得都多買一個.
有時候我們都是這樣的.

$^{801}$我自己也有這樣的隱憂.
見到有新手機出.
新手機出沒理由不買新的.
當然說服自己買一個新的.
但想想我這部又用得好.
又用得到.
而且它要做的功能.
基本上我全部都做得到.
我是不是需要買一個新的呢.
有時候我們就說.
我既然賺這麼多錢為什麼不買.
我一代一代的電話追上去.
反而我有錢用.
但如果按著衛斯理的說法.
這些其實不是你真的需要.
當你的電話真的不OK.
那你就換.
但如果你真的只差一代兩代.
其實跟你實際的需要是差很遠的.
所以如果是這樣的話.
你的錢應該是用來儲存.
去儲蓄.
而第三個它說的原則是什麼呢.
就是Give all you can.
盡你所能的去付出.
因為原來你努力去賺很多的錢.
接著用完你生活之後.
很多的錢剩下的拿來儲蓄的時候.
儲蓄來做什麼呢.
不是儲蓄來被蟲咬.
不是儲蓄來去生鏽的.
John Wesley說.
他說你就要將這些錢轉化.
而是去令到.
就是怎樣呢.
去幫到一些有需要的人.
主張這些基本生活所需之後.
剩下來的就要奉獻給.
教會或者窮人或者有需要的人.
去幫助這些人.

$^{841}$他們生活.
真的他們不行的人.
來幫助他們生活.
是能夠過到他們的日子.
所以John Wesley.
他自己的生活實踐是怎樣呢.
他自己就是.
當他在年輕的日子的時候.
在當時在英國.
他的年薪是30鎊的時候.
他當時的生活.
他需要花費28鎊.
他就奉獻了2鎊.
後來到他.
越來越名氣大.
越來越收入高的時候.
他說他每年的收入.
是有1400鎊的時候.
他仍然維持.
每一年的支出是28鎊.
而是將全部的錢捐出來.
這個就是John Wesley.
他自己實際去運用.
怎樣去面對他的錢財.
他將他的投資組合.
越來越增加增加.
在天上面的投資組合.
就正如前面說的.
怎樣能夠將他的錢.
有效地去幫助一些有需要的人.
這個就是.
耶穌第一個說的.
給我們聽過的那種理財法則.
所以今天第一個的小結就是.
你的財寶在哪裡.
你的心也在哪裡.
留意這段聖經裡面.
其實再看第21節.
也有兩個那裡.
所以在第19至21節裡面.

$^{881}$出現了六次那裡.
耶穌很想透過這個重覆.
去強調問.
在哪裡.
就是問你的財寶在哪裡.
你的投資組合在哪裡.
而帶出來一個很重要的思考.
就是你想一想.
你投資得越多的地方.
是不是你越緊張的地方.
是不是應該這樣.
譬如你舉例你的身家裡面.
有九成是買了股票的話.
你很難忍得住.
不每天看一下股市.
是不是.
因為你投資組合越大.
你越緊張它的變化.
或者那個狀況是不是安全.
所以耶穌挑戰我們什麼呢.
如果我們的投資組合.
是擺在天上積蓄財寶的話.
我們就緊張.
整天看著天上的財寶.
OK嗎.
是不是穩當.
我們就不會再這麼緊張地上的財寶.
耶穌就挑戰我們的心.
看清楚我們的投資組合.
的大部分或者全部在哪裡.
如果在天上的話.
你就安心了.
如果你還在地上的話.
你可能會為了地上.
經常擔驚受怕你的財產.
你的財富會損失的時候.
你的心常常在忐忑不安裡度過.
所以這是第一個.
我們去思考的一個法則.
我們怎樣去儲存呢.

$^{921}$儲存在哪裡呢.
第二個和大家去思考.
耶穌接著說的22-23節的段落.
眼睛是身體的燈.
你的眼睛若明亮.
全身就光明.
你的眼睛若分花.
全身就黑暗.
你裡面的光若黑暗了.
那黑暗是何等大呢.
路加福音第11章34-36節.
同樣有這個類似的段落.
說到我們的眼睛是我們身體的燈.
這裡說的是一個重要的問題.
就是耶穌提及眼睛是身體的燈.
因為這個眼睛就是一個.
透過我們的眼睛去看到事物.
也透過這個眼睛.
讓我們能否發出光呢.
是一個很重要的關鍵.
這裡說眼睛所散發的光線的強弱.
代表了我們的視力.
這個視力可以引伸成觀察力.
洞察力,理解力的好壞.
所以剛才我們讀的NIV版本.
就是你的眼睛是否健康.
翻譯就是你的眼睛若是健康的話.
就是這樣.
你的眼睛若是不健康的話.
就是後面的狀況.
不過我想再解釋多一點.
其實原來的文字是什麼意思呢.
我們的眼睛是代表我們身體的燈.
我們的燈究竟能否亮呢.
這裡說有兩個狀況.
第一個就是.
眼睛若是明亮.
全身就光明.
在原來的文字裡其實是怎樣翻譯呢.
所以你的眼睛是單一的.

$^{961}$你的全身就是光明了.
這裡說原來它的翻譯是單一的.
即是你的眼睛是單一.
簡單來說就是專一.
集中的.
另外一個翻譯就是.
直率或者真誠等等的意圖.
即是說原來我們.
將我們的眼睛能夠專注.
我們能夠集中去一件事.
在那件事的時候.
我們就能夠發出那種光明.
這個就是它的說法.
或者我們怎樣能夠看得到對比是怎樣呢.
你的眼睛若是分花圓文是怎樣翻譯呢.
原來的文字就是.
若你的眼睛是邪惡的.
你的全身就是黑暗.
原來後面是說一個很大的對比.
就是一個邪惡的眼睛.
而你的全身是黑暗.
所以這裡是說一個這樣的問題.
就是你的眼睛究竟是單一.
很正直很專一地看著.
令你能夠發光.
還是你的眼睛是充滿邪惡.
而這種邪惡是令你帶來黑暗呢.
即是說我們的眼睛究竟在看甚麼呢.
所以回到我們今天的題目.
今天的題目就是見淺開眼.
我想解釋一下.
見淺開眼這句說話是無效的.
為甚麼無效呢.
因為你想想見淺開眼的時候.
即是說你見淺之前你是開眼還是閉眼呢.
閉眼.
見淺開眼.
即是說見淺之前你是閉眼.
閉眼是怎樣能夠見到淺呢.
即是你閉著眼是不會見到淺的.

$^{1001}$所以這個邏輯上是有問題的.
見淺是不應該能夠開眼.
因為你閉著眼是沒有辦法見到淺.
相反.
見淺是甚麼呢.
是令你閉眼的.
耶穌在這段說的是一個狀況.
因為當你見淺.
而當你見到這些淺的時候.
而令到你出現了.
剛才說你的眼睛出現邪惡.
見淺原來會令到我們眼睛.
有一個作機冒險.
一個風險就是我們會出現邪惡.
而這種邪惡就會將我們全身帶來一個黑暗.
是這樣的.
所以應該是相反的.
見淺是令到你的眼睛閉塞了.
見淺是令到你的眼睛進入黑暗.
見淺是令到你的眼睛進入一個關燈的狀況.
就是這樣.
所以見淺閉眼才是合乎聖經的說法.
因為當你閉眼的時候.
你就會看東西就開始不準確了.
甚至乎要用剛才的字眼就是進入一個邪惡.
其實說的是什麼呢?.
這裡說的是什麼呢?.
就是單一,專一.
什麼叫做邪惡呢?.
就是由一去到多的時候.
由一去到多的時候.
你想要更多.
所以這裡說的是邪惡是什麼呢?.
就是當你見到淺之後.
而進入邪惡.
就是因為你進入一個想要更多的地步.
當你想要更多的時候.
就會出現一個貪婪的狀況.
為了錢什麼都可以做.
為了錢什麼都可以做的時候.

$^{1041}$這就是我們生命裡的危機.
最近大家有沒有看東奧?.
就是冬季奧運會.
有很多項目.
當說到今年東奧的時候.
其中有一單新聞提及兩個人.
那兩個人都是中美混血兒.
中國和美國的混血兒.
那兩個人呢?一個叫谷愛玲.
一個叫劉美妍.
不知道你們有沒有聽過這個新聞.
這兩個人都是屬於中美混血兒的時候.
他們參加了今年東奧的運動會.
當然最後谷愛玲拿了銀牌.
劉美妍拿了團體的金牌.
不過重要的是.
社會的討論不是拿銀牌還是金牌.
而是說他們究竟.
作為一個運動員去參加這個運動的時候.
他們是什麼心態呢?.
當中其中左邊那個就是谷愛玲.
網上的描述就是這樣.
叫做披中國戰袍的叛徒.
什麼意思呢?.
因為谷愛玲本身就是屬於美國國籍的人.
美國國籍.
在美國一路長大一路受訓練.
成為一個在美國生活的人.
當然她有她自己的天才和才華.
以致她在滑雪上有一個很出名的成就.
這個都是她值得擁有的.
因為她用了很多努力.
不過帶來的討論就是.
原來她在2020年的時候.
她代表中國拿了兩金一銀的成績.
她幫中國拿了兩金一銀的成績.
代表什麼呢?.
她為什麼能夠代表中國參加這個奧運呢?.
因為她掛著中國籍.
她唯有你是中國籍的運動員.

$^{1081}$你就可以代表中國去參加比賽.
而她本身是在美國生活的美國人.
也是美國籍的人在這裡長大.
不過她就去了為了中國去參加這個比賽.
而其中根據一些網上的報導.
其實最主要為什麼她會去代表中國參加這個比賽呢?.
因為有很大筆的錢是給她.
包括她的贊助.
包括她自己代表中國能夠拿到金牌的時候.
那些獎賞.
本來也不是一個問題.
你甘願成為中國籍其實也不是一個問題.
不過其中她說了一句話.
最令人不舒服的是什麼呢?.
因為有些人問她.
你現在是不是已經成為了中國籍的人了?.
因為你代表中國參選.
你是不是已經脫離了美國籍了?.
她的回應就是說.
我在美國的時候我就是美國人.
我在中國的時候我就是中國人.
她沒有回答她是什麼國籍.
她只是說當我去代表中國比賽的時候.
我就是中國人.
當我回到美國的時候.
我就是美國人.
為了得到很龐大的利益.
她就選擇了可以一時是中國人.
一時是美國人.
一時中國籍.
一時美國籍.
所以那些人為什麼不喜歡.
或者要提出一些疑問呢?.
就是因為這個問題.
原來錢可以令到一個人.
自己的國籍.
自己的身份都不知不覺沒有了.
基督徒也一樣.
基督徒可以為了錢.
有時候也會失去了我們基督徒的身份.

$^{1121}$在教會我就是基督徒.
出去工作.
我就是什麼?.
我就是什麼人呢?.
你想一下.
原來可以這樣的.
而右邊這位少女就叫劉美妍.
劉美妍她是參加滑冰比賽的.
她最近能夠帶來美國隊.
美國隊能夠拿到滑冰的冠軍.
拿了金牌.
其中為什麼報紙報道.
會將這兩個人並排呢?.
因為原來劉美妍.
她也被中國邀請過代表中國隊.
原來有很多的錢.
邀請她甚至網上的報道說.
比谷愛玲更多更高的錢.
是希望她成為中國的代表.
去參加奧運.
而結果她的家人拒絕了.
她爸爸就教她女兒拒絕.
因為她爸爸說我們來到美國.
為的是什麼呢?.
為了要有自由的選擇.
為了要為我們的理想而去做選擇.
為我們的理想去做一個自由的決定.
就是他們來到美國的原因.
所以今天你作為一個運動員.
你可以做自己的選擇的時候.
就不應該是為了這些錢.
去做你生命的選擇.
所以結果劉美妍就選擇.
繼續代表美國.
一個培育她能夠在滑冰之路.
能夠有一個傑出成就的美國.
去參與這個奧運比賽.
亭姐妹原來我們都是這樣.
我們今天挑戰我們的就是.
帶著什麼樣的眼睛去看事物.

$^{1161}$我們會不會被一些東西去遮蔽了我們的眼睛.
單一的眼睛是一種慷慨的眼睛.
因為只是能夠看到神.
在當中我們願意去為人去付出.
邪惡的眼睛就代表一個貪婪的眼睛.
貪婪的眼睛令到我們只是為錢而活.
今天第二個小結就是問.
究竟我們全身是光明.
還是全身是黑暗呢?.
第三個法則就叫做服事法則.
服事法則裡面說的是.
一個人不能夠服事兩個主.
他不是恨這個就愛那個.
就是中這個就輕那個.
不能又服事上帝又服事馬門.
第三個耶穌說到的法則就是.
究竟我們用什麼服事的心態去服事呢?.
我們對錢財原來都為我們帶來一個挑戰.
就是究竟它是一個怎樣的物質呢?.
這裡說的是一個是服事神.
一個是服事馬門.
馬門的意思就是錢財的意思.
所以耶穌就說你有兩個主人可以選擇.
一個就是神一個就是馬門.
馬門是一個怎樣的主人呢?.
有些圖片是這樣說的.
都幾常見到這句話.
Money is a good servant but a bad master.
錢財是一個好的僕人.
但是它是一個壞的主人.
錢財是一個好的僕人.
因為你適當去用這些錢.
適當去使用你的僕人.
可以令到你很多東西商得益彰.
還有更好的狀況和表現.
但是如果你讓錢財成為你的主人.
它就不會是一個好的主人.
因為這個主人會怎樣呢?.
就好像右邊那幅圖一樣.
當你願意以錢財作為你的主人的話.

$^{1201}$你就會被錢財去操控.
操控著你的生命.
以至到你為錢財釘在錢財的十字架上.
就是它會奴役你一生.
你為了錢財去努力.
努力賺很多很多.
奴役你一生之後.
到你死那一刻.
你看一看.
你有十個零.
你有二十個零又如何呢?.
又怎樣呢?.
因為你這些東西帶不上天.
多少個零都是沒有用的.
它不能夠豐富你的生命.
不能夠讓你生命裡面能夠有成長.
所以這裡說.
如果馬門是你的主人.
它不是一個好的主人.
哪一個是一個好的主人呢?.
這裡說.
當我們以耶穌作為你的主人的時候.
祂會是一個好的主人.
因為你的主人對你怎樣呢?.
這裡說這個主人會為我們付上重價.
因為你知道在古代來說.
就是要買一個僕人回來成為你的僕人.
其實你的主人要付上代價.
而今天主耶穌就為了我們的原因.
祂願意甘願為我們付上一個重價.
去買我們回來.
去買贖我們回來.
所以這個主人是很愛我們的.
耶穌是一個很愛我們的主人.
所以我們能不能夠讓耶穌成為主人呢?.
旁邊那句說話說什麼呢?.
當耶穌用了重價去買我們回來的時候.
祂要讓我們生命有改變.
讓我們成為什麼呢?.
God's masterpiece.

$^{1241}$是Master的piece.
成為這個主人的一個傑作.
原來當你選擇了耶穌.
作為你的生命的主人的話.
耶穌要讓你的生命成為一個傑作.
是由這個Master所造成的一個碎件.
這個人生命裡面的其中一個碎件.
但這個碎件是一件傑作.
耶穌就是這樣讓我們經歷一個選擇.
就是究竟我們選擇耶穌這麼好的主人.
為我們捨棄生命付上重價的主人.
還是選擇一個不會做主人的錢財做你的主人呢?.
而很有趣的是我們往往希望如雨紅掌.
兩者兼得.
所以耶穌挑戰我們一個問題.
就是我們往往如雨紅掌.
想兩個都可以.
又用主耶穌作為我們生命的主人.
又用錢財作為我們生命的主人.
但這裡告訴你.
當我們以主耶穌為我們的主.
又以馬門作為我們的主的時候.
其實你是以自己作為你的主.
因為你的重要的目的.
你是想將耶穌和馬門成為服侍你的兩個僕人.
要搞清楚啊弟兄姊妹.
當你兩樣都想得到的時候.
你只是想利用主耶穌.
利用錢財去讓自己成為那個主人.
我們見得到今天有一套新的戲.
叫尋秦記.
尋秦記大家都很想看.
希望有得看.
尋秦記裡面講什麼呢.
其實是講那個王.
秦始皇.
秦始皇之前.
未上位之前.
就是秦始皇.
而在秦始皇的時候.

$^{1281}$他就要成為那位的王.
他當時.
我記得在很多年前.
就不是這套尋秦記.
是很多很多年前.
有一個人叫羅家良.
唱過一首歌叫秦始皇.
不知道你有沒有印象.
師母當然沒有印象.
你不是那個年代.
那首歌說.
大地在我腳下.
國計長於手中.
哪個再敢多說話.
有沒有印象.
Wesley是同我同我年代.
啞得到頭.
就是說.
秦始皇是想做什麼呢.
在不是一人之下.
而是眾人之上.
所以呢.
什麼都在他的腳下.
什麼都在他的手裡.
原來呢.
我們很多人都是這樣.
以耶穌.
都成為我的腳下.
以錢財成為我的腳下.
今天我們不是尋秦.
我們不需要尋秦始皇.
我們今天要尋主記.
我們想想.
尋覓我們的主.
誰是我們的主呢.
這段經文24字挑戰我們.
誰是我們的主呢.
我們要專心去服侍誰呢.
這個打錯字.
服侍誰呢.

$^{1321}$我們要專心服侍誰呢.
這裡其實有一個很重要的挑戰.
就是我們生命的優先次序.
如何排序.
剛才說.
當我們心態.
秦始皇的心態就是.
兩個在上面.
兩個在下面.
如何排我們生命的秩序呢.
最近我看到.
有另一單的報導.
有一間公司在美國.
我也沒吃過這間餐廳.
不知道佛羅里達有沒有.
這間餐廳叫In-N-Out.
沒有吧.
這間餐廳叫In-N-Out.
原來呢.
最近被人投訴.
什麼呢.
賣出來的食物裡面.
那些食物盒上面.
有聖經的經文.
可以看到.
這裡中間那幅圖裡面.
應該是箴言.
箴言24章16節.
就是他們印出.
他們裝個包的那些盒.
或者那些飲品.
那些外面.
都會印上一些經文.
這個就是他們公司裡面.
成立的傳統.
而面對的就是.
有些人就說.
我不喜歡有印經文在食物裡面.
我是要來買食物的.
我不要有聖經的金句.

$^{1361}$在你的食物裡面.
於是很多顧客就嘗試.
去批評他們.
要求他們取消這個規定.
而這個公司的負責人.
這個Slider.
他很清楚.
他們說公司.
這個就是他們的.
家族的一個做法.
這個家族.
是一個信耶穌的家族.
所以他們認定了.
透過他們的公司.
透過他們的生意.
要表明了他們的信仰.
所以最後他們就說.
你們無論怎樣說.
都好 明天投訴.
不會因為你們的說話而屈服.
所以他說.
這些經文仍然會繼續.
在我們的食物裡面.
他說.
他會選擇.
Stand firm.
選擇堅持.
這個世界就告訴你.
那些就是你的顧客.
你願不願意.
為這些顧客去屈服.
你自己 因為當你願意屈服.
你自己的話 你可能.
就會賺更多的錢.
但是你願不願意為了這些錢.
去成為錢的.
僕人呢.
這個見證.
很鼓勵我們.
告訴我們原來.

$^{1401}$這麼大的企業 他們都可以.
選擇去堅持信仰.
不是服事主.
那我們又怎樣選擇呢.
這個服事法則裡面.
其實最後帶來的就是一個.
排序的問題 怎樣排序.
才能夠好的.
服事法則呢 應該是這樣 我的理解.
就是.
我們作為神的僕人.
我們知道我們.
的身份 我們願意.
服事主耶穌.
主耶穌在我們上面.
我們願意服事主耶穌.
聽得明白嗎.
服事主耶穌 而錢財呢.
是應該被我們.
好好運用的對象.
所以我們不會.
在錢下面.
我們不會服事錢.
應該錢是服在我們之下.
是被我們合當的.
運用 這樣排序.
我們作為神的僕人.
我們應該好好用主人.
交給我們錢的託付.
以致我們成為又忠心.
又良善的僕人.
這個才是我們用錢的方法.
用這些錢合意的時候.
以致我們這些錢.
讓我們生活得到滿足.
我們生活得到供應.
但同時間也讓有需要的人.
能夠得到幫助.
這個就是上帝要我們用錢的方法.
以致到最後.

$^{1441}$讓我們更加能夠.
積聚財寶在天上.
弟兄姊妹.
千萬不要讓錢財.
功高蓋主.
當它功高蓋主的時候.
你就會被它奴役一世.
求主幫助我們.
今天我們要思考.
三個很重要的法則.
一個就是要.
積蓄的法則.
積蓄財寶在天上.
第二個就是.
開眼的法則.
我們要保持我們的眼睛是明亮.
專一.
以致我們身體能夠發光.
第三個我們的服事法則.
我們要服事上帝.
而不是服事馬門.
這裡聖經說.
是輕馬門.
不是說我們不用馬門.
我們不是不用錢.
不過錢在我們眼裡.
是我們所用的上帝.
給我們的工具.
求主幫助我們今天去思考.
我們怎樣能夠.
面對我們今天擁有很多錢財.
能夠用得合宜.
用得蒙主喜悅.
用得我們能夠積聚財寶在天上.
我們一起祈禱.
天父我們求你幫助我們.
因為.
我們都很感恩.
因為你給我們生活裡.
有很多的財富.

$^{1481}$以致我們今天生活.
都舒適在一個很好的環境裡.
但求主也讓我們.
不只是積聚地上的財寶.
讓我們合宜去運用.
去使用.
以致我們能夠捐出這些財寶.
去幫助一些有需要的人.
也幫助福音的事工.
求主你帶領我們.
讓我們教會裡都能夠成為一個.
忠心的錢財的管家.
也能夠將榮耀歸給你.
也能夠使我們在天上蒙福.
我們就將我們的禱告.
交在你的面前.
是奉耶穌的名字而求.
阿們.
\newpage



\section{歌羅西書 1:9-14}
\label{sec:031652}
\textbf{如何代禱?如何接受代禱?}~~~~黃紹權牧師
\newline
\newline
連結: \href{http://www.church.com.hk/acms/content.asp?site=cdc\&op=show\&type=product\&code=031652\&layout=sermon}{\texttt{031652}} ~~~~ 語音日期: 2026-02-15
\newline
\newline
\hyperref[sec:031651]{\small{< < < PREV SERMON < < <}}
~
\hyperref[sec:index_chronic]{\small{[返順時目]}}
~
\hyperref[sec:031653]{\small{> > > NEXT SERMON > > >}}
\newline
\newline
歌羅西書 1:9-14
\newline
\begin{longtable}{cl}
\hline
\hline
章節 & 經文 (和合本修訂版)\\
\hline
1:9 & \begin{tabularx}{0.7\textwidth}{X} 因此，我們自從聽見的日子就不住地為你們禱告和祈求，願你們滿有一切屬靈的智慧和悟性，真正知道神的旨意， \end{tabularx} \\ \\ \relax
1:10 & \begin{tabularx}{0.7\textwidth}{X} 好使你們行事為人對得起主，凡事蒙他喜悅，在一切善事上結果子，對神的認識更有長進。 \end{tabularx} \\ \\ \relax
1:11 & \begin{tabularx}{0.7\textwidth}{X} 願你們從他榮耀的權能中，得以在一切事上力上加力，好使你們凡事歡歡喜喜地忍耐寬容， \end{tabularx} \\ \\ \relax
1:12 & \begin{tabularx}{0.7\textwidth}{X} 又感謝父，使你們配與眾聖徒在光明中分享基業。 \end{tabularx} \\ \\ \relax
1:13 & \begin{tabularx}{0.7\textwidth}{X} 他救了我們脫離黑暗的權勢，遷移到他愛子的國度裡。 \end{tabularx} \\ \\ \relax
1:14 & \begin{tabularx}{0.7\textwidth}{X} 藉著他的愛子，我們得蒙救贖，罪得赦免。 \end{tabularx} \\ \\
[1ex]
\hline
\hline
\end{longtable}
$^{1}$福臣慈悲讓我們在聆聽主的說話之前.
我們一起同心低頭禱告 讓我們轉圍一生.
父神多謝你 讓我們再次回到你的啟示當中.
在你獨生兒子耶穌基督所作成的救恩基礎之上.
讓我們能夠可以在當中去生根建造.
你讓你的僕人保羅將祂的說話記下來.
讓祂所透露的對教會每一個信徒.
他們的關懷 他們靈性上的長進.
所付出的心 身 力.
我們都祈求主讓我們一方面可以體會.
但更重要的是我們同樣要去擁抱同一份的心智.
來在上主理的當中去仰望 去事奉.
走一條合神心意的道路.
配得起稱為一個真正的基督徒.
而不是一個虛有其名掛名在這個世界上面.
蒙混過關的一個教徒.
所以祈求主讓我們再一次在你的說話啟示.
再次彰顯的時候.
讓我們每一個能夠再次同樣在你的教導當中.
準確去領受當中的教訓.
並且透過聖靈的引導.
能夠有同一份的心思意念.
好像基督一樣能夠活在上帝的喜悅當中.
求你的幫助帶領我們同心禱告.
奉靠主耶穌基督得勝名自救 阿們.
不是不想休息 很難休息.
因為教會只有一個傳道人.
希望將來可以多一些弟兄姊妹同工.
可以一同事奉在這個港道的崗位上面.
也都很開心.
這些弟兄姊妹就算去旅行.
都原來參加崇拜的.
會記掛著我.
記你的一些問候的卡等等.
今天我們會去看歌羅西書的第一章第九至十四節.
我給一個講題叫如何代禱.
如何接受代禱.
無獨有偶在過去這三個星期裡面.
也都有兩位弟兄寫了一個電郵給我.
談論到禱告的問題和他們的猶豫.

$^{41}$特別他們去討論到怎樣去為人代禱.
自己也都怎樣去代禱 怎樣去祈禱.
是不是好像其他的宗教一樣這樣去禱告.
今天我就不是特別為他們有這條問題.
我相信這個是聖靈的模樣.
在適當的時間就將這個題目擺在我們面前.
的而且確很多的宗教都有所謂禱告的事情.
但是聖經談論到禱告到底是一回甚麼事呢.
很多時候我們作為基督徒在這方面我們很少去探究的.
因為我們就用回其他宗教對於所謂祈禱或者去品神.
這種的思維來套入到去基督教或者我們基督徒的信當中.
這個我是很希望大家請姐妹首先要學會一件事.
切勿將其他宗教的拜神方式套入去我們相信上帝的這個信的過程.
我再次重複這是一個很嚴重的一個問題來的.
因為我見得最多的時候很多的信徒.
都他的禱告和其他的宗教的所謂禱告是沒有分別的.
甚至他們去敬拜相信上帝相信耶穌的思維方式.
都和去敬拜或者思想其他的宗教是沒有分別的.
我見得最多的是什麼信耶穌的禱告不信的耶穌的人又禱告.
他們的禱告的內容往往都是為了自己自身的困難而去禱告的.
但是很可惜因為他們不明白聖經裡面談論禱告的時候就亂禱告.
譬如基督徒很多時候就說哪一個人有什麼困難有什麼病患有什麼難處.
我們就可以帶到上主的庚前然後祈求上帝為他解決.
首先我想告訴大家這種禱告最大的問題的焦點在哪裡呢?.
就是你為自己自身的事情去禱告.
你整個禱告的焦點是在你個人的心緣或者是某些人的心緣.
但是這種禱告從來都不是在聖經裡面談禱告的問題.
或者不是談禱告的焦點.
有些人就會說牧師你怎麼禱告呢?.
我告訴大家牧師會告訴我很多他自己的屬靈上的挑戰或者肉身上的挑戰工作上的挑戰等等的挑戰.
不過我很誠實地告訴大家.
因為按照聖經的教導我從來都不會為你這些挑戰而得到完滿解決來禱告的.
於是有一個傳道人聽到我這樣說之後.
他就有一個這樣的回應.
他說你這個黃兆權牧師正因為是屬靈的孤寒鬼.
那怎麼孤寒啊?.
他說祈禱又不用花你多少錢.
不用錢的是不是?.
你都斤斤計較.
跟別人祈禱希望上帝為他解決這件問題.

$^{81}$甚至你跟他一起禱告.
所以很多傳道人很喜歡聽見信徒有什麼需要就馬上拉著他來禱告.
我不是說這個行為錯.
但是就在跟他一起禱告的過程當中.
就在禱告的內容的時候就說.
上帝你為他解決這個問題.
如果他病了就讓你的私人拯救的手觸摸他的病處為他禱告.
希望他快些康復.
他有癌症請你觸摸他的癌細胞殺死他的癌細胞.
如果他告訴你他失業了.
上帝求你馬上賜一份工作給他.
使他能夠在這份工作中榮耀神.
上帝你能夠得到更多的供應給這個家庭.
使他更有力去事奉.
為什麼傳道人喜歡做這些禱告呢?.
因為這樣的即時禱告很討人的喜悅.
你越病我為你禱告你快些不用病.
你就覺得牧者很關心我了.
如果他禱告真的我沒多久病好了.
他就說上帝的工作了.
但是我告訴大家.
這些禱告其實往往只是為了個人願望去禱告.
亦都是這些假傳道人.
他們欺騙人的其中一種手段.
我會怎樣禱告呢?.
如果你告訴我你病了.
我就會跟你說我們一起禱告吧.
你不是一樣嗎?.
先聽完內容.
我會跟你禱告.
上帝弟兄姊妹現在在病患當中.
求主你在這個病的過程當中.
能夠對生命的掌權更清楚掌握.
不是醫生掌控我們的生命.
不是我個人掌控我們的生命.
因為我們知道掌控生命的是上帝.
上帝你可以使用醫生.
求你在這個醫療的過程中.
讓我學會怎樣向上主你感恩.
怎樣在這個過程中.

$^{121}$更學會好好保養自己的生命.
珍惜有的生命和健康來為主你所使用.
上主倘若我這個病未必會好.
但願在這個過程當中.
讓我餘下的人生.
可以為主你作一個更有效的見證.
讓我也體會到上帝你與我同在的那種掙扎.
我會這樣禱告.
為什麼呢?.
因為如果我們的禱告只是亂來.
只是按照自己的心意而去禱告.
或者甚至代頭的話.
你知不知道很大可能性.
你這種的禱告是一種對上帝的干犯或者干擾.
一次半次我相信上帝都會饒恕的.
但如果你長期用這種思維在上帝裡面.
將一些你明知上帝不會做的事.
或者你根本不知道上帝在過程當中.
有什麼特別的心意.
是透過病患困難帶到你生命當中的話.
你強要上帝滿足你的願望.
其實你真是監臣乃后.
所以弟兄姊妹.
當如果我們去監臣乃后的時候.
那份罪比你不禱告更加的強烈.
我恐防上帝會見到你的禱告.
直接反映你的不信在他面前.
上帝更加討厭.
所以有些人說我不信主都向上帝禱告.
我相信神.
你知不知道上帝聽一個不信主的人的禱告.
只得一種的禱告.
就是認罪悔改的禱告.
因為認罪悔改是他整個信仰的路程的起點.
如果他連這個起點都沒有.
我禱告凡事就已經得到上帝的滿足.
對不起你信的那個不是耶和華神.
你信的那個不是耶穌基督.
你信的那個是其他的宗教.
包括最普遍的其他的偶像.

$^{161}$王大仙 菩薩 觀世音.
那些的神就是滿足你.
因為他們的背後就是魔鬼撒旦.
而魔鬼撒旦就是最喜歡滿足你的需要.
所以當我們去詳細去看到禱告這個題目.
絕對不是一件.
只不過閉上眼.
跪在地上很誠懇地在上帝裡面.
去陳述自己的願望.
禱告有一個很重要的層次的學習.
今天我們就透過這段聖經.
來到一切的思想.
歌羅西書第一章第九至十四節.
有聖經的請大家一起打開.
一起去讀一下這段聖經好不好.
不是很長而已.
請大家一起打開這段聖經去讀一下.
好 我們一起讀.
因此我們自從聽見的日子.
也就為你們不住地禱告祈求.
願你們在一切屬靈的智慧不勝常.
滿足知道神的旨意.
好叫你們行事為人對得起主.
凡事蒙他喜悅.
在一切善事上結果子.
漸漸地多知道神.
照他榮耀的權能.
得以在各樣的力上加力.
好叫你們凡事歡歡喜喜地忍耐寬容.
叫我們能與眾聖徒在光明中同得基業.
他救了我們脫離黑暗的權勢.
把我們遷到他愛子的國裡.
我們在愛子裡得蒙救贖.
罪過得以赦免.
歌羅西書在第一章的第一節至第二節.
就有一連串介紹到.
保羅與哥羅西教會的關係.
然後進入到第三至第八節裡面.
就連續用了三個的猶如.
來指出作為一個信主的人.

$^{201}$他應該的生命是猶如福音的本質.
猶如主耶穌基督一樣.
他們會按照主的道理而去行事.
即或在這個世界當中.
按主使意行事是令人討厭的.
他仍然不會因此而去到屈服.
他會怎樣做呢.
他會對一切的罪.
能夠在這個世界的上面幫助人脫離.
這個是他的使命.
因此他願意直接地向人指出.
他的罪的問題.
這一份是討厭的工作.
因為在這個世界裡面.
沒有人喜歡被人去指出自己的罪.
但是也都同時反映出.
如果一個真正信主的人.
他是願意面對自己的罪.
面對自己的錯.
而不會去抵賴.
當一個人能夠生命似福音的本質.
似耶穌基督.
他們的生命.
就會在第一章三至八節裡面所講到.
他會成長.
即是越來越成熟.
並且他會繼續繁衍.
英文用了兩個字很有意思.
They will keep growing and multiplying.
這兩個字我們也都曾經講過.
是創世記的用詞.
今天我們會再講的.
所以我稍後會和大家再次討論這個問題.
而來到第九節至第十四節.
在這個部分裡面.
我們就看到保羅更進一步告訴他.
怎樣去為到哥羅西教會.
他們這個這麼成熟的生命.
來獻上感恩和禱告.
上文我們已經講過了.

$^{241}$如果大家去看一看.
第一章第三節已經講到.
他能夠有一間哥羅西教會.
或者這班信徒.
他對上帝那份委身而向上帝感恩.
保羅沒有擦人鞋.
他只是向上帝見到信徒的生命像紙.
基督向上帝感恩.
他沒有因為自己和哥羅西教會.
沒有直接的關係.
而希望用這個機會去搭關係.
沒有.
所以在這裡請大家弟兄姊妹留意.
作為基督徒在這個世界.
是我們弟兄姊妹會互相幫助.
但不是因為你幫助了他.
你搭了個關係給他.
之後你就要他回報給你.
沒有啊.
保羅都沒有這樣做.
你能夠去幫到一個弟兄姊妹的.
你能夠幫到一個信徒的.
不是這個給恩典或者幫你的人.
而是向上帝.
在這個大前題之下.
進入到第九至第十四節.
在這個部分我準備分三個部分.
和大家一起更深入去看看.
代禱或者禱告的問題.
一方面我們可以怎樣去代禱.
另一方面我們接受了別人的代禱.
我們應該怎樣去回應.
先從第一部分去看.
第一章第九節.
我們看的就是.
究竟在什麼層面上作出代禱.
一章第九節這樣說.
因此我們自從聽見的日子.
也就為你們不住地禱告祈求.
願你們在一切屬靈的智慧不勝常.

$^{281}$滿心知道神的旨意.
正如先前我說過.
我從來不會只是為弟兄姊妹.
你有的困難.
而希望上帝經過禱告之後.
為你去除困難.
我不會這樣禱告的.
困難是上帝加給我們的.
當然如果你自己招惹的.
這是另一個問題.
但是上帝容許.
或者上帝加在我們身上的困難.
最重要的目的.
是要讓我們去明白.
上帝在我們生命當中要做的事.
是拉近和上帝之間的關係.
所以代禱的目的.
不是說我看到你有什麼困難.
然後為你代禱完事情得到解決.
然後我們一起在教會.
Hallelujah讚美上帝.
或許上帝是要你在這些困難上經歷.
我不敢說哪一個人的見證.
我只是說我自己.
如果大家知道.
早幾年前.
2023年的時候.
差不多COVID完了.
我整個COVID期間都沒有中招.
就是過了COVID之後.
我去探訪之後.
就中招了.
有些人就會覺得.
牧師你真是不幸運.
你知道我那次COVID病的時候.
困在自己的房間裡足足十天.
頭七天發高燒.
我想了很多事情.
我問了上帝很多次.
上帝呀.

$^{321}$為什麼我要在一個這樣的時間.
這麼critical的時間.
才去懷COVID呢.
我要答應了一個姐妹.
為她辦喪禮.
你叫我COVID.
我怎樣出去處理喪禮呢.
我去答應過她的事.
我不能夠不做.
基督徒最起碼一樣東西.
講得出一定做得到.
答應了別人的事一定要做.
就算要付出多少我都要做.
虧損自己都要去做.
因為這個是誠信的問題.
世人不會這樣想的.
哎呀.
算數了.
要我自己付出.
那個時候.
當我最辛苦的時候.
上帝說.
你不病.
你怎會體會到那個人.
尤其是在喪禮裡面.
你要去處理喪禮的那個人.
他的病的掙扎呢.
我讓你的心跟他緊貼一點.
明白他多一點.
我還記得那個姐妹臨死之前.
她問了我一個問題.
她說牧師.
我很害怕.
我是不是得救了.
你知不知道我怎樣回答這位姐妹.
我很明白你怎樣害怕.
但是.
你能夠問這條問題.
你就是在主的裡面.
因為這個是你最期望的事情.

$^{361}$講很容易的就是這個世界.
講很容易的.
你看看你家裡那些.
家裡人怎樣對你.
他講很多東西批評你.
講的東西很容易的.
叫他們做.
上帝說我放你進去這個情況.
因為你的心更加體會到那個掙扎.
我明白上帝的心事.
我靜候我快一點退燒.
不過我就癲癲癲癲的.
吃過量退燒藥.
吃到自己都癲.
不死得真是好彩.
over dosage.
有沒有人一日吃十粒.
太多extra strength.
弟兄姊妹.
禱告從來不是為了個人目的.
而去做的事.
這裡你看看保羅.
自從我聽見.
哥羅西教會你們那些事情.
從哪裡聽回來.
前面我們已經看過有個人物.
他是和保羅一起.
作同工獨人的爾巴弗.
爾巴弗是誰.
爾巴弗就是哥羅西教會本地人.
因為他去了羅馬.
認識了保羅.
然後和保羅身上的傳道信了主.
將那個信約帶回去.
哥羅西這個城市裡面.
打開了哥羅西教會的.
其中一個信徒.
由他的口上得知到.
哥羅西教會的成熟.
保羅聽到.

$^{401}$他不單只是一張三字裡面感恩.
而且他繼續為他們做甚麼.
禱告祈求.
你們覺得奇怪了.
他們都很好.
哥羅西教會已經很棒了.
他們的靈性.
還祈甚麼禱.
有甚麼祈求.
祈甚麼.
是否祈禱他們.
教會建堂更加大.
或者使更多人.
透過教會可以得到興旺.
生意興隆.
尤其是哥羅西這個城市.
是一個商業城市.
交通要塞.
掌握了這個小亞細亞.
亞洲的中心點.
所有交通要經過它.
才從這個希臘世界.
進入到亞洲這個世界.
保羅是否這樣祈禱.
不是喔.
他沒有讓他們更富貴.
保羅的祈禱是甚麼.
願你們在一切屬靈的智慧.
悟性上滿心知道神的旨意.
祈禱第一件事.
尤其是無論你自己個人禱告.
或者你在代禱上面.
我們要祈求的是.
要幫助我們更加成為一個.
有屬靈智慧的人.
而且能夠明白滿心了解.
上帝的旨意.
在希臘文的上面.
這兩個的陳述.
是兩件很重要的事情.

$^{441}$第一件事.
為了你們一切屬靈的智慧.
悟性上禱告祈求.
意思就是說.
為你能夠得到屬靈的智慧.
是談論到一個人的領受能力.
或者簡單一點來說.
他是一個能夠承載更多上帝的智慧.
上帝的教導.
上帝的心思意念的能力.
智慧是指甚麼.
是指你能夠了解上帝的事.
看透事情的事.
我用一個簡單一點的說法.
保羅就是祈求上帝.
為了哥羅西教會的重信徒.
有一個更大的器皿.
可以承載上帝的教訓.
就是祈求給你一個大一點的桶.
這個桶可以幫助你再更多上帝的教訓.
而載了之後.
你都可以明白上帝的教訓是甚麼.
所以第二部份.
就是講到滿心知道神的旨意.
這個說話.
就是談論到一個人實體上.
官能上的層面.
他是可以承載不單止.
而且能夠有領受能力.
這個就談論到一個人的faculty.
即是他的了解能力.
人的了解能力是透過幾種方法去得到.
聽,看,聞,常.
這四方面.
保羅祈求哥羅西教會.
能夠在實體上他們去領受到.
知道上帝給他們的教導是甚麼.
所以換言之.
保羅不單止祈求哥羅西教會.
從上帝裡面給一個更大的承載的器皿.

$^{481}$去承載上帝的教訓之外.
並且要給他們有能力.
載了一些東西之後.
都能夠體會到當中的教訓.
今天很多的基督徒.
是完全沒有這個承載的器皿.
根本不想知道上帝說了甚麼.
所以就不讀聖經.
上完課就不要回教會.
回教會是另一回事.
學會聖經.
這個器皿.
根本是一個上帝的子民應有的態度.
是一個聖情.
如果一個人不想回去認識上帝的話的人.
根本不想要一個大一點的桶.
或者器皿來載上帝的教導的話.
這個人不是基督徒.
所以你不要說牧師.
我回教會30年了.
我回教會20年了.
我就是一個好基督徒.
對不起.
如果你根本沒有興趣對上帝的話.
你的桶根本不想要一個大一點的.
說實際一點.
我們York Region.
最近有一個新的計劃.
1月1號起.
每人派一個大的環保桶.
那個環保桶可以載你和我下去的.
每個人都說不要.
他們說不要因為說阻擋.
我明白的.
有些人家的屋子這麼小.
那個桶其實我都在掙扎擺在哪裡.
我擺在哪裡都被人說那個桶很噁心.
他們叫我擺在哪裡.
我擺在哪裡都被人說.
但是你沒有那麼大的桶.

$^{521}$你每個星期的回收.
你就好像打樁一樣.
打在那個小桶裡.
一個不夠就兩個.
最慘的就是.
打了下來之後.
轉頭有些人就像搭羅漢一樣.
放在上面.
風一吹就吹到周到.
他們不想要這個器皿.
再多一點都不想.
載了下去之後.
怎麼樣.
上帝的說話載了之後.
載了就看了.
這是第二類的信徒.
是呀 上教會 上竹主 上學.
讀好聖經 同學牧師.
寫好筆記.
沒問題.
我很欣賞這樣的信徒.
但是不了解.
他們的屬靈.
只能夠存留在一本筆記簿.
或者一個手提電話裡面.
沒有體會的.
保羅祈求的哥羅西教會.
是你們已經好了.
但是你需要定時.
有一個新的.
換一個大一點的桶.
來去承載更多上帝的說話.
承載了之後.
要去了解.
這個就是上帝的功夫.
這個就是禱告的焦點.
而且焦點是.
滿心知道神的旨意.
不是你的旨意.
不是我的旨意.

$^{561}$不是牧師的旨意.
不是教會長執的旨意.
是上帝的旨意.
還有一個特點.
這兩個陳述.
在希臘文裡面.
都用了一個叫作現在進行式的陳述.
每逢在文法的上面.
用現在進行式.
就指到這件事是恆常這樣做.
原來保羅縱使聽見.
哥羅西教會.
是已經是一間這麼成熟的教會.
但是他為他去禱告.
不是蜻蜓點水講講就算數.
而是恆常地為他們禱告.
弟兄姊妹.
從這一點裡面我們去了解到.
禱告從來不是為個人目的.
不要忘記保羅寫這本書信的時候.
保羅是什麼狀態.
保羅是被囚禁當中的狀態.
他在這個禱告裡面.
沒有為到自己解脫.
這些他自身的困難而去禱告.
相反的我們見到他的禱告.
從來不是為解決問題而成為焦點.
相反的.
他是來幫助一間已經成熟的教會.
能夠更進一步.
在屬靈的承載力和領袖力上面.
更加增強.
弟兄姊妹.
我們在為人代禱或者自己禱告的時候.
這個是必須要的焦點.
或者可以說是一個定律.
如果你禱告只是為個人的事情去禱告.
我可以很肯定給你聽.
你的靈性.
仍然是存留在其他的宗教信仰當中.

$^{601}$你如果這樣繼續去禱告.
只是去得罪上帝.
真正的信徒的禱告.
是祈求上主.
在我看到或者遭遇到的各樣挑戰之下.
使到我能夠有一個準確的角度和悟性.
去詮釋上帝你在當中的意義.
而在禱告的過程當中.
祈求上帝與我們一起去經歷和同行.
這樣的信仰.
這樣的禱告.
就是有根有機的禱告.
你的代禱.
也是真真正正為到那個人的屬靈的需要.
讓他明白上帝的旨意作出的.
如果你只是要求事情得到解決的話.
我可以告訴你.
你信錯教了.
你信其他的教吧.
你信天主教也可以.
天主教也是為你去解決問題的.
他怎樣解決?.
用錢.
教廷的力量.
所以為什麼那麼多人喜歡去大教會.
密基密基.
密宣密宣.
或者浸密浸.
為什麼?.
他用錢來收買解決問題.
而不是在你的靈性上面栽培.
如果你的信仰.
只是仍然存留在這個層面.
我請你不需要在二馬衛環浸信會.
你只需要去這些教會裡面.
他們滿足到你的.
但是.
我要告訴你.
他不會帶你去永恆的當中.
他只會帶你去死地裡面.

$^{641}$當然在那些教會裡面.
仍然有真的信徒.
我認識的.
但是他在教會裡面很大掙扎.
祈禱其實是什麼事呢?.
我用一個比喻.
祈禱就好像一對情侶.
有一天拍拍下拖.
想想你和你老公.
你女朋友.
我們這裡不要有女朋友.
除了那些未結婚的.
你有一天和他一起拍拖.
行下山.
行下Trail.
行下的時候.
突然間下大雨.
又沒有帶雨傘.
在這個時候怎樣呢?.
馬上.
那個男孩就牽著女朋友的手.
然後和她一起快步奔跑.
遠遠看到那裡有一個避雨亭.
馬上走進去那個避雨亭裡面.
當你去到那個避雨亭的時候.
兩個人淋到全身濕了.
好像一隻落湯雞.
兩隻落湯雞一樣.
但是.
當大家在這個落湯雞的時候.
在這個避雨亭當中互相對望.
甜在心頭.
為什麼呢?.
因為你和你最愛那個人.
最信任那個人.
一起走過一條.
一起奔跑的路程.
落湯雞雖然可能會感冒.
淋濕過身可能很麻煩.
又不知道要等多久才停雨.

$^{681}$但是那份彼此的交心.
當他牽著你的手走去這個避雨亭的時候.
那份的甜蜜和親近.
就是你們的關係更進一步的建立的時刻.
祈禱或者代禱.
就是為了弟兄姊妹能夠和神的關係.
更加相交認識明白.
並且一起走這條道路.
所以我們要為什麼去代禱.
或者為什麼去祈求.
沒錯.
可能你遇見一些困難.
但是你縱使遇見這些困難的時候.
你祈禱的.
上帝在這件事的過程當中.
我有沒有更大的器皿.
明白你的教導.
勝在你的教導.
並且在當中明白主理的教導的意思.
看見你怎樣和我一起同行.
所以禱告不是只說當被給.
不是說一下就算數.
也不是期望阿彩有沒有得到成就.
而是一種和上帝建立關係.
明白主心意的一個.
很重要的屬靈功課.
好了.
當我們明白這個重點之後.
從第十節至第十一節.
普羅在這裡更進一步告訴我們.
代禱的目的.
或者代禱的起始點.
是以匹配主所喜悅的事為起點.
我們看看好不好.
我們一起讀好不好.
從第十節至第十一節.
很多人讀到這段聖經.
只會四個字.
力上加力.
禱告就是要力上加力.

$^{721}$最好上帝給我多點力.
我就天下無敵.
但是原來禱告的起始點.
剛才我們說了原則.
現在我們說起始點.
就是以匹配主所喜悅的事為起始點.
匹配主的喜悅的事為起始點.
聖經當中.
我剛才已經說過.
對一個未信主的人來算.
他的代禱.
只有一種代禱上帝會聽的.
就是他認罪悔改的禱告.
所以如果一個人說.
我信主是因為我祈禱上帝.
應驗了我的心思意念.
然後我相信他.
我告訴你.
他根本沒有信主.
因為上帝不會聽那些禱告.
有些人說.
我怎麼信耶穌.
有一次我因為遇見交通意外.
已經血流成河在我身上.
我跟著禱告.
上帝說你快點拯救我.
讓我快點康復.
讓我脫離這個險境.
突然之間救護車出現了.
消防員出現了.
鋸開車門.
把我拉出來.
上帝應驗了我的禱告.
我信主.
兄姊們.
如果你是這樣信主的.
我邀請你重新想清楚.
因為拯救你從車禍當中.
不是一定上帝.
而是你交了稅.

$^{761}$所以有消防員仍然可以當值.
來拯救你.
你說不是.
牧師你真的很狹窄.
那時候有個路人走過.
剛好他有手提電話.
看到車禍.
他打電話.
這不就是上帝的工作嗎.
為什麼你認為這是上帝的工作.
這不是一個人正常的工作嗎.
不要每件事都說一定是神做.
牧師你這樣說.
不是萬事都互相效力嗎.
萬事互相效力.
但是很多事情上帝用的時候.
不是你所指的這種情況.
一個公民的責任拯救了你.
你可不可能是這樣.
但是無論怎樣都好.
你都沒有脫離過你的罪性.
你只不過給一次的好處.
來收買了你自己.
不過你說這個是上帝.
所以很多人的信仰.
接著長期苛索上帝.
要求上帝給好處他.
來收買我們去信他.
上帝不需要的.
我講得俗一點這樣說.
你不認罪悔改.
就說我信了耶穌的話.
我用其他方法去信的.
上帝不稀罕你這種所謂信的.
從這段聖經第十節至第十一節裡面.
我們清楚看見.
因為神他在我們做任何事上.
經過禱告了解上帝心思意念.
有理解明白上帝的教訓的人.
他接著做事是對得起主.

$^{801}$而且這件事做出來.
是主所喜悅的事情來的.
從來一個人信主.
從來一個人禱告.
不是你祈禱之後.
接著我對得起我自己.
不是.
對得起主.
什麼叫做對得起呢.
很有意思這個字眼.
在希臘文裡面.
這個字叫做Asios.
可以講得叫做.
另一個文雅一點的翻譯叫做匹配.
如果你是讀會計的.
你就會知道有一樣東西叫做.
Matching principle.
做每一樣東西要對稱的.
你說是嗎.
只有會計才有.
不是呀.
做很多事都要有的.
譬如最簡單.
醫療.
醫療有個水平的.
早兩個禮拜因為我病得很嚴重.
去了急症室.
去急症室之後.
我從來沒有做過這麼多的檢測.
做了兩次CT scan.
一次X-ray.
接著驗血.
驗了六筒血.
有些人就說.
你看看多好.
其實不是好不好的問題.
這是一個Standard procedure.
他做每一樣東西.
他做一個醫療.
在急症他一定要做這樣的檢測.

$^{841}$好了 做完檢測之後.
他說我們還是找不到你什麼問題.
為什麼會封了喉嚨.
不過.
我馬上給你兩種藥吃.
第一.
他給了兩粒半類固醇給我吃.
其實我是不想吃的.
不過.
他給我就吃了.
然後再噴那些類固醇的puff.
因為我封了喉嚨.
他說那些可以撐開你的喉嚨.
丁姐妹.
過了急症室之後.
我去過醫生那裡.
家庭醫生.
我和那個家庭醫生說.
他沒有其他的治療方法給我了.
為什麼呢.
我說有沒有一些根治到的問題.
為什麼他不馬上處理我的情況.
或者可能我的喉嚨.
是不是正如你們一路所檢測到的.
我喉嚨的氣喉是比較狹窄呢.
需不需要做手術呢.
等等這些所有的問題.
我問我的家庭醫生.
為什麼急症室不做.
你知道我的醫生很有智慧地回答我嗎.
他說因為他是急症室.
急症室只要你不死就可以了.
他不是做根本治療.
只要你在危急的時候.
解決了你不用死就是這麼多東西.
你封了喉.
他就使你不封喉.
所以他要給你類固醇.
要給你puff去噴.
我後來才發現那些puff.

$^{881}$不應該每次每一瓶噴.
他給了兩瓶.
很慷慨啊.
醫院給你藥.
加拿大要買的那些藥.
一噴每瓶噴四次.
每二十分鐘噴一回.
總共噴三回.
你想一下我在一個半小時.
一個小時之內噴了多少次.
四三一十二.
每一瓶噴了十二次.
總共噴了二十四次.
正常這些puff是每次噴一次.
每個小時噴一次.
為什麼他要噴得這麼頻密.
急救嘛.
matching.
對應.
匹配.
你是要.
你就快呼吸不了.
所以我就要急救.
用這個情況去對應.
如果再對應不了.
他就插喉了.
但是什麼長期插喉.
不是長期插喉的.
你不再有呼吸問題.
拔喉的了.
是不是.
弟兄姊妹.
回來這裡看.
保羅的禱告.
要他們有悟性.
有智慧的悟性.
又要明白上帝的旨意.
是做什麼的.
要使得他們能夠.
繼續行事為人.

$^{921}$對應得到上帝.
所以一個真正的信徒.
其實很簡單認的.
你不用告訴我.
你是一個基督徒.
還是不是基督徒.
我在平時生活上看到.
你是不是一個基督徒.
還是不是基督徒.
你在教會可能對牧師很好.
對弟兄姊妹很溫柔.
說話都陰聲細氣.
可能你一回到辦公室的時候.
不停地罵人.
無理地罵人.
你做人做事.
能否在了解了上帝的智慧之後.
做事對得起主呢.
還有一個審判基礎.
不是你.
而是主的標準.
是上帝所喜悅的標準.
所以接下來.
為了要去給.
哥羅西教會的信徒明白到.
他們所作的事.
能夠真正是否.
合神心意主所喜悅.
對得起主.
於是乎有三個的分詞.
接著就從第十節的下半節.
開始這樣說.
他說在一切善事上.
結果指.
漸漸地多知道神.
照著他榮耀的權能.
得以在各樣的力上加力.
好叫你們凡事歡歡喜地.
忍耐寬容.
哪三個的分詞.

$^{961}$第一個分詞.
就是讓你們在善事上.
結果指漸漸地多知道神.
正確的翻譯是.
在一切事上.
不斷結果指.
我重複再說一次.
是不斷地結果指.
漸漸使你知道神.
你不斷有在生命上.
能夠去成熟起來.
而成熟起來.
是指上帝的質素.
不是指這個世界.
是指上帝的質素.
而在這個認識的過程當中.
這個漸漸地多知道神.
就是第二個的分詞.
這個漸漸地多知道這個字眼.
原來原文的字眼就叫.
不斷繁衍.
英文翻譯就叫Multiplying.
原來第一個結果指.
不斷地去.
Making fruit.
或者Bearing fruit.
和在這裡漸漸地多知道神.
不斷Multiplying.
和剛剛我們看完.
第三至第八節裡面所說的.
福音的本質是一模一樣.
多結果指.
不斷繁衍.
第三個的分詞.
就在第十一節.
十一節說什麼呢?.
召他榮耀的權能.
得以在各樣的力上加力.
是在上主的榮耀的權能裡面.
不是你和我得到什麼光榮.

$^{1001}$在上主的光榮能力之下.
這個能力越來越發展起來.
所以這裡所說.
得以在各樣的力上加力.
不是你有力上加力.
而是上帝的權能.
在各樣事情上面.
越來越有影響力.
在這三個的分詞之下.
造成了一個結果.
好叫你們凡事歡歡喜喜地忍耐寬容.
使得你們做一些事情.
可以歡歡喜喜地忍耐寬容.
什麼叫做歡歡喜喜地忍耐寬容呢?.
這個就是上帝的屬性.
弟兄姊妹.
你知不知道上帝忍了我們犯罪多久?.
忍了多久?.
由創世記的第三章引導我們現在.
不是的.
我只有二十三歲.
忍了我二十三年.
你當自己二十三歲.
上帝在你未出世之前已經知道你存在.
上帝忍了你很久.
但是忍你做什麼呢?.
是期望你會改變那一天.
哥羅西教會保羅為他們的禱告.
就是要讓他們能夠在禱告之下.
在代禱之下.
讓他們能夠體會到上帝這份心懷意念.
去活出一個匹配上主喜悅的生命.
所以從這段聖經裡面.
很清楚讓我們看見.
其實我們要恃的是恃福音本質.
恃耶穌基督.
恃回創世記裡面.
上帝原本做人的樣式.
什麼樣式?.
上帝創造這個世界是為了什麼而創造的?.

$^{1041}$有一句說話經常出現.
你看看創世記第一章第二章.
你洗到他怎麼樣?.
生養眾多片滿地面.
不過很可惜很多人只懂得自理者地.
你都沒有上帝的基礎.
你走來自理就搞砸這個世界.
所以就搞到現在BC省這樣.
一殺殺.
現在好像說九個還是十個.
我不是很記得.
最初是十一個.
現在又說九個.
後來十個.
殺足.
就是因為我們的政府.
我們的政府縱容人們隨便做變性.
牧師你又有歧視.
歧視變性.
不是.
你知不知道變性的過程當中.
影響人的思想有多嚴重.
人的思想很容易被扭曲.
被煽動.
失控.
我們政府討人喜悅.
最重要討你的喜悅.
你喜歡變性就變性.
現在搞出一單這樣的事.
陸續有來.
因為它不是做上主所喜悅的事.
當你不是做上主喜悅的事的時候.
這些的禱告是多餘的.
但是如果你的禱告是能夠幫助信徒.
能夠更加的活得像上帝.
理解得更像上帝.
好像上帝一樣有歡歡喜利.
忍耐寬容.
等我們回轉的這一天.
換言之.

$^{1081}$保羅的禱告.
就是為了哥羅西教會可以活得像.
我們很熟悉的一句話.
以基督的心為心.
肥勒比書第二章第五節.
當我們的生命能夠活得越像基督.
以祂的心為心.
好像福音本質一樣的話.
弟兄姊妹.
這個就是我們在.
代禱或者禱告上的起始點.
要活得匹配主所喜悅.
到最後尾第十二至十四節.
這個身份叫做.
「聖徒產業」的身份.
請大家留心我讀出.
叫我們能與眾聖徒在光明中同得基業.
祂救了我們脫離黑暗的權勢.
把我們遷到祂愛子的國裡.
我們在愛子裡得蒙救贖.
罪過得以赦免.
接受代禱之後的哥羅西教會.
應該作出一個怎樣的回應好呢?.
保羅在此再次以個人見證來反映出.
要向上帝獻上感恩.
雖然保羅是作出代禱的人.
哥羅西教會是受代禱的受眾.
但是保羅以自己的經歷告訴他們.
你們要向父神獻上感恩.
為什麼呢?.
因為我們能夠經歷過這一連串的改變之後.
我們可以得到與眾聖徒光明中同得的基業.
很多時候說到基業這個問題.
Inheritance.
有些人就說.
哎呀,所以信耶穌很有福啊.
有些什麼福呢?.
不說清楚了.
什麼福啊?.
那個福是讓我們能夠在光明當中的.

$^{1121}$成為一個聖徒的身份的.
光明當中.
這個光明就直接指到耶穌基督.
如果大家去看《約翰福音》第一章第一節第二節.
你就會知道.
光其實在聖經裡面很連貫地指到上帝或者耶穌基督.
上帝讓我們成為在耶穌基督的基業裡面的其中一個聖物.
就是讓我們能夠得到一個身份.
這個身份就好像移民了一個國家裡面.
你為那個國家能夠給你能夠認同的.
歸入其中一份的一份喜悅.
而感到光榮.
這一點也都反映到移民教會的一個問題.
很多信徒其實移民去外國.
其實從來不覺得那個國家是他自己的國家.
他覺得去那裡只不過是去掠取那個國家裡面給你一些好處.
如果大家記得早幾十年前那些人移民.
經常移民是為了什麼呢.
那裡有好的醫療.
有本護照.
讓我方便些去旅行.
不用拿簽證.
不過我告訴你你去不了俄羅斯的.
你拿著加拿大護照.
去哪裡呢.
拿著加拿大的護照.
去那些享受的國家.
為了這些東西.
他從來沒有覺得是一種光榮.
因為他根本上不投入一個國家裡面.
保羅在這裡告訴我們.
我們能夠享受到一個成為耶穌基督的整個國家裡面.
或者國度裡面的一份子.
你有沒有以此為榮呢.
而這份榮耀是怎樣得來的.
看清楚第十三節.
是他救我們脫離黑暗權勢.
把我們遷到他的國家裡面.
所以我們能夠蒙受愛子的救贖.
罪得赦免.

$^{1161}$在這裡強調了幾點.
第一點.
原來整個過程不只是耶穌基督一個人的做事.
是父神在做事.
父神有一個心願要將人拯救脫離罪.
而耶穌基督就有一份體貼上帝父神的心思.
甚至以死亦都願意去成就父神可以拯救人的目的.
而當耶穌肯這樣做的時候.
上帝父神亦都給他建立一個由他所建立的國度.
讓我們進入耶穌基督這個國度裡面.
而聖靈就在這個過程當中.
光照我們使我們能夠進入罪得赦免的光景.
維繫我們在這個罪得赦免的狀態.
所以從這個過程裡面給我們充分看到.
聖父聖子聖靈其實在幫助著哥羅西教會.
怎樣成為一個更加合神心意的教會.
而保羅為這件事繼續懇切禱告.
持續去禱告.
而希望他做到的.
能夠明白上帝的心思.
為什麼要明白上帝的心思.
之後會繼續詳細講.
但是你可以自己看.
如果你說牧師我看不明白.
你可以等我講到.
但是在這個過程當中.
上帝和耶穌基督彼此的配搭.
最後讓我們脫離罪惡.
聖經從來都沒有離開過.
聖信是脫離罪惡這個焦點.
只不過今天很多教會的人說.
你信耶穌.
可能教會的人很有愛心的.
幫到你.
他們沒有工作找到工作給你做.
他們很大脅的.
你沒有錢.
他可能有時短期借錢給你吃.
子女沒有書讀.
教會有基金可以支持他們讀書.

$^{1201}$香港還有新花款.
如果你是基督教聯會的會友.
堂會的會友.
可以免費給你機會.
埋葬在基督教永遠墳場.
這些東西他信主很多人.
但是從這段聖經裡告訴他.
我們為了進入到基督光明的國度裡.
而要進入是因為我們的罪可以得赦免.
因此.
每一個信徒在主的裡面.
都好像天父忍了我們這麼久.
也好像主耶穌基督為了父神拯救的心.
自己忍受死在十字架裡.
而今天我們作為信徒.
活在這個地上.
因此不要為了有問題有困難.
忍不住就要求神解決你的問題.
最精彩的人生.
基督徒的人生.
肅神子民的人生.
是有上帝和你一起經歷你的忍耐困難.
你今天的生命.
就能夠不折不扣地成為一個.
領受光明基業的眾聖徒的其中一位.
不知道大家有沒有留意到.
當一個人愛另一個人的時候.
多麻煩的事你都會做.
我曾經認識一個男孩.
他很喜歡一個女孩.
那個女孩不知是有心耍他.
作弄他.
還是真心的.
他跟她說.
他說我小時候去哪裡吃過一個麵包.
我很喜歡吃的.
現在可惜不在那裡.
如果可以就去吃.
你知道那個男孩.
要坐多久車去到那個女孩說的麵包店.

$^{1241}$買個麵包給他吃.
足足坐了兩個多小時車.
當他把麵包拿到那個女孩的庚前的時候.
你猜發生什麼事.
是不是那個女孩痛哭流涕.
以身相許呢.
那個女孩跟他說.
你為什麼買個麵包給我.
那個男孩跟他說.
你喜歡吃的.
你說過我掛在心裡.
我就去買給你.
他說你都傻的.
誰叫你買.
我正在減肥.
我不吃這些麵包的.
靈姐妹.
我希望你不是那個女孩的心態.
我希望你在屬靈上都不是這樣的心態.
上帝做了這麼多的事.
希望將你和我帶回祂身邊.
你就找其他的理由.
去否定上帝對你所做的事.
否定耶穌基督為你.
希望為父神的心願得到達成而去做的事.
其實最終的損失就是你.
誰可以得到去天國.
是不是上帝一定要你進天國.
上帝很想你進天國.
不過你不願意上帝是不勉強的.
你不甘心樂意是不勉強的.
禱告也是一樣.
如果你的禱告只是為了個人的目的而去禱告的話.
上帝沒有需要去成就你的.
你去找魔鬼和你成就就可以了.
或者他的爪牙.
包括天主教會.
或者很多其他的基督教會.
你會失落永恆光明在基督裡面的記憶.
禱告的目的是為人脫罪離開罪惡.

$^{1281}$然後回到上帝當中.
當你回到上帝之後.
再加力為你禱告.
持續為你禱告.
希望上帝的能力在你身上繼續作功.
希望直到你在這個世界上活出匹配主喜悅的事情.
過程不容易.
是會充滿艱辛的.
是充滿需要忍耐的.
是需要寬容的.
但是你不要小看這些忍耐和寬容.
你不要以為我們的神會無限量這樣做.
人都不會這樣做.
當那個男孩子知道那個女孩子這樣戲弄他之後.
馬上跟他分手.
原來你是這樣的心態的人.
如果這個男孩子說.
不要緊啦我很喜歡她.
跟她結了婚.
那你就自己找苦來受.
你一生的婚姻就在這樣的狀態裡面.
你看不看到普羅的禱告是什麼.
現在懂不懂禱告.
懂不懂怎樣為人代禱.
不用擔心.
有時候我們禱告聽起來好像很自私.
因為真的很自私.
但請你改過過來.
讓你的禱告能夠給你力量和上帝一起.
並肩下一條上帝所起源的道路.
走到去領受基督耶穌的國度.
這條的方向而發展.
一起來同心低頭禱告.
父神多謝你給我們在這段聖經裡面的學習.
我們知道禱告是一個看似簡單.
原來卻是非常之有深度的一個課題.
也是一個屬靈的操練.
需要我們在當中不斷地嘗試.
藉著禱告可以讓我們看到自己的軟弱.
看到自己的罪惡.

$^{1321}$從而讓我們可以去到上帝的面前.
再次支取你的力量讓我們回轉.
求你的帶領.
也要為我們為其他人的代禱.
也能夠更加的準確.
免得成為一個監臣乃後的一個狀態.
求你的啟明.
求你的聖靈的提醒.
免得我們去干犯上主你.
在禱告裡面犯罪得罪神.
求你也讓我們能夠透過準確.
明白上帝你在禱告上的教訓.
讓我們知道如何活出一個.
匹配主所喜悅的屬靈人的人生.
直奔永恆光明的國度.
禱告奉靠主耶穌基督得勝命之求.
\newpage



\section{腓立比書 3:13-14}
\label{sec:031653}
\textbf{奔向標竿人生 (普)}~~~~劉彤牧師
\newline
\newline
連結: \href{http://www.church.com.hk/acms/content.asp?site=cdc\&op=show\&type=product\&code=031653\&layout=sermon}{\texttt{031653}} ~~~~ 語音日期: 2026-02-15
\newline
\newline
\hyperref[sec:031652]{\small{< < < PREV SERMON < < <}}
~
\hyperref[sec:index_chronic]{\small{[返順時目]}}
~
\hyperref[sec:031654]{\small{> > > NEXT SERMON > > >}}
\newline
\newline
腓立比書 3:13-14
\newline
\begin{longtable}{cl}
\hline
\hline
章節 & 經文 (和合本修訂版)\\
\hline
3:13 & \begin{tabularx}{0.7\textwidth}{X} 弟兄們，我不是以為自己已經得著了；我只有一件事，就是忘記背後，努力面前的， \end{tabularx} \\ \\ \relax
3:14 & \begin{tabularx}{0.7\textwidth}{X} 向著標竿直跑，要得神在基督耶穌裡從上面召我來得的獎賞。 \end{tabularx} \\ \\
[1ex]
\hline
\hline
\end{longtable}
$^{1}$有一段聖經在腓立比書第三章十三到第十四節.
事實上是我們二月份的背經的經文.
我們會背的話,一同來背這個經文.
弟兄們,我不是以為自己已經得著了.
我只有一件事,就是忘記背後努力面前的.
向著標竿直跑,要得神在基督耶穌裡.
從上面召我來得的獎賞,Amen.
馬上這個禮拜要進入農曆新年了.
在這裡先再一次恭賀大家.
新春快樂,日日蒙福,Amen.
你知道我們華人都很注重祝福的.
所以過新年一定家家戶戶要掛個春聯.
意思就是祝福要臨到這個家庭裡面.
在家裡面我們喜歡掛福這個字.
不但正的貼,道的貼.
Make sure祝福一定要臨到.
在拜年的時候大家喜歡說吉祥話.
恭喜發財,年年有餘.
當然馬年更是金馬迎春,萬馬奔騰.
馬年是象徵著一個向著標竿直跑的人生.
我就想到在聖經當中.
保羅也常常是用奔跑來比喻他自己的生命跟福事.
事實上今天我們在腓立比書的經文.
就是我們二月份背經行動的經文.
如果你還沒有加入背經行動.
我鼓勵你趕快參加背經行列,Amen.
今年開始我真的感覺神在我們教會.
真的有一個屬靈氣氛的一個翻轉在這裡.
那我感覺神的聖靈也真的把.
深深的渴慕的心放在我們中間.
所以我們一定要抓住神所給我們這個時機.
好好裝備我們自己,Amen.
這個禮拜六我們有裝備日的課程.
下禮拜六我們有非常精彩的New Man Conference.
我們這次的主題是精兵崛起.
我鼓勵大家一定要趕快報名參加.
給你們要趕快報名參加.
因為讓我們在今年真的是成為一支裝備整齊的屬靈大軍.
Amen, Hallelujah.
你知道保羅形容他自己的生命跟福祉.

$^{41}$就是一個奔向標竿的人生.
這為什麼當他離開這個世界之前.
他在提摩代後書第四章第七節他也這樣說.
那美好的仗我已經打過了.
剛跑的路我已經跑盡了.
保羅就是形容他自己的生命福祉.
就像一個奔跑的路程一樣.
但是這種人生的奔跑.
可能不像你我所以為是那種百米衝刺的這種奔跑一樣.
反而好像是一種馬拉松式的賽跑.
你要知道百米衝刺跟馬拉松式的賽跑.
是完全不一樣的奔跑方式.
百米衝刺就是你一口氣要衝到底就好.
但是馬拉松式的奔跑.
跑得快固然很重要.
但更重要的是能夠堅持到底.
因為跑得好並不困難.
但要在這個信仰的道路上.
人生的道路上堅持到底.
卻不是一件容易的事.
特別是當我們在面對人生各樣的困境的時候.
特別是當我們看見我們的景況.
不但沒有好轉.
反而每況愈下的時候.
我們的心常常就是會灰心失望.
而且在信仰的道路上我們停滯不前.
那你說我們怎麼樣能夠在這條漫長的人生道路當中.
堅持到底我們能夠奔向標竿得著獎賞呢?.
我們一同來讀一段聖經.
在希伯來書第十二章第一到第二節.
兩節的聖經可是非常的寶貴.
我們一同來念好不好?.
我們既有這許多的見證人如同雲彩圍著我們.
就當放下各樣的重擔.
脫去容易殘累我們的最純心忍耐.
奔那擺在我們前頭的路程.
仰望為我們信心創始成中的耶穌.
他因那擺在前面的喜樂就輕看羞辱.
忍受了十字架的苦難.
便坐在神寶座的右邊.

$^{81}$從今天這個經文當中我們看見有三方面.
怎麼樣能夠讓我們有一個奔向標竿的人生.
第一個你一定要明白你不是孤單的.
跟你們說你不是孤單的.
在西班牙第十二章第一節這裡說.
我們既有這許多的見證人如同雲彩圍繞.
換句話說.
這條信仰的道路不是沒有人跑過.
你我所經歷的試煉困境.
我們心中的憂悶我們的煩躁.
也許形式有所不同.
但也不是說沒有人經歷過.
如果我們讀過詩篇第四十二篇.
第五節你就看見詩篇的作者在這裡對他的心說話.
詩篇的作者說我的心你為何憂悶.
為何在我裡面怎麼樣煩躁.
可見得他也曾經落入到一種憂悶煩躁的光景當中.
如果我們讀過以利亞的故事.
你還記得以利亞他在羅騰樹下求死的那個禱告.
在列王記上十九章第四節裡面.
他坐在樹下求死說耶和華罷了.
求你取去我的性命.
因為我不勝於我的列祖.
你也看見這麼一個偉大的先知.
他在神面前也有這種絕望的光景.
甚至要求神取去他自己的性命.
或者你也讀過主耶穌在克西瑪尼園的禱告.
馬太福音第二十六章三十八節.
耶穌的禱告說我心裡甚是憂傷.
幾乎要死啊.
我相信你很難再找到另外一個這種形容.
對人的沮喪那種絕望沮喪.
能夠更形容這麼恰當的一個經文.
我們看見耶穌也曾經經歷過這樣一種光景.
我們就看見歷世歷代的信徒.
他們都曾經在這條信仰的道路上奔跑過.
他們也曾經經歷過你我所經歷的憂悶煩躁.
挫折跟沮喪.
甚至一種在無可奈何當中的等候.
想到這裡應該讓我們心很得安慰不是嗎.

$^{121}$因為很多時候我們以為怎麼只有我們自己是這樣子.
只有我們是這個樣子.
甚至有時候連魔鬼撒旦也這樣控告我們不是嗎.
你看看你.
別人都好好的.
就是你這個樣子.
No No No No.
其實不是只有我們.
而是歷世歷代的信徒都跟我們一樣.
經歷過同樣的軟弱.
但雖然是如此.
歷世歷代的信徒也同時向我們做見證說.
神是怎麼樣信實的帶領他們跑完那當跑的路.
而且打完那當打的仗.
Amen.
事實上這為什麼希伯來書的作者在第11章.
就是在這個經文.
我們讀了經文之前.
他用了40節的篇幅.
來列出一連串的信心的英雄.
就是要來激勵我們.
鼓勵我們.
能夠繼續不斷的在信仰的道路當中.
我們能夠向前奔跑.
事實上如果你有機會讀到希伯來書第11章.
你讀到這個作者怎麼樣列出所謂的信心的英雄.
你就發現這些所謂的信心英雄.
跟我們都是一樣平凡的人.
當他們在面對困難的環境.
以及艱鉅的使命的時候.
他們一樣感受到那種壓力.
他們會心中有憂慮.
充滿恐懼.
就像神呼召摩西的時候.
你還記得神呼召摩西.
要摩西帶領以色列人.
要離開埃及.
400年奴役的綑綁的時候.
摩西沒有對神說.
神啊.

$^{161}$我等了你40年.
你怎麼到現在還呼召我.
沒有.
你還記得摩西怎麼回答神嗎.
摩西就說no no no no.
我是左口笨蛇.
請你另請高明吧.
你要知道當摩西說我是左口笨蛇的時候.
他一點不是像我們中國那種客套.
口裡說不行.
你很想那種.
你知道嗎.
no no no.
他真的覺得自己不行.
我真的左口不舌.
神啊我真的不行.
甚至到一個地步.
連神都生氣了.
能夠讓神生氣也很不容易.
最後連神都覺得他真的不行.
真的就是不行左口笨蛇.
好啦好啦.
我讓你哥哥幫助你.
讓你哥哥來幫你吧.
連神都覺得他不行.
可見摩西並不是像我們所以為.
那個很偉大的領袖.
很有自信心.
很會講話.
充滿雄心壯志.
一點都不是.
你還記得大衛嗎.
當神猜請薩姆爾要去高某大衛做君王的時候.
那個時候大衛只是家中一個去放羊的小孩而已.
沒有人看得起他.
事實上當薩姆爾到他的父親耶錫家.
要去高某耶錫的一個孩子做君王的時候.
薩姆爾對耶錫說把你所有的兒子都叫出來吧.
耶錫把所有兒子叫出來.
薩姆爾怎麼看.

$^{201}$怎麼都不在其中呢.
就對耶錫說.
你還有沒有其他的兒子啊.
耶錫才想.
喔還有一個啦.
在曠野放羊啦.
趕快把他找回來.
你看大衛平凡到一個地步.
連他爸爸都把他忘掉.
你說平不平凡.
夠平凡了吧.
所以我常說如果大衛生活在今天的世上.
他非常需要心靈內在一致.
連爸爸都把你忘掉.
我們看見不但舊約時代.
神所使用的很平凡.
新約時代不也是嗎.
提摩泰是個體弱多病.
優柔寡斷的人.
換了你我恐怕我們永遠不會選他做同工不是嗎.
在這條信仰的道路上.
我們所依靠的一點不是人的勢力跟才能.
我們所依靠的也不是人的智慧跟力量.
我們所依靠乃至神的靈.
這個靈是指在舊約時代.
能夠行神蹟騎士的靈.
也在新約時代.
給我們智慧能力的靈.
而這個靈今天.
也正住在你我的裡面.
歷史歷代的信徒都向我們說明一件事情.
那就如果像他們這樣平凡的人.
在那一種受限制的環境當中.
都能夠憑信心過一個得勝的生活.
那麼今天你我靠著神的恩典.
我們也必然能夠得勝.
所以你們在基督裡面你要明白我們不是孤單的.
我們不但有著許多的見證人圍繞著我們.
我們更是有主耶穌基督與我們親自與我們同在.
祂是昨日今日直到永遠不改變的神.

$^{241}$祂今天仍然能夠一致.
祂今天仍然能夠拯救.
祂今天仍然能夠帶來改變.
Amen Hallelujah.
所以第一個你要知道你不是孤單的.
第二個要除去心中的攬阻.
給你們除去心中的攬阻.
其實很多時候今天攬阻我們經歷神恩典的.
從來不是外在的環境.
而是我們心中的攬阻.
所以希伯來書的作者在這裡說.
我們要放下重擔.
脫去容易纏累我們的罪.
這裡的重擔跟罪.
其實都不是指的外面的環境.
或是黑暗的詮釋.
反而是指的我們心中.
那種攬阻我們跟隨神的靈.
你說在我們心中有什麼東西常常攬阻我們呢.
三方面第一個.
我們要求神除去灰心的靈.
跟著祂求神除去灰心的靈.
特別是當我們所期待的事情.
一直慈顏未得的時候.
你經歷過這樣的情況嗎.
你期待神成就的事情.
你期待神工作的事情.
一直慈顏未得.
很多時候當不順利的環境.
接二連三的靈到的時候.
常常就讓我們在挫折當中.
我們失去了盼望不是嗎.
就像彼得一樣.
你要記得主耶穌從死裡復活之後.
向彼得跟其他門徒顯現了兩次.
想想看.
如果你是彼得.
主耶穌向你從死裡復活顯現了兩次.
難道你不應該非常興奮.
歡喜快樂.

$^{281}$道去就傳揚主耶穌從此死裡復活囉.
不是嗎.
但彼得卻沒有這樣做.
反而他帶了其他的門徒.
一同到加利利海去打魚去.
哇.
他已經三年半沒有去打魚了.
為什麼現在從操就業去打魚了.
我告訴你因為他灰心了.
不但對神灰心.
而且對自己也灰心.
對神.
他看見自己跟隨耶穌三年半.
現在好像是一場空.
雖然主耶穌從死裡復活了.
但誰知道將來又會如何呢.
會不會又是另外一場空呢.
對自己.
他看見自己本來以為自己很剛強的.
可後來竟然在僕人跟卑履面前.
三次不承認主.
他一點沒有自己所以為的剛強.
灰心.
灰心使他失去勇敢的心.
繼續跟隨神.
這為什麼當主耶穌第三次.
向他們顯現的時候.
主耶穌特別對彼得說.
你來跟從我.
這就是回想到第一次.
主耶穌見到西門彼得同樣的呼召.
你來跟從我.
主耶穌再一次提醒他.
起初的呼召.
你來跟從我.
而且.
你要餵養我的小羊.
換句話說.
不要灰心.
神要繼續使用你的.

$^{321}$跟你媽說不要灰心神要繼續使用你.
你要繼續跟隨神.
來完成神的託付.
因為是個標竿為你設定.
有一個獎賞為你預備的.
第二.
要除去疲倦的靈.
跟你媽說要除去疲倦的靈.
你知道有時候我們在人生的各樣的困境跟挫折的時候.
我們會感覺到很疲倦.
不是灰心.
而是感覺累了.
一次一次的征戰.
一次一次的付出.
卻沒有看見所期待的結果.
我們感覺身心靈疲倦.
有時候不再想走下去.
就好像以利亞一樣.
當以利亞在別斯巴的曠野.
其實他在羅騰樹下求死的時候.
他其實在正面對.
已經經歷過很大的聖戰.
他不是在那個時候.
其實他是經歷過人生很高峰的時候.
但他打完這個聖戰之後.
殺了四百個巴黎的先知之後.
忽然間.
他聽見王后耶喜別要派人追殺他.
他就逃到別斯巴的曠野.
對以利亞而言.
他不是灰心.
也不是失望.
就是累了.
神啊我知道你是全能的神.
你的旨意也必定成就.
但我不行了.
我累了另請高明吧.
他不想再征戰了.
許多時候這也不是我們的光景嗎.
多少時候一次又一次的高山低谷.

$^{361}$一次又一次的輸輸贏贏.
感覺累了.
不想再奔跑了.
但我們看見神.
猜想天使怎麼樣對以利亞說呢.
另外一句上第十九章第七到第八節你看見.
這裡天使為他預備了飲食.
就對他說你起來吃吧.
因為什麼.
你當走的路甚遠.
換句話說.
這裡不是你的終點.
不要停留.
要繼續向前奔跑.
結果以利亞領受了從神來飲食之後.
第八節怎麼說呢.
他仗著飲食之力走了多久啊.
四十晝夜哇.
神的恩典真的很夠用.
Amen.
四十晝夜.
神的恩典真的夠用.
而且你也走得下去.
弟兄姊妹.
可能今天你也在面對生活跟事奉當中的上上下下.
可能今天你面對人生各樣的困境.
你也感覺疲倦了.
有時候你也不想再繼續.
可今天神的話要對你說.
嘿.
你當跑的路甚遠.
跟你媽說你當跑的路甚遠.
要將你下垂的手.
發酸的腿.
重新挺起來吧.
繼續向前奔跑.
我再次說是個標竿為你設立.
是個獎賞為你預備.
主的恩典永遠夠你用.
不但夠你用.

$^{401}$而且超乎你所求所想.
Amen.
Hallelujah.
第三.
要除去害怕失敗的零.
跟你們除去害怕失敗的零跟他說.
今天很多人在面對人生的道路的時候.
心中常常充滿害怕.
害怕什麼.
害怕失敗.
很多人在還沒有跨越出去以前已經心裡害怕.
不敢跨越.
特別如果從前曾經失敗過.
那更是不願意再去嘗試.
就好像有一次.
主耶穌在彼得的船上講到.
路加福音第五章第四節怎麼說呢.
講完了對西門說把船開到水深之處.
下網打魚.
其實難道彼得不知道打魚要去水深之處嗎.
他是漁夫啊.
很有經驗的漁夫啊.
他早就知道他當然知道.
這個捕魚一定要水深之處去.
但是他第五節他說.
我前一天晚上.
By the way.
我已經打過整夜勞力了.
夫子.
做過了啦.
什麼都沒有打著啦.
換句話說.
Been there, done that.
什麼都試過了.
什麼都沒有用嘛.
第二點今天多少時候.
我們的心也常常容易因為生活的困境.
過去的失敗.
在我們心中害怕.
讓我們不再依靠神來繼續向前.

$^{441}$我們看見彼得在說.
就依從你的話語.
他就聽從了耶穌的話.
再去水深之處打魚的時候.
第六節說結果這一次怎麼樣.
圈住許多魚.
而且怎麼樣.
往險些裂開.
Hallelujah.
這次真的不一樣啊.
第二點今天無論.
你面對前面的生活.
那個道路有多少未知的因素.
無論你的過去.
經歷過多少的失敗.
無論你心中有多少的害怕.
你千萬不要因為害怕而不再向前.
你要知道你的神.
仍然是你生命中的大牧者.
他仍然擁有你生命的答案.
你要憑信心再次向前跨越.
因為這一次會不再一樣的.
你們跟你們說這一次會不再一樣.
跟他說.
Hallelujah.
最後第三.
要仰望那為我們信心創始成中的神.
跟你們說.
仰望為我們信心創始成中的神.
因為希伯來書第十一章第一節怎麼說呢.
這裡說信是所望之事的實體.
是未見之事的確具.
可見信心不是虛無縹緲的東西.
信心是一種有確具的把握.
這個信心的確具是什麼.
這個信心的確具是在於我們所信的神.
所以說我們所信的神是為我們信心創始成中的神.
換句話說.
祂不但是那位開始在我們生命中工作的神.
祂也是那位要完成在我們生命當中工作的神.

$^{481}$Hallelujah.
祂不但開始祂的工作.
祂也必然完成這工作.
所以才叫做創始成中的神.
換句話說.
神做事絕對不會做到一半.
祂就說哎呀搞錯了.
神不會搞錯.
神開始的工作祂必定完成這工作.
Amen.
Hallelujah.
但雖然是如此.
我們需要明白神的時候不是常常像我們所以為的時候.
神的時候不是在我們的掌握當中.
這為什麼我們常常需要到神面前來仰望神.
我們需要到神面前來等候神.
等候當然不是一件容易的事情.
因為我們既不知道神的時候.
也不知道神的方法.
所以我們心中才常常憂慮.
不是嗎.
我常常就像一個守夜的人一樣.
我想在我們中間有些人在年輕的時候當過兵的經驗的.
當兵的人都有曾經守夜的經驗.
我請問你守夜辛苦不辛苦.
當然很辛苦.
所謂長夜漫漫.
真的很長.
那個時間是一秒一秒的度過.
但雖然長夜漫漫.
但對守夜的人來講心中是有個確距.
那個確距是無論如何天一定會亮.
Amen.
同樣的.
今天當我們在等候神的時候也是如此.
無論環境多麼令人絕望.
我們心中都有個確切的盼望.
那就神必按照他的旨意來完成他的工作.
我就想到主耶穌的榜樣.
我不知道你有想到過.

$^{521}$為什麼主耶穌在面對十字架的苦難的時候.
以及這麼多人給他這種羞辱的時候.
主耶穌能夠這麼平靜.
有沒有想到過.
那麼勇敢去面對十字架的道路.
當然你會說因為他是神.
他除了因為他是神之外.
有沒有其他的原因.
讓他能夠如此平靜去面對呢.
我相信那個秘訣就是在於主耶穌確實的知道.
有一個神得勝的時候是為他預備的.
當他在克西馬尼恩被捉拿的時候.
他知道那是一個受苦的時候.
所以他才說那句話.
時候到了.
受苦的時候.
時候到了.
賣我的人盡了.
而且當賓丁來捉拿他的時候.
耶穌沒有反抗.
連彼得想保護耶穌都被耶穌拒絕.
為什麼耶穌拒絕他.
因為耶穌知道這是一個受苦的時候.
他需要順服.
他不要反抗.
他需要順服.
因為這是受苦的時候.
但他心中他更知道.
有另外一個時刻是為他預備.
那是一個得勝的時刻.
那是個榮耀的時候.
那就是在第三天之後他要從死裡復活.
所以在希臘白書第十二章第二節裡面.
這裡才說.
他因那擺在前面的喜樂.
就輕看羞辱.
忍受了十字架的苦難.
便坐在神寶座的右邊.
我們看見主耶穌之所以能夠勝過這個苦難的環境.
正是因為他堅持的相信.

$^{561}$有另外一個時刻為他預備.
那是個榮耀的時刻.
那是個得勝的時刻為他預備.
所以他能夠輕看這個時刻的苦難.
Amen.
同樣今天當我面對人生的困境的時候.
你也要堅持的來相信.
無論現在的環境多麼惡劣.
你要相信有一個得勝的時刻是為你預備的.
Amen.
跟人家說有一個得勝的時刻為你預備的.
這裡面屬靈的道路是一條不斷向前奔跑的道路.
在這條道路上有高山有低谷.
有得勝的時刻也有挫敗的時刻.
但最重要的.
不只是眼前的成功跟失敗.
而是要知道你一生的道路在神手中.
就算經過死因的幽谷也不怕遭害.
因為知道神必定與你同在.
就像保羅一樣.
我喜歡保羅所說我只有一件事.
雖然在他生命當中很多事情在發生.
但保羅說其實只有一件事.
那就是忘記背後.
努力面前向著標竿直跑.
因為他知道是有神的標竿為他設立.
是有神的獎賞為他預備.
不只是保羅.
你我也一樣.
我們都是屬神的兒女.
我們一生的道路也都在神的手中.
神在我們生命當中是有標竿為我們設立.
也有獎賞為我們預備.
我們要效法保羅一樣的說其實我只有一件事.
奔向標竿.
得到獎賞.
我們一起禱告.
親愛的主我們再次來到你面前.
雖然在我們人生的道路當中常常高高低低起伏上下.
但主今天再次讓我們定睛仰望你.

$^{601}$真的像保羅說的其實只有一件事.
雖然太多的起起伏伏.
但其實只有一件事最重要.
能夠奔向你所為我們預備的標竿.
得著你所為我們預備的獎賞.
主求你今天早上幫助每位弟兄姊妹.
無論我們的環境是如何.
多麼的叫人憂慮.
多麼叫人恐懼.
多麼叫人絕望.
讓我們今天眼目單單定睛在你身上.
深深的知道你是我們生命的大牧者.
我們生命的答案都在你手中.
你是為我們信心創始.
也是承宗的神.
你開始的工作你必定完成.
主讓我們存著歡喜快樂的心.
無論環境如何.
定睛在你身上.
勇敢向前奔跑.
也真的知道主在你那裡.
有永遠的福樂為我們預備.
謝謝主.
阿們.
\newpage



\section{創世記 11:9}
\label{sec:031654}
\textbf{聖經人聖經事:分散在全地}~~~~黃克勤牧師博士
\newline
\newline
連結: \href{http://www.church.com.hk/acms/content.asp?site=cdc\&op=show\&type=product\&code=031654\&layout=sermon}{\texttt{031654}} ~~~~ 語音日期: 2026-02-05
\newline
\newline
\hyperref[sec:031653]{\small{< < < PREV SERMON < < <}}
~
\hyperref[sec:index_chronic]{\small{[返順時目]}}
~
\hyperref[sec:031655]{\small{> > > NEXT SERMON > > >}}
\newline
\newline
創世記 11:9
\newline
\begin{longtable}{cl}
\hline
\hline
章節 & 經文 (和合本修訂版)\\
\hline
11:9 & \begin{tabularx}{0.7\textwidth}{X} 因為耶和華在那裡變亂了全地的語言，把人從那裡分散在全地面上，所以那城名叫巴別。 \end{tabularx} \\ \\
[1ex]
\hline
\hline
\end{longtable}
$^{1}$《生命樹宣教網絡》.
《聖經人》《聖經時》.
每天都會和大家見面.
我們來看看聖經人士背後的上帝.
接著繼續住巴別塔.
這個塔是有名字的.
這一集開始有名字了.
「分散在全地」.
耶和華說.
看那他們成為一樣的人民.
都是一樣的言語.
如今還做起這事來.
起一個標記.
以後他們所要做的事.
就沒有不成就的了.
好 我們下去.
在那裡辯論他們的口音.
使他們的言語彼此不通.
咦 這樣做什麼呢.
這樣建好塔.
如果你看到上面.
頭那幾節.
接著一跳.
今天創世記一章十六節了.
其實耶和華說.
我們下去看看他們做什麼了.
我們創造的人.
又用方舟救回他們的人.
哇 現在很厲害.
很旺啊 下面.
好像在建東西.
看他們在建東西.
不如我們下去看看吧.
當然這些.
都是一些對於上帝.
擬人化的收持學的應用.
上帝哪用下去看他.
早就知道了.
不如我們下去看看吧.
那麼下去耶和華就說.

$^{41}$看吧 他們成為一樣的人民.
一樣的言語.
最重要是出現了兩次說.
一樣 一樣.
他們都由於這個城市.
或者這個塔.
這個都市化有一個特點.
聖經原來早就說了.
就會使人們成為一樣的.
所以你看今天的大城市.
大城市的人的世界觀.
價值觀.
一般衣著.
講的說話等等.
你發覺他們都有一個特色.
就是一樣.
一樣.
哇 面孔 面型 化妝.
頭飾.
他們穿什麼衣服.
全部都是一樣.
那麼這個一樣一樣.
又沒什麼不好的.
他們完全變成一樣一樣.
其實如果你說不好又不好.
譬如如果他們的世界觀.
如果他們的屬靈觀.
等等是偏離上帝的.
他們就一起偏離了.
如果有人帶著他的路.
走歪了就走歪了.
如果那幫人變成很暴戾.
就一起很暴戾了.
一樣的.
如今他們做了這件事.
就是說他們用一個品牌.
來將自己更加堅固了.
那麼以後他們要做的事.
就沒有不成就了.
就是說他能力比以前更大了.

$^{81}$沒有一不成就的事.
意思就是說.
他們很多時候會做的事情都會成功.
就沒有什麼難倒他們.
如果他們要做他們會做到.
這是真的.
因為大家的力量集中在一起.
這個在商業用一個詞叫synergy.
就是叫做協作.
一加一大於二.
一加一是三了.
因為大家兩個人合起來.
他比兩個分開的人的能力更加大.
他們一起協作起來.
就是協同效應.
synergy叫協同效應.
就會很厲害了.
那麼現在變成一樣的人.
一樣的語言更加厲害了.
其實好啊.
那就更大發展了.
但是上帝有點擔心.
因為他知道當人這樣的時候.
他們會容易走錯路.
離開正路的.
於是就說不如我們下去.
變亂他們的口音.
令他們言語不通.
又沒有說他們很不對.
預先做了一個介入.
這個介入就是令他們的口音變化了.
彼此言語不通.
講起不同的語言上來.
講不到溝通不到.
這個上帝很壞.
為什麼要這樣做.
其實他壞是用我們現在這個想法來判斷.
為什麼上帝要給一棵分別視角樹在伊甸園裡面.
如果沒有.
人就不會錯,不會死.

$^{121}$變成機械人.
現在是增加了言語不通.
增加了他們一起的難度.
譬如今天我們可不可以一起呢.
一起,我們又會英文又會什麼.
全世界都通.
即是全世界都通的.
再加上現在科技發展到翻譯技術這麼好.
全世界都通的.
不過是增加了他們一起的難度.
不令他們變成一樣.
有些不同.
於是他們沒有全世界一起能夠.
synergy這樣的協同效應而做出來.
如果要做壞事真的可以壞到極.
就變亂他們的口音.
在這裡干預一下.
於是就從那裡分散他們全地.
他們停工,為什麼.
因為講不通.
就不做那個城.
那塔還不做呢.
也不做.
當然今天又有些考古學發現高塔在哪裡.
以前有些馬雅文化都說那些塔很高.
品牌,地標,厲害.
你見到那些大廈.
以前香港都很proud of自己.
因為全世界十大最高的建築物.
有兩座在香港.
IFC和ICC,兩座很厲害.
當然現在被人趕過了.
他們要停工了,不做那城.
因為在那裡變亂了天下人的口音.
令到他們分散在全地.
所以那城名叫巴別.
留意一下.
變亂這個字是Bla.
Bla是混合,混淆,混在一起.
變亂是令到他有些東西又一齊的音.

$^{161}$有些東西又不一齊的音.
譬如我時常去柬埔寨.
我不學的,靠翻譯的.
開始我學一二三,拍照.
一二三,妹BiBi.
我只學妹BiBi.
後來我發覺很好笑.
他們走的時候.
走啦走啦,再見再見,走啦走啦.
走啦走啦,走啦走啦.
如果是平輩,走啦走啦.
講什麼呢?.
講厲害厲害,厲害.
普通話厲害厲害.
厲害就是走啦走啦.
我就發覺到厲害有沒有像中文的離開.
有些像離開.
厲害是混淆,混合.
Bla是變亂將他混合.
有些像有些不像.
談到好像談不到.
這個日文我不太懂.
不過後來見到時時有人寫或說.
大丈夫大丈夫.
好像一句音是大丈夫大丈夫.
大丈夫什麼意思?.
大丈夫即是不用怕不用怕.
很強不用怕.
咦?跟中文又像.
很多時候又像又不像才叫變亂.
於是口音有些像有些不像.
有些聽得明白有些聽不明白.
於是不如我們聽得明那些一齊.
不然都不知道怎麼辦.
他們就開始分散在全地.
找自己的地方再遷移.
那個城和那個塔就開始叫做巴別.
巴別他們叫做Babel.
Babel就是混雜了.
這個Bla是動詞來的.

$^{201}$混亂.
混雜是個名詞來的.
Babel就是那個地方.
那裡的口音很亂.
這個巴別塔就因為.
被一般基督徒覺得不好.
為什麼?.
因為他們建一個塔.
跟上帝比高.
錯了.
所以上帝就令他們變亂.
口音講不到.
其實未必是.
只是上帝介入.
不想他們用相同的口音.
於是他們變成非常集中.
和做出一個協同的效應.
於是他們要做什麼事.
都做到.
好就好.
如果壞就真的很差.
上帝都不想他們這樣.
不如令到他們變成很多不同的民族.
而且上帝是想人是怎樣的?.
是生養眾多遍滿全地.
這個令到他們口音的不同.
都是另一個方法.
令到他們遍滿全地生養眾多.
但是他們剛剛就想什麼呢?.
聚集在一起.
不要分開.
用一個地標和力量將他們吸引.
上帝於是不如令他們分開.
他們免得聚集在一起.
又容易集中力量.
如果壞就糟糕了.
我都不想再對於他們有任何毀滅的動作.
那就算了.
叫他們分開.
混淆了言語令到他們分開.

$^{241}$於是他們紀念那個塔.
就說那裡我們忽然之間亂了.
混雜.
Babel.
就是這個意思.
好了.
看看今天的金句.
因為耶和華在那裡混亂天下人的言語.
使眾人分散在全地.
那個城的名字就叫巴別.
神是想人生養眾多.
遍滿全地.
地球這麼大.
為什麼只去兩河流域呢?.
明白嗎?.
不要分散.
而且還建一些更加集中的東西.
建一個地標建一個城.
更加不好.
所以神就做這件事.
是分散變亂他們的口音.
他們就將那個城建到一半拋棄了.
為什麼呢?.
我們亂了說不了.
所以我們不再在那裡了.
那個地方的名字就叫巴別混雜.
不過後來.
由於巴別塔.
或者巴別或者巴比倫的.
祖先起源的地方.
被聖經後面一直指為是一個.
魔鬼撒旦勢力的代表.
所以認為是不好的.
特別是說.
塔頂通天那裡認為是跟神比較.
其實我以為不是.
只不過他們想集中.
而神就是想他們分散傳遞.
遍滿傳遞.
生養眾多.

$^{281}$這是其中一個.
我自己對這個變亂口音的解釋.
多謝大家.
我們今天到這裡.
我們一起禱告.
主要我們求你幫助我們.
你給了一個任務我們.
去修理看管伊甸園.
又叫我們生養眾多.
遍滿傳遞.
甚至是伊甸園以外.
但是在巴別這個士拿這個地方的時候.
看到人們喜歡聚集.
於是你就變亂他們的口音.
他們製造的品牌也不再存在.
要做一個塔做一個城.
不過不是上帝喜歡.
主我們求你幫助我們.
看清楚你的心意.
讓我們努力的.
去建造一個更美好的世界.
雖然這個世界越來越看來越差.
不過由於你讓我們繼續在世上存留.
幫助我們造炎造光榮耀利.
帶領人更加美好防護.
求主與我們同在.
加我們力量.
也讓我們記得.
主給我們最先的使命.
禱告奉靠耶穌基督的名字.
阿們.
\newpage



\section{約書亞記 5:1-15}
\label{sec:031655}
\textbf{在神聖使命中更新與順服}~~~~劉銳鋒牧師
\newline
\newline
連結: \href{http://www.church.com.hk/acms/content.asp?site=cdc\&op=show\&type=product\&code=031655\&layout=sermon}{\texttt{031655}} ~~~~ 語音日期: 2026-02-15
\newline
\newline
\hyperref[sec:031654]{\small{< < < PREV SERMON < < <}}
~
\hyperref[sec:index_chronic]{\small{[返順時目]}}
~
\hyperref[sec:031656]{\small{> > > NEXT SERMON > > >}}
\newline
\newline
約書亞記 5:1-15
\newline
\begin{longtable}{cl}
\hline
\hline
章節 & 經文 (和合本修訂版)\\
\hline
5:1 & \begin{tabularx}{0.7\textwidth}{X} 約旦河西亞摩利人的眾王和靠海迦南人的眾王，聽見耶和華在以色列人前面使約旦河的水乾了，直到他們過了河，眾王因以色列人的緣故都膽戰心驚，勇氣全失。 \end{tabularx} \\ \\ \relax
5:2 & \begin{tabularx}{0.7\textwidth}{X} 那時，耶和華對約書亞說：「你要造火石刀，第二次為以色列人行割禮。」 \end{tabularx} \\ \\ \relax
5:3 & \begin{tabularx}{0.7\textwidth}{X} 約書亞就造了火石刀，在哈爾拉勒山為以色列人行割禮。 \end{tabularx} \\ \\ \relax
5:4 & \begin{tabularx}{0.7\textwidth}{X} 約書亞行割禮的原因是這樣：從埃及出來的眾百姓，所有能打仗的男丁，出了埃及以後，都死在曠野的路上。 \end{tabularx} \\ \\ \relax
5:5 & \begin{tabularx}{0.7\textwidth}{X} 這些從埃及出來的眾百姓都受過割禮；但是那些出埃及以後，在曠野的路上所生的眾百姓卻沒有受過割禮。 \end{tabularx} \\ \\ \relax
5:6 & \begin{tabularx}{0.7\textwidth}{X} 以色列人在曠野走了四十年，直到那從埃及出來，全國能打仗的人都消滅了，因為他們沒有聽從耶和華的話。耶和華曾向他們起誓，必不容許他們看見耶和華向他們列祖起誓要給我們的地，就是流奶與蜜之地。 \end{tabularx} \\ \\ \relax
5:7 & \begin{tabularx}{0.7\textwidth}{X} 他們的子孫，就是耶和華興起接續他們的，都沒有受過割禮；因為在路上他們沒有受割禮，約書亞就為他們行割禮。 \end{tabularx} \\ \\ \relax
5:8 & \begin{tabularx}{0.7\textwidth}{X} 全國的人都受了割禮，留在營中自己的地方，直到痊癒。 \end{tabularx} \\ \\ \relax
5:9 & \begin{tabularx}{0.7\textwidth}{X} 耶和華對約書亞說：「我今日將埃及的羞辱從你們身上除掉了。」因此，那地方名叫吉甲，直到今日。 \end{tabularx} \\ \\ \relax
5:10 & \begin{tabularx}{0.7\textwidth}{X} 以色列人在吉甲安營。正月十四日晚上，他們在耶利哥的平原守逾越節。 \end{tabularx} \\ \\ \relax
5:11 & \begin{tabularx}{0.7\textwidth}{X} 逾越節的第二日，他們吃了當地的出產，就在那一天，吃了無酵餅和烘過的穀物。 \end{tabularx} \\ \\ \relax
5:12 & \begin{tabularx}{0.7\textwidth}{X} 他們吃了當地出產的第二日，嗎哪就停止了。以色列人不再有嗎哪了。那一年，他們就吃迦南地的出產。 \end{tabularx} \\ \\ \relax
5:13 & \begin{tabularx}{0.7\textwidth}{X} 約書亞靠近耶利哥的時候，舉目觀看，看哪，有一個人站在他對面，手裡拿著拔出來的刀。約書亞到他那裡，對他說：「你是屬我們的，還是屬我們敵人的呢？」 \end{tabularx} \\ \\ \relax
5:14 & \begin{tabularx}{0.7\textwidth}{X} 他說：「不，我現在來是要作耶和華軍隊的元帥。」約書亞就臉伏於地下拜，說：「我主有甚麼話，請吩咐僕人吧！」 \end{tabularx} \\ \\ \relax
5:15 & \begin{tabularx}{0.7\textwidth}{X} 耶和華軍隊的元帥對約書亞說：「把你腳上的鞋脫下來，因為你所站的地方是聖的。」約書亞就照著做了。 \end{tabularx} \\ \\
[1ex]
\hline
\hline
\end{longtable}
$^{1}$主持人:各位弟兄姊妹,平安!.
主持人:今日要和大家分享的題目是:.
「在神聖使命中,更生於信服」.
主持人:經文是記載在《約書亞記》的第五章一至五節.
主持人:不知道大家有沒有一些經歷,就是在一些很重要的時刻.
主持人:譬如說第二天,你要面對一個很重要的考試,或者是面試,或者是比賽.
主持人:在前一天,你不是很努力地去溫習,去練習,去預備.
主持人:而是去做了一些不是很相關的事情.
主持人:譬如說你去公園走走,又或者你約朋友去喝一杯咖啡.
主持人:又或者你是安靜下來,有一個很深度的祈禱.
主持人:這些事情,從結果上來看,是不會增加你對明天考試,或者面試,或者比賽的一些作用的.
主持人:因為你不會記多了一些東西,也不會練習熟習一點.
主持人:但是當你做完這些事情之後,你覺得很奇妙的就是.
主持人:你很焦躁的心能夠平靜下來.
主持人:你混亂的思緒也能夠開始變得清晰.
主持人:你就好像得到了一種很全新的內在的力量.
主持人:去面對明天那個很巨大的挑戰.
主持人:今天要分享的經文是約書亞記第五章.
主持人:在當中記載的就是這樣一個在大戰之前的奇特的準備.
主持人:他們要做什麼大戰呢?.
主持人:因為他們要準備進入迦南地的時候.
主持人:要面對第一場的戰爭,就是攻打耶利哥城.
主持人:大戰在即,他們有什麼的預備呢?.
主持人:本身要迎接戰爭的話,要做些什麼的?.
主持人:是不是應該很緊張?可能要部署很多東西,是不是?.
主持人:也可能要派一些人去偵察敵情,又或者我們要排兵佈陣.
主持人:又或者我們要多些的武器,去預備一些武器準備打仗.
主持人:但是神的命令是很奇怪的.
主持人:神是讓全軍停止前進.
主持人:讓所有的南丁接受一個可能令他們接下來幾天都走不了路的一個割禮的手術.
主持人:之後又不是馬上再去攻打耶利哥城.
主持人:而是集體去過一個逾越節,好像吃大餐那樣.
主持人:聽起來好像是自廢武功,是不是?.
主持人:打仗之前還行割禮?可不可以打仗之後先行割禮?.
主持人:可能身體軟弱,不太方便打仗.
主持人:但是為什麼要搞到自己最虛弱的時候來迎戰呢?.
主持人:今天我們一起去看一看,在這麼緊張的時刻,神為什麼要祂的子民做這些看起來不太合時宜的準備呢?.
主持人:其實這個背後可能揭示一個道理,就是所有神聖使命得以成功.
主持人:最核心,卻是最容易被人忽略的一個秘訣,就是不是我們為神能夠做些什麼.
主持人:而是神在我們裡面做了些什麼.

$^{41}$主持人:他們的身份,並且預備他們信服的心,以迎接神聖使命的挑戰.
主持人:他們有什麼神聖的使命呢?.
主持人:我們知道在出埃及的年代,摩西重臣領受了使命,就是要大以色列人離開埃及,出埃及.
主持人:到了約書亞這一代,新的一代起來了.
主持人:他們的神聖使命就是要在約書亞的帶領下,憑著信心去跨越若幹河,征服迦南地.
主持人:並且將其作為上帝的應許之地承受為業,這個就是他們新一代的神聖使命.
主持人:所以在這裡我們看到,以色列人他們停下來,其實是神對他們的特別要求.
主持人:有時這個停頓是很重要的,是為了我們內在的生命和神的關係上,可以做一個最徹底的預備.
主持人:所以我們常常都聽到一句說話就是,停下來是為了走得更遠.
主持人:有時我們真的需要停一停,以致我們能夠走得更遠.
主持人:這段經文裡面,我們一起來看一看,神怎樣去預備以色列人,怎樣要他們停一停.
主持人:有兩方面我想大家去思想,一方面就是更新,另一方面就是信服.
主持人:我們先看看更新,首先就是更新的身份,以葛禮重建聖若群體.
主持人:在第一節到第九節是說他們透過葛禮來更新他們的身份,以確立他們是屬神的子民.
主持人:第一節是說,約旦河西阿摩利人的諸王和靠海迦南人的諸王,聽見耶和華在以色列人前面使約旦河的水乾了.
主持人:等到我們過去,他們的心因以色列人的緣故就消化了,不再有膽氣.
主持人:這節經文很厲害,它描述一個很厲害的場景,就是以色列人所向披靡,阿摩利人和迦南人聞風喪膽.
主持人:這裡也提到,那些敵人的心因以色列人的緣故就消化了,不再有膽氣.
主持人:這句經文在和合本修訂版是這樣翻譯的,眾王因以色列人的緣故都膽戰心驚,勇氣全失.
主持人:也就是說,以色列人靠氣勢就已經贏了,令到敵人聞風喪膽.
主持人:是不是以色列人很厲害?當然不是,厲害的是神,而不是以色列人.
主持人:因為是神令到他們可以過到約旦河,是神在背後承受這個神蹟.
主持人:所以在經文裡面提到,聽見耶和華,就是說耶和華才是成功的關鍵.
主持人:是他們戰勝的關鍵,而不是以色列人厲害,而是耶和華厲害.
主持人:接著在第二節就提到,耶和華就吩咐約書亞和以色列人要行割禮.
主持人:對於神的吩咐很特別,就是約書亞沒有問原因,只是順服去執行.
主持人:我想問弟兄姊妹,如果你是約書亞的話,你會不會去問神為什麼?.
主持人:快要打仗了,大戰在即了,還要去行割禮?行割禮有什麼危險?.
主持人:身體軟弱,可能被人乘機去攻擊他們.
主持人:如果你是以色列人的領袖,你是約書亞,你會不會去跟神爭辯?.
主持人:你會不會去跟神討論為什麼會這樣?但是我們看到,約書亞沒有.
主持人:祂只是順服執行耶和華的命令,祂沒有問原因.
主持人:祂不問原因是因為祂不知道神為什麼要這樣做,還是祂很明白神的心意.
主持人:所以祂也不用問,照做就可以了?究竟是怎樣呢?.
主持人:接下來的經文其實是講了答案給我們聽的,在四到七節已經講明了.
主持人:約書亞其實是很清楚明白神的心意的,為什麼要行割禮呢?.
主持人:是因為從埃及出來的,能夠打仗的男丁,全部都已經死在曠野那裡了.
主持人:因為他們不聽從耶和華的話,他們的子孫,即是在曠野出生的那一代,他們是沒有受過割禮的.
主持人:意味著什麼呢?意味著新的一代,他們需要重新跟神立約,確認他們作為神子民的身份.
主持人:約書亞很清楚這一點,所以她照著去做.

$^{81}$主持人:為什麼割禮對以色列人來說這麼重要呢?.
主持人:因為割禮其實是代表著跟神立約的永恆的一個記號.
主持人:除了一個記號,也是他們身份的象徵,代表他們是屬神的子民,是一個很重要的身份象徵.
主持人:同時也是信仰的傳承,他們透過割禮將信仰一代一代傳承.
主持人:這個儀式其實是源自於《亞伯拉罕之約》,即是在《創世記十七章》所提到的.
主持人:男嬰出生第八日就要行割禮,這個象徵著加入神的選民群體.
主持人:是終身要奉行神的律法,這個也是區分猶太人和外邦人的神聖標誌.
主持人:在屬靈的層面上,割禮其實也是代表著將我們生命裡面一些屬肉體的罪性和屬血氣的自我除掉.
主持人:事實上摩西在他臨終之前也跟以色列文講過割禮的屬靈的意義.
主持人:在《新命紀》三十章六節,我們一起來讀一讀.
主持人:耶和華,你神必將你心裡和你後裔心裡的污穢除掉,好叫你盡心盡性愛耶和華,你的神,使你可以存活..
主持人:這是很重要的,要將他們心裡的污穢除掉,是一個除污和潔淨的過程.
主持人:今天是年廿八,大家清楚了嗎?我不知道你們洗了沒有.
主持人:我們每年做一次,可能有些人不止一次.
主持人:其實洗撻撻是很重要的,是去救迎新和除污潔淨的過程.
主持人:除了我們的家庭要這樣去潔淨之外,其實我們個人的生命也很需要潔淨.
主持人:所以我們需要潔淨,除污,才能夠被主使用,否則神很難使用我們.
主持人:在東南亞國家,牧師常常都需要去趕鬼,被邀請去趕鬼.
主持人:曾經有牧師被邀請去趕鬼,他答應了,但趕鬼的時候,那隻鬼很厲害.
主持人:他知道牧師做過些什麼,他一見到牧師就說,你整個禮拜沒有靈修了.
主持人:不單止說他沒有靈修,還跟他說,你這個禮拜將他整個禮拜所做的.
主持人:不太見得人的事,全部當眾說出來,你說真的很難面對這個處境.
主持人:所以牧師就敗北而逃,結果去趕鬼沒有趕到鬼,反過來被鬼趕走了.
主持人:如果我們的生命不潔淨,我們也很難去服侍神.
主持人:其實魔鬼也知道我們的軟弱,所以這個也成為一個破口,魔鬼去攻擊我們.
主持人:求主幫助我們,讓我們真的能夠成為聖潔的人,一個合用的器皿給神使用.
主持人:因為耶穌基督已經幫我們行了非人手所行的國禮.
主持人:正如歌羅西書所說,我們一起讀一讀.
主持人:你們在他裡面也受了不是人手所行的國禮,是基督使你們脫棄肉體成肉的國禮.
主持人:所以因著耶穌基督的緣故,我們已經被潔淨,除了污了.
主持人:我們已經脫去了肉體成肉,以致我們今天能夠坦然無懼來到神的面前.
主持人:可以被神使用,所以是一個很重要的潔淨過程.
主持人:他們在打仗之前,要首先潔淨自己,確立他們的身份是屬神的子民.
主持人:才能夠為神去爭戰.
主持人:第八節和第九節又提到,他們受了國禮之後做了另一件事.
主持人:就是住在營中,住了幾天,等到痊癒.
主持人:這個期間,耶和華對約書亞說.
主持人:我今天將埃及的羞辱從你們身上滾去了.
主持人:因此那地方叫吉甲,直到今天.
主持人:所以國禮更新的象徵,其實是從埃及的羞辱中得到釋放.

$^{121}$主持人:這句說話不是對過去的叛逆的總結,更加是對未來的信服的期許.
主持人:國禮是上帝和亞伯拉罕所立的永恆之約的核心記號.
主持人:所以以色列人出埃及之後,他們在曠野漂流了四十年.
主持人:新生一代未受過國禮,所以導致了這個聖約的外在表徵中斷了.
主持人:進入了迦南之後,約書亞奉耶和華的吩咐.
主持人:為全體的男丁行國禮,這個標誌著他們身份的重建.
主持人:新一代的以色列人不再是曠野的流氓,他們不再是野漂一族.
主持人:我們今時今日有港漂一族,當時可能是野漂一族.
主持人:但他們已經不再是野漂一族,他們是重新被納入亞伯拉罕的後裔.
主持人:是屬神的子民,這個神聖的群體.
主持人:亦都標誌著他們羞辱的除去.
主持人:在這裡提到一個地方叫吉甲,吉甲的意思就是滾離或者是除掉.
主持人:其實是意味著神聖的轉折.
主持人:提醒以色列人如何帶領他們在過去的失敗中重新站立起來.
主持人:另外就是標誌著信服的實踐.
主持人:國禮不是可有可無或者你可以選擇做或者不做的事情.
主持人:而是立約的強制性要求,一定要去做,不做會怎樣呢?.
主持人:在《創世記》告訴我們,不受國禮的男子必從民中被剪除.
主持人:所以信服是很重要的.
主持人:為何上一代的以色列人會倒閉在曠野,不可以進入迦南地呢?.
主持人:就是因為第六節所說的,他們沒有聽從耶和華的說話.
主持人:沒有聽從耶和華的說話,意思就是他們不信服.
主持人:不信服就會倒閉在曠野,這是很嚴重的後果.
主持人:所以信服對我們來說是很重要的.
主持人:為何在攻打耶利哥城之前要行國禮呢?.
主持人:這句有什麼重要的意義呢?.
主持人:不可以攻打完才來行嗎?.
主持人:會不會在真的行了國禮之後,敵人趁機,趁他們身體軟弱的時候來攻擊他們呢?.
主持人:我們這裡看到,神命令以色列人在敵人面前行國禮.
主持人:我們經常背詩篇二十三篇,有一句就是:在我的敵人面前,你為我擺設筵席.
主持人:我們這裡看到,在我的敵人面前,你為我行國禮.
主持人:是不是很不同呢?為何在我的敵人面前要行國禮呢?.
主持人:這個行動好像是自廢武功,或者是自己削弱自己的戰鬥能力.
主持人:看起來是這樣,但實際上不是,實際上是對信服的一個深刻的考驗.
主持人:雖然受國禮真的可能削弱了他們的所謂戰鬥能力.
主持人:這個其實是真的有機會被人攻擊的.
主持人:因為在創世記二十四章的時候,曾經都出現過類似的事情.
主持人:在事件城當中,這些居民因為行了國禮,他們的身體就變得虛弱.
主持人:雅各的兒子們就趁機攻擊他們.
主持人:但在約書亞記當中,我們看到神是親自去保護以色列人的.

$^{161}$主持人:迦南地的諸王是因為恐懼而不敢去攻擊他們.
主持人:他們已經聞風喪膽,根本不敢去攻擊他們.
主持人:因為看到神的保護.
主持人:所以在即將發動對耶律哥城的進攻之前.
主持人:神是沒有命令他們去收縮一些攻城的器械.
主持人:也沒有叫他們去做一些軍事的部署.
主持人:而是命令他們行國禮.
主持人:這是一個很重要的屬靈戰爭之前的預備.
主持人:我們看到在國禮之後,他們住在營中.
主持人:他們等到痊癒,雖然身體虛弱.
主持人:但他們卻可以經歷到神的同在.
主持人:在國禮完成那天.
主持人:約書亞也可以預見到耶和華軍隊的元帥.
主持人:並且被吩咐脫下他的鞋.
主持人:因為他所站的地方是聖的.
主持人:這就表明真正的爭戰的權柄.
主持人:不在於人多或武器有多先進.
主持人:不在於你是否活在聖約當中.
主持人:軍事勝利的前提是聖潔和信服.
主持人:不是武力優勢,不是你的軍隊有多大.
主持人:真正的德性是始於內在的更新.
主持人:而不是外在的軍事部署.
主持人:信服神有時會讓我們處於危險的境地.
主持人:但神的保護始終伴隨著我們.
主持人:在我們的生活當中.
主持人:信服神的命令可能寓意著我們要放棄世俗的安全感.
主持人:選擇一些看起來比較艱難的道路.
主持人:但正如以色列人所經歷.
主持人:我們看到神的應許和力量永遠成為我們堅固的後盾.
主持人:第一就是更新的身份.
主持人:他們要重新確立他們是神的子民.
主持人:除了更新的身份之外,他們還要過更新的生活.
主持人:他們要守禦節以迎接新的轉變.
主持人:在接下來的經文第十到十二節提到.
主持人:他們在吉格安營之後守禦節.
主持人:禦節對他們來說是一個很重要的節日.
主持人:這個節日是紀念神在埃及的時候.
主持人:如何拯救他們和如何供應他們的需要.
主持人:所以在出埃及記的時候.
主持人:就標誌著他們從埃及的奴隸當中得到釋放.

$^{201}$主持人:從而過了一個自由的生活.
主持人:而在約書亞記當中.
主持人:就標誌著以色列人從曠野的漂流開始進入應許之地.
主持人:開始他們新的生活.
主持人:所以有新的生活,他們就有不同的生活方式.
主持人:和以前是不同的.
主持人:神的供應方式亦都改變了.
主持人:我們見到神的供應方式.
主持人:從曠野的神蹟開始轉變為自然的土產.
主持人:他們以前是吃從天上降下來的嗎哪.
主持人:但是他們進入迦南地之後.
主持人:不再吃嗎哪了.
主持人:他們是吃迦南地.
主持人:就是這個應許之地的土產.
主持人:是豐富的一個土產.
主持人:這是一個轉變.
主持人:神的供應轉變了.
主持人:嗎哪的亭子其實是標誌著.
主持人:神是對以色列人的新階段的供應.
主持人:他們不再依賴超自然的食物.
主持人:而是開始享受應許之地的豐富.
主持人:這個轉變其實都提醒我們.
主持人:神的供應總是適應我們的需要.
主持人:在曠野的時候.
主持人:神是以嗎哪來供應以色列人的需要.
主持人:那時候因為他們沒有辦法自給自足.
主持人:在那個階段他們沒有辦法自給自足.
主持人:但是在他們進入應許之地之後.
主持人:神就轉向了自然的豐收.
主持人:這是他們能夠自給自足的一個新階段.
主持人:所以在我們的生活裡面.
主持人:神的供應也是這樣的.
主持人:他根據我們的需要和階段.
主持人:以不同的方式來供應我們.
主持人:所以嗎哪停止是一個必然性.
主持人:因為神要我們去成長.
主持人:而不是要我們停留在屬靈的嬰兒的階段.
主持人:所以這個好像是一個斷奶的行動.
主持人:我相信我們每一個人都有經歷過這個斷奶的階段.
主持人:斷奶看起來好像比較狠心一點.

$^{241}$主持人:但是這個也是比較必不可少的.
主持人:因為我們要成長 我們要成熟.
主持人:一定要斷奶.
主持人:弟兄姊妹想問你.
主持人:你斷了奶沒有?.
主持人:我這樣問好像很冒犯.
主持人:弟兄姊妹說我信主這麼久.
主持人:你還問我斷了奶沒有?.
主持人:但事實上在我們當中.
主持人:真的有可能我們信了很久.
主持人:但其實我們還是一個未斷奶的嬰孩.
主持人:我們還是仰望著嗎哪的供應.
主持人:我們還是靠恩典的時代.
主持人:我們還是停留在嬰孩的階段.
主持人:我們是一個老嬰孩.
主持人:求主幫助我們.
主持人:我們不要停留在那個階段.
主持人:我們要成長.
主持人:所以逾越節的慶祝.
主持人:其實提醒我們要感恩神的救贖和供應.
主持人:同時也期待祂對未來的應許.
主持人:逾越節不單止是歷史的紀念.
主持人:也更加是對我們個人救贖的提醒.
主持人:在今時今日我們有很多的節日.
主持人:我自己是一個不太過節日的人.
主持人:包括生日我也不過.
主持人:沒有什麼儀式感.
主持人:但是我有時候又發覺其實有節日也挺好的.
主持人:因為節日提醒我們要做一些事.
主持人:要紀念一些事.
主持人:要記得一些事.
主持人:所以節日對我們來說是很好的一個機會.
主持人:讓我們去記得一些很重要的事.
主持人:還有做一些很重要的事.
主持人:對以色列人來說.
主持人:逾越節是很重要的.
主持人:他們記得神的恩典和神的救贖.
主持人:在新約裡面.
主持人:在現在我們的年代.
主持人:我們也透過耶穌基督的犧牲.

$^{281}$主持人:我們能夠得蒙救贖.
主持人:我們可以從罪當中釋放出來.
主持人:不再作罪的奴僕.
主持人:我們也可以過新的生活.
主持人:所以逾越節的慶祝.
主持人:也教導我們要期待神對未來的應許.
主持人:正如以色列人進入應許之地一樣.
主持人:我們也期望神的國度的完全實現.
主持人:所以在慶祝逾越節的時候.
主持人:我們不單止是紀念過去的救贖.
主持人:也期待未來的榮耀.
主持人:在這段經文裡面.
主持人:我們看到有幾件事他們做了.
主持人:第一就是行國禮.
主持人:這個就是恢復他們聖約的身份 屬神的子民.
主持人:接著他們就守逾越節.
主持人:是紀念他們出埃及的時候 神的救贖.
主持人:是在重述救恩的歷史.
主持人:接著就是嗎哪子式.
主持人:這個是在說他們需要成長.
主持人:接著他們的生活方式都要有所轉變.
主持人:神的供應方式也都轉變了.
主持人:最後在接下來的十三到十五節的經文是說到元帥的顯現.
主持人:這裡是提醒他們神是親自掌權.
主持人:也都宣告征戰是屬於神的.
主持人:這一連串的事件其實是在展現一幅圖畫給我們看.
主持人:是在說以色列民他們的身份是越來越進深.
主持人:他們從被救贖的群體開始進深到是一個守約的子民的群體.
主持人:接著再進深到他們是承受產業的繼承者.
主持人:這個身份也是在表示他們和神的關係是更加親近的.
主持人:我們見到他們除了有更新的身份之外.
主持人:還有更新的生活方式.
主持人:除了更新之外,另外一樣重要的就是信服.
主持人:我們見到他們在當中有信服的行動.
主持人:他們交出主權,以神為元帥.
主持人:在十三到十五節,約書亞靠近耶利哥的時候.
主持人:約書亞就問他是來幫助他們還是幫助他們的敵人.
主持人:那個人就回答:我來是要作耶和華軍隊的元帥.
主持人:其實這個回答是顛覆了約書亞的觀念的.
主持人:本身一個人來當然很直接地問.

$^{321}$主持人:你是來幫我們還是幫我們的敵人?.
主持人:都很正常地問.
主持人:但是這個人的回答真的顛覆了他的思想.
主持人:他說我是來作耶和華軍隊的元帥,我不是來幫你的.
主持人:我是來做元帥的,我是帶領你的.
主持人:你要聽我的話.
主持人:這個是完全不同的角色.
主持人:耶和華軍隊元帥的出現其實是象徵著神對他的子民的領導.
主持人:在舊約當中,神常常都透過先知或者是祭司.
主持人:或者是一些領袖來領導他的子民.
主持人:但是在這裡,神是親自以元帥的身份來出現.
主持人:所以強調他對以色列人進入應許之地的直接領導.
主持人:約書亞有什麼反應呢?.
主持人:他的反應是謙卑和信服.
主持人:他立即就俯伏下拜,詢問神的旨意.
主持人:他的謙卑和信服是一個很重要的事情.
主持人:因為他作為以色列人的領袖.
主持人:他本來可以很驕傲和很自大.
主持人:但是他選擇了謙卑和信服.
主持人:他的反應告訴我們.
主持人:無論我們的地位如何.
主持人:我們都應該以謙卑的態度來面對神的呼召.
主持人:所以約書亞的信服不單止表現在他脫下他的鞋.
主持人:更加表現他很信服的態度.
主持人:謙卑俯伏地下拜.
主持人:在聖經裡脫鞋是象徵謙卑和聖潔.
主持人:所以脫鞋表明他所站的地方是聖地.
主持人:他願意以聖潔和謙卑的態度來面對神的呼召.
主持人:耶和華軍隊元帥的教導其實是強調神的主權.
主持人:在我們生活中我們常常都需要意識到神的主權.
主持人:將我們的計劃和決策交託給祂.
主持人:正如約書亞在元帥的面前放下自己的計劃和驕傲.
主持人:我們都需要在神的面前放下我們的計劃和驕傲.
主持人:讓神來作主.
主持人:在我們個人的生活中是這樣.
主持人:在教會的事工中也是這樣.
主持人:在建堂的事工中也是這樣.
主持人:讓神來作主.
主持人:讓神來帶領.
主持人:所以征戰不是我們為神而戰.

$^{361}$主持人:而是神為我們而戰.
主持人:我們弄清楚這一點是很重要的.
主持人:不是神來幫我們實現我們想做的事.
主持人:而是我們參與到神的計劃當中.
主持人:我們按照神的吩咐去做.
主持人:神為我們而征戰.
主持人:所以真正預備的是脫鞋謙卑.
主持人:放下我們自己的計劃和想法.
主持人:元帥的刀已經拔出來了.
主持人:我們只需要站在神的旁邊就可以了.
主持人:最後一點的總結.
主持人:今年我們教會的主題是「活出使命」.
主持人:盼望我們能夠在新的一年裡.
主持人:好像以色列人那樣進入征戰,可以彭順得業.
主持人:我們能夠在這個時代裡活出神給我們的使命.
主持人:在今天所講的經文裡.
主持人:其實是講到有割禮,是身份重建的一個很重要的步驟.
主持人:接著就是守節,逾越節.
主持人:是一個對舊恩歷史的回顧.
主持人:是提醒他們神是他們的救贖和供應者.
主持人:接著就是嗎哪的亭主.
主持人:嗎哪亭子就意味著他們需要成長.
主持人:不能夠停留在嬰孩的階段.
主持人:他們的生活都需要去更新.
主持人:接著就是元帥的顯現.
主持人:提醒我們需要將主權交託給神.
主持人:所以對我們有以下幾方面的提醒.
主持人:第一就是持續身份的更新.
主持人:我們需要不時停下我們的腳步.
主持人:透過聖禮,例如主餐和神的說話.
主持人:再次確認自己是被耶穌基督的寶血所賣熟的聖約的子民.
主持人:我們要脫去舊有的生命.
主持人:讓我們能夠有新生的樣式.
主持人:第二個要被提醒的就是.
主持人:要從被動到主動的成長.
主持人:信仰不能夠只停留在享受恩典的階段.
主持人:不能夠停留在嗎哪的階段.
主持人:神終是要我們成長的.
主持人:我們要承擔管理產業和傳揚福音的責任.
主持人:所以今天神要我們成為好的管家.

$^{401}$主持人:我們要成為管家去管理產業和傳揚福音.
主持人:實踐祂在這個時代給我們的使命.
主持人:最後提醒我們要認清爭戰的主權.
主持人:無論是屬靈的爭戰.
主持人:或者是個人的困境.
主持人:或者是事奉的挑戰.
主持人:我們要問的不是問神.
主持人:神啊,你可不可以幫我實現我的想法.
主持人:或者幫我實現我的夢想.
主持人:或者神啊,你可不可以幫我去克服這個困難.
主持人:不應該這樣去問神.
主持人:應該怎樣去問呢?.
主持人:應該問神啊,或者元帥啊.
主持人:你的戰略是甚麼呢?.
主持人:你的心意是甚麼呢?.
主持人:你的計劃是甚麼呢?.
主持人:我應該怎樣可以配合你?.
主持人:我應該怎樣可以參與在其中呢?.
主持人:這樣就不同了.
主持人:不是要神幫助我們成就我們心所想的東西.
主持人:而是我們參與到神的計劃,神的心意裡面.
主持人:我們與神同工.
主持人:所以神聖的使命是始於神在我們生命當中的更新.
主持人:成於我們對神的傳言的信服.
主持人:怎樣才可以完成這個神聖的使命呢?.
主持人:就是要更新和信服.
主持人:求主幫助我們,讓我們能夠在新的時代裡面可以活出使命.
主持人:我們一起同心祈禱.
主持人:持愛天父.
主持人:讓我們能夠更加見到你的應許和力量.
主持人:求你幫助我們在神聖使命當中能夠更新和信服.
主持人:以謙卑和感恩的態度跟隨你的指引.
主持人:在面對挑戰和困難的時候.
主持人:求你賜給我們持守信仰的力量.
主持人:相信你的應許永不落空.
奉主耶穌基督的成名祈求.
阿們.
\newpage



\section{雅歌 8:6-7}
\label{sec:031656}
\textbf{真愛律}~~~~賴若瀚牧師
\newline
\newline
連結: \href{http://www.church.com.hk/acms/content.asp?site=cdc\&op=show\&type=product\&code=031656\&layout=sermon}{\texttt{031656}} ~~~~ 語音日期: 2026-02-15
\newline
\newline
\hyperref[sec:031655]{\small{< < < PREV SERMON < < <}}
~
\hyperref[sec:index_chronic]{\small{[返順時目]}}
~
\hyperref[sec:031657]{\small{> > > NEXT SERMON > > >}}
\newline
\newline
雅歌 8:6-7
\newline
\begin{longtable}{cl}
\hline
\hline
章節 & 經文 (和合本修訂版)\\
\hline
8:6 & \begin{tabularx}{0.7\textwidth}{X} 求你將我放在你心上如印記，帶在你臂上如戳記。因為愛情如死之堅強，熱戀如陰間之牢固，所發的光是火焰的光，是極其猛烈的火焰。 \end{tabularx} \\ \\ \relax
8:7 & \begin{tabularx}{0.7\textwidth}{X} 愛情，眾水不能熄滅，江河也不能淹沒。若有人拿家中所有的財寶要換愛情，就全被藐視。 \end{tabularx} \\ \\
[1ex]
\hline
\hline
\end{longtable}
$^{1}$多謝五姊妹,大家好.
大家好嗎?.
好.
我們昨天過了我們的情人節.
其實應該翻譯做愛情節.
或者真愛節.
其實應該就叫做真愛節.
因為我覺得一年有一次.
來慶祝或者紀念.
我們人間最重要的一份感情.
就是愛.
愛是人間最重要的感情.
大家同意嗎?.
同意嗎?.
如果沒有愛.
這個世界會變得更冰冷.
如果沒有愛.
這個世界也會變得很黑暗.
冰冷就沒有溫暖.
黑暗就沒有光芒.
更加重要的就是.
可能連生命都沒有.
沒有生命力.
沒有生氣.
真愛在這個世界都很難找到.
因為這個世界是愛慾不分.
有很多時候.
愛和我們的慾是混在一起的.
我覺得廣東人很有趣.
你拿一樣東西給他.
你愛不愛?.
愛!.
就拿了過來.
我們不知道為什麼.
我們將愛變成了.
我喜歡,我想要.
這就叫做愛.
其實真正的愛不是拿過來.
真正的愛是給出去.
今天是情人節之後的第一天.

$^{41}$快要到農曆新年.
我的孫子十六年前出生.
他剛好是情人節那天.
也是剛好農曆新年的年初一.
我想很難碰到.
同一個日子落在同一天.
太多的人覺得愛是虛無縹緲.
很不切實際.
當講愛的時候就好像在天上的一些東西.
但是其實真愛不是一個哲理這麼簡單.
真愛是一些行動.
大家如果注意.
哥林多前書十三章.
那個愛的真諦.
其實全部都是行動.
全部都是action word.
就是行動的詞.
來表達.
我曾經在書店裡看到一些小書籤.
每一個都是愛事.
愛事是什麼呢?.
愛事是在你痛苦難過的時候握著你的手.
這是不是一個行動?.
愛是跟你分享一杯冰淇淋.
你咬一口我咬一口.
很浪漫.
愛是什麼呢?.
愛是用一整夜的時間.
為你排隊買一張車費.
愛是什麼呢?.
在我極度疲倦的時候.
我很累的時候.
仍然為你擠出一絲微笑.
大家知道有時候你很不想笑.
但為了它的緣故.
仍然擠出一些笑容.
這種叫做愛.
在剛剛這個週末.
我們剛剛從旅館過來.
我們有全部十五隊.

$^{81}$連同我們在內.
我們一起來參加啞歌夫婦營.
我們在啞歌夫婦營裡.
討論過啞歌書裡的四個定律.
四個大的定律.
這四個大的定律.
第一個叫做戀慕律.
然後贊悅律.
然後親密律.
然後解困律.
每個律都有五個小律.
所以我們已經說了二十個律.
讓他們回去.
我想今年也有很長時間可以練習.
而今早應趙牧師的要求.
我們就一起來崇拜.
我們就講啞歌書裡最重要的兩節經文.
就是第八章第六到第七節.
是講真愛律.
真愛律.
真愛的追尋.
其實在啞歌書裡.
兩節聖經講得非常清楚.
你能不能夠給啞歌書下一個很簡單的定義呢.
對不起.
對愛下一個簡單的定義.
愛其實是.
當然很抽象.
但是我們剛才講了.
愛是一個行動的詞.
愛我覺得它簡單的一個定義.
我看到一個是最合適的.
也最能夠讓我們行出來的.
應用的.
愛是什麼呢.
愛就是當對方的滿足和安危.
成為和我的滿足安危同樣重要的時候.
我是已經墮入了愛河裡.
我在愛的中間.
大家注意.

$^{121}$這裡有兩樣東西.
對方的滿足成為我的滿足一樣重要.
對方的安危和我的安危是同樣重要.
這兩樣如果加在一起.
一應用下去.
你就已經在愛的裡面.
譬如對方想吃一些東西.
你就說我偏不給你吃.
你覺得這個人對他是不是有愛呢.
除非是為了血糖高.
醫療原因.
如果沒有其他的原因.
你就是偏不給他吃.
那怎麼辦呢.
太多時候我們夫婦之間.
甚至朋友之間.
很容易就賭氣.
你喜歡嗎.
我就不給你.
我告訴大家.
距離愛還是很遠的.
譬如對方很晚都不回來.
或者一直都沒有打電話.
那你掛不掛心他呢.
如果你掛心他的安危.
我告訴你.
你是有他在你的心中.
也是有愛的.
他不回來.
你很擔心.
如果你心底裡有一絲的意念.
最好他不要回來.
我告訴大家.
危機的信號已經出現.
最好他從今以後都不要出現在我面前.
那就更加危險.
所以當對方的滿足和安危.
成為你的滿足和安危同樣重要的時候.
我告訴大家.
你已經是在愛的裡面.

$^{161}$我們現在說說真愛律.
有五個小律.
真愛律有五個小律.
這個在《雅歌書》第八章第六到第七節.
其實是在說整個《雅歌書》的一個總義.
一個結論.
一個約事.
求你將我放在你的心上如印記.
帶著你的臂上如銼記.
因為愛情如死之堅強.
窒恨如陰間之殘忍.
所發的電光是火焰的電光.
是夜幻的烈焰.
愛情縱水不能熄滅.
大水亦不能淹沒.
若有人拿家中的財寶來換愛情.
就全被藐視.
刻在心中如印記帶著臂上如銼記.
是什麼意思呢?.
這裡其實是在說印.
印就是印在他的心裡.
印在他的臂上.
心就是你內在的地方.
就是你感情最中樞的地帶.
而印在你的心裡.
就是說他已經刻在裡面.
不能夠磨滅.
然後第二個就是.
帶在你的臂上如銼記.
是什麼意思呢?.
帶在臂上如銼記.
就是讓人看見你對對方的愛.
就是在臂上就被人看見.
在心裡就是內藏於心.
所以這個其實就是佳偶的一個要求.
他要求什麼呢?.
他要求良人放他在他心裡.
他也要求良人.
也要帶他在他的臂上被人看見.
我們經常說有心心相印.

$^{201}$我們經常都聽見.
就是說心裡有對方.
而彼此之間兩個心是重疊在一起.
佳偶的要求.
我相信也是很多人一樣的要求.
譬如朋友.
譬如親人.
我們都很想我們所愛的人能夠放在我們心裡.
也將我們的愛表露在他們的生命當中.
約翰一書第三章第十八節說.
小子們,我們相愛不要只在言語和事頭上.
總要在行為和誠實上.
心是好的.
但要放在我們的行動上.
有一次有一個人來教會.
他是推銷保險的.
人家叫他信主.
他說我心裡信就行了.
我心裡信.
叫他洗禮.
他說不用了.
到他賣保險的時候.
有個弟兄就說.
我心裡跟你買就行了.
我不用出錢.
我心裡跟你買.
有心就無力.
有心但沒有外表的行動.
有個弟兄上完我的瓦戈.
他就寫在他的評價裡.
他說恩愛的內涵.
要有內藏的火熱之情.
也要有外顯的奔放流露.
我覺得他寫得很好.
內藏的火熱激情.
然後有外在的奔放流露.
我們需要有心,有行動,有口.
也要有實際的操作.
我們現在來看.
瓦戈裡面說的愛有五個特徵.

$^{241}$第一個,愛如死之堅強.
愛是堅強.
堅強是什麼意思呢?.
好像死一樣堅強.
你覺得死有什麼堅強?.
死有什麼堅強?.
其實死亡是無可抗拒的.
當死亡的神來敲門的時候.
你可不可以說我不開門.
我不去了,我不出去.
不可以的.
愛如死之堅強的意思是.
愛是無法逃避的.
我們每個人都要走到最後一站.
好像坐火車一樣.
坐到最後一站.
你說我不下車了,可以嗎?.
不可以的.
全部都要被你趕下車.
可能有些人早些下車.
有些人遲些下車.
但到總站都要下車.
愛的頑強就好像死亡一樣.
無可抗拒.
猶太人說世上有十二樣堅硬之物.
這是猶太人的智慧裡面所說的.
他說先是石頭.
但是鐵可以削石.
然後火可以溶解鐵.
到這裡我覺得都可以理解.
然後水可以將火熄滅.
這也是可以理解的.
然後他說.
雲可以蒸發水.
這就有點奇怪.
風可以將雲吹散.
人可以抵擋風.
恐懼可以使人消沉.
越講越覺得不太通.
不太有邏輯.

$^{281}$到最後那幾樣.
酒能驅除恐懼.
然後睡眠可以消除酒意.
這是用來講的.
但你知不知道他最後是說什麼呢.
他其實就是最後那兩句.
死亡比睡眠更持久.
即是你一睡下去就不醒.
那就沒有了.
然後死亡也不敵愛情的力量.
他最終是說什麼.
其實他講前面的都是為了鋪排最後那句.
死亡是不敵愛情的力量.
今天你說愛是很堅強的.
為什麼這麼多人的愛情或婚姻這麼脆弱呢.
什麼原因呢.
什麼原因我們的婚姻關係沒辦法持久.
沒辦法到永遠呢.
其實很簡單的一件事.
就是那個毀身.
那個毀身不夠.
毀身不足夠.
所以那個愛就會消失.
曾經有一個年輕的傳道人.
他要出道.
他第一次要主持婚禮的時候.
他就很緊張.
很緊張就請教一個老牧師.
如果我主持婚禮的時候忘記了東西.
下一步要做什麼我忘記了.
那你覺得要怎樣彌補呢.
老牧師就教他一招.
他說其實有時候都會忘記的.
如果可以就將你背過的聖經.
就將其中一節背出來.
保白,即是保住它.
然後你要找什麼,你要看什麼.
那個時候就會回來了.
這個年輕人就記得這個秘訣.
這一招.

$^{321}$他主持婚禮的時候真的忘記了.
忽然間他腦海裡就出了一節聖經.
他就隨口而出.
他說.
因為他們所作的,他們不曉得.
對著那對新人就這樣說.
當然這是一個笑話.
但是我又覺得真的.
有很多人步上聖壇.
他們所作的,他們真的不曉得.
如果他真的曉得的話.
他就不會這麼快就散了.
就不會這麼快就大家要離了.
所以要步上聖壇的時候.
我們又要知道我們在聖壇上所說的.
那個對方的誓言.
愛情婚姻不是一種買賣.
即是利益不足夠的時候.
我就退了,我就不做了.
愛情和婚姻其實不是一個普通的遊戲.
即是我不喜歡了,我沒興趣,我就不玩了.
不是,它是需要我們的毀身.
毀身就是以身相許.
無論怎樣都好,我們要從一而終.
這個毀身其實有很多樣的好處.
第一,它使兩方都有安全感.
即是不會經常害怕對方何時離開.
遇到第三者的時候.
就可以堅決地說不行.
不用考慮,根本不需要考慮.
愛神的毀身已經在這裡.
而可以因為這樣的毀身之後.
培養我們之間的感情.
而不是以利益或個人喜好為主的感情.
我相信這是夫婦之間很重要的一關.
這一關如果不守住就很難.
今晚我們會和這十三對夫婦再一次在主的面前重訂婚盟.
即是重一次再溫一下我們走上聖壇的時候.
我們到底說過些什麼呢?.
有時我們會忘記了.

$^{361}$所以愛的堅強在於什麼呢?.
在於我們對對方的毀身.
窒恨如陰間之殘忍.
突然之間你會覺得奇怪.
講愛講愛怎會講到窒恨呢?.
講到窒道和憤恨呢?.
和又講到陰間呢?.
你會覺得有點奇怪的地方嗎?.
窒恨如陰間之殘忍.
這個真愛的特徵是善妒.
即是說如果真的有愛.
在夫婦之間的那種愛會有一種躁意出來.
所以如果你的丈夫或太太對你有時有些吃醋.
其實是正常的.
你說正常的?我覺得是正常的.
識道的躁意是正常的.
什麼原因呢?因為他關心你.
因為他看重你.
在社交的場合裡面.
譬如有另一方注意到你的太太.
而且在獻殷勤的時候.
你如果一點都不關心一點都不著緊.
我覺得是有些問題的.
譬如當配偶問你.
你去哪了?你怎麼這麼晚回來?.
你又不要覺得你這麼注視.
其實他對你說這句話問這個問題.
其實是對你的關心.
他不是要限制你的自由.
而是他想知道.
而他也有權利去知道.
在社交媒體上.
有時我們丈夫或太太.
發現對方有異性的追求.
或是有些特別的注意感到不舒服.
我覺得這些都是正常的.
因為這裡說「疾恨如陰間之殘忍」.
在兩方的立約的關係裡面.
盟約的關係裡面.
會有時有這種「疾度之情」.

$^{401}$男女之間會不會有一種「柏拉圖式的愛情」呢?.
這個叫「Platonic Love」.
即是我純粹和你一種男女之間的友情.
不牽涉到男女的親密.
只是普通朋友.
可不可以呢?.
原則上是可以的.
不過實際操作上是有困難的.
因為愛到濃時.
或是朝夕相處到一個地步.
有時真的會把持不住.
所以還是要小心.
Don't play with fire.
因為「疾恨如陰間之殘忍」.
所以如果你明白愛的特質的時候.
有時你對夫婦之間的關係.
也可能明白對方為何會這樣表現.
有一個問題就是.
「疾恨如陰間之殘忍」和「愛是不疾度」.
有什麼關係和不同呢?.
哥林多前書不是說愛是不疾度的嗎?.
愛是不疾度.
這裡說到「疾度」也不得了.
那是什麼原因呢?.
是怎樣協調呢?.
其實這是對象的問題.
哥林多前書十三章的對象是所有的人.
但是雅歌書第八章的對象.
就是在婚盟裡立約的人.
就不同了.
愛是有疾度的成分.
第三就是愛是火熱的.
所發的電光是火焰的電光.
是耶和華的烈焰.
這句就將愛情和火放在一起.
愛情和電放在一起.
其實有很深刻的象徵意義.
火焰象徵愛情的強烈.
還有燃燒的能力.
我覺得是對的.

$^{441}$很多時候相戀的人.
大家建立了情感之後.
就有一種激情.
這種激情其實是正常的.
我經常說人彼此之間.
夫婦的愛有三個階段.
第一個階段就叫熱愛.
就是激情.
就是相戀和剛剛結婚沒多久.
但是過了一陣子.
你覺不覺得這種愛的火熱.
慢慢漸趨平淡.
我們中間有些都結婚二三十年.
三四十年.
我們夫婦已經過了五十年.
所以老夫老妻.
到中間那個階段.
不是叫熱愛.
就叫恩愛.
恩愛有沒有熱呢?.
有沒有火熱呢?.
有沒有?.
有.
但是那個熱沒有開始的時候.
那種激烈澎湃的熱.
你知道我們這個年紀.
我和我太太Esther.
我們是過了熱愛.
過了恩愛.
現在是戀愛.
同病相憐.
大家一起攜手步入我們最後人生的階段.
但是這個戀愛有沒有熱呢?.
有沒有?.
都有.
很可能就是人生漸趨平淡.
沒有那種激情的表現.
但是仍然是像電池一樣.
仍然可以運作的.
仍然可以讓你可以走的.

$^{481}$這種愛才會持久.
這裡是說一對夫婦之間的愛.
是耶和華的烈焰.
是火焰的電光.
是很激情的.
但是如果你看這段經文.
其實耶和華這個詞應該沒有出現.
我們和學本就將裡面的表達.
翻譯成神的名字.
如果你看英語的譯本NASV.
它其實是一種強調.
The very flame of the Lord.
它有說耶和華神.
但是新譯本和NIV都沒有耶和華神這個名字.
是非常猛烈的火焰.
和like a mighty flame.
意思是什麼呢.
其實我接受雅歌書裡面沒有提耶和華.
沒有提神的名字.
沒有.
除了這本之外.
還有以斯帖記也沒有.
沒有不等於愛不是從神那裡來.
只不過我們明白聖經的翻譯.
有不同的地方.
其實我們人在愛的進程裡面.
剛才我說了.
沒有開始那種澎湃並不重要.
但是你仍然要有一種熱力.
彼此之間仍然是互相扶持.
繼續往前走.
大家都知道.
電池如果你剛剛新買.
是很厲害的.
到一段時間.
如果有時有得有時有不得.
有時死了火.
要你去jump start.
這個電池很不靠得住.
死下死下.

$^{521}$這種電池就需要調整.
第三種就是電池死了.
一點電力都沒有.
如果我一點電力都沒有.
是不是要換個電池呢.
這個比喻就到這裡為止.
你不用換電池.
慕迪先生.
他曾經研究耶穌的生平.
有一次他要主持喪禮.
他想看看耶穌怎樣主持喪禮.
找來找去都沒有找到耶穌主持喪禮的步驟.
因為耶穌去的喪禮.
那個死人都復活.
所以如果你的電池沒有電.
不要緊.
你來到主的面前.
你說主啊.
你是那個讓人從死裡復活的.
從電池可以復生的.
你就將你的電池帶到主的面前.
所以沒有什麼喪事需要耶穌主持.
耶穌不會和你主持喪事禮拜.
每一個喪事禮拜來到主的面前.
都可以將死人復活.
我們無論怎樣都好.
我們的愛如果愛火已經慢慢熄滅的時候.
我們可以來到十字架的底下.
讓耶穌基督的愛再一次重新燃點我們.
第四 愛是排難的.
我們說愛是堅強的.
愛是窒道的.
然後愛是火熱的.
第四 愛是排難的.
愛情縱水不能熄滅.
大水也不能淹沒.
神話故事裡經常說白馬王子和美麗公主相見.
大家從此之後就樂也庸庸地.
過如下的一生.
實際的生活有沒有像神話故事裡所說的那麼順暢呢?.

$^{561}$沒有.
而雅歌書這一對年輕的戀人後來結婚之後.
他們都沒有對婚姻有一種特別的不切實際的憧憬.
他們都知道愛情裡是有困難的.
是會遇到大水的 會遇到縱水的.
所以他們不是一對不知天高地厚的男女.
他們也不是一對只憧憬羅曼蒂克的愛情.
而認為人生不會遇到困難.
我們昨晚在雅歌營裡處理解困律.
我們都覺得很特別.
聖經雅歌書提供了一些解困的原則.
我們在旅館裡前台的人離開我們聚會的地方很近.
他們經常聽到我們大笑.
笑聲很多.
我們有歡笑 也有眼淚.
昨晚我覺得笑聲很大.
很爆笑.
我們真的無法理解我們在討論解困.
其實應該是很嚴肅地去處理解困的問題.
我們可以存著一個很輕鬆的態度去面對我們人生的困境.
其實縱水和大水都是我們經常會遇到的.
每個婚姻都會經歷.
好像我們的小船經歷了人生的衝擊.
遇到這些人生的衝擊的時候.
你覺得如何能夠維持你的婚姻.
能夠經過驚濤駭浪呢.
我覺得就是你們兩個人要同心.
曾經有一對夫婦.
他們在南中國海的時候.
那艘船翻沉了.
他們兩個人就在那裡等待別人的援救.
差不多過了一晚.
後來救了他們出來的時候.
記者就訪問這個丈夫.
你在等待別人的援救的時候.
你心裡有什麼感受呢.
因為那裡有鯊魚出沒.
隨時可能沒有人知道.
也可能救不了也說不定.
這個丈夫就說.

$^{601}$我有很矛盾的感受.
第一,我很怕.
第二,我很開心.
很怕大家都明白.
記者就說.
為什麼你在這個時候這麼開心呢.
他說因為我太太從來都沒有這樣緊緊的抱著我.
所以他很開心.
所以是遇到困難的時候.
你才看到大家之間那種感情的流露.
所以從來沒有經歷到太太那種緊緊的抱著他.
他覺得這個困難也是值得的.
所以昨晚我們說.
其實如果困難你能夠兩個人處理.
你們的婚姻關係可以更上一層樓.
不單單是婚姻關係.
其實也可以是朋友和親人的關係.
第四,有四隻船.
其實人生要經歷四隻船.
第一隻船叫做友誼的船.
第二隻船叫做定情的船.
第三隻船叫做困難的船.
就是這段經文所說的.
你猜第四隻船是什麼呢.
我故意不放進去.
你猜猜第四隻船是什麼呢.
用英文來表達.
就是主權的船.
你的人生之所以能夠有突破.
你的人生之所以有前景.
有光芒.
就是因為你們將你們的主權交託給主耶穌.
就是Lordship.
Friendship, Courtship, Hardship and Lordship.
所以在濃情蜜意的關係裡當然是好.
但很多時候婚姻是需要配合排難解困的愛.
最後一樣我們來看看.
若有人拿家中的財寶來換愛情.
就全被藐視.
愛情有沒有價碼的.

$^{641}$有沒有啊.
所以愛情是無價的.
如果你看這段經文.
經文裡面所說的.
最後也是最高峰的.
如果有人拿家中的財寶來換愛情.
他是全被藐視的.
這個在希伯來文的表達是這樣的.
若有人拿家中所有的財寶來換取愛情.
他所得到的是什麼呢.
就是被藐視.
愛情是沒有價碼的.
愛情不是很便宜的.
大家不要搞錯.
它不是便宜,它是無價的.
沒有辦法可以定一個價碼在上面.
曾經有一個青年人叫王大夫.
他向自己的女朋友小玲求婚.
她說小玲,我很愛你,你嫁給我可以嗎.
雖然我沒有錢大莊那麼多錢.
有華夏,有跑車.
但我愛你是全心全意的.
小玲就回應了.
她說大莊我也很愛你.
不過你可不可以介紹錢大莊給我認識.
這個叫做有錢人終成眷屬.
這個世界是不是這樣的.
不是有情人,是有錢人終成眷屬.
有錢人可能辦的婚禮或找到的對象比較風風光光.
很合乎理想.
但是我告訴大家.
有錢人也有有錢人的苦.
有錢人也不是因為有錢就得到真正的愛情.
大家都可能聽到金錢不是萬能.
他可以買到物質享受但買不到心靈健康.
他可以買到山珍海味但買不到你的胃口食慾.
他可以買到各樣的保險但他不能保你內心平安.
他可以買高樓大廈但不能買到家庭美滿.
他可以買到一紙婚書.
可以大家一起結婚但他買不到真摯的愛情.

$^{681}$愛是無價.
如果用家中的財寶來換愛情是全被藐視.
是很可笑.
不知道大家之前有沒有看過一齣電影.
叫做The Indecent Proposal.
叫做不道德的交易.
這部電影我很早期坐飛機看.
是關於金錢愛情的抉擇和信任的一個很簡單的故事.
但他探討人性和道德界線和社會價值觀的一個預言.
我覺得這齣電影他拍的時候.
他不只是討論一些男女之外.
他是討論一個很深層的價值觀的矛盾.
他講一對夫婦Diana和David.
他們因為經濟的缺缺所以去了拉斯維加斯.
試試賭一下能不能贏回那些錢.
在那個時候他們遇到一個富翁.
這個富翁叫John.
這個富翁看中了Diana.
就和他做了一筆交易.
他用一百萬來買Diana和她相處一晚.
一百萬.
他們兩夫婦就在掙扎.
因為他們經濟的問題.
所以丈夫David就開始贊成.
真的住了一晚.
住了一晚之後.
David就一直問那晚發生了什麼.
那晚你們兩個做了什麼.
這個太太Diana就一點都不說.
兩個人後來就造成很大的矛盾.
然後他們拿到那一百萬之後.
David就覺得這一百萬都沒什麼用.
好像失去了他和太太的親密和信任.
所以在拍賣的時候.
他一直和另一個人鬥出高價.
然後那一百萬就財散人安樂.
沒有了.
錢也沒有了.
大家的婚姻彼此的信任也沒有了.
那晚不知道做了什麼.

$^{721}$也沒有實際知道那晚做了什麼.
但他們兩個人的感情已經崩潰了.
所以大家決定各走各路.
你覺得是不是很悲哀.
但這齣戲有一個很詩意的結束.
你看看右邊的碼頭.
就是他們一開始相識的地方.
他們兩個人其實大家都很失落.
大家都覺得人生竟然走到這個地步.
所以不約而同地就走回到他們原來相遇的碼頭.
大家一抬頭的時候看到對方.
這樣就結束了.
你猜最後會是怎樣呢.
等你們來猜.
等於一百萬來買愛情.
你會不會出呢.
你會不會接受呢.
這個其實是愛情和信任和道德的界線.
一個很特別的問題.
我覺得對人性的真實可以揭露出來.
真愛是不能賣的.
你用錢買,買不到.
正如神對我們的愛.
也沒有辦法你用錢去買.
主耶穌甘願為你和我寫生在十字架上.
如果你愛的電池現在缺能力的話.
你來到主耶穌基督的十字架面前.
所以一會兒我們會有一個分享.
我們會用會點一個房間.
我們會有一個聖餐.
當我們的婚姻感到無力的時候.
我們來到主的面前.
因為主耶穌基督就是愛的源頭.
神就是愛.
我們能夠相愛是因為祂先愛我們.
而我們的生命可以從神的愛那裡再一次燃點.
然後重新出發.
Happiness is a choice.
喜樂是一項抉擇.
同樣,Love is a choice.

$^{761}$曾經有書寫,愛是一個抉擇.
你決定要和對方繼續靠主維持下去.
愛是一個抉擇.
你願不願意用神聖的愛去愛你的配偶呢?.
去愛其他在你身邊的人呢?.
愛是一個抉擇.
仇恨也是一個抉擇.
你要選擇哪一樣呢?.
我們低頭一起祈禱.
天父,我們獻上我們的感恩在你面前.
有十五對連同牧師,師母一起.
去研討,去大家一起談論.
關於《雅歌書》裡面的婚姻和愛情的訊息.
求主你保守我們能夠走完人生的路程.
手牽手,心連心.
來走在主你託付的路上.
保守我們,帶領我們.
我們今天的時間交託在你的手.
奉耶穌基督的聖名祈禱.
阿們.
\newpage



\section{啟示錄 19:11-16}
\label{sec:031657}
\textbf{白馬之王}~~~~陳崇基牧師
\newline
\newline
連結: \href{http://www.church.com.hk/acms/content.asp?site=cdc\&op=show\&type=product\&code=031657\&layout=sermon}{\texttt{031657}} ~~~~ 語音日期: 2026-02-15
\newline
\newline
\hyperref[sec:031656]{\small{< < < PREV SERMON < < <}}
~
\hyperref[sec:index_chronic]{\small{[返順時目]}}
~
\hyperref[sec:031658]{\small{> > > NEXT SERMON > > >}}
\newline
\newline
啟示錄 19:11-16
\newline
\begin{longtable}{cl}
\hline
\hline
章節 & 經文 (和合本修訂版)\\
\hline
19:11 & \begin{tabularx}{0.7\textwidth}{X} 後來我看見天開了。有一匹白馬，騎在馬上的稱為「誠信」、「真實」，他審判和爭戰都憑著公義。 \end{tabularx} \\ \\ \relax
19:12 & \begin{tabularx}{0.7\textwidth}{X} 他的眼睛如火焰，頭上戴著許多冠冕；他身上寫著一個名字，除了他自己沒有人知道。 \end{tabularx} \\ \\ \relax
19:13 & \begin{tabularx}{0.7\textwidth}{X} 他穿著浸過血的衣服；他的名稱為「神之道」。 \end{tabularx} \\ \\ \relax
19:14 & \begin{tabularx}{0.7\textwidth}{X} 眾天軍都騎著白馬，穿著又白又潔淨的細麻衣跟隨他。 \end{tabularx} \\ \\ \relax
19:15 & \begin{tabularx}{0.7\textwidth}{X} 有利劍從他口中出來，用來擊打列國。他要用鐵杖管轄他們，並且要踹全能神烈怒的醡酒池。 \end{tabularx} \\ \\ \relax
19:16 & \begin{tabularx}{0.7\textwidth}{X} 在他衣服和大腿上寫著「萬王之王，萬主之主」的名號。 \end{tabularx} \\ \\
[1ex]
\hline
\hline
\end{longtable}
$^{1}$大家好,各位馬年快樂.
在這裡先去拜年,今天年廿幾呀?.
廿八,是嗎?.
糟了,我搞錯了,我知道今年沒有年三十的.
所以星期二就是初一,大家是不是放假?.
真的放假?.
美國來的,請假放假,OK.
自己放自己假.
如果你是馬年的話呢?.
誰屬馬年的?.
有喎,OK,屬馬年的,我們特別恭喜你,馬年快樂.
我跟弟兄姊妹聊天的時候,我發現年輕人是不會背十二生肖的.
有誰會背十二生肖的?.
有誰會背?.
這麼少人?.
下一年是什麼?馬年之後是什麼?.
羊嘛,是嗎?怎麼數?.
有誰不會數十二生肖的?.
這麼多,那些叫做年輕一點的.
很容易的,是嗎?.
鼠牛虎兔龍蛇馬羊猴雞狗豬.
是嗎?.
鼠牛虎兔龍蛇馬羊猴雞狗豬.
很容易的,單字這樣背,是嗎?.
好像,創出利民生,書士得殺王.
舊約目錄都是這樣背的,怎麼會不會背?.
如果你不會背,請你去華埠廣場轉一圈,好嗎?.
沒理由不會的,你是華人,是嗎?.
說到馬年,馬就是講速度和成功和那種衝勁.
我們希望新一年都是盼望有些突破,是嗎?.
個個都是事事如意的.
特別是我們新一年,新的馬年裡面.
我們有很多好的消息,是嗎?.
還有我們非常之有盼望的.
譬如,接下來星期六.
CMA就會上到多倫多那裡參賽.
那個卡拉OK的大賽.
我們希望他馬上成功,好嗎?.
希望他馬上成功.
上個星期,我去了一個教會.

$^{41}$有一個小的年宵.
跟團體一起去.
他們其中一個攤位是寫揮春的.
寫什麼呢?找來找去寫什麼呢?.
有人說寫成會.
成會搭馬,好嗎?.
萬馬奔騰,成會搭馬,好嗎?.
我寫的那些字.
很噁心,是嗎?.
那些漂亮的藝人寫的那些,我懶得放上來.
左邊那些,有人叫我寫馬上有錢.
因為Percy幫我作了個花名.
馬年,對了.
馬上,即是Mark Sir.
Mark Sir有錢,Mark Sir有錢,Mark Sir有錢.
太好了,太好了,太好了.
因為新年的緣故.
我也選了一節經文.
是跟馬有關的.
聖經裡面提到很多馬.
除了搭馬之外.
這段聖經是來自啟示錄.
啟示錄第十九章提到一隻馬.
馬上有一個君王.
我們不妨用經文來思想一下.
到底馬年裡面有什麼值得我們反省的.
特別生命的方向,生命的節奏等等.
這裡提到一隻白馬出現.
白馬上有一位君王坐著.
然後用了很多描述來告訴我們.
這個君王是值得我們敬拜和跟隨的.
但是我們打開啟示錄的時候.
特別是新征頭,我們要先弄清楚.
因為很多時候還沒接觸過聖經.
或者我們初初接觸聖經的時候.
我們不太熟悉聖經到底在說什麼.
和聖經裡面不同寫作的方法.
我們就會當作啟示錄做什麼.
就是通勝.
以前有沒有看過通勝.

$^{81}$有吧.
今年會不會行運.
或者將來,未來有沒有犯太歲.
要做些什麼等等.
那些紅色一本的,有沒有看過.
有啊,現在沒有一本了.
現在是下載在APP上看的.
當然那些是看上腦,騙你十年八年.
那些東西是不值得相信的.
但是很多時候我們都會當啟示錄.
就是預測未來,預先寫好的報紙.
告訴我們世界末日到了.
或者世界末了就快到了.
會有什麼事情發生.
然後我們就嘗試去做一些配對.
這裡提到一些符號.
提到一些異象.
和我們今天在報章或新聞上看到的異象.
是否符合呢.
如果是符合的話.
那一定是啟示錄裡面說的話.
一邊展示一邊揭曉.
然後告訴我們是對的.
好像預測未來那樣.
其實不是這樣的.
完全不是這樣的.
所以今天你聽了什麼都好.
啟示錄就不是通勝.
明白嗎.
啟示錄不是通勝.
不要當它是通勝.
當然它有說未來的東西.
告訴我們耶穌一定會再來.
這位君王一定會將新天新地帶來.
公義真理它將會帶來.
這個是我們盼望.
但是很多時候不是我們現代人這樣說.
或者那種印象.
就是預先寫好的報紙.
啟示錄有它的背景.

$^{121}$啟示錄是當時第一世紀的時候.
特別寫給一些小數的教會.
他們當時在羅馬帝國底下.
他們受到很多壓迫.
受到很多打壓.
受到很多壓力.
於是他們有時候灰心.
有時候他們失去盼望.
啟示錄就是用一個這樣的方法.
告訴他們.
鼓勵他們.
特別是用一些他們當時熟悉的image.
當時他們熟悉的符號.
象徵.
來表達出上帝仍然是掌權的.
所以是一本鼓勵的書.
是一本牧羊的書.
不是預測未來.
大家知道在第一世紀的時候.
我們經常說.
新約的背景.
就是當時羅馬帝國非常富強.
他們會用很多的方法來維持住.
帝國裡面的太平和秩序.
特別是啟示錄的背景.
那些教會.
他們很多時候在今天的東地中海的地方.
特別是今天的土耳其.
巴勒斯坦地等等.
他們受到很多的壓力.
因為羅馬帝國要去到.
維持住那個太平盛世.
為了維持住羅馬帝國所塑造出來的和諧.
就叫那些異見分子不要出聲.
或者打壓他們.
不准他們說任何話.
不准他們說反羅馬的話.
或者去挑釁羅馬帝國政權的東西.
所以在東地中海.
即是今天的土耳其巴勒斯坦地等等的地方.

$^{161}$他們就會推行一種叫做.
帝王的崇拜.
帝王崇拜.
即是國家領導人有個廟.
你就一定要去供奉他.
於是有很多的經濟壓力.
對當時的人來說.
因為很多時候他們的經濟和宗教.
是完全混為一談的.
不能分開的.
我們就可以分開.
以前是不能分開的.
你要去購物.
或者去做生意.
或者你要在商務裡面有順利.
你就需要去拜一拜那個國家領導人.
你要去那個廟.
向他供奉也好.
或者向他上香也好.
那你才會有人脈.
你才會有生意.
如果你不這樣做的話.
你就會失去很多的保護.
所以啟示錄就是給當時的教會.
來對他們講說話的.
因為當時的教會.
他們相信的是耶穌基督才是主宰.
如果有個國家領導人.
他們要去供奉他的話.
那就成為一種很大的.
生命裡面的張力.
那怎麼辦呢?.
我們會蝕底.
怎麼辦呢?.
我們受到那麼多壓力.
我們受到那麼多的迫害.
那怎麼辦呢?.
啟示錄就是用一個這樣的方法.
去跟他們展現出.
真正作王掌權的是哪一位.

$^{201}$所以啟示錄裡面.
我們看這些經文的時候.
我們需要有一個這樣的背景.
怎樣鼓勵到他們呢?.
我們先看看這段經文.
這段經文我們講過了.
是講馬的.
右邊那裡有一個空了的框.
當我們讀這段經文的時候.
請你想像一下.
到底他所描述的.
所形容的.
是會怎樣的呢?.
讀一讀好嗎?.
然後你讀的時候.
那些意象.
或者那些image.
你會想像到一個怎樣的畫面呢?.
啟示錄十九章十一節.
後來我看見天開了.
有一匹白馬.
這匹白馬騎在馬上.
是稱為誠信.
真實.
他審判征戰都憑著公義.
他的眼睛如火焰.
頭上戴著許多冠冕.
他身上寫著一個名字.
除了他自己.
沒有人知道.
他穿著浸過血的衣服.
他的名稱為上帝之道.
眾天君都騎著白馬.
穿著又白又潔淨的細馬布.
跟隨他.
有利劍從他口中出來.
用來擊打列國.
他用鐵匠管轄他們.
並且要串全能上帝.
烈怒的駕走池.

$^{241}$在他衣服和大腿上寫著.
萬王之王.
萬主之主.
你想像到一個怎樣的圖畫.
想像一下.
這些是文字的描述.
到底那個圖畫會是怎樣呢?.
那個圖像會是怎樣呢?.
如果你在AI那裡.
他出了一張這樣的畫給我.
你覺得是否合理?.
白馬.
口中有利劍.
然後眼如火焰.
穿著的衣服是甚麼?.
浸過血的.
有血沾上的.
這裡沒有寫萬王之王萬主之主.
大概應該想像到一個這樣的意象.
這種意象看下去的時候.
雖然現在是AI generate出來.
但當時大概應該看到一個這樣的畫面.
並且我們講到背景是怎樣的.
當時的人在羅馬帝國底下受壓迫.
在羅馬帝國的街上走走.
你大概會經常看到一個這樣的影像.
無論是石雕,浮雕,長畫等等.
你都會經常看到有白馬.
或者有馬匹戰馬.
因為在羅馬的文化裡面.
羅馬的帝國.
他們是怎樣建國的?.
是用軍事建國的.
好像今天用飛機大炮一樣.
每次一個地方被羅馬佔據了之後.
佔領了之後.
將那個地方打殘了之後.
那些軍隊就會海船回歸回京都.
或者各大城市.
就會有遊行,有儀式.

$^{281}$中央電視台就會宣佈一件事.
就是羅馬人多威風.
帝國是必成的.
帝國既然是必勝的話.
你每一個在帝國裡面的人.
就需要就範的.
你一定要跟他.
一定要拜他.
這個就是普遍羅馬的市民.
在街上會看到的形象.
但是啟示錄將這個形象顛倒過來.
因為他形容的是誰?.
他形容的不是羅馬將軍.
也不是羅馬帝王.
他形容的騎著白馬上面的.
是另一位.
這一位是怎樣的呢?.
他的名號你看到嗎?.
是誠信,真實.
Faithful and truthful.
還有公義有justice.
羅馬帝國不會說這些的.
因為羅馬帝國說什麼?.
誰權頭大.
就是誠信,真實和公義.
什麼叫公義?.
我打殘你就是公義.
什麼叫真理?.
我摒過金的你.
我就是真理.
這個叫真理.
這個叫誠信.
和我們今天的帝國都是一樣.
無論你說什麼國家都好.
你說英國也好.
中國也好.
美國也好.
什麼國家都好.
我們都是說權頭的.
但是這一位坐在白馬上面的啟示錄說.

$^{321}$這一位最終會得勝掌權一切的君王.
他的名號.
他的權頭,他的名號.
是誠信,真實和公義.
這個告訴我們到底哪一位才是真正的君王.
啟示錄是做什麼的?.
啟示錄是要揭穿一些東西的.
你看一些意象.
看一些image的時候.
就是去review那個reality.
就不是讓你看外殼這麼簡單.
好像洋蔥這樣慢慢慢慢撕開的時候.
你就明白裡面原來是爛的.
他就是想揭示告訴你在地上的帝國.
你慢慢撕開每一個層次的時候.
其實天下烏鴉是一樣黑的.
無論他打著什麼的旗號.
語言上面也好.
image上面也好.
其實本身裡面都是腐爛的.
因為這個是帝國的權利.
這個是人的腐敗的權利.
如果你說要找世間上.
有些我們相信,能夠相信,能夠去跟隨的.
就一定不是地上帝國的權柄.
千萬不要搞錯這件事.
啟示錄要一直問一個問題.
到底你的生命是向著哪一個方向?.
你的生命追趕的方向是什麼?.
追求的方向是什麼?.
今天我們都有很多這些熟悉的畫面.
無論剛才說到的一些帝國.
political power也好.
有很多熟悉的畫面就是.
在我們生命裡面.
我們都有不同的吸引我們跟隨的東西.
無論是一些我們的事業,金錢.
或者生命裡面的安逸.
我們想去control一些東西.
去控制一些東西.

$^{361}$啟示錄就是想去揭示這件事.
它們本身是不值得我們將生命交上的.
真正你要將你生命的賭注放在哪裡?.
放在這一位上.
其中真正白馬.
真正作王的這一位.
這一位是誰?.
當然就是耶穌基督.
這裡說到祂有許多的冠冕.
當時的王.
古代的王.
包括羅馬帝國.
佔領了一個地方的時候.
他們就會在那個地方.
然後拿了本地君王的皇冠.
然後就帶了上去自己那裡.
有些故事是這樣說的.
將軍去到一個地方.
羅馬王去到一個地方.
他們就會帶幾個冠冕在頭上.
代表什麼?.
是我已經戰勝了你.
表明我對你們的權柄.
你們是在我之下.
如果你熟讀啟示錄的話.
你就知道.
啟示錄裡面介紹過一些人物.
是有很多冠冕的.
我們現在看的是第十九章.
差不多結束了.
在頭十八章那裡.
已經介紹過一些人物.
是有不同的冠冕的.
我們之前說的啟示錄.
是介紹過幾個所謂叫做.
Unholy Trinity.
就是邪惡的三一.
就是有三大邪惡的化身.
是用不同的symbol去表達的.
第一個就是龍.

$^{401}$第二個就是從海裡面出來的怪獸.
第三個就是地出來的怪獸.
在第一個怪獸出現的古龍.
龍Dragon.
它出現的時候.
啟示錄就去到揭示.
這隻龍是很厲害的.
怎麼厲害呢?它頭上有七個冠冕.
戴著七個冠冕.
就是說它掌管了所有東西.
原來它的身份是誰呢?.
就是撒旦.
接著到第十三章.
有隻怪獸從海裡面出來的.
為什麼從海裡面出來呢?.
因為以前近東古代中東的人.
他們會覺得海是一些混沌的東西.
是不能戰勝的,很混亂的東西.
凡是從海裡面出來的東西.
十該其九都是不好的.
所以他們就說在海裡有一隻怪獸走出來.
而那隻怪獸就是他們平時見到的羅馬帝國.
他們所接觸的羅馬帝王.
外表是很光鮮的.
但是它所揭示的原來是這隻怪獸.
並且是在混沌出來的.
第十三章第一至十節就說到這個怪獸.
原來就是羅馬帝國.
它有十國七頭.
並且有十隻國.
每一隻國就有一個的冠冕.
有十個的冠冕.
並且它的名號是一個褻瀆的名號.
因為自稱為神.
自稱為霸主.
所以說到白馬上面有許多冠冕的時候.
它就說其實真正的冠冕在誰那裡呢?.
在耶穌基督.
在耶穌基督的頭上.
既然耶穌已經有許多的權柄.

$^{441}$所有眾多的冠冕都戴在祂頭上.
地上的權柄就會突然縮小.
它就不再是可以哈哈爸爸.
不再是有權柄的.
我們繼續看看其他異象.
浸過血的衣服.
這些羅馬市民經常都會看到.
因為海船回歸的時候.
將軍士兵手上總會有些什麼?.
血.
無辜人的血.
因為他們哈哈爸爸統治四方的時候.
回來的時候.
他們手上,衣服上就會有無辜人的鮮血.
你覺得這些很明顯.
很鮮明的一種形象.
這就是啟示錄所說的.
為什麼耶穌基督會有沾了血的衣服呢?.
你可以哈哈爸爸.
然後有無辜人的血沾上你的衣服.
但同樣也可以犧牲自己.
流自己的血.
以致你能夠拯救到其他人.
這就是十字架上的耶穌.
祂是有沾了血的衣服.
但是祂做的是用犧牲的血來得勝.
去擊敗邪魔.
去拯救我們.
又有什麼?.
有利劍.
這也是當時羅馬帝國最威風的事.
是他們帝國軍隊的標號.
兩人的利劍.
有武器.
他拿著利劍的時候.
他就可以東征西討.
但這裡說到這位騎著白馬的耶穌基督.
他也有利劍.
但他的利劍不是去廝殺.
他的利劍是什麼?.

$^{481}$上帝的道.
他的利劍是從口裡出來的.
即是說他所廝的福音.
用福音.
用十字架的福音.
來戰勝四方.
不是用利劍.
不是用廝殺.
也不是用暴力.
不是用恐嚇.
耶穌基督去征服世界.
改變世界.
是用真理,用愛.
去感動我們.
剛才我們提過.
受洗.
受洗的時候.
弟兄姊妹.
未下水池之前.
他們也會講見證.
有沒有誰聽過一個見證是.
我信主因為有人用槍指著我.
所以現在我要受洗.
每一個生命都是經歷過.
上帝對他的恩典.
我們不配這個恩典.
所以我們回應這個恩典.
是一個以愛還愛的道理.
這個道理絕對不是地上帝國的道理.
但這是天國的道理.
這就是耶穌基督.
是他可以改變世世代代.
千千萬萬的人.
以羅馬帝國是不能夠.
最後一個名號.
在衣服上,大腿上.
是有萬王之王,萬主之主.
一聽的時候.
用當時的背景.
一聽這些名號的時候.

$^{521}$第一個感覺就是.
這個一定是羅馬帝王.
因為羅馬帝王.
在他的錢幣或石碑上.
經常說他就是king of kings and lord of lords.
因為雄霸四方的時候.
還有誰可以被他打.
他就是萬王之王,萬主之主.
每一個都是這樣的代號.
羅馬帝王要這樣自稱的時候.
就是想人做什麼.
就是被他無條件的效忠.
You pledge your allegiance to me.
你向我無條件效忠.
這就是帝國地上要我們做的事.
但為我們死的是耶穌基督.
他說如果你應愛也會回應我的話.
我才是你值得無條件效忠.
因為我無條件地救你.
無條件地拯救你.
所以你看看.
啟示佬最後介紹這位白馬上的君王的時候.
他就介紹了一位信實,真實,公義的救主.
他的德性是最忠的.
他是值得你和我無條件效忠.
Unconditional allegiance.
對我們今天來說.
我覺得最重要的那句說話.
是這段第十四節.
第十四節說什麼呢.
他說這隻白馬來到的時候.
後面有很多白馬跟著.
看到嗎.
不是只有一匹.
後面有無數匹白馬跟著.
白馬跟著的騎上白馬的人.
後面跟隊的那群人.
穿著什麼衣服.
又白又潔淨的西馬布去跟隨他.
所以你聽到那些馬聲的時候.

$^{561}$他不只是一隻馬.
是萬馬奔騰.
但後面跟著的很多白馬.
我們第一個感覺就是後面白馬的戰士.
又白又雪白又有西馬衣.
這些大概應該是天使吧.
應該是Angels或Angelic Beings等等.
確實可以是.
但如果看回其他經文的話.
似乎告訴我們這些不是天使.
在這個意象裡面.
看到白馬上的軍王後面跟隊的那群騎著白馬的.
似乎就不是天使.
如果看17章第14節的話.
就提到有一些人.
是和這位被殺的高揚耶穌基督一起作戰的.
並且他們是一群怎樣的人呢.
是這樣形容的.
是夢召被揀選.
眾心與上帝福音的人.
這一群人就是跟著他作戰的.
以什麼作戰呢.
不是飛機大炮.
是以福音來作戰.
第十九章就提到教會.
怎樣形容教會呢.
教會就是身穿明亮潔白西馬衣的.
這一群人有西馬衣.
西馬衣代表就是聖徒的義行.
就是我們的righteous act.
我們的acts of kindness.
acts of goodness.
我們在生命裡面.
在家庭裡面.
在我們身邊的人裡面.
去到行的義行.
如果你這樣看的時候.
你會發現第十九章說的.
最終要來得勝的軍王耶穌基督.
跟著他的那群人.

$^{601}$眾天君跟隨他的.
原來就是中心的信徒.
就是基督徒.
他問一個問題.
就是告訴我們.
如果我們跟隨這位基督的話.
跟隨這一隻白馬的話.
我百不得我是其中一個.
是一個隊裡面其中一個.
是跟他的節奏行.
跟他的方向行.
是十四節.
似乎就是一個邀請.
invitation.
就是你是不是那個.
其中.
眾天君.
騎著白馬.
跟隨這位真正的軍王耶穌基督.
到底你的身份是什麼.
你屬不屬於這個天君裡面.
其中一員.
如果是屬於的話.
我們值得將我們的命交給他.
值得將我們的決定.
我們任何的意願.
我們盼望與利益跟他align.
跟他合作.
如果你這樣看的時候.
這本書啟示錄.
或者介紹這個白馬之主的時候.
其實就是帶給當時的人.
和我們今天的人.
心裡面有很多安穩.
我們現代的生活裡面.
特別這一代.
這幾年.
我們心裡面有很多拉力.
很多不同的張力.
有時候我們都不知道自己.

$^{641}$在跑去哪裡.
不知道在看什麼.
不知道在說什麼.
總之我們就很忙.
我們都沒有時間去想這些事.
對不對.
啟示錄就是叫我們.
停一停下來.
到底你.
跑去哪裡.
你走去哪裡.
心裡面有很多安穩.
近這幾天我自己.
在不同的情況.
很多人提醒我.
我自己之前都有說過.
我發現.
我最近.
很趕.
不知為什麼很趕.
有人提醒我的時候.
我發現其實我在趕什麼.
為什麼每件事都要這麼快.
為什麼每件事.
都好像很急很倉促.
所以.
有些人很關心我.
突然就跟我說.
其實你這麼快.
你會miss了很多東西.
miss了很多東西.
沿路的風景你全部miss了.
你自己的狀況.
可能你自己都不明白.
我一向覺得.
效率.
做事快做事趕.
是我引以為傲的.
對不對.
Someone has to get things done.

$^{681}$但最近我自己就.
被人提醒.
特別昨天.
在social media.
有一個朋友.
他post了一段很簡單的.
有關.
安息.
休息.
他說我們很多時候.
很快很趕.
停不下來靜不下來.
我們說話很快.
說話很倉促.
有時候傷害到人.
做事很快有時候急忙.
如果我們24/7這樣看下去.
其實主要的原因是什麼.
我那個social media的朋友就說.
主要的原因就是.
因為你不相信.
上街會工作.
你想控制.
我身後跑河.
到底我的生命的方向.
我的節奏.
是為了什麼.
到底我跟隨的是誰.
我想新幾十年裡面.
啟示錄要問的.
不只是要我們馬到工成.
不是每樣東西要急趕.
快速.
辦事要快速成功要快速.
我們不是.
Who are you following.
你再跟隨誰.
但我們跟隨這位主的時候.
我們有.
有那種意願.

$^{721}$我們是跟他的節奏.
因為如果這位君王.
是最終得勝的.
起碼我們的心境.
我們會有安穩.
如果是這樣的話.
在我們所有的關係裡面.
其實我們的說話.
我們待人接物.
其實都會有所改變.
對我們身邊的人說話.
可能我們會放慢.
放緩一點.
因為最終不需要我控制.
上帝仍然作供.
在我待人接物.
處理其他人的時候.
我可以和平很多.
因為這個世界有很多戰火.
很多爭戰.
但我知道這位白馬上的君王.
是必定會在來的.
我心裡面有安穩.
我都有勇氣.
跟著上帝的節奏.
有平安有喜樂.
可以同心低頭祈禱嗎.
可以低頭祈禱的時候.
想一段安靜的時間給大家.
在上帝面前你可以說.
最近到底我自己在幹什麼.
忙什麼.
求主讓我能夠安靜.
跟隨祂的節奏.
和祂一起事奉.
主耶穌基督.
你是那位真實.
誠信公義的主.
你是那位萬王之王.
萬主之主.

$^{761}$求主你帶領我們.
也讓我們跟隨你的節奏.
不是跟隨世界的節奏.
賜給我們勇氣.
賜給我們這份安穩.
讓我們跟隨你的時候.
我們有平安有喜樂.
並且有活力.
我們信靠你.
帶領著我們每一個個人的生命.
我們的每一個關係.
從我們自己和自己的關係.
身邊的家人.
工作,教會,社會.
跟著你這個.
信實.
真實的節奏.
去跟著你.
相信你.
讓我們的生命能夠更加經歷到你的特性.
帶領著我們.
如果我們太趕的時候.
讓我們在你的慈愛裡面能夠停下來.
讓我們知道你愛我們.
並且在背後不斷作供.
主叫我們的生命.
能夠不斷跟隨你.
我們巴不得每一個.
都是在天君裡面的一份子.
穿著你所潔淨的衣服.
事奉敬拜你.
祈禱奉傳明教.
恭祝各位馬上平安.
\newpage



\section{羅馬書 15:14-33}
\label{sec:031658}
\textbf{以行動活出因信稱義}~~~~袁惠鈞牧師博士
\newline
\newline
連結: \href{http://www.church.com.hk/acms/content.asp?site=cdc\&op=show\&type=product\&code=031658\&layout=sermon}{\texttt{031658}} ~~~~ 語音日期: 2026-02-15
\newline
\newline
\hyperref[sec:031657]{\small{< < < PREV SERMON < < <}}
~
\hyperref[sec:index_chronic]{\small{[返順時目]}}
~
\hyperref[sec:031659]{\small{> > > NEXT SERMON > > >}}
\newline
\newline
羅馬書 15:14-33
\newline
\begin{longtable}{cl}
\hline
\hline
章節 & 經文 (和合本修訂版)\\
\hline
15:14 & \begin{tabularx}{0.7\textwidth}{X} 我的弟兄們，我本人也深信你們自己充滿良善，有各種豐富的知識，也能彼此勸戒。 \end{tabularx} \\ \\ \relax
15:15 & \begin{tabularx}{0.7\textwidth}{X} 但我更大膽寫信給你們，是要在一些事上提醒你們，我因神所賜我的恩， \end{tabularx} \\ \\ \relax
15:16 & \begin{tabularx}{0.7\textwidth}{X} 使我為外邦人作基督耶穌的僕役，作神福音的祭司，使所獻上的外邦人因著聖靈成為聖潔，可蒙悅納。 \end{tabularx} \\ \\ \relax
15:17 & \begin{tabularx}{0.7\textwidth}{X} 所以，有關神面前的事奉，我在基督耶穌裡是有可誇的。 \end{tabularx} \\ \\ \relax
15:18 & \begin{tabularx}{0.7\textwidth}{X} 除了基督藉我做的那些事，我甚麼都不敢提，只提他藉我的言語作為，用神蹟奇事的能力，並神的靈的能力，使外邦人順服； \end{tabularx} \\ \\ \relax
15:19 & \begin{tabularx}{0.7\textwidth}{X} 甚至我從耶路撒冷，直轉到以利哩古，到處傳了基督的福音。 \end{tabularx} \\ \\ \relax
15:20 & \begin{tabularx}{0.7\textwidth}{X} 這樣，我立了志向，不在基督的名已經傳揚過的地方傳福音，免得建造在別人的根基上； \end{tabularx} \\ \\ \relax
15:21 & \begin{tabularx}{0.7\textwidth}{X} 卻如經上所記：「未曾傳給他們的，他們必看見；未曾聽見過的事，他們要明白。」 \end{tabularx} \\ \\ \relax
15:22 & \begin{tabularx}{0.7\textwidth}{X} 因此我多次被攔阻，不能到你們那裡去。 \end{tabularx} \\ \\ \relax
15:23 & \begin{tabularx}{0.7\textwidth}{X} 但如今，在這一帶再沒有可傳的地方，而且這許多年來，我迫切想去你們那裡， \end{tabularx} \\ \\ \relax
15:24 & \begin{tabularx}{0.7\textwidth}{X} 盼望到西班牙去的時候經過，得見你們，先與你們彼此交往，心裡稍得滿足，然後蒙你們為我送行。 \end{tabularx} \\ \\ \relax
15:25 & \begin{tabularx}{0.7\textwidth}{X} 但如今我要到耶路撒冷去，供應聖徒的需要。 \end{tabularx} \\ \\ \relax
15:26 & \begin{tabularx}{0.7\textwidth}{X} 因為馬其頓和亞該亞人樂意湊出一些捐款給耶路撒冷聖徒中的窮人。 \end{tabularx} \\ \\ \relax
15:27 & \begin{tabularx}{0.7\textwidth}{X} 這固然是他們樂意的，其實也算是所欠的債；因為外邦人既然分享了他們靈性上的好處，就當把肉體上的需用供給他們。 \end{tabularx} \\ \\ \relax
15:28 & \begin{tabularx}{0.7\textwidth}{X} 等我辦完了這事，把這筆捐款交付給他們，我就要路過你們那裡，到西班牙去。 \end{tabularx} \\ \\ \relax
15:29 & \begin{tabularx}{0.7\textwidth}{X} 我也知道去你們那裡的時候，我將帶著基督豐盛的恩典去。 \end{tabularx} \\ \\ \relax
15:30 & \begin{tabularx}{0.7\textwidth}{X} 弟兄們，我藉著我們的主耶穌基督，又藉著聖靈的愛，勸你們與我一同竭力為我祈求神， \end{tabularx} \\ \\ \relax
15:31 & \begin{tabularx}{0.7\textwidth}{X} 使我脫離在猶太不順從的人，也讓我在耶路撒冷的事奉可蒙聖徒悅納， \end{tabularx} \\ \\ \relax
15:32 & \begin{tabularx}{0.7\textwidth}{X} 並使我照著神的旨意歡歡喜喜地到你們那裡，與你們同得安息。 \end{tabularx} \\ \\ \relax
15:33 & \begin{tabularx}{0.7\textwidth}{X} 願賜平安的神與你們眾人同在。阿們！ \end{tabularx} \\ \\
[1ex]
\hline
\hline
\end{longtable}
$^{1}$羅馬書第39章.
今天要講羅馬書第15章14-33節.
我們剛才讀了經文.
今天的主題是什麼?.
以行動活出恩順清義.
我們知道在羅馬書裡有恩順清義的教導.
但我們也很需要將恩順清義的教導活出來.
我時時說我們基督徒就好像一個人的頭很大.
雙腳很小,走不出來.
其實我時時這樣說是我自己的寫照.
我四歲才懂得走路.
頭很大,我的綽號叫做大頭君.
但雙腳小得不得了.
走路時頭先大,雙腳跟著.
所以經常鍾埋牆,鍾埋台角.
我的臉也要開刀,要縫針.
還衝過玻璃門.
以前的玻璃不像現在的安全措施的玻璃.
是很危險的.
所以我們也要小心.
千萬不要因為頭很大,雙腳很小.
保羅在這裡就教我們.
怎樣能夠活出恩順清義.
心水清的弟兄姊妹就會說.
牧師,你上一次才講完第35講.
為什麼現在這麼快說是第39講?.
其實我在幾年前已經講完羅馬書第十四章.
也講了第十五章頭.
所以已經講了這三篇的道.
就是第36講,第37講和第38講.
所以今天就是第39講.
下一次就是40講.
整個羅馬書只有40講.
所以下一次就是最後一講.
如果認識我的弟兄姊妹都知道.
我是非常喜愛舊約.
但是羅馬書就是第一部令我喜愛新約的書.
其實整部羅馬書都是一個首尾的結構.
這個首尾的結構一共有三層,很簡單.
第一層是什麼?.

$^{41}$一頭一尾都是講到福音.
在羅馬書的一章一至五節.
保羅說是神的福音.
在羅馬書的十六章二十五至二十七節.
保羅說是他所傳的福音.
一頭一尾都是講福音.
第二層是文案.
在羅馬書的一章六節和七節.
保羅向羅馬教會文案.
然後去到羅馬書最後.
他再次向他們文案.
在羅馬書的十六章一至二十四節.
這個就是我下次要和大家講的篇道.
有人說牧師文案你都可以來講一篇道.
是什麼來的?.
我講這麼多道就是要和大家一起來讀經.
所以我時時都說.
千萬不要以為有些經文有很多人名.
又有很多經文好像文案那樣.
你就不去讀,跳了它.
其實所有聖經裡面的經文都是與我們有益.
都是給我們的教導.
所以千萬不要跳了不讀.
記不記得我講大胃傳的時候.
曾經講過一篇道.
是講大胃勇士的名單.
連名單我都可以講一篇道.
所以你不要說文案更加可以講一篇道.
這個整個結構最中心的地方是什麼?.
是講到恩順清義.
羅馬書的一章八節至十五章頭.
就是講恩順清義的教導.
恩順清義的真理.
而今天的經文.
羅馬書的十五章十四節至三十三節.
就是講出保羅教我們怎樣由他自己的事工.
他的計劃能夠行出恩順清義.
所以我們不妨再用放大鏡.
看看今天首尾結構最中心的那兩點.
就是首先是羅馬書一章八節至十五章頭.

$^{81}$是講到恩順清義的真理.
恩順清義的教導.
記不記得保羅講什麼?.
我們得教是否因為我們的行為?.
不是.
因為我們接受耶穌基督.
耶穌基督的義就披在我們的身上.
以致我們能夠與神和好.
但是否真是我們自己的聖潔.
我們自己的義?.
不是.
所以我們自己應該個人立場.
應該服從神.
應該自潔.
應該克服罪.
追求公義.
然後羅馬書第十二章保羅說.
之後我們要怎樣?.
將自己獻上當作活祭.
用我們的恩賜來事奉神.
然後羅馬書第十三章講什麼?.
我們與政府的關係.
我們要以善勝惡.
這樣來對待我們的政府.
接著在羅馬書第十五章第十四至三十三節.
就是這個首尾結構的下半節.
保羅說.
我教你們怎樣用行動活出恩順清義.
保羅打開他的心.
讓羅馬教會的信徒見到.
他怎樣活出他自己的教導.
怎樣活出恩順清義的真理.
所以今天這篇道的大綱.
一共只有三點.
很簡單.
第一點.
羅馬書第十五章第十四至十五節.
保羅說.
我教你怎樣能夠行出僕人的謙卑.
在羅馬書第十五章第十六至二十一節.

$^{121}$保羅說.
我教你們怎樣活出使者的身份.
在羅馬書第十五章第二十二至三十三節.
保羅說.
讓我教你們怎樣順服神的旨意.
所以我們第一點來看看.
就是怎樣能夠行出僕人的謙卑.
羅馬書第十五章第十四至十五節.
保羅說.
在第十四節第興門.
我自己也深信你們是滿有良善.
充足了諸般的知識.
也能彼此勸戒.
保羅不是高高在上來教訓人.
保羅絕對可以這樣做.
因為保羅是一個非常淑齡的人.
但他仍然謙卑自己.
他能夠見到弟兄姊妹當中有什麼可取的地方.
這個也很值得我們來學習.
我們很多時候覺得自己不可一世.
但也很需要看看弟兄姊妹有什麼可取之處.
保羅就是這樣.
他見到羅馬教會有什麼可取之處.
第一他們滿有良善.
羅馬教會不像加拉太教會和哥林多教會那麼多問題.
我曾經聽過一個牧師說哥林多教會的系列.
這個系列叫什麼.
問題多多哥林多.
很順口.
哥林多教會真的很多問題.
加拉太教會也是.
羅馬教會就沒有那麼多問題.
但始終他們不是沒有瑕疵.
始終他們不是完美.
但保羅仍然能夠見到他們有什麼好處.
他們是良善.
除了良善.
還有第二點.
他們也是有充足的知識.
良善是什麼.

$^{161}$是有正確的道德觀念.
充足的知識是什麼.
說的知識是屬靈的知識.
這裡是說羅馬教會是有正確的教義.
這是第二點.
還有第三點.
保羅說他們是彼此勸戒.
他們懂得彼此勸戒.
勸戒這個字.
Nizeteo這個字.
Nizeteo原文的希臘文.
就是透過警告.
透過教導.
將人們本來在錯誤的道路上.
帶領他們走回正確的道路.
這個就是勸戒的意思.
很多人認為.
當我們有心理問題的時候.
應該怎樣.
去找心理學醫生.
但保羅怎麼說.
我們應該用神的話語彼此勸戒.
而不是去看心理醫生.
保羅說最好的方法是什麼.
首先你們是良善.
你有正確的道德觀念.
然後有充足的知識.
用神的話語彼此勸戒.
所以當我們有心理問題的時候.
不是去找心理醫生.
是用神的話語彼此勸戒.
一個人身體出現問題.
當然是去找醫生.
但如果心理出現問題.
是用神的話語來勸戒.
這是保羅給我們的教導.
所以我們說出了第一點.
保羅走出僕人的謙卑.
他看到羅馬教會有什麼可取之處.
第二點.

$^{201}$他教我們怎樣活出使者的身份.
羅馬書15章16至21節.
保羅在這裡說出有三個身份.
也是我們的身份.
第一個身份是什麼.
是祭司的身份.
在羅馬書15章16節.
保羅說使我為外方人.
作基督耶穌的服役.
作神福音的祭司.
叫所獻上的外方人.
因著聖靈成為聖潔.
可蒙原南.
這裡服役這個字.
就是Latourgos這個字.
原文的希臘文.
Latourgos在羅馬書13章6節.
被翻譯成差役.
這個是一般的公務人員.
但是用在屬靈的方面.
就是帶領敬拜的神職人員.
就好像舊約的祭司一樣.
但是我們要注意.
保羅不是說他自己是舊約的祭司.
他是新約的祭司.
他是作神福音的祭司.
什麼是新約的祭司.
在彼得前書的2章5節.
彼得教我們.
要作聖潔的祭司.
藉著耶穌基督.
奉獻神所閱納的靈祭.
一個新約的祭司.
是要獻上神閱納的靈祭.
等一下我們就會說.
這個靈祭是什麼.
在彼得前書的2章9節.
彼得也說.
唯有你們是被揀選的族類.
是有君尊的祭司.

$^{241}$是聖潔的國度.
是屬神的子民.
所以我們是新約的祭司.
我們不是舊約的祭司.
舊約的祭司是做什麼的.
在舊約.
人民不能夠自己去親近神.
他們親近神.
是透過祭司幫助他們.
現在因為我們新約.
透過耶穌基督所做.
祂為我們釘十字架.
祂稱我們為義人.
以致我們能夠與神和好.
所以我們新約的人.
不需要透過祭司.
就能夠直接去到神的面前.
所以若果今時今日仍然有人.
說某些人就是祭司.
好像天主教說神父就是priest.
這是什麼.
是不認識到聖經話語的真理.
是誤解了聖經的話語.
我們每個基督徒都能夠直接去到神那裡.
我們不需要有一個中間人.
如果我們說仍然需要有舊約祭司的職份.
就是違背了我們新約信徒的身份.
保羅只不過是將自己比喻成為一個祭司.
他不是說他自己是一個舊約祭司.
他希望羅馬教會的信徒.
能夠明白他的身份.
能夠明白他的工作.
我們先來看一看有兩點.
第一點就是舊約祭司的身份.
舊約的祭司是什麼.
是神面前代表人.
怎樣代表人.
獻上祭物.
當人民犯罪不潔的時候.
祭司就會獻上祭物來替那些人民贖罪.

$^{281}$所以舊約的祭司是執行贖罪祭.
他是神和人當中的宗寶.
我們現在就不需要.
耶穌基督是我們的宗寶.
已經是為我們釘了十字架.
我們不再需要獻動物的祭.
來潔淨自己.
保羅說他是作福音的祭司.
不再是獻祭贖罪的祭司.
他是作福音的祭司.
羅馬書的十五章.
在羅馬書的第一章.
他說他自己是奉召為使徒.
特派傳福音.
所以我們見到.
新約的祭司是什麼.
是傳福音.
等一下我們就會見到.
這個就是保羅第二個身份.
就是傳道者.
我們講完了舊約祭司的身份.
和新約祭司的身份.
不同的地方.
我們不妨看看.
這個獻祭的祭物.
舊約的祭物是怎樣.
無瑕癡是最好的獻給神.
要選那些動物沒有單點.
是完美的才能獻給神.
你說新約呢.
保羅叫自己做外邦人的使徒.
他將外邦人獻給神.
你說舊約的祭物是聖潔的.
新約的祭物怎能聖潔.
怎樣才能聖潔.
就是要這些接受耶穌基督的人.
認罪悔改.
當他們接受耶穌基督的時候.
耶穌基督的義就披在他們的身上.
所以當神見到這些外邦人的時候.

$^{321}$就見到耶穌基督的義.
見不到他們的罪.
是不是就是聖潔.
所以新約來說.
保羅是帶領這些外邦人信主.
讓他們認罪悔改.
成為聖潔.
有耶穌基督的義披在他們的身上.
然後將他們獻給神.
以致他們能夠蒙神越立.
所以這個就是新約和舊約.
祭品不同的地方.
保羅和羅馬教會說.
我是經過你們要去西班牙.
希望你們能夠支持我.
希望你們能夠給我有多些機會.
接觸到外邦人.
以致我能夠帶領他們信主.
就好像舊約的祭司獻祭.
將他們獻給神.
蒙神越立.
說完了保羅第一個身份.
活出使者的身份.
第一個身份就是祭司的身份.
第二個身份是傳道人的身份.
是傳道.
在羅馬書的十五章十七節.
保羅說.
所以論道神的事.
我在基督耶穌裡有何誇獎.
這裡是一個轉捩點.
保羅由祭司的身份.
轉到傳道的身份.
他不單止將人的靈魂獻給神.
他也有宣講神話語的恩賜.
宣講福音的恩賜.
保羅說我有值得可誇的地方.
你說.
自己誇自己不是那麼好.
保羅絕對知道.

$^{361}$在哥林多前書的一章三十一節.
和十章十七節.
保羅說.
誇口的當指著主誇口.
你不是誇你自己.
你指著神指著主來誇口.
在著哥林多後書的十二章五節.
保羅說.
除了我軟弱以外.
我並不誇口.
你要誇的時候.
你就誇你的軟弱.
為什麼.
因為當我們是軟弱的時候.
我們就是最剛強.
因為我們是依靠神的力量.
不是依靠我們自己的力量.
所以當我們軟弱的時候.
我們誇口.
就是誇我們是依靠神的能力.
很多時候我們會見到.
有些畫很漂亮.
如果我們見到有幅很漂亮的畫.
你會不會誇.
這幅畫的筆真的很棒.
你會不會這樣誇.
又或者你會不會誇.
這幅畫畫在這張畫紙上.
這張畫紙真的漂亮.
你會不會這樣誇.
當然不會.
你誇誰.
你誇畫家.
你誇畫家真的厲害.
畫出這麼漂亮的畫.
所以我們都是一樣.
我們只不過是神的筆.
只不過是神的墨.
我們只不過是神的器皿.
千萬不要自己誇大.

$^{401}$一切都是神的作為.
這個就是保羅告訴我們.
每一件事其實都有兩個極端.
一個極端.
人們就會自己誇大.
但是另外一個極端.
就是假謙卑.
很多時候我都聽過.
有人誇一間教會.
說神真的祝福你教會了.
但是那些弟兄姊妹或者牧師會怎麼說.
唉,一般般.
你這樣說的時候.
你以為謙卑.
其實你在做什麼.
你在貶低神的工作.
你在貶低神的事工.
所以你千萬不要假謙卑.
所以不要自誇.
但是也不要假謙卑.
保羅在羅馬書15章18節說.
除了基督籍我作的那些事.
我什麼都不敢提.
只提他藉著我言語行為.
用神跡騎士的能力.
並聖靈的能力.
使外邦人馴服.
怎樣能夠使外邦人馴服.
是不是保羅的工作.
不是.
保羅說是耶穌基督藉著我做的.
兩次說耶穌基督藉著他.
一次說是聖靈的能力.
很多時候我們就會怎樣.
在自己的面上鐵金.
這個人是我帶領信主的.
這個是我的果子.
但是保羅說什麼.
不是我,是神.
在《使徒行傳》14章27節.

$^{441}$經文說.
神藉著他們所行的一切事.
並神怎樣為外邦人開了信道的門.
有沒有看到.
首先是神開了福音的門.
然後怎樣.
是神藉著他們傳福音.
才能夠成功.
在《使徒行傳》15章12節.
神藉著他們在外邦人中所行的神跡騎士.
首先是神給他們能力有神跡騎士.
然後是藉著他們將福音傳開去.
經文說.
保羅是有神跡騎士的能力.
這個我們也說過.
因為保羅是使徒.
所以他有神跡騎士的能力.
現在信徒呢.
我們已經沒有了神跡騎士的恩賜.
要聽清楚.
我被很多人罵.
這個袁牧師就是不相信有神跡.
我絕對相信神跡.
但是再沒有神跡的恩賜.
因為我們已經有整部聖經在我們的手上.
不需要神再透過神跡騎士的恩賜.
來證明某些人所說的詐神的話語.
因為全部聖經都在我們的手上.
我曾經說過一篇道.
究竟神跡醫治和方言的恩賜還存在嗎?.
你想知道為什麼我們再沒有神跡恩賜.
你回去聽聽那篇道.
經文說.
是神令到外邦人順服.
是神的旨意所有接受祂的人都會怎樣?.
順服真道,聽從福音.
有很多人斷章取義.
聖經馬書第十章說.
保羅說只要口裡承認心裡相信.
就能得救.

$^{481}$他們說我是一個得救的人.
為什麼?.
因為我口裡承認心裡相信.
看戲要看完整齣戲.
不能斷章取義.
經文說.
所有口裡承認心裡相信的人都會怎樣?.
順服真道,聽從福音.
所以我們的行為就表現出我們是一個真正相信的人.
我講羅馬書第十章的時候說過這一點.
所以你想聽回的.
你可以聽回那篇道.
口認心信就能得救嗎?.
保羅在羅馬書第十五章第十九節.
甚至我從耶路撒冷直轉到以利利谷,到處傳了基督的福音.
以利利谷是保羅時攻最西北的角落.
大概距離耶路撒冷是一千里.
保羅在馬其頓和阿該亞十八個月.
相信他在那個時候曾經探訪過以利利谷.
在使徒行傳的十九章.
保羅不單止是祭司.
不單止有傳道的身份.
他也有先鋒的身份.
在羅馬書十九章二十節.
他說我納了志向不在基督的名被稱過的地方.
傳福音.
免得建造在別人的根基上.
這是保羅的名言.
保羅說我不會建造在別人的根基上.
什麼意思?.
他有感動將福音傳到美德之民.
就是福音未接觸到的地方.
保羅不要去到一些地方那些人已經聽過福音.
他說我要去美德之民那裡.
在羅馬書十五章二十一節.
就如經上所記未曾聞之他信息的.
將要看見未曾聽過的將要明白.
這節的經文是引用了以賽亞書的五十二章十五節.
說尼賽亞來臨的時候.
所有未聽過福音的人都會聽到福音.

$^{521}$所以保羅說我就要這樣做.
我要將神的福音傳到美德之民.
那些未聽過福音的人的地方.
每個神的僕人都有不同的使命.
有些被神召去做什麼?.
打江山,打根基.
但是也有些人被神呼召去在根基上建築.
所有的牧師和教師我們是做什麼?.
在教會工作.
教會的人已經接受耶穌基督.
就是在已經打好了的根基上建築.
宣教士好像Annie Fong是開荒牛.
去打江山,去打根基.
一個教會兩個職責的人都需要.
我們需要有些人來打根基.
但也需要有些人在根基上建造.
在羅馬書15章23至24節.
保羅說根據當時的地理.
可以傳福音的地方他已經傳了.
所以他希望去到西班牙.
西班牙就是當時認為最西面的地方.
好像去到地極一樣.
所以保羅說我想去那裡.
他和羅馬教會的弟兄姊妹說.
希望你們能夠幫助我.
給我一些資助.
然後我就會去西班牙.
羅馬教會不是保羅建立.
所以他不會在羅馬教會逗留很久.
因為他說過他不希望在別人的根基上建築.
我說有三點怎樣活出恩信稱義.
我已經說了兩點.
第一我們要謙卑.
走出獨人的謙卑.
第二點我說要活出使者的身份.
保羅說他有祭司的身份.
有傳道的身份.
有先鋒的身份.
還有第三點.
他是教我們怎樣能夠順服神的旨意.

$^{561}$羅馬書15章22至33節.
你說你怎樣看到保羅說要順服神的旨意.
你看看羅馬書15章32節.
並叫我順著神的旨意歡歡喜喜的到你們那裡.
與你們同得安息.
保羅是一個非常之尋求神旨意的人.
一早在羅馬書的第一章第十節.
他已經說在禱告之間常常懇求.
或者照著神的旨意.
終能不平坦的道路往你們那裡去.
在第一章保羅已經跟羅馬教會的人說.
我很希望憑著神的旨意可以去到你們那裡.
然後在羅馬書的第二章第十八節.
保羅說既從律法中受了教訓.
就曉得神的旨意也能分別是非.
我們怎樣能夠知道神的旨意.
保羅怎麼說.
在神的律法那裡.
明白神的律法.
明白神的教訓.
我們就能夠曉得神的旨意.
有誰比保羅更加明白神的律法.
保羅是一個法利賽人.
所以他絕對明白神的律法.
他非常之明白神的旨意.
看著羅馬書的十二章二節.
保羅還說.
不要效法者的世界.
只要心意根深而變化.
叫你們測驗何為神的善良.
順傳可喜悅的旨意.
保羅叫我們做什麼.
要測驗神的旨意.
我們要明白神的旨意.
所以保羅在哥林多前後書.
已不所書.
歌羅西書提摩太前書.
在這幾部書一開始的時候.
保羅怎麼說.
奉神旨意蒙召作耶穌基督使徒的保羅.

$^{601}$寫信給你們.
保羅是一個非常之明白神旨意的人.
他願意活在神的旨意當中.
他願意服從神的旨意.
這個也是我們每一個信徒應該要做.
尋求神的旨意.
服從神的旨意.
說到神的旨意.
保羅一共給了我們七件事.
我們要注意.
第一件.
我們要明白神在我們生命當中的召命.
在使徒行傳的十五章十六節.
使我為外邦人作基督耶穌的服役.
保羅絕對明白神在他身上的召命.
要他做什麼.
要他服侍外邦人.
他就是耶穌基督的服役.
所以我們每一個基督徒.
我們應該知道神在我們身上的召命是什麼.
如果我們不知道神在我們身上的召命是什麼.
我們應該要怎麼做.
跪在神的面前來禱告.
直到神告訴我們.
他在我們身上的恩賜是什麼.
他在我們身上的召命是什麼.
我們也需要明白.
隨著我們生命的改變.
神在我們身上的召命都會不同.
所以千萬不要一成不變.
我們很需要繼續尋求神在我們身上的召命.
繼續活出神的旨意.
除了第一點.
我們說我們要知道神在我們身上的召命是什麼.
第二點保羅說我們要服從神的旨意.
在羅馬書十五章二十二節.
保羅說什麼.
我因多次被攔阻,總不得到你們那裡去.
原文希臘文攔阻這個字.
就是Encocto這個字.

$^{641}$是什麼意思.
切除,切入的意思.
你說我知道了.
神攔阻我的時候.
將我切除了.
那我就什麼都不用做了.
不是這樣的意思.
其實這個字和古時的堡壘有很密切的關係.
在古時的堡壘四周圍堡壘都有一條坑.
挖了一條坑來做什麼.
防止敵人的攻擊.
攔阻敵人的攻擊.
甚至在坑裡面可能裝滿水.
當敵人來攻擊的時候.
就會攔阻他們.
這個就是這個字根的意思.
在這裡保羅用的文法.
希臘文的文法.
是不完美的文法.
什麼叫不完美的文法.
就是繼續持續的進行.
就是神不斷地在攔阻保羅.
你說.
神能夠怎樣攔阻一個人.
神可以用君天動地的神跡其事.
你說有沒有例子.
有.
記不記得神叫若拿去傳福音給尼尼美.
若拿就怎樣.
阿賓叫狗就越叫越走.
叫他去東邊.
他去哪裡.
他就希望去他斯.
就是西班牙.
當時就認為是最遠的地方.
他希望有這麼遠走這麼遠.
走到天腳底.
結果怎樣.
上了一隻船.
被水手拋了他下海.

$^{681}$被一條大魚吞了.
去到尼尼美把他套出來.
你說這個不是神跡其事.
就很難令人相信.
你什麼時候見過一條大魚吞了一個人.
去到目的地把你套出來.
那個人還不能死.
絕對是神跡其事.
所以第一個方法.
神可以用神跡其事來攔阻.
還有第二個方法.
是改變歷史的流程.
神可以令到保羅在外邦人的工作.
忙到不得了.
分身不暇.
不能夠有時候去到西班牙.
雖然保羅已經想去西班牙有很久.
神就不斷在這裡挖坑.
令到他不能夠去.
但是保羅絕對願意馴服.
這個就是第二點.
我也要說.
有時神會攔阻我們.
但是未必等於神不准許我們做這件事.
可能是時間還未到.
一切都操縱在神的手裡.
甚至時間也是.
說了兩點.
執行神的旨意.
第一我們要知道神在我們身上的召命是什麼.
第二我們需要服從.
有時神會攔阻.
第三我們需要籌備.
保羅知道絕對是神的旨意.
要他向外邦人傳福音.
所以他知道自己的召命.
第二他也願意服從神.
這個就是第二點.
但是他仍然需要有計劃.
有人說在神的旨意我就不需要計劃.

$^{721}$神幫我計劃一切.
不是的,在《箴言》第十六章第九節.
說人心籌算自己的道路.
為耶和華指引祂的腳步.
是神指引我們的腳步.
但我們仍然需要自己來策劃,籌備.
這個就是《羅馬書》第十五章二十三節.
保羅說好幾年.
我切心想望到西班牙.
去的時候可以到你們那裡去.
西班牙在希伯來文就是「他斯」.
就是我剛才說若拿先知想去那裡.
但是就不成功.
當時就是所有人認為是最西面的地方.
包括保羅他也想去那裡.
去得到美德之吻.
很有趣,經文說.
保羅切切想去到羅馬教會.
他怎會知道有羅馬教會?.
在《使徒行傳》第十八章.
告訴我們有一對夫婦,記不記得?.
伯基拉和亞居拉.
他們在哥林多前書第十六章.
在哥林多認識保羅.
告訴保羅有關羅馬教會的事.
所以保羅在這六年來.
一直禱告希望去到他們那裡.
來更加認識這班弟兄姊妹.
早在羅馬書的一章十三節.
保羅已經說.
弟兄們,我不願意你們不知道.
我屢次正義往你們那裡去.
要在你們中間得奢果子.
如同在其餘的外邦人中一樣.
只是到如今仍有阻隔.
有沒有看到?.
神不斷來到攔阻保羅.
所以很多時我們雖然是籌備.
但未必一定能夠成事.
要看神准不准去.

$^{761}$你說,有些經文說是撒旦攔阻.
究竟是神攔阻還是撒旦攔阻?.
你看《使徒行傳》第十六章六節.
聖靈既然禁止他們在西亞西港.
他們就經過佛雷加,迦拉太一大地方.
有時是聖靈禁止,經文告訴我們.
還有鐵撒羅尼加前書二章十八節.
所以我們有意到你們那裡.
我保羅有一兩次要去.
只是撒旦阻擋了我們.
究竟是神阻擋還是撒旦阻擋?.
是神還是鬼?.
兩樣都是.
記不記得我和大家說大衛傳的時候.
我說大衛數點人民不為神所喜悅.
是誰叫他數點人民?.
在撒母耳上的二十四章一節.
說是神激動大衛數點人民.
但在歷代志上的二十一章一節.
同樣的事情,說是撒旦起來攻擊以色列人.
你說是誰?.
我講這篇道的時候都講過.
很多時神是容許撒旦來試探我們.
以致完成神的旨意.
所以兩樣都是對的.
在羅馬教會的統治之下.
條條大路通羅馬.
我們經常都聽見這句話.
所以也有大路通往西班牙.
但是沒有任何的顯示告訴我們.
保羅曾經到達西班牙.
但是這些都不重要.
最重要的是什麼?.
他是有計劃去到西班牙.
他為福音傳到西班牙他有禱告.
這才是最重要.
據我們所知.
福音是在公元大概二百五十年才傳入西班牙.
但是早在這個之前的二百年.
保羅已經日夜為了西班牙來禱告.

$^{801}$這才是最重要.
說開又說,我新近才發現.
原來在中國的滁州.
發掘到有一塊東漢的銅鏡.
上面刻著基督教的文字.
如果這個真的屬實的話.
最早基督教傳入中國是什麼時候?.
東漢的時候大概是公元八十六年.
還早過福音傳入西班牙.
神真的很恩待中國人.
這麼早福音已經傳入中國.
弟兄姊妹你們有沒有感動?.
你們有沒有感動將福音傳到未得之民那裡?.
說了三點,第四點是什麼?.
行神旨意都有它的優先次序.
在羅馬書的十五章二十五節.
「但現在我往耶路撒冷去,恭給聖徒」.
保羅知道神在他身上的召命是什麼.
保羅願意馴服神.
甚至他有計劃去做神的工作.
但最後原來有優先次序要改變他的計劃.
剛才說保羅已經由耶路撒冷去到以利利古.
就是距離耶路撒冷一千里的地方.
他在那裡去西班牙很近.
但突然得到消息,原來耶路撒冷撓饑荒.
聖靈感動保羅要籌備捐款.
將捐款帶回耶路撒冷.
如果真是這樣,他會怎樣做?.
他要再走一千里路走回耶路撒冷.
跟他去西班牙的路剛剛相反.
但神的優先次序不容許保羅不服從.
保羅願意服從,有聖靈感動他.
所以雖然他一直希望去西班牙.
他仍然將捐款帶回耶路撒冷.
在使徒行傳的十一章和十二章.
也有提到耶路撒冷撓饑荒.
在加拉太書的第二章.
彼得和保羅一起來約定.
保羅要傳福音給外方人.
彼得要傳福音給猶太人.

$^{841}$所以保羅去到外方人的地方.
聽見耶路撒冷撓饑荒,他會怎樣做?.
在外方人那裡聚集款項,捐款.
然後將這個帶回耶路撒冷.
在加拉太書的二章十節.
保羅說:他們叫我紀念窮人.
這也是我本來熱心去行的.
保羅絕對願意.
在羅馬書的十五章二十六節.
馬其頓和阿該亞人樂意.
籌出捐項給耶路撒冷聖徒中的窮人.
保羅收集了捐款.
包括加拉太教會的捐款.
帶回耶路撒冷.
捐項這個字.
原文希臘文Kainonia.
是很有意思.
也解釋就是團契.
真的很有意思.
當我們大家有團契的生活.
發現弟兄姊妹有需要的時候.
我們就下力去捐給他們.
來彼此幫助.
所以外邦人非常願意捐錢給猶太人.
這就像他們獻了猶太人的債一樣.
其實第一間的教會在哪裡?.
是耶路撒冷.
是猶太人的教會.
過了七年之後才派出第一批的宣教士.
去到安提阿.
來帶領第一批的外邦人來信主.
然後在那裡建立教會.
福音首先也是傳給猶太人.
舊約就是神的好消息.
首先是給猶太人.
救贖也是透過耶穌基督.
耶穌基督是以色列人.
所以也是先給猶太人.
然後第一批的宣教士是猶太人的使徒.
難怪耶穌基督在約翰福音四章二十二節說.

$^{881}$「救恩是出自猶太人的」.
所以是猶太人將這些屬靈的教導.
將福音傳給外邦人.
現在他們自己在物質上有需要.
有沒有看到外邦人也是願意.
供應給猶太人.
首先是猶太人將屬靈的訊息給外邦人.
現在他們願意將物質上的供應給猶太人.
我們在這件事上看到.
猶太人和外邦人非常之合一.
彼此幫助.
這也是給我們的寫照.
要合一的心我們不能夠完成神的工作.
我們要同心合意才能夠興旺福音.
保羅絕對明白這一點.
保羅明白神旨意裡面有他的優先次序.
除了優先次序.
還有呢.
行神旨意我們就能夠將祝福帶給人.
這個就是第五點.
在羅馬書的十五章二十九節.
保羅說:我也曉得去的時候必帶著基督豐盛的恩典而去.
你問保羅怎麼知道他去行神旨意的時候.
一定會將神豐盛的恩典帶過去.
這個是保羅自己生命當中的經驗.
無論保羅是卑賤或者他是富裕.
無論保羅是自由之身或者他在坐牢.
無論他是在吃豐盛的筵席.
還是他被棍子打.
保羅知道耶穌基督一定會祝福他.
保羅最後去到羅馬.
他是以囚犯的身份去到那裡.
他就在羅馬寫《肥納比書》.
他在《肥納比書》的第一章說什麼.
他遇到被囚困.
他有鎖鏈鎖住他.
還有人毀謗他.
然後在《肥納比書》的一章十八節他怎麼說.
無論怎樣.
基督究竟被傳開了.

$^{921}$為此我就歡喜.
並且還要歡喜.
他講了兩次歡喜.
歡喜什麼.
耶穌基督的福音被傳開.
哪怕他是坐牢.
哪怕他被人打.
他一樣是這麼歡樂.
在《肥納比書》的四章七節他還說.
神所賜給人意外的平安.
必在基督耶穌裡保守你們的心懷意念.
保羅的喜樂不在乎是環境.
很多時候我們見到環境不好.
整個人就沒有心機.
保羅不是這樣的.
也不是在物質上.
在監獄他什麼都沒有.
但是他這些喜樂是哪裡的喜樂.
是靈裡面的喜樂.
他知道當他拿這些籌款回去耶路撒冷的時候.
他的敵人就在等著他.
所以他在《羅馬書》的十五章三十一節.
他求神叫我脫離猶太不順從的人.
因為他知道當他帶籌款回去耶路撒冷.
那些猶太人不會放過他.
而歷史真的告訴我們.
當他拿籌款回去耶路撒冷.
他就會被人誣告.
他帶領外邦人進入聖殿的內院.
將他收監甚至將他殺死.
還有第六點.
就是行神旨意的能力的力量.
保羅說我的力量從哪裡來.
《羅馬書》的十五章三十節.
我藉著我們主耶穌基督又藉著聖靈的愛.
有沒有看到.
他說他的能力是藉著耶穌基督的能力.
在《肥納比書》的四章十四節.
我靠著那加給我力量的凡事都能做.
有什麼可以做.

$^{961}$凡事都能做.
因為他是靠著耶穌基督的能力.
當他在此已不所書.
他教我們如何穿上全副的軍裝.
在已不所書的六章十節.
他怎麼說.
我還有沒了的話.
你們要靠著主.
依賴他的大能大力.
作剛強的人.
當我們依靠耶穌基督.
我們就能夠得到大能大力.
我們就凡事都能做.
還有呢.
他也是藉著聖靈的愛.
在詩篇的一百四十三篇第十節.
詩人說.
求你指教我遵行你的旨意.
因你是我的神.
你的靈本為善.
求你引導我到平坦之處.
保羅他的力量是從哪裡來.
三位一體的神的兩位其中.
第二位格耶穌基督的力量.
第三位格聖靈的愛.
都在幫助他行出神的旨意.
還有第七點最後一點.
在羅馬書的十五章三節.
不單止保羅是靠神的能力.
他還要不住地禱告.
在羅馬書的十五章三十節.
勸你們與我一同竭力.
為我祈求神.
竭力這個字就是竭力.
其實在這裡原文的希臘文.
保羅只要用竭力就已經夠了.
已經讓我們看到有很大的掙扎.
但他加多個竭是什麼意思.
強化了這一點.
有很大的掙扎是很竭力.

$^{1001}$甚至這個字可以用在真賤的上.
為什麼?因為禱告就是屬靈的真賤.
我們都很需要禱告.
神的主權和人的禱告.
永遠都是相成相負.
缺一不可.
當我們禱告的時候.
神就是透過我們來完成祂的事功.
所以是很重要.
耶穌基督祂自己都需要禱告.
祂被施洗約翰.
施洗之前祂是有禱告.
祂在曠野四十天祂有禱告.
在登山變相祂也有禱告.
甚至當祂餵飽五千人.
餵飽四千人.
祂一樣是禱告捉者.
祂使道拉薩路起死回生之前.
祂也向神來禱告.
耶穌基督自己都需要禱告.
所以我們都是一樣.
保羅很多時候都會請人來幫祂禱告.
其實經文就是這樣.
很多時候我們都會覺得.
請人幫我禱告.
很丟臉的.
失威了.
或者是示弱了.
請人幫我禱告.
但是連保羅都請人幫祂禱告.
保羅為什麼請人幫祂禱告呢?.
第一,為了平安.
叫我脫離在猶太不順從的人.
羅馬書十五章三十一節頭.
還有為了神事公的成功.
捐項可蒙聖徒的月納.
十五章三十一節尾.
還有為了喜樂.
為什麼保羅那麼喜樂.
羅馬書十五章三十二節尾.

$^{1041}$歡歡喜喜的道你們哪裡去.
所以我們都很需要禱告.
甚至請弟兄姊妹幫助我們來禱告.
在結束的時候我們來看看.
保羅不單止是教我們恩順清義的教導.
他甚至在他自己的生命.
在他自己的事公.
教我們怎樣能夠活出恩順清義的教導.
以致我們能夠帶領更多的人.
能夠順主.
首先是什麼?.
他有一個謙卑的心.
他是活出使者的身份.
有什麼身份?.
是祭司的身份.
帶領人去到神的面前.
還有呢?.
傳道的身份.
他是指著神誇口.
傳講神的話語.
還有呢?.
是先鋒的身份.
將福音傳給微德之民.
他願意服從神的旨意.
一共有七點.
我不會一一地重複.
很多時候我們都需要監察我們自己的生命.
我們是不是做得太少?.
我們是不是因為我們逐世的生活太忙?.
當我們傳達神的福音的時候.
又好像不夠力量.
又好像很膽怯那樣.
其實保羅告訴我們什麼?.
用不用怕?.
不用怕.
因為我們是倚靠耶穌基督的力量.
我們是倚靠聖靈的愛.
不是我們自己去做.
願意我們能夠學習保羅一樣.
活出我們自己的信心.

$^{1081}$將榮耀頌讚歸給神.
讓我們一起低頭禱告.
不單止保羅給我們有恩順清義的神學.
他也是以身作則.
讓我們看到怎樣能夠活出恩順清義.
願意我們有謙卑的心.
願意我們能夠活出我們是者的身份.
做一個祭司,做一個傳道,做一個先鋒.
而且是跟隨你的旨意,信服你的旨意.
願意在我們當中.
如果不明白自己的召命是什麼.
都能夠來到你的面前.
跪在你的面前來禱告.
求問自己的召命.
願意我們能夠活出你的旨意.
能夠靠著耶穌基督的力量,聖靈的愛.
能夠帶領更多的人信主.
將一切的榮耀頌讚歸給你.
我們禱告奉主耶穌基督的命而求.
阿們.
\newpage



\section{哈巴谷書 3:2}
\label{sec:031659}
\textbf{復興:憐憫與回轉的盼望}~~~~羅泰然牧師
\newline
\newline
連結: \href{http://www.church.com.hk/acms/content.asp?site=cdc\&op=show\&type=product\&code=031659\&layout=sermon}{\texttt{031659}} ~~~~ 語音日期: 2026-02-15
\newline
\newline
\hyperref[sec:031658]{\small{< < < PREV SERMON < < <}}
~
\hyperref[sec:index_chronic]{\small{[返順時目]}}
~
\hyperref[sec:031660]{\small{> > > NEXT SERMON > > >}}
\newline
\newline
哈巴谷書 3:2
\newline
\begin{longtable}{cl}
\hline
\hline
章節 & 經文 (和合本修訂版)\\
\hline
3:2 & \begin{tabularx}{0.7\textwidth}{X} 耶和華啊，我聽見你的名聲；耶和華啊，我懼怕你的作為。求你在這些年間復興你的作為，在這些年間將它顯明出來；在發怒的時候以憐憫為念。 \end{tabularx} \\ \\
[1ex]
\hline
\hline
\end{longtable}
$^{1}$張華峰議員.
張國柱議員.
主席.
張國柱議員.
\newpage



\section{創世記 11:31}
\label{sec:031660}
\textbf{聖經人聖經事:要往迦南}~~~~黃克勤牧師博士
\newline
\newline
連結: \href{http://www.church.com.hk/acms/content.asp?site=cdc\&op=show\&type=product\&code=031660\&layout=sermon}{\texttt{031660}} ~~~~ 語音日期: 2026-02-06
\newline
\newline
\hyperref[sec:031659]{\small{< < < PREV SERMON < < <}}
~
\hyperref[sec:index_chronic]{\small{[返順時目]}}
~
\hyperref[sec:031661]{\small{> > > NEXT SERMON > > >}}
\newline
\newline
創世記 11:31
\newline
\begin{longtable}{cl}
\hline
\hline
章節 & 經文 (和合本修訂版)\\
\hline
11:31 & \begin{tabularx}{0.7\textwidth}{X} 他拉帶著他兒子亞伯蘭和他孫子，哈蘭的兒子羅得，以及他的媳婦，亞伯蘭的妻子撒萊，一同出了迦勒底的吾珥，要往迦南地去；他們來到哈蘭，就住在那裡。 \end{tabularx} \\ \\
[1ex]
\hline
\hline
\end{longtable}
$^{1}$今天和大家一起討論另一個在聖經裡極為重要的人物.
就是亞伯拉罕.
說到亞伯拉罕就要說到他的前身.
就是他的父親.
他的祖父.
聖經人聖經時由生命說書的網絡為大家提供.
今天的題目是要往迦南.
要去迦南.
要去迦南?.
我們想起亞伯拉罕.
神叫亞伯拉罕離開本地本族父家.
去我指示你的地方.
當然他出去時也不知道自己去那裡.
後來發現原來去迦南.
我們一想起去迦南.
就想起亞伯拉罕.
但這條聖經告訴我們.
原來不是.
最早要去迦南地的並不是亞伯拉罕.
而是他拉.
留意一下.
創世記31章A.
主要是31節.
創世記31章.
31節很長.
這個版本不夠.
上半節他拉帶著他兒子阿伯拉.
和他的孫哈蘭.
的兒子羅德.
並兒父阿伯拉的妻子薩萊一起出發.
看下一章.
出了迦勒底的吾爾.
要往迦南地去.
他們一起走到哈蘭.
住在那裡.
最後他拉沒了205歲.
死在哈蘭.
大家留意下這條聖經說什麼.
再來.
他拉帶著他兒子阿伯拉.

$^{41}$和他的孫子哈蘭的兒子羅德.
即是他拉帶著他兒子阿伯拉.
和孫羅德.
和他的兒父阿伯拉的妻子薩萊.
一起出發去迦南.
離開迦勒底的吾爾城.
要往迦南地去.
去到哈蘭.
住在哈蘭.
他拉一直活到205歲.
死在哈蘭.
故事開始.
他拉這個人.
會不會都是信耶和華神.
極有可能.
我們常說亞伯拉罕帶著家人去迦南地.
其實最早帶的不是亞伯拉罕.
當時未改名叫亞伯拉罕.
亞伯蘭是跟隨的.
跟隨誰?跟隨他的父親他拉.
Tegra他拉.
Tegra他拉的意思是住宅地.
他拉帶著他的兒子阿伯拉.
Abram.
我們現在常常讀Abram.
將B放在後面.
變成Abram.
阿伯拉.
Abram帶著.
Abram.
Abram是被高舉之父.
這個父親.
好ye 呀 爸爸好ye .
好叻ge 爸爸.
他拉的兒子出世.
很想他兒子將來做一個好叻ge 爸爸.
幾好.
不是偉人.
偉雄.
又偉大又做英雄.

$^{81}$這些名字很大壓力.
考試差點都覺得自己的名字不配偉雄.
都幾難.
將來有很多事業叫偉業.
都幾大壓力.
原來只不過是打份工.
給自己壓力.
亞伯蘭的意思.
他拉想他兒子做一個好叻ge 父親.
被高舉的父親.
好呀.
無論你富裕還是貧窮.
都可以做一個好好的父親.
亞伯蘭這個名字幾好.
他孫子哈蘭的兒子羅德.
哈蘭是登山者.
如果你喜歡行山.
不妨改名哈蘭登山者.
他兒子哈拉有亞伯蘭和哈蘭.
哈蘭早死.
留下兒子羅德.
一齊跟隨他.
羅德就是遮蓋.
一齊帶著義父.
亞伯蘭的妻子薩萊.
薩萊是公主.
很多女孩從小到大都是公主.
應該改名薩萊.
他們離開當時的加勒底的梧椅.
加勒底是Castile.
破土者.
加勒底的梧椅.
梧椅我們比較熟悉.
Wuer是火焰.
火紅的聖詩.
梧椅城很厲害.
現在考古學者.
一直發掘梧椅城.
在兩河流域.
有法拉底河和底加利斯河.

$^{121}$差不多出口出海的地方附近.
出海是商業城市.
有兩條河一路夾到很近的地方.
水頭充足.
又是交通要道.
梧椅發展得很厲害.
火紅的城市.
街道發展到現在.
街道有溝渠.
泥路也有溝渠.
溝渠可能是石頭.
溝渠還有渠蓋.
讓牛羊車輪掉下去.
這個城市好發展.
不知會不會掘到水管.
原來梧椅有很多二樓.
還有.
很驚訝.
由於水頭充足.
他們在高的地方把水引過來.
梧椅城很有可能有自來水的供應.
很驚訝.
有熱水嗎?.
很厲害.
帶他們離開迦勒底.
這是後來的巴比倫的地方.
梧椅城要去迦南.
迦南哪裡來的?.
去到地中海海邊那塊地.
如果和梧椅比.
和迦勒底比.
落後很多.
去迦勒底做什麼?.
去發展嗎?.
去發展也不會去到那麼遠.
又不懂.
因此很有可能.
神給一個呼召.
叫他拉帶他們出去.
後來說.

$^{161}$亞伯蘭聽到神的聲音出去.
即是他也聽到.
也同時願意出去.
否則「爸爸你去吧,我們在這裡」.
他就跟著去.
羅德就跟著去.
其實亞伯蘭當日也願意跟著一起去.
就往迦南.
Kanaan.
Kla'an.
Kla'an其實迦南的意思是「低地」.
低地.
由一個很好的地方去到一個名叫「低地」.
就是海邊.
就是低地去迦南.
他們去到一個好好的地方.
千萬不要以為四個人一起去.
四隻樂陀.
我想有四百人左右.
整個大家族一起出去.
否則又要有壯丁守衛.
又要有人煮飯.
又要有人推車.
又要有人養樂陀.
又要有人起屋服侍.
沒有幾百人去不到.
只有幾家人一起去.
去到一個地方叫迦南.
即是羅德的父親.
不是.
迦南有兩個兒子.
有亞伯蘭和迦南.
迦南有個兒子叫羅德.
而迦南這個人早死.
所以要離開就帶孫子一起去.
他們去到一個地方.
又叫中文又譯為迦南.
Ka'an.
Ka'an和迦南是不同意思的.
Ka'an是哈蘭的H.

$^{201}$Ka'an是哈蘭的地方名.
剛剛去到一個交界處.
路一直到迦南一個城.
然後分開兩條路.
他們去到那裡.
走左還是右呢?.
不知道.
從吳爾到迦南走了一段路.
看來這裡不錯.
不如找個地方先住.
再想想怎麼走.
於是就住在那裡.
他就活了205歲.
死在迦南.
今天要講的故事就是這樣.
有什麼特別呢?.
有的.
第一.
最先要去迦南地的.
不是亞伯蘭.
是他的父親塔拉.
我相信是神帶領塔拉和他的兒孫一起出去.
塔拉是一個人物.
就是他帶亞伯蘭出去.
當然不是亞伯蘭.
什麼時候說是亞伯蘭的.
是.
不過我相信亞伯蘭也有心去.
不然就像我剛才講的,亞伯蘭說你去吧,我不去了.
這裡吳爾這麼好,我不走.
我的存款在這裡,物業在那裡.
打都不走.
但他也願意跟他父親一起去.
於是就這樣去.
第一.
我們要多謝.
欣賞一下塔拉.
竟然離開這個地方.
本來他就住宅在.
吳爾.

$^{241}$他的名字叫住宅地.
住宅在吳爾.
不過竟然願意出去.
第二是要往迦南.
就是今天的題目.
去一個不太好的地方.
但這裡是神叫.
他就去了.
這一點真是很好.
神叫他就去了.
這一點很值得欣賞.
神叫他就去了.
這個要往迦南.
不過很可惜走到哈蘭.
走到一個交叉路.
神叫他由火焰之城.
這麼火紅的地方出發去一個低地.
他也願意去.
塔拉也厲害.
亞伯蘭也厲害.
不過要去的時候.
去到哈蘭一個交叉路.
他就停了.
就住在那裡.
第三點要和大家分享的.
就是塔拉去到這裡停了步.
這個如果我這樣說.
可能是塔拉不夠好的地方.
沒有繼續走.
去到人生交叉點.
那怎麼辦呢?.
先停一停.
人生交叉點.
有沒有經歷過人生交叉點?.
去到那裡.
不知怎樣走.
不如你會怎樣做呢?.
不如停步.
不要再走.
再下去真的不知怎樣.

$^{281}$左右不是右又不是,左右有壓力又有壓力.
前無退路,後有追兵.
那怎麼辦呢?.
不如停下來先看清楚.
一路去.
由火焰之城.
火紅的城市去.
要去一個低地.
很有心智.
去.
神叫我去的.
誰知來到人生交叉點停了.
不再去了.
住在那裡.
結果塔拉住在那裡.
活了205歲.
死在哈蘭.
真慘.
這是塔拉結局的描述.
塔拉活了205歲.
死在哈蘭.
這也是她的墳墓所寫的.
這個人塔拉在墳墓上寫著205歲.
死在哈蘭.
不是本地人士.
可能是這樣.
這裡差了一點.
她沒有繼續去.
但我想告訴你.
不知道.
因為聖經的記載是這麼少.
譬如她會不會去到交叉點.
交叉點可能是人生交叉點.
不過剛剛用哈蘭這個名字來表達.
去到交叉點.
可能病.
可能老.
可能腰骨很痛.
都會繼續去.
騎駱駝.

$^{321}$怎麼騎呢.
可能體力不計.
可能身體欠佳.
可能有很多原因.
她就想住在那裡.
不如阿仔亞伯蘭你自己和羅德去吧.
我身體受不了.
或者有很多原因.
我想在那裡休息.
如果你是亞伯蘭你會怎樣.
那好吧.
你繼續休息.
看醫生吧.
我們先走了.
如果亞伯蘭會.
很有可能.
聖經沒有說.
我只是作這些故事.
很有可能亞伯蘭是這樣.
爸爸既然是這樣.
我們在這裡住陪你吧.
看看怎樣.
交叉路.
不如我們在這裡等一等.
可能文說前路有很多強盜.
如果去會有死傷.
我有幾十個壯丁怕什麼.
都不想一個人受傷.
不如歇一歇.
很多種原因去到交叉點.
你停了.
結果他拉死在那裡.
205歲.
我們很難評價.
哈拉是否不敢不肯不想再前行.
而住在那裡呢.
不過有些事我想說的是.
有些事是沒有辦法的.
去到某些交叉路.
你是會停一停.

$^{361}$不過我想說的.
我鼓勵的就是.
不要停在那裡太久.
如果停就會住落了.
不肯動.
那你要繼續走.
神帶我從火鴻城市去一個低地.
我的使命就是去那裡.
神叫我去.
我就去.
遇到交叉點.
休息一休息.
我繼續去.
這是很重要的.
記載第32節可能就不同了.
如果哈蘭死在半路.
205歲死在半路.
這個塔拉都厲害.
死都去.
真的不錯.
死在哈蘭是否很差.
那又不是.
因為我不知道.
將來去到天上.
應該都會見到塔拉.
問問他.
塔拉你當時在哈蘭不走住在那裡.
可能他會告訴你.
我騎駱駝騎得久.
腰骨塌了.
我不行得了.
哎呀.
你都厲害.
肯由火鴻城市出來.
其實都幾好的.
在這個故事裡我們學到一些東西.
都幾好.
看看今句.
他們出了迦勒底的吾爾火鴻城市.
要往迦南地區低地.

$^{401}$走到哈蘭交叉點.
就住在那裡.
大家要留意.
第一.
如果神帶領你出來.
我雖然是在迦勒底破土.
即是說.
我開始建設我的第一步.
破土然後入住.
不過.
神就說你離開這個火鴻城市.
跟我去一個低地.
好啊我去.
去到交叉點.
就停了.
我希望大家不會去到交叉點停了.
是困難的.
我鼓勵你繼續前進.
無論你的交叉點是什麼.
我們祈禱.
多謝我們的主.
想不到你這麼精密.
你這麼厲害.
用一些地方名.
幾個地方名.
表達了我們人生的一個.
發展與路向.
我們每一個人.
都有自己破土的地方.
就是我們原先建立的地方.
最開始的地方.
然後.
那個地方發展得多好.
我們的人生發展得多好.
變成火鴻到不一樣.
在這個時候主叫我們.
放下這個火鴻的事業.
叫我們去一個低地.
在那裡重新來過.
其實我們的主耶穌也是這樣.

$^{441}$從天而來.
在馬槽裡出生在一個低地.
來到醜惡的人世間.
然後來這裡生活.
就去一個低地.
是有些工作有些任務.
遇到很多人生交叉點.
真的不知如何是好.
好像哈蘭一樣.
都不知走哪一邊.
不過幫助我們.
我們看到他拉住低了.
直到他死.
我們不知道什麼因素.
但幫助我們遇到人生交叉點.
我們繼續向前走.
去你要我們指示我們去的地方.
要去我們的迦南.
給我們這個力量.
給我們勇氣.
給我們有能力去踏出.
可以去到主指示我們的地方.
禱告奉耶穌名求.
阿們.
\newpage



\section{哥林多前書 12:1-14:40}
\label{sec:031661}
\textbf{成全的愛}~~~~林羅麗斯牧師
\newline
\newline
連結: \href{http://www.church.com.hk/acms/content.asp?site=cdc\&op=show\&type=product\&code=031661\&layout=sermon}{\texttt{031661}} ~~~~ 語音日期: 2026-02-15
\newline
\newline
\hyperref[sec:031660]{\small{< < < PREV SERMON < < <}}
~
\hyperref[sec:index_chronic]{\small{[返順時目]}}
~
\hyperref[sec:031662]{\small{> > > NEXT SERMON > > >}}
\newline
\newline
哥林多前書 12:1-14:40
\newline
\begin{longtable}{cl}
\hline
\hline
章節 & 經文 (和合本修訂版)\\
\hline
12:1 & \begin{tabularx}{0.7\textwidth}{X} 弟兄們，關於屬靈的恩賜，我不願意你們不明白。 \end{tabularx} \\ \\ \relax
12:2 & \begin{tabularx}{0.7\textwidth}{X} 你們知道，你們作外邦人的時候，隨事被引誘，受了迷惑去拜不會出聲的偶像。 \end{tabularx} \\ \\ \relax
12:3 & \begin{tabularx}{0.7\textwidth}{X} 所以，我要你們知道，被神的靈感動的，沒有人會說「耶穌該受詛咒」；若不是被聖靈感動的，也沒有人能說「耶穌是主」。 \end{tabularx} \\ \\ \relax
12:4 & \begin{tabularx}{0.7\textwidth}{X} 恩賜有許多種，卻是同一位聖靈所賜。 \end{tabularx} \\ \\ \relax
12:5 & \begin{tabularx}{0.7\textwidth}{X} 事奉有許多種，卻是事奉同一位主。 \end{tabularx} \\ \\ \relax
12:6 & \begin{tabularx}{0.7\textwidth}{X} 工作有許多種，卻是同一位神在萬人中運行萬事。 \end{tabularx} \\ \\ \relax
12:7 & \begin{tabularx}{0.7\textwidth}{X} 聖靈彰顯在各人身上，是要使人得益處。 \end{tabularx} \\ \\ \relax
12:8 & \begin{tabularx}{0.7\textwidth}{X} 有人藉著聖靈領受智慧的言語；有人也靠著同一位聖靈領受知識的言語； \end{tabularx} \\ \\ \relax
12:9 & \begin{tabularx}{0.7\textwidth}{X} 又有人由同一位聖靈領受信心；還有人由同一位聖靈領受醫病的恩賜； \end{tabularx} \\ \\ \relax
12:10 & \begin{tabularx}{0.7\textwidth}{X} 又有人能行異能，又有人能作先知，又有人能辨別諸靈，又有人能說方言，又有人能翻方言。 \end{tabularx} \\ \\ \relax
12:11 & \begin{tabularx}{0.7\textwidth}{X} 這一切都是由惟一的、同一位聖靈所運行，隨著自己的旨意分給各人的。 \end{tabularx} \\ \\ \relax
12:12 & \begin{tabularx}{0.7\textwidth}{X} 就如身體是一個，卻有許多肢體，身體的肢體雖多，仍是一個身體；基督也是這樣。 \end{tabularx} \\ \\ \relax
12:13 & \begin{tabularx}{0.7\textwidth}{X} 我們無論是猶太人是希臘人，是為奴的是自主的，都從一位聖靈受洗成了一個身體，並且共享這位聖靈。 \end{tabularx} \\ \\ \relax
12:14 & \begin{tabularx}{0.7\textwidth}{X} 身體原不只是一個肢體，而是許多肢體。 \end{tabularx} \\ \\ \relax
12:15 & \begin{tabularx}{0.7\textwidth}{X} 假如腳說：「我不是手，所以不屬於身體」，它不能因此就不屬於身體。 \end{tabularx} \\ \\ \relax
12:16 & \begin{tabularx}{0.7\textwidth}{X} 假如耳朵說：「我不是眼睛，所以不屬於身體」，它也不能因此就不屬於身體。 \end{tabularx} \\ \\ \relax
12:17 & \begin{tabularx}{0.7\textwidth}{X} 假如全身是眼睛，聽覺在哪裡呢？假如全身是耳朵，嗅覺在哪裡呢？ \end{tabularx} \\ \\ \relax
12:18 & \begin{tabularx}{0.7\textwidth}{X} 但現在神隨自己的意思把肢體一一安置在身體上了。 \end{tabularx} \\ \\ \relax
12:19 & \begin{tabularx}{0.7\textwidth}{X} 假如全都是一個肢體，身體在哪裡呢？ \end{tabularx} \\ \\ \relax
12:20 & \begin{tabularx}{0.7\textwidth}{X} 但現在肢體雖多，身體還是一個。 \end{tabularx} \\ \\ \relax
12:21 & \begin{tabularx}{0.7\textwidth}{X} 眼睛不能對手說：「我用不著你。」頭也不能對腳說：「我用不著你。」 \end{tabularx} \\ \\ \relax
12:22 & \begin{tabularx}{0.7\textwidth}{X} 不但如此，身上的肢體，人以為軟弱的，更是不可缺少的； \end{tabularx} \\ \\ \relax
12:23 & \begin{tabularx}{0.7\textwidth}{X} 身上的肢體，我們認為不體面的，越發給它加上體面；我們不雅觀的，越發裝飾得雅觀。 \end{tabularx} \\ \\ \relax
12:24 & \begin{tabularx}{0.7\textwidth}{X} 我們雅觀的肢體自然用不著裝飾；但神配搭這身子，把加倍的體面給那有缺欠的肢體， \end{tabularx} \\ \\ \relax
12:25 & \begin{tabularx}{0.7\textwidth}{X} 免得身體不協調，總要肢體彼此照顧。 \end{tabularx} \\ \\ \relax
12:26 & \begin{tabularx}{0.7\textwidth}{X} 假如一個肢體受苦，所有的肢體就一同受苦；假如一個肢體得光榮，所有的肢體就一同快樂。 \end{tabularx} \\ \\ \relax
12:27 & \begin{tabularx}{0.7\textwidth}{X} 你們是基督的身體，並且各自都是肢體。 \end{tabularx} \\ \\ \relax
12:28 & \begin{tabularx}{0.7\textwidth}{X} 神在教會所設立的：第一是使徒；第二是先知；第三是教師；其次是行異能的；再次是醫病的恩賜，幫助人的，治理事的，說方言的。 \end{tabularx} \\ \\ \relax
12:29 & \begin{tabularx}{0.7\textwidth}{X} 難道個個都是使徒嗎？難道個個都是先知嗎？難道個個都是教師嗎？難道個個都是行異能的嗎？ \end{tabularx} \\ \\ \relax
12:30 & \begin{tabularx}{0.7\textwidth}{X} 難道個個都是有醫病的恩賜嗎？難道個個都是說方言的嗎？難道個個都是翻方言的嗎？ \end{tabularx} \\ \\ \relax
12:31 & \begin{tabularx}{0.7\textwidth}{X} 你們要追求那更大的恩賜。我現今把最妙的道指示你們。 \end{tabularx} \\ \\ \relax
13:1 & \begin{tabularx}{0.7\textwidth}{X} 我若能說人間的方言，甚至天使的語言，卻沒有愛，我就成為鳴的鑼、響的鈸一般。 \end{tabularx} \\ \\ \relax
13:2 & \begin{tabularx}{0.7\textwidth}{X} 我若有先知講道的能力，也明白各樣的奧祕，各樣的知識，而且有齊備的信心，使我能夠移山，卻沒有愛，我就算不了甚麼。 \end{tabularx} \\ \\ \relax
13:3 & \begin{tabularx}{0.7\textwidth}{X} 我若將所有的財產救濟窮人，又犧牲自己的身體讓人誇讚，卻沒有愛，仍然對我無益。 \end{tabularx} \\ \\ \relax
13:4 & \begin{tabularx}{0.7\textwidth}{X} 愛是恆久忍耐；又有恩慈；愛是不嫉妒；愛是不自誇，不張狂， \end{tabularx} \\ \\ \relax
13:5 & \begin{tabularx}{0.7\textwidth}{X} 不做害羞的事，不求自己的益處，不輕易發怒，不計算人的惡， \end{tabularx} \\ \\ \relax
13:6 & \begin{tabularx}{0.7\textwidth}{X} 不喜歡不義，只喜歡真理； \end{tabularx} \\ \\ \relax
13:7 & \begin{tabularx}{0.7\textwidth}{X} 凡事包容，凡事相信，凡事盼望，凡事忍耐。 \end{tabularx} \\ \\ \relax
13:8 & \begin{tabularx}{0.7\textwidth}{X} 愛是永不止息。先知講道之能終必歸於無有；說方言之能終必停止；知識也終必歸於無有。 \end{tabularx} \\ \\ \relax
13:9 & \begin{tabularx}{0.7\textwidth}{X} 我們現在所知道的有限，先知所講的也有限， \end{tabularx} \\ \\ \relax
13:10 & \begin{tabularx}{0.7\textwidth}{X} 等那完全的來到，這有限的必消逝。 \end{tabularx} \\ \\ \relax
13:11 & \begin{tabularx}{0.7\textwidth}{X} 我作孩子的時候，說話像孩子，心思像孩子，意念像孩子；既長大成人，就把孩子的事丟棄了。 \end{tabularx} \\ \\ \relax
13:12 & \begin{tabularx}{0.7\textwidth}{X} 我們現在是對著鏡子觀看，模糊不清；到那時，就要面對面了。我如今所認識的有限，到那時就全認識，如同主認識我一樣。 \end{tabularx} \\ \\ \relax
13:13 & \begin{tabularx}{0.7\textwidth}{X} 如今常存的有信，有望，有愛這三樣，其中最大的是愛。 \end{tabularx} \\ \\ \relax
14:1 & \begin{tabularx}{0.7\textwidth}{X} 你們要追求愛，也要切慕屬靈的恩賜，尤其是作先知講道。 \end{tabularx} \\ \\ \relax
14:2 & \begin{tabularx}{0.7\textwidth}{X} 那說方言的，不是對人說，而是對神說，因為沒有人聽得懂；他是藉著聖靈說各樣的奧祕。 \end{tabularx} \\ \\ \relax
14:3 & \begin{tabularx}{0.7\textwidth}{X} 但作先知講道的，是對人說，要造就、安慰、勸勉人。 \end{tabularx} \\ \\ \relax
14:4 & \begin{tabularx}{0.7\textwidth}{X} 說方言的，是造就自己；作先知講道的，是造就教會。 \end{tabularx} \\ \\ \relax
14:5 & \begin{tabularx}{0.7\textwidth}{X} 我希望你們都說方言，更希望你們作先知講道；因為說方言的，若不解釋出來，使教會得造就，那作先知講道的就比他強了。 \end{tabularx} \\ \\ \relax
14:6 & \begin{tabularx}{0.7\textwidth}{X} 弟兄們，我到你們那裡去，若只說方言，不用啟示，或知識，或預言，或教導，給你們講解，我對你們有甚麼益處呢？ \end{tabularx} \\ \\ \relax
14:7 & \begin{tabularx}{0.7\textwidth}{X} 就連那有聲而沒有生命的東西，如簫，如琴，發出來的音若沒有分別，怎能知道所吹所彈的是甚麼呢？ \end{tabularx} \\ \\ \relax
14:8 & \begin{tabularx}{0.7\textwidth}{X} 號角吹出來的音若不清楚，誰會預備打仗呢？ \end{tabularx} \\ \\ \relax
14:9 & \begin{tabularx}{0.7\textwidth}{X} 你們也是如此；若用舌頭說聽不懂的信息，怎能知道所說的是甚麼呢？你們就是向空氣說話了。 \end{tabularx} \\ \\ \relax
14:10 & \begin{tabularx}{0.7\textwidth}{X} 世上有許多種語言，卻沒有一樣是無意思的。 \end{tabularx} \\ \\ \relax
14:11 & \begin{tabularx}{0.7\textwidth}{X} 我若不明白那語言的意思，說話的人必以我為未開化的人，我也以他為未開化的人。 \end{tabularx} \\ \\ \relax
14:12 & \begin{tabularx}{0.7\textwidth}{X} 你們也是如此，既然你們切慕屬靈的恩賜，就當追求多得造就教會的恩賜。 \end{tabularx} \\ \\ \relax
14:13 & \begin{tabularx}{0.7\textwidth}{X} 所以，那說方言的，就當祈求有翻方言的恩賜。 \end{tabularx} \\ \\ \relax
14:14 & \begin{tabularx}{0.7\textwidth}{X} 我若用方言禱告，是我的靈在禱告；但我的理智沒有效果。 \end{tabularx} \\ \\ \relax
14:15 & \begin{tabularx}{0.7\textwidth}{X} 我應該怎麼做呢？我要用靈禱告，也要用理智禱告；我要用靈歌唱，也要用理智歌唱。 \end{tabularx} \\ \\ \relax
14:16 & \begin{tabularx}{0.7\textwidth}{X} 不然，你用靈祝謝，那在座不通方言的人，既然不明白你的話，怎能在你感謝的時候說「阿們」呢？ \end{tabularx} \\ \\ \relax
14:17 & \begin{tabularx}{0.7\textwidth}{X} 你的感謝固然是好，不過不能造就別人。 \end{tabularx} \\ \\ \relax
14:18 & \begin{tabularx}{0.7\textwidth}{X} 我感謝神，我說方言比你們眾人還多； \end{tabularx} \\ \\ \relax
14:19 & \begin{tabularx}{0.7\textwidth}{X} 但在教會中，我寧可用理智說五句教導人的話，強過說萬句方言。 \end{tabularx} \\ \\ \relax
14:20 & \begin{tabularx}{0.7\textwidth}{X} 弟兄們，在心志上不要作小孩子。但是，在惡事上要作嬰孩，而在心志上總要作大人。 \end{tabularx} \\ \\ \relax
14:21 & \begin{tabularx}{0.7\textwidth}{X} 律法上記著：「主說：我要用外邦人的舌頭和外邦人的嘴唇向這百姓說話；雖然如此，他們還是不聽從我。」 \end{tabularx} \\ \\ \relax
14:22 & \begin{tabularx}{0.7\textwidth}{X} 這樣看來，說方言不是為信的人作標記，而是為不信的人；作先知講道不是為不信的人作標記，而是為信的人。 \end{tabularx} \\ \\ \relax
14:23 & \begin{tabularx}{0.7\textwidth}{X} 所以，全教會聚在一處的時候，若都說方言，偶然有不通方言的或是不信的人進來，豈不會說你們瘋了嗎？ \end{tabularx} \\ \\ \relax
14:24 & \begin{tabularx}{0.7\textwidth}{X} 若個個都作先知講道，偶然有不信的或是不懂方言的人進來，就被眾人勸戒，被眾人審問， \end{tabularx} \\ \\ \relax
14:25 & \begin{tabularx}{0.7\textwidth}{X} 他心裡的隱情被顯露出來，就必將臉伏地，敬拜神，宣告說：「神真的是在你們中間了。」 \end{tabularx} \\ \\ \relax
14:26 & \begin{tabularx}{0.7\textwidth}{X} 弟兄們，那麼，你們該怎麼做呢？你們聚會的時候，各人或有詩歌，或有教導，或有啟示，或有方言，或有翻出來，凡事都應當造就人。 \end{tabularx} \\ \\ \relax
14:27 & \begin{tabularx}{0.7\textwidth}{X} 若有說方言的，只可有兩個人，至多三個人，且要輪流著說，也要有一個人翻出來。 \end{tabularx} \\ \\ \relax
14:28 & \begin{tabularx}{0.7\textwidth}{X} 若沒有人翻，就當在會中閉口，只對自己和神說就是了。 \end{tabularx} \\ \\ \relax
14:29 & \begin{tabularx}{0.7\textwidth}{X} 至於作先知講道的，只可有兩個人或是三個人，其餘的人當慎思明辨。 \end{tabularx} \\ \\ \relax
14:30 & \begin{tabularx}{0.7\textwidth}{X} 假如旁邊坐著的得了啟示，那先說話的就當閉口不言。 \end{tabularx} \\ \\ \relax
14:31 & \begin{tabularx}{0.7\textwidth}{X} 因為你們都可以一個一個地作先知講道，使眾人都可以學習，使眾人都得勸勉。 \end{tabularx} \\ \\ \relax
14:32 & \begin{tabularx}{0.7\textwidth}{X} 先知的靈是順服先知的， \end{tabularx} \\ \\ \relax
14:33 & \begin{tabularx}{0.7\textwidth}{X} 因為神不是叫人混亂，而是叫人和諧的神。在聖徒的眾教會中， \end{tabularx} \\ \\ \relax
14:34 & \begin{tabularx}{0.7\textwidth}{X} 婦女應該閉口不言；因為，不准她們說話，總要順服，正如律法所說的。 \end{tabularx} \\ \\ \relax
14:35 & \begin{tabularx}{0.7\textwidth}{X} 她們若要學甚麼，應該在家裡問自己的丈夫，因為婦女在會中說話是可恥的。 \end{tabularx} \\ \\ \relax
14:36 & \begin{tabularx}{0.7\textwidth}{X} 難道神的話是從你們出來的嗎？難道是單臨到你們的嗎？ \end{tabularx} \\ \\ \relax
14:37 & \begin{tabularx}{0.7\textwidth}{X} 若有人自以為是先知，或是屬靈的，就應該知道，我所寫給你們的是主的命令。 \end{tabularx} \\ \\ \relax
14:38 & \begin{tabularx}{0.7\textwidth}{X} 若有不理會的，你們也不必理會他。 \end{tabularx} \\ \\ \relax
14:39 & \begin{tabularx}{0.7\textwidth}{X} 所以，我的弟兄們，你們要切慕作先知講道的恩賜，不要禁止說方言。 \end{tabularx} \\ \\ \relax
14:40 & \begin{tabularx}{0.7\textwidth}{X} 凡事都要規規矩矩地按著次序行。 \end{tabularx} \\ \\
[1ex]
\hline
\hline
\end{longtable}
$^{1}$謝謝敬拜對大靈來一個很好的敬拜.
剛才邱牧師問下什麼日子.
大家還是在說情人節.
昨晚希望大家都過一個很開心的情人節晚上.
昨晚我們第一次是周六和清崇分開的.
原因是終於把拍拍拖拖的情愛.
留在青年區裡面去分享.
反而在大堂也是借著情人節的日子.
分開的主題叫做「承傳的愛」.
為什麼會想起這個題目呢.
原因就是情人節在最近的商業社會之下.
吹噓得很厲害.
昨天在尖沙咀也是一個很熱鬧的地方.
但是實際上情人節不是我們說愛情的.
你要知道2月14號是紀念一位神父.
叫做聖雅倫,但我們喜歡叫做華倫天奴.
因為習慣了這個翻譯.
其實是他殉道的日子.
為什麼他殉道呢.
有很多傳說,但兩個通常說得比較多.
第一個就是因為當時打仗.
皇帝不想男丁有家室.
太掛念,不夠空巷.
於是禁止當時的青年男女叫士兵去結婚.
而華倫天奴卻幫了一群想結婚的人一起結婚.
於是他就死了.
所以可以看到他不是自己結婚的.
他是為了承傳別人的愛情而犧牲的.
第二個就是殉道版,就是福音版.
我小時候就是看這個版本的電影.
他就說到華倫天奴在傳福音的時候.
在當時的第一世紀,基督教是受逼不吉的日子.
於是被抓去坐牢.
去到監獄裡面,那些肉粥對他們不好.
他又繼續向別人傳福音.
於是當中的一個首領就說.
如果這個上帝這麼厲害,就叫他治好我的女兒.
原來因為他有一個失明的女兒.
在那個時間裡面,當然是代禱之後就有一件奇事.
這個女孩子就真的康復了.

$^{41}$於是他整家,包括這個肉粥,都是決定信主的.
亦都因為這個緣故,他善待很多在肉粥當時因為宗教的原因被抓的信徒.
甚至有很多都藉故將他們釋放.
所以這一次其實華倫天奴都是被釋放的.
但當他繼續去傳福音的時候.
最終他就被捕,並且判了他在269年的2月14日那天斬頭.
所以那個日子並不是一個浪漫順化愛情的故事.
而是一個為了福音緣故來去殉道的一個神父的故事.
但為什麼他發展了有情人節卡那些呢?.
在電影的時候拍得很好的.
他就說這個神父知道自己這次真的會殉道.
他曾經醫好了這個女孩,那個女孩亦都藉故很多時候都會去探望他.
他就寫了一封信來向祂道別.
因為他的署名就寫著From Your Valentine.
所以你看到現在的商業就是這樣.
所以很多時候都會寫給你的愛人,你就會借這個名字.
而他很氣憤,我感動的地方是那個女孩收到這封信知道他這次真的一定要死.
爸爸問都沒有辦法,他很傷心.
於是他就寫了一封信來向這個能夠醫好他的眼睛.
又帶他信耶穌的一個神父來致敬.
他就寫了一封信叫做To My Valentine.
其實真的有這個Valentine這個人出現的.
而且那套戲做得很好.
就是他叫了他身邊的基督徒.
或者曾經都是接觸過這個神父的.
就全部寫卡進去這個監獄裡面去代禱和去支持這個神的僕人.
當然故事都是悽慘的.
終於斬頭就死了.
因此如果真的去世未的時候.
這個節日是有點悲傷的.
但是也都很看得到的.
反而不是愛情.
而是一個人是為了別人的緣故.
是真的願意犧牲自己.
甚至真是在我們基督徒來講.
都是一個很大的提醒.
今天我們在這個世代裡面.
其實是受一點苦都覺得很苦.
你有沒有想過有一天.
其實你會為主而去犧牲自己的性命.

$^{81}$對於我們這班城市人在和平長大的這一代.
的而且確是一個好像離得很遠的故事.
不過說起愛的時候.
無論在什麼年紀什麼的人.
其實都是很需要和很渴求愛.
沒有愛好像一個人是很難有快樂和滿足.
我們相信我們每一個嬰兒的出生.
包括我們自己的出生.
都大部分應該是由父母的愛之下來長大的.
如果我們的父母真是很用心的.
無條件的去愛我們.
我們相信在愛的家庭長大的小孩子.
他們的心靈是強大一點的.
也都能夠真的叫健康一點的.
但是經過不同的時代的改變.
我們發覺現在的小孩子的心真的又脆弱很多.
是不是因為有離婚呢.
最初我們有這樣想的.
但是這班人在這麼多年.
一直去同一間教會去看那方面的事情的時候.
我們發現原來我們這一批的小孩子.
和上一批的小孩子最大的分別.
都不是父母是否真的恩愛.
你問我們那一代.
我們的父母絕對不恩愛的.
我家是暴力家庭來的.
打架是打到劈血的.
所以我完全感覺不到他們恩愛的.
不過我們那個年代.
是大部分的兄弟姐妹沒有這麼少人.
原來一個人在一個家庭裡面.
父母是一個老師.
我們會看著他們怎樣相處.
知道什麼叫做真是彼此相愛.
又或者就算他們有吵架.
他們怎樣復和.
對於我們來說.
都是一個很無形當中的學習.
但是怎樣都不夠兄弟姐妹多.
我們教會最厲害的就是柏祥 保佑姨.

$^{121}$他們是十兄弟.
十個人已經是一個小組了.
所以在兄弟姐妹之間的恩怨情仇就精彩了.
原因就是如果兄弟姐妹多.
你不可以一輩子只有兩個朋友.
今天可能和二哥很友好.
但是下一次突然間我們就會吵架.
但是下個星期又不知為什麼.
你又可以和三哥又會變成朋友.
總之我們的關係.
愛與被愛 公平與不公平.
有沒有被欺凌這些所謂.
其實我們在家裡已經有很多鍛鍊.
所以你明白.
為什麼現在才會去到學校的時候.
經常都會被欺凌.
是因為在家裡從來沒有操練過什麼叫做欺凌.
所以他們的能力真的脆弱很多.
但是當一個人在家裡長大.
然後學了很多愛與被愛的技巧.
都會很快就想出去找朋友.
因為朋友就是我們自己選擇的.
因為和哥哥和姐姐爭執.
說到最後 媽媽兩個都罵.
你們兩個是兄弟.
你們是兄弟姐妹.
你們不能選擇.
這輩子你們都是親人.
無論將來怎麼樣你們都一定彼此扶持.
我相信我們都聽我媽媽的.
但是我們很想自己找朋友.
中學的時候.
因為那些朋友是我們自己選擇的.
正如以前的人很不喜歡萬分瓦解.
就很喜歡自由戀愛.
但是你發覺.
自由戀愛帶出來的眼淚.
絕對不比萬分瓦解的時候少.
為什麼在朋友之間當中都會有困難呢.
在基尼神明會有很多醫治 禱告服侍的事.

$^{161}$你就會發現很多人是不是只是拍拖.
會覺得很傷心呢.
其實有很多人一想起小時候.
很多時候是你初中的時候的同學.
你覺得是被出賣.
那種感覺和傷口.
反而是直根在一個成年人的心裡面.
是很重要的.
很影響他往後.
如何和朋友和同事相處.
不過有一種愛我們少提的.
就是你有沒有發覺.
原來人類這個動物.
我們普遍當中是有一個良心的.
無論你什麼的膚色.
有讀書沒有讀書都好.
原來創造主.
其實吹那口氣吹入我們當中.
讓我們成為一個有靈活人之外.
其實我們普遍是有個底的.
我們的惻隱之心.
我們的同情之心.
是承托著整個的社會和人類的.
我們見到人有事.
然後百方之援去幫手.
那個愛其實骨子裡面.
在神創造的時候.
已經和我們分享了他的性情.
不過有沒有留意.
其實現在這個世代.
特別是當有了傳媒.
加上現在新媒體的時代.
這一種丁點兒的愛.
被魔鬼一直一直在消磨和偷走.
因為我們已經開始.
不是好像我們爸爸那一代.
人人為我我為人人.
是不中聽的.
現在加上電片.
這麼流行的時候.

$^{201}$令我們完全失去信任.
今天是的.
這個世界是凶險的.
我們要警醒的.
但頂尖我突然在預備的時候.
我都會去想一下.
其實我們是不是又要小心的呢.
今天如果我們在外國有難.
你試試打個電話回來找人求救.
可能你打很多個都沒人聽.
特別我們這些.
每逢見到自己的通訊錄沒有的.
我大部分都不聽的了.
一定按了它.
每逢吃飯前後都不用想你了.
一定有些人打電話給你.
但有沒有想過.
如果我們真的當到人有需要去求救的時候.
原來這個世界四望無人.
人在你旁邊真的看著你.
但沒有人肯出意換手.
其實我們在製造一個什麼樣的世界.
讓我們去存活呢.
當然情人節.
個個都期望有份愛情.
因為這個愛情.
是令到我們的荷爾蒙分泌.
是紊亂的.
於是我們就沒有理性.
純粹追求浪漫的感覺.
但當這種荷爾蒙分泌沒有了之後.
大家就會很清醒了.
於是就發生了這麼多.
失戀,扔鍋,各方面的事情發生.
但弟兄姊妹.
我們的心其實很脆弱的.
包括我們那個無形的心.
在這些恩怨情仇裡面.
慢慢慢慢.
你就發覺你的心.

$^{241}$開始有裂痕.
甚至是碎.
所以越年紀大的.
不好意思.
就是越歷多的人.
你發覺我們普遍都會理性一點.
冷酷一點.
不是因為他們有經驗或者成長了.
而是其實反而我們是脆弱了.
我們更加脆弱.
所以成年人要去認識新朋友.
去一個新的群體.
是很艱難的.
很多的考慮.
很多的計較.
因此我都明白.
當我們去傳福音的時候.
走了幾年幸福小組.
我們開始發現.
已經有的存倉.
就已經用完了.
原來我們真的只有那幾個朋友.
然後你發覺.
我真的很想為主瘋狂.
嘗試去接觸一下新的人.
其實我們發覺.
人家都不是很容易接受你的.
於是我們就只可以去做街暴.
不是街的暴力.
是街頭暴動.
為什麼?.
因為起碼拍了一張福音單張給你.
我傳了福音信不信是你的事.
還要當你最好的照顧.
好像是很艱難.
但頂智妹.
我們是不是真的認命?.
我們的生命是不是就是因為.
曾經有過不好的經驗.
於是我們的心就一直碎.

$^{281}$或者是為了想它不爛.
我們就讓它變黑就可以了.
完全是對於身邊的靈魂的需要也好.
人家真的有肉體的需要也好.
我們都是在一個存疑的態度裡面.
愛為什麼這麼困難和為什麼這麼不開心?.
其實我覺得很在乎這兩方面的取態.
究竟什麼叫愛呢?.
愛對你來說是付出還是接受?.
可貨是人人都有的.
但是有沒有想過.
有時候如果你愛是你付出的時候.
你已經覺得很開心.
其實你真的不是很需要計較對方.
的反應.
為什麼我說父母的愛.
是最近似天父的愛呢?.
因為無端端就說大肚子.
然後痛得媽媽都不認得.
生了那塊東西出來.
但是不知道為什麼.
就是抱在身邊的時候.
你那種愛就會湧現.
你為了他不睡覺.
你為了他一叫你就要醒.
然後你沒辦法罵他.
因為他都聽不到你說什麼.
你只有你遷就他.
他永遠不會聽你的.
我們這樣學做媽媽和學做爸爸.
這個愛是你自己.
總之見到那塊東西你就開心.
他吃得久你就開心.
就算他大便你也會覺得開心.
是不是有一點點的呢?.
你不在這個愛裡面.
你不會明白的.
但反過來.
有很多人覺得愛是因為.
我感受到別人愛我.

$^{321}$我覺得他愛我.
現在大部分的下一代.
都是在這個位置.
所以每個人都是等你來愛我.
當大家都是在計較.
在等的時候.
你就知道.
愛放在天秤上面.
就已經是褪色了.
某程度來說.
更加是變了陣.
因為愛已經不是.
從心底裡面想對方好.
愛是你希望對方好的.
我去服侍他.
我去幫忙他.
就算做父母.
我們去教養弟兄姊妹.
我們這些弟弟妹妹.
或者是子女.
其實你的心應該是愛的.
你是希望他好.
所以你才限制他打遊戲機.
限制他外出.
背後的目的.
都是應該想成全他.
幫助他.
讓他能夠變好一點.
所以你明白.
當這樣去看的時候.
我們會在一個很繁忙的世界.
但是現在所有繁忙的工作.
繁忙的負擔.
我們其實沒什麼時間去細味.
我們自己想要的那份愛是什麼.
往往只是覺得.
恐懼去付出.
怕吃虧.
怕收不回回報.
所以不如冷血一點.

$^{361}$冰冷一點就算了.
但是在婚禮裡面.
我們又重重複複.
常常都看到這篇愛的詩章.
永遠都在婚姻裡面出現.
我們曾經在國內.
很多基督徒都會用.
不會說那些說話是從哪裡來的.
這個愛篇在國內是很流行的.
他們不知道是正經的.
很多大學生都不知道.
但是他們很喜歡這一段說話.
因為裡面說的愛是非常吸引.
有什麼樣的愛是可以恆久忍耐的呢.
如果能夠恆久忍耐就不會離婚.
就是我忍忍忍忍忍.
忍到有一天我受不了.
於是我就回頭走了.
愛是要有恩慈的.
恩慈的意思是上帝的性格.
是因為祂的憐憫.
是祂能夠看出你的軟弱.
明白你一切的缺失.
祂明白你.
所以祂能夠接納你.
也不計較地用愛來代理.
這個叫恩慈.
然後原來愛不是亂來都可以的.
愛是有真理的.
所以一切不應該的事.
都是要自己去控制不做的.
你想一下.
如果我們彼此之間的相處.
朋友之間的相處.
都能夠凡事包容.
凡事相信.
凡事盼望.
凡事忍耐.
那我們就不會出現這幾年的事件.
很多朋友都不知道怎麼聯絡.

$^{401}$因為我們就是凡事不能相信.
凡事不能忍耐.
凡事不能包容.
愛是否真的可以永不自息呢?.
很多人覺得這些詩篇很美麗.
很羨慕.
但是做不到的.
屬於天上的.
可能只是弄個相架.
掛在牆上.
有空的時候去安慰一下自己.
但弟兄姊妹.
今天我們為什麼會在由哥林多前書的十二,十三,十四.
三章一起去看這段經文呢?.
因為.
我很想讓大家去明白.
這篇詩篇不是天上有.
這篇詩篇是保羅寫信給提摩哥林多教會.
要求他們要學習這份愛.
在基督教裡面這份愛不是遙不可及.
這份愛是我們曾經品嚐過的.
如果你沒有經歷過上帝對你這種恆久忍耐.
又有恩慈.
而且是有真理.
並且是永遠長存的愛.
就不會等到我們情願.
願意去悔罪歸正.
接受這份的救恩.
耶穌基督用百分之一百個人的身份.
活在地上.
傳揚信息.
並且在十字架上面.
他說要赦免他們.
因為他們所作的.
他們不知道.
這就是這份令人感動的愛.
很多時候我們在愛裡面.
我們都會有計較.
不是說計較我付出多少.
你付出多少.

$^{441}$起碼我們會計較那個人好不好.
值不值得我愛呢?.
我們有勸那些人.
他是這樣.
你不要理他.
他不值得你這麼愛.
我們覺得這是對的.
不應該濫用.
值不值得去愛.
但聖經在羅馬書第五章七到八節說.
耶穌基督在我們還做罪人的時候.
為我們死.
如果要等人類值得耶穌基督去釘十字架.
可能末後都找不到這個時間.
我們沒有一個人.
覺得自己是值得這位神.
為我們所承受的一切苦.
耶穌基督的示範讓我們知道.
什麼叫做愛.
愛是付出的.
是成全的.
成全這群人類.
能夠脫離魔鬼的合制.
能夠離開罪的權勢.
能夠回去成為天父的兒女.
一切的好處.
都是和耶穌基督無關.
這份愛是單單為了成全我們.
幫助我們.
讓我們得到最好的救恩.
因此.
約翰是一個很明耶穌基督的門徒.
無論他是在《約翰一書》.
甚至是《約翰福音》.
他都常常提及這些說話.
當我們為何要切實彼此相愛.
最重要是因為.
我們在主的知識裡就知道.
什麼是這種愛.
這種愛是有標準的.

$^{481}$我們人與人之間.
稱兄道弟.
弟兄姊妹.
我們要有一個標準去愛.
彼此相愛不是隨便愛.
是喜好去愛.
你聽我說.
我就愛你.
組長.
小心.
我叫你做什麼你就做什麼.
我就愛你.
但我叫你做什麼.
你都推推搪搪.
那我就不愛你了.
弟兄姊妹.
你們要彼此相愛.
像我愛你們一樣.
這是我的命令.
耶穌基督給我們的是命令.
愛要愛到像什麼呢.
像耶穌基督愛我們一樣.
所以無論你遇到在.
教子女方面的問題.
無論你在朋友之間的問題.
你不可以輕易說.
算了.
放棄了.
因為你沒有問.
上帝是不是算了.
上帝是不是放棄了.
很多時候.
特別是在世伯小祖常常說的.
牧養關顧的時候.
我們有兩個很吊詭的.
要不就是.
叫都不聽.
就仆心仆命.
令我們去愛他.
然後就遷就他.

$^{521}$陪伴他諸如此類.
但是這個愛卻沒有真理.
因為往往很多時候.
不是你只是遷就他.
你就可以成全他.
你遷就他.
你就成全他的惡.
中國有句說話叫做成人之美.
還有下一句.
叫做不成人之惡.
那個故事就是.
當有一個人.
我忘了名字.
許配了一個太太給他.
但是他知道他的太太.
原來和另一個人是相愛的.
所以他就成人之美.
就是退婚.
讓他們兩個能夠結婚.
但是君子也不能成人之惡.
你明知道那個人是惡的時候.
你不可以繼續的.
遷就他奉承他.
然後你就任由他一輩子的惡下去.
這些就要在語場片裡學.
譚蘭卿那些.
我不知道你認不認識這個人.
她那些惡家婆就是這樣.
永遠都是.
縱容前頭婆的兒子.
因為他是後母.
對前面生下來的兒子.
完全縱容.
令他變成一個廢柴.
每個人都覺得這個後母很好.
對他很好.
但是他一生就是.
放在他後母的溺愛裡.
所以弟兄姊妹.
愛是要有真理的.

$^{561}$耶穌怎麼對我們.
祂不會宣佈有罪的就當無罪.
有罪就是要承認自己有罪.
然後我們要悔罪.
這樣我們就能夠.
可以白白的.
得到這一份的救恩.
因此怎樣去學習.
這一份愛在我們整個的肢體裡.
事實上都要我們.
真是好好的去思考.
並且真是好好的來學習.
做父母.
怎樣能夠成全我們的小朋友呢.
究竟一輩子扶著他.
還是有時候要看著他跌倒.
但是都未必馬上出手幫呢.
今日亦都是香港很多.
已經比較有錢的父母.
真是要好好去考慮的.
為什麼我會將十二十三十四章.
一起去看呢.
因為保羅寫信給哥林多教會.
哥林多是一間什麼的教會.
他是一間大的教會.
是一間有錢的教會.
是一間非常有能力與才幹的教會.
但是當很多的恩賜與能力.
在一間教會的時候.
不應該是蓬勃的嗎.
興旺的嗎.
大家很以教會為榮.
很喜歡教會的嗎.
不是.
因為裡面是鬧得很厲害.
於是保羅就用了這三章的聖經.
來講述弟兄姊妹之間.
教會之間彼此相愛.
在講什麼的會事.
我用三文字來形容這三章.

$^{601}$表面是麵包.
在講什麼是恩賜.
中間是最重要的.
因為中間吃什麼.
好不好吃.
就是靠它了.
這就是愛.
最妙的道.
保羅說.
這個愛篇是最妙的道.
是要我們細味的.
然後在底的時候.
亦都要去限制.
有規有矩.
怎樣來去運用恩賜.
所以愛與規矩是不會相沖的.
在愛裡面.
亦都一樣有規矩.
因為特別在群體裡面.
沒有規沒有矩.
是沒有辦法真正彼此相愛的.
因此弟兄姊妹.
當我們去想的時候.
今天希望你從一個新的角度去看這個愛.
這個愛不是在講夫妻的.
這個在講的.
是屬於上帝的群體的.
教會之所以重要的原因.
現在有很多的迷惑.
經常說我信耶穌.
但我不信教會.
我是基督徒.
但我不去教會.
留心你身邊的朋友.
留心我們的下一代.
永遠是洗腦的.
洗腦他們覺得是很對的.
但這是被魔鬼.
慢慢慢慢偷走的真理.
什麼叫教會.

$^{641}$聖羅穆斯在上一篇已經講了.
教會是耶穌基督的身體.
是上帝充滿的地方.
耶穌基督渡成肉身.
就要來到這個世界.
讓人知道祂是神.
神是怎樣愛我們.
我們是透過耶穌基督的.
整個生命事跡和教導.
我們理解得到.
我們看得到.
我們明白得到.
今天耶穌基督已經升天了.
我們見不到耶穌在.
未信主的人.
怎樣可以見到這裡有個神呢.
其實就是所有的信徒.
被呼召一起追責.
這個就是教會.
教會不是建築.
建築不是教會.
只要我們搬走.
這裡就已經不是教會了.
但當基督在這裡追責的時候.
我們就建立了.
耶穌基督的形象.
透過這一群人.
人們就知道這裡有個神.
而且這個神究竟做過什麼事.
這個神和我們有什麼關係.
都是透過每一個弟兄姊妹.
不同的故事.
來呈現這個神的美德.
讓其他未信主.
未認識神的人.
可以有機會來找到上帝.
因此弟兄姊妹.
我們不是反對教會有錢.
教會有力.
教會有知識.

$^{681}$我們只是反對.
我們將上帝的榮耀奪取.
保羅說我們要追求恩賜.
很多屬靈的恩賜.
在保羅的《哥林多前書》裡.
教得很清楚.
我們看很多上一課都要教過.
但是有個問題.
為什麼有些弟兄姊妹.
或者某些教會出現了很多的鬥爭.
或者是古靈精怪的問題.
就是我們都是人.
我們有一個很軟弱的位置.
我們喜歡成功.
我們喜歡力量.
我們喜歡我們自己那種滿足感.
於是往往當我們能夠.
經歷到恩賜的時候.
我們就開始有一點迷失.
迷失的意思就是.
突然間覺得自己.
我按手為人祈禱.
他痊癒了.
我按手為人祈禱.
他就倒下了.
嘩!好勁啊!.
自己都說自己好勁.
但是我很想說.
two party 有責任的.
會中有責任.
為什麼?.
因為我們人都是習慣了拜偶像的.
弟兄姊妹.
我們禱告服事.
是給我們的組員出來學習的.
你們就很害怕.
但是我們就希望能夠訓練到.
給大家機會去學習.
因為你是可以為人祈禱的.
信徒皆濟私.

$^{721}$每一個信耶穌的人.
你就是上帝的代表.
你都可以為人祈禱.
你都可以得到上帝的恩典.
去透過你給其他人的.
我們在最初的時候.
都真的被同行很認真地質詢.
為什麼你們會這樣?.
祈禱,祝福這些東西.
應該是牧師做的.
怎麼可能給平信徒.
平信徒.
為什麼會給平信徒去做呢?.
也有一些其他教會來到.
都不OK的.
什麼叫不OK呢?.
他會選人的.
如果有牧師,有傳道人在.
他會選走哪一行的.
你明白嗎?.
因為他始終覺得什麼?.
牧師強一點.
就是這麼簡單.
區長強一點.
執事強一點.
他將人的Title.
不是一個服侍的崗位.
而是以為他有Power一點.
以為他強一點.
為什麼《大經史會》當中.
也有人使用「因事而迷失」?.
是因為會中造成的.
今天他向我祈禱之後.
我真的康復了.
於是以後我們有人病.
我們就會想.
還是找他吧.
我們在《醫治報道大會》.
就有這些奇景.
林美德牧師的團隊.

$^{761}$個個都很有經驗.
很多時候很坦白.
上台講的.
都不及別人前線做的經驗多.
但你不要管.
我們這些靈性姐妹.
個個都要在這個位置.
因為林美德在這個位置.
全部坐著.
死都不走.
有沒有見過?.
有啦.
什麼時候有動作呢?.
突然間那邊很嘩然.
因為師母為了一個.
脊骨扭歪的人.
手又回復了.
那邊很感動.
很嘩然.
於是就發現.
咦?.
有些人影.
慢慢就那邊排.
靈性姐妹.
我常常提醒大家.
林美德牧師做了三十多年.
他仍然叫做林美德牧師.
他沒有叫做林美德醫生.
你就明白了.
如果他真是.
逢人走過去.
他可以醫到你的時候.
他早就改了招牌叫做什麼?.
林美德醫生.
醫治的是上帝.
來《醫治報道大會》.
很多都是我們靈性姐妹.
我們是基督徒.
我們都這樣想.
難怪外面的人這樣想.

$^{801}$靈性姐妹.
上帝給予這個能力.
與恩賜我們.
是讓教會能夠服侍.
能夠彰顯祂的榮美和能力.
你能夠來到一個報道大會.
有人為你祈禱.
是因為上帝醫治你.
於是他覺得你的上帝.
原來有醫治的能力.
讓他能夠認識實.
但是我們人類真是.
往往.
樣樣都差多.
教會又差人多.
又差錢多.
最重要是差的地方夠多.
還有要夠靚.
就是這樣.
我們好像去參加球賽.
看哪個球賽靚.
舒服些.
便宜些.
多福利些.
我們去參加.
靈性姐妹.
這個是錯的.
教會是耶穌基督的身體.
祂要彰顯了.
耶穌基督的榮美.
上帝要充滿那個地方.
一切蛇蟲鼠蟻.
如果上帝在.
其實是會被驅走的.
我們很渴求經歷神.
我們很渴求神.
在我們教會當中.
你有沒有想過.
上帝臨在同時間.
是一個審判的日子.

$^{841}$因為當上帝臨在的時候.
一切的罪惡.
一切的邪靈.
是要完全離開.
無處可躲進.
所以小心.
小心些什麼.
小心了我們成為了.
拜偶像的群體.
我們害死了.
那些有恩高的弟兄姊妹.
因為我們的崇拜.
加重了他們.
會變質.
因為人真的很難.
一百元買不通.
可能一百萬可以買通.
初初是很謙卑的.
被弟兄姊妹這樣捧的時候.
是我們加強他們迷失的.
但當然.
自己也要很好的去忍受.
當你在恩賜.
在上帝的使用當中.
要小心.
耶穌基督是合一的主.
祂的身體是正常的.
一個正常的身體.
手臂不會特別大.
或者不會像千層腳一樣.
只有一隻腳很大.
你想想.
如果你的眼睛大過你的嘴巴.
看到你也會很害怕.
一個身體健康.
是講平衡.
各安各職.
眼睛的比例.
腳和身體的比例.
每樣東西都是正確的.

$^{881}$就叫做漂亮.
弟兄姊妹.
西班牙小組.
很有質素的被虛榮蒙蔽.
原因是.
我們很努力去做小組.
於是我們很期望小組分組.
原來可以分組成區長.
好像升職一樣.
現在更厲害.
竟然發展到.
他們可以自己做區務.
還等同一個全能的位置.
於是你會發覺.
謙卑服侍的心.
就會被大家喊你的名字喊得多.
就開始迷失.
每個人都喊著你.
不同的頭銜都喊著你.
那種被人追捧的小心.
耶穌基督在被釘十字架之前.
進耶路撒冷.
都是無數的人在喊.
禍殺啊 禍殺啊.
多少人來擁戴啊.
接著下一秒.
釘他十字架 釘他十字架.
就是那些.
弟兄姊妹.
在細胞小組裡的服事.
為何我們要叫細胞小組.
因為所有細胞小組.
一直一直組織下去.
在合一裡面.
這個區有它的特色.
傳統人帶甚麼特質出來.
但可能它是心.
這個是肺.
那個是腎.
你是眼.

$^{921}$可能另外有些人是聰明的.
它是腳.
我們不需要一樣的.
在合一裡面.
我們是互為肢體.
意思就是說.
我們只是住在.
耶穌基督身體裡面.
其中一個部門.
甚至我們廣泛去看.
在深水Po 區.
一百三十多間教會.
有大的教會.
有小的教會.
有些是英文的教會.
有些是其他華語的教會.
但是你要知道.
我們是同一間教會.
無論教會有多美.
教會有多不美.
但是.
那裡都是上帝充滿的地方.
如果我們全部都合一.
在基層的深水Po 區.
他就會知道教會是甚麼.
教會是一個講耶穌的地方.
對人很好的地方.
如果有一間教會出事.
所有教會都蒙羞.
經歷過嗎.
所以為甚麼我們這麼害怕.
報紙說有傳道人出事.
其實與你無關.
你不認識他.
你媽媽也不認識他.
但你都覺得很怕被家人看到.
所以姐妹.
我們真的互為之體.
一個身體受損.
另一個身體都會痛.

$^{961}$另外在保羅教裡面.
為甚麼不說沒有恩賜就好.
不准用恩賜.
那就沒事了.
那就OK了.
有很多人是這樣的.
很怕靈恩就是一刀切.
甚麼都不要鬥.
靜靜坐在那裡就行了.
不時都沒有.
那就不會發生事了.
但你別這樣說.
這樣也做了魔鬼的詭計.
因為上帝的大能.
上帝的能力.
沒辦法透過他的教會.
他的身體來彰顯.
這樣一樣是做了魔鬼的詭計.
神為甚麼要將恩賜賜下呢.
第一.
讓我們每一個人都知道.
我們只是一些功能或功用.
究竟今天有甚麼功能.
那只是一個功能.
你按工作就有功能.
不按工作的功能.
對你來說是沒有意思的.
我們要知道反而是.
我們是共同去事奉一位主.
今天神要用我在這個功能裡面.
叫弟兄姊妹登憶處.
這是沒有其他原因的.
然而.
如何去平衡自己的心態.
在這個事奉裡面.
上帝又祝福你.
你又不容易跌倒呢.
為甚麼保羅要寫得這麼仔細.
在十二章.
他寫的比喻很好.

$^{1001}$甚麼是需要恩賜.
有恩賜的是不是很威風.
有恩賜的是不是值得去膜拜呢.
你應該膜拜似恩賜的那位神.
你不應該膜拜他.
為甚麼.
保羅說.
甚麼人才需要化妝.
你的眼睛很大.
小妹.
你會畫眼線嗎.
不會吧.
因為我不畫眼線.
我除了結婚之外沒有畫過眼線.
因為畫完之後會發現自己很大氣.
眼睛太大了.
這樣搞不定.
甚麼時候需要裝飾.
不盡美的才要.
你的手臂粗到媽媽都不認得.
你一定穿有袖的.
如果你有個小盲腰.
你當然很想努力衝界線.
露齒.
你會不會像豬腩肉那樣來露齒.
不會吧.
當然要遮住.
保羅很清楚.
身體裡面.
如果別人看起來不體面的地方.
我們就會遮住.
我的頸很漂亮.
就給別人看.
我的鎖骨很漂亮.
就給別人看.
所以每個結婚的減肥.
都是希望能夠看到鎖骨.
因為你想讓別人看你漂亮的地方.
因此,弟兄姊妹.
恩賜是上帝對我們很大的憐憫.

$^{1041}$不同的教會組成的時候.
總有缺陷.
在一個聚會裡面.
總是有些缺陷.
上帝是很樂意將恩賜給我們.
去幫助我們.
很坦白.
我經常用敬拜隊做比喻.
因為你想想.
現在敬拜隊很厲害.
很多弟兄姊妹是越來越有能力.
我們最初開教會的時候.
彈的不是很會彈.
唱的也不是很會唱.
結他也不懂得彈.
怎麼搞都搞不定.
偏偏又要搞樂隊.
對於我這個傳統神學的弟兄姊妹.
我真的搞不定他們.
很難的.
一個人.
雙手都要搞很久的協調.
還跟著五,六個人.
十隻手.
你又層,我又層.
我又說你不對.
你又說我不對.
當年連樂理都不懂.
連講什麼key都沒辦法.
夾到同一個channel.
但是你問我.
是不是敬拜很差呢.
你問我,我覺得不是.
很有趣的.
我們想像一下.
上帝說祂將恩賜賜給我們.
讓我們能夠叫人得益處.
能夠幫助他的身體.
彰顯上帝的榮美.
是一回什麼事.

$^{1081}$如果我們這群.
哇,全部都已經很OK.
專業的夾到.
那其實上帝沒什麼做的.
天時很空閒.
哇,欣賞一下你們.
多努力,練得多好.
哇,唱得不錯.
OK,搞定.
弟兄姊妹又覺得聽音樂會.
哇,音響很漂亮.
聲音很漂亮.
哇,很好.
今晚的show真是好了.
你知道什麼叫看show.
我們是這樣看的.
但是.
如果在一個教會群體裡面.
雖然他們樣樣都很差.
我們當年那群小朋友.
做敬拜隊都要知道.
都不會很厲害.
但是他們一大敬拜.
上帝的臨在很大.
為什麼.
哇,這個彈琴不行.
好像看不清.
上天使就一下子走下來.
要什麼.
幫忙.
將恩賜與能力加在他身上.
然後結他那個又不行.
鼓又不行,亂敲.
唱的那些.
聽起來好像有些問題.
不對音.
但是.
大家因為很害怕.
其實當天的那群小朋友.
真的,每個人都很害怕.

$^{1121}$被捧上台.
你以為那麼容易嗎.
很害怕會怎樣.
死抓住上帝.
所以我常常鼓勵他們.
哇,練歌的時候.
有一額汗.
但是你沒有來.
他們一起敬拜的時候.
很感動.
就是因為.
他們軟弱.
上帝臨在很多.
因為上帝臨在.
在整個會場裡面.
下面敬拜的弟兄姊妹.
是很容易與上帝相遇.
簡單來說.
天使滿屋飛.
因為真的好無.
去幫助他們.
弟兄姊妹.
當你這樣去想的時候.
你就知道.
我們無論在什麼場合裡面.
無論你聽說到.
你大敬拜.
或者在外面服侍也好.
你知不知道.
在天上面.
天父是與我們一起同供的.
因為我們是代表他的身體.
當我們越軟弱.
越無能的時候.
其實更加大下上帝的恩告.
既然是這樣的時候.
憑什麼覺得驕傲呢?.
為什麼我們要自誇呢?.
為什麼我們要去吹捧某一些.
有恩賜的弟兄姊妹呢?.

$^{1161}$原來你沒有想過.
你可能害死他們.
當然我不是說.
你不可以讚賞他們.
但是記住.
我們不要為自己.
雕刻偶像.
因為這是犯十戒的.
當我們能夠從這麼謙卑的心.
去尊崇我們的上帝.
炫耀張整恩的教會.
或者我們這群朋友.
我們是基督徒.
我們都是同心的.
彼此相愛.
我們活出上帝的樣式.
其實不用怎麼傳福音.
因為你們就是福音了.
見到你們就是福音了.
你知不知道.
中國在三十年很大的壓制裡面.
當一開放的時候.
基督徒是翻了很多倍的.
沒有人想過.
以為基督教是會滅絕的.
但是你想不到.
基督徒的人數是多了很多倍.
在大爆發的時候.
是不能傳福音的.
其實現在都沒有了.
你不可以向十八歲以下的人傳福音.
你不可以在街上跟別人傳福音.
所以我們香港.
真的能夠做的.
我們就盡量在街上做.
我們就做外面的攤位.
就是這個意思.
我們有這個自由.
沒有的.
為什麼會無端端信徒.

$^{1201}$在這幾十年這麼痛苦的裡面.
竟然可以倍增.
就是因為.
可能就是一條村.
一間屋.
有一群人.
每個禮拜.
按時去聚會.
知道他們有些事唱著唱著.
又不知道他們在做什麼.
但是你不要管.
你總是知道這群人.
彼此相愛.
個個都願意為對方付出.
從不計較.
還有.
哇!村裡面有這麼壞的人.
為什麼會.
跟他們在一起之下.
又變好了.
於是人一定有好奇.
人是會自己走來.
找耶穌的.
弟兄姊妹.
我們有沒有這個胸襟.
你都是被上帝去使用.
並且我們都要立志.
原來我們的愛.
是要用來成全別人的.
不是成全自己的.
組長.
你帶走倍增.
你人歸主.
是好的.
但是不是用來成全自己的.
你要成全你的下一代.
要去訓練你的核心.
期望他一天代替你的位置.
世班小組這一個.
完全倒轉的領導概念.

$^{1241}$是很合乎聖經的.
大的就是要去服侍世.
什麼叫服侍世.
成全他.
你能夠有另一個組長.
能夠起來成為區長的時候.
你可以成全他.
讓他做.
你退下來.
支持他.
所以很多弟兄姊妹經常都很害怕.
師父你真的退休了.
哇!我沒有工作了.
有什麼好害怕的.
教會是神的.
還有.
我們見到下一代.
幹嘛在那裡堅持.
離開這個位置.
成全他們上屋.
一天站在這個位置.
沒有人上這個位置的.
區長.
區務.
組長們.
我期望你.
永遠成全的是什麼.
有人代替你的位置.
這個就是我們最開心的時間.
最光榮的時間.
弟兄姊妹.
你們做父母的.
難道不想你的子女比你厲害.
是不是.
做老闆又不喜歡.
做老闆就是.
用你的力量.
增加自己的收入.
用完可以棄.
但這個不是上帝說的愛.

$^{1281}$上帝愛我們.
是要成全我們.
讓我們更使他.
上帝不是將我們用完就棄.
弟兄姊妹.
讓我們明白耶穌基督的心.
因為你很喜歡的愛篇.
麻煩你讀完前頭的三節.
聖經不可以一截一截讀.
多倫多前書十三章的第一至三節.
講得很清楚.
就算你會講萬人的方言.
你有天使的話語.
你又有先知講道之能.
你就知道上帝的奧秘.
你是上帝的代言人.
So what.
沒有愛.
一切算不得什麼.
你只是那些bang bang 聲.
吵耳的鑼.
吵耳的筆.
在上帝的眼中.
一點都沒有價值.
所以弟兄姊妹.
什麼叫在主裡面合一.
在主裡面合一不是齊一.
而是我們要彼此相愛.
彼此成全.
我的貢獻出來.
是令到小組更加能夠.
去祝福其他的人.
我貢獻我的才幹出來.
在教會裡面去服侍.
令到更多人來這裡.
他們能夠舒舒服服的.
去認識這位主.
弟兄姊妹.
當大家都有這樣的心的時候.
你就知道.

$^{1321}$保羅沒有反對教會的弟兄姊妹.
追求屬靈的恩賜.
但是他的指向教導很清晰.
有真理.
他說你要追求的是.
造就教會的恩賜.
你要追求的.
做的事.
什麼都可以有自由做的.
但是.
原則是.
能不能夠造就人.
所以不要經常問規矩.
這間教會其實本身.
我開始那個年代沒有規矩的.
我們是沒有規沒矩的.
你想一下.
每個建築牧師.
我們這些就算開國元老那些.
有幾個高雅星早辰那些.
都沒有的.
牧師去帶團的.
他又幫大家拿行李.
拿到台灣的姊妹都接受不了.
為什麼牧師做這些事.
牧師說我還有力氣做的時候.
讓我做吧.
我沒力氣的時候就讓你做吧.
丁姊妹.
你要記住.
你有自由的.
但是背後請你想一想.
是不是造就那個人.
我們在最初做小組教會的時候.
頭一批的組長是很清楚的.
我們那時候每個月開一個組長會.
經常都說這句話.
就是將對方的好帶出來.
什麼叫做小組.
什麼叫做爹地媽咪.

$^{1361}$怎樣去做一個好的老闆.
好的團隊.
最簡單的.
不是將他的不好帶出來.
剛才保羅說的.
不好那些怎麼樣.
幫他蓋住他.
修飾一下.
讓他看出來漂亮一點.
但是他好的.
反而要讓他出.
推他上台.
你明白嗎.
這個胸襟.
不是很多領袖可以做得到.
做父母的.
如果你沒有這樣的胸襟.
你沒辦法看到你的子女比你厲害.
我和我老公經常都說.
女兒有時候做服侍都很辛苦.
在外面.
幸好是女兒給點面子.
我都和女兒說.
我經常都說這句話.
現在你就是林以樂的女兒.
但是我們很希望有一天.
我們是林盈熙的爹地媽咪.
明不明白.
今天你叫你做組長.
什麼組長什麼區務.
真的很開心.
但是你要知道有一天.
反過來你叫他.
什麼區務.
那就差很遠了.
我都過了兩年了.
沒有我的勵思牧區.
我非常高興.
我看到我們的牧區.
仍然有很好的領袖.

$^{1401}$弟兄姊妹.
希望在新的一年裡.
我們求主給我們有三個方向.
我們要學習成全的愛.
學習什麼呢.
第一件事求主給我們.
有他同一的體潔.
看到我們弟兄姊妹.
看到我們身邊的朋友.
他們的美.
香港人這幾年搞砸了.
我們大家都是一看就看壞的.
弟兄姊妹.
看看你旁邊那個.
他有好東西的.
知道嗎.
他有好東西的.
神的所有的創造.
最後那句都說什麼.
神看著是好的.
我們一定要記住.
這是上帝的心意.
多壞的祖源.
八次返一次.
六次返一次.
他有那麼多時間記著你.
他又不返小祖.
但是你怎樣對他.
看他是好的.
第二.
愛是要有真理的.
我們真的要很努力學的.
我都覺得很辛苦.
學了一輩子了.
怎樣用神的口.
意思是.
神叫我們要造就祂的說話.
我們媽媽一開口就罵你.
不罵你.
就算是稱讚你.

$^{1441}$所以我們都不是很習慣稱讚別人.
覺得稱讚別人有點毛骨悚然.
覺得稱讚別人有點虛偽.
但是現在這裡不同了.
這裡很玻璃的.
你小心你的說話.
你的說話是帶有能力的.
經常都拆毀了它.
其實你叫他乖一點.
或者叫他收拾一點東西.
他覺得你罵他了.
弟兄姊妹.
唉.
沒辦法.
愛嘛.
愛就要成全了.
做大的就要服侍小的.
努力去學.
可以這樣說.
想一想.
這樣說.
能否祝福他.
祝福不了他.
再想一想.
你叫他收機.
不准打電腦.
你說出來.
罵一罵.
都沒用的.
收回他.
想一想.
怎樣說可以令他收機呢.
祈禱吧.
求神讓我們有智慧.
但是還有一樣東西.
我們要求神讓我們有行動.
要有行動.
上帝其實是每一天.
都是在訓練我們的屬靈生命.
我們所遇的.

$^{1481}$無論是好還是不好.
背後都有上帝的好意.
讓我們真的能夠被鍛鍊.
能夠更加堅壯.
讓我們在這個末世的世代裡.
我們能夠擇樂得穩.
我們能夠有一個愛的群體.
互相扶持.
不要放棄.
我經常說小組.
不要放棄你的小組.
一個小組.
兩個小組.
就算你覺得不滿意.
不要緊.
最重要我們不要放棄.
我們努力地去學習.
慢慢慢慢的時候.
這個身體就越粗越靚.
越粗越fit.
不粗的身體就只有肥.
教會一樣.
你的小組也是一樣.
願上帝親自賜福給我們.
讓我們的使命人生.
明白上帝對教會的心意.
我們切實彼此相愛.
而且將對方的好帶出來.
Amen.
\newpage



\section{哥林多前書 13:1-8}
\label{sec:031662}
\textbf{給屬靈暴發戶的妙道}~~~~冼嘉琪傳道
\newline
\newline
連結: \href{http://www.church.com.hk/acms/content.asp?site=cdc\&op=show\&type=product\&code=031662\&layout=sermon}{\texttt{031662}} ~~~~ 語音日期: 2026-02-15
\newline
\newline
\hyperref[sec:031661]{\small{< < < PREV SERMON < < <}}
~
\hyperref[sec:index_chronic]{\small{[返順時目]}}
~
\hyperref[sec:031663]{\small{> > > NEXT SERMON > > >}}
\newline
\newline
哥林多前書 13:1-8
\newline
\begin{longtable}{cl}
\hline
\hline
章節 & 經文 (和合本修訂版)\\
\hline
13:1 & \begin{tabularx}{0.7\textwidth}{X} 我若能說人間的方言，甚至天使的語言，卻沒有愛，我就成為鳴的鑼、響的鈸一般。 \end{tabularx} \\ \\ \relax
13:2 & \begin{tabularx}{0.7\textwidth}{X} 我若有先知講道的能力，也明白各樣的奧祕，各樣的知識，而且有齊備的信心，使我能夠移山，卻沒有愛，我就算不了甚麼。 \end{tabularx} \\ \\ \relax
13:3 & \begin{tabularx}{0.7\textwidth}{X} 我若將所有的財產救濟窮人，又犧牲自己的身體讓人誇讚，卻沒有愛，仍然對我無益。 \end{tabularx} \\ \\ \relax
13:4 & \begin{tabularx}{0.7\textwidth}{X} 愛是恆久忍耐；又有恩慈；愛是不嫉妒；愛是不自誇，不張狂， \end{tabularx} \\ \\ \relax
13:5 & \begin{tabularx}{0.7\textwidth}{X} 不做害羞的事，不求自己的益處，不輕易發怒，不計算人的惡， \end{tabularx} \\ \\ \relax
13:6 & \begin{tabularx}{0.7\textwidth}{X} 不喜歡不義，只喜歡真理； \end{tabularx} \\ \\ \relax
13:7 & \begin{tabularx}{0.7\textwidth}{X} 凡事包容，凡事相信，凡事盼望，凡事忍耐。 \end{tabularx} \\ \\ \relax
13:8 & \begin{tabularx}{0.7\textwidth}{X} 愛是永不止息。先知講道之能終必歸於無有；說方言之能終必停止；知識也終必歸於無有。 \end{tabularx} \\ \\
[1ex]
\hline
\hline
\end{longtable}
$^{1}$開始的時候,我想問有沒有看過這個Blockgame?.
沒有?.
這個叫做Would you rather的Blockgame.
這個遊戲,首先有兩個選擇.
你會寧願選擇A還是B?.
我現在就問大家一個問題,Would you rather?.
假設現在大家是12歲.
如果有人告訴你,你會成為一個附庸.
百萬附庸,或者千萬附庸.
你現在12歲.
你會想現在立即做這個附庸.
還是十年後才做這個附庸?.
他給你兩個選擇.
現在立即做,還是十年後才做這個附庸?.
你現在是12歲.
十年後,有人選擇十年後.
當然有人選擇現在.
哈哈哈哈.
為什麼會選擇十年後呢?.
為什麼會選擇現在呢?.
為什麼呢?.
沒有對錯的.
當然我們有不同的選擇.
如果選擇現在那些.
就會即時變了一個身份.
那些叫做暴發戶.
即時做了暴發戶.
大家知道暴發戶是什麼嗎?.
不知道嗎?.
不知道不要緊,說一點.
暴發戶原來有英文名.
他叫做New Money.
是形容一些人突然間變得很富有.
不知道什麼原因.
可能是突然間賣了地.
突然間很富有.
這些很富有的人就有一個.
非正式統計的模式.
就是他很喜歡炫耀式消費.
什麼是炫耀式消費呢?.

$^{41}$他可能是買一些他根本不需要的東西.
但純粹想表現他有多麼有錢.
純粹想表現他有多麼有錢.
展示他有多麼有錢.
所以他有一個特徵.
就是炫耀式消費的觀念.
弟兄姊妹,有沒有遇過,有沒有見過?.
他們可能是非正式統計.
暴發戶可能會買一些名牌.
告訴大家他很有錢.
相信我們當中沒有暴發戶.
我們很謙卑.
就算有暴發戶我們也沒有炫耀式消費.
但原來昔日在哥林多的時候.
他們有一群暴發戶.
但不是金錢的暴發戶.
他們是屬靈的暴發戶.
什麼是屬靈的暴發戶呢?.
原來當時在哥林多的教會裡的那群信徒.
突然間在很短時間內.
有一些屬靈恩賜.
包括講方言,先知講道等等.
所以突然間他們這一群人.
很短期內很喜歡炫耀式消費.
他們的屬靈恩賜.
很喜歡展示他們的屬靈恩賜.
所以今天我給他們一個題目.
這一群屬靈暴發戶.
當時保羅如何處理這群屬靈暴發戶呢?.
當時在哥林多的時候.
哥林多是一個很新興的城市.
當時很富裕.
很多元文化.
很多哲學在裡面.
算是一個非常富裕的港口城市.
我們都知道.
保羅去宣教的時候.
他不是長時間留在哥林多.
他離開之後發現.
為什麼那群信徒.

$^{81}$慢慢地相信了一些其他的理論.
我們稍後再說.
他究竟相信了什麼.
使他有炫耀式的消費觀念.
相信我們今天.
剛剛說了.
弟兄姊妹都沒有屬靈暴發戶在當中.
但是今天的提醒.
我們都知道暴發戶是突然間暴發的.
沒有人知道什麼時候發了.
所以今天保羅這個提醒.
我們只願意記得.
如果有一天我們突然間.
不小心做了屬靈暴發戶的時候.
我們都記得他的提醒.
今天我們看哥林多前書十三章.
在十二章尾.
保羅說.
我指示你們一條更妙的道.
更妙的道.
所以引入第十三章.
這個道是什麼更妙的道呢.
我們一起看看.
不如一起讀.
今天我們看《環球聖經》譯本.
哥林多前書十三章一至八節.
如果我用人的方言說話.
甚至用天使的方言.
沒有愛.
我就成了曹擇的羅或次爾的拔.
如果我有先知信息.
也明白一切奧秘.
一切知識.
如果我有傳遞的信心.
可以移散.
卻沒有愛.
我就什麼也不是.
如果我把所有的一切周濟人.
如果我捨身被人焚燒.
卻沒有愛.

$^{121}$這對我毫無幫助.
愛是寬容.
愛是恩慈.
不嫉妒.
不自吹自擂.
不自高自大.
不做不合體統的事.
不求自己的益處.
不輕易被激怒.
不圖謀惡事.
不喜歡行不義.
而為真理歡喜.
凡是包容.
凡是相信.
凡是盼望.
凡是堅忍.
愛永不倒下.
而先知信息將被廢棄.
方言將會停止.
知識將被廢棄.
我們再一次低頭祈禱.
親愛的主我們.
求你讓我們謙卑來到你面前.
聽你的話語.
父啊我們沒有一個願意成為.
屬靈步伐戶.
我們沒有一個願意炫耀自己的生命.
但是父啊.
求你的話語成為我們的提醒.
成為我們的幫助.
願意我們懷著感恩.
謙卑的心去領受你的教導.
天父多謝你.
我們同心禱告奉主明求.
Amen.
剛剛保羅一開始.
他頭三節.
已經很連珠爆發地.
出現了五次的如果.
在原文那裡.

$^{161}$是有這五個如果.
看和合本.
所以我今天看《環球聖經》譯本.
比較貼近原文地.
他譯了五個如果出來.
所以當時.
保羅其實是很肉緊地.
指出哥林多教會的一些事情.
保羅說什麼呢.
他說如果我用人的方言說話.
甚至用天使的方言.
如果我用這些天使的方言.
如果我有傳遞的信心可以移散.
如果我將所有的一切周祭人.
如果我寫幾生被焚燒.
但是卻沒有愛.
一切都是徒然.
就是說.
如果縱使我有種種的恩賜.
種種的成就.
我就算多偉大去周祭窮人.
做了多少善事.
我就算欠了一條命.
但如果我沒有愛的話.
一切都是徒然.
到底當時哥林多教會出現了什麼問題.
發生了什麼事呢.
我們一起看看.
哥林多教會.
他們當時剛剛說.
有一群信徒.
他們相信一些屬靈恩賜.
是特別高超的.
包括講方言.
和先知講道.
這些恩賜他們覺得是很高超的成就.
除了這些想法之外.
還有什麼呢.
原來他們有很多的思想偏離了真理.
所以哥林多前書一直在寫.

$^{201}$保羅斥責他們發生了什麼事.
例如保羅說他們分黨派.
你們分黨派.
所以保羅說你們中間不可以這樣.
一個說屬保羅.
一個說屬基督.
不可以這樣分黨派.
當時那個情況是怎樣呢.
好像說.
如果我們當中有人跟Billy Graham報道.
有人跟堂頌榮.
Jack跟堂頌榮.
之後有人問你是跟誰.
我說我是LG的.
那些人就滴汗了.
當時的情況就是這樣.
表面上大家好像很熟齡.
很追求一些學問.
跟師傅學習.
但實質上是在比較大家.
是否有一個更高的地位.
還有倫理.
性倫理上的混亂.
保羅指責他們.
你有什麼理由可以這樣取你的繼母.
取了繼母.
還有什麼呢.
還有爭權奪利.
告上法庭.
第六章說.
還有什麼呢.
恩賜上比較.
不單是覺得自己的恩賜是超人一等.
還要跟其他的弟兄姊妹比較.
例如說怎樣比較呢.
他們覺得講方言.
和先知講道.
是特別高超的.
最頂級的恩賜.
其他的恩賜都是比較次等的.

$^{241}$例如今天我們有款待的恩賜.
我們有人喜歡煮東西.
有人喜歡敬拜.
神給我們不同的恩賜.
明明基督是一起擺上我們恩賜去服侍教會.
但當時哥林多的信徒.
他覺得講方言和先知講道的恩賜是最頂級的.
你煮飯,招待那些.
沒什麼.
絕對不是一個什麼.
當時就是這樣比較.
所以質度.
保羅說你們當中有質度紛爭.
還有什麼呢.
聚會裡面很混亂.
怎樣混亂呢.
聖經裡面說.
十四章.
保羅斥責他們說你們講方言的話.
方言是什麼恩賜呢.
其實是用一些舌頭發出的聲音.
先不說是真方言還是假方言.
用舌頭發出的聲音去祈禱.
去講這些方言的時候.
他覺得是特別高人一等.
但保羅說十四章.
如果你要用方言的恩賜.
你要造就人.
沒有人懂得聽你的恩賜.
你搞亂了整個聚會.
如果你在說一些人家不懂的東西.
你不是造就人.
沒有人幫你翻譯方言的話.
那你不要說.
上年我們去柬埔寨短孫.
因為是去機構服侍.
所以星期日我們可以去當地的教會.
當地的教會講柬文.
他有英文的翻譯.
所以我們就帶著耳機聽著APP.

$^{281}$其實是即時翻譯的.
但就透過一個耳機幫我們翻譯.
一開始也沒什麼的.
我們照樣敬拜.
英文歌敬拜.
聽英文翻譯給我們.
去到最後牧師就祈禱.
一聽的時候.
咦?為什麼講方言呢?.
我們等了一會.
不知道會不會翻譯了方言給我們聽.
我們等了一會.
已經沒有人翻方言.
我們就走了.
那個情況也挺搞笑的.
我們第一個反應就是.
方言祈禱又沒有人翻方言.
我們就在等.
一個搞笑的情況.
但當時在哥林多的教會裡.
講方言是一個他們認為.
很超班的屬靈恩賜.
為什麼呢?.
為什麼會形成這個現象呢?.
哥林多這個地方是一個新興城市.
也很富裕的城市.
裡面的人是很有知識的.
受很多的教育.
很有文化.
但為什麼會出現.
教會出現這樣的問題呢?.
原來當時最明顯是因為.
他們的思想裡出現了.
一個過分實現的末世觀.
過分實現的末世觀.
為什麼呢?.
我們經常提到諾斯底主義.
其實整個新約的背景.
也受諾斯底主義裡很多的影響.
但諾斯底主義在這個哥林多的時期.

$^{321}$是未成形的.
只是一個初形.
未完全破壞教會.
但你知道一個主義的思想.
並不是突然進入破壞教會.
它通常慢慢滲入.
當時為什麼那群哥林多信徒.
可以做出這麼奇怪的事情呢?.
就是因為諾斯底主義.
混合了福音的真理.
先說說諾斯底主義是什麼.
主要是強調什麼.
諾斯底主義主要是強調.
智慧很重要.
所以他們相信得著智慧有竅門.
而哥林多的信徒.
當時他們就相信要得著智慧.
就要做一個屬靈的人.
做一個屬靈的人就要有屬靈的恩賜.
屬靈的恩賜最高級別就是講方言和先知講道.
所以明明他的推論是怎樣來的.
因為他們相信要得著智慧.
要得著智慧就一定是屬靈人.
屬靈人怎樣證明自己是屬靈人呢?.
就要有屬靈的恩賜.
怎樣證明自己有屬靈的恩賜呢?.
就要講方言和先知講道.
最嚴重最麻煩的是怎樣呢?.
他們覺得如果他是屬靈人.
有屬靈恩賜的時候.
他講方言其實是跟天使說話.
他不是跟人說話.
他就覺得高人一等.
當他可以講方言的時候.
他就覺得我是天使的層次.
我不是人的層次.
所以這就是他背後的概念.
他跟天使說話.
跟天使說話的時候.
也代表著他相信.

$^{361}$這是歪理.
也代表著既然他能夠跟天使說話.
即是他是屬靈人.
他是屬靈人的時候.
在地上他所犯的.
他身體上所犯的事.
都不重要.
只要他還有這個恩賜.
就代表他是屬靈.
所以他整個思想就是這樣建構出來.
其實也解釋了為何可以如此淫亂.
為何他身體可以如此淫亂.
可以娶繼母.
可以亂來.
因為他覺得身體不重要.
他的靈已經跟天使是同一個層次.
所以這些混亂的思想.
加上福音.
混合在一起的時候.
進入了教會裡面.
這就是為何當時.
他們會有這些屬靈的超越感.
為何會有這些屬靈的超越感存在.
但是保羅當時就直斥他們.
他說你們這些恩賜.
真是算不得甚麼.
算不得甚麼.
如果你有這些恩賜.
但是沒有愛的時候.
只有個「曹」字.
保羅說你曹集的羅 次義的拔.
你引起大家的注意.
但是沒有意思.
甚麼是羅 甚麼是拔呢.
原來就是當時一些異教很喜歡用的工具.
去敬拜他們的神的一些工具.
或者一些儀式.
所以只有個「曹」字.
保羅說如果你有這些廣方言的恩賜.
但是你沒有愛在裡面.

$^{401}$你只有個「曹」字.
弟兄姊妹我們有沒有一些東西呢.
自以為自己很有成就.
很自豪 很自豪.
但是神說.
如果我們有這些成就.
有這些才剛.
如果沒有愛的話.
這些東西都是Nothing.
神說這些東西都是Nothing.
保羅給這些人一個屬靈治療的方案.
甚麼方案呢.
他們驕傲到極點.
但是保羅給他們甚麼方案呢.
保羅說用愛.
用愛.
愛是最大最Powerful.
一般我們講治療的時候.
有三種 有愛.
有內服 有外敷.
有情緒.
所以今天我用三樣東西.
保羅怎樣提出愛的治療.
三樣東西.
首先講內服.
我們要有愛的內涵.
弟兄姊妹 內涵是怎樣來的呢.
內涵我們經常說是培養出來的.
坊間世界都告訴我們.
內涵是培養出來的.
但是愛的內涵是培養不了的.
愛的內涵我們培養不了的.
為甚麼這樣說呢.
因為愛是需要有愛的種子在我們心裡.
如果我們有很漂亮的泥土.
我們不停淋水.
不停曬太陽.
不停施肥.
但是裡面沒有種子的話.
是不會有花出來的.

$^{441}$所以我們要有愛的內涵的時候.
我們要有神在我們生命裡面.
才種得出這種愛.
要有神在我們生命裡面.
用神這個種子在我們裡面做源頭.
第一 愛是寬容.
愛是寬容 愛是恩慈.
其實就是說神的屬性在我們裡面的時候.
就彰顯出這種愛.
不是靠我們自己.
從來都不是靠我們自己.
愛是寬容這個詞.
寬容的相反是急躁.
這個觀念在聖經裡面.
只是用在神身上.
寬容很多時候形容在神的身上.
神是寬容的神.
聖經裡面很多都說到神是怎樣寬容.
特別在有一句經文.
特別特別提醒我們.
神是怎樣寬容我們.
彼得厚書裡說.
主所應許的尚未成就.
有人以為他是擔言.
其實不是擔言.
乃是寬容你們.
因為神不願意一人沉淪.
乃願人人都悔改.
有人疑問.
為什麼神你又說要罰那個人.
你又不罰.
神你又說來救我.
你又不救.
特別在身體抱恙的時候.
特別在面對挑戰的時候.
我們經常都問神.
神你又說來救我.
為什麼你不帶我回天家.
但原來神.
一個最大最大的寬容.

$^{481}$去達到極致.
就是祂願意等候.
所以和合本說.
愛是恆久忍耐.
《環球聖經》譯本譯做寬容.
神願意寬容我們.
祂願意等候我們.
這個是神最極致極致.
接著又說到.
愛是恩慈.
恩慈這個字很奇妙.
《聖經》裡只出現過一次.
只出現過一次.
而我們一知道恩慈.
一定是形容神.
有恩典,有人事,有憐憫.
這個字很特別.
很特別.
恩慈的人可以怎樣.
可以善待一些對自己不好的人.
神透過耶穌基督來到這個世界上.
就是極致地表達了神這種恩慈.
無論別人怎樣對耶穌基督不好.
祂仍然願意主動地去拯救.
這個是神極致的恩慈.
神的恩慈是祂主動的施予.
縱使人是不配.
所以當時高林多的信徒.
他自高自大到極點.
因為我們看高林多前書的時候.
其實這個字出現了很多次.
不停被保羅說他自高自大.
但最奇妙的是.
保羅提出神的愛是恩慈.
無論他們怎樣自高自大.
神恩慈的愛沒有放棄過他們.
同樣神也是這樣對我們.
祂也是對我們恆久忍耐又有恩慈.
弟兄姊妹.
你相不相信.

$^{521}$人是沒有能力去克服自己的自高自大.
如果我是一個驕傲的人.
我沒有辦法叫自己不要這麼驕傲.
相不相信.
同樣我們沒有辦法去克服.
自己心裡面對別人的一些判斷.
很感恩我聽過一個弟兄的分享.
一個很真誠的分享.
很感恩弟兄會跟我這樣分享.
他說他未信主之前很驕傲.
他表面上會對別人笑哈哈.
但實際上他心裡面對別人很多判斷.
去小組的時候表面上笑哈哈.
但心裡面別人分享的時候.
會判斷別人的判斷.
心裡面他會一邊聽.
但心裡一直在說別人.
很感恩這個弟兄的分享.
因為他很真實的知道.
在認識神之前是什麼人.
神就是當他認識了神之後.
神的愛充滿了他的時候.
這些驕傲都放下了.
這個不是靠他的能力可以改變.
是因為神的愛充滿了他.
充滿了他,他就可以謙卑下來.
所以弟兄姊妹我們都要學.
讓神的愛充滿著我們.
內服.
接著說一下外服是什麼.
外服是講愛的行為.
我們知道.
我們不是靠修行.
愛的行動不是靠修行.
不是靠什麼行為.
德行去賺取任何東西.
但是因為愛是聖靈的果子.
聖靈的果子.
如果我們有愛在我們裡面的時候.
很自然就可以流露一些行為的模式.

$^{561}$但是不代表.
就算做到那件事不代表很有愛.
但如果有愛的時候.
很自然可以流露這些模式.
例如這些模式是什麼呢?.
愛的時候不做不合體統的事.
不求自己的益處.
什麼是不做不合體統的事呢?.
不合體統的事就是不光彩的事.
一些不光彩的事.
一些不應該做的事.
損害了別人尊嚴的事.
其實保羅想事情很有邏輯.
我們看到不做不合體統的事.
其實和後面那句.
不求自己的益處很有關係.
因為通常我們做不合體統的事.
都是為了自己的益處.
所以保羅說不要做不合體統的事.
不要為自己的益處.
我們很多時候無端端不會傷害別人的尊嚴.
不會傷害別人的尊嚴.
不會做不光彩的事.
除非那件事很多時候和自己的利益有關.
最極端的例子就是《飢餓遊戲》那套戲.
《飢餓遊戲》那套戲不是遊戲.
很殘忍.
那些人為了保住自己的性命.
在荒島裡面殺了一些他不認識的人.
去到最後誰死到最盡就留下性命.
人就是這樣.
為了自己的益處.
可以做一些不合體統的事.
還記得COVID那幾年.
我們不是我們當中.
其實是灣區一個這麼有文化.
這麼有教育的城市.
還記得COVID那幾年.
有人會因為爭口罩.
有人會因為一些Hate.

$^{601}$所以到處向人吐口水.
所以我們見到這個不合體統的事.
當人去到一個為了自己的利益的時候.
我們是想像不到他可以怎樣不合體統.
但是保羅說.
用愛,要有愛在我們裡面.
不可以做這些不光彩的事.
不可以餵著自己的益處.
用愛.
愛還有什麼呢?.
不圖謀惡事.
我們熟悉的和諧本說:不計算人的惡.
我看英文的NASB.
更加給我一個新的體會.
那裡說什麼?.
It does not take into account a wrong suffered.
意思是什麼呢?.
一個更深的層次.
不單止不計算別人做錯了的事.
不單止不計算別人的惡事.
更加不計算自己遭受了多少冤屈和不公平.
更加不計算自己有多少suffer.
而圖謀做惡事.
不會這樣做,不會.
意思是當別人對自己有惡意的時候.
不單止不計算那個人有多差.
更加不計算他對自己.
令自己有多委屈.
令自己有多慘.
為什麼呢?.
為什麼要這樣呢?.
因為原來當我們不計算自己有多委屈的時候.
我們就已防止自己進入一個自憐的狀態.
一個自憐,憤世嫉俗的狀態.
保羅說你不要進入這個狀態.
因為你進入這個狀態很容易得罪神.
一個自憐和憤世嫉俗的人很容易得罪神.
所以保羅說不要計算別人的惡事.
不要圖謀惡事.
就好像打球,踢球.

$^{641}$我們當中有沒有人踢球?.
有沒有人打球?.
當我們打球的時候.
如果對方那一隊又撞,又踢,又推.
犯規一樣.
我們會不會同樣這樣做呢?.
我們不會的.
他踢我,我不會踢他.
就好像一場打球的賽事.
我們不會因為對方攻擊我們自己.
而攻擊他.
同樣我們不會因為對方對自己的惡.
而作出一些惡事得罪神.
我們不會做一些惡事得罪神.
不知道大家覺得自己是否一個口齒伶俐的人.
大家覺得自己口齒伶俐嗎?.
通常跟別人辯論的時候.
是誰贏?.
我覺得自己不是口齒伶俐的人.
通常我跟我先生辯論的時候.
是很好笑的.
我經常說著說著就忘記自己在說什麼.
然後就會哈哈哈哈.
就完事了.
我明明很生他的氣.
說著說著就忘記自己在說什麼.
就完事了.
我們可以不計算對方的惡.
不計算對方的惡.
不是靠自己.
不是靠我自己的能力.
我沒辦法不計算對方的惡.
但神的愛在我們當中的時候.
就可以做到.
這不是死忍的.
千萬不可以死谷爛谷.
令自己不要計算對方的惡.
因為這樣忍的時候.
會冤枉到病.
冤枉到病.

$^{681}$冤枉一下就會一個大反彈爆出來.
所以我們不是靠自己去死忍.
在2023年做過一個統計.
就是全年在2023年.
總共有四萬五千多人死於因為槍的緣故.
可能是被人槍殺.
或者是用槍自殺.
四萬五千多宗.
有二萬七千多宗.
二萬七千多宗是用槍自殺.
一萬七千多宗是用槍殺人.
其實這個數字是很恐怖的.
但為什麼會有這些事發生呢?.
就是因為很多人.
他們被迫.
迫到一個點的時候.
他們的怒氣.
或者他們的屈怨.
就要麼假裝在自己身上.
要麼就假裝在別人身上.
所以弟兄姊妹.
哥林多前書保羅教我們.
不是屈屈埋埋.
叫我們不要計算別人的惡.
我們就死忍.
不計算別人的惡.
我們要回到根本.
用愛去融化我們生命裡的元氣.
用愛融化我們生命的元氣.
只有用愛才可以不計算別人的惡.
弟兄姊妹你覺得.
你容易不計算別人的惡嗎?.
你容易嗎?.
在我們未有神.
在我們生命裡面的時候.
我們有一個反射機制.
很自然有一個反射機制.
就是別人如果有箭射過來.
我們就很自然地射他.
未有神在我們生命裡面的時候.

$^{721}$有神在我們生命裡面的時候.
我們就很自然地射過來.
我們很自然地有個盾牌擋一擋.
但有時都是死頂.
有時是死頂.
擋到有點累的時候.
有點憤怒就射回去.
但我們期待今天更深的學習.
是用愛.
用愛融化這些元氣.
當有箭再射過來的時候.
我們祈禱求神幫助我們.
先用愛去融化這些元氣.
融化這些元氣的時候.
就不用再計算.
不容易的.
但我們靠禱告.
靠禱告,弟兄姊妹.
第三樣東西.
治療方案.
講心境.
好像今早Tucker講.
我們講emotional health.
怎樣呢?.
有愛在我們生命裡.
我們生命的態度.
會有少許不同.
特別是怎樣呢?.
我們愛不是麻木的寬容.
不是麻木的忍受.
我們是喜歡真理.
只是喜歡真理.
意思是甚麼呢?.
愛是不會愛到不分對錯.
愛不會愛到縱容人.
如果我們見到一個很慘的賊.
牽著一家大小一起去打劫.
很慘的小朋友沒東西吃.
但不會因為這樣而不報警.
一定要報警.

$^{761}$有愛就是有真理.
因為聖經說的.
真愛是不會恨不義.
只喜歡真理.
大家有沒有聽過一個term.
叫PUA.
PUA.
PUA本身這東西.
本來最初的時候.
是一些心理學家研究出來的.
一些溝通的技巧.
本身是中性.
但現在近年.
不知道近多少年.
現在已經形成一個比較.
不是那一樣東西.
不是溝通技巧那麼簡單.
就是PUA這個term.
是理解成一個情緒的操控.
因為PUA這個心理的技巧.
當被濫用的時候.
真的會造成一種情緒的操控.
所以現在當說PUA的時候.
就好像代表著情緒的操控.
說一下PUA是什麼呢.
主要就是說.
其實PUA在我們身上很多很多.
因為我們不停和人溝通.
不停和人溝通.
就有機會發生這些事情.
就有機會被人PUA你.
或者你PUA別人.
我們看一下.
PUA到底是什麼呢.
就是一種技巧.
去打擊對方的自信.
形成一些關係模式.
是務求怎樣呢.
可能令你依賴他.
情侶的話可能會跟你說.

$^{801}$你不漂亮的.
只有我要你.
然後你老闆可能跟你說.
你學歷這麼差 經驗又沒有.
是我才願意請你 用你這麼多年.
這些就是PUA的手法.
例如家裡怎樣PUA呢.
可能跟你說.
我也是因為愛你才打你 才關心你.
有些夫婦.
不應該說夫婦.
因為愛你才這樣對你.
我也不想罵你 發你脾氣.
不過我實在是太相信你了.
所以我要在你面前發脾氣 罵你.
等等等等.
他刻意令你相信.
他在你的生命裡很重要.
他刻意踩低你 抬高自己.
簡單來說就是一種PUA的技巧.
弟兄姊妹 知不知道我們基督徒 經常都被PUA.
基督徒經常都被PUA 想想.
想想好不好.
例如會怎樣跟你說.
你是基督徒 你要原諒別人.
你會怎樣跟你說.
你是基督徒 你應該愛的.
我記得我差點被人PUA.
我跟我第一個女兒的時候 作小產.
她去看醫生 看婦科.
婦科解釋完.
真的有機會流產.
但最後她回答了一句.
你是基督徒 你看得很開.
我那一刻就 嗯.
你說她算不算是PUA.
或者我有沒有被PUA.
就是在乎她是否經常重複用這些溝通方法.
而我有沒有被PUA.
就是在乎我的心態有沒有調節.

$^{841}$有沒有被這東西推倒了我.
所以這就是弟兄姊妹.
我們會不會容易被人PUA我們.
我們要捉住真理很重要.
捉住真理很重要.
神說愛是只喜歡真理.
不喜歡恨不義.
當我們要捉住真理的教導在我們心裡的時候.
別人就沒那麼容易PUA到我們.
別人沒那麼容易PUA到我們.
就是說你愛嘛 你的神說愛.
你不愛我 為什麼我問你借錢你不借給我.
為什麼你不原諒我.
我們在真理裡面.
我們的愛在真理裡面.
有真理就分懂什麼是不義.
一個恨不義的人是怎樣呢.
可能是當一個傷害自己的人.
若干年後他得到某程度上的惡果的時候.
我們的心會不會覺得很涼快呢.
還是我們見到的時候有憐憫呢.
就是reflect到我們有沒有不義.
我們是不是只喜歡真理.
弟兄姊妹我們要記得.
愛是在真理裡面.
所以我們要經常熟讀聖經.
經常熟讀聖經認識真理.
聖經裡面說凡是包容 凡是相信.
凡是盼望 凡是堅忍.
但是不離開真理.
不離開真理.
不可以死忍難忍.
但是離開了聖經的教導.
弟兄姊妹我們有什麼可以傳到永恆呢.
我們有什麼可以傳到永恆呢.
永生我們知道.
但是我們有很多關乎永恆的事情.
其實我們沒有能力去參與.
例如我們的永生.
我們沒有能力去令自己將來的永生更加漂亮.

$^{881}$更加豐富.
因為神說你將來上天階的時候.
那個樣子一定是最漂亮的.
所以我們不能靠自己去做什麼.
去令永恆有什麼改變.
但是聖經裡面說.
唯有愛.
唯有愛這件事可以在我們今生彰顯出來.
也是用到永恆的.
所以第八字是說愛永不倒下.
愛永不止息.
唯有愛這件事是我們可以主動.
不是被動.
也有永恆的價值.
保羅說一切的先知信息都要被廢去.
所有的恩賜都要被廢去.
唯有愛是達至永恆.
說一些見證.
我想我們很多弟兄姊妹都認識.
我們以前一位前教會的姊妹.
叫Vivian Fong.
一位姊妹突然間很短的時間被神接走了上天階.
在八天之內.
我們一班姊妹一起禱告.
我們在理性上知道她上天階.
情感上我們不捨得.
但當我們數算這位姊妹在我們生命裡的種種同行.
這位姊妹如何服侍神.
這位姊妹如何愛心的時候.
我們很體會到她的愛真的達至永恆.
不單止在這裡.
不單止在今生彰顯了.
她這種來自神的愛.
是真的達至永恆的價值.
所以弟兄姊妹說.
愛是真真實實的.
永不止息.
願意我們都不是憑血氣去愛.
用神的愛去愛人.
我們學習一起用神的愛去愛人.

$^{921}$保羅用了五個如果.
一開始的時候用了五個如果.
今天我很想用一個如果去提醒弟兄姊妹.
今年記不記得我們的主題是什麼.
Fever.
記不記得Fever是指什麼.
有提示.
Faith in Fellowship and Action.
神父召我們用信心用團契用行動.
去將我們的Mission.
將耶穌基督帶到我們的鄰居.
帶到我們的Community.
這是我們今年的主題.
所以今年我們發覺.
很多時候我們一起去推廣福音的事工.
我們接下來會有Mission Card.
我們會有短訊.
我們會一起去做一些Community Planning.
一些Outreach.
我們積極做福音事工.
積極做這些事的時候.
如果我們Fever大發熱心.
又獻出所有的恩賜.
卻沒有愛的話.
就什麼都不是.
我們就毫無幫助.
弟兄姊妹.
今天保羅提醒我們.
我們要用愛.
要有神的愛在我們當中.
我們不是追求恩賜.
不是追求資源上的增多.
我們求神的愛在我們生命裡面.
充滿我們.
很感恩今天有個姐妹又說多一個見證.
我們Mission不是用教導教回來.
是神的愛在我們裡面去改變我們.
記不記得上次我們說有隻雞腿.
大家派雞腿.
我知道我聽到好幾個弟兄姊妹分享.

$^{961}$我知道那些弟兄姊妹本身是Doer.
她不多說話.
也不多主動邀請身邊的人.
特別是不認識的人.
她是不會的.
但是今年Fever就是提醒她.
有一個姐妹去買菜.
她問隔壁的主婦.
你有沒有回教會.
我見到弟兄姊妹有很多突破.
因為一個Mission Mind.
願意今年我們一起去Fever的時候.
大發熱心去擺上我們恩賜的時候.
我們更加要被神的愛充滿.
弟兄姊妹我們一起低頭祈禱.
親愛的Joy多謝你.
多謝你告訴我們你是愛的源頭.
我們擁有一切的東西.
但如果沒有愛的話.
我們就什麼都不是.
多謝你透過保羅去提醒我們.
給我們一個更妙的道.
就是用愛去服侍.
用愛去奉獻.
天父啊.
我們生命裡面.
曾經有種種得罪你.
種種的至高至大.
好像哥林多教會的信徒.
我們都曾經有各方面的得罪你.
天父啊,多謝你用愛去囂觀去原諒.
多謝你對我們恆久忍耐.
以致我們有機會去接受這個福音.
天父啊,求你幫助我們.
特別是今年,亦都是以後.
我們Living Grace的弟兄姊妹.
可以用你的愛去服侍.
去帶更多的人去認識耶穌基督.
天父啊,多謝你.
多謝你聽我們禱告.

$^{1001}$我們同心禱告奉主耶穌基督的命求.
Amen.
\newpage



\section{詩篇 8:1-9}
\label{sec:031663}
\textbf{星空與塵土 (普)}~~~~張磊牧師
\newline
\newline
連結: \href{http://www.church.com.hk/acms/content.asp?site=cdc\&op=show\&type=product\&code=031663\&layout=sermon}{\texttt{031663}} ~~~~ 語音日期: 2026-02-15
\newline
\newline
\hyperref[sec:031662]{\small{< < < PREV SERMON < < <}}
~
\hyperref[sec:index_chronic]{\small{[返順時目]}}
~
\hyperref[sec:031664]{\small{> > > NEXT SERMON > > >}}
\newline
\newline
詩篇 8:1-9
\newline
\begin{longtable}{cl}
\hline
\hline
章節 & 經文 (和合本修訂版)\\
\hline
8:1 & \begin{tabularx}{0.7\textwidth}{X} 耶和華－我們的主啊，你的名在全地何其美！你將你的榮耀彰顯於天。 \end{tabularx} \\ \\ \relax
8:2 & \begin{tabularx}{0.7\textwidth}{X} 你因敵人的緣故，從孩童和吃奶的口中建立了能力，使仇敵和報仇的閉口無言。 \end{tabularx} \\ \\ \relax
8:3 & \begin{tabularx}{0.7\textwidth}{X} 我觀看你手指所造的天，並你所陳設的月亮星宿。 \end{tabularx} \\ \\ \relax
8:4 & \begin{tabularx}{0.7\textwidth}{X} 人算甚麼，你竟顧念他！世人算甚麼，你竟眷顧他！ \end{tabularx} \\ \\ \relax
8:5 & \begin{tabularx}{0.7\textwidth}{X} 你使他比神微小一點，賜他榮耀尊貴為冠冕。 \end{tabularx} \\ \\ \relax
8:6 & \begin{tabularx}{0.7\textwidth}{X} 你派他管理你手所造的，使萬物， \end{tabularx} \\ \\ \relax
8:7 & \begin{tabularx}{0.7\textwidth}{X} 就是一切的牛羊、田野的牲畜、空中的鳥、海裡的魚，凡游在水裡的， \end{tabularx} \\ \\ \relax
8:8 & \begin{tabularx}{0.7\textwidth}{X} 都服在他的腳下。 \end{tabularx} \\ \\ \relax
8:9 & \begin{tabularx}{0.7\textwidth}{X} 耶和華－我們的主啊，你的名在全地何其美！ \end{tabularx} \\ \\
[1ex]
\hline
\hline
\end{longtable}
$^{1}$其實之所以選詩篇就是因為今年我們應該一整年都會.
絕大多數的時間我們都會在晨禱一直讀詩篇.
這也是我們今年讀過的一篇詩篇.
詩篇第八篇.
你知道詩篇本身就是在舊約猶太人的手中.
他們是詩歌本.
他們是在會堂當中敬拜.
聖殿當中會堂當中敬拜上帝用的.
但是你知道嗎.
今天我們讀了第八篇.
直到第八篇才是整個詩篇當中第一首純粹的讚美.
讚美類型的詩.
第八篇是第一篇.
是大衛仰望浩瀚的星空.
看到璀璨的星星和月亮.
那他就發出了這樣的一個讚美.
那整個詩篇的佈局.
弟兄姊妹我們讀詩篇不僅僅要看內容.
我們要看到的佈局.
我說了第八篇才是讚美.
那前面後面是什麼.
你如果看的話.
詩篇第一篇第二篇.
基本上我們會歸類為.
它是整個150篇詩篇的序言.
但從第三篇開始一直到第七篇.
這是詩篇的.
這五篇叫做哀歌.
然後從第九篇到十三篇又是哀歌.
那前半段的哀歌主要是以大衛逃亡的時候.
大衛落難的時候.
大衛在苦境當中的時候.
在神面前發出的一個哀歌.
那後半段這五篇呢.
是有一群窮苦人.
貧乏人.
有這樣的一個群體.
那他們也是被逼迫的時候.
被壓迫的時候.
發出的一種哀歌.

$^{41}$表達他們的痛苦.
那在這中間呢.
夾著正中間就是我們剛才讀的這一篇讚美.
對上帝的讚美.
弟兄姊妹.
如果我們只讀哀歌.
我們就會覺得信仰當中充滿了痛苦和掙紮.
如果我們只讀讚美.
我們就誤以為信仰當中從來就不會有任何的痛苦和挑戰.
所以真正的信仰是什麼呢.
是在哀歌當中去讚美.
在黑夜當中.
在苦難當中.
在最高潮的時候向神發出讚美.
所以我常常說一句話.
我說世人和基督徒你知道有什麼不一樣嗎.
世人呢.
他們笑著笑著就哭了.
神的兒女呢.
哭著哭著.
就笑起來了.
哈雷路亞.
所以各位.
你可以哀哭.
你可以困惑.
你可以感到你軟弱.
但是不要忘記要抬起頭來.
不要忘記要舉目.
要仰望.
要讚美.
坦白說今天.
雖然我們來到教會會讚美.
雖然我們有敬拜團.
我們有詩班.
我們有很好的樂器.
音樂音符.
但是通常我們不太會讚美.
因為我們的生活型態往往是面朝黃土背朝天的一種生活型態.
我們的問題主要就是在於背朝天.
為什麼.

$^{81}$因為我們很現實.
你知道嗎.
我們的主要的文化是一種實用主義的文化.
我們面朝電腦面朝手機.
但是我們所做的一切都如此的跟這塊地.
跟這塊黃土地相連.
我們活得太接地氣了.
我們不太會仰天而活.
我們不太會向死而生.
中國人你知道嗎.
通常只有在一種情況下會仰天長嘯.
什麼時候.
活不下去了.
對吧.
太痛苦了.
你知道.
你看到一個中國人仰天的時候.
你就知道他日子過得很痛苦.
他不太容易的時候.
他才會仰天.
詩篇第八篇是仰望夜空.
因為提到月亮和星宿.
詩篇十九篇是姊妹篇.
是看到旭日冬生.
因為提到太陽如永世.
如新郎出洞房.
對吧.
我們剛剛讀完十九篇.
星期五的時候.
葛木是在晨禱當中的時候.
分享了康德的一段話.
打出來給大家看了.
康德說什麼呢.
他在這一本《實踐理性批判》當中.
他說有兩樣東西.
我對他們的思考越是深沉和持久.
他們在我的心靈中喚起的讚歎和敬畏.
就越來越逆酒迷心有加無已.
一是我們頭上璀璨的星空.
另一個就是我們心中的道德律令.

$^{121}$這是康德.
一位哲學家.
我們不認為他是基督徒.
但是當他看到璀璨星空浩瀚星海的時候.
他發出讚歎.
他發出敬畏.
再看一個人.
康德如果你不認識.
這個人你從大陸來的.
你一定認識.
叫做溫家寶.
前總理.
他也寫過一篇詩叫做《仰望星空》.
我也念給大家.
仰望星空.
我仰望星空.
他是那樣莊嚴而聖潔.
那凜然的正義.
讓我充滿熱愛.
感到敬畏.
仰望星空.
我仰望星空.
他是那樣自由而寧靜.
那博大的胸懷.
讓我心靈熙熙.
意味仰望星空.
我仰望星空.
他是那樣壯麗而光輝.
那永恆的炙熱.
讓我堅守信念奮發圖強.
仰望星空.
我仰望星空.
他是那樣遼闊而深邃.
那無窮的真理.
讓我苦苦追尋上下求索.
各位.
這些人都不信耶穌.
無論是哲學家政治家.
無論是這個世界的.
很有智慧的人.

$^{161}$或者很有詮釋的人.
但是這些人仰望星空的時候.
會發出讚歎.
會發出上下去求索.
苦苦去追尋.
somehow有一個.
比讓我他們的詮釋.
他們的智慧以上.
還有一個真理.
他們說他們苦苦地在那裡尋索.
他們尋找了嗎.
他們找得到嗎.
弟兄姊妹今天你我來到教會.
我想對大家說.
我們對星空的仰望.
我們對天空的仰望.
至少要超過超過康德和溫家寶好不好.
所以不要太接地氣.
個人建議.
強烈推薦你們能夠把手機調成飛行模式.
因為我們現在要飛了.
我們要起飛了.
不要再接地上的新聞了.
不要再接地上的短信了.
各位主日.
主日就是不要讓你接地氣.
讓你仰望的.
讓你離地上騰的.
弟兄姊妹.
我們很多的問題.
都是在於此地樂不思蜀.
甚至也原諒我說.
此地甚苦還是不思蜀.
還是不仰望.
還是不想天.
天上有一個家.
天上有一位上帝.
巴不得神借著今天這首詩篇.
來帶領我們仰望.
來帶領我們讚美.

$^{201}$阿們.
詩篇當中中間有一句話說.
並你所陳設的月亮星宿.
應該念星宿.
對吧.
各位那裡的想這個字.
原文其實是去思想.
去默想.
去深思.
不光是抬頭看.
你的腦袋也要去思想.
思想什麼呢.
思想神指頭所造的天.
思想神所陳設的月亮和星宿.
想著想著你就會發出讚嘆.
我算什麼.
想著想著你就會發現.
神顧念我們.
上帝眷顧我們.
我預備的時候突然想起一首歌來.
你們一定都會.
你問我愛你有多深.
怎麼樣.
你去看一看.
你去想一想.
月亮代表我的心.
弟兄姊妹.
你知道主對你.
主的情不移.
主的愛不變.
你想要知道主對你愛有幾分嗎.
去仰望吧.
去仰望去默想去深思.
在這個人算什麼的感嘆.
在聖經當中出現過幾次.
除了出現在今天我們讀的經文以外.
詩篇還有一次以外.
出現在約伯記.
約伯記的十七章第七節怎麼說呢.
他說人算什麼.

$^{241}$你竟看他為大.
將他放在心上.
你竟將你的心.
將他放在心上.
但是這個另一個譯本.
或者更貼近原文的意思就是.
你竟將你的心思集中於他.
但是我很喜歡和和本.
放在心上.
因為我們中國人叫做心上人.
你知道嗎.
你是上帝的心上人.
約伯說人算什麼.
請問約伯在講這句話的時候.
約伯是不是有點生不如死.
約伯遭遇了什麼.
疾病.
財產盡失.
兒女死亡.
朋友的不理解.
妻子的咒罵.
能想到的壞事全在他身上.
在這樣的一種狀態當中.
你知道上帝怎麼樣把約伯.
從這樣的狀態當中帶出來的.
你記得嗎.
你如果參加成人哲學你就知道.
到了約伯紀的三十八章.
四十一章到四十一章.
四章的聖經.
上帝親自向約伯顯現了.
這一位宇宙的創造者.
親自向這一個活得生不如死.
苟延殘喘的約伯顯現.
回應他的訴求.
回應他心中的痛苦.
你知道這四章在講什麼.
這四章一共上帝.
反問約伯了七十多個問題.
帶約伯去看他指頭所造的天.

$^{281}$他的手所造的宇宙萬物.
包括海洋的深處.
地球的尺寸.
光和黑暗的界線.
水的循環.
還有雪.
冰雹.
風雨.
露水.
霜凍.
天上的星球.
各樣的動物對不對.
講到什麼.
有獅子.
有馬.
有烏鴉.
有鹿.
有野驢.
有鴕鳥.
有鷹.
你就稀奇了.
這不整個一個國家地理雜誌嗎.
但是看完這一切的時候.
看到四十二章.
到了四十二章的時候.
約伯就怎麼樣了.
就是剛才我們唱的那首詩歌.
從前我風雲有你.
我痛苦不得了.
現在我釋懷了.
現在我就釋懷了.
我就放下了.
我就坦然了.
豁然開朗了.
哇哦.
原來看看動物.
看看天空.
看看神所創造的萬物.
有這樣的功效.
聖經當中還有沒有這樣的人.

$^{321}$亞伯拉罕生不出孩子來.
膝下無子.
愁啊.
上帝幹嗎.
來來來.
我領著你數星星.
數得過來嗎.
你的後裔也要像天上的眾星那麼繁多.
是吧.
當諾亞經歷了整個大地遭到審判之後.
心還有餘悸.
然後呢.
上帝就把他叫出來.
諾亞諾亞來看什麼.
看彩虹.
看rainbow.
看到彩虹了嗎.
你的心就放心好了.
這裡有一個彩虹之約.
永子妹.
當人在那裡擔心.
明天吃什麼.
穿什麼.
耶穌讓我們看什麼.
啊.
看飛鳥.
看什麼.
看野花.
你看地裡的花.
你看天上的鳥.
有信心了吧.
永子妹.
是不是從舊約到新約.
我們都看到.
當我們看到這一切神所創造的時候.
我們就應該能看到神的永能和神性.
我們就應該能夠看到天父的愛和供應.
大衛在這裡說.
耶和華.
我們的主啊.

$^{361}$你的名在全地和其美.
你將你的榮耀彰顯於天.
永子妹.
這究竟是多浩瀚的一個宇宙.
我給大家看一個一兩分鐘的一個小video.
好嗎.
能幫我打開嗎.
這是在高科技園區.
地球看不見了.
這是太陽系.
銀河系.
接下來我也不說了.
因為到了我的知識盲區了.
這是人類的天文望遠鏡所能夠看到的極限.
但並不是盡頭.
謝謝.
各位.
怎麼樣什麼感覺.
我算什麼.
你算什麼.
各位.
地球算什麼.
太陽算什麼.
對不對.
有這麼感覺吧.
我看了.
因為為著預備這篇講到.
我看了一些.
越看越上頭.
我看到底是多大.
太陽系在銀河系當中是多大.
你知道多大嗎.
有一句話我倒是記得很清楚.
雖然看不懂那些光年那些數字.
已經beyond my imagination了.
他說太陽系就好像一個quarter.
那銀河系就是整個美國.
整個美國的一個quarter就是太陽系.
我們求主幫助我們.
我們比大衛更先進.

$^{401}$我們有望遠鏡了.
我們有哈勃望遠鏡了.
我們應該比大衛發出的讚美和驚嘆.
要更深邃.
對吧.
你知道在1961年.
那個時候4月12號.
這是第一次.
是蘇聯的宇航員.
叫做加加利.
他是第一個上天環繞地球一周的.
回來之後.
我就想到凱撒的一句名言.
我來了.
我看見我征服.
什麼意思.
他繞著地球走了一周之後.
他說我上天了.
我沒看到上帝.
沒有上帝.
然後他回來之後.
他就被蘇聯授予蘇聯英雄的勳章.
次年美國阿波羅8號登月.
三位宇航員在月亮上展開了聖經.
甚至其中他們有人念了.
誦讀了創世記的第一章.
祈諸神創造天地創造這個創造那個.
這是美國跑到月亮上去.
在那裡讚美上帝.
你知道嗎.
我查了一下.
這幾十年來.
至少進入太空的美國宇航員當中.
下來之後.
有71個人領受護照.
想要全世界服侍主.
其中有14個人.
的的確確他們走上了就辭去了工作.
進了神學院.
做了傳道人.

$^{441}$上天下來.
有一個人拿到了勳章.
說沒有事.
從天上下來.
也有一些人說上帝啊.
你太偉大了.
做什麼宇航員.
我要服侍你.
你想要做哪一種.
有人會說.
今天科技都已經這麼發達了.
人都已經上登月了.
你還信上帝.
各位我們怎麼回答他.
人都已經登月了.
你還不信上帝.
這裡說.
他將他的榮耀.
彰顯於天.
而且他的名在全地是何其美.
主耶和華.
我們的主你的名在全地何其美.
這是第一節和第九節對應起來的.
這是重複.
這是這首詩篇的主旋律.
這裡講說.
他的名充滿了全地.
神的名是什麼.
神的名其實就是上帝的屬性.
上帝的作為.
在聖經當中.
名.
神的名從來不只是一個發音.
一個稱呼.
神的名代表了他的作為.
代表了他是一位怎樣的神.
他怎樣介入我們的生活.
怎樣介入人類的歷史.
所以其實在整本聖經當中.
都強調神的名的重要和寶貴.

$^{481}$在十節的第三節說.
不可妄稱神的名.
在新約的主禱文說.
願人都尊你的名為聖.
各位我們不是要保護一個發音.
我們是要尊崇他的所事和他榮耀的作為.
我問大家一個問題.
舊約十節說不可妄稱神的名.
我問你們.
請問我們可不可以稱耶和華的名.
我們可不可以稱呼他的名.
可以對吧.
聖經說不可妄稱.
但是沒有說不可以稱.
其實大衛在這裡也稱耶和華這個名字.
但是在舊約你知道嗎.
以色列人基本上是不稱的.
他們不稱就避免了妄稱比較安全.
但是你知道對於他們.
這位神是不可妄稱其名.
那就乾脆不要稱.
不敢叫他的名字.
所以對他們而言.
這位神是高高在上的.
卻不是一位親密無接的.
卻不是一個以馬內利.
與我們有親密關係的.
耶和華這個名字.
其實現代有人.
因為他不稱.
其實已經不知道發音是什麼了.
沒有母音流傳下來了.
所以有人說叫雅威.
可能叫雅威.
也可能叫耶和華.
我們不曉得.
我們也不去分析他的舊約的原文.
但是你知道.
耶和華這個名字.
是他們上帝跟神的選民之間.

$^{521}$很特殊利約的一個名字.
如果是上帝的話.
那個字叫Elohim.
神.
上帝.
這是正式的名字.
但是在親密的人之間.
會有一個尼稱.
會有一個親密的稱呼.
專屬於你的.
各位.
你和一個人如果關係比較好.
或者你跟一個人關係不怎麼樣.
會不會在稱呼對方的時候.
用什麼樣的一個名字去稱呼對方.
是有講究的.
對吧.
你對不同的人.
會稱呼他的時候.
或者你跟他關係不同.
在不同的處境當中.
你會使用不同的名字.
對不對.
如果你要叫他全名的話.
那是什麼時候.
我不提名字了.
我常常聽我們的團體裡邊.
有人稱呼你全名.
我就知道.
這個要責備一下了.
各位.
我們有的時候.
我們不敢稱呼其名.
是因為我們跟他的關係.
沒有到那個份上.
對吧.
但是上帝賜給我們他的名字.
他把他的名字傳於我們.
讓我們可以奉他的名行事.
奉他的名禱告.

$^{561}$讓我們可以呼求他的名.
表明他要跟我們建立一個美好的關係.
那好了.
各位.
基督徒會不會妄稱神的名.
有沒有妄稱神的名的危險.
昨天神州團體我們一起在聚會的時候.
我們討論過這個topic.
對不對.
昨天晚上.
各位.
當然了.
有很多美國的電影文化亂稱Jesus.
亂稱Christ.
但是我們先不說這些.
因為我們基督徒應該不會犯這麼低級的錯誤吧.
但是你知不知道.
我在默想的時候.
我就想到.
你是一個基督徒.
神的名字本身就在你的身上.
不是嗎.
你是基督徒.
你的名字本身就有Christ.
庸姊妹.
我們身上就好像我姓張.
我身上帶著一個張字.
如果姓王你身上就帶著一個王字.
你是王家的人.
今天你是基督徒.
這個名字本身就烙印在了你的身上.
我問你們.
你們有沒有一種時候.
你不太希望別人知道你是一個基督徒.
你想把這個名字.
把這個稱呼隱藏起來.
有沒有這樣的時候.
有啊.
什麼時候.
做壞事的時候.

$^{601}$闖紅燈的時候.
錄怒證的時候.
千萬別讓別人認出我是個基督徒了.
各位.
怎麼叫做王稱神的名.
什麼叫做尊主的名為聖.
不是一個發音.
不是我們嘴巴上說的哪句話.
而是這個名字如影隨形.
已經綁定在我們所有的日常生活當中.
在這些經文當中.
其實有一個字.
其實在這章經文當中有一個字.
在這章出現了四次.
叫做全.
在全地和其名.
全扣希伯來字.
在這章聖經當中一共出現了四次.
第一節第九節.
剛才我們知道了.
然後中間翻譯成萬物.
翻譯成一切.
這個字強調的不是數量.
而是強調了一個完整性.
無例外.
無一例外.
各位.
神的名在整個世界.
包含了你所有的生活領域.
神的名都在其間.
你知道這個字最多出現在哪裡嗎.
我查考的時候.
我發現說原來最多出現在創世記的第一章和第二章.
出現了29次.
所以神創造一切走獸.
一切的飛鳥.
有的時候翻譯成各樣.
然後治理全地.
充滿全地.
有時候翻譯成全.

$^{641}$你會發現神在創造的時候.
就用這個字.
然後神就說.
造完了就說什麼.
希伯來文叫Tofu.
豆腐.
好.
甚好.
神的名在這一切.
神的作為在這一切當中.
甚好.
弟兄姊妹什麼叫做.
他的名在全地都美呢.
我引用我們中國的兩句話.
一個在詩篇當中.
一個在禮記當中.
叫做普天之下.
莫非王土.
率土之兵.
莫非王臣.
什麼意思.
就是整個天下.
普天下.
都是我王的土地.
都是我王的領土.
然後這個地上的子民.
這個地上的百姓.
無一例外.
都是屬他的百姓.
天無二日.
人無二處.
就好像天上永遠不會有兩個太陽.
我們的生命當中.
也不應該有兩個主宰.
他的名在全地.
他的名在你的全領域當中.
都何其美.
何其美這個字.
在中文當中翻譯的.
當然很有美感.

$^{681}$但是原文.
其實這個字是非常的.
非常廣泛的一個意思.
威遠.
雄壯.
壯麗.
壯觀.
Majesty.
有一首詩歌.
英文詩歌很有名.
對吧.
Majesty.
怎麼樣.
Worship His Majesty.
Unto Jesus.
我有點忘了.
Majesty.
這個字在英文聖經是這樣子翻譯的.
所以我列了不同的譯本.
新譯本.
翻譯成尊榮.
現代譯本.
翻譯成偉大.
呂珍中譯本.
翻譯成壯麗.
當代譯本很有意思.
你的名怎麼樣.
向徹大地.
各位.
不同的譯本.
用了不同的字.
都想要把這個含義給表達出來.
弟兄姊妹.
簡單地說.
這個名.
在普天之下.
在整個大地上.
都充滿了大能.
請問.
主的名.

$^{721}$與你而言.
有多大的威嚴.
有多大的能力.
有一本書.
親愛的.
別把上帝縮小了.
你說到我們的問題.
我們生命當中的掙扎.
往往有兩件事情.
第一.
你不知不覺就把上帝給縮小了.
第二.
你不知不覺.
就把魔鬼放大了.
就把世界給放大了.
弟兄姊妹.
求主幫助我們.
讓我們在這裡讚美他的majesty.
讓我們讚美他的名的能力的時候.
你知道什麼叫做讚美.
讚美這個意思.
其實有很多.
有承認.
高舉.
頌揚.
榮耀.
彰顯.
都包含在讚美的意思.
這個詞語當中.
那什麼叫做讚美神的名.
剛才有一段歌詞.
我們來讚美神的名.
什麼叫讚美神的名.
就讓神的屬性.
神的作為.
在我們中間再次顯大.
再次被高舉.
顯為高.
顯為大.
你說神本身就很高.

$^{761}$本身就是至高的至大的.
但是我再說.
我們的信心不夠.
我們的環境的艱難.
常常就不知不覺.
就把這位上帝給縮小了.
那在讚美當中.
要再次讓這位神顯大.
各位.
你記得.
出埃及進迦南之前.
那十二個探子嗎.
其中那十個探子.
去了迦南帝之後回來說什麼.
他說那些雅納族人不得了.
他們好像是巨人.
我們呢.
詐猛.
哇.
我想這是一定有誇張的成分.
對吧.
他就把人家給放大了.
他就把他的艱難給放大了.
他就把自己給縮小了.
但是加勒約什雅呢.
上去吧.
我們足能得勝.
他們不過是我們的食物.
完全.
到底是我們是他的食物.
究竟他是獵物.
我是食物.
還是我是獵手.
他是食物.
各位.
好像出現了一個顛倒.
用此美讚美.
就是這樣的一個意思.
在新約聖經當中.
其實好幾處引用了今天的詩篇.

$^{801}$我們先讀這一節.
在馬太福音當中.
其實有一群專業敬拜神的人.
叫做祭司.
叫做文士.
當這些專業敬拜的人.
職業professional.
在這敬拜神的人.
看到耶穌的神蹟騎士的時候.
就很不屑.
但是.
有一種人叫做小孩子.
小孩子在店裡就在那裡讚美.
和善吶歸於大衛的子孫.
這些專業人士不願意了.
你們跑掉了.
你們是誰啊.
你們是小子.
這是小毛孩子.
你們為什麼跑到店裡來這樣子.
但是耶穌怎麼說呢.
耶穌就引用了今天的經文.
你從嬰孩和吃奶的口中.
怎麼樣.
詩篇是怎麼說的呢.
建立能力使仇敵和報仇的閉口無言.
這裡把建立能力替換成什麼.
完全讚美.
完全了讚美.
各位.
讚美神.
讚美我們的神.
就是我們在神面前得著能力.
得著屬靈征戰的能力.
勝過仇敵的路徑.
你是一個讚美上帝的人嗎.
你說我是啊.
我常常唱歌.
我也很喜歡聽讚美詩.
甚至我也學習在我的話語當中去讚美神呢.

$^{841}$我問大家一個問題好不好.
讚美和表揚這兩者之間有區別嗎.
有啊.
有一天你說張姐妹你好漂亮啊.
李弟兄你好聰明啊.
或者今天你作為一個大人對小孩說你好乖啊.
好可愛啊.
各位.
這跟我們讚美神有什麼不同.
我們對他人的.
也可以說是一種讚美.
我們對他人的讚美裡邊.
其實不見得包含了我們對自己的否定.
對不對.
你們當中有沒有這麼讚美別人的.
李姊妹你好漂亮.
我有點醜.
張弟兄你好聰明啊.
就顯出了我的愚笨.
我在你面前confess.
這個讚美叫confess.
confess我的愚笨.
你們有沒有這樣讚美過別人.
應該不太會吧.
除非你特別自卑.
各位.
不以悔改.
不以懺悔作為前提.
不以自我否定作為前提的讚美.
這個叫做表揚.
我們團體當中有好多木道友.
多年的木道友.
對聖經也很熟.
也蠻喜歡耶穌.
耶穌很棒啊.
做得不錯.
我說你又表揚耶穌了.
你不是讚美.
當大衛在這裡讚美.
你的名在全體和其美的時候.

$^{881}$他一定帶出來的一個叫做什麼.
我算什麼.
是一種自我的否定.
看到你的偉大.
我就看到了自己的渺小.
看到了你的聖潔.
意味著我就認識到了我自己的污穢.
我已在你面前承認我的污穢.
來讚美你的聖潔.
承認我的愚拙.
來讚美你的智慧.
承認我的無能.
我的墮落.
我的敗壞.
來讚美你的崇高和偉大.
其實我在這個週選回應詩歌的時候.
我一直在兩首詩歌當中掙扎.
一首叫做人算什麼.
別說人算什麼.
你僅顧念他.
他算什麼.
這是一首.
後來我就想說.
行我要唱另外一首.
你真偉大.
其實這兩首如果能綁定在一起多好啊.
我就在想.
一方面承認我的渺小.
但是硬幣的另一面呢.
就是讚美主的偉大.
你知道現代的心理學.
跟我們所信的福音最大的差別是什麼.
心理學認為.
今天現代人的主要問題.
就是他們的自我形象太糟糕了.
把自己看得太低了.
所以要想辦法建立他們的自信心.
自尊心.
抬高他們的自我評價.
來解決他們生命當中的問題.

$^{921}$但是福音.
福音恰恰相反.
福音認為人主要的問題.
你主要的問題.
就是你把自己抬得太高了.
看的太重要了.
福音的主要目標就是把你拉下來.
讓你在這位至高的神面前破碎下來.
當然我也要說.
他拆毀他就重建.
他拔除他就栽植.
他積傷他就殘骨.
福音就是要讓你經歷這樣的一個過程.
讓你在上帝面前徹底的破碎.
所以我給罪下了一個現代性的定義.
罪就是自我認知障礙.
一個不相信福音的人.
存在著嚴重的自我認知障礙.
他不知道他幾斤幾兩.
他永遠看不清他幾斤幾兩.
當你看到神的偉大.
神的榮耀.
你就回過頭來.
能夠更多地認識自己.
羅馬有一個古詩人叫做奧維德.
他說過這樣一句話.
他說.
我們對那些使我們受傷的事物.
相信的比較慢.
是這樣嗎.
話也可以反過來講.
我們對那些對我們有利有好的事物.
通常相信的比較快.
你說你今天好漂亮.
你馬上就信了.
各位.
如果今天有人批評你的時候.
你不太容易接受.
如果今天講完道.
有人說張牧師.

$^{961}$你這邊的信息有待提升.
你覺得我會怎麼反應嗎.
你有做牧師的潛力了.
謝謝你提出來.
但你知道我心裡怎麼想嗎.
你講講事實.
真的會這樣.
各位.
求主幫助我們.
約翰加爾文說.
不認識上帝.
你就永遠不可能真實地認識自己.
而不認識自己.
你也不可能完全地認識上帝.
求主幫助我們.
讓我們在這.
這篇聖經當中看到人.
我們自己究竟是誰.
在哈佛大學.
這是夢寐以求的長春藤.
有一個哲學大樓叫做艾莫森大樓.
這個哲學系的大樓是1900年建的.
當時有人建議在這個大樓上印上一個名言.
其實是一個古希臘哲學家的.
叫做普羅達格拉斯.
他說man is the measure of all things.
人是萬物的尺度.
但是當時的哈佛校長叫做艾略特.
他堅持用詩篇的第八篇的第四節.
人算什麼.
已經顧念他.
堅持把這節聖經.
所以對他的評價很高.
有人說就是艾略特締造了哈佛.
讀到這裡的時候我就想.
哇 今天的哈佛校長都進到愛潑斯坦案裡去了.
當年的校長可見美國的墮落.
各位人究竟是什麼.
人文主義者說人是萬物的尺度.
人是宇宙的中心.

$^{1001}$但是在大衛.
對於一個神的兒女而言.
人究竟是什麼.
其實人論不斷地受到很多的挑戰.
我今天要分享一點.
達爾文說.
人類與其他動物是有共同祖先的.
說人是什麼.
高級動物.
高級靈長類動物.
弗洛伊德說.
人是受潛意識的驅動.
被深層的衝動所左右.
在慾望和壓抑之間拉扯的一個矛盾體.
你活得很矛盾.
卡爾馬克斯說.
人是由社會經濟結構所塑造.
本質上你是勞動生產關係當中的一個成員.
在整個的勞動階層當中.
你是工人.
你是農民.
你就是經濟.
基本上你是個經濟要素.
雍子媚到底人算什麼.
這裡說.
人算什麼.
你盡顧念他.
世人算什麼.
你盡眷顧他.
雍子媚這裡用了兩個人.
一個叫做人.
一個叫做世人.
其實在原文它是用了兩個字.
第一個叫做以諾世.
第二個叫做adam.
亞當.
這兩個字到底表達什麼.
亞當.
adam.
這個字其實在希伯來文當中.

$^{1041}$它的字根是另外一個.
叫做什麼.
adamah.
adamah和adam之間是用同一個字根.
這個adamah就是塵土.
就是土壤.
所以表明亞當的本質是塵土.
我不在這裡一一念了.
雍子媚她思念我們不過是塵土.
表明我們是軟弱的.
我們是脆弱的.
我們是卑賤的.
你不是黃金.
你不是鑽石.
你是塵土.
那以諾世呢.
我也把這些經文copy下來.
這四章二十六節.
賽特.
賽特是誰.
該隱殺了亞伯.
然後就生了一個代替亞伯.
叫做賽特.
然後賽特也生了一個兒子.
名叫以諾世.
很重要一節經文.
那時候人才求告耶和華的名.
以諾世這個字.
就是在神面前去求告神的名.
所以這兩個字加起來是什麼意思.
我承認我的軟弱卑微.
承認我的本質如同塵土.
所以我就要在你面前來求告你.
弟兄姊妹.
今天特別是從啟蒙運動直到今天.
人不斷地膨脹.
人不斷地想要人定勝天.
人不斷地.
弟兄姊妹你們能不能感受到.
人作為受造之物的一種狂妄.

$^{1081}$剛才我跟大家講的那個例子.
那個阿波羅八號.
讀完創世記第一章回到地球之後.
你知道被告了.
你知道他們把NASA告了.
從地區法院一直.
這個地區法院第五巡迴上訴法院.
一直告到最高法院.
告他們什麼.
你們用的納稅人的錢.
你怎麼能夠上去念聖經呢.
你知道最高法院怎麼判的嗎.
外太空.
不是我們的管轄範圍.
不受理.
各位啊.
人的狂妄.
人的自戀.
我覺得我們中國人有很多的詞很自戀.
香港回歸的時候.
97年.
有一個字叫做什麼.
每個人都知道.
普天同慶.
他把天給你拉下來了.
某個領導人去世了.
天地同悲.
天地為之事事.
多麼自作多情啊.
弟兄姊妹.
天的目的只有一個.
就是彰顯神的榮耀.
任何一件事情.
不會普天同慶的.
只有普天頌揚.
榮耀歸主.
任何一件事情.
不會天地同悲的.
只有耶穌在十字架上死亡的時候.
天都變成黑暗.

$^{1121}$天都要流露出悲傷.
弟兄姊妹.
不要那麼自作多情.
承認我們的卑微.
承認我們的渺小.
甚至這個詩篇當中.
用了一個形象.
叫做吃奶的嬰孩.
來表達我們的軟弱.
弟兄姊妹.
嬰孩.
吃奶的嬰孩.
表達了人.
或者是人類當中.
最脆弱.
最無助.
最沒有自衛能力的.
一個形象.
對吧.
剛出生的孩子.
請問他有知識嗎.
沒有.
他有聰明嗎.
沒有.
他有能力嗎.
沒有.
弟兄姊妹.
他不是一個獨立自主的人.
他每一天.
每一頓飯.
每一次的吃喝拉撒.
都需要有外來的幫助.
嬰孩.
我們不過像那吃奶的嬰孩.
我們需要被顧念.
我們需要被眷顧.
我好喜歡這個眷顧這個字.
眷.
你知道在中文當中.
這個字是一個家族用語.

$^{1161}$我們叫做家眷.
家眷對應的什麼.
家長.
對吧.
我們有一位天父.
他是我們的家長.
我們是他的家眷.
他供應我們.
他出去掙錢來給我們花.
他供應我們每天的日用的飲食.
他供應我們情感的需要.
永子妹.
我沒有時間多講了.
聖經上甚至表達他的供應.
他的care.
我們很難相信.
他說連你們的頭髮.
你有數過別人的頭髮嗎.
我想我小的時候.
我媽今天在這.
我媽那時候年輕.
她讓我給她拔白頭髮.
現在拔不完了.
那個時候.
只有在這種親密的關係當中.
除了這種場景.
我真的想不到.
數頭髮.
但是上帝說他給我們數頭髮.
再回到剛才那個picture.
那個video.
那麼偉大的一位上帝來數頭髮.
求主給我們信心.
所以在他的眷顧當中.
一方面我們承認我們的卑微和渺小.
另一方面在他的眷顧當中.
我們看到我們的尊貴與榮耀.
這裡有一些經文翻譯上.
甚至有一點不敢翻譯.
第五節說.

$^{1201}$你叫他比天使微小一點.
口號叫做神.
其實那字的原文就是Elohim.
就是神.
但是好像不太敢講.
不太敢講.
但是原文的的確確是講到.
我們比神這位造物主微小一點.
然後神冠以榮耀,尊貴.
Glory and Honor.
第五節.
我們盡賜於神.
然後呢.
這些聖經說.
讓我們觀禮他手所造的是萬物.
一切的牛羊.
田裡的獸.
空中的鳥.
海裡的魚.
金星海盜的.
都就是萬物了.
萬物在我們以下.
上帝在我們以上.
弟兄姊妹.
在進化論當中.
人是什麼.
人是高級動物.
人是最聰明的靈長類動物.
人是食物鏈的頂端.
但是你本身還是一個動物的範疇.
你和你家的狗是同類的.
1988年有一部片子叫做和尚.
雖然不是基督徒拍的片子.
但是非常鼓勵大家看.
不講這個內容了.
和尚有六個記稿人.
其中有三個都信了耶穌.
其中有一個叫做謝選俊.
他是在參加唐崇榮牧師查經班的時候信主的.
他那個時候他就問問題.

$^{1241}$他問.
他說.
假如今天有一個最笨的.
IQ最低的.
天生有認知缺陷的.
有巨大的智力障礙的.
白痴一樣的人.
有一個最聰明的.
最有本事最厲害的猿猴相比.
他說有什麼區別.
有什麼區別.
弟兄姊妹.
其實這個問題真正的挑戰在什麼呢.
人的尊貴可不可以被量化.
人的尊貴.
人的價值可不可以被量化.
如果可以.
弱者就必被邊緣化.
如果可以.
人間就一定有三六九等.
弟兄姊妹.
讓我這樣說.
如果沒有基督教信仰.
你說來到美國.
美國是人權.
無論是美國文化還是中國文化.
永遠不可能.
有一切人與一切人的平等的尊嚴和價值的觀念.
絕不可能有.
中國人說偏見.
偏字怎麼寫.
一個人字旁一個扁.
我們都會把人看扁.
真的.
我就記得911的時候.
雙子座那個大樓崩塌的時候.
你知道嗎.
保險公司要巨額的去賠償.
但是那個公司裡邊的高管.
有企業的高管.

$^{1281}$社會的精英.
也有在後廚打工的.
其實這個非法入境的人.
沒有身份的人.
你知道他們的賠償的金額是天差地別的.
為什麼.
人家有價值.
你呢.
你沒太有價值.
弟兄姊妹.
我們的價值.
我們的尊貴.
到底是什麼來量化.
什麼來定義的.
我們的本質高於一切的動物.
為什麼.
因為神是按照他的形象創造了我們.
任何個體的缺陷都不能抹殺人的尊貴.
各位.
你知道你是多麼尊貴嗎.
有一節經文.
很harsh.
我覺得勝經有的時候就是很harsh.
我們讀一下好不好.
一二.
人在尊貴中而不省悟.
就如死亡的畜類一樣.
你個畜生.
畜生不如.
不如嘛.
對吧.
為什麼.
你沒有意識到你的尊貴嗎.
什麼叫尊貴.
你沒有看到人身上上帝的形象嗎.
弟兄姊妹.
如果你看不到這個上帝的形象.
一個目中無人的人.
他一定心中無神.
一個心中無神的人.

$^{1321}$他早晚把你給看扁了.
他一定目中無人.
聖經說.
起初神創造了天地萬物.
整個創造的最高峰.
就是神按照他自己的形象樣式所造的人.
是這個世界這個宇宙最尊貴的.
最複雜的生物人.
之所以尊貴.
是因為聖經告訴我們.
上帝把他的形象.
把靈魂放在我們的裡邊.
天地萬物都不可以與我們相比.
只有人有情感.
有愛.
渴望擁抱.
渴望友誼.
渴望關心.
有無私的奉獻和付出.
只有人可以創作和欣賞巴哈的音樂.
偉大的康塔塔.
偉大的歌劇小說和繪畫.
只有人可以發展和品嘗.
描述自己的喜樂憂傷.
然而也只有人.
會背叛自己的尊貴.
出賣造他者的形象.
就是神的形象.
求主幫助我們.
不要受到這個世俗的觀念去影響.
再次來到神面前.
承認說.
我們有身上帶著神的形象.
這些聖經最令人震驚的是.
就是上帝創造了僅次於他的人.
比他微小一點的人.
各位我們的人權.
我們的自由.
我們的尊貴和我們人生當中.
一切有價值的來源是什麼呢.

$^{1361}$就是你與上帝相似.
各位今天一隻老虎撲向你.
給我一個理由.
為什麼我要把那個老虎打死保護你.
而不是把你打死餵那隻老虎.
為什麼.
世界上有80億人口了.
老虎瀕臨滅絕了.
為什麼.
你知道嗎.
按照科學家DNA說.
我們96\%的DNA和老鼠.
不是老虎.
和老鼠相似.
各位啊.
如果你不相信神的話.
你和老鼠和老虎和狗沒什麼兩樣.
我99\%的情況下.
我覺得我都應該打死你這個人.
然後去餵那隻老虎.
保護物種多樣性.
其實我還想講一個點.
我時間可能不夠.
我很快.
這兩三個周前.
三個周前大約.
馬斯克在一個播客節目當中.
講了一個AI的預言.
對吧.
我覺得今天AI的發展.
已經衝擊到了人的定義.
他說了很多事.
我真的沒時間講了.
他基本上講.
基本上他的意思就是說.
將來人類所有能幹的事.
AI全能幹.
那人幹嗎.
人的意義是什麼.
人活著的價值是什麼.

$^{1401}$弟兄姊妹.
人的尊貴到底是什麼.
如果他在智力上.
他在知識儲備上.
都超過了人類.
那弟兄姊妹.
在計算能力之外.
我們還有什麼.
我們面對這樣的一個.
人工智能發展.
技術進步的.
這樣的一個社會環境.
弟兄姊妹.
我們忽然發現.
我們比一切.
我們所造的.
更尊貴的.
不是我們的計算能力.
因為你的計算能力.
連個AI連個計算機.
都比不上.
各位.
那是什麼呢.
再說.
是我們的靈魂.
是上帝的形象.
當機器人能做人類.
大部分都能做的事.
各位.
教育意味著什麼.
你把他送到長春藤.
意味著什麼.
你總不能.
指望你的孩子.
將來超過AI吧.
超過機器人吧.
那你只能把他.
當機器人來養了.
各位不是的.
讓我這樣說.

$^{1441}$讓我對.
我們中間這些.
做父母的說.
我們最重要的.
就是機器人.
無法學習.
無法做到的.
弟兄姊妹.
就是我們的信仰.
我們的尊嚴.
靈魂的自由.
神聖的愛.
上帝將他的形象.
將他的榮耀.
聖潔.
尊貴.
和他一切道德倫理.
的選擇的能力.
賜給了我們.
這就是我們之所以尊貴.
求主幫助我們.
讓我們.
求主恢復我們.
受造榮美的形象.
我給大家念這段文字.
這是尼爾頓.
這是一位清教徒的作家.
被譽為.
最後一位清教徒的作家.
在他的《施樂園》當中.
第四章對亞當和夏娃的.
一個描述.
亞當夏娃這兩位.
身形更為高貴者.
直立而高大.
如神般挺立.
披戴著與生俱來的尊榮.
在赤裸的威嚴中.
顯得彷彿萬物之主.
她們額頭坦然高舉.

$^{1481}$目光清明無懼.
流露出對自身尊嚴的意識.
對於至高造物主的順服.
一切受造之物.
都在她們腳下安然順服.
亞當為沉思而勇毅而造.
夏娃為柔美與溫柔的恩典而生.
她為主權與理性之相.
她為和諧與愛的光輝.
美吧.
這就是神的形象.
這就是神在我們生命當中要成就的.
所以成為基督徒.
弟兄姊妹.
讓聖靈繼續在我們生命當中做工.
將我們的心智改換一新.
讓我們每天脫下舊人.
穿上新人.
恢復有真理的人意和聖潔.
所以各位啊.
再次與神同工.
管理萬物.
有的時候管理萬物不容易.
我養狗.
你們有多少養狗的.
當你養狗的時候.
當你那狗狗先來家的時候.
無論是你adopt.
還是買了一隻小狗.
你做的最多一個字叫什麼.
sit.
他就不sit.
給你個treat.
sit.
各位啊.
連讓一隻狗坐我下都做不到.
我們.
但聖經上說我們要管理萬物.
我們和上帝一起管理萬物.
求主幫助我們.

$^{1521}$打開我們的眼睛.
讓我們看到.
我們是上帝最得意的作品.
今天在哪裡找一個.
神的形象和樣式的人呢.
來教會.
我們不是披了羊皮的狼.
我們是披著.
披戴基督生命的罪人.
Amen.
讓我們在教會當中回應主的呼召.
恢復主榮美的形象.
完成神治理的使命.
\newpage



\section{撒母耳記上 18:1-23:29}
\label{sec:031664}
\textbf{做個良朋益友}~~~~李思聰牧師
\newline
\newline
連結: \href{http://www.church.com.hk/acms/content.asp?site=cdc\&op=show\&type=product\&code=031664\&layout=sermon}{\texttt{031664}} ~~~~ 語音日期: 2026-02-15
\newline
\newline
\hyperref[sec:031663]{\small{< < < PREV SERMON < < <}}
~
\hyperref[sec:index_chronic]{\small{[返順時目]}}
~
\hyperref[sec:031665]{\small{> > > NEXT SERMON > > >}}
\newline
\newline
撒母耳記上 18:1-23:29
\newline
\begin{longtable}{cl}
\hline
\hline
章節 & 經文 (和合本修訂版)\\
\hline
18:1 & \begin{tabularx}{0.7\textwidth}{X} 大衛對掃羅說完了話，約拿單的心與大衛的心深相契合。約拿單愛大衛，如同愛自己的性命。 \end{tabularx} \\ \\ \relax
18:2 & \begin{tabularx}{0.7\textwidth}{X} 那日掃羅留住大衛，不讓他回父家。 \end{tabularx} \\ \\ \relax
18:3 & \begin{tabularx}{0.7\textwidth}{X} 約拿單愛大衛如同愛自己的性命，就與他立約。 \end{tabularx} \\ \\ \relax
18:4 & \begin{tabularx}{0.7\textwidth}{X} 約拿單從身上脫下外袍，給了大衛，又把戰衣、刀、弓、腰帶都給了他。 \end{tabularx} \\ \\ \relax
18:5 & \begin{tabularx}{0.7\textwidth}{X} 掃羅無論差遣大衛往何處去，他都做事精明。掃羅立他作軍隊的指揮官，眾百姓和掃羅的臣僕都看為美。 \end{tabularx} \\ \\ \relax
18:6 & \begin{tabularx}{0.7\textwidth}{X} 大衛打死了那非利士人，同眾人回來的時候，婦女們從以色列各城裡出來，歡歡喜喜，打鼓奏樂，唱歌跳舞，迎接掃羅王。 \end{tabularx} \\ \\ \relax
18:7 & \begin{tabularx}{0.7\textwidth}{X} 眾婦女歡樂唱和，說：「掃羅殺死千千，大衛殺死萬萬。」 \end{tabularx} \\ \\ \relax
18:8 & \begin{tabularx}{0.7\textwidth}{X} 掃羅非常憤怒，不喜歡這話。他說：「將萬萬歸給大衛，千千歸給我，只剩下王國沒有給他！」 \end{tabularx} \\ \\ \relax
18:9 & \begin{tabularx}{0.7\textwidth}{X} 從這日起，掃羅就敵視大衛。 \end{tabularx} \\ \\ \relax
18:10 & \begin{tabularx}{0.7\textwidth}{X} 次日，從神來的邪靈緊抓住掃羅，他就在家中胡言亂語。大衛照常彈琴，掃羅手裡拿著槍。 \end{tabularx} \\ \\ \relax
18:11 & \begin{tabularx}{0.7\textwidth}{X} 掃羅把槍一擲，心裡說：「我要將大衛刺透，釘在牆上。」大衛閃避了他兩次。 \end{tabularx} \\ \\ \relax
18:12 & \begin{tabularx}{0.7\textwidth}{X} 掃羅懼怕大衛，因為耶和華離開自己，與大衛同在。 \end{tabularx} \\ \\ \relax
18:13 & \begin{tabularx}{0.7\textwidth}{X} 所以掃羅叫大衛離開自己，立他為千夫長，他就領兵出入。 \end{tabularx} \\ \\ \relax
18:14 & \begin{tabularx}{0.7\textwidth}{X} 大衛所做的每一件事都精明，耶和華也與他同在。 \end{tabularx} \\ \\ \relax
18:15 & \begin{tabularx}{0.7\textwidth}{X} 掃羅見大衛做事精明，就更怕他。 \end{tabularx} \\ \\ \relax
18:16 & \begin{tabularx}{0.7\textwidth}{X} 但以色列和猶大眾人都愛大衛，因為他領他們出入。 \end{tabularx} \\ \\ \relax
18:17 & \begin{tabularx}{0.7\textwidth}{X} 掃羅對大衛說：「看哪，我將大女兒米拉嫁給你，只要你作我的勇士，為耶和華爭戰。」掃羅心裡說：「我不好親手害他，要藉非利士人的手害他。」 \end{tabularx} \\ \\ \relax
18:18 & \begin{tabularx}{0.7\textwidth}{X} 大衛對掃羅說：「我是誰，我是甚麼出身，我父家在以色列中算甚麼，豈敢作王的女婿呢？」 \end{tabularx} \\ \\ \relax
18:19 & \begin{tabularx}{0.7\textwidth}{X} 掃羅的女兒米拉到了當嫁給大衛的時候，掃羅卻將她嫁給了米何拉人亞得列。 \end{tabularx} \\ \\ \relax
18:20 & \begin{tabularx}{0.7\textwidth}{X} 掃羅的女兒米甲愛大衛。有人告訴掃羅，這件事在掃羅眼中看為合宜。 \end{tabularx} \\ \\ \relax
18:21 & \begin{tabularx}{0.7\textwidth}{X} 掃羅心裡說：「我把這女兒嫁給大衛，作他的圈套，好藉非利士人的手害他。」所以掃羅第二次對大衛說：「你今日可以作我的女婿。」 \end{tabularx} \\ \\ \relax
18:22 & \begin{tabularx}{0.7\textwidth}{X} 掃羅吩咐臣僕：「你們暗中對大衛說：『看哪，王喜歡你，王的臣僕也都愛戴你，現在你就作王的女婿吧。』」 \end{tabularx} \\ \\ \relax
18:23 & \begin{tabularx}{0.7\textwidth}{X} 掃羅的臣僕照這話說給大衛聽。大衛說：「你們把作王的女婿看為小事嗎？我是貧窮卑微的人。」 \end{tabularx} \\ \\ \relax
18:24 & \begin{tabularx}{0.7\textwidth}{X} 掃羅的臣僕回奏說，大衛說了這樣的話。 \end{tabularx} \\ \\ \relax
18:25 & \begin{tabularx}{0.7\textwidth}{X} 掃羅說：「你們要對大衛這樣說：『王不要甚麼聘禮，只要一百非利士人的包皮，好在王的仇敵身上報仇。』」掃羅的意圖是要大衛落在非利士人的手中。 \end{tabularx} \\ \\ \relax
18:26 & \begin{tabularx}{0.7\textwidth}{X} 掃羅的臣僕把這話告訴大衛，大衛就歡喜作王的女婿。日期還沒有到， \end{tabularx} \\ \\ \relax
18:27 & \begin{tabularx}{0.7\textwidth}{X} 大衛和跟隨他的人起來前往，殺了二百非利士人，將包皮足數交給王，為要作王的女婿。於是掃羅將女兒米甲嫁給大衛。 \end{tabularx} \\ \\ \relax
18:28 & \begin{tabularx}{0.7\textwidth}{X} 掃羅見耶和華與大衛同在，女兒米甲又愛大衛， \end{tabularx} \\ \\ \relax
18:29 & \begin{tabularx}{0.7\textwidth}{X} 就更怕大衛，常常與大衛為敵。 \end{tabularx} \\ \\ \relax
18:30 & \begin{tabularx}{0.7\textwidth}{X} 每逢非利士的軍官出來打仗，大衛做事比掃羅任何臣僕更精明，因此他的名極受尊重。 \end{tabularx} \\ \\ \relax
19:1 & \begin{tabularx}{0.7\textwidth}{X} 掃羅吩咐他兒子約拿單和眾臣僕要殺大衛，但掃羅的兒子約拿單卻很喜愛大衛。 \end{tabularx} \\ \\ \relax
19:2 & \begin{tabularx}{0.7\textwidth}{X} 約拿單告訴大衛說：「我父掃羅想要殺你，現在你要小心，明日早晨留在一個僻靜的地方藏起來。 \end{tabularx} \\ \\ \relax
19:3 & \begin{tabularx}{0.7\textwidth}{X} 我會出去，到你所藏的田裡，站在我父親旁邊，與父親談論到你。我看情形怎樣，會告訴你。」 \end{tabularx} \\ \\ \relax
19:4 & \begin{tabularx}{0.7\textwidth}{X} 約拿單向他父親掃羅說大衛的好話，對他說：「王不可得罪王的僕人大衛，因為他未曾得罪你，他所行的對你都很有益處。 \end{tabularx} \\ \\ \relax
19:5 & \begin{tabularx}{0.7\textwidth}{X} 他拚了命殺那非利士人，並且耶和華為全以色列大施拯救。那時你看見，也很歡喜，現在為何要犯罪，流無辜人的血，無緣無故殺大衛呢？」 \end{tabularx} \\ \\ \relax
19:6 & \begin{tabularx}{0.7\textwidth}{X} 掃羅聽了約拿單的話，就指著永生的耶和華起誓：「我絕不殺他。」 \end{tabularx} \\ \\ \relax
19:7 & \begin{tabularx}{0.7\textwidth}{X} 約拿單叫大衛來，把這一切事告訴他。約拿單帶他去見掃羅，他就像以前一樣侍立在掃羅面前。 \end{tabularx} \\ \\ \relax
19:8 & \begin{tabularx}{0.7\textwidth}{X} 此後又有戰爭，大衛出去與非利士人打仗。他大大擊敗他們，他們就在他面前逃跑。 \end{tabularx} \\ \\ \relax
19:9 & \begin{tabularx}{0.7\textwidth}{X} 從耶和華來的邪靈又降在掃羅身上，掃羅手裡拿槍坐在屋裡，大衛正用手彈琴。 \end{tabularx} \\ \\ \relax
19:10 & \begin{tabularx}{0.7\textwidth}{X} 掃羅想要用槍刺透大衛，把他釘在牆上，他卻躲開掃羅，掃羅的槍刺入牆內。當夜大衛逃走，躲起來了。 \end{tabularx} \\ \\ \relax
19:11 & \begin{tabularx}{0.7\textwidth}{X} 掃羅派一些使者到大衛的房屋那裡守著他，等到天亮要殺他。大衛的妻子米甲對大衛說：「你今夜若不逃命，明日就要被殺。」 \end{tabularx} \\ \\ \relax
19:12 & \begin{tabularx}{0.7\textwidth}{X} 於是米甲將大衛從窗戶縋下去，讓他走；大衛就逃走，躲起來了。 \end{tabularx} \\ \\ \relax
19:13 & \begin{tabularx}{0.7\textwidth}{X} 米甲把家中的神像放在床上，頭枕在山羊毛的枕頭上，用衣服蓋起來。 \end{tabularx} \\ \\ \relax
19:14 & \begin{tabularx}{0.7\textwidth}{X} 掃羅派一些使者去捉拿大衛，米甲說：「他病了。」 \end{tabularx} \\ \\ \relax
19:15 & \begin{tabularx}{0.7\textwidth}{X} 掃羅又派一些使者去看大衛，說：「把他連床一起抬到我這裡，我好殺他。」 \end{tabularx} \\ \\ \relax
19:16 & \begin{tabularx}{0.7\textwidth}{X} 使者進去，看哪，神像在床上，頭枕在山羊毛的枕頭上。 \end{tabularx} \\ \\ \relax
19:17 & \begin{tabularx}{0.7\textwidth}{X} 掃羅對米甲說：「你為甚麼這樣欺騙我，放我仇敵逃走呢？」米甲對掃羅說：「他對我說：『你放我走吧，我何必要殺你呢？』」 \end{tabularx} \\ \\ \relax
19:18 & \begin{tabularx}{0.7\textwidth}{X} 大衛逃跑躲避，來到拉瑪的撒母耳那裡，把掃羅向他所行的事全告訴他。他和撒母耳就去，住在拿約。 \end{tabularx} \\ \\ \relax
19:19 & \begin{tabularx}{0.7\textwidth}{X} 有人告訴掃羅說：「看哪，大衛在拉瑪的拿約」。 \end{tabularx} \\ \\ \relax
19:20 & \begin{tabularx}{0.7\textwidth}{X} 掃羅派一些使者去捉拿大衛。去的人見一隊先知受感說話，撒母耳站在當中領導他們，掃羅派去的使者也受神的靈感動說話。 \end{tabularx} \\ \\ \relax
19:21 & \begin{tabularx}{0.7\textwidth}{X} 有人把這事告訴掃羅，他又派另一些使者去，他們也受感說話。掃羅第三次派使者去，他們也受感說話。 \end{tabularx} \\ \\ \relax
19:22 & \begin{tabularx}{0.7\textwidth}{X} 然後掃羅親自往拉瑪去，到了西沽的大井，問人說：「撒母耳和大衛在哪裡？」有人說：「看哪，在拉瑪的拿約。」 \end{tabularx} \\ \\ \relax
19:23 & \begin{tabularx}{0.7\textwidth}{X} 他就往那裡去，到了拉瑪的拿約。神的靈也臨到他，他一面走一面受感說話，直到拉瑪的拿約。 \end{tabularx} \\ \\ \relax
19:24 & \begin{tabularx}{0.7\textwidth}{X} 他也脫了衣服，也在撒母耳面前受感說話，一日一夜赤身躺臥。因此有人說：「掃羅也在先知中嗎？」 \end{tabularx} \\ \\ \relax
20:1 & \begin{tabularx}{0.7\textwidth}{X} 大衛從拉瑪的拿約逃跑，來到約拿單面前，對他說：「我做了甚麼，有甚麼罪孽，在你父親面前犯了甚麼罪，他竟要尋索我的性命呢？」 \end{tabularx} \\ \\ \relax
20:2 & \begin{tabularx}{0.7\textwidth}{X} 約拿單對他說：「絕無此事！你必不至於死。看哪，我父做事，無論大小，沒有不告訴我的。我父親為甚麼要隱瞞我這件事呢？不會這樣的！」 \end{tabularx} \\ \\ \relax
20:3 & \begin{tabularx}{0.7\textwidth}{X} 大衛又起誓說：「你父親確實知道我在你眼前蒙恩。所以他說，『這事不要讓約拿單知道，免得他愁煩。』我指著永生的耶和華起誓，又指著你的性命起誓，我離死只差一步而已。」 \end{tabularx} \\ \\ \relax
20:4 & \begin{tabularx}{0.7\textwidth}{X} 約拿單對大衛說：「你心裡所求的，我必為你成就。」 \end{tabularx} \\ \\ \relax
20:5 & \begin{tabularx}{0.7\textwidth}{X} 大衛對約拿單說：「看哪，明日是初一，我必須與王同席用餐，求你讓我去藏在田野，直到第三日傍晚。 \end{tabularx} \\ \\ \relax
20:6 & \begin{tabularx}{0.7\textwidth}{X} 你父親若見我不在席上，你就說：『大衛懇求我允許他趕回本城伯利恆去，因為他全家在那裡獻年祭。』 \end{tabularx} \\ \\ \relax
20:7 & \begin{tabularx}{0.7\textwidth}{X} 你父親若說好，你的僕人就平安了；他若大怒，你就知道他決意行惡。 \end{tabularx} \\ \\ \relax
20:8 & \begin{tabularx}{0.7\textwidth}{X} 求你施恩於僕人，因你在耶和華面前曾與僕人立約。我若有罪孽，你就親自殺死我，何必把我交給你父親呢？」 \end{tabularx} \\ \\ \relax
20:9 & \begin{tabularx}{0.7\textwidth}{X} 約拿單說：「絕無此事！我若確實知道我父親決意害你，怎麼會不告訴你呢？」 \end{tabularx} \\ \\ \relax
20:10 & \begin{tabularx}{0.7\textwidth}{X} 大衛對約拿單說：「你父親若嚴厲回答你，誰來告訴我呢？」 \end{tabularx} \\ \\ \relax
20:11 & \begin{tabularx}{0.7\textwidth}{X} 約拿單對大衛說：「來，讓我們到田野去。」二人就往田野去了。 \end{tabularx} \\ \\ \relax
20:12 & \begin{tabularx}{0.7\textwidth}{X} 約拿單對大衛說：「願耶和華－以色列的神作證。明日約在這時候，或第三日，我一探出我父親的心意，看哪，若對大衛是好意，我怎麼會不派人來告訴你呢？ \end{tabularx} \\ \\ \relax
20:13 & \begin{tabularx}{0.7\textwidth}{X} 我父親若有意害你，而我不告訴你，送你平安地離開，願耶和華重重懲罰約拿單。願耶和華與你同在，如同從前與我父親同在一樣。 \end{tabularx} \\ \\ \relax
20:14 & \begin{tabularx}{0.7\textwidth}{X} 你要照耶和華的慈愛恩待我，不但我活著的時候免我死亡， \end{tabularx} \\ \\ \relax
20:15 & \begin{tabularx}{0.7\textwidth}{X} 就是耶和華從地面逐一剪除大衛仇敵的時候，你也永不可向我家斷絕恩惠。」 \end{tabularx} \\ \\ \relax
20:16 & \begin{tabularx}{0.7\textwidth}{X} 於是約拿單與大衛家立約：「願耶和華從大衛仇敵的手來追討。」 \end{tabularx} \\ \\ \relax
20:17 & \begin{tabularx}{0.7\textwidth}{X} 約拿單因愛大衛如同愛自己的性命，就叫他再起誓。 \end{tabularx} \\ \\ \relax
20:18 & \begin{tabularx}{0.7\textwidth}{X} 約拿單對他說：「明日是初一，你的座位空著，人必察覺你不在。 \end{tabularx} \\ \\ \relax
20:19 & \begin{tabularx}{0.7\textwidth}{X} 到第三日，就要走一段長路下去，去到你遇事那天所藏的地方，在以色磐石的旁邊等候。 \end{tabularx} \\ \\ \relax
20:20 & \begin{tabularx}{0.7\textwidth}{X} 我會向磐石旁邊射三箭，如同射箭靶一樣。 \end{tabularx} \\ \\ \relax
20:21 & \begin{tabularx}{0.7\textwidth}{X} 看哪，我會派僮僕，說：『去把箭找來。』我若對僮僕喊說：『看哪，箭在你的這邊，把箭拿來』，你就可以平安回來；我指著永生的耶和華起誓，你一定沒有事。 \end{tabularx} \\ \\ \relax
20:22 & \begin{tabularx}{0.7\textwidth}{X} 我若對孩子說：『看哪，箭在你的前方』，你就要離開，因為是耶和華差你去的。 \end{tabularx} \\ \\ \relax
20:23 & \begin{tabularx}{0.7\textwidth}{X} 至於你和我，我們所說的話，看哪，耶和華在你我中間作證，直到永遠。」 \end{tabularx} \\ \\ \relax
20:24 & \begin{tabularx}{0.7\textwidth}{X} 大衛就去藏在田野。到了初一，王要坐席用餐。 \end{tabularx} \\ \\ \relax
20:25 & \begin{tabularx}{0.7\textwidth}{X} 王照常坐在靠牆的位子上，約拿單在對面，押尼珥坐在掃羅旁邊，大衛的座位卻是空的。 \end{tabularx} \\ \\ \relax
20:26 & \begin{tabularx}{0.7\textwidth}{X} 這日掃羅沒有說甚麼，因為他說：「大衛或許有事，偶染不潔，還未得潔淨。」 \end{tabularx} \\ \\ \relax
20:27 & \begin{tabularx}{0.7\textwidth}{X} 初二，大衛的座位還空著。掃羅對他兒子約拿單說：「耶西的兒子為何昨日、今日都沒有來用餐呢？」 \end{tabularx} \\ \\ \relax
20:28 & \begin{tabularx}{0.7\textwidth}{X} 約拿單回答掃羅說：「大衛懇求我允許他回伯利恆去， \end{tabularx} \\ \\ \relax
20:29 & \begin{tabularx}{0.7\textwidth}{X} 說：『求你讓我去，因為我家在城裡有獻祭的事，我哥哥吩咐我去。如今我若在你眼前蒙恩，求你讓我去見我的兄弟。』所以大衛沒有來赴王的筵席。」 \end{tabularx} \\ \\ \relax
20:30 & \begin{tabularx}{0.7\textwidth}{X} 掃羅向約拿單怒氣大發，對他說：「你這頑梗悖逆之婦人所生的，我怎麼會不知道你選擇耶西的兒子，自取羞辱，也使你母親露體蒙羞呢？ \end{tabularx} \\ \\ \relax
20:31 & \begin{tabularx}{0.7\textwidth}{X} 只要耶西的兒子還活在世上一天，你和你的國必保不住。現在你要派人去，把他帶到我這裡來，因為他是該死的。」 \end{tabularx} \\ \\ \relax
20:32 & \begin{tabularx}{0.7\textwidth}{X} 約拿單回答父親掃羅說：「他為甚麼該死呢？他做了甚麼呢？」 \end{tabularx} \\ \\ \relax
20:33 & \begin{tabularx}{0.7\textwidth}{X} 掃羅向約拿單擲槍要刺他，約拿單就知道他父親決意要殺死大衛。 \end{tabularx} \\ \\ \relax
20:34 & \begin{tabularx}{0.7\textwidth}{X} 於是約拿單氣憤憤地從席上起來。他在初二這天沒有吃飯，因為他為大衛愁煩，又因為他父親羞辱了他。 \end{tabularx} \\ \\ \relax
20:35 & \begin{tabularx}{0.7\textwidth}{X} 次日早晨，約拿單按著與大衛約定的時候到田野去，有一個小僮僕跟隨他。 \end{tabularx} \\ \\ \relax
20:36 & \begin{tabularx}{0.7\textwidth}{X} 約拿單對僮僕說：「你跑去把我所射的箭找來。」僮僕跑去，約拿單就把箭射在僮僕的前方。 \end{tabularx} \\ \\ \relax
20:37 & \begin{tabularx}{0.7\textwidth}{X} 僮僕到了約拿單落箭之地，約拿單呼叫僮僕說：「箭不是在你的前方嗎？」 \end{tabularx} \\ \\ \relax
20:38 & \begin{tabularx}{0.7\textwidth}{X} 約拿單又呼叫僮僕說：「快去，不要站在那裡！」僮僕就撿起箭來，回到主人那裡。 \end{tabularx} \\ \\ \relax
20:39 & \begin{tabularx}{0.7\textwidth}{X} 僮僕不知道這是甚麼意思，只有約拿單和大衛知道這事。 \end{tabularx} \\ \\ \relax
20:40 & \begin{tabularx}{0.7\textwidth}{X} 約拿單把他的弓箭交給僮僕，吩咐他說：「你拿到城裡去。」 \end{tabularx} \\ \\ \relax
20:41 & \begin{tabularx}{0.7\textwidth}{X} 僮僕一去，大衛就從南邊出來，俯伏在地，拜了三拜。他們彼此親吻，一起哭泣，大衛哭得更悲哀。 \end{tabularx} \\ \\ \relax
20:42 & \begin{tabularx}{0.7\textwidth}{X} 約拿單對大衛說：「你平平安安地去吧！因為我們二人曾指著耶和華的名起誓說：『願耶和華在你我中間，以及你我後裔中間作證，直到永遠。』」大衛就起身走了，約拿單也回城裡去了。 \end{tabularx} \\ \\ \relax
21:1 & \begin{tabularx}{0.7\textwidth}{X} 大衛到了挪伯的亞希米勒祭司那裡，亞希米勒戰戰兢兢地出來迎接他，對他說：「你為甚麼獨自一人，沒有人跟隨你呢？」 \end{tabularx} \\ \\ \relax
21:2 & \begin{tabularx}{0.7\textwidth}{X} 大衛對亞希米勒祭司說：「王吩咐我一件事，對我說：『我差遣你，吩咐你的這件事，不可讓任何人知道。』因此我已告訴一些僕人到某處去。 \end{tabularx} \\ \\ \relax
21:3 & \begin{tabularx}{0.7\textwidth}{X} 現在你手中有甚麼？請你給我五個餅或是可以找到的食物。」 \end{tabularx} \\ \\ \relax
21:4 & \begin{tabularx}{0.7\textwidth}{X} 祭司對大衛說：「我手中沒有普通的餅，只有聖餅，只能給沒有親近婦人的年輕人。」 \end{tabularx} \\ \\ \relax
21:5 & \begin{tabularx}{0.7\textwidth}{X} 大衛回答祭司說：「我們確實沒有親近婦人，如同往常我出征的時候一樣。平常行路，僕人的身體都還分別為聖，何況今日豈不更使自己分別為聖嗎？」 \end{tabularx} \\ \\ \relax
21:6 & \begin{tabularx}{0.7\textwidth}{X} 祭司拿聖餅給他，因為在那裡沒有別的餅，只有那從耶和華面前撤下的供餅，就是換上熱餅的日子取下來的。 \end{tabularx} \\ \\ \relax
21:7 & \begin{tabularx}{0.7\textwidth}{X} 當日，有掃羅的一個臣僕在那裡，他留在耶和華的面前，名叫多益，是以東人，作掃羅的畜牧長。 \end{tabularx} \\ \\ \relax
21:8 & \begin{tabularx}{0.7\textwidth}{X} 大衛對亞希米勒說：「你手中有沒有槍或刀？因為王的事緊急，連刀劍兵器我都沒有帶。」 \end{tabularx} \\ \\ \relax
21:9 & \begin{tabularx}{0.7\textwidth}{X} 祭司說：「你在以拉谷所殺的非利士人歌利亞的那刀，看哪，裹在布中，放在以弗得後邊。你若要可以拿去，除此以外，再沒有別的了。」大衛說：「沒有甚麼可以跟它比的了！請你給我。」 \end{tabularx} \\ \\ \relax
21:10 & \begin{tabularx}{0.7\textwidth}{X} 那日大衛起來，躲避掃羅，逃到迦特王亞吉那裡。 \end{tabularx} \\ \\ \relax
21:11 & \begin{tabularx}{0.7\textwidth}{X} 亞吉的臣僕對他說：「這不是那地的國王大衛嗎？那裡的人跳舞唱和：『掃羅殺死千千，大衛殺死萬萬』，不是指著他說的嗎？」 \end{tabularx} \\ \\ \relax
21:12 & \begin{tabularx}{0.7\textwidth}{X} 大衛把這些話放在心裡，就很懼怕迦特王亞吉。 \end{tabularx} \\ \\ \relax
21:13 & \begin{tabularx}{0.7\textwidth}{X} 於是他在眾人眼前一反常態，在他們中間裝瘋作癲，在城門的門扇上胡寫亂畫，任由唾沫流在鬍子上。 \end{tabularx} \\ \\ \relax
21:14 & \begin{tabularx}{0.7\textwidth}{X} 亞吉對臣僕說：「看哪，你們看這人瘋了，為甚麼帶他到我這裡來呢？ \end{tabularx} \\ \\ \relax
21:15 & \begin{tabularx}{0.7\textwidth}{X} 我豈缺少瘋子，你們竟然帶這人到我面前瘋癲嗎？這個人可以進我的家嗎？」 \end{tabularx} \\ \\ \relax
22:1 & \begin{tabularx}{0.7\textwidth}{X} 大衛離開那裡，逃到亞杜蘭洞。他的兄弟和他父親全家聽見了，都下到他那裡去。 \end{tabularx} \\ \\ \relax
22:2 & \begin{tabularx}{0.7\textwidth}{X} 凡生活窘迫的、欠債的、心裡苦惱的都聚集到大衛那裡，他就作他們的領袖，跟隨他的約有四百人。 \end{tabularx} \\ \\ \relax
22:3 & \begin{tabularx}{0.7\textwidth}{X} 大衛從那裡往摩押的米斯巴去，對摩押王說：「請你讓我父母搬來，跟你們在一起，等我知道神要為我怎樣做。」 \end{tabularx} \\ \\ \relax
22:4 & \begin{tabularx}{0.7\textwidth}{X} 大衛領他父母到摩押王面前。大衛住山寨一切的日子，他父母也都住在摩押王那裡。 \end{tabularx} \\ \\ \relax
22:5 & \begin{tabularx}{0.7\textwidth}{X} 先知迦得對大衛說：「你不要住在山寨，要到猶大地去。」大衛就去，來到哈列的樹林。 \end{tabularx} \\ \\ \relax
22:6 & \begin{tabularx}{0.7\textwidth}{X} 掃羅聽見大衛和跟隨他之人的下落，掃羅正在基比亞，坐在山頂的柳樹下，手裡拿著槍，眾臣僕侍立在左右。 \end{tabularx} \\ \\ \relax
22:7 & \begin{tabularx}{0.7\textwidth}{X} 掃羅對左右侍立的臣僕說：「便雅憫人哪，聽著！耶西的兒子也能把田地和葡萄園賜給你們各人嗎？他能立你們各人作千夫長和百夫長嗎？ \end{tabularx} \\ \\ \relax
22:8 & \begin{tabularx}{0.7\textwidth}{X} 你們竟都結黨害我！我兒子與耶西的兒子立約的時候，無人告訴我；我兒子挑唆我的臣僕謀害我，像今日這樣，也無人告訴我，為我憂慮。」 \end{tabularx} \\ \\ \relax
22:9 & \begin{tabularx}{0.7\textwidth}{X} 那時以東人多益站在掃羅的臣僕中，回答說：「我曾看見耶西的兒子往挪伯去，到了亞希突的兒子亞希米勒那裡。 \end{tabularx} \\ \\ \relax
22:10 & \begin{tabularx}{0.7\textwidth}{X} 亞希米勒為他求問耶和華，給他食物，又把非利士人歌利亞的刀給了他。」 \end{tabularx} \\ \\ \relax
22:11 & \begin{tabularx}{0.7\textwidth}{X} 王就派人把亞希突的兒子亞希米勒祭司和他父親的全家，就是在挪伯的祭司都召了來，他們都來到王那裡。 \end{tabularx} \\ \\ \relax
22:12 & \begin{tabularx}{0.7\textwidth}{X} 掃羅說：「亞希突的兒子，聽著！」他說：「我主，我在這裡。」 \end{tabularx} \\ \\ \relax
22:13 & \begin{tabularx}{0.7\textwidth}{X} 掃羅對他說：「你為甚麼與耶西的兒子結黨害我，把食物和刀給他，又為他求問神，使他起來謀害我，像今日這樣？」 \end{tabularx} \\ \\ \relax
22:14 & \begin{tabularx}{0.7\textwidth}{X} 亞希米勒回答王說：「王的眾臣僕中有誰比大衛忠心呢？他是王的女婿，又是你的侍衛長，並且是你宮中受敬重的人。 \end{tabularx} \\ \\ \relax
22:15 & \begin{tabularx}{0.7\textwidth}{X} 我今日才開始為他求問神嗎？絕非如此！王不要歸罪於我和我父全家，因為這事，無論大小，僕人都不知情。」 \end{tabularx} \\ \\ \relax
22:16 & \begin{tabularx}{0.7\textwidth}{X} 王說：「亞希米勒，你和你父全家都是該死的！」 \end{tabularx} \\ \\ \relax
22:17 & \begin{tabularx}{0.7\textwidth}{X} 王吩咐左右的侍衛說：「你們轉身去殺耶和華的祭司吧！因為他們幫助大衛，知道大衛逃跑卻不告訴我。」但王的臣僕都不願動手殺耶和華的祭司。 \end{tabularx} \\ \\ \relax
22:18 & \begin{tabularx}{0.7\textwidth}{X} 王吩咐多益說：「你轉身去殺祭司吧！」以東人多益就轉身去殺祭司，那日殺了穿細麻布以弗得的，共八十五人， \end{tabularx} \\ \\ \relax
22:19 & \begin{tabularx}{0.7\textwidth}{X} 又用刀把祭司城挪伯中的男女、孩童和吃奶的都殺了，又用刀殺了牛、羊和驢子。 \end{tabularx} \\ \\ \relax
22:20 & \begin{tabularx}{0.7\textwidth}{X} 亞希突的兒子亞希米勒有一個兒子逃脫了；他名叫亞比亞他，逃到大衛那裡。 \end{tabularx} \\ \\ \relax
22:21 & \begin{tabularx}{0.7\textwidth}{X} 亞比亞他把掃羅殺耶和華祭司的事告訴大衛。 \end{tabularx} \\ \\ \relax
22:22 & \begin{tabularx}{0.7\textwidth}{X} 大衛對亞比亞他說：「那日我見以東人多益在那裡，就知道他一定會告訴掃羅。你父的全家喪命，都是因我的緣故。 \end{tabularx} \\ \\ \relax
22:23 & \begin{tabularx}{0.7\textwidth}{X} 你可以住在我這裡，不要懼怕。因為尋索你命的也要尋索我的命，你在我這裡可得保護。」 \end{tabularx} \\ \\ \relax
23:1 & \begin{tabularx}{0.7\textwidth}{X} 有人告訴大衛說：「看哪，非利士人攻擊基伊拉，搶奪禾場。」 \end{tabularx} \\ \\ \relax
23:2 & \begin{tabularx}{0.7\textwidth}{X} 大衛求問耶和華說：「我可以去嗎？我可以去攻打那些非利士人嗎？」耶和華對大衛說：「你可以去攻打非利士人，拯救基伊拉。」 \end{tabularx} \\ \\ \relax
23:3 & \begin{tabularx}{0.7\textwidth}{X} 大衛的人對他說：「看哪，我們在猶大這裡尚且懼怕，何況到基伊拉去攻打非利士人的軍隊呢？」 \end{tabularx} \\ \\ \relax
23:4 & \begin{tabularx}{0.7\textwidth}{X} 大衛又再求問耶和華，耶和華回答說：「你起身下基伊拉去，我必將非利士人交在你手裡。」 \end{tabularx} \\ \\ \relax
23:5 & \begin{tabularx}{0.7\textwidth}{X} 於是大衛和他的人往基伊拉去，與非利士人打仗，大大擊敗他們，奪取他們的牲畜。這樣，大衛救了基伊拉的居民。 \end{tabularx} \\ \\ \relax
23:6 & \begin{tabularx}{0.7\textwidth}{X} 亞希米勒的兒子亞比亞他逃往基伊拉到大衛那裡的時候，手裡拿著以弗得。 \end{tabularx} \\ \\ \relax
23:7 & \begin{tabularx}{0.7\textwidth}{X} 有人告訴掃羅，大衛到了基伊拉。掃羅說：「神將他交在我手裡了，因為他進了有門有閂的城，把自己關起來了。」 \end{tabularx} \\ \\ \relax
23:8 & \begin{tabularx}{0.7\textwidth}{X} 於是掃羅召集眾百姓，要下去攻打基伊拉，圍困大衛和他的人。 \end{tabularx} \\ \\ \relax
23:9 & \begin{tabularx}{0.7\textwidth}{X} 大衛知道掃羅設計陷害他，就對亞比亞他祭司說：「把以弗得拿過來。」 \end{tabularx} \\ \\ \relax
23:10 & \begin{tabularx}{0.7\textwidth}{X} 大衛說：「耶和華－以色列的神啊，你僕人確實聽見掃羅設法要到基伊拉來，為我的緣故毀滅這城。 \end{tabularx} \\ \\ \relax
23:11 & \begin{tabularx}{0.7\textwidth}{X} 基伊拉人會把我交在掃羅手裡嗎？掃羅會下來，正如你僕人所聽見的嗎？耶和華－以色列的神啊，求你指示僕人！」耶和華說：「他會下來。」 \end{tabularx} \\ \\ \relax
23:12 & \begin{tabularx}{0.7\textwidth}{X} 大衛又說：「基伊拉人會把我和我的人交在掃羅手裡嗎？」耶和華說：「他們會交出來。」 \end{tabularx} \\ \\ \relax
23:13 & \begin{tabularx}{0.7\textwidth}{X} 於是大衛和他的人約有六百名起身離開基伊拉，往他們所能去的地方去。有人告訴掃羅，大衛離開基伊拉逃走了，掃羅就停止出發了。 \end{tabularx} \\ \\ \relax
23:14 & \begin{tabularx}{0.7\textwidth}{X} 大衛住在曠野的山寨裡，在西弗曠野的山區。掃羅天天尋索大衛，神卻不將大衛交在他手裡。 \end{tabularx} \\ \\ \relax
23:15 & \begin{tabularx}{0.7\textwidth}{X} 大衛看到掃羅出來尋索他的命。那時，他住在西弗曠野的樹林裡； \end{tabularx} \\ \\ \relax
23:16 & \begin{tabularx}{0.7\textwidth}{X} 掃羅的兒子約拿單起身，到樹林裡去見大衛，使他的手倚靠神得以堅固， \end{tabularx} \\ \\ \relax
23:17 & \begin{tabularx}{0.7\textwidth}{X} 對他說：「不要懼怕！我父掃羅的手無法害你 ，你必作以色列的王，我必作你的宰相。我父掃羅也知道這事。」 \end{tabularx} \\ \\ \relax
23:18 & \begin{tabularx}{0.7\textwidth}{X} 於是二人在耶和華面前立約。大衛仍住在樹林裡，約拿單就回家去了。 \end{tabularx} \\ \\ \relax
23:19 & \begin{tabularx}{0.7\textwidth}{X} 西弗人上基比亞到掃羅那裡，說：「大衛不是在我們那裡，在樹林裡的山寨中，在荒野南邊的哈基拉山藏著嗎？ \end{tabularx} \\ \\ \relax
23:20 & \begin{tabularx}{0.7\textwidth}{X} 現在，王啊，請隨你的心願要下來，就請下來；至於我們，一定會把他交在王的手裡。」 \end{tabularx} \\ \\ \relax
23:21 & \begin{tabularx}{0.7\textwidth}{X} 掃羅說：「願耶和華賜福給你們，因為你們體恤我。 \end{tabularx} \\ \\ \relax
23:22 & \begin{tabularx}{0.7\textwidth}{X} 請你們回去，再確定一下，調查並看清楚他落腳的地方，是誰看見他在那裡，因為有人告訴我他很狡猾。 \end{tabularx} \\ \\ \relax
23:23 & \begin{tabularx}{0.7\textwidth}{X} 你們要看清楚，調查他藏匿的每一個地方，回來給我確實的報告，我就與你們同去。他若在境內，我必從猶大的千門萬戶中搜出他來。」 \end{tabularx} \\ \\ \relax
23:24 & \begin{tabularx}{0.7\textwidth}{X} 西弗人動身，在掃羅以先往西弗去。大衛和他的人卻在瑪雲曠野，在荒野南邊的亞拉巴。 \end{tabularx} \\ \\ \relax
23:25 & \begin{tabularx}{0.7\textwidth}{X} 掃羅和他的人去尋索大衛。有人告訴大衛，他就下到巖石那裡，留在瑪雲的曠野。掃羅聽見了，就在瑪雲的曠野追趕大衛。 \end{tabularx} \\ \\ \relax
23:26 & \begin{tabularx}{0.7\textwidth}{X} 掃羅在山的這一邊走，大衛和他的人在山的那一邊。大衛急忙躲避掃羅，掃羅和他的人正四面圍住大衛和他的人，要捉拿他們。 \end{tabularx} \\ \\ \relax
23:27 & \begin{tabularx}{0.7\textwidth}{X} 有使者來對掃羅說：「非利士人入境搶掠，請快快回去！」 \end{tabularx} \\ \\ \relax
23:28 & \begin{tabularx}{0.7\textwidth}{X} 於是掃羅不再追趕大衛，回去迎擊非利士人。因此那地方名叫西拉‧哈瑪希羅結。 \end{tabularx} \\ \\ \relax
23:29 & \begin{tabularx}{0.7\textwidth}{X} 大衛從那裡上去，住在隱‧基底的山寨裡。 \end{tabularx} \\ \\
[1ex]
\hline
\hline
\end{longtable}
$^{1}$昨天是甚麼節日?.
昨天有沒有情?.
好像有人喝醉了.
如果大家有看昨天的「靈命一郎靈修」.
講到耶利米蘇的經文是神用永遠的愛愛我們.
今天可能是連續有關係.
今天是關於 知不知道有沒有誰今早有看?.
彼此相愛.
昨天我們剛好在團契聚會內容裡面講這個題目.
大家分享祈禱,學習如何從約翰福音,約翰一書等經文.
提醒我們如何彼此相愛.
無獨有偶今早的靈修經文就是約翰福音13章34節.
主耶穌賜給我們一條新的命令.
叫我們甚麼?.
彼此相愛.
主耶穌說我怎樣愛你們,你們也要怎樣愛別人.
講到情人節.
講一下友情.
在大街生平裡面我們發覺有一個很重要的人物.
和他有一段非常美麗的友情.
可以見到弟兄彼此相愛的心.
講到朋友.
現在這個世界的朋友跟我們以前的朋友的定義可能很不同.
因為現在的媒體的發達.
我們這一代不知道,我們下一代的那些.
經常都說他有多少個朋友.
Facebook就說.
有沒有人用Facebook?.
有些.
經常都說想和你做朋友.
你要不要他做你的朋友?.
所以現在的現代人對於朋友的定義.
其實真的跟我們過往的認識可能很不同.
很多時候那些朋友.
雖然跟你做了朋友.
不過可能是最大令你痛苦的人也說不定.
因為現在發覺很多那些青年人的自我形象.
或者他們的心理的影響.
甚至有憂鬱.
很多跟現代的媒體有關.

$^{41}$媒體很多時候在上面說很多東西.
隨便說.
然後就會引致傷害到他們的自我.
自尊心,自我形象.
那個本來是朋友.
其實發覺原來不是朋友.
究竟我們怎樣能夠做一個真正的.
好的朋友呢?.
從大衛和約拿大身上可以見到.
做一個良朋益友.
在大衛的人生當中.
他也有很多朋友.
一會兒我們也會看一看.
不過我相信.
講大衛一定不可以少講.
他這位的摯友,這位良朋益友.
無論在他生命的安危上.
或者在他情緒的支持上.
情緒支持很重要.
更重要我們一會兒會看到.
就是他跟神的關係上.
所以當我們講到在主內邊.
或者我們學習.
聖經的命令彼此相愛的時候.
就是真的跟神不可以分開.
我們跟世人有不同的一種關係.
或者朋友的關係.
與在神內建立的一種關係是很不同的.
所以我們見到這個就是.
約拿單純是大衛的幫助.
可以算是他的良朋益友.
也可以說是一個屬靈的知己.
朋友也有很多種.
「真言」我們也有看過幾篇經文.
其中一篇是這樣說的.
「真言」講了很多關乎朋友的關係.
十八章二十四節.
我們一起讀這篇經文,請讀.
「濫交朋友的自取敗壞.
但有一朋友比弟兄更親密」.

$^{81}$比弟兄就是他肉身的兄弟.
比肉身的兄弟更親密.
但是如果有些壞的朋友.
所以可以說是自取敗壞.
因為那些濫交朋友.
好像今天有很多人.
他不能夠找到一些益友.
反而有很多是甚麼?.
損友.
有些是損友,有些是益友.
我們希望我們不要交到一些損友.
能夠有一些好的益友.
教會裡面的弟兄姊妹就很有福了.
因為這個實在是我們的本錢.
我們有很多機會能夠交到益友.
但我相信也都求神幫助我們.
不單是能夠交到益友.
自己能夠成為別人的益友.
這個我相信是很重要的.
所謂彼此相愛就是有彼有似.
所以我們希望能夠有好的良朋.
也都希望我們自己成為別人的益友.
我們怎樣去學習交友呢?.
也都是成為別人的幫助.
實踐愛心.
我們思想一下這個大衛與若拿丹之間.
有多方便我們一起去借鏡學習.
第一我們看見它是大方不疾度.
在聖經 三號二記上十八章.
上一次我們看《打哥利亞》就是十七章.
今天看十八章.
十八章第一至第三節.
我特意選《環球聖經》譯本.
因為它其中有一個翻譯.
是比較翻譯得好一點的.
就是第一節 我們一起讀完這三節.
大衛跟蘇羅說完話以後.
若拿丹與大衛結為生死之交.
若拿丹愛他如同自己的性命.
從那天起.

$^{121}$蘇羅把大衛留下來.
不讓他回父家去.
若拿丹愛大衛如同自己的性命.
就與大衛立約.
真的是愛人如己.
若拿丹愛大衛如同自己的性命.
《環球聖經》翻譯特別的是第一節.
結為生死之交.
所以若拿丹愛他如同愛自己的性命.
講了兩次.
所以由第一節和第三節.
我們可見這是一個很重要的重點.
所以這個友情是出於愛的.
愛好像愛他自己甚至比自己更重要.
這就是他們這種生死之交的關係.
為什麼我說大方不著道呢.
我們都知道.
其實愛.
在哥林多禪書第13章大家很熟悉的.
愛是什麼?.
愛是不著道.
不自誇不張剛.
所以我們看到.
在大衛和若拿丹之間.
有一種大方不著道.
什麼叫愛呢.
我們從這裡看見.
若拿丹自己是一個戰士.
若拿丹打仗厲害嗎?.
真的厲害.
因為他在.
撒米爾紀上14章前.
他已經帶領.
以色列的軍隊去打飛利士人.
你知道以色列和飛利士是世仇.
打了很多場仗.
上一次我們講.
和飛利士人打就是哥利亞.
所以呢.
其實呢.

$^{161}$這個飛利士人不斷和以色列人有爭戰.
在十四章的時候即是早幾章.
若拿丹.
已經是帶領軍隊.
進入飛利士人的營裡面.
是非常危險的.
他帶著這些拿兵器的少年人.
和他差不多一年多的.
大家很有膽識.
去到這個.
飛利士人的營中 我們讀讀這節經文十四章的第六節.
我們請讀好嗎.
若拿丹對拿兵器的少年人說.
我們不如.
過到未受割禮人的防營那裡去.
或者耶和華為我們施展能力.
因為耶和華使人得勝.
不在乎人多人少.
這是一句名言.
我們看到若拿丹有沒有勇氣.
非常勇氣.
若拿丹的勇氣從哪裡來.
從上帝來的.
因為他很明顯知道我們幾個人如何深入別人的營中.
但是他不怕 或者耶和華要在當中.
為我們施展能力.
因為他認識到一個重要的道理.
因為耶和華使人得勝.
不在乎人多人少.
所以他的勇氣.
是出於對神的認識.
和信心.
所以我們看到他是一個.
不單是一個戰士 也是一個愛神的人.
現在過了幾章了 去到第17章.
又來了.
菲律賓人向以色列人罵戰.
罵了多少天 大家還記得嗎.
40天 天天出來說有誰出來和我打.
OK.

$^{201}$誰出來和他打.
數落高人一頭.
出不出去.
沒有.
好了.
14章才有一個若拿丹.
若拿丹都不夠膽出去.
這才顯出大衛的特別地方.
所以40天想了很久 出還是不出.
雖然他真的知道.
他以求神 不過他的信心還未大到可以面對這個巨人的哥利亞.
但是無論如何 我們知道他是一位愛神.
而且是能夠有把握 也有戰爭經驗的一個人.
結果誰出去打哥利亞.
贏還是輸.
如果你自己.
很有本錢.
打仗很厲害.
也認識上帝.
但是現在不夠膽出去.
結果有一個.
不知道從哪裡來的.
又是少年人.
他說憑上帝的能力出去打 結果打贏.
請問你當時的心情會是怎樣.
開心還是不開心.
開心 真是很大方.
酸酸的 妒忌是很正常的.
因為.
明明自己本來是少年英雄.
所以在爸爸眼中 以色列軍隊當中 他也是很厲害的.
結果40天不敢出去.
就覺得很厲害.
現在竟然有一個.
比自己.
還沒有經驗的戰士.
又沒有裝備.
竟然.
憑上帝的能力打敗他.
這個敵人.

$^{241}$你的心裡面.
正常來說 我們的心裡面會覺得.
總之是不是我應該出去呢.
或者.
又有什麼厲害的 又有什麼不對的.
或者.
這次你厲害 看下次吧.
人那種的.
罪性 或者我們人的.
天然有的.
就是會嫉妒.
所以才說我們要對付.
即是說 人之常情會有禍.
否則我們不是罪人.
人家比你厲害.
和你一起畢業的.
現在他.
升得很快.
找到一份很好的工作.
我呢.
還是單身.
他已經有孫子了.
對不起 可能沒那麼快.
起碼有著落了.
不知道你會不會一樣的心情都會有.
有些酸酸的.
或者甚至覺得.
有什麼厲害.
讀書的時候我明明比他厲害.
誰知道出來工作的時候.
你剛好找到一份好工作.
這個就是人之常情.
不過我們知道是不對的.
當我們信了主之後.
因為如果有愛.
他就會克服這種的嫉妒.
有愛就說.
我看見他.
能夠成為神所使用的.
我看見他能夠可以.

$^{281}$雖然他效過我.
不過現在神使用他.
我看到神的角度很開心.
這種是要聖靈幫助我們.
要有一種新的生命才會活出來.
舊的罪性就是剛才那種情況.
所以我們發覺.
在聖經裡面說到.
什麼是愛.
其實不只是說他有什麼需要就幫幫他.
固然是.
我們很難想到什麼例子去說什麼是彼此相愛.
但其實原來是從心開始的.
從心開始你的心裡面能夠.
好像約拿丹教他愛他如同自己的性命.
真是難得.
因為他體重他.
和他覺得這個都是.
神真的很愛的人.
這個從這個眼光.
出發點就很不同.
所以在和合本他說心相契合.
但在環球異本他說成為生死之交.
我覺得這個形容.
約拿丹和大衛是實在更加貼切.
真是他們生死與共.
他們兩個人都看見神的能力.
也都看見神愛以色列國.
所以為了神愛以色列國.
他們的心就很歡喜.
無論這個是我也好或者大衛也好.
只要神的名興旺.
神德榮耀我的心就很歡喜.
這個就是那種.
不同的地方.
所以我們見到這兩個少年英雄很相似.
這裡就說.
耶和華使人得勝.
不在乎人多人少.
誰說的?.

$^{321}$大衛就這樣說.
大衛在十七章.
四十七節我們一起讀.
由使者眾人知道.
耶和華使人得勝.
不是用刀用槍.
因為.
爭戰的勝敗.
存在乎.
耶和華.
像不像?.
很像.
所以我們看到那個微妙的地方就在這裡.
所以他們兩個的說話都很相似.
兩個說話很相似,相似在哪裡呢?.
就是相似在.
他們對神的信靠和認識那裡.
所以 C.S. Lewis,路易斯,大家應該聽過他的名字,很出名的.
一位基督教作者.
他曾經這樣說,他說當兩個人.
在同一條道路上發現彼此.
在同一條路上發現彼此.
友情就開始發生.
友情開始的時候通常也會有這樣的對話.
這個作者這麼寫,很有趣.
什麼?.
你也是這樣想.
我以為只有我一個人這樣想.
如果白話一點就說,什麼?.
原來你也是這樣想的?我以為只是我這樣想.
就是說他們在一條路上發現有一個人跟他們自己的想法很相似.
友情就在那裡開始發生.
這個很相似,在這裡讓我們看到就是.
他們識英雄者,什麼?.
還英雄.
識英雄者還英雄,就是說你跟我一樣相似.
我們爭戰在乎.
不是在乎人多人少,不是在乎.
用刀用槍.
我們爭戰的勝利全在乎耶和華.

$^{361}$這個就是他們自己心相契合的原因.
如果你找良朋益友,找一些心相契合,跟你一起.
都有這個深愛上帝,愛神的角度,而且彼此去欣賞.
彼此可以大方不嫉妒,這個真是真愛的表現.
所以我們看到他主動跟他結盟,剛才我們看見他跟他立約.
主動結盟,歸為王子,你知不知道?.
當時他是訴羅王的兒子,他是王子.
大衛是誰?大衛是一個木桶.
所以一個王子主動跟一個木桶結交,結盟,立約.
這種真是很大方,自己降卑自己,很謙卑.
來到說我欣賞你,不在乎我的地位,不在乎我的財富.
不在乎我的能力,在乎我們大家都有一個同一個心.
這個真是好,所以這種真是一個真愛.
難怪他們一見異故,也是真正友情的基礎.
弟兄姊妹,今天你想成為別人的良朋益友.
你找一個良朋益友,我們都要學習做一個大方不嫉妒的人.
這不是天然,是需要神的幫助,是需要聖靈的更新.
還有我們有屬神的眼光,有屬神的眼光.
他這麼厲害,好啊,這種就是寶貴的地方.
所以反過來,蘇羅對大衛是相反的.
蘇羅見到大衛非常嫉妒,他的心懷疑殺了他.
找機會害他,這種不是真正的朋友.
所以我們見到蘇羅和約拿丹一個父一個子.
倉天共地對於大衛.
所以你想做誰?你想做蘇羅還是約拿丹?.
在乎我們的眼光,我們的心性.
我們能不能夠有這種屬神的愛心.
最近是冬季奧運會,大家有沒有看?.
最慘的那個Lindsey,Lindsey Vonn.
所以我們上個星期在開獎執會,有人放了一個影片給我們看.
說她怎樣慘跌.
這個不是冬季奧運會,這個是曾經有的夏季奧運會.
這個小姐是基督徒,叫Esther Kim,韓國人.
她20歲就代表美國可以參加跆拳道的比賽.
那時候她們準備出賽之前,誰代表美國呢?.
她有一個好朋友,跟她一樣厲害.
這個好朋友叫做Kay Po,是白人.
所以當他們要選拔第一,即第一選手,都有排名的.
代表美國出去做這項比賽的參選者.
他們要自己打自己,即是準決賽,即是甄選賽.

$^{401}$甄選出來,贏了的就可以代表美國出去.
在這個甄選賽的時候,這對好朋友一起練習,平時互相切磋.
結果Kay Po那天就傷了膝蓋,傷得很明顯.
所以如果Esther跟她打,她就知道一定贏她.
因為她們大家差不多,她傷了膝蓋是一個很大的忌諱.
所以Esther相信她能夠輕易擊敗她.
但因為她們的友情,還有她這種體育精神.
所以她最後決定棄權,自己放棄讓Kay Po她的好朋友出戰.
所以結果很多人都很詫異,有些新聞記者訪問她.
為什麼你會這樣,為什麼你不為自己爭取.
她的答案是,我把自己奧運的美夢丟掉了.
其實不是的,你們以為我把奧運美夢丟掉了.
其實她說,這是我把奧運的得勝機會交給我最好的朋友.
還有我覺得最有可能贏的人.
真的很偉大,這麼輕的年紀,真是少年英雄.
所以我們發覺後來美國特意給錢Esther和她的爸爸.
她爸爸是她的教練,訓練她女兒.
一起免費參加奧運的遊歷,以致支持她們自己的隊伍.
本來要自己出錢的,現在美國給錢她.
為什麼呢?因為她給了她一個獎項,叫Citizenship,即是好公民獎.
她表揚她這種大方和真正的.
所以後來那些報道說,She is the real champion,她是真正的贏家.
所以我們真的很難發現,人生當中都是你爭我鬥.
真正的朋友甚至能夠不只是不嫉妒,而是有共同的目標.
為了大局著想,來欣賞對方,成全對方.
這種真是愛.
愛容不容易學?真的不容易學.
不過也可以學,當我們有時有嫉妒的時候.
我們記得,愛是不嫉妒,不自誇,不張狂,不做害羞的事.
所以我們見到這個實在寶貴.
所以大衛和約拿丹,第一,讓我們看到.
他們這種友情建立在大家彼此欣賞,不嫉妒,反而大方.
第二,很明顯,這就是後來她的父親說她贏了.
第二方面,慷慨肯犧牲,愛就是犧牲.
耶穌基督愛我們,在哪裡看見?就是為我們犧牲.
我們看到第十一章第四節,一起讀,請讀.
約拿丹從身上脫下外袍,給了大衛,又將戰衣,刀,弓,腰帶都給了他.
然後第二十章第四節,請讀.
約拿丹對大衛說:「你心裡所求的,我必為你成就.」.
我們說,素羅是一個非常敏捷的人.

$^{441}$他就說:「好啊,我收你做我的女婿.」收誰?.
收大衛,收你做女婿.
不過其實是否真的想收他做女婿呢?不是.
我們儒學班的同學一定讀過這件事.
因為第一次我進去的時候,他問我.
「究竟素羅要大衛多少個羊皮?」.
我忽然間都記不住.
就是一百個羊皮,即是說在非利士人的包皮.
即是殺了一百個非利士人.
來交換將他的女兒米拉嫁給大衛.
其實為何要一百個羊皮呢?.
就是希望大衛去打的時候,被非利士人殺了.
所以其實他是害他的.
所以他不是真的想收他做女婿,是因為素羅嫉妒大衛.
結果大衛成不成功呢?.
不單止成功,還拿回來多一倍.
顯明了他的誠意,即是說我真的很想做你的女婿.
誰知道,答案是甚麼呢?.
答案是「對不起,我已經將他嫁了給別人了」.
「豈有此理」.
「你搗亂又叫我去一百個羊皮」.
「現在我拿了二百個回來給你」.
「你說沒有了,取消了這份婚姻」.
素羅就是一個這樣刻薄的人.
明明答應了別人的事,給不給呢?.
不但不給,還要害他.
相反約拿丹就不同了.
約拿丹說「你心裡所求的,我必為你成就」.
不是只有港字,剛才我們看到.
他將他的外婆給了大衛.
夠了,夠了沒有?.
還不夠,還給他甚麼呢?.
戰衣呀,刀呀,夠了.
還不夠?還給你一把弓,夠了.
還不夠?腰帶也給你,即是全部武裝都給你了.
你要記得,約拿丹是甚麼人呢?.
戰士來的.
你最能使用的工具,你自己家裡有沒有甚麼工具.
有些東西,我們有八樣工具.
用來用去都各有幾件,你有沒有覺得.

$^{481}$你的衣櫃裡面有幾十套衣服.
用來用去都那幾件.
因為有些東西是最重要的.
有些東西是戰士來說最重要的.
你會發覺,約拿丹這個真是慷慨,肯犧牲.
真是很慷慨的,這個真是很難得.
所以他很慷慨,不要緊,我給你.
很重要的東西,他給他.
為甚麼?因為他表達他的誠意.
其實不單止,不單止將外袍和戰衣給他.
約拿丹甚至說,你將來做皇帝.
我做你的在上,真是很難.
明明蘇羅死了,那誰做皇帝?我.
所以不單止連戰衣戰袍,連皇位都給他.
這個真是慷慨.
這種慷慨,相信真是一種犧牲.
如果我們說要交良朋益友,或者彼此相愛.
其實真是看到,究竟慷不慷慨,願不願意犧牲.
這也是一種很重要的指標.
有一次教會有一對兄弟,不是真是親生兄弟.
是主內兄弟,一個叫William,不是這個William.
一個叫Paul,今晚Paul好像沒有回來.
這裡有個Paul,Paul信主很久了.
William剛剛信主.
但是William以前好像昨天我們團契裡面扮飛仔一樣.
講粗口.
他有一個自小長大習慣了三句不對,兩句就是粗口.
他信耶穌,他決心知道這樣不好的.
所以他就戒,結果戒不戒得到?.
這麼容易嗎?講了幾十年.
他習慣了,講幾句,一興奮起來,他就爆兩句粗口.
一生氣起來又爆兩句,又是那些東西.
所以教會就比較好,比較留心.
回到自己工作崗位,周圍環境.
他又習慣了這樣.
他就跟Paul講,我知道很不對,不過習慣了這樣.
Paul就幫他想辦法,不要緊.
你現在跟自己立一個約,跟神立約.
你講一句粗口,你自己誠實,講一句粗口就罰5元.
那5元就用來奉獻.

$^{521}$所以你記住那個星期裡面,在工作也好,哪裡也好.
講了幾十句,你就拿50元回來,那個主人奉獻.
所以他就這樣做,結果第一個星期罰了100元.
然後第二個星期都是差不多.
第三個星期好一點點,不過都80多元.
第四個星期又回到100元.
所以四個星期之後沒有什麼進展.
他就跟Paul講,我真的不行了.
Paul說不要緊,我們繼續祈禱,求神幫助我們.
不過求神幫助我們之餘,我想個方法幫你.
現在開始你再這樣算,不過你不用自己再罰了.
你就回來告訴我.
他說好,又開始,又開始,過了一個星期回來.
他說我又失敗了,這個星期如果要罰就罰120元.
結果Paul怎樣呢?.
Paul就開了一張空頭支票給他.
只是Paul的錢,現在我給你支票.
這個星期我幫你出,下個星期你就用這張支票.
你說多少你就寫上去.
哇,現在他幫我出.
所以之後很大進步,過了一周55元.
厲害了很多.
那怎樣,誰罰?.
Paul罰了,可憐.
所以寫了55元放進去奉獻袋.
再過了一個星期,厲害了20元.
20元誰出?.
又是Paul出,又厲害了.
結果第三個星期,很開心的William跑去Paul那.
這個星期,0,我成功了.
有時候他自己逼自己都不行.
有一個人去幫他.
有一個人去成為他的accountability,負責任.
二來就是他這麼大愛心.
我怎樣可以濫用他的愛心呢?.
我不可以再令他失望.
所以他就鼓起勇氣,依靠著.
不停祈禱,不停提醒自己.
結果他就在這個軟弱上面勝過了.
親愛的弟兄姊妹,什麼是良朋益友?.

$^{561}$有時候我們會發覺.
原來這些,甚至為他們犧牲.
甚至為他們著想.
不是根根計較,甚至不是佔他們便宜.
而是他們真正為了建立對方而付出的犧牲.
這種真的很偉大.
求主都幫助我們.
我們都能夠去學習.
這種慷慨,肯犧牲的一種美麗的表現.
第三,我們看見在大衛和約拿丹的交情裡.
這段真的很寶貴.
就是多方常保護.
這段經文有很多.
我們一起讀完這七節經文.
我們一起讀吧,七節不是很多.
一起讀,請讀.
「素羅對他兒子約拿丹和中神博說要殺大衛.
素羅的兒子約拿丹卻甚喜愛大衛.
約拿丹告訴大衛說.
我父素羅想要殺你.
所以明日早晨你要小心.
到一個僻靜地方藏身」第三節.
「我就出到你所藏的田裡.
站在我父親旁邊與他談論.
我看他情形怎樣.
我必告訴你.
約拿丹向他父親素羅替大衛說好話.
說王不可得罪王的僕人大衛.
因為他未曾得罪你.
你所行的都與你大有益處.
他拼命殺那非理士人.
耶和華以色列中人大行拯救.
那時你看見甚是歡喜.
現在為何無辜要殺大衛.
流無辜人的血自己取罪?」.
這裡見到約拿丹怎樣保護大衛.
他用他的言語勸素羅.
又出計叫大衛不要回來.
你要躲起來.
然後他說我有什麼就告訴你.

$^{601}$我會在我父親旁邊與他談論.
勸他不要來害你.
結果第四節.
約拿丹向父親素羅替大衛說好話.
有沒有試過你幫人說好話?.
王不可以得罪大衛這個僕人.
你記不記得他對你有功.
記不記得對我們的國有功.
那時你不是很開心嗎?.
為何你現在要流無辜人的血自己取罪呢?.
這是不對的.
他真的很大膽.
他真是伸張正義.
他真是為朋友用言語保護他.
又用行動保護他.
所以連素羅都說他對.
第六節讀完它.
「素羅聽了約拿丹的話.
就指著永生的耶和華起誓說.
我必不殺他.
約拿丹叫大衛來.
把這一切事告訴他.
帶他去見素羅.
他就仍然侍立在素羅面前」.
所以真是勸服了素羅.
本來想殺他,結果怎樣?.
結果改變主意.
不殺他.
再離開他面前.
不過素羅就是這樣.
他經常都變.
有時好,一下轉眼又不行.
這是他的劣性.
所以素羅不是真的改變.
不過都是合理的.
素羅有時也會有良心發現.
有時也會有良心發現.
為什麼我這麼壞.
但是他的劣性沒辦法真的改.
他轉眼又來.

$^{641}$這種就不要跟他交朋友了.
因為是而又非.
不斷地轉來轉去.
但是約拿丹就不同了.
約拿丹真的用言語保護他.
顧及他的聲譽.
考慮到他的得失.
而且當他有危難的時候.
他挺身而出.
這些就是好朋友.
不是你好的時候.
那個奉承你.
是你落難的時候.
他挺身而出.
成為你的保護.
當然我們知道.
都是合理的.
是公義的.
不是什麼都幫.
有些人明明是做壞事.
你幫他這樣不對.
但是他其實是無辜的.
或者有時他不是故意的.
但是他落難的時候.
你作為朋友.
是多踩他一腳.
還是再去指責他.
還是挺身而出去保護他.
我們需要分開這兩個不同的.
在《箴言》裡面27章這麼說.
我們一起讀.
《箴言》27章第9節.
請讀.
高柔與香料使人心喜悅.
朋友誠實的勸教也是如此甘美.
所以作為朋友.
我們需要真誠的勸教.
才使他得益處.
如果他一錯就再錯.
你又不理他.

$^{681}$那就是害了他.
坦誠的告訴他.
這是你的問題.
需要改變.
不過是因為愛出於真誠的勸教.
而不是那種落井下石的心腸.
所以我們看到這種真是在患難的時候.
在被害的時候能夠挺身而出.
去竭力維護.
一同分擔.
這種才是真正的朋友.
在打棒球的歷史上.
有一個.
如果今天去Dodgers.
我們自己的家庭遊戲.
你會看到有一個銅像.
這個銅像是紀念這件事.
就是Jackie Robinson.
Jackie Robinson.
她是第一位參加美國職業棒球大聯盟的黑人.
所以我們知道.
黑人進入一些重要的美國的競賽.
或者某些社會地位是經過很大的奮鬥.
很大的掙扎.
而且實在是一個突破.
所以當時來說只有白人打棒球.
但她是第一位進入職業.
美國大聯盟的黑人選手.
所以她克服了皮膚膚色的障礙.
面對很多人對她的歧視和嘲笑.
她克服了.
成功了.
成為了職業選手.
不過你知道打棒球的事.
有時看狀態.
有時狀態好些.
有時狀態差些.
結果有一場在他們做主場的時候.
打的時候.
他因為犯了一個錯誤.

$^{721}$結果令到他們敗落了.
所以當時在場的觀眾.
特別是他的粉絲.
知道什麼是粉絲嗎.
現在我們說的朋友就是那些.
那些所謂的粉絲.
那些朋友喝倒彩.
本來是自己的朋友.
現在他一時失誤.
結果怎麼樣.
喝倒彩.
當時Jackie Robinson非常尷尬.
非常羞辱.
在場的中間四圍包圍幾千人看她的時候.
大聲喝倒彩的時候.
說多難受的情況.
結果在那裡差不多有十幾秒的時候.
他差不多想哭了.
結果有一位朋友.
這位朋友叫P.V. Rees.
他是白人.
他走出去.
是他好朋友.
他在當時去到球場的中間.
搭著Jackie.
然後和她一起低頭.
在眾觀眾面對.
很多觀眾喝倒彩的時候.
他和她一起面對.
結果隔了不夠十秒.
全場就縮靜.
很多人就反懲.
很多人就後來鼓掌.
就因為這個好朋友.
他挺身而出.
他願意來和她同擔羞辱.
面對眾人嘲笑.
他與她一同奮鬥.
所以今天就有這個銅像.
去表揚這件事.

$^{761}$就是Jackie Robinson和P.V. Rees的一種美好的情誼.
這個Dodgers.
我們家裡的遊戲.
所以我們看到這種真友情.
你發覺你和你的朋友.
當你落難的時候.
他們反過來踩你.
反過來杯葛你.
反過來拒絕你.
這些不是真正的朋友.
所以反過來我們應該怎樣學.
我們應該學習.
這種真是多方相保護.
來保護我們的朋友.
最後我們在這裡看見.
家賢多勉勵鼓勵.
這裡很有意思.
蘇羅是用惡言的.
相對的人.
《三國記》第20章.
我們來看看這裡的第十節.
先讀一下.
我們一起讀.
大衛對若拿單說.
你父親若用利言回答你.
誰來告訴我呢.
這個就是他剛才說的.
如果你的父親罵你怎麼辦.
結果怎樣.
是不是.
有沒有罵他.
我們來看看.
後來還有一章.
三十節至三十二節.
你看看蘇羅.
請讀.
蘇羅向若拿單發怒.
對他說.
你這橫耿勃逆之婦人所生的.
我豈不知道你喜悅耶西的兒子.

$^{801}$自取羞辱.
以致你母親露體蒙羞嗎.
耶西的兒子若在世間沒著.
你和你的各位必站立不住啊傻子.
現在你要打發人去.
將他捉拿交給我.
他是該死的.
三十二節.
若拿單對父親蘇羅說.
他為什麼該死呢.
他做了什麼呢.
他還維護他.
不過在這裡你看到蘇羅那種恐怖.
他發怒的那種惡樣.
那種惡言.
你有沒有看過媽媽罵兒子的恐怖片.
以前在香港真的聽到.
我真的聽到在街上.
媽媽罵兒子罵到用三字經罵他.
你們有沒有聽過.
很可怕.
這個蘇羅真是厲害.
我們再說這種時好時壞.
忽冷忽熱.
一時對你好.
一時惡到不得了.
真的不是我們要做的朋友.
箴言提醒我們.
不要和一些人做朋友.
包括什麼人.
包括喜歡說閒話.
洩漏密時.
回去自己看看.
箴言二十章十九節.
很喜歡說閒話.
別人告訴你密時.
他就洩漏出來.
他又說你壞.
挑撥離間.
這些真的不是好朋友.

$^{841}$所以這個箴言提醒我們.
不要和這些人做朋友.
第二不要和哪些人做朋友.
就是好飲酒吃肉.
只吃喝玩樂.
叫做酒肉朋友.
箴言二十三章二十節.
所以吃肉喝酒要遠離.
這是聖經教導的智慧.
第三就是這種.
反覆無常.
一時對你好.
一時對你好差.
其實是動機不良.
心術不正.
這是箴言二十四章二十一節.
回去看看.
他的心性很差.
所以有時裝出來對你好.
其實一時對你好.
一時對你好差.
這種反覆無常.
就表示他不是真正的.
這種動機不良.
心術不正.
他說聖經提醒我們.
要遠避邪死人.
第四就是好生氣的.
好生氣的人不可與他結交.
暴怒的人不可與他來往.
箴言二十二章二十四節.
所以聖經給了我們很多智慧.
其實這個反面來講.
就是我們應該怎樣追求.
與人做朋友.
如果你很喜歡說閒話.
那就要小心.
只是那些朋友的關係.
建立於吃吃喝喝那裡.
這些就可能會有些問題.

$^{881}$如果你發覺自己經常反覆無常.
你應該檢查你的內心.
有沒有這些的嫉妒.
有沒有這些的虛假.
有沒有這些的欺騙.
這些問題.
好生氣的.
就是留意我們自己的情緒.
我們的脾氣.
因為你這樣.
一定沒有人喜歡與你做朋友.
所以我們看到這個就是.
好朋友.
相反有一句.
是一位很好的基督徒的領袖.
這樣說:A friend is one who makes me do my best.
所以真正的朋友是鼓勵我做到最好的.
這個真的是很好的朋友.
所以加言多鼓勵.
你有沒有用你的說話.
不是去講人閒話.
不是去爆你的脾氣.
不是去反覆無常的來講欺騙人的東西.
而是成為一個幫助人做他最好的.
鼓勵人的說話.
在《撒慕爾記》上23章這樣說.
17節.
我們一起讀.
對他說:不要懼怕.
我傅蘇羅的手必不加害你.
你必作以色列的王.
我也作你的宰相.
這是我傅蘇羅知道了.
於是二人在耶和華面前立約.
這個真是很大.
剛才我們說的不單止是一種大方.
犧牲不嫉妒.
更加的看見他用他的說話肯定他.
他說:你將來必作以色列的王.
他說:我會幫助你.

$^{921}$所以他們就在神面前往下.
很認真在耶和華面前立約.
這是多美麗的鼓勵.
更好的我們看見他用神的話語鼓勵他.
第23章和這章前面的16節.
再讀好嗎?請讀.
「蘇羅的兒子約拿丹起身.
往那樹林去見大衛.
使他倚靠上帝得以見顧」.
這裡我們看見真是.
這就是主內的弟兄姊妹經常有的.
就是我們倚賴上帝.
又幫助人倚賴上帝.
所以我們看見當大衛低沉的時候.
本來他有信心.
但是約拿丹鼓勵他.
他說:使他倚靠上帝得以見顧.
這個真是很美麗.
所以大衛和約拿丹不單止是有友情.
更加是以神為中心的友情.
不單止有愛.
他更加是以神的愛去愛對方.
以神的話去兼顧對方.
所以我們看見他們在神面前.
在耳環面前立約.
這很明顯就是以神為中心.
他們起誓立約.
祝福對方.
然後鼓勵他.
所以我們看見他的鼓勵.
在20章裡面出現了很多次.
你回去自己看看.
在20章裡面我們看這幾節.
我們一起讀.
12節.
「約拿丹對大衛說.
願耶和華以色列的上帝為證」.
這個提到願耶和華以色列上帝為證.
第14節.
「你要照耶和華的慈愛待我」.

$^{961}$第16節.
「於是約拿丹與大衛結盟.
說願耶和華藉大衛的仇敵.
追討背弱的罪」.
這裡又提到甚麼叫以耶和華立約.
他以耶和華為我們當中的證據.
以耶和華的慈愛恩待.
以耶和華的公義.
你去結盟我們不會背弱.
耶和華為我們持守.
否則背弱的時候.
怨神重罰我們.
42節.
最後這節讀完它.
「約拿丹對大衛說.
我們二人曾指著耶和華的名起誓說.
願耶和華在你我中間.
並你我後裔中間為證.
直到永遠」.
然後23章第18節.
「於是二人在耶和華面前立約」.
真的很美麗.
你看見他差不多每句話都是以耶和華為中心.
以耶和華為證.
而且他說因為以耶和華為證.
我們不單止這一代.
就是我以後那一代.
甚至直到永遠.
結果大衛真的遵守這個結盟.
我們後來看見大衛真的遵守.
不單止厚待了素羅.
善待了約拿丹.
後來我們看見他更加善待他們的後裔.
所以這些真是好朋友.
真是以神為中心的好朋友.
這些真是愛.
是以神為中心的彼此相愛.
這些真是愛.
因為他以真理.
以真誠.

$^{1001}$來建立友情.
跟我們中國人的儀器有些不同.
有些是滴血為盟.
真的不同.
因為這種是以神為中心的友情.
所以如果我們成為一種能夠聆聽.
能夠鼓勵.
能夠體諒.
而且以神的話彼此堅固.
這種真是建立在神裡面的寶貴良朋益友.
我們看見大衛是一個非常重友情的人.
他以德報怨.
他是堅守承諾的.
他以犧牲贏得好友.
所以不單止約拿丹.
其實後來當我們讀信經.
他有好幾位好朋友.
包括了胡西和巴西萊.
都是生死之交.
但是約拿丹就是一個最特別的名證.
有沒有見過這個人?.
杏林子.
一位在台灣很出名的基督徒作家姐妹.
她有一篇祈禱文.
她的祈禱文是怎樣的呢?.
我們今天就作為講道結束.
她說:主,我喜歡朋友.
因為你知道孤單的路是多麼難走.
當我跌倒時.
我需要有人扶我.
受傷時有人為我.
寂寞時有人伴我.
快樂時有人分享.
痛苦時有人分擔.
受傷時有人為我.
寂寞時有人伴我.
快樂時有人分享.
痛苦時有人分擔.
得意時不嫉妒我.
失意時不離棄我.

$^{1041}$當我驕傲時有人提醒我.
當我犯錯時有人指正我.
若是我得罪了他仍能寬恕我.
我的主在我擁有別人的友誼之前.
先教導我如何成為別人的朋友.
這句最重要.
不以衣冠待人.
取人.
不以才世論人.
不以儒字斷人.
不以美醜閱人.
不閱己之長.
不揭人之短.
不誇己之功.
不忘人之恩.
教我言語柔黃.
態度誠懇.
心情開朗.
精神愉快.
因為這樣你才可以做人家好的朋友.
即使我不能為朋友做什麼.
也讓我默默付出我關懷的心意.
無盡的祝福.
好聽嗎?.
不得我們從大衛和約拿丹.
和很多朋友的美麗見證身上.
都能成為別人的良朋奕友.
這豈不是主耶穌?.
難怪是真愛.
難怪聖經說耶穌是我們最好的朋友.
我們低頭祈禱.
我們多謝主耶穌稱我們為朋友.
為我們捨命.
當我們越領受明白你捨身大愛.
我們就能夠學孝主實踐你的命令.
彼此相愛.
成為別人的良朋奕友.
也能夠讓我們的教會充滿著.
豐富主裡面的友情.
奉耶穌基督的名.

\newpage



\section{列王記上 1:15-31}
\label{sec:031665}
\textbf{應許消失了?}~~~~廖武洪牧師
\newline
\newline
連結: \href{http://www.church.com.hk/acms/content.asp?site=cdc\&op=show\&type=product\&code=031665\&layout=sermon}{\texttt{031665}} ~~~~ 語音日期: 2026-02-15
\newline
\newline
\hyperref[sec:031664]{\small{< < < PREV SERMON < < <}}
~
\hyperref[sec:index_chronic]{\small{[返順時目]}}
~
\hyperref[sec:031666]{\small{> > > NEXT SERMON > > >}}
\newline
\newline
列王記上 1:15-31
\newline
\begin{longtable}{cl}
\hline
\hline
章節 & 經文 (和合本修訂版)\\
\hline
1:15 & \begin{tabularx}{0.7\textwidth}{X} 拔示巴進入內室，到王那裡。那時，王很老了，書念女子亞比煞正伺候著王。 \end{tabularx} \\ \\ \relax
1:16 & \begin{tabularx}{0.7\textwidth}{X} 拔示巴向王屈身下拜，王說：「你要甚麼？」 \end{tabularx} \\ \\ \relax
1:17 & \begin{tabularx}{0.7\textwidth}{X} 她對王說：「我主啊，你曾向使女指著耶和華－你的神起誓：『你兒子所羅門必接續我作王，他必坐在我的王位上。』 \end{tabularx} \\ \\ \relax
1:18 & \begin{tabularx}{0.7\textwidth}{X} 現在，看哪，亞多尼雅作王了，你，我主我王卻不知道。 \end{tabularx} \\ \\ \relax
1:19 & \begin{tabularx}{0.7\textwidth}{X} 他獻許多牛羊、肥犢為祭，請了王的眾兒子和亞比亞他祭司，以及約押元帥，他卻沒有請王的僕人所羅門。 \end{tabularx} \\ \\ \relax
1:20 & \begin{tabularx}{0.7\textwidth}{X} 但你，我主我王啊，以色列眾人的眼目都仰望你，等你告訴他們，在我主我王之後誰坐你的王位。 \end{tabularx} \\ \\ \relax
1:21 & \begin{tabularx}{0.7\textwidth}{X} 若不然，我主我王與祖先同睡的時候，我和我兒子所羅門必列為罪犯了。」 \end{tabularx} \\ \\ \relax
1:22 & \begin{tabularx}{0.7\textwidth}{X} 看哪，拔示巴還與王說話的時候，拿單先知也進來了。 \end{tabularx} \\ \\ \relax
1:23 & \begin{tabularx}{0.7\textwidth}{X} 有人奏告王說：「看哪，拿單先知來了。」拿單進到王面前，臉伏於地，向王叩拜。 \end{tabularx} \\ \\ \relax
1:24 & \begin{tabularx}{0.7\textwidth}{X} 拿單說：「我主我王，你果真說過『亞多尼雅必接續我作王，他要坐在我的王位上』嗎？ \end{tabularx} \\ \\ \relax
1:25 & \begin{tabularx}{0.7\textwidth}{X} 他今日下去，獻了許多牛羊、肥犢為祭，請了王的眾兒子和軍官們，以及亞比亞他祭司；看哪，他們正在亞多尼雅面前吃喝，說：『亞多尼雅王萬歲！』 \end{tabularx} \\ \\ \relax
1:26 & \begin{tabularx}{0.7\textwidth}{X} 至於我，就是你的僕人，和撒督祭司、耶何耶大的兒子比拿雅、王的僕人所羅門，他都沒有請。 \end{tabularx} \\ \\ \relax
1:27 & \begin{tabularx}{0.7\textwidth}{X} 這事果真出於我主我王嗎？王卻沒有告訴僕人們，在我主我王之後誰坐你的王位。」 \end{tabularx} \\ \\ \relax
1:28 & \begin{tabularx}{0.7\textwidth}{X} 大衛王回答說：「召拔示巴到我這裡來。」拔示巴就來，站在王面前。 \end{tabularx} \\ \\ \relax
1:29 & \begin{tabularx}{0.7\textwidth}{X} 王起誓說：「我指著救我性命脫離一切苦難的永生的耶和華起誓。 \end{tabularx} \\ \\ \relax
1:30 & \begin{tabularx}{0.7\textwidth}{X} 我既然指著耶和華－以色列的神向你起誓說：你兒子所羅門必接續我作王，他必繼承我坐在我的王位上，我今日必這樣做。」 \end{tabularx} \\ \\ \relax
1:31 & \begin{tabularx}{0.7\textwidth}{X} 於是，拔示巴屈身，臉伏於地，向王叩拜，說：「我主大衛王萬歲！」 \end{tabularx} \\ \\
[1ex]
\hline
\hline
\end{longtable}
$^{1}$今天我跟大家去看這段經文.
可能大家覺得很特別.
為什麼在一個這麼開心的日子.
56歲生日.
又是年廿八.
快要過年了.
但是我要跟大家說的是一個權力的鬥爭.
是一場權力的鬥爭.
讀《遼王記》上第一章的時候.
我一直讀就好像在看一套宮廷權鬥的古裝韓劇一樣.
我相信在在很多人都有看韓劇.
裡面的角色有什麼?.
有世子.
有底下.
有大妃娘娘.
有左相.
有右相.
大家為了王位的繼承.
在那裡明爭暗鬥.
在那裡拉攏親信.
經文一開始就描寫大衛王已經老到.
其實剛才我們只是讀一段.
應該我們從第一節讀到最後一節.
整個故事看一次.
不過我就省略了前面和後面的東西.
不過原來經文一開始就描寫大衛王已經老到.
連蓋很多張被子都不覺得暖的狀態.
他身體退化.
老眼昏花.
手疼腳震.
有心無力.
以致退居幕後.
不問朝政.
好像一個老糊塗一樣.
沒有清楚去交代他的接班人.
以致令到整個國家的權力陣痛.
他的兒子亞多尼亞.
剛才讀的經文有提過這個人物.
亞多尼亞就乘機上演了一場權力的鬥爭.
就趁著父王老懵懂.

$^{41}$就想先下手為強.
封自己做世子.
王位的繼承人.
開始的時候.
我一路讀.
我會形容列王記上第一章是一個.
宮廷鬥爭的記錄.
但當我仔細去釘針.
裡面的經文描寫的時候.
如果大家有時間可以回去整個故事看一看.
我就發現.
列王記上第一章可以說是.
上帝應許存亡的一個啟示.
上帝的應許.
會不會因為人的問題而消失.
而落空的呢?.
上帝的應許.
又會不會因為人的態度而實現.
以及成就呢?.
就好像剛才陳執事講到.
五十六年裡面.
有起有跌.
有時候因為人的緣故.
我們看不清神的心意.
我們失去了神在我們教會.
或者在我們個人生命裡面的那個應許.
沒有實現也說不定.
但我們也看到有些人.
或者有些弟兄姊妹.
他們很努力.
很忠心.
很忠誠地去信任神的應許.
於是看到教會一路發展.
在這裡.
我們看到兩種生命的態度.
關係到神的應許.
落空.
抑或實現.
亞多尼亞的行動.
似乎在破壞上帝的心意和應許.

$^{81}$但是另一邊廂.
我們看到先知拿丹的敏銳.
和拔士巴的堅持.
卻去成就了上帝的心意和應許.
其實列王記上第一章.
不單止記錄了一場權力的鬥爭.
它也為我們去解讀.
上帝的應許.
如何會因為人的態度.
而被尊重或者被破壞.
列王記上第一章.
甚至可以說.
為整本列王記的內容留下伏筆.
列王記不是一本.
很光彩的英雄的故事或者歷史.
而是一部屬靈失敗者.
或者我們說失敗者的筆記.
一個記錄.
整本列王記記錄了一個又一個的.
王朝興起衰落.
裡面重複出現了一句評語.
如果大家有讀列王記上下的時候.
就看到這一句說話.
就是說.
他行耶和華眼中看為惡的事.
或者看為正的事.
整本列王記帶出一個很重要的訊息.
那個訊息是什麼呢.
那個訊息就是.
哪個人做皇帝做領袖做牧師都好.
不是最重要的.
最重要的是.
哪個人願意讓上帝作王.
他有沒有好好地去實現.
和成就上帝的應許和心意.
我分開兩方面來看這個故事.
第一方面我們看到.
說到神的應許.
人的自大.
和野心.

$^{121}$是會破壞上帝的應許的.
人的自大和野心.
破壞上帝的應許.
剛才我說過當大衛瑞瑞努爾.
無力又無心去領導國家的時候.
上帝應許他那個永遠的王朝.
正面臨一波政治的暗湧.
他的兒子亞多尼亞靜悄悄地.
掀起一場權力的鬥爭.
亞多尼亞是誰?.
大衛的兒子.
是第幾個兒子?.
原來亞多尼亞是哈琴.
原來大衛有很多個太太.
其中一個叫哈琴.
是哈琴在希伯倫為大衛所生的兒子.
大衛在希伯倫作王七年零六個月裡面.
原來有六個太太.
分別為他生了六個兒子.
而亞多尼亞是排第四的.
亞多尼亞從小在皇宮裡面長大.
他目睹爸爸的榮耀.
和宮裡面的鬥爭.
他也見證到他前面那三個哥哥.
那個皇阿哥.
就是誰呢?.
暗聯,基利亞和亞沙隆.
因為私慾,因為野心而墮落和失敗.
但是亞多尼亞沒有因為前車可鑒.
學到功課.
相反他驕傲自大.
原來作者用了一句很簡單的說話.
去描寫亞多尼亞在第六節.
這裡沒有打出來.
第六節.
他說他父親即是大衛.
數來沒有洗他幽門說.
你是作什麼呢?.
他岑俊美生在亞沙隆之後.
原來亞多尼亞生得和他哥哥.

$^{161}$大家如果記得亞沙隆都很帥的.
很英俊.
生得和他哥哥亞沙隆一樣.
那麼盡美.
容貌盡美.
可能因此而得到爸爸的寵愛.
因為沒有了一個帥哥.
亞沙隆死了.
有另外一個帥哥.
所以爸爸特別寵愛他.
他好像一個嬌生慣養.
被爸爸寵壞的皇帝一樣.
爸爸從來沒有好好地去管教他.
從小到大.
他想要什麼就有什麼.
他做什麼.
或者他有什麼行為.
大衛他爸爸從來沒有去過問.
你在做什麼.
或者為什麼你要這樣做.
爸爸的溺愛默默地塑造.
亞多尼亞那個自大的性格.
令他以為自己就是天之驕子.
皇位的繼承人.
當大衛顯得無能為力.
權力真空.
皇位的繼承權又未宣佈的時候.
作者就描述了.
描寫什麼呢?.
他就描寫亞多尼亞自尊.
說我必作王.
就為自己預備車輛馬兵.
又派五十人在他前頭奔走.
換句話說他帶著大隊人馬.
前面又有五十個人.
聲勢浩大地為他開路.
他們浩浩蕩蕩.
準備一場登基的大典.
亞多尼亞又拉攏了很多軍事.
政治.

$^{201}$和宗教的領袖做他的盟友.
包括大衛的以前很得力的元帥約克.
和他的王子.
還有祭司阿比亞他.
亞多尼亞.
不可能不知道父王的屬意.
他應該知道.
先知拿丹曾經向大衛.
表述過上帝的心意.
應許過所羅門繼承皇位的那位.
但是亞多尼亞的自大和野心.
令他拒絕上帝的心意.
他想先發制人.
去破壞上帝的應許.
他可能以為人定勝天.
我的意念高過上帝的意念.
我的計劃勝過上帝的心意.
他特意排斥那些知道上帝的應許.
和忠於上帝心意的人.
包括先知拿丹.
比拿瓦這些勇士.
和他的兄弟所羅門.
當然所羅門的母親拔士巴.
亦不例外被他排斥於門外.
亞多尼亞在耶路撒冷的郊外.
這些全部是前面的故事.
剛才那一段沒有說.
亞多尼亞在耶路撒冷的郊外.
隱羅傑水泉旁邊.
所犀列盆石那裡.
擺設了一場盛宴.
去宰殺牛羊肥毒.
請所有王子和猶大的神僕來參加.
為他自封為王來慶祝.
剛才我們唱歌celebrate celebrate.
在這裡慶祝.
當大家舉杯慶祝.
高呼願亞多尼亞皇萬歲的時候.
好了.
上帝的應許彷彿被人的驕傲自大破壞.

$^{241}$人的野心似乎取代上帝的心意.
那上帝的應許是否因此而消失呢?.
聖經裡面有其他的例子.
告訴我們.
人的自大和野心.
的確可以令上帝的應許消失落空.
上帝的心意似乎因為人的問題.
被破壞或者扭曲.
不要以為我們每個人坐在這裡.
每個人都為教會為神做很多事.
你知不知道可能你在破壞神的應許.
神的心意.
但聖經說.
人會失信.
但上帝會不會失信呢?.
其中一個事例.
巴別塔的事件.
是聖經裡面最早最經典的一個例子.
在《創世記》第十一章一至九節記載.
當時天下的人.
有一樣的口音和言語.
但是因為他們自大和野心.
以致自食其果.
當他們來到一個平原的時候.
本來好好的.
合而為一.
和諧共處.
安居樂業.
但是他們竟然異口同聲地說什麼呢?.
他說:來吧!.
我們要建造一座城和一座塔.
塔頂通天.
為要傳揚我們的名.
免得我們分散在全地上.
他們建造巴別塔的目的.
是想什麼呢?.
想通天.
通天的意思.
原來這個字.
原文可以解作一個.

$^{281}$很誇張很誇張的一個表達.
表示非常非常非常非常之高.
高到衝上雲霄.
這個字強調人的驕傲和野心.
以為可以用自己的力量.
去建造不得了的事.
Ricky 經常說的.
不得了的事物.
就是那個柱子.
另一方面.
古代近東的塔.
通常有很多石級的.
塔頂有一座神廟.
象徵人慢慢爬上去.
與天同等.
去挑戰上天的權威.
塔頂通天象徵人想憑自己的力量.
爬升到天上.
與神同等.
甚至想取代神的位置.
有時我們小心教會.
可能取代了神的位置.
而不是讓神在當中以馬來里.
基本上.
當日天下世人想建造一座城.
和一座塔.
目的是什麼?.
揚名六萬.
以為人的力量可以媲美上天.
與天比高.
人自大.
加上野心.
以為可以通天.
結果.
他們連在地上都站不穩.
被上帝變亂口音.
分散到全地.
另一個很經典的例子.
我們也很熟悉的.
就是蘇羅.

$^{321}$蘇羅王.
撒慕爾記上第九章和第十章.
記載上帝高立蘇羅.
成為以色列第一位的君王.
原來蘇羅在上帝眼中.
本來是一個合用的器皿.
有潛質.
成為一個好領袖.
去滿足以色列人的渴求.
渴求有領袖.
有皇帝管理他們.
在撒慕爾記上第十章有幾句話.
顯示上帝對蘇羅的應許是很清楚.
和很豐富的.
首先.
在撒慕爾記上第十章第一節.
上帝透過撒慕爾對蘇羅說.
「這不是耶和華高你作他產業的君嗎?」.
上帝是親自去揀選蘇羅的.
另外在撒慕爾記上第十章第七節.
他說「這兆頭臨到你.
你就可以趁時而作.
因為神與你同在」.
原來神也應許與他同在.
上帝也公開肯定蘇羅.
承認他就是被他所揀選的.
聖經在十章二十四節說.
「撒慕爾對眾民說.
你們看耶和華所揀選的人.
眾民中有可比他的嗎?」.
於是眾民大聲歡呼說.
願王萬歲.
但很可惜.
蘇羅的自大和野心.
扼殺了上帝對他的應許.
他親手將上帝的應許.
豐富的應許.
破壞甚至刪除了.
蘇羅沒有珍惜上帝的應許.
他不信服上帝的心意.

$^{361}$擅自獻祭.
又沒有完全滅絕阿瑪利人.
還為自己立了一個紀念碑.
所以上帝說.
在《撒慕爾記》十五章十一節說.
「我立蘇羅為王.
我後悔了.
因為他轉去不跟從我.
不遵守我的命令」.
最後.
撒慕爾的一句說話.
成為蘇羅失去上帝應許的一個總結.
亦成為我們每一個基督徒.
我們教會的一個警惕.
在《撒慕爾記》上第十五章二十二節.
我們都很熟悉.
背也不定.
耶和華喜悅.
凡濟和平安濟.
喜如喜悅人聽從他的話.
聽命聖於憲濟.
順從聖於公羊的自由.
因為人的自大.
人的野心.
上帝撤回對人的應許.
雖然蘇羅仍然坐在皇帝的寶座上.
大家都知道.
蘇羅其實有很多心虛的事情.
好像惡魔進入他的心裡面.
要找大衛彈琴.
以致平靜他的情緒.
但是上帝已經不再與他同在了.
曾言有一句說話.
我們都很熟悉.
驕傲在敗位以先.
狂心在跌倒之前.
雅各亦提醒我們.
上帝阻擋驕傲的人.
賜恩給謙卑的人.
我們是一間謙卑的教會.

$^{401}$上帝一定會賜恩給我們.
巴別塔的事件.
蘇羅的態度.
和亞多尼亞自封為王這件事.
都在警惕我們自大和野心.
可以使我們親手破壞.
甚至消除上帝對我們的應許.
但是歷史其實在不斷重演.
不斷教訓我們.
人自以為可以改變天命.
最近有一部電影.
說天命不可以隨便去改.
改變上帝的心意.
掌握自己的前途和命運.
到頭來是怎樣.
到頭來只是白費心機和努力.
譬如我們在歷史裡面看到.
那個亞歷山大帝32歲.
幾乎可以征服了全世界.
他看自己是神之子.
他東征西討.
不斷在證明.
他可以靠自己的力量.
去征服世界.
但是.
當他在光芒四射.
最燦爛的時候.
突然之間暴斃.
留下的只是混亂.
和40年的權力鬥爭.
最終.
他辛辛苦苦打回來的江山.
被他幾個將軍.
分裂為幾個王國.
如果大家有看希臘歷史的時候.
就知道.
今天國際政治局勢和軍事的對壘.
也都有相同的啟迪.
人覺得自己多厲害.
野心多大也好.

$^{441}$掌管歷史.
掌管人生命的上帝才是BOSS.
當人以為歷史會站在我那邊.
成功一定是屬於我們的時候.
通常.
就是我們跌倒的開始.
曾言提醒我們.
人心多有計謀.
唯有耶和華的籌算才能立定.
有人這樣說.
他說上帝的應許.
從來不是去燃點我們的野心.
而是獎勵謙卑順服的人.
我相信當.
稍幾班浸信會的頂子會.
牧師傳道同工.
謙卑的時候.
神就會獎勵大家.
好了 第二方面.
這裡更加精彩了.
人的敏銳和堅持成就上帝的應許.
人的自大和野心.
真的會令上帝的應許消失.
但是上帝其實有答應過我們.
祂的應許是永不落空的.
記得當所羅門獻殿的時候.
他很肯定地宣告.
在《列王記》上第八章.
如果我們再看下去一點.
耶和華是應當稱頌的.
因為祂照著一切所應許的.
賜平安給他的民以色列人.
凡值他僕人摩西應許賜福的話.
一句都沒有落空.
保羅提醒信徒.
在《光輪多號書》第一章二十節說.
上帝的應許.
不論有多少.
在基督都是實的.
所以值著他也都是實在的.

$^{481}$叫上帝因我們得榮耀.
希伯來書的作者也都很肯定地告訴我們.
在希伯來書第六章十七節.
他說上帝願意為那承受所應許的.
格外顯明他的旨意是不更改的.
就起了誓.
那麼上帝的應許和心意.
怎樣可以實現和成就呢?.
我用AI幫我做了一些漫畫.
AI做的.
有一天上游的水壩突然崩潰.
可能有很多人聽過這個故事.
洪水一直湧落下游的市鎮.
有個虔誠的基督徒.
見洪水一直湧過來他住的那三層的村屋.
情急起來.
迫切地祈禱.
去求上帝去救他.
他忽然之間記起.
以賽亞書第四十三章第二節.
大家知不知道那一節是什麼?.
那裡說你從水中經過.
我必與你同在.
你躺過江河.
水必不漫過你.
他記起這節聖經之後.
並且對上帝的應許深信不疑.
心裡面感覺很平安.
但是洪水繼續湧過來.
淹了他地下那層.
他立刻走上一樓.
他相信上帝會救他的.
果然.
他見到有人撐著一隻獨木舟撐過來.
撐獨木舟的人叫他快點上船.
他望一望獨木舟.
看見它又細又窄.
很不穩.
於是拒絕來者的好意.
寧願繼續等待.

$^{521}$因為他相信上帝會救他.
洪水開始湧到一樓.
於是他立刻走上二樓避難.
與此同時.
他見到有隻坐著幾個人的山板駛過來.
駛過來.
但是他見到船上的人很多.
很擠迫.
他擔心坐下去太重會翻艇.
於是他拒絕上山板.
繼續等待上帝救他.
不久洪水已經淹到二樓.
他無可奈何走上天台繼續等待.
很快他又見到有一艘漁船駛過來.
但是漁船又聲又臭.
心想不如多等一會兒.
或者會有一艘豪華遊艇來接他.
他仍然相信.
上帝會救他.
但是洪水很快.
連天台也淹了.
最後水淹眼眉.
他未見到豪華遊艇經過.
已經被洪水淹死了.
他來到天堂.
見到上帝.
他很不甘心地問上帝.
他說:上帝!上帝!.
你說過.
我可以從水中經過.
你必與我同在.
我的江河水必不慢過我.
但是你的應許都沒有實現過.
上帝很溫柔地對他說.
我已經派了很多艘船去救你了.
你沒有留意嗎?.
在《列王記》上第一章.
除了亞多尼亞還有兩個關鍵人物.
影響上帝的應許和心意.
他們是誰呢?.

$^{561}$先知拿丹和所羅門的媽媽拔士巴.
在這裡我們見到拿丹的敏銳.
配合拔士巴的堅持實現了上帝的應許.
成就了上帝的心意.
首先拿丹的敏銳和機智力挽狂瀾.
阻止了一場血腥的政治陰謀.
不僅挽救了所羅門和拔士巴的性命.
同時成就了上帝的應許和心意.
大家應該記得.
大衛借刀殺人.
謀殺了自己的戰士烏尼亞.
搶走他的太太拔士巴.
他同時犯了殺人和姦淫的罪.
之後拿丹很勇敢和機智地.
引用了一個比喻去斥責大衛.
引導他認罪悔改.
拿丹對上帝的心意很敏銳.
同時他敏銳於被上帝閱立的領袖.
因為犯罪會帶來的後果.
於是他就付諸行動.
去維護上帝對大衛家的應許.
在這個時候.
拿丹再次發揮他敏銳的本能.
他察覺到亞多尼亞的行動.
了解這些行動背後.
是對上帝的心意和應許的挑戰.
於是他馬上做事.
他去找拔士巴.
對他說:哈及的兒子亞多尼亞作王了.
你沒聽到嗎?.
我們的主大衛不知道.
現在我可以給你一個主意.
好保存你和你兒子所羅門的性命.
你進去見大衛王.
對他說:我主我王.
你不曾向婢女起誓說.
你兒子所羅門必接續我作王.
坐在我的位上.
現在亞多尼亞怎麼作了王呢?.
你還與王說話的時候.

$^{601}$我都會隨後進去證實你的話.
拿丹不單止提出警告.
也為拔士巴設計了一個很完美的計劃.
拔士巴先進去內宮見大衛.
提醒他當年的承諾.
就是你兒子所羅門必接續我作王.
坐在我的王位上.
這個承諾其實也是上帝對大衛家的應許.
然後拿丹在適當的時候進去.
去印證拔士巴說的話.
拿丹示範了一份敏銳的眼光和洞察力.
而且很機智地作出計劃和行動.
在第二次世界大戰發生之前.
英國首相邱吉爾.
早就看出了希特拉的野心.
當時英國大部分的政客.
都相信只有綏靖政策.
才可以換取和平.
但是邱吉爾不同意.
他不斷地提出警告.
他說希特拉不是要和平.
是要征服歐洲.
很多人嘲笑希特拉是一個好戰分子.
但是1939年.
當德軍入侵波蘭.
在1940年法國淪陷的時候.
大家如夢初醒.
證明了邱吉爾有先見之明.
他比整個歐洲的領袖和政客更加敏銳.
認清希特拉和納粹的狼子野心.
有人形容.
擁有敏銳眼光的人.
或者敏銳眼光的牧師.
傳道我們的執事.
通常.
認住這件事.
通常是勇於質疑大家公認的常識.
以及能夠在細節裡面.
看到大局的趨勢.
他會比其他人早一步.

$^{641}$看到人性的問題.
做牧師也是.
做領袖也需要這樣.
我再說一次.
擁有敏銳眼光的人.
通常勇於質疑大家公認的常識.
以為是的.
一向都是這樣做的.
越質疑.
能夠在細節裡面看到大局的趨勢.
他會比其他人早一步看到人性的問題.
拿丹沒有趨炎附勢.
沒有跟大隊.
沒有跟那些表面看起來很壯大.
很強勢的人來排隊.
為什麼呢?.
因為他對上帝的心意.
對人性.
對大局很敏銳.
又有一個漫畫.
也是AI做的.
小朋友可以看看.
有一個叔叔.
在海邊釣魚.
小朋友走過來.
很有興趣地看著叔叔釣魚.
叔叔的釣魚技巧.
很厲害.
很快就釣了一條魚.
旁邊的小朋友看到拍爛手掌.
叔叔看到小朋友興致勃勃.
於是準備將他釣到的魚全部送給他.
但是小朋友搖搖頭.
叔叔莫名其妙地問這個小朋友.
「小朋友為什麼你不要?」.
小朋友很不客氣地回答.
「我只想要你手上的魚竿」.
叔叔就問.
「你要魚竿幹什麼?」.
小朋友很聰明地說.

$^{681}$「這條魚雖然很多」.
「但是怎樣都會吃完的」.
「但是如果我有魚竿」.
「我就可以自己釣一輩子都有魚吃」.
有多少基督徒.
好像這個小朋友這麼聰明.
只要抓住最重要的上帝.
最重要的上帝.
堅持祂的應許.
結果一定比你眼前或者即時的滿足好.
畢氏巴曾經是大衛寵愛的妃嬪.
但是她因為和大衛通姦.
而承受犯罪的結果.
失去第一個為大衛所生的兒子.
是不是?大家都記得.
她內心一定感到很愧疚.
甚至可能一直不敢抬頭.
但是聖經記載.
在《撒慕爾記》第十二章二十四至二十五節.
聖經說.
「大衛安慰他的妻畢氏巴.
與她同寢.
她就生了兒子.
給他起名叫所羅門」.
接著又說.
「耶和華也喜愛她.
就藉先知拿丹賜她一個名字」.
大家知不知道所羅門有另外一個名字.
名字叫做耶底底亞.
原來如果考聖經問答錯.
可能很多人都不記得.
耶底底亞.
意思是因為耶和華愛她.
上帝透過大衛和先知拿丹.
安慰畢氏巴.
肯定她生的第二個孩子.
所羅門.
是耶和華所愛的.
後來大衛又配合上帝的心意.
答應畢氏巴.

$^{721}$所羅門將來要繼承王位.
延續上帝的心意和應許」.
畢氏巴可能感覺.
自己何德何能呢?.
我又是一個戴罪的婦人.
然而.
她知道她是一個蒙恩.
得到寬恕的人.
她沒有忘記上帝的應許.
但是她沒有用自己的方法.
去實現上帝的應許.
她只是默默地在等待.
和堅持.
堅持.
直到上帝透過拿丹建議她.
鼓勵她為上帝的應許付諸行動.
接著她默默地去信服.
進入內宮.
以一個很卑微的姿態.
聖經形容.
她屈身向王下拜.
然後很堅定又很溫婉地向王說.
「我主啊.
你曾向婢女指著耶和華.
你的神起誓說.
你兒子所羅門必接續我作王.
坐在我的位上」.
現在阿多尼亞這樣做.
那段經文大家都記得.
畢氏巴的堅持不是出於自大和野心.
而是她對上帝應許的信任和忠誠.
她知道如果大衛老殘.
大衛老懵懂.
加上阿多尼亞的自大和野心.
上帝的應許就會消失落空.
當她跟大衛說完這番說話之後.
拿丹也來到王的面前.
印證畢氏巴所說的話.
結果.
大衛記起上帝的應許和心意.

$^{761}$即刻行動.
宣佈所羅門就是王位的繼承人.
結果.
所羅門登基.
上帝的應許和心意得以成就.
不是出於人的自大和野心.
而是因為人對上帝的心意.
和應許的敏銳和堅持.
上帝的應許得到延續和實現.
所羅門代表大衛建造了聖殿.
而且成就大衛家永遠坐在王位上的應許.
聖經裡面有一個很出名的人.
以斯帖,我相信在座有些姐妹.
特別改自己王后的名字,是以斯帖.
以斯帖的故事是另一個很好的例子.
由於她的叔叔沒抵改的敏銳.
提醒以斯帖.
焉知你得了王后的遺憤.
不是為了現今的機會嗎?.
加上以斯帖.
死就死罷的堅持,記不記得?.
死就死罷的堅持.
終於挽救了她的同胞.
免於滅絕的危機.
延續了上帝透過她的選民.
成就普世救贖的應許.
說到人的堅持.
我又想起海倫海勒.
那個盲的,失明的作家.
他年幼的時候因為患病.
令到雙目失明.
雙耳又失衝.
他一生活在黑暗和沉默之中.
好像一個完全與世隔絕的人.
但是他對上帝的信心.
叫他選擇.
不向命運來低頭.
他開始很努力地.
去學習用手去捉摸.
用手指來聆聽.

$^{801}$甚至學會說話.
他的堅持令他上了大學.
後來還成為作家到處去演講.
又參與社會的運動.
推動傷殘人士的權益.
而且寫了很多本書.
激勵人心.
他有一句名言這樣說.
他說了一句.
很有名的說話.
他說人一生中最大的悲劇.
不是失明.
而是放棄.
最大的悲劇.
不是失明.
而是放棄.
也許我們在座當中.
你曾經相信.
或者抓住上帝的一個應許.
但是你等了很久.
未見到實現.
甚至感覺上帝不會兌現你的應許.
你開始覺得很累.
兌現他的應許.
你開始覺得很累.
很不耐煩.
你以為上帝的應許不了了之.
你想放棄.
但是這個時候.
請你不要忘記.
上帝不是人.
人會失信.
但是他永遠可信.
上帝的應許一定會實現.
保羅鼓勵我們.
我們縱然失信.
他仍是可信的.
因為他不能背負自己.
上帝的應許可能會遲來.
但是他的應許一定一定不會消失.

$^{841}$不會落空.
最後一個總結.
幾句說話.
上帝的應許會消失嗎?.
不會.
答案是大家都說不會.
但是會消失的有兩件事.
哪兩件事呢?.
第一.
就是我們對上帝應許的敏銳.
對上帝應許的敏銳.
第二.
我們對上帝應許的堅持.
對上帝應許的堅持.
亞多利亞的自大和野心.
將上帝的應許好像變成會消失.
會落空.
但是拿丹和拔士巴用他的敏銳.
和堅持成就了.
看來.
人看來.
會消失的應許.
我看到一句說話很有意思.
希望能夠成為教會和每一個.
牧者全道領袖.
弟兄姊妹的鼓勵.
他說上帝的應許.
不會過期的.
但是你的參與.
會影響他兌現的速度.
上帝的應許不會過期.
但是你的參與.
會影響他兌現的速度.
但願這句話.
成為.
稍縱整個教會.
你和我個人.
面對新的一年.
挑戰的一個提醒和鼓勵.
請我們一起低頭祈禱.

$^{881}$因為你透過.
亞多尼亞.
拿丹.
和拔士巴.
他們的行為.
來提醒我們.
如何去抓緊你給我們的應許.
我們願意.
來敏銳神理在我們生命.
在我們教會裡面的心意.
也抓住.
從起初.
到現在你給我們教會.
和給我們每一個已經相信你的人.
那個應許.
在這個世上.
生活的時候.
我們靠著.
你給我們的應許.
來快樂的前行.
讓.
神理在我們的生命裡面.
來得到榮耀.
我們這樣禱告.
奉耶穌基督的名.
\newpage



\section{創世記 12:4}
\label{sec:031666}
\textbf{聖經人聖經事:離開本地本族父家}~~~~黃克勤牧師博士
\newline
\newline
連結: \href{http://www.church.com.hk/acms/content.asp?site=cdc\&op=show\&type=product\&code=031666\&layout=sermon}{\texttt{031666}} ~~~~ 語音日期: 2026-02-07
\newline
\newline
\hyperref[sec:031665]{\small{< < < PREV SERMON < < <}}
~
\hyperref[sec:index_chronic]{\small{[返順時目]}}
~
\hyperref[sec:031667]{\small{> > > NEXT SERMON > > >}}
\newline
\newline
創世記 12:4
\newline
\begin{longtable}{cl}
\hline
\hline
章節 & 經文 (和合本修訂版)\\
\hline
12:4 & \begin{tabularx}{0.7\textwidth}{X} 亞伯蘭就遵照耶和華的吩咐去了；羅得也和他同去。亞伯蘭離開哈蘭的時候年七十五歲。 \end{tabularx} \\ \\
[1ex]
\hline
\hline
\end{longtable}
$^{1}$大家好,一齊又讀聖經.
看看聖經人物,聖經故事.
離開本地本族富家.
加字有些疊在花上.
離開本地本族富家.
一看就知道在說誰呢?.
亞伯蘭,亞伯拉罕.
是《生命說書》的網絡為大家提供的聖經短講.
今次的聖經故事繼續說亞伯蘭.
接續著上邊.
創世記十二章第一節.
耶和華對亞伯蘭說.
你要離開本地本族富家.
往我所要指示你去的地.
我必叫你成為大國.
我必賜福給你.
叫你的名為大理也要叫別人得福.
為你祝福的,我會賜福給他.
咒詛你的,我會咒詛他,地上萬族都要因你得福.
亞伯蘭照著耶和華的吩咐去了.
羅德也和他同去.
亞伯蘭出哈蘭的時候,七十五歲.
留意一下這裡.
這裡十二章第一節.
耶和華對亞伯蘭說.
你離開本地本族富家.
這裡的地點在哪裡呢?.
上一次是說.
爸爸他拉帶著亞伯蘭.
薩萊和羅德一起出去.
去到哈蘭.
本來想去迦南,人生交叉點哈蘭.
卡蘭就停了.
然後就一直住下.
直到他拉死了,二百零五歲.
我想問問你,如果你是亞伯蘭你會怎樣?.
這個好的父親.
亞伯蘭是一個好名字.
一個很偉大的父親.
一個好父親多數都是好兒子才會做好父親.

$^{41}$他就一直在那裡和他拉一起.
直到爸爸死了.
這個時間如果你是他你會怎樣呢?.
三個選擇.
繼續在那裡住下去.
但不太順,不是我無端端離開火鴻城市吳爾.
來這裡住是為了甚麼呢?.
如果去迦南,神指示我去的,那我就去.
如果我來到這裡,為了甚麼呢?.
不如有另一個選擇是甚麼呢?.
回家去吳爾,有那麼多親戚朋友在那裡.
不如回吳爾吧.
這個時間正在哈蘭人生交叉點的時候.
耶和華對亞伯蘭說.
離開本地本族父家.
往我所指示你的地方去.
這個時間是在哈蘭人生交叉點和亞伯蘭說的.
不要回去父家.
離開本地本族父家.
繼續去我指示的地方.
我指示的地方在迦南.
其實你承繼我對你父親說的話,去迦南.
其實也有跟亞伯蘭說的.
因為阿他一定跟他說.
我有一個感動,神叫我去迦南.
很落後的,海邊地方,低地.
沒有我們吳爾那麼火紅.
你去不去?.
亞伯蘭說我也有這樣的感動.
好呀,你去哪就去哪,一起去吧.
這個時間神特別呼召亞伯蘭跟他說.
我跟你爸爸說過.
現在又跟你說.
離開本地本族父家.
不要停在迦南.
我會在那裡.
我叫你指示的地方.
你會在那裡成為大國.
我會賜福給你.
你的名字會大.

$^{81}$你會叫別人得福.
你去那裡.
我要在那裡祝福你.
哇,在那裡祝福我,好呀.
而且沒有危險.
不危險.
為什麼?.
因為有我在.
為你祝福的,我會祝福他.
如果就在你,我就在他.
地上萬族會因你而得福.
你去那裡吧.
你會在那裡叫萬族得福.
在那裡叫萬族得福.
不是現在.
你如果聽我講,去到那裡.
我會叫你在那裡叫萬族得福.
願途危險,不怕.
我會保護你.
你去吧.
於是,四節.
亞伯蘭,好呀.
就照著耶和華的吩咐去了.
那羅德怎樣?.
羅德,你想回去吾爾還是跟我一起?.
羅德說,我當然跟你.
好呀,跟吧.
也和他同去.
亞伯蘭就出哈蘭.
這個人生交叉點的時候.
都不是年輕人了.
七十五歲.
就是這段聖經.
我們再看,第一.
離開前往.
離開本地本土富家.
往我指示你的地方.
我也有這樣的感覺.
如果是離開本地本土富家.
去我指示你的地方.

$^{121}$而那裡不知是甚麼.
不過你去那裡會很厲害的.
那就去吧.
好呀,很厲害.
我這個經歷我也試過.
就是差不多近二十年前.
當我離開我本身事奉了二十八年的宣道會.
的堂會,宣教,中國侍工.
神學.
和宣道院等等工作的時候.
我心裡面是很徬徨.
心裡面是很失落.
我當時的感覺簡直是這樣.
那些同事,其他宣道會牧師都說.
「黃牧師你離開宣道會嗎?」.
「沒有理由,我聽到也不相信」.
「因為你的額頭有個宣字在那裡」.
「做了二十八年是宣道會堂會信主」.
「進宣道會神學院」.
「然後出來就在宣道會工作一路」.
「只是在宣道會畢業之後工作了二十八年」.
「做過很多工作」.
「建立過宣道差會,希伯倫差會」.
「宣道差會,法國社協會」.
「宣道院,社會服務」.
「各種工作都做過」.
「你怎會離開呢?」.
當時我離開.
我的感覺是離開本地本族父家.
不過似乎是時間到.
神叫我離開.
又安排一些事情發生.
於是我就離開.
往神指示的地方去.
這裡還好一點.
因為它有講去迦南.
我離開的時候.
做什麼呢?我不太清楚.
不過當時我已經有兩個機構.
一個叫「開心樹社會服務」.

$^{161}$一個叫「生命樹宣教網絡」.
我正在做這兩項工作.
於是我就投身進去做這兩項工作.
這兩項工作就是我的迦南地.
現在迦南地擴闊了.
還有一個「獨人神學院」.
擴闊了.
可能會再擴闊一點.
就是會做一些人生教練.
或者輔導工作,聆聽人的心聲.
可能會越來越多.
會不會做長者的各樣服務.
增強一些呢?會的.
你見到去那裡是離開前往.
當時我最初離開.
我想是經歷了亞伯蘭很相似.
就是心裹裹跳.
不知道會怎樣的呢?.
是有一個地方去.
不過會怎樣的?.
去到那裡是否實好?.
中間會不會有強盜打?.
各樣不知道.
不過你去.
我會叫你成為大國.
我會賜福.
我別人因你得福.
我會祝福你的.
於是我就膽大地去.
別人會因你得福.
所以整個行程.
或者我去到的目的地.
無論什麼樹也好.
生命樹,開心樹還是什麼.
我都是想做一件事.
就是叫別人得福.
就是我人生最主要的目的.
亦都是亞伯蘭要走出去.
去迦南的目的.
就是叫地上萬族因他得福.

$^{201}$於是有第三點.
亞伯蘭說好啊.
我就照著耶和華的吩咐去.
不是自己想.
是照著耶和華的吩咐去.
其實我都曾經可以.
發生一個不義的事件.
火紅.
我試過和一兩間大機構.
在我博士畢業之後.
我博士是企業管理.
做過一些培訓.
我第一宗收錢的培訓.
是飛去上海.
和一間國際知名的美國大機構.
做過兩日三日的培訓.
是做什麼呢?.
是對於他們申請的二十幾個.
企業管理.
即是經理的trainee.
和二十幾個工程師的trainee.
組合在一起.
然後給他一個入職前的.
心理狀態的培訓.
那個培訓我就是.
將我當時比較拿手的.
生命裡很需要的五個自尊感.
自尊感是尊敬的尊.
自尊感.
這個self esteem.
trainee他們.
提升他們的自尊感.
我去.
飛機飛過去.
住好的酒店.
然後做三日的培訓.
錢是一萬元美金.
是我第一次的生意.
一萬元美金三日.
嘩!幾好啊.

$^{241}$一萬元美金.
八萬元港銀.
我就想.
如果一個月做一次就夠了.
我就想開始大展拳腳.
但當時其實我是一個牧師.
似乎神沒有叫我繼續做這個.
叫我收起那些東西.
專心跟我去迦南.
喂!不是吳爾嗎?.
迦南是低地.
其實是高地.
因為那裡很多人要因利得福.
於是我就照著耶和華的吩咐去.
當時我都不小.
不過沒有七十五.
都四十到五十了.
就這樣發生了.
所以我對於這段經文.
我都有相同的感受.
亞伯蘭厲害.
第一,神叫他離開去我的地方.
然後你去到那裡.
人會因你得福.
人因我得福.
不如我得福.
但他不自私.
神叫他就去了.
第三點是照著耶和華的吩咐.
他就出去了.
亞伯蘭將妻子撒娜和贈兒羅德.
帶了他們所有在哈蘭人生交叉點.
積蓄的財物.
所得的人口.
都帶往迦南地區.
他們就到了.
人生交叉點停一停.
原來神都有心意.
這裡才看到.
他們住在那裡.

$^{281}$雖然父親離開了.
但是他在哈蘭.
積蓄到很多財物和人口.
這個做生意的很厲害.
於是亞伯蘭就帶著他很多.
增加了在人生交叉點.
匯聚了新的力量資源.
然後出發去神為他預備的地方.
叫萬人得福.
這就是這件故事.
今句就在這裡.
創世記十二章第四節.
我們一起讀.
亞伯蘭就照著耶和華的吩咐去了.
照著耶和華的吩咐.
羅德和他一起去.
亞伯蘭出哈蘭.
這個人生交叉點.
吾毅然出去的時候.
年齡是七十五歲.
不是一個年輕人.
七十五歲.
又賺了這麼多錢在哈蘭.
那怎麼辦?.
繼續住下去.
再勞一個大筆.
然後衣錦還鄉.
回去火紅的地方做儲蓄.
吃老本.
退休.
他沒有這樣做.
照著耶和華的吩咐出去.
七十五歲才去.
願主祝福你.
如果你未到七十五歲.
都不要想著停下.
七十五歲出發.
往神指示你的地方.
只要你聽耶和華的吩咐就可以了.
神要使用你成為多人祝福的器皿.

$^{321}$一起祈禱.
天父多謝你.
並不是我們很厲害,很有用.
雖然我們都有些能力.
不過最重要是你.
在後面高舉我們.
是你支持我們.
是我們能夠站得住.
是因為你自己.
萬物從主而來.
我們這個人能夠站得穩.
是從主而來.
多謝主呼召給我們特別的機會.
我們可以去服侍一些人.
叫別人得福.
幫助我們.
好像剛才讀這本聖經.
原來神是想透過我們.
叫多人得福.
叫萬人得福.
原來我們在世上生活的目的.
是要叫人得著你的福氣.
給我們勇敢的做一個分享者.
勇敢的做一個傳福音的人.
勇敢的做一個敬拜的人.
帶領更多人在你面前得福.
主有什麼要做.
有什麼地方要去.
求你吩咐我們讓我們知道.
我們願意雖然年紀不小.
不過願意去.
是因為你自己的痛楚和賜福.
讚美你.
願你得著榮耀.
禱告奉耶穌基督的名.
阿們.
\newpage



\section{約翰福音 6:1-15}
\label{sec:031667}
\textbf{榮耀君王:要他作王}~~~~林鍵牧師
\newline
\newline
連結: \href{http://www.church.com.hk/acms/content.asp?site=cdc\&op=show\&type=product\&code=031667\&layout=sermon}{\texttt{031667}} ~~~~ 語音日期: 2026-02-14
\newline
\newline
\hyperref[sec:031666]{\small{< < < PREV SERMON < < <}}
~
\hyperref[sec:index_chronic]{\small{[返順時目]}}
~
\hyperref[sec:031668]{\small{> > > NEXT SERMON > > >}}
\newline
\newline
約翰福音 6:1-15
\newline
\begin{longtable}{cl}
\hline
\hline
章節 & 經文 (和合本修訂版)\\
\hline
6:1 & \begin{tabularx}{0.7\textwidth}{X} 這些事以後，耶穌渡過加利利海，就是提比哩亞海。 \end{tabularx} \\ \\ \relax
6:2 & \begin{tabularx}{0.7\textwidth}{X} 有一大群人因為看見他在病人身上所行的神蹟，就跟隨他。 \end{tabularx} \\ \\ \relax
6:3 & \begin{tabularx}{0.7\textwidth}{X} 耶穌上了山，和門徒一同坐在那裡。 \end{tabularx} \\ \\ \relax
6:4 & \begin{tabularx}{0.7\textwidth}{X} 那時猶太人的逾越節近了。 \end{tabularx} \\ \\ \relax
6:5 & \begin{tabularx}{0.7\textwidth}{X} 耶穌舉目看見一大群人來，就對腓力說：「我們到哪裡去買餅給這些人吃呢？」 \end{tabularx} \\ \\ \relax
6:6 & \begin{tabularx}{0.7\textwidth}{X} 他說這話是要考驗腓力，他自己原知道要怎樣做。 \end{tabularx} \\ \\ \relax
6:7 & \begin{tabularx}{0.7\textwidth}{X} 腓力回答他：「就是兩百個銀幣的餅也不夠給他們每人吃一點點。」 \end{tabularx} \\ \\ \relax
6:8 & \begin{tabularx}{0.7\textwidth}{X} 有一個門徒，就是西門‧彼得的弟弟安得烈，對耶穌說： \end{tabularx} \\ \\ \relax
6:9 & \begin{tabularx}{0.7\textwidth}{X} 「這裡有一個孩子，帶著五個大麥餅和兩條魚，但是分給這麼多人還算甚麼呢？」 \end{tabularx} \\ \\ \relax
6:10 & \begin{tabularx}{0.7\textwidth}{X} 耶穌說：「你們叫大家坐下。」那地方的草多，人們就坐下，男人的數目約有五千。 \end{tabularx} \\ \\ \relax
6:11 & \begin{tabularx}{0.7\textwidth}{X} 耶穌拿起餅來，祝謝了，就分給坐著的人，也同樣分了魚，都照他們所要的來分。 \end{tabularx} \\ \\ \relax
6:12 & \begin{tabularx}{0.7\textwidth}{X} 他們吃飽後，耶穌對門徒說：「把剩下的碎屑收拾起來，免得糟蹋了。」 \end{tabularx} \\ \\ \relax
6:13 & \begin{tabularx}{0.7\textwidth}{X} 他們就把那五個大麥餅的碎屑，就是大家吃剩的，收拾起來，裝滿了十二個籃子。 \end{tabularx} \\ \\ \relax
6:14 & \begin{tabularx}{0.7\textwidth}{X} 人們看見耶穌所行的神蹟，就說：「這真是那要到世上來的先知！」 \end{tabularx} \\ \\ \relax
6:15 & \begin{tabularx}{0.7\textwidth}{X} 耶穌知道他們要來強迫他作王，就獨自又退到山上去了。 \end{tabularx} \\ \\
[1ex]
\hline
\hline
\end{longtable}
$^{1}$歡迎大家來到我們的周六崇拜.
上年七月我們開始了一個新的講道系列.
我們要宣講的是新約聖經裡面的約翰福音.
我們宣講的題目是「榮耀的君王」.
今年一月我們離開了約翰福音第一章至第四章.
進入約翰福音第五章.
聖經學者告訴我們.
從約翰福音第五章至約翰福音第十章.
小路約翰帶我們進入了一個節慶週期的一段教導.
我們發現由第五章開始.
每一章的背景都有一個猶太人的節日.
約翰福音第五章背後.
他說有一個節日.
結果他都沒有跟我們說那個是什麼節日.
但我們後來發現那個節日是安息日.
第五章的背景是安息日.
第六章的背景是愉悅節.
第七章至第九章的背景是駐彭節.
然後去到第十章有獻殿節.
在這個節慶週期的當中.
耶穌通過了他的神蹟和教導.
以節期作為一個背景去啟示他自己.
其實耶穌所行的.
我們一直都知道不是一個普通的神蹟.
而是一個天國的記號.
當中耶穌行過的神蹟或記號.
包括有騎子行走.
這個是第五章的神蹟.
然後有五餅二魚是第六章的神蹟.
然後是水上行走是第六章的神蹟.
下星期回來我們會看那個神蹟.
最後有一個第九章裡面.
他使核子看見.
耶穌用這些天國的記號.
啟示他屬天的身份和他的使命.
使他們可以更加清楚的認識他.
然而他們越是清楚認識耶穌.
猶太的領袖結果就是越來越.
恨惡要逼迫他.
去到約翰福音第十一章.

$^{41}$第五十七節的時候.
祭司長和法利賽人.
是已經下令要捉拿耶穌.
要送他上十字架.
但是所有的一切.
他們的仇恨是出於約翰福音第五章.
我們說那裡有一個叫做黑色的安息日.
在那個黑色的安息日裡面.
耶穌行了那個騎腿行走的醫治.
然而猶太人越發想要殺他.
因為他犯了安息日.
然後他又稱上帝為他的父.
將自己和上帝捍衛同等.
這個是我們之前有三段.
整個三篇講章在這裡說的.
約翰福音第五章這個黑色的安息日.
和很激烈的自辯.
耶穌自己為自己辯護.
耶穌總算逃過一劫.
現在來到約翰福音第六章.
再沒有猶太人領袖的死亡追捕.
現在我們有的是猶太的平民百姓.
他們愛護擁戴耶穌.
於是耶穌一見到這班愛護他的平民百姓.
他心愛的羊.
他就憐憫非常.
於是耶穌行了一個滿載溫馨的奇妙神蹟.
通常我們這樣想的.
耶穌自己親自下廚為他們預備一餐飯.
那餐飯叫做五餅二魚.
以五餅二魚餵飽了五千個男人.
再加上女人和小朋友.
相信閒閒地都有二萬人.
耶穌親自下廚為他心愛他的羊.
去做一餐非常溫馨的愛宴.
總共餵飽了差不多有二萬人.
這個是我們要讀的故事.
有趣的地方就是當我們今天去看這個故事.
是去關於約翰的.
當我們仔細地閱讀約翰福音裡面.

$^{81}$對同一個故事這個故事的描繪的時候.
我們驚訝發現約翰福音第六章所描繪的.
不是一個溫馨戀愛的情節.
而是一幅叛亂謀反的圖畫.
這個是我們最終得出來的結果.
說回耶穌所行的五餅二魚.
耶穌所行過無數的神蹟騎士裡面.
約翰是選擇了只有七件神蹟.
來作他的天國記號.
即是耶穌所做的事情是沒有人可以進錄.
然而在沒有人可以進錄的神蹟騎士裡面.
約翰選了七個去表明.
耶穌是誰是他選擇的天國記號.
這七個神蹟包括有第二章.
加拿分燕水變酒.
這是第一件神蹟.
第四章大臣兒子復活了.
第五章我們上次看的騎腿三八進安息.
第四個神蹟就是我們今天要看的.
五餅二魚飽千人.
然後下個星期回來我們看.
第五個神蹟黑夜狂風水上行.
然後讓福音第九章還有一個生來失明的見光明.
最後拉薩路讓福音第十一章.
死了四天出盆墓.
我想你留意.
五餅二魚在眾多耶穌所行過的神蹟當中.
是唯一一個在四福音裡面都有記載的.
另外一個神蹟不是耶穌行的.
是行在耶穌身上的就是復活.
只有兩個神蹟在四福音裡面都有記錄的.
一個是耶穌的神蹟他復活.
但他是被復活所以他是被神蹟.
另外一個是唯一一個耶穌行過的神蹟.
在四福音裡面都有記載的就是五餅二魚.
所以你開始明白五餅二魚是多特別的一件事.
而在約翰福音裡面.
五餅二魚剛剛好排行的位置就是第四正宗.
在七個神蹟裡面的中間.
連貫了上半部和下半部.

$^{121}$除此以外.
五餅二魚也是約翰福音當中七個神蹟.
和七個我試的宣告的連接點.
我不知道大家知不知道.
在約翰福音裡面有很多的七件事.
我們通常一說他有七件神蹟.
但他除了七件神蹟之外.
他還有七個我試的宣告.
七個我試的宣告是比神蹟遲開始的.
但是到了我們第四件神蹟的時候.
是耶穌第一次在神蹟裡面說我試.
於是你發現五餅二魚是一個很特別很特別的神蹟.
就在這個神蹟的當中.
耶穌宣告我試生命的糧.
由此可見對於約翰而言.
五餅二魚是一個非常關鍵重要的神蹟記號.
不過你要真正明白五餅二魚這個天國記號.
我們先要回到第五章.
因為原來第五章和第六章雖然是兩個完全不同的故事.
但是竟然有微妙的關係是有連結的.
我們回到我們很熟悉的第五章.
我們說約翰福音第五章所記載的.
不是一個普通的神蹟.
而是基督十字架天國救贖的記號.
耶穌是誰?.
耶穌是上帝的羔羊來到羊門.
耶穌是永生的活水來到不事大的池水.
耶穌醫治了一個騎腿三十八年的人.
這個三十八代表了出埃及以色列人.
出埃及之後半逆不順流浪曠野三十八年.
沒有辦法進入上帝應許之地.
沒有辦法進入上帝的安息.
三十八代表了這群以色列人.
耶穌來要呼喚要救贖要復活.
要帶領靈裡面半逆癱瘓罪孽深重的以色列人.
進入上帝所應許永恆的安息裡面.
以色列人癱瘓不能進入上帝的安息.
是因為他們迷戀埃及的偶像.
於是耶穌做完神蹟之後.
最後和這個人說了一句話.

$^{161}$他說不要再犯罪.
不要再回頭不要再貪戀埃及的偶像.
這個世界的榮耀.
你卻要愛慕堅守上帝的安息.
你要奔跑朝向上帝永恆的天國.
於是耶穌做了這個很奇妙很榮耀的神蹟.
然而可惜的是施行神蹟容易.
全陽福音很艱難.
他做完神蹟他要解釋神蹟.
耶穌的神蹟和解說.
沒有得到猶太領袖的接納.
相反引來他們很憤怒的迫害.
他們說想要殺他.
和他們一個很不公平的審判.
他們把上帝耶穌放在被告席上.
他們自己卻不知道.
其實被人審判的不是上帝.
而是他們自己.
不是耶穌被審判.
而是他們被耶穌所審判.
但當時的耶穌很有趣.
他將計就計假裝好像很無奈地受審.
其實他是為要通過他自己的自辯的時間.
因為被告的時候.
猶太人會給你自辯的機會.
於是他通過他自辯的見證.
繼續向這些猶太領袖宣講天國的福音.
目的是要使他們得救.
上星期我們在讀耶穌自己的獨白見證.
在他很精彩的獨白見證當中.
耶穌提起天父為他在地上留下了多項見證.
所有見證都是出於天父.
天父留下的是第一.
施洗約翰的見證.
師父天父留下的第二.
是他自己臣職其事.
聖供他的教導所有的見證.
天父留下的第三個見證是聖經的見證.
最後在眾多聖經.
舊約聖經的見證裡面.

$^{201}$耶穌特別強調天父留下給他的見證.
是摩西的見證.
耶穌特顯要指出摩西的重要性.
摩西是走最尾的那個.
為什麼?.
因為摩西.
如果你讀通摩西你就可以得救.
而摩西正正是第五章和第六章的連結.
當你由第五章去到第六章.
你知道第六章你讀的是什麼?.
是摩西的見證.
第六章你讀出來的是五餅二魚.
五餅二魚說的是天降的生命糧.
天降的生命糧其實是從哪裡出來的?.
是從摩西舊約的摩西那裡出來的.
是嗎哪的神蹟.
如果你心水清.
你會覺得耶穌的五餅二魚強又強.
不過比起舊約的嗎哪.
是一個小規模的重演.
你明白我的意思嗎?.
耶穌餵飽了多少人?.
五千個男人.
加上女士們和小孩子.
頂多是二萬多人.
舊約上帝每一天要餵飽多少人?.
二百萬人.
而且他一餵就餵了四十年.
三十八年.
你怎麼算?.
怎麼算耶穌這個東西好像很厲害.
但只是一個很小規模的重演.
不用這麼緊張.
但這個小規模的重演.
同時也是一個大規模的盛傳.
我們說耶穌究竟最後餵飽了多少人.
我問你.
你就算了五千個男人.
再加上婦孺和小孩子.
就二萬多人.

$^{241}$還沒有.
因為他做了這個神蹟之後.
還有沒有剩餘的?.
有.
剩餘的麵包可不可以再派?.
可以.
剩餘的麵包有多少個?.
十二個籃.
每個門徒都有一個籃.
門徒最後拿著籃子.
拿來做什麼?.
自己吃嗎?.
吃到退休嗎?.
吃到見上帝嗎?.
還是拿著籃子.
籃子又可以變成五餅二魚.
這次比五餅二魚厲害很多.
因為籃子其實比五餅二魚大很多.
這個籃子裡面的麵包可以被傳到地極.
我想問問你.
結果耶穌行這個神蹟.
有多少人吃了?.
現在還在滾動.
還在繼續滾動的數字.
是不是?.
我們也吃了.
現在何止是200萬人.
不知道有多少人.
因為福音在那個神蹟裡面承傳.
所以是一個小規模的重演.
是不是?.
但是一個大規模的承傳.
因為那個神蹟其實還沒完.
今天仍然持續著.
繼續著.
Anyway.
回到《約翰福音》.
《約翰福音》第五章到第六章的時候.
耶穌的聽眾再不是猶太領袖.
第五章是很難,很難以的.

$^{281}$猶太領袖.
到第六章是平民的百姓.
還有祂的門徒回來.
因為第五章門徒是失蹤了.
現在耶穌的門徒回來.
最親祂的人.
還有一些愛戴祂的平民百姓.
整個故事其實是.
耶穌救贖,上帝救贖的延續.
為什麼呢?.
因為.
如果你記得我一直強調上帝愛三八.
三八就是那三十八年的騎腿人.
代表不是猶太領袖.
是所有猶太人.
是不是?.
所以之前耶穌傳的福音.
因為祂被人去告.
只是對著那群領袖.
現在不是一樣了.
現在的福音對象變成了所有的猶太人.
上帝愛三八.
甚至將猶太獨生子賜給他們,叫一切相信的不致滅亡.
反得永生.
因為上帝猜他的兒子到三八來.
不是要定三八的罪.
而是要使三八恩下得救.
現在祂有一個新的聽眾.
這個聖眾是什麼?.
就是三八.
於是我們去到第六章的時候.
我想你記得.
耶穌繼續將上帝的愛.
和這個福音向五千個男人.
再加上他的妻子兒女一起去宣講.
我想你知道最後五名義如天降生命樑的故事.
是一個愛的故事.
是一個上帝救贖罪人的故事.
我想帶你回去看看我們這個非常熟悉的故事.
我們先看馬可福音.

$^{321}$或者婦女福音的版本.
在馬可,馬太,馬可,盧迦的福音書裡.
都有記載同一個故事.
在三個故事裡面.
馬可的故事是最詳盡的.
不過也加了一點.
馬太和盧迦所加的東西在括弧裡面.
我想再和你看看那個故事的大綱.
婦女福音是這樣寫的.
他們要耶穌帶門徒去一個荒野的地方去探險.
因為來往來找他們的人太多.
連吃飯的時候都沒有了.
為什麼要走呢?因為太忙了.
於是他們坐船去加利利河.
起初盧迦告訴我們.
他們是想去一個白菜大的地方.
你先記住.
然後有很多人看到他們上了船.
他們在陸邊.
這個不是出海的問題.
那個湖很小.
他看到他們於是跟著他們.
各乘步行跑到他們.
由湖一邊去到另外一邊.
你想想他們由香港去九龍.
但是其實香港和九龍是連在一起的.
於是你發現你跑步比他還要快.
於是他們去到九龍那邊等耶穌上岸.
這就是發生的事情.
從各乘步行跑到那裡.
比耶穌還要快.
於是耶穌一下船.
慘了.
一看到一大群人.
然後就憐憫他們.
這是馬可所寫的.
馬太也是重複了同一件事.
耶穌看見他們就憐憫他們.
因為他們如同羊沒有牧人.
真是很棒.

$^{361}$於是耶穌做什麼呢?.
如同羊沒有牧人就開始教導他們.
這是馬可說開始教導他們.
但是馬太和路加也加.
同時也治好更多的病人.
然後天已經黑了.
那個地方偏僻.
於是門徒提醒耶穌.
現在天黑了.
我們應該放他們出去吃飯.
再黑一點就不行了.
是門徒主動叫耶穌.
去猜他們吃飯.
但是耶穌反問.
我們可以給他們吃.
然後他們說.
我們可以給他們吃什麼.
我們有多少餅夠他們吃.
就拿了五餅二魚出來.
是沒有小朋友這件事.
是約翰先告訴他們.
不要管了.
我們現在暫時看到.
然後他們就一群一群坐下.
有100有50.
就是這樣.
這是傅瑞福音給我們看到的.
我們很熟悉的故事.
我現在想來看看約翰的版本.
約翰這個版本起碼遲了二三十年.
我想你特別留意.
約翰第四個福音.
新加了什麼元素.
約翰的版本是這樣的.
有一大群人.
因為看見他在病人身上所行的神蹟.
就跟隨他.
然後耶穌就上了山.
耶穌現在逃難不是因為剛才太忙.
連飯都沒得吃.

$^{401}$是有另外一個原因.
就是那些人瘋了的一窩蜂追他.
是因為看見他在病人身上所行的神蹟.
於是耶穌上山.
山在舊約聖經裡面代表上帝的地方.
你越上山越近上帝.
所以上山代表我要親近上帝.
已經開始有一種張力.
地上有一股權勢.
壓迫地來.
於是耶穌要退上山.
然後再加上一個新的東西.
是我們以前不知道的.
原來當時猶太人的逾越節近了.
然後還有很多不同的地方.
約翰的版本裡面是耶穌主動先.
我們到那裡去買餅.
是門徒沒提出叫他們走之前.
耶穌已經先說.
我們去哪裡買餅給他們吃.
這個是約翰獨有的.
耶穌做主動.
其他的地方.
傅雷福音是門徒先作主動.
然後耶穌是回應的.
這裡是耶穌一早已經知道他想做什麼.
他一早已經做了神蹟.
神蹟不是後來的一個回應.
而是從頭到尾耶穌知道神蹟.
於是馬上說我們應該給他去哪裡買麵包.
因為下面那句我們讀到.
就是耶穌已經知道他想做些什麼.
然後又加上一句.
這話是要考驗皮靈.
他是原知道.
一早已經知道自己要做什麼.
一早已經想好那天就是做神蹟.
不是因為事情發展到那個地步.
沒辦法不做神蹟.
接著還有安德烈的說話.

$^{441}$帶了一個小孩子.
全部都是新的東西.
這個小孩子帶了五個大麥餅.
大麥餅又是新的東西.
之前就說五個餅.
現在就告訴我們是麥餅.
兩條魚.
人們看見耶穌所行的神蹟.
重複之前開始說人們看見耶穌的神蹟.
就說他真是那要到世上來的先知.
耶穌知道他們要來勉強他做.
全部都是新的東西.
就獨自又退上山.
你看到嗎?.
起頭和結尾都是和上山有關係的.
你讀完這兩個版本.
你有什麼感覺?.
我的感覺一讀完之後.
就應該將黑色和白色都相反.
同一個故事來的.
同一件事情來的.
然而馬可福音或者傅雷福音.
選擇了一個比較光明的角度.
一面去看這個故事.
而約翰福音就相反.
他是特別要展示這個故事的黑暗面.
所以我把背景轉了.
現在原來約翰是選擇同一個人來的.
是同一件事情來的.
但是約翰選擇了一個黑暗面.
為什麼?因為耶穌是真光.
光照在黑暗裡面.
於是將黑暗的東西全都暴露了.
當然耶穌最後要照明這個世界.
但是現在還沒照明之前.
他先將黑暗的東西展露了出來.
讓我帶你進入今天的經文.
約翰福音第六章第一節.
「這些事以後,耶穌渡過加里尼海就是提比亞海」.
他一開始的時候說「這些事以後」.

$^{481}$「這些事以後」就是一個新故事的開始.
一個新段落的開始.
這個段落是約翰福音第六章.
總共有七十一節.
正如我之前說的.
需要四五篇才可以講完.
所以今天所做的是頭盤來的.
是開場白來的.
他說耶穌在這個時候.
在這個「這些事以後」的時候.
他渡過加里尼海就是提比亞海.
如果你記得之前在約翰福音第五章.
耶穌在耶路撒冷.
這個故事我們不知道隔了多久.
但是現在已經轉了一個地理的地方.
耶穌再回到北部的加里尼.
接著第二節很重要.
第二節有一大群人.
因為看見他在病人身上所行的神蹟.
就跟隨他.
約翰沒有告訴我們究竟有多少人.
但是我們後來知道有五千人.
有一群人我們通常想有幾十人.
有一群人來到Crossway.
那群人不可能是五千人.
但是我們後來看到是五千人.
這群人為什麼要來跟隨耶穌.
因為他們看見神蹟.
這句話我們之前已經聽過了.
有些人看見神蹟就相信他.
接著下一句是什麼.
耶穌卻不相信他們.
這群人看見神蹟就要跟隨耶穌.
好不好?不好的.
尤其是你經過撒瑪利亞.
你知道要聽見耶穌的說話才相信才行.
這群人有沒有聽過?.
完全沒有聽過.
是不是來聽的?.
都不是.

$^{521}$是因為神蹟而來.
來到的原因也希望有更多的神蹟.
所以所有東西都不是好的預兆.
而且還有五千個男人.
厲害了.
有信仰學者提出.
加利利的村莊其實人數是偏少的.
而且比較分散.
他說你要聚集五千個人.
甚至五千個男人.
再加上五千或者二萬個人.
計算所有的家庭.
其實是很難的.
而這裡可以做到的原因.
是因為那個時候是逾越節近了.
跟這個很有關係的.
就是有很多朝聖的人.
其實經過那個地區.
他們由世界各地不同回去耶路撒冷朝聖.
因為是逾越節.
在中葡.
有人問你知不知道.
為什麼那邊特別多人?.
那個山好像有群人停在那裡.
幾百人.
可能現在上千了.
為什麼?什麼事?.
然後有人說有沒有聽過耶穌?.
我們最近祂很紅.
很流行.
很多人點讚.
我們今天碰到祂.
你去不去?.
痛不痛?.
我們去.
就去.
於是這樣才聚集了一群人潮.
一群五千個男人的人潮去找耶穌.
而耶穌當時的反應是什麼?.
耶穌當時的反應是上了山.

$^{561}$帶著門徒一起逃跑.
回到上帝的庚前.
有可能門徒不是很願意.
難得.
耶穌什麼時候有這麼好的生意?.
為什麼不做多一點生意?.
難得今天有這麼多客人.
耶穌卻拉著他們上山.
因為世界的誘惑就在面前.
OK.
在這個時候.
靠近猶太人逾越節的時候.
這是火上加油的時候.
大家都在等待救主君王來的時候.
然後有人聽到消息.
有一個這樣的人.
然後有人說見過他的行神蹟.
其他人可能幾千個裡面.
還有幾千是沒看過他的行神蹟.
很想他行神蹟來看看.
但是聽了很多.
於是這是一個火上加油.
等待救主君王來打羅馬的時候.
我想你回到六章第二節.
這是很重要的.
我們在《婦女福音》裡面是看不到的.
是約翰給我們一個很赤裸的主題.
主題就是有一大群人.
因為看見他的.
在病人身上所行的神蹟.
因為看見他的神蹟.
跟隨了他.
這個迷戀神蹟是一個很負面的主題.
所以我想如果有時候.
我上個星期還有一些新人來了.
來香港旅遊的.
都問起為什麼我對神蹟這麼負面呢.
我說很可惜.
很可惜你之前半年不在這裡.
因為如果你跟著我們看約翰.

$^{601}$其實一早在這裡跟我們說.
神蹟是一個很可怕的勢力.
千萬不可以因為神蹟看耶穌.
然後他重複又重複地警告我們.
神蹟是一個很恐怖的東西.
在第二章第一節開始說到.
那些人看見神蹟.
也同時那時候也是元月節.
已經一年前可能兩年前都可能.
第二章第二十三節第二十四節.
耶穌在耶路撒冷元月節的時候.
過元月節的時候.
有許多人看見他所行的神蹟.
就信了他的名.
一定要看下一節.
耶穌自己卻不信任他們.
或者卻不信他.
在原文是同一個字.
因為他認識所有人.
看通看透.
但如果你沒有了最後那一節.
因為有許多人看見他所行.
你信了他.
嘩正.
這是我們教會侍工的金句.
打出來.
每個星期就祈禱就可以了.
你看完下一節.
會不行的.
完全不行.
去到第三章.
約翰重覆一個警告.
他說有一個法利賽人名叫.
尼哥德姆斯猶太人的官.
這人夜裡來見耶穌.
對他說.
拉比我們知道你.
由上帝那裡來作老師.
我們知道我認識你.
我們知道你是上帝的對人.

$^{641}$因為你所行的神蹟.
若沒有上帝同在.
無人能行.
好像很好.
但耶穌接下那句.
就是你必須要重生.
這個來說因為神蹟.
所以來找耶穌的人.
耶穌.
讓我告訴你發生了什麼事.
你是一個黑夜人.
你是一個未重生的人.
你是一個不認識上帝的人.
Oh my god.
因為神蹟來找耶穌.
原來是反映了.
你根本不認識上帝.
如果你說因為聽到我的到.
所以我今天晚上來找你.
這樣就對了.
因為他們聽了見證.
見證的到.
所以留耶穌在撒邁尼亞多兩天.
這是完全合乎正經的.
但這個人因為看見神蹟.
來找耶穌.
很客氣的.
耶穌說.
你是黑夜人.
你是未重生.
你根本不認識我.
去到第四章.
我們記得那個故事.
迦納的第二個神蹟裡面.
他這樣說.
耶穌自己作過見證說.
先知在自己的家鄉.
先沒有人尊敬的.
什麼意思.
到了加里尼加里人.

$^{681}$他說歡迎他.
尊敬他.
因為他們也在耶路撒冷過節.
曾經看過他在節期裡面.
所做的一切.
一切是什麼?.
神蹟.
這群人他們歡迎他.
因為耶穌的神蹟.
因為他們見到神蹟.
因為他們盲愛神蹟.
所以他們歡迎耶穌.
但他們歡迎耶穌的理由.
是因為神蹟.
不是因為聽.
耶穌說.
這些人在我面前.
令我作為一個先知.
在家鄉沒有人尊敬.
什麼意思?.
先知的工作是.
傳講上帝的話語.
但如果有一個人.
他歡迎你.
不是因為你的話語.
是因為你的神蹟.
這個人其實就不尊敬耶穌.
不尊敬上帝.
不尊敬先知.
所以我們可以見到.
所有的這一切.
告訴你.
神蹟是一個.
多麼令人害怕的東西.
而現在我們面對的.
就是一個神蹟的故事.
正如我之前所說的.
這是一個.
叛亂謀反的圖畫.
他們要謀反羅馬.

$^{721}$因為他們發覺耶穌.
很厲害.
耶穌這個人.
本身就是一件武器.
一件可以翻天覆地的武器.
但他們叛亂的.
不單止是羅馬.
而是上帝.
因為去到最後的故事.
他們是要強迫耶穌.
做他們不願意做的事情.
他們要強迫耶穌.
做他們眼中看為對的事情.
他們要掌管耶穌.
接受耶穌.
做他們要做的事情.
他們要追求一個.
地上的國度.
要追求一個地上的.
君王地上的榮耀.
我想帶你去看看整個.
五坪二魚故事的大綱.
然後再回來.
繼續看我們故事的經文.
你發現聖經學者告訴我們.
《約翰福音》第六章第一節.
十五節.
也是呈現一個.
ABCDCBA前後呼應的架構.
很有意思的.
你看看那個AA.
AA都是在說漢建神跡.
於是漢建神跡.
耶穌上了山.
然後漢建神跡.
耶穌倒退山上.
但是有個分別的.
分別是漢建神跡的第一個神跡.
是中文你看不到的.
因為中文是沒有眾數.

$^{761}$和沒有單數的分別.
你看英文就看到了.
原來第一次上面那個A.
漢建神跡是眾數的.
漢建什麼神跡.
許多病人身上的神跡.
漢建醫治神跡.
你明白嗎.
是眾數的.
去到下面那個A.
最後那個.
高潮.
是單數的.
那個神跡是什麼神跡.
五千的壯男身上吃飽.
五千的壯男.
吃飽了.
因為耶穌施展了一個神跡.
這個神跡厲害.
這個神跡是一個軍事的神跡.
打仗最重要是什麼.
糧餉.
除了可以醫.
如果有個軍醫.
耶穌又是軍醫.
又可以主持糧餉.
我說大佬厲害了.
所以你去到下面的神跡.
神經病.
因為一個神跡.
漢建一個神跡.
最後要擁戴他為王.
不是我們唱的詩歌.
什麼Crown.
協會什麼什麼Crown.
喂大佬.
是叫他.
今天就和我們一起.
下耶路撒冷.
直搗黃龍.

$^{801}$黃袍加衣.
現在就是時候.
這個是他們心想的東西.
這個故事可怕了.
接著BB.
哪裡去買兩百銀錢的餅.
接著吃精.
碎屑.
裝滿了十二個男子.
上面B是不可能的任務.
下面是不可能的結局.
難以猜想的結局.
接著CC.
CC又有趣.
CC都在說.
五個墨餅.
一個是門徒向耶穌說話.
另一個是耶穌對門徒說話.
接著所有的焦點在中間.
耶穌崩餅.
拿起彼來.
捉謝了.
哪裡聽過.
最後.
盲猜.
所以整個故事我們知道.
所有的一切其實是.
流向十字架.
那五千個男人.
聽不聽得明白.
就是另外一件事.
完全.
但是我們知道.
我們今天作為基督徒.
我們看回這件事.
所有的東西.
和約翰福音第五章是一樣.
約翰福音第五章.
是一個關於基督十字架.
天國救贖的記號.

$^{841}$上帝要愛三八.
上帝要拯救三八.
這個故事一樣.
都是說耶穌受了十字架.
要拯救.
讓祂自己拯救.
犯了罪.
在罪孽深重的.
以色列人三八.
兩個故事是一樣的.
不過那五千個瘋子.
我想.
是沒有辦法看到這個故事的.
我們回到這個故事.
我們繼續讀第五節.
耶穌舉目看見一大群人來.
就對腓力說.
我們到哪裡去買餅.
給這些人呢?.
他說者說是要考驗腓力.
他自己原知道要怎麼做.
一早已經決定要做五餅二魚.
腓力回答他說.
就是兩百個銀幣的餅.
也不夠給他們每人吃.
一點點.
耶穌一舉目.
到了這裡我們也不知道.
舉目有多少人.
直到後來五千個男人.
二萬多人次.
誇張了.
耶穌看見這群人.
肥女也看見.
就對腓力說.
我們到哪裡去買餅.
給這些人吃.
這是耶穌的主動.
婦女福音.
耶穌不是主動是被動的.

$^{881}$但是這裡耶穌很主動.
因為耶穌一早已經決定了他要做五餅二魚.
餵飽五千人這件事情.
聖學者在討論為什麼要找腓力.
有沒有想過為什麼要找腓力.
因為有兩個可能.
第一個人來說.
估計腓力可能是負責採購和物流的.
猶大是管數的.
十二門徒各有專長.
我又不明白為什麼不找以前做稅吏的馬太.
利美做管數多好.
不要管了.
但是有人說他可能是負責採購和物流的.
所以耶穌問他.
第二.
是因為腓力是白菜大人.
如果你之前看.
他們開船的時候.
其中一個方向是朝向白菜大.
就是說他是地頭蟲.
於是耶穌問你這裡熟.
去哪裡買.
這說是要考驗腓力.
耶穌自己一早已經考驗他要什麼.
如果他要考過試.
他應該怎樣通過呢.
為什麼他不通過呢.
我們看看腓力的回答.
腓力說.
兩百個銀幣的餅.
都不夠吸貪門每人吃一點點.
一個銀幣等於當日.
普通員工一天的工資.
計算條數大概八個月的薪水.
少則 minimum pay 幾萬元.
多則十幾萬甚至幾十萬的人.
為什麼他要這樣說呢.
他說就算我們有.
我不知道他有沒有錢.

$^{921}$就算我們有八個月的薪水.
拿去拼命買麵包.
每人都吃到一點點.
其實他的意思就是說.
不要玩了耶穌.
不要玩了.
為什麼要做什麼.
就算我們有錢.
都未必買到麵包.
麵包店已經關門了.
或者一個小鎮的麵包店.
怎麼可能買到這麼多麵包.
我們要去十個小鎮.
我們十個小鎮.
我們怎樣做物流.
就算我們找到麵包.
我們怎樣搬運回來.
耶穌是拿來搞的.
如果你記得在腓力福音的時候.
他一早已經建議了.
不如給他們分散回家.
自己去謀生.
所以他心裡面的想法就是.
耶穌不要找這些東西來搞.
之前我們已經做了一個建議.
送他們走.
叫他們自己去謀生.
每人去謀生找一點點.
我們要幫他們謀生.
怎麼可能呢.
拉比不要這麼固執.
我想這個是他的心態.
說到這裡.
你應該有一種人神鬥的感覺.
耶穌已經發了板.
主動說我今天就是要餵飽這些人.
就算你沒有辦法.
如果你老闆這樣說.
你最多可以請教一下耶穌.
耶穌不是說.

$^{961}$我決定了你做.
你知道耶穌是很合理的.
如果你說.
拉比你想這樣做也好.
不過我真的不知道.
你可不可以請教拉比.
你覺得我們應該怎樣.
但他這裡不是的.
他是人神鬥.
他是跟他捏過.
或者我們可以問.
如果肥力沒有氣權.
他基本上是氣權.
說我不考試了.
如果肥力沒有氣權考驗.
如果肥力會嘗試去回答這條難題的話.
這條難題就是.
我們到哪裡去買餅給這些人.
就是成千上萬的人.
他的答案應該從哪裡來.
答案應該是從聖經來.
答案應該是由摩西來的.
這也是為什麼耶穌要考肥力.
因為肥力應該是熟讀摩西的.
如果你看回第一章的時候.
肥力找到拿彈業.
跟他分享耶穌的時候.
他是這麼說的.
他說拿彈業.
摩西在律法書上所寫的.
那一位我們遇見了就是耶穌.
你要說出這句話.
你要熟讀兩樣東西.
你要熟讀耶穌.
第二你要熟讀.
摩西在律法書上所寫的東西.
這就是耶穌在批評那些猶太領袖.
不相信他是摩西書上所寫的.
現在肥力自己親口說.
你就是摩西書上所寫的那一位.

$^{1001}$肥力你好像很熟悉摩西五經.
肥力你好像很熟悉耶穌.
原來他所謂的知識.
全部都是腦袋裡面的知識.
去到身體力行的時候.
他原來是不明白耶穌是誰.
如果你真的回到摩西五經裡面.
去找答案的時候.
說真的在摩西五經裡面.
唯一可以餵飽成千上萬的人.
就是天降的嗎哪.
是這麼簡單的.
如果你熟讀.
你頂到你跟耶穌說.
不會是嗎哪吧.
你也不好意思說.
但是老師教科書裡面只有一個案例.
可以餵飽成千上萬的人.
就是只有嗎哪.
是不是你想跟我們說.
嗎哪就是答案.
我們要去哪裡找這麼多麵包.
天上.
但是肥力想都沒有想過.
肥力不知道耶穌是誰.
現在不知道.
當時不知道.
後來都不知道.
後來是什麼時候呢.
後來是最後晚餐都不知道.
你發現肥力再出來的時候.
又被耶穌揭穿了.
耶穌對肥力說.
對他說.
這個十四章九節.
我與你們在一起就是魔教.
你還不認識我是誰嗎.
看見我的就是看見了父.
耶穌不斷重複地說.
你怎麼說張父遣給了我呢.

$^{1041}$他不明白他不知道.
耶穌是誰.
所以肥力這次考試失敗.
五餅二魚的時候失敗.
最後晚餐的時候也是失敗.
不過.
你不要看不起肥力.
因為按照初期教會的傳統.
其實我是初期教會傳統.
對肥力的身世之後.
不是很仔細知道.
但是我們知道.
在耶穌受死撲末.
升天聖靈降臨之後.
肥力是通過了他的考試.
我們知道這麼多.
就是他衷心勇敢.
全人奉上傳揚福音.
福音是什麼.
就是那個從天降下來的生命靈.
他每次傳福音的時候.
都想起之前的失敗.
然後他用心的傳揚.
我告訴你.
這個就是天降下來的生命靈.
肥力用他的餘生.
單單吃喝耶穌.
歡喜傳揚耶穌.
受苦殉道.
為耶穌生存.
他是釘十字架死的.
這個是肥力的最後勝利.
當時考試失敗無所謂.
Math考試一百分都沒用.
你最後要有最後勝利.
這個也是我們每一個人.
必須要有的人生方向.
回到約翰福音第六章.
考試失敗的又何止是肥力呢.
然後有另外一個人.

$^{1081}$也會考試失敗.
這次又點名了.
約翰福音第六章第八節.
有一個門徒.
就是西門彼得的弟弟安德烈.
他其實也是白菜大那裡的人.
安德烈走來對耶穌說.
什麼意思呢.
安德烈說.
耶穌我也要和你咬下手瓜.
我又來一場人臣鬥.
你沒理由這樣詆毀肥力.
我也是白菜大人.
所以我現在和他一起.
咬下手瓜.
這裡有一個孩子.
帶著五個大麥餅和兩條魚.
但是分給這麼多的人.
我們還不知道有多少人.
你去到最後才知道.
還算什麼呢.
這裡通常我們的理解.
是有一個很可愛的小孩子.
很想幫忙.
但是他不自量力.
自告奮勇地自己走來耶穌那裡.
奉上五餅如魚.
但其實更有可能.
當時的情景是安德烈.
招募了一個兒童小演員.
來演一場臨時的天國劇場.
然後又找了五餅如魚作為道具.
然後.
只有笑容.
這個小朋友拿著五餅如魚.
來到耶穌面前.
目的是要耶穌覺得不好意思.
叫耶穌知難而退.
要耶穌知道他個人的想法是.
如何的荒誕無稽.

$^{1121}$拿著這個東西.
小孩子我全部奉上.
你看我是五千個男人.
在這裡.
這個小朋友拿著的.
中文翻譯是挺搞笑的.
五個大麥餅.
它不是大的麥餅.
是大麥的餅.
這個不一樣的.
你想像一下.
大披薩.
它不是大麥餅.
是大麥餅.
大麥餅的原文是很小的.
所以其實是五個大麥所造成的.
大麥就是窮家.
窮人家庭才吃大麥的小餅.
很小的.
再加上兩條小魚.
沙甸魚那些.
所以加起來兩個人吃三文治.
一餐就差不多.
所以其實是很搞笑的.
他故意弄到耶穌覺得不要搞了.
一個小孩子.
五個小麥餅.
兩條小魚.
現在要面對著五千個男人.
壯男.
二萬個人.
誰知道耶穌將計就計.
你想搞我.
你想玩我.
耶穌的順水推舟.
和門徒我覺得是開個玩笑.
你試我.
你跟我鬥.
我就是用你給我的東西.
做你覺得不可能的東西.

$^{1161}$覺得無稽的東西.
以奇人之道還治奇人之身.
第六章第十節.
耶穌說.
你們叫大家坐下.
那地方的草多.
人們就坐下.
結果坐了.
因為他們是五十人一百人.
所以其實挺好數的.
最後數到男人就有五千人.
耶穌拿起餅來祝謝了.
分給坐在的人也.
分了與照他們所要的來分.
我想你想像一下當時的情景.
門徒當時那種困惑.
那種無奈.
抗議不停.
你想像一下你要五千人.
或者二萬人.
分組分五十一百人坐下.
都需要時間.
耶穌說你們去.
將他們分五十一百人坐下.
你搞了一個小時.
在一個小時裡面.
你做門徒的心情是怎樣.
我覺得跟他們說.
瘋的.
你跟其他人說.
你覺得.
人們說為甚麼叫我們坐下.
為甚麼要分五十.
為甚麼要分一百.
為甚麼呢 是不是有東西派.
不知道.
你做門徒覺得很無奈.
瘋的.
不知道老師瘋了甚麼.
一小時大家坐下.

$^{1201}$沒有東西吃.
那你怎麼死呢.
所以他們心情是極其的複雜.
困惑無奈.
老師這樣說.
抗議不停.
然後耶穌最後拿起餅來祝謝了.
分給他們的時候.
你仍然覺得瘋的.
為甚麼要做這些事呢.
如果我去拍一套戲.
我會.
其實我真的不懂得怎樣拍.
我覺得你怎樣拍呢.
是否用塊布遮住.
抱著然後不停地.
來來來.
我覺得你沒有可能.
很小的.
那五個餅很小.
那兩條魚很小.
一眼就看完.
為甚麼會.
砰砰砰會怎樣呢.
我覺得完全是拍不到的.
我都看過不同的耶穌傳.
都是覺得不太行的.
但是耶穌就行了這個神蹟.
最後那句話是.
照他們所要的.
分給他們.
as much as they want it.
要多少有多少.
可以再回來.
再添.
all you can eat.
自助餐.
放題.
不可思議.
在我們還沒有去故事結局之前.

$^{1241}$我想你看看這幅圖畫.
約翰幫我們畫了一個.
屬靈的圖畫.
在這幅圖畫裡面.
他們本來是去一個曠野的地方.
但是無端端曠野裡面突然間.
有一個大草原.
好像沙漠裡面有一個綠洲.
有一個大草原.
然後.
有一個極其豐盛的顏值.
雖然我們吃來吃去都是大麥餅.
窮家人的東西.
如果你見到一個這樣的神蹟.
你吃什麼你都覺得很有味道.
你覺得這是什麼魔術食物.
不可能的東西.
極其豐富的顏值.
當然有些人可能突然間質量很好.
我們不知道.
還有另外一個就是.
用大麥餅餵飽很多人的神蹟.
所以我只是告訴你們.
這幅新約的圖畫.
兩節而已.
其實會讓人聯想到三幅舊約的圖畫.
三幅舊約的圖畫.
分別是在詩篇裡面的.
在以賽亞書的.
和列王記下的.
我想和你們很快看看這三段的聖經.
第一個.
我想大家一想都已經知道了.
詩篇23篇.
「若是我的牧者我必不自缺乏」.
「他使我躺我在青草地上」.
明明是荒蕪之地.
突然間青草地.
「令我在可安歇的水邊」.
你馬上想到耶穌.

$^{1281}$其實第四章已經是祂生命的活傳.
又有水又有草.
你現在是羊.
不用吃麵包.
「但青草地上又有水不自缺乏」.
然後你繼續下去你就知道了.
然後祂要做什麼.
「祂使我的靈魂甦醒」.
「為自己的名引導我走而路」.
「祂使我的靈魂甦醒」.
另一個翻譯是.
「祂使我的生命復活」.
使我復活重生.
然後我可以為祂的國祂的義.
去盡心竭力.
你明白嗎.
其實你開始明白耶穌在做什麼了.
在整個故事裡面.
我們看到耶穌就是那個好牧羊人.
我今天沒有時間跳到第十章去看約翰福音.
但是你知道第十章.
耶穌就說自己是好牧羊人.
「我是好牧羊人」.
「好牧羊人為羊寫命」.
祂不是單單青草地溪水旁.
祂是要為羊寫命第十章第十一節.
但同時之間.
如果你記得在約翰福音第十章裡面.
耶穌又加了一句.
「我的羊聽我的聲音」.
「我又認識他們」.
「他們會跟從我的羊」.
於是我們必須要問一個問題.
在約翰福音第六章五名異儒神跡故事裡面的當中.
二萬個男女.
是不是都是耶穌的羊.
我們要問這個問題.
因為現在祂有溪水旁.
又有青草地.
但問題不是.

$^{1321}$因為我們後來知道這群人要的是地上的國度.
而耶穌怕了他們.
他們來追祂的時候.
祂就上山了.
這群五千起二萬幾人.
男女.
全部都是天國劇場.
天國記號裡面的臨時演員.
隨便招募.
所以不要以為有份在聖經出現.
很厲害.
原來是臨時演員.
我們可以這樣說.
如果這群人要真的成為耶穌的羊.
他們必須要真心的迴轉叛改.
他們必須要聽耶穌的話.
不可以再看耶穌的神蹟.
還有得救的.
我不是說現在他們做了臨時演員.
就以後做臨時演員.
他可以成為真正的羊.
不過他一定要聽耶穌的話.
以一生順靠跟隨耶穌.
追求天上的良善.
我們是一樣的.
我們怎樣可以成為耶穌的羊.
我們要真心的迴轉.
我們要聽耶穌的話.
不要整天在想神蹟.
要在這裡聽說話.
一生的跟從.
這幅是第一幅舊約的圖畫.
詩篇的圖畫.
第二幅圖畫其實更厲害.
不過我們沒有那麼熟悉.
第二幅圖畫是以賽亞書二十五章.
醒學者一想到在《約翰福音》裡面.
走不掉的一定是以賽亞書二十五章.
以賽亞書二十五章講到有一個筵席.
上帝為他子民所預備的筵席.

$^{1361}$是一個滅絕死亡的筵席.
很誇張.
以賽亞書二十五章第六節.
在這山上.
要發生在山上.
耶穌說必為他萬民擺設筵席.
有肥金和美酒.
就是滿有骨髓的肥金和精釀的美酒.
在這山上.
我必吞滅那前果萬民的面紗.
和那遮蓋列國的遮蓋物.
他已吞滅死亡直到永遠.
這是以賽亞書裡面一個救贖筵席.
一個很偉大的意象.
這個筵席其實結合了加拿大婚宴水面酒.
和五餅二魚飽千人兩個神蹟.
當然他沒有特別講麵包.
但耶穌那個神蹟太誇張.
令大家吃得太歡喜.
還在山上.
他上了山 人們跟著來.
在那裡擺設了筵席.
他又是加拿大婚宴水面酒.
所以所有東西.
死亡者說你逃不掉.
就是耶穌自己說.
我是生命的糧.
到我這裡來的絕不飢餓.
第六章三十五節.
信我的永不乾旱.
我是生命的糧.
第六章第四十八節.
第五十節.
這是從天降下來的糧.
只人吃了就不死.
最緊要那個筵席吃了之後.
會吞滅死亡的筵席.
人若吃這糧必永遠活著.
那個人就是我自己的肉.
耶穌很明顯.

$^{1401}$所有東西讓我們看見.
這個很奇異的偉大的意象成全了.
在那個傍晚裡面有五千個男人.
二萬個人次.
他們這班人其實是天國的臨時演員.
那一晚他們真的有得吃.
吃了耶穌天國的筵席.
我想問問他們吃了之後.
是否可以出世入生.
沒有.
你明不明白.
雖然你吃了一個奇異的魔法食物.
譬如你來我們教會.
如果你不是遵循耶穌.
你偷了我的聖餐吃也沒有用.
你不要以為吃了聖餐你下星期運氣好些.
沒有用的.
你也不要以為你吃了四十年聖餐一定上天堂.
完全兩回事來的.
那班人那天吃了耶穌的筵席.
但是除非他們迴轉.
他們沒有得救.
但是如果他們真心的迴轉.
這個救贖的筵席是有份的.
如果真心迴轉.
如果我們聽耶穌的話記得嗎.
不是看耶穌的神蹟.
是聽耶穌的話.
一生信靠一生跟隨耶穌為上帝而活的時候.
他們就可以跟這個筵席連上關係.
這是第二幅舊約圖畫.
最後第三幅一定要有的.
跑不掉的.
最後第三幅的圖畫是獵王基夏.
神人.
以利沙的神蹟.
跟耶穌那個是一樣的.
獵王基夏.
四章第四十二節.
有一個人.

$^{1441}$帶著初熟的果子.
以及大麥做的餅和新麥水.
神人說把這些給大家吃.
當時有個饑荒.
仆人說.
這些豈可擺在一百人的面前.
二十個大麥餅.
一百人吃怎麼夠.
以利沙說.
你只管給大家吃吧.
仆人就擺在他們面前.
他們吃了.
還有剩下的.
於是他用二十個大麥餅.
餵飽了一百個人.
已經是神蹟了.
耶穌用五個大麥餅.
餵飽了五千個壯男.
還有二萬人.
何其偉大.
整個故事告訴你.
耶穌是誰.
耶穌是一個更偉大更美善的以利沙.
He is the greatest and better Elijah.
耶穌是上帝終極的先知.
以利沙是先知.
耶穌是終極的先知.
於是我們必須要聽他的說話.
如果你有機會.
可能下星期我回來.
我帶你看看以利沙.
這群吃餅的一百人.
不是隨便在街上找人.
誰肚子餓.
吃餅的是誰.
先知的門徒.
這裡說的是上帝.
當他呼召這群先知.
先知的門徒來跟隨他.
有饑荒的時候.

$^{1481}$上帝供應他們的需要.
為什麼.
先知在列王記下的作用是什麼.
跟那些偶像假神撐過.
其實你想想跟我們有什麼關係.
當然有關係.
耶穌說.
你當求他的國他的義.
其他一切都加給你.
我的小羊你不要擔心.
只要你跟隨我.
我必照應.
是一個很美麗的圖畫.
不過不要說了.
我們下個星期回來.
可能再有機會看.
但這裡還沒完.
因為20個大麥餅.
餵飽一百人之後還可以剩下.
問題是.
耶穌那裡五餅二魚.
餵飽了一千上萬人.
還可以剩下嗎.
當然可以剩下.
回到總結那裡.
第六章第十二節.
他們吃飽的時候.
耶穌對門徒說.
要把剩下的瑞士收拾起來.
他們就把五個大麥餅的瑞士.
大家吃剩的收拾起來.
裝滿了十二個男子.
正如我剛才說.
神蹟還沒完.
因為這十二個男子.
會變成新的五餅二魚.
要供養以色列人十二個支派.
十二個支派.
耶穌全部都能拯救.
拯救到底.

$^{1521}$你不要以為耶穌.
錯過了機會.
沒有吃那頓飯.
那頓飯根本不重要.
是影子今天我們可以得到的.
才是實體.
不過有條件.
我們必須要真心回轉.
聽耶穌的話一生跟從他.
最後兩節.
人們看見耶穌所行的神蹟.
就說這真是那要到世上來的先知.
然後第十五節.
耶穌知道他們要來強迫他作王.
就獨自又退到山上去了.
這真是那要到世上來的先知.
是什麼.
是摩西.
是申命記第十八章十八節.
我必在你們弟兄中.
給他們興起一個先知像你.
我要將當說的話放在他的口裡.
他要將我一切吩咐的話.
全告訴你們.
你們要聽他的.
你發覺這群人其實他們很聰明.
如果耶穌在這裡行了神蹟.
我們吃了飯.
我們說多謝耶穌.
我們不會想其他人.
這群人正在想聖經.
這傢伙可以這樣做.
他是誰.
他一定是那個先知.
但他沒有留意下面那句說話.
下面那句說話就是.
我要將我的話放在他的口.
他所說的你都要聽.
他們就是不聽耶穌.
下一句轉個面.

$^{1561}$就要強迫他做他不想做的東西.
所以他們答應了一半.
你不是吃了飯就算了.
你要想一想耶穌究竟是誰.
但他想了.
答應了一半.
接下來他們要的.
就是做自己眼中捍衛對的事.
要強迫耶穌作王.
我們知道耶穌來就是要作王.
耶穌就是上帝所揀選.
所認祂的救主君王.
然而耶穌堅決.
祂不要做地上的君王.
耶穌要做天國十字架上的君王.
這裡是他們為什麼要迫耶穌.
做這個君王的理由.
但我們去到楊福音的高峰.
彼拉多對耶穌說.
你是猶太人的皇媽.
第十八章第三十六節.
耶穌回答說.
我的國不屬於這個世界.
我的國若屬於這個世界.
我的部下將會為我爭鬥.
只是我的國不屬於這個世界.
第十九章第十四節.
哪日是逾越節的預備.
彼拉多說.
看那你們的王.
他們就喊著.
把他釘在十字架上.
除了凱撒以外.
我們沒有王.
我們要好像凱撒一樣的王.
我們要好像耶穌一樣的王.
我們要好像凱撒一樣的王.
我們要像猶太人一樣的王.
要幫助猶太人.
但我喜歡凱撒的權兵.

$^{1601}$甚至我們可以這樣說.
其實以色列人很恨自己有十字架在手.
只不過他現在要借刀殺人.
因為十字架屬於羅馬人.
但有一天.
如果以色列人可以統領世界.
他們保留著十字架的威力.
可以威嚇世人.
把他釘十字架.
除了凱撒以外我們沒有王.
於是彼拉多把耶穌.
交給他們去釘十字架.
耶穌是天上降下的生命麵包.
十字架是他人生的目標和使命.
但我們回到約翰福音最後.
第六章第五十一節.
我就是從天降下來的生命的糧.
吃我的必永遠活著.
我的糧就是我的肉.
然後耶穌跟我們說.
不要為那毀壞的食物操勞.
要為那傳道永生的食物操勞.
就是肥你終於考過試那一刻.
人子要賜給你們的就是永生的食物.
我們的時候到了.
我們下個星期回來.
我們會繼續思想二元聖經.
我想用三句很簡單的說話.
去總結今天的訊息.
第一句.
不是賜人食物而是成為食物.
耶穌來不是賜人食物而是成為食物.
第二.
耶穌來不是於人有用而是作人自保.
不是讓我們可以洗洗它.
而是要成為我們的自保.
第三.
耶穌來不是要滿足我們的慾望.
而是要改造我們的慾望.
要我們不再追求地上的食物.

$^{1641}$而是追求天上的食物.
願耶穌成為我們每日的力量.
我們日用的麵包今天賜給我們.
我們日用的飲食其實是麵包.
那個主圖文.
我們日用的麵包是甚麼.
是耶穌.
今天賜給我們.
賜給我們耶穌.
願耶穌成為我們每日的力量.
生命的自保.
願耶穌改造我們.
使我們從此以後為那全道永生的食物去操勞.
我們今日祈禱.
五餅二魚.
一個何等孤單的神跡.
耶穌一個人.
要頂住五千個男人.
還再加十二個他自己憤怒的門徒.
二萬群情洶湧的猶太平民.
原來他們和猶太領袖都是一樣.
他們都是在追求一個地上的國度.
現以耶穌來.
要給我們一個不一樣的國度.
求主你幫助我們.
使我們明白這個神跡記號真正的意義.
讓耶穌成為我們的天良.
讓耶穌成為我們的自保.
讓耶穌改造我們的慾望.
使我們不再追求地上的食物.
而是追求使我們全道永遠的食物.
為此而操勞.
多謝你.
我們期待下星期再次回來思想.
祈禱奉主耶穌 祈禱冥求 阿們.
\newpage



\section{哥林多後書 4:7}
\label{sec:031668}
\textbf{大能的勇士 (普/粵)}~~~~湯傳哲牧師
\newline
\newline
連結: \href{http://www.church.com.hk/acms/content.asp?site=cdc\&op=show\&type=product\&code=031668\&layout=sermon}{\texttt{031668}} ~~~~ 語音日期: 2026-02-20
\newline
\newline
\hyperref[sec:031667]{\small{< < < PREV SERMON < < <}}
~
\hyperref[sec:index_chronic]{\small{[返順時目]}}
~
\hyperref[sec:031669]{\small{> > > NEXT SERMON > > >}}
\newline
\newline
哥林多後書 4:7
\newline
\begin{longtable}{cl}
\hline
\hline
章節 & 經文 (和合本修訂版)\\
\hline
4:7 & \begin{tabularx}{0.7\textwidth}{X} 我們有這寶貝放在瓦器裡，為要顯明這莫大的能力是出於神，不是出於我們。 \end{tabularx} \\ \\
[1ex]
\hline
\hline
\end{longtable}
$^{1}$很高興能夠在今天和大家分享主的信息.
很高興今天和大家分享神的訊息.
今天要和大家講的題目是大能的勇士.
今天想和大家分享的題目是大能的勇士.
剛才我們已經讀過那段經文.
剛才已經讀過一段經文了.
等一下我們要繼續往下看.
我們先做一個禱告.
天父謝謝你帶領我們.
因為你讓我們開啟我們的眼睛.
讓我們從心裡能夠聽見你的聲音.
知道你在呼召我們.
讓我們都可以成為你大能的勇士.
這樣禱告奉耶穌的名.
我們今天看的這段經文講到.
以色列人他們遇到大困難.
今天經文看到以色列人遇到大困難.
我們有沒有困難呢.
我不知道你們有沒有感覺到有困難.
我確實感覺到有困難.
我就真的感覺到有困難.
因為物價在上漲.
因為物價不停上升.
或許像我這種依靠退休金生活的人.
最能夠感覺到困難.
尤其是我這種依靠退休金生活的人.
最能夠感覺到困難.
從前去飯店吃飯十塊錢就一道菜.
以前去餐廳點一個菜只是十塊錢.
現在要十五塊.
現在已經升到十五塊.
從前買一磅青耕菜要一塊錢.
以前在超市買一磅青耕菜可能只是一塊錢.
現在兩塊多有時候三塊.
現在已經接近兩塊甚至三塊.
所以實在有困難.
所以實在困難.
保羅他怎麼說呢.
保羅怎麼說.
他說我們是瓦契.

$^{41}$古時候瓦契是最便宜的東西.
因為你只要到田裡去抓一把土.
你只要去田裡面去餵一些泥土.
用水合在一起捏一捏.
扣一些水用手捏一捏.
然後放在爐裡燒一下就成了瓦契了.
放在爐裡面燒一燒瓦契就搞定了.
保羅說基督徒就不一樣.
保羅說基督徒不一樣.
因為我們這個瓦契裡面有寶貝.
因為我們瓦契裡面有寶貝.
我們可以依靠神.
可以克服逆境度過難關.
克服逆境度過難關.
但很多時候別人看不出我們裡面有寶貝.
有時候我們自己都忘記自己裡面有寶貝.
有時候連基督徒.
我們自己都忘記我們心裡面有寶貝.
我們追求的是外面的這些身份,地位,家世,學問,錢財.
我們追求的是身份,地位,家世,學問,土地.
我們覺得這些外面都是很寶貝.
我們覺得這些外在的東西比較寶貝.
可是保羅說我們只是個瓦契.
但保羅說我們是瓦契.
真正的寶貝是我們心裡的信仰.
真正的寶貝是我們心裡面的信仰.
是耶穌基督.
是我們裡面的耶穌基督.
在經文說要顯明這莫大的能力.
這段經文說要顯出這個莫大的能力.
是出於神,不是出於我們.
所以這個意思就是我們要把在裡面的寶貝把它活出來.
所以是要將裡面的寶貝活出來.
可是很多時候我們都覺得我就是這樣,你要接納我.
但有時候我們覺得我就是這樣,你要接納我.
那都是自己,老自己活出來.
都是活出我們的腦.
我們沒有活出基督的新生命.
所以保羅的話就讓我想起舊約的人物基甸.
保羅的話就讓我想起舊約的聖經人物基甸.

$^{81}$所以我們一起來讀這個《釋詩記》第六章一節.
我們一起來讀《釋詩記》六章一節.
以色列人與其黨.
又行耶和華眼中看為惡的事.
耶和華就把他們交在米殿人手裡.
你看以色列人遇到困難了.
以色列人遇到困難.
經文說他們是犯了罪.
經文說他們犯了罪.
行耶和華眼中看為惡的事.
行神眼中看為惡的事.
然後神就把他們交在米殿人手裡.
所以神就將他們交在米殿人手裡.
基督徒常常會受苦的.
基督徒其實經常會受苦的.
有些苦不一定是你做錯了什麼.
有些苦不是因為你做錯了什麼.
很多時候我們受苦是我們做錯了事.
但有些受苦的確是我們做錯了某些事.
現在好像我們物價上漲.
現在我們遇到物價上升.
主要的原因就是川普在打貿易戰.
主要原因是特朗普打貿易戰.
而在美國總統是大家選出來的.
在美國總統是大家選出來的.
所以現在我們受苦跟我們自己有關係.
所以我們現在受苦跟我們自己有關.
那麼基督徒受苦.
很多時候是因為生活見證不好.
基督徒受苦有時是我們自己生活見證不夠好.
你為什麼受苦.
因為有家人對你有壓力.
為什麼家人對你有壓力.
因為你見證不好.
你不能把你家人帶來信主.
你見證不好.
你帶不到你的家人來信主.
或者不是見證問題.
是你根本都忘記了我們有個使命.
或者不是見證不好.

$^{121}$你根本忘記了我們是有使命的.
所以我們應該要帶領家人信主.
我們要帶領家人信主.
我們應該要改變社區.
我們要改變社區.
就不會在社會上有這麼多壞人.
社會不會有這麼多壞人.
經文說起初以色列人不明白為什麼.
經文說以色列人起初不知道為什麼.
所以神就派先知去警告他們.
先知是什麼人呢.
先知是什麼人.
用先知這個名字好像他可以預告未來.
用先知這個名字好像他們可以預告未來.
其實牧師有的時候就是先知.
其實牧師有時候就是先知.
牧師常常講到.
牧師時常講到.
說如果你不按照神的話去做.
如果你不按照神的話去做.
會有一些結果的.
先知不是算命.
牧師也不是會懂得未來.
牧師也不知道未來.
可是我們把神的信息傳給人.
可是我們把神的信息傳開去.
裡面就講到未來.
裡面就有講到未來.
我們看看這個先知到以色列人傳的什麼信息.
我們看看這個傳給以色列人是什麼信息.
我們一起來讀.
我是耶和華里門的神.
你們住在阿摩利人的地.
不可敬畏他們的神.
你們竟不聽從我的話.
以色列人犯了什麼罪呢.
其實這項罪一直都跟著他們.
其實這項罪一直跟著他們.
因為他們就是說自己一方面是神的子民.
他們一方面說自己是神的子民.

$^{161}$他們要跟隨神.
要跟隨神.
可是他們又拜巴黎.
但一方面又去拜巴黎.
跟那個世界上人沒什麼不同.
跟世界其他人沒什麼分別.
別人看不出他們是神的子民.
在你的生活裡有沒有這種情況呢.
神要你誠誠實實做人.
神要你誠誠實實的做人.
但是當金錢誘惑你的時候.
你的原則不定放在那裡.
但是金錢引誘你的時候.
你可能就放棄原則.
你看美國現在最近在整頓移民法.
最近美國在整頓移民法.
因為很多有居留權的人.
很多拿了居留權的人.
感覺到很大的壓力.
有壓力了.
因為很多人拿了居留權沒有住在美國.
因為很多人拿了居留權.
但他不住在美國.
一大半的時間都在自己原居地.
一大半的時間都住在原居地.
他們有不同的目的.
有些人是打算將來老的時候.
享受美國的福利.
有的是想著老的時候在美國享受福利.
有的人是為逃稅.
有的人為了逃稅.
還有其他原因.
有其他原因.
聽說中國有很多高幹把自己的孩子送到美國來.
連習近平的孩子也在美國.
連習主席的女兒也在美國.
所以現在川普說.
在美國作威作福.
把自己的女兒兒子送到美國來.
這才是我們不歡迎.

$^{201}$所以特朗普說.
這些在國內作威作福.
欺壓老百姓.
將子女送來美國.
所以華人中間有很多人有壓力了.
所以很多人覺得有壓力了.
或許你會說我沒有做這種事.
可能你會說.
但我沒有做這些事.
我沒有出賣我的良心.
但是你心裡是不是也在等候.
如果價錢高的話.
你就願意犧牲了.
如果價錢夠高的話.
我是不是就願意犧牲.
當時以色列人怎麼樣呢.
當時以色列人怎麼樣.
摩西帶以色列人出埃及.
摩西帶以色列人出埃及後.
約瑟爾帶領他們入迦南.
約瑟爾帶他們入迦南.
以色列人原來是遊牧民族.
以色列人本來是遊牧民族.
不懂得怎麼耕地.
不懂得耕種.
但是要定居在巴勒斯汀.
他們要學種地.
他們就要學耕種.
當地的居民有很多經驗.
當地的居民有很多耕種.
所以他們應該去學.
他們應該去學.
但是當地居民.
是把跟巴力崇拜.
和耕田連在一起.
但是當地的居民是拜巴力.
和耕田是連在一起的.
如果要田地出產豐富就要拜巴力.
如果要出產豐富就一定要去拜巴力.
以色列人要學種地.

$^{241}$就要跟著當地人一起去拜巴力.
以色列人如果要去學耕種.
就要去跟著拜巴力.
有一次我看在華盛頓看報紙.
有一次我看華盛頓報紙.
裡面記載一個房地產經紀人.
他裡面記載了一個房地產經紀人的文章.
他說我是不跟隨風水的.
他說我不看風水的.
但是風水就是一般選屋的常識.
其實風水就是一般選屋的常識.
和一般的住家原則沒什麼不同.
和一般住家原則沒什麼分別的.
可是他把算命跟風水連在一起.
但是他一旦把算命風水放在一起.
變成迷信了.
那就變成迷信了.
看起來這個經紀人很有見識.
看這個經紀人好像很有見識.
但是他後來就說有一個白人找他買房子.
但他跟著又說有一個白人找他買屋.
看見他是華人就問他.
這個房子風水怎麼樣.
見到他是華人就問他.
你覺得風水怎麼樣這間屋.
這個經紀人說這個房子你喜歡.
經紀人就說你看這間屋你喜歡的話.
空氣好陽光充足.
房間安排都合乎常理.
房間的安排都很好的話.
所以風水好的.
那就是風水好的.
他沒有澄清自己不相信風水.
但他沒有澄清他自己不相信風水.
可是別人可不以為這樣.
但別人不以為是這樣.
所以覺得要看風水的.
連外國人都這樣.
所以就覺得要看風水.
連外國人都是這樣.

$^{281}$那麼以色列人如果是想要豐收的話.
那以色列人如果要豐收的話.
那一定要按著當地人的方法.
他們做什麼呢.
他們就拜豐收之神.
以色列人為了種地.
以色列人為了耕種.
原來跟當地人學是沒有什麼好講的.
本來跟當地人做就沒什麼好講的.
可是耕種常識跟拜巴黎連在一起.
但是跟拜巴黎聯在一起.
這個信仰就受了打折扣了.
這個信仰就打了折扣了.
因為巴黎的信仰是跟淫亂連在一起的.
因為巴黎的信仰是跟淫亂一起的.
如果以色列人要去拜巴黎.
女的就要去扮亞斯塔魯女神.
女的就要去扮這個女神亞斯塔魯.
在當地廟裡做妙計.
在廟裡面做妓女的.
這個男的就去扮巴黎.
那麼男的就要去扮巴黎.
去廟裡面和妙計做愛.
去廟裡面和妓女進行性行為.
有沒有用呢?.
有沒有用?.
有啊!很多女的.
以色列人不會生孩子.
很多當地人是生不了小朋友的.
所以到廟裡去做妙計.
真的懷孕了.
去了廟裡面又真的懷孕了.
為什麼?因為可能那個.
不會懷孕的原因是男的不會懷孕.
可能原因是男的不行.
所以在人性中最弱的.
最脆弱的弱點.
人性最脆弱的弱點.
以色列人就有了藉口.
以色列人有了藉口.

$^{321}$這是為了學耕地.
這個就是學耕種.
去廟裡面過夜就好了.
去廟裡面過夜不就行了.
戴上面具扮一下假神.
誰都不知道誰是誰.
戴面具扮一下假神.
你都不知道是誰.
你還可以跟不同的男的女的去做愛.
還可以跟不同男的女的發生關係.
所以有很多以色列人就去犯罪了.
所以很多以色列人就去犯罪了.
說不定你不拜這些淫亂的神.
你不拜這些淫亂的神.
或者你不看那個馬年運程.
你不去看馬年的運程.
但是你的人生目標和不和.
是不是和那個拜偶像的人差不多.
但是你人生的目標.
是不是和拜偶像的人差不多.
是不是你把耶穌當菩薩來拜.
你有沒有當耶穌做菩薩這樣去拜.
你在心裡面想的是有求必應.
你在想有求必應.
就好像我們來美國念書.
就好像我們來美國讀書.
畢業以後就找工作.
畢業之後我們去找工作.
你就禱告神求神給你一個好工作.
去禱告求神給好工作.
然後留下來.
有了工作就申請居留權.
找到工作就留下來申請居留權.
那你心裡很有感恩.
說這是神的安排.
你很感恩.
你說這是神的安排.
我要事奉主.
我要愛主.
我現在要事奉主.

$^{361}$我要愛主.
我要為主而活.
事實上我們開始看世界上的物質享受.
但我們開始看世界的物質享受.
然後我們就買大房子.
我們買大房子.
買大樓大屋.
買大車子.
買大車.
買大電視.
我們看電視.
看股票.
看股票上落.
現在加一項看手機.
看手機.
互聯網.
看互聯網.
所以我們出賣我們的時間.
我們出賣了我們的時間.
精力.
然後我們說沒有時間讀經了.
接著我們說.
沒有時間讀經.
沒有時間禱告.
沒有時間祈禱.
不能去聚會了.
沒有時間去聚會.
沒時間關心別人了.
沒有時間關心人家.
我們不敢愛主.
不敢付出.
怕事奉太熱心被人家笑.
怕太熱心去事奉會被人家笑.
有的時候我們想突破.
我們想突破.
可是我們突破不了那個固定的生活模式.
但你突破不了這個固定的生活模式.
當一個信耶穌的人.
你的信仰出現危機.
當一個信耶穌的人.

$^{401}$你的信仰發生危機的時候.
神一定會出來糾正的.
神一定會出來糾正.
於是就有米電人來了.
所以就米電人就出現了.
就像現在川普突然間搞一個貿易戰.
聖經說這些米電人像蝗蟲一樣.
聖經說這些米電人好像蝗蟲一樣.
把以色列人的豐收都吃了.
把以色列人的豐收全部吃了.
以色列人怕米電人.
怕得不得了.
以色列人很怕這些米電人.
基甸就是怕米電人.
基甸就是怕這些米電人.
天使顯現的時候基甸在做什麼.
聖經說它在酒渣裡面打麥子.
聖經說基甸在酒渣裡面打麥子.
聖經說因為怕米電人.
有沒有想過在酒渣裡打麥子.
跟怕米電人有什麼關係呢.
你有沒有想過在酒渣裡面打麥子.
跟怕米電人有什麼關係.
原來是這樣的.
這個以色列人種麥子種葡萄.
以色列人種麥子種葡萄.
在收麥子的時候.
他們在曬古場上打麥子.
他們打麥子的時候.
就要去空曠地方去打這些麥子.
在收葡萄的時候.
放在酒渣裡面去採.
把它榨出水.
所以葡萄的時候就去酒渣裡面.
然後就採碎了以後.
就葡萄汁流出來.
他們去釀酒.
採碎之後葡萄汁就出來.
就去釀酒了.
可是基甸怕米電人.

$^{441}$但他們怕米電人.
所以他就跑到酒渣裡面去打麥子.
所以基甸就走去酒渣裡面去打麥子.
這米電人來的時候.
只想到在空曠的賣場上有麥子.
米電人來的時候.
他只會去空曠那個麥場那裡找.
可是他沒有想到.
米電是躲在酒渣裡面打麥子.
沒想到基甸就走去酒渣裡面去打麥子.
所以米電人.
他為什麼在酒渣打麥子.
因為怕米電人.
他沒辦法.
那為什麼基甸走去酒渣那裡打麥子呢.
是因為他怕這些米電人.
你看接下來他說這個話.
二十六節我們一起來讀.
二十六節我們一起讀.
於是基甸在那裡為耶和華築了一座壇.
起名叫耶和華沙龍.
就是耶和華賜平安的意思.
這壇在亞比爾這族的羅佛拉直到如今.
所以基甸為什麼建一座壇.
叫做耶和華賜平安.
所以基甸建了這座壇叫做耶和華賜平安.
為什麼起這個名字呢.
為什麼要起這個名呢.
起名都有原因的.
起名都是有目的的.
就是因為他心裡面期望平安.
就是他心裡面期望平安.
基督徒一定聽過基甸的故事.
你們一定都聽過基甸的故事.
我們知道聖經裡的基甸是個大能的勇士.
我們知道聖經記錄這個是大能的勇士.
他曾經拯救以色列人脫離米電人的逼迫.
他曾經拯救這些以色列人去脫離米電人的逼迫.
可是二十六節說他心裡期望耶和華賜平安.
第二十六節說他心裡面卻希望耶和華賜平安.

$^{481}$你是不是也期望耶和華賜平安.
基甸生活環境裡四周都是戰爭.
基甸生活裡面四周都是爭戰.
這個米電人來的時候帶著他的駱駝牛羊.
米電人來的時候帶著牛羊駱駝.
好像是蝗蟲一樣把以色列人的收成都吃掉.
好像蝗蟲一樣把以色列人的收成全部吃掉.
吃不完就帶走.
以色列人就躲到山上去.
以色列人就逼到要躲到山上.
這一群以色列人.
從前有神的帶領.
離開埃及的時候.
聖經說連狗都不敢向他們咬舌.
連狗都不敢向他們咬舌.
但是他們墮落了.
他們墮落了.
米電人成了他們的剋星.
這些米電人成為他們的殺星.
他們見到米電人就躲起來.
所以他們基甸心裡是充滿了懼怕.
所以他們心裡是充滿了懼怕.
我們的生活環境.
也是不是充滿了威脅,懼怕.
我們的生活環境是不是充滿了威脅,懼怕.
去年我去台灣.
台灣那個時候.
大陸天天跑到那附近去做飛彈試射區.
大陸天天去那裡搞飛機的威脅.
然後就說這個地方要打飛彈.
你不可以經過這裡.
這個地方要打飛彈你不可以經過這裡.
然後就把台灣四周都包圍起來.
就把台灣四周都包圍起來.
所以我們就像基甸一樣.
我們就好像基甸一樣.
我們渴望平安.
現在我們面對的敵人.
不只是像中國大陸那樣.
我們現在面對的不只是像中國大陸那樣.

$^{521}$家裡有矛盾.
我們家裡也有矛盾.
子女不聽話.
父母相處不和.
父母相處不好.
那個婆媳中間有問題.
婆媳也都有問題.
常常有衝突.
常常都有衝突.
所以我們不平安.
約翰福音第十四章二十七節.
記載耶穌的話.
約翰福音十四章二十七節記錄了這一段.
那我們一起來讀.
我們一起來讀.
我留下平安給你們.
我將我的平安賜給你們.
我所賜的不像世人所賜的.
你們心裡不要憂愁.
也不要膽怯.
所以耶穌應許門徒.
這不只是一個應許.
耶穌復活以後.
耶穌復活之後.
聖經說門徒把門都關了.
聖經記載門徒都將門關上.
因為怕猶太人.
可是耶穌向他們顯現的時候.
第一句話就是說願你們平安.
耶穌向他們顯現的時候.
第一句話就是願你們平安.
今天你有沒有平安.
你應該有平安.
你是應該有平安.
因為我們這個瓦器裡面有寶貝.
有耶穌基督.
你應該說在我裡面有平安.
你應該時常把你心裡的寶貝顯露出來.
你要時時讓這個寶貝顯露出來.
在你苦難中間還可以享受平安.

$^{561}$在你苦難中你仍然享受平安.
讓我們看看這個天使怎麼對機電的話.
怎麼說.
我們看一下天使怎麼對機電說.
你看它這個十六節第十二.
六章第十二到十六節.
第六章十二至十六節.
我們先一起讀.
我們一起讀.
大能的勇士啊.
耶和華與你同在.
你可以繼續看下面.
我們繼續讀.
機電說 主啊 耶和華若與我們同在.
我們何至遭遇這一切事呢.
我們的列祖不是向我們說.
耶和華領我們從埃及上來嗎.
還要繼續.
他那樣奇妙的作為在哪裡呢.
現在他卻叼起我們.
將我們交在機電人手裡.
所以機電和以色列人都受盡了米電人的苦.
可是天使對機電說.
大能的勇士.
天使對機電說.
大能的勇士.
這是很大的諷刺.
你有沒有什麼煩惱.
你有沒有煩惱啊最近.
最近的股票市場有的時候上有的時候下.
現在市場上上落落.
你是不是也心裡跟著七上八下.
你看股市是不是心都七上八下.
那個COVID肺炎好像已經過去了.
COVID肺炎好像已經過去了.
可是我前不久聽到有人染上了COVID.
或許你很小心.
你每次出去都戴著口罩.
你出去會戴口罩.
我都看見你們中間有人戴口罩.

$^{601}$見到有人都戴口罩.
而且人多的地方就不去.
人多的地方你都不去.
可是還是有人染上COVID.
現在那麼多年.
但是現在經過幾年之後還是有人染上COVID.
聖經說.
平安了.
誰與你同在.
聖經說平安.
神與你同在.
你是基督的精兵.
你是大能的勇士.
是大能的勇士.
你會想到底有沒有搞錯.
你在想有沒有搞錯.
跟我講.
你跟我說話嗎.
如果我是大能勇士.
為什麼那麼多不平安.
如果我是大能的勇士.
為什麼我還有那麼多不平安.
還有互聯網上癮的問題越來越嚴重.
還有上網這個問題.
你上什麼網.
你上的網是什麼網.
你的配偶上什麼網.
你的另外一半上網是什麼網.
你的兒女上什麼網.
你的子女上什麼網.
很多時候上網的情況很嚴重.
其實很多時候上網的問題非常嚴重.
可是我們都不敢說.
但我們不敢說.
因為非常羞辱神的命.
因為很羞辱神的命.
然而聖經說.
你是大能的勇士.
但是聖經說你是大能的勇士.
你心裡就有大問題了.

$^{641}$你心裡面就有大問題了.
我這麼軟弱.
為什麼我是大能的勇士.
我這麼軟弱.
我怎麼會是大能的勇士.
可是你看下面一節怎麼說.
我們看下一節怎麼說.
一起來讀.
你靠著你這能力.
去從米殿人手裡拯救以色列人.
不是我差遣你去的嗎.
這個話是什麼意思.
這句話是什麼意思.
天使說你靠著你這能力.
天使說你靠著你這個能力.
你有沒有看英文聖經講得非常清楚了.
如果你看英文聖經.
這句話更加清楚.
天使不是要基甸到什麼地方去得能力.
天使不是要基甸去哪裡拿能力.
他說你就靠著你已經有的能力.
他說你靠著你已經得到的能力.
就可以打敗米殿人.
這就是今天我們的主題.
這就是我們今天的主題.
我們都是瓦器.
看起來不起眼沒什麼力量.
看起來不起眼沒什麼能力.
但是在我們裡面有寶貝.
而且我們應該有平安.
還有能力.
有能力.
我們裡面有寶貝.
就像基甸一樣.
就好像基甸一樣.
他不知道自己有能力.
他怕敵人怕得要命.
他怕敵人怕得要死.
他要躲起來打麥子.
他甚至要躲起來打麥子.

$^{681}$可是上天使告訴他你有能力.
但是天使告訴他你有能力.
不是基甸有能力.
不單止基甸有能力.
你有能力我有能力.
你都有能力我有能力.
凡是相信耶穌的人都有能力.
凡是接受相信耶穌的人就有能力.
這就是瓦器裡的寶貝.
這就是我們瓦器裡面的寶貝.
天使對基甸的話很簡單.
天使對基甸的話其實很簡單.
只要你願意用你的能力.
只要你用肯用你這個能力.
按照神的命令去行就好了.
按照神的命令去做就行了.
不管你的問題有多大.
神的能力可以幫助你勝過逆境.
能夠應付你一切頭痛的事.
可以幫你應付一切頭痛的事.
也幫助你去傳福音.
也可以幫助你去傳福音.
你不必害怕.
只管把你心裡的寶貝顯露出來.
你只要將心裡面的寶貝顯露出來.
問題怎麼解決呢.
問題要怎麼解決呢.
首先就要像基甸一樣.
首先你要像基甸一樣.
承認自己很絕望.
承認我不能處理這個問題.
承認自己處理不到這個問題.
然後我們要相信有一個神.
我們要相信我們有一個神.
我們如果把問題帶到神面前.
如果將這個問題帶到神的面前.
祂的能力比我的大.
祂的能力比我大.
一定能夠解決我們的問題.
一定可以解決我們的問題.

$^{721}$再下來就是要委身依靠神.
接著就是要委身依靠神.
願意把自己交託在神的手中.
你要把你的問題清楚明白地寫下來.
你要把你的問題清楚寫下來.
願意把這個問題交給神.
然後找一個帶導的同伴.
找一個帶導的同伴.
互相守望 彼此監督.
互相去守望 彼此監督.
清楚明白來到神的面前.
清楚明白的來到神的面前.
求神來幫助.
求神來幫助我們.
可是我們現在的情況.
很多人不敢承認自己有問題.
很多人甚至不敢面對自己.
承認自己的問題.
知道有問題不願意承認.
知道問題但不敢承認.
那更覺得就是帶到神的面前.
神也沒辦法.
更加沒辦法去帶到神的面前來.
不相信神可以幫助你.
不敢相信神會幫助我們.
然後你必須把你心裡的寶貝顯露出來.
你必須把我們心裡的寶貝.
耶穌基督顯出來.
大能的力量才會出來.
大能的能力才可以顯出來.
什麼是攔阻基督徒.
把自己的力量埋沒呢.
什麼攔阻我們基督徒.
把我們已有的能力埋沒了呢.
有一次有人跟我談起信仰的問題.
有一次我跟人談起信仰的問題.
他說現實生活非常殘酷.
他說現實的生活是殘酷的.
我說了低頭沒有辦法.
我除了低頭我沒有其他辦法.

$^{761}$我們好像以色列人遇到哥利亞.
他是身高八尺.
他身高八尺.
全身都穿在盔甲裡.
全身都披在盔甲裡面.
他和我們罵陣.
我們一起讀罵陣的話.
我們一起讀這一句.
我今日以以色列人的軍隊罵陣.
你們叫一個人出來與我戰鬥.
誰敢出去.
沒有.
一個人都不敢出去.
沒有人敢出來.
然後大衛沒有盔甲.
大衛連盔甲都沒有.
本來嫂羅說我的盔甲給你戴.
嫂羅說我給盔甲你戴.
結果他一試說太大不能穿.
誰知一試太大穿不到.
他就穿著牧羊的衣服.
他只是穿著牧羊的衣服.
他連刀都沒有.
刀都沒有.
他就到河邊撿了五塊石頭.
只是去河邊撿了五塊石.
就靠他平常打獅子老虎用的彈弓.
就拿了他平常打老虎的彈弓.
然後就靠這五塊石頭.
就把哥利亞打死.
就靠這五塊石頭.
就將巨人哥利亞打死.
他靠著信心.
靠他會甩石頭的能力就得勝.
他靠著這個信心.
靠著他這個會扔石頭的能力就得勝.
就好像天使對姬甸說的.
正如天使對姬甸說.
你靠著你已經有的能力.
去打敗米殿人.

$^{801}$你靠著你已經有的能力.
去打敗米殿人.
這個今天耶穌已經復活了.
今天耶穌已經復活了.
耶穌復活的大能在我們身上.
每一個人你都體會過.
耶穌復活的大能.
你們每一個人應該都體會過.
耶穌復活的大能.
要不然你不會坐在這裡.
我們有改變的經驗.
所以我們悔改相信耶穌.
跟隨耶穌加入教會.
所以我們願意去悔改.
相信耶穌跟隨耶穌進入教會.
每次聽到講道.
你心裡都會受感動.
每次講道你心裡都有感動.
你會決志全心跟隨主.
美國有很多好講員.
每次講道我們聽的時候.
都會受感動.
我們每次聽道我們都受感動.
就把自己獻給主.
但是你回到家裡面去.
但是當你回到家裡的時候.
還沒有改變.
沒有改變.
覺得無力.
覺得靈性枯乾.
你的問題在哪裡.
經典的故事告訴我們.
我們的問題在覺得手中能力不夠.
我們的問題是.
我們覺得手裡的能力不夠.
所以我們不能得勝.
只好妥協.
所以我們不能得勝.
我們只好妥協.
你看另外一段聖經怎麼講的.

$^{841}$另外一段聖經.
我們一起來讀.
阿摩斯八章十一節.
人饑餓非因無餅.
乾渴非因無水.
乃因不聽耶和華的話.
在生活中間我們為什麼妥協呢.
在生活中我們為什麼會妥協.
我們覺得很軟弱沒有能力.
我們覺得軟弱沒有能力.
你有沒有試過用你心中的寶貝.
你有沒有試過用心裡面的寶貝.
如果你不用.
你怎麼知道你不能得勝.
你怎麼知道你不能得勝呢.
天使對機電說.
你靠著你有的能力.
就可以打敗米電人.
你靠著你已經有的能力.
就可以打敗米電人.
你的米電人是誰呢.
你的米電人是誰.
說不定夫妻相處就是你的另一半.
可能就是夫妻相處.
你的另外一半.
兒女管家就是你的兒女.
或者你的子女管教的問題.
這個弟兄姊妹相處.
就是你家裡的人.
弟兄姊妹相處.
你家裡的人.
這個婆媳相處.
就是你的婆婆.
你的媳婦.
就是婆媳相處.
你的婆婆.
你的媳婦.
然後再加上經濟壓力.
還有經濟壓力.
有的時候失業.

$^{881}$有時候失業.
有的時候物價上升.
物價上升.
股票跌價.
股票跌.
你的個性情緒都受到挑戰.
你的個性你的情緒都受到挑戰.
你要讓你裡面的寶貝.
基督的生命顯出來.
你要讓你心裡面這個寶貝.
耶穌基督的新生命去活出來.
你要信服神的話.
你的話.
你的能力就會顯出來.
你的能力就可以顯出來.
要用你自己已經有的能力.
你要用你已經有的能力.
一定會勝過你生命中的米電人.
一定可以勝過你的米電人.
生命中的米電人.
天使向基甸顯現.
告訴他神與你同在.
你用你已經有的能力.
去拯救以色列人.
你用你已經有的能力.
去拯救以色列人.
因為基甸怎麼辦呢.
那麼基甸怎麼做.
你看我們讀下面的話.
我們讀17至18節.
基甸說.
我若在你眼前蒙恩.
求你給我一個證據.
使我知道與你說話的就是主.
記下去繼續.
繼續.
求你不要離開這裡.
等我歸回.
將禮物帶來供在你面前.
主說我必等你回來.

$^{921}$於是基甸去預備了禮物.
然後基甸就去預備祭物.
這裡預備了這個祭物.
不是專門普通用來獻祭的東西.
祭物這個詞.
不是完全專用來獻祭的詞.
也不是待客的一個餐點.
也都不是一個待客的大餐.
而是一個犧牲的禮物.
是一份犧牲的禮物.
什麼是犧牲禮物呢.
什麼是犧牲的禮物啊.
就是對基甸來說是很大的犧牲.
對基甸來說是一個很大的犧牲.
他做了什麼呢.
他做了什麼.
他把一隻羊羔殺了.
他將一隻羊羔殺了.
用細面做無效餅.
用細面做無膠餅.
一隻羊羔多少重啊.
多重.
一隻羊羔幾重.
你平常買一磅羊肉多少錢啊.
平時買一磅羊肉多少錢啊.
他就是那麼多的.
那麼大的一個犧牲.
他就是那麼大的犧牲.
有一次解經學者說.
基甸為什麼要用這個祭物呢.
有解經學者說.
為什麼要用這個祭物呢.
他是要認識這個人.
到底是人還是天使.
他要認識這個人.
到底是一個人.
還是來自神的使者.
他是有一個很大的犧牲.
他有一個很大的犧牲.
他是自己.

$^{961}$人這方面要犧牲.
他人這方面自己要犧牲.
然後另外要看看.
神的使者怎麼樣接納.
另外一方面要看神的使者.
怎麼樣去接納.
如果基甸犧牲獻上.
這個使者就像普通人吃飯一樣.
把他獻的吃掉了.
如果使者像普通人吃飯一樣吃了.
那麼他就知道這個是一個普通的人.
他就會知道這個是一個普通的人.
不能夠顯出是神的使者.
不能顯出是神的使者.
他想不到這個使者竟然.
用火燒了他獻上的這個祭物.
想不到這個使者就用火燒了他這個祭物.
基甸什麼反應呢.
基甸有什麼反應啊.
我們看下一段神經.
一起讀.
基甸見他是耶和華的使者.
就說:哀哉主耶和華.
我不好了.
因為我敵面看見耶和華的使者.
我們和基甸一樣.
盼望平安.
我們好像基甸一樣.
我們盼望平安.
我們覺得自己軟弱沒有力量.
我們覺得自己軟弱沒有能力.
所以我們跟基甸有很多相同的地方.
神經特別強調.
還有一點跟我們一樣.
聖經特別強調.
有一樣東西跟我們一樣.
就是基甸是小信的人.
神告訴他.
你就靠著你的能力去拯救以色列人.
神就跟他說.

$^{1001}$你靠著你的能力去拯救以色列人.
不是我猜你去的嗎.
是我猜你去的.
神的話包括了很大的使命和應許.
神的話包括了重大的使命和應許.
因為基甸是小信的人.
他一直不相信神跟他同在.
你看看在第六章十三節.
他說什麼.
你看十三節.
耶和華若與我們同在.
我們何至遭遇這一切事呢.
他為什麼說若.
因為他不知道有沒有.
他不知道有沒有.
然後再看六章十七節.
再看十七節.
我若在你眼前蒙恩.
求你給我一個證據.
使我知道與我說話的就是主.
所以他真的是不信.
不敢相信.
他真的不敢相信.
然後第三次你又看他怎麼講.
第三次你看他怎麼講.
你若果聚著所說的話.
我手拯救以色列人.
我就把一團羊毛放在和牆上.
這就是我們很有名的.
這個辨屍辨肝的應許.
這個就是一個很有名的羊毛辨屍辨肝的應許.
從來沒有一個牧師會說.
你回家去試試看.
放一塊衛生紙在外面.
天天看是濕的還是乾的.
從來沒有牧師說.
你將在外面擺一塊衛生紙.
看它有沒有肝還是不肝.
尋求神的旨意的方法都錯了.
尋求神的旨意方法都錯了.

$^{1041}$但是他是要清楚知道.
是不是真的神在差遣他.
但是他想清楚明白.
是不是神親自差遣他.
所以他預備一份犧牲的禮物.
看看神用什麼方法來接納.
你有沒有這樣去試試看.
獻上一個很大犧牲的禮物.
看看神交給你的使命是不是從神來的.
這個基甸是確定了.
他知道是從神來的.
他就遵著神的話去做了.
他就遵從神的話去做.
四十記第六章二十五到二十七節說他做了什麼.
二十五到二十七節說他做了什麼.
我不再讀了.
就是簡單說.
就是簡單和我們說.
神要他再次獻上犧牲的禮物.
然後要他拆毀巴黎的祭壇.
要他拆毀巴黎的壇.
要車下木偶做柴.
要拆下這些木偶做柴.
然後將他的犧牲獻在這樣的壇上.
將他犧牲的禮物燒在壇上.
你看他以前就拜過巴黎.
他以前就是拜過巴黎.
就是他犯過奸淫.
跟別人一樣.
他犯過奸淫.
跟其他人一樣.
有一次我帶小組查經.
查考聖經人物約瑟.
約瑟在面對試探的時候有一個問題.
約瑟在面對試探的時候.
他有一個問題.
他到底要信服神呢.
還是信服在環境的引誘之下.
他要信服神呢.
還是信服環境的引誘呢.

$^{1081}$然後我就問說.
在我們的環境中.
你有什麼困難呢.
那我就問你們環境有什麼問題.
就好像約瑟面對的試探.
就好像我們面對約瑟的試探.
所以有個弟兄說.
有個弟兄就說了.
他是警察.
他出勤回來要報出勤多久.
他出勤之後.
他回來要報出勤的時間.
他的同事都是以少報多.
那些同事當然是以少報多.
好比出勤兩小時報六小時.
出勤兩個小時.
但是報了六個小時.
他這樣講了以後.
馬上有個另外的弟兄不是警察.
也這麼說.
另外一個弟兄不是警察.
但是都站出來說.
他說我出差的時候.
一定要以少報多.
我出差的時候.
都是要以少報多.
如果我不多報的話.
我就賺的錢很少.
我不報多.
那我賺的錢就少了.
而且我的同事都一樣的.
大家都這麼做.
我同事都是這樣.
大家都是這樣一起做的.
那他說怎麼辦呢.
那怎麼辦呢.
他說只有妥協.
沒有別的辦法.
只可以妥協.
沒有其他辦法.

$^{1121}$現在你查經了.
聖經教導我們不能妥協.
我們現在查經.
耶穌說叫我們不可以妥協.
你先要試試看.
是不是神的心意.
你要知道是不是神的主意.
如果是神在指示你的話.
如果是神指示你的話.
你就要去聽.
如果是.
大部分的人都是這樣說.
這是你牧師個人的意見.
但大部分人說.
這些只是你牧師的個人意見.
我有選擇權.
我可以不聽.
我有選擇權.
我可以選擇不聽.
你就沒有辨別清楚.
到底是不是神要你這麼做.
所以你根本沒有分清楚.
誰要你這麼做.
祭奠是禱告以後搞清楚了.
祭奠的情況是祈禱之後搞清楚了.
然後他就改變.
他去改變.
他配合神的要求.
他要配合神的要求.
其實神的意思太清楚了.
他要我們抵擋魔鬼.
絕對不妥協.
為什麼我們要妥協呢.
為什麼我們要妥協.
因為我們沒有信心.
我們怕同伴的壓力.
我們怕其他人的壓力.
我們怕有損失.
我們怕損失.
我們怕和別人不同會有壓力.

$^{1161}$我們怕和其他人不同.
我們會有壓力.
而祭奠就按照神的命令去做.
但基連他按照神的命令去做.
他是冒著死亡的危險.
他冒著死亡的危險.
他是憑信心去拆毀巴黎的祭壇.
他憑著信心去拆毀巴黎的壇.
對他來講也是非常不容易的事.
其實對他來講都是很不容易.
他選擇晚上做.
所以他選擇晚上去做.
第二天就有人發現了.
但第二天還是有人發現了.
就對祭奠的父親約阿施說這個話.
就對祭奠的爹地約阿施說.
你看我們一起讀.
三十節.
將你兒子交出來.
好賜死他.
因為他拆毀了巴黎的壇.
砍下壇旁的木偶.
你願不願意付代價.
去實現神的心意呢.
你願不願意付代價實現神的心意呢.
今天在美國實現神的心意一點都不難.
現在其實在美國.
要實現神的命令是不難的.
你不必付什麼生命的代價.
你不需要付上生命的代價.
很多人按照聖經的話去行事.
在社會上還有好名聲.
有很多人憑著聖經的原則去做事.
還會有好的人望.
很多人.
我以前見過這樣的執事.
拼命爭取要當執事.
我見過以前有些執事.
拼了命都要去當執事.
為什麼呢.

$^{1201}$因為他出去以後告訴別人說.
我是某某教會的執事.
他做了之後.
他出去跟別人說.
我是某某教會的執事.
美國人看他另眼相看.
美國人就會對他另眼相看.
但願我們事奉是單純的動機.
但願我們事奉服侍的動機是單純.
我們始終如一付代價事奉主.
我們始終如一付代價去事奉主.
因為你遵守神的命令.
是要得神的喜悅.
因為我們遵守神的命令.
是要得神的喜悅.
幾點遵守神的命令.
幾點去遵守神的命令.
連生命的代價他都願意付.
連生命的代價他都願意付出.
結果呢.
他發現沒有那麼危險.
結果他發現其實沒那麼危險.
很多時候我們不敢遵守神的命令.
很多時候我們不敢去遵守神的命令.
是魔鬼在嚇你.
說如果你聽神的話去做.
你命都沒.
就說你聽神的命令.
你連命都沒.
但是神的命令是.
遵守祂的命令就有福.
但神的應許是.
遵行祂的命令你必定有福.
你看看約翰福音第十四章二十一節的話.
耶穌在約翰福音十四章二十一節說.
大家一起讀.
有了我的命令又遵守的.
這人就是愛我的.
愛我的必蒙我.
父愛他.

$^{1241}$我也要愛他.
並且要向他顯現.
所以你到底要聽誰的話.
所以到底你要聽哪一個人說呢.
我們從機電的經歷學到什麼呢.
我們從機電的經歷我們學到什麼呢.
我發現機電每一次禱告都有一件事是一樣的.
我發現機電每一次禱告都有同一件事.
他要發現到底是不是神的意思.
他要清楚明白是不是神的意思.
弄清楚以後他就遵行到底.
他明白清楚之後呢.
他就一定遵行到底.
他不怕失敗.
有一個登山專家.
他連續三次攀登世界最高的聖母峰.
他連續三次去爬世界最高的聖母峰.
都失敗了.
第三次失敗.
大多數隊員都死在雪崩裡面.
第三次失敗的時候.
很多隊員都死於雪崩.
他回到英國.
他們回到英國.
大家歡迎他.
然後那裡面掛了一幅.
這個聖母峰的照片.
掛了一幅很大的聖母峰的照片.
他就對著那個聖母峰的照片說.
他就對著聖母峰的照片說.
他說聖母峰我以.
這個.
我還活著的登山隊員的身份.
聖母峰.
我以還活著的登山隊員的身份.
還有沒有出生的登山隊員的身份.
加上還沒出世登山隊員的身份.
你可以打敗我們一次.
兩次.
三次.

$^{1281}$但是你不能永遠打敗我們.
但是你不會永遠打敗我們.
有一天我們會登到山頂.
有一天我們會登上去山頂.
因為你不會長大.
因為我們每一次都會比上一次預備得更好.
但我們每一次都會比上一次預備得更加好.
馬太福音第17章21節的話.
馬太福音17章20節.
一起讀.
你們若有信心像一粒芥菜種.
受試對這座山雪.
你從這邊羅到那邊也必羅去.
並且你們沒有一件不能作的事.
所以我們都是瓦契.
我們都是瓦契.
都是小信的人.
畢竟我們有很多的弱點.
我們其實都有很多的弱點.
但是你要知道你心裡有寶貝.
雖然你信心很小.
小到像芥菜種.
你不要小看它.
但你不要小看它.
要讓你的信心顯露出來.
一定會發出能力.
積澱求神成就別人都覺得不可能的事.
積澱求神成就在人眼中看不可能的事.
神就按著祂的信心成就了.
神就按著祂的信心去成就.
祂只是一個瓦契.
祂一步一步把心裡的寶貝顯露出來.
但祂願意將心裡的寶貝顯露出來.
神成就大事是一個開始.
神成就大事只是一個開始.
積澱有一個目的.
祂要召集以色列人去打敗欺壓他們的米田人.
祂要召集以色列人去打敗他們的米田人.
然後祂要成就心理期望的和平.
祂要成就心裡的期望和平安.

$^{1321}$我們的目的是什麼呢.
我們要成為大人的勇士.
同心合意一起為主征戰.
同心合意為主征戰.
把和平帶給全世界.
然後以色列人聚集在一起.
有多少人呢.
有三萬兩千人.
你看神怎麼宣告的.
我們一起讀.
現在你要向這些人宣告說.
凡懼怕膽怯的.
可以離開基列山返回去.
我們都知道這個故事.
然後神又說.
這個趴在地上喝水的人回家.
神說趴在地上喝水的人回去了.
這個手捧著水喝的人留下.
手捧著水喝的人可以留下.
最後只有三百人.
最後只剩三百人.
基地不怕嗎.
你看第七章七到九節.
七章九至十節.
當那爺爺和華吩咐基地說.
起來下到米殿營裡去.
因我已將他們交在你手中.
倘若你怕下去.
就帶你的僕人普拉下到那營裡去.
繼續.
於是基地帶著僕人普拉下到營旁.
那個經文告訴我們.
基地怕不怕.
這段經文說.
基地怕不怕.
一定是怕的.
當然怕了.
要不然他為什麼帶著僕人普拉.
下到營旁.
要不然他不會帶著僕人普拉下去.

$^{1361}$既然神說怕的都回家.
既然神說怕的人應該回去.
為什麼基地沒有走.
因為他是願意為神冒險犯難.
因為他願意為神去冒險.
從他第一次見到神的實質就是這樣的.
他第一次見到天使.
他就已經是這樣.
他願意貫徹到底.
他會貫徹到底.
雖然他沒有十足的把握.
但他有一點點信心.
雖然他不完全相信.
他願意憑著一點點信心去試試看.
他願意憑著這一點信心去試試看.
你看他帶著肉餅.
帶著湯去見天使.
你看他帶著肉餅和湯去見天使.
他不知道這個天使到底是神派的.
還是一個普通的先知.
他不知道這個天使是神派.
還是普通的先知.
但是他存著信心去試試.
他就願意存著信心去試試看.
然後他確定是神.
他確定是神.
神要他砍了偶像.
要拆掉巴力的壇.
他知道會冒犯眾怒.
他知道他這樣會激起公憤.
可能會來殺他.
可是他憑著小小的信心.
試試看.
但他憑著一點點信心他願意去試.
結果神就保護他.
後來他不肯定是不是神要他帶以色列人.
之後他不肯定是不是神要他帶以色列人.
猜險他去拯救以色列人.
他又小幸去求神.
試試看羊毛會變濕還是變乾.

$^{1401}$他就憑著這個小小的信心去問羊毛這個濕.
他的方法不對.
但是神都是對的.
但是神每一次都是對的.
然後下去神要他帶三百人.
每個人帶著瓶子.
每個人帶著玻璃瓶.
裡面裝火把.
裡面裝著火把.
然後到洗臉人中間.
去到米電人中間.
高舉火把.
高舉這個火把.
然後他喊什麼.
我們一起讀.
然後他喊什麼.
十八節.
喊叫說耶和華禍機電的刀.
這就是一個信心的冒險.
這又是一個信心的冒險.
神的方法常常是讓人不理解的.
神的方法其實常常讓我們不了解.
我們平常說人多好辦事.
神說只要三百人.
但神說你帶三百人就行了.
平常打仗帶刀.
平時人家打仗是帶刀.
你們帶瓶子.
你們帶玻璃瓶.
連刀都不要.
刀都不用帶.
可是喊的是什麼呢.
但他們喊的是什麼呢.
這個耶和華禍機電的刀.
耶和華禍機電的刀.
平常打仗是敵鄰我暗.
平時打仗是敵鄰我暗.
現在神要他們敵暗我明.
但神要他們敵暗我明.
以前以色列人見到米地人.

$^{1441}$像老鼠看到貓.
以前以色列人見到米地人.
好像老鼠見到貓.
逃走趕快逃.
快點走.
他們現在要高舉火把.
就把臉照亮給人家看.
要把自己的臉照清楚給人家看.
然後吹著大喊.
米地人就亂成一團.
以色列人連刀都沒有動.
以色列人連刀都不用用.
就勝了.
就得勝了.
米地人就逃走了.
米地人就逃走.
神的軍隊打大勝仗.
神的軍隊大勝.
你願不願意順服神的方法.
看神怎麼用祂的方法為你打勝仗.
曾經有那麼一個作家.
他和他的未婚妻訂婚四年不敢結婚.
因為他收入太少.
因為他收入太低.
他有很多夢想.
希望去全世界旅行.
他希望寫郵寄.
報導這些去旅行的經過.
他想去寫網誌.
去記錄去外面玩的旅行.
可是他沒有把握.
他不敢成行.
有一天他坐在公園裡.
遇到一個老先生.
兩個人一起看松鼠從一棵樹跳到另外一棵樹.
這個松鼠跳過去了.
他沒有抓住他要去的地方.
他抓住下面另外一根樹枝.
然後從這個樹枝爬到他想去的地方.
然後老先生就說.

$^{1481}$我看過很多松鼠在樹上跳來跳去.
大部分的松鼠不會第一次就跑到他們要去的地方.
很多松鼠是第一次不會跳到想去的地方.
可是他不會因為自己.
跳不到自己想去的地方.
就不跳.
但他不會因為自己跳不到想去的地方.
就停止不跳.
這老先生就在說了.
這個老先生在說.
我想如果他們不想老死在一棵樹上.
他們必須要跳.
我想如果他們不想老死在這棵樹上面.
就一定要冒險.
所以這個作家兩個星期以後就結婚了.
所以這個作家兩個星期之後就結婚了.
他去了歐洲.
又環遊世界.
專門寫旅行報導.
他的夢想完成了.
每一次遇到他要冒險的時候.
他就想到那個松鼠的故事.
他就想起這個松鼠的故事.
如果不想老死在一棵樹上.
如果不想老死在這棵樹上面.
一定要冒險.
敢不敢冒險.
只要你願意.
你就是大能的勇士.
只要你願意.
你就成為大能的勇士.
那應該怎麼做呢.
那應該怎麼做.
第一.
你要知道你是神大能的勇士.
雖然你是瓦契.
但你心裡有耶穌.
有耶穌在你心裡.
生命不一樣.
有耶穌在你心裡.

$^{1521}$生命就不一樣.
你這個瓦契不是普通的瓦契.
你心裡有寶貝.
你就讓你心裡的寶貝發出來.
讓你心裡的寶貝發出來.
你已經有能力.
你能夠靠著你能力成為神大能的勇士.
你可以靠著這個能力成為神大能的勇士.
第二.
不要老覺得自己軟弱沒有能力.
神說你靠著你已經有的能力.
就可以打敗彌殿人.
神說你靠著你已經有的能力去.
就可以打敗彌殿人.
第三.
你要用你微弱的信心.
你要使用你微弱的信心.
雖然現在屆賽總那麼小.
雖然屆賽總那麼小.
只要有信心.
神必定成就在人看來不可能的事.
第四.
你要願意為主獻上犧牲的奉獻.
你要願意為主獻上犧牲的禮物.
看看神怎麼樣在小事上.
顯出你奉獻的成果.
看看神怎麼樣在小事上.
可以顯出你奉獻的成果.
第五.
你要有信心為神冒險.
第五.
你要有信心為神去冒險.
願意奉耶穌的名去斥責魔鬼.
你要奉耶穌的名去斥責魔鬼.
你要相信神一定會保護你.
會賜你平安.
會賜給你平安.
你環境雖然沒有改變.
你環境可能不會有大改變.
可會你會感覺到神的同在.

$^{1561}$但你會感覺到神的同在.
這世界上的物質文明.
這個世界上的物質文明.
讓我們聽不見神的呼召.
我們聽不見神的呼召.
我們只聽見錢的聲音.
我們只聽見享樂的聲音.
今天神在你心裡呼召你.
你是大能的勇士.
你雖然是瓦契.
你心裡有寶貝.
你雖然是瓦契.
但心裡有寶貝.
寶貝有莫大的能力.
是出於神.
你願不願意成為這支基督精兵的中間的一個.
你願不願意成為基督精兵的其中一員.
你願不願意成為神國的大能勇士.
你想不想成為神國大能的勇士.
那我們一起低頭禱告.
天父我們謝謝你.
從祭奠的身上.
我們看見我們應該做什麼.
求你讓我們有勇氣.
去實踐你的話語.
禱告奉耶穌的名.
\newpage



\section{提摩太前書 3:1-7}
\label{sec:031669}
\textbf{長老的崇高呼召,資格和品格 (普/粵)}~~~~張景恒牧師
\newline
\newline
連結: \href{http://www.church.com.hk/acms/content.asp?site=cdc\&op=show\&type=product\&code=031669\&layout=sermon}{\texttt{031669}} ~~~~ 語音日期: 2026-02-15
\newline
\newline
\hyperref[sec:031668]{\small{< < < PREV SERMON < < <}}
~
\hyperref[sec:index_chronic]{\small{[返順時目]}}
~
\hyperref[sec:031670]{\small{> > > NEXT SERMON > > >}}
\newline
\newline
提摩太前書 3:1-7
\newline
\begin{longtable}{cl}
\hline
\hline
章節 & 經文 (和合本修訂版)\\
\hline
3:1 & \begin{tabularx}{0.7\textwidth}{X} 「若有人想望監督的職分，他是在羨慕一件好事」，這話是可信的。 \end{tabularx} \\ \\ \relax
3:2 & \begin{tabularx}{0.7\textwidth}{X} 監督必須無可指責，只作一個婦人的丈夫，有節制、克己、端正，樂意接待外人，善於教導， \end{tabularx} \\ \\ \relax
3:3 & \begin{tabularx}{0.7\textwidth}{X} 不酗酒，不打人；要溫和，不好鬥，不貪財。 \end{tabularx} \\ \\ \relax
3:4 & \begin{tabularx}{0.7\textwidth}{X} 要好好管理自己的家，使兒女順服，凡事莊重。 \end{tabularx} \\ \\ \relax
3:5 & \begin{tabularx}{0.7\textwidth}{X} 人若不知道管理自己的家，怎能照管神的教會呢？ \end{tabularx} \\ \\ \relax
3:6 & \begin{tabularx}{0.7\textwidth}{X} 剛信主的，不可作監督，恐怕他自高自大，落在魔鬼所受的懲罰裡。 \end{tabularx} \\ \\ \relax
3:7 & \begin{tabularx}{0.7\textwidth}{X} 監督也必須在教外有好名聲，免得被人毀謗，落在魔鬼的羅網裡。 \end{tabularx} \\ \\
[1ex]
\hline
\hline
\end{longtable}
$^{1}$弟兄姊妹 主內平安 農曆新年快樂.
今天是主鎖定的日子.
我們在其中要歡喜快樂 向神獻上感恩.
今天是主鎖定的日子.
我們在其中歡喜快樂 向神獻上感恩.
讓我們一起低頭禱告.
讓我們一起低頭祈禱.
我們在天上的父神 願你掌管 帶領今天的主日信息.
讓我們無論講的 聽的 都在基督裡同盟造就.
我們今天要一起思想兩段非常相似的經文.
提前三章一到七節 和提多一章五到九節.
今天我們要一起思想兩段非常相似的經文.
提前三章一到七節.
提多一章五到九節.
這兩段的經文是關於監督的呼召 資格和品格.
這兩段的經文是關於基督監督的呼召 資格和品格.
為什麼要放在一起看呢?.
因為它們就像一幅立體化的兩個視角.
因為它們就好像一幅立體的畫的兩個視角.
共同為我們構畫出神對教會領袖清晰完整的藍圖.
共同為我們構畫出神對教會的領袖清晰完整的藍圖.
弟兄姊妹 當我們談到教會的領袖時.
弟兄姊妹 當我們談到教會的領袖的時候.
根據新約聖經的教導 是指三個不同稱呼.
但同一個職份的牧師 長老 監督.
他們負責教導 牧養和治理教會.
提前三章一到七節 是保羅寫給當時留在以弗所教會的童工提摩泰.
一段有關呼召設立教會監督的重要教導.
提多一章五到九節 是保羅寫給另一位留在克里特教會的童工提多.
另一段有關設立教會長老的重要教導.
我們今天會以提前三章一到七節為主要的經文.
但是會參考引用提多一章五到九節.
我們會主要是以提前三章一到七節為主要的經文.
但是會參考提多一章五到九節.
這兩段經文不僅是對當時以弗所教會和克里特教會的提醒.
也是對今日所有教會和屬神兒女的寶貴教導.
也是對今日所有教會屬神兒女的寶貴教導.
我們的教會T4C在呼召和選立牧師長老時.
我們的教會T4C在呼召和選立牧師長老的時候.
是以這兩段經文所列出的資格和品格作為標準.

$^{41}$而不是單看恩賜,口才或社會的成就.
弟兄姊妹現在讓我們一起來看提前三章一到七節的經文.
提前三章一節是講到崇高的護召.
保羅在提前三章一節開宗明義的說.
人若想要得到監督的職份就是羨慕善功.
這話是可信的.
提前三章一節講到崇高的護召.
保羅在提前三章一節開宗明義的說.
人若想得監督的職份.
就是事務善功,這話是可信的.
這裡的監督指的是教會中的領袖職份.
包括牧師,長老.
這裡所講的監督所指的是教會中的領袖職份.
包括牧師,長老.
羨慕善功是什麼意思呢?.
事務善功是什麼意思呢?.
弟兄姊妹.
保羅首先肯定了渴望成為教會領袖監督的動機是正當的.
弟兄姊妹,保羅首先肯定了渴望成為教會領袖監督的動機是正當的.
這不是為了權力或地位.
而是羨慕善功.
一種為神國度事奉的渴望和負擔.
而不是為了權力或地位.
而是事務善功.
一種對神角度事奉的渴望和負擔.
神許願有人願於委身.
承擔牧養群養的責任.
神許願有人願於委身.
承擔牧養群養的責任.
這份動機應當是為了服侍神和祂的教會.
渴望參與神國度的事奉.
這份動機應該是為了服侍神和祂的教會.
渴望參與神角度的事奉.
事奉神不是追求個人的榮耀.
而是回應神的呼召.
以服侍人來服侍神.
事奉神不是追求個人的榮耀.
而是回應神的呼召.
以服侍人和服侍神.
例子.

$^{81}$一位臨命成熟的弟兄看到教會的需要.
心裡有負擔.
例子.
一位熟齡成熟的弟兄.
看到教會的需要.
心裡有負擔.
願意接受神的呼召.
並付出時間和精力去教導,牧養和管理.
願意接受神的呼召.
並付出時間和精力去教導,牧養和管理.
這是一種值得鼓勵的羨慕善功.
這是一個值得鼓勵的羨慕善功.
提前三章二節講到核心品格.
做監督的必須無可指責.
只做一個婦人的丈夫.
有節制自守端正.
樂意接待遠人,善於教導.
必須無可指責.
這是所有品格的核心.
必須無可指責.
只做一個婦人的丈夫.
有節制自守端正.
樂意接待遠人,善於教導.
必須無可指責.
這是所有品格的核心.
提多一章六節.
提到若有無可指責的人.
提多一章六節.
提到若有無可指責的人.
提多一章七節.
也提到監督必須無可指責.
提多一章七節.
也提到監督必須無可指責.
弟兄姊妹,無可指責不是指完美無罪.
乃是指無可挑剔沒有把柄.
一個領袖的生命.
經得起公眾的檢視.
沒有人能指控他.
一個領袖的生命.
經得起公眾的檢視.

$^{121}$沒有人能夠指控他.
只做一個人的丈夫.
提多一章六節也提到.
只做一個婦人的丈夫.
提多一章六節.
也提到只做一個婦人的丈夫.
這裡強調的是在婚姻上的絕對忠貞和信實.
他必須是一夫一妻.
婚姻盟約的忠誠守護者.
他必須是一夫一妻.
婚姻盟約的忠誠守護者.
例子.
這意味著一個領袖的婚姻是公開的.
是光明正大的.
例子.
這個意味著一個領袖的婚姻是公開的.
是光明正大的.
他和異性的關係有清晰的界線.
避免任何形式的曖昧.
不恰當的獨處.
他和異性的關係有清晰的界線.
避免任何形式的曖昧.
不恰當的獨處.
讓自己的婚姻成為堅固的見證.
而不是給仇敵流破口.
讓自己的婚姻成為堅固的見證.
而不是給仇敵流破口.
有節制,自守,端正.
這裡強調的是一個領袖內在生命的成熟度.
有節制,自守,端正.
這裡強調的是一個領袖內在生命的成熟度.
弟兄姊妹.
這些品格不是為了贏得人的誇獎.
弟兄姊妹.
這些品格不是為了贏得人的誇獎.
而是為了活出基督的生命.
成為神無瑕疵的見證.
而是為了活出基督的生命.
成為神無瑕疵的見證.
並能有效地服侍神的教會.

$^{161}$樂意接待遠人.
提多一章八節也提到樂意接待遠人.
樂意接待遠人.
提多一章八節也提到樂意接待遠人.
這裡是指開放家庭.
接待客旅.
表現出慷慨,開放,樂意分享的生命.
這裡所指的是開放家庭.
接待客旅.
表現出慷慨,開放,樂意分享的生命.
例子.
節日假日時.
像現在農曆新年.
邀請新來的家庭或單身的姐妹一起用餐.
讓她們感覺到家的溫暖.
例子.
節日假期的時候.
邀請新來的家庭.
或者是單身的弟兄姊妹.
到你的家裡一起吃飯.
讓他們感受到家的溫暖.
例子.
善於教導.
不僅是能講道.
更是能清晰,準確,有智慧的傳講.
和解釋聖經的真理.
善於教導.
不僅是能講道.
更是能夠清晰,準確,有智慧的傳講.
和解釋聖經的真理.
用教導來餵養群羊.
保護他們免受異端的勤勞.
用教導來餵養群羊.
保護他們免受異端的勤勞.
例子.
當弟兄姊妹對《譽頓論》或《末世論》感到困惑時.
例子.
當弟兄姊妹對《預定論》或者是《末世論》感到困惑的時候.
教會領袖能夠用聖經耐心的,有系統的講解.
教會的領袖應該是能夠用聖經內心地,有系統地去解釋.

$^{201}$而不是避而不談.
或者用自己的意見來回答.
而不是避而不談或者用自己的意見來回答.
提多書一章九節更強調.
堅守所教真實的道理.
就能將純正的教訓勸化人.
又能把爭辯的人駁倒了.
提多一章九節講到.
堅守所教真實的道理.
就能將純正的教訓勸化人.
又能把爭辯的人駁倒了.
弟兄姊妹.
長老除了善於教導.
弟兄姊妹.
長老除了善於教導.
他必須堅守真道.
用真道建造勸化人.
他必須堅守真道.
用真道建造勸化人.
用真道抵禦並反擊錯誤的教導.
例子.
當有異端譬如成功神學說.
信主必發財.
例子.
當有異端譬如成功神學.
說信主必發財.
或否認基督神性的教導.
進入教會時領袖不能保持沉默.
或否認基督神性的教導.
進入教會時領袖不能保持沉默.
他必須是一個真理的守護者.
站出來用聖經的真理.
清晰地向揚情解釋什麼是對的.
他必須是一個真理的守護者.
站出來用聖經的真理.
清晰地向揚群解釋什麼是對的.
並有力地指出錯誤何在.
保護揚群不受異端的侵害.
並有力地指出錯誤何在.
保護揚群不受異端的傷害.

$^{241}$提前三章三節講到消極的品格.
不飲酒滋事.
不打人.
只要溫和.
不增進.
不貪財.
提前三章三節是講到消極品格.
不因酒滋事.
不打人.
只要溫和.
不增進不貪財.
弟兄姊妹提前三章講的是積極的品格.
現在提前三章三節講的是消極的品格.
弟兄姊妹提前三章二節所講的是積極的品格.
現在在提前三章三節講的是消極的品格.
很明顯地從積極的要做什麼.
轉到消極的不可做什麼.
很明顯地從積極的要做什麼.
轉到消極的不可做什麼.
不因酒滋事.
提多一章七節也提到不因酒滋事.
不因酒滋事.
提多一章七節也提到不可因酒滋事.
這裡不僅指不酗酒.
更指一切因酒精影響而產生的不當行為.
這裡不僅是指不醉酒.
更指一切因酒精影響而產生的不當行為.
不打人.
提多一章七節也提到不打人.
不打人.
提多一章七節也提到不打人.
這裡不僅是對身體的暴力.
也延伸到語言上的攻擊.
好爭吵和易怒的性情.
這裡不僅是對身體的暴力.
也延伸到語言上的攻擊.
好爭吵和易怒的性情.
例子在會議中因為意見不合就提高音量.
用尖銳的語調打斷對方.
甚至進行人身攻擊.

$^{281}$用尖銳的語調打斷對方.
甚至是進行人身攻擊.
譬如你根本不懂你這個想法太愚蠢了.
溫和這裡是指對待軟弱和受傷的人.
有耐心和溫柔.
溫和這裡是指對待軟弱和受傷的人.
有耐心和溫柔.
例子當一個教會肢體.
反覆為同一件事情軟弱跌倒時.
例子當一個教會肢體.
反覆為同一件事情軟弱跌倒時.
不是嚴厲指責.
而是溫柔的鼓勵陪伴他成長.
不是嚴厲的指責.
而是溫柔的鼓勵他陪伴他成長.
不增進.
不為無關緊要的事爭吵.
不為自己的面子爭辯.
不增進不為無關重要的事情爭吵.
不為自己的面子爭辯.
例子在同工會議上.
當自己的提案被否定時.
能謙卑接納以團隊為重.
例子在同工會議上.
當自己的提議被否定時.
能謙卑接納以團隊為重.
而不是堅持己見.
最後導致會議不歡而散.
而不是堅持己見.
最後導致會議不歡而散.
不貪財.
在錢財上清清白白.
不利用職位謀取私利.
不貪財.
在錢財上清清白白.
不利用職位謀取私利.
提多書一章提到.
不貪無意之財.
比不貪財更強調錢財的來源.
提多書一章七節提到.

$^{321}$不貪無意之財.
比不貪財更強調錢財的來源.
例子如果有機會為教會採購.
在採購時須透明謹慎.
主動申報利益衝突.
例子如果有機會為教會採購.
在採購時須透明謹慎.
主動申報利益衝突.
絕不讓人在錢財上對自己有絲毫的懷疑.
絕對不讓人在錢財上對自己有絲毫的懷疑.
提前三章四到五節講到家庭建政.
好好管理自己的家.
使兒女凡事端莊順服.
提前三章四到五節講到的是家庭建政.
好好管理自己的家.
使兒女凡事端莊順服.
人若不知道管理自己的家.
焉能照管神的教會呢.
人若不知管理自己的家.
焉能照管神的教會呢.
好好管理自己的家.
這裡指出家庭是神檢驗一個人.
領袖能力的試驗場所.
好好管理自己的家.
這裡指出家庭是神檢驗一個人.
領袖能力的試驗場所.
管理家庭的能力.
被視為管理教會能力的預習和證明.
管理家庭的能力.
被視為管理教會能力的預習和證明.
弟兄姊妹.
神首先呼召我們在家庭中活出信仰.
弟兄姊妹 神首先呼召我們在家庭中活出信仰.
最真實的屬靈生命表現.
不是在教會大堂的椅子上.
而是在家庭的餐桌上.
最真實的屬靈生命表現.
不是在教會大堂的椅子上.
而是在家庭的飯檯上.
使兒女凡事端莊順服.

$^{361}$這裡表明兒女在父親的帶領和教養下.
端莊順服.
使兒女凡事端莊順服.
這裡表明兒女在父親的帶領和教養下.
端莊順服.
如果她連自己的家都管理不好.
如何管理神的教會呢.
如果她連自己的家都管理不好.
怎麼可以管理神的教會呢.
弟兄姊妹 管理不是專制的控制.
而是以愛,真理和榜樣.
使兒女端莊順服.
弟兄姊妹 管理不是專制的控制.
而是以愛,真理和榜樣.
使兒女端莊順服.
例子 一位父親若能以愛和原則.
而非獨裁和放任.
例子 一位父親如果能夠以愛和原則.
而不是以獨裁和放任.
帶領家庭靈修處理子女的衝突.
教導他們敬畏神.
帶領家庭靈修處理子女的衝突.
教導他們敬畏神.
這就證明她具備了牧羊神的教會.
所需的同理心,智慧和權柄.
這就證明她具備了牧羊神的教會.
所需的同理心,智慧和權柄.
除了使兒女端莊順服.
除了使兒女凡事端莊順服.
提多一章六節.
對於兒女的要求更具體.
兒女也是信主的.
沒有人告他們是放蕩不服約束的.
提多一章六節.
對於兒女的要求更加具體.
兒女是信主的.
沒有人告他們是放蕩不服約束的.
這裡強調信仰的傳承.
兒女不僅是行為端莊順服.
這裡是強調信仰的傳承.

$^{401}$兒女不僅是行為端莊順服.
更是有真實的信仰,是信主的.
例子,在當今社會.
這意味著長老的兒女.
雖然也會面臨世界的試探.
例子,當今社會.
這意味著長老的兒女.
雖然也會面臨世界的試探.
但總體上他們的人生方向是尋求神.
而不是沉溺於網路,物質主義或混亂的性關係中.
並願意順從父母的屬靈權別.
提前三章六節講到屬靈的成熟度.
出入教的不可作監督.
恐怕他自高自大.
就落在魔鬼手塑的刑罰裡.
提到三章六節所講到的是屬靈的成熟度.
初入教的不可作監督.
恐怕他自高自大.
就落在魔鬼所受的刑罰裡.
初入教的不可作監督.
這裡是要求屬靈上的成熟度和生命的歷練.
初入教的不可作監督.
這裡是要求屬靈上的成熟度和生命的歷練.
信仰需要經歷考驗和時間的試煉.
以避免因徒難的地位提升而驕傲.
從而陷入魔鬼因驕傲而墮入的圈套中.
以避免因徒難的地位提升而驕傲.
從而陷入魔鬼因驕傲而墮落的圈套裡.
弟兄姊妹.
這並不是說初信者不能參與事奉.
而是不能立即擔任監督的重任.
弟兄姊妹.
這裡不是說初信者不能參與事奉.
而是不能立即擔任監督的重任.
事奉應當循序漸進.
從較小的責任開始.
事奉應該是循序漸進.
從較小的責任開始.
例子.
一位信主不久.

$^{441}$很有恩賜的年輕弟兄.
雖然火熱.
但生命還未經歷足夠的試煉.
例子.
一位信主不久.
很有恩賜的年輕弟兄.
雖然火熱 但生命還未經歷足夠的試煉.
若太快被選立為長老.
他很可能容易看自己過於所當看.
這是極其危險的.
若太快被選立為長老.
他很可能看自己過於所當看.
這是非常危險的.
當他面對批評或事奉有果效時.
就極易變得氣餒或驕傲.
對他和教會都是傷害.
當他面對批評或事奉有果效時.
就會極易變得氣餒或驕傲.
對他和教會都是傷害.
提前三章七節講到對外見證.
監督也必須在教外有好名聲.
恐怕被人誹謗就落在魔鬼的網羅裡.
提前三章七節.
也講到對外見證.
監督也必須在教外有好名聲.
恐怕被人誹謗就落在魔鬼的網羅裡.
在教外有好名聲.
領袖的見證不僅在教會內也在社會上.
在教外有好名聲.
領袖的見證不僅在教會內.
也在社會上.
他的誠信,工作態度,為人處事.
應得到外邦人的尊重.
他的誠信,工作態度,為人處事.
應得到外邦人的尊重.
這並不是說我們要討好世界.
而是要在世界中活出光合元的作用.
這不是說我們要討好世界.
而是要在世界中活出炎和光的見證.
讓人透過我們的生命看見基督.

$^{481}$讓人透過我們的生命來看見基督.
落在教外沒有好名聲會成為福音的阻礙.
並給仇敵攻擊的破口.
如果在教會外沒有好的名聲.
就會成為福音的阻礙.
並且給仇敵攻擊的破口.
例子:一位被考慮的長老人選.
如果在公司裡以苛刻對待下屬聞名.
例子:一位被考慮的長老人選.
如果在公司裡以苛刻對待下屬聞名.
或是在社區中因財務糾紛被鄰居議論.
或是在社區中因財政糾紛被鄰居議論.
即使他在教會裡表現良好.
也不符合在教外有好名聲的標準.
即使他在教會內表現良好.
但也不符合在教外有好名聲的標準.
現在是結論和回應的時候了.
弟兄姊妹我們看完了提前三章1-7節.
和提多一章5-9節.
弟兄姊妹我們看完了提前三章1-7節.
和提多一章5-9節.
這兩段經文為我們提供了一幅完整的領袖畫像.
他是一個內在生命成熟.
家庭建政有力.
真理上剛強穩固的人.
他是一位內在生命成熟.
家庭建政有力.
真理上剛強穩固的人.
他既是一個溫柔的牧者.
也是一個勇敢的真理守護者.
他既是一位溫柔的牧者.
也是一個勇敢的真理守護者.
當我們看見這一幅清晰的畫像.
我們可能會鬆一口氣.
當我們看完這一幅清晰的畫像.
我們可能會鬆一口氣.
心裡想太好了.
這標準是給牧師長老的.
與我無關.
心裡會想太好了.

$^{521}$這標準是給牧師長老的.
與我無關.
當我們看到這些高標準時.
可能會感到氣餒.
當我們看到這個高標準的時候.
可能我們會感到氣餒.
完全合格的監督.
主耶穌基督.
完全合格的監督.
主耶穌基督.
祂是無可指責的.
完全絕智.
善於教導.
完美管理神的家.
但是祂卻從來不自高自大.
祂是無可指責.
完全絕智.
善於教導.
完美的管理神的家.
但是卻從不自高自大.
也許你會禱告說.
主啊,我達不到這些標準.
或者你會祈禱說.
主啊,我達不到這些標準.
但是我相信主也必安靜地回應你.
我知道,所以我為你釘死在十字架上.
主也必定安靜地回覆你說.
我知道,所以我為你釘死在十字架上.
弟兄姊妹,我們不是靠自己的能力.
達到這些標準.
而是靠基督的恩典.
弟兄姊妹,我們不是靠自己的能力.
達到這些標準.
而是靠著基督的恩典.
聖靈在我們裡面動功.
使我們能夠活出基督的品格.
聖靈在我們裡面動功.
使我們能夠活出基督的品格.
讓我們不灰心喪志.
而是憑信心領受恩典.

$^{561}$靠聖靈的大能.
一天一天更像我們的主.
讓我們不要灰心喪志.
而是憑著信心領受恩典.
靠聖靈的大能.
一天比一天更像我們的主.
所以弟兄姊妹,讓我們再仔細地想一想.
弟兄姊妹,讓我們再仔細地想一想.
我們是否聽懂保羅教導話語裡.
更深一層的含義.
我們是否能夠聽懂.
保羅教導話語裡更深一層的含義.
我們每一個人.
都是神國度的子民.
都是神的管教.
我們每一個人.
在基督裡都是神國度的子民.
都是神的管教.
我深信.
今天神能藉著保羅的話語.
向我們每一個人說話.
我深信.
今天神仍然藉著保羅的話.
向我們每一個人說話.
對教會牧長團而言.
對教會的牧長團而言.
當我們呼召和選立新領袖時.
當我們呼召和選立新領袖的時候.
我們必須用保羅教導的監督資格為標準.
全面考察一個人的生命,家庭和真理根基.
而不是只看恩賜口才或社會的成就.
對潛在的領袖而言.
神的呼召是崇高的.
標準是高超的.
神的呼召是崇高的.
標準是高超的.
請問你是否裝備好接受神的呼召.
這並不是一個冰冷的考核表.
而是一個在基督裡成長的藍圖.
這並不是一個冰冷的考核表.

$^{601}$而是一個在基督裡成長的藍圖.
神正在教會中.
在你的日常生活中,在你的辦公室,在你的客廳,在你的書桌前.
塑造你成為神所合用的一個器皿.
神正在教會中,在你日常生活中,在你的辦公室,在你的客廳,在你的書桌前.
塑造你成為神合用的器皿.
對每一個基督徒而言.
我們都可以用這份藍圖來檢視自己的生命.
我是否在追求無可指責的品格.
在我的職場和社區中活出見證.
我是否追求無可指責的品格.
在我的職場和社區中活出見證.
我是否在忠心管理神給我的家.
帶領家人一同走在信仰的道路上.
我是否在忠心管理神給我的家.
帶領家人一同走在信仰的道路中.
我是否在真理上不斷扎根.
以至於不僅能夠明白.
而且還能夠向人分享甚至為真道辯護.
我是否在真理上不斷扎根.
以至於不僅能夠明白.
還能夠向人分享甚至為真道辯護.
藉著使徒保羅的教導向我們說話.
願神祝福他自己寶貴的話語.
讓我們一起低頭做結束的禱告.
讓我們一起低頭做結束的祈禱.
親愛的天父.
為教會領袖,牧師,長老和監督.
設立了如此清晰又崇高的標準.
為教會領袖,牧師,長老和監督.
設立了如此清晰又崇高的標準.
我們懇求你在教會中興起這樣的領袖.
也幫助他們在生活的每一個層面.
追求聖潔無可指責成為你榮耀的見證.
賜給他們渴慕善功.
為你神的國度承擔責任的心.
賜給他們渴慕善功.
為你的國度承擔責任的心.
並賜下恩典使他們的行事為人.
能與蒙昭的恩相承.

$^{641}$並賜下恩典使他們的行事為人.
能與蒙昭的恩相承.
阿們.
謝謝.
多謝.
\newpage



\section{詩篇 62:1-12}
\label{sec:031670}
\textbf{實在靠得住}~~~~陳永就牧師
\newline
\newline
連結: \href{http://www.church.com.hk/acms/content.asp?site=cdc\&op=show\&type=product\&code=031670\&layout=sermon}{\texttt{031670}} ~~~~ 語音日期: 2026-02-15
\newline
\newline
\hyperref[sec:031669]{\small{< < < PREV SERMON < < <}}
~
\hyperref[sec:index_chronic]{\small{[返順時目]}}
~
\hyperref[sec:031671]{\small{> > > NEXT SERMON > > >}}
\newline
\newline
詩篇 62:1-12
\newline
\begin{longtable}{cl}
\hline
\hline
章節 & 經文 (和合本修訂版)\\
\hline
62:1 & \begin{tabularx}{0.7\textwidth}{X} 我的心默默無聲，專等候神，我的救恩從他而來。 \end{tabularx} \\ \\ \relax
62:2 & \begin{tabularx}{0.7\textwidth}{X} 惟獨他是我的磐石，我的拯救；他是我的庇護所，我必不大大動搖。 \end{tabularx} \\ \\ \relax
62:3 & \begin{tabularx}{0.7\textwidth}{X} 你們大家攻擊一人，使他被殺，如歪斜的牆、將倒的壁，要到幾時呢？ \end{tabularx} \\ \\ \relax
62:4 & \begin{tabularx}{0.7\textwidth}{X} 他們彼此商議，要把他從高位上拉下來；他們喜愛謊話，口雖祝福，心卻詛咒。（細拉） \end{tabularx} \\ \\ \relax
62:5 & \begin{tabularx}{0.7\textwidth}{X} 我的心哪，你當默默無聲，專等候神，因為我的盼望是從他而來。 \end{tabularx} \\ \\ \relax
62:6 & \begin{tabularx}{0.7\textwidth}{X} 惟獨他是我的磐石，我的拯救；他是我的庇護所，我必不動搖。 \end{tabularx} \\ \\ \relax
62:7 & \begin{tabularx}{0.7\textwidth}{X} 我的拯救、我的榮耀都在於神；我力量的磐石、我的避難所都在於神。 \end{tabularx} \\ \\ \relax
62:8 & \begin{tabularx}{0.7\textwidth}{X} 百姓啊，要時時倚靠他，在他面前傾心吐意；神是我們的避難所。（細拉） \end{tabularx} \\ \\ \relax
62:9 & \begin{tabularx}{0.7\textwidth}{X} 人真是虛空，人真是虛假；放在天平裡就必浮起，他們一共比空氣還輕。 \end{tabularx} \\ \\ \relax
62:10 & \begin{tabularx}{0.7\textwidth}{X} 不要仗勢欺人，也不要因搶奪而驕傲；若財寶加增，不要放在心上。 \end{tabularx} \\ \\ \relax
62:11 & \begin{tabularx}{0.7\textwidth}{X} 神說了一次、兩次，我都聽見了，就是能力屬乎神。 \end{tabularx} \\ \\ \relax
62:12 & \begin{tabularx}{0.7\textwidth}{X} 主啊，慈愛也是屬乎你，因為你照著各人所做的報應他。 \end{tabularx} \\ \\
[1ex]
\hline
\hline
\end{longtable}
$^{1}$主持人:早晨!弟兄姊妹,首先代表總會祝大家新的一年,我們能夠身體健康,零命來到長進,很開心能夠在大家當中宣講上帝的聖言..
今天跟大家講這個詩篇,六十二篇.我想到從前工作的時候,我不是一出生就做牧師,我也有我的學習,工作的經驗等等,後來神就呼召我,在教會裡面來服侍..
我記得以前工作的時候,有一些同事很卸膊的,斤斤計較,不負責任的.那些同事其實背後都很不滿他,有時甚至浮面..
如果找他工作,譬如我提議找誰誰做事,他就會加一句,找誰誰做事,誰靠得住,你知道下一句是什麼嗎?.
原來你們身邊這麼多豬奶,我不知道,誰靠得住豬奶都會上樹,我不知道是豬公還是豬奶,我沒有去考究..
真的,我又沒有考究這句,為什麼不是牛公會上樹,是豬奶會上樹呢?.
不過肯定一件事,誰靠得住豬奶都會上樹,就講到可靠和責任,我們待人接物其實是一個很重要的生命質素..
今天我們宣講詩篇62篇,就是提醒我們,神是可靠的那一位,是依靠的那一位,是靠得住的那一位..
我們再有一個祈禱,我們求主來帶領,親愛的聖靈願你來幫助,成為我們的力量,成為我們的依靠,成為我們隨時的幫助..
讓我們藉著你的聖言,生命能夠得到堅固.我們這樣同心的祈禱,奉耶穌基督的名求,阿們..
這篇詩篇告訴我們,是大衛的詩交給靈長趙耶杜頓,他有幾個細節,然後就說,.
唯獨等候神,我的心才有安息,我的拯救從他而來.唯獨他是我的盤石,我的拯救,他是我的保障,我必不自動搖..
你們大家攻擊一個人,要把他殺害,如同拆毀傾斜的牆,張倒的壁,要到幾時呢?.
他們彼此相依,無要把他從磚位上拉下來,他們喜愛放話,他們口須祝福,心卻就坐..
我們知道詩篇其實是聖經裡面最長的經卷,150篇.我們透過篇收,有卷一,卷二,卷三,卷四,卷五,對照摩西的五經..
而62篇就放在卷二,那裡絕大部份都是大衛的詩,或者和大衛有關的詩..
你會看到如果我們這樣看這篇詩篇的時候,會代入這個詩人在面對的一些艱難.我們讀詩篇,當然可以說是博大精深的八幾篇..
有些研究的學者就將它分為十五類,有登基詩,有上行詩,有讚美詩等等,有些就將它分為六七類,不同的做法..
不過我們可以大致都分為兩類,一個就是哀歌,就是我知道上帝好,我現在面對艱難,上帝你拯救我呀,等著我,那就是哀歌..
我知道上帝好,上帝你幫助我,神真的救了我,我讚美主,那就是讚美詩,大概是這樣..
不過當然也有例外的,譬如哈利路亞詩就是另外一個例外,不過我們有時特別做牧師,教聖經要弟兄姊妹明白才行,對不對?.
他明白了,他大概看詩篇,他會有這樣的印象,他就投入到詩篇裡面,私人那種心靈的共鳴..
如果我們看這篇詩篇,很明顯大衛當時遇到什麼樣的境況?.
如果你代入去看回這個聖經的時候,他的兒子都叛變他,阿沙隆,甚至追殺他,他逃命的時候,為什麼?.
那些人想把我從高位裡面拉下來,是在面對一個敵人攻擊的時候,但是這個時候卻是告訴我們更加需要倚靠神的時候..
那個調是用耶道頓,這個人如果你看回歷代詞,他出現過的,是一個師班的領班,領唱,或者今天我們這樣大家明白一點,是師班長來的,領唱..
當然這個字都可以解成認的意思,可以代表了承認耶和華他是值得我們等候的神,承認他是唯一的神..
事實上我們改名字都有個期盼的,等於我的名字陳永就,我爺爺幫我改的,他可能都希望我所有的成就..
無獨有偶,我爸爸叫阿成,不知道是不是真的有關係,真的有些關係也說不定,所以從名字都會反映了一些事情在裡面..
無論如何我們都知道這篇詩寫的時候,詩人很想我們讀的時候能夠體會得到,承認得到上帝是可靠的那一位..
如果你讀這個和合本就很熟的了,我的心脈脈無聲,專等候神,我講到喜歡用新譯本的..
他翻譯成唯獨等候神,我的心先有安息,唯獨等候他,他是我的盤石,他是我的拯救..
這個唯獨,詩人一開始就這樣宣告,其實是貫出成首詩的,第一二節之外,我們再看下去一會兒,第四五節都是這樣講,第六節第九節都是這樣講..
有一種肯定,有一種絕對的意味,縱然現在我面對著敵人的攻擊,但是我內心裡面有一份安穩,外面是波濤洶湧的,.
但是我仍然宣告上帝是我唯獨可以依靠的那一位..
這個更加凸顯了他心裡面那份安穩和平靜,唯獨安息在這篇詩篇裡面重複地出現,.
是表達對上主那份信心,是默念他那份恩慈,是降服他的主權..
詩篇裡面講到確實,唯獨他有這樣的意味,確實,盛言,或者剛才講和合本益專心,專那個意思,真這樣的意思..
是一份在外面就算怎樣的環境也好,但是我內心有一份寧靜,是一個他依靠神的秘訣,就是他將這些事情交託給神..
唯獨安息貫穿了這篇詩篇,你再看下去很多的詞彙都是重複地出現的,一會兒我們讀下去..

$^{41}$盤石,拯救,保障,不致動搖.盤石就比喻為那份恆久和堅固,.
你對照下文就是那些傾斜的牆,張倒的壁,長倒眾人推,很脆弱..
但不是,這邊是很堅固,形成了一個對比.至於拯救,是詩人他親身這樣經歷,所以他才有這份信心,他有一份這樣的安靜..
不單止這樣,他也呼籲其他人好像他這樣,相信上帝的拯救,保障就有一種避難所的意味..
你看看面對敵人的那種攻擊,殺害,甚至推倒,詩人就不禁要發出一個疑問,他說要到幾時呢?.
我們讀詩篇很多時候這句騙語都會出現的,要到幾時呢?可能是我們艱難裡面去問上帝要到幾時呢?.
但如果你看這裡上下文,又好像不是,面對那些充滿了野心要加害他的人,.
詩人說你們不怕上帝嗎?你們還不收手,要到幾時呢?.
反而好像是勸了他們不要,你們不要做這樣的事情來得罪神..
因為真是令人心寒,你看看這班人,口就祝福,心就坐,口裡不一,口蜜福劍..
我想起文革的時候有幾句說話不知道大家有沒有聽過,語錄不離手,萬歲不離口,天天說好話,背後下毒手..
這個不是我講的,是周恩來講的,你上網找得到的,不是我講的..
語錄不離手,個個拿著本無語錄,萬歲不離口,天天說好話,背後下毒手..
你想想你身邊是這樣的人,這麼陰險的人,詩人怎樣說呢?要到幾時呢?你們不要再作惡了..
他不是那種粉粉不平,而是他反而勸人們不要再這樣,再這樣行惡..
面對逆境艱難,我們內心難免躁動的,我們往往選擇自己的方法,.
為什麼別人這樣攻擊我,我就打他,很正常,或者我走,也很正常..
但你看到詩人他選擇什麼方法呢?他是等候神,他是等候神..
這個是他和上帝有一份很親密的關係,他意識到上帝在他的生活當中..
我想意識神同在是非常重要的一件事情,我們作為牧者傳道,常常都勸勉弟兄姊妹要靈修,要親近主..
其實我們說屬靈操練只不過是一個習慣而已,我們有沒有和神有這樣的關係,而敏感祂在我們生活中的作為呢?.
有一個人他很肚餓,一個大男人他很肚餓,家裡沒有東西吃,他就走出去找東西吃,一出門就被他見到一個斬柴佬,一個巢夫..
有沒有東西吃?有沒有東西吃?沒有啊,我的柴,送給你吧.不要啊,不要啊,退卻了人家一番的好意..
接著再行前,又見到一個小販,賣那些煲煲,餐餐,碗碗碟碟,有東西吃我就沒有啊,這樣吧,送個煲,送個餐給你吧..
我不要啊,我要食物啊.繼續這樣來,很怒氣的來行前,終於見到一個小販有東西賣了,有食物了,有些生肉在那裡,無所謂啊,我送給你..
不過我沒有煲,沒有餐,沒有生火器具,生肉吃不到啊.他才發現他原來錯過了沿途很多的機會..
抄下頂尖,我們有沒有每日感到他敏感上帝他自己的臨在,意識神的同在呢?.
我們說靈修,操練其實不是一個很神秘的事情,如果你當為一個習慣,你就覺得很簡單,你早上起床要不要刷牙?.
我現在也是的,起床我刷完牙之後,梳洗個人比較清醒之後,我第一件事,我一定不看手機的,為什麼?.
因為一看已經有八,九十個WhatsApp了,很多的消息會進入你的腦裡,但是清晨的時候我要歌頌我的主,我要親近我的神,.
我就會默想一節聖經,或者一兩節有時候,看完之後反覆地來祈禱,然後坐在那裡恭恭敬敬,請說,獨人敬聽..
有時候內心裡面就會浮現一些我沒想過的事情出來,你不要小看那十五分鐘,其實不用多,你做得密,每天這樣做,持之以恆..
我們現在經常說養生,有沒有養一下心呢?安安靜靜地在神的面前,其實這個是不是也是耶穌葡萄樹和之子的教導呢?.
「尼了我不能作甚麼」,但是我們說葡萄樹,我們很困難的,為什麼?因為我們沒有誰種過葡萄,.
但是如果我轉一轉個比喻,這個「尼了我不能作甚麼」就是你手機和充電的關係,.
你晚上做什麼,每一晚,當然是充電才睡覺,對不對?你不充了電,你也會亂七八糟,你會下床走回去充電,.
我現在也是,我如果出街,我忘記帶錢包,我都不拿,但如果忘記帶手機,我去到地鐵站都回頭,.
為什麼?因為放了東西在裡面,你手機要充電,你有沒有和神有這樣的聯繫?.
「尼了我不能作甚麼」,正正私人和神有這一份親密的關係,他能夠辨識到,他聆聽到,跟隨他的神,.
他也告訴我們知道,唯獨神才是我們的盼望,第五到八節,.

$^{81}$唯獨等候神,我的心才有安息,因為我的盼望是從祂而來,唯獨祂是我的盤石,我的拯救,我的保障,我必不自動搖,.
我的拯救,我的榮耀都在於神,我堅固的盤石,我的避難所都在於神,.
眾民啊!你們要常常倚靠祂,在祂面前傾心吐意,因為神是我們的避難所..
第二個段落,你會看到第五和第六節,好像剛才第一二節重複的,.
不過如果你細查,你會發現第一節是講拯救,這裡轉了,唯獨是我的盼望,.
為什麼?因為第一節他經歷過上帝的拯救,現在還不行,我還在等候,但我有那份盼望,.
這個盼望其實是從神過去對他的拯救而來的,是一個有基礎的信心,.
他相信上帝會保護他,拯救他,事實上信心就是沮喪最有效的解毒劑,.
在神面前默默無聲,這個不是懦弱的表現,而是你對上帝那份信靠和把握,.
我必不自動搖,如果讀和合本,第四節很有趣的,我必不肯動搖,.
我第一次讀我覺得很怪,如果用中文翻譯是不對的,但新日本他沒有了,他照改了,免得在那裡太多事情,.
其實是一個曾帝來的,他第一段還有一點猶豫,去到第二段五到八節很清楚的,.
我必不自動搖,因為他經驗過了,我不會再動搖的了,由始至終我不會再被打倒的了,.
不是因為我,而是凸顯了上帝他那份可靠和對他那份信心,第七節更加講到榮耀,.
這個當然可以是上帝的榮耀,但用在人身上都可以,猶太人裡面這個都是代表名譽地位那些,.
詩人在這裡做了個對比,那些自以為社會上有名望的人做了個對比,.
但他說不要這樣,不要依靠這些,反而你要依靠什麼,又再出現一些重複的詞彙了,.
堅固的盤石代表那份穩妥,避難所可以遮風擋雨,詩人面對敵人的攻擊,身邊有人密謀加害,.
將他從專位拉下來,拉下來會怎樣,會失去現有的一切,但詩人又明白,.
其實現有的一切根本不是你們給我的,是神給我的,所以神才是我的保障,.
有這一份這樣的洞察力,從而他才對上主有這樣的一份信心,.
接著他說中文,如果你看和學本就是百姓,可以想像大衛和跟隨他的百姓說,.
你們都要依靠神,是肺腑之言,他呼籲他們要傾心吐意,.
依靠神不要將你的信心放在一些金錢物質甚至人的權位身上,.
傾心吐意就是你和神可以有一份很親密的關係,你不需要隱藏,.
你可以將你的軟弱或將你的憧憬跟上帝所說,基於這樣的理由,.
所以詩人才勸我們不要依靠這些東西,人的力量,你單單依靠神,.
依靠神就是不看任何時勢,任何環境,而為什麼他能夠有這份秘訣,.
這份秘訣就在於他和神之間那份親密的關係,.
詩人面對艱難,頹垣敗瓦,甚至身邊有人惡意加害,.
但他仍然選擇盼望,專心等候神,這個不是一個無奈或逃避,.
而是因為有一個正確的方向,他內心有一份踏實,而不急於去回應這個環境,.
這份安息和把握是因為他對神有這份盼望,我不知道你是否正在面對困難,.
你對神有沒有這份盼望呢?.
上兩個月我們去了日本開會,.
每四年一次的,整個亞洲區的播道會,有菲律賓,新加坡,香港等等,.
會聚首一堂,很感恩,今年林牧師領軍我們去了,.
我們香港都有近四十個參加者,我特意在富士山打卡,.
這個不是AI,真是我,到此一遊,除了宣講和工作坊外,.
我很享受的是他會將不同地區和國家的弟兄姊妹放在一起,.

$^{121}$弟兄歸弟兄,姊妹歸姊妹,.
我那組有新加坡,馬來西亞,菲律賓一起,.
有一個弟兄的見證令我深受感動,.
這位弟兄不是傳道人,他是一個社工,.
他告訴我他信主前的一些見證,.
他在日本從事社會工作,幫一些菲律賓來的比較低下階層的外勞,.
他很有心,大家都知人離鄉自然,去到那個地方當然很多需要,.
他想如何能幫到這些菲律賓的外勞,.
於是直頭去了菲律賓的一間大學進修,.
大家知菲律賓馬來西亞出名的其中一個地點叫垃圾山,.
那個地方那些人每日就去執垃圾,.
透過執垃圾而生存,.
你可想而知那個環境多惡劣,.
他就去這些地方做一些田野考察,.
做一些問卷調查,.
其實大家都明,這些問卷大概都是這樣,.
你生活得開心不開心,你生活得辛苦買辛苦多少度,.
跟住想要什麼福利,大概是這樣,.
他計劃這樣的研究希望幫,.
預備這些就去到垃圾山,誰知他迎面而來就有個小女孩,.
他就想和這個小女孩做個問卷調查,.
他估到的,問她開心不開心,當然不開心,.
你有什麼需要,當然多些福利,.
於是他就問,.
小朋友你生活得開心不開心,.
開心,我生活得開心,.
他呆了,不知如何問下去,.
跟住他說你為什麼開心,因為我有主耶穌,.
他呆了,.
這個小朋友沒理由騙他,.
他的影響大到,.
為什麼信仰的能力可以給一個在這麼艱難的小朋友,.
有這樣的希望,.
於是他改變了研究的題目,.
他後來去研究宗教,.
跟住研究基督教,跟住研究主耶穌,研究福音,.
最後信了耶穌,.
娶了菲律賓老婆,.
回到日本來服侍,.
她說回去見證我自己心裡很大的感動,.

$^{161}$親愛的弟兄姊妹,你遇到什麼艱難,.
神可不可以成為你的倚靠,.
神可不可以成為你盼望的力量,.
求主成為我們的幫助,.
先人告訴我們,.
唯獨神才能成為我們的倚靠,.
卑賤的人都是虛空,.
高貴的人都是虛假,.
放在天平上就必升起,.
加在一起,比空氣還輕,.
不要倚靠那些強暴得利的財物,.
不要妄想靠那些搶奪回來的錢財,.
如果財寶加增,.
你也不要倚靠它,.
神說了一次,.
我兩次聽過,.
能力是在乎神,.
主啊,.
此愛是屬於你,.
你必照各人所作的,.
報應各人..
當時人呼籲大家倚靠神之後,.
他就從另一個角度,.
突顯什麼不可靠,.
包括人和錢,.
第九節,.
卑賤的人和高貴的人,.
做了一個對比,.
這裡他不是貶低那些人,.
他的意思等於我們的成語,.
恭候將相,.
凡夫走卒,.
他是說什麼階級的人,.
在神面前,.
或者生命那種裡面,.
其實都好像一口氣,.
甚至比空氣還輕毫無重量,.
而世人以為金錢就可靠了,.
甚至不擇手段,.
貪戀那些不義之財,.

$^{201}$斯人指出,.
你不要這樣來依靠,.
那不是,.
我正途穩的呢,.
其實正途穩都不是你依靠的,.
你依靠的是上主,.
一個人如果有財有勢,.
就容易此世寧人,.
甚至更加貪婪,.
以欺控的方法,.
獲取那些不義之財,.
斯人就提醒,.
不要這樣,.
人的財富加增,.
你的勢力都加增,.
但你的盼望不要改變了,.
你的盼望仍然在於神,.
要常常依靠他,.
而不是依靠此世寧人的人,.
神說了一次,.
我聽了兩次,.
這個當然可以是數量上,.
但我們讀詩篇很多時都說,.
好像走樓梯那樣,.
有個層底的,.
好像走高了一級,.
一次,.
兩次,.
上去加強了,.
是什麼?.
是神的慈愛,.
有些日本現在就將它翻譯成,.
神堅定的愛,.
重點在於那份堅定,.
神會報應各人,.
這個報應,.
當然我們知道,.
在基督的台前,.
當耶穌再回來的時候,.
每一個都要交帳,.

$^{241}$作為信徒我們都要交帳,.
神會使人審判,.
報應各人,.
當然可以是那時候,.
但神自己行其事,.
可以很快也行,.
神是公義的,.
不是像人那樣,.
心口不一,.
面對這樣的神,.
先人就呼籲我們,.
可以耐心,.
專心,.
放心地來投靠,.
因為能力在於無限的神,.
當人敬畏神的時候,.
你就沒有其他的恐懼,.
神的力量足以幫助你,.
綽綽有餘,.
人心放膽來投靠,.
事實上我們每天面對職場的生活,.
我們有很多的抉擇,.
有很多的辨別在裡面,.
我們需要有一個這樣來依靠神的心,.
我曾經和一位弟兄聊天,.
很好笑,.
他說,.
職場那種的凶險,.
但以你就是一輩子被人扔到獅子坑一次,.
我每天上班都幾乎被人扔到獅子坑,.
真的,.
有時真的很凶險,.
我試過有一個肢體,.
他做的那些私募基金,.
在內地,.
叫他做某件事情,.
給他一張支票,.
寫了他的名字,.
簽了他的名字,.
沒有寫銀碼,.

$^{281}$但你要做一件事情,.
我都好奇,.
我最後問他,.
你有沒有收到?.
他說,.
當然不能收了,.
收了還可以回來,.
我們有很多這些抉擇,.
是需要我們和神有那個關係敏銳來做決定,.
是不是真的依靠他呢?.
提到依靠,.
我想起我舊的一位弟兄,.
他自己本身做一個部門的主管,.
做得不錯的,.
有業績的,.
誰知換了一個頂頭老闆之後,.
就將他調去一個毫無發展的部門,.
其實大家心中有數的部門等執的,.
他覺得很不公道,.
為什麼這樣對我,.
接著就明白了,.
原來他調了他那些人坐他的位置,.
他更加不憤氣,.
於是電郵總公司,.
但不得要領,.
這些內部當然信那個最大的安排,.
給你會怎樣?.
要不就打過,.
剛才這樣說,.
要不就走,.
不是,他領收的時候,.
甚至叫他等候,.
他真的很奇怪,.
好,他就馴服,.
他也是警衛神的,.
等候,.
他繼續稱職在他的崗位裡面,.
做回應該做的事情,.
沒有因為這樣而做其他的瞞騙的事情,.
為什麼兩個多月,.

$^{321}$三個月不夠呢?.
這個頂頭老闆突然被人解僱了,.
他都呆了,.
接著又接到電郵,.
又調回他原來的崗位,.
後來才知道辦公室消息,.
原來這個頂頭老闆,.
在公司搞婚外情,.
不行的,.
紀律不行,.
總公司說不行,.
你操手不行,.
結果馬上解僱他,.
他以為在那裡很委屈,.
在那裡等候的時候,.
原來上帝已經在做事,.
親愛的兄弟姐妹,.
我想讀這篇詩篇,.
我們共鳴的地方就是,.
面對這些事情,.
我們可以同流合污的,.
是不是?.
我們也可以粉飾疾俗,.
但詩人不是,.
詩人給了我們第三個選擇,.
他是專心倚靠上帝來做回應,.
事實上人真的很軟弱,.
你看看這段詩篇說,.
人放在天平上,.
家在一起,.
比空氣還輕,.
第九節那裡說到,.
其實那個虛空,.
和傳道書的虛空是一樣的,.
好像一陣風一樣,.
他提醒我們不要奢望,.
不要奢望在人那裡得到什麼,.
但是縱然人好像空氣,.
雄無那麼輕,.
但上帝仍然看重你,.

$^{361}$看重我,.
看重我們每一個,.
提到人那種脆弱,.
生命的那種短促,.
我就想起早兩個月,.
大家都知道,.
宏福縣大火,.
我們看到這件事的時候,.
心裡面真的流淚,.
不是說笑,.
整個城市都受到創傷,.
我們執委會很快做決定,.
我們說要成立一個鎮災委員會,.
於是立刻很快,.
幾天我們就給了一個電郵,.
中堂會,.
很快,.
兩個星期都不夠籌,.
二百多萬已經,.
我們都立刻舉辦祈禱會,.
叫過千人參加,.
我們本來開頭是想著,.
透過我們播道會的堂會,.
認識到哪些災民就幫他,.
是用這樣的方法想著,.
所以那個規模,.
不是我們想像中那麼大,.
誰知道當我們這些申請表,.
給一些災民拿到,.
而放在他們那個WhatsApp的群組,.
和那個一戶一社工的群組,.
結果一傳十,.
十傳八,.
八傳千,.
我們短時間幾天內,.
收到過千的申請表,.
我們真是連影印機都拍壞,.
有人直接上來我們辦公室問,.
是不是有錢派,.
我們不是想著這樣的,.

$^{401}$是上帝逼我們,.
逼我們要去回應那個,.
那麼大的需要,.
當時我們立刻想,.
我立刻上帝給我的領袖,.
不是只是派錢,.
我們真的要幫到這些有需要的人,.
我們要告訴他,.
是在艱難裡面,.
在眼淚當中,.
神愛你,.
神沒有忘記你,.
於是我們立刻改變了做法,.
不是只是派錢,.
而是有一張支票,.
給他上來我們教會當中,.
簽收之餘,.
我們為他祈禱,.
而事實上已經四散了,.
不是在大埔,.
有些去了香港濟,.
所以我們找廣福,.
有些去了洪水橋,.
所以我們找元朗恩福,.
有些在將軍澳的中轉屋,.
於是我們找書院道真,.
我盡量找了所有教會,.
打電話聯絡,.
怎樣怎樣安排,.
很多很多的工作,.
日以繼夜,.
但我們的同事都沒有怨言,.
因為我們很想告訴他們知道,.
在艱難裡面,.
教會關心你,.
很多的故事聽,.
有人性的醜惡,.
我曾經聽過一個申請者,.
因為我都派票,.
我下去做,.

$^{441}$他說那些親戚好日都不找,.
知道我們屋燒了,.
政府有錢給,.
打電話來問有沒有得分,.
多醜惡,.
他家裡竟然有人過身,.
但都有人性的光輝,.
有些淪社,.
你都聽到,.
叫大家走,.
然後失散了,.
竟然在教會重遇,.
真是相擁而哭,.
恍如隔世,.
結果牧師為你祈禱,.
我自己接觸的其中一個都是這樣,.
我都表明身份,.
我是牧師,.
簽名支票,.
我為你祈禱好不好,.
多數都樂意,.
我試過有一個祈完禱之後,.
他說我十幾日未試過這麼平靜,.
牧師你可否再為我祈禱,.
我當然樂意,.
樂意為他祈禱,.
然後祈完禱之後,.
他小時候上過一些基督教小學,.
我說你願不願意信耶穌,.
我和你決志,.
他說好,.
我想信耶穌,.
即場和他決志,.
你可以想像,.
如果你是那裡的災民,.
你看著那間屋,.
什麼都在燒,.
很想回去拿東西走,.
但結果什麼都拿不到走,.
這豈不是都是我和你的寫照嗎?.

$^{481}$你可以拿到什麼走?.
當你有一天離開這個世界的時候,.
你可以拿到什麼走?.
生命是這麼脆弱,.
但儘管是這樣,.
神祂是愛我們愛到底,.
將祂的獨生愛子耶穌基督,.
給過我們成就救恩,.
只要每一個人口裡承認,.
心裡相信的時候,.
就必得救..
\newpage



\section{創世記 14:20}
\label{sec:031671}
\textbf{聖經人聖經事:麥基洗德}~~~~黃克勤牧師博士
\newline
\newline
連結: \href{http://www.church.com.hk/acms/content.asp?site=cdc\&op=show\&type=product\&code=031671\&layout=sermon}{\texttt{031671}} ~~~~ 語音日期: 2026-02-08
\newline
\newline
\hyperref[sec:031670]{\small{< < < PREV SERMON < < <}}
~
\hyperref[sec:index_chronic]{\small{[返順時目]}}
~
\hyperref[sec:031672]{\small{> > > NEXT SERMON > > >}}
\newline
\newline
創世記 14:20
\newline
\begin{longtable}{cl}
\hline
\hline
章節 & 經文 (和合本修訂版)\\
\hline
14:20 & \begin{tabularx}{0.7\textwidth}{X} 至高的神把敵人交在你手裡，他是應當稱頌的！」亞伯蘭就把所有的拿出十分之一給他。 \end{tabularx} \\ \\
[1ex]
\hline
\hline
\end{longtable}
$^{1}$各位網友大家好.
又一齊來讀聖經.
聖經人聖經是生命宣教網絡提供.
今天去看一個聖經人物.
這個聖經人物厲害了.
就是無始無終.
無俗譜.
就是神的至高神的祭司.
叫墨基西德.
我們看看墨基西德.
仍然是亞伯拉罕創世記的故事.
創世記十四章.
第十七節.
亞伯蘭撒拜基大老馬和他同盟的王回來的時候.
所多馬王出來在沙微谷迎接他.
沙微谷就是王谷的意思.
又有撒冷王.
他的名字叫墨基西德.
拿著餅和酒出來迎接他是至高神的祭司.
這是什麼事呢?.
他為亞伯蘭祝福.
願天地的主至高的神賜福於亞伯蘭.
至高的神把敵人交在你手裡.
是應當稱頌的.
亞伯蘭把所得的拿出十分之一來.
給墨基西德.
這個就是要和大家分享的故事.
那是什麼呢?.
墨基西德是今天的主題.
其實這段時間在第十四章十七節之前.
說到有些王.
這些王其實都是土王.
一個個地區一個個城這樣叫王.
舉一個例子.
譬如大家都聽到這個佛祖釋迦牟尼.
他是一個王子.
其實不是叫王子.
他只不過是鄭泛國的王子.
其實是一個城主的兒子.
有很大的城和周邊的地區.

$^{41}$有個人是阿頭叫王.
有王還有國家名.
其實只不過是城邦之主.
當時周圍都是很多城邦.
有些叫國王有些叫國王.
有些叫王其實都是阿頭.
當時打仗.
有五個王和四個王.
這五個王聯合打仗.
結果四王戰勝了五王.
其中亞伯蘭是幫四王的.
亞伯蘭看他這樣.
你不要以為他文藝書生.
原來亞伯蘭都能夠領兵.
而且他自己有手下的兵源.
而且他自己這個族都很好打.
他自己手下很多人.
一個富戶.
手下的家丁都遙用善箭.
他就殺敗基大老馬這個王.
和其他同盟的王.
贏了回來.
索多瑪王就出來迎接他.
索多瑪王就想叫他.
你分些你贏了的東西給我.
你可不可以給一些你搶到的東西給我.
除了這個索多瑪王.
還有一個撒冷王.
Salam.
Salam是平安.
今天猶太人見到面都是用這個字.
不過音有少少不同.
現在他們見到面不說hello.
而是說Salam.
Salam是平安.
祝福大家平安.
其實就是撒冷.
撒冷是這個名字的來源.
Salam平安王.
他的名字叫麥羈洗德.

$^{81}$麥羈洗德.
麥羈我的王是.
Jadak.
Jadak是洗德.
麥羈.
麥羈Jadak.
其實這個麥羈Jadak是什麼意思呢.
麥羈就是我的王是洗德.
Jadak.
是這個意思.
這個Jadak是什麼呢.
不知道.
應該說我查考不到.
不過整個名字的洗德意思就是我的王是洗德.
這個洗德是廣東話的洗不洗德.
但不是這個意思.
是Jadak.
Jadak.
麥羈洗德.
我的王是Jadak.
這個麥羈洗德是撒冷王.
他是至高臣的祭司.
即是說耶和華的祭司.
至高臣的祭司.
在當地應該是很有名的.
很有身份地位的.
他帶著餅和走出來迎接亞伯蘭.
一個如此至高臣的祭司.
又叫撒冷王.
即是說他在屬靈上有很高的位置.
在地方也是一個王.
他見到亞伯蘭得勝之後回來.
帶著很多獵物回來.
得勝的牆上拿到的物資回來.
他帶著餅和走出來迎接亞伯蘭.
而且他還為亞伯蘭祝福.
「願天地的主至高的臣賜福給你」.
「至高的臣將你的敵人放在你手裡」.
「是應當稱頌的」.
讚美主.

$^{121}$是他使你得勝.
「至高的臣將敵人放在你手裡」.
「是應當稱頌的」.
「至高的臣將你的敵人放在你手裡」.
「是應當稱頌的」.
「至高的臣將你的敵人放在你手裡」.
亞伯蘭一聽到後.
留意這裡.
亞伯蘭把所得的戰利品拿出來.
拿了十分之一的戰利品出來.
給默基世德.
拿出十分之一獻給默基世德.
獻給默基世德.
即是當作獻給「至高的臣」.
十分之一亦是以色列人.
後來以色列人和猶太人.
在舊約的時候.
要獻上十分之一的來源.
當然在西乃山有很多指示.
關於獻上十分之一.
但最早的十分之一的奉獻給至高的臣.
就是在亞伯蘭出來的.
獻給默基世德.
默基世德是至高臣的祭司.
所以獻給他.
就是顯示獻給至高的臣.
得勝回來.
默基世德說.
「臣應該被稱頌」.
「至高的臣幫了你得勝」.
亞伯蘭二話不說.
就把自己戰利品的十分之一獻給默基世德.
即是我把這十分之一.
是我用血贏回來的獻給主.
這是最第一次獻上十分之一.
而默基世德就重要了.
當時是一個很重要的人物.
而這個重要的人物.
他更重要的是什麼呢?.
就是他被希伯來書.

$^{161}$認為是神的代表.
是很重要的人物.
有多少像是象徵尼賽亞.
讓他知道.
這件故事也很重要.
就是亞伯蘭一見到默基世德.
都獻上十分之一.
「至高臣的祭司」.
請你在每一次得勝之後.
你要記得這句.
「至高臣把敵人交在你手裡」.
不要想著我老鼠仔多大多厲害.
這次是我的頭腦.
我的能力將對方得勝.
不是的.
是「至高臣把敵人交在你手裡」.
請你記得要稱頌神.
獻上十分之一.
獻上禮物.
這句就是創世記十四章二十節.
「至高臣把敵人交在你手裡」.
是應當稱頌的.
神為你得勝.
他使你得勝.
應當稱頌.
那我要做什麼?.
拿出十分之一來獻給神.
不要忘記對神的感恩奉獻.
祂使你得勝.
祂使你有工作.
祂使你獲得.
而你今天能夠獲得.
是神給你的.
使你獲得.
不是我的學位嗎?.
不是.
不是我的能力嗎?.
不是.
不是我的手段嗎?.
不是.

$^{201}$是「至高臣」使你獲得.
請你記住這四個字.
「至高的臣」.
將「至高的臣」.
事事擺在你的心裡,將祂擺在第一.
祂會祝福你的.
當然你要跟隨祂.
獻上你感恩的禮物.
每一天.
每一個時刻.
帶著感恩的心都是神給我的.
好,祈禱.
多謝我們的主.
萬物從你而來.
我們所有的每一樣東西.
甚至我們的才華.
我們的才幹.
我們的能力.
全部都是你的.
如果不是你給我們,我們怎會獲得呢?.
就算有很多人的才華和才幹.
是比我們高.
但我們看到他並沒有得勝.
生活也很老倒,很艱難.
其實是因為神有祝福我們.
是因為你在我們背後.
幫助我們能夠做到.
幫助我們得勝.
幫助我們甚至勝過你的敵人.
完全是你的恩典.
是我們要多謝你的.
我們願意拿出感恩的心.
又拿出我們所有得勝的財物.
我們得到的戰利品獻給主.
不會說這些屬於我的.
其實是屬於你的.
我們就將它獻給主.
願你越立.
給我們一個敬畏的心.
給我們有一個感恩的心.

$^{241}$給我們有一個心.
是將至高的神.
時時放在我們生命的首位.
誰聽我們悲鳴的祈禱.
奉靠耶穌基督命求.
阿們.
\newpage



\section{路加福音 16:19-31}
\label{sec:031672}
\textbf{今生決定永生}~~~~梁成裕牧師
\newline
\newline
連結: \href{http://www.church.com.hk/acms/content.asp?site=cdc\&op=show\&type=product\&code=031672\&layout=sermon}{\texttt{031672}} ~~~~ 語音日期: 2026-02-15
\newline
\newline
\hyperref[sec:031671]{\small{< < < PREV SERMON < < <}}
~
\hyperref[sec:index_chronic]{\small{[返順時目]}}
~
\hyperref[sec:031673]{\small{> > > NEXT SERMON > > >}}
\newline
\newline
路加福音 16:19-31
\newline
\begin{longtable}{cl}
\hline
\hline
章節 & 經文 (和合本修訂版)\\
\hline
16:19 & \begin{tabularx}{0.7\textwidth}{X} 「有一個財主穿著紫色袍和細麻布衣服，天天奢華宴樂。 \end{tabularx} \\ \\ \relax
16:20 & \begin{tabularx}{0.7\textwidth}{X} 又有一個討飯的，名叫拉撒路，渾身長瘡，被人放在財主門口， \end{tabularx} \\ \\ \relax
16:21 & \begin{tabularx}{0.7\textwidth}{X} 想得財主桌子上掉下來的碎食充飢，甚至還有狗來舔他的瘡。 \end{tabularx} \\ \\ \relax
16:22 & \begin{tabularx}{0.7\textwidth}{X} 後來那討飯的死了，被天使帶去放在亞伯拉罕的懷裡。財主也死了，並且埋葬了。 \end{tabularx} \\ \\ \relax
16:23 & \begin{tabularx}{0.7\textwidth}{X} 他在陰間受苦，舉目遠遠地望見亞伯拉罕，又望見拉撒路在他懷裡， \end{tabularx} \\ \\ \relax
16:24 & \begin{tabularx}{0.7\textwidth}{X} 他就喊著說：『我祖亞伯拉罕哪，可憐我吧！請打發拉撒路來，用指頭尖蘸點水，涼涼我的舌頭，因為我在這火焰裡，極其痛苦。』 \end{tabularx} \\ \\ \relax
16:25 & \begin{tabularx}{0.7\textwidth}{X} 亞伯拉罕說：『孩子啊，你該回想你生前享過福，拉撒路也同樣受過苦，如今他在這裡得安慰，你卻受痛苦。 \end{tabularx} \\ \\ \relax
16:26 & \begin{tabularx}{0.7\textwidth}{X} 除此之外，在你們和我們之間，有深淵隔開，以致人要從這邊過到你們那邊是不可能的；要從那邊過到這邊也是不可能的。』 \end{tabularx} \\ \\ \relax
16:27 & \begin{tabularx}{0.7\textwidth}{X} 財主說：『我祖啊，既然這樣，求你打發拉撒路到我父家去， \end{tabularx} \\ \\ \relax
16:28 & \begin{tabularx}{0.7\textwidth}{X} 因為我還有五個兄弟，他可以警告他們，免得他們也來到這痛苦的地方。』 \end{tabularx} \\ \\ \relax
16:29 & \begin{tabularx}{0.7\textwidth}{X} 亞伯拉罕說：『他們有摩西和先知的話可以聽從。』 \end{tabularx} \\ \\ \relax
16:30 & \begin{tabularx}{0.7\textwidth}{X} 他說：『不！我祖亞伯拉罕哪，假如有一個人從死人中到他們那裡去，他們一定會悔改。』 \end{tabularx} \\ \\ \relax
16:31 & \begin{tabularx}{0.7\textwidth}{X} 亞伯拉罕對他說：『如果他們不聽從摩西和先知的話，就是有人從死人中復活，他們也不會信服的。』」 \end{tabularx} \\ \\
[1ex]
\hline
\hline
\end{longtable}
$^{1}$今天我們分享的題目是.
今生決定永生.
我們人生要做很多決定.
不知道你覺得你在人生裡面做過.
其中最重要的決定又是什麼呢?.
我們最近很喜歡問AI.
於是我也問了AI.
人生十個最重要的決定是什麼呢?.
當然AI的答案不能夠作準.
但是可能我們中間也做過以下.
十個人生AI認為是最重要的決定.
我們一起看看.
究竟是哪十個很重要的決定.
第一.
是你的老年怎麼過.
以及你要留下什麼.
我想問問有沒有人想過這件事的?.
謝謝大家.
大家舉了手也可能.
我不覺得一定是年長才想這些的.
我們一早預備好自己.
我後面看到有些很年輕的也舉手.
所以這件事跟年齡無關的.
多謝大家坦誠.
第九.
如何處理傷害,腐爛和失敗.
有沒有人也要處理過這件事呢?.
也有.
大家都很坦誠.
我們看看第八.
第八是如何對待金錢和時間.
這個應該多一點.
我們每個人都要的.
當然我們在迎新十二課.
也有門徒的金錢觀.
以及門徒的時間管理.
好.
在哪裡生活.
屬於什麼社群.
有沒有人想過的?.

$^{41}$沒有人想過.
大家都很認定是不是?.
是不是?.
有些人也想過的.
特別是過去幾年.
可能有些人也想過.
會不會移民.
住在哪裡.
好.
做什麼工作.
以及如何看待職業.
有沒有人想過的?.
應該很多人想過的.
是不是?.
現在看到Pastor Irene舉手.
你還有沒有在想的?.
多謝.
我們也知道.
不是開玩笑的.
和Pastor Irene玩一下.
她沒有想的.
她很開心的.
是不是?.
昨天說的好像是溫妮基.
如何看待職業.
我們也知道.
在迎新十二課.
我們也有門徒的工作觀.
是不是?.
好.
讀什麼書.
學什麼.
以及如何持續成長.
應該大家都一定有想過的.
是的.
又可以想一下如何持續成長.
是的.
又可以報一下教會.
學校.
是不是?.

$^{81}$是的.
好.
是不是生孩子.
以及如何做父母.
有沒有人想過的?.
是的.
我們當中的父母.
一定有想過的.
是的.
那就有了.
是不是?.
好.
你選擇愛誰?.
以及如何建立關係.
有沒有人想過的?.
有了.
是不是?.
一定有的.
是的.
好.
品格和價值觀.
有沒有人想過的?.
是的.
也有不少.
是的.
大家想一下.
其實排第一是什麼.
你猜猜.
排第一是什麼?.
其實排第一.
是信仰.
大家說得對.
是的.
原來就算AI.
你問他這個答案.
他也會說人生最重要的決定.
是你的信仰.
今天你的生命.
有沒有神人活著的意義是什麼.
雖然AI說這個是人生.

$^{121}$他覺得是最重要的決定.
但是其實我們要問一下自己.
其實我們有沒有真的想過的?.
是的.
你有沒有想過?.
你有沒有想過其實永生是怎樣的?.
又或者有沒有找過.
究竟我們的永生又是怎樣的?.
但是原來我們的永生是怎樣的?.
取決於我們今生的決定.
今天我們看一段經文.
是《路加福音》十六章十九至三十一節.
其實是延續上個星期.
陳牧師所分享的一篇經文.
其實陳牧師所分享的就是.
今天究竟我們.
如何運用神給我們的資源.
讓我們的生命結出果子.
而今天這段經文.
其實是上個星期陳牧師教導的一個延伸.
以及希望為我們生命所結出的果子.
我們有沒有為永生去想.
但是今天我們又如何去決定我們的永生呢?.
今天我們進入到《路加福音》十六章十九至三十一節.
我們去學習三個很重要的功課.
好.
我們看看今天我們學習的第一個功課是什麼.
今天學習的第一個功課.
是我們要反思生命的真正意義.
我們由第十六章的十九節讀到二十一節.
我們一起讀.
一二三.
有一個財主穿著紫色袍和細麻布衣服.
天天奢華燕落.
又有一個討飯的名叫拉薩路.
渾身生湯被人放在財主門口.
叫得財主卻紙上掉下來的零碎衝擊.
並且苟來舔他的瘡.
在這裡我們看到有兩個角色.
我們會首先一起去看看這兩個角色他們的生命是怎樣.

$^{161}$但是更加重要的是我們一會兒會看看.
這兩個角色的互動裡面.
對於他們怎樣看待人生意義.
又給我們有什麼反思.
我們首先看看財主.
財主曾經形容他身穿紫色的袍和細麻布衣服.
他是天天奢華燕落.
他身穿紫色的袍.
是反映當時他有很尊貴的地位.
而他穿著的細麻布衣服是很上好的料.
大家有聽過現在買完羊毛衣服.
有那些叫marino的羊毛.
或者你買一件cotton.
現在很漂亮的叫pima cotton.
有沒有聽過有的.
即是那些很微細很漂亮的羊毛.
或者那些很漂亮的棉質.
細麻布衣服就是那些.
一點都不粗.
很漂亮很滑的衣服.
他穿著很尊貴的衣服.
很漂亮質地的衣服.
過著奢華的生活.
其實紫色袍和細麻布衣服.
被形容為帝王一般的生活模式.
有沒有人覺得現在過著帝王一般的生活模式.
有的你也可能不敢舉.
是不是我明白的.
說笑而已.
但是這個財主.
他天天奢華燕落.
但是原來他的生命.
只是追求著自己的快樂.
和今世的享樂.
我們一會兒會看的.
他沒有理會過其他人的需要.
也沒有同情過其他人的需要.
這段聖經告訴我們.
他只是追求今生的享樂.
但是他也沒有想過.

$^{201}$將來會有神的審判.
但是我們又看看另一個角色.
這個角色聖經說是土犯.
就是乞丐.
叫做拉薩路.
他渾身生瘡.
被人放在財主的門口.
拉薩路這個名字來自希伯來文.
以利以謝的意思.
耶穌稱呼這個人為拉薩路.
可能特別想拿這個名字的象徵式意義.
以利以謝直譯就是神幫助我.
就是說這個窮人.
雖然是窮.
但是神是幫助貧窮的人.
土犯其實就是一個窮人.
這個窮人滿身生瘡.
可能患了麻瘋病.
而他是被人放在財主的門口的.
他可能是瘸腿的.
因為如果自己會走就不用被人放在門口.
這個窮人.
這樣子成為一個財主很大的對比.
他被人放在財主的門口.
但是我們看看這兩個這麼大的對比角色.
他們的互動又是什麼呢?.
在第二十一節告訴我們.
這個窮人被人放在財主的門口.
他只能夠得到.
財主從他的桌子掉下來的餅碎衝擊.
並且只有狗來舔他的瘡.
就是說縱然有人將這個窮人放在財主的門口.
這個財主沒有理會他的需要.
這個財主仍然天天奢華艷樂.
但是只會在桌子上的那些餅碎.
去讓這個人去吃.
而狗在當時猶太的傳統是以為是不潔的動物.
就是代表沒有人去理會這個人.
縱然被人放在財主的門口.
沒有人去理會他.

$^{241}$只有被以為不潔的動物來可憐他.
舔他的傷口.
這個有錢的人只是繼續追求自己的享樂.
有人被放在他的門口.
他眼看不到.
沒有同情過他.
沒有好好地照顧過他.
只有那些掉下來的餅碎去衝擊.
這個財主的硬心.
將這個貧窮人留在一個飢餓和無助的狀態.
因為這個財主只是追求著自己的享樂.
所以今天也帶出我們生命一個的反思.
就是究竟我們今天又在追求什麼.
今天我們有沒有去反思生命真正的意義.
今天我們追求的是什麼.
而我們追求的東西有沒有意義.
和永生有什麼關係.
我記得我初信主的時候.
雖然我信了耶穌.
但是我的價值觀是和這個世界一樣的.
所以我記得有一次我們回教會.
坐車回教會的時候.
當時我太太.
我們那時候還沒有結婚.
那時候蘇儀是我女朋友.
她就問我.
你正在讀法律.
因為我信耶穌的時候我正在讀法律.
你正在讀法律.
你將來想成為一個什麼樣的律師.
你猜我給她的答案是什麼呢.
(賺到錢).
為什麼你們這麼明白我.
當她這樣問我的時候.
我想成為一個什麼樣的律師.
我想成為一個最賺錢的律師.
我感覺到她對於我這個答案是很失望的.
但是對於我來說我覺得很理所當然.
我那時候讀法律也是.
都是想有穩定的收入和舒適的生活.

$^{281}$讀了法律做律師.
不做最賺錢的律師.
做什麼律師呢.
Helen 是不是.
做什麼律師呢.
所以她這樣回答我.
我真的感覺到她有點失望.
於是她就繼續問我.
她沒有放棄.
那你自己沒有興趣的嗎.
你最喜歡哪一科的法律.
你不喜歡去做.
你喜歡那一科的法律律師的嗎.
我要坦誠地承認.
我最喜歡的科不是最賺錢的科目.
最賺錢的當然是做商業律師.
其實我讀書的時候.
我最喜歡的是家庭法.
Family Law.
但是我知道做那些.
沒有什麼賺錢的.
於是她這樣問我的時候.
我就回答了一個.
可能很多父母都會跟你的兒女說一句話.
興趣?.
興趣可以當飯吃嗎.
但是她還沒有放棄.
她說那你信了耶穌.
你不想你以後做一個律師.
是做一個可以幫人的律師嗎.
你知道我怎樣回應她嗎.
所以今天我太太願意嫁給我.
也是恩典.
我嘲笑她.
幫人?.
開玩笑的吧.
不好玩的.
我信了耶穌.
但是我的價值觀是跟這個世界一樣的.
你問我.

$^{321}$我仍然在讀法律.
我信了耶穌.
但是我仍然在讀法律.
你問我想做什麼的律師.
哎呀.
另外兩位律師.
我當然想做最賺錢的那一種.
你還問我有什麼興趣.
還問我是不是想幫人.
幫人.
不好玩的.
但是我們今天要想一下.
生命真正的意義是什麼.
那時候的我.
我沒有想過.
這個世界追求錢.
我當然想做最有錢的律師.
我認識一位弟兄.
有一次跟他聊天.
跟他分享.
他說他未信耶穌之前.
都追求很多這個世界的東西.
他分享了兩件事.
第一件事就是.
他未信耶穌之前.
他追求什麼.
他追求去朋友堆之間.
去那些最流行.
最流行.
最貴的地方吃飯.
他就是追求那些.
總之朋友堆裡面.
哇.
哪家餐廳.
現在最流行.
最流行.
最貴的地方.
他就會去吃.
我就問他一件事.
喂.

$^{361}$那些最流行.
最貴的餐廳.
是不是真的很好吃.
他回答我.
Eugene.
好不好吃.
都是其次.
我說哇.
你去那些最流行.
最流行.
最貴的地方吃飯.
好吃都是其次.
那你追求什麼.
他說你不明白.
其實當時我的心態是追求什麼.
我真的想聽一下.
你去的地方不是為了好吃.
那是追求什麼.
他說追求就是能夠在朋友堆之中.
當他們分享現在最流行.
最貴的餐廳的時候.
我能夠很威風地在那些朋友面前說.
那家餐廳.
我去過了.
原來他要追求那種滿足感.
我不知道你有沒有試過.
跟別人聊天說.
那家餐廳很貴.
很難預約的.
然後你就說.
我去過了.
但是他說另外有一個更加大的滿足感.
我說什麼.
說來聽聽.
他說當朋友這麼說.
你不僅跟他說你去過.
如果你有那些關係.
你很帥氣地跟朋友說.
你也想去吃嗎.
我幫你訂位.

$^{401}$他說原來那種滿足感.
真的當時覺得很帥氣.
我不知道你有沒有覺得很帥氣.
你幫我訂位可以嗎.
幫我們付錢也好.
原來他以前是追求這些的.
但是他說信了主之後.
他在那家教會也做了很多關心基層的愛宴.
現在反而去參加一個.
跟基層一起吃飯的愛宴.
他覺得很滿足.
哇.
你變化也很大.
以前追求那些.
現在服侍和街坊吃愛宴.
你很滿足.
另外他以前很追求的是什麼呢.
就是追求賣帥的車.
你又說你賣帥的車.
你什麼心態.
他也是沒有的.
最重要的是什麼.
別人覺得我開的車出來帥不帥.
是的.
我很坦白我沒有這種掙扎.
因為為什麼呢.
最帥最貴的那些車.
最帥的那些.
我進不去.
我沒有這個掙扎.
你現在買了最帥的車.
是不是我也進不去.
是的我進不去.
所以我沒有這個掙扎.
是的.
但是他說他以前就是追求這些.
但是他信了主之後.
他現在專門買什麼呢.
買一些方便大駕的車.
他最喜歡就是開著他的七人車.

$^{441}$去參加弟兄姊妹.
開著他的七人車去送人回家.
你問我.
我也純粹看到你們是這樣的.
你知不知道.
尤其是晚上的時候.
是門女.
有時候門女晚上.
我看到你穿的帽子很好.
開著大車.
送人回你們不同區那裡.
或者是回自己同區那裡.
但是我看到這裡也是.
廣島東小隊有一輛車.
是不是.
九龍東小隊有一輛車.
然後另外不知道哪裡.
新界沙田有一輛.
你問我.
我覺得很感動.
你問我就是這樣.
原來我們信了耶穌.
我們的生命要反思生命的意義.
我們不再是追求那些世界.
告訴你.
令你很不滿足的那些.
因為那些不能夠帶給你永生的意義.
所以今天反過來.
你問我.
我又很為這位的肢體感恩.
看到他的改變.
看到他現在的生命.
他真的知道.
他要去幫助有需要的人.
所以今天到了第二個功課.
我們要去想.
就是我們怎麼去找機會.
為我們的永生去做決定.
我們從第二十二節讀到二十六節.
好不好.

$^{481}$這裡是在說.
一個財主和窮人.
在永生裡面的狀況.
好不好.
我們一起從二十二到二十六節.
我們一起讀.
一二三.
後來那土犯的死了.
被天使帶去放在亞伯拉罕的懷裡.
財主也死了.
並且埋葬了.
他在陰間受痛苦.
舉目遠遠地望見亞伯拉罕.
又望見拉薩路在他懷裡.
就喊著說.
我祖亞伯拉罕啊.
可憐我吧.
打發拉薩路來.
用指頭沾焦點水.
涼涼我的舌頭.
因為我在這火焰裡.
極其痛苦.
亞伯拉罕說.
兒啊.
你開回想你生前享過福.
拉薩路也受過苦.
如今他在這裡得安慰.
你倒受痛苦.
不但這樣.
並且在你我之間.
有深淵限定.
以致人要從這邊.
過到你們那邊.
是不能的.
要從那邊.
過到我們這邊.
也是不能的.
這裡告訴我們.
兩個都死了.
而這兩個角色.

$^{521}$在永恆裡面.
他們的狀況又是怎樣呢?.
一開始二十二節.
後來那土犯的死了.
那黑衣窮人死了.
但是他是被天使.
帶去放在亞伯拉罕的懷裡.
在舊約裡面.
我們也讀過.
以洛和以利亞.
是被神接走了.
天使的工作代表這是神的作為.
這個窮人被天使帶走.
代表在永恆裡面.
神對他的保守和眷顧.
在永恆裡面.
這個窮人在亞伯拉罕的懷裡.
亞伯拉罕的懷裡.
是代表上帝的拯救.
在《路加福音》十三章這段聖經說.
將來被神拯救的人.
是在神的國裡面.
和亞伯拉罕一同坐直.
代表他得著神的拯救.
這個窮人得著神的拯救.
並且在神國的福分.
和喜樂裡面有份.
在永恆裡面.
這個窮人.
他是能夠享受天家的福分.
永恆的福樂.
雖然在地上.
這個人無人紀念.
但是到了永恆.
他是得著神的拯救.
和神國的福分.
但是財主呢?.
財主這裡說他被人埋葬.
就是有人紀念他的.
但是在永恆裡面.

$^{561}$他的狀況又是怎麼樣呢?.
原來這個財主在永恆裡面.
是在陰間受苦.
陰間是死後惡人被神處罰的地獄.
我們發現這個財主.
和這個窮人去了不同的地方.
這個窮人去了神的國.
在亞伯拉罕的懷裡.
但是這個財主.
他去了陰間地獄受苦.
這個財主在世界的時候.
過著很奢華的生活.
但是在永恆裡面.
他就在陰間受苦.
這個財主在地上積聚的財富.
在地上積聚的影響力.
對他的永生毫無幫助.
今天我們也要提醒自己.
我們在世界上積聚的財富.
我們在世界上積聚的影響力.
對我們的永生毫無幫助.
所以這個財主在永恆的處境.
和他在世上的處境是剛剛相反.
這個窮人死了之後.
能夠在神國裡面.
在亞伯拉罕的懷裡.
但是這個財主他在地獄裡面受苦.
然後到了第24節.
我們看到這個人他就哭著說.
我祖亞伯拉罕可憐我吧.
打發拉薩路來.
叫他用指頭沾一點水.
涼涼我的舌頭.
為什麼?.
因為我在這個火焰裡面.
極其痛苦.
這個時候讓我們想起.
《路加福音》第三章這段聖經.
這段聖經說.
你們要結出果子來.

$^{601}$與悔改的心相稱.
不要自己心裡說.
由亞伯拉罕為我們的祖宗.
今天我們的生命要去為神結果子.
我們的生命要悔改回轉.
如果不是.
將來原來去到陰間.
那裡是永恆的火.
極其痛苦.
在以賽亞書也這樣形容.
在陰間的狀況.
因為他們的從事不死的.
他們的火是不滅的.
凡有血氣的都必憎恨他們.
當這個財主說.
亞伯拉罕你可憐.
走到這裡很熱很辛苦.
你叫拉薩路給一點水.
沾一沾我的舌頭.
亞伯拉罕的回應是什麼?.
他說二.
你應該回想你生前享過福.
拉薩路也受過苦.
如今他在這裡得安慰.
你倒受痛苦.
這個財主在世界上追求一切的享樂.
他死後很想神你施恩.
但是已經太遲了.
因為已經進入了那個不能逆轉的結果.
這個不能逆轉的結果.
就在二十六節告訴我們.
亞伯拉罕跟他說.
他說我幫不了你.
因為在你我之間.
有深淵限定.
以致人在這一邊.
就是享福的人.
和在陰間的人.
他說這一邊的人.
不能夠去到你們那一邊.

$^{641}$而你們這一邊的人.
都不能夠回到我們這一邊.
一個是不能逆轉的結果.
所以今天我們人要在我們.
還能夠去為我們永生做決定的時候.
為我們的永生.
好好地預備.
在我們的生命裡面.
有時候我們都會做一些決定.
會帶來不能逆轉的後果.
我曾經認識一位少年人.
當時他找暑期工.
他找暑期工的時候.
他被人吸引了去那些找換店工作.
其實當時他去的時候.
我們牧者.
我 Pastor Joyce 陳牧師都跟他說.
你不要去那些地方工作.
因為那些地方很多時候會涉及犯法的行為.
當時其實陳牧師很好.
他幫這位少年找了一份很好的暑期工工作.
人工對暑期工來說是十分之好的.
替他安排好.
安排他去上班.
他真的去上班了.
但原來他之後偷偷地.
又被人引誘了去找換店工作.
我們不知道.
有一天我收到一個電話.
他家人打來.
他說他被警察抓了.
因為被人告他洗黑錢.
原來他去找換店工作.
他的戶口被人利用了.
去洗黑錢.
接著很可惜的就是因為.
犯法的時候他剛剛過了18歲.
要接受的是成人的判刑.
以及成人的入獄.
一個的決定.

$^{681}$帶來了不能逆轉的結果.
所以我們今天都要問自己.
我們有沒有好好抓住機會.
去為我們的永生做決定.
有時候機會是我們要趁還有得抓住的時候抓住的.
我記得我奶奶經常跟我說.
以前打仗的時候一個故事.
我奶奶經常跟我說.
當香港淪陷的時間.
有一天我爺爺就被日本軍抓了.
只是抓了我爺爺而已.
但是接著他的親戚就跟我奶奶說.
你現在要走了.
你快點帶著你三個兒子.
當晚立刻走.
他說你今晚不走.
明天那些日本軍過來抓你們幾個了.
但是當時我奶奶很忐忑.
突然間老公被人抓了.
一個人帶著三個兒子.
我爸爸和我兩個伯父.
一個人帶著三個兒子可以怎樣.
但是親戚說你現在什麼都不要想.
你立刻收拾東西.
立刻就走.
我奶奶跟我說.
她很記得那一晚她跟對面有一個太太聊天.
對面那個太太也一樣.
她丈夫被日本軍抓了.
也是帶著兩個兒子.
但是接著她說.
我今晚決定走了.
但是對面那個太太說我能走到哪裡去.
我老公被人抓了.
我不行了.
我要等我老公回來.
我奶奶就說等不到你老公回來.
你今晚就跟我一起收拾東西一起走.
但是對面那個太太不願意.
那一晚我奶奶收拾東西.

$^{721}$帶著我兩個伯父和我爸爸離開了家.
她沒有再見過我爺爺.
我爺爺也是在戰爭的時候.
被日本軍殺了.
但是接著她跟我說.
原來她知道對面的太太她沒有走.
第二天.
原來日本軍就上來.
連她太太和兩個兒子都抓了.
我不知道最後他們的結果是怎樣.
不過很可能他們都是一家喪命.
原來我們是要趁我們還可以做決定的時候.
為我們的永生做好決定.
當我們為我們的永生做決定的時候.
我們考慮的是什麼呢?.
這個時候我很想分享一下.
有一次Pastor Ray 和我和陳牧師.
一個我覺得很反省自己內心的分享.
大家知道Pastor Ray 他很努力.
不僅是星期天在英語崇拜.
他其他日子.
他星期六又會有第二個崇拜.
然後他又會帶著很多他們的弟兄姊妹去做門徒訓練.
有一次他來跟我和Pastor John 聊天.
我們在聊天的時候他就跟我說.
那時候是暑假.
大熱天時.
他跟我們說.
他說他找了一個公園.
和他的弟兄姊妹做門徒訓練.
大熱天時.
暑假.
在公園那裡做門徒訓練.
於是我和陳牧師說.
Oh Pastor Ray.
So hot.
And you do discipleship in the park.
這樣.
但是接著Pastor Ray給我的回應 給我和陳牧師很大的反省.
他說 Hot.

$^{761}$Hell is hotter.
你說地獄更加熱.
就是說我們要想清楚.
當我們.
你問我他有什麼回應.
我和陳牧師真的很欣賞他.
不過星期六因為我們已經沒有位置再給他.
因為星期六我們是120\%使用率.
我們都沒有辦法再給他.
於是他找了一個公園去做門徒訓練.
但是你問我他給我們的回應很發人深省.
當我們真的想的時候.
我們有沒有想過.
如果我們沒有為我們的今生做決定.
我們的將來是怎樣.
Hell is hotter.
地獄更加熱.
所以今天當我們去想我們的永生的時候.
其實承接陳牧師上個星期的教導.
我們要問自己.
今天我們的人生意義是什麼.
剛才我們讀了一段經文.
我們的生命有沒有為神去結果子.
我們的生命有沒有改變.
當時聽的人特別有兩種人.
第一種是財主.
我相信耶穌想跟那些財主說.
今天神給你的資源.
你有沒有為神國好好做神國的管家.
但是另外一種人.
就是法理賽人.
法理賽人覺得自己很熟悉律法.
但是他們錯過了一件事情.
就是哪一個是救主.
我們待會讀第三個就明白了.
原來今天我們要認清.
哪一個是我們的救主.
並且我們的生命有沒有真正跟隨這位救主.
我們要去想.
我們要為我們的永生去做決定.

$^{801}$我們的生命有沒有真正結果子.
我們的生命有沒有真正去跟隨耶穌.
所以就帶入我們今天的第三點.
就是我們要竭力履行大使命.
我們由第27節讀到31節.
我們一起讀.
1,2,3.
財主說:我早啊,既是這樣.
求你打發拉薩路到我父家去.
因為我還有五個弟兄.
他可以對他們作見證.
免得他們也來到這痛苦的地方.
亞伯拉罕說:他們有摩西和先知的話可以聽從.
他說:我早啊,亞伯拉罕啊.
不是的.
若有一個從死裡復活的到他們那裡去.
他們必要悔改.
亞伯拉罕說:若不聽從摩西和先知的話.
就是有一個從死裡復活的.
他們也是不聽勸.
我們看這段聖經.
我們發現這個財主.
他在絕望裡面.
他知道他自己不能逆轉.
但是他仍然跟亞伯拉罕說.
我還有家人.
你可不可以叫拉薩路去我的家人那裡.
我的家人看到拉薩路.
他們會悔改的.
但是問題是什麼呢?.
其實神已經將他的獨生子.
差遣到了這個世界.
猶太人看到神的獨生子.
他們仍然不會去悔改.
當他們說:如果你要拉薩路去的時候.
他會悔改.
但是有一個比拉薩路更大的來到.
他們的生命都沒有回轉.
所以當我們去看這段聖經.
亞伯拉罕就說.

$^{841}$其實他們是有摩西和先知的話可以聽從的.
這個財主的家人應該跟他一樣.
天天過著奢華的生活.
沒有想過貧窮人的需要.
當時猶太的領袖覺得自己很認識聖經.
很認識律法.
但是他們看不到原來律法所指.
將來要來的救世主就是基督.
他們沒有接納到耶穌.
但是今天我們的生命.
我們要相信耶穌.
並且跟隨耶穌.
這個財主再跟亞伯拉罕說.
亞伯拉罕不是這樣的.
如果有一個從死裡復活.
去到他們那裡.
他們必定會悔改的.
那個從死裡復活的是誰?.
就是耶穌.
耶穌已經叫過死人復活.
自己也從死裡復活.
但是他們沒有接受過耶穌.
之前我們讀過一段聖經.
那些猶太人說.
你如果顯示一個神蹟給我看我就相信了.
然後耶穌跟他說.
除了若拿的神蹟.
再沒有其他神蹟.
我之前講過什麼叫若拿的神蹟.
若拿曾經被一條大魚吃過.
然後吐出來.
其實是預表耶穌從死裡復活.
有一個從死裡復活的人.
被神拆顯到我們當中.
但問題就是我們有沒有相信耶穌.
我們有沒有悔改.
以及真正跟隨耶穌.
所以亞伯拉罕作出了一個這樣的回應.
他說:若不聽從摩西和先知的話.
就是有一個從死裡復活的.

$^{881}$他們也是不聽勸.
今天我們覺得.
哇!神啊!如果你有神蹟奇事發生我就相信了.
神啊!如果見到這些奇蹟我就相信了.
但是其實我們已經有神的真理和教導.
但是今天如果我們不去跟隨神的真理和教導.
就算再有大的神蹟.
我們也不會相信的.
所以亞伯拉罕這個回應.
有兩個很重要的訊息給我們.
今天作為神的子民.
如果我們不忠心去走出神的話.
就算有一個從死裡復活的人來到.
我們也不會相信的.
但是反過來,當你願意去進行天父的教導.
你就能夠明白耶穌從死裡復活的意義.
所以今天這段聖經最後.
給我們一個很重要的呼召.
給我們一個很重要的教導.
是我們的生命要去履行大使命.
這個財主很想有人跟他家人傳福音.
但其實最重要的是他自己也沒有真的去行出生命.
所以今天原來我們要去行出大使命.
聖經說我們要去洗萬民作耶穌的門徒.
奉父子聖靈的名.
給他們施洗.
凡耶穌所教導我們的.
都教訓他們遵守.
所以今天我們每一個人.
我們都要努力去履行大使命.
不僅是我們信耶穌.
我們要跟隨耶穌.
我們也要令我們身邊的人相信.
並且跟隨耶穌.
最近在哈該牧區.
我們舉行了一個很別開生面的報道會.
叫做防騙講座報道會.
這個報道會是怎樣的呢?.
首先我們邀請了警方.
來跟我們分享一些防騙的信息.

$^{921}$我們收到了很多紀念品.
就是那張提子.
不知道大家有沒有看過.
那張很有趣的提子.
叫做提防騙子.
我覺得警方現在公共關係的講座做得十分好.
很好笑.
很幽默.
但是很好笑和幽默的背後.
有很好的信息.
教導我們不要被騙.
整個警方分享的過程裡面.
很多很輕鬆,幽默和搞笑的情景.
但是很多時候他們特別會分享一些影片.
他們分享一些影片的時候.
我經常聽到下面的人會這樣回應的.
哪有這麼蠢的.
早就告訴你了.
你又不知道.
譬如說有一個例子是什麼呢?.
他們有一條影片.
就是說有人會打給那些老人家的.
通常那些人打給老人家.
如果打給婆婆.
他就會說喂!媽!.
如果打給公公.
他就會說喂!爸爸!.
通常就會發生什麼事呢?.
我想問如果今天有人打電話給你.
而你是你媽媽.
你打給他的時候.
喂!媽!.
那你會怎麼說?.
我報警.
是的.
如果打電話.
喂!爸爸!.
那你會怎麼說?.
也是報警.
那個影片是這樣的.

$^{961}$那些人打給婆婆了.
婆婆就說喂!.
然後對面的家人就說.
喂!媽!.
於是媽媽就說.
是不是阿文啊?.
第一件事情就是如果有人這麼說.
你首先不要.
譬如說我打給Brian叔叔.
喂!爸爸!.
你是不是阿Ray啊?.
那你就將你的資料告訴別人.
你將你的資料告訴別人.
所以第一就是你不要將資料告訴別人.
但是接著那個影片裡面.
那個人就說.
是的!.
我是阿文啊!.
媽!.
我被警察抓了.
現在需要保釋.
那怎麼辦?.
要多少錢啊?.
要五千元啊!.
那怎麼給你啊?.
待會會有一個律師行的代表來找你的.
那你找你的時候.
你就將錢交給他吧.
接著電話就會交給另外一份.
你好!.
我是律師行的代表肥成.
是的.
接著他還問.
婆婆!.
請你將你的地址告訴我吧.
是的.
接著你就將你的地址告訴別人.
就是你剛剛告訴別人你兒子叫做阿文.
接著你就將你的地址告訴肥成.
過了一個小時之後.

$^{1001}$有一個穿得很斯文的人叫做肥成.
就會去到門口.
你是不是阿文的媽媽啊?.
我是律師行的.
你那五千元準備好了沒有?.
準備好了!.
那他就將那五千元給了肥成.
肥成是編出來的.
大家知道我叫肥成.
接著就走了.
但是到了晚上.
那個真的阿文就回來了.
媽媽!我回來吃飯了.
兒子!你出來了嗎?.
什麼出來了?.
你不是被警察抓去保釋了嗎?.
什麼保釋?.
那當然是被人騙了.
但是很多時候這些視頻.
我聽到後面很多人就會回應.
哪有這麼蠢的.
早就告訴你了.
你沒有聽過嗎?.
但是有時候我看這些視頻的時候.
我也覺得有時候我們的生命.
也真的被人騙了.
我們的永生也被人騙了.
為什麼呢?.
其實早就有人來告訴我們.
我們怎麼為我們的永生作決定.
那個就是耶穌.
耶穌早就來了.
並且告訴我們.
怎樣我們要為我們的永生作決定.
但是為什麼我們常常都.
好像聽不到這樣的呢?.
因為我們都被這個世界騙了.
我們都被這個世界告訴我們.
我們要追求不同的東西騙了.
以致到其實有人早就告訴我們了.

$^{1041}$所以你跟旁邊那個說.
你不要被人騙.
不要被人騙了你的永生.
是的.
當然錢也不要被人騙了.
不要被人騙了你的永生.
所以很重要.
今天我們要問自己.
當我們說我們要為永生作決定的時候.
究竟我們有沒有真的去認真想清楚.
我發現這個世界.
很多的決定我都會認真想清楚的.
我想問一個問題.
如果你接下來去旅行.
訂什麼酒店呢?.
你會不會先想清楚的?.
你會不會做很多的.
那些叫做什麼?.
就是很多不同的研究.
或者問人我資料的.
你會不會就這樣.
你去一個地方.
然後我跟你說你住Hotel A.
那你就去了.
但是有時候這麼簡單的東西.
我們都會用這麼多心思去考慮清楚.
為什麼我們沒有為我們的永生.
去考慮清楚.
還要耶穌已經一早來了告訴我們.
所以今天讓我們一起.
我們要為我們的永生.
作好決定.
在總結的時候我們請敬拜一堆.
其實這段聖經.
不是在說今天你是有錢還是貧窮的.
不是說今天如果你有錢你就不能夠去天國.
你貧窮才能夠去天國.
不是.
是在說我們的生命態度.
以及我們生命的決定.

$^{1081}$最重要今天無論我們富有或者貧窮.
我們有沒有決定相信.
並且真正跟隨耶穌.
讓我們今天一起.
去為我們的永生.
做一個重要的決定.
今天選了一個.
因為昨天是情人節.
我們也選了一首向神表達愛的詩歌.
這首詩歌不單單向神表達愛.
更加讓我們一起立志.
說我們願意將一生交給主.
因為這個是最美好的決定.
並且當我們這樣做的時候.
我們知道是我們一生最重要.
並且是最明智的決定.
要請大家打字起立.
我們就用這首詩歌.
去回應今天神給我們的訊息.
讓我們一起跟神說.
我願意將生命交給你.
並且一生愛你.
\newpage



\section{約伯記 39:19-25}
\label{sec:031673}
\textbf{無限馬力}~~~~曹子光牧師
\newline
\newline
連結: \href{http://www.church.com.hk/acms/content.asp?site=cdc\&op=show\&type=product\&code=031673\&layout=sermon}{\texttt{031673}} ~~~~ 語音日期: 2026-02-15
\newline
\newline
\hyperref[sec:031672]{\small{< < < PREV SERMON < < <}}
~
\hyperref[sec:index_chronic]{\small{[返順時目]}}
~
\hyperref[sec:031674]{\small{> > > NEXT SERMON > > >}}
\newline
\newline
約伯記 39:19-25
\newline
\begin{longtable}{cl}
\hline
\hline
章節 & 經文 (和合本修訂版)\\
\hline
39:19 & \begin{tabularx}{0.7\textwidth}{X} 「馬的力量是你所賜的嗎？牠頸項上的鬃是你披上的嗎？ \end{tabularx} \\ \\ \relax
39:20 & \begin{tabularx}{0.7\textwidth}{X} 是你叫牠跳躍像蝗蟲嗎？牠噴氣之威嚴使人驚惶。 \end{tabularx} \\ \\ \relax
39:21 & \begin{tabularx}{0.7\textwidth}{X} 牠用蹄在谷中挖地，以能力歡躍；牠出去迎擊仇敵。 \end{tabularx} \\ \\ \relax
39:22 & \begin{tabularx}{0.7\textwidth}{X} 牠嘲笑懼怕，並不驚惶，也不因刀劍退卻。 \end{tabularx} \\ \\ \relax
39:23 & \begin{tabularx}{0.7\textwidth}{X} 箭袋在牠身上錚錚有聲，槍和短槍閃閃發亮。 \end{tabularx} \\ \\ \relax
39:24 & \begin{tabularx}{0.7\textwidth}{X} 牠震顫激動，將地吞下 ；一聽角聲就站不住。 \end{tabularx} \\ \\ \relax
39:25 & \begin{tabularx}{0.7\textwidth}{X} 每逢角聲一響，牠說：『啊哈！』牠從遠處聞到戰爭的氣息，聽見軍官如雷的吼聲和吶喊。 \end{tabularx} \\ \\
[1ex]
\hline
\hline
\end{longtable}
$^{1}$多謝各位朋友,早安.
午安.
是不是還沒睡醒.
首先祝大家新年快樂.
龍馬精神,馬到攻城,好不好.
不是敷衍,都是真心祝福大家.
馬年,其實聖經裡面有很多地方講馬.
其中一處聖經裡面比較少人講.
但其實最詳細,最正面的經文.
就來自約伯記第39章.
這裡是最詳細,聖經裡面最正面描繪馬的.
小時候,人們就叫我做馬仔.
不說什麼原因.
我讀完中學之後.
大家都知道我成績是非常非常之差.
我就很孖孖呼呼的畢業了.
出來工作之後,我就做了室內設計.
沒人再叫我馬仔.
因為我已經成為了老闆的馬仔.
讀完中學之後,沒多久我就信了主.
在教會裡面,做事上面.
我都叫做一馬當先.
因為只要教會裡面有些特別的地方.
我都一馬當先去認頭去做.
所以我年輕的時候已經是教會裡面一個帶頭的人物.
到後期做了傳道人.
我的理念就是我們要討神喜歡.
我不會做一些拍馬屁,拍馬頭.
那些討人歡喜的事.
事奉,我就絕對不孖孖呼呼.
大家都看到,我都知道.
我不是一個孖孖呼呼的人.
不喜歡做一些走馬看花表面的事.
我們做傳道人要學習.
就要以主耶穌馬首是瞻才行.
雖然人老了.
但是我們都老馬有火.
希望大家馬年的時候.
就真的沒信主就馬上信主.
信了主的要馬上愛主.

$^{41}$大家都不知道有沒有經過一個唐人鋪.
在大學區附近.
我昨天有機會進去裡面吃了東西.
我就看到唐人鋪裡面有些農曆年的廢春.
上面有一句廢春.
我相信很多香港人都一定很喜歡聽.
馬上有錢.
喜歡不喜歡聽啊你.
不知道好事還是不好事.
馬上有錢.
如果是馬上我們可以很多東西.
馬上結婚.
獨裁者馬上下台.
甚至考車馬上考到車.
馬上入到籍.
雖然這些部分都可能是笑話.
但是神在約伯記裡面所說的有關馬.
就不是這麼好笑.
是一些很認真的東西.
當時約伯是失去了財產.
女兒 兒女們.
健康.
對他來說.
上帝所做的東西是超出了他所認識的東西.
於是他很希望.
好像法庭一樣跟上帝對質.
他要求上帝給一個合理的解釋給他.
但當然上帝沒有這樣做.
反而上帝在旋風裡面跟他說話.
他沒有解釋苦難的原因.
沒有交代背後究竟發生什麼事.
神反而帶他暢遊了他創造和維持宇宙的太空館.
他帶了他去飽覽太空館.
他帶他去看清楚這個宇宙有多麼的浩瀚.
是約伯你又認識多少.
又知道多少東西.
然後帶他去看動物園.
看看神怎樣令這些野生的生態.
在沒有人照顧之下.
仍然因為神可以繼續生存.

$^{81}$甚至再去看水族館.
看神給約伯看到一些神秘的大怪獸.
上帝要約伯看清楚.
在這個混亂苦難的人生裡面.
宇宙背後其實仍然都有一位智慧掌管的.
有權力的上帝.
而三十九章第十九到二十五節.
是屬於動物的段落裡面的高潮之一.
在這段經文前面提到山羊,鹿,野牛.
之後有提到鷹.
而這裡提到馬是最特別的.
因為馬和人類的文明,戰爭,權力是拉上關係.
而這些地方人以為可以控制.
但是原來上帝讓我們看到.
也是讓約伯看到.
原來人並不能夠駕馭.
這段聖經和大家在馬年的時候.
分享一個題目叫做無限馬力.
我們先看一段經文第十九到二十節.
第一件事就是馬力的源頭是來自上帝.
我們一起讀一讀這兩節經文.
馬的力量是你所賜的嗎?.
他頸上的鐘是你披上的嗎?.
叫他跳躍像蝗蟲嗎?.
他噴起之威嚴使人驚惶.
馬的力量這個字.
希伯來文這個字.
其實是指到能力,勇猛或者是勇氣.
這個字通常都是用在神又或者是勇士身上.
較為少用在動物.
但是這裡約伯記曾經有兩丑罕有地.
就將這個字放在動物那裡.
四十一章第四節或者希伯來文的版本是四十一章第十二節.
那裡是說到鱷魚或者翻譯做怪獸都可以.
而今天的這段經文是另一個罕見的例子.
那裡說到馬,即是今天的經文馬.
在聖經的其他地方如果要強調動物的力量.
例如獅子的力量,牛的力量.
他會用另一個字眼.
但是約伯記這裡他故意選用了這個字.

$^{121}$放在馬和巨獸的身上.
上帝將這個字套用在動物身上.
這樣約伯看得到,明白馬這種動物擁有的力量的源頭.
不是人可以給到,而是來自上帝.
也不是人類可以完全駕馭控制得到.
神說到馬的時候,大家不知道有沒有見過馬呢?.
馬的頸上有很多鬚,有些很長的毛.
中文譯做鐘這個字,鐘毛.
不單止是為了好看,更加其實有些實質的用途.
其中一個用途就是保護性質的用途.
當一些動物咬馬的時候,牠想咬牠的頸.
毛就幫牠擋住,令牠沒辦法很快咬到很致命的位置.
甚至乎當馬奔走的時候,有時候被樹枝,障礙物打到牠.
這些毛就可以形成一個保護的層面.
還有毛呢,馬的毛還可以幫牠排水.
當雨下的時候,牠不用擔雨傘.
因為馬的毛可以幫牠自動散水.
還有就是這些毛本身不單止可以散水.
甚至乎可以防止冬天的時候熱力流失得更快.
還有很多時候,當動物有昆蟲.
有些小動物滾到牠身上的時候.
只要牠拔掉毛,動物就會把牠的毒昆蟲全部拿走.
養過貓的人就會知道,貓是喜歡互相舔的.
大家有沒有印象?貓會互相舔對方的毛,牠們的身體.
馬也是一樣,馬也會舔對方的身體.
特別是牠們毛的地方,為什麼呢?.
因為這個是代表牠們的社交,社交的禮儀.
所以原來中毛有這麼特別的用處.
不過原來希伯來文這個字,毛這個字,媽這個字.
其實這個字的意思是振動,或者是抖震,甚至是雷轟.
這裡不單止是說一些實體的毛.
更加是描繪馬的昂首,那種威嚴.
頸項上的毛充滿著生命力.
這些生命力,當馬在走的時候就在震動.
而這些震動給人的感覺就是一種威嚴的氣勢.
上帝問約伯,是不是你讓牠披上這麼威武的中毛?.
這裡就給人的感覺就是,原來上帝親手創造馬之餘.
也幫牠披上這些沒有加工成份的這麼特別的傑作.
約伯你懂多少?你知道多少?你又明白多少?.
儘管人類可以騎馬,用馬,甚至某個成份可以控制馬.

$^{161}$但是馬與生俱來的那一種的本能,那一種的體魄.
其實是源於上帝.
上帝想告訴約伯,如果連一匹馬的生命火花和牠的威勢.
都超乎你的理解的話,整個宇宙,整個宇宙的智慧.
其實你又可以知道多少?你有什麼資格質疑我上帝?.
是啊,上帝就問約伯,是你讓這些頸項披上雷鳴般的威嚴嗎?.
這個是一個主權的宣告,馬的力量不是鍛鍊出來的.
是上帝給的,是上帝披上去的.
我們時時以為馬力就是等於智力.
我們以為馬力就是能力.
所以我們時時都把能力當成自己的資本.
但真正的動力其實是來自神所賜的定位和恩典.
當我們死撐,要生產自己沒有的力量的時候.
最終只會引來我們自己的崩潰.
反而當我們承認自己的無力的時候.
才開始學會去依靠神.
我想你另外注意二十節.
二十節那裡說,是你叫祂跳躍像蝗蟲嗎?.
祂噴氣的威嚴使人驚惶.
中文字用了二十多個字去描繪.
但希伯來文其實只有五個字就夠了.
神用了蝗蟲來比喻馬的跳躍.
跳躍這個字很特別的,希伯來文叫Vaash.
這個字的動詞是指震動或者是搖動.
你一聽震動搖動的時候.
剛才十九節的結束不是說馬的毛.
馬的毛的意思是什麼呢?.
就是指雷動或者是震動.
現在第二十節一開始第一個字.
跳躍這個字叫Vaash.
又是指震動又是指搖動.
這種在文學上叫做壓頭韻.
學者通常認為這兩個字本身在古代的用語上.
是震動性的聯想走在一起.
一個聽起來是聽覺上.
就是雷撼的聽覺的聲音,搖動的聲音.
一個就是視覺上的跳躍,蝗蟲跳躍的震動.
現在形容馬,換句話說.
連馬這種受造物身上都好像有神顯的一種狀態.
神顯通常都有這樣的狀況出現.

$^{201}$原來這種雷霆萬軍的狀態.
在馬的身上都可以見到.
約伯,你知道多少.
你連馬都駕馭不了.
你憑什麼去質疑上帝.
留意這裡用了一種帶有情感的未來式的表達.
我們指馬的跳躍就好像蝗蟲一樣.
大家有沒有見過蝗蟲.
知不知道蝗蟲是什麼.
蝗蟲原本就是一些草蜢的東西.
然後經過某一個季節或者變後變化.
牠就變成了蝗蟲.
蝗蟲的腿是特別粗壯的.
好像彈弓一樣.
科學家就研究.
他說原來蝗蟲可以跳到自己的體高十至二十倍這麼高.
當然馬是沒有可能的.
你什麼時候見到一隻馬去跳躍的時候.
牠可以跳高自己身體十幾二十倍.
是沒有可能的.
所以這裡重點不是高度問題.
而是那種爆發力.
那種興奮.
甚至乎是利用蝗蟲爆炸的跳躍力.
去帶出馬出征的時候.
牠那種激烈那種興奮.
那種充滿能量的動態.
不單止這樣.
下半節他就繼續說.
原來當馬興奮的時候.
很緊張的時候.
很奔跑.
跑得很快的時候.
牠會.
不是流口水.
牠是會呼氣.
而呼氣.
牠甚至在天氣冷的時候.
會噴一些氣出來.
其實我心裡面想.

$^{241}$不要說馬.
人都會噴氣.
天氣冷的時候你也會噴氣.
你也會噴一些煙出來.
但是馬是更加厲害的.
不知道大家有沒有被馬噴過一口氣出來.
有些都很難受.
不是那麼好受都不一定.
你要知道.
馬噴氣不是為了呼吸.
而是一種力量正在爆發的聲音.
馬沒有人教牠勇敢的.
勇敢是牠天生的.
這個就是上帝給的.
弟兄姊妹.
有時候在我們風暴的人生的經歷裡面.
我們以為靠我們的勇氣.
是自己可以生產這些勇氣出來.
神問約伯.
是你可以令到馬變成這樣的.
原來是神.
這種力量這種勇氣是來自上帝.
這種特殊的能力都是來自上帝.
不知道大家有沒有聽過.
力克胡捷這個人.
他出生的時候.
大家都可能聽過他的故事.
他是沒有四肢的.
他爸爸曾經站在醫院的停車場.
對著神說.
我保護不了我兒子.
我做不了一個正常的爸爸.
在那一刻.
他開始承認自己那種無能力.
沒有馬力的狀況.
所以他開始每天就為他的兒子祈禱.
而不是再靠自己的意志死撐下去.
當他承認自己.
當他承認自己沒有辦法單靠他自己的體能.
支撐整個家庭.

$^{281}$他才學會披戴著神所給的平安.
神所給的力量給他去面對竟然的事情.
後來力克成為了全球的報導家.
不是因為他爸爸多麼堅硬.
而是他爸爸學會了.
讓上帝披上他生命.
以至將上帝的能力駕馭在他生命裡面.
當人承認自己沒有馬力的時候.
我們才能夠真正經歷到.
來自上帝的那種無限馬力.
無限馬力不是人努力回來的結果.
而是神給我們的恩典.
當我們停止再扮演上帝的時候.
我們才能夠成為真正的受造者.
其實神不在乎你的馬力.
神哪有在乎.
神在乎的是你是否願意謙卑.
讓神為你披上他的威嚴神的能力.
但是你願不願意呢.
在馬年我們真的需要神的能力駕馭我們.
而不是再靠我們自己的能力.
還有第二樣東西.
我們是擁留神的能力的戰馬.
我們是一匹戰馬來的.
我們留意留意.
第二十一至二十三節.
我們一起讀出來好不好.
請讀.
他用蹄在谷中挖地.
以能力歡躍.
他出去迎擊仇敵.
他嘲笑懼怕.
並不驚惶.
也不因刀劍退卻.
箭袋在他身上.
噌噌有聲.
槍和短槍.
閃閃發亮.
二十一節就講到馬.
用蹄在谷中挖地.

$^{321}$拋地或者挖地這個字.
是指挖掘或者是陷潭.
這是一個什麼畫面呢.
這匹馬不是因為自己焦慮.
不是因為發脾氣而淡地.
我們小朋友小時候一不開心就淡地.
這些是小朋友喜歡玩的一個.
宣洩他情感的遊戲.
或者是那個表達.
但是馬為什麼要淡地.
不是他發脾氣.
不是他心情不好.
馬不停地在跑地.
震動.
噴氣.
淡地.
是因為代表他精力旺盛.
他等不及要衝出去.
他很想主人放他出去.
讓他可以自由地奔跑.
來為主人做事.
而以能力來歡騰.
就代表了嗎哪種興奮.
那種力量的澎湃.
他很自豪去展示自己的力量.
他不是去很懼怕.
他不是想躲避.
而是他很想奮勇前去迎敵.
就算前面拿著什麼兵器都好.
他是願意走到前面奮勇一戰.
他主動衝去刀劍那裡.
他主動迎向危險.
向戰場邁進.
你看看是不是很有氣勢.
甚至在二十二節那裡.
他說他嘲笑懼怕.
他並不驚惶.
也不會因刀劍退卻.
簡單來說就是他根本不怕死.
他無懼死亡.

$^{361}$當然經文不是說馬有什麼道德的思維判斷.
也不是說他腦海裡有什麼潛在善惡的分別.
而是用詩歌的題材.
誇張地去演繹.
去描寫他們那種戰鬥的本能.
進一步將馬體的恐懼變成一個笑話.
對他們來說刀劍恐嚇簡直是一個笑話.
他視恐嚇如無物.
戰馬是不會被恐嚇來擊倒.
刀劍的劍影是不會令他退卻.
我們要知道古代的戰場.
時時都聽到兵器碰撞的聲音.
勇士吶喊的巨聲.
還有受傷戰士慘叫的聲浪.
血聲來充滿整個環境.
也有金屬閃光 刀劍閃來閃去的東西.
你換成一個普通的動物.
一定會害怕到逃走.
就像我養貓一樣.
他聽到一點聲音.
貓就會閃了.
馬上會逃走 馬上會躲起來.
但是戰馬是不會這樣的.
為什麼呢.
因為他天生就是一隻戰馬.
這裡除了他的品種之外.
也當然他經過訓練.
反而更加能夠令他更加興奮.
去出戰 往前衝.
然後加上第39章第23節.
神說戰袋在他身上有聲音.
槍和短槍 閃閃發亮.
這裡就很明顯讓我們感覺到.
原來這些兵器的出現.
戰馬面對這些兵器出現的時候.
他完全毫無懼色.
沒有躲避 甚至主動去迎擊.
就好像一隻獵犬一樣.
獵犬為什麼跑得快一點.
你試一下你家裡的狗和獵犬賽跑.

$^{401}$哪一隻會跑贏.
為什麼有什麼不同.
因為他是一隻獵犬.
你那隻只是普通狗.
這個就是最大的分別.
為什麼這隻馬這麼不同.
因為他不是一隻普通馬.
他是一隻戰馬.
大家都知道約伯當時正在爐灰裡面.
來刮他的爛瘡.
或者大概他在躲避人.
甚至他在躲避神也說不定.
但是上帝就讓他見到.
那個不躲避刀劍的戰馬.
也讓他見到神給約伯認定的價值.
他不是一匹爛泥.
而是可以被神磨練過的正兵.
不知道你有沒有想過.
為什麼上帝要記載戰馬給約伯看.
或者為什麼要將戰馬講給約伯聽.
在古代近東的時候.
戰馬就是當時最高科技的武器.
當然今天的最高科技武器就是導彈.
可能是槍械.
甚至無人機.
但是以前馬就是最強的武器.
人以為自己可以透過訓練和戰術.
可以完全主宰著這些力量.
但是上帝在這段經文裡.
告訴約伯.
利用馬問約伯.
這種令人畏懼的生命力.
是你給還是我給.
約伯最痛的.
其實大概不是失去一切.
沒有錢.
沒有房子.
沒有子女.
我想這些都是很痛.
但不是最重要的東西.

$^{441}$他最熟悉的邏輯思維的破裂.
就是他最難明白的地方.
善有善報.
惡有惡報.
這個是全世界人最正確的邏輯.
約伯也是這樣想.
不過這個邏輯.
在他生命的經歷裡.
徹底崩塌.
他找不到這個邏輯.
但是上帝沒有給他答案.
上帝只是給他見到一匹戰馬.
這匹的馬衝向戰場.
在那裡嘶叫.
在那裡震動.
在那裡噴氣.
但是這匹馬有沒有問對錯.
有沒有問輸贏.
他只是活出被造的力量.
人類在戰爭裡面.
他會計算榮辱.
他會計算勝負.
但是上帝在創造的狂野.
生命界的裡面.
其實我們在這些生命界的裡面.
我們的判斷其實真的很有限.
上帝不是否定了公義.
上帝是在提醒約伯.
甚至是我們.
你以為自己明白因果關係.
但是你連馬的震動都沒有辦法賜予.
你有什麼資格來去批判.
是啊.
上帝不是要和約伯辯論.
誰對誰錯.
而是把約伯帶到一個更高層次的視野.
承認上帝的創造的裡面.
有一部分是人.
根本永遠都沒有辦法可以控制.
也沒有辦法可以用簡單的邏輯思維.

$^{481}$道德判決.
因果的東西.
來去明白套用在裡面.
當我們應用的時候.
弟兄姊妹.
我們也可以去提醒我們自己.
我們是神所造的什麼馬.
你想做的是普通馬.
還是戰馬.
我神學院有個院長.
他已經離開了世界.
我讀神學的時候.
他經常和我們神學生說.
他說你們出去神學.
出去做事之後.
他說你不是一個普通的工人.
你不是一個兵.
你是正兵.
正兵.
當然他是和神學生說話.
但是和基督徒的生命裡面都是一樣.
你想做的是一個什麼基督徒.
是一個普通的馬.
還是戰馬.
這個問題留給我們自己去回應.
有一次在火災裡面.
有一名消防員衝入火場去救人.
記者就問他.
你不怕嗎.
他就說當然怕了.
但是有人在裡面.
所以我要去救他.
這句話就好像戰馬的跑地一樣.
你說他怕不怕.
其實也可能怕的打仗.
不過是命大過恐懼.
當箭袋噌噌聲的時候.
當刀劍閃爍的時候.
神的精兵不是沒有懼怕的.
是選擇迎向呼召.

$^{521}$這個比懼怕更重要.
戰馬不是自己決定戰場.
戰馬也不是為自己而戰.
戰馬只是回應主人的方向.
我們要記得.
消防員衝入火場的時候.
是因為他知道背後有整個團隊.
去支援他.
當戰馬衝入戰場的時候.
因為他知道主人握住他的韁繩.
兄弟姐妹.
上帝握住我們的生命.
我們又有什麼值得懼怕呢.
最後第三樣東西.
我們是展現出馬到功成的進馬.
第一我們有馬力的來源.
第二我們是戰馬.
第三原來我們都是進馬.
我們讀這兩節經文.
24到25節請讀.
激動將地吞下.
一聽覺聲就站不住.
每逢覺聲一響.
他說啊哈.
他從遠處聞到.
噌噌的氣息.
剛才我們說了.
我們知道能力的來源.
也都清楚明白到能力的本質.
最後我們要昇華到生命的格調和結果.
我們是馬到功成的進馬.
24節提到震震這個字.
這是一個內在極為興奮的表達.
就好像肌肉一樣.
他不自覺地在震動.
大家有沒有試過.
通常有肌肉的才能做到.
我們瘦猛猛的做不到.
再加上激動這個字.
外在憤怒的情緒演繹.

$^{561}$帶出戰馬意志的高昂.
而將地吞下.
這個可以有一個誇張的表達.
意思就是他無論是在地上.
在翻起塵土.
就搞到這些塵土好像吞噬了大地一樣.
就好像他的威勢吞噬了大地一樣.
不過由於約伯在這裡.
這是神窮的約伯說的話.
而且是透過詩歌的題材表達出來.
所以極可能是一種文學的修辭都未定.
這是一個誇張的演繹.
他強調了馬的神勇.
馬的無懼怕.
他可以全程投入.
毫無懼惜地投奔在戰場.
這種不退縮.
這種氣勢到底是誰給的.
是怎麼來的.
而後半句更加說.
一聽覺聲就站不住.
這個意思真的非常重要.
因為他一等不及.
他一聽到覺聲打仗的時候.
他已經等不及他不要站在這裡.
他要出去.
所以在二十五節.
他就說每逢覺聲一響.
他說啊.
這裡真是一個很特別的字眼.
原來這個希伯來文字眼.
這個儀聲詞.
其實是模擬嗎哪種興奮的狀況.
用來表示萬一聽到.
這個發號司令的聲音.
號角的聲音.
他就立即很緊張.
很激動的樣子.
他很想作出一個回應.
而這個回應.

$^{601}$他就用了「啊哈」這個字眼.
這個讓我們學到幾樣.
很重要的提醒.
特別在馬年那裡.
第一樣東西就是.
我們有沒有敏銳的.
屬靈的觸覺.
駿馬和普通馬的分別.
就是在於他的敏感度.
他在遠方就能夠察覺到.
對手的動靜.
這個就是不同的原因.
為什麼今時今日.
大家都是人.
有些會厲害一點.
有些是有智慧一點.
有些有創業的能力.
為什麼.
因為可能他有.
不同的觸覺的能力.
馬 大家都是馬.
放十匹馬在裡面.
為什麼那些馬特別跑得快一點.
為什麼那些馬是不同一點.
因為他不是一隻普通馬.
他是一隻駿馬.
我們馬年.
不單止有一個體力的.
健康的身體更加重要.
這些都是重要的.
我們更加重要的就是.
要有屬靈的眼光.
在危機的裡面.
你能夠看到曙光.
在挑戰的裡面.
你看到上帝的手在工作.
這個就是人家看不到的東西.
你能夠看到.
這個就是屬靈的敏銳力.
今時我們怎麼能夠.

$^{641}$在我們的說話上.
在你處理的事情上.
在你的眼界上.
有屬靈的洞見.
有屬靈的慧眼.
有屬靈的觸覺.
這個就是分別.
你是駿馬還是一匹普通馬.
還有第二件事.
這段經文提醒我們.
我們有個毫不猶豫的行動力.
經文形容馬一聽見覺醒.
就沒辦法安靜坐在那裡.
馬到功成的秘訣.
是什麼大家知不知道.
功成就一定要馬到.
所以馬到才是一個最重要的因素.
剛才第一點的時候告訴我們.
神的能力已經給了我們.
第二點告訴我們.
我們的身份又確認了.
現在我們是回應呼召.
立刻行動的時候.
真正的成功不是靠運氣.
而是靠對神旨意的那種.
急不及待的熱情.
我記得我年輕的時候.
我回到教會的時候.
傳道目的者剛剛離任.
教會是剛剛一片震撼.
沒多久又來了一個傳道人.
當時我很小.
這個傳道人做了沒多久又離職.
當時我已經開始出來工作.
在那一刻.
大家不知道怎麼面對.
但是我們願意.
那些熱心的弟兄姊妹.
我們就回應上帝.
告訴上帝.

$^{681}$有什麼可以用我們.
為什麼我開始學結他.
就在那段時間.
我在那段時間開始學琴.
用到就用.
少年團契.
我去做導師.
我去彈結他.
我去彈琴.
那段時間就在那裡出現.
是急不及待想去幫忙.
想去做事.
就好像這段經文所說.
不能等了.
現在你知道有事了.
你要怎麼做.
坐在那裡就中.
應該要起來.
來一個急不及待的熱情.
還有第三樣東西.
跨越恐懼的海船心態.
你有沒有留意那裡用了「啊哈」這個字.
「吃啊」這個字.
希伯來文「吃啊」.
這個字在舊約裡面出現了十三次.
「吃啊」這個字.
是一種很興奮的喜樂的感嘆的表達.
是戰馬很興奮的歡喜的聲音.
當然很興奮是為主人而出戰.
但是更重要這個字眼.
是一種得勝的喜樂的字眼.
它更加含有就是它發出一個得勝的宣告.
基督徒馬到攻城.
不是指萬事如意.
而是指無論在什麼環境裡面.
我們都能夠靠主去跨城.
這個就是進馬的氣質.
環境越艱難.
鬥志就越高昂.
在疫情最嚴重的時候.

$^{721}$有一位基督徒的護士.
他連續上班花了十幾天.
醫院的走廊滿佈病人.
聽到的就是咳嗽聲.
然後都是咳嗽聲.
報警鐘聲不斷響起.
見到醫護人員走來走去.
他知道自己總有一天會受感染.
他家裡面有家人.
有小朋友有父母.
有人問他你不怕嗎.
其實他沉默了一下.
然後就說怕了.
但是我聽到號角的聲音.
他說他不是超人.
他也會累.
他也會哭.
但是當醫院的電話響起的時候.
當病房缺人的時候.
他就站不住.
就好像一匹馬一樣.
站不住.
這個不是衝動.
這個是回應.
弟兄姊妹.
約伯記三十九章那一匹的馬.
是一匹賤馬.
但是與此同時.
他也有一匹進馬.
他不能夠站著.
他要動.
他要說 啊哈.
那個不是一個狂妄的聲音.
那個是生命找到方向的興奮的聲音.
我們知道馬力是來自神.
我們明白我們不是爛泥.
是神使用我們的賤馬.
我們是賤馬.
我們更加明白.
角聲響起的時候.

$^{761}$我們不站著.
我們是一匹進馬.
與其他馬是不同的.
不是因為我們不怕.
不是因為我們很強.
是因為我們知道.
誰在掌管.
誰在前面.
真正馬道功成.
不是我們征服了多少領土.
不是我們贏了多少場硬仗.
不是我們教會多了多少人.
奉獻多了多少錢.
而是.
當神的號角響起的時候.
你和我有沒有退場.
弟兄姊妹.
當神響起一個號角.
在呼召人的時候.
你在哪裡.
有沒有退場.
馬年.
真是一個很重要的年份.
我們想求的是一條沒有戰場的道路.
還是求神賜給我們一顆賤馬的精神.
不要只顧著馬怨地上的泥土.
學一下這一匹進馬.
開始刨地.
神要我們學習謙卑.
甚至在最軟弱的時間.
為我們賜上無限的馬力.
我們準備好.
我們回應上帝.
這一句得勝的宣告了沒有.
我們一起禱告.
我們一起祈禱.
馬年的時間.
我們從聖經裡面.
我們要學習馬的特質.
主我們知道.

$^{801}$上帝你和約伯分享這段馬的特質的時候.
他想提醒我們或者提醒約伯.
我們連一匹嗎哪種威嚴的能力.
那種來源.
甚至是那種奇妙的東西.
我們都不懂得插頭.
都不了解.
在這個宇宙這麼浩瀚.
這麼多不明的東西.
我們實在沒有辦法去掌握很多東西.
但是主啊.
透過今天的經文.
幫助我們學會謙虛我們自己.
以致到我們來到你面前的時候.
我們才能夠有這個無限的馬力.
在我們身上擁留著的.
我們祈求你賜給我們有戰馬的那一種精神.
幫助我們學習.
做一個與眾不同.
非凡的生命的基督徒.
成為一匹進馬.
主啊幫助我們在馬年的裡面.
有這樣的立志.
多謝您收聽我們禱告的奉耶穌基督的恩主的命求.
阿們.
\newpage



\section{申命記 2:1-3:29}
\label{sec:031674}
\textbf{過去和現在的對話 (2) 聽歷史,講歷史,寫歷史}~~~~何淑斌傳道
\newline
\newline
連結: \href{http://www.church.com.hk/acms/content.asp?site=cdc\&op=show\&type=product\&code=031674\&layout=sermon}{\texttt{031674}} ~~~~ 語音日期: 2026-02-14
\newline
\newline
\hyperref[sec:031673]{\small{< < < PREV SERMON < < <}}
~
\hyperref[sec:index_chronic]{\small{[返順時目]}}
~
\hyperref[sec:031675]{\small{> > > NEXT SERMON > > >}}
\newline
\newline
申命記 2:1-3:29
\newline
\begin{longtable}{cl}
\hline
\hline
章節 & 經文 (和合本修訂版)\\
\hline
2:1 & \begin{tabularx}{0.7\textwidth}{X} 「我們轉回，從紅海的路往曠野去，正如耶和華所吩咐我的。我們在西珥山繞行了許多日子。 \end{tabularx} \\ \\ \relax
2:2 & \begin{tabularx}{0.7\textwidth}{X} 耶和華對我說： \end{tabularx} \\ \\ \relax
2:3 & \begin{tabularx}{0.7\textwidth}{X} 『你們繞行這山已經夠久了，要轉向北方。 \end{tabularx} \\ \\ \relax
2:4 & \begin{tabularx}{0.7\textwidth}{X} 你要吩咐百姓說：你們弟兄以掃的子孫住在西珥，你們要經過他們的邊界。他們必懼怕你們，但你們要分外謹慎。 \end{tabularx} \\ \\ \relax
2:5 & \begin{tabularx}{0.7\textwidth}{X} 不可向他們挑戰；他們的地，連腳掌可踏之處，我都不給你們，因我已將西珥山賜給以掃為業。 \end{tabularx} \\ \\ \relax
2:6 & \begin{tabularx}{0.7\textwidth}{X} 你們要用錢向他們買糧吃，也要用錢向他們買水喝。 \end{tabularx} \\ \\ \relax
2:7 & \begin{tabularx}{0.7\textwidth}{X} 因為耶和華---你的神在你手裡所做的一切事上已賜福給你。你走這大曠野，他都知道。這四十年，耶和華---你的神與你同在，因此你一無所缺。』 \end{tabularx} \\ \\ \relax
2:8 & \begin{tabularx}{0.7\textwidth}{X} 「於是，我們經過我們弟兄以掃子孫所住的西珥，從亞拉巴的路，經過以拉他、以旬‧迦別，轉向摩押曠野的路去。 \end{tabularx} \\ \\ \relax
2:9 & \begin{tabularx}{0.7\textwidth}{X} 耶和華對我說：『不可侵犯摩押，也不可向他們挑戰。他們的地，我不賜給你為業，因我已將亞珥賜給羅得的子孫為業。』 \end{tabularx} \\ \\ \relax
2:10 & \begin{tabularx}{0.7\textwidth}{X} 先前，以米人住在那裡，百姓又大又多，像亞衲人一樣高大。 \end{tabularx} \\ \\ \relax
2:11 & \begin{tabularx}{0.7\textwidth}{X} 他們跟亞衲人一樣，也算是利乏音人，但摩押人卻稱他們為以米人。 \end{tabularx} \\ \\ \relax
2:12 & \begin{tabularx}{0.7\textwidth}{X} 從前，何利人也住在西珥，但以掃的子孫把他們除滅，佔領了他們的地，接續他們在那裡居住，如同以色列在耶和華賜給他們為業之地所做的一樣。 \end{tabularx} \\ \\ \relax
2:13 & \begin{tabularx}{0.7\textwidth}{X} 『現在，起來，過撒烈溪！』於是我們過了撒烈溪。 \end{tabularx} \\ \\ \relax
2:14 & \begin{tabularx}{0.7\textwidth}{X} 從離開加低斯‧巴尼亞到渡過撒烈溪，這段時期共三十八年，直到這一代的戰士都從營中滅盡，正如耶和華向他們所起的誓。 \end{tabularx} \\ \\ \relax
2:15 & \begin{tabularx}{0.7\textwidth}{X} 耶和華的手也攻擊他們，將他們從營中除滅，直到滅盡。 \end{tabularx} \\ \\ \relax
2:16 & \begin{tabularx}{0.7\textwidth}{X} 「百姓中所有的戰士滅盡死亡以後， \end{tabularx} \\ \\ \relax
2:17 & \begin{tabularx}{0.7\textwidth}{X} 耶和華吩咐我說： \end{tabularx} \\ \\ \relax
2:18 & \begin{tabularx}{0.7\textwidth}{X} 『你今日要經過摩押的邊界亞珥， \end{tabularx} \\ \\ \relax
2:19 & \begin{tabularx}{0.7\textwidth}{X} 走到亞捫人之地。不可侵犯他們，也不可向他們挑戰。亞捫人的地，我不賜給你們為業，因我已將那地賜給羅得的子孫為業。』 \end{tabularx} \\ \\ \relax
2:20 & \begin{tabularx}{0.7\textwidth}{X} 那地也算是利乏音人之地，因為先前利乏音人住在那裡，亞捫人稱他們為散送冥人。 \end{tabularx} \\ \\ \relax
2:21 & \begin{tabularx}{0.7\textwidth}{X} 那裡的百姓又大又多，像亞衲人一樣高大，但耶和華從亞捫人面前除滅他們，亞捫人就佔領他們的地，接續他們在那裡居住。 \end{tabularx} \\ \\ \relax
2:22 & \begin{tabularx}{0.7\textwidth}{X} 這正如耶和華從前為住在西珥的以掃子孫，將何利人從他們面前除滅，使他們得了何利人的地，接續他們在那裡居住，直到今日一樣。 \end{tabularx} \\ \\ \relax
2:23 & \begin{tabularx}{0.7\textwidth}{X} 亞衛人先前住在鄉村直到迦薩；從迦斐託出來的迦斐託人將亞衛人除滅，接續他們在那裡居住。 \end{tabularx} \\ \\ \relax
2:24 & \begin{tabularx}{0.7\textwidth}{X} 你們起來往前去，過亞嫩谷。看，我已將亞摩利人希實本王西宏和他的地交在你手中，你要開始去得他的地為業，向他挑戰。 \end{tabularx} \\ \\ \relax
2:25 & \begin{tabularx}{0.7\textwidth}{X} 從今日起，我要讓天下萬民驚慌懼怕你，聽見你的名聲，就因你發顫傷慟。」 \end{tabularx} \\ \\ \relax
2:26 & \begin{tabularx}{0.7\textwidth}{X} 「我從基底莫的曠野派遣使者到希實本王西宏那裡，用和平的話說： \end{tabularx} \\ \\ \relax
2:27 & \begin{tabularx}{0.7\textwidth}{X} 『求你讓我穿越你的地，我走路的時候，只走大路，不偏左右。 \end{tabularx} \\ \\ \relax
2:28 & \begin{tabularx}{0.7\textwidth}{X} 你可以賣糧給我吃，賣水給我喝；只要讓我步行過去， \end{tabularx} \\ \\ \relax
2:29 & \begin{tabularx}{0.7\textwidth}{X} 就如住在西珥的以掃子孫和住在亞珥的摩押人待我一樣，等我過了約旦河，進入耶和華－我們神所賜給我們的地。』 \end{tabularx} \\ \\ \relax
2:30 & \begin{tabularx}{0.7\textwidth}{X} 但希實本王西宏不肯讓我們從他那裡經過，因為耶和華－你的神使他性情頑梗，內心剛硬，為要把他交在你手中，像今日一樣。 \end{tabularx} \\ \\ \relax
2:31 & \begin{tabularx}{0.7\textwidth}{X} 耶和華對我說：『看，我已開始把西宏和他的地交給你了，你要開始得他的地為業。』 \end{tabularx} \\ \\ \relax
2:32 & \begin{tabularx}{0.7\textwidth}{X} 「西宏和他的眾百姓出來迎擊我們，在雅雜與我們交戰。 \end{tabularx} \\ \\ \relax
2:33 & \begin{tabularx}{0.7\textwidth}{X} 耶和華－我們的神把他交給我們，我們就殺了他和他的眾兒子，以及他所有的百姓。 \end{tabularx} \\ \\ \relax
2:34 & \begin{tabularx}{0.7\textwidth}{X} 那時，我們奪了他一切的城鎮，毀滅各城的男人、女人、孩子，沒有留下一個倖存者。 \end{tabularx} \\ \\ \relax
2:35 & \begin{tabularx}{0.7\textwidth}{X} 只有牲畜和所奪各城的財物，我們都取為自己的掠物。 \end{tabularx} \\ \\ \relax
2:36 & \begin{tabularx}{0.7\textwidth}{X} 從亞嫩谷旁的亞羅珥和谷中的城，直到基列，沒有一座城是高得我們不能攻取的；耶和華－我們的神把它們全都交給我們了。 \end{tabularx} \\ \\ \relax
2:37 & \begin{tabularx}{0.7\textwidth}{X} 只有亞捫人之地，雅博河沿岸，以及山區的城鎮，你沒有挨近，這全是耶和華－我們神所吩咐的。」 \end{tabularx} \\ \\ \relax
3:1 & \begin{tabularx}{0.7\textwidth}{X} 「我們又轉回，朝巴珊的路上去。巴珊王噩和他的眾百姓出來迎擊我們，在以得來與我們交戰。 \end{tabularx} \\ \\ \relax
3:2 & \begin{tabularx}{0.7\textwidth}{X} 耶和華對我說：『不要怕他！因我已把他和他的眾百姓，以及他的地，都交在你手中；你要待他像從前待住在希實本的亞摩利王西宏一樣。』 \end{tabularx} \\ \\ \relax
3:3 & \begin{tabularx}{0.7\textwidth}{X} 於是耶和華－我們的神也把巴珊王噩和他的眾百姓都交在我們手中；我們殺了他，沒有給他留下一個倖存者。 \end{tabularx} \\ \\ \relax
3:4 & \begin{tabularx}{0.7\textwidth}{X} 那時，我們奪了他一切的城鎮，共六十座，沒有一座城不被我們所奪，這是亞珥歌伯的全境，巴珊王噩的國度。 \end{tabularx} \\ \\ \relax
3:5 & \begin{tabularx}{0.7\textwidth}{X} 這些堅固的城都有高的城牆，有門有閂，此外，還有許多無城牆的鄉村。 \end{tabularx} \\ \\ \relax
3:6 & \begin{tabularx}{0.7\textwidth}{X} 我們把這些都毀滅了，像從前待希實本王西宏一樣，毀滅各城的男人、女人、孩子； \end{tabularx} \\ \\ \relax
3:7 & \begin{tabularx}{0.7\textwidth}{X} 只有一切牲畜和城中的財物，我們取為自己的掠物。 \end{tabularx} \\ \\ \relax
3:8 & \begin{tabularx}{0.7\textwidth}{X} 那時，我們從兩個亞摩利王的手裡把約旦河東邊的地奪過來，從亞嫩谷直到黑門山， \end{tabularx} \\ \\ \relax
3:9 & \begin{tabularx}{0.7\textwidth}{X} 這黑門山，西頓人稱為西連，亞摩利人稱為示尼珥。 \end{tabularx} \\ \\ \relax
3:10 & \begin{tabularx}{0.7\textwidth}{X} 我們奪了平原的各城、基列全地、巴珊全地，直到撒迦和以得來，都是巴珊王噩國內的城鎮。 \end{tabularx} \\ \\ \relax
3:11 & \begin{tabularx}{0.7\textwidth}{X} 利乏音人所剩下的只有巴珊王噩。看哪，他的床是鐵床，按照人肘的度量，長九肘，寬四肘，現今不是在亞捫人的拉巴嗎？」 \end{tabularx} \\ \\ \relax
3:12 & \begin{tabularx}{0.7\textwidth}{X} 「那時，我們得了這地。從亞嫩谷旁的亞羅珥起，連同基列山區的一半和境內的城鎮，我都給了呂便人和迦得人。 \end{tabularx} \\ \\ \relax
3:13 & \begin{tabularx}{0.7\textwidth}{X} 基列其餘的地和巴珊全地，就是噩的國度，我給了瑪拿西半支派。亞珥歌伯全境就是巴珊全地，也稱為利乏音人之地。 \end{tabularx} \\ \\ \relax
3:14 & \begin{tabularx}{0.7\textwidth}{X} 瑪拿西的子孫睚珥佔領了亞珥歌伯全境，直到基述人和瑪迦人的邊界，就按自己的名字稱這些地，就是巴珊，為哈倭特‧睚珥，直到今日。 \end{tabularx} \\ \\ \relax
3:15 & \begin{tabularx}{0.7\textwidth}{X} 我又將基列給了瑪吉。 \end{tabularx} \\ \\ \relax
3:16 & \begin{tabularx}{0.7\textwidth}{X} 我給了呂便人和迦得人從基列到亞嫩谷，以谷的中央為界，直到亞捫人邊界的雅博河； \end{tabularx} \\ \\ \relax
3:17 & \begin{tabularx}{0.7\textwidth}{X} 還有亞拉巴和靠近約旦河之地，從基尼烈直到亞拉巴海，就是鹽海，以及毗斯迦山斜坡的山腳東邊之地。 \end{tabularx} \\ \\ \relax
3:18 & \begin{tabularx}{0.7\textwidth}{X} 「那時，我吩咐你們說：『耶和華－你們的神已將這地賜給你們為業；你們所有的勇士都要帶著兵器，在你們的弟兄以色列人前面過去。 \end{tabularx} \\ \\ \relax
3:19 & \begin{tabularx}{0.7\textwidth}{X} 但你們的妻子、孩子、牲畜，可以住在我所賜給你們的各城裡，我知道你們有許多牲畜。 \end{tabularx} \\ \\ \relax
3:20 & \begin{tabularx}{0.7\textwidth}{X} 等到耶和華讓你們的弟兄像你們一樣，得享太平，他們在約旦河另一邊，也得了耶和華－你們的神所賜給他們的地，你們各人才可以回到我所賜給你們為業之地。』 \end{tabularx} \\ \\ \relax
3:21 & \begin{tabularx}{0.7\textwidth}{X} 那時，我吩咐約書亞說：『你親眼看見了耶和華－你們的神向這兩個王一切所做的事，耶和華也必向你所要去的各國照樣做。 \end{tabularx} \\ \\ \relax
3:22 & \begin{tabularx}{0.7\textwidth}{X} 不要怕他們，因為那為你們爭戰的是耶和華－你們的神。』」 \end{tabularx} \\ \\ \relax
3:23 & \begin{tabularx}{0.7\textwidth}{X} 「那時，我懇求耶和華說： \end{tabularx} \\ \\ \relax
3:24 & \begin{tabularx}{0.7\textwidth}{X} 『主耶和華啊，你已開始將你的偉大和你大能的手顯給你僕人看。在天上，在地下，有甚麼神明能像你行事，像你有大能的作為呢？ \end{tabularx} \\ \\ \relax
3:25 & \begin{tabularx}{0.7\textwidth}{X} 求你讓我過去，看約旦河另一邊的美地，就是那佳美的山區和黎巴嫩。』 \end{tabularx} \\ \\ \relax
3:26 & \begin{tabularx}{0.7\textwidth}{X} 但耶和華因你們的緣故向我發怒，不應允我。耶和華對我說：『你夠了吧！不要再向我提這事。 \end{tabularx} \\ \\ \relax
3:27 & \begin{tabularx}{0.7\textwidth}{X} 你上毗斯迦山頂去，向東、西、南、北舉目，用你的眼睛觀看，因為你必不能過這約旦河。 \end{tabularx} \\ \\ \relax
3:28 & \begin{tabularx}{0.7\textwidth}{X} 你卻要吩咐約書亞，勉勵他，使他壯膽，因為他必在這百姓前面過去，使他們承受你所要觀看之地。』 \end{tabularx} \\ \\ \relax
3:29 & \begin{tabularx}{0.7\textwidth}{X} 於是我們停留在伯‧毗珥對面的谷中。」 \end{tabularx} \\ \\
[1ex]
\hline
\hline
\end{longtable}
$^{1}$今晚我們繼續.
哇 這個系列很長.
過去與現在的對話系列.
即是以前與現在說話.
我今天講題是.
聽歷史 講歷史 寫歷史.
不如這樣吧.
上次有Mentee Question.
這次我沒有放Mentee.
不過也有幾個問題跟大家玩玩.
好不好啊 開始了.
哇 誠經知識問題.
直接回答就可以了.
這條是溫暖的 比較簡單.
以色列人是誰的後代.
阿哥B對不對.
OK啊 好 大家很冷靜.
好 接著一條就深一點了.
好.
以東人是誰的後代.
A 是以素.
B 是以實瑪麗.
A 還是B 請A的舉手.
以素的舉手 舉高一點.
好像是這樣.
先不要搜索 以東人.
因為與今天的經文很重要.
是有相關的 一會兒你會聽到.
以東人是誰的後代.
A 以素的舉手.
有兩隻手 三隻.
阿南是按摩的吧.
OK 兩隻.
是以實瑪麗的舉手.
不舉的那些是什麼.
多一點啊 以實瑪麗.
很多啊 以實瑪麗.
答案是... 我先不說了.
一會兒會聽到的.
好 第二題.

$^{41}$第三題 摩壓人是誰的後裔.
好深啊 A 是羅德.
B 是拿克.
聽過沒有 這個名字.
很深啊 這次又弱了.
這次講生命計都很深.
阿南出的 別問我.
(笑).
A 舉手是羅德的後裔.
摩壓人.
一隻手.
兩隻手.
另外.
B 拿克 舉手.
R 痕 這個是R 痕.
這個不是手來的.
有點迷茫.
拿克是哪裡來的.
是舊約的.
創世記有的.
或者問一下大家.
聽過摩壓人沒有.
聖經有一個很出名的摩壓人.
是聖經一個書卷的名字.
猜不猜得到.
哦 沒有啊.
你早前講完.
阿南說他早前講完.
有個摩壓人.
很出名的 聖經一個書卷的名字.
是摩壓的.
是一個女子.
給了很多提示.
叫什麼名字.
摩壓的女子.
真的沒有想法.
路德是對的.
路德是很重要的.
耶穌的.
前人.

$^{81}$沒有路德就沒有耶穌.
好.
先不講特望.
一會兒你會聽到.
阿滿人是誰的後代.
B是羅德.
又回來了.
含.
A是含的.
哇 好多A哦.
含是誰聽過.
羅亞的兒子.
哪個兒子.
羅亞其中一個兒子.
B是羅德的.
沒有了.
阿滿人是誰的後代.
因為你看歷史書.
會經常出現這三個.
耳東人.
摩壓人 阿滿人.
我們先hold住.
今天講的主題是.
聽歷史.
講歷史.
寫歷史.
其實是一個進程.
我想我們每個人都會聽歷史.
但不知道大家有沒有講過歷史.
有沒有講過.
或者有人很喜歡歷史堂.
這次真的像歷史堂.
有沒有人喜歡歷史的 舉手.
我歷史很差.
我中學歷史很差.
現在開始喜歡.
我也是.
我中學歷史和中史都非常差.
我歷史和中史都不好.
但現在開始喜歡.

$^{121}$人長大了就開始喜歡.
這是一個過程.
由聽歷史 講歷史.
到寫歷史.
不如開始之前我們做一個祈禱.
親愛的主愛.
你是歷史的主宰.
當今天女兒要分享.
以色列人的歷史.
是你在他們當中的歷史的時候.
求神你開我們的耳朵.
你和我們說話.
在歷史裡面我們可以學到什麼.
在以色列人的歷史裡面.
我們會有什麼功課.
所以主人求你的靈充滿.
像剛纔的祈禱一樣.
充滿我們在這裡.
充滿我們每一個.
讓我們聽的 講的.
都能夠合乎你的心意.
按你的心意去宣講.
去聆聽 祈禱奉耶穌名求.
Amen.
好 很快.
歷史堂.
我們用最簡單的查經 識經方法.
五個W.
我們回顧一下申命記.
因為上個星期應該是.
申命記的第一章.
應該是主要講信心.
我不知道他有沒有講過一點點.
不如幾個用What.
Who, When, Where, Why.
我很快講一次.
申命記或Recap.
What 申命記.
我估計Amen也有講.
原文是第二部律法.

$^{161}$不是說有另一部律法.
是同一部律法.
只不過是再重新.
西乃山的律法.
摩西回顧過去.
和摩西臨終的遺言.
好像是一整篇遺書.
Who.
摩西.
是摩西寫的.
一直都是說摩西寫的.
摩西五經.
摩西說給新一代的.
以色列人.
When.
我不是說公元前或公元后.
是說當以色列人.
出了埃及.
在曠野四十年.
走了四十年.
又未開始進入.
應許之地.
中間這個位置就有申命記.
摩西和新一代的以色列人.
來一個回顧.
即是他們準備進入迦南地.
但未進入.
但去了曠野兜了四十年.
就是上一代那些.
都已經走了.
去到現在這一代.
在哪裡呢?在迦南以東的摩亞平原.
Why.
為什麼摩西要再.
第二次說這個律法解釋.
是要和新一代的.
以色列人發出.
勸告和警戒.
要勸他們要與神同行.
好了.

$^{201}$我們聽歷史.
不如我們一起讀一下經文.
大家不要睡著.
我們一起讀一下經文.
是二章的第一節.
我今天是說二章和三章.
我們轉回.
從紅海的路.
往曠野去.
是照耶和華所吩咐我的.
我們在西義山.
繞了許多日子.
他們兜了很多圈.
但不知道大家.
看經文的時候.
我一直都有一個問題.
明明他們又說跟著耶和華.
所吩咐做.
為什麼又要在這裡兜圈.
難道耶和華叫他們兜圈.
會不會這樣.
我們.
看聖經要看上下文.
我們回到.
第一章的結尾.
是一章的.
四十六節.
看不看得很像.
不如我讀一章那邊.
至於你們.
要轉回從紅海的路.
往曠野去.
和剛纔二章第一節.
差不多一樣.
但有不同是.
中間夾著紅色.
最後我都一樣.
於是你們在迦底斯.
住了許多日子.
但為什麼他們要兜圈.

$^{241}$原來中間有紅色.
第一章的結尾有說.
他們違背了耶和華.
第一章的結尾.
除了探子.
去看迦南地.
其實是說.
耶和華叫他們不要上去打仗.
他們就說要上去打仗.
所以神就說.
我離開你們.
他們就打敗仗.
這是一個很慘痛的經歷.
原來一章的.
四十一至四十五節.
就說了他們如何不聽神話.
原來是有結果的.
他們再兜圈.
但我這次.
第二章就比較好消息.
第三節我們一起讀.
耶和華對我說.
一起讀.
耶和華對我說.
你們要行.
這生的日子夠了.
要轉向.
不去.
終於等夠了.
兜了很多圈.
但停了很久.
終於神說.
Let's go 行得了.
這是整個.
以色列文這四十年的.
一個重要的轉捩點.
是神親自開聲.
跟他們說.
是跟摩西說.
你們要行.

$^{281}$這生的日子夠了.
要轉向不去.
起行了.
有時我們人生.
都會兜很多圈.
有時會不聽神話.
走錯路.
但神的憐憫大於一切.
神都會憐憫我們.
就說繼續向前走.
我們又說.
繼續聽歷史.
繼續向前走.
遇到什麼呢.
以色列人不是一開始就打仗.
他們遇到三個族群.
剛纔問過大家.
哪三個族群.
以東人.
摩爾人.
阿滿人.
在第二章裡.
神用篇幅.
第二章至第三章裡.
說這三個族群.
他們一開始起行.
神先說這三個族群.
原來是誰.
以東是誰的子孫.
以掃.
不是以實瑪利.
是以掃.
以掃是誰.
亞國的兄弟.
亞國的哥哥.
但他因為一碗紅豆湯.
換了長子名分.
就沒了.
原來他的後裔.
神也看管.

$^{321}$神跟他吩咐.
你要經過他們的境界.
但你不可以.
跟他們爭戰.
你不要跟他們打架.
不要跟他們爭戰.
你過去就可以了.
因為他跟你有血緣.
他是以色列的.
哥哥的後裔.
所以你們神和神.
沒有忘記.
神除了給地以色列人.
神也答應了.
以東人將西蟻.
交給他們.
所以西蟻是以東的.
是西蟻山.
所以神對自己.
納過的英雄不會忘記.
也不會忘記以東人.
另外摩厄人.
是誰的子孫.
看到嗎.
羅德.
是誰.
阿巴拉罕的誰.
親戚.
羅德.
原來是羅德的子孫.
聖經也是這樣說.
神說你不可以.
妖害摩厄人.
也不可以跟他們爭戰.
因為也有血緣關係.
神也答應了.
給他們阿爾.
阿爾是摩厄地的首都.
給他們.
神原來也留了地.

$^{361}$給他們.
所以叫你不要打擾他們.
第三 阿滿人.
又是誰的子孫.
又是羅德.
你會不會覺得.
又是羅德.
原來剛纔第二題第三題的答案.
也是羅德.
因為羅德在《創世記》裡.
摩厄人是跟大女生生的.
阿滿人是跟小女生生的.
這些亂亂的東西.
我就不說了.
原來聖經也是這樣說.
神說不可以.
妖害阿滿人.
也不可以跟他們爭戰.
因為也有血緣.
神也答應了.
把地給阿滿人.
原來神在這裡.
一開始不是說以色列人.
不是說什麼.
不是說迦南地.
是說以東人.
摩厄人和阿滿人.
再下去第二章.
是說什麼呢.
以色列人.
面對兩場戰爭.
其實經文很簡單.
在第二章和第三章.
第一章是說以色列人靠著神.
這次就有聽神話了.
靠著神.
神說你們起來.
就經過阿聯谷.
我已將阿摩釐人.
希什本王西雲.

$^{401}$和他的地交在你手中.
將他的地為業.
這是第二章.
另一個王.
耶和華又對我說.
不要怕他.
因我已將他和他的眾民.
並他的地都交在你手中.
你要代他將從前.
駐希什本的阿摩釐王.
西雲一樣.
很快原來以色列人.
起行之後.
渡過了以東摩厄阿滿人.
來了兩場戰爭.
兩場都贏了.
就擊敗了.
我們進入故事.
裡面.
申命記是摩西對著.
新一代的以色列人說.
去說剛纔的歷史.
如果你是.
新一代的以色列人.
你在摩西面前.
聽到這些.
你會覺得怎麼樣.
你會覺得.
不關我的事.
阿滿人我也不知道在哪.
我又沒見過.
那個什麼王.
那個什麼西雲.
阿摩釐王.
又不知道是希什本還是阿摩釐.
那個是誰 我沒見過.
那個什麼巴山王.
我又沒見過 關我什麼事.
是不是會這樣.
貓山王.

$^{441}$我貓山王會知道多些.
巴山王就不太知道.
貓山王就想知道.
但是你在新一代.
以色列人摩西跟你說這些的時候.
你會覺得.
都很重要的 我要聽.
還是怎麼樣.
大家說會很留心聽的 舉手.
如果你是新一代的以色列人.
你會很留心聽摩西說這些.
你舉手.
沒有.
大家都會覺得.
我站在這裡.
聽完什麼時候放學.
什麼時候放學.
什麼時候可以走.
又說第二次.
什麼重新誡命.
應該一講就講十界.
十界我都聽過.
又講這些都還沒到.
未到戲肉 為什麼不講十界.
為什麼用幾張 兩張.
的經文來講這些.
我們可能會這樣.
真的 我都.
為什麼要講這些.
直接講十界就行了.
是吧 我想在聖經裡面.
每一句說話 每一句經文.
記載都有意思.
如果我看回整個聖經.
我這裡不再引述.
我看回舊約 如果你搜索一下.
以東人 你搜索一下.
阿滿人.
這裡你搜索一下摩雅人.
神是講明.

$^{481}$叫以色列的後代不要跟他們.
打仗.
但是我可以告訴你 你搜索一下.
在士師記.
甚至在列王記.
你都會見到以色列人.
跟他們打仗.
我們有沒有聽到神的說話.
如果真的殺死所有摩雅人.
就沒有路德了 沒有路德 沒有耶穌.
神的計劃會不會.
這樣打彎了呢.
所以其實.
是很重要 我想摩西.
神要摩西講這些說話給.
新一代以色列民聽.
這是很重要.
因為上一代以色列民就是叛逆神.
所以他們不能進入迦南地.
但是神的心意是很想.
興起新一代的以色列人.
可以進入迦南地.
和可以真的.
秉承神給他們的.
誡命 去一代傳一代.
傳一代.
將他的律法傳下去.
今天.
我想起兩個星期前.
不是有聯合崇拜的.
有同福歷史.
大家用什麼態度去聽.
不用回答我.
我們會不會覺得.
教會的歷史都不關我的事.
我之前.
讀神學都要讀.
基督教內華史.
最初我會覺得.
關我什麼事.

$^{521}$今天關我什麼事.
但當我讀了之後.
我再去澳門 再去廣州.
再見教會.
我明白了.
不同宗派的歷史.
是發生了什麼事.
今天我見到不同宗派.
我可以知道.
原來為什麼會有這個宗派.
原來為什麼這個宗派.
原來為什麼救世君.
沒有浸禮.
原來為什麼救世君沒有聖餐.
原來是這樣的.
我相信.
我們今天都是.
用什麼心態去聽歷史.
可能你說.
我真的不喜歡歷史.
我一聽就睡覺.
就像催眠術.
很好睡.
但我想.
每一個歷史裡面.
其實是有血 有淚.
有裡面充滿故事.
好 我們再繼續.
講歷史.
剛纔我們說我們聽歷史.
有點距離.
好像不太關我事.
講歷史呢.
大家有沒有講過歷史.
和別人講歷史.
有沒有.
有 有.
弟兄多些.
姐妹不太講歷史.
不是用性別分.

$^{561}$我都不多.
但原來我不知道是不是人老了.
人大了.
講多了.
我發覺我都提醒自己.
不要講那麼多.
開始一口講歷史.
我們不如進入經文.
第三章.
一齊讀.
字是否有點小.
一齊讀.
那時.
我吩咐你們說.
耶和華你們的神.
已將這地.
賜給你們為業.
你們所有的.
勇士都要帶著兵器.
在你們的.
弟兄以色列人.
面前過去.
但你們的妻子.
孩子.
新畜.
我知道你們有許多的新畜.
可以住在我所賜給你們的.
各城裡.
等到你們弟兄在.
若淡河那邊.
也得耶和華.
所賜給他們的地.
又使他們.
得享平安.
與你們一樣.
你們纔可以回到.
我所賜給你們.
為業的地.
這裡我摩西吩咐你們.
摩西和新一代.

$^{601}$以色列人說.
我們已經下到第三章.
他說什麼呢.
其實摩西說什麼呢.
對著一群女變人.
摩西人和嗎哪西半支派.
他們是怎麼樣的.
他們留在若淡河那邊.
以東以西以南以北.
是以東他們說.
那裡就已經得地了.
但神說不是.
當時出埃及記就說.
你可以留下你的.
妻子孩子可以留在那裡.
但是.
你藍丁就跟我.
和以色列其他的支派.
就要去若淡河.
另一邊去打仗.
打完仗之後.
拿了得為業.
之後你纔回到家.
是這個故事.
在出埃及記的時候.
這個就是摩西回顧.
再講一次.
回顧這件事.
我讀下一版.
我特意沒有寫在哪裡出.
大家一起讀一次.
你們要追念.
耶和華的僕人.
摩西所吩咐你們的話.
耶和華.
你們的神.
使你們得享平安.
也必將這地賜給你們.
你們的妻子.
孩子和親戚.

$^{641}$都可以留在.
若淡河東.
摩西所給你們的地.
但你們中間.
一切大能的勇士.
都要帶著兵器.
在你們的弟兄前面.
過去.
幫助他們.
等到耶和華使你們的弟兄.
像你們一樣.
得享平安.
並且得著耶和華.
你們神所賜他們.
為業之地.
那時纔可以回.
你們所得之地.
承受為業.
就是耶和華的僕人.
摩西在若淡河東.
向日出之地.
你們的.
你猜猜這句經文是哪裡出來的.
是誰說的呢?.
猜猜.
是誰說的?.
摩西?.
約書亞?.
跟剛纔那句是不是很相似?.
內容是不是.
差不多一樣?.
這個呢.
是約書亞記.
一章十二節至十五節.
原來呢.
我們想想約書亞.
當申命記的時候.
他也是站在那裡.
聽摩西說話的那個.
但到約書亞記的時候.

$^{681}$這裡是說約書亞對裡面人.
加德人和馬爾納西.
半支派的人說的話.
就是約書亞.
由聽歷史.
變成一個講歷史的人.
他不只是站在那裡聽.
他已經是約書亞的時候.
當摩西交給約書亞的時候.
在他的年代.
他就成為一個講歷史的人.
他就說了出來.
我用AI問一下AI.
摩西和約書亞.
約書亞就是這樣承接了.
摩西交代他的.
他就說了出來.
我們再下去說多一個.
在申命記第三章.
我叫維延.
我也說過.
申命記都是摩西維延.
我們一起讀.
三章21至22節.
你們繼續讀下一段.
約書亞.
真正的是耶和華.
你們的神.
這個約書亞是23章.
是約書亞記的最後.
是約書亞的維延.
原來約書亞的維延.
和摩西的維延真的很相似.
我特意錄下了一句.
我看到的時候.
我發覺我很感動.
約書亞我真的很感動.
他將摩西當年.
他講了幾十年前.
吩咐約書亞講的說話.

$^{721}$他臨死的時候.
叫了所有領袖.
叫了長老.
族長 審判官.
官長 所有的人.
去到他床前面.
講的是.
摩西的維延.
他再提醒.
以色列民.
已經是再之後一代的以色列民.
因那為你們.
真正的是耶和華.
你們的神.
不知道大家看到這個.
有什麼感覺.
我覺得很感動.
是因為.
神的應許很實在.
隔了幾十年.
神都叫約書亞.
提再新一代第三代的以色列人.
這個應許.
不要忘記.
幫你們打仗的.
是神.
不要忘記.
不要用自己的力量去打.
你不要忘記.
對著敵人你很害怕.
你不需要害怕.
在約書亞記裡面充滿一句.
不要懼怕.
是約書亞臨走的時候.
當年摩西和耶和華.
如何提約書亞.
不要懼怕.
他今天臨死的時候.
就要和再新一代的以色列民說.
你不要害怕.

$^{761}$是耶和華.
我不知道大家覺得.
我覺得這句話.
可以拿得很實.
是我們的靠山.
因那為你們征戰.
的是耶和華.
你們的神.
我們呢.
我不知道你們有沒有徵戰.
未必我們不是打仗.
你會不會都工作上.
家庭上.
關係上都有些征戰呢.
我想鼓勵你們.
無論見到什麼.
你覺得不公義.
或者處理不了.
或者真的要打仗.
我們祈禱吧.
神會為我們征戰.
祂是我們的神.
當年這句話.
相信.
不只是當年的歷史.
是今天.
都是有力量的一句神的話語.
Amen.
那大家怎麼樣呢.
我問.
你呢.
我不是說你快要流遺言.
你們很久都沒有.
那你們有什麼故事呢.
你們有沒有人.
分享過歷史.
剛纔有些人說.
或者我轉個角度說.
有沒有人分享過.
你有沒有和其他人分享過你的見證.

$^{801}$神在你生命的作為.
有沒有啊.
有啊,有些人有.
我都鼓勵大家.
可能這就是這次的小功課.
向一個人.
家人也好,同事也好.
朋友也好.
分享神在你生命的.
一個見證,一個作為.
那我.
叫阿男放進去,功課.
好奇,是嗎.
好奇在寫,是功課.
我們每個人都有歷史.
剛纔都見了支筆.
那我們下去.
由支筆寫歷史.
我今天快一點.
第三點.
寫歷史.
寫歷史不是用筆寫的.
寫歷史.
是我們每個人都會偏寫歷史.
那我問一下大家.
有多少人.
會經常在ig post story.
有沒有人沒post過story.
在ig或facebook.
沒有吧.
我們每個人都會poststory.
你post什麼多.
喫東西.
什麼.
聽不到.
什麼.
one piece.
好.
喫東西又有,新聞又有.
鸚鵡,哦,寵物.

$^{841}$post很多story.
那我們現在ig或facebook.
social media story.
多久就沒了.
24.
都是24,沒變.
24小時沒了.
但其實.
我們每一個人.
你想或不想也好.
在意或不在意也好.
其實我們的生命.
就在寫歷史.
我們每一天的生活.
都在寫歷史.
而這個歷史.
不是24小時就不見了.
我們所作的每一樣東西.
我們說的東西,我們的言語行為.
我們待人處事.
其實我們都在一起.
寫歷史.
你說不是.
歷史當然是大人物才寫的.
約書亞,整本約書亞記.
我們這些small potato.
哪有歷史.
或者肥佬Trump,Trump就天天在寫歷史.
還要影響全世界的歷史.
我們又不是肥佬Trump.
我們一個小市民.
怎麼寫歷史呢.
我很記得.
雖然我歷史.
我中學歷史真的不厲害.
但我記得我中一.
我中一的時候.
第一堂的歷史課.
我老師就跟我說.
歷史.

$^{881}$是history.
而這個his是神.
是大街his.
是神的.
神的故事.
是神在他創造的人.
在這個世界裡面的故事.
他就說.
為什麼我們要讀history.
可能他要這樣說我們纔有心機聽.
他就說.
history repeats itself.
歷史是會重複.
所以我們在歷史裡面要學習.
要學習過錯.
不要再重犯.
但是我覺得.
我們認識神.
我們所認識的神.
history歷史是不是他寫的呢.
我們說.
掌管歷史空間.
時間的神.
是不是他一個人寫的呢.
我想大家都會覺得.
不是.
當然不是.
我覺得我們所認識的神.
他很想參與.
我們的故事.
他想和我們每一個.
寫我們的歷史.
寫我們的故事.
寫我們個人的故事.
寫我們這個群體的故事.
我這段時間.
一月去了幾個安息禮拜.
不好意思.
過年前講完衛燕又講安息禮拜.
安息禮拜一定.

$^{921}$我今個月都去了幾個安息禮拜.
參與了.
一定有一個環節.
是什麼.
新平點滴.
講回先人的.
安息禮拜基督教.
會講弟兄姊妹的.
新平歷史.
他們的兒子或女兒或孫子會講.
我剛剛上個禮拜.
參與了一個.
我覺得很特別的安息禮拜.
這個安息禮拜.
是有笑聲.
還有當.
這個弟兄.
最後.
他兒子講新平點滴的時候.
他邀請大家.
一同拍掌.
鼓掌.
為世伯的生命鼓掌.
我覺得很感動.
因為世伯都連長.
但是我聽到.
他講新平點滴.
我第一次有這個感受.
聽新平點滴.
我覺得很被鼓勵.
因為他有連長.
他九十多歲.
聽回整個香港歷史.
每一環.
都緊緊扣住.
到最後.
他連腦的時間.
去認識神.
他很認真.
他用信紙寫了一封信.

$^{961}$還給他的兒子看.
我對不對.
他對神很認真.
對他的信仰.
所以那一晚.
他的兒子說.
我們為了爸爸.
我們一同鼓掌.
為了他的生命.
他一生很精彩.
我不知道大家有沒有去過.
安息禮拜.
我自己做了傳道人.
會更加多.
這個給我一個很感動的安息禮拜.
原來是一個會影響人的生命.
就算他今天在世.
已經沒了氣息.
都會影響人.
我在世伯身上.
學到堅毅.
聽到堅毅.
我在世伯身上看到.
他有節儉的美德.
看到他對家人的愛.
他有很多兒孫滿堂.
很疼愛孫子.
孫子可能跟你差不多大.
很疼愛孫子.
很關心孫子.
很認真.
原來當我這樣聽.
做一個第三者.
我沒見過世伯.
我真的參加.
我都會被鼓勵.
我想.
如果將來我的安息禮拜.
我的家人會怎樣說.
我的生平點滴呢.

$^{1001}$不知道大家有沒有想過.
你想人家怎樣說.
你的生平點滴呢.
我們的生命的歷史是什麼呢.
我們的生命歷史是什麼呢.
\newpage



\section{箴言 4:23}
\label{sec:031675}
\textbf{祝福 •由心而發}~~~~潘梁其德傳道
\newline
\newline
連結: \href{http://www.church.com.hk/acms/content.asp?site=cdc\&op=show\&type=product\&code=031675\&layout=sermon}{\texttt{031675}} ~~~~ 語音日期: 2026-02-15
\newline
\newline
\hyperref[sec:031674]{\small{< < < PREV SERMON < < <}}
~
\hyperref[sec:index_chronic]{\small{[返順時目]}}
~
\hyperref[sec:031676]{\small{> > > NEXT SERMON > > >}}
\newline
\newline
箴言 4:23
\newline
\begin{longtable}{cl}
\hline
\hline
章節 & 經文 (和合本修訂版)\\
\hline
4:23 & \begin{tabularx}{0.7\textwidth}{X} 你要保守你心，勝過保守一切，因為生命的泉源由心發出。 \end{tabularx} \\ \\
[1ex]
\hline
\hline
\end{longtable}
$^{1}$我們過年常常都會講祝福的說話.
但是賀年我就心裡有些迷思.
什麼叫迷思呢?.
第一個難處就是我常常都不知道祝福別人什麼好.
最順口的就是恭喜發財.
但是我覺得發財好像基督徒一樣.
好像很是悔.
但是我要知道對方有什麼需要.
然後我祝福他.
我又會覺得有些難處.
我記得我剛剛結婚的時候.
我夫家有很多親戚.
很多阿姨,姨丈.
其實我都搞不清他們是誰打誰.
誰在做什麼.
誰已經結了婚.
誰是誰.
搞不清的.
那個情景通常都是.
你一進去拜年.
去到老爺,奶奶那裡拜年.
已經高堂滿座了.
嫂子來了.
很多人已經說恭喜恭喜.
然後我就說他是誰呢?.
搞不清.
然後他們一卵嘴的祝福我.
我都不知道祝福他什麼好.
我都不知道他有什麼需要.
所以我永遠都是祝福這四個字的.
身體健康.
不是.
又不是萬事勝意.
他祝福我.
我就說恭喜恭喜.
大家這樣說.
這樣我就灑過去了.
有什麼人需要我們祝福呢?.
第二個我的迷思就是.
其實我信不信我的祝福是有效的呢?.

$^{41}$我祝福他步步高升.
那他是不是真的會步步高升呢?.
我祝福他學業進步.
那他就會學業進步了.
不要說你自己不相信.
聽過的都不相信.
就算他真的步步高升了.
他會不會.
哎呀我今年可以加薪升職.
一定是過年的時候.
伯母祝福了我步步高升.
會不會呢?.
不會.
根本不會想起你.
因為你根本就不會記得.
誰祝福過你什麼.
你說的那個又不相信.
聽的那個也不太相信.
那你說來幹什麼呢?.
所以我今天就想一下.
今天的題目叫做「祝福由心而發」.
我們從聖經裡面去祝福的經文裡面.
去反思一下這個題目.
究竟我們需要什麼樣的祝福呢?.
究竟怎樣才能夠祝福我們呢?.
這個經文記載在一個我們很少讀的書卷裡面.
是全本聖經最短的書卷.
叫做《約翰三書》.
《約翰三書》當然是約翰寫的.
這本書是寫給誰的呢?.
他也是這麼說的.
他說「著長老的寫信給親愛的戈尤.
就是我誠心所愛的」.
很簡單.
開章明義你就知道.
收信人叫做戈尤.
是仕途約翰很疼他很愛他.
為什麼這麼疼他這麼愛他.
他做了什麼呢?.
聖經是沒有記載的.

$^{81}$不過有記載的是什麼呢?.
他怎樣稱讚他.
聖經裡面稱讚戈尤.
是會有真理在他的心.
不僅有真理還要行真理.
又有信心又有愛心.
他做了什麼呢?.
他就經常接待遠方來的弟兄姊妹.
可能那些是宣教士.
他就接待他們.
整本書就是鼓勵他們.
你們真的做得好.
繼續做.
記住繼續做.
就是這樣.
所以很短.
所以我們很少查經會查約翰三書.
不過約翰三書裡面有一句話.
有很多人.
如果要講約翰三書就會講這句.
就是這句是第一節.
因為他只有一章.
所以約翰三書的第一節是這樣.
第二節就是今天的金句.
你們讀吧.
好不好?.
一二三.
第一個要提的就是.
我祝福你.
祝福他.
但英文你會看到.
是什麼?.
祝福是什麼?.
I pray.
不是我們有什麼能力去祝福別人.
因為我祈禱神可以祝福.
阿蓋尤.
OK?.
祂祝福了他三樣東西.
第一樣是什麼?.

$^{121}$身體健壯.
很容易明白.
身體健康.
剛剛大家站在那裡.
誰說.
神啊你祝福我身體健康的.
有沒有?.
好啊.
這麼少?.
不是吧.
我覺得每個人都很著重身體健康.
不要緊.
聖經很落地.
祂都知道我們要身體健康.
第二樣祂祝福我們的.
約翰祝福阿蓋尤的.
就是叫他凡事興盛.
接著祂就祝福他.
靈魂要興盛.
很簡單的一句經文.
雖然簡單.
我們有些字都要解釋.
什麼叫做興盛?.
興盛.
英文叫做prosper.
prosperity.
發達,興旺,茂盛.
原文叫做You are dull.
為什麼我要想一想呢?.
我要想一想我怎樣記住祂怎麼讀.
我怎樣記呢?.
You are dull.
即是你很悶.
千萬不要那麼悶.
講到要很興奮一點.
You are dull.
這個原文.
當翻譯成為聖經裡面.
在不同的地方.
會有不同的意思.

$^{161}$包括亨通.
順利.
包括.
我看不到那些字.
因為下次那些字會大一點.
我這樣看不到.
成功.
不過最貼切的翻譯.
就是有一條很容易走很平坦的路.
那個就叫做You are dull.
那個就叫做豐盛.
就是說祝福他.
凡事都可以去到他的目標.
做的時候很容易.
沒什麼阻滯.
就叫做You are dull.
就是順利,亨通.
所以在不同的經文裡面.
譬如在創世記24章.
出現過You are dull這個字.
這個字是說.
亞伯拉罕年老的時候.
他的兒子沒有老婆.
於是他就派了一個僕人.
去遠方跟他娶老婆回來.
那個僕人就覺得.
這件事很難.
無端端找一個女人回來嫁你兒子.
你想一下那個.
mission impossible.
但是他真的去了.
真的按著神的帶領.
他真的見到尼伯加.
見到那個女孩又漂亮又好.
就很適合條件.
這句就是他見到尼伯加.
下一句他就這麼說.
那個人哪人.
定睛看他.
哪人就是那個僕人.

$^{201}$眼金金的看著尼伯加.
一句話都不說.
要曉得耶和華.
似他通達的道路沒有.
那個通達的道路.
就是說看看神給不給他.
很順利.
娶了他回來.
帶了他回去.
給他的主人.
做媳婦.
另外一節就是.
約書亞記.
約書亞記就是摩西.
臨死之前他吩咐約書亞.
叫他你要進入江南.
你要剛強壯膽.
其中有一句.
他勉勵他的說話.
你記住要守神的律法.
不可以離開你的口.
要有思想要跟著你走.
如果是這樣.
你的道路就可以.
亨通.
你的路就很容易走.
會很順利.
凡事順利.
這個就是.
輕王輕聖.
背後的意思.
ok.
第二個想和大家說.
的詞語就是靈魂.
靈魂其實我們.
虐莫都知道.
為什麼.
虐莫之餘也有些模糊.
什麼是靈?什麼是魂?.
這個詞語.

$^{241}$原文就是黑色那個.
叫做suhet.
如果你讀英文.
你心理學的東西.
有時.
在聖經裡面.
譯做我們的魂.
有時譯做我們的心.
我們的心.
有時候說要稱頌它.
那個就是我們的魂.
都是同一個詞.
有時會譯做思想.
有時譯做氣.
一口生氣.
一口令我們活起來.
那個.
都是這個詞.
簡單來說.
我們不分靈和魂.
靈魂.
要解釋.
我就給一個解釋.
你就是一個人.
你就望見Auntie Mary.
Auntie Mary你看到我的樣子.
你如果走過來和我握手.
你可以碰到我.
那個就是我的身體.
但Auntie Mary就不只是一個殼.
Auntie Mary有裡面.
我有我的思想.
我有我的情緒.
我有我的過去.
我有我的價值觀.
那個看不到部分的Auntie Mary.
就是我的靈魂所在.
就是說.
我們要我們的靈魂興盛.
我們看不到的部分都要興盛.

$^{281}$還有一個很特別的地方.
就是.
它包括了.
我們的思想,情緒,價值觀.
等等.
所有的東西.
它和我們的肉體不同.
我們的肉體會消失.
一天我們死了.
我們的肉體就消失了.
但是我們的靈魂.
是不滅的.
我們死了之後.
我們的靈魂會回到.
我們的天賦那裡.
大家都明白了.
我們再讀一次.
今天的金句.
1,2,3.
很有趣.
約翰.
祝福了他三樣東西.
A.
身體健康.
B 凡事興盛.
C 靈魂興盛.
是不是這樣.
應該是這樣.
三樣東西.
但他中間就加了兩個字.
正如.
他不是祝福你.
A,B,C.
他是這樣說的.
我祝福你A.
祝福你B.
正如你的C.
那是什麼意思.
就是說.
如果你有了C.

$^{321}$你就會有A.
還有B.
你靈魂興盛.
是和凡事興盛.
身體健壯.
這是基礎.
你有了靈魂興盛.
你就會身體健康.
你就會凡事興盛.
是不是.
有了C.
就會有A.
就會有B.
令我想起.
另外一段經文.
我不知道大家.
如何祝福自己.
祝福自己身體健康.
很多人都舉手.
祝福自己凡事興盛.
可能我們不是用這個詞語.
但是.
多數都是.
想那樣得那樣.
得心應手.
家和萬事興.
工作愉快.
步步高升.
順順利利.
我們可能是祝福這些.
我不知道.
但是我們通常.
祝福自己的事.
都是我們生命裡面.
要面對的問題.
我們覺得很難的事.
所以我們就想這樣祝福.
沒有那些很難的事.
沒有我們的煩惱.
就好了.

$^{361}$但是這段經文.
告訴我們.
你要A.
你要B.
你只要有C.
你有了C.
你就有A.
和有B.
OK.
我們.
我不說你們.
很少.
祝福自己的靈魂.
是的.
因為我們都是人.
我們都很顧念.
實體的東西.
忘記的東西.
正在經歷的東西.
很實在.
但是聖經裡面有這樣一句經文.
你們一起讀.
123.
原來神教我們.
我們很著重.
我們肉體.
實質的東西.
但是神提醒我們.
聖經裡面提醒我們.
看不到的東西最重要.
你有了C.
就會有A和B.
我們來看看.
是不是真的.
是不是這樣的.
究竟.
靈魂輕盛.
和身體健康.
是不是真的有關呢.
或者在這裡.

$^{401}$你的心裡面會問.
身體健不健康.
你知不知道.
知道的.
你沒有痛,沒有病,沒有癢.
沒有腫,沒有嘔,沒有暈.
你就沒有病.
你身體健康.
通常是這樣.
靈魂輕盛.
我怎知道我輕不輕.
我盛不盛.
我怎樣測度.
身體的健康.
有很多東西可以測度.
就算我沒有暈,沒有嘔.
我這個年紀都要年年驗身.
哎呀,你的血糖不行.
哎呀,你的體高體重不達標.
你的BMI不合格.
你的血脂不合格.
你的血糖不合格.
你的膽固醇太高.
是不是.
你的血壓太高.
你可能沒事的.
你感覺不到你有事.
但是有個標準.
你量一量.
驗一下血,驗一下尿.
照一下X光,照一下MRI.
就可能照出很多問題.
及早照出.
你就沒有問題.
就不要等事已經變成不可救藥.
你就去搞它.
身體.
我們除了感覺之外.
我們還有很多的指標.
去衡量.

$^{441}$我們是不是健康.
我們的靈魂.
怎樣衡量.
它是否輕盛.
你怎知道你的靈魂.
是否輕盛.
我們怎樣量度.
大家聽過.
FOS.
Index嗎.
有沒有人會聽過.
沒有嗎.
我都預了我.
因為是我作的.
我就想.
怎樣跟弟兄姊妹說.
他們的靈魂是否輕盛.
我怎知道自己的靈魂是否輕盛.
就彈了這句經文出來.
大家一起讀.
大家都很熟悉.
這是甚麼.
聖靈所設的果子.
Fruit of Spirit.
FOS Index.
大家反應真好.
竟然有人拍掌.
我以為我會被人聽到.
我們怎樣知道.
我們的生命.
我們裡面.
既然驗血又驗不到.
照X-ray又照不到.
照MRI都照不到.
我們的靈魂狀況.
那是要我們.
耐性.
我們去測驗.
我們裡面的心是怎樣.
心是怎樣.

$^{481}$你的心是充滿愛.
還是充滿恨.
你的心是充滿喜樂.
還是很重.
很憂傷.
你的心.
是很平安.
跟人的關係.
跟自己的關係.
跟神的關係.
都很舒服.
還是你經常坐立不安.
很焦慮.
騰來騰去.
好像有很多事未做.
很多事搞不定.
很怕靜.
你有沒有忍耐.
去面對你現在要面對的困難.
那些困難的人和事.
還是你一想起.
要上班見到那個人.
你就很煩躁.
你有沒有恩慈.
你對人.
有沒有包容.
有沒有憐憫.
還是你很冷酷.
你覺得我自己.
已經透不過氣.
還要我忍他.
還告訴我.
人人都有別人的難處.
已經夠多.
閉嘴.
還要我忍.
我們是否良善.
還是人家經常說.
媽媽你很兇.
爸爸你很兇.

$^{521}$你的下屬老闆你很兇.
通常人家不會告訴你的.
不過你在廚房聽到.
很壞的.
聽到人家說自己的壞話.
很惡.
我們很鮮酸刻薄.
我們很.
很mean.
我們.
有沒有信心.
還是我們的生命裡面充滿了很多恐懼.
外面可以很惡.
但心裡面可以很怯.
我們怕.
我們的工作沒了.
我們怕這樣怕那樣.
我們怕將來不知怎樣.
或者不敢去.
我們很怕.
我們的婚姻快不行了.
但我們不敢去.
連那個.
面對面唱情歌.
我都不敢去.
我怕去了之後.
不知道有什麼爆了出來.
我們是否溫柔.
還是我們是暴躁.
我們有沒有.
自控的能力.
我們有沒有節制.
還是我們經常情緒失控.
我們經常.
花錢失控.
我們經常飲食失控.
不應該吃那些東西.
拼命吃.
拼命飲.
不應該做的事.

$^{561}$我們知道不應該做的.
應該放下手機.
不要再打了.
但我們失控.
我們沒有了自制的能力.
FOS Index.
留意.
我不是說人生裡面.
不會有.
憂傷的時候.
有的,你家裡有人去世了.
難道你哈哈哈哈.
我是一個很好的基督徒.
我上上起下笑.
不是.
聖經裡面說.
哭有時笑有時.
但是如果.
那種情緒.
是長期的.
過份的.
那個負面情緒.
是長期過份的.
控制著你的生命.
控制著你的日常.
那.
你的靈魂就不會出現問題.
不是很豐盛.
人會怕的.
明天考試怕不怕.
考車牌,多怕.
那是不是你靈魂出現問題.
不是.
不過.
如果那種怕怕到.
三日睡不著,四日吃不下.
那就是過份.
原來我的生命.
是有一種很大的.
焦慮在.

$^{601}$影響著我和人.
對事的關係.
那樣的監察.
是你自己.
自治.
你的靈魂.
豐不豐盛.
如果是這樣.
我們的健康.
和我們的靈魂.
豐不豐盛.
扯不扯得上關係呢.
扯得上的,這個很容易明白.
所以先說.
大家都知道的,這句金句.
大家都讀過的,你又讀一次給我聽.
1,2,3.
開心就藥都不用吃.
憂傷.
可以入骨.
你這樣不開心.
你表面可以笑.
你裡面.
一直在哭.
聖經這樣說.
有沒有科學根據的.
有的,有醫學根據的.
宋醫西醫.
都是這樣說.
看書.
有樣東西叫.
七,我要看,因為不熟.
七情內傷.
中醫說的.
七種的情緒.
如果過度.
過分的傷.
過分的興奮.
過分的擔憂.
過分的憤怒.

$^{641}$如果過分.
這麼大.
很持久,不斷這樣的話.
就會傷及.
我們的內臟.
傷及我們的氣血.
根據中醫.
憤怒就會傷肝.
你抱喜.
原來開心.
在中醫來說是有問題的.
就會傷心.
你太過興奮.
很多人就在那裡出事.
是不是.
然後他又說.
你過分的思慮.
就會傷你的脾臟.
你過分的悲哀.
就會傷你的肺.
你過分的恐懼.
就會傷你的腎.
不是我說的,是中醫說的.
你說.
我相信西醫,西醫都是這樣說.
西醫說.
超過70\%的病.
是由情緒引起的.
你裡面的.
壞情緒.
你的緊張.
你的不安.
你的睡不著,你的憂慮.
是會引起.
很多疾病.
當然第一樣引起的就是情緒病.
抑鬱.
焦慮.
會引起腸胃病.
你會胃潰瘍.

$^{681}$你會搶耳機.
一緊張就要衝去廁所.
你會有心臟病.
你的血壓會飆升.
你會心肌緊塞.
心肌緊塞.
你的.
免疫力.
會降低.
因為你的情緒不好.
所以你就會經常感冒.
不會的.
剛好那個一刻.
你又來過.
你的免疫力.
影響.
你的敏感,你的濕疹.
很多的病.
都是由情緒.
引起.
所以.
聖經的話語是對的.
你有.
輕盛的靈魂.
你就有.
強健的身體.
我想這個.
大家都不需要.
多.
講解.
我想醫學上.
大家上那麼多YouTube.
不過再證實一次.
聖經說的話.
幾千年前說的話.
是對的.
我們來到下一個點.
究竟.
靈魂輕盛.
跟我們.

$^{721}$凡事輕盛.
有什麼影響.
是不是真的我靈魂輕盛.
我就會凡事輕盛.
我們來看看這個人.
這個人.
怎樣.
你覺得他怎樣.
開心嗎.
不開心.
很辛苦.
做人很辛苦.
四方八面的壓力.
又怕沒了工作.
有工作又要做到很晚.
又要交租.
媽媽又病了.
老婆又跟你吵架.
上司又不喜歡你.
下屬又罵你.
又不聽你的話.
回教會.
組長又不理你.
凱哲臣都不理你.
OK.
你的子女長大了.
本來是你的小棉襖.
現在就一言九鼎.
頂你.
簡直是.
諸事不順.
小人當道.
是非前身.
入不敷支.
得不到壓力嗎.
為什麼你這麼開心.
見到一個這麼慘的人.
你們快點照照自己有沒有恩慈.
就是在想那個是你自己.
我們覺得很辛苦.

$^{761}$年近衰晚.
YouTube就會上一些.
你今年流年好不好.
你的命.
犯太衰.
你改一下.
你的風水不好.
你改一下.
你說阿姨我們信耶穌.
我們不想這些.
真的.
你心裡有沒有說過.
我的命很苦.
我不明白.
為什麼我見到一個壞人.
拍完一次拖.
說了被人辜負.
第2個第3個第4個.
都是這樣非我.
有沒有搞錯.
我的命好不好.
我真的好不幸.
求你讓我凡事興盛.
求你讓我每一件事都順順利利.
不要讓我頭頭碰著黑.
人最大的問題.
就是覺得他的問題.
是由那些東西來的.
聖經上告訴我們.
我們的問題.
是為什麼不興盛.
是我們的心的問題.
阿姨你覺得教會不好.
沒有了這間教會.
我去過第2間我就好了.
那你就算了.
去到第2間.
又是這樣.
又有很多東西讓我看不順眼.
我就去第3間.

$^{801}$結果都是這樣.
初時都不錯.
為什麼後來這樣.
你知道為什麼嗎.
因為你去了.
搞到人家這樣.
你見到.
我老公真是壞.
不愛他離婚.
嫁過第2個.
這次學聰明了.
又嫁過.
伊利沙伯泰來漂亮嗎.
有沒有錢.
嫁了多少次.
8次.
真是多謝.
我忘了他嫁了多少次.
太多次了.
厲害.
為什麼.
他以為他婚姻的問題.
是他老公的問題.
他以為是他自己不幸.
他每次見到工作都和人吵架.
每個人都針對他.
他以為那些是.
他們的問題.
所以我們祈禱是怎樣.
求你叫我興盛.
叫我順利.
叫那些小人離開我.
叫我碰見好的上司.
叫我碰見好的下屬.
我們以為.
我們生命很多的不順利.
是因為什麼.
是因為姦會.
公司 老公 家人.
爸爸 媽媽.

$^{841}$是他們的問題.
於是我們覺得.
我們的命不好.
我們的運氣不好.
很多人.
未信耶穌的.
他們都覺得風水不好.
我就做一些事.
讓自己安樂一點.
對不起.
聖經反過來對我們說.
當我們裡面.
興盛.
我們外面.
就會興盛.
看看這兩個.
圖片.
就是一個很漂亮的花園.
鳥語花香.
那個是我們心靈的花園.
你心靈的花園.
聖經裡面說.
如果是很漂亮.
很整齊.
很茂盛.
這樣.
你就會道路行通.
是因為你裡面的心力.
改變了外面的環境.
相反.
如果你的心靈是.
雜草叢生.
又沒有人淋水.
你又不修護它.
破破爛爛.
你不理它.
你的生命.
就這樣.
如果你的生命.
是鳥語花香.

$^{881}$你就可以.
站直.
給兩隻拇指.
祝福到你自己.
而這個祝福.
就會帶去你的家.
帶去你的公司.
帶去你的教會.
帶去你的社會.
帶去你周圍的人.
因為.
因為什麼?.
因為我們的心靈.
豐盛.
老實說.
我們這裡坐在這裡.
每一個人.
生老病死.
每一個人都經歷過挫折.
每一個人生命裡面.
總被人欺騙過.
被人辜負過.
有人令我們失望過.
但是.
當我們心靈豐盛的時候.
我們會有那個韌力.
有那個信心.
和盼望.
去將.
那些不好的東西.
轉化成我們生命的智慧.
轉化成.
我們生命的力量.
相反.
同樣.
有很多人.
受了傷之後.
就不知道為什麼黑化了.
在他裡面有很多痕.
很多的怨.

$^{921}$令到.
我們的心靈不豐盛.
接著.
這幅圖片.
有多少人.
問一些這樣的問題.
明知我騙你的.
不是四個人.
其實這個是一個人.
是我們.
這幅圖.
代表什麼呢.
我們人生的生命裡面分開很多部分.
不是兩部分.
外面和裡面.
我分了四部分.
第一個部分就是我們的身體.
一個.
完美的組合.
就是我們的身體.
是跟著我們的魂.
我的魂是跟著我的靈.
是走的.
最先那個是什麼.
是你裡面有的東西.
外面的人沒有的.
那個就叫做聖靈.
當我們信了耶穌之後.
聖靈就在我們的心裡.
在你裡面.
一個.
理想.
豐盛的生命.
是應該這樣的.
聖靈行先.
做大靈.
我們的靈.
豐盛的.
願意跟隨神.
信任他.

$^{961}$聖靈說什麼.
那我就.
遵行你的旨意.
情緒OK.
忍一忍.
有平安有喜樂.
然後我的魂.
就會命令我的身體.
作出適當的反應.
所以.
聖靈說.
嗯.
星期日要送拜.
那個靈說.
好啊.
我們打開我們的心.
張開我們的口去敬拜我們的神.
好不好啊.
我有點不開心.
都要打開我們的口.
阿雲.
阿雲就說.
嗯 好吧.
打開你的口.
舉起你的手.
於是我的身體就.
我們就敬拜我們的神.
我們就經歷了.
那種從神而來的能力.
我們的道路.
就順利了.
相反是怎樣.
相反.
那些次序混亂了.
你的靈很弱.
誰說了算.
魂說了算.
你知不知道有很多事情要做.
每個項目都沒做完.
叫我回崇拜.

$^{1001}$這個禮拜不回.
下個禮拜也不回.
讓我先搞定所有事情.
努力工作.
阿雲就說.
堅持努力.
你壞了.
你死定了.
你怕不怕.
於是你就拼命在這裡做.
拼命在這裡做.
於是你的靈就越來越弱.
你就有很多的壓力.
你心裡面.
做得好的.
做得很好之後.
就怎樣.
下次要做得更好.
不然別人就不稱讚你.
人家就不升你職.
你就死定了.
現在經濟這麼差.
你很多的害怕.
你的生命.
就是一條不同的路.
再讀這一句.
我們結尾的金句.
很熟.
大家再讀.
123.
一生的果效.
是由你的心發出.
這個經文真的很熟.
熟到我們有點念口.
但是今天你再看的時候.
你明不明白.
什麼叫一生的果效.
是由你的心發出.
你的心就是你的靈魂.
你的靈魂.

$^{1041}$健不健康.
輕不輕盛.
你一生的果效.
你一生的是不是萬事亨通.
事事順利.
是不是如意吉祥.
是不是得心應手.
不是外面的東西.
不是你好不好彩.
不是你好不好運.
就算你很慘地.
遇見一個很壞的上司.
你裡面的愛.
你裡面神給你的能力.
可以勝過他.
可以改變他.
由你裡面出來的生命力.
可以勝過你周圍的環境.
你是影響世界的那個人.
不是被世界影響和踐踏的那個.
你是被世界影響和踐踏的那個.
記不記得那句金句.
親愛的弟兄.
我祝你.
不是我有能力.
是我祈禱.
神叫你.
凡事輕盛.
神的心意是想你輕盛.
輕和.
想你身體健康.
不過在此之先.
你要心靈健康.
你要靈魂豐盛.
你生命裡面要結出.
屬靈的果子.
因為那個是你道路亨通的理由.
ISIS是做什麼的呢.
很多人來到.
就將他們生命裡面的難處跟我們說.

$^{1081}$我們不會教他們轉工.
這間公司那麼差.
換個老公.
離婚.
轉教會.
那些是解決.
或者叫你解決外面的問題.
我們是會.
我們聽了那些問題.
我們去找.
不輕盛.
當然是心有事.
心有什麼事.
心受傷.
心想錯了.
心信錯了.
所以.
我們就跟他.
醫治釋放.
我們解決不了心.
其他的東西.
神會.
走在你前面.
帶領你.
你的心靈是健康的時候.
你的靈就會命令你的魂.
你的魂就會命令你的身體.
大家都走去你想走的那條道路.
而那條道路是輕散的.
是相信的.
是亨通的.
因為.
是神帶著你們.
我們一起起來.
我們一起禱告.
親愛的弟兄.
我願你凡事輕盛.
身體健壯.
正如你的靈魂.
輕盛一樣.

$^{1121}$親愛的姊妹.
親愛神的.
國民.
親愛的.
大家.
神的心意.
是你凡事輕盛.
神的心意是要.
祝福你身體健壯.
神的心意.
是要你靈魂輕盛.
我們的內心.
有哪裡需要除草.
有哪裡枯乾了.
有哪些負面的情緒.
在掌控著我們.
叫我們不斷.
不斷地撞板.
叫我們無力.
叫我們不失信心.
叫我們.
做不到.
我們想做的事.
神啊我們在這裡.
我們交代給你.
幫助我們.
靈魂輕盛.
幫助弟兄姊妹.
去察覺他們的內心.
讓他們信心知道.
這是神的心意.
要他們輕盛.
賜福給我們.
改變我們人生的命運.
從改變我們的心.
開始.
我們敬拜你.
我們讚美你.
願你的心意.
在我們每一個弟兄姊妹當中.

$^{1161}$成就.
祝福我們.
禱告奉耶穌基督.
得勝的命強.
我們一起說.
\newpage



\section{創世記 22:1-19}
\label{sec:031676}
\textbf{通過「考驗」成為敬畏神的人}~~~~黃浩君傳道
\newline
\newline
連結: \href{http://www.church.com.hk/acms/content.asp?site=cdc\&op=show\&type=product\&code=031676\&layout=sermon}{\texttt{031676}} ~~~~ 語音日期: 2026-02-14
\newline
\newline
\hyperref[sec:031675]{\small{< < < PREV SERMON < < <}}
~
\hyperref[sec:index_chronic]{\small{[返順時目]}}
~
\hyperref[sec:031677]{\small{> > > NEXT SERMON > > >}}
\newline
\newline
創世記 22:1-19
\newline
\begin{longtable}{cl}
\hline
\hline
章節 & 經文 (和合本修訂版)\\
\hline
22:1 & \begin{tabularx}{0.7\textwidth}{X} 這些事以後，神考驗亞伯拉罕，對他說：「亞伯拉罕！」他說：「我在這裡。」 \end{tabularx} \\ \\ \relax
22:2 & \begin{tabularx}{0.7\textwidth}{X} 神說：「你要帶你的兒子，就是你所愛的獨子以撒，往摩利亞地去，在我指示你的一座山上，把他獻為燔祭。」 \end{tabularx} \\ \\ \relax
22:3 & \begin{tabularx}{0.7\textwidth}{X} 亞伯拉罕清早起來，預備了驢，帶著跟他一起的兩個僕人和他兒子以撒，劈好了燔祭的柴，就起身往神指示他的地方去了。 \end{tabularx} \\ \\ \relax
22:4 & \begin{tabularx}{0.7\textwidth}{X} 到了第三日，亞伯拉罕舉目遙望那地方。 \end{tabularx} \\ \\ \relax
22:5 & \begin{tabularx}{0.7\textwidth}{X} 亞伯拉罕對他的僕人說：「你們和驢留在這裡，我和孩子要去那裡敬拜，然後回到你們這裡來。」 \end{tabularx} \\ \\ \relax
22:6 & \begin{tabularx}{0.7\textwidth}{X} 亞伯拉罕把燔祭的柴放在他兒子以撒身上，自己手裡拿著火與刀；於是二人同行。 \end{tabularx} \\ \\ \relax
22:7 & \begin{tabularx}{0.7\textwidth}{X} 以撒對他父親亞伯拉罕說：「我父啊！」亞伯拉罕說：「我兒，我在這裡。」以撒說：「看哪，火與柴都有了，但燔祭的羔羊在哪裡呢？」 \end{tabularx} \\ \\ \relax
22:8 & \begin{tabularx}{0.7\textwidth}{X} 亞伯拉罕說：「我兒，神必自己預備燔祭的羔羊。」於是二人同行。 \end{tabularx} \\ \\ \relax
22:9 & \begin{tabularx}{0.7\textwidth}{X} 他們到了神指示他的地方，亞伯拉罕在那裡築壇，把柴擺好，綁了他兒子以撒，放在壇的柴上。 \end{tabularx} \\ \\ \relax
22:10 & \begin{tabularx}{0.7\textwidth}{X} 亞伯拉罕就伸手拿刀，要殺他的兒子。 \end{tabularx} \\ \\ \relax
22:11 & \begin{tabularx}{0.7\textwidth}{X} 耶和華的使者從天上呼喚他說：「亞伯拉罕！亞伯拉罕！」他說：「我在這裡。」 \end{tabularx} \\ \\ \relax
22:12 & \begin{tabularx}{0.7\textwidth}{X} 天使說：「不可在這孩子身上下手！一點也不可傷害他！現在我知道你是敬畏神的人了，因為你沒有把你的兒子，就是你的獨子，留下不給我。」 \end{tabularx} \\ \\ \relax
22:13 & \begin{tabularx}{0.7\textwidth}{X} 亞伯拉罕舉目觀看，看哪，一隻公綿羊兩角纏在灌木叢中。亞伯拉罕就去牽了那隻公綿羊，獻為燔祭，代替他的兒子。 \end{tabularx} \\ \\ \relax
22:14 & \begin{tabularx}{0.7\textwidth}{X} 亞伯拉罕給那地方起名叫「耶和華以勒」 。直到今日人還說：「在耶和華的山上必有預備。」 \end{tabularx} \\ \\ \relax
22:15 & \begin{tabularx}{0.7\textwidth}{X} 耶和華的使者第二次從天上呼喚亞伯拉罕， \end{tabularx} \\ \\ \relax
22:16 & \begin{tabularx}{0.7\textwidth}{X} 說：「耶和華說：『你既行了這事，沒有留下你的兒子，就是你的獨子，我指著自己起誓： \end{tabularx} \\ \\ \relax
22:17 & \begin{tabularx}{0.7\textwidth}{X} 我必多多賜福給你，我必使你的後裔大大增多，如同天上的星、海邊的沙。你的後裔必得仇敵的城門， \end{tabularx} \\ \\ \relax
22:18 & \begin{tabularx}{0.7\textwidth}{X} 並且地上的萬國都必因你的後裔得福，因為你聽從了我的話。』」 \end{tabularx} \\ \\ \relax
22:19 & \begin{tabularx}{0.7\textwidth}{X} 於是亞伯拉罕回到他僕人那裡。他們一同起身，往別是巴去，亞伯拉罕就住在別是巴。 \end{tabularx} \\ \\
[1ex]
\hline
\hline
\end{longtable}
$^{1}$大英姐妹平安.
開始之前我們有個禱告.
一起來透過創世記22章1-19節.
看一下信心之父阿巴拉罕的生命.
我們知道今天你和我們對話.
不是要我們增加知識.
而是我們要更新我們的生命.
求聖靈在我們當中.
讓我們專心在你的話語裡面.
賜下一顆學無理願意受教的心.
禱告是奉主明求 阿們.
有一位年輕人.
他要申請成為宣教士.
一間警會的主考官就打給他.
叫他星期一凌晨三點來到警會面試.
年輕人都以為他聽錯.
再三確認 但警會都說是凌晨三點.
到了星期一晚.
那一晚很冷 只有十度.
天還沒亮.
他就準時三點到警會辦公室面試.
他一直等 等到九點 主考官才出現.
主考官就跟他說.
我們現在面試開始了.
第一我要考你英文.
請問你蘋果的英文是怎樣串和是甚麼.
年輕人當然馬上回答.
Apple 主考官說很好.
那我考你第二條問題 是數學題.
二乘二等於多少.
年輕人回答 四.
主考官說非常好.
你已經通過了面試.
明天我會在董事會上推薦你成為宣教士.
到了董事會.
董事會都質疑.
考官的考試內容是否有點兒戲.
但主考官卻稱讚這位年輕人.
說到他其實已經具備宣教士的資格.
第一 我叫他天寒地凍十度的早上來到辦公室.

$^{41}$他放棄了暖暖的被窩.
沒有埋怨 沒有反對.
按照要求來到辦公室.
我正在考驗他是否順服的人.
他通過了.
第二 我叫他早上三時來到辦公室.
他準時出現.
我正在試他是否一個有責任感.
會準時的人.
他都通過了.
第三 我九時才出現.
特意要他等六小時.
他照樣這樣做.
沒有任何發脾氣的跡象.
其實我是在測試他有沒有耐性和脾氣.
他都通過了.
最後 我要求他回答一些幼稚園都認識的問題.
他沒有投訴.
其實我是在測試他是否一個謙卑的人.
他都通過了.
其實我看過他的履歷表.
他的客觀條件.
他的服侍和神學訓練一向都OK.
我要測試他是否有成為神僕人的質素.
這個故事是一個宣教的差會.
如何考驗一個準宣教士.
是否符合成為神僕人的資格的一個方法.
雖然這個考驗好像有點兒戲.
但的的確確.
神最看重的是我們內在的質素.
而不是我們平時寫的那份履歷表.
而我們生命裡.
神也可能用不同的方法去考驗我們.
今天我們一起來看創世記二十二一至十九章.
一段大家都可能很熟悉的經文.
就是亞伯拉罕憲義煞.
神用了一個出人意外的方法去考驗亞伯拉罕的內心.
考驗往往在我們最不舒服的時候.
最不理解的時候.
看看我們是否看重神.

$^{81}$過於我們覺得重要的事.
我們來看看他的經歷.
亞伯拉罕的經歷.
讓我們去學習怎樣可以通過神的考驗.
在經文第一節說到.
這些事以後.
神考驗亞伯拉罕.
究竟這些事是包括什麼事呢?.
首先我們要明白.
創世記記錄亞伯拉罕的一生裡.
都圍繞著十二章當中.
神命令亞伯拉罕要離開本族本家.
他要成為大國.
要賜福給他.
地上萬族都因他而得福.
這個應許.
而這個應許.
很大程度上.
都是要透過他的後裔去實現.
而我們都知道.
在十二章至二十一章.
其實出現過幾次的後裔危機.
包括亞伯拉罕兩次說他太太撒拉是他的妹妹.
令撒拉差點被法老和亞比米勒玷污.
同時亞伯拉罕都多次想用自己的方法.
去實現神給他的應許.
當中他以為他的管家.
以利爾謝或下甲的以實瑪利.
是他的後裔.
但神多次強調.
撒拉所生的才是應許那一位.
而最終在二十一章的時候.
應許之子以實誕生.
神引證了他是信實的神.
會自己成就他的應許.
而神亦和亞伯拉罕說.
因為以殺身的才稱為你的後裔.
所以總括來說.
亞伯拉罕所有幫助神實驗應許的方法.
最終被神否定了.

$^{121}$而且亦只是為自己和他的後裔增添更多麻煩.
這些事以後.
神就考驗亞伯拉罕.
考驗即是試驗.
這裡不是說神不知道亞伯拉罕的情況.
要試驗他.
相反在聖經裡.
神要試驗一個人的時候.
是要讓當事人知道自己的情況.
同時成為後世的人的一個教材.
《出埃及記》是這樣說的.
不要懼怕.
神降臨時要試驗你們.
使你們敬畏他.
所以試驗的目的.
是要讓當事人知道什麼是敬畏神.
同時在歷代後世的讀者面前.
示範什麼是敬畏神.
而最終如果我們跳到後面的經文.
當亞伯拉罕通過考驗的時候.
天使都確認說.
現在我知道你是敬畏神的人了.
那麼神怎樣試驗亞伯拉罕呢?.
首先神就像一個朋友一樣.
叫亞伯拉罕.
然後亞伯拉罕就回應說.
我在這裡.
就好像以賽亞先知說我在這裡一樣.
他的回應表明了亞伯拉罕是願意信服.
願意接受差遣的一個態度.
然後神就說.
請你帶你的兒子.
就是你所愛的獨生子以撒.
去摩尼亞.
在我指示的山上.
將他獻為梵祭.
在這裡可能有人會認為.
和修本裡的你要帶你的兒子.
是一個命令.
但近代學者都認為.

$^{161}$這是一個請求的語氣.
所以我都譯為一個邀請.
他邀請說請你帶你的兒子.
去摩尼亞.
而梵祭.
就是將祭物殺死.
然後放在祭壇上.
完全焚燒.
代表了完完全全奉獻給神.
這個考驗給了你是亞伯拉罕.
你願不願意殺了.
和燒了自己的兒子.
我相信在座為人父母的.
可能會直接跟神說.
你不如燒了我算了.
何況以撒是亞伯拉罕.
等了幾十年.
是神應許的獨生子.
同時這裡要補充的是.
我們都知道神從來.
從來都不喜歡.
甚至恨惡人.
將自己的兒子或女兒.
獻做梵祭這件事.
恨惡到一個地步.
整本舊約聖經.
都重複提到.
使兒女驚火.
是神恨惡的事.
而夏文我們都見到.
其實神是會阻止亞伯拉罕這樣做.
所以一個如此有為神的信息.
和神如此珍惡的邀請或請求.
不如這樣我們問一個問題.
究竟神正在考驗亞伯拉罕什麼呢?.
可能我們可以在經文中找到一些端倪.
首先失去了以實瑪利的亞伯拉罕.
他究竟有多愛他的兒子以撒呢?.
愛到一個地步.
聖經第一次出現「愛」這個字.

$^{201}$就是《創世記》二十二章二節這段經文.
這裡透露了.
可能亞伯拉罕對以撒的愛.
已經蓋過了他對神的愛.
另一方面.
神一開始在第二節.
形容以撒是你的兒子.
就是你所愛的獨子以撒.
但是當神阻止亞伯拉罕.
下手傷害以撒之後.
就刪除了以撒是你所愛的這個稱呼.
這讓我們明白到.
神考驗亞伯拉罕是否敬畏神的內容.
就是邀請亞伯拉罕認清自己的真相.
你究竟重視我多些.
還是重視我賜給你的以撒多些.
究竟神賜下的禮物重要.
還是神本身自己重要.
所以我們看到.
雖然表面上神叫亞伯拉罕獻以撒.
實際上是要他獻上他的心.
以撒是他的命根.
是他的金波羅.
是他對神應許信心的基石.
但同時對亞伯拉罕來說.
他對以撒的愛.
可能已經蓋過了他對神敬畏的心.
所以神要考驗亞伯拉罕的是.
他所愛的他信靠的.
究竟是神本身.
還是神賜下的好處或應許.
弟兄姊妹.
在我們生活中.
是否容易出現這種愛神的賜福或好處.
多於愛神本身呢.
如果愛神的賜福多於愛神本身.
是一種以撒的話.
我們的生活會否久不久就出現呢.
好像神賜下我們這一代人.
有豐裕的生活.

$^{241}$這些是神的恩典.
手機由以前是奢侈品.
變成現在的必需品.
手機不單可以用來通訊.
亦可以用來娛樂.
但不知為何.
每天早上立志開個手機的聖經app.
或是看朋友傳出來的靈修小品的時候.
五分鐘不夠已經在看facebook ig.
十分鐘後連youtube都開了.
看盡天下的大事.
無怪乎現在的手機螢幕越出越大.
還要有兩個屏幕.
我們的生活的以撒都越來越大.
工作都可以成為另一種的以撒.
神賜下一份工作給自己.
可以發揮自己的能力.
收入亦有穩定的生活.
本來是一件美事.
但當我們越發揮自己所長.
我們就越投入工作.
同時慢慢將自己的價值觀.
安全感和工作捆綁在一起.
甚至可能覺得失去工作不行.
工作慢慢取代了神.
甚至事奉都可以成為我們心中的以撒.
服侍神本來是神的恩典.
我們可以透過服侍成長.
亦得到神賜給我們的滿足感.
但當我們不斷追求服侍帶來的滿足感.
我們開始緊張身邊有沒有人知道.
有沒有人賞識.
侍工慢慢取代了和神的關係.
慢慢我們的侍工果效.
成為我們自身的價值.
事奉的位置成為了我們的舞台.
其他人不可以碰.
我們的呼召 我們的恩賜.
慢慢變成我們的偶像.
因為侍工比神自己更重要.

$^{281}$阿帕拉罕愛以撒多過愛神.
同樣我們都可以愛子女勝過愛神.
當興趣班取代教會聚會.
當學業成績和課外活動取代了兒主阿灣那的時候.
可能我們都要問自己一條問題.
我們心裡的以撒是否浮現了出來.
弟兄姊妹 你們心裡的以撒又是甚麼呢.
事業 子女 還是苦心經營的成就.
當神說請你放下這些東西.
將焦點放在我身上的時候.
你的反應又會是甚麼呢.
面對一個好像不多合理的邀請.
阿帕拉罕在第三節作出了馴服的回應.
他做了一連六的動作.
清早就起來 預備了爐.
帶著兩個僕人和兒子 劈好柴.
起身往神指示他的地方去了.
他沒有像十八章那樣跟神討價還價.
反而在這六個連續的動作裡面.
我們看到阿帕拉罕好像沒有掙扎 沒有疑惑.
只有沉默.
還好像怕摩尼亞那裡找不到好的柴.
所以劈了柴 預備了爐 帶了兩個人去.
六個動作 一氣呵成.
讓我們明白到其實他已經下定決心要獻以殺.
雖然有學者認為這裡沒有描述阿帕拉罕的想法.
是因為要給讀者多一點空間.
去幻想當時他心裡有多麼矛盾.
有多麼悲痛欲絕.
但我個人認為.
其實他是醒悟到神考驗他的目的.
所以決定用信心和信服回應神的要求.
就算可能不太明白神的意思 仍然去做.
剛才說到神在考驗阿帕拉罕是否敬畏神.
而我們可以看到他帶著信心和信服的回應.
就是敬畏神的表現.
有什麼證據可以看到他的信心和信服呢.
我們有兩點可以看到.
第一 阿帕拉罕確信.
信實的神一定會落實他的應許.

$^{321}$不會反口.
在四至六節裡.
當他去摩利亞那兩天.
跟隨了主幾十年的他.
可能在回憶他和以撒的畫面.
以撒一出生的時候的喜悅.
可能在回想神的應許.
由一開始撒拉不辱.
到現在應許十二.
這幾十年來神的工作.
再想一想.
如果以撒死了.
神如何實現他的應許呢.
是否只是聽錯.
還是神反口要收回他的說話呢.
到了第三天.
阿帕拉罕舉目.
看到獻祭的地方.
他就和他的僕人說.
你們和隻驢留下來.
我和兒子去那邊敬拜神.
之後我們會回來你們那裡.
然後叫以撒揹著柴.
自己拿著火和刀.
二人同行一起上山.
值得留意的是.
和修本第五節.
阿帕拉罕說.
回到你們這裡來.
是沒有翻譯「我們」這個字.
近代的新漢語譯本和.
環球聖經譯本都翻譯出來.
而翻譯「我們」這個字的重要性.
除了因為原本有這樣的敘述之外.
更加重要的是.
反映了這兩天的路程裡.
讓阿帕拉罕更加確信.
神一定會信實地實現他的應許.
而以撒即是撒拉所生的.
是神實現應許那一位後裔.

$^{361}$所以以撒一定會回來那兩個僕人那裡.
因為二十一章.
神都和阿帕拉罕說.
因為從以撒生的.
才要稱為你的後裔.
當然當時的阿帕拉罕.
是沒有想過神會阻止他獻以撒.
但他總是覺得.
就算他獻了以撒.
神都會用什麼方法帶他回來.
落實神自己的應許.
所以阿帕拉罕才可以跟僕人說.
我們會回來你們那裡.
希伯來書的作者是這樣解讀這件事.
因著信阿帕拉罕被考驗的時候.
把以撒獻上.
他認為神甚至能使人從死人中復活.
意味著他得回了他的兒子.
阿帕拉罕有信心神會實現他的應許.
因為神甚至可以讓死人復活.
令他可以得回他的兒子以撒.
所以他可以和以撒二人同行一起上山.
而第二點我們看到他對神的信心.
就是因為他相信神的信實.
所以他相信神必然會看見和看顧他和兒子的需要.
所以用行動來回應.
證明他對神的信心.
其實以撒當時已經不是一個無知的小朋友.
起碼他知道獻祭的時候需要羔羊做祭品.
一般解經家都認為.
當時以撒最少十幾歲.
猶太的拉比甚至認為他已經三十七歲.
所以在第七第八節.
在上山的路途裡.
以撒見到爸爸手上只有火和刀.
覺得奇怪就問.
「爸爸,看看,火和柴都有.
但獻祭的羔羊在哪裡?」.
阿帕拉罕回答.
「兒子,神必定自己預備繁祭的羔羊」.

$^{401}$神必自己預備繁祭的羔羊.
我們可以視作一個信心的宣告.
這個宣告一方面表達阿帕拉罕.
深信神會預備獻繁祭的羔羊.
而另一方面也表達.
他深信神必然看見他和以撒的需要.
因為「預備」這個字可以翻譯為.
看見或看顧.
而「預備」這個字也貫穿了.
整段經文的重要字眼.
例如第十三節.
然後看見代替以撒的羔羊.
另外十四節.
耶和華以勒的以勒.
然後耶和華的山上必有預備的預備.
其實這四個字的語文都是同一個字.
阿帕拉罕相信整個獻祭的過程.
包括他清早起來準備一切.
然後和兒子現在上山.
其實神都看在眼裡.
神也會看顧他們的需要.
所以會為他們預備一切.
而且阿帕拉罕不只是宣告這麼簡單.
還用行動去證明他的信心.
九至十節.
當他和兒子上到摩利亞的山上.
他也用了一連六的動作去做回應.
捉壇.
把柴放好.
綁住兒子.
放在柴上.
伸手拿刀.
要殺他的兒子.
所以從經文我們可以看到.
一個敬畏神的人.
是會相信神是信實.
會實現他的應許.
同時也因為相信神是信實.
所以相信神在他的應許裡.
必然會看顧人的需要.

$^{441}$預備一切.
所以人可以用行動來證明他對神信實的信心.
來到這裡我們會問.
阿帕拉罕是信心之父才做到的.
我們這些凡夫俗子.
至少我是.
面對考驗的時候.
我們怎麼做.
但我相信.
阿帕拉罕並不是第一天已經成為信心之父.
因為就像一開始所說.
阿帕拉罕也有信心軟弱的時候.
兩次出賣老婆.
用自己的方法實現神的應許.
但這些失敗經歷.
沒有成為他的絆腳石.
反而成為了他確信神信實的資產.
同時在考驗裡.
神也給了阿帕拉罕時間去回憶這些資產.
另外後面的經文說到.
神必有預備.
在考驗裡.
神除了預備了代替以撒的羔羊.
還用了十二章至二十一章這麼多年的時間.
去預備阿帕拉罕這個人.
而今天.
同樣地.
神也在預備你和我.
問題只是我們怎樣將這些一點一滴神的工作.
轉化成為我們的資產.
我的想法是這樣的.
不要等考驗來到才尋求信心.
反而在平日的生活裡.
累積對神工作的回憶.
再讓這些回憶成為我們遇到考驗時候.
可以動用的資產.
而一個感恩的心.
在這裡尤其重要.
感恩除了讓我們知道神正在工作之外.
更加讓我們在困難的時候.

$^{481}$記得神是信實.
神仍然掌權.
當日阿帕拉罕要獨自面對神的考驗.
回憶神的工作.
感恩的是我們並不孤單.
因為神賜下了教會給我們.
信道群體的生活.
讓我們有機會分享神的工作.
彼此鼓勵.
聽到其他人的見證.
可以拉闊我們的視野.
提醒我們.
神不單在我自己身上有信實的工作.
也在歷世歷代的弟兄姊妹身上有這樣的工作.
幫助我們不會被自己的事情困死.
我參與的其中一個權契.
每次敬拜的時候.
帶敬拜的弟兄姊妹.
都會分享他們的近況.
而很多時候他們都會分享.
神在他們身上的工作.
甚至他們之前都有家庭見證的時間.
我很多次都因為弟兄姊妹.
他們分享到神的工作的時候.
得到鼓勵.
而我都深信那位預備分享的弟兄姊妹.
她的見證已經成為了她對神信心.
信實的一個資產.
所以鼓勵弟兄姊妹.
在分享代禱事項.
互相代禱的情況的時候.
定期都分享一件神在自己身上的工作.
藉此彼此祝福彼此鼓勵.
哥林多前書是這樣說的.
你們所受的考驗.
無非是人所承受得了的.
神是信實的.
他不會讓你們遭受到無法承受的考驗.
在考驗的時候.
總會給你們開一條出路.

$^{521}$讓你們能承受得了.
你們是包括你和我.
讓我們窩福家都可以一同成為.
彼此鼓勵互相祝福的群體.
回到經文.
在亞伯拉罕忍痛下手.
傷害自己的侄子之際.
十一至十四節.
神的使者突然在天上急急呼喚他.
亞伯拉罕 亞伯拉罕.
在這一刻.
亞伯拉罕都知道有些事情要發生.
所以立即回應.
我在這裡.
然後天使就阻止他說.
不要向這個孩子下手.
一點都不要傷害他.
現在我知道你是敬畏神的人了.
因為你沒有把你的兒子.
就是你的獨子.
留下不給我.
就像之前所說.
神考驗人並不是因為神知道人敬畏他.
是不是敬畏他.
而是讓當時的人知道.
什麼是敬畏神.
同時成為後人的榜樣.
而天使稱呼亞伯拉罕.
是敬畏神的人.
是因為他可以將他一生最重視.
最愛的獨生子.
都可以奉獻給神.
這個行為意味著.
眾言.
以撒是他等了幾十年的獨生子.
但是亞伯拉罕仍然將神擺在首位.
願意信服神的意思.
所以在十二字裡.
天使對以撒的稱呼都刪除了「你所愛的」.
即是保留你的兒子就是你的獨子.

$^{561}$因為亞伯拉罕都證明了.
他所愛的是神.
而他愛神也多於愛神賜給他的以撒.
十三,十四字.
雖然天使阻止了亞伯拉罕傷害以撒.
但是他的獻祭仍然要繼續.
就好像回應亞伯拉罕的信心宣告一樣.
當他抬頭一看.
就看見一隻公綿羊對角纏著一隻古木蟲.
亞伯拉罕當然馬上牽著那隻公綿羊過來.
代替以撒獻祭.
然後他把那個地方的名字改為耶和華爾.
直到摩西的時候.
還有人說在耶和華的山上必有預備.
新漢語的譯本.
在這裡有個備註.
「必有預備」直譯是「必看見」.
就好像之前所說.
「預備」這個字同時是看護或看見.
神必然透過他的預備.
看顧敬畏他的人.
而且人必然看見神的工作.
剛才敬拜的時候說到.
耶和華爾納神預備的東西是最好的.
為什麼耶和華爾納神預備的東西是最好的呢?.
因為神看見和看顧人的需要.
當亞伯拉罕起了名.
天使在十八章第二次講話.
神再次應許祝福亞伯拉罕.
和應許的原因是因為他所行的事.
沒有留下他的兒子.
十八節再補充主要的原因.
是因為他聽從了神的話.
如果我們對比這次應許的祝福.
和十二章的應許.
我們會發現.
因為亞伯拉罕這次的信服.
所以應許和祝福都是加倍了.
比如這次神說他親自起誓.
表明了應許有更加高的可靠性.

$^{601}$然後神說.
我必多多賜福給你.
我必使你的後裔大大增多.
再用天上的星.
海邊的沙.
作一個雙重的比喻.
顯示這個祝福有多大.
但我相信對亞伯拉罕最大的祝福.
是因為他的信服.
所以神也祝福了以撒.
如果我們看以撒的一生.
他也是一個很軟弱的人.
但面對一個這麼軟弱的人.
創世記二十六章.
神仍然在他面前顯現.
和賜福給他.
同時說到為什麼要賜福給他.
原因是亞伯拉罕聽從神的話.
遵守了神的吩咐.
亞伯拉罕信服帶來祝福.
祝福也帶到他下一代.
最後十九節.
亞伯拉罕回到他的僕人身邊.
然後一同起身回到別墅.
因為亞伯拉罕本來就住在那裡.
奇怪的是.
由第十節起.
亞伯拉罕伸手拿刀要殺他的兒子開始.
經文沒有再提以撒這個名字.
連這一節他們一起回去都沒有提到.
直到二十四章.
以撒要娶老婆這件事的時候.
以撒這個名字才再次出現.
可能因為矜持一役.
亞伯拉罕已經更加確定.
比起神賜下的好處和應許.
他更加愛神本身.
總括來說今天的經文我們看到.
神考驗亞伯拉罕是要讓他更加敬畏神.
亞伯拉罕之所以可以聽從神的話獻以撒.

$^{641}$是因為神一早透過幾十年的工作.
預備亞伯拉罕這個人.
同時也因為亞伯拉罕記得神是信實的神.
所以可以放下他的獨身子.
下手傷害他.
最終神阻止了這件事.
用公綿羊來代替以撒.
同時再次祝福亞伯拉罕.
甚至我們可以看到.
因為亞伯拉罕的信服.
軟弱的以撒因此都被神紀念和祝福.
也因為這樣.
亞伯拉罕被確認為敬畏神的人.
同時向我們顯示了.
記得神的信實是敬畏神的人的一個質素.
三十年前.
有一對弟兄和姐妹.
他們在外國讀完醫回來香港做醫生.
他們本來可以有高薪厚職的醫生事業.
有前途有安穩的生活.
按人來看這些都是神給他們很寶貴的恩典.
很寶貴的賜福.
但是沒多久神卻呼召他們去宣教.
當時對他們來說.
要放下香港的醫生職業.
不單止是要放下他舒適的生活.
甚至等同於放下神一直以來賜給他們的恩典.
同時要放下他們這麼多年努力讀書的一種身份.
但最終他們都決定放下一切.
去到南亞的工場.
由醫生成為宣教士.
其中弟兄很多年後在工場回來.
即是他們一起回來.
弟兄繼續為神工作.
最終成為第一位OMF華人總幹事.
幫助更多宣教士.
祝福不同國家的人.
我猜大家都知道他是誰.
他就是馮浩樓醫生.
在幾年前的一個見證分享裡.

$^{681}$馮醫生的太太鄭珍妮醫生都分享到.
她的女兒曾經在中六時問過她一個問題.
為什麼可以放下醫生的工作去宣教.
鄭醫生回應她說.
其實我一直都沒有放下過醫生的職業.
就算我現在都是一些宣教士的醫生.
同時其實當年在醫學院的訓練.
我現在都一直在用.
神的考驗不是為了收回給我們的東西.
而是為了我們可以成為敬畏神的人.
我們都向亞伯拉罕和馮醫生夫婦學習.
忠視神本身多過祂的賜福和應許.
成為敬畏神的人.
我們一起禱告吧.
透過亞伯拉罕獻耳撒.
我們看到什麼是敬畏神.
所以我們將我們每個人的心裡都交上.
我們的心都交上帶到你面前.
可能我們會重視我們的工作.
我們的恩賜甚至我們的人際關係.
求你賜下勇氣給我們.
將一切阻礙我們和你關係的事情一一消除.
又求主幫助弟兄姊妹.
讓我們不再只是找緊恩典.
更加要找緊這位賜恩典斯聯文的主.
讓我們都可以成為敬畏神的人.
願一切榮耀都歸給你.
我們的禱告是奉主明求.
阿們.
\newpage



\section{羅馬書 15:22-33}
\label{sec:031677}
\textbf{「無我」的事奉}~~~~秦民德傳道
\newline
\newline
連結: \href{http://www.church.com.hk/acms/content.asp?site=cdc\&op=show\&type=product\&code=031677\&layout=sermon}{\texttt{031677}} ~~~~ 語音日期: 2026-02-15
\newline
\newline
\hyperref[sec:031676]{\small{< < < PREV SERMON < < <}}
~
\hyperref[sec:index_chronic]{\small{[返順時目]}}
~
\hyperref[sec:031678]{\small{> > > NEXT SERMON > > >}}
\newline
\newline
羅馬書 15:22-33
\newline
\begin{longtable}{cl}
\hline
\hline
章節 & 經文 (和合本修訂版)\\
\hline
15:22 & \begin{tabularx}{0.7\textwidth}{X} 因此我多次被攔阻，不能到你們那裡去。 \end{tabularx} \\ \\ \relax
15:23 & \begin{tabularx}{0.7\textwidth}{X} 但如今，在這一帶再沒有可傳的地方，而且這許多年來，我迫切想去你們那裡， \end{tabularx} \\ \\ \relax
15:24 & \begin{tabularx}{0.7\textwidth}{X} 盼望到西班牙去的時候經過，得見你們，先與你們彼此交往，心裡稍得滿足，然後蒙你們為我送行。 \end{tabularx} \\ \\ \relax
15:25 & \begin{tabularx}{0.7\textwidth}{X} 但如今我要到耶路撒冷去，供應聖徒的需要。 \end{tabularx} \\ \\ \relax
15:26 & \begin{tabularx}{0.7\textwidth}{X} 因為馬其頓和亞該亞人樂意湊出一些捐款給耶路撒冷聖徒中的窮人。 \end{tabularx} \\ \\ \relax
15:27 & \begin{tabularx}{0.7\textwidth}{X} 這固然是他們樂意的，其實也算是所欠的債；因為外邦人既然分享了他們靈性上的好處，就當把肉體上的需用供給他們。 \end{tabularx} \\ \\ \relax
15:28 & \begin{tabularx}{0.7\textwidth}{X} 等我辦完了這事，把這筆捐款交付給他們，我就要路過你們那裡，到西班牙去。 \end{tabularx} \\ \\ \relax
15:29 & \begin{tabularx}{0.7\textwidth}{X} 我也知道去你們那裡的時候，我將帶著基督豐盛的恩典去。 \end{tabularx} \\ \\ \relax
15:30 & \begin{tabularx}{0.7\textwidth}{X} 弟兄們，我藉著我們的主耶穌基督，又藉著聖靈的愛，勸你們與我一同竭力為我祈求神， \end{tabularx} \\ \\ \relax
15:31 & \begin{tabularx}{0.7\textwidth}{X} 使我脫離在猶太不順從的人，也讓我在耶路撒冷的事奉可蒙聖徒悅納， \end{tabularx} \\ \\ \relax
15:32 & \begin{tabularx}{0.7\textwidth}{X} 並使我照著神的旨意歡歡喜喜地到你們那裡，與你們同得安息。 \end{tabularx} \\ \\ \relax
15:33 & \begin{tabularx}{0.7\textwidth}{X} 願賜平安的神與你們眾人同在。阿們！ \end{tabularx} \\ \\
[1ex]
\hline
\hline
\end{longtable}
$^{1}$各位弟兄姊妹,大家好.
今天是年廿八.
首先在此祝大家新年平安.
在我們思考上帝說話之前.
我們有一個禱告.
上帝,我們都需要在你面前.
去知取能力.
因為我們人有很多限制.
求上帝開我們的眼.
開我們的心.
向你敞開.
以致我們能夠.
看到上帝理想我們所做的.
帶領我們的腳步.
多謝你,我們祈禱奉耶穌的聖名.
阿們.
在上世紀三四六十年代.
是中國.
一個不容易的日子.
日本對中國發動全面的侵略.
抗戰結束之後.
亦都有內戰.
跟著就是一連串的動盪.
在中國最艱難的時候.
其實不少宣教士.
和中國的信徒.
都更加勇敢,剛強地.
站出來.
為主做見證.
在抗戰的時候.
有一個年青人.
回應了上帝的呼召.
他一生奉獻在全島工作.
他的名字叫邊雲波.
他當時是大學方重慶中央大學的學生.
雖然年紀輕輕.
但身體一直都不好.
因為十二歲那年.
日軍轟炸他的家鄉.
他曾經受傷吐血.

$^{41}$所以他一生.
都有嚴重的支氣官病.
他曾經有兩次機會出國讀書.
但他都拒絕離開戰亂的中國.
立志向苦難的同胞傳福音.
在1948年.
是內戰最激烈的時候.
邊雲波受上帝的感動.
未完成學位.
就毅然去到雲南,貴州這些貧瘠的山區.
向少數民族傳福音.
他被很多人質疑.
他很蠢.
但臨行之前.
他寫了一首長詩歌.
《獻給無名的傳道者》.
去表明他的心意.
這首詩也鼓勵了很多很多的信徒.
不知道在座當中有沒有人看過這首詩歌?.
有沒有?.
沒有?.
當中有幾句我相信值得和大家分享.
當黎明快要來臨的時候.
人世間便越顯得有厚.
艱難,幽暗.
但是.
李卓者向地說過.
是自己的手.
甘心放下世上的享受.
是自己的腳.
甘心到苦難的道路上來奔走.
選中這條不自由的道路.
並非出於無奈.
相反地.
他正是大膽地.
使用了自己的自由.
所以.
寧肯要淚水一河河地.
向內心湧流.
遙望著各國他的山頂.

$^{81}$就是至死也絕不退後.
1952年.
毓民波被雲南日報頭版點名批鬥.
被迫離開雲南.
1955年.
王明道在北京被捕.
作為王明道的同工.
毓民波亦被控反革命罪.
下放勞改.
32年之後.
到1987年.
他終於被平反獲得釋放.
當時已經是62歲.
但他仍然堅定地要傳福音.
而且比以前走得更加遠.
去到世界不同的地方.
繼續圍著奔走.
其實作為神的傳道人.
毓民波是將他一生.
身心奉獻作為活兒.
將自己完全交給神.
讓神帶他走神要他走的路.
在人看來,沒錯,他很蠢.
放棄了學位,前途.
放棄了在外國安全舒適的生活.
偏要選擇一條艱難的路.
為主的名,搞到身敗名裂,失去自由.
同樣是神的僕人,保羅.
其實他學上不是.
放棄了他自己優越的社會地位.
令人羨慕的前途.
他將整個生命交給主.
任由主去差遣.
保羅關注的是什麼呢?.
不是他自己的前路.
他能否有更好的生活.
他所關心的一生以來.
就是如何用他的恩賜和生命.
去見證基督,造就信徒.
沒錯,至死又絕不退後的編文波.

$^{121}$是無名的傳道者.
而表明活著就是基督.
死了就去益處.
一個為主的名,被自己同胞排擠.
最後被羅馬政權處決的保羅.
更加是一個無我的傳道人.
我們聽一聽經文.
我們會看到.
一個忠心僕人的生命是怎樣的.
我們今次會分三段去看這段經文.
我們首先看到.
做神的僕人首先要尊重神的心意.
讓神去帶領和事奉他的路.
第22到24節.
保羅和羅馬信徒說什麼呢?.
他說:因此我多次被攔阻.
不能到你們那裡去.
一落筆就說:因此.
很明顯,這是承接上文.
其實是第十九節.
第十九節說什麼呢?.
就是說:我從耶路撒冷直轉到以尼尼古.
到處傳入基督的福音.
這就是解釋為什麼他不能去到羅馬的原因.
這是地中海的地圖.
一個十字的地方就是保羅所身處的哥林多.
他寫羅馬書的時候就在哥林多寫的.
而在地圖的右下角.
會看到一條紅色的線.
是指示耶路撒冷.
而最上面紅色的線就是以尼尼古.
他從耶路撒冷去到以尼尼古一帶.
以尼尼古和哥林多之間.
就是馬其頓省和雅該雅省.
這些地方也是保羅福音足跡走遍的地方.
寫羅馬書的時候.
保羅已經事奉了二十二年.
他也五十多歲.
但仍然堅持傳福音給微德之民.
要完成耶穌基督給門徒的.

$^{161}$將福音傳到地極的使命.
當時他剛剛完成了在伊弗所的牧養.
然後輾轉到哥林多處理教會人士的問題.
也是在希臘半島.
在帖撒羅尼加和哥林多之間的一帶.
籌募一些捐款.
去接濟耶路撒冷貧窮的聖徒.
就算保羅很多次去羅馬探訪.
羅馬在哪裡呢?.
在意大利左邊.
大家可以看到.
紅點的地方就是羅馬.
他很多次都去那裡探望羅馬的信徒.
但因為種種的福祉他一直未能夠承行.
所以說我多次避難了.
不能夠去到這裡.
保羅說是避難了.
他是被動的.
事實上作為神的僕人.
保羅傳福音是聽命於上帝的大靈.
沒錯 祂有祂的計劃.
但很多時候聖靈都攔阻了祂的計劃.
在《攝行傳》十六章.
可能錯過了這一章.
記載到保羅被聖靈禁止在亞細亞港.
他打算去比推尼.
誰知耶穌的靈又不許.
神給他看到一個異象.
就是馬其頓人的呼聲.
結果他立即放下自己的計劃去馬其頓.
傳福音是耶穌給門徒的大使命.
保羅竟然看到亞細亞和比推尼的福音需要.
絕對是好事.
我相信保羅是不明白為何上帝阻止他.
但作為神的僕人.
總之收到指令就去做.
施奉其實是要突破自己的限制.
很多人其實都很怕計劃被打亂.
要臨時改動就會很煩躁.
但保羅不是按照自己的意願.

$^{201}$就算已經想好,做好準備.
他都要放下自己的大計.
去跟從聖靈的指示.
這就是神的僕人.
不過千萬不要以為.
服侍上帝不需要有計劃.
我們要有計劃.
但要將計劃交在上帝的手裡.
讓上帝按照他的心意去帶領,調整.
保羅說「但如今」.
在第23到25節.
大家剛才讀過的.
都出現了「但如今」.
「但如今」在這一帶沒有可存的地方.
甚麼意思呢?.
其實「但如今」就是說.
事情有不同的發展.
不是他的計劃.
這一帶沒有可存的地方.
是否在說.
這一帶所有人都相信耶穌.
不需要再存嗎?.
有沒有可能?.
不是啊.
而是保羅很清楚.
這一帶已經有信徒接觸傳福音.
他在希臘半島的福音工作.
是時候過一段路了.
保羅不只是滿足.
安頓在眼前的事工.
多年以來.
他一直向著一些.
未聽過福音的地方進發.
耶穌的大使命是甚麼呢?.
要門徒直到地極做他的見證.
而西班牙.
最左邊的.
叫做西班牙.
就是當時地中海最西的盡頭.
就是當時的地極.

$^{241}$所以保羅的目的.
就是要實踐最終的使命.
去到西班牙.
保羅多年以來都想去西班牙.
想去羅馬和西班牙.
但還是去不成.
可想而知.
他等了上帝開綠燈等了很久.
他去羅馬不是為了傳福音.
他是為了鞏固信徒.
和信徒相交.
他願意去那裡是甚麼呢?.
和他們彼此交往.
他希望和信徒有美好的團契.
讓他們成為保羅福音的夥伴.
讓這個地方成為一個宣教的基地.
保羅說甚麼呢?.
希望他去探望他們.
和他們有一個好的團契交往之後.
他們能為保羅送行.
甚麼是送行呢?.
送行不是禮貌上的道別.
say goodbye.
而是保羅知道.
雖然他自己去西班牙.
但福音是一個對公.
就算是最能打的屬靈強人.
報道家.
他需要甚麼呢?.
屬靈群體的同行和支援.
保羅去羅馬.
希望和羅馬的信徒建立一個互信.
親密的關係.
親密到羅馬的弟兄姊妹可以為他送行.
送行是甚麼呢?.
送行是指向物質和金錢的支援.
是一種實實際際的參與.
保羅希望和其他教會一樣.
羅馬的信徒也能和他建立一個福音的對公.
使他們能夠一起參與.

$^{281}$他再向西行的福音工作.
其實這是信徒合一的見證.
雖然保羅的心願多年以來都想做這件事.
但不是說去就去.
因為神有任務讓他完成.
在第二段落我們看到.
忠心事奉其實都要以教會合一為大前提.
第二十五節又出現了彈如金.
又有變數.
他要去耶路撒冷供應聖徒的需要.
為甚麼呢?.
因為馬其頓和瓦戈亞人已經捐了一些錢.
去給耶路撒冷的貧窮聖徒.
其實寫羅馬書的時候.
是在哥林多的港口.
由哥林多直接坐船向西行.
走著一些虛線的方向.
他就可以經羅馬去到西班牙.
但是上帝要他向東折返.
去耶路撒冷.
然後再調頭才去西邊.
就是西班牙.
不知為何亂了詞綴.
其實是不順路的.
有人計算過.
這樣一來一回.
保羅是要多走了二千里路.
為甚麼去耶路撒冷更重要呢?.
因為保羅籌到一筆錢.
是要將這筆錢送去耶路撒冷.
保羅提到欠債.
在第十五節的時候.
在第十六節提到欠債.
欠債是甚麼意思呢?.
其實保羅在羅馬書一開始的時候.
已經說了甚麼呢?.
無論希臘人,未開化的人.
所有的人.
他都欠了他們的債.
不知大家記不記得.

$^{321}$這是一種為福音.
負上的肅靈債務的心態.
其實中文翻譯的「債」字很有意思.
人字邊軸的「債」字.
是表示對虧欠的人.
所要負的責任.
保羅是用欠債的心態.
去對待主.
交給他的責任.
這種對人,對神的虧欠感.
成為了他事奉的驅動力.
向萬民傳福音.
就是他一生的義務.
所以他在哥林多前書說甚麼呢?.
若不傳福音我就有禍了.
保羅不單說自己欠了外邦人的債.
他還要外邦信徒知道.
原來他們也是欠了以色列弟兄的債.
固然他們是洛依.
其實他們也是欠債.
為甚麼說他們是欠債呢?.
如果大家記得十一章的時候.
保羅用日拿野橄欖給予外邦信徒.
其實神既然透過以色列人.
讓外邦人能夠分享福音.
外邦信徒是理所當然的.
要顧念耶路撒冷的弟兄的需要.
所以說.
因為外邦人既然分享了他們的靈性上的好處.
就當作是肉體上所需的供給他們.
就是這個意思.
雖然外邦信徒是有義務償還所欠的債.
在物質上支援.
但這不是出於欠債還債.
無拖無欠的心態.
而是要表達外邦信徒對耶路撒冷弟兄的愛,關懷和感恩.
但一直以來.
耶路撒冷教會的確有一部分人.
不願意接納外邦人做弟兄姊妹.
甚至抗拒他們.

$^{361}$有些猶太信徒要求外邦人.
一定要行割禮.
一定要接受猶太的傳統,宗教,文化.
歸化入猶太群體.
這樣才算是自己人.
這種以猶太為中心的心態.
其實成為教會合一的一個很大的難阻.
但保羅在真理上是堅定的.
他一直強調.
救恩不是建基於律法條文.
不是傳統的禮儀.
所有信主的人在主基督耶穌裡面都是一個身體.
如果這次耶路撒冷的信徒.
肯接納外邦信徒的捐獻.
這表示什麼呢?.
這是一個很重要的訊息.
即是耶路撒冷信徒願意打破隔膜.
彼此合一.
這是事關教會合一這麼重大的使命.
所以保羅去耶路撒冷有幾個很重要的目的.
首先他要紓解耶路撒冷貧窮信徒的生活困境.
也要幫助外邦信徒.
盡他們的本份.
但他們對耶路撒冷信徒的關心.
是向耶路撒冷的信徒主動伸出橄欖枝.
而保羅這個行動很重要的是打破兩個群體的隔膜.
促成大家的和解.
所以保羅說.
當我辦完這件事.
把捐款交給他們之後.
我就要路過你們那裡再到西班牙.
保羅不只是去.
保羅還說.
他們去是帶著豐盛的恩典去.
沒錯.
理想教會是不完美的.
無論在羅馬教會.
耶路撒冷教會.
都有很多很多自己的問題.
但是教會是神在救恩基礎上建立的.

$^{401}$神一定看著他的群體.
所以保羅很有信心.
我將帶著基督豐盛的恩典去.
在第三個段落.
我們看到一件事.
奉禱告其實是事奉最重要的力量.
三十到三十三節.
保羅說.
「弟兄,我藉著我們主耶穌基督.
又藉著聖靈的愛.
勸你們和我一起同心協力的為我祈禱」.
我們看到一個很謙卑的保羅.
很難想像.
那種心高氣傲和霸氣的保羅是什麼樣子.
保羅本來是一個極端仇視基督的法利賽人.
在大馬士革的路上.
他生命被耶穌徹底改變之後.
他的人生就以基督為中心.
所以我們到了保羅書信的時候.
我們時看到.
藉著我們主耶穌基督.
藉著聖靈.
在基督裡.
這些字眼是很常見的.
為什麼呢?.
因為保羅活著就是基督.
他是一個無我的僕人.
保羅其實大有學問.
他有影響力.
但是他知道一件事.
服侍從來不是靠自己的才幹.
他也看到自己是一個罪人.
更加是罪魁.
他更加不配做上帝的僕人.
他的服侍需要弟兄姊妹的支持和認同.
更加需要弟兄姊妹一同竭力的代禱.
「一同」這個字很有意思.
是由兩個希臘文所組成.
「同」和「戰爭」這個字.
保羅希望羅馬信徒.

$^{441}$和他一起為福音打仗.
所以第30到31節提出了幾件事.
是要求信徒「一同」為他祈求.
第一件事是什麼呢?.
使他脫離猶太不順從的人.
什麼意思呢?.
在耶路撒冷還有不少拒絕接受外邦信徒的死硬派.
他們這些人會成為這件好事的阻力.
所以保羅知道這是神的工作.
不是他的能力的工作.
第二件事是他這個服侍可以被聖徒越立.
其實保羅是沒有把握耶路撒冷的信徒.
一定會接受外邦人的捐獻.
要是神改變人心才行.
第三,他也要祈禱.
使我照著神的旨意歡歡喜喜到你們那裡.
和你們一起同得安息.
沒錯,保羅始終很希望能夠實現羅馬之行的計劃.
但他要祈禱.
不是一廂情願地要求上帝成就他的計劃.
他祈禱的是什麼呢?.
是要照著神的旨意去他們那裡.
如果這是神的心意.
他就能夠歡歡喜喜的去.
不是自己死裡的去.
保羅不是求神成就他想做的事.
而是求神讓他明白神要他做的事.
事實上是什麼呢?.
什麼是祈禱呢?.
祈禱不是教神做事.
祈禱是在祈禱裡面.
聽神怎麼說,想我怎麼做.
在三十到三十二和三十二節.
出現了兩次一同的意思.
保羅不僅是期望羅馬信號一同增賤.
也期望同得安息.
這個同字很重要.
保羅很重視教會的同心.
他不覺得自己是超級信徒.
他不覺得自己高人一等.

$^{481}$絕對沒有所謂平信徒,仕徒的分別.
他深信到羅馬和信徒互相勸勉的時候.
他就能夠同得安息,同得力量.
最後祝福羅馬信徒的是什麼呢?.
願賜平安的上帝和你們眾人同在.
平安 Shalom 這個字.
以色列人見面時時常說.
這是什麼意思呢?.
不是說新年所有人出入平安那種平安.
是說願你和上帝有一個相和諧.
相通和好的關係的意思.
這正正就是保羅寫書信給不同教會.
要幫助教會和神建立和諧關係的目的.
我們看回整段經文.
其實事奉有什麼重要原則呢?.
其實說白一點.
不僅是事奉.
我們對神,對人有什麼原則呢?.
第一個是什麼呢?.
無我的事奉是以神為中心.
不是滿足自己的心願.
事奉神不是一廂情願.
自己想怎樣去做就怎樣去做.
其實每個人都有自己的想法.
保羅也是有自己的計劃.
他想去特羅亞,比推尼.
但又想去羅馬.
每次都被聖靈帶去另外一個地方.
他聽神的命令事奉.
不過聽命事奉的先要條件是什麼呢?.
就是要和上帝有一個密切的關係.
你才明白上帝的心意.
如果我們和上帝沒有一個密切的關係.
和上帝疏離.
和上帝沒有溝通.
我們根本不可能知道上帝的心意.
在這樣的情況下我們去事奉.
可能只是自己做了一份義工.
在事奉的路上.
保羅給了我們一個很重要的原則.

$^{521}$就是聽神的指示.
按照神的步驟,方法,時間去做.
而不是我們覺得應該是這樣.
就應該是這樣.
上帝不需要靠我們做事.
而是我們需要靠上帝的恩典.
去參與上帝的聖功.
在事奉的路途上有很多的陷阱.
我自己也曾經失敗過.
不止一次,是很多次的.
我記得其中一次最深刻的印象.
就是我剛畢業的時候.
加入了基督徒的聖經群組.
有一年我們去到新加坡,馬來西亞.
巡迴報道,開讚美會.
特別是在新加坡的一次.
一個會場大約有千多二千人.
很開心.
大家唱詩歌.
一個二十多歲的初出茅廬.
自以為是的我.
在台上說話,大唱詩歌.
還獨奏.
自己選擇最適合自己唱的,最熟悉的歌.
唱完之後.
下面有些人留下來跟你說.
你真的唱得好.
真的以為真的好.
自己飄飄然.
回到香港的飛機上.
還在回味這次行程的成功,好享受.
但突然間心裡有一個完全相反的意念彈出來.
我是在事奉嗎?.
我是不是在威脅自己呢?.
我知道這個意念不是我自己有的.
我沒可能懂得想這些東西.
今天回想我自己當年是很幼稚的.
這個意念亦都問我.
你跟神的關係如何?.
你來到這個旅程.

$^{561}$這麼多天奔奔波波.
好像很勤奮.
有沒有靜下來祈禱靈修?.
什麼事奉來的?.
很享受.
事奉是來自己享受.
我事奉的是誰呢?.
一額汗.
很震撼.
不過我很感恩.
因為上帝肯管教我.
肯管教我即是不會丟下我.
會再用我.
在人生事奉的路途上.
其實有很多次跌跌碰碰.
人是有很多很多的挑釁.
所以我們真的要很謙虛.
在上帝面前.
看自己是什麼的傢伙.
第二個原則就是.
我們不是在服侍事奉.
我們是在服侍人.
本來眼裡不只是有事奉.
他著重的是人.
他本來對基督徒充滿仇恨.
他當年對基督徒殺氣騰騰.
要抓要鎖.
喊打喊殺.
但他成為信徒之後.
他在書信裡.
他流露到信徒那種很深很深的愛.
很溫暖的情.
他不僅看重信徒的成長.
耶路撒冷信徒那種的生活.
那麼傑具的.
他也很賞心.
為他們奔波籌款.
對於未得救的同胞.
如果大家記得的時候.
在《羅馬書》第九章.

$^{601}$他甚至說什麼呢?.
「我是骨肉之親」.
「我就是自己被就坐.
與基督分離」.
「我為了他們我也願意」.
保羅在神裡面.
從一個至高至大的人.
成為一個完全透明無我的人.
這是源於他將自己的真心獻上.
當作活祭.
這句話是他寫的.
他不是指我們要做這件事.
是他自己先做了.
是他自己心意被更新而變化.
所以今天我們有時候.
會否也想一想.
我們是否真的很熱心.
我們只是很熱心我們的事工.
為了做好一件事.
做成一件事.
達到我們預期的目的.
我們變得很所謂的Task-Oriented.
與弟兄姊妹相交.
全部是工作模式.
會不會呢?.
或者有沒有試過.
為了堅持自己認為好的做法.
令到群體撕裂,爭拗.
甚至造成傷害?.
我們的事奉.
會不會變成事奉一個Project?.
事奉一件事?.
或者.
好像當年那樣.
用事奉去滿足自己的好勝.
而不是在服侍上帝.
最慘就是在欺騙自己.
神所要的是謙卑,信服的僕人.
而不是要性格巨星.
第三個提醒就是.

$^{641}$重視團隊的事奉.
堅持同心合一.
其實保羅大有學問,大有能力.
但他從來不做個人Show.
他是邀請弟兄姊妹和他一起.
他需要討個夥伴.
他需要人送行.
信徒之間.
同心事奉有什麼條件?.
我覺得最重要的是好好溝通.
其實保羅不是羅馬教會的差派.
他在工作上沒有責任向他們匯報.
但你看看保羅.
一直以來.
他跟他們說.
這區已經沒有什麼可傳的地方.
我現在要去哪一邊?.
接著我想怎樣怎樣.
去西班牙.
又想來探望你們.
他跟將要接觸的弟兄姊妹.
講得清清楚楚.
不會有任何誤會.
這就是坦誠的溝通.
很重要.
做好對方的期望管理.
也要讓對方知道自己的期望.
這個我相信在事奉裡.
很多時候都會忽略的.
來到羅馬書的結尾.
書信結束就是問安的說話.
保羅剛才不是問了安嗎?.
但當我們下一次看十六章的時候.
你會看到保羅的問安.
不是公式化的圖說話.
十六章整章都是問安的.
保羅用了二十七節的篇幅.
除了問候不同群體.
還逐一點名三十五個人的名字.
和對他們的記掛.

$^{681}$長長的名單裡.
每一個跟他一起事奉的人.
他都非常賞心.
可以看到保羅是多麼重視.
跟弟兄姊妹之間的情誼.
多麼重視教會的合約.
今天我們都是跟保羅一樣.
信耶穌.
但生命有沒有像保羅一樣.
願意被更新,變化.
這是關鍵.
但願保羅對人,對神的態度.
都成為我們的榜樣和提醒.
好嗎?.
這是我們的禱告.
基督不單止在世上活過.
成為我們的榜樣.
上帝還興起您的僕人.
成為我們的榜樣.
以至我們周圍有很多見證人.
如同雲彩圍繞著我們.
我們不可以抵賴.
我們看不到.
我們不可以說我們不懂得怎麼做.
求上帝帶領我們.
讓我們在謙卑的生命.
能夠呈獻,服侍,拯救我們.
創造我們.
和我們一起走.
耶穌基督.
三一的真神.
我們祈禱奉耶穌的聖名.
阿們.
\newpage



\section{創世記 17:1}
\label{sec:031678}
\textbf{聖經人聖經事:作完全的人}~~~~黃克勤牧師博士
\newline
\newline
連結: \href{http://www.church.com.hk/acms/content.asp?site=cdc\&op=show\&type=product\&code=031678\&layout=sermon}{\texttt{031678}} ~~~~ 語音日期: 2026-02-09
\newline
\newline
\hyperref[sec:031677]{\small{< < < PREV SERMON < < <}}
~
\hyperref[sec:index_chronic]{\small{[返順時目]}}
~
\hyperref[sec:031679]{\small{> > > NEXT SERMON > > >}}
\newline
\newline
創世記 17:1
\newline
\begin{longtable}{cl}
\hline
\hline
章節 & 經文 (和合本修訂版)\\
\hline
17:1 & \begin{tabularx}{0.7\textwidth}{X} 亞伯蘭九十九歲時，耶和華向他顯現，對他說：「我是全能的神。你當在我面前行走，作完全的人， \end{tabularx} \\ \\
[1ex]
\hline
\hline
\end{longtable}
$^{1}$各位弟兄姊妹,祝你平安喜樂.
又過了一天了,又一起讀聖經.
聖經人聖經事由生命樹宣教網絡提供.
希望我們多讀聖經,多認識上帝.
今次是作玄傳人.
如果昨天你還記得,就是講墨基使德.
墨基使德就是亞伯拉罕打勝仗回來之後.
要拿自己所獲的十分之一.
獻給至高神的祭司墨基使德.
因此墨基使德再次在希伯來斯出現.
被認為是上帝.
最高的祭司.
代表.
早些.
表達象徵耶穌基督.
因此墨基使德是一個平安的王.
亞伯拉罕也獻十分之一給他.
即是說墨基使德是一個很高的位置.
今次就看作玄傳人,就是上帝要求亞伯拉罕.
作玄傳人,今次聖經人聖經事.
來到講亞伯拉罕,亞伯蘭,你知道要講很久了.
蒼生基第十七章第一節,留意一下.
你用一個故事的方式去想,你會容易記得.
當時亞伯蘭九十九歲的時候.
耶和華就向他顯現.
顯現.
如何顯現?馬上出現還是在夢中?.
這個顯現就不像在夢中,可能是見到異象.
對他說.
我是全能的神.
你當在我面前作玄傳人.
九十九歲的時候,耶和華向亞伯蘭顯現.
對他說.
我是全能的神.
你當在我面前作玄傳人.
我就與你納若,使你的後裔極其繁多.
我就與你納若.
使你的後裔極其繁多,綠色的詞就是解釋注釋來的.
我是全能的神,我不特別解釋了.
我不是大能的神,是全能的神.

$^{41}$即是說.
在我面前講.
沒有一件事我做不到.
很難啊,沒有的.
當然有很多論證,以今天上帝能不能夠做一塊石,他自己舉不起.
這些.
我們不回答這種問題,因為這樣.
是那個.
問題上面本身是不邏輯.
我們不回答這些不邏輯的問題.
好了.
我是全能的神.
你當在我面前作玄傳人.
這個「你」字藍色括著.
就是.
希伯來文原文是沒有的.
我是全能的神.
當在我面前.
玄傳.
應該是這樣的,人字是沒有的.
這些字加上去.
是.
翻譯者想幫我們更加容易明白.
當.
halak.
在,但它是命令式.
於是.
命令式就會變成.
當,中文的.
文法沒有命令式,加上要.
應該,應當.
當.
這個當是表示halak.
中文翻譯是一個命令式.
命令哪一個?.
只有一個,就是神和阿伯藍顯現,當在我面前作玄傳人.
於是加了「你」這個字.
其實我是全能神.
當在我面前.
玄傳.

$^{81}$這個「作」字.
其實.
haya,haya是成為.
「你當在我面前成為一個玄傳」.
「你當在我面前成為玄傳」.
這個玄傳.
tamin.
是沒有瑕疵.
正直的意思.
正直和沒有瑕疵有沒有關係呢?.
是有的.
沒有瑕疵即是說人完全正直,他說的話你就知道是這樣的意思.
沒有花架,你不用去想.
好像選擇一樣東西,有沒有花呢?.
其實他是沒有花的.
沒有瑕疵的.
即是說,他說的話你可以相信.
知道是真的.
其實就是正直.
你當在我面前正直.
這個意思.
正直就是不說假話,說直接的說話.
正直一點.
委婉一點可以嗎?.
BB靚不靚呀?樣子很難看,鼻子很扁,眼睛很小.
死啦,這個女兒.
將來怎樣嫁人呀?.
當人心這樣想.
正直就是哪裡靚呢?你女兒眼睛很小,鼻子很扁,將來很難看.
口面竟然那麼方形.
正直.
但是.
可不可以.
不說呢?.
可以的.
不違反正直的,例如,我女兒靚不靚呀?抱著這個BB.
哎呀,她的樣子很可愛.
睡得很舒服.
很可愛.
可愛說不出,那不要說了.

$^{121}$她睡得很軟.
睡覺BB,個個都好看.
OK,都不影響正直.
起碼不會說謊,眼睛那麼小,嘩,眼睛真是大呀,鼻子那麼高.
這樣不行.
正直意思.
是說.
對的說話,不會騙人.
這一句.
第一,很重要的一點就是.
要在要行走.
行走我面前成為正直.
我是全能的神.
行在我面前.
正直地行在我面前.
我是全能的神.
正直.
成為正直地在我面前行走.
來到神面前.
正正經經.
正正直直.
不要以為你騙到神,不要以為祂不知道.
正直.
這樣.
我就,這個「就」字.
應該.
是.
表示「你這樣我就這樣」,這是現代中文的用法.
「我就與你立約」.
其實是「我將會與你立約」.
Nathan,不要讀英文,英文會變成彌敦道.
Nathan.
我就與你立約.
「與你」這個意思是.
將來式來的,所以這個「就」的意思.
是將來式的.
I will make.
Nathan是I will make.
你在我面前正直成為正直地行走.
我將會與你立約.

$^{161}$意思是甚麼呢?第二個我想說的就是.
如果你不正直.
神是不會跟你立約.
神是與正直的人立約.
這個「約」是表示關係與對他的.
愛,對他的付出.
這個.
是「約」的意思.
會使你的後裔極其煩多,我的「約」就是使你的後裔極其煩多.
在神面前.
成為正直地行走.
我將會與你立約.
使你的後裔極其煩多.
在這裡我都先解釋一下這個「約」字.
立約.
很多時候,譬如我們誤解了.
我們很多時候是誤解了.
聖經裡面說的「約」.
甚至很多傳人都說「約」,就是神與我們立約.
其實就是神這樣對我們,要求我們這樣對他.
否則這個「約」就沒有了.
結婚.
兩個人,一男一女,大家立約.
很多時候,「約」就是一個合約,多了一個字在前面.
合約.
多了一個字在前面,盟約,不是.
其實神與我們立約只有一個字,「約」.
confident.
不是大家有一個agreement,有一個contract,不是那件事.
今天我和別人立約,買賣合約,我付多少錢.
他就給我一個房子,買賣合約.
我要付一個按金.
然後.
什麼時候付尾數.
完成了我要履行的部分,於是房子就要給我.
還有政府監管著.
如果我沒有付尾數,房子不會給我的.
或者他不能夠不給我的,後來反價,你多付20萬.
我才賣給你,尾數要加,不行的.
因為是大家的.

$^{201}$合約.
明白嗎.
但聖經的「約」不是這件事.
我要解釋一下這個「約」.
「約」是單方面的.
我和你立約.
神與亞伯拉罕立約.
納什麼約呢?我和亞伯拉罕立約,就是我要使後約極其繁多.
後面那些約,前面那些約,就是我要把這塊地給你.
你見到這些.
就是立約,我要令你的子孫多起來.
立約.
你一定要聽我的誡命,否則我就.
不給你這些了.
沒有的.
好像西乃山那樣.
十戒和其他的戒命.
如果你不守,我就追討你的罪,直到三四代.
守約思慈愛,如果你們守我戒命.
我就守約思慈愛.
對你的後約祝福,直到千代.
聽起來就好像.
如果他有這些給我,我要做我那份.
就是要守約,不守約就是我.
失落了我那部份,於是他也不守他的約.
不是,解錯了.
只不過是.
沒有接受到神,沒有跟隨到神,那樣他.
他會懲罰.
好像爸爸懲罰,但是.
那個約就是神已經納了,我要永遠地祝福我的子民.
他永遠是我的子民,雖然我現在鞭打他痛.
痛好幾代,不過.
他仍然是我的子民.
這個立約沒有改.
這個是單方面的,譬如好像現在很多人誤解了.
各位弟兄姊妹,現在這一男一女在主面前就有一個合約.
給錢我對方.
那於是隱喻了什麼呢.
如果一日有一方不愛對方.

$^{241}$沒有跟著這些事做,那另一方就可以都不愛他了.
因為他毀了約,所以另一方也可以毀約.
其實不是這個意思.
兩個人.
都自己在神面前向對方立約.
無論疾病健康,快樂不快樂,怎樣怎樣.
我都願意愛他.
那你是否願意?我願意,這是他對他的立約.
然後你是否怎樣,我又疾病健康,我都愛他,跟他一起,這是你對他的立約.
是兩個都同時立約,但那個並不是合約.
那個是約.
即是說對方不再愛我.
但是我仍然愛他.
因為我對他立約.
他對他立約,他對他立約.
那有一方不再愛了,在約上失去了,那這個人怎樣,我都好了,那我不用愛他了.
不是,我仍然愛他,這就是立約的精神,不是合約精神.
是立約精神,是信心和合約的分別.
我就將會跟你立約.
你呢.
不過第一步,我要求的是什麼.
我想跟你立約,不過我都要求你.
是自己在我面前.
去行走.
成為一個正直的人.
我就.
跟你立約了.
就令到你後裔多.
亞伯蘭就俯伏在地有句神說.
我與你立約.
神就跟他說,我與你立約,你要作多國之父,從此以後你的名字不要叫亞伯蘭,要叫.
亞伯拉罕.
我以納你作多國之父.
亞伯蘭俯伏在地,那是什麼意思,多謝,多謝你.
然後神就說,這個約是什麼呢.
就是你要做多國之父.
所以你要改名.
亞伯蘭呢.
亞伯蘭這個字是高舉的父.
高舉的父即是很厲害,很偉大的父親.

$^{281}$亞伯蘭.
但你要改個名.
是要做亞伯拉罕.
亞伯蘭.
亞伯蘭是多人的父.
由偉大的父親變成多人的父親.
亞伯蘭.
我以納你作多國之父,這就是第三樣說法,大家看看這三樣東西.
第一樣.
你要先做一個完全的人.
然後神就說,我將與你立約,我要這樣給你.
我想放一些東西在托盤上.
我都會先要求托盤是乾淨的.
明白放什麼進去?.
就是我要放進去,讓你由普通一個好的父親變成一個.
多國的父親.
這就是.
這段聖經所說的.
你的後裔極其繁多,國度從你而出,君王從你而立.
我要與你並你世世代代後裔,建立我的約.
你的世世代代後裔可能很差,不過我仍然建立我的約,做永遠的約.
是你後裔的神.
今天的金句就是.
創世記十七章第一節.
九十九歲向顯現的不要緊.
我是全能神,一起讀一下.
我是全能神.
你當在我面前.
作完全人.
我是全能的神.
你當在我面前.
作完全的人.
神在這裡要求我們.
行走得下.
行走.
成為一個.
正直的人.
這就是上帝對我們的要求.
神拿到一個很乾淨的器皿.
一個很乾淨的瓶.

$^{321}$他才可以把自己的東西倒進去.
就是這個意思,你要明白.
先做一個.
完全的人.
這個完全的人.
不要解釋為perfect,你要perfect,不是,他是說沒有瑕疵正直.
是一個.
是就說是,不是就說不是,好像主耶穌這樣說.
正直的人.
你走過去,不怕他後面暗箭射你,他不會陰你.
你和他一起,你知道他說的話是真的.
歡喜不歡喜,是他說出來的.
對與不對,是他說出來的.
神是喜歡這樣的人.
正直人.
不過很多時候我們是喜歡那些滑頭的人.
喜歡甜言蜜語.
自己都知道是假的,不過都相信.
其實是很傻的,神說我不喜歡的.
我喜歡人做正直的人.
正直的人我才和他有關係的.
一起祈禱.
主我們在你面前立志要做一個正直的人.
雖然我們並不是被要求完全.
無可指責,什麼錯都沒有,其實你都沒有這樣的要求.
你只是要求我們完全的意思是正直.
在你面前成為一個正直的人.
你就歡喜我們.
你就使用我們.
和我們納若.
叫一個就算是厲害還是不厲害的人,可以成為很多人的祝福多國的父.
亞伯拉罕這個名字就是這樣出來的.
主我們都在你面前願意成為亞伯拉罕.
不是自己夠好.
而是你.
教導潔淨我們.
讓我們可以成為多人的祝福.
隨聽禱告,奉耶穌名禱,阿們..
\newpage



\section{使徒行傳 11:19-30}
\label{sec:031679}
\textbf{落在異鄉,不為異客}~~~~駱瑞麟傳道
\newline
\newline
連結: \href{http://www.church.com.hk/acms/content.asp?site=cdc\&op=show\&type=product\&code=031679\&layout=sermon}{\texttt{031679}} ~~~~ 語音日期: 2026-02-22
\newline
\newline
\hyperref[sec:031678]{\small{< < < PREV SERMON < < <}}
~
\hyperref[sec:index_chronic]{\small{[返順時目]}}
~
\hyperref[sec:031680]{\small{> > > NEXT SERMON > > >}}
\newline
\newline
使徒行傳 11:19-30
\newline
\begin{longtable}{cl}
\hline
\hline
章節 & 經文 (和合本修訂版)\\
\hline
11:19 & \begin{tabularx}{0.7\textwidth}{X} 那些因司提反的事遭患難而四處分散的門徒，直走到腓尼基、塞浦路斯和安提阿。他們不向別人講道，只向猶太人講。 \end{tabularx} \\ \\ \relax
11:20 & \begin{tabularx}{0.7\textwidth}{X} 但內中有塞浦路斯和古利奈人，他們到了安提阿也向希臘人傳講主耶穌的福音。 \end{tabularx} \\ \\ \relax
11:21 & \begin{tabularx}{0.7\textwidth}{X} 主的手與他們同在，信而歸主的人數很多。 \end{tabularx} \\ \\ \relax
11:22 & \begin{tabularx}{0.7\textwidth}{X} 這風聲傳到耶路撒冷教會的人耳中，他們就打發巴拿巴到安提阿去。 \end{tabularx} \\ \\ \relax
11:23 & \begin{tabularx}{0.7\textwidth}{X} 他到了那裡，看見神所賜的恩就歡喜，勸勉眾人要立定心志，恆久靠主。 \end{tabularx} \\ \\ \relax
11:24 & \begin{tabularx}{0.7\textwidth}{X} 這巴拿巴原是個好人，滿有聖靈和信心，於是有許多人歸服了主。 \end{tabularx} \\ \\ \relax
11:25 & \begin{tabularx}{0.7\textwidth}{X} 他又往大數去找掃羅， \end{tabularx} \\ \\ \relax
11:26 & \begin{tabularx}{0.7\textwidth}{X} 找著了，就帶他到安提阿去。他們足有一年和教會一同聚集，教導了許多人。門徒稱為「基督徒」是從安提阿開始的。 \end{tabularx} \\ \\ \relax
11:27 & \begin{tabularx}{0.7\textwidth}{X} 當那些日子，有幾位先知從耶路撒冷下到安提阿。 \end{tabularx} \\ \\ \relax
11:28 & \begin{tabularx}{0.7\textwidth}{X} 內中有一位，名叫亞迦布，站起來，藉著聖靈指示普天下將有大饑荒；這事在克勞第年間果然實現了。 \end{tabularx} \\ \\ \relax
11:29 & \begin{tabularx}{0.7\textwidth}{X} 於是門徒決定，照各人的力量捐錢，送去供給住在猶太的弟兄。 \end{tabularx} \\ \\ \relax
11:30 & \begin{tabularx}{0.7\textwidth}{X} 他們就這樣做了，託巴拿巴和掃羅的手送到眾長老那裡。 \end{tabularx} \\ \\
[1ex]
\hline
\hline
\end{longtable}
$^{1}$各位兄弟姐妹,平安!.
很快又過了一年了.
今天是年初六.
也是祝大話.
大家以馬來黎.
馬年.
神與大家同在.
我們開始的時候,我們先來一個禱告.
最後,謝謝你.
因為你.
每個星期都給我們有你的懷願.
求聖靈.
這個時候就跟我們講說話.
來教導我們.
特別是.
將你.
將我們帶來澳洲.
一定是有你們的心意.
求你,不單止是.
讓我們能夠看見.
也能夠跟著去.
禱告,奉主耶穌基督的聖名.
阿們.
唐代詩人.
有一句很出名的.
說話,我想大家都很熟悉.
就是什麼.
獨在異鄉.
為異客.
每逢佳節.
陪思親.
短短這一句說話.
其實就是講出人.
在這個陌生的地方的時候.
那種孤單.
和疏離的感覺.
特別是在過年.
過節的時候.
不過當我讀這句說話的時候.
我有一個.

$^{41}$另一個感受.
有一句問號.
是不是真的樂在異鄉.
就是一定要成為.
異客的呢.
大家覺得怎樣看.
對很多.
移居海外的信徒來說.
環境陌生.
文化又很不同.
連這個信仰的生活.
都可能要重新適應.
但那種沒有歸屬感.
異客的感覺.
其實一點都不會陌生.
因為在《仕途行傳》那裡.
就已經記載了.
當時的信徒.
跟著猶太教的迫伯.
他們需要離開.
原來的地方.
到處去.
而這件事也記載在.
《仕途行傳》.
第七到第八章那裡.
當這個屍體犯.
殉道的時候.
就已經開始了.
信徒因為很怕.
那些所謂.
宗教狂熱分子.
正如那些保羅之類的人.
來執行屍刑.
所以被迫.
要離開耶路撒冷.
去到其他城市.
而其中一個很重要的地方.
就是安提阿.
而安提阿.
當時是.

$^{81}$羅馬帝國的.
第三大城市.
在現在的.
土耳其.
鄰近地中海那裡.
可以說是.
交通很便利.
商貿也很發達.
所以就匯集了.
各方人士.
都在那裡居住.
可以說是.
文化多元.
宗教氣氛也非常濃厚.
但是非常複雜.
在這個環境底下.
對於那些新的移民來說.
一定在適應上.
有很多張力.
從我們剛才.
讀的經文那裡.
你會看得到.
當一群被打散的信徒.
是怎樣一步一步.
走進這個新的城市.
走進這個新的人群裡.
是吧.
我會將這段經文分為四個部分.
來跟大家看看.
當時.
安提柯的信徒.
怎樣被神帶領.
找到他們自己的定位.
首先來看看.
第19節21節那裡.
經文第19節那裡.
將我們帶回去.
當時遭遇患難的情景.
我們看到.
這個屍體犯.

$^{121}$他被迫害.
被殺.
這樣.
信徒離開耶路撒冷.
去到.
菲利基,居比路和安提柯那裡.
這種經歷.
對當時的信徒來說.
真的很難覺得.
是神的祝福.
更加像失敗.
挫折.
甚至因為信仰.
他需要付出的代價.
是吧.
而四散這個詞.
有很大的.
不確定性.
因為信徒不是有計劃.
這樣離開的.
當然是非常匆忙.
帶著恐懼的心情出走.
他們在耶路撒冷.
根本上.
沒辦法再立足.
因為他們連去會堂.
也不可以.
那些人看到他們就會迫害.
但路加柯是從神的那裡.
來描寫這段歷史.
他暗示.
原來分散.
可以將福音.
突破原來的界線.
所以在經文裡.
特別指出.
那些人.
起初指向.
猶太人來到.
是港獨.

$^{161}$顯示了什麼.
顯示在潛意識裡.
其實他以為救恩.
是猶太人的專利.
而外邦人.
根本是不配.
或者不可能.
認識彌賽亞.
但二十節裡出現了一個轉折.
那裡是說什麼.
他說.
有幾個人.
就向希尼尼人.
傳講耶穌.
路加更加用了.
兩個字.
就是居比路.
希尼尼來人.
其實他是指.
那些已經相信了耶穌.
會說希臘話.
是僑居海外的猶太人.
而希尼尼人.
如果根據上下文.
來看應該是指.
希臘人.
而不是指當時那些.
所謂希臘化的猶太人.
現在我們新一點的譯本.
都已經寫了叫希臘人.
所以在這個.
這麼大集會的城市裡面.
如果我們在.
宗教的層面來看.
其實可以分為兩大類人.
一些就是猶太人.
而一些就叫.
所謂外邦人.
猶太人裡面.
當然有些會說.

$^{201}$希伯來話,亞蘭文.
和希臘話.
更加是裡面有信猶太教的.
有不信猶太教的.
現在加上.
有些信耶穌.
現在外邦的人群.
不用說.
一定會說希臘話.
或者加上自己本國的語言.
有些是.
信猶太教.
還有.
還沒信的.
後來當然有些信耶穌.
所以很複雜.
你看到.
這個言語文化.
和信仰.
交集的城市裡面.
對於新移入的.
移民.
居然能夠不斷地.
傳講耶穌.
以致福音.
就是這樣傳開.
所以路加.
在第21節.
他其實講了一個真正.
為什麼福音可以傳開的原因.
是什麼.
原來主與他們同在.
這個說法.
其實是很重要的.
在這裡告訴我們.
推動宣教.
向人傳福音.
建立教會.
不是因為環境是否穩定.
不是因為.

$^{241}$他們沒有苦難.
而是主與他們同在.
所以你看到.
其實分散沒有阻礙神的工作.
反而成為.
神旨意的展開方式.
看上來.
好像很失敗的歷史.
但在神的手裡.
卻成為救恩的管道.
所以.
特別在這裡提醒我們.
今天的進度.
這一段經文.
其實在邀請我們重新思考一個問題.
是什麼問題.
是當人.
我們的生活.
被打亂的時候.
當我們人生的需要轉向的時候.
我們是否只看到.
我們表面上的.
失敗呢.
還是看到神仍然掌權.
仍然在他背後.
做公呢.
分散.
並不是表示.
被遺棄.
離開熟悉的地方.
不是表示離開了.
神的心意.
有時候神的旨意.
慢慢地.
向我們踏出.
第一步的時候.
從那條路上.
向我們慢慢顯露出來.
大家有沒有想過.
我們怎樣移民.

$^{281}$我們移民的目的是為了什麼.
有些人是為了.
下一代讀書.
有一個好的學校.
有些人是為了.
生活更加好.
有些人是為了.
避開某種的壓力.
有些人.
可能都是跟著家人而來.
看小孩.
無論我們的答案是什麼.
你肯定不是為了.
傳福音而來.
很多人更加是.
來了這裡才信耶穌.
所以我們的心裡.
會不會覺得.
我們這樣離開了原來的地方.
偏離了神的心意.
對於那些已經信了耶穌的人.
想深一層.
可能不一定是.
正正是因為.
我們這種不輕易.
不經意地離開的人.
神就在你的心中.
將祂的心意.
顯露出來.
這也是為什麼.
路加沒有將他的焦點.
繼續放在這些人.
他們會失去什麼.
的身上.
反而問神在這個分散裡.
祂正在做什麼.
答案就是.
當我們向這個陌生的希臘人.
來分享福音的時候.
路加就看到.

$^{321}$聖靈.
已經在他們前面.
先行了.
同樣今天.
無論我們為了什麼.
事情.
來移民.
我們有沒有看到.
神在我們心裡.
的那種心意呢.
神已經在這裡先行了.
所以當這件事情.
從安提阿那裡.
傳回耶路撒冷.
教會那裡.
那些.
教會的領袖.
如果你是教會的領袖的時候.
你會怎麼想呢.
你會不會覺得.
這班信徒.
這麼快.
使人們來相信耶穌.
會不會有些.
好像.
不是.
合傳統啊.
會不會太快啊.
會不會有些什麼.
手腳啊.
但原來當他們明白.
原來聖靈已經在他們.
先行的時候.
聖靈已經在安提阿那裡.
做工的時候.
他們就馬上猜獻了.
巴拿巴.
看看他們到底在做什麼.
這個原因其實很簡單.
一方面因為耶路撒冷.

$^{361}$成為了早期.
教會的中心.
他們對於各地的信徒.
他們的發展.
其實是有.
關乎的責任的.
而安提阿那更加是一個.
大集會的都市.
還突然有一班.
人信了耶穌.
他們需要有一個.
很清晰的木樣的指引.
而另一方面.
我們想.
如果你有看到.
這一張的前上半部.
你會看到.
彼得當他向.
外邦人哥利柳.
來到.
宣講耶穌的時候.
他回去要和耶路撒冷的領袖.
來解釋.
為什麼這班人.
哥利柳會信耶穌.
顯示出.
其實教會對於.
這個士工.
向外邦人傳福音.
這個士工到底合不合乎.
神的心意.
其實他們都在摸著石頭來過河.
他們仍然都在.
摸索階段.
他們的心其實可能都未曾.
完全準備好.
但是聖靈.
就在他們前面.
先行.
所以在這個背景下.

$^{401}$耶路撒冷的人.
耶路撒冷的領袖.
他派遣巴拿巴.
去安提阿那.
來看一看.
其實是很順理成章的事.
說到巴拿巴.
巴拿巴兔子是什麼人呢?.
原來在第四章的時候.
都已經有介紹過.
我們看一看.
原來他是一個利未人.
他在塞浦路斯.
即是居比路那裡.
出生.
他本身也有.
希臘文化的背景.
對於所謂離散的人群.
他們的生活.
其實都有所了解.
而他的名字.
更加看得到.
是叫做勸慰旨.
即是說他本身.
其實他的名字已經帶有.
他有一種.
建立人,鼓勵人的意思.
但更重要的就是.
當保羅.
他被人懷疑的時候.
他挺身而出.
為他們做見證.
這些經歷.
其實使他.
成為一個跨文化的.
模樣的人.
是非常之適合的.
所以在第24節.
你看到路加怎樣形容他.
他說巴拿巴.

$^{441}$原本是一個好人.
被聖靈充滿.
但是有信心.
所以.
當他去到安提柯的時候.
他所看到的.
是什麼呢?.
和別人看到的.
是很不同的.
他看到的不是那些.
生活習慣,飲食條例.
與人不同的群體.
當然是.
他們當然不同.
但是他被聖靈感動.
他看到的是一個.
有神恩典同在的群體.
並且因此.
十分之喜歡.
想像一下.
當時的張力.
是怎樣的呢?.
如果一些猶太人.
一些外邦人.
大家一齊來崇拜.
崇拜之後.
一齊來吃飯.
哪個文化遷就哪個呢?.
很有問題的.
對於那些.
新去的.
牧者來說.
為了不出錯.
最簡單的是什麼呢?.
都是按傳統來做.
不要搞那麼多.
先做著.
慢慢才改.
但他不是.
巴伯克是展現出.

$^{481}$那種屬靈的洞察力.
他先去確認.
這些是不是神的工作.
他憑信心.
看到這些是神的作為.
而不是看到.
他們之間文化的差異.
巴喇巴的看見.
現在回想起來.
這幾年我都有參加.
承鄉短孫的工作.
其實深深體會出.
巴喇巴.
他看見神所賜的恩.
就歡喜這種感受.
鄉間的教會.
其實人不是很多.
但每一個人.
他背後都有他的.
被神帶領來到當地.
並且加入教會的.
這件事.
由他們身上.
你可以看到神的工作.
其實是無處不在.
真是令我大開眼界.
和巴喇巴有同感.
非常之欣慰.
神的工作竟然是在.
一些.
沒有什麼人道的地方.
來展開.
所以各位弟兄姊妹.
我們可能都來自不同的地方.
就算大家.
在中國來也好.
我們的生活習慣.
都可能很不同.
求主來到連問我們.
我們要有巴喇巴的那種眼光.

$^{521}$我們多看.
神的恩典.
我們多看人的優點.
我們不要只看到.
大家文化上的差異.
接下來.
巴喇巴.
他就顯出.
勸惠子的本色.
他就勸勉眾人.
要立定心智.
來恆久靠主.
他說.
恩典不是等於放任.
喜歡不等於.
沒有方向.
真正的牧羊.
是需要在.
辨識神的作為之後.
隨之引導信徒.
將這個焦點.
放在神的身上.
接下來.
在第25和26節.
他就更加做出一個.
非常之重要的行動.
做了什麼呢?.
就是去了大掃那裡.
找寶囊.
對於那些移民了的人來說.
他們的危機.
再不是.
受到逼迫.
而是什麼呢?.
感覺到會無根.
感覺到很飄.
很虛.
但有沒有想過.
為什麼他們要找掃羅呢?.
因為巴喇巴他也知道.

$^{561}$一間正在成長的教會.
一間這麼.
文化複雜的教會.
他需要的.
就是向下紮根.
才可以向上成長.
而巴喇巴他承認他自己的有限.
他不會獨自來承擔.
全負上身.
因為他知道.
神的工作.
更加不會是一個人.
去完成.
是需要同工的配搭.
和互補.
才可以做得成.
而剛剛好.
保羅他那時候.
正在好像沒什麼事要做.
而他的背景.
也是和巴喇巴.
他很相似.
而更加比巴喇巴好.
就是他對猶太的.
經典非常熟悉.
因此他更加.
是一個.
牧羊.
這個.
跨文化教會的一個更加適合的.
導師.
結果是什麼呢?.
結果當然是.
他們兩人一起事奉.
聖經說到.
他們在那裡一起事奉.
來教導很多人.
足足有一年.
而從另一個角度看.
信了耶穌的人又如何呢?.

$^{601}$信了耶穌的人.
是不是什麼都不用做呢?.
當然不是.
保羅都要去教他們一年.
耶穌的人當然是需要.
在真理上更加深入地認識.
耶穌.
以致路加在.
後面的章節.
更加會告訴我們.
安提阿.
成為了一個.
宣教的基地.
取代了耶路撒冷.
正是一個這樣的群體.
一些這樣的人.
這些信徒.
第一次被稱為.
基督徒.
在當時.
基督徒是一個.
什麼稱呼呢?.
可能是帶有貶義的.
因為.
在社會上.
可能那些人是嘲諷他們.
這些就是那些口口聲聲說.
講耶穌的人.
是沒有神的.
但正正是這個稱呼.
對我來說是十分之珍貴.
是不是?.
還有提醒我們.
基督徒的身份.
不是因為地理.
不是因為血統.
不是因為背景.
而是因為.
在我們的生命裡面.
人是看到耶穌基督.

$^{641}$所以耶穌基督的.
基督徒的稱呼.
不是什麼的完美.
什麼的好.
而是看到.
我們是願意.
讓基督活在我們的生命裡面.
我們是願意.
在這個分散裡面.
仍然是很忠心地來傳講耶穌.
我們願意.
在我們這麼多不同的.
裡面.
能夠對人有更多的.
包容和更多的.
恩惠.
我們願意.
與人家.
分享福音的時候.
成為別人的祝福.
所以其實這段經文.
是提醒我們.
教會.
能夠在外面.
站得穩.
是完全因為有人願意.
將福音與人分享.
有人願意.
以這個恩典.
去牧養教會.
有人願意.
為了教會的成立.
來付上代價.
更加有信徒.
能夠願意.
繼續來成長.
所以各位弟兄姊妹.
今天.
我們回到教會.
我們又怎樣來擺上我們自己呢?.

$^{681}$我們怎樣來回應.
上帝.
在我們身上的恩典呢?.
接下來.
講到安提柯教會.
被建成立之後.
當我們慢慢趨向成熟.
而經文的焦點.
慢慢就已經.
從他們的身上.
轉移開去.
轉移到那些需要的人身上.
經文講到.
當時是有幾位先知.
從耶路撒冷那裡.
來到.
其中有一位叫.
阿加布.
他藉著聖靈.
指出將來.
會有一個.
大的饑荒.
但值得我們留意的.
就是經文的重點.
不是講到.
這個饑荒是怎樣.
預言.
是應驗.
而是講到.
教會是怎樣去回應.
這個預言.
這個饑荒.
所以在第29節那裡.
你會看到一個募捐行動出現.
但是.
你留意一下裡面所寫的字.
定義這個字.
不是指到.
那些人被命令去做這件事.
是指到.

$^{721}$各人是照著各人的力量.
換句話說.
各人是按照神給他的恩典.
來做出捐獻.
所以這個.
募捐行動可以說是.
信徒之間的一個.
成熟的回應.
但更重要的就是.
方向.
你看到什麼.
安提拉教會.
本身就因為.
耶路撒冷那裡.
被逼迫而產生.
現在它是.
成為了回歸耶路撒冷.
的一個.
教會.
大家有沒有想過.
當時耶路撒冷的教會.
其實是受到.
雙重的打擊.
第一就是被猶太教.
逼害.
再加上饑荒.
其實他們是很需要.
外面的援手.
所以在30節那裡.
他們透過巴拿巴和蘇羅.
把這個捐獻.
交給長老.
他們就完成了一次.
跨地域跨文化.
跨處境的肢體的慣故.
這個也是.
使他們的教會.
能夠成為別人的.
祝福的出口.
這一幕是不是很感人呢.

$^{761}$一間新成立的教會.
一群新來的人.
竟然是可以.
負擔得起.
對武會的公益責任.
他們沒有想到.
他們自己需要什麼.
反而是在想.
在神國裡面.
是誰最需要他們呢.
原本.
他們是需要被人.
接納的群體.
現在卻成為了.
公認的群體.
再想深一層.
是不是只是.
說到捐獻這麼簡單呢.
其實也是.
代表了他們之間的.
一個合一的見證.
教會的合一.
從來都不是建立在地域.
制度上面.
而是在基督裡面.
雖然安提哥教會.
他們本身不是.
很富有.
很豐富.
但是他們學會了.
怎樣能夠將神給他們的東西.
與人來分享.
因為我們在基督裡面.
信主雖然是有先後.
但不會有高低.
我們信主雖然有.
不同的背景.
但都是互為肢體.
都是彼此.
來承擔.

$^{801}$英文有一句說話.
Sharing is caring.
當你能夠與人分享的時候.
你才會.
看到人的需要.
各位弟兄姊妹.
今天我們又怎樣呢.
我們願不願意.
也將神所賜給我們的東西.
與人來分享呢.
同時安提哥教會.
也幫我們重新來理解.
甚麼是宣教.
很多時候.
我們以為宣教.
一定要去很遠的地方.
但經文中.
宣教其實首先.
是需要對你身邊的人.
作出回應.
安提哥教會的信徒.
他沒有離開自己的城市.
卻是已經活出.
一種向外.
向人傳福音的心智.
在今天.
我們這個社會.
又怎樣呢.
我們這個社會.
正是經歷到不斷.
人口的流動.
有人離開.
有人也是走近.
很多未信的人.
可能已經成為了我們中間的.
鄰舍.
我們的同事.
我們熟悉的人.
在這種環境裡面.
能夠願意.

$^{841}$能夠看到別人的需要.
哪怕只是我們.
力所能及.
我們走出一小步.
我們已經踏出了宣教的道路.
記得有一次.
有位弟兄跟我分享.
他說.
其實我沒有想過.
我會參加宣教.
因為我年紀大了.
我英文又不好.
體力有限.
這是年輕人的事.
但後來發覺.
原來很奇怪.
神將一位.
馬來西亞的家庭.
搬到他附近.
最初大家只是打招呼.
關心一下.
慢慢建立了關係.
後來.
這個家庭有些問題.
他主動.
跟他講了一句說話.
我可不可以為你禱告.
這句很簡單.
亦是很窩心的說話.
他是打開了.
分享福音的門.
他告訴我.
原來宣教.
不一定要去到很遠.
神會.
將人帶在我們身邊.
只要我們願意.
來回應.
其實宣教就是這樣開始的.
換句話說.

$^{881}$在現代來說.
宣教的界限.
已經不是一刀切.
世界.
人口的流動.
比以前.
頻密很多.
那些不信耶穌的人.
已經默默地.
在我們身邊出現.
因此安提河教會.
給我們的提醒.
是很清楚的.
教會不一定要走得很遠.
但一定要學會.
用福音傳出去的.
那種心智.
一個真正.
被基督建立的群體.
最終一定.
能夠成為別人的祝福.
而今天.
我們這間教會.
亦都是一樣.
無論你是新移民.
你是舊移民.
或者你是已經在這裡住了很久的人.
只要你仍然活在神的恩典裡.
你就不是異客.
你是神家裡.
帶著使命的人.
我們中間有些弟兄姊妹.
可能已經升級了.
可能已經成為祖父.
母輩.
可能你覺得體力有限.
但我們可以用.
關心,我們可以用代禱.
我們用支持.
來參加神在你們.

$^{921}$這個地方的工作.
而成為神祝福.
人的管道.
我們不要隨便放棄.
最後總結一下.
雖然唐代詩人.
王維說過.
獨在異鄉.
為異客.
這句說話.
很真實.
很感人.
他說出了人在異地的時候.
那種無法抹去的.
那種孤獨感.
但當我們走完整個經文之後.
我們會發覺.
其實我們在基督裡.
我們不是孤單.
我們是被接納.
成為神家裡面的人.
也許我們都會不時來懷念.
過去.
我們不時都會感覺到.
我們還是不太習慣.
還是需要去適應.
但是老家在提醒我們.
神母.
讓我們完全適應.
完全安頓好.
才可以使用我們.
所以我鼓勵大家.
在這個星期裡面.
來嘗試.
做下面三步.
也許你已經在做.
但如果你沒有做的時候.
我們一起做好不好.
第一就是留意我們身邊的人.
特別是那些.

$^{961}$正在經歷變動.
感覺到孤單的人.
特別是那些新來的移民.
我們與他們連結.
我們主動和他們打招呼.
多點微笑.
他們建立良好的關係.
第三就是祝福他們.
我們可以在心裡面.
為對方禱告.
我們求神賜給他們平安.
當然有機會的時候.
我們更加送上溫暖的說話.
當然硬笑的福音.
是不太好的.
求主讓我們有智慧.
在適當的時候.
說出適當的話.
更加求聖靈.
來帶領.
我們就不是耳客.
我們是神.
使用我們.
來祝福人的常客.
我們是否願意呢?.
最後.
我有一副對聯.
作結束.
上卷就是.
樂在異鄉.
非孤客.
同奔天路.
是一家.
橫批就是.
以馬來尼.
在基督裡面.
我們不是孤單.
雖然天國很遙遠.
但我們是一家人.
在這個馬年.

$^{1001}$以馬來尼的神.
來走一條.
天路.
讓我們一起禱告.
主讓我們承認.
我們身在異鄉.
有時候都會疲倦迷失.
但你的愛.
是始終不變.
多謝你提醒我們.
你帶領我們來到這個地方.
我們成為了一家人.
求你在我們.
不安的心裡.
建立我們的信心.
在你的恩典裡面繼續發光.
讓我們不單只是.
為了生存.
更是為了見證你的榮美.
成為這個地方的祝福.
聽我們的禱告.
奉主耶穌基督的聖名.
阿們.
\newpage



\section{創世記 32:22-32}
\label{sec:031680}
\textbf{雅各之爭三部曲:與神之爭}~~~~林穎豪傳道
\newline
\newline
連結: \href{http://www.church.com.hk/acms/content.asp?site=cdc\&op=show\&type=product\&code=031680\&layout=sermon}{\texttt{031680}} ~~~~ 語音日期: 2026-02-22
\newline
\newline
\hyperref[sec:031679]{\small{< < < PREV SERMON < < <}}
~
\hyperref[sec:index_chronic]{\small{[返順時目]}}
~
\hyperref[sec:031681]{\small{> > > NEXT SERMON > > >}}
\newline
\newline
創世記 32:22-32
\newline
\begin{longtable}{cl}
\hline
\hline
章節 & 經文 (和合本修訂版)\\
\hline
32:22 & \begin{tabularx}{0.7\textwidth}{X} 他夜間起來，帶著兩個妻子，兩個婢女和十一個孩子，過了雅博渡口。 \end{tabularx} \\ \\ \relax
32:23 & \begin{tabularx}{0.7\textwidth}{X} 他帶著他們，送他們過河，他所有的一切也都過去， \end{tabularx} \\ \\ \relax
32:24 & \begin{tabularx}{0.7\textwidth}{X} 只剩下雅各一人。有一個人來和他摔跤，直到黎明。 \end{tabularx} \\ \\ \relax
32:25 & \begin{tabularx}{0.7\textwidth}{X} 那人見自己勝不過他，就摸了他的大腿窩一下。雅各的大腿窩就在和那人摔跤的時候扭了。 \end{tabularx} \\ \\ \relax
32:26 & \begin{tabularx}{0.7\textwidth}{X} 那人說：「天快亮了，讓我走吧！」雅各說：「你不給我祝福，我就不讓你走。」 \end{tabularx} \\ \\ \relax
32:27 & \begin{tabularx}{0.7\textwidth}{X} 那人說：「你叫甚麼名字？」他說：「雅各。」 \end{tabularx} \\ \\ \relax
32:28 & \begin{tabularx}{0.7\textwidth}{X} 那人說：「你的名字不要再叫雅各，要叫以色列，因為你與神和人較力，都得勝了。」 \end{tabularx} \\ \\ \relax
32:29 & \begin{tabularx}{0.7\textwidth}{X} 雅各問他說：「請告訴我你的名字。」那人說：「何必問我的名字呢？」於是他在那裡為雅各祝福。 \end{tabularx} \\ \\ \relax
32:30 & \begin{tabularx}{0.7\textwidth}{X} 雅各就給那地方起名叫毗努伊勒 ，說：「我面對面見了神，我的性命仍得保全。」 \end{tabularx} \\ \\ \relax
32:31 & \begin{tabularx}{0.7\textwidth}{X} 太陽剛出來的時候，雅各經過毗努伊勒，他的大腿就瘸了。 \end{tabularx} \\ \\ \relax
32:32 & \begin{tabularx}{0.7\textwidth}{X} 因此，以色列人不吃大腿窩的筋，直到今日，因為那人摸了雅各大腿窩的筋。 \end{tabularx} \\ \\
[1ex]
\hline
\hline
\end{longtable}
$^{1}$不知大家有沒有試過失眠呢?.
這個禮拜我有一天.
好像是我們那晚11點,12點睡.
我睡到3點多就醒了.
然後就眼光光一直到7點,8點都一直睡不著.
不知大家有沒有試過失眠.
我是沒有什麼原因.
但不知為什麼就3點醒了就睡不著.
沒有什麼原因的失眠.
但大家有沒有試過不是因為沒有什麼原因的失眠.
而是因為可能明天有一件事令你非常的害怕.
又卻是不可以不面對的事情呢?.
可能是等待一份醫療報告的結果.
可能是要向老闆承認剛剛這個星期你做了一件很嚴重的錯誤.
又或者可能是要見一個很久以前你得罪過,傷害過的朋友.
可能那種焦慮就好像一塊大石一樣.
壓在你自己的心頭裡面令你透不過氣.
今天我們來到《雅各之爭故事》的最後章.
過去一個月我們見到雅各前半生的奮鬥史.
在第一部曲他在家裡和哥哥耳掃到爭.
結果家庭破碎流亡異鄉.
在第二部曲他在職場和舅父拉班到爭.
在被欺壓當中他學會到盡忠.
最終他帶著妻兒,財富逃離哈拿.
去到今天第三部曲正正雅各處於一個前無去路.
後有追兵的狹縫之中.
他徹夜未眠.
後面我們上一講講到他和拉班和他的叔叔互相已經納了約.
不再互相侵犯,不再互相打擾.
所以他不可以再回去那裡.
但是前面就是他二十年前徹底得罪了的哥哥耳掃.
在這時候僕人回來報告了一個壞消息.
他說這個哥哥耳掃快要來了.
但他來的時候他帶了四百人.
我猜當時沒有任何一個可能性可以說服雅各.
但他的哥哥耳掃不是想要他的性命.
耳掃叫四百人做什麼呢?.
叫四百人一齊來這裡難道是拜年嗎?.
吃團年飯,吃盤菜,花車巡遊.
四百人的陣容我想在雅各眼中很明顯不是來歡迎派對.

$^{41}$而是一支討伐他的軍隊.
哥哥耳掃在二十年前那句.
等我爸爸死了 一死了我就馬上殺了你 雅各.
這句話彷彿在雅各的耳邊再一次響起.
雅各的一生都在找 都在計算.
他以為自己可以掌控一切.
但此刻死亡的陰影籠罩著他.
就是我們剛才讀經說雅博渡口的河邊.
這個在最黑暗的一夜.
雅各將要面對的.
不再是他和其他人的爭.
而是他生命徹底的轉化.
今日就讓我們透過雅各在雅博渡口的經歷.
反思我們的真實信仰生命.
和我們看看神究竟是怎樣叫人的生命徹底轉化.
和成長.
雅各面對著耳掃的四百大軍 雅各就怕得要死.
聖經形容他是心事懼怕.
而且他受煩.
上一講我們先講到其實雅各剛剛經歷完一個的靈性突破.
其實他終於明白.
原來生命的高低不在乎自己用多用力去抓住 多用力去爭.
而是在乎那個掌權的神.
不過.
不過今天來到在恐懼裡面.
我們再次見到雅各.
那個精打細算的老鵝.
再一次啟動了他的防禦機制.
第一我們見到雅各仍然都是很醒目的.
他就將他隨行的人口.
羊群牛群全部分作兩隊.
他心想如果他哥哥耳掃真的要來殺我.
整支軍隊 他殺了一隊 我還有另外一隊.
這個很聰明 叫做分散投資.
不要將所有的雞蛋放在同一個籃子裡面.
接著雅各做了什麼呢.
接著雅各做了一件很熟悉的事情.
他祈禱.
他彷彿這個是雅各的生命裡面記載.
他最謙卑最謙卑的一次禱告.

$^{81}$他說 耶和華神啊 我雖然是很不配.
但是你仍然將一切的恩典慈愛誠實.
賜予給我.
神啊 我是不配.
但你仍然賜給我.
眼前有一個很大很大的難關.
主啊 求你去救我.
這個禱告我們見到彷彿都是一個不斷的篇幅.
讓我們可能看到的時候我們都會想.
這個禱告聽起來真的是一個很全然的交託 是不是.
不過很諷刺的是.
當雅各一做完這個祈禱.
他做了什麼呢.
他祈完禱之後.
他又馬上回到他自己的籌算裡面.
他預備了極其豐厚的禮物.
很多羊 很多牛 很多駱駝.
並且他將這些禮物全部分成一批一批.
他說我預備這麼多禮物不夠.
我還要一次過送去.
我送完一批.
送完羊.
再送牛 再送駱駝.
等哥哥見到我的誠意.
等哥哥去消氣.
親愛的弟兄姊妹.
這個豈不是我們信仰裡面很真實的寫照嗎.
可能我們上個星期才剛剛經歷完.
神對我們的保守.
可能我們剛剛才在讀經裡面.
早兩天在靈修裡面 讀經裡面 詩歌裡面.
被神激勵 提醒.
當頭棒喝.
可能我們早兩天.
我們才經歷完這些事.
但是當我們再次面對困難.
面對恐懼.
我們內心的小信又再一次偷偷跑出來.
我們一面向神禱告.
我們求神掌權.

$^{121}$我們就說 主啊 我交託給你.
但一祈園禱轉頭.
我們就繼續用自己的方法.
用自己的手段 瘋狂向四周圍去找.
因為我們心底裡面.
我們仍然都相信不過.
神最後會保護我們.
我們雖然祈禱.
但我們又好像雅各一樣.
其實最後只是把神當作是最後的保險.
雖然我們有祈禱 但我們仍然最後很想靠自己的方法.
雅各以為只要禮物夠多.
只要送禮的策略夠好.
就可以買通這場危機.
保住自己的性命.
但他忘記了.
其實20年前.
正正就是因為他想用紅湯去買長子的命婚.
用羊皮的藝妝去買祝福.
他才能獲得這個田地.
到今天.
難道雅各還要用相同的方法去解決問題嗎?.
又或者不是雅各以為.
又或者是我們以為.
我們以為自己只要有夠多的錢財.
以為自己疏通得夠多的人脈關係.
又或者我們以為自己懂得夠多的知識.
個人的技巧技術.
我們就可以在人生裡面得著安全感.
什麼困境都難不倒我們.
所以不是雅各忘記了.
而是我們忘記了.
20年前.
10年前.
1年前.
2天之前.
正正是我們單單想靠著自己去找去爭取.
卻碰到焦頭爛額.
到今天.
難道我們還要繼續重蹈覆轍嗎?.

$^{161}$面對著充滿著人性軟弱的我們.
面對著容讓恐懼主導決策的雅各.
我想神很清楚知道如果不徹底去打破.
去拆毀這種靠自己掌握一切的安全感.
神很清楚知道如果無論是我們.
無論是雅各.
都永遠沒有辦法真正的成長.
去到夜間.
雅各將所有的妻兒家當都打發過河.
剩下他自己一個留在雅各渡口.
這是雅各一生最孤單.
也是他最赤裸最真實的時刻.
就正正在這個四周面黑漆黑一片.
雅各單人匹馬的時刻.
有一個人向雅各衝過來要和他摔交.
這個向雅各衝過來要和他摔交的人是誰?.
有沒有人知道.
這個衝過來雅各都.
和他摔交的是誰?.
今天如果我們有看聖經.
我們就知道.
荷西亞書講到原來這個和雅各摔交的.
是神或者是神的使者.
是天使.
但是當時如果我們代入一下雅各.
那個看著一個人衝過來的雅各.
他知不知道那個是天使呢?.
我們嘗試代入一下雅各.
當時雅各在想什麼呢?.
就是在想怎樣有什麼方法可以將人口生畜.
將我的妻兒甚至將我自己安置.
怎樣才是最安全呢?.
怎樣才可以逃過哥哥耳熟的刀劍呢?.
就在他剛剛在想這些的時候.
又剛剛將那些妻兒送了過河.
只有他自己一個人的時候.
有人衝過來攻擊我.
這個是誰?.
雅各心裡在想他在做這麼多的籌算的時候.
很害怕哥哥耳熟的時候.

$^{201}$有人衝過來.
雅各應該一定在想.
這個一定是耳熟.
又或者這個是耳熟派過來的殺手.
趁我一個人.
去殺了我.
心裡充滿著恐懼的雅各.
他再一次在這個時候硬著頭皮衝過去.
一出手抓住這個人.
去教力.
因為雅各很害怕.
一不留神他自己隨時沒命.
雅各這個一生都在用力抓住的人.
現在他就很用力抓住這個神秘人.
他用盡全力去推,去拉,去抗衡.
他不服輸.
他不願意被制服.
正如他的一生裡面他不願意馴服他的哥哥.
不願意馴服他的爸爸.
不願意馴服拉班甚至他不願意馴服神.
經年記載這場摔跤從夜間開始一直.
持續到黎明.
這個是一個極其漫長.
我們想像可以.
跟其力竭但是又膠著的一個搏鬥.
一直打到不知道什麼時候完結.
但是突然之間.
二十五節.
迎來了這個轉捩點.
那人見自己勝不過他.
就在雅各的大腿窩摸了一巴.
又或者好像新日本的翻譯是擊打了一下.
雅各的大腿窩.
正在摔跤的時候就扭了.
經文記載.
這裡究竟是矛盾還是諷刺呢.
從文字上面去到看.
我們看到.
天使見自己.
贏過雅各.

$^{241}$即是誰贏了.
即是雅各贏了.
天使贏不了即是雅各贏了.
是雅各勝利.
又正如這個其實都是天使自己所講二十八節所講到的說話.
你和神和人的交流你都贏了.
但是這個是從文字去看.
但若我們從現實去到看.
大腿窩扭了又或者好像新日本的翻譯.
大腿窩脫節脫了膠.
又或者這個大腿窩扭了的.
卻是雅各.
是那個得勝了的雅各.
臣藉著天使和雅各的摔跤.
再一次讓雅各用盡了自己的力氣去爭.
並且再一次讓雅各贏得勝利.
但是得到的是什麼呢.
得到這個勝利之後他得到了什麼呢.
和耳掃之爭.
雅各贏得的是被追殺被逃亡.
和拉班之爭他贏得的是被奴役被追趕.
今次他和神或者和天使之爭.
他贏得的是一個扭曲的大腿窩.
一個脫節的大腿窩.
大腿窩是人體其中一個最強壯的肌肉的支撐點.
是人站立在這裡是人發力的一個核心.
神的使者藉著摸或者藉著擊打.
令到雅各的大腿窩扭曲,脫節.
意味著神廢去了雅各引以為傲的力量.
那一刻劇痛傳來.
雅各站立也站不穩.
彷彿告訴他以後他也不可以再靠自己的力量.
去戰鬥,去逃跑.
這一幕其實很有意義.
原來神為了拯救我們.
有時不可以不拆毀.
和打破一些我們最引以為傲的東西.
可能是我們的健康.
可能是我們的財富.
可能是我們很引以為傲的事業.

$^{281}$可能是我們很重視的關係.
甚至是你自己的聰明才能,技巧.
神讓你們在這些自己引以為傲的事上經歷失敗.
讓你感覺到無力.
甚至讓你受到極大的傷害和痛楚.
其實神不是想害你.
而是神想讓你好像雅各一樣.
在無力之中學會一件事.
不再是抓住.
而是抱.
二十六節 當大腿扭傷之後.
雅各沒辦法再去摔跤.
他只能夠做一個動作.
死命地抱著天使不放.
他說你不給祝福 我不讓你走.
這個是雅各生命裡面一個很重要的轉折.
過去他想要祝福.
他是要用換的.
用騙的,用搶的.
現在.
他想要祝福.
他是在軟弱裡面.
在劇痛裡面.
在無力裡面.
單單去到仰望神的陰天.
他知道他想要祝福不再是一個交易.
而是乞求.
正如剛才我們在荷西亞書有提到雅各的故事.
荷西亞書用的字.
是雅各哭泣,懇求.
上一次他是站在拉班的面前.
理直氣壯地宣告神怎樣保護他.
如今.
他是跪在地上.
拉著他的眼淚.
忍著他的痛楚.
很卑微的.
去求神去保守他.
弟兄姊妹我們有沒有試過這個時刻呢?.
四周七黑一片.

$^{321}$你自己孤身一人.
彷彿沒有任何的朋友,沒有任何的親人去陪伴你.
連你的家人.
連教會的人都不明白你.
你過去一生最倚靠的,最引以為傲的,最可靠的.
都化為烏有.
你什麼都沒有.
這個時候還有人衝過來跟你下井下石.
要跟你相爭,徹夜去耗盡你的時間,心機精力.
然後還要反過手來.
把你打傷,把你打倒.
打到你傷一傷之餘還要揶揄你.
對喔,很厲害喔.
你贏了我喔,很厲害喔.
我們以為神看不到.
我們以為神不知道我們發生什麼事.
我們以為神聽不到我們的祈求.
是神不理我們.
甚至我們以為是神伸手攻擊我們,懲罰我們.
但是我們有沒有想過.
其實是神想要拯救我們.
是神想要從只懂得依靠自己的我們.
去拯救出來.
是神想要拯救我們,教會我們.
什麼叫做完全的降伏在祂的帶領之下.
所以神才把支撐著我們生命的每一條支柱.
通通拆毀.
讓我們徹底降伏於祂.
當雅各去求天使的祝福.
他就問了雅各一個問題.
你叫什麼名字?你是誰?.
神和天使豈不知道雅各叫什麼名字?.
神當然知道那個就是雅各.
天使當然知道那個就是雅各.
所以其實這個不是一個問題.
而是一個帶著審判的靈魂敲問.
二十年前雅各想要祝福他父親.
以撒問他一個同樣的問題.
「我的兒子你是誰?」.
雅各當時怎麼回答?.

$^{361}$昔日的雅各用謊言掩飾自己的身份.
「我是你的大兒子二叔」.
今日雅各想要祝福.
神再一次問他「你是誰?」.
「你叫什麼名字?」.
如今在神的面前.
雅各知道自己無法再去偽裝.
雅各最終坦言「我就是雅各」.
當雅各承認自己是雅各.
等於他承認自己是那個抓奪者.
昔日的欺騙者.
這是他在神面前最徹底的認罪.
他說「主啊 沒錯 我就是雅各」.
「我就是那個很卑鄙 很詭詐」.
「一直只想靠自己去找取祝福的雅各」.
就像我們剛才唱的詩歌一樣.
耶和華善愛其中一句說道.
「神令我們審查我們心裡的幽暗」.
然而就在他承認自己本相的一刻.
天使卻說「你的名字不要再叫雅各」.
「以後你叫以色列」.
「因為你與神 與人教力都得了聖」.
以色列這個名字有什麼意思呢?.
如果我們從一個最傳統的理解.
以色列在希伯來誌可以分拆成兩個部分.
一個部分是掙扎教力.
另一個部分是神.
意思就是一個很直接的理解.
就是一個與神掙扎,與神教力,堅持與神相爭的人.
正如這個就是神的使者對雅各改名的解釋.
如果大家記得.
第一講我們有說到雅各的名字是什麼意思.
雅各的名字是什麼意思呢?.
雅各的名字的意思就是「抓住」.
不單止是抓住.
還是抓住腳跟.
因為大家記得雅各出世的時候抓住他哥哥二叔的腳跟.
意思是什麼意思呢?.
抓住腳跟.
意思就是彷彿一個潛伏在暗處.

$^{401}$一個等其他人走過之後.
然後從後面抓住人的腳跟.
絆倒人的一個人.
一個很狡猾,很陰險的人.
但是今天他要變為以色列.
是一個與神教力,掙扎的人.
意思是他不再是一個站在別人後面.
去暗算別人的人.
如今他是要面對面.
用正面去面對其他人.
甚至去面對神的一個人.
但是以色列這個名字還有另一個拆解.
如果你依從當時的文化 原來希伯來人,阿摩尼人當時的名字.
如果在右邊的部分,後面的部分.
出現的字是神明的時候.
絕大部分是將神明視為主體.
而不是視為客體.
給大家一些例子.
在聖經裡面也有.
例如在《創世記十讀章》裡面.
「當夏革生爾實瑪利」.
「神的使者」就說這個名字的意思.
我們見到一邊是「聽見」一邊是「神」.
「神的使者」就說這個名字的意思.
是「神聽見」.
而不是這個兒子,這個人.
「你要聽神話」.
是「神聽見」.
不是「你聽神話」.
又例如.
「先知以西結」的名字.
這個名字是「神堅固」.
「神加力」.
很明顯這個就是先知的名字.
而很明顯不會有一個先知的名字叫做.
一個人,一個先知去到「堅固」或者「加力」給神.
還有一個例子就是在《史詩紀》裡面.
當基癲去到拆毀巴力的一些壇的時候.
一些祭壇的時候.
有人就叫基癲做「耶路巴力」.

$^{441}$一個就是「爭辯」.
一個就是後面「巴力」的意思.
意思是甚麼意思呢?.
聖經說是「等巴力和祂爭辯」.
而不是「等這個人和巴力去到爭辯」.
所以甚麼意思呢?.
以西結這個名字反過來是甚麼意思呢?.
以西結這個名字其實不是.
其實不是在說雅各要繼續同神相爭.
繼續要同神教力,掙扎.
反過來.
是要叫雅各放手.
叫雅各給神爭,給神教力.
昔日雅各和哥哥以素由出世爭到離家出走.
昔日雅各和拉班寇輔爭了二十年.
雅各和神的使者由夜間爭,爭,爭到黎明.
去到這一刻.
神對他說:你要叫以西結.
為甚麼?因為夠了,你放手,不要再爭.
讓我爭,讓我出力,讓我出手.
這就是以西列真正的名字的意思.
你不要爭,讓我爭.
而當神出手的時候.
神就令雅各的大腿窩扭曲.
並且在雅各的身上留下了他一生都有的一個「瘸腿」的記號.
三十一節記載.
日頭剛出來的時候.
雅各經過比魯伊納.
他的大腿就瘸了.
一個一世從此都跛腳的人.
為甚麼聖經有一而再再而三說雅各得了聖呢?.
一個跛腳的人怎樣算是贏呢?.
其實我們從第一講到第三講.
雅各的經歷一直在教導我們.
其實甚麼才叫贏?.
甚麼才叫勝利?.
甚麼是真正的勝利呢?.
原來真正的勝利不是爭贏甚麼.
原來真正的勝利不是打贏人甚至打贏神.
真正的勝利是人忠於肯輸給神.

$^{481}$真正的勝利是人忠於願意輸給神.
當雅各徹底降伏.
不再靠自己的血氣去爭.
而是僅僅去到抓住神的恩典的時候,他才真的贏了.
他贏得一個新的身份.
贏得一個真正的祝福.
當太陽剛剛日出的時候.
雅各跛下跛下走過雅博道口.
沒錯,他的身體是殘缺了.
但他的靈魂卻是完整了.
他不再是走得很快,甚至走路帶風的人.
那個很精於計算的生意人雅各.
他現在是一個身上帶著神恩典記號.
懂得依靠神的以色列.
騎腿這個記號在他下半生繼續每一天都提醒他.
陪伴著他,跟隨著他.
這個記號在他一生裡,他每走一步都提醒他.
沒錯,我真的沒有力.
我真的身體沒有力.
但是他已經不再需要靠自己的力.
他要靠的單單是神.
這個騎腿的記號也成為他見證的一部分.
見證證明他真的遇到過永生的神.
這是一個軟弱的記號.
卻又是一個真正特性的印記.
我再說一次,這是一個軟弱的記號.
卻又是一個真正特性的印記.
我們也有一個軟弱的記號.
又代表著特性的印記.
我們有沒有一個軟弱的記號.
作為基督徒.
我們有沒有一個軟弱的記號.
又象徵真正特性的印記.
有沒有?.
是什麼?.
有沒有?.
十字架,是不是?.
昔日耶穌基督.
在十字架上經歷死亡.
為我們成就了救恩.

$^{521}$沒錯,這就是一個很軟弱.
卻又有特性的記號.
不過其實這不是我最想問你的問題.
我最想問的是.
在你自己的生命裡面.
有沒有一個只是屬於你.
是你軟弱的記號.
卻又是代表著你真正特性的印記呢?.
這個印記可能是肉眼可見的.
又或者是你埋藏在心裡面.
深深處一段回憶.
可能是你過去曾經被人傷害過.
深深傷害過的疤痕.
又或者是你過去是你自己曾經傷害過其他人的經歷.
可能是你自己覺得自己生命裡面有些限制.
有些不完美.
有些做得不好的地方.
在我讀神學的時候.
我有一個神學老師教我.
上帝創造我們每個人其實都有自己的死穴.
我們總會有自己的限制.
總會有自己的軟弱.
並且有時候這些死穴是跟我們一輩子的.
當然如果有方法去幫我們跨過這些軟弱.
或者治好我們的傷痛.
心裡面的傷疤.
固然是好事.
但是有時候.
這些死穴.
這些限制.
不是我們想處理就處理得到的.
有時候.
上帝將軟弱.
將限制.
將死穴放在我們的生命裡面.
就是讓這些死穴.
去提醒我們.
原來我們真的不可以靠自己.
原來我們真的.
只可以持續不斷去依靠神.

$^{561}$去緊貼神.
這樣我們才會走人生前面的路.
我們才會向前走.
因為我們每一個人.
其實我們都有一隻跛了的腳.
一個被扭傷.
脫節.
甚至扭斷了的大腿骨.
這個月我們陪了雅各.
走過了他充滿著爭鬥的前半生.
在第一部曲我們見到雅各.
想用紅湯去交易親情.
結果家無寧人.
在第二部曲.
我們見到雅各.
想用樹枝去操縱財富.
結果發現全是神的供應.
來到今天的第三部曲.
我們見到.
雅各在恐懼之中.
他再次渴望自己去到新手.
捉住不同的救生圈.
但是神在雅博道口的黑夜.
親自的介入.
迫使雅各誠實坦然面對最真實的自我.
為了更深的建立.
神擊打雅各的大腿窩.
廢了他最強的支撐.
拆毀他最深層的保護網.
看似是傷害.
實質卻是恩典.
因為只有當雅各沒有辦法再靠自己去抓住.
逃跑或者戰鬥的時候.
他才懂得.
緊緊去抱住神.
我們也見到.
真正的德性.
原來.
是輸給神.
雅各變成了以色列.

$^{601}$不再是同人鬥爭.
而是讓神為他爭戰.
騎了的腿.
雖然叫到雅各.
往後都走得慢了.
走路的姿勢很難看.
但是卻是他一生裡最榮耀的勳章.
因為那個軟弱的記號卻是提醒他.
他不再靠自己的詭詐.
他是依靠神的恩典去站立.
親愛的弟兄姊妹.
親愛的朋友 今天你的雅各道口在哪裡呢?.
是不是在某些關係裡面 你仍然很想死命地去改變對方呢?.
是不是在一些工作或者經濟的壓力底下 你仍然很用盡各種的小聰明.
來轉型呢?.
是不是在你的內心深處 你仍然很害怕.
如果我不爭 如果我不走 我就會失去一切呢?.
又或許.
神容許困難臨到你.
讓你覺得好像跛了一隻腳一樣無力.
那是因為.
他想藉著痛.
藉著這個死穴.
讓你停止再掙扎.
從而明白.
神境外.
讓我們記得.
最大的德性不是我們掌控了環境.
而是降伏於神.
當我們願意放手 當我們願意不再和神出交 和人出交.
而是緊緊地抱著神的應許的時候.
我們就能夠像雅各一樣.
望見彼路伊勒的日出.
並且帶著恩典的記號.
繼續平安地.
向前行.
傳言降伏於神的雅各.
他最後不再畏懼.
他不再失縮於妻兒的後面.
不再是你們先走 你們先走 我害怕.

$^{641}$他最終是走在妻兒隊伍的前頭 他作帶領 他作承擔.
雖然他已經跛了.
但是.
他堅持七次.
他堅持七次完全的俯伏在地.
然後才迎向耳塑.
啊不是 經文說他不是迎向耳塑 經文說他迎向他的哥哥.
最終雙方沒有發生戰爭.
沒有發生爭執.
耳塑都沒有再和他計算過去的種種因緣.
耳塑作為哥哥.
而哥哥是跑過去迎接這個弟弟.
把他抱住.
是兩個人放聲的大哭.
雅各原本以為他自己要用生命去搏鬥的難關.
原來神早就已經為他化解.
原來他先前所有的籌算 所有的憂慮恐懼.
都是不需要的.
唯有.
那一晚的摔跤.
唯有那一晚他真的全然向神的降伏.
這個.
才是最重要的.
讓我們一起低頭祈禱.
因為你是雅各的神.
你更加願意成為我們這一班.
好像雅各一樣.
充滿著瑕疵的人的神.
主人我們承認我們很多時候都好像雅各.
我們面對心裡面的恐懼.
未知的將來的時候.
我們口裡面雖然經常說「教託」.
但我們心裡面依然瘋狂地去計算,去策劃,籌劃.
我們試圖用自己的方法去建立安全感.
我們很怕失去控制.
我們很怕去到承認自己的軟弱.
因為你有時甚至允許苦難,挫折.
淋到我們的生命.
就好像擊打雅各的大腿窩一樣.
你藉著痛楚去拆毀我們那個驕傲的內在生命.

$^{681}$求主讓我們在痛楚裡面不自絕望.
反倒我們學會去到緊緊抓住你的英血.
好像一個孩子抱著父親一樣不放手.
求你閱立我們生命裡面每一個傷痕,軟弱.
讓這些死穴被你轉化成為你恩典的記號.
求你除去我們那些很想抓住,抓奪的舊個性.
賜給我們一個以色列的新生命.
一個不再靠自己掙扎而降伏於你.
讓你出手,讓你用力.
單單靠著你,讓你為我們爭戰的生命.
當我們帶著軟弱離開這裡的時候.
願意閱立的日頭也照在我們的身上.
讓我們在軟弱裡面.
剛強壯膽靠著你繼續的前行.
我們這樣禱告,認罪,仰望.
是奉主耶穌基督的聖名致祈求.
\newpage



\section{約書亞記 1:1-18}
\label{sec:031681}
\textbf{求主給我看見異象}~~~~吳銘淇傳道
\newline
\newline
連結: \href{http://www.church.com.hk/acms/content.asp?site=cdc\&op=show\&type=product\&code=031681\&layout=sermon}{\texttt{031681}} ~~~~ 語音日期: 2026-02-22
\newline
\newline
\hyperref[sec:031680]{\small{< < < PREV SERMON < < <}}
~
\hyperref[sec:index_chronic]{\small{[返順時目]}}
~
\hyperref[sec:031682]{\small{> > > NEXT SERMON > > >}}
\newline
\newline
約書亞記 1:1-18
\newline
\begin{longtable}{cl}
\hline
\hline
章節 & 經文 (和合本修訂版)\\
\hline
1:1 & \begin{tabularx}{0.7\textwidth}{X} 耶和華的僕人摩西死了以後，耶和華對摩西的助手嫩的兒子約書亞說： \end{tabularx} \\ \\ \relax
1:2 & \begin{tabularx}{0.7\textwidth}{X} 「我的僕人摩西死了。現在你要起來，和眾百姓過這約旦河，往我所要賜給以色列人的地去。 \end{tabularx} \\ \\ \relax
1:3 & \begin{tabularx}{0.7\textwidth}{X} 凡你們腳掌所踏之地，我都照我所應許摩西的話賜給你們了。 \end{tabularx} \\ \\ \relax
1:4 & \begin{tabularx}{0.7\textwidth}{X} 從曠野和這黎巴嫩，直到大河，就是幼發拉底河，赫人的全地，又到大海日落的方向，都要作你們的疆土。 \end{tabularx} \\ \\ \relax
1:5 & \begin{tabularx}{0.7\textwidth}{X} 你一生的日子，必無人能在你面前站立得住。我怎樣與摩西同在，也必照樣與你同在；我必不撇下你，也不丟棄你。 \end{tabularx} \\ \\ \relax
1:6 & \begin{tabularx}{0.7\textwidth}{X} 你當剛強壯膽，因為你必使這百姓承受那地為業，就是我向他們列祖起誓要給他們的地。 \end{tabularx} \\ \\ \relax
1:7 & \begin{tabularx}{0.7\textwidth}{X} 只要剛強，大大壯膽，謹守遵行我僕人摩西所吩咐你的一切律法，不可偏離左右，使你無論往哪裡去，都可以順利。 \end{tabularx} \\ \\ \relax
1:8 & \begin{tabularx}{0.7\textwidth}{X} 這律法書不可離開你的口，總要晝夜思想，好使你謹守遵行這書上所寫的一切話。如此，你的道路就可以亨通，凡事順利。 \end{tabularx} \\ \\ \relax
1:9 & \begin{tabularx}{0.7\textwidth}{X} 我豈沒有吩咐你嗎？你當剛強壯膽，不要懼怕，也不要驚惶，因為你無論往哪裡去，耶和華你的神必與你同在。」 \end{tabularx} \\ \\ \relax
1:10 & \begin{tabularx}{0.7\textwidth}{X} 於是，約書亞吩咐百姓的官長說： \end{tabularx} \\ \\ \relax
1:11 & \begin{tabularx}{0.7\textwidth}{X} 「你們要走遍營中，吩咐百姓說：『當預備食物， 因為三日之內你們要過這約旦河，進去得耶和華－你們神賜給你們為業之地。』」 \end{tabularx} \\ \\ \relax
1:12 & \begin{tabularx}{0.7\textwidth}{X} 約書亞對呂便人、迦得人和瑪拿西半支派的人說： \end{tabularx} \\ \\ \relax
1:13 & \begin{tabularx}{0.7\textwidth}{X} 「你們要記得耶和華的僕人摩西所吩咐你們的話說：『耶和華－你們的神使你們得享安寧，必將這地賜給你們。』 \end{tabularx} \\ \\ \relax
1:14 & \begin{tabularx}{0.7\textwidth}{X} 你們的妻子、孩子和牲畜可以留在約旦河東、摩西所給你們的地。但你們中間所有大能的勇士都要帶著兵器，在你們的弟兄前面過去，你們要幫助他們。 \end{tabularx} \\ \\ \relax
1:15 & \begin{tabularx}{0.7\textwidth}{X} 等到耶和華使你們的弟兄和你們一樣得享平靜，並且得著耶和華－你們神所賜他們為業之地的時候，你們才可以回到你們所得之地，承受為業，就是耶和華的僕人摩西在約旦河東、向日出的方向所給你們的地。」 \end{tabularx} \\ \\ \relax
1:16 & \begin{tabularx}{0.7\textwidth}{X} 他們回答約書亞說：「凡你吩咐我們的，我們都必做；凡你差我們去的地方，我們都必去。 \end{tabularx} \\ \\ \relax
1:17 & \begin{tabularx}{0.7\textwidth}{X} 我們在一切事上怎樣聽從摩西，也必照樣聽從你。惟願耶和華－你的神與你同在，像與摩西同在一樣。 \end{tabularx} \\ \\ \relax
1:18 & \begin{tabularx}{0.7\textwidth}{X} 無論甚麼人違背你的命令，不聽從你所吩咐他的一切話，就必處死。你只要剛強壯膽！」 \end{tabularx} \\ \\
[1ex]
\hline
\hline
\end{longtable}
$^{1}$弟兄姊妹 新年快樂.
耶穌早晨.
讓我們一同為著我們聽和說的福事一同禱告.
讓我們能夠有一齊相聚的時間.
這是你呼召我們的.
讓我們今日都藉著在經文當中.
叫我們能夠有所領受.
叫你看見上帝你的帶領.
看見你藉著聖靈對我們的提醒.
奉著耶穌基督的聖名祈求 阿們.
弟兄姊妹 不知道你們有什麼事情是很害怕的呢.
有一次我們一群朋友出街.
突然有一個朋友說 小心點.
然後一群人就嚇得彈起.
我就直接衝了出馬路.
回過神來 大家就說 什麼呀 叫什麼要小心點.
那個朋友就說.
我純粹提一提你前面冷氣機滴水 叫你小心點.
大家就很害怕.
最特別是什麼呢.
最不特別是大家聽到這個警告的時候.
大家在想的事情其實都不一樣.
有人以為是有蟑螂.
有人以為有老鼠.
有人又說我以為是有個醉酒佬在衝出來.
甚至有人誇張到說我以為有鬼.
那我呢 我在怕什麼呢.
我在看什麼 我在想什麼呢.
我在這個鬧市旺角的街頭.
我突然在那一刻覺得.
他在叫的是因為有條蛇出來了.
因為我最怕就是蛇.
這件事令我發現了一樣東西.
就是原來有人這樣嚇你的時候.
是會令你直接聯想起一樣你最怕的東西.
根據這個科學的精神.
我一定要去證實一下這個理論是不是真的.
所以我用了同一個方法去嚇一嚇Katie.
結果是真的.
因為她最怕的是蟑螂.

$^{41}$我不知道你怕什麼.
或者你可以回去找人幫你試一下.
但是前提是關係要夠好.
不然下次就見不到面了.
但是通常我們一說怕什麼.
通常第一刻想可能是蟑螂 老鼠.
很多不同的動物或者是生物.
不過再想深一層.
可能這些恐懼都不是長期困擾著我們.
比起這些可能我們更加恐懼的是.
別離 恐懼憂柴憂迷的日子.
恐懼疾病 死亡.
甚至我們今天可以一起去想一下.
進入一個事奉當中的時候.
你又會怕什麼 你會擔心什麼呢.
今天我們講一下恐懼.
特別是怎樣面對進入事奉的恐懼.
說到恐懼的時候.
我們經常會想起一段金句.
今天主席也讀出了.
去鼓勵自己或別人的.
就是約書亞紀的一章九節.
這裡是整個第一章裡面.
剛強壯膽這個字其實出現了四次.
有三次都是在耶和華.
吩咐約書亞去接替摩西的工作的時候.
講到的 就是第六 第七和第九節.
而另一次就是裡面人家德人.
和馬來西亞那半個支派.
去回應約書亞的吩咐的時候去講的.
而根據聖經的敘事.
我們知道不斷重複的這些說話.
很可能是作者要告訴讀者.
一個段落的重點.
所以我們今天解決了三個問題.
第一 約書亞究竟在懼怕什麼.
第二 約書亞怎樣可以剛強壯膽呢.
最後我們要解決一個問題.
那什麼叫做剛強壯膽呢.
首先進入去講到約書亞的恐懼.

$^{81}$講到約書亞的恐懼.
在經文裡面其實沒有很直接地講了.
這個恐懼是什麼.
但我們都可以花點時間.
試試代入當時約書亞的處境.
看看究竟她在怕什麼.
第一個是約書亞接替摩西的工作的時候.
上帝很清晰地跟約書亞說.
要帶領以色列人走過約旦河.
甚至要進入迦南.
這個不單單是在曠野裡漂流生活.
而是更加可以想像到.
之後會有不斷的打仗出現.
作為一個游牧民族的新領袖.
要挑戰七個大邦大國.
不知道會不會在這個新崗位的工作內容.
令約書亞很擔心.
不知道大家有沒有試過.
圍著一些事奉的崗位很擔心.
記得我剛出來做傳道人的時候.
一開始心裡都滿懷信心.
覺得終於讀完神學.
要回應上帝的呼召.
進入全職的事奉.
但當教會真的將一個團契.
這個團契就是你負責.
那些年青人你就看著他們 教他們.
在這個時候就開始擔心了.
心裡每天上班的時候都會說.
不是吧.
有幾樣的.
第一就是 不是吧.
只有我一個大祖.
還要叫年青人分享心靈的軟弱.
怎麼做的.
又說 不是吧.
要我自己一個去處理那些年青人分手.
他們吵得這麼激烈 要我怎麼處理.
最後 最後不是吧.
我怎麼給意見那個家長.

$^{121}$怎麼教兒子.
我也是做完年青人不是很久.
不是很久.
只是小小的時間.
怎麼說呢.
可能你也會反過來問我.
不是吧.
你不是有上帝的呼召嗎.
不是吧.
你不是讀過神學的嗎.
不過做了這麼多年.
不是新手的時候才敢這樣說.
其實傳道人有很多事情要做的時候.
是神學院沒有教的.
所以有很多事情真的不懂得怎麼做.
可能當時約書人會不會也是這樣呢.
就是做了摩西的助手很多年.
但是對於打仗來說.
他是完全沒有經驗的.
也都沒有學過的.
所以對於自己這個新的工作崗位.
新的工作內容就覺得很害怕.
我怎麼懂呢.
不過如果你看出埃及記的時候.
你就會發覺應該不是這樣.
因為約書人在出埃及記已經出現過.
在出埃及記的十七章.
當時亞馬利人正在攻打以色列.
然後摩西就在山上舉手.
手一舉起.
以色列就會得勝.
垂低的時候亞馬利人就會得勝.
記不記得這個經過.
當時帶兵出戰的是誰呢.
就是約書人.
所以似乎約書人也有相當的軍事經驗.
或者作軍隊的首領.
這個未必是他擔心的.
會不會擔心的只是沒有了摩西.
他就真的要獨當一面.

$^{161}$這個可能是其中一個原因.
既然未必是對工作內容這麼多的恐懼.
有學者就會去推理.
會不會約書人的恐懼.
其實是對於十二支派能否合一.
能否團結的恐懼呢.
第一章有一個現象值得我們留意.
作者不斷強調約書人是神所揀選的.
是接替摩西的領袖.
包括他說如何和摩西同在.
就會和你同在.
正如你從前聽從摩西.
也會聽從你.
這些說話是不斷重複.
好像是反映著某一種的不穩定.
因為有學者推測.
可能約書人的壓力未必是軍事上的挑戰.
而是如何在內部能夠穩定.
特別是劉便仁 加德仁和馬來西那半支派.
為什麼要特別提他們呢.
在民數記32章和申命記第3章都提及這件事.
他們在過約旦河東之前.
曾經向摩西承諾.
我會和其餘的支派一起過河作戰.
直到他們得到產業.
為什麼要下這樣的承諾呢.
原因是他們在未過河之前.
已經拿到他們要的地.
他們說我們要定居在這個地方.
所以這個好像是一個隱憂.
現在摩西死了.
而在古代世界就是政權交接.
一個最容易出現.
那我就可以走數了.
那我就可以不理了的時間.
會不會有人會想.
我們都拿了地.
為什麼我還要過河和你冒這個風險打仗呢.
所以約書亞記1章12-15節特別提到.
約書亞要提醒他們.

$^{201}$摩西曾經和你們這樣說.
摩西曾經這樣吩咐你們.
這個不是一個溫馨提示.
出了一張面膜提醒你.
記不記得要去完西涌之鴨.
不是一個這樣的溫馨提示.
而是一個很重新約定.
記住這是一個很重要的立約.
所以我們推敲會說.
約書亞會不會也擔心.
他能不能夠坐穩領袖的位置呢.
始終一定不會在十二個支派裡面.
只有他一個領袖的才能.
總有其他的人都有那種領袖的素質.
可能他也會想他能不能帶領百姓的合一呢.
會不會有些人打著打著就說.
我站在這裡就算了.
我不會再走的了.
我不再打了.
同樣他也會想.
會不會因為他成為領袖.
這個領袖的交替的時間.
就會出現很多的分化.
很多的人.
大家的支派就各自有各自要做的事呢.
的確今天我們進入一個事奉.
也不是很願意.
也不是很輕易.
我們會接受他成為一個領袖的崗位.
所以部門的委員.
或者是團長職員會.
這些永遠都是最難邀請人去參加的.
我記得有好幾次邀請大英姐妹去事奉.
首先就說一下.
這個事奉的崗位是做什麼的.
我們有什麼想法.
有什麼方向.
有什麼意象.
然後大家都談得很投契.
也都覺得.

$^{241}$不是的 真的是這樣.
上帝應該要這樣帶領我們.
也都覺得自己能夠勝任去做到這件事.
但是一說到.
這個崗位都有個身份的.
有個名字的.
這個崗位叫做團長.
那些人馬上就像見鬼一樣.
馬上就.
不行不行不行.
千萬不要叫我做.
耍手擰頭.
經常都聽到大英姐妹說.
我是小薯仔來的.
我只可以做跟班的.
因為我們誰不知道領袖不容易做呢.
可能做決定的時候.
必定可能會觸動到一些人.
然後領袖就成天都要在這些原則.
和關係之中有很多的掙扎.
又或者大家都會說.
好啊 做啊 做啊.
但是真正行動的時候.
不見了人.
那怎麼辦.
領袖就要承擔起那個壓力.
可能又要出人又出力.
出誰人誰力呢.
就是自己和自己的家人.
老公就推他出來.
喂 幫忙吧.
你不要坐在那裡.
然後做都不是最難.
就是當你發現.
原來團隊可能都有人都在質疑你.
甚至在背後有很多討論.
傳回你耳邊的時候.
可能那個才是更令人氣餒的事.
或者當時約書亞就是背負起這些壓力.
所以當約書亞知道要過約但何.

$^{281}$然後他都要特別重申.
摩西和這些支派所立的約.
這些承諾.
然後當這些支派領袖去回應的時候.
說自己會忠誠於摩西的承諾.
特別要加一句.
你只要剛強壯膽.
彷彿要和約書亞說.
我們認可你是領袖.
你放心.
我們一定會繼續和你去拿迦南的地.
除了在團隊十二支派的恐懼.
對於一些學者.
甚至會從約書亞的年紀.
去推敲他所擔心的.
今天我們知道的是.
約書亞應該和迦勒是同輩.
雖然我們不知道約書亞接任領袖的年紀.
但我們知道當年他們進入迦南地規探的時候.
迦勒是四十歲.
進入迦南之後.
迦勒大概是八十五歲.
假設如果約書亞和迦勒的年紀相當.
我們就可以想到約書亞應該是八十歲的年紀.
去接任領袖.
學者就會想.
會不會對於約書亞來說.
他其中一個擔心就是.
自己的能力不及當年的勇氣.
也沒有當年的力量.
所以令他感到恐懼.
我想這些都是很合理的估計.
正如我一開始這樣說.
聖經裡面沒有一個很明確說.
約書亞的恐懼是什麼.
所以我們根本沒有辦法去確定.
可能你這一刻有點想反鄧.
即使你拿了十多分鐘說.
但說完之後你說沒有結果.
雖然我們沒有辦法去確定約書亞在怕什麼.

$^{321}$但我們可以從他的處境裡面.
去看到我們所恐懼的.
或者可能我們對事奉的內容的恐懼.
這個事奉太難了.
我根本不懂得做.
我做不到.
又或者我怎樣解決這麼多困難.
有這麼多問題.
這麼多炸彈要拆.
我怎樣搞定.
又或者說要安慰人.
最怕就是安慰人的工作.
或者可能那個恐懼來自於事奉的團隊.
就是我很怕意見不合的時候.
那些氣氛很尷尬.
那些責任會不會全部都落在我身上.
最後什麼都要我做.
又或者說我都曾經在團隊的服務裡面受過傷害.
我就不想再經歷一次.
我不想再受這些傷害.
同樣可能恐懼是來自於自己的能力.
就是覺得自己都有很多的掙扎.
怎樣去參與在服務的裡面.
怎樣去服務其他人.
帶領其他人.
又甚至可能覺得我都沒有這個信心.
或者我的屬靈生命是否真的足夠.
我真的信主不久.
或者我還有很多不足的地方.
上帝對約書亞的呼召是很獨特.
這個是我們要留意的.
雖然上帝不是呼召我們每一個人都要像約書亞一樣成為領袖.
但是面對著不同的事奉崗位的感動.
或者我們都面對著種種不同的擔心.
就算見到上帝對我們的吩咐.
但是我們都不敢去回應.
因為我們有很多恐懼.
但之後經文就要說的就是.
上帝怎樣讓到約書亞剛強壯膽呢.
上帝不是命令約書亞不准害怕.

$^{361}$而是要在他的吩咐裡有兩個讓他能夠剛強壯膽的原因.
第一個是因為上帝的應許仍然有效.
上帝向約書亞表明自己的計劃是沒有中斷的.
在一章五節裡說.
你一生的日子必定沒有人能夠在你面前站立得住.
這句說話我們聽的時候就覺得很安慰.
但對於約書亞來說.
他聽到的不是一個安慰.
是一個代表以色列人和上帝所納的約的部分.
因為在《申命記》裡上帝用同一句說話.
在七章二十四節裡和以色列人立約.
必定沒有一個人能夠在你面前站立得住.
所以這不是一個新的安慰 一個新的應許.
而是上帝和以色列人的約.
他要告訴約書亞是會繼續延續的.
就是摩西死了之後 上帝都沒有改變.
領袖好像在換.
但是上帝的救贖歷史不會因為這樣而中斷.
因為以色列人真正的領袖從來都不是摩西.
也不是約書亞 而是耶和華自己.
所以上帝第一件要告訴約書亞的事.
就是他和以色列所納的約是沒有終止的.
就是當日應許要賜給他們的地.
今天仍然有效.
所以要讓以色列人知道上帝仍然在掌權.
神仍然是守約的上帝.
神仍然會親自帶領他的百姓.
因此約書亞就可以靠著上帝去剛強壯膽.
上帝只是要將那塊地賜給他們.
要留意的是上帝不是說你要為自己去爭取.
人總要為自己去拿一些東西去爭取那塊地.
他不是這樣說.
而是說我只是要將那塊地賜給你.
在約書亞記裡面有一個非常重要的訊息.
就是以色列人的角色是甚麼.
以色列人的角色不是去捕捉一些.
上帝我做到這裡是這麼多的.
其他東西要靠你自己我做不來.
這不是我們認識的那個上帝.
也不是要靠自己的能力去創造一個怎樣的歷史出來.

$^{401}$他們是立約的群體.
在上帝的主權之下順服地行動.
所以神是賜那塊地給他們的主.
他們是承受產業的人.
同樣神是元帥.
這班以色列人是順服去出戰的百姓.
所以剛強壯膽從來不是說那種自信.
就是你要有自信你要對自己有信心.
不是這樣.
覺得要有那份自信可以為上帝做甚麼的成績.
反過來是你要對上帝的主權有那種信靠.
所以信心不是來自自己.
是來自對於上帝.
同樣今天事奉都不是說我一定要成功.
如果我這次搞這個聚會搞這些活動不成功.
那就完蛋了.
而是神若應許的時候.
上帝使用這件事就能夠成就.
以色列人不是靠自己去打那個應許之地.
他們是在上帝的應許之中前行.
所以今天我們也不是為自己爭取一個怎樣的地方.
也不是擔心自己做不做得出甚麼成績.
覺得自己會拖累上帝的事工.
或者認為要為上帝爭取到一個怎樣的地步.
因為這樣都是代表著.
我們以為上帝需要我們去補足一些他做不來的事.
反過來我們要剛強壯膽的.
就是要信服在上帝的大令裡.
做上帝要我們去做的事.
除了在第一個上帝說應許仍然有效.
第二個對他能夠剛強壯膽的原因.
是因為上帝和約書亞同在.
除了上帝的應許沒有改變.
上帝的同在也沒有改變.
他在一章第五節裡說.
我怎樣與摩西同在都必定與你同在.
我不會撇下你也不會丟棄你.
對約書亞來說這個同在是很實在的.
因為他正正就是摩西的助手.
他看著上帝怎樣與摩西同在.

$^{441}$所以今天同樣的應許.
是會與約書亞同在.
是他見過的 是他經驗過的.
很多時候我們在那個懼怕.
在那些很多擔心的裡面的時候.
都會將上帝在我們曾經在我們生命裡.
或者在我們身邊弟兄姊妹.
身邊做過的一些恩典的那些同在.
我們都忘記得一乾二淨.
而上帝就是要這樣去提醒約書亞.
上帝的同在是沒有間斷的.
就是以前是怎樣 現在都是這樣.
然後叫約書亞能夠剛強壯膽.
就是因著上帝的應許和同在.
約書亞就知道自己的角色.
就是要忠心信服在上帝的帶領裡面.
而不是要幫上帝去爭取什麼.
要成為一個怎樣的英雄.
因為忠心信服是他可以做到的.
所以他就能夠剛強壯膽.
好了 雖然我們知道了.
怎樣能夠剛強壯膽.
我們要解決最後一個問題.
就是什麼才是剛強壯膽呢.
怎樣是剛強壯膽呢.
在一章七到八節那裡這樣說.
只要剛強大大壯膽.
謹守進行我僕人摩西所吩咐的一切律法.
不可偏離左右.
使你無論往那裡去都可以順利.
在這個律法書裡面不可以離開你的口.
總要就夜思想.
很使你謹守進行這個書上面所寫的一切話.
原來神對剛強壯膽的解釋.
不是說你要情緒很穩定 很冷靜.
面對什麼的時候都知道不用怕.
不是這樣.
也不是說你要對自己有自信.
你的性格一定要夠強勢夠堅硬.
不是這樣.

$^{481}$而是在這一切的行動裡面.
你要忠於在上帝的律法裡面.
當你讀過《申命記》.
你就知道上帝有很多的律法.
和當時迦南七族的規定完全不同.
神只是要建立一個屬他的角度.
一個有公義有憐憫.
而且是一個神治的社會.
所以對約書亞來說.
要將這些規矩.
就是從前沒有的.
要一一好好記住.
甚至要就夜思想.
才不會遺漏.
其實記得都未必是最難.
最難是什麼呢.
最難可能是要實行出來.
對約書亞來說.
剛強壯膽的意思就是.
即使面對戰爭的壓力.
但仍然照著神律法去行事.
例如在耶利哥城以外.
已經過河了.
然後在耶利哥城門外.
他們要做些什麼呢.
神竟然吩咐這些以色列人要行割禮.
在人看來這是最不安全的時刻.
敵人隨時出來.
我們全部人都是受傷的狀態.
行動不便的狀態.
怎麼可能在那個時間打仗呢.
但在這個最不安全的狀態裡面.
他們就要馴服.
這個是剛強壯膽.
這才是約書亞真正需要的勇氣.
因為剛強壯膽不是要衝鋒陷陣.
而是在不安全或違背常識的情況下.
仍然要忠於上帝的話.
今天進入事奉裡.
我們都帶著很多恐懼.

$^{521}$可能會擔心自己的能力.
擔心能否升任服侍.
甚至擔心團隊能否合作起來.
一起服侍.
這些擔心都很真實.
可能都是當時約書亞所擔心的事.
但上帝叫約書亞要剛強壯膽.
不是說他要怎樣有自信.
要相信自己能夠做得到.
然後給自己一個擁抱.
不是說這件事.
而是約書亞知道.
真正他需要的是敬畏上帝.
他要有勇氣的.
是怎樣忠於上帝的吩咐和律法.
無論在什麼處境都要.
因為他的角色不是要成為以色列的英雄.
要完成上帝未能完成的.
將迦南地爭取給自己的民族.
約書亞真正的角色.
是要成為一個忠心的僕人.
按上帝的吩咐去做.
然後領受上帝為他們所預備的地方.
今天我們所擔心的.
又是否真是我們能夠處理得完呢?.
但今日經文提醒我們.
我們都不需要成為教會的英雄.
我們不需要帶著重擔.
覺得我做不到.
我只是在虧欠上帝的榮耀.
也不需要認為.
我們要為上帝做出一個怎樣的成績.
因為今日上帝要我們的耿耿.
只是忠心於他.
然後按著他的吩咐而做.
但你要留意的是.
今日不是說.
Kevin就說不用做了.
千萬不要這樣說.
我們說的是我們仍要做.

$^{561}$但為了什麼而做呢?.
今日我們仍然要傳福音.
但我們要知道.
拯救人的是上帝.
改變人心的是上帝.
同樣我們仍然要使人作門徒.
但我們要知道.
改變人的都是上帝當中.
同樣今日我們要忠心的服侍.
但我們要知道.
真正建立教會的主權.
是在主上帝的裡面.
所以我們只需要剛強壯膽.
按著上帝的吩咐盡力而為.
什麼叫盡力而為呢?.
就是可能沒有人看到.
但我們仍然做.
果效不大.
但我們繼續做.
因為我們做的原因不是被人看到.
不是為了有什麼結果.
而是我們要敬畏我們的上帝.
所以最後很想用耶穌的一句提醒.
成為我們今日的一個結束.
在《馬太福音》上有一章三十節.
耶穌這樣對他的門徒說.
我的額是容易的.
我的擔子是輕散的.
就讓我們放下.
只是上帝才能夠揹的擔子.
我們揹不起的.
我們放下它.
因為那個是上帝負責的.
我們擔起我們能夠應付.
耶穌基督給我們的額.
讓我們一同禱告.
因為你是憐憫我們的主.
你是明白我們需要的主.
所以你為我們去準備一個.
我們能夠應付的額.

$^{601}$就是忠心在你的律法.
忠心在你的教導當中.
求主讓我們有這個勇氣.
在我們想不跟著你的旨意而行的時候.
求主你去教導我們.
讓我們有勇氣去看見.
我們只要大大放膽去跟從你的時候.
其餘的一切就由你去負責.
上帝就求你讓我們有這個心.
奉著耶穌的性命祈求.
阿們.
\newpage



\section{羅馬書 5:12-17}
\label{sec:031682}
\textbf{不對稱的救贖}~~~~林玉英牧師
\newline
\newline
連結: \href{http://www.church.com.hk/acms/content.asp?site=cdc\&op=show\&type=product\&code=031682\&layout=sermon}{\texttt{031682}} ~~~~ 語音日期: 2026-02-21
\newline
\newline
\hyperref[sec:031681]{\small{< < < PREV SERMON < < <}}
~
\hyperref[sec:index_chronic]{\small{[返順時目]}}
~
\hyperref[sec:031683]{\small{> > > NEXT SERMON > > >}}
\newline
\newline
羅馬書 5:12-17
\newline
\begin{longtable}{cl}
\hline
\hline
章節 & 經文 (和合本修訂版)\\
\hline
5:12 & \begin{tabularx}{0.7\textwidth}{X} 為此，正如罪是從一人進入世界，死又從罪而來，於是死就臨到所有的人，因為人人都犯了罪。 \end{tabularx} \\ \\ \relax
5:13 & \begin{tabularx}{0.7\textwidth}{X} 沒有律法之前，罪已經在世上，但沒有律法，罪也不算罪。 \end{tabularx} \\ \\ \relax
5:14 & \begin{tabularx}{0.7\textwidth}{X} 然而，從亞當到摩西，死就掌了權，連那些不與亞當犯一樣罪過的，也在死的權下。亞當是那以後要來之人的預像。 \end{tabularx} \\ \\ \relax
5:15 & \begin{tabularx}{0.7\textwidth}{X} 但是過犯不如恩賜，若因一人的過犯，眾人都死了，那麼，神的恩典，與那因耶穌基督一人而來的恩典中的賞賜，豈不加倍地臨到眾人嗎？ \end{tabularx} \\ \\ \relax
5:16 & \begin{tabularx}{0.7\textwidth}{X} 因一人犯罪而來的後果，也不如賞賜，原來審判是由一人而定罪，恩賜乃是由許多過犯而稱義。 \end{tabularx} \\ \\ \relax
5:17 & \begin{tabularx}{0.7\textwidth}{X} 若因一人的過犯，死就因這一人掌權，那些受洪恩又蒙所賜之義的，豈不更要因耶穌基督一人在他們生命中掌權嗎？ \end{tabularx} \\ \\
[1ex]
\hline
\hline
\end{longtable}
$^{1}$各位弟子平安.
今天是農曆年初五.
在這裡恭祝大家新年進步.
身體健康.
上帝的恩典每天都臨到大家當中.
更加祝願大家.
學武神和信靠神的心.
每天都會增長.
讓我們首先一起低頭祈禱.
增長在新的一年 中國人的新年裡面.
主要是將我們每一個特慶節目交在您的手中.
求您的恩典時刻與我們同在.
就像今天在經文裡所說的.
您的恩典浩大 百百次給我們.
求您繼續來賜恩 去帶領我們生命的腳步.
願說的和聽的都同蒙您的月亮.
祈禱奉祖耶穌基督的名孝.
阿們.
我相信大家也有用過電腦.
牧師這句說話簡直是多餘.
手提電話更加不再說了.
每天都在手裡.
所以如果你用過電腦.
也說每天都在用手提電話的時候.
你大概都知道什麼叫做預設設定.
Default Settings.
大概都知道了吧.
原廠會幫我們設定好.
除非您更改電話或電腦的指令.
否則手提電話會根據原廠設計師的設定.
來運作手提電話.
例如手提電話已經預設有追蹤功能.
不知道您是否覺得.
我們的生命都好像裝了預設系統.
有些事情我們好像裝了魔一樣,不能不做.
我說的不是預定論.
而是保羅在這裡所說的就是亞當的系統.
是什麼來的?.
我們來看看保羅在《羅馬書五章》十二至十七節.
他告訴我們什麼叫做亞當的系統.

$^{41}$在十二節裡面他這樣說.
正如罪是從一人入了世界.
死又從罪而來.
於是死就臨到眾人.
臨到所有的人.
因為人人都犯了罪.
保羅在這裡所說的死.
不單只是我們每一個人去到某一個時間.
你的身體就會死亡.
你的肉體就會死亡.
其實他在這裡所說的都包含了我們生活中感受到的一份孤竭和孤乾.
在我們這個亞當系統裡面.
不知道大家覺不覺得.
感覺這個世界好像在一個靈活遊戲裡面.
競爭很激烈.
好像在玩淘汰賽.
不是你死就是我亡.
只可以留下一個.
於是大家為了生存.
我們要學辦公室政治.
我們要學人造我們的弱點.
當我們的上司或GM.
為了個人要表現向上爬.
他可能會犧牲他的下屬.
就好像這段經文所說.
一個人犯罪.
死就臨到眾人的現代版.
我們會發現在這個系統裡面.
無論我們多努力.
內心深處開始對人的信任.
對工作的熱情.
正在慢慢地枯竭和死去.
剩下什麼呢?.
就是對薪水的渴望和依賴.
令我們時時不自覺地繼承了我們原生家庭的遺傳.
例如可能是習慣了冷暴力.
習慣了過度的要求.
又或者用愧疚感來情緒勒索我們的下一代.
我們明明不太想成為像我們的父母那一種人.
不過不知為何在吵架的時候.

$^{81}$發現自己說了一些跟他們一模一樣的話.
就好像保羅在第十四節所說.
死就掌權了.
連那些不跟亞當犯一樣罪過的.
都在死的權下.
沒錯,死亡支配了全人類.
甚至那些不像亞當那樣犯了不服從上帝的罪.
那些人都不免要死.
我們都在一個錯誤的承繼裡面.
DNA全都進血了.
做完我們想做好人卻身不由己.
更甚的就是在我們的生命裡面有一份不安全的感覺.
有一份恐懼和焦慮感.
不知道你會不會覺得你很怕做錯事呢.
擔心自己走錯一步.
就好像骨牌效應那樣全軍覆沒.
好像亞當連累了整個世界.
在另一方面,如果我們真的做錯事的時候.
我們就遺傳了亞當的逃避.
躲起來否認.
又或者推到別人那裡說:不關我的事.
是他的錯而已.
在這個亞當系統裡面,我們越走越枯乾.
越走越沒力.
和人的關係越來越疏離.
越有隔膜.
彼此越來越沒了信任.
保羅說:我真是苦.
誰能夠脫離這個取死的身體呢?.
可以怎樣去脫離亞當的預設系統裡面呢?.
甚至他這樣說.
他說:但是過犯不如恩賜.
如果因為一人的過犯眾人都死了.
那麼上帝的恩典.
和因為耶穌基督因人而來的恩典之中的賞賜.
豈不加倍的臨到眾人嗎?.
聖經文給我們很大的鼓勵.
因為上帝藉著基督的修復力量.
遠遠大於亞當這個系統帶來的破壞力.
這個保羅稱之為基督恩典的系統.

$^{121}$是上帝對人一個不對稱的救贖.
固然有很多人因為亞當一個人的罪而死.
不過上帝的恩典更加豪大.
他藉著耶穌基督一個人白白賜給很多人的恩典.
也更加豐盛.
不知道大家在心底的深處.
特別是香港這段時間.
你會不會最怕失去工作呢?.
又或者達不到我們公司定下來的績效指標KPI.
又或者我們不知不覺間想達到別人對我們的期望.
在亞當的系統裡面不達標.
即是失敗.
即是等於你沒有價值.
不過在基督的恩典系統裡面.
保羅說上帝的恩典是加倍地賜給我們.
是白白賜給我們.
這意味著上帝對我們的接納,拯救和扶持.
遠超過由我們的工作做了甚麼.
以及由其他人對我們的評價來定義.
就算我們做得不夠好 上帝都愛我們.
當你意識到你的價值不是和你達不達標去掛鉤的話.
你突然間會發現自己可以輕鬆了.
而你會將你所有的精力.
將你的心神放在怎樣做好事上.
你反而更能夠展現上帝在你裡面.
上帝一早給了你的創意和天賦.
為甚麼會這樣.
因為你不怕輸.
你也輸得起.
或者你曾經在你家裡面犯過嚴重的過錯.
又或者你家裡只是處於破碎的邊緣.
保羅說「過犯不如恩賜」.
這句話對我們現代人來說.
就是上帝重建的能力.
大過我們曾經造成過的傷害.
沒錯 我們曾經是被定罪.
不過只要我們願意切換去到基督的恩典系統裡面的說話.
它能夠指教我們怎樣去修復和家人的關係.
給力量我們去道歉.
去寬恕.

$^{161}$甚至打破幾多人傳下來的詛咒.
現在我和成恩宣教師.
丙堂偉和師長一起正在上課程.
這個課程叫做「門徒生命轉化旅程」.
在教會裡面我們同時也在開辦這個課程.
因為我們一邊讀一邊教.
這個時候我們上的課是學習「針搖樹」.
靠著上帝的力量去修復和其他人的關係.
得到真正的彼此要恕.
其中有一個例子.
是一個女兒在很小的時候被她的親爸爸侵犯.
幾十年都過去了.
彼此當然是陌路人.
女兒也將這件事封塵在內心最深的地方.
但因為她想門徒的課程「門旅客」.
所以她要學習「針搖樹」.
在牧者的幫助下.
她重新去回顧這段傷害.
是真的不容易.
回想起來還是很痛.
這是一個很漫長的過程.
不過當事人真的很想去認真按照聖經的指引.
讓自己的心得到真正的平安和痊癒.
在整個漫長的過程裡.
她慢慢看到自己長大後如何報復父親.
在適合的時候.
女兒覺得自己再回想當日的事情時.
她沒有再有情緒反應.
而幫助她的牧者也發現這是合適的時候.
她就將她的父親出來見面.
彼此再說回當日的這件事.
她看到她的父親痛哭.
向自己的女兒道歉.
而女兒也向父親道歉.
因為她長大後.
用她可能的方法去報復父親.
在他們兩人彼此真正寬恕對方後.
他們之間的芥蒂也慢慢地清除.
大家可以去修復大家的關係.
大家聽完我講之後好像覺得很容易.

$^{201}$但其實真妖術真的不容易.
過程是很漫長和艱難.
所以如果大家在這之中.
你曾經受過很大的傷害.
請你不要貿然自己去處理.
請你去找教睦同工.
幫助你一步一步去達到這個目標.
真妖術真的不容易.
過程也很漫長和艱鉅.
不過在基督浩大的恩典幫助下.
人與人之間的關係其實真的可以修復.
並且開始重建大家的關係.
特別是信任的關係.
達到可以像聖經所講的.
真正的彼此相愛.
十七字說.
如果因為一個人的犯罪,過犯,死.
就因為這個人去掌權.
那些受弘恩又蒙所賜之義的人.
是否更要因為耶穌基督一人.
在他們的生命裡去掌權呢?.
各位靈姐妹.
你有沒有想過我們是可以.
在我們自己的生命裡去掌權,作王的呢?.
拿保祿的意思.
他說作王不是擁有特權.
而是擁有生命的主動權.
例如我們看看職場裡.
當所有人都在背後去議論某一個同事.
又或者為了升職去不擇手段的時候.
又或者你看到有些人有功勞自己領了.
有過失就推給其他人去受.
不過因為我們切換了基督恩典的系統.
我們領受了上帝給我們的這份義.
我們的安全感.
再不是去看我們老闆的面色.
又或者別人的面色.
我們的安全感.
是建立在上帝對我們的接納.
在那個時候.

$^{241}$你會選擇不參與卑鄙.
我們可以選擇成為鼓勵其他人的人.
或者是建立其他人的人.
這就是「作王」的意思.
我們要在我們的職場當中去作王.
因為周圍的環境.
我們身處的處境不能夠決定我們的品格.
是我們去帶領環境.
不是環境帶領我們.
是我們決定怎樣去處事.
怎樣去面對一些事情.
又或者在家庭裡面我們怎樣去作王呢?.
可能你結婚了就會知道.
當我們的另一半.
可能受了很大的壓力.
突然情緒爆發的時候.
又或者我們的小朋友挑戰我們的底線的時候.
我們再不是被亞當的系統去主導.
或者欺騙.
覺得自己是受害者.
於是跟著和對方去對罵.
因為我們裡面有一份.
上帝給我們的雄恩去撐著我們.
讓我們可以選擇很溫柔地去回應.
這就是「作王」.
我們擁有生命的主動權.
我們不再被情緒去統治.
而是可以靠著基督聖靈去幫助我們去主導我們的情緒.
將平安帶回我們的家.
我相信有小朋友的家長.
或者你都會試過.
其實我都試過.
在我們的小朋友去挑戰我們的底線的時候.
有時候我們情緒會失控.
我記得我都試過將書包.
扔進垃圾桶.
我都聽過有一個育兒專家.
他將他小朋友的書包從三樓扔到街上.
當然.
冷靜之後.

$^{281}$我們都會覺得很後悔.
為什麼我們會做出這樣的事.
因為那一刻.
我們就被亞當的系統所主導.
是我們要在家裡作王.
成為家人情緒穩定的基石.
我們不再被死亡亞當的系統去支配.
而是由基督恩典的系統所支配.
所以各位弟兄姊妹.
《天地經文》告訴我們有一個終極的釋放.
我們的人生不需要再哭哭啼啼或者死去.
因為我們不是為生存去掙扎.
而是在聖靈的幫助下.
在我們的生命裡面去作王.
如果我們此刻都覺得我們的生活很孤單.
信仰覺得是可有可無的時候.
請你記得.
順利給我們的恩典不是一張保險的單.
買了就放在這裡.
順利的恩典是一個引擎.
是一個動力的引擎.
是一個engine.
只要你祈求順利會給我們源源不絕的動力.
我們可以從亞當的死亡系統.
去切換到基督的恩典系統.
這是順利給我們不對稱的救贖.
我們沒有做過任何事.
所以不能等待交換.
是一個不對稱的救贖.
我們不再需要自己的努力.
去和上帝等待交換.
我們只要去祈求.
只要願意相信.
我們可以從舊的系統.
去切換到新的系統當中.
無錯,我們仍然會犯錯.
會搞壞一些東西.
不過上帝的恩典很豪大.
他時時會願意接納和寬恕我們.
並且賜機會給我們再來過.

$^{321}$再試過,再做過.
所以當你星期一.
我不知道星期一大家是否要上班.
或者今晚你回到家的時候.
覺得很大壓力.
和要遇到衝突的時候.
請你在心裡首先祈禱.
上帝幫助我記得我現在不是在亞當的系統裡面.
我們是受了洪恩的人.
已經切換到基督的恩典系統裡面.
所以上帝.
我們知道這是你給我們不對稱的救贖.
這個恩典是你賜給我們的.
幫助我在面對這件事情的時候.
靠你作王.
以致我可以榮耀你的名.
由我們同心祈禱.
親愛的上主.
我們的天父.
多謝你透過今天的經文告訴我們.
我們一直受了撒旦的圈套.
以致我們的生命仍然繼續在亞當的系統裡面.
很多時候我們生活得很累很枯渴.
以為人生就是一場不是你死就是我亡的靈活遊戲.
我們在職場裡面很沒有安全感.
在家庭當中不自覺地重複上一代的傷害.
就好像落入了一個改不了的指令.
這個預設系統裡面.
身不由己.
因為你的救贖是不對稱的.
雖然亞當一個人犯罪.
令到死亡是掌權.
但是你的恩典更加豪大.
我們多謝你告訴我們.
過犯不如恩賜.
我們的失敗.
我們的不達標.
我們過去所造成的傷害.
都不能阻隔你對我們的愛.
因為基督耶穌的恩典是加倍,是白白.

$^{361}$是足以覆蓋我們所有的缺點.
所以上主.
我們今天要在你的面前按一下.
去按這個切換制.
由我們切換到基督恩典的系統裡面.
要在我們的家庭.
在我們的職場裡面去作旺.
當我們在公司面對壓力.
又或者在家裡面對衝突的時候.
我們求聖靈去提醒我們.
我們不是再在亞當的系統當中.
成為一個受害者.
而是一個蒙受紅恩的人.
求你賜我們力量去拒絕卑鄙.
賜我們溫柔去化解衝突.
上主我們不再靠自己去死撐.
而是靠著你這個動力引擎.
讓我們每一天起身的時候.
都帶著這份不對稱的救贖.
在生命裡面和你一起作旺.
我們這樣透過藉主耶穌.
基督的名而求.
阿們.
\newpage



\section{馬太福音 4:1-11}
\label{sec:031683}
\textbf{如何對付那試探者——魔鬼撒旦}~~~~高國雄牧師
\newline
\newline
連結: \href{http://www.church.com.hk/acms/content.asp?site=cdc\&op=show\&type=product\&code=031683\&layout=sermon}{\texttt{031683}} ~~~~ 語音日期: 2026-02-22
\newline
\newline
\hyperref[sec:031682]{\small{< < < PREV SERMON < < <}}
~
\hyperref[sec:index_chronic]{\small{[返順時目]}}
~
\hyperref[sec:031684]{\small{> > > NEXT SERMON > > >}}
\newline
\newline
馬太福音 4:1-11
\newline
\begin{longtable}{cl}
\hline
\hline
章節 & 經文 (和合本修訂版)\\
\hline
4:1 & \begin{tabularx}{0.7\textwidth}{X} 當時，耶穌被聖靈引到曠野，受魔鬼的試探。 \end{tabularx} \\ \\ \relax
4:2 & \begin{tabularx}{0.7\textwidth}{X} 他禁食四十晝夜，後來就餓了。 \end{tabularx} \\ \\ \relax
4:3 & \begin{tabularx}{0.7\textwidth}{X} 那試探者進前來對他說：「你若是神的兒子，叫這些石頭變成食物吧。」 \end{tabularx} \\ \\ \relax
4:4 & \begin{tabularx}{0.7\textwidth}{X} 耶穌卻回答說：「經上記著：『人活著，不是單靠食物，乃是靠神口裡所出的一切話。』」 \end{tabularx} \\ \\ \relax
4:5 & \begin{tabularx}{0.7\textwidth}{X} 魔鬼就帶他進了聖城，叫他站在聖殿頂上， \end{tabularx} \\ \\ \relax
4:6 & \begin{tabularx}{0.7\textwidth}{X} 對他說：「你若是神的兒子，就跳下去！因為經上記著：『主要為你命令他的使者，用手托住你，免得你的腳碰在石頭上。』」 \end{tabularx} \\ \\ \relax
4:7 & \begin{tabularx}{0.7\textwidth}{X} 耶穌對他說：「經上又記著：『不可試探主—你的神。』」 \end{tabularx} \\ \\ \relax
4:8 & \begin{tabularx}{0.7\textwidth}{X} 魔鬼又帶他上了一座很高的山，將世上的萬國和萬國的榮華都指給他看， \end{tabularx} \\ \\ \relax
4:9 & \begin{tabularx}{0.7\textwidth}{X} 對他說：「你若俯伏拜我，我就把這一切賜給你。」 \end{tabularx} \\ \\ \relax
4:10 & \begin{tabularx}{0.7\textwidth}{X} 耶穌說：「撒但，退去！因為經上記著：『要拜主—你的神，惟獨事奉他。』」 \end{tabularx} \\ \\ \relax
4:11 & \begin{tabularx}{0.7\textwidth}{X} 於是，魔鬼離開了耶穌，立刻有天使來伺候他。 \end{tabularx} \\ \\
[1ex]
\hline
\hline
\end{longtable}
$^{1}$早晨 各位電影姐妹平安.
多謝師班帶領我們的禱告.
讓我們再一次在主面前獻上禱告.
天父上帝我們獻上感恩.
是你的恩典.
讓我們由歲首到年終都經歷生命的豐富.
求你繼續帶領賜福.
繼續向我們說話.
以讓我們生命裡面有你的話語.
成為我們的指引.
成為我們的燈.
成為我們的光.
以證我們每一步都不偏離.
看我們的耳朵.
讓我們留心聆聽.
看我們的眼睛.
讓我們見證到你的榮美.
猶如僕人口中的言語.
弟兄姊妹之間的心思意念.
蒙聖靈的啟迪.
蒙主的月立.
禱告奉耶穌基督的名求.
弟兄姊妹今天是大年初六.
所以在這裡祝賀每一位.
新年一日我們能夠經歷上主更多的恩典.
身心靈健壯.
能夠不單止龍馬精神.
而是在主的恩典當中.
我們的靈也是每天這麼精神的.
以致能夠領受上主更多的福樂.
更多的恩惠.
今天同樣都是教會大齋期.
第一個主日的開始.
大齋期是什麼意思呢.
大家可能都開始慢慢要去熟悉多一點了.
大齋期當中我們多一點默想.
你看到今天的詩歌.
土文.
我們都是在一個比較靜的時候.
多一點學習刻符.

$^{41}$學習詩社.
學習在一個簡樸的生活當中.
每年去到這裡.
很多弟兄姊妹都常常有很大的難處.
或者矛盾衝突.
明明是新年.
要慶祝的.
今年的大齋期開始.
大家知不知道是年初幾.
剛剛禮拜三.
初二.
中國人說新年了.
開年了.
就開始要學習獻璧.
學習走苦的路.
好像很多的衝突很矛盾.
每年都是這樣.
其實不像是今年.
都是大年初幾的日子.
差不多就是大齋期開始.
明明開始慶祝兩天之後.
教會就告訴你.
我們不要慶祝了.
我們要多一點去學習受苦.
我們吃東西要清淡一點.
禁食禱告.
好像很難拼在一起.
其實這個可能是上帝給中國人特別的恩典.
讓我們當中更深層的思考.
這個好像拼在一起的狀況.
其實就是信仰最實際的狀況.
信仰從來都不是順順利利.
開開心心.
快快樂樂.
好ye .
就過你一世.
肯定不是.
當中年長的弟兄姊妹.
或者信仰經歷得深的弟兄姊妹.
你會更明白.

$^{81}$一定不是.
今日經文正正就是講這件事情.
今日經文就是講耶穌受魔鬼的試探.
其實你應該很熟的.
因為每年去到這個時候.
都是講這些經文的.
可能我在這裡都講了很多次.
耶穌受魔鬼試探.
大家都很熟的.
這段經文第一個我要問的問題就是.
為什麼耶穌要受試探?.
你每個禮拜去教會.
剛才蘇雪教師帶我們去祈禱.
我們背主禱文的.
主禱文有一句是什麼?.
不叫我們.
就是了.
我們明明祈禱求上帝.
叫我們拿走試探.
不要見到.
這個是耶穌教我們的.
但是耶穌不是空口講白話.
你們沒有試探.
求上帝拿走試探.
是耶穌親自經歷過試探.
而祂為我們.
也教我們向上帝祈禱.
最好不要見到試探.
你明不明白那個內容?.
就是說.
試探是很難搞的.
是不容易的.
很多人將耶穌受試探這段經文.
就會解釋試探.
我們有分試探.
有分試煉.
其實你看看聖經.
沒怎麼分的.
他偶爾就用試探.
偶爾就用試煉.

$^{121}$明明同一個原文都是這樣用的.
其實告訴我們.
就是這樣.
人生就會經歷一些.
你沒有想過會面對.
而是魔鬼撒旦會出現的時候.
不要以為你信了耶穌之後.
魔鬼撒旦就不見了.
我信了耶穌之後.
有光環.
我發光.
成功了.
魔鬼見到我都退開.
避開.
有人講過.
越屬靈的人.
才會受魔鬼撒旦攻擊越多.
如果連魔鬼撒旦都不攻擊你.
很簡單.
你都屬於他的.
他不會攻擊你.
自己人.
聽得明白嗎.
很重要.
所以當你面對攻擊.
有試探來到.
不要以為有沒有搞錯.
上帝不祝福我的.
為什麼要給我試探.
而是你會發現.
這是上帝讓我們學習經歷的.
特別是馬太福音.
今天講的馬太福音第四章.
你如果回去看馬太福音.
從第一章開始.
即是那些很長的家譜之外.
然後就在講聖靈.
向瑪利亞顯現.
說耶穌要出世.
向約瑟顯現.

$^{161}$然後聖靈就帶耶穌去洗禮.
洗禮天就打開.
聖靈就降臨.
不斷是聖靈帶著耶穌.
由他未出世到他出世.
又帶著約瑟,瑪利亞逃難到埃及.
又回到拿撒勒.
不斷講聖靈和上帝使者帶靈.
而今天第四章這段經文.
一開始講得很清楚.
帶領耶穌去曠野受試探的是聖靈.
是聖靈帶耶穌去曠野.
是聖靈帶耶穌去到那個被魔鬼試探的地方.
所以上帝從來沒有叫我們.
告訴我們一輩子都不會見到試探.
當然我們祈禱求上帝幫我們.
希望我們不要見到試探.
後面還有一句.
是拯救我們脫離的惡者.
給力量我們贏他.
雖然那個期盼是很難.
對好事的期盼是要的.
有人病了.
走來跟你說你幫我祈禱吧.
我病了.
你不要說你病了可能病不好.
可能會死的.
不用祈禱.
我們信仰不是這樣的.
你都要為他祈禱.
雖然你知道真的為他祈禱.
上帝有恩典就醫他.
上帝要接他回去就是上帝的心意.
你將主權交給上帝.
但你都會將你期盼的去禱告.
所以不要今天聽完這邊就說.
牧師叫我回去就不用祈禱拿走試探.
甚至祈禱上帝給我多些試探.
沒理由的.
人當然是期望有一個所謂叫做.

$^{201}$順利一點好一點.
我們新年都祝大家出入平安.
身體健康.
事事順利.
都是期盼.
雖然我們都知道期盼不一定成就.
不過如果上帝有恩典成就.
就感恩.
所以我們將我們期盼表達.
是好事.
上帝我真的想生活好一點.
上帝我真的想有多些能力去服侍人.
上帝我真的想多些恩賜.
以至教會可以服侍你.
上帝我真的想全時間服侍你.
好想去讀神學.
你有這些祈禱你就祈禱.
上帝會為你開路.
這是很重要的.
這段經文裡面開始講魔鬼撒旦去試探耶穌.
大家都很熟的.
第一件你要記住一件事.
魔鬼撒旦不是一開始就叫耶穌.
你否定神.
你否定上帝.
沒有上帝的.
你不要跟.
魔鬼撒旦是用什麼呢.
是用一種.
就是要你懷疑一下上帝.
在一些狀況裡面叫你懷疑.
一開始頭兩個試探裡面都很清晰的.
就是講.
如果你是上帝的兒子.
就這樣這樣這樣.
還要是當時.
魔鬼出現的時候.
是當耶穌餓的時候.
耶穌餓.
餓很正常.

$^{241}$你肚子餓嗎.
你現在應該還沒有肚子餓.
過了半小時.
我再講長一點.
你就會肚子餓了.
你就去馬爾木斯廣場.
肚子餓了.
人總有餓的時候.
一定是.
這個餓當然是身體的餓.
但是有時候這個餓可能就是那種.
就是我們的渴求.
我們有一些安全感.
對生活.
對金錢.
對成功.
對認同.
那些渴想渴求.
很想被滿足.
很正常的.
不要以為覺得被滿足很奇怪.
人家想被滿足.
人家讚你.
當然開心.
多謝你.
讚我.
但是這種渴求.
是有時候會出現一些的縫隙.
因為當你渴求的時候.
自然就會有一個需要.
就是肚子餓.
很基本的需要.
在這裡耶穌就是肚子餓.
肚子餓吃東西很簡單.
最後最後先劇透.
這段經文最後.
天使就馬上來接受耶穌.
馬上有東西吃就放心.
不會餓到.
不要餓到耶穌.

$^{281}$雖然當時有禁食的操練.
但是這段經文開始的時候.
魔鬼就用了這個縫隙.
就是當人有一個渴望.
有一個需要.
有一個飢渴的狀況.
就會出現.
而這個出現一開始的時候.
是質疑.
是懷疑.
是教你懷疑上帝.
如果有人走出來叫你.
你不要信上帝.
上帝是假的.
你會說不是吧.
我是回教會的.
我信耶穌的.
我在良法堂的.
但是一開始的時候.
魔鬼叫你不是說你不要信.
而是魔鬼要你懷疑.
你如果是.
應該是這樣的.
如果上帝這麼愛你.
為什麼你的生活這麼苦.
如果上帝真的掌權.
為什麼香港會變成這樣.
如果上帝真的存在.
上帝真的這樣祝福你.
為什麼你的生活會這樣.
這個是魔鬼第一步.
也是很重要的一步.
對耶穌試探的內容.
如果你是上帝的兒子.
我們當然知道耶穌是上帝的兒子.
沒有人質疑.
並且天都打開.
天父就在上面叫下來.
謝是我的愛子.
這是不是在洗禮的時候很清晰.

$^{321}$但是魔鬼第一個就是挑戰耶穌這個身份.
我們作為基督徒.
我們這個身份.
我們這個屬於上帝的愛的兒女.
在上帝的家裡面追隨基督的人.
這個身份.
魔鬼是容易質疑這個身份的.
魔鬼是常常用這個身份去質疑.
以致我們開始慢慢懷疑上帝.
魔鬼不是第一句叫你不要信他.
你慢慢懷疑.
如果你真的上帝祝福你.
為什麼每個星期翻崇拜.
好像都沒什麼改變.
你還要是星期天沒有上班.
你賺少了錢.
你的生活變差了.
上帝真的祝福你嗎.
開始慢慢懷疑懷疑.
你就轉身離開了上帝.
這個要留意.
第二就是我們經常覺得.
我們熟悉聖經.
多背幾個金句就可以了.
讀聖經又很重要.
又不要懷疑牧師他們不讀聖經.
沒理由.
不過這段經文裡面有一個很清楚看見.
當耶穌講聖經.
人們再不是單靠食物.
跟著魔鬼來第二個試探的時候.
就講聖經.
還要引用詩篇九十一篇.
就是說你是上帝的兒子.
你跳下去吧.
經常這樣講的.
他要使者去托著你.
那段是聖經.
那段講的是.
上帝會命令使者用手托著.

$^{361}$免得你的腳碰在石頭上.
這段是詩篇九十一篇.
十一到十二節.
你可以回去看一下.
魔鬼都引用聖經去試探耶穌.
有些留心.
有人跟你講聖經.
不一定是對的.
有人跟你講聖經.
他可以是魔鬼.
不要以為講聖經的人一定是屬於上帝.
魔鬼都講聖經.
當然他歪曲這些道理.
當然他做很多歪理.
用這些方法.
但是很多人打了信仰的旗號.
結果是一些很不合理.
不合真理的表達.
你事奉很盡心.
可以了.
你回教會做好一點.
其他事情不用管.
生活不用管.
最重要是成功.
教會最重要是興旺成功.
所以無所謂.
用什麼辦法都可以.
現在外面的騙徒都是這樣騙人去做壞事.
我們用不同的方法.
總之人來到教會就行了.
我真的見到教會這樣的描述.
以致用一些我用這個詞叫做欺騙的方法.
引人回來教會.
我覺得這樣都可以嗎.
是的.
有時候人會這樣.
當我們不清楚明白魔鬼殺人的手段的時候.
我們成為了那個欺騙者.
試探者的時候.
有時候真的會拿著聖經.

$^{401}$都會令你離開上帝.
那怎麼辦呢.
當然聖經你相熟的話.
當然要看多些聖經.
但是同樣地不要再斷章取義.
不要再只是背金句.
背金句很重要.
不要只是背金句.
背了都要看前文後理.
知道要怎麼用.
這是很重要的.
每一年我都會.
很多人經常問我.
牧師你選一些金句.
好像新年一樣.
給我一些字.
給我一些祝福.
經常叫我給他一些金句.
我一給就一整段.
這麼大一段.
是這樣的.
沒理由只是看一節.
有時候我們容易斷章取義.
很危險的.
魔鬼都會用聖經去試探耶穌.
記住.
記住這件事.
這三個耶穌基督面對的試探.
是一場我們叫做生死戰.
因為是很實在的.
第一個描述是石頭變食物.
這個描述.
我們當然沒有石頭變食物的能力.
你有沒有.
沒有.
我們都要買東西吃.
撿石頭變食物就好了.
石頭變食物是什麼描述.
這是一個生存的焦慮.
去面對生命主權的描述.

$^{441}$希望你聽清楚.
是一個生存的焦慮.
面對生命主權的一個試探.
我們在香港生活的基本層面很高.
我們經常跟人討論.
我現在在中文大學.
其實都比較近深圳.
好多同事偶爾就回到羅湖吃飯.
甚至住在那邊.
他經常說在香港生活太難了.
香港謀生很艱難.
生存本身就是一種試探.
你怎麼生存.
在香港你不謀生哪有得吃.
大家都這麼辛苦.
挨餐餅才剛剛夠吃.
不過撒旦在這裡引誘我們.
將我們這個經濟的安全感.
擺在屬靈的原則之上.
耶穌有這個能力.
當然可以石頭變麵包.
其實有多難.
我們就不行.
對耶穌來說.
石頭變麵包有多難.
後期當他跟門徒一起的時候.
五餅如魚餵飽幾千人幾萬人.
所以對耶穌來說.
石頭變食物從來不是一個難處.
我經常開玩笑說.
耶穌用不著石頭變麵包.
不要肚餓了.
那就不肚餓了.
耶穌來的.
一句話我不肚餓了.
我一輩子都不用吃了.
當然他沒有這樣做.
他跟我們一起經歷人的需要.
但是這個描述.
魔鬼跟耶穌說.

$^{481}$你將石頭變麵包吧.
其實是挑戰我們.
是的你要謀生.
你面對現實.
你不要以為你是上帝的兒子.
所以說你若是上帝的兒子.
就好像懷疑你有沒有能力.
物理能力將石頭變麵包.
你是上帝的兒子.
你就有能力將石頭變麵包.
但是另一件事就是.
你是上帝的兒子.
其實你都被這些食物困住.
你真的是上帝的兒子.
你不用吃都可以.
你為什麼要吃東西.
你明白嗎.
很多時候我們都會這樣被試探.
基督徒都要謀生的.
基督徒不用上班嗎.
基督徒不用發工資嗎.
不會因為你是基督徒.
你的工資少一點.
你都過得活的.
不會的.
就算你是基督教的機構.
你都要發工資.
教會都要發工資.
都要生活的.
但是這裡不是說我們不生活.
我很喜歡這句金句.
我小時候就說金句.
我最喜歡的金句其實是這一句金句.
人活著不是單靠食物.
我就經常說下一句.
不是單靠就是都要靠.
所以我們都要靠食物.
我都喜歡吃東西.
大家都知道.
但是這裡不是單靠食物.

$^{521}$是很好的描述.
你不要以為靠食物就夠生活.
應該這麼說.
你不是只養著一條命就夠.
你要有價值地生活.
聖經裡面說耶穌是讓人得生命.
並且得的是豐盛的.
祂不是只想你多活幾多秒.
那口氣長多一點.
耶穌是想你生命豐盛豐富.
所以如果你生活只得食物是不夠的.
飽肚.
接著你就會餓.
餓之後你就飽肚.
如果你是這樣生活的話.
你和一般的動物.
其他的植物生物都沒有分別的.
你只是滿足到身體.
生活,生存的需要.
但是當這個生存變成焦慮.
變成你覺得我怎麼生存.
我怎麼活得好一點.
結果就是我要謀生.
結果你就離開了信仰的核心的話.
你就會出事的.
所以耶穌說不單靠食物.
是靠什麼呢.
是靠上帝口裡所說的一切說話.
你說上帝的說話能吃得下肚嗎.
上帝的說話能飽肚嗎.
是,人生都要繼續食物.
但是同樣地.
當上帝的食物能夠滋潤你的生命的時候.
你會不再一樣.
總之我鼓勵你.
在這個大齋期日子當中.
你真的每天都好好去靈修讀聖經.
我不知道你過去有什麼習慣.
如果你有好好靈修讀聖經.
如果你從來都沒有.

$^{561}$這段時間是很好的機會.
每天什麼都不要做.
先讀了聖經,祈禱.
就好像你肚餓,先吃了東西.
睡醒,第一時間做什麼.
空肚,先吃早餐.
先吃了早餐.
那個就是你的早餐.
如果每天你去到上帝面前.
你先飽食靈量的時候.
你發覺你的生命不一樣.
當你發覺原來不看聖經那天是會死的.
用這個字.
不看聖經是會死的.
我太太師母經常說.
不看聖經那天是會死的.
是真的.
你經歷過你就明白了.
所以我希望你真的好好地進入上帝的位置.
讓上帝的位置成為你生命的糧.
不是神學上說.
上帝是我們的糧.
是一個理念這麼簡單.
而是真的每天好好地去吃.
每天就好好地去看聖經.
祈禱和上帝親近.
上帝給的話語成為你每天的力量.
就好像吃飽飯.
很精力充沛地去做你的行動.
一模一樣.
所以這個是耶穌給我們的安慰.
也是提醒.
所以耶穌告訴我們.
食物不是不重要.
不過食物不能夠定義你的生命.
食物不能夠定義你的生命.
最重要是你找不找到你要追求的目標.
第二個試探是說到電上面要跳下去.
幾個課堂書.
二三個試探有不同的次序.

$^{601}$今天我們終於是馬來課堂這裡去描述.
第二個描述裡面說要在電上面跳下去.
其實是什麼呢.
耶穌被撒旦帶到電上面.
聖殿的頂.
就說你跳下去吧.
天使會保住你的.
這個試探很有趣.
第一這個試探是魔鬼叫耶穌做一些沒必要的事.
是不是.
肚餓吃東西都叫做必要.
走上天台跳下來也沒必要.
沒必要的.
沒人會突然間貪得意上天台跳下來.
沒必要的.
但是這個描述裡面是說一件事.
就是一種屬靈的試探.
一種在靈裡面的試探.
一種讓我們去看看.
究竟人生裡面我們追求的是宗教的表現.
或者表演還是我們真實的信教.
很多時候我們都是這樣的.
信仰變成了我們生命的一種投資.
一種交換.
上帝我這樣對你.
你應該這樣對我.
上帝我奉獻了.
你應該祝福我.
甚至起碼都祝福我.
戶口那裡不會少了錢.
股票不要碰著.
我是信徒.
上帝應該在我上面都發一下光.
發一下上帝的光.
人家見到我都威一點.
上帝令我有一個霉的.
明白我說什麼嗎.
很多時候我們都是這樣想的.
我們都覺得我們作為信徒.
追求信仰就應該有光環.

$^{641}$那種光環就是.
上帝給我很多的祝福.
以致人看到我很厲害.
做基督徒當然要做到很順利.
很威很厲害.
如果做基督徒做到很霉.
怎樣才算是耶穌.
那就見證不了上帝的榮美.
不行不行.
所以我一定要很光鮮.
很威很厲害.
這個就是屬靈裡面很大的陷阱.
也是魔鬼在試探裡面給我們很大的陷阱.
有時候我們都很想有屬靈裡面的宗教表演.
是不是.
真的有病.
我請牧師來祈禱.
將來就康復.
嘩真的厲害.
好像厲害一點.
我記得那時候在教會這麼多年.
梁伯棠也好.
在資源教會也好.
很多電影節目都說.
牧師來祈禱.
來探我.
甚至有時候覺得.
當年我還沒按立做牧師.
做宣教師的時候.
有些電影節目說.
教會不試探我的嗎.
我其實每個星期都探他.
不算的.
牧師來才算的.
比較厲害的一個祈禱.
真的跟著牧師去探他.
第二個禮拜就康復了.
我真的覺得.
是不是真的.
當然了.

$^{681}$上帝透過不同的人.
有不同的恩賜去發揮.
當然我做牧師為電影節目祈禱.
當然有很多人康復了.
當然也有一些是上帝的心意.
接他回去.
沒有擔保的.
但是我們很容易在想.
上帝厲害.
我也要厲害一點.
你不要將邏輯混亂.
當然了.
我們作為上帝兒女.
當然不要失禮他.
如果我們做到那部分.
有些部分我們做到的.
敬虔的生活.
好好服侍人.
有愛心.
你的信仰和你的行為匹配.
這些我們做到的.
但是我成功了.
我賺大錢了.
我發達了.
我能夠人生攀著高峰.
有財費.
有地位.
那些是很多不同的際遇狀況.
包括你努不努力的.
不是因為你是基督徒.
就等於你買那隻股票一定會升的.
不是教會投資小組買那隻股票基金一定會升的.
現在我在區會是執委會的副主席.
是投資委員會的主席.
所以我很大壓力.
至少黃牧師也知道.
其實不關我的事.
很多顧問.
不過如果區會買了那隻.
跌了怎麼辦.

$^{721}$我們會這樣想.
當然.
人家是希望買股票會升.
不然為什麼買股票.
投資就是為了賺利息.
聖經也教我們.
有財幹也不要收賣.
要用的.
所以收利息也重要.
這些都是好的教導.
但是如果我們以為.
我這樣盡了力.
做了我所做的事.
你發覺結果不好.
那就是上帝蒙羞.
所以上帝一定要祝福我.
讓我做的事順順利利.
風風光光.
還有人覺得很厲害.
其實背後你是否為了自己想享受多一點.
背後你是否為了自己的成功感強一點.
背後你是否為了自己的名聲大一點.
還是你真的想榮譽上帝的名聲.
頂住他的名脈.
我們那時候就是這樣.
這個試探其實是很深層.
他在說一件事.
就是我們好像一些.
其實沒有必要.
耶穌沒有必要走到天頂跳下來.
沒有必要的.
當然有人就說.
邏輯上耶穌是否真的這樣跳下來.
天使是否真的這樣捉住他.
其實問題是那件事不應該發生.
《聖經》裡面詩篇說的是上帝的保護.
就算有一天你會跌倒.
放心.
上帝的保護就會用天使去圍著你.
抱著你.

$^{761}$這個是上帝的愛.
但是沒有人叫你濫用.
我經常都有一個例子.
就算有人買一張膠布給你.
這個膠布是很厲害的.
一貼就馬上好了.
什麼傷痕都可以戒掉.
一貼就馬上好了.
那你會怎樣.
你當然會戴著它.
有需要的時候用.
你會不會說先試一下.
戒兩刀貼一下.
可以的.
戒到全身都傷了.
然後不斷貼.
為了證明這個膠布行不行.
不是這樣的.
這麼好的東西.
就戴著有需要的時候用.
因為那個不是需要.
還要.
其實.
為什麼無端端要發生這樣的事.
因為魔鬼撒旦的試探.
就是想你表演一下.
你是耶穌.
你這麼厲害.
你真的相信你的兒子.
你在天頂跳下來.
很多人看著你.
你就很威風.
耶穌.
你相信你的兒子.
你應該很威風的.
捧耶穌.
捧到他很威風.
很有趣的.
你看看聖經裡面.
捧耶穌捧到最威風的.

$^{801}$是魔鬼撒旦.
那些門徒反而覺得.
耶穌怎麼搞錯.
你做這樣的事.
耶穌你做這樣的事.
你不是應該很威風的嗎.
反而那些門徒不是很明白.
這個很有趣的地方.
小心.
記住.
留意.
我們很多時候都是這樣.
耶穌回應是.
不要試探主你的上帝.
真正的信仰.
不是看上帝為我們做了多少.
很驚天地.
很偉大.
很厲害.
很威風的事情.
而是我們在任何時候.
堅定的依靠祂.
不去懷疑.
這才是重要去聖歌試探的核心.
第三個試探.
在《瑪利福音》裡面說.
魔鬼將萬國的榮華.
給耶穌看.
然後就說.
這裡是萬國的榮華.
你拜我.
我就將一切給你.
你說.
魔鬼這個試探是不是真的很厲害呢.
萬國的榮華.
耶穌是上帝.
又是上帝的兒子.
祂什麼都有.
萬國榮華明明都屬於上帝.
不是屬於魔鬼的.

$^{841}$但實際上.
就真是魔鬼將萬國榮華.
展現在耶穌面前.
和頭兩個試探不同.
頭兩個試探是.
如果你是上帝的兒子.
石頭可以練練包.
他還是看著石頭的.
如果你在電裡面.
頂上面跳下來.
天使要保住你.
那些天使還沒出現.
當然沒有跳下來.
沒有發生那些事情.
但是第三個試探.
那些榮華.
萬國的榮華.
是立即展現在面前.
頂著面我們可能覺得.
是啊.
上帝會幫我.
我有能力.
我用心.
我一定可以勝過那些試探.
就是因為.
很多時候其實我們都沒有.
沒有真的面對那些事情.
很多人說.
如果我有什麼什麼.
我就會怎麼做.
簡單舉個例子.
有人跟你說.
你不要回教會.
我請你喝茶.
可能你會覺得.
喝茶而已.
不喝就罷了.
新年都喝那麼多了.
都生氣了.
不要喝了.

$^{881}$但是告訴你.
你如果今天決定不信耶穌.
我就給你十萬元.
你說.
扯的背後就是.
都看不到那十萬元.
說說而已.
如果人家真的就站在教會門口.
拿著一封利是.
你就有十萬元.
你拿了這封利是.
以後不踏進良法堂.
我就給你一封利是.
你說.
那個人傻的.
報警抓他.
拿了這封利是才抓他.
我們看到那些事情.
人生很容易去想.
特別在我們香港這個城市.
除了生活得很艱難.
但是我們很講效率.
很講速度.
很講回報.
如果魔鬼撒旦給你一條快速道路.
Fast Pass.
好像你去公園不用排隊.
你很快的.
很快的成功.
很快的能夠得到回報.
很快的能夠拿到你要拿的東西.
走捷徑的成功.
你願不願意去接受呢.
你說.
走一下捷徑.
我厲害而已.
但是如果走捷徑的成功.
他對著的.
告訴我們原來耶穌.
是要走上十字架的路.

$^{921}$大齋奇告訴我們.
我們開始去進入.
最後耶穌要走上十字架路的.
一個預備的日子.
耶穌沒有快速地成功.
耶穌沒有走出來.
做很多很偉大的事.
然後很快成功.
我經常都說.
如果我寫這個劇本.
耶穌這麼厲害出來.
醫病趕鬼.
有五千人做很多厲害的事.
人們都跟著他.
教會就成功了.
然後就改變這個世界了.
但是耶穌最後選擇的是.
慢慢逐步走上十字架.
是一個願意為人犧牲的路.
他沒有走捷徑.
弟兄姊妹.
在我們生命當中.
有時魔鬼就會用捷徑來吸引我們.
這個捷徑包括生命裡面的成功.
這個捷徑甚至包括信仰裡面.
我們以為的成功.
很快的.
在教會裡面走.
信仰我很虔誠的.
三個月就馬上完成某些課程.
然後就馬上出來做傳道.
馬上出來牧羊教會.
或者很快很快就覺得自己很成功.
很短時間裡面.
信仰呢.
生命是要慢慢浸淫的.
更重要的是.
原來這些捷徑呢.
是讓我們以為成功.
追求成功是我們生命的目標.

$^{961}$真的有的.
有朋友問我.
我想做牧師.
我說好啊.
三年夠不夠.
我說三年比較難.
那五年吧.
那我說五年.
讀完神學之後就五年.
讀神學幾年.
又讀多三五年.
那就十年吧.
我說都最少吧.
然後他就真的跟我說.
有沒有快一點的.
速成的.
我以前以為.
我以前教樂器.
大家都知道我彈結他.
我試過有電影節目跟我說.
我想學結他.
可不可以速成.
三堂學一首歌.
兩堂兩堂這樣.
當然我就不會教了.
兩三堂學到些什麼呢.
現在信仰居然都有速成.
居然有剛才說的問題.
是一個大學的同學問我.
我想做牧師.
我大發熱心.
那我三年夠不夠.
我想年輕一點服侍.
當然我感恩他想年輕服侍.
但他覺得不要緊.
只是學位而已.
甚至問我.
是不是有些神學院.
或者外面那些.
可以銜售.

$^{1001}$半年九個月就讀完了.
好像那些人想買個學位進修.
神學都要買個學位.
你真的做牧師.
真的做傳道宣教師.
你怎樣去牧羊鼎之門.
到時才學吧.
很想很快地成功.
我用這個字.
當然我不懷疑他服侍的心.
但為什麼他會這樣想.
當然可能覺得這個世界邪惡.
時日無多盡快.
對的 你就用不同的服侍方法.
你可以繼續大發熱心.
不一定要讀神學.
拿學位做牧者.
這是上帝呼召的一部分.
但如果上帝呼召你.
這個部分就要這樣生活成長.
沒理由期望要走捷徑.
更遑論我們每日的生活.
更遑論我們每日的賺錢.
上班 做生意.
我們很多時候都很想有捷徑.
最好想你告訴我.
這條路.
給我一個號碼那些.
以前很流行.
可能大家不知道我在說什麼.
怎樣可以很快地成功.
我們可能很想.
但其實實際上.
聖經裡面.
在這段經文裡面.
這個試探裡面告訴我們.
耶穌沒有選擇很快成功.
沒有選擇.
甚至沒有和魔鬼爭辯.
那些萬國明懷是我的.

$^{1041}$還給我.
是不是應該這樣爭辯.
耶穌都沒有.
耶穌只是很簡單 很直接.
和魔鬼撒旦說了一句.
那句話是什麼呢.
撒旦 退去.
因為經常記著.
要拜主你的上帝.
唯獨事奉他.
這一個就是.
耶穌最後最後一句說話.
在馬化音裡面.
讓魔鬼撒旦離開的經文.
魔鬼撒旦被耶穌趕走.
是因為耶穌說.
我的目標.
我的人生只有一個.
唯獨是對著我們的上帝.
對抗魔鬼撒旦.
有什麼方法呢.
怎樣可以在試探裡面.
去得勝呢.
一個很重要很重要的焦點.
就是你的人生是否正在對準上帝.
你眼裡面對準上帝.
身邊的人就不是.
身邊的事物也不能夠再一次.
去吸引你的目光.
不會再次偷看你.
不會讓你失去焦點.
很多時候人生失去焦點.
因為你不清晰.
成長人生的路都是這樣的.
信仰更加.
當你找到耶穌.
你很單單.
唯獨去事奉他.
這句話常常被我提醒.
特別是過去這幾年.

$^{1081}$我在大學裡面.
很多不同的事務.
在區會.
在大學.
很多會議.
很多工作.
很多時候我常常問自己.
究竟我是否很專心.
唯獨事奉上帝.
還是其實做很多其他的事.
不務正業.
我常常說笑.
突然問我.
牧師你又喜歡這件事.
你又會那件事.
為什麼.
我說我不務正業.
這是說笑.
正業是什麼呢.
基督徒的正業是什麼呢.
單單的.
唯獨的事奉上帝.
無論你是什麼角色位置.
無論你是牧師宣教師.
無論你是執事堂位.
無論你是弟兄姊妹.
你說我是平信徒.
你的焦點都是單單的去事奉上帝.
這個焦點很重要.
無論你搬椅子給隔壁的弟兄姊妹.
無論你讓座給那位有需要的朋友.
無論你回教會派單張.
無論你每一次奉獻.
把錢放進一個奉獻袋.
你都要告訴自己.
我是唯一單單服侍的對象.
第一個.
就是我們主上帝.
第一個就沒有其他的.
我不需要被人討好.

$^{1121}$我奉獻這麼少錢.
旁邊那個看到就很羞恥.
這麼多人看著我.
我拿多點錢出來奉獻吧.
教會登出來的周刊有名字.
我的名字這麼少錢.
登多點進去吧.
對不起.
這樣你奉獻多少錢都沒有用.
這樣都不是真心的奉獻.
希望你明白我說什麼.
我們要找到我們事奉的焦點.
當你看著上帝的時候.
你眼尾都找不到魔鬼撒旦的話.
你放心.
當然他走了之後會怎樣?.
會回來的.
聖經裡面其他福音書說得很清楚.
耶穌趕了魔鬼撒旦走了之後.
這裡說天使馬上事奉他.
接著魔鬼再等候機會再一次出現.
所以後面就說.
魔鬼就出現進入了猶大的心.
出賣耶穌.
魔鬼撒旦隨時繼續在這裡.
因為這個世界仍然是邪惡的世界.
當然耶穌還沒回來.
我們仍然生活在這個.
有魔鬼.
有撒旦.
有欺騙者.
有這位說方者在我們身邊.
你怎麼可以避開他呢?.
你怎麼可以看不到他呢?.
你怎麼可以.
他在你身邊怎樣去搖動你.
你都不被搖動呢?.
單單看準你的目標.
看遠一點.
看清晰你的目標.

$^{1161}$就好像一般生活.
走路.
開車.
看清楚你的目標.
現在我們很多時候低著頭.
拿著手機.
低著頭.
就算你找路.
按個導航.
接著你就看著那一步.
怎樣導航就帶你衝下山.
你就照樣走.
因為你都沒理會目標在哪裡.
目標在清晰是很重要的.
路有時彎彎曲曲.
有時跌跌碰碰.
這個是窄路.
有時是爛的.
不過放心.
身邊有弟兄姊妹一起.
更重要看準你的目標.
在這個大測期當中.
鼓勵弟兄姊妹.
每天好好地明修讀經.
好好地找回你生命的焦點.
看準.
看清楚.
唯獨去事奉我們的主上帝.
放心.
就算魔鬼撒旦在我們身邊.
很快就會離開.
很快就會避開你的敵黨.
而我們可以一起去勝過他.
請同心祈禱.
就是我們生命裡面.
有時真是面對難處.
在漫長的生命當中.
我們真是要和魔鬼打一場漫長的戰爭.
但我們不懼怕.
因為耶穌你成為我們最好的典範.

$^{1201}$你勝過魔鬼撒旦.
你將他趕走.
你讓他知道你是單單服侍上帝.
幫助我們.
有同樣的信心.
同樣的堅持.
看清楚我們生命的焦點.
更讓我們不被那些.
可能是人看來快捷成功的道理.
或者人看來一些宗教上的精彩的外衣.
甚至可能是一些生活裡面基本的需要.
而障礙了我們.
而我們真是能夠切切實實的.
好好地在你面前.
領受你的恩典.
每日就向著你去邁進.
儘管路有時跌跌碰碰.
有時高高低低.
甚至有時好似有些難過的位.
但我們不怕.
因為知道我們有一日就會走到這個目標.
求你幫助我們.
讓我們和眾弟兄姊妹一起去並肩去成長.
在你的恩典當中.
讓梁家堂眾弟兄姊妹.
能夠去體會生命的精彩和豐富.
以致我們走到最後那一日.
走到面見你那一日.
走到耶穌叫你回來那一日.
我們仍然不偏離.
仍然望著你.
以致被你讚賞我們成為良善忠心的僕人.
求你恩代.
求你賜福.
讓我們真是學懂.
去戰勝魔鬼撒旦的道理.
讓我們真是認清.
我們可以怎樣在你面前.
被你去提醒和鼓勵.
多謝主.

$^{1241}$聽我們禱告.
用耶穌基督的名堂.
阿們.
\newpage



\section{創世記 17:16}
\label{sec:031684}
\textbf{聖經人聖經事:從她得一個兒子}~~~~黃克勤牧師博士
\newline
\newline
連結: \href{http://www.church.com.hk/acms/content.asp?site=cdc\&op=show\&type=product\&code=031684\&layout=sermon}{\texttt{031684}} ~~~~ 語音日期: 2026-02-11
\newline
\newline
\hyperref[sec:031683]{\small{< < < PREV SERMON < < <}}
~
\hyperref[sec:index_chronic]{\small{[返順時目]}}
~
\hyperref[sec:031685]{\small{> > > NEXT SERMON > > >}}
\newline
\newline
創世記 17:16
\newline
\begin{longtable}{cl}
\hline
\hline
章節 & 經文 (和合本修訂版)\\
\hline
17:16 & \begin{tabularx}{0.7\textwidth}{X} 我必賜福給她，也要從她賜一個兒子給你。我必賜福給撒拉，她要興起多國；必有百姓的君王從她而出。」 \end{tabularx} \\ \\
[1ex]
\hline
\hline
\end{longtable}
$^{1}$各位姐妹,很快又一日了.
我們今天看的是.
「從他得一個兒子」.
「他」是指亞伯拉罕的太太.
舊名叫做薩萊.
新名叫做薩拉.
從他得一個兒子.
聖經人聖經事來自生命說書的網絡.
希望我們多讀聖經認識上帝.
看看他在人與事上怎樣管理.
以及背後我們的上帝是怎樣的.
從他得了一個兒子.
十七章十五十六節開始.
「神又對亞伯拉罕說」.
改了名了.
以前叫什麼?.
亞伯蘭是高舉的父.
即是好的父親,偉大的父親.
亞伯拉罕多國之父.
多人的父.
「你的妻子薩萊」.
薩萊.
薩萊不要再叫薩萊了.
薩萊的意思.
以前說過了.
是公主.
公主呀,很漂亮呀.
薩萊應該是很漂亮的.
公主呀.
我要叫她的新名叫薩拉.
薩拉.
薩拉就是多國之母.
多人之母.
配上亞伯拉罕叫多國之父.
多人之父.
多人的母.
我會賜福給她.
也使你從她得一個兒子.
我要賜福給她.
她要作多國之母.

$^{41}$就是薩拉的意思.
必有百姓的君王從她而出.
在她後面會出一個君王.
從她得一個兒子.
這是應許.
亞伯拉罕和薩拉會有一個兒子.
從她得一個兒子.
我是全能的神.
跟她說.
一直說到現在.
誰知她們不太相信.
亞伯拉罕就虎伏在地.
嬉笑說.
一百歲的人還能得孩子嗎.
我太太也老了.
薩拉已經九十歲還能生養嗎.
如果你看其他經文.
有三個人來到亞伯拉罕.
請坐先吃點東西.
今年你太太會生一個兒子.
明年就會生.
薩拉在帳篷後面笑.
其實她只是暗笑.
但其中三個人說.
你太太為何要笑.
薩拉說.
我沒有笑.
你是有笑的.
薩拉在帳篷後面笑.
九十歲還能生孩子嗎.
亞伯拉罕虎伏在地.
嬉笑說.
第一留意這個字.
亞伯拉罕是否不信神.
笑說.
一百歲還能生孩子嗎.
亞伯拉罕虎伏在地.
何謂虎伏.
我是全靈的神.
我要祝福你.

$^{81}$多謝.
虎伏.
你要生一個兒子.
不是呀.
多謝.
虎伏.
而心裡是喜笑.
不是嬉笑.
喜笑.
開心到笑出來.
開心到發夢.
喜笑.
虎伏在地.
開心到不得了.
很開心.
心裡說.
我都一百歲.
還能生孩子嗎.
這兩行不能認為是亞伯拉罕.
真的那麼不信.
薩拉在帳篷後面是不信的.
因為三位其中一位指責他.
你不信.
你偷笑.
即是不信我們說的話.
即是說我們不行.
薩拉說.
我沒有笑.
你真的有笑.
被上帝知道了.
上帝也沒有什麼不知道.
這裡神沒有說他.
亞伯拉罕虎伏在地.
這樣也行.
我一百歲.
我太太九十歲.
怎麼行.
但其實沒有.
為何有這句.
亞伯拉罕對神說.

$^{121}$但願以實瑪利會在你面前.
以實瑪利是誰呢.
就是神以前對亞伯拉罕說.
你會生一個兒子.
但薩拉說.
我九十歲.
哪裡能生.
如果以當時的慣例.
可以做的.
就是叫自己的婢女做妾侍.
做丈夫的妾侍.
於是為她生兒子.
由她而出.
所以就是她的兒子.
於是薩拉就用人的方法.
而當時有一個不夠好的地方.
亞伯拉罕也覺得.
我們也是這樣.
他接受了.
於是生了兒子.
跟婢女下嫁.
生了兒子.
叫以實瑪利.
屬於他的.
不過皮膚黑一點.
是埃及人.
下嫁是埃及人.
以實瑪利媽媽是埃及人.
是一個婢女.
但願以實瑪利會在你面前.
他也很認真的.
上帝說你會生兒子.
多謝多謝.
一百歲哪裡能生.
不過我已經有一個兒子.
生了.
多謝你已經給我一個兒子.
願以實瑪利希望他會在你面前.
你看看他.
有了有以實瑪利了.

$^{161}$好了.
大家要留意的就是.
我的字.
從他得一個兒子.
第一點想說的.
上帝有他的能力和吩咐.
得一個兒子我給你一個兒子.
誰知.
亞伯拉罕誤會了.
多謝主我已經有了.
就是以實瑪利.
你也祝福他.
願他會在你面前.
雖然皮膚黑了一點.
不過也很英俊.
少少亂七八糟也說不定.
不過你祝福他吧.
以實瑪利.
上帝說不是.
是你的妻子撒拉要給你.
生一個兒子.
是她自己生.
而且我要改名叫.
以撒.
我會和以撒堅定.
和你所納的約.
不是和以實瑪利.
不過第三點.
說的是螢光的.
至於以實瑪利.
神真厲害.
我真的很欣賞.
以實瑪利我不算.
很快就收回他的生命.
不是.
以實瑪利既然從你而出.
我也有口齒.
從你而出我也祝福.
我也英文你會賜福給他.
使他昌盛極其繁多.

$^{201}$他有十二個族長.
都成為大國.
他會好像.
我們現在看.
亞伯拉罕當時不知道.
因為他不知道以撒雅國.
雅國有十二個兒子.
他不知道十二支派.
上帝在這裡說.
以實瑪利雖然不是我說的那個.
是你用自己人手方法.
做那件事的那個.
不過既然從你.
亞伯拉罕而出.
我也認賬.
所以我會祝福.
以實瑪利.
我也會賜福給他.
使他昌盛.
很多人十二個族長.
我也要使他成為大國.
也真是大國.
因為今天你看看中東.
以實瑪利的後裔是什麼呢.
就是阿拉伯人.
以實瑪利的後裔.
阿拉伯人多不多?.
大不大個?.
當然大了.
沙地阿拉伯已經夠大了.
阿聯猶大不大?.
大到極了.
約旦呢?.
多到不得了.
敘利亞呢?.
敘利亞不是.
其他很多國家.
全部都是.
以實瑪利的後裔.
今天和以色列一起.

$^{241}$是大國.
巴勒斯坦人.
大部分是.
巴勒斯坦自治政府.
屬下的人民.
不是猶太人.
在中東.
以色列地方.
大部分都是.
以實瑪利的後裔.
阿拉伯人.
黑黑的.
其實是兩兄弟.
可以這樣說.
不過.
唉.
可惜呢.
不停相爭.
直到今天.
其實這個呢.
就是撒拉.
做了一件.
人手做的事.
我可以.
未生到.
不過我可以.
幫到上帝的手.
用這個辦法.
這樣.
其實呢.
就不好的.
好了.
今天的金句.
我必似呼吸他.
也要使你.
從他得而至.
這個問題呢.
就是有時候.
上帝給了一個.
應許我們.

$^{281}$不過呢.
我們用神學的說法.
就是叫.
Already but not yet.
神已經答應了給.
不過還沒到.
我們這些人.
好像小朋友一樣.
要一個玩具.
一粒糖.
要一杯果汁.
差不多要.
馬上要.
如果沒有.
就自己想辦法了.
這裡就是.
神說.
我會給你的.
你放心.
我會給你的.
哎呀.
我等不及了.
於是.
撒拉就說.
唉.
哪裡可以.
九十了.
怎麼生你呢.
我九十.
怎麼生.
玩吧.
不行了.
於是.
行.
我有辦法.
給我一個.
下甲皮類給你.
他這些.
做慣事的人.
很好生養.

$^{321}$應該.
一碰就生了.
於是呢.
生了一個.
二神瑪利.
其實都是福.
但都是禍.
不過神是.
按著他的.
應許.
說了就真的.
從你而出.
我也祝福.
結果呢.
後來.
沒有處理得好.
撒拉做了.
第二件.
不好的事.
是什麼呢.
將.
下甲.
和他的兒子.
二神瑪利.
趕走.
任由他在外面.
曠野.
少少水.
可能餓死.
慌死.
任由他了.
不理.
亞伯拉罕.
不知是否想發財.
怕老婆.
他太太說.
這樣就好了.
他在你手中.
你說了算.
結果被趕出.

$^{361}$不過神.
憐憫.
救了下甲.
和二神瑪利.
終於.
剪成兩個.
猶太人.
以色列人.
和二神瑪利.
其實呢.
就是對神的說話.
不夠信任.
而.
縱效禍根.
請你記得.
神的應許是確實的.
我必賜福給他.
他也要使你.
從他底的兒子.
就是以撒.
所以想起以撒.
就是上帝應許是確實的.
實現的.
信得過的.
一齊禱告.
多謝主.
因為你是一位.
信實可靠的主.
人忠言失信.
但你仍是可信.
因為你不能背負自己.
你說過的話.
你就會認.
我們都是倚靠你的應許.
而過每日的生活.
幫助我對你更有信心.
我們並不會去指責.
撒拉亞伯拉罕.
而是他們當時.
有他們決定的原因.

$^{401}$不過今天回看.
是他們當時的信心未夠.
而釀成今天後裔.
很多的災禍.
求你憐憫我們.
讓我們信靠你.
緊緊的跟隨你.
抓住你的應許.
而生活.
順天祈禱.
奉耶穌基督的名.
阿們.
\newpage



\section{歌羅西書 1:15-20}
\label{sec:031685}
\textbf{耶穌是永恆超級屬靈接合劑}~~~~黃紹權牧師
\newline
\newline
連結: \href{http://www.church.com.hk/acms/content.asp?site=cdc\&op=show\&type=product\&code=031685\&layout=sermon}{\texttt{031685}} ~~~~ 語音日期: 2026-02-22
\newline
\newline
\hyperref[sec:031684]{\small{< < < PREV SERMON < < <}}
~
\hyperref[sec:index_chronic]{\small{[返順時目]}}
~
\hyperref[sec:031686]{\small{> > > NEXT SERMON > > >}}
\newline
\newline
歌羅西書 1:15-20
\newline
\begin{longtable}{cl}
\hline
\hline
章節 & 經文 (和合本修訂版)\\
\hline
1:15 & \begin{tabularx}{0.7\textwidth}{X} 愛子是那看不見的神之像，是首生的，在一切被造的以先。 \end{tabularx} \\ \\ \relax
1:16 & \begin{tabularx}{0.7\textwidth}{X} 因為萬有都是在他裡面造的，無論是天上的、地上的，能看見的、不能看見的，或是有權位的、統治的，或是執政的、掌權的，一概都是藉著他為著他造的。 \end{tabularx} \\ \\ \relax
1:17 & \begin{tabularx}{0.7\textwidth}{X} 他在萬有之先；萬有也靠他而存在。 \end{tabularx} \\ \\ \relax
1:18 & \begin{tabularx}{0.7\textwidth}{X} 他是身體（教會）的頭；他是元始，是從死人中復活的首生者，好讓他在萬有中居首位。 \end{tabularx} \\ \\ \relax
1:19 & \begin{tabularx}{0.7\textwidth}{X} 因為神喜歡使一切的豐盛在他裡面居住， \end{tabularx} \\ \\ \relax
1:20 & \begin{tabularx}{0.7\textwidth}{X} 藉著他，神使萬有與自己和好，無論是地上的、天上的，都藉著他在十字架上所流的血促成了和平。 \end{tabularx} \\ \\
[1ex]
\hline
\hline
\end{longtable}
$^{1}$我一齊跟心頭說.
莫忘地從上帝你這條管子.
直接帶到我們當中.
讓我們能夠亦都彼此在肢體之間的建造.
好使一同在上主理的裡面仰望.
能夠合神心意地去活在每一天的生命生活當中.
再次獻上感恩.
無論先上先下的每一個追求主道的人.
我亦懇求上主理帶領我們.
禱告奉靠主耶穌基督得勝的明智求.
阿們.
在這個世界裡面凡事都講求關係.
但是關係這東西.
卻正正是今天這個世界當中.
很多人最弱的一環.
而在這個世界裡面很多時候.
被我們看見的現象就是.
無論是人與人關係甚至男女之間的關係.
以至於我們和這個世界的關係.
我們都越來越疏離甚至是破裂.
所以今天我們在這個世界當中我們發現到.
很多人就將關係這東西.
就將它開始貶抑.
在中文的言詞上面.
關係這個字眼很多時候都只不過是用.
講到兩者之間的連結.
但在英文的上面有一個很清楚的界定.
關係不單只是一個relationship.
而且關係也都有一種層面是比較低一點的.
叫做networking.
兩者的分別在哪裡呢?.
前者是一個真正人與人有一個緊密的連結.
他們是彼此互為影響.
並且是帶著一份彼此關懷去建造的.
的目的而去做的.
但networking就很不同了.
networking就只不過是一個聯繫.
在電腦世界的裡面我們很喜歡用到就叫做node.
人只不過是一個點.
每一個點會不會為我帶來一些好處.

$^{41}$或者帶來一些資訊.
使到我能夠受惠.
所以你們大家都有上網的.
有用手提電話的.
上網其實只不過是一個聯繫.
手機只不過是一個點.
我要看打球我要看奧運.
我要看某一些其他等等的一些活動.
這個手機或者電腦就只不過提供了一個點.
去收集這些資料.
你看完之後跟它是完全毫無關係的.
你甚至只是希望在當中.
得到一點點的你自己的盼望或者你的慾望的滿足.
所以當今時今日很多人將關係已經弄得不清楚的時候.
再結集一種networking的思維.
就中了魔鬼撒旦的詭計.
變成一種交易式的結果.
而這個交易式的結果我們就在創世記第三章.
第一至第七節的伊甸園的試探裡面.
以及在馬太福音第四章第一至十一節的當中.
耶穌被魔鬼試探的過程裡面就看得一清二楚.
魔鬼所提供的和在伊甸園裡面所提供的.
這條蛇所提供給人的選擇.
都是為了你個人的利益而去考量的.
所以當我們有任何的事情.
只要單單是由個人的出發點去考量的話.
就會造成了一種極端自私自利的狀態.
但是聖經裡面給我們看到的卻是另一幅圖畫.
尤其是今天我們所看的歌羅西書.
從第一章的第十五節至第二十節.
卻是耶穌帶來一個更加不同的一個重新建造.
而這個建造就是讓我們能夠與上帝.
與及人與人之間甚至與這個創造的一個連結.
所以我今天給了這個經文的講題就是.
耶穌是永恆超級屬靈的摺合劑.
或者你可以說萬能膠都可以了.
耶穌基督能夠提供了一個過去人性分裂.
與上帝關係分裂.
與及甚至與這個創造的分裂.
重新將它連結起來的摺合劑.

$^{81}$讓我們一起去看看耶穌怎樣.
透過祂一個獨特的工作.
祂的修補的工作.
能夠重新修補這個充滿破爛破裂的世界.
讓我們一起去看到上帝.
祂怎樣建立一個新天新地的開始.
讓我們一起回到歌羅西書第一章十五至二十節.
請大家等一會一起打開這段聖經.
我們也一起去讀多一次好嗎?.
讓我們熟悉一下這段經文.
因為不是很長而已.
我找到之後我一起去讀了.
一切都是借著他做的.
又是為他做的.
他在萬有之先.
萬有也靠他而立.
他也是教會傳喜之首.
他是原始.
是從死裡首先復生的.
使他可以在萬事上居首位.
因為父喜歡叫一切的豐盛.
在他裡面居住.
既然借著他在十字架上所流的血.
成就了和平.
便借著他叫萬有.
無論是地上的天上的.
都與自己和好了.
在上一堂的講道.
或者上一段的講道裡面.
我們看過了.
信徒和在牧者之間的關係.
不應該是一種好像老鳳上身的情況.
我們反而要經常要帶著一份感恩的心.
來去參與在我們一個屬靈的道路的上臂.
不然的話我們很容易.
造成了一個所謂忘恩負義的人.
不單止忘記了最初我們信主的.
原本的那個恩友和那個推動力.
更加我們否定了.
父神藉著耶穌基督.

$^{121}$怎樣那份體貼父神的心思意念.
而願意死在十字架上面.
滿足父神拯救的心願.
透過聖父,聖子,聖靈三方面的互動.
最終能夠幫助每一個.
認信耶穌基督的人.
能夠得到一份在光明當中的.
產業這份身份.
因此我們必須常常去記住.
父神怎樣去向他在我們身上的作供而帶著感恩.
這一份的推動力.
就成為我們每一個信徒.
怎樣繼續走一條基督徒的人生的路.
而在那個討論的過程當中.
我們也都談論到有關於禱告和代禱的問題.
叫信徒去.
不好意思,我先喝杯水.
在這個過程當中.
我們也都同時去探究到有關於禱告和代禱的問題.
很多時候我們對於代禱或者禱告.
我們只是為了自己個人的意願或者目的.
來求神去滿足.
但是當我們重溫保羅的介紹當中.
我們就會發現到.
原來保羅以身作則做了一個禱告的例子.
就是告訴我們知道原來禱告不是要為自己的意願來禱告.
而是為著上帝的意願.
來落實在我們身上而去發出的.
所以禱告很多時候.
今天的基督徒以為只是和普通其他宗教的禱告一樣.
見有問題我們就去禱告.
禱告這位神因為祂大能大力可以為我們解決問題.
解決了我們就等同於這位神是靈了.
有靈驗的能力.
我們所以就信祂多一點.
不靈的就不信.
這種的思想完全是一種.
我們稱之為民間信仰的一種信仰.
因為他們祈求相信一個神.
只不過是為他們帶來好處.

$^{161}$帶來一個風調雨順.
或者帶來利益獲取的機會.
從來沒有想過你去禱告去相信那個神.
是為那個神自己想祈望做的事而去出發去祈求.
這一點也讓我們看看.
今天我們作為基督徒.
為什麼今天的教會是這麼墮落的原因.
因為教會變成了和其他民間宗教沒什麼分別.
和其他廟宇的信仰沒什麼分別.
信了就回來求神.
求神的這種思維就等同於靈驗.
這個香火會鼎盛一點.
所以為什麼世教會沒人去.
因為沒那麼多資源.
沒那麼多靈驗的機會.
沒錢大教會可以多一點錢財.
可以使你更靈驗.
這種種種的思想扭曲了今天整個教會的世界.
而哥羅西的教會.
正正保羅就以這個成為起始點.
讓他們明白你往後所過的每一個屬靈的生命生活.
都是要為了上帝擺在教會裡面的目的而去努力的.
而在第一章的第十三節裡面.
我們更加看見保羅清楚指出.
我們的父神在他的工作之下.
我們可以得到一個很尊貴的身份.
入到一個新的國度裡面.
成為了在光明當中.
在主耶穌基督所建造的這個新天新地裡面的一份子.
將我們從一個黑暗的世界.
轉入到一個合神心意的世界裡面.
所以弟兄姊妹如果你的心或者信主的目的.
不是為上帝的心願而努力的.
你去到這個所謂新天新地對於你來講是會很悶的.
因為那裡不是滿足你的慾望.
那裡是滿足上帝的旨意的世界.
而在這個過程裡面.
聖子耶穌基督的拯救和他的赦免.
就將我們牢固地去連繫在上帝的身上.
所以今天當我們進入到第十五至第二十節裡面.

$^{201}$我們就會看到這段聖經很特別地用了詩歌體裁的形式來寫.
去介紹一位神所一直愛著的獨生兒子耶穌基督.
究竟他是怎樣的一位神的兒子.
而他的工作也都怎樣去滿足上帝.
重新建構上帝這個好像伊甸園裡面的那份人神關係.
以及人與人的關係.
所以我們看看耶穌怎樣將這些原本受罪所影響破碎的種種關係.
將它重新修復起來.
我們分開兩個部分去看.
第一個部分是從第一章第十五至十七節.
我們會看到耶穌透過他超越的屬性和他的身份.
將我們再次和創造永恆地去銜接起來.
我們去看看這幾節先講些甚麼.
第一章第十五至十七節.
請聽我讀出.
愛子是那不能看見之神的像,是守生的,在一切被做的以先,因為萬有都是靠他做的..
無論是天上的,地上的,能看見的,不能看見的,或是有位的,主治的,執政的,掌權的,一概都是借著他做的,又是為他做的..
他在萬有之先,萬有也靠他而立..
讀到這段聖經的時候.
你會覺得好像很哲學性.
詩歌體裁往往是一些比較有韻味,比較容易有些墨水的人都會覺得聽起來會有一種回甘的現象.
即是你們有沒有喝過好喝的茶,或者吃過甘草.
我最近吃過很多甘草,因為喉嚨不妥.
我今早都下了重藥,所以你見我的聲音可以重開.
吃重藥.
回甘的會聽起來,這段經文,回甘什麼呢?.
我們來看看.
在這裡所說的有關愛子耶穌.
他是一個不能看見之神的像,是守生的,就是一切被造的以先..
弟兄姊妹,當我們讀到這段聖經的時候.
第一樣東西最吸引我們的應該是一幅創造的圖畫.
為什麼是創造的圖畫呢?.
因為首先這裡所說到一切被造的已經很明顯是說到創造的問題.
第二樣東西,這裡說到不能看見之神的像.
現在透過愛子,因為第一個守生的,將祂顯現了出來.
其實這個陳述在創世記裡面,上帝創造人的時候.
是一模一樣的動機.
所以換句話來說,神讓人在祂的創造裡面被創造出來.
是將一個看不到的神的形象,擺在一個看不到的實體的創造裡面.
而人就好像愛子耶穌基督一樣,在這個看不到的形象的世界裡面.

$^{241}$將一個看不到的神,將祂顯露了出來.
而上帝當做人之後,就將人放在這個創造當中.
交託他去管理這個創造.
並且在這個創造裡面,要生養眾多,遍滿地面.
這些說法很熟悉.
因為之前第一章,第三節開始,我們已經說了福音和耶穌基督有這個特性,有這個屬性.
祂是讓我們在這個世界裡面做什麼呢?.
會成長的,會multiply的,會不斷衍生的.
將一個看不到的神,祂原本的動機,現在透過人做了出來.
但是人失敗了.
因為人的失敗,人本身其實是在做上帝的獨生兒子,愛子耶穌所原本應該做的工作.
所以愛子耶穌基督最大的不同之處,和上帝和聖靈不同的.
就是有一個實體存在過在這個世界當中.
祂將上帝的形象,看不到的形象,將祂說了出來.
而從第十五節至第十七節裡面,如果大家有機會去看一下其他的翻譯.
我們就會看到有一個很簡單的句式.
是你讀幼稚園已經會,不過你不知道原來你會.
就叫做verb to be,英文叫做is, am, are.
中文就叫做是.
而且是以現在式寫法,希臘文裡面現在式除了只是指時態之外.
同時間希臘文的現在式是指到一種永恆的真理.
就好像我們常常說太陽從東方升起.
你不會看到太陽從西方升起的.
這個是真理,不會改變的.
但是你不要跟我搏,不要說牧師我用鏡子倒轉看就反轉.
這個是你的問題,不是上帝設計的問題.
這是contrarian的想法,什麼都搏.
上帝不是這樣,上帝做的兒子就是將他真正看不到的樣貌形象實實在在顯露出來.
所以在第十五節開始就已經用這種結構.
耶穌基督是,Christ或者Son is.
或者在原文聖經裡面只是用了一個personal pronoun.
He,He是上文,耶穌基督是那個看不見的神的形象的實體.
而且他是手生的,是整個創造的手生.
所以當我在這裡介紹的時候就告訴我們耶穌是一個獨特的屬性.
他不單止是有上帝的形象,而且可以將上帝的形象化為實體去給我們看到.
而他看到這個這麼特別的屬性.
更加重要一點就是,原來往後所有創造都基於他作為一個模式.
所以創造的最高峰點,其中一個就是上帝做人.
將上帝顯現在創造的裡面.
你和我,如果你都仍然認為自己是人的話.

$^{281}$你的角色就在這個世界應該顯現上帝的形象.
不過很可惜就是罪導致人不能夠滿足了這個神的目的.
所以耶穌基督就要來到這個破爛了的創造裡面.
被罪怎樣去徹底將這些原本破爛了的關係將它修和.
上帝於是乎讓耶穌基督帶著人的樣式來到這個世界裡面.
將這個世界裡面很多人已經將人原本那個人的上帝的樣式已經模糊化甚至拆毀的.
變得不倫不類的狀態,將他再次修補過來.
所以為什麼耶穌要下場,我覺得直接一點,好像打球一樣.
他原本是做創造者,他根本設計就可以了.
好像一個球隊一樣,Eric喜歡踢球,你有沒有見過一個領隊下場踢球的.
放映東,他自己踢,打比賽的時候他當然設計好那場球對於對方是怎樣打.
然後他告訴球員怎樣打,在過程當中調兵遣將,他是做這個工作的,對不對.
他不會同時下場打的,除了放映東.
弟兄姊妹,為什麼要下場打,那班球員不行.
可能放映東行,所以他打了,我們不要高舉他人了.
但是上帝創造這個世界,人搞不定,下場搞不定.
上帝都告訴他,他又不照打,叫你好好治理這地.
你轉過來,就給一條蛇說兩句話,就已經動了你的計劃.
所以這個是個問題,所以怎樣.
耶穌唯有他自己真正下場打.
給你看一個人怎樣真正在創造裡面而活.
但是我們要知道耶穌來到這個世界的創造已經被摧毀到支離破碎的狀態.
他怎樣才可以將他原本的創造有份去決定怎樣設計,怎樣運作的創造被扭曲了.
可以重新修復,將原本見不到神的真正現象展現在這個世界裡面.
而展現在一個罪的世界裡面.
讓真正有上帝悔改認順的人可以回到上帝裡面.
所以在這個世界仍然活著時間去見證到上帝的實體現象.
原本就是基督徒的.
基督徒造成怎樣呢?我一會後再說.
但至少耶穌造一個真正在一個罪惡充滿的世界裡面.
連繫那些人,是會認順上帝的人.
有上帝生命的人,能夠讓他們可以回去和上帝的形象做一個連結.
所以在屬性上面,耶穌基督是展現上帝的形象在這個世界裡面.
並且招聚所有這些有神形象的人可以回來.
而他的身份是一個投生的,而且是怎樣呢?.
投生有什麼特色呢?我們再看第十六節.
因為萬有都是靠他做的.
原來他才是知道那個configuration的人.
我不知道大家喜不喜歡砌模型或者修理電器之類的東西.
我就很喜歡,我最不忍心的就是家裡的東西買回來壞了.

$^{321}$然後有人說要扔掉它,壞了一整個.
大家都知道我的車早前開在高速公路就壞了,燒了副引擎.
其實我沒有放棄過我的車.
不過當他拿去說要弄一萬元再加稅,我真的覺得很心痛.
我沒有放棄過弄它.
但為了燃眉之急沒辦法,要馬上買部車.
但我仍然找人去弄,找盡方法去弄.
終於找到人去弄.
當你看到一件東西壞了之後修復,那份喜樂是很難用筆墨形容的.
上帝就要讓這東西可以修整,所以他派了耶穌來.
但你找人弄,你也要找一個懂得弄的人.
我開什麼車大家都知道,不需要賣廣告.
我唯有找懂得弄這部車的人,為什麼要懂得弄這部車的人.
因為他知道這部車的設置參數是怎樣弄的.
每次我都弄車,很多東西我自己弄,油我自己換.
汽格我自己弄.
但每次弄每一樣東西的時候.
就算拆一粒螺絲,看上去是很簡單的東西.
扭一粒螺絲拆下來,弄完之後放些油,扭一粒螺絲上去,很容易.
不懂的人就說很容易.
你問我們教會的engineer,Alan,衛東這些電腦工程師.
每樣東西都有一個設置的.
你扭一粒螺絲放些油下來之後.
入油之前你上十粒螺絲,是不是扭下去就算了.
不是的.
你問衛東,他很喜歡換車胎的,冬天自己換的.
我真的很佩服他.
換車胎怎麼換,換冬季胎.
很容易,找個range插在四粒或五粒螺絲上扭鬆它.
然後車胎就掉出來,掉出來之後放回車胎上扭上去就可以了.
是不是這麼簡單.
你知不知道每粒螺絲有個磅數.
Alan一定知道,他做一些吊船,洗外牆.
每一粒螺絲有個磅數.
我那輛車正常來計算.
每一粒螺絲的磅數是每一粒螺絲要落81尺per pound.
我很難告訴你,每一磅要落81尺的磅數的扭力.
那個車胎在highway才會掉下來.
我試過付錢給dealership幫我換車胎,我以前試過.
車胎不夠三四公里,我聽到車胎有聲響我就覺得奇怪.

$^{361}$下車一看已經不見了一粒螺絲,然後拿回去給他修.
他說不見了一粒螺絲要你付錢,你修都不穩定要我付錢.
他說不是一粒,後面還有一粒又是不見了.
車胎全部是你換的,為什麼你現在要收我錢換兩粒螺絲.
兩粒螺絲我自己買多少錢?一元一粒,他要收我十七元一粒.
你真是謀取暴利,你不給他又不行.
你要找一個懂得configuration參數是怎樣做的人.
看不看到這段聖經,誰懂得這個參數.
找那個創造的主,愛子耶穌基督.
無論是萬有都靠他做的,無論是天上的地上的,看到的看不到的.
有位置的,有位置的,有主治的,掌權的.
他全部坎坷之道.
所以不單是屬性,是將上帝看不到的形象放在創造裡面.
而且他有能力,他是創造的mastermind,主腦.
他知道怎樣將創造歸回原有的狀態.
他為什麼要找一個這樣角色的人來到,或者愛子來到這個世界呢?.
因為當他這樣來到的時候,就能夠讓這個世界重新歸入一個正常的關係.
什麼叫正常的關係?英文叫做proper order.
有系統的,是按照每一件東西的組合來互相配搭,組合而成.
而互相的運作,所以大兄姐妹,如果我們仔細去看.
在以弗所書同樣保羅的書信裡面就講到.
每一個信徒在這個世界裡面,在教會當中是連結在耶穌基督身上.
做什麼事情?各按各職.
能夠各個肢體按照他的功能去發揮功用,互相配搭.
昨天我上網看了一條片,有一個國家有機械人,二三十個人一起在舞台上打功夫.
原來發覺那些機械人是有人遙控的,也不要緊.
最大的問題是他很多事情是做不到的.
《遙心一變》同一條片裡面就講到,在另一個國家裡面做了另一個機械人.
那個機械人的特點是什麼呢?.
他就是放一隻雞蛋在他面前.
他伸一隻手,機械人的手,拿起一隻雞蛋.
拿到自己面前,另一隻機械人的手,接過去.
去拿那隻雞蛋,然後再將雞蛋放開另一隻手去接的時候.
慢慢放回另一邊的盒子裡面.
介紹那個人就說,看起來就很簡單.
這是一個很複雜的行為.
他說由拿起不碾爛雞蛋,到拿到雞蛋,由一隻手轉去另一隻手.
互相那種配搭協調,而不會壓爛雞蛋.
到他轉手放回盒子裡的時候,不單止對準盒子,而且放下去的時候準確不需要雞蛋掉進盒子裡.
那個過程是非常之複雜.

$^{401}$弟兄姊妹,你知不知道神一早做了你和我.
不需要那麼複雜,祂已經做了這麼複雜的事情.
一隻手拿起雞蛋,另一隻手摺起來,放在另一邊.
因為祂知道參數是什麼.
不是只有表面一個機械人,懂得打功夫.
而是一些很仔細的,能夠了解到上帝心思怎樣去運作每一樣東西的人.
去模仿怎樣去做一個機械人.
回到屬靈的世界.
上帝派回一個老祖明白他整個創造的愛子.
因為這個世界就是由他創造而出的.
一切的東西,他之前已經知道怎樣做.
怎樣建立,怎樣去破除,全部在他管治之下.
當我們能夠在這個世界重新有一位設計師能夠參與在這個世界裡.
這位叫基督耶穌.
他來的目的,我們很多時候在教會說.
他死在十字架上,為我們死,簡單到不堪.
說話的人永遠都是這樣,覺得做事很簡單.
但是耶穌在做這個過程當中,為何要死在十字架上.
就是因為他知道要處理整件毀壞了原本創造的設計的問題所在,就是罪.
而罪要將它撕掉的時候,是一件很大花費能力的事情.
怎樣才能夠修補呢?.
就是要用很大的能力去處理罪所帶來的那種破壞性.
你說那種能力有多大?.
可能你想不到死在十字架上有多慘痛.
不過讓我介紹一樣東西,你一定會知道.
有一次我們在另一間教會講道.
講完道之後,就和一位弟兄坐在一起聊天.
介紹一下大家做什麼,我不用怎麼介紹,我做牧師.
我說弟兄你做哪一行的?.
他說我們在Oshawa一間車廠工作.
接著我就問,你在車廠做什麼工作呢?.
他說我做燒焊.
我說哇,燒焊這件事.
我小時候住樓下做鐵閘的經常都是燒焊的.
我說那你不是很辛苦嗎?因為燒焊是很殘硬的.
他說這份工作是挺辛苦的.
我說燒焊其實也不是很難的.
他說我也是找一支東西點燃就燒了.
他說不是啊,牧師.
我說不是啊?那你說多一點給我知道.

$^{441}$所以你不要說牧師經常和你聊天.
好像查家宅那樣,不是啊.
我想認識你多一點.
他說牧師你知不知道燒焊有一個特色.
我說什麼特色?.
第一件事,如果要燒焊那兩件要被拼接的東西.
一定要同一樣物質的東西.
我說什麼意思?.
燒焊,你焊鐵就要兩塊鐵.
焊的焊料也是要在鐵的材料上一樣.
你不可以焊鐵的材料也是要焊兩塊膠.
焊膠,大家知不知道焊膠?.
我也知道,我買了一個工具回來.
專門焊膠,因為有一次我家的車撞壞了一個泵把.
那我自己弄,弄壞了一個扣.
那你換一個泵把嘛.
焊膠,焊一個扣上去.
你牧師原來懂這麼多古靈精怪的東西.
他說如果你焊兩塊鐵呢.
那兩塊鐵也要同質地的東西.
你一塊是鋁,一塊是鐵.
兩者不可以這樣直接焊上去的.
我說是嗎?為什麼?.
不同質素嘛.
那些東西不同質,不同材料就焊不上去了.
我說原來是這樣的.
他說焊的時候焊的方向和模式也有不同.
這些我不懂了,他就說了一輪.
我就說你焊這些東西都很辛苦.
他說是啊,為了那兩塊鐵你們買的車.
大家知不知道你們的車不是一塊鐵弄出來的.
是很多塊鐵焊在一起的.
他說如果焊得不好.
車子走著走著那塊鐵會掉出來的.
所以焊了之後要有人檢查.
檢查那些焊口是不是完整,是不是合規格.
我說焊東西也有這麼大的學問嗎?.
他說還有.
你知不知道焊接的過程當中.
那些焊料,焊那條焊接的材料.

$^{481}$溫度是多少度?.
我說多少度?.
至少一千多二千度.
我說這麼熱?.
我說是啊,很熱的.
我說為什麼要這麼熱?.
你們開始發覺牧師什麼都問了.
問得這麼仔細.
他說因為不夠熱.
那些焊料和兩塊的金屬不能夠緊緊的融解.
要融為一體才可以將它接合.
不同的接合方式和焊接的模式.
就會影響到最後這兩塊鐵.
究竟可不可以承載這個力量.
原來兩塊鐵打對角焊接和打平焊接是不同的.
我說再問下去我可能是專家.
他說不是專家,你很多事情都不知道.
不過我告訴你,過程當中是很複雜.
各位聽眾,耶穌基督既是創造的主.
要將上帝看不到的形象活在創造裡面.
然後將創造破損,將它重新修補.
因為祂是總工程師,祂知道怎樣去焊接.
知道要怎樣接合.
而祂要焊接,因為祂看見問題是罪.
而怎樣將罪去排除,將人和神和人和創造放在一起焊接.
祂要付出很大的代價,幾千度的熱度去焊接兩塊鐵.
罪所帶來的可怕和傷害,耶穌就有多大的付出去完成.
所以罪帶來的世界是死亡.
所以耶穌要用死亡去補回人和神,人和創造之間這個破損的關係.
你明白為什麼耶穌要死嗎?.
你應該明白了.
因為破損的力度是用死來破損.
所以耶穌就用死在十字架上面,讓我們能夠重新焊接在上帝的創造裡面.
上帝的創造當中,我們原有的神的樣式.
從這一點給我們看見,為什麼耶穌說要將萬有重新在裡面去立,就是這個原因.
因為沒有了這位總工程師,就不能夠將這個原本破爛了的創造.
人與創造的關係,將它恰恰當當地去連接起來.
而這個連接不是短暫的,是永恆的.
不是用著先,頂著當先,而是一個徹徹底底的改造.
修補.

$^{521}$所以鄧小姐,當我們從保羅這個介紹裡面,我們就看見.
為什麼他要介紹耶穌基督,就是看不見神的形象.
為什麼祂是守生的,為什麼祂是創造的那位創造主.
完全的介紹,是讓你和我知道真正我們在這個世界裡應有的樣式.
而哥羅西教會,你們在這個地上,就應該匹配這個角色.
所以上次我們已經講過一個字眼,Cathars,猶如,就是好像耶穌基督一樣的屬性.
和福音的屬性一樣,你不單止是要生根繁殖,並且成熟起來.
而且要multiply,繼續的去衍生.
第二個部分,耶穌祂這個超越的屬性和身份,是將我們和祂直接連繫起來.
我們一起看看第十八至十節,我讀給大家聽,大家留心看著你的經文.
在祂裡面居住,既然借著祂在十字架上所流的血,成就了和平.
便借著祂叫萬有,無論是地上的,天上的,都與自己和好.
除了介紹到耶穌基督的屬性,祂的身份和祂的能力,去怎樣完成和維繫原本上帝的創造之外.
在這個時候,更因為耶穌能夠準確診斷出整個的問題所在,就是人不能夠在罪當中自拔.
因此主耶穌就將祂的焦點放在怎樣將人歸為人應有的類別.
而人應有的類別應該與耶穌基督有相同的心思意念,什麼心思意念?.
上文我們已經講過了,耶穌基督來到這個世界裡面,最重要的目的是讓人能夠好像祂一樣為了滿足父神的喜悅的旨意而去作上帝的事情.
所以人如果能夠像耶穌基督一樣得到和好的話,什麼叫和好?.
英文好講,和好就是你們兩個人之間的心思能夠再一次synchronize,能夠同心合意.
一個人能夠與另一個人和好,就是原本不能夠同心合意.
如果他能夠和好的話,就再次重新同心合意.
耶穌基督希望人能夠活在祂的新的修補了的創造裡面.
最重要的就是你要像我耶穌基督一樣能夠在當中活出上帝喜悅的心思.
上帝喜悅耶穌做的一件事就是上帝想拯救人.
所以祂以死都願意去做,作為基督徒的,屬神的子民的,同樣應該有這份心思意念.
為了上帝喜悅的事我們願意付出代價,而且是不惜為父神喜悅的事而付上代價.
聽清楚啊,這句話不是為教會,不是為牧師,不是為你自己,是為父神的心思意念而不惜付出代價.
為什麼我要這麼重複講呢?因為很多信徒誤會了教會想做的事就等於父神的旨意.
不是啊,不是啊,很多教會做的事不是父神旨意的事情.
不用說哪個天主教會已經是了,基督教會裡面更多.
在教會裡面,今時今日的教會,我直接講一句,只是做自己想做的事.
所以有一件事,在加拿大政府現在草擬將教會的charity position取締.
將教會的宗教取締,不准有免稅的能力,免稅的身份,甚至無法奉獻修具的權柄.
你知不知道在華人教會拆天這麼吵,為什麼?因為很多教會就在裡面用這種特權.
原本政府是對教會的尊重,但教會在這份尊重之下去濫用它的特權,來施傷受受.
我覺得如果政府拆了這個特權之後,而教會的信徒仍然是毀身的話.
我覺得這樣的教會是真的,我不怕再講一句給大家聽.
另外加拿大政府有個法例,叫C9法案,華人同工團契發起所有電郵聯署信.
叫所有弟兄姊妹都簽一份聯署信寄給政府,這份C9法案是講什麼?.
是說不准講一些有歧視性的說話,包括不准講同性戀是不對的事情.

$^{561}$你知不知道教會反應有多大?.
這些是不對的,應該豁免教會,不應該這樣去阻止教會去宣揚同性戀是錯的.
其實表面很風光的他們的說法,其實是在反映他們的內心.
因為他們講的東西,平時講一些合上帝心意的事情,他們根本不敢講.
你知不知道我在同工會寫了一句說話.
我說各位同工,你們這麼激烈地去反對C9法案,不如反過來想想.
C9法案如果順利通過之後,是不是你們港獨和你們所有信徒的見證.
以後不敢講同性戀是錯的事情?.
因為你怕被人告,你歧視是不是?.
我發了這個訊息,發了兩個星期了,一個同工都不敢回答我.
因為他們全部的心思就是要在教會這個空間裡謀取利益.
我說這番話是什麼呢?.
是要告訴我們知道今天的信徒的身份不能夠活出好像耶穌基督一樣.
是以上帝喜悅的心思來作為中心,而是以自己的心思來作為中心.
來去辦教堂,辦一些宗教的機構,叫做事工的東西.
沒力去見證上帝,上帝不喜悅的事情不敢講.
牧師我講了我會沒份工作,我沒份工作我沒得奉獻.
如果你講了之後你沒了份工作的話,那上帝的見證你見證了,上帝會看到.
哪一個看到的重要?你老闆看到的還是上帝看到的重要?.
馬上顯現你對上帝的信是有多少.
然後有些信徒說我小信,聖經都有小信,我是屬於小信.
那你不用信了,是不是?.
耶穌基督祂來到,從這段聖經裡面告訴我們.
祂是教會全體之首,祂是原始,我要特別提出這一點.
保羅講得很清楚,耶穌是教會全體之首.
不只是哥羅西教會,是所有屬神的子民的群體.
耶穌都是他們之首,所以你們馬首是瞻,馬年嘛.
那些童工又亂來,又說馬年要小心.
好像啟示錄那樣,那四頭馬末世已來臨.
我說廢話,末世根本就已經來臨了,你還講這些.
不要在那裡掩掩養養.
如果你要講馬,你應該講我們應該以耶穌基督馬首是瞻.
除此之外,祂是原始,什麼叫祂是原始呢?.
我讀一次英文聖經給大家聽,NET的版本.
He is the head of the body, the church, as well as the beginning.
你什麼都不懂,beginning都不懂,不過你一定懂得一個字眼叫meta.
因為在你的portfolio裡面可能買了很多meta的股票.
什麼叫meta?就是beginning.
一切所有的最根本的意思,沒有祂,其他不能存在.
耶穌在這個世界裡面是一切教會之首,而且是一切最根本最基本的元素.

$^{601}$所以祂從死裡首先復生,讓我們去看到能夠藉著耶穌基督可以同樣像耶穌一樣復生.
勝過死亡,因為死亡是很傷害,拆毀了神和人之間的關係,神和創造之間的關係.
剛才說過耶穌已經收復了創造裡面,而這裡就說到耶穌以死復生來讓我們和祂連結在一起.
如果你想知道這些主題,其實在哥林多前書裡面說得很清楚.
不過我現在不說哥林多前書,留給Joe說,哥林多前書是我其中一個大學的研究書卷.
嚇死Joe了,他說什麼我都知道,不用怕,不要緊的,你有你的得著,我從很多人身上可以學到很多東西.
弟子們,耶穌是個元首,祂從死裡首先復生,使祂可以在凡事上居守.
這裡是說給我們知道,耶穌再次讓我們有一個連結點提供出來.
而這個連結點是什麼呢?因為父喜歡叫一切豐盛在祂裡面居住.
這一句話就說給我們知道,耶穌基督要吸引所有能夠和祂連結在一起的人,都能夠活在上帝的豐盛裡面.
所以基督徒經常說要過得更豐盛,更豐盛,你有沒有連結在耶穌基督身上呢?.
耶穌告訴你怎麼連結,再看第二十節,既然藉著祂在十字架上所流的血,成為和平.
說什麼呢?我知道了,頭半句我明白了,耶穌寫在十字架上,成就了和平,什麼和平呢?.
這個和平的字眼其實和最後天上都與自己和好是兩個同義詞.
這個和平是指與上帝關係修和,因為人犯罪之後得罪的是上帝.
我再講多一次,人犯罪之後得罪的是上帝,所以死在這個世界上是上帝宣判的.
沒錯,是罪大,是上帝的宣判,你看創世記第三章說得很清楚.
上帝宣判了死,與人之間的關係分離了,但是耶穌將它連結在一起.
讓我們能夠與耶穌一起連結,一起與上帝和好,與上帝的心思意念合在一起.
所以我們可以藉著祂,就是耶穌,叫萬有地上的,天上的都與祂和好.
這句說話很熟悉,十六字剛剛講完,無論是天上的,地上的.
原來耶穌基督的創造,後來祂在這個世界的受死,復活,無非都是讓你和我這些肯真正認罪悔改的人.
與耶穌連結在一起,可以在上帝裡面和好,所以如果我們要與上帝和好.
我們必須要與上帝有和平的關係,我們必須要同時與耶穌基督和好.
如果你與耶穌基督沒有和好的,那怎麼能夠與上帝藉著祂可以同樣得到和平呢?.
再講多一次,和好,reconcile這個字.
它的意思,就好像我剛才所講到的,看摺兩塊金屬板一樣.
耶穌將這兩塊原本分離了的金屬板,上帝與我們連結在一起.
但是在和祂連結在一起之前,耶穌要在居間其中.
因此,我們能夠與耶穌基督和好的話,就會連接到祂與父神的和平.
透過這一連串的關係之後,耶穌這麼悽慘死在十字架上的工作就有功效.
能夠將我們不單止與創造連結在一起,而且透過祂我們可以與上帝連結在一起.
到了這裡,我們應該開始明白到保羅為什麼要在這裡將耶穌這麼清楚地介紹.
特別篇章地去介紹給哥羅西的教會知道呢?.
因為保羅已經在前面我們所講到的第一章,談論到有一位同工叫爾巴弗.
知道哥羅西教會已經很努力地去見證上帝.
但是保羅明白到如果你要繼續教會可以達至到基督耶穌的生命.
福音的本質是會能夠繼續更成熟,而且multiply的話.
他就要承接一個很重要的基礎,就是清楚掌握到耶穌基督是怎樣的屬性和工作.
希望藉此再下一城,所以他自己以身作則.

$^{641}$作為一個牧者,他自己來到在哥羅西教會的信徒面前一起作見證.
他沒有為他自己坐牢來禱告,但是他卻以上帝最喜悅的事獻上感恩.
感恩上帝一批這麼好的哥羅西的信徒,願意去繼續成長起來.
而不是停留在一個階段裡面.
又在這裡講到一樣東西,在我們教會裡面有不少弟兄姊妹已經受了浸.
我希望你不要滿足你自己只是受了浸.
我重複講多一次,你不要只是滿足你自己已經受了浸.
而沒有成長,沒有multiply,沒有成熟.
作為一個牧者,我希望你有一個好的屬靈開始之後.
要更加樂見上帝的旨意能夠承傳在你身上.
你的成熟就是你對我作為一個牧者的一個最大的鼓勵.
我看到你如果不生不死的狀態的話,我是很責心的.
求主幫助我們,真的好像哥羅西教會.
和有一位像保羅這樣的牧者,不只是私心滿載在教會當中.
藉著教會滿足一些宗教感的愉悅.
而是真真正正在你的生命生活當中見證出.
那個見不到的神,怎樣實際上在你的生命裡面顯露出來.
求主幫助我們,我們同心禱告.
天父多謝你讓我們能夠在這段聖經裡面學習.
讓我們看到保羅怎樣特別以篇章來介紹你的愛子耶穌基督.
祂的屬性和祂的能力.
祂怎樣為我們帶來一連串的修復工作.
祂怎樣能夠滿有權柄地去付出我們本來應我們去付的罪債.
但是因著祂的付出使得我們能夠和上帝你再次修和.
得以進入一個和平的關係,我們為此獻上感恩.
求主你的帶領,讓我們每一個屬你的子民都能夠真正合乎主你的心.
在你所擺我們所處的時間空間.
見證出一個見不到的神真正的形象.
禱告奉靠主耶穌基督得勝名自救.
\newpage



\section{歷代志下 20:1-3}
\label{sec:031686}
\textbf{絕境中的神操作 (華)}~~~~卓爾君牧師
\newline
\newline
連結: \href{http://www.church.com.hk/acms/content.asp?site=cdc\&op=show\&type=product\&code=031686\&layout=sermon}{\texttt{031686}} ~~~~ 語音日期: 2026-02-22
\newline
\newline
\hyperref[sec:031685]{\small{< < < PREV SERMON < < <}}
~
\hyperref[sec:index_chronic]{\small{[返順時目]}}
~
\hyperref[sec:031687]{\small{> > > NEXT SERMON > > >}}
\newline
\newline
歷代志下 20:1-3
\newline
\begin{longtable}{cl}
\hline
\hline
章節 & 經文 (和合本修訂版)\\
\hline
20:1 & \begin{tabularx}{0.7\textwidth}{X} 此後，摩押人和亞捫人，連同一些米烏尼人來攻擊約沙法。 \end{tabularx} \\ \\ \relax
20:2 & \begin{tabularx}{0.7\textwidth}{X} 有人來報告約沙法說：「從海的那邊，以東有大軍來攻擊你，看哪，他們在哈洗遜‧他瑪，就是隱‧基底。」 \end{tabularx} \\ \\ \relax
20:3 & \begin{tabularx}{0.7\textwidth}{X} 約沙法懼怕，就定意尋求耶和華，在全猶大宣告禁食。 \end{tabularx} \\ \\
[1ex]
\hline
\hline
\end{longtable}
$^{1}$弟兄姊妹平安.
在1940年5月呢.
有一場決定二戰走向的戰役.
在法國的海岸線就展開了.
當時有整整40萬的英法聯軍呢.
被德國的裝甲師層層包圍.
逼退到一個叫敦克爾克的沙灘上面.
那麼背後是茫茫大海.
前面是希特勒這個精銳的坦克軍團.
頭頂還有納粹空軍無情的轟炸.
根據當時英國軍方的數據分析.
在這種插翅難飛的絕境下.
他們最多隻能撤退3萬人.
這就意味著剩下的37萬名士兵.
若非戰死沙場.
就是成為俘虜.
親愛的弟兄姊妹.
在我們的人生旅途當中.
我們是否也曾經遇過.
或者現在就經歷這種敦克爾克的.
一個時候呢.
也許是面對大環境的衰退.
各樣的數據都在告訴你.
危機已經來臨了.
又或者是醫生遞上一個.
冰冷而殘酷的診斷書.
又或者是一個深愛的家庭關係.
瀕臨破碎.
不管你從哪個方向來看.
前有強敵後有追兵.
好像哪裡都找不到出口.
好像被逼到人生的那個沙灘上面.
那麼今天當現實的數據.
判了我們死刑.
當生活的困局把我們重重包圍.
我們該怎麼樣突破.
那個困局征戰得勝.
今天讓我們把時間倒轉回.
3000年前的猶大王國.
我們從歷代志下所記載的一段歷史.

$^{41}$我們去尋找那個破局得勝的鑰匙.
在歷代志下20章1到2節這邊說.
此後磨牙人和雅門人又有米烏尼人.
一同來攻擊約沙伐.
有人來報告約沙伐說從海外.
亞蘭那邊有大軍來攻擊你.
如今他們在哈西遜他們就是影擊敵.
這個故事的主角是猶大王約沙伐.
假象一下這可能是一個美好的早晨.
就像今天一樣.
這個國王剛喝完第一口咖啡.
結果大臣就臉色蒼白衝進來.
報告國王出大事了.
磨牙人雅門人連米烏尼人都來了.
三國聯軍已經殺到大門口了.
弟兄姊妹這個不是一般普通的邊境衝突.
那個叫復仇者聯盟已經來了.
以當時的猶大國的兵力.
打一個國家都打不贏啊.
更何況一次來三個.
這個在軍事專家的眼裡只有三個字叫死定了.
所以面對這種滅國等級的災難.
大魏的子孫約沙伐他第一個反應是什麼.
聖經告訴我們在歷代志下20章第3節.
說.
約沙伐怎麼樣.
便懼怕.
這三個字給我們極大的安慰啊.
聖靈沒有把屬靈的偉人描寫成不怕死死不怕的超人.
約沙伐他沒有什麼超級信心啊.
他沒有高舉雙手高喊說Hallelujah.
他的心裡大概只有我的媽呀.
把他嚇到靈魂出竅可能這個頭上的皇冠都歪掉了.
但親愛的弟兄姊妹其實懼怕.
是人類面對危機的一個原廠設定啊.
所以跟你旁邊的弟兄姊妹講.
會怕很正常.
所以不要覺得說自己信心不夠啊.
連神所重用的國王約沙伐都會怕呀.
你怕一下有什麼關係.

$^{81}$所以在跟你旁邊的弟兄姊妹講.
怕一下沒關係啦.
但請記住恐懼是人類的本能.
但信心是聖徒的選擇.
恐懼不是終點.
那個是一個信心的起點.
約沙伐之所以可以成為傳奇.
不是因為他不會怕.
而是在他怕得要死的時候做了一些關鍵的決定.
這些決定讓他從一個會怕的凡人.
變成了一個得勝的君王.
所以今天我們要從約沙伐的故事來學習三個.
絕境中的神操作.
當你的人生卡關當你完全不知道怎麼辦的時候.
這三招會幫助我們去打贏戰爭.
第一招叫承認無能就啟動大能.
跟你旁邊講承認無能啟動大能.
首先面對大軍壓境約沙伐做了一個看似政治自殺.
但是屬靈上是最強大的舉動.
在歷代志下20章3到5節.
約沙伐便懼怕定義尋求耶和華在猶大權力宣告進食.
於是猶大人聚會求耶和華幫助.
猶大各城都有人出來尋求耶和華.
首先約沙伐宣告進食.
尋求耶和華.
然後在歷代志下20章12節.
他講出了整本舊約聖經當中最辛酸.
但也最讓人想按讚的一句話.
在歷代志下20章12節.
這邊說因為我們無力抵擋著來攻擊我們的大軍.
我們也不知道怎樣行.
我們的眼目單仰望你.
請注意這裡有兩個關鍵詞.
第一個叫做我們無力抵擋.
no power.
我們打不過.
國防預算不夠.
導彈不夠.
精力不夠.
然後他說我們不知道怎樣行啊.

$^{121}$no clue.
我的腦袋一片空白.
完全沒招.
不知道該怎麼辦.
請問你啊.
這個是一位國王該說的話嗎.
按照世界的邏輯啊.
一個總統一個國王就算心裡怕得要死.
表面上要裝得胸有成竹嘛對不對.
就要說針對現階段的事態發展.
我方已啟動的標準化預案.
並完成動態風險評估.
目前已經達成跨部門資源整合.
整體態勢正處於高度受控的常態化管理階段.
反正講一大堆廢話.
然後自己已經偷偷打包要逃到國外去了.
但是約殺法.
他沒有.
他站在全國的百姓面前.
你注意十三節講的.
遊大眾人和他們的嬰孩.
妻子兒女都站在耶化面前.
表示說這個國王是擔在全國人民.
甚至小孩子的面.
他承認說我不行了.
只能仰望上帝.
為什麼這個示弱如此重要.
因為約殺法啟動了一個屬靈的定律.
信心不是說挺起胸膛對自己說我可以.
信心是跪下來誠實對上帝說我不行.
弟兄姊妹.
我們現代人最大的問題就是太會裝.
我們假裝自己穩如泰山.
然後每天都在喫胃藥.
我們說常常喜樂.
然後常常躲在房間裡面流淚.
我們以為信心就是要撐住場面不能讓別人看出我很軟弱.
否則上帝會覺得我很沒用弟兄姊妹錯了.
大錯特錯.
這個是一個很簡單的屬靈數學題.

$^{161}$當我認為我可以靠自己的時候那個算式裡面就不需要上帝了.
但是當我意識到我不能.
我破產了我崩潰了我無能為力了那個算式裡面就會有上帝的位置.
這就是為什麼上帝在哥林多後書對保羅說出那個震動天地的應許說.
我的恩典是夠你用的.
因為我的能力.
在人的軟弱上顯得.
完全.
如果今天我們自認無敵.
就變得自給自足.
但是在屬靈的世界裡面.
自給自足叫做死路一條.
因為只只離開葡萄樹結局就一個枯乾被火燒.
只有當我們開始承認自己軟弱的時候.
我們纔會開始依靠上帝啊.
我必須很誠實的告訴大家這個不是一夜之間就會學會的功課.
在我的人生當中我與軟弱的關係經歷了一個漫長的摔跤啊.
這大概可以分為五個階段.
第一個階段叫無知.
當我年輕氣盛的時候我覺得自己好極了.
我的判斷最有智慧我的眼光獨到我又聰明又屬靈又有恩賜.
簡直是上帝送給世界的禮物啊.
那時候的我不知道什麼叫謙卑啊.
只知道什麼叫驕傲.
第二階段叫發現啊.
後來我結婚啊.
弟兄姊妹婚姻是什麼啊.
婚姻就是一面4K高畫質的照妖鏡啊.
如果你想發現你自己都不知道的軟弱.
你去結婚就對了.
透過另外一半我突然發現自己有一大堆軟弱啊.
自我中心大呼大聲啊.
不體切妻子的需要啊.
那個鏡子就照得我無地自容啊.
然後進入第三個階段.
叫隱藏跟否認.
所以一開始我的本能就是藏啊.
當藏不住的時候我就變了我就否認了.
我哪有生氣啊.
我是講話丹田比較有力而已.

$^{201}$各種否認要拼命保護那個虛假的完美牧師形象啊.
到第四階段.
叫接納了.
終有一天我好累啊.
演戲太累了.
我纔會停止找藉口啦.
我在上帝跟家人的面前長嘆一口氣.
對啦這就是我啦.
我有這些限制我常常搞砸我有軟弱我需要幫助啊.
但弟兄姊妹這個不是終點.
當我們開始承認跟接納的時候.
會進到第五個階段.
叫誇耀.
最後我終於摸到保羅所說的.
所以我更喜歡誇自己的軟弱好叫基督的能力.
撫僻我.
誇耀自己的軟弱不是說我一直活在軟弱裡面.
而是誇耀上帝的大能現在成為我生命的力量.
讓我可以重新站立起來.
重新得力啊.
約沙伐也是啊.
當他在承認軟弱的那一瞬間.
上帝的能力就啟動了.
當我們開始承認自己的無能.
就是依靠上帝能力的開始.
想想看啊.
你開車.
用GPS導航.
什麼時候你最聽話會去聽導航的聲音.
絕對不是你覺得這條路我很熟的那個時候.
你開在很熟的路上你會把導航的聲音關掉就吵死了.
只有當你覺得.
迷路了.
你會把車停在路邊.
設定導航.
然後就開始乖乖的聽他每一個指令啊.
所以當約沙伐說我們不知道怎樣行.
我們迷路了.
接著他就說.
我們的眼目.

$^{241}$單.
仰望你.
我要看著導航.
我要聽著導航.
你可以想象當時的畫面.
在十三節描述全家大小連.
話都不會說的嬰孩都被抱出來.
他們什麼都沒有做.
沒有磨刀沒有練兵.
只是站在那裡.
眼睛就直直的.
看著聖殿.
看著上帝.
這個就像是一個被在外面被欺負的孩子.
他哭著回家找爸爸.
他甚至都不用說話.
他只要站在爸爸的面前用那個充滿眼淚無助.
但是卻專注的眼神看著爸爸.
那個眼神就在說爸爸我不行了.
我被欺負了.
你也看到了.
我告訴你.
很少有父親.
能拒絕這種眼神.
而弟兄姊妹.
這種徹底無助.
卻完全信靠的眼神.
就是上帝最無法拒絕的禱告.
親愛的弟兄姊妹.
如果你現在正處在一個不知道該怎麼辦的狀況.
我邀請你.
不要再試圖用自己的力量去硬撐了.
回到家.
關上房門做一個約殺法式的禱告.
你就看著神.
你說神啊.
在這個事情上面我徹底失敗了.
我管不了我的孩子我救不好我的工作我治不好我的病我看不到路.
我投降了.
但我把我自己.

$^{281}$交給你.
換你上場.
我只看著你.
當我們開始承認自己的無能為力.
神的大能就被啟動了.
然後再就是第二招.
交出主權神掌王權跟你旁邊講交出主權.
神掌王權.
當約殺法跟百姓承認自己無能之後.
現場就發生一件事情.
聖靈在會眾的當中領到一個立位人.
叫做亞哈西.
在歷代志下20章14到15節那時約華的靈在會中領到立位人亞薩的後裔.
馬泰亞的玄孫耶利的曾孫比拿亞的孫子撒迦利亞的兒子亞哈西.
他說猶大眾人耶路撒冷的居民和約殺法王你們請聽.
約華對你們如此說.
不要因著大軍恐懼驚慌.
因為聖拜不在乎你們.
乃在乎誰啊.
神.
亞哈西說聖拜不在乎你們乃在乎神這句話我們聽過無數遍.
但是我們往往忽略其中的意義.
亞哈西所宣告的.
是一份所有權轉讓書啊.
我講人話就是這個戰爭的主角換人了.
在五分鐘之前.
這個是約殺法的戰爭.
他是國王他要負責調兵遣將負責糧草他要負責輸贏.
所以他壓力山大.
因為責任在他肩上.
但五分鐘後.
神說停.
我接手了.
從現在開始我是主角你是觀眾.
各位啊這個就叫做聖拜不在乎你們乃在乎神啊.
那麼請問.
既然這個是神的戰爭.
所以該失眠的是誰啊.
是神嘛對不對但是我們的神不會失眠嘛.
但既然責任方已經變成神了.

$^{321}$戰爭不用我打了.
那我幹嘛失眠呢.
既然聖拜不在乎我.
我幹嘛胃痛呢.
想象一下如果今天我欠了一千萬債務.
討債公司天天上門恐嚇.
我每天愁得半死.
突然有一天有一個億萬富翁走過來拿出一張支票.
當著我的面簽好字交給債主然後他對我講.
這筆債算我的.
我已經跟他們談好了從今天開始這個債務.
轉到我的名下以後他們只會找我不會找你.
請問.
那一瞬間你會做什麼.
你會繼續哭嗎.
不會.
你會繼續發抖嗎.
不會.
你會馬上打電話.
叫一杯全糖的珍珠奶茶.
管壓壓驚.
為什麼.
因為責任轉移了嘛.
你自由了嘛.
這個就是雅哈西所宣告的好消息.
上帝已經買單了這個戰爭啊.
但這個是不是代表我們就回家睡覺躺平啊.
能不能躺平啊.
跟你旁邊講不要躺平啊.
事實上神給了一個非常唯妙.
很精確的指令.
歷來之下二十章十七節猶大和耶路撒冷人啊.
這次你們不要征戰.
要什麼擺陣戰者.
看耶和華為你們施行拯救.
這裡有兩個動詞界定了我們的責任跟神的責任.
第一個叫做不要征戰.
Do not fight.
這個是神的責任.
擊殺敵人改變局勢扭轉人心這個是神的工作不要插手.

$^{361}$也不要想用我們的血氣去幫忙.
然後要擺陣戰者stand firm.
這個是我們的責任.
我們不逃跑不崩潰.
站好那個鼠鄰的位置.
做好那個預備.
那麼什麼叫做擺陣戰者你想想看那個畫面很恐怖啊.
當這個敵人的戰車撞過來殺身震天神卻叫你站在那裡.
不準拔刀不準逃跑.
你只能張開眼睛看神怎麼修理敵人.
哇這個要很大的膽量我告訴你這比衝上去拼命還困難啊.
但這個叫做安息的勇氣.
神要我們在風浪當中站穩.
用行動證明我一步也不退我做好我的預備我相信神會出手.
但我不會越界去搶神的工作.
親愛的弟兄姊妹很多時候我們活得很累啊.
不是因為我們做的太少.
而是因為我們越界.
我常說我們基督徒有時候像一個領著基本工資的工讀生啊.
然後整天在操董事長的心啊.
這個不叫負責這叫自不量力啊.
我們沒有在劃分清楚什麼叫我的擺陣什麼叫神的征戰啊.
許多時候90\%我們的失眠都是我們試圖把神的征戰變成我們的征戰啊.
比如說在找工作的時候.
我的擺陣就是把履歷表修到連我媽都不認識的完美.
面試把牙縫的菜渣剃乾淨我準時出席穿著得體我持續進修我的技能.
這個是我的擺陣啊.
神的征戰就是面試官當天會不會剛好跟老婆吵架心情不好.
大環境的景氣會崩盤還是復出.
哪一扇門會關上哪一扇窗會開.
這個是神的征戰.
但我們的越界是什麼.
我面試完啊回家開始三分鐘刷一次email.
我們開始去搜面試官的臉書.
然後開始想哎怎麼還沒通知我啊.
我怎麼等了兩天還沒有消息我該不該再寫信.
去問啊.
弟兄姊妹那個是神的業務範圍啊.
我們負責精彩亮相.
神負責安排工作不要越界.

$^{401}$面對生病的時候我的擺陣是什麼.
準時喫藥嗎.
多喫花椰菜少喫鹹酥雞嗎.
保持喜樂的心嗎.
而神的征戰是什麼.
那個指數會不會降下來癌細胞會不會消失.
藥對我有沒有效我還能夠活多久.
這個是神的征戰.
當我們常常越界啊三更半夜不睡覺把身體感覺到一點點不對勁.
我們趕快上網查.
查完之後就覺得自己得了絕症啊.
親愛的弟兄姊妹焦慮殺不死癌細胞的但會殺死你的睡眠.
把征戰的主權交給上帝.
我們只負責把身體照顧好剩下的神會看著辦.
教育孩子也是個重災區對嗎我的擺陣.
就是按著真理教導他花時間陪他玩為他禱告他考0分的時候我忍住.
不要吼他.
神的征戰.
他的個性會不會改.
什麼時候會浪子回頭他將來會不會遇到好對象這是神的工作範圍嗎.
但我們的越界常常我們試圖用碎碎念來感化我們的兒女啊.
但請記得啊.
我們只負責撒種跟澆灌.
是神負責叫他生長.
所以弟兄姊妹我們搞清楚擺陣和征戰我們就不容易崩潰了.
我們最大的問題就是喜歡當一個碎碎唸的副駕駛.
這個神在開車我們在旁邊一直碎碎念.
前面紅燈踩剎車踩剎車前面右轉.
啊主你開太慢了要不要換我來開啊.
弟兄姊妹啊當憂慮來襲到我們心跳加速的時候.
請問自己一個問題就是我現在.
所擔心的這件事情.
是我的擺陣.
還是神的征戰.
如果是神的征戰.
就請你深呼吸.
然後對神說主啊這是你的業務範圍啊.
我不搶你的工作這場戰交給你了.
第三招.
讚美先行仇敵歸零跟你玩一場讚美先行仇敵歸零啊.

$^{441}$承認無人責任轉移接下來約沙法做了一個戰爭史上最荒謬的擺陣.
第二天清晨軍隊出發了通常走在軍隊最前面的是誰通常是拿盾牌的那種.
重裝步兵嗎或者是海軍陸戰隊嗎.
但是歷代志下20章21節記載.
說約沙法既與民商議了就設立.
歌唱的人送戰耶華使他們穿上聖潔的禮服走在軍前讚美耶華說.
當春瀉耶華因他的慈愛永遠長存啊.
弟兄姊妹啊.
走在軍隊最前面的不是坦克車而是師班啊.
這群人穿的不是防彈背心啊是那個絲綢的禮服啊.
他們手裡拿著樂器連一把菜刀都沒有他們是敬拜團啊.
約沙法把他們放在軍隊的最最最最前面啊也就是所謂的炮灰區啊.
想想看如果你是敵人的這個指揮官看著這個猶大軍隊塵土飛揚的列隊你會趕緊.
拿起望遠鏡看一下下一秒你就開始懷疑人生了.
是我眼花走在最前面那些人拿沒拿盾牌啊.
沒拿刀啊他們拿什麼拿麥克風啊.
哦還拿著吉他那個彈貝斯的那個男的怎麼還閉上眼睛那麼陶醉啊.
那個搖領鼓的女的會不會笑得太開心啊.
這個敵人的指揮官轉頭問他的副官這些這個國家是嚇到瘋了是吧.
這個是戰場不是ktv耶.
弟兄姊妹啊.
在軍事專家的眼中這叫自殺是種攻擊啊.
你想想看師班是沒有盔甲的對吧.
一件事過來.
我那個男低音就變成男高音了的.
但是約殺法明白一個連軍事專家都不同的奧祕.
敬拜讚美是最厲害的武器啊.
要知道我們人類的邏輯通常是這樣的啊.
打贏了啊開香檳切蛋糕.
然後讚美神對吧很合理看到結果才讚美.
但屬靈的邏輯完全相反.
是先讚美神.
神就動手瞭然後我們就打贏了.
大多數人是看見結果才讚美那個叫反應.
但約殺法教我們還沒有開打先讚美.
這個叫做信心啊.
弟兄姊妹讚美不是為了慶祝勝利.
讚美是為了帶來勝利啊.
請看二十二節記載一個神祕的時間點.
眾人方歌唱歌讚美的時候.

$^{481}$奚爾華就派伏兵擊殺那來攻擊尤大人的雅門人摩鴉人和希爾山人.
他們就被打敗了.
你注意方這個字啊when they begin to think.
就在那個do re mi那個第一個音符唱出來的時候.
那個第一個鼓聲敲下去的瞬間.
上帝就出手了.
他就派伏兵擊殺敵人甚至讓那個結盟的三國聯軍內訌互相殘殺.
弟兄姊妹為什麼讚美這麼厲害啊.
因為讚美就是信心的最高表現啊.
當我們在大能臨頭的時候快要死掉的時候還能夠唱歌.
魔鬼就慌了.
魔鬼想說這個人怎麼回事啊.
照理說他應該要哭啊他應該害怕他怎麼在唱詩歌啊.
他是不是掌握什麼我不知道的絕招啊.
弟兄姊妹.
這個是一個鼠鄰的奧祕.
讚美會讓仇敵混亂.
你的抱怨會引來魔鬼.
但你的讚美會建立神的寶座.
你的讚美在臨界就像轟炸機一樣.
要炸開那個鼠鄰的全屍.
炸開一個缺口.
讓神的大能介入.
那麼結局如何.
敵人全軍覆沒.
我們來看歷代志下20章24到25節說.
約沙法和他的百姓就來收取敵人的財物在屍首中建了許多財物珍寶.
他們撥脫下來的多得不可期待.
因為甚多隻收取了三日啊.
聖經說他們花了整整三天才撿完那個戰利品因為太多了.
那個地方原本叫戰場.
後來被改名叫做比拉伽谷.
歷代志下20章26節第四日眾人聚集在比拉伽谷在那裡稱頌耶和華.
因為那地方名叫比拉伽谷直到今日.
比拉伽的意思是什麼就是稱頌與祝福啊.
哇那個劇情大反轉.
本來是要滅國的悲劇結果因為敬拜讚美變成發大財的喜劇啊.
本來是來殺他們的敵人.
結果變成送快遞的.
把禮物送到他們的家門口啊.

$^{521}$苦難.
會在讚美當中變成一個祝福的容器啊.
在新約使徒行傳第16章當.
保羅和希拉被打個半死關在監牢裡面的時候.
半夜他們做什麼.
抱怨嗎不是.
他們禱告唱詩讚美神啊.
結果地大震動尖門全開鎖鏈脫落.
在教會歷史當中有個蜀靈偉人叫做蓋恩夫人.
他16歲的時候在父母的安排下被.
嫁給比他大22歲的蓋恩先生.
這一段婚姻對他來說是長達12年的折磨.
他的婆婆性格對他非常的嚴厲經常在丈夫面前挑撥離間.
甚至連他家中的僕人都可以去羞辱他.
而他的丈夫也無法理解他的靈性追求.
時常限制他的信仰生活.
更慘的是原本蓋恩夫人有美麗的外貌.
但是在22歲的時候他染上天花.
他的容貌盡毀.
雖然這個外在沒有什麼溫暖.
但他在內在裡面跟神有很親密的交通.
他發現當他停止抱怨.
轉而讚美的時候.
那個受委屈的痛苦就消失了.
所以後來他就提倡信徒.
不管你的階級或教育程度.
你都可以自己透過讚美跟禱告.
來直接經歷神的同在.
但不幸的是這被當時的羅馬天主教的權力中心.
就認為是煽動信徒可以越過所謂的聖職人員的地位.
所以他就被定為異端.
然後被長期監禁在地牢裡面.
但是就在那個陰暗潮溼與世隔絕的監牢裡面的時候.
他寫下了一個震撼後世的詩歌.
他說我是一隻關在籠中的小鳥.
遠離人群與天上的寬廣.
雖然我只能在籠中歌唱.
但我卻是為了我主而唱.
弟兄姊妹.
蓋恩夫人雖然她的身體被困住了.

$^{561}$但是她的讚美.
讓困境變成神用來得勝的環境.
而且她的生命影響了後世許許多多的信徒.
最後讓我們回到1940年的5月.
當40萬的英法聯軍被困在沙灘上面等待死亡的來臨的時候.
當時英國國王喬治六世下達一個命令.
就是5月26日為全國禱告日.
所以當天的照片記錄了全英國成千上萬的人都在教堂外面排隊.
他們不只是在禱告.
而是集體向神獻上敬拜.
當全國人民一起敬拜禱告的時候.
馬上帶來三個至今無法用邏輯解釋的結果.
首先希特勒.
在毫無預兆的情況下就下令.
裝甲部隊停止前進三天.
這個在軍事上完全不合理.
三天的停頓讓這個英法聯軍可以爭取時間撤退.
其次.
連海面都出現變化.
那個英吉利海峽即使在夏天呢.
強風也會讓船無法靠岸.
但當他們全國敬拜禱告之後這個海域出現了罕見的.
極度平靜.
這個讓800多艘用平民駕駛的救生艇漁船甚至遊船.
都能夠橫跨海峽.
第三個更關鍵的是突然那個低空.
就出現了濃厚的雲層跟霧氣.
就阻礙了德國的空軍進行精準的轟炸.
結果原本只能預測救出3萬人.
最後他們救回了33萬8千人.
歷史學家為這個爭論了幾十年不知道為什麼會這樣.
但對基督徒來講很清楚這是出於神的作為.
這個就是師班走在前面的見證啊.
親愛的弟兄姊妹你現在的心情如果很糟啊糟到我都不想來教會不想聽詩歌.
覺得上帝很遠.
請聽我講.
這正是你最需要唱詩歌的時候.
如果你要等到你心情好才唱那你要等很久啊.
你要現在就學會用讚美進行攻擊啊.
你可以找一首你最喜歡的詩歌然後開口宣告主啊雖然我看不見出路.

$^{601}$但我讚美你.
我不看環境我看你.
我相信這場仗你已經打贏了.
我要在敵人的面前為你唱一首得勝的歌.
當你的聲音響起.
神就要派伏兵擊殺你的仇敵.
擊殺你的憂鬱擊殺你的恐懼擊殺那些攔阻你的障礙.
因為神說勝敗不在乎我們.
而在乎.
神.
我們一起來禱告.
有時候我們真是落在一個光景的當中.
我們感到無助.
我們感到絕望.
我們不知道前面該怎麼辦.
主啊有時候在我們的生命裡面.
我們感到一個深深的無力跟疲乏.
我們甚至會說主啊你在哪裡.
但主你對我們說.
我一直都在.
主啊有些時候我們甚至不敢來到你的面前.
因為當我們看到我們自己的軟弱我們自己的罪.
我們害怕來到你的面前.
但主你說孩子們不要害怕.
不要害怕.
你可以來到上帝的寶座面前.
不是因為我們做得好或做得不好.
而是因為我們是神的兒女.
我們是神的兒女.
主啊當我們來到你面前的時候.
我們只要單單的仰望你.
雖然我們流著眼淚看著你.
主你都知道我們的困境.
主你都知道我們的需要.
主你就要心情為我們來掙扎.
主我們謝謝你.
因為我們深知.
不管在什麼樣的困境裡面.
不管我們在什麼樣的疲乏當中.
主你都會賜給我們能夠讚美的力量.

$^{641}$當我們開始讚美的時候.
當我們開始敬拜的時候.
主你所拆開的伏兵就出現了.
當我們開始敬拜的時候.
那個天敞開.
主你所拆開的那個天使就出現了.
你要帶領我們走過死癮的幽谷.
你要帶領我們走過那個絕境.
要進入到主你所預備的豐盛當中.
主我們謝謝你.
我們要向禰說我們要繼續來敬拜.
我們要繼續來讚美.
我們繼續來歌唱.
今天當我們歌唱的時候連石頭都要跳起來了.
那黑暗的泉石要破碎.
我們要看見主的光就在我們前面行.
主我們謝謝你.
將禱告奉祖耶穌基督的聖名.
Amen.
\newpage



\section{馬可福音 16:15}
\label{sec:031687}
\textbf{張瑪莉突破生命困境分享}~~~~張瑪莉女士
\newline
\newline
連結: \href{http://www.church.com.hk/acms/content.asp?site=cdc\&op=show\&type=product\&code=031687\&layout=sermon}{\texttt{031687}} ~~~~ 語音日期: 2026-02-22
\newline
\newline
\hyperref[sec:031686]{\small{< < < PREV SERMON < < <}}
~
\hyperref[sec:index_chronic]{\small{[返順時目]}}
~
\hyperref[sec:031688]{\small{> > > NEXT SERMON > > >}}
\newline
\newline
馬可福音 16:15
\newline
\begin{longtable}{cl}
\hline
\hline
章節 & 經文 (和合本修訂版)\\
\hline
16:15 & \begin{tabularx}{0.7\textwidth}{X} 他又對他們說：「你們往普天下去，傳福音給萬民聽。 \end{tabularx} \\ \\
[1ex]
\hline
\hline
\end{longtable}
$^{1}$我經常都覺得到我這個年紀.
還可以走走跳跳 這裡還挺清醒的.
其實真的全部都是恩典來的.
我來了幾天而已 感恩又沒有什麼jack leg.
所以都挺厲害的.
在這裡再和你們拜年 今天都說了是年初六.
希望每個人都是平平安安 喜樂 一家人齊齊整整.
很多馬馬馬的 龍馬精神之類的.
我經常都說 今年我的口號是馬上信主.
要快要快的 很開心的.
來到每次Cross Point Church 給我的感覺好像回家一樣.
我每次來到三藩市 我通常都是來這個教會分享.
今天的主題是牧師給我的 我覺得挺好的.
叫做「突破生命困境」.
其實我們人生裡面 每個人都會遇到很多很多事情.
我自己覺得我是一個比較阿Q精神的人.
除此之外 我想我是因為生命裡面有神.
我是1986年感恩 洗禮的基督徒.
我又不是傳道人 什麼都不是.
我只是好像你一樣 一個平信徒.
通常星期天我都是回教自己的教會.
我是在香港的港福堂.
在這裡我做了很多新年祝賀的事情.
這張是我自己的公司 我公司還沒倒閉.
三十多年還沒倒閉 我是做培訓的.
我不開班的 我主要是做一些企業形象.
有機會和你們在這裡 希望有機會可以分享.
「二馬內利」就是「神與我們同在」.
在這段時間 我覺得希望大家都是平平安安的.
就選了一大堆這些東西.
十字架在我們的心裡面.
很多時候我們真的不是需要掛著十字架.
這個就是你們的教會 我很喜歡你們的教會.
牧師他們開始決定要這個地方的時候.
我看 我沒見過一個教會的燈這麼漂亮的.
是不是?多有藝術感 真的很棒.
希望你們不要拆 很漂亮.
這個是這個教會的特色來的.
我想你去全世界的教會都未必有.
很棒 很棒.

$^{41}$那我為什麼今天來這裡不去.
我不知道你會不會打麻將 我現在準備學.
準備學打麻將 當然不是賭錢的.
我覺得要讓這裡繼續清醒.
可能我真的要學一下這個遊戲.
那你哪個教友是不賭錢的.
又喜歡打麻將的就找我.
我真的要學這個東西.
因為他老闆說這個東西一定要學的.
不然到時候你就知錯了.
我又很喜歡這段金句.
因為耶穌提醒我們 我們要去傳福音.
尤其我們自己是基督徒.
有沒有基督徒舉手看看 嗨 祝福你啊 祝福你啊.
真是感恩 其實我們基督徒其中一樣東西.
就是我們蒙福.
我們應該將這個福氣去傳出給其他人的.
因為很多人可能沒有機會聽福音.
所以來到這裡 今早是很早起床的.
你知道你要化妝 又練假眼陰毛.
很多事情要做 很多事情要做的.
所以搞了一段時間 很早就來了.
但是我自己覺得很感恩.
就是好像回家的感覺 有沒有.
這個是我屬靈的家.
我希望你有種回家的感覺.
這段話就是我們要傳福音 傳到地極.
我真是到處都去 譬如加拿大.
我去到Halifax那些地方.
去年我曾經11個月沒有停過飛.
兩次紐約 一次多倫多.
還是剛剛去到多倫多是下大雪的.
但是很漂亮 很漂亮.
這個是我自己覺得除了工作.
除了我家庭各樣東西之外.
是我自己心裡面很感恩.
是時時都有機會做的.
我是常常都出去的.
有沒有人知道我是香港小姐.
幸好這件不是錯事.

$^{81}$如果做錯事就慘了.
大家都知道我是1975年的香港小姐.
你找旁邊的手指腳指混在一起.
都不夠數.
其實原來我已經選了51年.
是57年 不是出世.
你想一下 到現在這個時間.
我覺得自己是一個很喜樂.
還有很蒙恩的人 真是很感恩.
說給你聽都不相信.
我是沒有老花的.
我現在和老朋友去飲茶.
經常都去我家看吃什麼點心.
我沒有老花 我又沒有什麼近視.
很棒的.
所以天父恩典.
待會你會知道天父為什麼這麼疼我.
讓我沒有老花.
這些就是以前的照片.
我經常都說你年紀夠大.
如果可能你還沒有認識我.
我其實經常都來這個教會的.
到你差不多年紀大.
就很多人給你東西.
所以剛才主席介紹我是院士.
其實這些你到差不多年紀.
你信了 不死就有的了.
他們會給你很多東西.
我又拿了很多獎.
十大傑青之類.
其實我由1975年選為香港小姐.
得到很多機會去事奉.
有很多機會參與社會不同的活動.
其實我很久已經沒有在香港做任何的公職.
但最近我是一位老師.
如果你有看香港的新聞.
HA就是Hospital Authority醫管局.
原來它的主席是我以前的老師.
在一次的活動裡面被它抓了出來.
現在我是紅十字會的管理員.

$^{121}$學到很多東西.
因為香港的醫療制度真的很好.
還有在紅十字會裡面我才知道.
原來全香港的血都是靠它那裡供應的.
我多無知.
我以為每間醫院都有.
今天我就選了三個故事.
但我不會說得很長.
因為我不是傳道.
我不是牧師那些.
但這三個故事.
我就請勁派隊上來.
因為他們今天會幫我.
你知道我年紀大了.
暫時未洗地參附.
不知道何時.
我去年還跟你們說.
我弄傷了膝蓋.
我是那些不太會認命的人.
醫生叫我鑽兩個洞.
然後微創.
然後就刮刮刮.
刮了半月板受傷撕裂的地方.
然後不知道擠些甚麼下去.
聽起來也害怕.
我沒有做那些事.
我就去針灸,拔罐.
然後放血.
然後物理治療.
然後甚麼機能再生.
所有你想得到的都做齊了.
你見我走上來.
我沒有咳嗽.
這就是天賦醫治.
所以不要怕.
有些甚麼我們求.
真的跟天賦說.
尤其是我們幫我們的神恩仔.
我現在抱了一個孫這麼多.
我經常都幫他祈禱.

$^{161}$我經常都說.
阿神你給我孫.
真是一輩子跟隨你到老都不偏離.
一開始我想唱一首歌.
暖和一下.
我選了一首歌是《活出外》.
好嗎?一起唱.
可以一起唱的.
擦不透 摸不著 卻看得見.
因為有你 因為有我.
甘心給你用心愛.
把心中這一份愛 活出來.
耶穌的愛 激勵我.
放開我的生命.
讓自己成為別人祝福.
耶穌的愛 點燃我.
心中熊熊愛我.
我們一起向世界 活出愛.
要同座唱一起啊.
有一份愛 從天而來.
比山高 比海深.
擦不透 摸不著 卻看得見.
因為有你 因為有我.
甘心給你用心愛.
把心中這一份愛 活出來.
耶穌的愛 激勵我.
放開我的生命.
讓自己成為別人祝福.
耶穌的愛 點燃我.
心中熊熊愛我.
我們一起向世界 活出愛.
我們一起.
我們一起向世界 活出愛.
我們一起.
去到世界不同的地方去活出來.
其實在很多的情況之下.
我們都有機會跟別人分享福音.
我讀書很少.
十歲才開始讀ABCD.
所以很多時候以前對著一些.

$^{201}$譬如他身分高一點.
或者是他教育程度高一點的人.
好像不知道應該說些什麼.
我經常覺得回到聖經裡面.
其實裡面有很多很多的智慧.
所以我自己很多時候都覺得.
能夠有時間.
我見到你很多的查經班.
我自己都真的要提醒自己.
因為經常都出了許傳福音.
還有其它的原因.
是少了自己的教會.
今天我選了三個故事.
希望你不要介意.
因為我不是去怎麼譯經 解經那些.
我只是覺得這三個故事.
是很能夠讓我分享今天這個主題.
我們人生裡面有時會挫敗.
有時越是在我們身邊的人.
我們最緊張,最愛的人.
原來他一個眼神.
有一句話.
或者他一個決定.
是打沉了我們.
我們好像回不了身.
還有怎麼樣可以去原諒,接納呢?.
很不簡單.
我選了這三個故事.
一個就是《約瑟的才醫》.
我想很多看音樂劇的都會看過.
《約瑟》是Jacob最小最小的一個兒子.
但是他爸爸很疼愛他.
有一點就是偏心.
所以帶來很多東西.
第二個就是蘇羅和大衛.
蘇羅是皇帝.
其實原來大衛是他的女婿.
但是這麼親這麼親的一個關係.
都會被人逼迫.
都會被人追殺.

$^{241}$第三個故事我就說耶穌釘在十字架.
祂那三年多全服飲.
有很多的神蹟奇事.
但是因為很多的妒忌.
各樣的東西.
最後是被人釘在十字架上.
但祂最後那句話我覺得很感動.
祂說「神啊 你原諒他們吧」.
因為他們現在所做的.
他們不知道.
我今天就和大家分享這幾件事.
這個主題其實就是.
那天和牧師師母他們.
喝茶,吃飯,聊天的時候.
就說其實我們人生裡面.
很多都是受到傷害的.
那個傷害有時可以很深.
有些時候可能你真的靠著神的醫治.
和安慰,寬恕.
我們很容易可以突破那種的困境.
好 那我開始了.
《有色的才藝》.
Joseph 剛才也說了他是最小的兒子.
他爸爸很疼他的.
就給了他一件很漂亮的衣服.
就是五顏六色.
我不知道五顏六色是不是一定很漂亮.
不過他就給了他一件很漂亮的衣服.
他那幾個哥哥就很妒忌.
為什麼爸爸那麼偏心.
只是喜歡他.
那些哥哥每天都要出去工作.
很辛苦.
這些文字可以.
如果你喜歡看這些可以看看.
因為都很小.
當那些哥哥出去工作的時候.
Joseph 就要送一些飯菜給他.
那些哥哥其實就已經很生氣.
有一次Joseph就傻傻的.

$^{281}$就跟他說.
其實你個個都差不多要敬拜他.
他做夢的時候就見到一些禾稈.
那些哥哥就更加不順氣.
人生是這樣的.
為什麼只是疼他.
為什麼不疼我.
就差到人生有時候真是難測.
那些哥哥就想將他扔到井裡.
就是殺了他.
其中有一個就有良心的.
就說不要了.
不要殺他.
不如賣了他.
就將他賣了他做奴隸.
就做那些工人就賣了他去埃及.
去到埃及他很幸運.
他就進了一個官員.
官員名詞無所謂.
就是護衛的官.
在他家裡他很勤奮很聰明.
就一直被他的老闆升職.
最後就做了一個管家.
但是原來人生裡面.
以為一波平了.
原來第二波又來.
不要緊.
靠著耶穌就成功了.
於是他是一個很順服和愛主的弟兄.
他在做事的過程裡.
原來這個大官.
就是他的老闆.
他的老婆就很差的.
就走去想非禮他.
還強姦他.
反過來.
宗善在走的時候.
就被太太拔了他的衣服.
因為他拔了他的衣服.
脫了他的衣服.

$^{321}$然後宗善就走了.
因為他知道神不會喜歡這件事.
很多時候有些引誘或者衝擊來到.
我們怎樣去站著走.
這個是很重要的.
於是就拔了他的衣服.
那個太太就告訴她老公.
惡人先告狀.
就說他想非禮我.
宗善想非禮我.
這個老闆當然生氣了.
就把他扔到監房裡.
神又真的很疼他.
他在監房裡.
原來宗善有一個特別的恩賜.
就是懂得解夢.
那裡有一個專門做酒的.
另外一個專門吃東西的.
這兩個人就和他一起坐牢.
他們兩個又做夢了.
不知道他們有這麼多夢法.
做夢的時候就問他.
祖師父你幫我解一下.
究竟這些夢是什麼意思.
他對著做酒的那個就說.
這個應該是三天後你就可以出監了.
另外一個專門做吃東西的.
就是那些善長.
就是吃東西的善長.
不是做慈善的善長.
就說你可能三天後會被人斬頭.
我不知道聽到他們會有什麼反應.
真的三天後這些事情全部發生.
原來他在埃及那裡.
法老王又做夢了.
他做夢之後.
所有的大神那些人.
全部和他無法解釋.
沒有辦法和他解一個夢.
這個做酒的那個.

$^{361}$因為放了監.
就回到皇帝那裡.
幫皇帝辦酒宴.
他說我好像記得有一個叫Joseph的.
他就好像可以解夢.
就去找Joseph.
Joseph去見到這個法老的時候就說.
你的夢是這樣的.
七年裡面.
是會全部很豐收.
風調雨順.
但是這七年裡面你就要好好地儲糧.
如果你不儲好糧的話.
七年之後就會有旱災.
連續七年的旱災.
如果你現在不儲好所謂要的糧食.
到時大家都很慘了.
這個法老真的聽他說.
七年過了之後就真的旱災.
幸好那段時間他聽Joseph說.
就儲起了很多糧.
於是乎他就升啊升啊.
就做了宰相.
你想一下一個小朋友.
被一些哥哥本來想殺.
他扔到井裡.
然後接著就去做奴隸.
奴隸裡面在監房裡面就遇到這樣的事.
然後接著還要是那些女人又這樣.
人生裡面有時候真的是你會.
一幫未平一幫.
你會一直這樣來.
但是我自己覺得就是.
他因為是一個很愛神很謙卑.
他就願意神賜給他智慧.
在最後他終於做了個宰相.
做了個宰相之後.
真是他自己的娘家.
就是他爸爸那些弟兄姊妹.
全部都差不多餓得死死的.

$^{401}$就要到處去找糧食.
就去到埃及那裡拿糧食.
第一次.
因為他是宰相.
他就已經給了很多糧食他們家裡.
好了 到第二次來的時候.
他就說你們有多少兄弟姊妹啊.
什麼什麼這樣說.
他說一個很小的弟弟.
但當然不是他 其實他是最小的.
但他對上還有一個.
於是乎他說.
那你帶那個最小的弟弟來吧.
當他帶那個最小的弟弟來.
第二次去拿糧食的時候.
我不知道為什麼.
就是Joseph會是.
將一些很名貴的東西.
放在他的糧食裡.
然後就說他們.
哦 你偷了東西走.
就全部抓了.
那到最後呢.
就是簡單一點說.
不說那麼長了.
那最後就是.
才和那些哥哥相認.
就是說其實我就是Joseph了.
但他沒有怪到他的哥哥.
在那裡這麼難受.
這些就是你看到這個杯子.
就是那些名貴的東西.
放在他的糧食裡.
然後就抓了他.
抓了他們之後.
接著去到最後.
完全是和氣收場的.
接著就爸爸兄弟姊妹.
全部都知道.
原來這個就是Joseph.

$^{441}$是當時他們曾經.
這樣去對待他的.
所以其實人生裡面很重要.
我覺得一件事就是Forgiveness.
就是如果我們是.
有一句電影對白我很喜歡的.
它說It's not that I can forget.
它說不是我可以不記得.
事情發生過的一定記得.
但是我們有時候可以做到什麼?.
原諒.
就正如很偉大的說.
神你原諒他們吧.
因為他所做的是不知道.
那我們只可以怎麼樣?.
不斷去為他祈禱.
我經常都說有時候遇到很mean的人.
不要瞄準.
我經常覺得.
就是不要整臭自己.
有一個秘訣告訴你.
就是如果你經常.
你想想如果將來我再老一點的時候.
我的樣子就.
多噁心啊你說.
但是如果你經常都笑.
隨便吧.
四兩撥千斤.
靠著耶穌的德性.
如果你經常都笑.
將來你年紀大一點.
你是笑文來的.
樣子是這樣的.
不是這樣的.
很可怕.
所以這是第一個故事.
我覺得提醒了我.
其實人生裡面.
如果要痊癒.
不是身體病這麼簡單.

$^{481}$很多人是心病的.
你知不知道.
心病就為需心藥醫.
那就是你真的要找到一個方法.
第二個故事我覺得是很特別的.
我最近不知道你們有沒有看.
看不看中文片的.
你們這裡.
我去了看.
《尋秦記》.
哪些人看了的舉手.
是啊.
去尋秦始皇.
很好笑的.
你有沒有聽過古天樂講英文.
很好笑的.
真的相信我.
我陪朋友看了幾次.
很好笑.
當中跟這個很相似.
就是怎樣呢.
原來那個叫韓紹龍.
就是那個古天樂.
他有個兒子.
兒子就是韓愈.
韓愈其實將來.
接管了秦始皇.
秦始皇就很生氣.
其實情況有點像這裡.
經常都想殺他們.
這家人.
但是韓愈他們一家人.
包括古天樂那些角色.
其實是他一直栽培.
一直栽培.
由這個秦始皇.
他叫他做師父的.
叫古天樂做師父的.
他叫做韓紹龍.
他一直栽培到他做到皇帝.

$^{521}$打完六國.
但是去到最後.
這個皇帝是想將他所謂的師父.
都要殺掉.
是一直追殺他們的.
追殺他們整家人.
你聽到之後你覺得.
嘩 真是人心何測.
怪不得經常都說.
我們要保守我們的心性.
要保守一切.
因為一切的果後從心而生.
我給你看看這個故事.
這些都是我在網上拿的圖片.
他真是很多次.
其實他們真是親戚.
還有一樣東西.
我覺得很特別的就是.
他說.
神所按的.
高的那個王.
其實.
素羅一直去追殺大衛的時候.
在一個山洞裡面的時候.
其實大衛是可以兩次.
去殺死素羅的.
但是他不殺的時候.
就是因為這句話.
他說這個王是神高的.
就是神是指定的.
是祝福他做皇帝的.
所以他明明有機會去脫難.
你明不明白.
殺了他之後他自己做皇帝.
但是他沒有.
我經常覺得有時候.
人的情操是很重要的.
還有我們作為基督徒.
我們的謙卑馴服.
要聽話.

$^{561}$是很不容易的.
有時候遇到一些事情.
你會覺得不是啊.
為什麼會這樣.
所以你看到後面那個就是素羅.
前面那個拿著一把刀素羅.
後面那個就是大衛.
其實是可以殺他的.
去到最後沒有做這件事.
是和氣收場的.
我覺得其實真是一個很簡單.
這裡有很多文字.
我不寫太多了.
他在那段時間.
其實大衛是很慘的.
一直被他追殺.
還要是他本身是親戚.
我自己覺得這段話很好.
他說我的上帝.
求你救我脫離仇敵.
其實有時候我們真的可以.
那個仇敵不一定是追殺我們.
他有很多的破壞.
很多的不公義的一些行為.
他說我要把.
請你把我安置在高處.
然後能夠脫離這些人的攻擊.
我們很多時候都是求父慈平安.
保護我們.
我經常到處去一些很不應該去的地方.
我還記得我在孟加拉.
有一次和選美會他們做副評.
說過了.
如果你要過去那邊看那個水災.
Oxford Bazaar那裡.
你一定要過一條橋.
不然你看不到.
我們真的祈完禱之後.
走遠過一條橋就倒下.
真的又是這麼巧.

$^{601}$是不是.
所以我自己覺得有時候.
我們真的好像埃塞俄比亞.
坐那些飛機.
轉兩個圈.
然後跌三個圈.
飛來飛去都飛不出那個海灘.
就是那個峽谷.
但是我們基督徒.
個個都是嘻嘻哈哈地笑.
但是有兩三個善長.
他們是未信主的.
下來的時候.
你知道嗎.
都是大家沒有seat belt的.
因為那個地方很窮.
坐那些飛機.
然後駕駛飛機的那個是傳道人.
要騎了一趟才上飛機的.
接著下來的時候.
看到他們真的手都握得快要瘀青了.
真的是害怕的.
但當我們知道.
原來我們將來會去哪裡.
基督徒知道的.
很清晰地.
如果他們真的信主.
接受了那個救恩.
我們是知道自己很清楚去哪裡的.
好了 第三個故事.
耶穌釘十字架.
我想大家都知道.
在他人生最後那三年多裡面.
他有很多的神蹟.
死人復活.
水變酒.
然後麻風.
很多很多很多.
這些全部都是對世人來說.
是一個恩典.

$^{641}$但是很可惜.
那些凡尼賽人看不過眼.
可能對於自己.
很多人是這樣的.
其實這個是一個自卑.
就是你沒有信心.
怕別人取代了你那個身份.
或者那個位置.
很多時候我們都會做一些事情.
回望人生其實很重要.
好像我現在這個年紀.
有一天你坐在搖椅上搖來搖去.
或者在輪椅上被人推的時候.
很多過往的事情會像flashback那樣.
所以做事我們有時候是.
對得起別人 對得起自己才可以.
這裡他講得很好.
耶穌到最後釘在十字架上面.
祂沒有埋怨任何人.
祂只是說 苦啊.
赦免他們.
因為他們所做的.
他們不知道.
如果你看回這三個故事.
很簡單的跟你說.
人生裡面其實就是這樣.
很多時候我們是.
心裡面覺得不平安.
或者有時候很有趣.
人到了某個位置爭權奪利.
為了自己的益處.
很多事情我們會做.
而是令到在身邊的人受傷害.
又破碎的.
所以我自己覺得.
希望我們每個人都不會這樣做.
耶穌基督釘在十字架上面給我.
我覺得是一個很大的愛.
一個犧牲的愛.
我們很多基督徒都是白白的可以得到的.

$^{681}$OK 我又想你們上來跟我唱歌.
這首歌我喜歡的.
今天整個主題我覺得.
只要你是原諒那個人.
你心裡面真的原諒他.
你就會得到那份平安.
這份平安不是世人可以給.
是天賦特別的給我們.
讓我們每一步走的時候很輕鬆.
好 我們唱這首歌《平安》.
廣東話的.
有灰暗 有傷困倦.
你讓我站穩 磐石之柱.
就算困楚艱難仍要著望.
我心藏著不死盼望.
耶穌是我牧者 永在我旁.
你將彼官都安慰我.
讓主你此刻平安.
不像世人所恥.
我的心不必再憂傷.
也不用膽怯.
主你此刻這平安.
是用重金換來的恩典.
求別相思 我得天賜.
受營費 若我得到赦免.
好 為了爭取時間.
我們唱一段好嗎.
但這首歌很好的 很多安慰.
我不知道你有沒有一個習慣.
我是背歌詞的.
我們待會慢慢再唱這首歌.
為什麼我會覺得背歌詞有兩樣東西.
第一 其實我當時的祈禱.
我都是用聖詩來祈禱.
你聽一下 它的邊霜令我們得到醫治.
它的刑罰可以令我們的罪得到赦免.
有時候我們不懂得祈禱的時候.
有一個音樂就很好.
我一個小秘密告訴你.
如果我讀書的時間.

$^{721}$所有的書本都是配音樂的.
我想我可能會考第一.
但很可惜不是.
你硬要用腦筋是不行的.
原來我和我女兒都是.
她考試是看著音樂劇讀書的.
即是那時候的音樂劇.
原來那種音樂是幫到我們的記憶的.
所以多些唱歌不怕五音不全有什麼所謂.
我最記得香港有一個牧師是利美學院的.
專門教敬拜的.
他知道我唱歌之後.
Mary不要緊 你就以你的拍子唱.
你知道是什麼意思.
幸好我還沒有奏音.
今天的主題 剛才聽完三個故事.
突破生命困境 人生傷害.
我們受傷了 怎樣可以得到醫治呢?.
我想和你說一點我自己的故事.
有些朋友可能是新朋友.
我是生於50年代,60年代的香港.
那時候很窮.
但是我在街上四年做街童之後.
被他們送了我進寶良谷.
在寶良谷裡面.
它不是一個有宗教信仰的機構.
那些姑娘,家長全部都是拜石頭.
拜白蘭樹等等.
但是原來難過的時間很快就過去了.
真的很快就過去了.
回望只是覺得感恩.
因為寶良谷是我娘家.
我經常都回去.
還有我最感恩的一件事就是.
那時候在寶良谷我可以出街.
又要拿出街紙.
我就在銅鑼灣所有不同的教會.
我都去過了.
那裡有靈良堂.
有很多不同的教會.

$^{761}$我經常都去那些教會坐.
它是給我星期天去.
我想是因為這樣.
所以我小時候就開始喜歡唱歌.
喜歡聽道理.
雖然有時不是很明白.
但是聖經說的.
它說神是會保護寄居的.
我們全部都是寄居的.
尤其是孤兒和寡婦.
神真的 我覺得祂很疼我.
OK 那我今天和你講一些.
我自己的故事.
每次我都是講媽媽的.
其實大約是八年前.
我自己講福堂是二千多人的教會.
總共有廣東話是四堂.
廣東話有四堂.
英文一堂 國語一堂.
其實我們總共有六堂不同的道.
就是崇拜的.
竟然在父親節.
我的主任牧師知道我沒有爸爸.
其實他可能都不知道.
知不知道我原來這麼憎恨他.
他說Mary 你經常出去分享.
你不如這個周活教會你來講.
我說 好啊.
我不知道是父親節.
是六月初.
我就答應了.
慢慢慢慢周遊列國.
周圍去完之後 回來了.
咦 父親節我去講?.
不是吧?.
我簡直是很少驚慌的.
我在台上從來不驚慌的.
我拿著咪高峰就說話的.
那些人來的.
星期五晚上睡不到.

$^{801}$直到星期六下午六點鐘.
不是 五點鐘開始.
一直講到七點鐘.
就有一堂要自己講.
就是講父親.
我可以講天父的愛的.
但是我不是.
我都說我不是牧師.
我怎樣都要講一些我自己.
和我爸爸的關係.
我有一個習慣是二十多年了.
在香港的太平山頂.
我經常上去行山的.
每天都走一個圈45分鐘.
就是要運動那些.
那天早上 星期六很早.
六點鐘我已經上了山頂.
插著耳朵一邊聽著聖詩一邊走.
走了五分鐘左右.
很清晰的.
耳朵有聲說.
「你爸爸呢?」.
我第一個反應就是.
「我天父的爸爸 你看我多棒」.
我自己和自己講話沒有出聲.
只聽到「你爸爸呢?」.
我就馬上說.
「我天父很疼我」.
他真的怎樣怎樣怎樣.
就是我自己心裡在想.
那就繼續走 你明不明白.
走多一會又來了.
我說「你的爸爸」.
「你的親生爸爸」.
我站在那裡.
靜呆了.
旁邊沒有人.
我不是傻的 你聽我說話就有紋路.
我說「糟了 什麼事呢」.
站在那裡.

$^{841}$因為真的自己上教會就要開始分享.
接著呢.
一點都控制不了.
好像開了水喉一樣.
我的眼淚不斷流.
全身好像本來的傷口.
全部劈開了.
被人全部撕開了.
那時我走整條路是要45分鐘.
我走了五分鐘.
一路哭哭哭哭哭.
哭了八個字.
還沒哭得完.
一路不停地抹眼淚.
我不知道自己在搞什麼.
我有個親生爸爸.
我從來都沒見過他.
又不知道他什麼.
你明不明白.
都不算是好人.
這樣對我.
我只是很多的憎恨.
很多的不開心.
很多的傷害.
在裡面一次過地哭.
好了 抹乾眼淚.
拿回車子.
回到公司我要工作.
就跟同事說.
「打個這樣的東西出來」.
都不知道說什麼.
「先弄出來」這樣.
我就開始有工作了.
到四點多五點鐘.
回到自己的教會.
唉 就上台說.
通常都是說我媽.
跟我說「好啊」.
「真的很感恩」.
「我跟她有十年時間」.

$^{881}$接著打了這個出來之後.
我已經說不出話了.
在台上就哭了.
我想一分多兩分鐘.
就覺得不可以全世界都看著你哭.
我現在其實隨時都哭得出來.
接著就說了.
一開聲第一句.
「我跟神認罪」.
「哇 神啊」.
「我全世界去傳福音」.
「為什麼我不可以原諒我這個爸爸呢?」.
可能那時候他還沒有信主.
我又經常說「寬恕」.
剛才說完「Forgiveness」.
說「愛」.
無盡的愛.
我是永不自釋.
我說「對不起啊」.
「我的天父」.
「我可以原諒他了」.
「我真的可以原諒我爸爸了」.
接著又在哭.
那天我想差不多所有的男士.
都重新立志.
都決志.
其實人生原來是這樣的.
我和我天父的關係.
因為我爸爸.
我沒有辦法感受到.
我真的完全是.
很理性的覺得.
「我知道有個天父就行了」.
但是自從那次之後.
我的靈明不再一樣.
我和你分享的一些東西.
是真的這麼奇怪的.
從來都沒有在外面分享過的.
我親愛的爸爸.
我現在終於可以寫了.

$^{921}$我親愛的父親.
這個是我爸爸.
我為什麼會找得到的呢?.
真的原來有些人是很有趣的.
原來他們知道我爸爸的名字叫張英超.
他是一個很出名的畫家.
這些是他的畫.
拍賣的.
然後呢.
他就做了一個.
我爺爺.
是不是叫公公不認識.
總之是我爸爸的爸爸.
的Thesis.
找了所有的東西出來.
於是.
原來我爸爸除了畫畫之外.
他做生意很有錢.
就是1952年在香港生了我之後.
就去了台灣住.
這張是我13歲畫的畫.
用那些普通的油彩.
就是那些上路作堂畫的畫.
這些是我畫的油畫.
很多很多.
我從小到大都畫畫.
原來是不是真的有遺傳的呢?.
沒有人教過我的.
那時候我還沒有跟老師.
那時候已經不斷畫很多很多畫.
我是一個畫家來的.
其實嚴格來說.
我畫得挺好的.
然後呢.
我爺爺呢.
就叫張珍侯.
我覺得這個是粵語片來的.
慢慢聽下去會笑.
他在上海非常有錢.
很出名的一個名人.

$^{961}$他的兒子就叫張英超.
就是我爸爸.
然後呢.
他是很喜歡攝影的.
所以他那時候在上海.
就開了很多這些.
那時候是哥達菲林的.
不是我現在的手機.
開了很多很多這些東西.
碰巧我也是攝影師來的.
這些全部是我的攝影來的.
我的東西還在法國賣著.
我是全世界到處展覽我的攝影的.
這些全部是我的攝影來的.
是不是真的這麼奇怪.
我爸爸的爸爸是喜歡攝影.
他攝影很厲害的.
我爸爸是畫家.
我又畫畫.
這些全部是我小時候做的東西.
到最近還好笑.
因為我看到知道了所有這些東西.
我有出世紙的.
因為我那時候要去日本.
就是同軍要去日本.
保倫他就查了我的紀錄.
原來我爸爸真是叫做張英超.
我媽是什麼名字.
我真名叫做什麼名字.
齊齊 我真名真是叫張瑪麗的.
於是我就開始找了.
這個人去哪裡了.
這個張英超.
看他還在不在世了.
我1975年選香港小姐.
原來他76年就死了.
我後來才知道.
我去到台灣.
然後他最後就差不多.
我有個姐姐.

$^{1001}$但都是回了天家了.
走了.
50年 52年.
這些是很私人的東西.
但我覺得現在我可以說.
這麼多年來.
我一直都在拆東西.
我放一些圖片在一起.
這間屋是我祖屋來的.
我爸爸的爸爸的屋來的.
是已經是.
如果你看這些文字.
是一個.
即是設計一種紀念品的.
不可以拆 不可以整的.
沒有人在裡面住.
上海很流行很多人.
是會進了你的屋.
5,6家人霸佔了一間大宅.
那些你是沒辦法趕他走的.
但我現在這間屋.
由頭到尾都是政府接管的.
是甚麼情況.
我不知道.
很有趣.
我就和熟悉的朋友說.
我想去看看這間屋.
竟然這樣出來了.
我說好.
我打電話去.
他們就說.
請問你有甚麼要求.
我沒有甚麼要求.
我想看看這間屋.
我就去看這間屋了.
你看這間屋多漂亮.
全部都是.
即是如果你搜尋上去.
張真侯的屋企就有的了.
他們帶我去看.

$^{1041}$幾個官員帶我去看的.
是不是漂亮得很厲害.
三層.
有十幾間房.
天花板全部畫了那些.
叫甚麼飛的那些.
那些殼那些.
然後我一路走一路走.
那時我還是受傷的.
上年.
即是2014年的12月.
一路這樣走.
現在你看到.
最高的那個.
呈了出來的那個就是我爸爸.
原來他也是漫畫會的主席.
真是.
我媽媽興本佳人.
我爸爸是才子.
有錢.
即是不只是那個才.
不是背著那個才.
所有的事情.
怎麼會搞到女兒.
進了孤兒院住足11年.
這個世界其實原來有很多事情.
那種傷痛是在我心裡.
我經常覺得.
是不是我很壞.
是不是我有些事情做得不好.
所以我爸就不要我媽.
即是小時候會想很多很多這些.
好像是自己的錯.
這些全部是他畫的東西.
我全部都查出來的.
那最有趣的就是.
我去洛杉磯.
如果你聽過我的見證.
我去洛杉磯有個阿姨的.
就是我在保良局的時候.

$^{1081}$在孤兒院的時候.
她很疼我的.
那我去探望她的時候.
有一次我剛好.
因為我找到我爸爸的照片.
那我就給她看.
問她知不知道這個是誰.
她92歲的了.
講國語的.
那我剛好有些主流的姐妹在那裡.
她就說.
啊 即是那個Sydney Barber.
我誒.
那我就馬上叫我姐妹幫我拍.
拍下來.
我問你怎麼知道這是Sydney Barber.
你怎麼知道我是她爸爸.
她誒.
我小時候14,15歲的時候.
即是她13,14歲了.
我媽媽是帶著她去拍拖的.
然後我爸爸是送了些畫給她的.
那段影片我拍了.
我拍下來.
在香港有些畫家都認識我爸爸.
那你想一下.
看回.
其實我真是應該是一個很有教養的.
在家裡學一下彈琴.
跳一下芭蕾舞那些.
但我為甚麼會在街上睡四年.
然後被人送去保養局.
原來人生裡面有很多東西.
我們真是沒得選擇.
很多東西來到我們那裡的時候.
你怎麼去做.
但每一次你的選擇.
你的決定.
其實是影響你的人生.
所以有時候我們自己要知道一樣東西.

$^{1121}$就是.
譬如我和我媽媽.
我以前真是很憎恨她.
你看我媽媽多漂亮.
興本佳人.
怎麼會生了些小朋友.
全部丟回香港.
然後自己一個人走了.
幾十年.
原來人在罪裡面.
是會越走越黑.
我爸爸也是.
你明不明白.
到他走了的時候.
76年走的時候.
身邊一個人都沒有.
原來人生裡面是真是有很多的挫折.
我和我媽媽在那十年裡面.
天父祝福我們.
有很好的時間.
我看著她.
一個很慈祥的老人家.
我現在找回我爸爸.
又找回我媽媽.
我想告訴你.
你是不是覺得這個世界.
這麼多東西都是這麼巧合.
我的朋友和我說.
快點去上海拿回這間屋.
我自己在笑.
我北京也買了一間屋很久了.
也是很漂亮的.
我自己覺得神.
我們經常說神是一個供應的神.
不是我們需要去求祂.
說你給些什麼我給.
不是那樣的.
其實剛才那首《平安》.
你聽的時候就知道.
祂真的給了你心裡面那種的富足.

$^{1161}$你知道每年香港小姐就選的.
一選就說美麗與智慧.
我心想美麗算了吧.
智慧很重要.
認識耶和華就是智慧的開端.
所以今天我想和你分享的就是.
今天的張瑪麗.
如果沒有神不是今天的我.
我心裡面只是有很多很多的感恩.
就是我感恩我現在所有擁有的東西.
我有一個兒子,一個女兒,有一個孫子.
每一件事我覺得無論怎麼樣都好.
可能你都在經歷著一些事.
可能在教會你都在經歷著一些事.
我覺得不需要怕 只要相信.
相信神是公義的.
祂會帶我們走.
我們心裡面有愛.
信望愛 我們以愛為首.
在愛裡面.
如果你看《哥林多前書》13章.
我是恆久忍耐又有恩慈.
一切你都會看到那個結果是怎麼樣.
不要計算人間的惡.
還有一樣東西就是很重要.
我們真的需要時時去敲神的門.
因為你敲神的門.
神的門一定開.
我們求 祂就會給我們.
我們找 會找到.
所以我自己覺得.
這世人裡面我自己真是.
你問我好像很壞.
由一無所有到我現在一無所求.
我覺得一切都是.
我只是現在不斷祈禱.
就是我的兒子,我的女兒.
我希望他們是一直跟隨神.
還有我的孫子.
我希望他一直在教會長大.

$^{1201}$什麼都不重要.
讀什麼名校之類的.
你知道很多媽媽是怎麼樣的.
我自己覺得.
我相信天父會祝福她好像祂祝福我一樣.
在這裡我就想問一下.
有沒有新朋友舉手給我看看.
你從來都沒有來教會的 有沒有?.
有 嗨 祝福你啊.
真是好了 還有沒有?.
新朋友.
我是每次都會邀請的.
因為我覺得我的經歷是我的經歷來的.
但是你今天來到教會.
我覺得你不是意外來的.
我知道有些新朋友可能你們害羞.
不過我在這裡希望你跟我做這個祈禱.
然後我們就跟你唱一首歌 好嗎?.
讓我們在聖利安修這裡.
Crosspoint Church.
大家可以遇上.
真是無窮的恩典和愛.
我相信很多朋友在這裡.
他們都是有一些經歷.
可能也有些傷害.
或者有很多不明白為什麼會發生的事.
但是主讓你愛我們.
你更加愛這個家.
你會保守這個家.
你保守我們每一個在這裡的弟兄姊妹.
是可以靠著你去得勝.
靠著你我們可以走一條.
我們覺得很有平安.
很感恩的路.
在這裡.
如果有朋友你是第一次來教會.
或者你還沒有接受耶穌基督.
我希望你在這裡.
輕輕跟自己說你願意.
或者你喜歡的人可以舉手.

$^{1241}$接受耶穌基督做你生命的救主.
我相信你會看到.
很多很多的奇妙恩典在你身上.
在你的家庭裡面發生.
甚至有些時候可能你根本上.
面對不到一些情況.
你的過去.
人可以不原諒你.
但天父知道.
祂會帶領你走前面的路.
願主賜福給我們每一位.
在新的一年馬年裡面.
平安健康.
主恩常在.
祈禱是奉著耶穌基督得勝的命求.
Amen.
\newpage



\section{約書亞記 6:1-27}
\label{sec:031688}
\textbf{同心打美好的仗}~~~~張智聰博士
\newline
\newline
連結: \href{http://www.church.com.hk/acms/content.asp?site=cdc\&op=show\&type=product\&code=031688\&layout=sermon}{\texttt{031688}} ~~~~ 語音日期: 2026-02-22
\newline
\newline
\hyperref[sec:031687]{\small{< < < PREV SERMON < < <}}
~
\hyperref[sec:index_chronic]{\small{[返順時目]}}
~
\hyperref[sec:031689]{\small{> > > NEXT SERMON > > >}}
\newline
\newline
約書亞記 6:1-27
\newline
\begin{longtable}{cl}
\hline
\hline
章節 & 經文 (和合本修訂版)\\
\hline
6:1 & \begin{tabularx}{0.7\textwidth}{X} 耶利哥的城門因以色列人的緣故，關得嚴緊，無人出入。 \end{tabularx} \\ \\ \relax
6:2 & \begin{tabularx}{0.7\textwidth}{X} 耶和華對約書亞說：「看，我已經把耶利哥城和耶利哥王，以及大能的勇士，都交在你手中。 \end{tabularx} \\ \\ \relax
6:3 & \begin{tabularx}{0.7\textwidth}{X} 你們要圍繞這城，所有的士兵繞城一次，六日你都要這樣做。 \end{tabularx} \\ \\ \relax
6:4 & \begin{tabularx}{0.7\textwidth}{X} 七個祭司要拿七個羊角走在約櫃前。到了第七日，你們要圍繞這城七次，祭司也要吹角。 \end{tabularx} \\ \\ \relax
6:5 & \begin{tabularx}{0.7\textwidth}{X} 羊角聲拖長的時候，你們一聽見角聲，眾百姓要大聲呼喊，城牆就必倒塌，各人要往前直上。」 \end{tabularx} \\ \\ \relax
6:6 & \begin{tabularx}{0.7\textwidth}{X} 嫩的兒子約書亞召了祭司來，對他們說：「你們抬起約櫃來，要有七個祭司拿七個羊角在耶和華的約櫃前。」 \end{tabularx} \\ \\ \relax
6:7 & \begin{tabularx}{0.7\textwidth}{X} 他又對百姓說：「你們向前去圍繞那城，帶兵器的要在耶和華的約櫃前過去。」 \end{tabularx} \\ \\ \relax
6:8 & \begin{tabularx}{0.7\textwidth}{X} 按照約書亞對百姓所說的，七個祭司拿了七個羊角在耶和華面前過去，他們吹著角，耶和華的約櫃在他們後面跟著。 \end{tabularx} \\ \\ \relax
6:9 & \begin{tabularx}{0.7\textwidth}{X} 帶兵器的走在吹角的祭司前面，後隊跟著約櫃走，號角繼續在吹。 \end{tabularx} \\ \\ \relax
6:10 & \begin{tabularx}{0.7\textwidth}{X} 約書亞吩咐百姓說：「你們不可呼喊，不可讓人聽見你們的聲音，連一句話也不可出你們的口，直到我對你們說『呼喊』的那日，你們才呼喊。」 \end{tabularx} \\ \\ \relax
6:11 & \begin{tabularx}{0.7\textwidth}{X} 這樣，約書亞使耶和華的約櫃圍繞那城，把城繞了一次。然後，眾人回到營裡，就在營裡住宿。 \end{tabularx} \\ \\ \relax
6:12 & \begin{tabularx}{0.7\textwidth}{X} 約書亞清早起來，祭司又抬起耶和華的約櫃。 \end{tabularx} \\ \\ \relax
6:13 & \begin{tabularx}{0.7\textwidth}{X} 七個祭司拿七個羊角，走在耶和華的約櫃前，他們吹著角；帶兵器的走在他們前面，後隊跟著耶和華的約櫃走，號角繼續在吹。 \end{tabularx} \\ \\ \relax
6:14 & \begin{tabularx}{0.7\textwidth}{X} 第二日，他們再把城圍繞一次，就回營裡去。六日都是這樣做。 \end{tabularx} \\ \\ \relax
6:15 & \begin{tabularx}{0.7\textwidth}{X} 第七日清早黎明時，他們起來，以同樣的方式圍繞城七次；惟獨這一日他們圍繞城七次。 \end{tabularx} \\ \\ \relax
6:16 & \begin{tabularx}{0.7\textwidth}{X} 到了第七次，祭司吹角的時候，約書亞對百姓說：「呼喊吧，因為耶和華已經把城交給你們了！ \end{tabularx} \\ \\ \relax
6:17 & \begin{tabularx}{0.7\textwidth}{X} 這城和其中所有的都要永獻給耶和華作當毀滅的，只有妓女喇合與她家中所有的可以存活，因為她隱藏了我們所派的使者。 \end{tabularx} \\ \\ \relax
6:18 & \begin{tabularx}{0.7\textwidth}{X} 但你們務必謹慎，不可取那當滅的物，免得你們受詛咒，取了那當滅的物，使以色列全營成為詛咒而遭受災禍。 \end{tabularx} \\ \\ \relax
6:19 & \begin{tabularx}{0.7\textwidth}{X} 只有金子、銀子和銅鐵的器皿都要歸耶和華為聖，放入耶和華的庫房中。」 \end{tabularx} \\ \\ \relax
6:20 & \begin{tabularx}{0.7\textwidth}{X} 於是百姓呼喊，祭司吹角。百姓一聽見角聲就大聲呼喊，城牆隨著倒塌。百姓上去進城，各人往前直上，把城奪取。 \end{tabularx} \\ \\ \relax
6:21 & \begin{tabularx}{0.7\textwidth}{X} 他們把城中所有的，無論男女老少，牛羊和驢，都用刀殺盡。 \end{tabularx} \\ \\ \relax
6:22 & \begin{tabularx}{0.7\textwidth}{X} 約書亞對窺探這地的兩個人說：「你們進那妓女的家，照你們向她所起的誓，將那女人和她所有的都從那裡帶出來。」 \end{tabularx} \\ \\ \relax
6:23 & \begin{tabularx}{0.7\textwidth}{X} 兩個作過探子的青年進去，把喇合與她的父母、兄弟，和她所有的帶出來，他們把她所有的親屬都帶出來，安置在以色列的營外。 \end{tabularx} \\ \\ \relax
6:24 & \begin{tabularx}{0.7\textwidth}{X} 他們用火焚燒了那城和其中所有的，只有金子、銀子和銅鐵的器皿都放在耶和華殿的庫房中。 \end{tabularx} \\ \\ \relax
6:25 & \begin{tabularx}{0.7\textwidth}{X} 至於妓女喇合和她父家，以及她所有的，約書亞保存了他們的性命。她就住在以色列中，直到今日，因為她隱藏了約書亞派來窺探耶利哥的使者。 \end{tabularx} \\ \\ \relax
6:26 & \begin{tabularx}{0.7\textwidth}{X} 當時，約書亞叫眾人起誓說：「凡興起重修這耶利哥城的，當在耶和華面前受詛咒。他立根基的時候，必喪長子，安城門的時候，必喪幼子。」 \end{tabularx} \\ \\ \relax
6:27 & \begin{tabularx}{0.7\textwidth}{X} 耶和華與約書亞同在，約書亞的名聲傳遍全地。 \end{tabularx} \\ \\
[1ex]
\hline
\hline
\end{longtable}
$^{1}$今天我們會讀到約書亞記六章一到二十七節.
經文是大家很熟悉的耶利哥城的攻破.
不知道大家當初聽攻打耶利哥這件事的時候.
會不會都覺得這是一件很奇怪的經歷.
當跟小朋友說這段故事的時候.
就會更加覺得很奇妙 很神奇.
六日繞城一次 七日繞七次.
之後大聲呼喊 城牆倒塌.
耶利哥城一下子被攻破.
說出來也不會有人相信有這樣的攻城方式.
但是一群以色列人竟然跟著去做.
並且最終得到勝利.
耶利哥城是一個不容易攻破的地方.
耶利哥是約書亞帶領進入迦南的第一場戰役.
而且耶利哥處於一個戰略的重鎮.
要攻佔迦南起了一個非常重要關鍵性的影響.
因此約書亞之前在第二章.
已經一早派人去窺探耶利哥城.
要了解一下城內的虛實.
不過很快探子就被人發現了行蹤.
因著得到基里拉俠的幫助.
他們兩個探子才能夠最終成功逃脫.
所以我們有理由相信.
耶利哥城是一個不容易攻破.
守衛又非常森嚴的地方.
當以色列人冰冷城下的時候.
他們就更加嚴陣以待.
就好像經文六章一節裡說.
耶利哥城門因以色列人就關得嚴緊.
無人出入.
這一場的確是一場不容易打贏的戰役.
在攻打耶利哥的描述裡.
聖經裡用了三個部分.
分別是征服之前對約書亞的吩咐.
征服的經過.
征服對約書亞的影響.
特別在五章十三到十五節.
上星期劉亦風牧師也分享過.
其實可以視為.
約書亞攻打耶利哥城的前奏或呼召.

$^{41}$我們不如一起再讀一次.
約書亞記五章十三到十五節.
預備一二三.
你是幫助我們呢?是幫助我們敵人呢?.
他回答說:不是的.
我來是要作耶華軍隊的元帥.
約書亞就俯伏在地下拜.
說:我主有什麼話吩咐僕人?.
耶華軍隊的元帥對約書亞說.
把你腳上的鞋脫下來.
因為你所站的地方是聖地.
約書亞就照著行了.
這一次是約書亞一次很特別的經歷.
他和耶華軍隊的元帥相遇.
約書亞見他手裡有拔出來的刀.
就對他發出提問.
你是幫助我們呢?是幫助我們敵人呢?.
有時當看不清楚狀況.
的確會容易誤判形勢.
因此約書亞問得很直接了當.
而耶華軍隊的元帥也答得很直接.
我來是要作耶華軍隊的元帥.
在這裡讓我們看到主和客.
很戲劇性的互換.
突然間對調.
最後更加令到約書亞毫無懷疑地.
向這位元帥俯伏下拜.
並且樂意聽候元帥的吩咐差遣.
是非常難能可貴的事情.
試想一下一個人突然出現在你面前.
然後你要將你的主權交給他.
由他帶領你.
這是一件非常不容易.
非常難得的事情.
試問有誰可以這麼快迅速地.
能夠轉換心態.
由我作主變成由他作主呢?.
並且甘心做僕人.
接受吩咐差遣.
相比於摩西的呼召的時候.

$^{81}$他諸多的推搪.
瓦割.
他遇見神的時候的頑強搏擊.
約書亞的順從更加顯得寶貴.
之後這個元帥對約書亞.
他說要把腳上的鞋脫下.
因為所站的地方是聖的.
不禁又想起摩西被呼召的情景.
在《出海及記》三章五節裡說.
不要近前來.
當把你腳上的鞋脫下來.
因為你所站之地是聖的.
在同樣的情景之下.
約書亞卻是不問緣由.
照著元帥所說的照著行.
一個完全信靠馴服的表現.
這次的經歷對約書亞來說是非常重要.
首先約書亞認定了神是他最可靠的支援.
是耶華軍隊的元帥帶領以色列人增戰.
不是靠約書亞本身的領導.
雖然約書亞是以色列人的領袖.
但他願意接受一位忽然出現的人物.
成為他更高的權威.
視他為從神而來的代表.
是很不容易的.
真的突然有人進入你的家.
說他就是那個家的主人.
你都會覺得很難接受.
但約書亞就欣然地接受.
欣然地去跟從.
其次約書亞也奉命要脫下鞋.
在這個時刻面對站在聖地上的聖爺的主.
被吩咐要分別為聖.
約書亞也順服地去遵從.
最後好像過去那麼多神所重用的僕人一樣.
他們都是先被分別為聖.
才再賦予任務.
我們的生命要先被塑造洗練.
才被神去使用.
賦予使命.

$^{121}$這個是恆久不變.
是我們要時刻牢牢謹記的真理.
可惜我們太多時候只顧著做一樣做一樣.
忙一樣忙一樣.
很多時候我們都只會重視工作的成效.
而忽略了生命的狀況.
就好像我記得在聖經裡面.
馬大瑪利亞的事件裡面.
耶穌去對馬大的一個提醒.
祂說 馬大 馬大 你為許多事思慮煩擾.
但是不可少的照一見.
瑪利亞已經選擇了上好的福分.
是不能奪去的.
所以我們必須要持續與神有關係.
持續與神連線 謙卑馴服.
不斷過土神喜悅的生活.
我們才能被神大大的使用.
這是我們要牢牢謹記的真理.
亦都是要我們常常被提醒.
常常被祂提醒的事情.
聖經接著記載到耶華對約書亞的吩咐.
就是要祂六日繞城一次.
第七日繞城七次.
後大聲呼喊 城牆就倒塌.
耶利哥就被攻破.
我相信大家都很熟悉的這個故事.
聖經裡說 耶華曉遇約書亞說.
看來我已經把耶利哥和耶利哥的皇兵大人的勇士都交在你手裡.
你們的一切兵兵要繞這城.
一日繞一次 六日都要這樣行.
七個祭司要拿七個楊角走在約櫃前.
到第七日你們要繞城七次.
祭司也要吹角.
他們吹的角聲拖長.
你們聽見角聲 眾百姓要大聲呼喊.
城牆就必塌陷.
各人都要往前直上.
這個是耶華神對約書亞的吩咐.
之後在第六節的經文.
其實是再重複三到五節的吩咐.

$^{161}$神命令的每一句說話.
凡是神向約書亞所說的話.
約書亞都說給祭司和兵兵知道.
其中不斷重複的騙語.
吹角成為了整個活動的重點.
一切的活動都是以吹角來開始和結束.
吹角代表著招聚一班以色列人.
齊心協力地去進行真戰.
這樣的響聲.
以及每一日繞行耶利哥城的隊伍.
一種好像巡遊典禮式的行進方式.
在音效和視覺上都極為壯觀.
試問當時那班耶利哥城的敵人.
看到這班以色列人圍著吹角.
他們的感想是怎樣呢?.
一定是覺得這班以色列人.
做一些很奇怪很出奇的事情.
很詫異他們的行為.
為什麼不是攻城而是圍圈呢?.
七日的時間相當於無敲節的七日.
以色列人在這個節期裡面.
七日要吃無敲餅.
象徵著他們傳言奉獻給神.
以色列人在這個節期裡面離開埃及.
並且親眼看見埃及的全軍覆沒.
第二次詳細技術預節的記載.
在《民數記》第九章.
預節之後接著就是.
從西乃山進入到漢夜去行進.
上星期牧師亦有提過.
聖經裡約書亞記第五章的預節.
其實和第六章很接近.
繞行耶利哥的技術.
和預節的記載很相連.
讓人想起第一個預節裡.
神擊敗以色列的仇敵.
使他們分別為聖歸給神.
這場戰爭亦成為了預節慶典的一部分.
不單止紀念第一次出埃及.
亦紀念在應許之地征服仇敵.

$^{201}$七日征服耶利哥的意義.
不單止是回憶起以色列人離開埃及的傳統.
更加提醒以色列人分別為聖歸給神的重要意義.
他們六日繞城一次 第七日繞城七次.
之後呼喊「城牆倒塌」.
耶利哥城就被攻陷.
一起完成這個壯舉.
能夠同心合力完成一件事是非常美好.
在教會服侍的上面.
我們都會有很多這樣的合作機會.
不過往往可能只是分工.
然後各自做自己負責的部分.
在這裡很明顯.
以色列人有一個很清楚的目標要一起完成.
大家只要跟足約書亞的指示.
就絕對沒有問題.
其實當中有沒有人曾經懷疑過.
約書亞的領導能力呢?.
他的領袖是否真的出於神呢?.
這種攻城的方法是否真的可行呢?.
約書亞如果身處於現在的社會.
可能這些問題都會不斷被人重複地問.
你是不是真的領袖對的?.
或者你是不是這個方法可以的?.
可能不斷地要去回答問題.
用約書亞的方法.
在今天可能會被人詬病為盲目跟從.
或者不夠客觀.
我相信現在的社會.
很難再用這種一言堂的方式來領導.
那種獨斷專行.
一個人說了算的方法.
將會被人批評為作風霸道.
不民主.
現代的領袖需要收集不同持份者的意見.
經過多番分析評估才做決定.
今天在教會要做決定.
也不容易用這種約書亞領導的方式.
太多人會質疑.
約書亞的領袖是不是他一廂情願的?.

$^{241}$或者有沒有事實根據?.
就算是多麼有份量的領袖.
幫他提出一個天馬行空的攻城方法.
我相信都沒有人願意去跟從.
在現今著重知識理解.
人人都可以讀經,研經,自行領受神話語的年代.
有多少人願意跟隨約書亞呢?.
現代人很喜歡用teamwork這個字.
teamwork這個字中文意思是.
團隊合作或者團隊的協作.
可以簡單說合作協作.
我試過問AI.
究竟合作和協作有什麼分別?.
大家分不分到這兩樣有什麼分別?.
AI就這樣說.
合作和協作是可以互換使用.
意思都是一起做事,一起完成某件事.
但比較細緻的區分.
合作是側重於互相配合的關係.
可以比較鬆散,分工清楚.
你做你的,我做我的.
方向一致就可以了.
協作有協同,協力的意思.
語感中強調一起動腦筋,一起做.
比較似大家坐在一起討論.
互相補位,一起完成一件事.
AI就很詳細地告訴我.
高效率的團隊合作的關鍵要素是什麼呢?.
就是要有清晰的溝通.
團隊成員公開地分享訊息.
確保每個人都了解目標,角色和進展.
要有共同的目標.
團隊就共同努力的方向達成一致.
並協調一致地進行工作.
信任和責任.
團隊成員彼此信任,盡職盡責.
並對自己的工作負責.
協作和互助.
團隊成員要互相幫助.
提供建設性的回饋.

$^{281}$並以尊重的方式處理衝突.
看完AI的描述.
不知道你心裡有什麼感受呢?.
你想想自己合作的團隊.
自己工作的團隊.
是否有過這樣的狀況呢?.
我看完之後的第一個反應就是.
有沒有可能發生呢?.
談何容易呢?.
要做到團隊方向一致.
彼此信任,盡職盡責.
互相幫助,有建設性的回饋.
彼此尊重.
真的不容易.
彼此信任.
互相盡力.
幫助,保衛.
是很美麗的圖畫.
但在我們現實狀況的上高.
上高有沒有真的有這樣的團隊呢?.
求主憐問我們.
若是都有一顆謙卑受教的心.
多聆聽上主.
亦都願意多聆聽他人.
謙卑馴服地去.
努力建立團隊.
回想起我過去曾經經歷過.
很多人和人合作的經驗.
最深刻的都是.
一起處理安排營會.
處理不同的大大小小的問題.
其中那種默契,保衛.
實在是難以形容.
能夠彼此信任,方向一致.
是很難得的團隊經歷.
是令到你很有力量去面對不同的挑戰.
就好像踢足球一樣.
清楚自己的崗位.
一個眼神,一個動作.
就可以無縫地連結溝通.

$^{321}$建立起默契.
互相配合彼此的信任.
向共同目標進發.
即是所謂的夢幻組合.
不知你們有沒有看球賽的時候.
見到那些夢幻組合.
那些很團隊的合作.
特別我有時看電視的時候.
見到那些穿插.
那些one two.
然後就進入了人家.
完全空門的狀況.
即是很美好的一個合作的關係.
很理想的狀況.
但實際的情況很多時候是.
我們各自只顧著自己.
表現自己.
怨隊友沒有配合好自己.
又或者怨隊友沒有顧及後方的狀況.
不斷地彼此埋怨.
製造紛爭.
互相懷疑.
不信賴我們的隊友.
其實.
這些人不單在球場上都是這樣.
其實在其他的狀況裡面.
在其他的團隊裡面.
這些情景都一而再地出現.
在教會裡面有很多的合作機會.
但亦都很多時候見到.
很多弟兄姊妹.
很多的爭執.
很多的猜忌.
很多的不合作的一些狀況出現.
求主連問我們.
叫我們能夠好好地建立團隊.
去打這場不容易的仗.
在第17到25節裡面.
一個很特別的描述.
講到毀滅和拯救的訊息.

$^{361}$重複地穿插排列.
梅花間竹地不斷地出現.
經文很長.
我讀給大家聽.
在當滅之物中拿取甚麼.
洗以色列仍成為當滅之物.
惹禍給他.
至於所有的銀子和金子.
以及銅鐵的器皿.
是歸上主為聖的.
他們要進入上主的寶庫裡去.
於是百姓呼喊.
祭司吹國.
百姓一聽到國聲就呼喊.
大大高呼.
城牆就倒下來了.
百姓都衝上城去.
人人爭相往前.
他們便攻取了那城.
他們便把城中的一切徹底滅絕歸神.
從男到女.
從少到老.
甚至牛羊奴都滅在刀下.
對偵查那地的兩個人.
約瑟說.
要進去那作妓女的女子家裡.
從那裡帶出那女子和那一切屬於她的.
就如你們向她所發的誓言一樣.
於是這些年輕的探子進去.
帶出那俠和他的父親.
他的母親.
他的兄弟.
和那一切屬於他的.
他們帶出他所有的親族.
安置他們在以色列營外.
至於那城.
他們用火焚燒.
也燒掉其中的一切.
只是銀紙和金紙.
以及銅鐵的器皿.

$^{401}$他們都給了上主殿中的寶庫.
而那俠那妓女和他父家.
並那一切屬於他的.
約瑟亞斯他們都存活下來.
他住在以色列中間.
直到今天.
因為他藏起了約書亞.
所差派去偵查耶利哥的使者.
經文很長.
不過很重複出現的.
是拯救和毀滅的訊息.
特別是多滅之物這個字.
在第十七節裡出現了很多次.
就是和毀滅原文是同一個字.
和第二十一節.
將城中一切徹底滅絕歸神.
的字根是相同.
在十七到二十五節裡.
和前面最大的不同的地方.
就是約書亞引介了滅絕令這個做法.
滅絕令的意思就是將物件.
完全獻給耶華傳言的毀滅.
曾經出現在摩西的申命記.
的二十章十六至十八節裡.
這樣描述.
但這些國民的城.
耶華利神既賜你為業.
其中凡有血凡有氣息的一個不可存留.
只要照耶華利神所吩咐的.
將這客人阿摩利人迦南人比利西人利未人.
希美人耶布斯人都滅絕正盡.
免得他們教導你們學習一切可真乎的事.
就是他們向自己神所行的.
以致你們得罪耶華你們的神.
摩西教導這班人進入迦南地的時候.
要將當地的迦南人滅絕.
在這裡你看到因為怕他們影響了以色列人.
以致他們得罪耶華神.
這個命令是神將神所賜給以色列人的地.
所有人口的聚居地都全都毀壞.

$^{441}$理由是當地的居民有機會影響以色列人.
讓他們行可真的行為.
就是和他們一齊敬拜當地神明的相關作風.
希望不要影響以色列人.
約瑟亞將耶利哥城所有一切都納入滅絕令裡面.
意思是一切的活物都要殺死.
值錢的物件要完全歸入神的家.
其中其他的東西就要銷毀.
雖然滅絕令意味著耶利哥是完全的毀滅.
但是在經文裡面卻給我們看到.
並不是意味著神毫無憐憫.
滅絕令只是神對這些國家計劃的一部分.
在其他經文裡面可以看到.
這些國家是可以有選擇.
正如申命記裡面20章第10節.
在申命記段落的前一點.
他說你臨近一座城要攻打的時候.
先要對城裡的民宣告和睦的話.
他們若以和睦的話回答你.
給你開了城.
城裡所有的人都要給你效勞.
服侍你.
若不肯與你和好.
反要與你打仗.
你就要圍困那城.
所以經文一再強調.
毀滅和拯救的訊息是同時出現.
在約瑟孝記的記載裡.
強調了拉俠和他的全家得救.
希伯來聖經用了86個字來描述拉俠的得救.
而耶利哥的毀滅就用了102個字.
故事最後也以拉俠來結束.
所以拉俠的得救和耶利哥的毀滅.
幾乎是同樣在段落裡很重要.
而第25節的結論裡也描述了.
拉俠是住在這些以色列人當中直到今天.
他和以色列人不分彼此.
已經是他們的一部分.
他不再作家南人.
已經成為了以色列人.

$^{481}$更加在耶穌的家譜裡出現他的名字.
另外經文也提到.
約瑟孝救活他們一家的理由.
就是因為拉俠隱藏了約書亞所打發窺探耶利哥的使者.
這裡用了「隱藏」這個字.
這個字在這裡再次出現.
還記得起初拉俠告訴探子.
他們要離開耶利哥之後.
逃到山裡去隱藏.
還有拉俠也隱藏了知道探子的消息.
在經文的段落裡.
當用隱藏這個字的時候.
特別強調到拉俠保護了探子誤導王的使者.
所以拉俠是拒絕他過去和家南人的關係.
最終他決心轉移效忠以色列人.
他決心和家南人一刀兩斷.
劃清界線.
基督徒的一生也是一場屬靈的真戰.
我們一直持續和罪惡的權勢相爭.
而神是絕對不容忍屬祂的子民被家南人同化.
甚至沾染他們這個以色人神聖的身份.
正如在第十八節裡描述.
只要你們要自己必須謹守遠離者多滅之物.
免得你們徹底滅絕歸神的時候.
從這多滅之物中拿取什麼.
使以色列人成為多滅之物.
惹禍及他.
一再強調我們不要成為被神滅絕的多滅之物.
我們要留心 要謹守 要分別為聖 跟從上帝.
我們選擇效忠耶穌 跟從主 作主的門徒.
我們就要撇下一切走上跟從主的神聖之路.
不再被世界的價值觀所影響我們.
正如保羅在羅馬書十二章一到二節去勸勉我們一樣.
我們一起讀吧 在下面.
一二三.
所以弟兄們.
將身體獻上 當作活祭.
是聖義的 是神所喜悅的.
你們如此事奉乃是理所當然的.
不要效法這個世界.

$^{521}$只要心意根深而變化.
叫你們察驗何謂神的善良.
順傳可喜悅的旨意.
所以我們的確需要跟罪惡.
跟罪惡的權勢劃清界線.
個分別為聖 跟從主耶穌基督的生活.
經文在最後二十六二十七節裡說到.
當時約書亞叫眾人起誓.
有興起重修者耶利哥城的人.
當在耶華面前受咒助.
他立根基的時候必喪長子.
安門的時候必喪幼子.
耶華與約書亞同在.
約書亞的聲名傳揚遍地.
約書亞最後要眾人起誓.
不可以再重修耶利哥城.
重修就必被咒助.
以色列人需要與耶利哥城劃清界線.
最後二十七節裡.
用耶華與約書亞同在.
約書亞的聲名傳揚遍地去落幕.
再次重申耶華與約書亞同在.
呼應一開始的時候.
耶華軍隊元帥的出現.
指示約書亞如何攻破仇敵.
最終約書亞完全信靠馴服.
並帶領以色列人最終得到勝利.
所以對於基督徒來說.
耶利哥的傾倒.
象徵著我們信靠神的能力.
如果我們沒有信心跟從.
什麼都不會發生.
但是我們信心還要加上信靠馴服.
基督徒就能夠克服人生所有困難.
五章十三節到六章二十七節裡提醒我們.
我們需要分別為聖.
各土神喜悅的生活.
神才能夠大大使用我們.
神才會賦予我們任務 使命.
我們一起努力去建立神的國度.

$^{561}$建立我們的團隊 協助教會.
協助我們身邊的團隊.
認識上帝 跟從上帝.
在網上看到一個故事.
話說在美國南北戰爭的時候.
有人問當時的總統林肯.
你們北方的軍隊.
禱告求上帝站在你們那邊.
但是南方軍隊也禱告求上帝站在他們那邊.
請問上帝到底會站在哪一邊呢?.
林肯回答.
我不是禱告求上帝站在我的這一邊.
我是求上帝使我能夠站在他的那邊.
我們的神是信實的.
只要我們站穩在神的那一邊.
他的大能一定可以幫助我們過得勝的人生.
為我們大家都學習約書的功課.
有顆謙卑受教的心.
願意各分別為聖 信靠馴服的生活.
最終同心完成神所託付給我們的事情.
不是要完成我們所想的事.
是要完成神所想的事.
我們一起祈禱.
求你憐憫幫助我們.
讓我們常常面對誘惑挑戰.
我們都能夠堅定站在你的那邊.
讓我們能夠立志過得勝的生活.
過土神喜悅的生活.
過跟從主耶穌基督的生活.
求你幫助憐憫.
當我們一起合作.
一起去建立團隊的時候.
求你幫助你在我們中間帶領著我們.
以致不是以我們的意思.
而是以你的意思.
\newpage



\section{創世記 19:36}
\label{sec:031689}
\textbf{聖經人聖經事:人的方法}~~~~黃克勤牧師博士
\newline
\newline
連結: \href{http://www.church.com.hk/acms/content.asp?site=cdc\&op=show\&type=product\&code=031689\&layout=sermon}{\texttt{031689}} ~~~~ 語音日期: 2026-02-12
\newline
\newline
\hyperref[sec:031688]{\small{< < < PREV SERMON < < <}}
~
\hyperref[sec:index_chronic]{\small{[返順時目]}}
~
\hyperref[sec:031690]{\small{> > > NEXT SERMON > > >}}
\newline
\newline
創世記 19:36
\newline
\begin{longtable}{cl}
\hline
\hline
章節 & 經文 (和合本修訂版)\\
\hline
19:36 & \begin{tabularx}{0.7\textwidth}{X} 這樣，羅得的兩個女兒都從她們的父親懷了孕。 \end{tabularx} \\ \\
[1ex]
\hline
\hline
\end{longtable}
$^{1}$大家好,我們又讀聖經了.
聖經人,聖經事,生命說,宣教網絡為你提供讀聖經.
今次就看看裡面的人和上帝在背後做些什麼.
是幫助我們認識人性更加多.
認識神性更加多.
今次看到有些人用人的方法.
其實今天這一段經文是很久以前曾經說過了.
不過是再重溫.
然後是由這個角度去看這段經文.
創世記十九章第三十節.
都是和亞伯拉罕有關係的.
亞伯拉罕坐在橡樹下.
有三個人迎面而來.
來到他住的地方.
亞伯拉罕看到有三個人跑出去.
請三位來我家住坐吃點東西再前行.
很恭敬.
馬上預備食物給他.
還找一隻美的羊來宰給他.
亞伯拉罕站在那裡服侍三個人.
為什麼會這樣呢?.
很有可能.
聖經沒有直接說.
不過似乎亞伯拉罕有一個感動.
是為了神感動.
他知道來的三個人是神.
於是他不敢怠慢跑出去.
如果願意來我家吃點東西.
加添心力再前行.
於是那三個其中一個就說.
明年今天你太太會生個兒子.
然後撒拉笑.
哈哈,誰會呢?.
是這樣的.
而這個時候.
亞伯拉罕已經和羅德分開了.
而他的兒子羅德.
羅德的父親是亞伯拉罕的兄弟.
所以亞伯拉罕也很照顧羅德.
羅德這個字.

$^{41}$其實是遮蓋.
因為怕住在索爾.
由於之後那三個人.
一個留下.
另外兩個就前行去索多瑪俄摩拉平原地方.
因為那裡有很多城市.
很壞,大家記得這段聖經.
於是忽然之間不再是那三個人.
而是改了.
又說我怎可以瞞住我的朋友亞伯拉罕呢?.
於是他告訴他.
他們是前行去毀滅索多瑪俄摩拉和平原上的城.
於是亞伯拉罕就說.
如果有五十個義人是否也毀滅城?.
四十個?三十個?二十個?.
大家記得這段聖經.
是羅德在說數.
這裡就改成耶和華和亞伯拉罕聊天.
真厲害.
於是耶和華就當他父上帝天父.
然後前行去做事的.
就是聖子和聖靈.
他們就去救羅德.
然後就叫他快點離開.
上山.
我都說我不想住在山上.
有個小城叫索爾.
不如讓我去那個小城.
那裡沒什麼害.
那邊緣地方.
那兩個就說好.
你去那裡.
你快點去.
如果你不去了.
我們做不到事.
因為開始有毀滅.
我留索爾在那裡的小城鎮給你住.
後來羅德想了想.
還是不要住在平原.
真的很危險.

$^{81}$羅德怕住在平原的小城.
索爾就和他兩個女兒.
由索爾上去住在山裡.
本來說不想住在山裡.
現在就住在山裡.
比較安全.
為什麼會這樣呢?.
不知道.
會不會是平原呢?.
因為死了很多人.
火燒,索多瑪,俄摩拉.
索爾很害怕.
又或者變成很凌亂.
不知道有沒有人來打劫搶劫.
很凌亂.
住在山上.
他和兩個女兒住在一個洞裡.
以色列如果有機會去.
那些石山.
古老石山.
多數都是石.
沒有樹,沒有草.
都是石.
很多洞.
一定有山洞.
大女兒對小女兒說.
我們的父親老了.
地上沒有人按照世上常規來這裡.
上來我們這裡.
那些山上的洞.
哪裡有人去呢?.
其實那裡很艱難.
水,食物都很艱難.
更加沒有人會上去那裡.
沒有人來.
那怎麼辦呢?.
來我們可以叫父親喝酒與他同寢.
這樣我們好從他存留後裔.
他們兩姐妹就做了這個決定.
留意一下.

$^{121}$這個決定.
第三十一節就說.
是大女兒對小女兒說的.
為什麼大女兒對小女兒呢?.
很有可能是因為爸爸老了.
呆呆的.
大女兒就要撐起頭家.
與小女兒搞定父親要吃的東西.
大家要的食物.
各樣東西.
釀些酒.
都是要的.
可能羅德以前在山下都有些錢.
不表示他不下山.
下山買東西吃都可以.
於是大女兒對小女兒.
大女兒有一個計劃.
就是沒有人來.
那怎麼辦呢?.
不如我們叫父親喝酒與他同寢.
這樣我們就從他存留後裔.
都沒有別的男人.
這個從他存留後裔.
我們好從他存留後裔.
這件事究竟是什麼來的呢?.
是我們好想有後裔.
還是我們好想我們的家族有後裔.
爸爸有後裔.
因為沒有男人來.
更加沒有別人來.
不要說男人.
沒有人來.
沒有男人我們生不了.
如果我們想要後裔.
一定要有男性.
又沒有.
我們好從他存留後裔.
不過這裡這樣看.
如果照這個字這樣看.
是我們叫父親同寢.

$^{161}$我們從他存留後裔.
不是為他存留後裔.
我們好為他存留後裔.
不是從他.
是我們從他這個男性存留後裔.
即是我們的後裔.
兩位大姐都好想有兒子.
都不想絕子絕孫.
所以不如我們和爸爸同寢.
從他存留後裔.
我們從他存留後裔.
於是那天叫父親喝酒.
不知道的.
大女兒灌醉了爸爸.
和他同寢.
何時躺下何時起來.
羅德都不知道.
第二天.
大女兒對小女兒說.
昨晚我和父親同寢.
今夜我們再叫他喝酒.
你就可以和他同寢.
這樣我們好從父親存留後裔.
不是為父親存留後裔.
我們存留後裔.
是這樣的.
於是喝酒小女兒和他同寢.
何時躺下何時起來.
父親都不知道.
其實羅德知不知道這件事?.
他不知道.
因為醉了睡到第二天醒來.
見到兩個女兒.
這麼早起床這麼晚起床.
他不知道昨晚發生過這件同寢的事.
羅德是懵懵地.
不知道發生什麼事.
他不了解.
這樣羅德兩個女兒.
聖經說.

$^{201}$兩個女兒都從他父親存留後裔.
不是為父親懷孕.
是我們兩個有後裔.
不想就此完.
當然有相當成份為了家族會完.
亦都說我們有小孩.
可能女性都歡喜小孩.
不過沒有辦法.
就是這樣.
這件事來說.
羅德可能不知道.
何時才知道?.
為什麼你們有後裔?.
懷孕為何你們會懷孕?.
沒有人理會.
可能這個時候他們才跟父親說.
其實是你的.
不是吧?.
怎麼搞的?.
羅德想不想我就不知道.
但是已經米已煮成飯.
羅德亦都不會說.
不如放下吧.
我都不想.
這樣不行.
在山上也沒有辦法.
羅德說.
好啊,你都很聰明.
還是說.
你這樣很不開心.
我不想的.
為什麼你們不問問我呢?.
不過已經做了.
可能不知道上帝是否歡喜.
怎麼搞的?.
羅德和兩個女兒都認識上帝.
因為他們真的見到.
索多瑪,阿摩拉和皮兒的城鎮毀滅.
那兩個是救他們出來的.
如果是三一上帝的話.

$^{241}$子和聖靈是救他們出來的.
他們應該認識上帝.
上帝喜不喜歡呢?.
不過已經懷孕了.
怎麼辦呢?.
難道叫他們拿尾巴也做不到?.
可能羅德都相當之遺憾.
怎麼會這樣呢?.
再看看三十七節.
大女兒生了兒子.
叫他的名字做摩亞.
摩亞的意思是他父親的.
真有趣.
直接改名叫父親的.
大女兒生了兒子叫做摩亞.
就是今天摩亞人的始祖.
小女兒又生了兒子.
叫他變阿米.
變阿米是Ben.
Ben是兒子.
兒子是Ami.
Ami就是阿米.
變阿米.
就是我民的兒子.
所以阿米就是我的百姓.
我的民族.
意思是.
由這個人.
將來會繁衍一個民族出來.
就是我民的兒子.
就是現在阿滿人的始祖.
大家都知道由聖經裡面.
摩亞人和阿滿人.
時時都和以色列人作對.
大家是敵人.
經常打仗.
摩亞阿滿經常來侵佔以色列.
和他打仗.
但其實你會看到.
羅德是亞伯拉罕的侄子.

$^{281}$大家都有相同的血脈源流.
因為亞伯拉罕和羅德的父親.
都是相同的父親.
本來是兄弟.
親戚.
摩亞阿滿本來是以色列的親戚.
但他們成為世代的仇人.
都很慘.
這裡有一幅地圖.
不知道大家在視頻裡看得清不清楚.
這裡是地中海.
V字形長的就是以色列.
耶路撒冷在這裡.
死海在這裡.
耶路撒冷在這裡.
這裡就是以色列.
今日約旦就在以色列旁邊.
中間有條街就是約旦河.
這邊是約旦.
這邊是以色列.
死海的頂部就是主耶穌洗禮的地方.
那裡洗禮的約旦河原來很窄.
大概在這邊吸一口氣.
拼命游過對面.
不用換氣都可以游得到.
很近的.
我想是十米之內.
那邊就是約旦.
這邊有以色列士兵.
那邊有約旦士兵.
大家都很悠閒.
約旦很近.
大家記得約旦的版圖是這樣的.
很有趣,很不規則.
一塊塊凸出來的.
這邊是以色列.
約旦的首都叫做安曼.
很近的安曼.
安曼是甚麼呢?.
安曼是今日叫法.

$^{321}$其實以往叫做阿曼.
阿曼地.
如果看舊約聖經.
阿曼就在這個位置.
在以色列的這一邊.
即是今日約旦的地方.
低一點就是摩亞.
如果是早期的以色列時代.
阿曼摩亞在以色列旁邊.
經常跟以色列打.
大家是兄弟,不過變成對抗.
這張圖可能清楚一點.
這是英文地圖,這是以色列.
這是約旦.
這是敘利亞.
所以我們去哥蘭高地.
就會看到敘利亞的境內.
這裡是黎巴嫩,約旦首都安曼.
Amen.
以前叫做阿曼.
阿曼摩亞.
就是這些地方.
大家會知道.
其實變成世代的仇人.
不是太好.
今天不叫反省.
反省不叫金句.
以免你覺得這樣是對的.
我今天選這段經文來說.
聖經人聖經事是因為早一陣子.
有人問我.
羅德兩個女兒.
其實神是什麼態度呢?.
我聖經沒有直接說.
不過我看其他經文.
我覺得神不太歡喜這件事.
就是沒有兒子.
我有理由與父親同寢生兒子.
這是正常的.
我覺得不是正常.

$^{361}$上帝未必歡喜.
是兩個女兒特別的大女兒歡喜.
反省的地方在這裡.
創世記十九章三十六字.
羅德與兩個女兒.
都從他父親懷了孕.
我有理由,有委屈.
沒有男人上來.
要生子立後.
這件事很正常.
只有一個男的就是我父親.
於是我與他生兒子.
也是家族.
以往很多人都是這樣.
以往很多時候.
甚至歐洲.
很多王族不想與其他平民百姓通婚.
所以一定要與王族的人結婚.
於是保存著王族純種血脈.
於是近親交配.
變成差勁.
我不知道.
但兩個女兒與羅德懷孕.
會不會兒子生出來.
有些遲鈍.
低低的,會不會呢?.
不知道,不過未必馬上會發生.
但這件事我不認為上帝會很歡喜.
很祝福.
照法在聖經看不到.
而他們是成為以色列的仇敵.
當成為以色列的仇敵.
即是表示上帝並沒有歡喜這件事.
但生子立後這件事.
天公地道,要做.
為何要怪他.
我只能這樣說.
所以不能鼓勵這件事.
所以不叫金句,叫反省.
不過我順帶一提反省.

$^{401}$CB23篇第一節.
耶和華是我的牧者.
我必不自缺乏.
這節聖經唸給女兒聽.
是不是這個意思呢?.
不是.
我想說.
這節聖經耶和華是我的牧者.
我必不自缺乏.
因為耶和華是我的目洋人.
I shall not want.
祂是我的牧者,我已經夠了.
其實原來的意思是.
耶和華是我的牧者,我已經夠了.
我沒有什麼特別需要.
如果神給我就給.
我沒有什麼特別需要.
即是說我有上帝就夠了.
祂要給什麼我,沒有給什麼我.
不要緊,我有上帝就夠了.
這篇聖經是給我們反省的.
如果耶和華沒有給我子女.
那怎麼辦?.
那就沒有了.
斷子斷孫.
我們的家庭就沒有了.
如果是這樣,那就沒有了.
上帝有這樣的心意,那怎麼辦?.
於是不通.
我嫁給老公,老公很久都生不了.
我們都沒有子女.
驗驗老公的精子不好.
我去生子立後,於是我去找別的男人.
這不能夠讓我們承認是可以的.
沒有了,那上帝既然沒有給我就沒有了.
你可以求.
或者上帝派一對獵人上來.
你可以求,上帝可能會給.
因為祂是我的牧者,祂會供給我們.
有祂,我就沒有什麼缺乏,就夠了.

$^{441}$我沒有什麼特別要求,因為有主德了.
可能上帝會派獵人上來,不知道.
又或者沒有,沒有了.
譬如哈拿去祈禱.
哎呀,我沒有子女,很慘.
求上帝給我一個.
如果祂真的給我,我就要他一生在聖殿中服侍.
上帝說好,我聽你祈禱,撒母耳出來.
又或者撒迦利亞和伊利沙伯.
哎呀,沒有啊,然後禱告.
神就說好,我給你,施洗有害.
如果神給就有了,沒有給怎麼辦?.
沒有了,很多人都沒有了.
那我要用我的辦法.
所以我今天想講的題目就是.
人的方法.
我要用我的辦法.
但其實未必是上帝的祝福.
未必是上帝的祝福.
可能今天摩艾雅滿都很困難.
而且整個亞同城的百姓作對.
神似乎並沒有特別的祝福.
雖然他們都認識上帝,並沒有特別的祝福.
今天約旦已經變成一個伊斯蘭國家.
並不認識.
也認舊約,也說是神的.
但信仰又走到另一邊了.
似乎也沒有上帝特別的祝福.
大家要記得.
耶和華有一個牧者.
I shall not want.
如果祂有給就給.
沒有給也好.
沒有給就沒有了.
不要用太多人的方法去做.
結果有時候未必好的.
神未必祝福,未必歡喜.
我們一起禱告.
求主幫助我們.
讓我們再次仰望.

$^{481}$記得你是我們生命的主.
又是我們的牧者.
一生都是你帶領的.
如果有你,我們就夠了.
有些事你賞賜就賞賜.
沒有賞賜,也有你的恩典.
也有你自己的原因.
我們就信服.
我們就等候.
我們就跟隨.
給我們有這個信心.
給我們有這個信服.
一路走的時候.
仰望你自己.
每一個時刻的供給與保護.
禱告奉耶穌基督名求.
阿們.
\newpage



\section{詩篇 104:24-35}
\label{sec:031690}
\textbf{總是記得創造主}~~~~江海薇傳道
\newline
\newline
連結: \href{http://www.church.com.hk/acms/content.asp?site=cdc\&op=show\&type=product\&code=031690\&layout=sermon}{\texttt{031690}} ~~~~ 語音日期: 2026-02-22
\newline
\newline
\hyperref[sec:031689]{\small{< < < PREV SERMON < < <}}
~
\hyperref[sec:index_chronic]{\small{[返順時目]}}
~
\hyperref[sec:031691]{\small{> > > NEXT SERMON > > >}}
\newline
\newline
詩篇 104:24-35
\newline
\begin{longtable}{cl}
\hline
\hline
章節 & 經文 (和合本修訂版)\\
\hline
104:24 & \begin{tabularx}{0.7\textwidth}{X} 耶和華啊，你所造的何其多！都是你用智慧造成的，全地遍滿了你所造之物。 \end{tabularx} \\ \\ \relax
104:25 & \begin{tabularx}{0.7\textwidth}{X} 那裡有海，又大又廣，其中有無數的動物，大小活物都有。 \end{tabularx} \\ \\ \relax
104:26 & \begin{tabularx}{0.7\textwidth}{X} 那裡有船行走，有你所造的力威亞探悠游在其中。 \end{tabularx} \\ \\ \relax
104:27 & \begin{tabularx}{0.7\textwidth}{X} 這些都仰望你按時給牠們食物。 \end{tabularx} \\ \\ \relax
104:28 & \begin{tabularx}{0.7\textwidth}{X} 你給牠們，牠們就拾起來；你張手，牠們就飽得美食。 \end{tabularx} \\ \\ \relax
104:29 & \begin{tabularx}{0.7\textwidth}{X} 你轉臉，牠們就驚惶；你收回牠們的氣，牠們就死亡，歸於塵土。 \end{tabularx} \\ \\ \relax
104:30 & \begin{tabularx}{0.7\textwidth}{X} 你差遣你的靈，牠們就受造；你使地面更換為新。 \end{tabularx} \\ \\ \relax
104:31 & \begin{tabularx}{0.7\textwidth}{X} 願耶和華的榮耀存到永遠！願耶和華喜愛自己所造的！ \end{tabularx} \\ \\ \relax
104:32 & \begin{tabularx}{0.7\textwidth}{X} 他看地，地便震動；他摸山，山就冒煙。 \end{tabularx} \\ \\ \relax
104:33 & \begin{tabularx}{0.7\textwidth}{X} 我一生要向耶和華唱詩！我還活的時候，要向我的神歌頌！ \end{tabularx} \\ \\ \relax
104:34 & \begin{tabularx}{0.7\textwidth}{X} 願他悅納我的默念！我要因耶和華歡喜！ \end{tabularx} \\ \\ \relax
104:35 & \begin{tabularx}{0.7\textwidth}{X} 願罪人從世上消滅！願惡人歸於無有！我的心哪，你要稱頌耶和華！哈利路亞！ \end{tabularx} \\ \\
[1ex]
\hline
\hline
\end{longtable}
$^{1}$各位親愛的姐妹們 午安 主佑平安.
今天是年初六.
在此也祝大家新年快樂 新年滿福 身體健康.
中國人有一句成語.
就是「一年之計在於春」.
2月4日開始已經是立春的時節.
正所謂「春回大地 萬象更新」.
這一句都是來自中國的傳統文化.
通常都是用來描述春天來臨.
大地回暖 氣候變得溫暖.
植物開始生長 萬物開始恢復生機.
「萬象更新」就是指宇宙間的一切景象.
一切事物都渙然一新 發生了變化.
這一句說話都是形容在春天的氣息當中.
自然界和社會上一切事物都顯得充滿活力.
煥發出新的生命力.
今天跟大家分享的這一首詩篇.
是稱頌神在創造和對受造物.
在供應上有很豐富的慷慨.
讚美祂去做出一個多元多樣.
錯綜複雜的世界.
一個肥沃富曉 生命力旺盛的世界.
整個世界都是神給一切生物的禮物.
詩人在口中的神是非常慷慨.
一切的美善都是從祂而來.
讓我們從神的創造去敬拜 讚美和禱告.
首先我們看看神所創造的其中一樣.
北極光.
多年前自己和朋友一起去的時候.
都遇不到 看不到 特意要去的地方.
現在從一些朋友家附近已經很容易看到.
為什麼這麼容易看到呢?.
一會兒再繼續說.
經文在第一至二節當中.
詩人去提醒我們.
這一篇其實是一個很輝煌的讚美詩歌.
去頌讚造物主的榮耀.
和祂偉大奇妙的創造.
在頌讚神的偉大的時候.
作者用了兩幅來自創造世界的圖畫.

$^{41}$去反映造物主各個的面向.
他用光作為開始.
就是你穿著光的外袍.
鋪展諸天的星芒.
在這個詩意的描繪當中.
神用光芒去遮蓋自己.
就好像披上了斗篷一樣.
亦都好像窗簾一樣.
展開了這個窮倉.
反映著祂的榮耀.
我相信對於我們身在香港這個城市.
對光是不陌生的 是很熟悉的.
因為在香港其實燈光是很光的.
我不知道大家如果去到有其他地方.
你會發覺相對比之下.
其實其他地方可能夜晚是很暗的.
香港夜晚都是很光的.
所以我們對光的敏感度都很強.
但是我們是否真的會看得見這個真光呢.
這個真光是可以穿越靈魂.
指引方向的光.
祂在哪裡呢?.
詩人其實就是提醒我們要看天.
看到神的偉大.
祂是愛這個世界.
祂的愛不單止是臨到人的當中.
在《約翰福音》三章十六節.
「神愛世人 甚至將祂的獨生子賜給他們.
叫一切順他的 不至滅亡 凡得永生」.
如果我們看原文來說.
可以說神不單止愛人.
其實神是愛整個的世界.
是「God loved the world」.
在原文是用了「Cosmo」的字.
如果以現在的宇宙學或宇宙論來說.
這個字就是「Cosmology」.
就是神是愛世界.
神所愛的是整個受造的世界.
然後詩人就是叫我們要去看地.
地就是三節至到第五節.

$^{81}$「在水中立樓閣的棟樑.
用雲彩為車鏈.
藉著風的翅膀而行.
以風為使者 以火焰為瀑翼.
將地立在根基上.
使地永不動搖」.
神是以雲彩作為祂的車.
這個車鏈其實是代表著帝王坐的車.
然後有風作為祂的使者.
將地立在根基上.
使地永不動搖.
為甚麼可以呢?.
因為祂是創造主.
祂可以控制大自然.
之後我們又再看看另一方面有關於海.
「你用深水遮蓋地面.
猶如衣裳.
珠水高過山嶺.
你的叱tak一發水便崩途.
你的雷聲一發水便崩流.
珠山升上珠谷沉下.
歸你為他所安定之地.
你定了界限使水不能過去.
不再轉回遮蓋地面」.
在這裡我們很清楚看到.
只有神有能力去定界線.
是這個地去定界線.
和海亦都是定界線.
不過我們現代的城市人.
大自然會受我們的控制.
我們習慣人去控制環境.
例如我們當中應該很多人都喜歡行山.
我自己都喜歡.
去接近大自然.
有人亦都喜歡行商場.
我都相信.
我們城市的設計.
不知道大家有沒有留意.
有很多商場.
無論大大小小其實有很多商場.

$^{121}$在我們城市裡面.
我們的設計其實.
可以說是用一個「快」字來構思.
我們的社區的行人天橋.
當我們去行行人天橋的時候.
很快便會將我們引渡到一些商場.
在行人天橋當中.
我們不像在路面行走.
大家不知平時如果要找路的時候.
我自己就是.
如果用Google Map要找路的時候.
行行在地面.
然後外面要入到商場就找不到路.
就沒有了那個訊號.
又要出到外面地面才會找到一個正確的方向.
在這裡可以說其實.
一個設計其實是在人的控制裡面.
去到商場裡面可能都會是.
當你都累的時候走到.
想坐下.
發覺其實都沒有甚麼椅子可以讓你坐.
因為他們是要做買賣的地方.
做買賣是怎樣.
不讓你慢下來.
你們要去吃東西坐下來消費.
你們如果想喝杯咖啡就去咖啡店.
坐下來慢慢這樣.
整個模式設計都會在一個控制的空間裡面.
如果我們去看神所創造給我們這些事物的時候.
這個有天地海這些大自然的景色.
這些自然的環境一直是存在著的.
反而當我們看回這些自然的環境的時候.
人就顯得很渺小.
我們就是從這些自然的環境裡面.
我們會認知到人其實不是世界的中心.
就好像我們行山的時候.
不知道大家有沒有試過.
行山的時候不小心扭傷或者跌倒.
其實沒有人會記得你的.
多年之後沒有人會記得你.

$^{161}$或者江海明傳道有一天行大霧山.
不小心扭傷跌倒.
但是如果這個山的山名.
我們數都數得出香港有些行山的地方.
大霧山 飛鵝山 獅子山 白仙嶺等等.
我們很多年後都會還記得這個山名.
因為它們屹立了很多年.
這些山巒是很偉大.
但是神是更加的偉大.
之後我們再去看一看.
詩人所說的有關山名.
就是在第十節至第十八節.
「耶和華使泉源湧在山谷 流在山間.
使野地的走獸有水鶴 野驢得解其鶴.
天上的飛鳥在水旁住宿 在樹枝上啼叫.
他從流角中囂冠山嶺.
因他作為的功效 地就豐足.
他使草生長 給鹿畜吃.
使菜蔬發長 供給人用.
使人從地裡能得食物.
又得酒能悅人心.
得油能潤人面.
得糧能養人心.
街尾的樹木就是尼巴嫩的香柏樹.
使耶和華所栽種的都滿了汁醬.
雀鳥在其上踏窩 子於鶴.
松樹是他的房屋.
高山為野山羊的住所.
岩石為沙番的藏處」.
我們在日常的生活當中.
我們經常都會記得一些很重要的日子或事情.
我們會包括在工作上要記得一些日子.
我們要講 Deadline.
要記得一些工作怎樣安排好.
或者我們的親朋戚友.
他們生日我們要安排約人慶祝.
否則就很難搞了.
真的忘記了很多事情.
有些什麼方法去加強人的記憶.
坊間有很多講座會教一些記憶的方法.

$^{201}$例如長期記憶的方法.
重複去背誦.
短期的記憶方法.
一首詩篇就是.
「神的靈是以智慧去創造萬物.
和去持續掌管,供應和更新」.
在詩篇一開始和結束.
都是以「我的心你要稱頌耶和華」.
這一首就是有一個首尾呼應的文學手法.
表示詩人對稱讚神是一生最重要的事.
是一生不可以不做的.
所以他不斷去提醒自己.
提醒自己究竟有什麼事情呢?.
就是神豐富的供應.
由大自然去供給.
就是「流國中囂冠山嶺.
使草生長 給六畜吃.
使菜蔬發長 供給人用」.
如果我們看清楚.
其實這一條是食物鏈.
就是神的慷慨.
他可以不做新畜甲.
其實對他沒有影響.
但是因為神的豐盛.
他要和不同的受造物分享生活.
神賞賜生命氣息.
除了讓受造物可以存活.
還有一個目的.
就是讓受造物享受生命的喜樂.
就好像小鳥可以很開心地提叫.
所以這裡去描述.
受造物其實是完全和直接去依賴神慷慨的賞賜.
這裡神是很近距離和他的受造物互動.
這裡就是去強調受造物去依賴神.
神也是同時賦予受造物能力.
例如雀鳥可以有能力去捉巢.
人類是可以耕作.
也可以架船.
所以在這個整個受造的世界裡.
我們會看到.

$^{241}$神也會對人有一個關注.
有一個關注就是在經文第十五節.
「有得酒能悅人心.
得油能潤人面.
得糧能養人心」.
就是神的慷慨.
我不知道大家有沒有養寵物.
我們知道可能有些其他的動物.
牠們吃飽了.
最基本的就是生存.
是為了生存.
而人吃食物不單止生存.
我們也要吃得開心.
現在的食物有很多種類.
很豐富.
以前我們吃牛.
還沒有這麼多種類的時候.
我們吃牛是很普通的牛肉.
現在多了很多種類.
已經是和牛.
和牛還要分不同的級數.
吃豬肉也是很基本的.
普通的豬肉.
但現在有很多不同的肉.
豬肉我們叫做「黑毛豬」.
人是很奇妙的.
有神的豐富供應.
如果另外有酒.
酒也是.
可以有葡萄釀酒.
有米酒.
有梅子釀酒.
其實這些都不一定需要.
其實對神沒有影響的.
如果沒有的話.
但神也是很慷慨地賜給我們.
在這裡去形容每一個受造物.
這些動物.
神都有為牠們提供棲息的地方.
有充分的水源.

$^{281}$和足夠的食物.
但其實我們會看到.
這裡沒有特別去將人的地位提升.
其實我們人和其他的受造物.
都在這個世界.
有各自的一個位置.
人類在當中和其他的受造物.
成為一個夥伴.
神是同樣為牠們.
為我們去預備生命.
住所和食物.
之後我們去看看.
在經文第十九節至二十三節.
「你做月亮已定季節.
太陽曉得何時西沉.
地派顯黑暗就有黑夜.
白獸便在林中出沒.
少壯獅子吼叫覓食.
搜尋神所賜的食物.
黎明時分.
他們便退回洞中休息.
這時人出去工作.
勞碌直到夜晚」.
這裡是講述時序和季節的變換.
我們知道整個地球有足夠的條件.
成為所有生物棲息的地方.
為不同的受造物提供不同的生活環境.
夜晚野獸出來尋找食物.
獅子的特性我不知道大家知不知道.
通常牠在夜晚行動.
牠是一個短跑高手.
如果我們健健康康的身體.
其實可以和獅子長跑.
但是短跑一百米就不行了.
這裡提及到.
連獅子夜晚在森林尋找食物.
神都會為獅子預備獵物.
所以在這裡到經文第24節.
為第10至23節作為一個總結.
夜禍所做的何其多.

$^{321}$都是你用智慧造成的.
遍地滿了你的豐富.
牠的供應很豐足.
這裡描述到有關海洋.
那裡有海又大又廣.
其中有無數的動物.
大小活物都有.
那裡有船行走.
有你所做的惡魚.
游泳在其中.
這裡提到在海洋當中.
有船和惡魚.
似乎我們看到在海洋的生物生活.
只有一些很大型的生物或魚類.
只可以在海裡生活.
而人就拿著船在海上航行.
這樣就顯得人很渺小.
如果從希伯來文.
聖經其他的經文包括《約伯記》.
我們會知道惡魚在當中比喻雲頓.
作者同時提到惡魚和船.
其實是說人類有勇無謀.
為什麼這樣說呢?.
因為在大海當中.
竟然拿著船在海面上航行有惡魚.
人類是讓自己身陷險境.
因為古代以色列人很少這樣做.
因為他們不熟悉大海的環境.
他們只會在加里尼湖打魚.
所以在新約的時候.
門徒在大海遇到風浪很害怕.
就要求耶穌平靜風浪.
如果說到熟悉大海的地方.
就是推羅黎巴嫩.
因為它是由地中海的東面向西面發展.
所以在這裡的經文提到船和惡魚.
究竟這個海和信仰有什麼關係呢?.
這個海其實都是在神的腳下.
因為在迦南宗教當中.
大海也是一個神明.

$^{361}$迦南有幾千個神明.
大海是其中一個.
是代表著雲頓的意思.
所以詩人提到大海也要臣服在神的權威底下.
神令到雲頓的海水.
成為了有生命的泉水和河水.
提供給大小活物在當中可以生存.
大海當中已經沒有任何的阻攔.
成為了提供船隻可以行走的通道.
詩人所認識的權能.
就是神掌握著.
所以惡魚也能夠在當中遊泳.
因為迦南的宗教.
惡魚是代表著恐怖的生物.
代表著製造混亂的怪獸.
但是在神的管轄面前.
這個獸也變成其中一個沒有殺傷力的生物.
因為神喜愛牠一切所做.
所以受造萬物的喜樂.
其實就是分享著神的喜悅.
一切都是仰賴著神祂按時的供應.
一切的受造物都是依賴著神祂的供應.
我們會想到神祂按時的供應.
就是以色列民他們要出埃及的時候.
神祂是給他們按時的賜予給他們的嗎哪.
詩人就是去表達出神的供應是全被的.
祂所創造的大自然是會按時出產豐盛的食物.
我們中國人有一句話.
「靠山食山 靠海食海」.
意思是如果你住在山上.
山上有很多的資源和食物.
提供你生存的資源.
如果住在海上 海產就是足夠供應和有謀生的需要.
如果現在我們想想要吃這些蔬果田園.
會想到什麼地方.
去哪裡 去新界.
新界有很多這些田園種蔬果的地方.
士多啤梨園 我們去過.
如果要吃海產的地方.
海鮮 想起什麼地方.

$^{401}$長洲 西貢 樓房山.
想起有這些地方.
所以都會按時.
神是按地方時間去供應.
大家在農曆年期間 有沒有逛花市.
有沒有買蓮花過年.
新年的蓮花有什麼代表.
考考大家.
蘭花.
水仙.
吉.
或者五代同堂.
桃花.
這些都是按著時節去提供需要.
平時就不會有.
平時過了新年之後就沒有了.
亦都是因為時節 氣候.
在之前的時間適合去栽種這些植物.
然後詩人亦都提到.
「你給他們 他們便拾起來」.
「你將手 他們飽得未食」.
「你掩面 他們便驚惶」.
「你收回他們的氣 他們就死亡」.
「歸於塵土」.
「你發出你的靈 他們便受造」.
「你使地面更換為新」.
這裡有一些字我們可以去反省.
有「你給他們」這個「給」字.
和「他們便拾起來」.
即是表示神願意給予人.
就可以拾起來.
拾起來即是我們會得到.
或者我們去收集.
如果神不給予.
人就不能去收集.
然後亦都有「掩面」.
有「神去將手 他們飽得未食」.
就是這個「將手」和「飽得」.
神去供應一切的需要.
因為祂去慈愛.

$^{441}$祂的手隨時向世人張開.
這讓我們想到主耶穌.
祂自己的雙手.
就是張開 被掛在十字架上.
都是為著世人的罪而張開.
為的是要我們飽得這個赦罪之恩.
之後經文第29節說.
「你掩面 他們便驚惶」.
掩面驚惶.
神如果要轉身離開.
不施予這個恩慈.
我們就會驚慌.
但是神是信實的.
祂答應拯救私恩到底.
縱然以色列人時常背叛神.
但祂也沒有掩面不救他們.
而是按著應許帶領他們進入迦南.
反而我們要去仰望神.
祂總是不向他們轉身離開.
主耶穌也在當中為我們世人.
作為了贖罪制.
另外私人提到.
「你收回他們的氣 他們就死亡 歸於塵土」.
收回 死亡.
神收回祂賜予人的那一口氣.
人就死亡歸於塵土.
神是賜生命的靈.
我們的生與死都是在祂掌握當中.
今天我們仍然享受著生活.
也是證明神對世人的憐憫和恩慈.
祂沒有收回我們的氣.
祂賜下祂的獨生子.
主耶穌自己用祂的身體.
去廢掉這個冤仇.
在十字架是使兩孝歸為一體.
與神和好.
我們都是因著主耶穌基督.
我們能夠與神和好.
我們也能與其他受造物和好.
使人告訴我們 神是一切萬物的擁有者.

$^{481}$接著在第三十節.
「你發出你的靈 他們便受造.
你使地面更換為生」.
這裡是頌讚神的創造和更新的恩典.
神發出的靈可以保守萬物的生命.
得著更大的能力.
使一切事物都能夠更新.
當我們去默想這段經文的時候.
我們盼望更深切地認識這位神.
祂賜生命更新的靈.
幫助我們重新得力.
我們會看見一切都是以神為中心.
以神為中心的讚美生命.
最後一段是三十一節至三十五節.
「願耶和華的榮耀傳到永遠.
願耶和華喜悅自己所做的.
你看地 地面震動.
他我生 生就無焉.
我要一生向耶和華唱詩.
我還活的時候 要向我神歌頌.
願他以我的默念為甘甜.
我要因耶和華歡喜.
願罪人從世上消滅.
願惡人歸於無有.
我的心啊 你要稱頌耶和華.
Hallelujah」.
最後在這裡.
第三十一節至三十二節.
施人去提醒我們.
面對著這位榮耀和權柄的神.
滿有智慧作為的神.
人和萬物都應該要讚美祂.
在這裡會看到地震和火山爆發.
很明顯都是神去發動.
不過這個現象.
我們可以被視為是神祂自己的顯現.
就好像在西乃山一樣.
是神祂要彰顯祂自己的榮耀.
不是在當中要危害聖命.
不過在這段經文裡.

$^{521}$一直都是去頌讚 讚美的時候.
有一個例外的地方.
就是施人在第三十五節.
突然間中斷了他的敬拜 讚美.
他轉為一個禱告.
是禱告什麼呢?.
就是「願罪人從世上消滅.
願惡人歸於無有」.
意思就是說.
如果沒有人類的話.
這個美好的世界.
就不會被毀滅 不會被摧毀.
一個舊約的聖經學者.
他可以說是暴露鴿蠻.
他提到罪人是什麼.
在這裡所說.
他說「不相信受造界的生命是神慷慨賜予的.
這些人無疑是想爭取更大的自主權.
不承認這個受造界是創造.
是由神創造的主所管理和托住的.
受造界是要用非常恭敬嚴肅的眼光.
去看待這位神.
他不容忍創造這個字眼遭到踐踏.
因為創造就是等同於禮物.
依靠和極度的慷慨」.
詩人唯一提到人類的地方.
就是人類去掠奪神的創造.
我們現在會看到.
國與國之間去搶奪天然資源.
人與人之間有沒有去掠奪其他人所擁有的東西.
人與其他受造物之間.
有沒有去掠奪他們生存的權利.
人與人之間.
有沒有因著自身的利益而去掠奪其他人的權利.
人與人之間有沒有因著去控制.
而是要控制其他人自身擁有神的恩典.
詩人提醒我們.
我們要以生命去讚美神為祭.
因為神是萬物的根源.
萬物都是依靠神的靈.

$^{561}$去掌管,保存,供應和更新.
下文我們也可以學習詩人所說的.
其實在她的讚美裡面.
她看得見神的榮耀.
在這裡和大家看另一個版本的一節經文.
這是來自《新普及譯本》.
「我的心你要稱重耶和華」.
這裡她所譯的就是.
「我身我心 要讚美上主 要讚美上主」.
這個不單只是心.
我們還要有一個身體力行 有一個行動.
我們怎樣去珍惜神所賜予給我們.
亦都去所愛神 所愛的受造物.
我們一起同心祈禱.
太后的天父 我們多謝你.
願意我們都能夠從詩篇裡面去學習.
我們去愛惜你所創造.
我們去愛惜大自然.
我們去愛惜你豐富的賜予.
你慷慨的恩典.
叫到我們得著這些慷慨的恩典的時候.
我們亦都懂得與人分享.
因為一切都是從你而來.
願意我們都能夠繼續將你自己的話語.
成為我們日常的提醒.
多謝你聽我們的祈禱.
我們祈禱交托 奉主明求 阿們.
\newpage



\section{哥林多後書 11:1-13:14}
\label{sec:031691}
\textbf{哥林多前後書 (6) 保羅以患難為榮}~~~~劉少康牧師
\newline
\newline
連結: \href{http://www.church.com.hk/acms/content.asp?site=cdc\&op=show\&type=product\&code=031691\&layout=sermon}{\texttt{031691}} ~~~~ 語音日期: 2026-02-11
\newline
\newline
\hyperref[sec:031690]{\small{< < < PREV SERMON < < <}}
~
\hyperref[sec:index_chronic]{\small{[返順時目]}}
~
\hyperref[sec:031692]{\small{> > > NEXT SERMON > > >}}
\newline
\newline
哥林多後書 11:1-13:14
\newline
\begin{longtable}{cl}
\hline
\hline
章節 & 經文 (和合本修訂版)\\
\hline
11:1 & \begin{tabularx}{0.7\textwidth}{X} 但願你們容忍我小小的愚蠢；請你們務必容忍我。 \end{tabularx} \\ \\ \relax
11:2 & \begin{tabularx}{0.7\textwidth}{X} 我以神嫉妒的愛來愛你們，因為我曾把你們許配給一個丈夫，要把你們如同貞潔的童女獻給基督。 \end{tabularx} \\ \\ \relax
11:3 & \begin{tabularx}{0.7\textwidth}{X} 我只怕你們的心偏邪了，失去那向基督所獻誠懇貞潔的心，就像蛇用詭詐誘惑了夏娃一樣。 \end{tabularx} \\ \\ \relax
11:4 & \begin{tabularx}{0.7\textwidth}{X} 假如有人來，傳另一個耶穌，不是我們所傳過的；或者你們另受一個靈，不是你們所受過的聖靈；或者接納另一個福音，不是你們所接納過的；你們居然容忍了！ \end{tabularx} \\ \\ \relax
11:5 & \begin{tabularx}{0.7\textwidth}{X} 但我想，我一點也不在那些超級使徒以下。 \end{tabularx} \\ \\ \relax
11:6 & \begin{tabularx}{0.7\textwidth}{X} 雖然我不擅長說話，我的知識卻不如此。這點我們已經在每一方面各樣事上向你們表明了。 \end{tabularx} \\ \\ \relax
11:7 & \begin{tabularx}{0.7\textwidth}{X} 我貶低自己，為了使你們高升，因為我白白地傳神的福音給你們，難道這算是我犯了錯嗎？ \end{tabularx} \\ \\ \relax
11:8 & \begin{tabularx}{0.7\textwidth}{X} 我剝奪了別的教會，向他們取了報酬來效勞你們。 \end{tabularx} \\ \\ \relax
11:9 & \begin{tabularx}{0.7\textwidth}{X} 我在你們那裡有缺乏的時候，並沒有連累你們一個人，因為我所缺乏的，那些從馬其頓來的弟兄都補足了。我向來凡事謹慎，將來也必謹慎，總不要連累你們。 \end{tabularx} \\ \\ \relax
11:10 & \begin{tabularx}{0.7\textwidth}{X} 既有基督的真誠在我裡面，在亞該亞一帶地方就沒有人能阻止我這樣自誇。 \end{tabularx} \\ \\ \relax
11:11 & \begin{tabularx}{0.7\textwidth}{X} 為甚麼呢？是因我不愛你們嗎？神知道，我愛你們！ \end{tabularx} \\ \\ \relax
11:12 & \begin{tabularx}{0.7\textwidth}{X} 我現在所做的，將來還要做，為要斷絕那些尋機會之人的機會，不讓他們在所誇耀的事上被人認為與我們一樣。 \end{tabularx} \\ \\ \relax
11:13 & \begin{tabularx}{0.7\textwidth}{X} 那樣的人是假使徒，行事詭詐，裝作基督的使徒。 \end{tabularx} \\ \\ \relax
11:14 & \begin{tabularx}{0.7\textwidth}{X} 這也不足為奇，因為連撒但也裝作光明的天使。 \end{tabularx} \\ \\ \relax
11:15 & \begin{tabularx}{0.7\textwidth}{X} 所以，他的差役若裝作公義的差役也沒有甚麼大不了。他們的結局必然跟他們的行為相符。 \end{tabularx} \\ \\ \relax
11:16 & \begin{tabularx}{0.7\textwidth}{X} 我再說，誰都不可把我看作愚蠢的；即使你們把我當作愚蠢人，那麼，也讓我稍微誇誇口吧。 \end{tabularx} \\ \\ \relax
11:17 & \begin{tabularx}{0.7\textwidth}{X} 我說的話不是奉主的權柄說的，而是像愚蠢人具有自信地放膽誇口。 \end{tabularx} \\ \\ \relax
11:18 & \begin{tabularx}{0.7\textwidth}{X} 既然有好些人憑著血氣在誇口，我也要誇口了。 \end{tabularx} \\ \\ \relax
11:19 & \begin{tabularx}{0.7\textwidth}{X} 你們是聰明人，竟能甘心容忍愚蠢人！ \end{tabularx} \\ \\ \relax
11:20 & \begin{tabularx}{0.7\textwidth}{X} 假若有人奴役你們，或侵吞你們，或壓榨你們，或侮辱你們，或打你們的臉，你們居然都能容忍。 \end{tabularx} \\ \\ \relax
11:21 & \begin{tabularx}{0.7\textwidth}{X} 說來慚愧，在這方面好像我們是太軟弱了。然而，我說句蠢話，人在甚麼事上敢誇口，我也敢誇口。 \end{tabularx} \\ \\ \relax
11:22 & \begin{tabularx}{0.7\textwidth}{X} 他們是希伯來人嗎？我也是。他們是以色列人嗎？我也是。他們是亞伯拉罕的後裔嗎？我也是。 \end{tabularx} \\ \\ \relax
11:23 & \begin{tabularx}{0.7\textwidth}{X} 他們是基督的用人嗎？我說句狂話，我更是。我比他們忍受更多勞苦，坐過更多次監牢，受過無數次的鞭打，常常冒死。 \end{tabularx} \\ \\ \relax
11:24 & \begin{tabularx}{0.7\textwidth}{X} 我被猶太人鞭打五次，每次四十減去一下； \end{tabularx} \\ \\ \relax
11:25 & \begin{tabularx}{0.7\textwidth}{X} 被棍打了三次，被石頭打了一次，遭海難三次，一晝一夜在深海裡掙扎。 \end{tabularx} \\ \\ \relax
11:26 & \begin{tabularx}{0.7\textwidth}{X} 我又屢次行遠路，遭江河的危險，盜賊的危險，同族人的危險，外族人的危險，城裡的危險，曠野的危險，海中的危險，假弟兄的危險。 \end{tabularx} \\ \\ \relax
11:27 & \begin{tabularx}{0.7\textwidth}{X} 我勞碌困苦，常常失眠，又飢又渴，忍飢耐寒，赤身露體。 \end{tabularx} \\ \\ \relax
11:28 & \begin{tabularx}{0.7\textwidth}{X} 除了這些外表的事以外，我還有為眾教會操心的事天天壓在我身上。 \end{tabularx} \\ \\ \relax
11:29 & \begin{tabularx}{0.7\textwidth}{X} 有誰軟弱，我不軟弱呢？有誰跌倒，我不焦急呢？ \end{tabularx} \\ \\ \relax
11:30 & \begin{tabularx}{0.7\textwidth}{X} 我若必須誇口，就誇我軟弱的事好了。 \end{tabularx} \\ \\ \relax
11:31 & \begin{tabularx}{0.7\textwidth}{X} 那永遠可稱頌之主耶穌的父神知道我不說謊。 \end{tabularx} \\ \\ \relax
11:32 & \begin{tabularx}{0.7\textwidth}{X} 在大馬士革的亞哩達王手下的提督把守大馬士革城，要捉拿我， \end{tabularx} \\ \\ \relax
11:33 & \begin{tabularx}{0.7\textwidth}{X} 我被人用筐子從城牆上的窗口縋下，逃脫了他的手。 \end{tabularx} \\ \\ \relax
12:1 & \begin{tabularx}{0.7\textwidth}{X} 雖然自誇無益，我還是不得不誇。我現在要提到主的異象和啟示。 \end{tabularx} \\ \\ \relax
12:2 & \begin{tabularx}{0.7\textwidth}{X} 我認識一個在基督裡的人，他在十四年前被提到第三層天上去；或在身內，我不知道，或在身外，我也不知道，只有神知道。 \end{tabularx} \\ \\ \relax
12:3 & \begin{tabularx}{0.7\textwidth}{X} 我認識的這樣的一個人—或在身內，或在身外，我都不知道，只有神知道— \end{tabularx} \\ \\ \relax
12:4 & \begin{tabularx}{0.7\textwidth}{X} 他被提到樂園裡，聽見隱祕的言語，是人不可說的。 \end{tabularx} \\ \\ \relax
12:5 & \begin{tabularx}{0.7\textwidth}{X} 為這人，我要誇口；但是為我自己，除了我的軟弱以外，我並不誇口。 \end{tabularx} \\ \\ \relax
12:6 & \begin{tabularx}{0.7\textwidth}{X} 就是我願意誇口也不算狂，因為我會說實話；只是我絕口不談，恐怕有人把我看得太高了，過於他在我身上所看見所聽見的； \end{tabularx} \\ \\ \relax
12:7 & \begin{tabularx}{0.7\textwidth}{X} 又恐怕我因所得的啟示太高深，就過於高抬自己，所以有一根刺加在我身上，就是撒但的差役來折磨我，免得我過於高抬自己。 \end{tabularx} \\ \\ \relax
12:8 & \begin{tabularx}{0.7\textwidth}{X} 為了這事，我曾三次求主使這根刺離開我。 \end{tabularx} \\ \\ \relax
12:9 & \begin{tabularx}{0.7\textwidth}{X} 他對我說：「我的恩典是夠你用的，因為我的能力是在人的軟弱上顯得完全。」所以，我更喜歡誇耀自己的軟弱，好使基督的能力覆庇我。 \end{tabularx} \\ \\ \relax
12:10 & \begin{tabularx}{0.7\textwidth}{X} 為基督的緣故，我以軟弱、凌辱、艱難、迫害、困苦為可喜樂的事；因為我甚麼時候軟弱，甚麼時候就剛強了。 \end{tabularx} \\ \\ \relax
12:11 & \begin{tabularx}{0.7\textwidth}{X} 我成了愚蠢人，是被你們逼出來的，因為我本該被你們讚許才是。雖然我算不了甚麼，卻沒有一件事在那些超級使徒以下。 \end{tabularx} \\ \\ \relax
12:12 & \begin{tabularx}{0.7\textwidth}{X} 我在你們中間，用百般的忍耐，藉著神蹟、奇事、異能顯出使徒的憑據來。 \end{tabularx} \\ \\ \relax
12:13 & \begin{tabularx}{0.7\textwidth}{X} 除了我不曾連累你們這一件事，你們還有甚麼事不及別的教會呢？這不公平之處，請你們饒恕我吧。 \end{tabularx} \\ \\ \relax
12:14 & \begin{tabularx}{0.7\textwidth}{X} 如今，我準備第三次到你們那裡去。我仍不會連累你們，因為我所求的是你們，不是你們的財物。兒女不該為父母積財，父母該為兒女積財。 \end{tabularx} \\ \\ \relax
12:15 & \begin{tabularx}{0.7\textwidth}{X} 我也甘心樂意為你們的靈魂費財費力。難道我越愛你們，就越少得你們的愛嗎？ \end{tabularx} \\ \\ \relax
12:16 & \begin{tabularx}{0.7\textwidth}{X} 罷了，我自己並沒有連累你們，你們卻有人說，我施詭詐，用心計牢籠你們。 \end{tabularx} \\ \\ \relax
12:17 & \begin{tabularx}{0.7\textwidth}{X} 我所差遣到你們那裡去的人，我何曾藉著他們中的任何人佔過你們的便宜呢？ \end{tabularx} \\ \\ \relax
12:18 & \begin{tabularx}{0.7\textwidth}{X} 我勸提多到你們那裡去，又差遣那位弟兄與他同去，提多佔過你們的便宜嗎？我們的行事為人不是同一心靈嗎？不是同一步伐嗎？ \end{tabularx} \\ \\ \relax
12:19 & \begin{tabularx}{0.7\textwidth}{X} 你們一直認為我們是在你們面前為自己辯護嗎？其實，我們本是在基督裡當著神面前說話。親愛的，一切的事都是為了造就你們。 \end{tabularx} \\ \\ \relax
12:20 & \begin{tabularx}{0.7\textwidth}{X} 我怕我再來的時候，見你們不合我所期望的，而你們見我也不合你們所期望的。我怕有紛爭、嫉妒、憤怒、自私、毀謗、讒言、狂傲、動亂的事。 \end{tabularx} \\ \\ \relax
12:21 & \begin{tabularx}{0.7\textwidth}{X} 我怕我再來的時候，我的神使我在你們面前蒙羞，並且又因許多人從前犯罪，行污穢、淫亂、放蕩的事，不肯悔改而悲傷。 \end{tabularx} \\ \\ \relax
13:1 & \begin{tabularx}{0.7\textwidth}{X} 這是我第三次要到你們那裡去。「任何指控都要憑兩個或三個證人的口述才能成立。」 \end{tabularx} \\ \\ \relax
13:2 & \begin{tabularx}{0.7\textwidth}{X} 對那些犯了罪的人和其餘所有的人，正如我第二次見你們的時候曾說過，現在不在你們那裡再次說：「我若再來，必不寬容。」 \end{tabularx} \\ \\ \relax
13:3 & \begin{tabularx}{0.7\textwidth}{X} 因為你們想求證基督是否藉著我說話。基督對你們並不是軟弱的，而是在你們裡面大有能力的。 \end{tabularx} \\ \\ \relax
13:4 & \begin{tabularx}{0.7\textwidth}{X} 他因軟弱被釘在十字架上，卻因神的大能仍然活著。我們在他裡面也成為軟弱的，但對你們，我們將因神的大能而與他一同活著。 \end{tabularx} \\ \\ \relax
13:5 & \begin{tabularx}{0.7\textwidth}{X} 你們總要省察自己是否在信仰中生活；你們要考驗自己。除非你們經不起考驗，你們自己豈不應該知道有耶穌基督在你們裡面嗎？ \end{tabularx} \\ \\ \relax
13:6 & \begin{tabularx}{0.7\textwidth}{X} 我希望你們知道，我們並不是經不起考驗的人。 \end{tabularx} \\ \\ \relax
13:7 & \begin{tabularx}{0.7\textwidth}{X} 我們祈求神使你們不做任何惡事；這不是要顯明我們是經得起考驗的，而是要你們行事端正，即使我們似乎經不起考驗也沒有關係。 \end{tabularx} \\ \\ \relax
13:8 & \begin{tabularx}{0.7\textwidth}{X} 我們不能做任何對抗真理的事，只能維護真理。 \end{tabularx} \\ \\ \relax
13:9 & \begin{tabularx}{0.7\textwidth}{X} 當我們軟弱而你們剛強時，我們也歡喜。我們所祈求的是：你們能成為完全人。 \end{tabularx} \\ \\ \relax
13:10 & \begin{tabularx}{0.7\textwidth}{X} 所以，我不在你們那裡的時候，把這些話寫給你們，好使我見你們的時候不用照主所給我的權柄嚴厲地待你們；這權柄原是為造就人，而不是為摧毀人。 \end{tabularx} \\ \\ \relax
13:11 & \begin{tabularx}{0.7\textwidth}{X} 末了，弟兄們，願你們喜樂。要追求完全；要接受鼓勵；要同心合意；要彼此和睦。如此，慈愛和平的神必與你們同在。 \end{tabularx} \\ \\ \relax
13:12 & \begin{tabularx}{0.7\textwidth}{X} 你們要用聖潔的吻彼此問安。 \end{tabularx} \\ \\ \relax
13:13 & \begin{tabularx}{0.7\textwidth}{X} 眾聖徒都向你們問安。 \end{tabularx} \\ \\ \relax
13:14 & \begin{tabularx}{0.7\textwidth}{X} 願主耶穌基督的恩惠、神的慈愛、聖靈的感動常與你們眾人同在！ \end{tabularx} \\ \\
[1ex]
\hline
\hline
\end{longtable}
$^{1}$各位電視節目平安.
我們繼續來看哥林多後書.
這次已經來到第六講了.
就由哥林多後書第十一章.
一路去到十三章.
就完結了哥林多後書.
保羅他以受患難為榮.
哇 有很多人說.
你誇的就是你的成就.
哇 你看我去到哪裡宣教.
我建立了多少間教會.
我做了什麼 做了什麼.
但是保羅不是.
他說我以為基督受苦為樂為榮.
一個很謙卑的僕人.
原來保羅因著基督的緣故.
去受了很多很多的苦.
而且他就勸勉弟兄姊妹.
他說你們不要失去了單純純一的信心.
你們最重要就是堅持.
你們的信心是足夠的.
信就是所望之事的實體.
是未見之事的一個確距.
我們的信心必須要足夠.
希伯來書第十一章叫做《信心的篇章》.
作者數了很多信心的偉人.
這些偉人走在我們前面.
就給我們看見一些榜樣.
我們要學他們這樣.
有純一的信心.
保羅講完信心之後.
他就以自己受過什麼的患難.
受過什麼的苦來數出來.
如果你看十一章.
由第23節起.
你就看見保羅這份單真的很厲害.
他說我比你們受苦是多的.
我多受勞苦.
多下監牢.
而且我受鞭打是過重的.

$^{41}$而且我時時在患難困所當中.
猶太人鞭打過我五次.
我每次是40減去一下.
就是39下.
他說我也是船沉了多少次.
還有就是我有一次.
是遇船壞三次.
一晝一夜在水裡面.
又屢次行路遭遇光河的危險.
盜賊的危險和其他外族人的危險.
城裡面的危險.
各樣的危險.
所以保羅就在這裡一路數一路數.
他是經過什麼危險.
他說在大患難裡面主救了他.
而且主是幫助他.
將他從敵人的手裡面救拔出來.
我們都知道保羅的經歷.
真是非常之豐富.
他就在這裡誇自己.
原來我受了這麼多的患難.
他以受患難為主的名受患難而誇口.
各位我們有沒有這樣做過呢?.
可能我們一見到患難.
我們就走甲不抖.
甚至我們會不認主.
因為以免我們避開那些的患難.
求主憐憫我們.
幫助我們.
讓我們光強.
好像保羅一樣.
然後去到十二章.
保羅就講到他所得到的啟示.
他說由於可能我所得的啟示甚大.
神就將一根刺加在我的身上.
我三次求過主.
可不可以拿去這根刺.
但是主都沒有答應.
保羅沒有提到那根刺到底是什麼.
於是很多佳經家就在猜.

$^{81}$到底保羅身上那根刺是什麼.
有很多人在猜.
其實就是他有眼疾.
所以有時在書信裡面.
他說你看看我寫的字多麼大.
因為眼睛不好.
所以寫的字是很大隻的.
他求主.
但是主沒有來.
主怎樣跟他說呢?.
主說我的恩典是夠你用的.
沒錯.
你有苦楚.
你有患難.
但是我的恩典更加夠用.
所以保羅他說.
我誇的就是誇自己的軟弱.
我什麼時候軟弱.
我什麼時候就光強.
保羅的態度真是令我們很得幫助.
一個這樣神的僕人.
他只誇自己的軟弱.
他從來沒有為自己誇口.
原來神是要在他的身上.
顯明自己的恩典夠保羅用.
他說免得我誇口.
「原來神要保羅你不要牙刷刷」.
因為我的恩典是承托著你的.
他說為基督的緣故.
我就以軟弱,急難,患難為可喜樂.
因為我什麼時候軟弱.
我什麼時候就光強.
保羅那種態度他承受苦難.
但是神給他一份的恩力.
可以托得住這些身體上的苦痛.
各樣在外面來的患難.
而叫他能夠為主做美好的見證.
有力的見證.
我們今天都可以成為耶穌基督有力的見證人.
所以保羅就是我們的榜樣.

$^{121}$他自己經過很多種各樣的苦楚.
但是在他身上仍然以主為樂.
顯出基督的向氣.
好 去到第十三章.
就是哥林多後書最後一章.
保羅說我第三次要去你們那裡.
保羅真的很愛哥林多這間教會.
所以很多學者都說.
保羅其實寫了三封信.
第一封信就是不見了的那封.
所以是一封未見過的信.
我們不知道的信.
然後第二封就是哥林多前書.
然後第三封就是哥林多後書.
然後保羅說我第三次要去你們那裡.
所以對他們有勸勉.
然後去到十三章的十一節.
保羅說我保羅沒了的話.
就要你們兄弟們都起來.
要做完全人.
要受安慰.
要同心合意.
要彼此和睦.
如此仁愛和平的神必常常與你們同在.
保羅語重心長地勸這班弟兄姊妹.
他說你們要起來.
要做一個完全人.
即是在行為上要有美好的見證.
要完完全全.
達到一個很高的標準.
我們千萬不要定一個很低的標準.
我們的標準要拉到很高的.
哇!這麼高的標準.
我努力去做.
我嘗試去做.
我靠著主的恩典我能夠做得到.
所以各位我們都要做完全人.
要起來.
要同心合意.
然後彼此和睦.

$^{161}$不要爭吵.
要和睦.
然後這樣仁愛和平喜樂的神.
就必在我們中間.
叫我們可以見證祂.
得著祂的恩典.
得著祂的祈福.
我們的神是一位祝福我們.
帶領我們的神.
所以保羅在這本書結束的時候.
是一個祝福.
是對他們一個祝福.
求主也幫助我們今天.
努力做一個完全人.
在我們的生命裡面.
我們都活出主的樣式.
見證我們主的榮美.
叫榮耀歸給我們天上的神.
願神賜福給各位.
\newpage



\section{以弗所書 4:14-16}
\label{sec:031692}
\textbf{建立健康的使命教會 (5) 以真理和恩典建立關係}~~~~鍾韵悠牧師
\newline
\newline
連結: \href{http://www.church.com.hk/acms/content.asp?site=cdc\&op=show\&type=product\&code=031692\&layout=sermon}{\texttt{031692}} ~~~~ 語音日期: 2026-02-22
\newline
\newline
\hyperref[sec:031691]{\small{< < < PREV SERMON < < <}}
~
\hyperref[sec:index_chronic]{\small{[返順時目]}}
~
\hyperref[sec:031693]{\small{> > > NEXT SERMON > > >}}
\newline
\newline
以弗所書 4:14-16
\newline
\begin{longtable}{cl}
\hline
\hline
章節 & 經文 (和合本修訂版)\\
\hline
4:14 & \begin{tabularx}{0.7\textwidth}{X} 這樣，我們不再作小孩子，中了人的詭計和欺騙的法術，被一切邪說之風搖動，飄來飄去。 \end{tabularx} \\ \\ \relax
4:15 & \begin{tabularx}{0.7\textwidth}{X} 我們反而要用愛心說誠實話，各方面向著基督長進，連於元首基督， \end{tabularx} \\ \\ \relax
4:16 & \begin{tabularx}{0.7\textwidth}{X} 靠著他全身都連接得緊湊，百節各按各職，照著各體的功用彼此相助，使身體漸漸增長，在愛中建立自己。 \end{tabularx} \\ \\
[1ex]
\hline
\hline
\end{longtable}
$^{1}$弟兄姊妹 早安.
跟大家拜個年.
祝大家主恩滿日.
龍馬精神.
龍馬精神 因為我家有一條龍和一隻馬.
所以我經常跟家人拜年都是龍馬精神.
這樣.
明天是人日 是不是?.
初七 是吧?.
所以祝大家有個平安的人日.
人人生日.
也是紀念天父給我們很多恩典的時候.
每天都是.
我先戴個頭盔.
我是一個科學不合格的人.
中三的時候要編去理科班或者文科班.
我freecam拜都是不合格的.
所以很自然就去了文科班.
但是我想說一說.
有一個很有趣的化學現象.
想跟大家分享一下.
這個化學現象我也是看回來的.
在Christianity Today.
今天基督教.
有一篇文章寫.
跳禍的鹽 seasoned salt.
他說有一種叫做鈉.
鈉就是sodium.
大家都知道.
很活躍的.
如果你是懂science.
chemistry那些.
不要挑戰我.
我說了我不合格.
sodium很活躍的.
在自然界裡面幾乎不可以單獨存在.
總是要跟其他的元素結合.
另一種元素就是chlorine.
鹿.
是有毒的氣體.

$^{41}$就好像在漂白水裡面.
你一聞聞那種很刺鼻的成分就是鹿.
但原來當很奇妙的.
當鈉和鹿結合在一起.
就變成了sodium chloride.
就是鹿化鈉.
也就是我們現在所稱為的食鹽.
沒有毒之餘.
還可以吃進肚子裡.
這個chemistry不合格的人突然覺得.
嘩!好奇妙啊.
當然大家中學有讀的時候.
大家都知道.
我其實不知道.
這個鹽還可以保存食物.
還可以去調味.
啊!真的太貼切了.
跟我們今天那一段的經文.
愛和真理其實有一點點像.
不要完全.
不要對號入座.
有一點點像鈉和鹿.
因為只是愛沒有真理.
愛就會變得很飄忽.
有時麻木.
什麼都要接受.
什麼都要混在一起.
但是只是有真理沒有愛.
嘩!就很功利.
很刺耳的.
帶來了傷害性.
所以如果沒有用愛心說出來的話.
就會reject人家.
推走人家.
甚至推離福音.
如果當真理和愛在一個人.
和一個教會裡面結合在一起.
真的會產生一個很漂亮很美的生命力.
所以主耶穌就這樣說.
你們要作世上的.

$^{81}$sodium chloride.
食鹽.
不是.
聖經是說鹽.
所以二弗所書第四章裡面.
保羅正正告訴我們.
一個成熟的教會.
就是要在真理和愛.
那裡一起去長大.
我們今天是建立健康的使命教會.
叫做第二部分.
我們之前說了活出敬拜人生.
今天我們會開始說有質素的相交.
group.
我們就看看保羅.
特別在第四章.
這兩至三次經文.
怎樣教導我們要成為成熟成長的一個教會.
這個基督的身體.
讓我們一起祈禱.
父啊!多謝你.
我們在新年裡面再次感恩.
多謝你給我們有健康.
有生命氣息.
每一天都是感恩.
我們可以祈禱.
我們可以坐在這裡.
我們可以有弟兄姊妹.
有你教會.
又是感恩.
我們有疾病.
我們有苦難.
但是有主你同在.
帶領我們經過.
幫助我們明白.
又是感恩.
求福你要選我們的國藩.
打開我們的耳朵.
眼睛.
心.

$^{121}$去聽你的話語.
多謝你.
奉主耶穌基督聖名祈求.
Amen.
(主持人:我們應該怎樣去活出來呢?).
(主持人:怎樣去與蒙召的因去相稱呢?).
(主持人:怎樣教會和生活真的去合一呢?).
(主持人:教會生活怎樣去合一呢?).
(主持人:和新生命倫理還有屬靈的挑戰等等).
(主持人:所以說如果第四章呢).
(主持人:我們在二弗所書可以說是一個轉捩點).
(主持人:從一個身份認知).
(主持人:知道自己是誰).
(主持人:到活出身份).
(主持人:到我是不是真真正正).
(主持人:活出我知道我是).
(主持人:神自己兒女那個的身份出來).
(主持人:在救恩和選選).
(主持人:神的選選教會得合一).
(主持人:我們就可以在合一中成長成熟).
(主持人:這個完全是一個).
(主持人:很長的旅程來的).
(主持人:但是也是主耶穌給我們一個福分).
(主持人:更加是我們有責任).
(主持人:或者一個獻情).
(主持人:一個offering).
(主持人:我們看看保羅).
(主持人:在這三節裡面跟我們說什麼說話).
(主持人:他說成長成熟).
(主持人:我們要在信仰上面長大).
(主持人:他說不要再作小孩子了).
(主持人:中了人的詭計和欺騙的法術).
(主持人:被一切邪說的風搖動飄來飄去).
(主持人:這節反映了).
(主持人:其實以弗所教會當時的處境).
(主持人:是很多異端邪說).
(主持人:很多錯誤的教導).
(主持人:例如洛斯底主義).
(主持人:猶太法思想).
(主持人:哲學 神秘主義的混淆).

$^{161}$(主持人:還有假教師的問題).
(主持人:所有都可以讓信徒誤入歧途).
(主持人:保羅說).
(主持人:如果你們聽他又被他動).
(主持人:是不成熟的表現).
(主持人:被各樣的風搖動).
(主持人:被人的詭計欺騙).
(主持人:和容易受錯誤教導影響).
(主持人:經文用「風」你就知道).
(主持人:因為風本身其實又看不到).
(主持人:又變幻莫測).
(主持人:天天都要看著天氣報告去穿衣服).
(主持人:有沒有風呢?).
(主持人:我去動物園都要看看).
(主持人:那個wind gust是多少).
(主持人:那個wind是多少).
(主持人:你說暖而已).
(主持人:但那個wind好勁的時候).
(主持人:我都不可以出去走那麼久).
(主持人:要看是吧).
(主持人:但是不知道的).
(主持人:我們廣東人很喜歡說).
(主持人:隨波逐流).
(主持人:跟著浪飄來飄去).
(主持人:就好像不穩定和很脆弱).
(主持人:當保羅說我們誤作小孩子的時候).
(主持人:其實這個小孩子不單單是幾歲大).
(主持人:英文是infant 嬰孩).
(主持人:當我們用嬰孩去說的時候).
(主持人:就知道是一個完全未成熟 剛出生).
(主持人:缺乏分辨力 缺乏理解力).
(主持人:所以希伯來書也說).
(主持人:他將屬靈的奶和乾糧去對比).
(主持人:因為baby才喝奶).
(主持人:大人也喝奶 我也喜歡喝奶).
(主持人:喝2\%好了 不要喝太多).
(主持人:下茶可以下5\%或者10\%).
(主持人:但不單單是喝奶去生存).
(主持人:嬰孩半歲).
(主持人:correct me if I am wrong).

$^{201}$(主持人:半歲就開始吃糊仔).
(主持人:4月半歲 差不多).
(主持人:要吃糊仔 要吃硬東西).
(主持人:不能保留30歲只喝奶 是不行的).
(主持人:這就是保羅圈面).
(主持人:我們不要只做嬰孩).
(主持人:我們要在信仰上成長).
(主持人:我們要走向成熟).
(主持人:這就是整個主題).
(主持人:所以當我們看到嬰孩).
(主持人:怎樣成長的時候).
(主持人:很有屬靈意義的).
(主持人:因為小朋友通常).
(主持人:好像還沒有很了解).
(主持人:我們說老順都還沒有生).
(主持人:其實二十多歲都可以).
(主持人:老順都還沒有生).
(主持人:不過往往會形容).
(主持人:小朋友會衝動行事).
(主持人:會受情緒影響).
(主持人:也不會經過想就說的).
(主持人:難道我想說句話會不會傷害到人).
(主持人:不會這樣的 是不是 小朋友).
(主持人:而且也未必很有自制能力).
(主持人:我這麼大的小朋友).
(主持人:也未必很有自制能力).
(主持人:自律也要繼續學習).
(主持人:所以父母或者長輩).
(主持人:我們通常跟小朋友說).
(主持人:我們都會說).
(主持人:小心啊).
(主持人:第一句就是小心啊).
(主持人:什麼都要小心).
(主持人:小心著涼 小心跌倒).
(主持人:小心看東西).
(主持人:小心不要看不對的東西).
(主持人:小心陌生人).
(主持人:很多小心).
(主持人:其實保羅的心腸就是).
(主持人:小心啊).

$^{241}$(主持人:不要再做小朋友了).
(主持人:不要在信仰裡再做嬰兒).
(主持人:提防陌生人).
(主持人:Say No).
(主持人:懂得Say No).
(主持人:在我們今天裡).
(主持人:我們這些會不會是).
(主持人:上YouTube很多的神學短片).
(主持人:很多情緒化靈修內容).
(主持人:很多資訊).
(主持人:但是我們會不會分辨).
(主持人:哪些是真理呢).
(主持人:我們會不會好像).
(主持人:熟還是不成熟).
(主持人:容易受情緒欺騙呢).
(主持人:古羅叫我們).
(主持人:不要走向其他的福音).
(主持人:專注堅持在基督真理裡面).
(主持人:就好像一艘船一樣).
(主持人:我看看有沒有).
(主持人:有的).
(主持人:check GPT).
(主持人:一艘船遇了暴風).
(主持人:但是無論船有多漂亮).
(主持人:其實都要看撈夠不夠穩).
(主持人:真理就是我們的撈).
(主持人:如果沒有紮根在基督裡面).
(主持人:環境一變).
(主持人:我們就飛走了).
(主持人:浪就衝走了).
(主持人:所以我們的問題).
(主持人:究竟我們知不知道自己信什麼).
(主持人:更甚的).
(主持人:我們知不知道自己).
(主持人:為什麼去相信我們所相信的事).
(主持人:Do we know what we believe).
(主持人:Do we know why we believe what we believe).
(主持人:我覺得這兩個問題).
(主持人:特別是去引導我).
(主持人:幫助我去).

$^{281}$(主持人:不單去看自己的信仰之餘).
(主持人:還是怎樣去幫我們的下一代去).
(主持人:Do you know what you believe).
(主持人:Do you know why you believe what you believe).
(主持人:所以保羅說).
(主持人:不要受到).
(主持人:好像從孩童那樣).
(主持人:要整天被受騙).
(主持人:欺騙).
(主持人:情緒影響).
(主持人:沒有節制).
(主持人:我們要在信仰上面成熟).
(主持人:都是回到最基礎).
(主持人:就是看聖經).
(主持人:查聖經).
(主持人:和弟兄姊妹一起研讀聖經).
(主持人:不要怕那些很複雜很難的).
(主持人:就skip了).
(主持人:很難不明白就skip了).
(主持人:不要緊的).
(主持人:上主約學).
(主持人:問Pastor Leo).
(主持人:問其他).
(主持人:我們去看《釋經書》).
(主持人:看有很優質的屬靈書籍).
(主持人:文章).
(主持人:聽podcast).
(主持人:看優質的影片).
(主持人:當然都要很分辨的).
(主持人:還有常在主耶穌裡面遵守祂的命令).
(主持人:因為祂就是道路真理生命).
(主持人:我經常用這個例子).
(主持人:就是我的).
(主持人:誰呢?).
(主持人:我大學的導師).
(主持人:大學團體的導師).
(主持人:他的妹妹初初進去銀行的時候).
(主持人:就有兩個禮拜什麼都不用做).
(主持人:只是數銀紙).
(主持人:數兩個禮拜的銀紙).

$^{321}$(主持人:當然是很久以前了).
(主持人:又沒有什麼機器).
(主持人:數數數數到熟悉).
(主持人:真銀紙是怎樣的).
(主持人:感覺).
(主持人:怎樣的看 怎樣的樣子).
(主持人:要很熟悉).
(主持人:因為假的東西實在太過厲害了).
(主持人:越來越做得不同).
(主持人:不認識真的).
(主持人:怎樣可以知道哪些是假的呢).
(主持人:保羅說).
(主持人:除了在信仰中要成長).
(主持人:他說在愛中要持守真理).
(主持人:今天說的).
(主持人:有質素的商家裡面第一課).
(主持人:所以這個篇幅可能會比較長一點).
(主持人:他說我們反而要用愛心去說誠實話).
(主持人:但一看「反而」這個字).
(主持人:Instead).
(主持人:他說Instead).
(主持人:要用愛心說誠實話).
(主持人:很鮮明的一個對比).
(主持人:他說信徒不會隨波逐流).
(主持人:是會選擇一條經過深思熟慮).
(主持人:即是有想過不是小朋友).
(主持人:亦無神影響的路).
(主持人:永遠讓自己在聖言當中).
(主持人:有了信仰有長大).
(主持人:吃乾糧會成長).
(主持人:才知道甚麼是真理).
(主持人:如果我們不讀).
(主持人:我們又怎知道甚麼是誠實話).
(主持人:甚麼是真理呢).
(主持人:在這個Amplified Bible).
(主持人:他說).
(主持人:Speaking the truth in love).
(主持人:即是誠實話).
(主持人:可能在中國人或廣東話).
(主持人:即是說真話).

$^{361}$(主持人:但其實這個Truth).
(主持人:是真理的意思).
(主持人:In all things).
(主持人:Both our speech and our lives).
(主持人:Expressing his truth).
(主持人:即是主耶穌的真理).
(主持人:基督的真理).
(主持人:聖經的真理).
(主持人:所以我們說).
(主持人:在愛裡算誠實話).
(主持人:包括在一切事上).
(主持人:我們都是言語生活).
(主持人:都是表達出神的真理來).
(主持人:這個真理).
(主持人:就不單單是).
(主持人:是不是我覺得對的東西).
(主持人:我就說).
(主持人:而是神的真理).
(主持人:而這個愛).
(主持人:究竟是什麼愛呢).
(主持人:是不是一味說).
(主持人:我是愛你).
(主持人:我為你好才這樣說的).
(主持人:這個我為你好).
(主持人:真的帶給很多人的傷害的).
(主持人:其實).
(主持人:大家都會).
(主持人:心裡面有所知道).
(主持人:我都是為你好).
(主持人:這樣跟你說).
(主持人:是不是有時).
(主持人:其實這個愛是什麼愛呢).
(主持人:我們最終是要回去).
(主持人:聖經那裡看).
(主持人:這個是什麼愛呢).
(主持人:這個是神的愛).
(主持人:神的愛是什麼).
(主持人:愛是恆久忍耐).
(主持人:又有).
(主持人:因此愛是).

$^{401}$(主持人:不嫉妒).
(主持人:不自誇).
(主持人:不張狂).
(主持人:不作害羞的事).
(主持人:不求自己的益處).
(主持人:不輕易發怒).
(主持人:不計算人的惡).
(主持人:不喜歡).
(主持人:不疑只喜歡真理凡事).
(主持人:包容凡事).
(主持人:相信凡事).
(主持人:盼望凡事).
(主持人:忍耐愛是永不止息).
(主持人:在心裡中讀).
(主持人:愛是恆久忍耐又有因此).
(主持人:究竟就會看到).
(主持人:究竟我的愛是不是神).
(主持人:叫我們這樣去愛).
(主持人:這個很大的反省).
(主持人:究竟我的愛是不是愛神愛人).
(主持人:和敬畏神的愛).
(主持人:我說為你好那些).
(主持人:其實是否自己想出來的).
(主持人:誠實說).
(主持人:I thought it's the right thing).
(主持人:但是不是真是聖經的真理呢).
(主持人:所以用愛心說誠實說這句說話).
(主持人:是很有深度).
(主持人:我們用愛心說誠實說).
(主持人:就不像嬰孩一般或者小朋友一般).
(主持人:會受情緒阻攔).
(主持人:或者要經過深思熟慮).
(主持人:基於基督真理裡面).
(主持人:我的一個決定).
(主持人:我要用愛心說誠實說).
(主持人:同時要為我們基督身體裡面).
(主持人:各個肢體一個愛的衛生).
(主持人:這個絕對不是說).
(主持人:只是說你做得不好).
(主持人:我愛你所以我和你這樣說).

$^{441}$(主持人:不是這樣 是一個衛生).
(主持人:我因為神愛我).
(主持人:所以我用敬畏神的愛去愛你).
(主持人:和一起去說真理).
(主持人:不代表自己懂得).
(主持人:也不代表自己沒有罪).
(主持人:所以這是教會成熟的關鍵).
(主持人:但也有張力).
(主持人:我們教會裡面不會缺乏聚會).
(主持人:也不會缺乏活動).
(主持人:不過有個問題我們要問).
(主持人:就是我們的相交有沒有質素).
(主持人:《聖經》上面「相交」這個字).
(主持人:是描述一個有深度 有方向 有生命力的相交).
(主持人:有質素的相交不只是去相處).
(主持人:是生命的相遇).
(主持人:不只是相處 是一個相遇).
(主持人:看飯 聊天 關心近況).
(主持人:很好的 不過我們不停留在那裡).
(主持人:我們會入得更深 入到生命).
(主持人:觸及一些我們感到很不舒服的東西).
(主持人:但正正因為這樣 我們的生命關係可以得到改變).
(主持人:不知道大家有沒有聽過這本書叫做《過猶不及》).
(主持人:其實左邊那本書是《過猶不及》).
(主持人:《Boundaries》這本書的新版).
(主持人:所以新版改名為《立戒線得自由》).
(主持人:有沒有聽過《過猶不及》《Boundaries》這本書).
(主持人:作者叫做Henry Cloud).
(主持人:Henry Cloud跟另一個作者作了一本書叫做《How People Grow》).
(主持人:《成長神學》).
(主持人:裡面說到真正的生命改變).
(主持人:是需要有恩典 有真理 還要有時間).
(主持人:恩典就是上帝和祂是誰 祂做了什麼事).
(主持人:祂給我們的恩惠 無條件接納我們的愛).
(主持人:這是恩典 是神給我們的禮物 給我們有安全感).
(主持人:真理就是神是誰 神做了什麼).
(主持人:讓我們好認識祂 也從而認識我們自己真實的情況).
(主持人:我們現在誠實面對不掩飾的情況).
(主持人:除了祂是誰 祂做了什麼).
(主持人:還有我真正是誰).

$^{481}$(主持人:而時間當然就是 要時間去浸淫).
(主持人:所以Henry Cloud說).
(主持人:如果只是有恩典 沒有真理).
(主持人:關係會很暖 很開心).
(主持人:見到大家會摟著你 好嗎).
(主持人:如果只是有恩典 沒有真理).
(主持人:就會變成溫柔 但是不改變).
(主持人:如果群體裡只有接納 只有包容).
(主持人:你OK的 OK的 做得很好).
(主持人:這些鼓勵的話當然是好).
(主持人:但是沒有進一步的時候).
(主持人:可能我們的生命就不會真正被更新).
(主持人:因為我們會感受到被愛).
(主持人:但是我們不會被挑戰).
(主持人:我們會感到很安慰).
(主持人:但是我們不會被醫治).
(主持人:我們會感到被接納).
(主持人:但是不一定會有改變).
(主持人:而這個愛不是主耶穌給我們的愛).
(主持人:如果只是有真理 沒有恩典又怎樣呢).
(主持人:其實就是sodium和chloride分開的時候).
(主持人:關係是很正確的).
(主持人:但是正確傷人).
(主持人:不是正確businessman).
(主持人:是傷害了人 傷人).
(主持人:相反 一個群體只是有標準 有對錯).
(主持人:只是指出問題 沒有接納 沒有憐憫 沒有耐性).
(主持人:結果會是怎樣呢).
(主持人:I feel so guilty and shameful).
(主持人:很羞恥 很羞愧 又得不到醫治).
(主持人:我會防衛 我不會躺開).
(主持人:我會離開 甚至沒有成長).
(主持人:所以真理沒有了愛 就會變成壓傷人的重擔).
(主持人:而不是the truth will set you free那種光).
(主持人:the truth will set you free 那種光).
(主持人:不只是有真理沒有恩典).
(主持人:我在想 其實做醫生 做老師).
(主持人:有時都突然想起).
(主持人:其實醫生是做什麼工作的).
(主持人:醫).

$^{521}$(主持人:對了 兩個 又醫病又醫人).
(主持人:又醫人).
(主持人:又醫病又醫人).
(主持人:很難很難).
(主持人:又醫病又醫人).
(主持人:是不是).
(主持人:如果你去看醫生).
(主持人:他只是告訴你).
(主持人:你OK的 沒事的 你回去).
(主持人:但其實你很大件事).
(主持人:即是掩飾不告訴你).
(主持人:what's going on的時候).
(主持人:得不到醫治).
(主持人:但如果他只是說).
(主持人:行了 快吃藥吧).
(主持人:你的情況很差 非常差).
(主持人:嘩 那你又突了一個突).
(主持人:不是吧).
(主持人:所以醫生是醫人 教師又如何).
(主持人:教師做什麼的).
(主持人:教學生 教書 亦都是教人).
(主持人:即是你教書本的知識上亦都是教做人).
(主持人:其實我在中學的時候喜歡).
(主持人:我或許都有分享過).
(主持人:所以如果聽過的).
(主持人:你又是當自己沒聽過).
(主持人:有位叫李老師 我喜歡的時候).
(主持人:他進來的時候).
(主持人:個個都害怕他).
(主持人:有個外號叫做女殺手).
(主持人:他進來的時候).
(主持人:他的眼睛已經說).
(主持人:我的眼睛可以看穿你們的心).
(主持人:你們不可以走出我的五指山).
(主持人:然後個個都很安靜 很害怕他).
(主持人:他的同學很頑皮 說話).
(主持人:佛抄).
(主持人:佛抄就不是我不在堂上聊天 不是的).
(主持人:他叫我抄一句).
(主持人:多言多語難免有過 禁止嘴唇時有自慰).

$^{561}$(主持人:中意不知是甚麼).
(主持人:嘩 這麼長).
(主持人:多言多語難免有過 禁止嘴唇時有自慰).
(主持人:你看我記到今天).
(主持人:不是我抄 不是我抄).
(主持人:但記到今天 你都知道有多深刻).
(主持人:原來是真言十章十九節 原來是聖經).
(主持人:其後我才知道這位老師是基督徒).
(主持人:嘩 五指山).
(主持人:嘩 吻過都聽到那些).
(主持人:他教書很好).
(主持人:但他經常叫你起身回答 我回答不了).
(主持人:因為沒有溫書 很害怕).
(主持人:但原來他是基督徒 很關心學生).
(主持人:沒有被他教過).
(主持人:黑面神 女殺手).
(主持人:有被他教過).
(主持人:在我大學二年級的時候).
(主持人:因為鼻咽癌回到天父那裡).
(主持人:嘩 感到很難過).
(主持人:因為這位李老師實在是).
(主持人:教書 教人 有真理).
(主持人:但有很多的恩典).
(主持人:因為和他接觸得多 就會知道).
(主持人:我在教師大學).
(主持人:加拿大一個師範讀一年的時候).
(主持人:老師都教了我們一樣東西).
(主持人:當我們和學生說話的時候).
(主持人:我們最好三文治).
(主持人:三文治?).
(主持人:有麵包 肉 麵包).
(主持人:麵包就是鼓勵).
(主持人:肉才說一下有什麼進步).
(主持人:What to improve).
(主持人:但要用麵包夾住它).
(主持人:這個形象令我記憶猶新).
(主持人:所以成熟的信仰群體).
(主持人:有恩典 有真理).
(主持人:而聖經裡面最好的人版是誰?).
(主持人:當然是主耶穌自己).

$^{601}$(主持人:道成了肉身 住在我們中間).
(主持人:充充滿滿有恩典 有真理).
(主持人:這個其實就是主耶穌給我們一個模範).
(主持人:我們怎樣去跟隨祂).
(主持人:祂的模式就是恩典 真理一起來).
(主持人:祂從來不是非黑即白的極端).
(主持人:祂如何運作最好的例子就是).
(主持人:上一年40天去一個靜修營的時候).
(主持人:其中一篇祈禱的經文).
(主持人:就是《約翰福音》第八章).
(主持人:那個婦人因為她犯姦淫的時候被抓到).
(主持人:文氏和花里賽人就帶她去主耶穌面前).
(主持人:他們都很想試探主耶穌).
(主持人:他們就說:耶穌 怎樣?).
(主持人:在《摩西律法書》裡面說的).
(主持人:這樣的女人用石頭打死祂).
(主持人:你說怎樣?).
(主持人:主耶穌見到有很多其他百姓).
(主持人:有他們的民事 有這位婦人在那裡).
(主持人:祂就彎下身 寫著寫著).
(主持人:那些人又繼續問祂怎樣怎樣).
(主持人:主耶穌就起來說).
(主持人:如果你們中間誰沒有罪).
(主持人:你們就先用石頭打他).
(主持人:接著祂又彎下身繼續寫).
(主持人:過了一會兒).
(主持人:慢慢一個一個由老至幼慢慢走開).
(主持人:我在想 老至幼當然是因為).
(主持人:年紀大一點 可能知道罪孽深重一點).
(主持人:小朋友未知自己 其實罪孽也深重).
(主持人:就剩下主耶穌和那個婦人站在一起).
(主持人:其實當我在這裡祈禱的時候).
(主持人:我心裡很難堪 很embarrassing).
(主持人:很humiliating).
(主持人:多麼的屈辱).
(主持人:做錯事還公然在那裡).
(主持人:我第一句就是Where is the man?).
(主持人:那個男人去了哪裡呢?).
(主持人:這個犯罪不是這個女人).
(主持人:Anyway這個不是重點).

$^{641}$(主持人:但是就帶到去 很屈辱).
(主持人:但是我又想).
(主持人:主耶穌 你知道會發生這件事).
(主持人:你都會讓它發生的).
(主持人:我就在這裡祈禱).
(主持人:其實就是主耶穌要給當時的百姓).
(主持人:甚至那些民事化理財人).
(主持人:根本要學習一個教訓就是).
(主持人:他們根本不是無罪或從未犯過罪).
(主持人:其實自己都有罪).
(主持人:而且這個方式是主耶穌用一個很大的憐憫).
(主持人:去憐憫這個婦人).
(主持人:我們都知道有很多的).
(主持人:addiction或者我們叫做addiction是甚麼).
(主持人:上癮的事或者我們的犯罪).
(主持人:很久都撇不清的事).
(主持人:其實我們自己的能力很難抽身的).
(主持人:但是有些極端的特別的方法).
(主持人:主耶穌就讓她在谷底的深淵抽出來).
(主持人:extreme measures).
(主持人:祂跟恆仁的婦人說).
(主持人:還有沒有人在這裡?).
(主持人:沒有了).
(主持人:沒有人定你的罪嗎?).
(主持人:沒有了).
(主持人:主耶穌說:那我都不定你的罪了).
(主持人:這豈不是恩典?).
(主持人:但是主耶穌再說一句).
(主持人:那你走吧).
(主持人:不要再犯罪了).
(主持人:這豈不是真理?).
(主持人:有真理有恩典 恩典給人有安全感).
(主持人:不被定罪 但是真理給人有方向感).
(主持人:我從此不再犯罪了).
(主持人:恩典讓我可以面對真相).
(主持人:真理引導我去走向改變).
(主持人:我都不被定罪 但是我從此不再犯罪).
(主持人:哇!我們聽過這個故事很多次).
(主持人:這個就是恩典和真理結合 除了綠化立之外).
(主持人:你想想一架火車要走的時候).

$^{681}$(主持人:只有一條很冷冰冰的路軌 沒有火車).
(主持人:就是知道有條路 但沒有火車走 沒有動力).
(主持人:但你又知道如果那架火車沒有路軌走會怎樣).
(主持人:就糟糕了).
(主持人:一定走了不知去哪了).
(主持人:不能遵循路向的方向).
(主持人:可以像真理和恩典 裡面就是火車軌和火車).
(主持人:要加起來才有帶來改變和真正的成長).
(主持人:普羅話用愛心說誠實話 就是有質素的相交的反照).
(主持人:人可以因為有恩典 承認軟弱 做真正的自己).
(主持人:但因為有真理 我們都願意被愛心提醒).
(主持人:被勸勉 被引導 不是互相討好 亦不是互相批評).
(主持人:而是在愛裡面 一起去面對真實 一起走向更新).
(主持人:這個恩典和真理 令到生命有改變).
(主持人:Henry Cloud 特別說到 要用時間去操練出來).
(主持人:群體需要有時間去操練出來).
(主持人:所以愛一開始 愛是恆久忍耐 就是要有時間).
(主持人:真正的屬靈相交 不是一次分享會就完成).
(主持人:是要一路恆久忍耐 長久同行 建立信任 經歷改變).
(主持人:玉米蛋花湯好飲 但老火燉湯的精華就在那裡出來).
(主持人:你去越南餐廳 第一口喝湯 嗯 開粉的).
(主持人:嗯 煲了十幾個鐘的).
(主持人:就知道了 既然我們對我們的舌頭都那麼嘴刁).
(主持人:為什麼不和我們的生命也嘴刁一下呢?).
(主持人:真理 生命 有時間).
(主持人:我們可以用三個簡單的問題問自己).
(主持人:幫我們檢視自己的群體生活).
(主持人:我們的相交夠不夠安全呢?).
(主持人:會不會被人敢做真實的自己說真話呢?).
(主持人:其實這就是恩典的溫度).
(主持人:Are we warm in grace?).
(主持人:我們的關係夠不夠真誠呢?).
(主持人:願意用愛心彼此提醒呢?).
(主持人:這個就是真理的亮度).
(主持人:夠不夠bright?).
(主持人:我們有沒有耐性陪弟兄姊妹).
(主持人:也讓弟兄姊妹陪我走人生成長的路呢?).
(主持人:這個就是時間的深度).
(主持人:多深呢?).
(主持人:若那三條都是yes的話).

$^{721}$(主持人:我們的相交就不只是一場聚會).
(主持人:而是讓神塑造我們生命的一個時間).
(主持人:真正的屬靈群體充滿恩典,持守真理).
(主持人:願意陪人走過的地方).
(主持人:在那裡不單止有愛).
(主持人:也會有被改變).
(主持人:這就是保羅想和我們說的話).
(主持人:而在持守真理裡面).
(主持人:我們凡事都要連於基督).
(主持人:各方面都要向基督長進).
(主持人:因為這個成熟不單止是知識多).
(主持人:我不是要be a better self).
(主持人:我不是要腦袋很大).
(主持人:我是要像主耶穌基督).
(主持人:我不是變得更好).
(主持人:但我要更像主耶穌).
(主持人:不是自我發展).
(主持人:而是基督成為我生命的中心點).
(主持人:就好像你一班人).
(主持人:我們去hiking).
(主持人:晚上去hiking真的不是太好).
(主持人:不過晚上去hiking一班人).
(主持人:然後就迷路了).
(主持人:大家都可能有不同的意見).
(主持人:不是呀 我剛才看到這裡).
(主持人:不是呀 我在那裡).
(主持人:我也覺得是).
(主持人:不是呀).
(主持人:直到大家都拿出一個準確的指南針).
(主持人:對準the true north).
(主持人:然後這就是方向).
(主持人:原來基督就是我們的北極星).
(主持人:當教會每一個人).
(主持人:都會追求更像主耶穌的時候).
(主持人:大家都一起追求).
(主持人:不要緊的 不是說先後 不是說多少).
(主持人:但都有一個同樣的心).
(主持人:我就是要更像主耶穌).
(主持人:我被主耶穌塑造 我肯開燙自己).
(主持人:我肯在外面持守真理).

$^{761}$(主持人:我肯在信仰成長).
(主持人:咦 那個教會合一成長就很自然發生了).
(主持人:我在這裡不多說很多).
(主持人:因為我們上次已經說了).
(主持人:活出敬拜人生).
(主持人:怎樣去先求祂的國 祂的義).
(主持人:怎樣在我們人生就是我們敬拜神).
(主持人:所以恩傳其實是).
(主持人:一個都很努力去學習成為屬靈辨識的教會).
(主持人:這個是凡事常連於基督).
(主持人:至於保羅 最後他就說).
(主持人:在基督身體裡 其實人人都要參與建造).
(主持人:如果我們在紅色字那裡).
(主持人:靠著祂 全身都連接得緊湊).
(主持人:全身 整個身體).
(主持人:百折 各按 各織 照著各體).
(主持人:這個就是各個肢體的功用).
(主持人:彼此相助 使身體漸漸增長 在愛中建立自己).
(主持人:基督的身體 我們弟兄姊妹互為肢體).
(主持人:這個形象 不用太多說).
(主持人:或者之後有幾堂都會說到).
(主持人:我們每個都是基督裡面的身體).
(主持人:每一個都會影響第二個).
(主持人:所以我們是不是很努力地在信仰中成長).
(主持人:都會影響彼此).
(主持人:最近的可能就是自己的家人).
(主持人:我們的生活群組).
(主持人:各領袖).
(主持人:整個教會).
(主持人:所以很重要的重點就是教會成長是全新的).
(主持人:開始自我).
(主持人:但是一切的 全新的 彼此的).
(主持人:因為我們的真理和恩典中建立).
(主持人:當然要跟別人有接觸才會學到).
(主持人:還有要有愛的).
(主持人:不是靠一個領袖 不是靠幾個同工).
(主持人:也不是靠節目 只是節目的活動).
(主持人:而是每一個肢體都有活動 有運作).
(主持人:不只是要動 不只是坐著).
(主持人:不是).

$^{801}$(主持人:祈禱本身就是一個很大的事奉 是不是?).
(主持人:在昨天看到新年嘉年華會).
(主持人:每一個部份都有弟兄姊妹).
(主持人:最開心看到的).
(主持人:很多東西都開心的 除了打球 玩遊戲).
(主持人:其實就是看到每一個部份).
(主持人:就有不同的弟兄姊妹參與).
(主持人:在其中又很賣力).
(主持人:又百折各按各職 因著他們的恩賜).
(主持人:這個實在是很開心的 看到).
(主持人:因為教會我也有份).
(主持人:我們不只是出席者).
(主持人:不只是來到 出席 完結了 拜拜).
(主持人:我們是參與者).
(主持人:我們不只是被慰養).
(主持人:我們也可以去服侍人 去供應給人).
(主持人:我們不只是等教會).
(主持人:你來到 要我就找我).
(主持人:而是我主動出擊 主動承擔).
(主持人:有什麼需要我可以參與?).
(主持人:我記得我看過一個演唱會).
(主持人:叫Matthew Harmonic Orchestra).
(主持人:管弦樂團 什麼都有的).
(主持人:有鼓 有弦樂 有管樂).
(主持人:全部是同管樂).
(主持人:但是我最).
(主持人:知不知道我最印象深刻的是哪一個樂器呢?).
(主持人:它不是打了很多).
(主持人:叮叮叮 幾下而已).
(主持人:那時候應該是在香港看的演唱會).
(主持人:它就是很準確地 等).
(主持人:來到了 叮叮叮 叮叮叮).
(主持人:可能真的不是很多).
(主持人:比起弦 同管 或者大鼓).
(主持人:叮叮叮就是那三個).
(主持人:開始是否很小).
(主持人:你一聽到叮叮叮 就不會忘記它).
(主持人:所以不是分大小 不是分多少).
(主持人:不是分動多大).
(主持人:而是我們彼此相助).

$^{841}$(主持人:願不願意我們各人都在裡面參與).
(主持人:想我怎樣可以在我們這個群體中營造).
(主持人:以致我們的教會成長 成熟和健康).
(主持人:保羅的四個標誌 信仰上要長大).
(主持人:不要只喝奶 要吃硬的東西).
(主持人:在真理裡面要有愛 要有恩典).
(主持人:這樣才在關係裡面有質素).
(主持人:我們每一天都要尋求主耶穌 尋求神).
(主持人:尋求祂的國和祂的義 方向正確).
(主持人:人人就參與建造).
(主持人:我們就會在新年講得最多的一句).
(主持人:身體健康).
(主持人:我們不是只求自己身體健康).
(主持人:我們求教會的身體都身體健康).
(主持人:以致這個健康的身體可以出去).
(主持人:但一些未認識主耶穌基督的人).
(主持人:認識真理和恩典).
(主持人:有真理沒有恩典 教會很硬).
(主持人:有恩典沒有真理 教會會很飄忽).
(主持人:沒有聯於元首基督 教會會很迷失).
(主持人:沒有彼此參與 教會會很虛弱).
(主持人:所以真正的健康使命教會).
(主持人:是愛中持守真理 在基督裡面長大).
(主持人:彼此中建造).
(主持人:弟兄姊妹 我們肯不肯?想不想?).
(主持人:我們一齊做吧 好嗎?Amen).
(主持人:讓我們一齊祈禱).
(主持人:多謝你給保羅).
(主持人:你給我們藉著保羅所說的).
(主持人:在主耶穌基督見到有真理 有恩典 更有時間).
(主持人:讓我們每日的生活裡 真是去模造我們).
(主持人:讓我們紮根在你自己的話語當中 你的真理當中).
(主持人:也讓我們順順感到你給我們的愛).
(主持人:讓我們都會操練 用敬畏你的心去愛神 去愛人).
(主持人:求主你叫我們這個光 不是在斗底的).
(主持人:求主你叫我們為世界的炎 不是沒味的).
(主持人:求主你給恩全家 成為一個健康的使命 有使命的教會).
(主持人:傳揚你自己的福音 帶出你的真理和恩典).
(主持人:多謝你 奉主耶穌基督的使命祈禱 Amen).
\newpage



\section{以斯帖記 8:3-9:16}
\label{sec:031693}
\textbf{神的救贖 人的責任}~~~~黃國棟醫生
\newline
\newline
連結: \href{http://www.church.com.hk/acms/content.asp?site=cdc\&op=show\&type=product\&code=031693\&layout=sermon}{\texttt{031693}} ~~~~ 語音日期: 2026-02-22
\newline
\newline
\hyperref[sec:031692]{\small{< < < PREV SERMON < < <}}
~
\hyperref[sec:index_chronic]{\small{[返順時目]}}
~
\hyperref[sec:031694]{\small{> > > NEXT SERMON > > >}}
\newline
\newline
以斯帖記 8:3-9:16
\newline
\begin{longtable}{cl}
\hline
\hline
章節 & 經文 (和合本修訂版)\\
\hline
8:3 & \begin{tabularx}{0.7\textwidth}{X} 以斯帖又在王面前求情，俯伏在他腳前，流淚哀求他阻止亞甲人哈曼害猶太人的惡謀。 \end{tabularx} \\ \\ \relax
8:4 & \begin{tabularx}{0.7\textwidth}{X} 王向以斯帖伸出金杖，以斯帖就起來，站在王面前， \end{tabularx} \\ \\ \relax
8:5 & \begin{tabularx}{0.7\textwidth}{X} 說：「王若以為好，我若在王面前蒙恩，王若認為合宜，我若在王眼前得喜悅，請王下諭旨，廢除亞甲人哈米大他的兒子哈曼設謀，要殺滅王各省的猶太人所頒的詔書。 \end{tabularx} \\ \\ \relax
8:6 & \begin{tabularx}{0.7\textwidth}{X} 我何忍見我本族的人受害？何忍見我同宗的人被滅呢？」 \end{tabularx} \\ \\ \relax
8:7 & \begin{tabularx}{0.7\textwidth}{X} 亞哈隨魯王對以斯帖王后和猶太人末底改說：「因為哈曼要下手害猶太人，看哪，我已把他的家產賜給以斯帖，也把哈曼掛在木架上了。 \end{tabularx} \\ \\ \relax
8:8 & \begin{tabularx}{0.7\textwidth}{X} 你們可以照你們看為好的，奉王的名寫諭旨給猶太人，用王的戒指蓋印；因為奉王的名所寫、用王的戒指蓋印的諭旨是不能廢除的。」 \end{tabularx} \\ \\ \relax
8:9 & \begin{tabularx}{0.7\textwidth}{X} 三月，就是西彎月二十三日，當時王的一些書記受召而來，按著末底改所吩咐的，用各省的文字、各族的語言，以及猶太人的文字語言寫諭旨，傳給那從印度直到古實一百二十七省的猶太人，以及總督、省長和領袖。 \end{tabularx} \\ \\ \relax
8:10 & \begin{tabularx}{0.7\textwidth}{X} 末底改奉亞哈隨魯王的名寫諭旨，用王的戒指蓋印，交給信差們騎上御用的王室快馬去頒佈。 \end{tabularx} \\ \\ \relax
8:11 & \begin{tabularx}{0.7\textwidth}{X} 王准各城各鎮的猶太人在一日之內，在十二月，就是亞達月的十三日聚集，在亞哈隨魯王的各省保護自己的性命，剪除，殺戮，滅絕那要攻擊猶太人的各省各族所有的軍隊，以及他們的妻子兒女，奪取他們的財產為掠物。 \end{tabularx} \\ \\ \relax
8:12 & \begin{tabularx}{0.7\textwidth}{X} 見上節 \end{tabularx} \\ \\ \relax
8:13 & \begin{tabularx}{0.7\textwidth}{X} 這諭旨的抄本以敕令的方式在各省頒佈，通知各族，使猶太人預備等候那日，好在仇敵身上報仇。 \end{tabularx} \\ \\ \relax
8:14 & \begin{tabularx}{0.7\textwidth}{X} 於是騎御用快馬的信差奉王命催促，急忙起行；敕令傳遍了書珊城堡。 \end{tabularx} \\ \\ \relax
8:15 & \begin{tabularx}{0.7\textwidth}{X} 末底改穿著藍色白色的朝服，頭戴大金冠冕，又穿紫色細麻布的外袍，從王面前出來；書珊城充滿了歡樂的呼聲。 \end{tabularx} \\ \\ \relax
8:16 & \begin{tabularx}{0.7\textwidth}{X} 猶太人有光榮，歡喜快樂，得享尊貴。 \end{tabularx} \\ \\ \relax
8:17 & \begin{tabularx}{0.7\textwidth}{X} 王的諭旨和敕令所到的各省各城，猶太人都歡喜快樂，擺設宴席，以那日為吉日。國中許多民族的人因懼怕猶太人，就自稱為猶太人。 \end{tabularx} \\ \\ \relax
9:1 & \begin{tabularx}{0.7\textwidth}{X} 十二月，就是亞達月十三日，王的諭旨和敕令要執行的那一日，猶太人的仇敵盼望制伏他們，但猶太人反倒制伏了恨他們的人。 \end{tabularx} \\ \\ \relax
9:2 & \begin{tabularx}{0.7\textwidth}{X} 猶太人在亞哈隨魯王各省的城裡聚集，下手擊殺那些要害他們的人。沒有人能在他們面前站立得住，因為各民族都懼怕他們。 \end{tabularx} \\ \\ \relax
9:3 & \begin{tabularx}{0.7\textwidth}{X} 各省的領袖、總督、省長，和辦理王事務的人，因懼怕末底改，就都幫助猶太人。 \end{tabularx} \\ \\ \relax
9:4 & \begin{tabularx}{0.7\textwidth}{X} 末底改在朝中為大，名聲傳遍各省；末底改這人的權勢日漸擴大。 \end{tabularx} \\ \\ \relax
9:5 & \begin{tabularx}{0.7\textwidth}{X} 猶太人用刀擊殺所有的仇敵，殺滅他們，隨意待那些恨他們的人。 \end{tabularx} \\ \\ \relax
9:6 & \begin{tabularx}{0.7\textwidth}{X} 在書珊城堡中，猶太人殺滅了五百人。 \end{tabularx} \\ \\ \relax
9:7 & \begin{tabularx}{0.7\textwidth}{X} 他們殺了巴珊大他、達分、亞斯帕他、 \end{tabularx} \\ \\ \relax
9:8 & \begin{tabularx}{0.7\textwidth}{X} 破拉他、亞大利雅、亞利大他、 \end{tabularx} \\ \\ \relax
9:9 & \begin{tabularx}{0.7\textwidth}{X} 帕瑪斯他、亞利賽、亞利代、瓦耶撒他； \end{tabularx} \\ \\ \relax
9:10 & \begin{tabularx}{0.7\textwidth}{X} 這十人都是哈米大他的孫子，猶太人的仇敵哈曼的兒子。猶太人卻沒有下手奪取財物。 \end{tabularx} \\ \\ \relax
9:11 & \begin{tabularx}{0.7\textwidth}{X} 那日，書珊城堡中被殺的人數呈報到王面前。 \end{tabularx} \\ \\ \relax
9:12 & \begin{tabularx}{0.7\textwidth}{X} 王對以斯帖王后說：「猶太人在書珊城堡中殺滅了五百人，又殺了哈曼的十個兒子，在王其餘的各省不知如何。你要甚麼，必賜給你；你還求甚麼，也必為你成就。」 \end{tabularx} \\ \\ \relax
9:13 & \begin{tabularx}{0.7\textwidth}{X} 以斯帖說：「王若以為好，求你允准書珊的猶太人，明日也照今日的諭旨去做，並把哈曼十個兒子的屍體掛在木架上。」 \end{tabularx} \\ \\ \relax
9:14 & \begin{tabularx}{0.7\textwidth}{X} 王允准這麼做。敕令傳遍書珊，哈曼十個兒子的屍體被掛了起來。 \end{tabularx} \\ \\ \relax
9:15 & \begin{tabularx}{0.7\textwidth}{X} 亞達月十四日，在書珊的猶太人又聚集，在書珊殺了三百人，卻沒有下手奪取財物。 \end{tabularx} \\ \\ \relax
9:16 & \begin{tabularx}{0.7\textwidth}{X} 亞達月十三日，在王各省其餘的猶太人也都聚集，保護自己的性命，擺脫仇敵得享平安。他們殺了七萬五千個恨他們的人，卻沒有下手奪取財物； \end{tabularx} \\ \\
[1ex]
\hline
\hline
\end{longtable}
$^{1}$大家早晨.
很高興又有機會回來和大家一起敬拜.
今天還是新年.
所以在這裡再一次參氣的向大家拜年.
恭喜大家今年新年.
這個馬年是龍馬精神.
在芝加哥我們知道.
每一年都會有農曆新年的巡遊.
我看過一些報紙就說.
今年的巡遊就在下個星期日.
不過我走了就來不了.
但去年我住在芝加哥都有去見過這個遊行.
都是很開心很好玩的一件事.
不過如果大家知不知道.
如果你在下個星期一.
你去到這裡的猶太人社區的話.
你又會見到另外一個巡遊.
這個就是今年的三月二號和三號.
星期一和星期二.
他們會慶祝猶太人一個節日.
叫做普爾節.
普爾節不是保禮.
不要看錯.
保禮的禮節是三點水.
這個禮節是黃字邊的.
他們又是會巡遊.
穿著很可愛的化曬裝.
穿著一些很可愛的衣服在這裡巡遊.
猶太人和華人很相似的.
就是他們每年都有很多節日.
而這些節日是有不同的意思.
可能有些是我們很熟悉的.
譬如是魚節.
是聖經舊約裡面有記錄過的節日.
普爾節就是另外一個在舊約裡面.
有記錄過的節日.
它的起源是在《儀式帖記》裡面所講到的.
就是說他們會在這個叫做亞達月的十四十五號.
就是為了猶太人脫離仇敵.
得到平安.

$^{41}$祭普爾原來是抽籤的意思.
就定了這兩天是做節日.
就要他們守節.
其實普爾節和魚節就是猶太人現在所守的節日裡面.
兩個在舊約裡面有紀念.
有記錄到的節日.
現在猶太人最重要的節日.
我聽我的一些猶太朋友.
我的兩個老師都是猶太人.
他們就說最重要的是他們的聖誕節.
就是在聖誕附近.
聖誕前後的一個節日.
但是除了漢尼卡之外.
就是普爾節,魚節也是他們很重要的節日.
而這兩個就是在聖經裡面有記錄下的節日.
今天我就想和大家.
用普爾節的故事.
看看神可以怎樣在我們裡面施行祂的救贖.
和這個節日的故事.
對我們今天又有什麼啟示.
神給了一個很重要的諾言給我們.
就是記載在詩篇裡面.
祂說:你留心聽我的禱告.
誰聽我懇求的聲音.
我在患難的那一天求告你.
你就必應運我.
在這裡我們看到一個formula.
有一條程式.
就是我們在患難裡面求告神.
神就會應運.
這是一個非常非常強力的應許.
給了我們信心.
我們可能現在看見這個世界好像很亂.
好像在患難當中水深火熱.
不知道怎樣.
但是神這個應許這個諾言.
是永遠存在的.
在我們的心裡面.
是令我們不用怕.
是可以很安心的面對各種挑戰.

$^{81}$希望大家可以絕對絕對的記住神這個應許.
不過這個應許裡面.
又有一個很有趣的問題.
就是神怎樣應運我們.
我們不妥了.
不行了.
或者我們病了.
或者我們沒錢.
被人lay off了.
或者我們有什麼法律的問題.
神怎樣可以拯救應運我們呢.
這個才是所謂的64 million dollar questions.
不是神會不會應運.
而是神怎樣應運.
在華人教會裡面.
我們很常見的有一個救贖觀.
或者一個祈求的觀念.
就是這個沙灘上的足跡.
這個土文.
這個故事或者這個祈禱.
大家可能都很熟悉.
就是有一個人在死了之後.
他就看回他生命裡面.
就是原來在沙灘上走的.
他就發覺主耶穌都一直陪著他走的.
所以在沙灘上就見到有兩個人的腳印.
但是他又見到.
他不碰很有困難的時候.
就有一對腳印.
於是他就問神.
他說主耶穌啊.
為什麼我不妥的時候.
你就不和我走.
我自己在這裡走.
主耶穌就說.
不是啊.
你不妥.
你有困難的時候.
是我抱著你走.
不是你自己走.

$^{121}$這個故事是一個很感人的故事.
很motivate人.
也是令到我們覺得.
原來我不碰的時候.
神就會抱著我走.
我不用做的.
我不用擔心的.
神會抱著我.
我不用.
我給他就行了.
我靠他就行了.
這個是一個很感動.
令到我們很安心的故事.
但是我要和大家講的就是.
這個是一個偏差了的救贖觀.
神不是一定這樣救我們的.
神可不可以這樣救我們呢.
是不是我們有困難的時候.
祂就抱著我們.
不用我們出力.
我們什麼都不用理.
祂抱著我們走.
我們就抱著祂就行了.
我有幾個孫子.
我的孫子很喜歡抱著我們.
有什麼事的時候.
你抱著祂就行了.
但是神是不是這樣.
我們有問題的時候.
就被我們抱著.
祂抱著我們走呢.
當然可以.
神可以這樣做.
絕對可以.
譬如有一個例子.
就是列王記下六章24到7章20節的例子.
講到在以利沙的時候.
亞蘭王攻擊撒瑪利亞.
突然間神就見到撒瑪利亞的軍隊聽見有聲音.
很害怕.

$^{161}$就突然走了.
那些以色列人什麼都不用做.
以利沙只是祈禱.
不用出去打.
那些人就走了.
這個故事裡面還有很有趣的.
就是有四個大麻風和軍長的故事.
今天不詳細講了.
大家可以回去看一下.
這個大故事裡面有一個小故事.
講到四個大麻風可以怎樣得到神的恩典.
而這個不信以利沙講的軍長又怎樣.
看見有得吃了.
本來我們機方沒得吃的.
但是看見那些亞蘭人突然間走了.
行了.
但是他們就有得看沒得吃.
結果就死了.
神可不可以這樣救贖我們.
什麼都不用做.
抱著我們走呢.
絕對可以.
但是.
神是不是只可以這樣做呢.
我就會說不是.
這個故事本身沒錯.
但是呢.
如果我們只是堅持神一定只可以這樣救我們.
當我們有困難的時候.
我們就什麼都不用做了.
就是在求神.
神就會救我們.
這個是我們心態上的偏差.
因為如果限了神可以這樣救我們的話.
我想問的就是.
那誰是神呢.
你當神是這個Ginny.
你擦一下神燈.
祂就會出來救你.
或者你當神是有求必應的黃大仙.

$^{201}$什麼菩薩.
還是這個是我們的主.
我們的神呢.
今天.
我就想和大家在以斯帖的故事裡.
這個普爾節的故事裡.
看一下神怎樣去救贖.
以斯帖這個故事是一個很有趣的故事.
剛才大家介紹有這個課程是講聖經裡面的敘事.
成本以斯帖其實就是一個長的故事.
一氣呵成的故事.
它的地點是波斯帝國.
是當時世界上最強的帝國.
它的地點是包括從非洲去到印度.
在中東一帶.
時間是約莫在公元前480年.
那時候如果你和聖經的舊約那裡看.
以色列人已經是一部分回去了.
是過了70年的被擄.
聖殿又建好了.
第二次聖殿又建好了.
但是還有很多以色列人.
猶太人是散居在各地的.
而以斯帖的時間就差不多和聖經最後的歷史記錄.
好像以斯拉.
尼希米是差不多時間寫成的書.
今天我們就可以在這個故事裡面.
看看神用另外一個方法怎樣去救贖.
和今天對我們有什麼意思.
以斯帖記是一本很特別的書.
如果你讀聖經你會發覺.
可能有一個問題就是.
為什麼聖經有一本這樣的書呢?.
它是聖經裡面唯一.
全本書沒有提過神這個字.
沒有說過有祈禱.
在故事裡面沒有提過祈禱這回事.
沒有任何的神蹟.
只不過是一個劇本.
如果你拍戲它就好像一齣戲.

$^{241}$一個劇本.
沒有任何的金句.
問題是.
為什麼這本書會放在聖經裡面.
成為聖經的正典呢?.
在猶太人裡面.
他們為什麼覺得這本書要放在舊約聖經裡面.
放在他們的聖經裡面.
變成一本他們不一定要讀的書呢?.
在今天經過這麼多次的宗教改革.
為什麼這本書還在呢?.
一本書沒有提過神.
甚至神這個字都沒有提過.
沒有提過祈禱.
沒有提過神蹟.
為什麼還在聖經裡面呢?.
因為這本書.
它講到的是神怎樣去拯救.
是不是有一張.
可能有一張slide不見了.
OK,不要緊.
這本書之所以成為正典.
現在解經家有一個分析就是.
這本書講到的是一個paradigm shift.
在舊約裡面一個規範轉移.
就是由神直接的介入.
轉到神在我們生活裡面介入.
你看舊約以前的神跡.
以色列人的歷史.
神每一件事都是直接介入的.
好像剛才講那個.
以利沙的時候那個例子.
他們被人攻擊了.
神直接介入.
發生了神跡.
於是以色列人得救.
但是以斯帖記沒有這樣的東西.
神不再是直接用神跡來介入人的歷史.
而是在日常生活裡面.
在生活裡面.

$^{281}$是介入我們生活的每一部份.
在那裡給我們帶領.
給我們保障.
這個故事就講神怎樣拯救以色列人.
怎樣去發生不同的事.
首先就是.
在故事裡面.
很有趣.
神不是叫敵人突然死光就算了.
他是說.
那時候的皇帝不可以取消.
他之前屠殺以色列人的預旨.
於是他就給了一條新的聖旨出來.
就說現在給你以色列人可以自衛.
猶太人就組織自衛隊.
反制敵人.
今天因為時間問題.
不可以將整個故事講完.
所以希望大家回去.
將以斯帖記從頭到尾一次過.
好像看小說一樣看完.
以斯帖記不可以平時查經讀一兩節.
一定要由頭到尾一氣呵成的看完.
因為是一個故事.
整本書就是一個故事.
你一定要由頭到尾.
好像追小說一樣追完.
你才知道它發生什麼事.
在事件最後.
猶太人慶祝他們得救.
慶祝這個普珥節直到今天.
這個故事大家回去看.
講到以斯帖做了波斯的皇后.
她的養父或堂兄得罪了奸臣.
奸臣哈曼想殺死猶太人.
以斯帖冒死去見皇帝.
但他沒有求皇帝什麼.
他只叫皇帝陪他吃飯.
在兩次吃飯裡.
有些奇怪的事發生了.

$^{321}$令到這個奸臣死了.
主死了.
猶太人就得救了.
詳細的事我等大家自己回去看.
我不詳細講了.
我想講的就是.
在這個故事裡有三個特徵.
對今天我們尋求神的救贖來說.
是有特別意思的.
第一個特徵在以斯帖裡見到是.
神提供了機會.
但是人有執行的責任.
和我們剛才講的.
沙灘捉人的故事很不同.
在那個模式裡.
我們希望什麼都不用做.
神抱著你走就行了.
但是在以斯帖裡見到.
是不同的.
神給我們機會.
我們自己也要做事.
第二個我們見到一個很特別的拯救方法就是.
在這個救贖的故事裡.
我們見不到有什麼英雄.
有什麼先知.
有什麼強人領袖.
第三個特徵就是.
這個拯救是在一個現有的機制裡進行的.
不再是一些我們想不到的東西.
或者是一些神蹟.
是超自然的東西.
而是在一些現有的機制裡進行的.
我們就看看這三點.
對我們今天有什麼應用的價值.
第一個特點就是.
神提供了機會.
人有執行的責任.
這次神沒有幫以色列人消滅敵人.
跟剛才說的那個故事.
列王記下的故事相比.

$^{361}$不同了.
列王記下裡.
神叫敵人聽見聲音就走了.
還有其他很多這些神蹟.
之前的故事都見到.
神直接出手.
打倒了敵人.
以色列人就得救了.
包括了在出埃及記的時候.
神降災.
直接殺了那些長子.
預節.
以色列人什麼都不用做.
坐在家裡.
吃個羊.
塗點血就行了.
不用做事了.
但這次不同了.
神沒有幫以色列人消滅仇敵.
沒有任何超自然的神蹟.
如果我們在這裡祈禱等神蹟.
可能我們已經被人殺光了.
神給了我們九個月的時間去準備.
如果你看這個故事.
你會發覺.
他們得到這個令羽說.
你們可以組織自衛隊.
你們是要自衛.
他們有九個時間去準備.
怎麼準備?.
他們要計劃.
有人說九個月之後.
我要把你們殺光.
你有九個月時間去預備.
你怎麼預備?.
你當然要想想.
如果要打到來.
我怎麼弄呢?.
我的屋子怎麼弄.
弄強一點呢?.

$^{401}$或者買好棍子.
磨好刀子.
要練拳.
健身.
你要自衛.
你不是手無筲箕之力.
你怎麼自衛?.
你當然要操練一下.
跑步.
舉重.
儲點糧食.
預算別人要來打你.
你當然要有準備.
至少要自衛.
到了時候.
你需要自己執行這個自衛拳.
神不會幫你一雷打下去.
拿刀衝過來的時候.
把你打倒.
你要自己出去打.
今天我們遇到困難的時候.
我們又想要什麼拯救呢?.
我們是不是想神.
求你把這個困難拿走呢?.
還是我們一樣要準備呢?.
我們是想要一個神跡的見證.
我有病了.
於是我就求神.
叫我病了.
不用看醫生.
不用吃藥.
突然間就沒有了.
我沒有工作.
於是叫神.
突然間有人三顧草廬.
敲你門.
什麼生.
什麼小姐.
我們聽過你的大名.
知道你很厲害.

$^{441}$重金禮聘你出山.
幫我們做事.
你自己不用找工作.
不用send resume出去.
不用工作呢?.
我們是不是將神.
當一個有求必應的Genie.
一個燈神.
擦爛燈神就出來搞定呢?.
自己什麼都不用做呢?.
我們怎麼準備呢?.
如果你沒有工作了.
你就坐在那裡.
沒有工作了.
神啊.
你幫我找工作吧.
你需不需要讀書.
寫好你的resume.
看看外面哪家公司有什麼呢?.
你可能是想要一個男朋友.
要一個女朋友.
你很寂寞.
是不是坐在家裡.
做宅男宅女.
只顧著玩遊戲.
那又如何呢?.
需不需要自己學一下.
去做好你的儀容.
扮好一點.
說話的時候說得合適一點呢?.
還是坐在那裡等神來.
當下幫你做個女人.
擺在你身邊.
你一醒.
這樣就行了呢?.
我們自己要準備.
神會供應.
但是我們也有執行的責任.
這是第一個.
在以色列的故事裡面.

$^{481}$我們看到的模式.
神給我們機會.
但是我們自己要做事.
第二個.
以色列裡面看到的模式是.
我們再也見不到英雄.
見不到什麼先知.
見不到什麼強人領袖.
以前以色列人來看歷史.
見到的模式就是.
一個英雄.
一個勇士.
或者一個先知.
或者一個君王.
他帶著一些人去打仗.
你看士師記.
你看武爾記.
你看列王記.
歷代至上.
歷代至下.
你也看到這樣的模式.
就是一個強人.
一個領袖.
帶著一幫人出去打.
他打贏就贏.
他打輸就輸.
但是在以斯帖記裡面.
我們見不到這個模式.
這次我們見不到一個領袖.
帶著他們去打仗.
你再看回那個故事.
以斯帖和末底改.
似乎完了之後.
就沒有出聲了.
沒有再做事了.
故事去到第六章第七章之後.
他們消失了.
以斯帖見到皇帝.
皇帝說以色列人可以自衛.
然後沒有了.

$^{521}$他們消失了.
再也見不到他們說什麼了.
沒有英雄.
沒有領導.
沒有一個領導.
當時的以色列人.
他們是散居在127個省.
從印度直到古實.
古實就是現在的埃塞俄比亞.
是非洲.
印度那時候的波斯帝國.
就去到印度的邊緣.
就是現在巴基斯坦.
差不多的地方.
有沙漠.
有森林.
有海.
有山.
有湖.
全部都散了.
你沒有辦法可以一個人.
那時候沒有Starlink.
Elon Musk那些衛星.
飛上去.
你沒有可能將一些訊息.
傳到不同的地方.
沒有可能一個領袖.
教一些以色列人怎樣打的.
他們唯有就是.
自己每一個地方.
想想我們自己在這個地方.
遇到問題.
我們怎樣解決.
沒有辦法一個人指揮.
也都需要因為他們的環境和安排.
來決定怎樣去自衛.
今天對我們有什麼意思呢.
我覺得這個model.
對今天教會來說.
是更加適用的.

$^{561}$就是我們再沒有處境化.
one size fit all.
解決問題的方法.
對不起.
我們是再沒有.
one size fit all.
解決問題的方法.
我們是需要.
每一件事都處境化.
根據當時當地的環境.
來做事.
在教會如果我們看回.
過去十幾年的歷史.
我們發覺很有趣.
十幾二十年前.
很多教會是一窩蜂.
去跟這套那套的.
跟Wick Warren.
Purpose Driven Church.
每個人都買一本書.
跟著芝加哥.
跟Willow Creek.
每個人都去跟他.
去他的seminar.
去那裡.
跟著Hew Song.
又去跟.
現在這些Mega Church.
我們見到很多都是沒落.
我不說神學上有什麼問題.
只不過見到.
當大家一窩蜂.
去跟一套東西的時候.
原來發覺不行.
我們每一個人.
每一個教會.
都需要跟著自己的.
那時候的處境.
來解決問題.
一個很重要的概念就是.

$^{601}$信徒皆濟私.
我們每個人.
都是直接跟神有關係.
我們不用特別跟著.
一些強人一些領袖.
說什麼就做什麼.
絕對不是叫大家.
不聽牧師.
不聽麗仙他們說話.
我們絕對尊重他們.
他們的經驗.
他們的教導.
都很有價值.
但是我們就不要去崇拜.
任何的那些強人.
特別是我們接觸不到那些.
外面的那些.
Detailed Evangelist.
那些不知道說什麼的.
走出來好像很有型.
其實背後很多骯髒的東西.
我們不知道.
我們大家要小心.
我們每一個人.
都可以直接跟神有一個關係.
是跟神那裡得到力量.
是得到接觸.
我們不要迷信.
什麼明星講完.
什麼宗教領袖.
我們是需要根據個人背景.
我們每一個人自己的背景.
當我們遇到問題的時候.
神會拯救我們.
但是會根據我們的背景.
來發展出拯救的方法.
講個秘密給大家知道.
我其實很多時候.
都是在大家教會裡面崇拜的.
雖然是在網上.

$^{641}$所以我知道大家幾個禮拜前.
有一個MBTI的workshop.
那時候我是很衝動的話.
如果我來得到就好了.
為什麼呢.
MBTI對我來說.
是特別有意思的.
可能大家.
有些人會知道.
我做醫生之前.
我是做資訊系統的.
我在UCLA那裡讀IT的.
IS的PhD的.
我的博士論文就是用MBTI來研究.
怎樣去訓練一些.
Managers去用那時候的software.
所以MBTI對我來說.
那個意思是非常之重要的.
見到大家已經過了四十年了.
這個工具還在.
見到大家是願意在MBTI裡面.
明白自己那個.
對我來說Cognitive Style.
就是你怎樣處理問題的方法.
是非常之重要的.
因為你不同的Cognitive Style.
你就有不同處理問題的方法.
我處理問題的方法.
和你處理問題的方法可能很不同.
絕對沒有辦法說.
誰對誰錯.
你要根據你的Personality.
你發展出來處理問題的方法.
所以很高興見到.
你教會會做這件事.
第三個.
神拯救的方法就是.
他會在現在的機制.
或者基礎建設的Foundation裡面進行.
他不再是神職的.

$^{681}$突然在上面走個雷.
砍了敵人就算了.
而是他在一個機制裡面.
現有的存在機制裡面.
來發展出拯救方法.
在當時去拯救以色列人的時候.
是在當時的法律的框架裡面進行的.
不是超自然的.
很有趣的就是.
為什麼要說這件事.
就是當時為什麼神就這樣說.
本來說要殺死那些以色列人.
不殺了.
我也不知道為什麼.
但是原來當時波斯有一個.
有一條法律就是說.
皇帝用戒指吸了印的東西.
就不可以取消的了.
皇帝第一次就說.
殺死那些以色列人.
他就不可以說.
現在不殺了就算了.
於是他就要出過第二條的命令.
就是說.
那些人可能還會殺你.
還會打你.
但是我准你們自衛.
所以他是在現有的機制.
法律基礎裡面.
幫以色列人找到一條出路.
今天呢.
神也可以用同樣的方法來拯救我們.
他會用已經存在的機制.
的基礎來拯救我們.
包括了什麼呢.
科學,醫學.
如果你有病.
我們不一定要求神.
你不用吃藥.
不用看醫生.

$^{721}$你就好了.
神可以這樣醫治.
但是神不需要這樣醫治.
如果神通過醫生.
通過藥物.
通過科學來拯救我們.
絕對依然是祂的掌管.
是祂的神蹟.
神可以通過法律來拯救我們.
如果你遇到什麼麻煩.
神可以通過現在有的法律來拯救我們.
神可以通過制度來拯救我們.
已經有的制度.
如果我們被人欺負.
我們去投訴.
這不是什麼不屬靈.
這是神通過現有的法律來幫助我們.
說一個好笑的例子給我們知道.
給大家聽聽一個故事.
真人真事在我身上發生的.
很多年前.
有一次和一個弟兄吃飯.
這個弟兄很慘.
他和我投訴.
他被人撞車在加州.
撞車的人沒有保險.
Uninsured motorist.
我朋友的車被偷走了.
他人受傷了.
看了身體療法.
但是對方不給錢.
於是他沒有錢給.
被人追債.
本來他買了Uninsured motorist的保險.
他的保險公司會賠償.
但是他發現一間教會.
一間很.
我不可以說是邪教.
但是很極端的教會.
就是告訴他.

$^{761}$基督徒不可以打官司.
你不可以去告別人.
你告別人就不符合聖經.
他就很慘.
他的保險公司就說.
你要我賠償給你沒有問題.
但是你要先file個claim去告那個人.
要file個claim去告那個人.
那個人不賠償.
我們的Uninsured motorist.
才會kick in.
會賠償給你.
這樣也很合理.
保險公司有他們做事的方法.
不是你說別人不賠償.
你問他拿錢.
這樣就不行了.
我朋友就很慘.
他的目的是告訴他.
他不可以去打官司.
不可以去告別人.
他的車壞了.
沒有錢買車.
他就被人追去chiropractor.
還是physical therapist的錢.
於是他就走來.
我也是剛好和他吃飯.
他就在我面前發呆.
我拗了頭.
我跟他說不是.
聖經說不可打官司.
保羅說不可以打官司.
好像這個情況.
不是很適合.
不是很適用.
保羅自己也打官司.
保羅被人告的時候.
上訴該殺.
為什麼要打官司.
結果怎樣呢.

$^{801}$我朋友就解決了問題.
但是解決問題的方法是什麼呢.
轉教會.
這個就是.
神是用現有的機制.
法律.
制度.
來幫助去解決問題.
還有一樣是朋友.
神可能通過朋友.
幫助我們解決問題.
總結就是.
我另外一個故事和大家說.
有一個爸爸和兒子說.
兒子你不會做功課.
不用怕.
我幫你.
於是有一天.
兒子就和他說.
老爸我不會做功課了.
你幫我做我就去打機.
老爸拗了頭說.
為什麼你不做功課.
叫我幫你去打機.
兒子就說.
你說幫我.
你幫我.
你幫我做.
我就去打機.
我們今天是不是這樣要神呢.
神幫我們.
於是我們什麼都不做.
我們就去打機.
他就抱著我們在山上走.
我們就抱著他.
很舒服的就算了.
如果你做爸爸的.
你是不是永遠都想幫你兒子做功課.
解決問題呢.
還是你只是想給他資源.

$^{841}$給他支持.
讓他學怎樣去解決問題呢.
到底我們想做這個巨嬰.
現在這個term很popular.
巨嬰.
我們想做神的巨嬰.
神抱著我們.
我們咬著他的奶嘴.
不用做.
還是我們想神拖著我們.
在沙灘裡面走呢.
到底我們作為基督徒.
我們希望得到的拯救.
是怎樣呢.
到底天父想我們怎樣呢.
這個我就希望大家可以想一想這個問題.
多謝大家.
\newpage



\section{馬太福音 5:27-30}
\label{sec:031694}
\textbf{重塑生命系列 (4) 色慾與貞潔}~~~~張健庭牧師博士
\newline
\newline
連結: \href{http://www.church.com.hk/acms/content.asp?site=cdc\&op=show\&type=product\&code=031694\&layout=sermon}{\texttt{031694}} ~~~~ 語音日期: 2026-02-22
\newline
\newline
\hyperref[sec:031693]{\small{< < < PREV SERMON < < <}}
~
\hyperref[sec:index_chronic]{\small{[返順時目]}}
~
\hyperref[sec:031695]{\small{> > > NEXT SERMON > > >}}
\newline
\newline
馬太福音 5:27-30
\newline
\begin{longtable}{cl}
\hline
\hline
章節 & 經文 (和合本修訂版)\\
\hline
5:27 & \begin{tabularx}{0.7\textwidth}{X} 「你們聽過有話說：『不可姦淫。』 \end{tabularx} \\ \\ \relax
5:28 & \begin{tabularx}{0.7\textwidth}{X} 但是我告訴你們：凡看見婦女就動淫念的，這人心裡已經與她犯姦淫了。 \end{tabularx} \\ \\ \relax
5:29 & \begin{tabularx}{0.7\textwidth}{X} 若是你的右眼使你跌倒，就把它挖出來，丟掉。寧可失去身體中的一部分，也不讓整個身體被扔進地獄。 \end{tabularx} \\ \\ \relax
5:30 & \begin{tabularx}{0.7\textwidth}{X} 若是你的右手使你跌倒，就把它砍下來，丟掉。寧可失去身體中的一部分，也不讓整個身體下地獄。」 \end{tabularx} \\ \\
[1ex]
\hline
\hline
\end{longtable}
$^{1}$好,弟兄姊妹,新年快樂!.
在新年的時候,願意我們一起領受在新的一年,更多領受神的恩典,願神的福氣與我們同在,願我們的福杯滿日..
好,我們一起先低頭祈禱吧..
在農曆新年的時候,給我們思想到我們在我們的親朋戚友當中,一起歡樂,一起聚在一起的時間,我們更加想起你給我們有教會的家庭,在主裡邊,我們都好像一班家人,我們在主裡邊一起去慶祝,一起去經歷,也是一起去領受你的恩惠..
願你帶領我們在新的一年,給我們有更多的蒙福,也有更多的能夠去行在你的旨意當中,奉耶穌基督的名祈求,阿們..
你知道農曆新年的時候,中國人必少不開吃東西,聚餐,團年飯,吃完過兩天就開年飯,這些是很開心的日子..
但是不知道你知不知道,剛剛過了的那個星期三是什麼日子?.
如果你在街上留意,特別你不是中國人,你會看到有些人的額頭有些灰弄黑了..
剛剛過了的那個星期三,就是Ash Wednesday,就是勞灰星期三..
一般來說,在基督信仰,無論是天主教或者禮儀教會,例如聖公會,路德宗,他們都有的..
Ash Wednesday,意思就是在主耶穌復活節之前40天開始,通常這個就稱為Lent Season,就是齋節期..
就是在這40天當中,我們去檢視自己的慾望,或者盡量不去在很多的聲色犬馬的生活當中停留,.
而去將我們的思想,將我們重新再一次聚焦,去思想我們在世界上是客女,而耶穌基督為我們寫了十字架,教導我們背起十字架去寫記..
所以這個Lent Season都是非常有意思的,因為它提醒到我們一樣東西,就是其實我們人活在世上,我們是一個物質的層面,我們是有很多的需求..
你需要吃,需要住,需要有溫飽,需要有安全感,各方面這些是在這個世界上物質的需求,這個是沒有錯的,神創造我們就是有這樣的需要..
但是問題是我們不是單單這些需要就夠,我們是屬靈的人,那時候其實是這個物質的需求,但是當人犯罪墮落之後,.
這些物質的需求往往不斷擴大,甚至乎侵佔我們整個人生的各方各面,以致那一切就成為了我們只有這裡這麼多東西,就是需要這些東西..
而來將我們跟神的關係,跟自己的內在靈性的反省都好像沒有再去重視了..
所以Lent Season其實就是每年一次,特別是復活節之前,就是壽苦節有一段時間,讓我們去反省人生在世的時候,其實是不是我們的生命就是被這些東西佔據了呢?.
我們的慾望究竟是專注在,是不是只是在這些物質方面,身體方面的需求呢?.
所以其實一直在講這一系列的,特別講到那個七個致命的罪的時候,其實這些很多時候在信仰的歷史當中,其實很多時候都是在Lent Season裡面講的,其實這些的題目..
因為其實就是提醒我們活在世上的時候,有時候我們實在太容易被這些東西佔據.而事實今天我們所講的第二樣的罪,如果英文來講就是Lust,就是情慾的罪或者色慾的罪..
這一方面,事實上我們要明白今天我們根本就是活在一個這樣的世界當中.我們現在活在一個世界,好像那些色情的文化是無所不在的..
現在我不知道是不是了,就是十多年前有人做過統計,就是在網絡那裡去尋求,就是搜索,網絡搜索,譬如Google去搜索一些東西的時候,最高第一排名的,你知不知道搜索什麼呢?.
就是搜索一些色情片,色情片段,色情片,或者看一些色情的東西..
這個是整個網絡差不多用了一半以上去做這些事情的,你想一下,網絡很多資訊不斷地傳來傳去,原來有很多是用來做一些Bandwidth,但是那時候做的可能沒有串流平台那麼多,.
沒有很多看Netflix,可能現在Netflix是No.1來的,但是姑勿論如何,就是這些的網絡的傳播,其實也代表了很多色情的事業,好像雨後春筍那樣..
因為以前很麻煩的,我記得我小時候,這些男孩子讀中學的時候,要看色情的東西要偷偷摸摸的,.
又要去買一本黃色雜誌,然後就躲起來,不想被人找出來,被老師沒收,躲在一起,當然我就沒有,因為我那時候信了耶穌的時候,但是我也知道有這些..
但是很麻煩的,有了Internet之後,特別是網絡的時候,現在很多人都不用,我記得有個姐妹她說,她說她人生當中發生一些事情是很令她驚訝的,.
或者很令她驚訝之餘很尷尬的,有一次就是她的兒子,她覺得兒子很乖,很純品,但是不知道為什麼去到teenager,.
有一次她就衝進去她的房間,不知道找兒子問什麼,一衝進去,她就看到那些三級的東西,嚇到她都不知道怎麼反應,.
她說究竟是馬上衝出來還是要學子去,她從來沒有想過會這樣發生,而且是因為進來家裡,.
其實我們現在就是活在一個這樣的時代當中,色情的刺激其實不單止是,有時候是你主動去找的,其實是推銷給你的,.
不知道你有沒有試過有些郵件或者一些網站,你去過之後不小心點了進去,點了進去之後,你發覺那些是不正經的東西,.
你馬上關了它,自此之後不斷會有這些色情的東西走出來,有時候你很尷尬,有一次開會的時候,有個同事的電話有個短訊響,.
隨手拿來看,一卡就彈了色情的廣告出來,其他同事在看開會,尷尬到整個臉紅了,可能其實他都不是故意的,.
但是他不斷地來推銷給你,現在我們就是活在一個這樣的潮流文化當中,其實是病態,當神做人的時候有各樣的需求,.
例如你肚餓的需求是絕對需要的,你想一下如果一個人不肚餓的時候,有些人有厭食症就是因為他沒有肚餓的需求,甚至將自己餓死,.
神創造人的時候,也有那個要生養眾多,遍滿全地,那時候其實有延續下一代的需求,這樣才有性慾望,.

$^{41}$如果不是,你想一下,每個人都沒有性慾望,未必會有男女生孩子,可能這個世界人口少了很多,.
事實上現在西方文化已經面對一個問題,就是出生率不斷下降,你想一下,如果要這個世界的人口維持的時候,.
每個家庭要生幾個小孩,維持而已,我不說多不說少,每個家庭要生幾個小孩,起碼都要有兩個,因為兩個人組成一個家庭,.
你要維持的時候,你下一代起碼都要有兩個,所以現在西方面對一個問題,很多家庭沒有小孩,或者有的只有一個,.
所以人口不斷地少,所以本來神造人的時候,人有一個男女結合,生兒育女去傳承下一代,.
這個是神造人有性慾需求的其中一個很正面的需要,但問題是我們現在活在一個世界,.
將所有神正面擺在我們裡面的慾望都扭曲了,你想像下一個情況,譬如吃東西肚餓,.
吃肉本身是一個神擺在你裡面的一個需求,但如果你是一個正常的人,譬如你每天有足夠的食物的時候,.
你會不會在街上見到人,譬如在超市買了一塊豬肉,在街上走的時候,你眼金金地看著人家的豬肉,.
看得眼睛都不眨眼,這樣說很誇張,或者你去到超市,擺了很多型型色色的豬肉牛肉出來,.
眼金金地看著雖然沒有錢買,但眼金金地看著人家的豬肉,會不會這樣,當然不會,但你想一下你有吃肉的,.
豬肉這麼噁心,或者牛扒,你有吃肉的,為什麼不會看著牛扒這樣看東西,為什麼一個男人看著一個女人,.
或者一個女人看著一個大隻佬,會色迷迷的這樣看,因為其實就是扭曲了,而且整個潮流文化都是這樣,.
還有電視整個節目都要推銷,要你用一些色迷迷的眼光去看另外一個對象,其實就是在推銷,這個是完全扭曲的,.
本來神所設計的慾望所有每一樣,如果擺在對的位置,其實是對我們絕對是好事,所以這個和禁慾是完全沒有關係的,.
禁慾的意思在以前中世紀的時候有些提出可能入修道院,禁慾那些,在宗教改革之後,我們基督教信仰其實就將這個觀念打破,.
所謂基督教的傳道人,牧者,我們會結婚的,例如天主教的神父,修士,他們就不結婚的,因為他們就提醒,.
因為性慾本身他覺得是一個不是那麼應該,一個神職人員不應該有的東西,或者要限制的東西,不是不應該有,是要限制的,所以他們就不結婚,.
但是我們基督教信仰就將性慾本身比較正面去看的,我們不覺得是完全要禁止的,只是要把它擺在一個對的位置,.
那怎麼擺在一個對的位置呢?首先你要明白他怎麼擺錯,究竟情慾本身,Lust,那個本質究竟是什麼一回事呢?.
那個本質其實是將人降格為物,將人物化,什麼為之將人降格為物呢?.
因為在神所創造的,特別是男女之間的慾望本身,就和人和食物之間的慾望,雖然大家的慾望是不同的,你有沒有想過?.
是不同的,我不說現在有什麼男和男,女和女,又可以變男又可以變女,那些我不說了,現在不是這個時間說這些東西,.
但正常來說,男女之間的慾望本身,和一般其他的慾望有什麼不同呢?你有沒有想過?.
你有很多其他的慾望,譬如你吃肉都是其中一樣,你想吃東西的慾望,或者懶惰都是一個慾望,後面我們都會說其中一個罪是貪食的罪,其中一個罪是懶惰的罪,.
但是男女之間的情慾的罪,和其他的慾望有什麼不同呢?你有沒有想過?.
其中一樣東西不同,就是男女之間的情慾,正常的那種,或者神設計的那種,不是看另一個對方是一個物件來滿足自己的,是互相有感情的,.
我們很直覺,或者很簡單你都會明白,我們說一個人,一個男和一個女,譬如他們有愛情,然後組織家庭,他們的慾望的滿足是有基本的愛情的基礎,.
我的意思是,我不是看你為一個發洩的工具來利用你的,而是大家彼此之間是有那個人與人之間那種互相彼此的看對方是一個人來的,你明白我的意思嗎?.
其實當人找不到他的滿足,找不到另一個人,或者他的滿足是在那些情慾上發洩的時候,其實就是將另一個對方變成了一個滿足自己的物質,滿足自己的存在,.
他再不是看對方是一個人來的,是一個發洩洩慾的工具來的,所以現在在AI時代,或者科技越來越發達的時候,所以有些人色情那些東西都不需要找人了,.
因為如果要發洩的時候,如果有相對同樣Quality,即同樣有品質的,可以滿足到他的性慾的,可能用機械人,或者什麼的,所以現在其實那些色情的事業,其實可能他們遲些也有AI時代可能也要轉型也不奇怪,.
但是簡單來說就是將人物格化,所以色慾本身的問題不是那個慾念本身,而是核心扭曲了,人本來有一個渴望是被愛的,人有一個渴望是跟人親密的,人本來有一個渴望是連結的,.
這個是男女之間的那個慾望本身背後的一個神所創造的基礎的底色來的,渴望親密,渴望被愛,渴望結連,.
但是這一切扭曲了你所渴望的親密,你所渴望的被愛,你所渴望的結連都沒有了,變成了一個只是滿足你的情慾上的一個需要,另外一個就變成了一個物,.
這種的關係是不會持久的,而扭曲了之後,人就將他人變成了滿足自己的工具,這個就是現在我們所處的一個情慾文化當中一個很現實也是很普遍的一個現象,.
所以有時候有很多年輕人他們覺得無所謂了,一夜情也好怎樣也好,我記得我初初來美國讀神學院的時候,有一次跟一個神學院的同學去聊天,他也是中國人來的,.
他說他初初來美國的時候,他來了美國久了很多,因為他讀大學的時候就來了美國,他說有一個很大的文化衝擊,他初初來美國的時候讀大學,.
他說有一個很大的文化衝擊,文化衝擊是什麼呢?他說那些人住大學的宿舍,那些男女關係很暖和,文化衝擊就是說他們竟然當一夜情,他們覺得沒什麼的,.
就好像打一場球一樣,就好像做運動一樣,其實背後就是因為他沒有了那種親密愛和連結,只是變成了一種滿足自己,將人變成了滿足自己的工具,.

$^{81}$所以耶穌的診斷是絕對正確的,問題本身不是那個慾望,是你的心,所以耶穌的診斷就說你的心是怎樣呢?.
指使我告訴你們,凡看見婦女動淫念的,這人心裡已經與他犯奸淫了.你說耶穌的標準很高,沒錯耶穌的標準是很高,但是標準之所以高,背後的原因就是耶穌看準,問題就是你的內心,.
你的內心扭曲了的時候,你看那些婦女為什麼會動淫念呢?就是因為你看他們,你不需要認識他們,你不需要跟他們有結連,你不需要跟他們有什麼親密關係,你只是看他們怎樣可以滿足你的慾望,.
那你就會色迷迷地動淫念,所以這個是那個心的問題,而這種問題其實反映到背後,其實在現代人,剛才所說的,充滿了整個潮流文化的時候,反映到現代人其實是有更深層的問題,.
在那個情慾的罪的背後,第一個就是孤獨,孤單,孤單本身是將很多人好像推向了這種在情慾當中來渴望連結的,第二個就是焦慮,.
人對於生命那種焦慮,或者那種沒有恐怖版的,當這種色情的發洩之後,好像就減輕了一些焦慮,但是那些都是很假的,或者都是很短暫的,真實的焦慮你是藉著一個真實和另外一個人的愛與被愛的關係當中,才可以真正的來舒緩,.
而第三個就是自我價值的破碎,當人找不到那個他覺得自己沒有價值的時候,可能他在那個生活,在那個男女關係都會很放縱,因為他覺得自己沒有價值,.
但是唯有在這些關係當中不斷的打滾的時候,很感覺到被需要,很感覺到我很滿足,很感覺到我的存在是有真實感,而另外一種就是控制慾,.
就是在一個這樣的關係當中的時候,你好像表面上我可以控制到,我想怎樣就怎樣,對自我的控制,對別人的控制,最後一種就是剛才所講的,就是習慣或者成癮,.
是無法自拔的,走不出來的,所以這一切不只是說我荷爾蒙分泌,令到我想有一些肉體上的滿足這麼簡單,其實背後那個真實我們內在那個屬靈狀況,我們的心如何呢,其實是影響的,.
而這一切的根源帶來很嚴重的後果,後果就是首先你對於屬靈的事是遲鈍的,你對於敬拜神,對神的追求的興趣可能都沒有,這個是很真實的,.
我曾經輔導過一個學生,我以前做過一批中學生的主要學老師,有一個學生原本他的信仰是很熱心的,很有心追求,後來我們都不知道什麼原因,.
每個禮拜都會見到他,但是你是知道的,你發覺他的人有些變了,譬如聽聖經那樣東西,以前很有追求,很發問,很熱心,現在好像都沒有什麼心,很嘔,你知道我說什麼,後來我們一直都不知道什麼原因,.
我跟另一個主要學老師,我們兩個人一起拍攝,都發覺有同樣的問題,都不知道什麼原因,為他祈禱,直至到後來那個學生有一次主動約我跟我聊天,我跟他單對單聊天,.
但是他脫離不了,尤其是他的女朋友都未信主,覺得這些無所謂的,他一直都脫離不了,所以他的心靈裡面就好像現在蒙了一個整個黑罩罩住了,.
所以他說我對神的各樣的興趣,各樣的東西,現在完全好像關了,完全關了,他跟我分享,他其實想離開這個罪,背後的那些就不說了,.
這個是很真實,而是他自己個人他知道,就算旁觀者我都看得很清楚,第二個就是那個關係的破壞,當你追求一種情慾將人控制的時候,人與人之間的關係是用了一些假的東西來替代一個真的東西,.
所以你的關係是會破壞的,對於現在其實在西方社會其中一個問題,剛才就說到出生率下降,其實不單止出生率下降,是有慾望去,這個是最近的一個統計我看過的,.
特別是那些男生的統計,主要是集中做男生的,有慾望去追女生,是有史以來是最低的,在現在這個時代,為什麼呢?有很多各樣的因素,我不說了,.
可能有些媽媽,如果有些兒子是單身的,你都會擔心,為什麼兒子想拍拖的慾望這麼低呢?現在是成為了一個普遍的社會現象,.
但我覺得可能,我不是說一定會發生,但在整個社會的普遍現象,可能整個社會的普遍現象就是很多假的東西取代了真實的關係,.
你上網找個聊天機械人聊天,很談心,可能他做得比你女朋友還要好,那還用拍拖嗎?更加就是你對情慾上的需要的時候,可能就看色情的東西就搞定了,諸如此類,.
假的東西取代了真實的東西,人與人之間很難建立真實的關係,而自我形象那種扭曲,剛才也說過,.
其實當你在情慾上放縱的時候,譬如有些人很隨便的關係的時候,很容易會跌進一個陷阱,就是他不會覺得自己的身體很有價值,.
就是很寶貴,有些東西你覺得很寶貴,我不想在婚姻之外跟別人污染了,.
但是如果這樣的關係,其實也會令人對自我形象扭曲,有很多心理輔導的個案都看到,.
有些人,譬如男女關係很混亂的時候,其實背後心靈裡有些人覺得自己沒有價值,不覺得需要尊重,不覺得需要受尊重,.
這種是很有關係的,而問題就是說,當這些男女關係失序的時候,人的自我存在感,藉著這些混亂的關係找自我存在感的時候,.
他的自我價值是不斷地降低的,他好像吸毒一樣,你覺得在這些男女關係找到滿足,或者覺得別人很重視你,.
覺得別人得到你所需要的關注的時候,但是你這種關注是假的,你越來越吸毒,其實背後心靈裡越來越輕賤自己,.
良心的麻木,剛才也說了一點,剛才說中學生輔導的個案,他說現在對於很多事情都麻木了,對於犯罪以前很敏感,.
現在對於犯罪都麻木,也對於真愛失去了耐心,你知道當你與人之間建立一個真愛的關係,是很麻煩的,.
當你與人之間建立一個真愛的關係,是很麻煩的,當你與人之間建立一個真愛的關係,是很麻煩的,.
這些一定是即食麵,不是即食麵,即食麵是吃不了的,但是當你有那些假的東西來替代的時候,你根本不想去真,你沒有耐心,.
男女單身配對的軟件,那些Dating App其中有一個指數去測試那個Dating是否可以成功,.
因為可能某個程度如果成功可能收到一些佣金,我不知道,但總之它有一個指數去測試怎樣成功,那個測試怎樣成功,就是如果他們Date三次,會否對方覺得可以繼續發展關係,.
但事實上來說,是否真實可否發展關係是否在三次裡面可以決定,你想想可能第一次就介紹一下你的名字,因為他們兩個是陌生人,.
其實因為現在很講求很快就要見效到一些東西,像打針那樣,那樣的時候對於真愛也會失去了耐心,而你沒有這些的時候,整個潮流文化不單止是在情慾那方面我們將人捆綁,.
甚至乎我們怎樣去認識建立真實的愛,剛才所講的那個背後是連結,是親密,是感受到愛的關係,這些所有都變了假的東西來取代,.
當這些假的東西取代的時候,必然會發展出那種對於情慾上,整個世界都是這樣,現在每個人都是這樣,都無所謂,就是這樣..

$^{121}$但聖經教導我們怎樣呢?聖經教導我們要尊重和學習追求一個聖潔的生活,特別是男女關係上的聖潔,神的旨意就是要你們成為聖潔,遠避淫行,要你們各人懂得怎樣用聖潔尊貴,守著自己的身體,不放縱私欲的邪情,像那不認識神的外邦人一樣..
楊保羅他很清楚的知道,我們基督徒是需要聖潔,但在這個聖潔當中,我們是特別對於情慾上的罪,特別要小心聖潔,因為這種罪本身是會不斷蔓延的,.
而我們要用聖潔尊貴守著自己的身體的時候,這個就是和那些不信神的人有很大的分別,所以他說不至於像那些不認識神的外邦人..
所以對於保羅來說,因為保羅活在的時代和我們現在的時代是一樣的,保羅活在的羅馬時代就是出了名的放縱,開派對,男女關係,很多那些雕像,或者你聽那些希臘神話,羅馬神話,其實羅馬神話是跟希臘神話的,.
就算那些神怎樣都勾引女人,關係都是這樣很亂的,保羅的時代和我們的時代是一樣的,但是保羅很清楚的知道,我們基督徒用聖潔來守著自己的身體的時候,.
但是保羅就說我們是有,因為我們是尊重,我們對於人與人之間真實的愛,生命的聯繫,我們看為很神聖,很尊重,如果我們沒有一個這樣的看法,對於人與人之間,特別男女之間的生命的聯繫是神聖,值得尊重的時候,.
這個是我們完全失去了神所設計的男女相交的慾望本身原先的計劃,第二點,我們要明白根本我們不是商品,我們是被買熟的人,.
而且你們不是自己的人,因為你們是用重價買來的,所以要在你們的身上榮耀神,保羅在哥林多前書就說得很清楚,就是說我們的身體是屬於聖靈住在我們裡面的電,你不要弄髒它,聖經也說一切的罪都是在身體之外的,.
但是犯奸淫是佔用自己的身體的時候,是佔用了聖靈住在我們裡面,我們的身體不是一個像啞活著的商品,但是我們是屬於神的,我們是需要為神來守住,你佔污了神的電的時候,你想一下,聖靈是住在你裡面的..
我們如果在這些罪上,如果是軟弱失敗的時候,聖經是教導我們,我們可以求神更新我們的內心,我們不需要被這些羞恥感來捆綁著我們,求神可以為我們做一個清潔的心,使我們重新有一個正直的靈..
這件事在世界當中是不容易做的,因為剛才所說,這個潮流文化,所有東西你都不想無端端點擊一下一個網站,不斷走一些色情的東西來,有時候男生就是這樣,上了癮都不知道了,上了癮都不由自主了,所以第一部分,我們要重新命名我們的學務,我們要想他們真正學務的是什麼呢?.
我們用一些假的東西來滿足,但你都要明白,你真正的學務不是那些假的東西,你真正的學務是和別人連結,你真正的學務是有一種親密的感覺,你真正的學務是那種被愛和愛,這些全部都不能夠從一個情慾的關係那裡得到滿足的..
你要有清潔的關係當中才能夠建立得到,那些多夫多妻的廢話,因為一個男與女的關係,你想一下,一方面你說我很愛你,不過我外面又有三奶,四奶,五奶..
你就知道根本是不可行的,我很愛你,不過有時候跟你在一起都挺窒息的,讓我一個星期有一天出去透透氣,那你就知道是不可行的,所以你要明白自己想學務的是什麼,而你學務的東西根本是不能夠從這些情慾那裡得到滿足的..
要從你的清潔那裡守住.第二就是劃定界線了,劃定界線是什麼意思呢?現在剛才所說的,根本什麼都會出現的,這個世界根本就是用這些來誘惑你..
有些網站,根本你是不應該再去的,或者我舉個例子,有一個傳道人,他這麼說,他跟一個他的禱告伴侶,兩個都是男的,他們兩個都是結了婚的,年輕傳道人..
他們有些事情要出外去,一個人去外面的時候,他們約好了他們出差的時候,每逢八點鐘就一定會打電話給對方,一起祈禱..
他們就會互相約定,問一個問題,有沒有女人在你的房間?因為可能很容易,你出差的時候,所以他們是定了一些界線,因為很多陷阱你定了界線,有些地方你不去,有些事情你不看..
是無所謂的,你說我小心一點就可以了,其實你不能小心的,你定了界線,你不去是最好的.對於色情這些罪,聖經有提過一個最明顯的教導是怎樣呢?就是要逃避的..
本身你說要鍛鍊自己的控制能力是不行的,是要逃避的.譬如有些男生女生,他們沒有必要的時候,他們可能還單獨走去一些單獨的地方,單獨的時間,.
但你不去,去的時候就很容易出事的,所以你要為自己去劃清這些界線,你知道那裡很大機會成為一個很大的試探的時候,你逃避是你最好的一個避免的方法,而逃避就是你首先劃清那些紅線,不要踩過去..
最後就是要走出隱藏,有時候可能我們跌了進去這些罪的時候,譬如有些人已經上癮了,或者看了些什麼,或者跌了進去這些罪的時候,你靠自己想改來改去都改不了,唯一這個是怎樣做呢?.
就是要找肝召,那些蟑螂在黑裡面,你怎麼打都打不了,你一召下去,你終於就可以搞定,找肝召的意思就是,譬如有些人他上癮了,舉例來說,他可能是要找輔導,.
你跟那些信任的人分享,一起來為他代討,起碼你是要給他肝召,你不是自己在黑暗當中去糾纏,剛才我所說的那個中學弟兄就是這樣,他自己在黑暗當中糾纏了很長時間搞不定,於是他跟別人分享的時候,那是他走出來的第一步,慢慢他就脫離了這個罪..
求神幫助我們,整體來說,現在我們是活在一個將情慾扭曲的時代的時候,我們究竟想要些什麼呢?我們的心有沒有靜呢?求神幫助我們..
\newpage



\section{創世記 21:13}
\label{sec:031695}
\textbf{聖經人聖經事:因他是你所生的}~~~~黃克勤牧師博士
\newline
\newline
連結: \href{http://www.church.com.hk/acms/content.asp?site=cdc\&op=show\&type=product\&code=031695\&layout=sermon}{\texttt{031695}} ~~~~ 語音日期: 2026-02-13
\newline
\newline
\hyperref[sec:031694]{\small{< < < PREV SERMON < < <}}
~
\hyperref[sec:index_chronic]{\small{[返順時目]}}
~
\hyperref[sec:031696]{\small{> > > NEXT SERMON > > >}}
\newline
\newline
創世記 21:13
\newline
\begin{longtable}{cl}
\hline
\hline
章節 & 經文 (和合本修訂版)\\
\hline
21:13 & \begin{tabularx}{0.7\textwidth}{X} 至於使女的兒子，我也必使他成為一國，因為他是你的後裔。」 \end{tabularx} \\ \\
[1ex]
\hline
\hline
\end{longtable}
$^{1}$第一節會平安.
眨眼又一日了.
生命樹又為大家提供聖經人,聖經事.
希望看到人與事的背後.
上帝是做些什麼.
神是怎樣的.
祂的想法是什麼.
「因祂是你所生的」.
這句話是上帝說的.
因為祂是你所生的.
這個也很重要.
我們先看一段聖經.
《創世記》二十一章.
「當時撒拉罕見哀琴人夏甲.
給亞伯拉罕所生的兒子.
氣笑對亞伯拉罕說.
『你把這侍女和她兒子趕出去.
因為這侍女的兒子.
不可與我的兒子以撒.
一同承受產業』」.
留意多一次.
你會看到黑色字,紅色字.
是聖經畫本的字.
綠色字是我的注釋.
留意多看一次.
當時撒拉罕見.
已經叫做撒拉了.
不是撒萊了.
當時撒拉罕見哀琴人夏甲.
即是她的婢女.
是撒拉的婢女.
「給亞伯拉罕所生的兒子氣笑對亞伯拉罕說.
『你把這侍女和她兒子趕出去.
因為這侍女的兒子.
不可與我的兒子以撒.
一同承受產業』」.
大家記得這聖經是怎樣的呢?.
就是神跟亞伯拉罕說.
「你會生她兒子」.
亞伯拉罕覺得他已經九十九,一百歲了.

$^{41}$怎會生得到呢?.
我太太撒拉已經老了.
起碼九十歲.
已經絕了,不能生了.
不過,既然神祝福.
她的「符福」就收到了.
不過,之後你看聖經就做錯一件事了.
撒拉見到要生兒子.
怎會生得到呢?.
這是撒拉出事的.
她以前的名字叫撒萊.
撒萊是公主.
撒拉是多人之母.
有辦法.
就是他們常用人的方法.
昨天說的題目是「人的方法」.
今天要看的又是「人的方法」.
就是出事在撒拉身上.
亞伯拉罕有沒有份呢?.
都有份,因為他接受這個方法.
撒拉說.
「你一百,我九十,怎會生呢?」.
「不過神說要生,不如這樣吧」.
「不要讓事情不發生」.
「不能生的,老公,那怎辦呢?」.
「我有一個很近身的婢女,多好啊」.
「比較黑的,是埃及人,叫夏格」.
夏格的意思是逃跑.
不知道是不是逃跑出來的.
被她買了做奴隸.
以來服侍她.
「我有一個我都很疼的婢女」.
「不如讓你做妾侍吧」.
「她那麼強壯,一定生得到」.
「又年輕」.
亞伯拉罕說「這樣好啊」.
其實撒拉先出事.
不過之後亞伯拉罕肯.
事情就發生了.
不過不是上帝給的.

$^{81}$因為兩個人在那個時間.
都不信上帝能夠讓一個.
一百歲和九十歲的人生兒子.
不行了,所以就用人的方法.
又是用人的方法,好像昨天那樣.
於是亞伯拉罕和夏格生了兒子.
生了兒子聖經就沒有在這段.
給亞伯拉罕所生的兒子棄孝.
就沒有說兒子的名字.
以為大家都認識.
我怕大家不記得.
亞伯拉罕和夏格的兒子的名字叫以實瑪利.
以實瑪利大家可能記得了.
以實瑪利是個好名字.
以實瑪利呢.
我不夠位置.
Ishmael.
你一見到L.
EL.
L是God.
Ishmael.
是神會知道.
神會了解.
其實也是好名字.
Ishmael.
也是好名字.
棄孝.
Jakob.
他的時式是PL.
PL就令到他這個字的解釋是嘲笑.
或者開玩笑.
如果是Kel.
即是另一個最常用的動詞.
的勢態就是笑.
但如果是PL.
這個Jakob就變成嘲笑.
或者開玩笑.
撒拉見到夏格的兒子以實瑪利棄孝.
小朋友笑會不會?.
當然會,不會有小朋友自己傻笑.

$^{121}$但這個樣子極有可能觸怒了撒拉.
其實沒什麼.
你怎知道他在說什麼.
不過棄孝是嘲笑.
極有可能夏格的兒子以實瑪利.
和Isha小朋友.
在家裡玩的時候.
可能Isha就笑Isha.
你這麼小,什麼都不懂.
你這麼瘦,好像猴子一樣.
瘦猴子,我不同我多大.
可能會笑.
媽媽撒拉見到很不開心.
你這個婢女的兒子笑我兒子?.
她又忘記了.
本來這個以實瑪利的出現.
是本來她叫阿巴拉罕.
和自己的婢女所生的.
本來以為家產是由她承繼的.
幫神生的兒子.
誰知道神說這個不算.
你阿巴拉罕和撒拉所生的我承繼.
哎呀,那怎麼辦?.
誰知道真的生了.
以實瑪利出來.
以實瑪利是大.
可能大他一兩歲.
可能這個埃及血統比較大.
不知怎樣.
大的哥哥笑弟弟.
不知什麼事.
不開心,棄笑.
於是撒拉和阿巴拉相當.
你把這個婢女趕出去.
這麼沒人情味.
後來她真的趕出去.
因為給水給幾個餅.
她就趕走了.
他們走了一段路.
到一個沒有路標的地方.

$^{161}$兩人打算死.
但神看到.
帶他們去一個水井救他們.
把他們趕出去.
很沒人情味.
撒拉用她的方法.
很沒人情味.
用方法令阿巴拉罕生兒子.
本來好端端.
現在自己生了.
她就不喜歡這個.
趕他出去.
笑我兒子.
因為這婢女的兒子.
不可以與我的兒子.
二十一同承受產業.
你把這個婢女和兒子趕出去.
因為她的兒子.
笑我的兒子.
這樣也有理由.
不過趕不到人走.
你教他吧.
你是主母.
但她現在用一個很重的辦法.
二十才是真正的兒子.
這個不是.
她不能夠和我們一起承受.
她不能夠承受我們的產業.
她不能夠和我的兒子.
二十分家產.
這樣.
很絕.
很惡.
阿巴拉罕因她兒子的緣故.
很憂愁.
這個兒子是誰.
二十瑪利.
怎麼辦.
趕他走.
不就是叫他死.

$^{201}$下甲都一場夫妻.
我妾來.
不就是叫他死.
很憂愁.
神對阿巴拉罕說.
其實神主動來.
跟他說話.
不用憂愁.
我心裡馬上很舒服.
為什麼呢.
原來我們的上帝.
一直看著我們發生每一件事.
看著.
全部知道.
跟阿巴拉罕說.
你不必為這童子和你的侍女.
下甲.
童子二十瑪利.
你不用為.
小孩子.
童子.
你不用為二十瑪利.
和下甲憂愁.
就她.
撒拉夏多國之母.
多人之母.
對你所說的話.
你都該聽從.
都該聽從.
應該聽的.
因為從以撒生日.
才稱為你的後裔.
撒拉沒有說錯.
沒有說錯.
對的.
不過我覺得不好.
沒有錯.
不過不好.
你都該聽.
因為.

$^{241}$由以撒生日才是你的後裔.
所以以撒開始才是你的後裔.
如果出去.
我會照顧她.
至於你的侍女的兒子.
我都必使他後裔成立一國.
沒事的.
不會死的.
我都會令到他變成一個國家.
國家會不會是一個人呢.
一個國家.
都不少人.
都有力量.
有勢力.
我會令到他後裔成立一國.
為什麼.
因為紅色字.
他是你所生的.
他都是你所生的.
嘩.
很感恩.
讚美上帝.
是什麼呢.
就是雖然他不是我.
令到你一百歲九十歲生的那個以撒.
但是由你所生的.
其實是你太太做錯了.
你都錯一半.
又不信我.
不過既然由你所生.
我都認數.
所以我會令到他成為一國.
因為他是你所生的.
我都會紀念.
嘩.
這位上帝真偉大.
讚美他.
守他的信用.
由你所生的.
好.

$^{281}$我都認.
雖然他不能夠跟著你成立一國.
照我的計劃.
不過我都會令到他成為一個國.
以實瑪利.
其實這個國大不大.
很大.
以實瑪利的後裔.
就是今天的阿拉伯人.
以色列人.
以撒雅各的後裔.
就是以色列人.
就是阿巴拉罕和以撒雅各.
下來的以色列人.
阿巴拉罕和下甲以實瑪利.
下來的阿拉伯人.
其實都是阿巴拉罕做祖先.
一個是希伯來人.
一個是阿拉伯人.
兩個後來變成世仇.
直到今天.
打了不少仗.
而這個阿拉伯人.
今天多數是信奉伊斯蘭的.
變了兩個.
和猶太教有很多的爭執.
很多的爭持.
變成敵對.
你會見到.
其實神說好.
我都會紀念.
因此你會記得.
一件事做了.
神第一.
他是看到他家的需要.
他安慰他.
幫助他.
和我都認.
我都承認.
我都跟手.

$^{321}$不過.
其實這件事不好.
因此這個.
其實都是人的方法.
而做出來的後遺症.
不過神都說.
既然如此.
我都認訴.
因為他是你所生的.
今句.
至於事類的兒子.
我必使他的後裔成立一國.
因為他是你所生的.
我見到這一句.
我心裡很感恩.
就是上帝說過的話.
他會承諾.
應許是實在的.
雖然不是他原先的.
他知道.
不過不是他原先的計劃.
既然如此.
你做了.
神都會在這裡.
做一些事.
幫你補救.
只要你繼續懂得跟隨他就好了.
如果起初的時候.
有信心就不會有下甲和以實瑪利.
這件事好多了.
到今天.
又不會出現阿拉伯人.
大家對抗對抗.
你說多好呢.
反而成為一個禍患.
亦都看到.
上帝在當中守他.
說話的承諾.
非常之感恩.
讚美主.

$^{361}$我們祈禱吧.
主我們看到.
神是祝福以實瑪利.
讓他成為一國.
主要原因就是.
神說你這個小孩.
都是由你所生的.
所以我都會看著他.
主我們很多時候.
自己是做得不好.
信心不足.
做錯了事.
或者做得不好.
不過神你是看到的.
你仍然會幫助.
會補救.
會帶領.
我們讚美你.
多謝主.
每一個時間的應許.
是多麼的確實.
你說的話.
你自己是形象的.
多謝主.
願你得著榮耀.
讚美你.
禱告奉耶穌名求.
阿們.
\newpage



\section{傳道書 1:2-7}
\label{sec:031696}
\textbf{它,他/她,祂}~~~~周朗先生
\newline
\newline
連結: \href{http://www.church.com.hk/acms/content.asp?site=cdc\&op=show\&type=product\&code=031696\&layout=sermon}{\texttt{031696}} ~~~~ 語音日期: 2026-02-22
\newline
\newline
\hyperref[sec:031695]{\small{< < < PREV SERMON < < <}}
~
\hyperref[sec:index_chronic]{\small{[返順時目]}}
~
\hyperref[sec:031697]{\small{> > > NEXT SERMON > > >}}
\newline
\newline
傳道書 1:2-7
\newline
\begin{longtable}{cl}
\hline
\hline
章節 & 經文 (和合本修訂版)\\
\hline
1:2 & \begin{tabularx}{0.7\textwidth}{X} 傳道者說：虛空的虛空，虛空的虛空，全是虛空。 \end{tabularx} \\ \\ \relax
1:3 & \begin{tabularx}{0.7\textwidth}{X} 人一切的勞碌，就是他在日光之下的勞碌，有甚麼益處呢？ \end{tabularx} \\ \\ \relax
1:4 & \begin{tabularx}{0.7\textwidth}{X} 一代過去，一代又來，地卻永遠長存。 \end{tabularx} \\ \\ \relax
1:5 & \begin{tabularx}{0.7\textwidth}{X} 太陽上升，太陽落下，急歸所出之地。 \end{tabularx} \\ \\ \relax
1:6 & \begin{tabularx}{0.7\textwidth}{X} 風往南颳，又向北轉，不停旋轉，繞回原路。 \end{tabularx} \\ \\ \relax
1:7 & \begin{tabularx}{0.7\textwidth}{X} 江河都往海裡流，海卻不滿；江河從何處流，仍歸回原處。 \end{tabularx} \\ \\
[1ex]
\hline
\hline
\end{longtable}
$^{1}$早晨 新年平安.
各位傳人堂的弟兄姊妹.
很開心 其實我答應了江牧師之後.
我答應了寄了在這本簿裡.
但我沒有留意這是甚麼日子.
今天是年初幾 年初六.
我到嫂子準備交資料的時候.
我就想起原來這是大家過年的第一次返聚會.
所以我就覺得開始有少少壓力.
所以我就特別花多少少時間準備了一些東西.
東西包括在稍後分享的資料裡.
還有一些實際的東西.
稍後會透過利是封.
稍後都會分派給每一位弟兄姊妹.
希望你們在新的一年有祝福.
新的一年大家剛剛休息足夠嗎.
見你想見的人見到嗎.
見的時候開心多些還是相反多些呢.
會不會是不見就好好的.
見到就有些特別.
不見的時候就活在想像中間.
但見到的時候就反而見到.
廣東話叫做眼淵.
有沒有聽到一些說話是比較性的.
即是說誰就這樣.
有沒有聽到這些說話.
你有沒有份發出這些說話.
我希望我們在教會的模樣之下.
我們會盡量減少這些.
今天想和大家分享的題目.
那裡有四個字.
我們一起讀出來那四個字好不好.
1,2,3.
他,他,他,他.
這個題目是不是很有趣.
原來我在這幾天做準備的時候.
原來我在大概十幾年前.
我在同一座大廈.
你們這一座大廈的三樓.
我分享過另外一個題目.

$^{41}$當時是一間院校.
原來是四個我.
我經常為了這些代名詞.
我呀,你呀,他呀,這些經常思考.
一會兒我再一路交代.
今天我們這四個他其實都不重要.
最重要是我們和這四個他的關係.
這裡我想帶出一樣東西.
因為我是來自於心理學和輔導的訓練.
關係是很重要的.
讓大家做一個熱身.
大家手機在附近嗎.
都碰到是吧.
我們先碰一碰手機.
不是叫你玩手機.
你碰一碰手機.
這部手機是你的.
不要拿錯人家那部.
這部手機是你買的嗎.
還是家人送給你的.
是你買的.
你自己選的.
款式是不是你自己選的.
你心愛的.
顏色是你心愛的.
你看著它.
你現在還有沒有當時那麼喜歡它.
有沒有逛過一些電話舖.
我昨晚吃飯逛過幾間電話舖.
有沒有心思思又想走過去.
又看到另外一部.
有沒有看到.
有沒有心動想換.
如果資料轉移不是問題的話.
是不是已經行動了.
所以我今天想問.
你和你的手機的關係是怎樣.
你還有沒有那麼喜歡它.
程度是多了少了還是一樣.
還是深厚了.

$^{81}$這就是我們和東西的關係.
就是我們接下來的第一部份.
就是遇上它.
一個死物它.
我們每個人心裡都為了很多不同的東西.
可能是車.
可能是物業.
可能是音響.
可能是一些衣服.
我剛才早來的時候聽到有弟兄姊妹聊天的時候.
就在討論一雙鞋.
但我就覺得討論裡面很豐富.
因為討論到不是鞋的價值.
而是鞋是用來舒服和開心.
用來襯衣服.
所以我就覺得原來那雙鞋的東西本身.
不是最重要的.
最重要的就是開心.
剛才開始前我和主席討論.
我心裡面的它.
其實我對物質就一般般.
唯獨是我這幾年就發現.
我和車的關係很密切.
我由有車牌一直都不擁有車.
基本上是不擁有車.
到我22年就買了一部二手車.
現在就去到第二部二手車.
原來我是會和車聊天的.
我是會稱讚它你今天很乖.
你每天都這麼乖就好了.
這是一些很奇怪的事.
我是自己會和它在心裡.
不會說出聲的.
但我都不介意.
有時我會告訴家人聽.
有時我在家人面前稱讚它.
因為其實再想想.
原來其實車是一個死物.
但原來這部車在這幾年裡.
為何22年要買車.

$^{121}$你想想22年是甚麼日子.
是疫情的甚麼.
是它最後準備離開.
但那時候我亦在那一年確診.
所以剛才我來的時候.
我是戴著口罩.
我一直都不脫口罩.
無論它怎樣好.
然後又有些感冒.
我都不脫.
我現在是尊重大家.
我脫了它.
一會我下去又戴回.
因為我確診的經驗不是很瘦.
我就因為這個不好瘦.
所以我就去愛惜我的家人.
我就想如果有部車我兼顧到的話.
我用部車提供一個比較衛生的環境.
不用他們在外面迫.
原來部車盛載了我和家人不同的時光.
所以原來這個死物它是一個甚麼.
是一個甚麼關係.
其實我買車的時候.
很多朋友跟我說.
你當時住大角咀.
你工作在長沙灣.
其實需不需要部車.
其實真的不需要.
其實養車的成本遠比搭的士貴.
所以我很多朋友就說.
我開頭都會聽.
但聽了我都是行動.
你們會不會.
有些東西你是知道一般般.
但你又行動.
你行動了好像知道有些不太好.
我們進入了一種忐忑.
又好像不好但又想.
有沒有這種情況.
但原來隔了一段時間.

$^{161}$我明白了.
原來我有兩個原因想有部車.
除了方便之外.
家人和原來我人生我看回.
我沒有甚麼懷疑沒有甚麼樂趣.
我沒有其他的嗜好.
原來我每一次工作完辛苦.
我上部車.
我可以有一個減壓.
這個就是我2022年到現在的經驗.
我有時晚上在家工作.
或者在辦公室工作.
困了十個小時.
原來我開車入一入油.
我會很舒暢.
這個就是我和車的關係.
今天我和大家帶了一些很特別的形式.
我希望新的一年.
第一次的訊息分享.
讓你有一個好的開始.
這個只是一個開始.
並不是內容很重要.
我希望今天透過.
有些活動你們參與.
可能有些反思.
這個就是我們東西.
人和環境的責任.
中間的關係.
經文是這樣說.
創世記二章十五節.
耶和華神把阿當帶來.
安置在伊甸園裡.
讓他耕種看管原子.
這裡就是說神揀選了我們人.
去管這個大地.
我們是比較特別的.
被揀選的.
所以在地上的很多東西.
都被我們掌管.
那麼這件事是一件美事.

$^{201}$還是一件壞事呢.
相信是一件好事.
但是我們不禁和這些東西糾纏.
所以這裡就是我們讀心理學.
或者我平時在院校教心理學.
常用的一些工具.
今天並不是一個心理學課堂.
但是我想中間有幾個東西.
有幾個理論和大家碰一碰.
讓大家感受一下.
這個就是我們心理學的其中一個.
很重要的人物叫做弗洛伊德.
弗洛伊德他在提倡潛意識.
潛意識就是說你沒有刻意去想.
但是它存在的東西.
大家有沒有做夢啊.
我試過有一天.
我發了很多訊息.
和別人有節日的祝福.
在幾百個祝福裡面.
其中一個就是我弟弟.
我有一個弟弟.
我們沒有一起住.
各自有自己的家庭.
我們平時不多聯絡.
因為他太忙.
他比我更加忙.
但是我發了一個訊息給他.
他奉旨不會回覆的.
因為他經常覺得這個哥哥.
做事比較不穩定.
他人生做了一份工作.
我人生做了數不出多少份工作.
所以他比較謹慎嚴肅.
但是竟然那天晚上.
我發夢見到他.
而他第二天竟然告訴我.
他又發夢見到我.
這麼甜蜜的.
兩個男人.

$^{241}$為什麼會有這個效果呢.
是不是那些一起出來的兄弟.
是不是有一些什麼特異功能呢.
其實在心理學很簡單.
就是那天.
他每天收千多個電郵.
他的工作.
裡面原來哥哥曾經找過他.
勾起了他下面一些.
不覺的潛意識.
惹起了之後.
晚上睡覺的時候.
就在惹起的小小東西裡.
再來一次回顧.
其實這樣就是潛意識.
但是你會發覺.
潛意識在冰山的哪裡呢.
在底部.
我想問下面浸著水的比較大.
還是上面的比較大.
下面比較大.
這樣才穩陣.
但是同樣地.
佛洛伊德就說了.
下面那個叫什麼.
我大家一起讀出.
下面那個.
本我.
最上面那個呢.
超我.
中間那個呢.
在真我這裡.
我邀請大家摸一摸臉.
會不會影響你的化妝.
中間這個就是真我.
原來我們永遠都在下面.
有些不知道什麼.
控制不了的慾望.
想法.
那個是小的.

$^{281}$小的大.
還是乖的大.
小的大.
叫他不要掣那個是超我.
超我在這裡.
超我比較小.
所以弄到誰呢.
弄到阿珍.
所以大家整天都在忐忑裡.
這個就是心理學的.
其中一個比較重要的理論.
我信了主不久.
我15年才回教會.
16年洗.
我反而在15年之前.
我就已經讀了很多這些書.
也都在市面任教相關的科目.
當時我還笑我一個老師.
我說這個老師比我年紀輕.
當時.
但他就說他是基督徒.
我就說你讀這些書這麼厲害.
為什麼會做基督徒呢.
不用信.
你靠自己都很厲害.
原來我反而跟了他讀了很多這些書之後.
我後來我05年讀心理學.
08年畢業.
我畢業之後.
反而我一步一步地走回基督的面前.
我覺得他說得對.
因為原來世人都犯了罪.
我們經常都有很多引誘試探在裡面.
在本我那裡.
但是我們如果要去.
每天去學習的時候.
我們更加提醒自己上面.
超我.
那我們就會明白.
怎樣去平衡自己.

$^{321}$做好每個禮拜每天的工作.
我們要在我們的環境.
很多引誘之下學習.
神怎樣供應給我們.
經文這樣說的.
《馬太福音》六章28,29節.
何必為衣服憂慮呢.
看清楚田野的百合花怎樣生長吧.
他們不勞動也不紡織.
但我告訴你們.
就算所羅門最榮耀時的村大.
也比不上這一些花中的一朵.
通常類似的經文都有.
上帝已經數算了每一隻小鳥.
每一隻麻雀身上的多少條羽毛.
所以我們其實.
基本上香港是很少人是.
因為凍親.
因為肚餓而影響了生活.
所以在香港.
我們應該不為這些擔心.
同不同意嗎.
好了.
所以我們要透過周圍.
很多好像很污穢的物質.
很引誘.
我們有時接觸了的時候.
我們會反思.
我是不是太貪心.
太過物質.
太過慾望呢.
但其實透過這些物質.
我們認識神的存在.
詩篇十九篇.
第一章是這樣說的.
天空屬雪神的榮耀.
窮倉傳揚他手的作為.
我當時說沒有開車.
沒有開車.
這個就是我當時.

$^{361}$這張相是我中年後拍的.
但原來我二十歲那一頭.
其實我是為了晚上進修.
我曾經有三年是開這一款車.
後來我中年的時候.
在街上遇到同一款車.
我忍不住走上去拍那張相.
原來我人生不停和車有些關係.
現在麻煩詩詩幫我去派.
今天我準備了四個利是封給大家.
現在就可以派了.
是這一款.
大家收到這個利是封之後.
裡面有一個祝福給你的.
就是一個經文.
基本上經文都不同.
全場只有兩個人會相同.
我有二十款.
我和家人一起在這幾天去整給大家.
大家現在收到.
去打開它.
我亦會有一首詩歌.
這些詩歌比較特別.
全部都是比較新派.
歌詞剛才有一張紙.
如果你們收到有歌詞.
你可以在一會播詩歌的過程裡.
間歇性我會講一些說話.
我希望剛才我分享了一些訊息.
然後你又有歌詞.
然後亦有你手上的經文.
大家可能做一些簡單的沉澱反思.
如果你能夠的.
你可以閉上眼.
你可以趁這個機會.
好像另類的靈修一樣.
我們趁這個機會和我們的天父溝通一下.
這是天父世界.
我心滿有安寧.
花草樹 瓊藏碧海.

$^{401}$訴說天父奇難.
去年新年的崇拜.
你在哪裡.
你不是在這個會場.
你在上一次的會場.
這是天父世界.
我們側耳傾聽.
天父透過執事門.
透過窮工門.
給你們找到一個新的位置.
你們在一個很美的環境.
繼續在祝福和夢想之下.
這些就是天父幫大家安排的.
這是天父世界.
在這個新的地方.
你和這個新的地方有什麼關係.
剛才我聽到窮工門.
關心廁所裡面的物資足不足夠.
每一個物資背後.
都有一些心思.
那些物資會令到你有一種福氣.
都是想你每一天活得更加好.
我最喜歡的食物是餐蛋麵.
你喜歡餐蛋麵嗎.
每一樣東西都有它的缺點.
無益.
為什麼天父容許餐蛋麵存在呢.
他提供給我的簡單的熱量.
稍微有些營養.
很方便地支持了我的工作.
當然我可以選擇更有益的食物.
不要為了餐蛋麵太內疚.
我一會兒就去吃.
有沒有為了一件很貴的衣服內疚過.
但這件衣服除了給你保溫.
去保護你.
它令到你很有自信心.
你穿著這件衣服去見人.
你去工作.
你去幫人.

$^{441}$那個效果計不到.
這件衣服在你的人生.
成為了其中一個很重要的步驟.
我相信天父給你有這樣的資源.
去買了這件衣服.
這件衣服你穿了之後.
能夠貢獻服侍的意義.
比我們付出的那個更高.
如果是這樣的這件衣服.
你就接受它吧.
而且你要穿著這件衣服.
再去做多些服侍.
祝福多些人.
我講的例子都是很不足.
你們自己思考一下.
你在世上和哪些物質.
哪些東西最糾纏.
它怎樣困擾你.
它會不會其實是一個祝福.
讓你有這個物質.
這個東西之後.
可以走得更加遠.
你信不信這個東西.
是天父透過一些天使安排給你的.
最後的心意你明白嗎.
好聽嗎這首歌.
現在這張白色的紙.
歌詞的背面有類似的畫面.
剛才我們完成了第一個部份.
是死物他.
我現在紅色寬了.
你會不會有筆.
如果沒有筆你舉手一下.
師士給你一支筆.
在裡面你可以寫出來.
你可以在心裡想一下.
剛才在整個過程裡.
無論是歌.
或者是我口頭的一些訊息的分享.
刺激到你.

$^{481}$或者是你手上的勵志峰的經文.
或者是剛才主席讀出來的一些經文.
今天不同的鋪排.
今早早餐吃的時候.
遇到一個夥計跟你聊天.
所有令你想起的.
你都可以記錄一點點.
在這一格裡面.
如果你不想著筆.
你在心中想一想.
這個就會是一個.
希望你今天.
在今年的第一個崇拜裡面.
希望透過這些反思.
對你今年有一個很好的提醒.
這個我就不限時間的.
我分享了那個做法之後.
我會繼續前進.
大家可以隨著你自己的情況去寫.
或者在心裡想一想都可以的.
現在我就繼續前進.
因為今天我的挑戰就是.
時間有限.
想說的東西太多了.
第二個就是男女仔撻.
這裡我絕對沒有機會去說很仔細.
不過這一版.
希望大家如果有興趣.
可以拍一拍那個畫面.
或者你怎樣都好.
你記住人生有八個階段.
人生八階.
如果你在手機的Google打人生八階.
按搜尋.
然後你選圖片.
你其實會見到這張圖片.
而這張圖片就是我放上去的.
即是綠色和紅色.
怎樣證明是我放上去的呢.
因為在紅色最低那格.

$^{521}$那個失望那兩個字.
我是特意打錯了字.
它原本不是這兩個字.
它英文是Despair.
中文翻譯做絕望.
但是因為我知道我放上去.
就到處都看到.
所以我特意打錯字.
好讓我有機會可以親口.
和聽這版的內容的出席者.
去解釋究竟.
它不是那個字背後.
原本那個絕望那個絕字這麼厲害.
其實只是有點失望.
今天我沒有機會去講仔細.
但是我相信大家看到這個畫面的時候.
不其然會怎樣的.
是不是按著自己的年紀.
然後看一次.
然後橫掃一下.
然後掃過去的時候.
大家看到綠色是好東西還是壞的東西.
紅色呢.
所以如果你見到的時候.
我希望你本身就處於綠色的區域.
但是如果綠色的區域失守了的話.
可能就會在旁邊紅色.
我們試一下往前.
往早年再看.
你應該總是有一個綠色的格子.
承托住你的人生.
可能是認同.
可能是勤力.
我相信很多香港人在勤力那格都撐得到.
這個理論.
黑心的老師就叫佛洛伊德.
所以剛才早前那個就是老師的理論.
現在就是學生的理論.
這裡我主要想帶出.
原來黑心的理論.

$^{561}$在我們讀心理學的人來說.
是很常見的.
我曾經有三份論文功課用過這個理論.
但裡面最重要的就是.
在頭那兩行.
裡面哪一個持份者.
哪一個人.
哪一個他.
會出現在這裡陪著他.
行頭那兩個階段.
就是父母論.
然後在三至六歲和六至十一歲.
又轉了持份者的和.
是哪些持份者進入了去陪那個人去成長.
老師.
學校那些校工.
校長.
還有同學.
開始有朋友.
好了.
十八十二十八.
十八三十五那兩格.
朋友.
開始有些.
發了一些粉紅色的夢.
可能有喜歡的人.
或者有人喜歡你.
開始有另一半的出現.
所以你看到他不停的和人扯上關係.
而去到三十五六十.
可能就再下一個周期.
可能我們有子女.
有下屬.
有學生.
或者有我們去幫的人.
可能要做一些整理.
這個理論帶出了這些.
不過我發覺坊間.
我做輔導員.
很多人都不開心.

$^{601}$其實他們搞錯了一些東西.
搞錯了什麼呢.
原來朋友.
親戚是不能選擇的.
朋友是可以選擇的.
或者同事.
無端端你當他做朋友.
那搞錯了什麼呢.
原來就搞錯了這四樣東西.
我很簡單和你說.
第一種.
叫做playmate.
是嘻哈的.
陪你吃飯.
玩.
陪你買衣服.
吹吹水.
不深入的.
大家都有這些朋友嗎.
有好過沒有.
有很重要的.
第二種就是功能性的.
公事上的.
讀書做功課的.
大家一起同一組的.
在一個團隊裡邊編制了.
我們一定要合作做好一些任務的.
那些是有具體的目的的.
這些很重要的.
我做補習的工作.
很多時候都和小朋友說.
你認識同學.
認識朋友.
是全級這麼多班都要認識的.
你不要只認識玩的那些.
你認識功課厲害的那些.
幫到你.
提醒你哪些重要不重要.
還要認識這麼多班.
因為你不知道哪個老師出卷的.

$^{641}$那些就是第二種.
但是我們很多時候不開心.
就是將頭兩種當作第三種.
我講一個秘密給你聽.
不要告訴別人.
那轉頭又怎樣.
肯定全部人都知道了.
接著又說他出賣我.
沒有的.
你自己先答錯了.
人家當你第一種.
你當他第三種.
就搞錯了.
通常在外面教書.
最難的就是第四格了.
第四格就是給了一條命.
你都可以的.
叫做Agape.
通常去到這裡.
父母都可能做到.
可能做到.
但是應該做到的.
所有你上網查的定義.
都是說神才做到的.
是無條件地付出.
所以我們都是.
你上了教會就有第四種.
如果你不信主的話.
第四種就可能很難得到.
我帶出了這些不同的朋友.
背後都是一些他.
男孩他 女孩他.
我想問.
在過去一年.
有哪一個他是最.
擾動你的呢.
可能是開心那邊的.
可能是擔心那邊的.
可能是激心那邊的.
我們對著這個人.

$^{681}$你心裡面有這個人.
我們要怎樣再跟他.
在今年新的一年裡相處呢.
我讀出.
哥林多前書十三章四至八節.
愛是寬容.
愛是恩慈.
不嫉妒.
不自吹自擂.
不自高自大.
不做不合體統的事.
不求自己的益處.
不輕易被捉弄.
不圖謀惡事.
不喜歡行不義.
而為真理歡喜.
凡是包容.
凡是相信.
凡是盼望.
凡是堅忍.
愛永不倒下.
其實剩下來的.
基本上剩下一樣東西.
就是一個字.
大家看到嗎.
愛字.
所以我們跟這些人相處的原則就是.
中國人是將心比己.
是這樣說.
在這裡金句都有說.
路加福音六章三十一節.
你們希望別人怎樣對待你們.
你們也要怎樣對待別人.
大家聽過了.
做不做到.
我呢.
二十年前的相.
所以你計到我兩個兒子.
現在幾多歲了.
我大兒子已經二十六歲了.

$^{721}$我小兒子二十三歲.
今年正式找到工作.
很感恩.
大學畢業了.
所以這裡我想說.
就有頭一個他.
死物他.
就是你會見到.
我跟食物的關係處理得不好的時候.
我會變成怎樣.
我會變成這樣.
二十年前.
幸好我認識了他跟我的關係.
我適當處理了.
我以前很喜歡喝奶茶.
加三匙糖.
早餐一杯.
下午吃飯這麼辛苦.
當然要喝.
下午茶這麼辛苦.
有人請.
當然要喝.
晚上喝多幾杯果汁.
這麼辛苦回來.
熱騰騰.
喝些凍果汁.
是不是很好.
物極必反.
過了龍.
所以當你重新找到.
跟他的關係的時候.
就可以有另外一個出路.
原來我無論這麼辛苦.
為了什麼.
可能我為了我家.
我們現在派第二個黎智峰.
我接下來會加速一點.
因為時間過得很快.
我趁派的時候.
大家知道有經驗.

$^{761}$收到就可以開了.
全部都是你自己.
可以拿走的祝福.
這個黎智峰有個福字.
這個福字是零型的.
我們有四方的朋友.
有沒有聽過.
識人好過識字.
適當地.
珍惜多一點的朋友.
可能是.
很必要的.
我其中一個兒子.
他由幼稚園開始.
到他大學畢業.
到他所有的工作.
原來都跟基督教有關.
但是起點就是.
誰介紹他讀那間幼稚園呢.
就是一個不信耶穌的人.
他說這間好.
那我就將那間.
給他介紹給他.
他就覺得.
這間很好.
他就說.
這間好.
那我就將我兒子放進去.
原來就開始了一條.
走基督教的路.
當時我們全家都不信主.
這個朋友我還有聯絡.
你呢.
有沒有一些朋友.
可能在銀行.
法律.
健康裡面給過意見你.
工作上面.
你呢.
你今年有沒有.

$^{801}$給他一個祝福的訊息.
是不是已經忘記了他.
那我先回答.
剛才我說的那個朋友.
我今年是我主動.
給了一個祝福他的.
但是他不肯見我很多年.
他是一個.
無價的一位女士.
她很漂亮.
但是我想.
她可能有些改變.
所以她這十年.
任何的場合.
她都不會給人見到他.
所以我們就在.
WhatsApp那裡做祝福.
原來他當時.
介紹了我兒子的一間學校.
那個祝福原來是這麼深遠.
所以當有朋友.
去請教你的時候.
你不要吝嗇好不好.
你懂的.
你照說.
你猜不到你這一句.
可能祝福了別人很多.
我同樣是說同一個兒子.
他原來幼稚園是這樣.
去小學都不是直接的.
都是很難才入到.
中學更加難.
大學都是.
我在另外一間不是基督教的大學.
交了那個訂金錢.
最後原來我們放棄了.
他竟然走回一間.
基督教的大學.
香港浸會大學.
結果在那裡.

$^{841}$更加又再奠定了他們一生.
那裡他在那裡.
學了很多東西.
是和他現在的工作.
和服事有關.
所以神.
他要行駛.
很多他的計劃的時候.
透過很多天使.
我剛才每一個的分岔路.
都有一些不同的天使.
給了一些意見我.
當然我要聽.
我們現在進入了這一格.
剛才在過程裡.
你想到哪一個他.
你都可以在這裡去想一想.
或者甚至乎寫下一點.
怕轉頭不記得.
我這裡快一點.
原來我們和神相遇.
其實神一直都在這裡.
因為他創造了這個天地.
他在先我們在後.
但問題是.
有多少人見過上帝.
有多少人見過耶穌.
是不容易的.
所以原來我們可能是.
透過很多不同的渠道.
很多的天使.
我們才感受到.
但無論如何.
馬可福音十章四十五節.
因為就是人子來.
也不是要受人服事.
而是要服事人.
並且要寫命.
作許多人的屬格.
祂真的來了.

$^{881}$道成肉身.
祂感受到我們感受的東西.
祂體會到我們所有物質.
人與人.
和祂和祂的父神的關係.
祂都沒有生氣過耶穌.
祂有生氣的.
祂都和祂說.
我都真是痛的.
按你心意.
祂都肯分享我的痛楚.
所以你都可以的.
但問題是.
你有沒有全心全意.
全智全力愛你的神.
馬可福音十二章三十節.
你有返崇拜.
你有沒有去小祖.
你有沒有查經.
你有沒有每日都禱告.
你有沒有靈修.
剛才主席原來.
都將我剛剛那段都說過.
不謀而合.
我想是神的安排.
所以你要學會和祂相處.
因為其實祂已經在等我們.
呼應剛才第一部份.
彼得前書四章十節.
作為神各樣恩賜.
我都要停一停.
你的恩賜是什麼.
你找到你的恩賜沒有.
插花.
煮東西吃.
關心人.
帶新朋友來.
人家有不開心去分擔.
這些都是恩賜.
我們各樣恩賜的好管家.

$^{921}$每個人都要照著.
所領受的恩賜彼此服侍.
將你的恩賜不要姑寒.
拿出來.
剛才這裡的改變.
這張照片.
看到我返教會前.
就是左邊.
我有一次返了教會.
不是很久.
應該是洗禮過期.
跟兩個兒子.
崇拜完之後和會友飲茶.
在茶樓拍的.
我特意叫他們擺這個姿勢.
有沒有分別.
原來返了教會.
跟神接觸之後.
祂會特別疼你的.
你一會兒在照片中會看到.
希望祂都疼你.
我們要派第三個風.
神透過什麼訊息去呼喚你.
除了叫你返回教會.
在你上次好險.
接不到的士的時候.
有沒有記得找一找祂.
祂有沒有跟著安排一輛的士給你.
有沒有找不到鎖匙的時候.
以為漏了.
跟著祂騎過領土.
跟著祂幫你找回.
可不可以不只是有事才找祂.
親自尋找我.
祂親自找你.
你有沒有親自找祂.
我和一個兒子.
第一份工竟然在教會.
第二份工在一間基督教的中學.
最後他現在是做自己的工作.

$^{961}$他接一些私人的工作做.
但他都是在一間基督教的音樂機構做義工.
這些會不會是我的安排.
我安排不到.
應該是神的安排.
神在你身上安排了些什麼.
除了安排了你返回這間教會.
我們今天是用腳走進來這間教會的.
你有沒有感恩.
我們走進來的.
時間的關係我繼續推進.
這裡遲了出.
大家有經驗了.
最後一個訊息.
最後一個部份.
你發覺剛才.
他都出來了.
最後那一格是什麼.
你幫我讀出來.
下一格是什麼.
是問號.
你想想.
我剛才說我什麼時候返回教會.
我2015年返回教會.
我2016年才洗.
剛剛才十歲.
有誰是信了主超過十年的舉手.
我知道你們.
你們都是師兄師姐.
但原來神在我剛才的照片.
還要找我.
原來是不遲的.
原來是要找我進神學院.
我之前讀了多少書.
他要我進神學院裝備.
我想他訓練了我畢業之後.
我應該還有些事要做.
所以我現在都繼續尋問.
究竟你想我做些什麼.
我們派發最後一個封.

$^{1001}$麻煩施事.
這個封我要特別講解一下.
這個封裡面有一張紙.
有底和面.
這張紙待會.
你可以填寫.
可能你沒有拿筆.
這個我想你真的拿筆.
下面的小圈.
就是你有興趣.
樂意的話就寫你的名字.
上面寫出一個訊息.
新一年.
我想找人關心我.
這件事.
我不知道什麼.
你自己寫.
新一年我想找人關心我.
寫在大雲那裡.
下面可以署名.
例如周朗.
但你說沒什麼.
我ok.
我有心有力.
下一張.
新一年我想去關心其他人.
兩邊都可以寫.
或者隨便寫一邊.
或者你整張都漏白.
不寫.
不要緊的.
寫完之後.
寫在封裡面.
全袋收起來.
全部交給教會.
我不會看的.
希望新的一年.
教會可以透過你的需要.
你需要被人幫.
或者你需要幫人.

$^{1041}$都不緊要.
我們是一起生活的弟兄姊妹.
所以希望新的一年.
第一次崇拜.
大家一起坦誠.
我用一首歌陪大家.
歌詞的過程.
我的訊息分享都去到尾聲.
這隊樂隊叫鹹蛋.
第四個峰.
是紫色的.
你拿一張紙出來.
你完成之後.
準備好.
詩詩就會收回你的峰.
我記得我第一年.
回教會的時候.
是沒有一次缺席.
很快被安排.
入了小組.
沒有一次缺席.
他們說要報名.
生命之道.
然後他們說要查經班.
我出席了.
然後他們在最後.
做報告的時候說.
環球聖經工會那年.
搞了40場講座.
在全香港不同的教會.
我有少少強迫行為.
我跟著他們的行程.
跟著走.
去了不同的教會.
聽他們的道.
那些道是有些深.
譬如其中有白維坤在講.
是有些深.
但我每次都出席.
他們的職員見到.

$^{1081}$弟兄你哪間教會的.
我懂得回答.
他以為我是執事.
傳道.
我還未洗.
他還以為你是神學生.
要出去做探訪.
教會安排了某些職事.
要幫幫忙.
你呢.
這些都要做.
剛才講漏了.
我沒有帶人回教會.
我記得我洗禮那天.
我帶了250人回教會.
我當結婚那樣搞.
把這個開心分享給人聽.
原來我不知道這樣.
是有點無禮.
那天衝擊到教會.
太多人.
還有其他受洗的弟兄姊妹.
不過不要緊.
感動就是要做這件事.
這首歌播多一分鐘就完.
大家嘗試合埋眼.
我最尾講多一個例子.
我認識一位前輩.
他是一個退休校長.
他得到一個癌病.
去到很嚴重的階段.
他得到一個癌病.
去到很嚴重的階段.
他得到一個癌病.
去到很嚴重的階段.
他得到一個癌病.
去到很嚴重的階段.
他得到一個癌病.
去到很嚴重的階段.
他得到一個癌病.

$^{1121}$去到很嚴重的階段.
他堅持一個星期打五日至六日乒乓球.
跟他打乒乓球那班人全部都不信耶穌.
他堅持一個星期打五日至六日乒乓球.
\newpage



\section{馬太福音 4:1-17}
\label{sec:031697}
\textbf{三項試探,一個中心}~~~~林豪恩先生
\newline
\newline
連結: \href{http://www.church.com.hk/acms/content.asp?site=cdc\&op=show\&type=product\&code=031697\&layout=sermon}{\texttt{031697}} ~~~~ 語音日期: 2026-02-22
\newline
\newline
\hyperref[sec:031696]{\small{< < < PREV SERMON < < <}}
~
\hyperref[sec:index_chronic]{\small{[返順時目]}}
~
\hyperref[sec:031698]{\small{> > > NEXT SERMON > > >}}
\newline
\newline
馬太福音 4:1-17
\newline
\begin{longtable}{cl}
\hline
\hline
章節 & 經文 (和合本修訂版)\\
\hline
4:1 & \begin{tabularx}{0.7\textwidth}{X} 當時，耶穌被聖靈引到曠野，受魔鬼的試探。 \end{tabularx} \\ \\ \relax
4:2 & \begin{tabularx}{0.7\textwidth}{X} 他禁食四十晝夜，後來就餓了。 \end{tabularx} \\ \\ \relax
4:3 & \begin{tabularx}{0.7\textwidth}{X} 那試探者進前來對他說：「你若是神的兒子，叫這些石頭變成食物吧。」 \end{tabularx} \\ \\ \relax
4:4 & \begin{tabularx}{0.7\textwidth}{X} 耶穌卻回答說：「經上記著：『人活著，不是單靠食物，乃是靠神口裡所出的一切話。』」 \end{tabularx} \\ \\ \relax
4:5 & \begin{tabularx}{0.7\textwidth}{X} 魔鬼就帶他進了聖城，叫他站在聖殿頂上， \end{tabularx} \\ \\ \relax
4:6 & \begin{tabularx}{0.7\textwidth}{X} 對他說：「你若是神的兒子，就跳下去！因為經上記著：『主要為你命令他的使者，用手托住你，免得你的腳碰在石頭上。』」 \end{tabularx} \\ \\ \relax
4:7 & \begin{tabularx}{0.7\textwidth}{X} 耶穌對他說：「經上又記著：『不可試探主—你的神。』」 \end{tabularx} \\ \\ \relax
4:8 & \begin{tabularx}{0.7\textwidth}{X} 魔鬼又帶他上了一座很高的山，將世上的萬國和萬國的榮華都指給他看， \end{tabularx} \\ \\ \relax
4:9 & \begin{tabularx}{0.7\textwidth}{X} 對他說：「你若俯伏拜我，我就把這一切賜給你。」 \end{tabularx} \\ \\ \relax
4:10 & \begin{tabularx}{0.7\textwidth}{X} 耶穌說：「撒但，退去！因為經上記著：『要拜主—你的神，惟獨事奉他。』」 \end{tabularx} \\ \\ \relax
4:11 & \begin{tabularx}{0.7\textwidth}{X} 於是，魔鬼離開了耶穌，立刻有天使來伺候他。 \end{tabularx} \\ \\ \relax
4:12 & \begin{tabularx}{0.7\textwidth}{X} 耶穌聽見約翰下了監，就退到加利利去； \end{tabularx} \\ \\ \relax
4:13 & \begin{tabularx}{0.7\textwidth}{X} 後來離開拿撒勒，往迦百農去，住在那裡。那地方靠海，在西布倫和拿弗他利地區。 \end{tabularx} \\ \\ \relax
4:14 & \begin{tabularx}{0.7\textwidth}{X} 這是要應驗以賽亞先知所說的話： \end{tabularx} \\ \\ \relax
4:15 & \begin{tabularx}{0.7\textwidth}{X} 「西布倫，拿弗他利，沿海的路，約旦河的東邊，外邦人的加利利— \end{tabularx} \\ \\ \relax
4:16 & \begin{tabularx}{0.7\textwidth}{X} 那坐在黑暗裡的百姓看見了大光；坐在死蔭之地的人有光照耀他們。」 \end{tabularx} \\ \\ \relax
4:17 & \begin{tabularx}{0.7\textwidth}{X} 從那時候，耶穌開始宣講，說：「你們要悔改！因為天國近了。」 \end{tabularx} \\ \\
[1ex]
\hline
\hline
\end{longtable}
$^{1}$弟兄姊妹平安.
今天是大齋期的第一主日.
其實每一年的大齋期第一主日.
我們都會讀耶穌的三個試探.
無論是甲年,乙年,丙年的經課.
都是耶穌的三個試探.
相信大家都熟悉的.
有一次我們學生團契在查經的時候.
就在討論這三個試探.
有同學認為其實耶穌的三個試探.
和自己沒有什麼關係.
為什麼這樣說呢?.
他說第一.
耶穌有能力吩咐石頭變食物.
但是我沒有這種能力.
第二.
我有畏高症.
你叫我站在墊頂上.
我還沒跌死都已經嚇死了.
第三.
我連自己都管理得不好.
要上個早堂我去不到.
要交個功課.
我都未必能夠準時交到.
你把萬國交給我.
那不就是找你煩.
你覺得他說得有沒有道理呢?.
如果是這樣.
耶穌的三個試探又有什麼意義呢?.
我就將這個問題問一下AI.
AI都很厲害.
它懂得以經解經.
它引用《約翰一書》二章十六節.
那裡的三項來解釋.
分別是肉體的情慾.
眼目的情慾.
和今生的驕傲.
它說分腐石頭變成食物.
就是肉體的情慾.
萬國的榮華就是眼目的情慾.

$^{41}$從墊頂跳下去就是今生的驕傲.
聽起來都頗有道理.
不過我們都可以和AI相確一下.
因為AI經常都有一句Disclaimer.
就是說我提供的資料可能有錯誤.
第一.
肚餓要食物.
其實是人之常情.
怎麼可以被貶稱為肉體的情慾呢?.
第二.
天地萬物本來就是上帝的.
耶穌看一看.
就算取回.
又怎麼可以叫做眼目的情慾呢?.
第三.
如果耶穌有能力在墊頂跳下去而不受傷.
為什麼又可以叫做今生的驕傲呢?.
你記不記得在今年一月二十五日.
有一位攀岩運動員赤手攀登台北101.
有沒有評論說他驕傲呢?.
好像不太聽到.
多數都是說他瘋狂.
那些媒體就去看看他的履歷.
然後就發現其實以他的履歷.
他的能力去攀登101只不過是小菜一碟.
對耶穌來說.
在墊頂跳下去.
而能夠全身而退.
那算得了什麼呢?.
這樣耶穌面對三個試探.
到底試探的是什麼呢?.
讓我們看看今天的經課.
帶我們從一個什麼角度去思考.
我們剛才讀過的四段經課.
第一段是舊約的經課.
舊約經課是創世記的故事.
我相信這個故事大家都非常熟悉.
那個故事就是在伊甸園裡面.
有兩個人.
然後有一天.

$^{81}$有一條蛇走來.
就跟其中一個人說.
喂.
園裡面是不是全部果子.
園裡面的果子.
是不是全部都可以吃呢?.
然後那個人就回答他.
他說.
其實全部都可以吃的.
不過有一棵是不能吃的.
吃了就會死的.
那條蛇就跟他說.
其實你有沒有想過.
是不是真的會死呢?.
可能不死呢?.
你知不知道.
你吃了那棵樹的果子.
你很聰明的.
你會好像上帝一樣.
你有沒有想過.
就是因為這樣他才不讓你吃呢?.
然後那個人就看那棵樹的果子.
咦.
好像又可以吃的.
又真的挺好看的.
又挺喜歡的.
於是就摘了來吃.
就給另外一個人吃.
吃完之後.
他就發覺自己赤身露體.
於是就用一些蠍去遮住自己.
其實這是什麼故事呢?.
這是一個騙案.
我們現在是不是有KK權騙案呢?.
原來很早很早的時候.
已經有伊甸園騙案了.
KK權騙案想騙什麼呢?.
騙錢是吧?.
那伊甸園騙案想騙什麼呢?.
騙情.

$^{121}$他騙情.
他騙什麼情呢?.
他騙感情.
他騙誰的感情呢?.
他騙那兩個人.
和上帝之間的感情.
那兩個人是誰來的呢?.
那兩個人在聖經故事裡面.
他們是人類的子祖.
其實他們是誰呢?.
你記不記得.
盧家福音.
有個續譜.
續譜的最後一句是什麼呢?.
不用找了.
我告訴你.
盧家福音的續譜最後一句.
就是說.
亞當是上帝的兒子.
那兩個人是誰呢?.
是上帝的兒女.
大家剛才讀福音書的時候.
有沒有發現上帝的兒子.
在福音書裡面又出現了.
魔鬼似乎是瞄準上帝的兒子.
是以上帝的兒子為攻擊對象.
我們試試看第四章三節.
他說.
你若是上帝的兒子.
叫這些石頭變成食物吧.
第六節說.
你若是上帝的兒子就跳下去.
因為主一定會托著你的.
第九節說.
你若苦復敗我.
我就將這一切賜給你.
你看到.
他是針對著上帝的兒子.
其實作為一個人.
道成肉身的耶穌.

$^{161}$需要食物.
都很合理的.
其實.
想玩一下笨豬跳.
都不是什麼事情.
你去旅行的時候都可以玩.
甚至.
他想收回萬國.
其實都不一定有問題.
問題是.
這些這麼自然的需要.
要滿足他.
究竟我們要付出什麼代價呢?.
我們今天有一些崇基的同學在這裡.
你們剛剛考完DSE沒多久.
應該還記得DSE中文裡面有一篇梵文.
叫做《如我所欲也》.
還有印象吧.
應該背過.
孟子說.
如我所欲也.
熊掌亦我所欲也.
二者不可兼得.
捨魚而取熊掌者也.
生亦我所欲也.
異亦我所欲也.
二者不可兼得.
捨生而取異者也.
耶穌的試探裡有一個是跟飲食有關的.
孟子這篇文章裡有一個例子是跟飲食有關的.
他說.
一缸飯.
一碗湯.
得到.
你就可以存活下去.
失去.
你可能會死.
但是如果呼呼喝喝這樣.
將這碗飯給人吃.
流浪者也不會接受.

$^{201}$如果踐踏之後給人吃.
連乞丐也不會接受.
道成肉身的耶穌會肚餓.
跟常人一樣有生理和心理的需要.
去滿足這些需要.
到底要付出什麼代價呢?.
是否什麼代價都可以接受呢?.
作為上帝的兒子.
在滿足這些需要的時候.
究竟有沒有底線呢?.
上帝的兒子.
究竟會怎樣對待上帝呢?.
在伊甸園的故事裡.
亞當和夏娃作為上帝的兒女.
他們是怎樣對待上帝的呢?.
他們被魔鬼成功離間了跟上帝的關係.
使得原本信任的關係變得猜忌.
於是做出背叛上帝的行為.
使得原來親密變成疏離.
最後要躲避上帝.
現在魔鬼又向上帝的兒子耶穌重施故技.
魔鬼能不能夠離間得到.
耶穌對上帝的信任呢?.
第四章四節.
耶穌這樣回答.
經上記著說.
人活著不是單靠食物.
乃是靠上帝口裡所出的一切話.
第二.
魔鬼能不能夠令到耶穌懷疑上帝.
而試探上帝呢?.
耶穌回答.
經上猶記著說.
不可試探主你的上帝.
第三.
魔鬼能不能夠令耶穌做出背叛上帝的行為呢?.
第四章十節.
耶穌回答.
撒旦,退去吧.
因為經上記著說.

$^{241}$當拜主你的上帝.
單要事奉他.
亞當和夏娃面對魔鬼試探的時候.
他似乎任由魔鬼帶著他去.
似乎看不到他有底線.
耶穌面對魔鬼試探的時候.
你看不看得到他的底線是什麼?.
當想起上帝的兒子應該是怎樣的呢?.
底線是什麼呢?.
我們很容易傾向制定一些行為守則.
例如有一百條要遵守的規則.
有一百條不可以違反的規則.
這樣就可以保障.
這樣很容易變成律法主義.
就好像小朋友經常會問.
為什麼他可以這樣做.
為什麼我不可以這樣.
為什麼他可以吃冰淇淋而我不可以這樣.
判斷好壞有時不只是行為本身.
而是和身份相關的.
大家有沒有聽過這個故事?.
有一天早上.
媽媽走進房間.
大聲叫.
兒子,起床了.
是時候上學了.
兒子用被子蓋著自己的頭.
說「我不上學,我不想上學」.
媽媽說「不行,你一定要起床上學」.
兒子說「旁邊的小明一點點不舒服都可以不上學」.
「為什麼我不可以?」.
媽媽說「你再不起床就來不及了」.
「難道你忘記了你是校長嗎?」.
即是行動是和身份相關的.
小明是可以一點不舒服不上學.
但是校長是不是不可以呢?.
我也不太清楚.
上帝的兒子是什麼呢?.
其實上帝的兒子不是獨立存在的.
你想想.

$^{281}$有稱得上為上帝的兒子.
就必然有另外一個身份存在.
那個就是上帝.
有上帝的兒子就必然有上帝.
所以才可以稱為上帝的兒子.
上帝的兒子應該是怎樣呢?.
這個問題不是獨立的問題.
是要放在上帝的兒子和上帝這兩者來思考.
於是問題就變成.
上帝的兒子和上帝的關係應該是怎樣呢?.
關係不是我們平時所說的那種感覺好不好.
關係不純粹是情緒式的.
從魔鬼對亞當夏娃的試探.
和對耶穌的試探來看.
關係是關乎誠信.
彼此之間有一個共識.
就是雙方是否可以用信任,誠實,不欺騙,不背叛來互相相待.
《新命紀》七章九節怎樣描述上帝呢?.
它用信實來描述上帝.
它說你要知道耶和華是你的上帝.
祂是上帝,是信實的上帝.
信實是包含可靠,穩固,誠實,忠誠不變和始終如一的意思.
上帝以信實來對待祂的兒女.
上帝的兒女會不會用信實來對待上帝呢?.
會不會這就是魔鬼想試探耶穌的地方呢?.
其實不單是上帝的兒女和上帝之間期望以信實相待.
其實在人類社會裡面我們都很渴望.
我們周圍的人和我們之間都是信實相待的.
我們不想我們和我們周圍的人之間的關係是背叛和欺詐.
詩篇32篇今天的經課帶我們回到人的層次裡思考.
我們看看詩篇32篇.
剛才聽過這篇詩篇是一篇認罪則赦的詩篇.
相傳是大衛王寫的.
這首詩篇和今天的福音經課有什麼關係呢?.
首先讓我們問一個問題.
大衛王是誰?.
大衛王是什麼身份呢?.
就是王.
那王是怎樣的呢?.
王應該是怎樣的呢?.

$^{321}$按照我們剛才的看法.
王不是獨立存在的.
如果你說自己是王.
你沒有指紋.
我們只是在取笑你.
有王就有指紋.
王應該是怎樣的呢?.
是在王和指紋的關係下思考的.
王和指紋的關係應該是怎樣的呢?.
其實不需要用愛和不愛的感性字眼去描述.
你試想一下.
你是一個指紋.
其實王愛不愛你.
對你來說並不重要.
重要的是什麼呢?.
重要的是.
王是否用信實來對待我.
是否誠實,守信,不欺騙,不背叛來互相相待.
他說的話是否算數呢?.
大衛王要認什麼罪呢?.
大衛王在這裡認什麼罪呢?.
他犯了什麼罪呢?.
撒慕爾記下十一章至十二章記載了大衛王.
和一個女人同房.
然後殺了一個男人.
在古代的君王.
後宮三千.
和一個女人同房是什麼罪呢?.
古代的君王.
操控生殺大權.
殺一個人是什麼罪呢?.
聖經指出.
那個女人是烏尼亞的妻子拔士巴.
而大衛殺的那個男人.
是拔士巴的丈夫烏尼亞.
烏尼亞是誰呢?.
我們看看撒慕爾記下十一至十二章是怎樣描述這件事情的.
第十一章第八節.
大衛就召烏尼亞從戰場上回來.
就跟她說.

$^{361}$你回家吧.
洗洗腳吧.
就是跟她說.
你回家餵養妻吧.
烏尼亞怎麼說呢?.
烏尼亞跟大衛說.
約桂和以色列和猶大的兵都住在棚內.
就是說.
我們和我們的兄弟.
現在在前線打仗.
我的上司約克和我的手下.
全部都在田野安營.
我怎麼可以回家吃吃喝喝.
和我老婆吃吃喝喝呢?.
我這樣做.
我怎樣對待我的兄弟呢?.
你看看這個人是誰?.
這是一個用誠信誠實.
不欺騙不背叛.
對待大衛王的人.
是不是?.
他是這樣對待大衛王.
對待他的國家.
對待他的同伴.
大衛王怎樣對待他呢?.
第十一章十四節十五節.
大衛王寫一封信.
叫他交給將軍約克.
然後信裡面寫什麼呢?.
你幫我搞定他.
還要不要讓人知道.
我相信沒有人會否定大衛王.
作為一個男人是有性需要的.
但是他有需要.
他就可以這樣.
叫別人的妻子來涉獄嗎?.
烏尼亞對王對國家忠誠.
這個戰士大衛王竟然暗殺他.
君王和子民之間的關係是這樣的嗎?.
大衛王犯的是什麼罪呢?.

$^{401}$淫亂?謀殺?.
這些只是罪名.
他犯的是他破壞了.
君王和人民之間的誠信.
他出賣和背叛了他的子民.
不知道你看到這個故事的時候.
你有沒有感到憤怒和悲傷呢?.
信實.
不欺騙.
不出賣.
其實是人類的共識.
如果你覺得憤怒和悲傷.
即是說你還有這種共識.
我們人間的社會.
有多少的災難,痛苦.
就是因為有人違反了這些誠信.
做了一些背叛和出賣的行為.
產生了多少的痛苦和悲劇.
我們試試問一個問題.
大廈維修公司應該是怎樣的呢?.
工程顧問公司應該是怎樣的呢?.
安全督導應該是怎樣的呢?.
物料供應商應該是怎樣的呢?.
這些所有的問題都不是獨立的問題.
你試試想想.
如果你將你住的那座大廈交給他們維修.
這些所有問題都是他和你之間的問題.
是在大家之間的問題.
那些維修商.
他們對於將自己的錢和生命交給他們的居民.
究竟存在著甚麼關係呢?.
他應該怎樣去對待住在裡面的居民呢?.
如果這些公司和職業人士.
破壞了他和居民之間的誠信.
欺騙了居民,背叛了居民.
所產生的悲劇和災難我們都看到.
所產生的痛苦很多居民都繼續承受.
一個人應該是怎樣的呢?.
一個人應該是怎樣做的呢?.
怎樣生活的呢?.

$^{441}$是需要放在一個關係裡面去思考的.
因為一個人和其他人.
是存在著這種誠信的基礎的.
我們是有這樣的期望的.
古人說.
人是有很多天然的需要.
賺錢養生當然是需要.
君子好財.
但都是隨之有道.
人還有很多其他的需要.
人需要有親密感.
人的性滿足亦都需要.
問題是.
當我們去滿足這些需要的時候.
有沒有問過我是誰?.
他又是誰呢?.
最近有一位著名的基督徒作家.
Philip Yancey 亦叫做楊肥力.
發表了聲明.
承認他和一名已婚女子.
保持了長達八年的婚外情.
他的聲明部份內容是這樣的.
我的行為違背了我對婚姻的一切信念.
亦都和我的信仰和諸術完全背道而馳.
給他的丈夫和我們雙方的家庭.
帶來了深深的痛苦.
我在道德和人性上都犯了錯.
我為自己造成的傷害感到無比悲痛.
我意識到我的行為.
會讓那些曾經信任我作品的讀者感到失望.
更糟糕的是.
我的罪羞辱了上帝.
我充滿了悔恨和懺悔.
除了上帝的憐憫和恩典.
我必無依靠.
其實.
一個男人和一個女人發展感情關係.
不一定有問題.
但是.
問題是這個男人和那個女人.

$^{481}$他們已經進入了一種承諾的關係.
即是這個男人和另一個人有一種承諾.
要求忠誠的關係.
那個女人又和另一個男人有一種承諾和忠誠的關係.
在這樣的情況下.
正如菲利普當時所說.
她做的事情.
讓那個女人的丈夫和自己的妻子都非常痛苦.
如果將她這段感情.
放在她和讀者的關係裡.
她就說.
讓那些曾經信任我的作品的讀者感到失望.
如果將她和這個女士的關係.
放在她和上帝的關係裡.
她就說.
我的罪羞辱了上帝.
那個行為的對和錯.
不因為那個行為.
而是她是誰.
做這個行為的是誰.
究竟.
我們知不知道自己是誰呢?.
我們能不能夠誠信地活出自己呢?.
其實是一個非常艱難的狀況.
是很不容易的狀況.
但是如果能夠這樣做的話.
是一個帶有拯救效果的狀況.
世界上有沒有人能夠這樣活出自己呢?.
使徒保羅說是非常罕有.
但是他在今天的書信經課裡指出有一位.
我們看看保羅怎麼說.
《羅莽書第五章十八節》.
他說「這樣看來.
因一次的過犯 眾人都被定罪.
照樣因一次的義行 眾人都被稱義得生命了.
因一人的背逆 眾人成為罪人.
照樣因一人的順從 眾人也成為義了」.
只要有人能夠這樣活出真誠.
這樣活出自己的誠信.
他做得到上帝的兒子這種身份.

$^{521}$活出的時候.
原來拯救的果效是普及全人類的.
我們難免會在人生裡掙扎.
有時會問我是誰呢?.
年輕的同學們.
你們都會不停問.
究竟我是誰呢?.
我們都很難完全清楚自己究竟是怎樣的.
我們也很難百分百地真誠地.
誠實地活出自己.
但是只要能夠去到某個程度.
就會產生某個程度的拯救效果.
試想想 在火災裡.
所有人都逃離火場.
但有些人知道自己是不可以逃離火場的.
當所有人都逃離火場的時候.
他們反而要衝進火場.
那些是誰呢?.
消防員.
他要知道自己是誰.
他知道自己是消防員.
他就不可以逃離 反而要衝進去.
為什麼呢?.
因為在火災裡的災民正在等待他.
火災裡的災民最渴望見到誰來呢?.
當然是消防員.
你們試想想.
原來消防員和災民是存在著這樣的.
共識和關係.
如果他能夠活出這種身份.
實踐這種誠信和共識.
其實就有人被拯救了.
我們不是消防員.
我們每個人都存在在.
和其他人有一種身份和共識的關係裡.
我們怎樣能夠信實地活出自己.
不背叛 不出賣呢?.
在去年12月.
天主教教廷冊封了一位阿根廷的評述圖商人.
作為地位僅次於聖人的「真福」.

$^{561}$這位商人為什麼會被冊封呢?.
他是不是有很多神蹟呢?.
不是.
情況是1961年的時候.
他的其中一間公司.
被美國的一個信託基金收購了.
信託基金就決定裁員1200人.
當時已經患了末期癌症的商人強烈反對裁員.
並且提出保留所有員工的改革計劃.
他的朋友表示.
在他接受癌症治療的時候需要輸血.
他工廠的工人排著隊去醫院輸血給他.
人數之多令到醫院大衛震驚.
一度以為他都是工會成員.
但是其實他是老闆.
老闆是怎樣的呢?.
老闆應該是怎樣的呢?.
其實老闆是存在在跟員工的關係裡面的.
當雙方都以信實相待.
雙方都得到拯救.
今天是大齋期的第一主日.
經過呼喚我們.
回到耶穌的故事.
在想我是誰呢?.
我跟誰存在著這些誠信的關係.
以致我要誠信地活出自己呢?.
願上主幫助我們.
\newpage



\section{馬可福音 9:14-29}
\label{sec:031698}
\textbf{變臉後傳 I 心弱力微}~~~~梁永善牧師
\newline
\newline
連結: \href{http://www.church.com.hk/acms/content.asp?site=cdc\&op=show\&type=product\&code=031698\&layout=sermon}{\texttt{031698}} ~~~~ 語音日期: 2026-02-25
\newline
\newline
\hyperref[sec:031697]{\small{< < < PREV SERMON < < <}}
~
\hyperref[sec:index_chronic]{\small{[返順時目]}}
~
\hyperref[sec:031699]{\small{> > > NEXT SERMON > > >}}
\newline
\newline
馬可福音 9:14-29
\newline
\begin{longtable}{cl}
\hline
\hline
章節 & 經文 (和合本修訂版)\\
\hline
9:14 & \begin{tabularx}{0.7\textwidth}{X} 他們到了門徒那裡，看見有一大群人圍著他們，又有文士和他們辯論。 \end{tabularx} \\ \\ \relax
9:15 & \begin{tabularx}{0.7\textwidth}{X} 眾人一見耶穌，都很驚奇，就跑上去向他問安。 \end{tabularx} \\ \\ \relax
9:16 & \begin{tabularx}{0.7\textwidth}{X} 耶穌問他們：「你們和他們辯論甚麼？」 \end{tabularx} \\ \\ \relax
9:17 & \begin{tabularx}{0.7\textwidth}{X} 眾人中的一個回答：「老師，我帶了我的兒子到你這裡來，他被啞巴的靈附著。 \end{tabularx} \\ \\ \relax
9:18 & \begin{tabularx}{0.7\textwidth}{X} 無論在哪裡，那靈拿住他，把他摔倒，他就口吐白沫，牙關緊鎖，身體僵硬。我請過你的門徒把那靈趕出去，他們卻不能。」 \end{tabularx} \\ \\ \relax
9:19 & \begin{tabularx}{0.7\textwidth}{X} 耶穌回答：「唉！這不信的世代啊，我和你們在一起要到幾時呢？我忍耐你們要到幾時呢？把他帶到我這裡！」 \end{tabularx} \\ \\ \relax
9:20 & \begin{tabularx}{0.7\textwidth}{X} 他們就帶了他來。那靈一見耶穌，就使他重重地抽風，倒在地上，翻來覆去，口吐白沫。 \end{tabularx} \\ \\ \relax
9:21 & \begin{tabularx}{0.7\textwidth}{X} 耶穌問他父親：「他得這病有多久了呢？」父親說：「從小的時候。 \end{tabularx} \\ \\ \relax
9:22 & \begin{tabularx}{0.7\textwidth}{X} 那靈屢次把他扔在火裡、水裡，要治死他。你若能做甚麼，求你憐憫我們，幫助我們。」 \end{tabularx} \\ \\ \relax
9:23 & \begin{tabularx}{0.7\textwidth}{X} 耶穌對他說：「『你若能』，在信的人，凡事都能。」 \end{tabularx} \\ \\ \relax
9:24 & \begin{tabularx}{0.7\textwidth}{X} 孩子的父親立刻喊著說：「我信；求你幫助我的不信！」 \end{tabularx} \\ \\ \relax
9:25 & \begin{tabularx}{0.7\textwidth}{X} 耶穌看見眾人都跑上來，就斥責那污靈說：「你這聾啞的靈，我命令你從他裡頭出來，再不要進去！」 \end{tabularx} \\ \\ \relax
9:26 & \begin{tabularx}{0.7\textwidth}{X} 那靈大喊一聲，使孩子猛烈地抽了一陣風，就出來了。孩子好像死了一般，以致眾人多半說：「他死了。」 \end{tabularx} \\ \\ \relax
9:27 & \begin{tabularx}{0.7\textwidth}{X} 但耶穌拉著他的手，扶他起來，他就站起來了。 \end{tabularx} \\ \\ \relax
9:28 & \begin{tabularx}{0.7\textwidth}{X} 耶穌進了屋子，門徒就私下問他：「我們為甚麼不能趕出那靈呢？」 \end{tabularx} \\ \\ \relax
9:29 & \begin{tabularx}{0.7\textwidth}{X} 耶穌對他們說：「非用禱告，這一類的邪靈總趕不出來。」 \end{tabularx} \\ \\
[1ex]
\hline
\hline
\end{longtable}
$^{1}$各位親愛的電影節目你們好.
我們又開始星期三的《茶經》了.
從歲守到年終.
神眷顧我們很重要的一句話.
就是我們天天仰望天天的彌大去抱擁.
剛剛在預備4月的《時詩記》.
第一章講到他們怎樣成功呢?.
成功有一樣東西.
第一點就是他們要不斷禱告仰望神.
如果我們每天禱告仰望神的時候.
我們必定得到那個力量.
我們今天會看到最後兩次登山變相.
我們來看後傳.
就是《深藥歷微》.
我們過去看到耶穌與三個愛徒.
彼得,雅各,約翰登山.
然後耶穌就變了他的形象.
就讓你看到主耶穌基督榮耀的真相.
因為耶穌下山後沒多久就要去耶路撒冷.
為人的員工被釘十字架.
受苦受死.
所以讓你們看到耶穌的真相.
讓他們有信心,有能力去面對更大的挑戰.
承擔更偉大的使命.
但是其實當他們下山的時候.
他們就已經遇到不同的挑戰了.
記不記得我們說.
上山休息就是為了下山面對挑戰.
有時候我們都需要停下來休息一下.
但我們休息就是為了要走更長的路.
第一個挑戰是什麼呢?.
就是對門徒.
下個星期我們就會看對耶穌的挑戰.
原來當時那些門徒.
不能夠醫好一個患了癲癇症被鬼附著的小孩子.
這件事的記載就在《馬太福音》第十七章十四到二十一.
但另外有兩本福音書都有記載的.
就是《馬太福音》第九章十四到二十九.
和《路救福音》第九章三十七到四十二節.
其中記載得最詳細的.

$^{41}$就是《馬太福音》第九章十四到二十九節.
我們就用這個作為一個基礎.
其他的兩個福音書作為一個參考.
辛苦了 伊安.
這次我們會讀得長一點.
就是《馬太福音》第九章十四到二十九節.
耶穌到了門徒那裡 看見有許多人圍著他們.
又有民事和他們辯論.
眾人一見耶穌就心稀奇.
就跑去問他的安.
耶穌問他們說:你們和他們辯論的是什麼?.
眾人中間有一個人回答說:.
夫子 我帶了我的兒子到你這裡來.
他被啞巴鬼附著.
無論在哪裡鬼捉弄他 把他睡倒.
他就口中留沒 咬牙切齒 身體枯乾.
我請過你們門徒把鬼趕出去.
他們卻是不能.
耶穌說:喂 不信的世代啊.
我在你們這裡要到何時呢?.
我忍耐你們要到何時呢?.
把他帶到我這裡來吧.
他們就帶了他來.
他一見耶穌 鬼便叫他重重的抽風.
倒在地上 翻來覆去 口中留沒.
耶穌問他的父親說:他得這病有多少日子呢?.
回答說:從小的時候.
鬼類似把他拎在火裡 水裡要滅他.
你若能作什麼 求你憐憫我們 幫助我們.
耶穌對他說:你若能信.
在信的人凡事都能.
孩子的父親立時喊著說:我信.
但我信不足 求主幫助.
耶穌看見眾人都跑上來.
就斥責那污鬼說:你這聾啞的鬼.
我吩咐你從他裡頭出來 再不要進去.
那鬼喊叫:使孩子大大的抽了一陣風.
就出來了.
孩子好像死了一般.
以至眾人多半說他是死了.

$^{81}$但耶穌拉著他的手扶他起來.
他 伯 就站起來了.
耶穌進了屋子.
門徒就暗暗的問他說:.
我們為什麼不能說出他去呢?.
耶穌說:非用禱告.
這一類的鬼還不能出來.
我們看到一段經文就是講耶穌下山.
馬上有些人來圍著耶穌.
為什麼呢?.
就說這些門徒沒有實力.
不行.
叫他趕鬼趕不了.
事實上門徒真的趕不了.
所以某個程度門徒真是啞口無言了.
我們看回原來第二天.
聖經裡面記載.
他們下山的時候.
就有很多人來迎見耶穌了.
在諾貝爾第九章第37節.
這段經文很清楚指出.
原來耶穌在晚上.
在山上登山變像.
他們整晚都睡不著.
所以那三個門徒打瞌睡.
翌日就下山了.
當他們下山的時候.
看到很多人圍著他們.
還有文士和法利賽人來跟他們辯論.
回到第九章第十四節.
「耶穌到了門徒那裡.
看見有許多人圍著他們.
又有文士和他們辯論」.
就是說耶穌看到那些門徒.
被很多人圍著.
去挑戰他.
質問他.
但是門徒在那一刻.
真是啞口無言.
因為當時耶穌是一個很矚目的人.

$^{121}$他們的門徒同樣都備受注意.
所以門徒.
他們曾經領受了主的使命和能力.
去醫病,趕鬼.
事實上他們曾經成功的.
在《路透社》第九章第一到第二節.
「耶穌叫齊了十二個門徒.
給他們能力,權柄.
制服一切的鬼,醫治各樣的病.
又差遣他們去宣傳.
神國的道,醫治病人」.
看見嗎?.
耶穌是給十二個門徒.
給他們能力,權柄.
他們出去.
兩個兩個出去.
事實上他們都怎樣呢?.
真是趕鬼,醫病.
他們回來很開心.
記不記得這段經文.
主耶穌說.
你不要圍著.
你只可以趕鬼治病.
開心.
他說最重要是怎樣呢?.
是因為你為神工作.
你會得到祝福.
這樣才是開心.
即是說門徒曾經成功的.
趕鬼治病.
但是這次他們是沒有能力去趕鬼.
結果他即時就會受到質疑.
被挑戰.
而挑戰他們最厲害的是什麼?.
就是當時那些宗教人士.
那些民事.
他們和閩門徒在辯論.
這是一個很大的感慨.
當你做得好.
當你成功的時候.

$^{161}$可能你是不被紀念的.
第一次的失敗.
可能那些人永遠都會記得.
記得你的失敗.
抹殺了你過去的成功.
這就是很無奈的.
甚至我剛剛出道的時候.
做傳道人.
去建立了兩間教會都是由零開始的.
但是我問一些要去開這些教會的人.
你知不知道誰開這間教會的?.
他們是不知道的.
原來可能在歷史裡面被抹去的.
今年是冥恩堂的第40年了.
很大感慨.
離開了10年了.
會不會過去這10年裡面來的人.
可能都不會知道有梁詠詩這個人的.
但無論如何.
人可能不記得.
但是神是會記得的.
這個很重要.
而另外一個很大的感慨是什麼呢?.
特別是在教會你做得好.
那些人可能會不記得.
但是你做錯了一些事情.
或者你講了很多編道給他們有得著.
但是可能你講了一篇的信息.
他們會覺得你針對他.
他就不妥你了.
但是電溫轉很無奈.
你想想每到新年的時候.
很多那些人就會講有些什麼運程.
可能他講十幾二十樣.
但是其實有很多樣都不準確的.
但是別人就會記得他一兩樣好像準確了.
於是就會相信他.
這個就是無奈了.
我們有沒有想一想他們辯論的內容是什麼呢?.
他說到那幫人來跟他辯論.

$^{201}$可能就說到.
你跟隨那個耶穌.
是不是真的好像他們所說的神呢?.
你們是不是真的有能力.
可以來講鬼.
如果真的有能力.
為什麼你講不到呢?.
那些鬼是不是比你厲害呢?.
你可不可以想像到當時的門徒.
因為真的講不到鬼.
還被人去挖苦去針對去挑戰的時候.
他們是何等的焦急呢?.
事實上他們百詞莫辯的.
因為事實擺在眼前.
他們無能為力.
所以他們真是啞口無言.
怎麼論下去呢?.
都不能夠說服那些挑戰者的.
所以有時候就有這種感慨.
但是不要忘記.
就在最適當的時間.
耶穌就出現了.
為什麼在山上的時候.
那三個門徒就說到.
哦.
在這裡真好啊.
這是事實.
但是為什麼我們說.
他們仍然要下山.
山上的退休.
就是為了下山要面對更大的挑戰.
和事奉.
主同樣享受在山上的寧靜.
與聖彌達去交往.
與摩西 以利亞來對談.
但是他不會忘記.
門徒當時的困境.
以及那個因為被魔鬼壓制.
而受痛苦的小孩子.
和他的家人.

$^{241}$所以他就去到他們當中.
正如我說到.
我自己也在著緊事奉.
縱使我放假也好.
如果不是放長假的話.
我也會繼續堅持.
和大家一起去查經.
這是很重要的.
特別當年我在閉恩堂事奉.
就算我放假去外地也好.
我也會繼續關心教會.
繼續寫講宗議秘的東西.
是不是?.
記住.
昔日摩西在山上四十晝夜.
與神友交往.
領受誡命.
但是下山的時候.
都是面對.
環境的百姓.
叫摩西上山四十天.
去領取誡命.
但是下山是面對什麼?.
那班子民.
拜金牛獨.
拜牛將.
對神不忠.
所以同樣主要下山的時候.
都碰到這個問題.
所以就發出這種.
哀嘆無奈.
我們看看《路經》第九章第四十一節.
「耶穌說:哀!.
這又不信又不謀的世代.
我在你們這裡.
忍耐你們要到幾時呢?.
將你的兒子帶到這裡來吧」.
耶穌說:你們這是.
又不謀又不信的世代.
主耶穌很有忍耐.

$^{281}$「我都要忍耐你們到幾時呢?」.
當時有很多人圍著耶穌.
和祂的門徒.
第九章第三十七節.
有許多人迎見耶穌.
這些人原來是看熱鬧.
想看神蹟.
但他們卻不願意毀身.
不願意悔改.
所以耶穌說.
看得多是沒有用的.
看到神蹟又如何呢?.
最重要是這班人的生命質素.
是否有擺上.
正如有些人去求耶穌顯神蹟.
耶穌怎麼說呢?.
路口福音第十一章第二十九到三十二節.
「當眾人聚集的時候.
耶穌開講說.
這世代是一個邪惡的世代.
他們求看神蹟.
除了若拿的神蹟以外.
再沒有神蹟給他們看.
若拿怎樣為尼尼美人成了神蹟.
人子也要照樣.
為這世代的人成了神蹟.
當審判的時候.
南方的女王要起來.
定這世代的罪.
因為她從地極而來.
要聽所羅門的智慧說.
看那在這裡有一人比多羅門更大.
當審判的時候.
尼尼美人要起來定這世代的罪.
因為尼尼美人聽了若拿所傳的.
就悔改了.
看那在這裡有一人比若拿更大」.
所以我經常說.
不要用神蹟去確認我們的信仰.
因為我經常說.

$^{321}$其他的信仰.
就算那些原始森林裡的部落民族.
他們拜他們的神.
都有能力行神蹟.
可以醫病.
所以我們不要將我們的焦點放在這裡.
譬如說車公廟的那些扣簽.
好像真的很吻合那個一年.
我說我不否認.
有這些靈驗的.
因為我相信靈界有一些力量.
那些人把他賣給魔鬼.
他都有這個能力去講一些預言.
所以不要用神蹟作為我們信仰的基礎.
是以神的話語.
有些人就叫耶穌顯神蹟.
耶穌說沒有啊.
除了若拿的神蹟沒有了.
你知道若拿是在漁福裡面三日被淘汰出來.
其實要知道耶穌多少個人死了.
但三日之後復活了.
那時候祂說你不要只看神蹟.
說到將來.
利利未人都要審判.
因為他聽了若拿的宣告.
他們悔改.
耶穌其實指責當時的人.
你看到其實我行了不少的神蹟.
但你還要去找.
你聽到我的宣告.
但你們不肯悔改.
所以你就要面對這個更嚴重的審判.
這就是當時耶穌所講的一段話.
在這個時候門徒被人用很難堪的說話.
來責難祂.
挑戰祂.
甚至冷嘲熱諷.
解救這個孩子.
他們無能為力.
所以耶穌出來幫他們.

$^{361}$留意這一段經文.
這是他們沒有期待或者沒有預期的.
我們看馬方第九章第十五節.
眾人一見耶穌.
都心稀奇.
就跑上去問他的安.
你有沒有留意到這一節經文.
只是在《馬可福音》裡面出現的.
是《馬可福音》獨特的記載.
意思就是說他們沒有預計到耶穌會出現.
甚至覺得很意外.
覺得很興奮,很刺激.
不是見到耶穌.
你這幫門徒還不死?.
但就在這一刻.
耶穌在適當的時間就出現.
給他們援手.
但記住今天神有祂自己的時間.
但神的時間不等於我們期望的時間.
但神的時間一定是最好的時間.
神知道什麼時候出手是最適合的.
正如你訓練你的孩子.
他們做功課.
有些事情不明白.
你不會馬上給他答案.
要他自己去試.
真的不行.
你才會來幫他.
但記住門徒肯定出乎意料之外.
開心,興奮.
因為見到耶穌的出現.
挑戰門徒的人.
就會出乎意料之外的失望和不快.
但無論那些人有什麼不同的反應.
主耶穌總是有祂自己的時間.
記住.
人歡迎渴望主來.
但不能夠因為這些渴望.
這種熱切期待.
驅使主早些出現.

$^{401}$有時人不想主耶穌出現.
拒絕.
但同樣這些拒絕,否定.
不能夠延遲主的來臨.
記住祂的時間是最好的.
當我們感到焦慮,疑惑.
主往往就是在我們最稀奇的時間出現.
請你記住這一點.
今天可能有很多事情較為妄言.
但記住神會在適當的時間.
是最好的時間來幫助我們.
我們看看那個爸爸.
他怎樣哀求主耶穌.
那些閒人挑戰者只是吃花生.
只是在看熱鬧.
但原來這位慈父的兒子是什麼呢?.
是獨生子.
在《路加文第九章》第38節裡有記載.
其中有一人喊叫說.
夫子求你看顧我的兒子.
因為他是我的獨生子.
看不看到?.
原來這個兒子是他的獨生子.
所以他很痛愛這個兒子.
但這個兒子竟然受了很多的苦.
於是他求耶穌幫助他.
但如果我們細心研讀聖經的時候.
我們看到主耶穌行神蹟的時候.
你有沒有留意到.
很少對被醫治的人有詳細的描繪.
只是輕輕的帶過.
很簡單.
耶穌無論醫治一個麻瘋病者.
或是那十個麻瘋病者.
都不會怎樣去描繪那些麻瘋病.
怎樣蠶食他們的身體.
怎樣痛.
怎樣辛苦.
怎樣被拒絕.
那個血流的婦人去摸耶穌.

$^{441}$是沒有講到這個婦人的痛楚在哪裡.
她怎樣被這個病人折磨.
啞巴鬼都是.
因為聖經不是要拍電影.
不會特別描繪這些慘恐怖的情況.
焦點只是放在主耶穌的能力身上.
但唯有這一次.
是有很詳細的記載.
為何會這樣呢?.
一會兒我們會解釋.
我們從聖經看到.
這個被鬼附著的.
患癲癇病的少年人.
是受著不少的苦楚.
被鬼附.
即是說被鬼附身.
是絕對不能夠自制.
失去了他自己的自由意志.
其實我也有幾次趕鬼的經驗.
但很多人都沒有聽過我講趕鬼的事情.
我不覺得這些是要去吸引人的.
不過下一次在《學是學非少年》那裡.
我都會講少少我曾經趕鬼的一些經驗.
因為剛剛跟他講到.
凌駕是有一定的力量.
但是被鬼附的人.
記住他是會失去了自由意志的.
啞巴鬼.
即是上了他的身.
說不出話來.
你有沒有經歷過有口難言呢?.
你有痛苦 你有痛楚.
但是你不可以向人表達.
那種是最傷痛的一件事.
不能夠表達自己的感受.
是很痛苦的.
這個人是被啞巴鬼控制著.
還有什麼呢?.
是被摔倒的.
摔在地上.

$^{481}$跌在水裡面.
又跌在火裡面.
更加被鬼要將他毀滅的.
我們看看《馬可福音》第9章第22節.
「鬼類似把他拿在火裡,水裡要滅他.
你若能作甚麼求你憐憫我們 幫助我們」.
聽到這個小孩子被鬼折磨得很厲害.
不能夠表達他的痛楚.
又不由自主地.
推到火裡面給火燒.
或者給水泡.
還有什麼呢?.
口中不斷地流口水.
控制不了自己.
咬牙切齒.
甚至會抽筋.
不單單這樣.
後來《基督經》記載.
他是身體枯乾.
即是說他營養不良.
健康極其差.
長期的受壓.
原來他自小就已經這樣了.
可以見到這個少年人.
在極差,極惡劣的環境之下.
很可惜.
那些門徒.
我很相信他們不止一個.
留在那裡的九個門徒一定是會同心的.
去一起的.
想用方法將這個鬼趕出來.
但是這次面對這麼嚴重的個案.
他們真的束手無策.
但是並不會絕望.
因為主耶穌的介入.
就令到這個小孩子的生命得到改變.
這個孩子的父親.
其實很值得我們欣賞和學習.
他絕對不介意在眾人面前.
向主下跪.

$^{521}$在《馬太福音》第十七章第十四節.
我們沒有打出來.
是講到甚麼呢?.
這個父親在這麼多人看著的時候.
他跪在地上求耶穌幫助他.
他苦苦哀求.
而且很謙卑的,很右麗的.
發出這個呼喊.
馬方第九章二十二節.
「你若能作甚麼.
求你憐憫我們 幫助我們」.
看不看到?.
他不是指令耶穌.
有時候我們祈禱好像是要.
學令我們的神一定要甚麼呢?.
就是要賜福我們.
今年是馬年.
有人就說馬上祝福.
我說不應該講這句.
你學令神馬上祝福你嗎?.
我會改一改.
你馬上靈修吧.
即是馬上祝福.
何時輪到你這樣跟神說話?.
所以你留意到.
他不是要求主一定要跟我這樣做.
他說神啊.
如果可以.
求你憐憫我,幫助我.
還要你留意一點.
他剛剛經歷到那些門徒.
沒有能力去趕跪.
但他對主仍然有這份信心.
不會因為主的門徒失敗.
而放棄了對主的呼吸.
其實這是我們今天.
很多基督徒需要學習的.
因為有時候我們會將信心.
放在一些牧者,傳道人身上.
但不得不承認.

$^{561}$有時候有很多牧者,傳道人.
會令我們失望.
有時候他們可能.
也有些做錯事的地方.
因為他們都是人.
於是他們說.
你看他們這樣.
我怎麼相信呢?.
我們應該將焦點.
放在神的身上.
今天雖然有牧者.
教會的領袖叫我們失望.
但我們不可以對主有失望.
記住.
人的失敗.
不等於神的失敗.
神仍然會繼續用他所揀選的人.
主並不會因為這一次門徒失敗.
放棄這個孩子.
更加不會因為門徒的失敗.
而放棄那些門徒.
所以就好像馬太福音.
在山上所見到的這句話.
好 馬太福音第十七章第八節.
「他們舉目不見一人.
只見耶穌在那裡」.
我們將我們的焦點.
集中在耶穌的身上.
你知道撒旦最喜歡用甚麼方法嗎?.
打落水狗.
要將人的信心打垮.
請你留意那小孩子.
一見到有甚麼反應?.
我們看看馬太福音第九章第二十節.
「他們就帶了他來.
他一見耶穌.
詭辯叫他重重的臭風.
倒在地上.
翻來覆去.
口中留沒」.

$^{601}$有沒有留意到原來一見到耶穌.
那隻鬼不是走.
是更加加害這個小孩子.
這個小孩子比平時受的苦更加嚴重.
這時候會不會因為這個緣故.
打擊了這個爸爸的信心?.
令到這個爸爸更加難受呢?.
記住苦難或者趕鬼.
未必即時能夠成功的.
我經常這樣說.
當一些人剛剛信主.
又或者當他要洗禮.
洗禮之前或者洗禮之後.
又或者他做了一些很重要的抉擇.
但是譬如要委身事奉神.
他往往就會受到很多的衝擊.
這些例證在我過去的墓會的時候.
都有不少這樣的情況.
有一個會友當年洗完禮.
星期天洗禮.
丈夫還送了一束花給他.
誰知星期一.
丈夫就逃走了.
因為他是水警的時候.
原來欠了很多賭債.
沒錢還.
丟下家裡 丟下兩個兒子.
那個會友很慘.
剛剛才洗完禮.
第二天 丈夫就走了.
還有一個家庭.
現在仍然有教會.
他決定要洗禮.
很難得他拜菩薩的婆婆.
都讓他來洗禮.
但當他快要洗禮之前.
他就發現了自己.
原來甲狀腺有事.
這時候就要吃一些藥叫「癲」.
但當他吃了藥的時候.

$^{641}$他剛剛懷孕.
但他說如果吃了藥.
就會令到小孩子有問題.
結果很無奈地要終止懷孕.
於是他就問為甚麼神會這樣.
他想放棄信仰.
但我跟他說.
為甚麼你不從另一個角度去看呢?.
為甚麼你會發現自己有問題呢?.
有問題就是因為有一個團友跟他很熟.
他在醫院裡面做護士.
覺得他的眼睛突出來.
或者有某些問題.
於是就跟他安排早些去看醫生.
才發現了這件事.
當發現了這件事的時候.
他就可以及早地來大約去處理.
如果他懷孕再多幾個月的時候.
那時候就有很大的困難.
但我想說的是.
原來當你要承擔一些更大的責任.
撒旦可能就會藉著這樣打垮你.
我自己當年開上水堂的時候.
都有這個經歷.
所以撒旦用不同的方法.
使我們的信心受損.
使我們知難而退.
所以我們更加需要他依靠我們的主.
你留意到很值得欣賞的就是.
這個慈父並沒有質疑耶穌的能力.
他就繼續求神憐憫幫助他.
因為主因應他的呼求.
於是再作出這個挑戰.
我們看一看.
《馬福音第九補章》23,24兩節.
耶穌對他說:.
「你若能信,再信的人凡事都能」.
孩子的父親立時喊著說:.
「我信,但我信不足」.
求主幫助.

$^{681}$這句說話很寶貴.
我想大家都很熟悉.
第24節我經常都會用的.
主啊,我信你.
但我信不足.
求你來幫助.
這個父親很坦誠地跟耶穌說.
我真是信心不足.
為什麼?一來到.
兒子就已經受到這麼多的苦.
怎麼辦?.
我們不應該怪責他.
因為這個兒子多年來.
他受著很大的折磨.
相信他曾經向不同的人.
甚至醫生來求助.
但都是徒勞無功.
一次又一次地失望.
他的信心越來越脆弱.
是很自然的事.
但竟然沒有掩飾自己那種缺陷.
而向神來求.
我信,但我信不足.
求主幫助.
今天我們有時候好像面對一些困境.
神好像沒有開路.
但你會向神求.
神仍然確信你是眷顧我的.
但我不得不承認.
我有不足的地方.
你幫我.
慈父不單單求主醫治他的孩子.
也為自己的信心向主來求.
求主加添他的信心.
使他能夠接受主更大的恩典.
所以記住.
信心原來可以求的.
羅馬書怎樣說?.
本於信,而至於信.
所以我們可以向神父求.

$^{721}$我們看看主耶穌的能力.
聖經說到耶穌看見眾人都跑來.
記不記得我剛才說.
很多人吃花生.
看熱鬧.
看演唱會.
甚至催錢.
想看看.
這個人又打到地上了.
死呀死呀.
看看現在怎麼辦.
但記住.
主耶穌並不是去表演.
不會讓人看.
你多些人來看.
好,現在我就做事.
不是.
耶穌不等候這些人.
然後即時斥責這個鬼.
馬芳第九章第25節.
耶穌看見眾人都跑上來.
就斥責那污鬼說.
你這聾啞的鬼.
我吩咐你從他裡頭出來.
再不要進去.
耶穌即刻斥責那隻污鬼.
這隻鬼很無奈.
一定要離開.
在這裡我們看到.
代禱的功效很重要.
你為所愛的人.
你帶他們去禱告.
促使他們更加經歷神.
你自己的屬靈生命.
都會得到突破的.
所以都鼓勵弟兄姊妹.
要多些為你們的牧者.
為我祈禱.
祝福我的事奉.
這是很重要的.

$^{761}$很簡單.
我接下來這個星期.
都要去不同的地方.
星期六在聶密普.
星期日又去到舍維爾.
然後又去到聶密普.
另外一個地方港島.
下個星期更厲害.
就會去到Exeter.
較遠的地方.
即是來回.
四個多五個小時.
如果沒有意外.
求神幫助.
有這份力量.
其實是需要你們為我禱告的.
主的祝福.
是超過了這個慈父的所想.
慈父只是求.
他的鬼離開.
但是主耶穌給了一個更大的保證.
就是說.
再不要進去.
封了這些鬼的後路.
不讓他再回頭.
不容讓這個孩子再受著搞擾.
他的恩典.
往往是超過我們所想所求.
主耶穌不單單給了我們拯救之恩.
更加有保守的恩典.
淋到這個孩子的身上.
我們有沒有為著今天我們失禱的.
不單單我們得著那份拯救.
是得著神的保守.
得著神的供應.
得著神的磨練.
得著神的眷顧.
而去感恩.
好了 我們看看.
當耶穌斥責這個鬼的時候.

$^{801}$嘩!又有什麼事情發生?.
《馬福地九章》第26節.
那鬼喊叫.
使孩子大大的抽了一陣風.
就出來了.
孩子好像死了一般.
以致眾人多半說他是死了.
主一呼喝.
那隻鬼就即時的離開.
但他還作出最後的一擊.
但那個孩子突然間抽得很厲害.
我們看到那些旁人.
就即時作出回應.
大部分人只是看表面.
就說:哎呀!死了!死了!死了!.
這個孩子死了.
作出這個即時的判斷.
就是說主耶穌說你沒用.
甚至害死這個孩子.
今天有不少的閒人都是這樣的.
正如有些人在某些明星裡面.
以前那些八卦雜誌.
這個明星是怎樣?那個明星是怎樣?.
與你何干?.
還有很多不是真實的新聞.
所以今天他們錯了.
有沒有留意到?.
多半的人說他死了.
一會兒我們看到這個小孩子沒有死.
這班人有沒有即時向耶穌道歉?.
對不起.
剛才看錯了.
不好意思.
沒有的.
這個就是很大的無奈.
今天有些人傳一些謠言.
但後來確認了原來他說的話是不真實的.
他會不會澄清?.
他會不會向當事人道歉?.
是不會的.

$^{841}$這就很可悲了.
我們看到主耶穌沒有理會這些錯誤的判斷.
是滿有愛心的.
扶起這個小孩子.
是用事實去否定了這些無知的回應.
主耶穌再一次彰顯了祂的力量.
撒旦一切行動都不能勝過主.
接著我們看到主耶穌發出一個很無奈的呼喊.
當主面對門徒的失敗.
旁人的嘲諷.
和那些只顧著看熱鬧的群眾.
但不肯毀身.
主耶穌也是這樣.
祂也發出無奈的呼喊.
剛才我們讀過了.
我們再看一次.
保安第九章第十九節.
耶穌說:愛不信的世代.
我在你們這裡要到幾時呢?.
我忍耐你們要到幾時呢?.
把他帶到我這裡來吧!.
留意這句話絕對不是針對他的門徒.
是針對當時的世代.
他再次引用舊約摩西對子民的形容.
舊約申命記三十二章第四到第五節.
和第二十節裡有這個記載.
他是盤石.
他的作為完全.
他所行的無不公平.
是誠實無礙的神.
又公義又正直.
這乖僻彎曲的世代.
向他行事邪僻.
有者弊病就不是他的兒女.
我要向他們演練.
看他們的結局如何.
偷換本是極乖僻的族類.
心中無誠實的兒女.
記不記得摩西都說.
你們是彎曲乖僻的世代.

$^{881}$行事邪僻.
所以沈華道真的很無奈.
不要忘記.
和耶穌所面對的處境是很一致的.
耶穌行過很多神蹟.
但是人們都沒有去委心比心.
當子民在抗疫的時候.
看到神的作為.
第一代的人特別經歷了十災.
如何過紅海.
如何擊盤石出水.
他們當時仍然在享受.
運柱遮蓋了熱切的陽光照射他們.
令他們不會受傷.
夜晚有火柱.
可以讓他們晚間都可以如常活動.
是不是?.
也都每一天在享受嗎哪的供應.
除了安息日.
但是安息日之前.
多拿一份又不會發臭.
但他們仍然不信靠神.
常常質疑神.
埋怨神.
甚至對神不忠.
所以今天我們會不會.
同樣犯了這個錯.
其實我們在經歷神很多恩典.
但我們都會埋怨.
而之前主耶穌指出.
這法利賽是民事不斷的挑戰.
祂都用「世代」這個詞.
去描繪祂.
這一段經文我很大的感受.
我們一起來看路加文第七章.
31到35節.
這裡很清楚的就是.
你隨時都會被人說你.
罵你.
如果你聽得多少人說呢.

$^{921}$他說到今天.
屍骸若寒.
去禁食.
吃蝗蟲野物.
那些人就說.
古老石山.
你都被鬼附了.
但主耶穌來的時候.
和那些瑞利妓女.
分享天國的福音.
和我們吃飯的時候.
就為耶穌貪酒貪食.
即是怎樣呢?.
你說甚麼你都會被人說.
但是很重要.
智慧之子以智慧為事.
很簡單舉個例子.
很多人都喜歡我那種講道的方法.
我一定是解經的.
然後就會應用在今天的生活裡面.
否則是沒有意思的.
但如果你講到今天的生活的時候.
可能你又會被人說你.
講政治呀.
當你完全不講這些的時候.
只是隔開肅若.
但那些人又說你離地.
我剛剛講車公.
車天車地.
有一個人留言.
說我幼稚.
講這些.
好好的傳福音.
後來我都回他的.
我用了歌羅西書去回他.
我們要用主觀的智慧.
將神的道向別人講解.
我就問他.
我在這段的分享都在講我們的信仰.
我們的神才是真實.

$^{961}$不是那些偶像的敬拜.
我們會不會是你的偏見.
以致你怎樣呢?.
聽而不聞呢?.
今天一定有不同的人對你有不同的評價.
但是求神幫助我們.
我們聽到別人對我們的挑戰.
我們不是馬上去反擊.
而是去反思究竟他們講得對.
還是他們誤解.
還是他們有偏見呢?.
好了.
如果我覺得我是對的.
在神面前無愧.
就繼續堅持.
但智慧之智.
以智慧行事.
所以耶穌感到很無奈.
這次連跟隨他的門徒都沒有信心.
如耶穌所說.
這是一個不信又彎曲薄謀的世代.
這會不會是現在這個世代的寫照呢?.
所以耶穌說我要忍耐你們要到何時呢?.
這個忍耐是包含了神的審判.
我們看看二賽亞書第一章第十三節.
你們不要用再獻虛浮的供物.
香品是我所憎惡的.
月初和安息日.
並宣召的大會也是我所憎惡的.
作罪孽又守嚴肅日.
嚴肅會我也不能容忍.
你聽到嗎?我要忍耐到何時呢?.
這班人表面有宗教的情操.
但生活卻是完全行不出應有的操守.
甚至我怎樣來對付你們呢?.
所以求神幫助我們.
期望我們的主不會在這一刻發出這種無奈的嘆息.
最後主耶穌的教導.
主耶穌不單單責備.
也會作出教導.

$^{1001}$這是很重要的.
門徒希不是一早已經領受了權柄.
去治病趕鬼嗎?.
但為什麼這次不能夠有效地趕出這個鬼呢?.
他們不明白.
所以他們向主來求問.
這種求問的心.
是值得我們來效法的.
於是主耶穌很清楚地來回應.
馬太福音第十七章二十六節.
「耶穌說:是因你們的信心少.
我實在告訴你們.
你們若有信心.
將一粒芥菜種.
就是對這座山說.
你從這邊來到那邊.
他也必來去.
並且你們沒有一件不能作的事了.
至於這一類的鬼.
若不禱告禁食.
他就不出來」.
原來主耶穌指出.
他們是心若力微.
他們的信心脆弱.
所以他們的力量自然會衰微.
信心和行動的果效往往是成為正比的.
有大的信心.
就能夠為主作大事.
當然信心並不是一下子就可以成功的.
是需要對神有更多的認識和信靠.
就好像一粒芥菜種那樣.
最終會成為一棵樹.
慢慢地成長.
我自己回看《命恩堂》.
今年第四十年.
自從我離開的時候.
已經由第一次只做了四個人.
去到二千多人的六間教會.
這個絕對是甚麼呢?.
不只是我的力量.

$^{1041}$而是弟兄姊妹的困操困力.
和對神的信心.
所以信心的加增.
其中的途徑是甚麼呢?.
就是禱告.
甚至有時要禁食祈禱.
我們留意.
耶穌那一刻.
有沒有禁食禱告呢?.
你等等.
等我禁食祈禱完.
我再來趕鬼.
沒有.
主耶穌不只是講一些理論.
但是你不要忘記.
主耶穌經常和神有親密的交往.
特別是剛剛登山變賬的時候.
更加行上單獨與神相處.
所以在這時候.
他的生命充滿了神聖的力量.
但我們不是這樣.
所以主耶穌教導我們.
要多些禱告.
甚至遇到一些很大的困難的時候.
挑戰.
更加需要禁食禱告.
禁食禱告.
請你記住.
不是說我黑虎幾身.
我去鞭策自己.
神啊你可憐我.
意思是我不再專注在吃東西上.
然後又怎樣呢?.
專注在神那裡.
當我越親近神.
越與神相近的時候.
我們的信心自然就會增強.
當你的信心弱的時候.
你的力量就未弱.
但當你的信心加強的時候.

$^{1081}$你就能夠有這個力量.
所以求神來幫助我們.
所以為什麼我特別提醒大家.
特別我們在英國.
在香港生活應該沒有這麼大的壓力.
會多些親近神.
倚靠神.
有沒有一些言讀神的話語呢?.
這是很重要的.
所以求神來幫助我們.
美國臨北戰爭的時候.
北軍就連連打敗仗.
主張繼續保持奴隸制度的南方節節勝利.
在這個時候.
林肯總統呼召全國的人.
禁食一天.
這就是1863年3月30日.
結果因為這個禁食.
整個局勢扭轉了.
最後大公告成.
廢除了黑奴的制度.
林肯的一句名句是什麼呢?.
「我不求上帝站在我們這邊.
我只能求祂幫助我們站在祂那邊」.
還有現在我們在英國.
記不記得第二次世界大戰的時候.
英國就面對德軍的狂烈轟炸.
全國快要覆亡的時候.
邱吉爾首相做什麼呢?.
呼籲全國的人禁食一天.
一起禱告仰望神.
結果英國的神蹟出現了.
我們可以化危為安.
當年我母親還未信主.
但是已經病入膏肓.
時日無多.
我都經歷過很多次禁食禱告.
神憐憫我媽媽終於信了主.
時斯撒母耳帶領以色列人.
在米斯巴舉行全民性的特別禱告會.

$^{1121}$在神的面前來悔改.
結果耶和華就將他從強敵.
非理士人的手上拯救出來.
他們在勝利之後就立了一個石碑.
上面就叫「以便以謝」.
「以便以謝」意思是什麼呢?.
到如今耶和華都幫助我們.
所以當人有極大困難的時候.
來到神的面前禱告.
最終是要用神來作主.
神蹟就會出現.
我不是強調禱告的威力.
不是我們去操控我們的神.
而是我們盡上去最大的努力.
我們盡力去禱告甚至禁食禱告.
然後我們跟神有更大的親近.
我們就會知道神的心意.
神有他最終的決定權.
所以可以轉化.
今天我們都可能會面對門徒所面對的困難.
一些難關很難衝過.
困厄很難得勝.
我們甚至會情緒低落.
失望.
但我們更加需要主給我們力量.
我們要去到主面前重拾信心.
甚至跟神說.
主啊.
我信.
但我信不足.
求你來幫助.
這就是主耶穌教導我們禁食禱告的意思.
我們必定要有信心.
邪惡最終會被擊敗.
我們現在看一些社會上的事情.
好像邪惡的人得著能力.
不斷去欺壓一些異人.
但是我們願意依靠.
仰望.
對神有信心.

$^{1161}$即使我們好像不太行.
但是靠著神的力量和勇氣.
去面對這個屬靈的爭戰.
或者有時候我們會受挫.
未竟全功.
就好像門徒在那一刻不能趕走那些鬼.
但主必定會加以援手.
如果我們對生命每一步都有把握.
十拿九穩.
那還需要信心嗎?.
不需要.
所以今天我們要竭力進心.
使我們生命得以更生.
當我們面對困厄.
不再是心若裂眉.
而是心強力壯.
最後.
主耶穌的宣告.
很多時候我們看到這裡.
以為補了段落.
當我們看下去.
但不是這樣的.
我們看看聖經裡面.
《路加福音》第九章43,44兩節.
眾人都猜疑神的大能.
耶穌所作的一切事.
眾人正稀奇的時候.
耶穌對門徒說.
你們要把這些話傳在耳中.
因為人子將要被交在人手裡.
很多時候我們被聖經的分段誤導了我們.
我故意用回我那本《死海古經》給大家看.
這段經文記載到耶穌教完門徒.
接著有一個新的分段寫著.
「預言受難」.
就說剛才我們所讀的那段經文.
記住原來當時不少人很驚訝耶穌所做的.
主耶穌得到別人的稱讚.
榮耀.
這個主很厲害.

$^{1201}$但耶穌立刻要轉移門徒的視線.
再一次告訴門徒.
我會受苦,受難,受死.
主耶穌不會轉移門徒的視線.
要他們集中在受苦,受死裡面.
不是假大空的.
主耶穌絕對不會讓人有盲目的樂觀.
盲目的跟從.
再一次提醒他們.
他們信念的彌賽會被出賣.
但事實上就是這樣.
是求神叫我們不要只看表面的榮耀.
要注意所付的代價.
好,我們繼續聽我們祈禱.
很願意神你的話語成為我們的提醒和幫助.
甚至有時我們面對困境的時候.
我們承認我們有時信心不足.
但我們就向主你呼求.
主啊,我信.
但我信不足.
求你來賜賜憐憫和幫助.
叫我們洞悉到魔鬼的詭計.
好時我們都會來打落水狗.
但讓我們在困境當中.
我們不會失去對你那份盼望和信心.
願意天父你自己親自向我們說話.
也祝福到每一個不單單參與查經的.
如果在網上他們在YouTube再看的時候.
都是祝福他們.
讓我們有所領受和得著.
多謝你.
聽下者面前的禱告交托.
拉奉主德聖的名.
阿們.
\newpage



\section{加拉太書 5:1}
\label{sec:031699}
\textbf{真自由 (華)}~~~~葛國光牧師
\newline
\newline
連結: \href{http://www.church.com.hk/acms/content.asp?site=cdc\&op=show\&type=product\&code=031699\&layout=sermon}{\texttt{031699}} ~~~~ 語音日期: 2026-02-22
\newline
\newline
\hyperref[sec:031698]{\small{< < < PREV SERMON < < <}}
~
\hyperref[sec:index_chronic]{\small{[返順時目]}}
~
\hyperref[sec:031700]{\small{> > > NEXT SERMON > > >}}
\newline
\newline
加拉太書 5:1
\newline
\begin{longtable}{cl}
\hline
\hline
章節 & 經文 (和合本修訂版)\\
\hline
5:1 & \begin{tabularx}{0.7\textwidth}{X} 基督釋放了我們，為使我們得自由。所以要站穩了，不要再被奴隸的軛挾制。 \end{tabularx} \\ \\
[1ex]
\hline
\hline
\end{longtable}
$^{1}$我們今天要講到真自由.
對我們來講 在北美的人來講.
我們不太能夠體會自由.
可是如果你看看在伊朗現在的人.
我們都知道最近伊朗因為示威的關係.
死了很多的人.
到底多少人.
現在其實也沒有辦法真實地統計 對不對.
但是這一場的革命.
這一場的示威從哪裡開始呢.
是從有一個女孩.
叫做Masa.
這一個女孩.
她就是為了要能夠爭取不要戴頭巾.
最後結果被殺了.
這一件事情引起了整個伊朗.
不單單是女性 甚至男性.
很多人都起來抗議.
因為在1979年之後.
伊朗的革命他們就規定之後.
伊斯蘭教的極端主義.
他們就要求所有的女性不可以不戴頭巾.
所以你知道.
剛纔你看到葉門也是如此.
你只看到那雙眼睛.
對於一個女人來講.
她再怎麼樣漂亮的頭髮.
她沒有辦法被看到.
所以這對女人來講這是一個很大的壓抑.
不僅僅如此.
很多的這些 我沒有時間細講.
我要講的意思就是說.
伊朗的人在那邊說我要自由.
對他們來講自由是這麼樣子的寶貴.
所以他們那時候喊出一個口號.
女性生命自由.
你看那英文就是寫著.
Women, Life, Freedom.
今天在我們的美國裡面.
在今天這個世代裡面.

$^{41}$很多的人在講自由的時候.
他認為什麼是自由啊.
什麼叫自由啊.
自由就是.
以前在1990年代.
臺灣有一個歌星叫做陳明依.
聽過嗎.
沒聽過.
但是她在一個廣告裡面.
替一個檸檬蜂蜜的一個飲料做廣告.
她在廣告裡面就喊出了一句話.
只要我喜歡有什麼不可以.
聽過嗎.
聽過沒有.
你可能沒聽過這個陳明依.
但是你聽過這句話沒有.
臺灣來的聽過沒有.
李明依.
李明依你們就認識了吧.
我給她改了姓.
只要我喜歡有什麼不可以.
這其實當時很流行.
也代表著現代的這一代的年輕人.
他們有一種.
他們覺得說只要我喜歡就可以了.
這一方面強調他的個人主義.
另外一方面也代表著.
現在很多的人對自由的一個看法.
就是隻要我喜歡有什麼不可以.
我想幹什麼就幹什麼.
中國大陸在2015年的時候.
有一個荷蘭實驗中學的一個心理學的老師.
他突然有一天提出辭呈.
辭呈上面只有十個字.
就是"世界那麼大,我想去看看".
這句話一下子也哄騙中國.
那他的意思是什麼呢.
他覺得在這一種名校.
在這種比較好的學校做老師.
在這種壓力的競爭的情況之下.

$^{81}$他覺得他想要有一個自由的生活.
你懂意思嗎.
世界這麼大,我想去看看.
他不想要再被綁在這個地方.
他不想在這個壓力的環境之下.
天天為了五斗米而折腰了.
那麼到底什麼是自由.
中國大陸後來又流行一句話.
最近的,不曉得你有沒有聽過.
他說"世上無難事".
怎麼樣,只要肯放棄.
其實這句話也說明了中國的內卷.
在這種內卷之下就產生了躺平.
就產生了一種覺得說.
不需要這麼辛苦了.
不要這麼努力了.
降低一點物慾吧,這好一點的.
降低一點的慾望,過一點簡單生活.
有腦子可靠,就靠腦子.
沒老爸可靠怎麼辦呢.
生活簡單一點,不需要那麼樣子的.
不一定要高興,生活就好好的過就好了.
"世上無難事,只要肯放棄".
什麼是真正的自由.
在美國在開國的時候.
你知道美國跟英國這個獨立戰爭.
有一句話非常的重要.
這句話叫做"不自由,勿凌死".
這句話是Patrick Henry在1775年.
3月23號舉行的.
Virginia的公約會議裡面.
他在St. John教堂裡面發表的演講.
在其中最後的一句話.
他就講出這句話.
他說什麼.
他說"Give me liberty or give me death".
他當時針對的就是.
面對到英國的.
因為他們是殖民地.
美洲是殖民地.

$^{121}$所以他們想要獨立.
他們不想再要交稅給英國.
不想要再受到英國的壓制.
所以他發出這句話.
這句話就激勵了.
當時候美國的這些移民.
他們起來發動獨立的戰爭.
我們看見不管是政治上面.
軍紀上面或者這個世界上面.
我們看見人們都渴望有自由.
對不對.
但是這些自由到底什麼是自由.
今天在美國.
我們自由到一個地步.
我想幹什麼就幹什麼.
我今天想變男的我就變男的.
我今天想變女的我就變女的.
我今天或男或女都可以.
我心情好的時候我就扮個男裝出去.
我心情不好的時候扮個女裝出去.
可不可以啊.
我穿褲子的時候我是男生.
但是我穿個裙子可不可以.
可以.
見怪不怪了.
第一次看到兩個男生在那邊親嘴的時候.
我都不敢看.
人家在那邊大大方方的親.
這就是美國.
以前我們一聽到大麻毒品.
我們說這個人完了這一輩子.
現在你家樓下就有人在抽大麻.
現在你可能走過一個地方就聞到大麻味道.
大麻合法化現在已經變成是趨向.
LGBTQ合法了.
這個世界已經搞不清楚到底是什麼叫自由了.
今天主動化怎麼說.
今天就單單這一節.
這一節非常重要的一節聖經.
加在太書裡面.

$^{161}$保羅在這裡第五章第一節說.
基督釋放了我們.
叫我們得以自由.
基督釋放了我們叫我們得以自由.
什麼叫做自由.
保羅在講這件事情講的這個自由.
他講到的是.
當耶穌基督為我們的罪定在十字架上面.
他把一切的律法的咒詛都除去了.
他把一切在我們身上的罪惡都赦免了.
我們今天應著耶穌基督.
我們信了耶穌之後.
我們就得救了.
我們就得釋放了.
我們就從罪惡的捆綁裡面釋放出來.
我們不再做罪的盧浦了.
這就是保羅講的基督釋放我們的真自由的意思.
親愛的同志們.
一個人可以有所謂的政治上面.
經濟上面講話你隨便講.
你要罵總統你要罵什麼人都可以.
在美國是這樣子是不是.
自由嗎.
有自由嗎.
但是可能你嘴巴很自由.
但是你的心非常不自由.
你被憂鬱所捆綁.
你被苦毒所捆綁.
你被這許許多多的這些東西捆綁在你的心裡面.
是不是.
所以在這裡.
他說基督釋放了我們.
教我們得以自由.
然後接著他說.
所以要站立得穩.
不要再被盧浦的惡挾制了.
所以什麼是真正的自由呢.
真正的自由.
第一個.
不再被律法挾制.

$^{201}$跟你旁邊說你不要再被律法挾制.
保羅在這裡講律法的時候.
記得我們整個加勒臺書在講這件事情.
我們需要了解這個當時候的背景.
因為當時候在十讀經在十五章記載著.
當時候.
猶太的這些基督徒從母會耶路撒冷教會.
他們說我們是母會來的.
母會來的.
這很大.
他來到了加勒臺教會.
加勒臺教會是一個外邦教會.
他說你們這些外邦人.
你們這些非猶太人.
你們信耶穌不夠.
你們還必須遵守摩西的律法.
你們必須要受割離.
保羅聽了之後.
他嚴厲地寫了這一封信.
加勒臺書你從頭讀到尾你可以看見.
這是保羅所寫的書上最嚴厲的一種.
他責備這些人.
他說你們今天所傳的是別的福音.
耶穌基督為我們罪釘在十項上面了.
我們不再是因為靠律法而得救了.
今天你竟然還教我們要遵循摩西的律法.
要受割離.
兄弟們.
這就是為什麼保羅在這裡說.
不要再被奴僕的惡挾制.
他指的是這個律法的惡.
我們看看惡.
看過牛嗎.
看看這個圖.
兩個牛要一起耕田的時候.
要放一個什麼東西.
要放一個惡.
所以保羅曾經用這個來做一個比喻.
信與不信不能同付.
那個意思是什麼.

$^{241}$就是你信的跟不信的.
這邊不僅僅指婚姻.
很多人誤會.
以為只有指婚姻.
也就是說.
你嫁娶都要在基督的裡面.
基督徒你要找基督徒.
不單止婚姻還指什麼.
任何這種親密的關係.
通通都適用在這裡.
合夥.
跟別人合夥.
你信的跟不信的合夥.
價值觀不一樣.
你要把這公司的營運十分之一奉獻出來.
他說十分之一.
開玩笑.
賺錢不容易.
還奉獻十分之一.
為這個問題你就吵不完了.
他說我們今天應該要有兩本帳.
你說基督徒怎麼可以有兩本帳呢.
我們應該是帳目清清楚楚的.
這些很多的價值觀不一樣.
所以信跟不信的不能同付.
那個二是什麼.
一隻牛往前走.
另外一隻牛怎麼樣.
我就不走.
可以嗎.
他把你綁住了.
所以二.
保羅在講這個二的時候.
在《使徒行傳》15章.
他講到這些猶太的基督徒.
猶太的基督徒來傳這個律法.
要得救.
必須還要受割離的時候.
他就講說.
他說現在你為什麼試探神.

$^{281}$要把我們祖宗和我們所不能負的二.
放在門徒的頸上呢.
就是說,這個二.
是律法所帶給我們的這個二.
你為什麼還要讓這些外邦信徒.
他們還要去揹負這個二呢.
不要再被這些律法來挾制了.
這個二.
本身就是一個捆綁.
猶太的拉比.
為了把摩西的律法.
把它歸納出來.
一共有613條.
這次我去看過.
我說那613條真是一個縱惡.
我連可力啉用30分鐘都講不完.
這613條的誡命.
那麼後來.
耶穌當年不是責備法利賽人嗎.
記不記得.
耶穌為什麼責備法利賽人.
還記不記得耶穌責備他們.
好,其中一個.
是什麼.
他們後來又發展了很多.
這些所謂口傳的律法.
律律加利律.
很多很反水的.
舉例來講.
他們說你在進餐之前.
你要先按照古人的遺傳洗手.
古人的遺傳洗手怎麼洗你知道嗎.
你看看他.
洗手要從手腕這個地方.
有沒有.
這樣洗.
從你今天回家試試看.
喫飯之前先去洗個手.
就照著這個方式去洗.
每一餐這樣子洗.

$^{321}$除了這個之外.
他還有規定.
比如說安息日.
你安息日的.
你只能夠走2000步.
請問現在一般.
你們那個手機裡面告訴你說.
一天要走多少步.
1萬步.
你們標準比較高.
我的印象好像只有5000步就可以了.
1萬步.
安息日不可以.
安息只能走2000步.
好了.
安息日呢.
你如果去治病.
去看病.
你只能夠治療到.
可以存活的階段.
你不能夠把他治療到.
完全康復的階段.
因此類的這些規定.
生活上面的細節奉獻十分之一.
連回鄉.
耶穌不是責備法利賽人說.
你什麼叫人家回鄉.
你回鄉喫過嗎.
有沒有喫過回鄉餃.
回鄉看過嗎.
回鄉很清對不對.
他說你們叫他連回鄉.
也要奉獻十分之一.
就是這麼細微的東西.
你算得清清楚楚.
這個叫做律法主義.
這個叫做遵行律法.
這個叫做被律法捆綁.
你說牧師啊.
那個跟我沒關係.

$^{361}$我也不是猶太人.
我現在也不照那樣做.
那我告訴你.
現代的律法是什麼.
什麼是現代的律法.
我是成長在一個傳統的教會.
是不是.
我成長在一個傳統教會.
我們成長在一個教會里面.
我們有很多的一些教會的教導.
這些的教導裡面告訴你.
這個可以做.
那個不可以做.
等等等等的這些.
親愛的同志們.
可能本來的用意都是很好.
但是不知不覺.
久了之後.
它就變成是一個形式.
它變成是一個傳統.
然後大家就照著這樣去做.
你不這樣做.
你就覺得你不夠聖潔.
你不這樣做.
你就不夠屬靈.
你懂這意思嗎.
牧師聚會當中.
我們教會也有.
來 請問.
我們有一個牧導.
是不是 牧導.
牧導做什麼.
牧導 你的瞭解是什麼.
牧師禱告.
是不是.
那沒牧師.
傳道.
那沒傳道呢.
長老.
好 那我請問你.

$^{401}$為什麼要牧導.
當然你如果去讀.
這個神學裡面有崇拜學.
你都可以講出一套東西來.
這個都是在教會歷史當中.
發明出來的.
發展出來的.
但我要問一件事情.
聖經有沒有告訴我們說.
我們在聚會的當中.
要有一個牧導.
有沒有.
沒有.
你覺得在能量堂有牧導.
你到別的教會去.
沒有牧導.
你去到英文堂也沒牧導.
英文堂是不是比較不樹林.
你懂我意思嗎.
可是我們可能在一個教會里面.
習慣之後.
我們覺得說一定要有一個牧導.
今天沒有牧導.
牧師沒有為我們禱告.
或者你的解讀是.
牧師要為你們禱告.
是不是.
很多時候我們把這些.
一些傳統的.
或者是人家告訴我們這樣做.
我們就變成.
你說我沒有教會背景.
我沒這些問題.
我們中國人有沒有一些的這種.
類似像所謂的.
律法的一些素質.
瑕疵.
有沒有.
有 舉例.
想想看.

$^{441}$很多.
不知不覺你從小到大.
你都被這樣教導.
是不是.
這些的傳統.
這些的遺文.
它好像是猶太人的那個律法一樣.
瑕疵著我們.
我們中國人就是從小就被教導說.
萬般皆下品.
唯有什麼.
讀書高.
所以你讀書好.
你這個進名校.
你是北大畢業.
你是臺大畢業的.
你在從小.
你一路上面.
都是被人家所尊敬的.
有沒有.
如果你大學也考不上.
你啥也不會.
你只會做生意.
你初中畢業.
你就去做生意了.
但你生意做得很好.
請問.
可不可以.
我認識一般我的學生.
溫州人都是如此.
可是在他們的心裡面.
永遠有一種萬般皆下品.
唯有讀書高.
這個中國人的一個觀念.
在瑕疵他.
所以今天他可以有機會讀神學了.
他也功成名就了.
他現在唯一差的就是什麼.
學位.
所以他開始來練神學.

$^{481}$他練完了本科.
他繼續要練碩士.
碩士練完之後練什麼.
博士啊.
加持加持加持.
我們不知不覺被這些捆綁.
主的話告訴我們說.
基督釋放了我們.
叫我們不要再被.
盧浦的惡所瑕疵了.
跟你旁邊說.
不要再被盧浦的惡所瑕疵了.
這個盧浦的惡.
也可以指的是.
罪.
耶穌在講到瑞士的一段話的時候.
耶穌這樣說.
他說.
我們來看看.
犯罪的就是.
罪的盧浦.
有沒有.
犯罪的就是.
罪的盧浦.
神的兒子.
天父的兒子若叫你們自由.
你們就真自由了.
這段經文.
裡面後面這樣說.
所以這裡告訴我們說.
如果我們是犯罪的.
我們就是.
罪的盧浦.
在我們裡面.
有沒有這裡所說的淫亂.
誤會.
邪情 惡欲.
貪婪 苦毒.
憤恨.
當我們在讀這些的時候.

$^{521}$可能你覺得 哇 保羅講的這些東西.
都是很糟糕的.
我不會有.
淫亂 誤會先不提.
邪情 惡欲你可能也不承認.
但是請問有沒有貪.
不要講貪婪.
我覺得這個比較嚴重.
會不會貪 有沒有貪心.
有沒有.
我通通都不貪的.
股票市場跟我沒關係.
薪水多少.
也沒有跟我沒關係.
我一點都不貪.
我人生就是百爛.
不會吧.
你心裡面還是會想.
多 多 多 多一點.
喫Buffet就是這樣的心態.
對不對.
我們會有貪心.
我們會不會有怒氣 會不會.
會不會 我從來不發怒.
我太太絕對不需要我發怒.
我眼睛一瞪她就乖乖了.
會不會發氣.
會不會發脾氣 會.
是不是 不是隻有弟兄會發脾氣.
男人會發脾氣 女人會不會發脾氣.
哎呀 那個蹬起來.
哇.
先生馬上就乖乖的了.
全家大小通通都乖乖的.
不講話.
我們不覺得這些問題啊.
我們不覺得這些問題.
因為從小我看到我爸爸媽媽.
我太太一發脾氣的時候 全家都乖乖了.
所以我從此就學會一件事情.

$^{561}$怎麼樣.
只要我一發脾氣 你就給我乖乖的.
是不是.
這個就是.
罪.
保羅在這裡講的.
這是罪.
腦恨 憤怒.
喧鬧.
哇 這一些東西.
這裡面還提了很多這一類.
我們可能不是那種.
手殺人放火.
但是我們都有這些東西.
保羅說這一些.
你要棄絕它.
不要再被它挾制.
我過去.
在國內一個神學班裡面.
在這邊教的時候.
有一個.
女學生.
年輕的一個孩子.
她來參加.
神學班的時候.
差不多才20歲左右.
但是她跟我講.
她怎麼信主的 我問了她.
她說.
她還沒有信主的時候.
那時候流行叫網吧.
有沒有.
她說她去網吧.
常常一去網吧.
就在那邊迷戀.
整天就泡在那邊.
她說她最高的一個記錄是.
她曾經有三天的時間.
不眠不休的.
就在那邊網吧裡面.

$^{601}$打電動.
親愛的丟子妹.
這樣子的一個人.
可以被釋放嗎.
剛纔主的話不是說.
基督釋放了我們嗎.
對不對.
主釋放了她.
很奇妙.
神釋放了她.
不僅僅讓她脫離網癮的捆綁.
而且.
後來她信了主奉獻出來.
做傳道人.
哈利路亞.
丟子妹你都沒有很興奮啊.
丟子妹你沒有很興奮啊.
這是福音.
這就是福音.
這就是耶穌基督的福音.
耶穌基督來要釋放被虜的.
要使受壓制的得自由.
Amen.
親愛的丟子妹.
耶穌基督來.
要來釋放我們.
這個自由纔是真正的自由.
今天有的人譏笑我.
我當年在工作的時候.
那麼我們工地裡面.
大概都要抽菸的.
都要喝酒的.
還有一個東西.
臺灣裡面還多一個東西.
都要喫檳榔.
每個人嘴巴都紅紅的.
我經常會碰到.
這些老闆.
這些的工人.
或者是這些遞煙給我.

$^{641}$遞酒給我遞檳榔給我.
我一概都說不不不.
有的人就這樣不會給你了.
他會用一些話來激你.
有一次有一個老闆.
塞了一個怪.
那叫什麼.
怪手叫什麼.
怪手知道.
挖土機的老闆.
他遞煙給我.
我說我不抽.
他講了一句話.
他用臺灣話講.
翻譯成普通話.
意思就是說.
他說不會抽菸的長不大.
不抽菸不會長大.
這句話很刺激人喔.
是不是.
可是.
他有抽菸的自由.
我呢.
我有不抽菸的自由.
他沒有不抽菸的自由.
今天很多的人.
沒有不抽菸不喝酒的自由.
沒有不吸毒.
不吸大麻的自由.
沒有不在淫亂裡面的自由.
他沒有辦法脫離手機的自由.
有沒有.
基督釋放了我們.
我們就得以自由.
天父的兒子教你自由.
你就得以自由.
所以信耶穌要帶來的.
一當一地有這樣的一個現象.
一個人如果信了耶穌.
他說他信耶穌.

$^{681}$他卻沒有改變.
他明明他裡面的耶穌不是活的.
約翰·威斯里.
曾經帶來英國大復興.
上帝使用他.
震動了英倫三島.
甚至這個福音傳遍了整個世界.
今天衛理工會.
巡禮會.
衛理中的教會遍佈世界各地.
威斯·約翰.
可是你想想看這個人當初.
他本來是一個.
在當時的教會里面.
他本來是一個.
那種乖乖牌.
他很努力地禱告.
每天就在起來早上禱告.
他也很努力地在那邊.
做一切的他所能做的.
去做好一個基督徒.
可是沒想到.
有一天當他在船上面.
看到有一批的.
摩拉維亞弟兄會的信徒.
當大風浪來臨的時候.
這群信徒一點都沒有懼怕.
還在那唱詩歌讚美主.
這時候.
威斯·約翰.
心裡面就.
發出了這樣子一個問號.
這一群人是怎麼樣子的一群人啊.
怎麼面對這樣子大的風浪.
他們不怕死嗎.
他自己怕得要死.
所以因為這樣子.
後來在1738年的.
一個早上.
8點30分.

$^{721}$他聽到.
有人在朗讀馬丁路德的.
《羅馬斯註釋》.
當他聽的時候.
神的靈就特別.
把那一段他所聽到的話.
就光照了他.
他形容說當時他突然感覺.
好像一陣的溫暖.
進入他的心裡面.
他立刻有一個把握.
有一個確信.
他知道耶穌基督是唯一的道路.
他知道耶穌基督就在他的裡面.
他從此知道.
他已經重生得救了.
Hallelujah.
兄子們.
從那一天開始.
威斯·約翰.
所努力傳講的一個信息.
就是音信稱義.
這一個.
神的僕人.
我要繼續講他.
除了基督.
當我們講基督要釋放我們.
什麼意思呢.
我們看另外一段聖經.
在哥林多後書第三章.
這一段聖經.
我們就讀前面的.
主就是那靈.
主的靈在哪裡.
哪裡就得以自由.
在這一節聖經裡面.
我們講到主的靈在哪裡.
主的靈是什麼呢.
主的靈就是聖靈.
聖靈在哪裡.

$^{761}$哪裡就得以自由.
今天當我們講到聖靈的時候.
教會存在兩種極端.
一個我稱它叫做.
牧原式的基督徒.
或者牧原式的教會.
去過墳墓嗎.
去過牧原嗎.
你進去的時候就感覺.
很沉重.
我們大家認罪.
我們現在大家.
安靜在主的面前.
哇.
那個氣氛很隆重.
然後唱詩歌的時候.
大家都好像沒有笑容.
手也不敢動.
兩隻手就放在那個地方.
拿個詩歌本在唱歌.
但是另外一個教會.
哇.
在那個地方炸.
然後呢.
整個聚會里面.
有人在那邊翻滾.
叫極端靈恩教會.
所以造成今天.
兩個極端.
但是現在又準備.
聖靈是誰.
聖靈就是主.
聖靈就是三位一體的神.
是不是.
所以他們同龍同等同尊.
神學一直正不正確.
正不正確.
為什麼當我們提到.
聖靈的時候.
我們就突然覺得.

$^{801}$哎呦.
好像聖靈這兩個字.
不要隨便提.
一提你可能就變聖靈恩派了.
劉子梅.
造成今天很多所謂的傳統福音派教會.
死氣沉沉.
一切主流的教會.
長老會 敬信會.
這些的主流教會 衛理工會.
都在衰退.
現在劉子梅.
主的靈在哪裡.
你就得以自由.
你要從這些的捆綁當中.
釋放出來嗎.
你要從這些成癮的當中釋放出來嗎.
你要從這些所謂不好的.
怒氣這些東西釋放出來嗎.
親愛的劉子梅 我們需要聖靈.
我們需要聖靈.
威斯利亞漢在講到.
他自己的經歷的時候 他講到這樣子.
他提到說.
他說有一天.
他跟60位當時候的尋道會的.
這些弟兄們.
一起在那邊禱告.
清晨三點鐘在那邊禱告.
禱告 禱告 禱告.
看見突然整個的.
聖靈大大的澆灌.
這些的尋道會的.
這些的弟兄們.
他們就在那邊大聲的禱告.
然後在那邊大聲的讚美神.
他們共同的在那邊發出聲.
神啊 你真是偉大.
你就在我們的中間.
永遠是我們的神.

$^{841}$他們就在那邊大聲的宣告.
這位神是那麼樣偉大.
親愛的劉子梅 從那一次之後開始.
神就開始.
使用這一批.
威斯利公會的這些人.
他們出去到處佈道.
當時候英國國教會.
非常的保守 非常的傳統.
把他們排斥在外 不讓他們在教堂裡面講道.
他們就到公園去講道.
到礦工去講道.
到礦坑裡面去講道.
哇 在這個街頭.
站在那個地方.
就是街頭佈道.
所到之處.
帶來了許許多多的罪人悔改.
甚至很多的礦坑.
那些工人.
臉上黑黑的.
聽到他講道的時候.
淚流滿面.
整個的臉是黑的.
只有兩道白的.
被那個淚水洗刷的.
兩道痕跡.
小偷 強盜.
聽到這些福音.
信了耶穌.
甚至有時候教堂裡面會看見.
有一群人就從外面走進來.
他們說我們要來信耶穌.
這些幫派的領袖.
帶著這些年輕人一起到教會里面.
信耶穌.
這豈不是我們所期望.
看見的教會嗎.
這豈不是我們今天.
所期望看見的基督徒嗎.

$^{881}$充滿著火力.
大聲的在那邊讚美神.
他們真是一群.
被聖靈充滿.
到處帶著聖靈的人也.
約翰·魏斯林那一年.
當他在帶領聚會的時候.
就有像今天一些.
靈恩教會里面.
所強調的東西.
比如說倒下去.
我沒有時間多講.
魏斯林·約翰在那邊講道的時候.
有好幾次.
他在講道的時候.
他開始就禱告.
開始呼召的時候.
纔剛剛講完道.
就看見一個倒下去.
一個倒下去.
而且不是.
他根本沒有下去為他們按手.
聖靈就這樣強烈的工作.
這一個一個的在會場各個角落裡面.
許多人就倒下去.
倒下去不是就倒下去.
有的人倒下去也沒幹什麼.
就是在那邊躺著.
好像以為被聖靈擊倒了.
有的人就在那邊流淚悔改痛哭.
有的人在那邊大聲的禱告.
最後喜樂的靈就從他那裡面出來.
這些人都遇見了聖靈.
對不對.
這是我們今天所需要的.
你渴慕聖靈的工作嗎.
但願聖靈在我們中間來工作.
\newpage



\section{撒迦利亞書 8:1-23}
\label{sec:031700}
\textbf{有神同在,恩福滿滿}~~~~黃馮潔霞師母
\newline
\newline
連結: \href{http://www.church.com.hk/acms/content.asp?site=cdc\&op=show\&type=product\&code=031700\&layout=sermon}{\texttt{031700}} ~~~~ 語音日期: 2026-02-22
\newline
\newline
\hyperref[sec:031699]{\small{< < < PREV SERMON < < <}}
~
\hyperref[sec:index_chronic]{\small{[返順時目]}}
~
\hyperref[sec:031701]{\small{> > > NEXT SERMON > > >}}
\newline
\newline
撒迦利亞書 8:1-23
\newline
\begin{longtable}{cl}
\hline
\hline
章節 & 經文 (和合本修訂版)\\
\hline
8:1 & \begin{tabularx}{0.7\textwidth}{X} 萬軍之耶和華的話臨到我，說： \end{tabularx} \\ \\ \relax
8:2 & \begin{tabularx}{0.7\textwidth}{X} 「萬軍之耶和華如此說：我為錫安而妒忌，大大妒忌；我為了它妒忌而大發烈怒。 \end{tabularx} \\ \\ \relax
8:3 & \begin{tabularx}{0.7\textwidth}{X} 耶和華如此說：我要回到錫安，住在耶路撒冷中間。耶路撒冷必稱為忠實的城，萬軍之耶和華的山必稱為聖山。 \end{tabularx} \\ \\ \relax
8:4 & \begin{tabularx}{0.7\textwidth}{X} 萬軍之耶和華如此說：將來必有年老的男女坐在耶路撒冷的廣場上，各人因年紀老邁而手拿枴杖。 \end{tabularx} \\ \\ \relax
8:5 & \begin{tabularx}{0.7\textwidth}{X} 城裡的廣場滿有男孩女孩在玩耍。 \end{tabularx} \\ \\ \relax
8:6 & \begin{tabularx}{0.7\textwidth}{X} 萬軍之耶和華如此說：在那些日子，即使這事在這餘民眼中看為奇妙，難道在我眼中也看為奇妙嗎？這是萬軍之耶和華說的。 \end{tabularx} \\ \\ \relax
8:7 & \begin{tabularx}{0.7\textwidth}{X} 萬軍之耶和華如此說：看哪，我要從日出之地、從日落之地拯救我的子民。 \end{tabularx} \\ \\ \relax
8:8 & \begin{tabularx}{0.7\textwidth}{X} 我要領他們來，使他們住在耶路撒冷中間。他們要作我的子民，我要作他們的神，都憑信實和公義。 \end{tabularx} \\ \\ \relax
8:9 & \begin{tabularx}{0.7\textwidth}{X} 「萬軍之耶和華如此說：你們的手要堅強；這些日子，你們已聽見先知的口，在萬軍之耶和華殿的根基立定、聖殿建造的日子所說的這些話。 \end{tabularx} \\ \\ \relax
8:10 & \begin{tabularx}{0.7\textwidth}{X} 那些日子以前，人得不著工價，牲畜也無人雇用；且因敵人的緣故，出入不得平安；因我使人與人互相攻擊。 \end{tabularx} \\ \\ \relax
8:11 & \begin{tabularx}{0.7\textwidth}{X} 但如今，我對這餘民必不像先前的日子。這是萬軍之耶和華說的。 \end{tabularx} \\ \\ \relax
8:12 & \begin{tabularx}{0.7\textwidth}{X} 因為他們要平安撒種，葡萄樹要結果子，土地必有出產，天也必降甘露。我要使這餘民享受這一切。 \end{tabularx} \\ \\ \relax
8:13 & \begin{tabularx}{0.7\textwidth}{X} 猶大家和以色列家啊，你們從前在列國中怎樣成為可詛咒的；照樣，我要拯救你們，使你們得福。不要懼怕，你們的手要堅強。 \end{tabularx} \\ \\ \relax
8:14 & \begin{tabularx}{0.7\textwidth}{X} 「萬軍之耶和華如此說：你們祖先惹我發怒的時候，我怎樣定意降禍，並不改變；萬軍之耶和華說， \end{tabularx} \\ \\ \relax
8:15 & \begin{tabularx}{0.7\textwidth}{X} 這些日子我也定意施恩給耶路撒冷和猶大家；你們不要懼怕。 \end{tabularx} \\ \\ \relax
8:16 & \begin{tabularx}{0.7\textwidth}{X} 你們所當行的是這樣：每個人要與鄰舍說誠實話，在城門口要按真正的公平來審判，使人和睦。 \end{tabularx} \\ \\ \relax
8:17 & \begin{tabularx}{0.7\textwidth}{X} 誰都不可心裡謀害鄰舍，也不可喜愛起假誓，因為這些事都為我所恨惡。這是耶和華說的。」 \end{tabularx} \\ \\ \relax
8:18 & \begin{tabularx}{0.7\textwidth}{X} 萬軍之耶和華的話臨到我，說： \end{tabularx} \\ \\ \relax
8:19 & \begin{tabularx}{0.7\textwidth}{X} 「萬軍之耶和華如此說：四月的禁食、五月的禁食、七月的禁食和十月的禁食，必成為猶大家的歡喜和快樂，以及美好的節期；所以你們要喜愛誠實與和平。 \end{tabularx} \\ \\ \relax
8:20 & \begin{tabularx}{0.7\textwidth}{X} 「萬軍之耶和華如此說：將來還有眾百姓和許多城鎮的居民要來。 \end{tabularx} \\ \\ \relax
8:21 & \begin{tabularx}{0.7\textwidth}{X} 這城的居民必到那城，說：『我們快去懇求耶和華的恩，尋求萬軍之耶和華；我自己也要去。』 \end{tabularx} \\ \\ \relax
8:22 & \begin{tabularx}{0.7\textwidth}{X} 必有許多民族和強盛的國家來到耶路撒冷尋求萬軍之耶和華，懇求耶和華的恩。 \end{tabularx} \\ \\ \relax
8:23 & \begin{tabularx}{0.7\textwidth}{X} 萬軍之耶和華如此說：在那些日子，列國中說各種語言的人，必有十個人強拉住一個猶大人衣服的邊，說：『我們要與你們同去，因為我們聽見神與你們同在了。』」 \end{tabularx} \\ \\
[1ex]
\hline
\hline
\end{longtable}
$^{1}$弟兄姊妹平安.
有神同在 恩福滿滿.
今年我自己對神的同在有很大的反省.
還有很多心思.
究竟神同在是什麼意思呢?.
究竟我們經歷神同在.
有什麼事情會發生呢?.
今天和大家花點時間去看看這個題目.
當然有神同在 一定有恩典.
一定有福氣 福氣是屬靈的福氣.
但是有時候這些恩典福氣.
是在逆境裡面出現的.
不是在你事事如意.
好像我們現在新年.
恭祝人家什麼呢?.
心想事成 萬事如意.
什麼都是順的.
但是其實有時候在艱難裡面.
經歷神的恩典更加寶貴.
你更加會珍惜.
更加會了解人.
還有更加可以給神大大的使用.
所以不要求一生都很順.
最好是無風無浪的.
無憂無慮的.
這個唱歌就唱得到.
但是真正實際的生活.
我和大家說.
是很多風浪.
很多掙扎.
很多挑戰.
和大家先來想一想.
直到現在你人生到了不管多少歲.
有誰是由你一出生懂事以來.
陪到你現在的?.
有的 因為媽媽好像Rita很有福.
她媽媽九十幾歲.
由她睜大眼做人的時候.
就看到她媽媽了.
有沒有人.

$^{41}$真的還有一些在你身邊的人.
一直陪著你走幾十年.
無論現在七十歲八十歲九十歲.
還有一個人在你身邊.
如果有的話.
你真是很蒙福很蒙恩.
好了 我們看看.
其實神很奇妙的.
他真的時常都希望.
親近我們人.
和多一些一起.
我們看看.
神從起初就樂意與人同在.
如果你讀《創世記》.
你會發覺神創造與人萬物.
去到最後那一天.
他就做人了.
先做了阿當.
做完阿當之後.
他們經常都見面.
神和他們經常都對話.
有時候如果不見了阿當.
他會問阿當你在哪裡.
就像我們如果養過子女養過孫就知道.
你做著做著或是去完街回來.
要找他們真的要找.
有時候他們不知道躲到哪裡去了.
但是神無論如何.
他很樂意與人一起.
就算當阿當犯了罪.
他真的躲起來.
自己知道很醜.
就做了一件樹葉的衣服.
當神叫他.
阿當你在哪裡.
當然那時候神已經體諒.
阿當一個人很孤獨.
就幫他做了一個配偶.
霞娃.
誰知道霞娃就.

$^{81}$唉 女人還是不要說那麼多.
有時候有人引你說話.
你就會踏上她的話題.
就說到沒雷公那麼遠.
甚至下到她的試探引誘裡面.
結果自己犯罪就好了.
還要回去找先生一起吃飯.
這個果很好吃.
你也吃一口吧.
誰知道那個就是神叫他們不要吃的果.
結果兩個人都躲起來.
但是神在伊甸樂園那裡.
風吹.
神在找他們.
阿當你在哪裡.
霞娃你在哪裡.
當人逃避神的時候.
神仍然是要去找人.
耶穌一樣.
祂說我來是要尋找拯救失喪的人.
所以神後來為他們兩夫婦做了一件皮衣.
皮的.
有皮做衣服一定是.
宰了羊 宰了牛的.
所以那裡有個救贖在那裡.
神為阿當 霞娃.
預備了贖罪的羔羊.
他們能夠穿著皮衣.
好過那些樹葉的衣服一陣風來就吹掉了.
而且.
當神祂做了人.
和人同在.
人是察覺到的.
就是說這個人是不同的.
亞伯拉罕.
以撒.
雅各.
約瑟.
他們這幾個族長.
無論他們去到哪裡.

$^{121}$無論他們種東西有一百倍收成.
或者他們說謊.
說太太是他的妹妹.
所以被人拿走了他的太太進去.
就是要娶他的.
無論怎樣都好.
你發覺.
神都仍然和他同在.
因為連那些外邦人.
看到他們這幾個族長.
都覺得他們是不同的.
他們說你有神同在的.
特別.
你知不知道.
亞伯拉罕不要說了.
他是信心之父.
以撒.
他被人欺負到真的沒地方站了.
你掘一個井.
是他爸爸以前掘的.
他找出來的.
那應該就是他用的.
誰知道這頭掘完的基拉爾人又來說.
是我們的.
又拿了他.
第二個井又是這樣說.
是我們的.
第三個井.
都想拿的.
但是神真的很賜福以撒.
他一路退讓.
退讓.
容忍.
你不給我嗎.
我又再去開發.
又有第二個井出來.
結果聖經說.
神與以撒同在.
神使得以撒的僕人種什麼都有.
一百倍收成.

$^{161}$所以那些敵人恨不得流口水.
快點來.
不要做敵人.
做朋友.
叫他幫他.
立一個約.
我不欺負你.
你也不要欺負我.
其實如果以撒聽到這句話.
你不欺負我.
你欺負了我很久了.
但是以撒說.
好吧.
我幫你簽了和約.
然後還請他吃了一頓飯.
你知不知道.
我說這些飯怎麼吃.
就好像詩篇23篇說.
在我的敵人面前.
你為我拜持嚴正.
吃得下.
他還吃完飯之後.
平平安安的送我的敵人走.
這種生命.
不是因為有神同在.
怎麼會做得出來.
我不和你做朋友.
不理你已經很好了.
我不去以惡報惡.
但你還能夠以善成惡.
還是祝福他們平平安安走.
有神同在的人.
旁邊的人看到的.
他會告訴你.
為什麼你這麼不一樣.
為什麼你好像和我們差一點.
你這麼平和.
為什麼你遇到驚險的事情.
你還是這麼淡定.
為什麼你病了.

$^{201}$別人就怨天尤人.
你還在猛感恩.
人們真的會問.
你為什麼會這麼特別.
因為有神同在.
可惜.
罪使人與神暫時分離.
這個情景大家很明白.
因為神是聖潔的神.
我們一犯罪.
不是神要不要離開你.
是你自己沒有那個面目.
很羞愧.
沒有面子見神.
就到處走 到處躲.
但神沒有錯.
祂不可以和罪一起存在.
有罪的地方.
神真的沒有辦法去容忍.
但是.
如果你去到新約.
你會發覺.
耶穌基督來.
祂不是一來就做一個聖人.
做一個超人.
一做就不是做Baby.
是做一個很偉大.
很壯的領袖.
祂不是.
祂來到是道成肉身.
與人同住.
所以道成肉身的意思是.
搭了一個帳篷.
在這個世界.
與我們一起住.
一起生活.
所以每次想到主耶穌的時候.
我覺得祂很偉大.
如果你住慣了一間很乾淨的房子.
我讓你.

$^{241}$譬如派你去.
你又很想去.
去做村教.
不要去得太久.
去一天.
去墨西哥.
你想想我們二十多年前.
在舊堂就已經開始去了.
我們有弟兄姊妹開Van.
下雨也照樣去.
我們真是厲害.
不改旗的.
那些墨西哥條灣的路.
沒有現在那麼好.
二十多年前.
震到你.
一有那些凹下去的地方.
那個Van一震.
你整個人頂上車頂.
還有那些車骯髒得不得了.
但我很欣賞這些弟兄姊妹.
真是風雨不改.
就是這樣用自己的Van.
出錢加油.
然後就帶著一班弟兄姊妹.
他們需要人車.
還有就是.
預備食物.
預備很多東西.
有一次還坐了一輛車.
那輛車爆胎.
在Freeway那裡.
真是多麼驚險.
又要一條一條的那些Lane.
要走到最近最近.
等人來到.
補胎即是換胎.
但他們一點怨言都沒有.
反而我自己和大家說.
真是很不好意思.

$^{281}$我第一次去到Tijuana.
你信不信.
我們去的是孤兒院.
那裡有宿舍.
但所謂宿舍都是睡地下的.
然後就聚會.
有個禮堂.
你知不知道.
我去到那個禮堂的時候.
他們都有排一些椅子在那裡.
我選不到一張乾淨的椅子坐.
真是好.
我覺得那時真的不知道為什麼.
那麼怕骯髒.
我看到那些狗好瘦.
小朋友跟著爸爸媽媽那些去.
他看到那些狗那麼瘦.
你知道美國的狗又肥又白.
養得很好.
他們真的好到怎樣.
不只是給熱狗的麵包皮給他吃.
是連裡面那條熱狗也給他們.
那些狗就跟著他.
不斷搖尾.
但當然真是很骯髒.
你不去洗手間.
去你也不想去.
但是我很欣賞那些人.
真是可以鍍成肉身.
我們就說很不方便.
又要戴遙控鏡頭.
又要找水洗.
那裡連乾淨水也沒有.
那些水很珍貴.
但那些年輕人不斷去.
幫人剪頭髮.
那些濕疹都跳出來.
跳到你身上.
有一個師母是做護士出身的.
她去到告訴我.

$^{321}$她說你知道為什麼一路跟那些人飛髮.
一路那些濕疹跳到我的身上.
但她說我不怕.
我只看到一塊一塊很可愛的臉.
無論小朋友大人.
他們沒有人跟他飛髮.
所以我們去幫他們飛髮.
一飛就全日了.
由長洲飛到我們走了傍晚.
我真是很欣賞那些人.
在美國這麼舒服的地方.
可以去到一些比較落後.
骯髒.
找不到一棵綠色的樹.
樹葉是黃色的.
一過了關你就發覺.
聖地牙哥和塞爾瓦那是完全兩個世界.
耶穌基督就是這樣住在.
我們這個這麼骯髒骯髒的世界裡面.
從嬰兒開始做.
一路童年 接著青年 成年.
這個是出於愛和真理給我們有恩典.
現在聖靈就時刻住在我們的心裡面.
真是的 你不孤單的.
你說我都見不到神.
都不知道神知不知道我的難處.
最近經常有人問這個問題.
你想一想.
他說我經常都犯錯.
經常都又來做那件不應該做的事.
我還是不是基督徒?.
我還是不是配稱為神的兒女?.
他就這樣問.
你們說呢.
如果你清楚認罪悔改一次.
願意接受耶穌做你的救主和生命的主.
那你就是基督徒了.
但你怎知道信了之後.
來了教會之後.
才發覺自己一樣好像以前那麼敗壞.

$^{361}$那些試探來到經不起那個引誘.
人家惹我生氣我都照樣反駁.
或者討個公道.
為什麼我沒有改呢?.
他就自己在那裡打個問號.
我還是不是基督徒?.
我配不配稱為神的兒女呢?.
我最記得何約瑟長老.
解釋他一個.
一點是我經常都記得的.
因為三十多年前和他一起在這裡同工.
我們最近收拾那些櫃子.
我發覺有好幾餅是何約瑟長老講道的.
他曾經講過一樣東西.
他說罪和軟弱是兩回事.
罪就是我們天生有罪性.
接著我們未信耶穌之前很放肆.
喜歡怎樣做就怎樣做.
任意妄為.
那些是罪惡.
但當我們去信耶穌的時候.
耶穌基督的寶血就洗乾淨我們的罪惡.
赦免了我們.
因為祂在十字架為我們流血.
而且祂還為我們祈禱.
祂說赦免他們 因為他們所作的 他們不曉得.
今天這些罪惡完全被耶穌清理了.
但如果魔鬼再來.
再來引誘你.
或者再來騙你.
說方說你不行了.
你又來做這件事了.
你要用一節聖經去擋住它.
什麼聖經呢.
《約翰一書》一章九節.
「他兒子耶穌的血也能洗淨我們一切的不義」.
耶穌基督的血可以洗淨我們一切的不義.
第九節就說了.
之前是這樣說的.
我們若認自己的罪.

$^{401}$神是信實的 公義的.
必要赦免我們的罪.
洗淨我們一切的不義.
所以今天我們的罪已經一次過.
被耶穌基督洗乾淨了.
那個的救贖是包括了過去 現今 將來.
所以很多人問.
如果一個人自殺.
他是基督徒 他自殺.
可不可以上天堂.
這是很爭議性的題目.
但當我們做了十幾二十年.
離婚和憂傷支持小組的時候.
我開始明白.
而且真正有人 真人地走進來我們那裡.
告訴我們.
我很不明白 我的主學老師.
靈性這麼高超.
他竟然跑到山頂那裡.
跳到懸崖那裡.
死了.
你叫我怎樣繼續去教會.
怎樣繼續去信耶穌呢.
為甚麼他都會走這條路呢.
後來這個「Church Initiative」這個機構.
它出了一些教材.
裡面很針對性就是說.
耶穌的救贖是很完全的.
過去 現今 將來.
所以當這個主學老師.
他做了一件這樣的事.
耶穌都會赦免他.
已經赦免了他.
而且是當他在心裡憂鬱.
或者是有其他的問題的時候做的.
耶穌更加了解.
所以當我們跟別人傳火的時候.
很多人是卡在過去一些事情.
你怎樣救他出來呢.
在我裡面安慰我們 鼓勵我們.

$^{441}$還有就是.
千禧年的神人永遠同在.
如果你讀啟示錄.
二十章 二十一章.
你真是很開心.
因為那裡說.
有新天新地.
那裡有神的障墨在人間.
而且沒有罪惡.
沒有黑暗.
沒有憂愁.
沒有疾病.
我有 我有.
還有沒有怨恨.
沒有罪惡.
那裡就是神和人永遠在一起的地方.
我想到耶穌基督.
他道成肉身的時候.
曾經在十字架上說了一句話.
「我的神啊 我的神啊 為甚麼要離棄我呢?」.
所以真是.
神有時真是會短暫地離開一個人.
當他背著很多罪的時候.
我們現在要來看一段聖經.
這段聖經是撒加尼亞書第八章.
都挺多節的.
但不要緊.
你知不知道我們看聖經 讀聖經.
有幾種讀法的.
第一 真是讀出聲的.
讀出來 一個個字.
不要勒的 一個個字讀.
第二 聽聖經.
眼睛不好.
聖經買到最大的字.
都沒有辦法看到.
那就要聽了.
第三 就是你一路讀 一路祈禱.
就是用裡面的字來祈禱.
譬如找到一句.

$^{481}$「不是依靠勢力 不是依靠才能.
乃是依靠神的靈方能成事」.
馬上就祈禱神 真是.
我們教會多大都好.
我們多少人都好.
我們多少同工都好.
我們不是依靠勢力 不是依靠才能.
是依靠神的靈 才有人信主.
神的靈 才有人真是樂意去事奉.
不是我的功勞.
不是那個群眾的功勞.
是聖靈.
好了 我邀請大家一齊來讀.
你讀的時候.
默想 盡量將它記在心裡.
弟兄們可不可以幫我們手讀這張slide.
一至三節 請.
耶和華講了一句說話.
如果你有機會數一下.
萬君之耶和華如此說在這章裡面.
出現了多少次你就知道.
這一章聖經全部都是耶和華自己講的.
現在是萬君耶和華.
即是說他是元帥來的.
他是為我們爭戰的.
他是得勝者來的.
他是君王來的.
他說現在我要為錫安.
即是為猶太人.
為耶路撒冷.
來到迴轉.
怎樣迴轉呢.
本來他離開了.
離開了才會有聖殿被毀.
城牆又沒有了.
火燒.
很多耶路撒冷都亂了.
因為被外邦軍隊毀滅.
因為裡面的人犯了太多罪.
但現在耶和華說.

$^{521}$我再次為我的城來火熱.
而且我要幫我這些漁民.
來懲罰外邦的國家.
後面有講.
其實神是憤怒他們犯罪.
但他說我是憤怒一點而已.
但你們罰他們這麼厲害.
即是說害這些猶太人這麼厲害.
神很憤怒.
而且最寶貴的是.
我現在要回到錫安.
要住在耶路撒冷中.
耶路撒冷就改變了.
不再是罪惡城.
稱為誠實的城.
萬君耶和華的山.
被稱為聖山.
甚麼都變成聖.
因為耶和華回到錫安.
請姐妹們幫我讀.
4到6節.
1,2,3.
萬君之耶和華如此說.
將來必有年老的男女.
坐在耶路撒冷街上.
因為年紀老馬就手拿拐杖.
城中街上必滿有男孩女孩玩耍.
萬君之耶和華如此說.
到那日這是在如聖的民眼中看為稀奇.
在我眼中也看為稀奇嗎.
這是萬君之耶和華說的.
你讀得很好.
這裡提到.
剛才那張slide.
就好像是過去的歷史.
預言.
但這裡是預言.
就是萬君耶和華如此說.
將來.
寫這本書的時候是未發生的.

$^{561}$可能現在我們都未曾完全發生.
他說將有年老的男女.
坐在耶路撒冷街上.
因為年紀老馬都是會手拿拐杖.
手拿拐杖不要覺得醜.
人到某個年紀.
或者是因為跌過.
因為有個病.
真的需要靠扶手.
但很好笑.
有時扶手變成一個.
安全杖.
好像很安全.
好像小朋友拿著安全杖.
那些被子.
不斷咬著拿著.
心裡就很安心.
我見過有弟兄.
經常拿著那支扶手杖.
但到他心急要做一件事時.
沒有拿著那支杖都可以做到.
即是我要走去那裡做一件事.
有時扶手杖不一定需要.
但到你需要時就很有用.
當你覺得不平衡時.
扶一扶它.
它就會令你平衡.
即是說.
當那天.
將來.
每個人都會.
不是七老八十.
八十不夠.
現在不知幾多老.
九老一百.
可能的.
現在全部人過九十.
甚至百幾歲.
一百一十幾都有.
你會發覺.

$^{601}$到時.
那些人.
是很長壽的.
還有在街上有很多男男女女玩.
你知不知道現象就是.
有長者的一代.
亦都有.
哇.
年輕的一代.
而且很多小朋友.
真是很興旺.
你沒有為教會祈禱感恩.
祈禱就是我們可以多些人來教會.
最後就多些年輕人來教會.
因為一有年輕人.
你就會有機會多些小朋友.
是嗎.
他們一定是一個個家庭來的.
我們祈禱求神祝福.
在這個末世的時間.
人認為很稀奇的事會發生.
就是教會很興旺.
以為老化.
不是.
原來老的還有很有用的.
真是托住.
支持住.
那些年輕人.
那些小朋友.
亦都是從小就認識神.
到老都不偏離神.
還有是彼此尊重的.
互相.
互相待託.
互相.
即是相愛的.
我真是發覺有時候做老人家.
做長者.
做得很辛苦.
因為我們一路是這樣讓給年輕人的.

$^{641}$他們要什麼嗎.
只要我們做到.
我讓給你.
但他們有沒有感恩的心呢?.
我們需要提醒一下他們.
尊重是彼此的.
鼓勵是彼此的.
支持也是彼此的.
這樣才會成為一個.
樂也容容的家庭.
或者是教會.
接下來.
弟兄們.
請你們再來.
七到九節.
請萬君之耶和華如此說.
我要從東方從西方救回我的民.
我要領他們來.
使他們住在耶路撒冷中.
他們要作我的子民.
我要作他們的神.
都憑誠實和公義.
萬君之耶和華如此說.
當建造萬君之耶和華的殿.
立根基之人的先知所說的話.
現在你們聽見.
應當手裡強壯.
這裡還有神應許什麼.
將來會從東方從西方救回他的民.
還有就是領他回耶路撒冷.
做他的子民.
他就做他的神.
這其實是一個約.
我們的神要和我們同在.
他是基於他是守約思前愛的神.
他不可以反悔.
所以他一直都是憑誠實和公義.
來做我們的神.
做我們的父.
還有神很注重什麼呢.

$^{681}$要建造聖殿.
你知不知道.
猶太人最驕傲的.
引以為傲的就是.
他們有聖殿.
他們有聖經.
就是那個律法.
有神的話.
所以我覺得這兩樣東西.
是他們很好的東西.
可以和神有關係.
建立很好的關係.
還有大家就有個集中點.
回來敬拜神.
回來學神的話.
好.
請姐妹們.
第十到十二.
哪日以先人得不著故駕.
生促也是如此.
且因敵人的緣故出入之人不得平安.
乃因我使眾人互相攻擊.
但如今我待者如聖的民.
必不像從前.
這是耶和華說的.
因為他們必平安撒種.
葡萄樹必結果子.
地土必有出產.
天也必降甘露.
我要舍如聖的民.
享受這一切的福.
哇 福氣來了.
之前這一節就說.
以前你被人欺負.
那些公價給你是很不公道的.
但是現在我要祝福你們.
你們不會像以前那樣被人欺負.
壓低你的工資.
甚至不給你工資.
做勞部哪有工資.

$^{721}$勞地沒有的.
但是現在我卻要祝福我的子民百姓.
讓他們撒種.
種什麼都有收成.
而且我又降雨給他們.
讓他們每一個都能夠享受一切的福氣.
今年的加州.
從去年年尾開始都很多雨水.
所以有沒有發覺.
每一次下完雨之後.
如果你有種東西.
那些葉又多了.
花又多了.
如果有果子出來.
那些果子都大過了.
還有很多汁.
因為有充足的雨水.
神正在賜福我們.
你知不知道呢?.
好了.
第八章十三到十四節.
應該輪到誰?.
弟兄了.
弟兄來了.
由大家和以色列家.
你們從前在列國中.
怎樣成為可咒詛的照樣.
我要拯救你們.
使人稱你們為有福.
細字讀完.
因為很有意思.
或譯什麼?.
使你們叫人得福.
你們不要懼怕.
手要強壯.
萬軍之夜如此說.
你們列祖惹我發怒的時候.
我怎樣定義和並不後悔.
所以這裡.
神應許一件事很好的.

$^{761}$將咒詛變為祝福.
還有你自己蒙福之後.
你要成為別人的祝福.
好了.
姐妹.
繼續.
現在我照樣定義私恩與耶路撒冷.
和由大家.
你們不要懼怕.
你們所當行的是這樣.
各人與鄰舍說話誠實.
在城門口.
按治理判斷.
使人和睦.
誰都不可心裡謀害鄰舍.
也不可喜愛喜賈誓.
因為這些事.
都為我所行護.
這是耶和華說的.
神是首若私持愛的.
很信實的.
祂也希望以後復興了的.
神的國.
以色列國.
現在屬靈來說.
我們宗教會也是.
當神復興我們.
我們就不會再你騙我我騙你.
各人是與鄰舍說實話.
和按真理.
神的話去判斷.
和使人和睦.
不是以牙還牙.
而是以眼還眼.
而且祂說誰都不可心裡謀害鄰舍.
你心裡有沒有謀害人.
沒有.
誰會謀害人.
基督徒一定不會.
但是你有沒有.

$^{801}$什麼都.
咒在心裡.
不開心就咒別人.
別人在你的車前面咳嗽.
你會怎樣說.
不要說了.
很多人都這樣說.
不過不要說了.
趕著去哪裡.
這裡說.
不要心裡謀害.
不是謀害.
什麼都有你的好處.
什麼好處都給你拿了.
那些其實就是貪心自私.
甚至有時候說一些謊話.
來掩蓋自己.
好.
接著.
18 19 20.
請.
萬君之耶和華的華臨到我說萬君之耶和華如此說.
四月五月禁食的日子.
七月十月禁食的日子.
必變為由大家歡喜快樂的日子和歡樂的節日.
所以你們要喜愛誠實與公和平.
萬君之耶和華如此說.
將來必有列國的人和多城的居民來到.
這一段很有意思.
神本來是每一年定一日做禁食的.
贖罪日.
那個日子.
但猶太人呢.
他自己就定多幾日的.
因為他們很注重那種儀式.
那種的.
就是禮儀.
可以去.
就好像表現表表明自己是很敬虔的.
我是經常都禁食的.

$^{841}$記不記得呢.
有這個新約呢.
有兩個人在聖殿裡祈禱.
一個在左邊.
我在右邊.
一個呢.
是法利塞人.
是不是.
他在神面前.
他祈禱說.
神啊.
你看我多愛你.
我做十分一奉獻的.
我又禁食幾多次.
一個禮拜.
是不是.
我又幫很多什麼.
一路數.
他一個清單呢.
就是他做了很多好事.
連禁食.
都包括在內.
你知不知道教會如果.
落在一個形式化的裡面呢.
大家去參與.
無論是來也好.
去事奉也好.
變了沒了意思.
因為是一個.
一個循例.
或者是.
每年都是這樣做的.
做為了什麼呢.
傳福音啊.
我們請一些微信的人回來.
但是回來又怎樣呢.
一次過.
就不見了人了.
很高興.
說有很多節目.

$^{881}$但是我們是不是一路繼續跟進.
去關心他們呢.
是不是真的用心.
不只是用外表.
佈置得多漂亮.
多多東西吃.
那些都還不夠.
何況禁食是一種謙卑自己.
在神面前的悔改.
如果做得不對.
就求神憐憫.
這些時刻來的.
你知不知道為什麼要訂那麼多禁食日子呢.
因為他們覺得.
如果這樣做.
神一定憐憫我.
神一定會賜福賜恩給我.
他拿了這些善行.
用來跟神交換.
就是說.
你看我多好.
我多多去教會.
我多多去事奉.
我多多去守這些節日.
你應該要賜福給我.
你一定要賜福給我.
代代都要得福,蒙福的.
人就變成在這樣的情況下.
用善行,功德去博取神的祝福.
神的恩典.
神說我不是要這些.
祂說你們不用再禁食.
哀哭.
你們要將這些日子變為.
大家歡喜快樂的日子.
歡樂的節日.
但是裡面說要喜愛誠實和平.
好 下面.
一起讀吧.
男女都可以喜歡讀.

$^{921}$21,22.
「這城的居民必到那城說.
我們要快去懇求耶和華的恩.
尋求萬軍之耶和華.
我也要去.
必有列邦的人和強國的民.
來到耶路撒冷.
尋求萬軍之耶和華.
懇求耶和華的恩」.
這裡更加有圖畫.
很多在外面來的強國,列邦.
全部都來歸向耶和華.
而且在那裡尋求神.
懇求神的恩典.
他們都不敢去懇求神的福氣.
只說給我恩典.
我們真是很軟弱.
我們真是很多罪惡.
你來憐憫我們吧.
他們大家都去親近神.
好 最後一章.
大家讀完它.
「萬軍之耶和華如此說.
這日子必有十個人.
從列國居民.
一個猶太人.
我們要與你們同去.
因為我們聽見神與你們同在」.
這裡更加奇妙.
我們開「幸福小組」.
如果有五,六個同工.
即是那些福伴.
能夠請到五,六個Best.
我們已經很開心.
一對一可以關心他們.
但如果你雙倍.
譬如你有六,七個人.
你請到十四個人.
更加感恩.
因為你有很多工作.

$^{961}$關懷,報導,跟進.
但這裡說.
「到有一日.
十個人從列國諸族.
拉著一個猶太人的衣襟.
說我們要與你們同去.
因為我們聽見神與你們同在」.
你說多美的圖畫.
你不用到處去找人.
不用去派單張.
要怎樣做.
但那些人自動走到你那裡.
十個人拉著你一個.
即是那種傳訪的機會.
那種將人帶到神面前的機會.
是多到這麼厲害.
一個帶十個.
如果我們教會.
今年裡面.
如我們禦宇堂能夠有這個目標.
不要一人帶十個.
一人帶一個也好.
但如果一人帶十個.
那些人是自己來的.
不用你每個月請吃飯.
或者要送禮物.
他才肯來一會.
而是自己追著你.
「我要來 我要來 我要來」.
這是一個最美的應許.
弟兄姊妹.
你的生命有沒有神同在呢?.
人家看不看得到.
在你生命裡面.
是不同的.
不一樣的.
我要飛快一點.
然後就去到我的總結那裡.
不要緊 是有出來的.
不會錯過的.

$^{1001}$好 這裡.
有神同在 恩福滿門.
人民都要享受長壽平安.
和興旺順利.
這個都應我們節.
個個見到面說身體健康.
長命百歲.
現在你可以這樣說.
很多人百歲.
還要事事如意.
很興旺.
你知不知道你這幾天出街吃東西.
如果老闆在餐館.
他一定做一樣東西給你.
芝麻球.
他說原來它是甚麼「碌碌金銀滿屋」.
他們都說這個好意頭.
無論大的芝麻球還是小的芝麻球.
他說就是等你「金銀滿屋」.
其實金銀.
沒錯是代表我們很興旺 很富有.
現在搶金搶得很厲害.
所以做金的生意.
他們真的很多機會.
但是其實真正的興旺.
是靈性的興旺.
聖經說.
我願你們身心靈都興旺 興盛.
周左邊為.
變成蒙福.
並要成為別人的祝福.
你今年有沒有一個祈禱.
求神讓我不要去咒罵人.
讓我去祝福人.
愛你的仇敵.
為逼迫你的人祈禱.
這是耶穌基督教我們的.
禁食的日子.
要變為歡喜快樂的日子.
不要只是去注重一些外表的形式.

$^{1041}$或者外表讓人覺得你很敬虔.
不是用那些.
是用心靈裡面.
你真的用心去對人.
不是只是用外表的一些善行.
去表明你怎樣去愛神愛人.
各國各族都要來尋求神.
和懇求私隱.
這個冬季的奧運會.
原來有兩個機構.
它做了一件事.
一個是美男浸秦會和一個機構一起做.
就是所有去參加奧運會的人.
都預備了一些單張給他們.
也有些特別的聚會.
如果你們那個時間不用比賽.
你們可以去參加.
還有一些個人的工作.
另外一個叫Rock Music.
他們舉行一些演唱會,音樂會.
也是吸引人去聽一些基督教的音樂.
他們就用這個方法.
他們說我們去不了全世界.
但是全世界某些人就來了.
米蘭那裡.
我就用這個機會去傳福音.
還有其實我很沒有時間去看.
但是我前天有機會看到一篇報道.
是一個牧師叫做李迪全牧師.
李迪全牧師.
他是講什麼呢.
我想他是一個日本的牧師.
他說全個奧運會最感動的一場.
就是日本有兩個商人的滑冰拿了冠軍.
你知不知道是日本破紀錄的.
從來沒有拿過商人滑冰的冠軍.
他說這兩個人其實是年歲差很遠的.
一個二十幾歲一個三十幾歲.
而且兩個都受過傷.
而且那個男的簡直要放棄.

$^{1081}$但是這個女的.
跟他一起七年去練習.
結果在一個短曲的比賽的時候.
他們失誤了.
失誤了輸了.
想著不行了.
但是你知不知道那個女的.
跟那個男的說.
今天我們失誤了.
所以我們更加要切換心態.
只要相信我們能夠做到的.
就一定可以.
於是就彼此信任.
他說真正的信任是什麼呢.
關係是什麼呢.
我接得住你.
接著他們就練習練習.
不要放棄再來做.
去到最後那個長曲就贏了.
而且打破了紀錄.
贏了很多人.
原因是什麼.
他說健康的關係是允許對方示弱的.
不是說你經常要以強人的姿態出現.
你可以告訴別人我現在很軟弱.
特別是那個跟他同工一起去比賽的人.
結果兩個人就贏了.
今天其實誰是我們的同伴呢.
剛才我問過你們.
有誰可以陪你一輩子.
如果有媽媽就最好了.
一百歲不死也是有媽媽的.
最好是媽媽.
但是有時候媽媽也會離開你.
先生也會離開你.
子女也要飛走了.
因為他長大了.
他不能常常陪著你.
有些子女很乖的.
常常陪著父母.

$^{1121}$不過無論怎樣也好.
真正是我們托著我們.
接得起我們的.
就算我們輸了.
失誤了.
他仍然是不離不棄你的.
只有耶穌基督.
我們的天父.
我們的聖靈.
他一生都陪伴著我們.
所以今天陸上和我們唱的「神同在」.
就是提醒我們在人生什麼時刻.
你都需要依靠神.
知道有人托著你.
你就不用怕.
自己不要手軟腳軟.
就是盡了你的力.
祂會擁抱你.
祂會托著你.
這就是一個信任的關係.
對神是完全的信任.
夫婦更加需要彼此信任.
不要他健康的時候你很愛他.
到他有病的時候.
你就飛了不知去哪裡.
逃避現實.
求神幫我們.
我們一起祈禱.
天父我們都是蒙恩的罪人.
都是軟弱的人.
因為我們就算信了主.
罪得到赦免.
我們仍然是事事會犯錯.
得罪你得罪人.
但是我們感恩天父.
你是不離不棄我們的.
感恩主耶穌.
你甚至肯替我們.
擔當所有這些過錯.
這些懲罰.

$^{1161}$叫我們得到平安.
得到醫治.
天父求你就繼續給我們有信心.
信靠你.
完全信靠你.
無論是得時不得時.
總要專心.
無論是在強壯的時候.
軟弱的時候.
仍然是要依靠神.
保持耶穌基督的心態.
溫柔謙卑.
和忠心良善.
奉耶穌基督得勝的命求.
阿們.
\newpage



\section{創世記 22:12}
\label{sec:031701}
\textbf{聖經人聖經事:我知道}~~~~黃克勤牧師博士
\newline
\newline
連結: \href{http://www.church.com.hk/acms/content.asp?site=cdc\&op=show\&type=product\&code=031701\&layout=sermon}{\texttt{031701}} ~~~~ 語音日期: 2026-02-14
\newline
\newline
\hyperref[sec:031700]{\small{< < < PREV SERMON < < <}}
~
\hyperref[sec:index_chronic]{\small{[返順時目]}}
~
\hyperref[sec:031702]{\small{> > > NEXT SERMON > > >}}
\newline
\newline
創世記 22:12
\newline
\begin{longtable}{cl}
\hline
\hline
章節 & 經文 (和合本修訂版)\\
\hline
22:12 & \begin{tabularx}{0.7\textwidth}{X} 天使說：「不可在這孩子身上下手！一點也不可傷害他！現在我知道你是敬畏神的人了，因為你沒有把你的兒子，就是你的獨子，留下不給我。」 \end{tabularx} \\ \\
[1ex]
\hline
\hline
\end{longtable}
$^{1}$各位電影姐妹大家好.
又來讀聖經了.
生命樹宣教網絡為大家準備的聖經人聖經事.
今次繼續說的是亞伯拉罕.
亞伯拉罕實在有很多東西要說.
從亞伯拉罕的身上我們會學到很多人性,神性.
特別是要認識多些神的東西.
人性就是知道原來人是這樣的.
反映到我們自己.
另外神性就是認識多些上帝.
原來祂是這樣的.
其實聖經比較常說是默想的文學.
什麼叫默想文學呢?.
就是看完之後.
讀完之後.
我自己經過默想.
它成為我生命的一部份.
這樣叫默想文學.
今天的題目叫做「我知道」.
「我知道」是上帝說的.
「我知道」其實是祂知道萬人的心.
好,先看看.
創世記二十二章第十二節.
天使說你不可再借童子身上下手.
一點不可害他.
現在我知道你是敬畏神的了.
因為你沒有將你的兒子.
就是你獨生的兒子留下不給我.
一讀這段聖經你就會知道.
亞伯拉罕要獻以撒.
當亞伯拉罕要獻以撒的時候.
實在這件事是很難的.
難在.
第一,神從來不歡喜人用人做祭物.
神很不歡喜的,很憎恨的.
第二,以撒明明是祂一百歲的時候給我的.
獨生一個兒子特別給我.
為什麼要我獻祭給祂呢?.
沒有理由的.
第三,這個兒子已經十多歲了.

$^{41}$少年了,我也很疼愛他.
第四,我如何向太太交代呢?.
真的非常難做.
很多種.
見到親戚朋友怎辦呢?.
在僕人面前怎辦呢?.
但經過一段時間內心深層的掙扎.
他終於認清如果是上帝叫我做的我就做.
這件事一點也不容易.
如果要我獻祭給祂,我就信外物.
但如果要我獻祭給祂,我就信外物.
牛羊拿幾隻都可以,沒理由要我獻祭給祂.
亞伯拉罕神早起身.
可能他太太還沒動,一早就預備讓僕人出發.
前一段經文沒有提及他驚動太太撒拉.
可能驚動太太撒拉不能出門.
神早就出發上山,去到神指示的地方.
近距離已經有幾位僕人在等我.
我和同志去那邊拜一拜就回來了.
那些柴通常是僕人揹的.
他就叫了以撒揹的柴.
以撒揹就不用僕人揹嗎?.
當然了,如果帶上僕人上去有問題.
他也想到,如果帶上去要殺以撒的時候.
可能僕人會攔阻,主人,主人,不可以,不可以.
不可以,他要護著少主,一定不行.
所以就不讓僕人去,自己去.
帶上兒子,兒子走到半路.
以撒也說:爸爸,獻祭的火和柴.
我揹著,火你拿著,刀你拿著.
那隻羊羔在哪裡呢?.
以撒也有些疑問,知道有多少QK.
這麼有趣的,又不讓僕人來.
但是以撒真的厲害,以撒平生不是做了很多事情.
沒有什麼像亞伯拉罕或者雅各那樣.
但是他是一個非常信服的人.
好了,繼續走,看看怎樣.
在神的山那裡有準備的了,這次上去了.
要下手害同志.
這位同志以撒,亞伯拉罕一百幾歲,我推測一百一十幾歲.

$^{81}$要捉著一個十幾歲的小孩,綁著他,放在壇上.
如果他要掙扎,你根本捉不到他.
一腳踢開你,跑了去,你追不到他.
但是他甘心樂意被爸爸綁著.
要殺了,要獻祭了.
然後天使就說不要在這個同志身上下手.
不可以害他一點,一個都不可以.
現在我知道你是敬畏.
第一點要和大家說一下這個螢光的字.
第一點要和大家說的是敬畏.
敬畏是什麼意思呢?.
Yarei,Yarei就是害怕.
是afraid的意思.
如果我是fear God, afraid of him.
這個afraid是比較上是敬畏,很害怕.
fear可以譯作fear,可以譯作afraid.
意思是害怕,害怕神嗎?不是很親愛嗎?.
以前的人真的害怕神.
我們都應該害怕神.
不過敬畏的意思,中文真是好,中文字真靚.
敬畏是因敬而畏.
因為我很敬重他,所以我害怕了.
不想得罪他,因為我很尊敬他.
敬而畏,因敬而畏.
我知道你是敬畏神的.
因為你沒有將你的兒子留下.
獨生兒子留下.
這樣都給神了,真是厲害了.
這個時候,不要動,不要動.
不可以下刀.
我知道你是敬畏我的.
那麥巴喊就醒了.
不要下刀,獻祭又如何呢?.
他就自然想,還有什麼可以獻呢?.
Brahma是看,是注意.
周圍看,不是這樣看.
有很多人這樣看,有沒有看到什麼?.
沒有,其實很多東西看他都看不到.
他沒有注意,他留意看每個地方.
竟然看到在樹叢中.

$^{121}$這個小樹並不是一棵,應該是山上的樹叢.
如果你一眼看下去,就隨便看,看不到.
最重要這兩個字.
就看到一隻公羊.
公羊.
山羊當然是兩個混在稠密的小樹中.
Brahma就取了那隻公羊來獻祭.
代替他的兒子.
在這裡有一個很重要的神後意義.
就是有一隻公羊.
公羊.
是一隻雄性,不是雌性.
雄性的公羊.
代替神預備了一隻雄性的公羊.
在小樹上掛著.
在樹上掛著.
掛在木頭上,親身擔當我們的罪.
牠被掛在木頭上.
勾著在小樹中.
於是就拿下來代替祂的兒子.
這個取代.
Tagar.
Tagar就是取代.
代替祂的兒子.
聖經是說,耶穌基督是義的代替不移.
於是我們被稱為義人,被稱為稱義.
是因為耶穌基督這個完全的義的上帝.
代替我們的不移.
這個代替,替代,取代,代替等等.
是福音很重要的含義.
就是義的代替不移.
這是彼得告訴我們的.
有一個屬靈的救贖的含義.
就是以撒很像獨生兒子.
不過在這個表現上.
上帝不肯把他的獨生兒子給我們.
這個很難得,令人很感動.
亞伯拉罕以撒已經表達了這一點.
神不吝惜他的獨生兒子.
為我們世人寫了.

$^{161}$另外用一隻公羊掛在木上來替代他的兒子.
所以在亞伯拉罕以撒的事上.
有兩個相同的救贖的神學表達.
一個是亞伯拉罕不吝惜自己的獨生兒子給我們.
另一個是用一隻公羊掛在木上來替代.
以撒給獻上的死.
這樣就很清楚了.
很多時候我們遇到困難的時候.
怎麼辦?怎麼辦?.
你就會發現.
原來上帝早有預備了為我們開出路.
一味向下望,望遠一點.
上帝為我們預備了.
我在Google上找到一張圖.
其實一隻羊的山羊的角比較長.
牠去吃草.
但有時候不夠吃,牠會吃葉子.
牠的小樹就伸高吃葉子.
吃著吃著.
那隻角勾住在樹枝的中間.
牠就吊住了,拿不回來.
牠是不會拆的,手伸不到就不會拆.
牠吊住了,很多山羊都會這樣.
吊住了.
如果久一點沒有人救牠.
牠就死在那裡.
或者有獅子來,牠就吃了牠.
熊來就吃了牠,狼也吃了牠.
吊住了.
掛在木頭上.
這幅Google圖有什麼不對呢?.
就是羊會在這麼空的樹枝掛著.
因為牠在吃葉子,所以這些樹應該會有葉子.
剛才吃了很多葉子.
不過吃得很矮,跟長頸鹿不同.
牠在上面吃,羊在下面吃.
這就是背景.
也畫了一個壇.
大概是這樣的意思.
掛著.

$^{201}$阿布拉罕就吸納地方起名叫耶和華爾勒.
直到今天還說在耶和華的山上必有預備.
耶和華爾勒.
如果看英文聖經就有點難了.
因為英文聖經通常把「Y」譯成「J」.
就變成Jirah.
Jehovah Jirah.
如果是英文聖經是這樣的.
不對的音,因為讀了「J」.
Jehovah Jirah.
就是耶和華爾勒.
「Jirah」這個字是什麼呢?.
是漢見.
而不是預備.
聖經說:直到今天還說在耶和華的山上必有預備.
我們經常說:耶和華爾勒,神有預備.
其實在原文上是不對的.
因為「Jirah」是漢見.
耶和華看到.
原來要知道我們所有的動作.
所有的奉獻.
神都看到.
當祂看到我們的需要的時候.
祂就為我們做一些事.
所以耶和華爾勒也說「預備」也是對的.
祂看到然後就預備.
祂看到我們已經乾了,需要水.
祂看到我們已經很餓了.
祂看到.
祂看到我們已經很餓了,需要水.
我只想知道你到底是誰啊?.
\newpage



\section{約翰福音 6:13-25}
\label{sec:031702}
\textbf{榮耀君王:不要害怕}~~~~林鍵牧師
\newline
\newline
連結: \href{http://www.church.com.hk/acms/content.asp?site=cdc\&op=show\&type=product\&code=031702\&layout=sermon}{\texttt{031702}} ~~~~ 語音日期: 2026-02-21
\newline
\newline
\hyperref[sec:031701]{\small{< < < PREV SERMON < < <}}
~
\hyperref[sec:index_chronic]{\small{[返順時目]}}
~
\hyperref[sec:031703]{\small{> > > NEXT SERMON > > >}}
\newline
\newline
約翰福音 6:13-25
\newline
\begin{longtable}{cl}
\hline
\hline
章節 & 經文 (和合本修訂版)\\
\hline
6:13 & \begin{tabularx}{0.7\textwidth}{X} 他們就把那五個大麥餅的碎屑，就是大家吃剩的，收拾起來，裝滿了十二個籃子。 \end{tabularx} \\ \\ \relax
6:14 & \begin{tabularx}{0.7\textwidth}{X} 人們看見耶穌所行的神蹟，就說：「這真是那要到世上來的先知！」 \end{tabularx} \\ \\ \relax
6:15 & \begin{tabularx}{0.7\textwidth}{X} 耶穌知道他們要來強迫他作王，就獨自又退到山上去了。 \end{tabularx} \\ \\ \relax
6:16 & \begin{tabularx}{0.7\textwidth}{X} 到了晚上，他的門徒下到海邊， \end{tabularx} \\ \\ \relax
6:17 & \begin{tabularx}{0.7\textwidth}{X} 上了船，要過海往迦百農去。天已經黑了，耶穌還沒有來到他們那裡。 \end{tabularx} \\ \\ \relax
6:18 & \begin{tabularx}{0.7\textwidth}{X} 忽然狂風大作，海浪翻騰。 \end{tabularx} \\ \\ \relax
6:19 & \begin{tabularx}{0.7\textwidth}{X} 門徒搖櫓，約行了十里多，看見耶穌在海面上走，漸漸靠近了船，他們就害怕。 \end{tabularx} \\ \\ \relax
6:20 & \begin{tabularx}{0.7\textwidth}{X} 耶穌對他們說：「是我，不要怕！」 \end{tabularx} \\ \\ \relax
6:21 & \begin{tabularx}{0.7\textwidth}{X} 門徒就欣然接他上船，船立刻到了他們所要去的地方。 \end{tabularx} \\ \\ \relax
6:22 & \begin{tabularx}{0.7\textwidth}{X} 第二天，留在海的對岸的眾人發覺那裡原來只有一條小船，而且耶穌沒有同他的門徒上船，是門徒自己去的。 \end{tabularx} \\ \\ \relax
6:23 & \begin{tabularx}{0.7\textwidth}{X} 另外有幾條從提比哩亞來的小船，卻停靠在主祝謝後給他們吃餅的地方附近。 \end{tabularx} \\ \\ \relax
6:24 & \begin{tabularx}{0.7\textwidth}{X} 這時眾人見耶穌和門徒都不在那裡，就上了船，往迦百農去找耶穌。 \end{tabularx} \\ \\ \relax
6:25 & \begin{tabularx}{0.7\textwidth}{X} 他們在海的對岸找到他後，對他說：「拉比，你幾時到這裡來的？」 \end{tabularx} \\ \\
[1ex]
\hline
\hline
\end{longtable}
$^{1}$歡迎大家來到我們的週六崇拜.
上個星期我們喫了一餐.
非常豐盛飽滿的年夜飯.
就是耶穌以五餅二魚.
餵飽五千個壯男的天國的筵席.
我們研讀的是約翰福音.
約翰福音第六章.
約翰福音敘事的角度.
和婦女福音.
即是馬可,馬太和奴家.
是有明顯的差異的.
所以同一個故事.
婦女福音選擇了一個比較光明的一面.
耶穌看見他們就憐憫他們.
而約翰福音卻展示了一個比較黑暗的一面.
這是一個真光耶穌照明黑暗的故事.
整個故事的中心.
你一看就知道了.
就是耶穌基督的十字架.
耶穌分餅.
耶穌拿起餅來祝謝.
很明顯的是最後晚餐.
所以這不是一個普通的神蹟.
而是一個天國的記號.
是一箇舊俗的記號.
是一個十字架的記號.
但是可悲的事.
甚至可以說是很可惡的事.
就是五千個男人.
猶太的男人.
可能是沿路去耶路撒冷.
去敬拜耶和華的五千男丁.
二萬多的男女老幼.
全部人都在看神蹟.
沒有人願意聽耶穌的說話.
猶太人的群眾是這樣.
耶穌的門徒也是這樣.
你記得在故事裡面.
有腓力和安德烈.
兩個不知為什麼.

$^{41}$合力去對抗耶穌.
就是不聽耶穌說的話.
耶穌叫他去東他就要去西.
就像我上個禮拜所說.
五男一女.
很可能是耶穌眾多神蹟裡面.
最熱鬧.
但是也是最孤單的一個神蹟.
有趣的是.
緊接下來的這一段.
又是另外一個神蹟記錄.
好像上次.
耶穌應該已經很失望.
如果我是祂.
我就不會再做神蹟了.
誰知道祂又再做多另一個神蹟.
這次說的是海上行走的神蹟記錄.
我想先和大家看看.
《約翰福音》第六章.
整張聖經的架構.
然後我們再看那段聖經.
究竟是怎樣配合整個架構.
《約翰福音》第六章.
可以大概分為五段.
你自己去讀.
你都大概分到這五段.
第一段很明顯的.
是「五名異魚為伴五千人」.
這是上個星期我們所講的神蹟記號.
接著是這個星期我們要讀的故事.
就是「黑夜狂風」.
海上行走的神蹟啟示.
接著第三是「耶穌不再上船追趕」.
這是一般我們覺得是一個過渡的經文.
前面是神蹟.
後面是教導.
這是一個過渡的經文.
一般來講.
聖經學者覺得內容沒有太多.
回來講究竟這段是怎樣去理解.

$^{81}$接著到第四第五.
很明顯是耶穌的教導.
第一是「我是天降生命的糧」.
這是天國教導的第一部分.
接著再講.
「喫喝人籽有永生」.
是天國教導的第二部分.
在這個分段的架構裡面.
有沒有覺得「黑夜狂風」.
「海上行走」與其他部分是格格不入的.
其他部分全部都和五餅二魚.
餵飽五千人有密切的關係.
但是「黑夜狂風」.
「海上行走」.
好像是一個非常獨立的個體.
很難和《讓福音》第六章.
上下前後是連上一個關係.
於是我們要問一個很重要的問題.
到底我們怎麼去明白.
《讓福音》當中的第五件神蹟記號.
它寫了七個而已.
如果寫了七個.
在耶穌眾多的神蹟記錄裡面.
選七個.
每一個都應該有它的道理.
我們現在問.
因為第五件神蹟是非常不顯眼.
好像和整體的格格不入.
於是人們都在問.
究竟怎麼去明白這第五個記號.
從故事發展來說.
二和一其實是相連的故事.
你看回馬太福音的記載.
都是這樣寫的.
兩件事是時序上真的這樣發生的.
於是你這樣寫是當然有它的道理.
但是我想說的是.
在《約翰福音》裡面.
這兩件神蹟記號.
不單止是時序上的相連.

$^{121}$可以這麼說.
它其實是兩件神蹟合而為一.
是非常之緊緊相連的.
這個是我經過了很久多方思想.
最後得出來的一個結論.
用一句赫連的話來表達.
就是所謂的好事成雙.
我們這裡有兩個神蹟記號.
好像一對龍鳳胎一樣.
所以好事就是一個女再加一個子.
一個是女生一個是男生.
看上去知道是不一樣的.
但是其實.
它們是龍鳳胎.
其實它們是成雙成對.
是一對一定要一對這樣看你才明白的.
當然你單獨看也行.
但是你要加上它們.
哦!那我明白.
為什麼這兩個是前後一起變成一個龍鳳胎.
可能到這裡有人會問我.
林牧師.
你說它們是龍鳳胎這兩件神蹟記號.
它們的連結點.
它的連接在哪裡呢?.
我告訴你.
它的連接就在十二的號碼上.
十二的號碼.
有沒有發現.
第一個神蹟記號的結束的時候是十二個男子.
那你說十二個男子為什麼是十二個.
因為耶穌有十二個門徒.
沒錯.
我們假設每個男子就是一個門徒.
然後下面那個故事.
下面那個神蹟記號的開始.
是十二個門徒.
上面完結的時候是十二個男子.
下面開始的時候是十二個記號.
五餅二魚的結局.

$^{161}$是有十二個男子的餅碎.
其實你一聽到這樣你就知道了.
其實它不是結局來的.
它是預表了故事還沒完.
還在下集.
為什麼呢?.
因為有餅碎.
就像我上個禮拜所講.
每一次分餅其實都有剩下的餅碎.
於是五餅二魚其實是一個.
沒有完結不斷更新的神蹟記號.
你明白嗎?.
你喫完之後又有剩下.
有時候我們在家裡打掃.
年廿八洗垃圾.
洗完之後又有塵.
洗完也沒有塵.
有沒有想過.
如果耶穌給你一個使命.
就是要派光餅子.
派完之後每次都會有剩下的.
完了.
派完之後又有.
五餅二魚不斷重複更新.
看完一集又有下集.
永遠無窮無盡的.
正如我上個禮拜所講.
這是一個不斷被分享的神蹟.
直到今天仍然繼續被分享.
繼續延續直到有一天.
終於完結了.
沒有餅派了.
就是耶穌回來的那一天.
世界沒了.
最後審判的一天.
那就再也沒有餅派了.
所以五餅二魚的結尾是十二個男子.
而海上行走的起頭是十二個門徒.
有沒有留意到海上行走.
是《讓福音》中第一件.

$^{201}$也是唯一一件.
行在門徒身上的神蹟奇事.
有沒有想過.
約翰寫了七個神蹟記號.
其他的神蹟記號.
都是行在某一個人身上.
或者某一群人身上.
你記不記得加拿大婚宴.
水變酒.
誰得益呢.
可能是那對新人.
大臣兒子復活了.
騎腿三十八年的人.
可以起身行走進入安息.
五餅二魚餵飽千人.
唯有這個是行在十二門徒身上的.
其他接下來的.
還有是一個失明的.
與生俱來失明的人.
可能是漢奴光明.
接著拉薩努的復活.
其他所有都是行在某人身上.
或者是一個大群人身上.
當然十二門徒是有份的.
但這個是唯一一個.
幫他們度身訂造的神蹟記號.
這個記號的意義是什麼.
這個記號的意義其實就是.
要成為他們將來的鼓勵.
和給他們一個認可.
為什麼我這麼說.
是因為很快十二門徒就要被耶穌所差遣.
差遣做什麼.
就是將上面的神蹟記號的餅碎.
生命的糧.
去全遍以色列.
以色列十二個支派.
再通過以色列人傳到地極.
於是他們有這個使命的時候.
上帝給一個認許他們.

$^{241}$一個很特別的認許.
就是告訴他們.
他們需要耶穌的同在.
如果耶穌沒有說要他們同在.
去都不要去.
但是耶穌答應他們.
祂將與他們同在.
天天與他們同在.
直到世代的終結.
但是很重要的地方就是.
他們要承擔了這個使命纔可以.
所以當你去想大使命.
或者耶穌的使命.
或者你基督徒人生的時候.
第一樣東西要想的就是.
你有沒有承擔過使命.
如果你沒有承擔過使命.
你不要以為.
耶穌天天與我同在.
真是好了耶穌.
Hallelujah Amen.
這些全部都是假東西.
因為你從來沒有承擔過使命.
那個認許跟你無關的.
但是如果你真的承擔過使命.
真真實實不是假東西承擔的時候.
那麼耶穌天天與你同在.
直到世代的終結.
這個認許就真真實實.
落在你的身上了.
所以我已經給了謎底.
大家去思想.
我們進入今天的故事.
聖經告訴我們.
如果我們將第五個記號.
就是第二部分.
再加上第三部分.
第三部分是一個過渡的經文.
剛纔我們覺得沒什麼用的.
但聖經說如果你將第五件神蹟.

$^{281}$再加上過渡的經文連在一起.
就變成第六章十六至二十五節的時候.
你再砌出來.
就成為了一個前後呼應.
ABCDCBA的架構.
這樣都可以.
真的可以用.
連那些過渡經文都可以用來做這個架構.
這個架構是很棒的.
首先最外層的是AA.
你知道這個架構真是有用的.
兩個地方都是黃家白龍去.
一個是門徒上了船.
接著是另外眾人上了船.
BB.
耶穌還沒到天已經黑了.
下面那個B.
耶穌沒有去第二天早上.
一個是黑夜一個是早上.
兩個耶穌都沒有去.
CC.
耶穌海面走.
尖尖靠近船.
接著下面的C.
耶穌接了上船.
就立刻到了岸.
但所有的東西最重要的永遠是最中間的.
中間你見到就是他們害怕.
接著耶穌叫他們不要怕.
怎樣由害怕轉接到不要害怕.
原因是因為耶穌說了一句話.
耶穌說.
I am.
我是.
如果你知道I am.
你聽開我講到你知道.
是在希臘文裡面叫Echo Amy.
Echo Amy有兩個翻譯.
第一個是是我.
你讀中文的時候.

$^{321}$它不是翻譯做我是.
它是是我.
但同時之間也可以翻譯做我是.
兩個都對的.
耶穌說是我的意思是什麼呢.
身份表明.
就是黑漆漆.
看不到.
耶穌說是我.
但是.
耶穌不單止是在講是我.
耶穌在講是我是.
我是其實是代表了.
耶穌說上帝在舊約裡面的名字.
記得在猜及紀第三章的時候.
耶穌上帝.
耶穌說上帝去找摩西.
摩西就說.
別人問起我你叫什麼名字.
我怎麼說.
上帝要猜諧摩西.
上帝就跟他說.
我是自有永有的.
這個是你的中文翻譯.
但是如果看回英文很簡單.
I am who I am.
耶穌上帝的名字就是我是.
不過他中文就覺得我是好像很難明白.
於是就將它解釋為我是自有永有的上帝.
也就是說上帝的名字是我是.
耶穌說我是.
耶穌說你不要害怕.
因為我是.
因為我.
也就是說上帝必與你同在.
為什麼你不要懼怕.
是因為耶穌會與你同在.
因為耶穌說上帝.
耶穌就是耶穌說上帝的化身.
會與你同在.

$^{361}$這個就是海上行走天國記號的意思.
但是我們進入今天的經文.
《讓福音》第六章第十六節.
到了晚上《讓福音》第十六章.
叛徒下到海里.
上了船要過加瓦隆去.
如果你看馬太福音的平衡經文的時候.
馬太是這麼寫的.
他們為什麼會下船呢.
是因為耶穌隨即推他們.
推他們.
就是叛徒上船.
做完五餅二魚.
那群人全部get不到.
只喜歡喫餅.
於是耶穌說.
散了散了散了.
耶穌要走路了.
然後是誰推他們下船呢.
是耶穌.
耶穌安排他們上船.
先到對岸.
等叫眾人散去.
然後疏散了群眾.
然後耶穌自己獨自上了山.
去禱告.
耶穌推門徒下去.
重要的地方是整個行程是耶穌精心安排的.
不是門徒自己說不如我們用海路走路.
不如我們去日本.
不是的.
耶穌幫他精心安排.
你們上船吧.
我遲點來找你們.
我自己先上山.
然後他自己單獨上山.
所有的東西是耶穌精心安排.
然後發生的事情.
我希望你們記住所有的東西都在耶穌掌控當中.
不是天有不測之風雲.

$^{401}$人有殺蟾之禍福.
不是.
是所有東西是耶穌掌管的.
於是我們繼續看.
第十七節.
天已經黑了.
耶穌還沒有來到他們那裡.
不然狂風大作.
風浪翻騰.
第十九節.
門徒耀老若行了十多里.
如果折返英里.
大概是三四個英里.
在海的中間.
加利利海.
總之兩頭不到岸就是問題了.
看見耶穌在海面上走.
他兩頭不到岸.
但是耶穌來找他們.
漸漸靠近他們就害怕.
這是約翰的描述.
約翰的描述相對來說是比較簡潔.
其中一個原因可能是.
楊福音面世的時候.
這已經成為一個家傳戶曉的故事.
小孩子都會背.
我想帶你去看馬可福音的版本.
因為馬可福音的版本是更加仔細一點的.
馬可福音是這樣寫的.
第六章第四十七節.
是同一個故事.
他到了晚上.
船在海中.
你記得兩頭不到岸.
耶穌獨自在岸上.
他看見門徒因風不順.
搖簍很苦.
就是耶穌看著他們.
見到他們是很不妥的.
天快亮了的時候.

$^{441}$他在海面上走.
往他們那裡去.
想要超過他們.
這個是馬可給我們一些很細微的地方.
第一中文翻得不好.
他說天快亮了.
如果要看英文.
其實是The fourth watch.
就是四更時分.
什麼是四更時分?.
四更時分對於意識人來說.
就是晨早三點到六點的那段時間.
所以這個不是真的很晚.
或者很早.
三點到六點這個時間.
如果你重新去想這個故事是這樣的.
耶穌叫他們上船.
上船的時候可能已經是傍晚.
天快要黑了.
還沒黑完.
到了晚上.
船已經在海中.
耶穌獨自在岸上.
晚上是幾點我們不知道.
可能是黑了.
可能七八點的時候.
不是香港七八點.
大家都知道.
那時候七八點就黑了.
七八點是以色人覺得是第一更.
第一更是六點到九點.
因為他們每一天是傍晚開始.
所以六點到九點是第一更.
九點到十二點是第二更.
十二點到三點是第三更.
接著三點到六點是第四更.
所以在第一更的時候.
耶穌已經在岸上.
什麼叫在岸上呢.
在山上是可以的.

$^{481}$但耶穌也可能已經從山上走下來.
在岸上.
那時候是第一更.
接著他第幾更纔去找他們呢.
第四更.
中間可能牽涉了六個小時.
你說耶穌六個小時幹嘛呢.
看戲.
在岸上會不會耶穌在那裡等了很久.
見到他們不行了.
但又不立刻救他們.
於是總之就等他們辛苦一下.
好了算了.
等了六個小時見到他們都不行.
於是去幫忙了.
但幫忙的時候又很奇怪.
耶穌的方法是.
在水面上走.
接著又想要超越他們.
耶穌拯救門後的方法很特別.
有沒有想過.
其實耶穌不需要在水面上走.
你說那怎麼辦.
興功嘛.
有沒有看過戲.
一跳就跳了過去.
為什麼要在水面上走呢.
所以水面上走其實是多餘的.
是專門的.
接著水面上走的時候.
又要好像載意要超過他們.
是看不清楚還是什麼的.
原來所有這一切.
那些細微點都是與《舊約聖經》有關的.
耶穌是專門做一場戲.
要表明祂的身份.
和祂來到世界.
祂的使命.
所以我告訴你們.
海面上行走這個神蹟記號.

$^{521}$其實是要帶我們回到《舊約聖經》裡面的三本書卷.
四個經文.
四段經文.
三本書卷四段經文.
哪三本書卷呢.
第一《約伯記》.
第二是《出埃及記》.
裡面有兩段經文.
最後是《以蔡亞書》.
我想很快快和你們看看.
究竟這幾段經文.
是表明了耶穌是一個什麼身份和祂來的使命.
《約伯記》第九章第八節.
這個是講上帝不是講耶穌.
不過耶穌很明顯是在映射.
這個就是關於祂的經文.
祂獨自鋪張諸天.
步行在海浪之上.
如果耶穌輕功過去了.
就沒有了這個步行在海浪之上的應驗了.
祂行大事不可猜度.
行其事不可勝數.
接著第十一節.
看啊.
祂從我旁邊經過.
祂不是走過來直接來.
祂是從我旁邊經過.
我看不見祂在我旁邊走過.
我沒有察覺到.
我想你知道一件事.
原來在舊約聖經裡面.
上帝每次榮耀顯現的時候.
都是在你旁邊察覺而已.
如果上帝直接去見你的時候.
你死了.
於是每次都是察覺.
所以你一見到.
有人在你旁邊察覺.
其實這個代表了上帝榮耀的彰顯.
所以耶穌故意在那邊察覺.

$^{561}$告訴你知道我是誰.
我就是上帝.
耶穌是誰.
你看了這段經文.
耶穌是上帝.
耶穌是誰.
祂是鋪張珠天.
步行在海浪上.
沒有人可以征服海的.
尤其是那個世界.
現在還是沒有人征服海的.
誰可以走在水中.
就是那個創造宇宙萬物的創造者.
和祂也是掌權者.
祂不僅創造了.
祂也是掌權宇宙所有的萬物.
這是第一點.
接著我要和你去看的是《雙一級紀》.
《雙一級紀》有兩段聖經.
第一段我們先看的是《雙一級紀》第十四章.
這是很熟悉的經文.
你看到這段聖經和耶穌行在水中.
那樣行在水面上的神蹟.
有多少相似的地方.
《雙一級紀》第十四章第十三節.
《五窩西說》不要怕不要懼怕.
要站穩.
耶和華用強勁的東風.
又有風.
在海水再一夜間退去.
又是海水又是夜晚.
以色列人下到海中走在乾地上.
接著破曉的時分.
又要破曉的時分.
耶和華使埃及的軍兵混亂.
你看到這裡有很多平衡的圖案.
不要懼怕.
強勁的東風.
在海中漁走在乾地上可以走路.
接著在破曉的時分.

$^{601}$上帝拯救上帝審判.
我們問耶穌是誰.
耶穌是拯救以色列人出埃及的耶和華上帝.
第二.
耶穌是審判埃及偶像假神的耶和華上帝.
因為最後破曉時分的時候.
不僅是拯救的時候.
也是審判的時候.
我想你大概記住這些東西.
接著我們要看第二段的插一級記.
第二段插一級記是.
插一級記第三十三章和三十四章.
在三十三章裡面.
如果你記得金牛族的故事之後.
摩西很大膽.
摩西想更認識耶和華.
摩西說求你.
顯出你的榮耀給我看.
上帝我知道你很噁心.
但是我想見你一面.
我想見到你的榮耀.
但是上帝說.
你沒辦法見到我.
我又繼續活下去.
於是上帝想了個辦法.
遮住他.
然後說我會在你面前經過.
又是經過.
經過是上帝彰顯他榮耀的一個代號.
於是去到第三十四章裡面.
第六節.
耶和華在他面前經過.
然後宣告.
耶和華.
有憐憫有恩慈的上帝不輕易發怒.
且有豐盛的慈愛與信實.
我們問耶穌是誰.
耶穌是上帝榮耀的彰顯.
第一個很重要.
耶穌是榮耀的彰顯.

$^{641}$第二.
耶穌是上帝慈愛與信實的化身和成全.
這些你不記得不要緊.
我們以前全部都讀過.
如果你記得.
這就是讓福音的主旨.
在我們讀那個Prologue的序的時候.
已經是這麼說.
已經有讀過這段聖經的了.
讓福音第一章第十四節.
道成了肉身駐在我們當中.
充充滿滿的有慈愛與信實.
就是這段聖經的慈愛與信實.
我知道你那本中文不是這麼翻的.
那個翻得不好的.
充充滿滿有慈愛信實.
我們有看見過祂的榮耀的彰顯.
正是負獨一兒子的榮光.
時間有限.
我們要看最後第四段的聖經議策下書.
議策下書第五十一章.
又是和耶穌行在水面裡面.
有很多平衡的地方.
第五十一章第十節.
上帝使海和深淵的水乾沽.
在海的深處開爐.
使舊熟的民走過.
不是你嗎.
我唯有我是安慰你們的.
這個我唯有我是.
如果你去看希臘文的翻譯舊約.
舊約是希伯來文.
但是它有個希臘文的版本的舊約的時候.
是雙重的Echo Amy.
就是雙重的我是我是.
希臘文是寫我是我是.
我是安慰你的.
所以又有很多的平衡.
耶穌是誰.
耶穌是帶領以色列人被老迴歸.

$^{681}$為什麼我這麼說呢.
因為.
議策下書講的是被老迴歸.
但是被老迴歸的模式.
又好像《雙埃及記》一樣.
所以要在海里面開爐.
然後行在乾地.
歷史背景是被老迴歸.
要解決罪的問題.
要使罪得到赦免.
但是畫出來的圖畫.
好像是第二次的《雙埃及記》.
於是我們問.
議策下書第51章告訴我們.
耶穌是誰.
我想告訴你們.
議策下書第51章告訴我們.
耶穌其實是受苦的僕人.
為什麼我這麼說呢.
因為議策下書第51章.
剛好夾在兩首《牧人之歌》之間.
如果你以前有聽我講.
你知道議策下書.
由40章開始到53章.
都是有四首《牧人之歌》.
或者是四個部分的《牧人之歌》.
議策下書第42章.
議策下書第49章.
然後議策下書50章.
就是這裡的前科.
然後議策下書52至53章.
那個受苦的僕人.
我們魚羊走迷.
上帝將我們的罪扔在耶穌的身上.
議策下書是舊約一本很美麗的書.
我記得你以前.
如果你記得我有講過議策下書的時候.
我說議策下書是舊約的雙城記.
議策下書開始的時候是講.
你們住在一箇舊城的耶路撒冷.

$^{721}$這個耶路撒冷本來應該是聖城.
但變成了偶像之都.
議策下書的結束的時候.
我們去到一個新城的耶路撒冷.
在那個新城裡面成為上帝之城.
在那裡我們再不拜偶像.
我們是敬拜真神.
然後我們問.
我們怎麼移民從上移到下.
怎麼從舊城移到新城.
我們要通過要跟隨一個.
上帝揀選喜悅的僕人.
那個叫做受苦的僕人.
那個就是耶穌.
於是我們跟隨耶穌.
祂為我們受死.
祂為我們的罪而受死.
祂拯救我們使我們聚得成眠.
但同時現在你要跟隨祂.
你不可以說多謝耶穌.
然後我繼續留在舊城.
那你就死定了.
你一定要跟著祂移民.
從舊城的偶像移民到.
新城裡面的獨因真神.
這個就是議策下書51章裡面.
那個受苦的僕人的形象.
我想帶你回到《讓福音》.
不要緊一會兒我會幫你總結的.
《讓福音》第六章.
我們要回到耶穌再水上行的事情.
他們就害怕.
耶穌卻對他們說.
我是.
我是耶和華.
剛纔那些聖經全部都跟這個有關係的.
於是跟他們說.
因為我是上帝.
因為我與你同在.
所以你不要害怕.

$^{761}$於是門徒接耶穌上船.
有了耶穌船就立刻到了他們所要去到的地方.
我們問耶穌是誰.
我們說耶穌其實就是.
耶和華上帝的化身.
祂就是我是.
但耶穌同時也是.
耶和華上帝差遣到地上.
承傳救贖的上帝的遺旨.
所以所有的東西都堆在一句話.
我是那裡.
有了耶穌的同在.
這是應用.
有了耶和華的同在.
我們可以得到拯救.
我們可以得到平安.
於是我們不再需要去害怕.
這是一個應用.
我想用一個圖表去總結.
海上行走神蹟記號對我們的啟示.
我們剛纔看了四段舊約聖經.
告訴我們以下這幾樣.
第一 耶穌是誰.
耶穌是創造萬物.
掌管萬有的耶和華上帝.
這是很重要的.
第二 這是通過約伯記的.
接著 耶穌是審判埃及.
拯救上帝子民的耶和華上帝.
第三 耶穌是彰顯耶和華.
榮耀耶穌是慈愛信實.
耶和華上帝是慈愛信實的承傳.
第四 耶穌是受苦的僕人.
祂為我們結束了那個祕虜.
讓我們取得赦免.
讓我們可以迴歸到上帝的面前.
第五 回到《讓福音》第六章.
暫時不開古.
這四段經文一同啟示了.
耶穌的身份和使命.

$^{801}$我知道大家可能不是很習慣去想.
但我要你們記住.
所有故事背後.
是以色列人.
猶太人怎麼去想這些經文.
他們看到那些圖案的時候.
他們會想起舊約很多的東西.
而每一箇舊約的圖畫.
其實都是在告訴我們.
耶穌是誰.
我想問問你們.
不知道你們有沒有留意到.
第二 第三 第四.
最終都是帶我們去到十字架的面前.
我們由第四開始.
耶穌是誰.
這個很簡單 耶穌是受苦的僕人.
祂承傳了舊俗.
帶我們回到上帝面前.
所有的東西和十字架有關係.
第三 耶穌是誰.
耶穌是彰顯上帝的榮耀.
我們現在在讀《讓福音》你就知道.
十字架是上帝榮耀的彰顯.
因為去到《讓福音》的後半部.
所有東西都是榮耀榮耀榮耀.
耶穌的祖國就是上帝的榮耀.
那個榮耀是十字架的榮耀.
所以你會發現.
當我們在讀《猜測記》的時候.
講到上帝的榮耀的時候.
我們知道最終的榮耀在哪裡會彰顯.
在十字架上面彰顯.
十字架同時也是上帝慈愛與信實的承傳.
上帝在音樂裡面去愛以色列人.
上帝答應他們要為他們除罪.
如何除罪.
上帝最後有沒有守住那位弱友.
如何信實地幫他們除罪的答案是十字架.
所以你一想到榮耀.

$^{841}$想到慈愛和信實.
所有東西都帶我們去到十字架.
第二 耶穌是誰.
耶穌是逾越節的羔羊.
因為這裡講的是出埃及記.
要審配埃及.
要拯救子民.
如何拯救.
是通過逾越節的羔羊.
通過除罪上帝的羔羊.
我想你們是這樣升的.
4 3 2 1.
1好像和十字架沒有一個直接的關係.
但是1很重要.
1告訴我們承傳救贖的是誰.
是創造萬物.
掌管萬有.
沒有東西是他掌管以外的東西.
這個很重要.
為什麼.
因為1和5是有直接的關係.
5是什麼呢.
我們告訴你.
1 2 3 4.
全部都是扔下去.
在5那裡.
5 讓福音第六章很有關係.
5就是耶穌將要賦予.
門徒一個福音的使命.
而當他們出於福音的使命的時候.
他們會有很多懼怕的時刻.
因為他們會進入一個黑暗的世界.
在他們懼怕時刻.
你怎麼可以不要懼怕呢.
答案就是.
因為猜他們的耶穌答應與他們同在.
猜他們的耶穌是誰.
是創造萬物.
掌管萬有的耶和華上帝.
所有的一切.

$^{881}$和福音的使命有直接的關係.
猜淺門徒進入世界.
我們問耶穌猜淺我們進入這個世界.
這個世界是一個怎樣的世界.
答案是一個黑暗的世界.
耶穌給我們一個使命.
這個使命是一個怎樣的使命.
是天國的使命.
同時這是一個光照在黑暗當中的使命.
而當光照在黑暗的時候.
這個黑暗是不會接受這個光的.
這個黑暗是拒絕的.
這個黑暗會恨惡這個光的.
所以在我們現在讀的神蹟記號裡面.
那個夜晚是很黑的.
然後有狂風大作.
然後在海面上也是巨浪鋪天.
於是在這個漫長的黑夜裡面.
在一個狂風大作的世界裡面.
在這個海浪翻騰的世界當中.
門徒需要耶穌的同在.
門徒需要耶和華上帝的同在.
以致他們不再懼怕.
當我在這裡畫這幅圖畫的時候.
說為什麼耶穌要做這個神蹟.
這個神蹟代表了什麼.
是因為耶穌會猜探他們進入這個世界的時候.
令我想起一部我也挺喜歡的科幻電影的名字.
那個電影是《星空奇遇記》裡面的一集.
叫做《Into Darkness》.
進入黑暗.
那個使命是一個什麼使命.
就是猜測他們進入一個黑暗的世界裡面.
這個正是十二門徒要擔起的福音使命.
使命是什麼.
要他們向黑暗的世界裡面傳揚.
世人都犯了罪.
虧缺了上帝的榮耀.
要告訴他們上帝的憤怒.
要告訴他們永恆的審判.

$^{921}$要在所有這麼黑暗的東西里面.
再告訴他們十字架的大愛.
我們現在是很厲害的.
我們跳過了前面那三四樣東西.
我們不講人的罪.
我們不講上帝的憤怒.
我們不講永恆的審判.
就這樣跟你講大愛.
你這樣你沒有大愛.
聖經裡面不容許我們這麼講.
為什麼會有這個神蹟記號.
為什麼耶穌知道.
他們會在海上大風浪.
原因是因為他們必須要傳講人的罪孽.
必須要傳講上帝的憤怒.
必須要傳講永遠的審判.
纔在這個罪與審判的憤怒當中.
去高舉耶穌的十字架.
他們人拒絕的不是所謂上帝的愛.
人拒絕的是你要告訴我.
我犯了罪 虧缺了上帝的榮耀.
你要告訴我上帝的憤怒.
永恆的審判 我接受不了.
但是這些是我們必須要傳講的福音.
於是耶穌說我必如你同在.
因為你要進入一個愛黑暗不愛光的世界裡面.
去傳揚這個福音.
《讓福音》第六章第十三節.
他們就把餅的碎屑收拾起來.
裝滿了十二個籃子.
這個代表了什麼.
代表他們承擔將來要有的福音使命.
要拿著這個籃子.
我不知道要怎樣去拍這段戲.
那個籃子可能很大.
然後耶穌可能變個法術.
孫悟空將它縮小.
你放在袋子裡.
以後會變大.
因為它代表了你將來的福音使命.

$^{961}$那些餅碎要傳給以色列十二個支派.
那些餅碎要傳到地極.
接著耶穌.
我想帶你們看第十七章和第二十章.
第十七章耶穌就要上十字架.
他說什麼呢.
他說你怎樣猜我到世上.
他在禱告上帝.
我也照樣猜他們我的門徒.
進入世界.
你去到第十七章的祈禱就包含我們了.
因為耶穌不僅是為他的門徒祈禱.
也為將來會信他的門徒的門徒祈禱.
除非你不認你是耶穌的門徒.
除非你說我只是來喫餅而已.
你不要指望我傳餅.
如果這樣你和耶穌就沒有份了.
但是如果我們喫餅之餘又是傳餅的話.
這個東西和我們有份.
上帝差遣我們做同一樣的事情.
於是你去到第二十章的時候.
耶穌三次講平安.
很有趣的.
耶穌第二十章第二十一節.
願你們平安.
負怎樣差遣了我.
我也照樣差遣你.
那平安是和差遣連上直接的關係.
正如我所說.
在《仰慕音》第二十章裡面.
耶穌連續三次和他們講願你們平安.
平安代表了什麼.
代表了耶穌受死復活所引來的.
屬天平安.
這個很重要的.
因為他已經受死已經復活.
他才和你講願你們平安.
平安代表了罪人與上帝之間的關係.
因為十字架的緣故.
現在可以修復可以和好.

$^{1001}$這個不是說我的心裡面.
我知道你最近心情不太好.
願你平安.
不是完全沒有關係.
這裡的平安.
耶穌說願你們平安的時候.
是一個差遣的動作.
耶穌說差遣門徒的時候.
差遣門徒進入黑暗世界的時候.
對他們的祝福.
耶穌說你不會像一個孤兒子進入去.
我會賜你們平安.
所謂平安就是他自己的同在.
是他同在的認認.
我要你帶著我的平安進入這個世界.
我的平安就是那些餅碎.
但我的平安同時也是我答應你們的祝福.
我答應你們的應許.
你不會像孤兒子進入這個世界.
我將要和你們一起.
於是.
因著屬天永恆的平安.
耶穌為我們成就的十字架.
屬天永恆的平安.
因著上帝同在的應許.
門徒現在可以全心信靠耶穌.
不再懼怕.
你明白嗎?.
所以不再懼怕不是一個轟動的.
不要那麼緊張,沒事的.
明天考試沒問題的,不要害怕.
冷靜點.
你讀得很好了.
你準備得很好.
完全沒有關係.
所有的不再懼怕是因為.
第一,耶穌成全了救贖.
第二,耶穌答應他必與你同行.
所以你不要懼怕.
我想帶你去看努加福音.

$^{1041}$因為努加福音在這裡.
給了我們兩個很寶貴的不要懼怕.
雖然努加福音這兩個不要懼怕.
是出於努加福音的中間.
就是耶穌還沒上十字架.
但你知道所有耶穌的故事.
都要去到最後十字架再讀回頭.
於是這裡耶穌說的不要懼怕.
也都與十字架絕對有關係.
努加福音第十二章第四節.
那最多隻能殺人身體的.
不能再做什麼的.
不要怕他們.
我提醒你們.
該怕那殺了以後.
又有權並把人.
應在地獄裡的.
是的我告訴你們.
正要怕他.
不要懼怕.
當你面對死亡的時候.
我會與你同在.
不要懼怕.
因為我已經為你成全了救贖.
不要懼怕.
因為你已經有永恆的財寶等待你.
你要怕的是.
有權備將你應在地獄裡面.
那個審判的上帝.
我想當時的門徒是完全聽不明白.
聽不入耳耶穌的那句話.
正如五名義義的時候.
安德烈.
肥力是完全不明白耶穌是什麼水.
但是我們知道.
在耶穌受死復活.
昇天.
聖靈降臨之後.
十二門徒裡面除了猶太之外.
因為猶太是不信.

$^{1081}$然後自殺了.
其他所有的門徒全部都是得到重生.
都是明白耶穌是誰.
也領受了耶穌的呼召.
愛慕耶穌的話語.
於是這段說話.
就是其中一段.
他們會想起的說話.
耶穌說你不要懼怕.
那個可以殺人身體的東西.
這個很重要.
為什麼呢?.
因為.
你知不知道十二門徒的下場.
你有機會去看十二門徒.
每一個除了約翰之外.
就是九十歲幫我們寫約翰福音之外.
其他十一個都是死的.
安德烈被人X了.
吊十字架.
肥力也是被人釘十字架.
為什麼?.
因為起初他聽不進耶穌的說話.
聽不明白.
但是到他們真的重新領受福音使命的時候.
他們聽明白了.
他們領受了呼召.
他們知道.
耶穌要派他們去到一個黑暗的世界.
傳揚上帝的福音.
也知道耶穌必與他們同在.
於是他們甘願接受.
他們自己十字架上的榮耀.
對於我們坐在這裡的弟兄姊妹.
我就覺得為上帝為福音殉道.
獻上生命的機會其實不是很大.
大家明白嗎?.
我不是說我們有沒有勇敢的問題.
其實真的會為上帝奉上生命.
沒有多少人.

$^{1121}$但是他輪不到你和我的.
然而在同一張聖經裡面.
我們正在讀的是《路加福音》第十二章.
有另外一個不要懼怕.
那個不要懼怕就和我們度身定做.
你走都走不掉的.
那個和我們息息相關.
那個不要懼怕從某個角度來說.
是另一種為主,為福音.
獻上生命的不要懼怕.
但是同時之間其實是同一件事.
那個不要懼怕從某個角度來說.
更加困難.
更加考我們.
是否真的願意為主,為福音.
獻上我們的人生.
你沒得選擇的.
你不可以說我們要死.
你都可以有機會死.
但是那個機率不是很高.
但是這個.
你今天走出去.
你不用走出去.
你已經可以開始活了.
這個不要懼怕是這樣的.
你們只要求他的國.
這些東西.
就必加給你們.
你們這些小羊.
不要懼怕.
你看到嗎.
不要懼怕.
因為你們的父.
落以把國賜給你們.
繼續.
你們要變賣財產.
為自己預備永不壞的前廊.
和用不盡的財寶在天上.
那裡賊不能近.
蟲不能住.

$^{1161}$因為你們的財寶在那裡.
你們的心也在那裡.
我要你們留意.
我們怎麼開始這裡.
是不要懼怕.
從不要懼怕開始.
你懼怕是做不到的.
所以很多人都是怕死的.
怕死是怕窮的.
我這麼做我怎麼生活.
耶穌說不要懼怕.
不要懼怕.
你信不信得過耶穌.
耶穌說你不要懼怕.
因為你要求他的國.
你不要懼怕.
因為上帝落以把他的國賜給你們.
上帝將祂自己的獨生子賜給你.
難道祂不會將你其他的需要.
去充充滿滿的給你嗎.
你不要懼怕.
你要放膽的去變賣你的所有.
將你的人生投資在天上裡面.
你就有財寶在天上.
然後耶穌跟你說.
你的財寶在哪裡.
你的心也在哪裡.
你的財寶不在那裡.
你的心也不在那裡.
你的心不在那裡.
你的靈魂也不在那裡.
你什麼都不在那裡.
你和那裡沒有關係.
是很可怕的一件事.
這是耶穌對我們的呼召.
不要懼怕的呼召.
這幾段幾十經文.
耶穌對我們發出永恆天國的呼召.
令我想起我神學院的一個同學.
他的名字叫做月池.

$^{1201}$我想好幾年前.
我也有重複的和大家分享過.
他們的故事.
我想再講一講.
他們的不要懼怕.
或者他們怎麼覺得上帝呼召他們.
要將他們的人生投資在永恆.
2009年我的同學月池和她的先生Wayne.
還有兩個當時是四五歲的女兒.
她由美國加州搬去巴布亞新幾內亞的一個小島.
那個小島叫BM.
大家記得嗎.
是一個部落.
那個部落BM這群人.
全世界只有兩三千人.
只有兩三千人是講那個語言的.
他們去那裡要住在那裡.
要學他們的語言.
然後為他們創造文字.
原因是要幫他們翻譯聖經.
你想想全世界只有三千人是講他的話.
沒有文字.
但需要創造文字.
為什麼這麼大工程.
因為需要將聖經給他們.
同年去的 上岸的有另外兩個年輕的家庭.
一個美國去的家庭.
另外一個是德國去的家庭.
他們一起到達 一起生活 一起學習.
三年之後.
他們都可以用流利的部落語言和他們溝通.
然後又為他們創造了文字.
不可思議在三年裡面.
然後2012年他們用了三個月的時間.
將聖經的故事基督的福音和當地人分享.
三個月.
因為之前開播的時候不講耶穌的.
不停地講舊約的創造.
阿巴拉罕那些不停地講.
去到最後.

$^{1241}$耶穌纔出現.
然後耶穌才上十字.
耶穌才復活.
這群人全部沒聽過那個故事.
三個月裡面他們聽了那個福音.
因為福音有人信主.
然後很快BM的教會就成立了.
2016年他們安立了他們第一批的長老.
2016年算一算 七年而已.
教會一安立了長老.
宣教士就將教會交給他們.
獨立在BM的群體裡面.
宣教士退居幕後.
繼續要翻譯聖經.
繼續要做輔助和訓練的工作.
但是誰在帶領教會.
就是BM那些人.
在這個小島上.
我的朋友和他一家.
居住了九年多 差不多十年的時候.
然後他們就搬了去臺灣.
現在他們在臺中.
做全時間訓練新一代華人的宣教士.
這個是他們的故事.
當中有很多人都還記得的.
那個記憶.
這個是一個奇異何等奇妙的故事.
然而每次你聽他們分享他們的呼叫的時候.
他們很強調一件事.
就是他們和其他人.
其他基督徒是沒有什麼分別的.
他們強調他們從來沒有聽過上帝的聲音.
不是上帝幫他們B-I-E-M.
不是和那些小孩讀書的時候.
有個什麼叫B-I-E-M.
沒有這些東西 沒有聽到上帝的聲音.
沒有見到任何的異象.
他說他純屬是按照聖經的教導.
聖靈的感動.
我們去慢慢去做我們覺得應該做的事.

$^{1281}$他每次他們分享他們的使命.
他們的呼叫的時候.
他也都鼓勵每一個基督徒都要去尋找.
並且活出他們自己獨特的福音故事.
這個是他們的故事.
但這個你應該有你的故事.
所以我現在想和你分享一段三分半鐘的分享.
2020年6年前1月份.
Wayne就是那個先生在臺北的一個分享.
主題很有趣.
主題就是我們其實沒有那麼特別.
三分半鐘而已.
但是很精淨的.
我想你留意每一句話每一個字.
不過他的分享是普通話的.
有字幕的.
希望大家看得明白.
三分半鐘而已.
我想你留意Wayne如何去鼓勵我們.
也都同時挑戰著我們.
我也知道說在不同的場合.
有的時候大家來聽.
講到或者是分享.
甚至買我的書或者看我的影片.
是想要了解我們的故事.
那我很謝謝大家.
就是我不是一個不願意講故事的人.
但是到了最後我想要講.
而且容許我講.
開始去找你們自己的故事.
成為一個有故事的人.
在聖經裡面在新約裡面.
你們知不知道耶穌在四福音裡面.
馬太福音 馬可福音 路加福音 約翰福音.
只有耶穌一個教導被記錄在這四個福音書裡面.
而且是在不同的場合.
耶穌有很多的東西是被重複的.
但是耶穌只有一個教導是在四個不同場合.
被記錄在四個不同的福音書裡面.
那這個是他這四個福音書的記載有點不太樣.

$^{1321}$但是我就做這種耶穌說誰想要得著.
或者是救著自己的生命的話.
他就會喪失生命.
那誰喪失生命的話.
他就會得著生命.
這個是耶穌.
這是唯一在四個福音書上面.
耶穌在四個不同場合.
而且被記載下來的.
我希望你們不會認為說一個Tip Talk.
我們這邊被邀請來的講員.
是特別有故事的講員.
耶穌為什麼會說一個那麼.
要嘛很哲學要嘛很極端.
什麼是要我喪失生命我纔會得著生命.
我相信耶穌是說.
我們的生命應該要做一些跟永恆連結的事情.
你的生命要跟.
不是我的.
你的生命要做一些跟永恆有連結的事情.
那你這樣的決定.
甚至這個世界會把你當成說.
你在把你的生命丟掉.
你為什麼要帶你的孩子去.
但是我要誠實說.
其實今天在這邊等一下來分享的講員.
我們都不是特別的人.
我們都不應該是特別的人.
我做的事情不應該特別.
因為我們回到聖經裡面.
我們看到的就是這樣子.
我們不應該.
我們也不是一個特別的人.
我們有的時候要打醒自己.
我們不要活在一個平行世界裡面.
好像我們跟著這個世界一直走一直走一直走.
只要這個世界按贊.
那我們就是在過一個.
這個這個這個纔是一個很特別的生活.
誠實來說.

$^{1361}$我們這三個講員或許是今天晚上最正常的人.
就是如果如果你的生命只是跟著這個世界走.
跟著這個世界的事觀一直走一直走一直走.
對不起.
我怎麼看聖經你的生命是非常的特別.
而且不是那種好的特別.
不是那種好的特別.
那所以我短短的跟大家分享真的不應該只是一個一瞬間而已.
而是我們把生命砸下去.
之後所累積的.
那我們什麼時候分享.
什麼時候可以貼在臉書.
什麼時候可以放在IG上面.
那就是為什麼有天堂.
那就是為什麼有永遠.
那就是為什麼有永恆.
生命到永生.
我真的再想多一次.
2009年.
三個宣教士的家庭.
六個大人.
九個小孩.
不愛惜自己的生命.
恨護自己在世上的生命.
是吧.
他在做什麼.
他不愛惜自己的生命.
他恨護自己在世上.
美國德國的生命.
來到巴布亞新基內亞的一個小島裡面.
要去到一個BM的部落民族裡面.
去做什麼.
學習他們語言.
全世界只有三千人能說他們的話.
那些東西而已.
要為他們創造文字.
用不著那麼嚴重.
為了要翻譯福音.
就是福音和聖經給他們.
他們和耶穌一起死了.

$^{1401}$就結出許多的死得來.
十年之後.
BM部落的民族.
居然有他們自己的聖經.
有他們自己的教會.
有他們自己的長老.
可以獨立的繼續傳揚他們的福音.
這是一個不可思議的故事.
不可思議.
有兩重.
第一不可思議是宣教士的奉賞.
更加不可思議的.
就是BM的教會竟然可以在十年裡面成立.
和獨立.
為什麼我這麼說呢.
因為相比之下.
我回到香港已經超過十一年了.
我又不用學新的語言.
我又不用翻譯聖經.
是吧.
然而我經常在想.
如果我今天離開香港.
我感覺是多年的努力.
就是付諸流水化為烏有.
他們可以光榮引退.
十年而已.
做完十年了.
而我就徒勞無功.
我想了很久.
他們搬家.
我又搬家.
他們努力.
我也努力.
他們奉賞.
我也奉賞.
為什麼他們這麼成功.
我這麼失敗呢.
我不明白.
我也不服氣.
然後呢.

$^{1441}$我就開始計算.
我就終於明白了.
要計算一下人口比例.
他們一家.
大小加起來.
他們經常很強調.
小孩都是他們的宣教士.
十五個人.
去到一個群組裡面有三千人.
比例是0.5\%.
我們現在搬回來香港.
一百五十萬人.
我們需要多少宣教士.
三萬七千五百.
三萬七千五百個宣教士.
你想清楚一點.
我不是說三萬七千五百個信徒.
宣教士.
還要精力充沛.
才華洋溢.
訓練有素.
九十後纔行.
大家已經過氣了.
很多人都在過氣了.
九十後三十多歲.
全人獻上.
奉上自己最黃金的十年.
這就是在發生的.
為什麼要做到這件事.
從這個角度來說.
你是不是覺得我們人手不足呢.
是不是我們人手大大大大大大大大大不足呢.
還是我們低估了.
我們這個小島上的使命.
我們需要三萬七千五百名.
全心全意.
全人全身奉上的宣教士.
精力充沛.
才華洋溢.
訓練有素.

$^{1481}$獻上他們的黃金十年.
黃金二十年.
黃金三十年.
我們不是要在這個小島裡面.
如果你聚集的是一群.
為學業為工作.
努力奔波.
來到這個小島上去尋覓伴侶.
追求夢想.
享受人生.
等待退休的兼職基督徒.
是沒有用的.
你明不明白.
我們需要一群.
單求上帝的國.
上帝的義的.
全職基督徒.
全職的宣教士.
讓他們覺得.
他們今天在這個小島上的原因.
不是去賺錢.
不是去建立家庭事業.
而是要將福音.
擺在這個小島上.
建立一間合乎上帝的教會.
我們需要一群這樣的人.
奉上一切.
為建立一所.
真正貼近聖經.
全部人生.
活出福音.
Part A Part B.
所有的一切.
是要建立一個這樣的基督教會.
你想要一個怎樣的教會.
你就需要怎樣的信徒.
如果你的教會.
沒什麼所謂的.
隨便找些人進來就可以了.
但是不可以這樣.

$^{1521}$我們要的是一個.
Impossible dream.
不可能的夢想.
而這個不可能的夢想.
同時是我們唯一的夢想.
The only possible dream.
我再想說.
我們除了人數不足以外.
我們還有一個外在的環境.
因為我們的外在環境.
極其的惡劣.
BM小島上的外在環境.
就是一個火山.
我們的外在環境.
就是一個裝假成真光之城的巴別巨塔.
你說誰比較可怕.
我們這個比較可怕.
他們活在他們的世界.
天天都想著生與死之間.
因為有一個活火山在他們旁邊.
我們這裡死了都不知道什麼事.
但我們的使命不能改.
除非你要去那個小島.
你在這個小島上面.
這個就是我們的使命.
上帝給我們的小島就在這裡.
上帝差遣我們進入黑暗的世界.
我想播幾分鐘.
一個月前我在那篇異象講道里面.
是講過一點點我們的使命.
在這個很惡劣的狀況下.
上帝要我們建立一個.
我們奉命.
沒得選擇的.
上帝給我們的.
要建立一個我所謂異逆流的教會.
當罪人遇上耶穌.
不是上升而是下降.
當蒙恩罪人跟隨耶穌.
他不是上升而是下降.

$^{1561}$而是一個反巴別的人生方向.
有沒有想過.
我們信主的時候是不是這樣想的.
我跟了耶穌.
應該是有個反巴別的人生方向.
還是想著耶穌會幫我步步高昇.
保羅比德約翰的人生就是這樣.
遇上耶穌之後.
遂利掌撒該的人生也是這樣.
要放棄所有跟隨耶穌.
以色列的老師.
尼哥德姆的人生也是這樣.
你開始明白我們的impossible dream.
一個不可能的夢想.
講真在香港有多少教會.
會經常關心提點.
他們當中的專業人士.
醫生 律師.
問他們有沒有考慮過辭職.
有沒有考慮過辭職.
去做一些更有天國意義的東西.
教會可以跟你共金服的.
我們可以減薪的.
我們不需要那些貢獻的.
但我們的人生一定要被上帝所使用.
我們為什麼要做這個不可能的夢想.
因為the glory of God.
上帝的榮耀.
上帝值得有一間這樣的教會.
上帝值得的.
因為他奉上所有給我們.
我們的教會應該稱得起.
我們的上帝 我們的救主.
而當我們去榮耀上帝的時候.
信不信由你.
聖經不停地告訴我們.
我們的靈魂的救贖益處.
就是最maximum的時候.
當你榮耀上帝的時候.
你也得到最大的救贖的益處.

$^{1601}$但我們不去追求上帝的榮耀的時候.
我們好像以為我們得到了很多地上的益處.
其實所有的東西都是虛空的.
耶穌告訴我們.
地上人追求的東西.
這些都是外邦人所求的東西.
你們需要這一切東西.
你們的天父都知道.
你們要先求.
我可以解釋 改寫.
單求 只求上帝的國 上帝的義.
這些東西都要加給你們了.
上帝說你們需要的這一切東西.
你們的天父都知道.
你們要求他的國 求他的義.
這些東西都要加給你們.
又如保羅在《羅馬書》裡這樣說.
萬事都互相效力.
叫上帝的人得益處.
你相信嗎?.
耶穌值不值得我們去傳言的信靠?.
我要結束了.
結束的時候我想跟你分享多一個朋友的故事.
原來在他們將近十年的小島宣教當中.
發生了一個小小的插曲.
就是我神學院的同學語詞.
兩度發現有癌症.
第一次的發現是他們剛剛學完語言.
預備可以首次 第一次去宣講福音的時候.
結果他們花了差不多一年的時間在臺灣去醫治.
然後再回到小島去跟他們分享福音.
然後幾年之後第二次他又復發.
那個時候是他們十年事工.
剛剛差不多要完成小島教會已經成立了.
獨立了的時候.
我想分享這四分鐘Wayne的分享.
不過這次是英文的.
我就幫大家做了一些中文的翻譯.
不過辛苦大家要去讀一下字幕.
(英文翻譯).

$^{1641}$新年快樂!身體健康!.
健康最重要!No!.
耶穌最重要!永恆最重要!.
大家明不明白?.
耶穌說:耶穌進前來對他們說.
受死復活的耶穌進前來對他們說.
所以你們要去使萬民作我的門徒.
奉聖父,聖子,聖靈的名.
為他們施洗.
凡我所吩咐你們的都要教導他們.
我就天天與你們同在.
直到世間的末了.
我們去祈禱.
我們是天父.
我們多謝你.
多謝你在五餅二魚之後.
有這個小小隻有十二個人經歷的神蹟.
因為馬可,馬太和約翰.
都將這個小小的神蹟.
只有十二個人經歷的神蹟.
記錄在他們的福音書裡.
因為這個故事.
和五餅二魚剩下的十二個男子.
掛上了終極的關係.
讓我們明白.
這個是耶穌給我們今天的認許.
當然認許的背後是差遣.
如果我們不領受這個差遣.
我們不可能迎接這個認許.
但是若我們在你的差遣當中.
若我們要去找我們.
與永恆連結的故事的時候.
我們所面對的所有的黑暗難處.
風浪當中.
就有你同在的認許.
幫助我們.
因為我們珍惜你的同在.
我們需要你的同在.
而唯一可以得著你同在的認許.
就是參與天國的使命.

$^{1681}$幫助我們明白這一切.
幫助我們知道今生的人.
今生的所有我們所擁有的東西.
時間,身體什麼的.
都是為了永恆的投資.
幫助我們深信這件事.
唯有這樣我們纔可以活出一個.
不一樣在人們的眼中.
是恨惡自己生命的人生.
幫助我們去相信你.
跟隨你.
我們這樣祈禱.
奉主耶穌基督明求.
阿姆.
\newpage



\section{約伯記 39:19-25}
\label{sec:031703}
\textbf{一馬當先,誰與爭鋒!}~~~~吳宗文牧師
\newline
\newline
連結: \href{http://www.church.com.hk/acms/content.asp?site=cdc\&op=show\&type=product\&code=031703\&layout=sermon}{\texttt{031703}} ~~~~ 語音日期: 2026-02-22
\newline
\newline
\hyperref[sec:031702]{\small{< < < PREV SERMON < < <}}
~
\hyperref[sec:index_chronic]{\small{[返順時目]}}
~
\hyperref[sec:031704]{\small{> > > NEXT SERMON > > >}}
\newline
\newline
約伯記 39:19-25
\newline
\begin{longtable}{cl}
\hline
\hline
章節 & 經文 (和合本修訂版)\\
\hline
39:19 & \begin{tabularx}{0.7\textwidth}{X} 「馬的力量是你所賜的嗎？牠頸項上的鬃是你披上的嗎？ \end{tabularx} \\ \\ \relax
39:20 & \begin{tabularx}{0.7\textwidth}{X} 是你叫牠跳躍像蝗蟲嗎？牠噴氣之威嚴使人驚惶。 \end{tabularx} \\ \\ \relax
39:21 & \begin{tabularx}{0.7\textwidth}{X} 牠用蹄在谷中挖地，以能力歡躍；牠出去迎擊仇敵。 \end{tabularx} \\ \\ \relax
39:22 & \begin{tabularx}{0.7\textwidth}{X} 牠嘲笑懼怕，並不驚惶，也不因刀劍退卻。 \end{tabularx} \\ \\ \relax
39:23 & \begin{tabularx}{0.7\textwidth}{X} 箭袋在牠身上錚錚有聲，槍和短槍閃閃發亮。 \end{tabularx} \\ \\ \relax
39:24 & \begin{tabularx}{0.7\textwidth}{X} 牠震顫激動，將地吞下 ；一聽角聲就站不住。 \end{tabularx} \\ \\ \relax
39:25 & \begin{tabularx}{0.7\textwidth}{X} 每逢角聲一響，牠說：『啊哈！』牠從遠處聞到戰爭的氣息，聽見軍官如雷的吼聲和吶喊。 \end{tabularx} \\ \\
[1ex]
\hline
\hline
\end{longtable}
$^{1}$馬年的時候我就用今年大概《輝春》,《電視》,《報刊》,《雜誌》.
最常出現的三句馬的祝福語.
特別是我看鳳凰衛視新聞的時候.
剛好習主席也用了這三句.
其中一句就是英國首相來到北京訪問的時候.
習主席勉勉宣布要在特朗普的巨棒底下.
帶領歐盟一馬當先過來.
不過習主席講了前面的部分.
後面的部分我就加上去.
大家祝福人有押韻,有屬靈的意義.
不如大家讀一次好不好.
龍馬精神,事主殷勤,馬到功成,主因豐盛.
一馬當先,靠主超前.
這三句祝福大家在不同的領域.
無論你是政界,商界,在家庭或者教會事奉裡面.
這三句祝福語都能夠給你力量.
馬很早就跟人類發生了聯繫.
我們看到原始人動月的壁畫裡面.
很多時候就畫了馬在草原奔跑.
所以中國人發現馬不單止可以吃.
它還有更重要的用途.
所以馬稱為六畜獅手.
比牛,比豬的價值更重要.
馬除了可以騎,旅行,去遠方.
要追回失去的牛羊.
馬的速度能夠有這樣的用法之外.
它還有一樣很特別的就是它作為戰馬.
由早期我們看到漢武帝.
因為得到中亞的寒血寶馬.
就打到新疆那邊.
中亞將一直影響著中國的匈奴人趕到匈牙利.
將突厥人趕到土耳其.
所以中國在歷史裡面少了這些少數民族的參謀.
但是馬在戰爭的用途.
除了中國之外.
在成吉思汗因為有蒙古馬.
贏了歐洲的高頭大馬.
將蒙古帝國擴展到歐洲的邊緣.
後來還有土耳其的突厥人.
那個鄂圖曼帝國.

$^{41}$他們也是靠著阿拉伯馬征服這個範圍.
這個帝國有六七八年這麼久.
雖然第一次世界大戰有飛機大砲有火藥.
但是在那個時候很奇怪.
那些人仍然用馬兵來衝刺.
去衝擊對方的戰壕.
我看過一個統計.
第一次世界大戰死去的戰馬有九十萬匹這麼多.
很厲害.
當然人死去是更厲害.
今天的戰爭我相信不需要用馬了.
用什麼呢?用電子馬.
用那些機械馬就比較多了.
中國的成語和馬有關的.
有大概二百六十八個這麼多.
大部分都是正面的.
除了人家說你是害群之馬.
你不要以為他恭維你.
或者說你指鹿為馬.
不是說你很聰明這樣的意思.
大部分是正面的.
馬和人類的關係很密切.
以致在我們鯉魚裡面.
很多時候用的馬是不能騎也不能跑的.
東南亞有個國家叫大馬.
其實他不大.
台灣以前有個執政者叫小馬.
其實他也很高.
還有選舉又有黑馬又有上馬又有落馬又有去馬.
還有造馬.
很多的東西.
總之數之不盡.
聖經裡面的馬也有.
聖經裡面大概有一百八十五節經文是講到馬.
舊約比較多有一百六十五節.
其中比較多集中在列王記下.
因為裡面講到很多戰爭.
新約只有二十節.
但是很特別.
二十節裡面有大部分是集中在最老美國卷書.

$^{81}$啟示錄.
啟示錄裡面講到戰爭用到馬.
其中一節很短提到有二萬萬匹戰馬.
即是二億匹馬.
你想想那個戰爭是非常之龐大.
聖經也肯定嗎哪種快速矯健和力大.
不過聖經叫人不要依靠馬.
有人靠車有人靠馬.
但是我們要提耀的明智.
神不喜歡人的腿快馬的力大.
祂喜歡那些敬畏祂和愛慕祂的人.
所以聖經在這裡做一個對比.
世間的人靠車靠馬.
但是我們要靠耀.
明明沒有車沒有馬.
神就出人意外的保守和帶領.
讓那件工作能夠完成.
我們今天這段經文就在約伯記.
全本聖經裡面.
連貫講到馬最長的一段經文.
就在約伯記三十九章.
我們今天來思想這一段.
我們會問為什麼在一卷講苦難的書卷.
會特別提到馬呢.
我們看看約伯真的很慘.
慘過梁天來苦過金葉菊.
那些年輕人不知道誰是梁天來不是梁啟超.
不知道誰是金葉菊不是皇上王那個什麼.
你去google一下就知道清代十大奇冤案.
裡面有這些名字.
你想想我列那些.
有誰同一個時間經歷那麼多.
最慘連他身邊的配偶都不扶持他.
質疑他.
你去死吧.
你還不放棄神.
他不保護你不賜恩於你.
你還相信他做什麼.
三個來安慰他的朋友卻是去定罪.
最初約伯都很有信心.

$^{121}$我指我的救贖主活著.
賞賜的是耶和華.
收取的是耶和華.
直至他被他三個朋友不斷刺激底下.
他開始恣意為著自己的冤屈來辯護.
在幾輪的辯護當中.
神就介入了.
神在旋風當中介入.
即是說神以大能的方式介入.
神就沒有直接給那個答案.
神是讓他看見更高的一種奧秘.
苦難是人間的一種奧秘.
因為我們看不到靈界的一幕.
因為我們是看到這個約伯記看不到的.
開頭那幾章.
這個世間邪惡的力量來到上帝面前挑戰上帝.
叫上帝撤離他所愛的人的保護.
不施恩給他.
我們在世間的人.
為什麼無端端會發生什麼事呢?.
我都按照律法,按照聖經去做.
我們因為看不到有些背後.
直至到有一天.
我們面對面去見到上帝的時候.
原來如此.
原來有一樣東西是這樣發生.
約伯記有一次經文.
安慰那些不識時移去世的人.
譬如嬰兒要前.
壯年嬰年早逝.
留下妻兒.
聖經說.
要叫這些人早點離去.
是為免他將來面對更大的禍患.
人生裡面有很多東西.
我們受時間空間限制.
我看不到我們整個人成幅圖畫的結構是怎樣.
直至到我們在神面前來看的時候.
原來當時發生這件事.
是對我生命有一個很重要的意義.

$^{161}$有一本書叫做.
天堂遇上的五個人.
就是寫星期二.
相約莫理的作者另外一本書.
當然他是以一個故事形式來展述.
他就說到.
我們人死後在天堂會遇上五個我們很可能不認識的人.
原來他對我們生命很有影響力.
有個人會走來說.
你知不知道你三歲的時候.
在馬路玩球.
球走了出去.
你在行人道兩個車縫之間.
就要追個球.
我當時駕車.
見到有個小朋友出來.
我一扭tie .
我就撞在一棵大樹.
死去.
三歲的小朋友不知道發生了什麼事.
原來他的生命是因為一個人扭tie .
而他有生命的延續.
其實我們人生裡面很多的遭遇.
真是很奇特.
我記得很多年前張國榮演過一部電影獲獎.
叫做霸王別姬.
他和國內一個明星.
那個男明星.
都很英俊.
張國榮的反串.
是一個京劇的名單.
他有一天在冰天雪地裡面.
行街的時候見到有一個棄嬰.
就是有人將小朋友.
不知道基於什麼原因.
就扔在路上的井旁.
他慈心一動就收養他.
他的僕人就說了一句話.
他說每個人的生命都有它的定數.
當然或者有姻緣等等.

$^{201}$你不要隨意去介入.
介入呢.
你搞擾了這個人生命的歷程.
或者影響你.
電影的故事就講到在文化大革命.
鬥死張國榮的.
就是這個他在風雪中收養的女兒.
你不去信某種宗教的姻緣等等的東西.
就是你那時候那一刻的慈心.
當然無可指責.
有時候當你一個人介入你生命的時候.
是會影響你的生命.
好像一灘水扔了一塊石.
一個漣漪一樣相關的關係.
我們生命上帝.
有時候叫我們遇到這個.
發生這件事有它不同的意思.
英國有一個女的Mystic.
就是靈修學者.
她說我們的人生其實很公平.
好像在下大雨底下所有人.
不過每個人站在一個位置.
伸手出去接的雨點都不同.
你接到這裡.
她接到另一處.
你生命領受的東西就變做不同了.
苦難在我們生命裡面.
我們只能夠從一幅大的圖畫裡面.
才能夠看見.
神跡就是風雨雷電.
神跡就是潮汐漲退.
神跡就是河馬.
鱷魚.
鴕鳥.
野山羊.
甚至馬.
來質問約伯和他們三個朋友.
你知不知道為甚麼.
一本萬造 各從其類.
各有特性.

$^{241}$馬就是在這種背景底下.
來出現.
其實除了讓他感覺到啞口無言.
在爐灰當中懺悔.
而且講了一句.
約伯記其中一句經典的說話.
從前我風雲有你.
如今親自見到你.
約伯作為一個義人.
或者說一個很資深.
很愛主.
捐很多錢.
在教會所有服侍.
捐很多錢去機構神學院.
原來我們都是道聽途說.
直至到我們被擊碎.
我們的生命被打碎的時候.
直接一些東西.
我們才能夠以一個柔弱的心靈.
去面對這個造物主.
和去重新了解自己.
了解神.
了解這個宇宙和人生的使命.
所以在這裡.
約伯的驚訝和讚嘆.
讚嘆神那種創造奇功的偉大奧妙.
神也要他像一匹戰馬一樣.
振作起來.
奮勇起來.
面對靈界囚敵.
作屬靈的征戰.
每年講到這個時候.
通常有很多信徒會問三個問題.
為什麼這個時候講生肖.
或者說我們講生肖會不會迷信.
還有有沒有聖經根據.
每次我都.
待會你們出去.
我那本小書.
每年都要做一個解答.

$^{281}$第一我們要看見.
我們傳福音保羅說.
要用諸般智慧.
主耶穌說.
誰是那忠心有見識的僕人.
能夠按時分量.
按時分量這四個字.
我很喜歡一本.
英文聖經的譯譯本.
譯得非常之好.
他譯做什麼呢.
Supplying the food according to different seasons.
春天的時候.
我煲血濕粥.
夏天的時候.
西瓜果凍.
秋天的時候.
雪甲糕.
冬天的時候.
蛇羹火鍋.
夏天蛇羹火鍋流鼻血.
你不是supplying the food.
according to different seasons.
按著不同的季節.
不同的群體.
就好像保羅在雅典.
阿略巴古所講的訊息.
我們要好像保羅.
過哥林多前書.
不傳別的.
只傳基督的十家.
只是認識一半的保羅.
當你看《士道行傳》.
保羅講未識之實.
你就知道.
保羅是超過.
過哥林多前書.
他所講的眼界.
所以我們是要.
最多最普遍講一個話題.

$^{321}$我在那裡.
不用陰陽五行.
不用術數.
不用曆法.
是以聖經的經文信仰來講.
第二.
基督徒對文化有三種態度.
隨俗,離俗,逆俗.
隨俗就是.
這個世間做什麼.
我們跟隨吧.
無相打壓.
這個當然未必對.
一種離俗就是.
斷言分開.
你長洲搶包山.
我那裡教會.
又弄一個基督教包山.
鬥去頂爛屎.
真是布什有講.
你那些會景巡遊.
我拿著十字架巡遊.
離俗的人會不會端午節.
不看爬龍舟.
不吃粽.
中秋節不拿燈籠.
不吃月餅.
其實有很多東西.
我們人好像魚.
生活在漁缸.
文化就是漁缸.
你沒辦法脫離文化.
你現在在教會聚會就是文化.
你坐的椅子.
地毯燈光.
你穿的衣服.
就是文化.
你出去吃的東西.
又是文化.
文化裡面的元素.

$^{361}$你利用一些中性成為福音的載體.
新年我不派利是.
不過人家派給我.
我就收.
其實利是現在基督教化.
全部寫好.
我們樓下有一個單張.
很好的.
應該全部拿去.
就是新年快樂.
劫下的利是封給人.
人家以為軟的東西要收.
誰知打開一張福音單張.
平時你給單張.
人家不要.
新年不要拿利是封.
就大吉利是.
多多都要.
拿回去再看.
所以你想想.
我們用豬般的智慧.
我們有很多傳福音的機會.
我們移風易俗.
改變舊俗的文化.
第三.
你說有沒有聖經依據.
聖經基本上.
每一種生肖動物都有.
除了一兩種.
或者是翻譯錯誤的問題.
聖經有沒有老虎.
有人說有.
好像祖堂一樣.
第幾章第幾節.
我想你一定是看印度版的英文聖經.
其實老虎大概在伊朗東面的地方生活.
印度的.
不是.
伊朗的東面.
印度的西面伊朗波斯那邊.

$^{401}$只有獅子.
那裡是一條線.
所以聖經.
你最遠的地域.
去到波斯的地方.
它就沒有記載老虎.
除非是所羅門王運的那些.
亞馬騮.
孔雀那些都有提到.
象牙那些.
不是在中東的地方.
但是唯一一本英文聖經.
有提到老虎的是哪一本.
印度版宣教士譯的英文聖經.
它凡是獅子.
那些人就譯做Tiger.
第一真是很道地化.
因為印度人.
如果你說以前來說.
他們認識老虎多過獅子.
但是原文來說.
那些人沒有一個字叫做老虎.
大部分都有聖經的根據.
特別我們剛才提到.
聖經裡面有馬.
總共有185節.
我們剛才的題目.
那個副題叫做.
成功人士的生命質素.
成功英文是success.
然後你就名成利就.
擁有權力與財富.
你擁有很多人脈關係.
一動的時候.
整個世界的政局都受影響.
世界的成功就是success.
但是我們基督徒就說功成.
Mission accomplished.
就是說交給你的職責任務.
你完成了.

$^{441}$中國人有句話很復古的.
很合乎聖經的.
就是功成身退.
你做完那樣工作.
你就退下來.
叫人只見到那個工作成果.
見到上帝沒有人見到你.
你是透明的.
中國也有好幾句話.
樣樣東西不見你在那裡做.
但是有你份在那裡參與.
功成Mission accomplished.
無論你做一個政界的領袖.
無論你是商界.
無論你在教會事奉.
或者家庭一個主婦.
你都可以有成就.
不單止只是你丈夫對你的供應.
給你一張信用卡.
無限地用.
而是你作為一個主婦的責任.
你成就了做物主給你的責任.
所以我們基督徒說功成.
世間說成功.
功成.
這個碼這七樣的素質.
我覺得是任何要mission accomplished的人.
去生命,屬靈,道德上.
都應該有這些屬靈的質素.
這七樣的質素.
第一,威嚴莊重.
作為一個要完成工作.
要事奉神的人.
是要怎樣呢?.
是有一個表裡一致的形象.
君子不重則不威.
孔子說.
君子就是一個領袖.
如果你不能莊重.
你沒有權威,沒有威嚴.

$^{481}$對你的下屬,同僚.
你沒有辦法去帶領.
在政界.
在宗教界.
我們看見.
當一個很出名,很出名的人.
突然捲入醜聞裡.
或者他生命有些瑕疵的時候.
被人揭發就要倒台.
國內某個宗教界.
一個本來應該獨身的人.
後來被人發現.
他三妻四妾.
兒女十幾二十個.
不單止他下台.
而令到國家的宗教政策.
影響基督教,其他宗教.
都是被收緊和檢點.
所以你看看那個.
一塊石頭扔進那個祠堂裡.
所帶來的震撼是多大呢?.
說回我們基督教.
都是很慚愧.
至少我知道有兩個寫書寫很多.
影響力很深遠的人.
他的書翻譯成很多的人.
宣佈說他不再寫書.
他不再講道.
因為他生命有瑕疵.
所以一個人.
你不要以為只是你的才幹得.
而是你的表裡一字.
生命的見證.
是很重要.
第二方面他充滿自信.
在谷中拋地.
一隻馬展示他那種力量.
十二年前有一部電影.
叫做雷霆戰區.
英文名叫做War Horse.

$^{521}$一個監製和導演.
是一個猶太裔.
很著名的電影人.
叫做Edwin Shepard.
因為他是一個猶太人的緣故.
我深信這部電影的背後.
一定去熟知約伯記這段有關戰馬的經文.
你看過那齣電影.
你知道那匹馬最初沒有白落.
遇不上一個識用他的人.
去泥地耕田.
被人用鞭打.
在泥濘地拖大炮上山.
避開肉裂.
直至有人賞識.
這匹馬的筋骨.
是屬於往前衝的戰馬.
他是第一次世界大戰.
九十萬匹馬中其中一匹死去.
你看最後一幕.
劇透了.
你看過也好沒看過也好.
真的很震撼.
他一路拖著鐵絲網.
鐵絲網又拖著木欄.
向著敵人振來衝去.
不單止他自己軍隊的人很震撼.
騎在上空的騎兵已經被人射了下來.
敵方也很震撼.
有匹瘋馬拉著一大串鐵絲網.
木向他衝過來.
那個場面很震撼.
這匹馬如果是一個人物.
他在那一刻攻城.
他在那一刻成功.
是這匹馬生命最榮耀最高峰.
就是那一刻.
是展現出來.
其實我們事奉神的人也是.
這種自信不是自私.

$^{561}$不是來自自我中心.
而是來自神的自我自信.
等待一個機會去發揮自己.
他又敢於承擔.
普通沒有訓練的馬.
見到刀刀劍劍.
又要騎上身.
很害怕.
但這匹馬是怎樣呢.
去迎接佩戴兵器的人.
承擔也是今天事奉祖師很重要.
最近一齣國劇很流行.
叫做《長安的勵志》.
提醒我們上班要準時.
不要遲到.
那個人就是遲到.
別人的同事將一件死差事扼了去.
每個人都笑.
你最適合做的.
又升官又發財.
其實死差事一定不能夠完成.
而且是代罪羔羊.
皇帝所愛的貴妃說.
很想吃貴妃笑那種荔枝.
你想想當時沒有高鐵.
又沒有空運.
又沒有高速公路.
要在廣東將剛剛熟或未熟的荔枝.
運上西安昔日的長安.
是根本不可能的事.
他就想用辦法.
頂硬上.
頂硬上這句話是很出名的.
是明朝一個名將袁崇煥說的.
好像他是東莞人.
他說這句話叫頂硬上.
即是說用肩膀承擔那件事去工作.
今天我們謝謝邦.
如果教會給一樣.
不要說死差事.

$^{601}$爛差事.
你敢不敢去承擔.
保證完成任務.
一定不成功就成仁.
這個承擔感覺是很好的.
曹操寫了一篇短文.
有四句話是真的很震動人心.
老騎伏匿志在千里烈士無憐壯心不移.
一隻老的馬在馬房等候人道毀滅.
但他的視野他的眼界是空懷千里.
一個曾經好做過豐功偉績的烈士.
壯心不移有個brave heart.
心仍然在跳動.
等待有無好似參宣最後一個時機.
臨了結之前抱住兩根碎柱.
在那刻參宣所殺的敵人.
比他一生所殺的更多和更厲害.
第四就是膽識過人.
他癡笑可怕的事.
並不驚惶也不因刀劍退回.
中國有句話雖千萬人吾往矣.
如果你知道這件是對的.
是神命定要做的事.
無論前面有多少反對.
耶利米神要他講一些不合時宜的說話.
神說你一定是由上到下.
君王到百姓都反對不同意你的.
但是我要使你成為堅城銅牆鐵柱.
耶利米冒險犯大能去講這些說話.
馬丁路德是一個小小教會的神父.
但是當他讀到聖經.
他明白上帝的真理.
他抗拒當時教會.
面對幾百個紅衣主教通通是學問淵博的人.
甚至是他的老師.
他的傳記很出名的一個題目.
就是他所講的說話.
Here I stand.
就是我站在這個立場.
除非你能夠以聖經駁斥我.

$^{641}$否則的話我不會改變我的信仰和微言.
今日如果在時代的氛圍底下.
時代的事包括教會裡面.
基督教的刊物,神學院裡面.
都未必和你的看法相同.
但是你在神面前領受是這樣.
你敢不敢不要說挑戰.
仍然建立在你原本的立場上.
第五,作好裝備.
他的箭袋有弓箭,長槍,短槍.
中國人有句說話.
十年磨劍,養兵千日,勇於一朝.
因此未必你一生的東西都有用.
但你不斷去升級自己.
到一個時候機會來的時候.
你不會好像十個童女那樣.
五個沒有準備好,燈沒有油.
只有那五個長期準備好的人.
能夠去迎接這個新郎.
第六,以衷致業.
他一聽覺聲就不耐站立.
事奉的起步點就是羨慕聖公.
你不羨慕,不愛神的工作.
你真的要停一停,想一想,去祈禱.
不要隨便去做.
因為你隨便去做就是獻煩火.
舊約裡面的煩火就是unauthorized sacrifice.
一個不跟律法,不跟憲制的標準.
隨意去憲制就擊殺在神的聖所裡面.
所以我們事奉神.
一定要羨慕神的工作,好好去裝備.
舊約二乘俠書第六章的撒拉弗三對翅膀非常之好.
一對遮臉.
叫人只看見他事奉的神,不看見他自己.
一對遮腳,保持生命的潔淨.
一對不斷拍翼,平衡自己.
不斷的高升.
我記得有一次在師班的訓練.
我說你們要學一個撒拉弗.
用三對翅膀去表達你們師班的詩歌敬拜.

$^{681}$不斷提升.
由不懂得看譜到懂得看譜.
不懂得一邊看著譜一邊唱字.
到有些合唱團受過訓練.
一張譜拿上來就唱了.
不是在譯Dolemi Fasso.
然後再看看這個音是怎樣唱.
有些真是習慣師班的人.
類似Ian Lok這些香港聲樂團.
拿起譜拉丁文就唱了.
為甚麼?他一開始就懂.
他唱了幾十年.
一邊累積,一邊keep up,一邊進修自己.
最後就是聽命順從.
聽到國聲.
他就有機會向前衝刺了.
今日教會和作信徒.
最起步最起步的一個功課就是信服.
我覺得這個很重要.
無論滿分,無論你基督徒生命的開始.
你要信服聖靈對你心的感動.
你要信服聖經話語對你的引導.
你要聽神的聲音.
你要弟兄姊妹之間相處.
彼此信服.
不要像智慧箭豬一樣.
唯你獨尊.
刺到四處跟你一起工作的人.
彼此謙卑自己.
而且就是信服於屬靈長者對你的勸導.
你起碼見成長的時候.
有很多的功經.
這些功經來自不想神的工作完成.
你們那些是豆腐爭工程.
一隻狐狸跳上去,你的城牆就垮了.
你起碼不要聽別人的說話.
你要攻城,你要成功.
一定要專心刺字.
一手拿工具,一手拿兵器.
丑的未來,各人對著自己的房屋建築.

$^{721}$丑的來,當國聲一響.
你就去那個國聲的地方.
會聚為神的家來爭錢.
貢獻這個箭馬有七種的質素,七德.
第一莊重,我們要事奉神.
必須表裡一致.
有一個形象上的見證.
因為我們是基督的戰訓.
我們要行事為人.
與所蒙的恩相稱.
第二要自信.
不要軟弱懼怕.
約書亞記開頭第一章已經說了很多的剛強壯膽.
因為約書亞要承接著摩西的工作.
他覺得他一比底下.
他覺得他自己比不上,很害怕.
第三要承擔,肩膀硬正.
有些教會的主一學是自己編寫的.
禮拜天要教主一學.
禮拜六晚那個人還沒寫好.
那些主一學的老師不知道怎麼搞.
或者有些教父的工作.
幾個月後問一下進度如何.
你又告訴我要做這件事嗎.
肩膀擔不住.
就算你不能夠完成.
有困難你要按時將進度報告.
教會是不是一定比外面的公司要求低一點呢.
其實不是.
外面的公司是賺一份工資.
討世間的老闆的喜悅.
在教會未必有人工.
你是討這個眾物主的喜悅.
是更加重要.
勇敢了.
這也是一個很重要的質素.
不畏艱難.
很好的裝備,有事奉的熱誠.
和能夠聽聖靈的聲音.
所以如果我們在任何的領域.

$^{761}$要成為一個功成的人.
不是成功,而是完成任務的話.
我們需要有這七樣質素.
我們怎麼去回應呢.
第一,感恩領受,善用所得.
各樣美善的恩賜都是從眾光之父而來.
我們將這些東西來善用服侍.
第二,不要驕傲了.
提放致詞.
詩篇有好幾段經文.
都將馬和倚靠神來比較.
國人不要看自己過於所當看的.
要看自己憑信心看得合乎眾道.
第三,要全然獻上.
我們就像馬的靈.
當有歸耶和華為聖那一天.
本來這幾個字只在大祭司的冠冕上有.
但是到那一天.
馬這種庸俗不及牛羊的動物.
都有歸耶和華為聖這個名牌的話.
就是說所有的人都可以分別為聖.
你都服侍神.
你想想多麼榮耀.
多麼特別的一個機會.
所以聖經說.
我們將身體獻上.
當作活祭.
我們全然的.
不是獻煩火.
當作活祭.
給神使用.
乃是理所當然的.
我剛才解了副題.
現在解正題.
十二年前我讀了一篇文章.
它的結論是.
此馬一出,誰與爭鋒.
它說的那匹馬是奧巴馬.
十二年前他是美國總統.
你看特朗普那麼霸道.

$^{801}$就知道美國的總統.
不單止在自己的地域裡管.
他管著全世界.
全世界的警察一樣.
搞亂世界一樣.
這個人一定不是基督徒.
如果他是基督徒.
讀過啟示錄十九章十一節.
你知道真正的此馬一出是什麼.
我看見天開了.
有一匹白馬.
騎在馬上的稱為誠信.
林牧師的名相同.
真實.
它審判和徵戰都憑著公義.
這匹馬是古今中外.
人類最榮耀的一匹馬.
因為上高是萬王之王,萬主之主.
如果我能夠做這匹馬就好了.
就算不是做這匹馬.
我們做一個.
將耶穌基督帶入耶路撒冷.
去成就救恩最後一步.
那個小奴區.
都是一樣.
因為斑馬,奴.
和這個奴區.
都是馬的科屬裡面.
馬有很多種不同.
在神在旋風當中.
向約伯質問的時候.
他講了幾句話.
約伯.
你則以聽者話.
要站立思想成妙作為.
你若聽從他.
道日亨通歷年福樂.
會不會神今年藉著這段經文.
提醒我們每一個基督徒.
我們雖然處於約伯生活的逆境困難.

$^{841}$包括去年黃福縣那些人.
如今的福利制度.
政府都願意走另一條路.
你想想如果好像經歷約伯.
又沒有保險.
又沒有政府為他走另一條路.
那不就更艱難.
但是你要聽這句話.
你仍然專心去仰望神.
你必道日亨通歷年福樂.
英文聖經這一段.
是譯得很好的.
就是.
Every day you are prosperous.
每一天你都是興旺的.
And happy for the rest of your life.
在你餘生的歲月裡.
你都幸福快樂.
嘩 多大的祝福.
道日亨通.
就是諸事順利.
歷年福樂.
不是今年 每年都福樂.
所以我希望大家用約伯這段經文.
作為我們的提醒.
最後.
由大家一起讀 好不好.
此馬奔騰.
不甘後人.
為主是瞻.
一馬當先.
老氣復溺.
志在千里.
道日亨通.
歷年福樂.
希望這幾句成為我們的.
今年座右銘.
願主賜福各位.
\newpage



\section{馬太福音 25:1-13}
\label{sec:031704}
\textbf{我的燈需要油}~~~~劉少康牧師
\newline
\newline
連結: \href{http://www.church.com.hk/acms/content.asp?site=cdc\&op=show\&type=product\&code=031704\&layout=sermon}{\texttt{031704}} ~~~~ 語音日期: 2026-02-22
\newline
\newline
\hyperref[sec:031703]{\small{< < < PREV SERMON < < <}}
~
\hyperref[sec:index_chronic]{\small{[返順時目]}}
~
\hyperref[sec:031705]{\small{> > > NEXT SERMON > > >}}
\newline
\newline
馬太福音 25:1-13
\newline
\begin{longtable}{cl}
\hline
\hline
章節 & 經文 (和合本修訂版)\\
\hline
25:1 & \begin{tabularx}{0.7\textwidth}{X} 「那時，天國好比十個童女拿著燈出去迎接新郎。 \end{tabularx} \\ \\ \relax
25:2 & \begin{tabularx}{0.7\textwidth}{X} 其中有五個是愚拙的，五個是聰明的。 \end{tabularx} \\ \\ \relax
25:3 & \begin{tabularx}{0.7\textwidth}{X} 愚拙的拿著燈，卻沒有帶油； \end{tabularx} \\ \\ \relax
25:4 & \begin{tabularx}{0.7\textwidth}{X} 聰明的拿著燈，又盛了油在器皿裡。 \end{tabularx} \\ \\ \relax
25:5 & \begin{tabularx}{0.7\textwidth}{X} 新郎遲延的時候，她們都打盹，睡著了。 \end{tabularx} \\ \\ \relax
25:6 & \begin{tabularx}{0.7\textwidth}{X} 半夜有人喊：『看，新郎來了，你們出來迎接他。』 \end{tabularx} \\ \\ \relax
25:7 & \begin{tabularx}{0.7\textwidth}{X} 那些童女就都起來挑亮她們的燈。 \end{tabularx} \\ \\ \relax
25:8 & \begin{tabularx}{0.7\textwidth}{X} 愚拙的對聰明的說：『請分點油給我們，因為我們的燈要滅了。』 \end{tabularx} \\ \\ \relax
25:9 & \begin{tabularx}{0.7\textwidth}{X} 聰明的回答：『恐怕不夠你我用的；你們還是自己到賣油的那裡去買吧。』 \end{tabularx} \\ \\ \relax
25:10 & \begin{tabularx}{0.7\textwidth}{X} 她們去買的時候，新郎到了。那預備好了的，與他進去共赴婚宴，門就關了。 \end{tabularx} \\ \\ \relax
25:11 & \begin{tabularx}{0.7\textwidth}{X} 其餘的童女隨後也來了，說：『主啊，主啊，給我們開門！』 \end{tabularx} \\ \\ \relax
25:12 & \begin{tabularx}{0.7\textwidth}{X} 他卻回答：『我實在告訴你們，我不認識你們。』 \end{tabularx} \\ \\ \relax
25:13 & \begin{tabularx}{0.7\textwidth}{X} 所以，你們要警醒，因為那日子，那時辰，你們不知道。」 \end{tabularx} \\ \\
[1ex]
\hline
\hline
\end{longtable}
$^{1}$各位弟兄姊妹,早晨,新年快樂.
我想我們每個人都有一個特性.
很喜歡聽故事.
我想在座的弟兄姊妹都不例外.
特別那些很感人的故事.
結局很出人意外的故事.
我們全部豎起耳朵來聽.
所以故事真的很吸引人.
耶穌是一個講故事的高手.
因為他講到很喜歡用比喻.
有一個人.
出去殺種,他把種子殺在田裡.
有四種不同土地.
有個爸爸,他有兩個兒子.
小兒子要分家產,提早分家產,反正遲些你都要給我.
就是很出名的.
浪子回家的故事.
這個故事真的太出名了.
很多世界出名的文學家.
都說這個故事是聖經裡面.
最出色.
最感人的短篇小說.
所以有很多電影的情節.
都是抄浪子回頭的.
一個家庭.
有些反叛的兒子.
不知怎樣跟家人吵架.
離家出走.
這個兒子在外面很多遭遇.
然後.
最後.
他竟然說.
想起.
家裡有個爸爸.
不如我回去看看他.
他就想好了台詞.
他說我不配做你的兒子,我做你的傭人.
我服侍你等等.
誰知這個爸爸竟然無條件接納他.
而且用很大很深.

$^{41}$很特別的愛.
來愛他.
我不知道在座各位弟兄姊妹.
你信耶穌.
有沒有因為這個浪子回頭的故事.
如果有教會請我去講報道會.
很多時候我都講.
浪子回頭的比喻.
因為耶穌的愛.
就是這個慈愛父親的愛.
今天我都想講一個故事.
十個童女的故事.
這個故事為什麼這麼出色呢?.
因為它令我們想很多事情.
這十個童女.
在等.
等什麼?.
等新郎哥出現.
猶太人的習慣很特別.
不是新娘最重要.
我們中國人結婚,每個人都說你去婚禮,我看看新娘有多漂亮.
其實你認識的新娘.
那天是特別漂亮.
有一次我在教會主持婚禮.
見到新娘.
其實我跟她很熟,見她在教會很久.
我都忍不住說.
原來你這麼漂亮.
可能我說錯話.
她說牧師姐你說我以前不漂亮.
今天才漂亮.
因為新娘就是那天.
穿最漂亮的婚紗.
化妝化得最漂亮,髮型最漂亮.
總之以最漂亮的面貌見人.
但是猶太人.
是新郎哥更加重要.
我覺得這個故事.
令你想很多事情.
其實.

$^{81}$你要問很多問題.
我今天都想問大家幾個問題.
第一.
新郎哥是誰?.
他們在等新郎哥.
請問你有沒有等?.
你知不知道新郎哥是誰?.
就是主耶穌.
主耶穌就是主角.
我們等什麼?.
等耶穌回來.
大家都是基督徒.
知道耶穌有一天會在空中降臨.
我們等了二千年.
祂還沒回來.
好像那十個童女一樣.
等到晚上都睡著了.
突然有人說.
新郎哥來了,全部就起來.
立刻去迎接.
迎接因為天黑.
還要點燈.
於是這個故事最關鍵的地方就出現了.
點燈.
是油燈.
一十個童女.
分了兩批.
聖經說有五個是聰明的.
有五個是比較蠢的.
我又要問大家.
你是聰明的還是蠢的呢?.
這樣都要問,我當然是聰明的.
但是聖經說.
比率是一半一半.
有一半是聰明.
但是很可惜.
有另外一半是愚蠢.
你們分兩邊坐.
我只是舉例.
我當余牧師那邊是聰明的.

$^{121}$因為牧師領主.
這邊沒有那麼聰明.
一半.
就是一半人聰明.
一半人蠢.
都挺可怕的.
為什麼他們是蠢呢?.
就是燈裡面沒有油.
沒有預備好油.
我記起我小時候.
很喜歡唱一首歌.
就是《我的燈》.
《需要油》.
《PowerPoint》就打了出來.
我想請詩琴幫你先彈一次.
然後我們一起來唱.
我不知道你是否年青.
我們都曾經年青過.
余牧師都曾經年青過.
我就早五十年前已經開始教會.
那些少年團契唱這首歌了.
你們可能二十年前,三十年前都唱.
請詩琴幫我們彈一次.
,很容易,我們一起唱,懂得唱的唱大聲一點,不懂得唱的一學就懂得唱了,好,預備起!.
是我燈永遠發亮堂.
我的燈需要氧吸諸上池下.
是我燈發光導轉作旺.
昌和散啦,昌和散啦,昌和散啦,茂茂黃金黃.
昌和散啦,昌和散啦,昌和散啦,漢娥娥.
讓我們都做聰明那批同類就沒有一半一半.
我們百分之百都是聰明的.
因為我們的燈裡面有油.
我們的燈為什麼需要油呢?.
我們就要明白在以色列當時的鄉村裡面.
他們用的是哪一種油燈.
我特意帶了一盞油燈來.
這盞油燈就是我去以色列的時候.
去拿撒勒,就是耶穌長大的地方.
就參觀一個拿撒勒村.
他就將那條村佈置到耶穌當年的光景.

$^{161}$就看到很多東西,有橄欖樹.
還有拿橄欖油的壓榨,是怎樣壓榨橄欖油.
又養羊,是怎樣養.
有些人就穿了牧羊人的衣服.
臨走呢,因為你要買門票進去.
每人送一盞這樣的燈.
所以你看到這盞燈很淺的.
這裡有一條燈芯.
你點燃它,這裡就裝油.
因為裝得多少很少.
所以一會兒就用完油.
如果你沒有預備後備油.
很快燈就熄了.
另外有個解釋.
因為當時的鄉村是很黑,晚上黑.
所以他們點的不是這盞燈.
燈在希臘文裡面都可以解作火把.
一個火把,那些光就強一點.
那些木啊,那些柴啊.
一樣是淋一點油下去.
那火把才點得著.
所以無論你是當.
這種普通人家的燈或者是火把都好.
油是必須要預備的.
因為新郎哥原來是要探全村裡面每一家人.
如果那條村有很多房子.
他就逐家逐戶去所謂「串門子」.
去跟人打招呼.
那些人就說「哎呀,恭喜新郎哥啊」.
請他吃點東西.
所以每家都坐一下.
所以不知道搞多久.
最後才去新娘的家.
接新娘去婚宴.
去擺酒的地方開始飲宴.
所以搞了很久.
那十個同類就等到睡覺.
突然有人說「新郎哥來啦」.
就全部起來.
誰知那些愚蠢的就發覺.

$^{201}$哎呀,糟糕了,原來我沒有帶油.
那盞燈很快就熄滅了.
所以他就問那些有油的.
「你可不可以借點油給我啊」.
有什麼可以借的嗎?.
如果我未信耶穌.
我將來回到天堂.
我假設給得起天堂守門口.
他就問「你憑什麼進來天堂?」.
「我是華僑基督教會的會友」.
「余牧師幫我浸的」.
「哎,我進得了天堂」.
彼得說「我不認識你」.
「你到底有沒有信耶穌的」.
你要接受耶穌做救主.
你要重生得救.
你才能進得天堂.
所以聰明的就預備訂油.
所以那盞燈時時發光,有光.
所以我就要問幾個問題.
請問什麼是油呢?.
你今天有沒有油呢?.
那盞燈是不是時時都亮呢?.
上個星期紐西蘭的所有教會.
大部份西人教會就舉行「Open Heaven」.
可能大家聽過.
就在Avena那裡.
一個大的體育館.
有萬幾人去的.
以前我和師母都去過.
不過現在年紀大.
而且很難泊車.
Shine TV有現場轉播.
7時半開始.
坐在家裡看電視.
我們都很感動.
真的有萬幾人.
就在Avena那裡.
很多年輕人就去前面.
一邊敬拜一邊跳.

$^{241}$因為他們盼望天開了.
聖靈降臨.
然後燒著他們.
讓他們火熱愛主.
大家一起為紐西蘭的教會祈禱.
為紐西蘭的教會復興祈禱.
為年輕人下一代祈禱.
為傳道人祈禱.
為求神興起更多工人祈禱.
所以這個就是.
有火點起靈火.
靈火再焚燒.
其實我們每個基督徒都很需要聖靈的火.
在我們的心裡燃燒.
所以你的燈.
第一.
有沒有點著.
火夠不夠旺.
然後聖靈有沒有在你的心裡工作.
所以我就希望今天.
我解釋這個比喻.
能夠問大家這些問題.
所以這個故事真是很精彩.
讓你可以反覆去思想.
所謂令你心動.
令你真是可以反覆去問.
新郎哥是誰.
你有沒有等他來.
你是否預備迎接耶穌再回來.
如果你真的見到耶穌.
那你會跟他上天堂.
還是我沒得上天堂不讓我進去.
你有沒有天堂的入場券.
很多人都說如果我浸了.
受洗了就有天堂入場券.
余牧師和我都做了牧師學了.
不一定的.
你不要以為你是會友就一定可以進入天堂.
有很多是掛名基督徒來的.
他信因為他想追求一個女孩.

$^{281}$那個女孩說你信耶穌.
嫁給你.
那不就多大話信.
是假的.
不是真悔改的.
因為他的行為沒有改變.
又是以前那樣.
所以我們真的要經歷主耶穌.
你真的要重生得救.
你跟耶穌建立一個緊密的關係.
然後你才可以見這個新郎哥.
所以聖經經常說.
教會就是新婦.
我們這群信徒.
這群弟兄姊妹.
基督說叫教會就是新婦.
所以新婦要裝飾整齊.
等候丈夫.
你有沒有裝飾漂亮呢?.
你像不像樣呢?.
你有沒有穿上婚紗呢?.
等新郎哥來.
嘩!原來他的新娘這麼漂亮.
然後我們又要問.
你是聰明還是愚蠢?.
你有沒有準備呢?.
你預備些什麼呢?.
所以下面我就繼續講有三樣東西.
第一.
我們都要學習等候的功課.
我們都在等待當中.
等呢.
沒有人喜歡的話.
我們會等車.
等人.
約了人去酒樓喝茶.
誰知等了很久都不來.
我在香港年輕的時候.
我們的團契裡面.
就有一個男團友.

$^{321}$他就住在北角.
再下去一點.
就是近側漁涌那裡.
那裡有一個地方叫做麗池.
美麗的麗.
池水的池.
是一個好像舞廳那樣.
他就住在麗池附近.
原來他真的是例牌遲到.
約了他去喝茶.
約了他去哪裡.
不是遲十五分鐘.
起碼都遲半個小時.
所以我們每個人都說.
你不要以為你住在麗池.
就例牌遲到才行.
下次我們不等你.
所以我們等人.
最近我在紐西蘭就要做一些小手術.
你知道沒錢去私家醫院.
那就排政府.
政府出錢.
醫生已經說了.
要排很久的.
等很久的.
半年也說不定.
九個月也說不定.
不知道什麼時候.
那沒辦法了.
想政府出錢.
那就排.
最後去年年底就有機會排隊了.
就做了一個小手術.
等.
我們都很不耐煩.
特別現在這個時代是一個.
instant的時代.
喝什麼.
即溶咖啡.
喝什麼即沖奶茶.

$^{361}$要快.
什麼都快.
特別我是一個很心急的人.
所以師母經常說我.
你不要這麼急.
什麼都要快.
去酒樓點了點心.
這麼久都不來.
催那些夥計.
為什麼我點了蝦餃,燒賣還不來.
等.
我們時時都在等.
我們看一段聖經.
所以你們要警醒.
因為不知道你們的主.
是哪一天回來.
求主讓我們知道.
幾更天會有賊來.
你就會說.
不是人到我那裡去.
你們知道的.
你們就要預備.
因為你們想不到的時候.
人子就來了.
想不到的.
我們真的等了這麼多年.
我們說耶穌會回來.
特別有很多所謂的先知.
有說過.
1970年代耶穌回來了.
那些全部都是假先知.
因為聖經講得很明白.
耶穌回來的那日子.
那時辰.
是沒有人知道的.
連子也不知道.
只有父知道.
連耶穌自己都不知道何時回來.
祂說是神先知.
神要祂何時回來.

$^{401}$祂何時回來.
所以我們不要亂猜.
但是我們一定要有一個很焦急.
在這裡等的心.
近了.
所以啟示錄第21章說.
主必快來.
主啊,我願你回來.
主必快來.
這裡.
因為我的眼睛不太好.
不如請大家一起讀這段經文,好不好?.
我身在高宗聖會,這世外還沒有過去,你這次都要承擔,天地都規矩,我的話若不規矩,但大日子如假先知,沒有人知道,而天生的事你也難知道,只有祂知道,唯有父知道.
真是原來沒有人知道的.
是天父才知道.
所以我們總是在等.
我們都要在這裡等.
所以第一個就是我們在這裡等.
第二點.
就是我們活在很多不知道的事當中.
所以他們這種未知.
現在我都年紀大了.
有時很多長者就說.
真的不知道何時回天家.
因為老人家有三樣東西等的.
等吃飯.
等睡覺.
然後又等死.
我們是三等公民.
等,甚麼都等.
如果你知道你何時死.
你會說這不就好了.
我知道我明年就死了.
不是的,如果你知道你明年就死了.
你都不知道怎樣辦.
哎呀,這麼快就來.
我就死了.
哎呀,死得了.
所以神很大恩典.
沒有人知道你何時離開世界.

$^{441}$是神的恩典.
我們有時咒人.
你沒有好死.
即是你死得不好.
我來了紐西蘭之後.
就見到我一個以前香港的團友.
他的丈夫是光緯牧師.
所以光師母葉壽娟.
就是我在香港教會的團友.
大家都知道光牧師很突然回天家.
所以師母很傷心.
他這麼快就走.
一點準備都沒有.
我說神的恩典.
如果他又病.
又進了醫院.
又插滿了喉.
又說不了話.
又不省人事.
你更辛苦.
有些人睡在醫院拖很久.
拖幾年都不走.
對於家人,親友來說.
是一個很大的負擔.
所以神是有恩典的.
雖然好像很突然.
你沒有準備.
不是說病病病就離開.
是突然.
但是神就叫我們經歷.
祂是信實的.
祂總有恩典.
所以很多事情是未知的.
我們就是活在未知的當中.
當然有些事情我們不知道.
我們很難估計.
但是我們就說凡事總有主的美意.
凡事總有主的安排.
主的時間總是最好的.
In God's time.

$^{481}$In God's way.
In God's method.
在神的時間.
神的方法.
神的處理.
我們Trust and obey.
只要信靠,信服我們的神.
我們讀一讀那段經文.
請大家一起讀.
「非常同歸於盡,無歸於盡」.
「全國各地的光榮」.
「當孫子發表了言義的誓言」.
「理論就知道眼前緊要」.
「天下人的夢境也要參與」.
「傳統的心動人心感激」.
「天官的道理」.
「保守的宗旨」.
「一切都完美了過後」.
「一切都已成功」.
「一切都完美了過後」.
「天下人的心動人心感激」.
「理論就知道眼前緊要」.
「天下人的夢境也要參與」.
有戰爭和戰爭的風聲.
現在我們今天這個世界就是.
經常說打仗.
特朗普又經常恐嚇人.
伊朗你還不和談.
我砌你的.
你有後果的.
這個世界很多戰爭還在發生.
俄烏戰爭.
中東那些緊張的局勢.
還有地震必然越來越多.
越來越頻繁.
越來越猛烈.
有人估計最近這幾年.
加起來的地震.
比過去五十年的地震還多.
即是地震那種密度.

$^{521}$最近不是聽到.
連台灣都有地震.
中國陽光又有地震.
廣東的陽光.
所以很多地方都有地震.
而且地震的強度.
比以前高很多.
動不動就說六級七級地震.
可能以往是三四級.
即是地震的頻繁密度.
所以耶穌說好像婦人.
生嬰兒的時候.
那種的陣痛.
是越來越頻繁.
越來越猛烈.
在座很多媽媽.
生嬰兒真的很辛苦.
所以我們要更加敬重我們的老婆.
因為她真的生嬰兒很辛苦.
陣痛那種的情況.
所以我們要知道.
不單止我們是活在未知當中.
還有第三.
有些事是無法借的.
余牧師可不可以借點愛心給我.
我不夠愛心.
問他借錢就可以.
借愛心是無法借的.
我的信心不夠.
有誰借信心來.
我不夠得心.
有沒有人借點得心給我.
哪有得借的.
是你自己努力的.
你都不愛主.
那怎麼借點愛呢.
你都沒有信心.
當我們說要一起有信心.
要光強.
要勇敢.

$^{561}$要站立得穩.
你都危危乎就快跌倒.
那怎麼站得穩.
所以有很多事是無法借的.
所以那五個愚蠢的.
問那些聰明的.
借點油來.
借了給你我便沒有.
因為不知道新郎哥何時來.
可能還要再等一段時間.
他才真的來到我們那裡.
所以很多事是無法借的.
我很感覺現在是末世的時候.
耶穌回來真的很近.
所以下面我很想今天跟大家提的.
麻煩下一章.
這些經文應該大家都熟的.
我就不讀了.
我們應該怎樣預備油.
第一 我們是光.
耶穌在登山寶訓說的.
「你們是世上的光,世上的焰」.
你是.
你不是去扮的.
你不是去做的.
你就是光.
你的生命是有光發出來的.
所以各位弟兄姊妹.
保羅在他的書信裡面經常提到.
你們在這個世上要好像明光照耀一樣.
燈不是放在斗底下.
而是放在燈台上.
就照亮一家的人.
這盞油燈如果我擺高一點.
就照得光一點.
擺低一點就沒那麼光.
如果我這個燈又大一點.
那就照得光一點.
如果中間找個高台我擺這個燈.
那都可能照到這個禮堂.

$^{601}$光光的.
所以我們就是光.
我們這個世界要做光啊各位弟兄姊妹.
你是炎.
你是防腐作用.
有很多食物.
加點鹽下去馬上很有味道.
沒有鹽.
沒有味道.
你今天有沒有影響力?.
人家很喜歡和你聊天.
因為你是一個很善良的人.
你心地善良.
你的言語很清潔.
很喜歡和你做朋友.
你不是那些經常罵人.
經常說粗話.
經常說一些很低俗的話.
你的言語好像嚴調和一樣.
你是光.
有光發出來.
光和暗是不可以並存的.
請問你晚上回家.
如果是你第一個進屋.
你第一件事是什麼?.
亮燈.
開了門之後亮了燈.
因為你不亮燈.
黑漆漆的.
什麼都看不到.
有光的地方.
黑暗就走了.
今天我們在這個世界.
要是光.
照亮你身邊的人.
在你的職場.
在你辦公的地方.
為耶穌做美好的見證.
第二.
在末世的時候有一些表現.

$^{641}$我想請大家一起讀一讀.
提摩太后書第三章一到五節.
請.
提摩太后書第三章一到五節.
忘恩負義.
心不勝怨.
無親情.
不解怨.
好施慘.
不能自仁.
心情苦惱.
不能自誠.
賣主賣友.
任以往.
自告自利.
愛而不愛.
有敬畏的內心.
卻未見證義.
正義.
提摩太后書三章一到五節.
是很重要的經文.
末世的時候.
有些表現.
就是人會愛錢.
專顧自己.
愛自己愛錢.
貪財.
狂傲.
狡猾違背父母.
我叫他三個愛.
兩個不愛.
第一愛什麼呢?.
愛自己.
然後愛錢.
第三.
愛艷樂.
就是今天我們這個世界.
「我」很大.
「我呀」.
知不知道「我呀」.

$^{681}$主一下自己的鼻子.
我請你吃飯.
多威風.
愛自己.
自己最大.
然後愛錢.
我看香港新聞.
剛剛過年.
金多寶.
如果一注獨中.
派三億港元.
所有馬會投注站.
每個投注站都排長龍.
大家想賭一下.
電視台專門訪問人.
為什麼你買寶.
我當然想中頭獎.
我去郵埠買那間大屋.
說到他好像真的中獎.
另外訪問一個女士就客氣些.
我沒錢供子女去讀大學.
如果我中了就供子女去外國讀大學.
這個好媽咪.
另外還有一個女士還客氣.
她說我隨便買下.
中其他獎幾十元都好.
頭獎輪到我.
每個人都有發財夢.
紐西蘭這裡都有.
很多個號碼.
一中了就很多錢.
發夢了.
然後第三個.
愛燕樂.
今日這個世界.
大家都喜歡享樂.
Entertainment.
所以有時我們牧師講道.
都要講些笑話.
吸引o下人.

$^{721}$如果唔係全部訓哂.
Entertainment.
我地都好想得啖笑開心o下.
三個愛.
但好可惜.
我地都有兩個唔愛.
第一.
唔愛良善.
今日呢個世界係一個好兇殘ge 世界.
世界各地都有好多槍殺案.
係雪梨好出名gwo 個海灘.
D 人開槍.
紐西蘭都有.
冇良善.
兇殘 狂暴 毀謗 違背父母.
唔聽父母.
有D 殺爹地媽咪ge 咁ge 罪案.
但係另外一個唔愛更嚴重.
我地唔愛神.
昨晚我帶一個普通話ge 弟兄查經班.
arm arm 查到哥林多前書最後一章第十六章.
保羅係哥林多前書十六章第廿二節gwo 度話.
若有人不愛主.
這人可救可助.
如果你唔愛主.
你是可救可助.
係提摩泰前後書裡面都話.
而家呢個世代.
特別提摩泰後書.
係保羅ge 好似遺囑咁樣.
因為好多賢經ge 人都話.
保羅最後寫ge 一封信.
就寫比佢ge 屬靈ge 真兒子提摩泰第二封信.
所以保羅特別係提摩泰後書第四章就話.
當跑ge 路我已經跑盡了.
所信ge 道我已經守住了.
佢知道自己好快會離開世界.
所以佢吩咐提摩泰.
佢話你係呢個末世ge 時代.
必須要愛神.

$^{761}$你要教導你ge 弟兄姊妹要愛神.
係啟示錄七封書信裡面.
第一封係寫比以弗所教會.
唔好失去jor 當初ge 愛.
起初ge 愛.
the first love.
你arm arm 信耶穌ge 時候你好愛主.
咁今日呢.
你參加完培靈會好似打jor 支強心針咁.
但係一兩個月之後又回復現狀仲慘過以前.
起初ge 愛.
我想請問去jor 邊呀.
我地失去jor 起初ge 愛.
呢個係第一封信寫比以弗所.
然後啟示錄七封信最後一封.
寫比老弟家教會.
耶穌點講呀.
佢話你唔好不冷不熱呀.
唔係我會將你係我口裡面吐jor 出黎.
look warm.
如果有人請你飲咖啡.
gwo 杯咖啡都涼涼地架.
你話有冇搞錯呀.
整杯咁ge 咖啡黎都唔係熱辣辣架.
不冷不熱.
一係你就叫杯凍咖啡.
一係就叫杯熱咖啡.
冇人飲一杯.
凍jor 擺jor 係度好耐架啦.
沖jor 出黎又未遞比你.
可能gwo 間咖啡好忙D 侍應又唔得閒.
冇拎埋黎你張台度.
你唔會鍾意架嘛.
所以各位弟子姐妹.
點解我地對神個愛.
有時候好似變淡jor 呢.
你問o下今年2026年.
你如果早十年前.
2016年.
你gwo 時係咪愛耶穌多D 呢.

$^{801}$你gwo 時係咪勤力D 傳福音呢.
奉獻俾教會都奉獻多D 呢.
而家退jor 休囉.
冇乜錢囉.
靠積蓄架咋.
咁所以你好多時候.
淡jor .
各位弟子姐妹我好想勸大家.
唔好做gwo 批愚拙ge 徒女.
同埋好落係提摩太后書呢度講.
三個愛 兩個唔愛.
千祈唔好.
愛自己 愛錢財 愛人.
兩樣ye 唔愛.
唔愛良善 唔愛神.
讓我地ge 燈時時拉著.
讓我地ge 燈預備好入.
請你低頭我地同心祈禱.
親愛ge 天父好感恩.
今日能夠黎到基督教華僑教會.
同弟子姐妹分享十個童女ge 比喻.
求聖靈都係我地內心裡面工作.
你親自黎感動我地.
向我地ge 心靈講說話.
讓我地都去選做聰明gwo 一批.
讓我地ge 燈係時時點著ge .
可以為你發光 為你燃燒.
求主你讓我地都預備好入.
我地時時都預備好入.
有聖靈係我地裡面.
聖靈係我地每個人ge 生命裡面工作.
讓我地唔好失去起初ge 愛.
又讓我地唔好不冷不熱.
好似老弟家教會咁樣.
讓我地熱切ge 愛你.
火熱ge 燃點起上來.
求主親自工作 幫助我地.
聽我地獻上ge 祈禱.
奉耶穌基督ge 命求.
\newpage



\section{希伯來書 3:7-4:11}
\label{sec:031705}
\textbf{創造之始的約:工作關係 (英/粵)}~~~~歐世達牧師
\newline
\newline
連結: \href{http://www.church.com.hk/acms/content.asp?site=cdc\&op=show\&type=product\&code=031705\&layout=sermon}{\texttt{031705}} ~~~~ 語音日期: 2026-02-22
\newline
\newline
\hyperref[sec:031704]{\small{< < < PREV SERMON < < <}}
~
\hyperref[sec:index_chronic]{\small{[返順時目]}}
~
\hyperref[sec:031706]{\small{> > > NEXT SERMON > > >}}
\newline
\newline
希伯來書 3:7-4:11
\newline
\begin{longtable}{cl}
\hline
\hline
章節 & 經文 (和合本修訂版)\\
\hline
3:7 & \begin{tabularx}{0.7\textwidth}{X} 所以，正如聖靈所說：「今日，你們若聽他的話， \end{tabularx} \\ \\ \relax
3:8 & \begin{tabularx}{0.7\textwidth}{X} 就不可硬著心，像在背叛之時，就如在曠野受試探之日。 \end{tabularx} \\ \\ \relax
3:9 & \begin{tabularx}{0.7\textwidth}{X} 在那裡，你們的祖宗試探我，並且觀看我的作為， \end{tabularx} \\ \\ \relax
3:10 & \begin{tabularx}{0.7\textwidth}{X} 有四十年之久。所以，我厭煩那世代，說：他們的心常常迷糊，竟不知道我的道路！ \end{tabularx} \\ \\ \relax
3:11 & \begin{tabularx}{0.7\textwidth}{X} 我在怒中起誓：他們斷不可進入我的安息！」 \end{tabularx} \\ \\ \relax
3:12 & \begin{tabularx}{0.7\textwidth}{X} 弟兄們，你們要謹慎，免得你們中間有人存著邪惡不信的心，離棄了永生的神。 \end{tabularx} \\ \\ \relax
3:13 & \begin{tabularx}{0.7\textwidth}{X} 總要趁著還有今日，天天彼此相勸，免得你們中間有人被罪迷惑，心腸剛硬了。 \end{tabularx} \\ \\ \relax
3:14 & \begin{tabularx}{0.7\textwidth}{X} 只要我們將起初確實的信心堅持到底，就在基督裡有份了。 \end{tabularx} \\ \\ \relax
3:15 & \begin{tabularx}{0.7\textwidth}{X} 經上說：「今日，你們若聽他的話，就不可硬著心，像在背叛之時。」 \end{tabularx} \\ \\ \relax
3:16 & \begin{tabularx}{0.7\textwidth}{X} 聽見他而又背叛他的是誰呢？豈不是跟著摩西從埃及出來的眾人嗎？ \end{tabularx} \\ \\ \relax
3:17 & \begin{tabularx}{0.7\textwidth}{X} 神向誰發怒四十年之久呢？豈不是那些犯罪而陳屍在曠野的人嗎？ \end{tabularx} \\ \\ \relax
3:18 & \begin{tabularx}{0.7\textwidth}{X} 他向誰起誓，不容他們進入他的安息呢？豈不是向那些不信從的人嗎？ \end{tabularx} \\ \\ \relax
3:19 & \begin{tabularx}{0.7\textwidth}{X} 這樣看來，他們不能進入安息是因為不信的緣故了。 \end{tabularx} \\ \\ \relax
4:1 & \begin{tabularx}{0.7\textwidth}{X} 所以，既然進入他安息的應許依舊存在，我們就該存畏懼的心，免得我們 中間有人似乎沒有得到安息。 \end{tabularx} \\ \\ \relax
4:2 & \begin{tabularx}{0.7\textwidth}{X} 因為的確有福音傳給我們像傳給他們一樣；只是所聽見的道對他們無益，因為他們沒有以信心與所聽見的道配合。 \end{tabularx} \\ \\ \relax
4:3 & \begin{tabularx}{0.7\textwidth}{X} 但我們已經信的人進入安息，正如神所說：「我在怒中起誓：他們斷不可進入我的安息！」其實造物之工，從創世以來已經完成了。 \end{tabularx} \\ \\ \relax
4:4 & \begin{tabularx}{0.7\textwidth}{X} 論到第七日，有一處說：「到第七日，神就歇了他一切工作。」 \end{tabularx} \\ \\ \relax
4:5 & \begin{tabularx}{0.7\textwidth}{X} 又有一處說：「他們斷不可進入我的安息！」 \end{tabularx} \\ \\ \relax
4:6 & \begin{tabularx}{0.7\textwidth}{X} 既有這安息保留著讓一些人進入，那些先前聽見福音的人，因不信從而不得進去， \end{tabularx} \\ \\ \relax
4:7 & \begin{tabularx}{0.7\textwidth}{X} 所以神多年後藉著大衛的書，又定了一天—「今日」，如以上所引的說：「今日，你們若聽他的話，就不可硬著心。」 \end{tabularx} \\ \\ \relax
4:8 & \begin{tabularx}{0.7\textwidth}{X} 若是約書亞已使他們享了安息，後來神就不會再提別的日子了。 \end{tabularx} \\ \\ \relax
4:9 & \begin{tabularx}{0.7\textwidth}{X} 這樣看來，另有一安息日的安息為神的子民保留著。 \end{tabularx} \\ \\ \relax
4:10 & \begin{tabularx}{0.7\textwidth}{X} 因為那些進入安息的，也是歇了自己的工作，正如神歇了他的工作一樣。 \end{tabularx} \\ \\ \relax
4:11 & \begin{tabularx}{0.7\textwidth}{X} 所以，我們務必竭力進入那安息，免得有人學了不順從而跌倒了。 \end{tabularx} \\ \\
[1ex]
\hline
\hline
\end{longtable}
$^{1}$大家早安,好開心能夠回來和你們在一起..
謝謝大家讓我有機會再次和你們分享神的話語..
我曾經做過兩個星期的一個系列,關於改革中的一些原則..
我曾經和大家分享神如何呼召我們成為傳道者,不僅僅是我..
而我今天開始會有一個新的系列..
我會從聖經的開始講一些說話..
我會從舊約的文藥開始講..
雖然聽起來不是一個很有趣,很有興趣的題目,.
但我可以告訴你,我們所明白關於這個約的時候,是幫助我們去明白耶穌基督為我們成就的一切..
如果我們不能夠正確掌握這個定義的時候,會讓我們產生很多在信仰和生活上的一些錯誤..
我們一起去禱告..
我們親愛的天上的父神,謝謝你再次讓我們回到你的家中..
我們求你幫助我們,讓我們明白你自己的話語,.
以致我們能夠活出你所呼召我們進入的生命,為你自己榮耀..
主耶穌基督,聖名已求,阿們..
大家都說我們已經進入了一個資訊的時代..
因為在我們觸手可及的資訊中從來沒有試過有這麼多..
如果我們能夠用聲音控制的時候,同樣地我們可以得到一些資訊..
在最近很冷的時期的時候,我借了我朋友一間旅渡假屋..
他們的熱水是不動的,是壞了的..
究竟我懂不懂得去檢測,檢修,或者去修理這個熱水呢?.
當然是不行的,但是所有我要知道的就是在我指尖之間..
當我有這個資訊的時候,就幫助我去檢視事情,.
以致我能夠將那個熱水修理好,.
以致在一個一月份很冷的晚上,我能夠有熱水沖涼..
聖經裡面其實有關於很多的資訊,幫助我們去了解和幫助我們解決一些問題..
而神的話語當中實在有很多的幫助..
但這不是說我們找到一個問題的時候,就去那裡為了解決問題去看..
而是我們要更加明白一個比較廣闊的概念和主題,以致我們能夠知道要怎麼去做..
以致我們能夠將我們的生命在根基上面建立得好..
在聖經裡面其實有一個重複的主題不斷的出現,也是對我們非常有重要的..
那就是文約.而這個文約其實有很多的資料讓我們去明白和知道..
它讓我們知道很多關於神的事情,也有關於我們自己的事情..
也都是在聖經裡面不同故事裡面的主題線,幫助我們去明白和體驗它是怎麼進行的..
如果你看到這個影像的時候,去到最底的時候,你會看到有很多人在當中得到幫助的人數的文..
就是因為我看的確是一個關於咖啡的內容,那我只要去到某個位置的時候我就可以停,然後就去看我需要看的內容..
如果我們要在這個短片裡面,在不同的地方停的時候,其實每個停的位置幾乎都是文約的位置..
這個文約的內容裡面能夠幫助我們體驗神在不同的歷史或計劃當中的一些做法..
究竟什麼是文約呢?.
文約是雙方之間的一個協議,類似一個合同的情況..

$^{41}$而在舊約的時代,文約是關於一個征服了國家的皇帝和一個被打敗的敵人的關係..
而這個文約就是規定了那個附庸國,他應該要做的事情..
例如他們要交稅,他們不可以再打仗打回那個皇帝,或者其他的事情..
但是有時候文約也是一個對等的兩個團體的關係..
這個是聖經裡面第二個的約,這個是阿當和夏娃他們結婚的時候的一個約..
這個也好像是他們的婚禮之中的一個誓言,他說這個是我骨中的骨,肉中的肉..
因為在當中,阿當不單止說夏娃是和他對等的,也都告訴大家聽,阿當說他會對夏娃是等於對他自己一樣的尊重和處理..
所以文約就告訴他們他們是誰,也是他們會做什麼事情..
在文約當中有不同的群體,也都當中有他們的義務或者誓約..
當我們看文約的時候,希望能夠讓大家看得更清晰的是在聖經裡面的不同的故事..
所以我希望藉著在當中看的時候,讓大家看到聖經當中不同主題的故事..
從我們這樣的追溯主題當中,我們就能夠看到關於耶穌,關於我們自己,和關於將來向天國的事情..
所以讓我們進入去看文約..
這個是當神去做人,去創造當中,他和阿當和夏娃所立的文約..
這個不是一個很常規的文約方程式,所以有些難的要去完全了解..
創世的二章十五到十八節就是我們要看的內容..
我會讀出來.創世的二章十五到十八節.耶和華神將那人安置在伊甸園,使他修理看守..
耶和華神吩咐他說,圓中各樣樹上的果子,你可以隨意吃,只是分別視我樹上的果子,你不可吃,因為你吃的日子必定死..
耶和華神說,那人獨居不好,我要為他做一個配偶幫助他..
在這個文約當中有分開三部分,我們要看這三部分是怎樣進入..
首先就是神創造了阿當在圓中工作..
所以首先那個工作是我們生命當中的一個基礎..
第二方面就是神在圓中種了兩棵樹..
其中一棵樹就是分別視我樹,當中的果子神已經吩咐阿當夏娃他們不可以去吃..
所以我們看見就是我們相信神也是在我們生命當中的一個很重要的基礎..
第三方面就是神知道阿當獨居不好,所以神就創造了夏娃..
所以關係是我們人生基礎的第三方面..
這三個基礎,工作,信靠神和關係..
所以第十七節是當中最接近文約的內容..
如果阿當和夏娃他們吃了那棵樹的果子,他們會死..
所以這個文約有一個很簡單的方程式..
信從神,你就會活下去..
如果你不信從神的時候,你就會死..
這個看起來好像就是這樣去操作..
大家都知道阿當和夏娃吃了這個禁果..
他們不信神..
他們就被趕出這個園子..
他們就開始死亡了,他們的靈性就死亡了..
他們的靈性死亡是因為他們和神遠離了..
他們和地圖之間的關係也都破裂了..

$^{81}$他們的工作成為了一個勞苦勞役..
懷孕生產也是一個很痛苦的事情..
最後他們也有肉身的死亡..
信從神,聽神的話,你就會活下去..
不信從神,你就會死亡..
我不太喜歡這個文約的觀點..
這個很接近真理..
但是當一件事情是接近真理的時候,.
會讓我們有一個很錯誤的想法出現..
讓我們將神的觀點,對神的觀點扭曲了..
因為如果你這樣去看這個文約的時候,.
你就會覺得神是一個好像在宇宙間的一個警察一樣..
他就在等待,如果你做錯了,他就會過來打你..
這個也不是一個很健康的父母教育子女的方法..
因為他不會幫助你去明白什麼叫愛,反而只是一個恐懼..
他也不會幫助我們去知道,.
有時候從另一個觀點來看,.
就是從撒旦的引誘裡面,.
讓我們能夠更加清楚看見什麼叫做文約..
因為撒旦引誘夏娃和阿當,.
對神的美善產生懷疑..
為什麼神不會將這麼好的果子給你們吃呢?.
因為他不是還要告訴你,你要去信從神,.
而是告訴你要相信神和相信祂的美善..
我想一個很重要的事情讓我們看到的就是,.
當我們去信靠神的時候,神就會供應..
很多歐洲的一些畫作裡面,.
都將這個果子形容為一個蘋果..
但是在猶太的傳統裡面,.
那個應該是無花果或者一個石榴..
其他的有可能會是一些橘子那類的食物,.
但是如果我們從我們家裡面去想的時候,.
就可能是一個很大很美的亞洲梨..
那你就會想一下,.
為什麼神不讓你們吃這麼美的亞洲梨呢?.
其實阿當和夏娃他們的思想就是,.
神其實已經為他們預備了..
所以這個也提醒了我,.
這是神在創造之開始的時候,.
祂的文藥的一個本質,.

$^{121}$是我相信神,祂為我預備,供應我的一切..
而這個說法是說要信從或者不信從,.
信從或者不信從就會死,.
這個是很不同的事情..
我們其實很清楚的就是,.
我們要去聽從神的時候,.
這個就是我們要做的事情..
但是仍然你會看到一個很不同的想法..
我相信有兩個在創造的時候的故事,.
我們可以去看..
我其實已經說了一些了..
首先就是當神去創造這個阿當的時候,.
祂就將阿當放在原子裡面..
這個是神對阿當一個很大的供應..
因為阿當不需要到處去覓食,.
他不需要去吃那些樹枝或者草,.
然後去想那些東西究竟可不可以吃,.
有沒有毒,各方面..
其實有時候我在想,.
那些人怎麼去知道那樣東西吃了會死呢?.
但是對於阿當來說,.
這不是真的,.
因為神將他放在原子裡面,.
是一個很大的供應..
但是我們去想的時候,.
阿當放在那裡的時候,.
完全不是這樣的事情..
而我相信那個伊甸園是一個很豐富的原子..
就好像你想想那些比較貴的超級市場裡面,.
有很多的漂亮的水果..
你會想到那裡有很多的蘑菇,.
有菜和各類的食物可以供你去食用..
而所有東西已經是預備好,.
可以隨時去收割..
也是神為我們預備了一個美好的一餐..
當然阿當要在當中去工作..
但是我們可以想一下,這個是神在原子裡面為他所豐富的供應..
而第二方面我們看到的是,.
神為阿當預備了夏娃..
神可能要讓我們知道這個很特別的供給..

$^{161}$大家可能看這幅圖就想起了Tom Hanks的其中一套電影叫做《荒島餘生》..
當他自己去到一個無人島的時候,.
因為沒有人在身邊,.
所以他將排球變成了他的朋友,一個人的樣子..
他叫他做Wilson..
神當然可以叫阿當做同樣的事情..
神也叫阿當到處看看他所創造的所有動物,.
給他們名字..
我們知道狗也是一個很可愛的朋友..
但是他們沒有一個動物是阿當適合的伴侶..
所以神為阿當做了一個最好的伴侶,就是夏娃..
而夏娃是神為阿當做最好的伴侶的供應..
同樣的,阿當也是夏娃一個很重要的供應供給者..
所以你從這個眼光去看的時候,.
你就會看到文藥其實有很多是關於供給供應的情況..
所以當阿當和夏娃吃了分別善惡樹上的果子的時候,.
他們並不是一個相信神的情況..
所以你看到一個最重要的影響就是,.
當他們因為這個罪的時候,.
他們就將神所給他們的供應破壞了..
所以當他們被趕離開伊甸園的時候,.
他們的生活就要靠他們的奴婦才能得到供應,.
因為神的供應已經斷絕了..
所以我們從這方面看的時候,文藥的那方面就很清晰了..
首先神是最偉大的那位,他是創造者..
而阿當和夏娃就在這個文藥當中是一個次一等,.
比較沒有那麼重要的位置..
而他們的義務和我們的義務就是相信神..
而神的義務就是要供應我們..
所以當我們思考這件事情的時候,.
究竟這個工作和關係是怎樣在文藥當中去明白呢?.
這個當然是我們有某種的方法去相信,去順從神..
我們也很盛言的,我們看見神給我們一些命令,.
叫我們關於工作的,也有很多的命令關於這個關係上..
但我也覺得可以從另一個角度更加容易去明白..
就是當我們看見我們是神所創造的那個位置的時候,.
工作和關係本身就來自於我們裡面..
這個是我們人本身裡面的一部分..
這個是我們部分,為什麼我們是一個人類的本質..
在創世的一章二十六字的時候,我們看見神是按著他的形象去做人..

$^{201}$因為神去工作去創造人,所以我們同樣的按照神的模式去工作..
有時候我們都好像看不到神是在一個關係當中的時候,.
但我們只要細心去想的時候,神是一個聖父,聖子,聖靈,.
在三維一體的時候,我們就能夠去明白到..
也都我們當我們看見我們是被創造有一個關係的時候,.
這個是我們可以去進行的事情..
所以我們一方面看見這是一個本質上,我們如何在關係上面..
所以我們很自然看到的,在我們裡面,神給我們一個工作和關係上的..
但也都另一方面我們也看到的就是,在文藥當中,神給我們是一個神會供應,會給我們的事情..
而在當中給我們一些的緊張..
所以我們看到我們可以工作為我們自己的生活維持..
我們可以更有智慧去工作,或者更努力去工作,以致我們能夠給自己更好的生活..
我們才可以很有創意的去工作,以致我們能夠為我們家人帶來更好的生活..
而我們的關係其實都是一方面的供應..
如果你是比較年輕的時候,你就明白到,是我們的家人如何去供應我們..
而這個是藉著關係而給予的供應..
所以也都讓我們明白到,神為什麼要供應那些孤兒和寡婦,因為他們沒有他們的家人和關係的那種供應..
所以我們同樣再看這個張力在當中,神去創造人的時候,他也都有工作和一個關係在那裡..
那些事情都是一件很好的事情..
而在這個創造的文藥當中,神是要求我們去聽從祂,以致祂能夠供應我們..
所以我們要去思想,我們的工作和關係,如何我們直接去依靠神,我們得到供應,而不是單單我們去工作和我們的關係,自己去努力獲取..
所以我相信這一點讓我們去依著這個點去尋找聖經裡面的一個關係..
當我們思想神一個特別的供應的時候,是超過我們能夠藉著努力去工作或者人際關係而獲得的..
而神這個特別的供應,是藉著我們對祂的相信而獲得的..
而我們的相信神,就是藉著我們的遵從神,有時候看上去好像都不是一個很有效的方法..
有一個例子就是我們的十一奉獻..
神提醒我們要將我們擁有的一部分,將它給神自己..
有時候我們說是一個十分之一的情況..
如果我們回想農業社會或者舊約的時候,這個就是一個初熟的果子或者農產品..
而對於他們當時來說,是他們整年所期待的一個修成,一個最好的一部分..
而當你想想我們今天的時候,如果你想吃士多啤梨的時候,整年都有士多啤梨,不是一個最需要初熟的情況..
你試想一下,如果當時我們要等待初熟的士多啤梨,又紅又大又靚,那神說你要將它奉獻給神..
而在馬拉基書三章十節的時候,神就告訴我們這個特別的供應,只要我們相信祂的時候,我們就可以得到..
使我加有量,以此試試我是否為你們搪開天上的窗戶,傾福於你們,甚至無處可用..
這個可能和我們思想上有些不一樣的地方,我們給的時候就會得到..
因為神說,當你相信我的時候,我就會給你,給得更多更好..
另一方面就是以蔡書五十章也給我們一個很好的例子..
這個也是神說對以色列人說,如果你繼續相信祂的時候,神會給他們一個特別的供給..
因為當時以色列人說我們在禁食當中,你會聽我們說禱告嗎?.
但神告訴他們,雖然你們禁食,但你不是做我告訴你們要做的事情..

$^{241}$而神要他們所做的時候,其實就是和他們剛剛相反的事情要他們去做..
因為神要求他們鬆開鴨..
將他們的餅分給飢餓的人..
將一些沒有地方住的人,將他們接到他們的家當中..
那些沒有衣服的,給他們衣服遮體..
而所有的事情其實就是要求他們的一個犧牲..
而這個是神呼召我們去做的事情..
如果我們相信神會這樣去做..
祂就會供應給我們,是超過我們所想的..
而以色列書這一段經文是一個很好的經文,大家可以再去看..
而我們看到最後的時候,神是將這個乾的地變成一個很豐富肥沃的土地..
在這段經文,11到12節是這樣說的..
和從修路徑與人居住的..
我們同樣看到,這個就好像神為他們供應的一個原子的圖畫..
所以當我們相信神的時候,.
神的供應是超過我們憑著我們的工作或者我們的關係而獲得的那種風俗..
而我們從耶穌基督身上也看到一個更清晰的圖畫..
因為耶穌基督是完全的信靠神自己..
因為藉著他的信心和對神的倚靠,.
他能夠為很多人帶來醫治..
藉著他的信心,他能夠餵飽很多人..
但是耶穌基督自己卻沒有得到神給他的那個供應..
取而代之的是他自己在十字架上被釘死..
但是我們也看到神給耶穌基督一個很特別的一個供給..
在地上,他被釘死和被羞辱..
而我們從菲勒比書這一段經文裡面看到神怎樣高舉他讓他得到這個榮耀..
在菲勒比書八到九節這裡這樣說:.
所以我們在耶穌基督身上和我們剛才所看的舊約裡面關於十一方面各樣的事情,.
我們看到神一個物質上的供應有時不是立即..
所以我們從這個我們看到神供應的第一個原則,.
就是神會在物質上供應,但是在祂的時間裡..
而神通常不會給出這些物質上的供應,直到天國裡面..
所以就像是耶穌基督一樣,也和我們一樣..
所以我們不可以期望或者期待神在今世裡面給我們滿足我們所有的物質上的需求..
但是當我們相信祂,我們就知道,尤其是在天堂,我們會得到更多的這些物質上的供應..
而我相信這一切是神為了要讓我們在信心上成長..
就好像我們最近要看那些奧林匹克的運動員時,他們不斷地運動去增強我們肌肉和忍耐力..
同樣的,當我們相信神的時候,雖然我們看不到這一切的供應的時候,能夠幫助我們信心長成..
我們有時也問,我們怎樣在信心裡成長,但是當我們又看不到神的供應的時候?.
這就是神提醒我們要回到創造的時候,神怎樣去完全的供應,是豐足的供應..

$^{281}$神所造的地圖是好的..
會按時降雨,也會有花朵盛開..
那些蟲產品會生長,雖然有季節的轉變,但是神會供足..
但是這個世界,這個地圖不是天國..
我們知道有時候農產品也不會收成..
我們也看到經濟會崩潰..
有時候那些壞人他們會贏了,甚至我們很努力的工作,他們會贏了這個場面..
我們看到到處的被毀壞,也看到一些惡的痕跡在我們生命裡,到處都看到..
所以我們要等待神完全的一個供應在天國裡給我們..
所以我們回想第一個原則,就是神的物質的供應不是臨時的..
第二個原則就是神的供應是會給我們一個屬靈上的需要..
所以我們看到什麼是我們最需要的屬靈的供應呢?.
就是我們最得赦免..
而我們知道耶穌基督已經為我們付了屬價..
而我們相信神的這個供應已經藉著耶穌基督的死為我們完全的成就了..
另外一方面,耶穌基督給了我們這個聖靈..
就是因為我們得了這個聖靈,我們能夠得到屬靈裡面的各樣的供應..
我們很喜歡保羅在這段經文裡告訴我們,神如何在屬靈裡祝福我們..
在哥林多前書一章四到七節說:.
又因你們在他裡面凡事富足,口才,知識都傳遞..
正如我為基督作的見證在你們心裡得以兼顧,.
以使你們在恩賜上沒有一樣不及人的,.
等候我們的主耶穌基督顯現.」.
這段經文是關於我們那個文曰裡面的所有方面..
但是我想藉著這些內容,.
去將文曰和我們如何得到,.
去告訴你,如何去成就,去告訴你..
在希伯來書裡已經告訴我們,.
我們如何去得到神的供應,.
以及供應是什麼,已經告訴我們了..
而這些經文是告訴我們,.
我們其實就是要休息,要安息..
而我們回到創世的故事的時候,.
在第七天神造物的公元畢了就安息..
我想到這個安息的時候,.
其實不是神自己本性上的東西,.
而是我們作為一個人的時候,.
我們是需要這個安息休息..
因為我們在肉體上需要這個安息..
而更加重要的是,.

$^{321}$神是要藉著這個觀念去讓我們明白,.
我們也要在這個宿命裡面去休息..
因為安息是我們對神的依靠的一個指示..
我很喜歡看我們明尼蘇達州的.
奧林匹克跨國運動員Jesse Diggins..
這個會是他最後一次參加奧運..
而他給我們看到的是,.
他完全將自己擺上,即是給出來..
我們當他看見他衝過這個終點線的時候,.
他整個人就跌倒在那裡..
他自己說,.
他對於他自己怎樣計算他的成功,.
不是那天得到什麼獎品,.
而是有沒有將自己完全給出來..
當他能夠將自己完全給出來的時候,.
他就覺得他自己可以安心下來..
這也讓我想起另一個明尼蘇達州的運動員,.
就是Lindsey Vaughan..
她年屆41歲..
她仍然在這次參加奧運,.
但很不幸地將她的腳斷骨..
她也表達了一個觀念,.
就是唯有她能夠將自己完全擺出來,.
她才能夠安心..
很有趣的是,.
我們作為一個嬌肚仇腎,.
我們被要求將自己完全擺上..
在這段經文11節裡說到,.
我們務必竭力進入安息..
這是盡了所能去得到安息..
所以我們要進入安息,.
也是需要這個努力..
因為第十節裡說到,.
我們不是要藉著我們的工作,.
因為不是我們的工作去補足我們一切,.
給予我們一切..
是因為耶穌基督成就一切,.
所以供給我們一切..
所以當我們相信耶穌的時候,.
我們就能夠得到安息..

$^{361}$這段經文也讓我們看到,.
我們心裡面的那種掙扎..
這段經文問到我們到底喜歡什麼,.
或者我們喜歡哪一個人呢?.
我們究竟要相信神自己的供應,.
還是我們相信自己能夠供給?.
以前我們有時候會問,.
究竟我們為什麼不去相信神呢?.
或者我們學習的東西看不見神的供應呢?.
或者我們希望供應給自己,.
但其實是從一個錯誤的方向去看..
就好像亞當和夏娃,.
他們看到一個果子覺得好吃,.
他們就去拿了,.
而他們沒有聽從神..
而這段經文挑戰我們,.
如何彼此去激勵..
鼓勵我們如何一起進入神所給我們的安息..
而我們有時候想到,.
我們教會有很多事情要做,.
也有很多事情要做..
我們很重要的事情就是,.
我們要彼此勉勵,.
我們要進入這個安息..
是進入到神已經為我們成就一切的那個安息..
而是我們能夠安靜地等待神,.
將來要給我們的安息,.
給我們一切..
我希望今天你能夠有這個安息在這裡,.
讓你去想到神能夠完全的供應給你..
讓我們一起祈禱..
我們親愛的天父,.
我們為你的美善獻上感恩,.
你如何如此豐足地特別地供應我們..
但我們必須自我承諾,.
有很多時候我們要靠自己的能力,.
為自己的供應..
而我們應該首先去依靠你..
所以我們要求你給我們特別的恩典,.
去相信你..

$^{401}$因為我們知道,.
你是值得所有的榮耀,.
當我們去依靠你的時候,.
我們就將這一切榮耀歸還給你..
奉主耶穌基督的聖名,.
阿們..
\newpage



\section{約翰福音 3:5-15}
\label{sec:031706}
\textbf{耶穌與尼哥底母談重生}~~~~駱柱成牧師
\newline
\newline
連結: \href{http://www.church.com.hk/acms/content.asp?site=cdc\&op=show\&type=product\&code=031706\&layout=sermon}{\texttt{031706}} ~~~~ 語音日期: 2026-02-22
\newline
\newline
\hyperref[sec:031705]{\small{< < < PREV SERMON < < <}}
~
\hyperref[sec:index_chronic]{\small{[返順時目]}}
~
\hyperref[sec:031707]{\small{> > > NEXT SERMON > > >}}
\newline
\newline
約翰福音 3:5-15
\newline
\begin{longtable}{cl}
\hline
\hline
章節 & 經文 (和合本修訂版)\\
\hline
3:5 & \begin{tabularx}{0.7\textwidth}{X} 耶穌回答：「我實實在在地告訴你，人若不是從水和聖靈生的，就不能進神的國。 \end{tabularx} \\ \\ \relax
3:6 & \begin{tabularx}{0.7\textwidth}{X} 從肉身生的就是肉身；從靈生的就是靈。 \end{tabularx} \\ \\ \relax
3:7 & \begin{tabularx}{0.7\textwidth}{X} 我說『你們必須重生』，你不要驚訝。 \end{tabularx} \\ \\ \relax
3:8 & \begin{tabularx}{0.7\textwidth}{X} 風隨著意思吹，你聽見風的聲音，卻不知道是從哪裡來，往哪裡去；凡從聖靈生的也是如此。」 \end{tabularx} \\ \\ \relax
3:9 & \begin{tabularx}{0.7\textwidth}{X} 尼哥德慕問他：「怎麼能有這些事呢？」 \end{tabularx} \\ \\ \relax
3:10 & \begin{tabularx}{0.7\textwidth}{X} 耶穌回答，對他說：「你是以色列人的老師，還不明白這些事嗎？ \end{tabularx} \\ \\ \relax
3:11 & \begin{tabularx}{0.7\textwidth}{X} 我實實在在地告訴你，我們所說的是我們知道的，我們所見證的是我們見過的，你們卻不領受我們的見證。 \end{tabularx} \\ \\ \relax
3:12 & \begin{tabularx}{0.7\textwidth}{X} 我對你們說地上的事，你們尚且不信，若對你們說天上的事，如何能信呢？ \end{tabularx} \\ \\ \relax
3:13 & \begin{tabularx}{0.7\textwidth}{X} 除了從天降下的人子，沒有人升過天。 \end{tabularx} \\ \\ \relax
3:14 & \begin{tabularx}{0.7\textwidth}{X} 摩西在曠野怎樣舉蛇，人子也必須照樣被舉起來， \end{tabularx} \\ \\ \relax
3:15 & \begin{tabularx}{0.7\textwidth}{X} 要使一切信他的人都得永生。 \end{tabularx} \\ \\
[1ex]
\hline
\hline
\end{longtable}
$^{1}$大家住內平安.
今日我講的仍然是上一次講的.
耶穌和彌哥帝母譚松生.
上週 週日我們有個聚餐.
剛好是年廿八洗垃圾.
今天已經進入新的一年.
所以年廿八是象徵著除舊迎新的.
今天已經是新的一年.
在此祝大家新年龍福.
今天是農曆年的初六.
如果你看到 過年如果是華人.
他們有一些會在門口貼著「五福臨門」.
五福是甚麼五福呢.
五福是講到壽 富 康寧 浩德和善終.
甚麼意思呢.
壽就是要活得長久.
富就是要過得好.
康寧不只是身體要好.
你的心情也要好.
你的心靈是有安寧的.
浩德除了壽 富和自己好.
你還要做好事.
最後才好好離開就是善終.
所以一句話就是.
盼望活得久 活得好 活得平安.
做人有得行 最後有得安穩.
在這五個當中最大的福是甚麼呢.
就是壽.
這就是五福之壽 講到壽.
今天我是講重生 從重生講到永生.
聖經所說的永生遠遠超過五福當中的長壽.
不單止活得久.
而且是得著真正從神而來的永遠的生命.
黎哥弟母見耶穌的時候.
黎哥弟母不是一個普通人.
她的年紀比耶穌更大很多.
是以色列人的老師.
可以說她是一個宗教學術的權威.
是法利賽人的領袖.
是耶路撒冷有名望的一個人.

$^{41}$但她卻很謙卑.
來到這個比她更年輕直直無名的耶穌面前.
所以你會看到在她身上.
黎哥弟母是很謙卑的.
她不是做那些客套的事.
她為甚麼來找耶穌呢.
就看到耶穌生命當中.
是有一個與眾不同的一件事.
這件事是她沒有的.
她就被耶穌吸引來見耶穌.
她在那種耶穌和法利賽人爭拗中.
他們的關係已經很緊張.
她就頂著這種壓力來見耶穌.
而且是晚上去見耶穌的.
為了避免很多的不必要的爭議.
耶穌知道這可能是黎哥弟母只有一次的機會.
所以祂沒有跟她說太囉嗦的事.
直接指向那個核心.
所以祂就說.
我實實在在告訴你.
人若不從生就不能見神的國.
重點有兩個.
一個就是從生.
一個是神的國.
不是那些知識的問題.
是生命上的問題.
我不是說知識不重要.
我們要有知識.
更重要的是生命.
從生是很重要的.
所以耶穌說你們必須要從生.
不單止是你要從生.
你們的時候就說到.
不單止是黎哥弟母.
而且是說到整個猶太宗教文化體系.
你們這一群人.
你們這一群法利賽人.
你們這一群教導的人.
你們需要從生.
你們以為你們是明白聖經的.

$^{81}$其實你們只是將聖經當作一種學術.
而且有了學術之後.
你們在這裡很驕傲.
成為一個驕傲的本錢.
其實他們真正要做的是要有生命.
他們需要有這個生命.
他們猶太人本來是一個追求真理的民族.
但是現在發展成形式化.
虛假的宗教體系.
他們是很虛假的.
他們從律法傳統當中.
沒有走出來.
反而是以那些律法那些規條.
成為自己驕傲的本錢.
當他們這樣做的時候.
他們就離那個真理越來越遠了.
最後走到的就是一條絕路.
一條死路.
因為正如聖經所說的.
字句是叫人死的.
正義是叫人活的.
他們就是守著那些字句.
所以耶穌就說你們必須要從生.
直接將他們那一套本質的東西說穿了.
就是將以色列人一千多年.
一直積累的那樣東西.
自以為屬靈的那種宗教成就.
全部瓦解了.
為什麼?.
因為你們必須要從生.
但是你們只是停留在那個知識那裡.
所以你們整個宗教文化.
缺陷是什麼?.
就是從生.
你們.
耶穌不是要幫他們修修補補.
或者調整他們修行的方式.
而是說你們這個舊的生命.
需要從生.
要重新開始.

$^{121}$不是幫你們再修修補補這個生命.
是要重新開始.
所以唯一的出路就是從生.
而且他們所想的那個神的國.
是和耶穌所想的不是那麼一致的.
他們說他們那種神的國是.
一種很狹窄的民族主義.
是憑著他們的血統而驕傲的.
但是耶穌所說的.
是你們不要憑著那個血統.
因為憑著那個血統.
你們是不可以進入神的國的.
而且聖父一早就對聖子說了.
我要將列國賜給你為基業.
耶穌基督他的國是用他自己的血.
買贖各國各族各方的人.
進入他的國的.
是萬民的.
不是那種狹窄的民族主義的.
是屬另一個國度.
不是憑著你的血統進入這個國度的.
是憑著那個生命.
是從上頭而來的那個生命.
你需要重生才有這個生命.
重生不是新約自由的.
舊約一早就說了.
我要除掉你的石心.
賜給你一個玉心.
這個就是一個重生.
所以舊約時代的人.
他怎樣去明白聖經呢?.
他是只有那些真正被聖靈帶領的.
他才看到神話語裡面真正那種活的意思.
那個正義.
所以耶穌就說了.
你們必須要重生.
耶穌不是要建立一個新的宗教.
而是揭開一個.
揭開了猶太人他們一直以為.
我們就是站在那個真理裡面.

$^{161}$但是他們最根本的真理就是重生的真理.
他們都忽略了.
他們都沒有捉住.
猶太人他們一直捉住的.
就是以為靠著行為.
靠著律法就可以得救了.
其實如果你讀創世記的時候你就知道.
人類的問題早在亞當的時候就開始了.
因為當亞當吃了那個分別善惡果的時候.
他們就困在那個善惡當中.
他們就用人自己的標準.
那個善惡標準.
努力去行善.
避免作惡.
希望可以更好地活下去.
但是無論你自己怎樣做.
多麼善惡.
結局就是死.
所以猶太人他們一直在那個分辨善惡樹那裡兜圈.
其實其他宗教亦都是.
其他宗教亦都是一直在用善惡的觀點.
用做好人做好事的觀點.
用則公德的方法去使自己有修為.
去達到這個免除死的災害.
但是始終他們都離不開那個善惡的兜圈.
其實真正能夠解決生命的問題的.
是耶穌基督.
只有耶穌基督來到這個地上的時候.
才叫人得生命.
人才可以真正有重生的生命.
不是繼續在那個分辨善惡樹那裡.
而是來到耶穌基督這棵生命樹那裡.
因為耶穌基督自己亦都說了.
我來就是叫人得生命的.
你們不應該一直在善惡那裡兜圈.
你們應該是因信清而因信得生.
這個真理是一早就講了.
耶穌基督不是只給標準給人.
而是把生命賜給人.
讓人可以去活出來.

$^{201}$所以耶穌對黎哥弟母說.
如果人沒有重生的話.
根本上不可能明白這些屬靈的事情.
你根本上就沒有真正明白聖經的意思.
亦都不能夠領會聖經裡面的精義是什麼.
因為聖經自己講了.
屬血氣的人是不可能領會神屬靈的事的.
所以若不是從水和聖靈生的.
你就沒有那個重生.
重生不是你所想的.
就再回到母親那裡再生一次.
不是這樣.
只有聖靈來到的時候.
在人生命當中去做潔淨的工作.
將一個新生命賜給人的時候.
你才可以得到這個生命.
你才可以真正明白.
你才可以真正將聖經的精義活出來.
所以從肉生生的就是肉生.
就是在亞當裡面.
從肉生生的就是肉生.
亞當是全人類的始祖.
代表著全人類.
亞當是有神的形象的.
這個形象傳給後代.
所以全人類都落在這個罪惡當中.
因為亞當犯了罪.
當亞當犯了罪的時候.
他不單止是在地位上墮落了.
他有那個罪.
而且他的人性的本性.
人的人性都受到影響了.
因為他墮落了.
整個人都受影響了.
他犯了罪.
罪是從亞當一人進入世界的.
死又是從罪裡的.
因為人犯了罪.
死就臨到了.
因為聖經一早就說了.

$^{241}$罪的功架就是死.
因為人犯了罪.
所以罪進入了這個世界.
死就臨到眾人了.
說到死的時候有三個層次.
第一個層次就是當人犯了罪的時候.
神和人就隔絕了.
這是靈性的死.
第二個死呢.
就是我們這個肉身有一天會扭壞了.
我們在地上都要面臨死亡.
還有一次死呢.
就是永遠與神隔絕了.
是永遠的.
就是將來的另外一次的死.
所以肉身生的人.
他有兩次的死.
一次的生就是肉身生.
媽媽生了我出來.
但是死就要死兩次.
這個肉身會扭壞.
以後還有一次要死.
人一直在亞當的生命的源流.
在亞當的源流.
亞當這條線一直在打轉.
就算他們怎樣努力也好.
他們都沒有辦法走出這個源流.
所以他們就面臨的.
無法解決的就是死亡的結局了.
從聖靈生的就在基督的裡面.
因為從靈生的就是靈.
因一人的信服.
因著耶穌基督的信服.
眾人就成為兒.
因著耶穌基督的信服.
我們因信稱義就成為神的兒女.
耶穌基督將他的兒披在我們身上.
我們就成為一個兒人.
所以我們的地位就改變了.
我們是神的兒女.

$^{281}$本質上我們這個生命.
就是一個新造的人.
所以有些基督徒就說.
我是亞當的後代.
本質就是罪人.
所以我們必然是犯罪的.
但是這裡說.
從靈生的.
你是一個新的創造.
是一個新的生命.
是整個人的本性已經改變了.
當然這個肉身依然會有壞.
所以在基督的裡面.
我們是有新生命的.
我們不要犯罪.
我們是可以有不犯罪的自由的.
我們是可以討神喜悅的.
我在訓學院有一個同學.
他未信主之前.
是那種很惹是生非的人.
專門跟人鬥的.
而且他做甚麼呢.
他自己開酒吧.
就走過天門的生意.
信了主之後.
他重新得救之後.
他就選擇離開那種黑暗的生活.
他的作息變了.
不再做夜貓.
性情變了.
面貌也變了.
我在他的錢包裡看到一幅照片.
就是他以前的樣子.
他以前的樣子就像一匹狼一樣.
那些眉毛是向上揚起的.
每天都跟人鬥.
臉是凶神惡煞的.
有狼性在.
但信了主之後完全不同了.
性情就好像一匹馬.

$^{321}$他很善.
不是馬善地人其.
他是願意為主的緣故去服侍人.
所以整個生命是改變了的.
因為他的生命的源流變了.
他從原先的亞當的源流.
進入基督的源流.
所以他有那個自由.
可以選擇去行善.
他有這個行善的力量.
他有這個生命.
有基督的生命在裡面.
所以一個基督徒.
他是生兩次死一次.
甚麼意思呢.
就是從父母生的.
他有這個育身.
第二個就是聖靈再生一次.
這個就是兩次生.
一般來說是有一次死的.
除非主耶穌基督來.
你就免了死.
你就只有兩次生.
就一直是永生.
所以基督徒是兩個生人.
兩次的生一次的死一般來說.
重生就是聖靈.
從亞當這個生命的源流.
帶進去基督這個生命的源流.
重生不是靠血統可以傳遞的.
很多父母.
他將最好的給他們的小孩.
最好的學校.
最好的生活.
最好的一些預備和資源.
給了兒子.
他們不單想兒子有個健康的身體.
但他們更盼望自己的子女.
自己的孫子孫女.
有個屬靈的生命.

$^{361}$但聖經這裡講清楚.
從肉身生的就是肉身.
從靈生的就是靈.
亞當的肉身生命傳給後代.
他們有個罪性.
但基督徒屬靈的生命.
不是靠著血統.
靠著父母的血源傳下來的.
屬靈的生命是需要你的子女.
去從神建立這個關係的.
不是說你自己生出來.
他們就自動.
就算是基督徒的父母也好.
他生出來就變成一個基督徒.
原則在哪裡.
原則就是從亞當生命生的.
那個就是肉身的生命.
從基督.
從聖靈生的.
那個是屬靈的生命.
所以從肉身生的就是肉身.
從靈生的就是靈.
所以一個是在肉身那裡.
每一個人都要來到基督這個源流那裡.
你才得到這個重生.
你才有這個屬靈的生命.
不是因為你的父母是基督徒.
你就得到這個生命.
人自己犯了罪.
按照公義來說.
他是要負責的.
他自己犯罪的功架就是死.
所以如果上帝不救我們的時候.
他沒有違背這個公義.
因為人自己犯了罪.
人就要為自己的罪去負責.
所以就算上帝不救我們也好.
他依然是公義的.
但是這位他有救恩.
他有恩典.

$^{401}$他有憐憫.
他有慈愛.
所以他按照自己的主權去使人得到重生.
得到新的生命.
耶穌自己說風隨著意志去吹.
你們不知道它從哪裡來往哪裡去.
凡是從聖靈生的也是如此.
當說到這個風的時候.
這個是相關的.
第一就是指到這個物質的風.
這個物理的一個現象.
第二個是另外一個意思.
就是指到靈和氣息.
如果你再看《創世記》的時候.
凡是將那個氣息那個生氣吹入人的鼻孔的時候.
他就成了一個有靈的活人.
意思就是這個生命就是從神的氣息來的.
是由神的吹氣得著這個屬靈的生命.
人活著是靠著一口氣.
氣一斷人就死了.
無論你的地位是多高也好.
你的職位是多好也好.
你多有錢也好.
都要面臨死亡這個問題.
風是隨著自己的意思吹.
不是隨著人的意思去吹.
風是按著上帝的意思去吹的.
因為風有兩個意思.
靈是隨著自己的意思去吹的.
當聖靈要拯救人的時候.
祂是感動一個人.
帶領一個人.
祂有自己的計劃.
不是按著人的計劃.
人的意思不可能操縱風.
也不可能操縱這個重生.
所以重生是由聖靈去按著自己的意思去做的.
聖靈按著聖靈從自己的工作去行的.
保羅他從來都沒有想過要信耶穌基督.
當他正去大馬抓基督徒的時候.

$^{441}$他沒有想過要信耶穌基督.
他反而是直黨耶穌基督.
所以保羅是沒有按著自己的計劃去信耶穌的.
他也沒有預備接受這個救恩.
什麼都沒有.
但是神的恩典就臨到他了.
所以重生不是人自己決定的.
是聖靈去決定的.
風隨著自己的意思去吹.
它從哪裡來我們不知道.
往哪裡去也不知道.
但是當風一吹的時候.
風吹到保羅的時候.
天上有大光一照保羅.
他就倒下了.
他沒有辦法去選擇.
所以不是人去揀選神.
是神去揀選人.
是完全出於神的恩典.
人揀神很可能會揀錯.
我記得我在廣州的時候.
我隔壁屋是做很大的一個水果檔.
我經常喉嚨痛就問他.
怎樣去揀選那些鴨梨.
怎樣去揀選雪梨.
他就說皮光肉滑.
我就按著皮光肉滑.
因為皮很光鮮的.
看上去是很順的.
肉質一定是比較好的.
他就將這個經驗告訴我.
當他教我去揀的時候.
揀那些皮光肉滑的.
大多數的時候都是很甜.
很潤的.
但是有時也會揀到那個黑心的.
因為皮光肉滑.
但是裡面已經變了質.
所以人揀水果都會揀錯的時候.
何況是一個墮落的人.

$^{481}$揀一個至高的神.
至聖的神.
你有沒有這個保障去揀得對呢.
所以不是人揀神.
是神揀人.
我在1981年出生.
我在1981年之前.
我從來沒有打過電話給我媽媽.
經他們找她說.
「你可以生我了」.
「你自己可以了」.
「我這裡也可以了」.
「不是兩個人商量好」.
這個是不可能的.
我媽媽說生我就生我.
同樣重生是怎樣的.
重生不是人去發起的.
是從天那裡降下來.
聖靈賜給人一個新的生命.
這個是重生.
不是人自己去決定的.
人會不會找神.
人亦都不會找神.
因為聖經一早就說了.
沒有明白的亦都沒有尋求神的.
人得救不是偶然的.
你想想我們得救的時候.
我們真的很多時候.
不知為何剛剛那段時間.
就是有基督徒向我傳福音.
剛剛那時候我就有時間去教會.
剛剛那時候就是那些訊息影響我的思想.
所以我們很多事情不知道.
你不知道.
我不知道.
可能身邊的人都不知道.
唯有神一早就知道了.
出於祂的簡淺.
出於祂的慈愛.
出於祂自己的主權.

$^{521}$因為神自己說了.
「我要憐憫誰就憐憫誰」.
「我要恩代誰就是恩代誰」.
聖靈不是跟著人去走.
聖靈是上帝的風.
祂有自己的時間.
自己的計劃.
自己的主權.
重生是神主動的.
當耶穌跟尼哥底姆說這番話的時候.
祂不是零散的.
祂是有一條很清晰的宣告.
祂說「人若不重生」.
「從肉身生的就是肉身」.
「從靈生的就是靈」.
「凡從聖靈生的亦都是如此」.
所以耶穌基督一直說一件事.
就是生命.
耶穌不是跟祂辯論.
尼哥底姆很有知識.
你是學術界的太斗.
我現在跟你說的不是這些知識的問題.
當然耶穌基督的知識比祂更高.
祂跟祂直接說的是生命的問題.
尼哥底姆代表著宗教知識.
亦都是代表著權威.
耶穌基督代表著從天而來的啟示.
神的啟示遠遠比人的理性知識要高.
代表著屬靈的生命.
所以重生.
屬靈的生命關鍵不是你知識有多少.
耶穌基督在這裡一直要提的就是.
你要重生.
你要有這個生命.
尼哥底姆聽了之後很震驚.
他是最高層次的知識學者.
他的知識架構很龐大.
他舊有的觀念很根深蒂固.
但是對重生這件事.
他可以說差不多完全不可以明白.

$^{561}$所以他就說.
怎麼能有這些事呢.
這件事怎麼可能呢.
耶穌就回答.
你是以色列人的先生.
還不明白這些事嗎.
意思就是.
你已經坐在摩西的位置上了.
你是整個以色列民族.
在你的教導之下.
你是應該明白的.
但是你自己不明白最重要的真理.
你本來應該是教導屬靈的真理的知識.
但是你最根本最基礎.
最重要的東西你都未捉得住.
猶太人他研究舊約律法一千多年.
他們的學識很淵博.
塔木德是他們舊約的注釋書.
可以說幾馬車你都拉不完.
但是他們沒有捉住這個核心.
你這麼有知識又怎樣.
但是你的知識沒有捉住.
真正這個信仰的核心是什麼.
所以我們今天的人.
很多時候我們已經是學富五車了.
今年是馬年.
可能你再多讀一馬車的書.
但是如果你沒有基督這個生命的時候.
再多的知識.
在神這個智慧這個生命的面前.
你是沒有價值的.
你是對真理的認識一竅不通的.
所以一個神學生.
如果你是只有幾個神學位的時候.
神學學位的時候是不夠的.
因為最基礎的就是基督的生命.
黎哥帝沒有他不明白.
所以一個人如果沒有基督的生命.
沒有聖靈的工作.
是不可能真正明白聖經的.

$^{601}$因為哥林書說了.
宿怯的人不能夠領會神.
他聖靈的事反倒為愚拙.
於是耶穌就跟他說.
你不明白.
耶穌跟他說.
我所說的是我們知道的.
我所見證的是我見過的.
你們卻不來到我這裡領受.
我們是指誰呢.
我們不單止是指耶穌祂自己.
也是指將來在教會中為主作見證的人.
包括那些為耶穌基督作見證的人.
傳福音的人.
他們所傳的不是人的思想.
這個福音是從前就隱藏的.
如今就顯明出來.
所以我們是包括我們今天傳福音的.
但你們卻不能夠領受.
猶太人他們被那種知識.
那種宗教所綑綁.
他們不能夠領受.
他們有舊約的啟示.
但本來舊約是指向耶穌基督的.
但他們卻拒絕舊約這個指向.
不來到耶穌基督的面前.
你們研究了很多的聖經.
但卻不肯來到.
聖經所指向的這一位耶穌基督的面前.
所以你再看後面的《約翰福音》時.
你就知道了.
你們查考聖經.
因為你們以為內中有永生.
給我作見證的就是這本聖經.
而你們不肯來到我的面前.
去領受這個生命.
聖經就好像一個證人一樣.
耶穌基督只是被見證過的那一位.
要你來到耶穌基督的面前.
你才可以得到這個生命.

$^{641}$但你們不願意來到我的面前.
你們不願意去領受.
人如果不來到耶穌基督的面前.
就得不到這個生命.
因為耶穌基督自己也說了.
祂是叫人得生命.
而且得得更豐盛的.
神一切的豐盛都有形有體.
在耶穌基督的裡面.
基督是道成了肉身.
祂本來就是真理的本體.
祂自己就是那個真理.
祂是神向人顯明自己的.
祂將這個真理顯明.
祂是從天降下來的.
將這個真理顯明給人.
伊斯蘭教最大的問題是什麼?.
他們拒絕了三位一體.
拒絕了道成肉身.
他們拒絕了十字架.
猶太人也拒絕耶穌基督.
所以他們同樣一個問題.
就是拒絕耶穌基督.
所以你和伊斯蘭教和猶太人說福音.
很難 他們聽不進去.
他們早就拒絕了.
為什麼?.
因為他們有自己的框架在.
你和他們說福音的時候.
他們當你是耳邊風.
他們認為你的信仰是不合理性的.
獨一的神怎會是三位.
神怎會是道成肉身.
所以靠著袋鼠這種的真理.
他們也是不接受的.
所以他們的系統.
他們的思維是封閉的.
他們就好像掩著耳一樣.
人是沒有辦法去打通.
打開他們的耳 打開他們的思想.

$^{681}$唯有聖靈來 才可以打開.
耶穌說.
「我對你們說地上的事.
你尚且不信.
若說天上的事 如何能信呢?」.
黎哥弟母說的是宗教的那種理性.
但是耶穌說的是天上的事.
不是地上的事.
祂說的是屬靈生命.
如果舊的框架不打破的話.
這種舊的思維不打開的話.
人是很難相信的.
天上的事.
不只是用頭腦那種知去領受的.
是要用信心去領受的.
它不只是那種學術的研究.
不是理性的推測.
而是聖靈啟示的.
聖經的說法是.
只有被聖靈光照.
被聖靈重生的人.
才可以真正去明白那個正義.
信不是完全明白了之後.
你才去信.
很多人說.
「等我明白之後 我就信」.
這句說法是不對的.
人通常是先自己做了.
決定要做什麼的時候.
再用那個理性去將它合理化的.
比喻一個糖尿病的人.
他除了自己身體的軟弱.
很想吃糖之外.
他就是被這種想吃糖的慾望.
去驅使他自己.
他就說「我就是要吃甜的」.
「我要找各樣的理由」.
「去支持我吃這個甜的」.
所以人通常的理性.
那種邏輯是被那種慾望去驅使的.

$^{721}$希伯來書一早就告訴我們.
我們因著信就知道.
所以不是先知道.
是要先信.
因信你才知道.
你需要先信你才可以知道.
爾賽納書一早就說了.
你們如果不信的話.
就不能明白.
所以是先信才可以明白.
在中世紀.
一個很偉大的神學家安息律說了.
信心尋求理解.
你要用信心回應神.
祂的話 祂的啟示.
在這個信心裡面你不斷地去追求.
去認識.
越信你就越知.
越知你就越信.
這個就是因信求知.
你要先信你才可以明白.
信和那種知不是對立的.
它們是彼此推進的.
就好像那個樹幹.
你信得.
你越信你就越知.
越知就越信.
所以是一圈一圈不斷地去增長.
你的信心也是這樣.
亞伯拉罕是因著信.
被夢召的時候.
他就專行出去了.
他要去哪裡.
他要往將來要得遺孽的地方去.
出去的時候他都不知道要往哪裡去.
他不知道.
他不是說我要找一部手機.
要有導航.
或者找一幅地圖.
我才做這個決定.

$^{761}$他是因信稱義的.
所以亞伯拉罕是傳著信心出發.
當他死的時候.
他都沒有得著所應許的.
但是他就遠遠地望見.
將來一個更美的家鄉.
所以他當自己是客裡的.
他就看到那個天上的城.
所以信心是讓人看到那個理性.
沒有辦法看得到的東西.
所以亞伯拉罕他的經歷告訴我們.
是因信稱義.
是你因著信你就知道了.
用信心的眼睛你就可以看得到.
靠著人那個被造的理性.
有限的理性和墮落的理性.
如果只是靠著我們這個墮落的人.
我們都墮落了.
怎可以明白天上的事呢.
經文講到天上的事.
那麼高.
但是我們人已經墮落了.
我們只可以靠著我們的理性.
沒有辦法去窮究得到天上的事.
所以只可以用信.
信就是你的理性的一種突破.
是一種飛躍.
理性和信心就好像一個是井底之蛙.
一個是在井口上的青蛙.
如果是在井底的那個.
它只可以看到井入邊的四圍.
它只可以測到井口是多大.
它可以通過推理分析和估計.
做一個結論.
在井口上的青蛙.
它看到的東西是不一樣的.
它說那個天不是你在地墮落了之後.
好像那個井底之蛙一樣.
看到是那麼狹窄的.
天是無邊無際的.

$^{801}$是比這個井口廣闊很多倍的.
人算得了什麼.
所以它所看到的.
和井底看到的是不一樣的.
信心就是一個突破.
當聖靈啟示人的時候.
信心就突破了這個井口一樣.
將人帶到那個井口之上.
看得更遠.
帶領人的心思.
帶領人更簡潔.
看得更寬.
引領人到高處那裡.
向天上去看.
我知道我們有些弟兄姊妹.
她不一定她的丈夫信.
她的丈夫有一個姊妹.
她的丈夫是做高管的.
技術人員.
從小就讀基督教的學校.
也讀過很多次不少的聖經.
但到現在都未信主.
她不反對姊妹過來聚會.
但她自己不信.
她先生原有的知識告訴她.
我這個理性是合理的.
我這一套是最好的.
你們那一套向我傳福音.
根本上不可能經得過我這個理性的頭腦.
所以向知識份子傳福音是很難的.
你要更多的耐心.
你要等候聖靈的工作.
因為他們固有的那種思想.
那種理念就好像牛皮燈籠一樣.
怎麼點都不明白.
人很難將它點明.
唯有聖靈來的時候.
才可以將它們點明.
將它們的思維打開.
因為福音大三章沒有告訴我們.

$^{841}$黎哥弟弟母有決志.
有那個回應.
但後來我們知道.
她和亞利馬太預備香料.
按照猶太人的禮儀去安葬耶穌.
這個.
所以我們會看到黎哥弟母.
她當時不一定馬上相信.
但後來她相信了.
因為當時她的理性太強.
她的宗教框架太深了.
但神沒有忘記她.
做聖靈工作的時候.
她就相信了.
後來她相信了.
因著主的緣故被排擠.
失去了她自己原有的地位.
失去了收入.
晚年為主受苦.
生活是靠執柴的.
那一夜的談話對黎哥弟母來說是值得的.
因為這一次的談話.
她就得著永遠的生命.
是很有價值的.
風隨著自己的意思吹.
聖靈隨著自己的意思去做.
所以神的工作.
當神做工的時候.
人就相信了.
神有祂的計劃和時間.
我們可以做什麼呢?.
我們傳福音只可以說我們要盡忠就行了.
那一刻只是靠神.
天上的事從哪裡來.
誰可以告訴我們.
一直說天上的事.
一直說地上的事.
地上就是這麼理性.
這麼感性的東西.
天上的事.

$^{881}$耶穌基督說.
除了從天上降下.
認救在天的人子.
沒有人升過天.
沒有人可以將天上的事告訴我們.
唯有耶穌基督除了從天降下.
認救在天.
天上的事只能夠從天上.
這一位耶穌降下來.
不是那些人升上去知道天上的事.
是耶穌自己親自降下來告訴我們.
耶穌基督在哪裡降下來.
耶穌基督是至高神的兒子.
祂從至高處來到我們當中.
祂成為人.
但祂成為人的時候.
祂也是神.
祂沒有失去祂的神性.
所以耶穌基督從天上降下來.
取了人的樣式.
有了人性.
但祂的神性依然是無限的.
祂在地上的時候有人性.
但祂依然是神.
祂的寶座在天上.
祂依然掌管著一切.
祂來到地上成為人.
祂依然坐在天上.
什麼意思?.
就是除了從天降下認救在天的人子.
祂來到地上.
祂依然坐在寶座上.
因為祂的神性是無限的.
祂用祂的權能托住萬有.
所以祂依然是掌管著天地.
祂的寶座依然在天上.
耶穌基督在地上.
祂是永恆之道,祂成了肉身.
祂降下降了下來.
但祂依然在天上.

$^{921}$這就是耶穌基督的獨特性.
祂的神性是無限的.
祂受人性的限制.
但祂依然在天上.
在地上的時候祂依然在天上.
耶穌基督說沒有人升過天.
但聖經一早就說了.
以洛升了天.
以利亞也升了天.
為什麼說沒有人升天.
這個天不是指普通的天空.
也不只是指一個屬靈的境界.
天,保羅說有三重天.
所羅門說天和天上的天.
天就是我們看到的大氣層.
一個天.
第二層的天就是宇宙.
第三層的天,保羅說有三層天.
就是神祂自己最高的寶座.
耶穌基督就是在最高的寶座降下來.
所以耶穌基督的寶座.
最高之處就是指神本體榮耀的所在.
所以說沒有人升過天.
我們人都是向下而升的.
但耶穌基督是降下來的.
沒有人憑著自己的意思.
可以到達神的最高處.
以利亞,以洛被接走的時候.
不等於他們去到最高之處.
神給他們一個合適的地方安置他們.
當耶穌基督的救恩未完成的時候.
沒有人可以真正進入到.
父那裡,最高那裡.
沒有人可以.
因為救恩未完成.
人都是有罪的.
沒有辦法到達最高神那裡.
唯有一位從最高之處降下來的時候.
完成這個救恩的時候.
才可以去到父神那裡.

$^{961}$我從父那裡來到這個世界.
又離開這個世界回到父那裡.
所以從最高之處降下來的時候.
成就了舊約律法預言的救恩的時候.
耶穌基督完成這個救恩的時候.
人才可以去到父神那裡.
救恩完成的時候.
耶穌基督說了一句.
救恩完成之後.
耶穌基督在救恩未完成之前就說過.
我在哪裡,你們也要在哪裡.
也就是說當救恩完成的時候.
祂要將人帶到神存在的地位.
要將這個永恆的生命賜給人.
唯有耶穌基督從天而降來到這個地上.
完成這個救恩.
人才可以真正去到父神那裡.
以利亞,以洛他們都不可以.
直到等到耶穌基督完成救恩的時候.
今天華人的文化當中.
你會看到現在很多的電視劇是講修仙的.
修仙就是講修行.
就是一級一級去上升.
講積那個功德的.
講練那個法術的.
從練氣,捉姬.
還有結丹,玄嬰.
一級一級直到化神.
他們的觀念是這樣的.
所以當他們的能力.
修行越高的時候.
他們的能力就越高.
他們的壽命也都越長.
所以這個是從地上那個人的思維.
那個思想.
幻想到升上天上的.
這個是宗教.
但是福音不是人從往上修行.
而去到達神那裡.
是神自己降下來.

$^{1001}$耶穌基督從最高處降下來.
為要將人帶入永恆的生命.
之後耶穌基督降下來的時候.
將天上的生命賜給人.
人才可以得獎.
耶穌基督被舉起在十字架上.
不是為了建立宗教.
不是再次為了建立那個宗教.
祂是要賜給人生命.
為了黎哥弟母更好地明白十字架的救恩.
耶穌基督用黎哥弟母熟悉的詩去講.
前面講到風.
現在祂講到摩西在曠野舉銅石.
那些被火石咬的快要死的人.
什麼都沒有做.
祂沒有那個行為.
也都沒有功勞.
因為祂沒有將自己的血或者靠著人.
去吸那些毒或者對人不好的東西.
吸出來.
也都沒有打過血清.
什麼都沒有做.
但是他們這群被石咬的人.
只是做了一件事.
就是仰望那個銅石.
一望就得到這個生命了.
這個不是行為.
這個是信.
所以得到這個生命.
不是靠你的行為.
是靠著信.
不是靠你有什麼功勞.
是靠著信服.
石.
舉起銅石.
我們讀石的時候.
一般看到石的時候.
就會想到魔鬼.
但聖經也都講到靈巧像石.
不是叫人像魔鬼.

$^{1041}$所以不是所有的石都是指著魔鬼.
「石」在這段經文的意思是.
被咒座的.
就是被咒座的一個記號.
聖經講明了.
凡掛在木頭上的都是受咒座的.
銅石.
如表耶穌基督.
耶穌基督要掛在木頭上.
是受咒座的.
耶穌基督本來是沒有罪的.
但是他為了我們罪的緣故.
擔當我們的罪.
願意被掛在上面.
被咒座.
成為這個咒座.
成為這個銅石.
所以耶穌基督用了這件事.
去和黎哥弟母說.
黎哥弟母很熟悉舊約.
她一說.
她可以明白得到.
摩西怎樣舉銅石.
人耳也照樣要被舉起.
就是說.
咒座要歸到她的身上.
耶穌基督是從神而來的.
他是沒有罪的.
他是從天上而來的.
但是他願意降低被咒座.
為了將福氣可以臨到人.
將死亡臨到他自己.
將永生的生命賜給信他的人.
這一切完全超過了黎哥弟母所想的.
所能夠明白的.
她明白的就是那種理性.
那種地上的事.
舊的律法知識.
和那種舊約最高.
他們研究舊約最高的一種成就.

$^{1081}$都沒有辦法達到耶穌基督所成就的.
這一種境界.
耶穌基督在這裡.
當這樣說的時候.
就等於建立了一個新的方向.
和新的標準.
你們以前是靠守傳統.
守律法.
靠著學問.
就以為可以得救的.
但是現在是靠著.
靠著仰望耶穌基督這個人子.
而得著這個生命.
得著真正的救恩.
你們才可以真正得救.
所以對黎哥弟母是一個新的方向.
新的一個標準.
所以在這裡.
耶穌其實對黎哥弟母說.
你是否真正願意.
除去舊的一套迎接這個本來就有.
你們以為新的一套.
舊約本來就隱藏了這個救恩.
你們願不願意除去這個舊的傳統.
離開這個舊的框架.
進入到我現在向你們這個啟示.
盛傳的舊約這個救恩.
這個從天上而來的啟示.
你們願不願意接受.
你們願不願意放下你們的一套.
黎哥弟母是眾多耶穌見過的領袖當中.
唯一一個可以信耶穌的.
最高的政治領導.
皮拉多他沒有信.
表面的王希律王也沒有信.
宗教當年的大祭司.
本年的大祭司戈埃法也沒有信.
唯有黎哥弟母他後來重申了.
他信了得著這個永遠的生命.
進入神的國.

$^{1121}$所以親愛的兄弟朋友.
我不知道你心目中有什麼.
是難阻你來到耶穌基督的面前.
讓你沒有辦法得著這種福氣.
如果你是願意得著這個贖天的福氣.
希望得著這個永永遠遠的生命.
遠遠比五福那種長壽更好的福氣.
唯有被聖靈重生.
在基督的裡面才可以得著這個生命.
因為聖經應許.
凡信耶穌是基督的.
是神的兒子的.
就可以得著生命.
願神賜福給你.
親愛的弟兄姊妹.
願神除去我們那些.
難阻我們明白真理.
專行真理那種固有的思想.
那種舊的模式.
求神真的幫我們打破.
願神真理的靈.
帶領我們更深入去認識耶穌基督.
得著這個豐盛的生命.
願智慧的靈引導我們.
在基督裡面得著贖天的智慧.
使我們能夠將身邊未信的.
親朋戚友.
特別是這段時間過年.
將他們的心意奪回.
使他們歸向主耶穌基督.
願神賜福給你們.
和你們的家人朋友.
願神祝福大家.
\newpage



\section{創世記 23:13}
\label{sec:031707}
\textbf{聖經人聖經事:人情世故}~~~~黃克勤牧師博士
\newline
\newline
連結: \href{http://www.church.com.hk/acms/content.asp?site=cdc\&op=show\&type=product\&code=031707\&layout=sermon}{\texttt{031707}} ~~~~ 語音日期: 2026-02-15
\newline
\newline
\hyperref[sec:031706]{\small{< < < PREV SERMON < < <}}
~
\hyperref[sec:index_chronic]{\small{[返順時目]}}
~
\hyperref[sec:031708]{\small{> > > NEXT SERMON > > >}}
\newline
\newline
創世記 23:13
\newline
\begin{longtable}{cl}
\hline
\hline
章節 & 經文 (和合本修訂版)\\
\hline
23:13 & \begin{tabularx}{0.7\textwidth}{X} 對以弗崙說，也給當地百姓聽：「你若應允，請你聽我。我要把田的價錢給你，請你收下，我就在那裡埋葬我的死人。」 \end{tabularx} \\ \\
[1ex]
\hline
\hline
\end{longtable}
$^{1}$大家好,又過了一天.
今天我們來看看《人情世故》.
這個題目很多時候我都會用聖經裡面的字眼.
不過今次《魏斷經文》有二十節.
是不是很害怕呢?.
很長,所以我選不到一節出來用.
我看到又是亞伯拉罕.
亞伯拉罕是一個很懂得《人情世故》的人.
很厲害,辦事很精明.
神祝福他,又給他一個很好的太太.
又給他一個很漂亮的媳婦.
祝福他很富有,怎樣怎樣.
其實你看到神的祝福不僅是給他地.
不僅是給他成為各國之父.
很多恩典.
不僅如此.
其實還給了亞伯拉罕一個恩典.
是他很有智慧,帶著家庭裡的家丁好好打.
四王戰勝五王,其中一個就是他.
他和四王一起去打仗,救回羅德.
打贏仗回來.
亞伯拉罕當日都很厲害.
又文又武.
還有他很敬畏神.
昨天說過了.
還有亞伯拉罕神祝福他.
他很懂得《人情世故》.
去到那裡,很多朋友都尊敬他.
今天字很密,因為我有二十節聖經.
這個故事很少人在教會說.
很少人特別說亞伯拉罕很厲害.
多數經文沒有說.
有時我特別和大家講一下.
亞伯拉罕對人情世故的掌握.
學一下.
很長的.
這個綠色是節數.
因為我沒有一節字.
擠在一起,字都細一點.
有二十三章.

$^{41}$記得有機會自己看看.
二十三章第一節.
撒拉享壽一百二十七歲.
撒拉即是亞伯拉罕太太.
是撒拉一生的歲數.
撒拉死在迦南地的基列阿巴.
就是希伯倫.
亞伯拉罕為她哀悼哭號.
基列阿巴是城.
阿巴之城.
阿巴是什麼意思?不太清楚.
基列阿巴,阿巴的城.
就是希伯倫,希伯倫是聯合.
亞伯拉罕為了和她一百二十七歲.
為了她一百年的太太哭了.
後來亞伯拉罕從死人面前起來.
因為她已經死了.
所以不再用她的名字.
這裡是當年的文化.
亞伯拉罕從撒拉的屍體旁邊起身.
從死人面前起來.
對當地人叫做黑人.
當地人的民族叫做黑人.
黑人叫做Kat.
為什麼叫做黑?中文不太清楚.
Ch讀Kat.
黑人.
黑人的意思是son of Kat.
就是黑,黑人.
其實黑人這個字是驚恐的意思.
這個民族叫驚恐.
不知道是不是他們經常很驚青.
黑人後來和以色列人打仗.
黑人叫驚恐.
跟當地人說.
我在你們中間是外人是寄居.
求你們在這裡給我一塊地.
好讓我埋葬我的死人.
使他不在我眼前.
我現在見到一條屍體真的很悲傷.

$^{81}$我要埋葬他.
我是一個外人是寄居.
可不可以給我一塊地.
第四節.
第一點就是.
我在你們中間是外人是寄居.
亞伯拉罕很厲害.
他會跟別人說.
各位我太太死了.
我想找個地方在你們這裡埋葬她.
不過我不是這裡本地人.
因此我要埋葬一個人.
我要問准本地人.
我要在這裡埋一塊地.
我想在這裡埋葬她.
表明她的心意.
不過她說.
我的身份是外人是寄居.
她的站位很清楚.
第一點我要說的.
學習.
就是很多時候.
我們去到某些場合.
某些情況底下.
說錯話.
做錯動作等等.
站錯位置.
到處都以為自己是老闆.
回家就是老闆.
一家之主你作主.
出來就未必是.
亞伯拉罕很懂得人情世故.
第一就是.
我的站位.
我只是一個外人是寄居.
雖然神告訴我.
這裡是給了我和我的後裔.
不過我不會說.
各位黑人大家留意.
雖然我來到這裡.

$^{121}$不過這個地方.
全片土地.
上帝已經賜給我.
屬於我.
將來也是這樣.
馬上被人打死.
我是一個寄居的.
是外人.
現在我太太死了.
我就想找個地方葬她.
很聰明.
我只是一個卑微的外人.
黑人回答亞伯拉罕.
我主請停.
這是一句禮貌的說話.
阿東就是主.
以色列人叫上帝不敢叫耶和華.
他們叫阿杜奈.
阿杜奈是主.
阿杜奈如果譯作英文字.
A字阿杜奈就要大佳.
這裡主是世佳.
即是我主.
尊敬的Sir.
尊敬的人物.
尊敬的先生.
請你聽.
你在我們中間是一位尊大的王子.
你在我們當中.
是像王子般受尊敬.
是一個很受尊敬的人.
只管在我們最好的墳地埋葬你的死人.
我們沒有一人不容你在他的墳地裡.
埋葬你的死人.
黑人就說.
亞伯拉罕Sir.
尊貴的亞伯拉罕先生.
你在我們當中是尊貴到像王子般.
沒有人會說你是外人.
你喜歡葬在哪裡就葬在哪裡.

$^{161}$你知不知道.
以往就算後來創世記說.
虐室死了都想什麼呢?.
你記得何時離開埃及回到祖家.
要帶我的屍體回去.
要葬在本地.
不想葬在那裡.
我不讓你葬在那裡.
因為你葬在那裡.
是自己人才能夠做的.
我們當你是自己人.
我們不會有人拒絕你在我們地埋葬死人.
第七節.
亞伯拉罕起來向那地的黑人下拜.
亞伯拉罕聽到馬上起身下拜.
多謝大家.
多謝對我的尊重.
多謝大家.
多識趣.
起身下拜.
對他們說.
如果你們有意叫我埋葬死人.
使他不在我眼前.
不知如何是好.
我要包住他.
等何時回家就拿走.
請聽我的話.
為我求所割的兒子以弗倫.
原來亞伯拉罕做事真好.
識人情世故.
一早已經打聽找到.
在以弗倫的地裡有一個洞.
很靚的.
用來做墳墓.
因為他以前是整個身體放進去.
然後化了一些骨合在一起再放進去.
他們的葬禮是這樣的.
他們的墳墓是這樣的.
所以山洞已經物色了好地方.
有時如果你們有意幫我的話.

$^{201}$我想求所割的兒子以弗倫.
我知道他有很好的洞.
其實亞伯拉罕沒有去過.
不過可能是打聽到那裡好地方.
有沒有去過不知道.
不過他不認識以弗倫.
以弗倫意思是象幼鹿.
當地人這個名字很有趣.
他像一隻幼鹿.
很活潑,很可愛的意思.
是小朋友的名字.
把田上的麥比拉洞給我.
第三他可以按足價賣給我.
做我在你們中間的墳地.
我不想你們說有塊地可以隨便用.
我是給足價錢的.
這是人情世故的第三.
屬還屬.
人情還人情.
跟著一句我們中國人叫什麼.
數目要分明.
不然真的沒有朋友做.
原來當時以弗倫正坐在黑人中間.
一班人跟他聊天.
黑人是眾子.
在很多黑人中間.
於是以弗倫在城門口出入的黑人面前.
對亞伯拉罕說.
到底是什麼?.
其實以往的官員,政府部門,城主,鎮長.
全部坐在城門口.
出入時看到誰.
一方面守住.
二方面在城門口做審判.
有什麼事爭執.
在那裡找到他們.
辦公室在城門口.
原來當時他們應該在城門口聚集.
亞伯拉罕跟他們聊天.
說我想找塊地.

$^{241}$大家一齊來開會.
原來以弗倫坐在那一班來開會的人當中.
以弗倫站起來說.
不如我主請聽.
尊貴的先生請你留意.
我送你這塊地.
連田間的洞也賜給你.
那塊田地當中的洞其實是山.
一塊梯田.
我送給你.
連同洞一起送.
在我同族的人面前都給你.
在我同族的人面前給.
大家可以作證知道.
可以埋葬你的死人.
亞伯拉罕很會做.
於是又在眾民面前向他下拜.
他很會人形世故,很會禮貌.
很多人沒有禮貌.
站都站不到,又不懂得站位,又不懂禮貌.
有些就喧賓奪主.
有些以為自己很厲害,很驕傲.
但是亞伯拉罕很謙卑,很有禮貌.
給足錢.
所以大家都很歡喜他.
他在以弗倫面前說.
你若英文請聽我的話.
我要把田地給你,求你收下.
我就在此埋葬.
這個「就」通常在聖經出現.
「就」是指future tense.
即將來我將會埋葬我的死人.
我要把田地給你,求你收下.
以弗倫回答亞伯拉罕說.
大家留意一下.
人就是這樣的.
尊貴的先生請聽.
值四百四百勒銀紙的一塊田.
在你我中間還算甚麼呢?.
只管埋葬你的死人吧.

$^{281}$這個以弗倫真會說話.
七尊敬的先生.
區區四百四百勒銀紙算甚麼?.
拿去吧.
不行不行,我一定要給.
不要緊,拿去吧.
我一定要給.
其實我不知道.
不過如果以當日來說.
四百四百勒銀紙.
一點都不少.
不過以弗倫很會說話.
尊敬的亞伯拉罕先生.
區區少少地.
四百萬不算甚麼.
四百四百勒銀紙.
在當時我認為不少錢.
這些在你我之間.
不要緊,拿去用吧.
為何說得這麼輕鬆,這麼闊綽.
因為亞伯拉罕說了幾次.
我一定給錢.
可能也不知道是否不值四百四百勒.
不過現在開大了數.
四百四百勒,可能即是四十四百勒.
搞定,現在四百勒.
亞伯拉罕照給.
亞伯拉罕聽從了以弗倫.
照他在黑人面前所說的話.
把買賣通用的銀紙.
平了四百四百勒及以弗倫.
平了四百四百勒.
即是說是真的量夠四百四百勒.
當時要來買賣用的銀紙給他.
一筆不少錢.
四百四百勒.
都真的要秤的.
銀.
一點也不少.
不過他這樣開大口.

$^{321}$亞伯拉罕既然又開大口.
好,我給足價.
你開個價來.
四百萬.
嘩,是不是.
照給.
這個是他識做的地方.
又照給.
於是麥比拉萬利前以弗倫的納法地.
及其中的洞及四位的樹木.
都定準歸約.
可能是他在黑人面前.
並城門出入的人面前買妥的.
這件事有個妥字.
有時我們人情世故.
算了,行行行,行行行.
其實不行的.
要認認真真.
君君真真.
給足錢.
人情還人情.
數目要分明.
於是他買妥的.
而且在眾人面前在城門口買妥.
即是說大家都知.
以後亞伯拉罕就將他妻子撒拉.
葬在加拿大萬利前的麥比拉田的洞裡.
這個萬利是什麼意思呢.
Mamre.
Mamre或者叫Mamre.
Mamre.
Mamre是力量與肥肉的意思.
就是後來的名叫希伯倫.
希伯倫.
C-H-A-T.
Cat Brown.
不是Hebron.
Cat Brown.
Cat Brown是聯合的意思.
於是就葬在萬利.

$^{361}$後來就叫希伯倫.
為什麼叫希伯倫呢.
因為就這樣說萬利.
以前的人都不知道.
所以寫到這個創世記實.
特別要加.
其實叫希伯倫.
後來叫希伯倫.
大家知道了.
這個就是亞伯拉罕做的事.
這件事.
從此那塊田和田間的洞.
就藉著黑人定准歸與亞伯拉罕作墳地.
定准.
亞伯拉罕作墳地可以嗎.
可以的.
任何亞伯拉罕的人作墳地都可以.
都買了這塊地.
定准.
好了.
結束了.
讀了二十二世紀.
這句是什麼呢.
這句很難唸.
我要把田價給你.
求你收下.
我在那裡埋葬了死人.
如果你英文聽我的話.
請聽我說.
很有禮貌.
給足田價.
是他的智慧.
特別是埋葬死人這麼敏感的事.
非常感恩.
在亞伯拉罕撞殺拉這件事上.
見到他人情世故很厲害.
這麼厲害都是上帝祝福他的.
所以你見到這個人這麼厲害.
你為他感恩.
上帝祝福他很厲害.

$^{401}$不要說他這麼厲害.
我沒有他這麼厲害.
於是生他的氣.
很多人都是這樣.
見人窮就多踩腳.
見人差就笑他.
我們去學習.
亞伯拉罕有寬闊的胸襟.
很有禮貌.
很君真.
很識人情世故.
大家都會歡喜你的.
我們一起禱告.
如果我們見到亞伯拉罕.
他真的很厲害.
他這麼厲害.
不會是他自己這麼簡單.
是有上帝祝福在他身上.
他很會說話.
很有禮貌.
很君真.
他很認識當地人.
是要辦妥這件事.
在眾人面前.
在城門口辦妥這件事.
各人都知道.
他會由外人成為本地人.
可以將死人埋葬在那裡.
所以我們見到這是你的恩典.
幫助我們學習亞伯拉罕處事小心.
人情世故.
禮貌君真.
讓我們在行事上無可指責.
求主聘我們這個職位.
我們懇求你幫助我們.
祈禱奉耶穌名求.
阿們.
\newpage



\section{撒母耳記上 17:21-26}
\label{sec:031708}
\textbf{成就不可能的任務 (粵/華)}~~~~許永勝傳道
\newline
\newline
連結: \href{http://www.church.com.hk/acms/content.asp?site=cdc\&op=show\&type=product\&code=031708\&layout=sermon}{\texttt{031708}} ~~~~ 語音日期: 2026-02-22
\newline
\newline
\hyperref[sec:031707]{\small{< < < PREV SERMON < < <}}
~
\hyperref[sec:index_chronic]{\small{[返順時目]}}
~
\hyperref[sec:031709]{\small{> > > NEXT SERMON > > >}}
\newline
\newline
撒母耳記上 17:21-26
\newline
\begin{longtable}{cl}
\hline
\hline
章節 & 經文 (和合本修訂版)\\
\hline
17:21 & \begin{tabularx}{0.7\textwidth}{X} 以色列人和非利士人都擺列陣勢，彼此相對。 \end{tabularx} \\ \\ \relax
17:22 & \begin{tabularx}{0.7\textwidth}{X} 大衛把東西留在看守物件的人手中，跑到戰場，問他哥哥好。 \end{tabularx} \\ \\ \relax
17:23 & \begin{tabularx}{0.7\textwidth}{X} 他與他們說話的時候，看哪，那挑戰的人，就是迦特的非利士人歌利亞，從非利士隊伍中上來，說了同樣的話，大衛聽見了。 \end{tabularx} \\ \\ \relax
17:24 & \begin{tabularx}{0.7\textwidth}{X} 以色列眾人看見那人就非常害怕，從他面前逃跑。 \end{tabularx} \\ \\ \relax
17:25 & \begin{tabularx}{0.7\textwidth}{X} 以色列人說：「這上來的人你看見了嗎？他上來是要向以色列人罵陣。若有人能殺他，王必賞賜他大財，將自己的女兒嫁給他，並在以色列人中免除他父家納糧服役。」 \end{tabularx} \\ \\ \relax
17:26 & \begin{tabularx}{0.7\textwidth}{X} 大衛對站在旁邊的人說：「若有人殺這非利士人，除掉以色列人的羞辱，他會怎樣呢？這未受割禮的非利士人是誰，竟敢向永生神的軍隊罵陣！」 \end{tabularx} \\ \\
[1ex]
\hline
\hline
\end{longtable}
$^{1}$- 姐妹平安! - 弟兄姊妹平安!.
- 今天已經是年初六了 - 今天是年初六.
不知道大家在這個農曆新年假期裡面做了些什麼呢?.
不知道大家在農曆新年假期裡面做了什麼呢?.
經歷了些什麼開心事呢?.
經歷了什麼開心的事呢?.
那大家覺得過這個農曆新年令你最開心的是些什麼呢?.
小朋友或者年輕人可能會說.
當然是逗利是.
小孩或者是年輕人可能會說拿紅包.
成年人就可能說有幾天假期可以休息一下或者去旅行.
成年人可能是有幾天假期可以去旅行或者休息一下.
或者有些人就會說可以吃到很多美食.
或者有人說可以吃到很多美食.
我因為幾年都不在香港.
我幾年也不在香港.
今年是第一年跟家人和親戚一起過農曆年.
今年是第一年跟親戚或者家人在香港過農曆年.
所以出席了不少的飯局.
開始的時候吃到好吃的菜都會有些幸福的感覺.
因為畢竟我都吃了五年的韓國菜.
畢竟我也吃了五年的韓國菜.
終於可以吃到香港的菜了.
連續幾天出席飯局之後.
每一餐都吃到飽到不得了.
每一頓飯也吃到非常飽.
感覺每次都吃到快想吐.
所以幸福感一下子就沒有了.
幸福感一下子就下降了.
其實關於過農曆新年的開心事.
我覺得可能還有一樣的小幸福大家可能不太為意.
還有一個小的幸福可能大家沒有為意.
讓我講出來看看大家同不同意.
讓我講出來看看大家是否同意.
就是過農曆新年的時候.
讓我們可以再一次有一個機會.
為新的一年訂下一些目標和願望.
就是每一個農曆新年.
我們可以為新的一年訂下一些目標或是願望.
因為就是我們才過農曆年.

$^{41}$就是因為我們才過農曆年.
所以我們有兩次新年.
我們有兩次的新年.
但是這個2026年轉眼就已經快過兩個月了.
所以這個2026年轉眼就已經快過兩個月了.
那麼起初我們訂立的新的目標或者願望.
起初我們過了一些新的目標或是願望.
可能在實踐方面就感覺有些不順.
在實踐的方面可能不是很順.
譬如你想改變一下你自己的一些壞習慣.
比如想改一下我們自己的壞習慣.
或者你很想建立一些好的習慣.
或者建立一些好的習慣.
但是過了兩個月一點進展都沒有.
因為兩個月一點進展都沒有.
甚至有些人都已經放棄了.
我年頭訂的目標就是做運動和去健身.
我年初訂的目標就是做運動還有健身.
那麼可惜兩個月過了就一點都沒開始.
可惜兩個月過了就一點都沒開始.
我老婆經常就說我.
你為什麼這麼懶一點都不想動.
我太太就跟我說為什麼這麼懶不想動.
我沒有說謊的她真的經常罵我.
我沒有說謊的她經常罵我.
所以過農曆年令我再次定立目標和願望.
所以過年的時候我就可以重新定下一個目標和願望.
希望新的農曆年我可以實踐我的目標.
在新的農曆年希望我可以實踐我的目標.
今天我想和大家分享的主題就是不可能的任務.
成就不可能的任務.
今天跟大家分享的題目是成就不可能的任務.
對基督徒來說最困難和成就不了的事.
可能未必是人生的目標或者願望.
對基督徒來講成就不到的事未必是人生的目標或願望.
可能是在信仰的實踐上面.
譬如過去幾年的世紀疫情或者是社會運動.
如過去幾年的世紀疫情或者是社會運動裡面.
大家都肯定經歷過那種無能為力的感覺.
大家肯定會經歷過那種無能為力的感覺.

$^{81}$在那段時間很多傳福音的施工都是停止了.
在那段時間裡面很多傳福音的施工也因此而停止.
過去幾年停止了的施工.
譬如宣教都好像很難或者沒有動力的重新開始.
過去幾年可能停止了宣教也可能沒有動力重新開始.
又譬如AI時代突然降臨.
我們很怕被人騙.
關心其他人和彼此相愛就變得格外小心.
關心他人或彼此相愛我們就可能格外小心.
今天我所分享的訊息是關於大衛對戰巨人哥利亞的故事.
今天我分享的訊息就是大衛對戰巨人哥利亞的故事.
我相信大家都對這個故事很熟悉.
我相信大家對這個故事也非常熟悉.
不過大家知不知道大衛是怎樣決定去接受這個所有人都認為不可能的任務呢?.
我大衛怎樣去決定去接受這個巨人哥利亞的任務呢?.
究竟需要什麼條件和因素去令到看似不可能的任務變成可能呢?.
需要什麼條件和因素去應對這個不可能的任務呢?.
希望透過今天的分享我們都好像大衛一樣.
在這個很不容易的世代成就神給我們基督徒的使命.
在這個不容易的世代成就神給我們基督徒的使命.
我會分享兩個成就不可能任務的關鍵因素.
我會分享兩個成就不可能任務的兩個關鍵的因素.
第一個因素就是需要願意信服隨時擔起這個使命的心.
第二個因素就是需要願意信服隨時擔起這個使命的心.
首先我將一些背景講給大家聽.
在今天的經文前面的17章1至11節提到的內容.
首先我跟大家分享一下經文的背景.
就是前面提過17章1到11節的內容.
大背景就是以色列的敵人非利士人要來跟以色列人征戰.
大背景就是以色列的敵人非利士人要跟以色列人征戰.
雙方的軍隊就在伯利恆附近的山谷對峙.
雙方的軍隊就在伯利恆的山谷裡面對壘.
我給大家看一下地圖讓大家有一個更立體的印象.
這是地圖,他們在以拉谷對峙.
另外有一張實景圖,可能還看得比較好.
還有一個實境圖.
其實說是山,最多來講只能稱為小山丘.
與其說是山,其實它只是一個小山丘.
要打起來,衝過去就是了.
古代戰爭,兩軍對壘的時候.

$^{121}$雙方都會派出戰士來對打.
古代的戰爭,雙方都會派出戰士互相對打.
看看誰會贏.
有點像我們中國看三國演義電視劇.
像中國我們看三國演義的電視劇.
都看到兩個將軍在對打.
你看到兩個將軍在對打.
最後的戰局通常都是打贏的將軍得到的.
最後的勝局都是打贏的那一方的將軍得到的.
還有這個故事是關於大衛對戰巨人哥利亞.
這個故事是大衛對戰巨人哥利亞.
哥利亞究竟有多巨大呢?多可怕呢?.
又給大家看一圖了.
給大家看一看圖.
聖經有特別說的.
他身高就是六爪零一虎口.
他身高是六爪零一虎口.
一爪就是指我們一般人的手肘的長度.
一爪就是手肘的長度.
六爪零一虎口大概是三米多.
不是三米多是一爪三米多.
算起來是九尺.
九尺.
大概九尺.
九尺多高.
九尺多.
其實我們最出名的籃球明星是姚明.
最著名的籃球明星就是姚明.
他也是七尺多.
他也只是七尺多.
哥利亞有九尺多.
哥利亞有九尺.
這個圖比較了大衛和哥利亞的高度.
其實差不多高一倍.
差不多高一倍.
很多人害怕.
哥利亞比大衛高一倍是超級恐怖的.
哥利亞比大衛高一倍是非常恐怖的.
還有聖經形容巨人的戰甲重五十七公斤.
甚至形容巨人的戰甲重五十七公斤.

$^{161}$武器也重七公斤.
這樣的巨人戰士誰敢挑戰他呢?.
誰有可能打贏他呢?.
所以十一節描述的以色列人.
裡面沒有人敢出去跟他打.
大家都很驚慌和極其害怕.
現在大家不要馬上審判以色列人.
尤其在這個對立很嚴重的世代.
詐騙很猖狂的世代.
我們怎樣去活出這個愛呢.
我們怎樣活出這個愛呢.
怎樣去講饒恕呢.
怎樣講饒恕呢.
跟你擁有不同政治立場的人.
你可不可以愛他呢.
跟你擁有不同政治立場的人.
你可以愛他嗎.
你可不可以關心他呢.
你可以關心他嗎.
可不可以傳福音給他們呢.
真的可以嗎.
其實很難.
真的很難.
我記得2014年我在香港讀神學的時候.
那時候香港正在發生佔領中環事件.
當時香港就發生一個佔領中環的事件.
那時候你跟我在政治立場上都有不同的意見.
當時同學在政治立場上都有不同的意見..
當時有人跟我的政治立場也有不同的意見..
不是跟我,是大家同學都彼此,彼此不是只有我..
同學之間也有不同的意見..
在課堂會激烈地討論..
在課堂裡面也激烈地討論..
當然我都深陷其中..
我也深陷其中..
有一次我跟一位同學就爭論到要老師來調解一下..
爭吵之後我跟那位同學就產生了隔閡..
簡單來說就是面左左..
就是過不去,對,面左左,翻譯不通..
就是看得不順眼,順不順眼..

$^{201}$回家後我反思自己,聖靈也責備和提醒我..
如果連將要成為傳道人都不能夠接納跟自己意見不合的人呢?.
如果將要成為傳道人的時候,我也不能夠接納其他人意見的時候..
那我還有沒有資格做傳道呢?.
還可不可以去遵守聖經的教導呢?.
於是我就去找那位同學..
我跟他道歉和分享聖靈的提醒..
他有相同的感受..
他也有相同的感受..
最後我們成為了班裡最好的朋友..
雖然在政治取向上,我們到今天還是不一樣的..
雖然在政治取向上,我跟他現在也是不同的..
但是我們一直都彼此關心對方..
所以有一些聖經教導,譬如說愛你的敵人,.
愛你不想愛的人,.
譬如有一些聖經的教導,.
愛你的敵人,愛你一些不想愛的人,.
要完成這些自己認為不可能做到的聖經教導,.
\newpage



\section{列王記下 2:1-18}
\label{sec:031709}
\textbf{雙倍恩膏,薪火傳承}~~~~溫艾衛傳道
\newline
\newline
連結: \href{http://www.church.com.hk/acms/content.asp?site=cdc\&op=show\&type=product\&code=031709\&layout=sermon}{\texttt{031709}} ~~~~ 語音日期: 2026-02-22
\newline
\newline
\hyperref[sec:031708]{\small{< < < PREV SERMON < < <}}
~
\hyperref[sec:index_chronic]{\small{[返順時目]}}
~
\hyperref[sec:031710]{\small{> > > NEXT SERMON > > >}}
\newline
\newline
列王記下 2:1-18
\newline
\begin{longtable}{cl}
\hline
\hline
章節 & 經文 (和合本修訂版)\\
\hline
2:1 & \begin{tabularx}{0.7\textwidth}{X} 耶和華要用旋風接以利亞升天的時候，以利亞與以利沙從吉甲往前行。 \end{tabularx} \\ \\ \relax
2:2 & \begin{tabularx}{0.7\textwidth}{X} 以利亞對以利沙說：「耶和華差遣我往伯特利去，你可以留在這裡。」以利沙說：「我指著永生的耶和華，又指著你的性命起誓，我必不離開你。」於是二人下到伯特利。 \end{tabularx} \\ \\ \relax
2:3 & \begin{tabularx}{0.7\textwidth}{X} 在伯特利的先知的門徒出來，到以利沙那裡，對他說：「耶和華今日要接你的師父離開你，你知不知道？」他說：「我知道，你們不要作聲。」 \end{tabularx} \\ \\ \relax
2:4 & \begin{tabularx}{0.7\textwidth}{X} 以利亞對以利沙說：「耶和華差遣我往耶利哥去，你可以留在這裡。」以利沙說：「我指著永生的耶和華，又指著你的性命起誓，我必不離開你。」於是二人到了耶利哥。 \end{tabularx} \\ \\ \relax
2:5 & \begin{tabularx}{0.7\textwidth}{X} 在耶利哥的先知的門徒來靠近以利沙，對他說：「耶和華今日要接你的師父離開你，你知不知道？」他說：「我知道，你們不要作聲。」 \end{tabularx} \\ \\ \relax
2:6 & \begin{tabularx}{0.7\textwidth}{X} 以利亞對以利沙說：「耶和華差遣我往約旦河去，你可以留在這裡。」以利沙說：「我指著永生的耶和華，又指著你的性命起誓，我必不離開你。」於是二人一同往前行。 \end{tabularx} \\ \\ \relax
2:7 & \begin{tabularx}{0.7\textwidth}{X} 有五十個先知的門徒同去，遠遠地站在他們對面；他們二人在約旦河邊站住。 \end{tabularx} \\ \\ \relax
2:8 & \begin{tabularx}{0.7\textwidth}{X} 以利亞捲起自己的外衣，用來打水，水就左右分開，二人走乾地過去。 \end{tabularx} \\ \\ \relax
2:9 & \begin{tabularx}{0.7\textwidth}{X} 過去之後，以利亞對以利沙說：「我未被接去離開你以前，你要我為你做甚麼，只管求。」以利沙說：「願感動你的靈雙倍感動我。」 \end{tabularx} \\ \\ \relax
2:10 & \begin{tabularx}{0.7\textwidth}{X} 以利亞說：「你求的是一件難事。我被接去離開你的時候，你若看見我，就必得著；若不然，就得不著了。」 \end{tabularx} \\ \\ \relax
2:11 & \begin{tabularx}{0.7\textwidth}{X} 他們邊走邊說話的時候，看哪，有火馬和火焰車出現，把二人隔開，以利亞就乘旋風升天去了。 \end{tabularx} \\ \\ \relax
2:12 & \begin{tabularx}{0.7\textwidth}{X} 以利沙看見，就呼叫說：「我父啊！我父啊！以色列的戰車騎兵啊！」以利沙不再看見他的時候，就把自己的衣服撕為兩片。 \end{tabularx} \\ \\ \relax
2:13 & \begin{tabularx}{0.7\textwidth}{X} 他拾起以利亞身上掉下來的外衣，回去站在約旦河邊。 \end{tabularx} \\ \\ \relax
2:14 & \begin{tabularx}{0.7\textwidth}{X} 他用以利亞身上掉下來的外衣打水，說：「耶和華－以利亞的神在哪裡呢？」打水之後，水也左右分開，以利沙就過去了。 \end{tabularx} \\ \\ \relax
2:15 & \begin{tabularx}{0.7\textwidth}{X} 在耶利哥的先知的門徒從對面看見他，說：「感動以利亞的靈臨到以利沙身上了。」他們就來迎接他，俯伏於地，向他下拜， \end{tabularx} \\ \\ \relax
2:16 & \begin{tabularx}{0.7\textwidth}{X} 對他說：「看哪，僕人這裡有五十個壯士，請你讓他們去尋找你師父，或者耶和華的靈將他提起來，投在某山某谷。」以利沙說：「你們不必派人去。」 \end{tabularx} \\ \\ \relax
2:17 & \begin{tabularx}{0.7\textwidth}{X} 他們再三催促，直到他不好意思，就說：「你們派人去吧！」他們就派了五十個人去，尋找了三天，也沒有找著他。 \end{tabularx} \\ \\ \relax
2:18 & \begin{tabularx}{0.7\textwidth}{X} 以利沙仍然留在耶利哥，他們回到他那裡，他對他們說：「我不是告訴你們不必去嗎？」 \end{tabularx} \\ \\
[1ex]
\hline
\hline
\end{longtable}
$^{1}$各位弟兄姊妹平安.
今天是年初六.
今年是馬年.
在農曆新春期間.
我都祝大家.
人強馬壯.
含義是說我們體力很強健.
人力很充沛.
更多人去信主.
我們每一堂崇拜都是壯無虛席的.
並且各位弟兄姊妹的靈命都是與日俱增.
然後我們大家都馬上喜樂.
今日的題目是.
雙倍恩高.
薪火全盛.
在列王記夏二章一至十八節.
剛才我們已經讀過經文.
先知是神的代言人.
即是作為神與人之間的代理.
會向人提出警告.
叫百姓迴轉.
甚至宣佈刑罰.
他的職責是要專心服侍上帝.
先知也會高納人為王.
和高納人為接班人.
舊約時代的以色列王國.
都是先知去建立.
去輔助和監督.
今日我們所看的經文.
出現了兩個很重要的先知.
就是以利亞和以利沙.
他們的事跡主要記載在列王記上和列王記夏.
他們是使徒的關係.
他們是怎樣遇上的呢?.
在列王記上的十九章說到.
耶和華叫以利亞去到大馬士革.
去告以利沙.
作先知接續他自己.
以利亞第一次遇見以利沙的時候.
以利沙在哪裡呢?.

$^{41}$以利沙原來在田中.
在耕田中.
之後以利亞經過他.
將自己的外袍.
搭在他的肩膊上.
以利沙明白到以利亞的意思.
以利沙和他的父母道別.
然後跟著以利亞服侍他.
以利亞很清楚.
耶和華要告他的接班人.
將先知的即時傳承給他另外一個人.
當以利沙見到以利亞之後.
以利沙也很清楚.
耶和華的呼召要臨到他身上.
以利沙信服神.
放下自己手上一直做的耕田工作.
跟著以利亞全身投入神的召命.
到了今天.
經文就是列王幾下第二章.
這段經文是記述.
以利亞被接升天的過程.
我們今天看看這段經文.
給我們有什麼啟迪.
有什麼看見.
我們一起同心圖告.
主耶和華和您這一刻要我們同在.
帶領著我們.
開我們的眼睛 開我們的耳朵.
讓我們聽得明白 看得見你.
主 讓今天的訊息.
能夠深深地進入我們心裡.
並且我們走出來見證你.
我們祈禱奉主耶穌基督的聖名自祈求.
阿們.
我們看看今天第一個重點.
就是不離不棄 堅定跟隨.
在列王幾下二章一至六節.
經文一開始的時候.
提到耶和華用旋風去接以利亞.
升天的時候將到.

$^{81}$於是兩使徒就由以法連山區的吉甲.
一路步行去到伯特利.
然後就由伯特利經過耶利哥.
去到約旦河.
以利亞在行程裡.
有三次要求以利沙留在原地去等他.
在第二 第四和第六節裡.
都有提到 但是每一次.
以利沙都堅決地表示.
他說 我指著永生的耶和華.
又指著你的聖命起誓.
我必不離開你.
以利沙的那個起誓.
表示他很堅決 很堅定.
表示他甘願地承擔起誓的責任和後果.
為什麼以利沙要堅持跟著.
以利亞去伯特利 去耶利哥和約旦河呢?.
其實很顯然是因為.
他知道神很快將他的師父接了.
他很珍惜和他的師父這段時間.
也很想繼續謙卑地跟隨他的師父學習.
以利亞很快就被接走了.
這個不單止是以利亞.
以利沙知道 經文說到.
連在伯特利 在耶利哥的先知門徒.
他們都很清楚知道這件事.
當以利亞和以利沙來到伯特利的時候.
在伯特利的先知門徒出現.
見到以利沙 他就跟他說.
耶和華今天要接你的師父離開你.
你知不知道?.
其後兩使徒來到耶利哥.
當地的先知門徒同樣問一樣的問題.
想要提醒以利沙.
關於以利亞很快會被接去的事.
以利沙兩次都同樣回答他們的問題.
「我知道 你們不要作聲」.
相信以利沙當時的心情很複雜.
他很傷心 因為師父要很快離開他.
他的心情也很不捨.

$^{121}$所以這句話大概就是說.
「我記得了 你們不用再提我了」.
先知門徒這兩句問他的說話.
就顯示神當日要將以利亞接去.
原來已經成為一個公開的秘密.
為什麼大家都知道這個公開的秘密呢?.
其實這個也不是什麼特別的公開秘密.
因為事情不單止以利亞 以利沙知道.
其實以利亞早就跟他們說過.
這個消息也告訴了先知門徒.
大概以利亞也同時宣告.
以利沙要在他被升天之後.
繼承他領導的地位.
以及繼承他作為先知的事奉.
以利亞每次經過一個地方的時候.
以利亞就叫以利沙停下來.
目的大概是要試驗一下以利沙.
是否一直有這個警覺.
警覺他自己會離開.
徒弟以利沙在這方面其實也很警覺.
當時也反映到他對於神給他這個事奉的堅定.
隨時後命去預備回應神給他的呼召.
記得我讀神學的時候.
我要去堂會實習.
有位牧者跟我說.
跟這位準傳道人說.
如果對真理的事不去堅持的話.
並且去寬鬆處理的話.
對其他人是一個殘忍.
為什麼會這樣說呢.
意思就是說.
其他人可能是有樣學樣.
以為這樣做是對的.
會令其他人很迷失.
甚至乎日後他們會跌倒.
所以我們對於真理一定要堅持堅定.
神是輕慢不得的.
我們看一段經文.
《迦拉提書》六章七至八節.
我們一起讀這段經文.

$^{161}$預備起.
不要自欺神是輕慢不得的.
因為人種的是什麼.
收的也是什麼.
信著肉體殺種的.
必從肉體收敗壞.
信著聖靈殺種的.
必從聖靈收永生.
如果教會,信仰圈.
服侍環境裡面.
有真理的偏差的時候.
跟錯了人.
而不是跟隨神.
而且又沒有人去指證出來的話.
那這個還是不是一個屬靈的群體呢.
真正的拆毀是什麼呢.
就是在屬靈上有錯誤.
一直存在著.
沒有人去指證出來.
就好像有一間房子.
有一個很大的裂縫.
住在裡面的人.
天天都看到.
他們習慣了.
那個裂縫仍然存在著.
而且也是慢慢越來越大.
可能還會影響到結構的核心.
領袖的角色.
發現了屬靈上的錯誤.
他就要去處理.
而不是要保護這個犯罪的肢體.
大家假裝在那裡睡覺.
大家一起叫不醒.
這樣只會令這件事沒有出路.
繼續陷入罪的當中.
反教會每一次見面的時候.
不是說你好我好.
大家好.
不是的.
我們是要背起十字架.

$^{201}$面對真實的自己.
要付出代價的.
耶穌為了我們.
也付出了很大的代價.
這份恩典是一個重大的恩典.
所以鼓勵大家.
在我們每天禱告的時候.
尤其是晚上禱告的時候.
我們要反省自己.
當天所做的事.
是不是通過陽光的測試.
為什麼呢.
就是說我們這件事.
是不需要在黑夜裡才去做的.
可以光明正大地去做.
信主基督徒也深信聖靈會提醒我們.
制止我們不要做得罪神的事.
而真正的建造有什麼呢.
是要讓人有短暫的不舒服.
但是方向是校正的.
各位親愛的領袖.
我們要將羊帶到神的面前.
叫他們堅持去跟從神.
一起去看聖經.
一起去明白真理.
去守望我們的羊.
而不是叫他們跟隨自己.
或者是錯誤地跟從自己之後.
也不去糾正他們.
不是這樣的.
也不是叫羊跟隨牧者.
或者跟隨區長,組長.
不是這樣的.
因為跟從人真的很危險.
因為人會犯錯.
人是不可靠的.
所以當人一旦走錯路的時候.
羊就會怎樣呢.
當然牠們會迷失方向.
會跌倒也不知道怎麼辦.

$^{241}$所以各位羊仔.
在座的羊仔.
羊仔我們同樣要跟從.
一定要跟從一個最大的.
最厲害的.
一個創造天地萬物的主.
大家確實明白清楚.
應該知道跟隨誰了.
我們看看第二個重點.
雙倍的靈.
神的應許.
在這幾下.
二章七至十二節.
當以利亞和以利沙.
繼續前行到約旦河的時候.
有五十個先知門徒.
一同跟著去.
離遠地站在他們的對面.
相信先知門徒要跟著他們兩使徒.
目的是要看看究竟有什麼事會發生.
他們兩使徒來到約旦河邊的時候.
雖然沒有渡海輪或者山板那些快艇.
不過以利亞有他的辦法.
他拿著自己的外衣.
將它捲起來.
然後擊打河水.
河水就分開了.
以利亞和以利沙就在這個江地裡.
走過去到約旦河的另一個岸邊.
這個水分開的神蹟.
也讓我們記起.
在摩西帶領以色列人出埃及的時候.
摩西同樣洗到紅海的水左右分開.
以色列人在江地裡走過去.
以利亞現在也是這樣選擇這樣做.
相信是要讓以利沙知道.
神從摩西開始.
在以色列人中間所做的揀選和拯救.
這些工作是沒有間斷過的.
神的作為是繼續在以利亞的事奉裡延續著.

$^{281}$以利沙大概也將師父以利亞所做的一切記在心裡.
一會兒當他要從約旦河回去另一邊岸的時候.
他也是將約旦河的水分開走過去.
以利亞和以利沙從江地走過去之後.
以利沙還記得神很快要將他施父接走.
所以這個徒弟寸步不離地跟著他施父.
以利沙總算這次通過了以利亞的考試.
以利亞接著就問以利沙.
他說:我未被接去離開你以前.
你要我為你做什麼?指管球.
這個指管球不是說以利亞有求必應.
意思是儘管開心見誠地說出來.
看看我能不能為你做些什麼.
事實上我們看經文說到.
以利沙所求的以利亞有條件地回答他.
以利沙馬上回答以利亞.
他說:願感動你的靈相配感動我.
可以見到他實在想過師父的問題.
他之前已經想了.
徒弟早就知道師父要去被接.
也知道師父要將先知的福事傳給他.
他見到自己未來的困難很大 很大的挑戰.
以利沙更知道他需要雙倍以利亞的屬靈能力才能應付.
確實以利沙日後所行過的神蹟.
按聖經記載是多過以利亞所做的.
有學者統計過.
以利亞行過的神蹟有八個.
而以利沙承接了師父的屬靈能力.
相配行了十六個神蹟.
以利亞如何回答他的徒弟以利沙呢?.
以利亞說:你求的是一件難事.
聖經不是說在神裡面沒有難成的事嗎?.
而事實上神賜給以利亞的屬靈能力.
或者是他的徒弟.
神賜給以利沙之後的屬靈能力不是雙倍的嗎?.
大家可以看得到.
也是和他同在.
為什麼以利亞這樣說這是一件難事呢?.
究竟以利沙所求的是有多難呢?.
當我們反思我們服侍上帝的時候.

$^{321}$可能會感到很難.
或者力不從心.
很多時候出來的果效.
跟我們想像有很大的落差.
不似預期.
有人會失望地敗退.
甚至會放棄服侍.
但有些人會求主賜給他更大的智慧.
更大的能力去應付.
忠心地去跟從主.
忠心地去服侍主.
神的差遣去跟從主.
不知道大家是屬於哪一類型的服侍者呢?.
我們經常會出去傳福音.
做關懷事工的工作.
探望我們年長的弟兄姊妹.
有些可能住在家居.
或者住在護老院.
陪陪他們.
期望他們的探訪可以給他們一點溫暖.
去應對他們的所改變.
但有些時候我們都會給一些關懷者.
態度有點冷漠.
甚至拒絕.
「你不用來探望我了」.
「我OK的,不用來了」.
可能我們都會感到心灰意冷.
覺得自己所付出的好像都沒有什麼回報.
上兩個星期.
2月8日.
我們有兩個成人牧區.
跨區去服侍.
在社區寫了「揮春街頭報道」.
當日一開始的時候.
有一位機構的牧師.
他分享了如何使用手上的揮春去傳福音.
因為揮春前面是祝福的語句.
後面是五色珠的內容.
令我們有很大很大的鼓勵.
他說.

$^{361}$跟別人去傳福音.
分享福音的時候.
神會賜給我們有雙倍的福氣.
即是Double Blessing.
第一個福氣是神賜給傳福音的人.
即是我們.
第二個福氣是神賜給聽福音又接受的人.
這樣就很好.
如果對方拒絕了我們所傳的福音.
第二個福氣就會回到我們身上.
即是我們有雙倍的福氣.
繼而我們就有力量去傳福音.
將神的福氣繼續送到我們街坊的心裡.
那天大概有40位弟兄姊妹一起出隊.
每位弟兄姊妹都很積極去傳福音.
一點都沒有灰心過.
我們使用了剛才所說的五色珠.
區長幫我們統計過.
大概當天我們都派了.
都是接近有福字的徽章.
有印刷的徽章 有寫的徽章.
都是接近有一千張.
我們都有170位街坊願意停下.
聽我們講整傳五色珠的福音.
並且願意讓我們做一個祝福的祈禱.
如果他們有決志的.
都為他們做一個決志的祈禱.
當我們服侍上帝的時候.
感到力不從心的時候.
或者效果不似預期的時候.
其實很常見.
這些經歷反映出.
我們屬靈的生命成長過程當中.
面對挑戰,挫敗,真實的狀況.
重要的是我們要保持一個信心.
持續的禱告.
依靠神的恩典和帶領.
讓神可以在我們身上施一個恩典.
神是著重我們的心智和我們的忠心.
而不是只看我們外表所服侍的情況或成績.

$^{401}$神會看到我們有心意和持續的愛心.
我們能夠透過服侍.
自己也會持續地學習忍耐和有無條件的愛.
以彌撒這一份「我必不離開你」的堅持.
也同樣體驗到對神的應許的抓緊.
讓我們反思到自己.
是否也具有屬靈的心智呢?.
在服侍當中是否也經歷到神的能力.
以及如何去接續為這個信仰傳承去努力呢?.
當我有時要開祖去探祖的時候.
聽到一些退休的祖源這樣說.
「我已經退休了」.
「已經沒有甚麼寄聖」.
「開始看聖經的時候」.
「我就很睏 要打瞌睡」.
然後我心裡就想.
其實很多學生一打開那本書.
都是一樣會感到睏的.
不關年齡是否退休的事.
也不一定是因為要看聖經才會睏.
這位退休的祖源不單這樣說.
他還說「我字也不識多隻」.
「上去看聖經的時候」.
「有很多不明白」.
「很多不理解的地方」.
「沒辦法了」.
「都是要這樣」.
其他的祖源就鼓勵他.
他跟他說.
「先不要說不行」.
「不可以」.
因為如果這樣說.
即是說還沒開始的時候.
就已經認定了自己是不行的.
我很同意這位祖源的鼓勵.
因為我們真的不要看輕自己.
因為在神是凡事都能的.
神會賜給我們有足夠的能力.
去給每一位追求真理的頂子妹.
當然聖靈也會指示我們去明白聖經.

$^{441}$雖然以利沙所求的事.
對於以利亞來說是有點難.
但神也是透過以利亞對以利沙作了指示.
祂說「我被接去離開你的時候」.
「你若看見我,就必得著」.
「若不然,就得不著了」.
這個就是說.
如果看見以利亞被接去的情景.
以利沙所求加倍的屬靈能力.
就必得著了.
即是說如果看不見以利亞被接去的時候.
以利沙所求的就得不著.
但我們會問.
以利沙不是已經跟著師父去到約旦河邊了嗎?.
他已經是寸步不離了.
以利沙還有可能看不見師父被接走的情景嗎?.
神要以利亞回答以利沙.
其實顯然也是要告訴以利沙.
祂先前一直這麼警醒地重視神所託付給祂的即時.
神其實已經預立了祂.
願意加倍賜予以利沙屬靈的能力.
來承擔祂的事奉.
但為了使得以利沙具體地知道神的心意.
神藉著以利沙看見以利亞升天的經歷.
告訴以利沙神必定與他同在.
賜予這個屬靈的能力.
雙倍臨到以利沙身上.
但我們要想深一層.
就算以利沙跟著以利亞來到約旦河邊.
以利沙也不一定可以看見以利亞被接去的情景.
記得在後面的經文裡.
雅蘭王派人到多丹捉拿以利沙的時候.
以利沙的僕人看不見以利沙所見到的滿山火車火馬.
以利沙求神開他僕人的眼睛.
讓他僕人才可以看見.
換言之.
以利沙之所以可以看到他和以利亞分開的火車火馬.
又見到以利亞被旋風接走的時候.
並不單單因為以利沙本人在現場.
而是完全因為神使得以利沙看到這個情景.

$^{481}$而那五十位先知門徒.
相信他們也看不見以利亞被接去的情景.
以利沙看到以利亞被接去.
其實是一個證據.
是神要向以利沙表明神願意給他雙倍的屬靈能力的證據.
經文提到有火車火馬將兩個人分開.
不論這是一個真實有戰車戰馬放出一些火焰.
或者是一個異象.
這些景象以利沙全部都看得見.
亦都交代了以利沙將會得到神的應許.
將屬靈的能力雙倍降臨在他身上.
我們看看第三個重點.
是得著能力繼承使命.
在耶和華記下二章十三至十八節.
在十三十四節.
以利沙被旋風接走的時候.
他的外衣掉下來.
而且這個外衣掉下來跟以利沙不太遠.
所以以利沙撿起了師父以利亞的外衣.
這個舉動亦表示.
以利沙願意繼承以利亞先知的事奉.
在列王記上十九章提及.
以利亞受神託父去告以利沙為先知.
去接續先知的事奉的時候.
以利亞就將他的外衣搭在以利沙肩膊上.
外衣亦與傳承有關.
這個一個新的以利沙就呈現出來了.
他以接續以利亞作先知的身份出現了.
所以如果有屬靈的長輩走過來.
將他的外衣搭在你的肩膊上的時候.
大家想一想是不是有點意思.
是不是有點古怪呢?.
當以利沙拿著這個外衣的時候.
他就說:耶和華,以利亞的神在那裡呢?.
這句話是一個禱告.
祈求神向他的僕人顯示他的能力.
而且這個能力是從他的師父的神耶和華而來的.
直至他證明他承接承繼以利亞的職務.
果然當以利沙一打這個水的時候.
同樣這個水也是分開.

$^{521}$以利沙就走過去了.
回顧以色列人的舊恩歷史.
我們曾經經歷過水分開.
人走在旱地裡的神蹟.
莫西大令以色列人過紅海.
約書亞也是大令以色列人走過約旦河.
以利亞擊打這個水.
就好像這個舉動是鼓勵以利沙.
上帝要將先知的工作.
屬靈的能力在以利沙身上延續下去.
而外衣更加是上帝給過以利沙.
接觸新的挑戰的一個記號.
當先知門徒見到.
以利沙在約旦河走過去的時候.
就知道耶和華和以利沙同在了.
原來在約旦河.
即是原先約旦河的50位先知門徒.
他們很遠的望著以利亞和以利沙去過對面岸.
現在只見到以利沙一個人回來.
以利亞不見了.
但當他們見到以利沙用外衣.
將河水分開的時候走過來.
他們就知道也承認.
感動以利亞的靈已經臨到以利沙身上.
亦表示以利沙要接觸以利亞的地位.
以後這班先知門徒就要被以利沙照顧.
於是他們就迎接以利沙.
俯伏在地上高.
表示他們心悅誠服地接納以利沙成為領導的地位.
十六至十八節.
由於其他人看不見以利亞被接去.
所以單憑以利沙一個人說的話.
這班先知門徒對以利亞未能忘記.
或完全消化到這件事.
建議以利沙派人去看看.
以利亞是否真的不在.
最終他們在山谷仔細尋找三天.
三天尋找都找不到.
徒勞無功.
這代表以利亞已經被接去.

$^{561}$亦肯定以利沙對他們所說的話是真實的.
而且帶有權威.
以利沙不但承接了師父以利亞先知的衣缽和能力.
更加帶有雙倍的屬靈能力去服侍上帝.
我們反思.
以利沙對於師父的離開.
他沒有因為這樣而迷失.
因為他知道自己的身份.
亦知道這個份要在裡面盡心盡力愛主.
做主喜悅的事.
繼承師父的使命.
同時亦知道神必定與他同在.
神會給他不同的印證.
我們去完成一個任務.
不要說是工作上的任務.
可能是學習上的任務.
或者是我們服侍上的任務.
原來都有分開三個層次.
第一個就是做了一件事.
第二個就是做到一件事.
第三個就是做好一件事.
做了一件事是什麼意思呢?.
這件事基本上已經完成了.
有人叫我去做,我就去做.
未必會用心去完成這件事.
為了做,要應付.
好像之前港鐵沙中線.
紅隧站的月台工程.
為了追趕完成的時間.
鋼筋剪短了.
引發質量問題.
這個就是打算做了一件事.
交代了,完成了.
遇到的期限就算了.
交差.
做到一件事是什麼呢?.
表示一件事成功完成了.
並且已經達到目標.
結果都出來了.
可能都有盡力去做.

$^{601}$但未必用心去做.
我記得讀書的時候.
放暑假要做暑期作業.
通常都不會一放假.
要回去就馬上去做.
絕對不會這樣.
一直玩,玩,玩.
玩到差不多要有新的開學.
一兩個星期前才開始做.
當然我就比較馬虎地去做.
快點講完它.
是做了一件事.
完成了這個功課.
但我沒有用心去做.
也在當中學習不到任何東西.
做好一件事是什麼呢?.
不單是盡力去做.
而且是用心去完成.
好像媽媽做一餐飯出來.
她都會用盡心思.
用盡力量去準備.
有一個很完善的一餐飯品質出來.
希望家人可以享受到一個豐富的晚餐.
期望各位服侍上帝的僕人.
不單止做了和做到這件事.
去交差.
和要做到神交託給我們的任務.
同樣我們都要做好這件事.
做一個又良善又忠心.
做一個無愧的工人.
這樣去服侍神.
以利沙在聖經記載.
他行過十六個神蹟.
多數和日常生活裡的供應有關.
有醫治的關係.
甚至乎以利沙死了之後.
這個神蹟繼續有出現.
以利沙行事較為溫和.
關注醫治,恩惠.
彰顯上帝的慈愛和豐富.

$^{641}$跟著師父所做的神蹟有些相似.
譬如洗禮人復活.
將水分開都有些相似.
亦都是師父給徒弟的一個榜樣.
但師父做的神蹟有些較為剛烈.
有些都有不同.
神用每一位的僕人.
不是因為他像哪一個.
而是因為他們很忠心.
將神放在僕人心裡的負擔.
每一位的屬靈前輩.
我們都值得去欣賞.
值得去尊重.
他們的生命有些會很溫柔地服侍.
有些可能會比較剛烈一點.
有些會安靜一點.
有些會默默地付出.
這些都是我們學習的榜樣.
繼承使命真正的意思.
是要將神交託給人的使命.
去薪火相傳.
繼續燃點下去.
一個傳一個.
生命影響生命.
使得服侍的心是不會熄滅的.
當教會要作出領袖傳承的時候.
通常都會產生一些不同的領袖風格.
領導的風格.
有些領袖有不同的屬靈的恩賜和才能.
按神的心意賜予他們的恩賜.
互相去配搭.
互相去服侍.
只要我們忠心堅定地去跟從神.
祝福神賜給我們的應許.
神必定會賜給我們加倍的恩高.
並且神會在新一代的領袖身上.
持續地去動功.
來彰顯神的大能的接觸和神的計劃.
這樣神的家就會得以繼續向前行.
一代一代地去傳承下去.

$^{681}$在社區裡面作主的登台.
作主美好的見證.
我們一起同心低頭禱告.
愛我們的主.
求你至下恩高和屬靈的能力.
使得我們能夠忠心地去服侍你.
成為一個無愧的工人.
彰顯你奇妙的恩典和能力.
又讓我們的生命能夠成為流通的管子.
將服侍的棒一代一代地傳下去.
讓筲箕灣浸信會成為明燈.
繼續去見證主的明.
主求你幫助我們.
好像以利沙一樣.
心貼心緊緊地去引從你.
無論前面有多少挑戰和困難.
我們都不輕易地說放棄.
我們要對你有信心.
我們要堅定地去依靠你的力量.
讓你的靈在我們身上運行.
使我們能夠行出你的旨意.
也求你賜給我們智慧和勇氣.
承接屬天的使命.
將你的辛苦一代一代地傳下去.
使更多人因為你的恩典.
得著拯救.
得著生命的更生.
將榮耀頌讚.
歸還給我們三一的真實.
我們祈禱奉主耶穌基督得勝命.
致祈求.
阿們.
\newpage



\section{哈該書 2:1-9}
\label{sec:031710}
\textbf{榮耀的殿}~~~~蕭麗琴傳道
\newline
\newline
連結: \href{http://www.church.com.hk/acms/content.asp?site=cdc\&op=show\&type=product\&code=031710\&layout=sermon}{\texttt{031710}} ~~~~ 語音日期: 2026-02-22
\newline
\newline
\hyperref[sec:031709]{\small{< < < PREV SERMON < < <}}
~
\hyperref[sec:index_chronic]{\small{[返順時目]}}
~
\hyperref[sec:031711]{\small{> > > NEXT SERMON > > >}}
\newline
\newline
哈該書 2:1-9
\newline
\begin{longtable}{cl}
\hline
\hline
章節 & 經文 (和合本修訂版)\\
\hline
2:1 & \begin{tabularx}{0.7\textwidth}{X} 七月二十一日，耶和華的話藉哈該先知傳講，說： \end{tabularx} \\ \\ \relax
2:2 & \begin{tabularx}{0.7\textwidth}{X} 「你要曉諭撒拉鐵的兒子猶大省長所羅巴伯、約撒答的兒子約書亞大祭司，和所有倖存的百姓，說： \end{tabularx} \\ \\ \relax
2:3 & \begin{tabularx}{0.7\textwidth}{X} 『你們中間存留的，有誰見過這殿從前的榮耀呢？現在你們看如何？在你們眼中豈不是如同無有嗎？ \end{tabularx} \\ \\ \relax
2:4 & \begin{tabularx}{0.7\textwidth}{X} 所羅巴伯啊，現在，你當剛強！這是耶和華說的。約撒答的兒子約書亞大祭司啊，你當剛強！這是耶和華說的。這地的百姓啊，你們都當剛強做工，因為我與你們同在。這是萬軍之耶和華說的。 \end{tabularx} \\ \\ \relax
2:5 & \begin{tabularx}{0.7\textwidth}{X} 這是照著你們出埃及時我與你們立約的話。我的靈仍要住在你們中間，你們不必懼怕。 \end{tabularx} \\ \\ \relax
2:6 & \begin{tabularx}{0.7\textwidth}{X} 萬軍之耶和華如此說：過些時候，我必再一次震動天地、滄海與乾地。 \end{tabularx} \\ \\ \relax
2:7 & \begin{tabularx}{0.7\textwidth}{X} 我必震動萬國，萬國的珍寶都必運來，我就使這殿充滿榮耀。這是萬軍之耶和華說的。 \end{tabularx} \\ \\ \relax
2:8 & \begin{tabularx}{0.7\textwidth}{X} 銀子是我的，金子也是我的。這是萬軍之耶和華說的。 \end{tabularx} \\ \\ \relax
2:9 & \begin{tabularx}{0.7\textwidth}{X} 這後來的殿的榮耀必大過先前的榮耀。這是萬軍之耶和華說的。在這地方我必賜平安。這是萬軍之耶和華說的。』」 \end{tabularx} \\ \\
[1ex]
\hline
\hline
\end{longtable}
$^{1}$各位鼎姐妹平安.
今天是年初六.
今年也是馬年.
在農曆新春期間.
我祝大家.
人強馬壯.
含義是說我們體力很強健.
人力很充沛.
更多人去信主.
我們每一堂崇拜都是壯無虛席的.
並且各位鼎姐妹的靈命都是與日俱增.
然後我們大家都馬上喜樂.
今天的題目就是.
雙倍恩高 薪火全乘.
在列王記下 二章一至十八節.
剛才我們已經讀過經文了.
我們先知道先知 先介紹一下.
先知是神的代言人.
即是作為神與人之間的代理.
會向人提出警告.
叫百姓迴轉.
甚至會宣佈他的刑罰.
他的職責是要專心服侍上帝.
先知也會高納人為王.
和高納人為他的接班人.
舊約時代的以色列王國.
都是先知去建立.
去輔助和監督.
今天我們所看的經文.
出現了兩個很重要的先知.
就是以利亞和以利沙.
他們的事跡主要記載在列王記上和列王記下.
他們是使徒的關係.
他們是如何遇上的呢?.
在列王記上的十九章說到.
耶和華叫以利亞去大馬士革.
去告以利沙作先知接續他自己.
以利亞第一次遇見以利沙的時候.
以利沙在哪裡呢?.
以利沙原來在田裡.

$^{41}$在耕田.
之後以利亞經過他.
把自己的外袍搭在他的肩膊上.
以利沙明白到以利亞的意思.
以利沙和他的父母道別.
然後跟著以利亞服侍他.
以利亞很清楚.
耶和華要告他的接班人.
把先知的即時傳承給他另一個人.
當以利沙見到以利亞之後.
以利沙亦都很清楚.
耶和華的呼召要臨到他身上.
以利沙就信服神.
放下自己手上一直做的耕田工作.
跟著以利亞全身投入神的召命.
到今天我們看的經文就是.
列王記夏第二章.
這段經文是記述.
以利亞被接升天的過程.
我們今天就看看這段經文.
給我們甚麼啟迪.
有甚麼看見.
我們一起同心圖告.
主耶和華在這一刻要我們同在.
帶領著我們.
開開我們的眼睛 開開我們的耳朵.
讓我們聽得明白 看得見你.
主 讓今天的訊息能夠深深入到我們的心裡面.
並且我們走出來去見證你.
我們祈禱 奉主耶穌基督的聖名致祈求.
阿們.
我們看看今天的第一個重點.
就是不離不棄 堅定跟隨.
在列王記夏二章一至六接.
經文一開始的時候.
提到耶和華用旋風去接以利亞.
升天的時候將到.
於是兩使徒由以法連山區的吉甲.
一路步行去到伯特利.
然後由伯特利經過耶利哥.

$^{81}$去到約旦河.
以利亞在行程裡面.
有三次要求以利亞留在原地去等他.
在第二 第四和第六節那裡.
都有提及到 但是每一次.
以利亞都堅決地表示.
他說 我指著永生的耶和華.
又指著你的聖命起誓.
我必不離開你.
以利亞的那個起誓.
表示他很堅決 很堅定.
表示他甘願去承擔起誓的責任和後果.
為什麼以利亞要堅持跟隨以利亞.
去伯特利 去耶利哥 去約旦河呢?.
其實很顯然是因為.
他知道神很快將他的師父接去了.
他很珍惜和他的師父這段時間.
也很想繼續謙卑地跟隨他的師父學習.
以利亞很快就要被接去了.
這個不單止是以利亞.
以利沙也知道.
經文講到.
連在伯特利 在耶利哥的先知門徒.
他們都很清楚知道這件事.
當以利亞和以利沙來到伯特利的時候.
在伯特利的先知門徒出現.
見到以利沙 他就跟他說.
耶和華今天要接你的師父離開你.
你知不知道?.
其後兩使徒來到耶利哥.
當地的先知門徒同樣問一樣的問題.
想要提醒以利沙.
關於以利亞很快會被接去的事.
以利沙兩次都同樣回答他們的問題.
他說 我知道 你們不要作聲.
相信以利沙當時的心情很複雜.
他很傷心.
因為師父要很快離開他.
他的心情也很不捨.
所以這句話大概就是說.

$^{121}$我記得了 你們不用再提醒我了.
先知門徒這兩句問他的說話.
就顯示神當日要將以利亞接去.
原來已經成為一個公開的秘密.
為什麼大家都知道這個公開的秘密呢?.
其實這個也不是什麼特別的公開秘密.
因為事情不單止以利亞 以利沙知道.
其實以利亞早就已經跟他們說過.
這個消息也是告訴先知門徒知道的.
大概以利亞也同時宣告.
以利沙要在他被接聖天之後.
繼承他的領導地位.
以及繼承他作為先知的事奉.
以利亞每一次經過一個地方的時候.
以利亞就叫以利沙停下來.
目的大概是要試驗一下以利沙.
是否一直有這個警覺 警覺他自己會離開.
徒弟以利沙在這方面其實也很警覺的.
當時也反映到他對於神給他這個事奉的堅定.
隨時後命去預備去回應神給他的呼召.
記得我讀神學的時候.
我就要去堂會實習.
有位牧者他跟我說.
就跟這位準傳道人去說.
他說如果對真理的事不去堅持的話.
並且去寬鬆去處理的話.
就是說對其他人是一個殘忍來的.
為什麼會這樣說呢 意思就是說.
其他人可能是有樣學樣.
以為這樣做是對的.
會令到其他人很迷失.
甚至乎日後他們會跌倒.
所以我們對於真理一定要堅持 要堅定.
神是興邁不得的.
我們看一段經文 《加拿大書》六章七至八節.
我們一起讀這段經文 預備起.
「不要自欺 神是興邁不得的.
因為人種的是什麼 收的也是什麼.
信著肉體殺種的 必從肉體收敗壞.
信著聖靈殺種的 必從聖靈收永生」.

$^{161}$如果教會 信仰圈.
服侍環境裡面有真理的偏差的時候.
跟錯了人 而不是跟隨神.
而且又沒有人去指證出來的話.
那這個還是不是一個屬靈的群體呢?.
真正的拆毀是什麼呢?.
就是在屬靈上有錯誤一直存在.
但是沒有人去指證出來.
就好像有一間房子 有一個很大的裂縫.
住在裡面的人天天都看到.
他們習慣了.
依然那個裂縫是存在的.
而且也是慢慢慢慢越來越大.
可能還會影響到這個結構的核心.
領袖的角色發現了屬靈上的錯誤.
他就要去處理了.
而不是說要保護這個犯了罪的肢體.
大家假裝在那裡睡覺.
大家一起叫不醒.
只會令這件事沒有出路.
繼續陷入罪的當中.
反教會每一次見面的時候.
不是說你好我好 大家好 不是的.
我們是要背起十字架.
面對真實的自己 要付出代價的.
GSO為了我們也要付上一個很大的代價.
這份恩典是一個重大的恩典.
所以鼓勵大家.
在我們每天去禱告的時候.
尤其是晚上禱告的時候.
我們去反省自己當天所做的事.
是不是通過陽光的測試.
我們這件事是不需要在黑夜裡才去做的.
可以光明正大地去做.
信主基督徒也深信聖靈會提醒我們.
制止我們不要做得罪神的事.
而真正的建造是什麼呢?.
是要讓人有短暫的不舒服.
但是方向是校正的.
各位親愛的領袖.

$^{201}$我們要將羊帶到神的面前.
叫他們堅持去跟從神.
和我們一起去看聖經.
一起去明白真理.
去守望我們的羊.
而不是叫他們跟隨自己.
或者是錯誤地跟從自己之後.
也不去糾正他們 不是這樣的.
也不是叫羊跟隨牧者.
或者跟隨區長 組長 不是這樣的.
因為跟從人真的很危險.
因為人會犯錯.
人是不可靠的.
所以當人一旦走錯路的時候.
羊會怎樣呢?.
當然牠們會迷失方向.
會跌倒 也不知道怎麼辦.
所以各位羊仔 坐在這裡的羊仔.
羊仔我們要跟的.
一定是跟一個最大的 最厲害的.
一個是創造天地萬物的主.
大家確實明白 清楚.
應該知道跟哪一個了 是不是?.
好 我們看第二個重點.
相配的靈 神的應許.
在這往幾下 二章七至十二節.
當以利亞和以利沙繼續.
前行到約旦河的時候.
有五十個先知門徒一同跟著去.
遠遠的站在他們的對面.
相信先知門徒要跟著他們兩使徒.
目的是要看看究竟有什麼事會發生.
他們兩使徒來到約旦河邊的時候.
雖然沒有渡海輪或者山板那些快艇.
不過以利亞是有他的辦法.
他拿著自己的外衣.
將它捲起來 然後擊打河水.
河水就分開了.
以利亞和以利沙就在這個岡地裡面.
就走過去 去到約旦河的另外一個岸邊.

$^{241}$這個水分開的神蹟也讓我們記起.
在摩西帶領以色列人出埃及的時候.
摩西同樣洗到紅海的水左右分開.
以色列人在岡地裡面走過去.
以利亞現在也是選擇這樣做.
相信是要讓以利沙知道.
神從摩西開始.
在以色列人中間所做的揀選和拯救.
這些工作是沒有間斷過的.
神的作為是繼續在以利亞的事奉中延續著.
以利沙大概也將師父以利亞所做的一切事.
都記在心裡.
一會兒當他要從約旦河.
回去另外一邊岸的時候.
他也是將約旦河的水分開走過去.
以利亞和以利沙由岡地那裡走過去之後.
以利沙還記得神很快要將他的師父接走了.
所以這個徒弟寸步不離地跟著他的師父.
以利沙總算這次通過了以利亞的考試.
以利亞接著就問以利沙.
「我未被接去離開你以前.
你要我為你做什麼 只管求」.
這個「只管求」不是說以利亞有求必應.
意思是儘管開心見誠地說出來.
看看我能不能為你做些什麼.
事實上接下來我們看經文.
說到以利沙所求的以利亞.
是有條件地回答他的.
以利沙馬上回答以利亞.
「願感動你的靈 雙倍感動我」.
可以見到他實在想過師父的問題.
他之前已經想過了.
徒弟早就知道師父要去被接.
也知道師父要將先知的福事傳給他.
他見到自己未來的困難很大 很大的挑戰.
以利沙更知道他需要雙倍以利亞擁有的屬靈能力.
才可以面對 才可以應付得到.
確實以利沙日後所行過的神蹟.
《按聖經》記載是多過以利亞所做的.
有學者統計過.

$^{281}$以利亞行過的神蹟有八個.
而以利沙承接了他施父的屬靈能力.
他雙倍行了十六個神蹟.
以利亞如何回答他的徒弟以利沙呢?.
以利亞說 你求的是一件難事.
《聖經》不是說在神裡面沒有難成的事嗎?.
而事實上神賜給以利亞的屬靈能力.
或者他的徒弟神賜給以利沙之後的屬靈能力.
不是雙倍的嗎?.
大家可以看得到 也與他同在.
為何以利亞這樣說 這是一件難事呢?.
究竟以利沙所求的有多難呢?.
當我們反思我們服侍上帝的時候.
可能都會感到很難 或者力不從心.
很多時候出來的果效.
跟我們想像有很大的落差 不似預期.
有人會失望地敗退 甚至會放棄服侍.
但有些人會求主賜給他更大的智慧.
和更大的能力去應付.
忠心地去跟從主.
忠心地去服侍主 和神的差異就這樣跟著.
我不知道大家是屬於哪一類型的服侍者呢?.
我們經常都會出去傳福音.
做一些關懷事工的工作.
探望我們年長的弟兄姊妹.
有些可能住在家居 或者有些住在護老院.
陪陪他們 期望他們能夠 我們的探訪.
可以給他們一點溫暖 去令他們有所改變.
但有些時候我們都會給一些關懷者.
那些態度有些冷漠.
可能甚至拒絕 你不用來探望我了.
OK 不用來了.
可能我們都會感到心灰意冷.
覺得自己所付出的 好像都沒有什麼回報.
上兩個星期 2月8日.
我們有兩個成人牧區 跨區去做服侍.
在社區裡面寫「揮春」 街頭報道.
當天一開始的時候.
有一位機構的牧師.
他分享如何使用手上的揮春去傳福音.

$^{321}$因為揮春前面是祝福的語句.
後面是五色珠裡面的內容.
令我們有很大的鼓勵.
他說 和別人去傳福音 分享福音的時候.
神會賜給我們 有雙倍的福氣.
就是Double Blessing.
第一個福氣就是神賜給傳福音的那個人.
也就是說我們.
第二個福氣是神賜給聽福音.
又接受的人.
這樣就真是很好.
如果對方拒絕了我們所傳的福音.
第二個福氣就會回到我們身上.
就是說我們有雙倍的福氣.
繼而我們就有力量去傳福音.
將神的福氣繼續送到我們街坊的生活裡面.
那天大概有40位弟兄姊妹一起出隊.
每位弟兄姊妹都很積極地去傳福音.
一點都沒有灰心過.
我們使用了剛才所說的五色珠.
區長幫我們統計過.
大概當天我們派了接近有福字的揮春.
有印製揮春 有寫的揮春.
都是接近有一千張的.
也有170位街坊願意停下.
聽我們講整傳五色珠的福音.
並且都是願意讓我們給他們做一個祝福的祈禱.
如果他們有決志的.
都為他們做一個決志的祈禱.
當我們服侍上帝的時候.
感到力不從心的時候.
或者效果不似預期的時候.
其實是很常見的.
這些經歷是反映出.
我們屬靈的生命成長過程當中.
要面對挑戰,挫敗,真實的狀況.
重要的就是我們要保持一個信心.
持續地禱告.
依靠神的恩典和帶領.
讓神可以在我們身上施一個恩典.

$^{361}$神是著重我們的心智和我們的忠心.
而不是只看我們外表所服侍的情況或成績.
神會看到我們有心意和持續的愛心.
我們能夠透過服侍.
自己也會持續地學習,忍耐.
和有無條件的愛.
以彌撒這一份「我必不離開你」的堅持.
也同樣體驗到對神的應許的抓緊.
讓我們反思到自己.
是否也具有屬靈的心智呢?.
在服侍當中 是否也一起經歷到神的能力.
以及如何去接續為這個信仰傳承去努力呢?.
當我有時要開祖去探祖的時候.
聽到一些退休的祖源這樣說.
「我已經退休了」.
「已經沒有什麼寄聖」.
「開始看聖經的時候」.
「我就很睏 要打瞌睡」.
其實我心裡在想.
其實很多學生一打開這本書.
都是一樣會感到睏的.
不關年齡退不退休的事.
也不一定是因為要看聖經才會睏.
這位退休的祖源不單這樣說.
他還說「我紙也不識多隻」.
「上去看聖經的時候」.
「有很多不明白」.
「很多不理解的地方」.
「沒辦法了 都是要這樣了」.
其他的祖源就去鼓勵他.
他就跟他說.
「先不要說不行 不可以」.
「因為如果這樣說」.
「就是說還沒開始的時候」.
「就已經認定自己是不行的了」.
「我很同意這位祖源的鼓勵」.
「因為我們真的不要抗拒自己」.
「因為在神是凡事都能的」.
「神會賜給我們有足夠的能力」.
「去給每一位追求真理的頂子妹」.

$^{401}$「當然聖靈也會指示我們去明白聖經的」.
雖然以利沙所求的事.
對以利亞來說是有點難.
但神也是透過以利亞對以利沙作出指示.
祂說「我被接去離開你的時候」.
「你若看見我 就必得著」.
「若不然 就得不著了」.
這就是說.
如果看見以利亞被接去的情境.
以利沙所求加倍的屬靈能力.
就必得著了.
若然 即是說如果.
如果看不見以利亞被接去的時候.
以利沙所求的就得不著.
但我們會問.
以利沙不是已經跟著他的師父去到約旦河邊了嗎?.
他已經寸步不離了.
以利沙還有可能看不見他的師父被接走的情境嗎?.
神要以利亞回答以利沙.
其實顯然也是要告訴以利沙.
他先前一直這麼警醒地重視.
神所託付給他的即時.
神其實已經預立了他.
願意加倍賜予以利沙屬靈的能力.
來承擔他的事奉.
但為了使得以利沙具體地知道神的心意.
神藉著以利沙看見以利亞升天的經歷.
告訴以利沙 神必定與他同在.
賜予這個屬靈的能力雙倍臨到以利沙身上.
但我們要想深一層.
就算以利沙跟著以利亞來到約旦河邊.
以利沙也不一定可以看見以利亞被接走的情境.
記得在後面的經文裡.
阿難王派人到多丹去捉拿以利沙的時候.
以利沙的僕人看不見以利沙所見到的滿山火車火馬.
以利沙求神開他僕人的眼睛.
讓他僕人才可以看見.
換言之 以利沙之所以可以看見.
他跟以利亞分開的火車火馬.
又看見以利亞被旋風接走的時候.

$^{441}$並不是單單因為以利沙本人在現場.
而是完全因為神要使得以利沙看見這個情景.
而那五十位先知們徒.
相信他們也看不見以利亞被接去的情景.
以利沙看見以利亞被接去.
其實是一個證據.
是神要向以利沙表明.
神願意給他雙倍的屬靈能力的證據.
經文提到有火車火馬將兩人分開.
不論這是一個真實有戰車戰馬放出火焰.
或者是一個異象也好.
這些景象以利沙全部都看得見.
亦都交代了以利沙將會得到神的應許.
將這個屬靈的能力雙倍降臨在他身上.
我們看看第三個重點.
得著能力 繼承使命.
在遼王紀下 二章十三至十八節.
在十三,十四節.
以利沙被旋風接走的時候.
他的外衣就掉下來.
而且這個外衣掉下去跟以利沙不太遠.
所以以利沙就撿起師父以利亞的外衣.
這個舉動亦表示.
以利沙願意繼承以利亞先知的事奉.
在遼王紀上十九章剛才也提及過.
以利亞受神的託付 去告以利沙為先知.
去接續先知的事奉的時候.
以利亞就將他的外衣搭在以利沙肩膊上.
外衣亦與傳承有關.
這個一個新的以利沙就呈現出來了.
他以接續以利亞作先知的身份出現了.
所以如果有屬靈的長輩走過來.
將他的外衣搭在你的肩膊上的時候.
大家想一想是否有點意思.
是否有點古怪呢?.
當以利沙拿著這個外衣的時候.
耶和華說:以利亞的神在哪裡呢?.
這句話是一個禱告.
祈求神向他的僕人顯示他的能力.
而且這個能力是從他的師父的神耶和華而來的.

$^{481}$直至他證明他承接承繼以利亞的職務.
果然當以利沙一打這個水的時候.
同樣這個水也會分開.
以利沙就走過去了.
回顧以色列人的舊恩歷史.
我們曾經也經歷過水分開.
人走在旱地裡的神蹟.
莫西大領以色列人過紅海.
約書亞也是大領以色列人走過約旦河.
以利亞擊打這個水.
就好像這個舉動是鼓勵以利沙.
上帝要將先知的工作.
屬靈的能力在以利沙身上延續下去.
而外衣更加是上帝給過以利沙接觸新的挑戰的一個記號.
當先知門徒看到以利沙在約旦河走過去的時候.
就知道耶和華和以利沙同在了.
原來在約旦河.
即原先在約旦河的50位先知門徒.
他們很遠的望著以利亞和以利沙去過對面岸.
現在只見到以利沙一個人回來.
以利亞不見了.
但當他們見到以利沙用外衣.
將河水分開的時候走過來.
他們就知道了.
也承認感動以利亞的靈已經臨到以利沙的身上了.
也表示以利沙要接觸以利亞的地位.
以後這班先知門徒就要被以利沙去照顧.
於是他們迎接以利沙 俯伏在地上.
表示他們心悅誠服地接納以利沙成為領導的地位.
十六至十八節大概由於其他人看不見以利亞被接去.
所以單憑以利沙一個人說的話.
這班先知門徒對以利亞未能忘記.
或者完全未能消化這件事.
所以建議以利沙派人去看看以利亞.
是否真的不在這裡.
最終是怎樣呢?.
他們在山谷裡仔細地尋找了三天.
三天尋找過了.
都不見 徒勞無功.
這就表示以利亞已經被接去.

$^{521}$也肯定以利沙對他們之前所說的話是真實的.
而且是帶有權威的.
以利沙不僅承接了師父以利亞先知的衣缽和能力.
更加帶著雙倍的屬靈能力去服侍上帝.
我們反思以利沙對於師父的離開.
他沒有因為這樣而迷失.
因為他知道自己的身份.
也知道這個份他要在裡面盡心盡力地愛主.
做主喜悅的事.
繼承師父的使命.
同時也知道神必定會與他同在.
神會給他不同的印證.
我們去完成一個任務.
不要說是工作上的任務.
可能是學習上的任務.
或者是我們服侍上的任務.
原來也有分開三個層次.
第一個就是做了一件事.
第二個就是做到了一件事.
第三個就是做好一件事.
做了一件事是什麼意思呢?.
這件事基本上已經完成了.
有人叫我去做 我就去做.
未必會用心去完成這件事.
為了做 我要應付.
好像之前港鐵沙中線.
紅隧站的月台工程.
為了追趕完成的時間.
鋼筋剪短了.
引發質量的問題.
這個就是打算做了一件事.
交代了 完成了.
達到的期限就算了 交差.
做到一件事是什麼意思呢?.
表示一件事成功完成了.
並且已經達到了目標.
結果也出來了.
可能也有盡力去做.
但未必用心去做.
我記得我讀書的時候.

$^{561}$放暑假要做暑期作業.
通常一放假.
我要回去馬上去做.
絕對不會這樣的.
一直玩著玩著.
玩到差不多要有新的開學了.
一兩個星期之前我才開始做.
當然我比較馬虎地去做.
快點講完它.
是做了一件事.
但我完成了這個功課.
但我沒有用心去做.
當然也在當中學習不到任何東西.
做好一件事是什麼呢?.
就是說不單是盡力去做.
而且是用心去完成的.
就好像媽媽做一頓飯出來.
她也會用盡她的心思.
用盡她的力量去準備.
有一個很完善的一頓飯的品質出來.
希望家人可以享受到一個豐富的晚餐.
期望各位服侍上帝的僕人.
不單做了和做到一件事.
去交差.
和要做到神交給我們的任務.
同樣我們也要做好這件事.
做一個又良善又忠心.
而且做一個無愧的工人去服侍神.
以利沙在聖經記載裡.
他行過十六個神蹟.
多數和日常生活裡的供應有關.
有醫治的關係.
甚至乎以利沙死了之後.
這個神蹟繼續有出現.
以利沙行事是較為溫和.
關注醫治,恩惠.
彰顯了上帝的慈愛和豐富.
跟著他施父所做的神蹟.
是有點相似.
譬如使人復活.

$^{601}$將水分開.
都有點相似.
亦都是師父給徒弟的一個榜樣.
但師父做的神蹟.
有些是較為剛烈的.
是有一點點不同的.
神用每一位的僕人.
不是因為他們像誰.
而是因為他們很忠心.
將神放在僕人心裡的負擔.
每一位的屬靈前輩.
我們都是值得去欣賞.
值得去尊重.
他們的生命.
有些會很溫柔地去服侍.
有些可能會比較剛烈.
有些會比較安靜.
有些是默默地去付出.
這些都是我們學習的榜樣.
繼承使命真正的意思.
是要將神交託給人的使命.
去薪火相傳.
繼續延典下去.
一個傳一個.
生命影響生命.
使得服侍的心不會熄滅.
當教會要作出領袖傳承的時候.
通常都會產生一些不同的領袖風格.
領導的風格.
因為領袖有不同的屬靈的恩賜和才能.
按神的心意去賜予他們的恩賜.
互相去配搭.
互相去服侍.
只要我們忠心堅定地去跟從神.
祝福神賜給我們的應許.
神必定賜給我們加倍的恩告.
並且神會在新一代的領袖身上.
持續地去動功.
來彰顯神的大能.
接觸和神的計劃.

$^{641}$這樣神的家就會得以繼續向前行.
一代一代地去傳承下去.
在社區裡面作主的登台.
作主美好的見證.
我們一起同心低頭禱告.
愛我們的主.
求你至下恩高和屬靈的能力.
使我們能夠忠心地去服侍你.
成為一個無愧的工人.
彰顯你奇妙的恩典和能力.
又讓我們的生命能夠成為流通的管子.
將服侍的棒一代一代地傳承下去.
讓「筲箕灣浸信會」成為明燈.
繼續去見證主的明.
主求你幫助我們.
好像以利沙一樣.
心貼心,緊緊地去跟從你.
無論前面有多少挑戰或困難.
我們都不輕易地說放棄.
我們要對你有信心.
我們要堅定地去依靠你的力量.
讓你的靈在我們身上運行.
使我們能夠行出你的旨意.
也求你賜給我們有智慧和勇氣.
承接屬天的使命.
將你的薪火一代一代地傳承下去.
使更多人因為你的恩典.
得著拯救,得著生命的更生.
將榮耀總贊歸還給我們三一的真實.
我們祈禱奉主耶穌基督.
得勝名自祈求.
阿們.
\newpage



\section{馬太福音 5:1-6}
\label{sec:031711}
\textbf{有福了!屬神的人 (1)}~~~~陳雅莉傳道
\newline
\newline
連結: \href{http://www.church.com.hk/acms/content.asp?site=cdc\&op=show\&type=product\&code=031711\&layout=sermon}{\texttt{031711}} ~~~~ 語音日期: 2026-02-22
\newline
\newline
\hyperref[sec:031710]{\small{< < < PREV SERMON < < <}}
~
\hyperref[sec:index_chronic]{\small{[返順時目]}}
~
\hyperref[sec:031712]{\small{> > > NEXT SERMON > > >}}
\newline
\newline
馬太福音 5:1-6
\newline
\begin{longtable}{cl}
\hline
\hline
章節 & 經文 (和合本修訂版)\\
\hline
5:1 & \begin{tabularx}{0.7\textwidth}{X} 耶穌看見這一群人，就上了山，坐下後，門徒到他跟前來， \end{tabularx} \\ \\ \relax
5:2 & \begin{tabularx}{0.7\textwidth}{X} 他開口教導他們說： \end{tabularx} \\ \\ \relax
5:3 & \begin{tabularx}{0.7\textwidth}{X} 「心靈貧窮的人有福了！因為天國是他們的。 \end{tabularx} \\ \\ \relax
5:4 & \begin{tabularx}{0.7\textwidth}{X} 哀慟的人有福了！因為他們必得安慰。 \end{tabularx} \\ \\ \relax
5:5 & \begin{tabularx}{0.7\textwidth}{X} 謙和的人有福了！因為他們必承受土地。 \end{tabularx} \\ \\ \relax
5:6 & \begin{tabularx}{0.7\textwidth}{X} 飢渴慕義的人有福了！因為他們必得飽足。 \end{tabularx} \\ \\
[1ex]
\hline
\hline
\end{longtable}
$^{1}$請姐妹平安.
在新幾年再一次恭祝大家.
主恩陪添 福氣延綿.
今天呢 剛才主席帶大家讀這段經文的時候.
你們有沒有覺得跟平時有點不同啊?.
他讀得很有力啊 你們有沒有覺得啊?.
可不可以回到上一頁的經文那裡?.
他帶我們讀的時候.
他是「謙和心靈貧窮的人有福了.
因為天國是他們的」.
他是怎麼讀的呢?.
「心靈貧窮的人有福了」.
啊 真是太好了.
主席一開始就已經幫我點出了.
是的 確實是「有福了」.
這個是我今天的主題.
「有福了 屬神的人」.
因為在原來的文字.
的確是真的這樣的.
是先說「Blessed 有福了」.
接著就是「心靈貧窮的人」.
「有福了」是什麼人?.
盼望我們今天.
你回到來.
雖然這段經文可能你已經很熟悉了 我相信.
甚至我相信當中我們有不少識輩.
就算你是未信主的人.
可能你都在不同的地方見過這些經文的了.
但是仍然帶著一顆謙卑的心.
盼望得著 祝福的心 好嗎?.
你要相信「屬神的人」是「有福了」.
你是「有福了」的 好嗎?.
新年的時候 我們總是會說.
迎春 接福.
我們都很喜歡「祝福」.
早上回來 見到大家穿得很漂亮.
很開心 都會彼此「祝福」.
用很多問候的說話 祝福的說話.
很希望大家在新的一年得著滿滿的福氣.
是的 我們農曆新年.

$^{41}$我們都很喜歡這樣彼此「祝福」.
亦都很渴望神賜福我們新的一年.
今天的經文 我相信雖然不陌生.
但是裡面蘊含豐富的意義.
所以為什麼你會見到「屬神的人」「上」這樣.
因為我覺得我預備的時候實在是豐富得太厲害了.
一堂根本講不完.
下個禮拜仍然是農曆新年.
所以仍然很想大家繼續一齊來接這個福.
所以我將那個經文八福.
斬為頭四福和尾四福.
如果你想看完尾的四福.
你下個禮拜一定要回來.
有下集的.
好 進入今天的經文.
《馬太福音》來到第五章.
我們今天讀的經文.
這是一個我們常說的「登山寶訓」.
耶穌在一個山上 或者一個平原那裡.
這樣去講這個訊息.
而「八福」就在這個登山寶訓的一開頭.
聖經就說 耶穌坐下就教訓他們 教導他們.
好 耶穌就見到很多的人.
他上了山 坐下來了.
就在門徒 到他跟前來.
他就開口教導他們.
所以你會看到.
其實《馬太》寫這一段經文的時候.
都是有刻意的鋪排的.
如果我們回頭看一點點 你會發現.
其實第四章正正就是耶穌呼召了門徒來跟從他.
那會不會就是《馬太》刻意的編排.
都讓我們知道 當呼召了你做門徒.
說你是耶穌的門徒.
你要跟從耶穌之後.
那接著你要做些什麼呢?.
那可能這個教導.
都是給昔日的門徒.
也都相信是給今天的我和你.
願我們即使再熟悉.

$^{81}$我們仍然都帶著我們很渴慕做門徒的心.
渴慕聽主的教訓的心.
一起來聆聽 好嗎?.
那這個教導 或者原來的文字.
教訓,教導 這個字眼.
是有一種重複,重重複複.
不斷教導的意味在裡面的.
就是說 你不要嫌煩.
這個是上帝想要給你的福氣.
願我們有聆聽的心來聽祂的教導 好嗎?.
那你們剛才看過了.
白福一開始是以有福了 這樣開始的.
有福了 就是心靈貧窮的人.
有福的人是什麼人呢?.
是心靈貧窮的人.
那如果我們比較熟悉的和合本的聖經.
我們會說是 不知道大家會不會背呢?.
虛心 我看到我姐姐姐妹說.
虛心的人有福了 原來是和合本.
其實虛心很難明白的.
所以新的譯本將它翻了.
再出一點的意思就是心靈貧窮.
也是原來文字真的意思是心靈貧窮.
心靈貧窮 就不是說到我們一看到貧窮.
通常這些字我們都不太想要的 是不是?.
特別新年樓樓不要說貧窮了 是不是?.
就是不是很意圖的 是不是?.
不過這裡的心靈貧窮不是說什麼都沒有.
很虛空 看到自己很低落.
它這個貧窮是在形容一份需要.
是在形容一份你心靈裡面有一個極大的需要.
聖經所用的這個字眼.
的而且確在原來是用來形容赤貧.
不是貧這麼簡單 是赤貧.
是一無所有 完全什麼都沒有.
這樣的字眼.
不過我們知道 它這裡說的不是生活貧窮.
它是在說心靈的貧窮.
當一個人感覺到 自覺到自己的心靈貧窮.
覺得自己心靈裡面有極大的需要.

$^{121}$有極大的渴望.
並且感覺到不可以靠自己.
我是無依無靠的.
要在上帝面前承認自己一無所有.
這個是心靈貧窮.
就好像有人形容 好像一個黑蓋.
一個黑蓋在上帝面前.
承認自己什麼都沒有.
Nothing 什麼都沒有.
除了上帝 別無依靠的這種狀態.
是關乎一個人在上帝面前怎樣看待自己.
上帝說 你如果自認你是一個心靈貧窮的人.
你有福了.
在英文版 有一個英文版的聖經.
叫做《The Message》 信息版的聖經.
它是這樣譯的.
它說心靈貧窮的人有福了.
它就譯為 當你走投無路的時候 你有福了.
很顛覆吧 走投無路 怎會覺得自己有福呢?.
但聖經是這樣說的 耶穌是這樣說的.
當你覺得自己走投無路的時候 你是有福了.
因為 為什麼呢?.
因為天國是他們的.
在這個英文版的聖經 它譯為.
當你感覺到你自己越少 上帝就越多.
祂的掌權 祂的rule在你身上也越明顯.
有福不是因為你什麼都沒有.
有福是因為當你什麼都沒有的時候.
上帝在你裡面很大.
這個是有福了.
你越少 上帝就越多.
當人能夠看見自己所領受的.
根本不是自己努力所得回來的.
都是你都要努力.
不過你再想深一層.
一切都是從神所領受的.
沒有可誇.
當人謙卑下來.
心靈貧窮到一個地步.
覺得自己沒有辦法靠自己.

$^{161}$只能夠單單依靠這位上帝.
你願意讓上帝居首位.
No.1 讓上帝第一.
而不是讓自己的自我放大.
這樣的人 耶穌說有福了.
有人這樣形容心靈貧窮.
或者虛心 是什麼呢?.
是在上帝面前無求.
無知 無有.
什麼都沒有 還有福.
無求.
不是真的說什麼都不要了.
無欲無求.
不要吃好東西 不要去街上.
不要穿衣服 不是這樣.
不是禁欲主義.
無求是在神面前完全放下你自己的心意.
無求是不為自己去求一點.
甚至乎不帶自己的企圖.
只是願意實踐神的旨意.
這是無求.
無知呢?.
我們經常說不要這麼無知.
這個無知不是說什麼都不知道.
很沒智慧 很蠢 很笨.
不是這個意思.
這個無知是即使你不知道.
上帝的旨意要怎麼走.
你不知道神將要在你身上實踐什麼計劃.
你不知道.
但是你仍然願意選擇信服.
你仍然願意留心神在你身上的作為.
這一種無知.
還有無有.
不是說買了房子 買了房子 買了車 什麼都沒有.
而是即使你擁有一切.
但是你認定這一切沒有一樣東西是屬於你的.
是上帝賜給你的.
你只是代為管理它.
而主若許可.

$^{201}$我們就全感恩的心.
去領受這些上帝給我們的福氣.
這個是無有.
如果人沒有由衷的謙卑.
沒有由衷自願首先降卑.
沒有凡事上先看到神顯為大.
只有表面的謙卑.
可能那些都是禮儀.
其實某程度上可能是一種偽善.
唯有真正的無求.
無知.
無有.
心明自己只不過是一個罪人.
我們何等需要神的憐憫.
憐恤.
何等需要神的恩典.
這是心靈貧窮.
是有福的.
「百福以心靈貧窮為起頭」.
有人就想 其實可不可以其他東西做起頭呢?.
我想你反過來背.
那也沒問題的.
都是有福的.
你無論抽任何一句出來講.
都是有福的.
不過似乎作者或者編輯者.
是刻意將心靈貧窮.
是有福的放在第一.
為什麼呢?.
因為如果人不首先自覺靈性貧乏.
不首先認定需要上帝.
不首先認定自己根本一無所有.
只能夠靠主.
這個世界上沒有一個人能夠信主.
唯有當你真正看見自己是何等需要主.
你悔改自認自己是一無所有.
在神面前我只不過是個罪人.
唯有這樣為起頭.
你才能夠真認識主.
唯有人看到自己的盡頭.

$^{241}$才會看到神的起頭.
這個是真正屬靈的福氣.
甚至有人說.
心靈貧窮這裡說的是.
你要甚至承認你自己破產.
承認自己在神面前破產.
反而能夠進到神的國去.
虛心的人.
心靈貧窮的人有福念.
天國是他們的.
你可以進入神的國.
接著耶穌就說到.
哀動的人有福念.
因為他們必得安慰.
哀動我們一聽到.
都會覺得是一些不太好的東西.
哭,悲傷,流淚,難過,傷心等等.
這些就是哀動.
不太想的.
最好不要有.
但是我們會發現.
哀動在原來的意思.
更加有一種隱藏,深切的悲痛.
不單是情緒上覺得難過.
其實某程度上是一種態度上的示弱.
在承認自己沒有力.
沒有力控制一切的事.
完全無助.
這一種哀動的表現.
我們可能會想.
難道耶穌會喜歡我們愁眉苦臉嗎?.
不是的.
祂常常叫我們賞賞喜樂.
我們祝福別人的時候.
都是說喜樂平安.
沒有理由祝福別人哀動.
是不是?.
是啊.
耶穌為什麼會說哀動的人是有福的?.
因為我們會發現.

$^{281}$在人生裡面根本沒有辦法.
你根本沒有辦法拿走哀動.
在我們的人生裡面.
有很多事情.
都會讓我們遇上哀動.
可能是為著世界所發生的事.
你每一天看到的新聞.
你身邊的人發生一些難過的事.
他們遇到不好的際遇.
或者是你自己.
無論生老病死,悲歡離合.
我們是沒有一個人能夠跳出.
這些人生的際遇.
生活困苦壓迫.
有失去心愛的事的時候.
這些都是人生一部分.
而且人人都必須經歷的.
傳道書說.
哭有時.
笑有時.
哀動有時.
跳舞有時.
可見人生根本沒有辦法.
永遠都沒有哀動.
生離死別會讓人難過.
悲歡離合.
孤單寂寞會讓人傷心.
我們甚至不是自己.
我們看到身邊的人受苦.
我們都會流下關懷和憐憫的眼淚.
這都是人之常情.
這種悲哀.
都是我們必須經歷.
無法逃避的.
但我又很感恩.
上帝說.
患懂得哀動的人有福了.
你沒有麻木到.
你還是一個活生生的人.
你有福了.

$^{321}$你沒有看見一切的事.
等於沒有看見.
你沒有不顯出憐憫.
不顯出關懷.
沒有不去同理別人.
你還未麻木.
還會哀動.
有福了.
我們可以確定.
如果沒有哀動.
就不需要安慰.
是不是?.
因著哀動.
我們能夠得著更大的.
就是蒙主的安慰.
經歷哀動的時候.
在哀動裡得主的安慰.
英文版.
這個英文信息版的聖經.
它說當你覺得自己.
失去最珍貴的東西的時候.
你是有福的.
只有這樣.
你才能夠被那個對你是最珍貴的擁抱你.
我們很難覺得.
我們失去最珍貴和最深愛的東西的時候.
會覺得自己有福.
但如果當你看到.
當你失去這一切的時候.
最珍貴.
對你來說最珍貴的東西.
會擁抱你.
那個最珍貴的東西是什麼?.
是我們的主耶穌.
唯有當你落在.
這個光景裡面.
你會深深感受主的安慰.
聖經應許必得安慰.
只要你帶著你的哀動.
來到主的面前.

$^{361}$是必得安慰.
你會擁有這個.
擁有這一份.
主給你最珍貴的禮物.
當我們難過,悲傷的時候.
就是我們最深切感受.
我們的天父願我們同在的時候.
詩人說.
我們流的每一滴眼淚.
在詩篇裡面.
詩人寫說.
我們流的每一滴眼淚.
神都用皮袋裝起來.
每一滴.
神都一一紀念.
神與你同在.
近日不知道大家有沒有看冬季奧運會呢?.
或者有些人說.
在傳媒上都看到很多這個信息.
被洗版了.
在冬季奧運會裡面.
當然是比賽.
奧運.
但是上演了異常動人的一幕.
這一幕是什麼呢?.
是一隊花式滑冰的日本隊.
日本的隊伍.
一隊雙人花式的隊伍.
在一個叫做短曲的項目裡面.
出現了異常,不尋常的失誤.
奧運運動員.
而且他們本來已經一直是世界冠軍.
他當時在短曲的項目裡面.
這個木原龍一.
他因為小小的失誤.
都不是小小的.
在奧運裡面極微小的失誤.
都是極大的失誤.
他沒有準確地接住他的搭檔.
沒有準確地接住三葡萄萊.

$^{401}$雖然三葡萄萊最後沒有受傷.
但是這一下的失誤.
足以令他們看來沒有辦法爭取世界冠軍.
就是在一結束這個項目的時候.
這個運動員龍一.
他跪在冰地上面嚎哭.
他沒有辦法接受自己的失誤.
很難過 很難過.
他不想他的伴侶.
他的同伴失誤而造成傷害.
雖然最終沒有.
但是他沒有辦法接受自己的失誤.
正當龍一跪在冰地上面.
全球的鏡頭都拍攝著他大哭的時候.
尼萊不僅沒有埋怨他 覺得他令他失去了一面金牌.
尼萊不僅沒有鄙視他 責備他 看不起他.
尼萊選擇了馬上上前安慰他.
擁抱他 鼓勵他 再一次堅固他.
告訴他 傳媒訪問的時候.
尼萊甚至乎說 我會願意再為龍一再畫一次.
他不僅沒有責備 沒有埋怨.
影響了他奪金的機會.
相反 他擁抱他的同伴.
他看見了龍一的哀動.
他哀動龍一所哀動的.
在第二天 還有一個項目叫「長曲」的項目.
其實是第一天就發生了失誤.
第二天就要再比賽了 計算過總分.
結果在第二天再比賽「長曲」的項目的時候.
他們竟然突破了驕人的成績.
破了世界紀錄的成績.
從來沒有人拿過這麼高的分數.
結果他們不僅追回了所有的分數.
而且最終獲得金牌.
我想說的是 得金牌固然很快樂.
但是我們能夠明白.
那些技巧,練習,體力固然是奪金的條件.
但是這一次能夠奪金最重要的因素.
是在失誤和失敗當中.
尼萊願意擁抱龍一的哀動.

$^{441}$看見他的傷心難過.
願意安慰他的心.
使他能夠重新振作.
這才是致勝的關鍵.
今天 我不知道我們心目中有沒有一些哀動的事.
不論是因為什麼事.
為自己,為別人,為世界.
你要相信主要安慰你.
但有時不是主不安慰我們.
是我們不肯受安慰.
我們會沉醉在自己的失誤裡面.
沉醉在自己的失敗裡面.
我們過不了的是我們自己.
只是繼續落在那個自憐和哀動裡面.
但我們沒有來到神的面前得安慰.
求主幫助我們.
你要將你的哀動帶到主的根前.
主應許你.
你必得安慰在我這裡.
你要在主那裡 你必得安慰.
所以哀動的人是有福了.
你要尋求主的安慰.
願我們肯受主的安慰.
接著 耶穌說.
謙和的人有福了.
因為他們必承受土地.
謙和總是給人一種感覺很懦弱.
很無能.
無能,懦弱的人才謙和.
但原來的文字.
或者和合本是譯作溫柔的人有福了.
謙和,溫柔不是說沒有用,沒有能力.
一無所有這麼簡單.
它是說一種剛柔合一.
原來文字是有剛柔合一的這種性格.
謙和.
這個字還可以用在哪裡呢?.
可以用在一些受過訓練.
願意信服主人.
願意接受管束.

$^{481}$這樣的字.
謙和不是說你沒有能力.
所以你要認輸.
它是說你即使大有能力.
但你願意接受神的管束.
願意信服這位真正的主人.
你願意控制自己的慾望,天性,衝動等等.
在神面前謙和下來.
是讓神.
大大讓神去管束我們自己.
所以謙和也可以說是.
對神信服上帝的旨意.
對人是一種以謙卑為基礎.
願意寬容.
願意不計較.
願意不予以報復.
哪怕即使在惡待.
甚至乎忍受侮辱的時候.
不是因為不想出聲.
不想爭吵.
而願意選擇謙和.
那是願意勇敢面對.
願意放大別人的優點.
收細對方的弱點.
不予以報復.
當神被侮辱的時候.
被人惡待我們的耶穌的時候.
祂都會發怒的.
但祂沒有自己心懷不平.
耶穌沒有為自己深深苦讀.
然後妒忌,嫉妒,嬲怒以致作惡舞.
我們可以忍受.
羞辱.
甚至不被尊重.
而不懷恨意.
只求主的公道.
只求主的公義彰顯.
這是謙和.
謙和嚴格來說.
具體化一點.

$^{521}$它是一種容量.
一個容器在你裡面很大.
你裡面有一個很大的容器.
承載得下.
當我們學校主的樣式.
主說你要來學我樣式.
耶穌說我心裡柔和謙卑.
你學校主的樣式.
你會發現.
你容得下那些人所捍衛用不下的事.
有一本小書.
我之前看過.
很多年前看過.
但我仍然很深刻.
它有一句說話.
這本書叫《宣教士的成功秘訣:放棄權利》.
在書的第一章.
它就說到.
我不記得是第一章還是序言.
已經說到了.
我一打開這本書.
沒幾句就說這句.
它說吃苦是容易的.
吃虧是難的.
當時我很多年前看《年少無知》.
我看的時候覺得.
我很不認同作者所說.
因為我覺得吃苦都已經很難.
怎會吃苦是容易的呢? 沒可能.
吃苦是難的.
吃虧是更難的.
我覺得應該是這樣.
我看完這本書.
當然我再仔細去看這本書.
當我仔細地看.
我終於明白作者在說什麼.
一個宣教士上到宣教的工場.
在宣教的工場裡面.
要讓未得知民,不認識耶穌的人.
去聽聞福音,去認識耶穌.

$^{561}$吃苦對他們來說可以說是基本.
甚至乎是必須.
他們往往更深的學習是吃虧.
如果是主.
如果是為主的緣故.
是主要我們吃虧.
是為主的緣故而吃虧.
我們還會不會覺得有福呢?.
很困難吧 很困難.
我不知道你現在是什麼狀態.
可能你看到這句你都會和我一樣.
吃苦已經夠難 最好不要吃苦.
還要我吃虧 說不過去.
但這個是謙和.
這就是聖經裡面所說的.
謙和.
你願意在上帝面前放下你應有的權利.
那些權利是你應有的.
哪怕你覺得是你應有的.
所以這本書有一個小的副題.
寫著「放棄權利」.
是的 確實是的.
願意放下你應有的權利.
不求回報.
不是求人的喜悅.
只討主的喜悅.
耶穌說當你不去為自己爭取的時候.
你不為自己去祈求的時候.
你必承受土地.
這個英文版又怎麼說呢?.
它說當你滿足於你就是你.
不多不少的時候.
那一刻你會發現自己成為一切.
沒辦法用金錢換取事物的真正主人.
是的.
在事奉的年日裡面確實.
我都覺得真的不會沒有苦澀的時候.
不要說是事奉的年日.
我相信大家都會認同.
我們人生裡面不會沒有苦澀的時候.

$^{601}$但如果你能夠安於上帝的手中.
安然接受主一切的安排.
我們會成為一個真正有福的人.
當我們願意以謙和的心.
放下那個我們自以為應得的權利.
應受的尊重.
我深信主必為你預備更美.
這是我這麼多年來我的信仰走過來的路.
當你願意伏下來.
不為自己去爭取的時候.
主必定會為你爭取.
來到今天最後的一個福.
有福的是饑渴無意的人.
饑渴無意的人有福了.
因為他們必得飽足.
一個人要感到肚餓,感到饑渴.
你才會去找東西吃.
心靈都一樣.
我們的靈命都一樣.
你要感覺到有饑渴感.
你要覺得你的靈裡面很餓.
你才會願意去尋求.
尋求什麼呢?.
饑渴你就會要無意.
我們來看看.
剛才我們說 前面饑渴無意有三個人.
有心靈貧窮的人,哀動的人,謙和的人.
你會發現作者是有刻意的編排的.
這三樣心靈貧窮,哀動,謙和.
都是一種饑渴.
都是一種倒空無有讓你謙卑你自己.
當你饑渴的時候.
你就去渴慕神的義.
神的義很廣闊.
但是在《馬太福音》這裡.
其實整個《馬太福音》第五章.
是重複出現了耶穌說.
怎麼說祂的義.
你要饑渴無意.
你要慕什麼義才行.

$^{641}$不是慕自己的義.
是慕耶穌的義.
耶穌說:我告訴你們.
如果你們的義不勝過民事和法理賽人的義.
是絕對不能夠進天國的.
當時的法理賽人覺得.
我們沒有殺人就已經是很容易了.
但是耶穌告訴你.
祂說不是.
凡向弟兄動腦的.
罵弟兄是廢物的.
說弟兄是白痴的.
都是.
都不能夠成就神的義.
求主幫助我們.
不只是履行表面的道德標準.
我沒有殺人了.
我沒有犯法了.
那是渴慕上帝真正的義.
我接近要結束了.
主比屬神國度的人的福氣.
是我們覺得很顛覆的.
是我們覺得和我們所認識的很不一樣的.
心靈貧窮.
哀動.
嫌和.
饑渴無意.
如果我們回去嘗試想像一下.
昔日在聽這個頻道.
這個訊息的人.
他們是什麼人呢?.
他們不是什麼達官貴就.
可能當中有小部份是.
但是大部份的都是.
漁夫.
農民.
牧羊.
他們本來就是貧窮.
本來就是常常哀動.
本來就是一無所有的人.

$^{681}$他們根本不需要再有人告訴他們.
什麼叫貧窮.
什麼叫要哀動.
他們根本就是經常都會哀動.
在當時的人來說.
貧窮,哀動,饑餓等等.
這些根本就是受咒詛的人.
根本就是沒有人會喜歡的.
到今天我們都不會想要這些.
但是耶穌卻說.
這樣的人.
有福了.
耶穌竟然邀請這些他想都沒有想過.
自己能夠成為不是咒詛的人.
不是社會最低層的人.
他是成為有福的人.
耶穌說這些人要成為有福的人.
今天.
耶穌也在同樣宣講這個好消息.
這個是福音.
這就是福音.
是顛覆我們所想的.
我們總覺得擁有越多才是最多福氣.
得著最多才是最好的福氣.
但是耶穌說如果你願意.
在祂面前倒空一切.
祂要將你填得滿滿.
過去的星期三我們踏入了預福期.
預福期就是.
我們在群組裡面有聲音出來.
一起去思想主耶穌基督的死而復活.
祂的捨己,祂的捨身.
我們的執事每一天都在群組裡面傳出.
一起去默想靈修的訊息.
有人說最近這兩個星期都是新年.
應該是很喜慶的.
好像很矛盾那樣.
有人會說.
怎麼可以又很喜慶.
但是又很默想受苦呢?.

$^{721}$這樣寫忌呢?.
是不是矛盾呢?.
不是.
從今天的經文可以看到.
無論你在什麼光景.
你自問是貧窮的,哀悼的,饑渴的.
主更加要與你同在.
在任何的光景.
我們都要得著這一份福氣.
再次謹祝大家.
蒙主的恩福,真正的福.
繼續在這個地上活出美善.
我們一同禱告.
你給我們的天賦.
再次感恩.
你從歲首到年終都看顧我們.
多謝主賜下屬靈的恩福.
不僅是世人所羨慕的.
更加是屬天裡面所羨慕的福氣.
求主你真的給我們.
幫助我們.
轉換我們的眼光.
羨慕這一份屬靈的福氣.
在你裡面.
成為心靈貧窮的人.
成為哀悼的人.
成為謙和的人.
成為饑渴無義的人.
因為.
這是屬靈的福氣.
這樣的人有福了.
求主幫助我們.
賜福教會.
賜福我們的弟兄姊妹.
奉耶穌基督名字祈求.
阿們.
\newpage



\section{加拉太書 1:1-3:29}
\label{sec:031712}
\textbf{加拉太書 (1) 因信稱義的真理}~~~~劉少康牧師
\newline
\newline
連結: \href{http://www.church.com.hk/acms/content.asp?site=cdc\&op=show\&type=product\&code=031712\&layout=sermon}{\texttt{031712}} ~~~~ 語音日期: 2026-02-25
\newline
\newline
\hyperref[sec:031711]{\small{< < < PREV SERMON < < <}}
~
\hyperref[sec:index_chronic]{\small{[返順時目]}}
~
\hyperref[sec:031713]{\small{> > > NEXT SERMON > > >}}
\newline
\newline
加拉太書 1:1-3:29
\newline
\begin{longtable}{cl}
\hline
\hline
章節 & 經文 (和合本修訂版)\\
\hline
1:1 & \begin{tabularx}{0.7\textwidth}{X} 我使徒保羅和所有跟我一起的弟兄，寫信給加拉太的眾教會。我作使徒不是由於人，也不是藉著人，而是藉著耶穌基督與使他從死人中復活的父神。 \end{tabularx} \\ \\ \relax
1:2 & \begin{tabularx}{0.7\textwidth}{X} 見上節 \end{tabularx} \\ \\ \relax
1:3 & \begin{tabularx}{0.7\textwidth}{X} 願恩惠、平安從我們的父神和主耶穌基督歸給你們！ \end{tabularx} \\ \\ \relax
1:4 & \begin{tabularx}{0.7\textwidth}{X} 基督照我們父神的旨意，為我們的罪捨己，要救我們脫離現今這罪惡的世代。 \end{tabularx} \\ \\ \relax
1:5 & \begin{tabularx}{0.7\textwidth}{X} 願榮耀歸給神，直到永永遠遠。阿們！ \end{tabularx} \\ \\ \relax
1:6 & \begin{tabularx}{0.7\textwidth}{X} 我很驚訝你們這麼快就離開那位藉著基督之恩呼召你們的神，而去隨從別的福音； \end{tabularx} \\ \\ \relax
1:7 & \begin{tabularx}{0.7\textwidth}{X} 其實並沒有另一個福音，不過有些人騷擾你們，要把基督的福音更改了。 \end{tabularx} \\ \\ \relax
1:8 & \begin{tabularx}{0.7\textwidth}{X} 但無論是我們或是天上來的使者，若傳福音給你們，與我們所傳給你們的不同，他該受詛咒！ \end{tabularx} \\ \\ \relax
1:9 & \begin{tabularx}{0.7\textwidth}{X} 我們已經說了，現在我再說，若有人傳福音給你們，與你們以往所領受的不同，他該受詛咒！ \end{tabularx} \\ \\ \relax
1:10 & \begin{tabularx}{0.7\textwidth}{X} 我現在是要得人的心，還是要得神的心呢？難道我在討人的喜歡嗎？我若仍舊想討人的喜歡，我就不是基督的僕人了。 \end{tabularx} \\ \\ \relax
1:11 & \begin{tabularx}{0.7\textwidth}{X} 弟兄們，我要你們知道，我所傳的福音不是按照人的意思； \end{tabularx} \\ \\ \relax
1:12 & \begin{tabularx}{0.7\textwidth}{X} 因為我不是從人領受的，也不是人教導我的，而是藉著耶穌基督的啟示而來。 \end{tabularx} \\ \\ \relax
1:13 & \begin{tabularx}{0.7\textwidth}{X} 你們聽說過從前我在猶太教中的行徑，我怎樣竭力壓迫殘害神的教會。 \end{tabularx} \\ \\ \relax
1:14 & \begin{tabularx}{0.7\textwidth}{X} 在猶太教中，我比本國許多同輩的人更激進，為我祖宗的傳統更熱心。 \end{tabularx} \\ \\ \relax
1:15 & \begin{tabularx}{0.7\textwidth}{X} 然而，那位把我從母腹裡分別出來、又施恩呼召我的神，既然樂意 \end{tabularx} \\ \\ \relax
1:16 & \begin{tabularx}{0.7\textwidth}{X} 把他兒子啟示在我心裡，讓我在外邦人中傳揚他，我就沒有跟有血有肉的人商量， \end{tabularx} \\ \\ \relax
1:17 & \begin{tabularx}{0.7\textwidth}{X} 也沒有上耶路撒冷去見那些比我先作使徒的，惟獨到阿拉伯去，後來又回到大馬士革。 \end{tabularx} \\ \\ \relax
1:18 & \begin{tabularx}{0.7\textwidth}{X} 過了三年，我才上耶路撒冷去見磯法，和他同住了十五天。 \end{tabularx} \\ \\ \relax
1:19 & \begin{tabularx}{0.7\textwidth}{X} 至於別的使徒，除了主的兄弟雅各，我都沒有見過。 \end{tabularx} \\ \\ \relax
1:20 & \begin{tabularx}{0.7\textwidth}{X} 我現在寫給你們的是在神面前說的，不說謊話。 \end{tabularx} \\ \\ \relax
1:21 & \begin{tabularx}{0.7\textwidth}{X} 以後我到了敘利亞和基利家一帶； \end{tabularx} \\ \\ \relax
1:22 & \begin{tabularx}{0.7\textwidth}{X} 那時，在基督裡的猶太各教會都沒有見過我的面。 \end{tabularx} \\ \\ \relax
1:23 & \begin{tabularx}{0.7\textwidth}{X} 不過他們聽說「那從前壓迫我們的，現在竟傳揚他原先所殘害的信仰」。 \end{tabularx} \\ \\ \relax
1:24 & \begin{tabularx}{0.7\textwidth}{X} 他們就為我的緣故歸榮耀給神。 \end{tabularx} \\ \\ \relax
2:1 & \begin{tabularx}{0.7\textwidth}{X} 過了十四年，我再上耶路撒冷去，巴拿巴同行，也帶了提多一起去。 \end{tabularx} \\ \\ \relax
2:2 & \begin{tabularx}{0.7\textwidth}{X} 我是奉了啟示上去的；我把在外邦人中所傳的福音對弟兄們說明，我是私下對那些有名望的人說的，免得我現在或是從前都徒然奔跑了。 \end{tabularx} \\ \\ \relax
2:3 & \begin{tabularx}{0.7\textwidth}{X} 但跟我同去的提多，雖是希臘人，也沒有勉強他受割禮； \end{tabularx} \\ \\ \relax
2:4 & \begin{tabularx}{0.7\textwidth}{X} 因為有偷著混進來的假弟兄，暗中窺探我們在基督耶穌裡擁有的自由，要使我們作奴隸， \end{tabularx} \\ \\ \relax
2:5 & \begin{tabularx}{0.7\textwidth}{X} 可是，為要使福音的真理仍存在你們中間，我們一點也沒有讓步順服他們。 \end{tabularx} \\ \\ \relax
2:6 & \begin{tabularx}{0.7\textwidth}{X} 至於那些有名望的，不論他們是何等人，都與我無關；神不以外貌取人。那些有名望的，並沒有加增我甚麼。 \end{tabularx} \\ \\ \relax
2:7 & \begin{tabularx}{0.7\textwidth}{X} 相反地，他們看見了主託付我傳福音給未受割禮的人，正如主託付彼得傳福音給受割禮的人； \end{tabularx} \\ \\ \relax
2:8 & \begin{tabularx}{0.7\textwidth}{X} 那感動彼得、叫他為受割禮的人作使徒的，也感動我，叫我為外邦人作使徒。 \end{tabularx} \\ \\ \relax
2:9 & \begin{tabularx}{0.7\textwidth}{X} 那些被認為是教會柱石的雅各、磯法、約翰知道神所賜給我的恩典，就跟我和巴拿巴握右手以示合作，同意我們往外邦人那裡去，他們往受割禮的人那裡去。 \end{tabularx} \\ \\ \relax
2:10 & \begin{tabularx}{0.7\textwidth}{X} 他們只要求我們記念窮人，這也是我一向熱心在做的。 \end{tabularx} \\ \\ \relax
2:11 & \begin{tabularx}{0.7\textwidth}{X} 後來，磯法到了安提阿，因為他有可責之處，我就當面反對他。 \end{tabularx} \\ \\ \relax
2:12 & \begin{tabularx}{0.7\textwidth}{X} 從雅各那裡來的人未到以前，他和外邦人一同吃飯，及至他們來到，他因怕奉割禮的人就退出，跟外邦人疏遠了。 \end{tabularx} \\ \\ \relax
2:13 & \begin{tabularx}{0.7\textwidth}{X} 其餘的猶太人也都隨著他裝假，甚至連巴拿巴也隨夥裝假。 \end{tabularx} \\ \\ \relax
2:14 & \begin{tabularx}{0.7\textwidth}{X} 但我一看見他們做得不對，與福音的真理不合，就在眾人面前對磯法說：「你既是猶太人，卻按照外邦人的樣子，不按照猶太人的樣子生活，怎麼能勉強外邦人按照猶太人的樣子生活呢？」 \end{tabularx} \\ \\ \relax
2:15 & \begin{tabularx}{0.7\textwidth}{X} 我們生來就是猶太人，不是外邦罪人； \end{tabularx} \\ \\ \relax
2:16 & \begin{tabularx}{0.7\textwidth}{X} 可是我們知道，人稱義不是因律法的行為，而是因信耶穌基督，我們也信了基督耶穌，為要使我們因信基督稱義，不因律法的行為稱義，因為，凡血肉之軀沒有一個能因律法的行為稱義。 \end{tabularx} \\ \\ \relax
2:17 & \begin{tabularx}{0.7\textwidth}{X} 我們若求在基督裡稱義，自己卻還被視為罪人，那麼，基督是罪的用人嗎？絕對不是！ \end{tabularx} \\ \\ \relax
2:18 & \begin{tabularx}{0.7\textwidth}{X} 如果我重新建造我所拆毀的，這就證明自己是違犯律法的人。 \end{tabularx} \\ \\ \relax
2:19 & \begin{tabularx}{0.7\textwidth}{X} 我因律法而向律法死了，使我可以向神活著。我已經與基督同釘十字架， \end{tabularx} \\ \\ \relax
2:20 & \begin{tabularx}{0.7\textwidth}{X} 現在活著的不再是我，乃是基督在我裡面活著；並且我如今在肉身活著，是因信神的兒子而活；他是愛我，為我捨己。 \end{tabularx} \\ \\ \relax
2:21 & \begin{tabularx}{0.7\textwidth}{X} 我不廢掉神的恩；如果義是藉著律法而獲得，那麼基督就白白死了。 \end{tabularx} \\ \\ \relax
3:1 & \begin{tabularx}{0.7\textwidth}{X} 無知的加拉太人哪，耶穌基督釘十字架，已經活現在你們眼前，誰又迷惑了你們呢？ \end{tabularx} \\ \\ \relax
3:2 & \begin{tabularx}{0.7\textwidth}{X} 這是我惟一要問你們的：你們領受了聖靈，是因律法的行為或是因聽信福音呢？ \end{tabularx} \\ \\ \relax
3:3 & \begin{tabularx}{0.7\textwidth}{X} 你們既然以聖靈開始，如今竟要以肉身終結嗎？你們是這樣的無知嗎？ \end{tabularx} \\ \\ \relax
3:4 & \begin{tabularx}{0.7\textwidth}{X} 你們受這麼多的苦都是徒然的嗎？如果真是徒然的， \end{tabularx} \\ \\ \relax
3:5 & \begin{tabularx}{0.7\textwidth}{X} 那麼，神賜給你們聖靈，又在你們中間行異能，是因律法的行為或是因聽信福音呢？ \end{tabularx} \\ \\ \relax
3:6 & \begin{tabularx}{0.7\textwidth}{X} 正如亞伯拉罕「信了神，這就算他為義」。 \end{tabularx} \\ \\ \relax
3:7 & \begin{tabularx}{0.7\textwidth}{X} 所以，你們知道：有信心的人才是亞伯拉罕的子孫。 \end{tabularx} \\ \\ \relax
3:8 & \begin{tabularx}{0.7\textwidth}{X} 聖經既然預先看見神要使外邦人因信稱義，預先傳福音給亞伯拉罕，說：「萬國都必因你得福。」 \end{tabularx} \\ \\ \relax
3:9 & \begin{tabularx}{0.7\textwidth}{X} 可見，那有信心的人和有信心的亞伯拉罕一同得福。 \end{tabularx} \\ \\ \relax
3:10 & \begin{tabularx}{0.7\textwidth}{X} 凡出於律法的行為都是受詛咒的，因為經上記著：「凡不持守律法書上所記的一切而去行的，都是受詛咒的。」 \end{tabularx} \\ \\ \relax
3:11 & \begin{tabularx}{0.7\textwidth}{X} 沒有一個人靠著律法在神面前稱義，這是明顯的，因為經上說：「義人必因信得生。」 \end{tabularx} \\ \\ \relax
3:12 & \begin{tabularx}{0.7\textwidth}{X} 律法並不出於信，而是說：「行這些事的就必因此得生。」 \end{tabularx} \\ \\ \relax
3:13 & \begin{tabularx}{0.7\textwidth}{X} 既然基督為我們成了詛咒，就把我們從律法的詛咒中贖出來。因為經上記著：「凡掛在木頭上的都是受詛咒的。」 \end{tabularx} \\ \\ \relax
3:14 & \begin{tabularx}{0.7\textwidth}{X} 這是要使亞伯拉罕的福，因著基督耶穌臨到外邦人，使我們能因信得著所應許的聖靈。 \end{tabularx} \\ \\ \relax
3:15 & \begin{tabularx}{0.7\textwidth}{X} 弟兄們，我照著人的觀點說，人的遺囑一經確定，沒有人能廢棄或加增。 \end{tabularx} \\ \\ \relax
3:16 & \begin{tabularx}{0.7\textwidth}{X} 那些應許原是向亞伯拉罕和他後裔說的，並不是說「和眾後裔」，指許多人，而是說「和你那個後裔」，指一個人，就是基督。 \end{tabularx} \\ \\ \relax
3:17 & \begin{tabularx}{0.7\textwidth}{X} 我是這麼說，神預先所立的約不能被四百三十年以後的律法廢掉，使應許失效。 \end{tabularx} \\ \\ \relax
3:18 & \begin{tabularx}{0.7\textwidth}{X} 因為承受產業若是出於律法，就不再是出於應許；但神是憑著應許把產業賜給亞伯拉罕。 \end{tabularx} \\ \\ \relax
3:19 & \begin{tabularx}{0.7\textwidth}{X} 這樣說來，為甚麼要有律法呢？律法是為過犯的緣故而加上去的，等候那蒙應許的子孫來到才結束，是藉著天使經中保之手而設立的。 \end{tabularx} \\ \\ \relax
3:20 & \begin{tabularx}{0.7\textwidth}{X} 但中保本不是為單方設立的；神卻是一位。 \end{tabularx} \\ \\ \relax
3:21 & \begin{tabularx}{0.7\textwidth}{X} 這樣，律法是與神的應許對立嗎？絕對不是！如果律法的頒佈能使人得生命，義就誠然出於律法了。 \end{tabularx} \\ \\ \relax
3:22 & \begin{tabularx}{0.7\textwidth}{X} 但聖經把萬物都圈在罪裡，為要使因信耶穌基督而來的應許歸給信的人。 \end{tabularx} \\ \\ \relax
3:23 & \begin{tabularx}{0.7\textwidth}{X} 但這「信」還未來以前，我們被看守在律法之下，像被圈住，直到那將來的「信」顯明出來。 \end{tabularx} \\ \\ \relax
3:24 & \begin{tabularx}{0.7\textwidth}{X} 這樣，律法是我們的啟蒙教師，直到基督來了，好使我們因信稱義。 \end{tabularx} \\ \\ \relax
3:25 & \begin{tabularx}{0.7\textwidth}{X} 但這「信」既然來到，我們從此就不在啟蒙教師的手下了。 \end{tabularx} \\ \\ \relax
3:26 & \begin{tabularx}{0.7\textwidth}{X} 其實，你們藉著信，在基督耶穌裡都成為神的兒女。 \end{tabularx} \\ \\ \relax
3:27 & \begin{tabularx}{0.7\textwidth}{X} 你們凡受洗歸入基督的都披戴基督了： \end{tabularx} \\ \\ \relax
3:28 & \begin{tabularx}{0.7\textwidth}{X} 不再分猶太人或希臘人，不再分為奴的自主的，不再分男的女的，因為你們在基督耶穌裡都成為一了。 \end{tabularx} \\ \\ \relax
3:29 & \begin{tabularx}{0.7\textwidth}{X} 既然你們屬於基督，你們就是亞伯拉罕的子孫，是照著應許承受產業的了。 \end{tabularx} \\ \\
[1ex]
\hline
\hline
\end{longtable}
$^{1}$各位網友平安,我們開始研讀加拿大書意見書.
加拿大書講一個很重要的真理.
就是因著信,我們可以得以稱義.
這是我們信仰的基礎.
是一個很重要真理的基石.
所以加拿大書意見書.
就叫做基督徒得自由的大憲章.
保羅他親自告訴我們.
因著信,我們可以被神稱我們為義.
就是我們宣告無罪.
因此這個是我們的憲章.
我們得自由的憲章.
一個的確據來的.
憲章是不可以改變的.
當你簽了一份的合約.
那份合約就開始生效了.
保羅就用六章聖經.
跟我們解釋這個他認識因信稱義的道理.
通常保羅在他的書信裡面.
就分前段和後段.
加拿大書一共有六章.
前段就是第一到第三章.
是比較教義性的.
而下一段就是第四到第六章.
通常我們就說保羅講一些實際行動.
對信徒一些的圈面.
讓我們怎樣去應用這個真理.
所以我就預備用兩次的時間.
就跟大家看完六章的加拿大書.
加拿大書如果用標題來分.
就可以說第一,第二章是保羅辯護這個恩典的福音.
他幫福音作辯護.
然後第三和第四章.
就是他解釋這個恩典的福音.
然後最後兩章就是應用第五,第六章.
應用恩典的福音.
保羅一開頭就說他這封信是寫給加拿大的各教會.
因為加拿大可以當為一個地區.
好像一個區域,一個省.
所以保羅好像不是寫給單一的教會.

$^{41}$就是這個地區的所有教會.
他們一起寫.
因此在很多學者的討論裡面.
就有兩個不同的理論.
到底保羅心目中所說的加拿大地區.
是指哪個地區呢?.
所以就有北加拿大論.
和南加拿大論.
這兩種的說法.
所謂北加拿大論.
就指比較北面的區域.
在《使徒行傳》裡面的記載.
就保羅在第二次的行程.
他才去這個地區的.
然後在第三次行程.
他又去.
換句話來說.
加拿大書就保羅去完第二,第三次行程之後.
他才寫的.
所以寫作的日期就比較遲一點.
第二種的理論就叫做南加拿大論.
如果你是聽從南加拿大論的話.
就保羅第一次宣教行程.
就去了這個地區.
故此順從這個理論的學者就說.
保羅寫這本書的日期是比較早一點的.
因為他去完第一次行程.
探完這個地區的教會之後.
他就開始寫這封信.
其實這封信很重要的地方.
就說到我們怎樣因著信稱賢.
大家都知道馬丁路德.
他發起宗教改革.
就因為他在修道院裡面.
自己讀加拿大書這本書.
他就發覺和當時天主教.
和當時一般神學的理論很不同.
因為當時的天主教會就賣那些贖罪券.
甚至有些神父去到街上就在那裡宣傳.
如果你買了贖罪券.

$^{81}$就好像叮一聲.
就把錢放在錢箱裡面.
你就可以把這些贖罪券救你已經逝世的親人,家人.
當你扔了這個贖罪券.
叮一聲.
你的親友的靈魂就升上天堂.
馬丁路德聽了之後.
覺得很不對勁.
什麼亂說都可以.
於是他就根據他對加拿大書很深刻的認識.
就開始貼出95條他對信仰的一些反省.
來公開挑戰當時的天主教會.
因此一場很大波浪的宗教改革.
Reformation.
就在德國興起了.
我曾經有一次去歐洲旅行.
當去到德國的時候.
就有機會參觀一些城市.
那個導遊就帶我們去參觀.
一些馬丁路德的歷史古蹟.
看到有一個地方有雙鞋.
鐵造的鞋.
是馬丁路德去過的一些腳蹤.
我還記得師母很興奮.
就站在鐵鞋上幫他拍張照.
就是說馬丁路德的Footstep.
我現在都走起馬丁路德的Footstep腳蹤之後.
我們今天都說我們的教會是根據聖經的教導.
前人有很多神學家為我們鋪了一條路.
我們就明白我們所信的教義是怎樣發展下來.
所以我們都在走先賢所留下的腳蹤.
所以我們要效法先賢的腳蹤來繼續走.
我們開始《加拿大書》第一章.
保羅每一封信都是用問安來開始的.
我是使徒或者奉召作使徒的保羅.
寫信給加拿大的教會.
然後我問你們安.
但是這封信問完安之後很出奇.
保羅竟然用真經來開始.
哇!問完安之後不是說一些東西.

$^{121}$而是說我很震驚.
我們中文和合本的翻譯就有少許不同.
這裡就說我奉命在你們中間.
就說很稀奇.
有些新的譯本就用震驚.
因為原文希利文這個字其實是很震驚.
很震蕩.
你們現在改了你們所信的東西.
因此這封信帶著一個辯護的口吻.
開始它的內容.
保羅就說我震驚.
你們現在好像走歪了路.
和我教你們的真理有些不同.
這樣不妥當.
於是保羅就開始提出.
要更正他們錯誤了的思想.
你們已經錯了.
現在馬上停止.
回頭走正確的路.
所以當我們叫人要悔改信耶穌.
我們做牧師的經常解釋什麼叫悔改.
就是你覺得你以前走錯了路.
如果這條路是錯的.
去不到的目的地.
快點回頭走過一條對的路.
正確的路.
所以保羅就用一個這麼強烈的口吻.
開始迦拉泰書.
他就說你們不要再搞錯了.
要回去你們正確的信仰.
就是耶穌基督為我們的罪而死在十字架上.
我們憑的不是自己的行為.
因為行為不能夠救人.
唯有從心裡面去信靠這位釘十字架.
愛我們的救主.
今天我們都說你信了耶穌沒有?.
是信耶穌.
所以有很多人說信耶穌最簡單.
又不用給錢.
又不用做什麼.

$^{161}$又不用讀什麼書.
你聽完之後.
你說我信.
這樣就得救了.
有些人就覺得會不會牙醉一點.
這樣就得救了.
所以其實神的恩典真的很大.
你信耶穌不用你有學問.
又不用你識字.
又不用你讀過什麼.
你說我信.
就是憑著信心.
你心裡面有個扭轉.
認自己的罪.
你悔改.
耶穌十字架的寶血.
就洗乾淨我們一切的罪.
所以在座如果有一些你還沒信耶穌.
我都很誠懇勸你.
快點做一個悔改認罪信耶穌的祈禱.
聖經告訴我們.
你若心裡信神.
叫祂從死裡復活.
你若信祂為你釘十字架.
洗淨你的罪.
你就必得救.
這就叫做「因信稱義」.
所以保羅就是這樣來勸勉弟兄姊妹.
而不僅「因信稱義」是神給保羅的一個任命.
因為保羅奉命做外邦人的使徒.
他要將這個真理帶給那些外邦人.
然後保羅又說.
我是為基督來傳這個福音.
所以我很奇怪為什麼你們中間有些人.
說我不是使徒.
對保羅有時候有些批評.
所以保羅好像也是辯護自己使徒的身份.
你們要接受福音.
我是向我顯現的主耶穌.
差我來將福音帶給你們這些外邦人.

$^{201}$而且加拿大書有一點就是.
如果你要信耶穌.
你要先做猶太人.
就是守猶太人的一些規條.
保羅是非常反對的.
他說聖經從來沒有要求我們.
我們要守國例.
我們要做猶太人做的事情.
因為救恩是伯伯給過我們的.
所以他就說.
因此神是提出「恩順清義」這個真理.
保羅就從經驗上去辯護.
就是第三章第一到第五節.
然後又從一些聖經去辯護.
就由第三章第四節到十四節.
然後又從常理去辯論.
就是三章十五到十八節.
他用這三方面去證明給你知道.
真是「恩順清義」.
我們既然認識律法的功用.
是叫我們知罪.
所以為什麼有時有些交通律法.
你超速警察可以抄你牌罰你錢.
你建一些禁區泊車.
警察會給你牛肉乾.
因為你知道這些是法律.
你知道自己犯了法.
所以律法是有用的.
就是知道你什麼叫做錯.
所以律法有一種提醒的作用.
你原來走錯了.
那怎麼做?.
快點悔改.
讓保羅這個頭腦很清晰的使徒.
寫了加拿大書前段的一到三章.
用很簡明的道理.
用一些辯護的方式.
用一些理論.
幫助我們對教義的明白.
我希望各位弟兄姊妹.

$^{241}$你同樣都明白.
願神賜福給各位.
\newpage



\section{創世記 24:7}
\label{sec:031713}
\textbf{聖經人聖經事:娶一個妻子}~~~~黃克勤牧師博士
\newline
\newline
連結: \href{http://www.church.com.hk/acms/content.asp?site=cdc\&op=show\&type=product\&code=031713\&layout=sermon}{\texttt{031713}} ~~~~ 語音日期: 2026-02-17
\newline
\newline
\hyperref[sec:031712]{\small{< < < PREV SERMON < < <}}
~
\hyperref[sec:index_chronic]{\small{[返順時目]}}
~
\hyperref[sec:031714]{\small{> > > NEXT SERMON > > >}}
\newline
\newline
創世記 24:7
\newline
\begin{longtable}{cl}
\hline
\hline
章節 & 經文 (和合本修訂版)\\
\hline
24:7 & \begin{tabularx}{0.7\textwidth}{X} 耶和華—天上的神曾帶領我離開父家和本族的地，對我說話，向我起誓說：『我要將這地賜給你的後裔。』他要差遣使者在你面前，你就可以從那裡為我兒子娶妻。 \end{tabularx} \\ \\
[1ex]
\hline
\hline
\end{longtable}
$^{1}$鼎師傅平安,又來看聖經了.
還是在看亞伯拉罕.
這次聖經人,聖經事,生命樹為你預備.
娶了一個妻子.
這是什麼呢?.
亞伯拉罕叫僕人.
為他的兒子要娶一個妻子.
誰娶一個妻子呢?.
本來娶妻子就算了.
不過聖經裡面的字.
亞伯拉罕說你幫我兒子娶一個妻子.
所以用一個不是兩個.
不是數目的問題.
是娶妻.
和娶一個妻子是經文的字.
我們看看怎麼發生.
這幾天講的東西都很長.
為什麼呢?因為是故事.
不是一節金句.
很長.
我們在聖經故事裡面學習.
二十四章第一節.
亞伯拉罕年紀老邁.
向來在一切事上也華都似福及譚.
你看到亞伯拉罕不單武功好.
四王戰勝五王.
他還很懂得人情世故.
又很懂得做生意.
很富裕.
去到做一個寄居的.
他沒有說神將這塊地給了我.
你們將來都是我管轄的.
我是一個寄居的人.
外人來的.
想找塊地娶我太太.
我煮你拿去.
這樣就很好.
開這麼高的價.
他聰明照給.
四百尺赫勒.

$^{41}$真是天價.
我覺得.
其實我不知道貴不貴.
不過應該貴.
四百尺赫勒都不少錢.
好了.
亞伯拉罕年紀老邁.
向來在一切事上也祝福他.
我剛剛講了很多福氣.
亞伯拉罕對管理他全業最老的僕人.
這個老僕人厲害了.
就是沒有說什麼名字.
很多時候上帝的僕人.
都沒有名字.
死了都沒有人知道.
沒有人紀念.
特別很多宣教而在外地死去的人.
經過一兩代之後.
沒有人再認識他.
歷史教歷史沒有他的名字.
不是和馬丁路德齊名嗎?.
沒有人認識.
這個老僕人都沒有人知道他什麼名字.
起碼沒有說.
亞伯拉罕對管理他全業的老僕人.
最老的那個.
就是說管理他全業.
大管家.
很厲害.
很信任的.
請你把手放在我大腿底下.
什麼意思?.
放手在大腿底下.
兩個男人.
其實放手在大腿底下.
當時.
就是一個我向你起誓承諾要做的事情.
你如果真的幫我做.
答應和我做的好.
你把手放在我大腿底下.

$^{81}$然後你發誓.
表示我真的承諾托住這件事.
請你把手放在我大腿底下.
起誓.
我要叫你指著耶和華.
是天地的主起誓.
幾樣東西.
第一.
不要為我的兒子娶這迦南地中的女子為妻.
你要往我本地本族為我的兒子以撒娶一個妻子.
娶一個妻子就是今天的題目.
這是第一點.
螢光字第一點.
是兩組.
第一是不要.
然後要.
不要為我的兒子娶這迦南地中的女子為妻.
隨便在本地選一個就容易.
又要我回到本地本族.
兩河流域.
他已經來了巴勒斯坦.
今天是巴勒斯坦.
以後是迦南地.
他要走回兩河流域.
即是他要向東走,向北,然後轉向東,去到兩河流域的根源.
然後跟著兩河流域走到底,出海才是吳爾城.
巴布蘭那裡出來.
嘩!走多遠.
最老的那個去了,我怕他中途躺下.
不知道怎麼辦.
叫他這樣做.
你不要.
你和我找我兒子娶個太太.
你和我找個媳婦.
叫個最老的僕人和我找個媳婦.
巴布蘭不自己去.
一百幾歲,一百幾歲.
起碼一百三十幾,近一百四十.
不去了,叫最老的僕人.
可能叫第二老那個去.

$^{121}$你和我找一個太太.
不要在本地找.
你要去我本地本族那裡找.
為我的兒子娶一個太太.
巴布蘭對她說.
如果去到這麼遠.
我想娶她,叫她嫁我少東.
以撒少爺.
她不肯來.
就要我們少爺,以撒過去.
給她見見面,那怎麼辦?.
看看樣子是否英俊,有沒有驚嚇.
有沒有樣子,怎樣都要先看看.
相親先.
如果女子不肯跟我來.
我就要帶以撒少爺去.
回去你原居之地.
亞伯拉罕第二句螢光的字.
不要,不要帶我兒子回去.
還要加一句在前面.
你要謹慎.
千萬要記得.
不要帶以撒回去.
她不屬於那裡的.
為甚麼呢?.
耶和華天上的神曾經帶我離開那裡.
本地本族父家.
跟我說.
你要在這裡.
賜這塊地那裡.
生子立後.
是這裡的.
所以你不可以帶我兒子回去.
我必差一點在你面前.
你就可以在那裡為我兒子娶一個妻子.
這裡說甚麼呢?.
不要,第二個螢光.
第一個就是不要,要.
第二個就是不要.
然後可以.

$^{161}$為我在那裡找一個.
這個不要就是我要保持這個承諾.
就是我要在這裡.
第三點要看的.
就是你可以這樣做.
就是為我兒子娶一個妻子.
在那裡.
那怎麼辦?我都不知道怎樣找.
不怕的.
神會差遣使者在你面前.
帶你怎樣做.
亞伯拉罕是有信心的.
看一個女子不肯跟你來.
我就使你的氣勢如雷無光.
因為不可以帶我兒子去那裡.
於是暴人就把手.
放在主人亞伯拉罕的大腿下.
就為這件事起誓了.
我覺得厲害的就是.
亞伯拉罕見到要為自己的兒子.
娶一個媳婦.
娶一個太太的時候.
他自己都很老了.
不方便走那麼遠.
就叫另一個很信任的老僕人.
去.
講到明.
不可以娶本地女子.
本地女子.
為什麼?.
因為我的後裔裡面.
特別這些我不想有本地人混進來.
我一定要離開本地本族父家.
不可以在那裡.
在這裡神要賜給我後裔.
在這裡.
其實我都沒有什麼本地本族父家的人.
所以麻煩你回去本地本族父家.
娶一個我們族裔的回來.
繼續繁衍.

$^{201}$不是一個迦南人.
是外族.
昨天才說.
我是外人.
是客女.
好像我們人生在世都是.
外人是客女.
在這個世界.
我要找一個同一個種族出來.
又有同一個信仰的人.
你都講得很清楚了.
亞伯拉罕是很遵守.
很尊敬上帝的說話.
所以不行.
你一定要跟我娶一個本地人.
如果他不肯.
那就算了.
這件事很厲害.
沒有人肯來.
那怎麼辦?.
那就算了.
不關你的事.
神有心意.
娶不到就娶不到.
羅德兩個女兒.
阿曼和摩訶的祖先就不是了.
沒有人想來山上.
那怎麼辦?.
我和爸爸同寢生子立後.
為我們自己.
自己的方法.
亞伯拉罕呢?.
如果你帶不到人回來.
他不肯來.
那就算了.
再去一次.
沒有了.
出WhatsApp請人.
沒有了.
Facebook.

$^{241}$沒有了.
如果沒有就沒有了.
不過.
我相信神會派他的使者為你開路.
你會和我的兒子娶到一個媳婦.
於是他就去了.
我本人的金句就是.
神說我要將這地賜給你做後裔.
所以你要回去.
以往的地方娶個妻子.
我在這個地方住.
但是.
我的後裔要回去我本地的人.
所以到今天.
你發覺以色列人.
基本上.
如果真的守住他們的信仰.
他們是不會和外人通婚的.
一定要和本地人.
因為和自己人.
自己的族裔.
才能夠維持他們純種的.
以色列人.
是這樣的.
他們很尊敬這件事.
當然今天很多都放鬆了.
亞伯拉罕其實是堅持著.
為什麼?.
因為.
神給了這塊地我.
但是我們的族裔要屬於他.
他做這件事.
你可以和我在那裡找你.
沒有?.
沒有就算了.
厲害.
很有信心.
我們一起禱告.
主求你讓我們像亞伯拉罕那樣有信心.
要回去那麼艱難的地方.

$^{281}$那麼遠的地方.
歷盡千辛萬苦.
路途遙遠.
萬水千山.
都要經過那些沙漠.
都要娶一個.
他本地同宗族的女子.
成為自己的媳婦.
主要都是他堅持著自己族裔與信仰.
要在這個新的地方發展.
以致他做這件事.
如果沒有.
沒有人肯跟著來.
就算了.
見到他有這個信心.
有這個交托.
是我們需要學習的.
讓我們在這些要與不要當中.
我們都做一個選擇.
好像學習亞伯拉罕一樣.
讓我們更多摸著你的心意.
更多信靠你.
祈禱奉靠耶穌基督明求.
阿們.
\newpage



\section{申命記 4:1-14}
\label{sec:031714}
\textbf{過去和現在的對話 (3) 聽誰說}~~~~李觀南傳道
\newline
\newline
連結: \href{http://www.church.com.hk/acms/content.asp?site=cdc\&op=show\&type=product\&code=031714\&layout=sermon}{\texttt{031714}} ~~~~ 語音日期: 2026-02-21
\newline
\newline
\hyperref[sec:031713]{\small{< < < PREV SERMON < < <}}
~
\hyperref[sec:index_chronic]{\small{[返順時目]}}
~
\hyperref[sec:031715]{\small{> > > NEXT SERMON > > >}}
\newline
\newline
申命記 4:1-14
\newline
\begin{longtable}{cl}
\hline
\hline
章節 & 經文 (和合本修訂版)\\
\hline
4:1 & \begin{tabularx}{0.7\textwidth}{X} 「現在，以色列啊，聽我所教導你們的律例典章，要遵行，好使你們存活，得以進入耶和華－你們列祖之神所賜給你們的地，承受為業。 \end{tabularx} \\ \\ \relax
4:2 & \begin{tabularx}{0.7\textwidth}{X} 我吩咐你們的話，你們不可加添，也不可刪減，好叫你們遵守耶和華－你們神的命令，就是我所吩咐你們的。 \end{tabularx} \\ \\ \relax
4:3 & \begin{tabularx}{0.7\textwidth}{X} 你們已親眼看見耶和華因巴力‧毗珥所做的。凡隨從巴力‧毗珥的人，耶和華－你的神都從你中間除滅了。 \end{tabularx} \\ \\ \relax
4:4 & \begin{tabularx}{0.7\textwidth}{X} 只有你們這緊緊跟隨耶和華－你們神的人，今日全都存活。 \end{tabularx} \\ \\ \relax
4:5 & \begin{tabularx}{0.7\textwidth}{X} 看，我照著耶和華－我的神所吩咐我的，將律例和典章教導你們，使你們在所要進去得為業的地上遵行。 \end{tabularx} \\ \\ \relax
4:6 & \begin{tabularx}{0.7\textwidth}{X} 你們要謹守遵行；這就是你們在萬民眼前的智慧和聰明。他們聽見這一切律例，必說：『這大國的人真是有智慧，有聰明！』 \end{tabularx} \\ \\ \relax
4:7 & \begin{tabularx}{0.7\textwidth}{X} 哪一大國有神明與他們相近，像耶和華－我們的神在我們求告他的時候與我們相近呢？ \end{tabularx} \\ \\ \relax
4:8 & \begin{tabularx}{0.7\textwidth}{X} 哪一大國有這樣公義的律例典章，像我今日在你們面前所頒佈的這一切律法呢？ \end{tabularx} \\ \\ \relax
4:9 & \begin{tabularx}{0.7\textwidth}{X} 「但你要謹慎，殷勤保守你的心靈，免得忘記你親眼所看見的事，又免得在你一生的年日這些事離開你的心，總要把它們傳給你的子子孫孫。 \end{tabularx} \\ \\ \relax
4:10 & \begin{tabularx}{0.7\textwidth}{X} 你在何烈山站在耶和華－你神面前的那日，耶和華對我說：『你為我召集百姓，我要叫他們聽見我的話，使他們活在世上的日子，可以學習敬畏我，又可以教導他們的兒女。』 \end{tabularx} \\ \\ \relax
4:11 & \begin{tabularx}{0.7\textwidth}{X} 那時，你們近前來，站在山下；山上有火燃燒，直沖天頂，並有黑暗、密雲、幽暗。 \end{tabularx} \\ \\ \relax
4:12 & \begin{tabularx}{0.7\textwidth}{X} 耶和華從火焰中對你們說話，你們聽見說話的聲音，只有聲音，卻沒有看見形像。 \end{tabularx} \\ \\ \relax
4:13 & \begin{tabularx}{0.7\textwidth}{X} 他將所吩咐你們當守的約指示你們，就是十條誡命，並將誡命寫在兩塊石版上。 \end{tabularx} \\ \\ \relax
4:14 & \begin{tabularx}{0.7\textwidth}{X} 那時，耶和華吩咐我將律例典章教導你們，使你們在所要過去得為業的地上遵行。」 \end{tabularx} \\ \\
[1ex]
\hline
\hline
\end{longtable}
$^{1}$今天是年初五.
仍然是農曆年期間.
所以也要跟大家說聲新年快樂.
主恩上載.
今晚我們繼續說我們《申命記》的第一個系列.
叫做《過去與現在的對話》的一篇訊息.
這一篇今天的訊息.
其實已經是這個系列的第三講了.
過年之前.
其實我們分別邀請了.
Pastor Eamon 和 Pastor Liz 來到我們中間.
他們和我們其實主要是集中分享了.
《申命記》的頭三章不同的內容.
《申命記》的頭三章的內容.
其實主要是在說什麼呢?.
就是在說摩西說回以色列人之前的一些歷史.
他和一群新的以色列人.
指的是那群在礦業出生的那一代.
叫做新一代的以色列人.
去回顧當日摩西自己.
如何被蒙召.
然後帶上一代的以色列人去出埃及.
然後他又在以色列人的中間做了一件事情.
就是為他們立了不同的領袖.
什麼千夫長,百夫長,十夫長那一類的.
而到他們一直走一直走.
就到了一個地方叫做加底斯巴尼亞的地方.
摩西就和一群以色列人說.
我們要照耶和華所說的.
上去德那地為業.
不過這群百姓聽到這件事之後.
他們就反而要求摩西要先派探子.
去窺探迦南整個地方.
最初他們提出這個要求.
說起碼等我們將來進到迦南地之後.
我們可以知道怎樣去走,走哪條路,去哪個城.
但是結果就是當探子回顧回來之後.
因為消息不太好.
有人又高又大,城牆又高又厲害.
這群百姓就害怕了.

$^{41}$繼而就走去埋怨摩西.
而這件事簡單來說.
長話短說就是令到耶和華.
要懲罰這群以色列百姓.
要他們在曠野裡一直漂流.
不過我很慶幸的是.
以色列人在這個漂流的日子裡.
他們仍然有耶和華的保守.
以致他可以擊敗那兩個阿摩利王.
叫做巴沙王和阿摩利王.
然後又記載了.
摩西因為以色列人其中兩個半支派.
在這裡因為打完仗,拿到地方的時候.
就主動要求摩西.
不如你分這些地給我們吧.
那麼摩西就真的.
應了這兩個半支派的要求.
就把那些地分了給他.
然後一直到第三章的尾段.
摩西也提到一件事.
就是耶和華原來親自跟摩西說了.
他說你不能進入迦南地的地方.
而是要約書亞成為一個.
以色列人的一個新的領袖.
繼續帶領這一班以色列人到約旦河.
過約旦河進去迦南.
繼續之後的路程.
那麼摩西就將整個過去以色列人.
由出埃及到他們當時.
一直居住在約旦河河東.
整個旅程所發生的所有事情.
無論是神怎樣恩代他們.
還是過去這班以色列人一些失敗的經歷.
摩西也一一地向這一班新的以色列人.
去陳明,去交代.
摩西用了三章的經文去回顧了一個歷史.
也都成為了一個很好的開場白.
繼而他就去到整段重溫歷史的真正重點了.
也是我今天要跟大家講的這一段的經文.
這一段第四章的內容.

$^{81}$不過開始之前讓我們先一同圖告.
天父上帝多謝你.
多謝你因為.
當我們去學習以色列人的歷史.
我們很多時候會想.
都其實不關我事.
但是主你卻同樣.
一些上一代以色列人發生的事.
你同樣藉著這些事蹟.
繼而去教導新一班以色列人.
同樣你今天去教導我們.
所以神我求你.
你幫助我們讓我們打開我們的耳朵.
讓我們知道.
聖經寫下這些東西.
是於我們有益.
是上帝要我們跟我們說話.
所以神我求你.
讓我們有謙卑的心去聆聽你的說話.
奉耶穌基督的聖名自祈求.
阿們.
好,那我們先去閱讀一下今天的經文.
今天選了的經文就是《申命記》的四章一至十四節.
看到嗎?.
好,我們一二三一起讀.
一二三.
以色列人啊.
現在我所教訓你們的律例典章.
你們要聽從遵行.
好叫你們存活.
得以進入耶和華.
你們列祖之上帝所賜給你們的地.
聖授為業.
所吩咐你們的話.
你們不可加添也不可刪減.
好叫你們遵守我所吩咐的.
就是耶和華你們上帝的命令.
耶和華因巴勒比爾的事所行的.
你們親眼看見了.
凡隨從巴勒比爾的人.

$^{121}$耶和華你們的上帝都從你們中間除滅了.
唯有你們專靠耶和華你們上帝的人.
今天全都存活.
我照著耶和華我上帝所吩咐的.
將律例典章教訓你們.
使你們在所要進去得為諜的地上遵行.
所以你們要謹守遵行.
這就是你們在萬民眼前的智慧聰明.
他們聽見這一切律例必說.
這大國的人真是有智慧有聰明.
那一大國的人有神與他們相近.
像耶和華我們的上帝.
在我們求告他的時候與我們相近呢?.
又哪一大國有這樣公義的律例典章.
像我今天在你們面前所陳明的這一切律法呢?.
你只要謹慎.
殷勤保守你的心靈.
免得忘記你親眼所看見的事.
又免得你一生這事離開你的心.
總要傳給你的子子孫孫.
你在何列山站在耶和華.
你上帝面前的那一天.
耶和華對我說.
你為我焦聚百姓.
我要叫他們聽見我的話.
使他們存活在世的日子.
可以學習敬畏我.
又可以教訓兒女這樣行.
那時你們更前來站在山下.
山上有火焰沖天.
並有昏黑密淪幽暗.
耶和華從火焰中對你們說話.
你們只聽見聲音.
卻沒有看見形象.
他將所吩咐你們.
當受的藥指示你們.
就是十條界.
並將這界寫在兩塊石板上.
那時耶和華又吩咐我.
將律例典章教訓你們.

$^{161}$使你們在所要過去得遺孽的地上遵行.
謝謝大家.
經文讀完了.
正如剛剛開場白所說.
其實當摩西回顧.
即是和班爾斯利人回顧完一些歷史之後.
今天這段經文其實是帶出了.
在第五章之後.
其實摩西就開始逐條逐條去講.
即是講回十界和不同的律例.
慢慢去講解律例.
他要重新律法的原因.
當然的.
其實我們之前如果有在這系列的頭兩篇來聽的話.
Pastor Eamon 其實有分享過.
申命記一方面其實是為了.
令到一班新一代的以色列人.
去知道耶和華過去在他們當中所做的一切的事.
因為這一班人其實很多新一代的以色列人.
其實他們是沒有親身經歷過的.
所以其實告訴他們.
讓他們知道原來上帝曾經對我們這個民族.
這一班人是有這些神蹟奇事發生過的.
讓他們知道.
讓他們原來知道他們今天能夠得以存活.
是有上帝的過去種種種種的恩典在裡面.
而其次.
如果大家之前曾經看過申命記的話.
你就會知道申命記的最後一章是在講摩西逆死.
上一篇讀Pastor Eamon 同樣有提過.
申命記其實是摩西的一個臨終的遺言.
所以我會想既然是遺言.
那很自然就會怎樣呢.
我們就會將自己認為很重要的一些事情.
要交代給我們的後人知道.
就好像那些保險櫃的密碼.
範圍竹放在哪裡那些.
臨死前都應該要告訴後人.
很正常的.
而對摩西的遺言.

$^{201}$他認為最重要的.
正正就是他們這一班後人.
要遵行耶和華的律法.
那為甚麼這麼重要呢.
剛剛我們讀的經文其實可以看到.
我們看回第一至四節的經文.
知道摩西在提醒以色列人.
他們要聽從.
遵行摩西所教訓他們的一切的律例典章.
為的就是要令他們能夠得以存活.
得以可以按上帝對這班以色列的列祖.
曾經應許過的事情.
能夠進入迦南地.
可以得地為業.
摩西在提醒這班以色列人.
能不能進入迦南.
其實不在於這班以色列人.
究竟打仗厲害不厲害.
還是他們本身是否軍備充足.
或者甚麼甚麼.
而是在取決於.
他們這班人有沒有聽神的話.
摩西特別提到.
有一個叫巴力比爾的一件事.
發生過的一件事.
巴力比爾其實發生過甚麼事呢.
在文述的二十五章.
說的是當時一班以色列人.
本身住在一個叫十庭的地方.
然後有摩亞的女子.
來到這班以色列人的裡面.
就令到這班百姓.
和這班摩亞的女子.
行淫亂的事.
繼而又聽了.
這班跟了的以色列人.
就跟隨了這班人.
跟了這班摩亞的女子.
一起去拜他們的神.
獻祭給他們的神.

$^{241}$還要吃那些祭物.
跪拜他們的神.
結果這件事我沒有帶出來.
後面其實說的是.
耶和華的嬲鹿淫到這班以色列人.
就有瘟疫淫到他們的裡面.
結果這一場瘟疫.
死了一共有二萬四千人.
摩西以巴力比爾的事件.
成為一個反面的例子.
去警告那班跟從偶像的人.
其實最終是會被剪除.
在他們的民族裡面.
在以色列的民族裡面會被剪除.
唯有專靠耶和華的人.
去到當時的今天.
他才可以繼續存活.
所以摩西其實千叮萬囑.
這班以色列人去說.
他所吩咐的.
又或者他所轉述上帝所吩咐的.
以色列人是要遵守這一切的說話.
他是不可以加添.
以及不可以生減.
他在提醒的是.
這班以色列人.
其實不要將自己的想法.
又或者世間的那些所謂文化傳統.
加在神的話上面.
也不要因為神的說話.
有些地方覺得不合自己的心意.
就隨意地刪掉.
因為其實我們在聖經之後都會知道.
在新約時代.
如果我們將神的話.
額外加添一些東西上去的時候.
其實正正就是那班法利祥人在做的事.
變成了律法主義.
要一班信徒純粹很謹守.
去那班他們定下來的一些傳統.

$^{281}$但是慢慢地就忽略了.
其實原來上帝最初給你的律法.
沒有加添版本.
原來後面有很多上帝在這些設立裡面.
是有很多祂的愛.
很多祂的憐憫.
很多祂的恩典在裡面.
而這班法利賽人.
反而就是在這些神的話上面加添了.
將律法成為了世人一個無形的枷鎖.
至於生減.
其實我想不難明白.
好像我們今天很多信徒都是這樣的.
聖經是合聽的就跟進.
不合聽的就放在一邊.
那是什麼呢?.
某情況我會說這個是一種放縱.
我們強調的是那個所謂的個人自由的重要.
但是其實某情況我們是在改上帝的說話.
去令到那個所謂的信仰.
成為了我們天然製造的信仰.
摩西指出了這兩件事.
加添和生減.
這兩件事.
都是會令到人離開了真正的生命之道.
這個不是真正的跟隨著上帝.
而當我看到這裡的時候.
我會甚至認為.
為什麼這些人會無緣無故去加或減神的話呢?.
我會認為的是.
其實那群人正正就是覺得.
上帝的說話就是有不夠的地方.
上帝做得不夠好.
所以我幫他加強一點.
上帝想得不完善.
這些多餘的我就幫他減掉.
這群人覺得自己比上帝強.
所以某情況他們的加添與生減.
其實是一份驕傲.
是一份自然.

$^{321}$是選擇放棄了去跟隨上帝的一個表現.
所以第一點.
摩西其實吩咐這群以色列人.
是一定要遵守律法.
要聽神的教訓.
他說的是因為這件事.
正正其實是關乎一個生死的事.
也可以說的是.
我們今天如果是我們這群人.
是我們在檢視我們其實生命是否尊主為大.
你可能可以去想.
究竟我們在上帝的說話裡面.
我們有沒有加添或生減呢?.
倘若有的話.
究竟上帝的說話大一點.
還是你的心意大一點?.
你可以用來成為你檢視自己的一個指標.
這是第一點.
第二點.
看五至八節.
第五至八節其實最主要就是說.
摩西帶出的就是.
當這群以色列人去遵行律法的時候.
其實地上的其他的萬民.
理應是會讚嘆上帝的智慧和公義.
也顯出了原來上帝是跟他們相近.
很接近的一個民族.
這群以色列人理應是這樣.
我們怎麼可以看呢?.
其實之前我有一篇道是講拉合的.
拉合的故事其實正正可以引證到這件事.
拉合這個女孩其實是一個外邦的女孩.
她當日.
當時初初遇見這兩個探子的時候.
她其實並未真的見過什麼以色列人.
但因為她聽到的是耶和華.
曾經怎樣幫這群以色列人.
出埃及過紅海.
擊敗了兩個約旦河河東的王.
就在說.

$^{361}$原來人心就會改變.
耶利哥城的人心都消化.
甚至這一位拉合.
他稱呼上帝這個耶和華.
是上天下地的上帝.
他在以色列人的際遇裡面.
他見到原來真的有神跟他相近.
所以第二點其實摩西提醒這群以色列人的是.
遵行律法是有神跟他靠近.
而這個也是見證神的一個很大的方式.
討論回我們吧,弟兄姊妹.
你的生命有沒有被人見證到神呢?.
又或者我這樣說.
你的生命是否讓人見證到.
我們信的上帝是一個獨一的真神?.
你的生命有沒有人.
因為你的生命一切的事.
你的行事為人.
繼而去讚嘆上帝呢?.
抑或你的生命.
只是被人見到上帝的一小小.
我也信耶穌的.
我也有教會的.
但其餘的行事為人呢?.
甚至其實你的行事為人.
是否讓人見到一個反效果的狀態呢?.
最近有一個很有趣的例子.
大家應該都曾經看過這個內容.
就是黃秋生的例子.
說的是他主演的一套戲.
叫做《不捨之罪》.
不知道什麼原因.
不能夠割續這個金像獎.
被人DQ了.
當他知道這個消息之後.
黃秋生就在自己的Facebook.
一個社交平台裡面.
他就發了一段經文.
就是《詩篇23》的其中一兩句.
重點也不是在於他的發文是什麼.

$^{401}$而是網民看到他發聖經的時候.
有人就很好奇.
就問:黃秋生你是不是信了主?.
黃秋生的回應也很有趣.
就是不是你想的那一種.
我想我們仙姑不論黃秋生是否信了耶穌.
因為這個不是我引用這個例子的重點.
我也沒有打算研究.
我也不認識他.
我著意的就是他的回應.
不是理想的那一種說話.
最初我也有想過.
究竟黃秋生說的是哪一種基督徒呢?.
但是當我繼續去想的時候.
我其實會覺得.
其實說哪一種也不是很重要.
而是這一句說話.
反映的.
是原來世界對基督徒.
是有一些黑板的印象.
有一定的對他們來說.
基督徒這幾個字.
你是基督徒.
他們就會有一些感覺.
有一些觀感在這裡.
甚至可能對於大部分的人來說.
那些形象的觀感.
也不是那麼正面.
所以我會說.
倘若聖經明明剛剛說.
我們遵行神的教訓.
理應是會令人讚嘆上帝的智慧和公義.
理應是會讓人見證到.
神在我們中間.
但是現實.
世界對基督徒.
卻是有另一種的詮釋.
那問題究竟出在哪裡呢?.
我想其實也不是很難明白.
就是我們基督徒.

$^{441}$或者大多數的基督徒.
其實沒有好好地遵行上帝的話.
以致在一些未信的人面前.
失去了所謂的見證.
剛剛也問大家.
你有沒有好好見證上帝呢?.
大家的反應也沒有什麼.
大家心知肚明.
我想其實.
不緊要的.
我們也會有失見證的時候.
我自己也會有.
但我想在新的一年.
剛剛六零年開始.
大家會不會又願意.
再一次重新立志.
好好地用你的生命.
去見證神呢?.
我們每次回來敬拜.
我們很投入.
那些歌詞我們甚至覺得很有共鳴.
但那份共鳴.
是不是只在唱敬拜詩歌的那幾句歌詞的時候呢?.
究竟我們的生命.
我們怎樣走呢?.
當我們舉手敬拜.
當我們很用力.
很渴求地去呼求上帝.
去敬拜上帝的時候.
究竟我們的生命是否在用同一件事.
去反映出來呢?.
這是第二點.
最後.
最後看的是第九至十四節.
剛剛說到其實.
摩西重新了當日在荷列山.
也就是在西乃山.
在說耶和華主動地跟這些以色列人立約.
還將十誡頒佈的情景.
再一次說出來.

$^{481}$摩西他再說.
其實表示的是.
這些律例典章.
其實不是摩西自己作的.
他再一次告訴新一代的以色列人知道.
這些我真是很真實地在上帝裡面得到啟示.
從而得來的.
所以.
當我跟你們說的時候.
麻煩你們.
學學.
怎樣懷著敬畏的心.
去尊重上帝的說話.
去遵守上帝的說話.
而同時.
摩西再一次表明.
他所說的是.
他之所以要不斷地再跟這些以色列人說當日.
耶和華究竟怎樣吩咐那些第一代以色列人那一切的事.
這些事.
不是摩西自己想做的.
是上帝託他做的.
上帝交代他一定要做.
而且.
不只是摩西要做.
這裡說的是.
世世代代的以色列人都要做.
甚至.
其實是.
每一個屬於上帝的人都要做.
畏的就是第十節了.
講的就是令到我們在存活在世的日子可以學習敬畏上帝.
還要教訓兒女都要這樣做.
因為.
為什麼要做呢?.
第一點說了.
因為遵行上帝的誡命是關乎生死的事.
所以要將這件事教訓你的子女.
要傳給你的子子孫孫.
其實是將生命傳遞下去.

$^{521}$那個真正的生命之道.
藉著要他們遵守律法.
繼而傳承下去.
對於以色列人來說.
甚至對於今天的信徒來說.
是一件極其重要的事.
我不記得是什麼情況下.
我曾經被人問過一個這樣的問題.
就是如果有一件事.
你可以傳給你的子女的時候.
你會選擇什麼?.
當時他的問題當然好像很輕鬆.
你可以選擇你的性格,樣貌,身材,財產,社會地位.
後面那兩樣我沒有.
什麼都可以.
他這樣說.
長話短說.
我當時的選擇就是.
跟隨耶穌.
我想的是我的子女能夠一輩子好好跟隨耶穌.
因為在我而言.
我覺得這件事比其他一切都更加重要.
我會覺得其他的事.
就算他本身沒有.
只要有耶穌帶著他.
我相信他們的人生.
最終總會過得豐盛.
雖然我不知道是哪一種的豐盛.
但起碼我仍然相信的是.
上帝會將最好的.
又或者叫做他們所需要的東西.
給他們.
在我們當中.
同樣有很多弟兄姊妹.
越來越多人有下一代.
如果今天我同樣將這個問題問你.
你會選擇什麼呢?.
要選擇將一件事.
傳給子子孫孫.
代表的是.

$^{561}$你有多麼看重這件事.
究竟今天我們.
你可能未有下一代.
有部分弟兄姊妹有的.
但同樣你都可以問問你自己.
究竟跟隨神又或者跟隨耶穌.
這件事對你來說.
是否最重要呢?.
摩西提醒我們.
我們要謹慎.
要殷勤保守我們的心.
免得我們忘記所經歷.
免得這件這麼重要的事.
會離開了我們的心.
會離開了我們的一生.
你見到我刻意highlight了.
即是變回黃色那兩個.
嚴格來說其實它是大了三個房間.
即是大了三格.
由於空間問題.
我又無謂弄太多了.
我特別將「忘記」這兩個字弄了.
我想人真正的問題其實不是不知道.
我想要跟隨耶穌.
對於我們這班人來說.
我們同樣不會不知道.
但是人真正的問題是.
很容易忘記.
就好像第一代的以色列人一樣.
他們是否不知道.
其實耶和華的能力有多厲害呢?.
不是的.
他們這班人走過了那個.
分開了那個紅海.
他們走過了那個.
你看見兩條鯨魚在旁邊走過.
不知道有沒有鯨魚.
總之他們見過魚.
他們見過的.
他們沒理由不知道其實上帝有多厲害.

$^{601}$但是.
這班人卻每每有困難來到的時候.
他們就忘記了.
耶和華曾經在他們中間發生過的事.
他們就忘記了.
原來我們要遵行耶和華所吩咐的說話.
以致他們會選擇用自己的方法.
或者一些過去的經驗.
去做一些決定.
去判斷一些事情.
他們最出名的.
做金牛毒這件事.
就是一個很明顯的例子.
所以在神的裡面.
其實我們不知道那些事情.
或許這樣說我會想.
其實上帝都會說.
他們所作的他們都不知道.
這樣就寬恕了.
但是問題.
我們真正的問題.
不是不知道.
是我們很容易就遺忘.
很容易就忘記了上帝的一些事情.
很容易就忘記了上帝.
在你生命裡面給過你的恩典.
忘記了過去二三十年.
坐在教會聽了那些教導.
或者這樣說.
我今晚說完你出去已經忘記了.
忘記了.
原來跟隨神.
跟隨耶穌.
是一件很重要的事.
你有沒有忘記.
相請敬拜對.
總結.
我今天把我的講題改了個名字.
叫做「聽誰說」.
很明顯就是想跟大家一起去反思一下.

$^{641}$這一刻我們自己的生命.
究竟我們在聽什麼東西去走呢?.
我們究竟在聽這個世界什麼聲音.
去做一些選擇呢?.
摩西很直接.
他提醒.
原來昔日有一群以色列人.
就是因為聽錯了那群摩亞女子的聲音.
就離開了神.
而我們剛才聽.
也有一些以色列人.
因為聽了探子的回報.
他們就選擇.
放棄了那個入迦南地.
得地為業的最佳時機.
還有一些以色列人.
因為忘記了耶和華的吩咐.
他就選擇按照他以前在埃及地.
的一些過去的經驗.
去做了那些金牛毒出來.
這些的忘記.
令到他們大大地得罪上帝.
以色列人的歷史裡面.
有很多的聲音.
這群人聽過很多東西.
而他們也曾經做了很多錯誤的選擇.
摩西要和這群新一代的以色列人.
去回顧這一切的歷史.
為的就是很想他們.
吸取這些前人的教訓.
走回一條對的路.
走回一條耶和華想他們走的路.
我們呢.
今天在這個世界同樣.
充滿無限的聲音.
最簡單的是我們買一樣東西.
上個網都已經要看看評論.
要看看評語.
要看看評分.
很多東西的.

$^{681}$現在AI又多一點了.
什麼都問一下AI.
然後又看看小紅書.
有不同的社交平台.
有不同的評論.
你身邊有很多不同的朋友.
同樣給很多不同的意見.
我們在聽什麼.
當我們去做一個決定的時候.
究竟我們在跟隨世界的不同聲音.
抑或.
你忘記了原來有一件事很重要.
跟隨耶穌.
究竟你看重的是什麼.
你看重的是世界的一件事是怎樣.
人們的評語是什麼.
還是.
你願意選擇.
降服.
去跟隨上帝.
\newpage



\section{創世記 12:1-4}
\label{sec:031715}
\textbf{多代同行•與神同行系列:當人生遇上呼召時}~~~~陳黎韻怡傳道
\newline
\newline
連結: \href{http://www.church.com.hk/acms/content.asp?site=cdc\&op=show\&type=product\&code=031715\&layout=sermon}{\texttt{031715}} ~~~~ 語音日期: 2026-02-22
\newline
\newline
\hyperref[sec:031714]{\small{< < < PREV SERMON < < <}}
~
\hyperref[sec:index_chronic]{\small{[返順時目]}}
~
\hyperref[sec:031716]{\small{> > > NEXT SERMON > > >}}
\newline
\newline
創世記 12:1-4
\newline
\begin{longtable}{cl}
\hline
\hline
章節 & 經文 (和合本修訂版)\\
\hline
12:1 & \begin{tabularx}{0.7\textwidth}{X} 耶和華對亞伯蘭說：「你要離開本地、本族、父家，往我所要指示你的地去。 \end{tabularx} \\ \\ \relax
12:2 & \begin{tabularx}{0.7\textwidth}{X} 我必使你成為大國，我必賜福給你，使你的名為大；你要使別人得福。 \end{tabularx} \\ \\ \relax
12:3 & \begin{tabularx}{0.7\textwidth}{X} 為你祝福的，我必賜福給他；詛咒你的，我必詛咒他。地上的萬族都必因你得福。」 \end{tabularx} \\ \\ \relax
12:4 & \begin{tabularx}{0.7\textwidth}{X} 亞伯蘭就遵照耶和華的吩咐去了；羅得也和他同去。亞伯蘭離開哈蘭的時候年七十五歲。 \end{tabularx} \\ \\
[1ex]
\hline
\hline
\end{longtable}
$^{1}$我們今天講的人物就是亞伯拉罕.
亞伯拉罕的原名叫做亞伯蘭.
後期神為他改名叫做亞伯拉罕.
而亞伯拉罕又被稱為信心之父.
當我一路預備這篇講章的時候.
我自己也會想.
亞伯拉罕其實何時開始認識耶和華是神的呢?.
我們說祂是一個信心之父的時候.
究竟祂的信心.
是否又是因為祂原生家庭的幫助.
以致祂有這麼大的信心呢?.
我就開始看聖經的時候.
聖經有這一段記載.
它說 古時你們的列祖.
就是亞伯拉罕和拿學的父親塔拉.
住在大河那邊 事奉別臣.
經文裡面 哪個是亞伯拉罕的父親?.
塔拉.
塔拉原居在吾爾.
吾爾在哪裡呢?.
這裡說的是大河那邊.
大河即是說 猶法拉底河.
而在螢光幕裡面.
你會看到有一個紅色的點.
就是吾爾的地方.
正正就是亞伯拉罕河的附近.
而經文裡說的.
塔拉在那裡 大河那邊 事奉別臣.
事奉哪個別臣呢?.
考古學家一路去發掘的時候.
就發現在吾爾這個地方的遺址.
有五座 還是六座 有幾座的神廟.
這些神廟都是用來敬拜月神的.
而這個神廟最高的那座.
有大概六,七層樓那麼高.
你今天可能會說六,七層樓不算什麼.
但是當時 亞伯拉罕的年代.
以當時的建築技術.
其實他們的房子只有一層 兩層.
所以你可以想像到.

$^{41}$當你遠遠地看著的時候.
其實整個城市都會看到月神這一座的神廟.
這一座神廟彷彿標誌著.
吾爾這個地方是用來敬拜月神的.
所以我們知道的亞伯拉罕.
其實不是因為家裡的緣故.
他知道耶和華是神.
因為塔拉其實也在這裡敬拜月神.
但是聖經也沒有用一個很詳細的篇幅.
去記載亞伯拉罕的信仰如何被建立.
不過反而用很長的篇幅.
去交代亞伯拉罕兩次的回應呼召.
為什麼聖經會這樣記載呢?.
因為亞伯拉罕兩次回應神的呼召.
這才是聖經想我們留意的重點.
這個重點就是人生是回應呼召的一個旅程.
我們知道人生是一個旅程.
但是這個旅程可以發生很多的事情.
而這個旅程也是我們回應呼召的旅程.
我剛才說過.
呼召不單止是神會呼召人全職服事.
神的吩咐同樣都是對人一個的呼召.
而在這個旅程當中.
當我們說回應呼召的時候.
其實就是神很想和我們在一起.
祂很想和我們建立一個很親密的關係.
而這份很親密的關係何時開始的呢?.
其實就是在神創造天地的時候.
這份關係已經存在.
我會形容這份關係就好像一對熱戀的情侶.
一對熱戀的情侶.
你問他要做些什麼.
去哪裡吃飯.
去哪裡喝東西.
卡拉OK還是去看電影.
不重要的.
最重要的是什麼.
最重要的是一切.
這份親密的關係最重要的是什麼.
是彼此知道對方在自己心目中.

$^{81}$在對方心目中有一個很重要的位置.
所以如果今天上崇拜.
坐在你旁邊的是你的家人.
坐在你旁邊的是你的情侶的時候.
你們有沒有一個親密關係呢?.
你可以看著他.
你問他.
你心裡有沒有我?.
如果你的親人轉頭的話.
你回去好好跟他長談.
是不是?.
一個親密的關係.
心裡面你會記掛著這個人的需要.
這個關係.
你不需要介懷大家做些什麼.
大家一起就好了.
而這份關係本來在創世的時候.
就已經存在.
但是當我們一路看聖經的時候.
我們發現.
由阿當夏娃吃禁果.
接著人的罪惡.
神他就呼召羅亞起方舟.
多多後來.
人又為了自己的明星.
去起巴別塔.
一次又一次.
歷史的片段告訴我們.
揭示人類那份驕傲.
人類的任性.
人類的忘本.
當我們有事相求.
又大又難的困難的時候.
我們就說.
神啊,你幫助我吧.
好奇,神啊,你來吧.
當問題解決了.
可能我們接著.
我們忘記了神是誰.
就算我們記得神是誰.

$^{121}$可能我們又會跟神說.
神啊,我想做自己的事情.
麻煩你先坐一邊.
不要打擾我.
我們一路一路.
人類的歷史告訴我們.
我們經常傷害神的心.
但是因為神對我們那份的愛.
祂很著緊我們跟祂的關係.
於是乎.
在歷史裡面見到.
神千方百計的.
用一些方法去重建.
我們跟神的親密關係.
到了亞伯拉罕的年代.
神就揀選了亞伯拉罕.
去成為重建這個關係的核心人物.
神要重建熟悉祂的群體.
神去跟亞伯拉罕說.
祂說:你要離開本地.
本族,父家.
往我所要指示你的地方去.
我必叫你成為大國.
我必賜福給你.
叫你的名為大.
你也要叫別人得福.
神去揀選亞伯拉罕.
祝福亞伯拉罕說.
我會叫你成為大國.
但是你想想.
亞伯拉罕是生於.
公元前2100年的時候的人.
在公元前2000多年.
其實那時候沒有國家的出現.
只是有城邦.
當這句話說.
你要成為大國的話.
亞伯拉罕根本沒有任何的概念.
什麼叫做大國?.
原來大國在原文裡面.

$^{161}$還有一個意思.
是形容一群蝗蟲的出沒.
大家有沒有見過蝗蟲?.
蝗蟲不是一隻出現.
是群居的.
而一群的蝗蟲.
我看一些昆蟲學說什麼呢?.
蝗蟲出現不會是一千隻.
不會是一萬隻.
一出來的話會多到數百億隻.
數百億隻.
就好像螢光幕.
你看到一群的蝗蟲出現的時候.
好像全部被蓋住.
太陽你都看不到.
掩蓋著整個城邦.
原來神祂要祝福亞伯拉罕.
祂要有眾多的子孫.
遍滿全地.
如果我說眾多子孫.
遍滿全地.
神今天新年的時候.
祂說要祝福大家.
你要有一百億個子女.
你想不想要?.
這麼多堂大家都搖頭.
一百億個子孫.
今天我們會想起.
嘩 改名字都煩了.
難道我叫他.
阿一 阿二 阿三.
數到數百億個嗎?.
但是我們要進到亞伯拉罕那個年代.
亞伯拉罕的時候.
他是務農為業.
那時候很需要有人手去耕種.
你有人手就有耕種.
有耕種就有收成.
有收成就是說你家裡有飯開.
所以是很重要的事.

$^{201}$並且當你子孫眾多的時候.
你就可以一路的.
將你家族的血緣.
一路的擴展 一路的承繼出去.
所以如果當時.
就是說新年的祝福.
我祝你兒孫滿堂.
當時的人個個都搖頭.
包括亞伯拉罕他都會搖頭.
「我要啊 我要啊」.
但是神就說.
如果你要這一份祝福的時候.
你首先要回應神的吩咐.
神吩咐什麼呢?.
「你要離開本地 本族 父家.
往我所要指示你去的地方」.
咦 這裡說.
神吩咐亞伯拉罕.
你要去神指示的地方.
即是去哪裡呢?.
我們再看一段經文.
這段經文就說 亞伯拉罕因著信.
夢照的時候就遵命出去.
往將來要得遺孽的地方去.
出去的時候.
還不知往哪裡去.
去哪裡 是還不知.
不知道去哪裡 總之你又去.
你想想我們今天.
當我們說要轉工的時候.
我們都會想一下.
轉了那份工.
那裡的工資.
是不是能夠應付到我日常開支呢?.
我們今天有人選擇會去移民的時候.
都會想一下.
我移居那個地方.
我有沒有能力好好地生活呢?.
但是亞伯拉罕.
他知不知道去哪裡呢?.

$^{241}$還不知.
他有一件事他知道的.
就是當他看聖經的時候.
他知道的就是他的太太.
撒拉呢.
她是不能生育.
所以他不知道去哪裡.
但是他知道的就是.
太太不能生育.
面對著神的祝福.
就說你要成為大國.
然後就說你要去.
我會用六個字形容神對他的吩咐.
就是高風險.
沒有可能的事.
高風險 因為我不知道去哪裡.
沒有可能 因為太太已經知道她不能生育.
比喻我們今天.
神跟我們說一些我們覺得是高風險.
沒有可能的事.
我們馬上會怎樣呢?.
老師問問題.
大家馬上低頭.
不要看著我.
最好不關我的事.
又或者說.
我想清楚了.
我擔心.
然後就 o下 o下.
聽不清楚.
然後就收線.
當沒事發生.
但是亞伯拉罕呢.
亞伯拉罕他盡明.
他說他願意用生命去經歷.
這個掌管萬有的神.
所以神揀選亞伯拉罕.
不是因為他英俊.
不是因為他能力夠.
他揀選亞伯拉罕.

$^{281}$因為神揀選的人.
是在問你的心願不願意.
當亞伯拉罕有一個願意的心的時候.
神就揀選他.
而我一直在預備這段經文的時候.
我會想.
神其實今天對我們有很多的應許.
不知道你有沒有數過.
我問AI 人工智能.
我說究竟聖經裡面有多少個應許呢?.
然後人工智能告訴我.
有七千多個應許.
然後我就會問自己.
我們今天有沒有亞伯拉罕的信心.
聖經說什麼.
我就聖經這樣說.
我就這樣相信.
我們有沒有亞伯拉罕.
對神那份的期待.
神啊.
你在聖經裡面.
你說過的那些應許.
我期待的是用我自己的生命.
我親自去經歷.
你那份的真實.
你對我的那份的愛.
那份的大能.
弟兄姊妹.
當我們看聖經的時候.
我們要知道.
聖經不單只是一些文字.
聖經更加是上帝那份的能力.
當人生遇上枯燥的時候.
亞伯拉罕遵命.
他就踏上回應神那份.
吩咐回應神呼召的旅程.
當他踏上這個旅程.
我們會看到他的起點在吾爾.
終點呢?.
終點叫做暫時保密.

$^{321}$到時機成熟的時候.
神就會告訴亞伯拉罕.
但在旅途裡面.
神有應許.
他會一直與亞伯拉罕同行.
幫助他生命有兩大突破.
第一個突破我會看到的.
就是不被世界定義的價值.
大家想想.
亞伯拉罕原居的地方在吾爾.
除了剛才我提到的.
就是在遺址發掘到的神廟之外.
其實考古學家.
一直發掘在遺址就發現了.
有一些在多年前.
亞伯拉罕年代的物件.
包括什麼呢?.
有一個石刻.
石刻大家看到照片裡面.
是一隻腳.
踩著下面有兩個.
好像人形的物體.
去托著這隻腳.
其實這個石刻代表什麼呢?.
是表達.
當如果你是戰勝一方的時候.
你就可以用腳踐踏戰敗者.
說的是當時的社會.
太弱留強.
一個榮辱的文化.
考古學家又發現.
又有一些數學的例板.
不是今天我們花學的學生去學習數學.
在亞伯拉罕的年代.
他們已經要學習算式加減乘除.
有學校.
而發現也有一些圓琴.
珠寶.
金飾.
各樣各樣的東西.

$^{361}$告訴我們這是一個非常富裕.
和文明的城市.
當你能夠攀附這個上流社會的時候.
你就能夠享受.
物質主義帶來給你的那份享樂.
帶來你一個.
好像人有一個超然的感覺.
一個很有優越感.
接著我就想起.
在十多年前.
當時我還在外面商界工作的時候.
話說.
當時任職的公司.
因為要控制成本.
就有一個安排.
就是說如果那些不需要見客的部門.
就要由中環搬去九龍灣這個地方.
十幾年前的中環.
仍然是像今天一樣.
是香港的核心商業的一個地點.
所以有很多專業人員在那裡發聲.
但是十多年前的觀塘.
九龍灣的時候.
是什麼環境呢?.
其實當時是一個新發展區.
是由工廠改建而成的.
很多後勤的部門.
都會搬到那裡.
當公司裡面的同事知道.
要搬辦公室.
在中環的同事就有很大的反應.
又投訴 又說要辭職.
好像以死相怯那樣.
那你問他.
為什麼你不搬去九龍灣 觀塘那邊.
他們就說.
你看我一身的行政裝束.
我的外套很貴的.
我的鞋是名牌的.
我在各條中環的街道走的時候.

$^{401}$這個是我的身份.
我是一個專業的高級的行政人員.
但是當我要搬到九龍灣觀塘的時候.
我要換來的.
就是一個叫做.
Business Casual.
就是一個便服的時候.
心裡面就好像我被降格了.
確實是的.
我們香港的生活.
是很急速.
很具競爭力.
我想我們每一個人.
都很想自己走在成功的道路上.
剛剛過去那兩天.
大家在街上有沒有看到很多人在排隊.
大家點頭都知道我在說什麼.
很多人都希望我會不會中獎.
我會不會拿到什麼什麼彩.
我們有時候都很想成功.
我很想有一個優質的生活.
我未必說我在群體裡面.
我要成為第一那個.
但是我們不想的.
就成為最後.
好像比下去那樣.
所以香港人很努力的去拼搏.
而我自己會想到.
就是我那時候在工作裡面.
我曾經見識過有一位高層.
我們就這樣稱呼他.
我們叫他做美女與野獸.
美女為什麼呢.
美女因為她一身的打扮很高貴.
都是拿起手.
大家看到什麼.
是我的手.
什麼都沒有.
但是你在她的手裡會看到.
哇 閃亮亮的東西.

$^{441}$所以我們稱呼她美女.
但是同時間.
為什麼我們會稱呼她為野獸呢.
就是當一些同事可能做一些事.
令到她覺得不稱心如意的時候.
她就會當眾去羞辱她.
跟她說.
你這麼蠢的.
你簡直是蠢材.
是這個團隊裡面的負累.
當時有些同事因為要開工.
要餵兩餐.
不敢出聲.
但是慢慢慢慢.
她就會開始認同別人對她那份說話.
對她那份詆毀.
接著就開始對自己沒有信心.
工作越來越差.
我們處身的都市.
我們的環境.
可能都充斥著這種榮辱文化.
但是我們要知道的就是.
在地上沒有一個人.
能夠真正明白到我們自己人生的價值.
唯有的是我們的造物主.
他才知道我們的真正價值是怎樣.
你想想外面的人.
他們怎樣去衡量我們的價值.
不就是這個社會價值.
今天看重些什麼的話.
你有在你身上的你就有價值.
但是我們的人生價值.
不是隨著這個社會改變去改變.
而是我們會問的就是.
我們的造物主.
當他做我們的時候.
他視我們為寶貴.
我們每一個都是按著神的形象去做的時候.
倘若我們身邊可能有些人說.
你這麼蠢的.

$^{481}$你這樣都不知道的.
哎呀 你沒有價值的時候.
我會告訴他.
其實你不識貨而已.
因為我是被神所造.
我是何等美麗.
而當我們說.
我們知道我們是被神所造.
我們是何等美麗的時候.
都讓我們學懂的就是.
在我們的群體相處的裡面.
耶穌祂曾經有一個教導.
祂說凡話底輕是魔力的 難免地獄的火.
魔力的意思就是愚蠢.
蠢鈍 沒有價值.
倘若我們這樣去視我們身邊的人的時候.
難免受地獄的火.
神看我們每一個都為寶貴的時候.
我們要懂得保護我們自己的心.
當這個世界充斥著很多.
可能是貶低我們.
詆毀我們 拒絕我們說話的時候.
我們要把焦點不是放在那些地方.
令我們變成一個腐毒的人.
我們的焦點要放在神的身上.
神告訴我們我們是被愛的.
我們是寶貴的.
今天好像也有些陽光.
完了崇拜 大家試試出去外面.
不要低頭只看手機.
抬高頭看看.
神為我們預設的天空大地.
是神因為愛的緣故為我們預備的環境.
我們打開聖經.
聖經其實是神寫給我們的情書.
神是歷世歷代的道.
祂告訴我們.
祂不斷地尋找.
用盡方法去重建我們和祂這份愛的關係.
讓我們把焦點放在神的身上.

$^{521}$神去吩咐我們.
神去呼召我們.
讓我們懂得用屬神的眼光去看自己.
去看成敗得失.
我們的安全感不是在地上任何的物質.
我們的安全感是因為神對我們的那份愛.
對我們的那份肯定.
或者我們有時做的決定.
或者我們的擁有不比別人多.
但是神說.
我們把焦點放在神身上的時候.
我們心靈是有一個極大的富足.
而第二樣東西.
我會見到神帶給亞伯拉罕那份突破.
就是他能夠超越既定的限制.
不知道大家有沒有試過.
我們自己知道神是很愛我們的.
我們被神視為寶貝.
我們都嘗試過去回應神.
神啊 你呼召我.
叫我幫你做些什麼.
或者你呼召我進到什麼嘗試裡面.
我有的.
但是走著走著.
我又覺得我自己好像沒有能力.
很多的不能夠.
接著我就停下來.
甚至我放棄了.
神你起初吩咐我所做的事.
我們又看看.
其實亞伯拉罕的人生.
他如何被神塑造.
以至他能夠突破這些.
他好像很多的關卡.
他不能夠跨過人生的很多限制.
我們看看神他第一次.
去呼召亞伯拉罕的時候.
經文記載他說.
他拉帶著他兒子亞伯蘭和他孫子哈蘭的兒子羅德.
並他兒婦亞伯蘭的妻子撒萊.

$^{561}$出了加拿大的誤議.
你想想.
神其實是呼召亞伯拉罕的.
但是在經文第一次呼召的時候.
誰帶著一群人出誤議.
是他拉.
他拉是誰.
他拉是亞伯拉罕的爸爸.
為何不是亞伯拉罕自己回應呼召.
自己走出去.
如果我們再看以下這些經文.
我們就明白了.
在11章26節.
他說他拉活到七十歲.
生了亞伯蘭 拿學 哈蘭.
這裡說的是他拉 即是亞伯拉罕的爸爸.
爸爸在七十歲之後.
相繼生了三個兒子.
有亞伯蘭 即是亞伯拉罕 拿學和哈蘭.
不過接下來.
經文也記載.
他拉共活了二百零五歲.
死在哈蘭.
亞伯蘭出哈蘭的時候.
七十五歲.
即是說.
亞伯拉罕他們有三兄弟.
接著爸爸二百零五歲.
在哈蘭死了.
他們在哈蘭之後.
爸爸死了.
亞伯拉罕.
他那時七十五歲.
他就離開了哈蘭.
繼續他呼召的路.
如果我們將二百零五減七十五的時候.
即是說.
亞伯拉罕是他拉一百三十歲所生的兒子.
即是說 亞伯拉罕他不是長子.
而在猶太人的文化裡面.

$^{601}$爸爸就是家主.
負責家裡所有的重要決定.
況且亞伯拉罕他做哥哥.
對於亞伯拉罕來說.
他的家庭崗位角色是什麼呢?.
就是.
我是小兒子.
有什麼事 爸爸決定.
爸爸決定完 哥哥又去做.
死跟著走就可以了.
所以亞伯拉罕縱然他很願意回應神.
但是他的心理上.
他很想倚靠爸爸.
他拉成為了亞伯拉罕的安全區.
而我們又會問.
為什麼他拉又會離開吾爾呢?.
原來他拉的家庭發生了一件事.
令他拉決定離開吾爾.
發生了什麼事呢?.
在11章26節記載.
他說 哈蘭死在他的本地.
加拿大的吾爾.
在他父親他拉至善.
他拉為什麼會離開吾爾呢?.
因為其中一個兒子哈蘭死了.
白頭人送黑頭人.
他的心很痛.
很沉重 很大的打擊.
於是他拉決定離開這個傷心地.
帶著他的家人去到哈蘭.
但是去到哈蘭之後.
經文繼續記載的就是.
「住在那裡」.
「住」不只是停下.
這裡還有「定居」的意思.
停在那裡定居.
家人停下來.
亞伯拉罕都停下來.
直到他拉離世的時候.
神的呼召才再次來到.

$^{641}$我們相比起亞伯拉罕的兩次呼召.
第一次 神跟他說.
「你要離開本地和親族」.
「往我所要指示你的地方去」.
這是第一次的呼召.
第二次 當神再呼召亞伯拉罕的時候.
他就說「你要離開本地,本族,父家」.
「往我所要指示你的地方去」.
多了一個字 是不是?.
那個是什麼字?.
「父家」.
「父家」就是他拉.
剛才提到的.
對於亞伯拉罕來說.
這是他的安全感.
對於亞伯拉罕來說.
他會覺得.
「我不行了」.
「你要我嘗試一些新嘗試」.
「一直以來我的成長」.
「都是有爸爸在我身邊」.
「沒有爸爸在身邊的時候」.
「我應付不了一些新的挑戰」.
亞伯拉罕的人生.
他需要成長的時候.
要有一些突破.
不過如果第一次.
神已經馬上跟亞伯拉罕說.
「你要生命成長 你要突破」.
「所以你要突破的時候」.
「你就要離開父家」.
我想亞伯拉罕那時候聽到.
他會有什麼反應呢?.
「什麼?你要我離開父家?」.
「為什麼?我都不覺得有什麼問題」.
「你為什麼要逼我去做一些我不明白的事情?」.
神就是一個這麼好的教練.
祂知道我們人生的進程.
祂知道我們需要突破的地方.
但是祂更加是一個很有智慧的神.

$^{681}$也很有忍耐能力的神.
祂知道何時出手才是最好的時機.
當你第一刻跟祂說的時候.
亞伯拉罕會不明白.
我的心都不接受的時候.
「你給我改些什麼?」.
「我只不過是行為上的改變」.
「我不會真的信服於你對我說的話」.
但是當亞伯拉罕停下來這四十年.
直至他拉離開這個世界的時候.
亞伯拉罕明白了.
人的陪伴很有限.
唯有神的陪伴才可以七成二十四.
一直陪伴著亞伯拉罕的成長.
繼續向前走.
我會想 其實亞伯拉罕這四十年的停低.
其實他也不好受.
因為他很清楚知道是神的呼召.
他也踏上了呼召的路.
但是他也因為自己的軟弱停下來了.
他會一直反思.
究竟人生我欠缺什麼.
直至到爸爸離世 他明白.
弟兄姊妹 我們的成長過程裡面.
有沒有一些話 神一次又一次地跟你說.
有沒有一些事件 不停不停地發生.
可能這些就是神繼續鼓勵你成長.
鼓勵你突破的一些信號.
告訴你 神是你最好的依靠.
當亞伯拉罕再次啟程的時候.
這一刻 神才會跟亞伯拉罕說.
你要去的地方 你去的目的地.
而當亞伯拉罕面對著再一次出發的時候.
他做了一件事.
就是每到一個新的地方 新的景點的時候.
他都在那裡祝壇獻祭.
可能我們會說 敬拜神而已 沒什麼特別.
是不是真的敬拜神而已 沒什麼特別呢?.
如果我用一個比喻去說.
大家可能就會明白.

$^{721}$當我們去旅行的時候.
導遊告訴你.
這裡就是這個地方的特色 一個景點.
大家會怎麼做?.
馬上在那裡排隊 拿個相機出來拍.
你拍照不是拍景點的.
你也會叫團友 過來過來過來.
我們一起在這裡拍張照.
留為紀念.
因為我們很想記下的就是.
在這個地方 有我所愛的人跟我一起經過過.
亞伯拉罕祝壇獻祭.
每一個新的地方 一路向前進.
去到一個新的地方 他就在那裡祝壇獻祭.
紀念神跟他同行的每一個片段.
亞伯拉罕的父家是對神的依賴.
今天 我們弟兄姊妹 我們的父家是什麼?.
我們的父家會不會有一些.
我們很牢固的思維模式的型類?.
我們的父家會不會是一些我們的壞習慣?.
我們的父家會不會是我們在家庭崗位角色裡面.
我們好像亞伯拉罕小時候.
我們經常很想去取悅別人的心.
而沒有問過自己 我可以怎麼去成長?.
神今天告訴我們.
祂是帶領我們生命成長的神的時候.
當祂的呼召來到 祂幫我們過關.
一關又一關叫我們成長.
我就回想起 我自己回應呼召的路程裡面.
其實我都好像經歷過亞伯拉罕迷惘的過程.
我記得那時候 當神的呼召來到的時候.
亞伯拉罕說他不知道他要往那裡去.
而我不知道是什麼呢?.
我不知道下一步 我應該要怎麼樣?.
然後有人跟我說.
他就說 你回應呼召.
你明天就回去馬上辭職 去讀神學.
讀完神學就全職事奉.
我說 啊 明天就回去辭職 讀神學 全職?.
不行啊 那時候我還在供樓.

$^{761}$我家人都還不知道我有回應那份心志.
我覺得我還沒有準備好.
接著我就不知道怎麼辦.
就好像亞伯拉罕那樣.
如果我答應了神 我不回應.
我又好像失信於神.
但是我勉強 我又勉強不來.
因為你問我 我的父家是什麼?.
可能就是我的生活質素.
可能就是我希望家人對我的一份認同.
接著在禱告裡面 神跟我說.
你用兩年時間去準備自己.
兩年時間去準備自己.
嗯 好 準備什麼呢?.
我想 當然是存錢.
因為供樓嘛.
我就說存錢.
還有找機會跟我的家人去表達我那份心志.
在這兩年的時候.
其實是在工作的環境.
我在銀行做一個部門.
就是幫忙去重組 部門重組.
部門重組是做什麼的?.
就是你要去提議一些方案.
讓公司去削減營運開支.
而最大的空間可以去削減開支的.
就是開刀 削人 人力資源.
對於我來說 我會想到這裡.
我說我很不想去做這件事.
因為我不僅是削減一個人頭.
其實每一個人都是背負著一個家庭的負擔.
我不想影響這麼多個人的家庭.
我跟神說 我不想去做這件事.
有一天 當我又是為了這些事情很惆悵.
因為第二天又要跟大老闆去表達.
告訴他我有什麼辦法.
獻計可以削減多少營運開支的時候.
當時突然聖靈跟我說.
祂說 你跪下祈禱.
接著我看一看.

$^{801}$其實那時候都很晚.
因為銀行的同事都做得很晚.
我自己部門的同事都走了.
但是我說 跪下祈禱?.
咦 我隔壁左右的部門的同事還坐在那裡很努力工作.
他們走過的時候.
看到我無緣無故跪下.
會不會打999?.
還是不知道我在做什麼?.
我說不要吧 很難看.
神很清楚地說.
你跪下祈禱.
接著我盡力.
我就在我自己的位置上跪下.
當我一跪下的時候.
我當然將我自己裡面很多很多的牢鎖.
我跟神說 神啊 你知不知道那些老闆真的無良.
神啊 我真的很不想去禁止什麼什麼.
我講完我自己要做的事的時候.
接著聖靈就跟我說.
祂說 你看不看到你的驕傲?.
你經常都以為你自己有點小聰明.
過五關斬六將 做了很多的事情.
祂說 你知不知道.
你經歷過由你出生到今天.
你人生進程的裡面 你經歷過很多又大又難的關卡.
誰牽著你的手過?.
接著到了那一刻.
我馬上跟神說.
神啊 我真是壞女兒.
我真是對不起.
我說我驕傲.
這次我再一次地在你面前.
我說我將我工作的裡面.
那份戰場.
我再交給你.
我說我不要我的小聰明.
但是我要的是.
神啊 你陪著我.
因為我很害怕.

$^{841}$我不知道如何繼續面對我的工作.
當我這樣祈禱完之後.
神給了我一本經文.
《二賽亞書》59章.
它說.
耶和華的棒臂並非縮短.
不能拯救.
耳朵並非發沉.
不能聽見.
神跟我說.
你不要再看你的手.
你看的是看神的棒臂.
當你看你的手的時候.
你在看你自己的能力.
當你看耶和華的棒臂的時候.
你看的是神那份的能力.
神對我那份的愛.
對我由始至終那份的不離不棄.
當神這樣跟我說的時候.
我只有一個反應.
就是願你旨意成就.
我說神 我交給你.
你繼續為我打這場美好的仗.
確實我看到神那份的信實.
在那兩年的裡面.
其實工作很艱難.
很大壓力.
但是神的信實.
祂說祂幫我打這場仗.
直至到我離開那個崗位.
進到去教會服侍的裡面的時候.
我竟然是.
我都不明白為什麼.
在公司的很多安排.
同事那份的信任下面.
我是一個員工.
都不需要去削減.
我知道是耶和華神.
祂的棒臂為我打這場仗.
再者的 原來神不是要我在這兩年的裡面.

$^{881}$只是去準備那些很現實的問題.
那些錢啊 怎樣跟家人交代.
神原來在這兩年的裡面.
最珍貴的是預備我的生命.
當我說要進到去教會服侍的時候.
教會的服侍.
不會比外面職場的工作來得容易.
甚至乎是更難.
因為教會的裡面.
我們說的是生命.
人的生命.
我們不能夠去控制.
他的成長步伐.
我們用愛心的服侍.
是不是弟兄姊妹.
立刻能夠成長去回應神呢?.
我們控制不了.
我們也不能去控制的就是.
我用十分地付出去愛.
很愛很愛弟兄姊妹的時候.
可能有時又會.
有些被拒絕.
被誤會.
很多的艱難在裡面.
又不容易跟別人分享.
但是原來.
在過去我跪下祈禱的裡面.
神今日仍然提醒我.
你不要用你自己的聰明.
你跪下.
你去禱告.
謙卑.
你看的是耶和華的旁碑.
祂一直去擁抱.
一直去帶領.
因為我在祂的裡面.
祂牽著我的手.
祂是我最親密最親密的神.
今日.
既然我知道神牽著我的手的時候.

$^{921}$我就可以很自由地去服侍.
我可以用我的生命.
去好好活下去.
然後將我的生命.
投資在永恆的裡面.
而亞伯拉罕.
他兩次的回應呼召.
確實都告訴我們.
我們的眼界.
不在地上面.
我們不要忘記.
我們天上面.
還有一個美麗的家鄉.
人生是一直的.
就是神不停不停地鼓勵我們.
跟我們有很多的吩咐.
有很多跟我們愛的說話.
祂告訴我們什麼.
祂的愛要引導我們.
塑造我們的生命.
人生的終點.
不是地上.
是在天上的時候.
就算亞伯拉罕.
他終結了他的一生.
那時候他只有一個兒子.
以撒.
神那份大國.
有眾多的子孫.
在他現實的眼裡他還沒看到.
但他仍然相信.
神是那個永活的真神.
神必然地用祂的全能性.
去成就祂自己的說話.
我們活在地上.
我們只要好好享受.
神跟我們同行.
然後我們等候.
到我們再見主面的時候.
神稱呼我們.

$^{961}$是忠心良善的好僕人.
我們就已經心滿意足.
所以弟兄姊妹.
今天亞伯拉罕.
他一生回應神的呼召.
告訴我們.
我們人生最重要的.
其實就是牽緊神的手.
跟神有一個很親密的關係.
無論我們說我們是高山低谷.
不要緊的.
因為神的傍臂與我們同在.
\newpage



\section{詩篇 121:1-8}
\label{sec:031716}
\textbf{繼續前行 (華/粵)}~~~~藍鷺平姊妹
\newline
\newline
連結: \href{http://www.church.com.hk/acms/content.asp?site=cdc\&op=show\&type=product\&code=031716\&layout=sermon}{\texttt{031716}} ~~~~ 語音日期: 2026-02-22
\newline
\newline
\hyperref[sec:031715]{\small{< < < PREV SERMON < < <}}
~
\hyperref[sec:index_chronic]{\small{[返順時目]}}
~
\hyperref[sec:031717]{\small{> > > NEXT SERMON > > >}}
\newline
\newline
詩篇 121:1-8
\newline
\begin{longtable}{cl}
\hline
\hline
章節 & 經文 (和合本修訂版)\\
\hline
121:1 & \begin{tabularx}{0.7\textwidth}{X} 我要向山舉目，我的幫助從何而來？ \end{tabularx} \\ \\ \relax
121:2 & \begin{tabularx}{0.7\textwidth}{X} 我的幫助從造天地的耶和華而來。 \end{tabularx} \\ \\ \relax
121:3 & \begin{tabularx}{0.7\textwidth}{X} 他不叫你的腳搖動，保護你的必不打盹！ \end{tabularx} \\ \\ \relax
121:4 & \begin{tabularx}{0.7\textwidth}{X} 保護以色列的必不打盹，也不睡覺。 \end{tabularx} \\ \\ \relax
121:5 & \begin{tabularx}{0.7\textwidth}{X} 保護你的是耶和華，耶和華在你右邊蔭庇你。 \end{tabularx} \\ \\ \relax
121:6 & \begin{tabularx}{0.7\textwidth}{X} 白日，太陽必不傷你；夜間，月亮也不害你。 \end{tabularx} \\ \\ \relax
121:7 & \begin{tabularx}{0.7\textwidth}{X} 耶和華要保護你，免受一切的災害，他要保護你的性命。 \end{tabularx} \\ \\ \relax
121:8 & \begin{tabularx}{0.7\textwidth}{X} 你出你入，耶和華要保護你，從今時直到永遠。 \end{tabularx} \\ \\
[1ex]
\hline
\hline
\end{longtable}
$^{1}$弟兄姊妹們,新年快樂!.
新年快樂!.
大家知道今天是正月初幾嗎?.
大家知道今天正月初幾嗎?.
初六了,對.
初六了,是的.
通常我們春節應該是我們最大的節日了.
然後我們有一個習慣就覺得.
過了春節才算一年真的開始,對不對?.
通常春節都是我們一年的開始.
過了春節才算一年的開始.
我們是很看重一些節日的禮節.
比如說我們過年要吃什麼.
等一下大家會嘗到不同地方.
弟兄姊妹他們帶來過年的傳統食物.
我們都很注重傳統.
今天弟兄姊妹帶來很多傳統的食物.
猶太人和我們華人有一點很相似的地方.
猶太人和我們華人有一些很相似的地方.
他們也是很看重一些節日的禮儀.
他們都很看重節日的禮儀.
猶太人他們成年的男子.
他一年要有三次帶著家人.
他們要從自己的家去到聖殿.
猶太人的男子有一年要三次帶著一家人一起去聖殿.
大家可以看到.
三次節日就是一次在愉悅節.
一次在祝蓬節.
一次在無殉節.
三次節日就是一次在愉悅節.
一次在無殉節.
一次在祝彭節.
那今天我要跟大家分享詩篇的經文.
它就是在講以色列人.
他從走在路上去到聖殿.
這一段路上有關的.
今天我們講詩篇裡面的一段經文.
就是講他們一路上去過節的時候是怎樣.
我想大家我們可以一起來讀這一篇的詩篇.
它其實不長只有八節.

$^{41}$我們大家可以用自己熟悉的語言來讀.
我們可以一起讀這幾節經文.
有八節的經文.
我們用廣東話也可以.
用國語也可以.
用英文也可以.
好我們一起來.
詩篇一百二十一篇.
我要向山舉目.
我的幫助從何而來.
我的幫助從遂天地的耶和華而來.
祂必叫你的腳動搖.
保護你的必不打頓.
保護以色列的也不打頓.
也不睡覺.
保護你的是耶和華.
耶和華在你右邊任被你.
白日太陽必不傷你.
夜間月亮必不害你.
耶和華要保護你.
免受一切的災害.
祂要保護你的性命.
你出你入.
耶和華要保護你.
從今時直到永遠.
大家有沒有hiking的經驗.
大家有沒有行山的經驗呢.
我有一次參加我們教會小組.
姐妹們的hiking小組.
然後她們很厲害.
她們輕輕鬆鬆就走了五六個mile.
有一次我參加一個姐妹的hiking.
她們很厲害.
輕輕鬆鬆就走了五六個mile.
但是我去了那一次以後.
就不敢再去了.
去了那一次之後我不敢再去了.
因為我去了以後.
走完那一次我回家躺了兩天.
去了那一次之後.

$^{81}$我回家睡了兩天.
其實大家知道嗎.
以色列人他們從自己家走到聖殿.
以色列人由自己家走到聖殿.
其實是一條很遠的路.
其實是一條很長很長的路.
大家可以看到地圖上面.
如果你住在加利利.
在加利利.
對 最北邊的地方.
你要走到耶路撒冷.
要走到耶路撒冷.
大家看到很多黃色的.
那應該就是山.
都是山路.
黃色的應該都是山路.
然後你想想.
她們在走這段路的時候.
她還帶著她的家人.
走這段路上面.
還要帶著家人 一家大小.
她還要帶上她獻祭給神的東西.
還要帶上獻祭的東西.
獻墨.
大家想一想那個場景.
古代的時候那個路.
不像我們現在的路.
沒有那麼好走.
想像一下那個場景.
以前的路沒那麼好走.
很多時候她們只有11路公車.
大家知道是什麼嗎.
很多時候只有11路公車.
大家知不知道是什麼.
她的兩隻腳只能靠腳走.
就是兩隻腳11路公車.
剛剛看到說耶路撒冷這一段路都是山.
所以她們其實走的這一段路.
是一條翻山越嶺的路.
所以她們走這條路.

$^{121}$是一段翻山越嶺的路.
我們剛才讀這個詩篇的時候.
它有四個字.
讀這個詩篇的時候有四個字.
它寫上行之詩.
上行之詩.
詩篇裡面其實有15篇的詩歌.
它都標誌了上行之詩.
詩篇裡面15首的詩篇.
都有上行之詩這個title.
這15篇是專門用在.
以色列人從家裡.
往耶路撒冷聖殿去朝聖的路上.
這15篇都是講.
由他們家裡去到聖殿裡面朝聖的路.
他們藉著詩歌互相鼓勵.
他們一路走一路唱.
一路藉著詩歌.
來準備好自己走向神.
一路準備自己去走向神.
以色列人朝聖的目的地.
他知道就是在聖殿.
他知道目的地是朝向聖殿.
聖殿就代表著神的同在.
聖殿就代表著神同在.
那是他們嚮往的目標.
這個是他們向上的一個目標.
不管他們在山下過著什麼樣的生活.
也不管一路走來遇到多少的難處.
不論他們在山下過著什麼樣的生活.
也不管向上有多少的難處.
他們一路向著那個目標前進.
一路就向上這個目標前進.
弟兄姊妹們.
我們的一生是不是就像這樣.
我們一生是不是都像這樣.
我們一生是在走向神.
我們一生都是走向神.
我們一生走向神的這個過程.
很像以色列人他們行路走向聖殿.

$^{161}$我們一生走向神的這個道路.
就好像以色列人走向聖殿.
有時候我們會路上有青苔.
容易打滑.
有時候路上有青苔容易劃傷.
有的時候有尖銳的石頭.
有時候有尖銳的石頭.
有時候雜草叢生.
沒有路要開出一條路來.
有時候很多雜草沒有路要開路.
有時候在我們走的路上.
會碰到毒蛇蚊蟲野獸的攻擊.
有時候有毒蛇蟲有野獸的攻擊.
有時候走路的人會覺得累口渴了.
有時候人會覺得疲累口渴.
去年的七月份.
我們一家人有去拉森的國家公園.
上年七月份我們一家人去了國家公園.
那個公園它有一條Trail.
是一年只開放很短的時間.
剛好我們到的那一天它就開放了.
在那裡有一個Trail.
是一年才開始.
剛好我們去到就開.
那我們就很開心.
我們想挑戰一下自己.
我們很開心想挑戰一下自己.
可是我們是高估了自己.
我們高估了自己.
然後我們低估了那個路的狀況.
也低估了那條路的狀況.
走到一半的時候.
走了一半的時候.
我們家小朋友他不想走了.
我們家的小朋友不想再走.
他們一直說我的腳好痛.
哎呀我的腳好痛.
因為那條路開始走的時候還挺寬的.
因為那條路開始走的時候很寬闊的.
但是走著走著非常的窄.

$^{201}$但慢慢走的時候越來越窄.
它一邊是山.
一邊是山.
一邊它是一個懸崖.
一邊是懸崖.
然後那條路只能走一個人.
那條路只可以走一個人.
走的時候你還累了不能停下來.
因為沒有地方可以坐下.
一路走的時候累到不可以停下.
因為沒有路可以坐下.
然後我們路上遇到一個年輕人.
在路上遇到一個年輕人.
他走到一半在路邊哭.
走到一半的時候他在路上哭.
一個女孩子.
他說他不要走了.
他說不要走了不要走了.
我們當時也很為難.
因為孩子走不動了.
那時候也很為難.
因為我們的孩子也走不動了.
那時候就對面來了一群.
他們是從終點往回走的人.
只是那個時候已經有人回頭了.
他們就跟我們說加油.
他跟我們說加油.
Almost there.
我們一聽.
Almost there.
我們一聽的時候.
哇就到了.
我們就鼓起勇氣.
再往前走一走就到了.
鼓起勇氣立刻再走.
結果還走了一個多小時.
但是到了以後我們就說.
還好他騙了我們.
到了之後幸好他騙了我們.
不是被騙他是鼓勵了我們.

$^{241}$他們鼓勵我們.
不然我們可能錯過了.
那個很美的火山的風景.
如果不是我們就錯過了.
那個很美的火山的風景.
以色列人在登山走向聖殿的過程.
以色列人在登山走向聖殿的過程裡面.
他們也是有時候會覺得累.
或者覺得很艱難.
或者也覺得困惑什麼時候才到.
或者覺得困惑什麼時候到.
那這些詩歌就鼓勵和提醒他們.
這些詩歌就鼓勵提醒他們.
什麼情況下腳都要抬起來.
一直向前.
向上走.
向山走.
直到目的地.
我們的信仰是一個行走的信仰.
我們的信仰是行走的信仰.
從舊約當中.
神就要亞伯拉罕離開你的本地本族本家.
從舊約當中.
神就要亞伯拉罕離開你的本地本族本家.
那我們看到這些經文當中.
看到經文裡面.
在以色列這個民族一直壯大的過程中.
神不斷地提醒祂的子民.
你是寄居的.
你要起來走.
你要起來要走.
我們之前有一段時間教會的主題.
講到的主題.
大家對馬太福音這一段經文也會很熟.
在希伯來書也提醒我們.
我們是客旅.
是寄居的.
希伯來書也提醒我們.
我們是客旅.
是寄居的.

$^{281}$我們的家不在這個世界上.
而是在天上.
我們這一生的過程.
是走向神的過程.
我們這一生的過程.
都是走向神的過程.
我們的目的地是在神那裡.
我們的目的地是走向神那裡.
神要亞伯拉罕走的時候.
祂一步一步的提醒指示他.
一步一步的提醒指示他.
今天神也要我們知道.
我們的家鄉在天上.
我們不要被這路上看到的東西所影響和干擾.
我們不要被路上見到的東西去干擾我們.
我們這一生的目標.
是離開世界走向神.
我們這一生的目標.
都是離開世界走向神.
我們過去的一年.
不知道弟兄姊妹們過得好嗎?.
過去一年大家過得好嗎?.
可能你這一年.
經歷了非常有紀念意義的事.
可能今年你過得很有意義的事情.
也有可能你經歷了很難忘的傷痛.
又或者你經歷了一些很難忘的傷痛.
我知道我們當中有的家庭有新生的寶寶.
我們當中有新的嬰孩出世.
也有人換了工作升職加薪.
有人換了工作升職加薪.
也有人經歷了失去親人和朋友的傷痛.
有人經歷了失去家人朋友的傷痛.
有人經歷病痛.
有人失去工作.
但不管是怎樣.
當我們回顧自己的這一年.
我們問自己.
當我們回顧的時候.
我們問自己.

$^{321}$我們的腳步走在哪裡?.
我們的腳步是走在上行.
要去見神,向往神的路上嗎?.
我們的腳步是走向上行.
要朝見神.
是向神的路上嗎?.
還是我們在向下走.
走世界的路?.
還是我們只是向下走.
走向世界的路?.
我們這一年走過人生的一個一個小目標.
這些目標和將來永恆有沒有關係呢?.
這些目標和將來的永恆有沒有關係呢?.
在過去的這一年當中.
我們的心裡有沒有嚮往過.
我們人生最終的目標?.
在過去一年裡面.
我們有沒有嚮往過那個最終的人生目標?.
當我們新的一年開始的時候.
新的一年開始的時候.
我們會做很多新的計劃和目標.
我們有很多新的計劃,新的目標.
我們也來問一問自己.
在這新的一年當中.
我的屬靈目標是什麼?.
新一年裡面我們的屬靈目標是什麼呢?.
這裡有一些問題可以給大家一些思考.
有一些問題可以給大家一起思考一下.
這一年我每天要花多少時間安靜和神獨處?.
我要看完幾本屬靈的書籍?.
我要讀多少聖經?.
我怎麼去關心和愛身邊的弟兄姊妹?.
我在教會的服事有什麼需要改變的?.
在家庭裡,在工作上,我要怎麼榮耀神?.
這些問題值得我們去好好地思考一下.
這些問題值得我們好好地思考一下.
把它們放在我們新一年的目標和方向.
放在我們新一年的目標和方向.
在這首詩歌裡面,詩人說.
這句很熟,對不對?.

$^{361}$我要向山舉目.
這首詩歌裡面一句很熟的說話.
我要向山舉目.
為什麼要向山舉目?.
在現代中文譯本裡面.
它是說我要向群山觀望.
在現代中文譯本裡面.
它說我要向群山舉目.
舉目觀望群山.
大家覺得山給你什麼樣的感覺?.
大家覺得山有什麼感覺呢?.
很壯觀.
有人點頭.
有人跟我一樣覺得山很可怕嗎?.
有人覺得山很可怕嗎?.
其實站在那個很多的山.
黑色一片的時候.
是很可怕的場景.
覺得很可怕嗎?.
當以色列人他們在走的時候.
他看到黑黑的一片山.
大家可以想像一下他的感受是怎樣的.
其實對那時候的人來說.
山不是代表很美的風景.
山不一定是代表很美的風景.
山通常是可怕的.
不可預測的.
有盜賊出現.
有野獸出沒的.
那為什麼詩人要我們看這些山.
不是越看越可怕嗎?.
為什麼詩人要我們看這些山.
不是越看越可怕嗎?.
其實詩人是在提醒我們.
當我們看到不可預測的未來.
我們需要幫助的時候.
我們的幫助從哪裡來?.
我們的幫助從造天地的耶和華而來.
這兩句它其實是.
一句由一部分人唱.

$^{401}$另外一句.
第二句是其他一部分人的回應.
這兩句經文其實一句是.
有一邊人唱.
然後另外一個是回應這樣唱法.
在一唱一和當中.
他們的信心就被重新燃起.
看著這些黑黑的山.
提醒走路的人.
要看這些山背後的創造者.
要看到山背後的創造者.
當我們的人生.
經歷到艱難的時候.
經歷艱難的時候.
其實是神提醒我們.
看到這些艱難背後掌管一切的神.
要看到艱難背後掌管一切的那個神.
親愛的弟兄姊妹.
我們習慣用眼睛看什麼?.
看人.
看環境.
你把眼睛放在誰的身上?.
我們把眼睛放在哪裡身上?.
放在世上的一切.
還是放在掌管一切的神身上?.
人在面對順利的環境的時候.
面對成功的時候.
通常看到的是自己的能力.
通常看到的都是自己的能力.
自己的聰明.
在艱難的環境當中的時候.
在艱難環境裡面.
我們只看到這個難處非常大,很重.
我們只看到這個難處很大,很重.
難以忍受.
很難忍受.
我們的眼睛什麼時候可以看到這一切背後的神?.
我們的眼睛什麼時候可以看到一切背後的神?.
祂是救我們脫離死亡的神.
有沒有弟兄姊妹知道這是什麼運動項目?.

$^{441}$這是什麼運動?弟兄姊妹知不知道?.
賽艇.
這是1715年英國開始這個賽艇的運動項目.
1715年英國開始了這個賽艇的運動.
之後就變成奧運會的大項目.
大是因為它人數很多參賽的.
大是因為人數多.
這個叫做八人賽艇.
但是大家仔細看圖.
它有幾個人參加比賽?.
大家看看有多少個人.
9個人.
很有趣.
這是唯一一個比賽項目.
那個比賽的選手.
他是看不到自己的目的地的.
他看不到終點.
他怎麼樣大家可以行動一致.
不會偏離方向.
他們怎樣可以行動一致不偏離方向呢?.
因為在這個賽艇上有一個指揮.
因為在賽艇上有一個指揮.
他是唯一一個可以看到目標的.
有他在.
他們就不會偏離方向.
有時候我們看我們的人生.
有時候我們看人生.
我們很難看得很清楚我們的前面.
很難看到很清楚前面的路.
有時候我們也覺得看不到希望.
我們看不到前面的希望.
但有一件事我們不要忘.
但有一件事我們不要忘記.
在我們的船上有一個指揮.
我們不需要真的看清楚前面的目標.
我們也不需要看路上的各種狀況.
我們也不一定要看路上各種的情況.
只要有這位指揮在.
我們只管看著他.
只要看著這個指揮.

$^{481}$我們的神就是我們人生旅途中的指揮.
我們的神就是人生旅途中的指揮.
只要祂在我們的人生當中.
只要有祂在我們人生當中.
祂與我們同在.
我們就能夠到達終點.
得到獎賞.
我們就能夠到達終點.
得到獎賞.
希伯來書的十二章二節.
希伯來書十二章二節.
仰望為我們信心創始成中的耶穌.
仰望這個詞在希臘文當中.
意思是說.
意思是.
將注意力從一切的人,事,物,環境,地方,時間.
所有事情轉移到耶穌身上.
牢牢地看著祂.
仰望耶穌.
不要看環境的艱難.
如果我們停下來去想我們的難處有多大.
我們只會把自己困住.
因為我們沒有辦法去面對.
我們沒有辦法想清楚神為什麼讓這些難處臨到.
因為我們沒有辦法可以想清楚為什麼這些難處要臨到.
我們也沒辦法靠自己去解決.
不要看難處.
只要看著耶穌.
仰望耶穌也就意味著.
不要去看順境當中的各種好處.
不要看這世界的好.
我們所擁有的一切.
免得我們的心被迷住了.
以至於我們看不見賜福給我們的神.
把我們的眼睛從各種東西上轉移開.
求主幫助我們.
求神幫助我們.
只看祂.
在第三節到第八節當中.
有一個詞是重複重複出現的.

$^{521}$在1,2,1,3到8節裡面.
有一個詞是重複又再重複的.
大家有注意到是什麼嗎?.
大家有沒有看到是什麼呢?.
對.
保護你.
是保護.
是 保護.
朝聖的人他在旅途當中會遇到挑戰.
朝聖的人在路途裡面會遇到挑戰.
他會遇到生命的危險.
詩歌在提醒他們.
詩歌提醒他們.
我們的神時刻在保護.
我們的神時刻都是保護.
今天也提醒我們.
不管我們的人生遇到什麼.
祂是時刻保護我們.
祂是時刻保護我們的神.
祂救我們脫離一切的災害.
是所有一切.
也就是說神對我們的保護.
是全時間.
全身心靈的.
全身心靈.
在這段經文當中.
在這段經文裡面.
祂提到日常的保護.
真正的保護.
祂提到日常真正的保護.
還提到一個非常重要的.
還有提到一個.
還有提到一個很重要.
跟我們的生命相關的.
這可能是我們所有人最大的擔心.
這可能是我們很多人很擔心的一樣東西.
但是生命是有限的.
生命是有限的.
就算人想盡一切辦法.
沒有一個人可以改變自己生命的壽數.

$^{561}$就好像人想盡一切辦法.
沒有辦法可以延長生命的壽數.
人生當中的災病.
人生裡面的災害.
意外.
是我們人力不可控的.
是我們人類不可以控制.
但是經文提醒我們.
神的保護可以免去一切的災害.
我們的性命不會被禍患所奪去.
我們的性命不會被這些禍患奪去.
它不是在說我們不用擔心沒有人會生病.
沒有意外會發生.
意思不是說我們不需要擔心.
我們沒有人病 沒有意外發生.
它意思是說.
我們不要為死亡擔心和害怕.
意思是說我們不需要為死亡而擔心擔憂.
有時候怕死比真死還可怕.
不管是什麼人.
今天你是基督徒.
也沒有辦法.
要擺脫你生命的結束.
擺脫你生命的結束.
誰都有那一天.
秦始皇開始找長生不老要.
他找到了嗎?.
秦始皇找長生不老藥找不找到?.
我不知道大家有沒有注意到這個新聞.
有沒有看過這個新聞?.
Brian Johnson他在2023年的時候.
進行了一個非常激進的實驗.
這個人在2023年有一個實驗.
因為他很有錢.
他有一個專門的醫療團隊.
來研究他怎麼樣可以讓自己永保青春.
來研究他怎麼樣可以保持青春.
那他就想出一個辦法.
他就想出一個辦法.
他和他17歲的兒子換血.

$^{601}$再把他的血換給他70歲的爸爸.
將17歲的兒子的血換給這個70歲的爸爸.
有沒有成功?.
失敗了.
不到半年就宣佈他的實驗失敗.
實驗當然失敗了半年之後.
聽說他現在又有新的研究方向.
聽說他有新的研究方法.
發現人很怕老.
人很怕老.
怕生命結束.
但因為罪.
誰都沒辦法改變這個結局.
沒有人可以有辦法解決這個結局.
但對我們相信神的人而言.
我們的生命不是在今生.
我們今生不是為了吃喝玩樂.
享受安樂度過.
我們今生不是為了吃喝玩樂.
安逸地度過.
我們今生是為了走到永恆裡去.
我們今生是為了走到永恆裡.
今生只不過是一個過程.
是我們一步一步走向神的過程.
如果一個基督徒.
在世上的日子結束了.
在世上的日子結束.
不管是什麼原因.
你的世上的日子結束.
不管是什麼原因.
在世上的日子結束.
只有一個解釋.
是神許可的.
你要回到祂那裡去了.
要回到神的那裡.
不是任何災害導致的.
因為死亡與我們而言不是終點.
而是歸途.
我們的性命永遠在神的手中.
我們需要有信心.

$^{641}$求主幫助我們明白這個真理.
求神幫助我們明白這個真理.
因為這是我們可以安然度過人生的最大安慰.
不知道大家有沒有弟兄姊妹看過這本書.
叫《深入非洲三萬里》.
有沒有看過這本書叫《深入非洲三萬里》.
它是記載著宣教士.
它記載著宣教士.
李文斯頓的傳記.
他是非常有名的宣教士.
他是很有名的宣教士.
他是廢除奴隸販賣的關鍵人物.
他是第一個進入非洲的醫生.
第一個進入非洲的醫生.
他把福音帶進非洲.
在書裡面講了他三次進入非洲的探險旅程.
在書裡面講到他三次進入非洲的探險旅程.
其中有一段記載到.
他和非洲當時非常有名的大酋長.
其中有一段記載到他進入非洲的一個大草場.
他們有一段對話.
當時那位大酋長他得了肺病.
他就問李文斯頓.
人死後有什麼盼望呢?.
李文斯頓說.
把心放在上帝的手中就有盼望.
於是他回答.
把心放在上帝的手中就有盼望.
李文斯頓把自己的一生放在上帝的手中.
我們把我們的心放在哪裡?.
我們的心放在哪裡?.
我們的心放在很多人的身上.
很多的事情上.
很多的事上.
還是我們把它放在上帝的手中.
這一篇詩篇提醒我們.
我們的神.
祂的大能遠超過我們可以想的.
祂的大能遠超過我們可以想像的.
也超過我們可以理解的.

$^{681}$祂能聽見我們聽不見的.
祂能看見我們看不見的.
祂可以注意到我們所注意不到的.
祂可以注意到我們注意不到的事情.
不管什麼時候.
祂對我們的保護從來都不懈怠.
神對我們的保護從來都不懈怠.
有時候我們什麼都沒有做.
神已經做了.
有時候我們知道是神做的.
但我們沒有看到.
但是祂默默地在做.
而是神默默地在做.
在我們一生當中.
神從來沒有停止保護我們.
祂知道我們面對的難處.
祂在默默地做.
上帝知道我們的痛苦.
神知道我們的痛苦.
祂一分一秒都不離開.
我們可以來讀這一節經文嗎?.
耶利米蘇29章11節.
我們一起來讀耶利米蘇29章11節.
耶和華說:我知道我向你們所懷的意念是持平安的意念.
不是降災禍的意念.
要叫你們抹後有子夢.
有這樣一位神.
我們的心還需要其他的歸屬嗎?.
小時候我最期待兩件事.
我小時候最期待兩件事.
一個是過生日.
一個是過年.
我從來沒想到.
我從來沒想過.
現在是我最怕的兩件事.
現在是最怕的兩件事.
因為小時候一直盼著長大.
我甚至跟神禱告說.
能不能讓我快點長大.
小時候很想長大.

$^{721}$甚至我祈禱跟神說.
可不可以讓我快點長大.
現在收回來不及了.
現在收不回了.
活著活著就會發現.
一年過得比一年還快.
活著的時候我們發現一年比一年好像過得快.
又是過去一年迎來新的一年.
又過去一年又迎來新的一年.
對世界來說我們是慢慢的老去.
對世界來說我們是慢慢老去.
但是對我們的目的地而言.
我們是離得更近了一步.
我們更進更近了一步.
做基督徒很幸福.
所以做基督徒很幸福.
因為我們的生命不是慢慢的老舊.
因為我們的生命不是慢慢的老去.
我們是一天新式一天.
我們一天新式一天.
我們的生命還有不斷更新.
重新開始的機會.
我們的生命還可以不斷更新.
重新開始的機會.
親愛的弟兄姊妹.
當我們走著新的一年的時候.
當我們走入新的一年的時候.
我們問一下自己.
我們的腳步要走向哪裡.
我們的眼睛要看到哪裡.
我們的心要屬於誰.
這個詩篇再一次提醒我們.
新的一年.
我們抬起自己的腳.
繼續靠主前行.
我們一同來禱告.
親愛的天父 我們謝謝你.
謝謝你讓我們新的一年開始.
我們在這個新的裡面.
我們靠著你.

$^{761}$我們可以不斷地更新我們的生命.
我們靠著你.
我們可以越活越新.
求你幫助我們在這新的一年當中.
我們立定心智.
我們走在你的裡面.
我們的眼睛單單只看向你.
我們的心只屬於你.
求你在我們這新的一年當中.
你做新的事.
你在我們的生命當中.
讓我們一同與你唱新的歌.
讓我們能活出那新的意象.
讓我們能夠繼續見證榮耀你的名.
繼續靠著你向前走.
謝謝主 榮耀歸給你.
禱告奉耶穌救命球.
阿們.
\newpage



\section{馬太福音 5:7-12}
\label{sec:031717}
\textbf{有福了!屬神的人 (2)}~~~~陳雅莉傳道
\newline
\newline
連結: \href{http://www.church.com.hk/acms/content.asp?site=cdc\&op=show\&type=product\&code=031717\&layout=sermon}{\texttt{031717}} ~~~~ 語音日期: 2026-03-01
\newline
\newline
\hyperref[sec:031716]{\small{< < < PREV SERMON < < <}}
~
\hyperref[sec:index_chronic]{\small{[返順時目]}}
~
\hyperref[sec:031718]{\small{> > > NEXT SERMON > > >}}
\newline
\newline
馬太福音 5:7-12
\newline
\begin{longtable}{cl}
\hline
\hline
章節 & 經文 (和合本修訂版)\\
\hline
5:7 & \begin{tabularx}{0.7\textwidth}{X} 憐憫人的人有福了！因為他們必蒙憐憫。 \end{tabularx} \\ \\ \relax
5:8 & \begin{tabularx}{0.7\textwidth}{X} 清心的人有福了！因為他們必得見神。 \end{tabularx} \\ \\ \relax
5:9 & \begin{tabularx}{0.7\textwidth}{X} 締造和平的人有福了！因為他們必稱為神的兒子。 \end{tabularx} \\ \\ \relax
5:10 & \begin{tabularx}{0.7\textwidth}{X} 為義受迫害的人有福了！因為天國是他們的。 \end{tabularx} \\ \\ \relax
5:11 & \begin{tabularx}{0.7\textwidth}{X} 「人若因我辱罵你們，迫害你們，捏造各樣壞話毀謗你們，你們就有福了！ \end{tabularx} \\ \\ \relax
5:12 & \begin{tabularx}{0.7\textwidth}{X} 要歡喜快樂，因為你們在天上的賞賜是很多的。在你們以前的先知，人也是這樣迫害他們。」 \end{tabularx} \\ \\
[1ex]
\hline
\hline
\end{longtable}
$^{1}$請姐妹願你們平安.
今天在早上回來的時候.
我覺得分外珍惜.
可以回來教會.
敬拜上帝.
能夠一起讚美神.
實在是恩典.
不知道過去的一個星期.
你所看到的.
所感受到的.
所聽到的.
一切身邊在你身邊發生的事.
你有什麼感受呢?.
或者會讓我們有傷心難過.
或許會讓我們有憤怒.
有疑惑.
或者會讓我們有很擔心.
很憂慮的時候.
或者也有.
讓我們仍然覺得感恩和平安的時候.
無論這一刻.
你的感受是什麼.
我盼望你仍然相信.
主與我們同在.
我們都將這一切的感受.
我們先來到上帝的面前.
去交託我們的主,好嗎?.
我們一起禱告.
親愛的主耶穌.
在這個破碎的世代.
在這個破碎的世界裡面.
你用你自己的破碎.
親自拯救這個世界.
你將你自己抹開.
將你自己的破碎.
成為世人的祝福.
主耶穌今日.
我們知道你仍然是以馬來你的神.
你仍然是與我們同在的神.
無論在我們的傷心難過裡面.

$^{41}$還是疑惑憤怒裡面.
還是擔心憂慮裡面.
你都與我們同在.
多謝主.
此刻我們仍然有感恩的力量.
此刻在患難中.
我們仍然能夠感受到從你而來的平安.
願主安慰我們的心.
堅定我們信靠你的心.
要我們不看環境.
單單定性仰望你.
並且深信.
你仍然是要將福樂賜給我們的神.
聽祈禱.
奉耶穌基督明智祈求.
阿們.
這個星期我相信.
大家都一定會更多去思考有關人生的問題.
或者可能你都會問很多有關苦難的問題.
但我們知道.
如果我們要思想人生.
我們要思想苦難.
我們不在主裡面.
不在主裡面去思想的話.
一切都是沒有意義的.
因為我們人生.
是我們主所賜給我們的.
唯有主.
有一切的答案.
願我們今天.
仍然渴慕聆聽主的道.
願意聽.
並且願意行.
行出主所教導.
生命願意經歷更新.
成為真正有福的人.
好嗎?.
還記得上星期我們一起來看八福.
在登山裡面的八福.
還記得頭四福.

$^{81}$講什麼嗎?.
我們來重溫一下.
耶穌說:心靈貧窮的人有福了.
因為天國是他們的.
哀動的人有福了.
因為會必得安慰.
天和的人有福了.
因為必定會承受土地.
幾何無意的人有福了.
因為他們會必得飽足.
我們今天再來看餘下的四福.
耶穌再繼續說:憐憫人的人有福了.
因為他們會必蒙憐憫.
清心的人有福了.
因為必定會見神.
締造和平的人有福了.
他們將會成為神的兒子.
為義受迫害的人有福了.
因為天國是他們的.
憐憫人的人有福了.
因為必定會蒙憐憫.
憐憫這個詞語我相信不難理解.
不需要特別去解釋.
你都明白什麼叫憐憫.
但在聖經裡面.
這個不單只是一個行動這麼簡單.
不單只是一個屬靈的特徵.
不單只是物質上具體的支援這麼簡單.
憐憫更加是與人共情.
與人共感.
我們的耶穌基督是賜給憐憫的主.
是主先憐憫我們.
以致我們可以憐憫別人.
一個人真正有憐憫心腸.
上帝都會應許他.
你都會必定蒙上帝的憐憫.
憐憫不單只是在旁同情一下.
憐憫不單只是見到對方.
他真的很可憐.
憐憫是進入對方的生命.

$^{121}$真正的同行.
真正的關懷.
真正感受對方的感受.
這個是憐憫.
在過去的幾天.
我親眼目睹一個又一個憐憫人的人.
我親眼目睹一個又一個.
願意展現這個屬天的價值的弟兄姊妹.
如果我們在群組裡面.
團體群組裡面.
你一定收到消息.
我們的弟兄姊妹.
遇上了嚴重的交通意外.
這四位.
三位的姊妹和其中一位姊妹的丈夫.
身著嚴重的交通意外.
去到不同的醫院的時候.
我見到我們的弟兄姊妹.
二話不說.
立即開車衝去牛津醫院.
我又見到一個又一個的姊妹來到醫院裡面.
當我們的姊妹失去至親.
痛苦煎熬的時候.
一個又一個的弟兄姊妹不住在醫院陪伴,關心.
全教會的弟兄姊妹不分晝夜.
在群組裡面發放代禱的訊息.
全教會一起同心禱告.
我見到大家自發成立關顧小組,關懷小組.
去幫助我們的弟兄姊妹.
有人負責去為姊妹送湯送飯.
有人負責每日打電話關心姊妹.
有人負責車出車入.
每一個都負上自己可以做的那一份.
我看到這一切.
除了很感動.
我更加看見他們走出了憐憫.
當我跟姊妹說.
「很多人在身邊,我們一起走」.
我告訴姊妹.
「是我很想多謝你」.

$^{161}$「是我很想多謝姊妹」.
「讓我們可以成為憐憫人的人」.
「我很想多謝我們的姊妹」.
「是她給我們機會與哀傷的人同行」.
「這是憐憫」.
當處理好一切要離開醫院的時候.
姊妹跟我說.
「我浪費了你們很多時間」.
我看著姊妹說.
「你不要浪費我們時間」.
「你是花了我們很多時間」.
不過我告訴姊妹.
「時間是用來花的」.
「上帝給我們的時間是用來花的」.
「是花在那些有意義和價值的事上」.
姊妹不是浪費了我們很多時間.
她是給我們機會把時間花在很有意義和價值的事上.
是有屬天價值的事上.
靚姊妹 你知道嗎?.
當你願意停下你手上的事.
為有需要的人禱告.
為有需要的人獻上一句關心問候.
當你願意付出時間.
去提供實際的需要的時候.
這些通通都是異常何等有價值的事.
憐憫人的人有福了.
你們都會別蒙憐憫.
我們的教會要成為一個憐憫人的教會.
我們的教會都要別蒙主的憐憫.
我們雖然軟弱.
我們雖然只是盡上我們所謂的綿力.
但我們的軟弱放在主的手裡.
成為別人的安慰.
讓我們還有能力可以憐憫別人.
之後耶穌說「清心的人有福了」.
因為他們必得見神.
清心是純潔無瑕.
是無罪惡,清潔,純潔,單純,專注於主.
不被世俗的纏擾.
心無旁騖 只是看見主.

$^{201}$如果要問一個清心的人是怎樣的呢?.
你去看我們的小孩子你就知道了.
我們的小孩子很單純.
每一個小孩子都是很簡單.
所以耶穌說我們要回轉像小孩子.
清心的人可以直接見到上帝.
與上帝溝通.
過去幾天.
當我去醫院探望其中一位受傷的姐妹時.
我問她「在交通意外時你的心情是怎樣的?」.
她回答我「很奇妙的傳遞」.
「無論是意外發生的時候」.
「還是我整個人要被送入醫院的時候」.
「還是我知道原來我都幾嚴重要做手術的時候」.
「還是手術之後我知道要慢慢復康的時候」.
「我沒有一刻覺得憂慮」.
姐妹說「整個過程我都覺得很平安」.
她說她一點憂慮都沒有.
而且我一直跟她聊天的時候.
我覺得她不是在說你真的感受到平安在她身上.
姐妹不是因為沒有遇到意外.
所以她覺得平安.
不是.
不是她知道不用做手術所以要平安.
不是.
而是無論遇上什麼樣的環境.
姐妹單單仰望神.
她將自己完全交託給神.
她很清心.
她沒有看眼前的環境.
她沒有看她將要面對的事.
她只是將自己交託給全能的上帝.
當我望著姐妹的時候.
我真的覺得姐妹有福了.
她是清心的人.
她對上帝這份清心,這份單純.
感動了在她身邊未信主的丈夫.
她的丈夫很擔心她.
很擔心她.
但當她看著姐妹的時候.

$^{241}$很單純仰望上帝.
當我去探訪她的時候.
我邀請她的丈夫.
有人跟我說跟她的丈夫傾福音是很困難的.
我去探訪之前有人跟我說.
你不能跟她的丈夫說耶穌的.
但當我去到醫院探訪她的時候.
我邀請她的丈夫一起為姐妹禱告.
他願意握著我們的手一起禱告.
當我跟這位丈夫說.
把握時間信耶穌了 好嗎?.
不要再等了.
她回答我:我會好好考慮.
我們也放在禱告裡面,好嗎?.
繼續紀念姐妹的丈夫.
是的.
因著姐妹的清心,很單純的仰望上帝.
她身邊的伴侶都看得見.
即使在意外中,大手術前.
姐妹真是清心就有福了.
得見上帝.
今天我們仍然要以信心的眼睛去仰望將來.
但我們知道到終末.
我們真的會必見上帝.
到終末的時候.
我們一定會看見神奪目的榮光.
清心的人有福了.
耶穌說:締造和平的人有福了.
因為他們會稱為「神的兒子」.
締造和平.
可能我們會想.
當有一個很不和睦的情景.
就需要去締造和平.
不過我相信聖經在這裡的呼籲.
不是當有事發生的時候.
而是每一天.
你有沒有想像過.
你是每一天在任何的情景都可以成為一個和平之子.
是活出一個和平之子的生命.
今時今日.

$^{281}$這是一個破碎的世界.
是一個關係撕裂的世界.
是一個人際間不再信任的世界.
是一個戰爭與衝突隨時發生的世界.
如果要使人和睦.
是要刻意的.
如果要締造和平.
你是要先攻克.
先要克性各種的仇恨.
克性你裡面的憤怒.
克性你的論斷.
克性你的拒絕.
你才能刻意的締造和平.
締造和平是要用心去做的.
是要一點堅持.
一點忍耐.
我們要由流心走.
耶穌這裡不是說喜歡和平的人有福了.
祂是說締造和平的人有福了.
是願意在仇恨當中散播和平.
願意將和平的種子散落在撕裂裡面.
最近在社交媒體裡面.
發生一件鼓勵仇恨的事.
我收到兩個女兒的學校發信給家長.
說有人在網絡上鼓勵中學生.
和不同的中學分成不同的隊伍.
在放學之後持刀去攻擊另外一間學校的學生.
我簡直覺得這是一個瘋狂的世代.
當這個世界不單止不是鼓勵和平.
是鼓勵仇恨.
鼓勵互相攻擊的時候.
我們要怎樣去面對這個世代.
弟兄姊妹.
這個世代的敵人已經夠多了.
你不製造敵人 敵人都會找上門來.
我們不需要再為自己製造更多的敵人.
當我們看到有不和睦的時候.
當我們看到有不和平的地方.
求主幫助我們不要再冷眼旁觀.
不要再在一旁覺得不關你的事.

$^{321}$當有人告訴你.
當你知道 當你發現有不和平的事的時候.
到底你是選擇多加一把嘴 多說兩句.
還是我們嘗試成為和平之子呢?.
我們願意嘗試讓對方知道.
我們看的東西都很有限.
我們願意承認我們不是全知.
我們願意在當中成為締造和平.
而不是加深仇恨的人.
和睦.
和睦的關係.
在這個世代已經變得離足珍貴.
而且是非常急需.
上幾個星期有一個弟兄跟我分享.
他說他在一個群組裡面.
看到一些差點要造成互相指責的言論.
當一刻 他嘗試去理解.
嘗試去理解對方的觀點.
知道對方可能都是一些情緒的發洩.
不是有心要針對或者傷害.
雖然不好受.
雖然看到這些言論不好受.
甚至覺得不合理.
但弟兄沒有選擇.
多加一把嘴 多反駁兩句.
而是他告訴我 他馬上選擇.
去私訊那些他覺得看到這句說話會不舒服的人.
立即安撫他們.
叫他們冷靜.
請他們停一停.
告訴對方 可能他們都是有些情緒發洩.
大家安靜點 冷靜點.
弟兄免去了無謂的爭吵.
有時候我們會說.
我們人人都不一樣.
人與人的相處總是很容易造出很多磨擦.
我們都彷彿會彼此傷害.
是的.
我們的性格.
我們的背景.

$^{361}$會造成我們彼此傷害.
但不要忘記.
主耶穌到今天.
仍然呼召我們彼此相愛.
當我們彼此相愛的時候.
眾人因此就認出.
我們是主的門徒.
今天我們說我們要活出.
基督徒的樣式.
活出與得救相稱的生命.
祂主幫助我們.
實在我們有軟弱 有做不到的時候.
但是.
主耶穌呼召我們彼此相愛的使命.
仍然是有效的.
仍然是可能的.
因為我們的主.
祂在《以賽亞書》說祂是和平之君.
祂是和平之君.
如果我們是和平之子.
我們就必稱為神的兒子.
因為我們的主是和平之君.
讓我們也來學習我們的主.
成為和平之子.
弟兄姊妹 這是一個選擇.
選擇要去做替做和平的人.
還是冷眼旁觀的人.
還是毫不關心的人.
甚至加深仇恨的人.
都是我們的選擇.
讓我們選擇替做和平.
因為替做和平的人有福了.
你會稱為神的兒子.
耶穌接著說.
為而受迫害的人有福了.
因為天國是他們的.
我相信這裡不是指簡單頂撞一下.
被人打壓一下這種.
而是在說為著你的信仰.
受嚴重的迫害.

$^{401}$當你堅持你的信仰.
堅持行屬天的價值的時候.
將要付上的代價.
一個拒絕隨波逐流.
不如這個世界妥協.
我們必然是必然會受這個世界的唾棄.
我們會被世界嘲弄.
許多時候.
我們都沒有要求別人和我們一起要行這個義.
但是只是單單你行義.
但是你的義顯出別人的惡.
你已經成為被攻擊的對象.
因為主耶穌的教導是何等顛覆這個世界.
因為當我們要行主耶穌的教導的時候.
我們就成為了這個不合理的世界.
他們裡面的規矩我們無法跟隨的時候.
我們就成為了他們的威脅.
我們的善顯出了別人的惡.
不過耶穌仍然說.
這樣的人是有福了.
耶穌在這裡突然之間有感.
由他們他們他們.
突然間變成了你們和我.
耶穌說:人若因我辱罵你們.
迫害你們.
捏造各樣的壞話.
誹謗你們.
你們就有福了.
要歡喜快樂.
因為你們在天上的賞賜是很多的.
在你們之前的先知.
人也是這樣迫害他們.
耶穌是直接跟你們說的.
不再是他們.
是你們.
耶穌彷彿看見.
正在聽這個訊息的人.
是昔日.
甚至是今日.
正在聽這個訊息的你和我.

$^{441}$當我們如果真的要立志成為一個.
為二受迫迫的人的時候.
我們面對的事情.
將會不是好受的.
但是耶穌重申.
即使有禍患擺在面前.
我們仍然是有福的人.
因為天上的賞賜是很多的.
主幫助我們.
不是只看這個地上所擁有的.
不是只看眼前所經歷的.
乃是看到天上將來的.
耶穌在第十章這樣說.
祂說:我猜你們如同羊進入狼群.
你們要靈巧掌蛇.
順良掌甲子.
我們是羊要進入狼群.
不過耶穌說.
你們要為我的名被眾人恨惡.
不過唯有忍耐到底的.
必然得救.
是的.
雖然要為二受迫迫.
不與這個世界的價值觀妥協.
不隨波逐流.
跟隨耶穌的教導去走.
但是忍耐到底的.
必然得救.
歷世歷代我們看到很多基督徒.
在不同的場景.
不同的國家.
不同的地域被迫害.
全球到這一刻.
仍然有超過三點一億人.
三點一億的基督徒.
他們是我們的弟兄姊妹.
有三點一億的基督徒.
仍然落在被極度迫迫的水平裡面.
即是平均全球每七個基督徒.
就有一個是今天活在迫害中.

$^{481}$他們會被殺.
他們會被迫迫.
他們會被捉拿.
他們是我們的弟兄姊妹.
耶穌不是不知道這些掙扎.
耶穌也不是不知道.
我們活在地上要面對的掙扎.
耶穌不會不知道.
雖然是這樣.
有迫不還難.
但是耶穌應許我們.
我們是有福的.
你要看到天上的.
並且要歡喜快樂.
不是為了這一刻.
是為了將來.
為了他日的盼望.
擁有這八種性格的人.
都是有福的.
如果我們仔細看這八福.
你會發現.
耶穌從第一個福.
去到最後的福.
頭和尾.
都有一個很重要的應許.
天國是他們的.
而頭四個福和尾四個福.
最後都是跟義有關的課題.
耶穌強調的是.
當我們今天按照這個世界的.
道德標準去行事為人的時候.
我們很容易會自以為自己已經做得很好.
但耶穌卻是要我們追求更超越的.
更高的.
耶穌要將焦點.
從表面的行為轉移到內心.
耶穌就應許天國是他們的.
如果我上星期也有提過.
如果你仔細回去看整個第五章.
你會發現耶穌的「義」.

$^{521}$祂所謂的「義」.
是整整一章都不停地重複.
耶穌說.
有人說不可以殺人.
這是道德的標準.
但耶穌說不單止不可以殺人.
你是要和好.
有人說不犯姦淫就可以了.
耶穌說不是.
是要不准動淫念.
有人說只要不背棄誓言就可以了.
耶穌說不是.
是要事就說事.
不事就說不事.
有人說.
走一里路就夠了.
耶穌說不是.
是要多走一里路.
有人說.
我們今天愛鄰舍已經夠了.
耶穌說不是.
是要愛仇敵.
這是耶穌的「義」.
你要饑渴慕這個「義」.
你要為這個「義」受逼迫.
是很顛覆的.
甚至我們說是很難理解.
但是不是不可能走的呢?.
不是.
生死苦難.
本來就是我們一生都要探討的話題.
當面對世界.
隨時每一分每一秒.
可能下一秒.
戰爭又再次來到的時候.
當苦難和失去隨時會在我們身邊發生的時候.
我們發現我們根本無法避免.
但最重要得到的福氣.
不是單單存活.
不是單單留著命.

$^{561}$很幸運沒事.
不是.
那是你要真正活出一個.
屬天價值的生命.
一個真正有福的生命.
耶穌給屬祂的人的福.
你再看一次.
在何等的一帶.
耶穌說「天國是他們的」.
他們會必得安慰.
必承受土地.
必得飽足.
必蒙憐憫.
必得見神.
必稱為神的兒子.
天國是我們的.
何等有福呢?.
這是屬神的人.
所得會享的福氣.
過去的這個星期.
確實讓人心痛又難過.
很多的問題.
很多的反思.
但如其來的噩耗.
讓我在上帝面前去求問.
當我屈膝在上帝面前.
去問上帝.
你要我們的教會學什麼功課.
上帝告訴我.
我們將要成為一個更合神.
更合神心意的教會.
面對苦難.
我們難免會問.
為什麼會有苦難.
不過我今天不打算解釋這個話題.
當我們的教會被迫一次又一次地.
去面對我們的弟兄姊妹.
落在重大的創傷.
和無法避免的意外的時候.
我們要問上帝什麼?.

$^{601}$我們到底有什麼領袖?.
約伯在約伯記這樣寫下.
「我知道我的救贖主活著.
活後祂必站在塵土上.
我這個皮肉滅絕之後.
我必在肉體之外得見神.
我自己要親眼見祂.
並不像陌生人」.
約伯經歷無法解釋的苦難.
他的痛 他的痛楚 他的苦楚 他的心情.
是無法能夠理解的.
但約伯仍然深信.
他說「我的救贖主活著」.
這是他的宣告.
在過去的一個星期.
當我不斷奔走不同的醫院.
每一次當我出車離開醫院.
回家後又再出車離開醫院的路程上.
我心中不住響起這首詩歌.
「我的救贖主活著」.
這首詩歌的歌詞是這樣說的.
它說「我知道我的救贖主活著.
祂是永活的主」.
當我在心谷迷失時.
祂領我走正義路.
我知道我的救贖主活著.
祂是永活的主.
當我在曠野孤獨時.
祂伴我鎖我的燈.
我知道我的救贖者永遠活著.
我心不再憂慮.
我要在每一個日夜之中.
去領受祂豐盛之愛.
我知道我的救贖者永遠活著.
我的靈不再沉睡.
當好覺響起的那一天.
我將要見祂榮光之面.
在苦難的當前.
我選擇了不再問上帝.
為何會有苦難.

$^{641}$而是要問上帝.
你想我們這個教會.
成為一間怎樣的教會?.
我們是一間只是學玩樂.
搞活動.
每個星期日大家回來聚會的教會嗎?.
還是我們真是一個能夠有福同享,有禍能夠同擔當的教會?.
當我們一再問上帝.
求主讓我們合而為一.
求主讓我們彼此相愛.
求主讓我們一同經歷上帝的時候.
我們一同與我們的弟兄姊妹度過哀傷.
就成為了連結全教會最重要的集體回憶.
就如約伯一樣.
當我們的教會經歷一切苦難之後.
約伯說:我從前風聞有理.
現在我親眼看見你.
是的.
弟兄姊妹我們的教會也要一樣.
當我們的教會一同經歷經過苦難之後.
我們要成為一間從前風聞有神的教會.
到現在親眼看見主的教會.
我們同心禱告.
親愛的主耶穌,我深信這一刻你知道我們的沉重.
你知道我們的哀動.
你知道我們的難過.
當我們看著眼前的事.
確實會讓我們無法釋懷.
但是主求主幫助我們.
讓我們專心看見你.
讓我們專心仰望你.
單單看見主.
主仍然與我們同在.
即使在混亂的世局裡.
動盪的世界裡.
破碎的世界裡.
主仍然掌權.
求主讓我們真心正正一同經歷你.
讓全教會能夠一同經歷你.
讓我們同心在你的面前去呼求.

$^{681}$主啊.
我們要在你裡合而為一.
我們要成為一間更合神心意的教會.
我們要成為一間親眼看見主的教會.
求主幫助我們.
我們聽祈禱 奉鼠明給.
\newpage



\section{歌羅西書 1:21-23}
\label{sec:031718}
\textbf{主耶穌已經開始了,我們又怎樣呢?}~~~~黃紹權牧師
\newline
\newline
連結: \href{http://www.church.com.hk/acms/content.asp?site=cdc\&op=show\&type=product\&code=031718\&layout=sermon}{\texttt{031718}} ~~~~ 語音日期: 2026-03-01
\newline
\newline
\hyperref[sec:031717]{\small{< < < PREV SERMON < < <}}
~
\hyperref[sec:index_chronic]{\small{[返順時目]}}
~
\hyperref[sec:031719]{\small{> > > NEXT SERMON > > >}}
\newline
\newline
歌羅西書 1:21-23
\newline
\begin{longtable}{cl}
\hline
\hline
章節 & 經文 (和合本修訂版)\\
\hline
1:21 & \begin{tabularx}{0.7\textwidth}{X} 從前你們與神隔絕，心思上與他為敵，行為邪惡； \end{tabularx} \\ \\ \relax
1:22 & \begin{tabularx}{0.7\textwidth}{X} 但如今，他藉著他兒子肉身的死，已經使你們與他自己和好了，把你們獻在他的面前，成為聖潔，沒有瑕疵，無可指責。 \end{tabularx} \\ \\ \relax
1:23 & \begin{tabularx}{0.7\textwidth}{X} 只要你們持守信仰，根基穩固，堅定不移，不致動搖，離開了你們從前所聽見的福音的盼望；這福音也是傳給天下一切被造之物的，我—保羅作了這福音的僕役。 \end{tabularx} \\ \\
[1ex]
\hline
\hline
\end{longtable}
$^{1}$在聆聽主的說話以先.
讓我們先同心低頭禱告.
仰望愛主.
父神多謝你.
讓我們能夠再次回到你的說話當中去學習.
因此我們知道.
在你的啟示裡面有很多事情.
和我們自己本身的性情.
以及甚至在教會多年所傳遞的神學.
有很大的衝突.
今天的教會.
活在一個以人為中心的一個神學基礎.
很多時候變得.
只是歪曲了上帝你的啟示.
以討人喜悅為中心的傳道.
但願主你給我們每一個.
屬你的子民都有一份智慧.
從天上而來.
藉著聖靈所帶領的引導.
使我們曉得怎樣.
真是將所聽到的錯謬.
能夠將它逐一排除.
讓我們在這些種種的神學毒瘤裡面.
中毒已深良久之後.
能夠得以再次釋放.
求你的帶領.
我們的同心禱告奉靠主耶穌基督得勝名自求.
阿們.
在說今天的經文.
以先做少少的回應.
因為我覺得有這個需要.
亦都是大英之媒很期望的.
在過去這幾個星期裡面.
特別是回應到在兩三個星期之前.
應該如果沒有記錯.
是兩個多星期之前.
我曾經說到.
有一堂講到有關於代禱與祈禱的.
主一崇拜裡面.
在網上不少的信徒.

$^{41}$有很大的壓力.
他覺得.
聽完黃牧師你的講道之後.
我們很有壓力.
我們聽了這麼多年.
這麼多傳道人.
在教會裡面所傳講的.
有關於禱告與代禱.
都是著重於真誠便可以了.
跟父神說話說甚麼都可以.
在聖經裡面.
是否這樣的啟示呢?.
卻偏偏.
跟我們這麼多年所聽的是自然不同.
我明白.
當將上帝的啟示.
按照聖經講述出來之後.
不少的信徒會覺得.
很難受.
因為這是很難接受的事情.
尤其是我談論到.
禱告一定要以上帝的旨意來禱告.
於是有很多信徒就覺得.
如果是這樣的話.
那我不敢禱告了.
如果這樣的情況.
幸好我不是黃牧師的教會會友.
否則我的禱告就不能夠繼續生存了.
為甚麼呢?.
他們往往的禱告是按照自己個人的意願來禱告.
甚至有些傳道人很懂得包裝.
你的意願就等同於上帝的意願.
或者你求.
將你的意願等同於上帝的意願.
甚至扭轉過來.
將上帝的意願變成你的意願.
這些是很危險的信仰.
基督教.
我們經常強調.
相信上帝.

$^{81}$以祂為主.
其他宗教不說的.
其他宗教說的就是信神.
這個神可以幫你.
這個神就像一個僕人一樣幫你.
但是基督教卻完全不是強調這一點.
基督教所強調的.
所相信的神是我們的主.
如果祂是我們的主.
我們就必須要尊重祂的主權.
所以我們去禱告的.
祈求上帝每一件事.
都是要按照上帝的為主來決定.
如果我們強行將我們自己的意願.
變成上帝的意願.
甚至要上帝滿足我們的意願的話.
我恐防你信的不是上帝.
並不是耶穌基督的父神.
而是相信一個你塑造出來的.
叫做基督教的神.
而這個神和其他異端宗教沒有什麼分別.
因為那個根本就不是耶和華神.
所以在這裡我為什麼這麼強調這一點呢.
我覺得有一位弟兄很好的.
他將我那堂講道.
用AI將它完全將它重新分析.
然後將那個主旨講出來.
黃牧師講的是要告訴今天的信徒.
知道我們信的東西有多少是離經背道的事情.
包括禱告.
弟兄姊妹你不是不需要真誠去父神禱告.
你是需要真誠地去到父神禱告.
但是你去到父神的面前.
你不是要求上帝父神滿足你.
才是父神.
你應該求的是父神.
你在這些事情上面.
你告訴我你的心思是什麼.
這個就是真正的禱告的意義.
而保羅就在歌羅西書裡面.

$^{121}$屢次又屢次做這個的見證.
因為他恐防那些討人喜悅的說話.
神學.
當他在哥羅西教會傳開的時候.
哥羅西的信徒.
會漸漸被這些這麼中聽的題目而吸引了.
甚至有些以為透過這些人為的思想.
神學塑造出來的敬虔的定義.
以為就是上帝定義.
傾虔是什麼.
所以當我們仔細去重新思想歌羅西書的時候.
我們應該更加學會怎樣去尊重上帝的主旨.
而不是我們的心意.
不然正在坐牢的保羅.
在寫歌羅西書的時候.
他就不應該說在歌羅西書裡面說的話.
而是應該說怎樣為我脫離羅馬的監禁而去禱告.
你們哥羅西教會應該怎樣幫助我.
保羅完全沒有這樣說過.
因為他知道.
由禱告以至於牧羊的上面.
是要塑造出來的每一個屬神的子民.
是他們真正信服在上帝的主權之下而活的信徒.
這個正正就是歌羅西書.
要去傳遞的一個很重要的訊息.
因為我們在這個地上.
活出有主的生命的時候.
我們就能夠與主一同去榮耀上帝.
所以為什麼教會今年的主題叫.
與基督同顯他.
即是上帝父神的榮耀.
就是這個原因.
在這個主題的背景之下.
今天我們來到第一章.
第21至23節.
21至23節說些什麼呢?.
我們逐步將它一層一層地剝開.
不過在開始的時候.
我想大家明白一句說話.
香港人有一句俗語.

$^{161}$這個俗語叫做.
執書行頭慘過敗家.
如果文雅一點這樣說.
就是叫做機會難逢.
切勿錯過.
有些什麼事要我們作為基督徒.
今天真的不要執書行頭呢?.
會不會是我們怎樣去更加謙卑.
在上帝面前.
跟從他的旨意而活.
這件事我們不要執書呢?.
還是今天我們只是憑順了主之後.
希望利用上帝的力量.
在這個世界謀取更多.
我們自己的慾望呢?.
這個正正就是今天.
這段經文開始要給我們.
要去更深入去思想的題目.
很多時候帶著一份.
不要錯過機會的心.
我覺得是正確的想法.
不要錯過任何機會.
但是機會往往是.
給那些有預備的人而去承接的.
今天很多的信徒.
信主的過程當中.
很多時候對於這樣.
真是缺乏了去好好緊握.
把握上帝這個拯救的機會.
即或你回了教會都好.
即或你受了洗都好.
我們的生命生活仍然.
表現出一種愛理不理的狀態.
甚至有些信徒覺得.
我信主只不過是剛剛.
我的朋友帶著我回教會.
或者這麼巧去到街邊.
有人派傳單.
我信了主而已.
但是缺乏了一樣東西.

$^{201}$如何在上帝裡面.
能夠真是爭取和上帝同活.
那個基礎我們缺乏了.
怎樣的人才可以.
回到上帝裡面.
這位聖潔的上主裡面同活呢.
無非一定要是對罪.
有一份很敏銳的觸覺.
以及厭惡.
信徒能夠在這個世界上.
真正活著的表現.
是因為我們對這個世界的.
罪惡厭惡.
而信了主的人.
因為有聖靈內注.
也都好自然對罪惡.
有一份的拒絕和厭惡感.
希望從罪的當中去掙扎出來.
而耶穌基督就做了這個工作.
即刻在我們之前上文.
歌羅西書的第一章的第九節至第十五節裡面.
我已經說了保羅是很刻意.
去介紹這位屬神的愛子耶穌基督.
他是怎樣來到這個世界.
而進入到第十五節至第二十節裡面.
更加清清楚楚地去介紹出.
主耶穌基督他那個特別的.
修和的工作.
怎樣去和創造與及他自己.
修和的關係建立在一個永恆的角度當中.
如果我們沒有和上主耶穌基督.
修和的話.
沒有藉著他的工作和創造再次修和的話.
我們也都沒有可能和上帝父神修和.
既然如此.
所以今天保羅沒有在這個題目上面停下來算了.
說了耶穌可以走了.
很多信徒是這樣傳遞的.
傳遞了說了耶穌拯救世人了.
死在十字架上面了.

$^{241}$你的罪得赦免了.
自由了.
喜歡做什麼就做什麼.
不是的.
保羅看到哥羅西教會他們的信是很清楚.
但是他仍然知道.
哥羅西教會是有他的進步的空間.
有一個好的基礎的同時.
保羅希望他們再下一城.
能夠更加走向成熟的地步.
所以在他繼續寫歌羅西書的時候.
他就以身作則地去將.
牧羊那個應有的風範顯露出來.
並且叫一切同心合意的信徒.
都可以一同.
怎樣去連結在基督耶穌上面.
不是說.
而是怎樣去達成到.
是耶穌基督和他真是連結起來的現象.
或者那個表現行為心思.
而這個也都成為了今天我們能夠讀到歌羅西書的時候.
我們所欣賞到保羅他的介紹.
怎樣讓歷世歷代.
與及他先前所寫的.
哥羅西教會這班的對象.
可以真正活在基督耶穌的封聖裡面.
領受一個基督國度裡面.
神子們的聖潔身份.
這個是前文的.
來到今天我們看的第21節至到第23節.
這一段落.
他要再進一步告訴我們.
聽耶穌既然他來到做了他第一步了.
也都做了完成了這個救恩的橋樑的建立之後.
哥羅西教會你們又怎樣呢.
而今天我們讀到這段聖經的時候.
我們也都怎樣的繼續參與在.
耶穌基督這個救恩的基礎的上面呢.
所以今天的講題.
主耶穌已經啟動了.

$^{281}$開始了這個救恩.
我們又怎樣的情況呢.
讓我們回到這段聖經一起去看看.
不是很長.
三節的經文.
我們一起讀一下這段簡單的經文.
希望我們能夠先熟悉一下.
這段經文的內文之後.
然後我們再仔細去看看.
這段經文所想講的內容是什麼.
我們一起打開之後.
我們一起讀一次.
有一節.
你們從前與神隔絕.
因著惡心.
心裡與他為敵.
但如今他借著基督的肉身受死.
叫你們與自己和好.
都成了聖潔.
沒有瑕疵.
無可責備.
把你們引到自己面前.
只要你們在所信的道上行心.
根基穩固.
堅定不移.
不至被引動失去福音的盼望.
這福音就是你們所聽過的.
也是傳語普天下萬人聽的.
我保羅也作了這福音的執事.
如果大家就這樣讀一次.
你就會覺得這段聖經沒什麼好說的.
我都知道了.
你是不是真的都知道這裡所說的話呢?.
我們試試分兩部分去看.
第一部分.
第21到22節.
在聖經裡面說到.
因為耶穌的緣故.
今日我再不是昔日的我.
我的英文寫的不錯.

$^{321}$因為耶穌的緣故.
我再不是昔日的我.
在第21節一開始這樣說.
用了一個很特別的語氣.
來介紹.
或者稱呼.
哥羅西教會在和合本裡面.
他這樣說.
你們從前與上帝界隔絕.
因著惡行心裡與他為敵.
和合本寫得平平無奇.
但是在原文希臘文聖經一開始.
他就用了一個強調式的陳述.
叫做至於你們呢?.
英文就有很多翻譯.
都將這個意思這個語氣譯了出來.
尤其是NET的版本.
一開始他說and you.
中文的翻譯應該是至於你們呢?.
這個是一個很強調的稱呼.
談論到至於哥羅西教會你們的狀態.
保羅在這裡說什麼?.
他說原來繼先前耶穌基督.
將你們已經提供了一個.
修和的工作關係之後.
你們的身份已經是有所改變.
怎麼改變?.
這裡教你寫得夠見證.
誰要受浸或者講福音.
要寫一個得夠見證.
學東西了.
保羅教我們.
教我們怎麼寫.
首先講一下你以前未信主之前的狀態.
所以如果你說.
牧師我準備參加浸禮.
你第一件事.
你在你的得夠見證交離年申請表的時候.
請寫清楚.
你昔日怎麼看你自己在罪裡面的狀態.

$^{361}$保羅在這裡這樣說.
你們從前與神隔絕.
因著惡行心裡與他為敵.
保羅第一件事提了出來.
告訴哥羅西教會知道.
你們昔日的狀態.
是與神有一個完全無關係的狀態.
這裡所說的隔絕.
什麼隔絕.
中文三四十年前可能明白.
隔絕其實是說毫無關係.
你不認識他.
廣東話你們與上帝以前都無親無故.
沒有什麼相關.
就是這個意思.
你以前當上帝沒刀.
這個就是英文NET的翻譯.
你們以前當上帝是一個.
strangers and enemies.
你當他是什麼水.
那你會明白.
什麼水都不認識他.
還有呀.
跟他專門唱雙和.
不是唱雙和.
唱對台.
他講什麼你就要反過來講.
這個就是我經常說的.
不信主的人.
就是一個contrarian.
什麼都跟人搏.
對的又跟人搏.
不對的又跟人搏.
總之要跟人不同.
明明人人說太陽從東方升起.
偏偏就說太陽不是從東方升起.
還要扭橫折曲.
扭一輪choice the topic.
然後是要跟你唱不同.
你以前是這樣的狀態.

$^{401}$保羅有沒有強烈.
他們有些什麼狀態.
與神又為仇又隔絕.
就是當上帝無道呢.
有一節我在這裡.
卻沒有任何東西看見保羅寫.
保羅只是說.
因著惡行心裡與他為敵.
咦.
跟保羅的風格不太像.
大家知不知道保羅是什麼性情.
保羅的性情就是.
人家有什麼錯.
他會直接指出人家有什麼錯.
不用說誰.
加拉太書裡說過.
保羅夠膽直接罵耶穌大弟子彼得.
因為彼得是假冒虔誠.
怎麼假冒.
沒有看見保羅的其他信徒.
他自己和不信的人一起吃飯.
一看見有信徒或保羅出現.
就馬上罵.
縮起來馬上走.
和所有未信主的人馬上有所分隔.
保羅馬上罵.
保羅性情是這樣的.
有時我都覺得.
我有少少沾染了保羅的性情.
弟兄姊妹.
做牧者不是說討人喜悅的話.
聖經要聲明做信徒.
不應該說討人喜悅的話.
但尤其是牧者.
就更加要點出信徒的問題.
所以我說道理是不好聽的.
越聽你越火滾越憎我都不奇怪.
曾經我去到一間教會.
有個姐妹跟我說.
牧師我們個個都很喜歡聽你講道理.

$^{441}$你聽了我一堂講道理就這樣說.
聽多兩堂之後你就會開始很憎我.
結局真是.
我講了十堂八堂之後.
接著她就憎到我死.
我講道理是很引人憎意的.
不過我講道理是為你的好處而去講.
你不喜歡聽不是不聽我講.
我沒有惡感.
但在聖經上如果你干犯的話.
是你有問題.
這是我作為牧者不想見的.
但保羅好像一改常態.
他不直接指出哥羅西教會有甚麼問題.
因為他們三不七七.
所以不想得罪他.
是不是這樣.
你就以為這樣.
但實質上保羅在這段聖經裡面所講的.
就不是針對人的罪是講甚麼.
而是列出一連串的指控.
這不是他寫這段聖經的主旨.
他寫這段主旨.
是要讓哥羅西教會學會怎樣.
參與在耶穌基督所已經成就的救恩這個基礎之上.
他怎樣再更加成熟起來.
所以講昔日的東西沒錯.
你是要知道你昔日自己的狀態.
保羅點到即止.
這個空間留給你自己.
每一個人自己去反思.
但你要反思完之後.
下一步你要了解的就是.
你怎樣去建造在基督這個救恩的基礎之上.
所以這個就是正正保羅我經常都講.
在他自己的神學裡面.
有一個很重要的核心神學.
叫做參與神學.
基督徒是需要參與在基督所建立這個救恩的基礎之上.
基督徒不是福音的消費者.

$^{481}$基督徒卻是福音的積極參與者.
我再重新講多一次.
基督徒不是福音的消費者.
而是福音的積極參與者.
什麼意思?.
基督徒不是我信了福音.
耶穌拯救了我.
接著我可以happy了.
我可以有上帝護體了.
見病醫病見鬼趕鬼.
我都想如果真的是這樣.
不過那些是異端的想法.
你信其他宗教的神就是這樣.
討好你幫你搞定他.
魔鬼最喜歡offer這些東西.
你看看登山上耶穌被試探了四十天.
魔鬼offer給耶穌多少?.
給耶穌多少好處?.
無論在實質的實體生活上.
吃.
現在的人喜歡吃.
魔鬼說不要緊石頭都變麵包.
吃.
宗教上給你做哥哥.
做head.
你上到這個聖殿的頂.
跳下來都不用怕.
因為有天使的手擋著.
你的腳不會跌倒.
宗教厲害.
神蹟都顯現了.
個個來蜂擁擁護你.
你成為教主了.
你看看.
名利.
你只要拜一拜我.
魔鬼說.
跟耶穌說.
拜一拜我.
你看見上面所有的城邦.

$^{521}$權力都給你.
魔鬼最喜歡滿足人.
基督教會裡邊個傳道人.
專門講滿足你呢?.
你就要小心.
魔鬼的代表.
真正假基督.
保羅在聖經裡.
強調的是.
哥羅西教會.
你既然昔日是和上帝.
成為陌生人.
甚至對敵的狀態.
但今天卻藉著主耶穌基督.
而得到打開一個新的門路.
這條的門徑.
讓你能夠和上帝怎麼樣.
我們再看第二節.
但如今.
他藉著基督的肉身受死.
叫你們與自己和好.
都成了聖潔.
沒有瑕疵.
無可指責.
把你們引到自己面前.
昔日你自己慢慢想.
但將來你怎麼做呢?.
這個就是主耶穌.
所希望做到的.
在這個地上我們要做的事.
而保羅把他坦白的說了出來.
耶穌的受死.
今天我們領聖餐.
這個餅這個杯.
餅就是肉身.
耶穌的死為我們承擔了這個罪的刑罰.
我們說過為什麼要用死來承擔.
因為死是罪帶來的最嚴厲的問題.
既然這麼嚴厲的問題.
就要用這麼嚴厲的方法.

$^{561}$將它來到花的力氣.
才可以將它改變.
扭轉過來.
目的為什麼要做這麼多事呢?.
叫你們與自己和好.
這個自己指誰呢?.
耶穌基督.
甚至上帝.
請大家弟兄姊妹留意.
和好這個字眼.
如果你有聖經.
是英文版本的.
你看看這個和好這個字眼.
是用了一個英文翻譯叫Reconciled.
和好.
即是原本不好.
現在就變回好.
首先我要強調和好.
不是世人定義的和好.
是在基督耶穌裡面定義的和好.
我再強調多說一樣東西.
和好不等於和平.
和好不代表我們所看的和平.
聽聞中東現在在打仗.
有沒有看新聞的?.
我看圖片也很厲害.
新聞片.
七百多支的導彈不斷飛起.
不斷飛進去不同的地方.
有很多信徒就發起了一個禱告.
尤其是我們教會.
有些教會的同工就說.
我要為世界和平禱告.
希望早日戰爭完結.
這是人的意願.
教會很喜歡用另一個字.
將和平重新演繹.
合一.
教會經常說合一.
什麼合一?.

$^{601}$在哪裡合一?.
不說的.
為了合一.
什麼都妥協了.
所以今天的教會.
不要吵了.
大家不要吵了.
有意見不要吵了.
合一.
妥協了.
尤其是教會.
會不會在教會的領袖當中.
有人犯罪呢?.
這句話我經常聽.
不是說我們教會這麼簡單.
我們教會以前試過了.
我所參與過的教會.
幾千人的教會也好.
當教會領袖.
犯罪之後.
為了合一.
不要說了.
各位姐妹.
這是用人的定義.
不要說就是合一.
不要說只是這件事沒有解決.
不是合一.
每個人心裡面是不是合一.
合一在哪裡?.
合一在骨氣那裡合一.
一起骨氣.
因為這件事沒有解決.
什麼叫做和好?.
在聖經這裡.
是上帝的定義.
所以請留心看清楚.
保羅怎麼寫.
叫你們與自己.
我特意再說一次.
自己是誰?.

$^{641}$保羅?.
不是.
耶穌基督.
因為上面說.
但如今借著基督肉身受死.
叫你們與自己.
這個是說耶穌基督受死.
為的是怎樣?.
與他來比鞠.
我們能夠和他和好.
原本不好.
變回做好.
這個定義在耶穌基督身上.
這個定義是怎樣?.
看完下面.
有一連串的描述.
告訴你和好的定義是什麼?.
是在當中成了聖潔.
沒有瑕疵.
無可責備.
而且可以帶你去到他面前.
弟兄姊妹.
這個修和的基礎.
或者和好的基礎.
是要在上帝所定的聖潔標準.
沒有瑕疵的標準.
無可責備的標準.
可以去到見到上帝的標準.
所以和好不是你好我好不出聲.
那些叫和好.
那些叫不出聲.
不出聲就不代表和好的.
可能心裡不斷在詛咒你.
不出聲而已.
沒有和好的.
但這裡強調的是什麼?.
是有聖潔.
沒有瑕疵.
無可指責.
誰沒有可指責呢?.

$^{681}$在這段聖經.
可能你未必這麼清楚.
但當我們遲一點.
當我們去讀到第二章的時候.
如果你喜歡的可以偷看一下.
或者很想知道多一些少.
在第二章的第16節至第19節裡面.
就談論到有關於一些異端的人.
在指控審判.
哥羅西教會裡面那些忠心信主的人.
他們用了很多人為的宗教條件.
定立了.
來指控哥羅西教會的人.
你們不虔誠.
不合格.
成為神的子民.
那些假的傳道就用了.
人為塑造出來的條件.
來量度哥羅西教會.
今天教會都是這樣的.
很多信徒集體來定立自己的標準.
來量度一個人的信.
究竟是不是虔誠.
不用說什麼.
說一個最簡單最簡單的例子.
今天的信徒.
絕大部分包括牧者.
認為虔誠的信徒.
一定要是回教會參加實體崇拜.
你不參加實體崇拜就不虔誠.
可能弟兄姊妹心裡都有這種壓力.
牧師今天我沒有參加實體崇拜.
我覺得很內疚.
你知不知道以前在我們教會有一對情侶.
這對情侶很有趣.
牧師不參加教會.
牧師之前就會來跟我說.
牧師我要請教.
我要請什麼教.
我未來兩個星期不參加教會.

$^{721}$為什麼.
我要去旅行.
我去旅行也不用請教.
請什麼教.
現在我們還好.
如果你不適身體不適.
你去實體崇拜.
所以為什麼我們搞這麼久.
大家明不明白.
我不是為了方便你.
而是我希望成全你對上主追求的心.
當你有這份心思意念在上帝裡面.
你不回實體崇拜不代表你不敬虔.
但現在的教會教牧童宮.
多少仍然是一定要推回教會的信徒.
回去教會實體崇拜.
你知不知道為什麼.
因為不回去實體崇拜的話.
收的奉獻會少了.
讓我說句不是陰毒.
這些表現是活該的.
整天搞mega church model.
讓教堂大一大.
沒有一千二千人聚會的教會.
零丁單薄的人.
大家沒有了那種氣氛.
覺得不夠虔誠.
不夠熱切.
不夠warm.
搞不起.
搞到每個人都要推信徒回教會.
兄弟姐妹.
首先我沒有否定實體崇拜的好處.
沒有.
你回去就回去.
我覺得回實體崇拜還有一個好處.
操練一下你的決心.
回教會.
特別是冬天.
由11月至3月.

$^{761}$你都不知道多想在床上的棉胎捲起來.
睡多一會.
你要掙扎.
從棉胎跳出來.
為了什麼?回教會.
是一個很好的學習.
鍛鍊.
但不代表因為你身體不適.
你回不了教會.
或者因為天氣的問題.
因為交通的問題.
回不了教會就影響你不敬虔.
回到這段聖經.
你要有這樣的條件.
聖潔無有瑕疵.
無可指責.
全部在上帝的標準裡面去看.
不是人為你不回崇拜.
實體崇拜你就不敬虔.
上帝看的是你追求的生命是怎樣.
你對上帝的尊重.
他的主權有多尊重.
來決定你的敬虔.
同樣道理.
禱告都一樣.
你對上帝沒有敬虔的心.
你只不過當上帝好像有求必應.
是你的滿足者一樣的情況.
你對他根本沒有尊重.
你就算回去崇拜祈禱多長都沒有用.
不是我說的.
是耶穌在馬太福音裡面.
論禱告講得很清楚的.
有些人在十字路口祈禱.
怕死人聽不到.
敬不敬虔?耶穌說這些假敬虔.
所以弟兄姊妹.
當我們如果能夠在上帝的面前.
可以示納的.
我們必須要有聖潔無瑕疵無雜備的評審.

$^{801}$過得到上帝.
當如果我們都能夠過得到上帝的話.
我們還用怕這個世界的人的基礎.
或者條件.
或者他們的神學.
來去指控我們不夠敬虔.
保羅在這裡.
為哥羅西教會打好一個很好的根基.
因為保羅其實聽到消息了.
經過以巴弗.
告訴他知道.
哥羅西教會在裡面有些人.
用很多的神學來扭曲上帝的說話.
用人的基礎.
用人的神學來取代上帝的啟示.
來去倒過來指控哥羅西教會.
這些我遲日會看的.
但是保羅在這裡告訴他.
過得到上帝你還怕什麼.
因為在這個世界上只有一個人.
是只能夠審判信徒是否敬虔的.
是否合規格回到上帝裡面的.
那個就是上帝他自己一個.
不是人.
不是哪一個.
連黃牧師都不是.
所以有些人聽完我講道之後.
他說黃牧師你講朕什麼會.
孫什麼會.
那些牧者都是這樣.
你自己都是浸信會的.
是啊.
我們教會都有這樣的人.
你不要以為我們教會沒有.
胡胡混混在那裡祈禱.
講極不聽.
希望沒有了以後.
我講道不只是講別人不好.
我們自己都不好.
怎樣改過呢.

$^{841}$保羅在這裡.
就是要讓哥羅西教會打好基礎.
不用怕將來有人會再來指控你.
使你覺得我們在主的裡面.
是不能夠走下去的.
基督徒在信仰的路上.
最多最多被攻擊的就是說.
你不像基督徒.
你知不知道當年我在.
滑鐵盧大學讀書的時候.
在圖書館.
很努力地讀書.
突然間遠遠聽見有個同學.
中國人來的.
很大聲地講粗口.
講廣東話粗口.
接著有個人在現場.
馬上用廣東話回答.
你還說自己是基督徒.
公眾講粗口.
弟兄姊妹.
當如果一個基督徒.
連上帝的基礎,聖潔.
還有瑕疵.
沒有瑕疵沒有可指責的時候.
都做不到的.
那個時候就應該值得我們反思.
哥羅西教會有沒有這樣呢.
保羅留個空間給他去想.
但至少保羅說.
如果你真的去決意相信上帝.
活出一個聖潔的生命.
無有瑕疵無可責備的話.
上帝的面前你都站得到的.
你不需要擔心.
在你裡面的人去指控你.
而這個也都成為了.
他們將來的整個.
屬靈的基礎的發展的一個.
很重要的基石.

$^{881}$免得很容易被人打殘.
今天太多信徒.
就是因為信了主之後很殘.
一信了主受了浸之後就散了.
你怎麼知道呢.
其實你看看你受浸的時候.
牧師或者教會送給你的那本聖經.
你就知道了.
還有沒有啊.
那本聖經是做什麼的.
是不是成為了.
你天天讀的聖經.
還是那個我們去攝.
灰塵的攝石.
攝了很多灰塵在上面的呢.
甚至可能你那本丟了也不奇怪了.
我曾經試過在我們街口.
有個收舊書的垃圾站.
見到有人扔幾本聖經在那裡.
我說誰扔聖經呢.
看一看.
什麼什麼姐妹競爭受浸紀念.
我就仁慈的想.
或者他去世了.
他的家人放了他在那裡.
我希望不是他自己放在那裡.
姐妹.
真正的基督徒是怎麼樣的.
我再重複講一次.
不是福音的消費者.
是福音的積極參與者.
他積極的在這個世界裡面.
去作見證.
活出一個聖潔無瑕疵無可指責.
能夠過得到神的生命生活.
而不是只是凍在那裡.
好像基督徒非基督徒的狀態.
就好像我們教會街口.
那個運動場的二樓的玻璃窗裡面.
你說是嗎.

$^{921}$牧師那裡發生什麼事了.
你知不知道我們教會.
對出街口的運動場二樓.
滑雪場那裡.
每天早上都有個人站在那裡.
很努力的站在那裡.
他站在窗邊一直看著你們.
讓每個人都很害怕.
有人監視著我們.
誰不知原來他是一個紙板人.
在那裡嚇嚇你們.
讓你們不要亂泊車.
很多基督徒就是這樣生活.
凍在那裡.
我是基督徒.
但是沒有生命.
不活的.
好像一個殭屍那樣站在那裡.
上個星期為什麼唱詩.
要唱《開到水深之處》兩次.
不要整天只是在碼頭那裡.
開到上帝恩典的海洋深處.
獲取更多上帝恩典.
保羅在這裡.
要讓每一個哥羅西教會的信徒.
以及今天讀到哥羅西教會.
這本書的信徒都知道.
我們是福音的參與者.
而不是福音的消費者.
接著來到第23節.
23節這樣說.
只要你們在所信的道上恆心.
根基穩固.
堅定不移.
不至被引動失去福音的盼望.
這福音就是你們所聽過的.
也是傳語普天下萬人聽的.
我保羅也作了這福音的執事.
你讀到這裡.
又是那三幅皮.

$^{961}$叫我們要有恆心.
要有堅固的信仰.
不要動搖.
明白了.
是不是這麼明白呢.
我小時候香港有個廣告.
叫十大洗腦廣告之一.
這個廣告說到一個小朋友.
放學回家之後.
打開冰箱.
看到有一支飲品.
那支飲品叫益麥多.
接著當他打開的時候.
準備去拿.
他外婆走過來跟他說.
不要喝啊.
他說為什麼不要喝啊.
有菌的.
我聽你姐姐說那支飲品有菌的.
喝了你會壞的.
有菌.
他就很不開心.
為了關掉冰箱.
想著走開.
接著走開沒多久.
轉頭他姐姐就放學進來了.
接著他進來之後.
打開冰箱.
馬上看到一支益麥多.
然後他就拿了出來.
撕開它喝.
接著臨入口那一刻.
那個小朋友.
就是他的弟弟就跟他說.
有菌的.
不要喝啊姐姐.
接著那句廣告的對白.
就這樣說.
當然有菌.
性乳酸菌.

$^{1001}$有益的.
你今天喝了沒有.
我現在喝了.
你今天喝了沒有.
我喝了.
看一下這裡.
最後那句.
保羅說什麼.
保羅說這個福音.
是你們聽過的.
有時傳語普天下萬人聽的.
我保羅也作了這福音的執事.
保羅我已經參與了.
去做這個傳福音的執事.
我保羅視福音為我參與的目的.
而不是消費福音的好處.
沒錯福音帶來有好處.
讓我們得以拯救.
回到上帝那裡可以無罪的.
和上帝活在一起.
但是福音的核心.
是要叫你和我在今天.
參與在傳福音的執事上面.
我保羅今天參加了.
你做了沒有呢.
活性乳酸菌.
阿婆不知道是有益的.
當然叫你不要喝了.
有菌的.
但原來是對身體腸胃有好處的.
或者我也需要喝一點.
弟兄姊妹.
什麼叫做根基穩固堅定不移呢.
英文翻譯原文用了兩個字.
叫做establish and firm.
什麼意思.
establish的意思就是說要成立的.
firm這個字是很特別的.
firm這個字眼.
如果用來作動詞就叫做堅固.

$^{1041}$如果用作名詞.
這個字眼除了講堅固之外.
可以指道一間公司.
原文用動詞.
所以我用動詞去解釋.
是要去經營的.
你要站得穩.
establish.
不只是開了一個詞.
所以大家有沒有見過以前.
大家大部分都是從香港來的.
我們這些老餅人士.
以前的店鋪.
特別是我住舊區.
店鋪面前有個很大的招牌.
寫著黃成基.
對不起我家店鋪的名字.
寫了個名字之後.
通常角落會寫著兩個西字.
成立.
由哪一年開始.
哪一年.
康熙第四年.
當然沒有了.
我作一個而已.
成立於1996年.
厲害勇敢.
97年之前你敢開檔.
這麼有信心.
成立於1900年.
更厲害.
民國都還沒開始.
是代表什麼.
開業.
信主有開始的開端.
沒問題的.
你需要有個開端.
你知道你的罪.
然後你自己回轉.
看到你對罪很憎惡.

$^{1081}$所以你要成立一個.
成為一個基督徒.
立志走在上帝的裡面.
活出聖潔無可指責.
無可責備的情況.
好.
但是你要繼續經營下去.
這是什麼神學.
參與神學.
講到悔改.
就是repenting.
不住認罪.
而是repented認過罪.
所以在《原文聖經》裡面.
這些字眼講得很漂亮.
告訴哥羅西教會.
以至今天你和我聽見的時候.
我們要在這個道上.
要做什麼.
不是說我信了主.
信了受針.
搞定了.
不是.
而是要繼續去.
在這個基礎上.
不只是去開檔.
而是要繼續的栽培.
經營下去.
所以廣東人都有句說話.
創業容易.
下面那句是什麼.
首業難.
基督徒的人生.
重點不是創.
開始和上帝接觸那一刻.
認了罪.
往後你怎樣去經營下去.
不要動搖.
怎樣才可以不動搖.
一路記住.

$^{1121}$那個福音的盼望成為動力.
這個不是我講的.
是保羅講的.
什麼意思.
你期望更加好.
昨日的我和今天的我.
應該不一樣.
因為昨日的我仍然犯罪.
今天的我和上帝和好.
明天的我怎樣.
你應該和上帝更加親密.
這個才是基督徒的人生.
哥羅西教會的建造.
是應該這樣的.
一個教會.
一個信徒.
能夠成長起來.
這個才是表現出你真正的信徒.
不然的話.
你仍然只不過是一個福音的消費者.
而不是一個福音的參與者.
保羅說.
我已經接受了這個邀請.
我去主動做這個福音的執事.
我去參與在這個福音上面.
不單止去傳福音.
我甚至將這個福音付諸實行.
能夠在我生命上面可以得到改變.
使得我更加確信上主.
我就做了.
你們怎樣呢.
弟兄姊妹.
我將這條問題.
交給你們.
保羅做了.
你有沒有.
同樣去參與做這件事情呢.
而不是只懂得喝益力多.
對不起 說了名字.
只懂得.

$^{1161}$你就喝了.
我沒有.
我喝了沒有呢.
不是的.
聖經本身的啟示不是難的.
這裡每個字你都懂得解釋.
為什麼做基督徒難呢.
就是因為我們不想離開罪.
那個阻力很大.
要順性去傳一些你喜歡聽的東西.
很容易的.
就好像那些牧者這樣說.
祈禱最重要.
真誠就行了.
跟上帝說什麼就行了.
上帝是大能大力的神.
一定會按照你的.
合神心意的禱告.
就一定會.
你就聽不到合神心意那個字.
不過你就是聽.
滿足你那個禱告的需要.
弟兄姊妹.
不要中招.
保羅做這麼多事情.
說這麼多事情.
你每個字都懂得看.
其實是要告訴哥羅西教會.
不要中招.
不中招之餘還要怎樣.
積極參與.
這樣你的生命.
才像基督的子民.
活在他的啟示當中.
活在他的心意之下.
讓我們一起同心低頭禱告.
祈求這份心志.
能夠在我們每一個人身上.
父神謝謝你讓我們在聖經裡學習.
你讓保羅.

$^{1201}$將他自己的心思意念.
以及他參與福音的意願.
清楚去作見證.
告訴昔日的哥羅西教會.
所有的信徒知道.
他沒有為自己的利益.
或者自己的盼望而祈求.
因為我相信.
只要我們的祈禱合上帝的心意的話.
沒有一樣東西.
上帝是不知道我們的需要.
求你的幫助大靈.
讓我們準確去找緊上主你的心思.
讓我們能夠在上主裡面.
能夠積極參與在.
福音的執事的崗位之上.
求你的大靈.
也都給我們有智慧.
懂識在這個世界各式異端.
所帶來的吸引和引誘.
求你的大靈.
讓我們有力量透過聖靈的鼓勵.
以及你聖經的啟示.
堅固我們的心.
禱告奉靠主耶穌基督得勝名自求.
阿們.
\newpage



\section{以弗所書 4:1-3}
\label{sec:031719}
\textbf{建立健康的使命教會 (6) 竭力同心同行}~~~~鍾梓昭傳道
\newline
\newline
連結: \href{http://www.church.com.hk/acms/content.asp?site=cdc\&op=show\&type=product\&code=031719\&layout=sermon}{\texttt{031719}} ~~~~ 語音日期: 2026-03-01
\newline
\newline
\hyperref[sec:031718]{\small{< < < PREV SERMON < < <}}
~
\hyperref[sec:index_chronic]{\small{[返順時目]}}
~
\hyperref[sec:031720]{\small{> > > NEXT SERMON > > >}}
\newline
\newline
以弗所書 4:1-3
\newline
\begin{longtable}{cl}
\hline
\hline
章節 & 經文 (和合本修訂版)\\
\hline
4:1 & \begin{tabularx}{0.7\textwidth}{X} 我為主作囚徒的勸你們，既然蒙召，行事為人就要與你們所蒙的呼召相稱。 \end{tabularx} \\ \\ \relax
4:2 & \begin{tabularx}{0.7\textwidth}{X} 凡事要謙虛、溫柔、忍耐，用愛心互相寬容， \end{tabularx} \\ \\ \relax
4:3 & \begin{tabularx}{0.7\textwidth}{X} 以和平彼此聯繫，竭力保持聖靈所賜的合一。 \end{tabularx} \\ \\
[1ex]
\hline
\hline
\end{longtable}
$^{1}$各位弟兄姊妹早晨.
今日講到的題目是.
「竭力同心同行」.
讓我們一齊先低頭祈禱.
愛我們的主.
我們來到你自己的庚前.
求你謙卑我們的心.
聆聽你自己的話.
多謝你招聚我們成為一.
求你幫助我們學習.
如何竭力同心同行.
成為一個合一的群體.
奉耶穌的名祈禱.
阿們.
我們會花一點時間.
一起看看以下這四方面.
第一.
我們所領受的究竟是一個什麼的呼召.
第二.
我們所被交託的.
是一個什麼的責任.
我們可以如何保持聖靈所賜予的合一.
而保持合一.
我們要付出什麼代價.
教會就好比一個交響樂團.
而聖靈所已經賜予的合一.
就好像一首已經被編寫好的交響樂曲.
這首樂曲其實是由主耶穌基督.
祂自己親自在十字架上.
編寫,創作.
而且已經創作完成.
然後就把這首樂曲交給聖靈.
在教會裡不斷地作出指揮.
所以在講到之前.
讓我先講一個故事.
讓我們透過這個穿插在講道中間的故事.
幫助我們學習怎樣的.
在神的交響樂團裡一起演奏.
有一個人.
他加入了座落在斯都菲爾社區裡的一個管弦樂團.

$^{41}$原來他本身是一個已經受過專業訓練的小提琴師.
經過多年的課堂深造學習.
又花了無數數以萬計的鐘頭去刻苦練習.
所以他對於他自己的音樂造詣是充滿了自信.
而且非常期待可以參與公開的演出.
他第一天來到這個管弦樂團參與排練.
心情是很興奮,很雀躍.
不過幾分鐘之後.
他開始覺得很失望,很失落.
因為每次當同管樂加入的時候.
他們的聲量實在是太大了.
蓋過所有其他樂器,所有音樂要發出的旋律.
至於弦樂在整個演奏過程裡面.
有經常的都愈拉愈快.
敲擊樂每次加入的時候都總是慢了個半拍.
整個感覺就是很嘈吵,很亂,很不協調.
就好像以下這個示範.
每一個人或者每一班樂師就好像一個自己的結黨.
紛爭中自己有自己的指揮,自己要彈奏的音樂.
在同一時間做自己要做的音樂.
就有點像這個情況出現.
演員,創作,你會不會喜歡聽這些音樂呢?.
你想再參加一個這樣的交流嗎?.
如果你是這個交流的成員.
你會不會覺得自己很差勁呢?.
在中場休息的時候.
這個小提琴家他在搖頭,嘆氣,自言自語地說.
這裡個個都真的很有才華.
不過為什麼這首寫得這麼好的交響的樂曲.
演奏出來的時候會變得這麼難聽.
指揮他聽到這句說話就帶著微笑.
而且很溫柔地對他說.
問題並不是在於才華.
而是在於他們沒有懂得去聽.
不單止要懂得聽自己.
更加要懂得聽別人.
要彼此的聆聽.
指揮的這句說話深深地印在小提琴家的心裡.
親愛的弟兄姊妹.
其實教會很多時候都好像這個管弦樂團一樣.

$^{81}$我們都是同樣已經得救.
我們都是同樣對信仰認真.
我們同樣個個都很有恩賜.
不過要學會怎樣一起同心同行.
一起活出和諧合一的信仰生活.
是需要時間,是需要謙卑.
是需要耐性,是需要刻意地努力.
一起來付出.
這個正正就是保羅在二弗所書第四章.
所要寫的內容.
他要寫第四章的主要原因.
因為他在前三章講完了一連串.
很豐富的聖經真理神學教導之後.
保羅在這裡就將焦點轉向實際的信仰生活.
對我們這樣說.
我勸你們.
行事為人要和所蒙的呼召相稱,對等.
信仰不單止是頭腦上的相信.
不單止是口頭上的宣認.
而是要實踐,是要操練,是要活出來.
不是一個人活出來.
是一班人在群體生活裡將信仰活出來.
所以懂得看譜並不代表.
並不等於懂得一起演奏.
既然我們都是被召進入神的這個交響樂團.
我們就要學會怎樣一起演奏.
在生活裡活出和諧和合一的群體生活.
保羅在未講怎樣透過行事為人來活出這個信仰之前.
他首先提醒我們其實我們有一個身份.
讓我們知道我們是誰.
他說我們是被召的.
這個呼召不單止是個人性.
而是群體性.
我們不單止是被召去做一個懂得單玩獨奏的基督徒.
我們是被召進入基督的這個身體.
進入神的這個交響樂團裡.
來一起協奏.
這個小提琴師慢慢明白了.
知道自己並不是被邀請來表演獨奏.
而是被邀請進入這個管弦樂團裡一起演奏.

$^{121}$大家都知道在一個管弦樂團裡有不同的樂器.
每一件樂器都有它自己所要負責演奏的部分.
每個部分其實是同樣重要.
有些負責協奏出主要的旋律.
有些負責協奏出旋律當中的和諧.
不過我們沒有一個人是單獨地為自己而存在.
沒有一個人是為自己而演奏.
It's not about you.
It's all about God.
將榮耀歸回給神.
二分所書第二章說 基督祂自己就是你和我的和平.
是我們的和平.
祂透過十字架的犧牲.
祂抹開自己的身體.
就是要建立一個合一的身體.
祂拆毀了中間隔斷的牆.
使兩下合而為一.
成為一體 一個的身體.
主耶穌不僅使我們和神和好合一.
祂更使我們彼此和好能夠合一.
祂不僅拯救了個別的個體.
祂更建立了一個信仰的群體.
我們稱之為一個新人類.
我們現在不再是外人.
我們不是外人.
我們不再是客人.
我們不是坐在這裡做客人.
我們是神家裡的人.
被焦聚歸入基督的這個身體.
彼此互為手足 互為肢體.
同被建造成為.
神藉著聖靈所居住的居所.
現在來到二分所書第四章.
我們知道聖靈所賜的合一.
這個觀念其實已經是一個.
講解得很清楚的屬靈現實.
不知道大家有沒有面對現實.
所以問題不是在問我們是否能夠合一.
而是問我們怎樣.
怎樣才能夠將這個合一的現實.

$^{161}$能夠呈現出來.
所以當我們停止.
將教會當成一群互不相干的陌生人的聚會.
彼此冷眼旁觀.
停止.
而是好像一家人一樣.
彼此相顧 看待 投入 參與.
彼此建立 彼此服侍.
合一就自然會在我們中間呈現出來.
所以保羅勸我們行事為人.
要和我們所蒙的呼召相稱.
意思是我們需要一種這樣的生活方式.
而這個生活方式.
是要和主耶穌所成就的和平.
所成就的合一相稱的 匹配的.
保羅說用和平.
就是耶穌已經建立的和平彼此聯絡.
竭力保守.
聖靈已經賜給教會的.
合而為一的心.
留意保羅並不是說創造合一.
製造合一或是去建造合一.
合一其實已經存在.
被主耶穌所成就.
被聖靈所賜予.
而且一直是由聖靈所維繫.
不過我們是要竭力去保守.
保守聖靈所賜的這個合一.
而保守的意思就是要保護.
看守 看緊它.
保護 守護它.
這是我們從主.
那位呼召我們主.
所交託給我們的責任.
合一其實是很珍貴.
價值連城.
所以我們要好好珍惜它.
不過合一同時也是很脆弱.
很fragile.
所以很需要我們竭盡.

$^{201}$我們一切的努力去保護它.
就好像一個人會竭盡.
他的努力去保護.
一件非常珍貴的家傳之寶.
不會讓他有任何的機會.
受到任何的破壞或傷害一樣.
在某一次的排列中間.
指揮突然停下來.
中斷了管弦樂團的排練.
對樂團來說.
樂譜上的音樂.
本身已經被編寫得非常完美無瑕.
沒有任何的缺陷或問題.
問題是當你們在演奏的時候.
你們個個都只顧著看自己的樂譜.
個個都只顧著演奏你自己認為對的音樂.
而完全忘記了看樂譜.
合一就好像樂譜上的音樂一樣.
其實一早就已經被編寫到.
裡面是有一份Harmony.
是有一個Unity.
是寫得很完美.
不過演奏的人需要透過排練,練習.
去竭力保守音樂裡面的Unity.
這個Harmony.
而教會當中的合一也是一樣.
同樣需要我們去很努力的.
去彩排,排練,練習.
以致我們能夠在群體生活裡面.
能夠活出在主裡面的和諧.
教會出現分裂.
甚少是因為某些教義出現分歧.
引起紛爭.
而是很多時候只是因為一些.
不經意,不自覺的自我中心.
自尊作祟,驕傲自大.
沒有耐性去彼此聆聽.
堅持自己的意見,看法,做法.
不願意去接納其他人的意見.
當被冒犯的時候.

$^{241}$又不願意去面對.
不願意去處理那些未解決的衝突.
人與人之間出現一些不必要的誤會.
心靈或者情感受到傷害.
選擇不直接溝通.
反而選擇在人的背後說三道四.
說人的壞話.
第一,需要你和我刻意努力去保守.
而我們需要經常看著誰?.
我們的元首主耶穌.
因為祂才是真正教會的頭.
我們也需要不斷的敏銳.
聖靈在我們心裡.
在我們群體中間的引導.
因為祂才是我們樂團真正的指揮.
只有當我們每一個人都專注.
留意主耶穌的引導,帶領,指揮.
我們才能拼湊出一幅完整.
能夠反映神榮耀的圖畫.
讓人看到的就是主耶穌祂自己.
保羅接著指出.
我們可以怎樣去保守聖靈所賜的合一.
並不是透過更新教會的章則.
並不是透過改變教會的架構.
而是透過建立基督的品格.
所以凡事要謙虛,溫柔,忍耐,用愛心.
彼此來互相寬容,用和平彼此聯絡.
這就是行事為人要和我們所夢的呼召相稱的實踐方向.
小提琴詩.
需要明白雖然它所演奏的部分是重要.
不過它所要負責的部分並不等同.
並不代表整個交響樂的全部.
有時它需要拉得輕一點,細聲一點,溫柔一點.
好讓其他的音樂部分能夠被聽見,被凸顯.
甚至乎是被襯托出來.
謙卑的人會將其他人.
看得比自己更加重要,更加有價值.
以至願意學習放下自己.
放下自己的自我中心.
所以主耶穌說若有人要跟從我.

$^{281}$就要捨己.
天天背起自己的十字架跟從我.
不知弟兄姊妹你願不願意學習捨己.
在群體裡學習捨己.
你都要跟從主祂自己.
謙卑的人明白.
自己的觀點並不是唯一的觀點.
不需要永遠都是自己正確.
不需要經常要人人聽取和接納自己的意見.
不需要經常要贏得人的認同和讚賞.
不需要經常要話事作主.
不需要經常霸佔主導領導的地位.
合一並不是統一.
而是要求同.
不過同時是存異.
因為神所創造的合一.
或者神三位一體之合一.
就是一種很特別的.
Unity in Diversity 的合一.
一種能夠同時容納多元.
同時能夠容納不同的合一.
合一是會死.
合一是會死於驕傲自大和自我中心.
所以每當自我自大堅持自己.
都要將自己凌駕在其他一切的惡氣之上.
交響樂樂裡的那個合一.
那個和諧就會完全被破壞被瓦解.
除了謙虛之外.
第二個最重要的基督品格就是溫柔.
因為溫柔並不代表軟弱無力.
而是一種可以受到控制的力量.
一種能夠以柔制剛的能量.
指揮從來都不需要大聲的破口大罵.
每當他需要作出矯正的時候.
他只需要用一些很溫和的手勢.
輕輕的來作出一些溫馨的提示.
能夠將演奏的那個旋律的聲量或者速度.
輕輕的很溫柔的改變過來.
溫柔是會選擇去建立.
而不是去欺壓令對方感覺到被壓被逼.

$^{321}$溫柔是會選擇去矯正.
不過同時是不會去羞辱.
不會令對方感覺到無地自擁.
溫柔是會選擇用柔和的說話.
去代替尖銳,尖酸,刻薄的說話.
是會選擇採取修復關係的行動.
代替報復的行動.
溫柔並不會忽視真正的問題.
不過是會用充滿恩典的方式.
來去正視那個問題和作出應對.
溫柔並不會迴避真相和不講真理.
而是會用保護守護合一的方式去正視真相.
自己的真相,其他人的真相.
和用充滿恩典的方式來表達真理.
在教會裡面我們不能用苛刻的說話去締造和平.
嚴苛的言辭或者會令人覺得很有那種震懾力.
很有威嚴,很有權柄,很有地位.
但卻是會破壞音樂裡面的和諧.
唯有溫柔反而能夠維繫音樂裡面的那種完整性,合一性.
Keep the music intact.
除了謙虛和溫柔之外.
第三個最重要的基督品格是忍耐.
因為我們都需要忍耐來繼續排練.
直至到協奏曲能夠和諧地被演奏出來.
有些樂手一直都不太掌握得到.
不太拿捏得到自己要演奏的部分.
整個樂團都需要不斷重複.
又重複,又重複來作出排練.
懂得忍耐的人.
知道要達到真正的和諧.
的而且確是需要時間.
保羅說我們要用愛心互相寬容.
其實這句說話後面假設了我們的成長以至成熟.
其實都是很需要一段很長的時間.
我們不可能在信主那一刻.
就「叮」一聲.
在那一剎那.
能夠由嬰兒搖身一變.
就成為一個已經長大成人的成年人.
變得成熟,變得完全.

$^{361}$因為在這個「Already but not yet」的追求成聖過程裡面.
我們每一個人都仍然會很「岩岩慘慘」.
三尖八角.
很不完全.
都會有機會令彼此受到傷害.
令彼此都感覺到很失望.
上星期Ada教父提到成長.
有三個很重要很重要的因素.
第一是因點.
第二是真理.
第三是時間.
我們不單止需要有因點的溫度.
是要暖的.
第二我們不單止需要有真理的亮度.
我們更加需要有時間的深度.
記不記得?.
而在這個時間的深度裡面.
忍耐扮演一個很重要的角色.
因為忍耐能夠容讓成長發生.
又同時能夠維繫合一.
是一個很關鍵的元素.
有忍耐的人會懂得說.
神祂仍然在做事.
神祂仍然在塑造祂,塑造你.
不過亦在塑造我.
當神繼續在你和我身上做事的時候.
我是會毀身繼續堅持和你同行.
不會放棄你.
因為神從來都沒有放棄過你和我.
所以合一是要靠有耐性地彼此同心同行.
以致這個合一能夠被維繫起來.
最後除了謙虛,溫柔和忍耐之外.
第四個亦是最後一個最重要的基督品格.
就是愛.
因為愛並不是對關係的一種感受,一種感覺.
而是對關係的一份毀身.
是一個決定,選擇,毀身.
所以愛是一種不離不棄的投入.
和參與.
愛是恆久忍耐.

$^{401}$以致會選擇迅速的寬恕,接納,包容.
不過同時相信它會被上帝慢慢的雕琢改變.
愛是會重視整體的合一.
而不只是重視個人的表現.
是會重視彼此之間的和諧.
而不只是重視個人得到其他人的讚賞.
若說沒有了愛.
謙虛會變成虛偽,虛假.
溫柔會變得軟弱,無力.
忍耐變成啞忍和無奈.
而我們的演奏只不過是一堆聲音.
是明羅,響拔而已.
因為沒有了愛 我們甚麼都不是.
所以當我們願意去愛的時候.
我們要問.
甚麼能夠保持教會的合一.
聖靈所賜的合一.
甚麼能夠建立自己的生命.
人類的生命.
建立基督的這個身體.
甚麼能夠幫助我們去活出主耶穌基督祂自己的生命.
因為祂就是我們的和平.
祂就是我們的合一.
曾經有一位姐妹她說.
其實我一直都不知道教會有多愛我.
直至有人在半夜來關心我,探我.
原來不需要長篇大論的道理.
不需要面紅耳熱的來辯論.
或者查經.
愛就是靜悄悄的.
透過行動.
將主耶穌基督的身體圍繫起來.
愛能夠將信仰變得有血有肉.
更加有生命力.
愛能夠將主耶穌基督所成就的和平合一.
活現出來.
不知道大家還記不記得.
我在講道之前曾經說過.
教會就好比一個管弦樂團.
聖靈所賜的合一就好比一首已經寫好的交響樂.

$^{441}$而這首交響樂是由主耶穌.
他親自在十字架上創作編寫和完成.
交給聖靈在教會裡不斷地指揮.
而我們.
每一個.
不是某一些人.
每一個都有責任.
透過用謙虛,溫柔,忍耐和愛.
就是主耶穌基督的.
這個生命素質.
來一起同心演奏.
有一晚.
這個小提琴家.
很想拉大聲一點.
來表現自己的表現.
不過最終他還是沒有這樣做.
因為他選擇整體的合一和諧.
多過要獨領風俗.
這個就是保持合一.
我們所要付的代價.
我們需要學習不斷地放下自己.
好讓主耶穌基督他自己能夠被看見.
被彰顯出來.
合一的代價其實是非常昂貴.
因為合一是主耶穌基督.
他用寶血重價的買贖.
透過打開他自己的身體.
來建造.
造成一個合一的身體.
買贖合一的代價固然是非常高昂.
你和我都負擔不到.
不過破壞合一.
所要付的代價就更加高昂.
我相信大家不想冒這個風險.
假如我們對價會合一掉以輕心.
我們其實就忽略了.
忘記了.
藐視了.
主耶穌基督在十字架上.
所彰顯的恩典和大愛.

$^{481}$結束了.
差不多了.
三十分鐘了.
在管弦樂團正式演出的那個晚上.
雖然演奏的都是同一班人.
他們都是用同一批樂器.
不過這次他們已經學會了.
看著指揮.
和學會了如何彼此聆聽.
所以當最後的一粒恩符.
被演奏出來.
慢慢在空氣之中消失的時候.
所有現場的觀眾都起立.
來歡呼喝彩熱烈的鼓掌.
為什麼?.
並不是因為某一個人的表現特別出色.
而是純粹因為這首音樂本身.
實在是太過動聽.
這就是神透過保羅給教會.
給你和我的意象.
當我們互相合一的時候.
這個世界就會聽到一個不一樣的聲音.
不是嘈雜,嘈吵,爭吵,爭競的聲音.
而是主耶穌所編寫和諧,合一的聲音.
所以接下來.
我想大家在聽以下這段館宴音樂的時候.
你自己去默想一下.
你可以合上眼.
默想一下你剛才所聽到的訊息.
然後去問自己一個問題.
就是神你帶領我來到恩泉.
來到神的這個大家庭裡面.
究竟你想我怎樣演奏?.
我現在是怎樣演奏你所給我的部分?.
而當我要願意繼續參與.
演奏主耶穌基督你所編寫的和諧,合一的樂章的時候.
究竟我缺乏的是什麼?.
我有沒有缺乏謙虛?.
我有沒有缺乏溫柔?.
我有沒有缺乏忍耐?.

$^{521}$我有沒有缺乏愛?.
請你安靜的.
在神的面前為自己祈禱.
有對夫婦慶祝結婚五十週年.
有人問他們婚姻成功秘訣.
他們回答說.
其實我們早就決定了.
我們不要贏了全家.
輸了那頭家.
我們要維護彼此的關係.
因為這個是更加重要的.
教會的合一也是這樣.
這種合一不是自然.
而且是無緣無故就會發生的.
所以我們要刻意選擇.
選擇和平而不是爭競.
在爭辯誰對誰錯.
選擇迅速去寬恕接納.
而不是彼此的排斥.
拒絕在任何人背後.
來對他說閒言閒語.
選擇用愛.
和用充滿恩典的方式.
來對他說誠實的話.
但是彼此誠實地去正視真相.
和帶出真理.
用來做一次低頭祈禱.
原始謙虛.
溫柔.
忍耐和愛的神.
使我們全然誠信.
幫助我們彼此去竭力.
效法主耶穌基督.
萬口一心地去定義榮耀神.
活出一個同心同恨.
合一和諧的信仰群體.
以致我們能夠與所蒙的呼召相稱.
求你幫助我們定睛地去仰望你.
因為你就是我們的指揮.
幫助我們彼此聆聽.

$^{561}$學習互相的建立.
讓我們的生命成為主.
所成就的和平合一的一個美好的見證.
同心祈禱.
奉耶穌的名.
阿們.
\newpage



\section{創世記 12:1-4}
\label{sec:031720}
\textbf{Rebbot 404 ID}~~~~林家棋宣教師
\newline
\newline
連結: \href{http://www.church.com.hk/acms/content.asp?site=cdc\&op=show\&type=product\&code=031720\&layout=sermon}{\texttt{031720}} ~~~~ 語音日期: 2026-03-01
\newline
\newline
\hyperref[sec:031719]{\small{< < < PREV SERMON < < <}}
~
\hyperref[sec:index_chronic]{\small{[返順時目]}}
~
\hyperref[sec:031721]{\small{> > > NEXT SERMON > > >}}
\newline
\newline
創世記 12:1-4
\newline
\begin{longtable}{cl}
\hline
\hline
章節 & 經文 (和合本修訂版)\\
\hline
12:1 & \begin{tabularx}{0.7\textwidth}{X} 耶和華對亞伯蘭說：「你要離開本地、本族、父家，往我所要指示你的地去。 \end{tabularx} \\ \\ \relax
12:2 & \begin{tabularx}{0.7\textwidth}{X} 我必使你成為大國，我必賜福給你，使你的名為大；你要使別人得福。 \end{tabularx} \\ \\ \relax
12:3 & \begin{tabularx}{0.7\textwidth}{X} 為你祝福的，我必賜福給他；詛咒你的，我必詛咒他。地上的萬族都必因你得福。」 \end{tabularx} \\ \\ \relax
12:4 & \begin{tabularx}{0.7\textwidth}{X} 亞伯蘭就遵照耶和華的吩咐去了；羅得也和他同去。亞伯蘭離開哈蘭的時候年七十五歲。 \end{tabularx} \\ \\
[1ex]
\hline
\hline
\end{longtable}
$^{1}$各位聽眾平安.
解釋一下reboot這個字.
我怕有些人不知道為什麼.
reboot即是重啟.
重啟.
404等下解釋給你聽.
首先去問大家一個問題.
有沒有留意.
我們現在慣常.
我們去朋友家.
或者去餐廳吃飯.
我們永遠第一件事.
做什麼.
什麼什麼.
聽不到.
password 什麼password.
YY password.
很常問的.
除了說聲hello.
有沒有WIFI.
WIFI呢.
我們很常問這個問題.
這樣反映什麼.
其實現在代表.
我們上網已經成為.
我們生活的一部分.
我們上網.
不再只停留在家.
或者.
圖書館.
以前是需要去上網.
除了在家.
也要去圖書館.
上網.
我們不需要在這些固定的空間.
我們才可以做到這件事.
現在只要我們.
手機.
數據電話足夠.
我們其實都可以任何時間.

$^{41}$任何地方我們都可以.
上網去看影片.
年輕人上網打遊戲.
有時候.
我們會.
利用這些科技.
我們都會去上網.
去找一些資料.
但其實我們.
上網都知道.
總會遇上這個曲.
404曲.
404曲代表什麼.
代表網絡錯誤.
原因.
可能你的網址.
串錯字.
多一個字少一個字.
又或者那個網頁已經被人spot.
甚至被人刪除.
你就找不到那個網址.
就變了404這個page.
404這個曲.
背後連著一個英文字.
就是Not Found.
Not Found即是找不到.
404 Not Found.
其實本身這個字.
亦都是一個潮語.
其實就是意思說我找不到.
甚至失蹤了.
可以說人是404 Not Found.
我找不到那個人.
失蹤了.
所以其實.
404 Not Found.
好像一個網址.
好像網址打錯了.
或者我們被刪除那條link.
那個網址出現了一些錯誤.

$^{81}$我們以為.
我們只有上網才會出現.
404這個曲.
但不單止上網.
會出現這個錯誤代碼.
其實.
連猶太人.
他在作為上主子民.
身份上面.
都會出現404這個曲.
同時.
我們可能.
都會出現這個曲.
我們嘗試去從今天的經文.
裡面的文本.
和編修.
這兩個的敘事世界.
我們去思想.
這段經文對我們有什麼提升.
又是時候要求.
大家打開聖經.
打開今天的經文.
創世記十二章一到四節.
邀請大家讀一讀.
創世記十二章一到四節.
大家打開.
創世記十二章一到四節.
一起讀.
創世記十二章一到四節.
一起讀.
創世記十二章一到四節.
一起讀.
創世記十二章一到四節.
一起讀.
創世記十二章一到四節.
一起讀.
創世記十二章一到四節.
一起讀.
創世記十二章一到四節.
一起讀.

$^{121}$創世記十二章一到四節.
一起讀.
創世記十二章一到四節.
一起讀.
創世記十二章一到四節.
一起讀.
創世記十二章一到四節.
一起讀.
創世記十二章一到四節.
一起讀.
創世記十二章一到四節.
一起讀.
創世記十二章一到四節.
一起讀.
創世記十二章一到四節.
一起讀.
創世記十二章一到四節.
一起讀.
創世記十二章一到四節.
一起讀.
創世記十二章一到四節.
一起讀.
創世記十二章一到四節.
一起讀.
創世記十二章一到四節.
一起讀.
創世記十二章一到四節.
一起讀.
創世記十二章一到四節.
一起讀.
我們看見.
上主第一個的命令.
就是吩咐亞伯蘭.
他要為他自己.
去一個地方.
而這個地方不是.
他現在所身處的地方.
也不是他們族群的地方.
更加不是.
他父親塔拉的地方.

$^{161}$更加不是他父親塔拉的地方.
而是去.
上主所要指示的地方.
那即是哪裡.
不知道.
上主這個命令.
其實對亞伯蘭來說.
其實有一定的難處.
第一我們看前少少經文.
就是在十一章三十一節.
亞伯蘭.
他身處的地方.
經文在族譜裡講到.
在迦勒底的吾爾.
我認為吾爾.
這個地方相等於.
古時的紐約.
倫敦和香港.
你想像到這三個城市.
究竟有什麼共通之處.
就是一個比較發達.
已經發展成熟的一個地方.
但是如今.
上主第一個的吩咐.
就是要亞伯蘭.
他去為他自己.
去一個原本.
他要離開原本.
是他出世之後.
所熟悉又發達的地方.
要他走離開.
他出世的地方.
你為到你自己去走.
就好像我.
由出世到結婚之前.
我住在葵涌與荃灣交界.
這個地方.
但我上年結婚之後.
有些弟兄姊妹知道我搬到大埔.
葵涌與荃灣.

$^{201}$由我出世到現在.
我很熟悉那些地方.
吃東西 生活 買什麼.
是非常方便的.
但我搬到大埔的時候.
其實對於我來說.
那一年的時候.
我就是要調節自己.
由一個原本很熟悉的地方.
去一個我不熟悉的地區.
其實我心裡有很多糾結.
其實我和我老婆爭論了很久.
我不是很想搬到大埔.
我想搬回荃灣.
荃灣始終是我的.
但在這些.
搬一個新的地區的時候.
本身就會出現一些心理.
和生活上的適應.
很奇怪.
從一個.
你想想.
我由荃灣葵涌搬到大埔.
都有這些問題.
更何況阿巴蘭呢.
上主吩咐.
他要從一個.
安穩的生活.
移到另一個地方.
而且這個地方不是一天半天.
可以到達的地方.
上主這個吩咐.
他沒有說.
在那裡.
也不知道需要的時間.
要多久.
只知道阿巴蘭.
你要去收拾包袱.
去一個不熟悉.
一個未知的地方.

$^{241}$對於我們來說.
一動不如一靜.
住得好好的.
為什麼要離開呢.
況且要離開.
本地本族富家.
其實反映另一件事情.
就是.
他自己需要主動放棄.
他家族的繼承.
家族產業的繼承權.
包括土地財產.
來到這裡.
我們其實都想像到.
體會到.
其實本身上主這個命令.
對阿巴蘭吩咐這個命令.
其實一點都不容易.
我們接著去看.
或者離開這個行動.
之後其實上主馬上說了.
說了一連串的說話.
就是他為阿巴蘭.
帶來了七重的祝福.
哎呀忘記按了.
帶來了七重的祝福.
第一就是成為大國.
第二就是賜福.
或者我們叫祝福給阿巴蘭.
使阿巴蘭的名為大.
使別人得福.
為阿巴蘭祝福的.
我們上主必會賜福或者祝福給他.
詛咒阿巴蘭的.
上主必會詛咒他.
地上的蠻族.
必因阿巴蘭得福.
短短七句的說話.
就已經出現了.
祝福這個字.

$^{281}$五次.
這些的祝福.
其實都是需要.
阿巴蘭有行動才可以兌現.
對阿巴蘭.
對阿巴蘭而言.
那七重的祝福.
其實更加沒辦法想像.
特別成為大國.
這個祝福.
對於一個75歲.
仍然沒有下一代的人來說.
這個祝福.
可能是一個天馬行空的事情.
我們看第四節.
經文很簡單地.
去交代.
阿巴蘭的回應.
當中其實沒有.
記載到.
阿巴蘭聽完命令之後.
他那種想法和感受.
也沒有任何.
記載關於他的回答.
他只是很直白地.
記載阿巴蘭.
遵照耶和華的.
吩咐去了.
我們可以理解.
阿巴蘭.
這個信心的回應.
同時.
都是他對上主.
信服的展現.
不過.
這種信心或者信服.
和404這個曲.
有什麼關係.
我們就要嘗試從.
另一個世界去理解.

$^{321}$偏修世界.
去理解這段經文另一重的意義.
什麼叫偏修世界.
偏修世界就是創世記.
成書那段時間.
那個年代.
已經不再是阿巴拉罕的時代.
已經是一個再更遠的時代.
我們要理解.
創世記成書的.
那段時間.
大多數學者.
認為那段時間就是.
以色列民被擄的時代.
直到波斯帝國期間.
即是創世記.
完全完成.
整本創世記是被擄時代.
直到波斯年間才有這本書出現.
弟兄姊妹.
我們還要記得.
如果我們這段時間.
其實以色列民已經亡國多年.
並且已經住在.
巴比倫這個地方.
有一段日子.
當時的以色列民.
其實已經在巴比倫這個地方.
落地生根.
有些更加會擁有.
自己的房屋.
甚至在被擄的.
以色列民之中.
有不少都是一些富裕家庭.
和宗教領袖.
雖然亡國被擄.
但是.
他們的日常生活.
其實不是我們一路以來.
覺得很悲慘.

$^{361}$其實不是這樣.
他們可以有房屋.
有自己的後代.
有自己的宗教生活.
不過我們參考.
舊約其他的先知書.
又或者以色列尼希米的記載.
上主定義.
要被擄的以色列群體.
要回到耶路撒冷.
重建聖城.
甚至.
重建聖文.
這個身份.
對於已經成家立室.
已經習慣了.
巴比倫生活的.
以色列民來說.
回到一個.
頹垣敗瓦的耶路撒冷.
做什麼呢.
那裡只剩下一群.
貧窮的同胞.
回去有什麼意思.
況且.
已經有一些住在巴比倫的.
以色列民其實已經不是.
第一代被擄的百姓.
而是一出生就已經住在.
巴比倫生活的第二代百姓.
耶路撒冷來說.
對於這群第二代的百姓.
其實是一個很陌生的地方.
對比現在.
現有的住處.
不是更好嗎.
起碼有房屋.
起碼我過得生活也不錯.
對於這一群.
被擄的以色列民.

$^{401}$他們好像忘記了一件事.
就是他們忘記了.
自己的身份和使命.
若然我們按.
創世記的編修時間.
角度出發.
創世記的讀者群體.
就是一群已經生活在巴比倫.
已經成家立室的.
以色列民.
特別創世記十二章.
一到四節.
今天這段經文作為.
創世記的核心.
編修者.
其實很想從.
阿伯南被呼召的序士.
去喚醒一群.
已經習慣了.
巴比倫生活的以色列百姓.
去喚醒他們.
自身的身份和使命.
當上主.
呼召阿伯南離開.
吾爾這個本地.
本族富家的地方.
伊朗迦南.
就好像當時.
被擄巴比倫的百姓.
他們要離開當地.
要回到耶路撒冷.
一個上主所.
應許的地方.
吾爾和巴比倫.
較相近的地方.
就是兩個充滿.
偶像敬拜的外邦文化.
並且是一個繁華.
富庶的地方.
對於被擄的.

$^{441}$以色列民來說.
巴比倫成為了他們的.
本地本族富家的地方.
回顧.
上主呼召.
阿伯南.
他要為自己去.
或者離開吾爾.
離開原本他的富家.
就是要將自己.
分別出來.
離開當地.
偶像敬拜的文化.
同樣地.
上主呼召以色列百姓.
要離開巴比倫.
回歸耶路撒冷.
我們一起去.
看以賽亞書52章.
11節.
以賽亞書52章.
11節.
好.
打開了的弟兄姊妹.
我想邀請大家一起去讀.
以賽亞書52章第11節.
預備 123.
(授課).
離開.
就是要將自己.
分別出來.
免得勝潔之民.
這個身份被濫污.
我們再.
回到創世記.
上主向阿伯南的.
說話之中.
在一開始提到這段說話.
剛才提到有兩個命令.
第一個是去或離開.

$^{481}$第二個.
就是要洗.
得福.
我們可以理解.
第二個命令.
上主要亞伯拉罕離開.
並且得到祝福.
但同時這個.
祝福.
背後其實賦予了一個.
使命.
就是要成為別人的.
祝福甚至是.
成為萬族的得福.
創世記.
這段敘事.
對於被擄的以色列民.
提醒第一.
他們是先祖亞伯拉罕的.
子孫.
他們是被上主所.
揀選的子民.
是一個分別為聖.
被上主祝福的群體.
第二.
先祖亞伯拉罕被祝福.
同時.
也被賦予使命.
就是要洗.
別人的福.
對於被擄的以色列群體.
同樣都要.
履行這個使命.
偏修者可能.
要喚醒這群.
被擄的以色列民.
他們本身的身份和使命.
他們現在所身處的.
巴比倫的確.
是繁華盛世.

$^{521}$他們還可以有如常的.
宗教活動.
但不要被現在的生活.
所困.
他們不是二等公民.
不要搞錯自己的身份.
他們真正的.
身份是上主.
聖潔的子民.
是要從巴比倫分別出來.
他們這個身份.
其實是被賦予了一個特殊的.
使命.
.
這群被擄之民.
其實要透過他們.
要回歸耶路撒冷.
重建上主的聖殿.
目的就是.
要成為萬邦之光.
我們再看另一段.
耳塞耳書的經文.
60章30節.
我們打開耳塞耳書60章30節.
的經文.
60章30節.
打開了的弟兄姊妹.
也邀請你們一起.
預備1,2,3.
.
其實.
兩段耳塞耳書的經文.
當然它也是對應當時的.
被擄的時代.
它有它的訊息.
它有它的主題.
但同時讓我們看到.
其實被擄的以色列群體.
其實他們都是要延續.
阿巴拉罕的使命.

$^{561}$就是要使到萬族得福.
使命的延續.
必須要回到上主.
應許先祖阿巴拉罕的地方.
即是迦南.
當上主.
吩咐被擄的族群.
要回到耶路撒冷.
就必須像阿巴拉罕一樣.
遵照耶和華的吩咐.
去了.
面對上主的吩咐.
作為上主的子民.
就必須要馴服.
不要停留在巴比倫.
這個本地.
本族.
富家這個地方.
回去.
發揮上主.
子民的身份.
成為列國的光.
弟兄姊妹.
我們雖然.
不是很需要.
在肉身上.
像阿巴拉罕一樣.
我們真的要整個人離開.
本地本族富家.
但是.
上主吩咐我們.
要離開一些.
讓我們基督徒.
身份變成.
404的事情.
例如一些心態.
一些價值觀.
當我們.
受上主的呼召.
其實我們被賦予了基督徒.

$^{601}$這個身份.
而這個身份.
就好像我們每個香港人.
必定會有的東西.
身份證.
我們一定會有.
身份證在身.
這張身份證代表.
我們是香港人.
這個身份.
當然.
我們基督徒的身份明證.
就是我們.
承認基督是主的時候.
我們成為.
神的兒女,是聖徒.
是基督的門徒.
這種.
是上主和我們之間的.
神聖關係.
我們被賦予了.
不過.
我們基督徒的身份.
很多時候和我們身份證一樣.
大部分的時間.
身份證.
不會掛在身後.
大部分時間.
我們把身份證收在哪裡呢?.
收在錢包.
不用的時候就收起來.
什麼時候用呢?.
什麼信用卡.
申請寬頻.
要登記資料.
甚至我們逛逛街.
被截停被查牌的時候.
我們才拿出來.
平時做什麼?.
把身份證放在錢包.

$^{641}$其實可能我們.
基督徒的身份也是這樣.
我們基督徒的身份.
也是這樣.
平時我們基督徒的身份.
就好像變成了404曲.
Not fine.
找不到.
為什麼會出現.
這種情況?原因可能就是.
停留在本地本族.
父家這個地方.
所謂本地本族.
父家不一定.
只是一個地方.
而是一些框住.
我們基督徒身份的事物.
這些事物可以.
來自我們的日常生活.
對於弟兄姊妹來說.
我們都很努力的成為.
一個基督徒.
但我們總會有軟弱的時候.
很常見的情況就是.
我們上班是一個樣子.
上教會是另一個樣子.
當我們每一天.
進到辦公室坐在同一個位置.
我們就自然進入了.
一個職場的模樣.
我們處理工作.
處理職場上的事.
我們自然就用.
我們職場上的規矩.
去處理.
甚至.
可能我們會用一些世俗的手段.
去處理事情.
從職場的角度.
用職場的規矩.

$^{681}$和身份.
去面對我們的世界.
只要不是犯法的事.
可以說是一個正常不過的事情.
處理手法.
然而.
我們基督徒的身份.
好像沒有帶在身上.
好像藏起來了.
好像我們身份證.
我們回到職場裡.
好像放在錢包裡.
這個基督徒身份就放在錢包裡.
藏起來了.
用的時候就拿出來.
我們的基督徒身份.
就會在我們日常生活中.
變成404 Not Fine.
在我們以前的工作環境裡.
公司總會有一些.
基督徒同事.
對於這些基督徒的同事.
不多不少都會想像到.
他們一定會有一些基督徒.
的素質.
而當其時.
我自己的上司.
也是一個基督徒.
而這位基督徒的上司.
有一個很特別的特性.
就是他的脾氣相當乾烈.
有時候.
同事有少少做錯事.
他都毫不客氣.
展現他的脾氣.
當其時在我心目中.
我留有一個疑問.
就是.
他不是基督徒嗎.
為什麼基督徒.

$^{721}$脾氣可以這麼強勁呢.
我不是指.
作為上司面對下屬做錯事.
要隻眼開隻眼閉.
而是可以用一些.
比較理想的方法.
去指出下屬做錯了什麼.
對於做錯事的下屬.
除了我們.
作為上司.
其實我們都不要忘記.
我們都是一個基督徒的上司.
我們不需要去用.
責罵這種方式.
去糾正同事的錯誤.
會不會上司.
可以多點鼓勵.
反而可以幫助到下屬.
或者同事可以做得更好呢.
做學生呢.
說一下做學生.
做學生最常做一件事.
一個很平常的事.
就是抄功課.
在座各位.
我相信個個都有抄功課.
包括我都有抄功課.
我以前讀書.
抄功課背後原因是什麼呢.
除了功課太多.
真的很現實.
多得來來不及.
又或者第三個原因.
功課很無聊不想做.
做都不想做.
有些大英姐妹.
有些年輕人.
可能會有這樣的情況.
不想做這麼無聊的功課.
早點回課室抄就算了.

$^{761}$甚至索性.
抄都懶得抄.
不交 一路拖.
拖到學期尾.
拖到整個學年完結.
我都不交.
學生總會面對這些問題.
但是有沒有想過.
如果我是.
在座各位大英姐妹.
自己審查.
年輕人我比較嚴厲.
你想想 我是一個基督徒.
我正在教會 但我就不交功課.
辭教.
試問.
這是基督徒的.
本份嗎.
以前在學校工作的時候.
有些同事會對.
一些經常不交功課的同學.
打了一個問號.
他們會問.
這個同學這麼有趣.
又說他自己.
上教會.
為什麼經常.
不交功課.
追了一年功課.
首先應該追作文.
追到嘔血.
追作文追到嘔血.
但他又說他自己.
上教會 是基督徒.
為什麼這個同學.
上教會上成這樣.
其實這些說話.
就是正正反映.
別人眼中的基督徒.
大英姐妹.

$^{801}$別人去問.
去問.
你是誰.
我們會怎樣回答.
我們很常會.
馬上回答的就是介紹自己的名字.
我是誰 我是蔡一華.
我是葉靜儀 我是張信義.
同時我為了.
更加鮮明讓自己.
讓對方認識我.
有一些鮮明的身份.
舉出來就是我在哪裡工作.
在哪裡的員工 我在做什麼職位.
學生的.
我就告訴他們.
我在哪間讀書.
東九會基.
會基東九.
我就讀了多少年級.
我讀了什麼科.
我選了什麼科.
大學生更加.
但是當我們.
回答這些問題的時候.
我們好像.
從來沒有想起過.
自己是基督徒這個身份.
我只想起我是誰 我叫什麼名字.
我在哪裡工作 在哪裡讀書.
我是什麼職位.
但是基督徒身份呢.
好像我從來回應別人.
我們從來回應別人的時候想不起身份.
就好像我們把身份證.
有需要的時候.
我才拿出來.
有需要的時候我才拿出來.
不用的時候.
我們放在錢包裡.

$^{841}$待待平安.
其實.
好像昔日的以色列群體都一樣.
他們其實知道他們自己是.
上主子民的身份.
在阿巴拉罕的自書.
不過他們自己已經習慣了.
他們另一個身份.
一個身處在巴比倫.
二等公民的身份.
但是原本.
在聖潔子民的身份.
他所要背負的.
上主所吩咐的命令.
就是延續.
阿巴拉罕要讓地上.
萬族成為光.
吸引列邦的使命.
他們忘記了.
上主就是叫他們.
不要停留.
現在身處的本地.
本族富家.
這個地方不是上主.
所賦予他們的身份.
而是要去或者離開.
前往上主.
所要指示的地方.
一個上主所.
應許的地方.
目的就是要去履行.
這個和這個身份.
匹配的使命.
對我們而言.
我們可能.
真是有需要的時候.
才記得起.
自己是基督徒身份.
甚至我們只是記得.
自己是被上主.

$^{881}$應許有永生這個身份.
而忘記了身份.
背後所賦予的使命.
我們不可能.
和阿巴拉罕一樣.
擁有相同的使命.
就是讓地上萬族得福.
但是我們卻是.
都有作為基督徒的使命.
就是怎樣.
讓別人.
在我們身上看見基督.
讓人透過.
我們的生命去見證.
福音大能.
我們現在打開.
加拉太書3章26-29節.
加拉太書.
3章26-29節.
加拉太書.
3章26-29節.
加拉太書.
3章26-29節.
加拉太書.
3章26-29節.
加拉太書.
3章26-29節.
加拉太書.
3章26-29節.
加拉太書.
3章26-29節.
加拉太書.
3章26-29節.
加拉太書.
3章26-29節.
加拉太書.
3章26-29節.
加拉太書.
3章26-29節.
加拉太書.

$^{921}$3章26-29節.
加拉太書.
3章26-29節.
加拉太書.
3章26-29節.
加拉太書.
3章26-29節.
加拉太書.
3章26-29節.
加拉太書.
3章26-29節.
加拉太書.
3章26-29節.
加拉太書.
3章26-29節.
加拉太書.
3章26-29節.
加拉太書.
3章26-29節.
加拉太書.
3章26-29節.
加拉太書.
3章26-29節.
加拉太書.
3章26-29節.
加拉太書.
3章26-29節.
加拉太書.
3章26-29節.
加拉太書.
3章26-29節.
加拉太書.
3章26-29節.
加拉太書.
3章26-29節.
加拉太書.
3章26-29節.
加拉太書.
3章26-29節.
加拉太書.

$^{961}$3章26-29節.
加拉太書.
3章26-29節.
加拉太書.
3章26-29節.
加拉太書.
3章26-29節.
加拉太書.
3章26-29節.
加拉太書.
3章26-29節.
加拉太書.
3章26-29節.
加拉太書.
3章26-29節.
加拉太書.
3章26-29節.
加拉太書.
3章26-29節.
加拉太書.
3章26-29節.
加拉太書.
3章26-29節.
加拉太書.
3章26-29節.
加拉太書.
3章26-29節.
加拉太書.
3章26-29節.
加拉太書.
3章26-29節.
加拉太書.
3章26-29節.
加拉太書.
3章26-29節.
加拉太書.
3章26-29節.
加拉太書.
3章26-29節.
加拉太書.

$^{1001}$3章26-29節.
加拉太書.
3章26-29節.
加拉太書.
3章26-29節.
加拉太書.
3章26-29節.
加拉太書.
3章26-29節.
加拉太書.
3章26-29節.
加拉太書.
3章26-29節.
加拉太書.
3章26-29節.
加拉太書.
3章26-29節.
加拉太書.
3章26-29節.
加拉太書.
3章26-29節.
加拉太書.
3章26-29節.
加拉太書.
3章26-29節.
加拉太書.
3章26-29節.
加拉太書.
3章26-29節.
加拉太書.
3章26-29節.
加拉太書.
3章26-29節.
加拉太書.
3章26-29節.
加拉太書.
3章26-29節.
加拉太書.
3章26-29節.
加拉太書.

$^{1041}$3章26-29節.
加拉太書.
3章26-29節.
加拉太書.
3章26-29節.
加拉太書.
3章26-29節.
加拉太書.
3章26-29節.
加拉太書.
3章26-29節.
加拉太書.
3章26-29節.
加拉太書.
3章26-29節.
\newpage



\section{箴言 4:10-27}
\label{sec:031721}
\textbf{保守你心勝過保守一切}~~~~周敬敏牧師
\newline
\newline
連結: \href{http://www.church.com.hk/acms/content.asp?site=cdc\&op=show\&type=product\&code=031721\&layout=sermon}{\texttt{031721}} ~~~~ 語音日期: 2025-08-03
\newline
\newline
\hyperref[sec:031720]{\small{< < < PREV SERMON < < <}}
~
\hyperref[sec:index_chronic]{\small{[返順時目]}}
~
\hyperref[sec:031722]{\small{> > > NEXT SERMON > > >}}
\newline
\newline
箴言 4:10-27
\newline
\begin{longtable}{cl}
\hline
\hline
章節 & 經文 (和合本修訂版)\\
\hline
4:10 & \begin{tabularx}{0.7\textwidth}{X} 我兒啊，要聽，要領受我的言語，你就必延年益壽。 \end{tabularx} \\ \\ \relax
4:11 & \begin{tabularx}{0.7\textwidth}{X} 我已指教你走智慧的道，引導你行正直的路。 \end{tabularx} \\ \\ \relax
4:12 & \begin{tabularx}{0.7\textwidth}{X} 你行走，腳步沒有阻礙；你奔跑，也不致跌倒。 \end{tabularx} \\ \\ \relax
4:13 & \begin{tabularx}{0.7\textwidth}{X} 要持定訓誨，不可放鬆；要謹守它，因為它是你的生命。 \end{tabularx} \\ \\ \relax
4:14 & \begin{tabularx}{0.7\textwidth}{X} 不可行惡人的路，不要走壞人的道； \end{tabularx} \\ \\ \relax
4:15 & \begin{tabularx}{0.7\textwidth}{X} 要躲避，不可經過，要轉離而去。 \end{tabularx} \\ \\ \relax
4:16 & \begin{tabularx}{0.7\textwidth}{X} 他們若不行惡，難以成眠，不使人跌倒，就睡臥不安； \end{tabularx} \\ \\ \relax
4:17 & \begin{tabularx}{0.7\textwidth}{X} 因為他們以邪惡當餅吃，以暴力當酒喝。 \end{tabularx} \\ \\ \relax
4:18 & \begin{tabularx}{0.7\textwidth}{X} 但義人的路好像黎明的光，越照越明，直到正午。 \end{tabularx} \\ \\ \relax
4:19 & \begin{tabularx}{0.7\textwidth}{X} 惡人的道幽暗，自己不知因何跌倒。 \end{tabularx} \\ \\ \relax
4:20 & \begin{tabularx}{0.7\textwidth}{X} 我兒啊，要留心聽我的話，側耳聽我的言語， \end{tabularx} \\ \\ \relax
4:21 & \begin{tabularx}{0.7\textwidth}{X} 不可使它們偏離你的眼目，要存記在你心中。 \end{tabularx} \\ \\ \relax
4:22 & \begin{tabularx}{0.7\textwidth}{X} 因為找到它們的，就找到生命，得到全身的醫治。 \end{tabularx} \\ \\ \relax
4:23 & \begin{tabularx}{0.7\textwidth}{X} 你要保守你心，勝過保守一切，因為生命的泉源由心發出。 \end{tabularx} \\ \\ \relax
4:24 & \begin{tabularx}{0.7\textwidth}{X} 要離開歪曲的口，轉離偏邪的嘴唇。 \end{tabularx} \\ \\ \relax
4:25 & \begin{tabularx}{0.7\textwidth}{X} 你的兩眼要向前看，你的雙目直視前方。 \end{tabularx} \\ \\ \relax
4:26 & \begin{tabularx}{0.7\textwidth}{X} 要修平你腳下的路，你一切的道就必穩固。 \end{tabularx} \\ \\ \relax
4:27 & \begin{tabularx}{0.7\textwidth}{X} 不可偏左偏右，你的腳要離開邪惡。 \end{tabularx} \\ \\
[1ex]
\hline
\hline
\end{longtable}
$^{1}$這錄音或錄影的版權是屬於基督教宣導會不覺堂所有.
我們的屬靈傳統之一是以真理教導為本.
歡迎你將此錄音或錄影的連結傳給親友.
懇請注意未經授權請勿轉錄或剪輯本錄音或錄影.
甚願好像《使徒行傳》二十章三十二節下所說.
「這道能建立你們,叫你們和一切誠誠的人同得基業」.
我們都很喜歡用真言的經文來鼓勵你身邊的弟兄姊妹.
或者甚至是一些未信的朋友.
我記得我小學畢業的時候會寫紀念冊.
當時的紀念冊仍然是一張張紙這樣去給.
現在的中學生已經不是寫紀念冊.
他們寫在校服上.
不過當時我寫紀念冊仍然是一張張紙給不同的同學老師.
他們寫完就回來.
你可能馬上還想起一些金句.
我記得其中一位老師將他的文字給我.
我一看他寫一句話說得合宜.
「就如金蘋果在銀網子裡」.
然後他還寫了那一張紙.
「真言二十五章十一節」.
原來這位老師教了我這麼久原來他是基督徒.
不過我們很喜歡用真言的說話.
可能剛才你看第四章二十三節的時候.
你要保守內心勝過保守一切.
都很可能是我們的金句.
不過我想說的是我們用真言來鼓勵人.
其實這不是一件錯事.
你不要誤會我.
不過當我們這樣用的時候.
是否等於我們完全明白當中的意思.
即使你明白當中的意思.
又是否等於你會切實進行呢.
今天我們來到這個時間.
大約是這個季題.
我們讀真言已經去到一半的部分.
當我們讀真言的時候.
當我們留心裡面智慧者對我們的說話.
而我們知道當我們聽完.
要讓我們的生命成長的時候.
有兩方面是很重要的.

$^{41}$第一你要聽.
不斷重複的「我而我而」你要聽.
第二就是你要行出來.
今天我們讀第四章的主旨主題.
其實都是提及我們在生命的真道上.
怎樣持續地去成長.
我們今天由第十節開始看.
不過如果由再看早一點.
即是由第一到第九節.
總括而言是父親對孩子的提醒.
當中第三節你看到.
他說當我即是父親.
在父親面前還是小孩.
是母親獨一嬌兒的時候.
他教導我.
說你的心要持守我的話.
遵守我的命令.
你就會存活.
這位父親提到他的父親.
即是爺爺.
怎樣勉勵他.
而現在這個父親也將同一番說話.
向他的孩子表達.
所以其實是有一種.
三代都一起領受智慧的意思.
弟兄姊妹.
我們信主.
在信仰的這條路上走.
其實你和我都知道.
我們都朝著這條.
最終而且是永恆的目標.
來去成長和前進.
而在這條信仰的道路裡面.
在這條路的中間.
我們會面對不同的事情.
你一路走.
你會一路面對.
作出不同的選擇.
而當我們作出選擇的時候.
我們就需要智慧.

$^{81}$在經文裡面.
父親告訴孩子.
因為他正正經歷過.
他曾經經歷過.
所以他提醒孩子.
要作出正確的選擇.
弟兄姊妹.
當我們從第十節開始來看的時候.
我們想不只是聽.
你知道.
而我們在實踐當中.
我們怎樣有正確的選擇.
今天從這段經文裡面.
有四個行動.
四個行動幫助你和我.
我們一起來提醒.
首先我們看第十和第十一節.
經文說.
我而要聽.
要領受我的言語.
你就必然連憶受.
我已指教你走智慧的道.
引導你走正直的路.
我們很多時候都會問.
為什麼有人會選擇.
走一條惡或者罪的路呢.
你想一想.
一看一看.
明明很清楚.
這條是正.
這條是邪.
這條是對.
這條是錯.
你總是覺得不知道為什麼.
為什麼有人會走錯那條路.
但其實當你有這樣想法的時候.
很多時候我們都是.
回過頭來想.
你回想回顧的時候.
你就會想.

$^{121}$為什麼我看不清楚.
為什麼他會這樣呢.
如果真的分得這麼清楚的時候.
如果那一刻真的分得這麼清楚的時候.
我想大多數人都知道.
哪一條是正.
哪一條是錯.
但當我們正在經歷那件事.
我們在那個期間.
在那個空間的時候.
正路和惡路.
很多時候都在同一個空間.
同一個時空裡面.
而當我們要作出選擇的時候.
我們就會面對很大的掙扎.
所以當你十字一開始的時候.
經文說要領受我的言語.
你就不言年益壽.
當然從字面解言年益壽.
那是指生命的年數會賜給你.
但如果我們連同往後的經文一起去看的時候.
就不是純粹提及生命的年數有多少.
而是說的是生命的質素.
和我們有沒有活出一個豐盛的生命.
所以當我們領受智慧的時候.
這個智慧幫助我們選擇.
走一條正直的路.
而當你選擇走正直的路的時候.
還有多一樣東西.
在第十三節經文提到.
要持定分魁.
不可放鬆.
要緊守,因為它是你的生命.
你明知每一個步驟都要step by step.
做完一步,做下一步,再做下一步.
這件事才叫做complete.
才叫做完.
如果你skip了一個step.
你可能會忽略了,跳了.
或者是沒有了一些很重要的東西.

$^{161}$不過你可能都會試過.
我發覺跳step原來是沒有問題的.
我們每一個人都上過數學課.
你知道跳step.
其實一兩次都不是不行的.
如果我們存在一種僥倖的心態.
我們不知道我們何時會面對一些錯.
或者不好的結果.
經文說持定,不可放鬆,緊守.
就是我們不要存在一種僥倖的心態.
以為OK,得過且過.
最後一節說因為它是你的生命.
我們知道生命重要.
所以我們不敢放鬆.
我不知道你的生命.
今天是一個怎樣的光景.
你圍到你的生命.
上帝很保守你.
又說即使我現在面對一個破碎的生命.
但我仍然努力靠著上帝的幫助.
逐步逐步建立起來.
父親在這裡提醒兒子.
你不要覺得你會有一種僥倖.
其實當我們回想我們信主的年日.
可能剛受洗的弟兄姊妹.
在靈裡面是一種很高峰.
我要上主一學去勤讀聖經.
我今年都不是錯.
信主信了很多年的.
你要知道我們仍然朝著那個生命之道.
一直成長的時候.
無論信主的年日長短.
我們都需要建立這份合神心意的生命質素.
而這份建立的過程是step by step.
你不要存著一份僥倖的心.
所以第一個行動.
求神幫助你去建立.
第二方面.
我們看第十四和第十五節.
第十四和第十五節我們一起讀.

$^{201}$不可行惡人的路.
不要走壞人的道.
要躲避不可經過.
要轉來而去.
當我們讀到這兩節經文的時候.
你很清楚知道裡面在說罪.
當一個人跌倒犯罪的時候.
往往都是出於一個很簡單的念頭.
即是不會第一步就跳下去.
不會第一步就犯一個很大的罪.
而是一種循序漸進式.
慢慢一步一步陷入裡面.
而聖經在這裡很清楚提醒我們.
當你意識到那個是罪惡之前.
就已經要讓自己避免走進那條罪惡的道路.
而這條罪惡的道路沒有分大小.
是罪就罪.
不會有中間的灰色地帶.
我們要提醒我們自己.
因為往往當我們走錯方向的時候.
那一刻我們都未必立刻意識到後果.
在這裡的弟兄姊妹.
我們未必個個都會開車.
不過你會坐車.
你坐車的時候.
你知道每一條路都有一個速度的限制.
可能你坐小巴.
你前面看到那個時速.
60,70超速它會響.
又或者每一條路的時速限制.
50,60,70,80.
諸如此類.
我就記得我第一次被人抄牌.
Speeding 捉快車.
其實都是唯一一次.
坦白說.
都是那一次.
很多年前我還在外國.
在教會回家的那條路上.
我記得那條路走60.

$^{241}$之前還沒有車的時候.
坐別人的車.
每個人都是這樣走的.
你知道不是走60.
一直走每個人都沒事.
到自己開車的時候.
試過一次,兩次.
都是沒事.
那次完了師班練習.
開車回家.
和一個弟兄一起載他回家.
一直走.
那條路沒有什麼車.
走60時速.
即是它的限制.
走著走著.
外國的路很多都是黑.
就是黑漆漆的.
走著走著路邊有一盞燈.
亮了.
隔壁的弟兄就.
Oops.
我就不告訴你.
我超速多少.
總之.
超速的是會被人捉的.
Speeding.
那次就被人這樣抄牌.
你犯錯之後.
我仍然很記得.
那個警察說.
你可以走了.
我的手震到.
然後腳震到.
那一下比彈風琴還要害怕.
我記得當時離開的時候.
我想應該全程都走四五十路.
你知道.
那一刻當你真的出事.
但其實你想想.

$^{281}$你已經試過很多次.
但你未必意識到那個後果.
我們說.
上得山多終如虎.
從意念到你真正犯錯.
或者跌倒.
中間都有一個空間.
或者一個時間.
你說不要緊.
那種僥倖的心.
沒有被人捉.
你就一步一步踏入去.
在香港有很多行山徑.
這個行山的警告牌.
我自己不是行很多山.
這個警告牌應該是我見到.
這個內容最大的其中一個牌.
他說懸崖危險.
切勿前進.
但當他動起這個牌的時候.
他不是在懸崖邊就動起給你看.
你見到那個牌後面.
還有一段距離.
才是那個懸崖.
其實他很早已經提醒你.
這裡已經是那個極限.
你要回頭.
他不會在崖邊才告示.
你跨越了就馬上掉下去.
目的是讓人有足夠的空間.
回頭.
所以當經民說不可行的時候.
經民繼續說.
你都要躲避.
你要轉來.
不論你在工作這件事上.
你可能會說因為方便的緣故.
我用了偏路.
用了一些偏方.
又或者你覺得我在行為上.

$^{321}$我不會影響到人.
甚至你可能會說.
當我沒有人的時候.
我在網絡上做了一些.
我以為沒有人知道的舉動.
在道德上越界.
做出婚姻關係裡不容許的情況.
弟兄姊妹.
這些都是一步一步陷入去.
你更加不要存在一種藉口.
那個藉口說.
不要緊.
我靠著神的恩典.
我為了弟兄姊妹的好處.
我試試看.
都是藉口.
剛才說.
無論信主的年日多或者短.
我們繼續讓自己生命成長.
無論信主多少年.
你不要說我能夠撐得住.
不要以為我有很委身的事奉.
我好像有一個保護罩.
我都信過我自己.
所以弟兄姊妹好好留心.
你身邊的弟兄姊妹.
導師,牧長給你的提醒.
神在很多方面.
其實祂都提醒你.
叫我們躲避,逃離.
不要進入罪惡的路.
要轉來.
這是第二個行動.
接著在第20至27節.
經文分別帶給我們兩個最後的行動.
第20節.
這幾節裡.
其實有幾個我們身體上的器官.
包括耳朵.
20節留心聽我的話.

$^{361}$眼不可使它們偏離你的眼目.
你眼要向前看.
雙目要直視前方.
口要離開歪曲的口.
轉來偏斜的嘴唇.
腳要收平你腳下的路.
你的道就必穩固.
然後不可以偏左偏右.
你的腳要離開邪惡.
你知道耳也好眼也好口.
這些都是我們很重要的感官.
而腳是一個行動出來的表現.
身體上所有的部份都需要回應和行動.
而你看什麼想什麼.
其實就影響你做出來的事情.
你和我都會上網.
你和我都會有不同的興趣.
你和我都會間中去找一些對你自己有.
你覺得很想繼續去看的東西.
通常你這樣想的時候.
你的社交媒體.
或者電腦就會不斷出現相關的資料.
早一陣子我有一位肢體.
他去了一個我從來沒去過的地方旅行.
因為他回來香港就有點阻滯.
我只是好奇.
查一下他回來的那班飛機.
究竟延遲了多少.
那間航空公司我沒坐過.
那家我沒去過.
那個地方我未必有興趣去玩.
我只是搜尋過兩次.
之後那兩天.
我的社交媒體全部都是關於那個國家的旅遊資訊.
都是講航空公司的宣傳.
我們看什麼.
我們的心就會.
焦點就會去了那裡.
所以你看什麼想什麼.
就很容易成為你生活裡的一部分.

$^{401}$不是上網.
可能你那段時間那期.
你經常都會聽某一個系列的詩歌.
當你一直聽的時候.
你很自然就會.
走一走街走兩步.
就會突然間這樣唱起.
我經常提詩班也是.
他們都會練歌的時候.
我猜他們會間中經常聽.
可能都是一樣.
一直聽得多.
可能不自然就會自然.
就會哼起.
所以是那種.
我去想我去聽.
其實就會引發到我其他.
我身體上面.
都會做了同一樣的事情.
當然在生活裡面.
我們都有很多不同的興趣.
例如行山運動攝影.
這些都是健康的.
不過你記得剛才我說.
在開始我們說的時候.
在我們這條信念的道路中間.
我們會面對不同的事情.
我們在中間.
或者會作出一些不同的選擇.
這些不同的選擇.
就讓我們究竟走在正.
還是走在罪惡的道路上.
這幾節的經文.
無論耳口眼腳.
提醒我們當我們聽見神的話.
而且領受的時候.
我們的心要存的.
要記得.
左耳入右耳出.
我們很簡單.

$^{441}$很容易去說.
但是如果當你知道.
這是你的生命的時候.
你就會好好地去感受.
所以你要留心.
這幾件事會不會出現偏差.
我不知道對於你來說.
這幾個感官.
可能你還有一個.
你覺得自己是弱一點的.
在各方面.
我的嘴巴怎麼說都不聽.
怎麼提自己都搞不定.
我的耳朵總是聽一點不聽一點.
或者我的眼睛.
我總是控制不了.
看一些不應該看的東西.
求神幫助你.
不會跌倒.
當這個部分.
經文繼續說.
要專心尋求的時候.
專一其實是提醒我們.
我們持定.
走一個正確的路.
並且你要時刻都要謹慎.
你留心.
不要偏離.
要並且遠離罪惡.
當然這個不是靠我們自己.
剛才你說.
我們自己都沒什麼信心.
可以搞得定.
但當你領受智慧的話之後.
求神幫助你.
一步一步你會明白.
你會作出.
對而有正確的選擇.
你心水清的時候.
其實還有多一個器官是還沒說的.

$^{481}$剛才那節金句.
我們一起讀.
你要保守你心.
勝過保守一切.
因為生命的泉源.
由心發出.
等一下我們會讀多兩次.
因為這節是今個月的金句.
我們知道心的重要性.
沒有了它.
身體的其他功能都不能夠運作.
你知道.
血液通過心臟.
去到身體的各個部位.
所以它是最重要.
你又想起以前.
古時打仗的時候.
通常他們會用盔甲去遮住心口.
今天他們可能會穿避彈衣.
那個避彈衣也是遮住心口.
又或者你記得勝經.
當保羅說.
神所賜的全副軍裝.
遮胸.
是那塊護心鏡.
都是一個很堅硬的東西.
心在經文裡面.
是超越了純粹一個器官的功能.
在當時的文化裡面.
心代表人的思想.
代表有情感有意志.
這些都是很核心的要素.
也都代表著生命的本質.
所以弟兄姊妹.
當你一路聽的時候.
你知道有很多事情我們需要謹守.
但是首先你要保守你的心.
因為心是所有事情的源頭.
我們所有的想法.
價值觀的來源.

$^{521}$英文聖經譯這一段.
保守這裡.
他譯作guard your heart.
guard.
即是你有保護.
有守護的意思.
之前那幾節.
剛才我們說.
耳,眼,口,腳.
都是從心而發的意思.
所以我們要從源頭開始處理.
我們怎樣檢視.
我們自己的情感和意念.
你怎樣做.
才能夠保守自己.
不被罪惡的意念.
勝過自己.
保守我們的心.
勝過保守一切.
你又會否誤會經文的意思.
不是說其他不重要.
或者其他不需要保守.
其他都需要.
不過更加重要的是.
要保守你的心.
因為生命的全源由心發出.
當我們領受智慧的時候.
我們的生命就被神更新和轉化.
所以是必須要有神的話在期間.
讓我們的生命流露出.
那份屬神的樣式.
相反如果你記得.
剛才我們提過的經文.
如果你沒有持定分悔.
你沒有躲避罪惡.
你沒有轉離.
你的心思就逐步逐步被罪籠罩著.
真言形容.
那是一個愚昧不義的生命質素.
所以更加留心.

$^{561}$我們在日常接觸.
我們聽看的事.
第四個行動.
當我們要去守護.
去保守的時候.
其實都要讓神引導你.
想起你自己的生命.
留心我們做每一件事的動機.
出發點是甚麼.
稍後主席會帶領我們.
唱今天的回應詩.
稍後那首回應詩叫做「眾生之神」.
詩歌有好幾節.
有六節的歌詞.
不過稍後我們唱的這首詩歌.
是一個形容一個人向神的祈禱.
幾節的歌詞都從不同的方面.
求神更新我們的心.
從意念到行動.
我們都立志按著神的引導而行.
即使在你的生命裡面.
或許正面對暴風.
或者正面對困惑.
但我們都願意聽得見神對我們的聲音.
而且我們能夠分辨出那個是神的聲音.
你記得四個行動.
建立生命的質素.
逃離不要走罪惡的路.
器官那幾樣東西你要專心不偏左右.
持定我們成長的目標.
保守求神幫助我們檢視我們的生命.
而我們留心我們生命的動機.
我們一起祈禱.
天父我們多謝你.
讓我們在領受你話語的時候.
叫我們的生命能夠得著.
被你的光照和帶領.
我們不敢輕看.
我們不敢自誇.
但求你賜給我們這份謙和的心志.

$^{601}$我們都時刻想起.
我們生命的每一步每一刻.
究竟我們在做什麼.
作為個人作為一個群體.
主要求你時刻都帶領著我們.
可能我們中間有弟兄姊妹正面對.
在生命上很多不容易的地方.
甚至可能在罪裡面被罪纏繞.
但當我們轉眼看著你的時候.
主要求你一步一步幫助我們.
你的同在和同行.
我們知道你一直都沒有離開.
多謝你.
願你垂聽我們的祈禱.
同心仰望.
奉耶穌基督的聖名.
阿們.
\newpage



\section{約翰福音 3:1-17}
\label{sec:031722}
\textbf{上帝愛世界}~~~~高國雄牧師
\newline
\newline
連結: \href{http://www.church.com.hk/acms/content.asp?site=cdc\&op=show\&type=product\&code=031722\&layout=sermon}{\texttt{031722}} ~~~~ 語音日期: 2026-03-01
\newline
\newline
\hyperref[sec:031721]{\small{< < < PREV SERMON < < <}}
~
\hyperref[sec:index_chronic]{\small{[返順時目]}}
~
\hyperref[sec:031723]{\small{> > > NEXT SERMON > > >}}
\newline
\newline
約翰福音 3:1-17
\newline
\begin{longtable}{cl}
\hline
\hline
章節 & 經文 (和合本修訂版)\\
\hline
3:1 & \begin{tabularx}{0.7\textwidth}{X} 有一個法利賽人，名叫尼哥德慕，是猶太人的官。 \end{tabularx} \\ \\ \relax
3:2 & \begin{tabularx}{0.7\textwidth}{X} 這人夜裡來見耶穌，對他說：「拉比，我們知道你是由神那裡來作老師的；因為你所行的神蹟，若沒有神同在，無人能行。」 \end{tabularx} \\ \\ \relax
3:3 & \begin{tabularx}{0.7\textwidth}{X} 耶穌回答他說：「我實實在在地告訴你，人若不重生，就不能見神的國。」 \end{tabularx} \\ \\ \relax
3:4 & \begin{tabularx}{0.7\textwidth}{X} 尼哥德慕對他說：「人已經老了，如何能重生呢？豈能再進母腹生出來嗎？」 \end{tabularx} \\ \\ \relax
3:5 & \begin{tabularx}{0.7\textwidth}{X} 耶穌回答：「我實實在在地告訴你，人若不是從水和聖靈生的，就不能進神的國。 \end{tabularx} \\ \\ \relax
3:6 & \begin{tabularx}{0.7\textwidth}{X} 從肉身生的就是肉身；從靈生的就是靈。 \end{tabularx} \\ \\ \relax
3:7 & \begin{tabularx}{0.7\textwidth}{X} 我說『你們必須重生』，你不要驚訝。 \end{tabularx} \\ \\ \relax
3:8 & \begin{tabularx}{0.7\textwidth}{X} 風隨著意思吹，你聽見風的聲音，卻不知道是從哪裡來，往哪裡去；凡從聖靈生的也是如此。」 \end{tabularx} \\ \\ \relax
3:9 & \begin{tabularx}{0.7\textwidth}{X} 尼哥德慕問他：「怎麼能有這些事呢？」 \end{tabularx} \\ \\ \relax
3:10 & \begin{tabularx}{0.7\textwidth}{X} 耶穌回答，對他說：「你是以色列人的老師，還不明白這些事嗎？ \end{tabularx} \\ \\ \relax
3:11 & \begin{tabularx}{0.7\textwidth}{X} 我實實在在地告訴你，我們所說的是我們知道的，我們所見證的是我們見過的，你們卻不領受我們的見證。 \end{tabularx} \\ \\ \relax
3:12 & \begin{tabularx}{0.7\textwidth}{X} 我對你們說地上的事，你們尚且不信，若對你們說天上的事，如何能信呢？ \end{tabularx} \\ \\ \relax
3:13 & \begin{tabularx}{0.7\textwidth}{X} 除了從天降下的人子，沒有人升過天。 \end{tabularx} \\ \\ \relax
3:14 & \begin{tabularx}{0.7\textwidth}{X} 摩西在曠野怎樣舉蛇，人子也必須照樣被舉起來， \end{tabularx} \\ \\ \relax
3:15 & \begin{tabularx}{0.7\textwidth}{X} 要使一切信他的人都得永生。 \end{tabularx} \\ \\ \relax
3:16 & \begin{tabularx}{0.7\textwidth}{X} 「神愛世人，甚至將他獨一的兒子賜給他們，叫一切信他的人不致滅亡，反得永生。 \end{tabularx} \\ \\ \relax
3:17 & \begin{tabularx}{0.7\textwidth}{X} 因為神差他的兒子到世上來，不是要定世人的罪，而是要使世人因他得救。 \end{tabularx} \\ \\
[1ex]
\hline
\hline
\end{longtable}
$^{1}$提醒姐妹平安.
讓我們再一次同心祈禱.
主讓我們多謝你.
讓我們經歷你的豐富.
在今天這個聖餐崇拜當中 我們一齊.
與你連結.
願你再一次進入我們心靈當中 成為我們提醒鼓勵.
看我們的眼睛 讓我們看見你的榮美.
看我們的耳朵 讓我們專心去聆聽.
更要以我們的手腳行動 去回應你的提醒和大愛.
有言報人口中的言語 弟兄姊妹請在心思意念.
蒙聖靈的啟迪蒙主的月立.
禱告奉耶穌基督的名.
阿們.
剛剛在福音書裡面 你聽到.
《約翰福音》三章.
相信很多弟兄姊妹朋友都熟悉的經文.
基督教一個很重要的核心信息.
是愛.
上帝愛世人.
這個描述就正正是今天.
福音書《約翰福音》三章十六節.
這一節很重要的金句.
我記得我信耶穌的開始是.
由中學階段在學校裡面認識信仰.
當我去到這間學校的時候 它是教會學校.
當時我沒有信仰背景 但要上聖經堂.
聖經堂要做的就是背金句.
要默書 要默金句.
除了我記得我經常都不合格 因為漏了一點點料料之外.
這一節金句.
很記得是第一堂老師就說.
整個基督教的核心就是在講上帝的愛.
而上帝愛世人就是這樣去展現.
這個也是我們今天和大家思考的重要題目.
因為上帝愛世人這個描述.
沒錯是基督教最流傳.
最廣泛 最被接受的一個描述.
就算你不是返教會 還不是信耶穌的朋友也好.
你都知道聖經裡面講的就是上帝愛人.

$^{41}$但這個愛 可能很多時候都會被我們誤解.
可能是我們一些習慣用批判思考去理解的人.
我們就會覺得這個愛往往就顯於過於感性.
甚至在這個充滿苦難.
荒謬和不公的世界面前.
就會容易顯得軟弱無力.
或者這句話 上帝愛世人.
已經成為了十分空泛的宗教的語言.
今天是大齋祈.
大齋祈是一個曠野的季節.
祂要求我們剝離了表面的繁華.
直視生命的幽暗和脆弱.
在今天大齋祈第二主日.
我們跟從耶穌的腳步.
一步一步走向耶路撒冷.
這是一個走向死亡的旅程.
當我們在這個背景之下.
在今天的經課當中我們談論上帝愛世界.
我們要談論的不是一種溫馨的情感安慰.
而是一場神學上的冒險和更新.
《約翰福音》三章一開始的時候.
講到有一位法利賽人名叫尼戈德姆.
是猶太人的官.
舊譯為尼戈底姆.
可能因為姆字 怕別人誤會他是女性.
所以就改了尼戈德姆這個翻譯.
今天經文始於一個深夜.
這位尼戈德姆是當時的知識分子.
他是法利賽人.
是拉比 即是一位老師.
他擁有嚴謹的邏輯和神學架構.
是一位極高學術地位和宗教權力的知識分子.
他在黑暗裡面走向了耶穌.
在現今我們身處的高速生活裡面.
很多人可能都同樣地處於某一種深夜黑暗當中.
對大時代的無力感.
對前途的迷惘.
以及對有知識體系崩潰的焦慮.
這個人在第二節說.
他在他的夜晚裡面就來到見耶穌.

$^{81}$他就跟耶穌說.
拉比 我們知道你是由上帝那裡來的師父.
因為你所做的神蹟.
若果沒有上帝同在.
沒有人知道.
尼戈德姆的開場白第一句就是.
拉比 即是老師 老師我們知道.
一開始他的描述就是有些東西我們知道.
但同樣地這是他的困境.
因為他試圖用他擁有的知識框架.
我們知道耶穌你是由上帝來做老師的.
這個框架去定義耶穌.
他追求一套解釋世界的理論.
試圖用嚴謹的因果邏輯等等去定義耶穌.
對知識分子來說.
理解就是掌控世界的方式.
我們同樣都是帶著我們知道的自信來到.
有些東西我們知道.
不過這些東西可能讓我們.
陷入了生命裡面更多的困境.
因為耶穌的回答就是這樣說的.
耶穌的回答是我實實在在告訴你.
人若不從生就不能見上帝的國.
耶穌隨即打破了他的框架.
耶穌說人若不從生就不能見上帝的國.
這個從生翻譯是由上而下的意思.
由上而下的生.
人要這樣才能見到上帝的國.
從生不是知識的增長.
耶穌在這裡用一個框架去描述.
從生是一種轉向.
耶穌指出認識上帝不是僅僅靠觀察或分析.
甚至學術研究.
而是靠著生命徹底的重整.
也讓我們明白.
生命追求的不是那些當時的宗教分子.
他們所追求的做師傅.
行大能騎士做偉大的事.
耶穌在這裡轉另一個框架去跟尼哥德默描述.
他說要的是要見到.

$^{121}$又或者是要進入上帝的國.
這些比其他一切都重要.
我們記得瑪利福音六章三十三節.
我們都熟悉的經文.
耶穌提醒我們要先求祂的國和祂的義.
這些東西.
即是其他憂慮,關心,擔心,需求的東西.
上帝就會加給我們.
所以首要的一種轉念.
是一種很不同的眼光.
尼哥德默就繼續這樣問.
如何從生?.
人老了,難道再從沒有福出來嗎?.
耶穌這樣回答.
我實在告訴你.
人若不是從水和靈生的.
就不能夠進上帝的國.
從肉生生的就是肉生.
從靈生的就是靈.
我說你們必須從生.
你們不要以為稀奇.
風隨著意思去吹.
你聽見風的響聲.
卻不曉得它從何而來.
往何而去.
凡是從聖靈生的.
都是這樣.
聖經裡面提到尼哥德默.
繼續用他熟練熟悉的邏輯去問.
如何從生?.
還要說從上而來的生.
難道回到媽媽的肚子裡嗎?.
耶穌的回應說.
從生的途徑是透過水和聖靈.
是唯一的途徑.
這個描述當日可能很多人不明白.
今天我們會明白多些.
因為我們透過水禮加入教會.
透過聖靈的洗禮.
讓我們生命不再一樣.

$^{161}$但是這個描述.
很明顯是超越當時尼哥德默的想像.
又或者超越我們的想像.
可能我們以為信仰應該透過理解和邏輯.
但是聖經裡面說的.
不是理解和邏輯不重要.
而是有更重要的眼光去看.
就是上帝的恩典和應許.
耶穌繼續用風去比喻聖靈.
風看不到.
但是看到祂的威力.
香港人很習慣.
我們經常打風就知道.
一打風就看到祂的威力.
在凡事追求精準計劃.
KPI的香港.
耶穌提醒我們.
上帝的愛和工作.
往往是在我們專業邏輯.
不可測的地方出現.
愛世界.
首先是要承認上帝的主權.
高於我們理性的預測.
因為耶穌繼續這樣說.
我實實在在告訴你們.
你們所說的我們知道.
我們所見證的我們見過.
你們卻不接受我的見證.
我對你們說地上的事.
你們尚且不信.
如果說天上的事怎麼信呢.
除了從天降下.
營救在天的人子.
沒有人升過天.
耶穌在這裡去責備尼哥德姆.
說我對你們所說地上的事.
你們都不信.
反映了很多時候我們的通病.
我們很容易透過過去.
我們一些理解.

$^{201}$一些習慣.
甚至在地上我們建構的邏輯.
政治等等結構.
我們的理念.
以至我們失去了原有.
信仰裡面有的.
其實不是應該這些邏輯.
而是一種對信仰的想像力.
對上帝應許那份信靠.
耶穌繼續有一個例子.
第十四節.
耶穌說摩西在曠野怎麼舉蛇.
人子也都必照樣被舉起來.
叫一切信他的人都得榮生.
耶穌引用舊約民數記.
銅蛇的典故.
當時以色列人在曠野裡面.
因為怨毒而死亡.
仰望銅蛇就能夠得到活.
我們今日的生活.
我不知道你覺得生活得怎樣.
或者我們都覺得.
我們好像在曠野當中.
充滿了怨恨甚至毒素.
耶穌強調.
上帝拯救世界的方法.
不是透過一個我們想像的.
完美的社會理論.
而是透過一個被舉起的受苦者.
這個對純上權力.
或者效率的現代人.
知識分子來說等等.
我們會覺得挺大的冒險和挑戰.
原來需要的不是甚麼理論.
甚麼的建構.
而是要接受這位被舉起的受苦者.
上帝愛世人.
甚至將他的獨生子賜給他們.
叫一切信他的不致滅亡.
反得永生.

$^{241}$這次金句中的金句.
剛才我小時候已經開始去背.
到今天仍然是我最喜歡的金句.
很直接了當地告訴我們.
上帝愛世人.
愛到甚至將他的獨生子賜給我們.
以致讓一切信的人就不再滅亡.
能夠得到永生.
很呃要,很精準,很好的描述.
上帝愛世人裡面的世人或世界.
這個字眼.
在羊羊荒裡面.
其實世界,世人這個字眼.
往往有少少負面的表達.
特別是指那些不認識上帝.
甚至敵對上帝的墮落體系.
我們熟悉的可能是魔鬼,撒旦.
又或者一些不認識上帝的人.
就是世人或世界這個字的描述.
上帝的愛不是一個他愛完美的世界.
而是一個反叛叛逆的世界.
愛不是一個可以被論證的命題.
而是一件上帝主動介入的事件.
對我們來說.
開始承認我不知道.
才是重新的開始.
上帝愛世界是愛那個真實卻是殘缺的世界.
而不是我們腦海裡面想那個完美.
很應該完美的理想國.
不是,愛不是這樣.
上帝所愛的不是抽象的人類.
上帝愛世人不是一句口號.
一句呼籲.
一個slogan.
告訴你上帝是愛人類的.
而是上帝實實在在的.
對每一個具體有罪在體系體制當中掙扎.
甚至可能不認識.
拒絕甚至抵擋上帝的每一個.
這個就是上帝愛的世人.

$^{281}$上帝的愛更加不是溺愛犯愛.
上帝的愛是透過賜下獨生子.
這一場代價極重的介入.
將世界從混亂裡面拯救出來.
聖經說因為上帝猜他的兒子降世.
不是要定世人的罪.
而是要世人因他得救.
這句經文對現在非常非常重要.
在我們世界裡面充滿著很多的指責.
投訴的文化.
很多審判的氛圍的社會裡面.
這段經文的宣告.
上帝猜顯他的獨生子來到世界.
原意不是要定世人的罪.
沒錯.
這個愛本身應該大有審判性.
因為當光進入這個世界.
黑暗拒絕光.
這個拒絕當然應該受審判.
但是上帝選擇的不是審判.
上帝選擇的是用愛.
耶穌基督道成肉身的愛.
不是說沒有了審判.
有一天我們回到天家.
我們都要被審判.
但是上帝不是直接因為我們的罪過去審判我們.
而是給了一個出路.
他親自用愛介入.
耶穌基督道成肉身成為人.
在我們當中去展現什麼是愛.
我們接受就能夠去面對這個審判.
如果上帝對這個這麼破碎的世界.
仍然都存在憐憫.
而不是單純去審判的話.
我們這些自己以為自己有能力.
知識,洞見力等等的人.
是不是應該放下那種旁觀者的審判感.
轉而學習去參與上帝的救贖.
復修,復和的愛呢?.
他們就照猶太人笨重的規矩.

$^{321}$把耶穌的身體用細麻布加上香料裹好.
在《約翰福音》的結尾裡.
尼哥德姆在日光日白之下.
走出來安葬耶穌.
《約翰福音》記載很有趣.
第三章我們今天經文裡.
尼哥德姆在深夜.
在沒有人見到的時候.
偷偷地去找耶穌.
很想去解決他生命裡.
追求人生價值的問題.
去到最後的結果的時候.
當耶穌受死.
耶穌要埋葬的時候.
他就在白日.
在眾人面前為耶穌安葬.
這位知識份子.
他由旁觀到毀身.
《約翰福音》就是記載這個過程.
上帝愛世人.
這份愛邀請我們用行動回應.
承認我們的專業.
我們的理性.
可能是沒有辦法去解決.
生命裡面的荒涼.
唯有從上而來給我們的啟示.
在各自的崗位上.
讓我們學習.
效法基督.
道成肉身.
我們如何介入.
如何去守護.
如何去看守.
如何去幫助身邊每一個.
很具體的世人.
那些可能是敵對信仰.
甚至很多人看不起的世人.
在大齋祈這個荒涼當中.
不再只是看見社會的崩壞.
看見很多負面情緒的狀況.

$^{361}$而是仰望那位被舉起的基督.
相信他的愛.
正正是要聖化這個城市.
當然是要透過我們.
而我們不再只是深夜裡徘徊.
而是在基督的愛裡.
得到重新的勇氣.
因為上帝愛世界.
所以我們沒有理由帶著絕望去生活.
因為上帝愛世界.
所以我們有責任去愛.
在大齋祈當中.
我鼓勵每一個.
多默想.
更重要多行動.
耶穌基督的教訓.
聖經裡面的教導.
很多時候需要我們去思考.
但更需要我們用行動去回應.
你見不見到你身邊的世人?.
你見不見到上帝所愛的世界?.
求主賜我們力量.
勇氣和盼望.
更多的愛.
一同去見證.
榮耀祂的名.
請同心祈禱.
天父上帝在這一刻.
求你親自向我們說話.
頂著我們這一刻去安靜.
去到上主面前.
如果我們腦裡閃出一些.
我們身邊真的不懂愛的朋友.
家人 同事.
又或者.
一些事物.
我們很大的情緒憤怒.
我們過去很想用邏輯理性去解決.
但今天.
主教導我們.

$^{401}$要去明白怎樣去愛.
要去實踐這份愛.
要去打破這些框架.
我們就去到主的面前.
求主憐憫.
我給你一點時間.
為你祈禱.
為你自己禱告.
為你身邊需要你回應的人.
和事物去禱告.
愛我們的恩主.
我們去到你面前.
求你施恩憐憫.
多謝你.
讓我們在世界當中.
去經歷生命種種.
當中的豐富.
當中的精彩.
讓我們嘗試機會去反省.
我們應該怎樣去活.
多謝你.
你透過昔日彌果德母去尋道.
再一次提醒我們.
我們生命尋求的是什麼呢?.
可能首要問的問題.
就是我們有沒有這份去尋求你的心.
主要是幫助我們.
幫助我們學到什麼叫愛.
在這個世代當中.
更需要愛.
但很多時候我們愛得沒有力量.
因為我們的愛太空泛.
我們的愛常常只是放在嘴唇邊.
或者在理念上面.
主要是幫助我們伸手去愛.
你愛的世界.
就是反叛的世界.
就是癡離破碎的世界.
你愛的世人.
就是那些跌跌碰碰.

$^{441}$常常都叛逆不聽話的世人.
但你為你所愛的世界,世人.
祈禱你道成肉身.
並且一步一步走上十字架.
將我們的手流出你的寶血.
等待我們擁抱回應.
還有幫助我們.
能夠成為知道如何活出愛的人.
能夠在領受你伯伯恩典愛的同時.
也願意去參與.
帶來希望.
帶來這個世界力量的愛.
幫助我們.
幫助我們能夠更加明白.
你的愛的意義.
更讓我們每一次再讀出.
上帝愛世人的時候.
不只是空泛的口號言語.
而是我們生命的力量.
並且是我們合圍的見證.
給力量我們這個群體.
盡心的去服侍你.
看見不同的需要而去伸手.
看見那些甚至連我們自己都很難接受的人.
和事和物.
但你仍然帶著盼望去愛.
只知道不容易.
但知道你的恩典夠用.
就求你給我們更多的恩典.
因為我們知道我們愛是上帝先愛我們.
求你賜恩憐憫.
求你讓我們繼續去反省.
繼續去思考.
特別在這個大齋祈日子當中.
好好地緊緊地去到你面前.
以致能夠得到你十字架的恩典.
在我們生命裡面彰顯十字架的榮美.
多謝主聽我們禱告.
奉耶穌基督的命求.
\newpage



\section{創世記 24:27}
\label{sec:031723}
\textbf{聖經人聖經事:耶和華引領我}~~~~黃克勤牧師博士
\newline
\newline
連結: \href{http://www.church.com.hk/acms/content.asp?site=cdc\&op=show\&type=product\&code=031723\&layout=sermon}{\texttt{031723}} ~~~~ 語音日期: 2026-03-18
\newline
\newline
\hyperref[sec:031722]{\small{< < < PREV SERMON < < <}}
~
\hyperref[sec:index_chronic]{\small{[返順時目]}}
~
\hyperref[sec:031724]{\small{> > > NEXT SERMON > > >}}
\newline
\newline
創世記 24:27
\newline
\begin{longtable}{cl}
\hline
\hline
章節 & 經文 (和合本修訂版)\\
\hline
24:27 & \begin{tabularx}{0.7\textwidth}{X} 說：「耶和華—我主人亞伯拉罕的神是應當稱頌的，因他不斷以慈愛信實待我主人。至於我，耶和華一路引領我，直到我主人的兄弟家裡。」 \end{tabularx} \\ \\
[1ex]
\hline
\hline
\end{longtable}
$^{1}$各位網友大家好.
今天跟大家讀聖經.
繼續看亞伯拉罕如何跟自己娶媳婦.
他就叫了一個僕人.
是他最信任服侍了他幾十年的老人家.
這個老人家.
我也很想認識他.
不過是沒有名字的.
只是那個人.
那人.
那僕人.
那人.
將來在天上見到.
我也想有機會認識一下這個人.
非常聰明.
應該沒有讀過書.
不過很值得學習他辦事.
以及對上帝的敬畏.
聖經人聖經事生命說書網絡為大家提供.
今天來到耶和華引領我.
希望我們更加認識耶和華.
除了認識人.
也認識背後神的工作.
今天就講到那個人定睛望一個女孩子.
那人就是一個僕人.
其實上集結尾是這樣的.
根據上一集.
昨日.
求你施恩給我主人亞伯拉罕.
使我今日遇見好機會.
這個僕人就是由瓦迦南地.
走向東北.
去到兩河流域的起源附近.
然後兩河流域向東南再走.
然後兩河流域伸入沃土.
去到兩條河的出海地方附近.
叫做吾爾城.
就是亞伯拉罕以前的家鄉.
想娶一個本地本家本族的女孩子.
做自己的媳婦.

$^{41}$其實這件事很難.
因為你那麼遠去找一個女孩子.
你肯不肯跟我回去.
嫁給一個你未見過的男孩子.
不過要走那麼遠.
如果以當日是回不來.
其實我看完這段經歷.
心想哪天會成功.
誰知顯出這件事後面.
真是亞伯拉罕的神厲害.
是給他遇見好機會.
他去到那個地方的時候.
這個僕人很厲害.
還有什麼呢?.
耶和華引領我.
雖然他那麼厲害.
他就把十隻駱駝.
放在水井旁邊.
他就在等黃昏.
很多女孩子沒那麼曬.
出來打水回家的時候.
他就看看.
他心裡有個計劃.
這個僕人真厲害.
就是我問哪個女孩子拿水喝.
他又肯給我水喝.
又肯給我的駱駝喝.
那麼好心的女孩子.
這個沒有錯.
又肯做事又殷勤.
於是他就見到一個女孩子.
他就問她.
可不可以給我水喝?.
她就馬上拿來給他.
她不是遞一個瓶子給他.
來一個礦泉水自己喝嗎?.
瓶子很大.
頂在頭頂.
很重的.
裝了水.

$^{81}$利伯嘉馬上拿著瓶子給他喝.
於是他拿不住.
這樣喝了水.
不單止你喝飽了.
還給你駱駝喝.
於是再打水.
再打水.
為什麼這麼多個再打水呢?.
因為有個槽給十隻駱駝喝.
十隻駱駝喝.
一隻駱駝喝你一瓶水.
你就打十瓶.
如果他每一隻喝你兩瓶水.
你就會打二十瓶.
如果有一隻駱駝喝三瓶水.
你打水打三十瓶.
還有很多僕人.
手乖起展.
手都避開了.
腰那一呎.
但這個女孩很開心.
很樂意.
這就是利伯嘉.
第一點.
定睛看.
那人就是僕人.
定睛看他.
一句話也不說.
要曉得搖滑.
似他通達的道路沒有.
他是定睛看利伯嘉.
不說話.
感受一下上帝曉得搖滑跟他說什麼.
就是說搖滑會跟別人說話.
會安排.
而且他是活的神.
另外這個僕人厲害的地方是什麼呢?.
如果你想聽到神的聲音.
你必須要一句話也不說.
大聲說話的人.

$^{121}$是聽不到別人說話的.
因為他在說.
跟著說「哎呀!神啊!」.
等等等等.
「你為什麼不帶我路?」.
「為什麼不告訴我你的心意?」.
「當然了,說了幾次你聽不到而已」.
「因為你在說話」.
所以一句話也不說.
要等候神就是小說話.
另外定睛看.
是一個環境.
看.
但自己要靜靜地聽.
然後你就會曉得搖滑的心意.
下頭也學足了.
他看到這個女孩子真是好.
於是他就拿了一個金環.
半赤赫勒而已.
兩隻一對金手足.
就是金環.
重十赤赫勒.
就是喀布勒女子.
「多謝你啊,這位小姐」.
「這樣幫我們守」.
「送一個金環給你,半赤赫勒」.
「送一對金手鐲給你」.
「就是十赤赫勒」.
「嘩!這麼有趣,為什麼金環半赤赫勒」.
「手足金足十赤赫勒呢?」.
「因為這個金環原來是戴在鼻子上的」.
後來經文看到是戴在鼻子上的.
戴在鼻子上不是很重.
經常側著頭不可以.
鼻子也會拉穿.
戴在鼻子上是漂亮的.
一對金手鐲是漂亮的.
一動手別人就看到.
這樣就很好.
而且還戴在鼻子上.

$^{161}$「你告訴我你是誰的女兒」.
「你住在哪裡的」.
「家裡有沒有地方讓我們今晚投宿一晚」.
就問這個喜人的底.
「你是誰的女兒」.
單刀直入,又不是吃一飯.
「你喜歡什麼顏色」.
「你的眼睛大不大」.
「喜歡聽什麼歌呢」.
「有空去哪裡玩」.
「你是誰的女兒」.
「因為如果我要成親」.
「我要問你爸爸」.
所以單刀直入.
「你是誰的女兒」.
「家裡有沒有地方讓我們住」.
「如果去家裡住一晚」.
「就什麼都清楚」.
女子回答得好.
「竟然真是遇見好機會」.
「神的大力」.
「我是墨家語拿學之子」.
「彼土利的女兒」.
留意,拿學是誰?.
拿學這麼熟悉?.
墨家可能不太熟悉.
太太,丈夫是拿學.
拿學是亞伯拉罕的爺爺.
同宗族,同家.
「我是拿學的兒子」.
「彼土利的女兒」.
拿學是亞伯拉罕的爺爺.
亞伯拉罕的父親叫他拉.
他拉的父親叫拿學.
現在是亞伯拉罕.
拿學,亞伯拉罕.
亞伯拉罕和彼土利是兄弟.
因為拿學是亞伯拉罕的父親.
拿學的兒子就是亞伯拉罕.
和彼土利或其他人.

$^{201}$「我是他的女兒」.
這個僕人居然是上帝的引領.
這麼大能,這麼厲害.
遇到他了.
「我家有糧取的地方住」.
那個人低頭向耶和華拜.
「耶和華神,你真是應當稱頌」.
不斷地,真是感恩.
「不斷地以慈愛誠實待我的主人」.
他代表他的主人.
所有遇到的都是耶和華.
「不斷地恩待我的主人」.
「我這個小人物不要緊」.
「耶和華在路上引領我」.
就是這個題目.
「耶和華在路上引領我」.
「耶和華引領我」.
「走到我主人兄弟家裡」.
彼土利竟然是這樣發生.
所以第三點.
記得我們和神來往.
要留心聽,不要說太多.
然後你會發覺神會做事.
祂會直接,很賣力地做.
想好辦法,想好問題.
然後低頭向耶和華下拜.
說些實實在在的感恩的話.
這是一件這樣的事.
和大家去看.
今日的金句.
「耶和華,我主人阿巴拉亨的神」.
「是應當被稱頌的」.
「因為祂不斷地慈愛誠實待我主人」.
「至於我這個小人物」.
「我非常感恩」.
「因為耶和華在路上引領我」.
信不信?.
信不信耶和華上帝會在路上引領你?.
會.
其實神從我們小時候開始.

$^{241}$預備我們這個人直到今日.
真的.
很多回望才看到.
以前很多艱難經歷.
回望才看到.
原來神在無造操練我這個人.
預備我這個人成為今天可以為主做的事.
你也是這樣的.
所以你一定要找到.
「耶和華,引領我」.
就是今日的題目.
記得神是活的.
是真的.
是會引領我們的路的.
我們一起祈禱.
讚美我們的主.
你真是配得稱頌的神.
竟然不遠千里帶領亞伯拉罕的僕人.
竟然回到他自己本身的父家.
回到他的兄弟彼土利那裡.
大家的父親都是拿學.
竟然是這樣.
爺爺是拿學.
竟然是這樣發生.
主人,這件事是你的安排.
看到你為以撒預備了一個很好的太太.
其實是說為亞伯拉罕預備了一個很好的媳婦.
看到她是很熱誠的.
很殷勤的.
很智慧的.
身體很健康的.
很健壯的.
完全都是主的祝福.
至於這個僕人.
我們讚美主.
這個無名的僕人.
那人原來是你自己手裡的工具.
你都很珍惜,很帶領這個工人.
他又很賣力工作.
主人,這是我們看到,學到的.

$^{281}$求主幫助我們.
小說話多聆聽你的心聲.
順天禱告.
奉耶穌基督的名禱.
阿們.
\newpage



\section{使徒行傳 12:1-24}
\label{sec:031724}
\textbf{神掌權 (華)}~~~~楊芹牧師
\newline
\newline
連結: \href{http://www.church.com.hk/acms/content.asp?site=cdc\&op=show\&type=product\&code=031724\&layout=sermon}{\texttt{031724}} ~~~~ 語音日期: 2026-03-01
\newline
\newline
\hyperref[sec:031723]{\small{< < < PREV SERMON < < <}}
~
\hyperref[sec:index_chronic]{\small{[返順時目]}}
~
\hyperref[sec:031725]{\small{> > > NEXT SERMON > > >}}
\newline
\newline
使徒行傳 12:1-24
\newline
\begin{longtable}{cl}
\hline
\hline
章節 & 經文 (和合本修訂版)\\
\hline
12:1 & \begin{tabularx}{0.7\textwidth}{X} 約在那時候，希律王下手苦待教會中的一些人， \end{tabularx} \\ \\ \relax
12:2 & \begin{tabularx}{0.7\textwidth}{X} 用刀殺了約翰的哥哥雅各。 \end{tabularx} \\ \\ \relax
12:3 & \begin{tabularx}{0.7\textwidth}{X} 他見猶太人喜歡這事，也去拿住彼得。那時候正是除酵節期間。 \end{tabularx} \\ \\ \relax
12:4 & \begin{tabularx}{0.7\textwidth}{X} 希律捉了彼得，押在監裡，交給四班士兵看守，每班四個人，企圖要在逾越節後把他提出來，當著百姓辦他。 \end{tabularx} \\ \\ \relax
12:5 & \begin{tabularx}{0.7\textwidth}{X} 於是彼得被囚在監裡，教會卻為他切切禱告神。 \end{tabularx} \\ \\ \relax
12:6 & \begin{tabularx}{0.7\textwidth}{X} 希律將要提他出來的前一夜，彼得被兩條鐵鏈鎖著，睡在兩個士兵當中；門前還有警衛看守。 \end{tabularx} \\ \\ \relax
12:7 & \begin{tabularx}{0.7\textwidth}{X} 忽然，有主的一個使者顯現，牢房裡有光照耀；天使拍彼得的肋旁，叫醒了他，說：「快起來！」鐵鏈就從他手上脫落下來。 \end{tabularx} \\ \\ \relax
12:8 & \begin{tabularx}{0.7\textwidth}{X} 天使對他說：「束上腰帶，穿上鞋子。」他就照著做了。天使又對他說：「披上外衣，跟我來。」 \end{tabularx} \\ \\ \relax
12:9 & \begin{tabularx}{0.7\textwidth}{X} 彼得就出來跟著他走，不知道天使所做是真的，以為見了異象。 \end{tabularx} \\ \\ \relax
12:10 & \begin{tabularx}{0.7\textwidth}{X} 他們經過了第一層和第二層監牢，就來到往城內的鐵門，那門就自動給他們開了。他們出來，走過一條街，忽然天使離開他去了。 \end{tabularx} \\ \\ \relax
12:11 & \begin{tabularx}{0.7\textwidth}{X} 彼得清醒過來，說：「現在我真知道主差遣他的使者，救我脫離希律的手，和猶太人所期待的一切。」 \end{tabularx} \\ \\ \relax
12:12 & \begin{tabularx}{0.7\textwidth}{X} 他明白了，就到那稱為馬可的約翰的母親馬利亞家去，在那裡已有好些人聚集禱告。 \end{tabularx} \\ \\ \relax
12:13 & \begin{tabularx}{0.7\textwidth}{X} 彼得敲外門時，有一個使女，名叫羅大，出來應門， \end{tabularx} \\ \\ \relax
12:14 & \begin{tabularx}{0.7\textwidth}{X} 認出是彼得的聲音，歡喜得顧不了開門，就跑進去報信，說彼得站在門外。 \end{tabularx} \\ \\ \relax
12:15 & \begin{tabularx}{0.7\textwidth}{X} 他們對她說：「你瘋了！」使女堅持真有其事。他們說：「那是他的天使。」 \end{tabularx} \\ \\ \relax
12:16 & \begin{tabularx}{0.7\textwidth}{X} 彼得不停地敲門；他們開了門，一見是他，就很驚奇。 \end{tabularx} \\ \\ \relax
12:17 & \begin{tabularx}{0.7\textwidth}{X} 彼得做個手勢，要他們不作聲，就告訴他們主怎樣領他出監；又說：「你們要把這些事告訴雅各和眾弟兄。」然後，他離開往別處去了。 \end{tabularx} \\ \\ \relax
12:18 & \begin{tabularx}{0.7\textwidth}{X} 到了天亮，士兵中起了不少騷動，不知道彼得到哪裡去了。 \end{tabularx} \\ \\ \relax
12:19 & \begin{tabularx}{0.7\textwidth}{X} 希律找他，找不著，就審問警衛，下令帶走他們處死。後來希律離開猶太，下凱撒利亞去，住在那裡。 \end{tabularx} \\ \\ \relax
12:20 & \begin{tabularx}{0.7\textwidth}{X} 希律向推羅和西頓的人發怒。他們那一帶地方是從王的土地供應糧食的，因此就託了王的內侍大臣伯拉斯都的情，一心來求和。 \end{tabularx} \\ \\ \relax
12:21 & \begin{tabularx}{0.7\textwidth}{X} 希律在所定的日子，穿上朝服，坐在位上，對他們演講。 \end{tabularx} \\ \\ \relax
12:22 & \begin{tabularx}{0.7\textwidth}{X} 民眾一直喊著：「這是神明的聲音，不是人的聲音。」 \end{tabularx} \\ \\ \relax
12:23 & \begin{tabularx}{0.7\textwidth}{X} 希律不歸榮耀給神，所以主的使者立刻擊打他，他被蟲咬，就斷了氣。 \end{tabularx} \\ \\ \relax
12:24 & \begin{tabularx}{0.7\textwidth}{X} 神的道日見興旺，越發廣傳。 \end{tabularx} \\ \\
[1ex]
\hline
\hline
\end{longtable}
$^{1}$謝謝弟兄姊妹給我這個機會來講道.
然後我也一起來禱告吧.
求你賜給我們臺下聽的弟兄姊妹們.
給我們的聖靈幫助我們能夠聽得懂普通話.
能夠理解今天的信息.
也幫助我能夠講得清楚明白.
奉主耶穌基督的名囚.
其實我也試著聽粵語.
如果我知道講的內容的話.
我覺得我能聽懂百分之七十.
所以也希望今天希望大家聽普通話的時候.
因為你們瞭解內容瞭解使徒行傳能夠聽懂百分之七十.
今天我們來講這個使徒行傳的第十二章.
在使徒行傳第十二章的時候.
這個教會的形式就急轉直下.
就感覺到真的是黑雲壓境.
因為聖經一開頭就這樣說.
那時希律王下手苦害教會中的幾個人.
用刀殺了約翰的哥哥雅各.
他見猶太人喜歡這事又去捉拿彼得.
你聽聽啊教會兩個最重要的領袖.
一個雅各已經被殺了.
然後希律又要去捉拿彼得.
對於任何一個組織來說.
是不是非常非常的恐怖啊.
你可以想象當時的場景.
希望大家也及時的把手機或者是你的聖經打開.
因為今天沒有PPT.
所以我們會不斷的引用經文.
大家可以隨時來看.
那首先在介紹這段經文之前.
我們要強調它的歷史背景.
我們要先要認識兩個重要的人物.
第一位就是雅各.
這位殉道的雅各到底是誰.
聖經說這個雅各是約翰的哥哥.
他就是十二門徒之一.
約翰就是寫約翰福音.
啟示錄約翰一二三書的這位著名的約翰.
現在被殺的雅各就是約翰的哥哥.

$^{41}$雅各也是耶穌的重要的同工.
是他核心的同工之一.
在耶穌登山變相的時候.
他帶的這個人就是彼得約翰和雅各.
但是我們要知道寫雅各書的那一位雅各.
並不是這位殉道的雅各.
寫雅各書的那位雅各是主耶穌的弟弟.
那個雅各後來也成為教會的非常重要的領袖.
那下面呢.
還有一個特別重要的人物就是西律.
其實這個西律在聖經新約當中出現過好幾次.
你知道嗎.
每一次幾乎每一次出現.
他都不是同一個人.
因為西律不是一個人的名字.
西律是家族名稱加統治的頭銜.
西律是王的這個稱號.
在新約的歷史當中.
主要有四個西律.
我們逐一介紹第一個西律叫大西律.
就是耶穌出生的時候屠殺嬰兒的那個西律.
其實他就是這一代人的這個家族的曾祖父第一代人.
第二代呢.
到了審問耶穌就是耶穌遭釘十字架的時候.
不是有一個審問耶穌的西律嗎.
那個西律是大西律的兒子.
是第二代的西律.
然後呢.
今天殺害雅各逼迫教會的這個西律就是叫做雅基帕一世.
他是孫子.
還有第四個西律就是後來使徒行傳.
最後有審問保羅.
還有一個西律那個是曾孫子.
他叫雅基帕二世.
你看在耶穌的年代.
其實這個西律王的確是羅馬政府的一個傀儡.
他只是這麼一個猶太的王.
但實際上他沒有任何的權力.
所以當年猶太人想要殺耶穌的時候.
他們還需要借羅馬的手.

$^{81}$你看他們要在本丟彼拉多手下受難.
是吧.
但是呢.
在十幾年之後.
當這個兒子雅基帕一世繼位的時候.
當時的羅馬皇帝叫做克勞迪.
他又把管理猶太地區的這個行政權力重新交給了西律.
所以現在這個西律就可以明目張膽的來逮捕來來處決這個猶太人.
他又掌握了這個行政權力.
國家機器又重新的交回到他的手裡.
所以路加就特意的描寫當時關押彼得的這個情景.
你想有多嚴酷.
先殺了雅各.
然後呢.
這個西律讓十六個兵丁看著彼得.
分為四班.
一班四個人.
那就是十六個人.
一組一排的人都在關押彼得.
然後彼得被兩條鐵鏈鎖著.
第六節他睡在兩個看守中間.
一邊一個看著他.
監獄外面還有看守.
顯然這就是死刑犯的待遇.
這就說明西律絕對就是準備.
在除教節一過.
馬上處死彼得.
但是神奇的事情就發生了.
就在馬上第二天就要提審他行刑的前一晚.
突然天使降臨.
天使就拍醒了彼得.
對他來對他說快起來.
然後這個鏈子就立刻從彼得的手上脫落.
彼得當時還在反懵呢.
彼得以為自己在做夢.
他什麼都做不了.
因為為什麼說他什麼都做不了.
你從天使的幾次談話當中就能明顯的感覺到.
天使說樹上帶子穿上鞋.
就是提醒彼得穿上褲子.

$^{121}$樹上帶子實際上就是穿上褲子的意思.
因為當時的人都需要擠一個帶子.
穿上鞋.
然後彼得就那樣做.
天使又說披上外衣跟我來.
還得指揮他披上衣服跟上了來.
然後彼得就恍恍惚惚.
因為他一直以為他自己在做夢.
然後呢.
當天使帶領彼得走到第一層牢房的時候.
門就開了.
接著他們就如入無人之境.
走到第二層依然開了.
最後只走到最後臨街的一道門.
天使就打開門.
護送彼得出了監牢.
走過一條街.
只走到一個安全的地方.
走過那條街.
離開監獄.
天使就離開了.
此時彼得才突然醒悟過來.
原來我沒有做夢了.
原來真的是主拆遣他的使者.
救我脫離細律的手.
然後彼得就站在空無人煙的這個大街上.
此時肯定是業已很深了.
他想我去哪呢.
他就說.
我可以去稱呼馬可的約翰家.
去他媽媽家瑪利亞那裡.
因為在當時馬可家.
已經就是基督徒祕密的聚會的地點.
當時就已經有很多人在那.
他們就是一個家庭聚會的一個家庭教會.
而這位被稱為馬可的約翰.
就這個馬可也是我們新約當中的一位重量級的人物.
他就是馬可福音的作者.
所以我們的這個聖經學者就告訴我們說.
實際上馬可福音背後的作者就是彼得.

$^{161}$因為馬可就是彼得非常得力的助手.
馬可福音就是彼得口述.
然後馬可記錄下來的.
所以你看彼得和馬可的這個深切的友誼.
是不是這段經文也透露了一些.
當時彼得沒有地方去他覺得到處都很危險的時候.
他想到了可以去馬可家.
所以我們就看到這段經文挺有意思的.
實際上他講述了兩個事件.
第一個事件是不是一個大事件.
這大事件是什麼.
西律逼迫教會.
這是一個大的歷史.
就是一個大的歷史背景.
當時的情況整個的教會都在面臨著巨大的逼迫.
然後這個大事件當中又有一個小事件.
這小事件是什麼.
就是彼得獲救.
為什麼說他是小事件.
你看聖經記載的這個細節多清楚.
天使三次提醒彼得穿衣穿褲穿鞋跟我來.
連這樣的細節都記錄下來.
但是那個大事件雅各被殺.
僅僅是一筆帶過.
所以一大一小.
但是他們當中給我們透露的信息是非常非常一致的.
這信息是什麼.
三個字就是我今天的講道題目.
叫做神掌權.
在大的歷史背景當中.
是什麼樣的情況.
君王的強權.
這是西律是不是.
宗教的壓迫.
為什麼是宗教的壓迫.
猶太教徒一直在迫害基督徒.
還有什麼.
輿論的敵對.
為什麼說輿論的敵對呢.
你以為當年初代教會在這個歷史當中.

$^{201}$老百姓當中.
老百姓都喜歡基督徒嗎.
完全不是.
大家注意看一下第三節.
第三節怎麼說.
西律見猶太人喜歡這事.
他殺了雅各之後.
他發現猶太人覺得這事很好.
這證明什麼.
證明民間的聲浪.
輿情.
大部分的老百姓.
是不支持教會的.
他們恨不得把教會全部剿滅.
因為在他們的頭腦當中.
這個基督徒就是宣傳異端邪說.
他們褻瀆神.
最讓他們受不了的是.
他們認為基督徒是迫害摩西的律法.
不守摩西的律法.
所以你就想想.
早期的基督教面臨的君王的強權.
這國家機器就用警察.
監獄首先壓迫你.
然後社會的主流宗教團體.
壓迫你.
普通的老百姓壓迫你.
你想想任何一個組織.
有可能在這樣的嚴酷的形勢下存活下來嗎.
但是教會不但存活下來了.
而且迅速的發展壯大.
而且兩千多年也沒有被任何人消滅.
否則就沒有我們今天坐在這兒.
是不是.
所以教會存在.
發展壯大.
一直活到今天這件事本身.
就證明了我講到的題目三個字.
神掌權.
它還證明了什麼.

$^{241}$教會就是神的旨意.
我們不管我們的教會當中有多少瑕疵.
有多少不完美的地方.
我們都不能否認一個事實.
我們存在這兒.
我們今天在這兒.
這是神的旨意.
因為神掌權.
所以無論時局如何變化.
神掌權.
你看陸家接著說.
我今天特意.
因為宣牧師是準備講到第十九節.
但是我多講五節.
我講到了二十四節.
所以我請Peter弟兄多讀了五節.
這後面的經文很有意思.
陸家接著說.
說西律王惱怒推羅西頓的人.
那一帶的地方是從王的地圖得良.
因此就託了王的侍臣勃拉斯都來求情.
一心求和.
這一段歷史是什麼意思.
我為什麼要講這個呢.
因為這個是前呼後應了.
前面講了西律.
那麼我們就要把西律的結局也講了.
是不是.
其實我現在講不光講聖經.
我還要講歷史當中的記載.
這個猶太曆史上最著名的歷史學家叫作約瑟夫.
他在猶太古史當中就記錄了這件事情.
他說這個推羅和西頓是怎麼回事.
他們是兩個港口城市.
非常小的港口城市.
他們這兩個城市是不產糧食的.
他們的糧食都是靠猶太地區出口進口給他們.
然後他們得罪了西律.
西律就怎麼樣經濟制裁他們.
就跟今天的政治很相似.

$^{281}$西律就經濟制裁.
他們就害怕了.
他們就派了使臣來求和來服軟.
然後西律.
這就給西律一個作秀的機會叫什麼一場政治秀.
因為任何一個政治家.
他為了就是掌權做政治表演.
政治秀是非常重要的.
因為他需要獲得人民振臂高呼呼喊他萬歲.
然後他就是君臨天下的那種效果.
他需要得到這種效果.
非常重要.
現代古代都一樣.
所以西律就抓住了這個時機.
你看有這樣的使臣來朝來求我.
然後當天正好是紀念羅馬皇帝的慶典.
這個史學家約瑟夫就說.
當日西律王穿上銀線製成的袍子.
坐在寶座上.
當太陽光照在他的衣服上.
就銀光閃爍.
閃閃發光.
百姓就被大大的震撼了.
他們就高呼說.
這是神的聲音.
不是人的聲音.
但是就在西律接受人們歡呼的時候.
他突然腹痛.
五天之後他就死了.
這是真實的歷史.
約瑟夫記錄的是歷史.
而路加在《史圖行傳》當中道出了真實的原因.
路加怎麼說.
西律不歸榮耀給神.
所以諸的使者立刻罰他.
他被蟲所咬.
氣就絕了.
最有意思的就是最絕的是這段經文的最後一句.
二十四節說.
神的道卻越發廣傳.

$^{321}$這就是我今天講到的主題.
歷史的主角不是西律.
不是猶太的宗教領袖.
也不是喫瓜群眾.
而是神.
下面這就是神掌權.
從歷史當中來看.
下面我們從另外一個角度來看.
神掌權.
是什麼角度.
從一個小的角度.
從個人的命運當中來看.
神掌權剛纔是從大的維度.
因為這段經文當中有一個特別讓人困惑的地方.
我相信大家仔細讀的話都會有同樣的問題.
那就是為什麼雅各被殺了.
而彼得卻獲救了.
難道說因為神更愛彼得不愛雅各嗎.
或者說難道是不是雅各做了什麼事情.
神就懲罰雅各.
經文根本沒有重點.
根本沒有意思要探討這個問題.
因為經文的重點告訴我們.
不在於誰得救.
誰沒有得救.
重點在於神掌權.
重點在於神的旨意高過人的期待.
因為我們人我們通常會把神的祝福理解為什麼.
大家想神祝福我們了.
我們一定會想說.
因為我們事事順利.
因為我們教會人人都很平安.
因為我們的教會很安穩.
我們教會最近沒有發生任何事情.
我們教會大家都和和睦睦的.
神祝福.
我當然不否認這的確是神的祝福.
但是這只是事情的一面.
事情的另一面是今天的聖經告訴我們.
神祝福更深刻的意義在於神的旨意成就.

$^{361}$你想到我們每個人的命運.
我們的個人的恩賜是不一樣的.
我們遇到的挑戰是不一樣的.
我們生活的境遇是不一樣的.
神所給我們的壽命年限.
身體是否健康.
其實都是不一樣的.
這並不意味著生活在聖經當中的人.
神就多愛他一點.
神就更祝福他.
如果要真的是這樣.
我舉一個最好的例子.
是不是耶穌的十二門徒都是被咒詛了.
因為他們中十一位都殉道.
就是今天獲救的彼得.
他今天獲救了.
我們都知道他最後是怎樣殉道的.
被刀釘十字架.
我們都知道保羅.
保羅經歷了多少神奇的事情.
甚至人們一度以為保羅是不死的.
保羅真的是被劇毒的毒蛇毒了以後.
保羅把蛇一下子甩出去.
自己完全沒事.
但保羅最後也是被處決了.
是不是.
我們怎麼樣解釋呢.
其實我們真的不能從表面上看.
好像平安就是神給我們每一個人的生命的目標嗎.
不是.
神真正的目標是讓他的榮耀.
他的旨意在我們的生命當中彰顯.
讓福音繼續前進.
想到我們個人的生活.
如果我們在病痛中.
在患難中.
你要怎樣榮耀神.
上帝要讓你在這樣的狀態當中榮耀神.
你覺得怎麼可能.
其實你想想完全有可能.

$^{401}$因為就比如說對我來說.
最讓我感到受觸動.
鼓勵我信心的.
是我看到那些在患難當中的弟兄姊妹.
仍然靠著神大有喜樂.
我就完全的真的被震撼到了.
因為生活順利.
富足.
信上帝的人可以做到.
不信上帝的人也有做到的.
拜佛的人也有做到的.
求神的人.
在廟裡求的人也能求到.
而真正我覺得在逆境當中.
獲得從天上的平安.
神的喜樂.
那種不是一般人可以做到的.
必然是上帝親自所賜予的.
其實小事是如此.
大事也一樣.
比如說我們的國家的家庭教會.
你知道在十多年前發展非常的迅速.
在十多年前.
在一線的城市裡面.
家庭教會人數居然都是可以破千人的.
有的家庭教會甚至有自己舉辦的幼兒園和小學.
我們想說.
上帝真的很祝福我們中國的教會.
是不是.
當然我們可以說神有祝福.
但是如今家庭教會被迫分散成更小的小組.
我們難道就覺得神不祝福了嗎.
恰恰相反.
你知道當基督徒在這樣的情況下的時候.
最直接的第一件好處就是.
基督徒必須非常認真地審視自己的信仰.
去教會這件事情.
絕不再是隨隨便便出於好奇和社交.
認識朋友.
履行公事等等任何這樣的原因.

$^{441}$如果你出於這些原因.
你是不會去教會.
你必須是從心裡切切地渴慕.
與弟兄姊妹一起去敬拜神.
難道這不是好事嗎.
還有因為分成了更小的團體.
以前是什麼事.
哎呀我做不了有牧師呢.
因為大家都有惰性.
這是人的本性.
但現在每一個人都要去牧羊.
每一個人都要去學著教導.
每一個人都要參與去事奉.
每一個人都要去決策.
其實這反倒成了培訓基督徒的臥土.
這難道不是好事嗎.
所以無論是我們的個人生活.
還是教會暫時遇到困難.
我們都不能夠很膚淺的.
很表面的泄氣說.
神不祝福我.
難道神不聽我的禱告.
不是這樣.
今天的經文告訴我們.
不要泄氣.
要知道神掌權.
最後我們從第三個角度.
我們看到歷史.
我們看到個人.
我們還要看到.
既然我們知道神掌權.
我們該如何做.
如何回應.
這就是今天經文告訴我們的第三點.
就是禱告.
你看經文說.
他們抓了.
殺了雅各.
抓了彼得.
教會卻怯怯的禱告.

$^{481}$在所有的教會施工當中.
我覺得最難推動的可能就是禱告會.
我們可能就覺得.
我們自下禱告就好了.
但是我們忘記了.
弟兄姊妹肢體連接在一起的禱告.
是何等的重要.
為什麼這麼說.
因為禱告不僅僅是求神.
成就我們的心願.
我相信當年為彼得禱告的那些基督徒.
不會說大家聚在一起.
就禱告這些事.
求神求你放了彼得放了彼得放了彼得.
讓他平安讓他平安讓他平安.
你覺得會是這樣嗎.
大家坐在一起就這麼禱告.
就重複這一句話.
絕對不是.
我相信他們是全方位的禱告.
我相信.
比如說如果你是當年禱告的其中之一.
你當時就處在那個情況下.
我想問問你會怎樣禱告.
你會怎樣為彼得禱告.
我相信你會這樣禱告.
求神保守彼得的信心.
即使在監牢裡也站立得穩.
你可能會禱告說求神賜給彼得平安和喜樂.
即便他面對身體的酷刑.
或者是即將被處決.
他心裡有從神而來的平安.
或者你肯定還會這樣的禱告.
神啊.
我們如今沒有了彼得.
沒有了雅各.
但求你給我們力量.
讓新起新的領袖.
求你使用不配的我們.
我們不管是年少也好.

$^{521}$還是年老也罷.
主求你讓我們親自的成為這個教會的一份子.
真正的一份子.
為他的興衰.
為福音的廣傳.
奉上我自己的力量.
所以這個禱告絕對不是說像求神拜佛一樣.
不停的唸叨.
只為了那一件事的禱告.
一定是全方位的一種禱告.
這個禱告.
跟不信的人拜拜的禱告.
還有一點特別不一樣.
我們基督徒的禱告.
因為基督禱告是通過禱告參與一場屬靈的征戰.
我們基督徒用禱告來全方位的封住每一個屬靈的破口.
為什麼這麼說.
你以為僅僅是希律猶太教徒在迫害基督徒嗎.
不是.
背後有一個黑暗的屬靈的勢力.
是背後撒旦黑暗的邪靈的勢力.
也裹挾著這些人間的傀儡.
不明就裡的人.
糊塗的人.
心裡黑暗被罪捆綁的人.
裹挾著這些人來逼迫教會.
所以我們基督徒的禱告是在全方位堵住這個魔鬼撒旦的破口.
不讓他在這個我們最薄弱的時候進攻進來.
也許我們外面會遇到很多的仇敵.
但是我們禱告求神讓我們裡面越來越合一.
最重要的就是合一.
如果你想象一下外面已經有了那麼大的壓力.
我們裡面再因為恐懼焦慮.
或者說一些自私的屬人屬血氣的虛榮的這些原因.
我們的裡面開始互相的產生不愉快.
結成小團體.
互相的不開心.
那麼教會立刻就散了.
所以我堅信.
當年教會的禱告是這樣的全方位的.

$^{561}$屬靈的征戰的打響屬靈戰爭的這樣的禱告.
為什麼呢.
還有一個原因.
我覺得是一個小小的細節.
就是當彼得真的被獲救的時候.
他就馬上來到了馬可家.
當他敲門的時候.
有一個使女羅大去開門.
羅大開門.
他激動的.
他居然聽到彼得的聲音.
他連門都沒開.
又往回跑.
跟大家說彼得回來了.
但是大家沒辦法相信.
他說你瘋了吧.
然後他說就是彼得.
大家說嗯肯定是他的守護天使.
他們覺得可能彼得已經死了.
很多的牧者講到這兒都會略帶.
就是輕微嘲諷的說一句.
你看他們的信心好小啊.
他們為這個事禱告.
他們都不相信神會成就.
其實我覺得恰恰顯明瞭他們的信心很大.
他們真的是為彼得獲救禱告.
但是他們不停在.
他們不僅僅為彼得獲救禱告.
所以我也是希望我們每一個基督徒.
是的.
我們是希望我們的病得醫治.
我們是希望我們的房子能賣一個好價錢.
我們能買一個好價錢.
我們是希望我們的孩子能找到好工作.
家庭幸福.
我們一定要為這些事情禱告.
因為我們就是神的兒女.
我們就活在這個世界上.
我們的喫喝拉撒還是要在這個世界上運行.
但是我的意思是.

$^{601}$我們的禱告不能以此為句號.
我們不能僅僅為這些事情禱告.
我們更多的要禱告.
上帝改變我們裡面的這個人.
這個本質.
改變我們看待苦難.
看待問題的那種態度.
給我們從天上來的屬天的價值觀.
讓我們真正有一個更新的生命.
然後當我們自然地把這個生命活出去的時候.
別人才能看到說.
噢你看他是不一樣啊.
他有這麼大的喜樂.
他有這麼多的智慧.
他遇到這件事情有這麼的平和.
他的愛心怎麼就確實比我們其他幾個人多呢.
這個時候神的榮耀自然就彰顯出來了.
神的旨意自然就在你的生命當中成就了.
所以這就是我今天要說的.
講了這麼多.
終歸就三個字神掌權.
教會是不會被消滅的.
無論強權是什麼樣的.
因為神掌權.
雅各殉道了.
彼得獲救了.
我們個人的命運高高低低.
不要被表面所迷惑.
因為神掌權.
所以如果我們知道的話.
我們就更應該切切地禱告.
是不是讓神的旨意和榮光在我們的生命中成就.
讓我們一同禱告.
在的天父.
歷史的主角不是君王.
不是輿論.
不是環境.
而是你自己.
無論是雅各殉道.
還是彼得獲救.

$^{641}$你的旨意都在成就當中.
你的道仍然興旺.
主啊.
求你賜給我們真實的信心.
不是相信事情一定照我們的心意發生.
而是堅定地相信你掌權.
無論我們身處逆境或患難.
都幫助我們忠心跟隨你.
專心仰望你.
也求你幫助我們成為禱告的教會.
在屬靈的征戰中站立的穩.
使你在我們個人家庭和教會中成全.
願一切榮耀都歸給你.
奉主基督的名求.
阿們.
\newpage



\section{啟示錄 1:1-2}
\label{sec:031725}
\textbf{序:看見啟示錄}~~~~趙崇健牧師
\newline
\newline
連結: \href{http://www.church.com.hk/acms/content.asp?site=cdc\&op=show\&type=product\&code=031725\&layout=sermon}{\texttt{031725}} ~~~~ 語音日期: 2026-03-01
\newline
\newline
\hyperref[sec:031724]{\small{< < < PREV SERMON < < <}}
~
\hyperref[sec:index_chronic]{\small{[返順時目]}}
~
\hyperref[sec:031726]{\small{> > > NEXT SERMON > > >}}
\newline
\newline
啟示錄 1:1-2
\newline
\begin{longtable}{cl}
\hline
\hline
章節 & 經文 (和合本修訂版)\\
\hline
1:1 & \begin{tabularx}{0.7\textwidth}{X} 耶穌基督的啟示，就是神賜給他，要他將必須快要發生的事指示他的眾僕人。他差遣使者指明給他的僕人約翰， \end{tabularx} \\ \\ \relax
1:2 & \begin{tabularx}{0.7\textwidth}{X} 約翰就將神的道和耶穌基督的見證，凡自己所看見的，都見證出來。 \end{tabularx} \\ \\
[1ex]
\hline
\hline
\end{longtable}
$^{1}$我們今天一起去看啟示錄.
首先我們會問一個問題.
為什麼要去看啟示錄?.
很多人一聽到啟示錄.
第一個反應是很難明白.
很多怪獸,很多符號,異象.
好像在說世界末日.
很多災難.
又好像和我們現在的生活距離很遠.
沒有什麼實際的關係.
我就想跟大家說一個關於一本書的故事.
大家有沒有聽過這個世界有一本書.
叫做世上最難明的書.
世上最難明的書.
大家知不知道是哪一本?.
亞裔人有額頭.
果然很博學.
這本書叫做Woynich's Manuscript.
中文也無謂翻譯.
它是一本15世紀的手稿.
裡面有文字.
又有插圖.
裡面有很多植物.
有很多人體的器官,身體.
還有很多天文的圖像.
看上去覺得很有意思.
很深奧.
但問題就來了.
就是到今天為止.
沒有人真正知道他寫什麼.
文碼學家研究過.
語言學家研究過.
現在我們人工智能.
人工智能也研究過.
結果都是一樣.
我們看到.
但是我們不明白.
很有趣的.
很多人每次說起啟示錄.
其實都用同一個語氣.

$^{41}$很難明白.
很多怪獸.
很多符號.
好像看到.
但不知道他想說什麼.
於是啟示錄就變成一本.
只是專家才看的書.
末日迷才看的書.
某程度上屬於一本關我什麼事的書.
彷彿是一本不想別人明白的書.
但我又想問.
啟示錄是不是一本不想讓人明白的書.
其實啟示錄一開始已經說得很清楚.
第一個字是啟示.
啟示這個字不是隱藏.
不是加密.
意思是要揭開.
要顯示.
要讓人看見一些事情的意思.
啟示錄和這本書最大的不同.
就是一本.
有人說這本是惡作劇.
永遠都解不明白.
但啟示錄就是一本神很想我們看得明白的書.
所以啟示錄本身不是隱藏了什麼秘密的天書.
而是一本教導我們作為信徒.
怎樣去看這個世界的書卷.
是一本告訴我們.
在這個世界很混亂.
很逼迫.
很多不合理的時候.
讓教會重新看得見這個世界是何等的事.
今天的題目不是解讀啟示錄.
也不是破解啟示錄.
而是看見啟示錄.
我希望我們能夠像啟示錄的作者一樣.
我們看得見耶穌.
看得見世界.
我們也看得見啟示錄當中神給我們的訊息.
好嗎?.

$^{81}$讓我們低頭祈禱.
親愛的恩主求主你都幫助我們.
打開我們的眼.
不單單是我們的肉眼.
更加打開我們靈裡面的眼.
叫我們看見.
而在今年的當中.
我們教會正式去研讀啟示錄.
求主你叫我們看見.
每一個異象的意思.
並且叫我們不需要膽怯.
求主你在當中.
要是這個世界有任何事要發生的話.
父神.
我們但願主你與我們同在.
也求主你讓我們看見.
不單單是苦難異象.
而是看見主耶穌基督.
看見主你的保守.
看見主你的拖帶.
看見主你的安慰.
因為父神你說過.
上天,下地所有的權柄.
都是屬於你.
就連我們都是屬於你.
求主你教導我們馴服.
我們祈禱奉耶穌你的名求.
阿們.
首先 如果你有拍照.
你就知道.
你用相機去拍攝不同的物件.
你是需要用.
不同的鏡頭.
同樣的道理.
要去讀不同的書卷.
我們都需要用不同的方法.
在教會的歷史裡面.
信徒去閱讀啟示錄.
一般來說.
有四個很普遍的方法.

$^{121}$而我們大部分的人.
其實也有接觸過其中的一個.
讓我很簡單地去講.
第一個方法.
叫做歷史方法.
意思就是.
將啟示錄裡面的內容.
對應第一世紀.
去到我們現在.
第二十一世紀當中.
所發生的歷史事件.
包括了可以是第一次.
第二次世界大戰.
可以是早兩年的新冠疫情.
好不好.
就是這2000年裡面.
所講的事情.
過去方法.
時間有點不同.
就是將啟示錄.
解釋成為.
剛剛好.
是主後70年.
前後所發生的事情.
也是當時.
教會要面對的事.
舉個例子.
當時所遇到的大災難.
不再是世界大戰.
而是在說.
70年左右.
當時耶路撒冷和聖殿被摧毀.
另外一個方法.
就是未來方法.
就是將啟示錄的內容.
理解成為.
不知道將來時間的事件.
而是從現在開始.
到將來何時發生我們不知道.
就連初代教會.

$^{161}$看都看不明白.
並且這個解釋的方法.
通常會以字面的意思.
去理解.
舉例你就明白.
其中一個例子正正就是.
聖經裡面提到.
假如.
某些時間.
這個世界會有大災難的話.
有一部分的信徒.
會被提.
不需要經歷死亡.
被提到天上.
就是這個學說的解釋.
第四個就眾議明白.
靈異方法.
很簡單.
就是將啟示錄.
不需要再知道時間.
不需要再知道事件.
只需要將它看成為一本.
正邪之爭的屬靈小說.
重點就是.
最後正.
一定會勝過邪.
舉個例子.
用這個方法看聖經.
不主在來.
這件事情.
只不過是象徵主耶穌的勝利.
未必真的會回來.
已經很簡單.
講到四種方法.
而這四種方法.
我相信我們有聽過.
網上不少.
你現在看到啟示錄的影片.
甚至以前我們小時候.
所看到的電視.

$^{201}$那種新奇趣怪的神秘節目.
其實都是用了.
其中一種方法.
為什麼呢?.
因為娛樂性高.
吸引力大.
充滿神秘感.
但問題就是.
這四個方法.
都未必是啟示錄.
想我們.
看啟示錄的方法.
啟示錄第一章第一節.
其實很清楚地告訴我們.
它想我們看見的.
是什麼.
第一章第一節.
經文這麼說.
耶穌基督的啟示.
就是神賜給他.
叫他將快要.
必要快成的事.
指示他的眾僕人.
他就猜顯使者.
曉遇他的僕人.
約翰.
好像很直接的經文.
但需要更多的解釋.
耶穌基督的啟示.
這六個字.
可以有三個意思.
首先就是.
關於耶穌基督的啟示.
就是講關於耶穌基督的事.
另外一個解釋.
可以解釋成為.
從耶穌基督而發出來的啟示.
第三點.
第三個意思.
可以解釋成為.

$^{241}$耶穌基督由耶穌基督.
去擁有的啟示.
只要我們再細心去讀.
這一節經文和下一節經文.
我們就會看到.
其實三個意思.
都會成立.
「周師神賜給他」.
說的是啟示.
不是耶穌.
自己想出來.
創作出來.
而是三一神裡的父神.
賜給耶穌基督.
現在已經屬於耶穌基督.
教託他去將快成的事.
指示耶穌基督的僕人.
教託他去將快成的事.
指示耶穌基督的僕人.
你見到彷彿是一個行動.
我們讀的是字.
但其實是講一個故事.
耶穌基督就差遣使者.
巧遇他的僕人耶穌.
耶穌沒有停在天上.
他是主動.
他自己差遣使者.
將屬於他.
關於他的啟示.
傳遞給約翰.
所知道.
我一直都沒有急著去解釋.
甚麼是快成的事.
因為不少人一聽到快成的事.
當然是很神秘,很刺激.
好像通勝一樣.
會讓我知道很多不同的事情.
很離奇.
但經文繼續解釋第二節.
約翰接收到的是甚麼?.

$^{281}$約翰說.
他所見證出來.
他所接收見證出來的事.
就是神的道.
和耶穌基督的見證.
神的道是主耶穌基督講過的說話.
見證就是主耶穌基督的生命.
十字架和他的復活.
凡自己所看見的.
都證明出來.
再加上約翰所看見的異象.
只是這兩節換句話說.
其實啟示錄想說的啟示.
不是說甚麼離奇大事.
從頭到尾他最想說的那件事.
是一切關於耶穌基督的事.
即是說啟示錄不是讓我們猜測下次何時再有大疫情.
讓我多買些口罩和廁紙.
不是讓我們預測這個世界何時有世界大戰.
啟示錄是很想我們看見主耶穌基督.
聖卷書都是屬於耶穌基督.
從耶穌基督而來最後再一次要讓教會看見主耶穌基督.
近幾年我相信大家都知道.
有很多香港人都移民.
去了很多不同的國家.
其中一個國家就是英國.
我最近看到一個關於信徒移民去英國的生活故事.
片段是說他去到英國之後.
天氣很不同.
食物又很不同.
生活又很不習慣.
下班有什麼可以做.
你不看足球就去酒吧.
從前的家人朋友又不在身邊.
身邊的人又完全不同.
你以前有份工作.
但你現在做的那份工作都不同.
什麼都不同的時候到底你還是你.
人自然都會問.
慢慢這個基督徒就問自己.

$^{321}$我還適不適應.
我還可以撐多久.
有試過移民的人都明白.
移民其實一點都不容易.
我用個比喻就好像一棵樹.
原本在溪水旁的一棵很可愛的小榕樹.
你掘它出來.
你擺它去撒哈拉沙漠.
那棵樹平時旁邊是有樹的.
現在身邊全部都是仙人掌.
於是這個信徒都問自己.
他也有問神.
神你想我怎樣.
現在來到這樣的境地.
將來到底又會是怎樣.
其實我會覺得.
神想我怎樣這個問題.
可能我們也有問過.
順境你可以問.
逆境我們也可以問.
神你想我怎樣.
當然有時可以發著脾氣問.
神你到底想我怎樣呢.
這個發脾氣的我就不解釋了.
我就解釋那一方面.
是出於真誠的心的那一種發問.
意思是如果我知道神想我怎樣.
一般來說.
那我就照著做.
是馴服.
是好事來的.
如果我照著做.
我就會假設將來.
應該OK的了.
我就照著做.
感受上.
如果我知道神想怎樣.
就好像我多多少少.
都看見將來是怎樣.
將來就不再是一無所知.

$^{361}$整個人平靜很多.
心也安定很多.
所以歸根究底.
神想我怎樣.
是出於人.
對於將來的不安.
對未來的擔心.
我們很想肯定一件事.
就是我們的將來.
是OK的.
大家都想肯定.
不過我想反問一個問題.
如果我跟你說.
你可以看到你的將來.
你看到的將來.
比你想像中還OK.
生活很穩定.
身體很健康.
前途一片光明.
精神融光煥發.
還青春了.
但是你看到所有OK的將來.
你找來找去.
你看來看去.
你都看不到耶穌基督在那裡.
那我就想問.
到底將來OK比較重要.
還是看到耶穌基督.
比較重要.
耶穌知道的.
如果我們將來.
看來多OK也好.
但是我們看不到耶穌基督.
其實真是很不OK的.
耶穌也知道.
如果我們的將來.
我們看到耶穌基督在那裡.
其實將來是怎樣.
彷彿我們都覺得.
應該是OK的.

$^{401}$因為我們看到耶穌在那裡.
所以弟兄姊妹.
我們要想一個問題.
就是我們到底想看到什麼.
水晶球.
我們去看到將來發生的是什麼事.
水晶球.
還是你想看到耶穌基督.
在將來裡和我們同在.
和我們同行.
啟示錄這本書很有趣.
它不是單單用眼睛去讀.
它更加是用眼睛.
用心眼去看.
去感受書裡的內容.
啟示錄.
我們今年會看到.
會有很多一幅一幅的圖畫.
展現在我們眼前.
而圖畫裡的細緻.
所描繪的事情.
就是要帶出當中的訊息.
比喻就好像你去到圖書館.
去到一個博物館.
你手上可以拿著一份長刊.
長刊會對所有的畫.
很多的壁畫都有很多的介紹.
但是你要感受得到.
你就始終都要走到.
那幅畫面前親眼去看.
你才能夠感受.
你才能夠明白.
經文說.
若翰便將神的道.
和耶穌基督的見證.
凡自己所看見的.
都見證出來.
這是第一章.
首尾呼應.
一條條到啟示錄最後的一章.

$^{441}$若翰亦親口說.
這些事.
他看完了.
他就說是我若翰所聽見.
所看見的.
這個若翰.
就是耶穌十二個門徒之一.
就是寫若翰福音的那個若翰.
就是和耶穌生活了三年的那個.
當耶穌啟示異象給若翰的時候.
當時若翰已經九十歲.
在當時.
幾乎叫人歲.
他眼有矇.
耳有聾.
聽都應該聽不清楚.
好奇問一句.
為何不在若翰.
精精神神.
青青春春的時候.
幾十年前.
就把啟示錄的啟示交給他呢?.
這樣不是應該講得更加清楚嗎?.
是要等若翰九十歲.
耶穌才把啟示告訴他.
是要等若翰九十歲.
耶穌才把啟示告訴他.
因為耶穌知道.
若翰什麼時候.
什麼處境.
才需要這個啟示.
若翰書寫啟示錄.
是公元大約九十年.
他來耶穌傳道.
已經十字架相隔了六十年.
六十年了.
十二信徒.
大部分都應該去世了.
基督徒已經在一個國家.
剛剛開始被大規模的迫害.

$^{481}$若翰看到的.
就是一個如日方中的國家.
迫害一群.
像螻蟻一樣的基督徒.
若翰自己也像螻蟻一樣.
被人關在一個.
愛琴海的希臘小島.
一聽到這幾個字.
很漂亮吧.
愛琴海 希臘小島.
好像旅遊一樣.
這個島叫做拔摩島.
現在看起來很漂亮.
不過當時不是這樣.
是一個荒島.
島上有一個礦場.
專門送囚犯過去.
打石頭.
把石頭打出來.
不是一樣.
他或許會想到.
我九十歲了.
我都未必有機會.
離開這個島.
應該打石頭.
打到自己中腦.
是在這個處境.
耶穌就知道.
是時候.
讓若翰.
看到啟示錄的訊息.
在若翰這個處境.
其實他所看到.
所有的事都很不合理.
他在地上過每一天.
都很不合理.
我覺得九十歲的他.
覺得他一輩子.
都不合理.
他所見到的事.

$^{521}$見到的人 現在的自己.
正常都會問.
為什麼會搞成這樣.
是在這個時候.
耶穌基督就讓他看到.
震撼人心的異象.
新天新地.
主會來臨.
因為耶穌知道.
他需要.
一個訊息.
就是在這個世界.
即使每一天.
你的際遇.
所有看到的事都不合理的時候.
你解釋不到的時候.
其實或許已經不再需要解釋.
當所有的事.
都很不合理的時候.
耶穌是想說.
喂 其實都OK的.
因為當所有的事都不合理.
很混亂的時候.
有一個主宰.
這個主宰不一定.
令到所有的事.
很快地合理.
但他可以確保.
我們的生命 我們的將來.
在他手中.
失控 脫軌.
若翰有.
貴屬.
若翰覺得人生不合理.
我九十歲.
會不會這種感覺.
不合理.
不公道 不知如何解釋.
其實我們或許都覺得.
不陌生.

$^{561}$我留意到一宗新聞.
是關於一些明星的報道.
他不是說明星有什麼新歌.
有什麼新戲.
而是說一些以前.
跟我年紀差不多.
現在跟我年紀相若的明星.
他們的子女都已經.
準備讀大學.
不過報道的重點.
不是他們的子女讀哪間名校.
讀哪間大學.
而是說他們的子女.
很明顯.
遺傳了.
明星父母的恩賜.
男的.
他很英俊.
女的也很甜美.
有一個你們應該認識.
整個家族姓謝.
姓謝.
以前謝伯伯.
在黑白片年代.
已經高大衛猛.
謝伯伯的兒子.
謝濟濟.
都盡得遺傳.
樣子很英俊.
最近謝伯伯的兒子.
謝孫仔.
在雪梨.
搭地鐵被人拍到.
報道說.
他當天.
頭髮蓬鬆.
不知道是不是刻意蓬鬆.
但即使如此.
蓬鬆的頭髮.
都難以掩蓋.

$^{601}$他那種晉浪的樣子.
和明星味.
原來他現在在澳洲.
讀音樂.
將來有可能做明星.
看到這段訊息.
最後的句號.
我的感受就是連續三代.
都那麼英俊.
真是不公平.
我每天梳頭之前.
都頭髮蓬鬆.
星味就一定沒有.
早兩個禮拜.
天氣熱汗味.
我就會說有一點.
這些例子.
說來可能覺得很好笑.
但我們要面對一個.
人生真實的問題.
就是人生可以很不合理.
認真一點說.
我的小兒子的情況.
我都可以去問.
為什麼是我不是他不是你.
有些不合理的事.
是出身的不公平.
為什麼別人家裡那麼溫馨.
我一出生.
我什麼都沒有做過.
我家裡已經破裂了.
有些不合理的事.
是性格和能力的差異.
為什麼別人天生就那麼樂觀.
我很羨慕.
為什麼我天生就那麼悲觀.
看那麼多東西我都是負面.
有些不合理的事.
可能是疾病和命運.
大家都那麼努力.

$^{641}$為什麼升職的是他.
不是我.
弟兄姊妹.
當我們遇到不合理.
總會有情緒.
希望事情快點過.
過了就沒事.
時間可以沖淡一切.
但就算事情真的過去.
我覺得我們得到的.
是一個很短暫的解決.
因為很快.
我們心知肚明.
類似的事情.
很快又會再發生.
不是有沒有問題.
是時間的問題.
唯有當我們看到這個世界.
有一個主宰.
而這個主宰.
跟我們說.
即使.
所有事情看上去.
有多麼不OK.
但是你的生命.
你的將來.
都在主宰的手中.
他保證我們.
不會失控.
不會脫軌.
啟示錄就是想我們看到.
一位主宰.
啟示錄想我們看到.
一個主宰的時候.
相信這個主宰是掌管著現在.
掌管著將來.
上天下地的一切.
所以無論什麼事情發生.
即使何等不合理.
但是我們的生命.

$^{681}$永遠都不會.
在主耶穌基督當中.
失控脫軌.
而這個主宰.
就是耶穌基督.
看見主耶穌是主宰.
對我們今天.
有什麼意思?.
其實我不知道大家知不知道.
這本書裡面.
聖經.
其實差點就沒有啟示錄.
給我們看.
沒有了.
即是啟示錄.
在聖經裡面其實是最後一本.
被歸納成為.
聖經正典的書.
有很大的爭拗.
曾經覺得不如啟示錄整本不要放進去.
我們不用看.
但是如果沒有了啟示錄.
我覺得很可惜.
太可惜.
因為新約裡面.
沒有一本書.
很多時候都好像那些道基督教導.
教我們.
你應該要怎樣做.
教會描述你應該怎樣生活.
你就討神喜悅.
讀下去.
好像上一堂又一堂課堂.
說我們應該要怎樣做.
大部份的書卷.
對於主宰來之後.
是怎樣呢?.
都是輕輕帶過.
給個例子你看.
彼得候書.

$^{721}$我們照他的應許.
盼望新天新地.
自居在其中.
新天新地.
蜻蜓點水.
輕輕帶過.
新天新地是怎樣呢?.
如果聖經沒有了啟示錄.
我想我們彷彿沒有了.
對永恆那種盼望的期盼.
唯有對主宰來之後.
那種新天新地.
要描繪得繪形繪聲.
臨離盡至夠具體.
才能為我們帶來.
這種期盼.
和.
現在.
深層的安慰.
啟示錄這樣說.
城裡面.
新天新地的新耶路撒冷城.
街道裡面.
當中.
有一道生命水的河.
在河的.
這邊和那邊.
有生命樹.
兩旁都有.
不是只有一棵.
平衡的.
樹上結了十二樣果子.
十二樣代表很多種類.
很完全很齊整.
每個月都會結果子.
你不用怕它枯.
它每個月都有.
有得給我們吃.
樹上的葉子乃是為醫治萬民.
這個圖像看到有生命樹的果子.

$^{761}$不會再有飢餓.
有生命樹的葉子.
不會再有疾病.
接著啟示錄去描述.
將來是怎樣.
神要拆去它們.
一切的眼淚.
不會有死亡.
也不會有悲哀.
哭號.
疼痛.
因為以前的事.
都過去了.
我們經常都想以前的事.
都過去了.
弟兄姊妹.
我們想不想活在一個不再有死亡.
不再有悲哀.
不再有痛苦的世界.
到了那一天.
我們可以和很多.
我們失去了的.
親人和朋友.
見面.
到了那一天.
生命有缺陷.
有殘缺.
長期患病.
妥當.
坐輪椅.
可以走路.
患有腦退化症.
可以回復昔日的微笑.
今天.
我們覺得.
我們很需要的事情.
在那一天.
不再需要了.
今天我們覺得很沉重的事情.
那一天.

$^{801}$再沒有這個重擔.
今天.
我們覺得自己不好的地方.
到那一天.
一切都更新.
舊的事.
都要過去被醫治.
我們無法完全想像.
新天新地.
確實是怎樣.
在教會裡.
用了很多不同的.
彩色玻璃.
去刺激我們的想像.
大家見到的這一個.
兩旁就是生命的樹.
讓我們見到.
上面有果子.
中間就有一條.
生命水的河.
紅色的.
是我們啟示錄其中一個異象.
大紅龍.
是告訴我們.
邪惡會被克服.
這個世界.
不再有死亡.
不再有痛苦.
羊象徵的是耶穌基督.
永遠作王.
回來.
我們現在每一個人.
其實我們都在和耶穌基督脫離.
我們並不能面對面.
與耶穌基督相見.
到那一天.
上帝的帳幕在人間.
所以弟兄姊妹.
我的結語是這樣.
啟示錄.

$^{841}$被我們看見的那一樣.
不是.
將來的圖畫.
不是.
我不是要大家看到將來的事.
而是.
開始讓這個將來.
影響今天的我們.
而在我們的現實生活裡.
我們.
已經能夠找到盼望.
和安慰.
因為將來不只是將來.
將來已經切入我們的現在.
啟示錄.
不僅是將來.
而是我們的現在.
我們要看見的是將來.
因為要看見將來.
才能夠看見今天.
今天的訊息.
很簡單和大家分享三點.
也是盼望.
我們接著在啟示錄看到的訊息.
就是看見主耶穌基督.
看見生命的主宰.
看見盼望和安慰.
讓我們低頭祈禱.
稱祂為恩主.
我求主你真的打開我們的眼.
打開我們的心.
矯正.
我們想看到的是什麼.
心所看慕的是什麼.
因為父神我們總會被自己.
蒙蔽.
憤心.
我們看見是這個世界的事.
但是父神.
我們的聖靈.

$^{881}$加力給我們.
叫我們看見我們應當追求的.
乃是肉眼所不能看見的事.
我們應當追求.
天上的事.
我們應當追求看見主理.
父神我們的肉眼.
現在或不能面對面看見你.
但願意聖靈在我們的心中.
讓我們的心.
時常看見.
感受到.
主耶穌基督來.
無論我們當下發生的是什麼.
無論世界發生的是什麼.
無論經歷的破碎.
無論失去的多少.
無論疾病何等的辛苦.
無論將來的生活.
甚至我們都不敢去想.
但父神求主你.
教導我們無論眼光.
去到任何時候.
任何的地方.
我們都看見主耶穌.
你就在當中.
我們就在當中.
我們祈禱奉耶穌的名及.
阿們.
\newpage



\section{約翰福音 4:46-54}
\label{sec:031726}
\textbf{信心成長進程}~~~~謝友德牧師
\newline
\newline
連結: \href{http://www.church.com.hk/acms/content.asp?site=cdc\&op=show\&type=product\&code=031726\&layout=sermon}{\texttt{031726}} ~~~~ 語音日期: 2026-03-01
\newline
\newline
\hyperref[sec:031725]{\small{< < < PREV SERMON < < <}}
~
\hyperref[sec:index_chronic]{\small{[返順時目]}}
~
\hyperref[sec:031727]{\small{> > > NEXT SERMON > > >}}
\newline
\newline
約翰福音 4:46-54
\newline
\begin{longtable}{cl}
\hline
\hline
章節 & 經文 (和合本修訂版)\\
\hline
4:46 & \begin{tabularx}{0.7\textwidth}{X} 耶穌又到了加利利的迦拿，就是他從前變水為酒的地方。有一個大臣，他的兒子在迦百農病了。 \end{tabularx} \\ \\ \relax
4:47 & \begin{tabularx}{0.7\textwidth}{X} 他聽見耶穌從猶太到了加利利，就來見他，求他下去醫治他的兒子，因為他兒子快要死了。 \end{tabularx} \\ \\ \relax
4:48 & \begin{tabularx}{0.7\textwidth}{X} 耶穌對他說：「若不看見神蹟奇事，你們總是不信。」 \end{tabularx} \\ \\ \relax
4:49 & \begin{tabularx}{0.7\textwidth}{X} 那大臣對他說：「先生，求你趁著我的孩子還沒有死就下去吧。」 \end{tabularx} \\ \\ \relax
4:50 & \begin{tabularx}{0.7\textwidth}{X} 耶穌對他說：「回去吧，你的兒子會活！」那人信耶穌所說的話，就回去了。 \end{tabularx} \\ \\ \relax
4:51 & \begin{tabularx}{0.7\textwidth}{X} 正下去的時候，他的僕人迎面而來，說他的兒子活了。 \end{tabularx} \\ \\ \relax
4:52 & \begin{tabularx}{0.7\textwidth}{X} 他就問甚麼時候見好的。他們對他說：「昨天下午一點鐘熱就退了。」 \end{tabularx} \\ \\ \relax
4:53 & \begin{tabularx}{0.7\textwidth}{X} 他就知道這正是耶穌對他說「你的兒子會活」的時候；他自己和全家就都信了。 \end{tabularx} \\ \\ \relax
4:54 & \begin{tabularx}{0.7\textwidth}{X} 這是耶穌從猶太回到加利利後所行的第二個神蹟。 \end{tabularx} \\ \\
[1ex]
\hline
\hline
\end{longtable}
$^{1}$弟兄姊妹早晨.
今日我們繼續我們的約翰福音講到.
記載主耶穌的幾個神蹟.
我們再次一起祈禱.
多謝我們天上的父.
我們的主耶穌.
我們的救主.
我們生命的主.
教會的元首.
讓我們得到你的救恩.
又讓我們在信心的路上.
繼續的成長.
求主你帶領我們.
讓我們知道.
怎樣繼續和你.
有一個親密的生命的聯合.
願你今日所教導我們.
是我們信心的成長的道路.
祈禱奉耶穌的名.
我曾經聽過一個故事.
不過我不知道這個故事是真是假.
這個故事就說.
曾經有一間教會.
它隔離是一間酒吧.
這間教會就.
經常有少少弟兄姊妹覺得.
不是很開心.
在正教會旁邊.
有一間酒吧.
雖然通常酒吧.
早上很少那麼早開門的.
多數是晚上開的.
但是晚上更加害怕.
因為我們有時晚上有聚會.
在酒吧隔離有少少危險.
我們新的堂子518那裡.
不是真的在隔離.
不過是隔了一條街.
都有時好像覺得.
是不是有少少危險呢.

$^{41}$不過我們都考察過.
那間酒吧都不是那些.
很多醉酒佬去喝的那種酒吧.
我們隔離街那間.
是一些比較正當的酒吧.
所以不用怕.
但是這間教會隔離那間酒吧.
就真的時時會有一些醉酒的人.
在那裡出沒.
於是那個教會.
在不安裡面.
就時時都很懇切地祈禱.
甚至有時祈禱的內容就是.
最好這間酒吧不要再存在.
在我們的隔離.
這樣的祈禱.
後來有一天.
那間酒吧就火燭.
燒了那間酒吧.
酒吧的老闆就去告這間教會.
說這間教會令到我們的酒吧火燭.
燒成灰燼.
損失很大.
那個教會就要辯護.
你會怎樣幫這個教會辯護呢.
這個教會辯護的方法就是說.
雖然我們有祈禱叫這間酒吧消失.
不過它後來火燭.
消失了.
是不關我們祈禱的事.
那有什麼重要呢.
關不關祈禱的事.
就是說那間酒吧的老闆.
相信這個教會的人的祈禱是有效的.
因為他們祈禱所以令到這間酒吧消失了.
但這個教會裡面的人就說.
你們消失了.
你們火燭是不關我們的事.
是不關我們祈禱的事.
換言之就是他們覺得.

$^{81}$祈禱是沒有效的.
我們有祈禱你消失.
不過你消失不關我們祈禱的事.
你會覺得那間酒吧消失.
和教會的祈禱有沒有關係呢.
這個是很巧妙的問題.
有很多弟兄姊妹都知道我在雪梨.
曾經在唐人街一間教會服侍了很多年.
這間教會是布州第一間浸信會的教會.
是1836年的時候成立的.
有不同的時代.
有時是很興旺的.
有時又試過稍微衰微一點.
我在那裡的時候.
在那個教會做到第十五十六年的時候.
有一個危機.
危機就是在教會的旁邊.
有一間越南餐館.
如果去過唐人街那裡.
可能你還記得.
當時有一間越南餐館.
後來那間餐館就賣了.
賣的時候我們就收到消息.
說買了這間餐館的人.
去申請一個酒牌.
我們收到這些消息.
我們就像剛才說的那間教會一樣.
有少少不安.
如果將來在我們隔壁那間.
不是再是普通的越南餐館.
而是一間酒吧.
唐人街其實都有很多這些酒吧.
有時夜晚的治安.
我們都有點害怕.
在唐人街我們這間教會.
不要說是夜晚.
甚至日間我們打開門.
都會有不速之客走進來.
日光日白都會有醉酒鬼走進來.
有時我們都有少少擔心.

$^{121}$因為在唐人街的酒吧.
真的會有醉酒鬼出沒.
我們收到消息說有這樣的申請的時候.
我們就組織了.
不是很多.
不可以整個教會幾百人去.
組織了我們大約15至20人的弟兄姊妹.
打算去雪梨市議會抗議.
因為我們知道某一天.
雪梨市議會就會開會.
要決定給不給酒牌的新老闆.
我們預備去旁聽.
而且我們還預備了一些抗議的牌.
預備去舉牌.
我們浩浩蕩蕩去到市議會的門口.
誰知道有一個保安員在那裡.
那個保安員看到我們拿著牌.
雖然我們還沒有進去.
沒有亂來.
但他看到我們的牌.
他馬上跟我們說.
你們不可以帶這些牌進去.
不可以舉牌.
如果我一看到你們舉牌.
我就會趕你們走.
我們就乖乖地收起牌.
沒有去舉牌.
但是後來.
我們沒有去追查究竟有沒有拿到酒牌.
但是始終現在在教會的隔離.
還是一間餐館.
不會是酒吧那種餐館.
我們有祈禱.
祈禱神給新的老闆拿不到酒牌.
他最終拿不到酒牌.
或者他拿到酒牌.
而不是開酒吧.
我們又不是很清楚.
究竟跟我們的祈禱有沒有關係.
這些都是一些比較難解答的問題.

$^{161}$在現在很小的基督教教會裡面.
真的有一些教會.
是不相信祈禱.
不相信神蹟.
當然有時候可以這樣說.
那個教會其實是相信祈禱的.
那個教會是相信神蹟的.
不過教會裡面有一些弟兄姊妹.
未必那麼相信神蹟.
又未必那麼相信祈禱.
都不奇怪.
所以我們要問自己.
靈福堂一定是相信祈禱的能力.
靈福堂一定是相信會有神蹟出現.
不過我們要問自己.
我作為靈福堂的一份子.
我相不相信祈禱的能力.
我相不相信神會施行神蹟.
大家都會聽過.
有些教會甚至已經不相信.
瑪利亞童貞女生子.
耶穌是這樣出生的.
耶穌是渡成肉身的.
耶穌是死後復活的.
有些教會已經不相信這些東西.
所以有些教會最終是會關門的.
我們都在雪梨見很多的教會.
他們關了門.
有時候有些教會關門.
我們都說我們很感恩.
因為他關了門.
我們可以去買了他.
或者是租了他.
因為他已經沒有在那裡主日崇拜.
但是我經常都仍然有個感覺.
耶穌仍然是很傷心.
見到他的教會是關門.
有一些的基督徒.
他們都仍然相信.
耶穌是童貞女生子.

$^{201}$仍然相信渡成肉身.
仍然相信耶穌是死後復活的.
不過中間有一些他們又不相信.
譬如他們可能不相信耶穌將水變走了.
他們不相信耶穌是拿五餅二魚.
可以餵飽五千人.
有些我們叫他做新派神學那些.
他們就說其實耶穌怎會需要五餅二魚來餵飽五千人.
因為其實秘密是有很多人自己帶了食物.
不過他們不拿出來.
是那個小朋友有一點天真無邪地拿出來.
所以後來大家分來吃的時候其實有很多東西吃.
就不用耶穌去行神蹟.
真是有些教會.
對於神的能力對於神蹟這些東西.
開始不相信之後.
我們就知道是一種骨牌效應.
一樣不信就倒了.
跟著一樣倒了.
跟著一樣又倒了.
在信仰上越來越多這些.
倒了的東西.
大家都聽過很多.
美國 歐洲 英國 包括澳洲.
很多教會就是這樣關門.
我們今天讀神蹟的記載.
可以說是很簡單.
簡單到怎樣呢.
大宗到第五十節就說哪人信耶穌所說的話就回去了.
正下去的時候他的僕人迎見他.
說他的兒子活了.
他就是這樣相信耶穌.
這個的官.
甚至不是猶太人.
在猶太人和神的立約裡面.
這些是外邦人是沒有份的.
所以他對於舊約的聖經.
猶太人的立約和神的立約.
都不是知道很多.
但是他會這樣去求耶穌.

$^{241}$然後耶穌說了這些話.
他就信耶穌所說的話就回去了.
就見到他的兒子是好了.
當然不容易的.
無論你信了多少年.
有時可能你的信仰裡面有一些納愧.
在那些納愧裡面.
你的信心慢慢漸漸流失.
你明白我的意思嗎.
在那些納愧裡面流失了信心.
這個官就因為單純地相信耶穌所說的話.
耶穌閱納他的信心.
真的醫好他的兒子.
而且這個不是屬於猶太人和神之間立約的外邦人裡面.
竟然在他家裡面發生神蹟.
你在曠野裡面見耶穌行神蹟.
和你在家裡面見到耶穌行神蹟是很不同的事.
而且他見到耶穌行的神蹟.
耶穌沒有去到他家裡面行神蹟.
原來我們信的耶穌是一位超越時間空間的神.
超越時間空間的主.
對於我們人的腦海來說.
這件事是很難想像的.
所以難怪有很多人不相信神蹟的發生.
尤其是現今科學化的時代 人工智能的時代.
很多人更加覺得我們何必相信耶穌會行神蹟.
這些奇怪的說法呢.
但是在這些時代裡面.
我們更加要持守我們的信心.
上個月我講約翰所記載的水變酒的神蹟.
大家記得我們講過神蹟最重要的意義和目的.
沒錯 耶穌真是關心有需要的人.
耶穌真是會幫助有需要的人.
甚至微小到他在婚宴裡不夠酒喝.
其實是很小事.
不夠酒喝就少喝一點.
就算沒有酒喝 喝水也可以.
我們經常喝水 不是經常喝酒.
都一樣可以.
但是耶穌一樣是關心他們的需要.

$^{281}$在今天記載的神蹟裡面.
耶穌是關心有病的小朋友.
耶穌真是醫好他.
耶穌真是關心父親對兒子的愛.
真是醫好這個小朋友.
有時我們覺得我們的信仰有些不容易的地方.
就是究竟神蹟 病得醫治.
我們有信心還是沒有信心.
信心夠大還是信心很微小.
和神的旨意 神的權能.
這些全部是一種很動力的互相關係.
這幾樣東西究竟哪樣是最重要.
究竟哪樣令到哪樣發生.
似乎不容易可以完全了解.
所以有時我們的信心.
可以分為很多類型的信心.
有時我們的信心純粹就是.
希望一些事情會發生.
所以很籠統來說.
我今天和大家分享的裡面.
將信心這個進程.
或者信心的類型.
可以分為四種或四個階段.
很多時我們第一個階段的信心.
可以說是無可奈何的信心.
唯有信.
還有甚麼可以做呢.
我們既然甚麼都做不到的時候.
唯有靠信心 唯有去祈禱.
好像這個父親他見到自己的小朋友病.
所有最出名的醫生都束手無策的時候.
他就無可奈何地唯有去祈禱去求耶穌.
他甚至都不是猶太人.
他甚至對耶穌的認識其實不是很多.
這個官他家離開加拿大.
根據研究是大約二十里路.
所以他一定是要.
雖然不是叫做長途跋涉.
都要有一段距離 一段時間去找到耶穌.
而且耶穌沒有一個登記的住址.

$^{321}$耶穌是周圍去.
第二章說他在加拿大的婚宴行神跡.
然後他又去了另外一個地方.
然後今天這裡記載他又回到加拿大.
他要這樣找到耶穌.
因為他已經是完全沒有辦法.
無可奈何地去找耶穌.
有時我們的祈禱 我們的信心.
會不會都是一種無可奈何的祈禱.
無可奈何的信心.
這個父親已經沒有其他的辦法.
第47字結尾形容.
耶穌求他下去醫治他的兒子.
因為他兒子快要死了.
病得他知道兒子快要死了.
才唯有去找耶穌.
所以我們都要問自己.
當我們遇到一些困難 有些危機.
有些束手無策.
有些是完全已經在我們.
能夠控制範圍之外的一些困難.
你是不是什麼其他都做齊了.
最後的方法 唯有祈禱了.
我相信我們每一個基督徒.
都有一些時間會這樣想.
真的要問自己.
祈禱是不是因為所有其他的方法.
我都試過了 都沒有辦法解決.
最後唯有祈禱了.
這種心態其實都很普遍.
但我們今天要學習的是.
要將這種想法都解除.
信心的祈禱不是因為其他的事我都做齊.
都沒有辦法 所以唯有祈禱.
但當然反過來說.
我們不可以說我只是祈禱.
由頭到尾只是祈禱.
什麼都不用做的.
只是祈禱就可以了.
這樣又要戒掉.

$^{361}$因為我真的認識一些基督徒.
他們生病的時候.
我自小我父母都是基督徒.
自小他們教我.
有病不要去看醫生.
有病不要去吃藥.
只是祈禱就可以了.
我又覺得這種都要戒掉.
由一個極端去到另一個極端.
任何極端多數都不是好事.
在信仰上面都是.
任何一種極端似乎都不是很好的.
所以應該要平衡一點.
你又要盡你自己的力去做.
你應該要做可以做的事.
同時祈禱.
不是說我做齊什麼都搞定.
但是都搞不定.
唯有祈禱.
這是一個信心的學習.
可以說這是一個.
很多時候我們初信的時候.
信心的初階.
初初未懂得祈禱 未懂得信心是什麼.
又未習慣要時時祈禱.
要時時為某些事祈禱.
我們可能用這種心態來祈禱.
第二種的祈禱我叫等候回應的祈禱.
因為神有應許到一些事.
我們抓住所謂抓住神的應許.
我們抓住神的應許的時候.
我們就要等待神的應許實現.
普遍來說.
應許可以分為三種.
第一種的應許就是.
最普遍的應許 最普世的應許.
這個世界上任何人都可以得到神的應許.
就好像《約翰福音》三章十六節說到.
我們信耶穌的時候不致滅亡 反得永生.
是全世界的人都可以得到的應許.

$^{401}$第二種的應許.
就只有信耶穌的人才可以得到.
不信耶穌的人.
對不起 這個應許是沒有份的.
如果你是信耶穌的.
這個應許是有份的.
譬如《羅馬書》八章二十八節說.
萬事互相效力 叫愛神的人得益處.
如果你是信耶穌的.
你有份.
不過說少少.
這個應許有很多人都誤解.
因為我信耶穌.
所以萬事互相效力.
令到我得益處.
我這樣信了耶穌.
就有好處.
耶穌就會給我好處.
好處是什麼呢.
是真的有好處.
不過我們當時誤解了好處的意思.
誤解了好處的內容.
究竟好處是什麼呢.
有時候我們以為.
我只要這樣祈禱 信神.
我想要得到的東西 我就可以得到.
無論是我考試一百分也好.
我拿到高薪厚職也好.
萬事互相效力.
我就會得到益處.
但其實這裡的益處是說什麼.
這裡的益處其實是說.
我們在耶穌基督裡面.
那個屬靈生命的成長.
越來越像耶穌.
這段經文說的好處是這一點.
你不相信的話.
你自己看回八章二十八節前的那幾節.
是很明顯說這一點的.
不是說你可以發大達.

$^{441}$可以賺到三百萬.
雖然有 我知道你們很多人都賺到三百萬.
第三種的應許就是.
很個別的應許.
就好像神給亞伯拉罕的應許.
咒詛你的 他必定會被咒詛.
你會成為萬國之父.
這個只是給亞伯拉罕.
沒有給我 我沒有份.
只有亞伯拉罕才有份.
所以在聖經裡面也從來沒有應許.
說神是會治好所有人的病.
聖經從來沒有這樣說.
甚至沒有說耶穌應許.
凡是在基督裡信耶穌的.
你的病就一定會治好.
聖經也沒有這樣說.
你想想如果耶穌真的這樣.
治好所有信祂的人.
那麼所有基督徒都不會病.
所有基督徒都不會死.
但是聖經又明明這樣說.
人人都有一死.
你信耶穌也好 你不信耶穌也好.
人人都有一死.
這句話很有趣.
有些人知道人人都有一死.
這句話是真的.
每個人都相信的.
你不相信是不行的.
你不相信都會死.
但是接著那句.
人人都有一死 死後且有審判.
這句話有些人相信 有些人不相信.
前面那一半你一定要相信.
後面那一半有些人不相信.
後面那一半有些人相信.
又未必是相信耶穌.
有些人會相信閻羅王.
他都會死後且有審判.

$^{481}$不過和基督教所說的審判不同.
這些耶穌沒有應許的.
有時我們都很想得到這樣的應許.
所以我們祈禱.
是祈禱神的旨意成就.
這個其實理論上每個人都知道.
但是當你真的遇到危機的時候.
當你真的自己有病.
或者你所愛的人有病.
你有時就不太想見到這樣的實際情況.
第三種的信心.
或者信心的進程去到這個階段的時候.
希望我們能夠真的有信靠耶穌的心.
真正的依靠祂.
所以剛才經文第四章第五十節說.
大人信耶穌所說的話就回去了.
因為他相信.
因為他相信他就很放心地走了.
沒有繼續叫耶穌快些給我證據.
證明你是會醫好他.
或者他當時沒有手機又沒有whatsapp.
更加不可以video call.
否則他可能會叫耶穌.
耶穌你等等.
我開了個電話現在和兒子視像通話.
讓我可以看到他現在真的康復了.
那我就相信了.
他又不用這樣.
他又不需要耶穌說.
耶穌你等等.
我拿張紙來.
無論當時是什麼紙也好.
葡萄紙也好.
你簽個約給我.
你簽名證明你說你會醫好我的兒子.
如果我回到家看到兒子還沒有康復.
我會去法庭告你.
因為你簽了名.
你說你會醫好我的兒子.
當時不能這樣做.

$^{521}$沒有手機.
但是這個人他不需要耶穌這樣給他證據.
他就相信了.
就好像說耶穌將水變酒的神蹟.
瑪利亞和耶穌說他們沒有酒.
但是瑪利亞又很聰明.
又不會指著耶穌做事.
她知道她不應該指著耶穌做事.
雖然我們現在的基督徒有時會覺得.
我們可以指著耶穌做事.
叫耶穌幫我們做些什麼.
瑪利亞很清楚.
所以我們說耶穌回答瑪利亞的時候.
那句說話好像很不孝.
好像很粗魯.
我的事與你無關.
我的時候還沒有到.
但是我們看到的就是瑪利亞磨練她的信心.
耶穌的時候什麼時候到.
她要有信心去等候.
然後瑪利亞更厲害.
她跟那些僕人說.
我兒子叫你做什麼.
你就照做.
那些僕人.
我又不知道那些僕人是怎樣的僕人.
總之那些僕人都很聽話.
瑪利亞說他叫你做什麼你就做什麼.
他們就真的去做.
那種是我們都會很羨慕的信心.
我們都很想得著的信心.
對耶穌的信心就是這麼單純.
祂說做我們就做.
當然對我們來說.
那種信心不會突然的.
是真的要培養出來.
去到那個階段.
我們學習到單純的依靠耶穌.
所以再說.
約翰記載這些耶穌所施行的神蹟.

$^{561}$最終的目的我們上次也說了.
最終的意義就是指向耶穌.
耶穌是用這些神蹟來彰顯祂的榮耀.
彰顯祂就是神的兒子.
以致人會信祂.
以致信了祂的人.
我們的信心會更加穩固.
更加的加深.
這些是神蹟的意思.
當然我們也說.
耶穌真的是關心人的需要.
所以耶穌真的會聽人的祈禱.
耶穌真的會聽你的祈禱.
信心去到第四個階段.
是一種具感染力的信心.
這位官回到家的時候.
看到自己的兒子治好了.
第五十三節說.
他便知道這正是耶穌對他說.
你兒子活了的時候.
他自己和全家都信了.
很簡單的描述.
他和他的全家都信了.
這就是一種有感染力的信心.
你的信心有沒有這種感染力.
還是有時候半信半疑.
先不要告訴別人.
看看是不是最終發生了.
最終是不是如我所願.
最終看看是不是神.
根據我的想法來應允了我的祈禱.
然後第二天才告訴別人.
如果是這樣的話.
也有一點半信半疑.
所以我們的學習信心.
是有成長的階段.
我也想分享一些我們買518教會.
其中一些過程.
過程是很長的.
要講得詳細一點.

$^{601}$今天沒有時間.
其中有一些我們學習信心.
學習祈禱的功課.
也想和大家分享一些.
在我們開始要決定.
是不是真的去買教會的時候.
在執事和教育同工.
我們執事會開會的時候.
我們大家都很認同.
很肯定一件事.
如果我們要開始這個.
所謂的構團的項目.
我們必須要以祈禱先行.
然後一路要以祈禱去支持.
然後才做得成.
這是我們很肯定的一個共識.
這裡也可以說.
我們的發展委員會和整個執事會.
都有這種信心.
要以祈禱先行.
以祈禱先行.
不只是開會之前有開會的祈禱.
開會的祈禱就是祈禱先行.
還沒開會就祈禱.
所以也可以叫做祈禱先行.
但是如果只是開會的時候.
那兩分鐘的開會祈禱.
又似乎不是真的可以去表達.
我們對神的倚靠和培育信心.
所以這種祈禱是不斷不斷的.
所以到了今天仍然都是.
因為我們還沒可以開工裝修.
另外那兩棟的建築物.
又還沒很肯定我們究竟現在是否夠錢.
去做這個裝修完的時候.
夠錢結帳.
現在都還沒很肯定.
雖然我們都知道.
之前弟兄姊妹的奉獻.
各種的籌款.

$^{641}$我們都是有一些存款.
但是這些存款是不是就已經夠了.
不用再籌款了.
不用再為建堂奉獻了.
其實我們還不知道.
不知道這樣說有沒有嚇到弟兄姊妹.
不知道夠不夠錢.
會不會爛尾樓.
這個就是信心的學習.
到我們和業主再談.
他說要二千萬澳幣的時候.
我們都有些害怕.
二千萬澳幣去哪裡找.
難道無緣無故天要掉下來嗎.
於是我們就想.
不如不要買全塊地.
我們買一半.
一半就是一千萬澳幣.
比較可行.
我們都幾乎要去碼了.
但是神很奇妙用他的方法告訴我們.
停止這件事.
你不要走了.
買半塊地的計劃是不行的.
後來我們真的發覺.
再想深一層.
真的不行.
如果我們只買一半的話.
是省了一千萬.
不過沒有停車的地方.
在雪梨的教會沒有停車的地方.
就沒有人來崇拜了.
所以原來真的不行.
原來周氏會再開會的時候就想.
現在好像地價跌了.
應該值二千萬.
我們要和業主再去談.
要他降低價格.
就說不如一千四萬.
都有些可能.

$^{681}$雖然都不容易.
都有些可能.
但我們明明知道那個業主開價二千萬.
他又用了一些錢去搞了很多東西.
然後再拿來賣的時候.
一千四萬.
他九成九都不肯賣了.
我們再祈禱再商量.
我們加一點.
憑信心加一點.
加到多少好呢.
我們的信心應該去到哪裡好呢.
想了想.
不如一千七百萬.
加了很多.
加了三百萬.
離開他要價都是三百萬.
中間位應該差不多了吧.
於是我們和業主喝了幾次咖啡.
喝完幾次咖啡之後.
他都不肯減價.
我們說我們沒有錢.
買不起.
他說不要緊.
你沒有錢.
我借給你.
不過我要收十厘息.
我們覺得他不知道是開玩笑.
還是獅子開大口.
過了兩天之後.
他通知我們.
不如不要十厘.
照市價收你利息.
他自己都覺得獅子開大口.
我們繼續和他商量.
又繼續祈禱.
大家可能都記得.
我們曾經有一段時間.
大家要這樣祈禱.
希望可以減價.

$^{721}$但我們經過了幾個月的祈禱.
經過了這幾個月和他喝咖啡.
請了他喝幾次咖啡.
他都不肯減價.
喝了這麼多次咖啡.
都不肯減價.
我們有一個感動.
我們求神幾個月.
叫神感動他減價.
神都沒有答應.
我們求他減價.
又求神感動他減價.
過了幾個月之後.
我們的感動就是.
我們求了神幾個月.
似乎不應該繼續求神.
繼續求神.
似乎不是正確的方向.
雖然我們說祈禱.
當然要持續地祈禱.
要不斷地祈禱.
但有時買地買屋.
你又不可以繼續這樣祈禱.
一年兩年這樣不行.
神給我們的感動就是.
我們不要繼續求神.
我們有信心.
不過是求神感動.
求神公認.
令到我們會找到二千萬回來.
買他的地.
很不容易.
我都很明白我們的發展委員會.
和我們整體的執事會.
所有的教務.
有份參與的弟兄姊妹.
都知道是很不容易的一件事.
但我們這個信心的學習.
神這樣帶領我們走這條路.
我們真的買了這塊地.

$^{761}$其他的故事不說了.
因為沒時間說這麼詳細.
但在我們學習的裡面.
我們不單在個人的生命.
而且在教會的事工上.
同樣我們要這樣學習這些信心.
我們一起祈禱.
多謝我們天上的父你.
知道人的需要.
尤其是你知道.
你自己教會的需要.
你的兒女的需要.
你知道主耶穌的門徒.
我們要操練信心.
我們要完全的信服.
帶著單純的信心跟從你.
但這種單純的信心.
我們要知道不是一種盲從.
不是一種隨便的說.
我們信就行了.
神你繼續憐憫幫助我們.
在耶穌將水變走的神跡.
我們見到指向耶穌.
亦幫助我們去培育我們的信心.
在今天所記載的神跡裡面.
神同樣教我們.
要我們培育信心.
甚至當時的大官.
他似乎只是一個外邦人.
但他有這種信心.
甚至成為我們作為基督徒的信心學習.
他在這裡懇求我們天上的父.
繼續帶領我們.
祈禱奉耶穌的名.
\newpage



\section{創世記 12:1-4}
\label{sec:031727}
\textbf{冒險的祝福}~~~~王家輝牧師
\newline
\newline
連結: \href{http://www.church.com.hk/acms/content.asp?site=cdc\&op=show\&type=product\&code=031727\&layout=sermon}{\texttt{031727}} ~~~~ 語音日期: 2026-03-01
\newline
\newline
\hyperref[sec:031726]{\small{< < < PREV SERMON < < <}}
~
\hyperref[sec:index_chronic]{\small{[返順時目]}}
~
\hyperref[sec:031728]{\small{> > > NEXT SERMON > > >}}
\newline
\newline
創世記 12:1-4
\newline
\begin{longtable}{cl}
\hline
\hline
章節 & 經文 (和合本修訂版)\\
\hline
12:1 & \begin{tabularx}{0.7\textwidth}{X} 耶和華對亞伯蘭說：「你要離開本地、本族、父家，往我所要指示你的地去。 \end{tabularx} \\ \\ \relax
12:2 & \begin{tabularx}{0.7\textwidth}{X} 我必使你成為大國，我必賜福給你，使你的名為大；你要使別人得福。 \end{tabularx} \\ \\ \relax
12:3 & \begin{tabularx}{0.7\textwidth}{X} 為你祝福的，我必賜福給他；詛咒你的，我必詛咒他。地上的萬族都必因你得福。」 \end{tabularx} \\ \\ \relax
12:4 & \begin{tabularx}{0.7\textwidth}{X} 亞伯蘭就遵照耶和華的吩咐去了；羅得也和他同去。亞伯蘭離開哈蘭的時候年七十五歲。 \end{tabularx} \\ \\
[1ex]
\hline
\hline
\end{longtable}
$^{1}$各位弟兄姊妹平安.
多謝李牧師的邀請.
讓我可以在今天這個崇拜的當中.
和弟兄姊妹一起來分享神的話語.
現在我們再次同心祈禱.
願僕人口中的言語.
弟兄姊妹的心思意念.
蒙聖靈的啟迪.
蒙主的月立.
禱告籍主耶穌基督的聖名而求.
阿們.
不知道弟兄姊妹有沒有一些投資.
在投資裡面有一個法則.
叫做高風險.
高回報.
反過來說.
低風險.
當然也是低回報.
如果有人告訴你.
低風險可以有高回報.
你就要留意那個是不是騙子.
這個世界.
沒有這麼大的蛤蟆隨街跳.
怎麼可能.
不需要冒任何的風險.
就可以得到一個好的回報呢?.
如果你想收穫.
付出代價.
就是必然的.
我想在當中可能有些弟兄姊妹在讀書.
在讀書裡面我們一定會經歷過的.
你考試要過關.
應該也是要很努力地去溫習.
不會就是明天考試.
你今天晚上不斷地去跟神禱告.
希望明天的試題你會懂得回答.
一定不可能是這樣的.
又或者在職場裡面.
你有機會.
現在可能少一點.

$^{41}$升職.
加薪.
離開一個你熟悉的工作崗位.
要去承擔新的挑戰.
所以呢.
加薪是應該的.
受多一點氣.
你的薪金也包括在內.
這個就是代價了.
也有些弟兄姊妹可能已經退休.
你覺得也做夠了.
應該好好地去享受生活.
不過呢.
也有一些弟兄姊妹可能會選擇.
退休終於可以做一些自己想做的事情.
可能會希望有技能解鎖.
以往沒什麼機會嘗試的.
你就去嘗試一下.
又或者去做一些服事.
是以前當你在職的時候沒有機會的.
當然你要去有新的技能.
去一個新的群體.
也是會付出一些代價的.
去到日常的生活.
我想最簡單的一件事情就是.
如果有一家餐廳.
很出名.
你很想去吃.
那你要付的代價呢.
就是麻煩去排隊了.
剛才在崇拜之前.
李牧師就帶我去.
他說買太古城的名物.
那你去到.
咦不用排隊的.
他說以前要排隊的.
我不知道大家知不知道是什麼.
就是那些蛋撻.
以前很多人排隊.
現在就沒有了.

$^{81}$你想享受.
當然你要付出.
可能就要排一兩個小時了.
問題來了.
作為基督徒.
我們需不需要付出什麼代價呢?.
耶穌基督.
已經付出了最大的代價.
為我們去犧牲.
那我們又怎麼樣呢?.
我們有沒有一些東西.
應該也是要來付出的呢?.
其實今天.
就要來挑戰弟兄姊妹.
神父呼召我們.
要去付出一些代價.
不過不是要拿了你的命.
呼召你.
離開一個.
安舒的環境.
你覺得很穩妥的環境.
然後.
進入祂為你預備一個更大的祝福.
這個祝福不一定是.
物質.
不一定是令你有更多的財富.
但是這個祝福裡面.
是可以讓你的心靈得到很大的滿足.
可以發現.
原來神在你身上的旨意.
也是實現神一個更加重要的計劃.
而你是有份參與在其中.
這是一個非常大的恩典.
因此今天要從.
一段可能大家都熟悉的經文.
就是神去呼召亞伯蘭.
我們知道.
祂後來有另一個名字.
就是亞伯拉罕.
神呼召亞伯蘭.

$^{121}$祂這個經歷.
就是給了我們一個很重要的學習.
去領受神的祝福.
就要願意冒險.
踏出一個安舒區.
我們回到經文.
就是《創世記》的十二章.
但是在這一章的經文開始之前.
即是在十一章是在說這些東西的.
在《創世記》的第十一章.
三十一到三十二節.
是這麼說的.
就是亞伯蘭的父親.
叫做他拉.
他做了什麼呢?.
他拉帶著他的兒子亞伯蘭和他的孫子.
就是哈蘭的兒子羅德.
而當時哈蘭已經離世了.
所以他拉帶著自己的孫子羅德.
即是亞伯蘭羅德.
和他的媳婦.
就是亞伯蘭的太太撒萊.
一起出了迦勒底的梧椅.
準備要去迦南的地方.
當他們來到一個位置叫做哈蘭的時候.
他們就停下住在那裡.
在聖經當中記載.
撒拉一共活了二百零五年.
他就死在哈蘭這個地方.
就是說他並沒有去迦南.
這個原本他們要去的地方.
在哈蘭那裡他們就停下來了.
接著在《創世記》的十二章.
就是我們今天看的這一段經文.
記載神去呼召和猜險亞伯蘭.
要離開他們已經住了很多年的地方哈蘭.
要去到原本他的父親.
被猜險被吩咐要去的迦南這個地方.
我們看這一段經文第四節的最後那一句.
我們就看到了.

$^{161}$當時亞伯蘭被呼召的時候.
是七十五歲.
我就不問大家了.
已經到了七十五歲了沒有?.
或者還有多久七十五歲了?.
如果你到了這個年紀.
神要呼召你做另外一件事情.
我想這個挑戰也相當大的.
在亞伯蘭的爸爸.
七十歲的時候.
應該這麼說吧.
亞伯蘭七十歲.
哈蘭就是一百四十五歲.
因為他們相差了七十五年.
亞伯蘭被神呼召要動身離開他的爸爸.
離開了之後.
哈蘭還繼續有命的.
因為在經文說.
他活了二百零五年.
我們加減數就算到了.
就是說當亞伯蘭離開之後.
哈蘭仍然繼續生活多六十年.
才離世.
神呼召亞伯蘭要離開本地本族父家.
要去神指示他的地方.
在第一節裡面是這麼說的.
呼召他離開一個我們叫做「安書區」.
放下他熟悉的環境.
已經擁有的安全感.
要忘記過去有多好的生活.
去到一個新的地方.
揭開新的一頁.
有新的開始.
其實就好像很多人選擇移民一樣.
全部都放下.
去到一個新的環境.
有新的開始.
本地本族父家.
那個捨棄的範圍.
是由外面一直走進去裡面的.

$^{201}$大家看到了.
本地.
如果我們用現在的說法.
可能離開你的城市.
離開這個國家.
然後將那個範圍再收窄一點.
就是你本族.
本族就是包括了你的親戚朋友.
或者你會加上弟兄姊妹.
這些熟悉的環境.
熟悉的人群.
離開.
最後.
最大的挑戰.
是離開父家.
離開自己最親的人.
關係最密切的人.
這樣的挑戰.
是一個一個更加難.
如果是舉家.
離開.
你還會有熟悉的群體.
但是如果是你自己.
幾個人離開.
其他你很熟悉的人.
熟悉的環境你全部丟棄.
實在是一件不容易的事情.
我不知道弟兄姊妹有沒有經歷過.
如果你是曾經移民.
又回來.
那你就經驗過.
你離開的時候的不容易.
回來的時候又好像熟悉.
但又好像陌生.
的一個這樣的情況.
都是另外一種的不容易.
當.
要離開一個土生土長的地方.
離開有很深厚感情的一個群體.
最後離開自己的老家.

$^{241}$失去了很多的保護.
失去了很多的保障.
那種不安全感.
就出來了.
這就是當時亞伯蘭要去面對的.
我自己也做了一些回顧.
去年.
就是我踏入目的.
三十年.
可能有些弟兄姊妹聽到的時候.
牧師你有沒有騙人.
三十年.
我算一算真的.
踏入目的三十年.
過去這三十年裡面我經歷了什麼呢?.
在這三十年裡面.
經歷了幾次離開.
安書區的挑戰.
而每一次的離開.
都不是我自己想的.
我用了一個詞彙.
叫做被動配合.
就是教會有這個安排.
有這個需要.
我就跟著大家的決定來做.
其實每一次的艱難.
比之前一次更加的大.
讓我簡單和大家分享一下那些經歷.
第一次的離開安書區.
就是離開一個.
我由神學院畢業之後.
服侍了超過十年的教會.
這個教會.
當然最初不是我自己的教會.
但是在那個教會裡面.
我經歷了一個傳道人.
成為了牧師.
因為這個教會安納我做牧師.
這個教會裡面見證我.
由一個單身漢.

$^{281}$接著變成一個已婚的人士.
見證著我成為爸爸.
第一個的孩子出世.
所以大家想到.
應該關係都很深厚.
人生都有很重要的片段.
在這個地方發生.
所以感受很強烈.
因為得到弟兄姊妹長輩的關心和支持.
所以照道理.
應該都會繼續在這個地方服侍.
但是經過了十年有多.
接近十一年的時間.
當時我自己成長的母堂.
出現了一個很大的困難.
就是所有的教牧童工都離開了.
所以看著我成長的一些長輩.
就跟我說.
「嘉輝,你可不可以回來教會呢?」.
他說「當日你去讀神學和呼召」.
「其實就是看到教會有需要」.
「今天教會很有需要」.
「因為一個教牧童工都沒有」.
當時我已經有兩個小朋友.
所以實在是一個很大的挑戰.
自己一個人很容易.
帶著家庭要面對這個轉變.
實在是不容易的.
但是想想.
當日去回應呼召的的確確.
是因為自己教會的需要.
所以就有這個回應.
到今天這個實在需要.
來到面前的時候.
結果自己又去掙扎.
但最後當然還是來到去信服.
結果就決定回到自己的母堂.
來擔任主任牧師.
但最大的挑戰是什麼呢?.
不單止是沒有教牧童工.

$^{321}$當我決定回去的時候.
教會就告訴我們.
教會的幹事和工友都辭職了.
因為當時教會有一個很大的困難.
所以當我決定回去的時候.
我媽媽都跟我說.
「你講真的?」.
「他其實不是關心我」.
「他只是擔心孫子而已」.
所以到後來我們真的回去的時候.
當然也經歷過一段不容易的時候.
但也看到神的恩典.
到回到自己的母堂.
一個既熟悉又陌生的環境.
因為畢竟由讀神學到服侍.
已經離開了十多年.
有很多人士已經轉變了.
但也有些人當然仍然是認識.
但一切都要從頭開始.
和之前服侍的教會已經經過了一段日子.
好像建立了一些東西.
要從頭再來過.
實在是一個很大的挑戰.
結果回到自己的母堂十年之後.
由最初沒有了所有的童工.
到我再次呼召來到的時候.
加上我之內教會有三個牧師.
也有傳道人.
也有幹事和工友.
教會在人看來好像是一個很好的景況.
但也在那個時間.
神的呼召.
透過自己神學院的母校就來了.
神學院的母校有些老師.
其實都是我自己昔日在神學院的同學.
他們就來找我和我說.
「家輝,神學院有一個崗位」.
「很想你回去幫忙」.
「因為那位童工準備退休」.
我就告訴他們.

$^{361}$如果神學院的母校叫我回去.
當然好像是義不容辭.
但手掌也是肉,手背也是肉.
離開自己的母堂.
回到母校.
在人看來.
怎樣的決定應該都不會是一個不對的決定.
繼續留下來.
很適合,這是我的母堂.
不會因為這個時間.
我覺得可以很安樂地離開.
但另一方面.
母校有需要,那怎麼辦呢?.
而最特別的就是.
回到母堂之後.
神再給我一個孩子.
所以就有三個孩子.
很小.
帶著更多的家人.
要去一個新的環境當中.
實在是不容易.
而且那段時間也是香港比較動盪的時候.
經過了禱告.
也跟家人分享.
跟母堂的長輩分享.
大家都會體會到.
這不是我自己很想做的.
但也是帶著大家的禱告和祝福.
前行.
但條件是什麼呢?.
你可以去神學院工作.
但要繼續留在教會當中.
結果有一段時間.
既要在神學院工作,是全職的.
也要繼續做教會的主任牧師.
當然有其他的同工一起去支援.
但那段時間只是一個短的時間.
到了我覺得沒有辦法去兼顧的時候.
大家就說.
好,我們放你.

$^{401}$但你仍然要做顧問牧師.
你的名字仍然要在這裡.
當然,這是我自己屬人的家.
有沒有名字其實也不重要.
最重要的是.
神要放我們在哪個位置.
結果.
因為這樣的緣故.
回到自己的母堂十年之後.
就回到神學院的母校工作.
本來以為.
那個應該是時逢的最後一哩路.
因為經歷了二十多年的牧養經歷.
大概已經開始數著手指.
什麼時候退休.
但結果只是回到神學院的母校.
服侍了六年的時間.
自己的總會.
因為前任的總幹事.
他說他一定要退休了.
其中一個原因就是因為師母也退休了.
所以他不可以讓師母自己一個人等那麼久.
還有一個更大的需要.
就是因為他要照顧孫子.
所以.
這些這麼大的理由.
所以教會也決定讓他退休了.
但是在這樣的環境當中.
其實那時候就是剛剛疫情開始了不久.
去找哪些人來擔任這個崗位呢.
而當時也是面對一個很大的挑戰.
就是大家現在正在經歷的.
就是教會的移民潮.
和教務同工的退休潮.
都在這段時間出現.
結果.
大家在.
尋覓.
邀請等待的當中.
沒有任何的人來回應.

$^{441}$以往的安排.
就是由同工自己去申請這個.
總會總幹事的崗位.
但是既然沒有人回應.
大家就說不如這次我們改變遊戲規則.
不用申請.
我們用提名的方式選舉一個人出來.
那我最初的想法就是.
他們有人選出來.
那就再去選吧.
但是結果呢.
去到這個提名截止之後.
他們告訴我.
我們只找了一個候選人出來.
那個就是你了.
(嘆氣).
那我就說上帝啊.
是不是有一群人聯合起來.
來戲弄我呢.
當然我不是說得這麼文雅.
當時有一個想法.
其實很多的同工都可以做.
為什麼是我呢.
我覺得我在神學院當中.
似乎能夠發揮的.
可以說是位置.
是更加的適合.
也是我總結了二十多年的服事當中.
我覺得雖然在神學院的服事有很多挑戰.
但是是最快樂的時光.
因為其實少了很多人事的問題.
只是處理學生的問題.
是相對容易的.
但是在這樣的安排當中.
最後就說.
雖然這個提名只是一群人的一個建議.
還要看所有整個會的領袖怎麼去看.
所以說如果領袖們都同意的話.
那我無話可說.
我只能夠相信.

$^{481}$大家的領袖應該沒理由會強過我一個人的想法.
那我會馴服.
結果在一個過百人的領袖會議當中.
大家來投一個信任票.
當然最後就是現在這個結果.
也是因為這個緣故所以要離開了一個.
認為是很快樂的事奉的崗位.
而我也相信當時我的估計是沒錯的.
為什麼呢.
因為今時今日在這個崗位當中.
對我來說最大的困難是什麼呢.
真的是離開了一個超級的按書區.
原來之前那些都不算.
因為仍然都在一個很實在的目的裡面.
但是現在在這個崗位可以讓大家知道.
我看法律文件和財務報表的時間.
是比看聖經還要多.
我經常收到的是什麼呢.
就是一疊又一疊的投訴信.
裡面附上很多律師撰寫的內容.
因為在一個這麼龐大的組織裡面.
我們有教會 有學校 有社會服務.
大家就可以想像有很多的困難會出現.
而這些困難不是大家用言語去解決.
凍結也是要用法律的途徑去處理.
所以我就會問.
每一次在看這些文件我就會問.
為什麼作為牧師我要去看這些東西呢.
我要花時間去做這些審核.
為什麼會是這樣呢.
我不明白.
離開了一個按書區.
要面對這些東西.
可以祝福到別人嗎.
我好像不覺得.
為什麼和我讀的聖經好像不一樣呢.
這些挑戰當不斷來的時候.
神就讓我知道.
其實有些時候都是一個選擇.
縱然我在這個位置裡面.

$^{521}$似乎不是在做很多目的的事情.
但是.
其實我知道我在做的每一個決定.
是讓人去看見這個教會.
在這個地方.
去怎樣去見證神.
因此我就自己有一個這樣的堅持.
除了在.
一至五很多的機會要去處理一些.
我不熟悉的工作之外.
在週末.
就像現在這樣.
哪間教會邀請我.
哪怕是去到離島.
我們有堂會在大澳的.
我都一定會去港島的.
有很多人會覺得你很傻.
你很空閒嗎.
我說我不是很空閒.
因為大家知道我一部分的時間.
其實是會經常會出外開會.
我飛都飛到埃及開會了.
我去大澳的港島有多出奇呢.
所以只要大家覺得有需要.
我就必定會去應邀港島.
為什麼呢.
因為每一次港島的時候.
我才會被提醒我是一個牧者.
我最重要的本業.
就是去宣講神的話語.
如果我連這個都沒有的時候.
我都不知道.
究竟為什麼我仍然要作為一位牧師.
因為大部分的時間.
我認為自己都是不務正業.
每一次離開熟悉的環境.
熟悉的工作.
其實有很多的焦慮.
有很多的擔心.
我想大家可以想像到.

$^{561}$這一種這樣的情景.
剛剛來到去踏入牧職的那些教牧童工.
其實一定會知道.
以前你都不是做這些工作.
一踏入去這個教會的牧職裡面.
才發現我再不可以.
用一個普通的頂梓梅的身份.
大家對我的期望是很不一樣.
哪怕是在謝飯之前的祈禱.
都會比其他人更加有力量.
這個就是我們要去熟悉.
一個本來我們不熟悉的狀況.
離開了安書區有很多信心的考驗.
有很多的不安.
很沒有把握.
遇到很多無法估計的困難.
這些全部都是必然的.
但亦都是因為.
可以這麼說.
危機四伏的裡面.
我們又同時間發現有很多新的可能.
所以當回顧三十年的牧職裡面.
發現每一次.
這樣的踏出.
帶來了新的經驗.
新的視野.
讓我去知道.
怎樣來去服侍神.
不是按自己的軌跡.
不是按自己的想法.
而是按照著神所定的方向.
希伯來書第十一章第八節這樣說.
「亞伯拉罕因著信.
夢照的時候就遵命出去.
往將來要得為業的地方去.
出去的時候.
他都不知道是去哪裡」.
信心的第一步.
就是願意離開一個已知.
進入未知的當中.

$^{601}$而單單是因為.
來到去信靠那位.
引領我們的神.
走出了一個安書的環境.
向一個不熟悉的領域進發.
需要鼓起很大的勇氣.
這個正正就是亞伯拉罕所需要的.
因為他不單止離開了本地本族父家.
他帶著他的家人.
去到一個完全陌生的地方.
不知道會發生什麼事.
在這樣的環境裡面.
遠行千駛.
一定是面對危險.
在經濟上也是很大的壓力.
也是很沒有保障.
這些壓力和困難.
即使是一個年輕人.
已經是不容易應付.
更何況.
大家記得剛才提過.
亞伯拉罕75歲.
他要來應付這一切.
是更加的難.
已經有差不多的年紀.
有自己的家室.
本來跟父家一起.
他享有很多很多的東西.
有很多的財富.
有服侍他們的僕人.
完全是沒有問題的生活.
但是現在離鄉別井.
重頭來過.
去到一個失去了.
穩妥生活的環境.
不知狀況的環境.
勇氣.
實在是需要很大很大.
勇氣聽來是一個很高尚的情操.
在風平浪靜的日子.

$^{641}$其實是說得很輕鬆的.
但是當考驗來到你面前的時候.
其實你就不想這麼有勇氣.
你希望我只不過是一個普通的人.
我想到的就是要避開一些危險.
想到的就是自保.
遇到危險的時候.
其實大家會怎樣來回應呢?.
我想不需要想一些很遙遠的地方.
想想.
過了幾年.
其實疫情不是過了很遠.
在疫情的時候.
當很多人搶物資的時間.
大家有沒有一起去搶呢?.
你怕不怕家裡沒有米呢?.
怕不怕沒有廁紙呢?.
我們過了的時候.
當然覺得沒有問題.
但是在正直的環境的時候.
其實我們跟其他人沒什麼分別.
大家去搶的時候.
我們也去搶一份.
哪怕是基督徒.
原來在這樣的日子當中.
我們都忘記什麼叫做勇氣.
真正的勇氣其實並不是無所畏懼.
我們也會有很多的擔心.
但是真正的勇氣.
就是在恐懼裡面.
我們仍然去選擇信服神.
不只是顧自己.
也會去顧念其他人.
而昔日的亞伯人.
就是展現這種勇氣.
因為他知道.
他去承受的不是自己的祝福.
他準備要成為其他人.
成為萬國的祝福.
在舊約裡面.

$^{681}$土地賦予了人一個很重要的身份.
因為地是屬於神的.
如果他去到那個地方是神要他去的話.
就證明他是神的子民.
但是亞伯人面對一個挑戰.
就是要離開那個本來是神給他的地方.
去另一個是神給他的.
當一個離開在那個過渡的時候.
他就失去了身份.
他會不斷去懷疑究竟我是誰.
去到外面的時候.
沒有人再認識他.
他的名聲,他的地位.
沒有人再去理會.
他建立了的影響力.
也要連根拔起.
所以的確是一件不容易的事.
我們看到.
有很多神所忠心使用的僕人.
他們回應呼召成為了宣教士.
其實他們正正就是要有這種勇氣.
就像剛才所說.
不是說無所畏懼.
去到一個陌生的環境.
語言又不通.
也不知道會遇上什麼事.
但是只知道一件事.
在那個地方有福音的需要.
所以他們就離開一個熟悉的環境.
離開很多支援的系統.
就會踏上宣教的歷程.
因此,兄弟姐妹.
我鼓勵大家多些關心宣教士的同工.
為他們禱告.
也要有實質的支持給他們.
否則他們所走的路.
實在是不容易的.
我們怎樣才可以有勇氣.
以致我們願意離開一個熟悉的環境.
繼續向前進發呢?.

$^{721}$當然,避缺.
當然要從認識神開始.
亞伯蘭.
他有勇氣.
就是因為他認識神.
並且信服神的安排.
在這個信服的過程當中.
他會去相信一件事.
神的應許.
神對他說過的.
一定是會實現的.
但是這個實現不是馬上出現的.
就不像現在因為馬年.
什麼人都說馬上就什麼.
對不起.
基督徒很少會看到馬上就什麼的.
我只想到一件事.
我們吃飯之前.
大家想馬上開餐也沒有說這麼多.
但是其他.
特別是神的祝福.
不是馬上就得到的.
可能要等.
等完再等.
都未必會馬上出現.
這就是一個信心的考驗和挑戰.
亞伯蘭放下這麼多的東西.
他失去了原本有的地方.
但是其實是神要為他和他的後代.
預備一個更大的土地.
讓他們成為一個大的國家.
得到更多的地土.
他失去了本族.
因為要離開自己的親友.
支援的系統.
但是神給他的應許.
就是讓他得到更多的子孫.
多到一個地步.
是無法數點變成了大國.
他失去了他的父家.

$^{761}$離開了他的爸爸.
離開了最大的保障.
但是神卻成為了他的父親.
這就是亞伯蘭的經歷.
要信得過的不是眼前的東西.
是信得過神所說的話.
以致他就前行.
我們怎樣才可以和神有這樣的關係.
怎樣才可以有這樣的信心呢?.
其實這也是過去我自己.
一直在問自己.
怎樣才知道我這個信服.
這個決定是真的對的呢?.
真是神想我要的呢?.
以往可能大家去想.
不外乎聽到牧師都是說.
讀經,靈修,禱告等等.
沒錯.
不過這些只是方法.
方法我們是有的.
而最需要最缺乏的是什麼呢?.
兄弟姐妹.
我們不是缺乏方法.
通常我們缺乏的是時間.
誰不知道要讀經.
誰不知道要靈修.
誰不知道要禱告.
但結果為什麼不做呢?.
因為我們覺得我們沒有時間.
沒有時間靈修.
沒有時間讀經.
沒有時間禱告.
不過有時間煲劇.
有時間去做其他事情.
我見過一些青年人告訴我.
哎呀!我很忙啊!.
我問他們忙什麼?.
因為都是在職的青年人.
很忙啊!.
因為下班後要去健身.

$^{801}$去健身.
我不知道我們丹管城堂的年輕人是否喜歡.
我見到教會的年輕人.
包括我兒子.
每個人走過去摸摸手臂.
都玩得很緊.
有時間去操練自己的身體.
但你告訴我沒有時間去做其他事情.
當然他們都很乖.
他們都有崇拜.
他們都有團契.
但他們擺放的時間在操練身體多.
還是擺放時間在與神建立關係上多呢?.
所以兄弟姐妹.
我們不是沒有方法.
我們只不過是忘記了.
我們用了藉口說沒有時間.
當沒有時間的時候.
什麼都可以不用做.
也沒有辦法做.
所以最初我們去想的.
並不是用什麼方法去建立.
與神一個緊密的關係.
怎樣才可以知道信服神的旨意.
其實就是我們先要預留起時間.
給神與我們一起.
而這也是我自己的堅持.
每一天當其他的同工.
還沒有回到辦公室的時候.
我會比他們早一點回來.
為什麼呢?.
因為我知道他們一回來的時候.
我就會很忙.
因為每個人都要來找我.
要簽名.
要問怎樣去做這些決定等等.
所以在一個最不受騷擾的時候.
那個就是我與神的時間.
如果我不預留時間與神一起的時候.
其實我知道那天我沒有任何事情可以做.

$^{841}$哪怕我有很多方法.
面對著我不熟悉的法律文件.
其實也不需要太過擔憂.
我首先現在可以找到AI幫我總結一下.
究竟那封信想說什麼.
再不熟悉的話.
我隨時都可以聯絡到我們的法律顧問.
他都很樂意去解答的.
都不用自己去想的.
但反而時間是其他人給不了我.
我要自己來知道哪一件事是更加重要的.
所以要留起時間來給神.
我也相信在聖經裡面.
有很多的人物同樣.
他們知道哪一件事是更加重要的時候.
就要留時間給那些人.
弟兄姊妹也一樣.
你會留時間給你的家人.
你會留時間去做一些你很想做的事情.
但你留下多少時間給神呢?.
神應許.
啞巴爛.
地上的蠻族都要因祂而得福.
當我們擺正了那些次序的時候.
我們知道我們的生命就不再一樣.
啞巴爛願意去冒險.
不只是因為要自己得到祝福.
他盼望要成為蠻族的祝福.
我們常常都會祈求神祝福我們.
我相信這也是我們心想的.
我們一定很想神祝福我們.
但我們有沒有想過.
神怎麼可以讓我們成為別人的祝福呢?.
當我們轉向去想想.
如何使其他人得到祝福的時候.
可能我們會更加願意.
又或者會鼓起更大的勇氣.
來離開一個安殊的環境.
這個祝福.
剛才一開始的時候說過.

$^{881}$不一定是在物質上.
但當你看到.
原來你願意地去踏出.
神就來使用你.
出現了一些改變.
我們就知道.
原來神透過我們.
可以使其他人得到這個祝福.
自己也會得到一個很大的滿足.
這個決定必然是冒險的.
在生活中的穩定是會受到挑戰的.
有很多未知的可能性會湧到你面前.
但之後你回望.
其實是得到一個更大的祝福.
我自己有一個很深刻的體驗.
這裡只能夠和大家簡單地說.
大家應該也會對弘福縣的大火有很深刻的印象.
在這個大火開始不久之後.
有很多街坊湧到我們會的學校.
因為學校就正正在弘福縣的前面.
當時我和整個會的一些長輩.
正在內地的教會交流.
而這間學校的校監.
就和我一起在這個旅程中.
我們收到校長很緊張地打電話來找校監.
校監說有很多街坊湧下來.
因為我們後面的屋苑著火了.
當時我們不知道是這麼嚴重.
因為我們不在香港.
我們只是聽到有這個消息.
因為同一團的團友的媽媽就住在那裡.
但幸好她收到一個訊息.
她媽媽的工人很機警.
一知道有這件事的時候.
但聽不到有火警鐘.
就帶著老人家走了下樓.
下樓之後就看到他們自己那座都燒了.
所以嚇了一跳.
當時校長很擔心.
他說有這麼多街坊湧來的時候.

$^{921}$我們不知道怎樣來處理.
我們說不要緊.
雖然是學校.
但我們學校裡面也有教會.
教會和學校的同工一起來接待那些街坊.
讓他們的家人也逃出來.
然後再過晚一點.
校長就說.
如果繼續這樣下去.
我們明天就不能上學了.
他問過教育局.
不可以就這樣自己決定不上學.
他一定要有校董會通過同意才可以.
但明天就可能不能上學了.
那怎樣召開這個校董會.
當時校監跟我說怎麼辦.
我說有什麼怎麼辦.
我說這個是災難.
我們不要想太多.
我就跟同學說明天一定上不了課.
總之大家要保持聯繫.
我們要去接待那些街坊是首要的.
如果教育局有人問.
有人去質疑挑戰等等.
麻煩直接來找我.
我一定會跟他們說.
當然校監也要做事.
我們在今時今日.
去到天腳底也能聯絡到.
所以很快所有校董都同意學校要停課.
我們自己就來決定.
當然後來教育局告訴我們.
不需要做這些.
因為這個情況特殊.
所以在附近的學校全部都要停課.
後來就已經處理了.
當然我們也跟校監說笑.
我說其實我做這個決定很輕鬆.
為什麼呢.
我說因為學校有什麼事.

$^{961}$坐牢的是校監.
所以我說你做了這個決定.
那麼多事他拔了.
總幹事你聽著整個過程.
如果有事的時候你不要丟下我.
但我們知道我們做了這件事.
其實是使得那些街坊.
得到一個短暫的地方.
可以安頓下來.
當然在後續也有很多不容易的情況.
但我們也盼望那個短暫的幫助.
是可以使得一些等待幫助的人.
不需要再太過徬徨.
所以因為這樣的經歷更加讓我知道.
叫我們在每一個位置.
其實都要有他自己的心意.
雖然災難的發生是沒有人想的.
但我們就有這樣的機會來服侍論社.
弟兄姊妹今天的世界.
教我們都是追求安定的生活.
可能這個也是我們很多人很想的.
但神的路.
往往都是叫我們帶著很多冒險和不確定.
但為的是什麼呢.
為的就是要讓我們可以承受祝福.
而其他的人因為我們這樣的緣故.
都會得到祝福.
但如果你不踏出這一步的話.
我們就看不見前面會發生什麼事.
所以願意我們都是可以得到.
從神的話語裡面給的提醒.
以致我們有那個勇氣.
因為我們知道無論什麼環境.
其實神都是和我們一起.
也讓我們去保持和神有一個緊密的關係.
以致我們聽得清楚.
祂給我們的呼喚.
也知道不是我們自己的意思.
而是祂的意思.
當我們去懷著使其他人可以得到祝福的使命.

$^{1001}$而不是顧自己的時候.
我們就會看到有很多很多的可能性出現.
求主幫助我們.
在一個充滿挑戰的當中.
我們不懼怕.
在有很多不確定的裡面.
我們能夠去確定的就是.
神是那位要去祝福而且永恆不變的神.
求主憐憫,示恩.
阿們.
\newpage



\section{創世記 24:56}
\label{sec:031728}
\textbf{聖經人聖經事:不要耽誤我}~~~~黃克勤牧師博士
\newline
\newline
連結: \href{http://www.church.com.hk/acms/content.asp?site=cdc\&op=show\&type=product\&code=031728\&layout=sermon}{\texttt{031728}} ~~~~ 語音日期: 2026-02-19
\newline
\newline
\hyperref[sec:031727]{\small{< < < PREV SERMON < < <}}
~
\hyperref[sec:index_chronic]{\small{[返順時目]}}
~
\hyperref[sec:031729]{\small{> > > NEXT SERMON > > >}}
\newline
\newline
創世記 24:56
\newline
\begin{longtable}{cl}
\hline
\hline
章節 & 經文 (和合本修訂版)\\
\hline
24:56 & \begin{tabularx}{0.7\textwidth}{X} 僕人對他們說：「耶和華既然使我道路亨通，你們就不要耽誤我，請讓我走，回我主人那裡去吧！」 \end{tabularx} \\ \\
[1ex]
\hline
\hline
\end{longtable}
$^{1}$各位電視姊妹,又讀聖經了.
時間過得很快,好像眨眼一天一樣.
眨一眨又一天一樣.
時間過得很快的.
新聞說書教網絡又為大家講聖經了.
聖經人聖經時最重要是看到背後上帝怎樣工作.
好了,今天的主題.
不要耽誤我.
不要令我遲延,不要阻礙我.
不要耽誤我.
好了,又是那件事.
就是阿伯喊取身好的續篇.
那個僕人,那人,聖經叫做.
將一個金環戴在鼻子上.
然後呢,怎麼戴啊?立即扭鼻子嗎?.
不是,當時很多女子都已經穿了個洞.
好像很多女子已經在耳朵穿了個洞.
隨時可以戴了.
穿了個洞,就戴了個金環.
另外就給她戴了一對金手鐲.
然後問:你是誰的女兒?.
家裡有沒有地方住?我去借宿一下.
有啊有啊,就跑回家了.
女子,創世記二十四章二十八節.
女子跑回去照著這隻蛇話告訴他母親和他家女的人.
很有趣,喜歡跟母親說.
特別是利伯加里卡,他講密家語拿學.
就是我爺爺那輩,都是女的先行.
然後跑回去,就將這隻蛇話告訴他母親,而不是父親.
有趣嗎?可能父親未必在家.
或者出去了,或者已經不在世.
他母親和他家女的人.
利伯加里有一個哥哥.
在這段經文,可能彼得利不是很主動.
或者不在家,或者有種種因素.
現在媽媽與哥哥在家裡說事.
利伯加里有一個哥哥,名叫拉班.
大家聽過了,拉班就是白色的意思.
白色是顏色來的,可能拉班是比較白的.
出世後,有一對兄弟在《聖經》裡很出名的.

$^{41}$就是以撒兩個兒子.
爭著出世的這個媽胎.
大兒子以素,很多毛,而且紅的人.
紅色的意思,叫做以東,以東就是紅色的意思.
以素就是毛的意思.
雅各爭著出世的意思.
拉班就是雅各舅父.
媽媽就是利伯加里.
拉班就是白色的意思.
看見妹妹走回來.
一走進來就看到她的臉.
這是金環的目的.
一看就看到她的金環,很富有.
很漂亮.
又看見她的金竹在她妹妹手上.
聽到利伯加里的妹妹這樣說.
那個人怎樣怎樣跟我說.
於是拉班就說,嘩!這件事很嚴重.
於是跑出去,往井旁去.
拉班雖然也是信耶和華的.
不過我怕他對金錢也挺有興趣的.
嘩!有錢的.
拉班就跑出去到井旁看看發生什麼事.
一看見他的女兒有金環與金手竹.
十四赫勒,很多錢,是金的.
跑出去看見駱駝等等.
對那個人說.
你借蒙耶和華賜福,請進來吧.
為什麼站在外面?.
嘩!很會說話,整家人都很厲害.
我已經收拾了房屋.
也為駱駝預備了地方.
拉班跑出來的時候已經吩咐工人.
預備好地方,有十隻駱駝.
有人來住,整頓一下.
這是貴客.
我們讀一下第三十二節.
那人進了拉班的家.
拉班卸了駱駝,用草料餵上拿水給那人.
和跟隨的人洗腳.

$^{81}$這裡叫拉班的家.
為什麼不是彼土利的家呢?.
拉班的家是哥哥做了主.
彼土利.
可能爸爸彼土利年紀大.
亞伯拉罕都這麼大年紀了.
彼土利應該不小了.
可能退休或不太活動了.
已經是拉班哥哥出來作主.
母親的家.
爸爸年紀大一點,太太年紀小一點.
於是哥哥拉班出來作主.
母親也作主.
爸爸彼土利年紀大也不少了.
跟亞伯拉罕差不多了.
拉班的家.
拉班作主.
可能很忙.
卸了駱駝,用草料餵上拿水給那人.
和跟隨的人洗腳.
可能那人很想說.
拉班先生我想說.
等等,我們餵了駱駝.
等等,我們快開飯了,先坐一會.
等等,我先用水洗腳.
很忙.
那人唯有去別人家作客.
忍一忍.
終於想說.
我們先開飯.
開完飯,坐下先吃.
怎知那人不再等了.
第一點想說.
我不吃,等我說明白我的事情再吃.
拉班說,好,請說.
到這裡才能說.
其實那人已經等了很久想說這件事.
他一心來這裡就是做這件事.
但又這樣,又這樣,又這樣.
那麼多事情,那麼忙.

$^{121}$到吃飯了,先不吃.
其實你猜那人肚不肚餓呢?.
我推測他也很餓.
不過不要給.
我們做事也是.
不要給第二重要的.
取代第一重要的.
第一重要是什麼?先搞好一件事.
然後才搞第二件.
這件事重要我先搞,然後第二第三.
第一重要是什麼?我想談一件婚事.
所以他說我不吃.
雖然我已經餓到鼓聲.
幾乎血糖都低,快暈了.
不過我先說.
我不吃,等我說明白我的事情再吃.
辦事要知所先後,輕重.
這是重要的.
我這麼老沒讀過書.
沒讀過管理學.
不過很厲害.
現在如果你願意.
以慈愛誠實待我的主人.
就請告訴我.
願以慈愛誠實待我的主人.
就請告訴我第二點.
如果你的女兒肯嫁給我.
你的妹妹肯嫁給我的主人的兒子.
嫁給我做我的少奶.
那麼你告訴我.
如果不行都告訴我.
光明正大,不要吃飯先想一想.
今晚明天再說,後天大後天.
那麼他沒有.
現在告訴我,行不行.
不行都不要緊.
我再向左向右再找再想.
就告訴我.
那麼你見到他很單刀直入.
告訴我,行不行.

$^{161}$那麼我才吃得下.
拉班和彼土利.
彼土利的意思是白色.
彼土利的意思是神毀壞.
這個利是EL.
是EL這個字.
EL是God.
彼土利是神毀壞.
或者神所毀壞.
神所毀壞.
神毀壞.
那麼你就回答他.
這件事既然是這樣.
爸爸可以出來說.
這件事是出於耶和華.
拉班也說是.
彼土利也說是.
這件事出於耶和華.
我們不能夠向你說好說歹.
我們不敢說什麼.
我不敢反對.
這樣的意思.
那麼就看看李伯嘉在你面前.
你就將他帶去.
照著耶和華所說的.
給你主人的兒子為妻.
爸爸拉響的僕人聽見他們這樣說.
就向耶和華俯伏在地.
不是向他父子俯伏在地.
一聽到響.
又俯伏在地.
見到耶和華.
這個僕人很敬虔.
當下僕人拿出金器銀器和衣服送給李伯嘉.
要給禮的.
又將寶物下訂.
真是過大禮.
將寶物送給他哥哥和他母親.
真有趣,又沒有送給爸爸.
是送給哥哥和母親.

$^{201}$真有趣.
僕人跟隨他的人吃了喝了.
那晚很開心.
住了一夜.
第三點.
住了一夜.
早晨起來.
僕人就說.
請打發我.
請容許我回主人那裡.
睡了一晚.
走了一個多月.
最少一個多月.
住了一晚.
成不成功?.
成功.
第二天就說.
我回去了.
不是吧.
如果是我.
住了幾晚.
先吃宿.
先歇一歇.
先呼吸一口氣.
每天騎駱駝.
你覺得好騎嗎?.
我也騎過.
我雙腳張開到.
辛苦到不得了.
突然雙腳還僵硬.
走很遠.
因為騎駱駝的駱駝背很闊.
一點也不開玩笑.
如果上面沒有椅子.
沒有艇.
真的騎不到.
很辛苦.
而且老人家.
這個僕人很老了.
不遠千里.

$^{241}$但是不留.
第二天睡了一晚.
可不可以讓我走.
我想回去了.
利伯加哥哥和他母親就說.
哎呀.
你一來到就帶他走.
讓女兒和我多住幾天.
最少十天.
讓我可以和女兒慢慢聊天.
僕人就說:耶和華既次給我通達的道路.
你不要耽誤我了.
請打發我走.
回我主人那裡.
你們不要耽誤我.
嘩!真是厲害.
你看到這個僕人厲害在哪裡呢?.
這麼遠.
千里迢迢來到.
只說我住一晚.
完成任務.
第二天起床就走了.
嘩!這樣也不拿多少早餐.
不吃幾天東西.
睡一覺.
不是吧?第二天就走了.
這個僕人就是這麼忠心.
你們不要耽誤我了.
不如讓我走.
於是十天都不等.
如果我呢.
我當時就說.
哎呀!這樣也很合理的.
讓女兒和母親多住幾天.
十天吧!不要緊.
十五天都可以.
二十天都不拘.
我在這裡辛苦一點等你們.
其實吃飯喝水睡一覺.
二十天不奇怪.

$^{281}$當時的人.
到耶穌時代.
記不記得他在水面走.
喂!那些人.
耶穌時代來喝酒.
走得遠.
都來七天的.
現在來第二天就走.
他真的覺得不行.
起碼十天八天吧.
但這個僕人完成任務深切.
不要耽誤我.
這個其實是我們要學的.
很多時候我們會被其他.
都很好的東西.
耽誤了我們最重要的一件事.
你不要耽誤我了.
讓我走吧.
於是就回去了.
他也不敢耽誤她.
多好啊!今代.
就是僕人說話.
既知我通達的道路.
你們不要耽誤我.
請打發我回去主人那裡.
請你記得.
人生的路上很短的.
不要兜兜轉轉了.
不要耽誤我.
我要完成我的任務.
盡快完成.
這個僕人很厲害.
如果今天機構工作.
或者辦事要請人.
公司請到這種僕人就厲害了.
可以是沒有學位.
沒有讀過MBA.
不是DBA.
不是Doctor of Business Administration.
不是管理學博士.

$^{321}$不是管理碩士.
不是管理學士.
什麼都沒有.
不過很有智慧.
很敬畏耶和華.
懂得聆聽.
然後很專心一致地做.
一做完就完.
馬上走.
很直接地當你做這件事.
想好了格局怎麼擺.
哇!真厲害.
這個僕人沒有名字.
不過很厲害.
他很專注.
專注在上帝的帶領.
專注在事情怎麼做.
他想好了對白.
第三.
一成功.
專注馬上走.
回程.
這個僕人真厲害.
我們辦事都是要這樣.
服事神都是要這樣.
另外看到.
神在背後不停地工作.
所以耶和華說.
記著我通達的道路.
我就不要耽誤了.
「哎呀!神給我的路很通」.
「反正通了,慢慢走」.
沒有.
我不耽誤.
馬上走.
就是神給我通達的道路.
馬上走上.
這樣就好了.
希望我們都學到怎麼辦事.
和仰望耶和華.

$^{361}$是賜我通達的道路.
我們祈禱.
多謝主.
我們看到這件事.
好像是一件人世間.
要娶媳婦的事情.
不過看到的是無名的僕人.
還是很老的.
原來是一個這麼出色的僕人.
他做得這麼好.
而且很敬畏神.
每一步都仰望你的帶領.
他又時時向人向耶和華.
俯伏敬拜.
他會欣賞上帝.
讚美上帝.
奇妙的作為.
他真是一個清心的人.
看到主自己的工作.
讓我們學習他.
更多認識你自己的能力.
讓我們勇於完成你所交付的任務.
在路上不耽擱.
就是完成任務.
瞄準著目標.
多謝主.
你賜他通達的道路.
我們同樣祈求通達的道路.
讓我們不耽誤的.
為主完成任務.
聽我們祈禱.
奉靠耶穌基督命求.
阿們.
\newpage



\section{利未記 1:1-2}
\label{sec:031729}
\textbf{從祭壇到心靈奉獻}~~~~羅泰然牧師
\newline
\newline
連結: \href{http://www.church.com.hk/acms/content.asp?site=cdc\&op=show\&type=product\&code=031729\&layout=sermon}{\texttt{031729}} ~~~~ 語音日期: 2026-03-01
\newline
\newline
\hyperref[sec:031728]{\small{< < < PREV SERMON < < <}}
~
\hyperref[sec:index_chronic]{\small{[返順時目]}}
~
\hyperref[sec:031730]{\small{> > > NEXT SERMON > > >}}
\newline
\newline
利未記 1:1-2
\newline
\begin{longtable}{cl}
\hline
\hline
章節 & 經文 (和合本修訂版)\\
\hline
1:1 & \begin{tabularx}{0.7\textwidth}{X} 耶和華從會幕中呼叫摩西，吩咐他說： \end{tabularx} \\ \\ \relax
1:2 & \begin{tabularx}{0.7\textwidth}{X} 「你要吩咐以色列人，對他們說：你們中間若有人要獻供物給耶和華，可以從牛群羊群中獻牲畜為供物。 \end{tabularx} \\ \\
[1ex]
\hline
\hline
\end{longtable}
$^{1}$當我去準備這次講道的時候.
就想起很小的時候的一些事情.
小到可能是讀幼稚園之前.
或者幼稚園的時間.
家裡的客廳裡面掛了一張耶穌的像.
掛在那裡.
我記得我小時候一起來.
我就去到那個像前向耶穌請安.
家裡人都覺得很好.
有時候我起來就提醒我這樣做.
我記得有一段時間是這樣的.
你不要管那些什麼聖像敬拜.
或者好像別人裝香這樣拜.
但我覺得這些情況就令我想起.
有一個習慣.
就是一起來就記得一個主耶穌.
當然我家人都是信主.
只是放在那裡好像一個裝飾.
不過我就好像一個請安的感覺.
其實我一想起就習慣了.
一起來就找主耶穌.
記得這位神在當中.
當我們講今天這段經文.
特別我們看李美記講獻祭的時候.
一想其實這些獻祭都很有意思.
因為這些獻祭是神所定的律例典章.
要讓人遵守又要實踐.
以致他們按時按後.
或者是不同的情境裡面都很好的獻祭.
包括他們應該要怎樣準備.
包括他們在什麼情境底下.
但是今天對於我們新約教會來說.
我們什麼時候叫獻祭呢?.
你說是不是翻崇拜叫獻祭呢?.
禱告是獻祭呢?.
但是無論如何有時候我們都不是定時定候.
祈禱都未必很定時地去禱告.
大家有參加40天祈禱.
今天祈禱了沒有?.
有時候不是說今晚臨睡前記得.

$^{41}$還沒做,立刻立刻先祈禱.
有時候我們都會這樣的情況.
或者有時候我們翻崇拜.
我們都是今天天氣又潮濕又冷.
不太想起床,可能沒回來.
或者有些緣故有些因素很容易撇開不做.
但是在舊約的儀式的獻祭裡面.
要求是很高的嚴格要求.
所以他們不可以失信於神.
他們要做好.
當然如果整體來說祭司代表百姓去獻祭.
在會目的獻祭或者在聖殿的獻祭更加重要.
但是他們平時自己去獻祭也可以.
當然去到祭司的面前.
今天和大家去看.
就是看到從舊約你還沒記.
特別一至七章裡面有五個的獻祭比較主要的.
是去幫助我們體會.
從一種祭壇裡面的獻祭.
提到我們心靈裡面的獻賞,奉獻.
甚至金錢的奉獻.
希望做一個比較一連串的去看.
我們更加明白當中更多的意義.
也會在新約的教會裡面.
究竟我們怎樣的態度.
怎樣去實踐這種的奉獻.
甚至是全人的奉獻.
金錢的奉獻等等.
所以今天的題目就是從祭壇到心靈的奉獻.
分開四方面.
一方面是奉獻是生命獻賞的延伸.
第二方面是奉獻是感恩的表達.
第三方面是獻金不是贖罪的手段.
第四方面留意這些字.
「金心獻賞」.
你記不記得一到七章是記載了五種的獻祭.
凡祭.
掃祭.
平安祭.
贖罪祭.

$^{81}$贖遣祭.
我都聽到有人回答了.
五個主要.
當然還有其他什麼「遙祭」.
什麼其他不同的祭.
但這個主要的五個祭.
雖然我們今天不是再獻一些牛羊.
不去收成.
但這個獻祭的精神.
是真的要繼續啟發我們今天奉獻的生活.
當然奉獻生活不單單是說金錢的方面.
金錢的奉獻.
只不過是財務上的運用.
但是生命,感恩,復和,敬拜.
這個真的是奉獻的延伸下來.
跟我們整個人,全人都有關係.
究竟神為什麼要百姓獻祭呢?.
如果在舊約裡面.
我會覺得有多方便.
第一.
大家知道當我們要獻祭的時候.
我們每個人.
無論是自己.
或者獻出來的祭身.
或者那些供物.
即是那些收成.
都要怎樣?.
都要潔淨的.
都要沒有殘疾的.
都要新鮮的.
這裡看到是教導我們聖潔.
因為神是一個聖潔的神.
所以我們去到神的面前獻祭.
我們用聖潔的一個.
我們整個人.
無論在心態上.
我們自己的行為上.
是一個聖潔的人.
去到神的面前.
所以崇拜的開始也是一個獻祭.

$^{121}$崇拜是獻祭.
開始裡面主席都會帶領我們.
獻上我們認罪的禱告.
我們潔淨我們自己.
提醒百姓.
神是聖潔的.
來到神的面前的人.
是要守潔心清.
第二方面.
獻祭當然是認罪和赦免.
這個聖潔之餘.
如果你真的有些地方得罪神.
得罪人的時候.
你都要在神面前.
藉著獻祭.
能夠挽回自己的罪.
能夠得到神的赦免.
又顯明了神是對人的憐憫.
所以這個就是其中的意義.
第三個意義就是感恩和敬拜.
獻祭也是回應神的恩典.
承認一切都是來自神的.
所以我們敬拜裡面.
我們都會不忘有感恩.
當我們主席和我們祈禱.
一起的時候都不會忘記.
這位創遷造地的神.
我們一切的好處.
都是來自神的.
我們從第一方面去看.
奉獻是生命獻上的現身.
首先講繁祭.
你未記得第一章是全部講繁祭.
我只是抽一支經文給大家看看.
第一章九節那裡說.
祭司要把一切燒在壇上作為繁祭.
是獻給耶和華的火祭.
是馨香的祭.
什麼叫做繁祭呢?.
繁祭的意思就是.

$^{161}$本身是完全焚燒的意思.
也有一個意思是上升的意思.
大家想想什麼叫做上升.
當你的祭品.
特別是這些全部都是動物.
它完全燒的時候.
會有什麼?.
燒起來的煙.
一路上升.
煙就會向上升.
所以這些祭品.
繁祭就是把動物燒盡.
所以所帶來的祭品是.
公牛,公面羊,山羊.
或者是單溝,戳甲等等.
這些動物和鳥類.
那做法剛才講了.
就是把傳言燒盡.
那皮呢?動物的皮就可以分拆出來.
可以給祭司.
那其他的就全部燒盡.
這個你看到沒有留一塊肉.
叉燒給你吃.
那就不會了.
那傳言就全部燒了.
這個就表明.
所有的東西都獻給神.
象徵是完全的奉獻.
和一個輕香的火祭.
你想想那些肉在燒有些香味.
你燒東西吃就知道.
那些味道你就不會流口水.
全部燒完.
最後可能是燒焦的味道.
這個就是當時的意思在裡面.
好了,去到新約裡面.
我們怎樣去對照呢?.
就是剛才我們讀過.
《羅馬書》十二章一節那裡說.
弟兄們,我以神的慈悲勸你們.

$^{201}$將身體獻上,當作活祭.
真是用身體才能夠表達.
剛才那個動物全件地燒掉.
我們心靈獻上就行了.
或者我們的財富獻上就行了.
身體獻上就表示你整個人.
是給了神的.
毫無保留的.
你再想想身體獻上.
就是身體都沒有了.
即是命都沒有了.
即是你就是這樣.
是全然毫無保留.
完全給了神.
是聖潔的.
是神所喜悅的.
當然這裡又是說聖潔.
即是又表明了.
剛才說了.
我們都要去神的面前.
要潔淨自己.
所以這個繁祭是提醒我們.
奉獻不是我們一部分.
不是我們某些的部分.
是整個人的生命.
毫無保留地獻上.
這個是當然我們整個概念上意識形態.
我們知道我們屬於神的.
我們是被神所使用的.
神賜給我們.
神要收取.
都是神自己的意思.
所以表明我們整個人是屬於神的時候.
你就會覺得.
你毫無保留地獻上.
是不會覺得.
是一件很困難的事情.
受靈的意思當然是說主我整個生命都是你.
你口頭上這樣說.
當然包括了我們的時間.

$^{241}$金錢.
我們的財幹.
我們的精力.
都是這樣獻給我們的神.
所以這個是生命的奉獻.
但是有時候.
我們都會在我們的生命裡面很多時候是工作.
在工作裡面我們會賺取一些金錢.
所以去到一個金錢的奉獻來說.
它代表了又可以用金錢.
今天的新約教會.
去代表了一種的態度在裡面.
通常我們的常年奉獻.
就會這樣給我們理解.
因為這個梵祭是恆常的奉獻.
是常常表達我們自己一種的心智的.
所以常年奉獻都可以說是一個.
我們來說.
一個金錢奉獻是常常表達我們一種的心智.
當然.
這個後面我再說.
究竟心重要一些還是金重要一些呢?.
待會再說.
但是我們要明白這件事.
好了,到了掃祭.
掃祭就是二章一節那裡這樣說.
有人獻掃祭為供物給耶和華.
用細面硝上油加上乳香等等很多條例.
掃祭是什麼意思呢?.
掃的本身就是穀物的意思.
植物.
字面的意思就是供物.
一個禮物給神.
這些是細面,乳香,無教餅,初熟的和水等等.
這樣獻給神.
方式都是燒.
燒在壇上裡面.
不是全部燒掉.
有些部分是給祭司的.
但是做的時候不可以有教.

$^{281}$也不可以加墨.
要加鹽.
但是主要的含義是什麼呢?.
你想想他們有這些的修成.
他們會經過什麼?.
一段時間的什麼?.
耕種.
聖經有說.
流淚,殺腫,歡呼,修葛等等.
就是說.
他們有些成果出來.
勞動了之後的成果.
就獻給神.
所以這個掃祭的意思就是說.
他們在生活上.
他們的生活,勞動,工作.
都要藉著生活,勞動,工作來獻祭給神.
去敬拜神.
所以這個就代表我們也是恆常的獻上.
是獻上我們勞力所得來的成果.
象徵我們日常生活裡面.
都要常常敬拜主.
哥林多前書十六章第二節說.
每逢七日的第一日.
每人要照自己的收入抽出若干.
保留起來免得我來時賢酬.
這裡當然是說一種金錢的奉獻.
當時可能七天就會有新金.
有些人是收周薪的.
有些是日薪的.
有些是月薪的.
但是現在說的就是每天七天.
就是主日的時候.
我們來到神的殿裡面.
我們都可以獻上.
你看到照自己的收入抽出若干.
意思是你的收入從哪裡來的?.
是你的勞力,你的工作.
這個在新教會裡面來說.
除了收成可以獻給神.

$^{321}$但是今天.
如果你帶一些禾菌,穀物來尖浸奉獻.
我們就很笨了.
不知道怎麼使用.
當然我知道有些內地教會是這樣的.
他們的農村教會.
真的是一籮冬瓜.
當是一種情況.
知道傳道人也不是很夠薪水.
但是希望收成裡面能夠幫助他們.
或者教會裡面有需要.
但是在今天一般的教會來說.
我們用了金錢代替了.
但是我不說金錢之先.
其實金錢可以代替很多事情.
就是錢可以做很多事情.
但是這個不是完全是這樣說.
但是我們要明白一件事.
你的勞碌的奉獻包括了什麼?.
包括你的事奉.
包括了像師班.
他們今天跟我們一起獻詩給神.
他是要有勞碌的.
禮拜幾唸詩?.
禮拜四晚唸詩.
每一首詩歌不是唸一次的.
是幾個禮拜一個月前慢慢唸下來的.
慢慢唸是慢慢唸的.
我講師班是因為有背後.
其實領師也是的.
剛才韋生執事也說他很怕領師.
也要很準備.
其實每個人都是的.
做主席.
做招待.
或者教會的事奉.
都是你們擺上自己的.
這個是很欣賞的.
是勞動裡面的.
有些人說我都沒有參與事奉.

$^{361}$當然了.
在這裡來說.
這些的金錢事奉.
又可以給我們提一個方向.
因為有些地方.
譬如這樣說.
教會有發展我們的擴展聖功的工作.
現在沒有了擴展聖功的獻金.
我們分開了.
不過大家未必知道的.
就是猜拳事奉獻.
和完善旁所的設備.
這兩樣都是繼續對教會的工作有發展的.
猜拳工作.
我們去到外面.
無論我們的旅費.
我們的幫助教會.
或者幫助猜拳工作.
都是弟兄姊妹參與之餘.
又要有金錢.
又要去奉獻.
又要去幫助別人的時候都需要.
教會另外一些.
剛才說的堂所的設備.
需要徵由所.
或者我們.
我們聽過有外展的工作.
有關具工作.
所以我們有外展備用金.
關固備用金.
這些都是一些工作.
但是這些東西未必是你有份去參與.
但是也有弟兄姊妹參與.
但是這種勞碌.
一種工作.
我們都是要去事奉的.
但是如果你說我沒辦法參與.
我這個奉獻都是美好的.
又參與又奉獻.
金錢當然更美好了.

$^{401}$這個都讓我體會到.
這些素祭帶來的精神.
是我們人要去.
讓我們的生活.
能夠去成為一個祭獻給神.
所以梵祭和素祭的.
兒童的地方是什麼呢?.
相同的剛才說了.
這些都是恆常去做的事情.
都是獻給神.
當然如果以實物來說.
梵祭就是獻動物.
素祭就是獻什麼呢?.
獻這些植物.
但是這兩個大家都是.
重要是在關聯人和神的關係.
是在表達人對神的關係.
獻上自己.
而梵祭是全人的奉獻.
是在說我們生命的奉獻.
整個人生命的奉獻.
而這個素祭.
是在強調我們日常生活裡面的奉獻.
所以你說我的生命已經獻給神了.
我是一個.
我是基督徒.
我知道生命是在於神的.
但是也是在生活每一個部分.
你是一個家庭主婦也好.
你是一間公司裡面的職員也好.
你自己不同的角色,身份也好.
你的生活也要獻給神.
所以新疆有很多地方都教導我們.
我們去服侍.
我們去工作.
我們參與一件工作.
我們面對我們的老闆也要好像什麼.
好像我們這樣服侍態度.
好像對神那種服侍一樣.
能夠尊重.

$^{441}$能夠擺上我們自己.
在生活上各層面.
都能有一個奉獻的心在裡面.
接著說奉獻是感恩的表達.
大家猜到是說哪個祭嗎?.
不是梵祭,不是素祭.
是平安祭.
平安祭我們看七章十二節.
就要把用油調和的無膠餅.
抹了油的無膠薄餅.
和用油調和的調勻的細麵造成的餅.
平安祭的意思當然是平安.
和睦的意思.
以色列人見到對面.
大家平安平安就是意思的字.
這個獻祭是獻什麼呢?.
是獻動物的.
是公牛,公羊和山羊.
這個方式有點不同.
跟梵祭.
梵祭是燒光的.
留一層皮給祭司.
這些只是脂油.
腰部獻給神.
即是脂肪.
有些肉.
我不詳細說了.
會給祭司.
其餘的肉.
會給獻祭的人和家人一起分享.
你會體會到.
這次獻上的動物.
很多人也有份分享.
大家真的可以吃了.
「表達了感恩,還願和自然獻上」.
強調了與人,神.
甚至大自然.
一個和睦的環境.
提醒信徒的奉獻.
帶著感恩和分享的意思.

$^{481}$這個分享帶來團契和睦.
所以平安祭.
簡單來說是對神恩典的回應.
這是感恩.
亦是人在困境中.
求神幫助我們.
以致得到神的拯救.
我們要用平安祭來還願.
表達信守承諾.
有時也會這樣.
有些人去奉獻傳道.
聽到神呼召他做傳道人.
他覺得不行.
除非你幫我處理某些事情.
或有些地方在負擔.
希望你幫我輕省這件事.
無後顧之憂奉獻傳道.
祈禱神真的幫他.
讓他不能離開.
獻上自己.
要還願,要感恩.
有時我們也會這樣.
神,你醫好我.
我一定會向你面前感恩.
這些都是從心裡的回應.
這也是正面的.
這不是一種交易.
是從心裡的回應.
神的恩典.
我們回應.
我們作出一些感恩的行動.
分享恩典也是.
奉獻不只是給神.
是跟人一起分享.
所以他分享給家人,祭司.
這是很美好的,很有心的獻上.
讓更多人為他獻上感恩.
所以你的生日感恩宴是可以的.
你的感恩分享給別人一起吃飯是開心的.
你的結婚週年感恩宴是可以的.

$^{521}$你儘管請我來.
是我為你感恩的.
其他都是.
我的團契也是很多的.
特別是長者團契.
年頭很多感恩的.
他們有感恩的聚餐.
在酒樓裡邀請牧童宮的弟兄姊妹.
一起參與,很開心的.
他們請我們一起去吃飯.
這種是看到他們感受到神的恩典.
很樂意去分享.
有時候有些頂尖女人團契裡.
有一個感恩的.
可能買些蛋糕,大家分享.
都是一個感恩的表達.
教會童宮裡常常都會有類似這樣的東西.
都是一個大家感恩,有什麼特別情況也好.
我們買些東西大家一起去吃.
都是一個感恩的分享.
當然,到了今天.
新約教會怎麼提醒我們呢?.
《過門多後書》九章十一節.
「你們必凡事富足,能多多施捨.
這就是說,當你這個感恩的行動去實踐的時候.
這是很好的事情.
因為你看到那個平安祭.
你宰一隻羊,不是燒到盡,上升給神.
而是只有神愛的油脂,脂肪,大家都不吃那些.
其他肉的就分給別人.
祭司也好,家人也好,朋友也好.
所以你看到,神很樂意大家彼此的分享.
你能夠去幫助別人,分享給別人.
其實就是在獻感恩祭給神.
這是很美麗的.
當然到了今天,教會裡面我們沒辦法.
每個人的感恩都整個圍在桌上吃飯.
吃到大家都吃得很飽.
我們很多時候用感恩奉獻代替.
教會有什麼?感恩奉獻.

$^{561}$你不只是在感恩禮拜才做感恩奉獻.
你隨時都可以做感恩奉獻.
這是很美麗的事情.
所以感恩奉獻是怎樣呢?.
同樣的,那些奉獻去到哪裡呢?.
這些奉獻仍然是幫助教會去支持教會的工作.
剛才也說了.
你去感恩去幫助別人.
你能夠凡事付足,多多施捨,都是幫助別人.
我們叫做教會的工作去幫助別人.
這些工作,我們繼續做傳福音的工作.
外展的工作.
關顧的工作,無論是會內會外的人也好.
報道的工作.
這些都是繼續去使用下去.
這些金錢會帶來什麼呢?.
是歡樂的.
因為這些工作帶來的果效,有救恩的喜樂.
這些工作帶來的有團契感恩的喜樂.
這是一個和睦感恩喜樂的祭.
所以感恩奉獻是美麗的.
好了,第三方面.
獻金不是贖罪的手段.
我講兩個祭.
就是贖罪祭和贖顯祭.
你還記不記得四章二十節說.
「祭司要為貪門贖罪,貪門就蒙赦免」.
贖罪祭是指罪.
字面意思是為罪的祭.
祭物都是動物.
公羊,山羊,公牛,綿羊.
視乎你的身份而定.
但是很重要的是血要灑在壇上.
陶在聖所.
肉是交給祭司享用的.
這個意思是罪是無意犯錯的罪.
但是血代表的是赦罪.
所以血是很重要的.
當然我們看到耶穌基督釘上十字架流出寶血.
就很明白這個意思.

$^{601}$血灑在壇上象徵了罪得赦免.
本來這些血是犯罪的人要流出來的.
但是不需要用動物代替.
舊約裡面強調了神的赦免.
提醒人要先潔淨心靈才是敬拜神.
所以有時他們是無意犯了罪.
他們有些地方不體統,得罪神.
他們就會有一個心靈覺得自己需要潔淨.
就可以有這個贖罪制.
我們在五章十五節還沒有記.
那裡是這樣說的.
「若有人在耶和華的聖物中無意中犯了罪.
有了過犯就要獻羊群中一隻.
沒有殘疾的公面羊給耶和華為贖遣制」.
贖遣制英文就是一個"guilty"這個字.
通常說法庭裡面定了你有沒有罪.
定罪有罪.
祭物就是面羊和山羊.
但是這個獻祭除了獻祭之外.
還要賠償.
因為牽涉了你對人的.
你犯這個錯誤是在損害別人的.
是令人有損失或者有損傷的.
這個要賠償.
不單是向神面前認罪.
你要做到賠償這件事.
這個當然某程度上是一種故意的犯罪.
需要賠償.
所以具體當然有些是可以計算的.
你偷別人的牛要還給別人.
還要加五分之一的情況.
等等這個獻祭.
主要的問題是你要承擔罪的責任.
以及要藉著賠償來修復關係.
贖罪和贖遣制.
很多時候是個人為自己的過犯去贖罪的.
但是有時候也會.
以色列人是全體.
百姓是祭司.
大祭司是代表以色列百姓整體是一個贖罪.

$^{641}$當然以色列人有贖罪日.
是一年一次的情況.
這些是群體性的.
百姓的共同是要承認神的主權在裡面.
要修復神與人的關係.
新約是怎樣教導呢?.
我們看希伯來書十章十二節.
「但基督獻了一次永遠有效的贖罪祭.
就坐在神的右邊」.
這個意思是經文說.
耶穌基督獻祭一次性.
完全,永遠有效.
祂的犧牲代替了舊約裡面.
反反覆覆的贖罪祭,贖遣祭.
祂的血也是用血.
能夠真正的赦免.
使人與神和好了.
當然祂的流血是代替我們自己.
不需要流血.
祂的死代替我們的死.
但今天你有沒有留意到.
教會哪裡是贖罪祭,贖遣祭.
有沒有奉獻「認罪奉獻」.
不好意思的.
我今天有個奉獻叫「認罪奉獻」.
我們沒有這樣的情況.
但更重要的是甚麼呢?.
剛才我提到的標題是甚麼呢?.
就是這個「獻金」不是贖罪的手段.
當時早期教會或者是天主教.
慢慢就流行了一種做法.
就是贖罪券.
在凡義之前他們才廢除了.
因為好像是用了「獻金」去贖回自己.
這些錯誤的罪.
但更離譜的是甚麼呢?.
買了這些贖罪券去犯罪.
這是很不合神的心意.
但從此至今.
新約教會神從來不是要人可以用金錢.

$^{681}$可以換取神的赦免.
這是無論如何都不可以的.
你自己得到赦免.
然後你感恩奉獻是另一回事.
但不可以換取這種手段.
但今天來說.
我們說不用了.
耶穌基督已經救贖了.
我們不用經常說認罪悔改.
我們更加需要.
我們知道的.
因為今天新約教會我們是怎樣呢?.
是包括了認罪悔改.
承認自己的罪.
接受耶穌基督的赦免.
這個可能是得救.
但平日我們仍然犯了不討神喜悅的事情.
得罪人,得罪神的事情.
我們都會覺得我們心靈有些污穢.
我們不能坦然面對神.
我們都需要在神面前認罪悔改.
這個當然我們有自己的認罪悔改.
教會裡面.
崇拜主席會帶領我們.
但大家都知道.
我經常說週五.
星期五早上認罪悔改的祈禱會.
真的是一個獻祭.
因為禱告也是一個獻祭.
在舊約裡面說.
這個禱告的獻祭.
就是贖罪的獻祭.
就好像新約現在這樣的實踐.
我們去禱告神.
我們去悔改.
無論是個人.
無論是群體.
都是一個美好的祭.
都是一個與神復和的祭.
我們是應該要參與的.

$^{721}$所以.
你平時在家裡為自己的罪獻上悔改認罪的禱告.
是很好的.
但你來參加祈禱會.
你會看到你是為教會一起.
一同有這樣的心智.
向神的面前認罪悔改.
這個我們都需要的.
另外一方面就是生命更新.
用聖祭的生活回應神給我們的救贖.
羅馬書剛才也說了.
傳人獻賞是聖潔的.
這個很重要.
他在說你獻賞的人要聖潔的.
這個才是合乎神的心意.
所以心意更新.
另外信心倚靠.
相信基督一次的獻賞.
我們就永遠得救.
你不需要再懷疑我們有沒有一個救恩.
而我們有救恩之後.
我們要好好做好我們自己.
當然這個是用我們的生命.
用我們的禱告來獻賞贖罪的.
金錢奉獻不是贖罪的手段.
第四點.
「甘心奉獻」.
「甘心樂意奉獻」.
有一次有個姐妹對教會來探訪的弟兄姊妹說.
她說當我丈夫崇拜的時候.
聽完道之後有所感動.
我就會馬上帶他離開禮堂.
因為接著就是奉獻.
不過大家放心.
待會兒就沒有奉獻了.
我們奉獻放在前面.
我順便說.
因為其實重要的是什麼呢?.
為什麼在一個崇拜裡面先獻賞.
是很重要的.

$^{761}$我們其實在一個禮拜第一次的崇拜裡面.
我們要準備好自己.
無論我們自己剛才說.
這些的賄罪的祭.
這些常年的奉獻.
或者我們這種感恩的禱告.
這種心態.
我們要帶到神的面前.
先獻賞.
所以崇拜開始是獻賞我們詩歌.
獻賞我們奉獻.
獻賞我們禱告.
接著才是領受.
無論是神的說話.
或者是其他祝福.
所以就是這樣的態度.
好了,說回這裡.
哥吾多後書9章7節說.
小眾的小收,多眾的多收.
各人要隨心所願.
不要為難,不要勉強.
因為神愛樂捐的人.
這裡聽到.
神真好.
講心不講金.
是的.
人的心多過體重.
你的數字.
大家知道耶穌稱讚的婦人.
是兩個小時前被紀念的.
因為她將自己所有獻給神.
神不是看人的金額多少.
是看奉獻的心是怎樣.
奉獻因為是關係的表達.
所以舊約的憲制,凡制,素制,平安制.
都是代表神和人的關係的表達.
這是很重要的.
因為你有神的心.
你才有獻祭.
沒有神的心,你就沒有獻祭.

$^{801}$金錢是一個外在形式.
真正的核心.
心沒有歸向神.
但有些是相反的.
我有錢就要奉獻.
心不在.
這個也不可以.
因為奉獻不是數字多少.
法利賽的人就很緊張.
聽聽那個奉獻箱叮叮噹噹有幾多次.
你叮一聲就大了一點.
說他當然奉獻了很多.
或者要看他們的人.
帶什麼修成.
一大捆,一大籮.
會這樣評估他們.
耶穌說不是這樣的.
提醒我們最重要的是心靈誠實.
但如果只有金錢奉獻.
沒有心靈的真實.
是一個很空洞的現象.
奉獻更加不是一種交易.
神是樂意的心.
但金心樂意奉獻.
其實就平衡來說.
我們有些情況.
不是所有人都可以參與很多地方.
但我們也可以用勞碌得來的錢財.
奉獻給神.
表達我們的心意.
這是美好的事情.
所以今天和大家看經文.
再加上新約的教導.
再講教會奉獻的全面一點.
我希望大家.
當大家去事奉神.
你的心態應該怎樣.
當大家去奉獻金錢的時候.
你的心態又會如何.
或者我們整個人生裡面.

$^{841}$有沒有一個全然獻上的心態.
就是在生活裡面.
分別為聖獻給神.
凡事要感恩.
多謝神.
和人分享美好的生命.
施比受更為有福.
以致我們更多的分享.
更多的得到神的恩典.
有更多的感恩行動.
這是很美麗的事情.
所以今天和大家去看.
奉獻不是一個例行的公事.
是和神關係的表達.
當我們就算是金錢的奉獻.
我們的心態是全然獻上的.
自己繁濟的意義.
感恩的分享.
好像平安祭一樣.
和別人分享.
看到別人有需要.
藉著感恩的奉獻.
更多的幫助別人.
然後就是潔淨.
敬拜自己的心靈.
這是贖罪贖獻祭在禱告裡面.
在悔罪裡面.
金錢不能夠換取的.
願神幫助我們.
我們同心祈禱.
更加藉著耶穌基督一次的獻上.
我們更多的心靈誠實的獻上.
以致我們更加知道.
我們的生活,行為.
我們的才幹,力量,智慧,金錢.
獻上的時候.
更加明白你的意義.
以致我們每一個獻上.
都是蒙神你的月亮.
蒙神你自己所揀選,使用.

$^{881}$在你的工作上.
求你親自帶領我們.
讓我們在學習裡面.
繼續被神你自己去使用.
以致我們成為一個.
獻上自己被主使用的教會.
禱告交託.
奉耶穌銘給.
阿們.
\newpage



\section{馬太福音 6:25-34}
\label{sec:031730}
\textbf{憂慮?何足掛齒?}~~~~石偉漢牧師
\newline
\newline
連結: \href{http://www.church.com.hk/acms/content.asp?site=cdc\&op=show\&type=product\&code=031730\&layout=sermon}{\texttt{031730}} ~~~~ 語音日期: 2026-03-01
\newline
\newline
\hyperref[sec:031729]{\small{< < < PREV SERMON < < <}}
~
\hyperref[sec:index_chronic]{\small{[返順時目]}}
~
\hyperref[sec:031731]{\small{> > > NEXT SERMON > > >}}
\newline
\newline
馬太福音 6:25-34
\newline
\begin{longtable}{cl}
\hline
\hline
章節 & 經文 (和合本修訂版)\\
\hline
6:25 & \begin{tabularx}{0.7\textwidth}{X} 「所以，我告訴你們，不要為你們的生命憂慮吃甚麼喝甚麼，或為你們的身體憂慮穿甚麼。生命不勝於飲食嗎？身體不勝於衣裳嗎？ \end{tabularx} \\ \\ \relax
6:26 & \begin{tabularx}{0.7\textwidth}{X} 你們看一看那天上的飛鳥，也不種也不收，也不在倉裡存糧，你們的天父尚且養活牠們。你們不比飛鳥貴重得多嗎？ \end{tabularx} \\ \\ \relax
6:27 & \begin{tabularx}{0.7\textwidth}{X} 你們哪一個能藉著憂慮使壽數多加一刻呢？ \end{tabularx} \\ \\ \relax
6:28 & \begin{tabularx}{0.7\textwidth}{X} 何必為衣裳憂慮呢？你們想一想野地裡的百合花是怎麼長起來的：它也不勞動也不紡線。 \end{tabularx} \\ \\ \relax
6:29 & \begin{tabularx}{0.7\textwidth}{X} 然而我告訴你們，就是所羅門極榮華的時候，他所穿戴的還不如這些花的一朵呢！ \end{tabularx} \\ \\ \relax
6:30 & \begin{tabularx}{0.7\textwidth}{X} 你們這小信的人哪！野地裡的草今天還在，明天就丟在爐裡，神還給它這樣的妝飾，何況你們呢？ \end{tabularx} \\ \\ \relax
6:31 & \begin{tabularx}{0.7\textwidth}{X} 所以，不要憂慮，說：『我們吃甚麼？喝甚麼？穿甚麼？』 \end{tabularx} \\ \\ \relax
6:32 & \begin{tabularx}{0.7\textwidth}{X} 這都是外邦人所求的。你們需要這一切東西，你們的天父都知道。 \end{tabularx} \\ \\ \relax
6:33 & \begin{tabularx}{0.7\textwidth}{X} 你們要先求神的國和他的義，這些東西都要加給你們了。 \end{tabularx} \\ \\ \relax
6:34 & \begin{tabularx}{0.7\textwidth}{X} 所以，不要為明天憂慮，因為明天自有明天的憂慮；一天的難處一天當就夠了。」 \end{tabularx} \\ \\
[1ex]
\hline
\hline
\end{longtable}
$^{1}$第一節末 主佬平安 讓我們一起祈禱.
成為我們生命的鼓勵和安慰.
願主你幫助我們 藉著你的說話更加堅固我們.
讓我們在你的信中更加堅穩.
願主你祝福我們 聽我們祈禱奉耶穌名字而求.
阿們 今天的題目是憂慮.
不知道大家有沒有想起憂慮的事情.
或者說很多憂慮的事情.
不過今天的題目不只是去到憂慮就完了.
還有何足掛齒.
何足掛齒即是不用經常掛在口邊.
有些事情事實上我們都不需要常常被憂慮影響我們的生命.
事實上今天我們讀的聖經都會說到這個問題.
不過首先我們都要問一個很真實的問題.
你是不是一個憂慮的人呢?.
憂慮的人很多事情都成為他的煩惱.
無論是錢,樓,身體,健康,老死等等.
所以我網上找到一幅圖 這種人叫Mr. Worry.
憂慮先生 原來憂慮先生真的有很多憂慮.
你看看他 包括了他的安全,測試,健康,財富,恐懼,失敗.
家庭,朋友,將來,責任,問題,人的關係,壓力.
住什麼地方都有很多憂慮.
所以原來憂慮都是一個很真實的課題.
所以我想耶穌都很關注我們面對著憂慮這件事情.
事實上憂慮都是一個世界性的問題.
中國有一個Mr. Worry 知道他是誰嗎?.
其實中國都有很多人叫Mr. Worry.
有一個很出名的 大家一說就知道他是誰 他叫紀仁.
大家認識他吧 很熟悉吧 平時都見過他.
事實上紀仁不是姓紀的.
在古代的時候有一個國家叫紀國.
紀國裡面有一個人 他常常都很多事情想.
以致出現一個問題 他發現他常常都很擔心.
擔心什麼呢? 擔心天會掉下來.
你們有沒有擔心的? 你們不擔心的嗎?.
你們讀了太多科學的東西了.
天掉下來就被砍 是吧?.
有時候流星可能都會有些混石群掉下來打到你.
會不會呢?.
事實上紀仁就是擔心天會掉下來.

$^{41}$所以他擔心天會掉下來 他就常常都睡不著覺.
因為他擔心天掉下來的時候.
他睡著了就不懂得避開天掉下來.
他又擔心吃東西的時候.
他怕吃著吃著天掉下來又會壓到他.
所以他就為這件事常常都煩惱.
當然他有朋友嘗試去安慰他.
他告訴他 天不會掉下來的 因為當時我不用這個版本.
因為是2000年前的解釋版本.
是不合現在的科學的 他的朋友的解釋.
無論怎樣 總之他的朋友就說.
天是不會掉下來的 你放心吧.
因為我們現在就是在宇宙裡面.
事實上他好像安慰到他.
不過其實他仍然都停止不了他的擔憂.
因此他又擔憂他又擔心地會掉下來 地陷.
你有沒有擔心地陷.
住在佛羅里達這麼久有沒有擔心地陷.
都有吧 都試過吧.
如果你看一些網上報紙報道.
有時某些地方好像最近都有一單.
好像在佛羅里達地陷下來.
有兩輛車掉進坑裡面.
不過他就說他整天都擔心地陷地陷.
又令到他生活很煩擾.
所以這個就成為了我們今天一個很出名的成語.
幾人就憂天.
事實上他不單憂天也憂地.
不單中國有幾人 事實上在以色列也有一些幾人.
有沒有想過以色列也有一些幾人.
今天耶穌裡面也有說一班幾人Mr.Worry.
以色列裡面的一班幾人Mr.Worry.
原來在門徒裡面也有這些人.
所以耶穌在登山寶訓裡面.
說了很多百福文士和法力賽人的教導.
跟著我們上次說到什麼錢財.
跟著轉頭就說到這個憂慮的問題.
今天我們讀的聖經裡面就在六章第25至34節.
如果你去看經文的段落.
你會發現其實可以這樣分開三個小段落.

$^{81}$因為這三個小段落都有相同的形式.
第25至30節一開始就說所以.
跟著我告訴你們不要為你們的生命憂慮.
跟著後面不要為你們的身體憂慮.
跟著第30至33節的段落再一次說到.
所以不要憂慮說.
最後第34節就說所以不要為明天憂慮.
所以你見到是用了三次.
所以不要憂慮不要憂慮不要憂慮的概念.
去說耶穌要說的這個經文的教導.
簡單來說我們可以這樣想.
第34節說所以不要為明天憂慮.
我們什麼時候才用所以這個字.
因為之前有因為.
所以34節的所以.
其中一個我們理解的方法.
就是因為31至33節那個因為.
因為31至33節所以有34節.
但是去到31節的時候.
原來又有一個所以.
所以這個所以是跟哪一個有關呢.
又跟前面有關.
前面的因為是什麼呢.
就是25至30節.
因為有上面說的一連串的時期.
即是說六章25節的所以跟什麼有關呢.
前面24節上面的部分.
是說什麼呢.
兩個星期前港獨說了.
考考大家的記性.
長者我會接受大家不記得的.
年輕的應該記得的.
Tony年輕的.
錢財嘛.
上次我們在六章19至24節.
就說到耶穌其實說到錢財的事.
我們怎樣去面對我們的財寶呢.
當中我就說了三個的法則.
就是儲蓄的法則.
就是我們要積蓄財寶在天上面.

$^{121}$開眼的法則.
讓我們的眼睛是專一.
而不是被很多東西蒙蔽了我們的眼睛.
第三個服事的法則.
就是我們要服事上帝.
要興旱滿門.
就是興旱錢財.
所以其實耶穌上次在說完.
這些關於錢財的事情的時候.
其實正正好就是因為很多時候.
錢財是其中一種.
最容易令我們帶來憂慮的一些東西.
所以耶穌在之後去到第六章25節.
就講完錢財之後.
就帶出錢財所帶來的憂慮的問題.
那當然我們憂慮的東西.
就不只是單純錢財的.
所以今天我們讀完這段聖經之後.
我們的應用就可以不單止在錢財上面.
我們面對錢財的憂慮.
我們這樣去處理.
我們更加是對我們整個生命裡面.
不同的憂慮.
我們可以用這種方式去處理.
今天的重點有兩個.
第一個就是憂慮重重.
剛才我們說到耶穌和門徒說.
要教他們怎樣做.
所以他說我告訴你們.
因為和門徒去說這件事情.
就說到很多的憂慮提醒他們.
他們會不會成為了憂慮之人呢.
不過我又要問一個問題.
憂慮的人有什麼特徵呢.
有什麼特色呢.
只是按著我們今天讀的聖經.
你有沒有想過.
耶穌怎樣定義什麼人.
是一些憂慮的人呢.
如果你看回整個段落.

$^{161}$其中第30節就解釋了答案.
第30節就說你們這些小信的人.
原來當耶穌說這群很多憂慮的人的時候.
原來就是說這群小信的人.
當他們小信的時候.
他們就是憂慮的問題.
產生的原因.
就是他們的Faith.
英文的翻譯叫Little Faith.
很少的Faith.
他似乎又好像認識神.
很信神.
不過有時候那種對神的信.
又真的比較微弱.
所以我們看到.
這裡耶穌就說這群.
他們要說你們這些小信的人是什麼呢.
就是那群的門徒.
那群跟隨耶穌的人就是你們了.
包括我們.
我們都是跟隨耶穌的人.
他說我們這些小信的人.
我們這些跟隨耶穌的人.
會不會我們都是一些小信的人呢.
而小信的人就是面對著.
我們第25至34節的一個憂慮的問題.
原來當我們小信神.
我們就會帶來憂慮.
我們反過來就問了.
什麼叫信呢.
小信你就帶來憂慮了.
Little Faith.
那你的Faith是什麼呢.
你的信是什麼呢.
希伯來書就給了我們一個答案.
就是希伯來書11章1節.
信就是所望之時的實體.
是未見之時的確具.
這個就是胡合本的翻譯版.
聽著了記得了.

$^{201}$信就是所望之時的實體.
是未見之時的確具.
是很明顯跟一些期望的事.
一些未見的事.
一些跟將來有關的事.
牽上一些的聯繫.
所以其實很多時候.
我們的憂慮是來自我們對於將來的看法.
我們對於明日的看法.
我們怎樣去看.
以致我們能不能解除我們的憂慮呢.
正正因為對於將來的事.
有些事我們未見到.
未捉到.
所以我們就要憑著我們的信去面對.
所以我的翻譯就是.
怎樣去翻譯憂慮呢.
憂慮就是所望之時的牢底.
是未見之時的失具.
你知道什麼是牢底了嗎.
這是廣東話.
希望能夠翻譯到.
如果有需要我今天不用翻譯.
希望能夠聽得明白.
牢底就是看到一隻一隻.
牙齒都僵硬了.
僵硬了都是另一個.
很難翻譯.
僵硬了.
沒什麼反應.
不懂得做反應.
我們對於未來的事.
我們有時候都牢底.
或者我們對於未見的事.
我們會進退失據.
好像沒有了把握那樣.
原來憂慮就是一個這樣的狀況.
所以我們在後面第六三十四字的時候.
就會讀到.
我們不要為明天憂慮.

$^{241}$因為很多時候憂慮.
是和明天有關.
是和我們未見的事.
所望的事.
是有一個直接關係.
當然我們既然有些.
我們未見到的事.
所等待的事.
可能會發生.
可能不會發生.
將要發生等等.
而這個過程裡面.
怎樣是帶來一個憂慮呢.
根據一些資料就說.
其實憂慮是一個什麼的狀況呢.
就是一個大腦的預演.
就是說原來有一些事情.
我們正在想.
可能是明天.
可能後天.
可能一個月後.
有些事情.
我們在腦裡面.
預備好去做一次.
我們將在腦裡面.
將一件將會發生的事.
自己用自己的方法.
自己的想法.
自己的態度去做一次.
而在腦裡面.
重複去預演.
以致我們成為了我們的憂慮.
最簡單就是.
你未溫書.
明天就考試.
你腦裡面就不斷重溫.
你沒有溫書.
你考試是不會的.
不會的話就會零分.
零分的就被媽媽罵.

$^{281}$其實你未考的.
不過你腦裡面.
想起這件事的時候.
想起你明天考試.
我就要被媽媽罵.
於是你想著想著.
唉不想想那麼多.
明天你又再想.
明天要考試.
我明天就要被媽媽罵.
所以重複重複.
其實在我們的腦裡面.
不斷預演重複去演習著.
同一件事情.
以致我們慢慢.
在我們心理上.
形成一些繃緊的狀況.
所以就成為了我們的憂慮.
所以人為什麼會失信.
或者為什麼憂慮.
為什麼會侵蝕到我們.
所以我嘗試按著.
我們今天讀的經文.
找出幾個原因.
幫助我們怎樣去認識憂慮.
第一個原因是.
為什麼我們會憂慮呢.
第一就是人想多了.
就是有些事情.
你想很多很多次.
那你開始想著想著.
你就會出現一種憂慮的狀況.
所以你見得到.
在短短我們讀的經文裡面.
只是由25節到34節.
10節.
你發現耶穌就在說.
憂慮重複這些字.
起碼出現了5個節數裡面.
都出現了憂慮的字.

$^{321}$證明了很多時候.
一次又一次去想.
就重複發生這件事.
重複去想這件事.
就很容易令到我們.
在那個思想上面.
就對這件事.
就會產生很多.
佔據你腦海的事情.
所以中國人有一句話叫.
日有所思夜有所夢.
為什麼會這樣呢.
因為你想得太多了.
次數太多.
多到什麼呢.
連做夢.
你的腦子都自動去運作.
因為我們在睡覺的時候.
由小腦運作.
小腦都不斷去運作的時候.
令我們想起那件事情.
所以其中一個.
很重要的情況就是.
人想太多就會這樣.
所以我自己一個觀察發現.
什麼人是沒有憂慮呢.
就是不用腦思考的人.
所以如果你分組的時候.
你說你沒有憂慮.
大家就明白了.
不要再問他問題.
因為他平時都沒有用腦思考.
因為你想得太多.
你不想就沒有憂慮.
不會在腦裡預演.
你都不會想明天發生什麼事.
那就沒有了.
事實上有時候.
有些人想得多到一個地步.
是什麼呢.

$^{361}$其實可能就會不開心.
中國還有一個很出名的人.
不知道你認不認識他.
他寫了一篇文章.
叫做《鱷羊流記》.
這個是我讀中學的時候.
要讀的課文.
必讀的課文之一.
我再講多一點.
你就知道是誰了.
《鱷羊流記》的作者叫做范仲淹.
大家就知道了.
范仲淹有兩句名句.
是什麼呢.
先天下之憂而憂.
後天下之樂而樂.
大家又認識了.
證明大家都在同一個年代.
讀同一本書.
原來范仲淹在《鱷羊流記》裡.
就是正正在講他自己的心態.
但是你留意.
他其中這篇文章裡面.
他寫到.
他說什麼呢.
居廟堂之高則憂其民.
處江湖之遠則憂其君.
進亦憂.
退亦憂.
然則何時而樂也.
什麼時候才可以開心呢.
他說其必曰.
先天下之憂而憂.
後天下之樂而樂.
什麼負.
還有一個負字.
我讀漏了.
原來呢.
我自己當日讀書的時候.
我覺得老師告訴我.

$^{401}$是很好的.
因為這代表了他那種.
愛民的心態.
什麼都為市民擔心.
為治理的百姓擔心.
但是我想一想.
其實他應該有抑鬱症.
會不會呢.
他無論在好的地方.
或者高處.
或者遠處.
無論進或退.
他都有很多憂慮.
很多擔憂.
以致先天下之憂而憂.
什麼叫先天下之憂而憂呢.
天下間所有的事.
人家還沒憂的時候.
我就憂了.
你想想.
天下間每個人都有憂的時候.
他還要比我們早憂.
你想想他有多少煩惱.
他腦裡有多少.
要處理的問題.
不斷在腦裡重複.
又重複.
所以你又問了.
他什麼時候才能後天下之樂而樂呢.
什麼時候.
死了之後.
應該他都看不到.
因為每一樣憂慮.
他都放在心裡的時候.
你想想他很難去快樂.
他很難去喜樂.
所以我都很擔心他.
可能都未必.
那麼開心.
原來人想太多.

$^{441}$就會成為.
令我們出現憂慮的一個很重要的原因.
第二個.
令我們出現很多憂慮的原因是什麼呢.
這裡說.
一個人想大了.
或者想過了.
想大了.
或者想過了是什麼意思呢.
這裡我們看到第27字說.
你們哪一個能直著憂慮.
使受掃多加一刻呢.
這裡原文裡面就說.
你們哪一個.
能直著憂慮.
使受掃多加一刻呢.
你能不能夠.
增加一刻在你的生命裡面呢.
就是說.
可能我們.
看你怎麼理解.
人定了你要什麼時候死呢.
你應該就是那時候死的.
中國人又有一句說話.
我怕我說錯.
炎黃.
要你三更死.
什麼什麼.
到五更.
你們懂就行了.
你們聽得懂就行了.
人生命的時間.
是有自己的定期.
上帝都給每個人.
有不同的生命定期.
所以耶穌就說哪一個能夠直著受掃.
令自己的生命能夠多一個小時.
的受命呢.
而事實上什麼叫想大了想過了呢.
就是人覺得可以.

$^{481}$人覺得可以將自己的受命.
多加一刻多加一日.
多加一個月多加一年.
如果沒有上帝容許呢.
是不可能發生.
而事實上真的有人是想.
做到這件事又是.
來自中國的例子.
大家認識這個人.
大家現在應該會看尋秦記.
可能你在不知道哪間戲院.
你會坐飛機去看.
可能Ivan他們會坐飛機去看.
尋秦記也說不定.
尋秦記裡面是說哪個皇帝.
秦始皇.
秦始皇最出名的其中一件事.
是什麼.
想長生不老.
想長生不老要找什麼.
找長生不老藥.
所以結果.
你在尋秦記那套戲裡面.
你發現他.
找不到長生不老藥.
因為在那個年代是沒有辦法.
找到長生不老藥.
當然在這套電影裡面.
大家看了沒有.
最後裡面.
其中有一段是說.
去到現代.
秦始皇去到現代.
能夠吃到美食.
然後他就問.
項少龍.
這裡有沒有長生不老藥.
他怎樣回答他.
整容.
一看就知道.

$^{521}$一試就試到你了.
秦始皇也想問.
來到現代有沒有長生不老藥.
我就問你.
現在在今天的社會有沒有長生不老藥.
又是沒有.
這些就是你們.
是不是現代人.
現代呢.
古代秦始皇的時候沒有.
今天有.
淘寶.
我淘寶一找.
真的有長生不老藥.
還有九十六個多.
是不是.
上面寫著.
不知道是什麼.
適合男女長生不老.
的仙丹藥.
淘寶有.
但是你敢不敢買.
你敢不敢買完.
九十六個多就可以讓你長生不老.
人民幣.
美金也很便宜.
十多元美金.
你發現.
你也會知道.
不可能的.
但是很多時候人.
可能覺得可以.
靠著自己想.
想得自己有能力.
去改變.
這個事情.
所以想大了也想過了.
過了.
而事實上在經文裡面.
使壽數多加一克.

$^{561}$還有另外一個翻譯.
另外一個翻譯是什麼呢.
使身量多加一爪.
使身體的長度.
多了一爪的長度.
一爪多少呢.
這裡告訴你.
大約十五英寸.
就是一爪.
這裡到這裡叫一爪.
當然每個人的身高不同.
但是普遍我們都會認為.
大約一尺半.
如果用這個翻譯.
你更加明白耶穌說.
我們有哪個能夠讓自己的身高.
身體.
多了一尺半呢.
有沒有可能呢.
今天我看網上.
很多時候明星.
站出來拍照好看一點.
他們會令自己.
身量多了.
就是靠增高鞋墊.
但是增高鞋墊.
你都要增高兩寸三寸.
四寸.
已經很過分了.
而事實上你想想.
穿了這個增高鞋墊.
你都很難再用高跟鞋增高.
就算計算.
都不可能有一爪的高度.
所以我們看到.
原來人很想.
用自己的方法可以做到.
但其實不是他能夠做到.
他想大了.
覺得自己可以.

$^{601}$他想過了籠以為自己可以.
所以呢.
不斷想如何令到不可以的成為可以.
這樣就出事了.
第三個原因.
令到為什麼人會憂慮呢.
就是想錯了.
想錯了.
是怎樣錯了呢.
就是認定了.
什麼都要靠自己.
才能夠解決問題.
所以就想.
因為你自己的人很有限.
自己的能力很有限.
當你只想靠自己又覺得.
又要做得到的話.
其實你就將自己在一個.
困境裡面.
但因為你自己.
認為你自己就能夠做得到.
而沒有想過.
其實是另外一些事情.
令你做得到.
是什麼呢.
包括了第一個重點.
就是想錯了什麼呢.
第一就是我們想錯了生命的優先次序.
第25節說.
生命不勝於飲食.
身體不勝於異常.
人經常在想一些.
在憂慮一些飲食的問題.
耶穌就說.
我們想錯了.
是因為我們生命的優先次序已經錯了.
因為生命和身體.
是更重要.
這是第一.
第二個是什麼呢.

$^{641}$你的生命和身體不是靠你自己.
第26節說.
是你們的天父.
是你們的天父養活他們.
而因為你們.
就是被飛鳥貴重.
所以天父就養活你.
原來是這樣的.
我們的飲食我們的異常.
不是靠我們自己.
去面對的.
原來真正供應萬物的.
是那位的神那位的天父.
如果你又記得.
耶穌怎樣教門徒祈禱.
記不記得我們教過.
我們說我們在天上的父.
我們認定耶穌教我們祈禱的時候.
我們經常說主禱文.
其實主禱文是什麼呢.
其實不是真正耶穌的祈禱.
是教門徒的祈禱.
我們要認清楚.
我們在天上的父.
剛才說到天父是怎樣.
尚且養活.
那些動物.
我們的天父會怎樣呢.
我們知道.
尊神的名之後.
我們日用的飲食.
原來不是我們的.
是神今日賜給我們.
是神伯伯賜給我們.
我們才能夠做到.
所以.
耶穌教我們的禱文已經說了答案.
說了答案.
我們不要想錯了.
我們生命裡的飲食.

$^{681}$靠自己去做得到.
當你想這個方向的時候.
已經錯了.
你的憂慮就會慢慢出現.
這也解釋了.
為什麼耶穌說.
我們的憂慮是什麼問題.
是小信的問題.
因為小信靠我們的神.
我們也不信靠我們的神.
我們怎能夠.
經歷到神給我們.
生活裡對明日的豐富.
所以.
人信不過我們的天父.
結果怎樣呢.
讓我們的憂慮.
在我們裡發芽成長擴大.
然後就將.
神的道塞住了.
將神的道塞住了.
為什麼我這樣說呢.
因為在《馬特福音》又這樣說.
《馬特福音》13章.
講述了殺種的比喻.
殺種的比喻裡.
講述了四種種子的處境.
其中有一種很清楚的.
殺在荊棘裡.
就是人聽了道.
後來有什麼.
世上的憂慮.
和前財的迷惑.
記不記得我們剛才說.
19至24節.
就是講述前財的迷惑.
耶穌剛剛說完前財的迷惑.
然後再講憂慮.
因為前財和憂慮.
很多時候有相關的東西.

$^{721}$而耶穌就很清楚.
第三種的.
殺種的比喻裡的處境是什麼呢.
原來這些前財的迷惑.
對於前財的貪戀.
以致以為靠這些.
能夠解決人生.
又要靠自己面對的時候.
產生了很多世上的憂慮.
原來就是將我們生命的道.
卡住了.
我們看不出果子.
活在一個荊棘裡面.
就是因為我們這種憂慮.
弟兄姊妹.
我們今天要認真去想一想.
我們是不是一個憂慮重重的人.
又這樣說.
憂慮本身不是病.
如果在醫學上.
問一下阿雷.
你們可以.
我上網找到.
憂慮本身在醫學上不是一個病.
但是憂慮本身.
當你容許它.
繼續發酵長大.
它就會成為病.
包括了什麼呢.
憂慮會引致anxiety.
成為你的焦慮.
焦慮是病.
而且.
另一個反向是什麼呢.
就是令你抑鬱.
我上網找了.
透過這個憂慮.
本來是正常的狀況.
但是.
令到.

$^{761}$焦慮就是.
令你的能量過高.
什麼叫能量過高呢.
當你擔心一些事情.
擔心的地方經常都很緊張.
你身體裡面.
很多能量都集中在.
你擔心的事情裡面.
所以你的人.
一個焦慮的人是坐不定的.
整個身體.
都受到影響.
經常都在一個很害怕很懼怕的狀況.
耶穌提醒我們.
原來我們的憂慮不認真去面對.
可以成為焦慮.
另一方面就成為了我們的depression.
成為我們的抑鬱.
就是令到我們沒有能量.
depression就是一個人沒有能量的狀況.
以致.
他生命裡面在一個低沉的狀態.
裡面.
經常都沒有動力去做事.
他說我沒有用了.
我沒有希望了.
因為他解決不了他的問題.
所以這個是一個重要.
對我們基督徒要思考的事情.
我們今天要是炎是光.
我們要為主.
去發光的時候.
如果我們處於一個憂慮.
以致成為一個焦慮.
或者抑鬱的時候.
其實是很可怕的事情.
分享一個見證.
就是我岳母.
最初認識她的時候就知道.
原來我岳母.

$^{801}$本來是.
生活健康都是正常的.
精神狀態都正常的.
什麼事情.
令到她後來有病呢.
其實就是因為.
我們本身住在屯門的.
有一些田地.
那些地方的.
政府就收地.
政府收地就要建西鐵.
剛好那個地方就紮在.
西鐵要走的路線裡面.
他就要收地.
收地的時候.
你地要收回來.
就不能說不同意的.
不過就看看.
怎樣賠償.
在這個收地事件裡面.
他們就按著.
花花草草.
怎樣賠錢等等.
其中這個過程裡面.
本來都是很正常的.
不過最大的問題就是.
岳母就說沒有地方住.
才是問題.
你賠多少錢給我都沒有用.
沒有地方住.
在這個過程裡面那些官員就說.
沒有的賠了錢給你.
就已經賠了給你了.
但是不行.
你怎樣都要給我一間公屋.
香港有所謂的公屋.
就是那些可能廉價的.
千多元港幣.
等於百多元美金.
就可以住了.

$^{841}$有幾百呎的地方給你住.
你沒有地方收了我的房子.
我應該給我一個地方住.
在這個過程裡面.
就不斷說你要給我一個地方.
但是他們一直都說不行.
搞了很久很久.
直到最後期.
終於他們安排妥協.
給了一間公屋.
給岳母上樓.
由收地到上樓期間.
因為很多次在這些位置.
很多的爭執.
爭取權利.
有很多.
岳母就經常想這件事.
經常想.
沒有地方住.
又想著想著.
終於上到公屋.
就覺得公屋的職員.
就是房處的職員.
是想趕他走的.
因為覺得你不應該有這間屋.
給你住.
於是就一直很擔心.
一直很擔心.
在田裡住的時候.
上到樓也是擔心.
於是就開始出現了.
他的精神病.
精神病到一個地步.
就是他試過在家裡.
關上門.
關上門的意思就是.
沒有人進來他的屋.
他會說有人進了他的屋.
搬動了一點點的東西.
動了一點點的東西.

$^{881}$甚至他試過有一個嚴重的階段.
就是他進去廁所洗澡.
他都說.
有房處的職員進了他的屋.
弟兄姊妹.
如果我們不好好認真去面對.
我們的憂慮.
我們可以讓他們成為一個很嚴重的病.
而憂慮就是因為.
我們想多了.
想大了 想錯了.
而事實上.
我們真的想得太多了.
如果你還記得.
Jenny有一次.
就說了一個研究.
不知道你記不記得.
她就說一些憂慮.
或者一些憂慮的處境.
有多少成.
有多少個百分比是會發生的呢.
或者.
有多少個百分比是不發生的呢.
多少.
92\%.
我證實了是錯的.
不是她錯.
是你們錯.
我們聽錯了.
因為我也為了今天這個訊息.
我都問Jenny拿回那個研究.
我就認真看回那個研究.
她就說那個研究其實是這樣的.
上次Jenny也有說.
不過她沒有說內容.
我今天就說多一點內容.
那個研究就是.
在賓州的州立大學裡面.
做了一個研究.
讓參加者在十天裡面.

$^{921}$來到當中.
要求這些參加的人士.
去寫下每天應該用四次的時間.
去寫下他們有什麼憂慮.
有什麼擔憂的事情.
而本身.
這些參加者.
是有一些廣泛性的焦慮症.
本身.
他們都很多事情要想.
所以就叫他們每天都寫下.
有四次的時間.
收集這些.
寫這些.
四次提醒他們.
要寫下他們有什麼憂慮的事情.
就是這樣每天寫下.
每天寫下.
寫下這些參加者所有憂慮的事情.
就是這樣.
寫下這些憂慮的事情.
於是那些研究員就嘗試.
看看他們憂慮的事情.
有多少會發生.
於是就在一路.
做一個研究的時候.
最後的發現是什麼呢.
答案就是原來這班參加者.
他們憂慮的事情.
不會發生.
是有91\%.
差了1\%.
準確一點.
上網看的.
是91\%.
擔心的事有91\%.
是虛驚一場.
跟著那9\%就是說有發生.
那9\%有發生的事.
那些研究人員就發現.

$^{961}$原來他們會發生.
一些憂慮的事情.
原來其中有三分一.
結果根本上就不是.
他們想的那麼惡劣.
原來我們想的很惡劣.
其實是我們想多了.
想大了想錯了.
所以原來有很多.
真的是這樣惡劣.
真正令到我們面對.
一些很難.
面對的環境.
很多時候是我們搞錯了.
所以最後就說.
大約有四分一的參加者.
他們的擔憂其實從來沒有發生.
即是有四分一的參加者.
他們所憂慮的那堆事情.
從來都沒有發生.
原來就是這樣.
這些憂慮其實是魔鬼給我們的陷阱.
魔鬼要讓我們墮入.
那個陷阱裡面.
我們被荊棘抓住.
以致我們更加信不過上帝.
讓我們成為小信的人.
所以今天.
第二個部分我要和你說的.
就是很重要.
不只是我們的長者需要聽的.
是所有人都要聽的.
就是退休生活.
第二點我和大家說退休生活.
我們的退休生活是怎樣呢.
應該是很開心的.
這個退休不是說.
我們年紀大了不用上班.
那種退休的東西.
我們退去我們的煩憂.

$^{1001}$退去我們的煩憂.
重新整理一下我們的生活.
我們怎樣能夠將一些.
我們經常搞擾我們的憂慮.
退走 減低.
以致不成為我們生命的傷害.
我們更加能夠堅定的信靠神.
我想都有四個.
退休生活的方法.
第一個.
讓你有優質.
優質生活變成.
優質生活.
優質一個好的生活.
的方法.
第一善用你的眼睛.
處理憂慮其中一個很重要的.
善用你的眼睛.
耶穌說你們看一看天上的飛鳥.
看清楚.
有些東西我們要睜開眼睛看清楚.
你看清楚.
那個法則.
你就會知道你有信心.
耶穌說你看看天上的飛鳥.
他也不種也不收也不在倉裡存糧.
天父尚且養活他們.
你們不比飛鳥貴種得多嗎.
原來說這些飛鳥.
從來都不需要種植.
不需要收割.
也不需要儲存倉庫.
原來天父都養活他們.
我們是不是.
應該會同樣被天父養活呢.
在早前.
大家記不記得早幾個禮拜.
科羅納淡是很冷的.
冷到有.
這裡多少度.

$^{1041}$因為我知道我只會用.
零度.
等於你們的三十幾度.
三十幾四十度.
你記不記得當時是怎樣.
所有伊灣拿.
都冷到不懂得動.
冷到死了.
不懂得動.
這個就是我家外面拍出來的.
整個伊灣拿都曬太陽.
但其實是動不了的.
有陽光.
但其實已經冷到動不了.
這班伊灣拿.
他們一直在冷.
結果.
我們駕車途中又看到.
第二天看到.
有一大量的鷹.
在湖邊吃這些伊灣拿.
他們正在開自助餐.
All you can eat.
這班飛鷹.
他們也不種也不收.
就是這樣.
我們吃自助餐也要付幾十元.
幾十元都不用.
All you can eat.
吃到飽.
第一姐妹.
睜大雙眼看周圍的環境.
你就知道天賦很奇妙.
讓這個大自然的系統.
能夠有上帝的供應.
讓我們不需要擔心.
我們的飲食.
不需要擔心我們的飲食.
第一姐妹.
在美國.

$^{1081}$你見過多少人是餓死的.
你有沒有認識一些人是餓死的.
美國人喜歡吃荷蘭豆.
因為飽死.
荷蘭豆.
美國人.
死因都是飽死多於餓死.
為什麼.
因為美國人吃得太多.
吃得太豐富.
太肥胖.
很多身體肥胖的病徵.
都是因為吃得太飽.
吃得太多.
吃得太盡.
所以在美國.
睜大雙眼一看.
就知道在美國.
飽死的人.
遠遠多於餓死的人.
所以我們為什麼需要擔心.
我們的飲食不充足.
睜開雙眼看一看.
事實上我上網看到.
原來在美國.
純粹餓死的人.
是很少見的.
通常.
如果真是.
食物缺乏而死.
那種食物缺乏.
不是因為沒有東西吃.
而是因為營養不良.
因為營養不良而導致.
他吃的東西.
不能幫助身體.
供應他需要而令他死亡.
或者老人家.
因為身體衰弱.
不能吸收營養.

$^{1121}$所以這樣死亡.
不是因為食物缺乏而死亡.
原來如此.
所以我們要睜開雙眼.
看清楚.
第二個法則.
多用腦.
想一想.
想一想不是.
什麼都不想就沒有憂慮.
而是用一個正確的方法.
去想一想.
這裡說.
何必為衣裳憂慮.
想一想野地裡的百合花.
想一想.
野地裡的草.
今天還在.
第二個重點.
多用腦想一想.
這些花.
這些草.
其實上帝都為他們.
預備一個最美麗的衣裳.
更何況我們呢.
今天我們最大的問題.
又是什麼呢.
我們是不是沒有衣服穿.
我們的問題又是相反.
我們太多衣服穿.
每天我們都是為了.
今天穿哪件衣服外出而煩惱.
哪件衣服.
襯哪條褲再襯哪雙鞋.
而煩惱.
有時每個星期日師母都會問我.
穿這雙鞋襯還是那雙鞋襯.
我就心想.
兩雙鞋都一樣.
沒有人看你的腳.

$^{1161}$你們會不會.
一回教會看著女孩子.
你的腳的鞋挺漂亮.
不會的.
我們水平線看人.
但我們今天的煩惱.
就是為了穿什麼衣服而憂慮.
但耶穌提醒我們.
我們怎會擔心.
我們沒有足夠的衣服.
沒有足夠的供應.
上帝怎會不給我們.
今天我們想想我們的物資.
往往是多過我們所有的.
如果我們用心去想想.
我們要扔的東西是遠遠多過.
是不是這樣說.
要做了.
就是要扔.
有些東西太多了.
不是天父沒有供應.
是供應過了龍.
你們要控制.
否則我們會有.
另一種問題出現.
所以我們要挑選一些衣服.
第三個.
我們想很快帶出來的.
如何令我們有優質的.
退休生活.
就是什麼.
就是不要說.
禁止你的口是什麼.
因為耶穌說.
所以不要憂慮說我們吃什麼.
喝什麼穿什麼.
因為原來憂慮的人.
很多時候在腦裡想得太多.
在腦裡逼得太多.
就開始說出來.

$^{1201}$說出來就經常掛在口邊.
當你掛在口邊越多.
就開始宣之於口.
慢慢成為真實.
本來不是的.
都變成了真實.
留意這裡.
喝什麼喝什麼穿什麼.
你看看剛才英文NIV的版本.
用了what.
三個what去表達.
什麼什麼.
而事實上憂慮.
其中一個說法.
在你腦裡出現了一個.
什麼想法的程序.
叫what if.
經常在腦裡出現.
發生什麼問題.
如果這樣我又想一堆解決方法.
所以當一個人掛在口邊.
what what what.
經常想就開始.
成為了他要想.
if怎樣怎樣.
以致他經常.
在憂慮裡的困擾.
所以我們控制我們的口.
讓我們不要經常.
擺太多這些東西.
事實上.
一會兒崇拜完了.
如果你現在已經在想.
我們一會兒吃什麼.
一會兒去哪間酒樓吃東西.
一會兒去吃什麼.
當你這樣想的時候.
你的心都已經不轉了.
所以很重要.
因為耶穌說.

$^{1241}$這個問題就是什麼人的問題.
這些是外邦人所求.
外邦人經常都會問what what what.
怎樣去想這些解決的方法.
說一個.
你知不知道.
每一次我回家.
去信箱的時候.
我最怕什麼.
最怕什麼.
IQ題.
最怕就是無信.
一個牧師無信.
怎樣算呢.
一個牧師無信.
無信的話怎樣做牧師呢.
而無信是什麼人.
是外邦人.
外邦人沒有那種信.
沒有問題就繼續擔憂.
但耶穌說我們門徒是什麼.
我們每天都有信.
打開信箱.
每天都有信在裡面.
你的生命是不是這樣.
你的信是不是一時有一時沒有.
我們每天都有神對神的信.
所以我們應該.
我們要知道我們耶穌說.
我們不要活得像外邦人那樣.
因為我們一切需要用的.
天父都已經知道了.
好了說了三樣.
用我們的眼睛.
用我們的腦袋.
控制我們的嘴巴.
最後是什麼呢.
最後先看右上角.
很用心去找這些圖案.
右上角.

$^{1281}$是什麼來的.
不求人.
你知道什麼叫不求人嗎.
不求人幫你拗痕.
我們要求什麼.
求神.
第四個很重要.
我們除去憂慮的方法就是.
我們不求人我們先求神.
所以耶穌說.
很重要是什麼.
我們要跟主教導的求.
跟著耶穌在.
主所教導的.
圖文裡面這樣去求.
所以.
這種祈禱是什麼.
很清楚說我們的祈禱是有先.
有後.
我們祈禱是有先有後.
第三節說你們要先.
求上帝的國.
和上帝的義.
後來就是.
這些東西都要.
加給你們了.
所以你又記得主圖文是怎麼寫的.
我們先求.
願人尊你的名為聖.
願你的國降臨.
願你的子兒行在地上.
耶穌就是這樣教我們先求.
上帝的國和上帝的義.
這些東西在哪裡呢.
在主所教導的圖文.
在後面說我們要用飲食.
神就會設計給我們.
免我們的債如同我們免了別人的債.
不叫我們陷入斯坦.
救我們脫離兇惡.

$^{1321}$原來就是這樣.
幫助我們能夠去調整.
我們的先後次序.
我們就減低憂慮.
因為當你有愁煩.
不知怎麼辦的時候你看著地上的東西.
地上的煩惱.
倒不如看看天上的.
那種豐富.
那麼你知道一定有.
在上帝裡面一定有.
所以到了差不多尾聲的時候.
我就會告訴你.
一個真正.
我覺得信耶穌的人.
應該有的兩句說話.
是什麼呢.
很容易記的,大家會背的.
就是先天下之樂而樂.
後天下之憂而憂.
那你又說牧師在亂說話.
先天下之樂而樂,後天下之憂.
先天下之樂是什麼樂.
喜樂.
我們一定要將我們的生命.
建在上帝給我們的喜樂裡面.
這個不詳細解釋.
因為在肥辣比書裡面.
Alice和Jenny已經說了很多次.
在肥辣比書裡面那種喜樂的生命.
要常常喜樂.
我們那種喜樂.
不是說我們只會玩.
要玩樂那種歡笑.
而是我們在神的基礎之下.
找到那種從神而來的喜樂.
那種Joy.
你又記得我和你說過.
Joy這個字.
我們可以怎樣拆開來理解.

$^{1361}$Jesus, Others.
接著才是Yourself.
你放在耶穌優先.
那你就能夠享受.
先天下之樂而樂.
活在地上如同活在天一樣.
你會滿有喜樂.
祝福,不會被困難所圍繞我們.
而後天下之憂而憂.
接著呢.
真的去到後面.
你才需要憂慮.
所以最後你什麼時候會面對憂慮呢?.
什麼時候會面對憂慮呢?.
是不會的.
又忘記了.
因為你一直都活著.
在先天下之樂而樂.
活在上帝的喜樂裡面的話.
你又怎會有憂慮呢?.
所以後天下之憂而憂是不會發生的.
等一下會認真去思考.
我們今天怎樣去.
面對我們的退休生活.
最後.
總結.
再用多一句經文去回應.
就是.
Janie 說過的.
肥辣比書.
一起讀吧.
肥辣比書四章六節.
應當一無掛慮.
只要凡事直著.
將你們所要的.
告訴上帝.
就是這樣.
肥辣比書裡面.
提醒我們一無掛慮.
為什麼呢?.

$^{1401}$先求上帝的國和上帝的義.
先和上帝傾心吐意.
先去告訴上帝.
這是我生命的需要.
上帝知道.
上帝會公認.
今天你信不信得過上帝呢?.
求主幫助我們.
日日有信.
我們一起祈禱.
這世界總有很多不足.
好像欠缺.
但其實.
上帝本來就是豐豐富富.
讓我們能夠轉眼仰望耶穌.
讓我們知道.
上帝的恩典是旺旺救我們.
讓我們能夠在.
困難,無助,缺乏.
困擾的時候.
專心仰望神.
讓我們知道主必定為我們.
開一條又闊又活的路.
求你帶領我們.
祝福我們能夠去實踐.
去應用今天所學的.
祈禱是奉耶穌的名字而求.
阿們.
\newpage



\section{雅各書 5:13-16}
\label{sec:031731}
\textbf{如何在生活中支取神的能力 (華)}~~~~劉彤牧師
\newline
\newline
連結: \href{http://www.church.com.hk/acms/content.asp?site=cdc\&op=show\&type=product\&code=031731\&layout=sermon}{\texttt{031731}} ~~~~ 語音日期: 2026-03-01
\newline
\newline
\hyperref[sec:031730]{\small{< < < PREV SERMON < < <}}
~
\hyperref[sec:index_chronic]{\small{[返順時目]}}
~
\hyperref[sec:031732]{\small{> > > NEXT SERMON > > >}}
\newline
\newline
雅各書 5:13-16
\newline
\begin{longtable}{cl}
\hline
\hline
章節 & 經文 (和合本修訂版)\\
\hline
5:13 & \begin{tabularx}{0.7\textwidth}{X} 你們中間若有人受苦，他該禱告；有人喜樂，他該歌頌。 \end{tabularx} \\ \\ \relax
5:14 & \begin{tabularx}{0.7\textwidth}{X} 你們中間若有人病了，他該請教會的長老們來為他禱告，奉主的名為他抹油。 \end{tabularx} \\ \\ \relax
5:15 & \begin{tabularx}{0.7\textwidth}{X} 出於信心的祈禱必能救那病人，主必叫他起來；他若犯了罪，也必蒙赦免。 \end{tabularx} \\ \\ \relax
5:16 & \begin{tabularx}{0.7\textwidth}{X} 所以，你們要彼此認罪，互相代求，使你們得醫治。義人祈禱所發的力量是大有功效的。 \end{tabularx} \\ \\
[1ex]
\hline
\hline
\end{longtable}
$^{1}$有一段聖經,在《雅各書》第五章十三到十六節,.
這裡說你們中間有受苦的呢,他就該禱告,有喜樂的呢,就該歌頌..
你們中間有病的呢,他就該請教會的長老來,他們可以奉主的名,用油抹他,為他禱告..
出於信心的祈禱,要救那病人,主逼叫他起來,他若犯了罪,也必蒙赦免..
所以你們要彼此認罪,互相代求,使你們可以得醫治..
一人祈禱所發的力量是大有功效..
我們就讀到這裡..
今天我們真的活在一個缺乏能源的時代,.
不但外面的環境我們看見有能源危機,.
在我們心靈的深處我們也常常在每天的生活當中有一種無力感,不是嗎?.
我們就需要能力在生活當中來做正確的選擇,.
我們需要能力來抵擋這個世界的引誘跟試探,.
我們更需要能力來活出一個得勝又豐盛的基督徒生活..
而這個能力從哪裡來呢?.
我想我們必須承認這個能力絕對不是從我們自己而來,.
因為我們人真的實在是非常軟弱..
我們的身體是軟弱的,.
年輕的時候愛到哪裡就去哪裡,.
愛吃什麼就吃什麼,.
但一過中年就發現好像過了保用期的機器一樣,.
到處充滿毛病,不是嗎?.
我們的意志是軟弱的,.
多少時候我們想做的事情我們做不到,.
所不願意做的事情卻偏偏去做..
我們的情感也是軟弱的,.
我們常常就活在愛與恨的掙扎之間,.
常常在憂慮恐懼的陰影下而生活..
所以我們必須承認作為一個人,.
能力不是從我們人而來,.
能力必須是從神而來,Amen?.
就好像我們在San Jose Light Rail,.
我們看見其實這個輕軌電車,.
其實這個電車本身是沒有能力的,.
它必須要接在上面的電纜上,.
這個輕軌電車想要啟動,.
一定要跟上面的電纜連結在一起,.
支取它的能源,車子才能夠啟動..
同樣的,今天在我們基督徒每一天的生活當中,.
我們也必須跟神的能力能夠連結在一起,.
我們才能夠在每天生活當中經歷神的得勝..

$^{41}$那你說我們怎麼樣能夠在生活當中支取神的能力呢?.
在今天這件事情當中,我們看見雅各教導我們,.
怎麼樣在不同的生活狀況當中來支取並且經歷神的能力..
首先第一個,在受苦的時候要禱告..
各位同學,在受苦的時候要禱告,跟他說..
當然雅各的意思不是說如果你沒有受苦就不需要禱告了,.
而是說如果做一個基督徒,.
我們都需要每天的生活當中來禱告的話,.
那麼在受苦當中你更需要來禱告..
因為在受苦中的禱告,最能夠使我們貼近神的胸懷..
因為以前在忙碌當中,我們的心向神是關閉的..
在忙碌當中,我們的耳朵也是下沉的,聽不見神的聲音..
可當我們在受苦的時候,忽然間我們的心向神是敞開了..
我們急於想聽神的聲音,我們急於想知道神的心意是什麼..
就像約伯一樣,在苦難還沒有降臨到約伯身上之前,.
他以為他很認識神,他以為他很認識自己..
只有當苦難來臨的時候,他的心才真正向神敞開自己..
他的話語變得很急切,.
他的內心也開始落入到懷疑跟掙扎當中..
一方面是因為他開始認識到一個以前他所不認識的自己..
其實你要知道,很多時候在這樣一個過程當中,認識自己也是一種痛苦的本身..
開始認識到自己並不像以前這麼完全了..
開始認識到自己其實也是何等的軟弱..
我們看見約伯在受苦當中的時候,他開始認識到自己的缺乏..
另一方面在懷疑跟掙扎當中,他也親自面對這位又真又活的神..
所以在約伯記督四十二章第五節非常出名的經文,我們一同來念..
我從前豐文有你,現在怎麼樣,親眼見你..
在受苦之前,他對神的認識是從遠處的跟隨..
在受苦當中,藉著禱告,他親眼遇見神..
從豐文到眼見,這是受苦的禱告所帶給我們最大的益處..
當然,在這過程當中,最難學的功課莫過於是等候..
因為我們人總是很急,不是嗎?立刻就要..
我們總是不耐煩等候..
我還記得有位基督徒向神禱告說,主啊,求你給我忍耐,現在就給..
我們不耐煩等候..
我們常常不明白既然神心意是要賜福給我們,為什麼不現在就賜福呢?.
如果神的心意是要來工作,為什麼不現在就來工作呢?.
為什麼要我們等候呢?.
事實上如果你仔細看聖經,聖經裡面充滿等候的例子..
當神呼召摩西到西乃山上,讓他在西乃山上進食四十晝夜之後,.

$^{81}$神才向他說話,神才頒布律法..
我們常常在想,那是不是四十天呢?.
有話快說嘛,那等四十天幹什麼呢?.
當主耶穌復活升天之後,神應許給門徒要降下聖靈在他們中間..
門徒們聚集在馬可樓等候禱告了多少天?.
十天啊,聖靈才降臨..
為什麼十天呢?神的心意要他們被勝算,為什麼不馬上降臨呢?.
多少時候你知道我們好像已經太習慣於這種現今世代所謂的速食文化..
我們不明白為什麼神要我們等候..
我們不明白是因為我們的眼目常常只看見事情的結果嗎?.
可我告訴你,神所看重的常常不是那個結果..
神所看重的是那個成長的過程啊..
跟例如說神看重成長的過程跟他說..
神看重在這個過程當中我們的心思意念可以被神更新嗎?.
神看重的是在這個過程當中我們的個性能夠被神塑造更加成熟嗎?.
神看重的是在這個過程當中我們的信心可以更加的堅定嗎?.
所以在等候的過程當中其實神的手正在工作..
我就想到在申命記第32章11到12節這段經文裡面講到一隻老鷹啊,.
是怎麼叫小鷹在飛翔..
在11節裡面我們看見一個圖畫,本來那隻小鷹很安穩的睡在鷹巢裡面,.
可能正在做夢,在做夢有一天它可以自己翱翔在空中的夢想..
沒想到忽然間老鷹就來啊,攪動鷹巢..
一攪動鷹巢,小鷹就被翻出鷹巢之外掉在空中,天旋地轉啊..
搞不清楚怎麼回事情,趕快就想展開小翅膀,想要揮啊揮啊揮啊..
沒想到就是力量不夠啊,開始往下跌了,眼看就要跌到地面了..
老鷹從下面一把接住它,把它放在自己的翅膀上,然後一同翱翔在空中..
哇,這個畫面看起來很美,但如果你是其中的小鷹啊,就不那麼美了..
天旋地轉的滋味,我相信我們都經歷過,怎麼回事啊?神啊,你還在不在啊?.
十二節講得很清楚,非常寶貴..
一同來念,這樣耶和華怎麼樣?獨自引導他,並無外邦神與他同在..
哦,好喜歡這句話..
這就是神怎麼引導我們的方式,這樣耶和華獨自引導我們..
就像一個老鷹教小鷹飛翔,在我們人生當中經過多少這樣的過程?.
一次又一次,直到我們成熟長大,合乎神的心意為止嗎?.
如果今天你真的明白神在你身上是一個命定的心意,.
而且神的手正在帶領你的話,.
那你就必然能夠在受苦的時候持續地禱告,Amen?.
以上書第六十四章第四節講得很清楚..
一同來念來,從古以來,人未曾聽見,未曾耳聞,未曾眼見,.
在你以外,有什麼神為等候他的人行事?Hallelujah!.

$^{121}$你可知道你的神是為等候他的人行事的神嗎?.
當你在受苦當中,就要等候神..
持續禱告,因為知道神正在工作,Amen?.
第二,在喜樂的時候要歌頌,跟你們說在喜樂的時候要歌頌..
當然雅各的意思不是說只要在喜樂的時候歌頌,不喜樂就算了,Alright?.
而是說,如果說你在喜樂的時候都不歌頌的話,.
那你怎麼可能在受苦的時候,在患難的時候還向神歌頌呢?.
事實上歌頌讚美,應當是基督徒的生活常態..
所以大衛在四篇一百一十九篇一百六十四節裡面,.
他說我因你公義的典章,怎麼樣?.
一天多少?七次讚美你,要注意,不是七天一次..
非常重要,今天很多基督徒是七天才一次來主日崇拜,來敬拜..
是一天七次..
什麼意思?不是叫你算今天現在第幾次,.
而是隨時要來敬拜讚美神,Amen?.
第二點,今天如果你真的明白神多麼愛我們,.
耶穌基督不但在過去為我們死在十日架上留下寶血塗抹我們一切的過犯,.
耶穌基督不但今天現在正在與我們同在,.
賜下智慧跟我們每天生活的力量,.
耶穌基督也在將來應許我們有永遠的生命,.
要接我們回天家與祂永遠同在一起..
祂是那位昨日今日直到永遠不改變的神..
如果你真的明白神的恩典是這樣的豐盛,.
我們怎麼可能只是安靜的坐在這裡,.
好像什麼事情都沒發生過一樣呢?.
每當我們思想到神的恩典的時候,.
我們就情不自禁會歡呼喜樂,歌頌並讚美神,Amen?.
這為什麼我常常鼓勵弟兄姊妹們在崇拜當中,.
一定要大聲來敬拜神嘛!.
傾全心全意來敬拜神,拍手來歌唱,對不對?.
舉手來讚美,開口大聲來禱告,.
有時候我鼓勵你們彼此勸勉一下,.
對旁邊說一些話,雖然有些人覺得很尷尬,.
但是我的意思不是叫你尷尬,.
我的意思叫你敞開你自己啦!.
很多人來到主日崇拜還是很封閉的,.
總是覺得口是緊閉的,.
雙手交叉在這裡看你們在做什麼,.
No,敞開你自己,永生神就在這裡,.
何等樂於觸摸你,.

$^{161}$何等樂於施恩典憐憫在你身上,.
只要你敞開你自己神的恩典,.
就能夠降臨在你身上,Amen?.
只有當你敞開你自己的聖靈,.
才能夠自由運行在你生命當中,.
你才能夠感受到神的工作,神的同在嘛!.
你知道每一次當我在非洲,.
跟我們分堂的弟兄姊妹在一起敬拜,.
我都很感動,.
看見這些非洲的弟兄姊妹,.
真的向神何等敞開啊!何等敞開啊!.
從老到少,從男到女啊,.
你看見大家都歡喜快樂,.
唱歌跳舞來敬拜讚美神,.
甚至在收奉獻的時候,.
都是唱一邊跳舞一邊到前面來奉獻給神啊!.
真的是看見他們每一個人來到身邊,.
那種歡喜快樂!.
我告訴你,.
我在非洲看見這些弟兄姊妹那種臉的表情啊!.
我真的從來沒有看見那種喜樂,.
從來在西方國家也看不見那種真情流露的喜樂,.
是那樣的喜樂!.
為什麼能夠那樣喜樂?.
就是因為他們每一次來到神面,.
那種全心全意的敬拜,.
不是因為他們很富足,.
其實在非洲的弟兄姊妹他們一無所有,.
非常的貧窮,.
可是當他們在神面全心全意敬拜的時候,.
那種喜樂就情不自禁的流露出來,.
Hallelujah!.
反過來說,.
今天我們在北美的華人,.
什麼都有,.
可是來到神的面前卻非常愁苦,.
重感的,.
很多,.
很愁苦,.
為什麼?.

$^{201}$因為我們沒有全心全意敞開在神的面前,.
我們就是依靠自己,.
習慣依靠自己,.
也許你說,.
因為他們是非洲人,.
牧師文化不同,.
中國人不來這一套,.
那很奇怪了,.
當你的孩子來到你面前,.
舉起手說,.
爸爸抱抱,.
你也沒說,.
少來這一套,.
中國人不來這一套,.
沒有啊,.
Right?.
你還不是說,.
抱起來了,.
你才不管別人怎麼看你,.
喔,.
又心動了,.
你根本不管別人啊,.
可見了這不是文化的問題嘛,.
而是一種學習跟操練的功課,.
可以經歷得到的嘛,.
是你一種學習嘛,.
我還記得我自己,.
第一次經歷敬拜讚美的時候,.
1991年,.
當我飛到漢城的,.
這個首爾參加,.
殘國的首爾參加,.
他們的敬拜讚美節慶的時候,.
三天的時間,.
在那裡看見,.
整天在我們敬拜,.
就是舉手在敬拜,.
揚聲在歡呼,.
拍手讚美,.
一開始很不習慣,.

$^{241}$真的非常不習慣,.
可是當我在,.
跟他們一起在拍手敬拜讚美,.
歌頌跳舞來敬拜讚美的時候,.
我告訴你,.
三天當中,.
我從來沒有流過那麼多眼淚的,.
可是我也從來沒有,.
那麼深刻的經歷神國,.
那三天在我整個生命翻轉的經驗,.
我就發現到過去,.
在我們的信仰經驗當中,.
我們多麼是盡力把自己包裝起來,.
保護起來,.
以至於我們的信仰,.
變得如此的理性啊,.
停留在知識的層面,.
我們沒有真正的經歷過神啊,.
可是什麼時候,.
當我們願意學習敞開我們自己的時候,.
我們才能夠經歷神的真實,.
所以我常常說,.
敬拜讚美的能力,.
不是當你在看別人敬拜讚美的時候,.
釋放出來的,.
敬拜讚美的能力,.
也不是在你聽別人敬拜讚美的時候,.
釋放出來的,.
敬拜讚美的能力,.
是你自己全心全意在敬拜當中的時候,.
才釋放出來的,.
Hallelujah!.
一旦你敞開你自己,.
你發現神的靈就能工作,.
Amen!.
我最記得我們教會一位弟兄,.
剛來到我們教會的時候,.
其實他還是個木道朋友,.
其實也不是真正的木道,.
他那個完全不信主,.

$^{281}$陪太太來聚會,.
現在很多我知道這邊丈夫是陪太太來聚會,.
那已經不錯了,.
願意進來是不錯了,.
我知道有些丈夫是在外面停車場,.
他們坐在車子上,.
讓太太來聚會,.
所以你今天進來已經很好了,.
那位弟兄就那個姿勢來,.
陪太太,陪太太,.
那我們敬拜很長嘛,.
他真的百般無聊啊,.
每次敬拜就不知道該怎麼辦,.
左腳這樣子,右腳這樣子,.
每天就不知道怎麼辦,.
有一次聽我講道的時候,.
就聽到說,.
如果你真的什麼歌都不會唱,.
開口練一下吧,張開口練一下,.
練歌詞總可以吧,.
他說好吧,.
既然說不會唱就練歌,.
他就那天決定敞開,.
好,練歌詞,.
沒想到一開口,.
剛好歌詞打了神羔羊,.
他一念神羔羊一開口,.
聖靈就感動他,.
開他的眼睛看見耶穌就是神的羔羊,.
為他的罪定死在十字架上,.
留下寶血洗淨他的罪,.
當場他接受耶穌做他個人的救助,.
Hallelujah!.
今天他是我們教育全職同工,.
Hallelujah!.
只要你敞開,.
神的靈就工作,.
當你喜樂的時候要歌頌,.
太多時候看見我們華人的信徒,.
什麼都不講的,.

$^{321}$都只是愛在心裡也口難開,.
喜樂也不開,.
什麼都不開,.
你要敞開一點,.
何等樂於施恩憐憫,.
真的,.
只要你敞開自己,.
不要再依靠自己,.
緊緊抓住,.
不用等到最後好像沒路可走,.
神的恩典今天就預備好在這裡,.
是樂意向你施恩的,.
第三,.
在生病的時候要尋求主的醫治,.
跟你說他生病的時候要尋求主的醫治,.
跟他說,.
你要知道我們所信的神,.
是為樂於醫治我們的神,.
在聖經當中,.
我們神的名字叫耶和華拉法,.
意思是什麼,.
耶和華拉法,.
神就是說,.
我就是你的醫治,.
意思就是,.
我就是你的醫治的意思,.
神樂於醫治我們,.
當主耶穌在地上的時候,.
如果你仔細讀四福音,.
耶穌在地上,.
基本上做兩件事情,.
無論到哪裡,.
只有兩件事情,.
第一個,.
傳天國的福音,.
第二個,.
醫治病人,.
無論走到哪裡,.
你看就是這樣,.
耶穌傳天國的福音,.

$^{361}$醫治病人,.
那你說,.
他是耶穌,.
你看見耶穌採檢十二個門徒的時候,.
叫他們出去幹什麼,.
一樣,.
做兩件事情,.
傳天國的福音,.
醫治病人,.
後來又檢選了七十個人,.
叫他們兩個兩個出去做什麼,.
一樣,.
傳天國的福音,.
醫治病人,.
直到今天,.
神把大使命給我們教會,.
大使命的內容是什麼,.
馬太福音講到傳天國的福音,.
要去使萬民做主的門徒,.
馬可福音講到,.
要醫治病人,.
守安病人,.
病人就好了,.
我們看見,.
主耶穌的事工,.
就是一個醫治的事工,.
聖靈我們看見,.
聖靈住在我們裡面,.
內住在我們,.
你知道聖靈給我們的恩賜,.
其中一個恩賜,.
也是醫治的恩賜,.
所以做一個基督徒,.
你一定要明白,.
我們的神是樂於醫治的神,.
當你我在生病的時候,.
醫藥也是神的恩典你知道嗎,.
我把醫藥叫做一種普遍的恩典,.
什麼叫普遍的恩典,.
就是像日光與水一樣,.

$^{401}$無論你是信不信耶穌的,.
都可以得到日光與水,.
普遍的恩典,.
醫藥也是普遍的恩典,.
但是對信耶穌的人,.
對神的恩典,.
還有一個額外的恩典,.
那是什麼,.
藉著禱告,.
得蒙醫治的恩典,.
所以做個基督徒你很有福,.
跟你朋友很有福我跟他說,.
你在生病的時候,.
不但是有醫藥的恩典,.
而且有禱告醫治的恩典,.
雙重的恩典在你身上,.
所以千萬不要用一個一個,.
兩個都要用,.
兩個都很重要,.
兩個都很寶貴,.
那怎麼樣能夠來經歷神的醫治呢,.
我們看見這段先講,.
第一個一定要出於信心,.
跟人說一定要出於信心,.
醫治最大的攔阻其實就在於不信,.
你還記得勝經講到,.
主耶穌在拿撒勒不多行神蹟,.
為什麼不多行神蹟,.
因為當地人什麼,.
不信,.
不是神不能,.
而是當地人不信,.
我們看見,.
醫治一方面是神的憐憫,.
另一方面,.
也是信心釋放出來的能力,.
所以你還記得在四福音裡面,.
好多次,.
耶穌在醫治完這些病人之後,.
就對病人說,.

$^{441}$是因為你的信心救了你嗎,.
信心太重要了,.
希伯來書第十一章第六節,.
我們一同來念,.
人非有信,.
就不能得神的喜悅,.
因為到神面前來人,.
必須信有神,.
且信他賞賜那尋求他,.
這講得很清楚,.
你到神面前不能只信有神,.
今天很多基督徒只是停留在這裡,.
信有神,.
那是不夠的,.
這裡說,.
不但要信有神,.
而且要怎麼樣,.
信他賞賜尋求他的人,.
你要相信,.
神樂意施恩典在你我身上,.
你說相信什麼,.
第一個,.
相信神能醫治,.
跟你們要相信主能夠醫治,.
跟他說,.
信心不是一種感覺而已,.
信心是一種選擇,.
選擇來相信,.
相信主能醫治,.
你還記得有一次,.
一個父親,.
他的孩子被鬼附,.
他就把這個孩子帶到門徒面前,.
門徒沒有辦法把這個鬼趕出去,.
他又把這個兒子帶到耶穌面前,.
他對耶穌說,.
你若能做什麼,.
請你幫助,.
很有意思,.
耶穌沒有回答說,.

$^{481}$我能或我不能,.
耶穌回答說,.
你若能信,.
在信的人,.
凡事都能,.
很有意思,.
這位父親在問說,.
你若能做,.
耶穌說,.
No,.
你問錯問題了,.
問題不是我是否能做,.
問題是你若能信,.
在信的人,.
凡事都能,.
所以這個父親就下一句話更有意思,.
他說,.
那我信我信,.
但我信不足,.
求主幫助,.
但是在原文裡面這個不足,.
其實我信不足在原文是我不信,.
換句話說,.
在我裡面真的無法相信,.
無法相信,.
但是父子因著你的話,.
我選擇來相信,.
我信我信,.
雖然我裡面是不信,.
我選擇相信,.
接著他說,.
我們也必須在很多環境當中也做這種選擇,.
很多時候在我們面對環境,.
我們裡面有許多的很難相信,.
無法理性,.
無法相信,.
裡面很多經驗無法相信,.
但是因著主你的話,.
我選擇相信你,.
信心的根基,.

$^{521}$不是我們的感覺,.
信心的根基必須在神所說的話語上,.
當你把你的信心建立在神的話語上的時候,.
你看見神的話語,.
神所應許的必定就要成就嘛,.
所以當他說我信,.
我選擇相信的時候,.
耶穌就醫治了他的兒子,.
Hallelujah!.
第一個,.
你一定要相信主能醫治,.
第二個,.
你要相信主肯醫治,.
各位要相信主肯醫治,.
不但是主能,.
而且主肯,.
有時候我們相信主能,.
但不見得相信主肯,.
總覺得我這麼軟弱,.
那些人可能神比較願意吧,.
那些比較屬靈吧,.
那神的恩典比較能夠靈到他,.
我可能不行,.
像我這麼軟弱,.
怎麼可能呢?.
我們相信神能,.
但不見得相信主肯,.
就像一個長大麻風的,.
你還記得當他來到耶穌面前,.
你知道當年在那個時代,.
長大麻風是被社會拒絕,.
要生活在城門之外的,.
是沒有人想觸摸,.
因為傳染病,.
那個長大風來尋找耶穌,.
我相信他很渴望得知,.
不過不覺得自己有什麼希望了,.
所以他在耶穌面前,.
主若肯,.
我就能得潔淨,.

$^{561}$雖然他這樣說,.
心裡恐怕不可能,.
我知道你不大肯的,.
不過如果你肯的話,.
如果你肯,.
我就一定得潔淨,.
我相信你能,.
但是我不知道你肯不肯,.
沒想到,.
耶穌不但說我肯,.
而且按手在他身上,.
伸手摸他,.
不但肯,.
而且憐憫恩典,.
伸手摸,.
沒有人要觸摸他的,.
耶穌卻伸手摸他,.
Hallelujah,.
耶穌不但能夠醫治你,.
耶穌也肯醫治你,.
第三,.
要堅持的相信,.
更何況堅持的相信,.
醫治的信心,.
必須是堅持到底的信心,.
就好像那個患血漏病的婦人一樣,.
她相信主能夠醫治,.
她相信主也肯醫治,.
但更重要的,.
她必須堅持到底,.
雖然當時你看見,.
眾人都圍擠耶穌,.
擁擠耶穌,.
根本沒有機會讓他來靠近耶穌,.
但她堅持到底,.
鍥而不捨,.
一直到她能夠觸摸到耶穌衣上的穗子,.
就是在那一剎那,.
耶穌感覺到能力從她身上釋放出來,.
所以才問旁邊的門徒說,.

$^{601}$有誰摸我?.
其實那些門徒說,.
什麼有誰摸你?.
大家都擁擠你,.
什麼叫有誰摸你?.
但他不一樣,.
因為能力被釋放出來,.
耶穌知道那是信心的觸摸,.
信心的觸摸,.
結果這位婦人就得到醫治,.
我們看見這是一種堅持到底的信心,.
是一種無論如何的信心,.
就無論如何我要抓住神的應許,.
無論如何我要抓住神的能力,.
直到神的應許成就,.
Amen?.
醫治的禱告,.
必須出於信心,.
第二,.
要彼此代求,.
跟你們要彼此代求,.
雅各書第五章十四節裡面講到說,.
若有病人要請教會的長老來,.
他們可以奉主名用於我,.
為他來禱告,.
我們看見這裡所說的長老是誰呢?.
我相信這裡所說的長老,.
一定不是指的教會一個職位而講,.
因為如果是這樣的話,.
我們教會沒有長老怎麼辦呢?.
這裡的長老應該是指的,.
在教會當中被指派委派,.
憑著信心愛心來服侍別人,.
這裡包括在我們教會的牧師們團隊,.
驅牧的團隊,.
小組長的團隊,.
禱告的團隊,.
我們看見這為什麼每一次在.
醫治聖餐主日當中之後,.
我們邀請在我們中間有生病,.

$^{641}$軟弱病痛的可以到台前來,.
我們禱告團隊樂意與你一同來禱告,.
我們盼望不但是主的愛,.
能夠永留在你身上,.
我們也讓教會肢體的愛,.
能夠永留在你身上,.
一同來永留,.
當你站在前面,.
禱告團隊或者牧者團隊,.
按手在你身上為你禱告的時候,.
是你的信心加上他們的信心,.
一同來領受主醫治的恩典,.
所以當你在病痛當中的時候,.
千萬不要一個人獨自禱告,.
你一定要來尋求主的醫治,.
讓教會的弟兄姊妹與你一同來禱告,.
也一同經歷神醫治的大能,.
雅各在這裡很清楚的教導我們說,.
當你受苦的時候要禱告,.
喜樂的時候要歌頌,.
生病的時候尋求主的醫治,.
換句話說,.
在人生每一種不同的情況當中,.
我們都要到神的面前來,.
藉著禱告經歷神的大能,.
今天我看見很多的基督徒在生命當中,.
最大的問題不是神不垂聽我們禱告,.
而是我們太少禱告,.
以至於我們無法經歷神的能力,.
保羅在哥林多厚書裡面很清楚教導我們說,.
基督在我們裡面不是軟弱的,.
基督在我們裡面是大有能力的,.
求主幫助我們在每天生活當中,.
常常迴轉到神的面前,.
讓我們平心與神連結在一起,.
在每天的生活當中,.
我們就能夠經歷神的真實跟神的能力,.
我們一起禱告,.
謝謝主,.
你不但是拯救我們的神,.

$^{681}$你也是賜下能力的神,.
讓我們在每天的生活當中可以靠你來得勝,.
讓我們迴轉到你面前,.
把眾弟兄姊妹帶到你面前來,.
無論我們今天在受苦的環境當中,.
或喜樂的環境當中,.
或在病痛的環境當中,.
讓我們在這個時刻不是依靠自己,.
讓我們重新迴轉到你面前,.
我們來依靠你,.
緊緊抓住你的應許,.
我們深深地相信主,.
今天你有恩典為我們預備,.
今天的聖靈今天親自地在我們心靈的深處,.
引導我們,.
帶領我們,.
讓我們來到你面前,.
再一次與你連結在一起,.
在每一種生活的景況當中,.
我們再一次知取你的能力,.
能夠向仇敵誇盛,.
向我們的病痛誇盛,.
也向這樣黑暗的權勢誇盛,.
帶領我們在每天的生活當中,.
經歷你的得勝,.
得勝而有餘..
謝謝主,將這些榮耀都歸給你,.
阿們..
\newpage



\section{詩篇 42:1-43:5}
\label{sec:031732}
\textbf{茅塞頓開}~~~~李業基牧師
\newline
\newline
連結: \href{http://www.church.com.hk/acms/content.asp?site=cdc\&op=show\&type=product\&code=031732\&layout=sermon}{\texttt{031732}} ~~~~ 語音日期: 2026-03-01
\newline
\newline
\hyperref[sec:031731]{\small{< < < PREV SERMON < < <}}
~
\hyperref[sec:index_chronic]{\small{[返順時目]}}
~
\hyperref[sec:031733]{\small{> > > NEXT SERMON > > >}}
\newline
\newline
詩篇 42:1-43:5
\newline
\begin{longtable}{cl}
\hline
\hline
章節 & 經文 (和合本修訂版)\\
\hline
42:1 & \begin{tabularx}{0.7\textwidth}{X} 神啊，我的心切慕你，如鹿切慕溪水。 \end{tabularx} \\ \\ \relax
42:2 & \begin{tabularx}{0.7\textwidth}{X} 我的心渴想神，就是永生神，我幾時得朝見神呢？ \end{tabularx} \\ \\ \relax
42:3 & \begin{tabularx}{0.7\textwidth}{X} 我晝夜以眼淚當食物，人不住地對我說：「你的神在哪裡呢？」 \end{tabularx} \\ \\ \relax
42:4 & \begin{tabularx}{0.7\textwidth}{X} 我從前與眾人同往，領他們到神的殿裡，大家用歡呼稱頌的聲音守節；我追想這些事，我的心極其悲傷。 \end{tabularx} \\ \\ \relax
42:5 & \begin{tabularx}{0.7\textwidth}{X} 我的心哪，你為何憂悶？為何在我裡面煩躁？應當仰望神，我還要稱謝他，我當面的拯救， \end{tabularx} \\ \\ \relax
42:6 & \begin{tabularx}{0.7\textwidth}{X} 我的神。我的心在我裡面憂悶，所以我從約旦地，從黑門嶺，從米薩山記念你。 \end{tabularx} \\ \\ \relax
42:7 & \begin{tabularx}{0.7\textwidth}{X} 你的瀑布發聲，深淵就與深淵響應，你的波浪洪濤漫過我身。 \end{tabularx} \\ \\ \relax
42:8 & \begin{tabularx}{0.7\textwidth}{X} 白晝，耶和華必施慈愛；黑夜，我要歌頌祈禱賜我生命的神。 \end{tabularx} \\ \\ \relax
42:9 & \begin{tabularx}{0.7\textwidth}{X} 我要對神－我的磐石說：「你為何忘記我呢？我為何因仇敵的欺壓時常哀痛呢？」 \end{tabularx} \\ \\ \relax
42:10 & \begin{tabularx}{0.7\textwidth}{X} 我的敵人辱罵我，好像敲碎我的骨頭，他們不住地對我說：「你的神在哪裡呢？」 \end{tabularx} \\ \\ \relax
42:11 & \begin{tabularx}{0.7\textwidth}{X} 我的心哪，你為何憂悶？為何在我裡面煩躁？應當仰望神，我還要稱謝他，我當面的拯救，我的神。 \end{tabularx} \\ \\ \relax
43:1 & \begin{tabularx}{0.7\textwidth}{X} 神啊，求你為我伸冤，向不虔誠的國為我辯護；求你救我脫離詭詐不義的人。 \end{tabularx} \\ \\ \relax
43:2 & \begin{tabularx}{0.7\textwidth}{X} 你是作我保障的神，為何丟棄我呢？我為何因仇敵的欺壓時常哀痛呢？ \end{tabularx} \\ \\ \relax
43:3 & \begin{tabularx}{0.7\textwidth}{X} 求你發出你的亮光和信實，好引導我，帶我到你的聖山，到你的居所！ \end{tabularx} \\ \\ \relax
43:4 & \begin{tabularx}{0.7\textwidth}{X} 我就走到神的祭壇，到賜我喜樂的神那裡。神，我的神啊，我要彈琴稱謝你！ \end{tabularx} \\ \\ \relax
43:5 & \begin{tabularx}{0.7\textwidth}{X} 我的心哪，你為何憂悶？為何在我裡面煩躁？應當仰望神，我還要稱謝他，我當面的拯救，我的神。 \end{tabularx} \\ \\
[1ex]
\hline
\hline
\end{longtable}
$^{1}$屬靈生命停滯不前.
可能是發生在個人身上.
有時候亦會發生在整個教會身上.
問題是當我們去到屬靈瓶頸地帶的時候.
這個時間應該做些什麼呢.
首先從錯誤的答案開始.
其中一個錯誤的答案就是.
什麼都不要做.
以前的人開用帆船來航海.
順風順水當然好.
自然不用做什麼.
就會順水而去.
但問題是逆風的時候又如何呢.
有人說逆風很簡單.
我下帆避下風頭.
等風向轉為順風的時候.
才再揚帆而去.
所以有些人可能覺得.
有時候處境困難停滯不前.
就什麼都不要做.
按兵不動先看好.
等候時機.
這個有時候的確是沒有錯.
但你想一想.
什麼都不要做.
就算下帆也是要做些什麼.
而且如果你本身逆水行舟的時候.
什麼都不做.
就只會不進則退.
相反.
另外一個情況就是.
如果你明明有一個機會.
天時地利.
你還是果足不前的時候.
不願意去突破.
現在的困境的時候.
就注定蘇州過後沒有單踏.
就會一事無成.
所以什麼都不做.
不是一個好的答案.

$^{41}$我想這裡.
大部分人也有.
人生的閱歷.
不過我不知道大家有沒有看過.
Pirates of the Caribbean.
Vegetales.
Vegetales大家有沒有聽過.
舉手如果看過.
不是很多.
如果大家在教會時間久一點.
尤其是參與那些兒童士工的時候.
我看到後面.
MARGIE舉手.
這些兒童士工就會見過.
這個其實是一個動畫系列.
大概三十多年前出世.
九零年代的時候.
是非常受歡迎的一個基督教兒童節目.
其中有一些歌手.
有一個叫Larry的唱歌.
會唱一些Silly Song.
那時候大家唱的時候.
都很開心.
其中一首歌叫.
We are the Pirates who don't do anything.
什麼都不做的海盜.
可能大家.
一頭霧水.
如果下次有機會去英文崇拜的時候.
聽到這些是做什麼的.
不要奇怪.
這些就是我們準備會用的題材.
這首Silly Song.
Silly Song的重點是什麼呢.
重點就是沒有重點.
即使是做海盜.
都知道不可能什麼都不做.
所以這首歌才叫做Silly.
就算做壞事.
做賊.

$^{81}$都要出點氣力.
不可能做一個宅男.
躺平.
所以同樣道理.
如果你知道.
你逐齡生命停滯不前的時候.
什麼都不做的時候.
你的問題是不會自己解決.
自己的.
今天我想和大家看兩段的經文.
一段是來自.
詩篇42篇和43篇.
另一段是來自啟示錄.
第二章第四.
第四到第七節.
寫錯了.
這兩段經文.
都是來自比較困難的文學.
題材的.
歷史敘事.
和書信的教導.
這兩件事大家看了很多.
尤其是福音書.
和保羅書信.
比較容易讀.
因為比較直接.
詩歌和啟示文學.
通常有很多.
引用的意象.
需要較多的文化背景.
才能夠理解.
今天就不會和大家好深入.
研究這些細節.
但我想從另一個文學的方法.
整個結構入手.
從這兩段經文.
可以看到一個.
三步突圍的策略.
我叫做回顧.
重整豁出去.

$^{121}$這是三步的.
突圍策略.
首先要回顧.
收集有關的資料.
有關的經驗.
客觀去判斷.
一直用那個方法.
有沒有優點缺點.
需要改變.
然後要重整.
reform.
下定決心許下承諾.
落實要做.
要做的改變.
第三要release.
豁出去.
知道自己已經盡了最大的努力.
去明白問題.
然後找到正確的方案.
就憑著信心.
藉著禱告將這個事情.
交給神.
這個簡單的.
三步突圍的步驟.
常識是做什麼都要這樣做.
所以生活當中.
很多實際的挑戰.
不只是靈性停滯.
你也會這樣做.
Iron Seam 做軟件.
做工程.
都是這樣.
如果有顧客投訴.
你寫的程式有什麼問題.
這就是三步解決的方法.
首先.
有什麼錯誤.
然後看看如何解決.
然後豁出去.
當你的屬靈.

$^{161}$生命停滯不前的時候.
這個三步.
突圍的策略.
可以幫助你.
茅塞頓開.
如果大家手上.
有手機有聖經.
鼓勵大家打開.
詩篇第42篇和43篇.
可以明白.
那個上文.
這兩篇.
一共有16節經文.
大多數聖經學者.
都相信這兩篇.
詩篇其實是一首.
三節的詩歌.
結構很容易.
從重複三次的副歌.
可以看出來.
其實從副歌.
也可以觀察到.
一個對唱的關係.
這首歌好像.
有自言自語的感覺.
在這三個.
這三個段落裡.
那個節數.
一二三裡.
你聽到的是私人的心聲.
心底的說話.
情感的流露.
但在副歌裡你聽到的是.
他的頭腦他的智慧.
理性回應.
可以說是他的頭腦.
他的智慧輔導他的心.
他的情感.
甚至你可以看到.
是神的靈透過他的智慧.

$^{201}$對他的情感對他的心說話.
42篇.
第一段落是一個回顧.
他的心裡.
回想起.
他以前的日子.
神啊我的心切無理.
如綠似無稽水.
我的心渴想神.
就是永生神.
我幾時得召見神呢.
我晝夜以眼淚當食物.
不住地對我說.
你的神在那裡呢.
我從前與眾人同住.
領他們到神的殿裡.
大家用歡呼稱頌的聲音守節.
我追想這些事.
我的心極其悲傷.
這個人聽來.
好像敬拜隊隊長.
敬拜隊隊長.
從前負責.
可能帶隊有個喇叭.
帶到人.
進入神的殿裡敬拜.
但如今.
人們嘲笑他.
你的神在那裡.
或者他可以流放到.
一個很遠的地方.
不能去聖殿.
甚至有可能聖殿已經不在了.
又有可能.
其實聖殿還在.
但是去到的時候.
就覺得神不在那裡.
總之教堂.
仍然每個禮拜.
有崇拜.

$^{241}$但台下冷清清.
一個六十人七十人的房間裡.
可能二十個人坐.
當他敬拜的時候.
感覺就好像.
對空氣說話一樣.
不再有.
歌唱讚美神的感覺.
但是這個時候.
他的頭腦理性智慧.
就以這首副歌.
來和他回應.
我的心啊.
你為何憂悶.
為何在我裡面煩躁.
應當仰望神.
我還要稱謝他.
我當面的拯救.
他的頭腦問他的心.
你為何相信這些一派胡言.
為何認為.
他的神已經不在了.
這些假話.
桃花依舊人面全非.
以前來過這裡的人.
現在都不來了.
為何說你的神不在這裡.
神就在這裡.
你停止你的抱怨.
打開你的眼睛去仰望神.
你要高聲去讚美他.
回顧.
這個步驟的重點.
其實最重要的是退後一步.
來看看那個實況.
來看看這個整全的大局.
有時候人的問題是.
我們過份.
專注於眼前.
被一些有限的經驗.

$^{281}$掩蓋了神.
在這個全人類.
人類歷史裡面.
保守的作為.
所以真正要回顧的時候.
不能被自己的人生經驗.
局限.
你說神不理你.
是不是呢.
你看外面有多少黑漆漆的地方.
你能夠在地球上.
有生命氣息.
已經是不簡單的事情.
你找到很遠的地方.
都找不到.
何來說神不理你.
但狹隘主觀的想法.
的確是我們人的天性.
尤其是年輕人.
年輕人.
人生經驗不多.
當然有很多抱怨.
不妥不妥.
但是就算是年長的.
有所有人生閱歷.
如果將自己的經驗.
變成了一種固執.
不再懂得.
向世界的事物學習.
不懂得向其他人學習的時候.
亦是一樣的.
當你只懂得延伸.
掩蓋失敗的經驗.
而不懂得.
從其他人身上.
學到一些成功的功課的時候.
你的人生經歷.
就是你的盲點.
我特別.
有時候和年輕人說.

$^{321}$或者和剛信主的.
新信主的朋友說.
不要讓.
有些人的人生經驗.
壓到你.
你怎樣可以不成為井底之蛙.
怎樣不會.
一味閉門造車.
學無先後.
達者為師.
你學一下.
你讀一下.
你便不會被一兩個人的經驗.
壓到你.
比你.
活多了三十年的人.
一般來說.
應該比你知道的更多.
但只是當他.
在人生當中不斷.
學習的時候才會.
他可能.
有些人信了主.
之後到現在都沒有.
揭過多少頁聖經.
讀過多少書籍.
他怎會真的.
比你懂得更多.
所以你應該觀察一下.
你所謂基督徒的父母.
有沒有讀過書.
你所謂的老師.
有沒有終身學習.
不認識歷史的人.
注定.
重蹈覆轍.
不認識歷史便注定重覆歷史.
所以回顧.
的意思是甚麼.
回顧的意思是開放你的思想.

$^{361}$不只是看著你自己.
每一個星期發生甚麼事.
回顧代表.
你不怕拿起一本書.
上網問問別人.
甚至人工智能也可以問問.
比以為自己不出戶.
便知天下事好.
你有回顧的胸襟.
的時候.
不代表你接受.
所有的意見別人說你也相信.
這個世界的確.
有很多噪音.
很多假話.
很多說話的確不是真的.
這就是.
開始我們.
看到第二個部分.
私人的心.
當他檢視自己目前的.
狀況的時候.
開始意識到那些問他.
你的神在哪裡.
他們憑甚麼說這些說話.
他開始意識到他需要做甚麼.
的重整.
我的神我的心在我裡面.
幽門所以我從約旦地.
從黑門嶺.
從米薩山紀念你.
你的瀑布發聲.
深淵就與深淵響應.
你的波浪洪濤漫過.
我身白晝.
耶和華必思慈哀.
黑夜我要歌頌祈禱賜我生命的神.
我要對神我的盤石說.
你為何忘記我呢.
我為何因受敵的欺壓.

$^{401}$時常哀痛呢.
我的敵人辱罵我.
好仗敲碎我的骨頭.
充滿不住地對我說.
你的神在哪裡呢.
之前我們在猜.
詩人究竟身在何處.
詩人是否沒有辦法回到聖殿.
因為他被流放異地.
或者聖殿已經被拆毀.
表面來看.
第六節好像是答案.
聽起來.
詩人好像在很遠的地方.
你看約旦地.
黑門嶺.
和米薩山.
地圖離開耶路撒冷很遠.
但你再想想.
地圖.
這三個地方互相都不太近.
所以我們反而懷疑.
這是否應該按照字面.
來理解.
這句說話呢.
你看黑門嶺在哪裡呢.
位於新約時代.
的加利利.
我們很熟悉耶穌長大的地方.
是一個高山.
海拔是九千多呎.
最高山.
在今天的敘利亞和.
黎巴嫩邊境.
靠近現在以色列佔領的.
戈蘭高地.
黑門嶺離開.
約旦都有一段距離.
我們不知道.
米薩山在哪裡.

$^{441}$但有些學者認為應該在.
黑門嶺附近.
所以似乎有一個感覺.
詩人是想借用.
這幾個地名主要表達.
是他離開聖殿耶路撒冷.
很遠.
但他可能是一個隱喻.
很遠 去到天腳底.
一個metaphor.
可以理解 再遠都不知道地名.
在這裡最遠都知道地名.
所以他是被.
一個屬靈放逐到很遠的地方.
當他去到.
第七節說到這個深淵的時候.
這個隱喻的理解.
看來更加可能.
深淵就與深淵.
響應你的波濤洪濤.
慢過我身.
約拿仙師.
在約拿書第二章.
記載他自己在大魚裡.
三日三夜在那裡禱告.
他禱告一開始就回憶.
他自己被人掟下水的時候.
泥流之際.
就是上到詩篇第七句.
所以如果你再看地圖.
看地理就更加沒有可能.
詩人又在約旦地.
又在高山上.
又在深水裡.
所以.
可能他只是.
用一些意象來表達.
他現在的心境.
其實詩人可能.
在耶路撒冷.

$^{481}$甚至在聖殿裡.
但是他覺得.
他與神的距離很遠.
他覺得他已經不在.
昔日充滿榮耀.
充滿活力的.
聖殿當中.
回到聖殿好像身在異鄉.
大海深處.
不止他的.
舊朋友不在.
他還被敵人包圍.
不停的奚落他.
你的神呢.
但是雖然如此.
他的理智仍然.
鼓勵他去重整.
他現在的情況.
他需要做些什麼.
他發現他需要向神祈禱.
他禱告裡面說些什麼.
首先他清楚肯定.
這個信仰.
神的愛沒有離開過.
這個地方.
神的愛從來沒有消失過.
百州耶和華必思慈愛.
黑夜我要歌頌.
祈禱賜我生命的神.
神從來沒有離開過.
人離開了神沒有離開過.
接著他因為有這個信念.
他就可以大膽的向神.
提出一個請求.
你可不可以現下形.
好像忘記了我一樣.
可不可以證明給我看.
欺壓我奚落我.
那些人是錯的.
我的主啊.

$^{521}$為什麼我必須要再被.
欺壓的時候在哭呢.
你幫我還擊.
可不可以.
他基本上是說什麼.
我想做些什麼我想要.
反駁但我也知道.
單靠我不行的.
我做任何的事情.
你啊我的主啊.
你首先要做一些事.
這個就是.
他學到的功課.
在神裡面的重整.
不是靠自己力量.
在神裡面的重整.
是要依靠神開始.
依靠神不是.
躺平在青草地上.
什麼都不想做一個.
宅男依靠神.
是很勤力的向神祈禱.
求祂出手.
讓我們可以加入.
神的行列裡面.
所以他提出這個重整的.
請求之後.
他的理智就再一次.
能夠唱出這首副歌.
我的心啊你為何憂悶.
為何在我裡面煩躁.
應當仰望神.
我還要稱謝他.
我當面的請求.
我的神.
最後他是43篇第三個段落.
他回顧.
和重整之後.
斯人就大膽地.
豁出去了.

$^{561}$release.
release的意思就是.
將他本來握在手裡的.
那個愁悶.
他覺得自己沒能力做的地方.
完完全全的放手.
拿到神的面前求神幫助.
所以在第43篇.
那個第一個段落.
裡面他能夠這樣說.
神啊求你為我申冤.
向不虔誠的國.
為我辯護.
求你教我脫離詭詐不義的人.
你是作我保障的神.
為何凋棄我呢.
我為何因受敵的欺壓.
時常哀痛呢.
求你發出你的亮光和信實.
好引導我帶我到你的聖山.
到你的居所.
我就走到神的祭壇.
到賜我喜樂的神那裡.
神我的神啊.
我要彈琴稱謝你.
我們之前在想.
這個詩人在哪個地方.
地理上.
我們剛才說.
可能不是肉身的流放.
而是一種心靈的流放.
第43篇.
我們這個第二段落.
看這個可能性更加有可能.
原來奚落他的.
那些敵人是一個.
不虔誠的國.
An unfaithful nation.
如果你舊約裡面看到.
不虔誠用單數的時候.

$^{601}$通常不是說外邦列邦.
是說誰呢 以色列民.
所以這個.
不虔誠的國.
Unfaithful nation 通常說.
是指以色列國.
是指以色列民.
他的同胞原本.
應該是神的子民.
他自己教會的人.
但現在反過來.
反抗他.
變成什麼樣子 詭詐不義.
不是來神裡面.
神的殿裡面敬拜聽道理.
反而在外面.
在休壇裡面敬拜偶像.
他還說.
這個私人樂隊領班.
你的神在哪裡.
給我看一下.
你的畫.
是不是靈驗 你的宗教過時了.
我們都很靈.
所以這個私人聽到這個.
當然在第二節裡面.
說到他自己非常愛童.
但是.
跳到第四節的時候.
他仍然可以喜樂.
為什麼.
我沒有寫出來 但第四行.
第五行就有了 為什麼.
關鍵就是中間第三節.
就是他知道.
怎樣可以從悲哀.
變成喜樂 就是向神禱告.
知道神.
真的最終會為他平反.
因此他所尋求的.

$^{641}$只有一件事.
神的亮光 神的信實.
聽到神的話語.
就已經足夠了.
即使沒有人來教他.
聽他唱歌.
沒有人來聽他講道.
即使沒有人他唱歌的時候.
跟著他唱.
他都完全不介懷.
因為他在神的真理當中.
已經看到神.
他不用很多人在他面前.
他知道神和他同在.
打開聖經他就知道神和他同在.
他最大的滿足.
最大的喜樂.
就是打開聖經 發現神的面.
神的榮光照耀著他.
這個喜樂滿足.
就已經讓他可以感覺到.
自己上了聖山.
還何必要計較.
那些已經離開 轉向其他.
神不虔誠的.
這些以前不忠信的以色列人.
那些以前教會.
所謂的成員.
嘲諷和欺騙的言語.
為何再被他困擾.
你最好的朋友.
是誰?.
你的琴.
可以和你一起向神敬拜的琴.
就是你最好的朋友.
所以這一段經文.
幾乎是整個旅程的總結.
所以使得他的智慧.
他的理性 他從神而來的靈.
再一次向他的心.

$^{681}$唱出這一首副歌.
我的心啊 你為何憂悶.
為何在我裡面煩躁.
應當仰望神 我還要.
稱謝他 我當面的拯救.
我的神 所以我們就可以在.
這個整個很短的一首.
三節的詩歌裡面看到.
回顧 重整.
和豁出去.
這三個特為的策略.
當詩人回顧.
他自己目前的情況的時候.
發現他有兩件事原本是出錯的.
第一 他太在意.
人家對他的嘲弄.
嘲諷.
他應該怎樣做.
不要理他們.
甚至不要當他們是朋友.
當他們是敵人.
為甚麼.
因為他們真的已經不是在神的裡面.
哪樣東西可以給他最大的喜樂.
不是別人的.
假話.
甚至不是他自己心裡面.
發過牢騷的心聲.
只是神的話語.
才是他真正想要聽到的聲音.
第二.
起初詩人似乎不是很肯定.
神真的會理會他.
會出手來幫助他.
但是他想深一層.
神為甚麼會不出手.
但神是要將我們放在一個.
患難 一個困難的處境當中.
要我們懂得去.
依靠他 要懂得向他求助.

$^{721}$讓我們知道.
你要甚麼平反.
你的喜樂不是靠你的力量.
靠人的力量 而是靠神的憐憫.
所以他.
能夠回顧 能夠重整.
他就能夠豁出去.
茅塞頓開.
讓我們再.
剩下一點時間看第二段的經文.
啟示錄一開始.
就記載了主耶穌基督.
透過約翰.
寫給小亞細亞.
七個教會的書信.
大部分色經家認為.
這些教會所面對的.
就是一個新約教會時代裡.
很典型教會見到的問題.
我們今天.
當然沒有時間去讀完七封信.
只是從第一封信.
寫給二佛所的信裡面.
我們其實都可以.
看到一個.
典範.
都可以看到.
當中怎樣去.
改變 亦都是一個回顧.
重整豁出去的一個步驟.
三部特圍的策略.
首先.
主耶穌透過這些說話.
幫他們回顧.
我知道你的行為牢固 忍耐 也知道你不容隱惡人.
你也曾測驗那.
自稱為是使徒卻不是使徒的.
看出他們是假的.
你能忍耐 曾為我的命.
勞苦而不困倦.

$^{761}$然而 我有一件事要責備你.
就是你把起初的愛心.
拋棄了 所以你要回想起.
你是從那裡墮落的.
說話的主.
他回顧的.
以弗所教會的行事.
有很多是值得讚賞的.
他們堅守聖經的真理.
他們斥責假的教師.
甚至為耶穌的名.
忍受很多的迫迫.
這一切的事上 他們都是.
努力不懈的.
但在回顧當中 他亦都.
指出一件事 他們是需要批評的.
需要被改進的.
就是以弗所教會.
最近好像拋棄了.
他們起初的.
愛心.
聖經學者有不同的看法.
對以弗所教會所拋棄的.
愛心有不同的看法.
有些人認為以弗所教會.
180度改變 完全失去.
對神的愛 變成愛世界.
如果你和尚文夏里.
和初期教會.
我們有關以弗所教會的歷史.
來看 這個就比較難理解.
更加有可能.
他們喪失了.
起初的愛 是他們互相.
弟兄姊妹之間的愛.
更加有可能 是他們對教會.
門外的那些以弗所成的人.
的愛 都變成很冷淡.
或者.
這些很多的虛假的教導.

$^{801}$很多的假話.
這些很多的患難的迫迫.
使得他們有點burn out.
有些就好像.
那個詩人 剛才我們讀.
這個詩人一樣.
看到一些不忠實的國民.
現在都放棄了.
都沒有力在愛了.
結果 曾經滄海難為水.
就轉入一種.
保護自己的姿態.
怕受傷.
開始什麼都不想做了.
躺平了 等別人去療傷.
主耶穌對他們說.
他們要重整 他們要悔改.
並且要悔改.
做起初所做的工作.
你若不悔改.
我要到你那裡去.
把你的燈台從原處拿去.
逆水行舟.
不進則退.
如果他們只是樂凡.
躺平等人療傷.
發生什麼事 就飄走.
大江東去南浮還.
應該怎樣做呢.
重整 反守為攻.
如果你害怕.
作出任何的犧牲.
的時候.
你就永遠沒辦法付出愛.
沒辦法將任何人.
帶到主耶穌基督的面前.
如果你.
不肯犧牲.
去付出這份.
從天而來的愛的時候.

$^{841}$贏的就是魔鬼.
如果有人好像.
之前那個詩人.
被挑戰的話.
你為神受了這麼多委屈.
神有怎樣幫過你.
神幫你做什麼事.
神在哪裡.
對他們的回覆.
最好是什麼.
你錯了 神的愛.
就在這裡.
神的愛就在我心裡.
神的愛就在我裡面.
神就在我裡面.
當你有這個能力.
大膽去說這句話.
神的愛就在我裡面的時候.
神就在我裡面的時候.
你就學到了.
這個詩人最終學到的事情.
第一.
根本不在乎.
你聽到你的對手說什麼.
他們根本就不知所謂.
神根本就在這裡.
第二.
你會更加向神求力量.
倚靠的力量.
因為你知道是神將這個愛.
給到你.
使你有力量去愛你身邊的人.
然後.
也能夠去豁出愛.
豁出這個去.
原因就是.
你知道神是公平的.
你知道神知道你在做什麼.
即使你面對挑戰.
即使你有些地方需要改.

$^{881}$神最終來說.
都會欣賞你.
付出的努力.
然而你還有一件可取的事.
就是你恨惡尼哥拉派的行為.
這種行為也是我所恨惡的.
凡有以德的.
都應當聽聖靈.
向宗教會所說的話.
我必將神樂園中.
生命樹的果子.
賜給他人.
這個就是神對我們.
的應許 為什麼我們能夠豁出去.
為什麼我們能夠明顯地.
蝕底 你想想做教會的人.
好像很蝕底 只是付出.
從來沒有收入.
因為我們知道當我們能夠.
付出本能付不出的愛的時候.
你就在經歷神.
你就是代表神.
如果力量不是神 哪裡來的.
你知道不是從你自己那裡來.
所以當你能夠.
去付出這些愛 不是在.
平等的療傷 而是真的走出去.
沒有的都給出去的時候.
你就知道神和你同在.
你就知道別人問你.
你的神在哪裡 你就可以說.
神就在這裡.
很清楚地知道 神對我們的愛.
是真的無盡無限.
我們天和地離開.
有多遠有多高.
神對我們的愛都是沒有辦法.
可以計算 所以我們有時候.
在一個患難當中 有時候困難的時候.
好像看到魔鬼.

$^{921}$撒旦的作為 真的也是.
非常之猖獗 我們甚至.
覺得是抵擋不住 我們甚至.
覺得是孤單 但我們知道.
最終來說我們其實不是.
被神刁棄的 神的愛.
一直在我們當中.
神一直沒有離棄我們.
所以我們 幫助我們真的能夠.
勝過這個世界的.
一切的誘惑.
這個他們的謊言.
他們的各樣的假話.
使得搖動我們的信心.
我們知道神和我們同在.
從今時一直到永遠 因為我們.
有主耶穌基督的話在影響我們當中.
當我們覺得是困難.
我們覺得無力再走下去的時候.
我們知道.
人的盡頭就是神的開始.
幫助我們真的重拾.
這個我們對.
人對神的愛.
使我們真的可以有一個新的超越.
更加的能夠.
彰顯神的榮美.
我們相信主耶穌基督因為.
天賦你的慈愛 聖靈的啟迪.
帶領和保守 必然.
和每一個聖徒同在.
從今天直到永永遠遠.
阿們.
\newpage



\section{詩篇 22:24}
\label{sec:031733}
\textbf{被看見的力量}~~~~司徒永富博士
\newline
\newline
連結: \href{http://www.church.com.hk/acms/content.asp?site=cdc\&op=show\&type=product\&code=031733\&layout=sermon}{\texttt{031733}} ~~~~ 語音日期: 2026-03-01
\newline
\newline
\hyperref[sec:031732]{\small{< < < PREV SERMON < < <}}
~
\hyperref[sec:index_chronic]{\small{[返順時目]}}
~
\hyperref[sec:031734]{\small{> > > NEXT SERMON > > >}}
\newline
\newline
詩篇 22:24
\newline
\begin{longtable}{cl}
\hline
\hline
章節 & 經文 (和合本修訂版)\\
\hline
22:24 & \begin{tabularx}{0.7\textwidth}{X} 因為他沒有藐視、憎惡受苦的人，也沒有轉臉不顧他們；那受苦之人呼求的時候，他就垂聽。 \end{tabularx} \\ \\
[1ex]
\hline
\hline
\end{longtable}
$^{1}$太好了.
我們今天可以相遇.
和大家分享.
都是我心裡覺得.
現在我們活著.
是需要很多鼓勵的.
我不知道大家會覺得.
我們的生活.
和我們的心情.
有沒有一些時候.
會去到一些困境.
我經常都有.
剛才姊妹提到我的身份.
亦都好知道.
其實當下有些行業.
都仍然是水深火熱的.
例如是零售業.
所以在我比較年青的時候.
除了喜歡講課之外.
當時我進入這個行業.
是五元一碗的涼茶.
其實有為的人.
是不會做這些行業的.
應該炒股票 買賣樓房.
那個年代我們叫這些生活.
星球人.
星球的意思就是.
每個星期賺一球.
一球即是一百萬.
那你賣多少碗涼茶才賣得到呢.
所以今天和大家給予一些鼓勵.
包括我仍然在工作裡面.
有很多挑戰.
在生活裡面.
你仍然會發現.
有很多事情是未解決的.
甚至見到有很多困境.
但原來發聲都不太方便.
會不會鬱到病.
有時我們都會在想.

$^{41}$我們的上帝 我們的信仰.
怎樣幫到我們呢.
如果利用民間的信仰.
都沒有鼓勵到你.
因為求那支籤都是中籤.
但其實中籤都是下籤.
因為據說是將下籤抽起了.
但另一個真是下籤.
我們是否看著這些社會現狀.
人云亦云.
我們的上帝在哪裡呢.
我記得在我還年輕的時候.
或者大家是否記得.
香港都走過一些困難的日子.
很明顯我自己都跌落當時的困難.
尤其是財困.
很記得那時的社會.
我自己認為.
這是很主觀的.
溫度.
我自己由上而下的溫度.
都比現在高一些.
覺得很多人很不開心.
二十多年前金融風暴.
很多人跌落情緒抑鬱.
真實的財困.
香港有一個偉大的發明.
就是燒炭.
中國人有四大發明.
粥粉麵飯.
第五大發明是香港人發明.
燒炭.
那些炭其實是燒雞翼.
但是人去到絕望的時候.
困境的時候.
很被氣絕.
於是自己都放棄自己.
我記得那時政府都很有心.
做一些系列的廣告.
去鼓勵我們.

$^{81}$記不記得這個廣告.
這個四眼仔.
還記不記得.
這個四眼仔又帶著憤怒.
又帶著失望.
是一種很不被解決.
不被看見.
一種很孤獨的心情.
那個背景的廣告.
我到現在都仍然記得.
就是他在街上踢球.
但是就好像今天這樣.
或者這兩天這樣.
天氣不似預期就下雨.
那個內心都是再次跌下去.
其實我的要求很卑微.
我想踢一下球.
難道這樣都很奢侈嗎.
於是就說了一句這樣的說話.
打球才下雨.
難道連上天都不喜歡我.
不知道大家有沒有.
有時都會跌下去這樣的情景.
就是很孤獨.
孤獨到一個地步.
自己看不起自己.
連上天都不幫你.
是一個很困難.
大家都沒有什麼出路的日子.
但是這些是不是會隨著.
我們的社會越來越發達.
我們就會越來越有解決呢.
你會發現不是的.
問題仍然在這裡.
在疫情的時候.
我看到這幅圖片.
我心裡面覺得很無助.
因為我們都怕那顆細菌.
怕到大家都躲在家裡.
但也有一些很謹守崗位.

$^{121}$但其實他都沒有辦法.
究竟裡面夠不夠強壯.
去應對那個災難呢.
尤其是這些災難很陌生.
在工作環境裡面.
不是我們互相鼓勵.
勤力一點就有出路.
但其實真實的情景.
我們以前是不是很相信多勞多得.
其實這只是一個迷思.
小弟很多年前寫過一本書.
叫做老闆要的不是牛.
原來做牛做馬.
可能是被人遺棄.
錦上路上有很多牛.
現在連豬都要搵食.
在山野裡面出現.
在市區裡面.
你會發現真實的社會現狀.
環境不是很深有好報之餘.
也不是盡力就有回報.
甚至你會見到一孩不如一孩.
貧窮的問題.
我們發現社會越發達.
但不知為何好像仍然沒有被看見.
有一次很多年前送一個很久沒見的朋友.
小時候很熟悉的.
但若干時候失聯.
聚會後送他回家.
去到深水Po 再行多兩條街.
其實他都不想我繼續進去.
但因為我們盛情難卻.
結果我都很驚訝.
去到那條窄路裡面.
一個沒有燈的樓房.
原來那裡真的有幾個小房間.
後來我才知道原來是tong 房.
我小時候已經好少少.
我知道有些廁所和洗澡的地方在外面.
整層的人都在用那個地方.

$^{161}$我好少少.
我的家和對面的家一起分享一個廁所.
但洗澡都是在家裡的露台或廁所.
人有三急衝到那裡.
竟然被人鎖住了.
因為裡面有人在用.
我們都會在那裡等.
但好奇怪.
不會有憤怒和不開心.
希望人們盡快出來自己進去.
是這樣的.
但多年之後我發現.
社會不是應該進步嗎.
原來是退步.
那個環境.
但你會發現有好多人.
有一部份人.
但真是所有的職聚.
連穿襪都穿不下.
肚滿腸肥.
有些人可以累積大量的財富.
但又會見到有些人仍然在困境裡.
我不知道我們一直說的物質貧窮.
或是我們生活的困境.
但都有好多人活在一個孤獨的空間.
心靈沒有被觸碰.
有一天我看一個報道.
說日本的年青人想到一個辦法.
把自己找出來.
大致上時薪一百元左右.
做什麼都可以.
當然全部都是正當的事.
真是好奇怪.
原來是有生意的.
這裡有個女士.
她說我想你陪我吃飯.
她很喜歡吃印度菜.
所以她穿了一套印度服裝.
我只想你不用說話.
坐在這裡陪我吃印度菜.

$^{201}$不需要說話.
可能她在生活裡有一種空虛感.
沒有被人理解或接受.
有些人只是陪人等人.
有兩年前我看足球.
美斯來香港.
我很興奮想看美斯踢球.
我發現我坐在美斯對面.
原來是美斯陪我看足球.
有些人空虛到一個地步.
很想有人陪他等人.
事實上那個人有沒有出現都不知道.
只是一種簡單的心靈訴求.
不是我們資訊爆棚嗎.
所以有時我們教會看聖經都很小心.
在小弟的教會仍然放置實體聖經.
因為很怕人滾手機的聖經app.
滾著滾著就有很多資訊進來.
社交群體的資訊進來.
又或者好像現在這樣.
不太中聽.
本來在看聖經的app.
滾著滾著就有其他資訊.
你會發現資訊是爆炸的.
人不是應該更加滿足嗎.
你會發現這裡落差很大.
所以是一種心靈的貧乏.
是一種生活的困境.
或者當你情緒甚至是身體出了狀況.
得如其來.
我一個很熟悉的家庭.
兒媳婦戀.
上幾個星期突然驗出血癌.
很絕望的.
另一個朋友在美國突然通知我.
她說她是第五期肝衰竭.
大約都是中年人.
又是有三個女兒.
得如其來是很絕望的.
我們是否要形容為嘔吐.

$^{241}$是沒有答案的.
我們的信仰在哪裡呢.
今天和大家回到聖經.
聖經怎麼說我們就怎麼相信.
我們的信德是從那裡來的.
我引用詩篇22章24節.
我們一起讀.
因為他沒有藐視憎惡受苦的人.
也沒有向他掩面.
那受苦之人呼求的時候.
他就垂停.
這個詩篇很特別.
是在舊約.
不過又被譽為十字架之詩.
如果大家看經文.
一開始就說這個很特別.
痛苦和讚美.
原來可以由困境絕望被棄絕.
去到讚美.
有出路.
所以詩人大衛在詩篇裡面.
他形容自己是一條蟲.
廣東話翻譯也很押韻.
第六節說我是蟲不是人.
有時發現是微不足道.
無能為力.
詩人大衛在這裡.
所形容的光景和他之前的一至二十二節.
完全不一樣.
又被藐視又被委屈.
甚至被犬類圍繞攻擊.
很絕望.
靠他的能力不能夠處理.
甚至他好像一開始的聖雁仔.
甚至埋怨.
究竟你是在那裡呢?.
為何叫十字架之詩呢?.
因為他除了突然間來.
我們的神不會藐視我們.
也不會厭惡受苦的人.

$^{281}$他會與我們同行.
他不會向我們掩面.
是否很積極?.
當我們求救的時候.
他會垂聽.
他會看見我們的需要.
但是原來第一節.
是一個一來就是我的神.
我的神為甚麼離棄我?.
當我們遇到生活的狀況.
遇到晴天霹靂.
遇到屙咀鼓勵不到.
不單止可能是發生在自己身上.
可能發生在社會那裡.
你見到宏福縣的事件.
不斷在沉澱的時候.
可能現在我們在接受之後.
我們不單止問為甚麼.
我們要問未解決.
否則不會有這樣的事件發生.
原來可能後面有很多結構性的.
依存的罪惡在裡面.
但是我們又可以怎樣被一起療癒呢?.
你會發現那個環境是痛苦到一個地步.
我的神為甚麼離棄我?.
這句說話是否很熟悉?.
這句說話是耶穌在十字架.
十架七弦裡面.
當他最痛苦的時候.
就算耶穌是神的兒子.
但是他被肉身的限制.
亦有救贖世人要完成.
所以我們現在在遇苦期.
你會發現我們需要思念.
和耶穌基督釘在十字架上的被棄絕.
這不是一個很優雅的形容.
這真是很真實的.
鼓勵大家去想一想自己困境的時候.
或者當下的你都在困境裡面.
又或者剛才我提到的.

$^{321}$在我身旁發生的事情.
無論是人我們都是無能為力的.
那個困境.
當你再在遇苦期裡面.
想起耶穌這句說話.
尚且他是道成肉身的.
上帝的兒子.
他經歷這種痛苦.
是否他不知道呢?.
其實他是在用詩篇.
二十二章第一節.
一開始我的神為甚麼離棄我?.
所以這篇詩篇.
也是很真實地說.
當時大衛遭逢的困境.
被誣衊.
被委屈.
被攻擊.
甚至裡面提到自己的骨都被截斷.
就是在這種光景裡面.
但為甚麼又叫做十字架之詩呢?.
因為耶穌用了大衛一開始.
真實狀況的呼求.
就是在二十二章第一節裡面.
但如果你看二十二章的標題.
它說用朝露的調.
彈的調要用朝露.
朝露即是早上.
早晨的那種露.
是否有些反差呢?.
露即是跳躍.
亦是早晨.
但亦是在痛苦裡面.
換句話說.
原來在痛苦裡面是有盼望的.
所以今天我們嘗試.
尤其是在禦苦期裡面.
我們可以看到耶穌在十字架上的痛苦.
但祂也賜我們盼望.
我們的信仰是有出路的.

$^{361}$至少祂告訴我們三件事.
耶穌在十字架上也告訴我們這三件事.
神沒有藐視我們.
神聾而眾知.
神知道我們發生什麼事.
神亦是掌管歷史的神.
神沒有掩面.
看到就當看不到.
我們在工作裡面有時看到很多挑戰.
看大不看小就當看不到.
在工作裡面也是.
有時老闆也是.
我作為別人上司的老闆.
也有上司就是我.
都要自我安慰.
一定會看到一些不合眼的事.
就當看不到.
我們這個神不是這樣的.
尤其是在十字架上.
祂也被棄絕那種呼求的感覺.
但祂一樣看著世人的需要.
我們的上帝是聆聽的上帝.
祂不會離棄我們.
尤其是在部節我們看到的盼望.
而這個盼望.
我們不妨去尋找.
我知道聽到這裡有些人說.
這些都是很抽象的.
你說神聽我.
什麼時候聽?.
你說誰聽?.
誰聽完也不見行動.
我想問問托爾斯泰.
不知道大家有沒有聽過這個故事.
托爾斯泰寫了《諧像瑪丁》.
同樣有另一個名稱.
除了叫《諧像瑪丁》之外.
還有《那裡有愛,那裡有上帝》.
托爾斯泰說.
愛簡單來說是通過三個故事.

$^{401}$愛不是用來說的.
愛是在行動裡面被感受.
只要我們能夠經驗愛.
就能經驗上帝.
《諧像瑪丁》托爾斯泰筆下裡說.
瑪丁很思念他太太.
很想親親太太.
因為太太離開了.
過世了.
於是他和耶穌有一個邀約.
約定了耶穌親自去捉摸他.
親自跟他聊天.
想通過跟耶穌聊天.
能夠思念他的太太.
於是那天瑪丁特別早起.
在等耶穌來.
當今天等耶穌來.
等的時候突然發現下大雪.
他在窗外望出去.
看到一個老人家.
這個老人家在雪地裡做什麼呢?.
就是在清理雪.
在清理雪的時候.
他看到這個老人家很生老.
又很冷.
他於心不忍.
自己暖暖的.
於是叫他進來.
給他一杯暖的奶茶.
暖暖他的身體.
接著這個老人家又出到雪地.
看一看手錶.
還早.
耶穌還沒有來.
在等候.
很敬虔的等候.
接著他又發現.
在雪地裡有一個女士.
抱著一個嬰兒.
走過.

$^{441}$又是很冷.
又於心不忍.
於是叫他進來.
給他什麼呢?.
首先把他忘妻的大衣.
很有紀念價值.
給他穿上.
因為覺得他不夠暖.
接著拿一個熱湯出來.
給他喝暖暖身體.
接著嬰兒開始哭.
拿一些奶奶出來.
補足一下.
接著很有暖意.
於是再上路.
很安慰.
已經去到中午.
看一看手錶.
耶穌又沒有來.
接著去到中午.
吃個下午茶.
見到外面開始日落.
有個女士賣蘋果.
賣蘋果有個攤檔.
小朋友挑皮.
走過去偷她的蘋果.
賣蘋果的老婦很不開心.
跑去罵她.
捉住她.
要報警.
馬丁又走出去.
又八卦.
勸架她.
不要吧.
給個機會給年青人.
這個年青人當然道歉.
跟著跟這個老婦說.
你給機會她.
跟著跟這個小朋友說.
她現在可以收檔了.

$^{481}$你順便幫她撿一下蘋果.
好啊.
兩人復和了.
又調解了.
接著回到家裡的時候.
已經是日落了.
耶穌有沒有來到?.
耶穌沒有來到.
跟著她看聖經.
她看聖經.
大家記不記得.
《馬太福音》那裡.
她一路看聖經.
看到聖經的《馬太福音》.
《馬太福音》那裡.
因為我餓了.
你們給我吃.
渴了你們給我喝.
我作客女.
你們留我住.
我赤身露體.
你們給我穿.
我病了.
你們看顧我.
我作監女.
你們來看我.
《馬太福音》第25章.
第35至36節.
當然這是耶穌的一個比喻.
跟著對談裡面問.
我什麼時候做過這些事.
跟著耶穌又再複述一次.
又是同樣這段說話.
在這個經文裡面.
或者看《蟹藏馬丁》.
耶穌有沒有出現過?.
耶穌出現了.
耶穌隱身了在這幾個人物身上.
又或者.
當然你可以倒過來.

$^{521}$你不用等候耶穌.
要看耶穌.
其實你已經可以做施與愛那個.
既是施與愛.
又可以去服侍愛.
服侍耶穌.
耶穌在哪裡?.
耶穌就在小子那裡.
在老父那裡.
原來正如托爾斯泰所說.
那裡有愛那裡就有上帝.
我想我們的社會裡面.
可能是欠缺了.
用愛來彼此觸碰.
觸碰那個困境.
或者觸碰那個孤獨的心靈.
耶穌用什麼方法呢?.
當然你可以在四福音裡面.
看到很多故事.
我隨意想到幾個.
耶穌去照門徒的時候.
都是看見門徒.
耶穌去傳福音的時候.
有很多人俗擁祂.
記不記得有一個瑞利薩戈.
是被很多人群嗤之以鼻.
因為覺得祂不是好人.
都是覺得祂不應該在裡面.
耶穌在人群裡面看到薩戈.
然後要去祂的家.
記不記得耶穌中午的時候.
去到井旁的日當空.
那裡有一個婦女.
是撒瑪利亞人.
當時猶太人和撒瑪利亞人.
有很多芥蒂.
猶太人更加藐視撒瑪利亞人.
但是耶穌和她對話.
為什麼那個婦女中午才來打水呢?.
我們可以推敲.

$^{561}$又是很多人看不起她.
所以不是人多的時候去.
是少人的時候去.
自然她自己都看不起自己.
其他人都看不起她.
因為她有五個老公.
現在那個也不是她.
很多人都覺得她很有問題.
成份不好 背景不好.
但是耶穌有沒有棄絕她?.
沒有 而且跟她說.
積著水來講生命之道.
甚至轉化了她的生命.
彼得即使三次不認主.
但是耶穌都沒有棄絕他.
甚至轉化彼得的生命.
所以彼得是殉道.
他說耶穌是釘十字架.
我不配 我要刀釘十字架.
一個被看見被愛的力量.
是可以很巨大的.
耶穌無時無刻.
就算自己被棄絕.
在痛苦裡面.
他仍然是看見人的需要.
只要我們謙卑自己.
好像十字架上的強度.
他的回應是什麼?.
今日你要同我在樂園.
那個不是doing來的.
是being來的.
我覺得服事當然是賺人熱淚.
我們要做什麼事工.
當然是很重要.
但是我們用什麼心態.
和隨時的幫助.
將耶穌那種愛.
能夠融入我們的生活.
無論是工作.
無論是社會.

$^{601}$無論是我們的人際關係.
家庭.
原來被看見.
或者我們簡單直接說.
被愛.
被愛的威力.
是超乎想像.
我不知道你最近那一次.
有沒有這種被翻轉.
重新的感受.
即使處境沒有被改變.
但你的心靈是重新上路.
這幅圖的背景.
其實都不是今年的.
是兩年前的.
或者是三年前.
那時候在疫情裡面.
困難,蜂擁,生意.
除了政府出手都沒有辦法.
我自己都跌入了心靈的幽谷.
我去了馬鞍山.
我也有這個習慣.
定期行山.
離開處境.
當我繼續做這個運動的時候.
那一次我去了馬鞍山.
如果你記得馬鞍山其實都很陡斜.
像陀螺一樣.
但為什麼我喜歡這些運動呢?.
因為去到那裡很絕望.
只有呼吸.
你不會想其他東西.
吃火鍋都要想.
下來再說.
只是呼吸.
跟自己聊天.
跟上帝聊天.
我將我的困難告訴他.
然後在我眼前.
我看到有些霧.

$^{641}$早上來的.
那條路看不清.
但我突然聽到聖靈跟我說話.
是很清楚的.
是廣東話.
祂說:使徒永富.
你已經很努力了.
不關你事.
跟我說:不關我事.
最後那句更加療癒.
祂說:你揹不起.
我再講多一次.
很簡單的這幾句.
使徒永富才對我說.
很欣賞我.
祂說:你已經很努力了.
第三句.
不關你事.
因為無論你多努力都改變不了.
第四句是什麼呢?.
你揹不起.
我們每天都很想改變世界.
化身成為超人.
但祂告訴我.
你揹不起.
我整個人降伏了下來.
雖然在眼前看到前路.
看不清.
但我知道一定有路.
我自己看不清.
原來這需要一種體會.
鼓勵大家.
愛的走動不只是宗教語言.
是屬靈的專利.
這本小書影響我很大.
我差不多20年前買涼茶的時候.
有朋友送給我.
《一碗清湯橋麥麵》.
在網絡上已經有電子版.
不需要買來看.

$^{681}$可能已經斷版.
裡面說到雪地裡的三濟那.
在一個橋麥麵的店舖.
晚上年三十晚十點打烊.
他們很緩慢地到達店舖.
店舖已經打烊.
太太有什麼幫到你.
可不可以做一碗橋麥麵給我們三濟那吃.
大家的反應會怎樣.
當然不會.
搗亂就早點來.
香港人現在九點半也找不到東西吃.
但這個橋麥麵店的老闆娘.
照舊透過腦再做這碗麵.
我想她的內心是有人吃我這碗麵.
這是我的天職.
重新透露.
為的是做這碗麵.
明年今天又再來三濟那.
是不是覺得很奇怪.
原來背後也有一個故事.
爸爸遭逢絕境困境去世.
媽媽很努力地長時間工作.
只有晚上才可以和兩個小孩.
一個四歲一個六歲.
又是年三十晚.
去吃一碗橋麥麵.
老闆娘和老闆聽到這個故事.
特別受感動.
老闆娘跟老闆說.
她老公這麼慘.
是不是應該做三碗給人家.
這個主意當然好.
但她老公說不要.
做三碗給人家.
令人覺得很難堪.
不過你放心.
其實我去年已經拿了一個特別大的碗.
好像火車頭一樣.
裡面放了三個麵.

$^{721}$怪不得今年又再來.
我劇透講到這裡.
來到這個位置也很感動.
是不是很樸素.
不需要那麼多計較.
這個鼓勵就是鼓勵家人.
後來他們走出了困境.
大兒子已經完成了學業.
他們三兒子特意回到橋麥麵鋪.
在雜房多謝老闆娘.
其實很陌生.
但原來就是這樣鼓勵他們好好活下去.
有沒有說過話呢?.
是沒有說過話的.
其實他們都不知道家人發生什麼事.
就是這麼樸素.
所以這個Being.
這個Being是可以承載我們的Doing.
不需要那麼多大道理.
所以德蘭修女也是這樣說.
她說偉大的愛就是從看到.
當然看到是否願意從見到.
來從做到.
能做多少就做多少.
剛才諧將馬丁也是這樣.
我讀的橋麥麵鋪的老闆娘.
一碗清湯橋麥麵.
多年前我去雜房玩.
是特意去找這間店舖是否存在.
據說已經轉了手.
據說是一個真實的故事.
求神幫助我們.
讓我們看得到.
也有很多事情未解決.
不過是有出路的.
大家有沒有興趣知道.
四眼仔二十多年後的今天.
在做什麼呢?.
是否在踢球呢?.
他是成為了健美先生.

$^{761}$再不是四眼仔那麼脆弱.
網上的都很幽默.
天不讓他去踢球.
做室內運動去健身.
我們的信仰是一個鼓勵不到的信仰.
痛苦神母讓我們每時每刻都可以彌撫得樂.
這不是一個超凡或免死金牌.
信耶穌就免於苦難.
沒有這回事.
但就像詩篇的二十二篇.
原來痛苦也可以有讚美.
就是神看到我們.
祂沒有掩面.
祂沒有離開我們.
祂讓我們跟到底 順到底.
神會刺盼望 刺出路給我們.
\newpage



\section{創世記 24:67}
\label{sec:031734}
\textbf{聖經人聖經事:得了安慰}~~~~黃克勤牧師博士
\newline
\newline
連結: \href{http://www.church.com.hk/acms/content.asp?site=cdc\&op=show\&type=product\&code=031734\&layout=sermon}{\texttt{031734}} ~~~~ 語音日期: 2026-02-20
\newline
\newline
\hyperref[sec:031733]{\small{< < < PREV SERMON < < <}}
~
\hyperref[sec:index_chronic]{\small{[返順時目]}}
~
\hyperref[sec:031735]{\small{> > > NEXT SERMON > > >}}
\newline
\newline
創世記 24:67
\newline
\begin{longtable}{cl}
\hline
\hline
章節 & 經文 (和合本修訂版)\\
\hline
24:67 & \begin{tabularx}{0.7\textwidth}{X} 以撒就領利百加進了母親撒拉的帳棚，娶了她為妻，並且愛她。以撒自從母親離世以後，這才得了安慰。 \end{tabularx} \\ \\
[1ex]
\hline
\hline
\end{longtable}
$^{1}$請見諒大家好,又來讀聖經了.
眨眼又一日.
生命樹宣教網絡.
又和大家一起讀聖經.
聖經人聖經時.
看看人與事背後的上帝.
他們和我們做了些什麼.
他是怎樣的.
更多認識他.
得了安慰就是今天的題目.
什麼得了安慰.
誰得了安慰呢.
這件事情故事還是和亞伯拉罕有關係的.
創世記二十四章第五十七節.
其實這裡是接續上次.
昨天說到什麼呢.
就是亞伯拉罕最信任的老僕人.
不過可惜這麼好.
是沒有名字的.
那人,那僕人.
不要耽誤我,我馬上要走了.
走了很遠千里迢迢.
回到亞伯拉罕的故鄉.
我估計要走一個月.
走這麼遠路.
去到一找到.
他不看別的.
只是完成他這件任務.
很直接的.
一完了之後找到.
住一晚.
第二天就說走了.
這麼遠去到.
不休息一頭半個月嗎.
我都說.
主耶穌時代的人去飲酒都會住一個星期.
何況他這樣他馬上走.
他很想快點完成任務.
很忠心的.
一路看著上帝怎麼帶領.

$^{41}$他就說我想走了.
不如打發我走.
你們讓我走吧.
說得好聽.
他們說.
沒有沒有.
利伯加的媽媽就說.
不如先住十天.
讓我和女兒住十天再聊聊天.
因為很難再見.
差不多不會再見了.
太遠了.
這個時候.
他們說我想回去.
這個人很堅持.
第57次.
他們說我們把女子叫來問問.
他問利伯加什麼意見.
就叫利伯加來.
利伯加來問他.
你和借人同去嗎?.
你是不是真的跟他走?.
他現在說要走了.
真的嫁給他.
認真的.
他走了.
其實這個問題.
同去即是現在馬上去.
這樣的意思.
因為.
禮物都收齊了.
難道現在說不去嗎?.
這樣還給他嗎?.
玩什麼?.
他們當然想著利伯加會.
哎呀!我也想住久一點.
很捨不得媽媽之類.
誰知利伯加就是撇撇脫脫.
我去,這是第一點要說的.
我去.

$^{81}$這件事不簡單.
因為他知道.
依夏一出去.
跟這個老僕人一去.
就已經不再回到故鄉了.
就和阿巴拉罕一樣.
不再回來了.
這些所有親戚朋友都沒有了.
不過我認為他回答得這麼直.
這麼簡單.
我去.
其實.
我相信他自己也想著看到.
上帝帶這個僕人來到.
然後在井旁相遇.
就遇到我了.
我給他倒水喝.
也給他倒在桌子上喝.
其實.
就是這樣而已.
然後在哪裡.
原來就是親戚來的.
於是.
他覺得這也是上帝預備的.
很有可能是這樣.
所以他很撇脫說我去.
其實李伯嘉不去也能嫁人.
一定能嫁人.
不過他說我去.
就是說.
他的心想.
這也是上帝的帶領.
他就很撇脫了.
我去這兩個字.
都是我們要學的.
很多時候知道整個心意淋到.
我不去先住多十天.
十天之後還是不捨得再住多十天.
然後再住多五天.
再住多三天.

$^{121}$再住多一天.
不如下個月吧.
其實他就很撇脫.
好吧.
我去.
於是他就預備了.
於是他們就打發妹子,李伯嘉和他的羽母.
妹子和羽母和阿伯克的僕人.
並跟從僕人的都走了.
打發不只是李伯嘉的人去.
還有一些妹子.
這個妹子.
就是跟著李伯嘉的人.
其實那些女孩是什麼人呢?.
如果看回這個妹子的希伯來文.
是可以翻譯為sisters.
就是很多姊妹.
是姊妹.
這個sisters.
究竟是真的家庭裡面的姊妹.
還是有些女的僕人呢?.
不知道,不過有些跟她一起的女孩.
有些跟她一起去.
並她的羽母.
這個羽母.
就是由希伯來文的字.
希伯來文的字.
就是.
「羽養」她.
Nursing 她.
Nurse 她長大的那個.
其實就是羽母.
即是說家裡是有錢人家.
帶著她,看著她長大的人.
又肯跟著她去.
於是就一起去了.
其實這樣也好.
新地方又不會一個人那麼孤單.
都走了.
第二點「都走了」.

$^{161}$想跟大家分享的.
就是有些時候.
你會發現到去到某一個位.
現在香港人很喜歡說去到某個位.
去到這個位.
這個位,去到這樣的位.
什麼位呢?.
即是說,去到這個時間.
這個遭遇.
那,都走了.
其實,很多時候.
我有個同學.
他最小的女兒結婚的時候.
他爸爸是牧師,我同學.
明天女兒結婚.
今晚,我同學的太太說她哭了.
牧師昨晚哭了,做什麼呢?.
不捨得女兒.
想不到我同學這位牧師哭了.
女兒出嫁了.
不捨得.
很多時候,兒子大了,去讀書了.
娶了媳婦了.
我不捨得兒子跟另一個人在一起.
這些不捨得,那些不捨得.
有時候去到某個位.
房子那麼大,不過家裡環境支持不住.
搬家小一點的房子,不捨得.
看著自己身體差一點,年紀大一點.
頭髮白一點,或者沒有一點.
很不捨得,以往常常向後望.
其實去到某些位,是要放手的.
都走了,第二點要跟大家說的就是.
有些時候,環境改變了.
關係改變了.
自己都改變了.
這些都走了.
走了就是走了,沒辦法了.
都要接受了.
如果你一直拿著不放手.

$^{201}$糟了,留多十天.
留多五天,女兒.
五天之後,三天,三天,兩天.
一天之後,一個月.
一月之後,不如女兒嫁個本地人.
禮物退回給他.
可能就是這樣.
所以其實.
在利伯家來說,第一點就是.
他知道常常帶你去.
他就去,我去,兩個字而已.
不說那麼多了.
而這些就是他的.
他的爸爸媽媽哥.
送他走,都走了.
就走了,好,放手.
Let go,放手,有些事情去到某個階段.
是會轉變的.
做什麼都走了,不要緊.
最重要就是給他們祝福.
於是為利伯家祝福.
然後送他走.
祝福什麼呢,我們的妹子.
願你作千萬人的母.
願你的後裔得著仇敵的城門.
哇,這麼厲害的祝福.
這兩句抄起來才行.
以後你嫁女兒的時候.
寫這兩句,願你作千萬人的母.
願你的後裔得著仇敵的城門.
意思就是.
第一句,願意你生養好多.
很重要的,因為以往的人都說了.
如果沒有子女.
覺得沒有賜福.
很不順利,很不好的,很差的.
所以一定要生,而且還生得多.
說到百子千孫.
以前中國人都是這樣.
百子千孫,現在一個都嫌多.

$^{241}$百子千孫.
願你作千萬人的母.
你的後裔得著仇敵的城門.
得著仇敵的城門之後.
你的對頭人都被你打低.
而且你佔領了他.
在他的城門那裡攻破了.
這個是很厲害的.
願意你生養眾多.
願意你事事順利.
勝利,就是這樣.
其實是個很好的祝福.
因為他們都走了.
記得了,有些事情.
去到差不多要放手.
就祝福,祝福而送走.
就不要攔住攔住.
打才走.
沒有鬥爭.
特別是一個家人.
既然走就祝福.
好了,繼續.
利伯加和他的侍女起來.
騎上駱駝跟著那一步人.
那一步人就帶利伯加走了.
這個時候.
以撒住在南地.
什麼叫南地呢?.
以色列是一條很長的.
北至南很長的.
中間點大概是耶路撒冷.
和死海.
因為耶路撒冷比較近死海.
再南一點.
希伯倫也是近死海的另一頭.
一個北一個南.
死海.
南地過了希伯倫.
再到死海以下.
就叫南地.

$^{281}$南地很清靜.
沒有什麼人住.
因為南地是曠野.
大部分都是曠野.
去南地.
以撒住在南地.
以撒長大後.
似乎沒有跟爸爸一起住.
當然有錢人家裡.
很多地方.
以撒住在.
比較南地曠野的地方.
住在那裡.
簡單清靜.
住在那裡.
以撒很一個人.
因為他的記載.
朋友不多.
家裡只有一個兒子.
沒有很多慶典.
其實還有以實瑪利.
是他的朋友.
不過他不常在附近.
以撒住在南地很安靜.
剛剛.
回來探望爸爸.
剛從.
比爾拉海內.
回來.
比爾拉海內.
是一個地方的名字.
比爾拉海內.
廣東話中文.
讀也有困難.
比爾拉海內是我第三點.
比爾是beer.
beer.
是井.
拉海.
Lakai.

$^{321}$拉海.
Lakai.
活的.
萊.
萊有時候.
讀萊.
萊.
的意思.
比爾.
Lakai.
萊.
比爾.
Lakai.
井.
拉海的意思是.
看顧我的永活.
者的井.
他們將.
名詞放在最早.
最主體的字.
井.
就是活著的.
的意思.
就是看顧.
我的永活.
者的井.
的意思.
有典故.
今天不說了.
明天再說.
天將晚以薩出來.
在田間脈上.
真厲害.
以薩的靈性不錯.
他很安靜.
他代表.
被殺的羔羊.
所以他也不錯.
起碼他不走.
信服.

$^{361}$以薩出來.
在田間脈上.
沉思.
在田間.
靈修.
默想上帝的作為.
沉思.
以薩當時比較沉.
心情比較低.
不知道.
是否有點情緒低落.
他沒有去到.
有抑鬱.
但不開心.
因為媽媽死了.
孤零零一個.
那麼.
好了.
舉目一看.
有些駱駝在來.
利伯嘉.
舉目看見.
急忙下了駱駝.
利伯嘉.
叫Ricard.
是個很厲害的女孩子.
無論身體健康.
起碼能夠.
弄三四十桶水.
給駱駝.
十隻駱駝.
一隻.
喝三罐水.
都可以.
三十桶水.
三十瓶水.
十個八個.
駱駝都喝八瓶.
四十瓶.
打水.

$^{401}$身體又健康.
又好靚.
聰明伶俐.
遠遠看見.
可能看出.
這個裝束有些不同.
衣服靚些.
或者有些.
有些不同氣質.
與風度.
一看.
莫非是這樣.
馬上看見駱駝.
急忙下了駱駝.
為什麼呢?.
如果是自己的丈夫.
沒理由自己.
站在駱駝上.
馬上下駱駝.
駱駝很聰明.
在田間迎接我們.
那個是誰呢?.
就是我們的主人.
馬上用拍子.
蒙上臉.
這是他們的禮貌.
第一次要見.
丈夫.
要蒙著臉.
中國有很多這樣的傳統.
蒙著臉.
在李伯嘉時代.
到洞房才看見.
新娘的樣子.
很美麗.
嚇暈了.
蒙著臉.
給丈夫驚喜.
夫人.
將所辦的一切事.

$^{441}$告訴二嫂.
遇到她.
是上帝為你準備的.
二嫂很開心.
便領李伯嘉.
進了他母親.
撒拉的帳篷.
媽媽撒拉的帳篷.
雖然媽媽死了.
但她的帳篷還在.
什麼是她的帳篷呢?.
每一個帳篷有什麼?.
很像.
以前的人住帳篷.
她不是住整個帳篷.
但長大後.
他們會分幾個帳篷.
好像不同的房間.
爸爸房間,媽媽房間,兒子房間.
媽媽不在.
但她的帳篷還在.
事情沒什麼改變.
原因是他們很掛念媽媽.
阿巴拉罕很掛念太太.
以撒很掛念媽媽.
因此.
媽媽的帳篷.
撒拉的帳篷.
還保存著.
她帶他進去撒拉的帳篷.
即是說.
她很看重這個人物.
將他帶到.
自己最愛的地方.
娶了她為妻.
是否一見到她.
就馬上去.
媽媽的帳篷洞房.
未必是.
沒特別說.

$^{481}$她說.
原來是我太太.
我都不知道.
不如你住在我媽媽的帳篷.
然後.
直接娶她為妻.
而且她很疼愛她.
以撒.
自從他母親不在.
這才得了安慰.
我只是補充一下.
以撒這個名字是什麼意思.
以撒是.
Isqar.
Isqar 怎麼讀成以撒.
其實是希伯來文.
Isqar.
也挺難讀的.
有時候.
用英文譯成.
Isaac.
怎麼譯成Isaac.
我也不太清楚.
不過譯成Isaac.
Isaac譯成中文.
就變成以撒.
I讀I.
是讀Y.
Y嘛.
但是.
I英文譯成中文.
就變成I.
可能是這樣.
其實以撒的希伯來文.
的字是.
他笑.
可能是撒拉生.
以撒的時候.
看到寶寶笑.
改什麼名字.

$^{521}$叫他笑.
以撒這個字.
就是他笑.
寶寶笑.
是不是.
Isqar.
自從他母親不在.
很傷心.
很不開心.
在田間哭了很久.
默想上帝.
好好沉思.
在這裡見到他太太.
李伯嘉.
得了安慰.
我很想大家.
能夠記得.
今天和大家分享的幾點.
溫習一下.
李伯嘉.
你是否願意跟他去.
神為我預備.
我去.
很好.
很撇脫.
等下先.
或者不去.
先通訊WhatsApp.
我去.
去到某個時間.
有些事都要放下.
時勢會改變.
祝福.
讓他出去.
第三.
有一位看顧我們.
永活上帝.
就是.
以撒住在那裡.
有典故.

$^{561}$明天再說.
以撒經常默想.
我都鼓勵大家多點.
田間默想.
今天和大家分享到這裡.
今天的金句就是.
以撒便領李伯嘉進與他母親.
撒拉的帳篷.
很掛念媽媽.
在那裡娶李伯嘉為妻.
愛她.
以撒自從母親不在.
得到安慰.
其實安慰是上帝.
特別為以撒.
預備的.
他喜歡的地方就是.
看顧活人的神.
的井.
神是看顧活人的神.
今天和大家分享到這裡.
我們一起祈禱.
讚美我們的主.
我們在這裡.
看到今天的聖經.
有很多好的榜樣.
第一次那個僕人不留戀.
馬上就走了.
去到甘縣只住一晚.
第二天就走了.
因為他見到任務完成了.
不想擔言.
李伯嘉就一聲我去.
他就起程了.
我們見到他的家人.
李伯嘉的家人就說.
既然你這麼有心要出發.
時間雖然改變.
環境改變.
女兒不在了,妹妹不在了.

$^{601}$要出發了,不過就為祂祝福.
我們見到以撒很安靜.
時時去田間默想上帝.
他在守住.
你是一位.
看顧人的神.
讚美你.
見到主撮合了這一對.
新郎新娘.
他們成為一對好的夫妻.
求主幫助我們.
頂住你自己的大力.
我們將自己很安穩的.
信靠的.
放自己在你的手中.
求主誰聽我們的祈禱.
幫助我們更有信心.
奉靠耶穌的名禱告.
在這裡我要補充一點.
因為.
有同事告訴我.
好像是前一次.
還是昨天.
我說錯了一個地方.
就是我說了拿學.
是亞伯拉罕的爸爸.
其實不是.
可能失魂.
其實拿學.
是亞伯拉罕的爺爺.
拿學生了其中一個兒子.
叫他拉.
他拉比亞伯拉罕更早.
帶亞伯拉罕和羅德.
一起離開.
本地本做富家吾爾城.
不過去到哈蘭.
住在那裡.
終於死在那裡.
亞伯拉罕繼續.

$^{641}$承繼父志.
出發.
帶著羅德去了.
迦南地神英雄的地方.
對不起說錯了.
同事提醒我.
拿學是亞伯拉罕的爺爺.
亞伯拉罕的爸爸.
拿學的兒子叫他拉.
讓大家知道.
對不起說錯了.
多謝大家.
\newpage



\section{箴言 5:15-23}
\label{sec:031735}
\textbf{快掙脫纏繞 當速速醒來}~~~~袁兆添傳道
\newline
\newline
連結: \href{http://www.church.com.hk/acms/content.asp?site=cdc\&op=show\&type=product\&code=031735\&layout=sermon}{\texttt{031735}} ~~~~ 語音日期: 2025-08-10
\newline
\newline
\hyperref[sec:031734]{\small{< < < PREV SERMON < < <}}
~
\hyperref[sec:index_chronic]{\small{[返順時目]}}
~
\hyperref[sec:031736]{\small{> > > NEXT SERMON > > >}}
\newline
\newline
箴言 5:15-23
\newline
\begin{longtable}{cl}
\hline
\hline
章節 & 經文 (和合本修訂版)\\
\hline
5:15 & \begin{tabularx}{0.7\textwidth}{X} 你要喝自己池中的水，飲自己井裡的活水。 \end{tabularx} \\ \\ \relax
5:16 & \begin{tabularx}{0.7\textwidth}{X} 你的泉源豈可溢流在外？你的河水豈可流到街上？ \end{tabularx} \\ \\ \relax
5:17 & \begin{tabularx}{0.7\textwidth}{X} 讓它們惟獨歸你，不可與陌生人同享。 \end{tabularx} \\ \\ \relax
5:18 & \begin{tabularx}{0.7\textwidth}{X} 要使你的泉源蒙福，要喜愛你年輕時的妻子。 \end{tabularx} \\ \\ \relax
5:19 & \begin{tabularx}{0.7\textwidth}{X} 她如可愛的母鹿，如優美的母羊，願她的胸懷使你時時滿足，願你常常迷戀她的愛情。 \end{tabularx} \\ \\ \relax
5:20 & \begin{tabularx}{0.7\textwidth}{X} 我兒啊，你為何迷戀陌生女子？為何擁抱外邦女子的胸懷？ \end{tabularx} \\ \\ \relax
5:21 & \begin{tabularx}{0.7\textwidth}{X} 因為人所行的道都在耶和華眼前，他察驗人一切的路。 \end{tabularx} \\ \\ \relax
5:22 & \begin{tabularx}{0.7\textwidth}{X} 惡人被自己的罪孽抓住，被自己罪惡的繩索纏繞。 \end{tabularx} \\ \\ \relax
5:23 & \begin{tabularx}{0.7\textwidth}{X} 他因不受管教而死亡，因極度愚昧而走迷。 \end{tabularx} \\ \\
[1ex]
\hline
\hline
\end{longtable}
$^{1}$《錄音》或《錄影》的版權是屬於基督教宣導會不各堂所有.
我們的屬靈傳統之一是以真理教導為本.
歡迎您將此《錄音》或《錄影》的連結傳給親友.
懇請注意,未經授權請勿轉錄或剪輯本《錄音》或《錄影》.
甚願好似《使徒行傳》二十章三十二節下所講.
這道能建立你們,叫你們和一切誠誠的人同得基業.
各位兄弟,愛與生命的源頭,今天我們來到你的庚前.
我們必須問自己幾個問題.
在我們的心中,我們還有沒有那份對你的忠誠和愛呢?.
我們在你的面前,我們願不願意放下自己,高舉耶穌?.
我們要問我們自己,我們看不看到耶穌基督給我們那份永恆的福分?.
幫助我們藉著我們今天的經文反思我們的生命.
也讓聖靈藉著你的話語更新我們的生命.
以上祈禱奉靠我主耶穌基督已經得勝的名字祈求,阿們.
因為當時已經在打八號風球,第一堂數五班要停運了.
然後風球不斷的升級到早上九點二十分.
天文台改發十號颶風信號.
當時我不知道你是否身處在安全的環境裡面.
我記得我身處家中外面是狂風暴雨.
我看到我家對面有一棟豪宅,那些豪宅是獨立泳池和健身室那種.
我看到對面的豪宅有個單位有一隻窗戶被強風扯開了.
然後一幅窗簾就扯了出來.
然後我再看清楚,同一棟豪宅另一個單位又有一隻窗戶被強風吹開了.
又一幅窗簾扯了出來.
是驚心動魄的景象.
但是夢神保守,雖然屋外面是狂風暴雨.
但是我自己回過頭來,我確實看到另一幅畫面.
我家中的袁冬甩悠然自得地在孵蛋.
她忠誠的丈夫曲奇守在她旁邊.
當時她眼前望到的就是會保護她的主人.
在冬甩身邊的就是那隻很忠誠的會餵她的老公.
因為她不用動的,原來愛情鳥如果感情好的話.
在孵蛋的時候老公會自動摸的.
不停地吞東西進肚子裡,然後就餵牠了.
就是這樣我們看到她在這個很艱難的情況.
好像很危險的環境裡面.
是那份愛,是那份忠誠的關係.
讓我們看到的小生命能夠在風雨飄搖的情況下.
安恕地在孵蛋.
同樣地我們看到今天經文開始的時候.

$^{41}$都是跟愛和忠誠是有關係的.
今天我們講到經文是出自箴言的五章第十五到二十三節.
是第一個段落.
然後我們會去到六章的一至十一節.
今天我們的主題是快爭脫纏繞.
當速速醒來.
今天經文你會發覺很特別的.
好像涵蓋了不同的範疇.
有講婚姻,有講借貸,擔保.
有講你要勤力工作這些的勸勉.
不過當我思考經文的時候.
我覺得裡面隱藏著對我們和上帝那份永恆關係的提醒.
好像提醒我們不要被世界在當中誘惑.
你不要被困在自我中心的思維.
你不要被短淺的視野阻隔了你和上帝的生命.
今天藉著這個經文我們會嘗試思考幾個主題.
第一,生命之道.
就是往眼外與衛生的命令直奔.
今天我們講的第一段經文是箴言的第五章十五到二十三節裡面.
有兩個很重要的命令.
飲池中的水是講述你要享受婚姻的親密關係.
第二個命令是愛你年輕的妻子.
是講述那份委身,愛,委身.
當我們發覺將這兩樣東西變成命令的時候.
很奇怪,命令你去愛,命令你去委身.
好像沒有彎著,沒有選擇.
經文似乎提醒我們.
世界有很多誘惑會令你和我失控的.
人與人之間的委身和關係.
是不斷被這個世界放縱情慾和高舉自己的文化衝擊.
同樣地,人與神的關係和委身.
那份文約也是備受挑戰.
第一段經文提醒我們.
愛和委身是命令.
是凌駕於我們的慾望和理性.
唯有守住這道命令.
我們能夠在世俗的洪流中義無反顧堅定不移的走諸道路.
今天講的第二段經文我給了一個題目.
叫做有種求生本能,叫做跪低.
今天講的第二段經文叫別人不要做擔保人.

$^{81}$別人借錢你不要做擔保.
以免別人還不到的時候.
你就備受牽連,陷入了債務的危機.
其實背後提醒我們.
很多時候我們做人都是自身難保.
不要不自量力.
更重要的是當你自身難保陷入困難的時候.
你要學習兩個動作.
第一,跪低.
第二,求救.
不過你有時會想.
其實我今天也沒怎麼欠債.
但是我們回到救恩這個層面的時候.
我會告訴你.
你和我在上帝的審判面前.
我們都是自身難保.
我們都虧欠了上帝的榮耀.
面對上帝我們要學習.
放下自己.
我們要跪低.
向神求救.
最後我們見到最後一段經文.
是關於懶惰人的提醒.
經文提醒我們不要只顧著睡.
不要面對眼前的危機.
視而不見.
去看看這個命令.
去看看那些螞蟻.
你要得智慧.
而我們在當中見到.
螞蟻是一種洞察環境的能力.
能夠回應環境.
預視冬天來臨.
並且行動為冬天作準備.
弟兄姊妹經文提醒我們.
我們要有一種目光.
我們要預視永恆的來臨.
我們要作準備.
不要失去那個豐盛的現值.
不要失去那個豐盛的生命.

$^{121}$而變得貧窮.
當我們進入今天的經文.
也是森嚴的第五章.
十五至二十三節的時候.
我們會相信這個經文是.
接續著第五章一至十四節的延續.
兩段經文都是講婚姻的.
前一段是叫我們遠離淫婦和情慾的試探.
到了今天這段的時候.
給我們一個很重要的命令.
如果你看到第十五節和第十八節.
兩個我在當中用紅筆.
也放大了一個字.
是一個命令的動詞.
你要喝你池中的水.
喝你井裡面的活水.
你要使你的泉源門風.
要喜愛你年輕的妻子.
這個講學自己的池水.
是講你享受婚姻裡面那種.
關於肉身的親密的關係.
而後面那一節.
喜愛你年輕的妻子.
是講你要在當中欣賞.
是愛.
只是愛一個.
這個是毀身.
只有享受婚姻.
在這段關係裡面毀身.
我們就能夠阻止到.
第五章第十六節的情景.
裡面講.
不要讓你那個泉流.
日流在外.
河水倒在街上.
是比喻這樣.
不要導致你的妻子離開你.
背叛你的婚民.
而走到街上.
穩定一段的關係.

$^{161}$在這裡我們發覺.
兩個重點.
不論是喝池中的水.
或者喜愛你的妻子.
你關於親密.
是愛和毀身的動詞.
都是用命令的字眼發出的.
不是勸你.
不是命令你.
不是給你參考.
這是命令你要做的.
你不能選擇的.
然後我們問.
上帝為什麼你做到這麼絕.
好像講到沒有彎轉.
好像很沒有自由.
不過如果我們看回前一段經文的時候.
他叫我們遠離淫婦和情慾.
再看回下面那段經文.
我而你為何迷戀陌生的女子.
為何擁抱愛邦女子的胸懷呢.
裡面用這個迷戀這個字.
是講那種失控的.
失控是愛到瘋 失控理性的.
原來人會迷戀陌生的女子.
然後被罪纏繞.
最後離開上帝 走向滅亡.
是的 這個世界很多的引誘.
多到讓你和我瘋狂的.
甚至我們說我們要控制自己.
我們知道不能做的.
不過我們忍不住.
我們在當中犯罪.
當中我只是舉兩個比較.
我們經常會遇到的引誘.
第一 在社會裡面.
鼓吹的那種追求自我掌控的驕傲.
我話事 我選擇 我的需要.
我要感覺良好.
在追求這些東西的時候.

$^{201}$我們會迷失在一個自我中心裡面.
我們看不到很多的東西.
第二 叫人專注在自己慾望之中的試探.
將慾望合理化.
有沒有見過聽你的廣告.
你想買就買 你想做就去做.
不用再想了 就是這樣.
滿足你的慾望.
見到那件東西 你喜歡買就買.
振興經濟.
正如我在之前的社關主任提到的訊息裡面.
我們說今天手機和社交媒體的文化.
其實就是在鼓吹這些元素.
我不知道你喜不喜歡手機.
我猜你和我都會經常滾手機.
滾手機的時候.
給我們一種很爽的感覺.
四個字 自我掌控.
看什麼我選擇.
和誰做朋友我選擇.
封鎖不封鎖他.
他煩我就封鎖他.
都是我選擇的.
甚至有時候你在家裡崇拜的時候.
什麼時候停過崇拜去廁所.
都是我選擇的.
另外我們會發現.
現在我們上的社交媒體.
這個演算法不斷將你和我喜歡的東西.
擺在我們面前.
去引誘我們要去拿.
我記得前陣子.
我要帶隊去柬埔寨.
短孫 短孫怎樣 你要機票酒店.
我當然在網上找.
如是者不找尤志凱.
你找完之後 天天在我的電腦和電話裡面.
說柬埔寨金邊.
這些機票便宜.
這間酒店又很漂亮.

$^{241}$那個東西又很划算.
不斷叫你買.
很划算.
如果喜歡.
打手機遊戲的弟兄姊妹.
你和我都可能知道.
手機遊戲背後.
有一個心理學團隊.
是心理學家團隊在設計這個遊戲.
針對你的慾望.
想想怎樣用音樂和畫面.
去刺激你繼續玩下去.
刺激你的慾望.
世界是很多讓人迷戀的慾望.
包括.
經文裡面提到的那些陌生女子的情慾.
所以我們會見到很多人在情慾上跌倒.
包括基督徒.
其實我們知道的.
很多時候我們知道理性上是不可以的.
不過我控制不了.
我的慾望太重要.
我要滿足我的需要.
然後當我說滿足我們的需要的時候.
情有可原四個字.
成為我們犯罪的遮醜布.
結果同居.
分外情.
成為性濫交.
成為了這個城市的日常.
到了今天.
當我前陣子看回香港統計處的數字.
我發覺.
現今我們這個社會裡面.
愛和衛生好像不是很重要.
根據我自己看回.
統計處的數字.
過去十年香港的離婚率仍然持續高距.
我不知道你知不知道現在離婚率是多少.
大概.

$^{281}$每兩對結婚就有一對離婚.
所以我們這裡有一半.
離了婚也是很危險的.
另外持婚和未婚的人口.
是持續增加的.
這些數字.
可能是告訴我們.
人對婚姻好像失去了信心.
面對這個高舉慾望和自由的世界.
我們看到.
神是用真言在當中.
用命令去回應.
就是那個愛.
衛生的命令.
不駕我們的理性和慾望.
不是說我考慮一下愛不愛.
你就要去愛.
我考慮一下還繼續不繼續衛生.
你是要去做.
唯有守住這些命令.
我們才能夠在世俗洪流裡面.
義無反顧的.
堅定不移的.
走主的道路.
當然我知道我們在講婚姻裡面的溝通.
和了解是很重要的.
不過當我們吵架的時候.
我們不認識對方的時候.
我們仍然有勇氣在當中.
選擇接納對方.
選擇溝通.
是愛.
愛的命令.
超越慾望和理性.
經文告訴我們.
當我們守住愛和衛生的關係的時候.
人就能夠享受到那份溫柔和溫暖.
生命就好像給活水那樣滋潤和豐富.
所以經常裡面你見到妻子是用池水.
活水在當中說滋潤豐富你的生命.

$^{321}$但我要講一件事.
很多人的處境不同.
可能有些人在婚姻裡面面對著很多的艱難.
這裡是講一個普遍的情況.
甚至乎在座有弟兄姊妹說.
其實我單身,這段經文不是很關我事.
不是啊,上帝設立婚姻是要藉著婚姻去預表一種愛與衛生的關係.
是告訴我們,人可以尋求上帝有一種愛與衛生和親密的關係.
弟兄姊妹,我不知道你們有沒有想過.
耶穌基督因著愛我們的緣故.
一早已經向我們傳言為生.
耶穌基督因著愛我們的緣故.
一早已經向我們傳言為生.
祂在十字架上毫無保留地將自己的性命和尊嚴.
為了我們全部不要.
愛與衛生.
我知道婚姻裡面總會有不似預期的情況.
同樣地.
我相信你相信了耶穌.
你和耶穌那份親密的關係也有不似預期的情況.
可能你和我會抱怨.
為什麼耶穌基督不聽我的祈禱.
為什麼上帝的命令你這麼難去遵守.
因為我好辛苦好多的掙扎.
或者我們會問.
為什麼上帝容許我在一個這麼差的環境裡.
我們會質疑.
不過本著愛與衛生的命令.
我們義無反顧地跟隨耶穌基督.
然後.
見證耶穌基督.
怎樣豐富你和我的生命.
今年我結婚第七年.
在這裡我必須承認.
我不是什麼婚姻戀愛專家.
我的感情史悶到飛起.
三言兩語說完了.
拍過一次拖結過一次婚.
我老婆就是我初戀情人.
然後就沒有了.

$^{361}$結了婚之後.
我和我太太豐富了我的生命.
我不是很去旅行的.
太太就抓住我.
帶我好像打衝鋒那樣.
去很多的地方.
不過我記得在裡面的經歷的時候.
我們都知道的.
旅行沿途景色花花世界.
但是我們最重要的就是.
兩個人一起衝去趕火車趕飛機.
一起走.
有時婚姻裡面.
總會有不似預期的事情會吵架的.
他很多時候命令我要約束自己的舌頭.
不要說一些和做一些.
傷害對方也不能回頭的事情.
有時我們會發覺.
在我們的關係裡面.
我們總有一些缺點.
沒有辦法能夠改變.
是搞不定的.
在這裡我必須.
懺悔和承認.
在洗碗的事情上.
我會洗鍋洗碗和洗鍋.
不過我總是忘記洗電飯煲.
這是我的陋習.
而每當我太太捉住我這個痛處的時候.
跟我說 袁紹添.
你又忘記洗電飯煲了.
我唯有看著家裡的鸚鵡布丁.
我像他那樣露出一個無辜的眼神.
和反應側著頭跟他說.
我太太就會笑到.
沒有動力繼續去罵我.
第一次明明是愛.
是衛生 是生命之道.
學習愛上帝.
愛你的伴侶.

$^{401}$珍惜和上帝.
或者和伴侶那份真誠的關係.
朝著愛和衛生的命令.
直奔.
有些認識我的弟兄姊妹知道.
我童年在一個龍蛇混雜的社區裡.
讀書和生活.
那個社區裡有很多閒雜的人等.
所以我父母就不放心.
我們小學上課.
都家長接送.
我很記得有一次.
我跟著父親回家.
當時經過一個貨倉.
父親通常走前面.
我是走後面.
突然有個醉酒佬衝出來.
他捉著我的手不讓我走.
當時小時候我很想掙扎.
想掙脫他的手.
但當時我就感覺到原來小朋友.
和大人的力量上有相當的差距.
我動不了我的手.
完全被壓制住.
我爸爸沒有發言.
當時我家裡有個家訓.
就是向人求救是弱者的表現.
是沒面子的事情.
不過當時我看著.
越走越遠的父親.
被人捉住我很害怕.
於是我終於開口大叫.
我要保命.
很感恩那個醉酒佬.
原來是爸爸附近的工友.
當時我爸爸就嚇他放手.
你放手呀.
你不要再騷擾我兒子.
那個人就放手.
原來有時.

$^{441}$跪下求救是可以保命的.
當然我要看誰.
你要找一個.
值得你在當中求救的人.
今日去到第二段經文.
是講僧賢的六章一至第五節.
在當中.
第六章一節是講我兒.
你若為朋友擔保.
替陌生人擊掌.
就在當中叫他們不要這樣做.
不要做擔保人.
因為當時很多人.
在當中.
不知道擔保人要付利息.
身邊的朋友還不了債.
他們就拍拍胸口.
然後發覺自己陷入財務危機.
經文提醒我們.
很多時我們其實是自身難保.
不要不自量力.
當我們看回經文的時候.
我們發覺裡面有兩個.
大概兩個很重要的命令.
大概.
跪下求救.
我兒因你的樂在朋友的手中.
只有一樣行財.
可以救自己.
你要在當中謙卑.
採取行動.
去懇求你的朋友.
亦都是你當時的債主.
去求他放過你.
謙卑這個字是屈服的意思.
希伯來文都幾精彩.
這個字是.
不停踩自己.
不停踩自己踩到地上.
如果他配合懇求一起用的時候.

$^{481}$我就用了一個比較現代化的詮釋.
跪下.
救自己.
這是一個命令.
或者你說今天我沒有欠債.
甚至乎我生活很富足.
外面還有很多人.
我身邊還有很多人欠我食飯和早餐錢.
但當我們回到上帝的救恩的時候.
其實.
你和我面對著上帝的審判.
我們是自身難保.
因為聖經說我們都.
虧缺了上帝的榮耀.
你和我都是欠了上帝的債.
所以死亡和審判.
成為了你和我的結局.
剛剛我又教完.
快要教完一季的西禮班.
長者西禮班.
長者西禮班裡面老人家都要做一件事.
就是拿一張紙.
我們叫認罪樹和靈界接觸表.
在當中將以前犯過的罪.
將以前拜過的偶像.
在耶穌基督面前列出來.
圈住他然後祈禱.
求耶穌基督赦免.
向神跪下.
敬罪悔改.
脫離上帝的審判.
這是唯一的方法.
是你我信仰的原點.
前陣子整理.
私人文件的時候.
我太太見到一封信.
有我的名字在那裡.
她未見過不知道是甚麼來的.
可能有些弟兄姊妹會知道.
我跟太太解釋.

$^{521}$大概當你讀大學的時候.
假如你學院的院長不小心發現你是全級考第一.
他就會寄一封信給你.
然後將你的名字刻在金屬牌上.
掛在辦公室的門口.
這封信提醒我.
我是一個很平庸的人.
當年我在大學讀書.
很努力讀到沒有朋友.
那時候為了追求成績.
壓力大到整個臉都暗瘡.
還暫迷.
當時有時要找輔導在當中.
幫我開解一下.
但事實上這封信.
對於我家人來說.
我家人是看都不會看的.
因為我另一個家人.
每次公開考試都幾乎會上電視.
他的獎學金.
當年的獎學金是多到可以.
買樓給首期.
這樣的比較.
追求不到的.
當時我覺得很累.
看著這封信.
我知道多努力.
很辛苦才拿到這封信.
但我維持不到這個成績.
當時我記得小學讀書的時候.
聽過一位耶穌.
於是我走回教會.
我跪下向耶穌求救.
我記得我走入教會崇拜的時候.
上帝的經文提醒我.
你不需要和別人鬥.
你不需要和別人爭.
你可以有耶穌基督的愛.
和永生.
放下自尊.

$^{561}$跪下向耶穌.
不是向AI.
不是向Google大神求救.
是向耶穌求救.
因為我們有一種.
尋求永生的本能.
叫做跪下.
不是向別人跪下.
是向耶穌跪下.
在今年的六月.
以色列向德克蘭的政權.
行動.
以F-35I為首的.
以色列空軍向德克蘭的政府.
高層和防空導彈帶來了.
毀滅性的打擊.
這場空襲的關鍵勝負在乎.
看得到和看不到.
你途中看到這架F-35I.
不單很難在雷達上.
被看到.
它在雷達上只是一隻小鳥.
那麼小.
憑著信息鏈的技術協助.
F-35I有一個很強的能力.
很快就能掌握到.
方圓幾百公里.
每一個角落的情況.
所以德克蘭革命衛隊.
藏起了防空系統和導彈.
一早看得一清二楚.
按下按鈕.
空對地的導彈就將這些東西摧毀了.
相反德克蘭革命衛隊的S-300.
和S-200防空雷達.
完全看不到F-35I.
最後只能夠站著任人打.
如是者當防空系統.
瓦解了之後.
就能夠如入無人之境.

$^{601}$去擊殺革命衛隊的高層.
和那些導彈部隊.
弟子們原來在.
真戰裡面.
看得到 看不到.
可以有截然不同的結局.
難怪聖經常叫我們.
看清楚 看呀.
弟子們你看不看到.
你看不看到你屬靈的危機.
你看不看到上帝要你走生命的方向.
還是你說.
其實我在上帝面前.
我什麼都看不到.
我什麼都聽不到.
箴言第六章第六節.
有一個很重要的命令.
懶惰人.
你去拆開螞蟻的動作.
就可以得智慧.
經文叫懶惰人.
睜大眼睛去看一下螞蟻的生活.
當中是有智慧在的.
有些人說聖經一定是叫我們.
學螞蟻這麼勤力.
不過我見到.
螞蟻在這裡更加有一個東西.
叫做危機意識的觸覺.
螞蟻知道冬天會來.
相反我們會見到.
懶惰人對危機視而不見.
你見到經文六節到第十一節.
經常說懶惰人在這裡睡.
經文說你不要睡.
你不要睡.
神叫懶惰人要醒.
打開他的眼睛.
否則他們豐盛的生命就會被奪取.
打開他的眼睛.
開啟那個永恆的觸覺.

$^{641}$這個觸覺是甚麼?.
對神的學務.
這個觸覺是甚麼?.
找智慧.
甚麼是找智慧?.
就是知道上帝的審判會來到.
所以我們見到.
真言第六章六節有個命令.
你要去怎樣?得智慧.
而這個智慧就是我們.
九章第十節.
我們一起讀.
敬畏耶和華是智慧的開端.
認識至聖者便是聰明.
今年港九培靈園經大會.
講完曾錦發牧師.
給了我們一個很好的提醒.
他提醒我們今天你和我.
很多人都很愛主.
不過今天很多基督徒跌倒.
不是因為他們不愛上帝.
很多都很愛上帝.
很多基督徒跌倒是因為.
他們不敬畏上帝.
敬畏上帝.
學務上帝.
剛剛過去的七月.
我參加了一個.
很多年青人都會聚集的活動.
那個活動在會展那裡.
當時我一早去了會展的會場.
天文台發出了黑色暴雨警告訊號.
意味著這個活動要取消.
於是我和天父說.
我很想進去看看.
決定繼續.
因為黑色暴雨的關係.
平時人流少了.
人家花了幾個小時排隊.
我花了三十分鐘就進了會場.

$^{681}$當時我看到很多.
打扮成動漫角色的年青人.
樣子很冷酷.
你看到年青人其實他們的樣子.
不會主動和別人交流.
不過我心想其實你扮得這麼漂亮.
扮得這麼好你也想別人注意.
於是我就鼓起勇氣.
其實我很害羞.
鼓起勇氣去找他們拍照.
和年青人.
用這個做漫畫.
你們扮得很漂亮.
我可不可以和你們拍一張照.
他們.
你找他們他們絕對會答應的.
雖然他們沒有笑.
但我感覺到他們很開心.
他們很渴望和別人有連結.
我知道他們骨子裡.
很想和別人.
有一份關係在當中連繫.
我在當中求神.
在當中開他們的眼睛.
因為他們真正需要的是.
和別人的連結和上帝的連結.
不過今天很多東西.
阻擋了我們的視線.
求神開了我們的眼睛.
我們今天.
都失去了永恆的觸覺.
年青一代都失去了永恆的觸覺.
求神幫助我們的下一代.
看到上帝.
求神幫助我們.
我們年輕人應該怎樣被關心.
然後回到永恆和愛上帝的關係.
回到今天的提醒.
生命之道就是往往愛與畏身的命令直奔.
在今天的第一段論經文提醒我們.

$^{721}$世界有很多誘惑.
憑我們的理智.
我們的計算.
我們抵抗不了.
愛與畏身很容易迷失.
很容易失控.
人與人之間的畏身.
面對著情慾文化的衝擊.
人與神的愛與畏身.
亦面臨著自由主義和個人主義的挑戰.
很多東西我們抵擋不了.
不過上帝的命令就是.
愛畏身.
義無反顧堅定不移的走主的道路.
今天第二段經文提醒我們.
其實你和我都是自身難保.
不要不自量力.
更重要的是.
你自身難保的時候.
你要懂得跪下求救.
找誰呢?.
向耶穌跪下求救.
因為在上帝的審判面前.
你和我都是自身難保.
面對上帝.
放下自己.
我們學下跪下向神求救.
最後一段經文提醒我們.
我們要有那種洞察能力.
我們要有那種永恆的觸覺.
敬畏上帝.
因為上帝的審判好像冬天一樣.
要來臨.
不要對上帝的審判視而不見.
不要讓我們的生命失去了那份豐盛.
變得貧窮.
特別是我們下一代的生命.
求神憐憫我們.
我們祈禱.
賜下婚姻.

$^{761}$你藉著婚姻.
預表什麼叫愛和毀身.
你親自賜下.
那位愛我們.
為我們毀身的愛子.
讓我們得到上帝的愛.
今天世界有很多東西在抵擋你.
叫我們離開耶穌基督.
是愛和毀身的命令.
幫我們持守住.
讓我們堅持走主的道路.
今天很多人提醒我們.
我們要維持自己的尊嚴和權利.
但是經文提醒我們.
在上帝的面前.
就算你有尊嚴.
就算你有權利.
你都是自身難保的.
幫助我們.
我們跪下求上帝的赦免.
求祂拯救我們的生命.
今天很多東西遮蓋我們的眼睛.
讓我們看不到上帝.
甚至我們對上帝的來臨.
和審判我們視而不見.
開啟我們屬靈的眼睛.
那份永恆的觸覺.
讓我們追求上帝.
敬畏上帝.
開啟我們下一代.
那個屬靈的觸覺.
讓他們見到.
他們真正渴望的.
是和永生上帝有連結的關係.
以上祈禱奉耶穌名求.
阿們.
\newpage



\section{路加福音 17:11-19}
\label{sec:031736}
\textbf{從醫治到拯救}~~~~梁成裕牧師
\newline
\newline
連結: \href{http://www.church.com.hk/acms/content.asp?site=cdc\&op=show\&type=product\&code=031736\&layout=sermon}{\texttt{031736}} ~~~~ 語音日期: 2026-03-01
\newline
\newline
\hyperref[sec:031735]{\small{< < < PREV SERMON < < <}}
~
\hyperref[sec:index_chronic]{\small{[返順時目]}}
~
\hyperref[sec:031737]{\small{> > > NEXT SERMON > > >}}
\newline
\newline
路加福音 17:11-19
\newline
\begin{longtable}{cl}
\hline
\hline
章節 & 經文 (和合本修訂版)\\
\hline
17:11 & \begin{tabularx}{0.7\textwidth}{X} 耶穌往耶路撒冷去，經過撒瑪利亞和加利利中間的地區。 \end{tabularx} \\ \\ \relax
17:12 & \begin{tabularx}{0.7\textwidth}{X} 他進入一個村子，有十個痲瘋病人迎面而來，遠遠地站著， \end{tabularx} \\ \\ \relax
17:13 & \begin{tabularx}{0.7\textwidth}{X} 高聲說：「耶穌，老師啊，可憐我們吧！」 \end{tabularx} \\ \\ \relax
17:14 & \begin{tabularx}{0.7\textwidth}{X} 耶穌看見，就對他們說：「你們去，把身體給祭司檢查。」他們正去的時候就潔淨了。 \end{tabularx} \\ \\ \relax
17:15 & \begin{tabularx}{0.7\textwidth}{X} 其中有一個見自己已經好了，就回來大聲歸榮耀給神， \end{tabularx} \\ \\ \relax
17:16 & \begin{tabularx}{0.7\textwidth}{X} 又俯伏在耶穌腳前感謝他。這人是撒瑪利亞人。 \end{tabularx} \\ \\ \relax
17:17 & \begin{tabularx}{0.7\textwidth}{X} 耶穌回答說：「潔淨了的不是十個人嗎？那九個在哪裡呢？ \end{tabularx} \\ \\ \relax
17:18 & \begin{tabularx}{0.7\textwidth}{X} 除了這外族人，再沒有別人回來歸榮耀給神嗎？」 \end{tabularx} \\ \\ \relax
17:19 & \begin{tabularx}{0.7\textwidth}{X} 於是他對那人說：「起來，走吧，你的信救了你！」 \end{tabularx} \\ \\
[1ex]
\hline
\hline
\end{longtable}
$^{1}$今天分享的題目是從醫治到拯救.
我不知道你有沒有試過收到很開心的禮物.
但是你收完之後.
我小時候這件事情就常常發生在我身上的.
因為其實我爸爸就很疼我們全家.
我爸爸也很疼我.
但是小時候爸爸做生意很忙.
所以我不是經常看到我爸爸.
但是每一次爸爸在家裡面的時候.
他也會用不同的方式讓我們去經歷愛.
我很記得有一次我快要生日了.
於是我媽媽就問我.
你想要什麼生日禮物?.
我就跟媽媽說我想要一架玩具飛機.
過了一段時間.
有一天吃完早餐.
我媽媽就說.
你進爸爸的房間吧.
爸爸要跟你說.
小時候不是常常看到爸爸的.
通常要進爸爸的房間.
也是我做了一些錯事.
所以我就說.
進爸爸的房間?.
我又回想一下最近有沒有頑皮過呢?.
有沒有做過錯事呢?.
帶著很擔心的心進了爸爸的房間.
進了爸爸的房間之後.
爸爸就跟我說.
他說媽媽跟我說.
你很想要一架飛機做生日禮物.
這架玩具飛機是送給你的.
那一刻我又擔心變為極度開心.
於是一開門那一刻.
哇!飛機啊!.
想著那架飛機.
然後飛走了.
小時候也不是常常看到爸爸回來跟我們吃飯.
因為他做生意很忙.
但是有一晚我記得他故意回來跟我們吃飯.

$^{41}$他回來的時候拿著一大袋禮物.
是什麼呢?.
原來是當時的兒童節目的錄音帶.
他說我今天帶了很多禮物給你.
一拿出錄音帶.
哇!跳飛機啊!.
有沒有人記得跳飛機?.
多謝你多謝你.
多謝你Tony.
多謝你多謝你.
我也明白的.
你記得跳飛機應該跟我同一個年代了.
然後不僅有跳飛機.
然後拿出來.
你看看.
醒目哥啊!.
如果你開始發現有些錯誤.
證明你跟我同一個年代了.
不是啦!不是醒目哥啊!.
是醒目仔啊!.
有沒有人記得醒目仔啊?.
有誠實的人在這裡.
然後我們很開心.
我仍然記得我們三姊弟一起跟我爸爸說.
不是醒目哥啊!.
是醒目仔啊!.
我們還說是醒目仔嗎?.
於是我就拿出了跳飛機.
跳著跳著跳走了.
又不記得多謝我爸爸.
我記得媽媽有一次跟我說.
她說為什麼你每一次收禮物.
你都不記得多謝你爸爸.
她說其實爸爸心裡面不是很想要你一句言語上的多謝.
但是原來爸爸是很渴望跟你有這一段的關係.
我媽媽給我這個教導.
到今天我都記得.
所以今天我仍然經常說多謝我爸爸.
因為爸爸做得很好.
他經常請我吃飯的.

$^{81}$他請我吃好東西的時候.
我就會說多謝爸爸.
但是我知道.
不只是一句公式.
是因為一段的關係.
但是當我們回想我們的屬靈生命的時候.
可能我們的屬靈生命也是這樣子的.
我們很多時候會懂得向神求.
但是我們又不懂得怎麼去回應.
很多時候我們的生命.
我們有不同的需要.
為我們的工作.
為我們的家庭.
為我們的感情.
為我們的健康.
很多時候我們都會很迫切的祈禱.
可能我們也會祈禱.
主啊!求你讓我進入這份工作.
主啊!求你醫好我的身體健康.
主啊!求你讓我這段關係可以和好.
我們也會這樣祈禱的.
但是當神真的開路了.
我們考試過了.
工作又找到了.
病又康復了.
問題又解決了.
但是接著我們會怎麼樣呢?.
我們很多時候都會繼續投入新的生活.
新的忙碌.
新的工作.
新的關係.
新的安排.
慢慢我們已經忘記了神了.
我們可能會少了回教會.
少了禱告.
少了親近神.
接著發現可能我們的生命.
我們只是有需要的時候才來找耶穌而已.
事成之後就各散東西.
我們發現可能我們的生命.

$^{121}$原來我們只是找神的幫助.
但是我們未曾和神進入一個真正的關係.
所以今天這個題目.
從醫治到拯救.
我們會去看一段很熟悉.
但是很值得我們生命反省的經文.
十個麻瘋病人得醫治.
但是只有一個轉身回來.
俯伏在耶穌面前.
歸榮耀給神.
我們今天我們的生命.
都可能經歷過神的醫治.
神的幫助.
神的帶領.
或者神的開路.
但是我們要問自己.
今天我們有沒有和神建立一段真正的關係.
所以今天很想藉著今天的經文.
我們一起去學習三個很重要的功課.
我們一起去彼此鼓勵.
今天第一個的功課.
是明白我們的生命需要主.
我們一起讀《路加福音》十七章十一至十三節.
我們一起讀.
一,二,三.
耶穌往耶路撒冷去.
經過撒瑪利亞和加利尼.
進入一個村子.
有十個長大麻瘋的迎面而來.
遠遠地站著.
高聲說.
耶穌夫子可憐我們吧.
當我們看這段聖經的時候.
我們首先看看這個背景.
耶穌往耶路撒冷去.
經過撒瑪利亞和加利尼.
在這裡帶出了兩個重點.
耶穌經過這個地方.
他看到一群被人忽略的人.
但是同樣這段聖經提醒我們.

$^{161}$耶穌去耶路撒冷是有著祂自己的使命.
耶穌去耶路撒冷是要在那裡釘十字架的.
但是耶穌為我們釘十字架.
就讓我們.
包括一群被忽略的人.
去得著救恩.
而耶穌在這裡遇見了十個麻瘋病人.
在當時來說是一個很真實的寫照.
因為當時的醫學並不昌明.
而乾旱的氣候也助長了很多皮膚病的生長.
所以麻瘋病是當時一個很普遍的問題.
在這裡你可能覺得麻瘋病跟我無關.
我不需要耶穌.
但其實麻瘋病人的需要不單單是身體的需要.
我們一會兒會看到麻瘋病人不單止身體有需要.
他的情感有需要.
他的情緒有需要.
甚至他的屬靈生命也有需要.
今天我們可能身體也有需要.
但今天可能說我身體很健康.
但我們也有情緒的問題.
我們可能會有情感的問題.
我們也有屬靈生命的問題.
我們需要主.
在這裡到了第十二節.
耶穌進入一個村子.
有十個長大麻瘋的迎面而來.
遠遠地站著.
說麻瘋病當時是被社會視為不潔的人.
他們要跟其他人分開.
要住在城外面沒有人住的地方.
在民數記五章二至三節裡面這麼說.
他說你吩咐以色列人.
使一切長大麻瘋的患陋症的並因死屍不潔症的都出營外去.
無論男女都要使他們出到營外.
免得污穢他們的營.
這營是我所住的.
看完這段經我們又發現.
身為一個麻瘋病人.
他有的不單止是身體的問題.

$^{201}$他有的更加是情感的問題.
因為他們是被人隔絕的.
他們有的是情緒的問題.
因為他們是被人歧視的.
而他們有的也是屬靈的問題.
因為他們是不能夠在聖殿敬拜神.
所以他們是被人厭棄.
沒有群體生活.
也沒有敬拜生活.
所以聖經說他們只能夠遠遠地站出來.
因為他們知道自己不潔.
而我們一會兒更加知道.
在這一群人裡面.
有一個人是撒瑪利亞人.
這個不單止是一個麻瘋病撒瑪利亞人.
更加是被猶太社群厭棄的人.
而這一群人當他們見到耶穌的時候.
第十三節告訴我們.
他們高聲說.
耶穌,夫子可憐我們吧.
有學者認為高聲說.
是一種向神呼求的格式.
但也跟當時可能律法的需要有關.
原來當時的律法說什麼呢?.
如果我們看《利未記》十三章.
四十五至四十六節.
原來當時如果有大麻瘋的人是要怎樣的呢?.
他說:身上有長大麻瘋災病的.
他的衣服要撕裂.
也要縫頭散髮.
蒙著上唇.
喊叫說:不潔淨了!不潔淨了!.
災病在他身上的日子.
他便是不潔淨.
他既是不潔淨.
就要獨居營外.
當我們去看這段聖經.
我們就明白他們當時可憐的狀況.
他們不單止已經身體有需要.
他們的情感有需要.

$^{241}$他們被人隔絕.
失去了群體.
失去了敬拜的生活.
不單止如此.
這段聖經更加說.
他們要怎樣呢?.
他們要衣服撕裂.
頭髮蓬鬆.
去到哪裡都要告訴別人.
我是不潔淨的!.
你試想一下.
如果你今天患了感冒.
然後我們在這個群體裡面說.
如果你患了感冒.
你要撕裂衣服.
頭髮蓬鬆.
到處去都要告訴別人.
我有感冒啊!.
不要走過來!.
這就是他們的苦況.
但是感冒是會康復的.
麻瘋病.
在那時候是不能醫治的.
所以我們就明白.
當時這些人的苦況.
他們是怎樣呢?.
被標籤.
被排擠.
被隔絕.
所以他們的出路是什麼?.
他們的出路就是唯有跟耶穌說.
耶穌!夫子!可憐我吧!.
通常在路加福音.
什麼人會稱呼耶穌為夫子.
就是他的門徒.
因為門徒經歷過耶穌的工作.
而這些麻瘋病人都稱呼耶穌為夫子.
他說主啊!求你的工作能夠怎樣?.
成就在我們的身體上.
所以其實這個可憐.

$^{281}$是表達一個人在絕望之中的呼求.
表達他生命需要主.
但是原來今天.
我們的生命改變的第一步.
都是明白我們的生命需要神.
我不知道你的生命.
會不會都曾經經歷過孤單.
或者經歷過無助.
經歷過被標籤.
甚至經歷過被拒絕.
可能甚至有時候.
是標籤和拒絕的不是其他人.
而是你自己.
有時候我們去到一個位置.
我自己都標籤自己.
我是沒有希望的.
我自己都拒絕了自己.
我跟自己說你是沒有希望的.
我曾經認識一個弟兄.
他回來教會一段時間.
然後有機會和他聊天.
我問他是怎樣回來教會的.
他說原來是一對夫婦帶他回來教會.
我說你是怎樣認識這對夫婦的.
原來是在茶餐廳認識這對夫婦.
他說有一次.
他在茶餐廳聽到有一對夫婦突然跟他聊天.
突然問他.
你好嗎?.
他說我不好.
我現在的生命很混亂.
他說我的感情一團糟.
我的工作快要爆煲.
我不知道我的生命有什麼方向.
在這個時候.
那對夫婦說不是的.
你有希望的.
然後他自己也說.
我看不到我有希望.
我看不到我有出路.

$^{321}$然後這對夫婦說不是的.
在神裡面你有希望的.
我就約你星期天回教會.
他真的在茶餐廳見到一對夫婦聊天.
邀請他回來教會.
但是有時候就是這樣.
但是當我們的生命.
去知道我們需要神.
原來就是我們生命改變的第一步.
那你可能說不是的.
我真的不像這個麻瘋病人.
但是我相信我們的生命裡面.
都曾經經歷過感覺無助.
或者經歷過有一些東西.
不在我們的控制裡面.
或者有一些東西.
我們不知道出路在哪裡.
可能我們都真的經歷過工作的張力.
甚至是工作爆煲.
我們又經歷過人際關係的張力.
帶來的情緒.
又或者我們經歷過.
家裡面有家人的需要.
甚至可能我們自己都經歷過.
身體健康的問題.
我們需要主.
但可能你說不是的.
梁牧師剛才你說的.
我全部都沒有的.
我很幸福.
我全家都很幸福.
但是我們的生命都需要主.
因為如果我們看回剛才那個麻瘋病人.
他其中一個的絕望.
是他不能夠敬拜神.
今天如果我們的生命沒有神.
我們都找不到生命真正的意義.
我們會不斷追尋這個世界.
告訴我們聽聽我們生命滿足的東西.
但是那些滿足不能夠帶給我們真正的滿足.

$^{361}$真正的滿足.
唯有我們回到神裡面.
我們就找到生命真正的意義和滿足.
所以很重要.
我們每一個都要承認我們的生命需要主.
當我們承認我們的生命需要神.
有很多東西不在我們控制.
很多東西不是我們自己搞得定的時候.
這個是我們生命改變的第一步.
而我們今天學的第二個功課.
是未見到結果已先.
我們仍然憑信心去信服主.
我們一起讀十四至十五節好嗎?.
我們一起讀.
一二三.
耶穌看見.
就對他們說.
你們去把身體給祭司查看.
他們去的時候就潔淨了.
來終有一個見自己已經好了.
就回來大聲歸榮耀與誠.
如果我們看這段聖經.
我很想我們去比較一下.
為什麼呢?.
因為這段不是在《路加福音》第一段.
有麻風病人得潔淨的記載.
原來在《路加福音》第五章.
有另外一段麻風病人得潔淨的經歷.
而我們去比較一下的時候.
我們就能夠明白這段聖經所教導我們的功課.
所以我們一起回到《路加福音》五章那段聖經好嗎?.
《路加福音》五章十二至十六節.
我們一起讀.
然後一起比較一下好嗎?.
我們一起讀《路加福音》五章十二至十六節.
我們一起讀.
一,二,三.
有一回耶穌在一個城裡.
有人滿身長了大麻風.
看見他就俯伏在地.

$^{401}$求他說:主日肯.
必能叫我潔淨了.
耶穌伸手摸他說:我肯.
你潔淨了罷.
大麻風立即就來了他的身.
耶穌祝福他.
你切不可告訴人只要去把身體給祭司擦汗.
又要為你得了潔淨.
照摩西所吩咐的獻上禮物.
對眾人作證據.
但耶穌的名聲越發傳揚出去.
有極多的人聚集來聽道.
也指望醫治他們的病.
耶穌卻去到曠野去禱告.
我們比較一下這兩段經文有什麼相似和分別的地方.
相似的地方是兩段經文裡面.
耶穌都依從摩西的律例.
去吩咐這些麻風病人要把身體給祭司擦汗.
但是跟之前不同的是什麼呢?.
我們會發現在之前這一段聖經.
耶穌有伸手摸病人.
而我們今天讀這段聖經.
耶穌沒有伸手去摸病人.
而我們去讀之前那一段聖經.
耶穌一摸病人的時候.
大麻風立刻就來了他的身.
但是我們今天看這段聖經.
大麻風沒有立刻離開這個人的身體.
這段聖經說他們去的時候就潔淨了.
就是說這段聖經跟之前那一段不同.
這段聖經那些麻風病人沒有立刻康復.
沒有立刻康復.
而他們是要帶著信心去見祭司的時候才康復的.
但是你試一下去幻想一下.
那個麻風病人的心情.
耶穌我求你來醫我.
然後你跟我說你現在去吧.
你去叫祭司看看你的身體.
然後我看一看我的身體.
不是啊.

$^{441}$那些大麻風還在.
你現在叫我去見祭司.
去到見祭司會怎樣?.
祭司會說你有沒有看過律例啊?.
你有沒有衣服撕裂啊?.
你有沒有頭髮蓬鬆啊?.
你有沒有到處告訴別人你是不潔淨啊?.
哇你這麼不潔淨還跑來祭司前面.
我現在就趕你走.
然後那些人就笑你.
還趕你走.
你想一下如果我還沒有康復.
你就叫我去見祭司.
我是不是真的會康復的?.
如果我康復不了.
我就被人笑.
還有被人趕走.
你能不能夠明白那個麻風病人的心情?.
但是在這裡他們沒有辯論.
耶穌說.
你們去吧.
你們去就將身體給祭司擦汗.
所以其實這一段聖經的關鍵句是.
他們去的時候就潔淨了.
這一次這個潔淨不是一個原地的潔淨.
不是那一刻就康復了.
那現在我可以去見祭司了.
耶穌叫他們見祭司的時候.
他們還是渾身麻風病的.
但是在他們信服耶穌的路上.
他們就經歷了醫治了.
所以原來信心不是感覺.
信心是願意按著主的話去行動.
所以我們要學一個功課.
就是當我們還沒有看到環境改變的時候.
我們是不是繼續用信心去跟隨主?.
之前那一段的聖經.
其實是想表明耶穌是接納那些麻風病人的.
耶穌是接納那些不潔的人.
所以用手去摸他們.

$^{481}$但是今天這一段的聖經是告訴我們.
這些人是經歷了一個信心醫治的大能.
就好像舊約一樣.
舊約乃萬得醫治是不是?.
先知說你去沖七次涼.
就得醫治了.
那是什麼意思?.
我還沒有醫治的.
我要跟著去沖七次涼才醫治的.
而今次這些病人的離去.
就是帶著信心去接受神的工作.
確實是神的工作.
因為第十五節說已經好了.
已經好了其實是一個被動的動詞.
就是說得到了醫治.
是別人醫治你.
不是自己得到醫治.
所以在這裡是很表明這個醫治是神的工作.
而有一個人他經歷了神的工作.
他又知道是神的工作.
他就回來歸榮耀彼神.
其實歸榮耀彼神不單止去宣認這是神的工作.
更加宣認這是神的身份.
所以今天我們要問自己.
當我們走在信心的路的時候.
當我們未見到有改變.
未見到結果的時候.
我們是不是繼續帶著信心去跟隨神呢?.
我在網上看過一個故事.
是說在十九世紀.
在十九世紀有一個叫George Muller的人.
他就在英國的布里斯托(Bristol)開孤兒院.
我剛剛其實回英國放假.
我也去過Bristol這個地方.
去探望一位很親愛的神學院老師.
不過我沒有機會去看過這些孤兒院的遺址.
這個George Muller在布里斯托開孤兒院.
他由幾十個兒童開始.
然後他又會再建多一間.
又再建多一間.

$^{521}$最後他建了五間孤兒院.
每一間孤兒院有幾百個兒童.
但是在過程裡面.
其實他也經歷過缺乏的時候.
經歷過困難的時候.
也會有人跟他說.
為什麼你還要堅持下去呢?.
但是他堅持下去.
因為這個是神給他的感動.
曾經有一個記載.
有一天他其中一間孤兒院裡面的職員跟他說.
我們今天已經沒有錢買早餐給那些兒童吃了.
於是George Muller就去祈禱.
祈禱之後突然間有人敲門.
那個人原來就在陣上面.
是焗麵包的.
他說我昨晚祈禱的時候.
神感動我告訴我們.
那些兒童需要麵包.
於是我就去焗麵包.
今天就送到來.
兒童是不是需要麵包啊?.
於是George Muller很感恩接收了這些麵包.
然後過了不久有人敲門.
什麼人敲門呢?.
原來是送牛奶的敲門.
他說我們運送牛奶的車.
剛剛在你門口壞了.
現在那些牛奶送不到去原地了.
擺在這裡是浪費的.
你可不可以要了它?.
就是這樣那天就有早餐給幾百個兒童吃了.
今天我們有沒有信心.
在我們還沒看到改變之前.
去跟隨著呢?.
我自己也有這樣的經歷.
在2007年的時候.
當時我自己的工作出現很大的張力.
我已經很想讓自己找一個出路.
所以當時我為自己想了一個出路.

$^{561}$就是不如我停止我的工作.
接著去讀神學.
我經常想就是.
如果我突然間不工作.
如果人們問我家人.
你兒子為什麼不工作?.
是很難說的.
但是如果我說我自己讀神學.
接著人們問我家人.
喂!為什麼你兒子突然不工作?.
接著家人就說.
Oh, he went crazy over Jesus.
這不就行了嗎?.
於是我覺得這個挺好的.
於是我就覺得.
就去讀神學吧.
我覺得這個計劃十分之好.
於是我就去找陳牧師.
陳牧師我想到一個很好的計劃.
我現在覺得我的工作有很多張力.
所以我想離開我現在的工作.
但是我又不知道有什麼好的計劃.
去離開我這個工作.
我一下子就去讀神學.
於是我讀神學的時候.
如果有些人問我家人.
為什麼你兒子現在不工作?.
Oh, he went crazy over Jesus.
這不就行了嗎?.
我想說我只是想去讀神學而已.
我不是想著去事奉的.
我想找一些東西去過冷河而已.
但是我想著.
應該牧師聽到有人想去讀神學.
應該很高興的.
於是我將這個計劃告訴陳牧師.
陳牧師就這樣跟我說.
他說這個實在是一個人想很好的計劃.
但是神怎麼想?.
於是我就回應他.

$^{601}$神怎麼想?.
我沒有想過.
我說這樣不行.
你將你的生命交給神.
你不是去尋求神的心意.
你現在就去尋求神的心意.
於是我就問.
那我怎麼去尋求神的心意?.
於是陳牧師就跟我說.
從今天開始你每個星期回祈禱會吧.
我說就是這樣?.
就是這樣.
那我有沒有出路?.
我說就是每個星期回祈禱會.
那時候我帶著這個信心.
陳牧師說你每個星期回祈禱會.
你就去尋求神的心意.
我一回祈禱會.
回了兩年.
這一兩年神都還沒告訴我.
我下一步是怎樣.
這一兩年其實工作有沒有張力呢?.
有的.
工作有沒有衝突呢?.
有的.
但是這一兩年.
我就是帶著信心去求問神.
很奇妙.
兩年之後.
就是在一個祈禱會裡面.
Pastor Pam為我的一個祈禱.
神讓我有這個全時間事奉的呼召.
是跟之前那個不同的.
之前我想讀神學.
是想著找一些事情做一下.
去過冷河嘛.
我那時候沒有想過全時間服侍的.
但是兩年之後.
當我真的將生命交給神.
神感動我的.

$^{641}$是一個全時間的服侍.
當然你問我.
在那一刻我跟神說.
沒有收入怎麼辦?.
但是神繼續去引領我.
我想說什麼呢?.
今天你看到.
還沒改變的時候.
你有沒有信心繼續跟隨主?.
今天好像那些麻瘋病人.
他們是憑著信心去的時候.
去經歷醫治的.
不是那一刻就得到醫治.
今天在我們的狀況裡面.
今天在我們生命的環境裡面.
可能你也覺得.
主啊我的出路是什麼?.
今天在人的環境裡面.
你也說好像還沒看到有改變.
但是今天.
我們是不是仍然憑著信心.
去信服跟隨主呢?.
當我們說我們憑著信心去信服跟隨主的時候.
就到了今天我們分享的第三個功課.
今天第三個功課.
是我們的生命回應.
並且真正進入救恩的關係.
我們從第十六節到十九節.
我們一起讀.
一二三.
這人是撒瑪利亞人.
耶穌說:結證了的不是十個人嗎?.
那九個在哪裡呢?.
除了這外族人.
再沒有別人回來歸榮耀與神嗎?.
就對那人說:.
喜害!.
走吧!.
你的信救了你.
如果你記得我一開始說的兩個故事.

$^{681}$我收到那個玩具飛機很開心.
但是我就跟著飛機飛走了.
或者收到那些跳飛機聰明的錄音帶.
昨天真的有一個弟兄.
還找回那些兒童節目的錄音帶給我看.
那些真的是珍藏.
大家記得我拿著跳飛機的錄音帶.
然後很開心地跳來跳去.
在那裡我媽媽給我的教導是.
爸爸是期待與你有一個真正的關係.
而今天我們學第三個功課就是一樣.
今天究竟我們是不是經歷神的幫助.
接著就完了.
但是其實今天我們有沒有真正.
與神進入一個救恩的關係.
今天我們是不是有一個真正跟隨耶穌.
一個親密的關係.
所以今天我們要學習一個很重要的功課.
我們的生命有沒有感恩呢?.
我們的生命有沒有歸榮耀給神呢?.
感恩不只是用口說.
說主啊多謝你啊.
很開心啊.
然後走了出去.
感恩是你承認祂真的是神.
如果你這樣去感恩.
你就會與神進入一個真正的關係.
祂是你的神.
你的生命被祂帶領.
同樣歸榮耀給神.
不是說完一切榮耀頌讚歸給我們的主.
歸榮耀給神.
當然我們很多時候說見證.
很多時候說感恩.
我們將榮耀頌讚歸給神.
但是真正將榮耀歸給神.
是我們的生命活出一個榮耀基督的生命.
當我們.
怎麼會活出一個榮耀基督的生命呢?.
當我們的生命跟隨耶穌.

$^{721}$與祂有一個真正跟隨耶穌關係的時候.
我們就活出一個榮耀基督的生命.
所以今天第三個功課很重要.
我們的生命要懂得回應.
並且進入一個真正救恩的關係.
如果我們看一段聖經.
這個人是撒伯尼亞人.
所以在這裡.
他是回來敬拜耶穌.
俯伏在耶穌腳前就是跪拜耶穌.
是在說這個人回到這裡去敬拜耶穌.
以耶穌為神.
但是吊詭的地方是什麼呢?.
這個人是一個撒瑪利亞人.
撒瑪利亞人是當時猶太人鄙視和欺絕的群體.
但是最吊詭的是什麼呢?.
就是十個人裡面只有這個撒瑪利亞人.
能夠真正認出耶穌的身份.
那些自稱是上帝子民的人.
認不出耶穌真正的身份.
所以耶穌在第十七節說.
他說:潔淨了的不是十個人嗎?.
那九個在哪裡呢?.
這個問題可能是耶穌問那個撒瑪利亞人.
但是這個問題更加是耶穌今天在問我們每一個.
自稱跟隨耶穌的人.
他同時要去那些不懂得回來感恩的人.
去反思他們的生命.
所以在第十八節.
耶穌說:除了這個外族人.
再沒有別人回來.
歸榮耀與神嗎?.
外族人在新約.
在這裡出現過一次而已.
就是用來描述那些外邦人.
但是這個字.
在聖殿原來是一個警報牌出現的.
聖經學者說原來以前的聖殿有一個警報牌.
那個警報牌是寫什麼的呢?.
就是寫著外族人.

$^{761}$那個警報牌是這麼寫的.
外族人不得進入聖殿的內院.
否則必被處死.
原來當時在聖殿裡面.
外族人是不可以進入聖殿內院.
去敬拜神的.
這是不是很大的諷刺呢?.
那些猶太人覺得不配進入聖殿內院敬拜神的人.
才是真正在敬拜上帝的人.
那些覺得自己進入聖殿內院敬拜神的人.
根本就認不出耶穌是神.
那些他覺得不配進入聖殿裡面敬拜的外族人.
反而正在敬拜真神.
所以今天我們要反思.
今天我們要反思什麼呢?.
我們是不是只有敬拜的外在.
但是我們的生命是沒有敬拜的內在.
今天我們每一個人都有敬拜的外在.
為什麼呢?.
你現在坐在教會裡面.
你進來教會.
你就有敬拜的外在了.
但是今天.
我們是不是只有敬拜的外在.
而沒有敬拜的內在.
接著你問.
什麼是敬拜的內在呢?.
我們看第十九節.
"就對那人說,起來走吧,你的信救了你".
這段聖經我們常常誤解了.
我們誤解了什麼呢?.
我們覺得.
他跟這個撒瑪利亞人說的話.
"起來走吧,你的信救了你".
這個信是什麼信呢?.
這個信就是相信耶穌能夠醫治我的信.
所以他相信耶穌能夠醫治我.
我的身體就能夠醫治了.
但是如果這個是撒瑪利亞人的信.
另外九個也有這個信的.

$^{801}$另外九個也是相信耶穌能夠醫治他的身體.
另外九個的身體也能夠醫治.
所以在這裡耶穌說的這個信.
不是信耶穌能夠幫助他的身體.
能夠醫治的那種信.
撒瑪利亞人的那種信.
是相信耶穌是主.
耶穌是神的信.
而這個信是他願意用生命去敬拜跟隨耶穌的信.
而這個信他的得救.
不只是身體得醫治.
是他得著真正的救恩.
所以原來這個撒瑪利亞人的信.
不只是信耶穌醫治我的身體.
他是信耶穌是神.
並且我一生敬拜跟隨.
我們也要問今天我們的信是什麼?.
今天究竟我們的信是信耶穌能夠幫助我.
所以今天耶穌能夠幫助我.
我就求耶穌了.
還是我們的信是信耶穌是神.
我們一生敬拜跟隨耶穌.
我在網上找到一個文章.
這個文章裡面說教會裡面有兩類人.
我們試試這個迎新十二課的形式.
因為待會會請你舉手.
迎新十二課也有A和B.
請你舉手.
他說教會裡面會有兩類人.
A類是什麼呢?.
經常經歷神的幫助.
考試過關.
感情挽回.
身體痊癒.
但是生命的主軸沒有改變.
仍然自己作王.
平時對神懇懶淡.
旨在出事時才來呼求.
A類.
我們看看B類.

$^{841}$B類也會求神幫助.
但是每次經歷恩典之後.
都會再回到主面前.
更多敬拜.
更多信服.
更多想認識他.
慢慢地這個人不只是得幫助.
而是被主得著.
我們試試迎新十二課.
覺得自己是A類的,你舉手.
多謝你的坦誠.
看到多謝你的坦誠.
我會找你聊天.
因為兩個我都認識.
覺得自己是B類的.
沒有舉手的是什麼類?.
是C類.
C類即是今生的C類.
我多謝剛才很誠實舉手的兩位.
我認識你的,我會找你聊天.
其實你沒有舉手.
我都想找你聊天的.
不過既然你舉手了.
我們就聊聊天.
再聊久一點.
未必一定是這樣的.
喝杯咖啡.
但很重要我們都要問自己.
我們是哪一類人?.
今天究竟我們只是回來.
希望得幫助.
抑或我們的生命.
真是想被主得著.
今天我們的生命.
是被依次那九個人.
生命得幫助,得供應.
但沒有真正進入救恩的關係.
抑或我們是懂得回來感恩的那個.
那個感恩不只是一個外在的感恩.
那種感恩是承認耶穌是神.

$^{881}$所以當我們承認耶穌是神.
我們就會將生命,時間和敬拜.
去獻給我們的主.
所以我們要問自己.
我們的生命是不是都經歷過.
神的幫助,神的依次.
神為我們開路.
但是我們是那九個人.
抑或我們是會懂得回來.
然後將我們生命的時間.
敬拜,事奉獻給主.
並且將神放在我們生命的中心.
所以這段聖經教導我們.
很重要我們要懂得感恩.
懂得感恩其實真的不只是口說.
懂得感恩是.
原來神在我們生命裡面的工作.
我們要將耶穌放在我們生命的首位.
以祂為神.
也都當我們說我們歸榮耀給神的時候.
不只是用口歸榮耀.
或者敬拜的時候.
很重要我們用口歸榮耀給神.
我們敬拜將榮耀歸給神.
但是我們的生命就有沒有活出一個.
榮耀基督的生命.
是唯有我們的生命去活出一個.
真正榮耀基督的生命.
我們才真正歸榮耀給神.
所以今天在總結的時候.
我們請敬拜隊.
這段聖經給我們三個生命.
很重要反思的問題.
第一個反思我們要問自己.
我們願不願意去承認.
自己的生命需要主.
今天我想我們每一個人都需要主.
就算今天你託梁牧師.
我很幸福 我一家都很幸福.
我們都需要主.

$^{921}$因為唯有回到主的裡面.
我們才找到生命真正的意義.
所以今天我們每一個人.
都要承認我們的生命需要主.
而這個是我們生命改變的第一步.
而當我們去承認我們的生命需要主.
我們要問第二個問題.
今天有時可能我們.
都未見到生命的狀況有改變.
我曾經聽過一個見證.
他說他在教會信了主.
然後信了主之後.
生意才來不好.
於是他就問.
但是他信了主了.
你這個神是不是真的?.
我信了你之後生意還好嗎?.
是不是的?.
是的.
但是我們要問自己.
今天我們未見到生命改變的時候.
我們是不是繼續願意憑信心.
去跟隨耶穌?.
但是更加重要的是第三點.
我相信在座我們都經歷過神的恩典.
經歷過神的幫助.
神的帶領.
我們有沒有回來歸榮耀給神.
就是說主啊.
我的生命要交給你.
由你帶領.
今天很重要.
我們要跟主進入一個真正救恩的關係.
今天你可能信了主很多年.
今天你也可能經歷過神.
在你生命裡面的幫助和帶領.
但是今天我們要用不同的角度.
去看這段話.
「起來,走吧,你的信救了你」.
這個信不是單單覺得.

$^{961}$耶穌能夠幫助我們的信.
當然這個信很重要.
當我們有需要的時候.
我們知道耶穌是我們的幫助.
但是這個撒瑪利亞人的信.
不只是相信耶穌能夠醫治他的身體.
是耶穌能夠讓他的全人得醫治.
是他的情感.
他的情緒.
他的關係.
他的家庭.
他的全人.
但是我們怎麼去經歷這個整全的救恩的生命.
唯有我們相信耶穌是神.
這個信不只是相信耶穌有能力幫助我們.
這個信是相信耶穌是主.
所以今天我們有沒有再一次.
願意將生命這樣交給耶穌.
神高揚佩特.
很想再一次藉著今天的回應詩.
我們一起向耶穌做一個生命的回應.
我們願意一起說.
主我願意我的信.
是承認你是主的信.
是願意承認.
我要一生跟隨你的信.
並且我就將一生交給你.
\newpage



\section{箴言 14:26-35}
\label{sec:031737}
\textbf{生命泉源}~~~~梁永善牧師
\newline
\newline
連結: \href{http://www.church.com.hk/acms/content.asp?site=cdc\&op=show\&type=product\&code=031737\&layout=sermon}{\texttt{031737}} ~~~~ 語音日期: 2026-03-01
\newline
\newline
\hyperref[sec:031736]{\small{< < < PREV SERMON < < <}}
~
\hyperref[sec:index_chronic]{\small{[返順時目]}}
~
\hyperref[sec:031738]{\small{> > > NEXT SERMON > > >}}
\newline
\newline
箴言 14:26-35
\newline
\begin{longtable}{cl}
\hline
\hline
章節 & 經文 (和合本修訂版)\\
\hline
14:26 & \begin{tabularx}{0.7\textwidth}{X} 敬畏耶和華的，大有倚靠；他的兒女也有避難所。 \end{tabularx} \\ \\ \relax
14:27 & \begin{tabularx}{0.7\textwidth}{X} 敬畏耶和華是生命的泉源，使人離開死亡的圈套。 \end{tabularx} \\ \\ \relax
14:28 & \begin{tabularx}{0.7\textwidth}{X} 君王的榮耀在乎民多；沒有百姓，王就衰敗。 \end{tabularx} \\ \\ \relax
14:29 & \begin{tabularx}{0.7\textwidth}{X} 不輕易發怒的，大有聰明；性情暴躁的，大顯愚昧。 \end{tabularx} \\ \\ \relax
14:30 & \begin{tabularx}{0.7\textwidth}{X} 平靜的心使肉體有生氣；嫉妒使骨頭朽爛。 \end{tabularx} \\ \\ \relax
14:31 & \begin{tabularx}{0.7\textwidth}{X} 欺壓貧寒人的，是蔑視造他的主；憐憫貧窮人的，是尊敬主。 \end{tabularx} \\ \\ \relax
14:32 & \begin{tabularx}{0.7\textwidth}{X} 惡人因所行的惡必被推倒；義人臨死，有所投靠。 \end{tabularx} \\ \\ \relax
14:33 & \begin{tabularx}{0.7\textwidth}{X} 智慧安居在聰明人的心中，在愚昧人的心中卻不認識。 \end{tabularx} \\ \\ \relax
14:34 & \begin{tabularx}{0.7\textwidth}{X} 公義使邦國高舉；罪惡是百姓的羞辱。 \end{tabularx} \\ \\ \relax
14:35 & \begin{tabularx}{0.7\textwidth}{X} 君王的恩寵臨到智慧的臣僕；但其憤怒臨到蒙羞的臣僕。 \end{tabularx} \\ \\
[1ex]
\hline
\hline
\end{longtable}
$^{1}$各位電影節目的同仁們,很高興在新的一年裡面第一次與大家聊天..
如果熟悉我的人都知道,我每年都會按照生肖講一篇道理..
今年是馬年,我應該講馬年..
但是你們教會很特別,我也很欣賞,每次都是指定經文,指定講題..
但是我又很乖,我一定會依著你們的指示去做..
所以如果想聽我講馬年話馬的話,可以在網上聽..
這篇新的講道,我很欣賞你們..
因為今天是3月1號,一個新的主題,是敬畏榮耀尊貴滿有能力的主..
你很簡單看一看,當你看新聞的時候,.
當國家領導人開會,人民大會堂裡面的人,個個都很認真去聽..
甚至在那裡做筆記,就不會像我們那些議員聽財政預算就去打盹..
如果我們都覺得一個人我們都要這麼尊重他的時候,.
何況我們說我們的神是滿有能力創造天地萬物的主宰,我們對他有沒有這份尊重呢?.
所以教會要求我說,這個生命泉源就是來自真言..
如果我們看真言的時候,要求我說的就是第十四章的二十六到三十五節..
一開始你有沒有留意到二十六二十七節的時候,都是講到要去敬畏神的..
敬畏耶和華的,大有倚靠他的兒女也有避難所..
敬畏耶和華的,就是他生命的泉源,可以使人離開死亡的網羅..
為什麼敬畏神就可以使人離開了這個網羅呢?.
將所有的已經給了他,其實原來他私下留下了一點點..
使徒就問他,是不是真的給了他?.
是啊,一講完這個謊話就跌在地上死了..
剛剛抬了屍體出去之後,太太就來了..
為什麼會這樣呢?.
你會明白的,結了婚的人就會知道了,拍拖就是一起走的..
但是結了婚之後,丈夫就先走的..
為什麼叫丈夫?就是一丈之夫,就先走..
太太來了,又照樣講這個謊話,結果他馬上就死了..
接著就是這節經文..
全教會聽見的人都很害怕,原來神是這麼威嚴的..
神是聖潔,不容讓小許的罪在當中..
所以他們就很害怕..
當我們離大海去敬畏神的時候,會產生什麼效果呢?.
有很多方面的效果,今天我們會去看..
第一,就是叫我們的生活上很嚴謹..
就是剛才曾賢所講,會叫我們脫離死亡的網羅..
如果我們離大海去輕忽了神,沒什麼所謂的..
我們自然就會輕忽我們自己的罪..
如果我們不懼怕神的審判,我們就會生活上鬆懈..
沒什麼所謂的,就犯錯了..

$^{41}$所以一個沒有信仰的人,他可以胡作非為的..
同樣有信仰的人,如果他對神沒有敬畏的心,.
他一樣可以離大海去做不應該做的事..
所以當我們肯離大海,倚靠神,對神有一個敬畏的態度,.
我們就會完全不同..
今天,經威耶和華,就是我們生命的泉源,.
可以使我們離開死亡的網羅..
這個就是箴言第十四章二十七節,.
就是我們今天要看的其中一節的經文,是不是?.
當然有時候我們牧者都需要反省的,.
因為我們常常都很強調,我們神是沒有恩典,.
慈愛的神,祂很願意寬恕我們..
即使我們如果犯了錯,原來我們都未祈求之先,.
神已經為我們預備了救恩的了..
只要我們一認罪,神就會赦免了..
我們很多時候都強調這些,我說這方面沒有錯的,.
但是我們就以騙概全,.
我們輕忽了神的榮耀,尊貴能力和聖潔,.
以致我們沒有這份敬畏的心..
所以當我們知道神是一個滿有權能,威嚴的神的時候,.
我們在生活上,我們絕對會嚴謹,.
不會隨便離開..
我們知道原來我們的神是最終的審判者,.
是真真正正的終審庭,.
祂的判斷是最正確的,.
你已經不能再上訴了..
這時候我們就不會輕率,我們不會只是看現在眼前的成敗得失,.
然後我們放棄了我們應有的原則,.
我們生活上就會很謹慎..
很簡單舉個例子,.
以前,我還在香港的時候記得,.
如果你隨地吐痰,會罰千五元的,.
我不是很明白有些人吐口水是有些陰影效果的,.
當你剛剛想吐的時候,.
見到有個食環署職員在那裡,.
你還敢不敢吐?.
馬上吞下去,.
因為會罰千五元的..
以前,我們說以前,.
有一個國家,你不是見到很多警察,.

$^{81}$不是很多閉路電視的,.
但你不敢隨便犯錯,.
我也去過這個國家,.
在哪個地方?.
新加坡..
新加坡沒有很多警察,.
但有很多便衣..
但你知道新加坡,.
很多地方都要罰錢,.
我較為節儉,.
但那次我在新加坡也買了一件衣服,.
這件衣服是很特別的,.
它說什麼呢?.
Singapore is a fine city,.
你有沒有留意到fine,.
這個字有兩個解釋,.
就是一個很好的城市,.
前面就是Singapore is a fine city,.
但後面有九個格,.
這九個格是說什麼呢?.
就是你犯了什麼錯,.
不依著交通規則過馬路,.
你會被罰,.
你來到隨地吐痰,可以被罰..
不知大家記不記得很多年前,.
有一個英國的年青人,.
在新加坡那裡,.
拿一個銀幣,.
劃花人的車,.
劃了一下,.
被人抓了,知不知道有什麼事?.
沒錯,要打藤,.
打七藤,.
但這一藤不是你媽媽打你的那些,.
是要檢查了你體格,.
然後才去打的,.
當年英國的外交部,.
都要為這個年青人求情,.
因為很嚴厲,.
但是我們有沒有想到,.

$^{121}$神都是很嚴厲的,.
有一天我們在神的審判台前的時候,.
我們要交帳的,.
你不要想著我所做的事,.
沒有什麼人知道,.
你想想現在所有人面辨識,.
香港要多放四億,.
雖然在淘寶買那些CCTV,.
是很便宜的,.
但是特區政府就買很貴的了,.
就是現在香港很多這樣的鏡頭,.
你做什麼事,隨時會被人看著,.
但是你不要忘記聖經說道,.
我們的神是無所不知的,.
你暗地裡要做的事,.
你以為沒有人知道,.
看那些不好的事,沒有人知道,.
但是你自己知道,.
但是不要忘記魔鬼和神都會知道,.
你可以欺騙到人,.
但是你欺騙不到神,.
有一天,.
你要在這個審判台前,.
你帶他去交帳的時候,.
你怎麼可以.
放縱你的邪情私慾呢?.
所以今天我們要知道,.
神不是單單慈愛,.
我們一起來再讀.
詩篇九十篇的第十一節,.
我們一起來讀出..
詩篇九十篇是.
摩西所寫的,.
他看到第一代出埃及的人,.
除了兩個之外,.
全部都死在曠野裡,.
就是他們對神沒有敬畏,.
這句話是怎樣的呢?.
神是讓我們知道.
你是會發怒的,.

$^{161}$你是為什麼事發怒,.
我們不可以得罪你,.
按著你該受的敬畏,.
曉得你的憤怒,.
神給我們很多機會,.
我們就要面對這個審判,.
不要想著我信了耶穌,.
就得到這個救恩,.
就可以隨便犯錯,.
很寶貴的一段經文,.
希伯來書第十章二十六,.
二十七和三十一節,我們一起來讀出,.
預備,一二三..
如果我們得知真道,.
但是故意犯罪,.
這個救贖,.
未必能夠臨到你身上,.
只有等候這個審判,.
所以說,.
落在永生神的手裡,.
是很可怕的一件事,.
但記住希伯來書,.
不是舊約,是新約,.
不要再誤會,.
舊約的神就惡些,.
新約的神就好些,.
是不變的,舊約的神一樣充滿慈愛,.
新約的神,.
同樣的充滿慈愛,.
但同樣是公義的,.
所以我們怎麼可以輕信呢?.
所以聖經說,如果我們忽略這麼大的救恩,.
怎麼能夠逃罪呢?.
所以我們就需要懷著恐懼,.
戰驚,.
作成我們得救的功夫,.
意思是什麼?.
我們要好好的過我們的生活,.
嚴謹地過我們的生活,.
敬畏神,.

$^{201}$知道神的審判是嚴厲,.
於是我們不會隨便,.
我們不會放縱,.
於是就能夠救我們脫離死亡的網絡,.
所以求神幫助我們,.
記住,.
神是一位威嚴,.
公義的神,.
甚至很愛我們,.
這份愛意志,.
祂會成為一個忌邪,妒忌的神,.
祂要求我們一生一心一意地跟隨祂,.
祂不容許我們一面敬拜神,.
一面敬拜霸凌,.
一隻手拿著世界,.
一隻手想拉著神,.
不是谷愛玲,.
不可以搭兩條船,.
只可以來到一樣,.
就是你選擇跟隨神,.
還是放棄世界的吸引,.
這裡我們是要去反省的,.
神絕對不容許我們身子是聖靈的殿,.
但也可以成為賊窩,.
絕對不可以,.
所以求神幫助我們,.
原來當我們敬畏神,.
是叫我們脫離死亡的網絡,.
我們的生活嚴謹,.
第二方面就是敬畏神,.
是可以祝福我們的兒女,.
甚至我們的下一代,.
在箴言第十四章二十六節,.
我們剛才讀,.
我們再一次讀,.
預備,一二三,.
敬畏耶和華的大有倚靠,.
他的兒女也有避難所,.
敬畏耶和華,.
他的兒女也有避難所,.

$^{241}$因為你做得好,.
那時候你就可以保護你的兒女,.
但還有一個原因,.
我們一起來讀詩篇,.
一百零三篇第十七節,.
我們一起來讀出,.
原來當我們敬畏神,.
尊重神,愛我們的神的時候,.
神不但祝福我們,.
還會祝福我們的子子孫孫,.
同樣一個人犯錯,.
離開神,.
其實他也會禍及他的兒孫,.
所以聖經很著緊一件事,.
就是家庭教育,.
這裡提醒每一個做父母的,.
我多次跟大家說,.
你不要只是著緊你的兒女的功課,.
可不可以考到一間好的學校,.
這個我們要重要,.
但最重要是你有沒有著緊他們的信仰,.
你有沒有好好的為他們祈禱,.
不只是迫他們回教會,.
而是分享你自己的一些領受,.
一些經歷,.
讓我看到這位神是何等的真實,.
你無法看他們一生一世,.
如果按正常情況,.
因為你都要離開,.
但如果你能夠讓你的兒女認識神的時候,.
神會看著他們,.
你何須擔心,.
所以以色列人很著緊家庭教育,.
你留意到以色列你要聽耶和華,.
我們神是獨一的神,.
你要盡心盡性盡力,.
愛耶華你的神,.
我今日吩咐你的話,.
都要記在心上,.
這是我們自己,.

$^{281}$我們要謹記敬畏神,.
但接著要做什麼,.
就是要殷勤教訓你的兒女,.
無論你坐在家裡,行在路上,躺下起來,.
都要談論,.
要繫在手上為記號,.
戴在額上為經文,.
又要寫在你房屋的門框上,.
並你的城門上,.
意思就是說,.
你有沒有隨時隨地和你的兒女分享,.
吃飯的時候,.
還是各自看手機,.
有沒有說說你自己的一些經歷,.
不只是死顧顧,.
分教會,讀聖經,.
你有沒有將你的領袖,.
和他們分享呢?.
所以他說到,.
你是要經常,.
但我們不是照字面的意思,.
要繫在手上為記號,.
戴在額上為經文,.
這就是在以色列拍攝的,.
他們真的在這些成人節裡面,.
會做這些事的,.
在手上真的戴著一個盒子裡面有經文,.
但聖經真正不是這樣,.
而是教導我們要,.
教導兒女,.
敬畏神,請大家一起來讀.
《新冥紀》第六章第一到第二節,.
我們一起來讀出..
這個敬畏的態度,.
不只是你自己要持守,.
你要教導他們認識神,.
同樣是一個尊貴榮耀的神,.
都教導他們要.
敬畏耶和華,.
就以致到.

$^{321}$他能夠.
尊重神,.
使到他得著祝福,.
很簡單舉個例,.
不只是理論,.
我不知道教會.
是不是會全奉獻袋奉獻的,.
我曾經.
在我自己的冥恩堂也說過,.
教會不是那麼多寡婦的,.
為什麼我們在.
主義學裡面.
收到很多小錢呢?.
你不用那些一毫兩毫,.
竟然.
在主義學裡面.
會找到出來的,.
五毫子,.
有時候可能我們就給點錢,.
去祝福他們,.
這樣的行動.
會令到你的兒女覺得.
神是很廉價的,.
你去喝茶給小朋友也不是這樣,.
這就是我們.
舉一個很簡單的例子,.
去反省,.
你如何教導你的兒女,.
去敬畏你的神,.
尊重你的神,.
我們崇拜的時候那種態度,.
會不會遲到,很輕率的,.
隨便,.
聽你說什麼,.
去看你在做什麼,.
所以為什麼那時候有第二代的.
基督徒,.
他們是不是會放棄信仰,.
其實做父母要反省,.
不要只將責任.

$^{361}$推在教會身上,.
所以我們要好好地教導.
你的兒女,.
真言有一句話是影響到我們,.
誤解了真言,.
那句話大家都很熟悉,.
教養孩童,.
使他當走的路,.
以致他到老都不偏離,.
這句是真言裡面.
我們很熟悉的,.
小朋友要教,.
當然要教,.
但當他長大了你要教嗎,.
當然要,.
真言你留意到,.
他教你怎樣教子女,.
但那些子女不是小時候的,.
很簡單,.
你要愛情專一,.
要為自己的發妻而歡欣,.
你會教他四歲五歲的小朋友,.
兒子,將來你娶老婆,.
要哄你老婆開心,.
不然你自己受苦,.
會不會,.
雖然這是真理,.
但不是那個時候教他,.
所以聖經裡面是很清楚的,.
特別在真言第六章裡面,.
我不讀這些經文,.
大家可以看,.
他教千萬要專一的,.
愛你自己的配偶,.
有婦之夫或者有夫之婦,.
有淫亂的時候,.
你是自尋死路,.
你有和別人的妻子,.
有不正常的關係,.
就算你給錢賠償,.

$^{401}$都賠不到,.
對方可能會很恨你,.
於是你就會有殺身之禍,.
事實上我們見到,.
有這些事實,.
有一個人,.
他的太太出軌,.
他的丈夫其實很好的,.
他很愛自己的妻子和孩子,.
但後來知道自己的太太,.
有第三者,.
他是一個三行工人,.
後來就用斧頭劈死了對方,.
但是很多的鄰居,.
親人都為他求情,.
說這個丈夫很好,.
我沒有說他這個.
殺害對方是正確,.
但你看到聖經就是這麼真實,.
所以聖經裡面提醒我們,.
就是要什麼呢?.
好好地教導我們兒女,.
敬畏神,這是第二方面,.
第三方面呢,.
就是真言提醒我們,.
我們如果能夠敬畏神的時候,.
我們的生活是更加和諧的,.
我們一起來再讀.
真言第十四章,.
今日經文第二十九節,.
我們一起來讀出,.
聰明的人.
不是那麼容易發脾氣的,.
那些一拍就著的人呢,.
其實是很愚忙的,.
我們再來看第十一和第十二節,.
這個也是同一段的經文,.
我們一起來讀出,.
因為王是有權力的,.
如果獅子嘯叫,.

$^{441}$在阿摩斯書裡面有提到這一句,.
The Roar of the Lion,.
獅子一嘯叫的時候,.
其他那些動物都害怕的,.
神來到嘯叫的時候,.
發出他的怒嘯,.
我們個個都應該面對很大的審判,.
所以我們更加要嚴謹,.
特別講到我們是要有那個.
溫柔的心,.
不要那麼隨便發怒,.
那時候呢,如果我隨便發怒的時候,.
就會帶來一些問題,.
我們一起來讀.
真言第十二章,第十八節,.
我們一起來讀出,.
如果我們講話太隨便,.
有時候就會.
令到那些人.
有不開心的反應,.
當年一對父子,.
爸爸叫他去交電費,.
誰知道他搞錯了,.
我經常都搞錯這些事情,.
那時候呢,.
那個爸爸說,死蠢,沒有用的,.
那個兒子聽見很生氣,.
拿刀劈,不要學了,.
拿刀去劈,.
那時候呢,.
那個爸爸說,死蠢,沒有用的,.
拿刀劈,不要學了,.
拿刀去劈他爸爸,.
這些其實是積怨,.
有時候一言可以輕幫,.
一言可以喪幫,.
所以我們要很小心,.
現在你想想,特別香港,.
我們看很多新聞,.
在地鐵,在公共場所,.

$^{481}$經常有些人來打架,.
當然他們內裡有很多的.
不安和煩躁,.
但有時候一句話就會令到對方很生氣,.
那我就回應,.
譬如,當年呢,.
我太太和我女兒聊天,.
她說,你爸爸很壞,.
口甜舌滑,.
我被他騙了,.
那你聽到你應該會有什麼反應?.
特別做丈夫那些,.
我就和我女兒說,.
女兒,我騙了你媽媽,.
不過現在你看到,.
說謊是不會有好後果的,.
你始終是會受苦的,.
我沒有直接去回應我太太,.
但我就和我女兒這樣說,.
那我太太又不知道.
可以再說些什麼,.
當然有時候我們要有智慧,.
不要隨便惹怒兒女的怒氣,.
《聖經》裡面是這樣說的,.
我們一起看,.
做父母那些,.
做兒女那些,未結婚那些,.
你們大聲地讀好不好,.
特別你爸爸媽媽在這裡,.
我們大聲地讀,.
《歌羅西書》第三章第21節,.
《歌羅西書》第21節,.
《歌羅西書》第三章第21節,.
我們一起讀,.
\newpage



\section{詩篇 150:1-6}
\label{sec:031738}
\textbf{全人的敬拜}~~~~陳崇基牧師
\newline
\newline
連結: \href{http://www.church.com.hk/acms/content.asp?site=cdc\&op=show\&type=product\&code=031738\&layout=sermon}{\texttt{031738}} ~~~~ 語音日期: 2026-03-01
\newline
\newline
\hyperref[sec:031737]{\small{< < < PREV SERMON < < <}}
~
\hyperref[sec:index_chronic]{\small{[返順時目]}}
~
\hyperref[sec:031739]{\small{> > > NEXT SERMON > > >}}
\newline
\newline
詩篇 150:1-6
\newline
\begin{longtable}{cl}
\hline
\hline
章節 & 經文 (和合本修訂版)\\
\hline
150:1 & \begin{tabularx}{0.7\textwidth}{X} 哈利路亞！你們要在神的聖所讚美他！在他顯能力的穹蒼讚美他！ \end{tabularx} \\ \\ \relax
150:2 & \begin{tabularx}{0.7\textwidth}{X} 要因他大能的作為讚美他，因他極其偉大讚美他！ \end{tabularx} \\ \\ \relax
150:3 & \begin{tabularx}{0.7\textwidth}{X} 要用角聲讚美他，鼓瑟彈琴讚美他！ \end{tabularx} \\ \\ \relax
150:4 & \begin{tabularx}{0.7\textwidth}{X} 擊鼓跳舞讚美他！用絲弦的樂器和簫的聲音讚美他！ \end{tabularx} \\ \\ \relax
150:5 & \begin{tabularx}{0.7\textwidth}{X} 用大響的鈸讚美他！用高聲的鈸讚美他！ \end{tabularx} \\ \\ \relax
150:6 & \begin{tabularx}{0.7\textwidth}{X} 凡有生命的都要讚美耶和華！哈利路亞！ \end{tabularx} \\ \\
[1ex]
\hline
\hline
\end{longtable}
$^{1}$大家好.
今晚我會飛到土耳其.
我會探索一個新的路線.
大家都聽過五月有弟兄姊妹跟我一起去土耳其.
你猜過去兩天發生的事.
我會飛到土耳其嗎?.
暫時沒有影響.
暫時很難說.
這世界真的發生太多事.
本來我打算恢復去以色列的聖地轉移.
暫時還未可以.
但需要為該地區禱告.
我想告訴大家.
很多時候有戰爭的時候.
我們都會被迫站邊.
看你幫誰.
其實是不需要的.
所有戰爭上帝都不喜歡.
因為你知道受害的人是甚麼人.
是無辜的人.
得益的是甚麼人.
是在上掌權的.
無論你是美伊這邊還是伊朗那邊.
哪邊都好.
我們不需要站邊.
因為我們屬於誰.
我們屬於天國.
對嗎?.
屬於天國.
我們可以談戰況,動態,地緣政治等等.
不過我們不需要站邊.
因為我們只屬於一位的國度.
一位天國的君主.
我們這樣去到風花雪月的時候.
其實已經出賣了我們自己的國度.
因為我們是不會受到戰火摧殘的人.
對嗎?.
我們就是那些袖手旁觀或隔岸觀火的一班人.
所以我們不需要站邊.
亦不需要亂拿聖經的經文來對號入座.

$^{41}$這些是我們傳統的教會經常有的弊病.
那裡有甚麼發生呢?.
中東有甚麼發生呢?.
我們就好像在聖經裡面有一些預早寫好的報紙.
這是非常不應該的.
應該學學怎樣去看聖經.
不要看錯聖經.
不要解錯聖經.
聖經沒有說這些東西.
所以大家不要弄錯這些東西.
但我們需要去作和平之志.
那你說聖經這些怎麼解釋呢?.
那你就要報讀我們的課程.
就是這麼簡單 祈禱和報讀課程.
我們先祈禱.
我們知道每逢有戰火摧殘每一個地方的時候.
你的心是很痛的.
因為每一個人都是按照你自己的現象所做.
特別是一些無辜的人.
婦女,兒童,孩童.
他們是被迫在戰火裡面生活.
主要是為了這些人去禱告.
我們為了這兩天美,意,聯,公,伊朗,伊朗還擊.
引致不只是單單一個國家的人受害.
而是整個地區,整個地帶,整個中東的地帶.
主要我們聽到很多傷亡的報道.
無論在杜拜,在一些Gulf States,在伊朗,在以色列.
在不同的地區.
我們見到有很多被戰火摧殘的建築物.
有很多人的生命受到很大的影響.
甚至是被摧殘.
主要我們求你憐憫.
憐憫這個地方因為裡面有很多很複雜.
政治的因素.
是一些近代政客或者不同的大國.
彼此去到爭拗,角力.
以致有很多平民百姓受到影響.
主要求你給我們一個屬靈的角度.
讓我們能夠成為和平之子.
縱使我們未必能夠身在那個地方.

$^{81}$但我們都為當地的居民去禱告.
主要求你去洗這個戰事.
不會繼續spill over.
不會繼續去escalate.
因為你愛世人.
主要求你讓當地的人能夠很快de-escalate到這場戰火.
以致人的生命不會受到摧殘.
主要求你教導我們在一個這麼混亂.
這麼多爭拗的一個世代裡.
教導我們明白我們真正的效忠.
是上帝是自己的國.
真正的效忠是我們的主耶穌基督.
他是那位和平之子.
並且將人與人之間和平的圍牆拆掉.
藉著他十字架的救恩.
將任何的人.
所有的人 猶太人,愛邦人,男與女.
遺老的,自主的 一切能夠在基督裡成人.
滿有慈愛的神.
並且你在這個亂世當中 你仍然是工作的.
求你幫助我們 讓我們在這個混亂的世代裡.
我們不是退縮.
我們更加願意以福音.
以我們的生命來見證你.
我們這樣祈禱 奉尊名禁.
好 今天我們繼續看詩篇 我們需要完結這個詩篇的系列.
我今天就選了詩篇最後一篇.
你記得詩篇有多少首詩歌嗎.
150首是嗎.
我們先看這首詩歌.
這個就是整個詩集.
最後最後一首的詩歌.
看上去其實非常簡單.
但對我們今天是非常重要的.
第一節這裡這麼說 他們要讚美耶和華.
在上帝的聖所讚美他.
在他顯能力的窮倉讚美他.
要因他大能作為讚美他.
按著他極美的大德讚美他.
要用國聲讚美他.

$^{121}$鼓瑟彈琴讚美他.
擊鼓跳舞讚美他.
詩而樂器銷的聲音讚美他.
大響的筆讚美他.
用高聲的筆讚美他.
凡有氣息的都要讚美耶和華.
你們要讚美耶和華.
在這裡結束的時候 詩篇告訴我們.
整本的詩集裡面有一個這樣的結論.
就好像你看一套戲.
看一個連續劇.
最後最後有一個大結局.
雖然我們知道有很多時候.
華文或者港產的.
大結局是很難的.
但是詩篇.
最後最後的大結局.
就是他總結了.
整個的詩篇.
到底想教導我們做些什麼.
就是什麼叫敬拜.
What is worship.
如果你問崇拜.
敬拜是一回什麼事呢.
這個問題是需要問的.
今天立到來的時候.
今天你去教會.
為什麼每一個禮拜你騰空一些時間.
一些可能.
你生命裡面可以成為優先的時間.
但是為什麼你不去做其他事.
你去敬拜上帝呢.
那你就要問.
什麼叫敬拜.
為什麼需要敬拜.
短短的六節.
作為整個詩集的結束.
就是嘗試去回答這個問題.
什麼叫敬拜.
簡單來說.

$^{161}$其實敬拜是什麼呢.
就是我們作為人去回應上帝.
的恩典和他的偉大.
我們作出一個回應的時候.
這個就是一個敬拜.
明白嗎.
這個就是我們的response.
所以我們首先是被動的.
上帝做了一些事創造我們.
創造這個世界.
一個美好的世界.
去救贖世人.
上帝首先做事.
然後我們作為一個被動的主體.
我們去回應他.
既然上帝是一個.
無限的存有.
infinite的being.
所以他配受的敬拜.
也是不能夠停止的.
這個詩篇一直告訴我們.
就是一個這麼重要的課題.
我們來到.
敬拜的時候.
和神的子民一起敬拜的時候.
我們就去回應上帝.
給我們的恩典.
特別是詩集.
聖經裡面的詩篇.
有一百五十首.
就是神的子民.
以色列人.
他們被流放.
然後他們回歸的時候.
他們去整輯出來.
在世世代代有的詩歌.
他們慢慢收集起來.
然後成為一百五十篇.
如果你留意詩集.
是有一個佈局的.

$^{201}$或者一個編輯的.
無論那個編輯是一個人.
或者一群人也好.
總之就是被流亡之後.
這些以色列人他們回歸.
然後他們就去將.
這個所謂正典的舊約聖經.
去到編輯出來的時候.
詩集就是其中一本書.
你看到這裡有很多像LEGO那樣.
其實很簡單.
整本的詩篇.
一百五十首.
是有五本書的.
五卷書.
有Book 1 to Book 5.
每一個.
一而就是一個門口.
告訴我們我們敬拜上帝的時候.
有一個Gateway.
有一個入口.
就好像你進去一個豪宅的時候.
開門打開.
然後你脫下鞋子.
然後你進去那間屋.
那個最堂皇的地方.
那個就是一和第二篇.
然後就分了.
五本書.
每本書最後的時候.
你見到Book 1最後的時候有一個D.
就是讚美.
所謂我們.
技術上我們叫它Doxology.
就是一個榮耀頌.
最後這個詩集.
Book 1.
收集了這些詩歌.
三至四十一篇.
然後最後就會有一個.

$^{241}$一節或者兩節.
就有一個所謂Doxology.
一個榮耀頌.
他就見到四十一篇.
十四節耶和華以色列的上帝.
是應當稱頌的.
從耿古直到永遠.
阿們.
Book 1就完了.
然後Book 2, Book 3, Book 4, Book 5.
三四五.
二三四五卷.
最後都會有一個.
類似的所謂D.
Doxology完結.
你見到譬如.
Book 3 第三卷.
其實就是.
七十三至八十九.
八十九結束的時候就有一個Doxology.
一百五十篇在哪裡.
一百五十篇就算是.
最後最後一個.
如果你學音樂的話.
你就知道什麼叫Coda.
Coda.
大家都不學音樂的.
Coda就是最後最後的時候.
那些交通樂團.
出到最多樂器.
所有樂器都走出來了.
這個就是.
一百四十六至一百五十篇.
我們剛剛看了就是一百第五十篇.
所以這些詩歌.
其實就是當時神的子民.
他們藉著不同的經歷.
有不同的創作.
他們就會.
收集起來成為今天.

$^{281}$我們有的.
舊約書卷裡面的.
詩篇.
好像我們.
現在少了.
現在沒有做這些東西.
以前我們那一代.
我們就有不同的詩集.
大家記得嗎.
那些比較成熟一點的信徒.
你記得嗎.
有生命性詩.
什麼青年性歌.
記不記得.
青年性歌有綠藍紅.
記得嗎.
記得.
或者.
有種主性歌.
有一本的.
為什麼要有那本呢.
以前那一代進到教會.
座位後面.
有個slot在那裡.
有個縫隙在那裡.
進來敬拜的時候.
拿起那本詩集.
主席就說我們今天.
唱第三首.
我們一起翻開第三首.
那裡有譜有詞.
是before.
有所謂projector.
記得嗎.
記不記得有projector.
這回事.
我只是說成熟的信徒.
年輕人不懂的.
上帝祝福你.
以前我們沒有.

$^{321}$所謂那些膠紙.
的projector.
之前我們就用那些詩集.
打開.
揭開它.
接著.
慢慢科技就發展.
有overhead projector.
記得嗎.
那時候我們敬拜的時候.
有個人坐在前面.
他就有張很多膠紙.
第一節他就放上去.
第二節就放上去.
有時候會放錯的.
有沒有做過這個位.
沒有.
有時候要寫上去.
或者印出來.
然後才有.
powerpoint.
開始有powerpoint.
是不是.
現在又有其他的科技.
我們又用其他的software.
等等.
你記得以前那套.
有一個集體的.
詩集大家都認識的.
所以我們就會.
唱那些詩歌.
詩集其實.
同一個.
過程.
去整理出來.
我記得以前.
在神學院的時候.
以前讀M.Dev的時候.
嘩.
回想起三十年前.

$^{361}$嘩 這麼老了.
三十年前.
那時候讀M.Dev的時候.
有些同學.
有些同學是very very much.
很著緊.
敬拜裡面是唱什麼歌.
他們是那些.
所謂改革中的教會.
那種傳統的.
保守的改革中教會很有趣.
他們的詩集只有150首.
就是這150首.
你打開.
你進去改革中傳統的教會.
打開詩篇的時候.
詩本 唱什麼歌呢.
只有150首.
一日復一日 一年復一年.
都是唱那150首.
只不過是將它配樂.
因為他們認為什麼呢.
他們認為上帝所啟示的話語裡面.
讓我們作敬拜.
其實就是.
只有這150首.
就不可以唱那些.
剛才Percy選的那些歌.
不可以的 他們覺得接受不了.
只能夠接受詩篇1至150篇.
那就配不同的樂曲上去.
那詩琴就很重要.
那詩琴就要彈.
還有一些更加保守的.
簡直是不用鋼琴.
還有一些鍵盤.
只有8個鍵.
就是.
我們今天唱第3首.
然後我們走去.

$^{401}$按第1粒進去.
主要是幫助大家.
不要走音.
就是這樣清唱.
有些同學真的是這樣的.
有些有一點拗頭.
就算我們所謂傳統的詩歌.
奇音典都不合格.
因為奇音典.
不在這150首.
Yeah.
所以不能唱奇音典.
太慘了.
所以任何的歌都不合格.
只能夠唱詩集裡面有的歌.
但我就有一點拗頭.
聖經裡面叫我們唱新歌.
何謂新歌.
剛剛我們唱那幾首.
You have to sing new song.
我覺得這個對聖經裡面.
執著的了解.
不是很對位.
聖經通常都是描述一件事.
Describe一件事.
它就沒有說去prescribe你某一些東西.
有很多地方有prescribe.
就是會規定你.
做一些事.
但敬拜上面.
通常它只是描述當代的人.
他們用什麼詩歌或者用了什麼音樂.
來去敬拜.
就沒有去dictate或者去決定.
世代代用什麼音樂.
或者用什麼詩詞.
所有的詩歌都是神學的application.
神學的應用.
我們對上帝的了解.
對上帝的認識.

$^{441}$那個所謂theology的神學.
我們應該怎麼應用在現在這個世代.
就算你寫了一首詩詞.
寫了一首詩歌或者寫了一首文章.
然後唱出來.
所以聖經就很少講到敬拜的時候.
會去.
決定.
我們要怎麼做.
所以我剛才講給你們聽你們覺得很好笑.
有些很傳統的.
連鋼琴都不用.
因為他們覺得勝經沒有鋼琴.
那我就心想.
那要怎麼樣.
我們就要.
好像耶穌的樣子.
是不是.
所以沒有說定死的.
但是是應用.
所以為什麼世世代代的神的子民.
是需要唱新的歌.
因為在這一代.
我們要回應這一代的需要.
在YouTube上面.
最近幾年特別是疫情之後.
有一班年輕人做得很好.
他們.
現在不知道還有沒有做.
年尾的時候他們就會做.
和我們一起經歷的.
在2024年.
這樣才算.
2024年和我們一起經歷的詩歌.
然後他們就把它們合起來.
唱在一起.
記不記得.
這些就是年輕的那些字.
他們就將.
去年.

$^{481}$經常唱的詩歌.
或者他們覺得很感動的詩歌.
他們就把它們串燒.
串在一起.
唱出來.
為什麼呢.
特別是香港背景.
香港的年輕一代.
他們的信徒.
他們的經歷就是.
很多時候.
觸碰到他們.
使得他們心動的.
這些詩歌就是和那些.
當時的社會的情況.
或者背景有關.
可能是一些哀歌.
可能是一些抒發他們的感情.
或者對上帝的一些呼喊.
他們就將它們合起來.
所以有些新歌是需要唱的.
We have to sing new songs.
詩篇第150篇.
這裡怎樣去到.
告訴我們什麼是敬拜.
你說如果不是.
去到決定我們用什麼音樂.
或者用什麼樂器.
用什麼詩歌去敬拜.
這裡是在說什麼呢.
是在說什麼有關敬拜的.
第一樣他說到.
瑞士我們要讚美敬拜.
第一節這裡說什麼.
他說你們要讚美耶和華.
在上帝的聖所讚美他.
在他顯能力的窮倉.
讚美他.
在他講到聖所的時候.
講到耶路撒冷.

$^{521}$聖所是什麼.
聖所不是人去的地方.
在近東的世界裡面.
所有的聖所.
都是上帝居住的地方.
是上帝的家.
上帝在那裡是坐著為王.
掌權在地上.
他去管理整個大地.
所以在聖所裡面.
去讚美他.
就是來到上帝的王權的面前.
去敬拜.
在他顯能力的窮倉讚美他.
這個就是近東.
古代近東人怎樣看這個世界.
這個講了好幾次.
他們認為這個世界是.
就是上面有很多水.
然後誰能夠有一個窮倉.
窮倉這個字.
在希伯來文其實是一隻鑊.
一隻很硬的鑊.
然後有一個鑊能夠.
頂住這些大水.
不會被它淹沒我們.
這個就是值得我們敬拜.
所以他說.
他顯能力的窮倉.
讚美他.
就承認這位主是比宇宙更大.
比我們人生更大.
是不是我們人生都會有高低.
但神比這一切都更大.
所以為什麼.
我們值得去讚美他.
所以這是第一個要回答的問題.
Who to worship.
第二個問題.
他想問的就是.

$^{561}$為何要讚美他.
What is the reason.
為什麼要去讚美他.
我們走進來的時候.
很多時候我們來到教會我們都不問這個問題.
因為第一個問題.
就是問泊車泊在哪裡.
希望你們也問這個問題.
有沒有遲到.
是不是要趕路來.
我們不是問Why do we worship.
Why do we get together.
為什麼要走在一起.
這裡告訴我們.
Because.
因他大能的作為讚美他.
按著他極美的大德.
讚美他.
這篇詩篇一開始就說.
你們要讚美我.
其實是命令式.
希伯來文其實就是Hallelujah.
Hallelujah.
剛才有沒有說這句說話.
唱這句說話.
Hallelujah其實就是希伯來文.
所以你去教會.
學了很多希伯來文.
Hallelujah.
你們要讚美我.
這裡在希伯來文.
你們要讚美我.
其實就是Hallelujah加Ya.
Hallelujah就是你們要讚美.
Ya就是Ya為的短寫.
所以你們要讚美我.
就是Hallelujah.
加Ya.
在這篇詩篇裡面.
十次的去到命令我們去讚美.

$^{601}$就是將我們的Attention.
將我們的注視力.
由我們自己.
轉移到.
天上這位掌權的神.
這個就是讚美.
敬拜的意思.
你進來的時候問Why.
這個不是一個純粹.
責任或者刻板.
要去做的屬靈作業.
我們要敬拜的時候.
我們要去問為什麼.
所以這裡講到.
要因什麼.
祂的大能的作為.
就是上帝創造這個世界.
創造我們的生命.
我們值得去讚歎.
值得被震撼.
大家都知道我最近幾個月.
我突然間很早起身.
我還是很早起身.
不知道為什麼.
Touch住那個鍵.
過去幾個月.
人可能老了.
我就五六點就起身.
我就會.
駕車去.
LA Fitness.
很多時候經過.
Lake Shore.
因為我去High Park.
在我駕車去Lake Shore的時候.
就會見到日出.
大家有沒有見過.
在Lake Shore見日出.
很少.
每次駕車經過.

$^{641}$我都覺得很讚歎.
很Wonder.
每一天的日出都很不同.
還有很美.
有幾次我忍不住.
我想停下車拍照.
不過當然沒有.
因為趕時間.
但欣賞那個創造.
欣賞到上帝的大能.
這個就是詩篇.
邀請我們做的.
在我們很忙碌的生活裡.
匆匆忙忙裡.
我們看看.
上帝自己的大能.
這是第一個原因.
第二個原因就是按照.
他的極美的大德.
Supreme Goodness.
因為這樣的緣故.
來讚美他.
什麼叫極美的大德.
其實詩篇整個鋪排裡.
已經講了.
上帝有什麼極美的大德.
剛才我們看了下面.
上面那兩個箭咀.
整本的詩咒告訴我們.
其實是一個故事.
150首的詩歌.
不是今天那些重主性歌.
他們是講一個故事.
整首的.
整本的詩咒一至到150篇.
其實是講一個故事.
開頭第三篇開始.
就是講到大衛的禱告.
因為他受苦在一個角度裡.
上帝應許了他.

$^{681}$然後建立他的角度.
然後一直唱下去.
你見到第89篇.
就是大衛的角度衰落.
因為流亡.
亡國的時候.
所以89篇是一篇很慘的哀歌.
所以.
1,2,3都很多哀歌.
很多的Lament.
因為神的子民受苦.
他們向上帝哀嘆.
去呼求的時候.
到89篇就更加更加坎坷.
因為89篇我給了一節給大家.
第38節.
是一個哀曲的禱告.
就是說上帝所建立的神的子民.
大衛的角度.
居然崩潰了.
沒有了.
89篇第38節.
但你的老奴.
上帝的老奴.
因為上帝的子民他們犯罪.
他們敗家.
然後整個的角度.
被毀滅.
因為你的老奴.
你的受高者.
就是君王,大衛等等.
當然最後是基督.
拒絕了他然後離棄了他.
主啊.
你從前憑信實.
向大衛起誓的慈愛在哪裡.
所以你可以想像.
在一個人生的失敗裡面.
整個國度亡國的裡面.
他們去哀嘆.

$^{721}$去作哀歌.
因為你知道以色列整本的舊約.
所謂The Beginning.
Of God's Grace.
就是什麼時候開始.
就是摩西.
摩西帶領以色列人出埃及.
然後以色列成為一個新的國度.
新的子民.
所以就將90篇這個.
神人摩西的禱告就放在.
卷四的開頭.
因為之前就是一個哀嘆的角度.
沒有了,但是到90篇.
摩西突然間好像.
上了舞台.
然後他就唱出第90篇.
哀歌.
就唱什麼呢?.
耶和華要到什麼時候呢?.
使我們早早飽得利的慈愛.
好使我們一生一世.
能夠歡喜快樂.
這個就是摩西.
有的禱告,帶領以色列人民.
第一次再去禱告.
所以你慢慢看到90篇開始捲四.
然後到捲五,就是這個國度.
有一個復興的盼望.
因為上帝的信實.
不是因為人的信實.
其實你這樣看的時候.
是一個大的敘事.
整篇詩篇.
就是在說這件事.
就是上帝的恩典.
在人的成功失敗.
然後再給他一個.
所謂second chance.
給他能夠有一個revival.

$^{761}$其實你看舊約的.
故事.
都好像我們故事一樣.
我們人生的故事裡面.
經常都會經歷這些.
89篇.
我們哀嘆.
上帝的信實在哪裡.
我們要等到什麼時候.
我們受苦的時候,上帝在哪裡.
可能我們很傷心,很沮喪,很挫敗的時候.
這個就是89篇.
但整個的詩集就是.
告訴我們,既然有一個上帝的恩典.
我們就不至於失去盼望.
何況我們今天有耶穌基督的死而復活.
上帝的子民.
最後.
是guarantee.
一個.
的victory.
一個勝利.
不是因為我們怎樣.
而是我們縱使人生都散了.
塌陷的時候.
縱使我們不順.
但上帝永遠是信實.
所以到第150篇的時候.
他就這樣結束,要讚美耶和華.
我們的生命最後的時候.
要讚美誰,不是我們厲害.
是上帝的恩典.
然後這篇詩篇就解答另外一個.
有關敬拜的問題.
三至五節就講到怎樣去敬拜.
你看到有多少樂器在這裡.
一,二,三.
可能七件吧.
是當時用的樂器.
重點不是.

$^{801}$一定要用這些樂器.
重點就是.
當時他們用各種.
拿了所有的工具.
所有不同的樂器.
他們去敬拜這位這麼信實的主.
今天我們怎樣去應用呢.
我們都需要拿了所有工具.
我們都應該要.
這樣去敬拜.
特別講到敬拜是一回甚麼事呢.
是一個動詞.
你見到有跳舞.
有擊鼓,全部都是有動作的東西.
這些我們華人.
基督徒的文化我們真的不懂.
因為我們就是.
好像碌木一樣.
或者我們表情.
亦都好像碌木一樣.
我們不懂怎樣表達喜怒哀樂.
因為我們從小到大.
我們的提升就不是鼓勵我們去做這件事.
是不是.
最重要是有Poker face.
樣樣東西都有Poker face.
無動於衷.
這樣才是強人.
才是好的.
但聖經裡面講到敬拜不是這樣的.
是有喜樂.
有喜怒哀樂.
要站起來.
要有動作.
要跳起來.
整個人去表達.
嚴肅不一定是屬靈.
屬靈的人.
都不一定是嚴肅.
明白嗎.

$^{841}$所以敬拜是一個認真的敬拜.
不一定是嚴肅的敬拜.
可能有這樣的空間.
是需要嚴肅的敬拜.
不過這裡講到.
是動詞.
是一個動詞.
身體是需要動.
你的心也需要動.
這個才是真正的.
整個人的敬拜.
到最後這裡講到.
誰應該去敬拜.
第六節.
凡有氣息的.
都要讚美耶和華.
凡有氣息的.
你看到有鋪排.
就算有最好的樂隊.
最好的樂器.
最好的音樂也好.
到最後第六節怎麼結束.
凡有氣息.
都應該要讚美耶和華.
氣息就是我們生命裡面.
最基本的禮物.
今天我們有口氣.
是上帝給我們最大的恩典.
對嗎.
到你有一天.
在CU的時候.
你就知道氣息有多貴.
氣息是很寶貴.
所以這篇詩篇說.
凡有氣息的都要讚美耶和華.
今天上帝給你一口氣.
你應當.
去讚美祂.
應當去敬拜祂.
在世間上有很多.

$^{881}$不同的權柄.
要求我們做這件事.
要去讚美祂 要去撐祂.
但是上帝的指紋.
那首音樂的節奏.
是不同世間的.
我們撐的.
唱的.
讚美的.
是第一位.
我們主耶穌基督.
我希望今天給你一個.
這樣的溫習.
什麼叫敬拜.
因為你記住.
敬拜的原因.
敬拜的意思.
以致我們聚集.
一起的時候.
就算有些誤離教會.
平常誤離教會.
他們進來的時候.
他們會知道這裡有不同.
因為耶穌基督.
值得敬拜.
耶穌基督是真神.
\newpage



\section{以弗所書 1:9-10}
\label{sec:031739}
\textbf{以「合一」迎主來 (華)}~~~~葛國光牧師
\newline
\newline
連結: \href{http://www.church.com.hk/acms/content.asp?site=cdc\&op=show\&type=product\&code=031739\&layout=sermon}{\texttt{031739}} ~~~~ 語音日期: 2026-03-01
\newline
\newline
\hyperref[sec:031738]{\small{< < < PREV SERMON < < <}}
~
\hyperref[sec:index_chronic]{\small{[返順時目]}}
~
\hyperref[sec:031740]{\small{> > > NEXT SERMON > > >}}
\newline
\newline
以弗所書 1:9-10
\newline
\begin{longtable}{cl}
\hline
\hline
章節 & 經文 (和合本修訂版)\\
\hline
1:9 & \begin{tabularx}{0.7\textwidth}{X} 照自己在基督裡所立定的美意，使我們知道他旨意的奧祕， \end{tabularx} \\ \\ \relax
1:10 & \begin{tabularx}{0.7\textwidth}{X} 要照著所安排的，在時機成熟的時候，使天上、地上、一切所有的，都在基督裡面同歸於一。 \end{tabularx} \\ \\
[1ex]
\hline
\hline
\end{longtable}
$^{1}$邱子貝 我想這個世界每一天都在變化.
最近發生什麼大事 你知道嗎.
每一個月的第一個禮拜是我們的聖餐主日.
我都會要講主在來的事情.
可能我們已經久了 時間也慢慢地.
就像聖經裡面有一些人機竅說.
主來 講那麼久了.
主耶穌講這主在來已經一千多年了 兩千年了.
主還是沒有來啊.
親愛的邱子貝 我們不要像這些人一樣.
是那個機竅 機諷神話語的人.
主說他要快來.
聖經裡面讓我們看見各樣的跡象都顯示主快來了.
如果你對這個聖經稍微有點了解.
聖經裡面在舊約以西節書38章39節.
講到有一個戰爭叫做哥格與馬格的戰爭 聽過嗎.
當然這一個主題是一個很大主題.
我今天不是主要在講這個.
我要講的就是因為最近伊朗的事情.
大家全世界都知道 是不是.
86歲的哈米尼被斬首.
跟著他一起的40位都是這些重要的軍官.
包括國防部長 包括參謀總長全部都被炸死.
他丟死你知道嗎.
其實你去看伊朗的影響力非常大.
包括葉盟 包括他周遭的這些的一些國家.
還有包括加沙的哈馬斯.
這些通通都是受他資助影響的.
可是你看這一段聖經裡面.
以西節書裡面讓我們看見一件事情.
就是聖經裡面預言說.
將來會在這個地區.
伊朗跟這些周邊的這些國家.
包括這個利比亞.
包括你看我把它圈出來的土耳其.
這些國家包括這裡面提到北方的急促的.
這個說法多數都認為說.
可能那個地方比較是指的是俄羅斯所在的地方.
這些對地方聯軍起來要一起來攻打以色列.
所以你可以看見中東可以說是在末世的預言當中.

$^{41}$一個非常重要的一個角色.
這個當然也牽扯到以色列以後的一個發展.
聖經上裡面提到說以色列要全家得救等等.
等等這些包括第三聖殿的這些的建立.
敵基督的出現.
這些整個末世的預言跟中東可以說是息息相關的.
哥哥馬哥的這個戰爭到了啟示錄第20章的時候.
又提到一次.
但是這兩個地方不太一樣.
這個地方是跟北方的這些國家聯合起來攻打以色列.
但是到了啟示錄20章的時候.
那邊就開始講了說是一個全部.
可以說是整個敵反對基督徒的.
反對這個以色列的所有的這些敵對的軍隊.
統統要集結起來.
要一起來攻打這個神的兒女.
不管他第一個跟第二個一樣不一樣或這些問題.
但是我們看見這些事情.
聖經裡面所說的這些預言正在一個一個的在實現.
主耶穌在馬太福音24章.
也清楚地講到整個末世的一些的徵兆.
我們一起來讀馬太福音24章.
第七節.
我們一起來讀.
預備 請.
民要攻打民.
國要攻打國.
多處必有悽惶.
地震.
好 這裡提到說民要攻打民.
國要攻打國.
COVID-19讓整個世界安靜了一段時間.
但是後來沒多久.
你就發現一些國家就開始打起仗來.
有沒有.
包括俄羅斯跟烏克蘭.
大家都戴口罩.
他們不戴口罩.
因為打仗還戴什麼口罩.
只能戴防毒面具.

$^{81}$你還記得嗎.
回想到四年前.
所以好像這個世界沒有辦法安靜下來.
到處都在征戰.
聖經已經清楚地讓我們看見.
末後的時候就是如此.
明天會更好嗎.
世界有一天會和平嗎.
會不會.
如果你從聖經來看是不會的.
所以聯合國他沒有辦法.
他註定要失敗.
因為這個問題的解決.
不是靠一個組織可以解決.
不僅僅是國家之間.
就在川普總統要攻打伊朗之前.
他發表了國情的諮文.
在國會發表.
據說是歷史上以來.
總統發表國情諮文最長的一篇.
但是當他在臺上講的時候.
下面怎麼樣.
你會很明顯地發現.
這個有鼓掌的聲音有沒有.
但是還有呢.
還有噓聲.
甚至還有對罵聲.
共和黨民主黨.
你可以看見美國這些年來.
政治的立場的分歧.
兩黨的對立.
也強化了人民之間的對立.
有沒有.
每一次的選舉.
都會帶來這樣子的一個結果.
這是很悲哀的一件事情.
曾經被美國人引以為豪.
也被世界人所推崇的.
所謂的民主制度.
也不過是如此.

$^{121}$我們不要把民主制度當神.
也不要把民主制度當做是神的旨意.
其實我們一再講.
聖經裡面講到老底家教會是什麼.
老底家原文的意思就是人民做主.
就是民主.
老底家是什麼樣的一個教會呢.
是主所責備的教會.
我們今天要的不是一個民主的國家.
民主的世界我們要的是什麼.
是一個上帝在掌管.
一個以基督為中心的一個國家.
或者一個世界.
Amen.
Amen嗎.
很多人還在思考中啊.
有問題嗎.
這個世界越來越紛爭.
但是丟子妹.
主的心意不是如此.
我們今天所讀的這兩句經文.
你看見一件事情.
神的心意是什麼呢.
他說要使天上地上一切所有的.
都在基督裡面同歸於一.
格魯西斯保羅在一章二十節也這樣說.
他要使所有的萬有.
無論是地上的天上的都跟自己和好了.
所以你會發現一件事情.
神的心意是要這個世界是合一的.
因為這是神的旨意.
都是照他自己所預定的美意.
叫我們知道他旨意的奧祕.
他的美意.
他的旨意.
神的旨意是什麼呢.
就是天上地上一切所有的.
都要在基督裡面有一天要同歸於一.
Amen.
但是丟子妹.

$^{161}$為什麼今天這個世界為什麼會這麼混亂.
我們看看上帝起初在創造的時候是怎麼說的.
上帝叫有光就有了光.
有暗就有暗.
是不是.
上帝叫光跟暗分明.
上帝看著是好的.
有沒有.
好的.
每一天照完就說好的.
好的.
上帝照的都是好的.
跟你旁邊說上帝照的都是好的.
最後一天當神創造完了之後.
上帝怎麼說.
上帝看著他一切所造的都甚好.
不僅僅是好是甚好.
親愛的丟子妹啊這個世界多好啊Amen.
好嗎.
最近南加州下了雨.
好不好.
你如果沒有下雨你就知道多不好.
是不是.
所以現在下雨多好是不是.
但是很多地方下雨不斷.
洪水都來了好不好.
現在這個世界到處都是這種極端的天候.
對不對.
本來到北極去看冰.
現在看什麼.
看滑冰.
所以整個世界亂套了.
整個世界不像起初上帝創造的時候那種光景了.
為什麼會如此.
大地好像都受了咒詛.
原來創世記第三章發生了什麼事情.
發生什麼事情.
亞當,夏娃犯罪是不是.
亞當,夏娃犯罪是怎麼犯罪的.
誰叫他們犯罪.

$^{201}$魔鬼引誘他是不是.
親愛的丟子妹你看到沒有.
所以魔鬼從一開始.
他就是要人分裂.
他就是要人離開神遠離神.
這是魔鬼的計劃.
跟上帝的計劃是完全不同的.
上帝的心意是合一.
對不對.
但是魔鬼的計劃是什麼.
是分裂.
所以親愛的丟子妹請你記得一件事情.
上帝的工作是合一.
魔鬼的工作是分裂.
哪個地方有分裂.
哪個地方就是魔鬼在那邊掌權.
哪個地方有合一.
就看見是上帝在那邊掌權.
Amen.
所以保羅在講的《以弗所殊》.
講到整個的一個藍圖.
你知道《以弗所殊》是在講到教會的藍圖.
他講到神的旨意.
講到神屬靈各樣的福氣的時候.
講完了之後他就講到這個恩典.
這個恩典呢.
他就講到這個恩典是最後要帶來一個什麼樣的恩典呢.
就是所有的一切.
天上地上所有的一切都在基督的裡面同歸於一.
同歸於一這個字原文的意思.
其實原文有17個字.
它是兩個字合在一起的.
前面Anna是"在"的意思.
所以它的意思是什麼呢.
就是"在"合在一起.
英文說把它sum up.
你們有沒有學會計的.
或者是你自己做簡單的計算.
用Excel.
你把這個禮拜支出多少錢.

$^{241}$或者這個收入多少錢.
你把所有的錢.
收入的錢通通把它列列列.
列完之後你最後用一個什麼.
sum.
sum是什麼.
總計.
總結.
是不是.
所以這個字的意思就是.
其中一個意思就是sum up again.
換句話說again表示.
這是本來神的心意就是如此.
神本來就是要萬有.
通通都在他的裡面是合一的.
夫妻是合一的.
家庭是合一的.
教會是合一的.
國家是合一的.
整個世界都是合一的.
甚至人與大地萬物通通都在基督的裡面.
是合一的.
這也就是將來千禧年將來這個新天新地.
神應許的時候你看見沒有.
整個人跟動物野獸.
通通都是和平的相處有沒有.
在你聖經裡面講到千禧年的時候.
親愛的弟兄姊妹.
那這是神的心意.
神要宇宙萬物都在基督的裡面合而為一.
所以將來請問弟兄姊妹.
我們到天堂的時候還會吵架嗎.
會不會吵架.
不會.
那為什麼在地上要吵架.
將來不會吵架嘛.
所以在地上好好吵一吵是吧.
親愛的弟兄姊妹請問在天上.
我們還會聽到打架嗎.
還會看見有人用槍炮彼此在攻打嗎.

$^{281}$不會.
那是神的果.
那是神的心意.
上帝起初在創造伊甸園的時候創造的世界.
伊甸園就象徵著今天的世界.
他所創造的起初的心意就是如此.
所以你看亞當在那邊叫大象的時候.
他叫獅子他叫猛虎.
他叫他的名字就叫他的名字.
他叫他來他就來.
這是起初上帝創造的意識.
請記得一件事情.
神的旨意神的計劃會不會被破壞.
不會.
雖然仇敵魔鬼想要破壞他.
但是神至終他的旨意一定成就.
你相信這件事情.
神的計劃必定要成全.
神在你身上的計劃.
神在你家庭的計劃.
神在他教會的計劃.
必定要成全.
這是神的心意.
所以親愛的兄弟姐妹.
反過來.
這個.
合一是一個祝福.
反過來呢.
分裂是宙主.
我們看見.
這個世界到處充滿紛爭.
根據統計.
美國的.
按照所謂的仇恨犯罪率.
根據這個統計.
它裡面這個表格.
你看最高的是51.5\%.
有沒有.
這個51.5\%是因為種族.
就是race.

$^{321}$ethnic.
根據因為這個關係.
而帶來的仇恨犯罪.
是佔到最高51.5\%.
那你看看.
特別我們洛杉磯.
紐約這些大城市.
各式各樣的不同的種族.
是不是.
不要說外面了.
我們華人教會也是大概至少兩個語言.
是不是.
所以種族的不同.
語言的不同.
很容易就造成了.
這種仇恨.
除了種族的關係.
另外23.6\%.
是因為宗教的關係.
你可以看見伊斯蘭教.
這個跟猶太教.
伊斯蘭教跟基督教.
有沒有.
這種仇恨的犯罪.
也是佔了第二高.
第三個是性傾向.
佔了18.1\%.
這個是美國左派政府搞出來的東西.
就是因為LGBTQ這個movement.
所造成的這種仇恨犯罪.
所以我們看見.
今天美國也是到處都是充滿著一個紛爭的一個情況.
家庭呢.
美國每一年夫妻離婚所帶來的影響.
產生了兩百萬的孤兒.
所以你看今天為什麼美國這麼多的一些好像定時炸彈.
好像不定時炸彈在那個地方.
有沒有.
很多孩子一下走進校園.
這個人也沒有犯罪記錄.

$^{361}$沒有前科.
但是你去看看這些的家庭.
很多時候是出了一些問題的.
所以家庭出問題.
夫妻關係的問題.
造成了許多的孩子也受累.
這是今天美國的家庭的很大的問題.
在這一切的當中.
親愛的同志們.
帶來的結果是什麼呢.
我們看看聖經上的話.
我們一起來唸好不好.
來 預備請.
凡一國自相紛爭就成為荒場.
一臣一家自相紛爭必站立不住.
一國自相紛爭就成為荒場.
所以其實今天美國最大的問題.
還不是經濟的問題.
美國最大的問題是什麼你知道嗎.
今天美國走在一個.
充滿了衝突對立的一個情況.
其實一個不好.
也很可能就.
也許就會再次發生內戰.
記不記得.
記不記得美國南北戰爭的時候.
我記得有一次陳導還提.
林肯總統當時要發佈解放黑奴宣言的時候.
他要面對到一個很大的一個很大的一個掙扎跟壓力.
為什麼.
因為當時候南北戰爭裡面對於解放.
對於黑奴的這一個問題是截然不同的.
甚至在他自己的這個他們這個北軍的裡面.
也有不同的聲音.
所以林肯尋求神的心意.
他清楚知道神的心意.
因此他勇敢的做了這一個這個解放黑奴的宣言.
神與他同在.
所以親愛的兄弟妹你會發現那個關鍵在哪裡.
關鍵就是我們是不是在神的旨意的裡面.

$^{401}$一個國家如果他們敬畏神.
這個國家神是祝福他的.
以耶和華那國為神的.
以耶和華的神為神的那國是有福的.
就說明美國本來是一個敬畏神的國家.
本來是我們的神很清楚.
現在你必須要講清楚到底你信的是哪個神.
這邊有穆斯林的阿拉.
這邊還有這個什麼什麼各式各樣的宗教.
有印度教有各式各樣的宗教.
所以你要說你是信神的.
你必須要講清楚你信的到底是什麼.
可是原本在美國沒有這個問題的.
一城一家自相紛爭必站立不住.
所以紛爭的結果就一定是.
最後是會走向滅亡.
一個家庭也是如此.
當然你會說哪有一對夫妻不吵架的呢.
但是請問講這句話的時候.
我們覺得吵架好嗎?.
不好,是吧?.
不是給自己找理由吧?.
吵吵吵,很多的夫妻吵到後來就怎麼樣?.
就sayonara.
是不是?.
拜拜.
有時候為一件小事情可以吵翻天,有沒有?.
有沒有?.
幸主時間久一點的回想一下.
有時候剛剛結婚的時候你為了一些什麼事情吵.
真覺得說很幼稚.
有沒有?我以前很幼稚.
跟你旁邊說不要那麼幼稚.
可是你就看見.
我們的老我就是如此的幼稚.
我們的老我.
我們的自我中心.
我們的舊人在我們的裡面.
他好像是一個嬰孩一樣.
意思就是他沒有長大.

$^{441}$老我要死掉.
對不起,應該講我們的那個新的人在我們裡面.
好像是嬰孩一樣,對不對?.
老我很大.
但是我們的新我呢.
神賜給我們那個新的我呢.
那個新的生命呢.
還是一個baby.
所以你永遠打不過這個老我.
這個老我強大到一個地步.
好像哥利亞一樣.
你根本無法打敗他.
親愛的同事們.
這是今天很多不信神的人他的問題.
沒有神的人根本都是老我.
是吧.
但是我們信了主的人.
需要天天的背起十架蛇脊來跟隨主.
那麼怎麼樣子.
我們才能夠成為一群合一的團體呢?.
剛纔我們讀的那三個聖經裡面.
你會發現都在強調一件事情.
就是在基督的裡面.
叫萬有地上一切所有的.
都在基督裡面同歸於一.
看到沒有.
在基督裡面同歸於一.
保羅在這裡也提到說.
藉著基督.
藉著他.
因著耶穌基督的緣故.
他拆掉了我們中間隔斷的牆.
使得我們與神和好.
然後我們彼此之間能夠彼此和好.
換句話說.
如果我們本來沒有耶穌基督.
我們跟神是什麼.
我們是跟神是不好的.
我們跟神是遠離神的人.
我們彼此之間是什麼.

$^{481}$我們彼此之間是不好的.
是敵對的.
所以這是神.
藉著他的兒子耶穌基督.
要藉著他定在十誡上面所做的這件工作.
就是每一次當我們在這裡.
守聖禪的時候.
請你記得一件事情.
耶穌基督定十誡.
不是一個很成就的一個古老的故事.
耶穌基督的十誡.
是人生當中我們可以靠著他.
來勝過一切困難.
一切的壓力.
一切的挑戰的根據.
也是因著耶穌基督的十誡.
也是我們能夠得勝老我.
勝過我們自己.
勝過我們這個老脾氣.
勝過我們這個救我的人.
的唯一的方法.
所以藉著在基督的裡面.
我們才能夠真正的合一.
什麼叫在基督裡.
丟子妹.
你在基督裡嗎.
在嗎.
還是.
禮拜天在.
天天都在.
哈雷路亞.
你在基督裡嗎.
我們需要常常問這件事情.
你在基督裡嗎.
這個答案.
這個題目.
這個問題.
我們可以分兩個階段來談.
第一個就是.
你是在基督裡.

$^{521}$還是你在基督外.
也就是說.
你是屬基督的.
你還是.
在基督以外.
是屬魔鬼的.
只有這兩個.
只有這兩個歸屬.
所以已經清楚了.
我們可以只有這兩個歸屬.
沒有別的.
是不是.
因為全世界都握在那.
惡者的手下.
所以你不是屬基督的.
你是屬誰呢.
你是屬魔鬼的.
那我請問丁永子妹.
你是屬誰的.
你是屬基督的嗎.
跟你旁邊說.
你要屬基督啊.
你的家要成為一個屬基督的家.
你的家要在基督裡面.
纔會有真正的和睦.
我曾經跟丁永子妹分享過.
我從小的我們家.
是沒有耶穌基督的.
我們是拜拜的.
所以我們家在我成長的過程當中.
就是風聲雨聲麻將聲.
聲聲鹿兒.
其實那個風聲雨聲就是吵架聲.
所以我小學四年級.
我就會打麻將.
假以時日.
我今天如果繼續這樣練功夫的話.
不得了.
還好小學五年級.
一個同學把我帶到教會.

$^{561}$真是神拯救我.
所以我成為我們家第一個信耶穌的人.
然後一步一步的到了我大學畢業.
其實我清楚知道神已經呼召我.
只是我沒有跟家人講.
然後在一個累贅當中.
神呼召我做傳道人.
就在大學畢業的前一段時間.
我就寫信跟我爸爸說.
我要去獨身學院.
好 美根羅斯德的秀.
兄弟們.
可能你聽我講我的見證.
講過很多次.
沒有關係.
你要常常講你的見證.
傳福音要先講你的見證.
你講講講.
你纔會常常提醒你.
我們是何等蒙恩的一個人.
你自己的故事常常講.
講到一個地步.
你非常熟悉.
隨時隨地.
你說我勢力背不起來.
你說我福音橋背不熟.
沒關係.
你的見證應該講得很清楚.
你的見證應該是怎麼樣.
信守戀來.
是不是.
匆匆忙忙.
遊刃有餘.
不要一上臺做見證.
哎呀.
結果匆匆忙忙.
下面就不要說了.
所以當我神呼召我出來.
很奇妙的.
上帝透過這一個方式.

$^{601}$讓我的爸爸被受刺激了.
刺激得非常的深.
後來我就請我們教會的牧師去探訪他.
那個牧師很好.
半年之後.
我父親信主了.
從我父親開始信主.
我的嫂嫂開始信主.
我的哥哥開始信主.
我的姐姐開始信主.
一個一個的.
我們家的人信主了.
所以從那個之後.
我們家不再是吵架聲.
麻將聲.
聲聲入耳.
我們家是歌唱聲.
歡笑聲.
聲聲入耳.
Amen.
丟子妹.
你的家人.
你的先生.
你的太太.
你的孩子還沒有信主.
你不要把它當做是無所謂.
我看到有一些父母親孩子.
大學之後離開家.
大學之後就不信主了.
父母親提起來的時候.
沒有什麼特別的感覺.
丟子妹.
如果你都不覺得你的孩子沒有信主.
這件事情是非常嚴重的一件事情.
你怎麼希望別人幫你的忙.
瞭解我意思嗎.
你真的要為他流淚禱告嗎.
希望主啊.
你把他帶回來嗎.
我發現很多父母親不care.

$^{641}$是因為他自己的信仰.
他也不care.
親愛的丟子妹.
如果你真知道這位主.
這麼樣子的重要.
一切的問題都在他的裡面.
沒有認識這位耶穌基督.
你再有怎麼樣成就.
你功成名就.
你再有多少錢.
都是枉然.
所以我們要好好的為我們的家人.
為我們還不信主的這些親戚朋友.
我們要為他們禱告.
代理他們信主.
這是一件很重要的事情.
我們今天講全家事奉主.
第一件事情.
全家需要信主啊.
你不信主還是什麼事奉主呢.
對不對.
所以你最重要的.
你先把他家人要帶到你主的面前來.
那麼當我們都在基督的裡面的時候.
你可以看見我們的家就不一樣了.
以前我們家每一次過年的時候.
大家聚在一起喫完飯之後.
那就精彩了.
小孩子玩他們的.
大人就打麻將.
我們家過年可以擺兩桌.
兩桌都不夠.
因為房子的關係.
所以只能擺兩桌.
所以大家輪番上陣.
你懂意思是啥.
現在不需要了.
現在大家喫完飯.
大家一起.
全家聚在一起分享唱詩歌.

$^{681}$親愛的兄弟姐妹.
不容易啊.
這個過程不容易啊.
一個老妖要能夠出頭啊.
真是不容易啊.
可是靠著上帝.
我們可以凡事都能.
在基督的裡面.
我們都能.
你若能信.
在信的人凡事都能.
抓住神的應許.
所以應許我們說.
當信主耶穌.
你和你一家都必得救.
Amen.
所以御主一家信了主之後.
他們全家就都很喜樂.
看到沒有.
第二.
在基督裡面的意思.
代表是一個關係.
就像我剛講的.
你不是一個禮拜天.
來到這裡.
表示你是基督徒.
那隻能說你是.
星期天的基督徒.
在基督裡的意思.
是住在基督裡.
我們來看看.
傳道書第四章十二節.
我們想念好不好.
來 請.
有人公聖孤身一人.
若有二人便能抵擋他.
三股合成的繩子不容易折斷.
前面兩句話我們都瞭解.
一個人孤身.
兩個人就能夠抵擋.

$^{721}$那請問三股合成的繩子.
這三股是哪三股.
所以這兩個人再加上一個誰.
如果用在我們夫妻關係.
所以有人把它發展出來.
這個叫做婚姻鐵三角.
有沒有.
其實我去查了一下.
這個婚姻鐵三角.
原本是一個心理學家.
美國心理學家.
他提出來的一個叫做.
愛情三陰論.
三陰論 陰術的陰.
他認為說一段完整且長久的婚姻.
是由三個核心元素組成的.
這三個叫做親密.
第二個叫做激情.
第三個叫做承諾.
他說如果這三個在你的婚姻當中都有的話.
這就是一個完美的婚姻了.
我覺得他在講的是一個烏托邦.
你現在就是美.
你想不想跟妻子intimacy的關係.
你想不想你們之間是有passion的.
夫妻之間充滿了強烈的吸引力.
浪漫感受.
每次在一起都很興奮.
像初戀一樣.
有沒有.
有沒有 想不想.
想嗎.
想 想歸想 做得到嗎.
你一看他那個表情要罵人的時候.
那個浪漫就沒了.
想跟他在一起的感覺都沒了 有沒有.
是不是 那怎麼辦.
所以你會發現所有這一些心理學.
如果不是在這個所謂的基督的.
用神的話語.

$^{761}$回到神的話語裡面.
你會發現很多是一些技巧或者一些教導.
或者是一些理論的東西.
或者是一些理想.
但是你很難做得到.
很難做得到.
我們怎麼樣子才能夠如此呢.
所以這個婚姻鐵三角告訴我們說.
這個婚姻鐵三角是什麼意思呢.
你看見當丈夫跟神之間的關係.
如果越來越接近的時候.
做妻子的跟神的關係越來越接近的時候.
一直往前走 往上走.
到達了.
你看見沒有.
你看那個虛線了沒有.
所以越往前跟神的關係越接近.
彼此之間的關係是怎麼樣.
也越接近.
瞭解這個意思嗎.
所以這個叫做婚姻的鐵三角.
也就是說你的婚姻當中要有第三者.
第三者 懂嗎.
這不是我說的.
這是由另外一個基督徒的婚姻專家提出來的.
這個叫做葛福偉牧師 也是我們葛家的.
不過他的last name叫做Kipferver.
怎麼會翻成葛.
他中文叫做葛福偉.
是在GooTV 現在非常在promote他.
他提出一個所謂的三方盟約.
他說這個第三者不是指的破壞關係的那個第三者.
這個第三者是指的聽我們禱告的神.
所以他的意思就是這裡所說的.
是哪裡所說的呢.
就是我們夫妻之間要藉著同心合一的禱告.
以至於讓第三者可以介入.
懂這意思嗎.
邀請這位第三者進入我們的婚姻當中.
親愛的兄弟姐妹.

$^{801}$這個原則聖經裡邊哪邊告訴我們的.
我們一起來唸好不好.
法帶於福音18:19-20 來預備請.
有兩個人在地上同心合一地求什麼事.
我們在天上的父必為他們成全.
因為兩三個人奉我的名聚會.
那裡就有我在他們中間.
你注意看.
前面提到兩個人同心合一地禱告.
神就怎麼樣.
神必為他們成全.
神有沒有說可能.
大概80\%.
有沒有.
他說什麼.
他說必.
必為你們成全.
當你們夫妻禱告.
當我們教會弟兄姊妹同心禱告.
當我們同心合一地來祈求的時候.
當我們的國家一起同心合一禱告.
神就會讓我們同心合一禱告的那個事情就成全了.
美國在諾曼第登陸的時候.
美國總統羅斯福提出來.
在廣播邀請全美國的百姓.
一起來為著這一次的登陸.
這一次的戰役禱告.
諾曼第登陸打勝仗.
扭轉了整個的這個世界大戰的結果.
親愛的同志們.
今天.
我們這個國家處在一個.
彼此對立的環境當中.
你叫大家禱告.
有的人還根本甩你都不甩.
是不是.
這是一個根本的問題.
這是今天美國根本的問題.
當我們不以這位神為我們的神的時候.
我們就沒有一個common ground.

$^{841}$我們就沒有一個共同的可以連在一起的.
你懂這意思嗎.
這就是使我們分歧的原因一個根本的原因.
那麼第二.
親愛的同志們.
教會需要同心合一.
教會同心合一.
你看見初代的教會.
耶路撒冷的教會.
他們同心合一.
聖經上告訴我們說.
他們同心合一一起橫切禱告.
然後聖靈降臨.
五星節降臨之後.
這些人所有的這些教會的丟子妹們.
他們順服使徒的教訓.
接著說什麼.
他們就在那個地方.
同一起彼此天天在一起彼此.
剝餅用餐讚美神得神的喜.
得眾民的喜悅.
主就將得救的人數天天加給他們.
菲律賓教會也是這樣子.
聖經上說他們從頭一天.
直到如今保羅提到他們說.
是同心合一的.
希望福音.
但丟子妹.
今天教會不幸的.
我們受這個世界的影響.
這叫世俗化.
許多的教會.
不是同心合一的禱告.
來決定事情.
而是看多數人的意見如何.
來決定我們要不要做.
是吧.
這叫做民主.
所以當我們投多數的人.
都覺得對的時候就是對嗎.

$^{881}$是不是.
請問耶穌被釘十字架.
是不是多數人要釘他.
有沒有.
有.
對不對.
所以什麼叫做多數人就對呢.
這不是聖經的道理.
只要這是神的旨意.
我們就同心去做.
所以同心合一禱告.
其實同心合一也跟能不能夠在一起禱告有關係.
夫妻不能夠同心禱告.
門都沒有.
我看到一些的夫妻.
什麼都可以在一起.
就是禱告沒法在一起.
我幫年輕人做婚前輔導的時候.
我都會要幫助他們做一件事情.
教他們怎麼禱讀.
兩個人一起學禱讀.
兩個人一起禱讀.
但是我發現有些人他們.
很難跟自己的另外一半禱告.
如果不能夠一起禱告.
代表你們的根基是非常非常的脆弱的.
同樣的.
如果同宮之間不能夠一起禱告.
或者是七嫂的禱告.
你再怎麼樣子很多的活動.
再很多的活動.
再很多的看起來很熱鬧.
最後我告訴你.
你會發現.
也不過是過眼雲煙.
丟子妹 我們要建立一個同心合一的教會.
昨天我們的國康師母.
就在這個講臺上面.
本來今天應該是宇康傳道帶領敬拜.
國康師母結果在敬拜練習的時候.

$^{921}$就突然昏倒了.
就整個人前面就筆直地倒下來.
就我們昨天去看她.
這邊有一個眼睛.
這邊的眼皮都黑了.
她的意識都還清楚.
所以可能是大概血糖過低之類的.
就暈倒了.
那我要講的是說.
本來她今天要來敬拜.
臨時出了這個事情.
馬上聯絡.
我昨天大概差不多十點多的時候.
接到宇康傳道給我的text.
那時候我正在有其他事情忙.
所以我看到的時候已經十一點了.
她就跟我說.
她今天不能來她還要照顧師母.
所以她說她已經把敬拜的事情都安排好了.
所以今天你看到我們鳳蘭.
鳳蘭在哪裡.
鳳蘭姐妹還有小雪師母.
他們都是臨時上陣的.
所以這就是一個團隊.
這就是一個同心合一的團隊.
我們盼望看見.
我們更多的是這樣子同心合一的團隊.
在我們教會里面出現.
不僅僅是我們國語堂同心合一.
三個堂同心合一.
我跟你講開再多的會不會同心合一的.
更多的禱告纔會同心合一.
我有機會在三堂聚會我都講這句話.
但是要從知道變成實踐.
這中間有一個很大的差距.
親愛的兄弟妹妹.
昨天我沒有直接看.
但是師母把張茂松牧師.
他們就在各地都舉行烈火特會.
昨天舉行的這場烈火特會.

$^{961}$全球有1.6萬人.
就是1萬6千人蔘加.
有的人在實體的.
有人在線上的.
一起參加.
我就聽到張牧師在講.
張牧師其實現在年紀已經大了.
而且又有癌症.
身體不好.
講話聲音已經有點沙啞.
他講說.
不知道講哪一句話.
講講.
然後就問弟兄姊妹.
弟兄姊妹的聲音.
其實我聽起來已經很大了.
他就說.
沒聽見.
然後下面就怎麼樣.
就更大聲了.
他再說一句.
沒聽見.
下面呢.
更大聲了.
弟兄姊妹.
不要看這個小動作.
你從這可以看見.
這是不一樣的教會.
是不是.
所以親愛的弟兄姊妹.
有時候.
你們不要怪牧師常常要提醒或者訓練.
這些東西就是操練我們是不是能夠同心合一.
如果我們所做的是神所喜悅的.
敬拜的時候.
讚美的時候.
大聲地讚美.
拍手來讚美.
不要好像是那個很宅男宅女的樣子.
你說我的個性就是如此.

$^{1001}$你說我的原生家庭如此.
那你信了耶穌幹嘛呢.
信耶穌不能夠改變你的原生家庭影響嗎.
信耶穌不能夠改變你的個性嗎.
上帝可以.
人若有願作的心.
這很重要.
當我們願意的時候.
神就會在我們身上工作.
當我們願意的時候.
神就會在我們的家庭工作.
當我們願意的時候.
神就會在我們的教會工作.
當我們願意的時候.
上帝就會在我們的國家工作.
因為神他是那位無所不能的神.
\newpage



\section{撒母耳記上 18:1-21:15}
\label{sec:031740}
\textbf{甚麼是我最終的依靠}~~~~李思聰牧師
\newline
\newline
連結: \href{http://www.church.com.hk/acms/content.asp?site=cdc\&op=show\&type=product\&code=031740\&layout=sermon}{\texttt{031740}} ~~~~ 語音日期: 2026-03-01
\newline
\newline
\hyperref[sec:031739]{\small{< < < PREV SERMON < < <}}
~
\hyperref[sec:index_chronic]{\small{[返順時目]}}
~
\hyperref[sec:031741]{\small{> > > NEXT SERMON > > >}}
\newline
\newline
撒母耳記上 18:1-21:15
\newline
\begin{longtable}{cl}
\hline
\hline
章節 & 經文 (和合本修訂版)\\
\hline
18:1 & \begin{tabularx}{0.7\textwidth}{X} 大衛對掃羅說完了話，約拿單的心與大衛的心深相契合。約拿單愛大衛，如同愛自己的性命。 \end{tabularx} \\ \\ \relax
18:2 & \begin{tabularx}{0.7\textwidth}{X} 那日掃羅留住大衛，不讓他回父家。 \end{tabularx} \\ \\ \relax
18:3 & \begin{tabularx}{0.7\textwidth}{X} 約拿單愛大衛如同愛自己的性命，就與他立約。 \end{tabularx} \\ \\ \relax
18:4 & \begin{tabularx}{0.7\textwidth}{X} 約拿單從身上脫下外袍，給了大衛，又把戰衣、刀、弓、腰帶都給了他。 \end{tabularx} \\ \\ \relax
18:5 & \begin{tabularx}{0.7\textwidth}{X} 掃羅無論差遣大衛往何處去，他都做事精明。掃羅立他作軍隊的指揮官，眾百姓和掃羅的臣僕都看為美。 \end{tabularx} \\ \\ \relax
18:6 & \begin{tabularx}{0.7\textwidth}{X} 大衛打死了那非利士人，同眾人回來的時候，婦女們從以色列各城裡出來，歡歡喜喜，打鼓奏樂，唱歌跳舞，迎接掃羅王。 \end{tabularx} \\ \\ \relax
18:7 & \begin{tabularx}{0.7\textwidth}{X} 眾婦女歡樂唱和，說：「掃羅殺死千千，大衛殺死萬萬。」 \end{tabularx} \\ \\ \relax
18:8 & \begin{tabularx}{0.7\textwidth}{X} 掃羅非常憤怒，不喜歡這話。他說：「將萬萬歸給大衛，千千歸給我，只剩下王國沒有給他！」 \end{tabularx} \\ \\ \relax
18:9 & \begin{tabularx}{0.7\textwidth}{X} 從這日起，掃羅就敵視大衛。 \end{tabularx} \\ \\ \relax
18:10 & \begin{tabularx}{0.7\textwidth}{X} 次日，從神來的邪靈緊抓住掃羅，他就在家中胡言亂語。大衛照常彈琴，掃羅手裡拿著槍。 \end{tabularx} \\ \\ \relax
18:11 & \begin{tabularx}{0.7\textwidth}{X} 掃羅把槍一擲，心裡說：「我要將大衛刺透，釘在牆上。」大衛閃避了他兩次。 \end{tabularx} \\ \\ \relax
18:12 & \begin{tabularx}{0.7\textwidth}{X} 掃羅懼怕大衛，因為耶和華離開自己，與大衛同在。 \end{tabularx} \\ \\ \relax
18:13 & \begin{tabularx}{0.7\textwidth}{X} 所以掃羅叫大衛離開自己，立他為千夫長，他就領兵出入。 \end{tabularx} \\ \\ \relax
18:14 & \begin{tabularx}{0.7\textwidth}{X} 大衛所做的每一件事都精明，耶和華也與他同在。 \end{tabularx} \\ \\ \relax
18:15 & \begin{tabularx}{0.7\textwidth}{X} 掃羅見大衛做事精明，就更怕他。 \end{tabularx} \\ \\ \relax
18:16 & \begin{tabularx}{0.7\textwidth}{X} 但以色列和猶大眾人都愛大衛，因為他領他們出入。 \end{tabularx} \\ \\ \relax
18:17 & \begin{tabularx}{0.7\textwidth}{X} 掃羅對大衛說：「看哪，我將大女兒米拉嫁給你，只要你作我的勇士，為耶和華爭戰。」掃羅心裡說：「我不好親手害他，要藉非利士人的手害他。」 \end{tabularx} \\ \\ \relax
18:18 & \begin{tabularx}{0.7\textwidth}{X} 大衛對掃羅說：「我是誰，我是甚麼出身，我父家在以色列中算甚麼，豈敢作王的女婿呢？」 \end{tabularx} \\ \\ \relax
18:19 & \begin{tabularx}{0.7\textwidth}{X} 掃羅的女兒米拉到了當嫁給大衛的時候，掃羅卻將她嫁給了米何拉人亞得列。 \end{tabularx} \\ \\ \relax
18:20 & \begin{tabularx}{0.7\textwidth}{X} 掃羅的女兒米甲愛大衛。有人告訴掃羅，這件事在掃羅眼中看為合宜。 \end{tabularx} \\ \\ \relax
18:21 & \begin{tabularx}{0.7\textwidth}{X} 掃羅心裡說：「我把這女兒嫁給大衛，作他的圈套，好藉非利士人的手害他。」所以掃羅第二次對大衛說：「你今日可以作我的女婿。」 \end{tabularx} \\ \\ \relax
18:22 & \begin{tabularx}{0.7\textwidth}{X} 掃羅吩咐臣僕：「你們暗中對大衛說：『看哪，王喜歡你，王的臣僕也都愛戴你，現在你就作王的女婿吧。』」 \end{tabularx} \\ \\ \relax
18:23 & \begin{tabularx}{0.7\textwidth}{X} 掃羅的臣僕照這話說給大衛聽。大衛說：「你們把作王的女婿看為小事嗎？我是貧窮卑微的人。」 \end{tabularx} \\ \\ \relax
18:24 & \begin{tabularx}{0.7\textwidth}{X} 掃羅的臣僕回奏說，大衛說了這樣的話。 \end{tabularx} \\ \\ \relax
18:25 & \begin{tabularx}{0.7\textwidth}{X} 掃羅說：「你們要對大衛這樣說：『王不要甚麼聘禮，只要一百非利士人的包皮，好在王的仇敵身上報仇。』」掃羅的意圖是要大衛落在非利士人的手中。 \end{tabularx} \\ \\ \relax
18:26 & \begin{tabularx}{0.7\textwidth}{X} 掃羅的臣僕把這話告訴大衛，大衛就歡喜作王的女婿。日期還沒有到， \end{tabularx} \\ \\ \relax
18:27 & \begin{tabularx}{0.7\textwidth}{X} 大衛和跟隨他的人起來前往，殺了二百非利士人，將包皮足數交給王，為要作王的女婿。於是掃羅將女兒米甲嫁給大衛。 \end{tabularx} \\ \\ \relax
18:28 & \begin{tabularx}{0.7\textwidth}{X} 掃羅見耶和華與大衛同在，女兒米甲又愛大衛， \end{tabularx} \\ \\ \relax
18:29 & \begin{tabularx}{0.7\textwidth}{X} 就更怕大衛，常常與大衛為敵。 \end{tabularx} \\ \\ \relax
18:30 & \begin{tabularx}{0.7\textwidth}{X} 每逢非利士的軍官出來打仗，大衛做事比掃羅任何臣僕更精明，因此他的名極受尊重。 \end{tabularx} \\ \\ \relax
19:1 & \begin{tabularx}{0.7\textwidth}{X} 掃羅吩咐他兒子約拿單和眾臣僕要殺大衛，但掃羅的兒子約拿單卻很喜愛大衛。 \end{tabularx} \\ \\ \relax
19:2 & \begin{tabularx}{0.7\textwidth}{X} 約拿單告訴大衛說：「我父掃羅想要殺你，現在你要小心，明日早晨留在一個僻靜的地方藏起來。 \end{tabularx} \\ \\ \relax
19:3 & \begin{tabularx}{0.7\textwidth}{X} 我會出去，到你所藏的田裡，站在我父親旁邊，與父親談論到你。我看情形怎樣，會告訴你。」 \end{tabularx} \\ \\ \relax
19:4 & \begin{tabularx}{0.7\textwidth}{X} 約拿單向他父親掃羅說大衛的好話，對他說：「王不可得罪王的僕人大衛，因為他未曾得罪你，他所行的對你都很有益處。 \end{tabularx} \\ \\ \relax
19:5 & \begin{tabularx}{0.7\textwidth}{X} 他拚了命殺那非利士人，並且耶和華為全以色列大施拯救。那時你看見，也很歡喜，現在為何要犯罪，流無辜人的血，無緣無故殺大衛呢？」 \end{tabularx} \\ \\ \relax
19:6 & \begin{tabularx}{0.7\textwidth}{X} 掃羅聽了約拿單的話，就指著永生的耶和華起誓：「我絕不殺他。」 \end{tabularx} \\ \\ \relax
19:7 & \begin{tabularx}{0.7\textwidth}{X} 約拿單叫大衛來，把這一切事告訴他。約拿單帶他去見掃羅，他就像以前一樣侍立在掃羅面前。 \end{tabularx} \\ \\ \relax
19:8 & \begin{tabularx}{0.7\textwidth}{X} 此後又有戰爭，大衛出去與非利士人打仗。他大大擊敗他們，他們就在他面前逃跑。 \end{tabularx} \\ \\ \relax
19:9 & \begin{tabularx}{0.7\textwidth}{X} 從耶和華來的邪靈又降在掃羅身上，掃羅手裡拿槍坐在屋裡，大衛正用手彈琴。 \end{tabularx} \\ \\ \relax
19:10 & \begin{tabularx}{0.7\textwidth}{X} 掃羅想要用槍刺透大衛，把他釘在牆上，他卻躲開掃羅，掃羅的槍刺入牆內。當夜大衛逃走，躲起來了。 \end{tabularx} \\ \\ \relax
19:11 & \begin{tabularx}{0.7\textwidth}{X} 掃羅派一些使者到大衛的房屋那裡守著他，等到天亮要殺他。大衛的妻子米甲對大衛說：「你今夜若不逃命，明日就要被殺。」 \end{tabularx} \\ \\ \relax
19:12 & \begin{tabularx}{0.7\textwidth}{X} 於是米甲將大衛從窗戶縋下去，讓他走；大衛就逃走，躲起來了。 \end{tabularx} \\ \\ \relax
19:13 & \begin{tabularx}{0.7\textwidth}{X} 米甲把家中的神像放在床上，頭枕在山羊毛的枕頭上，用衣服蓋起來。 \end{tabularx} \\ \\ \relax
19:14 & \begin{tabularx}{0.7\textwidth}{X} 掃羅派一些使者去捉拿大衛，米甲說：「他病了。」 \end{tabularx} \\ \\ \relax
19:15 & \begin{tabularx}{0.7\textwidth}{X} 掃羅又派一些使者去看大衛，說：「把他連床一起抬到我這裡，我好殺他。」 \end{tabularx} \\ \\ \relax
19:16 & \begin{tabularx}{0.7\textwidth}{X} 使者進去，看哪，神像在床上，頭枕在山羊毛的枕頭上。 \end{tabularx} \\ \\ \relax
19:17 & \begin{tabularx}{0.7\textwidth}{X} 掃羅對米甲說：「你為甚麼這樣欺騙我，放我仇敵逃走呢？」米甲對掃羅說：「他對我說：『你放我走吧，我何必要殺你呢？』」 \end{tabularx} \\ \\ \relax
19:18 & \begin{tabularx}{0.7\textwidth}{X} 大衛逃跑躲避，來到拉瑪的撒母耳那裡，把掃羅向他所行的事全告訴他。他和撒母耳就去，住在拿約。 \end{tabularx} \\ \\ \relax
19:19 & \begin{tabularx}{0.7\textwidth}{X} 有人告訴掃羅說：「看哪，大衛在拉瑪的拿約」。 \end{tabularx} \\ \\ \relax
19:20 & \begin{tabularx}{0.7\textwidth}{X} 掃羅派一些使者去捉拿大衛。去的人見一隊先知受感說話，撒母耳站在當中領導他們，掃羅派去的使者也受神的靈感動說話。 \end{tabularx} \\ \\ \relax
19:21 & \begin{tabularx}{0.7\textwidth}{X} 有人把這事告訴掃羅，他又派另一些使者去，他們也受感說話。掃羅第三次派使者去，他們也受感說話。 \end{tabularx} \\ \\ \relax
19:22 & \begin{tabularx}{0.7\textwidth}{X} 然後掃羅親自往拉瑪去，到了西沽的大井，問人說：「撒母耳和大衛在哪裡？」有人說：「看哪，在拉瑪的拿約。」 \end{tabularx} \\ \\ \relax
19:23 & \begin{tabularx}{0.7\textwidth}{X} 他就往那裡去，到了拉瑪的拿約。神的靈也臨到他，他一面走一面受感說話，直到拉瑪的拿約。 \end{tabularx} \\ \\ \relax
19:24 & \begin{tabularx}{0.7\textwidth}{X} 他也脫了衣服，也在撒母耳面前受感說話，一日一夜赤身躺臥。因此有人說：「掃羅也在先知中嗎？」 \end{tabularx} \\ \\ \relax
20:1 & \begin{tabularx}{0.7\textwidth}{X} 大衛從拉瑪的拿約逃跑，來到約拿單面前，對他說：「我做了甚麼，有甚麼罪孽，在你父親面前犯了甚麼罪，他竟要尋索我的性命呢？」 \end{tabularx} \\ \\ \relax
20:2 & \begin{tabularx}{0.7\textwidth}{X} 約拿單對他說：「絕無此事！你必不至於死。看哪，我父做事，無論大小，沒有不告訴我的。我父親為甚麼要隱瞞我這件事呢？不會這樣的！」 \end{tabularx} \\ \\ \relax
20:3 & \begin{tabularx}{0.7\textwidth}{X} 大衛又起誓說：「你父親確實知道我在你眼前蒙恩。所以他說，『這事不要讓約拿單知道，免得他愁煩。』我指著永生的耶和華起誓，又指著你的性命起誓，我離死只差一步而已。」 \end{tabularx} \\ \\ \relax
20:4 & \begin{tabularx}{0.7\textwidth}{X} 約拿單對大衛說：「你心裡所求的，我必為你成就。」 \end{tabularx} \\ \\ \relax
20:5 & \begin{tabularx}{0.7\textwidth}{X} 大衛對約拿單說：「看哪，明日是初一，我必須與王同席用餐，求你讓我去藏在田野，直到第三日傍晚。 \end{tabularx} \\ \\ \relax
20:6 & \begin{tabularx}{0.7\textwidth}{X} 你父親若見我不在席上，你就說：『大衛懇求我允許他趕回本城伯利恆去，因為他全家在那裡獻年祭。』 \end{tabularx} \\ \\ \relax
20:7 & \begin{tabularx}{0.7\textwidth}{X} 你父親若說好，你的僕人就平安了；他若大怒，你就知道他決意行惡。 \end{tabularx} \\ \\ \relax
20:8 & \begin{tabularx}{0.7\textwidth}{X} 求你施恩於僕人，因你在耶和華面前曾與僕人立約。我若有罪孽，你就親自殺死我，何必把我交給你父親呢？」 \end{tabularx} \\ \\ \relax
20:9 & \begin{tabularx}{0.7\textwidth}{X} 約拿單說：「絕無此事！我若確實知道我父親決意害你，怎麼會不告訴你呢？」 \end{tabularx} \\ \\ \relax
20:10 & \begin{tabularx}{0.7\textwidth}{X} 大衛對約拿單說：「你父親若嚴厲回答你，誰來告訴我呢？」 \end{tabularx} \\ \\ \relax
20:11 & \begin{tabularx}{0.7\textwidth}{X} 約拿單對大衛說：「來，讓我們到田野去。」二人就往田野去了。 \end{tabularx} \\ \\ \relax
20:12 & \begin{tabularx}{0.7\textwidth}{X} 約拿單對大衛說：「願耶和華－以色列的神作證。明日約在這時候，或第三日，我一探出我父親的心意，看哪，若對大衛是好意，我怎麼會不派人來告訴你呢？ \end{tabularx} \\ \\ \relax
20:13 & \begin{tabularx}{0.7\textwidth}{X} 我父親若有意害你，而我不告訴你，送你平安地離開，願耶和華重重懲罰約拿單。願耶和華與你同在，如同從前與我父親同在一樣。 \end{tabularx} \\ \\ \relax
20:14 & \begin{tabularx}{0.7\textwidth}{X} 你要照耶和華的慈愛恩待我，不但我活著的時候免我死亡， \end{tabularx} \\ \\ \relax
20:15 & \begin{tabularx}{0.7\textwidth}{X} 就是耶和華從地面逐一剪除大衛仇敵的時候，你也永不可向我家斷絕恩惠。」 \end{tabularx} \\ \\ \relax
20:16 & \begin{tabularx}{0.7\textwidth}{X} 於是約拿單與大衛家立約：「願耶和華從大衛仇敵的手來追討。」 \end{tabularx} \\ \\ \relax
20:17 & \begin{tabularx}{0.7\textwidth}{X} 約拿單因愛大衛如同愛自己的性命，就叫他再起誓。 \end{tabularx} \\ \\ \relax
20:18 & \begin{tabularx}{0.7\textwidth}{X} 約拿單對他說：「明日是初一，你的座位空著，人必察覺你不在。 \end{tabularx} \\ \\ \relax
20:19 & \begin{tabularx}{0.7\textwidth}{X} 到第三日，就要走一段長路下去，去到你遇事那天所藏的地方，在以色磐石的旁邊等候。 \end{tabularx} \\ \\ \relax
20:20 & \begin{tabularx}{0.7\textwidth}{X} 我會向磐石旁邊射三箭，如同射箭靶一樣。 \end{tabularx} \\ \\ \relax
20:21 & \begin{tabularx}{0.7\textwidth}{X} 看哪，我會派僮僕，說：『去把箭找來。』我若對僮僕喊說：『看哪，箭在你的這邊，把箭拿來』，你就可以平安回來；我指著永生的耶和華起誓，你一定沒有事。 \end{tabularx} \\ \\ \relax
20:22 & \begin{tabularx}{0.7\textwidth}{X} 我若對孩子說：『看哪，箭在你的前方』，你就要離開，因為是耶和華差你去的。 \end{tabularx} \\ \\ \relax
20:23 & \begin{tabularx}{0.7\textwidth}{X} 至於你和我，我們所說的話，看哪，耶和華在你我中間作證，直到永遠。」 \end{tabularx} \\ \\ \relax
20:24 & \begin{tabularx}{0.7\textwidth}{X} 大衛就去藏在田野。到了初一，王要坐席用餐。 \end{tabularx} \\ \\ \relax
20:25 & \begin{tabularx}{0.7\textwidth}{X} 王照常坐在靠牆的位子上，約拿單在對面，押尼珥坐在掃羅旁邊，大衛的座位卻是空的。 \end{tabularx} \\ \\ \relax
20:26 & \begin{tabularx}{0.7\textwidth}{X} 這日掃羅沒有說甚麼，因為他說：「大衛或許有事，偶染不潔，還未得潔淨。」 \end{tabularx} \\ \\ \relax
20:27 & \begin{tabularx}{0.7\textwidth}{X} 初二，大衛的座位還空著。掃羅對他兒子約拿單說：「耶西的兒子為何昨日、今日都沒有來用餐呢？」 \end{tabularx} \\ \\ \relax
20:28 & \begin{tabularx}{0.7\textwidth}{X} 約拿單回答掃羅說：「大衛懇求我允許他回伯利恆去， \end{tabularx} \\ \\ \relax
20:29 & \begin{tabularx}{0.7\textwidth}{X} 說：『求你讓我去，因為我家在城裡有獻祭的事，我哥哥吩咐我去。如今我若在你眼前蒙恩，求你讓我去見我的兄弟。』所以大衛沒有來赴王的筵席。」 \end{tabularx} \\ \\ \relax
20:30 & \begin{tabularx}{0.7\textwidth}{X} 掃羅向約拿單怒氣大發，對他說：「你這頑梗悖逆之婦人所生的，我怎麼會不知道你選擇耶西的兒子，自取羞辱，也使你母親露體蒙羞呢？ \end{tabularx} \\ \\ \relax
20:31 & \begin{tabularx}{0.7\textwidth}{X} 只要耶西的兒子還活在世上一天，你和你的國必保不住。現在你要派人去，把他帶到我這裡來，因為他是該死的。」 \end{tabularx} \\ \\ \relax
20:32 & \begin{tabularx}{0.7\textwidth}{X} 約拿單回答父親掃羅說：「他為甚麼該死呢？他做了甚麼呢？」 \end{tabularx} \\ \\ \relax
20:33 & \begin{tabularx}{0.7\textwidth}{X} 掃羅向約拿單擲槍要刺他，約拿單就知道他父親決意要殺死大衛。 \end{tabularx} \\ \\ \relax
20:34 & \begin{tabularx}{0.7\textwidth}{X} 於是約拿單氣憤憤地從席上起來。他在初二這天沒有吃飯，因為他為大衛愁煩，又因為他父親羞辱了他。 \end{tabularx} \\ \\ \relax
20:35 & \begin{tabularx}{0.7\textwidth}{X} 次日早晨，約拿單按著與大衛約定的時候到田野去，有一個小僮僕跟隨他。 \end{tabularx} \\ \\ \relax
20:36 & \begin{tabularx}{0.7\textwidth}{X} 約拿單對僮僕說：「你跑去把我所射的箭找來。」僮僕跑去，約拿單就把箭射在僮僕的前方。 \end{tabularx} \\ \\ \relax
20:37 & \begin{tabularx}{0.7\textwidth}{X} 僮僕到了約拿單落箭之地，約拿單呼叫僮僕說：「箭不是在你的前方嗎？」 \end{tabularx} \\ \\ \relax
20:38 & \begin{tabularx}{0.7\textwidth}{X} 約拿單又呼叫僮僕說：「快去，不要站在那裡！」僮僕就撿起箭來，回到主人那裡。 \end{tabularx} \\ \\ \relax
20:39 & \begin{tabularx}{0.7\textwidth}{X} 僮僕不知道這是甚麼意思，只有約拿單和大衛知道這事。 \end{tabularx} \\ \\ \relax
20:40 & \begin{tabularx}{0.7\textwidth}{X} 約拿單把他的弓箭交給僮僕，吩咐他說：「你拿到城裡去。」 \end{tabularx} \\ \\ \relax
20:41 & \begin{tabularx}{0.7\textwidth}{X} 僮僕一去，大衛就從南邊出來，俯伏在地，拜了三拜。他們彼此親吻，一起哭泣，大衛哭得更悲哀。 \end{tabularx} \\ \\ \relax
20:42 & \begin{tabularx}{0.7\textwidth}{X} 約拿單對大衛說：「你平平安安地去吧！因為我們二人曾指著耶和華的名起誓說：『願耶和華在你我中間，以及你我後裔中間作證，直到永遠。』」大衛就起身走了，約拿單也回城裡去了。 \end{tabularx} \\ \\ \relax
21:1 & \begin{tabularx}{0.7\textwidth}{X} 大衛到了挪伯的亞希米勒祭司那裡，亞希米勒戰戰兢兢地出來迎接他，對他說：「你為甚麼獨自一人，沒有人跟隨你呢？」 \end{tabularx} \\ \\ \relax
21:2 & \begin{tabularx}{0.7\textwidth}{X} 大衛對亞希米勒祭司說：「王吩咐我一件事，對我說：『我差遣你，吩咐你的這件事，不可讓任何人知道。』因此我已告訴一些僕人到某處去。 \end{tabularx} \\ \\ \relax
21:3 & \begin{tabularx}{0.7\textwidth}{X} 現在你手中有甚麼？請你給我五個餅或是可以找到的食物。」 \end{tabularx} \\ \\ \relax
21:4 & \begin{tabularx}{0.7\textwidth}{X} 祭司對大衛說：「我手中沒有普通的餅，只有聖餅，只能給沒有親近婦人的年輕人。」 \end{tabularx} \\ \\ \relax
21:5 & \begin{tabularx}{0.7\textwidth}{X} 大衛回答祭司說：「我們確實沒有親近婦人，如同往常我出征的時候一樣。平常行路，僕人的身體都還分別為聖，何況今日豈不更使自己分別為聖嗎？」 \end{tabularx} \\ \\ \relax
21:6 & \begin{tabularx}{0.7\textwidth}{X} 祭司拿聖餅給他，因為在那裡沒有別的餅，只有那從耶和華面前撤下的供餅，就是換上熱餅的日子取下來的。 \end{tabularx} \\ \\ \relax
21:7 & \begin{tabularx}{0.7\textwidth}{X} 當日，有掃羅的一個臣僕在那裡，他留在耶和華的面前，名叫多益，是以東人，作掃羅的畜牧長。 \end{tabularx} \\ \\ \relax
21:8 & \begin{tabularx}{0.7\textwidth}{X} 大衛對亞希米勒說：「你手中有沒有槍或刀？因為王的事緊急，連刀劍兵器我都沒有帶。」 \end{tabularx} \\ \\ \relax
21:9 & \begin{tabularx}{0.7\textwidth}{X} 祭司說：「你在以拉谷所殺的非利士人歌利亞的那刀，看哪，裹在布中，放在以弗得後邊。你若要可以拿去，除此以外，再沒有別的了。」大衛說：「沒有甚麼可以跟它比的了！請你給我。」 \end{tabularx} \\ \\ \relax
21:10 & \begin{tabularx}{0.7\textwidth}{X} 那日大衛起來，躲避掃羅，逃到迦特王亞吉那裡。 \end{tabularx} \\ \\ \relax
21:11 & \begin{tabularx}{0.7\textwidth}{X} 亞吉的臣僕對他說：「這不是那地的國王大衛嗎？那裡的人跳舞唱和：『掃羅殺死千千，大衛殺死萬萬』，不是指著他說的嗎？」 \end{tabularx} \\ \\ \relax
21:12 & \begin{tabularx}{0.7\textwidth}{X} 大衛把這些話放在心裡，就很懼怕迦特王亞吉。 \end{tabularx} \\ \\ \relax
21:13 & \begin{tabularx}{0.7\textwidth}{X} 於是他在眾人眼前一反常態，在他們中間裝瘋作癲，在城門的門扇上胡寫亂畫，任由唾沫流在鬍子上。 \end{tabularx} \\ \\ \relax
21:14 & \begin{tabularx}{0.7\textwidth}{X} 亞吉對臣僕說：「看哪，你們看這人瘋了，為甚麼帶他到我這裡來呢？ \end{tabularx} \\ \\ \relax
21:15 & \begin{tabularx}{0.7\textwidth}{X} 我豈缺少瘋子，你們竟然帶這人到我面前瘋癲嗎？這個人可以進我的家嗎？」 \end{tabularx} \\ \\
[1ex]
\hline
\hline
\end{longtable}
$^{1}$我們一連串思想薩姆爾記上的訊息.
特別提到大衛人生一些借鏡的地方.
在薩姆爾記上18-19章開始.
一直到整本書的結束.
其實是說到大衛在這段時間有極大磨難的日子.
因為他要逃避訴羅的追擊.
但其實不單止他有逃難的苦.
在這個過程中似乎他經歷了一種又一種失去.
甚至有時失去了方向.
更具體的是他失去了很多他曾經擁有的東西.
我們人生也有失落或一定有得有失.
可能我們也會有失去的時候.
失去一些本來我們擁有的人或物.
但我們知道這是人生有的過程.
與此同時我們見到大衛在這段磨難的日子.
很多失去的時候.
他寫了很多詩篇.
在這個時候他經歷神的同在.
是一個很明顯的階段.
他在靈性上過往已經紮好一些根基.
在少年的時候.
不過在這段時間可以說是紮根更深的時候.
因為他再不能依靠其他他擁有的人或物.
他發現他最深的倚靠是上帝自己.
所以他在這段時間寫了很多與神的關係.
如何感受到所有一切都沒有的時候.
上帝在他的身邊.
所以今天我們默許擁有很多東西.
希望求神憐憫.
不要讓我們經歷像大衛一樣的打擊和損失.
但是在這個時候.
當我們仍然有的時候.
或者我們都會經歷到失去的時候.
我們知道最終的倚靠是誰.
所以讓我們今早從三文治記載十八章.
至到二十一章這幾章經文裡面.
將我們其中一些經節看回當時大衛所面對的.
反省我們究竟甚麼是我們最終的倚靠.
第一方面我們今天要看的就是.
權力不是最終的倚靠.

$^{41}$為何這樣說呢.
在三文治記載十八章第五節.
講到當時大衛打完哥利亞.
也有約拿丹來支持鼓勵和讚賞他.
我們先看一連串的經歷.
我們先看十八章第五節.
我們一起讀這節經文.
蘇羅無論差遣大衛往何處去.
他都做事精明.
蘇羅就納他作戰士長.
眾百姓和蘇羅的臣僕無不喜悅.
好不好啊.
真是很好啊.
當時他受眾人認可.
蘇羅王知道他得勝了非利士人.
這是一個很大的功勞.
更加當他差遣大衛做事的時候.
大衛做事精明.
蘇羅非常賞識他.
所以一下子納他做戰士長.
他本來是沒有經驗的.
只是一個少年的木桶.
但一下子竟然身高要職.
成為了戰士當中的為首.
配不配啊.
配啊 因為他打贏了最大的哥尼亞.
沒有人敢打 也沒有人能打贏.
但只有他能勝過巨人.
所以每個人都服他.
眾百姓和蘇羅的臣僕無不喜悅.
這才重要.
因為他不只是蘇羅給了他一個高位.
而是其他人都很服他.
其他人都覺得這個真是好.
有他做我們的領導.
有他做我們的戰士長.
真好.
不過一轉轉到十九章.
我們看看十九章.
由第八到第十節.

$^{81}$我們一起讀 好不好?請讀.
此後又有爭戰的事.
大衛出去與菲律賓人打仗.
大大殺敗他們.
他們就在他面前逃跑.
從耶和華那裡來的惡魔.
又降在蘇羅身上.
蘇羅手裡拿槍坐在屋裡.
大衛就用手彈琴.
蘇羅用槍想要刺殺大衛.
釘在牆上.
他卻躲開.
蘇羅的槍刺入牆內.
當夜大衛逃走.
躲避了.
一下之間.
一個大的轉變.
真的扭轉了整個局面.
本來蘇羅和所有人都非常欣賞他.
他被升做戰士長.
有很好的權位.
有很好的地位.
有很好的在物質上和名利上的一大得益.
很可惜.
這個情況一下子又轉變.
不是大衛無能.
而是大衛形形的殺敗很多菲律師人.
本來應該更加再晉升才對.
但很可惜.
因為蘇羅沒有上帝在心中的人.
他被神棄絕之後.
惡魔進入他心.
結果他就大開殺戒.
窒道在心.
企圖將大衛殺死.
本來這麼重用他.
一下子變成這樣.
因為他那種罪惡的心性.
這種窒道和他已經離開了上帝.
沒有神的靈的管理.

$^{121}$反而有惡魔在他心中充滿他.
所以一下子讓大衛生命.
落入一個非常危險的情況.
結果大衛就唯有放棄.
逃走.
躲避.
撒梅記二十一章我們跳一跳.
我們看到大衛逃難的時候.
究竟他怎樣呢?.
這裡更諷刺.
二十一章第十.
十一節請讀.
那一日大衛起來.
躲避蘇羅.
逃到迦特王雅吉那裡.
雅吉的神僕對雅吉說.
「這不是以色列王大衛嗎?」.
雅吉的婦女跳舞唱和不是指著她說.
「蘇羅殺死千千.
大衛殺死萬萬嗎?」.
有人幫幫他.
不懂弄個手機.
看手機電話時有時會弄出聲音.
我們留意一下.
我們看到大衛逃避蘇羅.
去了哪裡?.
迦特在哪裡來?.
知道嗎?.
迦特是菲律賓五個城之中其中一個.
所以在菲律賓人這五個重要城裡.
其中一個有分封的王.
這個王叫甚麼名字?.
雅吉.
我們有沒有人叫雅吉?.
這個雅吉收納大衛.
不過是奇妙的.
原來菲律賓人都知道.
知道甚麼?.
知道現在做王的不是真正做王的.
因為神已經納了誰做王?.

$^{161}$納了大衛做以色列國的王.
所以其實蘇羅真的那麼憤怒.
連約拿丹都說將來你做王我要做你的宰相.
所以其實這個已經知道了.
通天的.
即是說上帝揀選了一個新的王將會代替蘇羅.
結果可以見得蘇羅的憤怒就在這裡.
所以連菲律賓人都知道.
結果他們說是囉 當時民眾都這樣說.
蘇羅就殺死千千 大衛就殺死萬萬.
請問殺死誰?.
(菲律賓人).
殺死我們這些人 確實是.
殺死菲律賓人.
所以其實我們見到大衛都出了名.
不單止在國內甚至國外.
不單止自己的軍人甚至敵軍都知道.
所以這個才厲害.
他就是一位戰士長 蘇羅所重用的.
甚至將來是以色列國 即是我們當時的敵國的皇帝.
從那麼高的地位 那麼高的權勢.
那麼高的聲譽 那麼高的所謂美好的前途.
他一下子成為了樂方而逃的人.
這真是一個很大的損失.
所以當時連那些菲律賓人都說.
阿吉王你知不知道這是誰?.
瘋瘋癲癲.
你看看這個以前說要殺死我們萬萬的那個.
現在來投靠我 投靠我們.
所以沒什麼好做的 做到要來敵軍投靠.
真是很慘.
這種是很大的失落.
不知道大家有沒有曾經試過這種失落.
或者可能以前在某個國家.
或者香港 或者越南 或者某些地方.
曾經做到很高位的移民來的時候要從低做起.
又或者可能本來做得很好的 很高職位的 很有前途的.
一下子被人lay off 就沒了一份職業.
現在我們昨天在幸福小組的電影節目聊天.
現在很多人都害怕究竟那份職業會不會保得住.

$^{201}$因為什麼呢?.
你們都知道吧.
連我們那些七八十歲八十九十歲都知道AI.
是的 真的不知道.
最近我的女兒都跟我說.
她說我們不知道我們county那邊的工會怎樣轉變.
有可能很多東西被AI代替了.
現在剛剛出來的大學畢業的更加徬徨.
有很多事找工作真的很難找.
因為她本來讀過科.
很好的 很有前途的.
現在一下子AI可以做到.
不用她 機械人搞定.
這是很大的打擊.
所以我們知道人生都是充滿著這些忽然間而來的改變.
甚至失落.
有很高的權位 很高的前途.
可能一下子轉變成為一無所有.
其實現在在pandemic之後也有很多轉變.
有很多的公司 小的small business.
也都關了很多門 我們都知道.
在美國申請破產的人是很高的.
上百萬計.
申請破產不是他沒了錢.
然後因為欠下太多債 無路可走.
唯有申請破產.
所以我們看到人生裡 整個世界裡.
我們看到真的有很多時候.
忽然之間在很高的權位可以失去.
在聖經裡約伯也這麼說.
約伯經歷過這個最深.
聖經約伯記三十章十五節怎麼說呢.
我們一起讀讀.
「驚恐臨到我 驅逐我的尊榮如風.
我的福祿如雲過去」.
我們最希望福祿受傳.
聖經也有提福祿.
不過福祿也可以如雲過去.
你是不是叫雲.
雲一下子就過去了.

$^{241}$無影無蹤.
好像很實在 忽然之間沒了.
「尊榮如風 風一下子來一下子就找不著.
風和雲轉眼消失 不能可靠」.
約伯深深體會到這一點.
所以尊榮和世界的權位也如此.
究竟有什麼我們可以抓緊的呢.
其實上帝不是反對這些.
有權勢 有地位 有學問 有前途.
神希望我們有.
我們祈禱也希望我們的子弟.
我們自己蒙這個祝福.
這個無可厚非.
不過其實上帝賜這些東西給我們.
不是叫我們成為我們人生的依求.
是叫我們善用祂來祝福人群.
即是說神給我們機會受教育.
有好的職業.
能夠有好的前途.
在這個世界上有影響力.
是為了去貢獻 去祝福.
去善用神給我們有的.
來服侍社會 服侍人群.
服侍我們周圍的人.
權位是有力的.
但不是濫權的.
善用權力就可以帶來祝福.
濫用權力就帶來災害.
剛好我們昨天大家都知道.
伊朗最高的獨裁者 跨台.
他在統治裡面有三四十年.
有很多很多的獨裁 獨權.
害了千萬人的靈魂.
有很多自己的國民被屠殺.
本來權勢可以去建立別人.
可以去服侍社會 服侍人民.
但如果濫用權力.
就會造成一個極大的痛苦災害.
不單是自己 國家 人民.
所以其實我們知道權力.

$^{281}$不是最終神要我們去依靠的.
我們有權 可能有好的地位.
好的前途.
這些都是應該放在神那裡給他使用.
在詩篇十六篇 大衛的金詩這樣說.
我的心啊 你曾對耶和華說.
你是我的主.
我的好處不在你以外.
以別的代替耶和華.
他們的仇苦必加增.
這真是大衛深入反省之後.
有一個自細明言.
我的心啊 你曾對耶和華說.
你是我的主.
我的好處不在你以外.
所以真是發現到其他所有的一切.
如果離開了上帝 在上帝以外的.
以別的代替了耶和華的.
他們的仇苦必加增.
這真是非常深刻的體會.
相反是同一篇詩篇.
如果大家回去看詩篇十六篇 金詩.
它說如果我將耶和華上擺在我的面前.
我的心就歡喜.
我的靈就快樂.
我的肉身就安然居住.
所以這是一種滿足 一種平安.
不再需要憂慮.
因為權勢之間會轉變.
但如果有這位賜給我們權勢善用一切的.
這位上帝在我們身邊.
雖然有時會失去.
好像約伯一樣.
神容許他失去.
但他真正的倚靠 最後發現.
真正的祝福 他發現.
都是在神自己那裡.
所以親愛的弟兄姊妹.
今天讓我們自己再反省.
希望我們不需要有這麼多失去.

$^{321}$不過我們現在可以有個聰明的調整.
究竟什麼是我們最終的倚靠.
是我們的權勢 是我們的地位.
是我們所擁有的物質.
還是這位賜給我們有這些善用的上帝.
Amen.
什麼是我們最終的倚靠.
我們發覺配偶不是最終的倚靠.
這裡就講到大衛的妻 叫什麼名字.
米鴿 她知道她的爸爸要殺大衛.
她就通水給大衛 大衛就逃走了.
讀完17節 17節請讀.
蘇羅對米鴿說.
你為什麼這樣欺負我.
放我囚敵逃走呢.
米鴿回答說.
他對我說.
你放我走 不然我要殺你.
是不是真的.
不是真的.
只是一個藉口.
所以米鴿從此就和大衛分開.
米鴿是大衛用性命換回來的妻子.
大家記得上次我們第一個女兒.
偶爾蘇羅反悔.
然後米鴿就愛上大衛.
說他少年英雄.
很羨慕他.
蘇羅就說.
你幫我殺菲律賓人.
拿100個羊皮回來.
結果大衛無死.
拿了200個.
所以他真的很勇敢.
圍著爭取這個和米鴿結婚的機會.
在《聖經》裡面說.
如果大家回去看我寫的書裡面有分析過.
就是大衛心裡不斷提及.
我要做王的女婿.
這就是他很大的動機.

$^{361}$所以做王的女婿.
是大衛很希望得到的.
所以他拼命.
冒生命的危險.
換取了妻子.
不過可惜我們發覺.
事不願為.
因為情況改變.
事實上他和米鴿之間也沒有一個很好的基礎.
因為大家不太了解對方.
更加重要的是他們和神的關係有很大的差距.
大衛打哥利亞的時候.
其實表面看.
米鴿看出來就以為他如何英勇.
如何厲害.
大衛的看法就不是.
大衛的看法就是知道是耶和華上帝給他的力量.
所以我們看到整個描述.
後來再記米鴿的事跡也好.
我們發覺大衛和米鴿之間.
有很大的差別.
就是對於上帝的認識.
可見這真是一個很關鍵性的.
去解釋為什麼後來米鴿完全被棄.
而大衛有另外的妻子的其中一個原因.
聖經裡面也提醒我們這樣說.
當然今天我們不是大衛當時的情況.
沒有三妻四妾.
不過一樣.
聖經也是這樣說.
這是耶穌的說話.
不單止說了一次.
我們看看路加福音十四章二十六節.
主耶穌這樣說.
我們一齊讀一讀.
「引到我這類來.
若不愛我.
勝過愛自己的父母.
妻子兒女弟兄姊妹.
和自己的姓名.

$^{401}$就不能作我的門徒」.
標準高不高?.
好像很高.
是不是有一點不可理喻?.
父母妻子兒女弟兄姊妹這麼重要?.
如果不能夠愛我勝過這些.
就不能夠作我的門徒.
當基督徒是不是真的要拋妻棄子.
離棄父母拒絕親友.
是不是這樣當基督徒?.
當然不是.
很明顯不是.
為什麼呢?.
因為耶穌自己不是.
耶穌自己非常愛護他的母親.
祂非常尊敬祂的父親.
也在祂的教導裡不斷提醒我們.
我們要愛人如己.
何況是家人呢?.
所以我們看到這很明顯的.
不是要我們放棄這些.
而是關鍵在哪裡呢?.
就是「若不愛我勝過」.
這才是關鍵.
可不可以愛父母愛妻子愛兒女?.
應不應該?.
但是如果他們代替了上帝.
即是說他們勝過了上帝.
在你的心目中.
這樣神說這個搞亂次序之後.
你就會有很多的受苦.
次序正確就有很多的祝福.
次序錯就有很多的受苦.
同一道理.
所以我們發覺這就是重點.
今天如果我們有夫妻.
神給我們的.
真是神的恩典.
不過我們知道.
我們必須要以主為中心.

$^{441}$以神比這些一切關係更優先.
這個適合了次序就得到神的祝福.
所以我們見到大衛在這段時間的失去.
他失去了一個保護他的妻子.
失去的那個無論是甚麼動機都好.
就是很愛他的米鴿.
但因為這個時候他樂方而途.
連他身邊的妻子.
這個朋友都失去了.
所以他心靈裡面有種非常的孤單.
當然我們再說.
米鴿未必能夠與大衛一樣的與神同心.
他家中還有些偶像.
當然那些偶像有不同的解釋.
不過無論如何.
他在後期我們發現.
一直以來他對神很少提.
米鴿很少提上帝.
亦都有少少藐視大衛愛上帝的心.
所以他只是看見表面的榮華富貴.
大衛的英勇.
但沒有看到他的內心.
那種倚靠神和愛神的心.
所以問題就在這裡出現.
事實上我們肉身的夫婦關係.
有一天都會過去的.
我們昨天都說了.
我們人生當中最長的時間的關係.
不是父母和子女.
是配偶.
所以在我們人生當中幾十年.
就是爸爸媽媽都可能是十幾二十年.
最多三十年.
但是夫妻如果能夠一生之久.
是真的最長的時間.
當然我們當中亦都有不少.
有些是單身的.
亦都有單親.
即是離開了.
所以一個人照顧子女.

$^{481}$亦都有單丁.
即是他的配偶過世了.
所以有單親的.
無論你是單親的.
還是配偶過世你是喪偶的.
現在是單丁的.
我們就會更加發覺.
原來最終的倚靠不是這些.
當然今天如果我們有配偶.
我們看重這個關係.
這個長久的關係.
不過次序一定要搞對.
我自己和師母最大深深覺得.
我們大家都一起愛主,事奉主.
就是成為我們最大最大.
連結關係的一種因素.
當然我們都是很欣賞對方.
很愛護對方.
不過次序一定不可以搞亂.
以上帝為優先.
我們就蒙福.
這個實在是真正的.
今年教會的主題就是提醒大家.
至於我和我家.
我們必定事奉耶和華.
中文翻譯為全家事奉主.
所以這就是以主為中心的家庭.
主為中心的夫婦關係.
在我的書裡有一個圖紙可以用.
我很喜歡用.
很多時候做婚前輔導.
或者和一些弟兄姊妹分享的時候.
都會提醒大家.
如果丈夫在一邊.
妻妾在一邊.
上帝在上面.
我們怎樣和神拉近關係.
就會影響我們和夫婦之間的關係的親密.
所以我們越近上帝.
夫婦之間的距離就越短.

$^{521}$越離開上帝.
夫婦之間的距離就會越大.
所以我們發覺.
如果你想你的婚姻蒙福.
婚姻美滿.
很重要的因素就是能夠大家向主.
你說我未信主的配偶.
幫幫他吧.
為他祈禱不要放棄.
我們還有很多弟兄姊妹未信主的配偶.
但是你要不斷提醒自己.
你想和他建立一個真實美好相愛的關係.
就是唯有大家都在主的裡面.
一起更加緊實.
或許他現在還未信主.
希望他近一點.
越近上帝.
你和上帝之間越近.
關係也會越近.
這不單是用在夫妻.
用在子女.
用在朋友.
用在其他人際關係上.
所以我們見到這是一個很重要的因素.
所以親愛的各位弟兄姊妹.
或許你是單親的.
或者你是單親的.
或者你是單丁的.
或者你是夫婦的.
但我們知道.
原來甚至夫妻的關係.
都不是我們生命當中最終的要求.
因為有一天這些東西都會失去.
Amen.
你們不敢說Amen.
這是事實.
第三我們看到.
當時大衛經歷的恩師不是他最終的要求.
少年人需要師父.
撒慕爾記上十九章.

$^{561}$當他逃難的時候.
他逃去哪裡呢?.
十八節這樣說.
一起讀好嗎?請讀.
大衛逃避.
來到拉瑪見撒慕爾.
張素羅向他所行的事述說一篇.
他和撒慕爾就往拿約去居住.
所以大衛一逃亡.
第一個想到是誰?.
逃去誰?.
先不在雅吉那裡.
先逃到撒慕爾那裡.
去到拉瑪.
因為這是撒慕爾居住的地方.
所以他就一五一十.
張素羅本來怎樣怎樣.
現在怎樣怎樣.
整個述說一篇.
他很需要跟別人說.
真的很需要跟別人說.
所以為什麼我們有時接到一些弟兄姊妹.
有時可能你們小組或者團契.
他們面對有些困難.
他們真的很想吐口苦水.
很需要找一隻耳聽一聽.
找個心共鳴.
所以大衛一想.
就想到這個屬靈師父.
這個高他為王的.
在十六章那裡的撒慕爾.
而且他不單止跟他個人關係.
他知道唯有這個以色列國.
這個是先知.
是一個最後的事先.
所以他就去到這個屬靈的領袖.
不是軍事領袖.
去到屬靈的領袖那裡.
來跟他吐苦水.
來請求他.

$^{601}$可能為他祈禱.
撒慕爾說我斷不停止為你們禱告.
還有就是向他尋求一些指引.
結果撒慕爾真的帶他去了拿約居住.
即是轉個地方.
可能素羅知道撒慕爾在那裡.
不要讓他找到.
去到二十章第一字.
我們一起讀.
大衛從拉瑪的拿約.
就是剛才那個地方.
逃跑來到約拿丹那裡.
對他說我做了什麼.
有什麼罪孽呢.
你父親面前犯了什麼罪.
他竟尋索我的性命.
所以後來大衛又從拉瑪的拿約.
再逃到另外去找約拿丹.
所以可見在這裡他離開了撒慕爾.
以後再沒有記載.
他跟撒慕爾之間有機會見面.
因為撒慕爾就死了.
所以撒慕爾年紀老邁就離開了.
所以這個他最終的依靠.
有可能當時二十章的時候.
大衛已經知道.
他不能夠再依靠撒慕爾.
所以他就回去找約拿丹.
再尋求約拿丹為他解圍.
所以撒慕爾可見.
對於大衛有一個很重要的影響.
可以說是他的屬靈師傅或恩師.
是高他做以色列王第一個.
告訴他上帝看人不像人看人.
神看人是看內心.
人看人是看外貌.
所以大衛就捉到這個很重要的.
得著撒慕爾的開啟.
靈竅得開.
聖靈也充滿他.

$^{641}$因為當時他就有從撒慕爾來的祝福.
所以我們看到他不單是以色列國的先知領袖.
也是大衛親身經驗過的屬靈師傅.
你們有沒有師傅?.
現在我們開門徒訓練班.
我們門徒訓練的學生.
是門徒訓練的徒弟和師傅.
老師也許會有.
或者以前做過團契導師.
在年青的時候.
現在他們在哪裡知不知道?.
你們有沒有跟他們聯絡?.
有些走了,有些失了聯絡.
有些香港更加不見了.
你們以前有沒有看過一套香港很久很久以前叫做.
我都不記得他的名字.
我只記得他叫草蜢仔.
草蜢仔叫做金桂祥.
你們誰有看過這套片?.
這麼少人看過?.
沒有人看過金桂祥?.
你們不是我的時代.
叫什麼名字?.
功夫就叫功夫?.
他是白人.
學完少林寺畢業之後.
打完那條牧人巷出來.
入世了.
去劫富濟貧做好事.
每逢遇到壞人.
或者危險關頭的時候.
那套片經常都是這樣.
就是他記得師父怎樣教他.
你們有沒有記得?.
阿祥記得.
你們都是叫阿祥.
所以師父對他很重要.
因為他經常記得當時師父怎樣教我.
他就用師父教他的應對當時的處境.
不知道現在有沒有得看.

$^{681}$什麼Netflix.
YouTube可以看看.
有少少懷舊.
講到懷舊.
不知道你們以前有沒有什麼老師很懷念他?.
我們以前畢業寫紀念冊的.
叫老師寫什麼紀念冊.
沒有了 都扔了.
我最記得我以前有一個中學老師.
我又愛又恨.
愛是因為他很用心教我們.
教我們英文.
當時我的英文普普通通.
但我記得他在中五會考最後一年.
他去我們中學教英文老師.
在那一年.
中四的老師教我的英文教壞了.
英文還差了.
中五的時候被這個老師教好.
結果我拿了優異的成績.
所以很愛他.
不過很恨他.
因為他經常用激將法.
很壞.
我平面這次考得好一點.
比上次好.
他說這樣也算好.
我看你也不會成功.
這樣?.
很像你.
可能激將法激到我更加用心.
我要厲害給你看.
他不只是這樣對我.
他對其他同學也是這樣.
所以他最喜歡罵人.
上課時我都被他罵.
當時我小學到中學都有機會參加朗誦比賽.
那次是英文朗誦.
拿了第三.
是我這麼多年來最高名次的一次.

$^{721}$各校制比賽拿第三.
每逢校制比賽.
音樂也好.
朗誦也好.
他們把名字放在學校外面.
這個同學拿了校制比賽的第幾名.
我回家看到我的名字很開心.
我還沒試過.
結果上課的時候.
這個英文老師英文朗誦.
他應該說「好ye 喎」.
他說「你是不是碰運氣的?」.
真是很生氣.
所以他又愛又恨.
所以我很深記得他.
另外在我人生當中的恩師.
是我的上學院院長.
他是劉瑤光牧師.
他去年回天教.
我寫了一篇紀念文章在紀念集.
題目叫做「良牧恩師」.
紀念他.
我紀念他關懷之恩.
因為我和師母進去上學院的時候.
我們要分開男女宿舍住.
已經結婚兩年了.
要分開住.
所以他就關心我們.
後來經過了九個月之後.
他就特意找了一個.
從來在神學院未有過的已婚宿舍.
給我們租了一間房間.
讓我們可以住在已婚宿舍.
不但關懷他也是造就之恩.
在我自己讀神學這麼多的過程裡.
劉瑤光牧師對於讀聖經解經影響我是最深的.
所以我多謝他造就.
推介之恩.
我來美國是他推介我來的.
所以我在橙園玲瓏堂.

$^{761}$最初一個教會事奉的時候.
教會問一個牧師是否認識人.
有沒有很好的牧師.
問了劉瑤光牧師推薦我和師母.
所以我們就這樣來美國.
最後他就是輔助之恩.
所以他一直以來在我們玲瓏堂事奉.
有時我們也有請教他.
我邀請了他.
唯一一個在兩次下令會做講員的.
就是劉瑤光牧師.
所以在他離開的時候.
我很懷念他.
所以寫了一篇文章叫做《良牧恩師》紀念他.
不過我們都知道.
今天我再想有事問他.
沒得問了.
只可以懷念 因為人總要過去.
在聖經裡面也是這樣說.
他說.
「你們不要倚靠君王.
不要倚靠世人.
他的氣一斷就歸回塵土.
他所打算的當日就消滅了」.
所以在聖經裡面說.
你不要倚靠君王.
不要倚靠世人.
為什麼?.
因為他的氣一斷就歸回塵土.
真是最好的最好的.
都有一天會離開的.
所以這裡表明了人的有限.
所以我們知道.
恩師也好 老師也好 導師也好 牧師也好.
不過不是我們人生最終的依靠.
這是大衛在生命當中深深經歷到的.
最後我們看到.
好友 也不是最終的依靠.
撒母耳記二十章第四十二節.
我們一起讀.

$^{801}$「若拿丹對大衛說.
我們二人曾指著耶和華的名起誓說.
願耶和華在你我中間.
並你我後裔中間為證.
直到永遠.
如今你平平安安地去吧.
大衛就起身走了.
若拿丹也回城裡去了」.
請問這句對話是誰和誰?.
若拿丹和大衛這樣說.
我們二人在耶和華面前起誓.
這是上次我們大家還記得的.
我們怎樣做良朋益友的時候.
我們講過很多經經的人.
他們立約彼此在神面前.
說願耶和華在我們中間.
並在你我後裔為證.
什麼時候說直到?.
直到永遠.
他的心是這樣直到永遠.
有沒有可能直到永遠?.
有可能就是後裔.
我們自己就不能直到永遠.
結果從這一刻開始.
這是最後記載.
大衛就起身走了.
若拿丹就回城裡去了.
從此二人就分道揚鑣.
所以我們上次剛剛講過.
大衛和若拿丹的友情非常寶貴.
非常重要.
但是這種關係也不能永遠維持.
最好的朋友有一天都要離開.
他不單止失去權勢.
失去妻子.
失去恩師.
最後連好友都失去.
這真是很大的失去.
大衛一生都有不少朋友.
有些永遠懷念.

$^{841}$但有些仍然令他很嘆息.
因為朋友有時好.
有時真的有些不好的朋友.
所以在詩篇41篇.
詩篇41篇這裡很諷刺.
我們一起讀一讀.
連我知己的朋友.
我所依靠.
吃過我飯的.
也用腳踢我.
真是很大的一種慨嘆.
有時朋友都會被他出賣.
這當然不是指若拿丹.
不過在大衛人生過程裡有很多好友.
曾經本來和他很好朋友.
後來令他很大失望.
有些甚至通水給蘇羅捉他.
有些亦忘恩負義.
本來對他很好的.
後來他忘記了大衛對他的恩情.
所以這裡讓我們看見.
知己的朋友真是好的.
不過亦有時令我們很大的創傷損失.
不知道你們有沒有好朋友.
曾經過往.
小學中學.
我很羨慕有些.
現在我們那些幾十年的老友在重聚.
幾十年的老友能夠維持關係.
真是很難得.
真是很珍惜.
我記得我自己中學時有好幾個死黨.
我們經常在學校出街.
一齊打橋牌.
一齊打球.
後來我考入政府做移民局.
有一個經常和我們在一起的又考入來.
大家還是同事.
所以那時很單純.
如果大家覺得那些友情.

$^{881}$我覺得最早期的友情很單純.
很珍惜.
不過就是因為太早期.
現在沒辦法找回.
特別是來了美國.
香港那些找不回.
現在有很多尋人法.
Facebook.
尋不到.
我打回以前有個小學同學.
坐在我旁邊的女同學.
對斯沒有什麼影響.
小學女同學.
我嘗試去找找.
那時經常經過她的屋.
找不到.
Facebook找不到.
有時找到.
不過真的發覺.
人生不斷的關係都會轉變.
我又記得我自己在政府工作.
最初在人大事務處.
要受訓練.
所以同班一起訓練十幾二十個人.
很老友.
因為大家進入一個新的行業.
大家一起受訓.
而且那時受訓.
早上要步操.
下午又要上課.
晚上又要一起做一些事.
有時要去練習智慧術.
有時要去游泳.
因為我們做移民局要學游泳.
又要去練智慧術.
有時要進入特別部門要練開槍.
所以大家有很多不同的事.
大家經歷過.
大家的步操.
熱到死.

$^{921}$一起磨那雙鞋.
磨到那雙鞋要見到自己的樣子.
要靚到.
所以我們畢業之後.
分配去不同的地方.
經常都有聯絡.
大家都聚會.
但是經過一年兩年.
就開始沒有了.
關係就是這樣.
後來我們神學院.
我記得神學院同班同學.
也是最難忘的.
可以說是最深的.
因為我們四年同窗.
要住在一起.
住在神學院.
由早上六點見到晚上九點.
四年.
有苦有樂.
因為我們在神學院.
睡三格床.
你們有沒有睡過.
三格床.
最慘的是睡中間的.
下面那個點來點去.
你又要動.
上面那個踏腳發夢的時候.
你又要醒.
真的很痛苦.
所以又開心又難忘.
不過很可惜.
大家一畢業的時候.
我們經常都說這句話.
我們分道揚鑣.
多謝疫情.
疫情就有WhatsApp.
以前沒有那麼多WhatsApp.
疫情WhatsApp幫我們組織.
前年我們的班會在WhatsApp可以見面.

$^{961}$一數三十年.
三十年沒有見面.
真是難得.
不過現在又是.
久不久有個通訊.
沒有什麼經常聯絡.
因為個個都忙.
個個都忙自己的事.
我不知道弟兄姊妹.
你有沒有都覺得.
有很多你的好朋友.
有些走就走.
離開就離開.
失聯絡就失聯絡.
其實都告訴我們.
人生就是這樣.
說到這裡.
牧師即是四大皆空.
事實上聖經裡面.
真是指示我們這個道理.
因為聖經告訴我們.
萬事皆空.
虛空的虛空.
為主事足.
這就是聖經的道理.
所以來到我們這樣的年紀.
七老八十.
我們會發覺.
當然有些年輕些好些.
我們希望珍惜現在有的.
如果神給你有權位有機會.
你知道這些是從神那裡來的.
不要將它成為你最終的要求.
但要善用它.
去服侍神要你服侍的對象.
如果你現在有家庭.
有妻有丈夫.
好好珍惜它.
因為不會永遠.
唯有在主的裡面.

$^{1001}$我們在天家裡面.
那個才是永恆的.
這個就是最重要.
所以繼續為我們微信主義的家人.
配偶 禱告.
繼續將他們拉近一點希望.
這個是我們的心願.
不過你要記得.
這個也不是我們最終的要求.
如果你失去了你的配偶.
你記得你永遠都不孤單.
因為有上帝.
如果我們今天.
朋友都很難找回.
恩師更加不在.
我們就知道.
我們能夠有一個.
我們真正最終的要求.
為什麼呢?.
因為萬福的源頭就是從神那裡而來.
所以無論是恩師也好.
無論是朋友也好.
無論是家庭.
我們的丈夫.
配偶 子女 權勢.
所有這些一切.
最終都是這位次因的主.
所以不要將神賜給我們這些.
代替了次因的主.
否則我們就會有很大的受苦.
相反次序擺正.
我們就會有很大的祝福.
因為縱視這些.
都會有一天失去的時候.
我們知道我們的心仍然安穩.
我們有我們最終的要求.
珍惜神給你現在有的機會.
有你的關係.
真的很珍惜.
但是不要成為我們最終的要求.

$^{1041}$讓我們的心僅僅是與主.
以他為我們人生的依靠.
\newpage



\section{創世記 25:16}
\label{sec:031741}
\textbf{聖經人聖經事:十二族長}~~~~黃克勤牧師博士
\newline
\newline
連結: \href{http://www.church.com.hk/acms/content.asp?site=cdc\&op=show\&type=product\&code=031741\&layout=sermon}{\texttt{031741}} ~~~~ 語音日期: 2026-02-21
\newline
\newline
\hyperref[sec:031740]{\small{< < < PREV SERMON < < <}}
~
\hyperref[sec:index_chronic]{\small{[返順時目]}}
~
\hyperref[sec:031742]{\small{> > > NEXT SERMON > > >}}
\newline
\newline
創世記 25:16
\newline
\begin{longtable}{cl}
\hline
\hline
章節 & 經文 (和合本修訂版)\\
\hline
25:16 & \begin{tabularx}{0.7\textwidth}{X} 這些都是以實瑪利的兒子們。他們的村莊和營寨按著他們命名；他們作了十二族的族長。 \end{tabularx} \\ \\
[1ex]
\hline
\hline
\end{longtable}
$^{1}$大家好,又和大家一起讀聖經.
《生命說書教網絡》.
聖經人,聖經事.
又和大家見面.
一起認識聖經人與事後面的上帝是怎樣的.
讀聖經最重要是認識上帝.
就是.
眨眼一日.
嘩,又來.
這次是十二續章.
大家記得我們上次一直說.
亞伯拉罕這個人和他的事.
亞伯拉罕已經過去了.
終於有一天他離開了世界.
歸到列祖那裡.
然後這次看十二續章.
十二個續,有續章.
哦,即是以色列十二支派.
留意一下.
《創世記》二十五章.
《創世記》挺好的.
挺有次序的.
是周到的.
說完亞伯拉罕.
就開始說.
他的兒子以實瑪利.
先說亞伯拉罕,再說以薩.
以薩然後他又要處理一下.
以實瑪利怎樣下去.
交代一下.
不是處理,是交代一下.
在他的寫作上處理好以實瑪利.
第十二節是撒拉的侍女.
撒拉是亞伯拉罕的太太.
她的侍女,即是他的婢女.
或者是買回來的奴隸.
由她的名字就知道了.
撒拉的侍女是一個埃及裔人.
下甲給亞伯拉罕所生的兒子.
就是以實瑪利.

$^{41}$撒拉是亞伯拉罕的太太.
就是元配.
這個侍女很有可能是一個婢女.
或者是女的奴隸.
買回來服侍主母.
什麼人士來的?.
埃及人.
埃及趁這個機會和大家解釋一下.
埃及這個字是有意思的.
Mithrae.
Mithrae呢.
TS是讀J.
Mithrae.
Mithrae是兩個海峽.
很有可能埃及去到出海的地方附近.
像兩個海峽對著.
Mithrae.
兩個海峽.
他是埃及人.
而他的名字本身就是下甲.
下甲這個字很有意思.
是逃跑的意思.
因此推測他可能在埃及.
不知道發生什麼事.
就逃跑了出來.
就改了這個名字叫逃跑下甲.
沒理由家人自己改他的名字叫逃跑.
逃跑.
或者可能他小時候很頑皮.
經常到處跑.
一說他就跑了出去.
可能是這樣.
所以他的名字叫做肯加拿逃跑.
這個下甲是給亞伯拉罕生了一個兒子.
叫做以實瑪利.
埃及人開始在亞伯拉罕的家.
服侍撒拉的時候.
開始認識神.
認識耶和華神.
於是亞伯拉罕跟他改了一個名字.

$^{81}$可能叫做以實瑪利.
Ishmael.
這個名字以前出現過.
我都介紹過了.
Ishmael.
其實這個L有一張EL.
Ishmael.
L就是神.
Ishma.
Ishmael.
Ishmael的意思是神會聽見.
神會聽到的.
大家一定要知道.
神是會聽你的禱告的.
亦都會聽你說話的.
所以小心說.
神是聽到的.
不要以為祂不知道.
不要迷倒.
不是的.
神雖然是你肉眼看不到.
但是你說的話神聽到.
甚至你心裡面說的話.
神都聽到.
Ishmael.
就是以實瑪利.
好名字.
第十三節.
還沒說.
第一點.
我想說的是.
有一個螢光的第一點.
以實瑪利其實是.
上帝是聽到的.
而且上帝是.
很信實.
很有信用.
跟亞伯拉罕說過.
你要出的後面成為大俗.
當時以撒拉就沒那麼.

$^{121}$信心不夠.
明明都說會生兒子的.
但是不會的.
我老公一百歲.
我九十歲.
怎麼生呢.
過世生到年齡.
生不到的了.
這麼多年都沒生到.
因此神沒給我們.
有辦法.
很進取的.
唯有在當時很流行的.
將他的侍女.
當時應該這樣做.
侍女下甲.
撒拉很疼愛的.
將她給亞伯拉罕做妾.
於是生了.
於是看到.
不能生育的是女兒.
而不是男兒.
生了以實瑪利.
認識上帝.
於是改這個名字.
神會聽見我們祈禱.
會聽見我們的心.
以實瑪利.
出來了.
以實瑪利.
神說.
以實瑪利不是我將來.
透露過他.
而興起一個大族的人.
不是的.
不過.
是由你所出的.
我都認.
我都算.
因此我會祝福他.

$^{161}$令他成為一個大國.
不過沒有想到.
將來以實瑪利.
就是今天的阿拉伯人.
剛剛與兄弟以色列人.
相爭了千百年.
在這個土地上.
在這個迦南土地上.
一直到今天.
應許之地.
叫做以色列地.
或者叫巴勒斯坦.
最近期新的幾百年的名字.
百年的名字.
就是巴勒斯坦.
以實瑪利又生了一群兒子.
這裡看也很有趣.
第二點我想說.
以實瑪利生了很多兒子.
接著他的家譜就寫在下面.
留意.
以實瑪利的長子叫做尼拜約.
名字我不解釋了.
尼拜約大兒子.
然後又有基達.
阿德別.
米比三.
米爾斯瑪.
杜瑪.
瑪薩.
哈大.
提瑪.
伊特.
拿菲斯.
基底瑪.
數一數.
有幾個兒子在這裡.
1.
2.
3.

$^{201}$4.
5.
6.
7.
8.
9.
10.
11.
12個兒子.
以實瑪利有12個兒子.
有沒有發覺.
和雅各有些相似.
12個兒子.
厲害.
以實瑪利.
12個兒子.
這麼厲害.
你知道12這個數目.
是聖經很歡喜.
表達上帝一個完整的數目.
神的心意一個完整的數目.
12個兒子.
似以色列十二之派.
就是.
第二點就是說.
其實神說我會祝福以實瑪利成為大國.
都會是一個不可以傾向的大族.
神說了真的會發生.
就有12個兒子.
這12個兒子.
第三點看看.
這12個兒子成為十二族.
這12個兒子以實瑪利種子的名字.
照著他們的村莊.
營寨.
逐漸.
他們各有自己的村莊.
後來發展到營寨.
比較大的區域.
作料第三.

$^{241}$第三點要說的是十二族.
是什麼來的呢?.
那12個兒子呢?.
這12個兒子是做了十二個地方.
和自己的家人.
一路繁衍出來很多人.
可以去到一個村莊的營寨.
做了十二族的族長.
好了.
十二個兒子做了族長.
就是governor.
chief的族長.
是管理那個村的營寨的十二族人.
那十二族人呢?.
咦.
十二族和十二支派不同的呢?.
是不同的.
用的字是不一樣的.
十二族這個族字呢.
這個族叫做Umar.
Umar.
意思是十二個民族.
十二族人.
十二家人.
不是家人.
就是十二個名字的民族.
十二個族人.
所以這個族字翻譯得很好.
十二個Umar.
十二族人.
但是是不是十二支派呢?.
十二支派呢.
以色列十二支派.
不是Umar這個字來的.
支派的意思是一支將.
一支將或者是管理或者權力.
也可以叫做支派.
就是那個group的人.
People group.
所以十二族的Umar是小一點的.

$^{281}$以及是指一個地方的一族人.
Umar.
支派呢.
是特別有一個字的.
就是一支棍.
好像十二支杖一樣.
上帝十二支杖.
好了.
四個十二支杖的族長.
以實瑪利享壽是137歲.
棄絕而死.
我剛才想說的.
就是十二族人和十二支派是不同的.
不過這裡顯示.
上帝是都祝福以實瑪利.
給他都像縮影版.
縮小版的以色列十二支派.
十二族人.
這樣神真的很好.
以實瑪利享壽137歲.
棄絕而死.
什麼叫棄絕而死呢.
棄絕就是壽終正寢.
去到那裡就很老了.
開始不行了.
好像沒電了.
棄絕而死.
即是死得很好.
歸到他列祖那裡.
這裡特別用了紅色字引大家注意.
歸到他列祖那裡.
這裡這樣說.
棄絕而死.
享壽137歲就完了.
句號.
但是他沒有.
他特別說歸到他列祖那裡.
他們以前的殯葬有一個墳墓.
亞伯拉罕買了一個山洞.
然後將撒拉先葬下去.

$^{321}$亞伯拉罕死了.
又葬在山洞裡.
就一齊了.
現在以實瑪利死了.
又葬在這裡.
歸到他列祖那裡.
就是他們歡喜葬在同一個山洞.
所以叫歸到他列祖那裡.
創世記的約瑟.
在埃及做完宰相之後離開世界了.
你們祭了我的身體.
回去應許之地故鄉的時候.
記得帶著我.
我知道有一天我們會回去的.
整個民族.
記得帶著我.
因為他要和列祖同睡.
這是他們的生命完整與終結.
那洞就越來越大了.
他要在山裡作.
如果去到耶穌時代.
羅馬時代.
他們會在第一年.
就是一條屍體包好之後.
就放在這裡露空.
不過用石擋住墳墓的門.
一個大山洞.
然後等它慢慢化了肉.
洗乾淨骨頭.
就放進盒子裡.
骨盒.
然後再作一個小一點的洞.
將骨盒放進去.
兒子又作大一點的洞.
將骨頭放進去之後.
就在骨盒裡放進去.
在爸爸旁邊.
在媽媽旁邊.
歸到列祖阿裡.
就這樣與列祖同睡.

$^{361}$能夠有這一句.
就是說上帝都將以實瑪利.
都算為是亞伯拉罕的宗族下來的.
是一家人.
所以將他叫做歸到他列祖阿裡.
而他子孫的住處在他縱弟兄的東邊.
從哈菲拉直到埃及前的樞爾.
正在亞述的道上.
其實他們住的地方.
如果是以實瑪利.
亞伯拉罕等等他們住的地方.
是南地.
比較是曠野一點的.
在耶路撒冷南很多.
我想他都要一個小時旅遊巴了.
住在那裡很荒涼很寧靜.
搭帳篷住的.
靠近地中海一點.
但其實看不到海很遠.
不過這一班以實瑪利的後裔.
全部是住在這個地方比較東邊一點.
比較東邊一點.
而部分就向南一點.
到了埃及前的樞爾.
又是一個曠野.
樞爾的曠野.
那裡有一條路繞上去.
到以色列的北.
就叫做亞述的道.
搬去住.
都不是搬去一個深山野嶺的.
都是一個通渠大邑.
都是有些地方.
是有些相對經過的大道上.
才有得來往.
就是說.
他們都有了一個地方.
是上帝給他的十二族人在那裡發展的.
而以實瑪利就歸到他列祖那裡.
都是相當好的結局.

$^{401}$有十二族人.
都是應驗了上帝的信實.
給他們都成為一個大族.
就是神應許以實瑪利的.
很感恩的.
看看今天的金句就是神信實.
我們認識神.
說過就是了.
Faithfulness.
就是以實瑪利種子的名字.
剛才讀了.
照著他們的村莊營寨.
他們有十二族.
而且他們那十二個兒子.
就做了族長了.
想不到以實瑪利.
神真的照著他的說話.
不用擔憂.
雖然撒拉趕他們走.
不過我都會祝福以實瑪利.
由於他是從里阿巴拉罕而出.
我都會祝福他.
都是成為一個大族的.
大家知道這件事了.
都有十二族人的.
雖然沒有什麼文名.
不過都有十二族.
上帝是信實的.
是看顧人.
以實瑪利的名義就是.
神會聽見.
我們記得一起祈禱吧.
天父上帝我們向你感恩讚美.
我們認識你的心裡面真的很舒服.
因為你是一位.
應許信實可靠的上主.
你是會聽到我們的心聲.
你聽到我們的祈禱.
你知道我們最.
困苦悲哀的地方.

$^{441}$你知道的.
你聽到我們的心聲.
你會聽見以實瑪利就是這個名字.
而他出世了.
雖然可能是開頭.
都有一些安排上不是最好.
不過既然他從里阿巴拉罕而出.
你就照樣在他身上要成就大事.
讓他有十二族的子孫.
一路成為大國.
今天是和以色列人一齊.
可以居住在一起.
不過很可惜他們彼此有很多爭鬥.
我們見到人有很多結果.
一路會生出來.
不過求儲憐憫.
你是會聽見的神.
願你得著榮耀.
禱告奉耶穌明求.
阿們.
\newpage



\section{民數記 5:1-31}
\label{sec:031742}
\textbf{神聖群體的健康學}~~~~李志輝傳道
\newline
\newline
連結: \href{http://www.church.com.hk/acms/content.asp?site=cdc\&op=show\&type=product\&code=031742\&layout=sermon}{\texttt{031742}} ~~~~ 語音日期: 2026-02-28
\newline
\newline
\hyperref[sec:031741]{\small{< < < PREV SERMON < < <}}
~
\hyperref[sec:index_chronic]{\small{[返順時目]}}
~
\hyperref[sec:031743]{\small{> > > NEXT SERMON > > >}}
\newline
\newline
民數記 5:1-31
\newline
\begin{longtable}{cl}
\hline
\hline
章節 & 經文 (和合本修訂版)\\
\hline
5:1 & \begin{tabularx}{0.7\textwidth}{X} 耶和華吩咐摩西說： \end{tabularx} \\ \\ \relax
5:2 & \begin{tabularx}{0.7\textwidth}{X} 「你要吩咐以色列人，把一切患痲瘋的、患漏症的和因屍體而不潔淨的，都送到營外去。 \end{tabularx} \\ \\ \relax
5:3 & \begin{tabularx}{0.7\textwidth}{X} 無論男女你都要送，把他們送到營外，免得他們玷污了他們的營，這是我住在他們中間的地方。」 \end{tabularx} \\ \\ \relax
5:4 & \begin{tabularx}{0.7\textwidth}{X} 以色列人就照樣做，把他們送到營外去。耶和華怎樣吩咐摩西，以色列人就照樣做了。 \end{tabularx} \\ \\ \relax
5:5 & \begin{tabularx}{0.7\textwidth}{X} 耶和華吩咐摩西說： \end{tabularx} \\ \\ \relax
5:6 & \begin{tabularx}{0.7\textwidth}{X} 「你要吩咐以色列人：無論男女，若犯了人所常犯的任何罪，以致干犯耶和華，那人就有了罪。 \end{tabularx} \\ \\ \relax
5:7 & \begin{tabularx}{0.7\textwidth}{X} 他要承認所犯的罪，將所虧負人的如數賠償，另外再加五分之一，交給所虧負的人。 \end{tabularx} \\ \\ \relax
5:8 & \begin{tabularx}{0.7\textwidth}{X} 那人若沒有至親可接受所賠償的，所賠償的就要歸耶和華，交給祭司；另外還要獻一隻贖罪的公羊為他贖罪。 \end{tabularx} \\ \\ \relax
5:9 & \begin{tabularx}{0.7\textwidth}{X} 以色列人一切的聖物中，所奉給祭司的一切禮物都要歸給祭司。 \end{tabularx} \\ \\ \relax
5:10 & \begin{tabularx}{0.7\textwidth}{X} 各人自己的聖物歸自己，給祭司的要歸給祭司。」 \end{tabularx} \\ \\ \relax
5:11 & \begin{tabularx}{0.7\textwidth}{X} 耶和華吩咐摩西說： \end{tabularx} \\ \\ \relax
5:12 & \begin{tabularx}{0.7\textwidth}{X} 「你要吩咐以色列人，對他們說：若任何人的妻子背離婦道，對丈夫不貞， \end{tabularx} \\ \\ \relax
5:13 & \begin{tabularx}{0.7\textwidth}{X} 有人與她同寢交合，這事瞞過她的丈夫，沒有被發現；她玷污自己，沒有證人指控她，也沒有被捉住； \end{tabularx} \\ \\ \relax
5:14 & \begin{tabularx}{0.7\textwidth}{X} 丈夫若生了疑忌的心，對妻子起了疑忌，認為她玷污自己；或是丈夫生了疑忌的心，對妻子起了疑忌，雖然她沒有玷污自己， \end{tabularx} \\ \\ \relax
5:15 & \begin{tabularx}{0.7\textwidth}{X} 這人要帶妻子到祭司那裡，同時為她帶十分之一伊法大麥麵粉作供物。不可澆上油，也不可加乳香，因為這是疑忌的素祭，是紀念的素祭，使人記得罪孽。 \end{tabularx} \\ \\ \relax
5:16 & \begin{tabularx}{0.7\textwidth}{X} 「祭司要使那婦人近前來，站在耶和華面前。 \end{tabularx} \\ \\ \relax
5:17 & \begin{tabularx}{0.7\textwidth}{X} 祭司要把聖水盛在瓦器裡，從帳幕的地上取些塵土，放在水中。 \end{tabularx} \\ \\ \relax
5:18 & \begin{tabularx}{0.7\textwidth}{X} 祭司要帶那婦人站在耶和華面前，使她蓬頭散髮，再把紀念的素祭，就是疑忌的素祭，放在她的手掌，祭司手裡捧著致詛咒的苦水。 \end{tabularx} \\ \\ \relax
5:19 & \begin{tabularx}{0.7\textwidth}{X} 祭司要叫婦人起誓，對她說：『若沒有人與你同寢，若你未曾背著丈夫做污穢的事，你就能免去這致詛咒的苦水。 \end{tabularx} \\ \\ \relax
5:20 & \begin{tabularx}{0.7\textwidth}{X} 但你背著丈夫，玷污自己，跟丈夫以外的人同寢。』 \end{tabularx} \\ \\ \relax
5:21 & \begin{tabularx}{0.7\textwidth}{X} 祭司叫婦人賭咒起誓，祭司對她說：『當耶和華使你大腿萎縮，肚腹腫脹時，願耶和華使你在你百姓中成為詛咒和咒罵； \end{tabularx} \\ \\ \relax
5:22 & \begin{tabularx}{0.7\textwidth}{X} 願這致詛咒的水進入你體內，使你肚腹腫脹，大腿萎縮。』婦人要說：『阿們，阿們。』 \end{tabularx} \\ \\ \relax
5:23 & \begin{tabularx}{0.7\textwidth}{X} 「祭司要把這詛咒寫在冊上，然後用苦水塗去， \end{tabularx} \\ \\ \relax
5:24 & \begin{tabularx}{0.7\textwidth}{X} 又叫婦人喝這致詛咒的苦水，這詛咒的水要進入她裡面，令她痛苦。 \end{tabularx} \\ \\ \relax
5:25 & \begin{tabularx}{0.7\textwidth}{X} 祭司要從婦人手中取那疑忌的素祭，把素祭在耶和華面前搖一搖，拿到祭壇前； \end{tabularx} \\ \\ \relax
5:26 & \begin{tabularx}{0.7\textwidth}{X} 又要從素祭中取出一把，作為紀念，燒在壇上，然後叫婦人喝這水。 \end{tabularx} \\ \\ \relax
5:27 & \begin{tabularx}{0.7\textwidth}{X} 祭司叫她喝了以後，她若玷污自己，確實對丈夫不貞，這致詛咒的水必進入她裡面，令她痛苦，她的肚腹就要腫脹起來，大腿萎縮；這婦人就在她百姓中成為詛咒。 \end{tabularx} \\ \\ \relax
5:28 & \begin{tabularx}{0.7\textwidth}{X} 這婦人若沒有玷污自己，是貞潔的，就要免受這災，並且能夠生育。 \end{tabularx} \\ \\ \relax
5:29 & \begin{tabularx}{0.7\textwidth}{X} 「這是疑忌的條例。妻子背離丈夫玷污自己， \end{tabularx} \\ \\ \relax
5:30 & \begin{tabularx}{0.7\textwidth}{X} 或是丈夫生了疑忌的心，對妻子起了疑忌，祭司要使那婦人站在耶和華面前，在她身上照這條例而行。 \end{tabularx} \\ \\ \relax
5:31 & \begin{tabularx}{0.7\textwidth}{X} 男人可免罪責；女人必須擔當自己的罪孽。」 \end{tabularx} \\ \\
[1ex]
\hline
\hline
\end{longtable}
$^{1}$陳克勤議員.
胡志偉議員.
陳克勤議員.
胡志偉議員.
陳克勤議員.
胡志偉議員.
陳克勤議員.
胡志偉議員.
陳克勤議員.
莫乃光議員.
陳克勤議員.
莫乃光議員.
陳克勤議員.
莫乃光議員.
陳克勤議員.
莫乃光議員.
陳克勤議員.
莫乃光議員.
其他事項.
\newpage



\section{民數記 5:1-31}
\label{sec:031743}
\textbf{神住在我們中間:聖潔的要求和實踐}~~~~何慧芬牧師
\newline
\newline
連結: \href{http://www.church.com.hk/acms/content.asp?site=cdc\&op=show\&type=product\&code=031743\&layout=sermon}{\texttt{031743}} ~~~~ 語音日期: 2026-03-01
\newline
\newline
\hyperref[sec:031742]{\small{< < < PREV SERMON < < <}}
~
\hyperref[sec:index_chronic]{\small{[返順時目]}}
~
\hyperref[sec:031744]{\small{> > > NEXT SERMON > > >}}
\newline
\newline
民數記 5:1-31
\newline
\begin{longtable}{cl}
\hline
\hline
章節 & 經文 (和合本修訂版)\\
\hline
5:1 & \begin{tabularx}{0.7\textwidth}{X} 耶和華吩咐摩西說： \end{tabularx} \\ \\ \relax
5:2 & \begin{tabularx}{0.7\textwidth}{X} 「你要吩咐以色列人，把一切患痲瘋的、患漏症的和因屍體而不潔淨的，都送到營外去。 \end{tabularx} \\ \\ \relax
5:3 & \begin{tabularx}{0.7\textwidth}{X} 無論男女你都要送，把他們送到營外，免得他們玷污了他們的營，這是我住在他們中間的地方。」 \end{tabularx} \\ \\ \relax
5:4 & \begin{tabularx}{0.7\textwidth}{X} 以色列人就照樣做，把他們送到營外去。耶和華怎樣吩咐摩西，以色列人就照樣做了。 \end{tabularx} \\ \\ \relax
5:5 & \begin{tabularx}{0.7\textwidth}{X} 耶和華吩咐摩西說： \end{tabularx} \\ \\ \relax
5:6 & \begin{tabularx}{0.7\textwidth}{X} 「你要吩咐以色列人：無論男女，若犯了人所常犯的任何罪，以致干犯耶和華，那人就有了罪。 \end{tabularx} \\ \\ \relax
5:7 & \begin{tabularx}{0.7\textwidth}{X} 他要承認所犯的罪，將所虧負人的如數賠償，另外再加五分之一，交給所虧負的人。 \end{tabularx} \\ \\ \relax
5:8 & \begin{tabularx}{0.7\textwidth}{X} 那人若沒有至親可接受所賠償的，所賠償的就要歸耶和華，交給祭司；另外還要獻一隻贖罪的公羊為他贖罪。 \end{tabularx} \\ \\ \relax
5:9 & \begin{tabularx}{0.7\textwidth}{X} 以色列人一切的聖物中，所奉給祭司的一切禮物都要歸給祭司。 \end{tabularx} \\ \\ \relax
5:10 & \begin{tabularx}{0.7\textwidth}{X} 各人自己的聖物歸自己，給祭司的要歸給祭司。」 \end{tabularx} \\ \\ \relax
5:11 & \begin{tabularx}{0.7\textwidth}{X} 耶和華吩咐摩西說： \end{tabularx} \\ \\ \relax
5:12 & \begin{tabularx}{0.7\textwidth}{X} 「你要吩咐以色列人，對他們說：若任何人的妻子背離婦道，對丈夫不貞， \end{tabularx} \\ \\ \relax
5:13 & \begin{tabularx}{0.7\textwidth}{X} 有人與她同寢交合，這事瞞過她的丈夫，沒有被發現；她玷污自己，沒有證人指控她，也沒有被捉住； \end{tabularx} \\ \\ \relax
5:14 & \begin{tabularx}{0.7\textwidth}{X} 丈夫若生了疑忌的心，對妻子起了疑忌，認為她玷污自己；或是丈夫生了疑忌的心，對妻子起了疑忌，雖然她沒有玷污自己， \end{tabularx} \\ \\ \relax
5:15 & \begin{tabularx}{0.7\textwidth}{X} 這人要帶妻子到祭司那裡，同時為她帶十分之一伊法大麥麵粉作供物。不可澆上油，也不可加乳香，因為這是疑忌的素祭，是紀念的素祭，使人記得罪孽。 \end{tabularx} \\ \\ \relax
5:16 & \begin{tabularx}{0.7\textwidth}{X} 「祭司要使那婦人近前來，站在耶和華面前。 \end{tabularx} \\ \\ \relax
5:17 & \begin{tabularx}{0.7\textwidth}{X} 祭司要把聖水盛在瓦器裡，從帳幕的地上取些塵土，放在水中。 \end{tabularx} \\ \\ \relax
5:18 & \begin{tabularx}{0.7\textwidth}{X} 祭司要帶那婦人站在耶和華面前，使她蓬頭散髮，再把紀念的素祭，就是疑忌的素祭，放在她的手掌，祭司手裡捧著致詛咒的苦水。 \end{tabularx} \\ \\ \relax
5:19 & \begin{tabularx}{0.7\textwidth}{X} 祭司要叫婦人起誓，對她說：『若沒有人與你同寢，若你未曾背著丈夫做污穢的事，你就能免去這致詛咒的苦水。 \end{tabularx} \\ \\ \relax
5:20 & \begin{tabularx}{0.7\textwidth}{X} 但你背著丈夫，玷污自己，跟丈夫以外的人同寢。』 \end{tabularx} \\ \\ \relax
5:21 & \begin{tabularx}{0.7\textwidth}{X} 祭司叫婦人賭咒起誓，祭司對她說：『當耶和華使你大腿萎縮，肚腹腫脹時，願耶和華使你在你百姓中成為詛咒和咒罵； \end{tabularx} \\ \\ \relax
5:22 & \begin{tabularx}{0.7\textwidth}{X} 願這致詛咒的水進入你體內，使你肚腹腫脹，大腿萎縮。』婦人要說：『阿們，阿們。』 \end{tabularx} \\ \\ \relax
5:23 & \begin{tabularx}{0.7\textwidth}{X} 「祭司要把這詛咒寫在冊上，然後用苦水塗去， \end{tabularx} \\ \\ \relax
5:24 & \begin{tabularx}{0.7\textwidth}{X} 又叫婦人喝這致詛咒的苦水，這詛咒的水要進入她裡面，令她痛苦。 \end{tabularx} \\ \\ \relax
5:25 & \begin{tabularx}{0.7\textwidth}{X} 祭司要從婦人手中取那疑忌的素祭，把素祭在耶和華面前搖一搖，拿到祭壇前； \end{tabularx} \\ \\ \relax
5:26 & \begin{tabularx}{0.7\textwidth}{X} 又要從素祭中取出一把，作為紀念，燒在壇上，然後叫婦人喝這水。 \end{tabularx} \\ \\ \relax
5:27 & \begin{tabularx}{0.7\textwidth}{X} 祭司叫她喝了以後，她若玷污自己，確實對丈夫不貞，這致詛咒的水必進入她裡面，令她痛苦，她的肚腹就要腫脹起來，大腿萎縮；這婦人就在她百姓中成為詛咒。 \end{tabularx} \\ \\ \relax
5:28 & \begin{tabularx}{0.7\textwidth}{X} 這婦人若沒有玷污自己，是貞潔的，就要免受這災，並且能夠生育。 \end{tabularx} \\ \\ \relax
5:29 & \begin{tabularx}{0.7\textwidth}{X} 「這是疑忌的條例。妻子背離丈夫玷污自己， \end{tabularx} \\ \\ \relax
5:30 & \begin{tabularx}{0.7\textwidth}{X} 或是丈夫生了疑忌的心，對妻子起了疑忌，祭司要使那婦人站在耶和華面前，在她身上照這條例而行。 \end{tabularx} \\ \\ \relax
5:31 & \begin{tabularx}{0.7\textwidth}{X} 男人可免罪責；女人必須擔當自己的罪孽。」 \end{tabularx} \\ \\
[1ex]
\hline
\hline
\end{longtable}
$^{1}$弟兄姊妹平安.
很開心在今天是新年年十三.
可以和大家一齊來敬拜主.
並且去聽神自己的話語.
其實今年的年初.
我們就教會的講壇.
選了門費系列的崇拜.
是會看民數記.
記不記得去年.
當我們選了啟示錄的時候.
我們都很驚慌.
因為一卷這麼難的經文.
怎麼說和怎麼學習呢.
但神又很施恩.
當我們逐章逐章地看的時候.
我們又發覺.
對上帝的心意,美意了解了很多.
同樣今年看的民數記.
有多少人看過.
一點都不容易看.
要講壇講也不容易.
不過我們都很相信.
只要神的話語一解開.
就能夠使我們有亮光.
並且讓我們能夠經歷生命的更新.
到現在三月一號.
其實我們已經講了民數記頭四章.
如果你有在這裡聽的時候.
你會知道.
頭四章都好像是一些技術的問題.
上帝吩咐摩西要數點以色列文.
數完十二支牌之後就去數利未人.
來到這裡第五章到第六章.
原來又是一個段落.
是強調聖潔是贖神子民的一個特質.
我們一定要在裡面好好學習.
其實信耶穌這麼久.
我猜大家都很喜歡神同在.
有神在就沒有人怕.
因為有創天造地的真神.

$^{41}$愛我們到底的真神.
和我們一起.
以至我們能夠得著最好的保護.
帶領和供應.
但是神和你一起住又怎樣呢?.
有沒有想過在家裡.
上帝在你家中和你一起住.
其實我們信了耶穌的人.
上帝就在我們心中.
身體是聖靈的電.
怎樣能夠活出那份聖潔呢?.
民數記的第五章.
在一開始的時候.
很清晰的一個訊息.
告訴我們.
神住在他百姓的中間.
所以在當時數點完人數之後.
然後上帝是吩咐.
他們好好地處理營地.
免得污穢了他們的營.
因為神在其中居住.
所以他有一些人.
要send out在營的外邊.
所以在第五章裡.
我們會有三個層面去看.
當神住在我們中間的時候.
我們要知道.
祂帶來了一個聖潔的要求.
而且神住在我們中間的時候.
祂會要求我們關係上都要有公義.
第三個就是神同住.
祂會監察隱藏的忠誠.
都是一個頗豐富的內容.
為了方便大家理解經文.
我們先讀第一個段落.
頭四節的經文.
我們一起讀.
這裡其實講得很清楚.
雷同頭四章.
都是耶和華吩咐摩西.

$^{81}$然後摩西照做.
最特別的是在第四節.
以色列人就這樣行.
使他們出到營外.
耶和華怎樣吩咐摩西.
以色列人就這樣行.
原來神的民是需要聽命於上帝.
神同在帶來聖潔的要求.
所以摩西就按著神的吩咐.
就要將那些長大麻風.
有漏症摸了死屍的.
都要出營.
長大麻風.
這裡說麻風不知道會不會痊癒.
其實按原文這裡是說.
一些類似腺界的皮膚病.
一些會帶來傳染的病症.
就需要離開營地.
因為整群的以色列民.
在將要前行的日子裡.
他們要入迦南.
要攻取耶和華所賜的美地.
所以整個營地都要保持聖潔.
所以有病傳染危險的就要離開.
有漏症的.
什麼是漏症.
就是譬如女人的血漏.
男人的淋病.
胃症.
這些都需要離開.
還有摸過死屍的.
死亡和次生命源頭的上帝相對.
死屍在當時.
因為沒有冷藏.
很容易腐化.
甚至會引起很大嚴峻的疫情災病.
所以這些對全群會帶來影響的人.
需要離開營地.
這個是衛生的條例.
更加是上帝的要求.

$^{121}$為什麼.
因為聖潔的神住在營中.
聖潔的上帝和污穢和罪惡是相對的.
所以有上帝存在.
那些污穢就要清除.
其實由摩西上西乃山四十晝夜.
領受神的曉喻的時候.
就會知道原來上帝願意與祂的子民同在.
甚至耶穌降世.
道成為肉身.
在我們中間有恩典有真理.
所以神會與百姓同在.
神的同住其實會帶來一個要求.
就是必須聖潔.
因為上帝是聖潔的.
當我們中間沒有聖潔污穢的時候.
上帝就不與我們同在.
所以聖潔的課題非常之嚴峻.
在新約裡面.
保羅亦有這樣的教導.
他提醒我們.
「豈不知你們是神的殿.
神的靈住在你們裡頭嗎」.
今天神不是住在會幕裡面.
祂住在教會.
祂住在熟悉祂的人心裡面.
所以問題再不是問神究竟在哪裡.
而是神如果住在我們中間.
神在我們裡面的時候.
我們的生命是否配得祂的同在.
這個很值得我們反省.
其實聖潔是甚麼意思呢.
聖潔不是完美.
而是按聖經所說.
是被分別出來.
聖數記五章裡面所說.
不潔的人其實只是要求他們暫時出刑.
當他們的身體回復潔淨.
甚至是摸過死屍七天之後.
他們就可以回到刑裡面.

$^{161}$所以聖潔其實不是要毀滅人.
不要人與他同在.
而是他透過聖潔的要求.
要人恢復聖潔.
以至上帝能夠與他同在.
喜約翰一書一章九節這樣提醒我們.
我們一起讀.
我們若認自己的罪.
臣是信實的,是公義的.
必要赦免我們的罪.
洗淨我們一切的不義.
原來我們真的不是聖潔的屬性.
我們常常會犯錯.
上帝給了一個出路.
就是當我們肯認罪肯悔改的時候.
上帝必然赦免.
又讓我們能夠忘記背後.
繼續向前跟隨神.
神要的潔淨不是趕俗.
神是要人恢復而不是棄絕.
好像我們幾年前遇到的疫情.
染疫的人要怎樣?.
要被隔離.
隔離那個是不是我們看不起.
不要那個染疫的人?.
不是,而是保護他.
可以能夠快快康復.
亦都保護全群.
不會被疫病污染影響.
聖潔不是完美.
聖潔是分別出來.
這個理解很重要.
如果聖潔等於完美.
就糟糕了.
為甚麼糟糕?.
因為我們人都是軟弱.
當我以為聖潔的要求是完美的時候.
我就會裝模作樣.
我會裝作自己很勁虔.
但其實真實的狀況.

$^{201}$上帝是會知道的.
很容易就變得假冒為善.
還有我們會犯錯.
我們會跌倒.
當真的犯錯.
得罪上帝的時候.
那份內疚,罪疚.
又會叫我們不能抬起頭.
而且當我們以為聖潔等於完美.
沒有錯失的話.
就會用很多外在的表現.
來衡量屬靈的程度.
人與人的相處.
就會有很多的屬靈比較.
其實這些都是不健康的.
根據《羅馬書》三章二十三節.
它將我們的狀況描述得很清楚.
世人都犯了罪.
虧缺了神的榮耀.
所以如果聖潔等於完美.
根本就沒有人能夠成為聖潔.
但是聖潔其實是一個甚麼意思呢?.
聖潔是被分別出來.
歸屬於神.
單單為神所用.
聖殿是要聖潔的.
為甚麼?.
不是因為它完美.
而是因為它屬於神.
祭司是要聖潔的.
不是因為他們沒有軟弱.
而是因為他們是被分別出來.
單單為神所用.
那麼基督徒呢?.
基督徒也要聖潔.
為甚麼?.
因為我們是屬於神.
根據《彼得前書》二章九節.
有一個很清楚的解說.
有一個身份的改變.

$^{241}$以前我們在神的恩典裡無份.
因為我們不是血緣.
是神的子民.
但是我們因著耶穌基督.
我們因著肯接受.
承認自己的罪.
求主赦免.
我們就成為了被揀選的一群.
而且我們可以成為.
將人帶到神面前的祭司.
有君王的尊榮.
並且因著神的同在.
我們是聖潔的國度.
去到哪裡就應該神的國.
臨到哪裡.
我們是屬神的子民.
是上帝帶領的特別的一群.
有存在的目的.
是叫我們宣揚那召你們出黑暗.
入奇妙光明者的彼得.
所以我們的存在是要顯示.
我們是屬於神的.
我們要聖潔.
是因為屬神的人.
可以跟神的聖情有份.
不是我們本身就聖潔.
不過因著我們相信上帝.
我們可以得起神的聖情裡的有份.
彼得後書一章四節.
我們又讀一下.
3. 神的神神已將一切.
關乎生命和虔敬的事賜給我們.
皆因我們認識那用自己榮耀和美德.
召我們的主.
因此祂又將有寶貴.
又極大的應許賜給我們.
叫我們既脫離世上.
從情欲來的敗壞.
就得與神的聖情有份.
我們真的會有很多的敗壞.

$^{281}$不過神已經將祂一切.
關乎生命和虔敬的事賜給我們.
叫我們能夠清楚認識.
那個用榮耀和美德.
召我們的主.
在跟隨祂的過程裡.
我們又得著寶貴又極大的應許.
以致就算我們原來多腐敗也好.
我們也可以在神的聖情裡有份.
是一個可以轉化.
可以追求更加像上帝的角色和身份.
神住在我們裡面.
我們要自潔.
以致神不會因為我們骯髒而離開我們.
只是這個段落已經給我們.
有一些很深刻的提醒.
今天是聖餐主日.
也是很適合我們去看這個經文.
講到那些能夠污染人心.
影響人生命的要離開營地.
今天的經文提醒我們.
我們其實會不會真的已經與罪一切.
已經習慣了.
以致我們有時候都沒有留意到.
自己內裡究竟是乾淨還是骯髒.
我們有沒有一些脾氣,習慣,錯誤.
或者一些隱藏的軟弱.
長期都在我們裡面沒有處理.
而這些事情其實是阻隔著.
我們與上帝和我們與人的關係.
神如果真的住在我們裡面.
住在我們中間.
我們就應該將那些應該潔淨的東西.
趕離我們的身心靈.
以致我們能夠與主的聖潔有份.
要知道聖潔是神與主的自然結果.
當我裡面上帝佔的位置越多的時候.
我們就越來越能夠活出主的榮美.
叫我能夠分別為聖.
當我裡面肯自潔.

$^{321}$我肯讓神來掌管的時候.
我就能夠更多的分別為聖.
更能夠為主捨棄.
更能夠為主堅持.
將所虧負人的如數賠償.
另外加上五分之一也歸於所虧負的人.
那人若沒有親屬可受所賠還的.
那所賠還的就要歸於服侍耶和華的祭司.
至於那為他贖罪的公羊在外.
以色列人一切的聖物中.
所奉給祭司的舉祭.
都要歸於祭司.
各人所分別為聖的物.
無論是什麼都要歸給祭司.
得罪的時候其實某程度都是干犯耶和華得罪神.
我們有時候常常以為.
其實我們只要和別人不和.
那些事情都很小事.
人和人相處就是這樣.
但是原來當我們每一次去欺負別人.
欺騙別人.
侵佔 傷害 苦待.
甚至有霸凌的時候.
都是對神的公義一個很大的冒犯.
也都會成為整個神群體一個玷污了的見證.
所以當人和人出現問題的時候不處理.
那個是一個很嚴重的狀況.
會令到按聖經說.
我們每一個屬神的人.
都在基督的身體裡面互為肢體.
當肢體和肢體之間不能夠復和.
不能夠相通的時候.
基督的身體會怎樣?.
好像塞血管一樣.
這裡不通 那裡不通.
然後基督的身體的運作.
就會成為一個很大的影響.
所以神在這裡吩咐.
虧負人的要有行動的悔改.
不單止口講悔改.

$^{361}$更加要有行動來顯示悔改的心.
要承認所犯的罪.
如果虧欠人的.
要如數賠償.
並且因著你之錯.
要加上五分之一.
額外表達那個補償.
當那些受害人如果已經不在了.
都是虧負他的.
我們那個賠償都依然要歸耶和華.
以致得罪神的事件.
能夠得著處理.
亦都讓人知道.
你是真正來悔改.
新約裡面有一個很好的例子.
路加福音十九章.
講到薩蓋.
薩蓋是甚麼水?.
薩蓋是瑞麗.
以錢為他人生最重要的追求.
但他甚至因著要.
他拿了瑞麗這個身份來做.
然後他自肥.
甚至訛詐.
很得人憎.
但當他遇了耶穌之後.
怎樣?.
他說主啊.
我把所有的一半給窮人.
我若訛詐了誰.
就還他四倍.
耶穌說.
今天救恩到了者家.
因為他也是亞伯拉罕的子孫.
其實真正的悔改.
是需要帶來實際的行動的.
神住在我們中間.
他關心我們彼此的關係.
真信耶穌的人.
是願意為錯誤付出代價.

$^{401}$以至關係能夠修復.
其實也很真的.
當你能夠悔改.
給予賠償.
人與人之間的問題就處理了.
然後甚至對方會因著你這個行徑.
認識你那個樂意施恩的上帝.
所以馬太福音也是這樣提醒我們.
在我們祭壇上高勸禮物的時候.
如果想起有弟兄向你懷怨.
就要將禮物留在壇前.
先去和弟兄和好.
然後來獻禮物.
因為你不處理對人的虧欠.
你獻的禮物其實等於沒有獻.
所以神也提醒我們.
快點好好想想.
究竟有沒有一些應該承認的錯誤.
你沒有承認.
有沒有一些需要償還的虧欠.
你沒有處理.
有沒有一些需要修復的關係.
你當沒有存在.
這些都是上帝要我們好好清楚處理.
第三部分就講到十一到三十一節.
這段落很長.
我鼓勵你回去自己看.
因為裡面有很多描述.
就講到神會監察隱藏的忠誠.
《悉經》書裡面提到這個叫做苦水的試驗.
或者叫做疑恨的掃祭.
就是處理當時婚姻夫婦關係上出現問題的時候.
然後上帝為他的子民設立一個處理的方案.
是什麼狀況呢.
就是妻子犯了姦淫的罪.
又或者她沒有犯到姦淫的罪.
但是丈夫起了疑恨之心.
為什麼這個不能處理呢.
因為以往如果犯姦淫有證人的時候.
無論男女都要處死.

$^{441}$但是沒有證據沒有證人.
不過丈夫有懷疑.
卻很憤恨處理不到.
然後這個試驗就是疑恨的丈夫.
要將妻子和祭物帶到祭司面前.
求祭司幫助處理事件.
其實這個安排對當代的女性來說.
是一個很大的保護機制.
我們現在香港男女較為平等.
不是男人說什麼就行.
但是當代的女性當她們被懷疑有淫行.
而又沒有證據的時候.
其實男人對她可以出現很多的折磨.
甚至很多的方法.
看聲書說古代的人.
當這些婦人被懷疑的時候.
有些人會要求她跳入河裡.
不死就是無罪.
用手要抬那些辣到不得了的東西.
你不受傷就是無罪.
又或者要你手要浸入一些熱水.
能夠不受傷就是無罪.
你說這些行不行.
其實真的很殘忍.
所以上帝設立了一個方案.
就是要丈夫如果真的有這個疑恨.
你其實都要獻祭.
因為你那個懷恨的心不對.
還不知道是不是你就生氣了.
所以他要獻授祭.
然後將太太帶到祭司面前.
祭司用聖水下了塵土.
然後再將一個女人要讀的奏座浸在水中.
然後就要那個女人棄世.
當她如果真的有淫行的時候.
她的大腿會瘦 肚子會脹.
但是如果她是無罪的話.
她就會懷孕.
然後上帝會為她洗脫罪名.
於是這個可以說是對當時一個無辜人的一個保護.

$^{481}$甚至也是一個能夠將罪突顯出來.
不隱藏的一個可行的方案.
當然看這裡的時候.
我都不列那些儀式.
因為我們真的不是在這裡學.
那個怎樣做那個規矩.
因為那個我們要掌握那個正義和原則.
在沒有證據的情況下.
上帝要求我們不用操控和不安.
來處理那個不能夠處理的狀況.
而是用神的方法 由神去判斷.
其實不單止在婚姻裡面.
甚至在職場上.
有些時候真的犯錯.
然後隱瞞.
但是上帝他是監察人心的.
希伯來書四章十三節.
這樣形容我們的上帝.
上帝什麼都知道.
所以我們以為犯罪的那個.
可能以為我們可以隱瞞配偶.
隱瞞同事 公司 教會.
但卻不能夠隱藏上帝.
如果你是被冤枉的.
你可能被懷疑.
你真的沒有做過.
但是神會知道 神會知道.
所以在過程裡面.
我們都是要仰望上帝.
彼得前書二章二十三節.
將到主耶穌被冤枉的時候.
他怎樣處理呢.
來去擠在我們的面前.
我們讀一下.
他被罵不還口.
受害不說威嚇的話.
只將自己交託那按公義審判人的主.
主耶穌也是這樣.
當然要澄清的要澄清.
要解說的要解說.

$^{521}$但是當沒有辦法來處理的時候.
就唯有交給主.
究竟我們可以怎樣實踐呢.
就著婚姻和教會兩個層面.
我有些這樣的整理.
其實在婚姻裡面會非常複雜.
上帝將婚姻設立.
叫男女由二成為一.
其實是一個愛的詞語.
是叫我們在那個關係裡面.
能夠得著保養固息.
和能夠彼此來配搭更完美的生活.
所以誠實是非常重要.
也都彼此要設立界線.
在裡面不是誰控制誰.
也都不適宜在猜測來懷疑對方.
所以如果真的在婚姻裡面.
我們是需要好好地建立.
你要去經營的.
如果真的出問題的時候.
其實需要選擇一些適合的時間.
切實地去溝通.
大家就著狀況.
來設定一些可行的界線和方案.
看看能不能夠處理到那個疑慮和問題.
萬一真的不行.
對方都不肯來談的時候.
其實就要尋求牧者的幫助.
不要在這裡硬撐或者在這裡啞忍.
至於教會群體的操練方面.
因為今天要說這麼多東西.
所以全部都是簡單版.
但是真的要做的其實是非常複雜.
我只是將原則說出來.
如果遇到有一些.
又是在這裡傳著傳著的留言.
在這裡說著說著的東西.
那又要怎麼樣呢.
就要停止留言審判.
學習以真理來彼此守護.

$^{561}$那些傳來傳去不知道真假的東西.
不要傳了.
首先如果知道教會裡面有這樣的情況.
要先為事情祈禱.
不要急於下結論.
如果我不是當事人.
我又不是合適去處理的那個.
我就不將這件事傳來傳去.
中國人都有句叫什麼.
謠言止於自者萬.
未證實是事實就不可以傳.
但是當我有責任的時候.
我就要按聖經教導的真理.
來去關心去處理.
目的是為了要挽回.
不是要定罪.
以至叫教會的群體沒受污染.
不會有一些流腫.
流到基督的身體裡面.
說就容易.
除非我們按照上帝教的方法.
用同一個標準來代人處事.
否則基督的身體都可以雜亂無章.
所以如果要上帝的聖潔.
在我們生命和教會群體裡面.
能夠流露出來的話.
我們其實需要上帝.
真的很大的憐憫提醒和幫助.
民數記在五章這裡.
有三個的呼籲.
神住在我們中間.
我們需要活出聖潔.
不要被污穢長期躲藏.
有一些需要處理的事.
是必須處理.
神重視群體的關係.
在人際關係上高.
要實踐公義.
是要修復虧欠.
不要欠別人.

$^{601}$我們教會有些時候.
難的就是認過錯.
認過錯就算了.
對受害的其實都有很大的艱難.
怎樣能夠做一些補償.
大家訂立一些對的相處之道.
和一些的界線.
都是非常之重要.
還有就是要拒絕留言和埋堆.
不要成為割舌的一份子.
要用愛和真理.
和一些合宜的方法去守護彼此.
因為我們的言行.
我們的舉措.
都和上帝的明星掛鉤.
我們行得不好.
除了我們被人指指點點.
自己承受罪的結果之外.
我們是污衊了上帝的明星.
所以需要我們非常之謹守.
我看這段經文的時候.
很受第一個段落.
以色列人的反應來感動.
摩西說什麼.
以色列人就跟著進行.
今天是新年年十三.
又是聖餐崇拜.
弟兄姊妹.
藉著今天的經文提醒我們.
要醒覺,醒察,除咎,迎新.
做一個合乎上帝喜悅的人.
以致我們不單止享受到神的同在.
我們更加能夠成為祂活活的見證.
剛才林牧師有提到.
世界變動得很複雜.
我們更加要好好地為上帝.
做上帝喜悅我們做的事情.
當我們能夠好好潔淨自己.
任主使用的時候.
神就在我們中間顯為大.

$^{641}$願以今天的經文.
能夠成為我們的提醒和幫助.
\newpage



\section{馬太福音 17:24-27}
\label{sec:031744}
\textbf{變臉後傳 II  掘銀交稅}~~~~梁永善牧師
\newline
\newline
連結: \href{http://www.church.com.hk/acms/content.asp?site=cdc\&op=show\&type=product\&code=031744\&layout=sermon}{\texttt{031744}} ~~~~ 語音日期: 2026-03-04
\newline
\newline
\hyperref[sec:031743]{\small{< < < PREV SERMON < < <}}
~
\hyperref[sec:index_chronic]{\small{[返順時目]}}
~
\hyperref[sec:031745]{\small{> > > NEXT SERMON > > >}}
\newline
\newline
馬太福音 17:24-27
\newline
\begin{longtable}{cl}
\hline
\hline
章節 & 經文 (和合本修訂版)\\
\hline
17:24 & \begin{tabularx}{0.7\textwidth}{X} 他們到了迦百農，收聖殿稅的人來見彼得，說：「你們的老師不納聖殿稅嗎？」 \end{tabularx} \\ \\ \relax
17:25 & \begin{tabularx}{0.7\textwidth}{X} 彼得說：「納。」他進了屋子，耶穌先對他說：「西門，你的意見如何？世上的君王向誰徵收關稅或丁稅？是向自己的兒子呢？還是向外人呢？」 \end{tabularx} \\ \\ \relax
17:26 & \begin{tabularx}{0.7\textwidth}{X} 彼得說：「是向外人。」耶穌對他說：「既然如此，兒子就可以免了。 \end{tabularx} \\ \\ \relax
17:27 & \begin{tabularx}{0.7\textwidth}{X} 但恐怕觸犯他們，你往海邊去釣魚，把先釣上來的魚拿起來，開了牠的口，會發現一個司塔特，可以拿去給他們，作你我的稅錢。」 \end{tabularx} \\ \\
[1ex]
\hline
\hline
\end{longtable}
$^{1}$今天會和大家去看「撈銀交稅」.
所以你會看到這個圖.
香港的弟兄姊妹一定會有熟悉的感覺.
但又不是很喜歡看到的 為什麼呢?.
我們收到綠色的大信封就知道.
是要填報稅表 是要交稅.
我們今天差不多來到登山變象的最後一次了.
我們看到耶穌基督登山變象.
就是讓門徒看到祂的真相.
就是說祂很快就要面對十字架的苦楚.
要受死 但是這個只是暫時.
祂真正的榮耀的本相就在登山變象那裡.
給你看到 但是記不記得我們講過上山.
他們說在這裡真好.
有時候我們有些地方我們覺得留在這裡就很好了.
但始終要回到自己熟悉的地方.
他們上山就要下山 要面對群眾.
上次我們看到當他們一下山的時候.
門徒就已經面對挑戰了 為什麼呢?.
他們不能夠趕鬼.
就被人質疑 你都沒有能力的.
你跟隨那個是不是真的行的?.
因為他們的信心少.
上次我們說的心弱力微 信心微弱.
你的力量自然就會減少.
最重要就是要好好的依靠神.
親近祂 靠著祂給我們有力量 然後去戰勝.
好了 挑戰完門徒 現在就去挑戰耶穌了.
就是有些人特意來挑戰耶穌的身份和地位.
就是問耶穌應不應該勒釘誰.
在這件事上給我們有什麼教導呢?.
我們先看這段經文.
《馬太福音》第十七章二十四到二十七節.
我們先從二方堂讀出來.
「到了加伯隆 有收釘誰的人來見彼得說.
你們的先生不勒釘誰麼?.
彼得說 勒 他進了屋子.
耶穌先向他說:.
「西門 你的意思如何?.
世上的君王向誰徵收官稅,釘誰?.

$^{41}$是向自己的兒子呢?是向外人嗎?」.
彼得說:「是向外人」.
耶穌說:「既然如此 兒子就可以免稅了.
但恐怕觸犯他們.
你且往海邊去釣魚.
把先釣上來的魚拿起來.
開了他的口 必得一塊錢.
可以拿去給他們作你外的稅銀」.
很簡單 只有四節經文.
那些人來問彼得.
「你的夫子勒不勒釘誰?」.
彼得說 「勒」 很快的.
然後去問耶穌 耶穌就問他.
究竟皇室成員是否需要納稅呢?.
應該不用納的.
但是耶穌說:「好吧.
免得他們有問題的時候.
你去外面釣一條魚.
然後在魚口裡挖一塊錢出來.
就可以交稅了」.
很短的經文.
但裡面有很豐富的屬靈的教導.
這段經文在《馬太福音》裡面.
記載了三個很獨特的神蹟.
你記得《四福音書》裡面.
只有一個神蹟是《四卷福音書》都記載的.
就是五餅二魚給幾千人吃飽.
但是《馬太福音》記載了三個獨特的神蹟.
這是其中的一個.
我們要先了解一下背景.
我們這段經文很容易.
將另一段經文混淆了.
甚至很多人就拿著這段經文來講.
就是我們要去信服掌權者.
因為該殺的物就歸給該殺.
神的物就歸給神.
所以要怎樣呢?.
就是要政教分離.
教會,基督徒千萬不要講政治.
就拿這一段來帶進去講.

$^{81}$我們看一看這段經文.
其實兩者是有分別的.
《馬太福音》22章15到22節.
當時法利賽人出去商議.
怎樣就著耶穌的話陷害他.
就打發他們的門徒和希律黨的人去見耶穌說.
「夫子,我們知道你是誠實人.
並且誠誠實實傳神的道.
什麼人你都不順情面.
因為你不看人的外貌.
請告訴我們你的意見如何.
納粹給該殺可以不可以」.
耶穌看出他們的惡意就說.
「假冒為善的人啊.
為什麼試探我.
拿一個上水的錢給我看.
他們就拿一個銀錢來給他.
耶穌說:這匠和這豪是誰的.
他們說:是該殺的.
耶穌說:這讓該殺的物當歸該殺.
神的物當歸給神.
他們聽見就稀奇.
離開他走了」.
這裡是很清楚講到.
是有兩班人就是法利賽人.
和希律黨的人.
聯手去對付耶穌.
希律黨的人是甚麼呢?.
都是猶太人.
不過他們是親當時的羅馬政府.
即是建制派.
其實和法利賽人原本是兩者互相之間.
是敵對對方的.
法利賽人有很強的民族主義.
所以一向都不喜歡被羅馬政府統治.
但現在有共同的敵人就是耶穌.
於是他們聯手去對付耶穌.
問耶穌是否應該納稅給該殺.
你留意這個稅字和丁稅的稅字.
在原文裡面是兩個不同的字.

$^{121}$納稅給該殺是當時羅馬政府統治下.
人民應該盡的責任.
正如我們在英國也好,在香港也好.
要納稅給當地政府.
你做事就要納薪俸稅.
在香港如果你有物業就要納物業稅.
這是必然的.
如果不納稅就觸犯了法例.
但今次我們看經文.
就是問耶穌是否會納丁稅.
甚麼是丁稅呢?.
如果按照《出埃及記》30章.
11到第16節的記載.
當時在抗疫的時候.
就做出人口普查.
好像香港也做人口普查.
原來20歲以上的男子.
每人就要交半射赫勒的銀紙.
好 剛才我講《出埃及記》30章11到第16節.
我只是抽14,16兩節.
就是這幾節出來看.
我們請Yvonne跟我們來分享.
這裡給我們看到原來他們是要納這個稅項.
這個稅項是做甚麼的呢?.
他們那時是出埃及的 是不是?.
那是沒有政府的.
這樣的時候當然不會有甚麼薪俸稅,物業稅.
但原來每一個20歲以外的人.
都要納這個釘稅.
這個錢是用來做甚麼的呢?.
就是維持匯募的費用.
而這個稅項你有沒有留意到.
連貧窮人都不可以不給的.
他說付出的不需要多出.
但貧窮的都不可以少出.
每人要給半射赫勒.
在主耶穌的年代.
就是20歲到50歲的人.
猶太的藍丁無論是居住在耶路撒冷.
猶太地或者在海外.

$^{161}$他們都需要來交這個釘稅.
就是維持當時聖殿的運作和經費.
往往他們如果要去逾節的時候.
就要去耶路撒冷首節.
他們就可以親自來繳交.
在逾節之前一個月.
巴勒斯坦地和外地各處都有專人.
負責收這個聖殿稅.
所以這一次.
詢問耶穌.
那些人絕對不是羅馬的官長.
所以沒有那些希臘黨的人.
是以色列人的官長所委派的人.
所以釘稅又稱為聖殿稅.
或者叫做贖命錢.
為什麼呢?.
是將神第十五節.
他為贖生命的禮物獻給神.
是含有感恩.
因為神拯救我們.
去呈現和承擔責任的意思.
不要忘記昔日的會幕.
聖殿是一個原來要花費很多很多錢的地方.
為什麼呢?.
他們每天清晨和黃昏的時候都要獻祭.
晚上要獻晚祭.
要獻上什麼呢?.
每天要獻上至少兩隻.
因為每一次獻祭都要獻一隻.
是一歲的羊羔.
你要留意.
要買的.
要錢的.
不便宜的.
同時要獻上酒.
麵和油.
並且是要燒香的.
而祭司穿的衣服.
那些掛飾.
是很貴重的.

$^{201}$祭司的袍是絕對不簡單的.
如果是大祭司更加要用金線去做的.
所以你要維持會幕和聖殿的運作.
是需要很多很多錢的.
因此的緣故.
就是要民間的人去獻丁歲.
如果他們不給稅的時候.
他們那些不是羅馬人.
是猶太人的官長.
即是民間的保長.
是可以扣留他們的財物的.
但是這個稅項要記住.
不是政府.
不是羅馬政府要他們交的.
所以某個程度如果他們不交.
是無奈的.
我想問問大家.
現在仍然是新蒸頭.
新年.
那些長輩是不是一定要派利是?.
這個不是法律規定的嗎?.
不派就要抓的嗎?.
沒有的啊.
如果長輩不派利是.
你最初心中不是這樣.
你不可以去追收他的.
這個不是必然的事來的.
你新年去拜年.
不一定要有手信的.
何必要有手信呢?.
但是我們習慣是會給的.
所以這個數額.
不是非常巨大的.
這個稅項.
窮人都不可以少出.
是大概他們.
平均兩至四天的工資.
請你記住.
一年只交一次.
平均是兩至四天的工資.

$^{241}$所以是可以承擔的.
如果人人都肯交的時候.
就能夠集極成球了.
就可以令到聖殿會幕有好的運作.
但是很可惜.
有不少人不肯承擔這個責任.
今天在教會裡.
往往都是這樣的.
很多人.
他們會參加教會的聚會.
甚至有很多要求.
有些人就說消費主義.
就要求這個要求那個.
我說不是啊.
不是消費主義.
因為消費主義.
是你給多少錢就期望有多少享受.
你去文華酒店吃東西.
喝一杯咖啡.
和去文華餐廳喝一杯咖啡.
價錢是很不同的.
但是服務都很不同的.
但是很多基督徒.
他只是期望教會做這個做那個.
但是他連基本的奉獻都沒有.
在英國現在的環境裡面.
都出現了這些問題.
聚會人多.
但是奉獻人少.
可能我們的經濟都不是那麼理想.
我們撕下撕下.
慢慢地錢就越來越少了.
我們就缺乏了這個信心的擺上.
甚至英國很多的座堂.
譬如教堂.
現在都會租給抗黨人的教會.
因為有一定的收入.
原來以前是禮拜堂.
已經運作不了.
運作要很大的經費.

$^{281}$所以今天我們基督徒.
我們都要反省.
我們看看究竟.
納還是不納呢?.
收稅的人就向彼得.
主耶穌會不會納這個稅項呢?.
《馬達方德十章》第24節.
到了加伯倫.
有收丁稅的人來見彼得說.
你們的先生不納丁稅嗎?.
你有沒有留意到.
他們沒有直接問主耶穌.
或者是對主的尊重.
而不好意思直接問祂.
很簡單 譬如婚禮.
一家大小來喝酒.
簽完名就進去了.
沒有放下禮金.
接待的人都知道好像沒有給人情.
但又不好意思直接問一家之主.
可能會問他認識的一家之主的孩子.
你爸爸或者媽媽有沒有忘記了給禮.
有沒有給人情.
這個很重要.
他們不好意思直接問.
於是就問他.
甚至教會的牧者有甚麼問題.
可能他們不會直接問牧者.
就會問那些執事.
執事會有很大壓力.
而彼得是甚麼?.
是耶穌的首席門徒.
所以就直接去問彼得.
但有另一個可能性.
是甚麼呢?.
因為知道彼得是一個口直心快的人.
經常衝口而出.
說話有時不經大腦.
記不記得彼得見到主耶穌行海.
就跟耶穌說.

$^{321}$主啊 如果是你.
我也走下來吧.
走下來幹甚麼?.
沒有意思的.
當耶穌在登山前主我們看了.
耶穌說要為人受死的時候.
彼得馬上拉著耶穌.
叫主萬萬不可.
怎麼可以做這件事呢?.
所以留意彼得有時是沒有心思熟慮的.
所以就容易從這些人口中.
找到一些把柄去攻擊.
有時很真實地說.
我也是一個腦界轉得快的人.
有好有不好.
所以很容易會得罪人.
有時沒有想清楚就說出來.
有時都很笨拙的.
所以彼得就說.
如果說不用想了.
就拿到把柄了.
就控告耶穌不尊重神.
不肯承擔責任.
沒有守神的律法.
請你記住.
很多時一些充滿了惡毒.
詭詐的言語.
和一些很謙虛.
很恭敬的話.
說出來其實差不多一樣.
很可能他說得很好.
但背後就是想陷害你.
想拿你把柄.
所以我們要細心分析.
正如剛才我們讀過的.
問耶穌應不應該立那個誓項的時候.
是怎樣呢?.
我們再看回《瑪達福音》.
第15到第17節.
我們剛才看過那一段經文.

$^{361}$就在《瑪達福音》第22章裡面.
我們一起聽聽Yvonne讀出來.
當時伯利賽人出去商議.
怎樣就著耶穌的話陷害他.
就打發他們的門徒.
和希律黨的人去見耶穌說.
夫子 我們知道你是誠實人.
並且誠誠實實傳神的道.
什麼人你都不信情面.
因為你不看人的外貌.
請告訴我們你的意見如何.
立誓給哥剎可以不可以?.
你有沒有看到他們問耶穌應不應該立誓.
前面說了很多話.
是讚賞,欣賞耶穌的.
夫子 我們知道你是一個誠實的人.
你一定是誠誠實實傳神的道.
甚至你不會看人的情面.
你不會理解他怕某些人.
你告訴我 應不應該立誓呢?.
是不是說得很好?.
但是前面第10節說什麼呢?.
他們去問耶穌是想陷害耶穌.
剛剛在過去這個星期說到.
車公廟求籤.
今次求了籤.
最後那句是.
「雖向人中類不仁」.
這句話是很認同的.
很多人都說車公廟很靈的.
那些人表面好像很仁義道德.
但是你不要忘記.
其實這些人很可能是心懷不軌的.
給就真的給.
他很爽快的馬上說.
立 怎麼會不立呢?.
但是他事先沒有徵求耶穌的回應.
他是不是越早帶袍呢?.
甚至沒有大沒有小呢?.
我很相信彼得.

$^{401}$他除了衝動之外.
其實對耶穌有一個很清楚的了解.
他知道主耶穌凡是人應該盡的本分.
應該承擔的責任.
主耶穌一定不會逃避.
不去做的.
而一切合理和有益於神的事.
主耶穌一定會樂意地帶入去承擔.
反之.
如果一些有礙於福音的事.
或者有不合法的成份.
不道德的事.
不需要深入研究.
主耶穌一定不開心.
和不會做這些事.
很簡單舉個例.
當年我在名人堂的時候.
甚至後來我離開了.
也很好.
我屬靈的女兒.
後來幫我找工作.
是我的幹事.
那些人打電話來請我去廣道.
他從來不會問我.
牧師你去不去?行不行?.
不會的.
他只是看我的時間表.
他有了我的時間表.
那個時間可以.
他也不會理會遠近的.
他就會和我去接廣道.
為甚麼呢?.
他知道我怎樣.
我是很樂意在港台上去服侍人.
不論遠近.
不論人多人少.
只要那個時間配合到.
我就會去的.
甚至他會知道.
那個可能是我生日.

$^{441}$可能是我孩子的生日.
可能是我結婚週年紀念.
但是也有時.
他也不會問我.
也會和我接廣道.
為甚麼呢?.
他知道我一定會.
正如今年我的生日.
我的結婚週年紀念也是星期天.
我照樣會去事奉.
所以這個答案是對主耶穌的認識.
也都是合於律法.
是忠於神.
和民族的熱情.
所以彼得馬上說「拿」.
但是接著.
彼得說完之後.
他就要去屋裡.
尋求耶穌的確認.
我們不知道.
究竟耶穌是聽到.
在門口的對話.
還是主耶穌的預知.
誰知道彼得進去的時候.
還未說話的時候.
耶穌就採取主動.
去反問約瑟.
藉著這樣來教導他.
分析「拿」和「不拿」的證據.
或者原因.
我們看看第17章第25節.
彼得說「拿」.
他進了屋子.
耶穌先向他說.
「世門你的意思如何?.
世上的君王向誰徵收.
關稅經稅?是向自己的兒子呢?.
是向外人呢?」.
你留意.
君王你看英文聖經就知道.

$^{481}$是眾數.
即是指普遍的國家.
不只是指羅馬政府.
世上所有的君主.
甚至納稅.
關稅是甚麼呢?.
關稅就是地方上.
直接收取貨物的稅銀.
即是很簡單.
有入口稅.
有些國家.
有些東西入口的時候.
就要打稅.
很簡單舉個例子.
為何有些人會走私.
新的iPhone上國內呢?.
可能在國內.
買的時候就要打稅.
你走私上去.
就節省了稅銀.
又或者你要有銷售稅.
你去日本或某些地方.
買東西的時候.
就要付銷售稅.
這就叫關稅.
丁稅是甚麼呢?.
剛才已經說過.
是給中央庫房的人的人頭稅.
很簡單.
香港現在都有.
你出境的時候.
就要收離境稅.
機場建設稅.
如果你不離境.
你不坐飛機的時候.
就不需要收稅.
我曾經在永善論壇說過.
很多人北上.
又或者很多國內的人下來.
影響香港的消費.

$^{521}$如果你想阻止.
這個消費很簡單.
每一次離境.
收50元香港錢.
離開香港境.
在羅湖.
在其他海岸裡面.
馬上減少很多.
一家大小上去.
就很多了.
國內下來的又要收稅.
就不會那麼多窮遊的人下來.
政府肯不肯去做.
所以原來.
這個丁稅是人頭稅.
兒子.
都是中數.
指皇室成員.
或者.
皇室的一些家屬.
彼得的回答是.
是很清晰的.
第十七章.
二十六節的上半節.
彼得說:是向外人.
耶穌說:既然如此.
兒子就可以免稅了.
彼得很清楚說.
當然不是徵收皇室成員的人.
他們不用納稅.
向外人去徵收.
斯特指紋納稅是什麼?.
是供養皇室成員.
讓他們過著安息.
超級享受的生活.
讓他們過著.
豪華的享受.
而皇室成員.
就有特別的身份.
皇子公主.

$^{561}$他們就不需要納稅了.
所以主耶穌很清楚表明自己的身份.
也是回應了.
登山變象的意義.
祂是至高神的兒子.
祂本身就是神.
為什麼要納這個徵稅呢?.
徵稅就是給聖殿運作的.
聖殿就是來事奉神.
祂本身就是神.
本身就是神的兒子.
所以不用納稅.
這個答案是正確的.
所以你有沒有留意到.
最近英國是怎樣的.
上星期你會看到.
安德魯皇子被控告.
但你其實留意到.
原來早就有鋪排了.
除去了皇室成員的銜頭.
貶為庶民.
即是成為一個普通人.
原來你看看.
因為皇室成員有豁免權.
是可以不用被起訴的.
正如一些外交人員.
都可以有豁免權.
好了.
這個時候.
所以安德魯被貶為庶民.
然後不讓他再在.
溫莎這一區住了.
做了這些步驟.
然後就怎樣.
現在就抓他 告他.
所以你留意到一定是.
皇室成員和政府方面.
是已經有了他的罪證.
但礙於是皇室身份.
所以就要先.

$^{601}$將他的身份除去.
即是我們講的.
「撕柴」.
這個就是這個意思.
耶穌基督本身就是神.
所以就不需要納這個稅項.
但是.
主耶穌.
為了人的著想.
雖然是神的兒子.
不需要納這個的丁稅.
但是為了別人.
他甘心狂復於.
律法之下.
我們看看第17章的27節.
主耶穌說我根本就不需要納這個稅項.
但是.
恐怕觸犯他們.
不如你就去.
找一個銀來交稅.
如果我們看中文聖經的時候.
觸犯這句話.
哇 耶穌都怕了.
這班有勢力的人士.
向惡勢力屈服.
原來觸露這個字.
原來的意思是甚麼.
是叛逆.
原來的意思是甚麼.
是叛倒的意思.
意思絕對不是.
所謂怕這班有權力的人.
於是就算了.
原來意思是.
不想叛倒這班人.
他不想因為他的拒絕.
這個是正確的拒絕.
他如果不交這個丁稅.
但就會引起.
其他人的誤會.

$^{641}$以為他輕視了神的律法.
不遵守.
聖殿的規範.
亦都令到聖殿.
不能夠有這個收入.
如果有些人效法.
聖殿就不能夠好地運作.
所以.
為了維持聖殿的運作.
他就要立.
這個稅項.
主耶穌這樣做的時候.
很真是有些委屈.
而且.
是有些荒謬.
但主耶穌為我們樹立.
一個非常非常好的榜樣.
為了別人的著想.
正如.
他要洗禮.
施主耶穌說.
我怎樣可以和你洗禮呢?.
他知道耶穌是不需要洗禮的.
因為洗禮的時候.
的意義就是我要承認我是一個罪人.
所以叫悔改的洗禮.
但施主耶穌說.
我怎樣可以和你洗禮呢?.
耶穌說我要盡諸般的義.
我要做一個榜樣.
於是他就堅持.
然後叫施主耶穌和他洗禮.
是不是?.
他不堅持自己的權利和自由.
他甘心.
放棄自己的權利.
這個甘願.
就被這些人看為.
凡人.
即是普通人來的.

$^{681}$所以就要交稅.
他為了別人.
不想叛倒他們.
或許.
那些人可能是無知.
幼稚.
但為了他們的得救.
為了他們的成長.
他放棄了自己的權利.
你不要忘記.
主耶穌經常很嚴厲的斥責罪惡.
但他又極力的.
避免一些不合理的得罪人.
我們需要.
持守這個態度.
你現在看到.
特別在香港.
很多人都不可以.
發聲.
你正確的發聲.
可能就會受到迫害.
你留意到.
你為公義發聲.
但原來你就會受到影響.
原來.
你一說話的時候.
可能就面對這些勢力的打壓.
於是很多人就怎樣?.
懼怕.
就不再作聲.
但主耶穌絕對不是.
懼怕這些勢力.
只是不想叛逃.
這是每一個成熟的信徒.
應該有的表現.
在《羅莫書》15章1到7節.
有興趣大家可以上網.
聽我講一篇《道》.
這篇《道》講到.
我們為別人著想.

$^{721}$願意放下自己.
就叫「恃強扶弱」.
又或者《哥林多前書》第9章12到23節.
大家都很熟悉.
這段經文就是「凡事都可行,但不都有益處」.
「凡事都可行,但不都榮耀神」.
之前講到.
應不應該吃拜偶像的食物.
他都講過這篇《道》.
就是為榮耀神而作.
大家可以在網上.
聽到的.
另外一些經文.
特別提醒我們.
在《哥林多前書》第8章9到13節.
這裡說我們不要做絆腳石.
請二汪同學讀出來.
「只是你們要謹慎,恐怕你們這自由,形成了那軟弱人的絆腳石.若有人見你,這有知識的,在偶像的廟裡坐直,這人的良心若是軟弱,豈不放膽去吃那祭偶像之墓矇?.
因此,基督為他死的那軟弱弟兄,也就因你的知識沉淪了.你們這樣得罪弟兄們,傷了他們軟弱的良心,就是得罪基督..
所以,食物若叫我弟兄跌倒,我就永遠不吃肉,免得叫我弟兄跌倒了.」.
這個意思就是說其實有些東西是合理的 是沒有問題的.
但是不要做一些東西 免得給人有誤會.
很簡單 剛剛上個月有一位姊妹 不是名堂的 今天來英國玩.
我幫她預訂了酒店住 我見她的時候 我只是在大堂見她.
「牧師 我們有些東西拿給你的 是國際的 不如你跟我上來拿吧」.
「我不會上去你的房間的」 可能只是逗留一兩分鐘 隨時給人有誤會.
如果剛好在房間出來的時候 一男一女 還見到有些認識的人.
「為什麼牧師在酒店的房間出來」 那就有很大的誤會了.
我說「不要介意 不如你上去拿吧 我在大堂等你 你拿給我吧」.
其實是沒有問題的 神可以造證 但是不想叛道人.
今天在這個社會裡面 我們大家都以自我作為中心.
甚至以自己作為重點 只有他這麼多.
我行我素 我沒有錯就可以了 我們很流行說這些.
這個也是我自己在少年時代的時候的想法.
但是我們的信仰是教導我們學習放棄自己 成全別人.
正如很簡單的 不只是休息的問題.
譬如我們的事 譬如我今天的查經 其實是錄影的.
因為我有時候要離開英國幾天 但是我不想停的.
義務的 但是我都會去做 但是更加重要的是什麼呢?.
我們不要做一些事情 但是我覺得沒有問題的.

$^{761}$我有所堅持 但是會叛道人 於是怎麼樣呢?.
如果我做一些事情 叛道人是很大的一件事.
會受到神的懲罰的 所以我們要很小心.
好嗎?第十八章第七節.
「這世界有禍了 因為將人叛道.
叛道人的事是免不了的 但那叛道人的有禍了」.
好 再看《路加福音》第十七章的第一到第二節.
「那叛道人的有禍了.
就是把磨石鑽在這人的頸項上 雕在海裡.
頑強於他把這小子裡的一個叛道了」.
我剛才在這個禮拜說車公廟那些籤好像很靈.
我就會看基督徒應不應該求籤.
或者我們怎麼去解讀 為什麼這些籤這麼靈呢?.
我更加會說我們基督徒應該是不可以參與這些的.
原來聖經說 如果做了一些事情是影響到別人.
特別如果那個是你 涉及一些你個人的利益得失的事.
或者和一些罪惡無關的事上.
我們應不應該盡量的容讓別人 幫助別人.
叫福音不會受到障礙 而不要成為絆腳石呢?.
當然 當我們面對一些絕對的真理.
關乎信仰重大的事情 我們的主耶穌基督.
並不會介意使人絆倒的.
你要搞清楚 在彼得的書信裡面.
曾經和大家查過彼得的書信.
耶穌基督是什麼?是盤石 但也是絆腳石.
聖經裡面說耶穌是絆腳石.
有些人拒絕耶穌就絆倒了.
正如祂和稅女妓女一起分享信仰 一起吃飯.
被別人來責備 被別人質疑.
說為什麼會和這些罪人在一起.
耶穌甚至在安息裡面趕鬼 治病.
就被人說他犯了安息的規條.
但是主耶穌不會因為這些責備 這些批評.
而拒絕別人 我就不做了.
等安息過了以後 才去治病 才去趕鬼.
不會的 要繼續面對.
所以求神讓我們有智慧去分辨.
什麼時候應該堅持 什麼時候應該要放棄.
我很喜歡一段的圖文 一張很古老的卡.
那裡寫著.

$^{801}$Her law grants me the sanity to accept the things I cannot change, courage to change the things I can, and the wisdom to know the difference.
那中文是什麼呢?.
這張是我當年夾在聖經的卡.
中文翻譯就說 神啊 求你賜給我靈色.
去接受那些改變不了的事情.
給我有勇氣去改善那些可以改變得到的.
但最重要是給我有智慧去分辨兩者的分別.
所以求神要幫助我們.
當然我想說 要成熟的生命.
你才會懂得怎樣去放棄.
剛才我說 我少年時代我不管那些導師.
或者那些人對我的批評.
在團派裡面 我都無所謂的.
十三,四歲的時候回到團體.
就算是姐妹也好 我都會搭搭肩膀的.
大兄姐妹有什麼所謂呢?.
但現在我都不會了.
但那時候我會理他那麼多嗎?.
我們是大兄姐妹有什麼問題?.
我又沒有不軌企圖 神又知道.
對方都知道了 理會別人怎麼看呢?.
但現在我不會這樣想.
就是因為神給我很多的教導.
所以求神幫助我們.
好了 去到我們這裡了.
就是「撈銀交稅」.
主耶穌說 好了 我們就納稅給他.
那耶穌怎樣來納稅呢?.
主耶穌吩咐彼得去海邊釣魚.
將釣到第一條魚.
從他的口裡面挖一個銀出來.
他的口裡面有一個銀.
就用來交稅.
縱使這件經文沒有再記載彼得怎樣去做.
但我們相信彼得一定會跟著主耶穌的指示去做.
也挖到這個銀 完成了這個使命.
但這件事值得我們去思考.
第一就是 為什麼主耶穌連交稅的錢都沒有呢?.
主耶穌在世上過著貧寒的生活.
這句經文是很特別的.

$^{841}$好 我們一起來看.
哥林多號書第八章第九節.
「他本是富足卻為我們變成貧窮」.
其實下面還有一句的 不過沒有打出來.
「我們本是貧窮.
但因為耶穌基督的貧窮而我們成為富足」.
是不是這句說話很玄呢?.
不是很明白的.
他本是富足 因為我們成為貧窮.
我們本來是貧窮的.
因為他的貧窮而我們成為富足.
我一講大家就會明白了.
每一個父母和子女說.
我原本是富足的.
我本來可以有很好的享受.
但因為生了你 結果我就成為貧窮.
你生出來 你一毛錢.
一個仙都沒有帶來這個世界.
但現在你多富足.
特別新年你賭多少利是?.
你賭這麼多利是 我付出了很多.
我要給很多利是 我貧窮 你就富足.
主耶穌因為我們寫在天上的榮耀而成為貧窮.
正正因為耶穌這樣的時候我們接受他.
於是我們成為富足 是不是?.
這裡就是主耶穌基督甘心的氣泉.
我們會不會時常埋怨自己不富足呢?.
但我很相信今天我們參與查經的.
或者聽這篇信息的弟兄姊妹.
一定是比耶穌當日的生活是富足得多的.
我們的物質享受是高過主耶穌的.
主縱使並不是沒有錢交稅.
但是他就用這個神跡取得的錢表明.
神有另外的安排.
他不需要自己親自來給錢.
間接引證了主耶穌原本就是神.
不需要來付這個聖殿稅.
是藉著餘口得這個稅.
當然 記住 並不是鼓勵我們不勞而獲.
只是一件很獨特的事情而已.

$^{881}$另外 也是給彼得一個考驗.
在山上的監聽到.
「這是我的愛子 我所喜悅的」.
然後他記不記得.
「你們要聽從他」.
這裡就是他要面對這件事.
主耶穌能夠將無變為有.
甚至可以點石成金.
但是為什麼要給他釣魚呢?.
這個並不是網魚.
網魚是一堆魚網出來.
這個是釣魚.
這個字在聖經裡只出現過一次.
不要忘記彼得是漁夫.
他以前是打魚的.
神就用他的專長去完成這個神跡.
所以今天弟兄姊妹.
你有什麼專長 有什麼特性.
你會不會用在神的事工上呢?.
問題是我們願不願意為主來獻上.
可能你未信主之前你已經學琴.
你已經懂得彈琴了.
那你信了主 你願不願意為教會做詩琴呢?.
你很喜歡插花的.
以前在家裡插一兩盆花 每個星期換.
但是你願意為教會來帶去怎樣呢?.
來帶去插花呢?.
我從小到大都喜歡說話.
神使用我 讓我喜歡來帶去講道.
釣魚還要怎樣呢?.
要有很大的耐性.
我到現在都沒有人試過釣魚.
直至這樣 也是彰顯神的權能.
要釣到一條魚.
還在第一條魚的口裡挖一個銀出來.
這個銀幣的價值原來是一射克納.
是夠兩個人來交稅的.
這個難度很高.
但是事情就是這樣完成了.
因此更加顯出基督的榮耀和祂的能力.

$^{921}$也都看到信仰裡面的那個吊詭.
主耶穌是榮耀的主.
但是可能祂身無分文.
雖然祂好像什麼都沒有.
但是祂樣樣都有.
祂滿有權能.
於是祂就有 是很寶貴.
所以值得我們來帶去敬佩.
最後我們看到主耶穌的情誼.
因為祂不單單為自己預備這個信.
祂甚至為彼得也預備了.
縱使彼得沒有這個要求.
這個錢銀是夠兩個人來用的.
是很美麗的.
所以祂和彼得有同道之誼.
也都有什麼呢?.
同伴之誼 有交往之誼.
當然彼得是要附上祂的努力.
但是因為祂的信服.
祂得到回報.
不單單物質 更加是祂信心的回報.
所以今天我們所信的神.
是耶華爾納的神.
祂會供應我們的需要.
只要我們用信心仰望祂.
信服祂就已經可以了.
在這裡可以說一個見證.
我預備這篇講章的時間.
剛剛預備到這一段的時候.
我就「叮」打了一個很短的訊息.
一個H字頭的.
就說有150鎊進了我的戶口.
我就以為是幾個星期前.
我港獨國安教會終於肯給我一個check in.
因為它沒有即時給.
國安教會都是H字頭的.
後來沒多久又收到一個訊息.
原來一位姊妹她發給我.
她說牧師 我入了一些錢給你.
當是新年的禮物.

$^{961}$我公司剛剛出了一些類似bonus.
我就給你.
我都很感動.
因為是這個愛心的支持.
我不是沒有錢.
但是很感動.
我沒有要求別人給我的.
但是她主動的給我.
我想一下 新年其實我都用了很多錢.
因為都有子女 有孫子 還有一些親戚.
但是我現在這幾年裡面是沒有收利是的.
我爸爸媽媽又離開了世界.
我岳父又離開了世界很久.
只有岳母在這裡.
但是我沒有在香港 是不是?.
這個時候只有出沒.
竟然會收到一封某個程度是利是.
很感恩.
原來我沒有要求.
但就預備.
你看 給得沒有要求.
但是主耶穌就來為她預備.
是多麼寶貴.
所以耶華已立.
後記是怎樣呢?.
《馬太記載》這個很獨特的記載是怎樣呢?.
不要忘記《馬太》本身的職業是甚麼.
是稅來的.
所以對交稅特別敏感.
所以聖明就透過它寫下了.
這件事是有原因的.
《馬太福音》大概完成的時間是.
公元80到90年.
即是那個時候聖殿已經被拆毀了.
那麼羅馬皇帝維斯柏.
就先規定將半赦赫勒的聖殿稅.
照樣交給另一個假神.
就叫做朱比塔的夫婦裡面.
所以猶太人就在一個困難裡面了.
他應不應該繼續立這個稅項.

$^{1001}$如果他們不立這個稅項.
就被視為背叛.
所以耶穌基督藉著這樣教導他們.
縱使你不願意.
但都要承擔這個公民責任.
事實上彼得在書信裡面作出這樣的教導.
所以都求神幫助我們了.
今日我們有機會去奉獻我們的福氣.
不要失去了這個福氣.
然後我們先者覺得可惜.
最後一個結論是甚麼呢?.
這件事讓我們看到.
主耶穌面對人對祂的身份的挑戰.
祂很有智慧的來回應.
祂用祂的權柄去表達祂的身份是怎樣.
更加重要的是.
教導我們有智慧去判斷.
哪些事我們要持守.
哪些事我們可以棄權.
目的就是不要叛倒人.
反而是建立人.
而且有祂的預備.
有祂奇妙的安排.
最重要的是我們是否用信心去信服.
好,我們一起來禱告.
天父多謝你.
因為你的話語何等寶貴.
雖然短短幾節的經文.
但都讓我們看到很多寶貴的教導.
神在幫助我們.
懂得怎樣去分辨甚麼是真心.
哪些是虛情假意.
教導我們懂得怎樣去有智慧去回應.
別人給我們的訊問,挑戰.
正如你透過彼得教導我們.
我們要常存準備的心.
是祂面對不同的人.
例如我們一些信仰上.
生活上的查詢甚至挑戰.
用溫柔智慧去回應.

$^{1041}$更加幫助我們懂得怎樣去分辨.
哪些事我們要持守.
哪些事我們可以妥協.
讓我們有勇氣面對我們可以改變的.
讓我們去接受不能改變的事.
讓我們去接受我們自己.
可以去面對這些情況.
多謝你為我們預備.
不僅為我們預備我們的救恩,得著拯救.
更加預備我們前面的路和我們的生活.
讓我們懷著這樣的心去領受.
願你祝福每一個參與查經的.
或者你有看YouTube的.
都幫助他們讓他們能夠有所成長.
聽下這面前的祈禱交托.
大奉主德勝國民 阿們.
\newpage



\section{約翰福音 15:16}
\label{sec:031745}
\textbf{呼召,回應,行動——帶來的祝福}~~~~Lydia Yeung姊妹
\newline
\newline
連結: \href{http://www.church.com.hk/acms/content.asp?site=cdc\&op=show\&type=product\&code=031745\&layout=sermon}{\texttt{031745}} ~~~~ 語音日期: 2026-03-01
\newline
\newline
\hyperref[sec:031744]{\small{< < < PREV SERMON < < <}}
~
\hyperref[sec:index_chronic]{\small{[返順時目]}}
~
\hyperref[sec:031746]{\small{> > > NEXT SERMON > > >}}
\newline
\newline
約翰福音 15:16
\newline
\begin{longtable}{cl}
\hline
\hline
章節 & 經文 (和合本修訂版)\\
\hline
15:16 & \begin{tabularx}{0.7\textwidth}{X} 不是你們揀選了我，而是我揀選了你們，並且派你們去結果子，讓你們的果子得以長存，好使你們奉我的名，無論向父求甚麼，他會賜給你們。 \end{tabularx} \\ \\
[1ex]
\hline
\hline
\end{longtable}
$^{1}$弟兄姊妹 早晨.
我今日Pastor Chloe.
經常叫她做Sue.
Pastor Chloe叫我有一些經文.
要有一些分享出來.
我就用了.
怎樣給神呼召.
呼召了之後.
怎樣給神我們怎樣去回應.
和怎樣去行動.
我想其實這三樣東西.
是很聯繫的.
你有呼召.
然後就會有怎樣去回應.
然後又怎樣去行動.
然後我就在這裡寫著.
還有帶來了一個祝福.
不知道大家會覺得.
真的有祝福嗎.
做這些事情.
你們覺得會不會有祝福.
有沒有.
是很有祝福.
很想和大家分享怎樣的祝福.
第一段經文是.
在《約翰福音》.
十五章十六節.
大家一起讀一讀.
其實在揀選.
主權.
永遠都在神裡面.
不是我們做了一些.
好與不好的事.
你可以去做.
不是一個交換的性質.
我經常看到.
揀選.
主權在神那裡.
我們會不會回應.
又是另一回事.

$^{41}$其實我們經常都聽到.
聽到神的呼召.
我們弟兄姊妹.
是不是很多時候都會聽到.
有沒有聽到.
不回答我.
我聽不到.
明明聽到.
都當聽不到.
我們怎樣去回應.
其實我覺得很特別.
當神願意呼召我們的時候.
祂讓我們結出美好的果子.
但我們怎樣去回應.
我覺得每一個人.
去回應神都不一樣.
好像在撒慕爾記上.
三章十節.
耶和華多次在夜間.
呼喚年幼的撒慕爾.
起初他都不知道.
後來叫了撒慕爾三次.
以利就知道.
耶和華和撒慕爾說話.
他就叫撒慕爾.
你要去回應神.
耶和華請說.
他的呼召就是在夜間.
我想每一個人都不一樣.
而我的呼召.
神怎樣去呼召.
我記得在2000年的時候.
有機會去做短孫.
去了黎巴嫩去做短孫.
那個情景就是.
去探訪一班孤兒.
我第一次見到這麼多孤兒.
我記得那次是派聖誕禮物.
你們也有.
派聖誕禮物的時候.

$^{81}$我看到很多很漂亮的眼睛.
那些孩子很漂亮.
但那時候我真的不太知道.
他們是孤兒.
但後來教會的弟兄姊妹和牧仔說.
這班是孤兒.
神就說你看看他們的眼睛.
你看看他們的眼睛.
我一看他們的眼睛的時候.
我的眼淚就不斷流出來.
知道他們是孤兒的時候.
我就在哭.
我就問神為什麼我會哭得這麼厲害.
神就給了我孤兒這兩個字.
直到現在.
我覺得每一個人聽到呼召真的不一樣.
我們要敏銳一點.
你真的會聽到神的呼召.
我還有一段經文.
大家也很熟悉的《約翰一書》一章十九節.
我想這段經文可能在聖經裡是一個又簡潔又深刻的一句.
我們愛因為神先愛我們.
這不是一個命令也不是一個要求.
而是一個很自然流露出來的.
當你去感受神的愛的時候.
你會很自然流露將這份愛出來.
為什麼我會這樣說呢.
因為我自己很實在經歷了神的愛神的拯救.
可能很多人不知道我之前的過去是怎樣.
很快地說一點.
我以前未信耶穌的時候.
我是一個非常抗拒去接受耶穌的人.
很抗拒.
我記得Chris的妹妹很早信耶穌.
她經常跟我傳福音.
你再跟我說我就不跟你哥一起了.
我覺得她很煩.
但我覺得神很特別.
祂給我一個很大的經歷.
我去了一次泰國.

$^{121}$我有邪靈上身.
弄了我很久.
弄了我病了很久.
邪靈上身我不說這麼詳細的事.
後來我經歷了神很實在的拯救.
我知道什麼是在死裡走出來.
我經歷了.
我就知道原來是這樣的.
原來神的愛可以這麼偉大這麼大.
當我經歷了神的愛的時候.
很想去拿這份愛出來.
很記得那時候神告訴我.
有機會出來服侍的時候.
我在想很多地方我都可以去.
因為我以前做生意.
我就跟神說到處我都可以去.
就是不要叫我去大陸.
因為很怕骯髒.
很怕見到Mickey Mouse.
不是美的那隻.
女士都知道我在說什麼.
很害怕.
但是不要跟神說.
你不想去做什麼.
偏偏就是這樣.
所以下一個視頻.
不是這個.
不是這邊.
是應該.
哪一邊.
這邊.
在20年前的時候.
其實21年前已經呼召我們去廣西南寧.
直到現在.
20年裡有很多恩典.
你看到我寫著.
去年我們有一個Fundraising Dinner.
我就寫了恩典滿滿20年.
今年神又再給我一個.
還要繼續寫什麼.

$^{161}$擴張神的國度.
我在想我做了20年.
還不需要吧.
都已經擴張了很多地方.
但是我記得.
我被神呼召去這個地方的時候.
你覺得害不害怕.
會不會害怕.
我在20年前.
就是2005年的時候.
我帶著我三個子女.
大女兒是15歲.
二兒子是10歲.
小女兒是9歲.
如果讓你帶著三個子女.
在那個20年.
不是現在2026年的中國大陸.
大家去過中國大陸玩了沒有.
現在是怎樣的.
很漂亮.
完全不一樣.
我記得那時候去.
我記得我奶奶真的罵我.
你為什麼要帶我的孫子去這些地方.
你在矽谷.
為什麼要帶他們去這些地方.
那時候我不懂得回應她.
我不知道怎樣回答她.
因為這是呼召.
她呼召你們夫妻.
為什麼要帶他們去.
她說的話都很實在.
我不懂得回答她.
但我知道神的呼召很強烈.
我們就去.
我記得去那一刻.
去到飛機場.
在飛機上降落那一刻.
我看著廿年前的廣西南寧.
那個地方是很偏僻.

$^{201}$很荒蕪的一個地方.
我看了下去.
是否真的要來這裡.
為什麼會這麼荒蕪.
好像去了一個.
好像還未發展的地方.
我當時就跟神說.
我很害怕.
可不可以退縮.
神很快就給了我一段經文.
在《瑪俄福音》五章三十六節.
經文裡面很簡單.
神經常給我一些很簡單的經文.
上一段的經文是什麼.
「我們愛,因為神先愛我們」.
這段更短.
「不要怕,只要順」.
我說這麼短的一段.
我一聽到這六個字.
我問神.
是不是你在跟我說話.
立刻想起聖經裡面.
當約書牙全心信服的時候.
神就會幫他.
給他一個有德性的策略.
當我們自己的眼對準神的時候.
其實我們不要上錯戰場.
神就在這裡給了我這段經文.
當我聽到神在跟我說話的時候.
我的心和情緒都立刻.
要像約書牙那樣信服.
來聽神的說話.
當我不害怕的時候.
我就明白.
我明白為什麼.
其實我經常都拿著.
這幾段經文.
在這二十年裡面.
我經常都這樣捉緊住.
「只要信,不要怕」.

$^{241}$「我們愛,因為神先愛我們」.
這幾段經文經常都在激勵我們.
怎樣去前行.
我們經常都說.
如果沒有神帶著我們去走.
我們是走不下去的.
在二十年前這些地方.
這麼荒蕪,荒僻.
沒有神的呼召.
如果你沒有一個呼召在你裡面.
其實你不能夠走下去.
但是當神告訴你.
你做的事是我差派你.
我真的看到在這麼多年裡.
神一直在開了很多門.
開了很多江河,開了很多資源.
開了很多不同的祝福給我.
我自己真的看到.
我就開始分享一些事功.
這個也是我們的口號.
我們的slogan.
「看顧在萬難中的孤兒寡婦」.
「以不離不棄的愛與他們同行」.
「相信愛可以改變生命」.
我在這麼多年裡.
神的開路之下.
我又開了孤兒院.
開了康復中心.
去服侍一些老癱的小朋友.
也有很多不同的地方.
大概在十年前.
現在已經到了緬甸.
不知道大家對緬甸有什麼印象呢.
KK園,真慘.
現在沒有了.
我聽聞收少了很多詐騙電話.
你們這邊.
我聽很多弟兄姊妹都這樣說.
好像少了很多.
之前每天起碼有十個八個.

$^{281}$現在真的少了.
其實不是這樣的.
緬甸不是你們想像中那麼差.
不過如果不差.
牧師經常跟我說.
如果這個地方沒有需要.
神又怎會叫你們去呢.
我覺得也是.
如果在十年前.
差不多十年前.
我們在2017年去.
未有這些事.
未有詐騙.
未開始.
可能幾年前我們去的時候.
現在你叫我去.
可能我真的不夠膽去.
但在那個時候有機會去.
我覺得是神很大的恩典.
給大家看看緬甸的呼聲.
給大家看看我們的活動.
我們興起了很多青線.
我們看到這麼多災難的國家.
最需要的是什麼呢.
神的愛.
他們很需要有耶穌基督.
當他們經歷到神的愛的時候.
我看到這一班年輕人.
真的很感動.
我可以說.
這個女孩是我們青少年中心的一個孩子.
她是一些難民.
他們長期打仗.
在他們的鄉下沒有機會讀書.
經常要到處去不同的地方.
因為打仗的緣故.
我們開了這個中心.
她來了我們中心三年.
那天我和她吃飯.
我問她.

$^{321}$你多久沒見過你的家人.
她說三年.
三年沒見過我的父母.
我的兄弟姊妹.
為什麼你不回去.
她說 第一 沒錢.
第二 很長途跋涉.
可能要回鄉下幾天.
我問 為什麼你都留在這裡.
你不想念你的家人嗎.
她說 我很想念我的家人.
但我在這裡 我有機會讀書.
我聽到她這一句.
她的成績很好.
我們這班年輕人.
我們有五十多個青少年.
全部在他們的成績裡是Top 10.
全部在學校有很好的成績.
我在想.
我不知道在這些落後的地方.
其實小朋友有機會讀書.
他們很珍惜.
聽完她之後.
我真的覺得我們是一個福氣.
有機會認識這班孩子.
是一個福氣.
每一年都有兩次百多個青少年.
一起來舉行這個營會.
我們的題目是什麼.
看到嗎.
All for Jesus.
這班孩子很熱心.
也很愛耶穌.
剛才你看到他們.
我記得我們有一個.
完了之後.
我們帶團隊去做Pair Walk.
去當地的地方Pair Walk.
你知道緬甸是一個佛教國家.
很多廟宇.

$^{361}$我們去到那些廟宇.
跟他們一起祈禱的時候.
其中一個大概十多歲的女孩.
她一直在哭.
我問她.
你以為在這些廟宇的地方.
有邪靈.
或者有些屬靈增賤的東西.
在為她祈禱.
然後再問她.
你為什麼剛才哭.
你知道她怎樣回答我嗎.
一個十六歲的女孩.
她說看到這裡有很多.
未信耶穌的人在拜佛.
很難過.
所以一直在哭.
聽完之後.
我們有沒有這樣的心.
我們自己有沒有這樣的心.
當你給了他們耶穌的時候.
他們經歷到神的愛的時候.
他們懂得怎樣為有需要的人祈禱.
我看到真的很感動.
很感恩.
所以我說在這二十年裡.
認識到他們是一個很大的福氣.
在去年.
不知道大家知不知道.
有一個地震很嚴重.
在三月底的時候.
我想Christ和我先生.
和我兒子都在那裡.
我們看到這麼大的地震.
死了很多的家人.
整條村200人.
死了100人.
我們就決定在這個地方.
再開一個孤兒院.
現在開了第一期.

$^{401}$這是第一期的興建座.
1月1日.
很感恩.
大家有機會可以進去祈禱.
孩子們.
這個就是.
剛才在外面.
有很多公仔.
我們還有一個program.
叫做「感於夢想」.
其實是我兒子Josh.
想出來的.
他經常跟我說.
媽咪我們服侍這麼多孩子.
將來一定會有一個.
會很厲害的.
我們為他們祈禱.
這個兒子叫做志剛.
我在三個月前.
真的很感動.
也很感恩.
我們可以打籃球.
我們的「感於夢想」有很多不同的活動.
有籃球.
有足球.
有跳舞.
畫畫 唱歌 什麼都有.
這個兒子很喜歡打籃球.
他在三個月前.
被緬甸的國家隊.
青少年的國家隊.
選了他.
我說兒子.
成就了你說的話.
我們真的有機會.
讓這些孩子.
不單要介紹耶穌給他們.
還要讓他們有自信心.
我也很想.
跟大家分享一個故事.

$^{441}$其中一個小朋友的故事.
他那邊打工.
然後他走了.
然後沒回來了.
然後我們四個都在那裡等著他回來.
每天都去等著.
然後一個多月.
兩個月.
差不多兩個月就我們知道了.
他不會回來了.
家他們的時候.
過得也不是很好.
每天要流淚吃飯.
就是.
那個.
爸爸媽媽的時候.
他們就會說.
你們爸爸媽媽都已經不愛你們.
所以拋棄你們了.
所以你在我們家也是沒用的.
反正不餓了嘛.
他們這樣會說.
因為他們不喜歡爸爸媽媽.
就是要找那個兔子的.
兔子吃的菜.
如果那個找不到的話.
別吃飯.
我去學校也是老師.
他們是認識爸爸媽媽的.
他們也知道爸爸媽媽不在了.
然後他們.
有時候會嘲笑我們.
那種每天都會.
哭著眼淚的時候.
我有時候我都想死了.
我都不想活在這個世界上了.
我覺得.
活在這個世界上了沒用.
就是沒有希望了.
就是我每天都會禱告.

$^{481}$跟爸爸媽媽說.
爸爸我不想在這個世界了.
可是如果.
我知道.
我知道你.
肯定會幫我們的.
然後我每天都哭著.
跟他說我不想在這個家了.
然後我想.
那種在一個.
那種有希望的地方.
然後他.
就聽我的話.
然後來到了英才家園.
然後在這裡的話.
溫暖.
我們得到很溫暖的家庭.
還有爸爸媽媽愛的感受.
我們也被愛到.
所以很開心.
還有我們能夠好好的.
可以學習.
去學習的時候也沒有.
沒有人嘲笑我們.
所以很開心.
最感動的就是.
以前母親節的時候.
我們如果.
我們在家鄉的話.
那個.
每一年.
到母親節的時候.
我們沒有媽媽.
所以.
一句都沒有說到.
媽媽我愛你.
還有母親節快樂.
然後來到英才的時候.
每一年.
母親節的時候.

$^{521}$我們可以說媽媽我愛你.
還有.
說到母親節快樂.
我們能夠見到媽媽.
養爸說.
我們都是你們永遠的.
爸爸媽媽.
我那時候很開心.
因為沒.
沒人說過我們.
我永遠都是你的.
爸爸媽媽.
我們一次都沒有聽到過.
我在英才的時候.
我在家鄉的時候.
我的心裡.
最大的夢想就是.
有一天.
如果我的爸爸媽媽.
一直都在旁邊.
在我的身邊.
陪我照顧好.
他們在看我長大.
我最喜歡的時候.
就是這個.
可是沒有.
所以養爸.
直接這樣說的話.
我心裡覺得很開心.
然後我以前也每天都會找到.
就是.
有一天會.
找到一個最愛我的.
爸爸媽媽.
演奏那未來.
從未被遺忘.
從今不再是一個.
與你同往.
不需再迷惘.
孤單感覺.

$^{561}$為你驅趕.
環彩百七的光.
燃亮心中盼望.
陪著你.
期待你.
一天天的絕蹤.
孩子.
長大.
孩子.
從今不再.
漂泊.
從今不再.
漆黑裡.
摸索.
其實我覺得很感恩.
就是在20年前.
神就告訴我們.
那個calling.
我自己的calling.
神給我的是孤兒.
當我們認識了這麼多的孤兒的時候.
讓我們可以成為他們的父母.
他們的父母.
詩恩跟我們說.
他說.
世上最好的就是.
養爸爸養媽媽.
我長大之後.
都要像他們一樣.
來服侍神照顧有需要的孩子.
詩恩是我們在緬甸雲彩家園的一個女孩子.
是神給我們的.
其中一個的女兒.
很感恩神.
給我們作為他們的父母.
這個身份.
讓我們可以看到這些孩子.
他們生命的翻轉.
來去榮耀.
榮耀八天父.

$^{601}$這個地方我剛才也說了.
我們已經可以入住.
所以希望.
更多的孩子.
可以來得到.
神的祝福.
我們將會.
另外可以收到二百多個小孩.
現在已經有.
幾十個.
有很多房間可以讓更多的小孩.
可以入住.
如果有機會.
都可以來支持這個項目.
2014年的時候.
因為打仗的緣故.
我們很多小孩都要逃難.
去兩個地方.
那時候我收到一幅圖畫.
不知道大家有沒有和你們分享過.
這一幅圖畫.
是一些小孩畫的.
你看到些什麼呢.
看到小孩畫些什麼.
看到嗎.
我自己真的覺得.
沒什麼理由.
小孩會畫.
有些子彈.
放下來.
他們在跪下祈禱.
其中一個班.
畫了出來.
我自己馬上拿了這幅圖畫.
看不看得清楚.
我自己馬上拿了這幅圖畫.
你看看這一邊.
全部的小孩.
很多人躺下了.
因為有子彈.

$^{641}$有些子彈放下來.
這個小孩在祈禱.
他們很想上學.
這是他們兩年前的圖畫.
剛剛我又收到一幅.
你看不看到這幅畫.
他們畫了些什麼.
上學.
原來對於一班.
在這樣的地方的小孩.
可以上學.
原來是一個恩典.
對於他們來說.
我走到這間學校.
我跟巴塞科里說.
這間學校很漂亮.
原來我們那邊.
肯定不會有這樣的學校.
這間學校非常之美麗.
我看到這幅圖畫的時候.
我覺得很開心.
小孩不用再在.
子彈之下生存.
很記得.
剛開始做孤兒院的時候.
一弟兄.
我跟他分享.
要做孤兒院.
他說我捐一千元給你們.
一千元你覺得可以做到什麼.
不是美金.
人民幣.
一千元美金和一千元人民幣.
基本上都不可以做到什麼.
當你願意去信服的時候.
當神召喚你.
你懂得去回應.
有行動的時候.
真的很感動.
也很感恩.

$^{681}$今天雲彩行動.
服侍超過一千個小孩.
我們看到.
神實在是很大的恩典.
對雲彩.
這個也是.
一些.
對弟兄姊妹的挑戰.
挑戰和呼召.
用優先的次序.
在我們的生命裡.
什麼才是優先的.
是你們的事業.
是你們的安全感還是神的國度.
忠心的託付.
如果今天主耶穌回來.
我們能夠坦然的.
這樣跟祂說.
主啊,我已經把你的託付.
放在第一位.
你們自己想想.
自己大家想想.
是不是神跟你說的時候.
你是怎樣回應.
在馬太福音六章三十三節.
你們要先求祂的國和祂的義.
這些東西都會加給你.
願我們都可以.
成為神捍衛.
我們一個忠心的僕人.
等到回去見.
耶穌基督的時候.
神跟我們說好.
你是一個又良善又忠心的僕人.
如果有機會你們自己.
可以.
回去看看.
這些能不能做到.
我們可不可以.
跟神說.

$^{721}$我們的優先次序是怎樣.
跟我們的耶穌基督.
我們看到.
在過去這麼多年.
這一群小朋友.
是很失落的.
但是.
我今天看到他們的笑容.
其實我已經覺得是.
一個很大的祝福.
希望大家可以.
繼續為我們這一群孩子.
去禱告.
要一起去.
很需要.
很需要大家的禱告.
也很需要大家的支持.
我今天都會擺了.
很多部份在外面.
如果你們看到有需要.
都可以出去.
助養或者做一個.
Monthly Support.
在外面大家可以看看.
我們一起去低頭祈禱.
主耶穌我們多謝你.
多謝你去使用.
這些行動成為很多孩子的祝福.
求神與他們同在.
讓我們知道.
怎樣去.
當聽到呼召.
我們怎樣去.
去作出一個回應.
回應之後.
怎樣去做這個行動.
因為我們看到.
當我們的心打開.
當我們心願意的時候.
你就是開道路開光學的神.

$^{761}$讓我們.
每一個弟兄姊妹.
無論在哪一個崗位.
我們都說我們在這裡.
請差遣我.
讓我們可以成為多人的祝福.
讓我們可以離開.
一些安書的地區.
可以去為神.
去更加作美好的見證.
結美好的果子.
多謝耶穌聽我們的禱告.
奉耶穌基督的名禱.
\newpage



\section{羅馬書 16:1-27}
\label{sec:031746}
\textbf{堅固人心的福音奧秘}~~~~袁惠鈞牧師博士
\newline
\newline
連結: \href{http://www.church.com.hk/acms/content.asp?site=cdc\&op=show\&type=product\&code=031746\&layout=sermon}{\texttt{031746}} ~~~~ 語音日期: 2026-03-01
\newline
\newline
\hyperref[sec:031745]{\small{< < < PREV SERMON < < <}}
~
\hyperref[sec:index_chronic]{\small{[返順時目]}}
~
\hyperref[sec:031747]{\small{> > > NEXT SERMON > > >}}
\newline
\newline
羅馬書 16:1-27
\newline
\begin{longtable}{cl}
\hline
\hline
章節 & 經文 (和合本修訂版)\\
\hline
16:1 & \begin{tabularx}{0.7\textwidth}{X} 我對你們推薦我們的姊妹非比，她是堅革哩教會中的執事。 \end{tabularx} \\ \\ \relax
16:2 & \begin{tabularx}{0.7\textwidth}{X} 請你們在主裡用合乎聖徒的方式來接待她。她在任何事上需要你們幫助，你們就幫助她；因她素來幫助許多人，也幫助了我。 \end{tabularx} \\ \\ \relax
16:3 & \begin{tabularx}{0.7\textwidth}{X} 請向百基拉和亞居拉問安。他們在基督耶穌裡作我的同工， \end{tabularx} \\ \\ \relax
16:4 & \begin{tabularx}{0.7\textwidth}{X} 也為我的性命把自己的生死置之度外；不但我感謝他們，就是外邦的眾教會也感謝他們。 \end{tabularx} \\ \\ \relax
16:5 & \begin{tabularx}{0.7\textwidth}{X} 又向在他們家中的教會問安。向我所親愛的以拜尼土問安，他是亞細亞歸於基督的初結果子。 \end{tabularx} \\ \\ \relax
16:6 & \begin{tabularx}{0.7\textwidth}{X} 又向馬利亞問安，她為你們非常辛勞。 \end{tabularx} \\ \\ \relax
16:7 & \begin{tabularx}{0.7\textwidth}{X} 又向與我一同坐監的親戚安多尼古和猶尼亞問安，他們在使徒中是有名望的，也是比我先在基督裡的。 \end{tabularx} \\ \\ \relax
16:8 & \begin{tabularx}{0.7\textwidth}{X} 又向我在主裡面所親愛的暗伯利問安。 \end{tabularx} \\ \\ \relax
16:9 & \begin{tabularx}{0.7\textwidth}{X} 又向我們在基督裡的同工耳巴奴和我所親愛的士大古問安。 \end{tabularx} \\ \\ \relax
16:10 & \begin{tabularx}{0.7\textwidth}{X} 又向在基督裡經過考驗的亞比利問安。向亞利多布家裡的人問安。 \end{tabularx} \\ \\ \relax
16:11 & \begin{tabularx}{0.7\textwidth}{X} 又向我親戚希羅天問安。向拿其數家在主裡的人問安。 \end{tabularx} \\ \\ \relax
16:12 & \begin{tabularx}{0.7\textwidth}{X} 又向為主辛勞的土非拿和土富撒問安。向所親愛、為主非常辛勞的彼息問安。 \end{tabularx} \\ \\ \relax
16:13 & \begin{tabularx}{0.7\textwidth}{X} 又向在主裡蒙揀選的魯孚和他母親問安，他的母親就是我的母親。 \end{tabularx} \\ \\ \relax
16:14 & \begin{tabularx}{0.7\textwidth}{X} 又向亞遜其土、弗勒干、黑米、八羅巴、黑馬，和跟他們在一起的弟兄們問安。 \end{tabularx} \\ \\ \relax
16:15 & \begin{tabularx}{0.7\textwidth}{X} 又向非羅羅古和猶利亞，尼利亞和他姊妹，阿林巴和跟他們在一起的眾聖徒問安。 \end{tabularx} \\ \\ \relax
16:16 & \begin{tabularx}{0.7\textwidth}{X} 你們要以聖潔的吻彼此問安。基督的眾教會都向你們問安！ \end{tabularx} \\ \\ \relax
16:17 & \begin{tabularx}{0.7\textwidth}{X} 弟兄們，那些離間你們、使你們跌倒、違背所學之道的人，我勸你們要留意躲避他們。 \end{tabularx} \\ \\ \relax
16:18 & \begin{tabularx}{0.7\textwidth}{X} 因為這樣的人不服侍我們的主基督，只服侍自己的肚腹，用花言巧語誘惑老實人的心。 \end{tabularx} \\ \\ \relax
16:19 & \begin{tabularx}{0.7\textwidth}{X} 你們的順服已經傳於眾人，所以我為你們歡喜；但我願你們在善上聰明，在惡上愚拙。 \end{tabularx} \\ \\ \relax
16:20 & \begin{tabularx}{0.7\textwidth}{X} 那賜平安的神快要把撒但踐踏在你們腳下。願我們主耶穌基督的恩與你們同在！ \end{tabularx} \\ \\ \relax
16:21 & \begin{tabularx}{0.7\textwidth}{X} 我的同工提摩太，和我的親戚路求、耶孫、所西巴德，向你們問安。 \end{tabularx} \\ \\ \relax
16:22 & \begin{tabularx}{0.7\textwidth}{X} 我這代筆寫信的德提，在主裡向你們問安。 \end{tabularx} \\ \\ \relax
16:23 & \begin{tabularx}{0.7\textwidth}{X} 那接待我，也接待全教會的該猶，向你們問安。城裡的財務官以拉都和弟兄括土向你們問安。 \end{tabularx} \\ \\ \relax
16:24 & \begin{tabularx}{0.7\textwidth}{X} 願我們的主耶穌基督賜恩典給你們各位！阿們！ \end{tabularx} \\ \\ \relax
16:25 & \begin{tabularx}{0.7\textwidth}{X} 惟有神能照我所傳的福音和所講的耶穌基督，並照歷代以來隱藏的奧祕的啟示，堅固你們。 \end{tabularx} \\ \\ \relax
16:26 & \begin{tabularx}{0.7\textwidth}{X} 這奧祕如今顯示出來，而且按著永生神的命令，藉眾先知的書指示萬民，使他們因信而順服。 \end{tabularx} \\ \\ \relax
16:27 & \begin{tabularx}{0.7\textwidth}{X} 願榮耀，藉著耶穌基督，歸給獨一全智的神，直到永遠。阿們！ \end{tabularx} \\ \\
[1ex]
\hline
\hline
\end{longtable}
$^{1}$羅馬書第四十章.
今天我們要講羅馬書第十六章.
我們知道羅馬書只有十六章的經文.
所以今天是最後的一章.
每講到每一個系列最後一章.
我都很不捨得.
羅馬書亦都不例外.
今日的題目叫做甚麼.
「堅固人心的福音奧秘」.
其實我在找題目的時候.
我很多時候想避免甚麼字.
奧秘,密碼.
有很多牧師,傳道人都希望放上這些字.
來引你去聽他的篇道.
其實我們不應該這樣做.
今日為何用到奧秘這個字.
就是在羅馬書裡自己講的.
保羅在羅馬書第十六章二十五節.
他說要照永古隱藏不言的奧秘.
堅固你們的心.
神的話語就是能夠堅固我們的心.
幫助我們能夠活出那美好的見證.
講到這裡又要稍為做一些廣告.
我講這麼多的系列.
都是希望跟大家一齊來讀經.
我經常都說.
你一個星期聽我五篇道.
然後讀相關的經文.
你可以兩年半讀完整部聖經.
而且對聖經是有很認識的了解.
因為我在每一個系列第一篇道.
叫做為何要讀這部書.
就是講到有很多的背景,很多的資料.
所以能夠幫助你認識這部書.
這樣讀經法還有甚麼好處.
是用字面來解經.
是用聖經來解釋聖經.
所以你會發覺是一致性的.
你有沒有讀聖經的時候.
你會發覺有很多矛盾.

$^{41}$但是當你用聖經來解釋聖經.
字面來解釋聖經.
你會發覺所有的矛盾迎刃而解.
我從來沒有發覺有任何的矛盾在聖經裡面.
很多人認為聖經裡面有很多的寓意.
很多的靈異.
他們甚至會引用保羅在《多倫多後書》的三章六節.
說「字句叫人死,精義叫人活」.
但是你不要將這本經文推得太遠.
你要看三章三節.
三章七節是在說甚麼.
那些刻在石板上的文字是在說甚麼.
是在說律法.
如果你是守著律法的字面.
死路一條.
我們要守律法的精義.
我們要透過新約接受耶穌基督.
我們這樣才能真正得到永生的盼望.
如果這個世上真的有形意,寓意解經的話.
那我就要問.
何時一節經文你是看它的實意.
何時一節經文你又看它的寓意呢.
有人說永遠都是指著將來.
是耶穌基督的救贖.
我同意.
但是不等於那節經文沒有了實質的意義.
我們的聖經和可蘭經.
和佛教的經書有甚麼不同的地方.
我們的聖經都是歷史.
是有考古學的印證.
我們的聖經是有神實實在在.
實質上的應許是應驗了.
是預言這是真真正正的神蹟奇事.
如果沒有了這些的話.
沒有了實質的應許.
我們這部聖經就是和其他宗教一模一樣.
沒有甚麼分別.
還有一個問題.
十個形意解經加寓意解經加.
同一節的經文解出來都不同的.

$^{81}$我們相信誰所說的.
我聽過有人說.
不是的牧師.
是有系統性來教我們怎樣解經.
我的問題就是.
如果真的有系統性的話.
那就不是每個人都懂得自己解經.
你是需要有神學的訓練.
聖經說甚麼.
每個人自己都能夠解經.
透過聖靈的光照.
所以我們真的要學習用字面解經.
用聖經來解釋聖經.
今天我們要講羅馬書的第十六章.
首先有少少的背景.
我講過羅馬書的第一章到第十五章的頭.
保羅講出恩信稱義的神學觀.
然後在羅馬書的第十五章尾.
他用他自己的事工.
用他的行動來活出恩信稱義的道理.
今天第十六章.
亦都是有關保羅的事工.
但他講到和羅馬教會弟兄姊妹的關係.
他講出了有很多人的名字.
他是問他們的安.
這些人都有參與保羅的生活和保羅的事工的人.
他向他們問安.
這張書給我們見到保羅對羅馬教會的愛.
彼此的信任.
亦都講出了很多的人名.
有人會說牧師.
人名你又講一篇道.
問安你又講一篇道.
是呀.
我要鼓勵大家來讀經.
你不要以為有名字有問安的東西.
你覺得很悶你就不去讀.
其實今日的經文好像禾稈cum 珍珠一樣.
在人名當中亦都有很多精彩的神學.
我們需要學習.

$^{121}$有很多精彩的教導.
記不記得曾幾何時我講過一篇道.
是大衛勇士的名單.
勇士的名單人名都可以講一篇道.
我當時講過.
若果你們聽了不覺得悶的話.
我就成功了.
今日這一篇道又是一樣.
若果你們聽了不覺得悶的話.
我又成功了.
其實羅馬書第十六章都有一個結構.
你們認識我的人就知道我很喜歡首尾結構.
首尾結構是ABCCBA.
就好像箭咀一樣指著中間的經文是最重要的.
今日羅馬書十六章就是一個平衡結構.
平衡結構是ABCABC.
A和A有連帶關係.
B和B有連帶關係.
C和C又有連帶關係.
今日我們十六章就是ABAB.
A是講甚麼.
是問安.
羅馬書十六章一至六十六節.
保羅問羅馬教會的安.
這些人名是有關羅馬教會.
然後到下面的A.
羅馬書十六章二十一至二十四節.
是問安的人.
這些是保羅的同工.
這些人名是保羅的同工.
不再是羅馬教會.
B的部分有兩點.
在十六章的十七至二十節.
是保羅給羅馬教會的忠告.
剛剛才問完安對這些人.
表示了大家的愛意.
保羅就馬上講真理.
如果你只是講愛沒有真理的話.
這些是假的.
所以我們見到很重要.

$^{161}$這裡忠告羅馬書十六章十七至二十節.
保羅給羅馬教會的忠告.
下面最後呢.
羅馬書十六章二十五至二十七節.
是最後給羅馬教會的祝福.
我們一起來先看看.
A問安的部分.
羅馬書的十六章一至到十六節.
一開始一章一節保羅說.
我對你們舉薦我們的姊妹菲比.
她是堅格理教會中的女執事.
菲比這個名字的意思.
就是明亮輝煌.
我時時都說.
在聖經裡的人名和她的性格.
是有連帶的關係.
我相信這個姊妹菲比.
她真的有外向的性格.
她真是一個人很明亮很開心.
保羅叫她把羅馬書帶到羅馬.
保羅在哥林多寫羅馬書.
哥林多在哪裡.
在現在的希臘.
羅馬在現在的意大利.
保羅在這封書上介紹了菲比.
這是早期教會一個很普遍的現象.
當一個弟兄或姊妹.
由一個城市到另一個城市.
來定居或工作的時候.
他時時都會帶一封介紹信.
證明自己是一個真正的基督徒.
在使徒行傳的十八章.
阿波羅就是有一封介紹信給他.
在哥林多後書的八章十八節.
提到有一個同行的弟兄.
也有一封介紹信.
其實在埃及沙漠裡有一個叫沙丘的地方.
考古學家發現了有很多這樣的介紹信在那裡.
看起來好像一個垃圾堆.
但居然找到有很多很有價值的文獻.

$^{201}$就在埃及的沙丘.
在羅馬書的一章二節.
保羅說:請你們為主接待他.
合乎聖徒的體統.
他在何事上要你們幫助.
你們就幫助他.
保羅叫羅馬教會的弟兄姊妹.
要幫助這位菲比姊妹.
為什麼?.
保羅說了有四個理由.
第一,她是一位姊妹.
是主內的好姊妹.
當有一個主內的弟兄或者姊妹來到我們當中.
我們怎樣?.
是不是趕人走?.
當然不是.
我們要招待她.
這是聖經的教導.
第二個原因.
她是一個執事.
她不單只是一個姊妹.
還是執事.
是堅格理教會的女執事.
保羅在哥林多寫羅馬書.
哥林多距離堅格理大概八至九里.
堅格理是哥林多的一個港口.
如果你要從哥林多出發到別處.
你就要從堅格理上船出發.
堅格理教會很可能是哥林多教會所建立的教會.
就好像我們CCAC.
我們在New Life.
就是我們的籍堂.
所以堅格理教會可能也是哥林多教會籍堂出去的教會.
她是一個女執事.
執事這個字.
Diakonos這個原文的希臘文.
就是執事.
慢慢演變成我們英文Deacon這個字.
這個字是沒有性別的.
沒有男沒有女.

$^{241}$是中性的一個字.
在福音書《使徒行傳》和《新約》的書信.
都有講到很多在教會服侍人的人.
但是在提摩泰前書三章八節之前.
沒有具體地講出這些人是需要有什麼資格.
有什麼質素.
在提摩泰前書的三章八節就有講到執事的資格.
而在提摩泰前書的三章十一節.
就有講到有女執事的資格.
有很多人說不是女執事.
這裡是執事太太的資格.
但是你有沒有發現在《新約》的書信裡面有講到長老太的資格.
我們這裡也有很多長老太.
你們有沒有見過聖經裡面有長老太的資格.
沒有的.
所以這裡絕對不是執事太太的資格.
這裡是女執事的資格.
除了菲比是一個好姐妹.
除了她是一個女執事.
她也是一個好幫手.
在《羅馬書》的一章二節.
保羅說:因為她數來幫助許多人.
也幫助了我.
這個菲比姐妹很喜歡幫助人.
幫助這個字是原文希臘文的prostatis這個字.
幫助人是可以是在經濟上的支持.
可以是物質上的支持.
究竟是什麼支持我們已經不知道.
如果是經濟上的支持.
這個姐妹就好像推也推拉的女底也.
記不記得這個姐妹.
在《使徒行傳》的十六章十四節.
她是專門買賣紫色布的一個商人.
很有錢的.
經文說這個菲比曾經幫助過保羅.
究竟怎樣幫助我們已經無從稽考.
還有她不單止是主內的姐妹.
是執事是好幫手.
她還是一個差役.
保羅叫她把《羅馬書》帶到羅馬.

$^{281}$這一部被聖靈默示的書帶到羅馬.
妳說差役只是打雜而已.
不是的.
在聖經裡面傳遞神話語的人很重要.
在希伯來書的二章二節.
新近的洛木斯才跟我們說了希伯來書.
二章二節怎麼說.
傳遞神話語的是天使.
將律例典章傳給我們世人.
透過天使.
所以我們千萬不要小看這些傳遞神話語的人.
所以我們都有責任來接待主內的弟兄姊妹.
去同情他們關心他們.
滿足他們一切的需要.
由羅馬書的十六章三節至十六節.
保羅一連串講出了有很多人的名字.
向他們來問安.
真的非常難得.
保羅不單止知道有這些人.
他還認識他們.
弟兄姊妹在我們粵語堂.
你知道有多少弟兄姊妹.
但是你認不認識他們.
這才是最重要的.
不只是知道有這些人.
我們還需要認識這些人.
保羅從來沒有去過羅馬.
他知道那裡有二十四個弟兄姊妹.
十七個是男的.
七個是女的.
他還問候兩個家庭.
還有很多人可能是他認識的.
但是他沒有說.
保羅間接直接與這些人的信仰.
與這些人的生活都有關係.
他甚至好像是羅馬教會的一份子一樣.
但是他從來沒有去過羅馬教會.
在羅馬書的十六章三節.
他說「問伯羈拉和雅居拉的.
按他們在基督耶穌裡與我同供」.

$^{321}$記不記得這對夫婦.
伯羈拉和雅居拉的名字.
在聖經裡一共出現六次.
其中有四次.
都是伯羈拉的名字先行.
伯羈拉是誰.
是太太.
現在在美國是很普遍.
Ladies first.
你當然是先擺太太的名字.
但是現在說的是古老社會.
是希臘的社會.
是羅馬的社會.
還是男輕女的.
居然伯羈拉的名字先行.
為什麼.
有很多人說伯羈拉可能.
他是很愛鄉很活潑.
是有很堅強的個性.
所以他就先擺.
我的確認識有很多主內的姊妹都是這樣.
但也有人說不是.
伯羈拉的社會地位高.
他可能是一個羅馬的貴族.
他的丈夫可能是一個猶太人.
是一個疊棚佬.
我們怎麼知道.
使徒行傳的十八章三節告訴我們.
這個雅居拉是疊棚的.
這叫夫貧妻貴.
我們不知道究竟答案是怎樣.
這只不過是我們的推測.
保羅在哥林多的猶太會堂認識這對夫婦.
因為大家都有傳福音的熱愛.
所以就變成了好朋友.
這對夫婦最初在羅馬.
後來我們知道歷史告訴我們.
羅馬王赫魯迪烏斯驅逐所有的猶太人.
如果這個伯羈拉真是猶太人的話.
他一定被驅逐.

$^{361}$而成為他的太太.
就算他是羅馬的公民.
就算是羅馬的貴族.
當你的丈夫被驅逐你會怎樣.
嫁雞除雞.
跟著他離開羅馬.
所以在這段驅逐的期間.
他們就去到哥林多.
在《使徒行傳》的十八章就認識了保羅.
兩年後他們去到爾忽所傳福音.
後來當這個赫魯迪烏斯去世了.
他們又回到羅馬.
所以保羅向他們問安.
保羅怎樣說.
在羅馬書的十六章四節.
也為我的命將自己的頸汗自知到愛.
這對夫婦真是不得了.
為了保羅將自己的頸汗自知到愛.
什麼意思.
伸條頸出來斬.
願意為保羅犧牲.
若果不是有這些人.
可能沒有了這一部羅馬書.
雖然是非比將這一部書從哥林多帶到羅馬.
但是如果沒有了伯基拉和阿基拉.
可能沒有了羅馬書也說不定.
但是你說不會啊牧師.
神的旨意一定會成就.
當然.
但是我們都要怎樣.
要多謝這些人.
所以保羅在羅馬書的十六章四節怎樣說.
很多時我們覺得保羅很偉大.
但是保羅提醒我們.
沒有了這些幕後的人幫我.
我不能夠成事.
所以你們要去多謝他們.
令我想起誰.
令我想起我們當中的Terry Ho.
他未受傷之前他怎樣.

$^{401}$每一個敬拜第一個最早回來的就是他.
現在當然是Benny長老.
最早回來的就是他.
敬拜完最後走的是誰.
又是他.
他去年傷了之後.
傷了手已經很久沒有回來.
我們都很需要為他來禱告.
為我們粵語堂所付出.
願意神繼續保守他.
讓他早日回來.
能夠和我們一齊侍工.
在羅馬書的十六章第五節.
保羅說.
「又問在他們家中的教會安」.
原來最早羅馬教會在哪裡舉行.
在伯記拉和阿居拉的家裡.
現在我們經常聽見國內有甚麼家庭教會.
原來最早期羅馬教會.
就在伯記拉和阿居拉的家裡舉行.
接著保羅向他自己初結的果子來問安.
在羅馬書的第十六章第五節說.
「問我所親愛的以拜尼圖安.
他在亞細亞時歸基督初結的果子」.
你說小亞細亞在哪裡.
是現在的土耳其.
記不記得在啟示錄的二章和三章.
耶穌基督寫信給七間教會.
有甚麼教會.
以弗所,善美娜,別加摩,推牙推拉,薩迪,.
菲拉鐵菲和魯底加.
這些就是在亞細亞的教會.
這個以拜尼圖就是在亞細亞.
他是保羅初結的果子.
注意保羅沒有說是我的果子.
給我們又會怎樣.
是我的我帶領他信主.
拍了針口.
但保羅沒有這樣.
他說是基督所結的果子.

$^{441}$為甚麼.
保羅在擦鞋?.
你說不是.
他自己說的.
羅馬書的十五章十六節.
要將外邦人獻給基督.
這個以拜尼圖就是保羅帶領他信主.
然後獻給基督.
所以是基督的果子.
在羅馬書的十六章六節.
保羅又問瑪利亞案.
他為你們多受勞苦.
在新約聖經裡有多少個瑪利亞.
一共有六個.
你們可以慢慢去數.
這個也是其中一個瑪利亞.
瑪利亞是一個非常常見很普遍的名詞.
保羅說甚麼.
他多受勞苦.
為了羅馬教會.
這個原文希臘文Kopiar這個字.
就是勞苦這個字.
這個姊妹做到汗流浹背一場劇.
我們香港人會說.
為了羅馬教會.
我們要多謝他.
還有呢.
他要向那些和他一起坐牢的弟兄.
來問案.
在羅馬書的十六章七節.
又問我的親屬與我一同坐牢的.
安多尼古和尤尼亞安.
安多尼古絕對是一個男人的名字.
但是尤尼亞可以男性可以女性.
但是說到和保羅一起坐牢.
那你想一下他是男還是女.
應該是男.
現在很多的世界都很開放.
但我都沒見過有男女同慾.
所以相信這個真的是一個男人.

$^{481}$保羅說這個是他的親屬.
親屬這個字又要解釋一下.
可以解釋做同族.
即是說是猶太人.
但也可以解釋做親戚.
真的是保羅的親戚.
這個字在羅馬書的十六章十一節和二十一節.
也都有用到.
如果我們注意的話.
保羅介紹這麼多人.
有很多都是猶太人.
但他沒有用到這個親屬這個字.
所以我相信這個真的是保羅的親戚.
等一下我們會見到還有兩個都是保羅的親戚.
在《使徒行傳》的十六章七節.
保羅說.
他們在使徒中是有名望的.
也是比我先在基督裡.
使徒這個字是什麼意思.
是被差派出去的.
有人說這些是什麼人.
經文說是有名望的.
是比保羅先在基督裡.
可能這些人他們是認識十二使徒.
或者甚至和十二使徒同工過.
保羅絕對不是白鴿眼.
他還向誰來問安.
一些做奴隸的信徒.
在羅馬書的十六章第八節.
又問我在主裡所親愛的暗白尼安.
在羅馬的皇室裡.
暗白尼這個名字是很出名的奴隸.
是一個奴隸的家族.
在羅馬有一個基督教的墳場.
如果你進去的話.
今時今日你會發現有一個很大的墓碑.
上面寫著暗白尼之墓.
這個真的非常特別.
世人的墳墓是怎樣的.
和你的名譽地位和你的身家成正比.

$^{521}$你的名譽越高你的身家越大.
你的墓碑越大.
但是在基督教裡.
就算一個奴隸也有一個冠冕堂皇的墓碑.
在羅馬書的第十六章第九節.
保羅又向在基督裡同工的爾巴奴.
和親愛的士大古文安.
士大古這個名字希臘文就是墨水或者谷水的意思.
這兩個人是誰.
我們現在已經無從稽考.
保羅這個人真的很特別.
他的傳福音是怎樣的.
由最低級的奴隸一直到最高的皇室.
他也是向他們來問安.
在羅馬書的第十六章第十節.
又問基督裡經過試驗的亞比利安.
這個亞比利是誰我們已經不知道.
但是我們知道這個人是經過試驗經文化.
其實是原文希臘文documos這個字的意思.
就是認可這個意思.
是被神所認可.
經過試驗被神認可.
是一個很高的榮耀.
這個亞比利這個人的名字.
有學者在羅馬皇室的支持者名單上找到這個人.
究竟是軍事上的支持還是經濟上的支持.
我們已經沒有辦法可以知道.
羅馬書的十六章十節.
保羅向第一個家庭來問安.
他說問阿彌多布加里人的安.
根據歷史的研究有學者發現.
這個阿彌多布可能是希律王阿基伯一世的兄弟.
也就是說是希律大帝的孫子.
記不記得希律大帝.
希律首先是誰.
是羅馬王差派來管理猶太人.
這個希律大帝曾經要殺死耶穌基督.
所以他要下令將所有兩歲之內的男嬰全部殺了.
而這個阿彌多布就是他的孫子.
所以這個阿彌多布是屬於希律的家族.

$^{561}$根據羅馬人的律法.
當希律王或者希律的家庭去世的時候.
他們的家庭就會屬於羅馬王.
保羅就是向這個阿彌多布的家來問安.
有沒有看到保羅傳福音.
是很廣泛的.
不單只是向一個個人傳福音.
他還注意到這個人的家.
所以保羅為什麼傳福音這麼有力.
他傳福音的方法是值得我們去學習.
不單只是向個人傳福音.
還向家庭傳福音.
還有更多王室的人.
在羅馬書16章第11節.
又問我的親屬希羅天安.
問那其數家在主裡的人安.
希羅天是誰.
Herodian.
是希律黨.
大家記不記得耶穌基督在世的時候.
那些猶太人的領袖有什麼人.
你說民事法理錯人.
這些我們最記得.
還有呢.
殺到那個人我們都記得.
還有就是希律黨.
這些人是希律的支持者.
其中有保羅的親屬.
保羅和他們有密切的關係.
還有保羅問候了第二個家庭.
就是那其數家.
當時在羅馬最出名的那其數是誰.
是克勞狄烏斯的秘書.
這個人對羅馬王有很大的影響.
後來克勞狄烏斯被謀殺.
然後尼祿登上王位.
記不記得尼祿.
這是一個很殘暴的羅馬君王.
他繼承了克勞狄烏斯.
當克勞狄烏斯死了.

$^{601}$那其數仍然倖免了一段時間.
但最後他還是被逼自殺.
當他自殺的時候.
他所有的財產他的家庭.
也都屬於尼祿王.
這個就是我剛才所說.
這個是羅馬帝國他們的律法.
保羅也都認識這個家庭.
所以我就說保羅傳福音的手法非常廣泛.
不是針對個人而是整個家庭.
還有當他傳福音的時候也是向上的.
有沒有看到他認識很多王室的人.
這些也都成為他對象.
為什麼.
因為王室的人是有權勢.
是能夠影響很多人.
這個真的很值得我們去學習.
就是要廣泛.
也都要達到比較高的地方.
以致能夠有影響力.
其實在羅馬書的十六章八至十一節.
除了以巴盧和士大古.
我們不知道是誰.
其他都是和王室的人有關.
也都令我想起什麼.
在肥納比書的四章二十二節.
記不記得保羅說.
在該殺家裡的人突突的問你們安.
有什麼人向肥納比教會來問安.
就是羅馬王室的人.
所以保羅當他傳福音的時候.
他很著重向上.
那些有勢力的人能夠有影響力.
另外就是那些很廣泛.
注意到每個人的家庭.
你不要以為保羅就是這麼死賄.
永遠就是針對王室的人.
還有三個姊妹.
這三個姊妹是誰.
首先就是土肥拿氏.

$^{641}$還有是土賦薩氏.
這兩個人都為主奴婦.
但是還有一個姊妹.
她是為主多壽奴婦的比識氏.
比識是什麼意思.
是波斯.
這個姊妹就是波斯人.
她是多壽奴婦.
就是又出現Copia這個字出現.
其實如果我們見到這三個姊妹.
之間也有不同.
首先就是這兩個姊妹.
土肥拿氏和土賦薩氏.
他們只不過是受奴婦.
但是比識氏是多壽奴婦.
比他們受得更加多.
還有如果你注意希臘文的文法.
你會發覺.
說到土肥拿氏和土賦薩氏.
是現在式的文法.
說到毗拉氏是過去式.
她受苦是過去的.
說不定這個姊妹已經退休了.
但是其他兩個.
現在仍然為主是奴婦.
還有保羅向一個很出名的人問安.
這個人間接和耶穌基督有關係.
在羅馬書的十六章十三節.
又問主蒙簡選的奴婦和她的母親安.
她的母親就是我的母親.
有人說我見到Sherman都突了一突.
這個是不是保羅的母親.
是不是這個奴婦就是保羅的兄弟.
不是.
意思是說什麼.
我和他們很熟.
她的母親就好像我母親一樣.
我們現在都會這樣說.
這個人為什麼這麼出名.
記不記得耶穌基督是釘十字架.

$^{681}$那些羅馬官兵叫他揹著自己的十字架走上山.
他走到一半就支撐不住.
在馬可福音十五章二十一節說.
有一個古利內人西門.
就是亞歷山大和老夫的父親.
從鄉下下來經過那地方.
他們就勉強他同去.
很背著耶穌的十字架.
原來這個老夫就是當年古利內人西門.
替耶穌基督背十字架那個.
古利內人是什麼.
是非洲人.
所以他是一個黑人.
經文很有趣.
是主內蒙揀選的老夫.
每個信主都是蒙主揀選的.
為什麼保羅要這麼說.
因為有特別的意義.
剛才說到他有個大哥.
亞歷山大.
很可能亞歷山大沒有被主揀選.
不是一個信徒.
但是這個老夫就是揀選神.
馬可在羅馬寫馬可福音.
他的讀者就是羅馬人.
我們知道.
馬可福音成書的日期在羅馬書之後.
極有可能馬可是讀了羅馬書之後.
才寫馬可福音的五章二十一節.
既然提到老夫.
讓我給你們看到連帶關題.
這個老夫昔日是耶穌基督背十字架那位西門的兒子.
接著保羅在羅馬書十六章十四節.
另外又說了五個人名.
但我們都不知道這些是誰.
我也不會讀這些人的名字.
接著保羅又向一個人問案.
這就帶出了當年一件冤獄案.
在羅馬書十六章十五節.
又問菲羅羅古和尤利亞.

$^{721}$尼利亞和她姐妹同阿林巴.
並與他們一處的眾城圖案.
這些人我們都不認識.
但值得一提的就是尼利亞.
在公元九十五年有一件轟天動地的事發生.
就是一對夫妻.
丈夫是格利敏.
妻子是竇米提拉.
他們都是基督徒.
就是因為這樣而被逼伯.
竇米提拉是誰?.
是羅馬王維斯柏先的孫女.
是當時羅馬王多米田的姪女.
因為他們是基督徒就被逼伯.
她的丈夫格利敏就被處死.
這個竇米提拉就被放逐.
就算她是羅馬王的孫女又如何?.
她嫁基督徒就把她放逐到荒島.
這個尼利亞就是他們的財政司.
接著保羅說什麼?.
保羅說我真的要給你們一個聖潔的親嘴.
在《新約聖經》裡面有很多這樣的親嘴.
都是同性的親嘴.
給我們現在會怎樣?.
哎呀!這麼噁心!.
搞機嗎?.
當時是很普遍的.
這些是聖潔的親嘴.
我已經講完了很多人名.
這些是羅馬教會裡面保羅向他們問安的人.
保羅問完安之後要被他們縱谷.
只是講愛是不行的.
如果你只是彼此相愛.
你不講真理是沒有用的.
你只是講愛只是害了對方.
保羅離開以弗所教會的時候.
他和長老怎樣說?.
在《使徒行傳》的二十章二十九至三十節.
我知道我去了之後必有空步的豺狼進入你們之間.
不外適羊群就在你們中間.

$^{761}$也必有人起來說活謀的話.
要引誘門徒跟從他們.
保羅亦警告加拉太教會不要隨從別的福音.
加拉太書的一章第六節.
在哥林多前書的四章十四節.
他警告哥林多教會.
我當你們好像我親生的子女一樣.
來警告你們.
在哥林多前書的十一章十四節.
保羅和哥林多教會說.
你們小心.
撒旦亦裝模作樣.
好像是光明的天使一樣.
提摩泰前書的第五章.
和提摩泰後書的保羅.
都有同樣警告提摩泰要小心假勞師假先知.
保羅是用愛心來忠告羅馬教會.
任何所謂的愛.
如果藐視真理包容罪.
都不是真正的愛.
我再說多一次.
無論所謂的愛.
如果藐視真理容認罪.
都不是真正的愛.
任何埋沒真理正去講愛去講團結的話.
這些都是假勞師都是假先知.
在羅馬書的十六章十七節.
保羅說.
弟兄們那些離間你們叫你們跌倒.
背負所學之道的人.
我勸你們要留意躲避他們.
有人就斷章取義.
你看這裡保羅說.
不要讓人離間.
我們要團結.
不惜一切來團結.
這裡在說什麼.
離間你們.
不去聽那些講真理的人.
不是說一定要教會合一.

$^{801}$只是講愛沒有真理是沒用的.
這些都是假勞師假先知.
保羅接著給了兩個假勞師的特徵給我們看到.
第一個是什麼.
他們是為了自己的目的.
他們的目的是什麼.
羅馬書的十六章十八節.
因為這樣的人不服事我們的主基督.
只服事自己的肚腹.
掛著自己的肚腹.
令我想起什麼.
許冠傑那首歌.
為兩餐什麼都肯祭前世.
是不是.
喂.
真的.
這些人為了自己的肚腹什麼都肯祭.
肥納比書的三章十七到十九節.
保羅同樣警告要小心假勞師.
這些人為了自己的利益.
為了自己的肚腹.
一樣的字.
在彼得後書二章一至三節.
彼得都說.
有假勞師假先知.
因為他們貪心.
騙做言語.
要在你們的身上來取利.
聖經從來沒有被.
罵假勞師罵得更加狠的華語.
更加嚴厲的教導.
對著假勞師有非常之嚴厲的教導.
還有呢.
這些假勞師的後果是怎樣.
在羅馬書十六章第十八節.
他們用花言巧語誘惑那些老實的人.
這些假勞師他們想做什麼.
希望看到什麼後果.
敗壞教會.
哥林多教會就是有人希望敗壞哥林多教會.

$^{841}$現在羅馬教會保羅一樣.
你們要小心.
你說大家都是基督徒.
怎會有人希望敗壞教會.
被撒旦利用.
撒旦就是最不值得教會合一.
所以撒旦一定會興起假勞師來敗壞教會.
在羅馬書十六章十九節.
保羅很高興.
因為羅馬教會他們馴服的聲譽.
傳遍了各地.
每個人都知道他們的信心.
在羅馬書一章八節.
保羅也有說類似的話.
保羅說你們要善上聰明.
要在惡上愚拙.
羅馬書十六章十九節.
什麼叫善上聰明惡上愚拙.
你們對真理要聰明.
你們要認識.
但你們不需要認識其他的東西.
那些假勞師說的話.
很多人都會說.
哎呀 我就要看別人說什麼.
好像John McCarter經常說.
你明知道這個坑渠臭.
你還走下去幹什麼.
我時時都說.
你怎樣辨別偽鈔.
最重要就是對真正的鈔票有認識.
其他的偽鈔.
五花八門.
你看了一千張偽鈔.
始終還有一千零一張你是沒看過的.
所以你如果走去看五花八門.
你輸定了.
你只需要認識那張真正的鈔票.
你就可以了.
這就是保羅所說.
在善上聰明在惡上愚拙.

$^{881}$接著保羅就給了羅馬教會有他的安慰.
在羅馬書十六章二十節.
他說賜平安的神.
快要將撒旦踐踏在你們的腳下.
願我主耶穌基督的恩.
常和你們同在.
我已經說了兩點.
第一是AB.
第一是羅馬教會的人保羅向他們問安.
第二是保羅給他們的忠告.
不可以只是講愛.
你們一定要講真理.
要小心假路線.
第三點又回來問安.
但現在不是羅馬教會的人是誰.
是問安的人.
是保羅的同工.
第一個是誰.
提摩泰.
羅馬書十六章二十一節.
與我同工的提摩泰.
提摩泰不單止與保羅同工.
還是同心.
在肥納比書二章十九至二十二節.
保羅與提摩泰的感情非常深厚.
在哥林多前書.
在鐵撒羅尼加前書.
和提摩泰前後書都有講到.
除了提摩泰還有誰.
還有路求.
在羅馬書十六章二十一節.
講到路求.
很可惜和合本的翻譯錯了.
說路求是保羅的親屬.
其實是錯誤的翻譯.
有兩個可能性這是誰.
第一個是在《術史圖行傳》十三章.
講他們開始安提阿時宮的時候.
其中有五個領袖.
一個是古利內人路求.

$^{921}$剛才講了古利內人是什麼人.
是黑人是非洲人.
所以如果真是這個路求.
他可能是一個黑人.
還有一個可能性.
是路加醫生.
我們熟悉的路加.
在《術史圖行傳》第二十章.
我們知道是路加寫的.
他說我們就是我和保羅兩個在哥林多教會.
因為我們知道保羅在哥林多寫羅馬書.
他說我們就是路加醫生與他一起.
所以路求可能是路加.
而路求和路加兩個名字真是通用的.
還有呢.
剛才講到保羅是有親屬.
耶孫就是其中有另外一個親屬.
在羅馬書十六章二十一節.
和我的親屬耶孫索西巴德問你們的安.
在《術史圖行傳》十七章.
我們知道保羅去到鐵薩羅尼加.
接待他的是誰.
是他的親戚耶孫.
這是保羅第一次訪問鐵薩羅尼加.
是保羅帶領耶孫得救帶領他來信主.
耶孫之後跟著他的親戚保羅.
與他同行一起做事工.
你說索西巴德又是誰.
這個人的宿舍就是索巴德.
這個人可能是住在伯利恆城裡的一個舊約學者.
他對聖經是有研究的.
在《術史圖行傳》二十章第四節有告訴我們.
還有一個人我們都很熟悉的.
不是知道這個人的名字.
是知道他寫字.
他是幫保羅寫書信的得鵰.
在羅馬書十六章二十二節.
「我者代筆寫信的得鵰在主裡問你們的安」.
在哥林多前書十六章.
加拉太書的六章我們都知道.

$^{961}$保羅因為他的眼睛不太好.
他需要有人來代筆.
需要有人來幫他寫他的書信.
這個得鵰就是那個人.
在羅馬書十六章二十三節.
「接待我也接待全教會的該由問你們安」.
我們知道這個該由是一個外邦人.
他是保羅向他傳福音.
是保羅幫他施洗在哥林多前書的第一章十四節.
他不單是接待保羅.
他還是接待了全教會.
最後兩個人就是以拉都和豁土.
在羅馬書十六章二十四節.
「城內管銀庫的以拉都和兄弟豁土問你們安」.
以拉都是一個有身份有地位的人.
他是城市裡的財政官.
這個人的名字是非常普遍.
所以相信未必是《使徒行傳》十九章和提摩泰後書四章裡所講的同一個人.
有一件事很有趣.
就是美國雅典古典研究學院在一九二九年在哥林多.
發掘到一塊大理石.
是用來鋪路的大理石.
上方刻著甚麼.
「以拉都公共工程專員自費補設了這條路」.
相信是同一個人.
這個人是有身份有地位.
而且是有錢的.
所以他自資補了這條路.
這個就是以拉都.
豁土就是一個主內的兄弟.
我已經講完了四個部分.
三個部分都講完了.
最後就是回到保羅給羅馬教會的祝福.
在羅馬書十六章二十七節.
保羅說「願榮耀因耶穌基督歸於獨一全智的神.
直到永遠也」.
有沒有看到.
都是一些讚美頌,榮耀頌.
就是我們所講的Doxology.
現在我們經版完結之後.

$^{1001}$我們有三一頌.
其實這些讚美頌用來做甚麼.
講出神的屬性,祂的作為,祂的能力,祂的智慧.
讚美祂.
其實詩篇裡面都有讚美頌,榮耀頌.
遲些我會和你們講詩篇.
就是講到詩篇有五個部分.
這個是猶太人的拉比都注意到.
有五部書.
去到最後都是一些讚美的話語.
和阿們,阿們.
有兩次的阿們.
大概的經文都是說.
「耶和華是應當被稱頌的阿們」.
這些最後的讚美詩就在詩篇的41篇.
72篇,89篇,106篇和150篇.
剛剛將詩篇分開五個部分.
就是好像和摩西五經很相似.
這個遲些我講詩篇的時候會和大家分享.
聖經裡面還有新約還有很多的讚美頌.
我不能夠一一讀出來和大家分享.
但我會留在PowerPoint裡.
如果大家有心的話.
可以下載來看.
我剛剛開始的時候和大家講了.
整個羅馬書的十六章都是一個平衡對偶.
但是羅馬書如果你注意的話.
第一章和第十六章都是有一個平衡的結構.
都有三點A,B,C.
A是什麼?.
說出來是福音.
在羅馬書的一章一節說是神的福音.
在十六章二十五節保羅說是我傳的福音.
這個是A.
B那點是羅馬書的一章二節說.
這個福音就是神從前藉著眾先知所應許我們.
然後在羅馬書的十六章二十五至二十六節說.
這個福音是永古隱藏不言的奧秘.
這奧秘如今被彰顯出來.
你說福音為什麼是一個奧秘?.

$^{1041}$已不鎖書第三章又說.
外邦人就是在教會裡面的外邦人.
是能夠和以色列人同蒙應許同為後治.
這個就是一個奧秘.
在舊約是沒有講到.
我亦都講但以利書的時候我講過.
舊約的先知看預言的時候.
他們是怎麼看的?.
見到耶穌基督第一次的來臨.
耶穌基督第二次的來臨.
這些聖經裡面有很多的預言.
但是就沒有了中間的部分.
是什麼?.
講到教會的時期.
我們現在這個時候.
為什麼?.
因為是一個奧秘.
是舊約的先知看不到.
那奧秘是什麼?.
以前沒有人知道的事.
現在透過新約彰顯出來.
這個就是奧秘.
所以福音是一個奧秘.
保羅給我們看見很清楚.
這個福音要來做什麼?.
第三點的C講什麼?.
羅馬書的一章五節和十六章的二十六節都講.
是要我們信服真道.
你不是說只是得救你就算了.
是不是?.
很清楚經文告訴我們什麼?.
當我們信服真道的時候.
就是表現我們是一個真正得救的人.
反映出來我們是一個真正得救的人.
很多人都說自己相信.
但是如果他們不是真的信服真道的話.
他們是不是真正的相信?.
不是.
所以我們見到.
這裡是有一個平衡的結構.

$^{1081}$而我今天的主題是什麼?.
羅馬書的十六章二十五至二十六節.
福音或者神的話語.
是能夠堅固我們的心.
使我們信服真道.
在結束的時候.
我們一起來看看.
現在教會真的很不安定.
不是說我們的教會.
普世的教會都不是很安定.
很多教會倒閉.
很多教會關門.
不是外在的影響.
不是外在的動亂.
是什麼?.
是內心對真理的掙扎.
是我們不能夠得到安定.
被撒旦引誘我們.
勝過我們.
以至搞到現在教會有很多很多的問題.
但是保羅說什麼?.
福音神的話語能夠安定我們的心.
在羅馬書裡面.
保羅給我們有很多很多的教導.
來安定我們的心.
就好像提摩太后書的一章十二節.
因為知道我們所信的是誰.
也深信他能保存我所交付他的.
直到那日.
神交給我們的話語.
我們是要保存.
不單止保存.
我們還要怎樣?.
我們還要信服真道.
要將這個真道活出來.
我們怎樣能夠活出來?.
靠著聖靈的帶領.
堅固我們的心.
將我們的信心.
將真理活出來.

$^{1121}$願神祝福他自己的話語.
讓我們低頭禱告.
主啊.
給我們一部更美好.
是被聖靈密使的書.
教我們甚麼是恩順情義.
而保羅在他自己的生命.
在他的時空當中.
讓我們看到他怎樣活出.
恩順情義的道理.
而在第十六章.
保羅讓我們看到他和羅馬教會的關係.
彼此的愛心.
彼此的相愛.
但是當愛沒有了真理的時候.
是沒有用的.
保羅給羅馬教會有他的忠告.
他最後勸羅馬教會.
他說神的話語能夠安定你們的心.
願你們都能夠信服真道.
願意我們都是一樣.
要繼續打那美好的仗.
繼續守著那真理.
你給我們的真道.
願意當你回來的時候.
我們都能夠被你稱為一個忠心良善的僕人.
將那榮耀還賺得貴給你.
我們今天禱告.
奉主耶穌基督的名義而後.
阿們.
\newpage



\section{以賽亞書 42:1-9}
\label{sec:031747}
\textbf{看哪!我的僕人}~~~~周弘毅牧師
\newline
\newline
連結: \href{http://www.church.com.hk/acms/content.asp?site=cdc\&op=show\&type=product\&code=031747\&layout=sermon}{\texttt{031747}} ~~~~ 語音日期: 2026-02-14
\newline
\newline
\hyperref[sec:031746]{\small{< < < PREV SERMON < < <}}
~
\hyperref[sec:index_chronic]{\small{[返順時目]}}
~
\hyperref[sec:031748]{\small{> > > NEXT SERMON > > >}}
\newline
\newline
以賽亞書 42:1-9
\newline
\begin{longtable}{cl}
\hline
\hline
章節 & 經文 (和合本修訂版)\\
\hline
42:1 & \begin{tabularx}{0.7\textwidth}{X} 看哪，我的僕人，我所扶持、所揀選、心裡所喜悅的！我已將我的靈賜給他，他必將公理傳給萬邦。 \end{tabularx} \\ \\ \relax
42:2 & \begin{tabularx}{0.7\textwidth}{X} 他不喧嚷，不揚聲，也不使街上聽見他的聲音。 \end{tabularx} \\ \\ \relax
42:3 & \begin{tabularx}{0.7\textwidth}{X} 壓傷的蘆葦，他不折斷；將殘的燈火，他不吹滅。他憑信實將公理傳開。 \end{tabularx} \\ \\ \relax
42:4 & \begin{tabularx}{0.7\textwidth}{X} 他不灰心，也不喪膽，直到他在地上設立公理；眾海島都等候他的訓誨。 \end{tabularx} \\ \\ \relax
42:5 & \begin{tabularx}{0.7\textwidth}{X} 那創造諸天，鋪張穹蒼，鋪開地與地的出產，賜氣息給地上眾人，賜生命給行走其上之人的神耶和華如此說： \end{tabularx} \\ \\ \relax
42:6 & \begin{tabularx}{0.7\textwidth}{X} 「我－耶和華憑公義召你，要攙扶你的手，保護你，要藉著你與百姓立約，使你成為萬邦之光， \end{tabularx} \\ \\ \relax
42:7 & \begin{tabularx}{0.7\textwidth}{X} 開盲人的眼，領囚犯出監獄，領坐在黑暗中的出地牢。 \end{tabularx} \\ \\ \relax
42:8 & \begin{tabularx}{0.7\textwidth}{X} 我是耶和華，這是我的名；我必不將我的榮耀歸給別神，也不將我所得的頌讚歸給雕刻的偶像。 \end{tabularx} \\ \\ \relax
42:9 & \begin{tabularx}{0.7\textwidth}{X} 看哪，先前的事已經成就，現在我要指明新事，告訴你們尚未發生的事。 \end{tabularx} \\ \\
[1ex]
\hline
\hline
\end{longtable}
$^{1}$我其實不是一個很熱愛大自然的人.
所以我很佩服那些很喜歡行山,看大自然的人.
特別是揹著很多工具.
望遠鏡,相機,長鏡頭.
去觀鳥.
不知道大家有沒有去過觀鳥.
看鳥鳥.
我記得很多年前.
我曾經去過米埔的濕地公園.
如果你有去過的話.
可能你有印象.
你走著走著.
其實那裡有些小屋.
那些屋你可以走進去.
看那些鳥鳥.
甚至因為裡面已經架設了很多望遠鏡.
你就看到很多的候鳥在那裡.
其實都挺漂亮的.
我是一個門外漢.
我不是很知道什麼鳥打什麼鳥.
搞不定的.
甚至怎樣搞那些工具.
那些望遠鏡都是一樣那麼不行的.
所以我以為.
當你看到那些鳥很遠.
好像挺漂亮.
我就將望遠鏡的鏡頭向著它.
希望近一點.
看到它的動態.
但是發現原來看不到的.
找不到鳥.
你明白嗎.
於是你就嘗試移動望遠鏡.
搞來搞去.
沒理由那麼難.
接著就知道原來你要調校焦點.
一邊移動一邊調校焦點.
終於對準.
接著就飛走了.
真的很煩.

$^{41}$所以在過程當中.
其實不是那麼容易的.
如果你要看到那些鳥鳥的形態.
已經很不容易.
更加難是我想像.
這件事我沒有做過.
觀星.
不是更難嗎.
黑漆漆.
你要找出那顆星.
我在年初和大家說.
異塞亞書的背景的時候.
我曾經說過.
異塞亞先知其實是活在一個.
面對外敵威脅的.
一個很動盪不安的大時代.
在異塞亞書36至37章.
曾經描述亞述王西拿基納.
來攻打耶路撒冷.
不過那次其實很感恩.
危機因為神的主動介入解決了.
不過解決了一次危機.
其實不代表什麼.
因為先知說了預言.
其實你們最後都要接受.
被擄亡國的結果.
對於我們來說.
可能在我們的生命歷程當中.
我們都曾經面對過不同的困難.
現在可能仍然在一個生命的低谷.
屬靈的黑夜當中.
如果在這麼黑暗的環境.
我們究竟怎樣能夠找到星光呢?.
當面對國家最黑暗的時候.
猶太人不知道究竟他們有沒有機會見到黎明.
就在這個時候.
神透過耳塞先知.
送了一首獨人之歌給他們.
這首歌的內容.
不是教他們怎樣去找到星光.

$^{81}$怎樣去找到光明.
其實不用像我官僚那麼辛苦.
弄來弄去都找不到.
甚至漆黑一片.
在茫茫的黑夜困難當中.
要找到那一點光明.
不用那麼辛苦.
耳催向這些猶太人宣告.
光明反過來是來找他們的.
所以不是靠他們的人力力量.
靠自己來找到光明.
所以今天和大家分享的.
《耳塞書》42章第一至第九節.
其實就是講這個內容.
好,我想先請大家和我一起讀一讀這段聖經.
預備起.
看吶,我的僕人,我所扶持所揀選.
心裡所喜悅的,我已將我的靈賜給他.
他別將功利傳給外邦.
他別勸揚,不揚聲.
也不使街上聽見他的聲音.
壓傷的爐位,他不截斷.
將殘的燈火,他不催滅.
他憑真實將功利傳開.
他別灰心,也不喪膽.
直到他在地上設立功利.
海島都等候他的訓誨.
創造諸天,鋪張窮倉.
將地和地所出的一併鋪開.
似氣息給地上的眾人.
又似靈性給行在其上之人的神耶和華.
他如此說.
我耶和華憑公義召你.
必參乎你的手,保守你.
使你作眾民的忠報,作外邦人的光.
開闔子的眼,領被囚的出牢獄.
領坐黑暗的出監牢.
我是耶和華,這是我的名.
我必不將我的榮耀歸給假神.
也不將我的稱讚歸給挑剔的偶像.

$^{121}$看吶,先前的事已經成就.
現在我將新事說明.
這事未發生先,我就說給你們聽.
前一個星期.
余姑娘就跟我們分享《以下書》的第四十一章.
那時她只是說了前半節.
一次到第二十節.
是表明耶和華是一個滿有恩典.
他是歷史的主宰.
其實當中的表述.
就有點像在法庭審案時.
控辯雙方的辯論.
先知告訴列國.
即經文四十一章第一節.
列國形容他們是「縱海島」.
告訴他們神是掌管一切的.
無論以色列或是異邦.
基本上全部都在祂的掌握之中.
即使由第二十一節.
前一個星期余姑娘沒時間說下半節.
下半場其實是說一些假神的質詢.
正因為這些假神全部都是虛空.
所以他們的辯解基本上是沒用的.
經過了四十一章辯論之後.
就來到今天我要跟大家分享的第四十二章.
轉而是講到神如何為人帶來安慰.
在第四十二章一開始就說.
「看吶,我的僕人」.
到底這個「僕人」其實是誰呢?.
「僕人」這個字在《倚創下書》裡.
第四十至五十五章.
一共出現了十九次.
而在很多人的心裡.
一想到這個「僕人」.
當然是寓表主耶穌.
很直接的連繫就放在那裡.
是寓表著耶穌.
不過如果我們看一看四十二章第十九節.
我今天又沒去到那一段那麼後.
講到經文這樣說.

$^{161}$「誰比我的僕人眼瞎,誰比我差遣的使者耳聾?.
誰瞎眼像那與我和好的,誰瞎眼像耶和華的僕人?」.
如果《倚創下》才有這樣的描述的話.
其實我們就很明白.
沒有理由說這個「僕人」是耶穌.
形容祂為「眼瞎耳聾」不應該是這樣的.
再加上有一些描述「僕人」這個字.
其實是複數.
所以它也可以說.
亞國甚至是整個以色列民族.
或者是那些散居列國的猶太人.
另外也有醒學者說.
「這些僕人有些地方是說一些特定的歷史人物也說不定」.
我全部列出來給你看.
所以「僕人」這個稱號.
你不要這麼快就黏在主耶穌身上.
可能在不同的地方其實有不同的解釋.
今天我和大家分享這一段經文.
醒學者叫它做「僕人之歌」.
由42章到53章一共有4首「僕人之歌」.
除了我今天和大家分享的第一首42章之外.
其他的已經列出來給你看.
如果我們回到第42章的第一至第四節.
剛才大家讀的時候.
我不知道你覺不覺得有點熟口熟面.
好像在哪裡看過.
答案是在《馬太福音》.
《馬太福音》就將它引用在耶穌身上.
我們試一下將《異賽亞書》42章.
和《馬太福音》12章相關的經文.
合起來看一看.
你就會知道原來是這樣.
看不看到?非常相似.
我們再看另外兩節.
到底耶穌這個僕人有什麼特質呢?.
在第一節中間的部分是這樣說的.
「我所扶持所謹選心裡所喜悅的.
我已將我的靈賜給他」.
這個描述又是熟口熟面.
因為它是在說耶穌受洗的時候的景象.

$^{201}$《馬太福音》三章十六至十七節.
你看看.
「耶穌受了洗.
隨即從水裡上來.
天忽然為他開了.
他就看見神的靈.
彷彿甲子降下.
落在他身上.
從天上有聲音說.
這是我的愛子我所喜悅的」.
接著到第一節的下半節.
最後那句.
「他必將公理傳給外邦」.
要預言這位主耶穌.
他有能力將公理傳給外邦.
這個托福是可以被完成的.
「公理」這個字在短短第一至第四節經文已經出現了三次.
所以可以說是整首《牧人之歌》的關鍵字之一.
「公理」.
「望翻」.
「以查書」.
不同的經文其實有不同的翻譯.
「公平」.
「公義」.
「審判」.
「冤屈」.
「辯論」等等.
神要告訴外邦人.
即是你和我.
即使我們在面對不同的困境.
生命裡面很大的艱難.
神說他的僕人仍然在.
神的僕人.
雖然要成就大事.
但這位主耶穌不是用了一個轟轟烈烈的方式來做事.
即使祂有能力.
祂有資格得到最大的榮耀.
不過祂選了一個卑微僕人的身份來世人中間.
42章第二節.
接著這樣說.

$^{241}$「他不謙讓,不揚聲,也不使街上聽見他的聲音」.
現今的世代.
聲大鑒惡.
自我抬舉.
叢林法則.
這些是現在這個世界的運作方式.
不過這裡三次提到這位僕人.
「不…」.
只是要凸顯了.
其實他和世人是不同的.
來到第三節.
這一節熟得不得了.
「壓傷的爐位他不截斷,將殘的燈火他不摧滅,他憑真實將功利全開」.
我們很容易就有一個解釋.
這個僕人真是體恤那些正在面對很艱難的處境的人.
這樣就解決了.
但是值得我們仔細想一想.
究竟壓傷的爐位是怎樣的.
相信大家都知道爐位草是空心的.
大家知不知道?.
它空心.
所以其實它是很容易被大風又或者很急的水流.
一過來它就斷了.
很容易斷的一些植物.
甚至如果有些動物要去河邊喝水.
一踩又斷.
所以它的構造是非常脆弱的.
而且爐位是不會好像其他植物一樣.
如果有一棵植物的枝斷了.
它就在旁邊生第二個出來.
但是爐位草是不行的.
斷了就是斷了.
是沒有辦法去修復的.
所以當這裡以剎先知說壓傷的爐位.
它當然不是說植物.
其實它是在說人.
它是要告訴讀者其實人是很脆弱的.
每個人都是很脆弱的.
很容易會跌倒受傷的.
在我們身邊.

$^{281}$甚至我們在面對很大的困局的底下.
很容易就會被那些排山倒海的問題.
將我們壓扁了.
除此之外.
以剎先知又有第二個描述.
將殘的燈火.
又值得想一想其實是什麼來的呢?.
在以前當然沒有電力供應.
沒有電燈.
晚上唯一的照明方法.
其實就是油燈.
如果忘記幫這盞油燈.
預備足夠的油.
當油差不多用完.
這盞燈很明顯越來越暗.
接著就閃一閃.
最後在它熄滅之前.
其實燈芯會出煙.
我們現在不用油燈.
不過我不知道你有沒有試過點蠟燭.
點到尾有個類似的觀察.
其實和壓箱的爐位一樣.
將殘的燈火.
都是在說人的狀況.
這一類人可能在最初的熟齡境況都不錯.
在那個階段其實生命可以發出光芒.
很積極參與服事.
很貼近神.
和上帝的關係都不錯.
不過隨著時間的流逝.
燈的油減少了就跌鑊了.
燈閃一閃.
原本得救的平安喜樂就不見了.
他們和主的關係.
其實只不過是剩下很微弱的光線.
大風少少一吹就熄滅了.
從靈性的角度來看.
如果這個火所剩無幾的話.
其實還哪有力量來面對困難和挑戰.
作為一個屋主.

$^{321}$晚上油燈差不多熄滅了.
開始出煙就吹熄它去睡覺.
不過以賽亞先知所說的這個僕人不是這樣.
他不會由那盞油燈來熄滅.
他不會摧毀我們剩下來的微小信心.
主耶穌的心意就是剪好一條燈芯.
幫盞燈加上油.
潑一下它.
令到火可以燒黃.
讓它重新發光.
甚至比以前更光.
我不知道你有沒有認識有些弟兄姊妹.
曾經在他們的生命下跌.
然後重新站起來.
有沒有.
那個生命是可以比以前更好.
因為上帝在他們的生命裡有恩典.
我作為牧者很多年.
神給我機會和不同的弟兄姊妹.
曾經走過在他們生命裡一些很大的難處.
我又一次又一次見證著.
神在這些弟兄姊妹生命裡的奇妙作為.
在很多沒有辦法的處境裡.
當然我見這些弟兄姊妹的時候.
一定會和他們分享聖經.
作為牧者不會就這樣和他們聊天.
然後拜拜.
一定會和他們分享聖經講什麼.
不過其實很多時候當我說完之後.
弟兄姊妹的回應就是.
我也知道.
甚至他們可以告訴我.
所以在那些時候.
我想說的是我不知道要說什麼.
沒什麼好說的.
唯有和他們一起來祈禱.
問題是你講完之後不會走了.
仍然在這裡.
那就祈禱吧.
因為我心裡的信念就是.

$^{361}$「百傷的蘆葦他不折斷,漿殘的燈火他不催滅」.
神是不會讓這些弟兄姊妹.
跌入心坑當中然後跌死的.
這不是我們的上帝.
之後奇妙的事就發生了.
下次見到他甚至不用下次.
在我們一起祈完禱之後.
當晚我就收到這些弟兄姊妹的訊息.
告訴我.
神怎樣回應了我們迫切的禱告.
不是說那個很大的問題整個解決了.
不過你將問題拆開細份細份細份.
似乎有些位置有點起色.
所以無論是什麼原因.
令我們那盞油燈開始閃爍.
屬靈生命開始出現了.
第四節告訴我們.
「這位僕人仍然不會灰心,不會喪膽」.
以下才知道.
故意用「不灰心」來回應「壓傷」.
為什麼呢?.
因為「灰心」和「壓傷」在原文希伯來文裡是同一個字根.
它同樣也用「不喪膽」來回應「張顫」.
原因是因為這兩個字的字根也一樣.
以賽亞先知清楚表明.
我們不用也不能夠靠自己來解決問題.
這位僕人最終必定會完成他的使命.
將真理傳開.
為的就是讓我們重拾救恩給我們的平安喜樂.
來到第五節.
以賽亞再一次提醒讀者.
「神是怎樣的?創造諸天,鋪張窮倉.
將地和地所出的一併鋪開.
似氣息給地上的眾人.
又似靈性給行在其上之人的神耶和華」.
他用了「創造」這個描述.
放在耶和華這個名字前面.
他不是耶和華行先,是反過來的.
是要強調神是那位掌管歷史主權的那位.
所以昔日的以色列和今天的你和我.

$^{401}$其實不需要懷疑.
接著去到第五節,第七節,下半節.
接著這樣說.
「他如此說,我耶和華憑公義召你,必參乎你的手,保守你.
使你作眾民的中保,作外邦人的光.
開闔子的眼,領被囚的出牢獄,領坐黑暗的出監牢」.
神在第六節,上半節.
很清楚地告訴讀者.
他要重申,其實他在第一節已經說過.
這個扶持和揀選.
因著神在這裡說他是呼召和說他是參乎.
這位僕人是要承擔起作眾民的中保.
請大家留意後面有個括弧.
「中保,原民是約」.
這裡說作中保的對象是眾民.
意思是什麼?.
我們很明白其實神在很早的時候.
創世記的時候.
已經和以色列人的先祖來到立約.
不過這裡要告訴我們.
這個聖約其實對象不只是以色列人.
還包含外邦人.
正因為這裡所說的僕人是耶穌基督.
他的犧牲帶來的救贖恩典.
是要給萬國萬民.
就正如大家很熟悉的.
《約翰福音》三章十六節.
「神愛世人,甚至將他的獨生子賜給他們.
叫一切信他的不至滅亡,反得永生」.
除此之外,接著在第六節的下半節.
也講到相類似的信息.
這裡說這位僕人是要作外邦人的光.
以賽亞當時宣講的對象是猶太人.
他們那時候是被亞述.
即是外邦人圍攻.
所以以賽亞告訴當時的猶太人.
這位僕人要作外邦人的光.
並且第七節還要說.
「開盒子的眼,令被囚的出牢獄.
令坐黑暗的出監牢」.

$^{441}$有沒有搞錯?.
我相信這些猶太人一定百思不得其解.
為什麼要看顧這些攻擊他們的外邦人?.
你不是應該幫我嗎?.
有什麼理由幫他們?.
不過如果我們看整套電影的話.
其實看完之後舊約.
講到關於以色列的民族歷史發展.
知道波斯王古烈.
是容讓這些被擄的以色列民.
從巴比倫可以回到故鄉.
我們就明白原來上帝在歷史裡.
是有祂奇妙的作為.
你想也想不到.
大家可能會比較仔細想想.
以色列民當時被擄的時候.
有坐過牢嗎?.
為什麼要出牢獄或者出監牢?.
聖經沒有講得很清楚.
但是如果我們看考古.
發掘出來亞述王西那基納.
有一些例版掘了出來.
那裡其實有提過.
在有一次戰役當中.
亞述奪取了以色列民四十六座的城邑.
擄走了超過二十萬人.
然後捉了他們去坐牢.
這份例版是這樣寫的.
所以我們不要說.
聖經跟歷史是不對的.
不是的.
我們嘗試了解多一點.
我們就明白.
好了 去到新約.
大家來到去看路加福音四章十八至十九節.
我不知道你有沒有印象這段聖經.
主耶穌在這兩節的經文說.
主的靈在我身上.
因為祂用膏告我.
叫我傳福音給貧窮的人.

$^{481}$差遣我報告被擄的德釋放.
乞眼的德看見.
叫那受壓制的德自由報告神.
閱立人的希臘.
印象主耶穌在路加福音裡.
引述以賽亞書的宣告.
這四十二章第七節.
其實多了一重的屬靈意義在其中.
我們嘗試仔細想想.
什麼叫做黑暗.
或者我嘗試用另外一個向度去問.
究竟你會怎樣定義黑暗.
什麼叫做黑暗.
黑暗如果用科學角度來說.
就是發出的或者反射的光.
比周圍少.
於是我們就看不清楚.
甚至是漆黑一片.
所以簡單來說.
不是有一種東西叫做黑暗.
一種物質叫做黑暗.
而是缺少了光就是黑暗.
如果黑暗是一種物質.
是一些東西.
原本黑暗我加多些物質.
那就黑暗多些.
不會這樣的.
在這裡.
以賽亞在第七節.
連續用了三個類比.
克子,被囚和黑暗.
其實是在說同一個指向.
就是在說那些在屬靈上.
缺少了光和神失聯的人.
當中當然是包括那些未信者.
因為他仍然在罪當中.
不過就算我們已經信了耶穌.
我們有些時候.
仍然會在情感某些處境.
甚至在靈性上.

$^{521}$我們覺得很孤寂無助.
我們不覺得上帝在身邊.
神你去了哪裡?.
為什麼你不出手?.
我們不覺得上帝看著我們.
這位僕人就是要將我們.
我們相信在在在當中.
絕大部分我們都信了耶穌.
這位僕人就是要將我們.
在失聯狀態的黑暗當中.
將我們拯救出來.
主耶穌在楊方八章十二節說.
「我是世界的光.
跟從我的就不再黑暗裡走.
我要得著生命的光.
面對著生命當中的地獄.
又或者靈性的倒退.
這就叫做黑暗」.
其實我們跟上帝的關係.
很多原恩惠導致我們倒退.
你想想.
可以因為我們在靈修.
用心讀聖經.
甚至是投入教會生活.
在事情上放軟了手腳.
廣東話叫做「比較敷衍」.
究竟我的信仰狀態.
是不是有點敷衍呢?.
這個就等聖靈提醒你.
又或者.
我們有些事情.
我們非要抓緊.
跟上帝吵架.
我就是不喜歡.
我不想.
我想用我的方法去做這個決定.
我想用我的生活方式去過每一天.
我就是不放得下.
又或者.
用養罪悄悄進入了我的生命當中.

$^{561}$諸如此類.
如果我們在這個處境裡.
有耶穌成為我們生命的光.
我們就能夠有方向.
祂就能夠幫助我們走出黑暗.
因為在以賽亞的年代.
有不少住在耶路撒冷的猶大民眾.
仍然在質疑上帝的能力.
都不知道行不行.
所以到了第八節.
先知又再次提醒他們.
我是耶和華.
這是我的名.
我必不將我的榮耀歸給假神.
也不將我的稱讚歸給雕刻的偶像.
這位叫做耶和華的神.
珍珠都沒那麼真.
祂才是那位掌管世界真正的力量.
因為在第九節繼續說.
先前的事已經成就.
現在我將新事說明.
這事未發生以先.
我就說給你們聽.
先前的事和新事.
在歷史進程當中.
其實就是分別指猶大.
這個國家裡的人.
將會被擄去巴比倫.
還有古列將要擊敗巴比倫解放猶大.
因為今天其實我們已經認識了主耶穌.
所以除了歷史的意義之外.
還加上一重的屬靈意義在其中.
以前因為罪的緣故.
我們成為了奴隸.
這是先前的事.
但是藉著主耶穌在十字架上.
為我們做的新事.
我們就被釋放.
當我們願意去回應的時候.
耶和華這位真神.

$^{601}$而不是第八節.
剛才說的那些假神偶像.
會得到稱讚.
神不單止要以色列民讚美祂.
祂更加期待願意看見萬國萬民.
包括外邦人來讚美祂.
在上一章的時候.
其實就是提到神的僕人以色列.
是一個不完全的僕人.
他會犯罪的.
可以說他不是很稱職的這個僕人.
但是以色列在神彰顯.
他在歷史裡面的作為.
上帝仍然用他.
擔當了一個不可或缺的角色.
可能是負面一點.
不過仍然有一個角色給他.
今天主耶穌這個僕人已經成就了救恩.
所以來到42章.
其實神對我們的期望就是.
我們同樣是不完全的.
我們這些不完全的僕人.
神盼望我們能夠繼續成為外邦人的光.
可以讓救恩傳到地極.
如果我們繼續看下去.
去到新約整個歷史的發展.
我們看到紹行傳一章第八節.
復活的主在祂升天之前.
他是差遣使徒直到地極作我的見證.
那裡所說的地極和見證.
其實你將這條時間線一直拉向後.
是當年以賽亞向亡國被老臣的指紋.
一個很遠很遠之後.
能夠實現的應許.
最後我想做一個很簡單的總結.
以賽亞書42章.
這首僕人之歌分成兩個段落.
第一至第四節要告訴我們.
耶穌基督這個僕人.
是選擇了一個非常卑微的身份來到世間.

$^{641}$在我們面對著很大的艱難的時候.
快要壓扁了.
或者我們靈性的那把火.
都是微弱得不得了.
所剩無幾的時候.
祂提醒我們不要靠自己.
因為這位僕人最終必定能夠完成他的使命.
將真理傳開.
為的是使人能夠得著.
其實已經給了你原先你應該有的平安喜樂.
不過不知為何走著走著就跌落坑渠.
去到第五至第九節.
祂說自己和外邦人有約.
是為了他們帶來光明.
因為我們已經認識了神.
如果在這一刻我仍然意識得到.
我在黑暗當中.
你意識不到就無話可說.
不過如果你意識得到.
你想從這個黑暗走出來.
求主幫助我們.
因為這位僕人能夠讓我們在生命裡重新發光.
求主幫助.
讓我們在彰顯他的作為上.
做好這個僕人的角色.
按照教會的年曆.
幾天之後.
不知大家有沒有留意.
2月18日.
其實就是今年的御苦期的開始.
御苦期另外一個名叫大齋期.
這是傳統上紀念耶穌基督受苦受死的節期.
當中一共有40天.
不計主日.
由聖灰日開始一直數.
數到復活節前一天.
40天.
是這樣計算的.
不計主日的原因.
就是因為每一個主日.

$^{681}$其實都是紀念主耶穌的復活.
好了.
盼望我們透過.
《以察書》42章.
第一首的僕人詩歌.
提醒我們.
主耶穌為我們擺上了的一切.
容讓我們在這段時間.
預備我們的心.
好好自省.
\newpage



\section{創世記 25:21}
\label{sec:031748}
\textbf{聖經人聖經事:祈求耶和華}~~~~黃克勤牧師博士
\newline
\newline
連結: \href{http://www.church.com.hk/acms/content.asp?site=cdc\&op=show\&type=product\&code=031748\&layout=sermon}{\texttt{031748}} ~~~~ 語音日期: 2026-03-01
\newline
\newline
\hyperref[sec:031747]{\small{< < < PREV SERMON < < <}}
~
\hyperref[sec:index_chronic]{\small{[返順時目]}}
~
\hyperref[sec:031749]{\small{> > > NEXT SERMON > > >}}
\newline
\newline
創世記 25:21
\newline
\begin{longtable}{cl}
\hline
\hline
章節 & 經文 (和合本修訂版)\\
\hline
25:21 & \begin{tabularx}{0.7\textwidth}{X} 以撒因他妻子不生育，就為她祈求耶和華。耶和華應允他的祈求，他的妻子利百加就懷了孕。 \end{tabularx} \\ \\
[1ex]
\hline
\hline
\end{longtable}
$^{1}$時間過得好快.
不是時間這麼簡單.
是每一個鐘頭都很快.
又第二天了.
今天又是生命樹.
和大家一起讀聖經.
看聖經的人與事.
和更多認識背後的上帝.
今次直接講祈求耶和華.
是今天的主題.
祈求耶和華.
我們看看.
創世記二十五章第十九節.
這裡由亞伯拉罕講下來.
亞伯拉罕離開世界.
撒拉離開世界.
亞伯拉罕離開世界.
以實瑪利離開世界.
亞伯拉罕的兒子以撒.
這裡開始講以撒.
以撒的後代.
是記在下面的.
上面講以實瑪利的後代.
這裡講以撒的後代.
亞伯拉罕生以撒.
這裡大家知道了.
以撒和誰呢.
以撒娶了利伯加.
以撒的名字是什麼意思.
記不記得.
以撒就是他笑.
可能是他出生的時候.
改什麼名字好呢.
然後小孩嘻嘻.
可能叫他笑.
以撒.
利伯加.
recap.
利伯加是誘捕者.
引誘.

$^{41}$然後捉他.
誘捕者.
可能是這個名字.
可能小時候的利伯加.
很擅長捉小鳥.
很擅長捉動物.
以撒娶利伯加為妻的時候.
以撒是四十歲.
四十歲當時是很年輕.
利伯加是巴旦亞蘭地的亞蘭人.
比土里的女兒.
是亞蘭人拉班的妹妹.
他們多數都要這樣講.
原因是沒有姓.
以撒,利伯加.
可能利伯加這個名字很多.
講利伯加是哪位.
因此要講一下.
他是哪裡的人士.
他是哪裡的女兒.
他是哪裡的妹妹.
這樣就解釋得清楚一點.
以撒因為.
第一件想和大家講的事.
以撒因為他妻子不生育.
原來利伯加也是不生育.
這樣就真得意了.
我自己對這第一點.
我自己的看法就是.
為何他不能生育呢.
上帝讓他不能生育.
如果講得難聽一點.
上帝令他不能生育.
上帝讓他不能生育.
他就沒有小孩.
當日沒有小孩很慘.
誰覺得慘呢?.
以撒覺得不慘呢?.
以撒覺得可惜.
而最覺得不開心的是太太.

$^{81}$因為在男女夫婦來說.
如果女兒無法生育.
無生小孩.
就覺得這當然是品性不好.
神不要祝福她.
或者講這些.
於是就有壓力.
很慘不生育.
這裡也是講妻子不生育.
於是以撒為他祈求.
第二點.
有一件難處.
上帝讓你遭遇的難處.
這個難處就會出現.
不過這個難處會令你更加認識神.
更加信靠神.
不要太抗拒.
為何這麼慘.
這麼不公平.
為何他有小孩而我沒有.
為何他的眼睛大我的眼睛小.
我們很多時候比較其實不好.
是上帝容讓那件事.
你要知道他的手是很美善的.
他的想法是很好的.
是善良的上帝.
重要的是他的美意.
你要接受這件事.
於是以撒為他祈求.
留意這個字.
「祈求」.
忘記撇了.
「阿花」.
「阿花」是「懇求」.
不單是祈求.
祈求就是禱告.
禱告阿主母仔.
求主讓她生一個.
「懇求」就是神.
懇切祈求.

$^{121}$可不可以給我們一個小孩.
很需要很想.
懇切祈求.
以撒為他「懇求」.
「懇求」多久?.
我不知道.
沒有說.
「懇求」.
你要知道.
有事發生很困難很艱難很痛苦.
你要做的就是「懇求」.
「懇求耶和華」.
「於是耶和華就應允他的祈求」.
「應允」這個字.
我也想解釋一下.
挺有趣的.
「祈求」就是「懇應允」.
「祈求」就是「應允」.
「應允」這個「呢」.
然後就是這樣.
「耶和華就應允他的祈求」.
即是「祂來求我」.
「耶和華祂來求我」.
「我就應允祂」.
就是應允祂求的東西.
「祂來求我」.
「我就應允祂」.
不是應允祂.
「向著祂的祈求」.
即是「耶和華臣是拿起祈求」.
「向著」或者「into」.
「耶和華是向著祂的祈求」.
即是「聽祂禱告」.
所以中文翻譯叫「應允」.
其實原來的字是「向著」.
「耶和華是向著祂的祈求」.
不是「背著」.
「背著」就糟糕了.
「背著」即是不聽.
「向著」即是聽.

$^{161}$「耶和華是向著耳塞的祈求」.
「於是他的妻子尼伯加懷孕了」.
第一.
遭遇一個很大的難題.
很大的痛苦.
很多感受.
記得「懇求耶和華」.
如果耶和華是向著我們的祈求.
事情就會成就.
交託給祂.
「懷了孕,孩子們在尼伯加的福中彼此相憎」.
兩個孩子胎兒在肚子裡打架.
彼此相憎.
大家爭.
怎樣爭呢?.
可能在肚子裡打架.
兩個孩子在肚子裡打架.
會怎樣呢?.
可能會起誓拳來腳往.
「揚」.
「揚」到尼伯加受不了.
很辛苦.
「哎呀!搞什麼?兒子」.
「為什麼兒子這麼活躍,這麼「揚」呢?」.
其實他不知道裡面有兩個兒子.
他以為兒子為什麼經常動.
手腳日夜地「揚」.
「為什麼會這樣?」.
辛苦到有點埋怨.
「哎呀!若是這樣,我為什麼活躍呢?」.
「哎呀!我真的想不到」.
「我很想要個兒子又沒有」.
「現在有了,真是「揚」到我死了」.
「不如如果沒有就算了」.
「「揚」死我了,為什麼活躍呢?」.
「不如死好過了」.
於是這個時間.
學到東西.
尼伯加嘗試以殺為他代討.
這次是他自己去求問耶和華.

$^{201}$這件事先是以殺懇切的代討.
然後是尼伯加自己痛苦到不得了.
於是幾乎把腳「揚」出來.
他就去求問耶和華.
求問.
問神「為什麼會這樣?」.
「很有趣,神有沒有回答他呢?」.
「沒有,不過告訴他」.
「有兩個國家在你的肚子裡」.
「有兩個民族要從你的身上出來」.
「那就是什麼?馬仔」.
「這族必強於那族」.
兩個是一個強一個弱.
「將來大的要服侍小的」.
「第三點要說的是命定」.
「神命定雖然人兩兄弟在爭」.
「但是神已經命定」.
「大的要服侍小的」.
於是事情就這樣發生了.
《創世記》二十五章二十一字.
我跟大家分享的金句.
「以殺因他妻子不生育」.
大家記不記得祈求是什麼意思?.
「懇求」.
懇求耶和華.
耶和華轉生.
英文是向著祂祈求.
於是尼伯加懷孕了.
請大家記住.
就像以殺一樣.
有困難有問題.
神懇求你.
「哀」.
到了尼伯加不明白.
為什麼肚子壓得這麼厲害.
「求問耶和華」.
又求問.
沒有去看婦產科.
「求問耶和華」.
這個要學習祈禱.

$^{241}$祈禱就是有時候神問你的事.
不是你問神的事.
有時候是你祈求問神的事.
神有回答.
但沒有直接回答你問題的答案.
沒有直接回答你問題的答案.
祂不回應你的問題.
祂只回答你最後一步的事.
最後一步.
現在喜歡用「matter」.
就是「之後的事」.
說「之後的事」.
是兩個胎.
是雙胞胎.
在裡面爭.
爭是彼此之間的關係.
是命定的.
而且大的要服侍小的.
是命定的.
我可以盡量懇求神會聽.
但在聽的時候.
神是不照著你的心意去聽.
不是你寫了紙給上帝的命令.
我以前在台灣宣教時.
有人說我已經下了字條給他.
下了字條給他.
即是我寫了兩個字.
「你做什麼啦?」.
是上層的官員下字條給下層.
上層高級的公司下字條給.
年輕的工作人員.
「我已經下了字條給你了」.
「你都不做了」.
這樣罵那個人.
「我已經下了字條了」.
「下了字條要做?」.
「你向下都可以下字條的」.
「向上跟神說可不可以下字條?」.
「我已經上了字條給你了」.
「你都不跟」.

$^{281}$「都不行」.
你儘管去懇求.
你儘管去到祂面前.
覺得很辛苦.
不過你的求問.
祂可能是將祂的心意告訴你.
不是照著你的答案.
不是照著你的要求.
不是照著你的命令.
是祂有祂的命定.
記得我們要信服.
祈禱.
多謝主.
你樂以聽人祈禱.
我們見到兩個禱告.
一個是以撒的懇求.
神就聽了.
也是按照你的應許.
以撒必須有後裔.
你就聽了.
於是就被李伯嘉懷孕.
誰知李伯嘉肚子癢得很厲害的時候.
他受不了太辛苦了.
他又求問神為何會這樣.
在這個時候神就告訴他.
你裡面有兩個在爭.
將來大的服是小的.
我命定了.
主我們見到你真是一個.
昔在今在以後永在.
是一個命定的主.
說有就有.
命立就立.
我們讚美主.
給我們的心潮向你信服你.
順天禱告奉耶穌名求.
阿們.
\newpage



\section{申命記 4:15-30}
\label{sec:031749}
\textbf{過去和現在的對話 (4) 入鄉除俗}~~~~葉向榮傳道
\newline
\newline
連結: \href{http://www.church.com.hk/acms/content.asp?site=cdc\&op=show\&type=product\&code=031749\&layout=sermon}{\texttt{031749}} ~~~~ 語音日期: 2026-02-28
\newline
\newline
\hyperref[sec:031748]{\small{< < < PREV SERMON < < <}}
~
\hyperref[sec:index_chronic]{\small{[返順時目]}}
~
\hyperref[sec:031750]{\small{> > > NEXT SERMON > > >}}
\newline
\newline
申命記 4:15-30
\newline
\begin{longtable}{cl}
\hline
\hline
章節 & 經文 (和合本修訂版)\\
\hline
4:15 & \begin{tabularx}{0.7\textwidth}{X} 「所以，你們為自己的緣故要分外謹慎；因為耶和華在何烈山，從火中對你們說話的那日，你們沒有看見任何形像。 \end{tabularx} \\ \\ \relax
4:16 & \begin{tabularx}{0.7\textwidth}{X} 惟恐你們的行為敗壞，為自己雕刻任何形狀的偶像，無論是男像或女像， \end{tabularx} \\ \\ \relax
4:17 & \begin{tabularx}{0.7\textwidth}{X} 或地上任何走獸的像，或任何飛在空中有翅膀的鳥的像， \end{tabularx} \\ \\ \relax
4:18 & \begin{tabularx}{0.7\textwidth}{X} 或地上任何爬行動物的像，或地底下任何水中魚的像。 \end{tabularx} \\ \\ \relax
4:19 & \begin{tabularx}{0.7\textwidth}{X} 又恐怕你向天舉目，看見耶和華－你的神為天下萬民所擺列的日月星辰，就是天上的萬象，就被誘惑去敬拜它們，事奉它們。 \end{tabularx} \\ \\ \relax
4:20 & \begin{tabularx}{0.7\textwidth}{X} 耶和華將你們從埃及帶領出來，脫離鐵爐，是要你們成為他產業的子民，像今日一樣。 \end{tabularx} \\ \\ \relax
4:21 & \begin{tabularx}{0.7\textwidth}{X} 「耶和華又因你們的緣故向我發怒，起誓不容我過約旦河，不讓我進入耶和華－你神所賜你為業的那美地。 \end{tabularx} \\ \\ \relax
4:22 & \begin{tabularx}{0.7\textwidth}{X} 我只好死在這地，不能過約旦河；但你們必過去得那美地。 \end{tabularx} \\ \\ \relax
4:23 & \begin{tabularx}{0.7\textwidth}{X} 你們要謹慎，免得忘記耶和華－你們的神與你們所立的約，為自己雕刻任何形狀的偶像，就是耶和華－你神所禁止的， \end{tabularx} \\ \\ \relax
4:24 & \begin{tabularx}{0.7\textwidth}{X} 因為耶和華－你的神是吞滅的火，是忌邪的神。 \end{tabularx} \\ \\ \relax
4:25 & \begin{tabularx}{0.7\textwidth}{X} 「你們在那地住久了，生子生孫，若行為敗壞，為自己雕刻任何形狀的偶像，行耶和華－你神眼中看為惡的事，惹他發怒， \end{tabularx} \\ \\ \relax
4:26 & \begin{tabularx}{0.7\textwidth}{X} 我今日呼天喚地向你們作見證，你們在過約旦河得為業的地上必迅速滅亡！你們在那地的日子必不長久，必全然滅絕。 \end{tabularx} \\ \\ \relax
4:27 & \begin{tabularx}{0.7\textwidth}{X} 耶和華必將你們分散在萬民中；在耶和華領你們到的列國中，你們剩下的人丁稀少。 \end{tabularx} \\ \\ \relax
4:28 & \begin{tabularx}{0.7\textwidth}{X} 在那裡，你們必事奉人手所造的神明，它們是木頭，是石頭，不能看，不能聽，不能吃，不能聞。 \end{tabularx} \\ \\ \relax
4:29 & \begin{tabularx}{0.7\textwidth}{X} 你們在那裡必尋求耶和華－你的神。你若盡心盡性尋求他，就必尋見。 \end{tabularx} \\ \\ \relax
4:30 & \begin{tabularx}{0.7\textwidth}{X} 日後你在患難中，當這一切的事臨到你，你必歸回耶和華－你的神，聽從他的話。 \end{tabularx} \\ \\
[1ex]
\hline
\hline
\end{longtable}
$^{1}$Hello,大家好.
其實每年來一次.
其實我是12年前就參加了東涌開東涌堂.
接下來就不是很多機會再在銅鑼灣堂那邊跟弟兄姊妹分享.
所以很多東西都已經忘記得七七八八了.
現在很多東西都忘記了.
但是每次來到都很歡喜,為什麼呢?.
因為我仍然可以用上帝的話語去服侍大家.
今天會和大家繼續看《申命記》.
頭四個題目都是在說上帝在上一代怎樣去下一代傳達他的心.
我想起我去年2025年我自己做了一件最大的project.
就是看了一本七八頁的書.
由年頭買了之後看到年尾,終於看完了.
不是中間太難看,我丟下了.
因為每一章的段落我覺得都很精彩.
是一本什麼書呢?.
這本書是說香港的.
叫做《香港村樂與鄉鎮,墜落與歷史》.
大家可能沒有興趣,你覺得是這樣的東西都看.
但是我很喜歡看這樣的作品.
說香港原來自己土生土長的地方原來有這樣的歷史出現.
那個作者叫做夏思怡,Patrick Hayes.
他本身是一個官.
1972年他就派來香港做李文輔的官員.
任內他做的事情就是去不同的新界的鄉村和那裡的居民聊天.
了解他們的需要.
然後到政府要推行任何政策的時候.
然後就給一些意見.
知道村民大致上是怎麼想的.
因為他本身就是一個歷史系出身的學者.
所以他就很能夠將這些情況表達出來.
也都很容易把這些檔案文件成為了他研究的材料.
我很佩服的.
因為他不只是工作的時候留在香港.
他退休之後都留在香港.
他的書是在香港寫成的.
他還是一口流利的語言.
如果大家在上網看那些影片.
他很多時候都會用廣東話去說.
看完之後我覺得.

$^{41}$長知識之餘.
不僅覺得自愧不如.
為什麼自己地方這麼多東西我都不懂.
這麼多東西都不知道.
近到我以前小時候住的深水Po .
原來深水Po 有發生過這樣的事.
怎麼來的我都完全不知道.
我對選擇去香港.
在香港入鄉隨俗的外國人我是特別尊敬的.
尤其是這幾十年來的話.
如果你是這幾十年才來香港的話.
其實懂得選擇的話就應該不會選擇留在香港.
市場你要找的職位.
你要找的工作等等.
應該不會選擇留在香港.
選擇內地應該好很多.
就好像這位夏思怡.
其實如果她要申請研究經費.
其實現在在大學裡面.
要申請關於香港題目的經費.
我想都不是這麼容易可以拿到.
香港生活又這麼貴.
廣東話又這麼難學.
大家知不知道廣東話很難學.
我問那些人.
原來他們覺得是一件很困難的事.
所以有時候見到YouTube.
彈很多懂得講廣東話的外國人.
我覺得又是特別的尊敬.
這麼厲害的.
願意肯落腳之餘.
還入鄉隨俗學我們那些話.
我覺得真是特別的尊敬.
所以我說入鄉隨俗其實是一件很受尊敬的事.
我想我們身邊有些外地的朋友來到香港.
落腳起爐的話.
我們都會對他們的選擇感到很佩服.
但是連結到今天我說的這段經文.
反而就是在《申命記》裡面.
他就要警告一班以色列民.

$^{81}$千萬不要入鄉隨俗.
不要跟隨之餘.
就是要除滅當地的風俗.
入鄉隨俗.
我們接著的經文.
我選了在《申命記》第四章第十五到三十節.
不過我再篩選一下.
有些相關的我們才拿出來看.
經文都好幾版.
不如我讀一版的經文.
各位弟兄姊妹你幫我讀一版一版.
我先讀.
大家跟著.
我們先來讀第一版.
第十五節開始.
因為上主在河列.
從火中對你們說話的那天.
你們並沒有看見任何形象.
所以你們要竭力保守自己的心靈.
免得你們敗壞自己.
為自己做各樣形象的偶像.
就像是男人女人的模樣.
或是地上各種走獸的模樣.
或是胸中各樣有翅膀的飛鳥的模樣.
請.
就算天上的黑暗.
不受任何人的對待貪污.
或是監獄.
但社會是上主的職責.
分派合乎他們的壓迫.
你們卻是上主所取的.
他把你們從埃及從鐵路裡領出來.
為要使你們作地產業的子民.
就像今天這樣.
你們自己謹慎.
你們要自己謹慎.
免得忘記上主你們的神.
與你們所立的約.
為自己做偶像.
就是上主你的神.

$^{121}$所吩咐所禁止的任何形象.
請.
主席.
今天我要叫天地向你們作見證.
你們在你們過約旦河去佔領的地上.
一定會迅速滅亡.
你們在那地上的日子.
必不長久.
一定會被滅絕.
上主必將你們分散在萬民之中.
你們剩餘的人.
在上主驅趕你們去的列國那裡.
必寥寥可數.
其實這個都很典型的.
在《申命記》裡面所寫的風格都是這樣.
重重複複.
講完一次就講一次.
或者將一個的訊息.
將它發很多次.
用不同的角度講一次.
都是這樣.
很簡單.
其實都不需要很多的解釋.
波西為什麼不容許一班以色列人.
在進入迦南地之後.
可以跟隨那裡的風俗.
入鄉隨俗.
沒有.
因為這裡反覆提及.
恐怕以色列人會跟隨迦南人的風俗.
學壞.
敗壞自己.
然後製造很多的.
模造很多的偶像.
作為神去敬拜.
以致上帝是法奴.
最終就會滅絕.
無法只在應許的地方長久居住.
然後被驅散.
分散去到不同的列國.

$^{161}$就再次.
本來帶你出來.
就從圍爐之間.
埃及帶你出來.
讓你得自由.
然後你因為不聽話.
然後你又再被驅散.
又再次成為奴隸.
這些都是大家反反覆覆.
將會在申命記裡面聽到的訊息.
可見這個是作為上一代.
上一輩的摩西.
新一代的以色列人一個很嚴正.
非常嚴厲的一個提醒.
語重心長的一個提醒.
我剛才說.
我入鄉隨俗.
一般來說都是一件好的事.
也是一件尊重當地人的一個表現.
也是說你是在表達一種友善.
你願意在那裡跟別人.
用對方的語言.
風俗跟他們相處的話.
其實是一個友善的表達.
也都這樣.
令到大家的相處都會融洽一點.
好像剛才我舉那些很多.
外地人來到香港去落腳那些的例子.
都是不同程度.
他們在這裡做的都是入鄉隨俗.
他們是會受歡迎的.
所以有些在裡面.
靠他們在香港的生活.
做YouTuber都會有足夠的收入.
這會為他們帶來好處.
其實是想深一層.
這是一個好的賣點.
就算說回我們昔日.
在清代的時候.
當外國的宣教士來到中國.

$^{201}$要進入中國內地傳福音的時候.
那些我們叫內地會的宣教士.
其實都是在做同樣的事.
大家可能都見過一些照片.
包括大德生.
或者是內地會的宣教士.
都是這樣的模樣.
不單只是穿著華服.
戴著薄帽.
你看見女的穿著長衣服等等.
都是一個這樣的裝束.
甚至他們去到哪個省.
學那裡的地方的方言.
不只是學國語這麼簡單.
不只是學華語這麼簡單.
是去到學那裡的方言.
甚至激烈到要留邊的男人.
很明顯他們是一個完全.
徹底地要進入中國文化當中.
意思是百姓會很尊重.
和很認同他們.
之所以說如果你要去到一個地方.
而不選擇去入鄉隨俗的話.
其實不是很發善.
甚至可能會令人覺得你很傲慢.
你覺得我的東西不好嗎.
為什麼你只是保留著你的圖片.
可能會令人覺得很傲慢.
不過這個正正和聖經所教導的不同.
聖經強調上帝的子民是要分別為聖.
分別為聖和列國的人是要不同的.
以色列文出來說.
上帝產業的子民.
是要和世俗分開的.
不是混在一起的.
所以回到我們今天的處境.
去到今時今日.
這件事又是可以怎樣去實踐出來呢.
又或者是應該那條界線是怎樣去定呢.
或者如果是仍然有界線的話.

$^{241}$那條界線是不是還是那麼明顯呢.
我們作為基督徒.
其實一直都在問這些類似的問題.
即是我們會面對那條界線的時候.
我們會受到很多的挑戰.
究竟摩西在生命的裡面.
給一班後輩的吩咐提醒.
又可以怎樣在今日.
在我們裡面可以展現出來呢.
我想和大家分享.
在經文裡面提到.
其實跟隨上帝的人.
在世上其實一直都會恆常地遭遇兩股張力.
我今天就很想講這兩股張力.
這股張力我相信是在摩西的時代.
在耶穌的時代.
甚至去到今日的時代.
我們都仍然要面對.
當我們要和世俗有這個分別為勝分出來的時候.
怎樣都是有這樣的張力走出來的.
第一個張力就是我們會有養學養.
有養學養.
剛才的經文其實講了很多次.
其實是將十界裡面的第二界.
不可以為自己做偶像.
空中地上地面下那些水那些.
其實是將這節的經文.
其實是將它擴張了.
擴張地來寫.
再舉很多的例子出來.
講明我們不可以做任何的形象出來.
當是神去敬拜.
這個很有趣.
即是和當時的文化.
當地的文化和加拿大的文化是完全的不同.
甚至和他們出來的那個埃及的文化.
亦都是不同的.
剛才講到在經文的開頭.
摩西去荷列山上領受上帝的律法的時候.
經文會告訴我們.

$^{281}$神在荷列山顯現的日子.
以色列人都是不見樣的.
都不見樣的.
只是聽到聲音.
他不會顯明有任何的形象.
讓以色列人可以描述出來.
可以繪畫出來.
甚至用來做個帽子.
做個像都沒有.
完全沒有這些.
所以有些學者去這樣形容.
他們認為耶和華的信仰.
即是我們跟隨猶太教傳下來的信仰.
是一個很著重聽覺的文化.
不是一個視覺的文化.
不知道大家是否認同.
他的理由就是.
因為神只是很想我們聽從.
我們用耳去聽上帝的說話.
用耳去聽覺去領受上帝的律法.
而不是我們要有很多視覺上的偶像.
或是有很多視覺的對象去膜拜.
重點就是.
人是否能夠聽從上帝的吩咐去做.
所以敬拜神的人.
神是要求他們要.
他們活出來的生命.
跟周圍的人不一樣.
那種分別就在這個地方.
不是在他們弄出來的像.
或是顯現出來的外表.
有什麼特別的不同.
重要的是他們是否願意去聽神的話.
我想這個想法.
我再深一層想下去的時候.
其實我覺得挺相似.
在基督教信仰裡.
其實一直貫徹著這種精神.
和大家回到一個場景.
在耶穌的時代.

$^{321}$福音書裡.
大家可能記得.
耶穌登山變像那一幕.
當有耶穌和摩沙希利亞一起顯現的時候.
彼得想做什麼.
大家記不記得.
少林踏棚大師.
踏三座棚.
每人起一座.
讓我們可以凝固了這個景象.
我相信因為他當時沒有手機.
如果彼得當年有手機.
他一定會很想拍下那個場景.
對彼得來說.
他看到這個景象對他來說是很神聖.
他只想凝固了這個景象.
所以他說踏三座棚.
大家保持著這個景象.
不過其實這個登山變像裡面.
最重要的是天父跟著說那幾句話.
天父說那幾句話.
大家在屏幕上看到.
這是我的愛子.
我喜愛他.
你們要聽從他.
也是的.
不是要著重看那幾個棚.
或者看我們幾個變得多光.
你是要聽從這位耶穌.
也是要聽從.
是這樣的貫徹.
我們是要聽從那個訊息.
不過這個就是張力所在.
因為眼前的景象.
很容易被我們更加側重.
我都去找一找.
究竟人是會喜歡用視覺多些.
還是用聽覺多些.
我不知道大家有沒有專家.
我問那些AI.

$^{361}$AI就告訴我.
人是會容易傾向更多是用視覺.
去捕捉很多周圍環境的資訊.
我不知道是不是一個真實的狀況.
但怎樣也好.
起碼對我的經驗來說.
很容易會側重我們看見的東西.
我們眼見的東西.
可能我們會相信視覺.
多於相信我們的聽覺.
對當年的以色列人來說.
他們看到迦南地那裡的人的風俗.
當他們拜那些神子.
那些偶像.
無論是來自於一些風調雨順.
又或者是一些保他們四景平安等等.
這些功能的偶像.
他們被這樣的情況吸引了.
他們就去跟隨.
按照很多不同的生物.
甚至天上的星獸等等的情況.
那個模樣去敬拜.
這個不只是一個.
對他們來說.
眼見的東西.
不只是一個屬靈對他們的影響.
也引申了很多東西.
引申了他們對上帝管治的反應.
又跟大家跳到另一個場景.
當以色列人看到.
周圍被非利人壓制的時候.
他們是城邦.
只不過每個城邦都有一個王去管治.
非利人這樣做可以.
不如我們這樣做.
他們也是這樣.
希望有養學養.
撒母耳記上第八章告訴我們.
這群以色列人的領袖.
對當時的士司撒母耳說.

$^{401}$你老了,你的兒子不行.
現在求你為我們立一個王.
治理我們.
最重要是最後一句.
好像列國一樣.
有養學養.
周圍的人是這樣做.
我們也要這樣做.
看到是這樣.
他們這樣.
那我們就跟著做.
如果是這樣的話.
可能到了我們今天.
這種傾向.
仍然會對我們帶來一些影響.
我們很容易被眼前的所見吸引.
以致我們有機會去到一個.
有養學養的狀態.
哪些是神科.
我們就要跟著那些科目去選.
哪些是筍工.
我們就要去那個方向走.
網絡的人設是這樣.
我又去學一學.
不要說這些情況.
教會也可能會這樣學.
哪些是增長的秘技.
哪些是人數增加的秘笈.
都會有我們想去學這些東西.
這個資訊量去到這麼氾濫.
這麼爆炸的年代.
其實有沒有辦法走得掉.
完全脫離有養學養的威力.
不容易的.
所以如果我們要按著我們今年的年題.
要多代同行的話.
我相信一件事就是我們.
作為上一代的.
絕對不能忽略我們的後輩.
我們的下一代.

$^{441}$正在處於一個什麼處境和環境.
正在面對什麼困難.
我想時代真的變了.
很多東西都不同了.
今天視覺影像一定是壓倒性地.
有一個這樣的位置.
我經常都會問.
為什麼這個年代我們還要.
用以往的方式去學習.
為什麼還要拿著筆寫字.
我都仍然要問為什麼要做這件事.
我們每個人都已經是.
無論在玩手機還是在玩電腦的時候.
拿著筆來學寫字有什麼特別的意思呢.
我都一直在想.
太多東西已經變了.
變到一個你不能回頭的情況.
影像真是實實在在.
是佔據著我們生活上面.
很多很多很多的空間.
我們如何在裡面能夠.
從當中我們有一個提醒.
當我看到我們面對這個張力.
去到今天是一個這麼大的.
壓倒性的位置的時候.
我就覺得我剛才說的.
神給我們的那個吩咐.
我們要去聽.
要從有樣學樣的裡面.
走出來.
我覺得反而有一些很有意思的東西.
如果我們繼續跟著.
周圍有樣學樣.
跟著不同的資訊.
有什麼彈出來我們去跟的話.
我們遲早會出事.
有一節經文.
我相信跟著大家的下一個月.
甚至再一兩個月內.
大家一定會在這裡的講台裡聽到.

$^{481}$新聞記第六章四至六節.
以色列你要停.
上主我們的神是獨一的上主.
你要盡你的全部心思.
全部的性命.
全部的力量.
愛上主你的神.
今日吩咐你的這一切話.
都要存在心裡.
繼續貫徹這個.
我們要去聽從主的那個吩咐.
大家看這段經文很熟口面.
因為耶穌都是用過這個經文去回應.
什麼是最大一家命等等.
這個系列裡面.
大家一定都會提及這段經文.
講這段經文.
所以上帝既然是希望我們.
真的要用一個聽從的角度.
去面對他的說話的時候.
其實不能夠忽略去操練.
我們的聆聽.
怎樣去聽.
十年前.
我開始了一個嗜好.
就是看小鳥.
幾乎每個星期休息的日子.
我都會去看小鳥.
在香港不同的角落.
找小鳥看.
看完之後.
拍得到的就拍.
這些都是我拍下來的東西.
但不要告訴你.
這些是十年裡面.
累積下來有這些.
不是每一次去都見到這麼多.
香港是不行的.
正正是因為.
小鳥不是在動物園裡面.

$^{521}$在籠裡面.
你一定會見到牠.
很多時候都見不到.
與其說我是去觀鳥.
其實我是去聽鳥多一點.
去到樹林裡面.
其實很多時候都見不到牠的樣子.
我其實都是聽到牠在唱歌.
聽到牠在叫.
牠可能從頭到尾.
我都見不到牠.
不過我知道牠在這裡.
聽牠的聲音.
牠的叫聲是一種特別的品種.
我就知道.
這種品種就是在這裡.
但很多時間我都是在聽.
不是在看.
原來發覺.
那個給我的反省.
原來不是我見到牠.
牠才存在.
原來我聽到牠的聲音.
牠都已經存在.
當我慢慢多去用我的聽覺.
去分辨那些小鳥的時候.
我發覺自己那種的集中力.
那種的耐性.
是真的多了.
是真的會多了.
提升了.
這個我覺得是給我一個.
學習去聆聽的其中一個技巧.
另一個聆聽.
我就說回我現在.
另一個學習.
我現在正在上一個證書課程.
叫做臨床目關.
即是去醫院裡面探病的院務.
所以會經常去醫院.

$^{561}$去病房裡面.
進去探病人.
我想說這個其實是一個非常.
非常好的一個聆聽的學習.
不要想著.
我去講很多東西給牠聽.
我靠近牠的床邊.
我就可以跟牠講.
三腹四肋全福音.
不行的.
基本上是.
很多時候我去到牠的床邊.
都是聽.
最大的發現原來是.
我發覺原來是.
仔細去聽一個人去分享.
牠會發覺生命真的可以很沉重的.
我是每次都.
我現在越來越發覺.
我每次不能夠聽超過三個人分享.
因為每個都是這麼沉重.
生命你會發覺.
當然如果牠要騙我.
當然是另一回事.
但我不覺得一個病了的人.
為什麼要騙我這麼多東西.
你發覺每個人的狀況.
面對的處境.
那種痛苦.
那種難受.
是可以很多方面發生.
去到那種情況.
我真的明白.
為什麼這麼多人.
不想聽別人分享.
我們真正耐心去聽.
其實生命真的很不容易.
但我同時也發現.
另一件事就是.
如果我們真的願意.

$^{601}$走到牠們旁邊去聽的話.
其實是一個很大的鼓勵和接納.
前兩個星期.
我去探望一個.
剛剛患了癌症.
開始做化療的媽媽.
走到牠的床邊.
其實牠一開始都不是很想說很多東西.
我問牠很多情況.
牠都說OK.
沒什麼的.
就這樣吧.
你聽的意思就是.
牠不是很想說.
不想說很多東西給你聽.
或者是覺得.
你不如找別人.
不要找我.
我也想著牠是這樣的.
我也會聽的.
牠也知道不是很想說.
但我臨走之前.
我跟牠說.
我是信耶穌的.
我會為你祈禱.
我不知道為什麼.
說完這一句之後.
牠的心好像突然有很大的觸動.
牠哭得很厲害.
突然哭得很厲害.
接著就說了很多很多東西出來.
原來原因就是因為.
牠也是信耶穌的.
不過牠沒有什麼.
信了耶穌之後.
沒有去上教會.
沒有跟上帝同行的生活.
但這一刻牠那種憂慮.
那些壓力.
那些很多很害怕的東西.

$^{641}$就一下子全部倒出來.
我發現原來一個人.
被人聽見.
被人去聆聽的時候.
可以帶來很大的幫助.
我相信這個媽媽.
也有一個很大的釋放在裡面.
可能沒什麼機會.
跟自己的家人不敢說.
但牠對著我.
當我說我也是基督徒的時候.
牠願意這樣去分享出來.
釋放出來.
聆聽.
接著第二個.
我們在世上所要面對.
和世界有分別的張力.
就是當我們在這裡生活得久的時候.
當我們住得久的時候.
人會求長壽.
我們家庭會求安居樂業.
政府要求長治久安.
很多東西我們都很想last.
很想好好延續下去.
多數都是我們共同的願望.
但大家有沒有想過.
要維持下去.
我們可能要有些辦法.
而這些方法.
未必所有都合上帝的心意.
摩西警告這些以色列人.
當他們在迦南地落腳.
久了之後.
對你有無有益.
就會出事.
因為他們要想辦法.
令自己繼續在應許的地方.
繼續生活下去.
這個地方本來是上帝給他們的.
是上帝應許的.

$^{681}$是上帝說這是流乃與物之地.
因為你生了兒子.
兒子又生了兒子.
在這裡住久了就會敗壞自己.
他們當然不是第一班人出事.
他們的祖先都是這樣.
創世記錄告訴我們.
他們一班人為了結集自己的實力.
他們要保持自己.
要顯示自己的能力.
於是建造一個巴別塔.
這個故事大家都聽過.
這是第一個例子.
為了令自己住得長久.
他們做一些事.
他們做一個巴別塔.
很容易.
我們為了讓事情繼續發生下去.
繼續延續下去.
我們會很想找一個方程式.
找一個秘笈去滿足.
我們內裡求生的慾望.
或者求長久的慾望.
就像彼得一樣.
這裡真好.
在登山變象的時候說.
這裡真是好.
我們都很想有這個慾望.
於是連彼得都會說搭座棚.
我們長住在這裡.
結果最終我們很容易.
因為為了要繼續延續下去.
繼續長久生存下去.
偏離了我們最初想要做到的那件事.
偏離了我們要想達到做的目的.
我們的目標.
人生裡面需要追求的目標和使命.
可能就是因為我們要維持下去.
為了生存.
我們很容易就失卻了.

$^{721}$在教會裡面都會的.
教會的一個狀況.
很容易因為我們要追求一種慣性.
追求我們的傳統.
每個聚會我們都知道要發生什麼.
大家今天回來.
你知道會有敬拜.
有奉獻.
有報告.
有道聽.
每個禮拜大家回來.
應該都是這樣的東西會發生.
這是一個可預期可計算的一個地方.
但往往就很容易.
當我們這樣去日復一日.
月復月.
每個禮拜都這樣的時候.
會讓我們很容易變成一個僵化了的群體.
而這個和教會本身的原意是有點不同的.
教會我們說叫Ecclesia.
本身是一個被呼召出來的聚集.
一群人是被上帝呼召出來去聚集.
去跟隨他.
是時刻在上帝的面前聽候於他.
以致可以成為這個世上的見證.
本身是應該.
不只是為了他們繼續的長久存在而存在的.
他們是有個使命.
帶著一個目標.
在上帝的吩咐下.
去到世上作見證.
所以我們是不可以.
也不應該將.
本身教會應該是一個很有活力.
很流動的一個movement.
變成一個動都不動.
只是為了去記述.
我們做過些什麼.
做過些什麼.
很多豐功偉績的monument.

$^{761}$前幾天我聽到一個.
一位牧師的分享.
我覺得很厲害.
我覺得是一個很厲害的一個情況.
一位韓國的牧師.
在韓國的釜山的YWAM那邊服侍.
他服侍了很多年.
服侍了二十多年.
十幾年前.
常常給他一個感動.
搞一個camp.
讓南韓的年輕人.
去接待那些.
北韓的脫北者.
那些後代.
可能他們離開了北韓.
去了中國之後.
再輾轉去到南韓生活.
做一些這樣的聚會.
讓兩韓之間的年輕人.
可以有一些這樣的連結.
這個營會很好.
讓大家彼此去認識.
也讓福音可以傳給他們.
一個這樣的營會.
反應越來越好.
一路做一路做.
但辦了十年之後.
同工一次的團隊裡面.
他們一起祈禱的時候.
突然間他們聽到上帝說.
要停了.
不要再搞這個營會了.
我想用智慧人的想法.
這個不是吧.
不是做完沒有人可以做.
是繼續有得做.
也不是已經做得太舊了.
也沒有人理會.
反應這麼好.

$^{801}$一路有這麼大的需要的時候.
為什麼要停呢.
但最終他們都選擇.
順服上帝的吩咐.
他們不知道為什麼.
總之他們真的停了營會的服事.
後來過了一兩年之後.
他們才發現.
原來上帝很想他們的服事.
有另一步去走.
不只是留在韓國.
他們要開始將多代同行的訊息.
向周圍不同的國家,地方,文化宣講.
後來他們多走幾年之後.
才發現上帝有這樣的鋪排和心意.
剛才說了.
在多代同行裡面.
我們一定不能忽略.
下一代他們現在面對的情況和心聲.
這是我猜測的.
如果我錯的話.
大家提醒我.
我猜測.
現在沒有很多的下一代.
會很有興趣或者很有熱誠.
進入一些傳統的體制裡面.
我不知道對不對.
如果我錯.
你糾正我.
不知道有沒有弟兄姊妹已經很立心.
有心智要做教會執事的呢?.
我不是呼召.
有啊.
如果你有的話.
那當然好.
但是我想.
這些體制.
很多的體制.
是不是這麼吸引我們的下一代.
願意參與在裡面呢?.

$^{841}$我想如果我又錯的話.
大家又糾正我.
我又覺得.
是不是一些套餐式的模式的東西.
大家都不會覺得很有興趣呢?.
一大堆頭.
一些運動.
大家一起要做什麼.
很多的模式.
很多的一套套餐.
大家一起去做的東西.
不知道大家又喜不喜歡.
如果我錯.
大家又糾正我.
我覺得以我接觸的下一代來說.
又好像不是很work.
我們的下一代未必是很喜歡跟隨很多的定規.
很多的一定要這樣做.
一定要這樣做才行.
很多的要很多的變化.
或者是很著重你的生命的本質.
特質的展現多一點.
最近我想大家都有看過這個女生的新聞.
剛剛拿了東奧女子滑冰金牌的劉美妍.
Alisa Liu.
她的故事如果大家都聽過.
簡單來說就是她在爸爸的栽培.
五歲就開始用數以萬計的美金.
栽培她去做一個滑冰的選手.
是一個很好的進步.
十三歲就已經參加世錦賽.
已經可以成為一個世界上頂級的選手.
但到了十六歲的時候.
她突然間就說我不玩了.
我不想再繼續滑冰.
原因就是我不想失去自己的生活.
我很想做一些一個年輕人應該做的事.
我不想活在大家的睹目和控制的底下.
於是她就去了旅行.
去了讀書.

$^{881}$去了讀大學.
總之就是做自己喜歡做的事.
有幾年的時間去做這些事情.
直到她有一次突然間.
感受到自己很想去滑冰旅途的那種熱誠.
回來的時候.
她才重投滑冰的行列.
她回去繼續訓練.
剛剛拿了東奧的金牌.
不想.
正正就是不想在大家的睹目.
在上一代的睹目的底下繼續去走.
所以我們都要學習的.
我們若果不想走在一個住久了.
我們會變得僵化的引誘的底下.
我們其實真的要學習更多的.
去欣賞上帝自己在裡面不同層面的工作.
無論大家其實可以多多的走出去.
我覺得大自然其實是一個很好的導師.
看看上帝的創造.
其實很重要的.
當我們認定這是天賦世界的時候.
這位是獨行其事的主.
當然我們說讀百科可以知道很多東西.
或者我讀地球科學.
我會知道很多東西的原理是怎樣來的.
規律是怎樣去運作.
但總有很多東西我們仍然不知道是很奇妙的創造.
它有它的秩序有機制.
不過仍然有很多的分別.
生命和自然都是很奇妙的莫測.
絕對不是很機械的運作當中.
有機會大家去看約伯記的第三十八到四十一章的時候.
上帝就正式用大自然的一大堆的例子去大約伯.
這幾章聖經是很好看的.
以致大到約伯都嘔了嘴.
沒有東西可以回應.
就在第四十二章.
不行啊.
這麼多講這些我完全搞不定.

$^{921}$大自然的東西我真的懂些什麼呢.
於是他說.
我真真正正見到原來你的作為是這麼奇妙.
因為你去掌管這一切.
只能夠謙卑下來.
我們需要除了聆聽神在自然界的工作.
也要聆聽神的聖靈在我們生命裡的引導.
上帝既然這麼愛我們的時候.
他一定會很想去引導我們.
與我們同行.
所以我們要學習去聆聽上帝的聲音.
以致我們能夠分辨.
哪條路是我們應該走.
不會走向慣性.
或者有樣學樣的道路上.
所以最後給大家的建議.
可能已經有些大家在做的.
我鼓勵大家繼續.
大家可能比較喜歡看聖經.
不過我鼓勵你用一個星期的時間去聽聖經.
聽聖經可能給你一個另外一個嶄新的體驗.
聽的.
教會最近也在推一元.
那個就是大家去聆聽.
可以去聆聽聖經.
另外除了聽上帝的話語.
大家也可以聽一位朋友的分享.
可能大家都會做.
聽一下他裡面的感受.
認認真真地聽一下一個人的感受.
這是一個很好的練習.
當然也可以每一天用五分鐘去聆聽一下.
聖靈怎樣向我們說話.
安靜下來.
無論是找一個你最合適的時間.
不用太久,五分鐘.
上帝足以讓上帝向我們.
聖靈向我們說話.
祝福我們每一位.
上帝叫我們去聆聽的時候.

$^{961}$總有很多方法可以讓我們聽得到.
請我們一同低頭祈禱.
主我們仰望你.
主我們多謝你.
多謝你的恩典.
多謝你的創造.
多謝你.
藉著你的律法.
叫我們可以去跟從.
去聽從.
主即使這個世界已經變到.
很多事情都回不了頭的時候.
主你仍然留下這些吩咐.
這些提醒.
叫我們不會有樣學樣.
叫我們不會走向僵化的道路.
幫助我們.
主啊,所以.
保羅這樣寫我們要心意更新而變化的時候.
你就這樣去幫助我們.
聆聽你的聲音.
聆聽你在我們裡面.
最微細的那個位置.
向我們發的聲音.
藉著你的話語的每一天.
都成為我們的糧食.
成為我們的呼吸.
成為我們的引導.
來幫助我們.
與我們同在.
奉耶穌基督的名.
阿們.
\newpage



\section{列王記下 2:1-14}
\label{sec:031750}
\textbf{多代同行•與神同行系列:主,我要你更多:以利沙的加倍之靈}~~~~唐郭麗婷牧師
\newline
\newline
連結: \href{http://www.church.com.hk/acms/content.asp?site=cdc\&op=show\&type=product\&code=031750\&layout=sermon}{\texttt{031750}} ~~~~ 語音日期: 2026-03-01
\newline
\newline
\hyperref[sec:031749]{\small{< < < PREV SERMON < < <}}
~
\hyperref[sec:index_chronic]{\small{[返順時目]}}
~
\hyperref[sec:031751]{\small{> > > NEXT SERMON > > >}}
\newline
\newline
列王記下 2:1-14
\newline
\begin{longtable}{cl}
\hline
\hline
章節 & 經文 (和合本修訂版)\\
\hline
2:1 & \begin{tabularx}{0.7\textwidth}{X} 耶和華要用旋風接以利亞升天的時候，以利亞與以利沙從吉甲往前行。 \end{tabularx} \\ \\ \relax
2:2 & \begin{tabularx}{0.7\textwidth}{X} 以利亞對以利沙說：「耶和華差遣我往伯特利去，你可以留在這裡。」以利沙說：「我指著永生的耶和華，又指著你的性命起誓，我必不離開你。」於是二人下到伯特利。 \end{tabularx} \\ \\ \relax
2:3 & \begin{tabularx}{0.7\textwidth}{X} 在伯特利的先知的門徒出來，到以利沙那裡，對他說：「耶和華今日要接你的師父離開你，你知不知道？」他說：「我知道，你們不要作聲。」 \end{tabularx} \\ \\ \relax
2:4 & \begin{tabularx}{0.7\textwidth}{X} 以利亞對以利沙說：「耶和華差遣我往耶利哥去，你可以留在這裡。」以利沙說：「我指著永生的耶和華，又指著你的性命起誓，我必不離開你。」於是二人到了耶利哥。 \end{tabularx} \\ \\ \relax
2:5 & \begin{tabularx}{0.7\textwidth}{X} 在耶利哥的先知的門徒來靠近以利沙，對他說：「耶和華今日要接你的師父離開你，你知不知道？」他說：「我知道，你們不要作聲。」 \end{tabularx} \\ \\ \relax
2:6 & \begin{tabularx}{0.7\textwidth}{X} 以利亞對以利沙說：「耶和華差遣我往約旦河去，你可以留在這裡。」以利沙說：「我指著永生的耶和華，又指著你的性命起誓，我必不離開你。」於是二人一同往前行。 \end{tabularx} \\ \\ \relax
2:7 & \begin{tabularx}{0.7\textwidth}{X} 有五十個先知的門徒同去，遠遠地站在他們對面；他們二人在約旦河邊站住。 \end{tabularx} \\ \\ \relax
2:8 & \begin{tabularx}{0.7\textwidth}{X} 以利亞捲起自己的外衣，用來打水，水就左右分開，二人走乾地過去。 \end{tabularx} \\ \\ \relax
2:9 & \begin{tabularx}{0.7\textwidth}{X} 過去之後，以利亞對以利沙說：「我未被接去離開你以前，你要我為你做甚麼，只管求。」以利沙說：「願感動你的靈雙倍感動我。」 \end{tabularx} \\ \\ \relax
2:10 & \begin{tabularx}{0.7\textwidth}{X} 以利亞說：「你求的是一件難事。我被接去離開你的時候，你若看見我，就必得著；若不然，就得不著了。」 \end{tabularx} \\ \\ \relax
2:11 & \begin{tabularx}{0.7\textwidth}{X} 他們邊走邊說話的時候，看哪，有火馬和火焰車出現，把二人隔開，以利亞就乘旋風升天去了。 \end{tabularx} \\ \\ \relax
2:12 & \begin{tabularx}{0.7\textwidth}{X} 以利沙看見，就呼叫說：「我父啊！我父啊！以色列的戰車騎兵啊！」以利沙不再看見他的時候，就把自己的衣服撕為兩片。 \end{tabularx} \\ \\ \relax
2:13 & \begin{tabularx}{0.7\textwidth}{X} 他拾起以利亞身上掉下來的外衣，回去站在約旦河邊。 \end{tabularx} \\ \\ \relax
2:14 & \begin{tabularx}{0.7\textwidth}{X} 他用以利亞身上掉下來的外衣打水，說：「耶和華－以利亞的神在哪裡呢？」打水之後，水也左右分開，以利沙就過去了。 \end{tabularx} \\ \\
[1ex]
\hline
\hline
\end{longtable}
$^{1}$今天我們要來向以利沙學習.
第一樣我想和大家分享.
我們要學習以利沙那份堅定的跟隨.
他堅定地跟隨他的師父以利亞.
更加堅定地跟隨神.
所以這個跟隨的焦點.
不只是在人以利亞身上.
而是更加在神的身上.
我們如何能夠看到這件事呢?.
當我們回看《鑷王紀相》第19章時.
其實之前鄭牧師也有和我們分享過.
當以利亞來遇見以利沙.
以利沙將他象徵著一個先知職份的外衣.
披在以利沙身上的時候.
以利沙就知道以利亞的意思.
就是他領受到神要將先知的職份.
從以利亞身上承傳下來.
來接受這個任務.
所以他知道神對他有一個呼召的時候.
他沒有和神推遲任務.
反而他很甘心樂意地.
安排自己離開安逸的生活.
來跟隨先知服侍神.
當我們看到經文記載他領受了這個呼召之後.
之後就沒有特別說他做了什麼.
直到今天我們要看的經文《鑷王紀相》二章.
一章十四節.
其實就去到一個時候.
就是當他的師父以利亞要被接昇天.
即是離開他的時候.
所以我們見到以利沙是一個溫和謙卑的人.
他接受了這個呼召之後.
他沒有直接說.
哇 以利亞你走開點 現在到我來了.
他沒有急於做第二個以利亞.
他只是一路默默地做一個助手.
跟著 看著以利亞如何去服侍神.
做一個先知 和他一起.
但其實去到這一刻.
他跟著以利亞 跟到這一刻.

$^{41}$以利亞知道自己.
神和他說我要接你昇天.
其實是什麼意思.
真的是交棒和接棒的時候.
真的來到了.
其實在這個時候.
以利亞就做了一件很特別的事.
以利沙一路跟著他.
他和他由吉甲這個地方 大家認識吉甲這個地方.
是心靈角的地方.
然後他們就經過了伯特利.
去耶利哥 然後去約旦河.
其實特別的事就是.
當他們一起走到每一個站.
以利亞都會和以利沙說.
他說 你就可以停留在這裡了.
你不用再跟我走下去下一站了.
你每次走一站就這樣說一次.
其實那時候除了以利沙之外.
還有其他的先知門徒是隨行的.
其實那些人都收到消息.
他就和以利沙說.
喂 你師父就要走了.
聽聞要接昇天了.
你是不是還跟著.
但以利沙每一次的回答都是很堅定的.
他和以利亞說.
我必不離開你的.
我不會揹負你.
我一定會跟著你的.
其實以利亞這個勸勉 叫他留下來.
其實都挺不尋常的.
其實明明是接班人.
跟著他去接班 怎麼會呢.
其實他每次都叫他.
你不用再走了 停留在這裡.
其實什麼原因呢.
其實都是一個考驗.
你跟了我這麼久.
你看到我所做的事.

$^{81}$其實你會不會再跟下去.
你會不會夠膽再做下去.
其實我們見到以利沙很堅定地和他說.
師父 我是不會離開你的.
不會被人 這個消息.
我不會被環境的轉變.
是你都讓我去離開.
好了 如果你是以利亞的徒弟.
這麼堅定地跟隨.
你會怎樣呢.
都挺開心的.
你不會說.
不聽我的話 明明叫你留下來.
我是你師父 我說什麼你就做什麼.
不是的 這個徒弟.
他不聽我的話.
就一直跟著他.
但是呢 我想以利亞很開心.
第一 因為有一個徒弟.
這麼忠心地跟著你.
其實都表示他很愛你.
他很願意一輩子跟著你.
他對你很忠心.
你作為一個師父.
你都會很開心.
是嗎 這樣的徒弟.
但是我相信.
以利亞更加開心.
甚至覺得老懷安慰的就是.
其實你知道一個偉大的先知.
其實他必定是一個很好的師父.
一個好的師父.
你最終要訓練人做什麼.
不只是跟隨他.
而是跟隨神.
所以當他見到.
其實當這個以利亞跟了他這麼多年.
知道這個任務.
承接這個先知任務是多麼艱鉅.
是怎樣艱鉅呢.

$^{121}$因為你看到.
那時候他在那個年代.
夜洗別是怎樣去推崇.
巴力 崇拜巴力.
然後呢 偶像崇拜.
然後又追殺那些.
先知追殺那些敬畏神的人.
然後又很多假先知.
然後其實那時候.
不過以色列十個支派.
其實那些人基本上全部都拜了偶像.
所以你說.
要像師父那樣來對抗這些.
皇權 黑暗的勢力.
還要去行.
師父都行很多神蹟騎士.
其實你想想 我可能都退縮了.
不要搞了 我好像都不怎麼做到.
但是一個平凡的以利亞.
他就是這樣堅持來到去.
走下去.
什麼原因讓他走下去.
就是因為他很清楚.
不是他自己的能力.
而是上帝對他的呼召.
上帝當日他知道他跟的.
當他那天去接受.
擺下自己的生活.
地上的一切來到去跟隨以利亞的時候.
他知道他不只是接受以利亞的邀請.
叫他去做一個先知.
而是他知道這是上帝的呼召.
叫他來承接他神的國的事情.
所以以利亞之所以能夠那麼堅定.
師父你說什麼我都是跟的.
除了跟你之外.
更加是因為我不能夠揹負.
上帝在我生命裡的託付.
上帝將這件事.
他的國擺在了他的心裡面.

$^{161}$所以他怎麼都沒有退縮.
無論他遇到什麼環境.
現在直情知道師父要走.
他都說我要繼續走下去.
所以我會覺得.
以利亞真的很安樂.
你找一個接棒的人.
他真的能夠成全這個使命.
他不只是跟著你.
環境多惡劣.
師父都走了.
但他仍然堅持能夠走在神給他的託付裡面.
他看到的不是隻是以利亞要走.
他看到的不是師父不在.
他看到的是神給他的使命.
所以他堅持堅持要來到繼續走.
就像我們今天受師的樣子.
有什麼可以堅持.
有什麼可以令我們不停地繼續服侍.
不是因為人.
是因為神自己.
他跟隨的是神自己.
是跟隨以利亞跟隨的神.
他跟隨的是差遣以利亞的神.
所以他能夠一路走下去.
其實他堅定不移地跟隨.
焦點要在神那裡.
不在人那裡.
其實對我們都有很大的提醒.
特別是我自己.
當我一路去預備的時候.
我更加對自己有個提醒.
如果作為一個牧者.
我帶領弟兄姊妹都很小心.
其實要有以利亞的心懷.
其實要將你牧養的人.
是帶去緊緊地跟隨誰.
是跟隨神.
不是跟隨我.
有時候我們都會說.

$^{201}$我叫你做事 你來聽我的話.
其實想真一點.
弟兄姊妹要跟隨的是神.
不是我.
所以我們經常要提醒自己.
我們要帶弟兄姊妹跟隨的.
亦都是神給教會的異象.
不是我自己的異象.
不是我自己的方法.
我要帶弟兄姊妹跟隨神.
以至到教會才能夠合一.
這就是神對我們今天的提醒.
另外 當我們看經文時.
裡面有說到.
當時有五十個先知門徒一起隨行.
不是隻有以利沙一個.
有幾十個人一起行.
但是行到最後.
其實只有以利沙一個.
是跟以利亞行到最終.
然後其他人是什麼呢.
其他人 聖經形容他們.
好像彭博特所說.
我們說這是世上.
信主的人都很多.
甚至很多人都自稱我們是主的門徒.
我們是先知.
但是走到最後.
真的能夠很清楚自己的抱負.
神給我們心裡的託付.
不管人的轉變 環境的轉變.
能夠堅持走到最後.
完成使命的.
其實是少數.
弟兄姊妹 你覺得呢.
你自己現在的狀況.
你會是站在遠遠地看著.
即是以利沙 以利亞那個.
還是你都是願意做以利沙呢.
或者你覺得我就是以利沙呢.

$^{241}$其實呢 我自己都有很多反省.
我都是在想.
究竟我不可以好像以利沙那樣.
其實呢 為什麼呢.
因為我想大家都聽過很多牧者.
和大家的分享.
因為做一個傳道人.
其實都不是很容易.
當我進來做之後.
現在到現在17年了.
我就發覺原來這份所謂的事奉.
功呢 就沒有我當初想象的.
我來服侍神啦.
我有這個抱負啦.
我就來擺上我自己的生命.
好像講得很偉大那樣.
其實當真的做下去的時候.
你發覺比外面的那份工作難做一百倍之餘.
那種人事呀 那種複雜呢.
其實呢 主啊.
即是天天都很難的.
真是不容易.
因為我以前那份工作都不容易.
其實這份更艱難.
其實有很多人呢.
即是如果你在同福岸看下去.
都有很多頂子妹.
在同福岸一段日子.
你都會可能知道.
其實教會都經歷了很多事情.
是吧.
經歷了很多不同的動盪啦.
人事的變動啦.
真是很多很多的事情.
其實每一次.
都不容易過的.
但是呢 完了之後.
我的頂子妹就說.
哇 其實你看到那些人.
那麼多雞飛狗走啊.

$^{281}$你有沒有想過你都離場的.
你認為我有沒有想過呢.
有啊 我真是大聲講得哭.
是 是有的.
沒有就真是騙你了.
為什麼我有.
因為我是一個普通的人來的.
是吧.
我有時候跟自己說.
啊 當我跟阿桑拉羅在那裡聊天.
我們是不是很傻呢.
即是派了兩份工.
即是進來同福堂.
還將雞蛋放在一個籃子裡面.
啊 真是是不是傻傻的呢.
在經歷這個時候.
然後我們有一個四字.
我們自己的真言.
就是我們經常說.
當我們遇到挫折的時候.
或者困難.
真是很煩啊.
然後我們就會說.
啊 主啊不如禰讓我們一起遠走高飛吧.
這個就是我們經常說.
我們一起進入大聖之後.
有沒有想遠走高飛啊.
其實真是經常想遠走高飛.
但是呢 當然.
每次我們這麼想.
其實這個不是一個輕易的決定.
我們都會想一想 問神.
哎呀 主啊.
其實是不是應該離開還是怎麼樣呢.
其實每次呢.
起碼到這一刻為止.
神每次都提醒我們.
祂說你要記得.
當日你是用什麼來決定.
放下你自己的工作.

$^{321}$全時間來到這個教會去服侍.
你要記得.
你當日是怎樣做這個決定的.
那是不是為了人呢.
是不是人家叫你進來做.
你為了回應那個人呢.
還是你自己.
啊 我想.
我覺得我很有恩賜啊.
我要進來發揮啊.
是不是這樣呢.
其實到最終我們回想.
都不是.
今天我們想 為什麼我們會在這裡.
是因為神放在我們裡面.
福音的託付.
因為我們是受過神恩惠的人.
我們知道呢.
那份救恩是多麼的好.
改變了我們的生命.
拯救了我們的生命.
所以我們很想.
主啊.
這個就是唯一人能夠得救.
能夠面對這個世界.
繼續走下去.
唯一的拯救.
唯一的道路.
所以我們很想.
將我們得到的.
告訴別人.
去分享.
這個全福音的使命.
另外就是.
神給我們.
對這個屬靈的家.
一份的愛.
那份要我們託付.
在這裡去服侍.
這裡的弟兄姊妹.

$^{361}$服侍這個神的家.
所以其實我們說.
哎 怎麼辦啊.
那好吧.
神禰既然有這個託付.
禰要幫助我們.
我們不是賴死不走.
而是.
當今天禰仍然將這個託付.
擺在我們心裡的時候.
求主幫助我們.
所以雖然經歷了.
很多不同的事情.
但是.
我們都一路堅持.
我們學的東西.
我們想學的東西.
不是說我們有多了不起.
而是我們想學習.
如何在.
走神的道路里面.
我們能夠按著.
一路一路.
是專注在神給我們生命.
託付那個負擔在我們心裡.
對神的國.
神的家的負擔.
然後我們學習.
不讓環境.
不讓周遭發生什麼事.
不讓人影響.
不讓消息影響我們.
而是專注專注專注.
在神給我們的負擔.
承擔的東西里面.
這個就是我們一路學習的.
我希望其實弟兄姊妹.
我們來到一個.
神的家的時候.
我們真的要學習.

$^{401}$因為人始終很容易.
將人去偶像化.
其實外面有很多不同的人.
KOL.
權威.
其實我們都會.
甚至屬靈的領袖.
我們都會有偶像化.
所以其實.
如果我們是看人.
跟人的時候.
這個人.
我不跟他不行.
沒了他好像.
連神都不在了.
其實真的要小心.
我們都很盼望.
我們不會是這樣.
也都是當.
有一天.
當我們自己來到.
退落裡的時候.
我們真的很盼望.
其實到了現在的年紀.
服事了這麼多年.
我們都會想.
其實這個是時候叫神.
祈禱神.
興起新一代的時候.
興起新一代來到接班的時候.
我們的祈禱.
是不是說.
神啊.
我是時候離開了.
你快點找個人.
頂著我的位置.
那我就可以走了.
沒人來頂著我的位置.
請不到人.
又走不了.

$^{441}$是不是挺有意思的.
其實.
是不是隻是這樣呢.
其實我們真的不是這樣.
我們真的盼望神興起.
我們的新一代下一代.
是能夠專注在神的身上.
祂很清楚.
我們經常問.
現在的童工.
你.
是不是回應神的呼召.
今天你留在這裡.
你今天在這裡做.
是不是回應神的呼召.
是的時候就可以了.
因為這個呼召.
可以讓你堅持下去.
在任何的環境下.
任何的轉變.
因為人是會流動的.
是不是.
一代也會過去.
新的一代又來.
所以人只能夠跟著神.
並且是神.
在我們心裡面的那份呼召.
讓我們能夠堅持.
走下去.
在任何的環境.
所以鼎姐妹.
我們一起來學習.
一起來實踐.
當我們見到.
爾利沙通過了師父的考驗.
來到約旦河.
爾利亞用他的外衣.
打水.
神蹟.
其實是地上的最後一個神蹟.

$^{481}$河水分開了.
一起過河.
走在乾地上.
來到這個關鍵的時刻.
我想跟大家一起讀.
這段關鍵的經文.
就是二章九至十四節.
爾利亞昇天之前.
他所跟爾利沙的一段對話.
我們一起讀.
一二三.
我願意向爾利亞說.
對爾利沙說.
我未曾被接去離開你.
你要我為你做什麼.
只管求我.
爾利沙說.
願感動你的靈.
加倍地感動我.
爾利亞說.
你所求的難得.
雖然如此.
我被接去離開你的時候.
你若看見我就必得著.
不然必得不著了.
他們正走著說話.
忽有火車火馬將二人隔開.
爾利亞就乘旋風昇天去了.
爾利沙看見就呼叫說.
我父啊.
以色列的戰車馬兵啊.
以後不再見他了.
於是爾利亞把自己的衣服撕為兩片.
他身上掉下來的外衣打水說.
掉下來的外衣.
會去站在約旦河邊.
他身上掉下來的外衣打水說.
耶和華.
爾利亞的神在哪裡呢.
打水之後.

$^{521}$水也左右分開.
爾利沙就過來了.
好了.
他要離開了.
爾利亞問.
爾利沙.
我要走.
我可以為你做些什麼.
彷彿就好像神透過.
爾利亞和爾利沙問.
你今天跟到我們這一刻.
你知道你接下來要接的東西是什麼.
有多艱難.
你需要些什麼.
來接這個任務.
你想我為你求些什麼.
其實.
如果是你.
你會求些什麼.
我在想.
如果我要接一個牧區.
或者我要接一個事工.
接一個小組的時候.
我可能會求.
主啊求你幫助我.
現在人數不夠.
下年你就將我小組的人數雙倍.
將教會的人數雙倍.
甚至三倍.
得救的人數加很多很多.
或者.
我覺得我自己不行.
主啊不如你.
將之前上手的能力.
雙倍給我.
讓我可以做得到.
我們可能會這樣求.
但是.
爾利沙求的是什麼.
感動你的靈.

$^{561}$加倍的感動我.
我們下一章.
下一章再下去.
就是說感動你的靈.
加倍的感動我.
他求一個.
他是求加倍的靈.
什麼是加倍的靈.
如果我們看回.
你還沒記得的21章.
講到長子.
以前長子是會承繼.
去爸爸雙份的產業.
長子.
就是一個加倍的承受的產業.
所以這裡加倍的靈.
可以有一個解釋就是.
爾利沙和爾利亞說.
我就是你屬靈的長子.
你就是我想繼承你的產業.
將你的力量.
將你的事業.
事功繼續傳下去.
但是.
當我們看《釋經書》的時候.
其實有一個更深的解釋.
就是說.
其實.
加倍的靈.
靈是什麼.
就是主自己.
主就是靈.
神就說.
其實我要的不是.
什麼人數.
成就能力.
而是說的是我是主.
當如果我要做你的事情的時候.
你要我做你的事情.
我要的是你自己的主.

$^{601}$所以他說求主.
我要你更多主.
我要你加倍地充滿我.
加倍地掌管我.
以至到.
你要加倍地與我同行.
以至到我能夠.
真的.
去做你自己的事情.
因為不是我能夠做.
是主.
你自己.
來去做.
所以這個就是.
以利沙知道自己很平凡.
你叫敵黨這麼多東西.
好像我們今天在一個亂世裡.
叫我們去做神的事.
很多人都說不要搞了.
做不到.
但是.
其實當日.
以利沙也是這樣.
面對一個很艱難的環境.
但是.
他就是這樣去做.
所以第二樣.
我想大家去學習的.
就是我們要學習.
我們不是說自己做到多少.
而是我們能夠學習.
以利沙這樣向神.
謙卑.
渴慕地來祈求.
然後他向神.
大大開口.
他求什麼.
願感動以利亞的靈.
加倍感動我.
主我大大向禰開口.

$^{641}$神啊.
禰自己來到.
我要的是禰.
所以頂子妹.
今天我們服事的時候.
我沒有跟神說.
我要的是禰.
我不是禰給我什麼能力.
什麼技巧.
給我一些成就.
我要禰 主自己.
所以呢.
這個就是這樣.
所以頂子妹.
其實現在.
有沒有一些事情.
你覺得很艱難的呢.
其實有沒有東西.
其實是神.
透過人或者去催逼你.
或者怎麼樣.
神其實放在你裡面.
但是你.
你一直都不敢走上去.
你覺得自己不夠能力.
你覺得自己做不到.
很多限制.
但是你要記得.
神只要你一個願意的心.
但是做事的是祂自己.
做事的是祂自己.
祂掌管.
那你覺得.
頂子妹你願不願意.
好像以利沙那樣求.
願意.
你覺得神會不會應允你.
你覺得神會不會應允我們.
我們一起來讀詩篇81篇第十節.
一起讀.

$^{681}$神對以色列說.
你們要向我張口.
我就充滿你.
神說你只要向我張口.
不然你有多少能力.
你就要知道自己有需要我.
所以你向我張口的時候.
我就大大的充滿你.
這個是神的應許.
大家記住.
所以當我們向神開口.
大大開口.
神你的令加倍來到.
感動我充滿愛.
意識到我能夠做你的事情.
和你一起同工的時候.
這個神就說.
我就應允你.
所以當我們去學習的時候.
與其我們退縮很多事情.
不如我們學習向神大大的張口.
其實這個以利亞和他的回應就說.
你求的難得.
你這樣說.
你要這件事.
要神自己.
不是指難不得.
只是說.
其實如果你是求的是神自己的話.
其實只有神可以決定給不給你.
因為這個屬靈的力量.
只有神可以給你.
不是我做你的師傅可以給你.
我可能只是給你一些經驗.
可能給你一些知識.
但是真的你能夠做事的能力.
神裡面的能力.
是神親自給你.
不在乎人.
是在乎神自己.

$^{721}$所以他說.
你這樣求.
就是難得的意思.
但是他說.
如果你親眼見到我.
你到去昇天的時候.
就是說.
神讓你打開了你的眼睛.
打開了你屬靈的眼睛.
見到屬天的事情.
這個就是神給你一個印證.
他應運你.
所以之後是怎樣呢.
以利沙親眼.
就見著他自己的師傅.
成旋風.
昇天.
他親眼見著.
那代表什麼呢.
然後他就拿起.
以利沙.
以利亞掉下來的那個外衣.
其實.
心裡面都很震驚.
有些人就說.
哇 好厲害.
他現在拿起那個外衣.
撲下去.
河水就分開了.
其實不是那種的心情.
我相信.
而是.
其實你都不知道行不行.
震驚一拿起那個外衣.
來撲水.
然後神給他一個很大的肯定.
就是.
若淡河.
河水分開.
這個是以利沙.

$^{761}$所行的第一個神蹟.
他施傅最後一個神蹟.
都是叫若淡河分開.
他承接他施傅的職份.
第一個的神蹟.
也是將若淡河分開.
所以我們見到.
神真的.
如何將以利沙.
完全改變成為他的器皿.
其實.
我們見到.
以利亞和以利沙的分別就是.
其實以利亞.
一走出來.
哇.
一個.
每個人都見到.
高高在上的先知.
他很厲害的.
在爭戰.
但是.
以利沙.
當我們看他之後所行的事.
其實你發覺他不是一個.
好像很外顯.
很高高在上的先知.
而是一個走入人群.
在人平常的生活裡.
來幫助人.
亦都是行神蹟的先知.
我想第三樣要大家學習的就是.
他憑著神給他的信心能力.
他就會行動.
然後他其實.
與其做這樣的.
哇.
爭戰的先知.
其實他就是做一個.
走入人群裡面.

$^{801}$在人.
你會見到一會兒.
我的表.
他所行的神蹟.
很多都在人平常的生活裡面發生.
那些叫需要.
在人平常的需要裡面.
顯出這些神蹟.
神的不平凡.
去祝福人.
你看他之後.
一個這麼平凡的人所行的神蹟.
其實.
今天沒有時間.
逐樣和大家講.
如果你去看聖經.
你看聶王幾下.
由二章開始看.
看到十三章.
其實你就會看到.
哇.
其實一個這麼平凡的人.
為什麼可以做到這麼多事.
他潔淨河水.
他醫治人.
他醫治.
他叫乃曼.
醫治了乃曼.
他又叫窮的寡婦.
得到油.
得到供應.
他讓人復活.
死人復活.
甚至當他死了之後.
神給他的神蹟其實就是.
他死了之後.
死的人.
碰到他的骸骨.
都能夠復活.
厲害.

$^{841}$但是.
其實.
厲害的是誰.
是神自己.
爾利沙本身就是一個神蹟.
他自己本身就是一個神蹟.
是不是.
其實弟兄姊妹今天.
我們都有一個神蹟.
我們能夠信神.
都有一個神蹟.
其實我們每一個人都可以被神用.
成為一個神蹟的器皿.
你願不願意.
有時候我們說我們經歷不了神.
其實就是因為我們.
不夠膽給神去用.
覺得我們做不到.
但是你要記住.
如果你願意被神用的時候.
你會經歷神蹟.
你自己都成為一個神蹟.
最後我們.
來總結一下.
爾利沙的一生.
我們剛纔講的三點.
就是我們要堅定地跟隨.
堅定地跟隨.
但是你要很清楚.
今天.
你跟隨神的人.
記住是跟隨神.
也都求神.
擺在我們每一個人的心裡面.
關於他自己的國的事.
神不是說.
我今天創造了你出來.
就擺你在地上.
你做你的事.
我有我的事.

$^{881}$神說不是.
你是我的創造.
你要記得.
我的事.
在你的眼裡.
記得有神國的事.
祂不是叫我們.
擺在地上上班.
放工.
喫飯.
玩.
去旅行.
這些我們都要做的.
但是你不要忘記.
最重要的是.
是神的國.
神的國在我們的心裡面.
所以求神幫助我們.
而當我們來跟隨祂的時候.
我們說主啊.
求禰賜給我.
從來以來的託付.
以我的這個人.
我這個可以承擔些什麼.
在禰的國裡面.
求禰讓我知道.
讓我一生能夠向著禰給我的標杆.
來到直跑.
然後我們說主啊.
我沒有能力.
我求禰.
禰的靈.
禰感動.
以利亞的靈.
加倍來感動我.
主我需要禰.
自己更多.
我需要的是禰.
來幫助我.
然後我們就憑信心.

$^{921}$兄弟姐妹.
我們有時候在這裡敬拜.
好熱烈啊.
整天都站著.
但是走出去的時候.
就做一個隱形的紀錄圖.
是不是這樣.
其實.
將神的能力.
將神的神蹟.
帶出去.
什麼是神蹟.
當我們去傳福音.
跟人說福音.
人信主.
是一個神蹟.
Amen.
今天我們每一個人都是一個神蹟.
我們能夠相信神.
我們的生命能夠完全被改變.
由走去死亡的路.
向永生.
是不是一個奇蹟.
是不是一個神蹟.
當我們用神的愛.
神的話.
我們為人禱告.
去關心人.
人有一個絕望的境況.
變成一個有盼望.
因為神的緣故.
是不是一個神蹟.
當我們去為人禱告.
有人生病.
有人在困難裡面.
之後禱告.
神醫治他.
是不是神蹟.
這些都是神蹟.
當我們將有限的金錢.

$^{961}$能夠願意跟有需要的人分享的時候.
讓他好像寡婦得到油一樣.
得到神的供應.
是不是神蹟.
兄弟姐妹.
這些都是神蹟.
其實我們.
其實講真.
這些事全部神自己都可以做的.
醫治人.
供應人.
他用不著我們做.
其實馬拉都是從天上掉下來的.
他今天說這個人沒錢.
他掉一塊錢在他面前.
直接掉下來也可以.
是不是.
對於神來說有多難.
但他為什麼要我們去做.
因為祂是我們的神.
我們是祂的代表.
今天我們全部都是祂在地上的代表.
大家知不知道.
知不知道我們自己的身份.
祂說不是我能夠做.
但你去為人祈禱.
我醫治他.
你跟人講福音救恩.
然後讓他得救在我的靈.
叫他能夠信我.
全部都是神的工作.
所以當我們為人祈禱.
人得醫治.
我們不用沾沾自喜.
你看我多厲害.
你看那個人得醫治.
sorry.
完全是神的工作.
但神要我們同樣經歷這個神蹟.
神要我們自己都成為那個.

$^{1001}$Amen.
\newpage



\section{雅各書 1:19-21}
\label{sec:031751}
\textbf{重塑生命系列 (5) 忿怒與溫柔}~~~~張健庭牧師博士
\newline
\newline
連結: \href{http://www.church.com.hk/acms/content.asp?site=cdc\&op=show\&type=product\&code=031751\&layout=sermon}{\texttt{031751}} ~~~~ 語音日期: 2026-03-01
\newline
\newline
\hyperref[sec:031750]{\small{< < < PREV SERMON < < <}}
~
\hyperref[sec:index_chronic]{\small{[返順時目]}}
~
\hyperref[sec:031752]{\small{> > > NEXT SERMON > > >}}
\newline
\newline
雅各書 1:19-21
\newline
\begin{longtable}{cl}
\hline
\hline
章節 & 經文 (和合本修訂版)\\
\hline
1:19 & \begin{tabularx}{0.7\textwidth}{X} 我親愛的弟兄們，你們要明白：你們每一個人要快快地聽，慢慢地說，慢慢地動怒， \end{tabularx} \\ \\ \relax
1:20 & \begin{tabularx}{0.7\textwidth}{X} 因為人的怒氣並不能實現神的義。 \end{tabularx} \\ \\ \relax
1:21 & \begin{tabularx}{0.7\textwidth}{X} 所以，你們要除去一切的污穢和累積的惡毒，要存溫柔的心領受所栽種的道，就是能救你們靈魂的道。 \end{tabularx} \\ \\
[1ex]
\hline
\hline
\end{longtable}
$^{1}$好 弟兄姊妹平安.
好了 今天我們一起繼續思想這個.
第三宗罪就是關於憤怒的罪.
那我們一起來低頭祈禱.
念父愛院 你教導我們.
主真的好像我們所剛剛讀過的聖經.
給我們一個溫柔的心.
領受所栽種我們的道.
因為唯有你的道在我們的心裡面.
才能夠結出美善的果子.
願你自己的話語.
再一次提醒我們.
讓你的聖靈影響我們中間來運行.
讓我們去聽到你自己的說話的時候.
我們自己親自的回應.
奉耶穌基督的名祈求 阿們.
好 今天我們講這個關於憤怒和溫柔.
其實不知道大家有沒有發覺.
現在你活在這個世界上 很多人是很憤怒的.
有沒有發覺.
就是我剛剛上幾個月前.
我回香港的時候 見有一個朋友.
他說 你出街都要小心點 類似.
我忘記了正確的語句.
他說 現在不同以前 現在很多人很兇的.
可能他提醒了我.
有時候不要那麼容易信任別人.
但是真的 現在我們活在一個時代當中.
很多人都好像很憤怒.
其中一個我覺得憤怒的原因.
是因為大家都很在社交媒體那裡去拿資訊.
你不知道有沒有發覺.
在很多社交媒體當中.
其實它是將你的情緒不斷的去擴大.
它的運作背後的原理.
就是你喜歡聽什麼 它就不斷的去餵給你.
所以其實慢慢慢慢你的回音谷是越來越大聲.
但是在人所接受到的資訊裡面.
最大聲 最容易影響我們的情緒.
很多時候都是那些負面情緒.

$^{41}$那些憤怒 那些嫉妒別人.
是最容易牽動到別人的情緒.
現在在一個這樣的社交媒體 一個反彈的時代.
就是不斷的去催促人用憤怒來回應很多事情.
很多時候特別是要即時的回應的時候.
很多時候的憤怒 可能你想清楚你就不會那麼憤怒.
或者你想清楚可能都不是這樣的.
但是如果你即時回應的時候 那些憤怒就發作.
最可怕的一件事就是如果你在社交媒體發憤怒.
就沒有什麼代價的.
譬如你見到一個真人 人與人之間.
你談了一些事情 大家一言不合 那怎麼辦.
和他吵架 或者甚至反面 你會不會這樣做.
很少這樣做的 除非真的很忍無可忍的 為什麼.
因為你要付代價的 你亂和別人吵架.
沒有人願意和你做朋友的.
你亂和別人吵架 在公司沒有人願意和你同一個隊的.
你付很沉重的代價 你都會克制一下自己.
但是如果你在社交媒體是沒有代價的.
甚至你越是表現激烈的言論.
可能很多人都很喜歡 還給你一個讚.
這個時候 現在我們就活在一個很憤怒的言論.
而更加在很多政治的操作那裡.
現在不知道你有沒有發覺.
很多選舉 或者關於政治的那些評論.
其實是沒有議題的 不知道你有沒有發覺.
特別選舉的時候 很多時候都不談論.
究竟有些什麼應該做不應該做.
或者做這個政策有什麼好處有什麼壞處.
很少有這些討論的.
很多時候就是誰誰是我們的敵人.
我們不可以給他什麼.
誰誰他上了說他是希特勒.
他是獨裁者 世界末日.
差不多就是這樣.
就是去黑化對手來拿選票.
但是其實本身是沒有任何的正面的解決方案的.
譬如你說 現在這個世界很大問題.
譬如很多非法移民 舉例來說.
我不是在選邊站.

$^{81}$但是我只是問一個問題.
你說不贊成要去那麼多非法移民來的時候.
其中一個做法可能就是要遣返他們.
但是有些人支持的.
支持的時候 支持很多非法移民繼續在這裡的時候.
那又怎麼樣呢.
你明不明白 他沒有任何方案的.
因為赤字各方面又越來越多的時候.
你支持的時候 那個方案又是怎麼樣呢.
又沒有說的.
是不是不理呢.
諸如此類 很多這些.
現在那個社會的分化越來越嚴重.
當然其實這個跟剛才所說的政治.
不是 剛才所說的社交媒體的操作是完全有關係的.
但是也跟背後的政治的操作是很有關係的.
人與人之間 本來其實談政治.
其實是在談一個公共領域的一些事務.
究竟大家講道理而已 你對或者你不對.
都可以講個理由出來 大家評評道理.
大家看看誰多人支持一點 那就是那個了.
其實在一個公共領域那裡.
應該就是這樣來去運作的 是不是.
但是現在其實那些就不是了.
現在就講身份政治.
就是身份政治你知道是什麼.
就是說你是什麼身份的人 你就投什麼的政黨.
或者你的身份一定要支持我們的.
就是我們是為你爭取權益的.
但是其實爭取什麼權益呢 是不是爭取得到呢.
那些是完全之後就沒有人講的了 是不是.
所以會在一個非常大的一個矛盾當中.
而這些矛盾是刻意地去操作.
事實上這個也是.
你平心而論 你想一下一個事情就是這樣.
如果一些平心而論可以討論的事情.
就必然大家會想到.
或者起碼都會談到有一點正面的方案出來的.
就是譬如有一個問題.
如果大家都是平心而論來討論的時候.

$^{121}$慢慢慢慢談一談.
就算你不是說完全找到一個解決方式.
起碼都有一點正面的東西應該可以會出來 是不是.
但是現在的問題就不是這樣的.
現在很多時候 你有一個問題出現的時候.
就只是罵罵罵那些跟你不同意見的人.
因為其實是不打算解決問題的.
在那個罵的過程當中其實是不打算解決問題的.
其實有點像什麼呢.
有點像黑幫 幫派那些的糾紛.
那怎麼辦呢.
罵罵罵 看誰用廣東話俗一點叫「搭馬」.
招收哪些人站在我那邊的時候.
打起來的時候就打贏了.
就是這樣來解決問題的.
所以現在那個政治的操作就是這樣.
其實人與人之間的衝突也是這樣的.
就是我沒有實際的數字去找到.
但是我不知道你有沒有發覺.
現在在街上的人開車都暴躁了很多.
有時候很兇的.
有時候那些人惡言之外你都懶得駁斥他.
為什麼呢.
怕他發瘋.
Wallage那些有套電影可能你都沒有看過.
那個是有一個Wallage的開車憤怒.
那個起因就是那個女主角.
當時趕著上班送兒子上學.
然後有一個心理不太正常的人.
接著就打他兩巴掌.
接著他就追著他的車.
接著就搞到很大件事.
但是這個世界現在是這樣的.
很多的憤怒.
因為人與人之間的憤怒不懂得怎麼處理.
而事實上經過那些社會的媒體.
社會媒介.
政治的操控.
因為你看到很多人都是你這個世界的壞人.
或者你在身邊很多人都是壞人的時候.

$^{161}$你還用這麼客氣.
因為你跟好人 跟好人就客氣.
你跟個個都是壞人的時候.
我為什麼這麼客氣.
我客氣可能還會被你欺負.
所以這個是有些人的處事方法.
一有什麼不對的時候.
那又怎樣.
先發惡.
發惡完之後.
其他人就不敢出聲.
有時候都可以的.
有時候得到他所以得到.
就是這樣.
現在我們就是活在一個這樣的時代當中.
在一個這樣的時代當中.
事實上真的都挺困惑.
你想一下如果一個小孩.
在一個家裡的成長.
從小到大.
每樣東西都是靠惡.
靠大聲.
那小孩長大之後會怎樣.
就是這樣 靠惡靠大聲.
但是現在這個社會整個都是這樣.
那怎麼辦.
而事實上求神真的要幫助我們.
有時候我們首先第一個要處理的時候.
首先要處理我們自己.
因為整個社會很多東西我們都處理不了.
但是我們怎樣去看自己.
或者我們怎樣去明白到.
原來憤怒真的會將人去摧毀的時候.
我們自己怎樣要警惕.
因為就好像以賽亞仙主義所說過.
我是嘴唇不潔的人.
我會在嘴唇不潔的民眾.
以賽亞仙主義很明顯.
他神呼召他的時候.
他明白到這個世界的恐怖.

$^{201}$他都不敢掉以輕心.
而同樣都是一樣.
如果今天我們要說.
說憤怒是一個很嚴厲的罪的時候.
但是我們就要明白.
其實我們現在就是活在一個.
很憤怒的一個世代.
很多人是很容易發怒的時候.
那我們自己也要小心.
首先說一下憤怒的本質.
基本上我問一下你們大家.
你有沒有想過在你自己的時候.
在什麼情況你會很憤怒的呢?.
其實就是覺得好像被委屈.
被屈的那種權利被侵犯.
本來是屬於你的東西.
或者是一些你自己的東西.
而被侵犯了.
這個就是憤怒很容易發生的一個原因.
就是感覺到權力被侵犯的時候.
譬如感覺到權利或者尊嚴被侵犯的時候.
譬如第一種就是.
我覺得被人忽視.
又不夠尊重我.
這樣會生氣的.
是不是?.
譬如舉例來說.
你可能在.
我作個例子.
你在公司開會.
剛剛在說你的方案.
你都還沒說到一半.
接著有個同事一句說過來.
這些是廢話.
沒有用的.
你可能會很生氣的.
是不是?.
因為你感覺到你不被尊重.
你感覺到你不被尊重.
或者你感覺到有人忽視我.

$^{241}$而事實上聖經那裡都有的.
聖經那裡講到以西帖書.
講到末底改怎麼得罪哈孟.
因為哈孟的時候.
他就是一人之下萬人之上.
皇帝重用的宰相.
每個人見到他都卑躬屈膝.
就是和他打招呼.
末底改就不理睬他.
是不是?.
他不和他卑躬屈膝.
他就很生氣.
生氣到要殺人全族.
殺他一個不夠.
殺他全家都不夠.
殺他全族.
他那種都是很生氣的.
所以以西帖他記那裡講到.
哈孟那種其實是他感受到.
他的自尊被傷害.
第二種就是感受到威脅.
感受到威脅其中有一個.
很多時候會發生的就是.
你感覺到一些東西失控.
不在你控制之外.
那你就生氣.
是不是?.
就是你感受到威脅.
其實這個有時候.
如果你家裡有養寵物都會有的.
就是譬如我那時候小孩子的時候.
我媽媽就教我.
那些狗吃東西的時候.
你不要搞牠.
是不是?.
因為牠可能會咬你的.
因為感受到威脅.
牠不是威脅.
牠感受到牠的食物可能會給你拿.
那牠會咬你.

$^{281}$而事實上.
人覺得自己受威脅的時候.
又會發怒.
是不是?.
或者有些東西不受你控制的時候.
你也會發怒.
就是你明明已經計劃.
可能你也試過了.
有些東西你很精心設計.
計劃了一二三四五這麼多的步驟.
但是去到最後就完全打岔了.
那你完了之後你也會很生氣.
或者在過程中被人打岔的過程.
你也會很生氣.
是不是?.
第三種.
就是有些未處理的傷害和腐毒.
這種我覺得就比較隱藏一點.
但是其實是會有的.
我們也有句話叫做.
賓口人忌賓口碗.
那賓口人為什麼會忌賓口碗呢?.
有沒有想過?.
就是因為可能其實.
牠有些傷害.
而你剛剛說的那些東西.
其實本來就不是說牠的.
但是可能就剛好擦中.
其實大家很容易想到的.
就是本來譬如人家碰你一下.
其實你不會覺得痛的.
但是如果那裡受過傷.
或者傷口還沒有好.
碰你一下你立刻痛到跳起來.
同樣都是一樣的.
有些傷痛本身你還沒有處理.
譬如那個人以前跟你有過一些過節.
你表面上說我寬恕了他.
我和他和好了.
我現在沒事了.

$^{321}$但其實又沒有處理.
又不是的.
到時可能他說了一點點的話.
你就即時就好像.
他殺父仇人那樣來攻擊他.
是不是?.
就是因為那些沒有處理的傷害.
很多時候這些我們都有的.
有時候你為什麼會有一些這樣的反應.
對某些人或者某些事.
可能有些事都是一樣.
你以前出現過類似的事情.
所以你會反應很大的.
如果這樣的時候.
你說憤怒帶來什麼後果呢?.
有沒有想過?.
憤怒帶來什麼後果呢?.
憤怒帶來的後果.
其中一方面就是.
我們再不能夠用正確的視野去看事情.
我們失去了正確的視野.
聖經那裡說.
「人的憤怒並不成就神的義」.
這裡很有趣的.
雅各說.
他說「人的憤怒並不成就神的義」.
意思是什麼呢?.
意思反過來說.
就是怎麼說呢?.
有些人覺得憤怒.
他自己的憤怒成就了神的義.
是不是?.
如果他不是這樣想.
雅各就不會這麼說.
事實上.
你想到其實也挺扭曲的.
你發完憤怒.
還覺得你的發憤怒是理直氣壯.
可以成就了神的義.
你說會不會有這樣的事情發生呢?.

$^{361}$當然有.
如果沒有的話.
聖經也不會這麼寫.
事實上你就看到.
其實發憤怒本身.
有時候會使我們的視野扭曲了.
我們覺得自己.
其實發憤怒其中一個.
很理直氣壯的原因.
就是什麼呢?.
就是你覺得自己是對的.
有時候你不覺得自己是對的.
你也心虛.
你也不會發憤怒.
也不會發得這麼盡.
是不是?.
就是其中一件事.
發憤怒你覺得自己是對的.
譬如剛才所說的那些事情.
你也覺得自己是對的.
就是我被人忽視.
我的尊嚴被受損.
那些事情我明明想好了.
現在被你搞砸了.
那你就發憤怒.
就是你完全覺得自己是站在一個對的位置.
發憤怒本身.
使我們的視野.
只是站在自己的觀點去看事情.
有時候.
當你靜下來去想一想.
從多方面去想一想的時候.
你沒有這麼容易生氣.
或者你會沒有.
可能生氣完你也沒有這麼生氣.
為什麼呢?.
因為可能.
也不是的.
可能其實我當時也有些錯的.
那你就沒有這麼生氣.

$^{401}$是不是?.
但是如果你越想越生氣.
我明明都是對的.
那個事情這樣這樣.
你就越想越生氣.
是不是?.
因為你的視野是沒有辦法矯正的.
所以在尊嚴那裡所說.
人有見識就不輕易發怒.
為什麼呢?.
因為你有見識的人.
其中一個方面就是.
你會看多方面的時候.
你不只是看自己.
那你就沒有這麼容易生氣.
所以這個第一個屬靈的後果的時候.
經常發怒的人.
你要想一想.
其實你的視野都扭曲了.
很自以為宜.
覺得自己是對的.
覺得自己一定是對的.
第二個屬靈後果.
當然是關係上受損.
是不是?.
關係上受損.
但是這個關係上的受損.
其實不僅是人與人的關係受損.
人和神的關係都是受損的.
人與人的關係受損.
你也明白的.
不用多說了.
就是生氣之後.
發完氣之後.
發完脾氣之後.
跟那個人有時候說話.
都好像面左左.
都會有的.
這個人與人的關係受損.
這個我也明白.

$^{441}$但是你有沒有想過.
其實發怒本身.
也是人與神的關係受損.
所以主耶穌說.
如果你去獻祭的時候.
想起有人向你懷怒.
你就放下你的祭物.
先去跟他和好.
然後再獻祭.
其實就是這個意思.
因為如果你跟人的關係.
在憤怒當中沒有處理.
神不予納你的祭物.
所以也都影響了人與神的關係.
所以有時候聖經那裡說到.
主禱文也有說.
我們謠說別人的過犯.
叫神也都謠說我們有過犯.
那個也是有關係的.
我們要跟神的關係正常的時候.
就不可以被憤怒來控制我們.
第三個當然就是.
靈命的僵硬.
這個就是用苦毒的人無法禱告.
也都無法饒恕.
這個也是很明顯的.
如果你心裡面還有很多苦毒的時候.
你講感恩其實是不合適的.
好像亞國書也有這樣說.
你心裡面充滿了苦毒.
一個泉源有些苦毒的.
你又怎麼能湧出那些讚美的水呢.
其實是不行的.
剛才也說過你不肯謠說的時候.
神也都提醒我們.
我們不謠說別人.
神也都不謠說我們.
主禱文那裡教的.
所以整體來說.
憤怒本身是將人的屬靈的生命摧毀的.

$^{481}$屬靈的生命如果有很多的憤怒不處理的時候.
這些就好像放一些毒藥.
如果你種植物的時候.
放了一些不合適的肥料.
有些不合適的肥料變成某些植物的毒藥.
你會發覺為什麼慢慢地那棵東西不斷下肥.
越來越癢.
癢到葉子又黃了.
然後慢慢地死掉了.
就是因為那些憤怒.
就是我們心靈裡面的毒藥.
所以聖經所說.
溫柔其實就是對於憤怒的一個相對的.
但是另一方面來說.
我覺得勝經所說的溫柔本身.
其實也是憤怒自己的藥.
憤怒這個毒是很難解的.
但是聖經說你用溫柔就可以把它解了.
溫柔的意思原本.
溫柔的希臘文.
原本就是在說一些馴服的野馬.
如果你想象到憤怒是一些不受控制的力量.
溫柔就是把那些野馬的力量被馴服了.
被馴服了.
被修復了.
而主耶穌自己當然是很溫柔的.
如果你說聖經中有誰是很溫柔的呢?.
有個評價就是摩西.
在舊約的時候.
神對摩西的評價.
他說摩西是一個很溫柔的人.
但是新約的時候.
其他就沒有什麼評價.
但是新約的時候主耶穌自己說.
祂自己對自己的評價.
就是我心裡柔和謙卑.
所以主耶穌也是一個很溫柔的人.
但是你想想主耶穌的溫柔.
是不是代表祂沒有憤怒的呢?.
是不是呢?是不是代表祂沒有憤怒呢?.

$^{521}$不是的.
當然主耶穌的溫柔不代表祂沒有憤怒的.
祂有發憤怒.
我在聖殿那裡趕走那些賣賣的人.
那次就發憤怒了.
那次真是.
你看你不要以為主耶穌很溫柔地趕走他們.
你想想.
主耶穌.
那時候聖經形容推翻別人的桌子.
真的跟別人反桌子.
然後怎麼樣.
還要拿一條繩子.
好像鞭一樣鞭.
我不知道鞭那個地.
還是鞭那些動物.
還是鞭人.
我就不知道了.
拿一條繩子好像鞭一樣.
你想想.
很溫柔地鞭.
根本是不可能的.
所以那個是一個很激動的.
很激烈的一個行動.
很激烈的行動.
所以.
主耶穌是沒有.
也有祂的怒氣的.
但是主耶穌的怒氣和憤怒那種是有什麼不同呢?.
或者再看看.
主耶穌也有責備那些返來賽人的.
祂跟別人舌劍唇槍的時候.
祂也罵別人.
說你毒蛇的種類啊.
誰教你脫離將來的憤怒呢?.
祂罵別人毒蛇的種類.
其實這個意思就是說.
其實這些是說得我們很低賤的.
都是很嚴厲的.
那麼.

$^{561}$那麼那個分別在哪裡呢?.
我就問你了.
主耶穌也有發怒啊.
但是我們也會發怒啊.
那麼如果發怒.
發怒本身如果你看為一種力量的時候.
是一種心靈裡面的一種.
湧現的一個.
對於一些事情的反應的一種衝擊的時候.
那麼主耶穌的發怒和我們的發怒有什麼分別呢?.
聖經就沒有教我們不要發怒.
但是很大的一個分別呢.
就是主耶穌不是出於那個自我的防衛機制而發怒的.
事實上這個是跟我們的發怒有很大的分別.
我們當然.
如果譬如你看到一些很不合神心意的事情.
你很生氣.
那麼這個都可能是可以合理的.
可能這個就是發意怒.
可能神都想你看到一些不合心意的事情.
你不是好像見到就當見不到.
覺得沒所謂.
就是神都不想我們這樣的.
但是那個問題就是說.
回想我們自己的時候.
因為如果我們今天說憤怒是一個罪的時候.
最大的分別在哪裡呢?.
最大的分別就是.
很多時候我們的發怒是出於自我的防衛機制的.
你想一下.
別人說一些話.
如果你是沒所謂的.
那麼你就不會生氣的.
但是如果說一些話.
你覺得你受到傷害的時候.
你要反抗的時候.
有時候有些人就馬上表現上已經很激烈了.
就是馬上可能撲口大罵或者怎麼樣.
但是即使有時候我們有教養一點.
或者有修煉一點.

$^{601}$你不是那麼激烈.
馬上有很大的跟別人反駁.
但是可能你心裡面都一樣很生氣.
一樣會有的.
是不是?.
因為歸根結底就是.
我們心裡面每個人是有一個很強的自我中心.
而這個自我中心也都是有一個很強的保護自己的機制.
有些什麼不合自己的心意.
有些什麼覺得在你的掌控之下你覺得失控.
有些什麼是覺得屬於你的.
你不可以被人拿走.
你的尊嚴或者什麼.
不可以被人拿走.
有些什麼你是覺得順理成章.
只是應該是屬於你的.
而你各方面被人不同的程度上冒犯.
你就會有反應.
而這個反應都是出於以自我為中心.
這些就是聖經所說的那種憤怒.
憤怒本身是一個情緒.
情緒本身就好像一條河流一樣.
河流可以流一些乾淨的水也行.
流一些骯髒的水也行.
因為它是一個channel 一個河流.
有些是合乎神心意的情緒也有.
不合神心意的情緒也有.
譬如說愛,恨,憤怒.
這些都是情緒的一部分.
每個人都有的.
神不是說叫你信了耶穌之後.
禁你一情六慾.
連憤怒也沒有 像一個得道高僧那樣.
看破世塵 怎麼戳你都沒事.
我們做基督徒不是做這一種.
色即是空 空即是色.
什麼反應都沒有那種的.
而是我們仍然都有我們的情緒.
但是在我們的情緒的時候.
我們就懂得怎樣去疏導.

$^{641}$我們要懂得分辨.
哪些情緒是合神心意的.
哪些情緒是不合神心意的.
而憤怒本身有一些是合神心意的.
就好像主耶穌.
祂也為了那些不義的事情憤怒.
但是問題就是對於我們來說.
很多時候我們活在一個罪惡的世界.
我們自己也是一個犯罪的人.
很多時候我們那些情緒.
同樣是憤怒的情緒.
很大部分不是合神心意的.
很多時候都是出於.
自我中心的防衛機制而來.
而這樣的時候.
所以聖經所說的溫柔.
就是說是要一種修復它的力量.
溫柔本身不是說那些力量是沒有的.
譬如好像水那樣.
你知道那些.
譬如有些地方要建水壩.
其實它是幫助那些水.
讓它可以正面的將那些衝擊的力量.
來發電的.
同樣一樣我們心裡面那些情緒.
是需要在神裡面來去純正的.
去改正的.
而這樣說的時候.
溫柔有什麼特質呢?.
溫柔是不是一個人說話.
都是很女性化的.
不會大聲的音配都是低兩度的.
是不是這樣就代表溫柔呢?.
有沒有試過有人被人很溫柔的去罵過?.
有沒有?.
你不要笑,其實是有的.
當然是有的,什麼沒有.
有些人是很溫柔的.
其實他在罵你的時候是很刻薄的.
是不是?.

$^{681}$所以那個不只是語氣的問題.
溫柔不是說一個語氣.
溫柔其實是說一個內在的素質.
一個內在的素質.
首先就是剛才所說的憤怒是一個情緒的時候.
所以我們的情緒是需要受聖靈來導引的.
我們的情緒本身.
有時候我們自己沒有刻意去反省的時候.
好像覺得要生氣就生氣,要笑就笑,要怎樣就怎樣.
特別是在西方的教育制度之下.
小孩子成長就盡量要少給他這些限制.
讓他自由發展,他想怎樣就怎樣.
讓他做回一個真我.
這也是現在一個社會當中非常大的一個方言.
這個方言是什麼呢?.
就是要活出真我這個方言.
你有沒有聽過?.
活出真我這個方言.
其實有意無意之間.
無論你在潮流文化.
無論電影,電視.
其實都很普遍的.
人覺得最高的價值是要活出真我.
活出你真正的自我.
特別是你上了一把年紀.
你都知道這些其實很多都不是太有用的.
但是年輕人是很相信這套東西的.
現在我們會在一個時代裡.
覺得要authentic,要真實.
所以你的情緒本身.
任何的情緒發洩.
只要是真性情.
都可以接受的.
不是身邊的人接不接受.
是你自己接受了自己的真性情.
你明白我什麼意思嗎?.
就是這樣.
但是.
其實聖經不是這樣教導的.
聖經教導不是說你所有的東西.

$^{721}$只要是真性情對得起自己.
就都可以不受控的.
其實不是這樣的.
反而我們的真性情.
就是被聖靈管治的性情.
因為聖靈去管治你們.
就是你的賞望.
或者你希望什麼.
用聖靈來管理你所希望的東西.
你喜歡什麼不喜歡什麼.
用聖靈來調節你.
譬如你以前很喜歡.
譬如我舉例.
有個人很喜歡跳舞的.
舉例來說.
他經常都想著跳舞的.
但是當他信了耶穌之後.
可能聖靈去管理他.
他可能那個心思意念.
不是那樣東西.
他反而更加著重.
可能他現在更加愛慕去敬拜神.
更加愛慕去親近神.
無論是聽詩歌,讀聖經諸如此類.
我們所愛慕的東西.
我們讓聖靈去管治.
同樣的憤怒的性情.
其實當被聖靈去管治的時候.
就成為了溫柔的一個特質.
而第二樣.
因為溫柔的人.
是回應而不是反應.
回應就是respond.
反應就是react.
兩個是不同的.
你知道很大分別.
react就是有些人.
一件事都是react.
你知道什麼意思嗎.
根本他講得白一點.

$^{761}$他的process(過程).
就是很即時本能地去反應.
而他不受約束.
甚至乎可能都skip了.
譬如我們小朋友說.
家裡的教導說.
你講話要用腦.
有沒有聽過.
我們小朋友.
意思就是說.
我們每個人講話都用腦.
怎麼會不用腦怎麼講話.
但是你明白背後的意思就是說.
你要process(過程).
你要想想有什麼後果.
你講完這句話.
你經過process(過程)之後.
那個就不是react(反應)了.
那個是respond是回應.
是不是.
就不是反應了.
你沒有經過process(過程)的.
那些就是react(反應).
就是這樣的意思.
事實上.
很多時候我們都需要去.
有所節制.
所以聖經那裡.
剛才我們讀的那段聖經.
不是說要快快的聽.
慢慢的說.
是不是.
慢慢的說慢慢的動腦.
那個慢慢的意思.
其實就是教我們.
你process(過程)一下.
就是處理一下.
你不要unfilter(清除).
就是毫無間隔的.
unfilter(清除).

$^{801}$就是這樣來到.
什麼都出來了.
但不只是講語言的.
其實就是在講我們那些憤怒.
其實有時候.
是很多時候即時的出來的時候.
其實我們要處理它.
有時候不需要這樣即時的出來.
第三個就是用真理說話.
但都是帶著恩典的.
用真理說話都是帶著恩典的.
當然我們待人處事.
我們都是要有一個真誠.
和一個誠實.
一個真誠和一個誠實.
但是有一個真誠和一個誠實.
去說話的時候.
也都不代表說.
你可以拿了面試金牌.
或者拿了尚方寶劍.
就可以斬殺誅邪.
到處去開機關槍掃人家那些不對的東西.
是不是.
因為其實我們雖然是講一些真實的說話.
但是都要帶著恩典.
聖經那裡說我們要說話.
要帶著鹽用鹽調和.
用鹽調和.
鹽是怎樣的呢.
我們中國人就.
你生妃子喝點鹽水.
即是怎樣.
下火嘛.
是不是.
鹽是下火的.
當然聖經那裡未必是有這樣的意思.
但是帶著鹽調和.
我想某個程度.
就是說我們的說話是有一種溫柔在其中.
我們有時候講真實的說話的時候.

$^{841}$我們都要選擇適當的時機.
是不是.
也都選擇用適當的表達.
更加重要的其實就是.
有時候我們的憤怒本身.
不是那麼容易靠自己就能夠處理得到.
如果我們那麼容易處理得到憤怒的時候.
就不會.
剛才所講的這個世界上越來越多.
World Wage(世界工資).
越來越多那些公路憤怒那些.
就是因為很多人都處理不了自己的憤怒.
怎樣處理.
就是憤怒有點像你的壓力煲.
就是很多年前.
我都忘記了多少年前.
不是有個什麼巴士.
巴士伯伯就在罵那個年輕人.
不是瘋了嗎.
那個視頻不是很多人看過的嗎.
就成了一個名言.
巴士伯伯就說.
你有壓力我也有壓力.
是不是.
事實上很多的憤怒本身.
是真的因為人有很多的壓力在裡面.
但問題就是壓力本身.
你是處理.
如果你不疏導它.
你不斷地去壓抑它.
是越壓抑就越嚴重的.
可能壓抑到點一爆出來的時候.
那就爆煲了.
所以憤怒本身是需要處理.
不是靠著壓抑.
那福音怎樣來幫助我們來處理呢.
第一就是十字架讓我們看到.
我不再需要自我防衛.
這一點是很重要的.
剛才所說的憤怒本身.

$^{881}$很多其實是自我防衛機制的反應.
你覺得.
譬如你要去保護自己.
你就會憤怒.
但是如果十字架那裡.
怎樣來告訴我們不需要防衛機制呢.
譬如我舉個例子.
譬如有個事情或者一個人.
說些話令到你覺得.
啊 我很不被尊重.
舉例來說.
哇 接著你就很生氣.
哇 為甚麼他可以這麼不尊重我呢.
為甚麼會這樣呢.
就是你很生氣.
但是如果在十字架的面前的時候.
多過真實的去認識到.
原來我最關注的尊重是從哪裡來的呢.
是從哪裡來的呢.
是從神那裡而來.
不是從人那裡而來的.
人對於我的不尊重.
No big deal(沒大不了).
真的No big deal(沒大不了) 你想想.
有時候你有意無意.
可能你都忽略了一些人.
都會有的.
有時候人忽略你.
你就覺得很大件事.
他為甚麼不尊重我.
但有時候你忽略人也有的.
你就發覺其實.
在十字架那裡其中福音幫助我們.
就是看到有一個正確的價值觀.
我們所找到的.
所要求的比如尊重.
要我們覺得.
我的價值在哪裡.
我自己在神面前是一個怎樣的人.
這些是這個世界動搖不了的.

$^{921}$有時候可能你聽到一句話.
你覺得人好像踩低你.
你會很受傷害.
然後你可能很生氣.
但你想想.
Who care(誰在乎).
有時候不需要這麼介懷.
有時候人怎樣看不是最重要.
但我知道我在神面前仍然是有價值的.
人可以看低我.
但是神仍然是這麼寶貴我.
是不是.
福音就是這樣帶給我們一個正確的.
來將我們所需要的那些寶貴.
其實拆解.
因為我們唯有在神那裡有一個正面的.
來看到自己是一個怎樣的人.
有一次有一個傳道人分享.
有時候有人很介懷.
覺得別人這樣說好像將他.
本來他這麼好的.
但好像將他降低了.
但傳道人自己說.
其實都無所謂.
其實他說得我這麼差.
其實我還差一點.
不過我在神面前.
神知道他不知道.
事實上是真的.
是不是.
我們還差一點.
你說我差.
其實我還差一點.
不過你不知道.
不過神就知道.
那又怎樣.
沒有關係.
神知道我還差一點.
都接受我.
是不是.

$^{961}$當你這樣的時候.
慢慢慢慢你那些很多防衛機制.
你又拆了.
那些不是一朝一日可以拆的.
因為你是經年累月.
在你做人的過程當中.
不斷的建立.
建了很多的.
又護城河.
又橋頭堡.
什麼都有.
一有什麼就開槍.
但是我們就是慢慢慢慢.
用神來拆.
我們自己裡面的防衛機制.
因為我們正確的在神面前來看自己.
而第二樣就是真實寬恕的力量.
剛才都講過.
有些是舊的傷.
你就被人碰一碰.
接著就發出來.
那歸根結底就是.
其實你還沒有真正寬恕那個人.
或者那些事情你仍然都放不下.
這個真的要求神幫助我們去寬恕.
而寬恕的力量.
真的唯有從福音而來.
人不斷的告訴你.
放下吧 寬恕吧.
你不寬恕.
最終都是難為了你自己.
其實每個人都會講這些.
你都可能跟別人講過.
但是到你自己的時候就不行了.
所以不關事.
我們真的唯有在神那裡.
才可以得到真正寬恕的力量.
就是求神來軟化我們的心.
讓我們看到我們那種寬恕是心裡面真實的寬恕.
慢慢你的傷口就好了.

$^{1001}$你的傷口以前是很深的.
好了之後.
可能真的仍然有疤痕.
但問題是你的疤痕是好了.
即使是有疤痕.
但是你都不會痛的.
是不是.
你的疤痕是比較噁心.
但起碼你不會痛.
但問題是有時候我們沒有處理.
不僅沒有去到疤痕的地步.
而是真的心裡面有些傷口.
所以一提起就有些東西.
好像一戳一戳它.
然後馬上那個反應.
超乎那個應該有的反應.
你就知道其實那些是舊傷.
我們真的需要去醫治.
我們真的需要在神面前求醫治.
所以當然也是聖靈在我們裡面.
結出溫柔的果子.
讓聖靈在我們裡面帶領我們.
無助我們.
這個是在裡面做.
不是在外面做.
最後就說一下怎樣去練習.
操練.
第一就是我們真的要學一下.
遲一點才說.
就是聖經也教我們要快快的停.
慢慢的說.
慢慢的動腦.
就是我們要遲一點才說.
有時候你即使很生氣的時候.
你不要即時回應.
你慢一點才回應.
這個也是一個練習.
有時候你過了那個不回應之後.
可能你之後就覺得不需要回應了.
有些東西.

$^{1041}$你可能也試過.
有時候即時很生氣.
馬上很想回應.
然後過了之後.
算了.
有什麼好回應的.
這些沒什麼.
是會這樣的.
就是你不要讓你的反應主導了你.
你的情緒是需要受控制的.
第二就是要分辨受傷和事實.
分辨受傷和事實.
這個也是很重要的.
有時候.
剛才也說過有些相關的.
就是那些反應本身.
不成正比的一些憤怒出來.
其實不是事實.
不是事實.
譬如舉例來說.
你聽了一句話.
你聽了之後很生氣.
你覺得這個人是在傷害你.
或者在攻擊你.
那你就很生氣了.
但其實他說的話可能不是事實.
不是在攻擊你.
只是可能是你聽為一個攻擊.
但其實可能.
或者問在場和其他三個人.
其他在聽的.
都不覺得他在攻擊你.
可能只是你覺得自己被攻擊.
然後你就很生氣.
你有沒有遇過這樣的情況.
也有的.
那問題就是.
事實本身不是這樣的.
但那個憤怒出來.
其實如果事實不是這樣的時候.

$^{1081}$就一定有其他原因的.
是不是?.
就是沒有無緣無故的愛.
也沒有無緣無故的恨.
當然有其他原因的.
那其他的原因.
就可能背後是有一些傷口.
可能根本那個人一看.
另外一個人已經很不順眼了.
他說的話的時候.
就覺得他一定是攻擊我.
那我就對號入座的時候.
他接收的就變成了他攻擊.
那他就馬上反應過來.
就馬上還火.
是不是?.
所以你要分清那個受傷和事實.
那個受傷本身.
我們剛才也說過.
就真的要求神來幫助你.
去處理那些傷口才可以的.
要不然的話都搞不定的.
就是你的世界.
就是你看的東西.
很多不是事實的東西.
就是很難搞的.
就是有時候去.
就是譬如一個輔導人.
就是一個輔導的角色.
一個輔導者的角色.
其實其中很大方面.
就是幫助那個受輔導的人.
搞清楚什麼是真什麼是.
什麼是真實什麼是不是真實的.
而事實上很多我們所建構的一些回應.
很多的東西.
為什麼會有這樣的回應.
究竟是不是真的事實呢?.
可能有時候就算是一個事實.
有個人說的話可能是得罪了你的.

$^{1121}$但是你都可以想清楚.
那個事實也都不是這麼表面.
那個事實可能是.
他都是無心之失.
有時候他說的話是很難聽.
但是他都不是真的要針對我的.
我平時都跟他OK的.
他只不過那時候說了一句話.
真的很難聽而已.
這個都是一個事實的另一方面來的.
就是你要搞清楚那個事實究竟是怎樣.
不需要說.
他一說這句話.
這個人就一定是跟你十冤九仇.
說要跟他血海深仇.
不是.
你要搞清楚那個事實才行的.
是不是?.
第三就是我們都是要學習寬恕.
所以聖經那裡所說的就是.
身體卻不可不要犯罪.
不可含怒到日落.
也不可給魔鬼留地步.
這裡說因為你不寬恕.
你被魔鬼留了地步.
魔鬼就最喜歡撩你這些.
不寬恕人的那些空隙的.
越撩越大.
洞越來越大.
然後魔鬼就鑽頭進來的了.
就是事實上人與人之間.
教會的事奉當中.
教會的關係當中.
很多都是.
這些空隙就是因為人不肯寬恕.
而魔鬼就留了地步在那裡.
剛才也說過了.
人有見識就不輕易發怒.
寬恕人的過失.
便是自己的榮耀.

$^{1161}$是不是?.
就是我們要寬恕別人的過失.
其實這個是我們的.
privilege(優惠)來的.
我們那個.
你想想你可以寬恕別人的時候.
其實這個是神給你的一個特權.
就是很多人不能夠寬恕.
其實他是很貧乏的.
你反而來說.
一個你豐富的人你可以寬恕.
因為這個是你的尊榮.
最後一點就是.
我們當然說話的時候.
都要學習一些溫柔.
就是少一點用.
你經常都是這樣的.
是不是?.
就是你少一點用.
都說了.
你就是這樣的人.
就是因為這樣的話.
就很容易被人覺得是人身攻擊的.
你說一件事就好了.
譬如你說.
那些菜很鹹.
你說這些菜很鹹就好了.
每次都煮得這麼鹹.
每次都覺得是攻擊的.
如果你說這次的湯很鹹.
那就好一點了.
但是每次都這麼鹹.
接著就有較壞的.
接著就說你煮吧.
是不是?.
就是這樣.
因為我們說事情的時候.
就不要去到人身攻擊的層次.
不會這麼容易挑起那些憤怒.
所以.

$^{1201}$另外反過來說就是.
我覺得這次很鹹.
所以這個都是兩方面.
就是學習.
就是你表達你自己.
就不要將矛頭.
以至手指指向別人那裡.
\newpage



\section{馬太福音 5:13-16}
\label{sec:031752}
\textbf{鹽和光}~~~~伍宇衛長老
\newline
\newline
連結: \href{http://www.church.com.hk/acms/content.asp?site=cdc\&op=show\&type=product\&code=031752\&layout=sermon}{\texttt{031752}} ~~~~ 語音日期: 2026-03-01
\newline
\newline
\hyperref[sec:031751]{\small{< < < PREV SERMON < < <}}
~
\hyperref[sec:index_chronic]{\small{[返順時目]}}
~
\hyperref[sec:031753]{\small{> > > NEXT SERMON > > >}}
\newline
\newline
馬太福音 5:13-16
\newline
\begin{longtable}{cl}
\hline
\hline
章節 & 經文 (和合本修訂版)\\
\hline
5:13 & \begin{tabularx}{0.7\textwidth}{X} 「你們是地上的鹽。鹽若失了味，怎能叫它再鹹呢？它不再有用，只好被丟在外面，任人踐踏。 \end{tabularx} \\ \\ \relax
5:14 & \begin{tabularx}{0.7\textwidth}{X} 你們是世上的光。城造在山上是不能隱藏的。 \end{tabularx} \\ \\ \relax
5:15 & \begin{tabularx}{0.7\textwidth}{X} 人點燈，不放在斗底下，而是放在燈臺上，就照亮一家的人。 \end{tabularx} \\ \\ \relax
5:16 & \begin{tabularx}{0.7\textwidth}{X} 你們的光也要這樣照在人前，叫他們看見你們的好行為，把榮耀歸給你們在天上的父。」 \end{tabularx} \\ \\
[1ex]
\hline
\hline
\end{longtable}
$^{1}$再一次歡迎各位.
有健康的身體 健康的靈命.
來一齊享受在主裡面的溫暖.
剛才讀的經文.
提供了今天訊息的內容.
「炎和國」是今天的講題.
求神藉著今天的訊息.
幫助我們認識每一天.
我們的生活必須有明確的目的.
和好的行為.
讓神得到榮耀.
一直到行完我們在這個世界的路程.
和主耶穌面對面的那一天.
希臘人研究到死海的鹽.
是有特別的用途.
所以他們稱這些鹽是神靈鹽.
Divine Salt.
因為希臘人用鹽醫治很多的病.
他們就用兩條蛇.
和他們醫學的神靈.
叫做漢默士.
作為他們醫治病人的標誌.
各位有沒有留意到.
很多醫療部門都用這個標誌.
羅馬人認為沒有什麼比陽光和鹽更加重要.
這個想法是因為受到兩種不同的信仰的影響.
就是羅馬人拜的太陽神.
希臘人拜的漢默士神.
就是剛才給大家看到的漢默士神.
所以羅馬人認為沒有什麼比陽光和鹽更加重要.
因為第一世紀的羅馬時代的信徒.
生活在這樣的社會文化背景裡面.
所以耶穌就用鹽和光來引出靈命上的教訓.
剛才讀的經文.
描寫出今天訊息的兩個重點.
第一個重點.
信徒是地上的鹽.
第二個重點.
信徒是世上的光.
求聖靈幫助我們明白這兩個重點.

$^{41}$請我們低頭祈禱.
請主耶穌求你的靈幫助我們在短短的時間內.
能夠在我們心裡面動功.
讓我們明白你給我們的教導.
奉主耶穌的名.
阿們.
今天希望用兩個重點來解釋清楚訊息的內容.
第一個重點.
有用的鹽.
用途有兩個.
第一個.
調味和諧.
比喻人和人之間過一個和諧的生活.
第二個防止腐爛.
防止變壞腐爛的人心.
第二個重點.
明亮的光.
在自然界裡面.
光可以放出熱能,電能,化學能.
幫助萬物生長.
在屬靈界裡面.
被主耶穌的光照的信徒.
在靈明上面得到掌盡.
這個重點也都包括兩方面.
光照的特質.
真光照入人心的時候.
啟示到人的生命有目的.
人的生命有價值.
第二方面.
光照人前.
憑著我們全福音.
使到光照在人間.
現在說有用的鹽.
第一個用途.
調味和諧.
鹽加食物之味.
比喻信徒有美好的見證.
就加添了教會的溫暖.
信徒必須藉著祈禱,讀經,默禱.
使到內心純潔.

$^{81}$行事為人.
得到神和人的喜悅.
這樣的信徒.
自然會使到教會成為一個溫暖的家庭.
吸引人來聽福音.
和信上帝.
耶穌教導門徒.
我們要有鹽的聖能.
和睦地一起努力事奉神.
榮耀神.
鹽本是好的,若失了鹹味,你們怎能用它調味呢?.
你們中間要有鹽.
彼此和睦.
馬太福音第五章.
亦記載耶穌這樣教導我們:使人和睦的人有福了.
因為他們必稱為神的兒子.
這樣我們才會彼此稱為弟兄姊妹.
我們稱為弟兄姊妹才有意思.
否則我們稱弟兄姊妹是枉然的.
信徒外在有好的儀態.
有些似企鵝的可愛作風.
吸引了很多動物來做他們的朋友.
信徒加上內在有好的涵養.
好像鹽的作用.
能夠促進和睦相處的氣氛.
自然會吸引其他人來加入這個溫暖的教會.
一同行走這條長遠的世界的路程.
和睦相處,彼此相愛是溫暖教會的實質.
約翰一書第四章約翰解釋實質的重點.
我們應當彼此相愛,因為愛是從神來的.
凡有愛心的都是由神而生,並且認識神.
有愛,你就是認識神.
《喜樂時光》是二零二五年出版的一本非常暢銷的著作.
作者是研究如何避免衰老的醫學教授.
他指出喜樂的要素和性格培養的要點.
喜樂的要素要有智慧,認識耶和華是智慧的開端.
要有愛心,要積極做事.
事事存在感恩的心,和藹仁慈待人.
性格要怎樣培養.
作者說我們要慢慢成長.

$^{121}$我們要有一個慷慨的性格.
適應歡樂或者憂傷的環境.
因為有神在看顧.
歡喜和外人交往為了傳福音的緣故.
作者指出喜樂的要素和性格培養的要點.
如果你們需要的話,我可以寄給你們.
作者特別強調,他說喜樂不可以裝在罐內.
隨時拿出來就好像很喜樂.
當你在困苦,憂傷的時候.
要記得喜樂的來源是由神傳來.
這樣就可以幫助你追求喜樂的心態.
我相信聖靈必定會引導神的兒女.
在人生的路程上常常享受到喜樂的時光.
不論在家裡或者在教會.
最先使外面的人從人與人對話裡.
見證到是不是一個溫暖的家庭.
最主要有兩個態度.
好像雅各在《雅各書》第一章的教導.
快快地聽,慢慢地說.
著名的文學家和哲學家Mortimer Adler.
他寫了一本非常非常流行的著作.
《如何講如何聽》.
各位,我不覺得我自己懂得講.
我又不覺得我自己懂得聽.
所以我讀這本書.
希望我能夠改進.
有一些主要的事.
特別大家留意.
要有一個有意義的對話.
使別人覺得這個教會是和諧的弟兄姊妹成立的教會.
我們對話的時候要有意義的對話.
我們要看看我們講話的時候的環境在哪裡.
講話的內容是什麼?.
有沒有意義的內容?.
講話的態度是怎樣?.
是快點講慢點聽,還是快點聽慢點講?.
講話的目的是什麼?.
是要罵人還是要向人請教?.
等等等等.
有用的鹽,第二個用途是防止腐爛.

$^{161}$防止食物腐爛.
有時沒有冰箱,要用鹽來保存食物.
耶穌教的門徒要先用鹽潔淨自己.
才能夠驅除信徒中間的毒.
保羅和雅各都講出那些毒是什麼?.
他們的喉嚨是敞開的墳墓.
舌頭弄鬼扎,嘴唇有毒蛇.
嘴唇有毒蛇的毒氣.
舌頭是喋喋不休的惡物.
充滿致命的毒素.
這些可以讓我們檢查,我們自己有沒有毒在我們裡面.
史唐寺更加給我們一個例子.
他說那時那些不順從的猶太人.
煽動外邦人使他們的心仇恨那些弟兄.
他們不但自己有毒,還煽動其他人.
雖然這些都是第一世紀雅各和保羅的教導,當時的信徒.
2026年的信徒也應該聽一下這些教導.
除去教會裡面的黑暗.
正如彼得教導信徒說.
不是君尊濟私有了被召出黑暗進入光明的國度的恩典和經驗.
你們就有經驗有責任去驅除黑暗.
保羅列出黑暗的事,不義,陰險,好鬥.
人家講一句,自己馬上講十句,要鬥贏人家.
欺詐,毀謗人,傲慢,自誇,narcissistic,這是英文.
幸災樂禍,滿懷嫉妒,好說柴賢,博弈父母,不守信用.
憎恨神,這個是非常黑暗,沒有人知道的,他自己憎恨神沒有人知道的.
凌辱人,不但自己去行,還喜歡別人去行.
這是羅馬書第一章二十八節到三十二節.
黑暗的事,人以為在黑暗中做的事沒有人可以知道.
但是神會知道,真光會驅除黑暗.
在希伯來書第四章,神的道是活的,是有效的.
比一切兩刃的劍更鋒利,甚至可以切入,剖開靈與魂,關節與骨髓.
並且能夠辨明心中的思想和意念.
黑暗的事,心中的思想和意念,神的道,神可以知道.
現在講明亮的光.
第一方面,光照的特質,真光照入人心,可以啟示出人生是有目的的.
以及人的生命是有價值的.
第二個方面是光照人前.
大給世人希望是我們教會的口號.
我們需要知道傳福音的要點.

$^{201}$然後才知道怎樣大希望給世人.
主耶穌對眾人說:我就是世界的光.
跟從我的必不在黑暗裡走,卻要得著生命的光.
你跟我已經得著生命的光.
明亮的光,我們要清楚我們已經得著生命的光.
我們已經行走在這個光裡面.
我們可以大給世人希望,一起在這個光裡面行走.
詩篇描寫,在黑暗裡面走的人.
以為沒有人知道他做的詭詐的事,人心詭詐壞到極處.
但是神知道,神要人更新,改變壞的心意.
當神悔改,走向光明,有一個清潔的心,正直的靈.
如果人離開世界的時候.
還不知道他是不是會進入這個永恆真神的光明的世界.
他去的地方就是黑暗的世界.
我們有責任將希望講給這些人聽.
耶穌指出,信徒是世上的光.
好像建在山上的城是無法隱藏的.
又好像人點了燈,不會放在量器底下.
而是放在燈台上,就照亮一家人.
我們的光應當照在人前.
所以要知道我們的生命是有目的和有價值的.
第一個光照的特質是帶著一個有目的的生命.
那個生命才有意義.
在主的真光照之下.
就好像我們聽到主耶穌對我們講.
我們的光應當照在人前.
第一個光照的特質就是愛人而己.
生命旋律這本書是描寫人生好像音樂的旋律.
有歡樂的笑聲,也有哀痛的哭泣.
正如所羅門王對人生的感嘆.
他說哭有時,笑有時,哀痛有時,跳舞有時.
但是仍舊要有目的.
牧長帶領我們信徒.
以找到主的大愛為目的.
就好像每天那些大鵝帶小鵝.
找到主給牠們吃的主人.
溫暖教會.
眾信徒相愛.
一直走到天家,就好像那些大鵝和小鵝.
漫步走到主人家.

$^{241}$我們很喜歡餵鵝.
所以每一年那些鵝生了嬰兒的時候.
就回來我們的正門敲門.
要找餵牠們的主人.
將來我們要回到天父那裡.
因為祂給我們的大愛.
另一個光照的特質是帶著神.
賜給價值給我們的生命.
這個特質就是榮耀尊貴.
並你所陳設的月亮星宿.
人算什麼你竟顧厭他.
世人算什麼你竟眷顧他.
你使他比神微小一點.
賜他尊貴為冠冕.
詩篇第八章第八篇是這樣說.
我們就好像聽到主耶穌對我們說.
神賜我們榮耀尊貴為冠冕.
既然詩人這樣寫出神做人都有價值.
就不要低估我們自己.
不要低估我們自己的光彩.
每一個人都有光彩.
好像企鵝一樣.
將光照在人間.
知道自己有價值的信徒.
他會努力奔跑前面的路.
一直到獲得榮耀尊貴的冠冕.
就是這本書的內容.
更加要彰顯出這個價值.
憑著熱情和堅定的信心.
在永恆的盼望.
快跑並且光照其他的跑道.
明亮的光的第二方面就是光照人前.
耶穌指出明亮光的作用.
人點燈不會放在陽氣底下.
而是放在燈台上就照亮一家人.
照樣你們的光也當照在人前.
讓他們看見你們的好行為.
又頌讚你們在天上的富.
我們可以將自己從黑暗走向光明這個經驗.
講給別人聽.

$^{281}$讓別人都可以見到照在人前的光.
藉著傳福音讓人可以認識主耶穌.
得到生命的光.
耶穌在世界的時候教導我們傳福音.
傳福音的心在生活上會彰顯出來.
聖靈必定會帶領我們晝夜讚美救主.
因為祂必賜給我們應許的祝福.
主怎樣教導我們傳福音呢?.
有好的行為讓人頌讚天上的富.
要彼此相愛讓人見到溫暖的家庭.
講述主的大愛讓世人追求主的救恩.
如果我們要有效地傳福音.
必須要在每一天的生活中時時實習這三個功課.
請各位弟兄姊妹時時記住.
你和我都是有用的鹽 明亮的光.
但願Fanny Crosby所作的歌唱出我們的心願.
請各位弟兄姊妹時時記住.
你和我都是有用的鹽 明亮的光.
但願Fanny Crosby所作的歌唱出我們的心願.
請各位弟兄姊妹時時記住.
你和我都是有用的鹽 明亮的光.
但願Fanny Crosby所作的歌唱出我們的心願.
請各位弟兄姊妹時時記住.
你和我都是有用的鹽 明亮的光.
請各位弟兄姊妹時時記住.
你和我都是有用的鹽 明亮的光.
但願Fanny Crosby所作的歌唱出我們的心願.
請各位弟兄姊妹時時記住.
你和我都是有用的鹽 明亮的光.
但願Fanny Crosby所作的歌唱出我們的心願.
請各位弟兄姊妹時時記住.
你和我都是有用的鹽 明亮的光.
但願Fanny Crosby所作的歌唱出我們的心願.
請各位弟兄姊妹時時記住.
你和我都是有用的鹽 明亮的光.
但願Fanny Crosby所作的歌唱出我們的心願.
\newpage



\section{以弗所書 6:1-4}
\label{sec:031753}
\textbf{在基督裡的親子關係}~~~~葉慧敏傳道
\newline
\newline
連結: \href{http://www.church.com.hk/acms/content.asp?site=cdc\&op=show\&type=product\&code=031753\&layout=sermon}{\texttt{031753}} ~~~~ 語音日期: 2026-03-01
\newline
\newline
\hyperref[sec:031752]{\small{< < < PREV SERMON < < <}}
~
\hyperref[sec:index_chronic]{\small{[返順時目]}}
~
\hyperref[sec:031754]{\small{> > > NEXT SERMON > > >}}
\newline
\newline
以弗所書 6:1-4
\newline
\begin{longtable}{cl}
\hline
\hline
章節 & 經文 (和合本修訂版)\\
\hline
6:1 & \begin{tabularx}{0.7\textwidth}{X} 作兒女的，你們要在主裡聽從父母，這是理所當然的。 \end{tabularx} \\ \\ \relax
6:2 & \begin{tabularx}{0.7\textwidth}{X} 當孝敬父母，使你得福，在世長壽。這是第一條帶應許的誡命。 \end{tabularx} \\ \\ \relax
6:3 & \begin{tabularx}{0.7\textwidth}{X} 見上節 \end{tabularx} \\ \\ \relax
6:4 & \begin{tabularx}{0.7\textwidth}{X} 作父親的，你們不要激怒兒女，但要照著主的教導和勸戒養育他們。 \end{tabularx} \\ \\
[1ex]
\hline
\hline
\end{longtable}
$^{1}$是的 鼎姐妹 午安.
剛才Warner提到今天是4月13日.
我都在新春期間祝賀大家.
新的一年靠主得令.
就好像進馬奔跑一樣 靠主得令.
好 鼎姐妹 我不知道對你來說.
你覺得今天的家庭裡面.
最常見的問題是什麼呢?.
會不會是同住在一間屋裡面.
但是父母和子女之間越來越沉默呢?.
又會不會是一家人坐在同一張飯桌裡面.
但是每個人都只是看著自己的手機?.
情緒化的反應取代了溫柔的對話呢?.
對我們一群在香港來到英國的弟兄姊妹來說.
我們的下一代很多都是接受英國的教育.
會不會我們都發現和我們的下一代出現了一個代溝.
令到父母和我們的孩子好像活在兩個世界一樣.
在這些張力之下 父母覺得教養很難.
兒女又覺得被誤解.
家庭覺得越來越疏離.
上次我和大家分享《已不鎖書》五章的時候.
我都分享過 保羅提到要看一個人的屬靈生命.
不是看他在教會做什麼.
而是看他在家裡是怎麼樣的.
上次就看到夫婦的關係.
今天就進入另一對的關係.
父母和子女的關係.
希望透過《已不鎖書》六章一至四節.
我們都一起去學習.
可能弟兄姊妹你會說 其實我沒有子女.
但我都很想你去留心聽今天的講道.
聽完你就會發現.
你在裡面都有一些使命要去做.
我們一起祈禱.
順著我求你去幫助我們每一位.
你讓透過今天的講道.
你幫助我們每一位弟兄姊妹體現.
你要我們怎麼做好.
無論是父母還是子女的角色.
主要我們更加看到在教會裡面.

$^{41}$我們怎麼去彼此的影響.
主求你在當中去幫助人們.
我們將會透過教堂獨斷名求 阿們.
我們一起讀一下這四節經文.
《已不鎖書》六章一至四節.
這次我們主要讀中文.
我們一起完成這四節經文.
1,2,3.
「你們作兒女的.
要在主裡聽從父母.
這是理所當然的.
要孝敬父母使你得福在世長壽.
這是第一條大英雄的誡命.
你們作父親的不要惹子女的氣.
只要照著主的教訓和警戒養育他們」.
今天的講道主要是分兒女篇和父母篇.
我們先看保羅對兒女的呼籤.
他說「你們作兒女的要在主裡聽從父母.
這是理所當然的.
要孝敬父母使你得福在世長壽.
這是第一條大英雄的誡命」.
短短三節經文.
是整段家庭的規範裡面第二組的關係.
接住了夫妻之後.
獨人和主人之前.
雖然篇幅比夫妻少.
但是色經同樣的重要.
因為保羅是將五章十五至二十一節.
在聖靈裡面我們要過一個智慧的生活.
具體地落實在家庭最基本的關係.
有幾個詞語很想和大家一起去探討.
保羅說到「作兒女」.
這個不單單是一個稱呼.
而是說到是一個關係.
是指到仍然依賴父母.
未完全獨立的人.
保羅正在針對那些仍然在父母權力之下.
仍然在成長學習階段的孩子.
仍然需要被教導,被引導的兒女.
但是更加重要的是.

$^{81}$保羅是向教會裡面的兒女說話.
這一點是很重要的.
初期的教會聚會裡面.
就不會像現在這樣.
兒童有兒童主日賀.
青少年有青少年聚會.
當時的信徒家庭是一起敬拜.
一起聽道.
所以保羅寫這篇信的時候.
是假設兒女是教會的一部份.
兒女亦都是聽道的人.
兒女是需要回應神話語的人.
一向橫教會都有一個觀念.
或者有一個做法.
就是兒童侍工不用特別請一個童工去做.
那通常是誰去做的 弟兄姊妹?.
通常就是請一個牧師.
然後就問他:你太太擅不擅長做兒童侍工?.
如果擅長的話.
其實那句話就是:我們想你買一送一.
這個觀念對我們今天很重要.
兒童不是供屬品.
是在神家裡真正的成員.
他們可能未完全明白信仰.
但是已經在教會裡面被牧養,被引導.
在過去的事奉日子裡面.
其中一個女孩子給我留下一個深刻的印象.
她是從媽媽的肚子裡已經回教會了.
我想問這裡是否有很多是從媽媽的肚子裡.
已經回教會的?.
可以給一些示意嗎?.
都沒有的 沒有理由的.
今天特別青少年和我們一起.
這個女孩子是從媽媽的肚子裡已經回教會了.
媽媽和爸爸都是很多活動的人.
星期天回教會是其中一項活動.
還有其它各方面的活動.
所以在這個小朋友裡面.
他的成長不是每個星期都回教會的.
但是去到四年級的時候.

$^{121}$有一天他媽媽和他說.
「我們今天不回教會,我們今天去哪裡玩?」.
這個小朋友就因為在教會學習了一些真理.
他那天就鼓起勇氣和媽媽說.
「不如你送我回教會」.
「然後你自己去玩」.
「你玩完你一點鐘再回來接我行不行?」.
面對著這個要求.
家長們你怎樣選擇?.
玩還是回教會?.
很難回答.
約好了.
最後這個媽媽就選擇和女兒一起回教會.
但是回來就馬上和我說.
「我女兒今天和我說了一些這樣的東西」.
其實下一句可能就是心裡想.
「是不是你教她的?」.
但是我覺得對這個女孩子來說.
不是人去影響她.
是神的話語去吸引她.
以致她有這個勇氣去和她父母說.
「我很想回教會認識耶穌」.
早兩天在哈勞基小組裡面.
有弟兄帶我們去看一本.
已故的楊木菊牧師的書.
書名叫《亂世風言:耶穌比喻的現代言》.
當天弟兄就帶我們去看浪子的比喻.
我在這裡不詳細說.
大家都知道浪子的比喻就是.
小孩子拿了父親的錢.
然後出去玩完.
沒有錢又回來.
楊牧師他在那個比喻裡面提到.
大兒子對他的弟弟回來的回應.
他表達哥哥是一個局外人.
他更加在書中提到.
其實教會都可能有很多局外人.
這些局外人是誰?.
就是弟兄姊妹你們回來認識主.
一路成長.

$^{161}$然後去到結婚,生孩子.
你帶著你的下一代回來.
這一幫第二代信徒.
他們從小就在教會.
但是他們是很容易成為教會的局外人.
我們從來不知道什麼時候.
我們的下一代回教會不再跟著父母回教會.
是因為他們自己很想回教會.
我們也不知道在什麼時候.
他們不再藉著父母的生命去認識上帝.
而是他們親自和主有一個關係.
我們和神的關係是父子的關係.
是不是啊? 弟兄姊妹.
我們的下一代和神的關係.
是不是都是父子的關係啊?.
楊牧師他就提到.
上帝只有子就沒有孫子.
今天 但是我們很多時候.
會不會都將我們的下一代.
視為只是回來的一個副產品呢?.
求主去幫我們.
究竟我們很想我們的下一代.
是因為我們而回教會.
還是他們都是一幫願意回應上帝恩典的人?.
第二個字是很想看「在主裡聽從」.
「聽從」是一個強烈的動詞.
意思是將自己放在權柄之下去聽.
並且按著所聽的去行.
在保羅的書信裡面.
這個字通常用來形容就是.
聽從基督,聽從福音,聽從使徒的教導.
所以保羅不是單純說家庭倫理.
而是說兒女對父母的信服.
是他們對基督信服的一部分.
「在主裡」不是在形容父母.
是在形容兒女的動機.
不是因為父母永遠是對的.
所以他們去聽從.
不是因為父母完美.
所以他們聽從.

$^{201}$而是因為兒女屬於主.
兒女知道自己屬於主.
所以他們聽從.
所以兒女去聽從.
其實這個回應是一種敬拜.
是一種門徒操練.
也是一種屬靈的姿態.
這裡和五章二十一節.
他說「全敬畏基督的心彼此信服」.
那個教導完全一致.
所以問題不是父母值不值得我們去聽.
而是我屬不屬於主.
不過弟兄姊妹在這裡.
也有一樣東西要說得很清楚的是.
聖經從來沒有要求兒女要信服罪.
信服暴力.
或者信服任何操控他們的人.
傷害他們的人.
在主裡面的聽從信服.
永遠不會違反主的心.
主的真理,主的愛.
如果父母的要求是要我們去違反聖經.
違反良心.
傷害自己或者其他人的話.
這些不是保羅所教導的聽從信服.
在這種情況裡面 神不會勉強你強行去支持.
而是要你在主裡面尋求幫助,保護和慈惠.
可能就是和教會牧者,道士.
成熟的弟兄姊妹去求助.
因為真正在主裡面的聽從.
是信靠神的心意.
而不是縱容任何可以摧毀生命.
扭曲神形象的權柄.
第三個字眼.
「這是理所當然」.
「理所當然」這個字.
可以譯作「公義的合乎」.
這句話其實有三重意思.
第一重意思是.
合乎一般的倫理.

$^{241}$在希臘和羅馬的世界裡面.
孝敬父母是維持社會秩序的基礎.
不聽從父母被視為敗壞的象徵.
我們可以想像一個小朋友去過馬路.
父母叫他停一停.
望去左又望向右.
小朋友完全可以自己衝.
但是大家都會知道.
聽父母的話是理所當然.
因為這樣是安全又合理.
第二重意思是合乎神的律法.
保羅隨即引用了十誡第五誡.
顯示了孝敬父母不僅是一個文化的要求.
也是神的命令.
即是認同父母不僅是一個社會的常識.
而是遵行神的命令.
第三重意思是合乎神創造家庭秩序.
父母是神所設立的權柄.
責責教導,保護,引導兒女.
而兒女也在秩序裡面學習信服,學習信任.
學習領受信命.
就像我剛才說小朋友過馬路的例子一樣.
信服父母不僅是服從一個人的權威.
而是在神所設計的秩序裡面.
學習信靠和被引導.
所以理所當然這個字不是說父母永遠是完美.
而是說神設立的秩序.
這個秩序包括社會倫理,神的律法,家庭創造的秩序.
當我們願意聽從,信服父母的時候.
我們就在進行社會上的行為.
當我們願意聽從父母的時候.
我們就是在敬畏神,進行神的名利.
我們就是活在神創造的家庭秩序裡面.
被保護,被培育,被引導.
換句話說.
兒女的信服就是指導日常生活裡面.
活出怎樣去信靠神.
但是這些是不是真的這麼合理呢?.
為什麼神要將祝福和這份信服放在一起呢?.
第二至第三節是在說孝敬父母.

$^{281}$這是第一條帶應他的誡明.
保羅引用了《出埃及》的20章12節.
這樣說:「要孝敬父母使你得福在世長壽」.
並且他補充這是第一條帶應他的誡明.
孝敬在舊約的意思是信服,尊重,照顧父母的需要.
孝敬父母是一種贖靈生命的表現.
反映了一個人對神的敬畏.
為什麼說是第一條帶應他的誡明呢?.
在十誡裡面只有第五誡清楚地附帶這個應許.
其他誡明雖然都會有祝福.
但是第五誡特別強調.
因為家庭秩序的重要性.
尊重生命源頭的重要性.
這條誡明對整個社會的影響.
而保羅如何重新應用這個應許呢?.
在《出埃及記》裡面他提到.
「使你在耶和華所賜你的地上得享長久日子」.
但是保羅就將應許普及化.
說成「使你得福在世長壽」.
他沒有再提「那地」這個字.
因為在新約的神的指紋裡面.
不再被局限在一塊地圖當中.
他將應許變成了普世性的原則.
這個應許也不是贖靈化的長壽.
有一些猶太文獻會將長壽贖靈化.
但是保羅沒有這樣做.
他所說的是一般貴文學的原則.
他說的是「安神自慰而行」的人.
人生更加有秩序.
更加無保護.
更加得人尊重.
不是保證沒有苦難.
而是尊重生命源頭的人.
生命會更加整全.
在舊約路得記記載.
路德信從了他的婆婆.
撈米回到耶路撒冷.
雖然前途很不明朗.
但是他的信服帶來了神的祝福.
最終成為了耶穌加譜的一部分.

$^{321}$這裡提醒我們.
聽從父母或者屬靈長輩.
是信靠神的自慰.
孝敬父母是信服,尊重,照顧,敬畏神的表現.
信服帶來了祝福.
這個祝福不單只是影響自己.
也可能影響整個屬靈的歷史.
就好像路德一樣.
當保羅說完兒女的責任之後.
他就將焦點轉向父母.
他這樣說.
「你們作父親的,不要惹他的氣.
只要照著主的教訓和警戒養育他們」.
在當時的世界裡面.
這句話是極具顛覆性.
他是說「你們作父親」.
為甚麼特別說到「對父親說話」呢?.
在第一,第二節明明是用「父母」這個字.
但在第六節,他清楚指出「你們作父親」.
原因最少有三個.
第一個就是.
在當時的社會,父親是握有極大的權力.
在猶太和希羅文化裡面.
父親掌握著家庭裡面的絕對權柄.
在羅馬的法律裡面.
父親幾乎擁有完全的決定權.
決定他們小朋友的婚姻,教育和未來.
有權去決定初生的嬰兒是否繼續活在這個世界上.
也可以決定是否將他們的兒女埋葬成為奴隸.
正因為權力集中在父親手中.
他們最需要被福音去重塑.
但是今天的環境.
會不會是姐妹,媽媽的權力都很大呢?.
這個我們可以都消失.
第二個原因為何他特別提到「作父親」.
是本來要重新定義父親的角色.
福音並非為父權背書.
而是為家庭帶來了一個新的秩序.
父親不再是權力的象徵.
而是屬靈的牧者.

$^{361}$以基督為榜樣養育兒女.
保羅呼召父親由掌控者轉為牧養者.
由權力的中心轉為愛的中心.
第三個原因是父親對孩子的影響極深.
父親的態度,言語,反應,行為.
會深刻地影響下一代.
怎樣去理解權柄,怎樣去看自己的價值.
怎樣去想像天父的愛.
所以保羅特別對父親說話.
不是因為父親更加重要.
而是因為他們的影響力更大,責任更大.
更需要被基督的愛去根深.
在兒童主教學的服侍裡.
我們回到教會很多年的弟兄姊妹.
可能都會有一個印象.
通常是姊妹做主教學.
還是弟兄做主教學.
通常是姊妹做定型的.
兒童主教學通常是全女班.
但是我記得有一年我在香港服侍的時候.
有一屆的年度一級的小朋友.
全部都是男孩子 十個.
頭一兩年還可以找一些姨姨去照顧他們.
因為還很可愛的 一,二年班.
去到三年班已經DEE跟我說 搞他們不行.
去到五年班的時候.
我就發出了一個邀請給兩個弟兄.
我問他們 你們願不願意來承擔這班小朋友的生命.
他們兩個就說 好啦 我們輪流來.
他們來跟小朋友玩 他們來給小朋友玩.
他們來教聖經.
但是更加重要的是.
他們兩個的生命就是在影響小朋友.
小朋友就是從他們裡面看見.
一個弟兄應該怎麼去做.
他們除了在自己的父親裡面去學習.
也可以在教會的叔叔,哥哥裡面去負責.
去學習怎麼去做一個弟兄.
所以兒童主義學我們都很需要弟兄加入的.
大家如果弟兄有感動都可以站下圖.

$^{401}$保羅在這裡不是要鞏固父親的權力.
而是要將焦點可以從我可以做些什麼.
轉到向我應該成為怎樣的人.
他不是叫父親要維持權威.
而是要效法機動.
不是叫父親要掌控家庭.
而是叫父親去牧養家庭.
說完了為什麼保羅特別提到你們作父親之後.
保羅就提了兩個命令.
一個消極的命令和一個積極的命令.
消極的命令他說不要惹兒女的氣.
惹兒女的氣不單止是指令到你的小朋友不開心.
是指用不合理,不按情理,不用愛的方式去對待孩子.
令到孩子心裡面激怒,被壓碎.
惹兒女的氣這個字在舊約七十事譯本裡面.
多數用來形容耶和華.
以色列人激怒耶和華的意思.
就是做了一些違反關係,違反愛,違反約的行為.
而在《歌羅西書》三章21節保羅用了另一個字.
在中文聖經都是講到惹兒女的氣.
但是他這個字有更豐富的意思就是.
令到小朋友心裡面感覺到辛酸.
令到小朋友失去了盼望.
令到小朋友覺得自己永遠都做得不夠好.
這裡是呼應《爾弗所書》四章26至27節講怒氣的危險.
怒氣如果惹怒了我們的孩子處理.
撒旦就會利用這個機會去破壞父母和子女之間的關係.
所以保羅是叫父親要避免一切會摧毀我們下一代心靈的行為.
包括過度嚴厲,不合理的要求.
濫用權柄,不公平,責備,囉嗦,羞辱.
不視小朋友的需要,用情緒控制小朋友.
用比較摧毀小朋友的自信.
保羅的意思是兒女雖然要順服父母.
但是兒女並不是父母的副產品.
兒女都是要有尊嚴,有權利,有神的形象.
父親不能夠操控,剝削或者壓榨他們.
在現在的社會裡面父母面對壓力都很大.
工作的目錄在香港的時候相信那個工作的壓力真的很大.
來到這裡相對鬆一點.
但是父母會面對另一種的壓力就是.

$^{441}$沒有工人,家庭支援少了,很多事情要做.
無論是媽媽,爸爸都要幫忙做家務.
弟兄你們有沒有幫忙做家務的?.
沒有反應.
有的,有姐妹代弟兄點頭.
情緒會覺得很疲累.
教養的資訊過多.
其實這些都是一個壓力.
但是保羅提醒我們壓力不能夠成為我們傷害.
我們孩子的理由.
我們的小朋友不是父母的情緒出口.
不是父母未完成夢想的現身.
也都不是權力的試驗場.
保羅接著講一個積極的命令.
他不單單不是叫父親不要做些什麼.
他還叫父親要做些什麼.
他這樣說.
「只要照著主的教訓和警戒養育他們」.
「養育」這個字是《義佛所書》五章二十九節.
用來形容基督保養護膝教.
意思是溫柔,細心,持重,以愛為動機.
以對方的成長為目標.
在兩若之間 這個字更加有培育到達應有目標的意思.
但是保羅提醒我們.
女兒不是王子,不是天才.
不是父母未完成夢想的容器.
養育不是將孩子塑造成父母想要的樣子.
而是幫助他們成為神創造他們.
要成為的人.
「教訓」這個字可以指為一個全人教育.
品格塑造,正面教導,生活榜樣.
教他們應有的價值觀.
父親的生命就是小朋友的教材.
「警戒」這個字是口頭勸喻.
溫柔糾正,設立界線.
防止小朋友走向危險.
不是情緒的爆發 而是愛裡面的提醒.
主的教訓和警戒其實都是基督為中心的教育.
其實原文是用「照著主的」這個字.
意思其實應該是以主的教訓和警戒養育他們.

$^{481}$意思就是教導的內容是來自主.
教導的方式都是來自主.
教導的目標是帶小朋友去認識主.
教導的動力是因為主的.
父親不是單靠自己.
在家庭裡面代表基督,反映基督,指向基督.
在我服侍的過程裡面.
曾經遇過好幾個小朋友.
回到聚會的時候有情緒.
但是那個情緒未必是在聚會裡面發生.
是他突然間覺得很傷心在哭.
有幾個男孩我都很深印象.
第一個男孩他有一天哭.
哭的原因是 後來我和他私底下再談.
哭的原因是因為他媽媽對他有一些要求.
而那個要求是要去養弟弟.
接著他就問了我一個問題.
待會你們都可以想一下這個問題怎麼去解答.
他說:其實我永遠都比他大.
我今天養他,我明天又養他,後天又養他.
那究竟我何時不用養他?.
我永遠都是他哥哥.
其實他那一刻說出來的時候.
現在說給你們聽你們覺得很好笑.
但是其實這個就是做大哥.
經常被要求你是大哥.
你養一下弟弟那種的壓力.
那誰養我?.
通常媽媽就說:在你弟弟和妹妹還沒出生之前.
全世界都在養著你.
但是那個人忘記了.
他帶著的是幾十年要養他的弟弟.
第二個都是男孩.
他和媽媽,爸爸相處有很多阻力.
他很想打遊戲.
但是他媽媽,爸爸對他很多要求.
很嚴格的.
考試成績很好.
又學很多十方全能.
游泳,吹笛,養療癥.

$^{521}$他很想打遊戲.
因為打遊戲這個事件就發生.
掛無線網絡,收起電話那些.
他就對著我哭.
他覺得其實是不是真的不可以做?.
其實都很困難.
一直為他們祈禱,一直為他們祈禱.
去到他初中的時候.
應該是中三.
有一天我看他的IG.
我就發現為什麼有個女孩跟你這麼親密的呢?.
我就問他:這是你女朋友嗎?.
他回答我:是的.
接著我就問他:你媽媽知道沒有?.
你爸爸知道沒有?.
接著我就說:你的IG.
其實有教會的叔伯兄弟,兄弟姐妹在那裡.
那些姨姨會很快告訴你媽媽,爸爸知道的.
你最好自己快點報案.
哎呀!我不知道怎麼說.
很難說.
不知道,他們罵我什麼的.
一個初中的小孩當然是怕這些.
第二天回教會.
他媽媽就跟我說.
喂!我想跟你說些什麼?.
我懷疑我兒子拍拖.
那我回答:知道還是不知道?.
我就說:你不如跟他好好談一談.
了解一下發生什麼事.
不要罵他.
你兒子又帥又十項全能.
有女孩喜歡很正常的.
我跟他媽媽這樣說.
接著他媽媽就說.
都是呀!我兒子這麼帥.
我就叫他忍著,不要罵他.
接著他媽媽星期天跟我說.
她就回家約個兒子.
我們星期六吃早餐.

$^{561}$怕不怕?如果你是那個初中生.
怕得要死.
但是星期六.
他就知道其實我不說就沒機會說了.
我現在就要和盤托出說給你聽.
是呀!媽其實我正在拍拖.
接著那個媽媽心裡面.
就是忍著裡面很多事情想說.
她只是說話.
多謝你告訴我.
很欣賞你告訴我.
不過你要留意123這樣.
接著很快就分手了.
分手了.
分手了她就主動告訴媽媽.
這個是她初中的經歷.
接著這個小朋友一直在教會成長.
去到她高中的時候.
有弟兄姊妹就和我說.
可不可以搞一些親子的講座.
不是找父母來講.
而是找走過這段路的年輕人來講.
我就想起這個男孩子.
我就問他可不可以來分享.
分分享給他.
他答應我.
其實我都沒有叫他要講回這件事.
我就任他講 你想講什麼就講什麼.
接著他真的在講座裡面.
對著眾家長講自己當日拍拖.
接著怎樣給媽媽.
要和媽媽交代.
接著在交代裡面.
他分享到我當初都很害怕.
但是我覺得媽媽對我的回應.
是我覺得很被信任.
這個是他和父母破冰.
破冰的那一刻去形容.
自此之後他每個星期日.
他都為他全家去預備豐富的晚餐.

$^{601}$接著以前就是被人mung Wi-Fi線.
出去煉水的時候.
就去KFC坐多一個小時.
因為KFC只能上網一個小時.
就打多一個小時機才回家.
但是之後 當破冰之後.
他會和媽媽講.
那時還在讀中學.
他和媽媽講.
我現在游完水和同學去KFC坐一會.
我今晚回教會.
我今天怎樣怎樣怎樣.
好像和女朋友報告行蹤.
弟兄姊妹對我來說.
我覺得很感恩.
因為媽媽和爸爸原本又很嚴厲.
很高要求.
到他們願意按著神的方式去愛他的孩子.
去容忍,去等候他的孩子.
現在這個小朋友已經大學生.
求主去幫我們.
我們都很需要幫助我們的下一代同樣這樣經歷.
保羅的教導其實和當時的文化是完全相反的.
當時的文化剛才說.
父親是一個權力的中心.
但是保羅提醒父親是一個愛的中心.
保羅當時的文化講到父親是可以懲罰的.
保羅的教導是父親你要默養.
保羅當時的教導是父親可以要求你去順服.
但是保羅的教導是父親要以基督的方式去引導.
保羅要基督徒的父親去記得.
你們做父親的職份是從我們天上的父而來.
神怎樣以恩典,忍耐,真理,慈愛對他的孩子.
你們都要這樣去對待你的孩子.
弟兄姊妹當我們聽完神對兒女,對父母.
對我們作為教會的膚淺.
我們作為教會的一群.
我們的孩子都是在看著我們.
我們有愛他們就去學習.
我們的教會裡面是有衝突他們完全收到.

$^{641}$我們的教會裡面是願意分享他們的道德.
我們都知道就是靠自己我們沒有辦法做到.
但是福音告訴我們我們不孤單.
今日這段經文有三個很重要的真理.
如果你忘記剛才說什麼不要緊.
你記得以下三個真理.
兒女的信徒是敬拜主的一部分.
父母的養育是一個屬靈的使命.
不是私人的情緒.
教會是一個屬靈的家庭.
是讓下一代學習基督的地方.
這段經文不是要我們成為完美的父母,完美的兒女.
而是要我們回到那位完美的天父面前.
因為我們都有很多缺失.
我們都不是完全的父母.
但是我們有一位完全的父親.
他願意以恩典,忍耐,真理,慈愛去養育我們.
他從來沒有按我們的過犯去待我們.
他沒有因為我們的軟弱去放棄我們.
他沒有因為我們的失敗去羞辱我們.
他反而以他的兒子耶穌基督的愛.
重新塑造我們,更新我們,接納我們.
今天神不是要我們死守律法.
而是回到祂的身邊.
願我們作為兒女的在主裡面學習,聽從.
因為我們有一位值得信靠的天父.
願我們作為父母的以主的教訓和警戒.
去養育我們的兒女.
因為我們自己就是這樣被主去養育.
願我們整間的教會,我們這個社會.
成為一個能夠讓下一代看見基督.
經歷基督,學習基督的福利家庭.
由主去幫助我們.
在主裡面重新得力.
今晚回到家的時候.
不妨試一下用一句溫柔的語氣.
和你的家人去說一句話.
或者用一分鐘有耐性的去聽他們說話.
讓福音從今天開始慢慢改變我們的家庭.
求主去幫助我們一起祈禱.

$^{681}$天父我們求你在裡面去幫助我們.
我們都是一幫民族人.
我們都是一幫罪人.
我們都是得罪你的人.
但我們多謝你.
你願意這樣去接納我們.
你願意不放棄我們.
以致我們能夠經驗你很深的愛,很深的恩典.
但是求你去幫助.
幫助到作為兒女的.
我們真的去學習怎樣去聆聽父母.
這份聆聽是因著耶穌的恩典.
作為父母的.
我們願意真的去幫助我們的孩子.
有著你的真理,你的教導.
去幫助我們的孩子.
去成為合理心意的人.
我們來到英國.
我們追求的不再是很多特徵的事,特徵的成績.
我們追求的真的很需要.
你讓我們看見怎樣去無阻我們下一代的生命.
你幫助我們的教會.
都能夠讓我們的下一代看見更多的恩典.
你更多的圖在怎樣和這間教會一起走過.
主求你就去看顧人們.
我們多謝你.
我們將會報告教堂和主名卡.
\newpage



\section{以賽亞書 44:1-23}
\label{sec:031754}
\textbf{告別虛空}~~~~周弘毅牧師
\newline
\newline
連結: \href{http://www.church.com.hk/acms/content.asp?site=cdc\&op=show\&type=product\&code=031754\&layout=sermon}{\texttt{031754}} ~~~~ 語音日期: 2026-02-28
\newline
\newline
\hyperref[sec:031753]{\small{< < < PREV SERMON < < <}}
~
\hyperref[sec:index_chronic]{\small{[返順時目]}}
~
\hyperref[sec:031755]{\small{> > > NEXT SERMON > > >}}
\newline
\newline
以賽亞書 44:1-23
\newline
\begin{longtable}{cl}
\hline
\hline
章節 & 經文 (和合本修訂版)\\
\hline
44:1 & \begin{tabularx}{0.7\textwidth}{X} 「我的僕人雅各，我所揀選的以色列啊，現在你當聽。 \end{tabularx} \\ \\ \relax
44:2 & \begin{tabularx}{0.7\textwidth}{X} 那位造你，使你在母腹中成形，並要幫助你的耶和華如此說：我的僕人雅各，我所揀選的耶書崙哪，不要害怕！ \end{tabularx} \\ \\ \relax
44:3 & \begin{tabularx}{0.7\textwidth}{X} 因為我要把水澆灌乾渴的地方，使水湧流在乾旱之地。我要將我的靈澆灌你的後裔，使我的福臨到你的子孫。 \end{tabularx} \\ \\ \relax
44:4 & \begin{tabularx}{0.7\textwidth}{X} 他們要在草叢中生長，如溪水旁的柳樹。 \end{tabularx} \\ \\ \relax
44:5 & \begin{tabularx}{0.7\textwidth}{X} 這個要說：『我是屬耶和華的』，那個要以雅各的名自稱，又有一個在手上寫著：『歸耶和華』，並自稱為以色列。」 \end{tabularx} \\ \\ \relax
44:6 & \begin{tabularx}{0.7\textwidth}{X} 耶和華－以色列的君王，以色列的救贖主－萬軍之耶和華如此說：「我是首先的，也是末後的；除我以外再沒有神。 \end{tabularx} \\ \\ \relax
44:7 & \begin{tabularx}{0.7\textwidth}{X} 自從古時我設立了人，誰能像我宣告，指明，又為自己陳說呢？讓他指明未來的事和必成的事吧！ \end{tabularx} \\ \\ \relax
44:8 & \begin{tabularx}{0.7\textwidth}{X} 你們不要恐懼，也不要害怕。我豈不是從上古就告訴並指示你們了嗎？你們是我的見證人！除我以外，豈有神呢？誠然沒有磐石，就我所知，一個也沒有！」 \end{tabularx} \\ \\ \relax
44:9 & \begin{tabularx}{0.7\textwidth}{X} 製造偶像的人盡都虛空，他們所喜悅的全無益處；偶像的見證人毫無所見，毫無所知，以致他們羞愧。 \end{tabularx} \\ \\ \relax
44:10 & \begin{tabularx}{0.7\textwidth}{X} 誰製造神像，鑄造偶像？這些都是無益的。 \end{tabularx} \\ \\ \relax
44:11 & \begin{tabularx}{0.7\textwidth}{X} 看哪，他的同夥都必羞愧。工匠不過是人，任他們聚集，任他們站立吧！他們都必懼怕，一同羞愧。 \end{tabularx} \\ \\ \relax
44:12 & \begin{tabularx}{0.7\textwidth}{X} 鐵匠用工具在火炭上工作，用鎚打出形狀，用他有力的膀臂來錘。他因飢餓而無力氣；因未喝水而疲倦。 \end{tabularx} \\ \\ \relax
44:13 & \begin{tabularx}{0.7\textwidth}{X} 木匠拉線，用筆劃出樣子，用鉋子鉋成形狀，又用圓規劃了模樣。他仿照人的體態，做出美妙的人形，放在廟裡。 \end{tabularx} \\ \\ \relax
44:14 & \begin{tabularx}{0.7\textwidth}{X} 他砍伐香柏樹，又取杉樹和橡樹，在樹林中讓它茁壯；或栽種松樹，得雨水滋潤長大。 \end{tabularx} \\ \\ \relax
44:15 & \begin{tabularx}{0.7\textwidth}{X} 這樹，人可用以生火；他拿一些來取暖，又搧火烤餅，而且做神像供跪拜，做雕刻的偶像向它叩拜。 \end{tabularx} \\ \\ \relax
44:16 & \begin{tabularx}{0.7\textwidth}{X} 他將一半的木頭燒在火中，用它烤肉來吃；吃飽了，就自己取暖說：「啊哈，我暖和了，我看到火了！」 \end{tabularx} \\ \\ \relax
44:17 & \begin{tabularx}{0.7\textwidth}{X} 然後又用剩下的一半做了一個神明，就是雕刻的偶像，向這偶像俯伏叩拜，向它禱告說：「求你拯救我，因你是我的神明。」 \end{tabularx} \\ \\ \relax
44:18 & \begin{tabularx}{0.7\textwidth}{X} 他們既無知，又不思想；因為耶和華蒙蔽他們的眼，使他們看不見，塞住他們的心，使他們不明白。 \end{tabularx} \\ \\ \relax
44:19 & \begin{tabularx}{0.7\textwidth}{X} 沒有一個心裡醒悟，有知識，有聰明，能說：「我曾拿一部分用火燃燒，在炭火上烤餅，也烤肉來吃。這剩下的，我豈要做可憎之像嗎？我豈可向木頭叩拜呢？」 \end{tabularx} \\ \\ \relax
44:20 & \begin{tabularx}{0.7\textwidth}{X} 他以灰塵為食，心裡迷糊，以致偏邪，不能自救，也不能說：「我右手中豈不是有虛謊嗎？」 \end{tabularx} \\ \\ \relax
44:21 & \begin{tabularx}{0.7\textwidth}{X} 雅各啊，要思念這些事；以色列啊，你是我的僕人。我造了你，你是我的僕人，以色列啊，我必不忘記你。 \end{tabularx} \\ \\ \relax
44:22 & \begin{tabularx}{0.7\textwidth}{X} 我塗去你的過犯，像厚雲消散；塗去你的罪惡，如薄霧消失。你當歸向我，因我救贖了你。 \end{tabularx} \\ \\ \relax
44:23 & \begin{tabularx}{0.7\textwidth}{X} 諸天哪，應當歌唱，因為耶和華成就這事。地的深處啊，應當歡呼；眾山哪，要出聲歌唱；樹林和其中所有的樹木啊，你們都當歌唱！因為耶和華救贖了雅各，並要因以色列榮耀自己。 \end{tabularx} \\ \\
[1ex]
\hline
\hline
\end{longtable}
$^{1}$鼎姊妹,午安,好似呢,過完年啦,今日係農曆年幾多,有冇人講得出呀?.
十二呀,咁都係正值正月啦,都恭祝大家馬年進步,身體健康.
如果大家每一年,係年初ge 時候有留意新聞呢.
年初,當你打開電視,你總係會見到有一位藝人,九十幾歲,年年都俾人訪問佢.
點解呢?佢就會扮成唔同形象去黃大仙,即係裝頭炷香,你一定見到佢.
如果你有留意下新聞ge 鏡頭,唔止年紀大,其實都有相當多年輕ge 面孔.
爭住去做同一件事,今年我就見到新聞訪問一位00後.
00後ge 年輕人,佢就話早一晚三點幾就已經走去排隊,為著要去裝頭炷香.
究竟我地ge 信仰,基督信仰,同信黃大仙有D 咩分別呢?.
德國神學家馬丁路德,佢就曾經咁樣講,佢個說話,英文係咁樣翻譯.
Religion that does nothing, gives nothing, costs nothing, suffers nothing, is worth nothing.
翻譯出黎就係話,一個咩都唔使做,乜ye 都唔使付出,唔使成本,唔需要承受痛苦ge 宗教信仰.
係無價值,佢咁樣講,我今日要同大家分享以賽亞書44章.
前文43章就講到以色列民要面對逆境,但係神ge 子民仍然將會蒙救贖.
所以就叫佢地唔使驚,因為神一定會保護佢ge 子民.
問題就係呢班上帝ge 子民究竟有冇睇得清楚佢地所信ge 耶和華上帝究竟係點樣ge 呢?.
我地先睇一睇頭五節經文,如果你有機會查經翻查下其他中文ge 譯本.
呂鎮忠譯本同詩高譯本都分別將第一節「我的僕人雅各」呢句翻譯ge 前面就多jor 一個「但」或者「但是」呢個連接詞.
意思即係點呢?其實華本無譯到出黎,但係如果我地有呢個連接詞ge 話.
我地就知道44章第一節同前面43章其實係有關連緊接住43章.
寫開古卷話比我地聽原來44章一節係緊接連於43章28節.
雖然神係43章28節咁責備佢ge 子民,所以我要辱沒聖所ge 首領使雅各成為咒詛,使以色列成為辱罵.
不過鏡頭一轉去到44章第一節,上帝仍然堅持佢gwo 個不變ge 應許.
佢話唔會因為百姓ge 活逆,然後我將你地整個猶大都毀滅jor 佢.
佢仍然要向佢地施恩,所以神係44章第一節上半節就講.
我地撥人雅各,我所揀選的以色列.
睇唔睇到呢度佢用緊重疊ge 方式一次又一次叫佢地個名.
跟住去到第一節下半節同埋第二節接住就咁話.
現在你當庭,造作你,又從你出胎,造就你,並要幫助你的耶和華.
如此說,我地撥人雅各,我所揀選的耶穌倫,不要害怕.
耶穌倫係上帝對以色列ge 匿稱.
如果我地睇返七十事譯本ge 翻譯,譯佢做親愛的以色列.
神係gwo 位造作並且造就以色列之神.
咁你問,造作同造就有乜分別?.
新譯本就將呢兩個字分別譯做創造同埋塑造.
就好似一對好愛錫自己孩子ge 爸爸媽媽.
佢唔單止將小朋友帶到黎世間.
佢仍然唔會話生仔就唔養仔.
佢一定會承擔起教養佢地成長ge 責任.
縱使神係要嚴懲呢D ge 以色列民.

$^{41}$但係唔代表要佢地萬劫不復,死緊.
因為佢地係被我,即係神所揀選.
係第一節同埋第二節重覆jor 兩次.
去到第三節到第四節,仲為佢地勾出jor 一個相當靚ge 願景.
因為我要將水囂冠口渴ge 人,將河囂冠乾旱之地,我要將我ge 靈囂冠你ge 後裔,將我ge 福囂冠你ge 子孫.
他們要發生在草中,像溪水旁ge 柳樹.
不停咁囂冠囂冠囂冠.
上帝希望人睇見佢ge 作為,不停咁囂冠囂冠再囂冠.
聖靈ge 囂冠,人ge 生命將要被改變.
聖靈ge 降臨,我相信大家好熟識.
係《哲行傳》第一章第八節,gwo 度ge 技術就係咁話.
但聖靈降臨係你地身上,你地就必得著能力,並要在耶路撒冷,猶太傳地和撒瑪利亞,直到地極,作我ge 見證.
但係當我地繼續睇落去,其實你望清楚《哲行傳》係整個聖靈ge 工作上面係點樣技術ge .
當然有正面ge ,正面就係話會幫助呢一班ge 使徒黎到行神即其事,去傳福音,係咪呀?.
不過請你唔好忘記呢度所寫ge 第五章第三節,咩黎架?.
阿拿尼亞夫婦因為欺控聖靈,結果係即時撲倒死亡.
去到第十六章,記載保羅最初係想去亞西亞港島ge .
不過聖靈禁止佢話唔准去,之後冇幾耐佢又想去呢個ge 被推離傳道ge 時候呢.
點知耶穌ge 靈卻不許,所以呢唔單止係服事,其實係我地ge 生命都係咁樣架.
唔係一帆風順架,各位弟兄姊妹,可以好多難處架.
不過神係呢個難處ge 當中,其實佢唔係想一掌推倒我地.
呢個從來都唔係上帝ge 原意,而係佢好想我地能夠係呢D 困難裡面.
仍然睇到究竟神係做緊D 乜ye ,佢想教我地乜ye 功課.
以至我地克服jor 呢D 障礙之後,我地可以辨識得到上帝ge 心意,然後跟住去行.
去到第五節,跟住就話,這個要說,我是屬耶和華的,那個要以雅各的名自稱.
又一個要親手寫歸耶和華的,或亦在手上寫歸耶和華,並自稱為以色列.
參考呢度寫比你ge 異賽亞書49章第16節,講到當其時ge 主人對於屬於佢地ge 奴隸.
普遍ge 做法,就會係奴隸ge 手上寫上,或者作一個印,作為一個標記.
不過呢度你睇清楚,寫標記ge 人唔係主人.
如果我地參考埋呂鎮忠譯本,就會清楚一D .
講到係咁樣講,將歸永恆主寫在自己的手上,並將以色列者名字作為自己ge 頭銜.
寫呢個名唔係人地逼你,係自動自覺將歸耶和華的標記寫係自己ge 手上.
神ge 心意唔單止係要叫以色列人得福,仲有見唔見到呢個哪個又一個.
即係話一個又一個再一個ge 人,好多人,佢地最終都會自願黎到歸於神.
歸於係以色列信仰ge 外邦人呢,呢度話比我地聽.
最終都會被納入整個神所祝福ge 群體當中.
所以你見到呢,我係呢一個ge 小段落gwo 度呢,係用jor 一個開引號同埋刪引號.
框住jor 以色列呢三個字.
跟住去到第六節ge 上半節呢,係咁樣話ge .
耶和華以色列的君,以色列的救贖主,萬君之耶和華,如此說.

$^{81}$以賽先知就用jor 幾個尊稱黎到去稱呼神.
藍色ge 字,以色列的君,以色列的救贖主,萬君之耶和華.
一方面佢係要強調神同埋以色列民gwo 個好特別ge 關係.
另外仲有話佢希望表明呢一個ge 神其實係滿有能力同埋佢必定能夠成就佢ge 心意.
好啦,佢跟住呢就同讀者講出jor 呢.
神三個好重要ge 特質,第六節ge 下半節.
「我是首先的,我是末後的,除我以外再沒有真神」.
「我是首先的,我是末後的」其實係啟示錄你會見到.
啟示錄ge 開頭一章八節,結束ge 時候二十二章十三節.
其實都有相類似ge 敘述係咪呀?一章八節.
主神說:「我是阿拉伐,我是俄梅加,是色在今在以後永在的傳能者」.
去到二十二章「我是阿拉伐,我是俄梅加,我是首先的,我是末後的,我是初,我是終」.
阿拉伐呢,呢度話比我地聽,係希臘文ge 第一個字母.
作為一個開端,其實就意味住佢係萬物之前,萬物ge 存在都係源於佢.
佢係地上面一切ge 勢力係黎自佢.
佢唔使靠任何人ge 努力.
就正如約翰福音一章第三節gwo 度講「萬物是藉著他做的,凡被做的沒有一樣不是藉著他做的」.
仲有「我是末後的」意思就係任何ge 力量都無辦法攔阻佢ge 旨意被成就.
所以無人能夠凌駕於佢之上,甚至與佢平排.
佢係至高無上ge ,即係佢係老闆呀.
一個老闆做任何ye ,唔使問你架,唔使問我架.
馬丁路德呢,比jor 我地一個好好ge 提醒.
佢話呢「我地對神ge 睇法太過人性化」.
你諗下呢句「我地對神ge 睇法太過人性化」.
其實聖經係詩篇五十篇二十一節都曾經比過我地相類似ge 責備.
佢話咩你,即係人呀.
行了這些事,我,即係講緊神啦.
患閉口不言,你仲想我同你一樣.
呢個正正就係人ge 想法.
我地好多時好似將望遠鏡調轉jor 黎睇.
將明明好有能力ge 上帝扁到好細.
你用個望遠鏡,加塊鏡片,睇下個太陽.
幫我望下望下,變成一支好微弱ge 蠟燭.
呢個就係我地黎架啦.
所以好值得我地反思ge 就係.
呢D ye 大家聽過曬,你都識講架.
但問題就係,神係我地心裡面ge 地位.
究竟係咪同佢係宇宙裡面ge 地位一樣咁高.
你諗o下,我再講多次.
神係我地心裡面ge 地位,究竟係咪同佢係宇宙ge 地位一樣咁高.

$^{121}$你可以好容易答,梗係啦.
聖經講架嘛.
不過當有事情發生,有D 切身ge 問題出現jor .
擺係我地面前.
上帝ge 地位就開始嗡嗡up up .
搖下搖下,我地自我的想法.
就會取代jor 至高無上ge 上帝ge 地位.
好期盼,今日呢段聖經能夠提醒我地.
重新認定,呢個先至係至高創造天地ge 主.
佢係gwo 位掌管一切ge 神.
跟住去到第七節.
呢度係好似一個法庭ge 辯論咁樣.
控方就咁樣話.
自從我設立古時ge 文,誰能將我宣告並且指明.
又為自己陳述,讓他將未來ge 事和必成ge 事說明.
正因為冇人係神ge 對手,冇得傾.
所以去到第八節,佢就繼續鼓勵呢班以色列民.
你們不要恐懼也不要害怕,我豈不是從上古說明指示你們.
並且你們是我的見證,除我以外豈有真神.
聖賢沒有盤石,我不知道一個.
神係呢度反覆強調.
除jor 我之外,再沒有真神.
四十三章又講,四十四章又講.
所以佢ge 獨一性其實好緊要.
神從來都唔係我地生命ge 一部分.
聽清楚呢句.
神從來都唔係我地生命ge 一部分.
我地ge 生活會包括我地ge 家庭.
我地ge 工作,或者你讀緊書,係你ge 學業.
你ge 朋友,你ge 社交圈等等.
當我信耶穌gwo 刻呢.
我就係咁多楷裡面呢就加多一瓣叫做神.
然之後呢我就好努力將呢一楷呢.
就不停咁去塞落去我生命裡面唔同ge 範疇.
究竟我地係咪咁樣ge 呢.
哦,信jor 耶穌嘛,咁梗係要擺落去啦.
我舉個例呀.
如果有一D 你唔係好熟ge 朋友呢.
黎你屋企探你.
咁因為你唔係成日見佢啦.

$^{161}$係咪呀,真係好似普普通通ge 朋友.
你知道佢唔會賴死唔走係你屋企.
係咪呀,最多坐夜D je .
咁點都會講拜拜咁就離開.
所以呢我斷估呢你個朋友呢.
都唔會似熟賣熟.
然之後呢去到你屋企呢就.
自把自為開你D 櫃.
然之後呢就lor D 零食黎食.
開你個雪櫃lor ye 飲.
會唔會咁架.
我相信係唔會架.
不過我地ge 上帝唔係gwo 個唔係幾熟ge 朋友.
佢入到黎我地屋企.
其實佢係可以隨意開任何一道門.
將一D 唔合乎佢心意ge ye 黎到lor 走.
因為佢係屋主黎架.
所以佢係我地ge 生命ge 當中.
唔係淨係一忽.
係我地生命裡面ge 全部.
呢個就係神ge 獨一性呀.
生命之主嘛.
其實就係咁ge 意思.
當我地ge 生命.
能夠認定ge 時候呢有呢份ge 認定.
第八節.
以下先知點樣講呀.
你們不要恐懼也不要害怕.
係兵荒馬亂.
好艱難ge 時間.
你切切祈禱.
我好想弟兄姊妹記得我地祈禱ge 對象係邊位呀.
我地同邊個人祈緊禱.
我地同上帝獨一ge 真神祈緊禱.
我地同呢一位ge 盤石真神傾緊計呀.
黎到第三個段落.
第九至第二十節.
耳塞描述到有個鐵匠又有個木匠.
佢地點樣製造偶像ge 過程.
第九節係咁樣話ge .

$^{201}$製造挑剔偶像的盡都虛空.
他們所喜悅的都無益處.
他們的見證無所看見無所知曉.
他們變覺羞愧.
虛空呢個字打jor 俾你架啦.
其實呢就係同創世記第一章第二節.
地時空虛混沌.
冤冕黑暗.
神的靈運行在水面上.
呢個空虛gwo 個字呢係一模一樣ge .
咁所以呢耳塞書係呢一度呢.
呢個虛空呢個字呢其實可以引申為呢.
冇意義冇幫助.
甚至呢好似第九節gwo 度所講.
毫無益處.
咁解.
大家只要立一立呢一段經文.
好長呀所以我唔打算同大家讀.
你就會明白點解作者.
講到敬拜偶像係一件幾咁荒謬ge 事.
創世記一章二十七節.
就話比我地聽.
神係照住自己ge 形象做人係咪.
不過當人去做偶像.
調返轉頭ge 時候呢.
gwo D 鐵匠呢gwo 段經文話佢好努力呀.
費寢忘餐搵個鎚黎係咁打係咁打.
直至到筋疲力竭.
木匠呢就為jor 搵gwo D 製造偶像ge 材料呢.
就踏遍整個ge 森林.
然之後呢甚至係親手栽種.
培植一D 貴價ge 木材.
精心量度設計製作.
不過諷刺ge 係第十五節呀.
這樹人可以人可用以燒火.
他自己取些烤火.
又燒著烤餅.
而且做神像跪拜.
做挑剋ge 偶像向他跪拜.
同一棵樹ge 木頭.

$^{241}$人竟然呢lor 部分呢lor 黎當柴燒.
另一部分呢就用黎挑剋偶像.
有冇搞錯係咪呀.
死而所思.
既然個木頭連自己避免被火焚燒.
即係自保呀.
都做唔到.
人點解仍然係要以佢為神.
希望佢幫你希望佢救你呀.
但係人呢就好似第十八節gwo 度咁樣講呀.
如果唔諗下其實邏輯上面根本就唔通ge .
係咪呀.
當我地將呢一個做法同創世記.
神照住自己ge 形象.
黎到做人拍埋一齊黎睇ge 時候呢.
我地就會見到一個非常之強烈ge 對比.
顯出人靠住自己ge 力量去做製造一個ge 神.
係幾咁愚昧可笑.
敬拜偶像呢.
唔單止好似第九節咁樣講到.
盡都虛空都無益處.
拜完之後都有後果.
拜極都係咁驚.
心裡面都係唔平安.
都係仍然覺得羞愧.
十一節咁講.
對於基督徒黎講呢.
我相信呢我地多數.
弟兄姊妹都唔會跌入去呢.
一個拜有形ge 偶像.
崇拜裡面ge .
唔係咁.
不過呢第十節仍然問我地.
誰製造神像.
誅造無益ge 偶像.
當佢咁問ge 時候呢.
我相信我地就好容易答呀.
心裡面就話.
誒梗係唔係我啦咁樣.
我冇份咁.

$^{281}$不過當我地黎到諗真D ge 時候呢.
我地會唔會跌jor 落去一個.
無形ge 偶像崇拜裡面呢.
唔係gwo 舊像.
gwo 舊木頭.
係無形ge .
我有三個問題想邀請大家去諗下.
第一.
有冇追求更多呀.
有冇追求更多.
追求更多ge 物質.
更好ge 居住環境.
放一個長D ge 假期去遠D .
然之後呢.
同屋企人去一個.
以前呢度未去過呀嘛.
咁去第二個國家玩下啦.
等我ge 仔女呢.
係年少年紀細ge 時候.
學多幾樣ye .
傍身呀嘛係咪呀.
總之乜ye 都要多D .
如果你話我唔係咁ge .
我冇咁物質主義.
但係我地仍然可以追求更多ge 權力.
名譽.
成功.
你話唔係咁講呢D ye ge .
不過呢.
我希望人地肯定我多D 喎.
追求別人ge 肯定.
甚至更高ge 層次.
希望能夠有更大ge 自由.
更加遠大ge 理想.
其實呢D ye 可能係好ye 黎架喎.
不過如果我地賦予上述呢D ye ge 價值.
係超過我地ge 上帝佢睇重ge ye 呢.
任何上述我提過ge 事情或者物質.
都可以成為我地ge 偶像.
例如.

$^{321}$我地睇我地ge 孩子ge 學業.
入到個咩學校.
能唔能夠讀到繼續讀上去.
入唔入到大學.
咩大學.
讀咩科.
比起我地睇重去建立佢同上帝ge 關係.
帶領佢地歸向耶穌基督.
更加重要.
又或者.
我需要問我自己.
我究竟係咪一心想成為一個好成功ge 牧者.
一個成功ge 牧師.
呢一個都可以跌jor 落去偶像崇拜裡面.
第二個值得我地反思ge 問題.
就係我地ge 信任同埋安全感.
究竟建立係D 乜ye 上面.
例如.
我地就黎退休啦.
官孫等於姐妹年紀慢慢大.
要諗退休.
究竟我計唔計得掂條數.
好重要架喎.
我地相信我地ge 健身計劃.
能夠確保我地身體健康.
於是我會花上好多ge 操練我身體ge 時間.
同我親近上帝ge 時間.
有一個明顯ge 落差.
咁呢.
我地又跌jor 落去偶像崇拜gwo 度.
係我地生命當中.
神最危險ge 對手.
其實就係我地叫做我呀.
gwo 個神.
主耶穌曾經警告過我地.
唔好過於依賴自己ge 能力.
希望以自我黎到掌控我地ge 未來.
記唔記得路加方十二章.
無知ge 財主gwo 段經文呀.
返jor 教會咁耐應該聽過架啦.

$^{361}$主耶穌提醒我地呀.
如果我地ge 生命有過多ge 憂慮呢.
其實另外一邊係表達緊乜ye 呀.
即係話我地對神ge 信任呢.
其實係出現jor 問題.
如果你呢樣又擔心.
gwo 樣又擔心.
冇ye ge .
即係其實我地唔係幾信得過神je .
第三個.
我地需要反思ge 問題就係.
究竟我地係服侍緊邊一個主人.
如果你未退休.
當然啦我地每一日.
返工都好長呀係咪呀.
其實都幾疲累架.
花上大量ge 時間同埋精力去應付老闆同埋公司ge 要求.
好啦如果我地更加有年長ge 年邁ge 父母.
又有年幼ge 小朋友.
呢個生命階段就仲難添係咪呀.
我地要去照顧我地ge 家人.
我唔會講話我地將我地ge 屋企人當做神咁擺.
我地通常都唔會ge .
不過呢如果佢地佔據jor 我地生命ge 首位ge 話呢.
其實佢就成為jor 我地ge 主人.
我地就會受制於滿足佢地ge 期望.
因為仔女要做乜做乜.
因為我地要陪佢做乜做乜.
然之後呢就忘記jor 上帝.
讓佢取代jor 上帝ge 位置.
又是跌jor 落去偶像崇拜呢一個ge 陷阱裡面.
所以簡單黎講.
凡係係我地心裡面取代神ge 位置ge 人.
字物全部都可以成為我地所製造ge 偶像.
結果就好似二十節咁講.
拜偶像ge 人呢.
佢係以灰為食.
心中昏迷.
使他偏邪.
他不能自救也不能說.

$^{401}$我右手中豈不是有虛方魔.
佢得到ge 只係有虛方藍色gwo 個字.
即係話呢做ge ye 呢全部都係陶冶.
冇用架.
去到最後一個段落.
二十一至到二十三節.
以前阿先知呢就用jor 一個好短短ge 救贖之歌呢.
黎到作為結束.
係二十一節.
神刻意重覆又是講jor 兩次.
你是我的僕人.
佢係要強調以色列文呢一個好獨特ge 身份角色.
正因為神係同以色列文有關係.
所以佢提醒佢地.
要紀念這些事.
請佢地唔好忘記原本佢地係一個.
耳仔又聽唔到.
耳聾眼睛又睇唔到.
眼辣ge 一個僕人.
但係佢地仍然係得夢拯救.
上帝係要改變佢地ge 生命.
歷史話比我地聽.
以色列文呢係無可避免.
係會面對亞述同埋巴比倫ge 侵略同埋被擄.
呢件事係一定會發生.
無得走雞架.
但係呢個正正就係上帝.
去施行管教以色列文ge 一個方法.
今日到我地讀以剎書.
神係要我地好清楚認定.
沒有別神.
佢係gwo 位獨一ge 主.
佢叫我地要遠離偶像.
告別虛空.
就好似23節最後gwo 句咁講.
要見證耶和華的榮耀.
呢個教導唔單止係針對以色列.
其實亦都係同我地每一個信耶穌ge 人講.
因為我地已經得著jor 耶穌俾過我地ge 救贖恩典.
22節我圖沒了你ge 過犯.

$^{441}$我圖沒了你ge 罪惡.
所以即使我地要面對生命裡面好大ge 艱難都好啦.
其實恩著我地同以色列文都係一樣.
都係耶和華ge 僕人.
我地同神有關係.
係ge .
就係因為我地同佢有關係.
佢就向我地呼籲叫我地歸向佢.
希伯來書12章第6節.
你一定聽過.
唯獨佢愛ge .
佢必定管教.
作為一個簡單ge 總結.
係今日我同大家分享ge 耶穌亞書44章裡面頭5節.
神話佢係會持守佢不變ge 應許.
所以唔單止係係以色列文身上.
係我地ge 生命裡面都係咁樣.
佢好希望我地即使面對生命裡面好大ge 艱難.
我地仍然見到佢做緊ye .
我地見到佢做緊ye .
求佢幫助我地去克服呢D ge 挫折.
呢D ge 障礙.
以致我地明白佢做緊乜ye .
辨識佢ge 心意.
以致我地一個又一個好願意ge .
歸向佢.
結果係蒙祝福.
不過gwo 個過程對唔住呀.
唔可以跳過架.
最好中間gwo 節冇jor .
跳跳去後面咪好囉.
對唔住呀我地ge 上帝唔係咁樣.
跟住去到第二個段落.
第六節到第八節.
先知挑戰我地要思考.
究竟神係我地心裡面gwo 個ge 地位.
究竟係咪好似佢係宇宙裡面ge 地位一樣咁高.
佢叫我地再次認定.
佢先至係gwo 位至高創造天地ge 主.
佢要掌管一切.

$^{481}$佢係獨一ge 神.
跟住去到九節到第二十節.
佢就講到敬拜偶像gwo 種ge 虛無虛假.
佢要我地反思ge 就係.
究竟我地係咪不停咁追求多D 多D 多D .
又或者我地ge 信任同埋安全感ge 基礎究竟係邊度.
究竟我地係服侍緊邊個主人呀.
去到最拉尾gwo 個段落.
以下就以一個救贖之歌提醒我地.
我地同神係有關係ge .
呼籲我地歸向佢作為結束.
有一個新約聖經ge 學者Antirite.
佢又咁樣講.
Isaiah gives us the God whom we didn't want.
but who amazingly wanted us.
以賽書向我地揭示jor 一個我地唔係幾想要ge 神.
我地唔想佢咁樣個波.
不過佢係好奇妙咁去想要我地.
想得著我地.
雖然人好努力咁樣塑造自己ge 偶像.
但係呢D ge 偶像只不過係虛空.
同神ge 榮耀冇得比架.
所以請我地記住佢ge 手係木頭上面.
黎到展開.
釘上.
就係要向人.
向你同埋我.
去展示佢ge 美麗同埋榮耀.
Antirite跟住就咁樣講.
比個原本ge 講法翻譯出黎係咩呢.
佢話呢個其實係一個好唔同ge 神黎架.
呢個係一個好危險ge 神.
一個顛覆性ge 神.
一個好似呢.
盲jor ge 黑衣咁ge 神.
佢手上面呢.
係傷痕累累ge .
黎到我地面前ge 神.
佢係一個透過火.
透過風.

$^{521}$透過餅同埋酒.
用佢ge 血肉之軀.
黎到我地面前ge 神.
佢同我地講.
唔係你揀選jor 我.
而係我揀選jor 你地.
楊康十五章十六節上半節.
係咁樣講.
然而我地真係能夠睇得見.
我地所相信.
或者我地信jor 好耐架啦.
呢位耶穌究竟係一位點樣ge 主.
\newpage



\section{出埃及記 35:30-36:1}
\label{sec:031755}
\textbf{比撒列的自白}~~~~杜崇基牧師
\newline
\newline
連結: \href{http://www.church.com.hk/acms/content.asp?site=cdc\&op=show\&type=product\&code=031755\&layout=sermon}{\texttt{031755}} ~~~~ 語音日期: 2026-03-01
\newline
\newline
\hyperref[sec:031754]{\small{< < < PREV SERMON < < <}}
~
\hyperref[sec:index_chronic]{\small{[返順時目]}}
~
\hyperref[sec:031756]{\small{> > > NEXT SERMON > > >}}
\newline
\newline
出埃及記 35:30-36:1
\newline
\begin{longtable}{cl}
\hline
\hline
章節 & 經文 (和合本修訂版)\\
\hline
35:30 & \begin{tabularx}{0.7\textwidth}{X} 摩西對以色列人說：「看，猶大支派中戶珥的孫子，烏利的兒子比撒列，耶和華已經題名召他， \end{tabularx} \\ \\ \relax
35:31 & \begin{tabularx}{0.7\textwidth}{X} 又以神的靈充滿他，使他有智慧、聰明、知識，能做各樣的工， \end{tabularx} \\ \\ \relax
35:32 & \begin{tabularx}{0.7\textwidth}{X} 能設計圖案，用金、銀、銅製造各物， \end{tabularx} \\ \\ \relax
35:33 & \begin{tabularx}{0.7\textwidth}{X} 又能雕刻鑲嵌用的寶石，雕刻木頭，做各樣精巧的工。 \end{tabularx} \\ \\ \relax
35:34 & \begin{tabularx}{0.7\textwidth}{X} 耶和華又賜給他和但支派中，亞希撒抹的兒子亞何利亞伯能教導人的心。 \end{tabularx} \\ \\ \relax
35:35 & \begin{tabularx}{0.7\textwidth}{X} 耶和華使他們的心滿有智慧，能做各樣的工，無論是雕刻的工，圖案設計的工，用藍色、紫色、朱紅色紗，和細麻作刺繡的工，以及編織的工，他們都能勝任，也能設計圖案。」 \end{tabularx} \\ \\ \relax
36:1 & \begin{tabularx}{0.7\textwidth}{X} 比撒列和亞何利亞伯，以及一切心裡有智慧，蒙耶和華賜智慧和聰明，懂得做聖所各樣用途之工的人，都照耶和華所吩咐的去做。 \end{tabularx} \\ \\
[1ex]
\hline
\hline
\end{longtable}
$^{1}$各位弟兄姊妹平安.
我們今天很特別邀請了比薩烈上來.
和我們分享他的見證.
我們一會兒就請他出來.
各位弟兄姊妹.
我是比薩烈.
其中我有份做的工作就是.
大家應該知道是甚麼吧.
對了 就是約櫃.
其實不止約櫃.
其實所有會幕裡面的帳幕.
都是我有份做的.
所以首先介紹一下自己.
我是猶大之派中護爾的孫.
也是烏利兒子比薩烈.
或者你都認識我爺爺護爾.
當年約書亞和亞馬力人爭戰的時候.
摩西上山祈禱.
摩西很累的時候.
亞倫和護爾就搬了一塊石頭給他坐下.
一邊一個扶著他的手.
直到以色列人大獲全勝.
阿穆斯邀請我來到沙圳這裡.
去分享我的事.
他在《出埃及記》.
其實我的名字是親自在《出埃及記》裡面出現的.
神就說他是提名召我的.
又說他是以神的靈充滿我.
使我有智慧,聰明,知識.
能夠做各樣的工作.
所以我今天想親自和你們去作見證.
原來神不單是呼召傳道人.
神也是呼召你和我.
因為神的呼召.
不只是叫我們做所謂神聖的事情.
他是叫我們能夠被他聖靈的充滿.
按他的心意去做各樣的事情.
我的名字在希伯來話裡面.
其實就是在神的任避下的意思.
根據聖經學者的分析.

$^{41}$這個名字的含義.
應當是解釋為.
我的帳幕是父神.
這個帳幕也可以解作避難所.
反映了神將人藏於帳幕的觀念.
這個名字對我來說是特別的貼切.
因為我要成為神帳幕裡面的帳人.
以色列人能夠在埃及被奴役.
多數人都缺乏了一個精巧的工藝.
在抗疫的環境裡面.
能夠完成到這麼精細.
跨越工種的這種技術.
確是神所賜的智慧和組織能力.
我們才能夠做到的.
回想當我們離開埃及的時候.
前面是荒涼的抗疫.
因此需要時刻尋找水源.
而食物的供應只能夠依靠嗎哪.
眼見百姓對神常常都是三心兩意.
照人來看.
他們的首要任務.
最重要是在抗疫裡面能夠生存下去.
或者求旅途的平安.
然而神對摩西的吩咐卻是.
當為我做聖所.
使我可以住在他們的中間.
弟兄姊妹這句話將神的心意揭開了.
神不單止拯救我們出埃及圍爐之地.
更要將我們塑造成一群.
以神同在為中心的子民.
會幕的目的不是一個建築的成就.
而是預表神的親自臨到.
在新約就成就在耶穌基督.
藉著聖靈在教會裡面運行.
而我所要做的.
就是會幕,法櫃和旗上的櫃蓋.
會幕裡面一切的器具.
公桌和公桌的器具.
純金的燈台.
一切的器具.

$^{81}$香壇,梵祭壇.
和壇的一切器具.
洗鑿盤和盤座.
供祭司積分用的一些精緻的禮服.
如果你今天看每日讀經釋義.
我們看到《出埃及記》三十九章.
特別是說到亞倫的聖服是很特別的.
而亞倫祭司的聖衣和衣服的安排.
那些膏油和聖所用的芬芳香.
都是很清清楚楚地.
神指示應該怎樣做的.
無論是尺寸,材料,顏色,器具,次序.
都不是隨便這樣做的.
因為神要我們知道祂是聖潔的.
祂的敬拜不是混亂的.
祂住在我們中間是恩典.
也是秩序.
弟兄姊妹.
當我聽見神說要以神的靈充滿我的時候.
我心裡面是非常感動的.
聖靈充滿.
原來就是神賜智慧,聰明和知識.
而完成祂的吩咐.
在今天來說.
就是在我們的家庭,職場和教會裡面.
能夠活出一個合神心意的能力和品格.
我知道在舊約的時代.
聖靈是特別揀選那些特殊任務的人.
不過這個特殊任務不只是限於先知和君王.
而我竟然作為一個工匠.
都能夠被提名.
可見神是很著重這個使命.
今天神也很著重你的使命.
可能你不是在教會裡面有甚麼領導的角色.
但可能你在家裡面作為一個家長.
你在職場上作為一個上司.
甚或你在你的社區,社團當中.
你都作為一個帶領者的時候.
原來神都給了一個特別的召命給你.
這個召命就是要你所在的地方.

$^{121}$能夠顯出你是屬於神的.
你是有神的性情.
你是有神的價值觀.
但知派阿和利亞伯跟我一起做工.
包括了整個被神賜智慧的群體.
所以我按照耶和華的吩咐.
完成了所有的工作.
原來聖靈的充滿.
是神將人從自己去作主.
帶回去神作主.
不過很老實說.
我很明白今天在教會和家庭裡面.
如果只靠人的力量服事的話.
很容易有三種情況出現.
第一種情況出現.
就是我們可能只是靠責任感.
去推動我們繼續的事奉.
我們看到有些人都很忠心.
在教會裡面甚麼都肯做.
但久而久之就會發現自己.
其實是在頂硬上.
勉強自己在崗位裡面服事.
又不敢對別人說不想做.
結果失去了事奉神的喜樂.
而且他越事奉的時候.
就感到越疲倦.
甚至他會覺得別人不明白他.
神也好像要他不斷不斷去付出.
弟兄姊妹.
責任感可以推動我們走一段路.
但不能夠保守我們完成到底.
只有聖靈充滿才能夠使人.
不是在一個被任務驅動.
而是一個被愛驅動.
不是被焦慮推著他去走.
而是能夠被主引導而行.
第二種情況.
有些人的事奉只靠他的血氣.
血氣的意思就是憑自己一股的動力.
或者是一個的野心.

$^{161}$血氣的意思就是他很多時候.
表現是很焦急的.
很多時候一用血氣的意思的時候.
就很想在其他人面前拿到主動權.
或者拿到領袖的角色.
或者是要壟斷整個決策裡面的討論.
甚至如果差起來.
他會用一些比較論斷的方式去壓制別人.
其實在教會有時我們都會將屬靈的事供.
變成了人的成就.
當事供只講求成效.
而放低了團隊的合一的時候.
當講求你輸我贏.
勝過在事奉中建立的生命.
或者當講求對錯.
而誰是被牧養的時候.
我們就會忘記神朝聚我們一起事奉的初心.
聖靈充滿的重要.
就是祂不單止給我們有能力.
祂更加先掌管我們的心思.
祂使你在忙碌裡面仍然能夠柔和.
在分歧裡面仍然能夠有節制.
甚至在衝突裡面.
你仍然能夠去尋求和睦.
換言之.
聖靈的充滿不是讓你更強勢.
而是讓你更像主耶穌.
第三種推動我們去事奉的情況.
就是我們只靠我們的熱情.
我亦都見過一些人開始是很火熱的.
他很願意投入.
但是熱情過後.
始終都要面對現實的生活.
小朋友要照顧.
長者要照顧.
自己的工作都要向人交代.
身心都有時疲倦.
需要休息.
所以如果事奉只是建立在熱情之上.
人就很容易.

$^{201}$不是突然停擺.
而是靜悄悄地退場.
所以聖靈充滿.
是使人從一時情緒的高漲.
轉到我們每天與神同在.
其實回想起來.
我們不是每一刻都能夠士氣高昂.
我們每一刻都可以被祂扶持.
又或者我們工作的時候.
不是每一件事都很順利.
但是每一次當我們倚靠神的時候.
祂就能夠使我們有更新的生命.
弟兄姊妹.
讓我用神對我的話作出總結.
神用祂的靈充滿我.
結果就是我能夠更有智慧.
更有聰明.
更有知識.
能夠做各樣的工作.
意思就是.
聖靈充滿使基督徒能夠在平凡之中.
活出一個不平凡的中心.
在家庭 在職場 在教會.
不是說我能做到多少.
而是我做的事.
是不是更合乎神的心意呢?.
但是神並沒有叫我一個人去完成這個會望.
神對摩西說.
我又派彈支派中.
阿希撒末的兒子.
阿荷利亞伯與他同工.
弟兄姊妹請你留意.
神的安排不是說你很厲害.
所以你做完所有事情.
神的安排是說.
我將恩賜分給眾人.
你們在同一個使命中彼此配搭.
我就是來自猶大之派.
阿荷利亞伯就是來自彈支派.
背景 性格 做事的節奏都很可能不同.

$^{241}$但是神將我們放在一起.
是為了教我們一件事.
對工的根基不單單是講默契.
更加是要講共同的意象.
共同意象的能力是超越個人性格的差異.
因為大家都是被聖靈充滿的.
在牧羊領袖當中.
我們要有一個共同的意象.
簡單來說就是我們要跟同一把尺.
那把尺就是神的吩咐和神的同在.
過於單單是尋求團隊的共識.
弟兄姊妹容許我用一個工作的例子去講.
如果一個工場沒有同一把尺.
人人都按自己的眼光去量度的時候.
最終每一件製成品都不會合規格.
同樣如果牧羊領袖沒有了共同的意象.
事工就會被個人的喜好拉扯.
有些人很著重活動.
有些人很著重人數.
有些人很著重形式.
有很多人著重效率.
表面搞到教會好像很忙.
但是裡面可能會失去了焦點.
回牧之所以能夠建造.
不在於我一個人技巧多高.
而是在於神給了大家一個很清楚的規格.
好讓人知道我們在建造的時候.
是神同在的居所.
而不是人所做的金牛毒.
因此我想將共同意象講成一個具體的對齊.
你們領袖團隊可以試一下用來檢視一下.
究竟我們做每一件事.
是否能夠引導人更去敬畏神.
更去親近神.
更去馴服神呢.
如果我們將事工做得很成功.
卻將人帶來神的同在.
這樣我們就不在建造就在拆毀了.
原來回牧是為敬拜而去建立的.
同樣教會也是為塑造門徒而存在.

$^{281}$如果我們用業績來衡量成功.
很快就會犧牲一個個的家庭關係和靈明.
共同的意象必須保護一件事.
就是人比事重要.
生命比成效重要.
在工場上我們會趕工.
趕工的時候最怕就是將製成品做壞了.
而需要整批來重造.
教會也是一樣.
如果領袖之間失去了信任.
事工再快也只是加速了內耗.
留意我們港合一不是故意拖慢了整個決策的過程.
乃是避免弟兄姊妹因未明白而造成傷害.
所以容許我再一次說.
共同意象要成為真正的童心.
必須有聖靈的充滿.
聖靈使領袖願意放下自己.
願意聽 願意讓 願意承傳.
願意將我的做法放在主的心意之下.
真的很鼓勵我們.
我也很感激他所說的話.
所以我繼續分享.
牧師最近看了一本書.
這本書叫做《類代》.
Generations.
是由兩位夫妻檔的牧師所寫的.
就是Pastor Hal和Pastor Leah.
是新加坡神之心教會的創堂牧師.
話說在2018年的時候.
當時我和Alky和阿傑一起去神之心教會參觀.
或者去認識他們的事工.
當時已經有一個很深刻的印象.
他們的年青人個個都走出來見證.
見證神如何使用這間教會成為青少年的一個家.
讓這個家感受到歸屬感.
我甚至聽到一個見證.
一個年青人.
那個教會在新加坡.
有一個年青人甚至在馬來西亞過境.
然後坐三個小時車回來這間神之心教會.

$^{321}$可以看到這一班年青人.
多麼以教會為他們的歸屬感.
我又聽到一個見證.
其中一位年青人因為他的家人是一個破碎的家庭.
所以他家人沒有什麼可以給他.
支持他去讀書.
神之心教會見到他家境這樣的情況的時候.
特別給他一個獎學金幫助他完成.
大家也都知道.
新加坡是一個很求競爭的一個地方.
所以他們對學業非常的緊張.
神之心教會同樣是幫助這些年青人.
在崇拜之後開放那個地方.
給那些年青人可以去溫習.
也有些哥哥姐姐去幫助他們補習.
讓他們的功課能夠追得上.
更加原來在這班教會裡面.
有很多年青人沒有爸爸.
在他的家庭裡面因為父母離異的緣故.
沒有爸爸.
所以在父親節的時候.
這班年青人是最傷心難過.
這個Pastor Howard牧師.
每年就找到這些年青人.
寫一張卡去鼓勵這些沒有爸爸的年青人.
再一次的鼓勵他們.
說你要勤力 你要努力.
你在神的愛當中你必定能夠成長.
同樣我也聽到不少的見證.
是說到這班教會.
怎樣去一代一代的人.
去興起投入在事奉.
當我去這班教會的時候.
其中一個最深刻的印象.
就是看他在教會裡面.
那個控制音響的那個.
就不是那麼大年紀的.
其實我們後面的這位弟兄.
也是剛剛浸泥加入的.
其實他年紀也很大了.

$^{361}$為什麼呢.
因為他控制音響的都是十幾歲的年青人.
最重要的就是他不只是一班人控制.
他是幾代的人控制的.
一個二十多歲的年輕人.
去幫助一個十六七歲的年輕人.
那個十六七歲的年輕人又幫助一個十一二歲的年輕人.
那個十一二歲的年輕人就去幫六七歲的小孩.
讓他們很小的時候.
就有機會投入在事奉的行列當中.
所以什麼叫做「類代」呢.
「類代」就是讓他們去建立.
一群一群不同年代的領袖.
他特別提醒.
「類代」這個字其實是有聖經根據的.
是出自於《爾賽亞書》五十八章十二節.
他說你要建立拆毀類代的根基.
或者可以說是立起代代相承的根基.
教會是要不斷興起不同世代的領袖.
服侍破碎的青少年.
為他們預備一個家.
所以「類代」不只是說繼承的.
「類代」是同時擁有一層層的領袖.
一起去服侍.
「類代」不代表是取代.
「類代」是代表增援.
他就用了一個例子.
不知打籃球.
大家就看看有一隊球隊打籃球.
他有一個專有名詞叫Deep Bench.
Deep Bench的意思就是.
他有很多作後備的.
那些後備球員其實球技不會比正選差.
但是他就不斷在比賽的過程中.
他將這些人因為他有新力軍.
他可以換進去成為增援整隊球隊的攻擊力.
所以我們看到Deep Bench.
就是說有很多儲備的領袖在當中.
其中神之心教會有一樣東西給我很大的印象.
就是他有三批人服侍.

$^{401}$剛才我們說他不同世代有領袖.
敬拜隊是有三個部份的敬拜隊.
一個部份是在服侍.
一個部份是休息的.
還有一個部份有時他們帶去一些外邊的特會的時候.
他們就會用他們這一班的領袖出去大敬拜.
我心想.
嘩 如果筲箕有這麼好就真是發達了.
其實現在我們的敬拜隊已經人多了.
不過其實都辛苦的.
因為都經常要戴.
沒得休息的.
哪裡有得休息的呢.
當然我這樣說的時候.
都是鼓勵.
我都期待我們的教會都能夠多代同行.
因為不只是靠我們上一代的.
我們都要多代的同行.
回想起筲箕.
回看我們教務團隊都很感恩的.
有不同年齡層的教務團隊.
我呢就屬於最早的一批.
牧師我都不想提大家的.
不過我經常都說.
過多幾年都退休的了.
但是我們還有一批.
阿傑啊 Samuel啊那些.
四十出頭的那些.
都是很當打的.
很厲害的.
我們還有那些三十歲出頭的.
我們有可風啊 曉怡啊.
另外我們當然顏牧師又小那些.
他做得可以做很久的.
我們其實是deep bench.
我們都很厲害的.
我想告訴大家.
我們真的要突破.
我們要突破就是.
我們的領袖不是只有一層的.

$^{441}$我們都要有不同的年代的.
同工啊 迎接未來的挑戰.
我真的覺得未來的挑戰.
是我們都想不到的.
發生的事我們都想不到.
牧師本來想著放長假.
為什麼我要放長假呢.
我其實是想著去中東那裡.
探訪那些宣教工場的.
誰知道昨晚六點多收到一個訊息.
飛機都不開.
那個地區的機場全部關了.
大家知道為什麼嗎.
因為美國和以色列打伊朗.
而我去的國家一看地圖.
全部都是中間的地方.
伊朗在那裡.
我去的國家伊拉克也在那裡.
我又去約旦.
約旦就在以色列和伊朗中間.
然後敘利亞.
全部都是這些受戰火波及的國家.
所以今天早上要馬上取消假期.
所以心裡面是很不容易過.
但是其實我少了假期.
或者少了這個行程事小.
最重要是求主保守.
是工場上的工人.
最重要是保守.
在這場戰爭裡面盡量沒有那麼多人傷亡.
這個才是重要.
其實局勢是風雲變色的.
我們都不知道前面會發生什麼事.
但是重要的是神要使用我們.
使用我們在這個世代裡面去見證祂.
不過在我們見證神的過程裡面.
我們都會有一些要克服的地方.
我嘗試舉出三個處境.
我猜大家都可能會聽過.
第一個處境就是在教會裡面.

$^{481}$有一些父母是很投入事奉的.
但是那些孩子有時候就感覺到.
教會搶走他的父母.
這個原因是什麼呢.
原因就是他感受到父母經常想著事奉神.
所以就無暇兼顧這些孩子的成長.
更加嚴重的就是因為這些父母本身.
他生活很匆忙.
甚至在工作上受很大的壓力.
以致他有時候在小朋友面前.
表達出那種躁動.
甚至父母之間都會有些爭執.
久而久之那些小朋友不是遠離教會.
就是把信仰看成一個破壞家庭的東西.
弟兄姊妹 多代同行在這裡的第一步.
就是教會的共同意象必須清楚地宣告.
我們所建造的不是事奉的機器.
我們所建造的是生生不息的門徒.
和敬虔的家庭.
最近我和韋傑傳道一起去拜訪楊啟厚博士夫婦.
他們兩位提醒了我們.
他們說銳意家庭門訓是很重要的.
記得我們不是要成為別人心目中的樣板.
而是要成為弟兄姊妹的同行者.
一家人吃飯其實是很重要的.
記得上次我和大家說家庭祭壇很重要.
我們一定要有家庭崇拜.
楊博士說其實一家人一起吃飯也很重要.
我不知道你們一家人一起吃飯.
一個星期有多少晚.
我自己數著數著可能只有三晚.
因為你也知道做教務.
平時很多晚上都要開會.
或者有其他的事情.
所以我很珍惜這三晚的時間.
可以和我們家裡和善恩和太太一起去分享一下.
分享一下我們家裡的事情.
分享一下我們自己的事情.
也可以有機會彼此紀念彼此代禱.
這個其實是很重要的.

$^{521}$不要看輕這個分享.
原來很多家庭裡面的困難.
在這個分享裡面能夠去解決.
今天我們特別要留意.
事奉的人不要忽略了家庭.
我們有優先.
我經常都記得幾個同心圓.
第一個同心圓當然就是我們以愛神為我們最重要.
但是愛神愛哪一個圈.
就是愛我們的配偶.
接著就是愛我們的家人.
接著就是愛我們的教會.
接著就是愛我們社區的人.
你會發覺這個同心圓是一直向外延展的.
這個也提醒了我們.
什麼是我們生命裡面的優先.
所以弟兄姊妹.
我再一次提醒.
當我們很著重事奉神的時候.
不要忘記.
我們的家人.
我們的配偶.
都是我們一個很重要的人.
是要我們關注的.
在第二個的情況裡面.
我們發現在教會有這樣的情況.
就是兩代對教會未來的看法是會有不同的.
可以說是在傳統與更新之間常常拉扯的.
年長的人他是很愛教會的.
他是持首尾好的傳統.
但是年長的人有一個特性.
就是比較習慣熟悉和安全的方式.
但是年輕一代很多時候.
比較對教會的現狀是不滿的.
他們期待革新.
但是當壓力累積的時候.
就很容易出現了這樣的對話.
年輕人就說.
你都不明白我們.
而年長的人就說.

$^{561}$你都不尊重我們.
弟兄姊妹.
其實多代同行.
我們是需要修復和和解的.
這個是很重要的一件事.
你想想你家裡你和父母有不同意見.
一定是有的.
你小時候就沒得選擇.
父母說什麼你就要聽.
當你長大了的時候.
就不一定全聽他的.
但是我們要怎麼做.
你發不發覺.
多數都是要我們遷就他.
如果我們做小的.
我們一定要遷就.
而這個遷就不是代表我們輸.
這個遷就是代表我們愛他.
當然有些事情是不能遷就的.
在真理的事上我們不能遷就.
但是有很多事情是關乎做法的.
是關乎文化的.
其實我們可不可以讓一讓呢.
我們需要修復和和解.
在神之心教會裡面有一個見證.
就是說這個爸爸因為吸毒.
和被人找到毒品.
就在監牢裡面坐了很長的時間.
他因而忽略了他自己的兒子.
這個兒子最初是很生這個爸爸的氣.
覺得這個爸爸不配做他的爸爸.
但是後來這個兒子在神之心教會裡面.
他感受到神的愛.
他也感受到神要他去原諒他的爸爸.
所以他就寫了一封信給他的爸爸.
就是說其實他已經在主裡面原諒了他.
他也希望他爸爸能夠在主裡面.
繼續去成長.
因為他在監獄裡面信了主.
爸爸的回信對這個兒子就是說.

$^{601}$我很詫異.
你竟然能夠原諒我.
其實我在這個監獄裡面.
我已經學習事奉神了.
我學習彈結他去敬拜神.
你說這件事是不是一件很美好的事呢.
父子能夠冰釋前嫌.
是因為主的愛在當中.
原來神在他們各自雙方都做了奇妙的事情.
今天我們要在共同的意象下.
同守著一個核心的價值.
就是我們要敬畏神.
我們要尊重人.
我們要以愛相待.
有時候我們在方法上是有不同的.
但是我們都能夠彼此的接納.
彼此的學習.
如果領袖能夠同心.
就能夠將會眾從不同立場的爭論.
帶回神的家.
大家願意能夠彼此聆聽.
尋求最好的出路.
第三個情況.
可能在教會裡面都會出現的.
就是年輕人想服侍.
但是卻怕被比較.
怕做錯.
怕永遠都只是一個頂替.
次好的東西.
有些年輕人他不是不願意.
而是他害怕.
有時候教會有一種氛圍.
就是覺得你不可以做錯事.
但是其實在教會裡面.
有些人確實是要被剁磨的.
經驗實在是不夠的.
你有沒有空間讓他可以做錯呢.
不是刻意要破壞.
而是有時候他真的做得沒那麼成熟的.
沒那麼圓滑的.

$^{641}$或者是覺得他的看法很極端的.
但是是不是因為這樣.
教會就打壓他.
不讓他說話呢.
所以今天我們要特別留意這件事.
有些人經常說.
年輕人是教會未來的領袖.
但是這句話.
其實本身是帶著一些問題的.
這個問題是什麼呢.
那就是說年輕人.
不是今天的領袖.
還有另一個問題很嚴重的.
就是說年長的人.
是過去的領袖.
而不是現在的領袖.
其實你想一下這個其實是很小心的.
所以在神之心教會裡面.
他特別說.
青年人就是現在的領袖.
如果青年人就是現在的領袖的時候.
我們有沒有一個空間讓他去犯錯呢.
我們有沒有一個空間讓他去嘗試呢.
弟兄姊妹.
我們看到香港過去幾年.
曾經有一個很大的社會的運動.
有很多的年輕人離開了他們自己所屬的堂會.
他們叫做離堂者.
他們失去了自己的身份.
他覺得教會不能夠明白這個世代的需要.
所以他們離開.
同樣在這幾年裡面有一間教會.
他重新的成立.
他是由一位神學院的老師.
走出來去成立的一個教會.
是專幫助這些離堂的青少年人.
青年人或者職青去重建的.
所以這間教會已經有七年的歷史了.
很快的七年的歷史.
神使用這間教會去幫助很多離堂的信徒.

$^{681}$他們回到他們這個地方.
是有一個信仰的生命.
弟兄姊妹.
類代同行必須將這個承傳成為一種文化.
不是將事供交出去就算完成.
而是要將生命價值方法.
負擔和權定一起去交托.
千萬不要想著等你準備好了.
我才給你事奉的機會.
而是我要先給你一個舞台.
我能夠陪你走過一起成長.
因此類代同行必須將承傳.
變成一個是大家都能夠預期的.
大家都能夠交接的制度和節奏.
如果這樣的話.
年輕人就不單單只是一個旁觀者.
或者一個幫手.
而是能夠被承傳的承擔者.
而年長的人也不只是苦苦的支撐.
能夠放心將下一代的人帶出來.
弟兄姊妹.
會幕的建造不單只是說做得精巧.
更加能夠說我們能夠傳下去.
技巧如果不能夠傳承.
工程只是做了一次.
生命如果不承接的話.
事工只能靠少數人的維繫.
在工場裡面有一個想法.
如果你可能也有這樣的經驗.
如果你打算做一件事.
快快手手搞定它.
就你自己下手做就是最快.
但是你要教人做.
你要幫人去懂得做.
其實你要多花幾倍的心機,時間,力量去做.
但是就是這樣.
我們願不願意多花點時間,心力.
去走多一里路呢.
如果那個師傅想做得快.
他自己做完所有事情.

$^{721}$徒弟就永遠學不會了.
因為快只能得到一個很快的成果.
但是唯有慢才能夠孕育到人才.
而教會真正要的.
從來都不只是將活動辦完.
而是要將門徒栽種出來.
如果你回燒浸一段日子.
你會覺得燒浸的事工應該都相當多.
其實燒浸最基本的東西.
可以說是套餐.
除了回崇拜之外就是回小組.
但是回小組之外.
其實教會都有很多機會讓我們去學習事奉和成長.
我們有很多特別的小組.
我們有印尼侍工.
有街頭報道隊.
有老人院侍工.
其實我們教會都有很多不同的工作.
可以一起去事奉.
如果所有的工作.
我們都是想找那些最熟手的人去做.
而不能夠讓一些新人加入的時候.
始終我們做這個工作都不能夠承傳下去.
今天.
保羅教導提摩泰說.
你在許多見證人面前.
聽見我所教訓的.
也要交託那忠心能教導別人的人.
提摩泰後書二章二節.
有些人說每一個人都需要一個保羅.
一個巴拿巴和一個提摩泰.
保羅就是我們屬靈的導師.
巴拿巴就是我們同心合意的夥伴.
提摩泰就是你要栽培的門徒.
不知道你在你的身邊有沒有保羅.
有沒有巴拿巴.
有沒有提摩泰呢.
在教會裡面.
神給我們有牧區的領袖.
無論是我們有組長.

$^{761}$我們有區長.
其實都是神給我們一個機會.
祂成為我們一個導引我們的人.
去幫助我們怎樣在信仰生命裡面去成長.
但同樣我們又有巴拿巴.
巴拿巴就是我們的組員.
就是可以和我們同心合一去事奉的人.
但有一樣東西你要注意.
在你的信仰生命裡面.
有沒有提摩泰呢.
就是你要刻意找一些人.
可能他的屬靈生命給你錢的.
你去幫助他.
你守望他.
你與他同行.
你知不知道往往你去做這件事的時候.
是你屬靈生命成長得最快的.
如果你沒有了下面這層提摩泰的話.
很多時候你就有一個危機.
就是你覺得OK了.
我信了主這麼多年.
我知道了.
那些講道.
牧師講完上一句.
我都知道下一句.
不過這句不適合我.
你猜不到我.
今天你猜不到我.
所以其實真的.
當我們信得主久的時候.
如果我們沒有了提摩泰.
沒有了我們要建立的下一代人.
其實是令到我們的生命.
都會有一種停滯不前.
所以我們怎樣能夠傳承呢.
我們有幾個具體的策略.
第一個就是我們有影子的同工.
意思是什麼呢.
就是我們上一代邀請下一代的人.
一起同工.

$^{801}$這個不一定是年紀的上一代下一代.
這個是給我年資錢的人.
或者給我信主年日錢的人.
上一代能夠邀請下一代的人同工.
而這個下一代的人.
不是臨時來跟一下的.
而是他能夠在事奉中.
是有跟隨實習和檢討的機會的.
讓他們能夠從上一代身上.
看見怎樣預備.
怎樣決策.
怎樣面對衝突.
怎樣牧養人心.
在沙鎮裡面.
我覺得代土隊其中一樣東西.
是可以幫到我們看到.
有代代的傳承的.
在過往的日子.
代土隊都好像被人覺得是精英分子.
一般人都是望而卻步.
覺得很難的.
但我相信神也要興起新一代的人.
都是對禱告有這個心智的.
他都能夠投入在代土的行列當中.
我們要打破一種精英的主義.
嘩 要很厲害才有代土隊.
不是.
當你有心.
你願意.
你將生命獻情的時候.
你在當中又成長又學習.
又或者在事奉裡面.
我們可以有拍檔.
商人的服侍.
就是我們有拍檔.
我們怎樣能夠讓新手.
在一個安全的範圍裡面嘗試.
讓成熟的人能夠在旁示範扶持和修正.
又或者我們說.
剛才我所說的.

$^{841}$我們有沒有輪替的制度呢.
我們有沒有真的有三班人.
一班人是服侍.
一班人休息.
一班人是對外的工作呢.
或者我們有沒有清楚的交接呢.
交接不是只是做些什麼流程.
或者聯絡人的名單.
更包含為什麼我要這樣做.
這個侍工的核心價值.
哪些事情是不可以妥協.
哪些做法是可以調整的.
所以今天我們要留意.
成全最難的.
從來都不是技術.
而是一個心態.
我們要將聖靈的充滿.
放回中心.
將共同的意象落實到生活.
今天我們有兩個行動.
我可以回應的.
第一我們要將共同的意象.
化成一個可以共同檢驗的話.
例如我們要建造一個.
以神同在為中心.
以門徒生命為優先.
以多代同行為文化的教會.
所以我們每次去討論我們的工作.
我們領袖就要問.
這件事能不能夠更盛載神的同在呢.
是否能夠更去建立門徒呢.
是否能夠促進多代同行呢.
共同意象不是寫在紙上的.
而是用來做決定.
第二個行動就是我們要求主.
聖靈充滿.
我們要先被主更新.
才去服侍.
服侍神不是靠血氣或者一股熱情的.
而是要常常回到主的面前.

$^{881}$求主潔淨我們的動機.
教準我們的方向.
更新我們的愛心.
賜下智慧.
所以你能夠在家庭和教會.
都能夠活得合乎神的心意.
你可以嘗試在這個星期裡面.
實踐一個具體的行動.
或者你與你的配偶兒女或者同工.
進行一次十分鐘的對話.
我們可以分享一下近況.
我們或者承認要調整的事情.
和彼此代禱.
那我們一起祈禱吧.
主啊.
求你以聖靈充滿我們.
使我們能夠被你掌管.
在家庭裡面有愛和節制.
在教會當中有合一和智慧.
求你使牧羊的領袖在你裡面同心.
得著共同的意象.
以你同在為中心.
以門徒生命為優先.
以多代同行為我們的文化.
也都求你使上一代願意承傳.
下一代願意承擔.
使我們一起走在你的面前.
建造你的家.
我們這樣在你面前祈禱.
是奉主耶穌基督得勝名字祈求.
對.
\newpage



\section{創世記 25:24}
\label{sec:031756}
\textbf{聖經人聖經事:果然是雙子}~~~~黃克勤牧師博士
\newline
\newline
連結: \href{http://www.church.com.hk/acms/content.asp?site=cdc\&op=show\&type=product\&code=031756\&layout=sermon}{\texttt{031756}} ~~~~ 語音日期: 2026-02-23
\newline
\newline
\hyperref[sec:031755]{\small{< < < PREV SERMON < < <}}
~
\hyperref[sec:index_chronic]{\small{[返順時目]}}
~
\hyperref[sec:031757]{\small{> > > NEXT SERMON > > >}}
\newline
\newline
創世記 25:24
\newline
\begin{longtable}{cl}
\hline
\hline
章節 & 經文 (和合本修訂版)\\
\hline
25:24 & \begin{tabularx}{0.7\textwidth}{X} 到了生產的日期，看哪，腹中是對雙胞胎。 \end{tabularx} \\ \\
[1ex]
\hline
\hline
\end{longtable}
$^{1}$很快有一天我們又看聖經《仁與士》.
然後記得生命樹為你呈獻經文.
認識背後的上帝.
果然是雙生.
昨天就說.
怎麼我的肚子裡面壓得這麼厲害.
神就說.
是兩個來的.
兒子們.
是兩個國來的.
裡面將來.
大的伺小的.
不是小的伺大的嗎?.
通常以前是小的伺大的.
是對的.
在那裡應該是這樣.
但是.
大的伺小的.
我命定了.
果然是雙生.
今天的題目.
生產的日子到了.
《創世記》五章二十四字.
福中果然是雙子.
嘩.
兩個的.
這裡好笑.
留意一下.
我們以前.
都聽過耳掃,啞角.
留意一下這段經文.
我研究裡面的字.
希伯來文的字.
原來有少少跟我們讀到的中文有些不同.
「先產的」即是哥哥.
哥哥是身體發紅.
通常我們說哥哥出來身體是紅的.
其實留意一下這個字是「發紅」.
發紅不是出來紅色.
嘩好紅.

$^{41}$是身體發紅.
Edmond Lee.
Edmond Lee是紅色.
但這個紅色.
不是紅色那種紅.
不是英文red那種紅.
不是紅色那種紅.
都可以說是紅色的.
但這種紅就是.
哎呀,你的臉怎麼這麼紅.
可以嗎?.
是膚色很紅潤.
是膚色紅潤的意思.
這個長者出來.
整個嬰兒的膚色很紅潤.
血氣很足.
看起來是個很健康的小朋友.
我相信耳掃也是.
很強健的.
嘩紅的.
很紅潤.
膚色很好.
血色很好.
另外「渾身有毛」.
「渾身有毛,如同皮衣」.
「貪污就給得起名叫耳掃」.
我解釋了這些字有點亂.
「身體發紅,渾身有毛,如同皮衣」.
「貪污就給得起名叫耳掃」.
「渾身有毛,如同皮衣」.
是聖經中文的話本.
原來在這裡.
「渾身有毛,如同皮衣」.
這個開頭的「有」.
藍色字有括號.
是表示中文字有.
但是原文.
希伯來文.
如果身上有希臘文.
是沒有的.

$^{81}$因此「渾身都是皮衣」.
這裡有括號.
「渾身皮衣」.
「衣」字應該是紅色的.
留意一下.
「渾身有毛,如同皮衣」.
是「渾身都是皮」.
「皮」和「衣」是分開的.
我們叫「皮衣皮衣」.
「皮衣褸」是一件.
原來用的字是「渾身皮」和「衣」.
不過,真是慘了.
這個「皮」字是譯得不對的.
「皮」其實不是「皮」.
是「毛」.
「渾身都是毛」.
不過用了「皮」這個字.
可能是「皮毛皮毛」.
其實是「渾身都是毛」.
「渾身都是毛」.
「皮衣」就是「毛衣」.
「衣」字是「榮耀」.
或者是外袍.
外袍,斗篷等等.
全身包著.
耳掃是怎樣的呢?.
「身體膚色十分紅潤」.
「而且渾身是全身都是毛」.
不過中文寫了「皮」字.
其實是「毛」.
「魚瞳」就在這裡出現了.
不過在這裡沒有提及.
其實是「魚瞳一件衣」.
「魚瞳是類似一件榮耀的外袍」.
即是說「毛好靚」.
好像穿了一件「毛衣」.
「毛衣」不是「油毛衣」.
你不要想錯.
是穿了一件「靴」.
穿了一件「皮草」.

$^{121}$這個小朋友全身紅潤.
而且渾身是全身都是毛.
而且毛好像穿了一件「毛球」.
「嘩!好靚啊!」.
於是就被稱為「耳掃」.
「耳紗」.
「耳紗」就是「多毛的」.
這是形容詞.
「耳紗」是多毛的.
這就是「耳掃」.
隨後又生了「耳掃」的兄弟.
「兄弟」即是一家人的「兄弟」.
「兄弟」這個字讓大家知道.
叫「Ark」.
「Ark」即是「兄弟」.
這個「兄弟」是手拿著「耳掃」的腳跟.
手拿著是「捉實」.
「抓住」.
「Ark」是「捉實」的腳跟.
「嘩!不放手」.
「Ark」是「捉住」的腳跟.
所以他們為「捉實」這個動作.
改了這個「兄弟」的名字.
叫他做「啞角」.
叫「啞角」,譬如普通話.
當時用普通話來做中文.
所以譯音應該是普通話的「啞角」.
「啞角」是三個音來的.
「啞角」和中文「啞角」有點相似.
不過「角」沒有中文字.
「啞角」,「啞」是對的.
「角」唯有譯成「啞角」.
「啞角」是抓住腳跟的人.
或者想「抽人家後腿」.
想做「取代者」.
「這個兒子捉住哥哥的腳跟出世」.
「當然想取代哥哥做大」.
「捉住腳跟不讓他出」.
就叫「取代者」.
「啞角」就是「捉住腳跟的人」.

$^{161}$為什麼要捉住哥哥的腳跟呢?.
因為想「抽人家後腿」.
想「取代他」.
所以有兩個兒子,「耳素」和「啞角」.
利伯加生下兩個兒子時.
耳素是六十歲.
剛剛六十歲.
六十歲開始成為大人.
開始壯年.
「啞角」的命名是對的.
從後面一直看下去.
「啞角」一生都在這個世界上.
不停地「抓」.
「抓」銀,「抓」貨,「抓」美女.
抓住,捉實,捉實.
甚至捉實天使不讓他走.
要給他祝福才可以讓他走.
捉實,捉實,捉實.
捉實拿了哥哥的獎指名分等等.
他真的沒有改錯名字.
所以他經常都心心不忿.
有些東西得不到.
「我要拿到」.
反而大哥「耳素」渾身事無.
十分強壯.
沒什麼所謂.
他為人大情大性.
「啞角」就「我要拿,我要取代」.
看下去《創世記金句》.
二十五章廿四節.
生產的日子到了.
福中果然是雙生.
一早就說了.
「哎呀,醫生醫生」.
「為何肚子嗝得這麼厲害?」.
「現在可以照到,驗到」.
「以前不懂就這樣了」.
「他會弱一點」.
「你吃些安靜的東西吧」.
「休息一下吧」.

$^{201}$神說「為何嗝得這麼厲害?」.
「他們在爭,是孖仔來的」.
「他們在爭將來大的要服侍小的」.
果然是雙生.
神說「這樣就這樣了」.
「明定」.
祈禱.
多謝我們的主.
你明定了每一個人.
明定了誰是誰的兄弟.
明定了誰是誰的子女.
明定了誰是誰的父母.
今天我們看到的是.
你為了尼伯加以撒.
預備了一對孖生的兒子.
在孖生的兒子裡面.
他們大家相爭.
正如今天我們和很多人相爭.
在家裡又相爭.
在教會裡又相爭.
在公司又相爭.
在學校又相爭.
在社會又相爭.
甚至我們搭車也相爭.
主要我們每樣東西都相爭.
可能這個爭奪.
要捉實就是我們人生的答案.
其實你早就知道.
幫助我們.
我們看到自己的影子.
看到雅各的時候.
好像看到自己是一個.
要捉實多點東西的人.
很多時候忘記了.
要捉實我們的主最為重要.
求你改變我們的心.
讓我們有更寬闊的心胸.
有你就夠了.
讚美你.
祈禱奉耶穌基督的名字.

$^{241}$阿們.
\newpage



\section{羅馬書 16:1-16}
\label{sec:031757}
\textbf{我們這一家•坐言起行}~~~~嚴增森弟兄
\newline
\newline
連結: \href{http://www.church.com.hk/acms/content.asp?site=cdc\&op=show\&type=product\&code=031757\&layout=sermon}{\texttt{031757}} ~~~~ 語音日期: 2026-03-01
\newline
\newline
\hyperref[sec:031756]{\small{< < < PREV SERMON < < <}}
~
\hyperref[sec:index_chronic]{\small{[返順時目]}}
~
\hyperref[sec:031758]{\small{> > > NEXT SERMON > > >}}
\newline
\newline
羅馬書 16:1-16
\newline
\begin{longtable}{cl}
\hline
\hline
章節 & 經文 (和合本修訂版)\\
\hline
16:1 & \begin{tabularx}{0.7\textwidth}{X} 我對你們推薦我們的姊妹非比，她是堅革哩教會中的執事。 \end{tabularx} \\ \\ \relax
16:2 & \begin{tabularx}{0.7\textwidth}{X} 請你們在主裡用合乎聖徒的方式來接待她。她在任何事上需要你們幫助，你們就幫助她；因她素來幫助許多人，也幫助了我。 \end{tabularx} \\ \\ \relax
16:3 & \begin{tabularx}{0.7\textwidth}{X} 請向百基拉和亞居拉問安。他們在基督耶穌裡作我的同工， \end{tabularx} \\ \\ \relax
16:4 & \begin{tabularx}{0.7\textwidth}{X} 也為我的性命把自己的生死置之度外；不但我感謝他們，就是外邦的眾教會也感謝他們。 \end{tabularx} \\ \\ \relax
16:5 & \begin{tabularx}{0.7\textwidth}{X} 又向在他們家中的教會問安。向我所親愛的以拜尼土問安，他是亞細亞歸於基督的初結果子。 \end{tabularx} \\ \\ \relax
16:6 & \begin{tabularx}{0.7\textwidth}{X} 又向馬利亞問安，她為你們非常辛勞。 \end{tabularx} \\ \\ \relax
16:7 & \begin{tabularx}{0.7\textwidth}{X} 又向與我一同坐監的親戚安多尼古和猶尼亞問安，他們在使徒中是有名望的，也是比我先在基督裡的。 \end{tabularx} \\ \\ \relax
16:8 & \begin{tabularx}{0.7\textwidth}{X} 又向我在主裡面所親愛的暗伯利問安。 \end{tabularx} \\ \\ \relax
16:9 & \begin{tabularx}{0.7\textwidth}{X} 又向我們在基督裡的同工耳巴奴和我所親愛的士大古問安。 \end{tabularx} \\ \\ \relax
16:10 & \begin{tabularx}{0.7\textwidth}{X} 又向在基督裡經過考驗的亞比利問安。向亞利多布家裡的人問安。 \end{tabularx} \\ \\ \relax
16:11 & \begin{tabularx}{0.7\textwidth}{X} 又向我親戚希羅天問安。向拿其數家在主裡的人問安。 \end{tabularx} \\ \\ \relax
16:12 & \begin{tabularx}{0.7\textwidth}{X} 又向為主辛勞的土非拿和土富撒問安。向所親愛、為主非常辛勞的彼息問安。 \end{tabularx} \\ \\ \relax
16:13 & \begin{tabularx}{0.7\textwidth}{X} 又向在主裡蒙揀選的魯孚和他母親問安，他的母親就是我的母親。 \end{tabularx} \\ \\ \relax
16:14 & \begin{tabularx}{0.7\textwidth}{X} 又向亞遜其土、弗勒干、黑米、八羅巴、黑馬，和跟他們在一起的弟兄們問安。 \end{tabularx} \\ \\ \relax
16:15 & \begin{tabularx}{0.7\textwidth}{X} 又向非羅羅古和猶利亞，尼利亞和他姊妹，阿林巴和跟他們在一起的眾聖徒問安。 \end{tabularx} \\ \\ \relax
16:16 & \begin{tabularx}{0.7\textwidth}{X} 你們要以聖潔的吻彼此問安。基督的眾教會都向你們問安！ \end{tabularx} \\ \\
[1ex]
\hline
\hline
\end{longtable}
$^{1}$大家平安.
我一直在想 大家早上很清楚了.
現在又過了十二點 午安好像太晚了.
多謝Julia幫我們讀一至十六節.
很不容易讀的經文 很多個字 很多個名字.
不知道大家讀這段經文時會否跳過它 跳過了它.
兩年前我選了這段經文在晚課中分享.
如果你有聽過 如果你還記得 那就辛苦你了.
相信是Julia透過她的話語對我們每一個人說話.
讓我們一起祈禱.
Deeper 我多謝你讓我們一起敬拜親近你.
在這一刻我們要去聽你的話語 求主你的靈幫助我們.
讓我們能夠帶領我們明白你的話語 明白你的心意.
主啊 你知道說話不配戴你的話語而帶著能力.
求主你藉著你的話語幫助我們.
因為我們不明白你的心意的時候 要讓我們的生命能夠被更新.
主啊 我們懇求你在這一刻與我們同在.
讓我們經歷你.
多謝你聽我們祈禱奉耶穌基督的聖名.
今日是羅馬書第39講.
不知道大家有沒有留意程序表 第39講.
我一直都是晚課 所以之前聽得一點也不多.
但大家聽了第38講的時候.
不知道大家對羅馬教會的認識有多少呢.
或者可能你回到38個星期的窩福 你對窩福這個教會的認識有多少呢.
可能38個月 甚至38年也不一定.
我這裡看一下 有38年的.
究竟我們對窩福這個群體的認識有多少呢.
今天我給這個主題 我們這一家.
我不知道大家看到我們這一家 你會想起些什麼東西.
會不會是法斯拉的一家.
還是窩福的一家呢.
除了我們三堂崇拜的弟兄姊妹.
其實我們星期天還有恩廟的群體.
有時他們會在我們當中與我們一起敬拜.
還有旺角的The Cave 不知道大家有沒有去過.
沒有去過的話 可能你已經去過了.
The Cave 同樣在星期天 主要是有崇拜的.
有青少年的崇拜 新天新地.
為了讓大家認識他們多一點.

$^{41}$我就找了這班少年人 為了他們的崇拜 定下了崇拜守則.
和大家分享一下.
這些字很小的 厭硬的.
當中說到手機 要將手機放入書包.
我不知道大家的手機怎麼處理.
應該教他們靜音.
還有他們說可以吃東西.
不過唱詩祈禱的時候 讀經就不可以吃.
這裡我們是不可以吃東西的.
他們說讀主題經文的時候 要一起讀 要投入.
今天就不需要大家讀了 太難讀了.
第四個就很精彩了.
盡量不要光明正大地打哭鬧.
我不介意的 大家可以打哭鬧.
吸多一點氧氣 可以精神一點.
他說遲到不要超過十分鐘.
我們是大人 所以我們不要遲到 我們要準時.
你看電影都不會遲到的.
沒有理由上帝在 你會遲到的.
另外他說要積極參與.
不要太吵 可能他們會聊天.
不要騷擾人 要尊重 要彼此尊重.
還有他說要聆聽他人的話.
他的意思是講道 講道的時候要聽.
所以就請大家聽.
裡面這些少年人的崇拜手法.
可能有一些都值得我們學習.
這裡就是窩福這個家的大大小小.
接下來我們要看的是羅馬教會的一個家.
保羅在第十五章 即是上一節第三十八講.
在上一節第十五章結尾的時候.
他交代了他去耶路撒冷的行程.
亦都交代了他計劃去經過羅馬.
然後再去西班牙傳道.
最後他就加了一句祝福.
總結了他的書信的內容.
來到第十六章最後一章.
我們看到充滿了很多的名字.
首先在第十六章的第一節 二節.
保羅就去推薦菲比這位姊妹.

$^{81}$保羅用了兩個銜頭去形容這位姊妹.
她是堅格利教會的女執事.
堅格利是哥林多這個地方.
是一個東面的口岸城市.
現在菲比來到羅馬這個地方.
相信她都是將保羅這封信.
羅馬書這封信帶給羅馬教會.
大概菲比是一個又忠心又可靠的領袖.
她要去羅馬處理事情.
於是由她帶這封信去羅馬教會.
就是最佳的人選.
銜頭保羅說她是贊助人.
其實第十六章二節最後一句.
應該是願文應該這樣翻譯.
她說她是很多人和保羅的贊助人.
贊助人這個字當時是一個特別的名詞.
就是指那些在物質上.
和用她在社會上的地位.
去幫助其他人的那些人.
可能保羅在哥林多的時候.
菲比就是這樣幫助她.
支持她傳道的工作.
所以保羅就請羅馬教會去接待和幫助.
來到羅馬的菲比.
首先我們看到菲比不是單單送信人.
她也不是單單需要別人接待幫助的人.
當我們想到在十五章剛才提到.
保羅很希望羅馬教會成為.
她去西班牙傳道的幫助者和基地.
就好像安提拉就是保羅去傳道的時候的基地.
在這樣的背景的時候.
而菲比的出現就成為一個很好的見證和榜樣.
在二千年前就有這些貴族.
這些有錢有頭有臉的信徒.
他們出錢出力熱心的去服侍.
大概今天我們不一定要是貴族.
也不一定要很有錢.
可能當中有些你是.
但不一定是這樣的.
也可以成為別人的贊助人.

$^{121}$你可以想想怎樣成為現代的菲比.
或者成為窩福的菲比.
又或者你是接待和幫助菲比的人.
在哥林多和保羅一起服侍就不止只有菲比.
其實在二十一到二十三節還有一些名單.
不過就留到第四十講 留到下一次表傳道講.
而我們就來到第三到第十五節那些名單.
可能你會問保羅書信裡面最後文案的名單有什麼特別有什麼好講呢.
其實在保羅的書信裡面.
十三卷書信裡面只有三卷會提名去問安.
在歌羅西書裡面提名一位靈法.
在提摩塔後書裡面就有三位.
但在羅馬書裡面我們看到.
剛才有沒有數過.
26個名 是26個.
是不是很特別.
那麼特別一定有些東西值得我們去研究.
所以不要睡著了 先聽完.
其實這個名單給我們看到羅馬教會的組成.
在這裡我們可以看到有不同的人的種族文化.
和有不同的社會階層.
在這裡我們看到名字紅色代表女性.
首先在十六章三節就提到一對猶太人的夫婦.
伯基拉和雅居拉.
可能大家比較有印象.
在史道行傳第十八章就提到這對夫婦.
由羅馬去到哥林多.
在哥林多和保羅相遇.
和保羅一起工作.
不是只是傳道 是工作.
當然也有一起傳道 傳福音.
之後他們一起去了以弗所.
然後又回到羅馬.
而雅居拉這個字是一個拉丁名字.
可能顯示了他們在羅馬生活了很多年.
又或者已經融入了羅馬的文化.
就好像今天這樣.
很多人都有英文名字.
或者有一個比較特別的.
有一個法文名字 或者是日文名字.

$^{161}$我們也看到有些外國人.
可能也會說一個中文名字.
來到香港或中國這個地方工作.
而這裡所問候的對象.
還有在他們家裡聚會的弟兄姊妹.
當時在羅馬的情況.
就不像今天這樣.
我們有這麼大的地方.
可以同一時間坐這麼多弟兄姊妹一起聚會.
所以羅馬的信徒.
基本上他們會在不同的家庭裡.
去聚會.
而可以提供這些地方的人.
可能都是比較富裕的.
就像今天可能你的小組團.
去某個弟兄姊妹家裡聚會.
你可能也會選一個比較大的地方.
可能他們相對比較富裕.
而雅居拉夫婦他們是做帳棚生意的.
所以他們應該比較富裕.
也有條件移居到其他地方.
除了這對猶太人的夫婦.
在第七節的安多尼古和尤利亞.
他們都是猶太人.
因為他們是保羅的親戚.
而這兩個名字分別是希臘名和拉丁名.
到了第十一節的希羅天.
意思是屬於希律的人.
和希律有關.
但他們都是保羅的親戚.
都是猶太人.
其餘可能是猶太人的.
包括了第六節的瑪利亞.
瑪利亞是很普遍的猶太人的名字.
但原來它也是很普遍的希臘名.
所以我們不能百分百肯定.
是不是猶太人.
接著在第十節的阿比利.
是當時在羅馬.
猶太人經常會用的一個希臘名.

$^{201}$所以很可能也是猶太人.
而十三節的魯夫.
這是一個拉丁名字.
但這個魯夫可能是馬可福音的魯夫.
即是古利內人西門的兒子.
古利內是北非的地方.
其實有很多猶太人聚居.
加上西門是一個很典型的猶太名.
所以應該也是猶太人.
以上所說的.
其餘.
除了以上所說的.
其餘的都是外邦人.
有拉丁名字.
有希臘名字.
而第十二節的比昔.
其實是一個波斯名字.
而且大部分這些名字.
都是當時奴隸的名字.
另外也有一些和帝國皇室有關係.
因為當時的皇室的人.
他們幫他們的奴隸改名字.
也是用不同的名字.
所以當我們看到這些名字的時候.
我們就可以分得出.
哪些原來和當時帝國皇室有關係.
例如在第十六章九至十一節.
和十五節裡面提到的人.
相信也是和當時的皇族有關.
在三至十五節的名單裡面.
我們可以看到羅馬教會的組成是很多元的.
包括有不同的文化種族.
有猶太人.
而且更多的是外邦人.
也有女性.
當中也有些是希臘化.
也有些是被羅馬的文化影響.
也有不同的階層.
好像剛才提到的.
白居拉和阿基拉做生意的人.

$^{241}$也更多的是低下階層的人.
特別當時奴隸是很普遍的.
在羅馬有三分二的人口是奴隸.
當大家想起奴隸.
其實不是大家所想的那些.
很苦力的那些.
當時的奴隸是做什麼的呢?.
教師 會計師 醫生.
他們都是被俘虜成為奴隸.
我不知道我們當中有沒有教師.
有專業人士.
有醫生 會計師.
恭喜你 你可以做奴隸.
另外在這十幾節的經文裡.
我們找到大概有五個家庭的教會.
有五個地方被弟兄姊妹聚集.
一起去敬拜.
這裡其實只是一個剪影.
雖然不可以完全代表.
當時羅馬教會信徒群體的情況.
但是我們看到一幅圖畫.
有關羅馬教會弟兄姊妹的圖畫.
我們可以想像一下.
這群不單單是基督徒.
他們同時是猶太人.
也有外邦人.
猶太人可能有不同的派別.
好像耶穌的門徒當中.
有憤慮黨的西門.
又有親建制的瑞利馬太.
外邦人就更加來自不同的文化.
不同的種族背景.
還有來自不同社會的階層.
這裡我們可以看到很立體.
看到當時的信仰群體的情況.
比我們今天的教會更加多元.
但是大家可以想像.
可能產生的困難問題.
也會更加複雜.
除了羅馬教會的組成.

$^{281}$也讓我們了解一下羅馬教會的處境.
羅馬是位於帝國的中心.
所謂天子腳下.
被教會帶來不同的挑戰.
首先是猶太人和羅馬人.
以至和帝國的關係.
其實自從耶穌的年代.
到保羅寫羅馬書的時候.
當中出現了兩次猶太人被驅逐離開羅馬城.
史諾恆第十八章也記載過.
在格魯吊·克羅迪亞做皇帝的時候.
猶太人被逐出羅馬城.
所以伯吉拉和雅居拉就去到哥林多.
歷史告訴我們這件事和基督徒有關.
可能就是當時猶太人和基督徒發生了衝突.
以致所有猶太人都被驅趕.
另外在帝國的眼中.
基督教和其他宗教很不同.
其實當時有很多不同的種族.
他們都有自己的宗教.
這是帝國所容許的.
而且這些宗教基本上都是多神的宗教.
即是當他們面對要敬拜君王的時候.
他們只會把君王的像多放一個.
對他們來說沒有什麼大挑戰.
但是對猶太教來說.
因為他們是一神宗教.
所以對他們來說就很不同.
不過帝國很多時候都豁免他們.
讓他們不需要去敬拜君王.
而基督教初期的時候.
也是被視為猶太教的一個分支.
所以他們都可以被豁免.
不過後來當基督徒和猶太人的關係緊張的時候.
於是這個豁免就沒有了.
而且基督教會和其他的宗教很不同.
其他的宗教都是單一的種族.
但是基督教會裡面.
剛才讀過很多不同種族的人.
而且還有很多奴隸.

$^{321}$於是你可以想像到.
當時的政權見到這樣一個群體的時候.
這個不像宗教.
這個是政治.
所以我們可以想像到當時的教會.
他們面對來自各方各面的挑戰和壓力.
剛才提到猶太人被逐出羅馬城.
包括猶太的基督徒.
一個由猶太基督徒建立的教會.
突然間那些猶太基督徒離開的時候.
那些外邦的信徒.
他們就需要承擔所有的角色和責任.
久而久之.
他們的習慣各方面就形成了去猶太化.
但是去到格魯迪皇帝死了之後.
這個禁令就取消了.
猶太人又再返回羅馬城.
包括猶太的基督徒.
你可以想像當時的情況帶來一些衝擊.
因為有很多的改變.
在之前大家還聽到的第十四十五章.
講到因為猶太人的飲食和節期的問題.
產生了很多的衝突.
其實當中不只是影響信徒關係合一這麼簡單.
而且這個是福音的根本.
在這樣的情況之下.
就影響這個群體不能聚焦去承擔福音的使命.
這就是保羅寫羅馬書的原因.
羅馬書不是一個純粹講教義的書.
而是針對當時羅馬教會處境的一本書.
為什麼要講這麼多之前我們所謂的教義.
其實是要讓讀者知道.
福音是怎麼樣的.
不是靠自己.
不可以自疑.
我們沒有一個人可以誇口.
愛邦人對著猶太人不可以誇口.
猶太人對著愛邦人其實都不可以誇口.
在這樣的情況之下他們才可以面對處理這些問題.
今天我們讀這段經文其實不只是圍繞這些名單.

$^{361}$當中保羅講到是關於民安.
經文一開始就說請向某某民安.
但第一件事我們要弄清楚.
究竟是誰去民安呢?.
民安這個動詞其實是用了第二身的訟數.
也就是說不是保羅第一身去民安.
保羅是請那些信徒.
羅馬書的讀者去民安.
是不是只是替保羅代他去民安呢?.
又或者只是顯示保羅在羅馬教的人物關係呢?.
應該沒有這麼簡單.
接著我們會繼續看下去.
我們解決了誰是民安.
原來是那些信徒.
民安的對象.
這裡不只是一些名字.
而是當中保羅對很多提名的信徒都有介紹和描述.
其中最重要的就是伯吉拉和阿居拉.
他們是保羅的同工.
甚至可以為著保羅犧牲自己的性命.
可想而知他們的關係有多密切.
但值得我們留意的不只是這樣.
還有就是保羅說.
大概是因為這對夫婦對外邦教會的貢獻.
大家有沒有留意.
保羅在這裡似乎刻意將猶太的基督徒和外邦的教會.
外邦的基督徒連在一起.
另外和保羅一起同工甚至坐牢的.
有名望的使徒安多尼古和尤利亞.
還有和他同工的伊班奴.
至於瑪利亞,土斐娜,土庫什和比瑟.
就形容他們是辛勞的姊妹.
甚至是非常辛勞.
另外也有幾位保羅稱他們是親愛的.
第五節的以拜尼圖.
就是屬於亞細亞當時最早信主的那一批.
第三節提到的老夫的母親.
保羅當她是自己的母親一樣.
大概她曾經照顧過保羅.
或者在保羅的工作中支持過他.

$^{401}$這些都是保羅個別認識的一些信徒.
但不只是他們.
還包括了五個家庭教會.
保羅也請讀者向他們問案.
保羅也記掛著他們.
這裡所說的除了保羅和他們的關係.
還有保羅如何欣賞他們.
也指出可能當時羅馬的信徒.
他們彼此未必很認識.
特別是來自不同家庭教會的成員.
看到這麼多名單.
我們可以思考一下.
保羅用什麼眼光,用什麼角度去看他們的肢體呢?.
而保羅的舉薦對讀者有什麼作用呢?.
如果你是羅馬教會的一員.
你聽到保羅這樣形容這些弟兄姊妹的時候.
你對他們會有什麼印象呢?.
會不會覺得更容易去接觸他們呢?.
雖然他們未必是你熟悉甚至你認識的.
又或者會不會本來你對他們有一些偏見.
不過聽完保羅所說的時候.
你會多一份尊重甚至多一份尊敬.
而可以去接觸他們呢?.
讀了這段經文的時候.
我就想到我會怎樣去形容.
我自己團體小組的弟兄姊妹呢?.
又或者跟我一起去事奉的弟兄姊妹.
我會怎樣去想他們,怎樣去形容他們呢?.
你又會怎樣去看你認識的弟兄姊妹呢?.
香港人是很厲害的.
很critical的.
整塊白板你看不到白色.
你只看到黑色的一點.
我們怎樣去看呢?.
求上帝幫助我們都懂得去欣賞.
相信剛才那二十六個名字.
他們不是聖人.
他們不是完全人.
他們的「人」沒有一個是.
但是保羅看見的.

$^{441}$是不是和我們看見的不同呢?.
在第十六章十六節.
保羅最後總結了一句.
他說你們要以聖潔的吻彼此問安.
這裡是彼此問安.
不是代保羅問安.
是羅馬教會他們要彼此的問安.
親吻問安是按照當時的文化和習慣.
可能今天還有些地方也是這樣.
我們也可能看見過.
但是我們的習慣可能就是握手.
這裡提到為什麼會是聖潔呢?.
保羅這裡不是要講衛生.
保羅也不是在講道德.
聖潔的意思是分別出來為上帝.
怎樣為上帝問安呢?.
我自己有一個經歷.
在兩年前.
大概是這個時間.
崇拜完了之後.
離開了.
在走廊.
遇到一個我認識的姐妹.
但不是經常交往.
但她突然停下來.
跟我打招呼問候我.
她想了解我的近況.
我感受到她非常真誠.
她問我有什麼可以為我祈禱.
我看到她如此真誠.
我就很認真地告訴她.
我有什麼需要去祈禱.
大概.
聖潔的問安.
是要用心和時間.
不只是Hi and Bye.
或者我們經常見到.
吃飯了嗎?.
人家還沒回答.
你已經走了.

$^{481}$可能你會問.
問安而已.
何必說這麼多.
是不是想多了?.
正如我們開始時提到.
保羅書信裡的問安.
這裡是很特別的.
保羅要強調的是要建立.
羅馬信徒當中.
那個群體彼此的關係.
而問安可能是實踐.
保羅教導的第一步.
破冰的第一步.
實際上.
彼此這個字.
在第十二章開始.
出現了十次.
完完全全一樣.
這個字出現了十次.
同一個字根多了一次.
十二章的五節.
就好像大家看到的那張slide.
開始是說我們彼此為肢體.
所以保羅說我們要彼此相親相愛.
我們要彼此真心去尊重人.
尊重對方.
我們要彼此同心.
再一次說我們要彼此相愛.
我們不去彼此評斷.
我們要彼此造就.
我們要彼此同心.
效法基督耶穌.
我們要彼此接納.
彼此勸戒.
今天我們說到彼此問安.
是不是偶然的呢?.
用同一個字出現了十次.
大概保羅告訴我們.
我們要彼此去建立.
而彼此問安就是第一步.

$^{521}$也是需要繼續去做.
去維繫和見證我們合一的關係.
為何保羅如此緊張羅馬教會要彼此接納.
要同心合一呢?.
第一就是要去榮耀上主.
好像第十五章五至六節.
之前也提及過.
你們要彼此同心.
效法基督耶穌.
使你們可以同心同聲.
榮耀主耶穌基督的父上帝.
大家有沒有想過.
同心同聲榮耀上主是怎樣的呢?.
今天我們一起來崇拜.
你敢不敢受到我們同心同聲呢?.
還是我們只是來到同一個地方.
不過我們只是個別去敬拜呢?.
彼此接納同心合一.
第二個目的.
是要承擔和實踐福音使命.
之前也提及過.
第十五章保羅打算去西班牙傳福音.
他希望羅馬教會能夠成為.
他去西班牙傳福音的一個基地和幫助.
但是要羅馬教會準備好承擔這個使命之前.
就必須處理很多他們當中的紛爭.
一些不和,一些嚴規.
用今天的說話來說.
可能就是撕裂.
彼此建立同心合一.
我們才會有空間去聚焦.
去承擔,去實踐上帝給我們福音的使命.
每一間教會都有很多問題.
窩福也有很多問題.
你常常聽到我不知道是什麼.
不過最近教會的問題.
其中一個問題是廁所.
窩福的廁所有很多問題.
常常都要整.
當整好之後.

$^{561}$你一定可以感受到.
因為你完了崇拜之後.
發覺要排隊.
我們感恩的是.
我們沒有因為對整廁所不同的意見.
而影響到我們的關係.
但是如果有一個教會.
他為了這些事情.
設施各方面.
卻忘記了要聚焦於傳福音.
那麼這是不是理想的教會?.
更重要的是.
這是不是上帝心意的教會?.
我們總結一下.
今天我們讀了一個經文.
讓我們看到羅馬教會的組成.
是由不同種族文化.
不同社會階層的信徒組成.
他們帶著不同的背景.
不同的習慣.
經歷.
而成為基督的門徒.
在相處的時候.
難免會有一些張力.
有些緊張.
保羅要他們學習彼此接納.
來到最後的時候.
保羅說你們要開始彼此問案.
接著就要建立那九個彼此.
目的就是要榮耀上主.
要承擔使命去見證福音.
今天這段經文對我們有什麼意義呢?.
希望大家不要只記得.
二十多個名字.
聽過一次.
以後可以跳過.
不是的.
我們要想想.
我們可以怎樣應用和學習呢?.
今天香港教會包括窩福.

$^{601}$其實在過往的日子.
儘管我們未必去到撕裂的情況.
但我們確實經歷了關係上的緊張.
之後的移民潮.
無論是教會領袖.
或是信徒的離開.
這一切的衝擊和改變.
我們不能當作沒有發生過.
就像羅馬教會所面對的.
雖然我們的處境不是完全一樣.
但我們同樣可以學習.
和回應保羅的呼籲和教導.
其實保羅不是呼籲.
他是命令.
我們要積極的建立彼此.
可以由民安開始.
既然我們現在少了很多人.
大家是否感受到?.
我們更加要把握這個時機.
去穩固這個基礎.
成為真正同心合一的群體.
做到同心同聲去榮耀上主.
同時間我們求上帝.
讓我們有一個更加廣闊的心.
去接觸和接納跟我們不同的人.
無論是還留在香港的.
或是剛剛來到香港的.
他們跟我們不同.
但一樣的是他們都是需要福音的人.
求上帝讓我們看到.
那些莊稼已經在門口.
莊稼熟了是會壞的.
我們需要積極的去承擔.
這個時代的福音使命.
還有.
求上帝讓我們利用愛.
去接納一些可能將來會回流的弟兄姊妹.
英國有一位牧者.
在一個很大的環教會.
在倫敦.

$^{641}$他估計未來可能有20-25\%.
香港去到的弟兄姊妹會回流.
我們是否準備好去接待他們呢?.
去歡迎他們呢?.
幫助他們再次融入這個群體.
羅馬教會的合一.
還有準備他們面對這個時代很快到來的逼迫.
可能對我們來說都是需要準備.
所以就讓我們這個群體成為一個家.
去建立我們的彼此.
以致我們更加能夠去準備好.
去承擔這個福音的使命.
以及能夠預備好迎接接下來的挑戰.
初期教會.
他們彼此親嘴問安.
其實就成為了他們在聚會時候的第一個環節.
在今天我們說到尾聲的時候.
我都很希望去邀請大家.
我們坐言起行.
我邀請大家.
其實剛才早堂就坐滿了.
就比較困難.
現在就有點空檔.
我邀請大家方便的話.
你可以稍而欲報.
你去看看身邊我們互為之體.
彼此為之體的.
你跟他彼此問安好不好.
你不需要跟他親嘴問安.
你可以跟他握手.
你可以跟他激掌激拳.
甚至有些你可以跟他擁抱.
大家不要吝嗇.
多謝大家給予反應.
如果大家沒有反應.
我就會很尷尬.
但這只是很少很少的熱身.
最後我們不要忘記.
我們同樣可以成為.
或者幫助現代的非比.

$^{681}$讓我們一起祈禱.
在這個時候我們稍為安靜的時間.
我們回想這段經文.
究竟上帝對我們說什麼話呢.
我們要怎樣去回應呢.
我們一起祈禱.
我們的一家之主.
天父上帝.
這一家是屬於你的.
除了我們求你幫助我們去學習.
保羅對其他肢體那份的欣賞.
和讓我們能夠實踐.
保羅的教導.
其實這就是你的教導.
你要我們彼此相愛.
不去彼此的評斷.
你要我們彼此去造就.
彼此同心的效法基督.
幫助我們能夠彼此接納.
彼此勸戒.
和讓我們懂得開始彼此問安.
是聖潔的問安.
最後幫助我們.
因為我們彼此是遺肢體.
最後讓我們去建立這份彼此.
讓我們準備好去承擔.
去實踐.
求主你擺在我們面前的福音使命.
幫助我們能夠走出去傳遞愛.
讓我們將主你給我們的生命活出來.
我們去見證你.
又讓我們有預備迎接將來的挑戰.
主你將我們這個群體.
這個家交在你的手中.
求主你繼續與我們說話.
多謝你聽我們祈禱奉耶穌基督的聖名.
\newpage



\section{路加福音 15:4-7}
\label{sec:031758}
\textbf{一比九十九}~~~~姚永光傳道
\newline
\newline
連結: \href{http://www.church.com.hk/acms/content.asp?site=cdc\&op=show\&type=product\&code=031758\&layout=sermon}{\texttt{031758}} ~~~~ 語音日期: 2025-08-17
\newline
\newline
\hyperref[sec:031757]{\small{< < < PREV SERMON < < <}}
~
\hyperref[sec:index_chronic]{\small{[返順時目]}}
~
\hyperref[sec:031759]{\small{> > > NEXT SERMON > > >}}
\newline
\newline
路加福音 15:4-7
\newline
\begin{longtable}{cl}
\hline
\hline
章節 & 經文 (和合本修訂版)\\
\hline
15:4 & \begin{tabularx}{0.7\textwidth}{X} 「你們中間誰有一百隻羊，失去其中的一隻，不把這九十九隻留在曠野，去找那失去的羊，直到找著呢？ \end{tabularx} \\ \\ \relax
15:5 & \begin{tabularx}{0.7\textwidth}{X} 找到了，他就歡歡喜喜地把羊扛在肩上。 \end{tabularx} \\ \\ \relax
15:6 & \begin{tabularx}{0.7\textwidth}{X} 他回到家裡，請朋友和鄰舍來，對他們說：『你們和我一同歡喜吧，我失去的羊已經找到了！』 \end{tabularx} \\ \\ \relax
15:7 & \begin{tabularx}{0.7\textwidth}{X} 我告訴你們，一個罪人悔改，在天上也要這樣為他歡喜，比為九十九個不用悔改的義人歡喜還大呢！」 \end{tabularx} \\ \\
[1ex]
\hline
\hline
\end{longtable}
$^{1}$《錄音》或《錄影》的版權是屬於基督教宣導會不各堂所有.
我們的屬靈傳統之一是以真理教導為本.
歡迎您將此《錄音》或《錄影》的連結傳給親友.
懇請注意,未經授權請勿轉錄或剪輯本《錄音》或《錄影》.
甚願好似《使徒行傳》二十章三十二節下所講.
這道能建立你們,叫你們和一切聖聖的人同得基業.
主耶穌很愛我們,以致到我們今日.
無論我們在哪一個處境,我們在任何背景.
我們今天能夠一齊聚集去敬拜,去讚美神.
又分享神的話語,都不是偶然.
都是出於主耶穌一份很深的愛.
今日和大家分享的經文是聖經新約路加福音第十五章四到七節.
是主耶穌講的一個比喻.
這個故事很簡潔,講到有一個主人.
他有一百隻羊,雖然走失了一隻.
他仍然會去尋找.
不過他這個尋找很特別.
他的尋找是要直到,找到.
並且將這隻羊帶回到家裡面.
他才有那份喜樂.
今日嘗試從聖經當中看看主耶穌.
他怎樣看一和九十九之間的關係.
羊有一個特性.
羊是需要牧放的.
羊是需要給牠有一個空間.
牠去覓食,牠可以自由活動.
不可以常常困住牠.
牠是需要有空間.
牠是需要有一個活動.
牠是需要有自由.
不過羊同時間又很容易走失.
羊又很容易被一些事情.
或者被一些吸引.
甚至乎被其他的東西.
引致牠偏離了牠原本應該要走的方向.
羊是這樣的.
你不可以困住牠.
是要給牠有自由的空間.
不過羊又有一個特性.
牠很容易會不自覺地走失.

$^{41}$不過在經文中說.
一切都是由第四節的第一句開始.
是的,縱然主人有一百隻羊.
不過牠走失了一隻.
剛才我們也說過.
牠仍然會先放下牠所擁有的.
是要找回那一隻.
是要織牠找到.
又將牠帶回家裡.
一切都是由第一句開始.
誰有一百隻羊.
這個「有」字來得很寶貴.
這一百隻羊是有一個主人的.
包括那隻走失了的.
我們不知道是什麼原因.
不知道牠為什麼會走失了.
不過牠縱然走失.
牠是沒有被忘記的.
牠仍然是屬於這一位主人.
並且這個主人用盡方法去尋找牠.
這個尋找不是等牠回來.
不是看看牠遲些會不會自己找回來.
不是.
這一隻走失的羊.
是,牠走失了.
不過牠沒有被忘記.
牠仍然是屬於有一個主人.
這個主人很著緊牠.
無論如何都要找回牠為止.
找到要直到找到牠為止.
大家認不認得這部是什麼機器.
果然用過.
是入票機.
我記得很多年前.
我去到銀行的位置.
有幾張票.
我打算去入.
我對於這方面反應比較慢.
我對於入卡入票手腳比較慢.
有一次吃午餐的時候.

$^{81}$趕著拿著幾張票.
其中要入一張.
排隊了.
到我的時候.
我比較笨拙.
手腳比較慢.
在猶豫之間.
我已經聽到後面有聲音.
是催促的聲音.
你知道已經是手腳笨.
所以有些壓力.
更加淪盡.
在電光火石之間.
我拿著一張票.
大家知不知道發生什麼事.
很聰明.
希望那個不是你.
竟然後面那位幫我入到他的位置.
說時遲那時快.
不過好彩.
那張是我要入的.
其實我講的是什麼.
我們現在身處的環境.
我們是很急的.
我們是不能等的.
不要說不能等.
我們停下來也不願意.
時間是很寶貴的.
不過這個主人.
那一百隻羊.
所有都是屬於他的.
就算一隻走失了.
他仍然願意先放下.
他所有的.
無論如何都要花時間.
花心思.
花力氣.
去尋找那一隻迷失了的.
不是等他回來.
不是坐著等他.

$^{121}$是透過身邊很多的人與事.
都要將他帶回來這個家當中.
是要直到他回到這個家當中.
這一份付出.
願意付出時間.
心思.
和力量.
是這個主人對一百隻羊的關心.
找到了.
他們歡歡喜喜.
將羊肝在劍上.
找到了之後.
很特別.
他找到這隻羊的時候.
他是將這隻羊放在肝在劍上.
肝在劍上這個字很特別.
這個肝在劍上.
肝字所指的是.
是用雙手要抬起它.
是很重的.
是用雙手抬起它.
是很重的.
這個主人見到這隻走失了的羊.
他即時做的是.
將它抱起.
放在他的肩上.
是很重的.
他所指的重不是這隻羊的重.
他將它放在他的肩上的意思是.
將這隻走失了的羊.
是他有的重擔.
他曾經背負的一些重擔.
甚至很多他纏累著他.
他仍然不能夠放下的一些負累.
又或者一些很積壓著他.
令他不能夠自己的一些事情.
這些的重.
主人一見到他的時候.
沒有任何說話.
就是將他抱起.

$^{161}$放在肩上.
這隻羊從此就有新的開始.
所指新的開始.
他不知道什麼原因走失.
或者他是貪玩.
或者他被其他東西吸引了.
或者他有很多難言之隱.
在走失過程他經歷過的種種.
他所做過的種種.
這個主人不需要他講一句話.
就一手將他抱起.
整個的重量都放在他的肩上.
這個小的重擔.
從此就可以交在這個主人當中.
他有新的開始.
我記得我兩個兒子小的時候.
我們經常會和他們去樂園玩.
大家知道什麼是樂園嗎?.
就是真的很大的那些.
有機動遊戲那些.
你知道他們小時候.
去樂園我們通常都會和他們約法三章.
就是講明去一定要聽話.
最重要就是要跟著爸爸媽媽.
不要走失.
每一次都要講.
講三次因為重要.
每次都是這樣講.
你覺得我兩個兒子聽不聽?.
聽不聽我講?.
當然聽 不過不做.
不聽就沒得去.
講到我都未曾講第二句.
他已經說知道了.
要跟著要拖著.
要跟著哥哥.
哥哥又要看著弟弟.
又要跟著爸爸媽媽.
不要走失.
我都未開口講第二句.

$^{201}$已經背了出來.
有一次我們去一個很大的樂園.
還要在外地.
一樣是這樣講.
不過在開始玩之前.
出發之前在樂園裡面.
出發去玩之前.
通常我們都會和兩個小朋友講.
記住一定要跟著哥哥.
哥哥一定要看著弟弟.
你們一定要在爸爸媽媽身邊.
不要走失.
走失了很大件事.
不過你如果走失的話.
通常我們會找一個地標.
通常都是最高的建築物.
或者最能夠看得到的.
我們就會和我兩個兒子講.
你如果真的走失了.
你認定這個位置.
我們一定在這裡等你.
最重要你跳得回來.
如果你真的走失了.
你離開這個位置.
我們一定在這裡等你.
有一次去了一個很大的樂園.
一開始玩都沒有半個小時.
終於就走失了.
最小的那個.
當然很擔心.
有人去找他.
有人就站在原位.
等他回來.
十幾分鐘之後.
終於回來了.
見到他回來的時候.
我就走前.
護底身.
一手抱起他.
放在我的肩膀上.

$^{241}$我沒有問他什麼原因.
為什麼你走失了.
我當然知道.
我當然猜到了.
不過最重要的是.
他記得我們講過.
你真的走失了.
你只要回來這裡.
我們一定等你.
他找到了.
把他放在肩膀上.
最重要的是.
他是歡歡喜喜的.
他見到走失的樣子.
他問他為什麼會走失.
我要花這麼多的氣力.
我要花這麼多的心思.
主人一見到走失的樣子.
抱起他的時候.
把一切他背負的.
都已經承擔了.
他即時的感覺是歡歡喜喜.
這個歡歡喜喜.
代表著主人說.
夠了.
好了.
你終於回來了.
其實主人一直都在他身邊.
一直在等他.
不要繼續向前走.
只要他肯停下來.
回頭看一看.
其實有一位很愛他的主人.
一直在尋找他.
一直在等他.
只要他肯停下來.
回頭看一看.
是有人很愛他.
很關心他.
很著緊他.

$^{281}$一直在尋找他.
直到找到他為止.
找到了.
歡歡喜喜.
這個歡喜來得很重要.
這個歡喜不只是說.
我失去的找回了.
這個歡喜更加代表著.
他那份寒蟲.
那份著緊.
和他對這隻走失的羊.
夠了.
回來真好.
接著說.
他尋找這隻羊.
直到他找到.
還要回到家裡.
回到家裡的時候.
他第一時間.
主人說.
我已經找到了.
不過重點是主人說.
我失去的.
他將走失的羊.
帶回到家裡的時候.
他第一句要說的就是.
主人說我找到了.
他是我的.
和你們一樣.
都是屬於我的.
這個說我.
來得很重要.
各位朋友.
各位姐妹.
這個主人.
不讓任何人.
對這隻走失的羊.
有任何一些懷疑.
他找到他了.
他跟眾人說.

$^{321}$我已經找到了.
他根本就是屬於我的.
是和你們一樣的.
這一份的接納.
這一份的保護.
我相信這隻走失的羊.
在他的內心裡面.
他是會有很深的感受.
我想的感受意思就是.
如果早點回來就好了.
早點回來就好了.
並且這個主人.
跟他的朋友和鄰舍說.
你們眾人都要和我.
一同的歡樂.
一同的歡喜.
接著.
最後第七節.
他說我告訴你們.
一個罪人悔改.
在天上也要這樣為他歡喜.
一個罪人悔改.
很明顯這個故事.
是一個比喻.
比喻羊.
就是這位創造天地的主.
獨一的真神.
所深愛的每一個人.
每一個都是屬於他的.
不過在當中卻是有走失的.
走失的原因很多.
因著人自由的緣故.
因著人自由選擇的緣故.
因著人種種的經歷走失.
不過這個主人說.
你們一個罪人悔改.
在天上歡喜.
歡喜.
天上的歡喜.
不單單是我失而復得.

$^{361}$這個主他找回了.
他失去了.
當然這是值得高興的.
不過這個歡喜.
再歡喜.
是一個很深的歡喜.
這個深的歡喜意思是.
這隻羊不需要再經歷.
種種令他很不堪的境況.
這隻羊不需要再經歷那種疲累.
這隻羊不需要再經歷.
種種的捆綁和枷鎖.
甚至乎內疚.
這個歡喜再歡喜.
是因為他們清楚知道.
這隻羊在走失的途中.
他所經歷的.
是很不容易的.
更加重要的是.
因為他知道.
這隻羊如果繼續走失的話.
他的結果是更加難堪.
所以這隻羊.
當他回來的時候.
他的主人歡喜再歡喜.
一個罪人悔改.
天上為他歡喜.
比九十九個不用悔改的義人.
歡喜還大.
在主耶穌的眼裡.
不是一比九十九.
在主耶穌的眼裡.
是一和九十九.
都是這麼寶貴.
每一個都是這麼寶貴.
每一個都是他所珍惜.
他所珍愛.
罪人悔改.
值得開心的.
聖經有一句話.

$^{401}$我告訴你們.
一個罪人悔改.
在天上也要這羊為他歡喜.
比為九十九個不用悔改的義人.
歡喜還大.
說的是罪人.
更加貼切的是.
罪人是要悔改.
小朋友在做什麼?.
玩水.
洗手.
我發覺在過往的日子.
我最難叫小朋友做的一件事.
特別是之前那幾年.
就是吃飯前要洗手.
不是.
是要吃什麼之前都要洗手.
吃什麼之前都要洗手.
是很難跟他們說的.
為什麼這麼難跟他們說呢?.
因為他們之前是沒有這個習慣的.
他們沒有這個習慣.
最多就是吃飯要洗手.
但都是這麼容易.
但早幾年我們是吃什麼之前都要洗手的.
不理你究竟是零食,小食,正餐.
早不知你去過什麼地方.
回來你都要洗手的.
是很困難的.
有時難到一個地步.
我都不知道怎樣說服他們.
但有些小朋友是很乖的.
很聽話的.
但他們就是做不到.
就算做到也好.
都是一時一時這樣.
你知道早幾年是很擔心的.
各位朋友.
小朋友你要他吃什麼都要洗手.
其實原因在哪裡呢?.

$^{441}$原因是因為他看不見.
所知看不見的是.
其實小朋友走過這麼多地方.
接觸過這麼多東西他雙手.
其實是有污糟的.
其實雙手很容易是沾染了不潔的.
雅館是一個短短的路程.
都已經雙手可以沾染了很多不潔的.
甚至乎在一些物件上.
表面好像很漂亮.
表面好像沒什麼特別.
不過其實認真一點.
都是有很多污垢在的.
最困難是因為他看不見.
所以做家長要很著緊.
看不見不是代表沒有.
看不見不是不存在.
不要等到真的有問題的時候.
你才會知道原來是要洗手的.
有些不潔有些污垢.
是連自己都看不見的.
不過卻是存在的.
走過一段路而已.
接觸一點物件而已.
都可以沾染到一些不潔.
我們的人生走過一趟很長的旅程.
無論我們所做的.
我們曾經想過的.
我們口中所說過的話.
我們所去過的地方.
包括一些網站.
很多很多其實都很認真地需要去正視.
聖經有一句說話是這樣說的.
在以賽亞書的五十三章九節.
那裡這樣說.
「我們都如羊走迷,各人偏行幾路,.
也或使我們眾人的罪孽都歸在他身上.」.
其實人就像羊一樣.
是需要有空間的.
是需要有一個曠野的地方去覓食的.

$^{481}$他是需要自由的.
他是需要有一個活動的空間的.
不過很容易就被這個世界很多的事情吸引了.
牽引了.
甚至乎是跟隨了.
聖經那裡說.
「我們都如羊走迷,各人偏行幾路.」.
不過羊離開了主人是很危險的.
羊離開了主人是很危險的.
羊是需要有一班的同行者.
可以大家互相的照應.
互相的守望.
一起去跟著這位主人.
去行走他們的路程.
約翰一書一章九節.
那裡說「我們若認自己的罪,神是信實的,是封而的,必要赦免我們的罪,洗淨我們一切的不義.」.
是,有些污垢我們用水可以洗得清.
或者再用清潔劑可以洗得到.
再用強力一點的當然不是洗手.
都可以擦得掉.
不過有一些的罪.
特別是人性的罪.
靠人是處理不到的.
有一些的罪錄印在人的心裡面.
是擦不走的.
有一些的罪咎.
曾經說過的話.
曾經去過的地方.
曾經做過的事情.
這些是纏繞著一個人的內心.
是一個重擔.
甚至有時候有些的習慣很想來到他跳出.
不知為何,我就是不看不行.
我就是不去不行.
不關我的事.
是我瞬間跳出來.
有一些的習性.
靠自己是處理不到的.
特別是一些的罪.
但我們很幸運.

$^{521}$我們有一位很愛我們的主.
這個主耶穌.
祂愛罪人愛到一個地步.
祂甘願成為人的樣式.
來到人世間.
不單止傳授這份拯救的福音.
祂更加在十字架上受死.
以祂無罪的聖潔.
來代替眾人有罪.
不過主耶穌在十字架上.
受死埋葬之後.
祂復活.
祂復活就帶著一群真心相信祂.
又跟隨祂的人.
能夠得到一份永生的盼望.
主耶穌一直尋找罪人.
祂就在我們當中.
祂就在我們中間.
不過主耶穌很好.
祂不會勉強我們每一個人.
祂透過身邊很多的人.
要將我們尋回帶回.
最後是需要我們願意接受祂的擁抱.
接受祂將完完全全整個放在主耶穌的肩上.
這個完完全全讓主耶穌放在肩上.
意思是一個完全的馴服.
信靠.
教託.
承認自己的過犯.
我們一生去跟從.
主耶穌尋找每一個走失了的人.
尋找到一個地步是要即到.
找到.
並且將帶回到這個大家庭當中.
各位未相信耶穌的朋友.
或許你在人生的路都走過一段日子.
都聽過福音的訊息一段時間.
又或許是第一次聽福音的訊息.
不過很想告訴你.
是有一位很愛你的神.

$^{561}$這一位主耶穌.
祂甘願來到人世間.
去尋找每一個走失了.
不過原本是屬於祂的人.
縱然是走失了.
不過你沒有被忘記.
祂就在你的身邊.
只要你回頭望一望.
願意接受祂.
將你背負著的重擔.
將你不能夠靠自己處理的一些罪.
將你不能夠靠自己內心的一些掙扎.
種種的疲累.
你都可以接受這位主耶穌.
由祂承擔.
天下人間唯有主耶穌可以承擔人的罪.
並且靠著祂的恩典.
一起回到這個大家庭當中.
去歡歡喜喜的享受一個新的開始.
主耶穌知道你的境況.
主耶穌知道你現在內心的掙扎.
主耶穌更加知道你一直所尋求的.
我鼓勵在在的每一位.
如果你是相信了耶穌一段時間.
不過礙於種種的原因.
你離開了.
你沒有繼續來祂依靠神.
你仍然是靠自己走你人生的路.
主耶穌說夠了.
你就回到家的當中.
再一次接受祂.
再一次你對祂緊緊的去依靠祂.
主耶穌說夠了.
你疲累了.
無論在現場或在網上的弟兄姊妹.
如果你是曾經相信.
但卻沒有緊緊的去跟從.
甚至乎是遠離了神.
我鼓勵你.
神向你作出呼喚.

$^{601}$再一次去到祂的身邊.
再一次接受祂的擁抱.
再一次回到這個大家庭.
一起去走這條天路.
如果有這樣的弟兄姊妹.
我很想邀請你.
可以舉起你的手.
一瞬間我會為你祈禱.
如果有這樣的弟兄姊妹.
你就向這位主耶穌舉起你的手.
你對祂遲疑嚇人.
我就回到你的身邊.
我不再靠自己來過我的人生.
我要再一次的緊緊的去認定.
你是我生命的主.
我緊緊的去跟隨你.
還有沒有?看到了.
在這裡我也作出第二個呼喚.
在他們當中.
未曾相信耶穌的朋友.
你未曾決志相信耶穌的朋友.
有很多身邊很愛你的人.
很愛你的家人.
很愛你的同事.
很愛你的朋友.
甚或你都不知道為什麼.
我就是來了神的殿去敬拜.
我可以告訴你.
這些種種都是因為.
有一位很愛你的主耶穌.
祂一直去跟隨你.
祂一直在你的身邊.
不過祂卻是容讓.
你有一份的空間,自由.
去接受祂,去相信祂.
不要再靠自己走人生路了.
不要再靠自己.
去背負這些各樣的重擔.
不要再靠自己.
去處理一些沒辦法跳出來的問題.

$^{641}$一些的習慣,一些的罪.
包括是曾經想過的.
做過的,所去過的.
唯有主耶穌可以潔淨我們的罪.
唯有主耶穌可以讓我們得著拯救.
唯有主耶穌可以背負我們各樣的重擔.
在這一刻.
我很想邀請未信主的朋友.
作出一個決定.
去相信這一位一直尋找著你.
很深愛你的主耶穌.
相信耶穌很簡單.
簡單是因為.
相信耶穌的重點不在我們身上.
相信耶穌簡單是因為.
所有的事情,所有的罪.
所有要處理的,所有要背負的.
都已經由主耶穌一一去承擔了.
不過信主耶穌也有一定的困難.
因為信主耶穌是需要放下自己.
來真心地去跟隨祂,相信祂.
信主耶穌.
我們可以作出一個祈禱.
這個祈禱就是.
我們真心去相信.
主耶穌就是那位獨一的真神,神的兒子.
我們接受祂在十字架上為我們成就的救恩.
我們邀請祂進入我們的內心.
去掌管我們的人生.
這些一切最重要的是.
我們要承認.
我們是一個罪人.
我們是需要祂.
在這一刻.
如果有這樣的朋友.
你願意今天來決志去相信耶穌.
我邀請你.
可以在座位上輕輕地舉起你的手.
讓我們能夠知道.
可以為你來祈禱.

$^{681}$如果你願意相信主耶穌.
就接受祂.
成為你的救主.
掌管你的一生.
讓祂能夠去擁抱你.
承擔你一切.
你只要相信,承認,接受.
如果有這樣的朋友.
很鼓勵你.
舉起你的手.
看到.
如果有這樣的朋友.
請你舉手.
主耶穌愛你.
主耶穌就在我們當中.
主耶穌說夠了.
不要再自己去背負各樣的重擔.
就將一切都交託.
將你的生命交給這位主耶穌.
祂會一生看顧著你.
祂的應許是永不改變.
如果再有這樣的朋友.
我很鼓勵你.
你就在你的座位上.
輕輕地舉起你的手.
讓我們能夠知道.
為你祈禱.
有我的朋友.
你可以放心.
我們的主耶穌很愛我們每一個.
無論我們的過去.
無論我們的現在.
重要的是.
我們就憑信心.
回到這位.
在我們的主的身邊.
有沒有這樣的朋友?.
如果有的話.
很鼓勵你.
這個舉手.

$^{721}$是一份回應.
這個回應是你和主耶穌的回應.
是你去回應主耶穌那份愛.
是你回應主耶穌一直尋找你的呼喚.
是你去回應主耶穌.
這份應許.
還有沒有這樣的朋友?.
如果有的話.
很鼓勵你.
不要再等了.
神一直在尋找你.
神一直在等你.
你就憑信心.
去接受這位主.
還有沒有這樣的朋友?.
我們同心祈禱.
無論在現場或在線上.
剛才舉手的每一位.
無論是再一次回到主耶穌身邊的弟兄姊妹.
又或者是剛才舉手.
決志相信耶穌,跟隨耶穌的朋友.
我們很為你們感恩.
很為你們歡喜.
你可以讓他們作出一個跟隨耶穌的禱告.
這個跟隨耶穌的禱告.
是你和耶穌之間的一個關係.
你真心去相信耶穌.
是神的兒子.
你接受耶穌在十字架上成就的救恩.
你承認自己的罪.
也求主耶穌的饒恕.
由誠心去邀請主耶穌.
進入我們的內心.
成為我們生命的主.
我們同心祈禱.
親愛的主耶穌.
我相信你是神的兒子.
我承認我是罪人.
求你饒恕我的罪.
願主耶穌進到我的內心.

$^{761}$掌管我的一生.
奉主耶穌基督的聖名祈禱.
阿們.
\newpage



\section{使徒行傳 13:1-12}
\label{sec:031759}
\textbf{黑暗敵擋光明}~~~~禤凱榮牧師
\newline
\newline
連結: \href{http://www.church.com.hk/acms/content.asp?site=cdc\&op=show\&type=product\&code=031759\&layout=sermon}{\texttt{031759}} ~~~~ 語音日期: 2026-03-08
\newline
\newline
\hyperref[sec:031758]{\small{< < < PREV SERMON < < <}}
~
\hyperref[sec:index_chronic]{\small{[返順時目]}}
~
\hyperref[sec:031760]{\small{> > > NEXT SERMON > > >}}
\newline
\newline
使徒行傳 13:1-12
\newline
\begin{longtable}{cl}
\hline
\hline
章節 & 經文 (和合本修訂版)\\
\hline
13:1 & \begin{tabularx}{0.7\textwidth}{X} 在安提阿的教會中，有幾位先知和教師，就是巴拿巴和稱為尼結的西面、古利奈人路求，與希律分封王一起長大的馬念，和掃羅。 \end{tabularx} \\ \\ \relax
13:2 & \begin{tabularx}{0.7\textwidth}{X} 他們在事奉主和禁食的時候，聖靈說：「要為我分派巴拿巴和掃羅去做我召他們做的工作。」 \end{tabularx} \\ \\ \relax
13:3 & \begin{tabularx}{0.7\textwidth}{X} 於是他們禁食禱告後，給巴拿巴和掃羅按手，然後派遣他們走了。 \end{tabularx} \\ \\ \relax
13:4 & \begin{tabularx}{0.7\textwidth}{X} 他們既蒙聖靈差遣，就下到西流基，從那裡坐船往塞浦路斯去， \end{tabularx} \\ \\ \relax
13:5 & \begin{tabularx}{0.7\textwidth}{X} 到了撒拉米，就在猶太人各會堂裡宣講神的道，也有約翰作他們的幫手。 \end{tabularx} \\ \\ \relax
13:6 & \begin{tabularx}{0.7\textwidth}{X} 他們走遍全島，直到帕弗，在那裡遇見一個術士—猶太人的假先知，名叫巴耶穌。 \end{tabularx} \\ \\ \relax
13:7 & \begin{tabularx}{0.7\textwidth}{X} 這人常和士求‧保羅省長在一起。士求‧保羅是個通達人，他請巴拿巴和掃羅來，要聽神的道。 \end{tabularx} \\ \\ \relax
13:8 & \begin{tabularx}{0.7\textwidth}{X} 只是術士以呂馬(他的名字翻出來就是行法術的意思)敵對使徒，設法使省長遠離這信仰。 \end{tabularx} \\ \\ \relax
13:9 & \begin{tabularx}{0.7\textwidth}{X} 掃羅，又名保羅，被聖靈充滿，定睛看他， \end{tabularx} \\ \\ \relax
13:10 & \begin{tabularx}{0.7\textwidth}{X} 說：「你這充滿各樣詭詐奸惡，魔鬼的兒子，一切正義的仇敵，你還不停止扭曲主的正道嗎？ \end{tabularx} \\ \\ \relax
13:11 & \begin{tabularx}{0.7\textwidth}{X} 現在你看，主的手臨到你身上，你會瞎眼，暫時看不見日光。」立刻迷濛和黑暗籠罩著他，他到處摸索，求人拉著手領他。 \end{tabularx} \\ \\ \relax
13:12 & \begin{tabularx}{0.7\textwidth}{X} 省長看見所發生的事就信了，因對主的教導感到驚奇。 \end{tabularx} \\ \\
[1ex]
\hline
\hline
\end{longtable}
$^{1}$我們一連幾星期都在看《市堂行傳》.
今天來到《市堂行傳》第十三章.
差不多來到一半了.
不知道大家來到這裡.
一邊看《市堂行傳》的故事的時候.
對早期教會的發展來到現在.
你有什麼觀察.
有什麼感覺呢?.
我自己的感覺就是.
來到這裡似乎給我感覺.
當時的教會是正在回應究竟他們在遭遇什麼事情.
正在回應當時的情況.
好了.
例如呢.
當教會有窮困人的時候.
有心有錢的信徒.
賣光他們的田產.
販物公用.
當他們受到逼迫的時候.
環境所逼.
離開耶路撒冷.
走到其他地方那裡.
同一個時間.
將福音帶到其他的地方.
好了.
當有人想聽福音的時候.
例如哥尼劉記得嗎?.
那麼.
彼得就根據他的要求.
走到跟他們講福音.
剛才我所講的全部都是好東西.
不過始終教會那時候.
都是因應遭遇的事情.
而作出回應.
作出反應的.
今天的經文告訴我們.
在這裡這一刻.
教會有一個很重大的不同.
教會從一個回應環境狀況的教會.
而變為一個很主動的教會.

$^{41}$開始第一次.
第一次由聖靈引導.
有意識的.
來宣講福音.
這個故事.
告訴我們基督徒.
很多我們需要知道.
或者當我們面對屬靈的征戰的時候.
我們應該怎樣來參與.
大家沒有聽錯.
是征戰.
我們今天對戰爭一點都不陌生.
四處都是戰爭.
今天我們所看到的.
是一個我們看得見的世界裡面.
看得見的戰爭.
但聖經告訴我們.
在看得到的屬靈世界.
我們是有很多我們看得見的戰爭.
現在是正在發生的.
其實.
如果你看回聖經.
當罪第一天進入世界的時候.
你可以說.
良善與罪惡之間.
神與撒旦之間.
一直都是存在這個的征戰.
從喪世的第一天開始.
邪惡的能力.
邪惡的勢力.
一直都用盡它的方法.
來阻撓天國的發展.
好了.
讓我們通過今天的經文看一看.
撒旦的工作.
或者我們當面對撒旦的工作的時候.
我們應該怎樣作為基督徒去面對.
我們一開始看的時候.
今天所講的故事發生在一個地方叫做.
安提阿.

$^{81}$安提阿是當時一個羅馬第三個最大的城市.
那個城市你可以說是一個很現代化.
一個大都會.
大融爐.
融匯了當時有希臘的文化.
羅馬,猶太,敘利亞.
很多不同的民族.
不同的文化在當中的.
而在安提阿這個地方.
耶穌的信徒.
是第一次被稱為基督徒.
也都是這個地方.
大家留意.
在安提阿.
而不是耶路撒冷.
是第一次.
這個教會是有跨地區的宣教的活動.
所以.
安提阿你可以說是在教會的發展歷史上.
是一個極其重要的一個城市.
而今天我們所看的故事.
我們看到一群人的聚集.
這個聚集不是普通的一個教會的會議.
而是一個很重要的時代的分水嶺.
什麼是分水嶺呢?.
剛才說過了.
以前福音或者教會.
都是在猶太人的搖籃裡面成長.
但是由那一次開始.
福音就走出猶太人的搖籃.
而開始擴散.
將福音帶到地極.
你可以說安提阿.
就在那個行動當中.
就是基督的宣教歷史的開始.
在這裡值得一提.
當時這幾個人.
你可以說是代表著教會的領導層.
我想看看他們的組合.
如果我們看這幾個名字的話.

$^{121}$發覺這群人真的是來自五湖四海的.
來自不同的社會階層.
不過.
而因著耶穌基督.
縱然有很多的不同.
仍然的合而為一.
我們看看這幾個領袖的名字.
第一個.
巴拿巴.
巴拿巴根據這裡所說.
就是塞浦路斯的猶太人.
好了,西面呢.
西面有個名字也叫做尼格.
意思是指黑色.
有人就說應該是來自非洲的黑人.
另外一個就是古里內人路球.
是來自北非.
另外一個就是馬念.
馬念就很特別了.
馬念就跟昏風王希律一同長大.
從小玩到大的.
希律是誰呢?.
大家都知道了.
這個希律呢.
就是當時.
殺了施洗約翰.
斬了他的頭的那個希律.
也是戲弄耶穌的希律.
一起玩大的一個朋友.
很明顯這個馬念不是普通人.
必定是一個貴族.
必定是屬於當時的一個精英階層.
素羅呢?.
素羅也就是保羅.
是當時受過很高等教育的一個法利賽人.
可以說是希伯來人當中的希伯來人.
單從這幾個名字.
我們可以看到這個領導的團隊是什麼?.
來自我們可以說是.
跨民族,跨文化,跨社會階層的一個組合.

$^{161}$如果這樣看的話,他就是坐在這裡開會.
十中十聯合國在這裡開會,是不是?.
不過大家想一想.
這個就是正正.
耶穌基督在十字架上成就救恩而應該有的果效.
所有人與人之間.
原本應該存在的隔閡.
原本應該有的那一面牆.
耶穌基督已經為我們所有人與人之間.
原本不應該有,但卻是存在的牆.
完全的拆毀.
所以不再分猶太人還是外邦人.
也不再分東方人還是西方人.
不管你是有錢還是貧窮.
也不管你是尊貴還是卑微.
所有的人.
只要我們相信耶穌基督.
在基督裡面都合而為一.
所有的不同都不再存在.
全因我們在耶穌基督裡面有同一的信心.
這一個就是神的國.
大家記不記得.
耶穌基督說過一個比喻.
撒菜種的比喻.
撒菜種是最小的.
但當它大了成為一棵樹的時候.
天上的飛鳥都會飛到這棵樹那裡.
在枝上飛來縮在樹上.
這一個比喻當中的演變.
很明顯就表達了神的國度的擴展.
當聖經說.
所有的雀都.
所有天上的飛鳥來縮在它的枝上.
同樣.
我不管你是什麼人種.
不管你是什麼的界別.
不管你是什麼的階級.
只要我們有耶穌基督的信心的時候.
就能夠進到天國.
就好像天上所有的飛鳥都聚在這棵樹上一樣.

$^{201}$這個情景就發生在早期的教會那裡.
早期的教會的人做什麼呢?.
聖經告訴我們.
在第二節.
他說「他們事奉主和禁食」.
當他們這樣做的時候.
聖靈就告訴他們.
你要為我分派這兩個巴拿巴和蘇羅.
去做我要他們做的事情.
我想大家想一想.
你猜當聖靈來.
叫他們猜派巴拿巴和蘇羅.
你猜這群門徒在做什麼的時候?.
在什麼場景,什麼場合裡面.
聖靈來吩咐發出這個命令呢?.
大家留意.
當時他們不是在開猜拳會議.
當時也不是在討論什麼教會擴展的策略.
他們在做什麼呢?.
很簡單.
只是事奉和禁食.
簡單來說.
就是敬拜神.
我不知道大家有沒有教過收音機.
現在不用了.
現在按個鍵.
我們老一輩教收音機.
你還沒教台前聽到什麼聲音?.
聽到沙沙聲.
現在沒有這個年代了.
聽到沙沙聲.
但是當你教對頻率.
你就會聽到頻道播放的節目.
很明顯.
我想問你.
你知不知道廣播這些電台的運作方式是怎麼樣的?.
他們會不會是說.
看看有沒有人教對頻道在聽我的節目.
沒有的時候收工不用做事.
播放也浪費力氣.

$^{241}$我一看到有人的時候我才播放.
是不是這樣呢?.
當然不是.
那個電台的廣播是不斷的播放他的節目.
聽不聽到就在乎我們有沒有教對頻道.
同樣聖靈其實一直都是跟人跟信徒說話.
只不過我們教不對頻道聽不到.
有什麼令我們聽不到聖靈的聲音?.
多的是.
我們生活上的忙碌.
我們的憂慮.
我們的娛樂.
我們的野心.
我們的驕傲.
我們的懶惰.
有一大串的東西.
可以讓我們聽不到聖靈的聲音.
不過這班信徒不同.
這班信徒會通過禁食.
他將世界的噪音調校到最低.
他通過祈禱.
通過事奉.
將調校的頻率調校到聖靈的頻率.
當他調校對的時候.
他就聽到這個很清晰的命令.
你要為我分派巴拿巴和素羅.
去做我召他們要做的事情.
很清楚.
我想大家留意一下.
這兩個人.
巴拿巴和素羅.
其實一直已經在事奉.
他們一直都在教會裡面.
擔任著很重要的領導的角色.
教導,事奉,建立教會.
在做很多的事情.
他們不是坐在教會裡面.
在等事情做,等事情發生.
不是.
神很多時候差派他的僕人.

$^{281}$是一些他已經願意站出來.
已經在事奉.
這班人已經是事奉的參與者.
而不是旁觀者.
不是翹起雙手等人.
看著別人做事.
當別人做完事之後.
就隨口說兩句來評論一下.
頂梓梅.
這個.
我們大家去思考的問題就是.
我們今天究竟是旁觀者.
還是教會的參與者.
如果我們想回應神的愛.
如果我們想領受更多神的恩典.
頂梓梅.
做一個參與者.
教會實際上需要更多的參與者.
來讓教會得著建立.
所以感恩.
第一代的基督徒.
早期的教會.
很多這些忠心.
願意參與.
充滿恩賜的信徒.
以致我們今天可以在福音的果子裡面.
去成長.
我再看下去.
在這個教會領受了聖靈的命令之後.
這班信徒就怎樣?.
禁食,禱告.
然後按手在他們頭上.
然後打發他們去.
我想大家留意一下.
我們聽起來很簡單.
但這對教會來說絕對不是一個容易的決定.
為甚麼?.
為甚麼不是一個容易的決定?.
我們要明白這兩個巴拿巴和蘇洛.
我相信是當時一個很重要的教會支柱.

$^{321}$很能幹.
是做到事的領袖.
但是教會不會因為這兩個最做到事的那兩個人.
證明你不是,你派別人行不行?.
派了兩個最能幹的人.
我們走了之後我們教會怎辦?.
教會會倒塌的.
沒有,你們要發覺.
他們沒有這樣的憂慮.
很明顯.
這個決定反映了告訴我們.
這班信徒他們的福音事業.
不只是看著自己的當.
是超越地區,直到地極的一個福音事業.
所以我們可以很合理的來做一個定論.
就是說當時第一代的基督徒.
他們做過一些正確的事情.
做對事.
他們很忠心的事奉神.
建立基督的身體.
信服聖靈的帶領.
所以聖經講得很清楚.
神就天天加他們的人數.
讓他們的音的範圍一直擴展下去.
神是很大的賜福.
很大的恩典.
不過現實.
如果大家有留意的話.
無論教會什麼時候做對事的時候.
撒旦走了的時候就會搞事.
有沒有留意?.
每一次教會有成功,有增長.
有果子結出的時候.
撒旦就會有去阻撓的工作出現.
在當時發生的事情就是.
當素羅和巴拿巴跟著聖靈的帶領.
去到一個地方叫做撒拉米.
然後在那裡的會堂傳講福音.
接著有一個羅馬政府的高官叫做方伯事求保羅.
很有心聽到這兩個人來講耶穌的時候.

$^{361}$邀請他們到他家裡來講.
這個方伯事求保羅的名字很蹺口.
要解釋一下.
方伯是指羅馬政府的代表的一個稱號.
事求保羅就是他的名字.
方伯事求保羅是當時那個區.
你可以說是最高的羅馬政府代表.
聖經是說.
是這個人主動邀請素羅和巴拿巴去見他.
不是因為想在他們身上拿什麼好處.
或者拿保護費.
而是很真誠的.
想聽誠的話語.
好了.
當神在某一個地方.
開啟了福音的門的時候.
大家留意.
撒旦很有趣.
就在那個位置.
捉了一道牆.
讓人不能夠清楚明白到這個福音.
撒旦就會開始有阻撓的工作.
除了這個事求保羅.
聖經在這個故事也提到另外一個叫做巴耶穌.
這個巴耶穌也叫做以呂瑪.
聖經說這個以呂瑪是一個行法術.
但是一個假先知.
他是猶太人.
聖經也告訴我們.
這個事求保羅和巴耶穌或者以呂瑪.
是經常在一起的.
大家想一想.
過不覺得奇怪.
兩個完全不同世界的人.
一個猶太人一個羅馬人.
一個是行法術的.
另一個是羅馬高官.
為什麼兩個完全不同的人會經常在一起.
我可以估計大概不是因為他們是好朋友.
而是因為什麼.

$^{401}$相互利益.
你拿他的好處他就拿你的好處.
兩個人就在一起.
這個羅馬的軍官我相信是看中了猶太人.
猶太人對了.
你幫我和猶太人建立關係.
又或者用他的法術來得到某些好處.
對於行法術的那個更加清楚.
一個高高在上的政治人物.
必定可以在身上拿到一些聲譽上.
物質上的好處.
不過.
當他見到這麼好的一個利益來源.
現在有威脅.
因為這個蘇羅和巴拿巴所講的.
聖經的真理.
會威脅到他和這個高官的關係.
我們看一看.
當他感受到這個威脅的時候.
聖經就說.
這個以女馬敵黨使徒.
要叫方伯不信真道.
雖然我估計當時這個以女馬.
大概去敵黨使徒的原因.
就是為了自己的好處.
但是我們需要明白.
如果我們從剛才那個世界觀來看的話.
每一列良善與罪惡與邪惡.
不斷的交戰當中.
所以我們相信.
這個以女馬的行動背後真正的敵人.
是撒旦.
而撒旦那個對人的搞擾.
從古到今可以說是沒有改變.
你們去到那幾招.
所以我在這裡想和大家看看.
五個撒旦最常用力阻撓教會.
或者迷惑人心的五個最常用的伎倆.
我想用五個D.
來代表大家容易記的D.

$^{441}$我想看第一個D.
第一個D是distort.
蒙蔽人的心.
撒旦會用很多的掩蓋.
或者扭曲福音.
令人無法明白福音真正的意義.
看不到福音的寶貴.
大家很常聽到.
這個世界哪有神.
基督教是西方用來控制我們的思想的一個工具.
永生.
死後永生.
你不要和我說死後.
我現在說的是現在.
現在的生活才是最重要.
現在我們已經是一個科技的世界.
我們上到月球甚至去到火星.
你還這麼迷信.
蒙蔽人的心.
第二個是distort.
分散人的心.
我很忙.
不要叫我回教會.
星期日.
五天工作六天工作.
星期日我要睡晚一點.
我要去飲茶吃飯.
拍拖什麼的都很煩.
不要叫我回教會.
很忙.
人總是有很多的事情.
娛樂焦慮.
世俗的追求.
來將人的心轉移.
對屬靈事物的追求.
第三個deceive.
欺騙人的心.
這個世界有很多虛假的宗教.
這個世界有很多半真半假的真理.
或者有誤導性.

$^{481}$好像很有道理的屬靈哲學.
這些都是謎畫或者是欺騙人的心.
讓人覺得我的屬靈需要已經得到滿足.
不再需要追求其他的信仰宗教.
最近大家都知道.
一個宗教領袖.
哈梅內伊被刺殺.
有人有些陰謀論.
就是說這個哈梅內伊可能是專登被人殺的.
為什麼.
因為在他們那個宗教的觀念裡面.
殉道者是怎樣.
殉道者有很多獎賞.
必定可以上天堂.
也有很多獎賞.
所以你是殉教者的時候有很多的好處.
我在想.
到他睜大眼的時候.
沒有這些好處.
不能上天堂.
是另外一個完全的世界.
我相信那個震驚的驚惶.
很大.
大得來.
沒有彎轉.
因為那個已經是永恆的目的地.
這個世界有很多虛假的宗教.
欺騙人心.
第四個D.
分裂教會.
這個是撒旦很常用的一個招數.
外面的東西搞不定教會.
在教會裡面煽動一些分裂的事情.
當人見到.
我在家裡已經掛了很多壓力.
我工作的環境掛了很多壓力.
你回來教會還要禁這麼多人事的問題.
我不回教會了.
人總是看到基督徒的問題,行為.
而拒絕接受福音.

$^{521}$所以這個很好用.
分裂教會.
第五.
Delay.
拖延人的心.
遲一點才信耶穌.
現在讓我多玩一會.
到適當的時候我就信了.
我記得我一次傳給一個國內的朋友.
他跟我說.
牧師.
我都覺得很需要信耶穌.
不過我現在不可以信.
我說為什麼你不可以信.
他說因為我是共產黨員.
共產黨員不可以信耶穌的.
我說那怎麼辦.
我說我退休我就可以信耶穌了.
我退休.
我就跟他失了聯絡.
我真心的盼望.
他現在應該退休了.
我真心的盼望.
他仍然可以有這個機會信耶穌.
不過我想跟大家挑戰的就是.
如果你是以上那五種.
陷在撒旦的伎倆裡面.
不能夠將你的眼目.
投放在耶穌基督的真理上面的時候.
大家要小心.
不要被撒旦的伎倆籠奴了你.
我們要找緊我們能夠有的機會.
將我們的人生投放在耶穌基督那裡.
如果我們是基督徒的話.
留意留心撒旦常用的伎倆.
千萬不要被他利用為.
成為他伎倆裡面的一個棋子.
我們要有足夠的認知.
防備提防撒旦在教會裡面的攻擊.
好了.

$^{561}$你說掃羅跟巴拿巴怎樣應對呢?.
好了.
經文這裡說.
這個掃羅.
也都叫做保羅.
大家留意一下.
為什麼特別在這裡呢?.
作者就強調掃羅.
另外一個名字是保羅.
從這裡開始一直下去都是用保羅的.
只會有這樣的改變.
在這個地方為什麼這麼特別呢?.
這個地方沒什麼特別.
但是在這一處.
你會發覺保羅漸漸好像是由一個.
事工的背景.
慢慢走到事工的前線.
甚至在燈光的底下照著他.
聖經告訴我們保羅被聖靈所充滿.
然後去責備巴耶穌.
祂說你不要再混亂主導.
然後跟祂說.
你要得到一個懲罰.
什麼懲罰呢?.
懲罰就是暫且不見日光.
大家覺得奇怪嗎?.
有沒有似曾相識的感覺?.
保羅自己也試過.
保羅自己也試過三天都看不到東西.
而巴耶穌也是一段時間都看不到東西.
我那時候想會不會是保羅憋著一口氣.
我那幾天這麼辛苦.
你也要受一下我那天的痛苦.
是不是這樣呢?.
當然不是了.
為什麼我肯定不是呢?.
因為聖經說的.
他在說這番話之前.
聖靈被聖靈充滿.
就是說那次所說的說話.

$^{601}$是聖靈的意思.
為什麼聖靈要讓巴耶穌盲了眼呢?.
我相信很有象徵意義.
大家想一想.
當時巴耶穌或者以禮馬.
他在做什麼呢?.
聖經說他混亂主導.
什麼意思呢?.
就是說一些東西令信徒看不到真理.
看不到真理.
他的行為就是讓其他要看到真理的人看不到真理.
想看到東西的人讓他們看不到東西.
神就讓他們真的看不到東西.
這是不是一個很諷刺的懲罰?.
以其人之道,還治其人之身.
聖經沒有告訴我們.
這個以禮馬會不會好像保羅那樣.
三天之後看東西了.
悔改,認罪,重新上路.
聖經沒有交代.
但是聖經所說的信息我覺得是相當的清晰.
正如保羅在其他地方這樣說.
我們不要自欺.
神是興蠻突突的.
人種的是什麼.
收的也是什麼.
所以我們要信著聖靈撒種.
就必從聖靈受永生.
所以我們或吃或喝.
無論做什麼.
都要為榮耀神而行.
所以弟兄姊妹,這是一個很重要的提醒.
對於我們每一個人來說.
我們今天活在世上.
究竟是按照我們自己的感受.
按照我們的情慾.
還是按照聖靈來生活呢?.
老實說.
如果我們單是按照自己的情感.
按照自己的情緒來生活.

$^{641}$我告訴大家,很容易的.
而且很爽的感覺.
為什麼?.
舉個例子.
你對我差嗎?.
我一定會在世衛上面弄死你.
你看不起我嗎?.
我一定在其他人面前唱衰你.
你處處都給我面子嗎?.
行行行,我以後什麼事情都站在你身邊.
很簡單.
按照情緒來生活.
按照我們的喜怒哀樂.
來做出行為上的抉擇.
很容易,很爽,很輕鬆.
不過聖經告訴我們.
有一個更好的生活方式.
什麼生活方式?.
就是一個榮耀神的生活方式.
這就是使徒保羅在這裡所展現的一種心態.
當你面對對福音的抗拒.
對福音的扭曲.
巴耶穌以裡罵的時候.
保羅就來面對這個邪惡.
糾正這個扭曲.
以致神的榮耀不受虧損.
當他這樣做的時候.
這個示求保羅.
看到這個所做的事.
很稀奇主的道就信了.
大家留意.
這個追求神的人.
看到這個神蹟就信耶穌.
但是留意.
他是稀奇主的道.
是主的道,神的話語首先讓他的心有所改變.
你可以說神蹟是印證了福音的訊息.
但是是神的道才能夠改變人的心.
神蹟,歧視往往是將人指向真理.
但是神蹟,歧視從來不是信心的終點.

$^{681}$所以大家要很清楚.
示求保羅也是因為這個經歷成為了第一次教會宣教旅程的第一個信主的外邦人.
因此福音就進入羅馬政府的權力架構.
而保羅也不再是以前逼迫教會的訴羅.
而是成為教會發展一個很重要的人物.
沒錯,撒旦從來都不停止攻擊,阻撓神的工作.
但是我們可以很放心.
因為撒旦的工作從來都不會得到實驗.
我想大家想像一艘船.
一艘船快要泊岸.
在一個深夜大風大霧的時候快要泊岸.
很快就要撞到這個岩石.
相當危險.
突然間燈塔就有一個強光劃破夜空.
照在這個水面上.
大家想像一下如果你是那個船主你會怎樣.
馬上安心了.
原本想著隨時發生意外.
現在看到了.
知道怎樣行走了.
這個蟻,女,馬.
就是這個例子的風浪和那個霧.
讓這個船家看不到怎樣走這條路.
而保羅就是這個燈塔.
保羅通過他的生命來讓聖靈好像光一樣劃破夜空.
穿過這個濃霧.
讓船家看到這個方向.
就好像聖靈的光.
一投射到這個世球保羅的心上的時候.
他就看得到真理.
他就可以從死亡的道路上走回永生的道路.
雖然撒旦的工作不斷阻撓神的工作.
但告訴大家.
黑暗從來不能驅走光明.
相反是光明可以驅走黑暗.
各位姐妹.
我們今年是粵語堂成立的26週年.
我想大家認真想一想.
神擺我們每一個信徒在一間教會裡面.
不是偶然.

$^{721}$神有祂的心意.
所以我們來到教會不單是適合朋友.
有更多的社交生活.
我們去飲茶,去購物.
有人陪我們.
這些重要的關係是相當好的.
但神告訴我們.
我們更重要的一個神的心意.
就是我們要在教會裡面建立教會.
投入在教會的事工裡面.
讓福音得道廣傳.
讓我們在過程裡面得逐造就.
這就是耶穌基督的心意.
所以讓我們把握這個機會.
想一想.
我們都希望成為這個燈台.
讓聖靈通過我們.
來將祂的光照射到這個黑暗的世界.
讓我們可以成為神一個忠心的兒女和僕人.
請我們從低頭祈禱.
26年不是一段短的時間.
但是神啊.
你在昔日如何帶領我們.
今日我們求神你仍然的賜福保守.
讓我們走在這個教會的路上.
我們個人的成長路上.
有你的恩典.
有你的帶領.
有你的賜福.
讓你的榮耀得著高舉.
讓我們的生命得著造就.
多謝你讓我們一起祈禱.
奉主名求.
阿們.
\newpage



\section{羅馬書 5:18-21}
\label{sec:031760}
\textbf{失敗中的英雄}~~~~梁錫強醫生
\newline
\newline
連結: \href{http://www.church.com.hk/acms/content.asp?site=cdc\&op=show\&type=product\&code=031760\&layout=sermon}{\texttt{031760}} ~~~~ 語音日期: 2026-03-08
\newline
\newline
\hyperref[sec:031759]{\small{< < < PREV SERMON < < <}}
~
\hyperref[sec:index_chronic]{\small{[返順時目]}}
~
\hyperref[sec:031761]{\small{> > > NEXT SERMON > > >}}
\newline
\newline
羅馬書 5:18-21
\newline
\begin{longtable}{cl}
\hline
\hline
章節 & 經文 (和合本修訂版)\\
\hline
5:18 & \begin{tabularx}{0.7\textwidth}{X} 這樣看來，因一次的過犯，所有的人都被定罪；照樣，因一次的義行，所有的人也就被稱義而得生命了。 \end{tabularx} \\ \\ \relax
5:19 & \begin{tabularx}{0.7\textwidth}{X} 因一人的悖逆，眾人成為罪人；照樣，因一人的順從，眾人也成為義了。 \end{tabularx} \\ \\ \relax
5:20 & \begin{tabularx}{0.7\textwidth}{X} 而且加添了律法，使得過犯增加，只是罪在哪裡增加，恩典就在哪裡越發豐盛了。 \end{tabularx} \\ \\ \relax
5:21 & \begin{tabularx}{0.7\textwidth}{X} 所以，正如罪藉著死掌權；照樣，恩典也藉著義掌權，使人因我們的主耶穌基督得永生。 \end{tabularx} \\ \\
[1ex]
\hline
\hline
\end{longtable}
$^{1}$弟兄姊妹主內平安.
除非你是剛剛從太空回來.
都知道中東的局勢是非常緊張.
就好像烽煙四起.
被人誤以為是二次世界大戰之後.
真是一個最危險的時期.
沒有說錯的 其實.
今天我們在一個安定的生活當中.
我們很多人覺得不關我們的事.
好像很遠 除非你最近買了機票.
要經中東回香港 可能要改一下.
但也不是的 其實影響很大.
我們知道 今天看新聞說.
原來英國那些天然氣的儲備只剩下一天半.
一天半 大家省點用.
那就要用很高價在國際的市場去買.
那一定會影響我們的生活.
那這些不是很奇怪.
因為戰亂這件事 從創世記開始到我們今天.
人類似乎是沒有進化過.
仍然是沉在一個罪裡面.
人與人之間的爭鬥 罪與國與國之間的爭.
直到好像主耶穌說要去到.
主耶穌基督再回來的時間才可以結束.
所以很多時候看回過去的歷史.
我們就好像一面鏡一樣可以提醒我們.
過去發生的事 今天我們仍然是在發生的.
為什麼呢?.
最近我和一些年輕人去談一下.
他們選那些GCSE選科.
當然年輕人有不同的理想 有不同的興趣.
有些人想選這一科 有些人想選那個科.
但是有一個共同點的.
大部分人跟我說 我一定不會選歷史.
一定不喜歡歷史.
我差不多問了十個 有九個半都是這麼說的.
我沒有看到哪一個會喜歡讀歷史.
我很奇怪 其實我自己本身就很喜歡讀歷史的.
可能我都受到我爸爸的影響.
因為我爸爸在大學的時候是讀中國文學史的.

$^{41}$所以我覺得為什麼會這樣呢?.
我想可能和傳統教學的方法有關.
尤其是香港 大家知道都是填鴨式的教育.
就是要你去記熟那些年月日.
發生了什麼事 考到試.
就是填那個填充或者答到就可以了.
其實你從來沒有想過.
究竟這個歷史和我們今天有什麼關係.
你不會有進入這個思維的 是不是?.
就好像大家去考這個Life in UK 是不是?.
1066年發生了什麼事 知不知道?.
每個人都知道 是不是?.
不知道就快點讀了.
就是Battle of Hastings 就是Hastings的戰役.
是影響了英國後來的一生.
不是要你記熟1066發生了什麼事.
這個不是重點 重點就是1066之後.
究竟英國成為一個什麼世界.
成為今天一個什麼世界.
這個就是我們讀歷史的主要原因.
如果真的要讀歷史 不是記熟那些年月日.
而是真的用一個我們所謂反思.
和一個批判性的角度去看.
一個反思 一個批判性的方法去看歷史.
我們不只是說發生了什麼事.
而是為什麼這些事會發生.
和這些事發生之後對我們有什麼影響.
為什麼說這些東西呢?.
因為我們今天看《舊約聖經》.
都是一本歷史的書.
這些事是已經發生了.
我知道我們早期教會已經將《舊約》分了類.
頭那五本書就是摩西五經 我們叫做律法書.
這個也是猶太人的律法典章來的 五本書.
接著以今天我們《聖經》的排位.
就是12本的所謂的歷史書.
由約書亞記開始講到以色列文.
去到迦南地方去征服迦南地的經過.
到今天我們所講的士師記.
就是去到他們在迦南地定居之後.

$^{81}$他們的墮落令到外族人的侵略.
而上帝是要興起去史詩去拯救他們.
這個上次我們講《基典》也有講到這個.
所以《聖經》其實你可以當作一本歷史的書去看.
有很多的故事 尤其是《舊約》.
有很多引人入勝的故事 我們看得津津有味.
主要是學生也很喜歡看的.
也可能有《聖經》的問答遊戲 我們也很喜歡去參加.
但是究竟我們要問 上帝在跟我們講什麼話?.
這個可以很深去研究這件事.
但是其實很簡單的一件事 就是剛才講的.
人性就是由抗世紀到今天.
看不到有什麼轉變 我們仍然活在罪裡面.
這個就是很明顯看得出來的.
最近我們也講關於逆轉.
我們叫reverse或者顛覆.
上帝的福音 他的教導是顛覆了我們人的思維.
所以為什麼很多人不信基督教 覺得不公平.
為什麼那個好人會遭遇這麼多的苦難.
那個壞人剛剛坐牢出來 他可以中六合彩.
你明白嗎?真的是真事.
試過有人說拿回六合彩 不可以的 他真的中了.
我們覺得我們的神不公平.
因為我們的思維就是好心得好報.
我們要做多點好事 多點福氣.
無論是中國 無論是外國.
那個倫常的思想就是這樣的.
但是上帝來的就是顛覆了我們的思維.
所以我們就很迷惘 我們就不明白.
上兩星期我們聽陳全道講的八福就是顛覆了我們的思維.
於是我們今天不再詳細討論.
今天講到的題目已經知道了 失敗中的英雄.
失敗怎麼會是英雄呢?.
一個英雄又怎麼會失敗呢?.
是 英雄莫問出處 可能他有前失敗.
但是你不會將他的失敗的東西搬出來講的.
你不會將不好的東西搬出來講的.
但是聖經就要強調 恩格斯所講的參宣.
他是一個失敗中的英雄.
就好像上次我們講的機電 他是一個軟弱中的勇士.

$^{121}$勇士又怎麼會軟弱呢?.
聖經就是要顛覆我們的思維 但是我們不明白.
所以在20世紀的時候 我們上世紀.
我們開始有這個所謂成功神學的一個理論.
大家可能都聽過 Prosperity Gospel.
基督教是一個分支 我不可以說它是異端.
但其實它講很多東西是違反了我們傳統聖經的教導.
它教導了我們就是 它基於什麼呢?.
基於靈恩的運動 很著重屬靈的恩賜.
很著重福氣 很重視行為.
很重視 跟著混合民間的倫理思想.
即是積福和好心得好報那種的心態.
但是聖經不是這樣的.
它們很注重什麼呢?.
就是說 如果你是很有信心的話.
你就 上帝會給你很多福氣.
你就會很健康 有很多財富.
如果你有病的話 你貧窮.
祈禱上帝都不答應 因為你信心不足.
我真的聽過有牧者是這樣說.
但不是這樣的意思 聖經從來不是這樣的.
成功神學是不對的 在這一方面.
甚至他們發展成就是 如果你奉獻越多.
你的福氣越大.
我真的聽過有個教會有個主席出來.
呼籲人們奉獻是這樣說的.
那窮寡婦只有兩個小錢 就沒有福氣了.
是不對的 聖經不是這樣的.
成功神學將這個豐富的生活.
用物質來衡量 是上帝祝福的證據.
但是上帝的教導絕對不是這樣.
上帝的教導是什麼呢?.
是逆轉我們的思維.
上次也說過了 福音來到是怎麼樣?.
福音來到是將原本你是軟弱.
你困苦 甚至絕望的人.
因為藉著信靠福音 得到一個新希望 新方向.
這個就是上帝要跟我們說的.
上帝告訴我們 就是你在困境那裡.
你可以因為上帝的大能 上帝的救贖.

$^{161}$你變成一個祝福.
而這個祝福不是用物質來衡量.
是你和上帝的關係.
是你怎麼去面對你的困難的時候.
仍然感覺到上帝跟你同在.
那個才是真正的福氣.
所以保羅在哥林多後書說過.
他說 因為我什麼時候軟弱 什麼時候就剛強.
大家記得這句話.
似乎是不通的.
你是剛強就是剛強 軟弱就是軟弱.
為什麼你什麼時候軟弱 什麼時候剛強呢?.
這個就是顛覆了我們的思維.
當你軟弱的時候.
你更加剛強地抓住上帝的手.
知道自己沒辦法勝過軟弱.
你知道自己根本就沒辦法可以對抗這個世界.
你只能夠是剛強起來抓住上帝的手.
所以保羅後來在哥林多後書六章裡也說.
他說 今天我們似乎是憂愁 其實是喜樂.
似乎是貧窮 其實是富足.
似乎是一無所有 其實是樣樣都有.
大家回去思考一下.
你今天覺得憂愁.
其實不是的 有喜樂在裡面.
問題是你看不看到喜樂在那裡.
問題是你用哪把尺去量度.
你以為自己貧窮嗎?.
其實你已經有很多東西了.
你以為你一無所有 其實你樣樣都有.
這個就是顛覆了我們的生命的教導.
今天所說的參宣也是.
他是顛覆了我們的教導.
參宣勇士 參宣是一個英雄.
但是參宣絕對不是一個好人.
不是一個成功的人 他很失敗 很失敗.
《時事記》有四章是講關於參宣的.
大家回去北韓看一下.
今天我們不是去研經 去看每一節.
參宣發生了什麼事.

$^{201}$但是參宣那個的生活 其實是非常差的.
他做的事真是 不要說他是英雄.
普通人都不應該做的.
在《時事記》十三章那裡開始.
現在過了機電之後 過了很多年.
聖經又再講了.
它說以色列人又行耶和華 眼中看為惡事.
留意那個「又」字 又失敗了 以色列人.
今天我們的生命都是一樣.
上帝一次一次提醒我們 拯救我們.
我們又失敗了 是不是?.
結果怎麼樣? 之前是交給了米電人.
這次是交給一些更加惡的人 就是非理士人.
非理士人就是孔武有力.
頭很高大的 打仗很厲害的.
以色列人就交給了這些非理士人手中.
多少年了? 40年.
上帝的公義來到的時候.
就將以色列人這些不聽話的人.
交給了非理士人40年.
但是上帝有憐憫一面.
公義和憐憫 上帝一直都在做事的.
他沒有放棄以色列人.
他就興起一個新的事實.
這個就是參選.
上帝什麼時候去呼召參選呢?.
上次基地 上帝走去找一個使者和他談.
這次不是.
是參選還沒出生之前.
上帝已經揀選他 和他媽媽說話.
慈禧記十三章三至五節說.
耶和華的使者向那婦人顯現.
這婦人就是參選的媽媽.
對她說過 看啊 以前你不懷孕不生育.
如今你必懷孕生一個兒子.
但是有條件.
條件就是你要謹慎.
清酒,烈酒不能喝.
任何不潔之物都不可以吃.
這個都不是很奇怪.

$^{241}$你懷孕都不應該喝酒 不應該吃髒東西.
但是接著怎麼說呢?.
看啊 你必懷孕生一個兒子.
不可以用剃刀剃他的頭.
不可以剪頭髮.
因為這孩子一出母胎就歸給上帝作拿西爾人.
他必開始拯救以色列脫離非利士人的手.
今天我們不要花太多時間去詳細討論什麼是拿西爾人.
拿西爾人是舊約一個猶太人.
其中一個族人 也就是某一種人.
他是一個見證出來.
是有上帝的記號告訴你.
他是跟從上帝作一些刻苦.
一個好像修道那種的生活.
他吃的東西是清淡.
他們不會喝酒.
他們不會碰那些髒東西.
尤其是死的物質.
死人的東西或者動物.
他們是有個特徵.
就是他們不剪頭髮的.
所以拿西爾人是不會剪頭髮的.
你看到那個人剪頭髮 他不是拿西爾人.
他只是模型.
這個是他的特徵.
上帝是很清楚告訴你 他是拿西爾人.
接著上帝的任務就是說.
他會拯救以色列人脫離非利士人的手.
上帝揀選人就是這樣.
上次他揀選基地 大家記得了.
基地是個什麼人 一個武膽鬼.
他們個個在前方對紅尾殿人打仗.
他就躲在大後方那裡弄走.
就是在工廠裡面.
但是上帝的使者來到他那裡.
第一句怎麼跟他說呢?.
他說大能的勇士.
基地 你搞錯了 點錯箱了.
我不是大能的勇士.
但是上帝是知道的.

$^{281}$上帝揀選人不是說你做過些什麼.
做了多少好事.
而是上帝認為你可以為他做多少事.
明白嗎? 這是很重要的一件事.
上帝今天揀選我們在教會那裡做事.
很多時候我們就說我不懂 沒有經驗.
上帝說不是.
我揀選你是因為我覺得你可以為我做更多事.
一定是的.
宣教士有意象走去埃塞俄比亞那裡做宣教.
以前去非洲都沒有去過.
不懂那裡的土語.
或者不喜歡吃那裡的食物.
那為什麼上帝揀選他呢?.
因為他可以為上帝做更多事.
上帝不是揀選以前在非洲住了幾十年.
很熟悉的地方 跟那些人很熟悉.
又懂那些語言 精通六國非洲的土話語言.
上帝不是揀選那些.
上帝不需要揀選那些.
是很不同的.
這個是顛覆我們的思維的.
我們今天不是的.
我們去建功我們要怎樣?.
我們就要弄好那張履歷表.
那個履歷表.
要告訴別人我做過什麼 我做過什麼 我懂什麼.
慌死人不知道的.
你要在這麼多人當中突圍出來.
你的履歷表要弄得很漂亮.
要給別人看的.
上帝揀選人不是這樣的.
不要說履歷表.
現在的小朋友都是.
幼稚園而已.
去建幼稚園也要有一個Portfolio.
知不知道什麼是Portfolio?.
我那個年代還沒聽過這個字.
原來現在很重要.
有作品集.

$^{321}$這個就是他兩歲的時候已經會畫這個圖畫了.
三歲就會作曲了.
這個小朋友叫大雄.
開玩笑的.
不放過他.
你不明白.
我們是要這樣的.
上帝揀選人不是你做了多少事情.
而是他認為你可以為他做多少事情.
這個是最重要的.
所以他揀參孫 揀基甸.
他就是這樣的情況.
參孫更加.
他什麼都沒做過.
他還沒出生就已經.
那上帝是不是揀錯人了?.
似乎是的.
但上帝怎麼會揀錯人呢?.
上帝是自有永有.
他是否在時間以外.
他知道參孫將來會是一個什麼樣的人.
他多好多壞.
上帝不知道嗎?.
為什麼上帝要揀選他?.
很奇怪吧.
這就顛覆我們今天的思維.
參孫不是一個好人.
他是一個徹徹底底失敗的人.
如果大家從參孫的生平去看那幾章聖經.
就很清楚了.
不是好人來的.
就是那個三大原罪.
大家記不記得?.
眼目的情慾.
肉體的情慾.
今生的驕傲.
這三招參孫都中了.
十四章一開始.
參孫成人的時候.
他怎樣呢?.

$^{361}$參孫就看上了一個非利士的女人.
他回來跟他的爸爸媽媽說.
他說我看中了一個非利士的女孩.
現在請你幫我安排.
我可以娶她為妻.
以前娶她為妻一定要父母之命.
參孫都不敢說.
走去私奔就算了.
所以他都問他父母.
但是不行的.
他入迦南地的時候.
上帝說得很清楚.
不可以跟外族人通婚.
非利士那些人不是信上帝.
他們拜偶像的.
他們不是屬神的人.
他爸爸媽媽怎樣說呢?.
他說你在猶太人族.
以色列族那裡那麼多女孩你不喜歡.
你是要選外族人的女孩嗎?.
那個未受國禮的非利士人的女孩嗎?.
他爸爸未受國禮.
他們不是屬神的人.
參孫怎樣回答呢?.
很簡單 我喜歡她.
喜歡就大了.
他不用怎麼辯護.
我喜歡.
現在年輕人就是這樣.
我喜歡.
你咬他不行的.
因為他喜歡.
你不喜歡 我喜歡.
參孫就是這樣.
這個人性沒有變過的.
幾千年前都是這樣.
參孫說我喜歡.
父母怕了他.
喂 父母記得的.
這個是應許.

$^{401}$是他會拯救以色列人.
從非利士人的手上拯救出來.
這個是大能的勇士將來.
這個是拿西爾人.
他父母應該知道的.
喂 現在怎樣.
你去跟非利士女人結婚.
日後怎樣拯救你族人.
從非利士的手上出來.
這是矛盾嗎?.
沒辦法 他堅持一定要.
參孫很驕傲的.
順利給他能力.
他是很大力的.
大力士就是參孫.
大家如果在聖經文大比賽.
舊約的大力士.
大家馬上舉手說參孫.
知道啦 個個都知道.
新約都有大力士的.
這個大力士是提多.
拿得多 知道嗎?.
OK 所以大家記得了.
好 參孫有一次就在這裡.
整個獅子都被他撕裂.
很威風的 不可一世被人知道.
弄個謎語出來.
叫人猜他做過什麼.
詳細的我不說了.
總之他是一個炫耀自己的人.
有一個專門搞惡作劇的人.
他不喜歡那班菲律賓人.
他怎樣呢?.
就要殺了他.
他不是.
他把三百隻狐狸的尾巴綁在一起.
用火燒著那些狐狸的尾巴.
多殘忍.
然後趕他去那些菲律賓人的田地.
燒著他們.

$^{441}$就像那些頑童在搞惡作劇一樣.
參孫是個好人?.
不是好人.
失敗?.
徹底的失敗.
他不要說做拿世爾人.
他根本做一個上帝的人.
他在上帝眼中應該是失敗.
這樣看法.
去到十六章一節.
這個又奇怪.
參孫.
一句而已.
參孫到了加薩.
在那裡看見一個妓女.
就與她親近.
原來參孫是個好色鬼.
你明不明白.
聖經就是這樣說的.
很簡單的.
不說他為什麼無端白事.
又沒有說後文.
什麼都沒有.
就是你們.
你放工走過波蘭街就找到個妓女.
他是那種.
想起要去做.
他被情慾.
這樣的肉體的情慾.
是沖昏了頭腦.
根本他就不想後果.
也不想.
什麼.
參孫就是這樣的人.
失敗?.
很失敗.
最失敗是什麼呢?.
十六章去到四至五節那裡.
「這時以後參孫又愛上另一個女孩」.
夠不夠?.

$^{481}$那個其實已經有些故事了.
已經分開了.
後來有第二個.
這個就叫做「大利拉」.
這個就是女主角出場了.
這個夜王的女主角.
這個真是壞的.
這個其實就是被非利士人買通了他.
要在參孫那裡.
試探她究竟為什麼這麼大力.
因為個個都怕了她.
她為什麼這麼大力呢?.
這個女孩就用很多很多的方法.
去引誘她.
哄她說真話.
參孫都不說了.
故事結局是怎樣呢?.
男人的結局是怎樣呢?.
通常都是衰給女人的.
結果就說了出來.
她將她的秘密說了出來.
就說其實我是靠我的頭髮.
你剪光我的頭髮我沒有這個能力.
上帝的能力離開了.
這個這麼大的秘密.
是影響她一生.
影響她民族以後的歷史.
她都說了出來.
肉體的情慾你說是多麼的重要.
多麼的重要.
不要笑這件事.
大家看新聞都知道.
美國的檔案裡面有多少名有錢人的照片.
這麼多年的罪行.
這些根本是沒有停過的.
人就是這麼軟弱.
這個當然大家熟悉參孫的故事都知道.
他出計謀剪了她的頭髮沒有力氣綁住她.
菲律賓人甚至挖瞎她的眼睛.
沒有她了 因為想折磨她.

$^{521}$綁住了她.
上帝的能力離開了她.
參孫成為一個普通的人.
手無縛雞之力的人.
因為沒有了那種能力.
因為她的能力不是從她自己而來.
很多時候我們以為能力是從我們自己而來.
我賺到錢 因為我本事.
我讀書厲害 或者我人緣好.
認識的人多.
不是 能力是從上帝而來.
很簡單 上帝拿回你一點點健康.
你就什麼能力都沒有了.
大家都很清楚這件事.
參孫也是.
上帝拿回一點點 她就沒有了能力.
被人抓住 眼睛都挖瞎了.
回到剛才看的那個《羅馬書》.
如果不是說完之後.
為什麼沒有說今天的那個《羅馬書》呢?.
《羅馬書》剛才五章那一段.
其實是一個很高深的神學.
今天不是想說這件事.
這個是說到耶穌基督是第二個亞當.
幕後的亞當 可能大家聽過.
因為看回那一段就是.
因為罪是因為一個人犯罪 就是亞當.
所以就令到我們有罪.
但是拯救救恩是因為耶穌基督來到.
是一個第二個亞當.
來到是令到我們有一個新的生命.
這個是保羅的神學思想.
但是我想大家去留意.
剛才那個經文其中一句.
他說 只是罪在哪裡增加.
恩典就在哪裡越發豐盛 很奇怪.
罪在哪裡多 恩典在哪裡多.
這個也是逆轉我們思維.
罪越來越多 那不就是什麼?.
就是上帝不理我們.

$^{561}$就是我們失敗了.
懲罰來了 是不是?.
按道理是這樣的.
但是不是.
保羅在《羅馬書》說的是一個逆轉性的福音.
他說你罪越多.
反而你更加看到上帝的恩典豐盛.
這個是真的.
其實很簡單的.
你患了很多重病.
很多的百病纏身.
你去看醫生 醫生你都治好了.
哇 這個真的是大恩典.
你越多病痛 你就越可以治好.
不是看到上帝的恩典嗎?.
不是看到醫生的大能嗎?.
一定看到的.
如果你是沒有罪的話 你沒有病痛的話.
你沒有病痛不用來找我的 是不是?.
那根本就不會看到我有多厲害可以治好你.
明白嗎?.
一樣原理.
所以罪在哪裡多.
恩典就越多.
這提醒了我們今天.
就是因為參宣失敗過很多次.
才能夠成為一個真正的英雄.
因為往後我會說他最後的結局.
你明白嗎?.
如果他一直都是一帆風順的.
沒什麼的 普普通通的.
你看不到他是一個怎樣的英雄.
最多就說上帝給了你一些能力.
你很大隻 你可以打贏很多菲律賓人.
但因為他跌到谷底 他犯很多罪.
他反而更加捉實上帝的手.
到最後他悔改.
他求上帝給他能力.
成為以色列民族的真正英雄.
罪越多 恩典越大.

$^{601}$這和成功神學是倒過來的.
成功神學不是這樣說的.
說你們不要犯罪 上帝恩典都走了.
不會的 世人都犯了罪.
我不是鼓勵你犯罪.
但接受一樣東西 世人都會犯罪.
所以失敗不是一個結局來的.
在聖經上帝的角度裡面.
很多時候我們看出戲或者看本書.
那個人失敗了 結局結束了.
那就教訓我們以後不要這麼失敗了.
那些故事就是這樣教訓我們.
但不是 原來失敗才是開始.
是上帝的救恩的開始.
沒有失敗我們不需要上帝的救恩.
沒有痛苦我們感受到上帝的同在.
因為一樣原理.
聖經是從來不會逃避說失敗這件事的.
不逃避這件事的.
你說聖經是一本偉大的書.
它寫的東西肯定是鋪得最好的.
寫到上帝的人民或者先知,使徒.
最漂亮的一面 是不是?.
但不是 聖經從來不會避開失敗.
因為它要我們看回這個歷史.
就是從失敗中看到上帝的那個恩典.
羅亞是一個義人 一家得救.
洪水來到的時候沒有被消滅.
但羅亞後來怎麼樣呢?.
喝醉酒犯了很多倫理的歪理的事.
亞伯拉罕是一個信心之父.
但他有沒有信心軟弱的時候 有.
雅各 舊約的那個雅各更加不用說.
是一個貪得無厭 是一個妒忌深重.
要爭長子命份 對子女偏心的 不是好人.
大衛有很多好的地方 也有他軟弱的時候.
聖經不會逃避這些東西.
就算是耶穌基督都有軟弱的時候 為什麼?.
因為他要告訴別人 耶穌都是有百分之百的人性.
耶穌有哭 耶穌有怕過.

$^{641}$耶穌不想拿個苦杯來喝的.
如果可以的話就拿走 祂都有猶豫過.
這就是聖經不會逃避人的那種軟弱.
彼得很清楚 彼得大家很熟悉的.
門徒彼得有時候覺得他有些精神分裂.
有時候他很厲害 有時候他很軟弱.
有一次在《馬太福音》十六章那裡.
耶穌問那些門徒 你說我是誰?.
那些人說你是先知 你是仁子 神子.
就是很多東西 但是說不出重點.
彼得就很清楚 他說你是基督.
是永生上帝的兒子.
彼得一個敢站出來這樣說.
這句話沒什麼 是不是?.
要夠膽 這句話是死罪.
顛覆國家的死罪 知不知道?.
為什麼?因為你說你是基督.
就是你是尼徹 就是那個猶太人的王.
你這個是違反了羅馬的政權.
猶太人不可以有自己的王.
你也是違反了猶太人的那個律例.
因為你是褻瀆了上帝 你說你是基督.
是上帝的兒子 所以彼得很大膽的站出來說.
耶穌就說對了 他說若拿的兒子西門.
他那時候叫西門 不是彼得 他說你有福.
他說現在我幫你改名 你的名字以後叫彼得.
彼得那時候是改名的.
他不是一出生就叫西門彼得的.
他一出生就叫Simon 不是Peter.
彼得什麼意思?彼得就是石 判石.
他說教會會建築在你這個石那裡.
陰間的權兵不可以勝過這個教會.
天國的鎖匙我給你 哇 彼得的權力多大.
他的能力多大.
但是過了幾個禮拜.
彼得三次不認耶穌 大家知道是吧?.
其實他說這番話的時候 他沒理由不知道.
可能他不知道他知不知道.
他是人他就不知道 他是神他就知道.
你明白 但這個不是重點.

$^{681}$重點就是彼得他都是失敗.
但是上帝也是委任重任給他.
我們很多時候都是經歷過失敗.
一下子彼得又說你是神的兒子 他很清楚.
一下子別人問他認不認識這個耶穌.
他說不認識的 我路過的 跟他不是很熟的.
他怕死.
但是上帝揀選的就是這些人.
還有一個叫馬可 大家聽聽馬可.
大家聽馬可就是因為馬可福音 是吧?.
馬可不是十二個使徒之一.
但是早期教會都是很活躍的.
大家看馬可 最熟悉他的經歷就是在《仕途行傳》.
都不是說很多的 《仕途行傳》.
保羅第一次去旅行的時候 不是旅行.
去宣教的時候 他帶了他去巴拿巴.
就是安提拉教會委任他們兩個去的.
他帶了幾個助手 一個就是馬可.
巴拿巴表弟馬可 另外一個是誰呢?.
另外一個一定是路加的 路加一定有份.
因為路加寫《仕途行傳》的.
他們一起去 走了一個圈 在附近就回來了.
去到安提拉教會返回職位 就回到猶殺了.
到第二次去旅行了 之前發生一件事就是.
這個馬可有跟著他們一起去.
但是沒有跟著他們一起回來 中途離開了隊.
《聖經》沒有說得很清楚為什麼他離開.
猜一猜吧 可以猜到的 可能怕辛苦.
怕辛苦啊 知不知道?.
奇怪吧?不奇怪的 真的很辛苦的.
以前這樣坐船去.
稍後我們會介紹一下退休會.
就是我們今天夏令營 我們今年那些好東西.
早期我們的退休會真的是很退休的.
真的是辛苦的.
黃牧師的年代 7點鐘起床要零收的.
沒有一口好吃的.
幸好那間是接近布里斯圖的退休會.
那些人完了之後就去布里斯圖五月花買些燒鴨回來.
不至於離開隊 不至於離開.

$^{721}$但是馬可真是搞不定了 中途離開了.
就是沒有團隊精神了 知不知道?.
齊上齊落的嘛 要進去賣瓶的嘛.
沒理由中途就走了去.
所以保羅就很不喜歡這件事.
第二次旅行的時候 就是去全島的時候.
巴拿巴亞又要帶著馬可去.
就給他一個機會 讓他學一下.
保羅說不行 一次不中 百次不勇.
你上次中途走了去.
搞亂了我們的計劃 不准帶他去.
結果怎麼樣呢?.
保羅和巴拿巴亞這兩個前線人員.
說這件事似乎是很小事 就吵架了.
結果就是 你可以說是鬧翻了.
因為結果他們不在一起了.
保羅就帶了另外一個人去西拉.
這次就走陸路.
巴拿巴亞照樣帶著馬可走海路去塞浦路斯.
因為他鄉下那裡 走了另外一邊.
那又怎麼樣呢?.
很奇怪 失敗了一個英雄.
《聖經》沒有把每件事講得很清楚.
後來怎麼樣怎麼樣.
但是有一樣東西我們留意到的.
就是保羅最後在提摩泰後書.
很多年後 提摩泰後書就是保羅孔晚期寫的.
他跟提摩泰說 他說現在有路階在我身邊.
他說你以後來找我的時候.
提摩泰你來的時候 你帶著馬可來.
因為他在服侍上於我有益.
完全改觀了 他特別嚴厲的說.
很清楚的叫提摩泰你下次來 帶著馬可來.
點名啊 帶著馬可來.
就證明馬可這段時間有一個很大的改變.
《聖經》沒有講的 但是我們都看到了.
後來馬可寫了《馬可福音》.
不是一個淡散的帥哥 不是那些二打六.
他是《馬可福音》的作者.
在服侍上與保羅有益.

$^{761}$你說這是多大的稱讚.
上帝就是去呼召一些失敗中的英雄.
馬可經歷過膽怯 經歷過失敗.
但後來保羅在眼中覺得對他有益.
對福音有益 對上帝的國有益.
所以人的故事很多時候因為你失敗就講完.
但是上帝的故事在你失敗之中.
你才剛剛開始.
這就是顛覆我們今天的思維.
罪在哪裡多 恩典就會多.
神今天不是要揀選一個完美的人.
神是在人最破碎的時候 最軟弱的時候.
要我們伸出手抓住他.
上帝就是這樣縱容我們.
當參宣去到最低的時候 最谷底的時候.
盲了雙眼 沒力氣 被人綁住.
求生不得 求死不能.
還被人帶到菲律賓人的廟裡面.
被人取笑他.
當他是一個動物園的動物給人看.
這個就是你們所謂的耶和華的大能勇士.
那個慈師 居然是這樣的.
他不殺他原因就是這樣 要取笑他.
結果去到《時西記》16章28節.
就是參宣最後一次去呼求上帝.
大家都很熟悉這個故事.
就是參宣去到這個菲律賓的廟宇裡面.
他說累了 叫人擺他在牆邊 柱邊那裡扶著.
跟著他祈禱 主耶和華求你顧念我.
他徹底的迴轉.
他知道以前做的事全部是失敗.
他說上帝就一次一次就夠了.
求你賜我力量使我向菲律賓人報仇.
報那個挖我雙眼的仇.
當然他不只是一個私人的仇恨.
因為菲律賓人是欺負了以色列人很多年了.
上帝要清洗參宣去拯救菲律賓人.
但是什麼都沒有做過.
是不是?.
那些情慾上的壞事做到盡.

$^{801}$真是哪有一個英雄呢?不是的.
但是最後怎麼樣呢?.
上帝怎樣使用了參宣呢?.
參宣祈完禱之後抱著廟宇的兩根柱子.
左手抱一根 右手抱一根.
然後將他拉下來 上帝給他一個能力.
把那個廟宇塌下來 聖經怎麼形容呢?.
他說押著那些菲律賓人的領袖那些人.
比參宣在死的時候所殺的菲律賓人.
還多過他活著的時候殺的那些菲律賓人.
猜不到是不是?真是猜不到.
是很可惜的 要死那一刻才可以成為一個真正的英雄.
但是上帝的那個計劃就是這樣 是我們猜不到的.
我們沒辦法可以猜到上帝是怎麼樣.
只能夠有一樣東西就知道.
就是你不要在你最谷底的時候.
在人生最無助的時候 你仍然一定要抓住上帝.
只能夠這個故事告訴你 就是要這樣的了.
如果不是的話 那個結局就是在那裡結束的了.
你的失敗是上帝的拯救開始 記住.
今天如果你覺得你是很失敗的.
覺得你的人生很糟糕的 你每樣都是不順境的.
我告訴你 這是上帝的恩典的開始.
我一定要抱著這個心態 不是結束.
如果你覺得是結束 你是一個失敗者.
如果馬可因為這樣 保羅不帶我去了.
以後我回去鄉下耕田就算了.
不會有馬可復任的 是不是?.
一定不會有的 所以耶穌說.
我來到不是照義人 是照聚人 很清楚了.
你不要以為英雄就是很聽上帝的說話的.
很乖的 每樣都做得很好的.
當然我們想做到這樣 但不是一定的.
不是一定的 順利來到就是要照聚人 不是照義人.
所以人們說教會的人都是這樣.
又多是非 又這樣那樣 經常都說的.
當然我們不同意 這是他們的藉口.
但也要檢討一下自己究竟做得對不對.
但是明白一件事就是.
就是因為我們是聚人才要回教會.

$^{841}$我們每個都是義人 何必進來呢?.
個個都健康 醫院倒閉了 不用開醫院 是不是?.
是很簡單的道理來的.
上帝要找的 呼召的不是完美的人.
不是要有一個很好的履歷 很好的 portfolio.
不是這些 上帝要找的是一個破碎的人.
上帝要找的不是成功的人.
因為如果你成功的話 你會驕傲.
上帝不會找你的.
上帝要找一個軟弱愚拙的人 等到驕傲的人.
以為自己有實力的人去羞愧.
保羅覺得很清楚.
上帝從來是找失敗的人.
上帝不是要找一個無瑕疵的所謂義人.
而是找一個願意被上帝醫治的罪人.
這個就是分別 這個就是逆轉的福音.
明白了沒有?.
所以保羅很清楚就是.
我的能力在軟弱的人是更加顯得完全.
如果你不軟弱的話 你怎麼顯得完全呢?.
你沒有病的話 醫生的能力怎麼會在你身上顯出作為呢?.
一樣原理.
上帝的能力在軟弱的人顯得更加完全.
我給大家低頭.
相信我們每個人都有些失敗的經歷.
我們不敢面對 不敢和別人說.
更加不敢交給上帝.
為什麼呢? 因為我們覺得自己不配 被人取笑.
覺得自己做得不夠 做得很差.
但上帝說 你的失敗中 我仍然這麼愛你.
你的失敗中 我的恩典在你身上顯得更加完全.
今天我們願不願意將我們的失敗 過去不好的事.
我們的痛苦交給上帝呢?.
讓上帝為我們編織一個新的故事.
你帶我們來到你面前.
藉著參雙的故事 保羅的教導.
你提醒了我們 我們都是失敗中的人.
因為罪在我們新世界裡.
自從創造天地開始到今天.
人是在罪裡沒辦法走得出來.

$^{881}$唯有一定是抓住上帝的手 抓住你的拯救.
今天我們不要用成功神學的定義來定自己是否成功.
父王真正的成功就是我們能夠認識主耶穌基督.
來接受這個救恩 這個就是我們的成功.
這個是上帝無限的恩典 繼續帶領我們.
當我們有絕望的時候 我們不要放棄.
因為父王今天揀選的是軟弱的人.
父王說你要我們為你做得更多 走得更加遠.
將來的路不容易走 但父王有你與我們同在.
過去可能我們經歷很多失敗.
但父王說不用跟我們說 不用怕.
只要將一切交託給你.
求主繼續帶領我們在未來的日子當中.
有主每一天都是與我們同在.
每一天靈裡可以更生.
更加愛主 更加像主 更加與主同行.
將一切永遠都交給上帝.
我們祈禱奉祂的聖名之祈求.
\newpage



\section{加拉太書 4:1-6:18}
\label{sec:031761}
\textbf{加拉太書 (2) 基督使我得自由}~~~~劉少康牧師
\newline
\newline
連結: \href{http://www.church.com.hk/acms/content.asp?site=cdc\&op=show\&type=product\&code=031761\&layout=sermon}{\texttt{031761}} ~~~~ 語音日期: 2026-03-04
\newline
\newline
\hyperref[sec:031760]{\small{< < < PREV SERMON < < <}}
~
\hyperref[sec:index_chronic]{\small{[返順時目]}}
~
\hyperref[sec:031762]{\small{> > > NEXT SERMON > > >}}
\newline
\newline
加拉太書 4:1-6:18
\newline
\begin{longtable}{cl}
\hline
\hline
章節 & 經文 (和合本修訂版)\\
\hline
4:1 & \begin{tabularx}{0.7\textwidth}{X} 我說，雖然那承受產業的是整個產業的主人，但在未成年的時候卻與奴隸毫無分別， \end{tabularx} \\ \\ \relax
4:2 & \begin{tabularx}{0.7\textwidth}{X} 仍是在監護人和管家的手下，直等他父親預定的時候來到。 \end{tabularx} \\ \\ \relax
4:3 & \begin{tabularx}{0.7\textwidth}{X} 我們也是一樣，在未成年的時候，被世上粗淺的學說所奴役，也是如此。 \end{tabularx} \\ \\ \relax
4:4 & \begin{tabularx}{0.7\textwidth}{X} 等到時候成熟，神就差遣他的兒子，為女子所生，且生在律法之下， \end{tabularx} \\ \\ \relax
4:5 & \begin{tabularx}{0.7\textwidth}{X} 為要把律法之下的人贖出來，使我們獲得兒子的名分。 \end{tabularx} \\ \\ \relax
4:6 & \begin{tabularx}{0.7\textwidth}{X} 因為你們是兒子，神就差他兒子的靈進入我們的心，呼叫：「阿爸，父！」 \end{tabularx} \\ \\ \relax
4:7 & \begin{tabularx}{0.7\textwidth}{X} 可見，你不再是奴隸，而是兒子了，既然是兒子，就靠著神也成為後嗣了。 \end{tabularx} \\ \\ \relax
4:8 & \begin{tabularx}{0.7\textwidth}{X} 但從前不認識神的時候，你們是給那些本來不是神的神明作奴隸； \end{tabularx} \\ \\ \relax
4:9 & \begin{tabularx}{0.7\textwidth}{X} 現在你們既然認識神，更可說是被神所認識的，怎麼還要轉回那懦弱無用的粗淺學說，情願再給它們作奴隸呢？ \end{tabularx} \\ \\ \relax
4:10 & \begin{tabularx}{0.7\textwidth}{X} 你們竟又謹守日子、月份、節期、年份， \end{tabularx} \\ \\ \relax
4:11 & \begin{tabularx}{0.7\textwidth}{X} 我為你們擔心，惟恐我在你們身上是枉費工夫了。 \end{tabularx} \\ \\ \relax
4:12 & \begin{tabularx}{0.7\textwidth}{X} 弟兄們，我勸你們，要像我一樣，因為我也像你們一樣。你們一點沒有虧負我。 \end{tabularx} \\ \\ \relax
4:13 & \begin{tabularx}{0.7\textwidth}{X} 你們知道，我因為身體有疾病才有第一次傳福音給你們的機會。 \end{tabularx} \\ \\ \relax
4:14 & \begin{tabularx}{0.7\textwidth}{X} 雖然你們為我身體的緣故受試煉，卻沒有輕看我，也沒有厭棄我，反倒接待我如同神的使者，如同基督耶穌。 \end{tabularx} \\ \\ \relax
4:15 & \begin{tabularx}{0.7\textwidth}{X} 你們當日的好意哪裡去了呢？那時若辦得到，你們就是把自己的眼睛挖出來給我，也都情願。這是我可以給你們作證的。 \end{tabularx} \\ \\ \relax
4:16 & \begin{tabularx}{0.7\textwidth}{X} 如今我把真理告訴你們，倒成了你們的仇敵嗎？ \end{tabularx} \\ \\ \relax
4:17 & \begin{tabularx}{0.7\textwidth}{X} 那些熱心待你們的人，不懷好意，是要隔絕你們，好使你們熱心待他們。 \end{tabularx} \\ \\ \relax
4:18 & \begin{tabularx}{0.7\textwidth}{X} 在善事上，時刻熱心待別人原是好的，卻不只是我與你們同在的時候才這樣。 \end{tabularx} \\ \\ \relax
4:19 & \begin{tabularx}{0.7\textwidth}{X} 我的孩子們哪，我為你們再受生產之苦，直等到基督成形在你們心裡。 \end{tabularx} \\ \\ \relax
4:20 & \begin{tabularx}{0.7\textwidth}{X} 我期望現今就在你們那裡，可以改變我的口氣，因為我為你們心裡難過。 \end{tabularx} \\ \\ \relax
4:21 & \begin{tabularx}{0.7\textwidth}{X} 你們這願意在律法之下的人，請告訴我，你們沒有聽見律法嗎？ \end{tabularx} \\ \\ \relax
4:22 & \begin{tabularx}{0.7\textwidth}{X} 因為律法上記著，亞伯拉罕有兩個兒子，一個是使女生的，一個是自由的婦人生的。 \end{tabularx} \\ \\ \relax
4:23 & \begin{tabularx}{0.7\textwidth}{X} 那使女所生的是按著肉體生的；那自由的婦人所生的是憑著應許生的。 \end{tabularx} \\ \\ \relax
4:24 & \begin{tabularx}{0.7\textwidth}{X} 這是比方：那兩個婦人就是兩個約；一個婦人是出於西奈山，生子為奴，就是夏甲。 \end{tabularx} \\ \\ \relax
4:25 & \begin{tabularx}{0.7\textwidth}{X} 這夏甲是指著阿拉伯的西奈山，與現在的耶路撒冷同類，因為耶路撒冷和她的兒女都是為奴的。 \end{tabularx} \\ \\ \relax
4:26 & \begin{tabularx}{0.7\textwidth}{X} 但另一婦人就是在上的耶路撒冷，是自由的，她是我們的母親。 \end{tabularx} \\ \\ \relax
4:27 & \begin{tabularx}{0.7\textwidth}{X} 因為經上記著：「不懷孕、不生養的，你要歡樂；未曾經過產難的，你要高聲歡呼；因為沒有丈夫的，比有丈夫的有更多的兒女。」 \end{tabularx} \\ \\ \relax
4:28 & \begin{tabularx}{0.7\textwidth}{X} 弟兄們，你們是憑著應許作兒女的，如同以撒一樣。 \end{tabularx} \\ \\ \relax
4:29 & \begin{tabularx}{0.7\textwidth}{X} 當時，那按著肉體生的迫害了那按著聖靈生的，現在也是這樣。 \end{tabularx} \\ \\ \relax
4:30 & \begin{tabularx}{0.7\textwidth}{X} 然而經上是怎麼說的呢？是說：「把使女和她兒子趕出去！因為使女的兒子絕不能與自由婦人的兒子一同承受產業。」 \end{tabularx} \\ \\ \relax
4:31 & \begin{tabularx}{0.7\textwidth}{X} 弟兄們，這樣看來，我們不是使女的兒女，而是自由婦人的兒女了。 \end{tabularx} \\ \\ \relax
5:1 & \begin{tabularx}{0.7\textwidth}{X} 基督釋放了我們，為使我們得自由。所以要站穩了，不要再被奴隸的軛挾制。 \end{tabularx} \\ \\ \relax
5:2 & \begin{tabularx}{0.7\textwidth}{X} 我—保羅告訴你們，你們若受割禮，基督就對你們無益了。 \end{tabularx} \\ \\ \relax
5:3 & \begin{tabularx}{0.7\textwidth}{X} 我再指著凡受割禮的人確實地說，他有義務遵行全部的律法。 \end{tabularx} \\ \\ \relax
5:4 & \begin{tabularx}{0.7\textwidth}{X} 你們這要靠律法稱義的是與基督隔絕，從恩典中墜落了。 \end{tabularx} \\ \\ \relax
5:5 & \begin{tabularx}{0.7\textwidth}{X} 至於我們，我們是靠著聖靈，憑著信心，等候所盼望的義。 \end{tabularx} \\ \\ \relax
5:6 & \begin{tabularx}{0.7\textwidth}{X} 因為在基督耶穌裡，受割禮不受割禮都沒有功效，惟獨使人發出仁愛的信心才有功效。 \end{tabularx} \\ \\ \relax
5:7 & \begin{tabularx}{0.7\textwidth}{X} 你們向來跑得好，誰攔阻了你們，使你們不順從真理呢？ \end{tabularx} \\ \\ \relax
5:8 & \begin{tabularx}{0.7\textwidth}{X} 這樣的勸導不是出於那召你們的。 \end{tabularx} \\ \\ \relax
5:9 & \begin{tabularx}{0.7\textwidth}{X} 一點麵酵能使全團都發起來。 \end{tabularx} \\ \\ \relax
5:10 & \begin{tabularx}{0.7\textwidth}{X} 我在主裡深信你們必不懷別樣的心；但騷擾你們的，無論是誰，必須承受懲罰。 \end{tabularx} \\ \\ \relax
5:11 & \begin{tabularx}{0.7\textwidth}{X} 弟兄們，我若仍舊傳割禮，為甚麼還受迫害呢？若是這樣，十字架絆倒人的地方就沒有了。 \end{tabularx} \\ \\ \relax
5:12 & \begin{tabularx}{0.7\textwidth}{X} 恨不得那騷擾你們的人把自己閹割了。 \end{tabularx} \\ \\ \relax
5:13 & \begin{tabularx}{0.7\textwidth}{X} 弟兄們，你們蒙召是要得自由；只是不可把這自由當作放縱情慾的機會，總要用愛心互相服侍。 \end{tabularx} \\ \\ \relax
5:14 & \begin{tabularx}{0.7\textwidth}{X} 因為全部律法都包括在「愛鄰如己」這一句話之內了。 \end{tabularx} \\ \\ \relax
5:15 & \begin{tabularx}{0.7\textwidth}{X} 你們要謹慎，你們若相咬相吞，恐怕要彼此消滅了。 \end{tabularx} \\ \\ \relax
5:16 & \begin{tabularx}{0.7\textwidth}{X} 我說，你們要順著聖靈而行，絕不可滿足肉體的情慾。 \end{tabularx} \\ \\ \relax
5:17 & \begin{tabularx}{0.7\textwidth}{X} 因為肉體的情慾和聖靈相爭，聖靈和肉體相爭，這兩個彼此敵對，使你們不能做所願意做的。 \end{tabularx} \\ \\ \relax
5:18 & \begin{tabularx}{0.7\textwidth}{X} 但你們若被聖靈引導，就不在律法之下。 \end{tabularx} \\ \\ \relax
5:19 & \begin{tabularx}{0.7\textwidth}{X} 情慾的事都是顯而易見的；就如淫亂、污穢、放蕩、 \end{tabularx} \\ \\ \relax
5:20 & \begin{tabularx}{0.7\textwidth}{X} 拜偶像、行邪術、仇恨、紛爭、忌恨、憤怒、自私、分派、結黨、 \end{tabularx} \\ \\ \relax
5:21 & \begin{tabularx}{0.7\textwidth}{X} 嫉妒、醉酒、荒宴等類。我從前告訴過你們，現在又告訴你們，做這樣事的人必不能承受神的國。 \end{tabularx} \\ \\ \relax
5:22 & \begin{tabularx}{0.7\textwidth}{X} 聖靈的果子就是仁愛、喜樂、和平、忍耐、恩慈、良善、信實、 \end{tabularx} \\ \\ \relax
5:23 & \begin{tabularx}{0.7\textwidth}{X} 溫柔、節制。這樣的事沒有律法禁止。 \end{tabularx} \\ \\ \relax
5:24 & \begin{tabularx}{0.7\textwidth}{X} 凡屬基督耶穌的人，是已經把肉體與肉體的邪情私慾同釘在十字架上了。 \end{tabularx} \\ \\ \relax
5:25 & \begin{tabularx}{0.7\textwidth}{X} 我們若靠著聖靈而活，也要靠著聖靈行事。 \end{tabularx} \\ \\ \relax
5:26 & \begin{tabularx}{0.7\textwidth}{X} 不要貪圖虛名，彼此惹氣，互相嫉妒。 \end{tabularx} \\ \\ \relax
6:1 & \begin{tabularx}{0.7\textwidth}{X} 弟兄們，若有人偶然被過犯所勝，你們屬靈的人就要用溫柔的心把他挽回過來；自己也要留意，免得也被引誘。 \end{tabularx} \\ \\ \relax
6:2 & \begin{tabularx}{0.7\textwidth}{X} 你們各人的重擔要互相擔當，這樣就會成全基督的律法。 \end{tabularx} \\ \\ \relax
6:3 & \begin{tabularx}{0.7\textwidth}{X} 人若沒有甚麼了不起，還自以為了不起的，就是自欺。 \end{tabularx} \\ \\ \relax
6:4 & \begin{tabularx}{0.7\textwidth}{X} 各人要省察自己的行為；這樣，他所誇口的只在自己，而不在別人。 \end{tabularx} \\ \\ \relax
6:5 & \begin{tabularx}{0.7\textwidth}{X} 因為人人必須擔當自己的擔子。 \end{tabularx} \\ \\ \relax
6:6 & \begin{tabularx}{0.7\textwidth}{X} 在真道上受教的，要把一切美好的東西與施教的人分享。 \end{tabularx} \\ \\ \relax
6:7 & \begin{tabularx}{0.7\textwidth}{X} 不要自欺；神是輕慢不得的，因為人種的是甚麼，收的也是甚麼。 \end{tabularx} \\ \\ \relax
6:8 & \begin{tabularx}{0.7\textwidth}{X} 順著肉體撒種的，必從肉體收敗壞；順著聖靈撒種的，必從聖靈收永生。 \end{tabularx} \\ \\ \relax
6:9 & \begin{tabularx}{0.7\textwidth}{X} 我們行善不可喪志，因為若不灰心，到了適當的時候就有收成。 \end{tabularx} \\ \\ \relax
6:10 & \begin{tabularx}{0.7\textwidth}{X} 所以，一有機會就要向眾人行善，向信徒一家的人更要這樣。 \end{tabularx} \\ \\ \relax
6:11 & \begin{tabularx}{0.7\textwidth}{X} 你們看我親手寫給你們的字是何等的大！ \end{tabularx} \\ \\ \relax
6:12 & \begin{tabularx}{0.7\textwidth}{X} 那些想要炫耀外表的人才勉強你們受割禮，無非是怕自己為基督的十字架受迫害。 \end{tabularx} \\ \\ \relax
6:13 & \begin{tabularx}{0.7\textwidth}{X} 他們那些受割禮的，連自己也不守律法；他們要你們受割禮，不過是要拿你們的肉體誇口。 \end{tabularx} \\ \\ \relax
6:14 & \begin{tabularx}{0.7\textwidth}{X} 但我絕不以別的誇口，只誇我們主耶穌基督的十字架；因這十字架，就我而論，世界已經釘在十字架上；就世界而論，我已經釘在十字架上。 \end{tabularx} \\ \\ \relax
6:15 & \begin{tabularx}{0.7\textwidth}{X} 受割禮或不受割禮都無關緊要，要緊的就是作新造的人。 \end{tabularx} \\ \\ \relax
6:16 & \begin{tabularx}{0.7\textwidth}{X} 凡照這準則行的人，願平安憐憫加給他們，和神的以色列民。 \end{tabularx} \\ \\ \relax
6:17 & \begin{tabularx}{0.7\textwidth}{X} 從今以後，不要有人再攪擾我，因為我身上帶著耶穌的印記。 \end{tabularx} \\ \\ \relax
6:18 & \begin{tabularx}{0.7\textwidth}{X} 弟兄們，願我們主耶穌基督的恩與你們的靈同在。阿們！ \end{tabularx} \\ \\
[1ex]
\hline
\hline
\end{longtable}
$^{1}$各位電影姐妹平安.
今天我們來到加拿大書的第二講.
就是加拿大書的下半部.
保羅在他每一封信的下半部.
都是生活渲染一個實際實用的一個部分.
就是加拿大書的四,五,六這三章.
我給一個題目給他.
就是耶穌基督叫我得自由.
自由是很寶貴的.
有很多人極力要爭取自由.
民主自由.
言論自由.
生活有自由.
不受約束.
我不知道你有沒有被一些東西捆綁住.
可能我們都被一些壞的習慣.
被一些舊的東西.
將我們好像綁緊了.
綁著我們一樣.
但是耶穌說祂釋放我們.
叫我們在基督裡面.
我們得到自由.
那保羅為什麼這麼注重實踐呢?.
因為有很多時候我們聽道.
聽了道之後就不記得了.
我們有時候一群牧師,傳奴人走來聊天.
你猜你講完一篇道.
那些弟兄姊妹會將那篇道存留在他們裡面多久呢?.
一般的統計就說.
聽完一般道.
其實接著走出教會.
去吃午飯的時候.
或者如果是星期六參加周老崇拜.
接著晚上去小組.
或者回家吃晚餐.
就已經沒有了一半了.
然後到了第二天.
又沒有了五分之三.
然後再過一天.
基本上就沒有了.

$^{41}$所以有些講完之後.
很不會做就說.
我上個星期講的題目.
大家還記不記得啊?.
根本沒有人記得的.
所以最近有一本書是講港道的.
叫做《道猶在意》.
就是神的道好像仍然在你的耳朵裡面.
這本書講什麼呢?.
就是聽道和講道的神學.
原來聽道也是一門學問.
也是一個神學.
你懂不懂聽道呢?.
很多人都研究the art of listening.
聆聽的藝術.
原來聆聽是一門藝術.
我們說你要懂得聽別人說話.
因為你不是聽他說的話的表面.
你開始會觸摸到.
原來他說這些話的背後.
其實他當時的情緒.
心理的狀況.
或者他受到什麼影響.
以至他這樣說.
哇 我很討厭那些人.
很大聲說話的.
可能就是因為曾經有人很大聲嚇過他.
所以他對這件事很不滿意.
很難忘記.
各位.
今天我不知道你每一次去聽道.
你帶著一個什麼心態.
唉 就看看你們這些牧師傳的人說什麼.
很多我都認識的了.
又是三幅被 說來說去都是那些沒什麼新意.
我說得比你好.
可能很多時候我們帶著批判的心去聽道.
所以我們沒有一個謙卑的心.
沒有一個柔和的心.
那又怎樣領受神的道呢?.

$^{81}$所以求神幫助我們.
不要只是做一個會聽道.
會批評的人.
讓我們聽完道.
是藏在心裡.
聖經說 我將你的話藏在心裡.
免得我得罪你.
所以求主也幫助我們.
讓神的道在我們生命裡面產生功效.
發生力量.
我們能夠實踐出來.
那我說說保羅在第四章說什麼呢?.
就是說被救贖的人.
我們都有異類的名份.
你知不知道你是神的子女.
多麼的尊貴.
奴隸和子女是有分別的.
因為當時在希臘的時代.
很多奴隸的.
但是如果是子女.
你可以承受家族的產業.
特別是長子.
所以長子的名份.
在聖經裡面很重視的.
因為你就承受父親雙倍的祝福.
承受家族遺留下來的產業.
而且聖經告訴我們.
我們是異類的身份.
我們就需要記得.
我們是活在基督裡面.
不要枉費了保羅.
浪費了在我們身上那種奴婦.
以致我們好像回復舊狀一樣.
還仍然被奴僕的扼押克制.
因此保羅就用亞伯拉罕.
他有兩個兒子.
兩個兒子的媽媽.
來做一個例籍.
亞伯拉罕第一個生的兒子.
是他一位妾侍夏甲所生的.

$^{121}$這個兒子的名字叫以實瑪利.
這個以實瑪利其實就是以後阿拉伯人的祖先.
第二個兒子就是神應許的兒子.
他說我要賜你一個兒子.
他必成為大哥.
而且你要領受很多的祝福.
這個兒子就是以實.
以實瑪利的媽媽就是夏甲.
夏甲後來被亞伯拉罕的太太薩拉趕了出去.
他很辛苦.
去到一個地方都沒有水喝.
在那裡求死.
他抱著一個很小的嬰兒.
又被家裡趕走.
好像沒有路可走.
神仍然眷顧他.
開了他的眼睛.
他就看到沒有遠的地方有口水者.
於是可以打水.
就不用渴死了.
以實因為他是應許的兒子.
所以就得到父親的祝福.
於是他建立了一個大國.
承受了亞伯拉罕的產業.
各位弟兄姊妹.
你知不知道你是神的子女.
有很尊貴的身份.
有時我自己祈禱時.
詩上中阿伯父說.
你怎麼這麼疼我.
我是你的兒子.
在我未祈禱以前.
爹地已經知道子女要什麼.
甚至我們未開口祈禱.
他已經知道.
他已經答應我們.
他已經預備了祝福賜給我們.
所以各位我們要持有我們神兒女的名份.
領受上帝給我們的祝福.
故此我們可以說.

$^{161}$這個應許的兒子.
是代表神給我們很多的恩典承受產業.
第五章.
第五章不要被奴僕的騙.
來仍然克制我們.
開頭我舉例提到.
你好像僕人一樣.
記不記得浪子回頭的故事.
那個在外面花了錢.
沒東西吃.
這個很慘的浪子.
他就想起爸爸.
於是他說我要回家.
就作定編台詞.
爸爸我不配做你的兒子.
我要做你的奴弟.
我們都服侍你.
爸爸老遠就走出來.
抱著他.
把最漂亮的袍子戴在他的身上.
把戒指戴在他的手上.
說你是我的兒子.
各位原來我們不是奴僕.
我們是神的兒女.
如果我們仍然接受割禮.
我們就代表我們仍然想舊的厄賢俗下去.
因此我們要將自己排除起.
這些奴僕的惡之外.
我們要信服神.
並且是信服神.
在我們身上給我們的帶領.
律法都包起愛人如己.
這句說話上面.
所以耶穌有律法師問他.
舊律法的總綱是什麼?.
耶穌就說.
你要愛神.
其次也相仿.
就是愛人如己.
今天我們有沒有學習.

$^{201}$愛人如己這個功課?.
我們要順著聖靈而行.
我們就不放縱肉體的程慾.
然後第五章的二十二節是很出名的.
就是聖靈所結的果子.
其實果子是單數的.
The fruit of the spirit.
沒有s的.
沒有雙數的.
就是說一個果子.
它有九種特質.
有些人就說有九個果子.
其實一個果子.
它有九種的特質.
大家一定要背了.
這九個聖靈所結的果子.
有時候我也背得不好.
所以我們要知道.
Love joy peace.
人愛喜樂和平.
溫柔信實節制.
然後就是.
溫柔信實節制.
調轉了.
應該忍耐恩慈良善.
然後信實溫柔節制.
人愛喜樂和平.
溫柔忍耐恩慈良善.
信實溫柔節制.
耶穌基督就叫我們在各方面都要成長.
所以有些門徒訓練.
就是怎樣訓練我們在這九個特質裡面.
我們都不斷地成長.
當然有些人某幾方面是長處強項一點.
有些地方是我們的弱項.
但是求神加恩典給我們.
給我們全面的成長.
最後第六章.
第六章講到.
一開頭就說弟兄們.

$^{241}$如果有人偶然被過犯所成.
你們光強的人.
就應該將他挽回過來.
而且你們要小心.
免得自己也被引誘.
如果有人偶然被過犯所成.
跌倒了.
挽回過來.
主教會很注重挽回那些失腳的弟兄姊妹.
然後神就說.
我們各人的重擔.
是要互相來擔當.
有些叫做「重擔」.
這個字在原文裡面.
就真的是超重的東西.
很重的東西.
然後各人就擔自己的背包.
有時有些擔子.
自己要揹自己的背包.
每一個軍人.
都有一個行軍的時候的背包.
裡面有水有各樣東西.
然後你又拿你的步槍.
所以那個背包是各個人揹.
但是有時長官就吩咐.
你們四個人幫忙抬大炮.
你們四個人要搬些什麼.
於是四個人就共享責任.
一起去擔一個重擔.
很重的擔子.
還有告訴我們.
我們要知道.
不單止是一起去承擔這些重擔.
然後保羅就提醒我們.
人種的是什麼.
收的就是什麼.
如果我們順著情欲殺種.
收的就是敗壞.
如果我們人若順著聖靈去殺種.
我們收的就是在基督裡面.

$^{281}$能夠得到永生.
然後保羅在第十四節說.
我斷不以別的誇口.
我只誇耶穌基督和祂的十字架.
保羅高舉耶穌基督的十字架.
然後在結束的時候.
保羅就說.
從此以後人都不可以再搞擾我.
因為在我身上有耶穌基督的印記.
很寶貴的一句話.
你不可以搞擾我.
因為我身上有基督的印.
我是屬基督的.
耶穌的印象在我的身上.
請問今天別人從你的身上見不見到耶穌.
我記得以前唱的一首歌就是.
使別人由你或者由我見到耶穌.
求主都幫我們.
讓我們在我們的身上彰顯耶穌基督.
顯基督偉大.
使別人看見耶穌在我的身上.
願天父賜福給各位.
\newpage



\section{歌羅西書 1:24-29}
\label{sec:031762}
\textbf{事奉主的喜痛交雜}~~~~黃紹權牧師
\newline
\newline
連結: \href{http://www.church.com.hk/acms/content.asp?site=cdc\&op=show\&type=product\&code=031762\&layout=sermon}{\texttt{031762}} ~~~~ 語音日期: 2026-03-08
\newline
\newline
\hyperref[sec:031761]{\small{< < < PREV SERMON < < <}}
~
\hyperref[sec:index_chronic]{\small{[返順時目]}}
~
\hyperref[sec:031763]{\small{> > > NEXT SERMON > > >}}
\newline
\newline
歌羅西書 1:24-29
\newline
\begin{longtable}{cl}
\hline
\hline
章節 & 經文 (和合本修訂版)\\
\hline
1:24 & \begin{tabularx}{0.7\textwidth}{X} 現在我為你們受苦，倒很快樂；並且為基督的身體，就是為教會，我要在自己的肉身上補滿基督未盡的苦難。 \end{tabularx} \\ \\ \relax
1:25 & \begin{tabularx}{0.7\textwidth}{X} 我照神為你們所賜我的職分作了教會的僕役，要把神的道傳得完滿； \end{tabularx} \\ \\ \relax
1:26 & \begin{tabularx}{0.7\textwidth}{X} 這道就是歷世歷代所隱藏的奧祕，但如今向他的聖徒顯明了。 \end{tabularx} \\ \\ \relax
1:27 & \begin{tabularx}{0.7\textwidth}{X} 神要讓他們知道，這奧祕在外邦人中有何等豐盛的榮耀；就是基督在你們心裡成了得榮耀的盼望。 \end{tabularx} \\ \\ \relax
1:28 & \begin{tabularx}{0.7\textwidth}{X} 我們傳揚他，是用諸般的智慧，勸戒各人，教導各人，要把各人在基督裡完完全全地獻上。 \end{tabularx} \\ \\ \relax
1:29 & \begin{tabularx}{0.7\textwidth}{X} 我也為此勞苦，照著他在我裡面運用的大能盡心竭力。 \end{tabularx} \\ \\
[1ex]
\hline
\hline
\end{longtable}
$^{1}$本節目由我們在聆聽主的說話而設.
我們一起同心低頭禱告.
我們的得救為上帝你的心願.
而去甘願受苦死在十字架上.
讓我們在這個時候再次為了這件事獻上感恩之言.
讓我們每一個都能夠好好地學會上主這份事奉的心志.
讓我們能夠在我們今天屬靈人生的路上所行所想所作.
都是滿有明白,體貼上主心思而去實踐的事情.
禱告奉靠主耶穌基督德性的名字而靠.
阿們.
移民很多人都為了去找工作而去酬賬.
我在加拿大這幾十年的日子裡面.
其實在工作上雖然是主要以事奉主為我的基礎.
但是有時候因應著財政的需要.
也都試過有很多不同的在加拿大的工作機會.
在英語世界裡面工作.
其實我想第一個挑戰就是對於這邊的文化或者語言上等等的種種挑戰.
我記得我有一份在英語世界的工作.
就是做一個移民中心的一個全日的工作.
interpreter.
大家順帶講一講.
interpretation和translation是兩種工作.
interpreter除了要將人說的話轉成另一個語言之外.
同時要捉到那個語氣.
這個是不容易的工作.
以及那個反應的快速是很大的需要.
所以當時為了要做這份工作.
我就要去考一個牌照.
要去考一個全日牌.
全日牌在加拿大特別安省還有兩個牌你要考的.
一個就是一般的全日牌.
另一個全日牌就是專門為了家庭暴力去做這個牌照.
考完這個考試才可以去真正做一個即時全日者.
我記得我第一份找的即時全日工作就是一個移民中心.
是一個所謂社福界的一個移民中心裡面的全日員.
但是全日不是這麼簡單的一件事.
不只是一個人去到懂得一個語言就可以的.
有很多的訓練在職訓練.
就在第一個星期基本上就很忙碌的.
除了要去觀摩人家其他的全日怎樣做之外.

$^{41}$因為我就做普通話一國.
和廣東話翻譯英語雙向翻譯.
所以那個即時的反應是很大的需要.
我記得我第一次做的時候.
全日沒多久就十多二十分鐘左右.
我真的要睡一個多小時才可以恢復元氣.
因為用腦是用得很厲害的.
因為這麼要求高的工作.
那時候我在這個移民中心就有一個導師.
專門去訓練我去做這件事.
他用了一個星期左右去訓練.
但是訓練到最後一個星期.
他給的評語就是說.
其實Alan就不用再訓練了.
因為他好好預備好做這份工作.
我想或者可能我的語言能力足夠他覺得.
但是當去到最後一天的時候.
但是當說了沒多久之後.
我們那個全日中心的主管就走出來.
他用了一個宣告.
是不是宣告我畢業呢?.
是不是有獎金,有獎品?.
原來不是.
我想和大家一起宣佈一個消息.
今天是我們一位同事的最後一天的工作.
我心想我剛剛來一個星期.
接著就有一個同事要走.
是誰呢?.
我似曾相識的一個字眼.
叫做「I would like to announce a bitter joy」.
我很熟悉這個字眼.
如果用中文翻譯來算.
這個字眼就是指一個悲喜交集的情況.
原來這位訓練我的導師.
他訓練我這一星期.
原來也是他人生當中最後一個星期的工作.
原來他訓練完我最後一天.
接著他就正式退休.
為什麼是一個叫做「bitter joy」呢?.
其實最主要的一個原因就是.

$^{81}$一方面公司會少了一個好的傳譯員.
因為他懂得三種語言.
除了英文.
所以其實我們公司差不多很多工作.
都背在他身上.
沒有一個懂得三國語言的人.
其實那個分分鐘可能要請三個人來頂一個人.
尤其是這些社福界的就沒什麼錢的.
就是薪水按到你低一低.
工作就多一多那些情況.
那這位同事離開其實是一種很大的損失.
但是回心一想.
當一個人踏入去退休的日子.
其實也是他人生另一個的片段的開始.
所以其實退休就不是不用做.
退休只不過是你的角色有改變.
你的人生的步伐可能會有一個改變.
所以為什麼要稱之為Bitter Joy呢?.
其實就是說到.
雖然是一件不開心的事對於一個公司.
但是對這位同事是一件很大的喜事.
其實這位同事後來因為他年紀大.
有很多電腦的東西.
那時候他後來退了休.
我也有去過他的家.
他其實我不知道是不是政府那麼好安排.
就是他退休之後的那個星期他就搬屋.
搬去老人屋.
我去到那個老人屋其實是很漂亮的.
是自己住的一個單位.
其實是一個新建的老人宿舍.
是很精緻.
所以這個字眼一直都擺在我心裡.
人生裡面其實有很多的事物都是一個.
我們稱為Bitter Joy.
悲喜交集的一個人生經驗.
也可以這麼說.
我去到這個公司的工作.
不多不少都是去承擔了他一些之前的工作.
但是這種世代相傳的過程.

$^{121}$其實好像真的一定要經過了這個悲喜交集的過程.
才能夠有一個交棒的時間.
為什麼我分享這個經驗呢?.
其實橫教會經常特別我們這個年紀.
我就很不亂的情況.
要退休的還有十年左右.
身邊很多的教務同工都比我大.
因為我們這些年輕人出來事奉.
貼近那些老人家多.
我屬於老人家那一批.
但是又不是他的年紀.
經常聽見他們說要接班.
繼承教會那個承傳的工作.
我覺得的而且確是這樣的.
但是教會有沒有真的預備好.
在這些交接的過程裡面.
可能會有一些正如我上司所說的.
所說的叫做bitter joy.
悲喜交集的轉接過程呢?.
還是我們無論交那個也好.
又不肯放手.
接那個又接不到手.
在這個過程當中反而不是悲喜交集.
而是大家互相嘔鬥呢?.
這個情況我想是很值得我們去思想.
而這個亦都是今天我們所看的.
聖經的經文在《歌羅西書》第一章.
從第24節至第29節裡面.
保羅正正談到一個他事奉的轉接.
保羅已經年紀大了.
已經是甚至下在監獄裡面.
身處在羅馬城的當中.
但是他沒有辦法去繼續牧養一間.
雖然不是他自己親手所建立的教會.
但是至少亦都是透過保羅認識主的一些童宮.
譬如我們所熟悉的以巴弗.
他作為一個在哥羅西這個城市.
一個這麼重要的小亞細亞.
或者亞洲的重要交通商業中心.
而去到羅馬認識了保羅.

$^{161}$甚至認識了基督.
信了主的信徒.
他回到哥羅西城去開展.
或者參與在這個教會的事奉的過程裡面.
保羅其實同樣要明白到.
他的事奉不是只是單靠自己的直接參與.
其實是真的需要好好地去交接.
給下一代在哥羅西的教會.
而以巴弗以及他的童宮.
其實正正就是需要面對一個這樣的轉接的過程.
而當我們去看歌羅西書今天這段的聖經.
亦都反映出保羅這份對教會那種傳承的矛盾.
和那種的喜悅.
怎樣兩者之間能夠可以去看到上帝的心意.
所以我將今天這段聖經的講題.
定之為事奉主的喜樂交集.
來成為我們今天這段聖經的思想的主要核心的題目.
讓我們作為屬臣的子民.
亦都抱著一種好像和保羅一樣的心思意念.
來去承傳教會.
這個應當將主導彰顯.
牧羊的角色能夠好好地做好.
讓我們一起回到這段聖經歌羅西書第一章第24至29節.
讓我們一起再讀一次.
雖然Karen已經幫我們讀過一次了.
讓我們再讀多一次.
因為看來這段經文裡面所講的東西有些似曾相識.
但又好像有一種若即若離不太掌握的內容.
讓我們一起去再讀一次.
請大家一起打開歌羅西書第一章.
我們從第24至29節一起讀.
好,找到之後我們一起讀.
把神的道理傳得傳備.
這道理就是歷世歷代所隱藏的奧秘.
但如今向他的聖徒顯明了.
神願意叫他們知道.
這奧秘在外邦人中有何等豐盛的榮耀.
就是基督在你們心裡成了有榮耀的盼望.
我們傳揚他,使用諸般的智慧.
勸戒各人,教導各人.

$^{201}$要把各人在基督裡完完全全地引到神面前.
我也為此勞苦.
照著他在我裡面運用的大能盡心竭力.
我們回顧一下上文我們講些什麼.
上文其實當我們去看回那段保羅的內容.
主要是想告訴我們.
其實真理本身的啟示不是很困難的.
不是說真的很刁鑽.
只有讀過很多書的人才會明白.
不是的,其實那個道理是很簡單的.
但是為難的就是大部分的情況是人不願意離開罪這個決心.
於是乎當我們缺乏了這份決心要離開罪的時候.
我們就會發現到我們的信就開始變質.
尤其是今天我們的相信.
究竟是不是直根於認罪悔改這個基礎.
還是以其他聖經以外的原因來所謂信了耶穌呢?.
其實這方面是很值得我們去深思的.
所以今天有些弟兄姊妹說牧師我是不是真的信了主.
我覺得他問的這條問題其實是一個很深思熟慮的一個考量.
如果一個人求求其其的信了主.
就說我做很多事有神保守有危險.
上帝要我離開這些危難困難.
這個神很靈所以我信他.
其實後者究竟是不是真的認罪悔改.
相信主而生命有真正的改變呢?.
這個是很值得我們思想的.
而保羅也都給我們看見.
在這個事奉的過程裡面.
他不介意去直接將這個認罪悔改這個基礎告訴信徒知道.
而且他甚至去承擔了一個執事.
福音的執事的角色.
來去將這個福音準確傳開.
而這個也是我們今天所看的這段經文的背景.
保羅在一個這樣的大前提之下.
究竟他怎樣去真正執行這個執事的工作呢?.
用今天的用語就是說作為教會的領袖.
甚至教會的每一個信徒.
我們是怎樣去真正執行上帝那個的救恩的啟示.
福音本身那個的傳遞.
怎樣才做得真是準確呢?.

$^{241}$這個也是我們這一段聖經裡面所看見的進一步保羅的申述.
如何作好這個福音的執事的這個行列.
已經參與了.
沒有說等啊求求其其過日子算數.
真是肯承擔起這個執事的責任了.
但是這個福音的執事的責任是怎樣執行呢?.
首先說一些結構上面的一些的陳述先.
如果在聖經的經文的文法的結構上面.
從第一章的第24節一到去到二章的第五節本身是一個整體的記載來的.
而它以一個叫做交叉式的一個陳述記事方式.
我已經說過很多次了.
這種無論你叫做交叉式的記載方式.
或者叫做夾心餅.
或者英文叫做Kiasic Structure.
或者叫做Sandwich Structure.
它是將傳統猶太人寫東西陳述的記載.
以夾心餅的形式去寫.
就是它一層一層去展述.
然後再將那個展述的結果說出來之後.
又再將那個展述的結果引申怎樣發展出來.
這個夾心餅的記載有一個好處的.
就是當讀者讀的時候就能夠容易去記載.
並且在介紹的時候就會更加容易去記得內容.
那你說牧師大約是什麼意思我都get不到.
那我就用唱流行歌的歌詞來做一個比喻.
大家有沒有留意現在絕大部份的流行歌曲.
都是將同一堆的歌詞重複唱多一次.
為什麼是這樣?.
因為這些歌我們現在稱之為洗腦歌.
它這樣唱的時候.
就能夠將那個歌詞的內容.
重複地刻在你的腦海裡更加容易.
Chiasic Structure其實就是類似這樣的意思.
不過今天我們不看整個的結構.
我們只看頭一半.
我們怎樣看頭一半呢?.
我們會分為三個部份去看的.
第一個部份.
第一章第二十四節.
一開始保羅說.

$^{281}$當他要執行這個福音的執事的身份的時候.
或者這份工作.
或者說得更加神學性的用詞.
事奉的時候.
他就道出了原來是需要經歷多重的為難.
但是他是很歡喜快樂地去到接受這些為難.
並且這些經歷能夠使他有真正事奉的喜樂.
我們一起看看這節經文說些什麼.
第二十四節這樣說.
「現在我為你們受苦,賭博歡樂.
並且為基督的身體.
就是為教會.
敲在我肉身上.
補滿基督患難的欠缺」.
我們做《釋經》研究的人.
在《聖經》裡面有.
大約十多段的《聖經》.
我們稱之為.
The Clux of Interpretation.
什麼意思呢?.
Clux這個字.
Clux就是一些死結纏繞著.
在《釋經》上最難解釋的《聖經》.
大約有十多段左右.
當然有很多可能大家都覺得很難.
但是在《釋經》上真是很難解釋的十多段經文.
第二十四節就是其中一節.
如果一般的信徒就這樣讀.
OK 搞定了 我明白的了.
你是不是真的明白?.
這是另一回事.
但是如果說《釋經》上.
說傳道的上.
很多牧者也都讀到這段《釋經》.
本身歌羅西書已經很少人說的了.
但是這句這麼難解釋的經文.
讀了其實實質上你仔細看看.
你其實真是知不知道他想說什麼.
你未必想知道.
未必真是掌握到他在說什麼.

$^{321}$我細心讀多一次.
或者慢一點讀多一次.
你自己去想一下.
究竟這節經文想說些什麼.
他說.
現在我為你們受苦.
賭博歡樂.
並且為基督的身體.
就是為教會要在我肉身上.
補滿基督患難的欠缺.
你真是明白嗎?.
我們逐節逐節看看他說什麼.
逐句逐句.
他說我現在是為你們受苦.
這句話說出來.
有些亂解經的人就會怎樣.
這裡教導我們.
聖經做傳道人要受苦.
所以我們做會友.
或者做教會執事.
教會領袖的.
就要使你的牧者有這麼苦得這麼苦.
這樣你才能夠滿足到保羅的要求.
還有講後面那句.
賭博歡樂.
你越虐待他.
對不起.
你越對他不好.
他就越歡樂.
這句話如果是這樣解釋.
就將所有傳道人演繹成為被虐狂.
他們是神經的.
他喜歡被人虐待.
是不是這個意思.
我這樣說你都會說不是.
但不是又為什麼呢這句話.
為什麼保羅說我為你受苦.
不是呀第二個解釋就是.
保羅想高他自己有多苦.
所以為什麼有些傳道人又出這一招.

$^{361}$叫做苦肉計.
搞到自己有這麼苦得這麼苦.
使教會的人特別會有.
覺得牧師很苦.
牧師很慘.
所以我們要對他好.
天主教很喜歡使這一招.
我認識很多神父.
很多牧師都是這樣.
講到自己有多可憐.
然後就怎樣.
讓會友覺得用憐憫牧者的心來接待牧者.
我曾經認識一個天主教教會的教區主教.
他曾經經常說我很窮我很窮.
我見他又不像很窮.
為什麼.
因為我每天見到他的時候.
經常有些人駕著路斯萊斯.
Benz 寶馬.
雖然我們有沒有人駕這些車.
沒有吧.
來接他去.
有一次我跟他說.
為什麼你出出入入.
又路斯萊斯又Benz又寶馬呢.
他說沒有呀.
會教友來接我去打球.
吃東西喝茶之類.
又有些傳道人是這樣的.
搞到自己好像很可憐.
然後怎樣.
然後就要會友來對他.
多多施下.
這樣去醒他.
其實我不怕.
告訴大家.
有信徒也是用這種方法.
來去以為.
私封神.
我曾經說.

$^{401}$曾經試過有信徒這樣做.
對我.
牧師你很窮的.
所以呢.
我給你一些錢.
你千萬不要給教會.
你給我我也是奉獻給教會.
不要呀你收回你自己.
你自用.
我說不要呀.
你給我我也是給教會.
然後他就.
我生你的氣.
你不給我機會去接待你.
真是慘呀.
我不怕.
我做你們牧師.
我不是很窮.
但我不是很富有.
足夠有餘.
但也不是.
富有到怎樣情況.
這個聖經究竟想說什麼.
本來在說我自己有多苦.
但我又覺得很開心.
來到去.
用謙卑.
來去羞辱信徒.
或者用謙卑來去施壓.
給信徒呢.
尤其是保羅和.
哥羅西教會沒有關係.
其實沒有直接關係.
是不是用這種方法來去.
希望.
後來的經文講到.
可以塞回一些腳.
賺回一些.
奉獻.
是不是這樣.

$^{441}$所以這句聖經很難解.
還有.
不是難解在這兩節.
我們看後面.
並且為基督的身體.
就是為教會.
要在我肉身上.
補足基督患難的欠缺.
究竟想說什麼.
我全部字都懂解.
並且為基督的身體.
全部字我都懂.
原來基督的身體.
就是教會.
教會就是基督的身體.
要在肉身上.
補滿基督患難的欠缺.
哇.
保羅你好厲害.
你可以用肉身.
去補足基督患難的欠缺.
即是耶穌死.
受害都不夠.
要找你來在肉身上.
補足基督.
難不難解.
我繼續說下去.
你會覺得.
我根本不明白這句聖經.
想說什麼.
我知道你很難解.
所以.
我用了一些時間.
我將這句聖經.
用原文重新.
翻譯一次.
希望我的翻譯.
能夠幫到你更加準確.
去知道這句聖經.
說什麼.

$^{481}$保羅.
他這樣說.
他現在.
我滿心喜樂.
即使.
我一直畏著.
哥羅西.
或者宗教會.
你們的人.
深陷多重為難.
或者不幸.
或者難堪.
但我的肉身.
仍然可以.
充充實實.
或者抖擻.
為的是什麼.
是為耶穌基督這個身體.
就是你們教會.
叫你們從原本.
缺乏的.
都可以在.
主耶穌基督.
為難當中.
再次得到好像我保羅一樣.
的充實.
可能我這樣翻譯完之後.
你都會覺得牧師.
很刺鳴多一點.
但其實說的是什麼.
那我就再用廣東話.
翻譯一次給大家聽.
那你就會明白了.
保羅.
我現在真是.
很開心.
可以怎麼樣.
為你們.
能夠付出我很多.
的難堪.

$^{521}$或者.
不想做的事.
而我.
因此也都可以.
繼續在體力上面.
有充沛的能力.
使到.
你們這個基督的身體.
就是你們宗教會.
宗信徒.
叫你們可以在.
原本.
不能夠.
充足地生活.
的狀態之下.
得到堅固.
好像主耶穌基督.
透過他的為難.
使到你們得到.
重新堅固.
簡單一點來說.
保羅想說的.
內容就是.
想告訴.
哥羅西教會的信徒.
知道.
他自己.
所經歷的事情.
不是自己拿來.
而是想去.
滿足耶穌基督.
使到.
交託給他.
使教會在牧養上面.
得到堅固.
這個目的.
用個例子.
你就會容易明白一點.
如果大家做過父母.
我想一定會明白.

$^{561}$就是.
照顧小孩.
你知不知道照顧小孩.
其實他還沒有出生之前.
你已經在照顧他.
用了很多心力去照顧他.
不過.
我們將時段特意.
放在他出生.
之後的頭幾個月.
你知不知道.
我已經很幸福.
我兩個小孩.
出生之後.
他們很快.
會調節到.
日與夜.
大約兩三個月.
什麼意思.
兩三個月內.
你.
日與夜都要照顧住.
他的需要.
他沒有分日夜.
所以.
日間要餵奶.
晚上也要餵奶.
晚上餵奶不是一般的餵奶.
不是一般的時間.
媽媽你還記不記得.
晚上.
小孩.
剛剛十點鐘.
餵飽奶之後.
睡到明天六點七點鐘.
沒有的.
他兩點鐘就來叫醒你.
兩點鐘過後.
四點鐘沒多久.
六點鐘就來叫醒你.

$^{601}$六點鐘之後.
八點鐘就來叫醒你.
一路就是這樣.
你慶幸的就是.
由兩個鐘變成三個鐘的分別.
三個鐘變成四個鐘的分別.
四個鐘的分別就變成五個鐘的分別.
你希望就是這樣的日子.
你知不知道.
那個難處是多少.
只是不能睡.
什麼都難受.
我記得那時候.
頭一個月.
師母說很累了.
頂不住了.
不如你來弄弄他.
你愛你的妻子.
你就要起來來弄弄她.
你愛你的妻子你就要起來來弄弄她.
你明天又要上班.
又要開會.
又要準備講道.
又要做程序.
這樣的為難.
就是保羅這裡所說.
所以保羅的書信裡面.
很喜歡用.
父母的情況.
生產之苦.
類似這些陳述.
去講到.
基督怎樣願意.
毀身來去.
付出一些.
其實未必是一般人.
都想願意付出的事.
為了家會.
就好像父母為了.
兒女的成長.

$^{641}$所以在所不貴.
睏都要起來來.
餵奶.
餵奶不止.
你有沒有試過半夜沖涼.
和他.
因為他賴了.
你全身髒.
你抹不乾淨.
你就要捧盆水.
去洗手間.
跟著和他沖完涼.
擦好身又怕他冷.
擦好身之後.
他又來肚餓.
你整晚都沒睡.
我兩三個月搞得定.
但有些搞一年半載.
那你怎樣頂.
所以你知不知道.
我們有第二個小孩的時候.
第一樣真的祈禱.
我不知道師母有沒有祈禱.
希望他像女兒一樣.
快點不用我們餵奶.
這種真是.
逆人性.
很難受.
如果用和合本的翻譯.
就叫受苦.
講得好像很嚴重.
其實這個字眼.
可以翻譯為多重為難.
很多種東西混在一起.
使我很難過.
難過不是sad.
你的生活有很多困難.
或者不幸.
這個字可以解釋為.
有小孩就不是不幸.

$^{681}$是慶幸.
有些想生都生不了.
難堪難過.
但你為了他的好處.
為了教會的好處.
保羅願意去付出.
而他的付出不是他付出了.
去補足基督有什麼缺陷.
基督已經成就了這個救音.
他將餘下的工作.
交給我們.
去幫助.
哥羅西教會的成長.
所以從這段聖經裡.
難的.
不是經文的字眼.
怎樣陳述.
難的是我們怎樣去理解.
保羅那份願意付出的心志.
弟兄姊妹做基督徒不是坐在那裡.
受了浸.
教會有很漂亮的彩色玻璃.
教會有冷氣.
冬天幸好有暖氣.
不是的.
教會如果在一個.
undesirable situation裡面.
宗信徒.
包括牧者.
有沒有願意一起去面對.
去處理這些挑戰.
簡單的說.
教會為了.
搞網上直播.
我廢寢忘餐.
我今早去洗臉.
還在想.
有什麼方法.
我現在連Starlink都想了.
衛星的上網.

$^{721}$當然了.
這些東西要花錢的.
不到我們的話.
那種.
每一個弟兄姊妹.
衛冬和我一起去想.
弟兄姊妹在想.
那個怎樣搏那個怎樣搏.
那個怎樣搞.
都是為了眾弟兄姊妹.
能夠有一個.
好的崇拜而去考量.
是我們有高清重播.
影片還是精心.
切割有字幕.
有圖片.
始終如果有弟兄姊妹.
想願意和我們一起崇拜.
我們考慮到這些東西.
這種心態.
是我們作為教會應該有的.
而保羅在這裡就講到.
這個為了.
信徒弟兄姊妹.
在上帝旨意之下.
能夠更加成熟.
就好像一章第22節裡面所講.
更趨成熟的.
屬靈生命.
而我們每一個人都要作出.
一些取捨.
而這些取捨可能未必是.
順著我們喜歡做的事.
或者我們的preference.
我們的喜好而去做.
就好像父母.
個個都想睡晚一點.
是不是.
最好.
公司上班.

$^{761}$就是十點鐘開工.
我睡到九點四十五分才起身.
收工最好早一點.
三點多.
餘下的時間.
我可以早點回家.
看電視.
很多人都是這樣想.
在家裡不用煮飯.
有人煮飯.
開飯.
桌子都不用我set up.
下來就像少爺少奶.
然後就可以.
吃完之後還不用洗碗.
個個都想做這些事.
是不是.
但不是的.
你為了家庭為了你的子女.
為了你的家人的好處.
願意去騰出一些時間.
去做工作.
一起去煮飯.
一起去開桌.
一起去吃飯.
保羅也是一樣的心態.
為了教會.
整個群體的好處.
補足他們.
自己屬靈生命的不足.
然後.
他願意擔當.
這個福音的.
執事的身份.
來幫助他們.
能夠成長.
即使保羅.
坐下監或者受下苦.
又何嘗.
不可以呢.

$^{801}$這個就是.
保羅所傳遞的訊息.
講到對父母的心態.
中國人有句說話.
叫做養兒方知父母心.
如果一個真正的.
好的兒女.
是不需要到自己養兒的時候.
都已經曉得.
明白父母的心思.
這個就是一個有智慧的人.
如果要到養兒才知道父母心.
都不錯的了.
都OK.
好過不知.
好過知都不對父母.
都是不理不睬.
作為信徒.
我們明白.
上帝.
為什麼有一個這樣的兒子.
叫耶穌.
願意寫在十字架上.
他沒有需要做的事.
但為了父神的.
心思意念.
拯救罪人脫罪.
這條通道能夠打開.
他願意受苦.
保羅.
他的事奉在這裡告訴我們.
他只不過是.
仿效.
耶穌基督.
願意.
為教會.
為基督的身體的成長.
好處而付出.
作為信徒.
第二部分.

$^{841}$第一章第25至27節.
25至27節這樣說.
我召神為你們.
做一次我的職份.
作了教會的執事.
要把神的道理.
得傳遍.
這道理.
就是歷世歷代.
所隱藏的奧秘.
但如今.
向他的聖徒顯明了.
神願意.
叫他們知道.
這奧秘在.
外邦人中有何等.
豐盛的榮耀.
就是基督在你們.
心裡成了.
有榮耀的盼望.
在這一部分裡.
保羅.
想說的是.
如何能夠.
完美地實現.
神的話.
就是將上帝的.
奧秘.
解釋清楚.
給哥羅西教會.
知道.
是否有幸可以讀到.
同樣我們要知道.
透過保羅這番說話.
去明白神的奧秘.
是甚麼東西.
在這裡再一次.
保羅第25節說到.
他是照著神所賜給他的職份.
來作教會的.

$^{881}$執事.
保羅在這裡說得很清楚.
他作為教會.
的執事.
他不是透過.
人的意願.
即教會的意願.
或者他自己的意願.
來做事.
他是要靠著甚麼東西.
是照著神你為我所賜的職份.
簡單來說.
保羅.
是需要在.
上帝的.
督導之下.
來完成上帝的.
心願.
要怎樣.
保羅第25節.
這道理就是歷世歷代.
隱藏的奧秘.
但如今要向.
他的聖徒顯明了.
原來保羅.
他要擔當的工作.
就是要將.
上帝的心願.
或者上帝.
在歷世歷代.
所說的話全部地.
準確地去傳遞給他們知道.
而這個訊息.
隱藏了很久.
歷世歷代的裡面.
都隱藏了.
甚麼叫歷世歷代.
歷世歷代這裡.
包括了兩方面.
第一方面就是指到.

$^{921}$時間上面.
很長時間都隱藏了.
而另一方面就是.
指到在世世代代.
裡面都很多人未明白.
即是在.
年代世代裡面.
都隱藏了.
為甚麼要隱藏.
因為這裡說到這是一個奧秘.
但如今要讓你知道.
保羅要做這個功夫.
那你就會在這裡問.
牧師那甚麼叫奧秘.
因為接著27節又再講奧秘.
你講清楚甚麼叫奧秘.
因為神願意叫你們知道.
這奧秘在外邦人當中.
有何等豐盛的榮耀.
這就是基督在你們心裡.
成了有榮耀的盼望.
又講奧秘.
這是甚麼奧秘來的呢.
首先這段聖經我們.
會明白到多一點點的資料.
就是原來.
哥羅西教會是一間.
充滿外邦人的.
教會.
不代表他沒有猶太人.
他有猶太人.
因為其實.
我們仔細看看之後的經文.
開始講很多猶太人.
的律法的問題.
這在第二章裡面.
但當我們仔細看看這段聖經.
原來有很多外邦人.
現在上帝要將一個奧秘.
要說有些榮耀.

$^{961}$是可以來到你們身上.
那是有甚麼奧秘.
奧秘.
這個字眼.
很多人就講到.
是一些很神秘的.
很神秘的.
甚至涉及一些.
可能神蹟奇事.
但如果當我們.
仔細去查考一下.
這個字眼奧秘.
其實是和舊約.
的神學有直接的關連.
最明顯不過的.
就是在.
但義理書的第二章.
如果大家.
剛剛聽完但義理書的時候.
你應該還.
記得但義理書第二章.
講些甚麼.
大家眼仔碌碌.
現在馬上就反聖經了.
第二章講些甚麼呢.
幸好我有一部.
laptop.
或者我有部tablet.
然後按了給牧師.
也不知道我正在反聖經.
不要緊.
你讓我知道也不要緊.
但義理書第二章講甚麼.
但義理書講到.
但義理.
當他剛剛在.
這個.
新巴比倫帝國.
裡面.
被栽培.

$^{1001}$新一代的.
一些皇室官員.
雖然在被擄.
但是新巴比倫.
利用一種.
懷柔手段.
將一些皇室猶大國的皇室.
包括但義理.
帶入皇宮訓練.
好好招待他.
用這種方法去軟化他們對於.
復國或者為國家.
或者仇恨巴比倫.
這種心態能夠減低.
去做這番功夫.
在那個時候.
尼布甲.
尼薩王即是新巴比倫王.
也是老略.
但義理的.
這位君王.
他就利用.
一個機會希望去.
殲滅.
他國裡面一些.
渾水摸魚.
在那裡浪費.
公帑的公務員.
這些公務員做什麼.
做.
他有名的術士.
哲士.
解神話.
天象.
這些的.
總稱叫做哲士.
這些有學識的人.
懂得天文地理的人.
希望將這班人.
完全將他們剷除.

$^{1041}$於是乎.
尼布甲尼薩王.
有個.
發夢的經歷.
記不記得.
在這個發夢的經歷裡面.
尼布甲.
尼薩王.
他找這班的哲士來說.
我有個夢.
使得我很困擾的.
你要為我解.
這班哲士就說.
好啊.
那你告訴我內容講什麼.
有一個夢.
用今日講就真是玩野.
是不是.
又要我解夢.
你又不講個夢.
連個夢都講給我知.
我忘記了.
我們不要討論這件事.
總而言之這件事.
我現在忘記了.
奧秘.
這裡好像有這個字眼.
隱藏的奧秘.
這班哲士.
不知解釋.
三番四次推推躺躺.
記不記得.
尼布甲尼薩王說.
我早知你們在玩野.
在推推躺躺.
你再解釋不到.
講個夢內容給我知.
我就全部殺了.
記不記得.
接著怎樣.

$^{1081}$那班人就.
騰騰震.
接著在這個時候.
怎樣.
聖經沒有新馬仔.
你不用擔心.
我騰騰震.
接著.
接著怎樣.
尼布甲尼薩王.
你們再解釋不到.
我就找我官員一個內侍神.
叫阿略.
麻煩幫我將這班人.
殺了.
當阿略走出去.
皇宮的時候就遇見.
但以理.
但以理於是見到阿略.
方方失失.
接著阿略見到他.
但以理在這裡.
傾傾下的時候就說.
阿王現在這件事.
很大件事.
又要殺了這班人.
要我做劊子手.
你知道殺人不是好事.
是犯罪.
雖然很多人覺得.
殺人不是犯罪.
至少在上帝.
裡面殺人.
不是一件好事.
怎樣搞好.
但以理不用怕.
等我回去祈禱.
但以理接著回去祈禱.
上帝接著說.
什麼事.

$^{1121}$你不用擔心.
我將這個夢的奧秘.
隱藏了.
告訴你.
這個背景.
最後導致.
阿略.
不需要殺.
不需要殺任何.
術士節士.
又能夠透過但以理.
在這個過程當中.
將夢的內容.
以及夢的意思.
告訴尼泡加利薩王.
逐步帶尼泡加利薩王.
進入去.
相信這位創始成中的上帝.
那時沒有耶穌.
他信耶穌.
信了這位創始成中的上帝.
但在進行的過程當中.
你看不看到上帝.
當找一個人.
能夠將.
一件奧秘.
隱藏的奧秘.
將它告訴他.
解釋清楚之後.
做出什麼結果.
一連串的請求.
尼泡加利薩王.
生命的改變的開始.
一班節士.
他們得到.
不需要被殺的押運.
不要忘記阿略.
他不需要去參與.
一些污穢的工作.
一些污糟殺人的行動.

$^{1161}$但爾尼.
和他的朋友.
在這件事上.
亦都讓.
尼泡加利薩王.
看見他與別不同.
他能夠解釋意象.
勝過其他.
術士十倍有多.
我以前講過.
十倍即是完美無瑕.
弟兄姊妹.
保羅在這裡.
再一次告訴.
哥羅西教會.
知道我.
來到的工作.
是上帝交了給我.
將世世代代.
隱藏了的奧秘.
什麼奧秘?請求的奧秘.
講了給你們知道.
所以為什麼在第27節講到.
這個奧秘是可以叫.
連外邦人.
包括.
哥羅西教會裡面每一個.
都能夠可以享受這個.
豐盛的榮耀.
一會兒我們會再.
將來會講這個榮耀是怎樣解釋.
不過.
從這個.
介紹裡面.
就讓我們知道.
作為.
傳道的人.
聽清楚.
作為傳道的人.
不只是牧師.

$^{1201}$你也是.
你是要將上帝的.
奧秘.
請求的訊息.
來對清楚.
介紹.
展釋.
給上帝.
安排去到你面前的人.
讓他們清楚.
知道這個.
舊因的意義.
所以.
弟兄姊妹當我們去仔細看.
這個所謂奧秘的意思.
就不是一些.
神蹟奇事.
玩兩手魔術.
不是講這件事.
其實奧秘.
想講一樣.
我們很熟悉的.
一個字眼.
就是上帝和我們.
傾蜜計.
怎樣傾蜜計.
你有沒有和人傾蜜計的經驗.
傾蜜計就是你兩個人.
靜靜地在這裡.
交談.
你知道.
我知道.
聽眾都知道.
上帝將這個.
蜜計跟你傾完之後.
然後將主旨.
傳遞給人知道.
簡單來說.
這是上帝對你一種.
無限的.

$^{1241}$推心置腹的.
相信.
你不會和一個不信你的人.
傾蜜計.
有時有些大兄姐妹.
私人下.
有些事想和牧師.
師母傾私底下傾.
當然我不會將你的事.
講出來.
因為上帝沒有叫我將你的事.
講給人知道.
但是他信得過你的.
亦都你不會.
講出去的.
他細心和你講.
所以這裡保羅所講.
上帝所談論.
到歷世歷代.
所隱藏的奧秘.
現在.
這個部分.
上帝和我傾完之後.
我告訴你們.
就是連愛邦人都能夠.
勒囊括在上帝.
豐盛的恩典.
之下.
在這裡我先做個註腳.
為何要讓.
連愛邦人都可以在上帝.
豐盛的恩典.
之下.
這裡中文還說.
何等豐盛.
英文的聖經這裡說.
這是一個.
the hope of glory.
是讓你們能夠.
在這個盼望裡面.

$^{1281}$可以得到的.
得到甚麼.
27節裡面講到.
glorious riches.
of the mystery.
在這個裡面.
你們可以得到最豐盛的東西.
人生最得到最豐盛的東西.
就是能夠再次得到.
上帝的創作.
成為其中一個部分.
救恩.
其實就是讓人.
回到上帝裡面.
能夠和上帝在一起.
伊甸園其實都是一樣.
伊甸園不是上帝.
弄件玩具出來.
放兩個人去玩.
伊甸園其實是一個.
人和神的.
交通的.
私人場所.
你們有沒有這樣想過.
所以為甚麼上帝知道.
當下他犯了罪之後.
有沒有留意.
上帝去伊甸園裡面.
找他們聊天.
因為私人交通.
信得過他.
但結果就是上帝信他們.
這班人信不過上帝.
現在上帝就要將.
這個奧秘.
講為宣傳.
公諸同好.
所以鄧英姐妹.
當我們從這段聖經裡面.
我們去看到上帝.

$^{1321}$藉著保羅的一連串.
的陳述.
他這份的使命.
其實是要將上帝.
多年來.
請求的奧秘.
將他擴展傳開去.
而他現在在逐步.
做出來給大家知道.
怎樣信呢.
怎樣能夠在.
這個請求裡面有份呢.
所以接著下去.
第二章裡面所講的東西.
第三章裡面所講的東西.
第四章裡面所講的東西.
一步一步帶他們.
回到上帝這個原本.
做人無罪的.
這個樣式.
當中.
所以歌羅西書有些人.
經常講第三章第四章裡面講很多.
什麼基督徒的生活.
道德的生活.
其實不是道基督教化.
聖經從來不是講道基督教化.
而聖經只不過是講一個.
新做的人的生命的.
樣式是怎樣.
明白嗎.
從這個角度來看.
歌羅西書保羅的.
傳道是這樣建立.
一間教會.
傳承這個救恩下去.
他為了.
傳承這個救恩下去.
於是乎他就要.
有些東西真是要.

$^{1361}$面對要有取捨.
而這個就在我們最後第28節.
和29節裡面所講的內容.
當如何落實.
這個實踐.
這個真真實實的.
傳遞救恩的.
事奉是怎樣能夠.
落實去能夠完成呢.
請我們大家一起.
打開第28至29節.
我為大家讀出好嗎.
第28至29節.
我們一起讀好不好.
我們傳揚他.
是用諸般的智慧.
圈戒各人.
教導各人.
要把各人在基督裡.
完完全全地.
引到神面前.
我也為此.
勞苦.
照著他在我裡面.
運用的大能盡心竭力.
這裡你所講的東西.
你傳道都會明白的.
所以我要傳揚基督.
要有智慧.
這些好像.
屬靈口號的東西.
講就很容易.
我們看看真是怎樣落實.
第28節.
保羅現在去到.
講到他要參與這份事奉.
是要用.
諸般的智慧.
甚麼是諸般的智慧呢.
諸般的智慧就是諸般的智慧.

$^{1401}$如果我將.
諸般的智慧用一個.
你能夠容易明白的陳述.
就是.
你就要想盡方法.
留心著.
並且去仔細.
觀察著之後.
分析情況.
那你就會開始明白多一點.
諸般的智慧.
用盡你.
所有的心力.
心神.
知識.
來分析了解情況.
讓我用臭仔做一個例子.
帶小朋友.
你真的要用諸般的智慧.
由一個BB仔.
他哭一聲.
媽媽就已經知道他想要甚麼.
你問問那些.
帶小朋友的.
媽媽.
我也有這樣的經驗.
帶小朋友一段時間之後.
經常覺得有點幻聽.
經常聽到BB在哭.
是不是我女兒在哭.
是不是我兒子在哭.
其實他沒有哭.
他還上學了.
那種.
你長期很.
敏銳地去.
體會他有甚麼需要.
他這樣的哭法你就知道他肚餓.
他這樣的哭法.
他又要換尿片.

$^{1441}$他這樣的哭法.
可能他不舒服.
弟兄姊妹.
諸般的智慧.
我講的不是.
求神給我一點智慧.
come on give me some power.
不是這樣的.
是用盡上帝給你各種的能力.
你是準確瞄準.
看到分析.
理解到狀況.
牧羊教會.
我不怕老實說.
我作為牧者.
我經常是這樣的.
做甚麼.
看你的狀況.
牧羊不是陪你去喝茶.
當然我一會可以去喝茶.
是不是一會去喝茶.
是不是慶祝你生日.
我?.
我們有很多個生日.
不是去喝茶.
不是去遊輪.
可以.
我再強調.
牧師不是不可以和你玩.
不是不可以和你去喝茶.
那些不是唯一的事情.
你要去做.
真正的牧羊.
就好像爸爸媽媽.
經常留意著子女.
就好像師母.
經常留意著兒子.
其實我都沒有留意.
他好像.
當然.

$^{1481}$穿著背心在家.
涉涉就來.
是不是.
要去和他討論.
穿衣服的好處.
雖然他仍然不穿.
沒有辦法.
他勉為其難穿了.
你走了他又脫了.
但你做了.
你的角色.
你要去留心.
留心不止.
你要.
勸戒個人.
不是我講是保羅講.
教導個人.
這兩樣是什麼呢.
勸戒呢.
不用我說了.
警戒勸戒有個特色.
你要留意著他.
當你留意到他的時候.
你就要講.
我媽媽和我以前的上司.
在香港工作的時候.
我生第一個小朋友的時候.
我很記得我的上司.
很好的上司.
Ellen我免費.
多放你幾天假.
你知道我們只有.
九天假期在香港大學工作.
那個年代.
免費放幾天假.
我有產假.
你放假吧.
不過我要.
提醒你一句.
小朋友.

$^{1521}$You are the master, they are not the master.
他是英國人.
他說.
你是做主的人.
跟著我說什麼意思.
When the time.
he or she needs to be taught.
you need to.
tell them right away.
我媽媽都講這個.
教小朋友要在.
那一刻他飯菜的時候.
才有用.
過了教就沒有用.
第一他不記得.
第二他不記得你教他的時候.
他就覺得你冤枉他.
他錯那一刻你就要教.
用豬般智慧.
看著他需要.
然後一見到他不妥就勸戒他.
這些弟兄姊妹.
這份工作.
不是好做的.
成日看著他.
你知不知道我以前在教會.
我打我做.
做好之後要印.
印程序表聽起來是很簡單的東西.
印程序表有多難.
多難.
按個按鈕就可以.
30張40張就出來了.
用什麼.
按個按鈕就很容易.
按一個按鈕之後.
那個結果可能很大件事.
最大件事就是.
如果你一按個按鈕你不理.
不觀察著他.

$^{1561}$你就會印40張.
錯的或者50張錯的程序表.
你印的時候.
你印多幾張就要拿出來.
看看機有沒有問題.
上面的東西是否正確.
一見有錯就馬上按按鈕.
停機.
程序表是很簡單的.
你知不知道.
大家有沒有去過.
印銀紙的.
發鈔的工廠.
梁牧師.
哪裡有這樣的東西.
有的.
你去美國.
菲律賓有間.
印美金的工廠.
他會告訴你怎樣印.
雖然他不讓你現場看著.
看看機旁邊怎樣印.
遠遠播放影片給你看.
原來印鈔的那些人.
印十多二十張就拿一張出來.
找個放大鏡逐張看.
你知不知道為什麼.
因為他想知道.
那些銀紙有沒有印錯.
有沒有印歪.
一有錯就按按鈕.
停機.
不想錯繼續下去.
做牧者.
其實做這個工作.
有誰喜歡被人.
指出他的錯誤.
做牧者.
做牧者其實有個.
有個叫做.

$^{1601}$我們叫做什麼.
工作的工商.
什麼工商.
習慣了整天說別人.
有些問題.
所以要整天調教自己.
要學會欣賞別人的.
好處.
不然就很容易變成一個.
整天指責別人的人.
第二.
教導國人.
教導是什麼意思呢.
我認識的.
牧師教.
教導國人就是你逐步.
去給指示他.
怎樣一步一步去做.
作為一個.
牧者他的工作.
就是要.
將信徒.
不懂的東西.
有系統地去教導.
所以就有主有學.
有個別的.
陪談.
見面.
牧師不是八卦.
不是想知道你的東西.
但是想跟你一起.
走這條路.
做信徒弟兄姊妹.
之間也是一樣.
你是想做弟兄姊妹.
一起走這條天路的時候.
走得對.
他不懂的你教他走.
但如果信徒不肯學.
不肯受教.

$^{1641}$其實不是牧者的問題.
而是.
信徒的信的問題.
除此之外.
保羅在這裡再講多一樣東西.
要把國人在.
基督裡完完全全地.
引到上帝的面前.
這句說話.
是那個教導.
用智慧去警戒.
之後的.
最終的目標.
牧者的.
最終的目標.
不是信徒跟著好聽.
我話成為我的擁躉.
也不是要求.
那個信徒或者會友.
成為了教會的.
其中一個事奉者之後.
就使到他盡.
不是.
而是最重要的就是使到他們.
能夠在基督的.
面前.
能夠.
過得上帝.
的審核.
所以這裡再講要把國人在基督裡.
完完全全地.
引到上帝的面前.
就是在上帝面前.
都站得住腳.
所以.
作為.
弟兄姊妹的.
教會領袖的.
我們的目標都是為了.
這件事使到你.

$^{1681}$走到去.
上帝面前.
上帝說.
OK了.
是吧.
所以有些信徒我現在發覺.
特別是牧祖會二十幾三十年之後.
你講什麼.
他都跟你拼.
你就不要再跟他講那麼多.
你講完上帝的說話.
就算了.
他想跟上帝交通.
交通不交通都沒有辦法.
因為他根本就覺得.
他是神.
他要求就是他自己能夠.
站得住他根本沒有想過.
他是不明白你是.
想他在上帝面前站得住.
當做到.
這個最終目的.
之後.
保羅再次重申.
我也.
為此勞苦照著.
他就是上帝基督.
在我裡面運用的.
大能盡心竭力.
保羅在這裡講到.
我做那麼多.
事情.
我用了.
上帝給我的能力.
不是去透過事奉.
來滿足.
我在你面前多威風.
多好處.
我心想牧師你懂了.
我跟你無話可說.

$^{1721}$不是 我不是完全懂.
但如果你的心想.
覺得我用那麼多時間那麼多力量.
去解說之後.
原來你以為我是.
為了我去誇耀自己.
其實保羅在.
哥林多前書裡面.
有介紹過.
保羅在哥林多前書.
裡面講到.
他從來不會誇口自己有多厲害.
他只誇.
自己的軟弱.
什麼叫誇.
就是要講.
他都給人看到.
他自己的不善之處.
而不是講自己有多厲害.
所以事奉主的人.
或者甚至基督徒.
他一生用盡.
他自己的力量.
其實不是去高他自己多虔誠.
問我多厲害.
我不怕.
老實講 這句不是誇口.
我會預備一堂.
現在大約是用.
通常三天左右.
我其實.
不需要預備 我坐出來.
看完一段經文可以講給你聽 講兩三個小時.
講的東西有質有數.
我不需要做很多功課.
為什麼要用.
三天時間去預備一堂.
弟兄姊妹.
因為在這三天.
裡面.

$^{1761}$我和這段聖經.
有充足的時間.
去交往對談.
被他教我之後.
然後我去.
體察你的需要.
去按你的需要.
去給你傳道.
但我不是.
去到張兄.
我多厲害.
我不需要.
我如果要拋書包.
我想.
我要拋的東西很多.
我不需要在這裡拋.
但是保羅在這裡講得很清楚.
我為了這一切所做的東西.
都是為了你們.
能夠可以在上帝裡面.
站得住.
不是為我多厲害.
弟兄姊妹我已經講了這麼多.
大家看到.
保羅的這個牧羊工作.
是一個長期的.
作戰.
是一件很長時間.
的一個.
對信徒的一個牧羊.
不是談花樣搞幾堂崇拜.
搞幾個福音.
搞一些暴徒會.
請一些名人來講道.
或者搞一些古靈精怪的.
事奉節目.
然後就教會.
哎呀多興旺.
沒什麼意思的.
如果你始終都不能夠在上帝的面前.

$^{1801}$能夠真的站得住腳.
我搞這麼多東西都是沒用的.
可能你覺得你自己.
好像很虔誠.
但這些虛假的虔誠.
卻不能夠帶你去上帝面前.
得到上帝的肯定.
所以從這段聖經開始.
第24-29節我們可以看到.
保羅其實從開頭.
到尾他的事奉.
都是充滿著對.
哥羅西教會的.
擔心.
就好像一個父母一樣.
從初生的嬰兒的時代.
他已經願意放棄.
他自己的睡眠的時間.
來為他的嬰孩.
來餵奶.
換尿布.
甚至怎樣.
跟他在.
成長的過程當中.
有很多的.
願意放棄自己原本的喜好.
為的是.
兒子兒女能夠得到成長.
我.
到最後.
做一個.
自身的一個故事.
作為一個總結.
我這樣說.
崇拜學講道學是錯的.
因為不應該告訴你原的.
講完這句你可能沒有心聽.
不過我講一件事.
可以幫到大家.
對保羅的事奉.

$^{1841}$多一重理解.
大家知道.
我很喜歡.
吃.
其實我很喜歡吃.
一些酌料.
什麼酌料包括.
燒鴨.
燒鵝.
我很喜歡吃燒鵝.
我很喜歡吃雞.
所以師母都知道.
所以我們經常在家.
買雞.
買雞做什麼.
蒸雞.
你知道.
現在的雞很貴.
超級市場.
一隻黃油走地雞.
16.99元.
生蚵蚵一隻.
16.99元.
都是小隻.
蒸了一隻雞.
很多時候我蒸雞.
蒸完之後.
我蒸雞一定要.
薑蔥蓉.
我很喜歡吃薑蔥蓉.
然後是斬雞.
你懂不懂斬雞.
很多人不懂斬雞.
你給他斬完那隻雞.
遍體鱗傷那隻雞.
不過不要緊.
我斬好了.
砌成一隻雞.
好像酒樓一樣.
準備一會兒吃.

$^{1881}$但是家裡有幾個人.
有些要蛋白質.
特別很多需要蛋白質的.
雞最多蛋白質.
最多蛋白質.
除了雞蛋.
我就.
輕輕想著.
吃兩塊雞肋骨.
等他們吃完.
吃完我才.
吃.
誰不知.
回家想著他們吃完.
一看.
沒有了大半隻.
去了哪裡.
我就問師母.
為什麼沒有了大半隻.
你吃了沒有.
我問她吃了沒有.
不是問她有沒有吃.
你猜誰吃.
這間屋吃這麼多.
我不說誰.
總之不是師母.
不是我.
當然有些人.
這樣吃不像樣.
這隻雞已經.
沒什麼肉.
黃油雞.
我說.
沒有了大半隻.
我斬雞.
斬得很精細.
將雞胸.
起肉.
沒有骨.
切好一件件.

$^{1921}$雞胸肉去了哪裡.
沒有了.
還有.
剩下什麼.
還有剩下.
一隻雞腿的後端.
近腳的位置.
一隻雞腿我都破開.
幾件.
還有一堆.
用皮包著骨的雞胸骨.
包括雞頸.
還在碟上.
我當然知道誰吃了.
討論完之後.
我就馬上跟她說.
我說.
你知不知道家裡有四個人.
我不介意.
你吃的.
做出來都是給你們吃的.
我也不介意.
給你先吃好一點的.
雖然我很想吃雞胸肉.
雖然我很想吃雞腿.
但我都留著.
等你們吃完我才去吃.
但你要懂得去考慮.
一家人.
平均分.
吃多少.
弟兄姊妹.
這些教導.
你猜會有什麼反應.
當然會說.
我不能吃的嗎.
你不是不能吃.
是吃多少.
是怎樣照顧到你的家人.
使一家人都去享受.

$^{1961}$這一隻.
用時間.
蒸得好吃的雞.
當然這堂講道.
不是講蒸雞不是講切雞.
不過我想告訴大家.
作為.
想家人好的.
永遠希望.
將最好的留給家人.
一個淑齡的家庭.
保羅.
將最好的留給.
西教會.
他不止留給西教會.
他留給以弗所教會.
他留給迦太教會.
他留給所有能夠.
接觸到的教會.
他受苦吃少一些.
有什麼何妨.
保羅的心智.
事奉就是這樣.
他講這麼多不是高於他自己.
是你.
有沒有能夠成全這份心.
我們一齊同心低頭禱告.
父神多謝你.
讓我們有一位這麼好的.
牧者保羅.
能夠將你在.
奧秘裡面.
隱藏了多年的奧秘.
將它清楚陳述出來.
讓我們作為信徒.
可以在當中.
明白主你的心思意念.
怎樣為我們.
付出一些本來.
你根本不需要.

$^{2001}$付出的事.
但你為了父神的喜悅.
他願意付出.
即或對自己.
未必是這麼喜歡的事.
但因著.
上帝你喜悅的事得到成就.
主耶穌基督.
滿心歡喜.
保羅同樣為了.
基督的身體為了教會.
亦都付上自己.
一切上帝你給他的.
智慧能力.
為的是讓教會能夠在.
上帝裡面站立得穩.
站得住腳.
所以他願意.
即使付出.
難堪的經歷.
與他的生命結連.
但他仍然願意.
因為這是主你所喜悅的事.
而這一份的受苦.
亦都為保羅帶來.
教會成長的喜樂.
弟兄姊妹.
給我們一些經驗.
在歷史裡面.
重覆保羅的心思意念.
今天我們作為信徒.
有沒有在上主的裡面.
同樣有所領受呢?.
求上主你幫助我們.
有這份同一樣的心思意念.
事奉著你.
祈禱奉靠主耶穌基督得勝名治異球.
阿們.
\newpage



\section{創世記 26:35}
\label{sec:031763}
\textbf{聖經人聖經事:心裡愁煩}~~~~黃克勤牧師博士
\newline
\newline
連結: \href{http://www.church.com.hk/acms/content.asp?site=cdc\&op=show\&type=product\&code=031763\&layout=sermon}{\texttt{031763}} ~~~~ 語音日期: 2026-02-25
\newline
\newline
\hyperref[sec:031762]{\small{< < < PREV SERMON < < <}}
~
\hyperref[sec:index_chronic]{\small{[返順時目]}}
~
\hyperref[sec:031764]{\small{> > > NEXT SERMON > > >}}
\newline
\newline
創世記 26:35
\newline
\begin{longtable}{cl}
\hline
\hline
章節 & 經文 (和合本修訂版)\\
\hline
26:35 & \begin{tabularx}{0.7\textwidth}{X} 她們使以撒和利百加心裡愁煩。 \end{tabularx} \\ \\
[1ex]
\hline
\hline
\end{longtable}
$^{1}$各位網友大家好,我們又一齊讀聖經了.
《生命述書》的網絡和我們一起看聖經人,聖經事.
希望你在人與故事的背後認識上帝更多.
今天講四個字.
其實通常都是講兩個字.
就是「心煩」.
「心裡愁煩」即是心煩.
今天看看哪個心煩.
是《創世記》二十六章三十四,三十五字,兩節,很短.
「以數四十歲的時候娶了黑人比利的女兒尤迪」.
「和黑人以倫的女兒巴神梅為妻」.
「這兩個女士常常使老爺奶奶以薩和利伯家心裡愁煩」.
先介紹一下她們.
「以數四十歲」.
這麼大年紀了,娶妻不可以不大.
當然有些可能十多二十歲去娶.
不過她四十歲.
四十歲其實是不大的.
因為當日阿寶罕百多歲.
以薩也百多歲.
當時的人很多也百多歲.
如果四十歲只是人生走了三分之一.
還有很長時間.
「以數四十歲」.
如果我猜當時來說.
就算遲婚了.
第一點.
「娶了」這個字很有趣.
希伯來文用的字是Lakak.
Lakak.
Lakak是什麼來的呢?.
就是.
真的「娶了」,大家留意一下.
中文是「取了」.
這個字是「娶了」.
其實這個「娶了」是「拎了」的意思.
「拿在手裡拎了」.
有點像中國南方人.
「你兒子多大?」.
「三十多歲了」.

$^{41}$「三十多歲還不叫他拎回人家」.
這是老一點的廣東話.
「拎回人家」.
「老友,大家差不多一年多了」.
「你太太過世了」.
「不如拎回人家」.
拎回人家,很有趣.
原來也有這個字.
Lakak.
以素「拿了」.
意思是.
不是經過三思六禮,名門正取.
「拎了」.
「不如大家一起歡喜」.
「一起拎吧」.
「拎了」.
好了.
他是什麼人呢?.
不是他們本族.
你知道.
以薩娶妻是亞伯拉罕.
特意叫他老婆人回故鄉那裡娶回來.
就是不想本地人.
誰知以素大情大性.
恐慕有力到處玩.
認識了兩個本地黑人.
本地有很多族.
迦南族很多族.
拜不同的神.
就娶了兩個黑人的女孩.
不是黑色.
黑字.
Ketit.
黑人其實是當地一個民族.
Ketit 意思是恐懼,懼怕.
民族名叫懼怕.
有些人認為.
聖經裡有黑人出現過很多次.
但其實在歷史裡找不到黑人.
是不是作家,創世記.

$^{81}$後來有些考古學者發現.
真的有黑人的記載.
而且黑人當時是一個部族.
後來一直擴大.
甚至成為一個很強大的國家.
後來又被滅了.
找回他們的泥板,文件等資料.
真是有黑人這一族.
黑人裡面有一個人叫比利.
比利就是Biree.
Biree就是「我的井」.
這個男士的爸爸叫「我的井」.
他的女兒叫「尤迪」.
「尤迪」是Yahudith.
「Yahud」是猶太.
「迪」是女孩.
「Yahudith」就是猶太女人.
「尤迪」的意思是猶太女人.
但是它變成一個英文的名字反而很多人都認識.
甚至我認識一個人是用Judith這個名字.
Judith這個英文名字.
Judith這個英文名字就是猶太女人的意思.
那麼是不是猶太人呢?.
是與不是.
其實「猶太」是地方名字.
那裡地方的一個女士.
就像香港人.
不過我的名字也叫香港人.
奇怪一點.
當時沒有什麼名字.
這個女子是猶太的女人.
還做了名字.
你要知道猶太不是猶大.
猶太是地方來的.
阿巴喊.
以撒.
然後雅各.
雅各有十二個兒子.
十二個兒子其中一個就是猶大.
這個名字.

$^{121}$猶大一直數下去就是大衛.
大衛再數下去就是耶穌.
在主耶穌的十二門徒裡面.
又有一個人叫猶大.
主耶穌自己的弟弟.
四上的弟弟又叫猶大.
寫猶大書的那個.
一個就是賣他的猶大.
死了.
有一個弟弟叫猶大書.
寫猶大書的猶大.
這裡不是猶大是名字.
猶大是地方來的.
還有另一個黑人叫以倫.
Elon.
Elon是指重哲尤.
強大的意思.
爸爸的名字叫重哲尤.
很有趣.
這個女子叫巴神末.
Bosmith.
就是香料的意思.
耶穌四十歲的時候.
就拿了兩個女人做太太.
拿了兩個女人.
用現在的詞應該這麼說.
拿了兩個女人做太太.
就這兩個.
拿了就拿了.
拿了.
就是這句了.
他們嘗試以撒和利伯加心裡愁煩.
是什麼意思?.
這兩個.
媳婦令到老爺奶奶心煩.
這樣也有話說的嗎?.
有的.
第一點.
就是他.
以數決定.

$^{161}$拿了兩個.
不是自己種族的太太.
這個在猶太.
在當日他的家族裡.
是一件不受歡迎的事.
因為你知道以撒承繼著.
阿巴拉罕的父親.
是要找一個本族的女子.
都還沒找到.
他自己就去拉了兩個女子回來.
入門也要在這裡住.
已經很煩了.
那怎麼辦?.
而且還令他們心裡愁煩.
他們兩個女子令他們心裡愁煩.
為什麼會這樣呢?.
主要原因就是.
很簡單.
有文化的差異.
黑人本地的民族有他們拜的神.
我們拜耶和華.
我們敬拜耶和華.
他們拜不同的神.
通常是有不同的世界觀.
有不同的價值觀.
如果你想認識一個人深一點.
你跟他聊天.
其中一定要聊到的就是.
這個世界從哪裡來的呢?.
世界是自然爆出來的.
Big Bang Theory.
一下子出現.
慢慢慢慢.
然後在海水裡有很多.
不同的百多個化學物質.
撞啊撞啊撞啊撞啊.
然後竟然間偶然.
只是發生一次.
撞出一個氨基酸.
然後撞啊撞啊撞出十個氨基酸.

$^{201}$然後十個氨基酸就沒了.
於是就組成了一個細胞.
這個單細胞加上細胞膜.
然後又有DNA等等.
變成一個單細胞的生物.
單細胞生物一路分裂.
變成多細胞的生物.
其他的動物就出現了.
一路由低等到現在的人類.
嘩!這樣也相信的.
那麼這個大爆炸論是怎麼來的呢?.
偶然發生.
好了.
這個世界觀.
就影響到他不同的價值觀了.
既然世界偶然出現.
那麼我們人其實也是偶然出現.
在這個世界的產物.
既然偶然出現.
講什麼人生意義呢?.
隨便偶然過了就算了.
我們基督徒又不是這樣看的.
所以不是有創造主.
創造主的創造是必定有目的的.
不偶然的.
嘩!如果這兩個人結婚.
很艱難的.
去到最後的時候會有很多爭執.
兩個人都喜歡一點.
生了個兒子.
哎呀!複雜了.
用什麼世界觀和價值觀來教他呢?.
這樣就很複雜了.
大家搶了.
星期天當然是去教會.
不是!星期天這麼難.
當然是和父母喝茶.
然後去玩.
不是!去教會.
這次兩個價值觀不同了.

$^{241}$那些誠實,公義,正直等等.
價值觀不一樣.
那麼生活上也很艱難.
就會心裡受煩.
又或者這個黑人女孩.
為什麼會和伊蘇這麼大情大性地在一起.
就是大家都這麼豪放.
大家都這麼豪放.
於是可能在家裡抱著一塊.
嘩!抱抱抱抱.
一會兒這個黑人女孩.
一見到伊莎.
這不是聖經說的.
一見到伊莎.
哎呀!爹地.
抱著又親一下.
嘩!不是吧.
於是.
哎呀!弄得很不開心.
好了,當這對父母.
伊莎,李伯嘉心裡受煩.
自然臉會怎樣.
長了.
還有什麼顏色.
黑色.
黑臉神地拉長臉.
當你這樣的時候.
那兩個女孩會不會也是這樣.
自然又會黑臉拉長臉.
哎呀!不要見那麼多了.
你爸爸和媽媽很麻煩.
如果萬一李伯嘉問一個很差的問題.
伊蘇!伊蘇!想問問你.
如果有一天你要選擇.
你會愛你老婆還是愛你媽媽.
哎呀!千萬不要問這個問題.
如果我掉下海.
你救我還是救你老婆.
千萬不要問這個問題.
這樣很麻煩的.

$^{281}$很麻煩的.
因為問這個問題.
會令對方多數想.
救不救你.
救他吧.
如果老婆問.
伊蘇!如果掉下海.
你救我還是救你媽媽.
你只會逼著她先回答我媽媽.
那你就是不愛我了.
不是嗎?不是嗎?.
於是.
嘩!複雜.
所以其實有很多不同的價值觀.
世界觀,文化衝突.
各樣面色.
一家人.
相見好自然.
下一句是什麼?.
同住就很難了.
於是.
大家心裡愁煩.
拉長臉.
黑了.
於是相處得很差.
不同文化兼容不了.
所以如果是一家人.
不要說家人.
不要說不同族.
相同香港人在一間教會裡.
都會不順眼.
為什麼?不同家庭.
不同背景.
同事.
同學.
弟兄姊妹.
都見到心煩了.
很多問題.
所以要寬容一點.
喜樂一點.

$^{321}$尊重一點.
多一點空間.
這樣就好了.
就是這樣的遭遇.
其實.
耳掃也很難做.
是因為他.
娶了這兩個回來.
令到.
爸爸媽媽.
很難.
於是整家人都心裡愁煩.
雅各有沒有話說?.
聰明的雅各當然不說話.
走遠一點.
是嗎?.
少衝突一點.
就心裡愁煩.
大家要小心.
如果遇到不同的文化.
次文化.
不同的家庭背景等等.
都要有多一點空間.
彼此接納與尊重.
否則就拉長了面.
心裡愁煩.
祈禱.
求主祝福我們.
讓我們看到.
在《創世記》記錄了很多.
上帝揀選的家庭裡.
很多恩怨情仇會發生.
而這些只是在家庭裡發生才會心煩.
如果在別人的家庭發生.
我們沒什麼心煩.
在自己家庭裡發生.
很心煩.
求主幫助我們.
更加喜樂.
更加寬容.

$^{361}$更加耐心又平安.
倚靠你.
就一切交託給主.
就自己發發樂樂地相處與生活.
更多的空間與尊重.
誰聽祈禱.
奉耶穌名求.
阿們.
\newpage



\section{馬太福音 9:14-17}
\label{sec:031764}
\textbf{變!定唔變?}~~~~王正傑牧師
\newline
\newline
連結: \href{http://www.church.com.hk/acms/content.asp?site=cdc\&op=show\&type=product\&code=031764\&layout=sermon}{\texttt{031764}} ~~~~ 語音日期: 2026-03-01
\newline
\newline
\hyperref[sec:031763]{\small{< < < PREV SERMON < < <}}
~
\hyperref[sec:index_chronic]{\small{[返順時目]}}
~
\hyperref[sec:031765]{\small{> > > NEXT SERMON > > >}}
\newline
\newline
馬太福音 9:14-17
\newline
\begin{longtable}{cl}
\hline
\hline
章節 & 經文 (和合本修訂版)\\
\hline
9:14 & \begin{tabularx}{0.7\textwidth}{X} 那時，約翰的門徒來見耶穌，說：「我們和法利賽人常常禁食，你的門徒卻不禁食，這是為甚麼呢？」 \end{tabularx} \\ \\ \relax
9:15 & \begin{tabularx}{0.7\textwidth}{X} 耶穌對他們說：「新郎和賓客在一起的時候，賓客怎麼能哀慟呢？但日子將到，新郎要被帶走，那時候他們就要禁食了。 \end{tabularx} \\ \\ \relax
9:16 & \begin{tabularx}{0.7\textwidth}{X} 沒有人把新布補在舊衣服上；因為所補上的會撕破那衣服，裂口就更大了。 \end{tabularx} \\ \\ \relax
9:17 & \begin{tabularx}{0.7\textwidth}{X} 也沒有人把新酒裝在舊皮袋裡，若是這樣，皮袋會脹破，酒就漏出來，皮袋也糟蹋了。相反地，把新酒裝在新皮袋裡，兩樣就都保全了。」 \end{tabularx} \\ \\
[1ex]
\hline
\hline
\end{longtable}
$^{1}$兄弟姐妹平安.
今天有點冷.
跟你旁邊那個打聲招呼.
歡迎一下他好不好 讓他溫暖一下.
今天我們要講一個.
變還是不變的訊息.
究竟變還是不變呢.
不過對於我們來說.
人性兄弟姐妹我們都知道.
其實我們都很抗拒.
隨便要去改變.
大多數人.
都不喜歡改變.
我們有.
習慣原有的環境.
有習慣一些熟悉的方法.
習慣一些.
向來的做法.
所以你要變.
大家都不會喜歡.
就好像你們今天回來的教會一樣.
你坐開的位置.
你都是坐在那個位置.
你回來如果你坐開的位置.
你見到有人坐了你就很不高興.
不會.
那就不至於要打他.
所以我們.
人是很習慣了.
一向的做法的時候.
我們要去變.
是從來都不容易的.
因為改變是代表著.
未知.
也代表著很沒有安全感.
因為你不知道.
改變了之後會怎樣.
所以就沒有安全感.
亦都因為未知.
因為沒有安全感.

$^{41}$甚至會很害怕.
好像我控制不了.
我掌握不了.
所以這些都是令到我們.
很難去隨便改變.
亦都很抗拒去改變.
所以就算我們知道.
需要改變.
但卻是想要穩定.
不願冒險.
最後都只會.
原地踏步.
我講個例子.
給你聽你就明白.
我明明知道.
我需要減肥.
但我卻是想要.
吃這個.
麻辣火鍋的話.
那我最後.
都只會怎樣.
減不了肥.
所以人就算我們知道.
需要改變.
但卻是不想改變.
我們只想穩定.
最後只會得到.
原地踏步的結果.
但對於我們來說.
究竟我們的信仰生活.
是否都是這樣呢.
因為我們很多時候.
可能我們都會很喜歡.
將我們的一些習慣.
等同了真理.
又或者我們將一些.
傳統的做法.
等同於一些屬靈的標準.
所以很多時候.
都會出現一句說話.

$^{81}$以前不是這樣做的.
以前不是這樣.
這個是.
常常會遇到.
而今天我們看的這段經文.
正正是講到一場關於.
改變的衝突.
施禁日後.
他的門徒來問耶穌.
我們和法理賽人.
都經常禁食.
但主耶穌的門徒就怎樣呢.
就倒不禁食.
他們不禁食的.
這個問題.
其實他背後是.
想表達一樣東西.
就是要問耶穌.
你為什麼不按我們一向的做法去做呢.
所以這場是.
關於變和不變的.
一場對話.
而耶穌的回答.
是不單止回應禁食的問題.
更加是.
揭示了神的國度要臨到.
新的時代已經開始.
而今天.
弟兄姊妹.
我們都應該問問自己.
如果當神在我們生命中.
帶來更新.
究竟我們會選擇.
拒絕.
還是迎接呢.
可不可以幫我問問你旁邊的人.
你是會拒絕.
還是迎接呢.
當我們看回經文.
當主耶穌.

$^{121}$在山上宣講完.
登山寶訓.
在第五章的時候.
當他下山.
施行了一連串的神蹟.
他去醫病.
其中他醫好一個癱子.
他更加宣告.
他的罪已經得到赦免.
然後他平靜風浪.
趕鬼入豬群.
更加呼召人跟隨他做門徒.
他呼召了瑞利瑪泰.
更加他和罪人一起吃飯.
當時耶穌做了這一連串的事.
其實是做了幾件.
對當時的人來說.
那些文士,法利賽人.
甚至可能包括斯金一項.
的門徒來說.
他們看到的就是這些行動.
都是很震撼性.
這些行動每一樣都衝擊著.
當時他們思想的框框.
所以沿路上就不斷有人.
去挑戰耶穌.
文士就質疑耶穌.
說了一些淺妄的說話.
因為他說.
他醫好那個癱子的人.
他跟他說.
但是他說的說話.
就是說你的罪得赦免.
因為對於猶太人來說.
從來只有神.
才可以有權柄.
去赦免人的罪.
而當時耶穌這樣說的時候.
就是說了一些淺妄的說話.
就是他其實在說.

$^{161}$他自己是什麼?是神.
所以當時文士會質疑耶穌.
他說淺妄的說話.
而法利賽人更加挑戰的是.
耶穌說為什麼會跟那些.
罪人和那些稅吏一起.
吃飯呢?.
而在這個時候.
當他們吃飯的時候.
有斯金一項的門徒和法利賽人.
就走來挑戰耶穌.
問他為什麼他的門徒.
不禁食?.
其實禁食這件事.
當時的律法其實只是要求.
在贖罪日禁食.
但是那些很敬虔的.
猶太人他們是會.
多加一些.
禁食的日子.
特別是法利賽人.
一個星期更加會有兩天禁食.
這個其實.
慢慢成為了他們一個贖人的指標.
你是敬虔的.
你一個星期裡面應該起碼有兩天.
禁食.
禁食是一種宗教操練.
是要和祈禱和.
佩改同時進行.
所以那些門徒居然不禁食.
是不對的.
所以他們這一刻要去挑戰耶穌.
為什麼呢?因為他們認為.
耶穌他作為一個老師.
他應該為他們的門徒.
去負責.
你的學生.
竟然不禁食.
竟然禁不禁田.

$^{201}$你作為老師是有責任的.
你是應該為他們沒有禁食.
去負責.
所以在這裡我們會看到.
民事,法利賽人和施浸日漢的門徒.
他們發出的挑戰.
其實都反映了他們只是.
遵守一種舊的標準.
那個舊的標準是什麼呢?.
他們以守律法,以人的傳統.
來作為一個禁田的標準.
是依靠行為.
來去稱義.
他們仍然停留在這個舊的標準.
裡面.
但是等一會.
當我們看到這裡的時候.
我們也會看到.
這些民事,法利賽人和.
施浸日漢的門徒.
他們都很律法主義.
他們都.
很以人的傳統來作一些禁田的標準.
他們仍然停留在.
靠行為來去稱義.
但是等一會.
其實我們在教會.
我們都很多時候會.
有一些我們自己的神學.
有一些我們自己.
心裡面的規則.
我在網上.
找了一些.
我中間是刻了.
有兩個字的.
那個是教會奇怪規則.
有一些規則是這樣的.
首先第一件事就是.
崇拜的時候.
男女是要分開坐的.

$^{241}$夫婦都要分開坐的.
有一些教會.
我看到我們就沒有分開坐.
也有一些教會的規則.
就是不准看電影.
為什麼呢?.
因為電影很多時候.
有暴力.
又有色情的成份.
我最近看了一部電影.
2D級的.
整部電影都是粗口的.
我很稀奇的.
就是以前看電影.
那部電影.
如果有一句粗口就變成3級.
但是現在竟然.
整部電影都是粗口.
就變成2B級.
所以是不是不應該看電影.
還有就是.
不應該參與婚宴.
為什麼呢?.
因為婚宴裡面會玩新娘.
玩新娘的時候.
通常都會.
有一點.
色情的成份.
所以.
也不適宜去參與.
這些婚宴.
很稀奇的就是.
耶穌也有參加婚宴.
可能那時候他們沒有玩新娘.
也不可以.
喝酒.
因為酒能使人放蕩.
所以不可以喝酒.
還有就是.
崇拜不可以拍掌.

$^{281}$給別人.
因為唯有神才值得.
給別人拍掌.
還有就是.
不准吃豬紅.
因為不可以.
吃血.
但是我想問大家.
有沒有人吃過豬紅?.
只有這麼少人.
敢舉手.
因為我平時.
是吃鴨血的.
不准吃豬紅.
也有些教會.
是不准染髮的.
因為染髮.
是不敬虔的表現.
我在一間台灣教會.
他們的分享.
我看到在牧師房門口.
貼了一張標語.
告示.
他們教會的師母.
訂立了一個規則.
規則是這樣的.
若為牧師好,請停止餵食.
下面還有兩句.
願主賜福你.
所以我說.
這些規則是不對的.
我們回教會.
很容易.
好像花旗菜人.
文士.
甚至師浸日行的門徒.
我們也有很多規則.
有很多傳統.
我們演變成.
一種敬虔的標準.

$^{321}$其實在經文裡.
問題不在於.
禁食的本身.
而在於當神.
去做新的事的時候.
我們人是否仍然.
要去停留在一些舊的標準裡.
這段經文的重點.
不是要去否定禁食.
不是說禁食不對.
而是要指出耶穌的彌臨.
是標誌著一個.
暫生的舊恩時代.
人不再需要去靠.
守律法.
靠行為去稱義.
而是要進入一個.
新的舊恩時代.
所以耶穌被質疑.
被挑戰不是因為祂錯.
而是因為祂所做的.
超越了當時這班文士.
法利賽人.
這班施浸日行的門徒.
他們所理解.
特別我們認識法利賽人.
他們很喜歡在律法上.
加上很多傳統等等.
他們加上這些傳統.
作為對信徒奉行生活的一種標準.
這樣做就會更健全.
他們更加用律法.
作為審判.
人的一些標準.
但最大的問題是.
他們將自己都不能遵守的規則.
要求別人遵守.
我們其實.
都很清楚一件事.
就是靠律法.

$^{361}$講規矩.
來叫人稱義.
這件事是.
達不到的.
如果不是.
我們都不需要耶穌基督的舊恩.
剛好今天下午有家長小組.
家長小組講的內容.
有點像.
我們很多時候幫小朋友.
訂立很多規矩.
但訂立完這些規矩.
我們往往見到小朋友都守不到這些規矩.
這是我們的人性.
是我們的罪性.
所以我們明白.
耶穌來到.
是進入一個嶄新的.
一個舊恩的時代.
人不再需要靠.
守律法和靠行為來稱義.
而是靠恩訊耶穌基督來稱義.
所以今天.
我們都要問的是.
我們是不是都是這樣呢.
當我們見到一些新的東西.
改變的時候.
我們會有什麼反應呢.
如果當你今天見到神.
祂使用一個.
可能你看不過眼.
你不習慣的人.
又或者當教會.
用一種我們大家都很不熟悉的.
方式來去傳福音.
又或者.
當年興一代.
有一些很不同的敬拜形式.
甚至去唱一些新的.
詩歌的時候.

$^{401}$大家會有什麼反應呢.
我們的反應是不是馬上說.
這樣也可以的.
還是我們真的會說.
會去想的是.
神是不是正要去做一些.
新的事呢.
因為我們屬靈生命的危機.
不是我們只是在說.
我們犯不犯錯那麼簡單.
而是我們很多時候用很多.
我們一些個人的想法.
一些舊的框架.
來去拒絕了神一些新的作為.
所以我們都要思想的是.
我們怎樣迎接神.
新的作為呢.
特別當.
我們經常看著外面的這條.
榮光街.
社區改變.
我們以往的做法.
以往要去祝福.
這條榮光街的居民的時候.
可能我們要祝福.
是一群基層人士.
但是當.
社區重建之後.
可能住在這裡的大部分.
都是一些中產的家庭.
那我們是不是.
還會用.
過往去做.
基層人士的工作的方式.
來去對待.
來去祝福這一群.
中產的人士呢.
又或者當我們要去.
看的是風險工作的.
改變.

$^{441}$還記得呢.
上一代說以前教會.
就會派很多東西.
其中一樣是派什麼.
派奶粉.
如果我們現在教會.
派奶粉.
有沒有人要.
可能有人要的.
我們樓下吃團餐的人.
都可能會拿吧.
所以福音工作.
不同了.
如果我們今天純粹.
只是派單張.
那些人不願意收.
所以我們要想.
究竟我們的福音工作.
應該怎樣去改變.
來去回應這個時代的轉變.
我們人的需要都改變了.
弟姐妹.
我們不能夠再用.
以前做年青人那一套.
來繼續做年青人的工作.
以前的年青人呢.
我做拿西爾導師.
的時代呢.
請他們吃東西.
和他們一起踢球.
和他們一起去玩.
他們就.
很開心地回團體.
有空打電話問他.
他們就.
都很乖地回團體.
但是現在.
整個時代不同了.
你再打電話找他.
他都不會理你.

$^{481}$你WhatsApp他.
已讀不回.
時代的轉變.
很快.
還有就是.
我經常會體現一件事.
我們的年青一代.
平板電腦大.
他們在看.
Netflix.
那些大.
所以他們在聽.
K-pop大.
所以他們在聽.
這個世界最頂級的.
娛樂長大的.
所以.
我們會問.
究竟年青的一代.
他們回到教會.
我們現在的敬拜.
是否真的能夠滿足.
年青一代的需要.
這是我們很需要.
去想.
我特別要.
提一提這幅圖.
是這幾天.
在網絡上.
不停有人分享.
這幅圖.
來自一篇文.
那篇文是討論.
在教會裡面.
老年人經常嫌歌太嘈.
年青人.
嫌歌太老.
今天我們已經很稀奇.
我們有個.
不算太老的年青人.

$^{521}$唱著.
我願將你.
但是.
這個分享.
很反映一件事.
很多教會仍然停留在.
年長的和年輕的.
仍然有很多的.
矛盾.
有很多的爭拗.
所以今天我們要去.
想的是.
不是說哪個對哪個錯.
而是說究竟.
當新的工作.
新的作為臨到的時候.
我們應該要怎樣回應呢.
這幅圖就是.
我們寫著土瓜灣.
比利街.
的項目.
當年在重建之後.
是有一千.
一百五十火.
規模是區內最大.
弟兄姊妹.
我們怎樣迎接神.
新的作為呢.
然後一項的門徒.
他們對耶穌的門徒.
不禁食.
感到困惑和不解.
因為他們都以禁食.
來表達對神的.
禁忌.
但竟然耶穌和他們的門徒.
不單止他們不禁食.
他們更加和罪人和罪理.
吃飯.
當時的環境是這樣的.

$^{561}$首先.
他們認為.
如果你是敬天的.
你是應該要禁食的.
現在耶穌和他的門徒.
他們不禁食.
你應該要禁食.
你沒有禁食還要吃飯.
這是不對的.
再加上.
你去吃飯的時候.
竟然是和罪人和罪理吃飯.
這是不對的.
再加上是不對的.
所以你明白.
當時約翰的門徒.
他們見到耶穌的門徒不禁食.
會感到困惑和不解.
但耶穌的回答是.
新郎和陪伴的人.
同在的時候.
陪伴的人豈能哀動呢.
這是一幅婚宴的圖畫.
在猶太人的文化當中.
婚宴是一個歡喜快樂的時候.
沒有人會在.
婚禮上禁食.
這對我們華人來說.
應該這樣說.
對人來說.
我們都可以理解的道理.
我想問問大家.
有沒有誰你去喝酒的時候.
你會選擇禁食.
我們沒有人會這樣做.
當時的環境就是這樣.
耶穌這樣回答.
其實是稱呼自己.
是一個新郎.
這同時.

$^{601}$是另一個令到當時.
這班文士.
他們聽到是一個很震撼的.
自我宣告.
因為在舊約裡面.
神是以色列的丈夫.
所以現在耶穌說.
新郎就在這裡.
意思就是說.
神親自臨到人間.
這個時代已經開始.
祂說祂自己就是神.
所以對當時的人一聽.
是另一個很震撼的宣告.
而主耶穌的表達.
是說禁食.
其實是象徵著.
哀傷和等候.
但現在不是一個等候的時間.
現在是一個歡慶的時候.
僅此而已.
耶穌沒有否定禁食.
祂說日子將到.
新郎要離開他們.
所以換句話說.
其實祂在說一個很重要的事.
不是不可以禁食.
不是不可以哀傷.
不是不可以等候.
而是說現在不是那個時間.
所以屬靈操練的意義.
是必須要配合.
時間和作為.
我們信耶穌不是.
要去守一些系列的規矩.
而是我們要與主同行.
我們要知道神在做什麼.
看神在做什麼.
神做什麼我們就做什麼.
神想我們做什麼我們就做什麼.

$^{641}$有時候可能神要我們.
安靜要悔改.
我們就去安靜悔改.
有時候神要我們.
要去慶賀我們就去喜樂慶賀.
屬靈生命的成熟.
不是說我們一直都在做.
同樣的一件事.
或者某樣的事情.
而是我們要分辨究竟這一刻.
神在這裡.
正是要做什麼.
弟兄姊妹.
今天你的生命.
是停留在.
一個哀傷的模式.
還是活在.
與主同在的.
一種喜樂裡面呢.
當我們真正去認識.
耶穌是新郎的時候.
信仰就不再只是.
一個責任.
而是一種關係.
經篇是說我們要去跟隨耶穌.
而不是跟隨.
律法人的傳統規條.
而是一份.
我們與神之間的同行的關係.
所以教會的.
更新是很需要.
看見需要.
很需要把握時機.
我記得當年.
我在土鎮上了一個.
門徒訓練.
我們用的一本書是.
不再一樣.
Chris是我其中一個的同學.
在.

$^{681}$不再一樣裡面.
我很深刻.
到今天我還記得.
有一章.
有一章裡面是這樣說的.
說到馴服的行動.
為什麼.
我這麼深刻呢.
因為他特別提到當我們看見.
神的工作的時候.
我們是立刻以信心.
去採取行動.
哪怕是一個微小的嘗試.
所以當.
其實這個提醒是說.
當我們真的看見.
神正在工作.
我們是應該要去用.
信心來去.
採取行動.
不用擔心.
那個嘗試可能會有什麼問題.
最重要的是我們要看見的是.
神的工作.
如果當我們看見.
神真的.
祝福著.
一群很特別的人.
或者可能一個新的群體的時候.
或者.
神在我們當中有新的事工出現的時候.
神真的正在工作的時候.
我們應該是以信心.
來去採取行動.
哪怕是一個微小的嘗試.
傳統的教會.
我們經常都很.
介懷一件事.
很介懷什麼呢.
很介懷我們做了一個事工.

$^{721}$然後就.
最後做不成.
做不成我們就要接了它.
其實弟兄姊妹.
如果.
神真的在這個階段要做一個事工.
而這個事工也會有一個結束的日子的話.
我們其實都應該是要以.
信心去採取行動.
哪怕只是一個微小的嘗試.
所以這裡說的最重要的東西是什麼呢.
最重要的是我們要看到需要.
要擺好時機.
你去會映神.
然後耶穌他用了.
兩個新舊的比喻.
他說新布呢.
就補在舊衣服.
另一個就是新袖就裝在.
舊皮袋裡面.
我沒有刻意.
去跟別人說.
我今個禮拜講道的經文.
也沒有跟別人分享.
我不知道我身邊的人.
我遇到的人他們是不是.
看到教會的.
崇拜周刊.
同時間在不同的環境.
有人跟我提這個比喻.
前晚.
因為我太太幫我兒子.
扯開褲子.
因為要補兩個洞.
我順便跟她說.
你也幫我補一條西褲.
因為也有些洞.
我還跟她說.
你可不可以幫我.
補褲袋.

$^{761}$因為我的褲袋.
不知道為什麼買回來的時候.
已經溶掉了.
我又沒有留意.
我發布幫我.
扯開褲袋.
我老婆馬上跟我說.
你沒有看聖經嗎.
聖經不是說.
新報不可以補舊報.
會扯爛的.
昨天.
因為我.
上午就去.
處理了一個.
安息禮拜.
吃便飯的時候.
有一個弟兄.
突然之間.
很興奮地跟我們說.
我最近聽了一個分享.
他用了一個比喻.
這個比喻說得很好.
聽他說什麼呢.
他說.
在說一些經文.
新報補在舊衣.
新手裝在舊皮袋.
他用了一個比喻.
是這麼說的.
如果你的電腦.
你換了一個硬件.
但是你裝的軟件.
全部都是舊的.
你都會變成一個舊的硬件.
你的軟件都沒有變.
所以他說.
這個比喻真的很好.
我沒有跟別人分享.
不知道為什麼身邊的人.

$^{801}$跟我說.
其實我們要了解這兩個比喻.
其實新報補在舊衣.
是會扯爛的.
因為新報補上去.
是會縮水的.
他不能夠.
牢固地釘在舊衣裡.
所以反而.
當他縮水的時候.
就會扯爛了舊衣.
而酒.
只在山羊.
或者綿羊製成的皮袋裡.
新酒裝在皮袋裡.
因為新的酒.
仍然在發酵.
所以發酵的時候.
會產生氣體向外.
向外施加壓力.
新皮袋就會.
有彈性.
皮袋受壓的時候.
是沒有問題的.
但是舊皮袋.
因為皮袋已經失去.
彈性的時候.
新鮮發酵的酒.
裝在舊皮袋裡.
就會使舊皮袋爆裂.
就會漏出酒來.
在這裡我們去看.
似乎在表達.
新和舊.
很難合在一起.
這個可能是我們經常會想.
當我們說新的時候.
舊的我們會覺得.
你只是不要我們.
你只是不需要我們.

$^{841}$新和舊我們經常會有.
矛盾和衝突.
一說新的.
就是不喜歡舊的.
覺得舊的不好.
但在這裡他不是真的在說.
新舊難合.
他在說.
不要以舊的方式.
靠行為得救.
來套用在新的教導.
就是靠恩典得救.
這個是耶穌說的.
比喻最重要的意思.
所以兩個比喻.
都是在說同一個重點.
就是新和舊.
是不能隨便去混合.
耶穌不是要改良.
猶太教.
而是要帶來一個.
承傳與更新.
我們要明白.
這個靠律法.
靠行為清移.
是一直都怎樣.
做不到的時候.
所以耶穌基督.
他的彌臨是帶來一個.
承傳與更新.
就是讓人是怎樣.
靠恩典去清移.
因為我們若果再.
以律法.
和行為去清移.
其實我們不能夠.
清移.
這正正是為什麼.
有耶穌基督的救恩.
最重要的一個原因.

$^{881}$所以福音不是舊的制度的.
一個補丁或補償.
或補充.
不是在原有的宗教框架.
加上一塊布.
而是福音說的是.
全新的救恩時代.
如果我們只是想在.
原有的生活方式上.
加一些耶穌.
但他不願意.
用生命去更新.
我們的結果.
就是那塊舊衣服.
就會撕裂.
那個舊皮袋.
就會破裂.
所以今天對我們很多基督徒來說.
我們最大的掙扎就是這樣.
我們很想跟隨主.
但我們價值觀.
不願意改變.
我們都很想經歷神日.
但我們生活的優先次序.
都不肯改變.
我們都很想突破.
但我們的心態.
又不肯改變.
新布補在舊衣裡.
只會裂得更大.
所以今天我們要知道.
神不是要修補我們的人生.
神是要更新.
我們的人生.
所以在整段經文.
作為最後的一句.
亦是整段經文的.
高潮.
為了將新酒裝在.
新皮袋裡.

$^{921}$兩樣就都.
保存了.
新酒就代表神的國.
一個新的生命,新的能力,新的恩典.
新皮袋就代表著一個.
願意被更新的.
生命.
所以今天各位姐妹.
我們需要的是成為一個.
更新的生命.
新酒本身是好的.
問題不在酒裡,而是在皮袋裡.
神今天仍然要.
設下新酒.
新的恩典,新的呼召.
新的意象.
土鎮是一間63年的教會.
但是神仍然有新的恩典.
新的呼召,新的意象.
給我們.
而問題是我們願不願意成為一個.
新的皮袋.
新的皮袋有幾種特質.
首先第一.
它是柔軟的.
它是願意被塑造的.
第二,它是有彈性的.
它是能夠承受成長的張力.
第三,它是能夠敞開的.
它是開得到的.
它是不會固守舊有的框架.
各位姐妹.
神不是要我們否定過去.
我們過去做得很好的.
我們所做的一切.
神不是要否定我們過去.
而是要帶領我們進入一個更豐富的將來.
我們的教會可以更新.
我們的家庭也可以更新.
你的個人的屬靈的生命.

$^{961}$也可以更新.
所以我們不要害怕.
要改變.
我們要確定的是.
這一切的改變.
都是來自神.
所以今天.
我們回到.
最初的問題.
變還是不變.
各位姐妹.
真理是永遠不變的.
神的愛也是不變的.
福音是不變的.
耶穌基督的昨日和今天.
一直都是永遠一樣.
祂昔日怎樣帶領土瓜人浸信會.
今天也仍然繼續.
帶領著土瓜人浸信會.
但是我們的生命.
我們需要不斷地被更新.
因為當神的國度臨到.
我們若果堅持停留在舊的框架.
就會錯過了.
祂新的作為.
所以要讓我們成為.
一個新的皮袋.
讓神的新酒在我們裡面.
去發酵.
讓我們的生命充滿彈性和喜樂.
改變未必是容易的.
我們都明白.
從來都不是一件容易的事.
但是我們去.
馴服神的更新.
一定會帶來祝福.
所以心願我們.
沒有成為一個守舊的人.
而我們能夠迎接神新時代的人.
我們懇求主.

$^{1001}$去幫助我們.
在我們裡面讓我們勇敢地.
去為道主義去改變.
最後我想和大家做一個練習.
你不需要回答我的.
不過你可以自己想想.
你可以去想一想.
十年前的你.
和現在的你.
你可以想一想.
十年前的你.
你最喜歡去的地方.
是哪裡呢.
而現在的你.
你最喜歡去的地方.
又是哪裡呢.
是不同的地方.
還是同一個地方呢.
十年前的你.
你最喜歡吃的東西.
今天你是否仍然喜歡呢.
還是你已經喜歡了.
吃其他的食品呢.
十年前的你 你最喜歡做的事.
十年後的今天.
你是否仍然喜歡做呢.
還是已經喜歡了.
其他的事情呢.
十年前你最討厭的物件.
十年後的今天.
是否都不同了.
還是已經沒那麼討厭.
還是可能你正在討厭其他東西呢.
如果當你發現.
十年前的你.
和十年後的今天.
是有改變的話.
在這裡告訴你一個很重要的訊息.
就是.
弟兄姊妹.

$^{1041}$你是可以改變的.
好不好都幫我和你.
旁邊那個說一聲.
你是可以改變的.
這裡有一段很好的時間.
差主耶穌我們去感激讚美你.
降世為人.
為我們成就救恩.
你讓我們不需要再去.
靠六法.
靠行為來去.
清儀.
讓我們可以靠恩順你.
而能夠去清儀.
主我們要去.
獻上感恩.
主我們今天來到你的面前.
求你幫助我們.
開放我們的生命.
讓我們真的能夠.
確保我們的生命都能夠去更生.
以致我們真的在.
不單止我們的教會生活也好.
我們教會.
坐落在土國環榮光街.
這個地方也好.
我們都能夠用一個更生的生命.
能夠可以.
祝福.
我們的弟兄姊妹.
更加可以祝福.
我們身處的這個社區.
主我們懇求你幫助我們.
可能拿走我們一些.
我們的頑固.
我們的固執.
讓我們更加願意的.
是學著你.
以致我們真的能夠可以.
祝福到更多的人.

$^{1081}$主願意親自去幫助.
賜福我們.
使用我們.
主我們多謝讚美你.
我們感其道交托.
是奉靠我主耶穌基督的聖名.
\newpage



\section{約翰福音 4:5-5:42}
\label{sec:031765}
\textbf{活水之道}~~~~李均熊牧師博士
\newline
\newline
連結: \href{http://www.church.com.hk/acms/content.asp?site=cdc\&op=show\&type=product\&code=031765\&layout=sermon}{\texttt{031765}} ~~~~ 語音日期: 2026-03-08
\newline
\newline
\hyperref[sec:031764]{\small{< < < PREV SERMON < < <}}
~
\hyperref[sec:index_chronic]{\small{[返順時目]}}
~
\hyperref[sec:031766]{\small{> > > NEXT SERMON > > >}}
\newline
\newline
約翰福音 4:5-5:42
\newline
\begin{longtable}{cl}
\hline
\hline
章節 & 經文 (和合本修訂版)\\
\hline
4:5 & \begin{tabularx}{0.7\textwidth}{X} 於是到了撒瑪利亞的一座城，名叫敘加，靠近雅各給他兒子約瑟的那塊地。 \end{tabularx} \\ \\ \relax
4:6 & \begin{tabularx}{0.7\textwidth}{X} 雅各井就在那裡；耶穌因旅途疲乏，坐在井旁。那時約是正午。 \end{tabularx} \\ \\ \relax
4:7 & \begin{tabularx}{0.7\textwidth}{X} 有一個撒瑪利亞婦人來打水。耶穌對她說：「請給我水喝。」 \end{tabularx} \\ \\ \relax
4:8 & \begin{tabularx}{0.7\textwidth}{X} 因為那時門徒進城買食物去了。 \end{tabularx} \\ \\ \relax
4:9 & \begin{tabularx}{0.7\textwidth}{X} 撒瑪利亞婦人對他說：「你是猶太人，怎麼向我一個撒瑪利亞女人要水喝呢？」因為猶太人和撒瑪利亞人沒有來往。 \end{tabularx} \\ \\ \relax
4:10 & \begin{tabularx}{0.7\textwidth}{X} 耶穌回答她說：「你若知道神的恩賜，和對你說『請給我水喝』的是誰，你早就會求他，他也早就會給了你活水。」 \end{tabularx} \\ \\ \relax
4:11 & \begin{tabularx}{0.7\textwidth}{X} 婦人對耶穌說：「先生，你沒有打水的器具，井又深，哪裡去取活水呢？ \end{tabularx} \\ \\ \relax
4:12 & \begin{tabularx}{0.7\textwidth}{X} 我們的祖宗雅各把這井留給我們，他自己和兒女以及牲畜都喝這井裡的水，難道你比他還大嗎？」 \end{tabularx} \\ \\ \relax
4:13 & \begin{tabularx}{0.7\textwidth}{X} 耶穌回答，對她說：「凡喝這水的，還要再渴； \end{tabularx} \\ \\ \relax
4:14 & \begin{tabularx}{0.7\textwidth}{X} 誰喝我所賜的水，就永遠不渴。我所賜的水要在他裡面成為泉源，直湧到永生。」 \end{tabularx} \\ \\ \relax
4:15 & \begin{tabularx}{0.7\textwidth}{X} 婦人對他說：「先生，請把這水賜給我，使我不渴，也不用到這裡來打水。」 \end{tabularx} \\ \\ \relax
4:16 & \begin{tabularx}{0.7\textwidth}{X} 耶穌對她說：「你去，叫你的丈夫，再到這裡來。」 \end{tabularx} \\ \\ \relax
4:17 & \begin{tabularx}{0.7\textwidth}{X} 婦人回答，對耶穌說：「我沒有丈夫。」耶穌說：「你說沒有丈夫是對的。 \end{tabularx} \\ \\ \relax
4:18 & \begin{tabularx}{0.7\textwidth}{X} 你已經有過五個丈夫，你現在有的並不是你的丈夫。你這話是真的。」 \end{tabularx} \\ \\ \relax
4:19 & \begin{tabularx}{0.7\textwidth}{X} 婦人對他說：「先生，我看你是一位先知。 \end{tabularx} \\ \\ \relax
4:20 & \begin{tabularx}{0.7\textwidth}{X} 我們的祖宗在這山上敬拜神，你們倒說，應當敬拜的地方是在耶路撒冷。」 \end{tabularx} \\ \\ \relax
4:21 & \begin{tabularx}{0.7\textwidth}{X} 耶穌對她說：「婦人，你要信我。時候將到，你們敬拜父，既不在這山上，也不在耶路撒冷。 \end{tabularx} \\ \\ \relax
4:22 & \begin{tabularx}{0.7\textwidth}{X} 你們所敬拜的，你們不知道；我們所敬拜的，我們知道，因為救恩是從猶太人出來的。 \end{tabularx} \\ \\ \relax
4:23 & \begin{tabularx}{0.7\textwidth}{X} 時候將到，現在就是了，那真正敬拜父的，要用心靈和誠實敬拜他，因為父要這樣的人敬拜他。 \end{tabularx} \\ \\ \relax
4:24 & \begin{tabularx}{0.7\textwidth}{X} 神是靈，所以敬拜他的必須用心靈和誠實敬拜他。」 \end{tabularx} \\ \\ \relax
4:25 & \begin{tabularx}{0.7\textwidth}{X} 婦人對他說：「我知道彌賽亞—就是那稱為基督的—要來；他來了，會把一切的事都告訴我們。」 \end{tabularx} \\ \\ \relax
4:26 & \begin{tabularx}{0.7\textwidth}{X} 耶穌對她說：「我就是，正在跟你說話呢！」 \end{tabularx} \\ \\ \relax
4:27 & \begin{tabularx}{0.7\textwidth}{X} 正在這時，門徒回來了。他們對耶穌正在和一個婦人說話感到驚訝，可是沒有人說：「你要甚麼？」或說：「你為甚麼和她說話？」 \end{tabularx} \\ \\ \relax
4:28 & \begin{tabularx}{0.7\textwidth}{X} 那婦人留下水罐，往城裡去，對眾人說： \end{tabularx} \\ \\ \relax
4:29 & \begin{tabularx}{0.7\textwidth}{X} 「你們來看！有一個人把我素來所做的一切事都說了出來，難道這個人就是基督嗎？」 \end{tabularx} \\ \\ \relax
4:30 & \begin{tabularx}{0.7\textwidth}{X} 他們就出城，來到耶穌那裡。 \end{tabularx} \\ \\ \relax
4:31 & \begin{tabularx}{0.7\textwidth}{X} 就在這個時候，門徒求耶穌說：「拉比，請吃吧。」 \end{tabularx} \\ \\ \relax
4:32 & \begin{tabularx}{0.7\textwidth}{X} 耶穌對他們說：「我有食物吃，是你們不知道的。」 \end{tabularx} \\ \\ \relax
4:33 & \begin{tabularx}{0.7\textwidth}{X} 門徒就彼此說：「難道有人拿甚麼給他吃了嗎？」 \end{tabularx} \\ \\ \relax
4:34 & \begin{tabularx}{0.7\textwidth}{X} 耶穌對他們說：「我的食物就是要遵行差我來那位的旨意，完成他的工作。 \end{tabularx} \\ \\ \relax
4:35 & \begin{tabularx}{0.7\textwidth}{X} 你們不是說『到收割的時候還有四個月』嗎？我告訴你們，舉目向田觀看，莊稼熟了，可以收割了。 \end{tabularx} \\ \\ \relax
4:36 & \begin{tabularx}{0.7\textwidth}{X} 收割的人已經得工錢，為永生儲存五穀，使撒種的和收割的一同快樂。 \end{tabularx} \\ \\ \relax
4:37 & \begin{tabularx}{0.7\textwidth}{X} 『那人撒種，這人收割』，這話可見是真的。 \end{tabularx} \\ \\ \relax
4:38 & \begin{tabularx}{0.7\textwidth}{X} 我差你們去收你們所沒有辛勞的；別人辛勞，你們享受他們辛勞的成果。」 \end{tabularx} \\ \\ \relax
4:39 & \begin{tabularx}{0.7\textwidth}{X} 那城裡有好些撒瑪利亞人信了耶穌，因為那婦人作見證，說：「他把我素來所做的一切事都說了出來。」 \end{tabularx} \\ \\ \relax
4:40 & \begin{tabularx}{0.7\textwidth}{X} 於是撒瑪利亞人來見耶穌，求他在他們那裡住下，他就在那裡住了兩天。 \end{tabularx} \\ \\ \relax
4:41 & \begin{tabularx}{0.7\textwidth}{X} 因為耶穌的話，信的人就更多了。 \end{tabularx} \\ \\ \relax
4:42 & \begin{tabularx}{0.7\textwidth}{X} 他們對那婦人說：「現在我們信，不再是因為你的話，而是我們親自聽見了，知道這人真是世界的救主。」 \end{tabularx} \\ \\ \relax
4:43 & \begin{tabularx}{0.7\textwidth}{X} 過了那兩天，耶穌離開那地方，往加利利去。 \end{tabularx} \\ \\ \relax
4:44 & \begin{tabularx}{0.7\textwidth}{X} 因為耶穌自己作過見證說：「先知在自己的家鄉是沒有人尊敬的。」 \end{tabularx} \\ \\ \relax
4:45 & \begin{tabularx}{0.7\textwidth}{X} 到了加利利，加利利人都歡迎他，因為他們也上耶路撒冷去過節，曾經看過他在節期間所做的一切事。 \end{tabularx} \\ \\ \relax
4:46 & \begin{tabularx}{0.7\textwidth}{X} 耶穌又到了加利利的迦拿，就是他從前變水為酒的地方。有一個大臣，他的兒子在迦百農病了。 \end{tabularx} \\ \\ \relax
4:47 & \begin{tabularx}{0.7\textwidth}{X} 他聽見耶穌從猶太到了加利利，就來見他，求他下去醫治他的兒子，因為他兒子快要死了。 \end{tabularx} \\ \\ \relax
4:48 & \begin{tabularx}{0.7\textwidth}{X} 耶穌對他說：「若不看見神蹟奇事，你們總是不信。」 \end{tabularx} \\ \\ \relax
4:49 & \begin{tabularx}{0.7\textwidth}{X} 那大臣對他說：「先生，求你趁著我的孩子還沒有死就下去吧。」 \end{tabularx} \\ \\ \relax
4:50 & \begin{tabularx}{0.7\textwidth}{X} 耶穌對他說：「回去吧，你的兒子會活！」那人信耶穌所說的話，就回去了。 \end{tabularx} \\ \\ \relax
4:51 & \begin{tabularx}{0.7\textwidth}{X} 正下去的時候，他的僕人迎面而來，說他的兒子活了。 \end{tabularx} \\ \\ \relax
4:52 & \begin{tabularx}{0.7\textwidth}{X} 他就問甚麼時候見好的。他們對他說：「昨天下午一點鐘熱就退了。」 \end{tabularx} \\ \\ \relax
4:53 & \begin{tabularx}{0.7\textwidth}{X} 他就知道這正是耶穌對他說「你的兒子會活」的時候；他自己和全家就都信了。 \end{tabularx} \\ \\ \relax
4:54 & \begin{tabularx}{0.7\textwidth}{X} 這是耶穌從猶太回到加利利後所行的第二個神蹟。 \end{tabularx} \\ \\ \relax
5:1 & \begin{tabularx}{0.7\textwidth}{X} 這些事以後，到了猶太人的一個節期，耶穌上耶路撒冷去。 \end{tabularx} \\ \\ \relax
5:2 & \begin{tabularx}{0.7\textwidth}{X} 在耶路撒冷，靠近羊門有一個池子，希伯來話叫畢士大，旁邊有五個柱廊； \end{tabularx} \\ \\ \relax
5:3 & \begin{tabularx}{0.7\textwidth}{X} 裡面躺著許多病人，有失明的、瘸腿的、癱瘓的。等候水動， \end{tabularx} \\ \\ \relax
5:4 & \begin{tabularx}{0.7\textwidth}{X} 因為有天使按時下池子攪動那水，水動之後，先下水的人無論患甚麼病都能得痊癒。 \end{tabularx} \\ \\ \relax
5:5 & \begin{tabularx}{0.7\textwidth}{X} 在那裡有一個人，病了三十八年。 \end{tabularx} \\ \\ \relax
5:6 & \begin{tabularx}{0.7\textwidth}{X} 耶穌看見他躺著，知道他病了很久，就問他：「你要痊癒嗎？」 \end{tabularx} \\ \\ \relax
5:7 & \begin{tabularx}{0.7\textwidth}{X} 病人回答他：「先生，水動的時候，沒有人把我放在池子裡；我正要去的時候，別人比我先下去了。」 \end{tabularx} \\ \\ \relax
5:8 & \begin{tabularx}{0.7\textwidth}{X} 耶穌對他說：「起來，拿起你的褥子走吧！」 \end{tabularx} \\ \\ \relax
5:9 & \begin{tabularx}{0.7\textwidth}{X} 那人立刻痊癒，就拿起自己的褥子走了。那天是安息日， \end{tabularx} \\ \\ \relax
5:10 & \begin{tabularx}{0.7\textwidth}{X} 所以猶太人對那被治好了的人說：「今天是安息日，你拿褥子是不合法的。」 \end{tabularx} \\ \\ \relax
5:11 & \begin{tabularx}{0.7\textwidth}{X} 他卻回答他們：「那使我痊癒的人對我說：『拿起你的褥子走吧！』」 \end{tabularx} \\ \\ \relax
5:12 & \begin{tabularx}{0.7\textwidth}{X} 他們問他：「對你說『拿起褥子走』的是甚麼人？」 \end{tabularx} \\ \\ \relax
5:13 & \begin{tabularx}{0.7\textwidth}{X} 那治好了的人不知道那人是誰，因為那裡人很多，耶穌已經躲開了。 \end{tabularx} \\ \\ \relax
5:14 & \begin{tabularx}{0.7\textwidth}{X} 後來耶穌在聖殿裡找到他，對他說：「你已經痊癒了，不要再犯罪，免得你的遭遇更壞。」 \end{tabularx} \\ \\ \relax
5:15 & \begin{tabularx}{0.7\textwidth}{X} 那人就去告訴猶太人，使他痊癒的是耶穌。 \end{tabularx} \\ \\ \relax
5:16 & \begin{tabularx}{0.7\textwidth}{X} 所以猶太人迫害耶穌，因為他在安息日做了這些事。 \end{tabularx} \\ \\ \relax
5:17 & \begin{tabularx}{0.7\textwidth}{X} 耶穌就回答他們：「我父做事直到如今，我也做事。」 \end{tabularx} \\ \\ \relax
5:18 & \begin{tabularx}{0.7\textwidth}{X} 為了這緣故，猶太人越發想要殺他，因為他不但犯了安息日，而且稱神為他的父，把自己和神看為同等。 \end{tabularx} \\ \\ \relax
5:19 & \begin{tabularx}{0.7\textwidth}{X} 於是耶穌回答，對他們說：「我實實在在地告訴你們，子憑著自己不能做甚麼，惟有看見父所做的，他才做；父所做的事，子也照樣做。 \end{tabularx} \\ \\ \relax
5:20 & \begin{tabularx}{0.7\textwidth}{X} 父愛子，將自己所做的一切事指示給他看，還要將比這更大的事給他看，使你們驚訝。 \end{tabularx} \\ \\ \relax
5:21 & \begin{tabularx}{0.7\textwidth}{X} 父怎樣叫死人復活，賜他們生命，子也照樣隨自己的意願賜人生命。 \end{tabularx} \\ \\ \relax
5:22 & \begin{tabularx}{0.7\textwidth}{X} 父不審判任何人，而是把審判的事全交給子， \end{tabularx} \\ \\ \relax
5:23 & \begin{tabularx}{0.7\textwidth}{X} 為要使人都尊敬子，如同尊敬父一樣。不尊敬子的，就是不尊敬差子來的父。 \end{tabularx} \\ \\ \relax
5:24 & \begin{tabularx}{0.7\textwidth}{X} 「我實實在在地告訴你們，那聽我話又信差我來那位的，就有永生，不至於被定罪，而是已經出死入生了。 \end{tabularx} \\ \\ \relax
5:25 & \begin{tabularx}{0.7\textwidth}{X} 我實實在在地告訴你們，時候將到，現在就是了，死人要聽見神兒子的聲音，聽見的人就要活了。 \end{tabularx} \\ \\ \relax
5:26 & \begin{tabularx}{0.7\textwidth}{X} 因為父怎樣自己裡面有生命，也照樣賜給他兒子自己裡面有生命， \end{tabularx} \\ \\ \relax
5:27 & \begin{tabularx}{0.7\textwidth}{X} 並且賜給他施行審判的權柄，因為他是人子。 \end{tabularx} \\ \\ \relax
5:28 & \begin{tabularx}{0.7\textwidth}{X} 你們不要對這事感到驚訝，因為時候將到，凡在墳墓裡的，都要聽見他的聲音， \end{tabularx} \\ \\ \relax
5:29 & \begin{tabularx}{0.7\textwidth}{X} 並且要出來：行善的，復活得生命；作惡的，復活被定罪。 \end{tabularx} \\ \\ \relax
5:30 & \begin{tabularx}{0.7\textwidth}{X} 「我憑著自己不能做甚麼。我怎麼聽見就怎麼審判，而我的審判是公平的，因為我不尋求自己的意願，只尋求差我來那位的旨意。」 \end{tabularx} \\ \\ \relax
5:31 & \begin{tabularx}{0.7\textwidth}{X} 「我若為自己作見證，我的見證就不真。 \end{tabularx} \\ \\ \relax
5:32 & \begin{tabularx}{0.7\textwidth}{X} 另有一位為我作見證，我也知道他為我作的見證是真的。 \end{tabularx} \\ \\ \relax
5:33 & \begin{tabularx}{0.7\textwidth}{X} 你們曾差人到約翰那裡，他為真理作過見證。 \end{tabularx} \\ \\ \relax
5:34 & \begin{tabularx}{0.7\textwidth}{X} 其實，我所受的見證不是從人來的；然而，我說這些話是為了使你們得救。 \end{tabularx} \\ \\ \relax
5:35 & \begin{tabularx}{0.7\textwidth}{X} 約翰是點亮的明燈，你們情願因他的光歡欣一時。 \end{tabularx} \\ \\ \relax
5:36 & \begin{tabularx}{0.7\textwidth}{X} 但我有比約翰更大的見證：父交給我去完成的工作，就是我正在做的，為我作證是父差遣了我。 \end{tabularx} \\ \\ \relax
5:37 & \begin{tabularx}{0.7\textwidth}{X} 那差我來的父也為我作了見證。你們從來沒有聽見他的聲音，也沒有看見他的形像。 \end{tabularx} \\ \\ \relax
5:38 & \begin{tabularx}{0.7\textwidth}{X} 你們並沒有他的道存在心裡，因為你們不信他所差來的那一位。 \end{tabularx} \\ \\ \relax
5:39 & \begin{tabularx}{0.7\textwidth}{X} 你們查考聖經，因你們以為其中有永生；而這經正是為我作見證的。 \end{tabularx} \\ \\ \relax
5:40 & \begin{tabularx}{0.7\textwidth}{X} 然而，你們不肯到我這裡來得生命。 \end{tabularx} \\ \\ \relax
5:41 & \begin{tabularx}{0.7\textwidth}{X} 「我不接受從人來的榮耀， \end{tabularx} \\ \\ \relax
5:42 & \begin{tabularx}{0.7\textwidth}{X} 但我知道，你們沒有愛神的心。 \end{tabularx} \\ \\
[1ex]
\hline
\hline
\end{longtable}
$^{1}$弟兄姊妹平安.
當我按木之後.
我就開始在崇基禮拜堂講道.
其實第一篇道就是三年前.
講舊約的那段經文.
叫做無水之道.
今天我就不講舊約了.
我就講福音書.
大福音書是講到井旁的婦人.
我想對大家來說都是很熟悉的一段經文.
熟悉到很多的細節.
很多的理由都已經聽過.
為什麼耶穌會見那個婦人.
或者證悟這些東西.
所以為了令大家不要太過.
從一個熟悉的角度去理解這個故事.
我想從另一個角度跟大家.
去講一下《約翰福音》這個記載.
我今天會選擇一個視角.
就是一個跟婦人同住同一個城市.
也就是在撒瑪利亞地區裡最佳城的居民.
這個角度.
透過他們的理解去看整個故事.
所以我是會用一個第一身的角度去講.
所以我不知道會不會令到全日的弟兄姊妹覺得很困難.
不過在開始之前我都可能要介紹一下.
其實這篇經文有很多的特點.
剛才讀經員用了很長時間去讀經.
不知道大家有沒有發現.
其實這段經文是四卷福音書裡最長的一段對話.
所以你看到有很多節的經文.
其次我們大齋期.
特別是甲年的大齋期.
我們叫EIA.
甲年的大齋期.
通常的福音經課都會選用馬太福音.
所以我們有時會叫做馬太年.
不過在大齋期裡.
特別是大齋期的第二,三,四,五的主日裡.
就會選用了約翰福音的經文.

$^{41}$大齋期的第一主日.
通常都會選用耶穌受試探的經文.
約翰福音就沒有耶穌受試探.
大齋期的第六主日就是宗之主日.
有時都會用約翰福音.
今年就會用馬太福音.
但是二,三,四,五主日都是用約翰福音的經文.
如果你回家可以留意一下.
或者你再看看的時候.
它都是說耶穌與某人相遇的記載.
上個星期就是大齋期的第二主日.
就是耶穌與尼哥德姆相遇.
今天就是與井婆的婦人相遇.
下個星期就會與一個生來合眼的人相遇.
到了大齋期的第五主日.
就會說拉薩路復活.
就是耶穌與拉薩路和他的家人相遇.
所以其實很明顯編寫大齋期福音經課的人士.
是很有心思的.
特別如果你將它串連的話.
上個星期是說尼哥德姆在夜間探問耶穌.
今天大齋期的第三主日.
就是耶穌在正午的時候.
在井旁與一個曾經有四個丈夫.
現在與他在一起的那個不是他丈夫的.
撒瑪利亞婦人相遇.
上個星期是德高望重的尼哥德姆.
今天是一個可能名譽不太好的女士.
下個星期就會在西羅亞池旁附近.
與一個生來合眼的人相遇.
到了大齋期的第五主日.
就在墳地與拉薩路和他的姐姐相遇.
這個編排其實在一些教會傳統.
就是預備復活節的洗禮.
就是由不同人與基督的相遇.
如何培養信心.
由尼哥德姆懵懵懈懈.
到拉薩路的復活.
所以今天這首詩歌是由灰到活水.
其實是說洗禮的進程.

$^{81}$在這幾個大齋期的主日裡面.
正正告訴我們信仰的歷程有很多種.
在過程中我們發現基督.
而且得到永遠的生命.
所以大家可能完結之後.
或者接下來的星期天聽到的時候.
就試試可否串連這四個主日裡面的經課.
看看我究竟在哪個屬靈的旅程的階段.
今天我們就從撒瑪利亞最佳居民的角度.
去重新再說《約翰福音》四章五至四十二節的記載.
要上下身.
朋友 我不是耶穌.
我不是門徒.
我也不是當天耶穌跟他說話的撒瑪利亞婦人.
我是最佳城裡面的其中一個居民.
就是那些經常被視為背景人物的人.
可能你有時也會有一種感覺.
就是你不是主角.
在一個世界大棋裡面.
一個普普通通不起眼的普通居民.
所謂路人甲乙丙.
但平凡如我都可以做我的主角.
所以我今天想從我的角度去告訴你.
那天在井邊發生的事.
在井邊發生的事.
當初改變了的不只是那個女人.
還令我們全城的人都捲入在這個改變裡面.
現在就讓我從我的角度.
去說這件令我們全城震撼的事.
我們城裡面的每一個人都認識她.
你們口中所說的撒瑪利亞婦人.
又或者井旁的女人.
不知道為什麼使徒忘記了她的名字.
但是我們最佳城的居民.
會叫她那個女人.
我們每次提起她都說那個女人.
有些人甚至用一些更加不好聽的稱呼去稱呼她.
說真的.
她是我們在城裡面最不想遇到的人.
感情史複雜.

$^{121}$名聲不好.
說她兩句沒有人會幫她說好話.
但其實我們都知道.
其實她是一個苦命的女人.
在我們生活的那個年代.
一個女人沒有男人的依靠.
根本就生存不了.
她好像一般的年輕女士一樣.
很早就結婚.
但她同時都很早做了寡婦.
但她沒有像一般女士一樣守寡.
很快又結婚.
可惜她再婚的那個男人對她很不好.
整天說她腿不好.
說她嚇死自己前夫.
還用很多不好聽的說話.
最後還用一些很小的事來幽了她.
在我們生活的那個年代.
幽妻其實很容易.
寫封信就可以了.
可以用很小的事來幽了老婆.
之後我想她都試過幾次的婚姻.
但我就不知道她結過多少次婚.
我想最後可能她對男人沒有什麼期望了.
於是現在就跟了一個肯包養她的有夫之婦.
叫做生活有些保障.
不用去做那一行.
知道我說什麼吧.
但她就成為了我們最佳城居文茶與飯後的娛樂.
成為了我們的話柄.
其實現在回想起來.
我們沒有憐憫她的不幸.
反而拿她出來說.
在她的傷口裡加鹽.
說真的有點不對.
不過這個最佳城居文對她的看法.
她自己都知道.
所以她總是選擇.
在正午最熱最沒人出現的時候.
去打水.

$^{161}$打水就每天都要.
大部分人就會黃昏日落.
沒那麼曬沒那麼熱.
在我們這些半沙漠地帶.
比較舒服的時候就出去打水.
所以你黃昏的時候.
就會見到很多人在井旁打水.
沒有人會在正午日頭最曬的時候去打水.
但正因為這樣.
她就選擇那個時候去打水.
不過那一天完全不同.
她沒有再避開人.
她跑回來找我們.
還很大聲地說.
「你們來看看 我遇到一個人」.
我們當時的反應是.
「什麼 你又搞了些什麼事出來」.
「跑成這樣 即是怎樣」.
但是她雙眼發光.
那份興奮.
那份自由.
那份毫無掩藏的勇氣.
我們從未見過.
起碼從未在她身上見過.
她以前是避開世界的.
但現在她衝著世界來做見證.
「朋友 我不知道你有沒有見過一個突然變得很自由」.
「自由到你不認得他」.
「原本你以為很熟悉」.
「但突然間覺得很陌生的人」.
那一刻我們就是這樣的感覺.
她氣喘喘地跟我們說.
「有一個人 他把我以前的事說出來」.
「但他不嫌棄我 他接納我 他理解我」.
「難道他就是基督?」.
你知不知道.
當她說我以前所做的事.
我們全城的人都知道她在說什麼.
而她口中的基督.
不單止猶太人.

$^{201}$我們這班撒瑪利亞人都很期盼上主.
要派來在我們這麼不堪的處境.
在現在我們這麼不堪的處境.
做我們的救主.
這方面這個女人就特別奇形成的.
不知道為什麼.
或者可能你都可以想像到.
她就在我們基里森山的聖殿那裡.
經常去敬拜 經常去供奉.
基里森山的聖殿沒有耶路撒冷那麼偉大.
規模那麼大.
但也同樣地沒有耶路撒冷的聖殿那麼多規矩.
所以她這樣的女人都可以去貢獻祭物.
所以她很樂意去基里森山那裡獻祭.
但她突然間說.
遇到基督.
有一個人她是基督.
還要將她以前所有的事都說出來.
但當她說的時候.
你發覺她再沒有羞愧.
不再閃避.
面容帶著一些光.
以前的她怕人說.
今天的她就好像打開了自己的生命那樣.
全部說給人知道.
我們就開始竊竊私語.
究竟發生什麼事.
究竟她遇到誰.
誰可以令那個人變成這樣.
我心裡的問題是.
她說遇到那個基督究竟是怎樣的.
究竟那個基督為何能夠令她那麼大改變.
於是我們一班村民就決定親自去看看.
到底她說的是什麼.
是真還是假.
我們其實不是完全相信她.
這個人那麼多問題.
突然間變得那麼熱情.
明天又不知道會變成怎樣.
但她那份被觸動的真實.

$^{241}$又確實是無法減去的.
於是我們就說.
不如我們自己去看看.
然後我們就見到她.
那個坐在井旁邊的耶穌.
祂看著我們.
並不是好像猶太人.
平時看著我們這班撒瑪利亞人那種厭惡的眼神.
而是用一種好像對著朋友.
好像很熟悉的目光看著我們.
我後來也有問那個女人.
我問她.
你見到一個猶太男人跟你說話.
你不害怕嗎.
她跟我說.
當她走到井邊發現那個男人的時候.
其實她是想立刻拿著水罐離開.
誰知那個男人突然問她.
可不可以給我一些水喝.
於是她才停下來.
一來出於好奇.
二來這個男人有求於她.
她被人拒絕得多.
所以就想幫她.
於是這個女人就問了她一句.
你既然是猶太人.
為何要向我這個撒瑪利亞婦人要水.
誰知我跟她說話.
就越說越出奇.
這個耶穌這個猶太人.
那個女人跟他說話的對話.
後來她跟我說就越來越離奇.
都不知說到哪裡去了.
一開始耶穌就跟她說.
可以給她活水.
她就以為這個人 耶穌在說的.
是不知道哪裡有個水泉.
那些水是會流的 所以是活水.
會流的水就是活水.
她以為她知道哪裡有個水泉.

$^{281}$於是她就跟他說.
如果耶穌告訴她.
哪裡有這樣的活泉.
她就不用那麼辛苦.
每天為了避開別人走那麼遠.
在正午的時候來這個井打水.
然後這個女人就問耶穌.
可不可以告訴我這個活水 活泉在哪裡.
但耶穌就突然叫她.
你叫你的丈夫過來.
於是那個婦人就跟我說.
她就打了個嗒.
她沒有細想之下就回了她.
我沒有丈夫.
誰知耶穌就跟她說.
你以前有五個丈夫.
現在跟你在一起的不是你的丈夫.
將這些事很仔細的說出來.
聽到這裡.
我就見到她是毫無掩飾自己的過去.
以前她會避開我們.
見到我們會感到很羞恥.
但今天她的羞恥變成了她的見證.
故事沒有變.
但說故事的人卻變了.
就好像她將水罐留在井旁.
她將過去的事.
過去的羞愧.
留在井旁.
留在耶穌身邊.
而她裡面就好像有個活水泉流出來.
你問我為何知道她將水罐留在井旁.
因為耶穌離開了罪加後.
她有一次來問我借水罐.
我問她水罐去了哪裡.
她說當日我將水罐留在那裡.
當時耶穌來到村 只顧著聽耶穌說話.
完全忘記了留下水罐.
我問她為何會留下水罐.
她說當初真的忘記了.

$^{321}$因為她跟耶穌說話後.
她的心很亂 很興奮.
突然間好像心一扇窗戶打開了.
有一陣風吹進來.
令到她的人完全忘記了.
但她離開了之後.
當她再找水罐的時候.
她就覺得這個好像一個記號.
這個記號提醒她.
她來井旁是要來裝水.
離開的時候.
卻是生命裡有耶穌.
有基督賜給她的活水.
以前她靠罐 靠自己背水回家.
現在她可以不再靠那些東西來解渴.
特別是她心靈裡的饑渴.
所以後來她跟我說.
忘記和釋放其實是一起發生的.
耶穌的釋放.
是發生在人最自然的反應裡.
你以為是忘記了.
其實可能是放下了.
但我問她.
你還問我借什麼水罐.
她說我還要洗衣服 煮飯 打水 做人.
不過我現在不是靠水罐來維生.
因為我是靠基督.
在我生命裡的活水來過每一天.
但我還要那個罐.
因為我生活還有很多責任.
活水不會取代我的責任.
我仍然要打水 煮飯 做人.
我會還那個罐給你的.
但我會繼續說我如何遇見基督.
她這樣說的時候.
我好像明白了一些事.
遇見基督不代表我們逃離我們的生活.
而是更新了我們的生活.
這個基督會賜活水.
使我們心裡有存緣.

$^{361}$我們心裡的饑渴可以得到滿足.
但水罐代表我們日常的責任.
兩者是並行的.
後來耶穌就在我們村裡住了幾天.
祂說了很多事.
不過祂有一句話令我不會忘記.
祂說「救恩是全世界的」.
祂曾經在哪裡說過「神愛世人」.
這個「世人」就是代表整個世界.
而救恩也是世界的.
不再分猶太人.
不再分撒瑪利亞人.
不再分男人 女人.
好人 壞人 正常人.
失敗的人 成功的人.
祂看穿了我們.
但沒有一句責備我們.
祂只有接納 只有邀請.
原來基督不是高高在上.
不是說在哪個山我們去見祂.
不是說在耶路撒冷我們接觸到祂.
不是說在基里森山我們接觸到祂.
而是祂親自來到我們的身邊.
在我們每天都要去的井旁.
和我們一起去感受疲倦 口渴.
祂明白那個女人的苦.
祂甚至問她要水喝.
祂用自己的需要 使我們放下自己的防衛.
使我們正視自己的需要.
我好像明白了些什麼.
這個女人為什麼改變得這麼快.
因為她見到的耶穌.
真是基督.
是我們每個人都需要的.
她那份自由.
她面上那份光彩.
不需要再掩飾的勇氣.
這些可能就是基督賜給她活水的果效.
耶穌說的活水不是只給她一個.
是給我們城裡每一個人.

$^{401}$給我們世界上每一個人.
我們城裡很多人都以為我們比這個女人好.
比這個女人乾淨.
比這個女人幸福.
比這個女人安定.
但其實反而是這個女人.
更加早嘗到真正的恩典.
我們雖然擁有禮儀傳統對神的知識.
但她擁有一樣東西.
她是真誠地承認自己的口渴.
真誠地面對自己的需要.
而耶穌最先回應的.
就是那些真心承認自己心靈乾渴的人.
起初我們因為她的見證.
因為這個女人所說的話.
我們去找耶穌.
後來我們親自見到耶穌.
親自聽耶穌說話.
我們只能說一句.
祂不單只是祂的救主.
更加是我們全城的救主.
是人類的救主.
是全世界的救主.
祂是基督無條件地接納我們.
釋放我們.
這個女人的生命被翻轉.
為的是令我們全城的人的生命都被翻轉.
她留下的水罐好像在提醒我們.
我們城裡每一個人都背著自己的水罐.
耶穌解去她的饑渴.
其實也解去了我們的饑渴.
我們明白到.
原來改變的不單只是她一個.
耶穌可以改變我們每一個人.
無論我們是什麼角色.
今天朋友聽完我這麼說之後.
你會覺得自己是什麼角色呢?.
你覺得你是旁觀者.
故事以外的人.
聽別人見證的人.

$^{441}$但我跟你說.
每一個人都可以在他生活當中.
在他每天去打水的井旁遇到基督.
他有時甚至會問你可不可以給他一杯水喝.
他就是基督.
他帶著他的需要來找你.
因為他知道你的需要.
他都經歷過你的需要.
那時候你會怎樣回應他呢?.
你會不會接待他.
讓他去揭露你的生命.
他揭露你的生命.
不是去審判你,批評你.
而是要你正視你的生命有什麼需要.
而他很直接告訴你.
我就是你所找的活水.
我就是你所找的基督.
我就是你所需要的就在你面前.
就像當日那個女人一樣.
她遇到基督.
又帶我們一起去遇到這個基督.
今日在你生活的井旁.
你有沒有遇到基督呢?.
或者你遇到的你又認不認得到他呢?.
我們一起祈禱吧.
主啊,你在我們生活當中與我們相遇.
在我們需要的裡面.
在我們饑渴的裡面.
在我們日常責任的裡面.
你就無聲無息地出現了.
我們很多時候會逃避.
我們很多時候會掩飾.
我們很多時候不願意面對我們的生命.
我們的過去.
但是你卻告訴我們不要緊.
你甚至問我們要水.
讓我們知道我們人是有需要的.
我們需要面對自己的困難.
自己的不足.
你又很願意賜下你的活水.

$^{481}$使我們生命可以豐盛.
求主你賜給我們這個活水.
讓我們心裡有從你而來的泉源.
我們每一天有力量地去過.
不代表我們不再需要打水.
它是代表我們打水的時候.
可以帶著自由.
帶著喜樂來做我們每天的日常.
又叫我們想想.
像那個女人那樣.
當她知道她遇到的是基督的時候.
她就將水罐留在井旁.
這樣告訴我們什麼呢?.
一個很微小的舉動.
讓我們看到她願意放下她的過去.
放下她的不安.
甚至可以說放下她的羞恥.
以致她能夠自由釋放.
因為在你裡是沒有分種族,階級,性別.
你願意叫每一個人.
心裡都有活水.
正如像罪加成的居民發言那樣.
你是世界的救主.
我們也享受與你的關係.
我們這樣禱告 是奉主耶穌明智祈求.
阿們.
\newpage



\section{以弗所書 4:25-32}
\label{sec:031766}
\textbf{建立健康的使命教會 (7) 只說造就人的話}~~~~陳燕華牧師
\newline
\newline
連結: \href{http://www.church.com.hk/acms/content.asp?site=cdc\&op=show\&type=product\&code=031766\&layout=sermon}{\texttt{031766}} ~~~~ 語音日期: 2026-03-08
\newline
\newline
\hyperref[sec:031765]{\small{< < < PREV SERMON < < <}}
~
\hyperref[sec:index_chronic]{\small{[返順時目]}}
~
\hyperref[sec:031767]{\small{> > > NEXT SERMON > > >}}
\newline
\newline
以弗所書 4:25-32
\newline
\begin{longtable}{cl}
\hline
\hline
章節 & 經文 (和合本修訂版)\\
\hline
4:25 & \begin{tabularx}{0.7\textwidth}{X} 所以，你們要棄絕謊言，每個人要與鄰舍說誠實話，因為我們是互為肢體。 \end{tabularx} \\ \\ \relax
4:26 & \begin{tabularx}{0.7\textwidth}{X} 即使生氣也不要犯罪；不可含怒到日落， \end{tabularx} \\ \\ \relax
4:27 & \begin{tabularx}{0.7\textwidth}{X} 不可給魔鬼留地步。 \end{tabularx} \\ \\ \relax
4:28 & \begin{tabularx}{0.7\textwidth}{X} 偷竊的，不要再偷；總要勤勞，親手做正當的事，這樣才可以把自己有的，分給有缺乏的人。 \end{tabularx} \\ \\ \relax
4:29 & \begin{tabularx}{0.7\textwidth}{X} 一句壞話也不可出口，只要隨著需要說造就人的好話，讓聽見的人得益處。 \end{tabularx} \\ \\ \relax
4:30 & \begin{tabularx}{0.7\textwidth}{X} 不要使神的聖靈擔憂，你們原是受了他的印記，等候得救贖的日子來到。 \end{tabularx} \\ \\ \relax
4:31 & \begin{tabularx}{0.7\textwidth}{X} 一切苦毒、憤怒、惱恨、嚷鬧、毀謗，和一切的惡毒都要從你們中間除掉。 \end{tabularx} \\ \\ \relax
4:32 & \begin{tabularx}{0.7\textwidth}{X} 要仁慈相待，存憐憫的心，彼此饒恕，正如神在基督裡饒恕了你們一樣。 \end{tabularx} \\ \\
[1ex]
\hline
\hline
\end{longtable}
$^{1}$弟兄姊妹 牧者平安.
繼續在以弗所書裡的經文.
由我們一起先得頭祈禱.
親愛的天父 你是全能的神 又是愛我們的天父.
配得我們的讚美敬拜.
多謝你眷顧我們 又眷顧我們的教會.
此刻我們祈求天父你挪開我們心裡面種種的雜念,煩惱.
親自對我們每個人說話.
讓你的話語成為我們去面對工作,生活和事奉你的力量.
我們這樣禱告奉靠主耶穌基督的名 阿們.
根據Gemini AI給我的資料.
一個標準男性成人的健康身體裡.
有大概30兆個細胞.
一兆是什麼概念呢.
就是一之後有12個零.
而同一個人體裡大概有39兆個細菌.
這個資料和我以前讀書時候認識的很不同.
因為當時是說人體裡的細菌比細胞的數量是幾倍的多.
於是我再去搜尋一下AI給我的資料的來源.
我發現原來我以前知道的東西.
已經在2014年因為在以色列的研究院裡所做的研究推翻了.
他們在完成這個研究之後.
發現其實一個標準男性成人的身體裡.
細菌的數量和細胞的數量比例其實是幾乎差不多的.
我為什麼要和大家分享這個資料呢.
因為我們從這個資料裡都可以理解到一件事.
一個健康的身體裡不是完全沒有細菌的.
同樣一間健康的教會也不是完全沒有問題的教會.
細菌在人體裡是有益有害的.
如果那些有益的細菌有什麼問題發揮不了功能.
那些平衡就會失去了.
那些不好的細菌就會影響我們健康.
但是要保持這個身體健康的方法不是說要將所有細菌清除的.
同樣我們在教會裡如果要保持教會健康.
不是說要將那些我們覺得有問題的人趕走.
也更加不是將我們覺得有問題的事掃入地氈底.
要身體健康人是要在日常生活裡去留意和實踐.
使我們身心健康的那些操練.
同樣教會裡我們要健康的話.
每個信徒我強調是每個信徒都要去學習和實踐.

$^{41}$使神的教導在日常生活裡去進行主導.
其中一種的學習和實踐就是我們今天經文裡所提醒我們的.
就是要留意我們的言語我們的說話.
為什麼呢?因為我們的言語都是帶著能力的.
我們千萬不要小看我們所說的話.
不是有一句話一言可興邦一言可喪邦嗎?.
我們不要否認我們自己的說話.
是可以帶出正面或者負面的能力.
所以我們的言語真的要有界線.
如果我們要發揮我們的言語裡面正面的能力的話.
我們只可以說有益有建設的說話.
這樣我們弟兄姊妹的相交就真的可以是有質素的相交.
而我們的教會就可以健康的成長.
今天經文一開始就說「所以」.
「所以」是一個連接詞來的.
是連接著上文的東西.
上文說什麼呢?.
上文就是保羅提醒教導的信徒.
如何在基督的真理基督的愛裡面去建立我們自己的屬靈生命.
以致我們令到耶穌基督的身體即是教會可以健康的成長.
以致我們可以做到神設立教會的目的.
什麼目的呢?就是要我們去將使人得救的福音傳揚.
這是我們的使命.
以致到最後我們能夠歸榮耀給神.
在第25節的經文裡面我們看到保羅提醒我們肢體之間我們說話.
我們要怎樣呢?氣絕謊言.
跟我們接觸跟我們往來的人我們之間說的話都要說誠實話.
在這裡所說的誠實話我們就要去留意了.
誠實話這個誠實跟我們在四章十五節裡面所說的用愛心說誠實話的誠實是不同的.
雖然在英文來講十五節的誠實跟這裡二十五節的誠實都是用the truth去翻譯的.
但其實原文是兩個不同的字.
不過雖然不同的字去到字根都是同一個字根.
究竟有什麼分別呢?十五節用的誠實是指真理.
而在二十五節所說的誠實是講真實.
它的相反詞是不虛假的.
當然不虛假真實的話一定包括真理.
究竟保羅為什麼這麼強調我們要講誠實話呢?.
因為首先我們不要忘記我們自己的身份.
我們是光明之子.
撒旦是說謊的人的父.

$^{81}$如果我們說謊不屬於神是屬於撒旦.
而我們的身份要很清楚知道我們是藉著耶穌基督的拯救.
以至我們能夠活在光明裡面的光明之子.
而聖經也提醒我們光明所結的果子是一切的良善公義和誠實.
是真實不虛假的.
還有我們不要忘記我們的言語真的有能力.
一句真心的話在適當的時候說出來.
就算你不是用來去讚人去鼓勵人.
裡面就算是包括一些指責提醒人的話.
只要是真心的.
只要是適當的時候去說的話.
人家都願意去聽去接受.
因為這個是對自己生命有益的.
反而無根無據的謊言.
你說得多漂亮都好.
你怎樣去讚人家都好.
如果對方是信以為真的話.
實在不會對他有益處.
甚至有機會害了他.
當我們願意去講誠實話的時候.
有時他那個正面的功效.
真的超過我們自己可以去想像的.
我記得我幾年前在玻利維亞認識一個弟兄.
他是在國內溫州來的.
我在去年大約四月去阿根廷參加一個宣教大會的時候.
在那裡再和他碰頭.
見到他和他的太太.
那個時候他告訴我.
他們剛剛因為生意的問題.
從玻利維亞搬去智利那裡做生意.
在這個宣教大會期間.
講完那個訊息.
那個弟兄就很受感動.
於是他也回應了.
願意奉獻自己全時間的禮福事實.
在大會裡見到我的時候.
他就來問我.
也告訴我問我意見.
我當時很自然的反應就是.
有弟兄肯奉獻自己.

$^{121}$去全時間奉獻神.
當然是鼓勵.
我說了很多鼓勵的話.
也答應我一定會為你祈禱支持你的.
在那天大會結束之後.
我就回到酒店.
答應別人為他祈禱.
那天晚上我就坐下來.
安靜的為弟兄祈禱.
在祈禱的時候.
總是想起我早上跟他說的鼓勵的話.
好像缺少了一些東西.
但我又想不到缺少了什麼.
總之就是有這樣的感覺.
第二天早上他繼續去大會的時候.
我見到他.
又再留意到他的太太在他身邊.
於是當時腦袋就想.
我知道我要跟這個弟兄說什麼了.
我就等到下午休息的時間.
我就鼓起勇氣.
就去跟這個弟兄聊天.
我就跟他說.
首先我也跟他說.
不過我跟這個弟兄說.
我發覺神早就給了你一個恩賜.
就是你的太太.
我鼓勵你跟太太一起奉獻去服侍神.
兩個一起去做神學的裝備.
我深信你的太太是神預備給你的.
成為你很好的幫手和陪伴.
接著我就繼續想說一些東西.
不過真的要鼓起勇氣.
心裡面也不停的去禱告.
我就直接告訴這個弟兄.
他在未來的裝備.
未來的事奉裡面.
他要留意些什麼.
有些什麼我相信他是需要幫的.
他要找他的太太幫他.

$^{161}$是非常直言的去說出一些東西.
心裡面是不斷的禱告.
為什麼呢?.
因為我很怕他不高興.
我更加不想因為我說的話純粹是出於我自己的意念.
我對他的認識而影響了他做決定.
他在大會結束之後.
他回到智利.
他馬上就報讀神學.
而且他還邀請我做他推薦的牧師.
他在大概五月尾左右就開始讀神學了.
是跟他的太太一起讀.
他去到去年年尾的期間.
這麼多個月.
他們要考試.
有時候考了不好又要重考.
每一個難關.
他都會發訊息來叫我跟他一起祈禱.
我自己的感覺就是.
弟兄很勤力去讀.
而且比我想像比我預計的更加勤力去讀.
他去到去年年尾.
在他回到智利的華人教會裡.
已經有實習講道機會一次.
而在今年在玻利維亞的華人教會裡.
都預留了一些時間.
讓他在那裡實習講道.
為什麼我要這樣好好的去祈禱.
好好的鼓起勇氣來跟他講呢.
我深信如果我沒有好好的為他祈禱.
好好的安靜去想一下.
這個弟兄他有什麼恩賜.
他有什麼的屬靈需要.
我是要去提醒他的話.
而是以一個牧者的身份好習慣性.
好自然的去鼓勵一個弟兄去參加.
接受神的訓練.
然後去全時間服侍神的話.
我相信如果我沒有做這些事.
我所講出來的話.

$^{201}$效果可能是很不一樣.
弟兄姊妹我們講誠實話.
實在是需要勇氣.
但我們更加需要去禱告.
因為禱告幫助我們讓神掌管我們內心裡面.
有什麼動機呢.
我們內心裡面有沒有什麼偏見呢.
也幫助我們不會輕率講一些我們自己覺得對的話.
能夠去禱告才去講誠實的話.
也幫助我們讓神去預備對方的內心.
去感受我們從心底發出的那份真誠和那份愛.
經文26到28節提醒我們.
作為光明之子.
無論在言語和行為.
我們都要有清楚的界線.
就是不可犯罪.
犯罪就是指得罪神.
這個可以說是我們一切心思意念.
言行舉止的底線.
我們不可以越過.
約瑟曾經對著引誘他的波提佛的老婆.
去講的一句話.
他說我怎能作者大惡得罪神呢.
這一句話正正就是在任何環境.
任何壓迫威脅試探之下.
都不可以被我們的心思意念.
言行舉止越過的一條底線.
我們的言語在什麼時候會越過這一條的界線.
然後去犯罪呢.
其實很多東西都可以影響我們.
愛在的人和事可以影響我們.
愛在的環境可以影響我們.
人家一兩句說話都可以煽動我們.
事實上歷史都讓我們看到.
人類都是很容易被說話煽動的動物.
不過最難抵抗的.
就是從我們內心的情緒出來的那種影響.
特別是經過長久壓抑的情緒.
我們可以去到如果不能夠將它再壓抑的時候.
它一爆發可以一發不可收拾.

$^{241}$令我們完全忘記了什麼叫做底線.
什麼叫做界線.
想想一對夫婦如果他要吵架.
兩個人都要講最後那句話.
才可以將裡面的憤憤不平發洩出來的話.
你說他們會怎樣呢.
從早吵到晚,從晚吵到天亮.
這個情況就好像箴言所說的.
火缺了柴就必熄滅.
就是說如果在兩個人裡面.
有一個很願意不要被自己的情緒支配而收口的話.
那就沒有柴燒了.
那些爭端,吵吵那些東西全部都可以慢慢停止.
阿葛他就曾經教導我們.
他說每個人都是強調每一個人.
要快快的聽,慢慢的講,慢慢的動腦.
他甚至解釋了為什麼要這樣.
他說因為人的腦氣並不能夠實現神的義.
人的腦氣都是血氣的情緒.
如果我們拿這些腦氣,拿我們血氣的情緒.
要去做神的事.
神都不會讓這些事成就.
更何況我們想用腦氣來去處理,去解決其他的事.
怎會成功呢.
所以簡單一句.
憑著人的腦氣可以說是一事無成.
當生氣的時候,又怎麼可以持守著街線,不去犯罪呢.
保羅就提醒我們.
不可含怒到日落,不可被魔鬼留地步.
猶太人的文化裡面,他們每一天是從日落開始的.
所以當保羅說不要含怒到日落.
意思就是我們千萬不要將我們的腦氣帶到新的一天.
所以在我們來說,真的不要讓腦氣留在我們心裡面過夜.
要快快的處理,不要讓它積壓下來.
腦氣藏在心裡面越久.
我們越有機會將這些腦氣積壓成心裡面的苦毒,憤怒和惱恨.
而這些東西就是魔鬼最喜歡用來控制人的武器.
這些就好像在《雅各書》裡面告訴我們.
是從地獄來的火在我們的心裡面燃燒.
不斷的燃燒,燃燒到有一天我們受不了.

$^{281}$就會從我們的口中出來.
那些就是讓鬧,毀謗,一切惡毒的言語.
英文有一句話其實是很貼切的.
它就說「Nip in the bud」,意思就是提醒我們.
對一些你早知有不好影響,不好的結果的事情.
要趁著它剛剛萌芽的時候就要壓殺了它.
不要讓它有機會.
所以盼望弟兄姊妹我們的內心在我們還沒有將腦氣積壓出來.
我們的口還沒有爆發那些惡毒的言語之前.
我們會快快去處理我們的腦氣.
那怎麼去處理呢?.
我們可以找人傾訴,我們可以找人跟我們一起祈禱.
甚至有些事情可能不可以跟別人說.
我們就在神面前去傾訴.
讓神去幫助我們,讓神聽完之後來安慰我們.
甚至給我們醫治悲傷.
我們也可以學大衛.
如果回到《詩篇》109篇裡面.
我們看到大衛將心底裡最惡毒的東西.
都在神面前表露無遺,完全傾倒出來.
《詩篇》109篇第一到第五節.
是大衛向神申訴他的冤屈,他的苦況.
經文裡面大概就是告訴我們.
他說那些惡人詭詐的人張開嘴巴攻擊他.
用謊言,用怨恨的說話來攻打他.
以惡報善,以恨報愛.
接著去到第八到十五節的經文.
大衛就去揍他的敵人.
要將他對敵人內心積壓的苦毒憤怒和惱恨.
全部發洩出來.
他想他的敵人是怎樣呢?.
年歲短小,子女要做孤兒,老婆要做寡婦.
子女到處行乞,最好就是後繼無人.
他的名字要在這個世界完全消滅.
所以簡單一句,大衛對他的敵人的咒詛是什麼呢?.
短命,妻離子散,絕子絕孫.
是不是很嚴厲,是不是很惡毒呢?.
但是我們看回第20節的經文怎麼說.
他說這就是那些與我作對.
用惡言議論我的人從耶和華所受的報應.

$^{321}$沒錯,大衛作出很嚴厲,很惡毒的咒詛.
但是他同時願意將他個人的報復.
公義的審判,懲罰的主權.
馴服的交到神的手裡.
憑著我們認識的神.
我們知道祂是慈愛.
我們知道祂是體恤我們每一個做祂兒女的怨弱.
我們知道祂是很願意去聆聽.
我們在祂面前傾出心裡最隱藏的惡毒.
最隱藏的怨鬱.
但是我們要學大衛這樣.
願意馴服的交在神手裡.
我們個人報復的權利.
將公義的審判,懲罰的主權.
完完全全交在神的手裡.
雖然我們的內心可能真是極度憤怒.
極度冤屈,卻不犯罪.
經文一直到第29節就告訴我們.
一句壞話也不可出口.
只要順著隨著需要講造就人的好話.
讓聽見的人得到益處.
壞話原文的意思是指腐爛.
就好像食物由好變腐爛.
說話又怎樣變腐爛呢?.
這裡的意思是指有一些說話.
真的可以讓生命由好變爛.
這些包括帶著惡意的說話.
而好話是指充滿善意的說話.
可以令生命越來越變得好.
我們廣東人很容易明白.
在這裡說壞話說好話.
我們廣東人很容易明白.
凡是說好話.
我們的心當然想為一些人或一些事情有幫助.
我們就多說幾句好話.
凡是說壞話.
心裡也可能想要重傷一些人.
又或者要破壞一些事情.
所以保羅的提醒是要說好話.
不要說壞話.

$^{361}$他甚至提到一句也不可出口.
一句都嫌多.
為什麼呢?.
因為他告訴我們.
一句壞話會使到聖靈擔憂.
這裡所說的擔憂不是指憂慮.
是傷痛,悲傷.
還是很沉重的那種傷痛,悲傷.
我有時在想.
當我們信主之後.
其實我們也有靈性軟弱的時候.
我們同樣會做一些神不喜歡,得罪神的事.
但聖經沒有用這種令聖靈擔憂來形容我們所犯的事.
對聖靈的影響沒有這麼嚴重,沉重的傷痛.
為什麼一句壞話有這麼大的能力呢?.
如果我們回到二弗所書第四章.
第一到第三節的經文裡看到.
保羅提出.
聖靈是賜給教會有合一的.
信徒在教會裡的責任就是要去努力保守這個合一.
因為只有合一的教會.
才有條件說我們要健康的成長.
才有資格說我們要完成神給我們的使命.
所以合一對教會有多重要呢?.
怎樣去合一呢?.
第一到第三節經文有說.
上個星期宗傳道的那篇道理.
你們要再聽聽.
我不會重複了.
怎樣在基督裡面成為合一的教會.
我們回頭看.
為什麼一句壞話.
可以有這麼大的能力.
使聖靈悲傷呢?.
因為我們都是主耶穌基督用的寶血去拯救回來的.
我們每一個人的身份是有聖靈的印記.
做憑據的.
不是就說我有這個身份.
我們是有那個身份證的.
所以我們在等待主基督再回來的日子.

$^{401}$應該在基督的愛裡面一起去建立.
建立我們自己的屬靈生命.
一起建立教會.
以致它能夠健康成長.
但是當我們不能夠節制我們的嘴唇.
我們在教會裡面.
在弟兄姊妹之間所說的壞話.
就會破壞合一.
聖靈怎會不會傷痛呢?.
一間不再合一的教會會是怎麼樣的呢?.
我在去年2025年.
一直都在網上幫一個在南美洲的姊妹做輔導.
這個姊妹我叫她做A姊妹.
最主要因為她跟另外一個姊妹.
我叫她做B姊妹.
她是有一點誤會.
而B姊妹就在教會的群組裡面拉黑了她.
大家都明白拉黑是什麼意思.
當時的牧者就要調停了.
所以她也在群組裡回應了一些真理的教導.
和一些公道的說話.
就是因為這樣.
牧者也被B姊妹拉黑了.
這個教會有兩個群組.
一個群組就單單是教會弟兄姊妹的群組.
他們用來每天靈修的時候大家分享訊息.
我也在這個群組裡面.
另外一個群組除了弟兄姊妹之外.
還有多年來一直加入的慕道的朋友.
或者是親戚,家人.
所以這是一個更大的群組.
很不幸的是B姊妹就用了第二個群組去拉黑大家.
所以事情就被很多人知道了.
看到的就是神家裡面的事.
我就這樣形容.
真是家醜外弱.
知不知道在這個群組裡面.
可以叛倒多少生命.
可以拆毀多少生命呢.
這一張就是不合一的教會.

$^{441}$其中的一些照片.
因為他們的事情已經被教會所有弟兄姊妹知道了.
所以在教會裡面的弟兄姊妹各有各的想法.
不同的說法.
所以都影響了弟兄姊妹的關係.
更加影響了他們對牧者的信任.
於是教會出現另外一張照片.
分裂在裡面.
過去的一年裡面.
教會用了整整一年的時間來修補弟兄姊妹的關係.
希望可以重建對牧者的信任.
其實我自己在旁邊和他們同行的日子裡.
我是很傷心的.
因為我看到的就是教會將很多的精神,精力.
奈何在裡面怎樣去重建,怎樣去重修.
又有另外一張照片出來了.
他們本來有三福傳福音的訓練.
到了一半就要停止了.
而他們這個姐妹小組裡面.
他們有一個探訪.
他們關懷社區的事工都要停止了.
這一張又是不再合一的教會有的照片.
試想想,當聖靈看到這些照片.
他怎麼會不擔憂,不傷痛呢?.
除了在我們中間要除掉我們那些.
令我們說壞話的壞情緒.
我們更加要懂得.
隨著需要去說造就人的好話.
《馬太福音》裡面有一節經文提醒我們.
因為心裡面所充滿的.
是會從口裡出來.
我們裡面存在的是善良的東西.
我們出來的就會是善.
我們存在的是一些惡念的話.
我們出來的東西自然就是惡.
所以我們內心裡面究竟充滿著什麼呢?.
可以幫我們在說話的時候的動機和心態.
都是正確,都是善良.
以致真的可以說到好話去見到別人呢?.
今天經文最後一句說到.

$^{481}$要仁慈相待全憐憫的心彼此饒恕.
正如神在基督裡面饒恕了你們一樣.
這句話是保羅給我們的.
要我們就是用這句話.
來保持在弟兄姊妹之間的那個相交.
令到我們的相交是有質素.
是互相建立的相交.
我記得有一次我的大兒子載我去一個地方.
他開著開著車.
前面那輛車的司機開得很慢.
他就跟得很貼.
很不開心又不是很耐煩.
他就一直按著按鈕.
想前面那個開快一點.
我看到的時候我就跟兒子說.
兒子,你當前面那個是你媽媽開車.
他反應慢,你不要催他.
這樣就真的有效.
我兒子真的好像沒有那麼不耐煩.
放慢了,又不要跟得那麼貼.
嚇到前面.
我就在想.
聖經裡面叫我們要彼此的饒恕.
正如神在基督裡面饒恕我們一樣.
我們是什麼人呢?.
我們那麼有限.
我們有沒有那種能力.
好像神去寬恕罪人那種的心.
來寬恕其他人呢?.
我深信我們很難做到.
我們很難做到這一句.
正如神在基督裡面饒恕了你們一樣.
不過,我們是可以凡事.
嘗試將自己放在別人的位置.
去理解一下他們的情況.
理解一下他們的角度.
為什麼他們會這樣做.
為什麼他們會有這樣的反應.
當我們願意這樣將自己放在別人的位置的時候.
我們至少會多了一些同理心.

$^{521}$而這個同理心.
真的可以讓我們能夠從心底裡去發出.
對人好一點.
對人體恤一點.
對人寬宏大量一點.
內心有這樣的動機和心態.
我深信我們出來的說話.
一定會是溫柔了一點.
謙卑了一點.
也都會是詛咒別人的.
弟兄姊妹還記不記得.
五合一的教會有的照片是怎麼樣.
我們很希望因全家呈現給聖靈看到的照片.
是教會合一.
是彼此相愛,是彼此詛咒.
讓聖靈滿足,讓聖靈安慰的照片.
要做到這樣的話.
我們要真的在我們當中.
一句壞話都不要說了.
我們只是按著適當的時候.
按著需要說一些詛咒別人的說話.
除掉我們一切心裡面的苦毒,惡行.
那些令我們有壞情緒的東西.
不要讓它們馬上在我們心裡面.
成為了我們出來的惡毒言語.
怎麼除掉呢?還記得嗎?.
找人傾訴,向神傾訴.
或者學大衛那樣.
你有多狠毒的話,你向神說.
不過,要將所有的主權交給神.
神要怎麼報應他.
是神的工作,是神的事.
還有就是我們盡量凡事.
站在別人的位置去理解別人的角度.
讓我們多一些同理心.
我們對人就會寬宏大量一些.
口說出來的都是好的說話.
一間健康的教會不是完全沒有問題的教會.
所以當我們看到教會裡面出現問題的時候.
我們不需要覺得失望,不需要覺得絕望.

$^{561}$因為我們知道這些的存在是正常的.
但我們一定要去保持教會健康.
每一個人都有一份的本份去做.
學習和遵從神的教導.
日常生活裡面去將神的教導活出來.
留意我們之間的相交,這是很重要的.
因為我們的言語真是帶著能力的.
不可以輕看我們的說話.
要用勇氣,要用禱告去講誠實話.
我們的言語都有界線,無論什麼環境.
我們都不要越過犯罪這個界線.
我們的言語如果是很想去發揮正面的能力的話.
就記住,只可以講建造人生命的說話.
讓別人聽到的時候得益處.
我們的教會如果肯這樣去遵守的話.
我深信我們可以保持著去成長.
我們才有條件,有資格.
去說健康的教會去成就神給我們的使命.
求神幫助我們,讓我們真的能夠又健康.
又可以去承擔神的使命去完成.
我留下兩個問題給大家.
希望你們在這個星期裡的靈修.
會好好在神面前去點試一下.
我究竟在什麼時候最容易犯上壞話的時候呢?.
我怎樣去節制自己不要說壞話呢?.
我怎樣去操練對別人更有同理心.
讓我說出來的說話更加柔和謙卑.
更加可以去建造別人的生命呢?.
\newpage



\section{詩篇 100:1-5}
\label{sec:031767}
\textbf{朝聖路上——歡欣敬拜}~~~~林廣潤牧師
\newline
\newline
連結: \href{http://www.church.com.hk/acms/content.asp?site=cdc\&op=show\&type=product\&code=031767\&layout=sermon}{\texttt{031767}} ~~~~ 語音日期: 2026-03-08
\newline
\newline
\hyperref[sec:031766]{\small{< < < PREV SERMON < < <}}
~
\hyperref[sec:index_chronic]{\small{[返順時目]}}
~
\hyperref[sec:031768]{\small{> > > NEXT SERMON > > >}}
\newline
\newline
詩篇 100:1-5
\newline
\begin{longtable}{cl}
\hline
\hline
章節 & 經文 (和合本修訂版)\\
\hline
100:1 & \begin{tabularx}{0.7\textwidth}{X} 普天下當向耶和華歡呼！ \end{tabularx} \\ \\ \relax
100:2 & \begin{tabularx}{0.7\textwidth}{X} 當樂意事奉耶和華，當歡唱來到他面前！ \end{tabularx} \\ \\ \relax
100:3 & \begin{tabularx}{0.7\textwidth}{X} 當認識耶和華是神！我們是他造的，也是屬他的；我們是他的民，是他草場的羊。 \end{tabularx} \\ \\ \relax
100:4 & \begin{tabularx}{0.7\textwidth}{X} 當稱謝進入他的門，當讚美進入他的院。當感謝他，稱頌他的名！ \end{tabularx} \\ \\ \relax
100:5 & \begin{tabularx}{0.7\textwidth}{X} 因為耶和華本為善；他的慈愛存到永遠，他的信實直到萬代。 \end{tabularx} \\ \\
[1ex]
\hline
\hline
\end{longtable}
$^{1}$多謝宋健牧師給我機會來到雪村跟大家分享.
當我想起要在雪村分享的時候.
腦海裡就出現了一段經文.
就是剛才宋健牧師讀出的經文.
詩篇第100篇.
詩篇第100篇是我自己回想.
2018年當我去到以色列的時候.
2018年這裡有沒有人去過以色列?.
沒有人去過以色列嗎?.
哇 也有人去過.
去到以色列的時候.
當時是我第二次去以色列.
是20年前去.
2018年.
就是2008年我第一次去.
2018年去到以色列.
跟沙浸,跟堅浸,跟天樂的弟兄姊妹一起去的時候.
特別是去到哭場的面前.
我突然間有一個感動.
要跟這群弟兄姊妹讀出詩篇第100篇.
詩篇第100篇裡面.
剛才宋健牧師也讀過了.
短短五節的經文.
我給他的題目就是《朝聖路上歡欣敬拜》.
當你熟悉詩篇的經文.
詩篇第93篇至100篇.
都是在說耶和華是王.
耶和華是王在當中敬拜,讚美神的時候.
我們怎樣來到聖經裡面說.
我們用心靈和誠實去敬拜神.
怎樣你能夠來到神的面前.
你很歡欣去敬拜呢?.
我不知道今早當你一路唱這三首敬拜的詩歌的時候.
你有沒有覺得從心裡面很想去讚美神呢?.
神的偉大,神的聖潔,神的公義,神的全能,神的同在.
令到你心裡面很想高聲大聲的去讚美神呢?.
詩篇第100篇一開始第一節的時候.
聖經裡面就提醒我們.
我們要普天下當向耶和華歡呼.
我不知道你有幾時曾經向神歡呼呢?.

$^{41}$過去幾天我的兒子帶著我們去塔斯曼尼亞.
有沒有人去過塔斯曼尼亞的?.
都有啊 不是那麼多人去過.
住在香港去到塔斯曼尼亞.
看到很大的草原.
我們在車裡面就哇.
看到很多的牛群羊群.
我們就哇.
我不知道你們什麼時候從心裡面哇一聲去讚美神呢?.
詩篇裡面告訴我們.
普天下當向耶和華歡呼.
當你收到心愛的禮物的時候.
你就哇一聲.
過去的星期五 今天是星期日.
過去星期五我們就和我的女兒慶祝生日.
原來這裡慶祝生日唱生日歌的時候.
那些工作人員點蠟燭 唱生日歌.
然後怎樣呢?.
舉手 說什麼呢?.
Harry 是不是啊?.
Ray 是不是啊?.
我不知道一句兩句三句.
是幫助慶祝生日的人.
當你拿到冠軍的時候.
當你大學畢業的時候.
你掟四方帽子的時候.
你有沒有覺得很歡欣很快樂呢?.
香港年輕人大學生.
他考到五星星的時候.
那種興高采烈.
拿到成績單的時候那種歡呼.
當你失業很久.
當你找到工作的時候.
可能你都很興奮很雀躍.
很歡欣.
聖經的學者告訴你和我.
詩篇第一百篇.
就是昔日聖殿建好的時候.
很多以色列人湧入敬拜.
去到聖殿的時候.

$^{81}$他們就唱詩篇第一百篇.
當昔日耶穌時代.
猶太人進入到聖殿的時候.
進入到會幕的時候.
他們就唱詩篇第一百篇.
短短五節的經文裡.
我們看到很重要的訊息.
聖經裡短短五節.
如果你細讀.
有七個的當字.
當向耶和華.
當讚美他.
如果中文譯本裡.
讓我們看到你們要.
無論是當或者要.
提醒今天我們每一個相信耶穌的人.
我們應該要這樣做.
我們要這樣做.
我們必須要這樣做.
人能夠甘心樂意.
真心實意.
去敬拜神 去讚美神.
是因為第三節的經文裡提醒我們.
你們當曉得耶和華是神.
當曉得原文的意思.
是認識 是承認 是相信 是接受.
當人經歷到這個信仰的真實的時候.
我們就能夠很歡欣的去敬拜神 讚美神.
當你聽到弟兄要轉會.
為什麼他認同這間教會呢?.
在這間教會裡他怎樣與神相遇呢?.
當他分享見證的時候.
為他感恩 是因為他與神相遇.
我相信這個相遇.
不單是頭腦的認識.
是他生命裡經歷到這個信仰裡的真實和具體.
可能你也認識舊約裡有一個很重要的人物.
叫做大衛.
大衛他面對巨人哥利亞的攻擊的時候.
聖經裡告訴你和我.

$^{121}$大衛很清楚說.
你攻擊我是靠刀槍 是靠外在的物件.
但是今天我去面對巨人哥利亞的時候.
他是靠著萬軍的耶和華.
因為他從小在牧場裡放羊的時候.
面對羊群 面對狼群攻擊的時候.
他靠的不是自己手上的武器.
他靠的是耶和華與他同在.
讓他能夠經歷到.
這個神是一個可信可靠的神.
舊約裡有很多的例子.
昔日的但爾里他放在獅子坑裡.
他毫無懼怕.
因為他相信耶和華上帝會保護他.
昔日的約書亞渡過約旦河.
他沒有錯 經摩西帶領他經歷過紅海的分開.
但是他要帶著以色列人渡過約旦河的時候.
完全不一樣.
是他吩咐抬著藥櫃的祭司.
走進這些河水的時候.
當腳踏進這個約旦河的河水.
河水才慢慢退卻.
約書亞仍然相信.
耶和華上帝必定會與他同在.
摩西,保羅和西拉在這個監牢裡.
他仍然沒有被四縛的圍牆困住.
他和西拉在監牢裡.
仍然很歡喜快樂地歌頌讚美神.
聖經裡的銷行傳給我們看到.
地大震動 所有的監門打開 鐵鏈都打開.
各位同學 今天你所相信的神是一個什麼神呢?.
詩篇第100篇裡說.
今天我們能夠歡欣的來敬拜.
是因為人能夠相信.
人能夠認識這位神的實在和寶貴.
今天你所相信的神是一個什麼神呢?.
你所相信的神 你認不認識祂呢?.
我不知道你家裡有多少本.
不同版本不同譯本的聖經.
還是剛剛過年的時候.

$^{161}$年廿八 中國人說會洗邋遢.
會清掃一下塵埃.
你有沒有試過每天都打開一本聖經去研讀呢?.
你口裡稱呼這個神.
但是你試問一下 每天24小時裡.
你有多少與神同在呢?.
你的工作 你的家庭 你的教養子女.
你的開車.
你有沒有經歷到神的真實呢?.
還是你好像粵語長片裡有一句話.
我不知道你有沒有聽過 有沒有看過.
有事就終無厭 無事就夏迎春.
當你面對困難的時候.
你就會切切的禱告期望 抓住神.
當困難過了的時候.
你就會過回你自己人生裡的目標和方向.
所以詩篇第一百篇裡.
第一樣東西提醒我們.
今天我們如果要敬拜神 要讚美神的時候.
我們要認識我們自己和神的關係.
無怪耶穌在《瑪達福音》第七章裡說.
凡稱呼我主 我主的人都不能進天國.
唯有遵行祂的話語的人.
我們才能與神同在.
弟兄姊妹 今天你認不認識你所相信的神呢?.
詩篇第一百篇的經文裡.
就提醒今天的你和我.
今天我們都是神所 都是祂所做.
都是屬於祂 都是祂的民 都是祂草堂的羊.
意思是什麼呢?.
意思是今天我們每一個相信耶穌的人.
我們都是神所創造 是神所揀選.
是神所牧養的人.
你有沒有被神揀選?.
你有沒有被神每天去牧養你呢?.
宗教改革家 神學家嘉義文說了一句話.
如果今天我們不認識自己.
我們就沒有辦法去認識神.
所以聖經裡 哥林多後書第五章裡.
一段很熟悉的經文.

$^{201}$若有人在基督裡 他就是新造的人.
舊事已過 都變成新的了.
今天的你和昨天的你有沒有改變呢?.
你的生活模式 你的習慣.
你和神的關係 你和人的關係有沒有更新呢?.
當我預備訊息的時候.
我回望我自己人生的路.
我找了三張我人生裡不同階段的相.
我不知道你有沒有試過.
找一下你人生裡不同的相.
我就找了三張.
由我很小的時候出生.
到我人生年輕的時候 到今天.
我就覺得我只能說一句話.
全部都是神的恩典.
你要說神的恩典 數之不盡 多而又多.
我從來都沒想過我成為傳道人.
我從來都沒想過今天可以對著這麼多人去講道.
沙浸以前的教會.
由一百多人 二百人.
開始去到我離開的時候.
十年我離開的時候是二千八百人.
當我在堅浸的時候.
我去的時候是一千四百人.
當我離開的時候差不多接近三千人.
我就覺得神給我很多的學習.
很多的突破 很多的經歷.
在信仰的路裡.
神為我預備了一位屬靈裡啟蒙的老師.
這位老師在我的教會裡.
成為我生命裡一個很重要的屬靈啟蒙的人.
當這位老師我入學的時候.
他第一次見我.
見了我一個小時.
見了我一個小時之後.
他給了一個評語.
他說林廣潤.
過去一個小時我好辛苦啊.
哇 為什麼這麼辛苦啊.
我說我問一句你回答一句.

$^{241}$你叫什麼名字啊.
林廣潤.
結了婚沒有啊.
結了婚.
你結了婚多久啊.
幾年.
你有沒有兒子有沒有女兒啊.
一個女兒.
你女兒叫什麼名字啊.
榮欣.
哇 如果你跟一個人聊天一個小時.
有問有答.
你覺得是不是令你很苦悶很抽筋呢.
所以他說如果入學之前他面試我.
他就一定不會收我這個神學生.
在過去三年多的神學裡.
他就成為我一個很重要的啟蒙老師.
幫助我 幫助我家裡.
幫助我整個家庭.
我就覺得這個都是神給我一個很大的恩典和祝福.
回望不單止信仰的路.
在我的職場的路裡.
神都很多的恩典.
我不知道今天你的職場在哪裡.
我有五年在學校裡工作.
有十年在當時的電話公司裡工作.
接著我就放下十幾年的工作.
進入到神學院裡接受裝備.
當我神學畢業的時候.
過去三十多年.
神帶我去了三間不同的浸信會去事奉.
很特別.
我離開了轉職又進入到這間教會.
離開轉職又進入到這間教會.
我人生三十多年就是事奉了三間教會.
但是每間教會我都離開又進入.
今年我又回到天樂浸信教會裡面.
做義務的牧師.
我想沒有一間教會.
俗語說會請一些吃回頭草的羊.

$^{281}$我就覺得這三間教會都樂意我回去幫忙.
我覺得這也是神給我在職場,牧會裡一個神的恩典.
神讓我經歷的.
就好像詩篇第23篇裡說.
祂帶領我去到青草地,溪水旁,可安歇的水邊.
面對困難,面對患難.
祂仍與我同在.
詩篇第100篇的作者裡.
開始的時候呼籲普天下的人.
都要來到神的面前去敬拜,去讚美.
是讓我們能夠認識我們自己.
又認識我們的神.
不單止認識我們人是神所創造.
第四節的經文裡說.
我們要稱謝進入祂的門.
讚美進入祂的院.
門是在說聖殿.
院是在說聖所.
在現代中文譯本裡的表達.
當人要進入到聖殿.
當人要進入到聖所.
當人要敬拜,讚美,稱頌神的時候.
我不知道我們應該存著什麼心態.
聖經的作者在這裡告訴你和我.
我們應該要有的心態.
是恭敬和謙卑.
當我去到以色列的時候.
去到一間教會.
大門關上了.
但旁邊就開了一個小的門.
每一個進入到聖殿去參觀的人.
一定要烏身鑽進去.
是表示進入到聖殿的時候.
那種謙卑,恭敬的心態.
你每個星期都回來說宣去崇拜.
當你回來的時候.
你有沒有渴望的心去敬拜神呢?.
你有沒有渴望今天.
神藉著詩歌,藉著禱告.
藉著神的話語和我說話呢?.

$^{321}$你有沒有願意來到神的面前.
經歷到神與你同在呢?.
當你很想經歷到神的同在的時候.
你有沒有認識.
我們的神究竟有什麼特性呢?.
詩篇第一百篇就提醒我們.
因為耶和華本為善.
善的意思就是好的意思.
如果你熟讀《創世記》的經文裡.
神創造六天的創造.
每一次的創造.
神都說是好的.
但去到創造人的時候.
神說甚好.
原來神對人的愛.
是超過今天你所想所求.
神對人的慈愛.
是我們想都沒想過.
想都沒想過.
我人生沒想過會來雪梨.
我來過一次.
昔日我兒子在埃德尼讀書.
我去過埃德尼兩次三次.
但雪梨我就來過一次.
之後就沒想過會來到.
但我沒想過可以來到雪梨.
我人生裡更加沒想過.
可以去到塔斯曼尼亞.
在香港經常說塔斯曼尼亞有多好啊.
說塔斯曼尼亞一定要吃生蠔.
說裡面的蠔養得很好.
我人生吃過蠔是很好吃的.
我不知道你人生裡.
有沒有想過一些你從沒想過.
從沒想像過.
從沒經歷過.
神的信實.
是因為我們是草場的羊.
是祂所創造.
神的信實.

$^{361}$是超過今天你和我所想所求.
因為祂的善是信實的.
是慈愛的.
信實是祂堅定不移.
祂會愛我們愛到底.
祂對我們的慈愛永遠不變.
神對你的愛有沒有變過呢?.
神的信實有沒有離開過你呢?.
我想在《舊約》裡面有一個很重要的人物.
這個人物叫做約瑟.
我不知道你有沒有留意《約瑟》.
《聖經》裡面有很長的篇幅.
去描述約瑟的生命.
他的經歷.
在約瑟自從他爸爸很年輕的時候.
送了一個彩衣給他之後.
他就面對著被他哥哥排斥.
被他兄弟出賣.
去到家裡圍爐.
他被這個主母誘惑.
結果陷害他坐牢.
在監牢裡他為人解夢.
在解夢裡他不被紀念.
在他人生裡跌落最低谷的時候.
他沒有忘記神與他同在.
所以如果你讀《創世記》完結第50章的時候.
第50章第20節裡面約瑟說了一句話.
你們的意思是要害我.
但是神的意念是好的.
是要成全今天的光景.
我不知道今天.
如果要你回想一下.
你信仰裡面的歷程的時候.
你有沒有覺得神的慈愛.
神的信實是何其的豐厚.
鼓勵你每次回來教會不是完了就離開.
當然你要開車15分鐘,20分鐘或者半小時.
但是當你完了崇拜的時候.
我知道下個星期就有一起用膳的時候.
你們一起享用食物,愛言的時候.

$^{401}$你也分享一下你信仰裡面的經歷,體會.
你過去一個星期,一個月裡面.
神有什麼恩典,祝福臨到你呢?.
不只是為了吃而留在這裡.
是為了彼此能夠在信仰裡面.
有分享,有經歷.
2018年當我去到以色列的時候.
很遠參觀很多的禮拜堂.
很遠的時候就聽到很多美妙的樂曲.
當我們帶著這二十多人.
一起進入到這個參觀的禮拜堂的時候.
原來知道我們可以自由的唱詩.
結果在那一刻裡面.
沒有準備之下,沒有練習之下.
我們就鼓勵所有的弟兄姊妹.
在香港是有沙浸,是有堅浸,是有天樂的弟兄姊妹.
他們過去都不認識的.
他們從來都沒有一起合著唱過歌.
但是去到以色列的時候.
去到這間禮拜堂.
我們就一起唱一首詩歌.
沒有講台,沒有喇叭.
沒有任何音響的設備.
但我不知道當你一直聽的時候.
有沒有聽到一些迴聲的呢?.
我覺得這也是神很特別的恩典.
從來都沒有合作過.
從來都沒有練習過.
但我們一起唱的時候.
我們就能夠將最高的榮耀.
給我們的神.
Sing Hallelujah to the Lord.
唱讚美我們的神.
在整個以色列的行程裡面.
神不單止讓我們一起學習.
神給我自己個人很特別的經歷.
剛才說2018年是我第二次去以色列.
最後的時候.
神竟然讓我遇到一個20年前.
成為我們帶領我們在以色列裡面.

$^{441}$一起去經歷,一起體會的老師.
他來自台灣.
但是他一直都留在以色列.
讓我遇到他.
他又認得我.
結果我們倆就拍了一張合照.
這張合照我覺得很特別.
張實賢老師20年仍然在以色列裡面.
做著很多帶領的工作.
和人分享整個以色列的行程.
神給我很特別的恩典和經歷.
人要認識神.
是因為神的慈愛,神的信實.
沒有改變過.
香港2024年裡面.
我的女兒和我看一套很出名的.
很修德的電影.
我不知道你有沒有看.
就是《九龍城寨遺城》.
當時對我來說.
是一個不單止回憶.
更加看見一個當時的社會.
當時的文化.
當時的處境.
是一個三不管的地方.
裡面的黑暗.
裡面居住的環境.
那種問題.
但是在很多困難裡面.
很多問題裡面.
神從來都沒有忘記過.
有九龍城寨這個三不管的地方.
很多人被忽略.
很多人不被愛.
但是神在很遠很遠的地方.
在英國倫敦.
感動一位20歲年輕的姐妹.
叫做潘寧卓來到香港.
當她進入城寨的時候.
在她的文章裡面分享.

$^{481}$一聞到這些味道.
她就知道.
這就是神要帶我來的地方.
一個在倫敦來的外國女子.
22歲.
她來到城寨.
將福音傳開.
去關心這些三不管的人.
吸毒的人.
戒毒的人.
沒有人關心.
沒有人愛的人.
潘寧卓帶著神的恩典.
在當中來建立這個群體.
將福音來帶給這個有需要的群體.
樹仁大學頒了一個榮譽博士學位給她.
表達她這種無私的愛.
在當中去建立這個群體.
無怪史托德穆斯.
在他的表達裡面.
他說今天我們的信仰.
不單止我們要讀好聖經.
你擁有很多本聖經都沒有用.
不單止我們熟讀聖經.
不單止我們鑽研聖經.
史托德穆斯說.
聖經是要我們遵從.
所以弟兄姊妹.
今天我們能夠開聲高歌的來頌揚.
是因為我們心裡面認識我們自己.
又認識我們所相信的神.
華福總幹事董家華牧師.
在他寫一本書.
叫做21世紀的指命門徒的時候.
他說今天每一個基督徒.
我們要和我們所相信的耶穌.
建立一個很親密的關係.
是一生之久的關係.
不單止我們和神建立關係.
我們更加要遵行耶穌給我們一個整全的使命.

$^{521}$詩篇第100篇.
大衛藉著短短五節的經文裡面.
提醒今天的你和我.
我們當向耶和華歡呼.
我們要敬畏讚美我們的耶和華.
當我們要懂得去歡呼去讚美的時候.
我們要認識我們自己.
我們要和神建立一個永恆的關係.
我們要認識我們神的特性.
所以在結語的時候.
我選了一節經文和大家讀.
不如我們一起讀出.
善治夏伯谷書第十三章十七,十八字.
我們一起讀,好嗎?請!.
雖然無花果樹不發旺.
葡萄樹不結果.
橄欖樹也不效力.
田地不出糧食.
圈中絕了羊.
棚內也沒有牛.
然而我要因耶和華歡恩.
因救我的上帝喜樂.
希望詩篇一百篇和夏伯谷書都放在你心裡.
我想你們都很熟悉.
但是希望這個信仰不只是頭腦的認知.
是我們生命裡面一個很寶貴的經歷和體會.
讓我們同心低頭來禱告.
還有當我們每天有機會打開你自己的話語的時候.
給我們一個謙卑受教的心.
不是我們熟悉當中每一個文字.
而是我們能夠願意期望打開你的話語的時候.
我們能夠聽到上帝向我們心靈裡面來說話.
因為你的話就是我們腳前的燈.
就是我們路上的光.
因為你的話就是我們每天生命裡的能量.
讓我們能夠在你的話語裡.
我能夠更加的認識我們自己.
又認識你.
幫助我們能夠成為一個聽到又行到的人.
讓我們的信仰不只是停留在頭腦的認知.

$^{561}$而是活在我們生命裡每時每刻.
在我們的家庭.
在我們的工作.
在我們的服侍.
在我們和人際關係的裡面.
我們能夠見到你與我們同在.
幫助我們.
讓我們的言語.
讓我們的行為.
讓我們的信心和愛心.
和我們所領受的救恩相稱.
祝福我們雪孫的每一位弟兄姊妹.
讓這個信仰成為我們的生命.
成為我們家庭裡一個很重要的角色.
幫助我們又祝福我們.
多謝你讓我們今天早上.
在你的裡面.
我們一起有敬拜.
一起有你的話語的分享.
聽我們的禱告.
奉耶穌基督你的聖名.
阿們.
\newpage



\section{約書亞記 2:1-24}
\label{sec:031768}
\textbf{抓著應許}~~~~吳銘淇傳道
\newline
\newline
連結: \href{http://www.church.com.hk/acms/content.asp?site=cdc\&op=show\&type=product\&code=031768\&layout=sermon}{\texttt{031768}} ~~~~ 語音日期: 2026-03-08
\newline
\newline
\hyperref[sec:031767]{\small{< < < PREV SERMON < < <}}
~
\hyperref[sec:index_chronic]{\small{[返順時目]}}
~
\hyperref[sec:031769]{\small{> > > NEXT SERMON > > >}}
\newline
\newline
約書亞記 2:1-24
\newline
\begin{longtable}{cl}
\hline
\hline
章節 & 經文 (和合本修訂版)\\
\hline
2:1 & \begin{tabularx}{0.7\textwidth}{X} 嫩的兒子約書亞從什亭暗中派兩個人作探子，說：「你們去窺探那地和耶利哥。」於是二人去了，來到一個名叫喇合的妓女家裡，在那裡睡覺。 \end{tabularx} \\ \\ \relax
2:2 & \begin{tabularx}{0.7\textwidth}{X} 有人告訴耶利哥王說：「看哪，今夜有以色列人到這裡來窺探此地。」 \end{tabularx} \\ \\ \relax
2:3 & \begin{tabularx}{0.7\textwidth}{X} 耶利哥王派人到喇合那裡， 說：「你要交出那來到你這裡、進了你家的人，因為他們來是要窺探全地。」 \end{tabularx} \\ \\ \relax
2:4 & \begin{tabularx}{0.7\textwidth}{X} 但女人已把二人藏起來，卻說：「那兩個人確實到我這裡來過，他們從哪裡來，我卻不知道。 \end{tabularx} \\ \\ \relax
2:5 & \begin{tabularx}{0.7\textwidth}{X} 天黑、要關城門的時候，他們就出去了。他們往哪裡去我也不知道。你們趕快去追他們，就必追上。」 \end{tabularx} \\ \\ \relax
2:6 & \begin{tabularx}{0.7\textwidth}{X} 其實，這女人已經領二人上了屋頂，把他們藏在她擺列在屋頂的亞麻梗中。 \end{tabularx} \\ \\ \relax
2:7 & \begin{tabularx}{0.7\textwidth}{X} 那些人就往約旦河的路上追趕他們，直到渡口。追趕他們的人一出去，城門就關了。 \end{tabularx} \\ \\ \relax
2:8 & \begin{tabularx}{0.7\textwidth}{X} 二人還沒有睡之前，女人就上屋頂，到他們那裡， \end{tabularx} \\ \\ \relax
2:9 & \begin{tabularx}{0.7\textwidth}{X} 對他們說：「我知道耶和華已經把這地賜給你們了，並且我們也都懼怕你們。這地所有的居民在你們面前都融化了。 \end{tabularx} \\ \\ \relax
2:10 & \begin{tabularx}{0.7\textwidth}{X} 因為我們聽見你們出埃及的時候，耶和華怎樣在你們前面使紅海的水乾了，並且你們怎樣處置約旦河東的兩個亞摩利王，西宏和噩，把他們完全消滅。 \end{tabularx} \\ \\ \relax
2:11 & \begin{tabularx}{0.7\textwidth}{X} 我們一聽見就膽戰心驚，人人因你們的緣故勇氣全失。耶和華－你們的神是天上地下的神。 \end{tabularx} \\ \\ \relax
2:12 & \begin{tabularx}{0.7\textwidth}{X} 現在我既然恩待你們，求你們指著耶和華向我起誓，你們也要恩待我的父家。請你們給我一個確實的憑據， \end{tabularx} \\ \\ \relax
2:13 & \begin{tabularx}{0.7\textwidth}{X} 要救活我的父母、兄弟、姊妹，和所有屬他們的，拯救我們的性命脫離死亡。」 \end{tabularx} \\ \\ \relax
2:14 & \begin{tabularx}{0.7\textwidth}{X} 那二人對她說：「我們願意以性命來替你們死。你們若不洩漏我們這件事，當耶和華將這地賜給我們的時候，我們必以慈愛和誠信待你。」 \end{tabularx} \\ \\ \relax
2:15 & \begin{tabularx}{0.7\textwidth}{X} 於是女人用繩子把二人從窗戶縋下去，因為她的屋子是在城牆邊上，她也住在城牆上。 \end{tabularx} \\ \\ \relax
2:16 & \begin{tabularx}{0.7\textwidth}{X} 她對他們說：「你們暫且往山上去，免得追趕的人遇見你們。要在那裡躲藏三天，等追趕的人回來，你們才可以走自己的路。」 \end{tabularx} \\ \\ \relax
2:17 & \begin{tabularx}{0.7\textwidth}{X} 二人對她說：「你叫我們所起的誓與我們無關， \end{tabularx} \\ \\ \relax
2:18 & \begin{tabularx}{0.7\textwidth}{X} 除非，看哪，當我們來到這地的時候，你把這條朱紅線繩子繫在縋我們下去的窗戶上，並要叫你的父母、兄弟和你父的全家都聚集在你家中。 \end{tabularx} \\ \\ \relax
2:19 & \begin{tabularx}{0.7\textwidth}{X} 凡離開你家門往街上去的，他的血必歸到自己頭上，與我們無關；凡在你家裡的，若有人下手害他，他的血就歸到我們頭上。 \end{tabularx} \\ \\ \relax
2:20 & \begin{tabularx}{0.7\textwidth}{X} 你若洩漏我們這件事，你叫我們所起的誓就與我們無關了。」 \end{tabularx} \\ \\ \relax
2:21 & \begin{tabularx}{0.7\textwidth}{X} 女人說：「就照你們的話吧！」於是她送他們走了，就把朱紅繩子繫在窗戶上。 \end{tabularx} \\ \\ \relax
2:22 & \begin{tabularx}{0.7\textwidth}{X} 二人離開，到山上去，在那裡停留三天，直等到追趕的人回去。追趕的人一路尋找，卻找不著。 \end{tabularx} \\ \\ \relax
2:23 & \begin{tabularx}{0.7\textwidth}{X} 二人回來，下了山，過了河，來到嫩的兒子約書亞那裡，向他報告他們所遭遇的一切事。 \end{tabularx} \\ \\ \relax
2:24 & \begin{tabularx}{0.7\textwidth}{X} 他們對約書亞說：「耶和華果然將那全地交在我們手中了，並且那地所有的居民在我們面前都融化了。」 \end{tabularx} \\ \\
[1ex]
\hline
\hline
\end{longtable}
$^{1}$弟兄姊妹早晨.
讓我們為著我們講道和聽道的服侍一同禱告.
因為是你的恩典去呼召我們.
願上帝今天藉著聖靈對我們說話.
藉著經文對我們說話.
奉著耶穌基督成名祈求 阿們.
兩個星期前剛剛開始講這個約書亞記.
今次不知道大家還記不記得上次大概在講什麼呢.
上一次我們講到關於恐懼.
在講恐懼很多種類的包括動物很多不同的東西.
而講到一些深層一點的恐懼.
對我來講應該是離別.
我很怕講拜拜.
因為有那種好像失去了的那種感覺.
所以就算來到雪梨已經生活了五年了.
每一次在機場跟家人說拜拜的時候.
都是眼濕濕都很不捨得.
都是會被醫毛制一段時間.
對於離別那種恐懼.
這個是我跟輔導員聊天的時候最近發現的.
就是源於我會說其實我很少一些.
就是以前的一些朋友.
中學的朋友我都沒有.
因為我讀完中學就不會再回中學的那些人.
很多同學其實都慢慢會開始沒有聯絡.
後來我才發現原來這個都是我面對恐懼的一個方法.
因為我很怕別離.
所以我會選擇逃跑.
我心裡面經常都會覺得反正大家升學之後各自生活.
其實就會慢慢覺得大家都很不同.
就會見面少了.
所以不如我快點不見.
這樣好像那個關係就會斷了就斷了.
我就會沒那麼傷心了.
原來人有不同的方法去面對恐懼.
經常都會有不同的反應.
心理學通常把它分為四類.
第一樣就是剛剛說的那種逃跑.
就是見到那個恐懼的時候就要離開那個恐懼的場景.
有些人不是的.

$^{41}$有些人選擇對抗.
就是直接對抗面對那個恐懼.
就是我們要繼續做朋友.
我們應該要繼續一起然後維持這段關係.
可能會是這樣.
甚至有些人會呆了.
反應不來.
就是腦袋一片空白.
什麼都想不到.
就是我都不知道說什麼.
很多東西想說但是不知道怎麼說.
同樣有些人會討好.
就是透過去取悅.
威脅那個來源.
來保護自己.
就是可能明明很多不捨得的.
但是永遠在嘴裡.
不要緊不用擔心我沒問題.
當人面對著恐懼的時候.
都會用不同的方法去保護自己.
可能有人逃跑.
有人對抗.
有人呆了.
有人會討好.
而很多時候.
這些反應其實沒有對和錯.
只是人面對情感的一種方法.
而回到我們的經文裡面.
上一次我們說到約書亞記的第一章.
就是說上帝怎麼回應約書亞的恐懼.
在約書亞要成為.
以色列人的領袖的時候.
有很多讓他可以擔心的事情.
包括可能是新的領袖角色.
打這些未知的仗.
又甚至是整個民族.
怎麼樣的團結.
這些事情全部都是約書亞所擔心.
但是上帝對約書亞說.
你要剛強壯膽.

$^{81}$因為前面要面對的問題.
其實這些是上帝的工作.
約書亞只需要記住一件事.
去忠於上帝的吩咐去做就可以了.
但是有時候我們說就是這樣說.
恐懼其實不是說放下就放下.
擔心就是擔心.
害怕就是害怕.
人有很多時候.
不是因為不知道真理.
然後去恐懼.
而是因為恐懼.
本身這件事.
有一個很真實的感覺.
所以問題就來了.
當恐懼真的出現的時候.
人除了逃跑.
除了對抗.
除了呆了站在那裡不動.
又或者可能是.
用一些方法去討好之外.
還有沒有另外一種的可能性呢.
約書亞記第二章.
就讓我們去看到.
拉合這個角色.
在恐懼裡面.
做了一個完全不同的選擇.
我們要明白拉合的選擇.
首先要知道拉合經歷了什麼.
需要讓我花一點時間.
講一下拉合這個故事.
剛剛主席幫我們讀到中間的部分.
但是前前後後在說什麼呢.
就是當時約書亞派了兩個探子.
去到耶利江城那裡.
去窺探一下這個地方是怎麼樣的.
探子就去到一個叫做.
拉合這個妓女的家裡面去停留.
大概就是因為妓女的家.
應該是很多人可以在那裡.

$^{121}$所以好像比較容易掩人耳目一點.
而耶利江王知道.
但是也知道了.
有以色列的探子去到進了城.
甚至他知道他去了拉合的那個家.
所以就派人要去拉合那裡去搜查.
拉合就將探子藏在屋頂的那些乾草裡面.
然後就對王的使者說.
我見過他們的.
不過他們走了.
王的使者就馬上追出去.
城門就關了.
在這個時候.
拉合其實已經知道.
耶和華已經要將這塊地賜給這幫以色列人.
所以拉合去救這些探子.
甚至請求這些探子.
就是說你可不可以拯救我和我的全家.
探子答應了他.
就是說只要拉合不洩露探子的事.
他就會保他一家的安全.
之後拉合就用繩子將探子吊出去.
然後探子再躲起來三天.
回去約書亞那裡.
而拉合也都按照吩咐.
將所有的家人聚集在自己屋子裡面.
也都將紅繩掛在窗外去等待救援.
這個就是拉合的故事.
我們今天要去討論兩件事.
第一件事就是拉合的恐懼.
就是他承擔著什麼樣的風險呢.
同樣第二個我們更加要問的問題就是.
拉合能夠克服這些恐懼的原因又是什麼呢.
我們最後第二個部分就是.
在約書亞記第二章的序事裡面.
我們其實很容易去看到拉合的那份信心.
這麼大的風險或這麼多的事情.
但他就是這樣去做了.
但如果你仔細去閱讀經文的時候.
你就會發現其實他的選擇都是極度危險.

$^{161}$甚至是風險很大.
首先第一個他要面對的是一個叛國的風險.
就是當耶利哥王知道以色列有探子進城的時候.
已經派人去到拉合的家裡去找的時候.
在古代我們會知道.
其實當你收留了敵國的探子.
其實這件事跟背叛已經沒什麼分別了.
所以一旦被揭發.
其實他拉合就會被視為是叛徒.
會被處決.
這個是他第一個風險.
他還要承擔全家人性命的風險.
拉合不是說你記得要救我.
救我就可以了.
他要求探子在以色列攻城的時候.
要保他全家全個家族.
這個意味著他必須要說服他的家人.
要將生命交付在這兩個探子的承諾當中.
如果探子失信沒有救他們的時候.
拉合不單止害了自己.
甚至他也害死了整個家族.
可能家人要聚集在家裡等救援的時候.
都會不期然地說.
你真的相信這兩個陌生人和他們的神?.
這個不是我們聽過或信過的一個上帝.
為什麼你會信得過他?.
同時他們也可能不知道.
以色列人究竟什麼時候才會來攻城.
什麼時候我們才會得到救援.
所以在經文的表達裡面.
我們相信他們大概要等兩個星期.
兩個星期.
兩個星期什麼地方都不敢去.
因為說離開了那間屋.
你就不會獲救了.
所以你一定要留在這間屋.
而本身這件事.
已經會令到城裡面很令人懷疑.
為什麼?.
為什麼整天都不出家呢?.

$^{201}$又不是過年.
然後一群姨媽姑姑在家裡做什麼呢?.
不單止這樣.
他還要承擔被暴露.
被人發現的風險.
探子吩咐他們還要做多一件事.
就是要將一條紅繩掛在窗外.
作為一個記號.
雖然這條紅繩是代表拯救的一個記號.
但同樣意味著他要公開表明.
他和以色列人站在同一條線.
但你可能會說.
Kevin沒那麼誇張吧.
為什麼掛一條紅繩.
就知道他站在哪一邊呢?.
原來繩這個希伯來文是Tequa.
Tequa這個字.
和盼望其實是一個.
很相似讀音的一個字.
所以我們沒有在迦南文化裡面.
看到有說要掛一條紅繩.
在外面的一個禮儀.
所以不其然都會令人想像.
就是為什麼要掛一條紅繩呢?.
紅繩和接下來攻擊我們的人.
有沒有什麼關係呢?.
所以其實這條繩.
都有一定的顯眼的程度.
算是拉闔的一個立場聲明.
也都令他成為城裡面.
一個很可疑的人.
但最重要的是什麼呢?.
最重要的是.
如果那些探子他們說得是.
他們一落千金的時候.
其實沒有這些風險.
可能你覺得都OK都值得.
因為真的可以救到全家.
但我們要知道是什麼呢?.
知道是上帝曾經吩咐過.

$^{241}$以色列人在申命記七章裡面.
去說你們要滅絕了迦南的七族.
沒有一個能夠留下.
所以其實儘管兩個探子.
去承諾任何的事情.
但對拉闔的一家來說.
其實他們能不能夠得救.
都不是這兩個探子說了就算.
最終都要看上帝的心意.
所以其實對拉闔來說.
探子的承諾.
其實可能都只是一張空頭的支票.
第二處.
如果你是拉闔.
拉闔是不是應該選擇一個.
更安全的一個選擇?.
他大可以討好耶利哥王.
將探子交出去.
甚至可能因為這樣.
他可以得到很多賞識.
去保到自己的安全.
但他沒有這樣去做.
他寧願冒著叛國.
連累家人.
甚至暴露自己的風險.
他都仍然選擇去保護探子.
這個就是我們要回答.
第二個問題的原因.
第二個問題是.
拉闔為什麼都願意承擔這麼大的風險?.
拉闔是怎樣能夠克服這些恐懼?.
剛剛主席幫我們讀了.
二章九到十三節.
要知道拉闔.
怎樣能夠克服這些恐懼的原因.
我們要看回這段經文.
拉闔之所以能夠承擔這些風險.
是因為他對上帝.
有一個很清楚的認知.
第九節開始.

$^{281}$他在那裡說給我們聽.
拉闔已經看清了上帝的心意.
拉闔第一句的宣告.
我知道耶和華已經將這塊地賜給你們了.
在希伯來文.
賜給這個字不是一個未來式.
就是可能會賜給你.
或者會賜給你.
所以我現在都是五十五十.
然後我買你大.
所以我就相信你.
不是一個這樣的環境.
賜給這個字是一個完成式.
所以對拉闔來說.
耶和華已經將這塊地賜給你們了.
他見到耶和華已經做了這件事.
這是第一個我們見到拉闔.
見到上帝的心意.
不單只是這樣.
拉闔雖然是一個外邦的女子.
卻已經看見上帝要完成的事.
這個是因為他見到.
上帝在歷史裡面的工作.
拉闔提到兩件很關鍵的事.
在第十節.
他說四十年前.
在紅海的神蹟.
和近期對阿摩利那兩個皇帝所做的擊殺行動.
所以他都明白原來發生了什麼事.
這兩件事分別是代表著.
以色列人在進入曠野生涯的開始和結束.
所以拉闔透過這兩件歷史事件.
看見上帝的能力 看見他是世界的王.
而上帝要帶領以色列人所做的事.
其實不單只是拉闔知道.
是整個迦南地都知道.
在九和十一節裡面我們會見到.
拉闔描述當時迦南人的狀態是怎樣.
其實他們很害怕.
因為他說它是驚恐到好像掉了在我們身上.

$^{321}$心是消化的.
這個是聖經裡面描述真戰的術語.
而同樣作者藉著拉闔的口.
說出上帝說的話應驗了.
因為摩西在出埃及記十五章的時候.
裡面都有一個這樣的預言.
就是他會提到迦南的居民心會消化.
現在拉闔告訴他們.
迦南的居民真的心會消化.
甚至在申命記十一章上帝都應許說.
這些恐懼要遍佈滿全地.
所以這個也再次證明了.
上帝的工作有進行的.
他沒有食言的.
他繼續做他要做的事.
而另一方面我們看到在之後的經文裡面.
迦南人都在用不同的方法去面對上帝的恐懼.
還記不記得有哪幾個不同面對恐懼的一些形式.
包括可能是對抗.
迦南人都有對抗的.
艾城就是這樣.
艾城和以色列要交戰.
同樣迦南人也好像呆了.
站在那裡不懂得給反應.
心都消化 勇氣全失.
甚至有人會用計謀希望討好.
所以欺騙人就假扮.
我們很遠離的.
你可不可以放過我們.
我們可不可以求和.
但是唯獨拉闔.
他是一個不一樣的迦南人.
面對著上帝攻擊這份這麼大的恐懼.
他選擇了效忠在上帝的裡面.
他希望藉著這個機會.
可以令到全家都得到這份請求.
所以拉闔這樣認順.
他說你們的神是上天下地的神.
上天下地這個稱呼.
是指著上帝有絕對的權威.

$^{361}$在迦南人的世界觀裡面.
他覺得這個世界是有不同的神明.
各有各的權力.
各有各的管轄範圍.
但是拉闔在這裡看到的.
不是其中一個有某個管轄範圍的上帝.
而是天上地上地下的主.
擁有一切權柄的上帝.
雖然拉闔仍然聲稱.
這個上帝是你們的神.
但其實他已經承認了.
這個上帝不單超越了迦南的巴力偶像.
但同樣他是超越萬有的.
這是令到他下了這麼大的決定.
要去背離耶利哥.
不惜一切都要投靠真正的上帝.
去進入這個上帝的國度當中.
拉闔告訴我們.
原來真正改變人的選擇.
往往不是如何令恐懼消失.
而是當一個人開始敬畏上帝.
當他見到誰才是真正掌權的時候.
恐懼就會被重新排列.
所以問題是你比較害怕誰呢?.
其實我們生活裡.
都經常聽到或經驗過類似的事情.
舉個例子.
有時和電影姐妹分享.
就是聊天.
為什麼會移民到雪梨裡面呢?.
總有父母會說因為小朋友.
因為他們所以我們移民.
因為很怕小朋友.
要在一個壓力太大的教育環境下長大.
很怕影響他們身心的健康.
所以有些父母下了一個很大的決定.
就是要離開自己熟悉的地方.
去移民到另一個國家.
移民其實都很恐怖.
各位移民過的應該都知道.

$^{401}$有時要離開家人離開朋友.
適應新的語言新的文化.
甚至可能你從來之前做的工作.
建立的那些工作上的能力.
或者職位.
來到可能都沒有了.
要重新再建立過.
重新再適應過.
如果你只是為了自己.
可能有很多人都覺得不要搞那麼多東西.
繼續這樣就很好.
但是因為心裡更大的那種需要.
然後很怕那些孩子.
在那個壓力太大的地方成長.
影響他們的性格.
影響他們心理的健康.
所以縱然有那麼多未知.
很多東西可以值得怕.
但是都不要管了.
再這樣去做這件事.
所以我們都見到.
當一個人見到更重要的事情發生的時候.
他就會開始不再怕.
不再怕以前一直在怕的那些東西.
在信仰上都一樣.
當一個人真正知道.
誰才是那個值得敬畏的上帝的時候.
他就能夠有勇氣.
在那個恐懼裡面仍然忠於上帝.
那一刻的信心.
那個表現是怎樣的呢?.
就是他知道誰是那個真正值得敬畏的.
所以他就不再怕那些其他的東西了.
儘管風險很大.
但他仍然忠於上帝.
所以他就從那個.
本身都會被毀滅的外邦人.
一個局外人.
變成在這個拯救計劃裡面.
有份的局內人.

$^{441}$甚至我們在大衛的家譜裡面.
都看見這個外邦女子的蹤影.
今天你的恐懼是什麼呢?.
又有些什麼是你面對.
而那些會阻礙你去跟從上帝的恐懼呢?.
對我來說除了離別.
一個經常都會影響我事奉的恐懼.
是我很怕覺得自己不夠好.
我很怕有相反的意見.
很怕其他人覺得我很差.
經常都因為有沒有人出席團契.
又或者是團契結束後.
被一個很忠誠的評語者說.
今天的團契其實還可以的.
我就會想很久.
他說了個「幾」字.
「幾」即是不好.
究竟有多不好呢?.
五分的不好還是八分的不好呢?.
這個就會影響我的心情很久.
然後又會弄來弄去.
覺得自己為什麼會做得這麼不好.
為什麼只得到一個「幾」的回應.
為什麼會這樣呢?.
回想起來就會想到.
應該因為小時候在讀寫障礙裡面.
在一個要求背誦寫墨書的香港舊約裡面.
我努力永遠都換取不到成績.
成績永遠都和我的努力沒有成正比.
所以總是覺得自己很差.
總是覺得其他人很怕.
他們如果看到我原來這麼差.
那會怎樣呢?.
因為這種恐懼.
所以我會很敏感其他人的看法.
我都會用不同的方法.
去對抗這些恐懼.
可能我會說服其他人.
會不斷地不斷地說.
因為我覺得爭取贏的那場架.

$^{481}$就代表我自己是夠好的.
又或者我會逃避.
我會對其他可能一些意見.
聽到那邊有些聲音.
我就會避開.
因為我都不懂得怎樣去處理.
覺得被拒絕的那種感受.
我覺得在裡面.
我不懂得回應.
我還是不走過去.
甚至有時候都會很想放棄.
就覺得大家喜歡怎樣就怎樣.
你們說了算.
團體是你們的.
挺好聽的.
但慢慢時風就開始有那種.
打份工.
總之朝九晚五.
你找到我就找到我.
但是找不到我.
我都不管了.
因為那種灰心是覺得.
我反正怎樣都不會做得夠好.
不知道你有沒有這種類似的狀況.
在你的生命裡面.
在你的生活當中.
可能我們和三千多年前.
拉合的恐懼是完全不同的.
但其實我們都有同一個掙扎.
就是究竟我們要怎樣.
我們要馴服在我們自己的恐懼裡面.
還是我們選擇敬畏我們的上帝呢.
的確這是一個選擇.
弟兄姊妹.
因為我們沒有辦法.
活在一個完全沒有恐懼的世界.
我們做什麼都不需要付代價.
不需要擔心.
不會是這樣的.
但我們可以學習.

$^{521}$好像拉合一樣.
在恐懼裡面.
仍然忠於在上帝的裡面.
因為當一個人真正知道.
誰才是天上地下的君王.
我們就會開始不再害怕.
那些以前害怕的東西.
所以我開始.
經常都要問這個問題.
如果我只是怕其他人的看法.
而我不去做.
我相信是上帝要我去做的事的時候.
其實我是在怕誰呢.
誰才是我的主呢.
當然其實說完這句話.
做完那些事.
怕不怕呢.
怕到發抖.
但當我想到.
上帝才是真正明白.
或者能夠了解我的心.
那個最終審視.
我事奉的那個人.
上帝的時候.
我就慢慢去明白.
其實我帶著這些恐懼去做.
但都不能夠因為怕人的評論.
然後我就不去忠於我的上帝.
所以我可以有我的立場.
因為我不是為自己去爭取什麼.
我是在回應上帝看為好的那些事.
但同樣我可以退讓.
因為我不是為了我的面子.
而去做什麼.
今天你的恐懼是什麼呢.
你又怎樣去面對那些阻礙.
來跟隨上帝的恐懼呢.
讓我們一同禱告.
憐憫人的上帝.
我們知道你明白我們的困難.

$^{561}$你憐憫我們的需要.
憐憫我們有軟弱.
有恐懼的時候.
求主你讓聖靈去開啟我們的眼界.
叫我們去看見我們在恐懼的是什麼.
為什麼這麼恐懼.
但同樣求你讓我們去看見.
你的偉大.
你是那位天上地下的君王.
叫我們能夠在你的裡面.
有著恩典能夠站起來.
同樣求主你讓我們.
在仍然未看見.
承諾未兌現之前.
但叫我們仍然有勇氣.
有信心去經歷.
等待.
求主在當中繼續祝福.
和帶領我們每一個.
奉主耶穌的聖名祈求.
阿們.
\newpage



\section{約翰福音 17:18}
\label{sec:031769}
\textbf{上帝為何不差派天使作宣教士?}~~~~羅錫堅牧師
\newline
\newline
連結: \href{http://www.church.com.hk/acms/content.asp?site=cdc\&op=show\&type=product\&code=031769\&layout=sermon}{\texttt{031769}} ~~~~ 語音日期: 2026-03-08
\newline
\newline
\hyperref[sec:031768]{\small{< < < PREV SERMON < < <}}
~
\hyperref[sec:index_chronic]{\small{[返順時目]}}
~
\hyperref[sec:031770]{\small{> > > NEXT SERMON > > >}}
\newline
\newline
約翰福音 17:18
\newline
\begin{longtable}{cl}
\hline
\hline
章節 & 經文 (和合本修訂版)\\
\hline
17:18 & \begin{tabularx}{0.7\textwidth}{X} 你怎樣差我到世上，我也照樣差他們到世上。 \end{tabularx} \\ \\
[1ex]
\hline
\hline
\end{longtable}
$^{1}$大家都知道我幾年前就加入了播道差會成為宣教士.
我心裡面就越來越有個感觸.
宣教工場好缺乏宣教士.
好缺乏.
我而家都請緊好多宣教士.
而且仲未請到.
因為大部份基督徒都唔想做宣教士.
那些想做宣教士的又未必個個口齒伶俐.
口齒伶俐的又好似財多身子弱成日病.
唔病的可能年紀大有曬所有抗體.
好似我咁.
年紀大就有限制.
怪唔知上帝在《以賽亞書》第六章第八召.
我可以差遣誰呢.
誰肯為我唔去呢.
因為冇乜人想去.
你知我諗ye 好多時諗出位.
我就諗一個好特別ge 問題.
點解上帝唔差派天使來做宣教士呢.
天使做宣教士可以解決好多問題.
有好多優點.
第一.
天使千千萬萬又唔會托手肘.
第二天使口齒伶俐唔會講錯ye .
你聽過邊個天使講錯ye 未.
墮落ge 三分一就唔計.
而家係上帝身邊服侍ge 三分二就唔會講錯ye .
天使會唔會病.
唔會病.
天使又唔會老.
天使又唔使訓覺.
又唔使食飯.
唔使坐飛機.
坐飛機好多時間.
天使仲飛得快過飛機.
又唔會浪費時間.
天使唔使籌款.
咩都唔使洗.
係咪用好慳時間.
所以天使做宣教士其實係非常之完美.

$^{41}$大家贊唔贊成.
有邊個贊成請你舉手.
有冇呀.
有幾個.
其他gwo D 你唔贊成.
咩原因唔贊成先.
可唔可以解釋下點解你唔舉手.
但我諗來諗去諗到最撈尾.
都要承認其實上帝係唔會派天使做宣教士.
天使係唔適合做宣教士.
咁點解唔適合做宣教士呢.
我地睇o下.
第一個原因.
天使係冇得夠見證.
你同我傳福音.
我地講到我沉淪.
上帝救贖我.
咁人地聽起上黎有共鳴.
天使冇.
大家聽過一首歌未.
Amazing Grace.
奇異恩典.
何等甘田.
甘以得道赦免.
呢個作者John Neutron.
原本係反賣奴隸.
後來信jor 耶穌做jor 牧師.
成個改變到.
我罪得以赦免.
史特普羅都話.
我係罪人當中我係個罪魁.
原本我係迫害教會.
而家我成為何等樣ge 人.
佢一個改變令到聽者有共鳴.
但天使做唔到呢點.
第二天使冇選擇.
佢未經過被召考慮.
掙扎馴服毀身過程.
我地試過.
或者有人話.

$^{81}$朗斯里好似冇乜.
成日見你笑都未見過你掙扎.
其實我有.
你知唔知.
當播道會猜獻禮gwo 陣時.
我同我太太跪左係度見唔見到.
江牧師gwo 陣時係我地差會ge 主席.
何志迪牧師就係我地.
香港派駐跨地域合作宣教計劃.
ge 其中一位代表香港.
同埋一班播道會ge 總會.
差會ge 董事.
gwo 陣時我71歲.
播道差會未試過猜派一個71歲ge .
第一次出去工場宣教.
我都唔知有邊間差會會猜派一個71歲.
你71歲去到台灣就好危險.
唔係台灣危險.
又gwo D 路凹凹凸凸高高低低成日撲撐gwo D 危險.
咁其實我地要經過考慮.
經過掙扎.
經過信服.
經過毀神.
咁幾個月前我又兼任埋多一份義工.
就係差會ge 總幹事.
因為我地ge 前總幹事辭職.
咁我如果每個月就來回差會同埋台灣.
咁就成日出出入入.
咁所以你地見到我好似多jor 呢.
其實有道理ge .
因為我有兩份工.
我ge 生活好充實.
咁呢個其實都係一D ge 掙扎.
我記得以前同大家講過.
我55歲返香港gwo 陣時同太太講.
唔係佢要求我.
我65歲返到加拿大買部靚車.
買部Cadillac.
就唔係好似我地結婚gwo 部車頭車尾都出煙gwo 種.
咁如果你買電動車就唔會出煙.

$^{121}$咁但係到而家都未兌現.
見到太太都唔係幾好意思.
咁但係呢個都係要經過掙扎信服ge 過程.
呢個選擇好緊要.
上帝唔想逼我地.
上帝俾你有個選擇.
所以有D 信徒係教會坐係度唔做ye .
其實上帝唔開心ge .
但佢俾你有個選擇.
其實仲有一點.
上帝差派你去做宣教士唔差派天使呢.
因為你比宣教士更加重要.
你比天使更加重要.
無錯.
點解你又會問.
我比天使重要呀.
我其實呢就成日坐係度.
天使一走就有兩隻翅膀.
你知唔知聖經話天使係一個服役ge 靈.
係一個Spirit.
係Servant.
係服侍上帝ge 靈.
你同我就係咩.
我地係上帝ge 兒女.
上帝創造我地.
係用佢ge 形.
按佢ge 形象黎做我地.
即係按佢上帝自己ge 形象.
佢ge 本質黎做我地.
咁我地係上帝ge 兒女.
我地蒙恩救贖之後.
我地直情係.
係與基督同一同坐係寶座上面.
將來我地要審判天使.
所以如果皇宮裡面你就係王子.
你就係王子.
雖然你阿媽就同你講.
你著jor 龍袍都唔似太子.
但唔似太子都係太子.
係咪呀.

$^{161}$你唔似上帝ge 兒女都係上帝ge 兒女.
所以上帝差派你去.
就差派一個好重要ge 代表去.
咁呢個上帝要差派一個好重要ge 代表去.
咁即係梗係好重要ge 職份.
天使呢個宣教士一個好重要ge 職份.
咁仲有第四點呢.
就係呢個係一個好有恩典ge 呼召.
即係上帝比你去呢.
係其實呢個恩典.
你唔好以為.
我出去同阿江牧師去派單張.
呢個就俾面上帝啦簡直.
平時我唔做.
我完jor 之後就同阿朱先生去飲茶.
而家就出去派單張咁樣.
咁呢個司祿普羅就講得好清楚.
佢話當我係母腹裡上帝就將我分別出黎.
又用佢ge 恩.
佢ge Grace 佢ge 恩呼召我gwo 一位.
既然樂意將自己ge 兒子啟示給我.
使我可以在外族人中傳揚他.
普羅就Catch到呢D .
呢個係恩典.
上帝今日叫你係傳恩堂成為一個會友.
成為一個出去服侍.
或者你歡迎大使.
或者你派單張.
或者你彈琴.
你打鼓.
你領詩.
廣道呢D 全部都係恩典.
唔係你有咩ge 本能.
又唔係你點樣去俾面上帝.
係上帝恩道.
所以因為咁ge 緣故.
司祿普羅珍惜呢個呼召.
比自己ge 性命仲珍惜.
所以佢話.
我並不珍惜自己ge 性命.

$^{201}$只求跑完我ge 路程.
完成我從主耶穌所領受ge 職份.
為神恩惠的福音作見證.
普羅就睇到.
因為咁ge 緣故.
佢又放低佢所有ge ye 去傳福音.
呢D 普羅好勁.
佢出身顯赫ge 家庭.
有錢有勢.
佢係羅馬公民gwo 個身份.
佢gwo 個老師就係全羅馬帝國最出名ge .
學問高深.
酬老精靈.
本來就大好前途.
但佢放低所有ge ye .
去做宣教士.
佢話我比人打過幾多下.
我係水裡面.
係大海裡面漂浮過幾多次.
我坐過幾多監.
受過幾多苦.
仲有喎.
普羅出去傳福音.
其實有陣時佢都唔夠食.
唔夠衫著.
記唔記得佢坐監獄.
叫阿提摩大打件衫比佢.
所以其實佢唔珍惜自己ge 生命.
佢睇到上帝呼召佢做宣教士gwo 個恩典幾奇妙.
有個故仔就講到耶穌.
呢個故仔黎ge .
唔知係咪真ge .
講到耶穌復活jor 就升天.
升天呢.
天使長Gabriel就問耶穌.
主耶穌你係世上已經做哂所有ge 救贖工作.
咁你點樣讓全世界ge 人知道呢個救恩.
相信你得救呢.
咁耶穌話.
咁我就訓練jor 幾個門徒.

$^{241}$11個.
原本12個就冇jor 一個.
仲有11個.
咁呢我就將呢個球交jor 比佢地.
由佢地去傳比其他ge 人.
其他ge 人又傳比其他ge 人.
咁就傳到全世界.
就係咁架.
咁Gabriel就覺得危險D .
喂.
呢11個都好似唔係好掂.
你咁樣叫佢地.
咁如果佢地唔傳呢.
或者佢傳錯jor 呢.
咁或者佢傳完下 甩jor 碌呢.
咁你有冇一個plan Bge .
有冇一個計.
有一個叫做.
我地叫做contingency plan.
係咪呀.
咁中文呢.
翻譯係咩呀.
即係話呢我地有個後備計劃.
如果呢個唔得呢.
就會唔會有一個plan B呢.
咁耶穌話冇plan B.
再靠佢地.
弟兄姊妹.
上帝冇plan B.
上帝就係靠你傳福音.
上帝唔會搵天使.
個理由已經講得好清楚.
咁就搵你.
咁唔搵你搵我.
所以最後晚餐gwo 陣時.
耶穌就講呢句話.
就剛才阿偉同我地讀呢段經文.
我知道大家聽到呢.
就最開心ge 經文短.
因為長gwo 陣時又阻時間.

$^{281}$而且又唔記得.
咁呢個呢我就做個好好榜樣.
你怎樣猜我到世上來.
我又怎樣猜他們到世上去.
其實呢句話後面呢.
係好多ye .
gwo 陣時係最後晚餐.
耶穌知道大家就快要.
要俾人釘死.
而且一個猶大又背叛佢.
咁佢.
佢一個禱告.
呢個禱告係全個.
我諗全個新約最勁ge 禱告.
約翰福音十七章你返去睇下.
佢對上帝話.
你怎樣猜我係世上來.
黎到呢度.
黎到呢度做咩呀.
黎到呢度釘緊十字架.
黎到呢度完成呢個救贖.
我又照樣猜佢地去世上去.
去做咩呀.
就係將呢個救贖計劃講出去.
傳出去.
叫個個人都聽到.
呢個.
就係.
我地ge 使命.
呢個就係點解我地全恩堂.
每年三月都有一個宣教月.
猜全月.
就係我地要.
將耶穌ge 福音.
帶到世上去.
咁我地知道.
點解上帝唔會叫天使做宣教士.
點解上帝冇plan b.
咁而家我地就諗.
咁究竟我地全恩堂弟兄姊妹.

$^{321}$今日都好似人才.
仔仔喎好多人係哂度.
我地點樣去.
點樣去傳福音呢.
點樣去做呢個宣教呢.
我首先要俾大家一個好清楚ge .
ge 認識一個.
宣教ge 概念呢個concept.
讀書ge 最緊要.
舊時我地.
我讀書有好多年代.
我讀書gwo 時你阿媽未出世.
起碼其中一兩個.
咁呀.
都唔知係咪可能出一世.
咁.
gwo 時個教授個老師話.
你最緊要係明白個concept.
個觀念.
你搞掂呢個concept得.
如果唔係你代代都唔得.
三個concept大家記住.
就包你受用不盡.
第一個呢.
就係教會.
係一個出去ge 群體.
出去ge .
好多教會內聚.
我地四面牆係度.
入黎gwo D 三寶人我地秒下佢.
咦呢個我唔識.
佢唔知有冇細菌帶入黎呢.
或者講你乜ye .
或者伯藏呀或者偷我地ye 呀.
即係我地.
好多教會呢.
係好內聚.
自己keep係度.
你知唔知個世界.
好孤單.

$^{361}$日本每一年.
有二萬人死亡係.
孤獨去世冇人知道佢死.
每一年有二萬人.
咁而家好多國家.
包括香港台灣日本.
南韓呀甚至可能.
外國都有一組叫出租人.
即係陪伴ge .
我就扯佢黎.
同你做ye 喎唔係同你打工.
唔係同你清潔.
佢黎呢就陪你ge je .
咁有D 人好孤單我想陪佢.
搵個人陪我食飯.
佢時中計.
佢同你食飯就唔同你傾計.
咁佢就.
唔會比你意見.
但係就陪你.
或者你話我想搵個人陪我睇戲.
日本最近呢.
就有好多孤獨ge 人呢.
咁佢地呢就新聞報道.
佢地呢就D .
佢地去.
幕後買D 奶呀.
有D 公司.
專登搵D 會同你傾計ge 人.
帶D 奶送比你.
咁佢每日一次呢就期待呢個呢.
就係gwo 日ge .
呢個就係gwo 日ge 高潮啦.
佢.
佢有個人黎送一盒奶比我.
同我傾一陣計.
有好多孤單ge 人.
唔單止成人孤單.
細蚊仔都孤單.
香港中小學生自殺率.

$^{401}$每一年都提高.
咁我地全英團搬到北角呢.
就發現有好多貧窮ge 人.
北角好多.
即係生活窮困ge 人.
多過柴灣.
我相信全香港.
都有好多孤單ge 人.
就唔係淨北角有.
咁所以我地全英團.
我地應該要做一個出去ge 群體.
我地睇到gwo 度ge 需要.
我地要.
我地要去.
接觸佢地.
耶穌就咁啦.
佢見到好多ge 人困苦流離.
佢就憐憫佢地.
因為佢好似羊無目人一般.
我地要存一個憐憫ge 心.
睇到周圍ge 人.
我地好欣賞我地全英團.
我地聯絡埋.
北角ge 教會.
唔單止北角 其他地區.
一齊去樹仁大學.
去教廣東班 廣東話班.
我地係度有開長者班.
開長者團契.
又開呢個.
兒童ge 活動.
又有廣東話班.
職場廣東話班等等.
咁其實呢D 就係出去ge 群體.
每一個信徒都係要出去ge .
因為耶穌就咁講.
所以你地要去.
go 使萬民作我的門徒.
奉父子聖靈ge 名.
給他們施洗.

$^{441}$我吩咐你地一切都要.
教導佢地遵守.
咁我就尚與你地同在.
直到世界ge 終結.
呢個係馬太福音.
28章19到20節.
有一本書叫做.
Transforming Mission.
六百幾頁英文.
中文有日文都係六百幾頁.
英文版呢 柯牧師有.
呢本書好勁.
呢本書就係被.
Christianity Today.
話係.
20世紀最勁ge .
100本書之一.
唔單止係基督教書籍.
20世紀人類出所有ge 書.
最緊要ge .
100本之一.
佢有一點佢講到.
六百幾頁英文你點睇呀.
柯牧師咁樣.
我講一句你就catch到個重點.
佢話你睇馬太福音.
唔單止由頭睇到尾.
你要由尾睇返到頭.
個尾就係馬太福音28章.
你透過呢個.
耶穌基督ge 大使命.
調返轉頭睇馬太福音.
你會睇到完全唔同ge picture.
完全唔同ge 景象.
原來馬太福音.
係一個宣教猜拳ge 書.
其實聖經好多書.
整本聖經就係宣教猜拳.
教會.
就係宣教.

$^{481}$但好多教會.
播道會都好多教會係無猜拳報.
三分一有.
大部分信徒.
都冇諗到宣教.
亦都冇支持宣教.
我地冇諗到.
我地係代表上帝.
出去出面.
咁你話宣教去到咁遠.
我唔去.
你唔需要去到日本.
唔需要去到台灣.
你上進中心出面.
就係你ge 工場.
你去到邊度就傳到邊度.
我自己都一個例子.
我二十八歲去讀書.
我二十三歲結婚.
廿四歲我同我太太.
見到gwo 度冇乜.
福音純正ge 教會.
我地就.
會一埋幾個人開一間播道會.
咁呀.
而家就變jor 六間七間.
自己開花結果.
咁再過一段時間.
我地開一間浸信會.
咁再過一段時間.
我地又同一個獨立教會開jor 個.
開jor 個分堂.
咁然後上帝先呼召我做傳道人.
你唔需要呼召我去.
放低工作讀神學做傳道人.
我未讀神學之前就開架喇.
咁港道係我喎.
師餐又係我喎.
當時七零年代.
香港播道會就唔準.

$^{521}$D 評信徒港道師餐.
肯定師餐唔準.
咁我又犯jor 規.
加拿大ge 播道會就自由好多.
Free Church.
咁上帝又冇罰我.
即係我個point就係.
每一個信徒.
都係宣教士.
一個教會唔係.
就咁去做宣教.
有個宣教部門.
有個巫師呀韋紅呀咁樣.
總之個宣教部門去做下.
成個教會就係宣教教會.
呢個宣教應該係你ge DNA.
如果你唔知道宣教做乜呢.
呢個係你ge 問題亦都係教會ge 問題.
但係我地全因堂講jor 咁多.
清楚講宣教如果你都唔知.
呢個就係你ge 問題.
所以我地要知道呢個.
係宣教.
就係我地教會.
ge being就唔係doing.
教會.
唔應該內向.
唔應該躺平.
宣教唔應該靠幾個人靠一個部門.
而係每一個信徒都要出去.
第二點.
宣教係一個差憲ge 群體.
我地要時刻.
培養.
我地ge 信徒成為宣教士.
我記得.
我.
好多時我係報道會有兩間英文.
堂我係gwo 度講到.
其中一位長老叫Dave Leung.

$^{561}$佢係美國出身ge .
華裔ge 人.
佢就.
結婚係香港結婚.
佢太太係湖北人.
佢太太講普通話就好正.
廣東話都好正.
生下三個仔.
有一次問個仔.
細個你大概想做咩.
arm arm 聽arm arm 聽.
我覺得呢句話後面.
睇到佢爸爸媽媽.
係培育佢將來做神的朋友.
其實佢係咪做牧師.
上帝會唔會呼籍佢唔知道.
但係我地教會都要留意.
有D 有潛力ge .
有心智ge .
好好地培育佢支持佢.
差憲佢.
呢個就係差憲ge 教會.
就係咁被差憲.
史獨行傳十三章.
安提柯教會中有幾位先知和教師.
就係巴拿巴.
古萊利人ge 路球.
馬念.
同埋蘇羅即係保羅.
佢地事奉主並且禁食ge 時候.
聖靈師.
要為我把巴拿巴和蘇羅.
分別出來.
去作我呼召他們所作的功.
於是佢地禁食禱告.
為兩人按手.
猜派佢地出去.
教會唔係偶然ge .
猜一兩個人出去就掂.
全因堂都好叻.

$^{601}$猜派jor 我同我太太出去.
全因堂亦都支持好多ge 教會.
但我地希望將來.
全因堂我地愈來愈大ge 時候.
我地愈來愈多人ge 時候.
我地就從今日開始.
禁食禱告.
究竟我地當中有邊D 人.
我地將佢分別出黎.
猜派佢出去.
我地唔係一個.
做宣教ge 教會.
我地係一個beeing.
我地唔係do宣教ge 教會.
我係be一個宣教ge 教會.
第三.
教會係一個公認ge 群體.
我地需要公認ge 宣教士.
你知唔知道.
一個軍隊一萬人.
一個軍隊一個division.
一個師一個陸軍一個師.
師長一個師.
一萬人裡面真係.
打仗zar 槍.
向前線打仗ge .
係三分一.
其他三分二.
就係要做通訊.
醫療.
gwo D ge .
而呢個一萬人裡面.
一萬人後面gwo 度.
就要有國家ge 資源.
因為三軍未動糧草先行.
你諗唔諗到.
我地派一個一萬人.
ge 軍隊上前線.
真係zar 槍.
搏鬥ge 係三分一.

$^{641}$其他六七千.
就係要.
讓佢幫佢地.
提供佢地ge 援助.
叫佢地知道去邊度打點樣打.
而呢個軍隊後面.
成個國家要支持.
咁呢個ge .
給予ge 群體係.
係好重.
腓立比書.
gwo 時保羅坐緊監.
咁佢話 腓立比ge 弟兄們.
你地都知道.
係我傳福音ge 初期.
離開馬其頓ge 時候.
除jor 你之外.
我冇收到任何教會ge 功級.
我係帖撒羅尼迦ge 時候.
你們還是一而再.
把我需用ge 送來.
la  剛才講到保羅咁勁ge .
但佢唔夠錢.
佢即帳盤.
即帳盤呢份ye .
都唔容易做.
我就未試過.
咁.
但係保羅要落到咁ge 場地.
即帳盤.
我同我太太.
係台灣.
我地ge 波道差會會計就好小心.
佢有一次問.
問我太太.
你地兩個早餐.
兩個夾埋食150蚊台幣.
夠唔夠飽.
佢從來未試過問.
星期食夠唔夠飽.

$^{681}$通常個早餐費都幾高.
但唔少蚊我兩個.
150蚊除4.
就等於大約30幾蚊港幣.
咁我地慳.
不過我地都未到即帳盤ge 地步.
想做都唔識就行.
咁但係其實我地.
真係需要公認宣教士.
咁好欣賞傳音堂.
傳音堂呢我地ge .
你睇我地ge .
我地ge 壁報呢.
我地今年ge .
我地呢一年ge 預算係.
34萬幾.
咁呢個係超過我地.
整個堂會ge .
經常費ge 十分之一.
超過我地.
咁呢個係好ge .
但係如果更多就更好.
因為我地係一個宣教ge .
ge 教會我唔係淨係做.
宣教我地本身就係.
宣教每一個人.
每一個細胞都係宣教.
你.
你戀愛ge 時候你對個女朋友.
話我每個細胞都愛你.
你對上帝話我每個細胞.
都宣教咁就arm .
時先生你係咪每個細胞.
都宣教.
唔使答我.
時先生望住我.
好嚴肅咁.
好咁你話.
牧師講得咁清楚.
呢個ge .

$^{721}$出去ge 群體.
吸引ge 群體.
咁我點樣做呢.
我有咩ye 可以做呢.
咁而家我地講啦.
究竟係全印堂在在每一位.
有乜ye ye 你可以做.
第一加入我地教會ge 宣教.
行列.
第一為宣教士禱告.
我真係好需要為宣教士禱告.
你ge 禱告就係我地ge 生命線.
我地.
我地全印堂支持.
8位9位ge .
宣教士.
9個單位.
佢每一個都要我地禱告.
我地係台灣.
係日本係澳門.
其實呢三個地方.
都係硬土.
台灣人好食好虔誠.
拜拜.
咁.
咁即係信耶穌ge 人.
真正重新得救.
得2\%.
日本得1/4\%.
係咪硬土.
澳門呢唔知點解間間教會.
澳門有個我記得幾年前.
同一位牧師經理牧師.
佢做jor 幾十年佢gwo 陣時已經六十幾歲.
gwo 陣時都成十年前.
佢應該年紀大過我.
我間教會做幾十年.
都得幾十年因為唔知點解D 人.
唔郁ge D 人又唔信耶穌.
信耶穌又唔郁.

$^{761}$即係好一個比.
殺彈籠罩得好犀利.
我地一定要派D 尖兵去打.
經濟上ge 支持.
係美國.
舊年有一對年青ge 宣教士.
去jor 台灣.
而家我每個月同佢地一齊.
佢地每年ge 經費.
要籌超過一百萬港幣.
佢有教會支持佢地.
佢地一期四年.
係承諾去支持佢.
咁呢個.
就係好感恩.
美國播道會.
美國播道會就係係.
1835年之前黎中國ge .
宣教教會.
佢地美國播道會ge .
差會有六百幾個宣教士.
三十萬人.
六百個宣教士.
香港接近三萬人.
香港播道會接近三萬人.
我地宣教士ge 比例就少過佢地好多.
我地得十幾位.
唔夠二十人.
所以我地其實真係有好多ge .
好多ge 發展ge 空間.
第三.
大後方ge 參與.
我地講過.
宣教士出去真係需要後面ge 支援.
咁你可以做咩ye 呢.
大後方ge 參與.
你可以.
參觀.
我地.
傳恩堂ge 宣教士工.

$^{801}$去搵傳恩使命團.
就係我地ge .
組織出黎做宣教.
搵佢傾計.
有邊個係傳恩使命團ge 成員.
請你舉手比大家認識下.
邵芳.
彬彬.
邝牧師仲有呢.
仲有彬豐仲有冇呀.
大家去搵佢傾計.
佢地好樂意同你傾.
你關懷個別ge 宣教士.
呢個好緊要.
佢地係.
異地其實都幾孤單.
我地有宣教士係柬埔寨.
兩個月前呢.
泰國同柬埔寨打仗.
佢係佢出面冇幾耐.
就見到D F16戰鬥機.
炸D 炸彈落黎gwo 度.
我地有宣教士係中東.
宣教士係中東.
而家中東係立立面.
D 飛彈唔知邊個飛.
個個都唔郁得.
飛機都唔郁得.
宣教士就係呢個環境.
台灣又好好.
除jor 地震之外.
地震颱風就冇乜其他D ye .
咁.
咁所以其實.
佢地需要我地去.
去關懷.
第四.
前線.
往前方.
希望下次全因派出黎.

$^{841}$gwo 個就唔係好似我咁大年紀.
而家在在當中.
有D 都年紀輕好多.
你輕我十年gwo D 就算係年輕.
咁.
咁或者會唔會上帝真係.
差派你出去.
咁你唔使驚.
因為我都出到去你都出到.
係人類ge 歷史裡面.
上帝有兩次.
ge 賦予人類好重要.
ge 任命兩次ge je .
第一次呢.
就係伊甸園.
創世記.
第一章第28節.
神就賜福給他們.
對他們說.
要繁衍增多.
蒼滿者地征服他.
要管理海.
利的魚 空中的鳥.
同埋地上所有ge 生物.
嘩好勁喎.
比個命令你 比個任命你.
比個authority你.
管理曬全世界所有ge ye .
空中ge 飛鳥.
海裡面ge 魚.
地上ge 走獸 所有環保.
全部你搞掂 全部係.
under你 勁唔勁呀.
咁大家知道我地就和jor .
啦係咪呀 伊甸園都未.
出黎就繁jor 盡.
咁成個世界都淪陷jor .
咁以至呢就產生出羅馬書.
第8章21.
第二節gwo 個光景.

$^{881}$被造ge 萬物盼望自己得著.
釋放脫離敗壞ge 奴役.
因為人ge 犯罪.
以致帶黎敗壞成為jor .
我地ge 奴役得著神而.
呂榮耀ge 自由我地知道.
被造ge 萬物直到如今.
都係痛苦深入.
咁其實第一條任命呢.
就用得唔好啦係咪.
過jor 咁跟住耶穌基督.
復活節ge 時候呢.
佢升天之前比我地第二個任命.
就係剛才我讀gwo .
gwo 兩字經文.
所以你們要去使萬民.
作我的門徒奉父子.
聖靈的名給他們施洗.
上帝點樣派耶穌到世上來.
耶穌亦都點樣派你同我.
到世上去.
呢個任命就緊要第一個.
因為第一個管理呢個地球.
管理呢個地球.
佢上帝都冇叫你.
管理宇宙ge 即係地球.
管唔到宇宙.
地球其實呢個天地會廢去.
將來會冇jor ge .
即係唔係永久ge .
但係我地搵人.
做耶穌基督ge 門徒呢就係永恆ge .
個價值留到永遠.
將來每一個得救ge 靈魂呢.
都係天國裡面度過永恆.
每一個唔得救ge 靈魂呢.
都係地獄裡面度過永恆.
所以第二條任命就緊要.
過第一條ge 任命.
咁今日.

$^{921}$全因堂ge 猜拳月裡面.
我請阿 Keith 同我整jor 一幅相.
就係我地.
全因堂ge 宣教.
ge 單位.
你都睇到.
其實如果你係逼報版度.
你就睇到更多.
睇到我地ge 預算.
我地而家籌到幾多.
我鼓勵大家真係睇下.
我地同心合意ge .
將全因堂成為一個.
出去ge 猜獻ge .
同埋供給ge .
教會我地一齊低頭祈禱.
張教主我地多謝你.
你自己ge .
救贖你ge 十字架上面.
ge 待死其實係成為我地ge 激勵.
我地願意.
承擔呢個使命呢個榮耀ge 使命.
就係出去我地要.
將你ge 福音傳開.
求你ge 令感動我地帶領我地.
叫全因堂成為一個.
宣教ge 教會我地咁.
祈禱奉耶穌基督的名.
\newpage



\section{創世記 27:2}
\label{sec:031770}
\textbf{聖經人聖經事:不知哪天死}~~~~黃克勤牧師博士
\newline
\newline
連結: \href{http://www.church.com.hk/acms/content.asp?site=cdc\&op=show\&type=product\&code=031770\&layout=sermon}{\texttt{031770}} ~~~~ 語音日期: 2026-02-26
\newline
\newline
\hyperref[sec:031769]{\small{< < < PREV SERMON < < <}}
~
\hyperref[sec:index_chronic]{\small{[返順時目]}}
~
\hyperref[sec:031771]{\small{> > > NEXT SERMON > > >}}
\newline
\newline
創世記 27:2
\newline
\begin{longtable}{cl}
\hline
\hline
章節 & 經文 (和合本修訂版)\\
\hline
27:2 & \begin{tabularx}{0.7\textwidth}{X} 他說：「看哪，我老了，不知道哪一天死。 \end{tabularx} \\ \\
[1ex]
\hline
\hline
\end{longtable}
$^{1}$眨眼又一日.
日子過得這麼快.
真是不知到那天死.
這句話是以撒說的.
不知到那天死.
《聖經人聖經事》由生命說書網絡.
跟大家一起讀聖經.
知道人性與神的工作感覺.
二十七章以撒年老.
一說就一條條到以撒年老.
以撒一生在聖經佔據的內容很短.
不是居於很多大作為.
不過承接著阿巴拉罕與雅各.
是他的重要性.
以撒年老.
眼睛昏花不能看見.
就叫了他大兒子以素來說.
「我兒(以素孫)我再借女」.
眼睛昏花.
「卡哈」的意思是很微弱.
眼睛很弱.
「眼睛昏花」就是眼睛很弱.
弱到怎樣?.
下一句就說了.
是弱到不能看見.
可能看到光,少少影.
誰都認不出.
他叫了他大兒子以素來說.
「兒子」.
他叫兒子先當然是親切.
很疼愛他.
因為以撒疼愛以素.
尼伯加疼愛雅各.
「兒子」.
以素說了.
「我再借女」是希伯來文一個字.
Henay這個字其實不是我在借淚.
Henay這個字是「去」的意思.
有點像甚麼呢?.
很多人在說「兒子,兒子在哪?」.

$^{41}$兒子聽到就「去」「去」「去」.
如果是希伯來人或道義式的人.
他們會說「Henay Henay」.
「在這裡,在這裡」或「這裡,這裡」.
「看一看,看一看」.
所以這個「看」不是「behold」.
不是「看一看」.
是「這裡,這裡」.
很多東西要搬,誰來幫忙撿椅子.
「Henay」「這裡,這裡」.
很多人有個巨獸「這裡,這裡」.
這就是「我在借淚」的意思.
這亦在《耳塞俠書》裡用過.
「我在借淚請差遣我」.
「我可以差遣誰,誰肯為我們去呢?」.
耳塞俠說「我在借淚請差遣我」.
「我在借淚」就是「Henay」.
意思是「Henay」是說.
「誰可以為我做些事?」.
「Henay」是這個意思.
「我在借淚」.
耳塞俠就說「Henay」.
耳塞俠又回應「Henay」.
這個字「我如今老了」.
他說「我如今老了」.
這個綠色字是希伯來文有.
而中文沒有翻譯.
而耳塞俠就說「Henay」.
然後耳塞俠又說「Henay」.
「你看看,在這裡,多老」.
「我如今老了,不知那一天死」.
這個「那一天」用藍色括弧.
即是希伯來文是沒有的.
「看看,很老了,不知死」.
即是「不知死」.
其實他是說「不知死」.
即是「不知到死」是甚麼.
其實是說「那一天」.
不過沒有用這個字.
「不知到幾時死」.

$^{81}$「現在你拿你的器械」.
「就是戰狼公往田野去為我打獵」.
「兒子,在這裡,你看看,多老」.
「不知幾時死」.
「現在你拿你的器械」.
「就是戰狼公往田野去為我打獵」.
「然後給我吃東西」.
「打獵」在中文來說是一個動作.
但這個「打獵」這兩個字.
「為我去打獵」有三個希伯來文的字.
第一個是「Jude」.
「Jude」即是「打獵」.
「Jude」即是「獵物」.
「為我打獵獵物」.
然後「Jude」即是「給我打獵獵物回來做食物」.
然後「煮我喜歡的味道」.
「煮我喜歡的味道」.
「然後拿來給我吃」.
「我不知幾時死」.
「你知不知道自己幾時死呢?」.
「不知」.
「不知又怎樣?」.
「你就要把握自己」.
「數數自己的日子」.
「好叫我得著智慧的心」.
「怎樣去生活」.
「因為我們都不知自己幾時死」.
「自己一眨眼以為自己很小」.
「真是不知幾時死」.
「我就開始去想現在做什麼」.
「以實有冇諗呢?」.
「有」.
「和我打獵做餐好ye 食」.
「成為我呢生最後ge 記憶」.
「請D ye 食食」.
「當咁樣你拿來給我吃」.
「使我在未死之先給你祝福」.
「你咁樣自己給我食」.
「是如此才能阿伯」.
「阿伯就是咁樣令到我滿意」.

$^{121}$「咁我就會係未死之前」.
這是動詞.
「死」是名詞.
中文第幾句是動詞.
「Move」.
「就係我死」.
「將要死」.
動詞有Tense.
「將要死之前」.
「就為你祝福」.
他已經有計劃.
「我唔知幾時死」.
「不過睇黎已經好老喇」.
「唔係幾好得ye 又睇唔到」.
「隨時就一下」.
「睜唔大眼就死jor 喇」.
「食番餐先」.
「和我打D 好ye 」.
「你食得我滿意」.
「我為你係死前祝福」.
「係我死前祝福」.
「做D ye 整得我滿意」.
「我就同你祝福」.
有少少脅迫ge 意思.
「就為你祝福喇」.
「呢個祝福唔係普通ge 祝福」.
「呢個祝福係遺命」.
「透過上帝祝福你」.
「呢個祝福係真係會成架喎」.
「因為以撒係跟住亞伯拉罕講比佢聽ge ye 」.
「然後耶和華點樣賜福我ge 後裔」.
「所以佢咁樣祝福係傳住亞伯拉罕ge 傳承落黎架喎」.
「所以係會發生ge 係因為耶和華」.
「咁我地ge 神就祝福喇」.
「講下祝福先」.
「祝福好多時候完jor 崇拜喇」.
「好請牧師祝福」.
「然後呢個牧師就上去喇」.
「好請領受上帝ge 祝福」.
「咁就祝福喇」.

$^{161}$「我唔係呢個神學ge 」.
「有D 仲係講好請黃牧師頒佈上帝ge 祝福」.
「嘩!我都唔知講唔講好」.
「因為呢個唔係我ge 神學黎架」.
「因為聖經係話從來都係」.
「為大的給為小的祝福」.
「所以呢其實成個崇拜裡面最受尊敬ge 」.
「唔係咁講呀」.
「為高ge 牧師囉佢帶領架嘛」.
「咁而家散會喇我地要走喇」.
「牧師為我地祝個福」.
「咁為我地求個福然後離開喇」.
「咁我唔祝福囉」.
「咁其實個福呢」.
「係個牧師自己ge 祝願與禱告」.
「但願」.
「我ge 但願黎ge je 」.
「如果上帝使乜但願呀」.
「個福氣跟住你」.
「跟住你使乜但願呢」.
「即係話係我ge 祈求祈願黎ge 」.
「所以我唔係頒佈緊上帝ge 祝福」.
「我真係聽過一個會友」.
「一個弟兄竟然夠膽問牧師」.
「問我真係當時未識點答」.
「佢問牧師呀」.
「點解你個禮拜崇拜完jor 祝福」.
「你祝完福我都冇福ge 」.
「嘩真係尷尬至極架」.
「點解祝完福都冇福ge 」.
「真係唔知點答」.
「後面我先至睇到」.
「原來係為大的及為小的祝福」.
「就好似爸爸媽媽送個細路仔出門喇」.
「喇今日搭校車小心o下」.
「自己搭車返去搭地鐵小心o下」.
「咁睇住路呀平平安安去平平安返」.
「阿媽就話小心食ye 健健康康」.
「讀書勤力」.
「爸爸媽媽ge 祝福」.

$^{201}$「咁當然而家D 細ge 都話」.
「爸爸媽媽我都祝福你身體健康」.
「又賺多D 錢買多D 雪糕比我食」.
「咁都好多細ge 祝福」.
「其實我地都開心幾中意ge 」.
「不過其實聖經係話」.
「為大的及為小的祝福」.
「咁好啦咁你明白啦」.
「但呢個祝福就唔係啦」.
「係上帝透過同阿伯拉伯約」.
「然後會去到佢ge 後裔ge 」.
「好啦,留意」.
「我想講ge 就係」.
「Heine,睇下我而家老到咁啦」.
「不知到哪一天死」.
「就係今日ge 金句」.
「咁就做到金句」.
「不知到哪一天死」.
「係ge ,我好想你學識數算字ge 日子」.
「有機會讀下詩篇第九十篇」.
「求你俾我智慧ge 心」.
「叫我數算自己ge 日子」.
「好叫我得著智慧ge 心」.
「我一生ge 年日係七十強壯到八十」.
「不過係勞苦愁煩轉眼成空」.
「我門便如飛而去」.
「所以求你教導我點樣數算自己ge 日子」.
「好叫我得著智慧ge 心」.
「咁就年齡啦,仲做到D 乜ye 」.
「哎呀,我好想做gwo 樣ye 呀」.
「咁樣做囉」.
「諗辦法行緊佢啦」.
「嘩,要讀好多書架喎,讀囉」.
「哎呀,我好想做咩工作呀」.
「咁你又唔讀書」.
「咁你又唔做ye ,又唔行埋去」.
「咁點會發生?」.
「哎呀,神都冇俾我呢個機會」.
「俾咩je ?」.
「特別一個開放ge 社會機會俾所有人」.

$^{241}$「你行埋去得架啦」.
講埋呢個例子完啦.
我識一個弟兄當時後生.
咁呢就我地做佢導師ge .
不停咁樣追佢.
做咩呢?.
呢個弟兄好ye 喎.
所有聚會佢都黎到ge .
點解呀?.
點黎到呀?.
朝頭早去旅行.
八點鐘站合會又黎到喎.
咁晏晝有特別聚會又黎到.
唔使返工呀?.
夜晚黑友陪靈會又到.
「嘩,唔係呀?你唔使返工咩要呀?」.
「咁你返幾時呀?」.
「我返夜晚九點半ge 」.
「有時返十點ge 」.
「o下,點解咁呀?」.
原來佢係做夜更ge 保安.
咁如果夜更保安.
佢仲有喎.
佢係好好呢份工我都想做.
佢係守一個係偏遠地方無人到ge 配電站.
有一間大屋裡面就全部都係gwo D 變壓器.
由高壓電變做普通平民家庭用ge 電.
咁呢就守住個度唔好比人有搞震.
咁一個人架咋.
咁呢訓jor 無人知呀幾好呀.
所以佢朝頭早出黎.
旅行精神都好飽滿架.
咁呢就訓住賺呀.
咁佢呢就當時係咁ge .
咁話喂.
「弟兄,咁你成日做呢份ye 做到幾時呀?」.
「成世都做保安?」.
「咁係呀,我想做傳道人架」.
「o下,唔係呀?做傳道人?」.
「咁你要讀書得架」.

$^{281}$「咁呢起碼做傳道人都要中學畢業」.
「先入到個學士架喎」.
「咁你中學畢業未呀?」.
「未呀」.
「咁你點呀?」.
「讀書囉要」.
咁佢又話「唔係呀,好辛苦架喎」.
「咦,咦,咦」.
又做緊gwo 份工作.
好多年「喂,靚咪架」.
於是成日追佢終於呢.
佢就讀完jor 中學真係讀jor 神學喎.
咁你即係要知自己行先架.
你早D 唔使初陀歲月.
好多時候呢,不知到哪一天.
真係用好我地D 時間得著智慧ge 心.
數算好自己日子.
咁呢就耳索開始數喇.
知道自己隻眼就黎唔睜得開架喇.
唔,食一餐.
食埋個仔最後D 餐就祝福喇.
就係咁樣啦,幾好呀.
咁幫助我地有智慧數日子呀.
我地祈禱.
求主幫助我地有智慧日子數算好.
有智慧ge 心.
求主幫助我地有智慧ge 心數算自己日子.
比我地知道.
我地每一日都係珍貴ge .
都係你比ge .
讓我地有計劃去做我地ge 人生.
有時未必成功.
但我地總係要行埋去嘗試一下.
去完成神比我地ge 使命與工作.
去完成我地自己心裡面所嚮往ge .
求主比我地有行動.
先數算自己日子.
因為我地都唔知道哪一天就會離開.
願主祝福.
加利.

$^{321}$祈禱奉耶穌名求.
阿們.
\newpage



\section{創世記 1:1-3:24}
\label{sec:031771}
\textbf{女人的使命}~~~~林羅麗斯牧師
\newline
\newline
連結: \href{http://www.church.com.hk/acms/content.asp?site=cdc\&op=show\&type=product\&code=031771\&layout=sermon}{\texttt{031771}} ~~~~ 語音日期: 2026-03-08
\newline
\newline
\hyperref[sec:031770]{\small{< < < PREV SERMON < < <}}
~
\hyperref[sec:index_chronic]{\small{[返順時目]}}
~
\hyperref[sec:031772]{\small{> > > NEXT SERMON > > >}}
\newline
\newline
創世記 1:1-3:24
\newline
\begin{longtable}{cl}
\hline
\hline
章節 & 經文 (和合本修訂版)\\
\hline
1:1 & \begin{tabularx}{0.7\textwidth}{X} 起初，神創造天地。 \end{tabularx} \\ \\ \relax
1:2 & \begin{tabularx}{0.7\textwidth}{X} 地是空虛混沌，深淵上面一片黑暗；神的靈運行在水面上。 \end{tabularx} \\ \\ \relax
1:3 & \begin{tabularx}{0.7\textwidth}{X} 神說：「要有光」，就有了光。 \end{tabularx} \\ \\ \relax
1:4 & \begin{tabularx}{0.7\textwidth}{X} 神看光是好的，於是神就把光和暗分開。 \end{tabularx} \\ \\ \relax
1:5 & \begin{tabularx}{0.7\textwidth}{X} 神稱光為「晝」，稱暗為「夜」。有晚上，有早晨，這是第一日。 \end{tabularx} \\ \\ \relax
1:6 & \begin{tabularx}{0.7\textwidth}{X} 神說：「眾水之間要有穹蒼，把水和水分開。」 \end{tabularx} \\ \\ \relax
1:7 & \begin{tabularx}{0.7\textwidth}{X} 神就造了穹蒼，把穹蒼以下的水和穹蒼以上的水分開。事就這樣成了。 \end{tabularx} \\ \\ \relax
1:8 & \begin{tabularx}{0.7\textwidth}{X} 神稱穹蒼為「天」。有晚上，有早晨，這是第二日。 \end{tabularx} \\ \\ \relax
1:9 & \begin{tabularx}{0.7\textwidth}{X} 神說：「天下面的水要聚集在一處，使乾地露出來。」事就這樣成了。 \end{tabularx} \\ \\ \relax
1:10 & \begin{tabularx}{0.7\textwidth}{X} 神稱乾地為「地」，稱聚集在一起的水為「海」。神看為好的。 \end{tabularx} \\ \\ \relax
1:11 & \begin{tabularx}{0.7\textwidth}{X} 神說：「地要長出植物，就是含種子的五穀菜蔬，和會結果子、果子裡有種子的樹，在地上各從其類。」事就這樣成了。 \end{tabularx} \\ \\ \relax
1:12 & \begin{tabularx}{0.7\textwidth}{X} 於是地長出了植物：含種子的五穀菜蔬，各從其類；會結果子、果子裡有種子的樹，各從其類。神看為好的。 \end{tabularx} \\ \\ \relax
1:13 & \begin{tabularx}{0.7\textwidth}{X} 有晚上，有早晨，這是第三日。 \end{tabularx} \\ \\ \relax
1:14 & \begin{tabularx}{0.7\textwidth}{X} 神說：「天上要有光體來分晝夜，讓它們作記號，定季節、日子、年份， \end{tabularx} \\ \\ \relax
1:15 & \begin{tabularx}{0.7\textwidth}{X} 它們要在天空發光，照在地上。」事就這樣成了。 \end{tabularx} \\ \\ \relax
1:16 & \begin{tabularx}{0.7\textwidth}{X} 於是神造了兩個大光體，大的管晝，小的管夜，又造了星辰。 \end{tabularx} \\ \\ \relax
1:17 & \begin{tabularx}{0.7\textwidth}{X} 神把這些光體擺列在天空，照在地上， \end{tabularx} \\ \\ \relax
1:18 & \begin{tabularx}{0.7\textwidth}{X} 管理晝夜，分別光暗。神看為好的。 \end{tabularx} \\ \\ \relax
1:19 & \begin{tabularx}{0.7\textwidth}{X} 有晚上，有早晨，這是第四日。 \end{tabularx} \\ \\ \relax
1:20 & \begin{tabularx}{0.7\textwidth}{X} 神說：「水要滋生眾多有生命之物；要有鳥飛在地面以上，天空之中。」 \end{tabularx} \\ \\ \relax
1:21 & \begin{tabularx}{0.7\textwidth}{X} 神就創造了大魚和在水裡滋生的各樣活動的生物，各從其類，以及各樣有翅膀的鳥，各從其類。神看為好的。 \end{tabularx} \\ \\ \relax
1:22 & \begin{tabularx}{0.7\textwidth}{X} 神就賜福給這一切，說：「要繁殖增多，充滿在海的水裡；飛鳥也要在地上增多。」 \end{tabularx} \\ \\ \relax
1:23 & \begin{tabularx}{0.7\textwidth}{X} 有晚上，有早晨，這是第五日。 \end{tabularx} \\ \\ \relax
1:24 & \begin{tabularx}{0.7\textwidth}{X} 神說：「地要生出有生命之物，各從其類，就是牲畜、爬行動物、地上的走獸，各從其類。」事就這樣成了。 \end{tabularx} \\ \\ \relax
1:25 & \begin{tabularx}{0.7\textwidth}{X} 於是神造了地上的走獸，各從其類；牲畜，各從其類；和地上一切的爬行動物，各從其類。神看為好的。 \end{tabularx} \\ \\ \relax
1:26 & \begin{tabularx}{0.7\textwidth}{X} 神說：「我們要照著我們的形像，按著我們的樣式造人，使他們管理海裡的魚、天空的鳥、地上的牲畜和全地，以及地上爬的一切爬行動物。」 \end{tabularx} \\ \\ \relax
1:27 & \begin{tabularx}{0.7\textwidth}{X} 神就照著他的形像創造人，照著神的形像創造他們；他創造了他們，有男有女。 \end{tabularx} \\ \\ \relax
1:28 & \begin{tabularx}{0.7\textwidth}{X} 神賜福給他們，神對他們說：「要生養眾多，遍滿這地，治理它；要管理海裡的魚、天空的鳥和地上各樣活動的生物。」 \end{tabularx} \\ \\ \relax
1:29 & \begin{tabularx}{0.7\textwidth}{X} 神說：「看哪，我把全地一切含種子的五穀菜蔬和一切會結果子、果子裡有種子的樹，都賜給你們；這些都可作食物。 \end{tabularx} \\ \\ \relax
1:30 & \begin{tabularx}{0.7\textwidth}{X} 至於地上一切的走獸、天空一切的飛鳥，並一切在地上爬行的，有生命的動物，我把綠色植物賜給牠們作食物。」事就這樣成了。 \end{tabularx} \\ \\ \relax
1:31 & \begin{tabularx}{0.7\textwidth}{X} 神看一切所造的，看哪，都非常好。有晚上，有早晨，這是第六日。 \end{tabularx} \\ \\ \relax
2:1 & \begin{tabularx}{0.7\textwidth}{X} 天和地，以及萬象都完成了。 \end{tabularx} \\ \\ \relax
2:2 & \begin{tabularx}{0.7\textwidth}{X} 到第七日，神已經完成了造物之工，就在第七日安息了，歇了他所做一切的工。 \end{tabularx} \\ \\ \relax
2:3 & \begin{tabularx}{0.7\textwidth}{X} 神賜福給第七日，將它分別為聖，因為在這日，神安息了，歇了他所做一切創造的工。 \end{tabularx} \\ \\ \relax
2:4 & \begin{tabularx}{0.7\textwidth}{X} 這就是天地創造的來歷。在耶和華神造地和天的時候， \end{tabularx} \\ \\ \relax
2:5 & \begin{tabularx}{0.7\textwidth}{X} 地上還沒有田野的草木，田間的菜蔬還沒有長出來，因為耶和華神還沒有降雨在地上，也沒有人耕種土地。 \end{tabularx} \\ \\ \relax
2:6 & \begin{tabularx}{0.7\textwidth}{X} 但是，有霧氣從地上騰，滋潤整個土地的表面。 \end{tabularx} \\ \\ \relax
2:7 & \begin{tabularx}{0.7\textwidth}{X} 耶和華神用地上的塵土造人，將生命之氣吹進他的鼻孔，這人就成了有靈的活人。 \end{tabularx} \\ \\ \relax
2:8 & \begin{tabularx}{0.7\textwidth}{X} 耶和華神在東方的伊甸栽了一個園子，把所造的人安置在那裡。 \end{tabularx} \\ \\ \relax
2:9 & \begin{tabularx}{0.7\textwidth}{X} 耶和華神使各樣的樹從土地裡長出來，可以悅人的眼目，好作食物。園子當中有生命樹和知善惡的樹。 \end{tabularx} \\ \\ \relax
2:10 & \begin{tabularx}{0.7\textwidth}{X} 有一條河從伊甸流出來，滋潤那園子，從那裡分成四個源頭： \end{tabularx} \\ \\ \relax
2:11 & \begin{tabularx}{0.7\textwidth}{X} 第一條名叫比遜，它環繞哈腓拉全地，在那裡有金子。 \end{tabularx} \\ \\ \relax
2:12 & \begin{tabularx}{0.7\textwidth}{X} 那地的金子很好，在那裡也有珍珠和紅瑪瑙。 \end{tabularx} \\ \\ \relax
2:13 & \begin{tabularx}{0.7\textwidth}{X} 第二條河名叫基訓，它環繞古實全地。 \end{tabularx} \\ \\ \relax
2:14 & \begin{tabularx}{0.7\textwidth}{X} 第三條河名叫底格里斯，它流到亞述的東邊。第四條河就是幼發拉底。 \end{tabularx} \\ \\ \relax
2:15 & \begin{tabularx}{0.7\textwidth}{X} 耶和華神把那人安置在伊甸園，讓他耕耘看管。 \end{tabularx} \\ \\ \relax
2:16 & \begin{tabularx}{0.7\textwidth}{X} 耶和華神吩咐那人說：「園中各樣樹上所出的，你可以隨意吃， \end{tabularx} \\ \\ \relax
2:17 & \begin{tabularx}{0.7\textwidth}{X} 只是知善惡的樹所出的，你不可吃，因為你吃它的日子必定死！」 \end{tabularx} \\ \\ \relax
2:18 & \begin{tabularx}{0.7\textwidth}{X} 耶和華神說：「那人單獨一個不好，我要為他造一個配偶幫助他。」 \end{tabularx} \\ \\ \relax
2:19 & \begin{tabularx}{0.7\textwidth}{X} 耶和華神用泥土造了野地各樣的走獸和天空各樣的飛鳥，都帶到那人面前，看他叫甚麼。那人怎樣叫各樣的動物，那就是牠的名字。 \end{tabularx} \\ \\ \relax
2:20 & \begin{tabularx}{0.7\textwidth}{X} 那人就給一切牲畜、天空的飛鳥和野地各樣的走獸都起了名。只是亞當沒有找到配偶幫助他。 \end{tabularx} \\ \\ \relax
2:21 & \begin{tabularx}{0.7\textwidth}{X} 耶和華神使他沉睡，他就睡了；於是取下他的一根肋骨，又在原處把肉合起來。 \end{tabularx} \\ \\ \relax
2:22 & \begin{tabularx}{0.7\textwidth}{X} 耶和華神就用那人身上所取的肋骨造了一個女人，帶她到那人面前。 \end{tabularx} \\ \\ \relax
2:23 & \begin{tabularx}{0.7\textwidth}{X} 那人說：「這是我骨中的骨，肉中的肉，可以稱她為女人，因為她是從男人身上取出來的。」 \end{tabularx} \\ \\ \relax
2:24 & \begin{tabularx}{0.7\textwidth}{X} 因此，人要離開父母，與妻子結合，二人成為一體。 \end{tabularx} \\ \\ \relax
2:25 & \begin{tabularx}{0.7\textwidth}{X} 當時夫妻二人赤身露體，並不覺得羞恥。 \end{tabularx} \\ \\ \relax
3:1 & \begin{tabularx}{0.7\textwidth}{X} 耶和華神所造的，惟有蛇比田野一切的走獸更狡猾。蛇對女人說：「神豈是真說，你們不可吃園中任何樹上所出的嗎？」 \end{tabularx} \\ \\ \relax
3:2 & \begin{tabularx}{0.7\textwidth}{X} 女人對蛇說：「園中樹上的果子，我們都可以吃； \end{tabularx} \\ \\ \relax
3:3 & \begin{tabularx}{0.7\textwidth}{X} 只是園子中間那棵樹的果子，神曾說：『你們不可吃，也不可摸，免得你們死。』」 \end{tabularx} \\ \\ \relax
3:4 & \begin{tabularx}{0.7\textwidth}{X} 蛇對女人說：「你們不一定死； \end{tabularx} \\ \\ \relax
3:5 & \begin{tabularx}{0.7\textwidth}{X} 因為神知道，你們吃的日子眼睛就開了，你們就像神一樣知道善惡。」 \end{tabularx} \\ \\ \relax
3:6 & \begin{tabularx}{0.7\textwidth}{X} 於是女人見那棵樹好作食物，又悅人的眼目，那樹令人喜愛，能使人有智慧，她就摘下果子吃了，又給了與她一起的丈夫，他也吃了。 \end{tabularx} \\ \\ \relax
3:7 & \begin{tabularx}{0.7\textwidth}{X} 他們二人的眼睛就開了，知道自己赤身露體，就編織無花果樹的葉子，為自己做成裙子。 \end{tabularx} \\ \\ \relax
3:8 & \begin{tabularx}{0.7\textwidth}{X} 天起了涼風，那人和他妻子聽見耶和華神在園中來回行走的聲音，就藏在園裡的樹木中，躲避耶和華神的面。 \end{tabularx} \\ \\ \relax
3:9 & \begin{tabularx}{0.7\textwidth}{X} 耶和華神呼喚那人，對他說：「你在哪裡？」 \end{tabularx} \\ \\ \relax
3:10 & \begin{tabularx}{0.7\textwidth}{X} 他說：「我在園中聽見你的聲音，我就害怕；因為我赤身露體，我就藏了起來。」 \end{tabularx} \\ \\ \relax
3:11 & \begin{tabularx}{0.7\textwidth}{X} 耶和華神說：「誰告訴你，你是赤身露體呢？莫非你吃了那樹上所出的，就是我吩咐你不可吃的嗎？」 \end{tabularx} \\ \\ \relax
3:12 & \begin{tabularx}{0.7\textwidth}{X} 那人說：「你賜給我、與我一起的女人，是她把那樹上所出的給我，我就吃了。」 \end{tabularx} \\ \\ \relax
3:13 & \begin{tabularx}{0.7\textwidth}{X} 耶和華神對女人說：「你怎麼會做這種事呢？」女人說：「那蛇引誘我，我就吃了。」 \end{tabularx} \\ \\ \relax
3:14 & \begin{tabularx}{0.7\textwidth}{X} 耶和華神對蛇說：「你既做了這事，就必受詛咒，比一切的牲畜和野獸更重。你必用肚子行走，終生吃土。 \end{tabularx} \\ \\ \relax
3:15 & \begin{tabularx}{0.7\textwidth}{X} 我要使你和女人彼此為仇，你的後裔和女人的後裔也彼此為仇。他要傷你的頭，你要傷他的腳跟。 」 \end{tabularx} \\ \\ \relax
3:16 & \begin{tabularx}{0.7\textwidth}{X} 又對女人說：「我必多多加增你懷胎的痛苦，你生兒女時必多受痛苦。你必戀慕你丈夫，他必管轄你。」 \end{tabularx} \\ \\ \relax
3:17 & \begin{tabularx}{0.7\textwidth}{X} 又對亞當說：「你既聽從你妻子的話，吃了那樹上所出的，就是我吩咐你不可吃的，土地必因你的緣故受詛咒；你必終生勞苦才能從土地得吃的。 \end{tabularx} \\ \\ \relax
3:18 & \begin{tabularx}{0.7\textwidth}{X} 土地必給你長出荊棘和蒺藜來；你也要吃田間的五穀菜蔬。 \end{tabularx} \\ \\ \relax
3:19 & \begin{tabularx}{0.7\textwidth}{X} 你必汗流滿面才有食物可吃，直到你歸了土地，因為你是從土地而出的。你本是塵土，仍要歸回塵土。」 \end{tabularx} \\ \\ \relax
3:20 & \begin{tabularx}{0.7\textwidth}{X} 那人給他妻子起名叫夏娃，因為她是眾生之母。 \end{tabularx} \\ \\ \relax
3:21 & \begin{tabularx}{0.7\textwidth}{X} 耶和華神用獸皮做衣服給亞當和他的妻子穿。 \end{tabularx} \\ \\ \relax
3:22 & \begin{tabularx}{0.7\textwidth}{X} 耶和華神說：「看哪，那人已經像我們中間的一個，知道善惡，現在恐怕他又伸手摘生命樹所出的來吃，就永遠活著。」 \end{tabularx} \\ \\ \relax
3:23 & \begin{tabularx}{0.7\textwidth}{X} 耶和華神就驅逐他出伊甸園，使他耕種土地，他原是從土地裡被取出來的。 \end{tabularx} \\ \\ \relax
3:24 & \begin{tabularx}{0.7\textwidth}{X} 耶和華神把那人趕出去，就在伊甸園東邊安設基路伯和發出火焰轉動的劍，把守生命樹的道路。 \end{tabularx} \\ \\
[1ex]
\hline
\hline
\end{longtable}
$^{1}$今天是一個特別的節日.
弟兄姊妹知不知道.
今天正正是三八.
和旁邊的姐妹祝她.
婦女節快樂.
我們今天講婦女節很開心.
通常是想在女人的世界.
諸如此類的.
其實三八的婦女節.
本身不是一個輕鬆的節日.
都是經過很長時間的鬥爭.
我們叫做兩性的鬥爭.
才成立了一個這樣的節日.
來慶賀一些成功.
成功什麼呢.
是成功爭取了很多婦女的權益.
是二十世紀這段時間發生的.
我們在香港的婦女.
她們可能不是很知道.
原來我們已經享受了很多.
我們之前上一代的女性.
是沒有的福利.
其實為什麼兩性的問題.
會成為一個社會的議題呢.
如果大家有些讀開歷史.
或者有些興趣看神話的都知道.
很多原始的社會.
本身沒有是社會的.
也就是說都是女人說了算.
原因就是唯有女人很肯定.
哪個小孩是自己生的.
因為那時候沒有DNA.
驗不到的.
如果看那些劇集介紹的.
就是那些男人排在窗口.
第二天搞完就走了.
根本沒有人確定.
哪個小孩真是屬於自己的.
因此沒有是社會.
因為小朋友是一定認得知道.

$^{41}$或者是女人會知道.
哪個是自己的子女.
在那個時候.
也沒有特別是需要用到機械.
勞動力來做生產的.
都是撿到魚.
撿到蘋果這樣喫的.
很原始的社會.
但是去到已經開始有些禮制.
制度的時間.
就發現男人的體力差很遠.
打獵.
去耕田.
開始覺得因為財富來源的轉移.
慢慢慢慢.
特別是中國人所定下來的.
所有的婚姻制度.
都似乎已經是男尊女卑的時代了.
因此在犯罪之後的世界.
男女好像永遠都在角力和增進裡面.
所以出現了世界各地.
特別是在基督教一直傳開去的時候.
更多的女性就會起來.
來爭取她們自己的所謂的權利與福利.
這個浪潮由18世紀.
原來是打到20世紀.
纔開始見到我們現在的初型.
其實你問我自己在60年代出生的.
我們當年都很清晰.
男女其實是不平等的.
經歷了整個會議都聽過.
我們講見證就知道了.
如果能夠真是有書讀的.
男女都可以平均有這個機會去讀書.
其實不是很多年的事.
當然如果你問下面的項目.
他們都不知道多想女生不用讀書.
另外就是工作機會能夠同等.
也都不容易的.
很多時候都是先請男性.

$^{81}$或者是差一點的工作纔去想女生的.
如果去到真的同工同酬.
某程度上我小時候.
那個時候還是一個爭鬥的時候.
不知道大家.
原來籤默都是挺厲害的.
我們都有很多的創舉.
包括我們是全港第一對的夫婦.
教母同工是同職同工.
這樣來入職的.
他們也都有很多壓力的.
因為特別是贈信會.
兩夫婦一起投票怎麼辦.
兩個是一票.
諸如此類很多的.
其實是在一些不公平的狀態裡面的.
現在here and now.
你們看教會的服事.
教會的狀態都不是完全的.
以前很多時候.
當我們在夫婦團契.
夫婦小組的那些年代.
大部分的姐妹是自動讓位的.
就是尊敬我們的丈夫.
因為上上都希望丈夫做個頭的.
所以如果你問女的傳道人要做牧師.
甚至要做主任牧師.
某些教會在現在還是在爭鬥的.
我們都很開放.
主任牧師的位置.
我們是由女士來處理的.
但是整個社會運動.
當然是包括了很多公民的權利.
投票的權利.
甚至能不能夠去領受遺產.
你知道自己給兒子給女兒.
現在我們寫的遺囑是很簡單的.
但是在香港還有一個不平等的.
我們的丁屋是什麼的.
全男全女的.

$^{121}$事實上是的.
還有一樣東西.
可能大家沒有介意.
也沒有想過的.
有沒有想過你生了個子女.
你跟誰姓.
是不是.
就算離了婚.
很多都還是跟著爸爸姓的.
所以有時候我們在馬來芝加那些一修.
結了一次婚.
我們經常以為那些同工是執錯資料的.
搞錯了.
三個孩子三個不同姓的.
原來還是跟著爸爸姓的.
在中國比我們香港更加先進.
現在有個新興的做的狀態.
因為長年的獨子.
長年的獨女.
你明白嗎.
當大家結婚的時候.
祖宗十八代都是靠你的.
那怎麼辦呢.
現在有個狀態就是.
獨子娶了個獨女.
然後他們就一定會生育眾多.
第一個就是跟著爸爸姓的.
第二個是跟著媽媽姓的.
第三個是跟著爸爸姓的.
如此類推.
就確保兩邊都有子子的延續.
所以其實他們真的挺先進的.
因此如果你看到在懷崇拜.
可能他們兩兄弟來到.
原來他們是不同姓的.
你不要以為他們離婚.
他們可能都是父母相愛.
但是生了兩個姓的小孩出現的.
因此在整個社會里面.
人類有出現的.

$^{161}$兩性的競爭.
其實是從來沒有停止過的.
我相信在你的家裡.
你都會感覺到這種氣氛.
但是在兩性裡面應該怎麼去相處呢.
特別今天這個訪談是講明白神的旨意.
而這麼難得在婦女節的時候.
因此你就知道為什麼我會定題目.
就是講女人的使命.
先小人後君子.
我們現在講聖經裡面的創造美好的原則.
大家聽完之後回家就不要開火了.
這個事實上是真的.
在聖經裡面.
他講得很清楚.
神是按照自己的形象來做男做女的.
因此我希望信主的弟兄姊妹.
我們都應該從心底裡面真的尊重.
以及知道我們無論是男生還是女生.
都是有神的形象而做的.
所以裡面是有尊貴的.
但是某程度來說.
我們發覺男女真的是不同的.
我們仔細去看.
你就知道整個創造的過程是不同的.
也因為在創造的過程裡面.
我們很清楚找到.
男人經常都會問神為什麼要做人.
這個可能要問上帝了.
但是很清楚的.
女人是從哪裡來的呢.
女人為什麼會出現的呢.
其實在創世記裡面是講得很清楚的.
不會混淆的.
希望姐妹們也深呼吸一口.
去接受創造的命定.
我們來看三聖經文.
我們很簡單的看一看.
在上世紀的第一章.
就講了上帝怎麼去創造這個世界的次序.

$^{201}$大家唱歌都唱得多.
如果懂得唱歌就知道.
第六天是做了地上的走獸和人類.
七天就完成了上帝的創造.
七天就完成了上帝的創造的工作.
但是在第一章的尾段很清楚的講了.
上帝是照自己的形象.
做男做女之外.
上帝是定義要賜福給他們的.
所以上帝去創造人類.
本來是享受上帝為他們預備的所有福氣的.
並且賦予人類一個獨特的任務.
第一就是什麼.
生養眾多.
遍滿全地.
這個對很多女人來說好像不太合適.
另外第二就是要什麼.
除了生育之外.
目的也是為了什麼.
要治理大地.
治理海里的魚.
所有上帝所創造的東西.
其實上帝就挑選了人類.
來做一個管理員.
去打理所有的東西.
也是超級的.
問題是.
第一章這麼簡單的去講.
這個做男做女.
其實是一個過程.
唯有在第二章才補充了.
做男做女的過程發生了什麼事.
經文很清楚.
上帝做了人類的時候.
是將一個塵土.
吹了他自己的靈氣進去.
神是靈.
神將自己吹進了你和我當中.
在人類當中.
他就成為了一個有靈的活人.

$^{241}$有靈的活人.
不是男女.
是一個有靈的活人.
但是怎麼會發展到會出現呢.
聖經很仔細地去講.
相信這個有靈的活人.
是在幫助管理大地.
也與神同行.
因為上帝和他好像是朋友一樣.
問題就是.
在這段時間裡面.
在15至17節.
神說了.
神將他安置在伊甸園裡面.
也命令他要做修理和看守.
然後對他有了吩咐.
吩咐就是樹上所有的果子.
是可以喫的.
不過留下了第一條命令.
就是唯一的命令.
他不可以喫一棵叫做.
分別善惡樹上的果子.
講明喫了之後的後果.
上帝從來在聖經的教導.
律例典章.
都是講得很清楚的.
所以讀聖經讀清楚一點.
你不可以喫.
你喫了的時候.
你是必定死.
OK.
講得很清楚.
你相信這個有靈的活人.
他應該和上帝講得很清楚.
時間過去.
在第18節的時候.
就發現一件事.
上帝就發現.
那人獨居不好.
我都是要為他.

$^{281}$做一個配偶.
來幫助他.
又記得在創世記的時候.
每逢上帝創造完.
特別是天空飛鳥.
諸如此類.
他都有一句話.
經常是一樣的.
就是「各從其類」.
所以相信這段時間.
我們會看到.
昆蟲,飛鳥,魚兒,走獸.
都在各從其類.
也就是說他們可能有生育.
纔會變多.
唯獨這個有靈的活人獨居.
所以上帝就看到.
那個人獨居是不好的.
於是就為他.
打算要做一個配偶.
所以這個女人的出現.
也是挺特別的.
這個出現是一個.
在創造的過程裡的尾段.
覺得獨居不好.
所以而發生的.
然後去到經文的22-23節.
也都將那個過程.
很仔細地講了出來.
她是在那個人身上的肋骨.
拿出來做了一個女人.
而且做了這個女人.
是上帝親自.
將她帶到這個男人的面前.
就這麼分了.
然後這個男人.
就講了一句第一句的說話.
就是人第一句說的話.
他說什麼.
這是我骨中的骨.

$^{321}$肉中的肉.
可以稱她為女人.
因為她是從男人身上.
取出來的.
那你就發現.
很難聽的.
同意嗎.
很浪漫,一點都不浪漫.
特別我們那個年代.
有多少人聽過.
我們上一代在介紹.
這個是我的女人.
有沒有聽過.
有啦,是不是.
如果好聽.
現在都會這麼用.
就是當時的女人不喜歡聽.
我起碼會說.
我女人這樣.
所以後來才發現.
現在說什麼.
我的太太.
我的老婆.
如果是說英文.
我的Honey.
總之你會覺得.
這是我的女人.
這個是我女人的感覺.
反而不好的.
不過你又想一下.
阿當.
這個有靈的活人.
他活了不知一段多長時間.
他一直在眼見的.
從來都沒見過人.
我養了一隻狗之後.
養了第二隻狗之後就很明白.
我第一隻狗叫栗子.
他們跟我一起生活了大半年.
他只見過我們.

$^{361}$他沒見過其他狗.
因為他太小.
我們很少帶他去街.
牛要去街.
發現他很有趣.
他見到狗.
跟狗是好玩的.
不是的.
他見到狗很惡的.
不停吠.
都不知發生什麼事.
見到人.
人人都怕人.
人家當然怕他.
所以我們很少帶他出街.
有時冬天用沙發蓋著自己.
看電視的時候.
這麼靜.
去哪了,栗子.
他跟我們一樣.
鼠咒在被窩裡.
跟我們一樣蓋著被子.
不知他在幹嘛.
看著看著.
我懷疑.
他應該不知自己是隻狗.
直至我們買了樽子出現的時候.
我們就更加肯定這個想法.
當然他見到他不停吠.
但是慢慢慢慢的時候.
他應該有視覺看到.
是兩類東西.
不是同比例的.
你明白嗎.
然後他們發覺他們兩個同喫一樣東西.
我們在桌上喫東西.
慢慢慢慢.
他相信已經知道.
他跟我們是不同類.
所以你們想想.

$^{401}$當見到夏娃的時候.
他那種驚喜.
是想到的.
哇.
有一隻這樣的東西.
你們明白嗎.
不是豹,不是獅子.
而是他明明的.
這一隻是什麼呢.
但是他很有好感.
很有歸屬感.
因為他們很相似.
因此這一首我們覺得.
唉,很辱華的詩篇.
你在什麼狀態去想.
這個是從男人心裡一見到他的時候.
就說這是我骨中的骨.
是我肉中的肉.
於是他為他定名.
就叫做女人.
女人這個字.
是絕對有那種歸屬感.
他明明是斑馬.
明明是獅子老虎.
但他們不是他的.
但是這個女人.
是男人從心底裡面說.
這個女人是我的.
你明不明白那種情懷.
男人看著我好像不知道發生了什麼事.
回到你追求你的時候.
第一眼見到他的時候.
你有這種感覺.
哇,這個就是我的了.
那種感覺.
所以那種歸屬.
其實是在創造裡面.
上帝是做男做女.
一個很巧妙和很奇妙的建築.
因為令到兩性.

$^{441}$有一個相互依附的關係.
我需要你.
你也需要我.
我願意歸屬於你.
你也願意歸屬於我.
這個就是做男做女.
也是上帝定下.
這個夫妻婚姻關係裡面.
一個最重要的元素.
就是這對男女.
他們應該那份相愛.
不只是我喜歡你對我好.
你喜歡哄我或者你遷就我.
而是我們彼此需要彼此相熟.
當這樣的結合的時候.
纔是真正的合一.
然後目的就是為了在生理上能夠生育下一代.
並且能夠同心合意的.
來管理上帝給他的大地.
在這個關係裡面.
原本男人是真的從心底裡擁有一個女人.
這個擁有是代表了.
不只是我愛她.
而是我對她所付出的承諾與責任.
如果你問弟兄.
開車的.
他怎麼對車.
他就明白了.
他們知道車是他老婆的.
如果你愛惜你的車.
等於愛惜你老婆.
就是類似的情懷.
你說我不開車.
你總有一樣東西.
你的手機.
你的相機.
或者你玩那些機械人.
那些奇怪的東西.
總之那個是你所愛你所屬.
你就會對他很厲害.

$^{481}$你不可以隨便借給別人用.
你很緊張他是否安好.
這個就是在男女被創造的時候.
男士很正常地對著他的女人.
應有的那種關懷與情懷.
並且是那種衛生和責任.
同樣女人呢.
女人無端端帶著他去一個物體的前面.
不知道是什麼.
然後那個人就說你是屬於他.
你看到夏娃是沒有反對的.
那種是我願意歸屬於你.
也是在上帝創造的裡面定下的方程式或者密碼.
所以我相信每一個結婚.
或者嫁給他的時候.
你決定的時候的女士.
你都知道你是願意將你自己的一生.
你都是毀身於對方的.
你根本就不會介意他是你的老公.
你也不介意灌呼成.
當你決定結婚.
你就想著你自己叫什麼太太.
這個是願意好像是低一點.
但其實你不覺得是低的.
你覺得是被愛.
終於我有個叫歸宿了.
而這種美麗的情懷是重要的.
因為這樣才能夠建立一個真正相愛的家庭.
然後下一代纔會明白什麼叫做愛.
在愛裡面除了被愛之外.
你還知道什麼叫負責任.
很多青年人就說我不想為了負責任而去愛.
錯了.
什麼叫愛.
愛就是負責任.
如果你不負責任的話.
你不要說愛字.
你根本談不上什麼叫愛.
所以在這樣的氛圍裡面.
本來人類應該是在一個相愛的裡面.

$^{521}$互相毀身的裡面.
你都是去快快樂樂地生活下去.
當然事情不是這麼開心.
去到第三章的時候.
我們就見到人類犯罪.
而引致到現在我們活在這個有罪的世界.
兩性為什麼會有這麼多的鬥爭.
但是當我們這次再去第三章的時候.
不好意思是婦女節.
試一下從婦女的角度去挑戰一下.
過往一直覺得都是你的女人.
最麻煩就是女人又貪美又貪喫.
搞到喫了禁果.
你要知道是夏娃喫了禁果的.
是她抵不住這個蛇的引誘.
但是我很想大家去想一下.
究竟是哪一個聽上帝的命令.
園中所有果子是不可以喫.
喫了的時候你必定死.
是誰?.
是亞當.
當大家變成有兩個的時候.
究竟亞當如何地詮釋或者教導這件事.
那我們就會存疑了.
因為所有的女士心中就有數.
往往可能你老公說的話.
在我們聽的來說.
就是說一些不說一些.
你明不明白.
於是往往老公第一次.
這些很溫柔很聽話.
大部分女生我相信.
老公說什麼就跟著他做什麼.
但是當你發覺跟了他就會死.
的時候你就會發現.
那些女士或者媽媽就越來越兇了.
是不是真的這樣.
你不要說清楚一點.
或者是不用你說.
我做完它.

$^{561}$我們相信在家常常見到這些場面.
夏娃究竟多理解這個禁果的事.
以至到蛇很輕易地挑戰了她.
你不一定死.
和必定死又不同了.
第二.
她挑動了原來你不喫禁果的目的.
是背後可能上帝有些好東西.
不想你們知道.
於是女人就真的耳朵軟.
被人挑戰一下她就拿了來喫.
跟著.
這個不重要.
這個第一步.
第二步.
她愛她老公的.
我們也在想管理大地.
為什麼夏娃一個人去做事.
噹噹整天睡覺.
在吞泡.
但無論如何也好.
她也愛老公的.
她摘了喫了.
咦 正啊.
好喫啊.
於是她就拿回去給老公喫.
你看到那個好意嗎.
這裡是一個很大的機會.
弟兄.
如果今天你老婆突然跟你說.
哎呀 死了.
我喫了麵包喫了老鼠藥.
你會怎樣.
第一件事.
不是罵她吧.
沒有沒有.
第一件事你可能喝水.
跟著可能就帶她去急診室.
如果她已經九八抱.
你已經召救護車.

$^{601}$你應該很緊張很擔心.
整件事.
聖經不是這樣看的.
因為你老婆喫了禁果.
你知道上帝跟你說她是必定死的.
就是說你竟然沒有理這件事.
然後夠膽死就是.
她給你喫你又喫一份.
所以你問我.
無論夏娃是聽了還是不聽了也好.
或者阿當自己也不是很立意要進行上帝這個命令.
於是結果她也喫了.
喫了之後我覺得身體一定有變化.
她一定要知道自己有大麻煩.
如果不是.
當上帝再找她的時候.
她沒理由會躲起來.
或者不敢見上帝.
在這件事裡面.
我也想強調我們所信的上帝.
雖然想法很分明.
但是上帝的慈愛往往是超過我們所想所求的.
你想想回到家裡.
看到一灘水在那裡.
玻璃又打爛了.
媽媽進屋的時候.
你就知道火都來了.
當然是立刻問誰打爛水.
你一定是爆發的.
但是上帝既然是上帝.
他怎麼會不知道這一對男女.
已經違反了他的命令.
喫了禁果呢.
他一定知道.
但聖經沒有說到上帝一來.
就是犯賊.
罵人的.
上帝仍然在找.
阿當阿當你在哪裡.
然後.

$^{641}$阿當阿娃雖雖格格.
害怕害羞害羞地走出來的時候.
神都沒有罵他.
神給了他一個很大的機會.
他問.
為什麼你會避開我.
如果你是聰明的小朋友.
你又相信你媽媽是一個.
有教養又有愛的媽媽.
在那一刻你還不識做.
你當然是自己扯自己的兒子.
或者是乍哭.
不好意思媽媽打爛了你的東西.
什麼類似的.
多生氣的媽媽.
都要保持自己的教養.
都會放一口氣.
都好心好氣跟女兒說.
下次不要這樣.
弄傷手之類.
但是很可惜.
她們看不到上帝的愛.
又或者我不知道.
她們對上帝的認識有多少.
結果纔是我們兩性.
出現了衝突的主因.
就是當上帝問的時候.
阿當這個負責人.
他是老闆.
夏娃是助手.
親自示範.
他說什麼.
都是你做給我的女人.
搞成這樣.
你很自然.
你見到夏娃就跟了.
就是你做的那條蛇.
搞成這樣.
自此你就知道.
我們人類世界.

$^{681}$為什麼進入了.
現在這樣的兩性衝突裡面.
就算我們所有福利.
什麼什麼平權都沒用.
回到家裡這個場景.
是經常發生的.
所以罪的恐怖性.
是令到神和人的關係.
受到破壞之外.
其實人與人之間.
就再沒有歸屬負責任.
有什麼事我承擔.
沒了肩膀.
兄弟姐妹.
在這個這麼恐怖.
或者這麼淒涼的反面之後.
人類當然就進入了犯罪的過程.
也都是被趕出這個伊甸園.
但是在當中整個的狀況.
兩性這麼多代一代一代的詮釋.
我相信我們現在香港的女士.
在我們這個年代.
可能六七十年代.
你還會聽到你媽媽.
在你臨嫁的時候會說些什麼.
我就嚴重一些.
我其實在讀書的時間.
我媽媽常常都會這麼說的.
就是男人靠得住.
豬奶都會上樹.
所以女兒.
你讀書一定要讀得好.
你一定要靠自己.
女人不要再靠老公了.
不然將來你就黑食了.
所以某程度.
你就會明白為什麼.
這一代的女子.
為什麼這麼努力讀書.
這麼強烈.

$^{721}$這麼想說話.
因為某程度上一代的咒詛.
一代一代傳到現在.
因此弟兄.
我想提醒你.
在整個創造的故事上.
上帝真的立了你們為家庭的頭.
不能走的.
夫妻關係在上帝的定義裡面.
你們是在前線的.
但是這個權並賦予了.
那種歸屬的感覺.
你的女人是屬於你.
她做錯事.
都是關你事的.
這個做前線就是有這個責任的.
所以下面怎麼背鍋的.
就是上面那個老鼎要辭職的.
因此.
在屬靈的事裡面.
彼得更加提醒我們.
他說.
你要知道女人是軟弱的.
比你男人軟弱.
但是.
她是跟你一起承受生命之恩的.
如果你不敬重她.
或者你們的關係在一個糾結的裡面.
你們的禱告都會有難了.
哇.
挺嚴重的.
原來在屬靈的裡面.
不是你一個聖潔.
你一個愛主就足夠的.
而是你們是二人成為一體的.
你另外一半的屬靈的狀況.
他能不能跟你同心.
是足以影響你的禱告.
有沒有被上帝來應允的.
因此.

$^{761}$為了自己的好處.
為了他的好處.
我們姐妹和旁邊的弟兄說.
請你做一個負責任的男人.
做得.
就是阿頭呢.
就要擔得起這個阿頭的責任.
但是在於我們女人又怎麼看呢.
因為今天講女人嘛.
在聖經裡面.
有一個這麼明顯的故事.
令到我們不得不承認.
女人是從哪裡來的.
女人真的是從男人那裡來的.
當然我們創造是上帝.
但是那個設計的過程是這樣的.
女人有沒有目的呢.
女人一出生就有目的了.
因為女人的被造.
是為了幫助你的男人的.
OK.
這個是最重要的.
因為你要幫助她一起生育下一代.
並且承擔起她的責任.
就是這個人類的責任.
就是治理大地.
當然.
現在今時今日在這個世界裡面.
我們發覺整個世界不可以走的.
這樣會死人的.
先不要不可以的.
OK.
背後是因為這個犯罪.
引來除了兩性的角力之外.
也帶來了咒詛.
和一個審判.
一個惡果.
對男人來說.
你要知道就是.
你需要終身勞苦.

$^{801}$才能夠在地裡面得到.
喫得飽的.
所以男人.
一定要肩負起這個責任.
就是去讓你的家人.
因為你的勞動.
為你的努力.
起碼能夠喫得飽.
而女人的受阻.
是雙倍的.
因為可能她真的先喫了.
上得很公平的.
神說她會加添我們懷胎的苦楚.
我們生產兒女是必受苦楚.
這個原來生孩子會痛.
其實是.
因為犯罪.
原來沒有犯罪.
其實生孩子是不痛的.
真的挺棒的.
但是怎麼辦呢.
我們知罪就信了主.
耶穌十字架.
赦免我們.
我們恢復了.
難道我們也要這樣捱嗎.
不是的.
這一句經文.
是女人都要教的.
有女的就要提醒她.
如果女人常傳信心愛心.
又聖潔自守.
就必在生產上得救.
很重要的.
這個對於女人的應許.
是解藥.
在毒液裡.
這個是解藥.
但是第二個惡果就糟糕了.
她說你必怨惡你的丈夫.

$^{841}$你的丈夫必管轄你.
這裡不關男人的事.
你要記住.
無論你怨惡什麼.
那樣東西都管轄你.
你明白嗎.
喫白粉的.
他沒辦法控制.
是白粉控制他的.
如果你很愛一個人.
別說你老婆.
可能你子女.
她的子女根本就是控制你.
她任何的表現.
都會觸動到你的情緒.
所以當你怨惡一個人.
意思是什麼.
如果你所有滿足的源頭.
都是歸功於他.
除他以外.
別無拯救.
那你就知道.
你一切的苦楚.
都將會由他而來.
我們叫這個依附關係.
一生我們.
從出生到現在都要處理的.
就是這一種依附關係.
帶來的所謂的苦果.
當你被他管轄的意思是什麼.
你再沒有自由了.
你見到很多在青春期.
青年期特別現在.
他們雖然說被父母愛惜.
但不知為什麼他們很缺乏愛.
很坦白.
今天有很多家長.
不好意思.
我們又要有一個智慧和聰明.
我們要禱告.

$^{881}$為什麼我們.
僕心僕命愛我們的子女.
但是他們這麼缺乏愛.
以至於他們的拍拖.
是很快的.
很快拍拖.
然後叫他分手之後.
一年後再拍拖.
是沒有人理會的.
他們沒有這個能耐.
簡單來說.
他隨便見到有另一個的時候.
他就會跳下去.
是因為他裡面.
好像吸毒一樣.
他是很需要這件事.
所以他沒有自由.
不能再選擇.
沒有理性.
隨便見到他需要一件事就拿.
因此在這個過程裡面.
令到很多婚姻失敗.
失戀.
其實對男生也是有禍害的.
你明白嗎.
如果是女生撲過來.
不是多少個男人可以頂得住.
但往往錯誤的婚姻.
就會帶來很多傷害.
於是男女之間在拍拖這件事.
都有很多的苦楚.
現在在疫情之後.
已經有很多人不結婚了.
不結婚而已.
不要搞錯 大英姐妹.
他沒有說過他們沒有男女關係.
只是沒有承諾.
沒有承諾究竟是瀟灑還是眼淚.
也就是說你都不敢要求對方.
為你有任何的承諾.

$^{921}$你明白那個狀態和心情嗎.
其實你問.
這並不是一個幸福的感覺.
但他們很流行自我安慰.
所以現在有個詞大家要學一下.
叫做單身優越.
什麼意思.
以前新年最怕去拜年.
你還沒結婚.
有人問你什麼時候嫁.
十年了你還問.
所以很多人很不喜歡去拜年.
但現在反過來.
他們大家催眠自己.
單身優越.
結了婚又離婚.
離婚又單身.
大家都是這樣的心態.
令自己好像更加堅強.
但是不是真的單身會優越一些呢.
如果你問.
現在是講聖經.
創世本身的原意.
是讓男和女成為一體.
以至有上帝賜給你的福氣.
但我們不得不承認.
我們現在面對很多單身的女性.
包括她找不到她值得偽身的男人.
第二.
可能她已經在婚姻的失敗裡.
第三.
可能.
多恩愛也好.
也會有一個先走的.
所以也有在這樣的狀況下被逼單身.
但我覺得在正常之下.
既然你要知道女人的使命是為了什麼.
幫助你的男人.
如果你發現你沒有男人出現的時候.
你有什麼心情呢.

$^{961}$我和很多單身的姐妹.
談這個問題.
當然她們說什麼.
都是抒發不了她們的情懷.
她們總是說師母你怎麼和我離婚了.
我也是離婚了.
但一個過來人.
嫁了的人.
我很想和你說.
如果你每一件事在主的帶領裡面.
但不知為什麼.
你覺得上帝硬是不把你的Mr.Right.
不把你的亞當.
擺在你面前.
我只是衷心說一句.
恭喜你.
因為上帝沒有派工作給你做.
姐妹們.
你有沒有從這個角度想.
回到聖經.
本來你創造出來是為了幫助你的男人.
這是你的使命和工作.
竟然上帝那麼愛你.
不用你做.
保羅說.
如果你不需要的時候.
你就可以單單為主而活.
何等美好呢.
你不需要和他商量.
也不需要牽手.
而且在現在這個時代.
實際上女性有自己的生存能力.
不需要像以前那樣.
一定要嫁了才能喫飽.
現在自己可以餵飽自己.
然後我們特別在香港.
很方便.
約人出來玩也是很好的.
所以你見到.
很多場合你都會見到這些情景.

$^{1001}$無論什麼年紀都是一群女人.
嘻嘻哈哈 喫喝玩樂.
其實這件事也是值得我們深思的.
每一個人總有機會是什麼.
有單身的一天.
如果我們太過依附一個男士.
來生存.
也是我們生命的Bug.
是要我們去面對的.
和我相熟的戰友.
由那麼年輕一起事奉的.
都很知道的.
我用了很多年時間.
去學習獨活的樂趣.
因為從小很多兄弟姐妹在一起.
也是父母的家庭.
結婚後有個老公是基督徒.
所以有很多夢想都在他身上.
也很想家裡永遠像24小時在一起.
但是不是.
原來上帝最終要我們學習.
好好愛自己.
好好預備自己.
也能夠享受獨活的樂趣.
如果你能夠這樣.
你才知道什麼叫切實彼此相愛.
兄弟姐妹在小組也是.
一定要組長幫你.
一定要組長提醒你.
一定要組長照顧你才上堂.
你未懂得獨立的.
你還是去生活.
你享受不了這份與主同行的樂趣.
你只是與你的組長同行.
這不是健康的.
我們人多麼相愛也好.
也不可以一起斷氣的.
所以所有的兄弟姐妹.
我們開始有些.
大家都是路由卡的.

$^{1041}$我們真的要預備的.
我們總有單身生活的一天.
我們能夠成為一個獨活的人的時候.
反而我們能夠彼此祝福.
每一個是男的.
是女的.
神都有在生命當中各自的美意.
上帝要我們去幫助男人.
不等於沒有你了姐妹.
你可以仍然有你的天下.
因為上帝仍然會使用你的.
但是如果夫妻都能夠獨活的時候.
我們就能夠盛載上帝給我們的使命.
然後我們又能夠真的靠著主.
學習彼此相愛.
彼此相屬.
在這個時間或一起或分開.
你都會經歷到上帝在我們的生命當中.
所為我們預備一切的福祿.
因此我們要留意保羅的提醒.
他說如果你們現在是有妻子.
或者老公纏著你的.
特別是退休的男人.
總是纏著老婆的.
保羅就說.
你不要求解脫.
你一定要繼續照顧她.
因為她需要你.
但是這麼幸運.
你沒有老婆.
或者沒有老公.
需要纏著你的時候.
他說什麼.
那就不要強求了.
你沒有空嗎.
那我就安排我自己的生活.
你明白嗎.
鼎姐妹當你這樣做的時候.
我們就學習在神的面前.
只是做一個忠心的牧人.

$^{1081}$去領受上帝所委派在我們自己生命當中的使命.
你有你的使命.
你的老公或者老婆都有他的使命.
你的子女都有他們自己的使命.
但當我們在一起的時候.
我們一家人.
我們有一家人的使命.
我們在同一間教會.
我們有同一間教會的使命.
當我們能夠明白的時候.
結婚就不是一件很恐怖的事.
或者單身.
都可能是一件很有樂趣的事.
但是當我們已經是有一半的鼎姐妹.
我在這裡鼓勵你.
一個很重要的功課.
或者我們需要學習的.
就是請你每天都應該為你的丈夫.
或者為你的太太去祈禱.
正如為我們的子女祈禱一樣.
不要只是抱怨他.
你想去改變.
為了改變他.
是上帝的.
所以請你去為他禱告.
為他禱告有什麼指標呢.
或者我們為他禱告做什麼呢.
這些書是有很多的.
我覺得是阿們的.
祈禱唯一的方向.
就是讓他成為合神心意的男人.
或者是合神心意的女人.
這是最終極的目標.
因為我們不知道什麼是最好.
我們每天的興趣都會有改變.
你覺得我老公要變成這樣.
但到時你又有第二個改變.
他又不知道變成怎樣了.
說真的.
大部分結婚的時候你都滿意的.

$^{1121}$不滿意你為什麼要嫁給他.
姐妹們.
但我聽到我們的姐妹們經常都說.
他經常都很目獨.
很悶的.
但你跟他拍拖時他不是這樣的嗎.
那倒是.
那你一直慣著他是這樣的.
那你為什麼會想著他會變呢.
當然神可以令到我們有改變.
但是不是必然的.
因此我們祈禱最好的.
最有智慧的.
總之就是求神讓.
無論我們的子女.
我們的老公或者老婆.
他們能夠成為合神心意.
就一定是無死的了.
大家都是合神心意的男人女人.
一起去生活的時候.
神就活在我們當中.
而最重要的.
在上次我們分享已經講了一個叫成全的愛.
我覺得在我們兩性之間的相處.
都需要有這個胸襟的.
每一個人的創造都是上帝的傑作.
我們幫手是幫什麼.
就是幫他的好能夠真的可以帶出來.
今天是女人.
我用女人的角度去看.
是的我們需要真心的去愛我們的丈夫.
但是愛有些是不真心的.
是因為他對你好所以你愛他.
他對你不好你是不會愛他的.
但是這個真心的愛就好像上帝的愛一樣.
其實你愛他.
就是因為你愛他.
沒有什麼原因.
但是在當中要愛不愛得到.
某種程度來說我們需要花時間去了解他.

$^{1161}$當你不瞭解他的時候就有很多的錯誤.
於是個人的期望的差誤.
就有很多不滿足的出現.
但當你瞭解明白的時候.
你的體諒你的愛.
就容易流通.
實際上是需要照顧的.
因為沒辦法.
我們乖的就怪我們奶奶.
如果你老公被奶奶媽媽照顧.
你怎麼也要接了一半女人的工作.
就做了她媽媽.
照顧她或者是幫助她.
但是這個不是終極的.
當我們成為伴侶的時候.
有時候陪伴和一起同行.
都是能夠發展我們彼此之間.
真正的認識和了解.
愛是要練出來的.
無端端如果你只是愛那件事.
你就要想想家裡多買了多少無謂的東西.
就是因為看的時候是很棒的.
覺得很好用的.
不知如何買回來在家打開或者用一次之後.
所有的缺點你才能看得清楚.
兄弟姐妹人都是.
結婚的時候你可能覺得很棒.
但是婚後一起生活的時候.
一起事奉的時候.
你就發覺原來有很多東西不無奇.
但是聖經說.
愛是永不止息的.
當你在愛的裡面的時候.
慢慢來磨.
慢慢來配合.
你就會能夠享受到.
上帝為我們預備的所謂的婚姻.
甚至在小組裡面也是一樣.
初初都是有認識期很開心的.
接著就來衝突期.

$^{1201}$但衝突期不要放棄.
當你真的讓神的愛貫穿的時候.
一個小組經歷一段衝突的時間.
能夠合一.
那種美好和甜蜜.
是令你真是哈利路亞的.
在這個金句裡面我只是選了這一句.
去提醒所有的姐妹.
為什麼要這麼辛苦.
為什麼要這麼困難呢.
沒有的.
因為哪一人獨居不好.
所以才需要你的出現.
當你承接上帝給你的使命的時候.
我常常說的.
Husband Ministry.
這個可能就真是女人第一個的使命.
當我們回到創造的故事.
有很多的安慰.
因為無論世界的價值觀.
怎麼去定義男生女生.
或者一個夫妻之間的關係.
那個權力.
很多時候是因為賺錢多少而定位.
但這一切都不是上帝的心意.
上帝創造男和女是平等的.
同尊同榮.
大家都是按照上帝的形象所做的.
所以無論你是男生還是女生.
請你以自己的角色.
自己的性別為榮.
喜歡自己.
因為你知道上帝看你是甚好的.
然後你請你學習的.
好好地去愛自己.
和好好地去愛神.
當在親密關係.
好朋友關係.
弟兄姊妹關係.
有磨擦的時候.

$^{1241}$怎麼可以擁抱衝突.
在我這30年的服侍裡面.
我只能這麼說.
你去愛主多一點.
唯有你愛主多一點.
你就能堅持下去.
無論是你的老公.
你的組長.
你的組員.
甚至你的教會.
你所有的制度.
往往沒辦法是不滿足你的.
但是當我們去愛主多一點的時候.
上帝的愛就成為我們的支撐.
我們就能夠可以去活出這一份的愛.
也簡單地說.
我們每一個人都忠於自己在上帝面前的領袖.
我們能夠靠著神的恩典.
我們一定能夠活出自己生命當中的美好.
因為上帝在創世的時候.
已經將我們每一個人的使命.
都能夠命定在裡面.
我們在罪的裡面.
在私慾的裡面.
可能看不到的.
我們需要有另一些的人.
我們的弟兄姊妹.
我們的配偶.
我們的親人.
他們能夠真的有一個承傳的心.
幫助我們.
以至我們的恩賜可以被髮掘.
可以被挑動.
我們裡面一切的好.
能夠拿出來.
所以希望你能夠.
看著你的另一半.
看著你的子女.
甚至看著你的組員.
你都要有個心態.

$^{1281}$你看到他是一個炭.
好像沒用一樣.
但是黑漆漆一樣.
但是原來他裡面可能是翡翠.
問題是我們有沒有能力.
將他裡面的好能夠帶出來.
然後將一切不好的東西.
慢慢地讓它消退.
如果我們能夠活在這樣的愛當中.
活在這樣彼此相愛的環境裡面.
我們相信我們的人生是一定幸福的.
\newpage



\section{歌羅西書 1:15-23}
\label{sec:031772}
\textbf{在基督裡}~~~~俞大鳴牧師
\newline
\newline
連結: \href{http://www.church.com.hk/acms/content.asp?site=cdc\&op=show\&type=product\&code=031772\&layout=sermon}{\texttt{031772}} ~~~~ 語音日期: 2026-03-08
\newline
\newline
\hyperref[sec:031771]{\small{< < < PREV SERMON < < <}}
~
\hyperref[sec:index_chronic]{\small{[返順時目]}}
~
\hyperref[sec:031773]{\small{> > > NEXT SERMON > > >}}
\newline
\newline
歌羅西書 1:15-23
\newline
\begin{longtable}{cl}
\hline
\hline
章節 & 經文 (和合本修訂版)\\
\hline
1:15 & \begin{tabularx}{0.7\textwidth}{X} 愛子是那看不見的神之像，是首生的，在一切被造的以先。 \end{tabularx} \\ \\ \relax
1:16 & \begin{tabularx}{0.7\textwidth}{X} 因為萬有都是在他裡面造的，無論是天上的、地上的，能看見的、不能看見的，或是有權位的、統治的，或是執政的、掌權的，一概都是藉著他為著他造的。 \end{tabularx} \\ \\ \relax
1:17 & \begin{tabularx}{0.7\textwidth}{X} 他在萬有之先；萬有也靠他而存在。 \end{tabularx} \\ \\ \relax
1:18 & \begin{tabularx}{0.7\textwidth}{X} 他是身體（教會）的頭；他是元始，是從死人中復活的首生者，好讓他在萬有中居首位。 \end{tabularx} \\ \\ \relax
1:19 & \begin{tabularx}{0.7\textwidth}{X} 因為神喜歡使一切的豐盛在他裡面居住， \end{tabularx} \\ \\ \relax
1:20 & \begin{tabularx}{0.7\textwidth}{X} 藉著他，神使萬有與自己和好，無論是地上的、天上的，都藉著他在十字架上所流的血促成了和平。 \end{tabularx} \\ \\ \relax
1:21 & \begin{tabularx}{0.7\textwidth}{X} 從前你們與神隔絕，心思上與他為敵，行為邪惡； \end{tabularx} \\ \\ \relax
1:22 & \begin{tabularx}{0.7\textwidth}{X} 但如今，他藉著他兒子肉身的死，已經使你們與他自己和好了，把你們獻在他的面前，成為聖潔，沒有瑕疵，無可指責。 \end{tabularx} \\ \\ \relax
1:23 & \begin{tabularx}{0.7\textwidth}{X} 只要你們持守信仰，根基穩固，堅定不移，不致動搖，離開了你們從前所聽見的福音的盼望；這福音也是傳給天下一切被造之物的，我—保羅作了這福音的僕役。 \end{tabularx} \\ \\
[1ex]
\hline
\hline
\end{longtable}
$^{1}$多謝各位聖徒,我們先一同祈禱.
讓我們能夠來到這個聚集一同見拜你.
求你幫助我們全心全意跟隨你.
全心全意聽你的話語.
認識我主耶穌基督.
和祂一切的榮耀.
跟耶穌基督一同祈禱.
我們這幾個禮拜都是全心全意系列.
我們講道主要的經文是基於二弗所說一章三至十四節.
那裡特別提到信了主耶穌之後的多重恩典.
甚至信主之前都是有恩典的.
我們希望是深入思想上帝給我們這些恩典之後.
也都會激勵我們對上帝是有全心全意的回應.
我們能夠認識主耶穌.
首先是由於上帝揀選.
未信主之前上帝已經揀選了.
我們每一個人都是這樣的.
在創世之前.
就是在耶穌基督揀選了我們.
使我們在祂面前成為聖傑.
沒有瑕疵.
不是真的使我們變成了沒有瑕疵.
只是上帝通過耶穌來看我們.
耶穌是真的沒有瑕疵的.
但是我們跟在耶穌後面.
上帝先看到耶穌.
以致上帝看不到我們的瑕疵.
然後在二弗所說一章五節提到.
預定我們藉著耶穌基督是得兒子的名分.
我們已經說過.
做上帝的後子.
要承受的是永恆的福氣.
信耶穌的人當然不止有永恆的福氣.
還有今生的福氣.
我們的福氣是由今生一直直到永恆.
所以當我們有機會跟別人分享福音.
就不要只說信耶穌是得兒子的名分.
我們一定要說清楚.
如果你不說清楚今生有福氣.
別人就會告訴我們.

$^{41}$但我退休才相信.
臨死前才相信.
反正你也是信耶穌想得永生.
不是的.
所以我們要說清楚.
我們信耶穌是真的會得永生.
但永生是由今世開始的.
由現在開始的.
我們活在世上就已經有福氣了.
一直到永生.
上星期因為突然有病進了醫院.
所以請了主任牧師來港讀.
他講到上帝的救贖.
雖然上帝揀選了我們.
但是上帝仍然要付出代價.
祂要祂的獨生子來到世上.
藉著祂無瑕無痴的血.
為我們的罪.
付出罪的贖價.
使我們願意相信.
跟隨耶穌的人.
來成為上帝的兒女.
能夠被看為沒有瑕疵.
今天我們要思想的.
應該是二哥所說的一章十節.
要照著所安排的.
在時機成熟的時候.
使天上地上一切所有的.
都在基督裡同歸於一.
這首經文最後一句是不容易理解的.
什麼叫同歸於一呢?.
是什麼意思呢?.
難理解的原因.
是因為在原文希臘文裡面.
用了一個比較少用的字.
使翻譯聖經的人.
都很麻煩.
不知道怎麼翻譯好.
比較現代的中文聖經譯館.
我們可以參考一下.

$^{81}$讓我們參考環球聖經協會的翻譯.
這裡這樣說.
等到祂的計劃時機成熟.
就使天上地上的萬有.
都歸於基督.
這是環聖的翻譯.
跳得太快了.
現在我們來看新漢語.
新漢語是一個比較新的譯本.
當所安排的時候滿了.
就使天上的地上的一切所有的.
都同歸於基督.
以祂為首.
今天主要參考了兩個近代中文譯本.
原來同歸於一是歸入耶穌基督裡面.
絕對不要拿道德經說的.
一生二,二生兩,二生四,四生萬.
不是那個一.
不是我們華人那些舊的理解的一.
那個一是說我們歸回宇宙裡面.
我們不是歸入宇宙裡面.
佛教都說歸入宇宙裡面.
與天地合而為一.
我們歸於一是歸入耶穌基督裡面.
我們今天會從另外一段經文.
在《歌羅西書》那裡去思想.
有關基督裡的意思.
在基督裡面.
為何選擇《歌羅西書》呢?.
原因是以弗所書與《歌羅西書》有近親的關係.
以弗所書講到教會是基督的身體.
《歌羅西書》強調基督是教會的頭.
這兩封書信是由同一位弟兄.
叫做推基古送去給教會.
以前送信不是找郵差.
要找專人去送.
古代沒有郵政局.
專人去送.
兩封信由同一位弟兄去送給教會.
結論是甚麼呢?.

$^{121}$這兩封信應該是差不多時間寫的.
並且我們從這兩封信的內容裡面看到.
其實內容都十分類似.
在《歌羅西書》.
短短九十五節的經文.
只是分了四章.
九十五節的經文.
大概有一半的經文都在以弗所書裡面提到.
在保羅的書信中.
有時他講一些東西.
他沒有講到原因.
他不會講原因.
為何現在要講這些東西給大家聽.
他不會講出原因.
但是他也不會是無的放置.
他應該一定是和當代的背景有些關係.
甚至和當代的教會.
那個地方教會有些關係.
他才特別講這些東西.
保羅在第一章似乎很突然地.
高舉耶穌基督的特性.
講到他很厲害,很高.
當然他是很高.
但為何他在其他書信不講.
而在《歌羅西書》講呢?.
其實很可能是和當時一個時代的哲學思想有關係.
那個哲學思想正在影響.
哥羅西地方教會的信徒.
所以他特別要講.
那個思想叫做「諾斯底主義」.
因此我先介紹一下.
有關諾斯底主義的一些信念.
我們了解當時的背景.
然後我們再思想.
保羅為何特別強調耶穌基督.
以及在基督裡面.
諾斯底主義的第一個思想就是.
信仰是一個遙遠,至高.
及人所不認識的獨一個體的神格觀念.
而這個神格有不同的名字.

$^{161}$其實這個思想是否有點像我們的神呢?.
我們在舊約裡面也看到.
我們信的至高上帝.
有時聖經也叫上帝有不同的名字.
有時叫全能的上帝.
有時叫耶和華爾辣.
也有些不同的稱呼.
第二件事就開始有些問題了.
他從獨一個體最高的神.
流出了很多神明.
衍生了很多神明出來.
有點像中國人的玉皇大帝.
下面有很多不同的神.
也有點像基督從聖靈聖父而出.
基督是從最高的神流出來的.
但是他們相信流出的神明.
是會逐漸遠離了最高的神.
慢慢地他們失去了一些神性.
他們的神性質慢慢地變得不太穩定.
有關神和耶穌基督這兩方面.
有些類同.
有關墮落.
諾斯底認為人的墮落的原因.
是因為蘇菲亞從人體裡面流出之後.
是從人體裡面流出而不是反叛上帝.
蘇菲亞是誰呢?.
蘇菲亞其實是智慧的意思.
人慢慢地失去了智慧就墮落.
跟我們的行為不是有直接的關係.
第四方面就是說流出的神.
它慢慢地變質.
墮落了,去到物質世界了.
那怎麼辦呢?.
它寄居在人那裡.
原來流出的神明沒有肉體.
所以它寄居在人的裡面.
它透過救贖.
它如何救贖自己呢?.
透過覺醒.
你失去了蘇菲亞,失去了智慧.

$^{201}$你找蘇菲亞回來就行了.
要覺醒.
要得到救贖.
就要透過救贖者幫助.
你自己搞不定也要找人幫助.
有些像佛教.
知道為什麼要念阿彌陀佛嗎?.
其實就是找阿彌陀佛救你.
你自己搞不定就叫阿彌陀佛救你.
最後一個特點就是.
世界是大有缺憾.
或者是因錯誤而產生的.
我們現在說的是創世這個世界.
但本質上是善良的.
這個世界其實是幻影.
有些像我們現在的電腦裡面叫.
虛擬世界.
Virtual Reality.
不是真的.
Virtual 即是模擬出來的.
這個觀點其實有點像佛教.
一切都是空的.
色即是空.
不是說好色的色.
它所說的萬物都是空的.
你看到的東西是假的.
所以當時教會很可能受了這些影響.
有關上帝觀,基督的觀念.
甚至是救贖論,犯罪.
全部都有影響.
如果這種思想.
學說在哥羅西教會裡面流行.
信徒很容易會偏頗.
原來只靠耶穌才行.
要找Sophia來幫忙.
所以保羅就寫了這封信去幫助他們.
認識耶穌基督是至高的.
歸於耶穌基督裡面就不需要其他東西.
我們現在思想今天的經文.
今天的經文只有九節聖經.

$^{241}$但我們可以分為三段來思想.
每三節為一段.
第一段十五節到十七節.
是講基督與世界的關係.
「愛子時乃看得見上帝之像」.
首先講到基督有完全的神聖.
我們讀《藏世紀》就知道.
亞當和夏娃是怎樣做的.
是上帝按照自己的形象去做.
但基督是由上帝裡面出來的.
希伯來書一章三字講得清楚一點.
「祂是上帝榮耀的光輝,是上帝本體的真像」.
這裡講的榮耀的光輝.
在希臘文的意思是從那裡放射出來.
放射出來不是說月亮反映太陽光.
基督不是反映出上帝的光輝.
是直接從上帝裡面射出來.
好像X光一樣從原子裡面放射出來.
由於基督是從上帝裡面出來.
讓我們了解祂是三位一體.
上帝中的其中一位.
也讓我們可以完全通過耶穌認識上帝.
但上帝在外就有特定的次序.
在舊約時代.
上帝藉著祂的說話和事件.
讓以色列人去了解祂的聖情.
因此在舊約我們看到上帝頒佈律法.
讓以色列人知道上帝是獨一的.
除祂之外就沒有其他上帝.
上帝也不准以色列人扮偶像.
因為偶像很多時候人做偶像真的很喜歡.
模仿動物或者飛鳥.
再加入一些人的想像力.
例如在巴比倫時代.
他們用獅子.
獅子本來沒有翅膀.
但上帝加了翅膀給他們.
這是巴比倫的標記.
在現實世界沒有這種動物.
羅馬人很喜歡用鷹做標記.

$^{281}$用鷹來做他們的標記.
上帝也會藉著一些特別的事件.
讓人去了解祂的聖潔和神聖.
在撒母耳記下.
第六章記載了一件事.
很多基督徒也不太明白.
大衛王想將上帝的藥櫃.
運入耶路撒冷亞.
這是上帝藥櫃的模型.
大衛當時用牛拉車來運藥櫃.
在路途中.
牛失前提.
不平穩.
藥櫃可能跌下車.
有兩個人扶手扶送.
其中一個叫烏薩.
他立即用手扶著藥櫃.
阻止藥櫃跌下來.
但是被上帝擊殺.
這件事很多基督徒也不太明白.
覺得上帝好像不講道理.
烏薩追於好意.
用手扶著藥櫃.
阻止藥櫃跌下來.
為什麼還要殺了他呢?.
其實上帝是要藉這件事.
使以色列人認識上帝的神聖.
其實大衛用牛拉車運藥櫃.
本身是做錯了.
藥櫃本身不應該用車來運載.
但很可能大衛記得.
藥櫃曾經被非利士人擄走.
後來非利士人周身危.
藥櫃在哪個城市呢?.
那裡的人又生痔瘡又生什麼病.
於是他們就找一輛牛拉車.
專門不找人去管理牛.
看看牛會不會將藥櫃拉回以色列人的地方.
果然是牛不用人教.
兩隻都是小牛.

$^{321}$從來沒有拉過車.
去到分岔口的時候.
牠懂得選擇一條路就是回以色列人的地方.
於是非利士人知道.
我們受的災難就是由於上帝整蠱我們.
把藥櫃放在我們身上.
我們就周身危.
他們用牛拉車是因為他們不懂規矩.
以色列人猶太人就應該懂規矩.
不可以用牛拉車.
為什麼?.
你看到藥櫃有兩支棍在旁邊.
要移動藥櫃應該由祭司去抬.
不應該找牛拉車.
所以上帝藉著這件事擊殺污殺.
讓以色列人有所學習.
知道上帝的神聖.
在新的時代.
上主藉著愛子耶穌基督來顯明自己的性情.
普羅繼續提到主耶穌和物質世界的關係.
15節提到是守生的.
是一切被造的已先.
17節又重覆說.
他在萬有之先.
萬有也靠他而存在.
15節講完17節又講.
這是希伯來人慣用的文學手法.
叫做首尾孚應法.
他用這個方法讓你留意中間的節.
一頭一尾包著.
中間的節度是最重要的.
16節是最重要的.
16節強調什麼?.
強調上帝藉著耶穌基督創造整個宇宙.
宇宙不是無端端就在這裡.
而是上帝創造的.
看得到的物質宇宙是上帝創造的.
看不見的呢?.
看不見的可能包括科學家所說的.
暗物質和暗能量.

$^{361}$在普羅時代當然沒有這些概念.
暗物質,暗能量.
他說的看不見的是什麼世界呢?.
是屬靈的宇宙世界.
聖經不是啟示很多有關屬靈世界的事.
但是這裡提到.
靈界裡面有些什麼呢?.
他是有權位的,統治的.
或是執政的,掌權的.
他不是說現在這個世界那些.
總統,總理,國防部長那些有權的人.
他是說在屬靈世界裡面.
原來在屬靈世界裡面都有管理階層的.
有不同的管理階層.
如果我們參考《但義理書》十章十三節.
二十節到二十七節.
我們看到天使之間有爭戰的.
可以有爭戰的.
不過這些.
無論我們看得到也好.
看不到的存在也好.
都是藉著基督而做的.
所以基督才是最權威的.
如果有人問.
神為什麼創造這一切呢?.
一個很簡單的答案就是.
十六節最後那幾個字.
藉著他.
為著他而做.
基督是不是很厲害呢?.
有時候我們認識的耶穌.
好像有點軟弱.
去到拉薩路的墳墓前又會哭.
在馬西馬尼園臨上十字架前又會祈禱.
又求天賦上帝.
可不可以將我的苦杯拿去啊?.
最後還好像看到.
耶穌基督死了在十字架上.
你可能看到一個軟弱的基督.
但實際上的基督不是這樣.

$^{401}$保羅將基督的真實身份告訴我們.
他是一位創造者.
是上帝愛之.
是整個物質宇宙和靈界都是靠他而做的.
不單是靠他而做.
為他而做.
並不是一個錯誤.
上帝不是不小心搞壞了.
造了一個世界.
不是這回事.
不是因為錯誤而創造了這個世界.
這是洛斯底派的人說的.
好了.
十五至十七節是講基督和這個世界的關係.
接著我們看第二段.
十八至二十節是講基督和教會的關係.
什麼是教會呢?.
教會在希臘文的意義.
是指一群被呼召出來的人.
即是你和我.
我們在一起就是教會.
基督成為肉身來到世界上.
就是要完成上帝建立天國的計劃.
教會就是一群被基督呼召出來.
救述出來.
去敬拜上帝的子民.
跟舊約一早有些圖畫一樣.
上帝藉著摩西拯救了一班子民出來.
最重要的事就是要敬拜上帝.
所以他們出了埃及之後.
接著差不多十幾章聖經講他們做什麼.
做回眸.
做回眸的器具.
然後就敬拜.
所以這是最重要的.
成為上帝的子民.
最重要的就是敬拜上帝.
救贖我們出來成為上帝的一家人.
主耶穌曾經對彼得這樣講過.
我要把我的教會建造在這盆石上.

$^{441}$陰間的權柄不能勝過他.
耶穌基督建立教會.
就是要勝過.
教會是勝過所有權柄.
我們不用怕看得見的邪靈世界.
當主耶穌從死裡復活之後.
我們還可以看到.
其實耶穌基督是勝過死亡的權勢.
因此保羅用身體來做比喻.
基督是教會的頭.
而教會就是基督的身體.
此外他也可以說是我們新人類的長子.
因為他是首先第一個出來.
從死裡復活.
所有其他的信徒死了之後.
去了哪裡呢?.
根據路加福音23章43字和肥立比書1章23字說.
是去到樂園與主同在.
今天如果我死了.
我回樂園了.
所以我不怕死.
癌症打我不死.
上帝什麼時候讓我死.
我回樂園而已.
無所謂.
然後呢.
等耶穌基督一起再回來.
聖經預言了耶穌基督會再回來.
但是呢.
事有未到.
到時我們會有一個新的身體.
我們會有新的身體.
在哥林多前的15章23節.
保羅就這樣說.
國人是按照自己的次序復活.
初熟的果子是基督.
因為他首先復活.
初熟的果子.
然後在他來的時候.
是那些屬於基督的.

$^{481}$這個是聖經告訴我們的真理.
19節要解釋一下.
上帝寬起一切的豐盛.
在他裡面居住.
什麼叫一切的豐盛呢?.
其實是指上帝一切豐盛的特質.
有關上帝的特質.
全部都在耶穌基督裡面.
換句話說.
耶穌基督不是次一等的上帝.
他從上帝裡面出來.
有上帝所有豐盛的特質.
他是完完全全和上帝一樣的人.
保羅一再強調這一點.
我相信是為了針對.
洛斯底主義的信仰,思想.
保羅強調宇宙不是邪惡的.
是上帝藉基督而做.
也是為他而做的.
教會也不是邪惡的.
雖然教會的信徒.
我們這些人有些不完美.
但是是基督的身體.
基督是教會的頭.
由於基督的完美.
上帝還使得全宇宙都再與自己和好.
這一切都是靠主耶穌.
在十字架上所流的血而達成的.
這裡給我們看到.
換句話說.
始祖犯罪.
亞當和亞娃犯罪.
帶來一個嚴重的後果.
就是連整個宇宙都被他們拖下水.
還使得人和上帝失去了和好的關係.
所以基督所完成的救贖.
使得全宇宙與上帝和好.
使得我們信祂的人與上帝和好.
成為上帝的兒女.
最後一段.

$^{521}$二十一至二十三節.
再靠近一點.
先講基督和世界的關係.
再講基督和教會的關係.
二十一至二十三節就靠近一點.
講上帝和信徒的關係.
你我都有份.
二十一節是講從前.
我們未信主之前的人的狀態.
二十二節是講現在的情況.
我們信主之後的情況.
二十三節是講我們將來的盼望.
由於對將來的盼望而產生的行為.
曾經有一個比較感人的故事.
講到有一個兒子不長進.
經常出去作惡.
被父親逐出家門.
從此以後不准回家.
但過了一段時間之後.
母親病重,快要死亡.
她未死之前.
有一個心願想見兒子一面.
丈夫無可奈何,沒辦法.
就去找兒子.
回來見母親最後一面.
兒子回家的時候.
母親也快要死亡,奄奄一息.
兒子在床邊大哭.
整段過去真是頑皮.
跪在床邊捉著母親的手.
父親又在另一邊捉著母親的手.
這位母親聽到兒子認罪之後.
轉世之前.
她最後用一些不太有力的力氣.
將兒子的手和父親的手拉在一起.
使父子和好.
耶穌被釘在十字架上.
張開雙臂.
似乎一隻手是拉著天父的手.
另一隻手是拉著我們.

$^{561}$使我們這些原本與神隔絕的罪人.
敵人.
能夠再次恢復和上帝的和睦關係.
很多人以為是原罪使我們與神為敵.
但是二十一節說.
是我們自己的行為和思想遠離上帝.
而成為了祂的仇敵.
但是藉著主耶穌基督作為我們的主.
上帝通過耶穌看我們.
將我們看成是聖見沒有瑕疵.
無可指責.
不是靠我們的行為.
真的能夠做到聖見沒有瑕疵無可指責.
而是靠著祂兒子肉身的死.
使我們來到上帝的面前.
但是我們不能夠以為.
現在我們順手拿了免死金牌.
從此之後可以任意犯罪.
不是的.
如果你真的相信那些人告訴你.
一次得救永遠得救.
得救之後你喜歡做什麼就做什麼.
殺人放火無所謂.
一樣死了之後就入天堂.
對不起聖經沒有這樣說.
二十三節說得很清楚提醒所有信徒說.
只要你們要持守信仰.
根基穩固.
還要堅定不移不自動搖.
不要隨便亂信那些什麼諾斯底主義.
現代的什麼主義都好.
不要隨便亂信.
所以我們不要相信一些人說一次得救就永遠得救.
就可以去為所欲為.
這些是錯誤的說法.
信心和行為可以說是前輩的兩面.
真正的信心是要藉行為證明出來的.
今年我們遲些會看另外一本書.
《雅各書》.
《雅各書》就很強調這一點.

$^{601}$成為他的兒女他的家人.
我們得蒙救贖.
原來我們要歸入基督的身體.
即是普世教會.
不是單單歸入地方教會.
是普世教會.
我們一切都是靠著主耶穌基督在十字架所流的血而促成的.
今天我們會手煮餐.
我們記念主耶穌為我們所做的事.
讓我們會全心全意去愛耶穌.
持守信仰.
堅定不移.
一起祈禱.
我們能夠認識您一切的豐盛.
以及不能夠測度到的厚愛.
求您幫助我們認識您更多.
愛您更多.
在信仰的道路上充滿信心.
堅定不移.
以我們的生命全心全意地事奉您.
奉耶穌基督的名祈禱.
阿們.
\newpage



\section{以弗所書 6:13}
\label{sec:031773}
\textbf{成就了一切,還能站立得住 (華)}~~~~王亞辰牧師
\newline
\newline
連結: \href{http://www.church.com.hk/acms/content.asp?site=cdc\&op=show\&type=product\&code=031773\&layout=sermon}{\texttt{031773}} ~~~~ 語音日期: 2026-03-08
\newline
\newline
\hyperref[sec:031772]{\small{< < < PREV SERMON < < <}}
~
\hyperref[sec:index_chronic]{\small{[返順時目]}}
~
\hyperref[sec:031774]{\small{> > > NEXT SERMON > > >}}
\newline
\newline
以弗所書 6:13
\newline
\begin{longtable}{cl}
\hline
\hline
章節 & 經文 (和合本修訂版)\\
\hline
6:13 & \begin{tabularx}{0.7\textwidth}{X} 所以，要拿起神所賜的全副軍裝，好在邪惡的日子能抵擋仇敵，並且完成了一切後還能站立得住。 \end{tabularx} \\ \\
[1ex]
\hline
\hline
\end{longtable}
$^{1}$我們一起來讀一下今天的主題經文.
是在以弗所書六章十三節.
還有菲利比書三章十三十字節.
我們一起來唸 來.
所以要拿起神所賜的全副軍裝.
好在磨難的日子 抵擋仇敵.
並且成就了一切 還能站立得住.
下一段經文 來.
弟兄們 我不是以為自己已經得著了.
我只有一件事 就是忘記背後努力面前的.
向著標竿直跑.
要得神在基督耶穌裡從上面.
召我來得的獎賞.
從今年年初開始 教會就是蠻有動力.
不論是在各個牧師小組.
或是各樣的士工當中.
真看到有活力.
又有這個聖靈的火在我們當中.
其中最讓我們感動的.
就是全教會的背誦聖經行動.
因為非常多的弟兄姊妹踴躍的參與.
那我們看到能夠一起來背誦聖經.
真的叫做患難中見真情.
因為背聖經實在是挺挑戰的.
特別像我這種記性本身就不是很好的人.
我就選擇了.
閉著眼睛背起來.
錄下來 記到APP上面.
因為我就發現如果在家背好了.
然後開車到教會.
到停車場.
下來遇到幾位弟兄姊妹聊天.
再到這個攤位.
根本就已經忘光了.
所以我就好 我就用這個閉眼睛.
哇 我跟你講 閉眼睛背.
我都背了三十幾遍才通過.
但是你知道嗎.
我不知道有沒有因為我是牧師.
所以放水 希望沒有.

$^{41}$我真的是一字不漏.
可是我就發現每一次這樣子在背聖經的過程.
神的話語越咀嚼.
就越有味道.
越來就越有感動.
而神就透過這樣一段經文.
對我有深度的一個感觸.
所以今天我們的信息其實.
如果你沒發現的話.
就是一月和二月的經文裡面的內容.
以弗所述六章十三節.
真的是我從小就在念.
可是這次我特別有感觸.
這邊說要拿起神所賜的全副軍裝.
不但是要拿起軍裝是神所賜的軍裝.
不但是神所賜的軍裝.
是要全副的軍裝.
還不能只是半套而已.
為什麼要穿戴全副軍裝呢.
因為會有磨難的日子.
好在磨難的日子我們可以度過.
為什麼要穿全副軍裝.
因為有屬靈征戰才能夠抵擋仇敵.
為什麼穿上全副軍裝.
因為這邊說要成就一切.
弟兄姊妹你知道嗎.
神在我們生命當中有極豐富的計畫.
極榮美的計畫.
祂期待我們能夠去成就一切.
你這一生不是注定一事無成.
你這一生上帝要你成就一切.
因此要穿戴全副的軍裝.
可是這一次當我在一直背聖經的時候呢.
最觸動我的卻是最後一句話.
成就了一切還能夠站立的住.
這段時間我們聽見有一些神所.
重用的僕人跌倒了.
一些很有影響力的機構.
也出了問題.
成就了一切卻是沒有辦法站立的住.

$^{81}$當我在讀這段經文的時候.
我向神深深的呼求.
吶喊 抓.
但願我不但能夠因著你成就一切.
更能夠靠著你還能夠站立的住.
Amen.
事實上神國的道路.
在各種不同的季節.
都有不同的挑戰和屬靈征戰.
而屬靈的詩題常常是在意想不到的時候.
我們來看一下這個圖.
就是一個神國度路不同季節的一些挑戰.
這是我和我的好朋友AI一起畫出來的圖.
我們就看在一開始神呼召你的.
你回應神呼召之前呢.
恐懼與自卑會導致你遲遲不前.
為什麼 我不行 我不知道未來會如何.
別人比較厲害.
而當我們進入了磨難的日子.
有時候血氣與短視.
導致我們偏行幾路.
磨難的日子有時候會讓我們偏行幾路.
磨難就是有磨也有難.
磨就是要熬 難就是要挨.
但是因為血氣的關係.
我不要再忍耐了 我受不了了.
就磨不下去.
因為短視讓我們想要快速的成功.
地上的成就.
我不要那個人瞧不起我.
讓你就挨不住了.
求神賜給我們專注力和耐心.
能夠定心永恆的呼召.
而當我們看到遇到仇敵的攻擊的時候.
最怕的就是生命中有破口.
而苦讀和貪婪.
會成為仇敵能夠趁虛而入.
在思想中把我們半跌得很慘的破口.
特別苦讀與貪婪.
所以因此我們要不斷的來向神.

$^{121}$和在神面前說我們要保守我們的心.
勝過保守一切.
而在成就之後呢.
是什麼樣的挑戰.
就是自負與偏離.
遠離了真神.
我們可以把這個圖先收起來了.
當我們成就一切之後.
你成為了一神之下萬人之上的大英雄.
你成為了所有人歌功頌德的對象.
你成為了一個絕對真理一般的一個存在.
就算你唱個走音別人都說.
哇你好像在唱藍調或是R and B一樣.
這個時候你會變得驕傲自負.
而這個時候問題就出現在你的眼目.
開始只定睛於已經得著的.
卻不再定睛於主耶穌.
自己從上面要你來得著的呼召.
其實說實在話也不是真的成就了一切.
才會有這種狀態.
多少時候當我們人有一點點成就.
我們就容易進入這樣的一個態度.
那麼解藥呢就是二月份.
腓立比書三章十三十字節這邊說的.
我不是以為自己已經得著了.
我不是以為我已經得著了.
我只有一件事.
就是忘記背後.
忘記背後一切的榮耀.
忘記背後一切的羞辱.
忘記一切的好與歹.
我要努力面前的向著標竿直跑.
要得神在基督耶穌裡.
從上面召我來得的獎賞.
不是我現在已經得到了獎賞.
不是我現在開始進入.
我自我的虛榮吹捧的獎賞.
永遠只有一件事.
要得基督耶穌裡.
祂施加的大能和救恩.

$^{161}$要釋放出來的極榮耀永恆的獎賞.
因此我在讀完這個經文的時候.
我真的向主說主啊.
幫助我能夠度過磨爛的日子.
主要幫助我能夠勝過仇敵的攻擊.
主要幫助我能夠.
經歷得著一切你召我來得的獎賞.
而在這一切當中我還能夠站立的住.
事實上聖經當中也記載了.
許多成就一切.
站立得住.
或是成就一切後來站立不住的一些故事.
讓我們今天也能夠用這些故事來成為借鏡與學習.
我們先來看一下聖經當中有哪些站立不住的例子.
首先就是大胃王.
哇 大胃王是我們多麼喜歡的一個人物啊.
他從小就滿有聖靈的同在.
他勇敢戰勝哥利亞.
他專注於神的呼召.
他還能夠饒恕 掃羅.
他又竊慕神 歡迎 約會 設立祭祀之分.
尤其是最寶貴的一句經文.
是他對上帝的一個愛慕的情書.
我只有一件事.
就是一生一世住在主的殿中.
瞻仰祂的榮美.
哇 這樣的大胃.
竟然在他好像成就了一切的時候.
卻跌倒了.
他與巴士巴犯了姦淫罪.
他借刀殺人 殺了伍裡亞.
而且不但如此 他不是只是跌倒一次啊.
他第二次跌倒.
他在詩篇51篇的時候 他悔改.
說啊 請你不要掩面不見我的面.
可是在他好像真正的成就了一切.
在耶路撒冷登基之後.
他又站立不住了 弟兄姊妹.
他發生了什麼事情 你還記得嗎.
他吩咐約鴨 要去數典人數.

$^{201}$這個時候大胃王已經變得驕傲自負.
你可能說為什麼數典人數有問題.
上帝不是也要叫摩西數典人數嗎.
但是這個時候的數典人數.
大胃王他是在使用國家的資源.
並不是在做神要他做的事情.
並不是用在治國與愛神的上面.
卻用在炫耀國力自我欣賞的虛榮之中.
以至於以色列因著他的跌倒.
遭到了瘟疫的咒詛.
你不覺得很諷刺嗎.
大胃王在巨人面前站住了.
在掃羅需要饒恕的這樣的一個關卡中.
他站住了.
可是在他順境的時候.
他卻脫去了鼠鱗的軍裝.
因此跌倒了.
另外一個人物是所羅門王.
所羅門王是歷史上最有智慧的人.
我不知道你是不是也跟我一樣.
從小在教會聽到所羅門王的故事.
都是好事 沒有人說他不好.
為什麼 因為他很有智慧.
所以我們可能一想所羅門王.
他有智慧 所以他做的一定是對的.
我沒有他聰明 所以我就乖乖聽著就事了.
哇 結果可是我發現.
所羅門王他真的就是完美到底嗎.
哦 很有智慧的人.
真有智慧.
每次想到他娶了一千個老婆.
我說哇 這個不可思議.
我一個老婆都搞不定.
一千個老婆 這是何等的智慧.
他選擇政治聯姻啊.
我在想這個太brilliant了.
政治聯姻就世界太平.
每個都生個小孩.
哇 那就搞定了 沒問題了.
brilliant.

$^{241}$然後他又與各國貿易買賣.
包括埃及的交易.
哦 這個太棒了.
把世界各地的資源都引到了神的國度.
好得不得了.
他後來又把各國賓妃的偶像假神.
都帶到以色列來.
還設立秋壇讓他們去拜假神.
哇 這個太有智慧了.
這有智慧 這是對嗎.
你說對啊.
不對啊 怎麼可以呢.
把假神帶到以色列.
還幫他設立秋壇.
這是一個什麼情況.
弟兄姊妹 你知道嗎.
其實索羅王早就跌倒了.
你知道嗎 在申命記十七章這邊說到.
是摩西的這個律法.
不可回埃及買馬匹.
因為不能再走回頭路了.
可是他和埃及做貿易.
不可為自己多立賓妃.
恐怕他的心偏邪.
哇 他不只是多立賓妃啊.
立了一千個賓妃.
並且這邊說不可為自己多積金銀.
他的王宮 自己住的地方.
比聖殿還要奢華無數倍.
為什麼一個充滿智慧的王.
最後會跌成這個樣子呢.
也許他心裡想 這些律法很好.
但是摩西是寫給那些沒有智慧的王所寫的.
我可是生活智慧王 我沒問題.
這些律法是寫給那些邪惡的.
不愛神的王設立的.
我可是從小就是神眼中的乖乖牌模範生啊.
我還這麼謙卑的說 上帝啊 我需要你的智慧.
所以他可能因為這樣的關係啊.
反而他開始慢慢一步一步的跌倒了.

$^{281}$弟兄姊妹 敬畏耶和華是智慧的開端.
但是當索門王不再敬畏神的時候.
他的智慧雖然仍然能夠讓他.
富足 甚至成功.
但是卻得不著富足 又不加上憂慮的祝福.
弟兄姊妹 索門王犯的錯不是普通的大的錯啊.
你知道他設立的這個秋壇.
最後整本舊約啊 都在處理這個問題.
所以呢 當他成就了一切之後.
權力與榮耀 是他親看了樹林的軍裝.
另外一個王是西西伽王啊.
西西伽王也是一個不得了的好王耶.
列王記下十八章五節這樣描述的.
西西伽王倚靠耶和華以色列的神.
在他前後的猶大獵王中 沒有一個及他的.
哇 是非常棒的王 他除掉秋壇 打碎偶像.
恢復聖殿 重新設立玉月節.
他又經歷了戰勝亞述王的神蹟.
並且我們都知道他經歷了上帝給他多十五年的壽命.
但是在他晚年的時候啊.
哇 你看這麼好的一個王.
當他成就了一切 他卻為了自我的虛榮心.
使巴比倫的使者在這個拜訪的時候.
他把國庫 武器庫 所有財寶都給他們看.
這就成為了後來巴比倫攻打 俘虜猶大國的伏筆.
這問題出在哪裡啊.
如果你有好朋友 他來你家做客.
頂多待在客廳吧.
真正的好朋友 再好的.
來房間參觀一下remodel廁所 這很好了.
肯定不會把你銀行的存摺來給他看吧.
希西加王在晚年的時候成就了一切 卻站立不住.
這是一些寶貴的例子.
首先就是迦勒 85歲依然求戰.
他永不忘記神的呼召.
當時其實約書亞帶領的以色列已經打遍了大部分的江山了.
身為這個第一戰將 他絕對可以優先來選擇一塊肥沃又安全的土地.
結果他真的第一個出來選了.
可是你記得他選什麼樣的地嗎.
他選那塊還沒有打下來的西伯倫啊.

$^{321}$那是一個充滿了亞納族的巨人.
而且是當年12個探子所描述的.
對方是巨人 我們是詐謀的那塊地.
沒有想到迦勒竟然說求你將耶和華那日應許我的這山地給我.
為什麼呢.
因為他到最後仍然不忘記起初的呼召.
約瑟是另外一個人.
他在權力巔峰的時候 他仍然穿戴神的品格.
哇 身為埃及的宰相 他真的是一根指頭可以調動全世界啊.
但是在這樣的一個時刻.
他出賣他的十個兄弟出現的時候.
他為何可以選擇不報仇 反而饒恕呢.
我相信因為他已經選擇要保守他的心 勝過保守一切.
他選擇定義不要苦讀怨恨 成為生命中的破口.
這何等難得的事情.
因為他要報仇太容易了 太容易了.
在權力巔峰的時刻 穿戴神的品格真是不容易.
但是也是因此 他能夠真的成就了神在他生命中最終極的命定.
成為埃及的宰相不是一切 能夠與家人和好.
並且讓以色列能夠成為一個合一的族.
之後成為一個神所使用的國度.
約瑟是一個極關鍵的人物.
弟兄姊妹 因著他持守 他經歷了.
他不但站立得住 他經歷了神在他生命中最榮耀的呼召.
撒母耳是另外一個例子.
他是誠信 清廉 至忠 活出神的聖傑.
他是一個非常特別的角色 不但從小就開始服侍神.
有人說他是聖經當中唯一的一位同時兼備先知.
祭司君王的角色 並且他是唯一一個掃羅王和大衛王都要乖乖聽話的人.
這樣的人不可思議 多麼有影響力.
可是在他人生的尾端 他的一句話是震撼了人心.
他把所有人召集到廣場 就在公開的問說.
我奪過誰的牛 我搶過誰的驢.
我欺負過誰 虐待過誰.
從誰手裡接受過賄賂 因而眼瞎呢.
若有 我必償還.
而在眾人面前 他可以如此的坦然.
百姓也齊聲說 你未曾欺負我們.
未曾虐待我們 也未曾從誰手裡授過什麼.
哇 這個是很寶貴的一個榜樣.

$^{361}$他的誠信與聖潔.
使他成就了一切 還能夠清廉.
另外一位成就一切 還站立得住的就是保羅.
新約的保羅 他真的是殉道前.
一直都守住真理 直到得著獎賞.
但今天我們在這裡 不是來品頭論足.
說誰比較屬靈 更不是批評論斷.
決定誰成功 誰失敗.
每個聖經的故事 也包括每一位.
現今的屬靈的巨人神的器皿.
生命中的成敗或聖衰啊.
其實都是我們值得學習與借鏡的教材.
但是我們千萬不要開始論斷 定罪.
或是妄下結論.
這些站立不住的人 神會親自來管教.
神會親自來管理.
那是神與他們之間的事情.
反而我們該做的呢.
應該是為這些人來饒恕 來帶導.
也是透過這些機會來自省.
並且重新定睛仰望.
畢竟今天的站立 並不代表明天不會跌倒.
而昨天的站立不住.
並不代表你今天不能再次起來為主而活.
Amen.
聖經當中也充滿了跌倒之後 又站立起來的例子.
我們就來看到.
大胃王就是極寶貴的例子.
我們來看一下 我和我朋友做的另外一張圖.
我們看到大胃王在事奉的巔峰的時候.
哇 多麼屬靈 迎接約櫃 設立帳目與祭司.
來 好像就是建立大胃的帳目啊.
可是 他後來首次的跌倒.
八十八事件 借刀殺人.
謙卑悔改 寫下了十篇五十一篇.
重新讓信仰根基回歸神自己.
可是很不幸的 他二次的跌倒.
他擅自數點人數 強調軍力 強調數字與成就導致.
瘟疫讓以色列非常的受苦受難.
難道這樣就是宣告他再也不能被神使用了嗎.

$^{401}$難道這樣就是宣告他是一個失敗的人嗎.
但是我們看見 當他再一次.
你說厚臉皮也好 死皮賴打也好.
他再一次悔改 謙卑回來 竊慕神的國.
這一次才是他開始預備聖殿.
開始大規模的編排立位人與祭司班次的時刻.
所以弟兄姊妹 上帝是能夠再一次給我們機會.
給我們恩典的 不論你的昨天如何.
你昨天再成功 不代表你今天不會跌倒.
但是你昨天 就算你跌得體無完膚.
羞愧 羞於見人.
並不代表神不能繼續來使用你.
只要我們願意迴轉來到神面前.
殊不知 這是你離神在你生命當中.
從基督耶穌裡 從上面招你來的獎賞和護照.
最接近的時刻 Amen.
那麼我們要如何.
所以我們要如何能夠站立得住.
讓我們今天彼此來提醒來勉勵.
首先 我們要有問責的環境.
好不好意思說 問責的環境.
英文就是accountability.
以弗所述六章十八十九這邊說.
保羅說 並要在此警醒不倦.
你是需要警醒的.
然後呢 就是一個問責的一個圖畫了.
為眾聖徒祈求.
你對弟兄姊妹 你有問責 有責任.
並且也為我祈求.
弟兄姊妹 你對我能夠站立得住 也有責任.
所以這邊講得很清楚.
我們需要彼此守望.
需要彼此來問責.
因為我們知道 各樣的試探太多了.
不論是財務上面.
不論是在情慾上面 在權利上面.
我們不但需要有清楚的界線.
和保護的原則.
我們更需要有人這樣子的.
彼此的問責 才能夠成就一切啊.

$^{441}$才站立得住.
所羅門的問題就是.
He's above the law.
這是極大的問題.
那我們怎麼樣能夠知道.
我們是有一個問責的環境.
有一個謙卑的聆聽的耳朵.
讓別人來問責你啊.
請問最近一次有人.
指出你的錯誤是什麼時候啊.
你可能說 就是今天早上啊.
就我的旁邊那位.
不是問責 只是責而已.
其實我告訴你 這是神很大的祝福啊.
有人在你旁邊 耳提面命啊.
還是可以跟你嘮嘮叨叨一下.
這其實真的是一個幸福.
如果你已經很久沒有人.
好像來指責你.
你好像大家都覺得我都不錯.
所以沒什麼人指責我.
但是聖經講得很清楚.
當人都說你好的時候 怎麼樣.
你就有禍了.
所以弟兄姊妹.
如果有人指責你.
還為那個指責你的人氣啞啞的.
你今天應該轉換個心情.
不然我就有禍了.
所以首先我們需要有問責的環境.
我們需要有禱告的夥伴.
第二呢 榮耀永遠是屬於神的.
我們一起來說 榮耀永遠屬於神.
請問你是否有將榮耀.
一次又一次的歸給神呢.
這個事情我們可以在.
跟啟示錄的二十四位長老來學習.
這二十四位長老是不得了的.
你看這幾天年來.
跟隨神的基督徒.

$^{481}$跟隨多少億人.
可是有這二十四位長老.
在啟示錄最後的最大的.
高揚的敬拜的時候.
只有這二十四個人.
他們應該是最頂尖的人物.
可是看到聖經上的描述.
他們每一次 每一天來到寶座前.
他們就將他們頭上的冠冕怎麼樣.
放在高揚的寶座前.
Casting the crown.
將這個冠冕.
榮耀的冠冕放到祖人面前.
為什麼呢.
原來我們發現.
這些二十四位長老.
他們真正所愛慕的不是冠冕.
他們所愛慕的是.
主基督高揚他自己.
那我們怎麼樣可以知道.
我就發現其實很特別.
因為他們每天.
每天來到主面前.
都把crown放到主面前.
你有沒有發現.
如果他把crown放到主面前.
隔天不是就沒有了嗎.
怎麼還有crown.
原來他放在主面前.
我猜啦.
耶穌最後還還給他們.
還帶回去.
隔天再來.
哇.
所以弟兄姊妹.
如果我們將榮耀歸給神.
神會將他榮耀跟我們分享的.
但是我們如果抓著不用的話.
抓著不願意將榮耀歸給神.
那問題就開始了.

$^{521}$那我們怎麼樣知道.
我們人能夠將榮耀歸給神呢.
我在自己思想的時候.
我就想到說.
問大家一個問題.
你是否有注意到.
你身邊默默服侍的人.
你是否能夠將功勞.
歸給這些默默服侍的人呢.
還是你一天到晚就開始想.
啊 我很重要.
我很了不起.
然後這些旁邊人你還嫌他.
嫌東嫌西的.
還覺得他不夠聰明.
還成為你的麻煩.
求主幫助我們.
能夠成為一個.
將榮耀歸給神的人.
不要自己一直來竊取搶奪的人.
接下來就是要來穿戴軍裝.
我們今天.
應該大家都很清楚很熟悉了.
要穿戴軍裝就是神的話語.
真理公義福音信德救恩.
神的道.
要為眾聖徒祈求.
我們都知道.
但是我們真的有在穿戴嗎.
還是我們一月份.
把聖經經文背完了成就了一切.
然後就把它忘記了.
什麼樣能夠讓我們真的是穿戴軍裝.
請問你是否每天都期待神對你說話呢.
或者你是否成為能夠.
領受神的話語.
然後去分享的人呢.
我就.
發現一個渴望神對你說話的人.
其實隨時隨地.

$^{561}$上帝都樂意對你說話.
因為飢渴慕矣人有福了.
因為他們必得飽足.
所以如果你真的竊慕.
不論任何人講到.
你都會有收穫.
不會說你今天一進來.
蛤今天不是留牧師.
光頭大鬍子.
好吧來補眠吧.
今天如果你沒有覺得神對你說話.
不是我的問題.
如果你真的有一個渴慕神的一個心.
任何的場合的敬拜讚美.
你都能夠被神大大的觸摸.
甚至是清唱的時候.
你都可以被神大大的感動.
所以我們要穿戴軍裝的關鍵是.
要渴慕神對你說話.
最後就是要保守我們的心.
我們來說保守你的心.
真言四章二十三節這邊說.
你要保守你心勝過保守一切.
因為一生的果效是由心發出的.
弟兄姊妹你可以想像嗎.
如果約瑟他那時候選擇蒙蔽他的心.
哇這個整個聖經舊約的歷史就寫不下去了.
但是因著他保守他的心.
神的旨意就成就了.
弟兄姊妹.
神在你生命中有美好的旨意.
但是今天很重要的.
你要來選擇保守你的心.
這不是為了別的.
因為一生的果效由心發出.
主就幫助我們.
讓我們不是以為自己已經得著了.
讓我們保守我們的心.
永遠定精在基督耶穌裡.
為我們得著的那個護照的獎賞.

$^{601}$今天不論我們親愛的弟兄姊妹.
你是在磨難當中呢.
還是你在一個征戰之中.
或者你現在正在成就還不錯的時刻.
求主讓我們都是選擇要來穿戴軍裝.
讓我們能夠站立得住.
並且能夠看到神的旨意在我們當中.
至終可以成就.
讓我們一起禱告.
在我們生命中有美好的計畫的神.
今天幫助我們.
讓我們就來到恩典的座前.
能夠成為在磨難的日子.
有些弟兄姊妹正在磨難的日子.
我們可以有勇氣.
我們可以有耐心.
我們可以定精於你.
我們當中也許有在極大的征戰當中.
幫助我們不要有血氣.
保守我們的心.
或是我們當中有些是蒙神大的祝福.
做了在一些影響的位置上面的弟兄姊妹.
主要幫助我們能夠永遠謙卑.
不要自負.
主要定精在天上的獎賞.
讓我們成為一個柔軟的人.
Amen.
\newpage



\section{提摩太後書 2:3-4}
\label{sec:031774}
\textbf{突破知易難行的瓶頸}~~~~招世超牧師
\newline
\newline
連結: \href{http://www.church.com.hk/acms/content.asp?site=cdc\&op=show\&type=product\&code=031774\&layout=sermon}{\texttt{031774}} ~~~~ 語音日期: 2026-03-08
\newline
\newline
\hyperref[sec:031773]{\small{< < < PREV SERMON < < <}}
~
\hyperref[sec:index_chronic]{\small{[返順時目]}}
~
\hyperref[sec:031775]{\small{> > > NEXT SERMON > > >}}
\newline
\newline
提摩太後書 2:3-4
\newline
\begin{longtable}{cl}
\hline
\hline
章節 & 經文 (和合本修訂版)\\
\hline
2:3 & \begin{tabularx}{0.7\textwidth}{X} 你要和我同受苦難，作基督耶穌的精兵。 \end{tabularx} \\ \\ \relax
2:4 & \begin{tabularx}{0.7\textwidth}{X} 凡當兵的，不讓世務纏身，好使那招他當兵的人喜悅。 \end{tabularx} \\ \\
[1ex]
\hline
\hline
\end{longtable}
$^{1}$今個星期不知大家有何看法和感受呢?.
雖然我是在香港.
但大家也知道這個星期.
以色列和美國的軍隊閃電襲擊伊朗.
將伊朗的獨裁者和大部份的領袖消滅.
當我聽到這個消息時,我又想起.
其實是今年二零二六年一月份的時候.
美國的特種部隊.
大家記得,深入威尼瑞拉的首都.
進入總統府.
將威尼瑞拉的總統和總統夫人.
劫持回美國受審.
我在此不去批評這個行動是否合法.
是否正確.
不過讓我想起以色列的軍隊和美國的軍隊.
他們在戰術和實力上.
是令到全世界嘩然的.
很難想像竟然能夠做到這樣的情況.
現在打仗和以前很不一樣.
但由於我們見到以色列軍隊和美國軍隊.
這些表現真是令我們嘩然震驚.
我就想起美國的第一任總統.
華盛頓,喬治·華盛頓.
他曾經說過一句話.
"Discipline is the soul of an army".
我將它翻譯成中文.
紀律是軍隊的靈魂.
可否跳一跳,我有沒有,應該有.
紀律是軍隊的靈魂.
能夠讓少數人變得強大.
讓弱者獲得成功.
並且贏得所有人的尊敬.
我就想起這句說話.
與此同時這句說話.
也勾起了我對基督徒其中一個想法.
其實這不是我對基督徒的想法.
不知道大家有沒有忘記.
我們在基督耶穌裡面.
其中一個身份是什麼.
除了神的子女.

$^{41}$耶穌基督的門徒.
我們在基督裡面其中一個身份.
是耶穌基督的精兵.
大家記不記得.
大家有沒有聽過.
很多時候我們不記得.
原來我們正在參與一場屬靈的戰爭.
保羅怎麼說呢.
保羅在提摩太后書第二章第三第四節說.
你們要和我同壽苦難.
好像基督耶穌的精兵.
凡在軍中當兵的.
就不會讓誓無前身.
好腳勒招他當兵的人.
喜悅.
成為基督的精兵.
對很多基督徒來說.
包括你和我.
很多時候只是一個觀念.
你聽過的.
你知道的.
但對我們來說.
似乎是一個觀念.
知道.
但怎麼做出來呢.
知易行難.
所以保羅說的我們都知道.
做耶穌基督的精兵.
是要吃苦的.
你受訓練.
去操練的時候.
不容易的.
要能吃苦受磨練.
那你才能夠完成任務.
除此之外.
要做耶穌基督的精兵.
我們必須要專心.
要專注.
要守紀律.
守規矩.

$^{81}$否則的話.
我們就很難成為一個.
強大的軍隊.
可以打勝仗.
聖經說叫他招我們當兵的人.
是誰呢.
就是我們的教主耶穌基督.
開心快樂.
他也是我們的統帥.
問題是什麼呢.
今天我們發覺有很多教徒.
不單是忘記了.
我們是基督精兵的身份.
忘記了.
於是我們不知不覺地成為了逃兵.
我們在屬靈的戰場上慢慢消失.
忘記了要打仗.
忘記了我們在屬靈的戰爭當中.
所以很多教徒我們發覺.
被卡住在知而不行.
你知道的.
不過沒有走出來.
知道神的話語教導是什麼.
但是做不到.
我們卻卡在屬靈的瓶頸裡面.
我們的屬靈生命於是變成了.
呆滯不前.
不進不退.
卡在那裡.
什麼都不能動.
失去了一個屬靈的動力.
甚至我們連在世上生活.
都失去了方向.
我不知道這是不是大家的情況.
我嘗試去了解.
嘗試去體會.
今天我們基督徒所面對的挑戰.
如果是這樣的時候.
很自然的是什麼呢.
我們就隨從自己的心意.

$^{121}$我們有什麼需要.
我們有什麼慾望.
就被這些東西帶著我們去走.
每天的生活.
譬如很簡單的.
大家知道早睡早起的.
有沒有聽過.
我們小時候家裡都說的.
老師也說的.
早睡早起身體好.
我們知道要有自律的生活.
維持一個健康的人生.
聽過的知道的.
有多少人做到.
早睡早起經常做運動.
做運動好不好.
好.
有多少人經常做運動呢.
但是.
很多時候我們就缺乏了那種毅力.
和那個自律.
我們說到今天這個世界很亂.
大家現在經常聽.
不用牧師講你都知道.
你感覺到的.
你感受到的.
這個世界很亂.
為什麼這麼亂.
太多謠言.
太多假新聞.
太多那些不盡不實的報道.
言論在我們當中滿天飛.
我們知道這些是有這樣的情況出現的.
但是我們還聽到很多人大聲疾呼.
講什麼.
要誠實.
我做了一個有誠信的人.
特別美國的社會是講trust(信任)的.
講credit(信用)的.
你有沒有多少誠信.

$^{161}$你的信用卡有多少錢.
看看你過去的表現是怎麼樣的.
是不是.
很多時候我們聽到很多這樣的呼籲呼聲.
大聲疾呼在那裡喊叫.
但是實際上.
有多少人.
真的可以言行一致.
這個我覺得是很值得我們今天基督徒去思考的一個問題.
因為我發覺就算在教會裡面.
有很多基督徒甚至教會的領袖.
說謊當吃生菜.
說一套做一套.
我不希望在我們教會裡面見到這樣的情況.
但是我知道.
有很多的教會.
有這樣的情況出現.
時不時大家聽到.
教會裡面有些醜聞發生.
這個反映出教會屬靈生命的那種呆滯不前.
教會屬靈生命的那種軟弱.
我們現在教會.
以前教會在社會當中影響力很大的.
牧師是很有權威很有尊嚴的.
現在挑戰教會挑戰牧師的權威.
因為很多傳道人失敗.
因為很多教會裡面太多偽善的人.
所以為什麼年輕人不來教會.
這個是第一點.
這個是不是值得我們去思想的事情.
當我去思想這個題目的時候.
我發覺今天會不會在教會裡面你寬恕牧師這樣說.
很多只是星期一來聽道的人.
而不是來學行道的人.
於是很多教徒.
就被困在知而不行.
或者是知而難行.
這個屬靈的瓶頸裡面.
於是我們卡住在那裡.
不上不下.

$^{201}$那你說什麼實踐大使命.
你說什麼做基督徵兵.
你說什麼去打屬靈的仗.
沒有辦法.
所以這幾個星期跟大家分享.
如何突破屬靈的瓶頸.
今天希望能夠幫助大家.
突破一個知而行難.
或者知而不行.
這個屬靈的瓶頸.
如果我們想突破這個屬靈的瓶頸.
我覺得第一件事是什麼呢.
我們要培養一個習慣.
什麼習慣.
尋求明白真理的習慣.
聖經裡面怎麼說呢.
法國書第一章二十二字.
指示你們要行道.
不要單單聽道.
自己騙自己.
因為聽道而不行道的人.
就像人對著鏡子.
看自己本來的面目.
看到之後隨即忘了他的相貌如何.
唯有詳細察看那傳遍使人自由之律法的.
並且事常如此.
這個人既不是聽了就忘乃是實行出來.
就在他所行的事情必然得福.
我希望我們在座的弟兄姊妹.
都有做靈修.
有沒有做靈修.
我希望你們已經有做靈修的習慣.
但我想說的是.
就算你有做靈修的習慣.
天天看聖經祈禱.
不等於你明白真理.
你同意我這麼說嗎.
更加不等於我們有尋求明白真理的習慣.
首先.
如果我們會做靈修.

$^{241}$怎麼做靈修呢.
讀經,祈禱.
我希望你們每天都有做.
每天讀聖經,每天祈禱.
你會發覺很容易忘記.
讀完就不記得了.
是不是這樣.
我相信如果我們有這樣的情況.
一個很基本的原因是什麼呢.
可能我們誤解了什麼叫做做靈修.
做靈修其實不是一個祈禱,讀經的宗教活動.
很多時候我們只會當做靈修是讀經,祈禱.
於是你讀完聖經,祈禱.
你覺得是了這件事了.
做完了.
那當然很容易忘記了.
其實做靈修的意思是.
我們與上帝建立一個親密的關係的時間.
把這段時間我們與神在一起.
讀經,祈禱.
是一個宗教活動.
但是不止於,不限制於只是一個宗教活動.
靈修其實是要我們進一步與神建立一個更加親密的關係.
讓我們認識神.
讓我們更加能夠體驗到神對我們的愛和恩典.
以及學習我們如何去愛神的一個操練.
所以靈修對一個基督徒來說.
就不應該是一個很死板的,死固固的.
每天就讀經,祈禱.
完了一件事,就合完聖經就走人了.
做靈修其實是我們與神在一起.
見到上帝在這裡.
然後看看上帝有什麼要跟我們說.
我們有什麼問題就向上帝問.
大家有一個很活潑的互動和對話.
因為只有這樣你才能建立關係.
但很多時候我們做靈修的時間.
起床,趕著上班,趕著帶小孩上學.
於是讀經,祈禱,搞定了一件事.
了了一個做基督徒的責任就完了.

$^{281}$第一件事我想告訴大家.
如果你想要明白真理的話.
你必須首先知道靈修是與神見面的時間.
那你就不要趕快去讀聖經.
趕快去祈禱.
在你做靈修的時候.
首先第一件事是什麼?.
要靜下來.
靜下來,為什麼要靜下來?.
讓我們意識到神的同在,神在這裡.
因為你和神見面.
你想和神溝通.
你想有屬靈的操練.
直至到你能夠靜下來覺得.
神在這裡,與你一起.
神在這裡,與我們一起.
然後你才開始.
那你有什麼要跟我們說,上帝?.
今天你想教我們什麼?.
提倡了什麼?.
然後我們有什麼事情聽不明白.
我們不知道,想不通.
為什麼會這樣?.
有沒有去問神?.
你做靈修的時候.
有沒有神討論問題?.
有沒有神去辯論.
你一些不同意聖經的說法?.
很多時候我們沒有,是嗎?.
有沒有?.
我希望你有.
但是我知道很多基督徒是沒有的.
讀完聖經,完成了.
啊,這個聖經安慰我.
這個聖經教導我什麼.
那就放進去.
放了多久不知道.
可能蓋上之後就馬上忘記了.
所以聽師傅去做靈修.
你要明白那個意思.

$^{321}$做靈修其實我稱之為第五類接觸.
大家聽過第三類接觸.
就是外星人與UFO齊齊出現.
第四類接觸是什麼?.
就是外星人將人帶上去.
去太空船裡面做研究各樣.
第五類接觸是什麼?.
人主動.
不過不是與外星人接觸.
人主動set aside時間與神接觸.
這是第五類接觸.
我們initial.
因為我們愛上帝.
我們知道應該這樣去親近神.
我們initial提出這個方法.
我們用時間去認識神.
知道祂多些,明白祂多些.
這是第五類接觸.
所以聽師妹,如果你有做靈修的話.
你做靈修的時候.
有沒有意識到神的同在?.
如果沒有的話.
我希望從今天開始.
你再做靈修的時候.
首先你要安靜下來.
求神在你心中向你顯現.
你意識到神的同在.
然後才開始去讀經,去祈禱.
學習怎樣去分享,去看聖經.
與神去討論.
你不需要做聖經老師.
因為那個不是學聖經的時間.
那個是你與神互動,溝通的時間.
多些去問神問題.
在做靈修的時候.
多些與神聊天.
不要讓我們的靈修時間.
變成一個死板的宗教活動.
我再說,做靈修是一個活潑.
與我們天父,爸爸見面.

$^{361}$溝通,互動的一個時候.
我昨晚剛剛從香港回來.
回到十點多了.
又要教宗了.
所以我昨晚只睡了三個小時.
我想說的是什麼,不是說這些.
這次我回來很開心見到我媽媽.
其實是一個主要目的,見一下我媽媽.
我媽媽九十多歲了.
如果大家看過我那個牧人心語的話.
耳朵聽不到,很少說話.
不過我想告訴你.
每次當我見到她,我捉住她的手.
我掃著她的背.
我跟她說話,我知道她聽不到.
不過我看到她很開心.
她好像明白我說什麼.
雖然聽不到,但好像明白我說什麼.
我忽然間有一個感覺.
好像我跟我媽媽那種溝通.
是超乎人的言語.
是盡在不言中的一個交往和溝通.
我發覺我們跟天父上帝聊天.
做靈修也好,怎樣去稱呼也好.
都有類似的情況.
很多時候我的耳朵聽不到神跟我說話.
聽不到祂的聲音.
不過我知道神在跟我說話.
有沒有這樣的經驗?.
我們的肉眼看不到神.
因為神是個靈.
但是我們知道神現在跟我們同在.
有沒有這樣的經驗?.
有沒有感受到神真正在我們身邊?.
有時候我們有些困惑.
在聖經的時候,聽到或者跟別人聊天的時候.
有些我們不明白.
神沒有跟我們解釋.
但忽然間想通了.
忽然間神給我們一個新的體會,新的看見.

$^{401}$我們才能夠明白神的心意.
不知道大家可否觸摸得到,可否體會得到.
原來我們跟神之間的相處是可以這樣的.
弟兄姊妹,如果我們能夠突破了.
即是我們想突破一個知而行難.
或者知而不行的純粹瓶頸.
我覺得我們首先要清楚知道.
神的心意是甚麼.
我們要明白甚麼是真理.
神到底想我們怎樣做.
當我們從聖經,從神的話語裡面領受到.
咦,這些東西真是來自上帝的.
清清楚楚,沒有含糊的.
是上帝真的給我們的教導是這樣.
神真的想我們這樣做.
弟兄姊妹,神要我們做的.
神吩咐我們要這樣走的.
我們怎可以不走?.
你怎可以不聽?.
甚麼令我們這麼大膽?.
至高的上帝吩咐我們做些事.
跟我們說些事,我們就當他沒到.
有沒有這樣的情況? 弟兄姊妹.
求神幫助我們,給些提醒.
所以為甚麼我剛才說.
如果我們要突破純粹瓶頸.
首先我們要培養一個.
尋求明白真理這個習慣.
如果我們能夠知道這是真理.
是真理,是來自上帝的.
我們信耶穌的人怎可以不聽?.
甚麼令我們這麼大膽?.
我們怎可以不去做?.
為甚麼不可以?.
可以的,很多文字就是這樣.
因為我們沒有去尋求明白真理.
好好認識聖經.
於是神啟示的真理.
就變成一個讓我們似懂非懂.
似明非明的教導.

$^{441}$既然你好像聽得懂,又好像聽不懂.
我們不自覺地將真理模糊化了.
你聽到我的意思嗎?.
似懂非懂,似明非明.
將真理模糊化了.
既然真理模糊化了,也不是很清楚.
既然不是很清楚.
就沒必要這麼緊張去做.
這會不會就是我們今天的問題?.
所以保羅教我們.
要詳細察看傳遞使人自由的律法.
詳細察看使人自由的律法.
而且要事尚如此.
甚麼意思呢?.
就是我們要時時去研究去默想.
去尋求明白神的吩咐,神的旨是什麼.
要將這成為一個經常的習慣.
不只是讀經祈禱.
要搞清楚到底神說了甚麼.
神要我們做些甚麼.
甚麼是真理.
甚麼是神要我們做的事.
這是我們突破屬靈生命成長的瓶頸的一個基礎.
如果你相信上帝的話.
為突破屬靈生命成長的瓶頸做好準備.
從牧師的角度,你明白嗎?.
但為何那麼多基督徒就是沒有這樣做.
就是沒有進一步去尋求明白神的真理.
我也在想,我做牧師是因為我希望幫助大家.
我想在聖經裡面有個比喻.
耶穌說的比喻.
大家猜到是哪個比喻嗎?.
是殺種的比喻.
大家記得勝經裡面有殺種的比喻嗎?.
耶穌說的故事.
那故事是說甚麼的呢?.
殺種的比喻,馬太十三章.
「凡聽見天國的道而不明白的.
惡者就來,把殺在他心裡的種子奪了去」.
大家明不明白聖經在說甚麼?.

$^{481}$這些人可能是那星期來教會參加崇拜聽道.
但是無心裝載.
我不是說大家,不好意思.
我是說聖經裡面說的那些人.
那星期來教會聚會,但是無心裝載.
不求甚解,聽完就算數.
沒有用心去尋求明白到底在說甚麼.
聖經說的就是殺在路旁的種子.
然後有些殺在石頭地上.
「殺在石頭地上的,就算人聽了道立刻歡喜領受」.
這是另一種人.
意思是甚麼呢?.
有些人不單止聽見神的話和聖經的教導.
覺得不錯,挺好的.
「只因心裡沒有根,不過是暫時的」.
這就像我剛才所說,沒有進一步去探究.
聽見聖經說得不錯,教導不錯,有道理.
但是沒有進一步去探究.
甚麼意思呢?.
未能將聖經所教導的真理變成我們自己的信念.
未能將聖經的啟示變成我們的價值觀.
於是神的道會在我們裡面生根.
這些人,聖經怎樣說.
「一旦為道遭受患難或迫害,立刻就跌倒了」.
沒有生根,沒有去探究.
還有一些人是甚麼呢?.
「殺在荊棘裡的,就算人聽了道後來.
有世上的憂慮,錢財的迷惑.
把道擠住了,就結不出果實」.
甚麼意思呢?.
是這些人聽了道後.
不單是那些領導更進一步.
還覺得有道理.
聖經說得好.
不單覺得好,願意接受.
但是當我們面對現實生活的壓力.
或者有錢財權力的迷惑的時候.
這些人選擇不去實踐聖經的教導.
遵從神的吩咐.
轉而去順從自己的感覺,慾望,需要.

$^{521}$所以他們屬靈的生命就呆滯不前.
被塞住了,沒有去實踐.
於是結不出果子.
最後終人.
「實在好吐淚滴」.
就算人聽了道明白了.
現代中文聖經的翻譯.
是指那些人聽了信息後並領受了.
不單是明白,而是領悟了.
即是他們不單是明白和接受.
而是領悟了其中的真理.
於是後來結出果實.
有100倍的,有60倍的,有30倍的.
這班人去探究,去明白,去知道.
於是結出果子,因為願意去行道.
我又覺得很奇怪,有些人100倍,有些人60倍,有些人30倍.
很明顯你會發覺,這些人行道的結果是不同的.
有些人多些,有些人少些.
修成的果實是有參差的.
這反映出每個人對神話的領受程度是有不同的.
我們屬靈的生命也不一樣,但無論如何.
這些人是結果子的教徒.
是有生命的教徒.
沒有被卡住,在瓶頸的教徒.
他們不斷有突破.
有些人多些,有些人少些.
有些人快些,有些人慢些.
但這些是繼續在主裡面成長的教徒.
如果大家記得的話.
上次我說突破不求人的瓶頸的時候.
除了這些東西會阻礙我們生長之外.
還有提到肉體的程慾和世俗的衝擊.
包括聖經所說的憂慮,錢財.
或者其他慾望的迷惑.
再加上不求人那種驕傲的心態.
於是我們把神的道卡住了.
我們塞住了讓我們屬靈生命成長的通道.
所以最後,弟兄姊妹.
求神幫助我們.
要培養出不單單是讀經,祈禱的習慣.

$^{561}$要培養出尋求明白,真理去行道的習慣.
如果我們清楚明白神的旨意.
又知道應該怎樣在世界裡面生活.
弟兄姊妹,我們怎樣可以不聽?.
我們怎樣可以不做?.
我們怎樣可以忤逆神的旨意?.
我再說多一次,好不好?.
如果我們清楚知道.
這些是聖經所說的,是來自上帝的.
是神要我們去做的真理.
是聖經給我們的教導.
我們怎樣可以不去行?.
我們怎樣可以不聽?.
可不可以?.
你覺得可不可以?.
當然我們做了,是吧?.
當然可以,沒有什麼說不可以的.
但是不應該,不應該.
不過很可惜,今天有很多基督徒.
包括教會的領袖,不敬畏上帝.
今天有很多基督徒,包括領袖,不敬畏上帝.
That is why.
我聽到有人說,牧師.
我們現在生活在一個欽典時代,有沒有聽過?.
我們生活在一個欽典時代.
我所相信的上帝是滿有慈愛私恩的神.
祂會饒恕我們一切的過犯.
祂會私恩赦免我們一切的罪惡.
大家覺得這是不是聖經的教導?.
你覺得這是不是聖經的教導?.
耶穌怎樣說?.
路加分第十三章.
路加十三章第四節說.
重戰西雷亞.
有一幢樓叫西雷亞樓倒塌了.
什麼原因不知道.
砸死了十八個人.
那些人知道這些人一定比住在耶路撒冷的其他人更加邪惡.
所以神要刑罰他們.
耶穌怎樣說?.

$^{601}$「我告訴你們,不是的」.
不,不是的.
「你們若不悔改,都要如此滅亡」.
耶穌說得很嚴重.
事實上聖經講得很清楚.
「上帝萬不以有罪的為無罪」.
但我們竟然當神沒有到.
我們不敬畏神.
耶穌說如果我們犯罪不去悔改.
不去凝聚神的刑罰.
是必然會臨到我們的身上.
所以主佑弟兄姊妹.
求聖靈幫助我們.
求聖靈照明我們心中的黑暗.
讓我們能夠看到.
自己在什麼事情上得罪了神.
我們不是不知道聖經講了什麼.
很多時候我們不是不明白神想我們怎樣去做.
但我們有意無意的當聽不到當看不見.
我們沒有去實踐神的吩咐和命令.
我們明明得罪了上帝.
但不知為何我們都無所謂.
It's OK.
這是為什麼教會今天這麼軟弱的問題.
可不可以一起祈禱.
天父上帝.
祈求主佑弟兄姊姊聖靈幫助我們.
看到我們對你們的冷漠.
對你們的吩咐和教導.
無心壯哉.
不求甚解.
天父上帝.
祈求你饒恕我們.
只喜歡來聽道查經.
但沒有盡心去思考.
去探索你話語的意義.
我們把你的教導.
只是當作一些不錯的道理.
好像孔子,孟子教對我們的生活哲學一樣.
我們沒有看成是.

$^{641}$來自至高上帝不變的真理.
祈求你饒恕我們.
我們在天上的父母.
祈求你赦免我們.
選擇隨從西族.
過意進行你的吩咐和教導.
祈求你赦免我們.
忙於面對生活的壓力.
多過關心自己.
有沒有去為你爭戰.
完成你交給我們在世上的任務.
天父上帝.
祈求你幫助我們.
能夠培養出尋求明白真理的習慣.
並且不是單單聽道.
更學懂怎樣去行道.
怎樣去實踐真理.
阿們.
我未講完.
不過我覺得我們需要.
在聽道的時間.
如果神的靈感動你.
隨時隨地去祈禱.
我們怎樣能夠行道.
怎樣才能夠實踐.
本來有兩段經文跟大家分享.
《雅各書》第一章二十五節.
「這個人既不是聽了就忘記.
而是實在地行出來.
就在他所行的事上必然得福」.
然後《馬利福音》第七章.
亦都這樣說.
「所以凡聽見我者話就去行的.
好比一個聰明人把房子蓋在盤石上.
雨淋水沖風吹撞著那房子.
房子總不倒塌.
因為根基立在盤石上.
凡聽見我者話不去行的.
就好比一個無知的人.
把房子蓋在沙土上.

$^{681}$雨淋水沖風吹撞著那房子.
房子就倒塌了.
並且倒塌得很大」.
這兩段經文告訴我們.
行道和遵行神的旨意.
是會帶來祝福的.
兩段是這樣說的.
首先帶給我們祝福的是甚麼呢.
在我們所做的事上.
無論神叫你做甚麼.
在我們所做的事上.
會蒙神賜福.
當我看到這件事的時候.
我就想起.
這些在這裡比較長時間的.
弟兄姊妹都知道.
當我們為總統分堂去購物物業的時候.
曾經有過一段風波.
我們曾經為了應該向哪間銀行貸款.
引起一些很嚴重的紛爭.
當時其實我們已經給了.
購買物業的Earnest Money 按金.
但是聖靈照著聖經提醒我.
教會的合一比教會購堂重要.
即是如果我們買了禮拜堂.
但是導致教會分裂.
這個不值得.
我相信這個是不符合上帝的心意的.
在那個時候.
於是就得到教會信託委員.
我們的Trustee 支持這個看法.
於是我立刻叫停了購買的合約.
我們不買了.
如果你還記得那段歷史的時候.
後來我們就買了現在Paris Hilton這個地方.
這個地方比原先我們想要賣的那個地方.
差不多大一倍.
那個好像都不夠二萬呎.
而且那個地點位置和停車的地方.
都比以前那個好很多.

$^{721}$所以讓我能夠想起這個故事.
聖經說進行神教話語還有一個好處.
就是什麼呢.
當我們遇到人生的風暴的時候.
我們能夠站立得住.
不知跌倒.
我們教了二十五年.
在二十五年裡面.
我們首先經歷的是科技泡沫的爆破.
九零年初的時候.
然後經歷金融風暴.
之後世紀疫情.
還有舊在寇唐,建唐.
在時光的發展上面.
遇到很多大大小小的衝擊.
我們不是風平浪靜.
順風順水的.
在我們發展的過程裡面.
遇到很多的挑戰,難處和危機.
我們嘗試努力按照聖經的教導.
去面對和處理這些危機和風暴.
起碼我自己一個人.
原來不只是我自己一個人.
去面對這些挑戰和難處.
No.
原來天父上帝一直在我們的身邊.
他一直在這裡.
我們心急.
我們覺得為什麼會這樣.
為什麼你不做些什麼.
原來神說.
小安無措,未到時候.
當時間到了,神一出手,馬上搞定.
我經歷很多這樣的情況.
當神一出手的時候.
問題不單單可以去處理好.
而且這些危機真的變成轉機.
成為神加添給我們的祝福.
為什麼?.
好像我剛才說的.

$^{761}$因為我們學習.
我們嘗試將教會建立在盤石上面.
這個盤石是什麼?.
就是耶穌基督.
去向我們所啟示的真理.
聖經的教導.
我很希望以後我們教會.
還會繼續面對很多的難處和挑戰.
我們還有危機的.
但只要我們行聖經.
實踐神的話.
聽神說.
我們不單單能夠順利地過渡.
並且我們會見證神的恩典.
比過去的更多和更大.
所以看完節目.
我們知道要行道,是不是?.
知不知道?.
應該知道.
我們大家知道要行道的.
但如果你發覺.
但我好像還未行.
牧師,你說這麼多.
我現在的光景.
我現在的情況.
仍然被卡在知而不行的瓶頸當中.
還有沒有其他原因呢?.
有的.
還有其中一個原因是什麼?.
害怕.
恐懼.
這個也不是我說的.
聖經怎麼說的?.
馬太方用的字眼是.
雨淋水沖風吹.
你看到雨淋風吹風吹.
經常都有的.
聽起來好像不是很嚴重.
不是很重要.
就是一般的風風雨雨.

$^{801}$其實不是的.
如果你看原文的時候.
原文聖經裡面用的字眼是什麼?.
是暴風雨.
是河水泛濫.
大家在電視上看過沒有?.
有暴風雨,河水泛濫.
或者是海嘯.
我們看到那些殺傷力.
很害怕.
害不害怕?.
真的令我們害怕.
丁子梅,我們有心理準備.
當我們純粹行聖經的時候.
當我們進行神的吩咐的時候.
我們真的會遇到一些.
令我們覺得很恐懼.
很震驚的一些事情.
最低限度會令我們覺得不舒服.
譬如聖經教導我們什麼?.
要饒恕你的敵人.
要愛仇敵.
要接納傷害我們的人.
行不行?.
沒有說不行的.
不過我不喜歡.
我不開心.
我不快樂.
我覺得我不行.
或者你上班.
工作的時候.
你拒絕虛報工作的業績.
或者你不去隱瞞顧客.
將真相告訴顧客.
你的同事會將你邊緣化.
會不會?.
甚至會帶來嘗試對你的不滿.
再嚴重一點可能就炒魷魚.
將你解僱.
在家裡面.

$^{841}$我們說愛心.
用愛去包容.
不要經常想著怎樣去操控家人.
或者要贏.
家人都會給予自己競爭.
但是當我們嘗試用愛去包容的時候.
我們會被誤會為溺愛.
或者是無差別縱容.
所以早日聽之會.
不要忘記剛才聖經給我們的應許.
如果我們願意去進行神的吩咐的時候.
實踐神的話語.
我們必然在所行的事上是得福.
是害怕的.
但是你要相信神會祝福我們.
聖經用的字是什麼?.
必然得福.
不單是得福.
是必然會得福.
所以我們要有心理準備.
剛才說過.
當我們立志要行聖經.
要去行道的時候.
我們會遇到像是河水氾濫.
好像暴風雨.
甚至是海嘯一樣的衝擊.
但是神應許我們.
我們總不會倒塌.
我們能夠站得住的.
我們不會跌倒.
我們怎樣才能夠實踐聖經的真理和教導呢?.
我覺得我們要學的是什麼呢?.
將一些抽象的聖經真理和教導.
轉變成一些可以行走的具體步驟.
我們不需要要求自己一步登天.
忽然之間能夠做到的.
譬如大家很熟悉聖經是什麼呢?.
要常常喜樂.
不絕禱告.
凡事借因.

$^{881}$懂得背吧.
你是不是常常喜樂?.
有沒有不絕禱告?.
有沒有凡事借因?.
你知道的.
但是做了多少?.
我們應該怎樣做呢?.
我們大家都想的.
誰不喜歡常常喜樂.
不絕禱告.
凡事借因呢?.
好.
那你們去想吧.
怎樣能夠將它變成一個具體的行動呢?.
譬如說.
每天上課之前.
我意思是睡覺之前.
我剛才叫你寫下一些感恩的事項.
每天寫下一件令你感恩的事項.
這是一個具體的行動.
養成一個習慣.
以前我在香港坐地鐵.
以前我經常讓座給別人.
忽然有一個年輕人讓座給我.
告訴我.
牧師你夠老了.
看你的樣子.
都要讓座給你.
怎會讓座給我呢?.
我當然讓座給師母.
怎會讓我自己坐下.
師母就站在那裡.
誰知道.
我讓師母坐下.
那個坐在旁邊的年輕人.
又立即站起來.
你坐在這裡吧.
一來覺得我真的老了.
第二樣我覺得.
很小事一個人來的.

$^{921}$有什麼大不了.
寫下吧.
行不行呢?.
回去試試吧.
我只是舉例而已.
第二樣.
學習隨時去祈禱.
當我們不和別人說話的時候.
可不可以多些和神傾談呢?.
大家有沒有和神傾談呢?.
我上次說到.
不知道有沒有和你們說過.
神很喜歡和我們傾談.
我們不需要有禱告事項.
有代禱需要.
吃飯前才去祈禱.
隨時隨地.
任何時間任何地點.
想到什麼就說什麼.
你和你的家人.
和你的家人是這樣傾談的嗎?.
有想好嗎?.
除非是很大件事.
否則見到面就說話.
為什麼我們和上帝不可以這樣呢?.
你要培養一個這樣的習慣.
第三.
我教你一個很重要的功課.
學習去笑,行嗎?.
大家懂不懂得笑呢?.
我發覺今天很多人不懂得笑.
每個人拉長臉.
皺起眉頭.
可能身體壓力太大也說不定.
我記得我年輕的時候.
我經常對著鏡子.
什麼笑容是最漂亮的呢?.
哪個角度看會好一點呢?.
現在拍照也是.
這邊好一點.

$^{961}$這邊沒有那麼胖.
這邊好像有點腫.
一定會.
對不起.
大家明白我說什麼嗎?.
大家懂不懂得學笑呢?.
原來有人告訴我們.
微笑可以讓你放鬆你的心情.
而你的笑容可以令人放下戒心.
當我們和別人聊天的時候.
你時時帶著笑容的時候.
你會給人一種很舒服.
很暢快.
很和平的感覺.
但是要練習的.
因為我們很多時候笑不出來.
經常撫著口臉.
壓力太大.
所以對自己也有幫助.
學學怎樣笑得漂亮點.
學習怎樣向人微笑.
時間關係我講這麼多例子.
弟兄姊妹.
如果我們真的喜歡回來教會聽道.
我們喜歡回來教會查聖經.
我們的屬靈生命是不會轉化的.
不會有很大的改變.
你會多了一些聖經的知識.
除非我們能夠行道.
行道有沒有挑戰?.
有沒有難處?.
Definitely.
一定有.
所以我剛才講過兩次.
我們必須有心理準備.
求神賜給我們勇氣.
在行道的時候.
可以去面對挑戰.
和被別人針對.
甚至被別人敵視的困難和恐懼.

$^{1001}$但如果我們愛上帝.
我們真的相信神.
這個神是真的.
我們敬重祂.
弟兄姊妹.
我們應該樂意去聽神的話.
學習將神的吩咐和教導.
變成一些可以實踐的具體行動.
張弟姊妹.
今日你離開這個地方之前.
我希望你記住一件事.
我們每行出一小步.
就是突破我們知易行難.
這個成長瓶頸的一大步.
從小開始.
先想一些你可以做到的事.
慢慢加下去.
你就越來越能夠有力量.
實踐聖經的教導.
求神一令.
親自向你的心講話.
改變你的信仰生活習慣.
讓你能夠更加經歷到神.
當你越來越經歷到上帝的時候.
無論外面的風雨有多大.
你有能力去面對.
並且成為上帝榮耀的見證.
將人吸引回來教會.
你看我周幾多空位.
外面有很多需要聽到這些訊息的人.
弟兄姊妹.
你和我就是那個媒介.
可以將人帶回來.
求神一令我們去教會.
我們一起祈禱.
天保上帝我將這事情交在手中.
孩子能夠想到要講的都講了.
我講得不清楚的.
求聖靈你去保祝.
求聖靈你在我心裡面跟我們講話.

$^{1041}$讓我們明白自己的光景.
到底我們信耶穌那麼久.
我們是哪一類人.
我們有沒有被卡住在一個.
知而不行的瓶頸當中.
以致我們的經歷越來越少.
求天父你恩代我們.
賜力量給我們.
突破這個贖利成長的瓶頸.
奉耶穌基督的聖名.
\newpage



\section{出埃及記 17:1-7}
\label{sec:031775}
\textbf{誰為妳受擊打?}~~~~林玉英牧師
\newline
\newline
連結: \href{http://www.church.com.hk/acms/content.asp?site=cdc\&op=show\&type=product\&code=031775\&layout=sermon}{\texttt{031775}} ~~~~ 語音日期: 2026-03-08
\newline
\newline
\hyperref[sec:031774]{\small{< < < PREV SERMON < < <}}
~
\hyperref[sec:index_chronic]{\small{[返順時目]}}
~
\hyperref[sec:031776]{\small{> > > NEXT SERMON > > >}}
\newline
\newline
出埃及記 17:1-7
\newline
\begin{longtable}{cl}
\hline
\hline
章節 & 經文 (和合本修訂版)\\
\hline
17:1 & \begin{tabularx}{0.7\textwidth}{X} 以色列全會眾遵照耶和華的吩咐，從汛的曠野一段一段地往前行。他們在利非訂安營，但百姓沒有水喝。 \end{tabularx} \\ \\ \relax
17:2 & \begin{tabularx}{0.7\textwidth}{X} 百姓就與摩西爭鬧，說：「給我們水喝吧！」摩西對他們說：「你們為甚麼與我爭鬧呢？你們為甚麼試探耶和華呢？」 \end{tabularx} \\ \\ \relax
17:3 & \begin{tabularx}{0.7\textwidth}{X} 百姓在那裡口渴要喝水，就向摩西發怨言，說：「你為甚麼把我們從埃及領出來，使我們和我們的兒女，以及牲畜都渴死呢？」 \end{tabularx} \\ \\ \relax
17:4 & \begin{tabularx}{0.7\textwidth}{X} 摩西就呼求耶和華說：「我要怎樣對待這百姓呢？他們差一點就要拿石頭打死我了。」 \end{tabularx} \\ \\ \relax
17:5 & \begin{tabularx}{0.7\textwidth}{X} 耶和華對摩西說：「你帶著以色列的幾個長老，走在百姓前面，手裡拿著你先前擊打尼羅河的杖，去吧！ \end{tabularx} \\ \\ \relax
17:6 & \begin{tabularx}{0.7\textwidth}{X} 看哪，我要在何烈的磐石那裡，站在你面前。你要擊打磐石，水就會從磐石流出來，給百姓喝。」摩西就在以色列的長老眼前這樣做了。 \end{tabularx} \\ \\ \relax
17:7 & \begin{tabularx}{0.7\textwidth}{X} 他給那地方起名叫瑪撒，又叫米利巴，因為以色列人在那裡爭鬧，並且試探耶和華，說：「耶和華是否在我們中間呢？」 \end{tabularx} \\ \\
[1ex]
\hline
\hline
\end{longtable}
$^{1}$各位弟子們平安.
邀請大家再次低頭祈禱.
去招待我們在你的聖殿敬拜.
我們知道部分的弟兄姊妹妹.
因為身體或其他原因.
雖然現在很渴想在殿中.
但未能夠.
要透過科技在網上參與.
無論我們在哪個位置去崇敬你.
求你閱立.
我會對你的渴想.
對你的敬拜.
由於讓你的訊息不是從人而來.
而是你透過聖靈.
在我們心裡去動工.
打開我們的耳朵.
打開我們的心.
來領受你要對我們的說話.
成為我們生命的力量.
指引和方向.
由於講的和聽的都同往你的願.
祈禱奉主耶穌 基督的名求.
阿們.
各位弟子們平安.
今天我們要分享的經文.
剛才少爺姐妹都為我們讀出.
發生在一個地方叫利菲頂.
在希伯來文之中.
利菲頂這個名字的原意.
很漂亮很漂亮.
意思是休息的地方.
想想.
以色列.
我們記得它出埃及之後.
走了一段很長很長很長的路.
而每一站無論是起行拔營.
去哪裡.
在哪裡停留扎營.
它們一直都是遵照耶和華的吩咐.
大家記不記得我曾經.

$^{41}$去和大家分享這個文掃記的時候.
在整個以色列.
它在曠野的路程裡面.
運柱和火柱成為.
耶和華上帝同在的證明.
當運柱和火柱移動了.
它們就會移動.
停的時候它們也會停.
所以經文去到這裡告訴我們.
它們一直是遵照耶和華的吩咐.
於是它們來到利菲頂安邑.
在它們心目中.
這應該是一個可以躺臥的青草地.
有安潔水源的地方.
很正常和合理的期待.
不過實際並不是這樣.
在這個休息的地方竟然沒有水喝.
各位靈姐妹.
當你遇到不似預期的狀況的時候.
你的心情會是怎樣呢?.
而你又會有什麼反應呢?.
我在想像不像我們人生的寫照呢?.
我們以為結了婚.
就會有休息和安穩.
結果就有很多的家務.
還有想不到的磨擦.
我們以為升職.
就會有安穩和安息.
結果就是更大的壓力和責任.
各位靈姐妹.
不知道你在現在的人生階段.
你最覺得自己休息不了.
心裡面有一個很大的渴想.
Hunger.
你的Hunger for.
你要渴想的是什麼呢?.
會不會是你努力付出之後.
得不到你期待的回報?.
還是有一份你看不到未來的焦慮呢?.
當你的需要得不到滿足的時候.

$^{81}$當你的期待不似預期的時候.
你會怎樣呢?.
我們在經文裡面見到.
百姓的反應是向摩西爭鬧和埋怨.
他們覺得自己的要求很合理.
「上帝,我服從你.
我聽到你說來到這裡扎營了.
為何連水都沒有?」.
他們將上帝的供應看作理所當然.
「你不是告訴我.
你是我們的牧者.
我們必不自缺乏嗎?」.
沒錯,我們相信,我們大家都同意.
百姓的需要是很合理.
因為這個是人類基本生存的需要.
他們不是妄想.
這個是他們真正的需要.
但我們要留意的是.
他們的態度就很不合理.
他們埋怨摩西.
「你為何要將我們從埃及領出來.
讓我們和我們的子女.
以及我們的生畜都在這裡渴死?」.
他們還懷疑.
耶和華是不是還在我們中間?.
這句說話的意思是甚麼呢?.
原來上帝的存在和真實.
必須要透過即時滿足我們的需要.
來證明上帝.
如果上帝不按照我的時間表.
不給我想要的東西.
他就不夠資格做我的上帝.
當我們將上帝看成是有求必應.
又或者祂是我們的公人.
我們就會失去對上帝應有的態度.
一份敬畏祂的心.
慢慢我們會變得很傲慢驕傲.
變得心癢.
將人和神的關係調轉了.
不過自己也不太察覺.

$^{121}$認為上帝應該要聽我的話.
應該聽我的禱告.
滿足我們的需求.
慢慢我們口裡說.
「信靠順服上帝」只不過是一個口號.
而不是我們要努力實踐的目標.
這真的很危險.
我們要謹慎.
我們要留意.
摩西在其中的反應非常真實.
「你們為何要和我爭辯?」.
「你們為何要試探耶和華?」.
摩西覺得他很無辜.
很委屈.
甚至覺得百姓是無理取鬧的.
「當初不是你們嗎?」.
「不是你們向上帝呼冤」.
「渴望上帝拯救你們離開埃及做奴隸的地方嗎?」.
「現在橫行霸道」.
「你們忘了順理成章地讓法老」.
「心甘情願地放你們離開埃及嗎?」.
「當你們離開埃及的時候」.
「你向身邊的埃及人拿甚麼給你」.
「只是希望你們快點離開」.
「當是過紅海的時候你們親身經歷的神蹟」.
「日落在曠野沙漠無水一無食物的地方」.
「上帝不是一一供應給你們嗎?」.
「就在上一章而已」.
「你們說沒有肉吃」.
「於是就有成千上萬的鵪鶉」.
「來飛離你們的營地」.
「這麼快就忘記了?」.
「你們再次來試探上帝?」.
不過百姓群情洶湧.
甚至想用石頭來打死摩西.
這是人經常在壓力和焦慮之下.
有的失智和失控的表現.
摩西做對了一件事.
他不再與百姓爭拗.
他轉向呼求上帝.

$^{161}$他意識到自己不是這群人的老闆.
他只不過是一頭領頭羊.
沒錯 他走在前面.
不過他只不過是領頭羊.
他只不過是一個中介人.
沒錯.
是上帝差派他領導以色列人離開埃及.
但實際上.
他只不過是上帝與以色列人之間的中介人.
他不是上帝.
他根本沒有能力在沙漠這個甚麼都沒有的地方.
去變出水給百姓喝.
各位弟兄姊妹.
無論你是在家庭,職場或教會擔任甚麼職能,職事.
又或者你是一個領袖.
請你記得.
你不需要亦不能代替上帝去解決大家的困難.
當你感到很大很大的壓力.
大到你想躲起來想逃避的時候.
當你覺得很無辜很委屈的時候.
記住學摩西.
將問題和困難.
帶到呼召你事奉的上帝面前.
因為上帝才是大家的「波士」.
上帝垂聽了摩西的禱告.
他對摩西有一個吩咐.
藏著一個很深奧的恩典.
他要摩西拿著那支曾經擊打尼羅河.
象徵著審判.
亦象徵著領導的那支手掌.
去擊打河裂的盆石.
在第六節上帝說.
我要在河裂的盆石那裡.
站在你的面前.
原來在古代站在一個人面前.
通常都是一位被告.
他這位被告對著審判他的法官.
站在那裡接受審判.
百姓埋怨上帝 試探上帝.
本來是要承受懲罰和擊打的.

$^{201}$不過上帝卻選擇站在受罰者這個位置上.
去代替百姓.
當摩西拿著那支審判的手掌.
去擊打盆石上的時候.
其實是打在上帝的身上.
《普臨海·哥林多》前書十章.
講述以色列人在曠野這幾十年的經歷.
特別提到這一段的歷史.
在第四節保羅說.
那個盆石就是基督.
讓我們知道.
上帝給世人的是一份對稱的救贖.
我們跟百姓一樣犯了罪.
不過受懲罰被擊打的.
他是上帝主耶穌基督.
上帝甘願站在被擊打的位置.
去做代罪的羔羊.
《儀場亞先》在《儀場亞書》第53章第5節這樣說.
他為我們的過犯受害.
為我們的罪孽被壓傷.
因他受的懲罰.
我們得到平安.
因為他受的鞭傷.
我們得到醫治.
盆石被擊打之後.
流出活水.
解除了百姓的乾渴.
各位聽眾.
你明不明白.
你體不體會得到.
上帝對我們這一份心恩厚愛呢?.
如果我們真的捉得住.
上帝給我們的這份愛.
我們就不會向人去拿水喝.
什麼意思?.
其實很多時候我們底層的渴望.
我們很乾渴的地方.
其實是想得到.
其他人對我們的認同.
欣賞和肯定.

$^{241}$甚至有時候我們都會乞求對方.
給我們這份接納.
不過主耶穌說.
我已經為你受了擊打.
在十字架上所流出的寶血.
是我對你愛的證明.
無論你是一個怎樣的人.
又或者無論你怎樣看待你自己.
我永遠都會接納你.
你是我的愛子.
我愛你.
沒錯就是這一份承諾.
帶給我們平安.
能夠填滿我們心靈底層的饑渴.
各位弟兄姊妹.
你捉不捉得住上帝這一份的愛呢?.
今天是婦女主日.
我相信當初設立婦女主日的目的.
除了要紀念往往被忽略的婦女貢獻之外.
其實我們更加要反思的是.
我們身邊有很多人.
無論男女.
是否只是處於某一種隱形的饑渴當中呢?.
很多人在家庭.
在職場.
以及在教會當中.
時時被要求像盤石一樣的穩固.
不斷要供應.
但是很少人去問他.
你口不口渴?.
亦很少人關心他們有什麼需要.
其實我們每一個人.
無論什麼身份.
什麼地位.
什麼性別.
什麼種族.
都需要別人的尊重.
和公道的對待.
什麼是尊重和公道的對待呢?.
不是坐在高位上指責別人.

$^{281}$為什麼沒有水喝?.
而是學校上帝主動走到受擊打的盤石位置上.
與其成為那個只懂埋怨指責試探神的人.
我鼓勵大家願意成為那個站在盤石上的.
難難.
超難.
所以求上帝給我們一份願意的心.
賜給我們力量.
去見證別人的辛勞.
去主動去分擔別人的重擔.
讓他們不再在孤單裡面來口渴.
各位弟兄姊妹.
上帝帶我們來到沒有水的利菲頂.
不是玩我們.
而是要幫我們去看清我們屬靈的狀況.
去幫我們屬靈的生命去除雜質.
修剪那支阻礙我們結果子的部分.
記得在《楊芳十五章》這樣說.
上帝是修剪去栽培葡萄樹的人.
而我們就是枝子.
當有些地方不結果子的時候.
祂會修理乾淨.
因為當我們在乾渴的時候.
在缺乏的時候.
我們才會發現.
我們對上帝的信心.
是否只是建立在祂給我們物質的供應上面呢?.
我們是否看清楚.
我們是否能夠捉住上帝給我們那份不對稱的愛呢?.
所以弟兄姊妹請你記得.
當你覺得最口渴.
最無力.
最想放棄的時候.
請你望一望主耶穌.
望一望那塊盆石.
主耶穌基督.
在那裡有位上帝.
祂正站在你的面前.
去甘願為你受擊打.
好讓你的生命可以流出活水.

$^{321}$上帝的活水能夠流入你的心靈裡面.
無論你向人去埋怨也好.
你指責其他人也好.
你試探上帝也好.
都不能夠解除我們心靈的深處底層的乾渴.
唯有來到上帝的面前去呼求祂.
呼求祂給你活水來喝.
呼求祂幫你捉得住.
祂要給你這份心恩厚愛.
耶穌在《約翰福音》裡面這樣說.
祂說「誰喝了我所賜的水.
就永遠都不口渴.
我所賜的水要在他裡面成為泉源.
直湧到永生」.
喝了耶穌所賜的活水.
你會有源源不絕的動力.
在這個好像沙漠曠野的人生旅程當中.
我們不再懼怕.
不再憂慮.
也不再迷惘.
可以活出美好的人生.
並且會有滿足的喜樂.
讓我們同心祈禱.
親主你知道我們每一個人的狀況.
你知道我們每一個人此刻所面對的困難.
上主懇求聖靈在我們心裡面動功.
來感動我們.
以致我們不單止在頭腦裡面知道你愛我們.
而是在心靈的深處感受到.
你對我們的心恩厚愛.
當我們能夠捉住你對我們的愛的時候.
面對世間一切的困難.
挑戰.
迷惘.
讓我們焦慮和沒有安全的境地.
我們都因為捉住你對我們的愛.
而有一份安全感.
有一份動力去面對一切的挑戰.
最願意摩西成為我們的榜樣.
面對很多艱難的裡面.

$^{361}$我們來到你的面前.
去呼求你的拯救.
你必然開我們的眼睛去看到出路.
去看到我們應該要走還是站.
要堅持還是放棄.
因為你是指引生命路途的那位上帝.
主我們記得.
我們的好處不在你以外.
你會帶領我們經過死暗的幽谷.
雖然眼前會有很多令我們覺得很害怕.
很迷惘.
很不知怎樣算不敢跳過去的死暗幽谷.
但我們知道只要我們.
聽得清楚你對我們的吩咐.
那一條會是美好的道路.
求上帝你就這樣和我們同在.
賜我們信心.
去行每一段生命的路程.
我們這樣祈禱過.
藉主耶穌基督的名而求.
阿們.
\newpage



\section{創世記 27:13}
\label{sec:031776}
\textbf{聖經人聖經事:我兒,咒詛歸我}~~~~黃克勤牧師博士
\newline
\newline
連結: \href{http://www.church.com.hk/acms/content.asp?site=cdc\&op=show\&type=product\&code=031776\&layout=sermon}{\texttt{031776}} ~~~~ 語音日期: 2026-02-27
\newline
\newline
\hyperref[sec:031775]{\small{< < < PREV SERMON < < <}}
~
\hyperref[sec:index_chronic]{\small{[返順時目]}}
~
\hyperref[sec:031777]{\small{> > > NEXT SERMON > > >}}
\newline
\newline
創世記 27:13
\newline
\begin{longtable}{cl}
\hline
\hline
章節 & 經文 (和合本修訂版)\\
\hline
27:13 & \begin{tabularx}{0.7\textwidth}{X} 他母親對他說：「我兒，你所受的詛咒臨到我身上吧！你只管聽我的話，去牽小山羊來給我。」 \end{tabularx} \\ \\
[1ex]
\hline
\hline
\end{longtable}
$^{1}$弟兄姊妹大家好.
又來讀聖經.
聖經人士.
我們由當中去看背後上帝奇妙的作為.
以及看人性有時很軟弱.
有些事也不錯,有時很軟弱.
生命說書的網絡為大家提供.
今天是我兒就坐歸我.
一講我兒,你知是父母講的.
這句說話是利伯家講的.
我們看看利伯家發生了什麼事.
我兒就坐歸我.
繼續講的是.
亞伯拉罕到以撒.
以撒不是很多事情.
以撒最多事情發生.
當然是他被父親亞伯拉罕帶上山.
耶和華以勒.
另外就是他臨終之前.
吃夜美,又要祝福.
亞伯拉罕欺騙哥哥出事.
現在看看.
27章第五節.
以撒對他兒子以素說話.
利伯家也聽見了.
以素接著出了去田野打獵.
要得夜美帶來.
利伯家就對他兒子亞伯說.
我聽見你父親對你哥哥以素說話.
其實這是什麼事呢?.
如果你還記得昨天.
以撒對兒子說.
兒子啊.
兒子就說:嘿呢,我在這裡.
認他.
兒子就說:嘿呢.
兒子就認他一樣.
看看我,已經很老了.
我都不知道什麼時候死.
你和我去打獵.

$^{41}$就打些獵物給我做食物.
弄好它.
我就為你祝福.
其實我都說.
以撒祝福就祝福.
祝福之後才叫他打獵.
那就沒事了.
叫他打獵回來吃了才祝福.
給這個條件他真是害事了.
其實有時候都教訓到我們做事直接些.
喂,先搞好這件事.
先不要說吃.
先不要吃.
兒子進來,我先幫你祝福.
不死了,我怕自己不知道哪天死.
幫你祝福.
祝福完之後.
打點東西給你吃.
他都馬上去打.
所以沒所謂了.
誰知就出事了.
他就出去打獵了.
以素往田野去打獵.
誰知以撒與他兒子以素說的話.
以撒的太太尼伯加都聽到.
馬上和尼伯加疼愛的小兒子雅各說.
兒子,我聽到你爸爸和哥哥說話.
你把野獸帶來做成美味給我吃.
我好在未死之先.
在耶和華面前為你祝福.
這個祝福我上次說了.
這個祝福很重要.
這個是在耶和華面前為你祝福.
即是上帝介入在這件事裡面.
另外以撒知道和承接著.
亞伯拉罕對他納的約.
這個祝福很重要.
現在,我兒.
你要照著我所吩咐你來聽從我的話.
這是第一點想跟他說的.

$^{81}$兒子,現在你照著我說的話做.
即是說,尼伯加出主意了.
出主意做什麼?.
要介入在這件事裡面.
要做介入.
攔阻爸爸為哥哥祝福.
要為弟弟祝福.
這件事你會見到.
這個一字,尼伯加的介入.
其實是不好的.
為什麼不好呢?.
第一,產業不是自己的.
是老公的.
老公要給兒子.
為什麼我要阻止他呢?.
我疼愛小兒子.
我想老公給小兒子,不要給大兒子.
其實這是不公道的.
好了,你聽我說.
你去羊群那裡.
給我拿兩隻肥山羊羔來.
為什麼不拿綿羊羔呢?.
羊肉不同,山羊羔可能比較鬆.
而且羊肉比較像野獸.
所以就找山羊羔.
不過要肥,因為要羔才好吃.
沒有羔,只有羔才好吃.
我便照你父親所愛的.
給他做成適合他的味道.
你便拿到父親樂裡給他吃.
使他在未死之先給你祝福.
所以這個計謀.
是他馬上一閃就做了.
其實李百嘉的計謀很厲害.
很聰明的女孩子.
很有辦事能力的女孩子.
就嫁了,你見她打水的時間很殷勤.
她從小到大已經很厲害.
雅各對他媽媽說.
這計謀也好,不過.

$^{121}$他媽媽說不行,為甚麼呢?.
他哥哥耳掃渾身都是毛.
他一出生就渾身都是毛.
我身上光滑,剛好沒有毛.
如果父親摸著我.
就必定認為我是欺控人.
我反而遭到咒作不得祝福.
他以為我欺騙他,其實是在欺騙他.
很有趣的.
他哥哥說不行,因為父親摸我的手便知道.
如果他說我欺騙他.
他就在我中斃,根本不會有祝福.
還遭到咒作.
這句說話大家可以在這裡學到.
當你欺控人的時候.
你未曾得到欺控而帶來的好似福分.
你是先遭到咒作.
所以要留意.
這裡也教到我們.
怎知利伯加說.
第二點.
我言,你遭的咒作歸到我身上.
你只管聽我的話.
去將羊羔給我就得.
兒子,這些事讓我處理,大人的事你不用處理.
咒作歸我.
媽媽幫你揹,我幫你頂.
你只管聽我的話.
這些說話.
其實神是聽到的.
在聖經裡很多時候.
人說這些說話.
神是計數記得的.
如果真的咒作會歸到他.
譬如彼拉多說.
你們要釘耶穌十字架.
如果流無辜人的血.
我查不到他有什麼叛亂的罪.
如果流無辜人的血怎麼辦?.
當時的大祭司,祭司,文士,法利卓人和百姓.

$^{161}$說這件事歸到我們和我們子孫身上.
我們頂.
哇,上帝聽到的.
以後他們很慘的.
發覺流人血的事歸到他們身上.
都很慘,神真是聽的.
讓我頂.
這些都是媽媽.
你得福得,我落地獄都濟.
其實尼伯加很疼愛雅各.
忽然間心急起來用這個藉口去騙爸爸.
其實有沒有別的路可以走呢?.
如果知道以撒這麼疼愛自己.
不如尼伯加走進去以撒那裡.
老公,我剛剛走過帳幕聽到你這樣說.
祝福耳掃,不如分一點給小兒子雅各.
好啊,一起祝福吧.
可能都會有一點.
他不理會,如果他不知道紅豆湯事件.
其實祝福大兒子是對的.
雙倍祝福,長子名分.
他就說不要緊,就著就著,我頂.
你去聽我話做.
於是雅各就拿了兩隻肥山羊羔.
他母親就照父親所愛的口味做成了美味.
第三,尼伯加真的使得那些計劃.
想著名額.
尼伯加又把家裡所傳大兒子.
耳掃上好的衣服給小兒子雅各穿上.
為何要找耳掃上的衣服給他穿呢?.
真是不同的.
我相信雅各的衣服應該柔軟一點.
舒服一點.
因為他一直在家裡,是宅男,媽媽很疼他.
耳掃是出去打獵.
出去打獵的人穿的衣服會不會很柔軟,絲質?.
當然不會.
會不會有很漂亮的棉質?.
當然不會,當然是穿麻衣,粗身.
頂得,刮到不會損,掉到地上不會爛.

$^{201}$而且當日洗燙都不會很厲害.
於是就會有田野泥味,草味.
拿耳掃身上漂亮的衣服給他穿.
又用山羊羔皮包在雅各的手上.
特別猜到他會摟他親嘴.
所以包上頸也很厲害.
在光滑的地方鋪上羊毛,羊皮.
並且將美味和餅.
餅和食物變成一餐飯.
交給他兒子雅各.
於是就這樣發生了.
你要留意,這張聖經令我們有所感動.
在《創世記》這段人與故事裡.
你會見到尼伯加愛子深切.
其實雅各也承接了他母親尼伯加.
是屬於有鬼才,很多妓女能取得好處.
其實母子倆合謀欺騙父親和丈夫.
而且他們是否知道會招來咒詛呢?.
是知道的.
雅各這個兒子知道,媽媽這個兒子也知道.
不過媽媽說:不怕啦兒子.
我兒如果有咒詛就歸到我身上.
這個很疼愛的兒子也不是理由做出欺騙父親的事.
因此這個兒子是不對的.
聖經照記載在這個重要的族長家庭.
以色列族長家庭.
族長的意思我們叫做族長時代.
阿巴拉罕以薩雅各是族長時代.
這個族長時代的污點.
是尼伯加用妓女去欺騙以薩.
給大家了解這件事.
你只要聽他拿羊羔給我.
咒詛歸我.
其實咒詛真是歸他.
因為他很快會見到.
雖然那一刻先勝一招.
但之後兩個兒子不對.
以薩,尼伯加,以素雅各四人.
都變成很淒涼.
大家關係破裂了.

$^{241}$這是一個馬上的咒詛.
還有其他咒詛可能內心很不舒服等等.
記得要避開這些咒詛.
很重要的.
正正直直祈禱.
求主耶穌幫助我們.
你教導我們.
是就說是.
不是就說不是.
若再多說.
一定是有一個來自邪惡的理由.
主我們在這個時間.
求你幫助我們有一個勇氣.
可以直言.
可以說出.
可以去仰望上帝.
其實無論怎樣.
你有你最好的大靈.
你有你最好的賜福.
免得我們用自己的手段之後.
不獲得你的祝福.
這樣更慘.
主求你自己真的幫助.
讓我們躲避離開這些咒詛.
尋求你自己真正的祝福.
誰聽禱告奉耶穌基督的名.
阿們.
\newpage



\section{尼希米記 1:1-4}
\label{sec:031777}
\textbf{生命無牆,好快收檔 (1)}~~~~顏沛恩牧師
\newline
\newline
連結: \href{http://www.church.com.hk/acms/content.asp?site=cdc\&op=show\&type=product\&code=031777\&layout=sermon}{\texttt{031777}} ~~~~ 語音日期: 2026-03-08
\newline
\newline
\hyperref[sec:031776]{\small{< < < PREV SERMON < < <}}
~
\hyperref[sec:index_chronic]{\small{[返順時目]}}
~
\hyperref[sec:031778]{\small{> > > NEXT SERMON > > >}}
\newline
\newline
尼希米記 1:1-4
\newline
\begin{longtable}{cl}
\hline
\hline
章節 & 經文 (和合本修訂版)\\
\hline
1:1 & \begin{tabularx}{0.7\textwidth}{X} 哈迦利亞的兒子尼希米的言語如下：亞達薛西王二十年基斯流月，我在書珊城堡中。 \end{tabularx} \\ \\ \relax
1:2 & \begin{tabularx}{0.7\textwidth}{X} 那時，我有一個兄弟哈拿尼，同幾個人從猶大來。我問他們那些被擄歸回、剩下殘存的猶太人和耶路撒冷的情況。 \end{tabularx} \\ \\ \relax
1:3 & \begin{tabularx}{0.7\textwidth}{X} 他們對我說：「那些被擄歸回剩下的餘民在猶大省那裡遭大難，受凌辱；耶路撒冷的城牆被拆毀，城門被火焚燒。」 \end{tabularx} \\ \\ \relax
1:4 & \begin{tabularx}{0.7\textwidth}{X} 我聽見這話，就坐下哭泣，悲哀幾日，在天上的神面前禁食祈禱， \end{tabularx} \\ \\
[1ex]
\hline
\hline
\end{longtable}
$^{1}$弟兄姊妹早晨.
係呀今日ge 講題有D 特別ge .
係生命無牆跟住做咩呀.
好快收檔呀仲要係第一喎咁即係講緊呢.
我會有第二次ge 續集ge 所以今日呢係一個開場白黎ge .
好啦咁呢弟兄姊妹你有冇玩過塵塵碟呀.
玩過呀舉手黎睇下非常之刺激好興奮架喎所以你下一個禮拜呢破冰ge 時候.
可以玩呀玩ge 時候呢滑心都黎一黎手都震埋呀因為呢你個人呢要非常之.
手定心定眼定仲要腦袋要定.
仲要呢你睇清楚呢上面邊一層呢gwo D 木呢可以慢慢咁推出黎呢你要.
確保你唔好整lum 佢呀嘛係咪但係呢往往呢玩到最底gwo 層.
得返中間gwo 條木ge 時候呢係難度最高架因為呢你係上面呢.
一個角落你都會點呀?好驚呀,好驚佢lum 架,不過呢最興奮最刺激ge 時候到呢,就係乜ye 呀?.
就係點話?大家都唔係幾興奮架,就係一ye gwo 個層次做咩?.
lum 落黎gwo 下呢,就全世界ge 人類就做咩架?gwo 下就會尖叫啦,大嗌啦,嘩一聲啦.
不過呢,係尼希米時代ge 時候,成場lum jor 落黎,係一D 都唔興奮架,就好似我地頭先讀ge 經文一樣喎,就係尼希米記ge 一章,一字到四字,我地呢再慢慢睇一睇呢段經文講乜ye 先.
佢話呢,亞達舌西王二十年,gwo 個係基斯留月,就係講ge 係波斯帝國主前四四六年ge 十一至到十二月啦,佢話呢,我係書山城,書山城堡ge 裡面,gwo 時有一個ge 兄弟叫做哈拿尼,同幾個人呢,係呢個ge 猶大黎啦.
我呢就問佢地啦,gwo D 被擄歸回,靖下殘存ge gwo D ge 猶太人,做jor 耶路撒冷ge 情況,究竟而家係點架?.
弟兄姊妹你要記得呢,當時呢已經係第三批ge 猶太人,佢地被擄歸回返屋企ge 時候啦,然之後佢地就話啦,gwo D 被擄歸回,靖下ge 漁民,係猶大省gwo 度遭遇乜ye 呀?.
大能,跟住仲咩呀?受凌辱,跟住呢耶路撒冷ge 城牆係做咩呀?最底gwo 行,被做咩呀?拆毀jor ,跟住城門做咩事呀?被火燒喎,下一版再睇埋先.
跟住尼希米聽見者話,就坐低哭泣悲哀幾日啦,跟住呢佢係天上ge 神面前禁食禱告,於是尼希米一路禁食禱告,禱告,禱告,禱到幾時呢?.
就係二章一節啦,就係亞達薛西王ge 二十年gwo 個呢叫做彌撒月,就係主前ge 四四五年三至到四月,即係講緊佢一路禁食祈禱,祈jor 四個月.
尼希米聽到白色壽福,佢好傷心,佢一個好難過ge 時候,佢坐低係度喊悲哀幾日禁食祈禱,一路祈,祈祝四個月,點解呀?就係因為當時城牆已經拆毀jor 啦,城門已經被火燒.
但係對於尼希米黎講係一件好大件事黎,係一件好嚴重ge 事,係以色列民受辱,係一個好大ge 羞恥,當時呢班人已經冇尊嚴冇地位ge 時候.
弟兄姊妹呀,你要知道呢,佢地而家係被擄做咩事呀,佢地而家返屋企呀,可以返屋企啦,終於返返自己屋企耶路撒冷ge 時候,唔使再流離啦,唔使做苦工啦,唔需要睇人面色啦,唔需要逆來順受,唔使寄人籬下,我終於可以返屋企.
但係返屋企,我有屋企呀,但係冇城牆,城門又冇埋,被火燒,講緊乜ye 呀,即係等於你有屋企,但係你屋企冇門架喎,又冇鎖架喎,乜ye 人都可以走入去你屋企gwo 度,即係講緊咩呀,D 人可以偷曬你D ye 啦,甚至乎佔據jor 你屋企ge 位置啦,霸佔jor 你屋企ge 地盤.
耶路撒冷冇城牆,冇城門,點樣防衛呀,一攻ge 時候打lum 晒D 人呀,全部可以死晒呀,擄掠晒佢所有ge ye 啦,耶路撒冷可以移為平地,所以城牆,你睇下呢幅圖,圍住成個耶路撒冷.
城牆係表達緊一份好安全,一份好大能ge 同在同埋保護,基本上係360度圍住成個耶路撒冷係好安全同在ge 地方,城門係要lock架,係要上鎖架,係表達緊一份ge 保障,乜ye 人出入都好啦,都要安檢架喎.
非常安全,係一個民族住係耶路撒冷,係一份生活非常之滿足ge 一個自豪感黎ge ,係一個滿足感黎ge ,就好似我地北上ge 時候,弟兄姊妹你有冇北上呀,有冇北上呀,好鍾意北上ge ,舉下手黎睇下先,三樓冇人認我,其實暗中已經去jor .
就好似我地北上ge 時候,過關咪要安檢ge ,咁人又要檢啦,跟住貨又要檢啦,一次仲記得你知啦,我有架買sung ge Benz架嘛,即係其實我都識zar 車ge ,不過係買sung gwo 架車姐,買sung ge Benz一擺落去gwo 條輸送帶gwo 度就過啦,跟住之後個人咪係呢度安檢架嘛.
突然嗶一聲,發生咩事呢?就我竟然傻豬豬,唔記得jor 將我有個包包呢,除jor 落黎呢,都要安檢架嘛,所以當你嗶一聲ge 時候,全世界ge 人,全部就會望住你啦.
但係耶路撒冷,城牆被拆,城門被火燒,佢地當時失去jor 一個好大ge 同在同埋保護,耶路撒冷殘缺,外敵好自然就入侵,自出自入,被擄ge ,而家可以返屋企,但係點解返到耶路撒冷呢個屋企,繼續會被仇敵黎到去操控,仍然冇改變過呀.
佢地繼續過緊一個被擄ge 生活,神ge 心唔係咁樣,尼希米好清楚知道,如果唔再重建呢個城牆ge 話,百姓就會失去jor gwo 份同在同埋保護.
就有句金句咁樣講,盜賊來,跟住做乜ye 呀,盜賊來,跟住做乜ye 呀,無非要做乜ye 呀,偷竊殺害毀壞呀,再黎多一次,大聲多一次,準備起,盜賊來,無非要偷竊殺害毀壞.
一個冇城牆ge 人,講緊信仰根基唔穩陣ge 人,都唔知自己信緊乜ye ge 人,就算知道自己信乜呀,但我同神ge 關係好遠好疏離呀.
或者一個對自己都唔認識自己ge 人,講多次,一個對自己都唔認識自己內裡面最deep真實gwo 一面ge 人.
受敵就可以係佢身上面最軟弱ge 地方偷竊殺害毀壞,可能係我地ge 眼睛唔知見到D 乜ye ,睇見色情ge ye ,睇見網絡,可能係我地ge 腦袋裡面裝有好多負面扭曲ge ye ,一D 唔乾淨ge 幻想.
可能係我地心窩裡面有好多錢財ge 迷惑,物質裡面ge 私慾,甚或係我地成長ge 裡面,長大ge 時候好多驕傲,自卑,比較,妒忌,嫉妒呢D ge ye ,甚至關係裡面ge 破口呀.
婚姻裡面ge 上司與下屬ge 婆媳ge 關係,同輩裡面ge 關係,甚或係教仔呢,一個太太一個丈夫兩個教仔ge 價值觀唔同ge 時候就可以鬧大交啦.
只係一個小小ge 破口,我已經唔係講緊一個好大ge 窿,只係講緊一個小小ge 破口,罅惡者就已經有位可以攝到入黎,佢就可以自出自入,個家就會被破壞啦.

$^{41}$頭先講緊,當人失去jor 呢個城牆城門,帶黎關係,帶黎生命裡面ge 破口,如果教會失去jor 呢一個城牆城門ge 時候,冇jor 神ge 同在臨在保護ge 時候,教會就會分黨分派啦.
教會就會出現分裂ge 情況,因為教會都係罪人ge 地方,罪人亦都係帶住罪性,就我頭先講,如果人冇jor 城牆城門,帶黎生命關係裡面ge 破口,罅惡者同樣可以偷攝jor 人生命裡面最軟弱ge 地方,彼此攻擊.
當教會今年開始新ge 季度黎啦,新ge 一頁展開啦,我同你係生命ge 裡面都要堅持呢八個字,呢八個咩字呀,同你隔離gwo 陣講一次先.
渴慕令火點燃生命,我地大力D 有力ge 宣告話,呢八個字準備起,渴慕令火點燃生命,弟兄姊妹,呢八個字唔係教會ge 口號,呢八個字唔係教會ge slogan,呢八個字係我地屬靈生命ge 表達.
呢八個字係我地作信徒生命裡面呈獻比主ge 禮物,就係乜ye 話,渴慕令火點燃生命,當我地一宣告ge 時候,上帝佢一定會好喜悅我地呢個生命ge 擺上.
我地上帝一定好喜悅我地人人都渴慕呢個靈火,你有冇諗過,當你試過,大熱天時ge 時候,當你行山呀,當你跑步ge 時候,嘩好頸渴呀,好口乾ge 時候,發現已經缺水呀,快要暈低ge 時候,你gwo 刻最大ge 渴慕係乜ye 呀.
你gwo 刻好想有D 冰凍ge 水係我前面,然之後你就會點呀,咩話,佢話堆jor 佢,都要開蓋ge ,開蓋跟住之後就點呀,咕吱咕吱好大力咁樣快D 飲快D 飲,我好想將你心裡面ge 賞望可以得到手,呢個就係渴望呀嘛.
渴慕聖靈ge 火做乜ye 呀,頭先詩歌已經有唱,練正我地,消去我地生命ge 雜質,我地渴慕靈火做乜ye 呀,因為每一日靈修ge 時候,神ge 說話好似光呀,當我地一路讀神ge 說話,原來光就會點呀.
進入到我生命裡頭呀,所有黑暗ge ye 就要離開啦,當我地每日靈修ge 時候,我渴慕靈火主呀,今日你downloadD 乜ye 俾我呀,我好想親近你,好想經歷到你,原來呢個光成為我地生命ge 指引.
每日讀,每日讀,認真讀喎,唔係褻讀呀,每日認真讀ge 時候,靈裡面就會更新,改變我地gwo D ge 壞習慣,當我地渴慕靈火ge 時候,每日祈禱呀,祈禱呢,廿四成七都可以祈禱呀,唔係坐係度食飯先祈禱呀,廿四成七隨時隨地都可以祈禱,你有冇點名,為你緊張ge 人,未信主ge 人,你想起咪一個人,即刻就為佢祈禱呢.
當你渴慕靈火ge 時候,你會快跑就去傳福音,因為你好想人黎到去信主,你唔想人呢,好似呢,如釀酒迷咁,個人偏行幾路,你有生命ge 自由,但係出面好多人都冇呀,你會好緊張佢地ge 生命.
當你渴慕靈火ge 時候,你更加唔會諗到教會有任何ge 缺陷,你就每個月立定心智做乜ye 呀,呢個家要有量,十分一ge 奉獻呀,當你渴慕靈火ge 時候,你每個禮拜都會好珍惜,返黎小組ge 時間,好想見到你,好想見到佢,因為主裡面一家人可以彼此守望,建立gwo 份真正同行ge 關係.
當你渴慕靈火ge 時候,生命藏住神ge 說話,神就係光,光進到生命裡面,黑暗行開啦,那惡者冇佢ge 位置呀,呢度冇佢ge 地盤呀.
當我地話,今年唐慶ge 時候,廖牧師上台黎到宣告,佢俾我地一份ge 鼓勵,要lor 住先知那丹ge 信心,拔事巴一樣ge 忠心堅持,去喂住上帝ge 說話同埋應許,去鼓勵我地,我地要係迦納地啦.
個問題係,我地有冇憑信心去領受神說話ge 大能,zar 住神ge 應許,我地實踐呀,係經歷上帝gwo 份ge 豐盛呢.
而我會問,如果我地連神ge 應許,真理同在,係乜都唔知ge 話,又或者我地根本冇每一日努力去讀經,主呀,你今日有咩說話同我地講,咁我點樣去喂住神ge 應許呢.
其實,我地要以真理為基礎,我地要以真理為基礎,重建返我地ge 生命,我地要以真理為基礎,重建返整個ge 教會.
所以,弟兄姊妹,你有冇發現,一月一號開始,我相信每一個牧者,企係個台上面,都同你講呢句說話,你靈jor 修未呀,問o下你隔離gwo 個.
你靈jor 修未呀,跟住再加多句,每次裝備一出ge 時候,頭先第二季裝備出jor 啦,又同你鼓勵你,一人做咩呀,一人一裝備呀,最少同你隔離gwo 人講一次先啦.
肯定少少講一次呀,一人一裝備呀,加多句,咁你報jor 未呀,咁姐係報jor 未姐,哎呀,我聽到未囉.
弟兄姊妹呀,我又成日聽到有一句說話架,都係你地教我架,我入黎ge 時候就已經聽到啦,你地成日話,有感動就咩ye 呀.
但係我往往呢,成日睇到呀,聽就聽到啦,有感動就行動啦,但係咁咁下呢,個人可以失蹤jor 喎,感動完呀,跟住之後呢,咦,個人呢,又話服侍ge ,去jor 邊呀.
點解呀,因為有感動,但你背後都冇神ge 同在去印證呢個感動究竟係咪真ge ,當你有感動ge 時候,其實你都有需要真理ge 印證架.
你都要去問個神,係咪你呀,係咪我行出去呀,以至到呢個回應先見過呀,唔係咁容易黎到去放棄,用真理作為基石,重建自己生命ge 城牆.
重建教會ge 城牆,城門係要上鎖架,城門係有lock架,lock住D 乜ye 呀,lock住一D 唔合神心意ge ye ,唔讓神gwo D 唔鍾意ge ye 走入去我地教會裡面.
走入去我地ge 生命裡面,呢個咪就係冷惡者,偷竊,殺害,毀壞囉,城牆,城門除jor 保護之外,仲有另外一個好重要ge 意思,就係一條界線.
耶路撒冷係一個神同在ge 地方,又可以叫做大衛之城,又可以叫做錫安,所羅門王係耶路撒冷黎到去建殿,成個猶太人ge 宗教同埋政治ge 中心,猶太人直情睇呢,成個耶路撒冷係世界ge 中心黎ge .
因為以色列書五章五節gwo 度講到,主耶和華如此說,這就是耶路撒冷呀,我曾將祂安置在列國,列邦做乜ye 呀,都係祂ge 四圍,耶穌同樣係耶路撒冷gwo 度釘十字架,賣妝,復活基督教ge 教會.
亦都係係耶路撒冷gwo 度開始黎到成立,呢個福音係係耶路撒冷gwo 度一度傳,傳到地極,耶路撒冷更加係將來新天新地ge 預表,將來永恆ge 國度,所以又叫做新耶路撒冷.
對於呢一班被擄ge 百姓,佢地返屋企呀,係耶路撒冷呀,係城牆,係城門ge 保護呀,我帶住神ge 同在呀,有子女,主女子女名份ge 身份呀,我可以係gwo 度安然居住呀,我可以一齊黎到去敬拜讚美,我可以每一日聽到以色列黎到去讀經,牢記神ge 說話,我唔再被擄啦.
原來呢份ge 保護係帶黎戒善,有城牆有城門同無城牆無城門呢係好大ge 分別呀.
將屬神不屬神ge 分開,將屬天屬地ge 黎到去分開,將神鍾意ge ye 同唔喜悅ge ye 厭惡ge ye 黎到去分開,將屬靈同屬肉體ge ye 黎到去分開,乜ye 係神要ge ye ,乜ye 係神不要ge ye .
你有冇問過神點解筲箕灣浸信會要係西灣河街175號仲有培理學校係我地當中,你有冇問過自己主呀,你創造我係呢個地上,係呢個世界,係呢個家庭,究竟主呀,你俾我ge 心意係D 乜ye 黎呀.
當我預備ge 時候,我發現原來我地缺少jor 一個看門口ge 人,我地缺少jor 一個看門口ge 人,看門口ge 人呢,企係個城門gwo 度架,係點呀,日以繼夜守住條界線,守住那惡者你行開,你不得進來,唔好再係人生命裡面偷竊殺害攻擊.
有一段經文,一個大祭司,耶和耶大,當佢設立呢一個祭司利美團隊ge 時候,佢話,有一班人要做咩呀,守門口,有一班人特意要黎我守門口,佢話,閒雜人ge 不能進來,外邦人不能進來,不聖潔ge 不能進來,唔係神要ge ,請你唔好進來呀.
原來屬神ge 人係要守住條界線,原來屬神ge 人,我同你都係守住門口gwo 個人黎架,我同你都有個城牆,神ge 靈臨在ge 時候,要保護我地ge 生命,保護教會ge 時候,我地原來都係守門口ge 人黎架.
講個故仔比大家聽,係一個小村莊ge 裡面住jor 兩個家庭,一個家庭係爸爸ge 屋企,一個家庭係女兒ge 屋企,一個左邊,一個右邊,佢地屋企之間好簡單ge je ,就係圍上呢個圍欄,就大家知道啦,我有界線ge ,呢邊係爸爸ge ,呢邊係女兒ge .
爸爸呢一個好守規矩ge 人,佢好尊重界線,所以呢,佢入面個花園呀,所有修成呢,都種得非常之好呀,但係呢個女就唔係啦,佢好鍾意唔守規矩,成日越界,攝去爸爸gwo 個圍欄側邊呢,就偷jor 爸爸裡面gwo D ge 修成.
有一日呢,爸爸呢同對呢個女兒ge 行為呢,成日點解我番企呀唔見jor 呀,點解我種ge ye 突然間會唔見jor ge 呢,跟住對呢個女兒ge 行為覺得好不滿呀,點解你會做D 咁ge ye ,但係呢個爸爸ge 心諗,我又唔想同個女有任何拗撬,有任何衝突呀,於是呢個爸爸好聰明呀,佢竟然諗jor 個計仔喎,佢竟然係個圍欄旁邊呢,種jor 好多棵大樹呀.
種jor 好多棵大樹ge 時候,個女兒一望ge 時候,醒定呀,呢條界線黎架啦,唔好過黎呀,跟住之後呢,地不知,你知啦,D 樹會點架,陽光普照,越山越高,跟住好多gwo D ge ,嘩,真係完全好豐盛啦,勁多樹葉啦,跟住呢,遮曬D 陽光,於是爸爸ge 果園呢,就變得越黎越暗.

$^{81}$於是爸爸gwo D 收成呢,竟然就係因為咁失收啦,於是佢為jor 保護條界線,就失去jor 自己ge 利益,於是爸爸忍唔住同個女,都係要我地坐低,講一講先.
原來當彼此尊重返界線,守住界線呢,原來先可以減少衝突同埋損失,當彼此尊重界線ge 時候,守住界線,原來先可以保護大家,減少個衝突同埋損失.
建立界線係因為愛,建立界線ge 基礎係因為愛,係因為愛帶黎彼此ge 尊重保護,相反,如果失去jor 愛ge 時候,人就會越界,人就會冒犯人.
弟兄姊妹呀,你知唔知點解會冒犯呀?係因為人裡面ge 自私呀,還有人裡面ge 貪婪先會越界,界線唔係冰冷,界線唔係同人割蓆,而係一份屬靈ge 智慧去彼此保護,去尊重互信,最終ge 目的只得一個,就係建立一個健康ge 基督ge 身體.
所以我地頭先讀ge 今句啟示錄廿一章十七字,我地一齊大聲ge 黎對讀出,準備起!.
凡不潔淨的,和可恨可憎,與虛荒之士的人,都不得進那城,只有明志,寫在羔羊生命冊上的,才得進去.上帝特意係耶路撒冷,百姓返屋企,返屋企啦,回歸啦,城牆城門,有個好大ge 保護呢班ge 以色列民..
當彼此時守界線ge 時候,呢班以色列民處在其中,係上帝揀選ge 一群,係聖結ge 一群,係祭司ge 一群,係上帝喜悅ge 一群,而家返屋企啦!.
但係點不知聖經話比我地聽,佢遭大能受凌辱,你會唔會同我都問一同一條問題,我返屋企等於無返屋企,咁我返屋企,你做乜ye 呢?我點解要返屋企呀?而家城牆被拆,城門被火燒,我帶黎ge 後果結果呀,都仍然一樣遭大能受凌辱,我呢個仇敵自出自入ge ?.
失去jor 神ge 家,失去jor 保護,失去jor 呢條ge 界線,當你再睇呢段ge 經文,尼希米記打後ge 經文,呢班以色列民已經無往館呀,已經無jor 條界線ge 底下,去同外邦ge 人通婚,去拜偶像,安息人,係敬拜主,佢地走jor 去做生意呀,所以整個民族都已經失去jor 成個信仰裡面ge 見證..
尼希米好傷心,好難過,悲哀,禁食,禱告四個月,對於尼希米黎講,被擄歸回其實係一個膚字話,被擄歸回明明係神ge 應許呀,被擄歸回,神係為gwo 班以色列民好架嘛,帶佢地返屋企,要保護佢地,但點解返到屋企,而家繼續被擄呀?.
尼希米覺得好羞恥呀,點解會發生呢D 事,我心裡面好痛,當我地話,我地要一齊經歷渴慕靈火,弟兄姊妹,你有冇諗過邊一個係最唔想我地渴慕靈火呀?邊個呀?邊個最唔想我地渴慕靈火?.
那惡者呀,係一個好短ge 時間裡面,我看見那惡者,佢點樣偷竊jor 人生命裡面最軟弱ge 地方,就黎到去彼此攻擊,原來當一個人怒氣蓋頂ge 時候呢,從來冇見過呢一個ge 樣子呀,當情緒失控ge 時候,原來噴出黎gwo D ge 言語係可以好恐怖架,.
當我係度諗,究竟gwo 個人裡面藏jor D 乜ye ye ,係佢心底呢,只要一有爆炸點ge 時候,嘩,佢裡面入面ge 污穢,全部透過個語言噴曬出黎,又有D 時候睇見,講者無心,下一句咩呀,聽者有意呀,講者無心,聽者有意,於是就誤會jor 對方啦,.
呢個小小ge 破口就出現jor ,由彼此黎到去受傷,我又看見你gwo 幅圖畫,原來我地生命ge 裡頭,我地成日都做jor 小導演,小導演呢,最叻ge 一件事,就係編劇,我地自編,自導,自演,同呢個事實係不符,但你又唔會fact check架喎,就以為係真,.
你知唔知道呢個係思想陷阱黎架,係叫做妄下判斷呀,我地又中jor 罵惡者ge 詭計,自己傷害自己呀,甚至乎你以為人地都係咁呀,當一個人好介意別人點樣去睇你ge 時候,.
不管人ge 樣子,甚或是人同你話話gwo 個聲音,你都會零射敏感,你會零射在意架,你就以為別人有惡意黎到去對待你,有時候又看見,因為一D 個人ge 意思,或者我好想爭取一D ye ge 時候,.
就可以利用jor 自己ge 權力,去越過jor 所有ge 規則,城牆城門本身係比我地帶黎保護,係有界線,gwo 個目的係為jor 我地成隊主裡面ge 家人,一齊行出黎作見證,.
而呢D ge 事情ge 發生,用屬靈ge 眼睛看見,那惡者真係好古惑,那惡者同樣看見連人,係連我地自己都唔認識自己生命最軟弱ge 地方,就係我地生命ge 破口就lor 黎黎到去彼此攻擊,.
但係就係因為我地連自己ge 破口都唔知,我地仍然以為自己係arm ,其實係我地盲jor ,係我地根本睇唔清楚自己ge 本相,總係以為別人有問題,其實係自己都看不清自己生命ge 軟弱破口已經比那惡者lor jor 黎彼此攻擊,互相破壞,.
就好似呢個好靚ge 紅蘋果一樣,見到ge 時候BlingBling眼,你見到ge 時候食完好想快D 咬佢一啖啦,當你一咬一啖正喎,好甜呀,好爽呀,再咬多啖啦,原來你發現,發現乜ye ,中間有蟲,所以金玉皮外可以乜ye 架,敗水其中,內裡面生蟲呀,.
但係呢,我問主,我覺得有一個好奇妙ge 地方,神呀,點解咁短ge 時間,而家幾多月呀,三月咋,一月,唔好話一月啦,可能上年年尾ge 時候,我地已經講,主呀,我地要追求,開悟靈火,點燃生命,點解一個咁短ge 時間,容讓呢D ge 事情,砰砰砰咁樣,突然間爆炸咁樣出現ge 呢,.
我一度諗,如果唔浮面呀,又點樣見到人生命ge 破口呢,那惡者咪只會越呢越埋,但當浮面ge 時候,神工作緊ge 時候,人看見生命裡面ge 破口,弟兄姊妹呀,你知唔知呢個破口叫咩名,呢個破口叫做玻璃心,.
我同你都一樣,都有呢個玻璃心,係我地最脆弱,最敏感ge 地方,而浮面要處理ge 時候,就係主裡面呢,一切都要見光死呀,唔容讓那惡者再係我地玻璃心ge 裡面呢,偷竊,殺害,毀壞..
弟兄姊妹呀,你明白點解呢班人,明明被擄,而家返屋企,仍然再被擄,係佢地ge 心靈再被擄,其實我地都係一樣,我地係主基督ge 名下,我地係舊恩,我地已經有呢份永生ge 禮物,仲話聖靈已經係我地ge 生命裡面ge 印記,去俾我地封聖,.
應該好自由呀,應該好喜樂架,應該係大雷ge 遮蓋ge 裡面,我地唔使再坐監呀,係籠奴ge 裡面,係罪ge 捆綁ge 裡面,已經得釋放喇喎,點解,真係好想問,點解,我地成日都係唔同ge 原因ge 裡面,甚或是我頭先講過,玻璃心自己捆綁自己呀,自己捆綁自己呀,.
我地ge 心仍然被腦運,當你問o下你自己,你ge 心有冇被捆綁呀,你有冇有D 人ge 樣子好恐嚇你呀,你有冇被一D 聲音黎到佔奪jor 你腦裡面,嘩,你應該係好清潔架,你唔需要被呢D 煩惱黎到影響你,你有冇諗o下,你自己做jor 小導演,做jor 小劇場,你自編自導自演,.
係你自己認為係,唔係呀,同事實不符呀,可能你裡面有好多唔滿意ge ye ,可能你裡面有好多罪ge 核制,金錢物質ge 核制,慾望ge 核制,甚至乎尊嚴大過一切,你問o下你自己,有冇被呢D ye 黎到去擄奪我地生命裡面ge 自由,.
弟兄姊妹呀,你有冇諗過仇敵呢,係度按按腳,係度剝花生,拍爛手掌,正呀,見到呢班人,嘩,仍然被擄,信完耶穌得一樣,仍然被擄ge ,佢係度拍爛手掌,佢點會想見到我地合一復和呀,.
弟兄姊妹呀,當我地話進入新ge 季節,進入一個新ge journey,一個新ge chapterge 時候,渴慕靈火,我地最期盼去到終點,就係成個教會要復興呀,但係復興唔係齋up 架,復興唔係淨係騎呀騎呀騎呀騎,我知道神有工作,但係復興ge 裡面,最重要一個位係我地自己呀,.
如果我地生命唔改變ge 時候,點樣復興呢?如果我地生命唔改變唔更新,我地呀,同自己都不復和呀,當你諗o下,當我地自己同自己都不復和ge 時候,連自己玻璃心都唔知呀,連自己被那惡者黎度偷笑殺害毀壞都唔知,以為自己arm 曬啦,咁我地生命點復興呀?.
當我地同人ge 關係不復和ge 時候,小組又點復興呀?當我地每一個信徒係主裡面呀,如果我地唔守住教會裡面ge 城牆,遵守教會裡面所set出黎gwo D 章則規矩ge 時候,越界,唔係正確ge 軌道,神ge 軌道上行,教會點復興呀?.
當我地諗,去到總結ge 時候,我想同你分享一個,我ge 朋友同我分享一個見證,佢有一日呢,係佢facebook messengergwo 度,佢收到一個form7ge 同學,佢message佢,即係講緊幾廿年前ge 事,起碼廿幾年前ge 事,.
呢一個同學係特登係facebookgwo 度搵返我呢個朋友出黎,點解呀?因為佢有ye 要講,原來佢呢個朋友同我呢個朋友去say sorry,佢文字裡面表達jor ,係一個form7ge field trip裡面,.
係我ge 朋友沖涼ge 時候,佢lor 晒佢D 衫,佢表達緊對自己呢個行為,覺得好遺憾,亦都係表達jor 一個遲來ge 道歉,當我問返我朋友,你仲記唔記得呢件事呀?.
佢話我已經唔記得啦,不過佢記得另外一件事,就係當佢沖涼ge 時候,比人影低jor 佢ge 相,佢當刻覺得好欺凌,被人欺凌,被冒犯,更加有gwo 份ge 創傷係度,但我又係度諗,.
如果唔係神光照呢位同學ge 話,其實lor jor 件衫je ,仲係少不更易ge 時候呀,所以gwo 時一件好笑ge 事,點會帶你自責,遺憾,要係facebookgwo 度等jor 幾十年,搵返呢一個ge 朋友出黎,要同佢say sorry呢?.
講真呢廿幾年,話短唔短,話長唔長ge 日子,但我好相信,當神光照佢gwo 刻開始,佢ge 心靈已經唔好受,係呢咁長ge 日子裡面,佢點樣過,直到佢搵到我ge 朋友,同佢表達歉意,呢句誠意表達ge sorry,.
佢好清楚,當佢講jor 出黎,就係見光死,佢內心ge 捆綁,係呢廿幾年裡面,佢完完全全真係捆綁得釋放,呢個同學,係當時ge 日子,係佢生命裡面,係生命無長,生命冇jor 城牆城門ge 時候,.
佢冇好好保護大家彼此gwo 份ge 友誼,當時ge 唔成熟,幼稚,自私,佢越過界線,佢冇返人,失去jor gwo 個見證,尼希米好難過,好心痛,佢話我要重建呢個耶路撒冷ge 城牆城門,佢要守住屬天ge 界線,因為耶路撒冷係神同在ge 地方,.
呢一班ge 百姓係聖結架被揀選,上帝係我地唯一ge 忠心,尼希米話我要重建城牆,佢要保護呢一班百姓到底,你仲記得嗎?上年五十五週年,陳康牧師已經係台上面宣告jor ,我地一群乜ye 人黎,係耶和華ge 軍隊,.

$^{121}$我地呢一班弟兄姊妹係乜ye 黎,係耶和華ge 軍隊,同你隔離gwo 人講一次,我地邊個黎,係耶和華ge 軍隊,我地大力D ,重力D ,再講多次,我地係耶和華ge 軍隊,弟兄姊妹,你知唔知耶和華ge 軍隊已經係迦南地,你知唔知你而家係邊?.
你而家係一個豐盛ge 地方,係一個呢迦南地ge 裡面,呢班ge 百姓做咩?成為耶和華大能ge 軍隊,係聖潔同在,係神ge 遮蓋,保護ge 底下,去到今年,五十六週年,仲記得廖牧師講乜ye ?.
去台上面宣告,我地要記得,係迦南地ge 裡面,我地要打仗架,我地要屬靈真戰,但係真戰係必,乜ye 呀?係必得勝呀,真戰係必得勝,但係點樣得勝呀?去鼓勵我地,要重現呢個城牆,每日黎到去讀聖經呀,以神ge 說話,成為我地ge 英雄,神比jor 我地ge 英雄,.
而呢一個係豐盛ge 英雄,同你隔離gwo 人講一次,係咩黎ge 話?豐盛ge 英雄,但係我地要lor 住呀,我地要用信心去持守住呢個恩典,往前直上,以至到我地ge 生命要結出果子..
弟兄姊妹呀,過往ge 日子,教會ge 經歷,教會ge 故事過去啦,我地已經係泥沼gwo 度上返黎啦,已經要往前直上,我地而家係度做緊乜ye ?再講多次,我地要成為耶和華ge 軍隊,一齊講一次,要成為耶和華ge 軍隊,成棟ge 筲箕灣浸信會比jor 你啦,係度啦..
我地今日要做咩?重建呀,重建呢個敬拜ge 祭壇,重建呢一個祈禱ge 祭壇,重建呢一個我同你都有份ge 宣教ge 使命呀,再講多一次,我地係度做乜ye 呀?重建呀,重建呢一個敬拜ge 祭壇,重建呢一個祈禱ge 祭壇,仲有一個重建我地生命裡面要彰顯出黎ge ,每人都有份ge 傳福音呀,一個宣教ge 使命呀..
但係我地而家最重要ge 一樣ye ,係重建返我地ge 生命,係重建返我地ge 生命,唔好再俾那惡者再偷竊jor 我地生命裡面ge 自由,自己攻擊自己,自己攻擊其他人,冇呀,上帝已經係度啦,堵塞我地一切ge 缺口..
所以你明白,生命無牆,我地會點樣呀?好快收檔,同你隔離gwo 陣講一次先啦,生命無牆又會點呀?我地一齊起立好嗎?當你諗到你生命ge 破口係乜ye ?你ge 眼睛,你ge 耳朵,你ge 嘴巴,你ge 心,你ge 腦袋,你知道啦,惡者係度偷竊緊我地生命裡面最自由ge 地方,可能我地心裡面有好多罪惡係裡面,.
被捆綁,心靈深處好多地方係我地唔想俾人治,但係弟兄姊妹呀,你知道嗎?今日上帝佢要將釋放ge 能力要充滿係我地裡面,今日上帝要臨在呢個地方,要釋放宗教會呀,佢唔容約呢個地方係聖潔ge 地方,被冷惡者再次黎對偷竊,殺害毀壞,今日要重建自己,今日要放低過往一切ge 重擔,.
你誠誠實實咁樣去面對自己ge 真相,唯有你自己照一照鏡,你心靈裡面照一照鏡,你就確知道你裡面係幾唔掂架啦,當你照一照鏡,倒空你自己,主呀,你光照我,主呀,你光照我,主呀,你話俾我聽,我生命裡面最軟弱ge 地方係乜ye ?.
我想做返一個自由人,我想做返一個豐盛ge 人,我唔想再破碎呀,我想經歷抓住原來上帝ge 英渠係咁美好架,昔日ge 故事已經過去,無論你生命裡面有幾多裂縫一樣,你有冇諗過,光盡到你裡頭ge 時候,所有ge 裂痕都已經被上帝愛ge 同在依照,已經縫補jor 你生命裡面所有ge 缺陷,黎啦,我地一齊去面對..
一齊去主力聯絡禱告.
\newpage



\section{路加福音 17:20-37}
\label{sec:031778}
\textbf{隨時預備,警醒迎接神的國}~~~~陳傳華牧師博士
\newline
\newline
連結: \href{http://www.church.com.hk/acms/content.asp?site=cdc\&op=show\&type=product\&code=031778\&layout=sermon}{\texttt{031778}} ~~~~ 語音日期: 2026-03-08
\newline
\newline
\hyperref[sec:031777]{\small{< < < PREV SERMON < < <}}
~
\hyperref[sec:index_chronic]{\small{[返順時目]}}
~
\hyperref[sec:031779]{\small{> > > NEXT SERMON > > >}}
\newline
\newline
路加福音 17:20-37
\newline
\begin{longtable}{cl}
\hline
\hline
章節 & 經文 (和合本修訂版)\\
\hline
17:20 & \begin{tabularx}{0.7\textwidth}{X} 法利賽人問：「神的國幾時來到？」耶穌回答：「神的國來到，不是眼睛看得見的。 \end{tabularx} \\ \\ \relax
17:21 & \begin{tabularx}{0.7\textwidth}{X} 人也不能說：『看哪，在這裡！』或說：『在那裡！』因為神的國就在你們心裡。」 \end{tabularx} \\ \\ \relax
17:22 & \begin{tabularx}{0.7\textwidth}{X} 他又對門徒說：「那些日子將到，你們渴望能看見人子的一個日子，卻看不見。 \end{tabularx} \\ \\ \relax
17:23 & \begin{tabularx}{0.7\textwidth}{X} 有人要對你們說：『看哪，在那裡！』或說：『看哪，在這裡！』你們不要出去，也不要追隨他們。 \end{tabularx} \\ \\ \relax
17:24 & \begin{tabularx}{0.7\textwidth}{X} 好像閃電從天這邊一閃直照到天那邊，人子在他的日子也要這樣。 \end{tabularx} \\ \\ \relax
17:25 & \begin{tabularx}{0.7\textwidth}{X} 可是他必須先受許多苦，又被這世代所棄絕。 \end{tabularx} \\ \\ \relax
17:26 & \begin{tabularx}{0.7\textwidth}{X} 挪亞的日子怎樣，人子的日子也要怎樣。 \end{tabularx} \\ \\ \relax
17:27 & \begin{tabularx}{0.7\textwidth}{X} 那時，人又吃又喝，又娶又嫁，直到挪亞進方舟的那日，洪水就來，把他們全都滅了。 \end{tabularx} \\ \\ \relax
17:28 & \begin{tabularx}{0.7\textwidth}{X} 同樣，就像在羅得的日子，人又吃又喝，又買又賣，又耕種又建造， \end{tabularx} \\ \\ \relax
17:29 & \begin{tabularx}{0.7\textwidth}{X} 到羅得離開所多瑪的那日，有火與硫磺從天上降下來，把他們全都滅了。 \end{tabularx} \\ \\ \relax
17:30 & \begin{tabularx}{0.7\textwidth}{X} 人子顯現的日子也要這樣。 \end{tabularx} \\ \\ \relax
17:31 & \begin{tabularx}{0.7\textwidth}{X} 在那日，人在屋頂上，東西在屋裡，不要下來拿；人在田裡，也不要回家。 \end{tabularx} \\ \\ \relax
17:32 & \begin{tabularx}{0.7\textwidth}{X} 你們想想羅得的妻子吧！ \end{tabularx} \\ \\ \relax
17:33 & \begin{tabularx}{0.7\textwidth}{X} 凡想保全性命的，要喪失性命；凡喪失性命的，要保存性命。 \end{tabularx} \\ \\ \relax
17:34 & \begin{tabularx}{0.7\textwidth}{X} 我告訴你們，在那一夜，兩個人在一張床上，一個被接去，一個被撇下。 \end{tabularx} \\ \\ \relax
17:35 & \begin{tabularx}{0.7\textwidth}{X} 兩個女人一同推磨，一個被接去，一個被撇下。」 \end{tabularx} \\ \\ \relax
17:36 & \begin{tabularx}{0.7\textwidth}{X} 兩個人在田裡，一個被接去，一個被撇下。 \end{tabularx} \\ \\ \relax
17:37 & \begin{tabularx}{0.7\textwidth}{X} 門徒回答他說：「主啊，在哪裡呢？」耶穌對他們說：「屍首在哪裡，鷹也會聚在哪裡。」 \end{tabularx} \\ \\
[1ex]
\hline
\hline
\end{longtable}
$^{1}$好 今天我們就進入到這個題目.
這個題目就是在末世有關的.
關於什麼呢?.
我們立刻就進入去看.
那裡隨時準備.
警醒迎接神的國.
我不知道大家在一生中有些什麼.
常常多作準備.
是一件什麼事.
為什麼要持續.
我常常都聽到很多一生的事業的東西.
例如女士就說減肥是一生的事業.
是一生的.
有些人就說做運動也是一生的事業.
有些人就說扮靚也是一生的事業.
我們很多是一生的事業.
當有時我聽到女士說減肥是一生的事業的時候.
同時我又聽到有些人說.
減肥也是我的終身事業.
但是吃東西是我一生的興趣.
真是很困難.
有另一面.
剛剛我有個組長去了東京跑馬拉松40多公里.
做運動是一生的事業是很明白的.
但是有些人說做運動就好像拍拖一樣.
頭幾次很熱情.
之後就開始找藉口.
希望你拍拖不是經歷這樣的事.
扮靚是一生的事業有人就這樣說.
他說護膚是一場同時間的賽跑.
雖然我跑得慢.
但是我裝備多.
裝備多就是搞不定.
當然這個也是一個笑話.
我們很多東西真是需要很有恆心去堅持做.
有些做不到.
都明白的又不是這麼大件事.
不過有些東西我們的確要堅持去走去做.
因為那些東西不是我們喜歡或者不喜歡.
那些東西不是這樣.

$^{41}$就會產生很大的反效果.
例如消防員.
問你什麼時候準備好出去滅火救人.
當然不是當你看到收到電話你才準備.
其實是天天的準備.
時刻的效命.
因為危險隨時發生.
我看過報告.
看來也是相當有興趣.
原來一個消防員當他輪班的時候.
他不是在等著做什麼.
玩一下牌.
然後吃一下東西.
吃一下榴槤.
吃一下粽子.
如果是我就做得到.
我看完之後我發現我是肯定做不到的.
第一.
他早上需要好好安排.
接更檢查工具.
然後他就要消防訓練.
多層樓宇.
拯救就是走來走去.
然後他就要在裝備維修保養.
外勤巡查.
然後他需要出外檢查街井.
然後樓宇通道.
不斷在準備.
不僅如此.
還要有體能訓練.
會打排球和跑步.
好了.
我想問問大家知不知道為什麼.
歡迎在加拿大回來的朋友.
是你們.
昨天Carmen也提醒你們.
因為與我無關.
大家上網有父母在聽.
有人問為什麼消防員這麼喜歡打排球呢?.
大家有沒有想過?.

$^{81}$有沒有想過?.
我真的不明白.
有一個他說他很明白.
解釋了.
不知道是不是真的.
他說消防員打排球是第一個運動量高.
跑步運動量高.
第二.
受傷機會低.
因為消防員如果打籃球.
你知道我是男人.
我不知道你是男人.
我打籃球做事是不可以輸的.
我除了死給你看.
被你擊倒.
可能他們會抓住你.
踢你.
踢球.
所以這是最安全的分兩隊.
他們真的這麼說.
但是要運動.
要保持.
然後要有課堂的時間學習.
然後繼續維修.
車輛保養.
然後繼續等候.
後命在裡面等交接.
那個不是等那一刻來才準備.
是天天準備.
常常準備.
當然今天我們不僅是說這件事.
這件事當然重要.
但是我們說這件事更加重要.
就是末世的時候我們怎麼準備.
是不是末世的時刻.
才準備之前要做什麼.
今天這段經文.
就是一些宗教領袖.
法利賽人.
當時跟耶穌的對話.

$^{121}$他問耶穌天國是怎樣.
神國降臨是怎樣.
是不是政治性.
是不是領域性.
但是耶穌顛覆他們的想像.
將焦點轉向上帝.
不是他日.
是今天我們要準備什麼.
而今天上帝又是怎樣掌權.
今天我看到大家.
有些人又再來.
有些人沒有小組.
很開心.
所以我今天多送兩點給你.
送五點好不好.
開心嗎?.
好像不是很開心.
好不好.
說到五點好不好.
真的笑得很開心.
今天跟隔壁說一聲好不好.
今天有五點學習.
五點學習.
學習一.
一開始的對話.
一二十二節我立刻讀.
一二三.
法利賽人問.
神的國何時來到.
耶穌和教主.
神的國來好多時.
看在你們心裡.
看在哪裡.
因為神的國就在你們心裡.
法利賽人問一個問題.
神的國很有型有底.
他在問何時降臨.
人也不得說.
耶穌的回答.
是不是這樣的呢?.

$^{161}$我們很明白.
在當日在歷史.
或者不單止是他們的歷史.
在以往世界歷史.
無論是波斯.
無論是秦始皇.
無論是巴比倫.
或者是亞述.
希臘.
羅馬帝國等等.
都是有型有底地域性的.
外在化的.
地緣政治.
不過耶穌顛覆他的思想.
他說當日神國.
神國有兩種.
一個是末世.
神國會降臨.
第二個.
他說神國不只是末世才來.
如果我們記得耶穌出道第一句說.
天國近了.
你們應當悔改.
天國近.
天國已經降臨.
即是有兩種.
第一.
祂降臨在我們內心.
第二.
祂會一致過末世的末端.
會噴一聲出現.
好了.
所以第一件事我們正在學習的.
天國已經降臨在你內心.
是超越外在的尋求.
所以我們不是說地緣化.
政治化等等.
我們很明白的.
因為就算這幾年我們見到.
在俄羅斯打烏克蘭.

$^{201}$即是很地緣政治.
特朗普打委內瑞拉.
打伊朗.
我們也不用說也知道.
我們會想起很多軍事.
F35.
不同的裝備.
航空母艦.
突然間進入國家去打等等.
不過我們不是說這些.
天國不是這種形態.
耶穌顛覆我們的思想.
祂進到人的內心.
是我們要願意開放.
讓上帝進到我們內心.
並且在我們內心更新改變.
用這些例子你就更加容易拿捏到.
什麼是神國降臨.
前幾天我和廖牧師.
我們早上吃早餐.
有些mentoring section.
我們去了一間茶餐廳.
附近而已.
很多間茶餐廳.
我覺得香港茶餐廳真的很寶貴.
因為香港可以沒有任何一件事.
除了茶餐廳.
否則香港五分之一人應該沒有東西吃.
很重要的.
我們坐下來吃東西.
那間茶餐廳看到我們教會.
看到的位置.
你嘗試接觸那麼多間.
有哪一間看到我們教會.
什麼都有關係的.
我和廖牧師聊天.
突然間有個侍應走來.
他說你是不是就是那間教會.
指著我們.
樓上說是啊是啊.

$^{241}$幸好我們說了好話.
大家在外面說教會.
說神吧.
簡短一點.
他說他也是基督徒.
幾年前信了主.
聽到真是好.
他很興奮.
可能他知道我們在說話.
我沒有問他.
接著他說.
信了主幾年前.
在這裡經歷.
我就問他.
你為什麼信主.
誰帶你回教會的.
這個問題正不正常.
正常.
你問一問你旁邊那個.
誰帶你回教會的.
當然是親人或是朋友或是同事等等.
很正常.
接著他這樣說.
沒有啊.
天父帶我的.
大家很神聖.
我問一個地域性的問題.
形有體的問題.
我很有趣很想問.
什麼意思.
大家很神聖.
他說我是香港人.
我從小到大都有人跟我傳福音.
不過我不相信.
不過幾年前.
我生命發生了狀況.
所以時機到了.
他真的說時機到了.
大家說英文.
時機到了.

$^{281}$他知道他要信主.
於是他自己就去教會.
於是他信了主.
我想說什麼呢?.
其實神的國常常都想降臨.
重點是我們是否讓神國降臨.
明白嗎?.
明白吧?.
所以他一直說.
我覺得很好.
當然他有說他的處境.
我不在這裡多提了.
這種的太明奇妙.
神國就在內心.
而重點是我們拒絕.
還是邀請上帝的國繼續在這裡成長發酵.
掌管我們並且帶我們生生育月.
又另一個例子更加明白.
這個初信.
有一個信了很久了.
我們大家都知道.
我們教會門女門徒.
生命成長轉化旅程.
門女.
是門徒在這裡成長.
其中我們現在到了不同的教會.
都幫他的牧者開始有門女.
其中一個教會比我們更大.
差不多兩千人的教會.
有一個執事公開分享了.
我去分享.
很有能力的執事.
做事很厲害.
美國做完事.
香港做事.
家裡管理得非常好.
大家可能覺得.
一定是一個男士.
其實是一個女士.
現在女士很厲害.

$^{321}$家裡她說了算.
我認識她老公.
家裡女士說了算.
今天在香港一點都不意外.
她說了.
我分享的見證我自己也很入心.
她說她也知道.
我也知道她有能力.
她的頭腦也很清醒.
她這樣說.
當家裡有爭執.
有看法不同.
她的腦子轉得很快.
她理據很足.
她說真的贏了.
她說了其中一次.
的時候我自己覺得.
神國真的在裡面.
正在生長.
她說有一次也是.
她的女婿是西人.
來香港探望她.
和孫子.
上教會那天.
怎知道他們很慢.
孫子還沒換衣服.
再這樣就會遲到.
於是她就加入.
幫孫子搞定了.
搞定我們.
你知道西人和中國人.
西人.
我的兒子我的女兒.
我教不是你搞的.
她就說你為什麼要搞我.
我要搞這件事.
然後她就說.
我照顧小孩.
你也不知道在哪裡.
當然是用英文說.

$^{361}$差不多的意思.
然後她女婿立刻退縮.
然後回教會.
回到教會她又是很有能力的人.
看到教會有些事情做得不好.
她就說我應該這樣這樣.
但是在崇拜裡面聖靈提許.
聖靈提許我給智慧給能力.
但是你缺少一件事.
就是謙卑溫柔和忍耐.
我們知道的.
能力是屬靈恩賜上帝給我的一點也不難.
但是屬靈的果子是要自己結出的.
要結出來並且向上帝的聖靈改變.
於是她又.
完了崇拜就跟她的女婿.
又跟教會裡面說.
用詞不好.
道歉.
她說道歉這個詞在腦子裡面從來沒出現過.
聽完之後真的很感動.
神光在我們內心究竟可不可以見面見到.
不是似乎上帝已經說神國在我們裡面.
重點是我們讓不讓神國.
在我們裡面繼續掌控掌管.
在日常生活中活出神國的價值.
在職場家庭人際公共生活中活出愛.
公義謙卑和誠信.
這是其中一個很重要的神國氣息.
所以第一點我們值得問.
我們會不會每天早上都是投民主主要.
今天我面對不同的會議.
不同的家庭的需要.
或者我作為一個長者照顧家庭.
我可以有什麼要上帝在我的領域掌權.
你是否願意每天都這樣禱告.
讓上帝可以用上帝創造我們的生命力.
並且我們的中間神國.
是否已經在展現那種喜樂平安.
有時我們見到一些基督徒.

$^{401}$好像每天都去葬禮.
完全裡面沒有那種生命喜樂.
真是很奇怪.
聖經說的不是這樣.
什麼是這樣呢?.
第一點.
究竟神國是否你的內心.
我們看看我們身邊的人的心.
你是看不到的.
因為有肉有骨有生.
你看看你的心.
也不是看出來.
是你感受出來.
第一點.
第二點.
繼續下去他們的交談就很有趣.
神國不單是我們的內心.
不過在他日末日的時候.
又很真實會見到祂的彰顯.
我們一起繼續讀下去.
22,25節.
123.
他又對門徒說.
日子將到.
你們巴不得看見人子的一個日子.
卻不得看見.
人將要對你們說.
漢娜在那裡.
漢娜在這裡.
你們不要出去.
也不要跟隨他們.
因為人子在他降臨的日子.
保障閃電從天這邊.
一閃直到天那邊.
此時他必須先受許多的苦.
又被這世代棄絕.
這裡說的是末世中的末世.
我先解釋一下.
聖經有多處地方提到末世.
這是其中一個絕跡.

$^{441}$而末世其中有些警號.
跡象.
是這30年才出現的.
末世何時可以倒數.
耶穌死而復活開始進入末世幾個階段.
而末世裡面有一個末端.
我跑馬拉松跑40公里.
你會衝線.
衝線的末端.
末端會有些警號.
而警號是這幾十年才出現的.
例如文打文 國打國.
當時不可能有世界打到大戰.
現在科技容易了.
第二次世界大戰.
第三次也可能來了.
科技做到.
例如他會說到氣候或者污染.
甚至海裡面的生物.
四分一三分一死亡.
現在的污染.
你走著走著也不知道死了多少.
例如他說到買賣是不需要用錢的.
用密碼的.
現在誰還用錢.
有的很少.
當然拿個電話.
我老婆什麼都在裡面.
DOOT.
Wechat什麼的.
DOOT了.
還有.
現在還有很多晶片進入身體.
買東西很容易.
是這幾十年出現的.
大聖經2000年說了.
有這些跡象.
所以.
如果末日的末端真的很快出現.
又或者福音傳到地極.

$^{481}$以前沒有科技.
現在真的AI.
所有文字十年內.
我們聽到不同的民族幾千個語言.
十年內都會聽到福音.
是這十年發生的.
如果真的的話.
他就說末世原來會有異端出現.
耶穌會來了.
看看這裡.
看看那裡.
然後會有開始迷惑人.
他說你不要出去.
不要跟他.
並且.
如果耶穌再回來.
死而復活.
再末世的末端回來.
一出來是所有人都看到的.
從這邊到天下都看到.
現在什麼時候.
一樣東西出現.
會全球看到.
是這個時候才行.
我不知道是不是.
不過肯定科技行.
YouTube都出了.
小紅書都出了.
以前是完全想不到你說什麼的.
你明白我說什麼嗎.
不過重點是你要小心.
因為很多的異端會出現.
今天的異端多到不得了.
很顯明很普世性的異端.
會看到.
就算我們會中.
或者我們神學院教很多的牧者.
我教博士的學位.
很多牧者都曾經跌入異端.
曾經跌入異端.

$^{521}$真的不少.
例如現在我們也看到.
耶穌說見證人會.
他說耶穌是受造物.
天使長米加納.
很多人相信.
很多人去了.
否認三一神.
摩門教.
後期聖徒教會.
他說耶穌和撒旦是兄弟.
天父的長子.
當然他除了聖經以外.
有些典章摩門經他們會跟進.
很多異端除了聖經以外.
他有些東西要跟進.
要小心.
又或者有些說耶穌已經來了.
剛剛來了.
先知,載臨,運動等等.
你不要覺得很遠.
我們教會曾經完了崇拜留下.
都有東方閃電派傳單.
結果我們牧者就下去.
送客.
不單止這樣.
我有一次講不到.
有一位姐妹約我見面.
一坐下.
她跟我說.
陳牧師.
有些事要告訴你.
什麼事?.
她說耶穌已經回來了.
是嗎?.
然後她說.
還是女人.
我也在找.
你明不明白?.
我也想到.

$^{561}$所以你要小心.
如果聽到我說得怪怪的.
有兩件事要做.
一是馬上去報告總會.
說這個人糟糕.
第二是要走了.
要很小心.
現在我們也知道.
這種形態.
其實他們很多都是好人.
不過他們被迷惑了.
他們很多我認識的.
有些人是君子.
但是我們要小心.
他們信錯了.
不是三一真神.
他不是用神人異性.
耶穌是全神.
也是全然人.
他明白人的痛苦.
並且有些人說.
我是使徒.
我今天有些經典.
我說話跟聖經同等權並.
要很小心.
先講耶穌暗暗在後.
我不知道你有沒有留意新聞.
剛剛日本幾天前.
出了一個新聞.
定了統一教會.
不只是異端.
也是邪教.
他在日本騙了200億人民.
即是港幣10億.
這個教派很奇怪.
他的始創人是50年代.
在韓國.
韓國很多數字.
信主你要小心.
這個叫文仙明.

$^{601}$50年代出生.
他認自己是誰呢.
他是耶穌.
他老婆厲害.
他老婆是聖靈.
很多人信.
你覺得是不是會笑.
我也想笑.
不過竟然.
他說他才是真正的父.
真正的母.
他是無原罪.
他所生是無原罪.
所以跟很多人發生性關係.
他有配婚.
並且還有一樣很有趣.
他在敬拜的時候.
拜豬頭.
不知道是不是很喜歡吃豬骨.
我覺得你不要叫基督教.
我不理你.
但你踩過來.
我們一定要劃清.
所以第二件事.
我們要小心.
我們怎樣才可以維繫到.
不會信錯.
一定要在正統的教會.
要學習.
要群體裡面成長.
有很多異端並且扭曲.
潛入教會.
第一個問題問.
你懂不懂分辨.
你怎樣分辨.
找誰幫你分辨.
第二.
當你遇見這種被迷惑.
他很多都是好人.
你怎樣幫他.

$^{641}$我們教會都不少.
曾經過去被迷惑.
接著回來.
當然我不會告訴誰.
也不需要.
但是我們常作準備.
很好的.
自己搞清.
也幫助別人.
好了.
第三件事.
我覺得是最難防的.
第一個其實在防.
天國在哪裡.
是不是有形態.
不是.
在神的心.
在人的心裡.
第二.
接著說.
我們要防備很多異端.
不是在這裡那裡.
第三是最難防.
第三是甚麼呢.
我們一起看.
在二十六到三十字一起讀.
一二三.
狼牙的日子怎樣.
人子的日子也要怎樣.
宏水就來.
把他們全都滅了.
杜羅德出所當那的雅人.
降下來.
把他們全都滅了.
人子再來的時候.
很突然.
狼牙的時候.
其實警告了他們一百年.
都不回應.
水一來就完了.

$^{681}$接著說羅德的時候.
發生.
走的時候.
火和硫磺降下就滅了.
說甚麼呢.
這一點.
我們在平常中.
我們要察覺有危機.
當時又吃又喝又買又賣.
好像沒事一樣.
我們要小心.
好像溫水煮蛙.
在這個世代走著.
以為沒事.
其實危機就在面前.
他用這個作比喻.
非常之重要.
他說當這些人.
毫無警覺.
也不知道審判臨到.
突然來到的時候.
就如夢初醒.
在那個時刻.
羅亞和羅德.
不是每個人都是.
窮凶極惡的.
當然有很多惡人.
在那裡描述.
但也有不少的人.
沉溺在日常的事務.
對神的警告.
完全麻木.
置之不理.
這是最大的屬靈危機.
往往都不是明目張膽犯罪.
就是將這些東西正常化.
看不到.
將神邊緣化.
不公不義.
當沒有發生.

$^{721}$失去永恆的警覺性.
有些日常生活我們知道.
很多時候我們覺得.
習慣了其實都怪怪的.
例如我自己很喜歡看電影.
不知道你喜不喜歡看電影.
我已經看了三部電影.
只是上個月.
我想再看多一點.
有些不好看.
有些好看.
但是看電影.
如果你喜歡看.
看電影之前.
都會有段短片.
卡通片.
一分多鐘說.
你要關掉你的手提電話.
你光的時候就弄暗.
別人看電影不要吵著人.
如果你看電影多.
你就知道.
整天看電影中間都會有電話響.
為什麼?.
他都不理你的短片說什麼.
我理你都傻了.
今天耶穌說很多話.
可能沒人理會.
還有些更有趣.
電話都響了.
按熄它.
之後不會聽.
喂?.
我在看電影.
其實這段是沒有內容的.
一會兒說再見就關掉.
你不要聽一樣.
有時我又不肯.
有時我電話突然響.
我都馬上關掉.

$^{761}$我都不記得.
你都不要.
喂?.
剛才有人電話響.
現在又到我.
很怪的.
很多自我的重要性.
又有一個.
我叫黑色幽默.
這些英國人最喜歡的.
這個網上看的很有趣.
他說有班人去打獵.
兩個兩個分組.
一班人兩個兩個去打獵.
接著早上回來了.
最後一組.
他回來了.
只有一個人回來.
明明是兩個.
那個人回來.
就背著一個很大的鹿.
你知道打獵鹿很厲害.
你背著鹿很厲害.
但是你的夥伴.
亨利呢?.
怎麼不見了?.
我們打一打.
他暈了.
可能有一點中風.
我們的同伴就說.
你為什麼不背著亨利回來.
背著鹿回來.
真是的.
我背著亨利回來.
鹿就會被人偷.
鹿值錢的.
我背著鹿回來.
亨利是沒有人理的.
我一會兒再去.
先背著有錢的東西.

$^{801}$其實是一個幽默的笑話.
這個世代就是這樣.
犧牲著一些人.
不斷.
無論那個社會價值觀.
那種賺錢的方法.
其實是在死亡中賺錢的.
遊戲機大麻毒品賭錢.
還有合法化.
賭完馬賭.
不要緊.
再賭大一點.
賭完人.
很多人一直在破產.
沉溺.
繼續這裡好像沒事.
很危險.
他說這種狀況.
看不到這種慢慢邊緣化.
沉醉在世界價值觀當沒事.
我剛才也說了.
有下半截的時刑.
在這間附近的茶餐廳.
接著我問他一件事.
喂.
你做侍應又要回教會.
很難的.
你星期日怎麼辦.
他說.
我是這樣的.
我每個星期日.
他住元朗.
他回這裡.
我說你為什麼回這麼遠.
那個侍應說.
我就跟老闆講明.
我星期日至少給我兩個星期.
可以回崇拜.
不上班.
為什麼.

$^{841}$回崇拜成長很重要.
我必須要講一件事.
對一個侍應來說.
我相信每一元都很重要.
我們有些專業人士.
賺少一點.
但當然有些專業人士.
做好工作.
六天還要賺錢.
我覺得是很離譜的.
我講下一個侍應.
我現在還在找一份.
希望可以回到三個崇拜.
因為我不僅要崇拜在教會成長.
我要服侍主.
神國在我們裡面.
不僅是自己成長.
我們要服侍主.
不是世界在說話.
他說話很有趣.
我自己很欣賞.
當然我聽了一段時間.
我突然間心想.
其實你看到教會.
星期六有崇拜.
不過我當然不講.
搶羊不做.
不做這件事.
我為他祝福也很感動.
但我知道我繼續跟他說.
他應該快要被解僱.
因為講了很久跟他聊天.
他要工作.
第三件事我們要檢視行程表.
是不是為了又吃又喝又做.
等等等等地上的事.
而忘掉了上帝.
靈修崇拜家庭服侍等等.
我們怎樣保持醒覺.
這一點是很難回應的.

$^{881}$因為有些人不是大奸大惡.
就是在平常中.
在紅樓裡迷失了.
怎樣可以在裡面.
不僅是我們的人生不是活在當下.
是有永恆.
和有未來.
這種眼光活在今天.
是第三點我們需要學習.
我們這樣的警醒.
其實是需要付出代價.
和劃分時間.
這是根本對我們一個很重要的動力.
亦是蒙恩進入神國必然的表現.
第四點.
是一個很重要的高峰.
耶穌想和我們解決一件事.
31到33節.
當那日人在房上.
東西在屋裡.
不要下來了.
人在田裡也不要回家.
你們要互相羅登的妻子.
凡喪要好處生命的.
必喪掉生命.
凡喪掉生命的.
必救好生命.
他開始講的形態.
重點不是你聽到多少.
你明白多少.
或者你感動多少.
重點是你要作出決絕的抉擇.
不要回頭的跟從.
他說那日在房上東西房屋裡.
不要回去拿了.
在田裡也不要回家.
回想羅登的妻子.
大家明白嗎?羅登的時候突然間.
火同硫磺降下.
神仙警告了一家叫他走.

$^{921}$但是他走著走著.
不是不走,是走著走著,突然間望回頭.
命都沒有了.
他提醒我們要行動不要眷戀世俗.
不要眷戀你覺得認為跟從性的違反.
不要回頭.
正在走的了,這是很貼切的.
因為很多時候我們走到一半.
突然間覺得.
先看一看.
好東西.
我又要天國.
我又要世界.
通實.
唯識.
沒有這東西.
你要選擇.
耶穌強調屬靈的原則始終不變.
我們有兩個生命.
一個今世的生命,一個永恆的生命.
我們要保存今世的生命.
眷戀回頭就失去永恆.
其實兩個可以沒有抵觸的.
當我們好好擺位.
活好永恆的生命.
今生是準備著永生.
不回頭.
說到這裡.
我想起一些例子.
雲峰縣的確是很傷痛.
對全香港人都是.
兩個多月.
報道會做了一件事.
很想回應,很想讓他們見到.
耶穌的愛.
還有我們可以做一些事.
如果大家記得,我們在十月尾.
我們開始籌集資金.
我們總共籌集了六百多萬.
每一號都給了他們.

$^{961}$是嚴厲之急.
有些上學要馬上買鞋.
有些老人家沒了輪椅馬上做.
總共一千九百戶.
我們探到一千四百戶.
每一戶我們按照他們的人.
多少老人家,有些過了生.
我們有些希望幫到的.
小小的金錢.
可以後視.
有些就是剛才說的.
買到學費等等,我們分好.
分好不單止我們分好,逐個聯絡.
我們坐在這裡聽他們說.
如果沒有時間,不想說,沒問題.
我只是坐在這裡問.
你們安排好了沒有?.
有些就說了很久,很傷心.
有些說完拿張支票,再回來說.
我們很願意.
我記得有很多場面.
有兩個我想特別在這裡分享.
第一個,發生事後那幾天.
我們馬上去了.
我們過渡會的墓者,我們買了很多墓者.
其中一個,我很記得.
有一個長者,老人家.
他的衣服.
他跟我說,他知道我是牧師.
這件衣服.
其實是我三天前.
離開榮福院的同一件衣服.
我只是說.
我要跟他買這件衣服.
他說不用了.
這件衣服也夠了.
我突然想起.
其實人需要多少東西呢?.
其實一件衣服可能也夠了.
我們一開衣櫃.

$^{1001}$一開我們的鞋櫃.
女士的衣櫃.
永遠都欠一件衣服.
永遠都欠一件東西.
對不起,我只是這樣說.
男士永遠遊戲機都欠一個遊戲.
即是上網.
永遠都還有東西想要.
這個老伯伯.
今天在我們分行,已經順道住在這裡.
我介紹了他去分行.
因為他臨時住屋在啟德.
另一個場面.
我也很想說.
我們有幾張椅子.
他們坐在這裡聊天.
有一個人抱著他的背包.
很緊,不放.
我說你放下吧.
慢慢說吧,很安全.
他說不用了,我不放.
其實我知道.
然後他也說.
我那天也是拿了背包.
如果我打算拿多一點.
我都好像街坊一樣.
你今天見不到我.
我也有點想走.
他說我要的其實不多.
這個背包我就夠了.
其實我們今天要些什麼呢?.
捐亂些什麼呢?.
背包可能就夠了.
這個世界有很多東西吸引.
在經濟學上說.
unlimited desire.
即是你無窮無盡.
你的渴求.
但聖經告訴我們.
知足才是上路.

$^{1041}$第四樣.
他問我們.
我們有沒有羅德妻子式的捐亂.
我們要多一點,大一點,漂亮一點,好一點.
是不是真的我們想要的呢?.
如果我們不是有這樣的歷程.
是不是這樣就等於.
我們沒有人生呢?.
還是說上帝在我們生命中.
有不同的計劃呢?.
有沒有一些罪惡世俗的價值觀.
你依依不捨的.
不奮氣的.
令你無法轉向神呢?.
又或者被這件事.
困住了.
走不出來.
有沒有願意.
堅守信仰永恆的產業.
付出代價呢?.
可能是你時間.
可能是你分辨出來.
或者是你.
甚至是你金錢去服侍.
去參與這個神國的事工.
那裡又是什麼呢?.
第四個.
我們要作出決擇.
並且一個不回頭的決擇.
求主幫助.
最後第五.
也是今天的終極.
相當恐怖的一幕.
我們來看看.
34到37節.
一起讀好不好?1,2,3.
我對你們說.
當那一夜.
要取去一個.
撇下一個.

$^{1081}$要取去一個.
撇下一個.
他講的一個場面.
就是當.
末世中的末端.
神國回來的時候會怎樣呢?.
那一夜兩個人在一張床.
即是夫婦.
有一個真正的信主.
被提去.
我們叫被提.
一個撇下.
突然間你回想.
床大了這麼多.
沒得遇.
又或者一個.
推磨.
即是兩個人在磨墨紙.
弄得幼細.
本來很輕散的.
突然很重.
一個不見了.
到時沒得遇.
他已經告訴你.
所以最後這個.
是一個非常艱難.
的預警.
是一個分離的審判.
你多親密的關係中.
都會區分.
不是因為你是好朋友.
不是因為你是家人.
區分只是因為.
你跟隨主神國在你裡面.
才可以發生的一件事.
當時那些宗教領袖.
覺得他亞伯拉罕子孫.
在一個家族在一個民族裡面.
我們就得拯救.
沒有人.

$^{1121}$可以帶你走.
除非你跟上帝建立.
這種關係.
說到這裡我想起.
有一兩個場面.
我兒女都長大了.
我記得他們小時候.
六歲我兒子.
三歲我女兒.
一起去日本玩.
我兒子坐不穩.
有點像我.
有些角度會弱.
他就周圍走.
當時我都會.
推著嬰兒車.
推我女兒.
我兒子就永遠不坐嬰兒車.
就喜歡推嬰兒車.
我女兒就永遠不會不坐嬰兒車.
所以他們是最佳拍檔.
我女兒就坐嬰兒車.
我兒子就喜歡推嬰兒車.
我兒子就推得很快.
以為那輛是弗拉利.
他是一個.
弗拉利的賽車神.
在日本推來推去.
閃來閃去.
像東京一樣很多人.
我記得那次.
不知道在找路還是在看什麼.
很短可能幾十秒.
一看.
賽車手不見了.
乘客都不見了.
大件事.
那怎麼辦.
在日本不懂英文.
他們不講英文.

$^{1161}$很難的.
廣東話一定不懂.
國語更加不會.
當時是廿年前.
怎樣生存.
大人都難.
他現在沒有電話.
很害怕.
找啊.
結果找不到.
找不到根本不講.
找不到就一講就哭.
找到.
很開心.
我跟兒子說.
你由開車到哪裡.
你都要看回頭.
看到我的樣子才好.
至少我看到你.
他也害怕.
所以不坐車.
但我想告訴大家.
在門口的時候.
你是找不到的.
沒得找.
你最愛的那些.
都找不到.
永遠隔離.
他說.
屍首在難.
鷹也在難.
鷹是找腐屍吃.
你見到鷹就見到.
哪裡的腐爛是顯然易見.
那些訊號明顯.
見到.
不過你理不理.
你理不理.
很顯然易見.
見到.

$^{1201}$問題是.
會不會發生.
是何時發生.
怎樣發生.
所以第五個問題.
很重要.
你是否意識到.
屍首已經充滿大地.
那些記取.
到處都是.
審判的時候.
臨到.
審判分離.
末日來臨的時候.
神會根據你的信仰和生活狀態.
是否反映神國在你內心.
作出分別.
取去一個.
撇下一個.
重點是你是哪一個.
我再說一次.
取去一個撇下一個.
耶穌說了.
重點是你是哪一個.
我經常形容聖經是.
open book exam.
但大家有考試.
考試.
以往小學都是通常.
背完就去默書.
大學不玩這東西.
開給你看.
你怎樣形容.
你見到.
怎樣形容.
很重要.
大家不要搞錯.
這經文不是要我們恐懼.
或者只是停留在恐懼.
相反.

$^{1241}$要我們按著這種內容.
活出有盼望的反思.
究竟為甚麼而活.
生活是否正正反映.
神國在我們的生命中.
面對這麼多未來.
不確定可能會更加.
混亂的世界.
現在經常說這個叢林.
的.
生活形態.
怎樣.
可以知道天國.
近了.
怎樣忠心警醒.
活出每一天.
結語.
我們請敬拜他.
耶穌今天的訊息.
非常嚴肅.
和迫切.
天國已經隨著.
他的說話和記號.
我們開始見到很多的景象.
當然.
我們沒辦法.
計算那一天.
但是我們.
卻可以決定我們今天.
怎樣生活.
我再說一次.
耶穌說了會發生.
不過我們今天.
已經要決定.
我們怎樣活出每一天.
叫我們準備好末世.
總結.
我們學了.
五件事.
第一.

$^{1281}$學了天國內住性.
神國.
應該在我們裡面.
第二.
知道再來的突然性.
如果我們今天不是這樣活著.
或者不是我們.
我們旁邊的人.
你覺得可以慢慢愛的人.
你跟他說.
閃電一樣.
第三.
生活的警覺性.
小心跌入世俗的價值觀.
幾種意識形態.
溫水煮蛙.
對神麻木.
第四.
抉擇的徹底性不回頭.
羅亞妻子.
第五.
審判的個人性.
我們只可以因為主.
和我們的關係.
可以活到永恆.
我們一起祈禱.
我想在這裡特別為大家禱告.
不知道這段是福音.
這段是訊息.
有什麼停留在你內心.
有什麼你很想突破.
有什麼很想聖靈幫助你.
可以越走越有力.
越走越有盼望.
如果你是很想主.
我明白你的心.
我很想跟隨.
你所說的而行.
如果你是其中一個.
我鼓勵你這一刻.

$^{1321}$將你的手放在你的心.
我為你禱告.
重點不是我們知多.
重點是我們活出如何.
所以在這時候.
我想特別為你禱告.
神我們多謝你.
親愛天父.
讓我們看見天國的奧秘.
和人子必定會再來的真實.
求你救我脫離.
屬靈的麻木.
和對世界的眷戀.
甚至對自己裡面.
自我.
比屬靈的真實更強.
神求你.
救我脫離這些.
賜我謹整的心.
天天在生活中.
尊你為王.
以永恆的眼光.
活出今日.
直至到見你面.
主我們多謝你.
你的提示.
必然帶來我們生命的更新.
讓聖靈自由在我們生命工作.
我們歡迎你.
求你幫助.
讓我們有著力量.
快步踏前.
跟隨你.
多謝你.
聽我們禱告.
奉耶穌名求.
\newpage



\section{出埃及記 3:1-22}
\label{sec:031779}
\textbf{呼召的關係}~~~~江芷敏傳道
\newline
\newline
連結: \href{http://www.church.com.hk/acms/content.asp?site=cdc\&op=show\&type=product\&code=031779\&layout=sermon}{\texttt{031779}} ~~~~ 語音日期: 2026-03-08
\newline
\newline
\hyperref[sec:031778]{\small{< < < PREV SERMON < < <}}
~
\hyperref[sec:index_chronic]{\small{[返順時目]}}
~
\hyperref[sec:031780]{\small{> > > NEXT SERMON > > >}}
\newline
\newline
出埃及記 3:1-22
\newline
\begin{longtable}{cl}
\hline
\hline
章節 & 經文 (和合本修訂版)\\
\hline
3:1 & \begin{tabularx}{0.7\textwidth}{X} 摩西牧放他岳父米甸祭司葉特羅的羊群，他領羊群往曠野的那一邊去，到了神的山，就是何烈山。 \end{tabularx} \\ \\ \relax
3:2 & \begin{tabularx}{0.7\textwidth}{X} 耶和華的使者在荊棘的火焰中向他顯現。摩西觀看，看哪，荊棘在火中焚燒，卻沒有燒燬。 \end{tabularx} \\ \\ \relax
3:3 & \begin{tabularx}{0.7\textwidth}{X} 摩西說：「我要轉過去看這大異象，這荊棘為何沒有燒燬呢？」 \end{tabularx} \\ \\ \relax
3:4 & \begin{tabularx}{0.7\textwidth}{X} 耶和華見摩西轉過去看，神就從荊棘裡呼叫他說：「摩西！摩西！」他說：「我在這裡。」 \end{tabularx} \\ \\ \relax
3:5 & \begin{tabularx}{0.7\textwidth}{X} 神說：「不要靠近這裡。把你腳上的鞋脫下來，因為你所站的地方是聖地」。 \end{tabularx} \\ \\ \relax
3:6 & \begin{tabularx}{0.7\textwidth}{X} 他又說：「我是你父親的神，是亞伯拉罕的神，以撒的神，雅各的神。」摩西蒙上臉，因為怕看神。 \end{tabularx} \\ \\ \relax
3:7 & \begin{tabularx}{0.7\textwidth}{X} 耶和華說：「我確實看見了我百姓在埃及所受的困苦，我也聽見了他們因受監工苦待所發的哀聲；我確實知道他們的痛苦。 \end{tabularx} \\ \\ \relax
3:8 & \begin{tabularx}{0.7\textwidth}{X} 我下來是要救他們脫離埃及人的手，領他們從那地上來，到美好與寬闊之地，到流奶與蜜之地，就是迦南人、赫人、亞摩利人、比利洗人、希未人、耶布斯人之地。 \end{tabularx} \\ \\ \relax
3:9 & \begin{tabularx}{0.7\textwidth}{X} 現在，看哪，以色列人的哀聲達到我這裡，我也看見埃及人怎樣欺壓他們。 \end{tabularx} \\ \\ \relax
3:10 & \begin{tabularx}{0.7\textwidth}{X} 現在，你去，我要差派你到法老那裡，把我的百姓以色列人從埃及領出來。」 \end{tabularx} \\ \\ \relax
3:11 & \begin{tabularx}{0.7\textwidth}{X} 摩西對神說：「我是甚麼人，竟能去見法老，把以色列人從埃及領出來呢？」 \end{tabularx} \\ \\ \relax
3:12 & \begin{tabularx}{0.7\textwidth}{X} 神說：「我必與你同在。這就是我差派你去，給你的憑據：你把百姓從埃及領出來之後，你們必在這山上事奉神。」 \end{tabularx} \\ \\ \relax
3:13 & \begin{tabularx}{0.7\textwidth}{X} 摩西對神說：「看哪，我到以色列人那裡，對他們說：『你們祖宗的神差派我到你們這裡來。』他們若對我說：『他叫甚麼名字？』我要對他們說甚麼呢？」 \end{tabularx} \\ \\ \relax
3:14 & \begin{tabularx}{0.7\textwidth}{X} 神對摩西說：「我是自有永有的」；又說：「你要對以色列人這樣說：『那自有永有的差派我到你們這裡來。』」 \end{tabularx} \\ \\ \relax
3:15 & \begin{tabularx}{0.7\textwidth}{X} 神又對摩西說：「你要對以色列人這樣說：『耶和華－你們祖宗的神，就是亞伯拉罕的神，以撒的神，雅各的神差派我到你們這裡來。』這是我的名，直到永遠；這也是我的稱號，直到萬代。 \end{tabularx} \\ \\ \relax
3:16 & \begin{tabularx}{0.7\textwidth}{X} 你去召集以色列的長老，對他們說：『耶和華－你們祖宗的神，就是亞伯拉罕的神，以撒的神，雅各的神向我顯現，說：我實在眷顧了你們，眷顧你們在埃及的遭遇。 \end{tabularx} \\ \\ \relax
3:17 & \begin{tabularx}{0.7\textwidth}{X} 我也曾說：要把你們從埃及的困苦中領出來，往迦南人、赫人、亞摩利人、比利洗人、希未人、耶布斯人的地去，就是到流奶與蜜之地。』 \end{tabularx} \\ \\ \relax
3:18 & \begin{tabularx}{0.7\textwidth}{X} 他們必聽你的話。你和以色列的長老要到埃及王那裡，對他說：『耶和華－希伯來人的神向我們顯現，現在求你讓我們往曠野去，走三天的路程，為要向耶和華我們的神獻祭。』 \end{tabularx} \\ \\ \relax
3:19 & \begin{tabularx}{0.7\textwidth}{X} 我知道若不用大能的手，埃及王不會放你們走。 \end{tabularx} \\ \\ \relax
3:20 & \begin{tabularx}{0.7\textwidth}{X} 因此，我必伸出我的手，在埃及施行我一切的神蹟，擊打這地，然後，他才放你們走。 \end{tabularx} \\ \\ \relax
3:21 & \begin{tabularx}{0.7\textwidth}{X} 我必使埃及人看得起你們，你們離開的時候就不至於空手而去。 \end{tabularx} \\ \\ \relax
3:22 & \begin{tabularx}{0.7\textwidth}{X} 每一個婦女必向她的鄰舍，以及寄居在她家裡的女人，索取金器、銀器和衣裳，給你們的兒女穿戴。這樣你們就掠奪了埃及人。」 \end{tabularx} \\ \\
[1ex]
\hline
\hline
\end{longtable}
$^{1}$大家好.
早晨 早晨.
時間真的過得很快.
今天又已經是夏令的時間.
我記得上一次說到好像是冬令時間.
我們剛剛過完新年.
又進入這個大齋期.
很快就到什麼節日了.
復活節了.
而且還是UV的十週年.
教會很快就十歲了.
其實最近我經常都在想.
教會快要十週年了.
有什麼東西可以送給教會.
獻情給神呢.
如果我自己有個子女是十歲的話.
我都會跟祂好好慶祝.
更何況是十週年.
我想都會很想跟祂展望一下未來.
所以今天這一篇講道.
很想跟大家一起去思想.
我們在這個十週年裡面.
可以獻情一份怎樣的禮物給我們的神.
今天我們會看《出埃及記》的.
三章 十九章和三十二章.
是有點貪心.
不過大家可以預備好你的聖經.
隨時打開我們一起去看.
我都會在PowerPoint裡面.
都有展示到一些經文的.
大家不用擔心.
《出埃及記》三章.
字有些小.
不過不要緊.
我會跟大家講一下故事.
《出埃及記》三章一至六節.
其實是一個上帝呼召摩西的經文.
摩西其實是第一次遇見神.
神呼召他帶領以色列人出埃及.
其實那時候.

$^{41}$摩西是在曠野那裡.
牧放羊.
帶著他外父的羊.
一路走一路走.
去到曠野的另一邊.
那座山叫做荷列山.
就在那裡神主動跟他顯現.
你見到有光.
在荊棘裡面著了火.
但是那些荊棘是燒不爛的.
所以摩西覺得很奇怪.
走過去看看什麼事.
在這個時候.
神就叫他.
摩西 摩西.
摩西就說.
我在這裡.
神就跟他說.
你不要走過來.
因為你要脫掉你腳上的鞋.
你站的地方是一個聖地.
然後神再介紹自己.
祂說.
我是你父親的神.
是亞伯拉罕的神.
以撒的神.
雅各的神.
摩西聽完之後.
他就遮住自己.
他很怕見到神.
我們要在這裡問.
這兩幅圖.
究竟哪一幅像一至六節所形容的.
剛才的經文.
通常都是白天放羊.
沒理由是晚上的.
但是下面那幅圖.
好像有一樣東西沒有畫到.
就是摩西有沒有遮住自己的臉.
是沒有的.

$^{81}$所以我覺得上面.
黑漆漆的圖比較像.
不過不知道為什麼是晚上.
不要緊.
我們看這兩幅圖.
其實都有一個共同點.
就是他也有畫到.
摩西脫掉了鞋.
其實我經常看這段經文.
我經常在想.
我們回家也脫下鞋子.
是不是走進去.
就是一個比較聖潔的地方呢.
又不是.
為什麼上帝在這裡要描述.
摩西要脫掉鞋子.
因為他站的地方是聖地呢.
聖.
就是說聖潔.
聖潔究竟是什麼意思.
聖潔是分別出來的意思.
就是從世俗裡面分別出來.
究竟怎樣去定義.
哪裡是世俗.
哪裡是聖潔呢.
有什麼原則呢.
其實聖潔也是一個.
頗抽象的概念.
在摩西五經.
其實經常提到.
關於聖潔的經文.
或者他的條款.
不過我今天不是想解釋.
什麼是聖潔.
因為要講很久.
聖潔在這裡.
每一次都出現在.
上帝呼召人的時候.
他總是呼召人的時候.
要叫人去自潔.

$^{121}$為什麼呢.
我們有看《士師記》.
《先知書》.
譬如好像《耳塞》.
上帝呼召他的時候.
都叫他自己要潔淨.
為什麼上帝要給一個.
這樣的一個動作.
人去聽他的呼召呢.
其實上帝每一次提到.
這個聖潔的要求的時候.
他是劃定了.
什麼的界線是為之聖潔.
就好像摩西那樣.
你站在這個地方.
其實你天天都在這裡放羊.
但是我就將這個地方.
今天分別為聖.
你要脫下鞋子.
才可以站上去.
我們就看看.
其實摩西天天都在.
這個地方裡面.
他天天都放羊.
為什麼上帝要將這個地方.
變為聖潔呢.
如果我們想一下.
我們天天工作的地方.
就好像我們的辦公室.
其實一點都不聖潔.
還想快點逃離那個地方.
或者現在你的工作地方.
變成了你的家.
你都不會覺得.
為什麼有一個.
這麼特別的景象.
或者想像不到.
上帝會在那裡呼召我們呢.
其實上帝.
好像在跟摩西說一樣東西.

$^{161}$你現在要聽我說.
你來到這個地方.
不再是你平日.
很熟悉放羊的地方.
今天這個地方是我主事的.
今天這個地方.
是我要去呼召你的.
上帝將這個對話.
搞清楚那個.
誰是主事的那個.
每一次上帝呼召人的時候.
都要搞清楚這個次序.
是我才是主.
是我給你這個聖潔的界線.
我一直在寫的時候.
就想起主人.
和貓,狗,寵物的關係.
其實貓和狗.
牠們都有自主權的.
其實你在家裡面.
牠都可以選擇親不親近你.
而且可能有時候關係掉轉了.
你要服侍牠.
又要鏟屎.
又要摸牠.
要給牠足夠的關注.
不過你始終是主人.
因為你是買東西給牠的.
有病的時候.
是你帶牠去看醫生.
而且你是要餵東西給牠吃的.
你和牠有一些不同的操練.
或者你要給牠一些規矩.
當那些規矩一出來.
聽到指令的時候.
這些寵物.
這些貓和狗.
就知道你就是那個主人.
寵物是要聽主人的話.
我想起這個關係.

$^{201}$當然不是我們直接代入.
我們就是寵物.
我們比寵物尊貴得多.
或者自由得多.
希望大家不要誤會我這個例子.
但我正想說出.
當上帝給這個指令的時候.
我們要聽.
我們要服從.
那個關係就變得清楚了.
牠是主.
我們是跟從.
上帝給了我們很大的自由意志.
當祂跟我們講這些指令的時候.
如果我們選擇去跟從祂的話.
我們選擇去跟從祂的話.
我們就和祂恢復一個.
從屬的關係.
好像摩西一樣.
摩西和上帝.
恢復一個從屬的關係.
祂知道誰是主.
當然出埃及記.
不單止是講呼召摩西一個人.
呼召是一個進程.
慢慢祂由呼召摩西.
變成呼召整個以色列的群體.
就在出埃及記十九章那裡.
上帝又再一次去呼召以色列人.
耶和華神就跟摩西說.
你現在去到百姓那裡.
叫他們明天要自潔.
叫他們洗衣服.
然後第三天我就會降臨在西奈山上面.
其實這個西奈山和荷里山都是同一座山.
不過是山脈和山峰的關係.
來到這一座山的裡面.
上帝要為他們定界限.
祂說你們要謹慎.
不可以上去這個山.

$^{241}$也不可以摸這個山邊.
只要你摸這個山就會致死.
第十三節.
祂說也不可以伸手摸他.
因為必被石頭打死.
說用箭射頭.
無論是人還是畜生都不得活.
這樣的聖潔的界線是誰定的?.
是上帝定的.
這次是不是跟呼召摩西的時候一樣?.
是不一樣的.
所以每一次被呼召的時候.
那個聖潔的界線很不一樣.
其實上帝每呼召一個人.
那個聖潔的界線不一樣.
其實不是我們要深究.
為什麼他這次被呼召.
他要做這個聖潔的儀式.
其實是每一次那個聖潔的標準來到的時候.
我們就知道上帝就是那個主.
我們就是要聽祂說.
祂就是被呼召的那個上帝.
所以在這裡我們又再看到一次.
那個呼召是我們都要擺回適合我們的位置.
當上帝的呼召來到的時候.
我們這個群體就要知道.
誰是我們的主.
擺回一個適合的姿態.
這段經文不單止是說.
我們要擺回一個適合的姿態.
擺回一個適合的位置.
去知道上帝是上帝.
更加是調整我們的眼光.
第三章 神與摩西的對話.
繼續說.
祂說:我確實看見了我百姓在埃及所受的困苦.
我也聽見了他們因受監宮苦袋所發的哀聲.
我也確實知道他們的痛苦.
這段經文是上一次跟大家講到的時候也有提過.
不過三章這裡再重覆一次.

$^{281}$上帝是那位親近我們的上帝.
摩西聽到上帝的呼召.
說他要去.
上帝要差派他去法老那裡.
並將他的百姓以色列人從埃及領出來.
摩西第一個反應是什麼?.
第十一節 摩西說:我?.
我是什麼人?.
竟然可以去見法老?.
還要將以色列人從埃及那裡領出來?.
如果大家對聖經有些了解.
會知道摩西究竟有什麼背景.
摩西以前是埃及的王子.
他在那裡成長.
但是他很看不過眼這些埃及人去欺負他的同胞.
所以他打死了那個欺負他同胞的埃及人.
然後他就逃難去到米甸做牧羊人.
在曠野那裡一留就留了40年.
其實他已經和外面這個世界脫節了.
他不知道外面是發生什麼事.
埃及已經離他很遠了.
那時候沒有新聞 不會像我們現在這樣.
今天發生什麼事 明天全世界都知道.
他還要面對一個不可以撼動的法老.
帶以色列人出埃及 憑什麼呢?.
但是上帝的回應是什麼?.
我們十二節看到是什麼.
他說:我必與你同在.
我必與你同在.
上帝沒有回答他那個問題.
我究竟是什麼人.
上帝也沒有說:你行的.
我為你打氣 你一定行的.
上帝沒有說這些.
上帝是說:我與你同在.
我與你同在.
而且第十四節繼續上帝也說.
我是誰? 我是自有永有的那個.
你看到這些經文全部都是「我我我」.
摩西就問上帝:我是什麼人?.

$^{321}$我有什麼可能? 有什麼quality?.
上帝沒有回答他這個問題.
不過上帝也繼續用「我」這個字去開頭.
說:我必與你同在.
我是自有永有的那個上帝.
I am who I am.
其實翻譯成英文也不知道是什麼.
或者翻譯成原文.
也不知道究竟他這個名字是什麼意思.
I am who I am.
好像沒有什麼可以定義上帝.
我們回顧古晉東一些文化,一些神明.
他們很多時候都用一個名字去顯示.
他自己有什麼功能.
他有什麼屬性.
就好像太陽神.
就好像風暴神,月亮神這些女神.
你一聽他的名字就知道他要對付什麼.
或者他有什麼功能.
其實的確這樣的神明都比較容易捉摸.
因為一說就知道他是代表什麼上帝.
什麼神.
但上帝就在這裡用「I am who I am」.
「我是自有永有」.
他去回應摩西.
他在這裡否定人為他下一個定義.
否定那些關注他的名字.
否定那些對他定義的定義.
很有趣.
上個星期就是年宵.
在年宵裡面.
我們教會也有擺一個Outreach的攤.
那個攤就是玩遊戲和請人喝咖啡.
很有趣.
有一個U Chicago的學生.
他一直喝咖啡和Ken說.
Ken就介紹他.
這個咖啡可以喝到什麼什麼水果味.
因為我們沖手沖.
很多時候說這個咖啡有什麼朱古力味.

$^{361}$有什麼荔枝味.
我們喜歡喝咖啡的人.
慢慢學懂這些.
品嘗這個咖啡.
但那個學生很有趣.
他就說.
為什麼咖啡不可以是咖啡味.
很有趣.
今天Ken就和我分享.
他也是.
為什麼咖啡不可以是咖啡味.
為什麼要用其他味去定義.
這個咖啡是什麼味.
神就好像這樣.
他沒有給其他東西去定義.
祂就是這個世界沒有意思.
祂已經在這裡的上帝.
祂就是創造的上帝.
祂就是和我們一起的上帝.
而且上帝說.
你不要再看你自己了.
你整天說我我我.
你現在看一下.
我究竟是哪一個.
上帝在這裡和他說.
不要再看你不行的東西了.
不要再看著自己的限制.
你的焦點應該是要看著.
我這個呼召你的神.
其實我也很明白摩西.
他的掙扎.
因為神叫他做那件事.
和他的實力差距很遠.
你完全覺得沒有可能.
的確上帝有時候.
叫我們做的事是沒有可能.
其實我又想起一件.
很沒有可能的事.
其實我們生命裡.
都出現過很多不同的上帝.

$^{401}$還是大靈.
但我覺得最沒有可能的.
就是上帝帶領Ken.
來到芝加哥這裡讀書.
因為我們以前讀書.
都是身無分文.
也覺得自己的英文很有限制.
在我們十年前讀神學.
就是2016年的時候.
Ken有一次無緣無故打給我.
那時候我們還沒有拍拖.
我和他都沒有什麼關係.
但因為我純粹住在長洲.
他遺留在長洲.
他有需要就打給我.
在那個星期六.
他發燒.
因為星期六就要看急症.
在香港.
看急症當時是要100元.
但他的戶口又不夠100元.
所以按不到那張100元出來.
所以他就打電話給我.
問我借錢.
現在這些電話就要很小心.
但那時候他就問我借100元.
我都覺得他很慘.
又發燒.
又要看急症.
又沒有人在他身邊.
家人又在深水Po 住.
我就馬上去長洲醫院.
借錢給他.
我那時候都很難想像到.
一個身無分文.
身無分文但一兩千都有.
但他連100元都沒有.
但他竟然來到芝加哥讀書.
回望我們過去.
他完成了課程.

$^{441}$現在還差一篇文.
但我們用了很多錢.
我們回想.
其實這些都是上帝給我們的恩典.
上帝也是和我們同在.
和恩典很夠用的上帝.
我們計算一下.
我昨天問他.
你猜一下我們這三年.
其實應該用了多少錢呢.
Trinity的網站就說.
我們一年要用五萬美金.
兩夫妻.
加上租屋和在那裡讀書.
我們讀了兩年半.
我想怎樣都有十二萬.
我一直都沒有想過.
上帝可以這樣帶領我們來這裡.
還有我們剛剛來到的時候.
好像整桌人在orientation dinner.
吃飯.
我一句英文都聽不到.
好像自己來錯了地方.
很多的質疑.
到今天可以在這裡服侍.
我都覺得完全是上帝的恩典.
如果當日我們要去計算.
計算自己有多少限制的話.
我相信我們不能夠站在這裡和大家分享.
如果我們只是要看著數字的話.
我們是不需要向前走.
但上帝就是那個呼召我們的上帝.
就是要叫我們不要只看著自己的限制.
要看著這個自有永有.
和祂會供應的上帝.
第二點就是.
我們不單止要看自己的不足.
我們更加要將我們的焦點.
放在這個自有永有超過一切的上帝身上.
第三.

$^{481}$呼召其實不是一日完成的一件事.
是一個進程.
就好像出埃及的三章去呼召摩西.
然後到十九章.
上帝是呼召整個的以色列群體.
上帝說.
你們要聽我的話.
你要去作我祭司的國度.
作聖潔的國民.
上帝在十九章那裡是這樣呼召以色列人.
作祭司的國度.
為聖潔的國民.
其實這個進程.
不單止是由一開始摩西第一次見上帝.
就聽到這個呼召.
是一點一滴不斷去塗造摩西.
塗造以色列人.
摩西不斷在神面前去掙扎.
他說其實我真的很不濟.
我又不懂說話.
你派另一個人吧.
摩西是一個很將自己的軟弱.
呈現給上帝看的一個人.
上帝亦都在這一段的關係裡.
沒有走開去聽摩西.
去說他自己多不濟.
你看其實不單止是三章.
第四章都一直在說.
摩西和上帝不斷地去談判.
我們都看到十九章.
以色列人是怎樣回應上帝的呼召.
他不像摩西那樣有很多來回.
他是一個很簡潔的一句說話.
很簡潔的一句說話就是十九章八節.
如果大家有聖經可以打開.
百姓只是和上帝說.
凡耶和華所說的.
我們都要遵循.
只是一句說話就回應上帝的呼召.
其實這是一個很淑寧的群體.

$^{521}$馬上說「Amen」的一句說話.
不過其實我們看到.
十九章之後.
就是上帝去呼召摩西上山.
和他說這個法版.
和說怎樣做一個聖潔的群體.
在他上了山之後.
摩西上了山之後.
百姓見到摩西遲遲未下山.
就和亞倫說.
亞倫起來.
我們要做一個神像.
可以在我們面前引路.
剛才說那些凡耶和華所說的.
我們都要遵循.
已經拋諸腦後了.
亞倫更加去聽這些百姓說.
他說以色列.
你們快點回去拿些金來.
我們要一起去做一個金牛獨.
這個領袖.
是人家說什麼.
他就做什麼.
他完全沒有想過.
然後做了金牛獨之後.
亞倫還說.
以色列.
這個就是領你們出埃及的神.
他一點都沒有糾正.
這一班走錯路的以色列人.
以色列人原來是想要一個.
看得見摸得到.
其實不一定是耶和華神.
你再看看三十二章的經文.
他也當摩西是帶領他們出以色列的神.
只要讓他們看得見摸得到.
有足夠的安全感.
這個就是他們的神.
簡單來說.
誰幫到我誰就是我們的神.

$^{561}$以色列人就是這樣.
用他們的用語去包裝偶像.
其實他們不需要一個神.
他需要一個工具.
如果今天他們有AI.
可能都覺得AI是他們的神.
什麼都問得到.
什麼都解答得到.
原來他們所說的.
只是一個口號式的敬拜.
我再看看摩西和神來回這麼多個圈.
去承認自己的軟弱.
承認自己的掙扎.
這個關係來得更加真實.
呼召是一個關係的開始.
或者我可以形容它是一個坦白的開始.
當我想起我怎樣信主的時候.
我都是一個很軟弱的狀態.
或者我是和神坦白.
我有什麼得罪他.
就好像剛才那首詩歌《真相》.
我慢慢揭開一些包住了我的保護機制.
或者一些我很想欺騙自己欺騙人的東西.
我打開了這些.
見到最真最軟弱的自己.
我們來到神的面前.
那一刻是我和神最親近的時間.
由摩西被呼召的那一天.
他就開始和神建立這個關係.
這個關係不是我們看得很完美.
是一個跌跌碰碰.
一仆一碌的關係.
他也有做錯事得罪神.
但是他和神的關係就很真實.
就好像在三十三章也有說.
神和摩西的關係是怎樣呢.
我們一起去讀《馬雀埃及記》三十三章十一節上.
耶和華與摩西面對面說話.
好像人與朋友說話一般.
好像朋友一樣.

$^{601}$我們再看《申命記》三十四章十節.
以後以色列中再沒有興起先知像摩西的.
他是耶和華面對面所認識的.
這個面對面的認識.
是一個很真實,很坦白的關係.
神體重的不是我們可以做些什麼.
神體重的是我們和祂的關係.
所以這個呼召是叫我們進入和祂一個這麼親密的關係.
是要坦坦白白的.
在我身邊.
我覺得最親密的一個關係.
當然是我和我老公.
但亦都是我和我媽媽.
在這個關係裡面.
我們經歷過很多磨擦.
但是在這些磨擦裡面.
我慢慢覺得我們觸摸到對方究竟是一個怎樣的人.
我和我媽媽小時候很疏遠.
因為她很年輕生我.
為了謀生.
我十歲以前都是和我媽媽一起住.
因為爸爸媽媽經常都不在身邊.
所以我要經常學會看別人的眉頭眼額.
很努力做一個大人症.
很懂得看別人需要些什麼.
但是同時在這個保護自己生存機制的底下.
我對我父母慢慢積累了很多的埋怨.
到我青春期的時候.
我開始回家和他們住.
我就經常都爆發.
去控訴他們之前為什麼要這樣.
一開始說完.
大家都是大哭一場.
我媽媽其實都很愛我,很疼我.
不過她都不懂怎樣去接到我的情緒.
她都覺得很委屈.
因為她已經很盡力地去愛我.
其實她只是十八歲就生我.
她自己都未搞得定.
到我十八歲的時候.

$^{641}$我才慢慢開始明白.
原來做一個大人要這麼多事要顧.
所以我慢慢明白她多一些.
加上我十九歲開始信耶穌.
我們回教會.
我們一家人試著牽著手去祈禱.
每一次想起過往那些關係的傷痛的時候.
我們都會在哭.
但是我們願意彼此去聽.
去坦承我們心裡面的感受是怎樣.
我覺得和我媽媽開始越走越近.
不是真的physically.
而是emotionally.
我知道她為什麼要這樣做.
我開始盛載她的關係.
開始慢慢地越走越近.
是基於那個坦白和神的恩典.
她不斷地光照我們.
神的恩典是很好.
她不斷地讓我們一家人去修補這個關係.
讓我們變得可以彼此擁抱.
我不知道大家有沒有經歷過你的友誼.
都是試過在一些吵架,坦白之後.
大家的感情變得更加牢固.
弟兄姊妹更加容易吵架.
弟兄姊妹之間有沒有試過坦承.
令到你覺得更加close.
上帝在這裡也和我們說.
我們需要和祂有一個坦白的關係.
今天我們去問UV十週年.
我們有什麼可以獻情給我們的神.
在這段經文裡面再一次提醒我們.
上帝是很想和我們恢復一個關係.
在祂的呼召裡面.
我們要回到和祂重熟的關係.
祂才是我們生命的主.
祂才是UV的主.
祂是那個話事的那個.
而我們的眼光又再一次對正這個呼召我們.
和超越一切的上帝.

$^{681}$祂就是那個自由永有超越一切的上帝.
一切都是從祂開始.
UV也是.
而祂也邀請我們和祂進入一個更深.
更坦白的一個關係.
不知道大家在這幾點裡面有什麼想法呢.
或者今天如果上帝來到我們當中.
問我們願不願意再更深一步和祂建立這個坦白的關係.
不單止每個星期日回來.
是在我們平日裡面和祂建立一個更坦白的禱告.
那個同行的關係.
讓我們一起去祈禱.
主啊求你幫助我們.
當我們在這個你如我乍的社會裡面.
在我們營營役役的生活裡面.
我們有時都好像忘記了.
你就是我們生命的主.
主啊求你親自幫助我們.
回復這個和你重熟的關係.
你就是我們的主.
你就是UV的主.
也求你幫助我們.
不要經常面對困難的時候.
就想起自己的限制.
我們能夠多一些去想像.
和多一些去凝望著你自己的能力.
當我們講到教會要去expand的時候.
當我們要去找一個更適合的地方的時候.
當我們要面對很多財務上的困難的時候.
主啊求你幫助我們.
能夠去望著你這個自由永有超越一切.
而且供應很足夠的上帝.
主啊我們都懇求你.
挪去我們心裡面那個金牛毒的敬拜.
我們真的要去來到你的面前.
坦坦白白的去展開我們自己.
讓你去對付我們的生命.
將我們最軟弱的一面呈現給你.
主啊求你.
閱立我們的禱告 閱立我們的敬拜.

$^{721}$奉主名求 阿們.
\newpage



\section{約翰福音 1:1-6}
\label{sec:031780}
\textbf{在恩典裡面對新的挑戰}~~~~潘智翔牧師
\newline
\newline
連結: \href{http://www.church.com.hk/acms/content.asp?site=cdc\&op=show\&type=product\&code=031780\&layout=sermon}{\texttt{031780}} ~~~~ 語音日期: 2026-03-08
\newline
\newline
\hyperref[sec:031779]{\small{< < < PREV SERMON < < <}}
~
\hyperref[sec:index_chronic]{\small{[返順時目]}}
~
\hyperref[sec:031781]{\small{> > > NEXT SERMON > > >}}
\newline
\newline
約翰福音 1:1-6
\newline
\begin{longtable}{cl}
\hline
\hline
章節 & 經文 (和合本修訂版)\\
\hline
1:1 & \begin{tabularx}{0.7\textwidth}{X} 太初有道，道與神同在，道就是神。 \end{tabularx} \\ \\ \relax
1:2 & \begin{tabularx}{0.7\textwidth}{X} 這道太初與神同在。 \end{tabularx} \\ \\ \relax
1:3 & \begin{tabularx}{0.7\textwidth}{X} 萬物都是藉著他造的，沒有一樣不是藉著他造的。凡被造的， \end{tabularx} \\ \\ \relax
1:4 & \begin{tabularx}{0.7\textwidth}{X} 在他裡面有生命，這生命就是人的光。 \end{tabularx} \\ \\ \relax
1:5 & \begin{tabularx}{0.7\textwidth}{X} 光照在黑暗裡，黑暗卻沒有勝過光。 \end{tabularx} \\ \\ \relax
1:6 & \begin{tabularx}{0.7\textwidth}{X} 有一個人，是從神那裡差來的，名叫約翰。 \end{tabularx} \\ \\
[1ex]
\hline
\hline
\end{longtable}
$^{1}$隨著新年又是一個新的開始.
今日我們都想講一個題目.
就是如何在這個恩典裡面.
去面對一個新的開始.
和一個新的挑戰.
我想從一節的經文開始.
面對新的挑戰.
希望我們可以從耶穌基督的恩典裡開始.
在約翰福音是這樣說.
從他的風承裡.
我們都領受了恩典.
而且欣賞佳人.
很喜歡風承這個詞.
Planaroma.
這個詞的原文是指Complete.
即是它是完整的.
是完美的.
我們因為從出生到現在.
我們未經歷過什麼叫做Complete.
未經歷過什麼叫做.
我是足夠,我是不缺乏的.
是完美的.
所以我們就不太明白什麼叫做風承.
因為我們很少人是風承的.
我們從小到大都有各種不同的貧窮.
這些貧窮會讓我們覺得缺乏.
我們缺乏很多的東西.
可能從小到大.
缺乏父母的認同.
缺乏一些人對我們的認同.
尊重,讚美.
各樣各樣的東西.
甚至長大後會發覺.
某個階段會缺乏錢.
我們會缺乏很多的東西.
所以我們不知道什麼叫做風承.
但在經文裡面.
它唯一一個是風承.
就是從它的風承裡面.
去嘗試去領受這些恩典.

$^{41}$所以希望在新的開始.
新的挑戰的時候.
我們回到主耶穌的風承裡面.
去領受一下.
其實祂是風承的.
祂是什麼都不缺乏.
是我們自己覺得缺乏.
所以我們就投射了一個貧窮缺乏的.
一個影像在我們身邊的環境.
但如果我們從它的風承裡面.
去領受.
就是一幅不同的影像.
所以這個風承是非常重要.
我很喜歡在寫約翰福音的時候.
其中有一個寫commentary的神學家.
是這樣形容約翰福音.
他說約翰福音.
他給它起了另一個名字.
叫做風承福音.
The Gospel of Abundance.
他是這樣解讀約翰福音.
他說約翰福音的典題.
其實他覺得是這一節的經文.
他覺得接下來.
整本約翰福音耶穌的每一個事件.
其實都是在.
具體彰顯.
它的風承.
在什麼層面流露出來.
所以這個風承.
不是一個很抽象的.
它是一個很具體化的.
我們可以生活中.
經歷的東西.
oh sorry.
就好像這節經文所說.
歌羅西書二章九節.
因為神本性一切的風承.
祂所有最好的東西.
都是有形有體.

$^{81}$所以這份風承.
不是抽象你觸摸不到.
祂是很有形有體的.
居住在基督裡面.
我們在祂裡面.
可以得著這種風承.
所以這種風承.
是神不單止給了基督.
由基督可以.
分享給我們.
你要在祂裡面.
不是在祂外面.
我們經常在基督以外.
想尋求這種風承.
但是你花了一生人的時間.
只要在基督外.
你得到的最後都是缺乏.
都是貧窮.
一點風承的感覺都沒有.
因為我們一開始追求風承的方式.
根本就是錯的.
這種風承是神在基督裡面.
留給我們.
一切的風承有形有體.
神在基督裡面解釋.
這種風承是怎樣流露出來.
例如我們在新年的時候.
每一次到了新年.
我們總是有一種焦慮.
什麼焦慮呢.
通常到了某個年紀.
你會發覺有一種焦慮.
叫做時間的焦慮.
你總是覺得.
過了一年我又打碎了.
你又覺得時間不多了.
你就覺得很焦慮.
所以對時間的焦慮.
是很貧窮的一種地方.
因為我們從小到大有期望.

$^{121}$應該幾多歲完成什麼.
你眼看著.
又進入中年危機.
在我這個層次.
牧師也有牧師的期望.
快50歲了.
50歲的牧師.
應該是無憂一千人教會才可以.
有這樣的期望.
應該要擁有兩座堂.
但你又覺得很焦慮.
眼看著時間到了.
所以每一次到了新年.
我們有什麼時間的焦慮.
所以在2025年的時候.
在網上見到有一種.
大家知道有些人喜歡原富.
炫耀你的財富.
包包,LV.
是一個最新的.
炫富的方式.
叫做炫耀你的時間財富.
有沒有人夠膽在網上.
說我有大把的時間.
我用之不盡的時間.
沒有人夠膽炫耀.
因為一說到時間.
大家都覺得貧窮.
我們時間是不夠的.
隨著一年一年的過去.
我們對於時間的焦慮會越來越嚴重.
所以時間.
是否經歷了基督的豐盛.
是重要的.
主耶穌基督在約翰福音第一個神蹟.
彰顯了他對時間的掌控.
水變酒.
水要變酒.
是要經歷時間的.
在聖誕節的時候.

$^{161}$有個小組來我家.
那個小組很喜歡喝酒.
帶了一瓶.
他們說很好喝的日本酒.
這瓶酒十二年.
剛剛好.
上等的好酒.
他們喝得很開心.
好的酒需要時間的醞釀.
唯有主耶穌.
可以瞬間將水.
變成最好的酒.
在舊約的聖經.
有希西佳的故事.
希西佳年輕的時候很厲害.
帶領整個以色列.
背起一個偶像的敬拜.
回到神身邊.
但年老的時候得了一個病.
他很痛苦.
時間不夠.
他向神求.
求神可以給我多一點時間.
二賽亞幫他求.
神就給二賽亞.
多給他十五年.
希西佳不太相信.
憑什麼神可以給我十五年.
二賽亞說.
你不相信可以向神求一個訊號.
求什麼訊號呢.
神就給他一個訊號.
你想求.
太陽的位置.
向這邊移.
意思就是.
以前的人.
你想我鐘突然.
從兩點變成三點.
還是從兩點變回一點.

$^{201}$所以神給他一個訊號.
告訴他一個很重要的真理.
就是誰在掌控時間.
是神自己掌控時間.
如果神說.
要多給你十五年.
就是十五年.
所以希西佳和主耶穌基督.
在第一個神蹟.
都彰顯了其中一種.
很有形有體的風聲.
是我們都在求的風聲.
是什麼.
時間的風聲.
我們都希望新的一年.
我們不會再陷入一種時間的焦慮.
覺得時間不夠.
我們要覺得時間是夠的.
神祖.
我們是為著永恆去做.
我們在地球這短短的幾十年.
只有一個.
如果你喜歡打遊戲.
知道什麼叫新手村.
新手村就是你遊戲的第一關.
讓你熟悉一下這個遊戲.
熟悉完畢業.
出了新手村.
Explore the whole world.
才是真正神給你的.
所以將來有神的國來臨.
那個才是你出新手村.
真正進入你真正風聲生活的那一天.
現在幾十年.
在地球你是在學什麼.
學怎樣去Run這件事.
怎樣去信靠神.
怎樣讓自己變成一個不一樣的人.
當神覺得你Ready的時候.
你就是你畢業的時候.

$^{241}$所以這個時間是夠用的.
是神Set好的.
新手村的時間一向都有限.
但神做我們是無限的.
是神所做.
所以這個觀念來到主耶穌那裡.
那個時間是很不同.
所以新的一年.
在去年進入新的一年的時候.
我都希望給自己一個Expectation.
新的一年我希望可以放下.
一些我對時間的焦慮.
嘗試一個不同的方式去處理.
我怎樣看自己的時間.
為自己定了一個目標.
就是2026.
我要開始一種比較慢活一點.
我要改變一下我的節奏.
節奏的節奏再慢一點.
不怕沒有時間.
這樣的節奏去過生活.
所以在2025我計劃2026的時候.
我就迫自己.
說今年.
我一定要每一個月.
讓自己有一個主日.
是不用講道.
是可以純休息.
所以今年我編了.
其實我原先是編了我有12個星期不用講道.
今年.
我開頭很開心.
希望今年可以慢活.
誰知道後來我發覺高興得太早.
隨著時間的轉移.
我發覺我那12個星期.
一個一個被人拿走.
這個打電話.
不行你幫幫忙.
從12個變成11個.

$^{281}$11個變成10個.
10個變成9個.
9個變成8個.
8個變成7個.
現在最新update是剩下6個.
我希望可以守住6個.
起碼可以讓自己有休息的時間.
所以我不知道新的開始.
新的挑戰.
你如何計劃.
你接下來一年的時間表.
你的節奏如何計劃.
我希望你在風盛裡計劃.
不是陷入一種.
時間焦慮裡.
去亂計劃.
因為時間焦慮你會亂計劃.
你會塞滿它.
萬念時間不夠.
搞得不好.
跟自己說.
在基督裡你有足夠的時間.
完聖所有的事.
在基督裡的時間觀.
是跟你想的不同.
只要神的時間到.
你發覺所有事都是.
水到渠成.
是真的水到渠成.
在基督以外你可以花20年.
去打拼一個目標.
在基督裡可以一個月搞定.
為什麼水變走.
是質的變呢.
為什麼十幾年才釀成的.
好酒主耶穌基督說.
我一句話就可以了.
因為神是掌控時間.
只不過我離開了.
基督的時候陷入.

$^{321}$時間的焦慮.
我們就用我們想到的.
時間觀去衡量.
我們這一生應該怎樣過.
其實不是你想像中那麼痛苦.
你回到基督的恩典.
裡面和它的風盛.
你看你的時間表.
所有的事都會不同.
所有的事你會發覺.
水到渠成.
水到渠成是最舒服的.
是最經歷到恩典.
所以希望大家可以.
新的挑戰不要陷入焦慮.
不夠缺乏貧窮的當中.
去計劃你的未來.
回到基督裡面.
知道它的恩典.
它的風盛包圍著你.
用這個去計劃.
下一個我們說說.
面對新的挑戰.
恩典怎樣幫助你.
從它的風盛當中.
我們領受最大的是恩典.
而且是恩上加恩.
我很喜歡這個字.
不是只是恩典.
而是恩上加恩.
我很喜歡有些神學家.
寫楊枯的解讀.
這個恩上加恩的意思.
就是一層一層的恩典.
疊加上去.
2025你經歷了很多恩典.
疊加上去.
所以不是一種的恩典.
它指的不只是.
我們都知道我們信主是恩典.

$^{361}$它指的不是one time grace.
就是你信了主你得救.
它是指daily grace.
你每一天所需要的.
新的恩典不停地疊加給你.
這些恩典可以幫助我們.
第一個處理過去的失敗.
創傷.
花一點點時間.
不要帶著一些baggage.
去你的future.
尤其是新的一年.
有些事該放下就要放下.
有些事沒處理過就要花一點時間去處理.
處理了之後.
你比較輕鬆一些走新的路.
所以新的挑戰.
我們希望可以處理一些過去的傷痕.
和一些創傷.
還有恩典可以幫你面對今天的挑戰.
恩典也幫你面對.
未來的未知和焦慮.
我經常和小祖聊天的時候.
我常常聽到.
最大的焦慮就是.
基本上我身邊所有人都很害怕.
今年我們會被AI取締.
有沒有這樣的焦慮.
你不再重要.
你現在做的事.
你的崗位可能都會沒有了.
我最近看了一篇新的文章.
AI有沒有可能取代牧師呢.
答案是有可能的.
所以是很焦慮的.
有一天你是不再被需要.
回到基督的風格.
去處理這種焦慮.
這種對於未知.
覺得自己.

$^{401}$很害怕自己會被取代.
會沒有價值.
會不被神愛的焦慮.
第二個我想.
簡單說說.
我都知道在.
你們印的衣服裡.
都寫得很清楚.
Living Grace.
你們的路線是.
Being Before Doing.
這個是重要的.
我們要在恩典裡去服事.
我們看主耶穌.
給我們一個好的示範.
在瑪麗福音3:16.
主耶穌受了浸.
就從水上來.
天忽然為他開了.
他就看見神的靈.
彷彿甲子降下.
落在他身上.
從天上的聲音說.
這是我的愛子.
因為這是一個很重要的點.
就是主耶穌.
在他還沒開始.
出來服事之前.
先發生了這件事.
然後才開始.
他的服事.
當主耶穌受浸.
他聽到父神跟他說.
你是我的愛子.
是我所喜悅的話的時候.
主耶穌是一件服事都沒有做.
他只是一個木匠的兒子.
他沒有任何的公共事務.
沒有參與過任何的服事.
沒有施行過任何的神蹟.

$^{441}$所以在這件事.
在我們面對新的服事的時候.
你一定要讓這件事.
這個身份.
成為你核心.
守著的東西.
如果你守不住這件事.
接下來會發生什麼事呢.
到了瑪麗福音第四章.
很清楚.
一開始.
先要在曠野.
經歷魔鬼撒旦的試煉.
到了第四章.
聖靈就帶主耶穌.
去曠野禁食40天.
然後魔鬼出來試煉.
而魔鬼的試煉.
從頭到尾.
都是target這個身份.
他跟主耶穌說.
如果你真是神的愛子.
你做這件事.
證明給我看你是神的愛子.
我們是會中招的.
我們是會.
透過服事.
證明一些東西給神看.
我好像真是你所愛.
我們很容易.
最容易.
我們用服事的果效.
表面看到的果效.
來convince自己.
好像還有價值.
好像真是神你所愛的child.
我是神所愛的兒子.
我是神所愛的兒女.
但是你發覺.
我們是倒轉了.

$^{481}$主耶穌是先知道.
我是神的愛子.
然後開始他的服事.
我們是相反.
透過服事來證明.
原來我真是神的愛子.
就是中了撒旦魔鬼給我們的temptation.
所以.
這件事是重要的.
如果你不知道在沒做任何事之前.
你的whole being.
你真正的身份.
就是神.
用主耶穌在十誡上的.
恩典換來的.
愛子.
是他所喜悅的.
before you do anything.
你已經是他喜悅的.
你守不住這個身份.
我們一去到曠野.
經歷的時候.
你就發覺你全部中招.
那個聲音會出來.
你憑什麼覺得神喜歡你.
你看你.
有沒有參與那一刻.
你看你有沒有參與服事.
很多這些聲音.
然後你就開始有個desire.
證明一下自己.
原來我真是需要一個神蹟.
很多時候發生你會求神蹟.
當神蹟沒有來的時候.
你會灰心喪志.
灰心喪志不是因為你很愛那個人.
我幫你求了神蹟.
為什麼神不聽.
重點是因為神沒有聽你求的神蹟.
你覺得你不再是.

$^{521}$神所喜悅的.
這個才是.
我們最深的恐懼.
這個恐懼.
是這一生.
都會伴隨著你.
不斷的攻擊你.
所以撒旦魔鬼.
每一個攻擊都是針對你這個身份.
所以如果你一放棄.
你又覺得你有個desire.
想要證明這件事.
那你就中招了.
你一定要這個身份在哪裡守住.
基督裡面很簡單.
我是主耶穌基督.
所揀選所救贖的.
我在祂的恩典裡面.
我就是祂的喜悅.
純粹是祂揀選.
不是因為我本身有什麼.
值得的東西.
這件事你守得住.
你就發覺你的福事.
是完全不同的.
心境去面對.
新的挑戰.
新的福事.
是完全是另一種畫面.
你看到的東西都不同.
你經歷的東西也會不同.
所以很重要學主耶穌.
不要陷入撒旦魔鬼騙你.
最後一個就是.
我希望在敬拜新的一年.
你都可以在基督耶穌的.
恩典裡面去敬拜.
約翰芬說了一句敬拜.
很好的定義.
時間將到如今就時了.

$^{561}$真正的敬拜是怎樣呢.
他就是用心靈.
這裡說的心靈.
我用英文翻譯成原文.
其實應該是在靈裡面.
然後和真理.
真理和靈.
然後拜祂.
注意這個S不是大寫.
是小寫.
是指你的靈.
在你的靈裡面用真理去敬拜祂.
因為富有者容易被人拜祂.
神是一個靈.
所以拜祂必須在靈裡面.
然後用真理去拜祂.
很多時候我們明白.
用真理去拜祂是簡單的.
你認識真理.
你知道真神是誰.
所以你去拜祂.
但在靈裡面就不容易解釋.
因為很多時候我們是真理停留在這個位置.
然後我們就用這個位置去敬拜.
所以很多時候敬拜.
進不了一個你覺得很.
很解脫或者很復活的狀態.
因為我們從來都不是在.
真理沒有進到你的靈裡面去敬拜.
而真理進到靈的位置.
是需要一個過程.
這個過程不是一天到位.
而是你要開始.
去練習這個過程.
這個過程是什麼呢.
大家記不記得.
主耶穌這段話是在什麼情況下講的.
他是和撒瑪利亞婦人.
聊天的時候講.
他和撒瑪利亞婦人聊天的時候.

$^{601}$他講了一樣很特別的東西.
就是主耶穌表現了.
一種是正常的人.
看起來是一個EQ很低的表情.
正常的人看起來是一個EQ很低的表情.
就是他專門.
篤了撒瑪利亞一個很不想提的東西.
他和撒瑪利亞婦人說.
你叫你老公來找我.
撒瑪利亞婦人說.
我沒有老公.
他說你說得沒錯.
你已經有五個老公.
現在那個都不是你老公.
大家知道對撒瑪利亞婦人來說.
這個是他最不想面對的一條傷疤.
是他最敏感的地方.
他為什麼中午12點.
這麼曬.
避開所有人去打水.
就是因為在2000年前.
經歷過五段挫敗的婚姻.
你整村人都出了名.
知道不受歡迎.
他知道避開所有人的閒言閒語.
一個人走去打水.
他最不想提的東西.
但是主耶穌一句話就提了出來.
他還要避開.
避開主耶穌都不放過.
我知道.
以前那五個都不是你老公.
我們都不是.
這是他最不想提的東西.
我最不想提的東西.
所以你要問.
為什麼要讓記載這件事呢.
如果你要看整段故事.
如果你刪除了那幾句.
都可以.

$^{641}$主耶穌見到撒瑪利亞婦人.
已經跟她說.
如果你喝我給你的水.
就是活泉不會渴.
跟著跟她說.
這種是真正的敬拜.
是活水.
為什麼要加那句.
主耶穌要點她的那句.
你會發覺很特別的地方.
如果主耶穌沒有點出那句.
你就不會明白.
什麼叫在靈裡面.
用真理去敬拜.
怎樣進入到靈裡面.
去真理敬拜.
什麼為之心靈誠實.
就是你必須要透過你生命中.
最黑暗.
最不想面對最真相的那一點.
順著那一點.
進入到靈裡面.
然後真理.
開始進入那個位.
真理進不到那個位.
是因為我們將所有我們不想面對.
最脆弱的位壓進了我們的潛意識裡面.
你關了門.
然後你停在.
你有意識的那幾個百分比裡面.
然後你想真理停在這裡去敬拜.
永遠都不行.
你要打開這道門.
隨著那個你最不想面對的位.
一個很深的地方.
真理必須進來.
主耶穌到你的愛接納.
必須進到你最脆弱的位.
所以主耶穌說完這句.
撒瑪利亞夫人就聽得明白.

$^{681}$於是你看到她的表現.
告訴你她聽得明白.
她馬上跑回村莊見所有人.
見人就說.
她知道我所有東西.
她知道我曾經有五段失敗的婚姻.
你不是之前不敢提這件事.
現在衝進來每個人見人就說.
她突破了.
她知道什麼叫最釋放的敬拜.
就是你不再掩飾你最差的那一點.
你最脆弱的那一點.
然後你從那一點進去.
主耶穌對你接納的那份真理.
恩典的真理.
進到你最深的那個位.
然後就在那裡敬拜.
這個正神想給你的敬拜.
我在兩個星期前就發生一件事.
大家看到我今天走路都不是很厲害.
因為最近.
我一月底就去了一趟.
我成長的地方.
梁牧就哄我去.
因為他們之前有一個宣教士在比利時.
我讀中學的時候在比利時.
我已經離開了25年沒有回去.
梁牧就哄我可不可以陪他去一趟.
他哄了我去.
我就哄Shirley陪我去.
我就說.
你想不想看我成長的地方.
Shirley肯之後.
我就哄Allen和Mary陪我一起去.
在我去比利時之前.
我的腳有痛風.
痛風痛了幾天就好了.
好了之後我覺得沒事.
我就去比利時.
我非常放肆地吃.

$^{721}$因為很久沒有吃那邊的東西.
吃完之後.
回來的時候就停在LA三天.
在LA三天.
什麼都沒有做.
全程放肆地吃.
吃完回來就發覺.
他復發中招.
因為太放肆.
中招之後我又很乖.
馬上禁食.
禁食之後他又康復了.
康復之後又到了週末.
週末他又回來.
教一些人吃東西.
然後又開始放肆.
放肆之後他又復發.
復發之後我又開始禁食.
禁食之後他又康復.
康復之後剛好又遇到過年.
大家知道過年發生什麼事.
又放肆.
難得過年.
我記得那天星期六晚.
我們有新年準備.
訂了很多甜食.
我吃了很多甜食.
我星期六晚完了崇拜.
星期日有一個小組要來我家.
我接了一個小組.
然後星期一又有另外一個小組來.
所以我星期日星期一都要接.
有兩組人.
而且都幾大組.
但我星期六晚吃得很厲害.
我回家星期日起床就發覺.
大件事我又開始禮了.
昨天不節制.
今天就禮了.
Shirley一早就走去.

$^{761}$買海鮮.
準備煮東西給大家吃.
我起床就發覺.
好像腳有點痛.
今天怎麼辦.
剛好我吃了安眠藥.
其實剛發的時候.
如果我吃安眠藥.
至少可以幫我解痛.
對我有效.
Shirley回來的時候.
我跟她說.
我昨天吃多了.
我們家對面的.
目標幫我買了一張Advil.
然後她就.
開始罵我一輪.
誰叫你不節制.
罵完一輪就說幫你買不理你.
她又不幫我買.
然後她就說那群人快來了.
我沒時間理你.
然後她下去處理所有的東西.
然後我一個人就留在床上.
在二樓.
我就陷入一種很委屈的狀態.
然後我就說這是黃金的機會.
這是我最脆弱的時刻.
我就邀請主耶穌.
來談談.
我想你幫我一起去處理.
這種很委屈.
很憂傷.
一旦有事.
身邊的人幫我去對面買一張藥都不肯.
我被人放棄.
不值得被愛.
沒有價值.
這種所謂的情緒.
我想趁這個機會.

$^{801}$這是我最脆弱的一面.
主耶穌黃金機會.
快點進來處理.
我邀請了聖靈.
主耶穌和我一起進去處理.
我大概就.
當然主耶穌就和我一起去處理.
處理了大概十分鐘之後.
我就發覺.
突然間在靈裡面.
我是第一次聽到.
聖靈很清楚的聲音.
很美的聲音.
我沒聽過這麼美的聲音.
在你靈裡面有一個很美的聲音.
我就想到一首詩歌.
那首詩歌就是.
這首詩歌是我.
近十年.
都沒聽過Playsteam.
選來唱的詩歌.
我已經完全忘記了這首詩歌.
但那天我處理了.
最脆弱的那個位置.
就是這首詩歌.
一個很美的聲音.
在裡面湧出來.
那一刻我終於明白.
什麼叫敬拜.
是真的.
進到你靈裡面.
開始湧出來的敬拜.
聖靈有一個工作.
羅馬書第八章.
聖靈的工作是將.
父神的愛.
湧入你的靈裡面.
在那天.
我很喜歡這個詞.
聖靈將父神的愛.

$^{841}$在你的靈裡面.
這樣的感覺真的灌出來.
那種敬拜.
完全是.
委屈的心情.
完全煙消雲散.
隨之而來.
是你被神的愛.
包圍著.
你整個人.
非常有力量.
我就立刻起來.
腳仍然是痛的.
我就按下去.
然後那群人來的時候.
他們看到我的狀態.
是一個聖靈.
充滿愛的狀態.
去招呼每一個人.
那天應該是.
Pastor Alan和Mary.
和我聊到一兩點.
才離開.
整個群組接待完.
星期一那個群組又搞了很久.
兩個群組順利招呼完.
那兩個群組的人.
沒有一個看得出我們之前一地雞毛.
他們看到的是.
我一個充滿愛心的聖靈.
充滿耐性.
充滿能量的狀態.
那個狀態是哪裡來的.
就是之前你.
被主耶穌提到一些.
你很不想面對的事.
然後讓主耶穌進到那個位置.
真的進到你的靈裡面.
然後那種敬拜湧出來.
就是一種最真實的敬拜.

$^{881}$在恩典裡面的敬拜.
就是將主耶穌.
身上加恩的恩典.
不是停留在這裡.
是停留在你靈裡面.
最深的那個位置.
怎樣進到最深的那個位置.
就是每一次.
有些人和事發生.
會引發你最受傷.
你最不想碰那個位置.
那個就是你的黃金的機會.
所以任何時候.
神藉著你身邊的人和事.
開始引發你那個位置出來的時候.
不是神要罰你.
也不是旁邊的人.
故意想騷擾你.
是因為神想給你一個黃金的機會.
是神邀請你.
可不可以讓我的真理.
開始進到你的那個位置.
然後在那裡.
我與你見面.
然後在那裡有一個很真實的敬拜.
所以新的一年.
希望我們的敬拜從這裡開始.
進深一點.
變成一個不一樣的敬拜.
你的投入.
變得不一樣.
我們一起來祈禱.
早一地來到你面前.
我們謝謝你.
你的豐盛帶來了恩上加恩的恩典.
是源源不絕的恩典.
在我們面對新的一年.
新的挑戰.
新的開始.
我們需要從這種豐盛的恩典.

$^{921}$裡面開始.
幫助我們面對新的一年.
我們有很多的焦慮.
新的挑戰我們有很多的焦慮.
我們有很多的恐懼.
幫助我們回到基督裡面.
去領受這一份豐盛.
這一份恩上加恩的恩典.
開始有能力去面對.
幫助我們.
我們的福事在恩典裡面.
我們的敬拜都回到恩典裡面.
去面對我們最脆弱.
最敏感的位置.
然後我們有一個最健康最活潑的靈.
去敬拜你.
求主你親自帶領我們.
今年我們的敬拜福事.
都上升到一個不同的等級.
我們今天的祈禱感恩.
奉主耶穌基督的名求.
Amen Amen.
\newpage



\section{腓立比書 3:18-19}
\label{sec:031781}
\textbf{重塑生命系列 (6) 貪食與節制——當慾望失去界線}~~~~張健庭牧師博士
\newline
\newline
連結: \href{http://www.church.com.hk/acms/content.asp?site=cdc\&op=show\&type=product\&code=031781\&layout=sermon}{\texttt{031781}} ~~~~ 語音日期: 2026-03-08
\newline
\newline
\hyperref[sec:031780]{\small{< < < PREV SERMON < < <}}
~
\hyperref[sec:index_chronic]{\small{[返順時目]}}
~
\hyperref[sec:031782]{\small{> > > NEXT SERMON > > >}}
\newline
\newline
腓立比書 3:18-19
\newline
\begin{longtable}{cl}
\hline
\hline
章節 & 經文 (和合本修訂版)\\
\hline
3:18 & \begin{tabularx}{0.7\textwidth}{X} 因為，我屢次告訴你們，現在又流淚告訴你們：許多人行事是基督十字架的仇敵。 \end{tabularx} \\ \\ \relax
3:19 & \begin{tabularx}{0.7\textwidth}{X} 他們的結局就是滅亡。他們的神明是自己的肚腹；他們以自己的羞辱為光榮，專以地上的事為念。 \end{tabularx} \\ \\
[1ex]
\hline
\hline
\end{longtable}
$^{1}$好,弟兄姊妹平安.
我們一齊先低頭祈禱.
特別為這個世界來禱告.
天父啊,你是這個世界的主.
你在天上面坐著為王.
我們想起詩篇第二篇那裡.
講到萬民都是憎鬧.
我們都是謀算虛妄的事.
但是你在天上面來去癡笑.
因為最終你的受高者都會掌權.
我們看到現在這個混亂的世界.
人與人之間的仇恨也都有爭戰.
邪惡的勢力好像都是不斷地來去擴張的時候.
天父啊,有時候屬你的子民活在世上.
有很多我們所見到的事情.
我們見到不合神理的心意.
但是我們真的承認.
我們自己個人能夠所做的都是很少.
但是我們為著這個世界來去禱告.
因為我們知道.
無論這個世界的各樣的困難.
各樣的不是合神心意的.
最終你自己的旨意都會達成.
因為你也都答應過在最終的有一天.
在天下萬物都是在基督裡面同歸於一.
在你的救贖的旨意當中.
沒有一件事是例外的.
沒有一件事是在你救贖的旨意之外的.
願你的旨意成就.
也都願我們在這個世界的時候.
我們為著這個世界來去求平安.
我們知道最終我們的禱告.
在神面前必然會蒙英韻.
因為我們就是按照基督的心意而求.
願你保守.
很多現在在美國當中也都有很多的.
出去打仗的士兵 很多的軍人.
也都保守在戰亂的地方.
在伊朗這個國家也都有很多的平民百姓.
我們求你保守他們每一個.

$^{41}$願意無辜的人的傷亡.
都是減到最低.
我們願你的旨意成就.
我們為著我們今天的崇拜.
交給你的手中.
我們就是將自己的路放在你面前的時候.
讓我們看到你在這個世界的心意.
也都讓你看到你在我們生命當中的心意.
我們這樣禱告奉耶穌基督的名祈求.
阿們.
不知道你知不知道.
現在其實你們每一個人所享受的.
比古代的皇帝更加多.
就是你可能完全都沒有想過.
但是其實你所每日的享受.
真的是以前皇帝也都未必有.
你想一下你熱的時候怎樣.
有冷氣.
想吃冰淇淋的時候怎樣.
一開冰箱就有了.
那些食物.
很多是很新鮮的.
差不多是即日運到的.
就算不是即日.
都是兩三天之內運到的.
你想一下.
其實只是這幾樣東西.
皇帝都沒有你這樣的享受.
接著你還可以怎樣.
你如果要去哪裡呢.
坐飛機 坐車.
以前的人就算皇帝要微服出巡.
他坐的那些馬車.
我們廣東話說.
「凳到肺都甩埋」.
即是那些馬車沒有避震的.
你想一下那些馬車.
只有那幾個木鹿在那裡推.
你見到很威風.
你坐在上面那裡.

$^{81}$蹬下蹬下.
可能又熱又性.
繞著帳篷蹬.
蹬兩日都未必去到.
他要微服出巡的地方.
但是現在你坐飛機.
全世界到處都去到.
這些只是一些很小的事情.
我們現在的生活水平.
享受是很多人夢寐以求.
皇帝都想像不到的享受.
你想一下以前楊貴妃吃荔枝怎樣.
嘩 真是厲害.
勞師動眾.
不知道楊貴妃想吃多少粒新鮮的荔枝.
想運到來不用變黑的荔枝.
都要搞到這麼厲害.
成為了中國歷史當中一個.
不斷流傳的一個故事.
現在你吃荔枝.
嘩 不單是夏天吃.
冬天都可以吃荔枝.
很多事情是想像不到的.
當然在這個世界.
有很多地方有很多國家都很缺乏的.
都是的.
不過今天我們的講道.
當然是講回我們現在的處境.
在美國這個地方.
所以在美國這個地方本身.
其實我們是不愁缺乏的.
我們是豐富 豐富到過盛.
以致到所以今天所講的.
貪食這個罪本身.
就好像聖經所講.
有時候我們好像將我們的滿足.
我們的肚腹.
在肥辣秘書那裡講.
他們的神就是我們自己的肚腹.
即是我們廣東話都有.

$^{121}$我們肚餓的時候我們叫做什麼.
叫做什麼.
祭五臟廟.
你有沒有聽過.
當然這個都是以前人那些說話的方式.
現在你說祭五臟廟不代表你拜神的.
但其實某個程度都是一樣.
你為什麼要祭五臟廟.
就是說我們的肚腹就是我們的神.
是不是.
事實上我們現在就是活在一個.
這樣的過渡的一個文化當中.
很多東西我們是不愁缺.
我們是過渡.
而在一個這樣的過渡文化當中.
是我們飲食方面過渡.
我剛剛幾天之前.
有看過一個很短的一個小小的紀錄片.
它講到都很有趣的.
就是說原來在七八十年代.
美國很多人都不是黏肥的.
很多人都是正正常常.
不是肥仔肥女很大的身形.
是在七十年代尾八十年代開始.
整個美國現在那個黏肥的很大的嚴重的問題.
其中一個原因.
當然有食物的因素的.
就是因為那個時代後來其中有一個的政策.
是美國的農業的政策.
就是政府補貼了很多那些種玉米的人.
然後就製造了很多Corn syrup出來.
就是那些玉米的甜.
那些Corn syrup就差不多便宜到不用錢.
因為政府補貼他種玉米的.
就變成大量的湧入去各樣的食物製加工那裡.
變成很多東西就很甜.
變成吃甜了那些人就肥了.
這個是其中一個原因.
這個原因就是整個社會政策的因素.
但是其中他也都說了另外一樣東西.

$^{161}$就是說.
以前在七八十年代之前.
是很固定人是一天吃三餐的.
有時候都沒有三餐的.
有時候是兩餐的.
那些小孩子都是回家的時候.
那些父母都不會給他零食吃的.
等一下就會吃飯了.
我們中國人也有這樣.
等一下開飯不要吃飽了.
等一下要開飯.
以前是這樣的.
但是現在就是怎樣?.
不停的.
就是不斷的零食各樣的東西.
就不斷的來的.
而且那些食物的份量是多了大了很多.
以前那些汽水是12盎司而已.
現在那些16,18,24都有.
是不是?.
這個是太多了.
而我們就是活在一個很過度的一個.
就是太多不受控的,多得太多的一個時代.
食物是其中一樣.
娛樂也是.
以前沒有那麼多娛樂的.
很多時候,我記得我以前小孩的時候.
很多時候就不會跟父母或者不會跟大人說.
很悶的.
為什麼?.
你有沒有跟父母說過.
小孩都會說很悶,會不會?.
不會的,為什麼?.
因為如果你跟父母說很悶.
下一句說話會是怎樣呢?.
你幫我做這件事.
幫我洗碗.
就是這樣的.
所以你小孩你最悶.
你自己找些東西來做,找些東西來搞.

$^{201}$你不要跟父母說你很悶.
不然就叫他幫我去街上買一瓶醬油回來.
我們那個時代是這樣的.
我們現在沒有那麼多過度的娛樂.
我們反而是在生活當中找些東西來充實自己.
但現在就不是了.
現在的娛樂是多得過度.
你知不知道.
譬如你家裡有Netflix.
Netflix就是那些串流平台看電影的.
你知不知道Netflix的電影多得.
基本上.
我不是說你看.
你只是搜索出來都搜索不完.
譬如Netflix在它的資料庫儲存.
譬如它有幾十萬套電影.
基本上你搜索都搜索不完.
你就算打一些關鍵字.
它出了幾十套給你.
其實它背後可能同一個關鍵字.
可能背後還有二三百套.
它都不會展示給你.
因為你太多了.
我不是說你看.
我只是說你搜索.
都搜索不完.
所以其實就是活在一個.
不斷有太多娛樂的東西.
令到你覺得一停下來.
以前你有空閒的時間你會怎樣.
你會怎樣.
看書.
有些人說做夢也是一件好事.
可以思考一下.
讓你的思想有些.
其實你想一些東西.
都是在整理一些思緒.
現在怎樣.
你會發覺沒有這些時間.
一有時間就滾手機.

$^{241}$滾手機我之前可能說過.
但我最近再看過有一個.
是真實做過一個研究.
其實原來一滾手機.
因為它是英文來的.
會brainwalk.
即是說真的會.
腦袋會遲鈍.
簡單來說就是變蠢.
真的原來滾手機會令你變蠢.
這個是做研究的人.
他不是做社會學研究.
他是那個Neuroscientist來的.
是那些腦內科的神經醫生來的.
因為他scan那些人的腦.
即是他研究那些腦的結構.
因為現代人這樣來不斷刺激.
只是同一個很小的腦的位置.
是會令到其他的位置會退步的.
他說對於成年人還好.
退步到一個地步.
如果成年人能夠戒得甩.
都有機會可以恢復.
即是可以回頭.
但是對於那些發育未健全的teenager.
或者是二十五歲之前.
在二十五歲之前的人.
這樣腦發育是irreversible.
是回不了頭的.
因為已經被它破壞了.
所以很多這一代的小朋友.
基本上那個人是很擔心的.
他應該是UCLA的腦外科研究員來的.
所以他很擔心.
他說整個社會現在是這樣的時候.
根本是等不及.
說再等五年之後才立法.
根本等不及.
五年之後一個generation就沒有了.
他說現在已經沒有了第一個generation.

$^{281}$他說短視頻這個發明.
已經沒有了第一個generation的小朋友.
他說再等五年又沒有第二個generation.
他說根本已經等不及.
現在是一個過渡的娛樂.
就是過渡的消費.
過渡的消費也是一樣.
因為我們有多餘的.
美國就是出了名的一個消費社會.
美國是.
所以Trump那些關稅戰.
為什麼可以打得成呢.
因為全世界都賣東西來美國.
如果你是買東西.
不斷地買東西的時候.
全世界的人都賣東西給你.
所以他才有籌碼跟人打關稅戰.
就是用這個來跟人談判.
如果你買很少東西.
就等於你去街市買菜一樣.
你買了雞碎那麼多.
又跟人講價人家鬼踩你嗎.
是不是.
但是如果你大核就不同了.
什麼都可以買的.
美國就是一個這樣的消費大國.
就是我們不斷地消費.
不斷地消費.
我有一次跟一個朋友聊天.
我們那時候大家都是來讀書的時候.
我在神學院讀書.
他就在別的大學讀書.
他就說一件事.
他說有時候他的同學很蠢.
花錢花得很蠢.
他沒有想過那樣東西值不值得.
然後有錢就花去用信用卡買了之後.
然後又不夠錢又要分期付.
很蠢的.
他說一件事我從來沒有想過.

$^{321}$他說其實你沒有想過.
其實是專登的.
專登不是那個人專登.
他說是美國社會專登的.
因為他專登不教那些小孩怎麼花錢.
然後就等他們花錢.
等他們花錢之後.
美國的經濟就會好.
因為美國的經濟是一個消費型的經濟.
所以這個想法又是真的.
學校從來都不會教小孩怎麼花錢.
怎麼省錢.
怎麼有金錢觀念.
沒有的.
一個正常的教育體系的時候.
這件事從來都不講.
我去到大學都不講.
基本上差不多.
你想想.
是哦.
他說的那件事我從來都沒有想過.
他說可能是美國整個國策.
專登搞那些人.
搞到這樣每個人都是花錢.
更加是過度的刺激.
這個和剛才講的有一點相關.
現在那些刺激.
不單是你玩那些刺激運動的刺激.
現在是那些電子刺激.
電子的刺激.
其實滾手機也是一個刺激.
你不斷地滾的時候.
它根本是不斷地.
讓你的腦子分那些多巴胺出來.
是刺激你.
就好像你不斷地.
有些人不斷地看電影.
不斷地吃薯片.
吃爆谷.
看足兩個小時.

$^{361}$吃足兩個小時.
是不是.
就是那些刺激它.
因為它一邊吃.
它已經變成機械性動作.
它根本停不了.
它是不斷地刺激著自己.
而滾手機其實.
它是想製造一個不能夠停的.
事實上這個是很真實的.
你有沒有發覺滾手機的時間.
和你做其他事的時間.
好像是兩個時鐘一樣.
你有沒有發覺.
你做其他事的時間.
或者舉例你看書.
或者看聖經的時間.
為什麼五分鐘十分鐘.
好像無限延長一樣.
這麼久都還沒完.
但滾手機.
一拿手機出來.
然後怎樣.
一看一看.
打算玩五分鐘而已.
現在有時間.
滾五分鐘.
一看.
滾了一個小時.
是這樣的.
為什麼.
因為它的設計.
就是不斷地刺激到你的腦.
來分泌一些多巴胺出來.
而這一種.
剛才已經說過了.
是會冰凍的.
是會令到你的人變蠢的.
其實這個是和你吸毒的人.
令到腦再不能夠有基本的運作.

$^{401}$沒有了是一樣的.
現在我們就是處於一個.
這樣的過渡的一個時間.
而在一個這樣的社會當中.
我們吸收我們身體過多的需要.
過度的講究的時候.
就去執著.
很享受執著自己的一些享受.
過度的滿足.
其中一種就是.
我們等不到.
我們很想很快即時地要滿足.
就是你的耐性越來越短.
上一次說憤怒.
其中一個憤怒的原因.
我覺得會不會是.
和人的耐性越來越短.
也都可能有關係.
現在這些東西.
很多東西都不能等了.
即時不給你.
你覺得.
怎麼會這樣.
現在什麼都即時有了.
而這個過於專注的時候.
我們就很容易去上癮.
很容易上癮.
很多東西會令到你上癮.
而這個世界設計很多東西.
是令到你會上癮的.
而這樣.
我們就明白到.
哥林多前書所說.
凡事我都可行.
但不都有益處.
凡事我都可行.
但無論那一件.
我總不受他的克制.
這個是變成了一個.
非常重要的提醒.

$^{441}$沒錯我們是什麼都可以做的.
但是現在的問題就是.
第一做的東西.
是不是有益.
這個是一個問題.
即是等於你什麼都可以吃的.
但吃的東西有沒有益.
是一個問題你要問.
第二個就是.
現在更大的問題就是.
我們能不能夠說.
我們做的東西不受他克制呢.
做的東西不受他克制呢.
這個是一個很大的問題.
而現在的特別說克制這件事.
現在的數碼的演算法.
其實是不斷地令到你被他控制的.
因為他知道怎樣去吸引你的注意力.
整個演算法的計算過程就是.
他要知道你喜歡什麼就不斷地給你.
這個大家都知道的.
你看過一個貓仔狗仔的視頻的時候.
很多貓仔狗仔的就不斷地給你.
因為他知道你喜歡看那些.
你喜歡什麼他就不斷地給你.
甚至乎你看Netflix都是一樣的.
你舉例來說看那一套.
這一類的科幻片.
他就不斷地.
同一類型的科幻片.
就不斷地發給你.
有時候我都覺得.
我不是想看這麼多這些的.
但他就不斷地發這些給你.
而事實上當這樣的時候.
他就要令到你的專注越來越狹窄.
他不斷地要取你的吸引力.
亦都是令到你所看的世界.
亦都是越來越小.
你想想譬如我看過一套科幻片.

$^{481}$我不想每次都看科幻片.
我想看文藝片可不可以.
我想看紀錄片可不可以.
他就不給那些給我.
他不斷地繼續看那些科幻片給我.
我的世界就越來越小.
我的視野就越來越少.
你YouTube都一樣.
你聽過一個時事新聞的評論.
譬如舉例來說是罵特朗普的.
他就不斷地把所有關於罵特朗普的東西給你.
你本來對於他都沒有什麼意見.
久而久之都覺得特朗普不死都沒有用.
這個壞人來的.
就是這樣.
而這樣的時候.
其實他是在刺激你的情緒.
很多時候.
最多人在我長大的一個社會.
或者氣氛的時候.
都有人教我們.
就是你不要那麼容易.
去給你的情緒來控制你.
你做有時候你要冷靜一點.
有時候你的東西你要心思熟慮一點.
你用你的理性.
用你的思想來分析.
不要只給情緒.
我們叫做情緒用事.
但是現在看我們的電話.
我們的資訊的時候.
甚至乎那些電視的主流的媒體.
報新聞那些都是一樣的.
他只是想有一個目的.
就是要刺激你的情緒.
那個新聞的主題.
也只是有一個目的.
是引起你覺得很大反應的.
報新聞都是要這樣報的.
所以很多假新聞出現的.

$^{521}$因為他根本就不是想要報一個事實給你聽.
不是要客觀分析給你看.
只是想要牽動你的情緒.
為什麼.
因為如果你的情緒受到控制的時候.
你就會被他拉著鼻子走.
如果你是很客觀的分析事物的時候.
有時候你說也未必是這樣.
你就未必完全聽他的.
但是如果他是主導到你的情緒的時候.
他想要帶你去哪裡.
你就帶他去哪裡.
我記得我也說過有一個真實的.
有一個訪問.
其實是娛樂圈有一個明星.
就和有一個主持人做一個訪問.
開玩笑地說.
但其實是真事.
那個娛樂圈的明星就說.
你知不知道有一段時間.
我相信地球是平的.
現代人還相信地球是平的.
怎麼搞的.
就算不讀書你也會的.
他說為什麼呢.
他說就是因為我有一次看了一個YouTube.
他說地球是平的.
說得有情有理.
然後我看完之後.
他之後不斷發那些地球是平的影片給他.
就是Conspiracy Theory.
我說什麼政府怎樣欺騙人諸如此類.
很多東西串在一起.
越看越覺得有情有理.
後來他就相信地球是平的.
當然後來他這樣說出來.
他就知道被人欺騙了.
後來就走出來.
但你想想這麼荒謬的事情.
也能欺騙人.

$^{561}$有些事情其實並不是這麼荒謬.
是逆真逆假的.
他就抓住你了.
而不去分析究竟事情.
這個就是現在.
因為在一個網絡的媒體的時候.
我們會在一個這樣的時代.
每個人看一個三分鐘五分鐘的短片.
就覺得自己是專家了.
三分鐘或者五分鐘解釋.
比你厲害一點.
十分鐘.
十分鐘解釋以色列和哈馬斯的恩怨情仇.
你就覺得成為以色列專家了.
哇 中東國家地緣政治.
哇 完全理解了.
Come on.
你看了十分鐘的短片.
就覺得自己是一個專家這麼厲害.
根本.
你就知道根本就不是.
但很多人的情緒就被這樣牽動了.
牽動了的時候.
你看世界很多事情其實是很扭曲的.
不只是你其實我也是一樣.
所以我們就是活在一個這樣的時代當中.
我們的資訊那樣東西.
其實整個演算法.
其實簡單說一句.
它就是要控制你.
令你上癮.
令你不能停下來.
一句話是很真實的.
不用錢的東西是最貴的.
是不是.
哇 你說不用錢的.
哇 我每天滾手機滾兩個小時.
你叫我去看一套戲我看兩個小時.
我都要給錢入場.
是不是.

$^{601}$哇 現在的電影票都不便宜.
我看手機不知道要滾兩個小時.
滾八個小時都可以.
都不收我錢.
那是不是很划算.
所以就說不用錢的東西是最貴.
為什麼它不用錢呢.
因為它是控制.
令到你不斷地使用它.
而無法自拔.
它已經主導了你各樣的.
你的想法.
你的世界.
你所去看的東西.
它得了這樣之後.
還貴過你給十多元入場看一套戲.
是不是.
不用錢的東西是最貴的.
而這個現在就是所有.
所以為什麼你停不下來.
你有沒有試過.
沒有了電話.
可能幾個小時或者半天.
覺得好像這半天或整天.
錯過了很多東西.
如果你沒有這樣的感覺.
我恭喜你.
你還不是很上癮.
我真的聽過有個大學生中學生.
他們說現在的手機.
我睡覺都要放在床頭.
我睡覺都放在床頭.
你睡了覺你都不會知道.
但心理上根本撕不掉.
所以去到哪裡第一件事.
吃東西當然是手機先吃.
看風景都是手機先看.
但是你知不知道原來.
我也看過一個研究.
你以為很多人都是這樣想的.

$^{641}$譬如你去旅遊.
那些人去旅遊就想著.
拍很多照回來做紀念.
很多人都有這樣的想法.
但原來我看過一個研究.
其實如果你以這樣的想法.
去旅遊的時候.
你只顧著拍照.
其實你是不記得看過的東西.
你外判了給相機.
然後上載了.
你就以為你記得了.
其實反而是你的腦是沒有記憶的.
因為你已經外判了給它.
那個不再屬於你.
所以你的記憶就要在Facebook那裡.
五年前你做過什麼.
三年前見過什麼朋友.
如果它沒有提醒你.
你和它那種很歡樂的時光.
好像都現在要依賴給它.
而事實上.
我們就會是一個這樣的演算法.
來嘗試將你的人生.
好像它就成為了你的主導.
你的人生再不是你自己主導.
是它主導你.
所以這是一個危機.
所以這個貪食的危機本身.
不是說只是吃到很胖為止貪食.
而是說有些東西已經過量到一個地步.
我們根本就無法抽身出來控制.
它變成了我們的神.
它變成了我們的神的時候.
所以這個危機帶來了.
對於我們自己來說.
有什麼影響.
有什麼轉變.
第一我們控制自己的能力是越來越少的.
控制自己的能力是越來越少的.

$^{681}$事實上你發覺很多東西.
你不由自己控制.
譬如你想我控制我自己.
這個時間我想做某樣東西.
譬如舉例來說.
這個時間我想和朋友聊天.
或者打電話去談心.
或者約朋友出來.
但是你要發覺控制不了.
有時候大家的生活節奏.
似乎是被手機主導.
聊電話聊著聊著.
嘟嘟有個訊息來.
不聊了有其他東西看.
諸如此類.
我們對於自己人生的掌控是越來越少.
在一個表面上是很多工具給我們用的世界.
但是不是我們用的工具.
現在是有些工具主導了我們自己的人生.
你想一下都是很暖濃的.
你開車是想著那個車是一個工具.
帶你去某一個地方.
但是現在就變成了.
你什麼時候上車上了車去哪裡.
現在由它來控制.
你想一下其實這個和綁架是沒什麼分別的.
是不是.
而事實上我們真的被綁架了.
你看什麼你做什麼.
甚至乎什麼潛移默化的影響了你喜歡什麼.
其實是完全被綁架了.
那個失去注意力是一個很大的問題.
失去注意力是一個很大的問題.
現在我也是之前.
也是兩三個禮拜前.
就說到那個應該是一個哥倫比亞大學的學校的報告.
你知道哥倫比亞大學是一個長春藤的名牌大學.
就是名校來的.
他就說那些學生沒有能力去讀書.
沒有能力去完成一本書.

$^{721}$那些教授教書已經再不能夠填一本書給他看了.
因為那些學生就是失去了那個專注力.
就是剛才所說的那一代已經沒有了.
他再沒有辦法坐下來完成一本書.
再沒有看書的能力是沒有了.
其實很多的專注能力都是令到我們失去的.
你想一下多難來去禱告呢.
就是因為很多時候我們的專注能力失去了.
當我們去禱告的時候.
因為你慣性給東西刺激你的腦.
剛才所說的brainwashed就是那個原因.
就是因為你的腦就好像吸毒那些人.
吸毒那些人你玩短視頻的時候.
其實是沒有資訊的.
他就不斷地刺激你 哇好笑啊好笑啊.
有隻貓仔在那 哇好可愛好可愛.
其實是沒有資訊的.
他就不斷地去刺激你.
當你一沒有了這種多巴胺在你的腦的時候.
那件事就做了不久.
三分鐘已經是受夠了.
因為你的短視頻三分鐘之後就要下一個了.
所以很多東西你是失去了專注能力的.
你看書當然不用說了.
祈禱不用說了.
讀經不用說了.
這些是完全沒有了的.
所以現在人的靈修生活.
或者現代人的屬靈操練這方面是一個很.
我們現在是面對著一個.
整個世界好像跟我們反對的一個時代.
當然我這樣說的時候不是危言聳聽.
因為我們真的活在一個這樣的時代當中.
如果你真的要讀經靈修的時候.
如果你是上了癮.
你首先真的要戒毒來開始的.
這個不是誇張的.
因為你是中了毒的.
你三分鐘的專注力之後就沒有了的時候.
你要慢慢慢慢戒毒出來.

$^{761}$而事實上其中一個.
等一下我會說一下我們需要怎樣.
所以久而久之我們對於屬靈的事情變得很遲鈍.
你的心靈是不敏銳的.
看很多事情是很表面的.
因為所有東西都是別人給你的資訊.
你自己沒有心靈的洞察力.
因為第一你根本沒有靜下來的空間去思考事情.
第二也是你跟神的事情的時候.
其實我們人是很特別的.
好像跟神相交很多時候是要沉澱在深一個層次的.
無論是尋求神等候神禱告.
你不是傻傻滾在表面上做了搞定交貨.
然後mark個tick.
你是需要沉澱.
但是你一直很多事情都沒有沉澱下去.
沉澱下去的時候就變得越來越遲鈍.
所以久而久之我們的生命的焦點都轉移了.
轉移去了哪裡呢?.
肯定就不在神那裡.
轉移去了自己那裡.
其實都不知道有多少轉移去了自己那裡.
即是你說我不去神那裡去回自己那裡.
但其實都不知道有多少轉移去了自己那裡.
因為根本已經變成了一個.
講得白一點.
就是那些zombie.
你知道現在這個zombie.
你知道什麼是zombie嗎?.
就是那些喪屍.
喪屍片你知道什麼時候流行嗎?.
原來以前是沒有喪屍的.
即是最低限度不流行.
是不是完全100\%沒有我就不敢講.
但最低限度不流行.
是在80年代之後才開始有一兩套喪屍片.
然後慢慢慢慢喪屍的類型就開始流行.
喪屍是一種什麼樣的生物呢?.
喪屍是一種.
它是仍然會動的.

$^{801}$仍然會吃東西的.
仍然會有反應的.
打它會有反應的.
但它只是反應來的.
它不會自己想事情的.
它一個群體一個喪屍去那裡.
全部喪屍就會群一起去的.
喪屍是一種這樣的一個存在來的.
為什麼這個喪屍片會在年輕人開始流行呢?.
其實背後又有一個文化的背景.
就是因為它發覺.
其實很多人覺得其實慢慢慢慢.
其實就是好像喪屍化.
我們就好像活在一個喪屍的社會當中.
而生命的焦點轉移.
喪屍就不是在自己那裡.
根本都不知道去了哪裡.
就是沒有了焦點.
其實去到一個地步已經沒有了焦點.
你看著手機就像喪屍一樣滾呀滾呀滾呀滾呀.
滾了幾個小時就有了焦點.
就是這樣的一個情況.
我們看回主題表演的時候.
主題表演是怎樣去面對慾望的呢?.
最開始我們看到主題表演.
它出來傳道的時候.
它就在四十週夜.
在禁食四十週夜在狂野那裡.
那裡當然有不同的神學的解讀.
主題表演為什麼一開始出來傳道的時候.
要去那裡.
怎樣代表四十週夜.
怎樣代表新的出埃及.
它的四十週夜.
或者怎樣對比摩西在西大山的四十週夜.
諸如此類.
有很多神學的解讀.
我這裡不是想講那些.
但其中一樣東西我覺得都是非常重要.
如果以人性來說.

$^{841}$主題表演它在出來傳道之前.
首先它要在神面前.
它不讓這個世界的東西來主導.
究竟它是一個怎樣的.
它整個人生整個的事奉.
不是讓這個世界主導.
它去到狂野那裡.
甚至去到其中一個試探.
魔鬼去試探它.
將石頭變成食物.
它其中一個慾望都是餓嘛.
吃東西嘛.
這個時候它是拒絕的.
而有時候適當的禁慾.
其實是需要的.
適當的禁慾是絕對需要的.
現在很多時候那些.
都講到好像很壞.
有任何禁慾的東西.
都好像很醜化.
好像講到中世紀那些.
說是將自己虐待.
來鞭策自己那些.
這個完全是扭曲了.
在有教會或者有我們的信仰.
任何一段的教會歷史.
都必然有禁慾的成分.
在那個屬靈操練裡面.
沒有一個時代的教會歷史.
的屬靈操練.
是沒有禁慾的元素在其中.
你聽得懂我說什麼嗎.
因為根本那個就是必須.
就等於你煮任何的菜.
必然會離不開有鹹味.
有甜味 有酸味 有辣味.
這幾種味道.
是一定會有的.
任何的菜做出來都會有.
同樣的在任何的屬靈傳統.

$^{881}$任何的靈修操練那裡.
必然要有對付肉體這一方面.
是必然要有這個元素的.
否則你的屬靈生命是很膚淺的.
在所有的教會歷史.
任何的屬靈操練裡面.
必然都會有要對付你的肉體.
要捨己.
要不要讓你的慾望去主導你.
你要限制它.
適當的限制它.
以致到你再一次重奪.
你是慾望的主人.
不是慾望是你的主人.
這個就是主耶穌在四十晝夜禁食的示範.
去到另外一個試探.
很多人知道是食物的試探.
但是去到另外一個試探.
在聖殿跳下來.
其中有些人的解讀.
都是這個是另外一種的慾望.
就是成名的慾望.
你看大衛高拍飛.
在上面飛下來.
那就是成名了.
諸如此類.
所以聖經那裡說到.
怎樣去對付慾望不受限制.
節制是其中一樣.
很重要是聖靈的果子.
節制是聖靈的果子.
我們要重奪生命的主權.
節制本身就是那個tie 盤.
你的車你要開的時候.
你要懂得知道那個tie 盤在哪裡.
節制本身就是那個tie 盤.
節制在這裡說到的時候.
其實你可以自己去想.
因為你有些什麼.
你覺得不受控制的時候.

$^{921}$你需要去節制.
但是如果簡單用.
我們剛才所說的事情來說.
第一就是.
你要知道要停.
有時候你停不了那些東西.
你更加要為自己去畫一條紅線.
要停.
要為自己畫一條紅線要停.
譬如小朋友.
我用小朋友的比喻.
小朋友玩手機.
玩到不會停.
我們大人怎麼辦.
有些家長不理會.
就算了.
我們也不想小朋友這樣.
我們大人怎麼辦.
可能要限制時間.
我過了幾點就停了.
不能上網了.
舉例來說.
就要停.
對於小朋友來說.
他知道要停.
但是他停不了的時候.
他需要有一些外在的幫助去停.
剛才所說.
可能父母在家裡.
就限制了時間.
幾點之後就停了上網.
但是你不是小朋友.
所以我們自己有時候要.
都要為自己定一些界線.
然後要對自己殘忍.
我們很少聽.
對敵人仁慈又對自己殘忍.
我們覺得對自己殘忍這些東西.
最好就不要.
但其實是需要對自己殘忍.

$^{961}$因為你不會停的時候.
你不殘忍就繼續.
所以如果你知道那件事不能停的時候.
你定了這樣的時候.
你硬了在心上.
譬如你說.
我幾點之後.
起碼有這幾個小時是空閒時間.
可能你真的要把電話收起來.
然後就跟自己說.
放在抽屜裡關起來之後.
我看都不看.
碰都不碰.
跟自己說這兩個小時.
無論如何.
我都不可以開這個抽屜.
一開始來說.
你都是又述又引的.
可能你心裡也跟自己說.
不用這樣.
我不用放在抽屜裡.
我擺出來充電都可以.
充電我都不看.
你明白嗎.
就是不殘忍.
不殘忍就是放鬆.
放鬆就放鬆.
其實你是做不到的.
有時候你對自己殘忍就是這樣.
然後你畫了一條這樣的線.
你就要這樣去做的時候.
慢慢慢慢你才能戒.
現在是在說戒毒.
這個絕對不是誇張.
因為那些真的跟毒品是一樣的.
那些杜巴胺不斷刺激腦分泌.
剛才那個神經科學的醫生.
他已經講到明.
他說根本是有證據.
令到你腦的結構被他改變了.

$^{1001}$就好像吸毒的人吸了毒.
為什麼吸毒的人對於那些.
譬如我不知道你有沒有遇過.
如果有些吸毒的人.
他平時是一個很好的人.
但是吸了毒之後為什麼這麼沒有底線.
又偷錢.
又甚至做一些很不好的事.
去傷害別人.
去拿錢那些都會做的.
就被吸毒控制了.
當然差是這樣.
但是其實有研究的.
因為那些吸過毒的人.
那些毒品影響了他其中腦的.
那個Moral Reasoning.
就是那個道德判決那部分.
以致到即使他未吸毒之前.
是一個很有道基督教養的人.
對於那些道德判斷.
因為他吸了毒影響了他.
他判斷的能力是降低了很多.
以致到所以他吸了毒之後.
所以他甚至打人去搶人錢.
他都會做.
同樣一樣就是.
現在那些控制我們的東西.
其實是有時候令到我們上癮.
你不對自己殘忍.
吸毒戒毒是很殘忍的.
你是戒不了的.
第二就是那個能夠延遲滿足.
能夠延遲滿足.
有些東西你很想去做的.
但是你要訓練自己.
Retrain自己.
那樣東西是你做了一些.
你認為要做的東西之後的一種獎勵.
那樣是一個Retrain的過程.
是一個重新訓練的過程.

$^{1041}$你明白甚麼意思嗎.
譬如舉例來說.
我們小朋友都是這樣被人訓練的.
譬如很想看電視.
或者很想打機.
父母就跟他說.
你先做完功課.
不做完功課就不要開電視.
你不用指意.
沒得看的.
那我們就很有動力.
那就做功課.
乖乖的.
做完就看電視.
而且看電視那種.
真的成為了一個獎勵.
其實是Retrain了我的Pattern.
Retrain了我.
去重新訓練了我的生活節奏.
我每次放學回家的時候.
我就先做功課.
我不需要每一次就有人逼我去做功課.
因為在那個過程當中.
我延遲滿足.
也去重新訓練.
去Retrain自己.
甚麼是Reward.
Reward就不是經常有.
就好像那些吃薯片.
那些Reward.
就不是像現在那些剛才說的痴肥.
不斷地吃不斷地吃.
因為他覺得那個不是Reward.
不斷地吃不斷地吃.
但我們小朋友那時候.
是做了某些東西才有.
譬如Cookies.
或者糖.
你看為Reward的時候.
你對於那樣東西.

$^{1081}$Reward的東西令他上癮.
那個也是會減低的.
這個是另外一個策略.
所以有些東西.
你是要延遲它.
而你應該要擺在前面.
把Retrain自己做到Reward.
第三當然就是.
不要給感覺做主.
讓神做主.
這個也是一個很重要的一點.
甚至乎還沒有現在網絡這些控制的時候.
都會有這樣的情況.
人很多時候會給感覺做主.
舉例來說.
譬如說讀經祈禱.
讀經祈禱可能做了五分鐘.
之後就覺得很悶.
都沒有領受.
這個感覺.
很多人就讓這個感覺做了主.
然後就算了.
不做了.
是不是?會這樣的.
但是其實.
很多屬靈的操練.
不是看感覺的.
你知道這件事要做的時候.
即使你的感覺很悶.
這樣真的沒有什麼得著.
你跟自己說.
這樣如何?.
繼續再做.
你有沒有這樣去Train自己?.
是要這樣去Train自己的.
保羅說.
我洩功克己身叫身服我.
I beat my body and make it my slave.
即是我打自己的身體.
令到他做我的奴隸.

$^{1121}$他的意思就是這樣.
而我們是沒有這樣鞭策自己的.
我們對於自己是很仁慈的.
對於敵人是很殘忍的.
但是保羅說我對敵人很仁慈.
對自己很殘忍.
不要給感覺做你的主.
很多時候.
很多是感覺.
感覺本身是一個.
給我們去經驗這個世界的東西.
但是最終.
我們是一個怎樣的人.
不是我們的感覺去無助我們做一個怎樣的人.
是你是一個怎樣的人去無助你有什麼感覺.
你明白我的意思嗎?.
你是一個怎樣的人.
你怎樣去看這個世界.
就會有很多不同的感覺.
你會愛會恨.
譬如你看前面這個人是一個壞人.
你又不喜歡他.
你就會有這樣的感覺.
因為你這樣看這個人.
所以同樣都是一樣.
你要成為一個怎樣的人.
不是讓你的感覺去Train你.
而是要Train自己去做一個怎樣的人.
然後你的感覺才會受你去Train.
而這些是真的需要我們.
首先你要知道.
現在我們就是活在一個世界.
只是讓感覺來控制著你.
這個所有都是講感覺的時候.
所以無怪乎有時候我們就像剛才所講的.
變了一個Zombie世代.
變了一個喪屍世代.
沒有任何回應的能力.
就像上一堂所講.
我們只是有React.

$^{1161}$我們沒有Respond.
任何東西都是React.
React是怎樣呢?.
多一下你的手指.
或者以前做那些小朋友做的那些實驗.
你打你的膝蓋就自然會一腳踢上來.
那些是React.
但你不是Reaction.
而福音其中一個很重要的.
就是幫助我們找到真正的生命之糧.
真正我們追求的滿足.
是由神而來的滿足.
往往其實我們沒有真正的滿足的時候.
這個世界的東西才能夠吸食我們.
所以我們要在神那裡得到滿足的時候.
我們真的就讓神來去無助我們.
剛才所講的幾樣東西.
有建議.
第一就是定界線.
第二就是真的要有一些安靜的空間.
你要有一些安靜的空間.
你要收起你的手機.
開郵輪做React.
其中一樣東西好.
就是沒有手機.
這個是其中一個最好的.
你去到其他地方.
有手機一樣是不行的.
那個就沒有手機.
還有要學習一些屬靈的操練.
這個屬靈的操練沒有時間講.
我們遲些主日學會教一些屬靈操練的東西.
其實這個是很重要的.
用真實的關係取代了虛擬的刺激.
真實的跟人聊天和在網上跟人聊天是不同的.
真實的跟人出來見個面吃個飯.
和在網上那種接觸都是不同的.
我們要有些真實的去建立取代那些虛擬的東西.
求神幫助我們.
\newpage



\section{創世記 27:19}
\label{sec:031782}
\textbf{聖經人聖經事:我是你的長子以掃}~~~~黃克勤牧師博士
\newline
\newline
連結: \href{http://www.church.com.hk/acms/content.asp?site=cdc\&op=show\&type=product\&code=031782\&layout=sermon}{\texttt{031782}} ~~~~ 語音日期: 2026-02-28
\newline
\newline
\hyperref[sec:031781]{\small{< < < PREV SERMON < < <}}
~
\hyperref[sec:index_chronic]{\small{[返順時目]}}
~
\hyperref[sec:031783]{\small{> > > NEXT SERMON > > >}}
\newline
\newline
創世記 27:19
\newline
\begin{longtable}{cl}
\hline
\hline
章節 & 經文 (和合本修訂版)\\
\hline
27:19 & \begin{tabularx}{0.7\textwidth}{X} 雅各對他父親說：「我是你的長子以掃。我已照你吩咐我的做了。請起來坐著，吃我的野味，你好為我祝福。」 \end{tabularx} \\ \\
[1ex]
\hline
\hline
\end{longtable}
$^{1}$大家一看這個主題.
「我是你長子耳掃」.
你猜不猜到是誰說的呢?.
是不是耳掃說的呢?.
如果承接著這個故事.
當然不是.
是雅各說的.
騙他爸爸耳殺.
「我是你長子耳掃」.
《生命說宣教》網絡的聖經人聖經士.
和大家一起讀聖經.
希望你可以多點支持宣教.
請多點人一起來看我所說的話.
都預備得挺辛苦挺好的.
挺好的.
每天五元.
支持我們宣教工作.
150元一個月.
我們送一個聖經人聖經士給參加者.
希望不會看我樣子很害怕.
如果看得害怕.
覺得吃驚.
可以只聽聲音.
好.
開玩笑.
聖經人聖經士.
「我是你長子耳掃」.
《創世記二十七章》第十八節.
雅各到他父親那裡說.
「我父親」.
他說.
「我在這裡,我兒,你是誰?」.
去到他爸爸那裡.
雅各就說.
「我父親」.
挺有趣的.
我不知道.
我沒有特別研究過.
這是否當時人的說話習慣.
或當時習慣叫「usage」.

$^{41}$是否這樣慣用.
我們通常見到爸爸.
去找爸爸.
「爸爸」.
叫「爸爸」.
很少叫「我的爸爸」.
是否這樣.
「我父親」.
如果你說「我父親」.
有時會生疏.
心裡有些隔膜.
或叫親切一點.
「爸爸」.
「父親大人」.
向你請安.
「什麼事?」.
「爸爸」.
即是這樣.
他說「我父親」.
「是我的爸爸」.
這麼尊敬.
因為他要騙他.
可能拿著東西成王成孔.
所以說「我父親」.
耳塞說「我在這裡」.
「你是誰?」.
是這樣的.
你會見到.
很有趣.
耳塞問「我在這裡」.
然後「兒子」.
他認得是兒子.
「你是誰?」.
他聽得出嗎?.
我認為耳塞聽得出.
這是啞角的聲音.
是兒子的聲音.
不是大兒子的聲音.
不過很有可能.
兒子較少找父親.

$^{81}$生疏一點.
「我父親」.
「騙我在這裡」.
「騙你」.
「兒子」.
「你是誰?」.
看不到.
但他心目中不聽聲似啞角.
啞角說「我是你長子耳塞」.
這是第一次.
耳塞查過很多次.
「你是誰?」.
爸爸查過.
但他明騙他.
「我是你長子耳塞」.
「我已經照你所吩咐我做」.
「請你起來坐著吃我的野味」.
「好給我祝福」.
說得不錯.
不過他又說得多了一點.
「你起來吃我祝給你吃的野味」.
他爸爸說「你去打獵」.
「找獵物做糧食」.
耳塞的個性可能是.
「爸爸,打獵了」.
「野豬,起來吃吧」.
沒他說得那麼文雅.
「請起來坐著吃我的野味」.
「好給我祝福」.
耳塞又查他.
「兒子,為何你那麼快」.
別欺負看不見的人.
我心水很清.
盲俠聽聲很叻.
心水又很清.
「可兒,你如何找得這麼快」.
「剛出去沒多久」.
「這麼快搞掂,還整好買」.
他又說「是因為爺爺」.
「這次麻煩了」.

$^{121}$「你把這件謊言拉到耶和華落水」.
「斷不以有罪的為無罪」.
「你拉到耶和華落水」.
「耶和華,你的上帝」.
「你的上帝」是尊敬的說法.
很少說「我的上帝」.
因為爸爸在.
「你的上帝」是使我遇見好機會得著.
一出去有隻野豬撞過來,然後躺下.
其實野兔跑得很快,不易捉.
使我遇見好機會.
第一,你是誰?.
第二,為何你那麼快?.
第三,我兒子走近一點.
我摸一摸你.
第三次查他.
知道你兒子耳熟不識,有懷疑.
什麼因素?.
因為他一聽到聲音.
「咦?弟弟啞角,不是大兒子」.
然後問「什麼事?」.
「哦,我已經整好野味了」.
「o下?那麼快?」.
心水很清.
不過他這樣說.
查多次,你走近一點.
我摸一摸.
幸好李伯家族知道.
啞角靠近父親耳索.
耳索摸著他說.
「咦?聲音是啞角的聲音」.
所以他一進來.
「爸爸,我的父親」.
他已經知道是啞角.
「手卻是耳素的手」.
耳索也說得很清楚.
心知是可以,然後再查他.
撕掉他的毛,看他捱不捱痛.
「哎呀!」.
如果撕掉他的毛.

$^{161}$可能啞角也會捱痛.
「哎呀!」.
大家都很狡猾.
聲音是啞角,手卻是耳素.
他分辨不出來.
因為他手上有毛.
像哥哥耳素的手一樣.
於是便給他祝福.
又說,第四次查他.
「你真是我兒子耳素?」.
他回答「是的」.
「你遞給我,我很欠我兒子的惹味」.
「然後為你祝福」.
他便遞給他.
「吃了」.
吃完後拿酒給他喝.
這招真是絕.
給他喝的酒濃一點.
喝完後有點嗡嗡的.
可能沒那麼時間查他.
所以這是第四句.
「你真是我兒子耳素?」.
他查了第四次.
「他父親耳殺對他說」.
「我兒,你上前來與我親嘴」.
他很聰明.
吃飽後,臨祝福前.
便說「你與我親嘴」.
親嘴表示親密.
他們經常都會親嘴.
有可能是祝福前想親嘴.
表示先擁抱兒子.
其實同一時間耳殺都非常聰明.
「你上前來與我親嘴」.
是第五次查他.
上前來與父親親嘴.
父親一聞他的衣服.
父親親嘴時先聞他的衣服.
一聞.
聞到衣服上的香氣.

$^{201}$有兩個可能.
第一,他的衣服有點土土地.
惹味.
不是惹味.
有點土土地.
有點惹味.
有點草的味道.
甜惹的味道.
亦可能他喜歡某種味道.
這可能是他大兒子喜歡的味道.
為他祝福.
「我兒的香氣如同耶和華」.
「似福之田地的香氣一樣」.
即是甚麼意思.
一聞是惹草的味道.
植物的味道.
可能他兒子身上有這些味道.
「碌地」也說不定.
甜惹的味道.
可能有泥土的味道.
可能有草,植物的味道.
一聞.
所以其實以撒都很聰明.
心很清.
查了兒子五次.
算上親嘴順便聞他的時候.
於是就去祝福了.
這個祝福有五個願.
很重要.
第三個願是甚麼呢.
「願上帝賜你天上的甘露.
地上的肥土.
並許多五穀生就」.
所以第一個願是.
「願你食得飽.
而且富足富有」.
食得飽.
第二.
「願多民事奉你.
多國跪拜你」.

$^{241}$即是你要統治多國.
自己立國並統治多國.
第三.
「願你做你兄弟的主.
連你母親的兒子都向你跪拜」.
就是在兄弟之中.
你做一家之主.
第四.
「誰咒詛你.
願他受咒詛」.
即是「願你蒙耶和華的保護」.
第五個是為你祝福的.
「願他也蒙福」.
就是「願你蒙福順利」.
其實這五個願都頗重要.
於是就為他祝福了.
接著我們看看今日的金句.
27章19節.
「雅各對他父親說.
『我是你的長子耳熟』」.
已經整理好了沒有?.
請你留意.
這一句不要學.
因為這一句是欺騙他父親的說話.
亦看到他不是無意.
是有意的.
有預謀的.
而且這些欺騙.
還拉了耶和華下水.
耶和華被他遇到好機會.
其實是很差.
順溝了這麼久.
原來他的靈性很弱.
由尼伯加至雅各來看.
他的靈性都不是很強.
真是很可惜.
有時我們看到一些人.
靈性好像很好.
說出來很好聽.
但其實他背後是怎樣.

$^{281}$骨子裡真是不知.
或者真的很好.
去到某些位置.
他就會失腳.
我這樣說即是你很行.
不是,我都去到某些位置.
可能會失腳.
因為當我聞到野味的味.
可能我就很開心.
其實每個人都有他的弱點.
雖然以撒察覺到.
檢查了五次.
不過都被兒子騙到.
其實是他有心騙他.
想好了辦法.
「我是你長子二嫂」.
其實最深層是甚麼呢.
雅各很想做長子.
老實說.
如果長子的位置是你的.
就是你的.
不是你的.
如果上帝要給長子名分的.
相應祝福給你.
祂會給.
這就是我們要信靠主.
有祂最好的安排.
大家一起祈禱.
天父上帝我們信靠你.
對我們有最好的安排.
安排可能每個人都不同.
不過我們總覺得上帝安排別人比我們好.
總覺得別人的命比我們好.
為何他那麼好.
我又沒有.
心裡很疑惑.
甚至很多時候見到作惡的很肥.
會為作惡的心懷不平.
不過其實知道主有你的作為.
你是做物之主.

$^{321}$是看顧之主.
我們仰望你.
既然主那麼大力安排我們.
我們就順著你的心意.
其實每個時候.
你都為我們有最好的管理.
我們信靠你的掌管.
禱告奉耶穌名求.
阿們.
\newpage



\section{約翰福音 6:22-40}
\label{sec:031783}
\textbf{榮耀君王:永生天糧}~~~~林鍵牧師
\newline
\newline
連結: \href{http://www.church.com.hk/acms/content.asp?site=cdc\&op=show\&type=product\&code=031783\&layout=sermon}{\texttt{031783}} ~~~~ 語音日期: 2026-03-07
\newline
\newline
\hyperref[sec:031782]{\small{< < < PREV SERMON < < <}}
~
\hyperref[sec:index_chronic]{\small{[返順時目]}}
~
\hyperref[sec:031784]{\small{> > > NEXT SERMON > > >}}
\newline
\newline
約翰福音 6:22-40
\newline
\begin{longtable}{cl}
\hline
\hline
章節 & 經文 (和合本修訂版)\\
\hline
6:22 & \begin{tabularx}{0.7\textwidth}{X} 第二天，留在海的對岸的眾人發覺那裡原來只有一條小船，而且耶穌沒有同他的門徒上船，是門徒自己去的。 \end{tabularx} \\ \\ \relax
6:23 & \begin{tabularx}{0.7\textwidth}{X} 另外有幾條從提比哩亞來的小船，卻停靠在主祝謝後給他們吃餅的地方附近。 \end{tabularx} \\ \\ \relax
6:24 & \begin{tabularx}{0.7\textwidth}{X} 這時眾人見耶穌和門徒都不在那裡，就上了船，往迦百農去找耶穌。 \end{tabularx} \\ \\ \relax
6:25 & \begin{tabularx}{0.7\textwidth}{X} 他們在海的對岸找到他後，對他說：「拉比，你幾時到這裡來的？」 \end{tabularx} \\ \\ \relax
6:26 & \begin{tabularx}{0.7\textwidth}{X} 耶穌回答他們說：「我實實在在地告訴你們，你們找我，並不是因見了神蹟，而是因吃餅吃飽了。 \end{tabularx} \\ \\ \relax
6:27 & \begin{tabularx}{0.7\textwidth}{X} 不要為那會壞的食物操勞，而要為那存到永生的食物操勞。這食物是人子要賜給你們的，因為父神已印證了。」 \end{tabularx} \\ \\ \relax
6:28 & \begin{tabularx}{0.7\textwidth}{X} 於是他們問他：「我們該做甚麼才算是做神的工作呢？」 \end{tabularx} \\ \\ \relax
6:29 & \begin{tabularx}{0.7\textwidth}{X} 耶穌回答，對他們說：「信神所差來的，這就是神的工作。」 \end{tabularx} \\ \\ \relax
6:30 & \begin{tabularx}{0.7\textwidth}{X} 於是他們對他說：「你行甚麼神蹟，好讓我們看見而信你呢？你到底要做甚麼呢？ \end{tabularx} \\ \\ \relax
6:31 & \begin{tabularx}{0.7\textwidth}{X} 我們的祖宗在曠野吃過嗎哪，如經上寫著：『他從天上賜下糧食來給他們吃。』」 \end{tabularx} \\ \\ \relax
6:32 & \begin{tabularx}{0.7\textwidth}{X} 於是耶穌對他們說：「我實實在在地告訴你們，那從天上來的糧不是摩西賜給你們的，那從天上來的真糧是我父賜給你們的。 \end{tabularx} \\ \\ \relax
6:33 & \begin{tabularx}{0.7\textwidth}{X} 因為神的糧就是那位從天上降下來，並且賜生命給世界的。」 \end{tabularx} \\ \\ \relax
6:34 & \begin{tabularx}{0.7\textwidth}{X} 於是他們對他說：「主啊，請常常把這糧賜給我們！」 \end{tabularx} \\ \\ \relax
6:35 & \begin{tabularx}{0.7\textwidth}{X} 耶穌對他們說：「我就是生命的糧。到我這裡來的，絕不飢餓；信我的，永不乾渴。 \end{tabularx} \\ \\ \relax
6:36 & \begin{tabularx}{0.7\textwidth}{X} 可是，我告訴過你們，你們已經看見我，還是不信。 \end{tabularx} \\ \\ \relax
6:37 & \begin{tabularx}{0.7\textwidth}{X} 凡父所賜給我的人，必到我這裡來；到我這裡來的，我總不丟棄他。 \end{tabularx} \\ \\ \relax
6:38 & \begin{tabularx}{0.7\textwidth}{X} 因為我從天上降下來，不是要按自己的意願行，而是要遵行差我來那位的旨意。 \end{tabularx} \\ \\ \relax
6:39 & \begin{tabularx}{0.7\textwidth}{X} 差我來那位的旨意就是：他所賜給我的，要我一個也不失落，並且在末日使他復活。 \end{tabularx} \\ \\ \relax
6:40 & \begin{tabularx}{0.7\textwidth}{X} 因為我父的旨意是要使每一個見了子而信的人得永生，並且在末日我要使他復活。」 \end{tabularx} \\ \\
[1ex]
\hline
\hline
\end{longtable}
$^{1}$歡迎大家來到我們的周六崇拜.
三個星期前我們來到《約翰福音》第六章.
是整本《約翰福音》裡面最長的一章聖經.
共有七十一節.
在整章聖經裡面.
我們可以把它分為五個段落.
第一個段落是.
「五餅二魚餵飽五千」.
是《約翰福音》中的第四個神蹟記號.
接著是「黑夜狂風海上行走」.
是《約翰福音》中第五個神蹟記號.
接著有一段是一個過渡的經文耶穌不在.
接著有一群人上船去追趕.
然後帶我們進入天國教道的第一部分.
「我是天降生命的糧」.
接著還有另外一個天國教道的第二部分.
是「吃喝人子就有永生」.
就這樣看來你會發現第二部分.
「黑夜狂風海上行走」.
和其他的部分有點格格不入.
其他部分都是和耶穌五餅二魚的神蹟有關係.
但是中間第二個好像和它完全沒有關係.
如果你上次兩個星期之前在這裡.
你會知道經過了仔細的研讀的時候.
我們會發現原來五餅二魚和海上行走.
不單在時序上是相連.
他們更是一個神學上.
即是聖經上的龍鳳胎.
正所謂好事成雙.
五餅二魚的結尾是十二個男子.
而海上行走的起頭是十二個門徒.
海上行走在《約翰福音》中.
是唯一一件行在門徒身上的神蹟騎士.
其他神蹟騎士都是眾人或者某一個人.
唯有這個是耶穌為門徒度身訂造的神蹟記號.
原因是什麼呢?.
是要給他們將來事公上的一種鼓勵和一個認許.
因為很快耶穌就要差派這十二個門徒.
在耶穌受死活升天.
聖靈降臨之後.

$^{41}$耶穌不在了.
當然聖靈來了.
耶穌與他們同在.
耶穌要差派他們進入一個黑暗的世界.
Into Darkness.
為了要將他們所承受的餅碎.
每一個門徒都拿了一籃的餅碎.
那個餅碎是耶穌生命的糧.
要將這個生命的糧傳給以色列的十二支派.
又再通過以色列人傳到地極.
在這個 Into Darkness 進入黑暗世界的使命當中.
他們極其需要耶穌的同在.
而耶穌自己也答應他們.
「必天之與門徒同在,直到世界的終結」.
於是在海上行走的神蹟裡面.
出現耶穌的同在.
成為了他們將來很大的鼓勵.
但我們經過了這兩個神蹟記號之後.
我們進入今天的經文.
今天的經文總共有十九節的聖經.
我將它分為七個部分.
我想你留意每一部分.
我給它的分題.
這七個部分分別是.
逾越鐵桿粉 過海 君王.
大家知道什麼叫鐵粉.
這個不是普通的人.
有一群人是逾越的鐵粉絲.
他們過海 君王.
待會我們讀到那些經文.
你就會比較明白.
第二 拉彼耶穌進會堂宣講永生道.
第三 信神所差的就是做主公.
第四 我父賜天良生命給世界.
第五 我是生命良 信我不饑渴.
第六 見我還不信來的不驕氣.
最後 君子從天降 信子得永生.
你看了我的分題.
或者你看了整個十九節的故事大綱.
如果我問你.

$^{81}$你覺得故事的核心思想在哪裡.
你看不看到.
不如這樣吧.
我們將整個核心放在最中間那裡.
故事的核心思想是在中間那裡.
我父賜天良.
這個良食是從天而降的.
這個良食給了這個世界新的生命.
永恆的生命.
我想問你.
如果你看到這個故事的中心思想.
令你想起什麼呢.
我一看完之後.
我就想起楊福音三章十六節.
金句中的金句 上帝愛世人.
甚至將他的獨生子賜給他們.
叫一切信他的不致滅亡反得永生.
我父賜天良.
上帝將他的兒子.
差到這個世界裡面.
將生命賜給這個世界.
上帝愛世人.
甚至將他的獨生子賜給他們.
如何知從天而賜.
叫一切相信的不致滅亡反得永生.
生命給這個世界.
原來這段聖經.
又是代表了上帝對以色列人.
那個叛逆子民的大愛.
是上帝對世界裡面.
離棄了他與他隔絕的人.
的一種愛 呼喚愛的表達.
如果你熟習約翰福音前面的五章聖經.
你會知道在約翰福音第四章裡面.
上帝愛撒瑪利亞.
於是耶穌來到撒瑪利亞.
呼召那個不單止撒瑪利亞的婦人.
還有其他跟婦人聯想關係的.
一群撒瑪利亞的人.
接著去到第五章.

$^{121}$你記得嗎 上帝愛三八.
耶穌行了一個神的記號.
他選擇了一個三十八年行不到的.
其實是以色列人的代號.
所以上帝愛三八.
去到約翰福音第六章.
我可以說上帝愛十二.
因為他有十二個門徒.
有十二個餅碎.
十二個餅碎是給以色列人的十二個支派.
所有的東西代表了上帝.
要將這個救恩帶到以色列民的當中.
所有的一切.
都是在講上帝的大愛.
上帝的拯救.
有了這個大圖畫.
我們現在可以進入今天的福音故事.
約翰福音第六章第二十二節.
二十三節.
第二天.
留在海的對岸的眾人.
發覺那裡原來只有一條小船.
而且耶穌沒有和他的門徒上船.
是門徒自己去的.
另外有幾條從提比利亞來的小船.
就停靠在主.
祝謝後給他們吃餅的地方附近.
這個說第二天.
第二天之前那天發生什麼事.
五餅二魚.
之前那天還有一樣東西.
凌晨的時候.
四更.
如果你記得四更是三點至六點的時候.
耶穌行在海上.
之前那天發生的事.
現在的第二天早上.
他說留在海的對岸的眾人.
發覺.
如果你看他的原文.

$^{161}$他的英文都是這樣的.
看見.
他沒辦法看見.
門徒就是搭了一條船走.
然後耶穌沒有去.
他看不到.
因為沒有閉路電視.
所以這個看見的意思是.
他們發覺.
可能他們第二天早上起床的時候.
是訪問了當地的人.
當地人是親眼目擊到.
說起昨天晚上.
耶穌去了哪裡.
我們只看見門徒坐船去了那邊.
大概那邊.
耶穌沒有上船.
所以我們不知道.
但是他們找不到耶穌.
於是他假設耶穌.
他怎麼也去了那邊.
在同時之間.
第十三節.
看見當日.
他們來到之前分餅的地方.
發現海邊又有幾隻小船.
可能是今天早上才到的.
停在那裡.
之前下船的地方.
於是他們就預備要上船.
第十四節.
這時眾人看見.
耶穌和門徒都不在那裡.
就上了船.
其實只是搏一下.
因為你真的不知道.
耶穌有沒有過海.
你假設耶穌過了海.
於是你搏一下.
往加伯隆去找耶穌.

$^{201}$第一個要你知道的就是.
我們說眾人.
這個眾人不是五千人.
這個眾人不是二萬人.
像是耶穌餵餅那些.
有五千有二萬.
不是.
因為只有幾條小船.
只有幾條小船.
每條船可能十幾個人.
所以你加起來.
可能最多都是幾十個人.
這幾十個人.
我叫他們做忠實粉絲.
他們不是普通人.
他們是鐵幹粉絲.
他們不知道為什麼.
就是要找耶穌.
為了要找耶穌.
他們搏一下.
他們不知道耶穌真的過了海.
他們就上船往加伯隆去.
那時大家可能都知道.
加伯隆是耶穌加利利.
事公的總部.
雖然我們昨天不知道.
耶穌有沒有過海.
但是門徒過了海.
耶穌可能找其他方法.
我們不知道.
但是我們搏一搏.
於是去到第25節.
在海的對岸真的找到祂.
說找到了祂.
對祂說拉比你.
何時到這裡來的.
上帝保守.
難得有情人.
無論如何都要過海找耶穌.
結果居然在對岸找到耶穌.

$^{241}$他們一見到耶穌.
他們很激動地叫耶穌拉比.
你有機會回去看《讓福音》前面.
門徒叫耶穌就是叫拉比.
什麼意思.
會不會這班人其實很盼望.
可以做耶穌的入室弟子.
於是一見到.
不是叫老闆.
叫拉比.
是不是他們都想從此以後.
加入戰團.
和十二個門徒一齊捏過.
做耶穌的跟班.
他們問耶穌一個問題.
他說你何時到這裡來.
我想他們心裡面.
很多神學家都說.
他們更想問的是.
你怎樣來到這裡的.
照我們所知.
昨天只有一條船.
而且你沒有上船.
究竟你如何來到這裡.
他們不知道耶穌如何來到這裡.
如果他們有人告訴你.
昨天凌晨四更.
三點至六點的時候.
耶穌水上行走.
帶我們過來的.
我想他們一定立刻熱血沸騰.
在我們還沒看耶穌.
怎樣去回答他這個問題的時候.
因為他問耶穌你什麼時候來到.
我們見到耶穌回答他的時候.
我想交代一下.
這個故事還有一些背景.
是我們需要知道的.
如果你繼續讀下去.
第26節不停地讀.

$^{281}$讀到第6章59節.
你會發現原來他們一群人.
找到耶穌的地方.
是在會堂裡面.
他們不是在曠野.
不是在岸邊找到耶穌.
他們是在加伯隆的會堂裡面.
這些話是耶穌在加伯隆會堂裡.
教導人時候說了.
耶穌在加伯隆的會堂裡面.
講到教導人.
因為教導的對象包括了十二門徒.
包括了在加伯隆裡面的猶太人.
猶太人才會去會堂.
還有的就是跟著他過海.
立定志向要跟隨的幾十個鐵幹粉絲.
這是新人來的.
他們就加入了這個戰團.
更重要的是.
故事還有另一個更大的背景.
背景是逾越節近了.
千萬不要忘記逾越節近了這個背景.
來追星的眾人.
全部都是逾越節的鐵幹粉絲.
大家明白嗎.
他們不是單單很看得起耶穌.
他們不是單單覺得耶穌很帥.
於是想跟隨耶穌.
他們跟隨耶穌的原因是因為.
逾越節近了.
耶穌很重要.
但是逾越節更加重要.
他們想要的是逾越起義.
反清復明.
反清復明其實是反羅馬.
復以色列.
這件事對他們來說很重要.
講到最後.
耶穌只是他們眼中.
他們手中的一個棋子.

$^{321}$但是很有威力的棋子.
但是這個棋子對他們整個夢想很有幫助.
於是他們必須要找到耶穌.
我想問你.
你猜耶穌知不知道.
看不看得出這群人的動機和野心.
當然知道.
耶穌一眼就看透了他們.
於是耶穌基本上沒有回答他們的問題.
耶穌馬上單刀直入.
第26節.
耶穌回答他們說.
他們問的問題是什麼.
他們問的問題是.
你什麼時候過來的.
你很傻.
耶穌單刀直入.
耶穌看透他們的心.
耶穌說我實實在在的告訴你們.
你們找我並不是因見了神蹟.
而是因吃餅吃飽了.
這是什麼意思呢.
這個很重要.
你中文是看不到的.
就是.
如果你這樣讀中文.
你就說你們找我不是因為見了神蹟.
你以為那神蹟是五命而餘的.
不是.
他連神蹟都不認識.
他只知道因為吃飽了.
不是的.
原來這個神蹟是福數.
看不看到.
是福數.
這是什麼意思呢.
很多基督徒.
很多牧師.
聖經學者.
都有一個誤會.

$^{361}$我覺得這裡.
他們以為這群人.
只是希望再吃多一餐免費午餐.
於是他們長途跋涉.
坐船到江.
來到嘉伯龍.
這個不是過九龍.
不是這麼簡單的.
為了做這麼多事情.
去吃多一餐免費午餐.
你一聽就知道沒這麼簡單.
完全不是.
在《讓本音》第六章第二節裡面.
五命而餘的故事裡面.
也可以告訴你.
大群人起初因為看到.
耶穌在病人身上所行的神蹟.
這個神蹟是福數.
他們就跟隨耶穌.
當時沒有這麼火.
什麼時候開始火呢.
他們看見另外一個單數的神蹟.
那個單數就是.
吃了五命而餘之後就火了.
於是他們看到五命而餘的神蹟.
就強迫耶穌.
強迫耶穌要做王.
所以你看到他們起初注目的神蹟.
由福數變成單數.
但是他們的熱情大大的升級.
因為那個單數的神蹟.
比之前福數的神蹟更加重要.
更加讓他們覺得.
就是祂了.
所以你看回上面那裡.
耶穌說.
我實在告訴你們.
你們找我.
並不是因為前面那些神蹟.
你們找我是因為後面那個神蹟.

$^{401}$因為你們都有吃餅.
吃完餅你就瘋癲了.
所以你知道這群人.
其實耶穌知道他們來找我.
不是因為他醫治的神蹟.
耶穌知道他們來找他.
是因為吃餅的神蹟.
也是因為吃了餅.
他們想叫耶穌做王.
但是耶穌逃脫了.
不要緊.
有一群瘋癲的粉絲.
鐵粉絲.
過河來找他.
是因為他們想繼續叫耶穌做王.
如果你可以想像.
我們拍一套戲.
那些幾十個眾人.
幾十個鐵粉絲.
昨晚吃了餅.
手中還有些餅碎.
拿回家不能入睡.
門徒有餅碎.
他們都帶了一些進袋.
回家磨磨.
這個餅碎是什麼東西.
這個人怎麼可能有.
這樣的法力.
五千個男人吃得飽飽.
我現在還在舔餅碎.
Who is this person.
這個人究竟是什麼.
他來這個世界.
究竟要做些什麼.
於是睡不著覺.
清早就起身.
立定志向要找他.
這幾十個眾人.
絕對不是為了免費午餐而來的.
他們是盼望耶穌做王.

$^{441}$他們希望陪伴耶穌的左右.
他們希望輔助耶穌成就千秋霸業.
然後和耶穌一起共享榮華.
你明白嗎.
這個叫鐵粉絲.
打破鐵鞋無壁處.
得來全不費工夫.
就是要找他.
怎樣都要找他.
我今天如果失去了他.
我也不知道以後能不能找到他.
於是趁著他還在附近.
怎樣都要找他.
這一切耶穌都看到了.
耶穌都看透了.
於是你聽耶穌的終極回答.
耶穌終極回答.
怎樣去回答呢.
耶穌說.
不要為那毀壞的食物操勞.
而要為那存到永生的食物操勞.
這食物是人子要賜給你們的.
因為父上帝已經認成了.
你馬上要想一想.
你知道所講的食物.
不是代表你日用的飲食.
因為這班人不是傻子.
他不是因為要免費午餐.
這班人有更巨大的野心.
你明白嗎.
這個食物代表了地上的國度.
耶穌知道他們來找他不是要食物.
耶穌知道他們要的是地上的國度.
耶穌知道他們要的是地上的榮耀.
於是耶穌和他們說.
你不要為地上的國度.
不要為地上的榮耀操勞.
你要為存到永遠天上的國度.
天上的榮耀操勞.
這個很重要.

$^{481}$我想問你.
耶穌說你要不要操勞.
要.
你要不要努力奮鬥.
要.
要不要獻上所有.
當然要.
但是你要小心.
你奮鬥努力犧牲獻上的方向.
有兩個選擇.
一個是必定扭壞的方向.
你們現在的方向就是必定扭壞的.
死硬的.
另外一個是必不扭壞.
存到永生的方向.
OK.
耶穌說你們不要追求地上.
要追求天上.
接著下半節他說.
這食物是人子要賜給你們的.
因為上帝已經印證了.
但你記得食物不是代表食物.
下半節說什麼.
他說這食物.
你現在可以代入去.
這國度.
你們追求的國度.
這榮耀.
這食物.
這國度.
這榮耀.
是人子.
要賜給你們的.
人子是什麼.
人子就是人子.
No,人子是《大義利書》第七章裡的人子.
耶穌不是經常叫自己人子的.
有時.
然後他會發牌.
人子是誰.

$^{521}$人子是《大義利書》第七章裡的.
在耿古常在者面前.
天父那個.
因為是義將.
領受了永恆天國的人子.
他說所有天國的榮耀在我手中.
我要賜給你們.
因為天父已經.
印證了.
我就是.
那個天國的.
掌權者.
領受者.
天父如何.
印證耶穌.
天父如何印證耶穌.
有很多人.
很多基督徒.
很多牧師.
學者就會說.
你看看.
之前耶穌行那個神蹟的時候.
祂.
祝謝了.
然後才分餅的.
如果你看那個.
婦女福音.
即是Synoptic Gospel.
馬康馬太.
勞嘉勒 還加了一句 祂望天.
望天祝謝了.
然後再分的.
芝麻開.
什麼意思 望天.
上帝幫幫我.
五千人.
不知道行不行.
我就不行了.
看你怎麼看.
於是祂說.

$^{561}$祂望天.
祂祈求上帝.
上帝認穩他.
這樣上帝就印證了.
這個印證.
只是那個小小神蹟的印證.
但是如果你看約翰福音.
你知道上帝.
給了耶穌一個更偉大的印證.
一個更終極的印證.
在約翰福音第一章裡面.
你記得.
那些人去找施洗約翰.
施洗約翰為耶穌作證.
我讀一讀給你聽.
第一章第三十二節.
施洗約翰又為耶穌作見證說.
我曾經看過聖靈彷彿甲子.
降在.
降下.
停在他的身上.
停在耶穌的身上.
聖靈降下.
這個是上帝的印證.
我先前不認識他.
可是那猜我來用水施洗的對我算.
你看見聖靈降下來.
降下停留在哪個人的身上.
他就是那個將來.
要用聖靈施洗的.
我看見了.
我作證了.
他是上帝的兒子.
他就是基督.
這個印證是什麼.
是他受洗的時候.
聖靈的印證.
接著去到約翰福音第三章三十四節.
又說上帝所猜來的.
說上帝的話.

$^{601}$因為上帝所賜給他的聖靈.
是沒有限量的.
在舊約聖經裡面的聖靈.
去一去.
就會收回的.
但是這個不收的.
因為他.
聖靈降下停留在他的身上.
與他同在直到永遠.
毫無限量.
聖靈是上帝高納耶穌的記號.
你開始明白.
上帝怎樣印證耶穌.
他在說的不是.
簡單之前擦餅的那一刻.
上帝印證了他.
他是告訴他.
上帝就是上帝所.
耶穌就是上帝所.
揀選的君王救主.
所以耶穌可以給你的不是食物.
耶穌可以給你的是天國.
耶穌可以給你的.
是無盡的榮耀.
是永恆的榮耀.
這是一個天國的印證.
這是一個榮耀的印證.
於是我們繼續看.
於是他們問耶穌.
耶穌叫他們.
要操勞不好地上的事.
要操勞天上的事.
於是他們問耶穌.
我們該做什麼才算做上帝的工作.
這個工作是眾數來的.
是福數來的.
耶穌回答他們說.
對他們說.
信上帝所差來的就是上帝的工作.
這個工作是單數來的.

$^{641}$你中文是看不到的.
你看英文或者看原文就很清楚了.
我想你記得耶穌之前說什麼.
耶穌之前說你們要操勞.
你們要努力.
你們要奮鬥.
於是他們再追問.
我們應該做什麼才是做上帝的工作呢.
如果你問我.
我覺得這個問題是很正常的.
和有道理的.
耶穌叫他們為上帝事操勞.
於是他說我要做什麼才是做上帝的工作.
為上帝操勞呢.
這個問題本身是很正常.
很有道理.
但是很多福音派的牧師.
聖經學者.
一聽到他們問要做什麼.
反應很大.
就說哇.
律法主義出來了.
治疫病毒發作了.
每次跟你講天國就說.
我該做什麼.
瘋的.
不是的.
於是耶穌就回答他.
不是你做什麼.
是你要信我.
於是他們在想像.
耶穌出來.
錯了錯了錯了錯了.
然後背了一句耳筆鎖書.
你們得救是本乎恩恩著信.
這不出於你們自己.
而是上帝所賜的.
也不出於行為.
免得有人自誇.
耳筆鎖書第二章第八節第九節.

$^{681}$講完.
得救本乎恩恩著信.
不是你的行為.
什麼都不要做.
耶穌就加了一句.
這就是上帝的工作.
單數.
上帝的工作只有一件.
就是信上帝所差來的那一位.
就是信耶穌.
於是他們說.
是的.
恩信稱賢.
他們想做很多福數的事情.
耶穌跟他們說.
得救不是靠你做事的.
得救是靠你信我.
你恩信稱賢.
你聽聽.
分派的牧師和學者指出.
他說你看看兩個的分別.
那些人只想著做事.
那是福數.
耶穌提出的是信我.
是一件事.
只有一件事.
就是你依靠我做你的救主.
不是靠你的行為.
上帝的工作其實不是要求你做什麼.
因為你什麼都做不到.
你做完之後.
都是骯髒的東西.
根本你沒辦法拯救.
於是你們的工作就是要相信.
耶穌為你已經完成了那個工作.
耶穌為你上了十字架.
耶穌為你完美地承傳了律法.
你的工作就是要相信耶穌.
你什麼都做不到.
如果你問我.

$^{721}$我會這樣說.
恩信稱賢當然是對的.
依靠耶穌.
當然是對的.
唯一的出路就是依靠耶穌.
但我們都要問一個問題.
問題就是.
這個真是耶穌在這裡想要表達的東西嗎.
耶穌那句話.
信上帝所差來的就是上帝的工作.
是真的叫他們恩信稱賢嗎.
耶穌是不是真的在跟他們說.
你們太內捲了.
要學習躺平.
要依靠恩典.
不要靠行為.
不要動.
懂不懂躺平.
你們就是了.
你們太過律法主義了.
經常要做太多事情.
於是耶穌跟他們說.
不可以再內捲.
要躺平.
不准認真讀聖經.
不准認真聽教.
睡覺.
躺平.
現在就躺.
當林牧師舞蹈.
對的了.
不可以積極.
不可以靠行為.
這個叫信耶穌.
神經病.
耶穌說上帝的工作.
上帝要你們做的工作就是.
信耶穌所差來的那一位.
什麼叫信耶穌.
信上帝所差來的那一位.

$^{761}$我想帶你們去看一個.
福音書故事裡面.
我其中一個最喜歡的故事.
拿起福音第十章.
有個人跑來問耶穌說.
良善的老師.
我該做什麼才能承受永生.
因為永生是一個故事.
最後結局.
時間關係我們跳到結局.
耶穌看著他就愛他對他說.
你去變賣你的所有.
然後就有財寶在天上.
然後來跟從我.
如果你有機會回去看這個故事.
你會發現這個故事裡面.
信靠.
動詞.
或者是信心.
名詞.
是從來沒有出過的.
但是整個故事.
從頭到尾都是在說什麼叫信耶穌.
你明白嗎.
這個字不存在的.
信靠這個字.
信耶穌這個字沒有出現過的.
信心這個字沒有出現過.
但是整個故事.
都是充滿著福音.
整個故事啟示給我們知道.
什麼叫信福音.
你要知道福音其實是一個專有名詞.
福音說的是一個有關於君王的好消息.
所以信福音就是信君王.
君王是誰 君王是耶穌.
整個故事信沒有出現過.
但是整個故事都是跟信耶穌有關係.
這個人走來問耶穌.
我該做什麼才能.

$^{801}$承受永生.
答案不是.
什麼都不用做.
什麼都不可以做.
信耶穌就好了.
耶穌有沒有這麼累.
因為有沒有感悔.
沒有.
答案也不是變賣所有捐給窮人.
我們不是說道德主義.
我們不是說律法主義.
正確的答案.
唯一合乎聖經的答案.
唯一正確的答案是跟從耶穌.
你要永生.
你怎麼可以有永生.
跟從耶穌.
但是耶穌一眼望透他.
他知道除非這個年輕的財主.
大家知道他是年輕財主.
變賣所有.
放下他的錢財.
他不能跟隨耶穌.
但是令他得著永生的.
不是變賣所有.
令他得到永生的是跟隨耶穌.
耶穌說你要永生.
你就必須要有我.
因為我就是永生.
你有我就有永生.
你跟隨我就是永生.
這是一幅.
真信耶穌.
得著永生的圖畫.
什麼叫信耶穌得永生.
就是這幅圖畫了.
Jay Packer在《簡明神學》裡.
曾經這樣說.
他說我們所謂的信心.
他是這樣定義的.

$^{841}$是一個人全心的反應.
全心對福音的回應.
包括了這個人的心思.
心靈.
意志.
情感.
都包在內.
說真的什麼都包在內.
全人全心對上帝的回應.
然後他說.
適時的改革中.
神學家這樣分析.
真信心.
假信心有很多種.
真信心有三個層面.
第一個.
信心是認知.
notitia.
就是由上面開始.
但其實是最底層那邊.
拉丁文的意思是.
你需要有熟悉福音的內容.
你需要認知.
你知道福音.
你需要有知識.
有人跟你解釋了.
林牧師跟你解釋了什麼叫福音.
接著下一層.
或者他中間那層叫.
Ascensus.
就是.
你同意了.
我跟你說.
世人都犯了罪.
規矩是上帝的榮耀.
你現在聽了.
你有了知識了.
但你同不同意.
我不同意.
是世人.

$^{881}$有些人是犯了罪.
規矩是上帝的榮耀.
但不是我.
第二.
他說.
你要有知識.
notitia.
你要知道福音的內容是什麼.
接著.
你要明白到.
福音的一部地方.
你是同意.
你同意.
我覺得.
耶穌說的是對的.
但還有一層.
fiducia.
最後那個是.
信靠.
毀身.
這個人.
第一.
他聽了福音.
他知道福音的內容.
第二.
他開始認同福音的內容.
你說我是罪人.
我同意.
你說耶穌是獨一無二的教主.
我也同意.
但最後.
福音對這個人有一個呼召.
這個人必須要有.
fiducia.
他要有對福音的信靠.
對耶穌的信靠.
和毀身.
這個毀身是全人的信靠.
是全人的投入.
全人的跟隨.

$^{921}$你想想.
這個年輕的財主的問題就在這裡.
他和耶穌談永生.
他有知識了.
耶穌跟他說了.
或者他之前讀舊的聖經.
已經知道.
永生是什麼.
然後耶穌跟他說.
你要變賣你所有來跟隨我.
他也同意的.
為什麼.
因為他悠悠愁愁地離開.
他不同意.
他何必悠悠愁愁.
他悠悠愁愁的原因.
就是因為他同意耶穌所說的話.
但他做不到第三件事.
第三件事就是他要毀身.
他毀身不了.
公規一規.
撒旦也是這樣.
你問撒旦.
你知道撒旦知道福音是什麼嗎.
當然知道.
甚至我問你.
撒旦你覺得你認不認罪.
可能撒旦也認罪.
是的.
我是專門和上帝合謀.
我就是滿逆的.
但是你問撒旦.
你願不願意降服一人.
不行.
於是我們回到.
楊福音第六章第29節.
他們問耶穌.
我們該做什麼才算做上帝的工作.
耶穌的答案是.
上帝的工作只有一件.

$^{961}$就是來跟隨我.
上帝的工作只有一件.
就是你們要信靠我.
你們要跟從我.
跟從上帝參拜來的那一位.
那一位我一直和你說.
就是我.
如果這班人在耶穌的面前.
他們怎麼可以得到真正的信心.
第一.
他們要聽耶穌的話.
因為耶穌從頭到尾在解釋給他們聽.
什麼是天國的福音.
耶穌要他們得到Notitia.
接著耶穌要他們用心的思想.
耶穌的教導.
耶穌告訴他們.
那些有罪的人才會來找醫生.
無罪的人是不需要找醫生.
那你就要回去思考.
你是有病還是沒有病.
你是有罪還是無罪.
直到他們說.
耶穌我回來了.
我有罪.
接著耶穌說.
好,來跟隨我.
第三個.
全人的跟隨.
全人的行走.
耶穌引領他們的道路.
這個叫信靠.
於是耶穌對他們呼召.
要求他們來跟從我.
這個也是耶穌對我們今天.
每一個人的呼喚和要求.
我們要不要永生.
要永生我們就要來跟隨他.
你說這班人聽不聽得明白耶穌說.
耶穌說你要信我.

$^{1001}$你要跟從我.
你聽不聽得明白.
其實他們聽得明白.
是現代的福音.
福音派的牧師聽不明白.
他們完全聽得明白.
但是他們就是聽得明白.
耶穌叫他們全人奮身.
他們心想我們還不行.
我現在沒有辦法全心信你.
我也沒有辦法全人奉上.
於是他們再問耶穌.
你要給我一些東西我可以拿.
他們不是反對.
他們明白耶穌叫他們做的事.
他們明白耶穌叫他們放下所有去跟隨他.
但是他說你玩到這麼大.
玩到這麼大你要給我一些東西我可以拿.
於是第六章第三十字.
於是他們對他說.
你行什麼神蹟.
好讓我看見而信你呢.
你到底要做什麼呢.
有些人覺得他們是老人癡呆.
昨天不是做了一個神蹟給你看.
還要問神蹟.
昨天那個神蹟是做一次的.
你現在叫我全人奉上.
昨天那個神蹟我跟隨你三天是可以的.
如果要我跟隨你一輩子.
赴湯蹈火在所不辭.
你是不是應該給我一個更大的神蹟.
他們根本不是老人癡呆.
他們沒有忘記昨天五餅二魚那個神蹟.
但是現在耶穌挑戰他們.
提高了賭注.
賭注是全人奉上.
於是他去問耶穌.
你叫我全人奉上.
你一定要給我更大的神蹟.

$^{1041}$這個故事情節.
其實你在《約翰福音》第二章已經看到了.
耶穌出來清潔潔淨那個聖殿.
玩到很大.
那些人又不是馬上阻止他.
他問玩到這麼大.
你是不是應該給我們一些東西看.
讓我們知道你的身份.
如果你真的上帝揀選過君王.
那我就沒得說了.
他們想用神蹟來證明耶穌的身份.
所以你想.
耶穌挑戰他們的信心.
他們反過來挑戰耶穌的身份.
但我想你留意一件事.
他們要的不是要多聽耶穌講的道.
這個很有趣.
因為你回去撒瑪利亞的時候.
他們是想聽耶穌講多些道.
你想我知道你是不是上帝的君王.
你講多些道.
我用道去肯定你的身份.
但他們要的耶穌.
他們要的是更大的神蹟.
在《約翰福音》裡面.
神蹟代表了一個人對地上國度的追求.
地上榮耀的追求.
唯有話語才是對上帝天國.
上帝榮耀的追求.
你明白嗎.
所以當他們要求耶穌.
不是聽更多的道.
而是做更大的神蹟.
你知道他們人生的方向已經完全錯了.
這種對耶穌的反抗.
這種不順對耶穌的挑戰.
其實一直持續到《馬開福音》.
到十字架的面前.
我想你大力看《馬開福音》的結宗.
他們仍然在叫耶穌行神蹟.

$^{1081}$《馬開福音》第十五章第二十九節.
耶穌已經盯著十字架.
「經過的人羈嘯他說」.
「你這一拆毀聖殿三日又建造起來」.
「救救你自己從十字架上下來吧」.
「眾祭司和民事也是嘲笑他說」.
「他救了別人不能救自己」.
「上以色列的王」.
「基督現在你從十字架上下來」.
「好讓我們看看你就順利」.
他說如果你現在做一個神蹟給我們聽.
然後你還沒死.
做多一個神蹟.
就是這個而已.
你能下來.
你這樣也不死.
我們就順利.
但你知道這個所謂的順其實不是真順.
如果耶穌真的下來.
他說我現在順利.
這個順是一個邪惡的順.
一個以自我為中心的順.
因為他們希望耶穌為他們成就地上的國度.
他根本不是普普在耶穌面前.
他們貪愛世界.
他們高傲恣意.
他們不明白耶穌不可以下來.
因為如果耶穌下來.
罪人沒辦法得著救贖.
如果耶穌下來.
天國的福氣.
永生的禮物.
沒辦法臨到人間.
唯有上帝為我們預備.
蒙他揀選的受苦僕人.
才可以將天國的救恩帶進這個人間.
但以色列人一直要的.
不是十字架上的耶穌.
回到楊卜音第六章第三十節.
於是他們對耶穌說.

$^{1121}$你行什麼神蹟.
好讓我們看見而順利呢.
到底要為我們做些什麼呢.
第三十一節.
他給耶穌一個提示.
我們的祖宗在曠野吃過馬奶.
如經上所寫.
他從天上降下糧食來給我們吃.
他給耶穌一個提議.
他說你可不可以重複這個神蹟.
一天重複三次.
再重複四十年呢.
我們一想五餅二魚餵養五千人.
我們覺得很厲害.
但你記得我講那篇道的時候.
我已經跟你說.
這個對比當年出埃及.
相差很遠.
出埃及的人是二百萬人.
吃一餐耶穌.
耶穌頂窿兩萬人.
你怎麼比呢.
接著還是一天三餐.
吃足四十年.
於是他們跟耶穌做了一個提議.
你昨天做的那個真的很不錯.
但你現在要我們全力.
是不是應該給我們多一點.
不如這樣吧.
你試一下看會不會賭公.
今天試一下一天三餐.
如果可以的話.
可不可以再持續多做一會兒.
那我們就可以信你了.
這個是他們的要求.
我想你看看耶穌的回答.
耶穌的回答.
第六章第三十二節.
於是耶穌對他們說.
我實實在在地告訴你們.

$^{1161}$那從天上來的糧.
不是從摩西賜給你們的.
那從天上來的真糧.
是我父賜給你們的.
這個中文又是沒辦法了.
因為這些文法的問題.
你這樣看中文就看不到.
就是過去式現在式.
當然你可以的.
不過令到他太過困難了.
其實他在這裡說.
那從天上降下的糧.
不是摩西賜給你的.
摩西賜給你的是一個過去式.
在原文是過去式.
那從天上增糧.
是我父賜給你的.
這個賜給你的是現在式.
摩西的bread program.
摩西有一個麵包計劃.
已經過了.
已經收檔了.
接著耶穌再說.
不過.
現在有新promotion.
我父有一個現在式的麵包計劃.
你明白嗎.
他是這麼說的.
我們可以說.
你知不知道從前.
舊約裡面的摩西麵包計劃.
只是一個影子.
現在我們promotion中的.
推廣中的.
才是終極的實體.
楊寶欽第一章第17節.
慈愛信實.
創作慈愛信實影子的律法.
是藉著摩西所頒佈的.
慈愛信實的實體.

$^{1201}$是由耶穌基督來的.
我沒時間回去跟你說.
但你明白我的意思.
舊式的那個program.
他現在還在問耶穌.
可不可以重做摩西的神蹟.
耶穌說.
那件事已經過了很久.
那個program早就沒有了.
不過我告訴你.
現在有一個更厲害的.
我可以將你重新再翻譯一次.
就是這樣的.
那個從天降下來的糧食.
以前那個.
是摩西所賜給你的.
不過已經成為了過去式.
不過我告訴你.
其實當年不是摩西賜給的.
是我父上帝賜給的.
第二.
摩西死了.
但我父還沒死.
我父不會死.
接著我告訴你.
停了這麼久.
我父現在又有一個新的麵包計劃.
現在進行式.
現在我父開始.
賜下天糧.
接著告訴你.
現在的天糧不是從前的碼拿.
現在的天糧是升級了的碼拿.
是碼拿2.0.
因為他加了一句說話.
他說True Bread.
真.
真糧.
他說現在我要給你的.
不是從前摩西賜給的.

$^{1241}$你為什麼還在問舊的東西.
我有一個真正的麵包.
是現在要給你的.
於是耶穌繼續.
因為上帝的糧.
就是那位從天上降下來.
並賜生命及世界的.
就是他自己.
我們.
你不要笑他們.
因為我們已經知道.
最後的結局是什麼.
我們知道耶穌就是那個糧.
耶穌說那個升級了的碼拿.
那個碼拿2.0.
那個天上降下來的.
真糧真碼拿.
過去的影子.
就是我自己了.
就是我從天上降下來.
但他不知道耶穌從天上降下來.
耶穌說的麵包.
就是他自己.
但他們聽不聽得懂.
所以第34節.
他們對耶穌說.
主啊.
請上上把這糧給我們.
這樣又可以.
這樣又相信.
但其實是一種誤會.
他以為耶穌和摩西一樣.
是一個麵包師傅.
他以為耶穌現在.
可以按一個按鈕.
跟著天上可以降下.
一個升級了的碼拿.
這個我找了一個.
叫做super low.
超級麵包.

$^{1281}$吃了超級麵包.
就是碼拿2.0.
升級了.
所以做超級的軍人.
於是這班人誤會了.
他以為耶穌只是一個新的摩西.
當然耶穌是一個新的摩西.
但耶穌何止是摩西.
他於是說.
摩西現在死了.
拜託你按一下按鈕.
現在降下來.
我們開始吃超級麵包.
我們不需要再回到碼拿.
他說主啊.
請上上把這糧給我.
一場誤會.
這裡中文的翻譯.
又出了另外一個問題.
這些人怎麼稱呼耶穌.
主啊.
在原文.
可以翻譯做主.
也可以翻譯做先生.
女人明白了.
因為在楊福音裡面.
很多人喜歡說耶穌叫先生.
你記得那個薩瑪利亞婦人.
也是起初叫耶穌先生.
這個薩瑪利亞婦人.
也是說了很多差不多的東西.
牛頭不答馬嘴.
女人在第四章裡面.
耶穌和她說活水的時候.
請把這水給我.
使我不渴.
耶穌在說天降活水.
她聽不明白.
那時候的婦人叫什麼.
先生.

$^{1321}$在原文是同一個字.
我不知道為什麼中文.
現在你看不起女人.
女人就先生.
這群人就叫主啊.
要不兩個都主啊.
要不兩個都先生.
這群人其實不認識耶穌.
但是那個神奇女俠.
神奇女俠.
很厲害.
半個小時一個小時之內.
再不是先生.
先生去先知先知去基督.
基督說我是最後.
你是世界的救主.
Can you believe it.
這群人.
不信不明白.
從頭到尾只叫耶穌.
做先生.
他們聽不明白.
但你要原諒他.
他以為耶穌是神和世.
於是耶穌畫圖畫出牆.
第三十五節.
耶穌對他們說.
我就是生命的糧.
I am the bread of life.
我是生命的糧.
到我這裡來的.
絕不饑渴.
信我的永不乾旱.
這段聖經.
如果熟讀.
有人服飲.
應該會想起兩段經文.
我想大家都記得.
耶穌說.
你吃我這個麵包你不再飢餓.

$^{1361}$你記得耶穌怎樣跟他說.
凡渴者水活水的會再渴.
誰渴了我所賜的水.
你永遠不渴.
很平衡.
於是第四章和第六章.
加起來又不渴又不飢.
所有東西.
全部被滿足了.
一個是何等偉大.
彼此對應的福音.
這是你第一段.
你會聯繫到的經文.
但有另外一段更加重要的聖經.
是沒那麼容易去.
聯想的但非常重要.
耶穌說.
我是生命的麵包.
在這個詞句裡面.
有兩個名詞一個叫生命.
一個叫麵包.
我問你.
兩個名詞當中哪個比較重要.
想一想.
我是生命的麵包.
請問生命比較重要.
還是麵包比較重要.
當然是生命比較重要.
所以這段聖經.
帶我們去到生命在約翰福音裡面.
第一次的出場.
第一次的出場在哪裡.
在第一章.
第一章第四節.
他裡面有生命.
這生命是人的光.
光照在黑暗黑暗.
沒有勝過光.
黑暗沒有明白光.
不明白要拒絕但又沒有辦法勝過光.

$^{1401}$難怪耶穌是生命的麵包.
因為他裡面.
一早已經告訴你他裡面有生命.
如果耶穌變了麵包.
麵包裡面當然有生命.
但我想你再留意一件事.
這個麵包裡面不單止有生命.
還有光.
他說這生命就是人的光.
那個生命裡面.
耶穌的永生裡面.
又有真光.
所以.
生命裡面有真光.
麵包裡面又有真光.
所以耶穌要給我們一個真光.
生命麵包.
行不行.
不單止是生命麵包.
是真光的生命麵包.
有少少籠統.
但你要明白這個故事.
如果麵包裡面有真光就麻煩了.
因為當真光麵包.
來到黑暗世界裡面.
你知道黑暗世界的人.
是不會喜歡吃真光麵包的.
是不是.
都不知道為什麼這些學校叫真光.
在黑暗世界裡面.
為什麼叫真光呢.
於是.
你發現下一遍.
那些人就不信了.
為什麼 因為真光麵包.
真光麵包是不合口味的.
第六章第36節.
可是.
我告訴你們.
你們已經看見我.

$^{1441}$你們還是不信.
為什麼你們不信.
因為那光是真光.
來到世上照亮所有的人.
他來到自己的地方.
自己的人不接納他.
因為光照在黑暗裡面.
所以世人不喜歡吃真光麵包.
還沒完.
耶穌繼續宣講天國的福音.
我真的要記住.
耶穌現在在說道.
耶穌說道的目的.
不是要定他們的罪.
耶穌說道的目的是光照在黑暗當中.
呼喚.
黑暗去迴轉.
OK.
第六章第37節.
凡負所賜給我的人.
必到我這裡來.
到我這裡來的.
我總不挑起.
是什麼意思.
我們福音派的神學家.
牧師一看到.
凡負所賜給我的.
到我這裡來.
預定.
揀選.
凡負所賜的才會來.
反過來說.
凡不是被負所揀中的就不會來.
有什麼辦法.
牧師的朋友叫幾個不來.
他們沒有被上帝揀選.
有什麼辦法.
如果你們這樣想的話.
我想你過分消極也過分被動.
而你這個.

$^{1481}$想法.
你這個預定揀選的想法.
是不切合.
耶穌在會堂當中.
的主題.
你一定要記得.
耶穌不是在裝神學講座.
然後一早打了來.
預定論揀選.
耶穌不是在做神學講座.
耶穌是在說福音的.
講道.
耶穌不是要他們.
消極不順.
沒有被揀選.
下次不用回來.
不是.
耶穌不是要他們消極不順.
而是要他們積極的迴轉.
因為你留意下半句.
到我這裡來的人.
我總不丟棄他.
簡單來說.
如果你現在在現場.
在聽耶穌講道.
你突然之間.
有良心發現.
知道了.
原來你不是站在天父那一邊.
舊約的以色列人.
有兩種人.
一種是手約一種是伴約.
但那些犯罪伴約的人.
不知道自己是伴約.
所以才要迴轉迴轉.
如果你現在在聽耶穌的講道.
你開始發現原來我不是.
在天父那一邊.
那你應該怎樣.
你應該快快的迴轉.

$^{1521}$下面那句.
你到我這裡來.
你今天要到我這裡來.
我會不會丟棄你.
他不是說如果你不是選選.
對不起你來錯了我也不會收你.
我查一下你的名字.
你的名字很明顯不在這裡.
對不起.
耶穌說煩到我這裡來的.
我也不會丟棄.
耶穌不是在講預定選選.
耶穌在講的是.
悔改迴轉.
我想帶你看之後.
第四十五節.
下個禮拜我們回來才看的經文.
在先知書上寫著.
他們要蒙上帝的教導.
但聽了上帝的教導.
而學習的.
都到我這裡來.
OK.
耶穌在引用的已錯下書.
第五十四章.
為什麼要引用.
已錯下書第五十四章.
不是叫你認命.
上帝有沒有教導你.
你的心有沒有就向上帝.
沒有 對不起你不會來我這裡.
不是要你認命.
而是耶穌告訴我們.
今天救贖已經成全.
呼喚大家.
快快來到我這裡.
來到我這裡的.
總不丟棄.
我再講一次下個禮拜回來再講多一點.
這段聖經.

$^{1561}$先知所講的.
是已錯下書.
已錯下書五十四章.
已錯下書五十四章是已錯下書.
五十三章之後的章.
已錯下書五十三章講什麼.
受苦的僕人.
救贖的成全.
已錯下書第五十四章.
一開始的時候是這麼說的.
你者不懷孕.
不生育.
的要歡呼.
你者未曾.
經過.
燦爛的要歡呼.
揚聲的歡呼.
這是要何說的.
講了什麼 講了以色列人被擄回歸.
講了以色列人好像一個不能生育的婦人.
但是現在可以.
生育了 因為她可以.
有人得救.
已錯下書第五十四章第二節.
要擴張你的著墨之地.
伸展你.
居所的慢指.
不要退回.
要放長你的繩子.
堅固.
你的房子.
所有東西要擴張.
為什麼 因為救贖已經成全.
耶穌是呼喚他們.
今天就是救贖.
的時候回轉回轉回轉.
最後三節.
因為.
第三十八節.
我從天上降下來.

$^{1601}$不是要按自己的.
旨意行 而是要.
進行差我來的.
那位的旨意.
差我來的那位旨意就是.
他所賜給我的要.
要我一個.
也不失落.
並且在末日使他.
復活.
在這兩節聖經裡面旨意出現過.
三次.
在最後一節第四十節裡面.
又會再出現過.
多一次 耶穌說我不是要.
按自己的旨意.
中文翻譯是意願.
英文原文也是旨意.
我不是要做我自己的.
旨意行 我要遵守那.
差我來的旨意.
已經令你想起.
赫西瑪利元那一幕了.
天父上帝的旨意.
是什麼 天父上帝的旨意.
是他賜給我的.
我一個都不會.
失落 並且在.
末日使他復活 我們知道.
末日每個人都會復活的 信不信耶穌的人.
都會復活的 所以這裡說的復活不是.
復活被定罪.
使他復活的時候.
不被定罪 使他復活的時候.
可以得享永生.
耶穌定義.
要行天父的旨意.
令我們想起.
在楊福音第四章.
他去撒瑪利亞的時候.

$^{1641}$耶穌和門徒說.
我有食物吃膳你們不知道的.
我的食物就是真行.
行我差我來的.
那位的旨意 完成.
他的工作.
耶穌將完成.
上帝的旨意 行上帝的旨意.
給予做食飯.
是一個維持生命的東西.
如果我不.
做上帝的旨意.
我就死了.
我必須要 我活著.
只有一個目的 就是行上帝的旨意.
於是耶穌從天降下來.
要遵行上帝的旨意.
要完成上帝的工作.
於是在楊福音裡面.
第四章 耶穌親身.
特意來到.
撒瑪利亞 他要焦聚.
上帝的羊群.
然後在楊福音第六章.
現在耶穌在做什麼 他來到.
加伯隆的會堂裡面.
不是給神學講座 是呼喚.
迷羊迴轉.
oh my god 然後我們一定要.
知道耶穌不是在那裡.
說你死了你死了你死了.
耶穌在那裡是.
教導他們 叫他們迴轉.
然而我們知道.
耶穌單單尋找.
焦聚 呼喚.
接納是不夠的.
如果耶穌只是來到.
尋找焦聚 呼喚 接納.
那些羊 都是沒有辦法.

$^{1681}$令他們不失樂.
末日復活的時候 讓他們進天國.
唯有十字架.
才可以拯救到眾人.
使他們進入永生.
使他們不失樂 使他們復活的時候.
得著永生.
十字架在哪裡.
在摩西裡面的影子.
楊浦陰第三章.
第十四節 摩西在.
曠野怎樣舉蛇.
引子也必須.
照樣的地舉起來.
使一切信他的人.
都得著永生.
這一節是.
永生第一次的出場.
在整本楊浦陰裡面第一次的出場.
接著那一節是什麼 上帝愛世人.
都是有永生的.
但是第一次的出場.
永生在哪裡 永生在十字架上.
又說摩西.
怎樣舉蛇 引子要被舉起來.
永生在哪裡.
引子被舉起來的那一刻.
最後一節 今天最後一節.
第四十節.
他說 因為我父的子義.
少使每一個人.
見了子.
而信的人得永生.
這個希臘文的見了.
這個見了是說.
定睛仰望的意思.
不是 愁眼了.
不好意思 不小心看一看你.
都沒有穿衣服 不是.
是他們定睛仰望.

$^{1721}$耶穌的.
那一位 當你定睛.
仰望十字架的時候.
就是你得著永生.
的時候 災難末日.
我要使他復活.
進入永生 好 我時間夠了.
我們下個禮拜.
回來才繼續.
聽耶穌的講道 那篇講道還沒完.
很長的戲.
不過我想提醒大家.
我也要提醒我自己.
耶穌整篇講道的核心.
和動力.
上帝愛世人.
甚至將他的獨生子.
給他們叫一切相信的不至滅亡.
反得永生.
我們每一次在這裡相聚的時候.
我們每一次的敬拜.
和講道都是一樣.
不是跟你玩神學講座.
不是跟你講.
乜乜論乜乜論.
是要你明白福音.
從那個.
思想 從知識上面.
知道 讓你知道.
福音和你有關係.
讓你認同 世人都犯了罪.
虧缺了上帝的榮耀.
接著呼喚你.
全人的投入.
信靠.
委身給這個.
唯一可以拯救你的.
耶穌基督.
整篇道的核心.
是我父.

$^{1761}$賜天良生命.
給世界.
你有沒有.
來到耶穌基督的十字架面前.
正正仰望.
我們獨一的救主呢.
如果你沒有.
你要細心的思想.
因為祂是你唯一的拯救.
如果你有.
你記得那首詩歌吧.
當轉眼仰望耶穌.
當你轉眼仰望耶穌的時候.
再救主.
恩典的大光宗.
世俗事必要顯為虛空.
這個不是.
為了唱歌而寫的歌詞.
這個是一個基督徒的人生.
如果我們真正的.
轉眼.
定睛仰望耶穌的時候.
我們必定會.
轉眼定睛仰望.
永恆的天國.
因為世俗的事.
都會顯為.
虛空.
於是我們不只是.
看了耶穌得救了.
繼續去這個世界追求世界.
不是我們看了耶穌之後.
我們亦都轉眼定睛仰望.
永恆的.
國度.
我們開始要追求.
一個不扭壞.
傳到永生的.
榮耀和財寶.
所以我想用兩句說話.

$^{1801}$總結和延伸.
今天這個故事.
因為今天這個故事.
其實是沒有人信主的.
但我想像.
如果當中有人信.
我相信之後可能會有人信.
他們的人生會怎樣呢.
所以我現在有兩部分.
純粹三句說話.
眾人熱烈追.
耶穌不住跑.
耶穌望天上.
眾人看地下.
眾人不信他.
耶穌呼回轉.
這個很難搞.
人們.
如果你當他是一個神學講座.
眾人不信他.
耶穌就.
和他.
論揀選.
耶穌不是論揀選.
耶穌是呼回轉.
我們一定要明白.
這是耶穌的心情.
接著帶我們去到中間的十字架.
我瀑次天良.
生命及世間.
如果這個故事有下文.
是上面故事的反轉.
會是怎樣呢.
眾人若回轉.
耶穌不丟棄.
不是我錯.
耶穌自己說的.
你來的話我必不丟棄你.
耶穌神高揚.
眾人逃死亡.

$^{1841}$最後眾人求天良.
眾人再不望地下.
耶穌必同在.
你明白嗎.
我們沒有這個故事.
但是我們要延伸思想.
這個故事.
願這個成為我們的故事.
我們要祈禱.
耶穌在.
《讓福音》裡面的道.
真是高深莫測.
有很多難明的地方.
或者我們欣賞自己的神學.
傳統.
我們見到以色列人.
和他們的對話很容易先入為主.
你為何又要做事.
信耶穌就好.
接著我們將信耶穌變成一個.
不咸不淡.
完全無味的東西.
於是我們.
圍爐的在取暖.
說我們個個都信了耶穌.
說我們個個都恩順清義.
其實根本沒有人明白什麼叫信耶穌.
在.
過去的.
幾百幾千年裡面.
通過聖經.
通過神學.
聖學者的反思.
他們告訴我們.
綜合整本聖經的教導.
信不是單單.
頭腦上的認知.
甚至不是情感上.
做一個抉擇.
做一個祈禱那麼簡單.

$^{1881}$信是一個寸刃.
的投入.
這也是.
為什麼那個年輕的財主.
他若不變賣.
所有跟隨耶穌他不能信耶穌.
耶穌不是.
出一個難題給他.
耶穌是在教導他.
什麼叫信耶穌.
得永生.
耶穌看著他就愛他.
看著我們.
也愛我們.
耶穌看著這班.
加伯隆的.
鐵粉絲.
他也愛他們.
他呼喚他們.
若來我這裡.
一個也不凋棄.
但你必需要.
操勞.
天上的事情.
你必需要仰望.
我的十字架.
多謝你.
給我們這麼寶貴的福音.
願我們.
全心的領受.
懷著感恩的心.
一起迎接接下來的聖餐.
我們今天祈禱.
奉主耶穌彼得名求.
\newpage



\section{箴言 6:20-35}
\label{sec:031784}
\textbf{管教的責備是生命的道}~~~~莊明偉牧師
\newline
\newline
連結: \href{http://www.church.com.hk/acms/content.asp?site=cdc\&op=show\&type=product\&code=031784\&layout=sermon}{\texttt{031784}} ~~~~ 語音日期: 2025-08-24
\newline
\newline
\hyperref[sec:031783]{\small{< < < PREV SERMON < < <}}
~
\hyperref[sec:index_chronic]{\small{[返順時目]}}
~
\hyperref[sec:031785]{\small{> > > NEXT SERMON > > >}}
\newline
\newline
箴言 6:20-35
\newline
\begin{longtable}{cl}
\hline
\hline
章節 & 經文 (和合本修訂版)\\
\hline
6:20 & \begin{tabularx}{0.7\textwidth}{X} 我兒啊，要遵守你父親的命令，不可離棄你母親的教誨。 \end{tabularx} \\ \\ \relax
6:21 & \begin{tabularx}{0.7\textwidth}{X} 要常掛在你心上，繫在你頸項上。 \end{tabularx} \\ \\ \relax
6:22 & \begin{tabularx}{0.7\textwidth}{X} 你行走，她必引導你，你躺臥，她必保護你，你睡醒，她必與你談論。 \end{tabularx} \\ \\ \relax
6:23 & \begin{tabularx}{0.7\textwidth}{X} 因為誡命是燈，教誨是光，管教的責備是生命的道， \end{tabularx} \\ \\ \relax
6:24 & \begin{tabularx}{0.7\textwidth}{X} 要保護你遠離邪惡的婦女，遠離外邦女子諂媚的舌頭。 \end{tabularx} \\ \\ \relax
6:25 & \begin{tabularx}{0.7\textwidth}{X} 你不要因她的美色而動心，也不要被她的眼皮勾引。 \end{tabularx} \\ \\ \relax
6:26 & \begin{tabularx}{0.7\textwidth}{X} 因為連最後一塊餅都會被妓女拿走；有夫之婦會獵取寶貴的生命。 \end{tabularx} \\ \\ \relax
6:27 & \begin{tabularx}{0.7\textwidth}{X} 人若兜火在懷中，他的衣服豈能不燒著呢？ \end{tabularx} \\ \\ \relax
6:28 & \begin{tabularx}{0.7\textwidth}{X} 人若走在火炭上，他的腳豈能不燙傷呢？ \end{tabularx} \\ \\ \relax
6:29 & \begin{tabularx}{0.7\textwidth}{X} 與鄰舍之妻同寢的，也是如此，凡親近她的，難免受罰。 \end{tabularx} \\ \\ \relax
6:30 & \begin{tabularx}{0.7\textwidth}{X} 賊因飢餓偷竊充飢，人不藐視他， \end{tabularx} \\ \\ \relax
6:31 & \begin{tabularx}{0.7\textwidth}{X} 但若被抓到，要賠償七倍，他必賠上家中一切財物。 \end{tabularx} \\ \\ \relax
6:32 & \begin{tabularx}{0.7\textwidth}{X} 與婦人行姦淫的，便是無知，做這事的，必毀了自己。 \end{tabularx} \\ \\ \relax
6:33 & \begin{tabularx}{0.7\textwidth}{X} 他必受損傷和羞辱，他的羞恥不得消除。 \end{tabularx} \\ \\ \relax
6:34 & \begin{tabularx}{0.7\textwidth}{X} 丈夫因嫉恨發怒，報仇的時候絕不留情。 \end{tabularx} \\ \\ \relax
6:35 & \begin{tabularx}{0.7\textwidth}{X} 他不接受任何賠償，你送許多禮物，他也不肯和解。 \end{tabularx} \\ \\
[1ex]
\hline
\hline
\end{longtable}
$^{1}$《錄音》或《錄影》的版權是屬於基督教宣導會不各堂所有.
我們的屬靈傳統之一是以真理教導為本.
歡迎您將此《錄音》或《錄影》的連結傳給親友.
懇請注意,未經授權請勿轉錄或剪輯本《錄音》或《錄影》.
甚願好似《使徒行傳》二十章三十二節下所講.
這道能建立你們,叫你們和一切誠誠的人同得基業.
我們會詳盡講解,可以見到「真言對犯淫亂」這個課題看得很重要.
讓我們在聆聽聖經的話語之先,我們先一齊低頭有祈禱.
七因主,我們將到以下這個時間都交上給你.
願意你與我們同在,幫助我們去聆聽你的話.
求主你幫助我們,讓我們心門真的去打開.
能夠切切實實去聽到你的說話.
亦都在生活當中去進行,去謹守.
願主你幫助,願聖靈你同在.
我們這樣祈禱,交託奉主耶穌基督的名求,阿們.
都好想大家去打開同樣的聖經.
我們今日的經文是在「真言」的第六章.
經文其實比較長,經文的第六章和第七章.
大家如果有唐人聖經都可以去打開看一看.
一開始第一節講到.
「我兒,要遵守你父親的命令,不可離棄你母親的教誨」.
在六章二十節的經文.
其實是一個我們比較熟悉的宣誡.
大家如果一路看「真言」的時候.
會發覺好像有點熟悉.
原來開始的時候,在「真言」第一章.
已經有類似的經文.
但是你仔細去看的時候.
動詞是轉變了,去用了「遵守」.
去用了「遵守」這個字.
意思更加是去謹守.
要小心翼翼去遵行.
用這個詞代表程度提高了.
要我們留心去遵守誡命.
遵守甚麼的命令呢?.
其實整個大段落是在講犯淫亂的罪.
是婚姻以外的性關係.
要遵守就是不去行淫亂這個命令.
而其實這個詞「遵守」和第二十節,即後幾節.
你看到聖經那裡看到.

$^{41}$後幾節「保護」這個字其實是有接近的意思.
其實英文的翻譯可以用guard字這個字.
同一樣的字去形譯這兩個字.
這裡是有講要去保護你遠離邪惡的婦女.
原來經文是想連在一起去看.
即是當我們二十節那裡說.
我們切切實實去守住誡命的時候.
這個誡命,這個命令就去保守你.
去保護你遠離邪惡.
免你去禍患不用承擔罪的惡果.
很清楚其實神想我們去遵守誡命的最大目的.
最大的目的是想去保護,去保守我們.
就是這裡六章二十至二十四節成為一個段落.
但其實今天講到的經文.
在後面有同樣一個差不多的段落.
就是七章第一至五節.
這裡是另一個平衡的段落.
它都是以遵守神的誡命言語.
以誡命去保護你遠離這些惡事,惡人.
一頭一尾括著另外一個段落.
我們先去集中看看這兩個段落.
經文提到我們應該怎樣去對待這個誡命呢?.
有兩句都是提到.
我們要將這個誡命,就是神的教訓放在心上.
要牢記的,記實神的教訓誡命.
而另一句呢,它說刻在你的心板上.
意思都是類似的,背的,要去記的.
縱管純粹在腦裡面在回想著神的命令誡命.
神的誡命就可以不斷在心裡面,腦裡面出現.
就可以達到神要去達到的目的.
就是真正去保守,保護我們遠離邪惡.
其實現在這個世代,將一些東西記實,牢記.
已經不是一件很自然會去做的事了.
大家同不同意嗎?.
因為實在太方便了.
實在太方便,我自己都要悔改.
自從有網絡,甚至AI.
只要你按一個按鈕,東西就全部出來了.
你要找的資料,你要找的經文.
不同的版本,語言,余文,網上一應俱全.

$^{81}$自然而然,你去背誦經文,傳記在心裡.
就會少做了.
這樣的話,其實神怎能夠透過你所學過.
去默想過,你記得過的經文.
去在你重要的時刻,重要的關頭去提醒我們呢?.
所有這些我們即時可能在網上找到的經文.
其實仍然在外面的,不是屬於你的.
不是在我們的心裡面.
我們真的要做回應該做的.
不時背誦經文,重複地去看.
我們要刻意去做,傳記在心,聖言在心.
神知道我們不容易將這些誡命傳記在心上.
就指出一些外座的方式.
讓我們上上見到,不斷提醒我們.
他這裡說放在頸上,一直掛著它.
一直在提醒自己.
將誡命,將神的道掛在身上.
在面對試探的時候,你可能見到.
你可能手觸碰到,就提醒到了.
另一亦提到,放在指頭上.
可以是很早期,以色列人真的有一種穿戴.
是繞著手指提醒自己的穿戴.
但亦都可以像猶太人一樣,就好像這幅圖.
你見到他手上的前臂,是纏著一個叫做經盒.
這個經盒是綁到中指那裡.
就纏在手上的這條帶,你經常都會見到.
你戴著,經常都會見到.
常常在提醒自己.
其實現今的人都會用這些不同的方式.
可以去提醒到自己.
譬如教書店,店舖,有製作一些手腕.
有些繩帶,有寫一些經文,有些圖文.
是讓自己時時刻刻都見到.
我們的確是要常常被神的誡命不斷提醒.
這些誡命對於我們是有積極保護的作用.
你看看六章二十二節的經文.
他這裡說的是行走,是躺臥,是睡醒.
基本上說的是二十四小時,全天候的.
神要護著你.
我們所相信的神,我們知道祂不休息,不打頓.

$^{121}$祂就有能力保護我們.
很特別的,他這裡說的是當我們睡醒的時候.
祂必與你談論.
其實「致」的意思是大聲的,熱情的,與你談話.
是講到神的誡命,祂的智慧.
祂會與我們大聲熱情的去說.
是要讓我們聽到,讓我們可以遵從.
我想起今屆的港九培靈研經會,曾錦發牧師.
他在講道裡面講自己的一個例子.
他講到當時他太太在開車.
曾牧師就坐在他旁邊,正在睡覺.
但是就在那一刻,有一把很清晰,很響的聲音叫曾牧師起身.
然後叫太太要停車.
然後要停車.
終於趕得及,剎停了那輛車.
一輛大貨車緊緊的,在他們很近的位置.
在他這輛私家車面前,打橫全速的駛過.
神就用他的方法保護了,救了他們夫婦兩個.
我們未必在生活當中聽到神會很大聲保護.
但是我相信大家總有經歷過神特別去提醒.
神去保護你的時候.
我記得有一次,我自己本身都算是一個好脾氣的人.
但是那一次,有一樣事情激得我很生氣.
其實聖經說不可含怒到日落.
但是搞不定,整晚都睡不著.
就在那裡亂七八糟.
早上起床,其實已經想著要去找這個我不滿的人去理論.
但是早上看到一節的聖經.
都是箴言的經文.
它就說,遠離紛爭是人的尊榮.
如往人都愛爭吵.
我就不敢再去找這個人了.
神那一刻叫我要遠離.
我們經常以為,誡命是叫我們不要做一些事情.
好像學校的訓導老師不准我們做這樣,不准我們做那樣.
感覺是負面的.
但是其實經文這兩大段落講得很清楚.
禁止我們去做的原因.
叫我們遵守誡命的原因.
是去愛護我們.

$^{161}$神是去保護我們.
要保護我們遠離邪惡.
惡事,惡人.
雖然剛剛這兩段經文還沒講到很多人亂的事.
這麼亂的罪.
但其實整段都是圍繞著這個很重要的課題.
就是在夫妻關係以外的性關係.
甚至其他箴言都有提得多這個課題.
如果你有聖經,你看第五章.
它有講到.
今天這段經文第七章又是整章都有去講.
為什麼講這麼多呢?.
為什麼這麼重要呢?.
因為確確實實是和我們生命有關的事.
是和我們生命有關的事.
今天講到的題目.
管教的責備是生命的道.
在這裡,六章二十三節的經文.
管教原文的意思其實是一些教訓,誡命,指引.
責備這個字.
就是有去改正的意思.
即是先有一些誡命,有指引.
告訴你什麼是正確的道路.
但是有做得不好的時候.
就要去接受改正.
這樣為之生命的道.
為之生命的道.
是和生命有關,和性命有關的.
七章二節同樣講到.
遵守我的命令就是存活,就得存活.
要遵守命令.
不要和人發生在婚姻外的性關係.
否則結果會是死亡.
其實這個箴言都多次重申.
第二章就是這樣講到.
與有夫之婦犯淫亂可以陷入死亡.
第七章同樣有這樣講到.
所以說遵守神的誡命.
才可以去存活.
我們知道聖經講得很嚴重.

$^{201}$這種行淫亂的罪.
新約同樣看得很嚴重.
譬如我們所熟悉的前書的經文.
說淫行是得罪自己的身體.
這個是聖靈的電.
是神的電.
所以經文要我們遠避淫行.
要我們第一句就說要遠避淫行.
今天這段箴言同樣叫我們.
切切實實要留意,要避開.
但不是像剛剛多前書所講的.
要遠避淫行.
卻是更加提醒我們.
是要從一開始就避開.
從一開始就避開.
的確如果我們已經洗濕了頭.
已經被吸引了.
就不是那麼容易避開.
所以經文提醒我們.
一開始就避開.
不要自己深入了泥沼.
難以自拔.
譬如這個經文首先提到舌頭.
首先提到舌頭.
淺眉的舌頭.
由嘴滑舌.
可能我們原本不是有心去行淫亂.
但對方未必是這樣想.
對方可以是一心想纏著你.
尤其是現在大家常用電子通訊應用程式.
譬如WhatsApp等等.
對方很容易就能夠接觸到你.
甚至他可以在WhatsApp.
或者其他一些通訊的程式.
就跟你講很親密的話.
或者是一些淺眉討好你的說話.
引誘了你.
你就慢慢進入了這個網絡裡.
經文提醒我們要避開.
WhatsApp或者其他通訊方式.

$^{241}$都不是可以很隨便地去給人.
尤其是當我們越來越講到一些深入的話題.
或者不合宜的話題的時候.
就要懂得停.
要stop.
他提了說話.
然後他提到視覺.
就是女性的眉色.
眼皮.
經文特別提到要小心.
不要被對方的眼睛所吸引.
那個時期有些婦女.
是有做各樣眼部的化妝.
男士就這樣被吸引到了.
不單止是眼部.
這裡整體講女士的眉色.
讓男士看見動心.
動心這個字其實是什麼意思呢.
是講想佔據.
想據為己有的意思.
其實這個字就是和《出埃及記》裡面.
十誡裡面.
貪戀淪卸妻子.
這個貪戀這個字.
是同一樣的字.
見到.
望到.
貪戀的心就萌芽.
這字整體就是講及視覺.
就是講視覺.
其實現在這個世代.
比起舊的時代更加複雜.
更加複雜.
由於有互聯網.
很多不理想的資訊是氾濫.
你不是刻意避免的話.
很多這些引誘人的圖片.
影片就會入了眼目.
其實是影響你的思想.
你的定力.

$^{281}$可能對著真人的時候.
都更加容易跌倒.
雖然這段經文一直是講.
男士受誘惑犯罪.
但是現今這個世界.
其實女士.
同樣受到很多的引誘.
網絡上.
在廣告版上.
甚至有陌生人直接去找你.
要留心是從一開始.
從一開始.
經文提醒我們從一開始.
避免這些圖片.
視像.
其實這一兩年.
新聞我們都見到很多.
無論男或女.
被懲騙的個案都很多.
要小心被別人的說話.
跟你講的說話.
或者別人的外表.
去吸引了你.
記得經文說.
行淫亂會是生命有關的事.
是生命有關的事.
去到第26節.
繼續講有夫之婦獵取生命.
都是影響你的生命.
之後就說我們不要玩火.
我們不要玩火.
這裡講人若兜火在懷中.
他的衣服豈能不燒著呢.
人若走在火炭上.
他的腳豈能不燙傷呢.
基本上這個比喻.
是要講犯淫亂的嚴重後果.
是不可逃避的.
後果很嚴重.
你想想你在胸口捧著一堆火.

$^{321}$在胸口捧著一堆火.
衣服一定會燒著.
你走在炭火上.
你的腳一定會燙傷.
是在講必然性.
一定會發生.
接著經文亦都說.
與鄰舍的妻同寢.
亦都是這樣.
這樣講犯淫亂的.
就好像玩火.
是可以隨時奪去生命.
有生命的危險.
是走向死亡的道路.
亦都與神的懲罰有關.
凡親近他的難免受罰.
是涵蓋所有犯淫亂的人.
玩火的比喻.
那個後果是避免不了的.
而同樣神的刑罰.
同樣都是避免不了的.
弟兄姊妹犯淫亂.
是生命有關的事.
是玩火的事.
是會燒著的.
你想想如果你見到一堆火.
你都不會撲過去.
你都會走開.
由於這種罪惡的嚴重性.
經文就跟我們講.
要一開始就避開.
避開一些不合適的對話.
避開不合適的圖片.
視像等等.
不應該進入眼目的.
我們都將它避開.
不要讓自己陷入在網絡裡面.
接著這個段落就要講到.
承擔的種種不好的後果.
他先講到是賠償.

$^{361}$說賠是七倍.
必賠上家中一切的財物.
在舊約出埃及紀律法裡面.
其實有各種的律法.
是講到若果導致別人損失.
就要賠多少.
通常是講賠兩倍.
至到五倍不等.
這裡賠償是講要賠七倍.
其實不是一個實質的數字.
他想講一個完全的數.
七是一個完全的數.
想講的是賠償你的所有.
賠償所有.
如果用舊約的律法去講.
就是賠上你自己的田地,家當.
包括你自己的自由.
就直接成為了屬於其他人的奴隸.
一個犯淫亂的人.
不折不扣成為了一個欠債的人.
並且是永遠無法自己去清還這些債項.
這個明明是屬神的人.
現在他首先是欠了神.
剛才講到身體是聖靈的殿.
他怎樣可以去補償這些債項呢.
其次這個人也是欠自己配偶.
或者對方配偶的債.
自己或者行淫的對象.
在婚姻的約裡.
其實就是在神和神人面前所訂立的婚盟.
神是見證者.
而雙方是立約者.
納至在婚約裡面對對方忠誠.
行淫亂就是背棄了這份神聖的約定.
於是第三十四,三十五節.
就說到和他人的妻子行淫.
對方這個丈夫就會因責恨去發牢.
不留情地去報復去報仇.
其實我們都是在偶爾看到一些新聞.
都有報一些慘劇.

$^{401}$一個人被罪之後.
可能是丈夫或者太太出軌之後.
這個人被罪.
就引發更大的罪惡.
他就去報復.
心態是不會顧及後果.
在現代社會我們都屢見不善.
而經文最後這句話.
你送上許多的禮物.
他亦都不肯和解.
這裡的禮物是各樣不同形式的賠償.
其實禮物這個字更加是一個法庭常用的詞.
即是你在法庭上聽到一個很大的賠償金額.
但是都不能解決.
不會和解.
對方不接受.
行淫亂作是嚴重地得罪了自己.
或者配偶的行為的一個罪惡.
經文亦都說這個行為是得罪自己.
得罪社群的行為.
得罪自己.
他說這個行為其實是無知的.
是毀了自己.
無知這個字基本意思就是失去了理智.
沒有用腦去思考.
其實根據剛才經文所說.
行淫亂的.
他先會成為一個很欠債的人.
欠神的債項.
欠其他人的債項.
我們將來亦都要面見神.
行淫亂亦都會惹來對方不留情的報復.
而且亦都會在社群當中不被接納.
充滿羞辱.
所以經文說去行這件事的時候.
其實基本上失去了理智.
毀了自己.
其實我們看回現實的情況.
我們都知道的確有些犯淫亂的人.
家人不再原諒他.

$^{441}$子女不再理會他.
甚至有失去了工作.
當你以為和淫亂的對象可以繼續在一起的時候.
很多時犯淫亂的最終都不是和對象會再在一起.
所以說犯淫亂罪的人可以說.
是因為這件事就真真正正賠上了一切.
很多是之前辛辛苦苦經營的.
無論是自己的家庭.
無論是他的事業.
是他的明星等等.
就這樣賠上了他.
有很多例子成為前車之鑑.
所以可以說這樣去做的.
其實經文說真的好像喪失了理性.
沒有經大腦去思考.
沒有顧後果的一些行為.
經文說他怎樣去毀掉自己.
他更加說會受損傷.
受羞辱.
這裡的損傷可以是身體被打.
剛才說對方的配偶可以不留情面的報復.
而損傷也可以是遭受神的懲罰.
實際上如果在舊約和以色列女子行淫.
是會被判死刑的.
是會被判死刑.
現在很多地方沒有這種法例.
但聖經也有說過神那種好像直接的去懲罰.
災禍臨到,疾病等等.
聖經強調神是會審判人的神.
犯淫亂的罪可能現時未受到審判.
但始終是要面對審判.
聖經提醒我們不要存僥倖的心態.
神祂是知道一切的神.
最終我們都要見祂.
而且這裡說同樣是一個羞辱.
犯淫亂的人在神面前的罪行是一個羞辱.
而且在社群當中也不被接納.
犯淫亂的惡名也會一直被記住.
一生都有這個污名.
在人面前的羞辱.

$^{481}$在神面前的羞辱.
其實看了這麼多犯淫亂的人.
其實他多方的得罪.
他得罪了神.
得罪了對方的配偶或者自己的配偶.
得罪了社群.
更加他得罪了自己.
就是會有嚴重的後果.
而且這個後果會一直存在.
經文很多方面講到的時候.
其實就是叫我們不要陷入犯淫亂的罪裡面.
要我們時時刻刻去作警惕.
去提醒我們.
但如果我們有弟兄姊妹已經在這種罪裡面.
那怎麼樣好呢?.
按著聖經的真理並不是沒有出路.
天父說要我們謹記聖經裡面提到我們新約約翰一書.
一章九節的經文.
我們若認自己的罪神是信實的是公義的.
必要赦免我們的罪洗清我們一切的不義.
重點是什麼呢?.
先是知罪認罪和真正脫離這個罪.
神按著他的恩慈憐憫就會去赦免.
如果我們有弟兄姊妹在這個罪裡面的時候.
求神去光照去知罪去認罪去悔改.
正如今天講到的題目.
管教的責備是生命的道.
責備就是去改正.
真正去改正這個才是生命的道.
才能夠存活.
如果有需要的弟兄姊妹.
也可以找相熟的教母.
即是大家相熟的教母.
聖經清楚告訴我們.
行淫亂是有各種的惡果.
真正是毀了自己也毀了他人.
求神去讓跟隨他的人都能夠自己自潔.
不參與在其中不用承擔各種罪的損傷,羞辱.
弟兄姊妹.
今日的真言給到我們清晰的提醒.

$^{521}$神要我們遵守誡命.
其實是於我們有益.
神是想保護我們遠離這些邪惡.
出發點是想愛護我們.
他想保護我們.
而這是一件跟我們生命有關的事.
正如見到火沒理由撲過去.
神提醒我們及早避開.
包括跟人不恰當的對話.
包括一些不良的資訊,訊息,圖片,視像.
最後犯淫亂的確是有很多的惡果.
亦遭神的審判懲罰.
願神幫助我們都不要讓自己在其中.
若然有犯錯的亦都及時悔改.
去得到神的赦免.
求神憐憫幫助我們遠離這些罪惡.
我們低頭一起有祈禱.
參與天賦上帝我們將到我們這裡各人都交上給你.
願主你讓到我們時刻將你自己的誡命.
你的教導,教訓傳記在心.
以至於我們可以常常記得你提醒我們的地方.
你的目的是想保護我們,保守我們.
免除邪惡,免除惡果.
求主你也讓到我們去謹慎,警醒.
去避開一切這些隱憂,誘惑.
求主你讓到我們能夠在適當的時刻.
真的能夠即刻的去停止.
願主你去與我們同在.
常常去看顧著我們.
亦都願主你去幫助.
若然我們有當中在這些罪惡裡.
你的憐憫,你的愛去吸引我們去迴轉.
以至於我們能夠去悔改,歸向你.
去得心,去走在生命的道上.
主我們這樣多謝你.
願主你常與我們同在.
我們這樣祈禱交託.
奉耶穌基督的名求.
\newpage



\section{撒母耳記上 18:1-20:42}
\label{sec:031785}
\textbf{多代同行 (2) PeerWalk (1) Godly Friends, Godly Wisdom}~~~~李觀南傳道
\newline
\newline
連結: \href{http://www.church.com.hk/acms/content.asp?site=cdc\&op=show\&type=product\&code=031785\&layout=sermon}{\texttt{031785}} ~~~~ 語音日期: 2026-03-07
\newline
\newline
\hyperref[sec:031784]{\small{< < < PREV SERMON < < <}}
~
\hyperref[sec:index_chronic]{\small{[返順時目]}}
~
\hyperref[sec:031786]{\small{> > > NEXT SERMON > > >}}
\newline
\newline
撒母耳記上 18:1-20:42
\newline
\begin{longtable}{cl}
\hline
\hline
章節 & 經文 (和合本修訂版)\\
\hline
18:1 & \begin{tabularx}{0.7\textwidth}{X} 大衛對掃羅說完了話，約拿單的心與大衛的心深相契合。約拿單愛大衛，如同愛自己的性命。 \end{tabularx} \\ \\ \relax
18:2 & \begin{tabularx}{0.7\textwidth}{X} 那日掃羅留住大衛，不讓他回父家。 \end{tabularx} \\ \\ \relax
18:3 & \begin{tabularx}{0.7\textwidth}{X} 約拿單愛大衛如同愛自己的性命，就與他立約。 \end{tabularx} \\ \\ \relax
18:4 & \begin{tabularx}{0.7\textwidth}{X} 約拿單從身上脫下外袍，給了大衛，又把戰衣、刀、弓、腰帶都給了他。 \end{tabularx} \\ \\ \relax
18:5 & \begin{tabularx}{0.7\textwidth}{X} 掃羅無論差遣大衛往何處去，他都做事精明。掃羅立他作軍隊的指揮官，眾百姓和掃羅的臣僕都看為美。 \end{tabularx} \\ \\ \relax
18:6 & \begin{tabularx}{0.7\textwidth}{X} 大衛打死了那非利士人，同眾人回來的時候，婦女們從以色列各城裡出來，歡歡喜喜，打鼓奏樂，唱歌跳舞，迎接掃羅王。 \end{tabularx} \\ \\ \relax
18:7 & \begin{tabularx}{0.7\textwidth}{X} 眾婦女歡樂唱和，說：「掃羅殺死千千，大衛殺死萬萬。」 \end{tabularx} \\ \\ \relax
18:8 & \begin{tabularx}{0.7\textwidth}{X} 掃羅非常憤怒，不喜歡這話。他說：「將萬萬歸給大衛，千千歸給我，只剩下王國沒有給他！」 \end{tabularx} \\ \\ \relax
18:9 & \begin{tabularx}{0.7\textwidth}{X} 從這日起，掃羅就敵視大衛。 \end{tabularx} \\ \\ \relax
18:10 & \begin{tabularx}{0.7\textwidth}{X} 次日，從神來的邪靈緊抓住掃羅，他就在家中胡言亂語。大衛照常彈琴，掃羅手裡拿著槍。 \end{tabularx} \\ \\ \relax
18:11 & \begin{tabularx}{0.7\textwidth}{X} 掃羅把槍一擲，心裡說：「我要將大衛刺透，釘在牆上。」大衛閃避了他兩次。 \end{tabularx} \\ \\ \relax
18:12 & \begin{tabularx}{0.7\textwidth}{X} 掃羅懼怕大衛，因為耶和華離開自己，與大衛同在。 \end{tabularx} \\ \\ \relax
18:13 & \begin{tabularx}{0.7\textwidth}{X} 所以掃羅叫大衛離開自己，立他為千夫長，他就領兵出入。 \end{tabularx} \\ \\ \relax
18:14 & \begin{tabularx}{0.7\textwidth}{X} 大衛所做的每一件事都精明，耶和華也與他同在。 \end{tabularx} \\ \\ \relax
18:15 & \begin{tabularx}{0.7\textwidth}{X} 掃羅見大衛做事精明，就更怕他。 \end{tabularx} \\ \\ \relax
18:16 & \begin{tabularx}{0.7\textwidth}{X} 但以色列和猶大眾人都愛大衛，因為他領他們出入。 \end{tabularx} \\ \\ \relax
18:17 & \begin{tabularx}{0.7\textwidth}{X} 掃羅對大衛說：「看哪，我將大女兒米拉嫁給你，只要你作我的勇士，為耶和華爭戰。」掃羅心裡說：「我不好親手害他，要藉非利士人的手害他。」 \end{tabularx} \\ \\ \relax
18:18 & \begin{tabularx}{0.7\textwidth}{X} 大衛對掃羅說：「我是誰，我是甚麼出身，我父家在以色列中算甚麼，豈敢作王的女婿呢？」 \end{tabularx} \\ \\ \relax
18:19 & \begin{tabularx}{0.7\textwidth}{X} 掃羅的女兒米拉到了當嫁給大衛的時候，掃羅卻將她嫁給了米何拉人亞得列。 \end{tabularx} \\ \\ \relax
18:20 & \begin{tabularx}{0.7\textwidth}{X} 掃羅的女兒米甲愛大衛。有人告訴掃羅，這件事在掃羅眼中看為合宜。 \end{tabularx} \\ \\ \relax
18:21 & \begin{tabularx}{0.7\textwidth}{X} 掃羅心裡說：「我把這女兒嫁給大衛，作他的圈套，好藉非利士人的手害他。」所以掃羅第二次對大衛說：「你今日可以作我的女婿。」 \end{tabularx} \\ \\ \relax
18:22 & \begin{tabularx}{0.7\textwidth}{X} 掃羅吩咐臣僕：「你們暗中對大衛說：『看哪，王喜歡你，王的臣僕也都愛戴你，現在你就作王的女婿吧。』」 \end{tabularx} \\ \\ \relax
18:23 & \begin{tabularx}{0.7\textwidth}{X} 掃羅的臣僕照這話說給大衛聽。大衛說：「你們把作王的女婿看為小事嗎？我是貧窮卑微的人。」 \end{tabularx} \\ \\ \relax
18:24 & \begin{tabularx}{0.7\textwidth}{X} 掃羅的臣僕回奏說，大衛說了這樣的話。 \end{tabularx} \\ \\ \relax
18:25 & \begin{tabularx}{0.7\textwidth}{X} 掃羅說：「你們要對大衛這樣說：『王不要甚麼聘禮，只要一百非利士人的包皮，好在王的仇敵身上報仇。』」掃羅的意圖是要大衛落在非利士人的手中。 \end{tabularx} \\ \\ \relax
18:26 & \begin{tabularx}{0.7\textwidth}{X} 掃羅的臣僕把這話告訴大衛，大衛就歡喜作王的女婿。日期還沒有到， \end{tabularx} \\ \\ \relax
18:27 & \begin{tabularx}{0.7\textwidth}{X} 大衛和跟隨他的人起來前往，殺了二百非利士人，將包皮足數交給王，為要作王的女婿。於是掃羅將女兒米甲嫁給大衛。 \end{tabularx} \\ \\ \relax
18:28 & \begin{tabularx}{0.7\textwidth}{X} 掃羅見耶和華與大衛同在，女兒米甲又愛大衛， \end{tabularx} \\ \\ \relax
18:29 & \begin{tabularx}{0.7\textwidth}{X} 就更怕大衛，常常與大衛為敵。 \end{tabularx} \\ \\ \relax
18:30 & \begin{tabularx}{0.7\textwidth}{X} 每逢非利士的軍官出來打仗，大衛做事比掃羅任何臣僕更精明，因此他的名極受尊重。 \end{tabularx} \\ \\ \relax
19:1 & \begin{tabularx}{0.7\textwidth}{X} 掃羅吩咐他兒子約拿單和眾臣僕要殺大衛，但掃羅的兒子約拿單卻很喜愛大衛。 \end{tabularx} \\ \\ \relax
19:2 & \begin{tabularx}{0.7\textwidth}{X} 約拿單告訴大衛說：「我父掃羅想要殺你，現在你要小心，明日早晨留在一個僻靜的地方藏起來。 \end{tabularx} \\ \\ \relax
19:3 & \begin{tabularx}{0.7\textwidth}{X} 我會出去，到你所藏的田裡，站在我父親旁邊，與父親談論到你。我看情形怎樣，會告訴你。」 \end{tabularx} \\ \\ \relax
19:4 & \begin{tabularx}{0.7\textwidth}{X} 約拿單向他父親掃羅說大衛的好話，對他說：「王不可得罪王的僕人大衛，因為他未曾得罪你，他所行的對你都很有益處。 \end{tabularx} \\ \\ \relax
19:5 & \begin{tabularx}{0.7\textwidth}{X} 他拚了命殺那非利士人，並且耶和華為全以色列大施拯救。那時你看見，也很歡喜，現在為何要犯罪，流無辜人的血，無緣無故殺大衛呢？」 \end{tabularx} \\ \\ \relax
19:6 & \begin{tabularx}{0.7\textwidth}{X} 掃羅聽了約拿單的話，就指著永生的耶和華起誓：「我絕不殺他。」 \end{tabularx} \\ \\ \relax
19:7 & \begin{tabularx}{0.7\textwidth}{X} 約拿單叫大衛來，把這一切事告訴他。約拿單帶他去見掃羅，他就像以前一樣侍立在掃羅面前。 \end{tabularx} \\ \\ \relax
19:8 & \begin{tabularx}{0.7\textwidth}{X} 此後又有戰爭，大衛出去與非利士人打仗。他大大擊敗他們，他們就在他面前逃跑。 \end{tabularx} \\ \\ \relax
19:9 & \begin{tabularx}{0.7\textwidth}{X} 從耶和華來的邪靈又降在掃羅身上，掃羅手裡拿槍坐在屋裡，大衛正用手彈琴。 \end{tabularx} \\ \\ \relax
19:10 & \begin{tabularx}{0.7\textwidth}{X} 掃羅想要用槍刺透大衛，把他釘在牆上，他卻躲開掃羅，掃羅的槍刺入牆內。當夜大衛逃走，躲起來了。 \end{tabularx} \\ \\ \relax
19:11 & \begin{tabularx}{0.7\textwidth}{X} 掃羅派一些使者到大衛的房屋那裡守著他，等到天亮要殺他。大衛的妻子米甲對大衛說：「你今夜若不逃命，明日就要被殺。」 \end{tabularx} \\ \\ \relax
19:12 & \begin{tabularx}{0.7\textwidth}{X} 於是米甲將大衛從窗戶縋下去，讓他走；大衛就逃走，躲起來了。 \end{tabularx} \\ \\ \relax
19:13 & \begin{tabularx}{0.7\textwidth}{X} 米甲把家中的神像放在床上，頭枕在山羊毛的枕頭上，用衣服蓋起來。 \end{tabularx} \\ \\ \relax
19:14 & \begin{tabularx}{0.7\textwidth}{X} 掃羅派一些使者去捉拿大衛，米甲說：「他病了。」 \end{tabularx} \\ \\ \relax
19:15 & \begin{tabularx}{0.7\textwidth}{X} 掃羅又派一些使者去看大衛，說：「把他連床一起抬到我這裡，我好殺他。」 \end{tabularx} \\ \\ \relax
19:16 & \begin{tabularx}{0.7\textwidth}{X} 使者進去，看哪，神像在床上，頭枕在山羊毛的枕頭上。 \end{tabularx} \\ \\ \relax
19:17 & \begin{tabularx}{0.7\textwidth}{X} 掃羅對米甲說：「你為甚麼這樣欺騙我，放我仇敵逃走呢？」米甲對掃羅說：「他對我說：『你放我走吧，我何必要殺你呢？』」 \end{tabularx} \\ \\ \relax
19:18 & \begin{tabularx}{0.7\textwidth}{X} 大衛逃跑躲避，來到拉瑪的撒母耳那裡，把掃羅向他所行的事全告訴他。他和撒母耳就去，住在拿約。 \end{tabularx} \\ \\ \relax
19:19 & \begin{tabularx}{0.7\textwidth}{X} 有人告訴掃羅說：「看哪，大衛在拉瑪的拿約」。 \end{tabularx} \\ \\ \relax
19:20 & \begin{tabularx}{0.7\textwidth}{X} 掃羅派一些使者去捉拿大衛。去的人見一隊先知受感說話，撒母耳站在當中領導他們，掃羅派去的使者也受神的靈感動說話。 \end{tabularx} \\ \\ \relax
19:21 & \begin{tabularx}{0.7\textwidth}{X} 有人把這事告訴掃羅，他又派另一些使者去，他們也受感說話。掃羅第三次派使者去，他們也受感說話。 \end{tabularx} \\ \\ \relax
19:22 & \begin{tabularx}{0.7\textwidth}{X} 然後掃羅親自往拉瑪去，到了西沽的大井，問人說：「撒母耳和大衛在哪裡？」有人說：「看哪，在拉瑪的拿約。」 \end{tabularx} \\ \\ \relax
19:23 & \begin{tabularx}{0.7\textwidth}{X} 他就往那裡去，到了拉瑪的拿約。神的靈也臨到他，他一面走一面受感說話，直到拉瑪的拿約。 \end{tabularx} \\ \\ \relax
19:24 & \begin{tabularx}{0.7\textwidth}{X} 他也脫了衣服，也在撒母耳面前受感說話，一日一夜赤身躺臥。因此有人說：「掃羅也在先知中嗎？」 \end{tabularx} \\ \\ \relax
20:1 & \begin{tabularx}{0.7\textwidth}{X} 大衛從拉瑪的拿約逃跑，來到約拿單面前，對他說：「我做了甚麼，有甚麼罪孽，在你父親面前犯了甚麼罪，他竟要尋索我的性命呢？」 \end{tabularx} \\ \\ \relax
20:2 & \begin{tabularx}{0.7\textwidth}{X} 約拿單對他說：「絕無此事！你必不至於死。看哪，我父做事，無論大小，沒有不告訴我的。我父親為甚麼要隱瞞我這件事呢？不會這樣的！」 \end{tabularx} \\ \\ \relax
20:3 & \begin{tabularx}{0.7\textwidth}{X} 大衛又起誓說：「你父親確實知道我在你眼前蒙恩。所以他說，『這事不要讓約拿單知道，免得他愁煩。』我指著永生的耶和華起誓，又指著你的性命起誓，我離死只差一步而已。」 \end{tabularx} \\ \\ \relax
20:4 & \begin{tabularx}{0.7\textwidth}{X} 約拿單對大衛說：「你心裡所求的，我必為你成就。」 \end{tabularx} \\ \\ \relax
20:5 & \begin{tabularx}{0.7\textwidth}{X} 大衛對約拿單說：「看哪，明日是初一，我必須與王同席用餐，求你讓我去藏在田野，直到第三日傍晚。 \end{tabularx} \\ \\ \relax
20:6 & \begin{tabularx}{0.7\textwidth}{X} 你父親若見我不在席上，你就說：『大衛懇求我允許他趕回本城伯利恆去，因為他全家在那裡獻年祭。』 \end{tabularx} \\ \\ \relax
20:7 & \begin{tabularx}{0.7\textwidth}{X} 你父親若說好，你的僕人就平安了；他若大怒，你就知道他決意行惡。 \end{tabularx} \\ \\ \relax
20:8 & \begin{tabularx}{0.7\textwidth}{X} 求你施恩於僕人，因你在耶和華面前曾與僕人立約。我若有罪孽，你就親自殺死我，何必把我交給你父親呢？」 \end{tabularx} \\ \\ \relax
20:9 & \begin{tabularx}{0.7\textwidth}{X} 約拿單說：「絕無此事！我若確實知道我父親決意害你，怎麼會不告訴你呢？」 \end{tabularx} \\ \\ \relax
20:10 & \begin{tabularx}{0.7\textwidth}{X} 大衛對約拿單說：「你父親若嚴厲回答你，誰來告訴我呢？」 \end{tabularx} \\ \\ \relax
20:11 & \begin{tabularx}{0.7\textwidth}{X} 約拿單對大衛說：「來，讓我們到田野去。」二人就往田野去了。 \end{tabularx} \\ \\ \relax
20:12 & \begin{tabularx}{0.7\textwidth}{X} 約拿單對大衛說：「願耶和華－以色列的神作證。明日約在這時候，或第三日，我一探出我父親的心意，看哪，若對大衛是好意，我怎麼會不派人來告訴你呢？ \end{tabularx} \\ \\ \relax
20:13 & \begin{tabularx}{0.7\textwidth}{X} 我父親若有意害你，而我不告訴你，送你平安地離開，願耶和華重重懲罰約拿單。願耶和華與你同在，如同從前與我父親同在一樣。 \end{tabularx} \\ \\ \relax
20:14 & \begin{tabularx}{0.7\textwidth}{X} 你要照耶和華的慈愛恩待我，不但我活著的時候免我死亡， \end{tabularx} \\ \\ \relax
20:15 & \begin{tabularx}{0.7\textwidth}{X} 就是耶和華從地面逐一剪除大衛仇敵的時候，你也永不可向我家斷絕恩惠。」 \end{tabularx} \\ \\ \relax
20:16 & \begin{tabularx}{0.7\textwidth}{X} 於是約拿單與大衛家立約：「願耶和華從大衛仇敵的手來追討。」 \end{tabularx} \\ \\ \relax
20:17 & \begin{tabularx}{0.7\textwidth}{X} 約拿單因愛大衛如同愛自己的性命，就叫他再起誓。 \end{tabularx} \\ \\ \relax
20:18 & \begin{tabularx}{0.7\textwidth}{X} 約拿單對他說：「明日是初一，你的座位空著，人必察覺你不在。 \end{tabularx} \\ \\ \relax
20:19 & \begin{tabularx}{0.7\textwidth}{X} 到第三日，就要走一段長路下去，去到你遇事那天所藏的地方，在以色磐石的旁邊等候。 \end{tabularx} \\ \\ \relax
20:20 & \begin{tabularx}{0.7\textwidth}{X} 我會向磐石旁邊射三箭，如同射箭靶一樣。 \end{tabularx} \\ \\ \relax
20:21 & \begin{tabularx}{0.7\textwidth}{X} 看哪，我會派僮僕，說：『去把箭找來。』我若對僮僕喊說：『看哪，箭在你的這邊，把箭拿來』，你就可以平安回來；我指著永生的耶和華起誓，你一定沒有事。 \end{tabularx} \\ \\ \relax
20:22 & \begin{tabularx}{0.7\textwidth}{X} 我若對孩子說：『看哪，箭在你的前方』，你就要離開，因為是耶和華差你去的。 \end{tabularx} \\ \\ \relax
20:23 & \begin{tabularx}{0.7\textwidth}{X} 至於你和我，我們所說的話，看哪，耶和華在你我中間作證，直到永遠。」 \end{tabularx} \\ \\ \relax
20:24 & \begin{tabularx}{0.7\textwidth}{X} 大衛就去藏在田野。到了初一，王要坐席用餐。 \end{tabularx} \\ \\ \relax
20:25 & \begin{tabularx}{0.7\textwidth}{X} 王照常坐在靠牆的位子上，約拿單在對面，押尼珥坐在掃羅旁邊，大衛的座位卻是空的。 \end{tabularx} \\ \\ \relax
20:26 & \begin{tabularx}{0.7\textwidth}{X} 這日掃羅沒有說甚麼，因為他說：「大衛或許有事，偶染不潔，還未得潔淨。」 \end{tabularx} \\ \\ \relax
20:27 & \begin{tabularx}{0.7\textwidth}{X} 初二，大衛的座位還空著。掃羅對他兒子約拿單說：「耶西的兒子為何昨日、今日都沒有來用餐呢？」 \end{tabularx} \\ \\ \relax
20:28 & \begin{tabularx}{0.7\textwidth}{X} 約拿單回答掃羅說：「大衛懇求我允許他回伯利恆去， \end{tabularx} \\ \\ \relax
20:29 & \begin{tabularx}{0.7\textwidth}{X} 說：『求你讓我去，因為我家在城裡有獻祭的事，我哥哥吩咐我去。如今我若在你眼前蒙恩，求你讓我去見我的兄弟。』所以大衛沒有來赴王的筵席。」 \end{tabularx} \\ \\ \relax
20:30 & \begin{tabularx}{0.7\textwidth}{X} 掃羅向約拿單怒氣大發，對他說：「你這頑梗悖逆之婦人所生的，我怎麼會不知道你選擇耶西的兒子，自取羞辱，也使你母親露體蒙羞呢？ \end{tabularx} \\ \\ \relax
20:31 & \begin{tabularx}{0.7\textwidth}{X} 只要耶西的兒子還活在世上一天，你和你的國必保不住。現在你要派人去，把他帶到我這裡來，因為他是該死的。」 \end{tabularx} \\ \\ \relax
20:32 & \begin{tabularx}{0.7\textwidth}{X} 約拿單回答父親掃羅說：「他為甚麼該死呢？他做了甚麼呢？」 \end{tabularx} \\ \\ \relax
20:33 & \begin{tabularx}{0.7\textwidth}{X} 掃羅向約拿單擲槍要刺他，約拿單就知道他父親決意要殺死大衛。 \end{tabularx} \\ \\ \relax
20:34 & \begin{tabularx}{0.7\textwidth}{X} 於是約拿單氣憤憤地從席上起來。他在初二這天沒有吃飯，因為他為大衛愁煩，又因為他父親羞辱了他。 \end{tabularx} \\ \\ \relax
20:35 & \begin{tabularx}{0.7\textwidth}{X} 次日早晨，約拿單按著與大衛約定的時候到田野去，有一個小僮僕跟隨他。 \end{tabularx} \\ \\ \relax
20:36 & \begin{tabularx}{0.7\textwidth}{X} 約拿單對僮僕說：「你跑去把我所射的箭找來。」僮僕跑去，約拿單就把箭射在僮僕的前方。 \end{tabularx} \\ \\ \relax
20:37 & \begin{tabularx}{0.7\textwidth}{X} 僮僕到了約拿單落箭之地，約拿單呼叫僮僕說：「箭不是在你的前方嗎？」 \end{tabularx} \\ \\ \relax
20:38 & \begin{tabularx}{0.7\textwidth}{X} 約拿單又呼叫僮僕說：「快去，不要站在那裡！」僮僕就撿起箭來，回到主人那裡。 \end{tabularx} \\ \\ \relax
20:39 & \begin{tabularx}{0.7\textwidth}{X} 僮僕不知道這是甚麼意思，只有約拿單和大衛知道這事。 \end{tabularx} \\ \\ \relax
20:40 & \begin{tabularx}{0.7\textwidth}{X} 約拿單把他的弓箭交給僮僕，吩咐他說：「你拿到城裡去。」 \end{tabularx} \\ \\ \relax
20:41 & \begin{tabularx}{0.7\textwidth}{X} 僮僕一去，大衛就從南邊出來，俯伏在地，拜了三拜。他們彼此親吻，一起哭泣，大衛哭得更悲哀。 \end{tabularx} \\ \\ \relax
20:42 & \begin{tabularx}{0.7\textwidth}{X} 約拿單對大衛說：「你平平安安地去吧！因為我們二人曾指著耶和華的名起誓說：『願耶和華在你我中間，以及你我後裔中間作證，直到永遠。』」大衛就起身走了，約拿單也回城裡去了。 \end{tabularx} \\ \\
[1ex]
\hline
\hline
\end{longtable}
$^{1}$好,頂指滿滿.
好,歡迎大家回來崇拜.
事隔了一個星期.
一個星期之前.
再一個星期之前.
我剛剛講完上一篇.
又講另一篇了.
今天來到一個新的月份.
即是三月.
我們又來到一個新的系列.
這個系列其實是教會今年全年的主題.
我們講多代同行.
如果大家有印象.
我們之前在一月份.
應該是講了第一個多代同行的系列.
那個系列主要是講與神同行.
這一次是講如何與同輩同行.
所以我將這個系列的名字改了一個名字.
改了叫做多代同行二支.
Peer Walk.
其實是怎麼來的呢.
就是我有一點點認識這個Peer Talk這個字.
因為我在想的時候.
我就會想一般人的情況下.
如果兩個人會有Peer Talk這件事的話.
通常這兩個人的關係都是較為密切的.
可能是你的夫妻關係.
又或者可能是你的閨密.
你就會跟他有Peer Talk這件事.
通常都是有一些較為深厚的友誼在裡面.
兩個人呢.
通常基本上都應該是會有一個較為.
比較信任對方的一個關係在裡面.
所以當我們以今一個系列.
我們說我們要講學如何與同輩.
與朋友去同行的時候.
同樣都是在講我們如何可以與他們.
建立一些更深厚的關係.
以致我們與人的關係.
特別是在我們的群體之間.

$^{41}$我們的關係不只是留於一個表面的關係.
而是可以真真正正地做到.
我們能夠彼此同行這件事.
一起去走這個信耶穌的這條路.
好,那講到做人的同行者.
又或者我改一改這個字眼.
同行者就好像很荒謬.
有些文言文.
寫一些字書面語應該這麼說.
我嘗試改一改.
我說不如講.
講如何做人一個真誠的朋友這件事.
在我預備這篇信息裡面.
講到要做這件事的時候.
我想起我之前.
在大概我沒記錯是二三年年底.
我曾經講過一篇道.
提過一件事.
我不知道大家有沒有印象.
就是我說.
其實朋友這個字眼.
在今天這個社會裡面.
其實越來越模糊.
因為科技的發達.
我們結識一些所謂的新朋友.
其實好像是一件很容易的事.
大家只需要在一個社交平台裡面按個鍵.
你已經可以有一個新朋友了.
而在那篇道裡面.
我還記得我曾經建議過大家回去看看.
你看看你自己的社交平台上的朋友清單.
究竟當中其實有多少個.
你會真的稱呼為你一個真正的朋友呢.
我不知道大家有沒有回去做這件事.
但是無論如何.
當我們今天說要學.
如何成為別人的一個真誠的朋友的時候.
成為別人的一個同行者的時候.
或者我會想其實重點.
不是在別人身上.

$^{81}$因為別人是怎樣.
其實我們並不能夠控制這麼多.
也就是說就算他肯.
因著你們的友誼去作出一些改變.
都不是一時三刻就可以見到的一件事.
反而我會覺得更加著重的那個位置.
理應是在我們自己裡面.
究竟我們自己有沒有做好自己的本分呢.
我們有沒有令到那些人覺得.
我們其實不是那些社交平台裡面的.
那個朋友清單的其中一個.
那就算了.
而是覺得你是一個真的值得相交的好朋友呢.
那麼盼望我們今天一起去看看.
有些什麼聖經去教我們.
開始之先我們一起去祈禱.
天父上帝多謝你.
多謝你透過你聖經裡面的說話.
有很多不同的故事.
不同的教導在裡面.
讓我們知道究竟.
當我們在這個世上要學著你.
去成為地上的天聖.
與人同行的時候.
我們有很多東西可以參考.
就求神你今天藉著我們一起去看聖經裡面.
我們能夠明白這一切的事情更多.
求主你親自的去到我們當中.
你的聖靈臨在教導我們.
打開我們每個人的心.
讓我們明白你的心意.
聽我們禱告教託.
奉耶穌基督的聖名臨自祈求.
Amen.
好那麼今天想介紹.
上一整篇道.
最主要想介紹給你們認識的.
就是其中一個的聖經人物.
他的名字叫做約拿丹.
那麼聖經裡面其實記載約拿丹這個人的事.

$^{121}$其實不是很多.
就是經文不算很多.
大概其實都是只有四張經文.
那麼要看的話其實大家都很快就可以看完.
不過我今天沒可能跟大家讀四張經文.
在這裡讀就好像太長了.
讀完都可能夠時間了.
所以為了節省時間.
今天跟平時有一點不同.
就是我不預大家要一起讀經文.
反而是我會在之後的講解裡面.
我們跟大家一起高呼一下.
究竟約拿丹在聖經裡面的不同事蹟.
看看約拿丹這個人有什麼特別值得我們學.
而為免我們當中有人完全不認識約拿丹這個人物.
你看到我寫了一些簡介出來.
是約拿丹的一些小小的背景.
那麼約拿丹的英文就是Jonathan.
而他希伯來語的一個意思.
就是這個名字的意思.
就是耶和華賜予.
就是給耶和華所賜的一件事.
那麼他是以色列人的王.
以色列第一個皇帝.
蘇羅的長子.
而他很有趣的是第二個皇帝的大衛的朋友.
那麼約拿丹也是一個很深得民心的人.
在一次的戰爭裡面.
根據聖經記載的一個內容.
其實當時以色列人跟一幫菲律賓人打仗.
在過程裡面.
其實蘇羅曾經要求過那幫以色列的百姓去起誓.
他就說.
如果我們不等到晚上.
向那幫菲律賓人報仇.
就吃些什麼的話.
一定要受到詛咒.
那麼約拿丹就因為當蘇羅要求以色列百姓去發誓的時候.
他不在場.
他不知道.

$^{161}$所以在他在森林的過程裡面.
他用了他的匠.
他就點了一些蜂蜜去吃了.
結果蘇羅因為自己起過這個誓.
那麼他真的覺得需要殺了約拿丹.
不過當時因為約拿丹在一幫百姓裡面已經有些功績了.
已經覺得他有些名聲了.
所以這幫百姓很主動地為他求情.
而這個求情也真的使蘇羅覺得.
那好吧.
那就不用他死了.
所以我剛剛說為什麼說約拿丹心得百姓就是這個意思.
而整個故事記載.
剛剛說了.
約拿丹的故事其實就是十四章到十八至二十章.
大概是這四章的經文.
主要是集中在《撒慕爾記》上.
那麼為什麼我會選用約拿丹來分享同輩同行的這個話題呢.
就是因為.
其實大衛與約拿丹他們之間的友誼.
是在整個舊約的聖經裡面.
其中一段最偉大的友誼.
甚至他們之間的感情.
是說到超越了一般男女之間的愛情.
這個說法不是我說的.
而是聖經很實在地記載了.
在《撒慕爾記》的一章二十六至二十七節.
當這一句經文.
大家看到剛剛說到《撒慕爾記》上.
而這裡是《撒慕爾記》下.
說的是當時大衛知道約拿丹的生命最後是戰死的時候.
大衛就曾經用一首哀歌去紀念他的老朋友.
就是大家看到的經文.
很多人讀這一段經文的時候.
很有趣的.
有人會認為.
大衛與約拿丹之間.
其實是不是一種同性戀的愛情呢.
他們會將大衛與約拿丹.
說成聖經其實是支持同性戀的一個例子.

$^{201}$不過.
如果我們看原文的話.
這裡說的愛情.
其實原文不過是在說愛.
就是love.
純粹是一個這樣的字眼.
大衛在哀歌裡面.
用過於婦女的愛情.
其實是在說一份過於對一般女性的愛.
這一份愛不只是在說男女之間情情塔塔的愛.
而是包括了母愛.
後面有很多媽媽.
也包括了在子女對父親中間的那份愛.
所以大衛其實想表達的是.
約拿丹與他之間的友情.
是超乎了一般倫理關係裡面.
特別是家人之間的那種親密的愛.
也就是在說大衛與約拿丹之間.
他們的關係是比親人更加親密的意思.
為什麼對大衛來說.
約拿丹這麼特別呢.
為什麼他們之間會有一份這麼深厚的感情呢.
約拿丹做過什麼.
以致大衛在約拿丹死了之後.
會這麼悲傷.
要作哀歌這麼嚴重呢.
讓我們繼續看下去.
我想大家會明白多一點.
第一點.
我們先看撒母耳記上十八章的一至四節.
這段經文其實是說大衛與約拿丹在聖經裡面.
根據聖經的記載.
他們第一次的見面.
而這段經文一開始就說.
約拿丹的心與大衛的心深相契合.
約拿丹愛大衛如同愛自己的性命.
為什麼約拿丹會這麼欣賞大衛呢.
如果大家在這裡揭前前後後.
其實聖經沒有很詳細地說.
不過我嘗試和大家再看多十八章之前的經文.

$^{241}$讓我們可以稍微了解一下.
為什麼他會這麼欣賞大衛.
我們先看撒母耳記上的十四章.
這裡的經文其實是說當時蘇羅在帶領一幫以色列人和非利時人去征戰.
而他們兩邊的軍隊分別在旁邊的兩個山峰.
他們就在那兩個山峰上去扎營.
戰況其實一直處於一個膠著的狀態.
你又攻不進 我又守不住.
總之就是在兩個山峰上膠著了.
而這段經文就說.
當時約拿丹和那個幫他拿兵器的少年人.
就想著兩個私下走到非利時人的防營裡面.
約拿丹和這個拿著兵器的少年人說.
我們不如去到那幫未受國禮的防營裡面.
或者耶和華為我們施展能力.
因為耶和華使人得勝 不在乎人多人少.
然後這個拿兵器的少年人就回答他.
我就一切聽你的心意去走 我們可以去.
我必定會跟你去 跟你同心.
長話短說 故事大家可以自己繼續看.
結果這兩個人在非利時人的營中.
他們真的突襲成功.
進去殺了大概二十個人.
然後因為他們的突擊之後.
這個防營的營裡面.
展現了一片混亂的狀況.
也令到本身在另一座山的那幫紮營營的以色列人.
就覺得有機可乘.
繼而他們就攻進去.
令到以色列人在這個戰役裡面獲勝.
這是第一件故事 關於約拿丹的罐仔.
然後我們再看 薩姆爾記上十七章的經文.
這段經文大家應該較為熟悉.
大家看到圖片就會知道其實是大衛打哥利亞.
就是巨人哥利亞的一個故事.
當時同樣 以色列人又跟非利人爭戰了.
而在非利人當中.
有一個叫做哥利亞的巨人.
他就在對陣的裡面叫囂.
就跟那幫以色列人說.

$^{281}$你們找一個人出來跟我戰鬥.
而在過程中他也辱罵了以色列的軍隊.
以色列的那幫軍人.
以色列軍人看到他真的高很多.
自然就很害怕.
也因為這樣沒有人回應哥利亞的要戰.
沒有人敢應戰.
當中包括了蘇萊他自己.
而當時那個年紀尚輕的大衛.
他本身純粹去到軍營.
純粹是為了送食物給那三個在重軍的哥哥.
在過程中他聽到.
哥利亞這樣在罵以色列的軍隊.
是萬軍耶和華的軍隊.
他就主動請纓說我要迎戰這個哥利亞.
面對著這個哥利亞 大衛沒有害怕.
他還跟哥利亞說.
你來攻擊我是靠那些刀槍和銅擊.
我來攻擊你是靠萬軍耶和華的名義.
就是你所怒罵帶領以色列軍隊的上帝.
今天耶和華必將你交在我的手裡.
我必殺你.
斬你的頭.
又將菲律士軍兵的屍首.
給空中的飛鳥地上的野獸吃.
令到普天之下的人都知道.
以色列中有一個上帝.
又要令到這班人都知道.
這個耶和華是會令人得勝.
不是用刀用槍.
因為真正的勝敗.
從來都是在乎耶和華.
結果大家應該聽過的故事都會知道.
大衛真的擊敗了那個巨人哥利亞.
也同樣令到以色列人在這個戰事裡面.
可以擊退這班菲律士人.
兩場的戰役.
我們見到的是約拿丹和大衛.
兩個人都有一些一樣的特質.
就是他們都非常勇敢.

$^{321}$還有他們兩個都非常敬畏耶和華.
他們面對著跟菲律士人爭戰.
他們兩個人一樣不怕.
他們確信的是耶和華才是令人得勝的那個.
以至他們分別願意做那個突擊的先鋒.
做那個單挑巨人的勇士.
所以剛剛我們在第十八章那裡說.
約拿丹的心與大衛的心相契合.
其實不是無緣無故.
因為約拿丹其實見到的是大衛和自己一樣.
他們都是非常敬畏耶和華.
是一個與耶和華同行的一個人.
所以約拿丹基於這一份愛.
這一份看見大衛與自己一樣.
那份對耶和華的敬畏.
他就願意跟大衛立一個盟約.
我們要知道.
當時其實大衛不過是一個鄉村的木桶.
就是他負責牧羊.
他不是已經是一個跟隨作戰.
一向不是那一類的.
他是一個小木桶.
而約拿丹是什麼呢?.
約拿丹是蘇羅王的長子.
是王子.
甚至其實他理應是那個王位的繼承人.
不過約拿丹他沒有看到他出生.
他反而在剛才第十八章的一至四節裡面.
他將自己的外袍.
將他的戰衣.
將他的刀工.
腰帶都給了大衛.
將這些本身屬於王子身份擁有的一些物品.
一些隨身物.
他將它給了大衛.
他在提升大衛的身份.
他在提升大衛在以色列百姓眼中的身份.
所以聖經說約拿丹愛大衛如同愛自己的性命.
他在說的不是真的表面的一個交情.
不是真的剛好你和我都是以色列人.

$^{361}$你和我都剛好打非利士人.
我們做朋友吧.
不是.
我們看到他將自己擁有的物品給了大衛.
說的是他將大衛的益處.
他看他的益處.
和在看我自己是一樣的重要.
所以第一點我們在約拿丹身上學到的是甚麼呢.
當我們跟別人說我們要建立一個真誠的友誼的時候.
我們是不會計較大家的身份地位的差距.
我們會為對方去想.
我們不會計較哪一個吃虧.
哪一個付出多助.
反而是我們會很願意為他付出.
今天我們交朋友是怎樣的呢.
我們未必是這樣.
我們很多時候都會看我們有沒有共同的興趣.
信耶穌的就可能會看有沒有共同的信仰.
會看談不談得來.
會看大家的身份地位會不會差距很遠.
有些人就未必.
可能我的身份地位和你差很遠.
為了某些利益我是特意和你做朋友都有的.
但是我想我們同樣都可以去思考一下.
當我們今天去做別人朋友的時候.
我們有沒有學下若拿丹一樣.
如果我真的覺得那個人很值得我交朋友的時候.
我會不會計較他的身份.
我會不會想其實你配不配做我朋友.
或者我配不配做你朋友呢.
會不會想這麼多這些呢.
或者大家可以想一想.
這是第一點.
第二點.
剛剛也提到大衛擊敗了非利士人的巨人哥利亞.
繼而就令到這群以色列人打勝仗這件事.
這個勝仗其實就令到以色列裡面.
這麼多不同的城的人都很開心.
不過好死不死的就是.
當索羅王領軍回到這個國家的時候.

$^{401}$回到以色列的時候.
這群婦女出來的時候.
她們就跳舞在那裡唱歌.
唱不唱就唱了一句.
索羅殺死千千大衛殺死萬萬.
索羅就因為這一句話就開始妒忌大衛了.
而如果在薩姆爾記上18章之後打落的經文.
一直去到20章的內容.
如果我們可以去數一數.
索羅因為這一份妒忌.
其實前前後後他有六次想殺死大衛.
而大衛的好朋友約拿丹.
也都知道他爸爸的確很明確地表明他想殺死大衛.
在第19章的1-6節.
同樣都不讀了,大家可以自己看.
我都會recap給大家看.
在第一節其實他已經講了.
就是在說索羅和約拿丹.
和他這麼多個神僕.
很明確地表明我們要殺死大衛了.
而第二第三節也都解釋了.
約拿丹因為本身喜愛大衛.
他本身和大衛結了盟.
他就主動地提醒大衛.
你要小心.
而他轉過頭.
自己也是為了大衛向他爸爸索羅王求情.
就跟索羅王說.
大衛曾經為了耶和華.
為了以色列人拼命地殺敵.
他做的一切事情其實是有益於你.
索羅王.
他提醒索羅王.
他其實最初自己都曾經因為大衛做的事情.
他很欣賞他.
他為他做的事情很開心.
但是如果現在突然之間調轉過頭.
你就說要殺大衛.
你是在流無辜人的血.
是會叫自己在神面前犯罪的.

$^{441}$索羅王聽到他的兒子約拿丹這樣說.
最後他決定在這一次裡面.
他的確選擇放過了大衛.
所以你見到他最後寫著.
「我必不殺他」.
他還要跟耶和華起誓這樣說.
不過他的誓是有少許漏的.
因為過了沒多久.
真的沒多久.
就在四節的經文之後.
在第十節裡面.
不是這一段.
第十節裡面沒有打經文.
索羅王又因為妒忌大衛.
他又再一次想殺他.
而大衛這一次.
都很清楚自己知道.
索羅王真的定義要殺他.
他連夜地逃跑.
他逃去了一個叫做拉瑪的拿約.
找撒慕爾.
還躲藏在那裡.
不過之後索羅王也因為聽到他躲藏在那裡.
就派人去抓大衛.
繼續想著要取他的性命.
於是大衛又走了.
在那個叫做拉瑪的拿約的地方又走了.
他這次跑去哪裡呢?.
他跑去找約拿丹.
他找約拿丹.
他其實一見到約拿丹的時候.
他第一件事就不是說你爸爸殺我.
他也有說的.
不過他是用一個質問的態度.
去質問約拿丹.
他跟約拿丹說.
究竟我做了什麼?.
我做錯了什麼事?.
在你爸爸面前我犯了什麼罪?.
以至他一定要殺了我.

$^{481}$這一次約拿丹本身是不知道.
他爸爸仍然在追殺大衛.
所以他當下.
他仍然在跟大衛說.
他幫訴羅王說話.
他說訴羅王不會對我有隱瞞的.
如果他真的還想殺你.
我一定知道.
但是由於當時的對話.
大衛很堅定地告訴他.
我真的跟死相差一步.
約拿丹聽到的時候.
他選擇的是.
他願意跟大衛合謀.
他親自去再試一試訴羅王.
是否真的有殺機在裡面.
最後.
的確一如大衛所說.
約拿丹他真的知道了他爸爸訴羅王.
要殺死大衛.
甚至其實在他.
去試訴羅王的過程裡面.
訴羅王其實當下.
氣得連約拿丹也想殺死.
所以他很肯定.
訴羅王是真的想殺大衛.
於是約拿丹按照原本他們合謀的時候.
他們約定過的.
彼此立了一個約定.
他就立刻通知大衛.
叫他立刻走.
剛剛說的故事大概在.
《撒母耳記》上的十九至二十章.
大家可以回去自己看一下.
整個故事大概詳細是什麼.
因為我也飛了不少.
在這幾章的經文.
即是剛剛十八至二十章.
我們看到的是.
訴羅王一次又一次.

$^{521}$想要殺大衛.
但是作為他的兒子約拿丹.
他不惜要跟爸爸.
有衝突要翻臉.
甚至好像剛才說.
令到自己的生命都有危險的時候.
他仍然選擇著.
要幫他的好朋友大衛說好話.
他又順守他們最初的約定.
順守那個諾言.
去通知大衛要著草.
所以說的是.
當我們作別人的一個真正的朋友的時候.
其實不只是在那個所謂的.
順境的時候吃喝玩樂.
就是朋友.
我想大家都有一定的年紀了.
你作為的.
我想你在朋友裡面.
你都知道如果單純吃喝玩樂的.
那些叫做走肉的朋友.
只是走肉而已.
你沒有走沒有肉.
他就未必會來的.
真正的朋友是會在那個困難的裡面.
在他被人誤會的時候.
在他可能在背後被人用言語攻擊.
甚至實際的攻擊.
你會為他挺身而出.
你會選擇站在他那邊.
盡量去幫他.
但前提當然.
說的是在真理和公義的基礎上面.
而不是盲目地死撐.
就好像若拿丹聽完大衛這樣說.
他都不會就這樣說.
喂呀 老伯想殺你.
我相信你.
我馬上回去幫你擋住他.
他不是這樣的.

$^{561}$他都會聽完大衛說.
他都會選擇回去試一試數落王.
究竟若拿丹你現在這樣說的東西.
是不是真的呢.
而不是大衛說什麼都相信.
他都會試.
他都要看一下.
其實若拿丹是不是在公義裡面.
在真理裡面.
仍然是值得相信的那個人.
最後一點了.
最後一點說的是.
當若拿丹知道了他爸爸.
訴羅定義要殺死大衛的時候.
剛剛說他暗中通知了大衛逃走.
而在他們離別的時候.
他們也有一個時間可以互相見面.
在離別的時候.
大衛見到若拿丹的時候.
他主動的就是去俯伏在地.
去拜了三拜.
兩個人親嘴.
彼此哭泣.
大衛哭得更加慘了.
若拿丹就跟大衛說.
我們兩個人曾經指著耶和華的名氣誓.
願耶和華在你我中間.
並你我的後裔中間維持.
直到永遠.
如今你平平安安地去吧.
於是大衛就走了.
若拿丹也回到了他的城裡.
在這個離別的時間.
若拿丹提到他們曾經指著耶和華的名氣誓.
誓的內容包括大衛和若拿丹自己.
也包括了他們兩個人的後代.
若拿丹說的曾經指著耶和華的名氣誓.
就是在他們的交談過程中.
剛剛說他們要試探數落王.
他們的對話中其實曾經提過.

$^{601}$在二十章的十三至十五節.
他就說.
他願耶和華與你同在.
如同從前與我父親同在一樣.
你要照耶和華慈愛恩待我.
不但我活著的時候免我死亡.
就使我死後耶和華從地上剪除你受敵的時候.
你也永不可向我家絕了恩惠.
所以這裡說的是.
他不僅說不要令自己沒事.
假若當大衛真的像耶和華所高納他的時候.
理論上他應該像聖經所說他會繼位.
但是當時他繼位的過程中.
他就說當日他只要有一天能當上皇帝的話.
你記住要保存我的性命.
就算我死了.
你也要保存我後代的性命.
他所說的是這一點.
這裡很有趣的是.
其實若拿丹當時還沒有試探他的父親.
就說這樣的話代表什麼呢?.
其實這裡他說這番話代表若拿丹.
其實當時已經認定了.
耶和華要將整個以色列國交給大衛.
所以他才要求大衛將來作王的時候.
你恩待我和我的後代.
所以當若拿丹如今和大衛離別的時候.
他再一次提起這份恩待後代的約定.
他繼而祝福大衛平平安安地離去.
他所說的是.
他提起這件事.
就盼望將來大衛真的如最初的約定那樣做.
而聖經在之後也沒有特別再提及.
究竟若拿丹和大衛之間有什麼交集的位置.
他們沒有再見面.
有什麼事情發生沒有說.
直到剛才撒母耳記下知道他死了之後的事情.
但我們在整個若拿丹和大衛之間的故事裡面.
我們看到的是其實他們兩人這份友誼.
從來沒有一刻不在耶和華的裡面.

$^{641}$由他們最初因為各自對耶和華的敬畏而惺惺相惜.
到他們選擇結盟.
及後他們一切的起誓立約.
甚至離別.
他仍然都在耶和華的裡面.
所以頂智慕我們今天.
其實我不知道大家平時怎樣對朋友.
當我們作人員一個真誠的朋友.
特別是我們一班有信仰的人.
你有沒有為你的朋友上上去.
祈禱紀念他們呢.
抑或大家那一刻合不合適一起玩就算了.
我想不在於他信不信耶穌.
即是那個朋友信不信耶穌.
而是當他很多時候朋友之間會有彼此相交.
會有一起聊天的時間.
他跟你分享的事情.
究竟你會不會為他祈禱呢.
還是大家聊完,八卦完就算了.
又或者可能純粹給一些所謂的建議.
那個建議可能是來自我們做人的一些小智慧也好.
或者一些常識也好.
是不是講完就算了.
當我們有信仰的時候.
我們這個群體.
究竟我們跟別人做朋友.
跟這個世界外面的人.
他們去跟其他人做朋友.
其實有沒有分別呢.
我不知道大家有沒有.
但是在約拿丹裡面我們看到.
他將他跟大衛的一切的事情.
常常放在耶和華裡面.
所以都盼望我們.
當我們作為有信仰的人.
一班基督徒的時候.
或許我們可以做多一點.
或許我們可以多走一步.
你不需要真的抓住你.
按手有利有勞的那種祈禱.

$^{681}$你知道他的事情.
你背地裡為他祈禱.
甚或你純粹想起他的名字.
突然之間好像很久沒見.
連你記得他 你就為他祈禱都可以.
但是我們有沒有做呢.
總結.
今天我們大概很簡單地.
高估了在聖經裡面.
大衛與約拿丹之間的那份友情.
他們的一些小事情.
一些不同的事蹟.
特別是約拿丹付出的方面比較多.
我在預備的過程裡面.
都上網找了一下資料.
我覺得其中一個挺好的.
就是他們都在說約拿丹這個人.
他用了一段經文去概括這件事.
當然他的篇另外是一篇很大篇的道.
我不會重複整篇道給大家聽.
但我覺得他很有趣的是.
他用了尼加書六章八字.
去概括了約拿丹.
我將他按著我對於約拿丹這個故事的理解.
我都fine-tune了一點.
尼加書六章八字.
行宮二號憐憫全謙卑的心與你的神同行.
我想大家很熟悉.
特別是過去幾年不斷地聽的.
那怎麼走出來呢?.
在約拿丹說服素羅王.
不要流無辜人的血.
其實我們看到的就是.
就算事情.
那個有多親的父親.
他所說的事情.
他很想做.
但當他知道這件事.
是不合乎上帝的心意.
是沒有公義的時候.

$^{721}$他選擇去說服他的父親.
所以我們看到的是.
約拿丹是一個要堅持按上帝公義行事的人.
而約拿丹幫大衛逃離素羅王的迫害.
即使他自己有生命危險.
他仍然要選擇的就是.
要保持他心裡的那份良善.
他不單看大衛不過是一個無名小卒.
或者當時有少少地位.
就選擇放棄他.
他仍然保持著那份良善.
去憐憫素羅王無辜要殺大衛這件事.
同樣約拿丹貴為一個王子.
貴為一個王位的繼承人.
他願意跟大衛結盟.
甚至剛剛看到.
我們見到他某程度上已經認定了.
耶和華揀選大衛做以色列王.
他沒有因為這樣.
就覺得利益受損.
我放棄大衛了.
他沒有.
他仍然謙卑馴服耶和華的心意.
而最後我們見到的是.
無論當初約拿丹去突擊那班非利士人.
甚或在他之後跟大衛的相交裡.
他們一切的事情.
一直都將每一個部份放在耶和華裡.
我們見到的是.
約拿丹是一個一直選擇同神同行的人.
所以當我們今天這個系列.
我們說我們要學怎樣跟朋友同行的時候.
在約拿丹的生命.
在約拿丹怎樣對待大衛.
他的確是一個很好的榜樣.
讓我們可以學習.
我盼望的是.
當我們今天聽完.
原來做朋友是可以這樣.
又或者可能你聽過這些.

$^{761}$不過到你走的時候.
你發覺原來做不到.
那怎麼辦呢?.
讓我們效法約拿丹.
就算我們可能有性命危險.
他仍然堅持很真誠地對待大衛.
他很真誠地為著大衛的好處著想.
盼望我們都一樣.
可以學成為別人眼中的約拿丹.
做人家一個真誠的朋友.
做人家一個覺得你是一個值得信賴的朋友.
而不是那些社交平台.
朋友圈裡面的隨便一個.
讓我們一同祈禱.
是上帝求你的愛先充滿我們.
因為當我們去到.
我們說我們要學怎樣用你的愛.
去到與身邊的人同行的時候.
的確我們這個世界有很多很複雜的人.
我們都會有那份風險.
可能遇著不好的人.
但是上帝求你.
你親自去到我們當中.
你用你的愛填滿我們.
以至我們不是靠我們自己的力量.
去逼那些東西去跟別人做朋友.
而是當我們很清楚知道.
上帝在我們心靈裡面的時候.
我們就有這一份愛.
我們就能夠藉著你給我們的這一份愛.
去與人同行.
成為別人那個真正的朋友.
成為別人那個真正的同行者.
求上帝幫助我們.
祝福我們.
倘若我們過去有受到傷害.
因而卻步.
因而我們說我們都是不好.
與人保持距離.
倘若我們真的有這樣的想法的時候.

$^{801}$求主你再一次去到.
更新我們.
你再一次讓那份更大的愛.
你那個偉大的愛.
那個完全的愛.
進到我們的生命當中.
醫治我們.
釋放我們.
叫我們再一次有力.
去愛別人.
去成為上帝你那個美好的見證.
求主你這樣祝福.
帶領我們每一個.
幫助我們.
奉耶穌基督的聖人名字祈求.
Amen.
\newpage



\section{馬可福音 5:21-34}
\label{sec:031786}
\textbf{耶穌,麻煩快點}~~~~莫樹堅牧師
\newline
\newline
連結: \href{http://www.church.com.hk/acms/content.asp?site=cdc\&op=show\&type=product\&code=031786\&layout=sermon}{\texttt{031786}} ~~~~ 語音日期: 2026-03-08
\newline
\newline
\hyperref[sec:031785]{\small{< < < PREV SERMON < < <}}
~
\hyperref[sec:index_chronic]{\small{[返順時目]}}
~
\hyperref[sec:031787]{\small{> > > NEXT SERMON > > >}}
\newline
\newline
馬可福音 5:21-34
\newline
\begin{longtable}{cl}
\hline
\hline
章節 & 經文 (和合本修訂版)\\
\hline
5:21 & \begin{tabularx}{0.7\textwidth}{X} 耶穌又坐船渡到對岸，有一大群人聚集到他身邊；他正在海邊。 \end{tabularx} \\ \\ \relax
5:22 & \begin{tabularx}{0.7\textwidth}{X} 有一個會堂主管，名叫葉魯，也來了，一見到耶穌，就俯伏在他腳前， \end{tabularx} \\ \\ \relax
5:23 & \begin{tabularx}{0.7\textwidth}{X} 再三求他，說：「我的小女兒快要死了，求你去為她按手，使她痊癒，可以活下去。」 \end{tabularx} \\ \\ \relax
5:24 & \begin{tabularx}{0.7\textwidth}{X} 耶穌就和他同去。有一大群人跟隨他，擁擠著他。 \end{tabularx} \\ \\ \relax
5:25 & \begin{tabularx}{0.7\textwidth}{X} 有一個女人，患了經血不止的病有十二年， \end{tabularx} \\ \\ \relax
5:26 & \begin{tabularx}{0.7\textwidth}{X} 在好多醫生手裡受了許多苦，又花盡了她所有的，一點也不見好，反而更重了。 \end{tabularx} \\ \\ \relax
5:27 & \begin{tabularx}{0.7\textwidth}{X} 她聽見耶穌的事，就夾在眾人中間，從後面來摸耶穌的衣裳， \end{tabularx} \\ \\ \relax
5:28 & \begin{tabularx}{0.7\textwidth}{X} 因她想：「我只摸到他的衣裳，就會痊癒。」 \end{tabularx} \\ \\ \relax
5:29 & \begin{tabularx}{0.7\textwidth}{X} 於是她的流血立刻止住，她覺得身上的疾病好了。 \end{tabularx} \\ \\ \relax
5:30 & \begin{tabularx}{0.7\textwidth}{X} 耶穌頓時心裡覺得有能力從自己身上出去，就在眾人中間轉過來，說：「誰摸我的衣裳？」 \end{tabularx} \\ \\ \relax
5:31 & \begin{tabularx}{0.7\textwidth}{X} 門徒對他說：「你看眾人擁擠著你，還說『誰摸我』呢？」 \end{tabularx} \\ \\ \relax
5:32 & \begin{tabularx}{0.7\textwidth}{X} 耶穌周圍觀看，要見做這事的女人。 \end{tabularx} \\ \\ \relax
5:33 & \begin{tabularx}{0.7\textwidth}{X} 那女人知道在自己身上所成的事，就恐懼戰兢，來俯伏在耶穌跟前，將實情全告訴他。 \end{tabularx} \\ \\ \relax
5:34 & \begin{tabularx}{0.7\textwidth}{X} 耶穌對她說：「女兒，你的信救了你，平安地回去吧！你的疾病痊癒了。」 \end{tabularx} \\ \\
[1ex]
\hline
\hline
\end{longtable}
$^{1}$各位早晨.
很多謝大家讓我有機會一起崇拜.
多謝謝牧師.
所以我經常說「謝謝謝牧師」.
對的,多謝各位掌摺.
我在播道會的幕會.
就因為播道會的規矩有年紀的限制.
所以幾年前過了這個年紀.
就退休去了加拿大,跟家人和兒女一起.
師母就說一件事,最不習慣.
因為我在香港的長時期.
就有飛鴻姐姐幫忙.
就不用煮飯,不用做家務.
現在去了加拿大的生活.
她說沒工人用了,又要煮飯.
這麼多年都沒煮,現在要學習煮.
然後就要做家務.
這是她的經歷.
但我自己就說.
「你沒有,但神的供應我有」.
「你就是我的工人了」.
但這個見證是不可以隨便說的.
因為師母不在崇拜當中.
所以可以自由地分享.
但你見到她千萬不要說這件事.
神的供應是我們很需要上帝的.
因為人生太多問題.
太多周邊環境,世界局面的問題.
動盪的問題,個人的問題,健康的問題.
生活的問題,教會的問題,事務的問題.
如果能夠解決就不是問題了.
很多事情是未能解決的.
我們就嘗試祈禱吧.
「耶穌,你說你很厲害,你幫幫我吧」.
但經常我們都希望耶穌快點幫忙.
我便說,祈禱給耶穌有個七字真言.
是廣東話文化中學會了的.
這七個字是甚麼字呢?.
不用急,最重要是快.
所以今天的講題就是「耶穌,麻煩你快點」.

$^{41}$當然是很文雅的.
「耶穌,麻煩你快點」.
因為是由powerpoint同時在國語堂講的.
所以沒有「耶穌快點」.
「耶穌,麻煩你快點行不行?」.
這段經文今天要講的是記載在馬可福音.
就是耶穌不夠快.
有人不自知症離開世界.
我們要看看這裡究竟在講甚麼.
我先讀聖經給大家聽.
今天這段經文記載在馬可福音第五章.
第五章這裡的記載就是這樣.
五章二十一節開始.
「耶穌坐船渡到對岸,有一大群人聚集在祂身邊.
祂正在海邊,有一個管會堂的人名叫葉魯.
他來了,一見到耶穌,就俯伏在他腳前,再三求他說:.
「我的小女兒」.
這個稱呼是女兒.
「快要死了,求你去為她按手,使她痊癒.
可以活下去」.
耶穌就和她同去.
繼續下一章.
「有一大群人跟隨她,擁擠著她.
有一個女人患了經血不止的病有十二年.
在好些醫生手裡受了許多苦.
又花盡了她所有的,一點也不見後,反而更重了.
她聽見耶穌的事」.
這個聽見即是聽見少少.
沒有見過,沒有經歷過.
不過聽人說好似得得地.
「聽見耶穌的事,就夾在眾人中間.
從後面來摸耶穌的衣裳.
因為她想」.
這個是自然其想.
不知道為什麼,沒有什麼根據.
「我只要摸到她的衣裳,就會痊癒.
於是她的血綠黑止住.
她覺得身上的疾病好了」.
再下一章.
「耶穌頓時心裡覺得有能力從自己身上出去.

$^{81}$再去眾人中間轉過來說.
誰摸我的衣裳.
回頭對她說.
你看眾人擁擠著你.
還說誰摸我.
你不是痴痴地吧」.
語氣就是對比耶穌有一個心腸要做一件事.
一定要這個婦人出來.
去認識這個救主.
不是停留在一個身體的醫治.
要認識救主.
但這個門徒跟隨了耶穌.
那個眼光完全和耶穌有相距.
一個凡.
不過耶穌沒有理會.
要看造者是的女人.
那女人知道自己身上所成的事.
就恐懼戰驚.
來匍匐在耶穌跟前.
將實情傳告訴她」.
耶穌的回應.
耶穌對她說.
「女兒」.
同樣一個葉老說自己的女兒的稱呼.
「女兒.
你的信救了你」.
平安地回去吧.
「你的疾病痊癒了」.
疾病痊癒了.
特別是指她身體痊癒.
所以相對來說.
你的信救了你.
這個救.
就在前文後理最合理的解釋.
就是你能夠有生命的方向.
靈魂得救.
聚得赦免的得救的意思.
你認識了救主了.
再下一章.
「耶穌還在說話的時候」.

$^{121}$那就糟糕了.
「一個就醫好了」.
那之前那個怎麼辦?.
「有人從會堂主管的家裡來說.
你的女兒死了.
何必還勞駕耶穌呢?」.
耶穌不理會他們所說的話.
就對會堂主管說.
「不要怕.
只要信」.
「於是他帶著彼得,雅各和雅各弟弟約翰.
同去.
不許別的人跟著他」.
再下一章.
「他們來到會堂的主管家裡.
耶穌看了一片吵鬧.
看到一片吵鬧.
被有人大聲哭泣哀號.
就進到裡面對他們說.
為什麼大吵大哭呢?.
孩子不是死了.
是睡著了.
他們就嘲笑耶穌.
耶穌把他們都趕出去.
帶著孩子的父母.
和跟隨他的人進了孩子所在的地方.
就拉著孩子的手對他說.
大利大古米.
翻出來就說.
女孩.
我吩咐你.
起來.
那女孩立刻起來走動.
她已經十二歲了.
他們就非常驚奇.
耶穌切切祝福他們.
不要讓人知道這事.
又吩咐給他東西吃.
在馬可福音裡面.
一直在介紹耶穌是怎麼樣的一個耶穌.

$^{161}$在這裡.
寫作就有點特別.
通常記載耶穌.
有一個人有需要.
耶穌做了一些事.
然後就得到解決.
一件事單元過.
但是剛剛讀的段落.
就將兩件事放在一起.
因為其中一件事.
就令到後面那件事.
好像太遲了.
將兩個故事放在一起.
就有特別的訊息了.
兩個故事.
無論是相同.
類同的要點.
或者是不同的要點.
都顯赫帶出訊息來.
我們就看看有幾方面我們要留意.
按著經文要求我們明白.
耶穌是一個大能的耶穌.
牧師你又講耶穌.
牧師不講耶穌講什麼呢.
難道講吃飯嗎.
這裡經文要求我們.
無論我們信了耶穌.
或者是誤導尋索人生.
什麼叫信耶穌.
這個都是要求我們.
重新溫習理解.
重新新鮮的再看一下.
這個耶穌究竟有多厲害.
究竟有多厲害.
耶穌在這個十二年患經血不止.
其實是一個絕望的境況.
一路的武律醫治.
記載這個故事.
錢又用光了.
醫生又欺騙了他.

$^{201}$不僅不好還繼續嚴重.
在一個絕望的境況下.
另外一個十二歲的女孩.
死了不用說了.
怎麼絕望.
絕望到死了.
在這裡介紹耶穌.
就是一個耶穌扭轉人的絕望.
他就有希望.
今天看到這段經文.
就這樣描述耶穌.
昔日是這樣的做法.
今天對我們說.
重新對我們說.
今天的耶穌都可以將我們個人.
我們的生活.
我們的世界的絕望.
帶到去盼望.
我想像不到.
他怎麼可以.
非常好的問題.
那就要謙虛一點.
怎麼能想像耶穌.
可以介入醫好這個.
連醫生都醫不好.
十二年病患血漏的人.
怎麼可以碰一碰就行呢.
我想像不到.
對呀.
如果你想像得到.
你就自大了.
為什麼可以想得透.
這個十二歲的女孩子.
死了可以重新復活過來呢.
我真的想不透耶穌.
為什麼會做得到這樣.
你想不透是正確的.
你如果說自己想得透.
你就自大了.
都是那一句.

$^{241}$今天這麼多的問題.
世界這麼多的問題.
教會這麼多的問題.
教會不只是說同福堂.
同福堂根本沒問題.
我只是說隔離間的教會.
這麼多的問題.
人這麼多的問題.
世界這麼多的問題.
社會這麼多的問題.
不說今生的問題.
來生怎麼處理.
這麼多問題我真的想不透.
耶穌你怎麼可以.
怎麼搞得定.
幸好不用你去想.
如果要你和我想就想不到.
這裡只是很直接地說.
很直言地說.
耶穌就是這麼厲害.
將人的盼望轉化為希望.
將黑暗變成光明.
絕路變為有前途.
這個是好消息.
還有一個好消息.
誰有份.
你說我吧.
像我這麼厲害就行了.
這個故事有相同的地方.
兩個都稱為女兒.
兩個都是十二有關.
一個十二年病患.
一個十二歲.
兩個相同都是絕望的裡面.
那個十二經血不停的.
你好像說.
她沒有馬上死亡的危險.
但是絕對是很接近死亡.
另外十二歲的女孩子.
就直情死亡.

$^{281}$所以都是絕望.
這個是相同點.
帶出兩個故事的相同.
同時有一個非常不相同.
來特別做對比的就是.
第一個是婦人.
婦人在當時大男人的文化.
是沒有地位的.
連經文的名字都沒有提.
直直無名.
還有當時的文化.
宗教文化法律裡面.
這樣流血法.
是不潔淨的人.
不可以去敬拜上帝.
不可以有正常社交生活.
任何一個潔淨健康的人.
碰到一個這樣病患.
叫做不潔淨的人.
自己都變為不潔.
然後不能夠敬拜上帝.
難了自己敬拜上帝.
難了停止了自己正常社交生活.
所以沒有人想碰她.
她也知道這樣去碰到人是不合法的.
沒有錢.
窮困.
被人欺負.
在醫生手下受了很多苦.
然後就說錢都沒有了.
所以又窮又孤單.
又沒有正常社交生活.
又沒有敬拜上帝的機會.
什麼都沒有.
無名無分.
但是耶穌救不救她?.
耶穌關不關心她?.
絕對是.
相對就是葉老會堂.
當時這個宗教中心的領導人.

$^{321}$雖然不是專業的宗教人士.
但是管得會堂都有些身份.
有些名望.
是男人.
甚至有個名字提出來.
然後有人尊敬.
耶穌幫不幫?.
幫.
如果我們假象地說.
如果我這麼好.
就值得耶穌幫.
這是你的錯謀.
如果你說.
是我這麼謙虛.
那些有錢人上帝就很討厭.
這是你的錯謀.
這段經文直接說一件事.
提出來.
對我們說.
耶穌的愛.
耶穌的幫助.
耶穌的願意.
耶穌所愛的人.
所有的人.
無論你是什麼背景.
無論你有什麼問題.
經文裡面沒有一個問題.
大得過耶穌可以幫的.
沒有一個問題.
是太遲.
耶穌幫不到的.
今天世界的動盪局面.
今天的經濟的問題.
政治的問題.
軍事的問題.
看報紙回來.
你都看得到.
內裡還有更多.
你不知道的東西.
耶穌能不能搞定?.

$^{361}$能不能搞定?.
聖經就說.
耶穌連生命都能夠.
從死裡復活.
大能的耶穌.
就是這麼厲害.
耶穌其實不是.
其實有些東西你搞不定.
眼見的問題.
你可能搞得定.
不過我童年有創傷.
有缺陷.
我父母又不夠愛錫我.
使得我這個傷痕.
這個創傷.
帶著在我整個人生當中.
是沒得醫治的.
這段經文直接回答你.
就試試耶穌.
我真想不到耶穌.
怎樣可以重新再回頭.
在我的經歷裡面.
創傷失望裡面.
能夠醫治得到我.
都是那句.
想到你就做耶穌.
就是想不到.
耶穌說我很願意.
很願意幫助你.
無論是你經歷人的.
帶來的失望.
帶來的創傷.
過往你以為沒得救了.
沒得醫治了.
今天這段經文直接跟你說.
Yes.
大能的耶穌.
可以醫治幫助.
我們每一個人.
再向著光明再起步.

$^{401}$OK.
好.
那是什麼意思?.
接著下一個的真理.
在這段經文說了.
信心.
堅定的信心.
接著那一章.
我馬上會解釋什麼是堅定的信心.
牧師.
講道理是講耶穌.
我不知道聽了多少年.
講完耶穌當然講信心.
都是那一句.
如果牧師不講耶穌.
不講信心.
講什麼呢?.
所以都是信心.
經文.
在這段經文特別顯赫的講信心.
在耶穌.
知道是太遲.
人認為太遲.
當然他不是認為太遲.
救這個十二歲的女孩子.
然後就.
管會堂的人.
就出來.
不用煩耶穌.
不用麻煩他.
你的女兒死了沒得救.
然後在這三十六節裡面.
就說.
耶穌不理會他們所說的話.
就對會堂主管說.
不要怕只要信.
在我的經歷裡.
很多信徒誤解了.
濫用了這句話.
你不用理會.

$^{441}$信就行了,不用想了.
人群做什麼你就做.
不經大腦就行了.
上帝沒給你腦子嗎?.
怎麼會不想呢?.
耶穌一直叫信徒.
叫門徒去想事情.
怎麼會不想呢?.
所以這裡一定不是叫他們不要想.
只是信.
這個經文的編排.
在第二個部分.
就回頭說耶路的女兒.
然後耶穌對著耶路說.
不要怕只要信.
這個經文要求我們回去理解.
第一個十二年患經血不止的婦人.
那個信心.
同時今天.
同時當時是要求耶路.
有這樣的信心.
所以我們回頭看看.
這個婦人的信心是什麼信心.
經文裡面說.
她很有人生的問題.
解決不了.
但是一點點聽見耶穌.
很薄弱.
然後自己就奇想.
如果耶穌真的這麼有能力的話.
好像我聽聞的話.
我碰到他衣服的一角一角.
那些能力就能夠解決.
我人不能夠解決的問題了.
你看起來好像很美麗.
很浪漫.
但是當時裡面其實就是一個迷信多於信心.
當時的文化宗教的規條.
沒有教導任何一個這樣的舉動.
不過她自己想.

$^{481}$很堅定.
很薄弱.
不過很堅定.
很determined.
很專一的要做.
然後她是排除萬難的.
當時的環境.
很多人擠擁著耶穌.
哇!耶穌啊!你真好啊!耶穌啊!耶穌啊!.
如果好在當時沒有閃光燈.
沒有簽名.
然後就說耶穌啊!.
自拍!耶穌!跟我自拍!.
沒有啊.
她把人推開.
別人笑她.
別人看低她.
別人覺得她骯髒.
別人覺得她臭.
別人覺得她沒錢.
別人覺得她沒地位.
很低微.
不管.
所有周圍的人怎麼看都不管.
她就專一的.
我要碰到耶穌.
我要試一試耶穌.
就是這麼簡單.
不是很偉大.
不是很多神學知識.
不是很多理論.
一個簡單.
窮困沒地位的女性.
碰一下耶穌就行了.
耶穌給她的評價是什麼呢?.
女兒.
你的信救了你.
你認識我是救主.
不是停留在身體的意志.
你認到我是救主.

$^{521}$這個就是這樣的信心.
排除萬難.
有時候教會要你.
探探哪個關心哪個.
好忙啊.
我要去旅行啊.
我要吃飯啊.
我要休息啊.
是你不想做.
想做的事.
排除萬難.
再回想一下.
大家年輕拍拖.
追求伴侶的時候.
排除萬難.
媽媽說不准出街.
你就從窗口偷出去.
不准拍拖.
你就說要去補課.
回學校.
排除萬難.
那些都是不對的.
精神排除萬難.
我想說什麼叫排除萬難.
為耶穌而活.
相信祂.
很難的.
我很多事務.
你去排除萬難.
專心找一找耶穌.
你沒信過耶穌.
你在這裡尋道.
今天是一個好的時刻.
我的信心就不知道了.
不是很多.
聖經又不認識.
不過我聽過你好像可以的.
我也想試一試.
如果耶穌是假的.
你試一試.

$^{561}$對你沒有損失.
耶穌你幫幫我.
我願意相信你.
是假的.
對你沒有影響.
但如果假若耶穌是真的.
很明顯今天我想告訴你.
耶穌是真的.
在座當中這麼多人.
信了耶穌這麼久這麼多年.
他都會告訴你.
在經歷裡.
在認知裡.
在心理平安裡.
都說耶穌是真的.
你就試一試.
在這裡.
有自己的責任.
堅強的信心.
就算少許的驚心.
都是自己的責任.
三十六節.
已經有聲音了.
不用了耶穌.
他死了.
太遲了.
三十六節.
耶穌不理會他們.
在整段經文裡.
這幾句不理會.
我覺得很好.
也是非常精準的.
有些版本.
耶穌聽了.
耶穌聽了也不是錯.
不過就是譯不到精神.
這個原文的字.
聽的字.
就好像我們英文.
有個字根叫Acoustic.

$^{601}$聲音.
Acoustic.
Engineer的東西叫Acoustic.
原文聽的字.
就是用這個字根.
不過在這裡的Acoustic字根.
加上有個前設詞.
所以就將這個重心.
轉變過來.
耶穌是聽見.
不過好像沒有聽過一樣.
聽了這個絕望.
灰心.
太遲.
沒用的了.
的聲音.
但是不被影響.
這個字.
所以這個中文不理會.
是很有味道的.
今天是我們的責任.
如果說要信耶穌的話.
責任是不理會.
不應該聽的.
指引和人的聲音.
在座當中.
如果有一些資歷更深的信徒.
建立這個堂.
買這個堂.
沒可能啦.
金錢機遇都沒可能啦.
沒可能的聲音.
是我們的責任去扔開它.
而見到當然上帝已經在做事.
現在見到了.
接著前面還有很多上帝要我們生活的.
無論是個人.
無論我們是管理.
無論是管理社會.
或者管理機構.

$^{641}$或者自己的家庭.
有這個聲音.
絕望.
沒可能.
那些人是沒可能改的.
我也沒可能改的.
太遲了.
這是我的責任.
我們自己經民今天要求我們的責任.
三個字是什麼?.
三合會.
不是不理會.
不要理那些聲音.
但是那些聲音真的很大.
當然大啦.
不大怎麼會咆哮到你呢?.
恐嚇到你呢?.
你死定了.
你沒用的.
不要怕只要信.
是自己的責任去不理會.
然後在經文裡面.
耶穌雖然沒說過不理會.
但是行為上就是.
當門徒都不明白.
信徒跟隨耶穌的信徒都不明白.
耶穌這麼多人.
你還要找那個人出來.
要教導他認識你.
完全體會不到耶穌的心意.
但是耶穌不理會.
然後去到會堂主管的家裡.
大吵大鬧.
人又嘲笑他.
然後耶穌的動作是什麼?.
趕他們出去.
無論不理會.
按照真理而行.
然後趕走那些不應該有的障礙.
排除萬難.

$^{681}$只去碰到耶穌.
這是你跟我信心的責任.
所以信心不是很模糊.
很感受感性的.
「我信啦」.
這裡是很清晰地有頭腦的決策.
趕開不應該聽的聲音.
趕走不應該有的指引.
環境和人的聲音.
很多是與上帝的心意很有相距.
是我們的責任.
去不聽不理會.
然後去到耶穌那裡.
甚至是碰到一個崩.
那個力量已經足夠解決人不能夠解決的問題.
你試一下.
你心中祈禱.
耶穌我聽完你這麼厲害的話.
我也想試一下.
或者我們信了耶穌.
就像門徒一樣已經跟隨了耶穌.
但是還沒有體會到耶穌的心意.
深度的心意.
今天就馬上跟耶穌說.
我改.
你提醒了我.
我改.
我就要自己紀律.
自己discipline.
自己有這樣簡單排除萬難堅定的信心.
最後.
真正的信心是帶有行為.
不斷的行為.
不是一朝一夕.
是不斷的更新.
簡單的說.
耶穌以身作則.
祂的生命的行為是什麼呢?.
我們剛才說到.
這個澱不潔的血流的人.

$^{721}$自己成為不潔的.
耶穌.
你在這麼多人當中.
有一個人將會摸你.
不是你不知道.
是不是?.
經文裡面雖然沒有明顯說.
但是耶穌沒有理由不知道將會有人摸他.
但是你都容許不潔的去接觸你.
但是你又沒有害怕過變成不潔.
你絕對不是不潔.
你甚至使不潔淨的人成為潔淨.
就不會被不潔淨的人使你不潔淨.
耶穌.
雖然你不是中國人.
又不懂得功夫.
但是你閃避一下.
應該都可以吧.
雖然你沒有學過功夫.
聖經沒有記載.
但是宇宙萬物之在.
Master of the Universe.
怎會不是功夫Master呢?.
你怎樣都一兩招都可以.
但是你站在那裡.
專門容許不潔淨來接觸你.
只是你去接觸不潔淨.
你說是不是真的這樣?.
這個故事的微妙寫法就是.
還有另外一件事.
12歲的女孩.
碰了死屍.
觸摸死屍.
都是成為不潔的.
當時的文化.
但是耶穌專門要去.
還要專門拉著手.
叫她得醫治.
其他聖經說.
耶穌的大能說一句就行了.

$^{761}$都不用去.
你去到說一聲就行了.
你不用去.
你回家吧.
搞定了.
不是搞定撿屍.
是搞定得醫治.
說一句就行了.
隔山打牛的功力都有.
他這次專門去到.
碰了她.
拉起她.
起來.
立刻得醫治.
立刻起來.
立刻復活.
在這樣的理解下.
經文要求我們.
經文要求.
即背後是聖靈要求我們.
今天看見這個真理.
耶穌的醫治.
是去到疾病當中.
疾病不會使到醫治的人成為疾病.
要我們潔淨的人.
經歷過耶穌的潔淨醫治恩典.
要伸手主動出去接觸.
不潔淨的人.
是我們的使命.
個人的使命和教會的使命.
是不是.
是吧.
明白了嗎.
我下班了.
誰是今天不潔淨的人.
很多不潔淨的人.
不是我們每個人都直接去的.
不過我想提議.
有些事情我們可以先做.
教會內外.

$^{801}$特別是教會內部.
打大門做生意.
什麼人都可以進來.
有很多人是麻煩人.
很多麻煩人.
麻煩人是用最客氣的說法.
令人討厭.
很多人都不想和他聊天.
他又麻煩又怨憤.
又貪心.
叫他聚會.
叫他學習.
叫他服侍.
又不見他.
又多多推搪.
叫他來吃飯.
教會聚餐.
他見到他.
聚餐還要帶一個盒子來.
他拿走了.
這類的人.
例如這樣說.
是我們的責任身手.
用耶穌的愛去撫摸.
醫治他.
從這樣學習開始.
耶穌叫我們有生命行為的改變.
好像耶穌.
不怕接觸不潔淨的人.
好像光去到黑暗.
黑暗就被趕走.
就不會耶穌的光.
消消恥變成黑.
所以我們真正有信心.
求上帝幫助我們.
求上帝幫助我們在計劃當中尋找.
主動尋找.
教會.
由教會內開始.
接著到我們的生活圈子.

$^{841}$平時有些社交能力很低.
那些英文叫social oddity.
又不懂說話.
說話得罪人.
樣子又不好.
最少是令人不喜悅的.
神幫助我們.
今日這段經文.
讓我再重新溫習一次.
耶穌是多麼偉大.
多麼慈愛.
多麼愛我們.
我們都經歷到.
接駁耶穌的幫助與拯救.
是拿出一個簡單.
派翠萬難堅定的信心.
我一定要耶穌.
然後耶穌改變我們的生活.
耶穌改變我們得醫治的人.
去醫治其他不潔淨的人.
\newpage



\section{加拉太書 5:1}
\label{sec:031787}
\textbf{得自由之路 (華)}~~~~葛國光牧師
\newline
\newline
連結: \href{http://www.church.com.hk/acms/content.asp?site=cdc\&op=show\&type=product\&code=031787\&layout=sermon}{\texttt{031787}} ~~~~ 語音日期: 2026-03-08
\newline
\newline
\hyperref[sec:031786]{\small{< < < PREV SERMON < < <}}
~
\hyperref[sec:index_chronic]{\small{[返順時目]}}
~
\hyperref[sec:031788]{\small{> > > NEXT SERMON > > >}}
\newline
\newline
加拉太書 5:1
\newline
\begin{longtable}{cl}
\hline
\hline
章節 & 經文 (和合本修訂版)\\
\hline
5:1 & \begin{tabularx}{0.7\textwidth}{X} 基督釋放了我們，為使我們得自由。所以要站穩了，不要再被奴隸的軛挾制。 \end{tabularx} \\ \\
[1ex]
\hline
\hline
\end{longtable}
$^{1}$我們上次的信息裡面也講到.
加勒太蘇第五章第一節講到.
基督釋放了我們 使我們得自由.
自由非常的寶貴.
對我們來講 我們大概不太能夠體會.
現在全世界的人都在關注伊朗.
美國跟以色列每一天都有進展.
可能有些人天天在看.
到底這個情況會怎麼樣.
不管你喜歡不喜歡 不管你贊成不贊成.
這一場戰爭影響的一定是全世界.
是不是啊 有沒有感覺.
有沒有感覺.
油價有沒有感覺.
有嗎 哎呀 馬上有感覺了.
是不是 所以這個跟我們太有關係了.
但是你知道嗎.
有一些人他們比我們更關心.
他們昨天在洛杉磯舉行示威.
他們在那邊抗議.
他們在那邊慶祝.
他們慶祝美國跟以色列.
這個去攻打這個伊朗.
因為他們認為他們忍受了47年.
他今天他們終於有機會.
他們可以回到他們自己的母國了.
這些居住在.
美國的這些伊朗人.
在全美國各地他們舉行遊行.
他們在那個地方慶祝這件事情.
當然他們慶祝的不是他們自己的國家被轟炸成這樣.
但是我想大家了解這個意思是不是.
因為他們長期以來在這個伊朗的政權之下.
他們沒有得到自由.
他們被壓迫的中間.
所以今天他們有機會可以得到自由了.
對他們來講自由是何等的寶貴.
跟你旁邊說自由何等的寶貴.
可能我們今天我們沒有像他們這樣子被壓制.
但是我們被不同的東西所壓制.

$^{41}$所以這一個基督釋放了我們使我們得以自由.
這個信息不單單是對伊朗的人說的.
這個信息是對所有的人所說的.
所以今天我們要來看繼續的來看.
今天這段經文我想我們特別要問.
從這個哥林多後書這一段聖經裡面.
其實應該是要從第十三節.
不過太長了.
我就等一下我會再提.
那麼保羅在哥林多後書第三章這一段聖經.
其實從第三章的第一節一直開始.
他就講到說我們今天新約的時候.
我們的職分跟舊約的職分是不一樣的.
我們今天都是新約的執事.
執事是什麼呢.
執事就是minister.
不是deacon 是minister.
換句話說你我凡是經歷了耶穌基督的救贖之後.
我們都成為新約的minister.
minister在這個英國來講就是牧師了.
那我們今天講說我們在我們美國來講.
我們講minister跟pastor.
一個是沒有暗牧的一個是有暗牧的.
但我想講的重點不是在這裡.
那聖經上讓我們看見每一個被耶穌所救贖的.
我們都是新約的執事.
minister是在做什麼.
是在服事神的.
可是保羅在這裡講到.
他說他們的心幾時歸向主.
帕子就幾時出去了.
他在說什麼.
他這邊是在引用舊約裡面當摩西.
他在山上領受律法的時候.
四時晝夜他在神的面前親近神對不對.
就是以色列人可不可以去見神.
不可以.
在舊約的時代以色列人如果見了神.
神如果向他們顯現.
如果他們看見神會怎麼樣.

$^{81}$會死的.
所以摩西一個人到山上四時晝夜在那邊親近神.
神把那個律法交給他對不對.
結果當四十天之後.
他得到了律法.
那個摩西的師伴是不是.
這個就是舊約聖經裡面闖出黎明生.
這五卷書 創世記.
出埃及記.
立位紀.
民宿紀.
森明紀 這五卷叫做什麼.
叫做摩西五經懂嗎.
俠義上對於來講.
聖經裡面講的律法.
就是指的摩西五經.
猶太人稱這個摩西五經叫做Tora.
Tora.
那廣義來講.
整個聖經都是可以.
這個律法都可以指的是整個聖經對不對.
但一般來講.
以色列人講的律法.
指的是摩西五經或者是舊約.
那丟著唄.
那麼保羅在引用這段故事的時候.
引用這個歷史的時候.
他提到什麼呢.
他說摩西是那個舊約的執事.
他遇見神的時候呢.
他的臉上就反照出神的容光.
所以當摩西走下山的時候.
以色列人不敢直接看他 為什麼.
因為臉上有容光.
那個光他們沒有辦法指使.
接著保羅就說了.
那個摩西啊.
代表舊約的執事.
他的容光是會漸漸退去的.
可是呢.

$^{121}$新約的執事.
那個容光卻不會退去的.
因此.
他就提到了.
他說他們的心幾時歸向主.
帕子就幾時出去了.
為什麼有帕子.
那個帕子是什麼.
那個帕子就是當摩西.
他下山的時候.
對不對.
以色列人不敢看他 對不對.
所以怎麼樣.
他要有一個帕子.
但是.
保羅說他的帕子什麼時候出去.
那個帕子代表什麼呢.
那個帕子就代表律法.
好 這是這一段聖經原來的意思.
等一下我們會引申出來.
所以第一個.
如果我們要在基督裡面得自由.
我們要從律法當中.
得到釋放.
跟你旁邊說要從律法當中得釋放.
我們這段聖經一直在講.
保羅在迦納泰書講說.
基督釋放我們.
就要我們得以自由.
所以不要再做.
怎麼樣.
律法的盧浦了.
換句話說.
當我們強調.
我們還在靠著舊約的律法.
就像這些猶太人一樣.
他們還在那邊想要靠著守律法.
想做行割禮.
想做這些東西.
來得神的喜悅的時候.

$^{161}$這就是保羅所告訴他們的.
你們不要再做律法的盧浦了.
我們今天會不會有律法的盧浦.
會不會做律法的盧浦呢.
會不會.
今天新約的時代的基督徒.
會不會做律法的盧浦呢.
會.
每個禮拜天都來敬拜.
是吧.
過去你是長老會背景的.
請告訴我.
請舉個手.
請問長老會怎麼敬拜.
至少以台灣的長老會.
敬拜的時候要規規矩矩的.
坐在那邊安安靜靜的.
唱詩歌.
唱傳統詩歌.
所以當台灣敬拜讚美風吹起.
全台灣的教會都開始敬拜讚美.
現在敬拜鼓啊吉他什麼都上去的時候.
長老會開始面臨挑戰了.
他們的年輕人開始流失了.
有些教會開始發現不對了.
開始開了青年人的崇拜.
有一個台灣的牧師叫.
也是姓葛.
葛牧師.
他的兒子做了一些很好聽的歌.
其中一首歌叫什麼.
葛肇熙.
所以他的教會裡面年輕人很多.
為什麼.
因為他把.
他自己就是一個做詩的.
做現代的這些敬拜讚美詩歌.
所以在他們的青年崇拜裡面.
鼓啊吉他什麼通通上來了.
年輕人就很喜歡.

$^{201}$你懂我意思嗎.
但我講這句話的意思不是說.
老年人你不可以喜歡.
OK.
好吧.
如果你旁邊有臉長的.
你跟他說你也要喜歡.
北美的教會.
大部分的會有從台灣來.
從香港來的.
都大概是三四十年以前.
從台灣從香港來的.
很多人還帶著過去.
台灣香港的那種敬拜的觀念.
或者那種傳統習慣.
來到這裡的時候沒有改變.
所以我們還是有帕子.
當我們敬拜的時候.
我們帶著一層帕子.
我們看著別人舉手.
我們覺得說哎呀.
這是靈恩教會啊.
當我們拍手的時候.
這是靈恩教會啊.
我有一次到一個很多年前.
到一個.
華人教會去講道.
講到時候不知道怎麼講.
我怎麼會談到.
反正我談到舉手了拍手.
講完之後他們教會的一個神學生.
就跑來上來跟我講.
他說他偷偷跟我講.
他說牧師啊.
你今天講道很危險哦.
我說什麼事情啊.
他說可能以後.
你大概不會再來這邊講道了.
我說什麼事情啊.
這麼嚴重啊.

$^{241}$他說你今天講道.
說我們敬拜讚美說要拍手啊.
要大聲的拍手唱歌啊.
他說我們教會.
說這是不行的.
啊.
丟子妹有這麼樣子保守的教會啊.
除去帕子.
如果我們的眼睛有這個帕子.
有這個傳統.
有這個律法在我們的裡面.
你沒有辦法.
敞著臉來見這位神.
你破帕子遮住你了.
還有很多時候.
我們會被一些傳統的教導.
或者傳統的習慣.
而不一定.
而不是聖經的原則哦.
這些東西就一直在影響我們.
有一次我到另外一個.
北美的一個教會講道.
講完之後有一個老弟兄跟我分享.
他說.
他說葛牧師啊.
你今天講道說我們.
敬拜.
可以倒.
可以舉手.
可以拍手.
他說啊.
我實在沒有辦法舉起那個手來.
我說為什麼.
他說因為.
以前我在文革的時候.
我每一次舉手我就想到.
想到什麼你知道我就不用講了.
所以在那一個年代背景長大.
他覺得說舉手.
是他一個.

$^{281}$心中的一個痛你懂意思吧.
所以今天來了教會裡面又要舉手了.
他舉不起來.
就像早期.
我希望我們教會現在沒有.
早期有一些國內來的弟兄姊妹.
來了教會裡面之後.
在小組.
查經.
他就說哎呦.
我們在大陸那一套.
今天教會也來搞這一套.
他小組研習毛語錄了.
給你旁邊說.
去掉這些的.
傳統捆綁.
去掉這些律法的捆綁.
去掉這種律法形式主義的捆綁.
現在你為什麼.
你不是在像以色列人一樣活在那個舊約的律法裡面.
你必須要這樣這樣這樣這樣這樣你才叫做敬拜神.
你必須要如此如此如此你才是事奉神.
我們今天很多時候事奉也是如此是吧.
我們叫什麼叫事奉.
服飾.
就是什麼叫服飾.
我們常常覺得我們這樣子的服飾.
我們今天我們要在教會裡面做什麼樣子的一個服飾.
我們覺得我們服飾應該是這樣子.
聖經.
並沒有.
把這些.
形式的東西叫我們傳流下去.
我們北美的教會很多都是.
按照.
剛才我講的.
三四十年前以前的台灣香港教會.
的模式在運作.
教會有團契有沒有.
你現在到全美國各地.

$^{321}$教會裡面是以小組為單元的.
是少是又少.
大部分都是以團契為主的.
但是現在你回台灣去到香港去.
你會發現.
你已經變成少數民族了.
如果你還習慣團契的話.
我今天沒有時間來講.
團契小組的差別等等這些.
這個我們在小組上訓練沒有講過.
但丟著沒.
我們按照過去的傳統方式我們的做服飾.
所以一個教會呢.
去參加團契.
團契安排聚會每個禮拜有不同的聚會.
今天請這個講員明天這個辦一個辯論會.
後天再辦一個講座.
好我們聚會弄得這個這個很.
好像多彩多姿是不是.
再加上來個趴啦那更好.
但丟著沒我們忘掉了到底團契存在是為了什麼.
我們每一年教會裡面辦要辦一次的陪靈會.
這是標準格式.
要辦一次步道會.
你走你走到北美各地的教會去.
你會發現大概都是如此.
我請問.
為什麼一年要辦一次陪靈會.
為什麼.
為什麼教會一年要辦一次步道會.
現在臺灣教會沒有在辦步道會.
來做步道.
你懂意思.
步道會都是為著收割.
辦幸福小組最後來幾個人什麼迷人講座.
那個叫步道會對不對.
是做什麼收割.
我們傳統的教會裡面呢.
我們前面什麼都沒做.
了不起印個單張發起邀請別人來.

$^{361}$然後就期望那一次的步道會.
今天我請個遠志明來.
今天我請個唐崇隆牧師來.
我能夠請張陸家牧師來.
我要請這些有.
好來然後怎麼樣呢.
就靠他一腳把你踢進去.
踢進天堂嗎.
就認為.
這種做法太多年不work.
但是我們很多時候我們習慣這樣的模式.
因為我們也看不到別的模式.
所以我們就一直在這樣子的當中服侍.
一組張陸服侍.
中指妹.
這就是你的帕子.
也沒有辦法經驗到這位神.
是那位做心事的神.
是那位又真又活的神.
所以.
保羅在這裡講了說.
他們的心幾時跪向主.
帕子就幾時出去了.
律法主義.
律法刑事.
很多時候.
教會裡面.
處理事情也是如此.
我們把大家找來開會.
然後開會開會.
開完會之後.
大家異軍討論完之後做決定.
好贊成這樣子方案的A方案的舉手.
贊成B方案的舉手.
好贊成B A方案的人多.
所以我們就去做了.
請問丁永哲妹.
我們跟美國外面的club有什麼差別.
我們跟外面的社團.
開會有什麼差別.

$^{401}$唯一的差別就是.
我們開始的時候做了一個禱告.
That's the only difference.
神的心意在哪裡.
聖經的原則在哪裡.
我們幾乎很少去談這些東西.
所以帶了一個結果教會.
十年二十年.
也看不見上帝的作為.
你懂我意思嗎.
所以我們說.
把你的心.
這個眼睛的這個帕子挪去.
今天我們眼睛沒有這個帕子.
但是我們的心裡有這個帕子.
是不是.
這些帕子.
這些過去的.
習慣過去的方式.
遺文.
律法.
傳統.
這些.
就可能成為我們跟神之間的攔阻.
好第二個.
第二個就是偶像.
偶像不一定是.
這個.
有物質的.
聖經怎麼說呢.
因為你的財寶在哪裡.
怎麼樣.
你的心就在哪裡.
請問你的財寶在哪裡.
Bank of America.
West Fargo.
或者你的.
你的財寶在.
股票市場.
可能你的股票.

$^{441}$可能你的財寶在哪裡.
在洛杉磯有兩棟房子.
在德州有一棟房子.
在佛羅里達州有一棟房子.
在紐約也有一棟房子.
所以你的心怎麼樣.
Everywhere.
of America.
這是一個事實.
耶穌講的一點都不錯.
是不是.
所以錢財.
為什麼在這個地方上下文講的說.
耶穌說你們不可以有兩個主.
你不然就是敬拜馬門.
要不然你就是事奉什麼.
就是事奉主.
所以你只能選擇一個.
你要事奉主.
你還是要事奉馬門.
現在有沒有準備.
今天很多時候.
其實我們基督徒.
比不信的人.
面臨的挑戰更大.
為什麼.
外面不信主的人.
他追求金錢.
也沒有人去罵他說.
你怎麼追求金錢.
對不對.
他也不需要禮拜天來聽牧師講.
這麼一篇道說.
什麼什麼什麼什麼跟什麼.
什麼新洲的偶像.
我大大方方的賺錢.
是不是.
而且我還教你怎麼賺錢.
我賺錢多.
你看網上面教你怎麼賺錢的.

$^{481}$這些視頻一大堆.
對不對.
偏偏我們基督徒來了教會裡面.
哎呀金錢不要成為偶像.
其他留什麼.
我要再聲明.
是不是基督徒都應該窮光蛋.
不是.
基督徒可不可以有錢.
可以.
但是基督徒有錢.
而不要變成錢的奴隸.
懂這意思嗎.
基督徒可以有錢.
但是錢不是你的.
你的.
你不是錢的主人.
錢的主人是誰.
是主.
這是保羅在這裡講的目的.
在.
汐谷灣區流行的一首打油詩叫做灣區工程師之歌.
這個工程師之歌這樣說.
一進汐谷心裡發毛.
二手舊車東奔西跑.
三十出頭白髮不少.
四尺小屋要價天高.
五彩屏幕鍵盤敲敲.
六神無主終日辛勞.
七夕流浪之旅難找.
八萬頭垮.
股票套牢.
九點回家只想睡覺.
十萬家當房還搶不到.
最後這個打油詩的結論就是.
千辛萬苦虛無縹緲.
不知道在幹什麼.
今天很多的人人生就是如此.
是嗎.
不知道人生在幹什麼.

$^{521}$一股腦的大家都要到汐谷去.
因為那是發財的地方.
要做工程師一定要到汐谷去.
要能夠過好生活一定要到美國來.
是嗎.
像劉子梅.
世界上的人他們沒有方向.
因為他們沒有神.
但是基督徒你我.
我們應該要有人生的方向.
為什麼.
因為.
我們的神.
他對我們的人生有旨意.
所以不要讓金錢不要讓任何的東西.
來成為我們的偶像.
影響到我們來跟隨這位主事奉他.
好了.
那麼.
最後我再提一個.
這我也曾經提過.
我說這個偶像.
很多時候對我們北美華人來講.
是我們的兒女.
一切都是為了兒女.
省思儉用.
打拼做什麼.
最後都是為了兒女.
是吧.
為兒女一點錯都沒有.
可是如果兒女變成我們的偶像.
那就不是神的心意.
有一次有一個農夫.
一個農村的一個基督徒夫婦.
他們平常並不太在意兒女的屬靈生活.
有一天當他們夫婦兩個要去教堂去禮拜.
他們就告訴他的小兒子說.
他說牧師講的你聽不懂.
所以你不用去.
趁這段時間.

$^{561}$你把牛放出去養一養.
去吃草.
好了.
教堂完了之後.
這對夫婦就去做禮拜.
做完禮拜回來之後.
竟然看見小兒子.
只有把那個母牛牽去吃草.
卻把小牛關在這個這個牛棚裡面.
這父親看得很生氣.
他就把他找來了.
準備好好的修理他一頓.
結果到了傍晚小兒子回來了.
父親就開始問他了.
你怎麼只帶母牛去吃草啊.
你為什麼留下小牛在餓肚子呢.
小兒子從容不迫.
從從容的說.
這都是你教的.
因為禮拜天只有大人能夠到教堂.
小孩子都要留在家嗎.
親愛的豬子妹.
我們做父母親的榜樣很重要.
我覺得這個故事反映出來.
是我們今天華人教會是相反.
今天我們看到很多基督徒啊.
甚至不信主的人.
他把小孩子送到阿娃娃來.
然後怎麼樣.
自己走人.
有沒有.
所以我聽到有人說.
哎呀.
我們小孩子就在你們阿娃娃長大的.
就在你們小朋友長大的.
那我說你怎麼不來了.
他把教會兒童世公當作baby seat.
你懂意思嗎.
或好一點.
他覺得教會可以教他好的做人處事的道理.

$^{601}$小孩子在教會不會變壞.
你懂這意思嗎.
豬子妹.
難怪今天很多年輕人在在在在教會長大.
去了大學之後.
他一去不復還.
因為他看到父母親的價值觀.
其實那個是最真實的.
你說你愛豬嗎.
你說你看中什麼.
但是你好像當我要去這個準備考試的時候.
或者我要去踢soccer的時候.
你說沒關係.
去踢吧.
以前我們小孩子.
小的時候都會去.
師母都帶他們去踢soccer.
但是碰到有一次.
他們大概會有.
比較多的時間.
會在禮拜天要去比賽.
我們就決定不要讓他去參加這個soccer team.
豬子妹.
你不要小看.
你的一個決定.
就在告訴你的孩子.
你的價值觀是什麼.
如果有別的事情可以做.
禮拜天就不用去了.
你就在告訴他一個價值觀.
就是有別的事情.
比去禮拜天做禮拜更重要.
求助幫助我們.
所以在我們心中有沒有這些的偶像.
在我們心裡面有沒有這些不合神心意的價值觀.
在我們心裡面.
甚至有沒有一些的誤會.
一些罪惡.
成癮.
各式各樣的這些癮.

$^{641}$手機的癮.
我覺得這是現在最大的問題.
是不是.
所以我們跟神之間的關係都被他影響.
因此.
那麼保羅在哥林多後生裡面講到.
要除去這個 拿去這個帕子.
那接著他是說什麼呢.
我們一起來練這段聖經.
我們再來練好不好.
來 預備 請.
主就是那靈.
主的靈在哪裡.
哪裡就得以自由.
我們眾人既然敞著臉.
得以看見主的榮光.
好像從鏡子裡返照.
就變成主的形狀.
請記得這裡說.
他說主的靈在哪裡.
哪裡就得以怎麼樣.
OK 好.
然後他前面說主就是那靈.
這裡讓我們看見一件事情.
主就是那靈.
主聖靈.
他們沒有分家的.
他們是合一的.
親愛的兄弟姐妹.
我們也不要把他分家.
所以其實我很不喜歡.
我們這樣分說你是靈恩派的.
聽我準備.
當你講這句話的時候.
那請問你什麼派的.
你什麼派.
你懂我意思嗎.
當我們講他是靈恩派的時候.
什麼意思.
你可能自己講都不知道.

$^{681}$所以教會裡面在分這些東西的時候.
無不知不覺的就把我們對聖靈還可以分聖所謂的靈恩派,非靈恩派.
我們豈不是都應該是有聖靈的嗎.
是不是.
沒有聖靈你怎麼重生呢.
沒有聖靈你怎麼樣子能夠浮世呢.
沒有聖靈的充滿.
我們怎麼樣子來浮世成傳福音呢.
對不對.
我們何等需要聖靈.
主就是那個靈.
主的靈在哪裡.
也就是基督在哪裡.
哪裡就會得到自由.
如果基督在你的心裡面.
主的靈在你的裡面.
你的心就會得到自由.
可是接下來他就提到了那個方法.
或者說那個途徑.
他的意思就是說.
今天新約的時代.
聖靈要來做一個工作.
不再是靠著守律法了.
聖靈今天重生了我們.
因此我們得到自由.
我們可以敞著臉.
來看見主的榮光.
好像從鏡子返照.
我今天早上是有去看鏡子.
但是鏡子看不到後面.
她看到我的後面.
說牧師你的頭髮怎麼翹起來了.
趕快去幫我用點水把它撲一撲.
現在你還看不到沒關係.
丟著沒.
你看鏡子的時候你看到誰.
你看到誰.
當你去看鏡子的時候你看到誰.
看到你自己是不是.
古代的鏡是銅鏡.

$^{721}$所以呢.
並不像我們現在這個鏡子這麼清楚.
但是保羅用這個比喻來講說.
當你看鏡子的時候.
鏡子就會返照出來.
是不是.
好.
那他講的這個鏡子是指什麼.
就是當你來看主的時候.
就會把主的樣子怎麼樣.
返照出來.
親愛的丘子妹.
所以保羅在這邊講到.
就是聖靈來釋放了我們.
讓我們跟神之間的關係.
是可以這麼樣子的敞開.
是這樣子的可以自由的.
沒有律法的限制.
沒有條文的限制.
沒有policy的限制.
自由的來到神的面前.
來敬拜他.
那你看看.
這個.
大衛在.
大衛在這個詩篇.
二十七篇.
第四節.
我們來念這一段經文好不好.
來預備請.
有一件事.
我曾求耶和華.
我仍要尋求.
就是一生一世.
在耶和華的殿中.
瞻仰祂的榮美.
在祂的殿裡求問.
大衛.
他說有一件事情我曾經尋求.
我曾經尋求就是我曾經禱告.

$^{761}$然後他說什麼.
他說我仍要.
I still.
我仍要我還要尋求.
表明他不是一次的尋求.
他一再的尋求.
他尋求的這件事情就是.
一生一世住在耶和華的殿中.
怎麼樣.
瞻仰祂的榮美.
我們都知道.
當你看什麼.
你就會像什麼.
當你看什麼.
其實.
你就在影響你的思想.
我們的思想非常重要.
真言第四章二十三節說.
一生的果效.
是由心發出的.
所以我們的心如何想呢.
我們的心很容易受到.
我們所看的影響.
中國人講說靜珠者赤.
靜默者怎麼樣.
黑.
所以.
小孩子去交朋友的時候.
如果你看見他的朋友.
都是那個穿著那個.
女孩子穿著都是很暴露的.
男孩子穿著打仗都是看起來就像那個.
小混混的.
請問你做父母親擔不擔心.
擔不擔心.
沒關係現在年輕人都是這樣.
都是這樣子.
都這樣子嗎.
你也可以找到都不是這樣子的.
你懂意思嗎.

$^{801}$所以我們會替小孩子考慮.
應該選擇什麼樣子的朋友是嗎.
是不是.
因為他跟什麼樣子的孩子.
還記得孟母三遷的這個例子嗎.
有沒有.
哎呀天天都在那邊這個.
在那邊砍肉啊.
在那邊賣菜啊.
那好了小孩子從小長大.
他就怕說以後他也變成怎麼樣.
來一斤豬肉五塊錢.
所以孟母要三遷.
是不是.
我們看見我們的環境影響我們很大.
你看什麼你在想什麼.
這個.
柴契爾夫人.
這是據說.
因為.
這個這不是一個.
說一個很肯定的.
但一般說認為說.
這是他說的一段話.
也有人說是英國的一個詩人說的一段話.
anyway.
這不重要.
他說這一段話什麼呢.
他說什麼話.
啊.
他說.
注意你的思想.
他們會變為言語.
注意你的言語.
他們會變為行動.
注意你的行動.
他們會變為習慣.
注意你的習慣.
他們會變為性格.
注意你的性格.

$^{841}$他們會變為你的命運.
換句話說這一切最後變成是我們的命運.
會變成影響我們的一生.
是什麼.
是從我們的思想開始的.
亞歷山大.
那個征服了歐亞非山州的亞歷山大.
你知道嗎.
他受到荷馬的一個史詩.
就是荷馬的史詩.
叫做以利亞德.
這個以利亞德呢.
裡面.
有一個主角叫做阿基里斯.
所以.
亞歷山大非常喜愛.
這個荷馬的這個詩.
他讀了再讀.
他一再的讀.
他把這個阿基里斯當作是偶像.
他渴望像他一樣.
他變成是征服.
然後像英雄一樣.
然後像英雄一般的.
能夠榮耀的歸來.
所以這種的追求.
就使得他以以利亞德.
作為他的軍事指南.
率軍東征.
並且他還在特洛伊.
你們如果去過希臘的.
在特洛伊.
在那邊祭拜.
甚至據說.
亞歷山大在出征.
這個波斯之前.
特地跑到特洛伊.
去那個地方向阿基里斯去致敬.
而且在南征北討的期間.
他甚至隨身攜帶.

$^{881}$一本亞里士多德校定的以利亞德.
所以你看看這一本.
荷馬的史詩.
如何深深的影響他的思想.
促成了他.
南征北討.
征服了歐亞非三州.
我們的思想在受什麼影響呢.
今天我想我們很多的人.
都不可避免的受到這些.
social media 的影響.
根據統計.
美國大概.
美國人平均.
每一天.
看這些所謂的social media.
五個小時一分鐘.
平均.
當然了看什麼東西.
可能有的人說.
哎呀我是工作關係.
沒關係.
不必.
我沒有要大家來解釋.
我只知道大家看見.
現在已經變成是.
大家看手機.
看這些social media.
已經變成是一個很普遍的現象.
一天五個小時.
這是平均值.
有的人更多.
然後呢.
一天當中呢.
大概.
去打開手機來看.
186次.
你們數學好一點的幫我算一下.
假設你在醒著的時間.
186次那是多久看一次.

$^{921}$所以你會發現今天很多的人.
開車也看手機有沒有.
很危險是吧.
好.
然後呢.
開會的時候呢.
也看手機.
還有呢.
還好.
目前我沒有看到足不崇拜看手機的.
如果有看到大家舉報一下.
好吧舉報一下.
提醒他不要再看手機了.
你被手機捆綁了.
吃飯的時候看手機.
是不是.
有什麼重要事情重要到.
兄弟們我們已經被他捆綁了.
你不承認沒關係.
但是你就是被捆綁了.
現在就準備.
所以.
我們受到這些的影響.
這些東西就是.
他在塑造我們的思想觀念.
你喜歡那個pink.
black.
你的腦袋裡面就是啊.
就是pink black.
對不對.
你看的是韓劇.
你天天想的招式我想的都是韓劇.
我現在不知道是哪一個韓劇了.
我只記得最早的時候.
那個已經是大概15年前.
為了讓我的孩子們學.
學中文.
我到中國去.
培訓的時候呢.
我就去給他們買大腸精.

$^{961}$因為那時候正在紅那個大腸精.
那麼所以就給他買那個dvd.
vcd dvd.
那裡面是有那個中.
是中文發音的.
因為是從中國買的.
所以他中文發音.
就為了他讓他學中文.
ok 好嗎.
沒有別的意思.
ok.
而且大腸精.
據說我自己沒有去看過.
據說是比較.
稍微還健康一點.
是嗎.
有人在點頭.
謝謝你的support.
ok.
現在我們父母親很多時候.
給小孩子一個手機.
太早你是害了他.
我聽到有家長.
他給小孩子.
用那個最原始最原始的.
那個smart phone.
就是像那種.
摺疊式的.
就是.
你如果真的是要打電動玩具.
他只能夠有那個什麼.
那個蛇的那種有沒有.
非常基本的.
這其他啥都沒有.
那個是刻意的.
他刻意這樣做.
他刻意不要讓孩子.
太早去接觸這些東西.
我們今天面臨的挑戰很大.
特別做父母親的.

$^{1001}$所以.
這些都在影響我們的思想.
都在塑造我們的人格.
塑造他將來的命運.
今天你看見很多的孩子.
為什麼會拿著個槍就衝進.
衝進這個教園裡面去殺人呢.
那不是一天兩天變成這樣子的.
他可能打電動玩具不知道打了多少年了.
他已經殺人如麻了.
是不是.
所以他到了那個時候的時候.
他就是他就覺得.
今天有機會來了.
家裡面又有槍去了.
就乾了這些事情.
求主幫助我們.
讓我們.
從基本的功課開始.
我們要能夠.
變成主的樣子.
從哪裡開始呢.
從每一天建立.
早起的靈修導入開始.
怎麼把它關了.
剛才.
這個照片是我在.
這已經是三十年前.
我到中國去培訓的時候.
那麼我刻意的.
因為我知道.
那時候中國的教會.
那時候我們都流行唱一首歌.
叫中國的早晨五點鐘.
傳來祈禱聲.
所以我就去看.
我就問他們說.
我說你們都到早上起來.
你們到哪去禱告啊.
他們就說我們到大深山裡面去.

$^{1041}$我就去了.
跟另外一個傳道人.
我們就一起去了.
我們去看了之後.
就把這個照片照下來.
那個他們是早上四點鐘在那邊禱告.
因為我們到了關西.
所以有培訓.
他們是在培訓.
所以五點鐘的時候.
他們四點鐘.
這一群人他們自己禱告完的時候.
五點鐘再到培訓的地方.
他們在一起的.
在一起的晨禱.
走遍大家南北到中國那個年代.
早上五點鐘都是如此.
都是起來晨禱的.
但是曾幾何時中國開始城市化.
曾幾何時中國的環境條件好了.
開始外面的社會.
外面的教會的幫忙越來越多的時候.
結果中國的教會開始漸漸的.
墮落了.
後來我開始.
再一次的傳講.
中國早晨五點鐘起來禱告的這個信息.
很多老一輩的.
傳道人他們知道.
聽了之後他們真是流淚.
因為他們知道我在講什麼.
年輕一輩不知道.
聽說.
但是他們的老的一輩知道.
張國當年的復興就是因為禱告.
就是因為早上起來.
可能這些人大致都不是一個.
聖經都沒有讀過一遍.
可是他們知道起來禱告.
在那樣困難的環境當中.

$^{1081}$是罪孽者這樣子的禱告.
支撐他們面對逼迫困難的當中走下去.
是在這樣子的禱告當中.
他們建立跟神親密的關係.
教會因此復興.
弟兄姊妹今天我們在北美.
不要中國的早晨五點鐘.
北美的早晨幾點鐘.
我們教會是六點鐘.
可能對很多人來講都很挑戰是吧.
可是現在弟兄姊妹.
你聽牧師講到.
在黎良堂今年第十四年了.
你可以說牧師啊.
你有沒有點新鮮的道可講.
對不起這就是新鮮的道.
一個道如果你沒有去做到.
這個道永遠對你來講都是新的.
Amen.
不管你今天做了執事.
做了長老做了小助長做什麼.
你更需要起來禱告Amen.
早期我們教會裡面同工就任時候.
其中有一項要同工們立約的.
就是說要建立靈修的生活.
要固定參加教會禱告會.
為什麼要這樣規定.
因為我們是在服事.
我們不是在外面世界在做事.
我看過太多的同工服事到後來burn out.
因為沒有跟神友好的關係.
看到太多的教會因為服事的人本身屬靈的光景不對.
所以肉體血氣都出來.
開會的時候吵架.
開會的時候大家意見不合.
然後最後教會分裂.
就是因為.
因為我們太疏忽到.
我們這個人何等重要.
沒有人就沒有工作.

$^{1121}$你有再好的再大的教堂.
你有再多的錢財.
如果你沒有神的同在.
一夜之間神都可以把它拿去.
誰跟你旁邊說.
從明天開始.
清早起來靈修吧.
Amen.
如果你可以.
你可以來加入教會的.
教會六點就開始有宣讀聖經.
六點半就開始成禱.
那你說哎呀.
這個時間可能不方便.
或者有其他原因沒關係.
如果你不能夠參加.
你自己要有啊.
你自己又沒有.
又不來參加.
那你怎麼樣.
你就是擺爛.
千萬不要擺爛.
兄姊妹.
千萬不要放棄.
我不放棄你.
我講這麼多年.
我希望看見不是只有八十個人.
我看見是我們所有的弟兄姊妹們.
每一天都建立起跟神一個穩定的靈修生活.
Amen.
最後.
這裡講到說.
瞻仰主的容美.
在他的殿裡求問.
然後因此呢.
就帶來一個結果.
哥林多後世說什麼.
說變得容上加容.
就如同怎麼樣.
主的靈變成的.

$^{1161}$也就是說.
當我們這樣子長期的跟主親近的時候呢.
我們就越來越像誰呢.
就越來越像耶穌基督了.
丟這邊看看你旁邊像誰啊.
你說夫妻啊.
我們說夫妻臉夫妻臉.
為什麼會夫妻臉呢.
哦.
天造的一雙.
地設的一對.
是嗎.
不是.
因為他們兩個人經常在一起常常在一起.
你看我我看你慢慢就怎麼樣.
就越來越像了嘛.
是不是.
對不對.
那你看兩個不像的.
啊對不起啊.
這應該是一個自然的結果.
當你跟主親近的時候.
你一定是越親近他怎麼樣.
就會有主的.
柔美在你的身上彰顯出來.
有沒有.
求主幫助我們.
要能夠建立.
要能夠建立這樣子的關係啊.
其實不是單單靠每一天早上.
那半小時或者一個小時的靈修.
一定要操練到一個地步.
就是時時刻刻的.
活在主的面前.
或者說.
常常的.
住在.
主的裡面.
那是一個什麼樣子的一個生活呢.
那是一個隨時隨地的禱告.

$^{1201}$清淨神.
常常晝夜默想主的話.
這個操練不容易的.
兄弟妹.
在我們聚會之前啊.
你試試看下次你來.
不要多.
我只要求你五分鐘好不好.
你提前五分鐘來.
坐在那邊.
像剛才我在那邊安靜啊.
直到主席在講說.
兄弟妹.
現在我們要準備聚會了.
請大家安靜.
我都還聽到有人在講話.
親愛的兄弟妹.
你看我們多不容易安靜是不是.
我們的心很多的時候.
都聲音在我們旁邊.
但親愛的兄弟妹這個功課太重要了.
你知道嗎.
連佛教徒現在他們都追求這個東西.
在台灣多年前就曾經有人說.
他願意付十萬台幣的錢.
去參加一個五天的瑞翠.
那個瑞翠教你什麼呢.
就叫你冥想打坐.
很多有錢的生意人他願意付這個錢.
為什麼.
因為他太忙了忙到他身體出狀況.
心靈也不安定.
所以他覺得他能夠有那一種的.
他希望能夠有這樣子一個冥想的操練.
親愛的兄弟妹.
這是我們今天基督早就教導我們的.
要我們住在他的裡面.
這個操練非常非常的重要.
有一次在中國一次的培訓當中.
都是各地來的通工.

$^{1241}$傳道人.
在那一次聚會當中我講完信息之後.
我就帶他們做一個操練.
我就帶他們唱了這首歌.
除你以外.
怎麼樣.
在天上我還能有誰.
除你以外.
在地上.
然後我唱完這首歌之後.
我就帶他們一起禱告.
我就跟同學們講.
我說.
我們福祉很多.
我們很多的聚會.
我們很忙碌.
但是你愛主嗎.
我說盼望.
我們有一段時間你都不要禱告.
你就是安靜在主的面前.
你知道嗎.
安靜了一段時間之後.
我說弟兄姊妹.
你有感動的開口禱告.
許多的弟兄姊妹.
這些同工們.
開口哭著禱告.
向主承認說主啊.
我們太忙了.
我們忙到許多時候.
我們忘了跟你的關係.
主赦免我們.
親愛的弟兄姊妹.
求主幫助我們.
那不可少的只有一件事情.
大衛說我過去曾經尋求.
我如今還要尋求.
就是一生一世住在他的殿中.
瞻仰他的榮美.
向他求恩.

\newpage



\section{使徒行傳 9:1-21}
\label{sec:031788}
\textbf{不一樣的人生:新方向}~~~~李黎秀兒師母
\newline
\newline
連結: \href{http://www.church.com.hk/acms/content.asp?site=cdc\&op=show\&type=product\&code=031788\&layout=sermon}{\texttt{031788}} ~~~~ 語音日期: 2026-03-08
\newline
\newline
\hyperref[sec:031787]{\small{< < < PREV SERMON < < <}}
~
\hyperref[sec:index_chronic]{\small{[返順時目]}}
~
\hyperref[sec:031789]{\small{> > > NEXT SERMON > > >}}
\newline
\newline
使徒行傳 9:1-21
\newline
\begin{longtable}{cl}
\hline
\hline
章節 & 經文 (和合本修訂版)\\
\hline
9:1 & \begin{tabularx}{0.7\textwidth}{X} 掃羅不斷用威嚇兇悍的口氣向主的門徒說話。他去見大祭司， \end{tabularx} \\ \\ \relax
9:2 & \begin{tabularx}{0.7\textwidth}{X} 要求發信給大馬士革的各會堂，若是找著信奉這道的人，無論男女，都准他捆綁帶到耶路撒冷。 \end{tabularx} \\ \\ \relax
9:3 & \begin{tabularx}{0.7\textwidth}{X} 掃羅在途中，將到大馬士革的時候，忽然有一道光從天上下來，四面照射著他， \end{tabularx} \\ \\ \relax
9:4 & \begin{tabularx}{0.7\textwidth}{X} 他就仆倒在地，聽見有聲音對他說：「掃羅！掃羅！你為甚麼迫害我？」 \end{tabularx} \\ \\ \relax
9:5 & \begin{tabularx}{0.7\textwidth}{X} 他說：「主啊！你是誰？」主說：「我就是你所迫害的耶穌。 \end{tabularx} \\ \\ \relax
9:6 & \begin{tabularx}{0.7\textwidth}{X} 起來！進城去，你應該做的事，必有人告訴你。」 \end{tabularx} \\ \\ \relax
9:7 & \begin{tabularx}{0.7\textwidth}{X} 同行的人站在那裡，說不出話來，因為他們聽見聲音，卻看不見人。 \end{tabularx} \\ \\ \relax
9:8 & \begin{tabularx}{0.7\textwidth}{X} 掃羅從地上起來，睜開眼睛，竟不能看見甚麼。有人拉他的手，領他進了大馬士革。 \end{tabularx} \\ \\ \relax
9:9 & \begin{tabularx}{0.7\textwidth}{X} 他三天甚麼都看不見，也不吃也不喝。 \end{tabularx} \\ \\ \relax
9:10 & \begin{tabularx}{0.7\textwidth}{X} 那時，在大馬士革有一個門徒，名叫亞拿尼亞。主在異象中對他說：「亞拿尼亞！」他說：「主啊，我在這裡。」 \end{tabularx} \\ \\ \relax
9:11 & \begin{tabularx}{0.7\textwidth}{X} 主對他說：「起來！往那叫直街的路去，在猶大的家裡，去找一個大數人，名叫掃羅；他正在禱告， \end{tabularx} \\ \\ \relax
9:12 & \begin{tabularx}{0.7\textwidth}{X} 在異象中看見了一個人，名叫亞拿尼亞，進來為他按手，讓他能再看得見。」 \end{tabularx} \\ \\ \relax
9:13 & \begin{tabularx}{0.7\textwidth}{X} 亞拿尼亞回答：「主啊，我聽見許多人講到這個人，說他怎樣在耶路撒冷多多苦待你的聖徒， \end{tabularx} \\ \\ \relax
9:14 & \begin{tabularx}{0.7\textwidth}{X} 並且他在這裡有從祭司長得來的權柄，要捆綁一切求告你名的人。」 \end{tabularx} \\ \\ \relax
9:15 & \begin{tabularx}{0.7\textwidth}{X} 主對他說：「你只管去。他是我所揀選的器皿，要在外邦人、君王和以色列人面前宣揚我的名。 \end{tabularx} \\ \\ \relax
9:16 & \begin{tabularx}{0.7\textwidth}{X} 我也要指示他，為我的名必須受許多的苦難。」 \end{tabularx} \\ \\ \relax
9:17 & \begin{tabularx}{0.7\textwidth}{X} 亞拿尼亞就去了，進入那家，把手按在掃羅身上，說：「掃羅弟兄，在你來的路上向你顯現的主，就是耶穌，打發我來，叫你能再看得見，又被聖靈充滿。」 \end{tabularx} \\ \\ \relax
9:18 & \begin{tabularx}{0.7\textwidth}{X} 掃羅的眼睛上立刻好像有鱗一般的東西掉下來，他就能再看得見，於是他起來，受了洗， \end{tabularx} \\ \\ \relax
9:19 & \begin{tabularx}{0.7\textwidth}{X} 吃過飯體力就恢復了。 \end{tabularx} \\ \\ \relax
9:20 & \begin{tabularx}{0.7\textwidth}{X} 立刻在各會堂裡傳揚耶穌，說他是神的兒子。 \end{tabularx} \\ \\ \relax
9:21 & \begin{tabularx}{0.7\textwidth}{X} 凡聽見的人都很驚奇，說：「在耶路撒冷殘害求告這名的不就是這個人嗎？他不是到這裡來要捆綁他們，帶到祭司長那裡去嗎？」 \end{tabularx} \\ \\
[1ex]
\hline
\hline
\end{longtable}
$^{1}$我們一起有一個很簡單的祈禱 好嗎?.
天父 我們藉著你如此耶穌基督的名譽.
現在我們就能夠坦然進到你的面前.
天父 我們知道你在這裡.
我們很恭敬的將到以下的時間交託給你.
懇求你藉著你的靈.
感動引導我們.
讓我們能夠得到生命的餵養.
也都得到生命的改變.
我們恭敬在你面前 等候你.
知道你會工作.
奉耶穌基督的聖名 阿們.
人生有時會有不同的轉變.
不知道你的人生有沒有很大的轉變.
以致你現在有不一樣的人生.
聖經記載有一位人物.
他自從在大馬色路上.
遇見主耶穌之後.
他的生命就有很巨大的改變.
這位是哪位?.
哎呀 真是的 沒錯.
這位就是保羅.
今年我會嘗試用一個系列.
跟大家分享.
就是從保羅看不一樣的人生.
今天是第一講.
今天的題目就是.
不一樣的人生.
新的方向.
我們知道 大家知道.
蘇羅又叫做保羅.
蘇羅就是希伯來的名字.
而保羅就是羅馬的名字.
如果我們讀《士徒行傳》到第十三章.
我們就看到保羅正式向外幫人宣教.
十三章打後.
所以你會發現從《士徒行傳》十三章開始.
蘇羅又叫做保羅.
因為保羅是羅馬的名字.
更適合使用.

$^{41}$所以今天我的講道.
如果你見到蘇羅又見到保羅.
其實是同一個人.
我們有一段聖經.
我們先來看看.
在《士徒行傳》九章.
我先讀給大家聽.
從第一節到九節.
請大家慢慢思想神的話.
九章第一節.
蘇羅仍然向主的門徒.
口吐威嚇兇殺的話.
去見大祭司.
求文書給大馬式的各回堂.
若是找著信奉者道的人.
無論男女都准他捆綁帶到耶路撒冷.
蘇羅行路將到大馬式.
忽然從天上發光四面照著他.
他就撲倒在地.
聽見有聲音對他說.
蘇羅!蘇羅!你為什麼迫罷我?.
他說:主啊!你是誰?.
主說:我就是你所迫罷的耶穌.
起來!進城去!.
你所當作的事.
必有人告訴你.
同行的人站在那裡.
說不出話來.
聽見聲音卻看不見人.
蘇羅從地上起來.
睜開眼睛.
竟不能看見什麼.
有人拉他的手.
領他進了大馬式.
三日不能看見.
也不吃也不喝.
接著十七節.
一路去到二十一節.
我邀請大家一起讀.
十七節開始.

$^{81}$請讀.
阿那尼亞就去了那家.
把手按在蘇羅身上說.
兄弟蘇羅.
在你來的路上.
向你顯現的主.
就是耶穌.
打發我來.
叫你能看見.
又被聖靈充滿.
蘇羅的眼睛上.
好像有輪立刻掉下來.
他就能看見.
於是起來受了洗.
吃過飯就見葬了.
蘇羅大馬式的門徒.
同住了些日子.
就在國會堂裡宣傳耶穌.
說祂是神的兒子.
凡聽見的人都驚奇.
耶路撒冷殘害.
求告者明的.
不是這人嗎?.
並且他到這裡來.
特要捆綁他們.
大道祭司長那裡.
不一樣的人生.
新的方向.
我們先看到第一樣.
是有經歷.
蘇羅的人生.
他有新的方向.
因為他在大馬式路上.
真實地經歷著.
主耶穌向他顯現.
他的人生有新的方向.
成為了另外一個人.
保羅有什麼經歷呢?.
我們一起來找找.
保羅有什麼經歷呢?.

$^{121}$我們先看到.
有光照耀.
第三節.
聖經說蘇羅行路.
將到大馬式.
忽然從天上發光.
四面照著他.
所以才看見.
有大光照耀他.
接著.
與主對話.
第四節.
他就跌倒在地.
聽見有聲音對他說.
蘇羅!蘇羅!.
你為什麼逼迫我?.
他說.
主啊!你是誰?.
主說.
我就是你所逼迫的耶穌.
所以我們看到.
保羅.
與主耶穌有真實的對話.
這是一個對話.
我們再看.
有主的吩咐.
第六節.
起來進城去.
這就是主耶穌.
和保羅說.
起來進城去.
你所當做的事.
必有人告訴你.
所以.
保羅的經歷裡.
主耶穌很清楚地吩咐他.
很清楚地吩咐他.
我們再看.
主對他說.
第七節.

$^{161}$同行的人站在那裡.
說不出話來.
聽見聲音.
卻看不見人.
這個經歷.
是保羅與主耶穌之間的.
一個經歷.
是一個很真實.
個人的經歷.
保羅的生命有很大的改變.
他的生命有很新的方向.
因為他有什麼?.
因為他有很真實.
與主耶穌.
個人的經歷.
我們看他的轉變.
第一第二節.
和十九二十一節.
是一個極大的對比.
我們看第一節.
到第二節.
聖經說.
素來仍然向主的門徒.
口吐威嚇兇殺的話.
去見大祭司.
求文書給大馬式的國會堂.
若是找著信奉者道的人.
無論男女.
都准他捆綁帶到耶路撒冷.
十九節.
看見轉變了.
吃過飯就見狀.
蘇羅和大馬式的門徒.
同住了些日子.
就在國會堂裡宣傳耶穌.
說他是神的兒子.
凡聽見的人都驚奇.
說什麼?.
在耶路撒冷.
撐開求告者名的.

$^{201}$不是這人嗎?.
看見他轉變了嗎?.
所以第一第二節.
蘇羅還是殘害教會.
反對逼迫基督徒.
但到了十九二十一節.
他卻與基督徒同住.
甚至在國會堂裡.
傳講見證耶穌是神的兒子.
二十一節.
凡聽見的人都驚奇.
耶路撒冷撐開求告者名的人.
不是這人嗎?.
弟兄姊妹.
我們信主.
或者我們生命遇見主耶穌之後.
我們的人生都應該有一個.
不一樣的人生.
都應該有一個新的方向.
你和我的人生是否有新的方向.
你知不知道在乎什麼?.
為什麼有些人沒有新的方向呢?.
是在乎你和我有沒有真正經歷過主.
有沒有真正經歷過主.
每個人的經歷都不一樣.
未必需要像保羅那樣.
有沒有光照過來.
開燈就有了.
未必像保羅一樣.
當主耶穌向他顯現的時候.
是有大光直照他.
未必是這樣.
但我們信主的時候.
在黑暗罪惡裡.
你記得你怎樣信主嗎?.
是主而明亮的神聖.
這樣照入我們的內心心田.
讓我們知罪悔改.
每一位應該是這樣.
帶給我們在罪裡面得到赦免和釋放.

$^{241}$保羅和主耶穌的經歷.
剛才提過.
就是他和主有直接的對話.
有直接的對話.
信主之後的你和我.
有沒有恢復和神建立一個.
個人真實的關係.
有沒有啊?.
你和主有沒有真實的關係?.
有啊.
意思就是說.
你和神能不能溝通?.
有沒有溝通?.
有啊.
神有沒有藉著聖靈.
向你的內心說話呢?.
弟兄姊妹.
在一個和神的關係裡面.
一定有和神對話和溝通.
就是說.
我們向主耶穌禱告.
有沒有聆聽到主藉著聖靈向你說話.
甚至在這個關係裡面.
神透過他的聖靈會感動你.
會吩咐你甚至.
該做的事.
該走的路.
或者甚至責備管教你.
有沒有經歷過?.
應該有的.
這是一個正常的關係.
同親關係一定有溝通.
你認識一個朋友.
只有你說話.
從來沒有聽他說過話.
這是溝通嗎?.
這不是溝通.
你又亂說.
這是一個正常的關係.
人和人之間.

$^{281}$一定有相從往來的關係.
這叫溝通.
所以人和神更加是這樣.
所以我們要一同去看看我們的生命.
到底你和神之間.
有沒有溝通到的關係.
神有沒有在禱告裡面.
在和他的關係裡面.
來引導你呢?.
我們和神的經歷.
不僅僅是在物質各方面.
神的供應.
是不是很好?.
神有沒有供應你?.
不然就坐不下了.
但不僅如此.
和神之間的經歷.
不僅僅是神保守我們.
和家人平安.
神有保守我們嗎?.
有,但不僅僅如此.
不是的.
最重要的是.
當然這也是神給我們的祝福.
但最重要的是.
我們和神有沒有一個能夠溝通對話的關係.
你有沒有聆聽到他的心意和吩咐.
有沒有尋求過他的旨意.
有沒有尋到.
生命有沒有因為和神的結連.
這個溝通關係.
而被改變有新的方向.
這個來得重要.
這個改變.
這個新的方向.
可能是一些很大的改變.
我們聽過很多見證.
我們都認識有一位.
以前是黑社會.
做第二把交椅的.

$^{321}$你知道黑社會.
拿著一把刀在街上不斷砍人.
他真的做過那些.
真的做過那些.
後來信了耶穌之後的他.
生命改變了.
得了神的赦免.
他生命改變了.
真的離開黑社會了.
離開黑社會也不容易.
不是你要離開就離開.
你沒離開的時候就你砍人.
你離開的時候就別人砍你.
但是他真的願意離開.
而且離開得到.
後來還讀得成書.
後來還做了一間學校的校長.
厲不厲害.
真的厲害.
你校長是誰.
他以前是黑社會的.
現在變成這樣.
就因為信了耶穌.
這個可能是大的改變.
人生有很大的新方向.
我以前也認識一位姊妹.
我在美國做工做了很多年.
特別在不同的學校我也做過.
在幼兒園,小學,中學.
我不是做老師的.
我做行政的.
我幫校長的.
其中一位老師.
是我的同事.
她是白人.
她一直結婚很多年都沒有小朋友.
我跟她也很好聊.
有一次聊天.
她告訴我.
我英文名叫Barbara.

$^{361}$她說Barbara我告訴你.
我發生的事.
我說好啊.
她說她未信耶穌之前.
是一個很任性的女孩.
白人來的.
很任性的女孩.
很濫交.
很放縱.
很多不同的男朋友.
而且每晚都去不同的男朋友那裡過夜.
你說多放縱呢.
信了耶穌之後.
剛信的時候.
她仍然未能完全脫離這種罪.
因為很多年了.
已經成為了習慣.
但已經信了主.
聖靈在她裡面.
她說她有一個祖上.
突然醒了.
她其中一個男朋友還在她旁邊.
她醒了的時候.
坐在床上.
聖靈責備她.
你在做什麼.
你在這裡做什麼.
所以有時候我們得罪神的時候.
神會在那個時刻.
可能聖靈問你.
你在做什麼.
聖靈問她.
你在做什麼.
她當場醒了醒.
她信了主了.
不過靈命很弱.
剛信.
她覺得很羞恥.
忽然之間.
從來未試過的羞恥.

$^{401}$立刻起來.
離開那個地方.
離開她所有的男朋友.
向神認罪.
聖靈幫助她.
她回到教會.
後來一路離開這些.
人生有一個新的方向.
離開這些.
甚至在教會裡.
神為她預備.
另外一位她的丈夫.
一位很好的基督徒.
接納她 原諒她.
他們有一個新的開始.
她告訴我.
Barbara 因為這樣.
我一生都不可以有疑慮.
這是我的後果.
犯罪帶來會有後果.
這是我的後果.
這些新的方向.
我們身邊很多不同的見證人.
從很大的改變裡有新的方向.
但新的方向不只是這樣.
我又不是黑社會.
我又不是淪落成這樣.
這個新方向.
其實也在於我們每天.
當我們和神溝通對話的時候.
透過祈禱.
透過讀神的話語.
神會改變我們.
有沒有?.
有沒有禱告?.
有沒有祈禱?.
有沒有在聖經被神改變?.
真的有.
流眼淚.
沒有人見過.

$^{441}$不是做給人看的.
不斷流眼淚.
你給我傻傻的.
因為神向的說話.
聖潔的靈向的說話.
我們污穢的自己就反照出來.
所以在那裡.
可能新的方向是這麼簡單.
不是大奸大惡的人.
才有新的方向.
你和我.
其實每天都需要新的方向.
讓神的話神的靈.
改變我們對事對人的看法.
離開自我.
離開神不起悅的性情.
有沒有一些性情.
你覺得很一般.
都很想離開.
有沒有?.
每個人都可能有.
離開神不起悅的.
不只是人不起悅.
人不起悅我們都要檢討一下.
但不只是人.
是神不起悅的性情.
和自我.
和行為.
需要.
有新的方向.
當主耶穌向訴羅顯現的時候.
訴羅問了一個問題.
有沒有留意?.
訴羅問主.
主啊.
你是誰?.
我相信我們都可能.
需要問自己一個問題.
到底.
你覺得主耶穌是誰?.

$^{481}$在你過往的生命經歷裡.
到底主耶穌是誰?.
你有沒有在你生命裡.
你和我.
能夠真實地經歷過祂呢?.
我們很好.
我們有機會參加很多教會.
很好的聚會.
我知道你在節目裡.
參加很多團契.
有小組有團契.
有不同的聚會.
真的很好.
有這麼多機會.
很好.
我們參加很多不同的聚會.
甚至我們參與很多的事奉.
都是很好的.
但是不等於.
你真的經歷過神.
同意嗎?.
重要的是.
你和我有沒有真實地經歷.
你自己才知道.
有沒有真實地經歷.
聖靈引導你進入真理.
最重要的是.
羅馬書說.
我們的心.
聖靈和我們的心同正.
我們是神的兒女.
所以聖靈如果在你的心裡.
已經進入你的生命.
你就一定經歷過聖靈.
因為他的工作是這樣的.
是會引導你進入真理.
所以你有沒有經歷過.
聖靈引導你進入真理呢?.
有沒有經歷過.
聖靈改變你的生命.

$^{521}$而走人生新的方向呢?.
有沒有呢?.
我們很希望.
五十年不變的樣子.
是不是?.
有時候見到一些人.
很多年沒有見過.
一樣的樣子.
為什麼?.
吃什麼防腐劑?.
真的好.
但是.
我盼望我們的生命.
不要五十年都不變.
我相信耶穌.
數一下手指.
十年,二十年,三十年,四十年,五十年.
你的生命怎麼樣?.
也有一點點改變.
已經不錯了.
完全沒有.
我們就擔心了.
有沒有改變?.
這個很重要.
聖靈有沒有因為神的名.
引導我們走二路的時候.
引導你進入祂的真理呢?.
真理的聖靈.
我們需要更多.
鼓勵你和我.
更多的獨自來到神面前.
學習一下.
你知不知道聆聽聖靈的聲音.
要學的?.
不是每個人都完全懂的.
我最初也不懂的.
但是要學的.
要操練的.
要從留意開始.
從留意.

$^{561}$從更多去到神的面前.
學習一下.
留意一下.
聆聽一下.
聖靈的聲音.
真實地去經歷一下.
聖靈在你身上的工作.
未信主的朋友.
或者網上的.
我們需要更多的經歷.
經歷這位神.
尋找祂.
經歷祂.
你和我的人生.
才有新的方向.
對於信了主的弟兄姊妹.
其實我們很想認識神.
我們每個人都想認識神更多.
你知不知道怎樣才能認識神更多?.
就是你越經歷祂.
你就越認識祂.
你都沒有經歷.
我信了主三十年.
三十年之前有.
現在沒有了.
其實我們的靈性不是很正常.
你都不經歷祂.
你怎會認識祂呢?.
人和人之間都是這樣.
你都不和祂交往.
你怎會認識祂這個人是怎樣呢?.
和神都是這樣.
所以我們一定要.
越去經歷這位神.
怎樣經歷呢?.
禱告多些讀字去祂那裡.
和祂聊天.
有沒有和神聊天?.
不只是代求那些.
那些靈計量.

$^{601}$那些叫代禱時間.
我也有的.
另外分開時間.
和祂聊天有沒有?.
只是去尋求.
和祂讀字聊天.
這樣你才認識祂.
接著留意聖靈.
在我們生命裡怎樣.
帶來感動我們.
這樣你要去做.
如果聖靈感動你.
你整天都不做.
聖經說.
我們會麻木的.
聖靈會為你擔憂.
感動來來去去.
你都不動.
聖靈擔憂.
慢慢麻木了.
所以當聖靈感動我們的時候.
我們最聰明.
就是第一時間學一下.
真理的聖靈感動.
那些一定是合真理.
我們就學一下.
很小的事都好.
就去做.
那你就慢慢習慣.
明白聖靈的感動是什麼.
聖靈的聲音是怎樣.
怎樣引導我們進入真理.
這個叫操練.
我們看到.
保羅的人生有新的方向.
因為他真的經歷神.
你和我的人生都盼望有新的方向.
要怎樣?.
要真的經歷神.
我們再看第二方面.

$^{641}$第二方面我們看到.
有反省.
第八節.
蘇羅從地上起來睜開眼睛.
卻不能看見什麼.
有人拉他的手.
領他進了大馬式.
三日不能看見.
也不吃也不喝.
保羅從地上起來.
幾多日?.
三日.
不能看見.
也不吃也不喝.
到底保羅三日在做什麼?.
做什麼?.
請記住我說.
有啊.
十一節.
主對他說起來.
往直街去在猶大的家裡.
訪問一個大數人名叫蘇羅.
他靜什麼?.
禱告.
主對他說這個他.
是指亞拿尼亞.
主耶穌吩咐亞拿尼亞去找保羅.
保羅當時在做什麼?.
主說了.
他靜什麼?.
禱告.
我們相信.
三日.
當保羅盲了眼的時候.
他不吃不喝.
正在禱告.
禱告什麼?.
不說了.
不過.
很合理的推論.

$^{681}$他一定在反省.
因為後來有改變.
主耶穌在大馬式路上向保羅顯現.
對於保羅其實.
是一個很嚴重的反省.
很嚴重.
請問保羅是誰?.
保羅是誰?.
第三節有沒有說?.
有啊.
聖經第三節保羅說.
保羅自己說.
我原是猶太人.
生在基里加的大數.
長在這城裡.
在加瑪列門下.
按著我的祖宗.
嚴謹的律法受教.
熱心事奉神.
與你們眾人今天一樣.
我也曾逼伯奉者道的人.
直到死地.
無論男女.
都所戴下監.
這是保羅被抓的時候.
自己的分數.
保羅的出身.
他的出身是怎樣的?.
他出身是一個非常嚴謹.
宗教背景的人.
他熱心的事奉神.
完全效忠於猶太教.
殘害基督徒的一個人.
他自以為是對的.
直到什麼時候?.
直到復活的主耶穌.
在大馬錫堪的顯現.
他親眼看見主耶穌.
親耳聽見主耶穌.
他才知道.

$^{721}$原來他過往所做的一切.
是錯的.
對於保羅.
你可以想像.
是一個當頭一棒.
甚至當時如果你是保羅.
你會很難接受.
他真的需要安靜反省.
難怪他又不吃又不喝.
他真的需要反省.
反省什麼呢?.
到底發生了什麼事呢?.
到底我以前.
以保羅而言.
以前我做過什麼呢?.
我下一步應該如何行呢?.
聖經說他正禱告.
三天的禱告反省.
之後的保羅.
人生有沒有改變?.
巨大的改變.
有新的方向.
弟兄姊妹你和我當不當反省呢?.
當不當在禱告裡反省呢?.
當有事情發生.
很多事情發生在我們人生.
有開心的.
有很多不開心的.
不如意的.
當事情發生的時候.
我們很容易.
立刻想想如何解決.
正常.
但是或許.
我們可以把握這個機會.
去想一想.
反省一下.
特別在禱告裡反省一下.
我自己有沒有責任呢?.
到底.

$^{761}$神如何看待我呢?.
這件事發生.
到底.
神如何看待這件事的發生?.
有沒有試過這樣想?.
這樣想是不錯的.
有一次我有機會和一位姊妹聊天.
那位姊妹很好.
我們聊到.
人和人相處不是很容易.
你覺得容易嗎?.
不是很容易.
事就容易處理.
但人就很難處理.
所以我和那位姊妹聊了很久.
她說人和人相處不容易.
很容易會有誤會.
而且不知道如何處理.
我問那位姊妹.
那你會怎樣處理?.
她真的很好.
她說我會反省一下.
我會反省一下.
是不是自己的問題.
有沒有這樣反省?.
如果有事發生.
和人有問題.
當然是她的問題.
她要改就要悔改.
可能是的.
你有沒有呢?.
她說我會反省一下.
我真的會去祈禱.
看看聖靈會不會在禱告裡.
給我有什麼說話.
有什麼引導.
我想如果更多人像她.
我相信天下就更太平.
神的家就可能更和睦.
但我相信一個懂得自我檢視.

$^{801}$或者自我反省的人.
就一定很有進步.
自我知道.
有時你有沒有發現和人相處.
可能是我們自己也不明白.
怎麼說都不明白.
你有沒有發現?.
怎麼說都不明白.
我們那次團體分享.
那些叫做「萌塞」.
不懂得寫廣東字.
叫做「萌塞」.
挺好的.
怎麼說都不明白.
如果我們能夠多點去反省自己.
現在發生的事.
或者任何事.
想一想.
我的責任在哪裡呢?.
到底神怎麼看我呢?.
不是人怎麼看你.
神怎麼看我呢?.
到底神怎麼看這件事的發生?.
有什麼意思讓它發生?.
我可以從中得到什麼屬靈的好處呢?.
如果這樣想.
可能我們要健康很多.
我們的生命.
或者進步更多.
在禱告裡面.
我們學習更多反省.
這個不是很多人一起最會的.
做得到的.
有時也會的.
神會透過不同的人去幫助我們.
但亦很重要的是.
你獨自去到神的面前.
人生between神和你.
這個是反省.
保羅透過三天的反省.

$^{841}$他的人生就有很大的改變.
我們再看.
最後一方面就是.
有迴轉.
第九節.
十九節.
吃過飯就見撞.
蘇羅和大馬式的門徒.
同住了西日誌.
就在國會堂裡宣傳耶穌.
說祂是神的兒子.
凡聽見的人都驚奇說.
在耶路撒冷.
撐開求告者名的.
不是這人嗎?.
他迴轉.
保羅透過三天禱告的反省.
他人生有什麼方向呢?.
迴轉.
看到迴轉了嗎?.
經文解釋了.
他迴轉.
三天的不吃不喝.
反省相當認真.
你和我有沒有遇到一些事.
曾經用三天的時間.
甘食禱告.
尋求神的旨意和反省.
我都沒有試過.
一天我試過.
真的甘食一天.
真的尋求一些很重要的事.
但我都未試過三天.
保羅很厲害.
三天的他.
相信保羅就在神面前找到答案.
他有迴轉.
你怎樣知道他迴轉了?.
你知不知道他迴轉了?.
知道.

$^{881}$你怎樣知道?.
有什麼表現?.
第一.
十九節.
聖經記載他肯吃飯.
吃飯是我們每天都吃的.
不是的.
他三天不吃不喝.
後來他吃飯.
代表他明白了.
接受了.
需要一個接受過程.
反省你知道嗎?.
特別是反省完.
你和一些事改的時候.
那時候更加吃不下飯.
但他吃得下.
他接受了.
知道下一步怎樣走.
你才吃得下飯.
十九節.
肯和大馬式的門徒同住.
同住了些日子.
以前他是怎樣的?.
以前捉拿他們.
甚至殺他們.
現在竟然肯和他們同住.
原來大家都知道是弟兄.
二十節.
在國會堂裡宣傳耶穌.
說他是神的兒子.
真是一百八十度的改變.
可以說是一個戲劇性的迴轉.
本來反對耶穌.
現在竟然去傳講耶穌.
是神的兒子.
連周圍的人都無法相信.
像戲劇一樣.
在耶路撒旦殺害.
求告這個名字的人.

$^{921}$不是就是他嗎?.
為什麼現在變成這樣?.
這個就是迴轉.
迴轉是什麼?.
弟兄姊妹.
過往所做的.
經過禱告反省之後.
知道是不對的.
是神所不起悅的.
就怎樣?.
U-turn.
我的開車的人經常都U-turn.
U-turn.
即是什麼?.
不做以前錯的那些.
重新開始做對的那些.
是神所起悅的那些.
又不是為了應付別人.
在聖靈光照底下真是懊悔.
真是反省.
真是改革.
其實是極不容易的.
是要付上代價.
甚至付上生命的代價.
保羅是一個熱心的猶太教徒.
極有前途的年輕人.
得到祭司長的看好.
當時猶太教的人很尊敬他.
甚至取得祭司長的權柄和命令.
去捉拿基督徒.
他當時很大權.
在大馬式路上遇見主耶穌之後.
他整個人改變.
生命改變.
方向改變.
其實保羅的迴轉.
是需要很大的勇氣.
他需要很大的勇氣.
去面對當時所有的人.
包括什麼?.

$^{961}$包括所有的祭司長.
所有的祭司長和猶太教的人.
還要面對當時的基督徒.
起初不相信他.
不信任他.
誰會相信這樣的人會改變呢?.
連基督徒都不信任他.
以前的保羅都是捉人坐牢的.
現在的保羅是怎樣呢?.
被人捉坐牢.
甚至信道.
成為信道的人.
這是什麼?.
這是福音的大能.
有時迴轉是不容易的.
需要付上代價.
很需要謙卑.
但這是福音.
有時如果不迴轉.
不能成就神在你和我生命上.
更大的計劃.
保羅的迴轉不單是付上代價.
而且他跟隨著.
你留意到.
他是非常有彈性和靈活.
他迴轉是非常快.
可以說我們開車來說.
是一個sharp的U-turn.
三天盲了眼.
跟著聖經記載他和門徒住了三日子.
就開始在會堂宣講耶穌是神的兒子.
是一個完全清楚快速的迴轉.
但我們跟隨著.
有時都需要很有彈性.
很靈活.
真理的立場就堅守不變.
是的.
但若過往有些靈媒只是做法.
或者現在明白了真理的做法.
就立刻迴轉.

$^{1001}$很有彈性.
例如我們開車.
有試過走錯路U-turn嗎?.
有嗎?.
有.
你們這麼聰明.
沒有人.
只有幾位.
我相信.
每個人都.
GPS有時都要U-turn.
真的有U-turn.
例如我們開車.
當我們走錯路的時候.
你發現不太對勁.
為何不太相似呢?.
真是走錯路.
你會怎樣做呢?.
你會否說.
不要緊.
只是錯了.
多走十分鐘.
慢慢走.
多走廿分鐘.
有試過嗎?.
當然沒有.
立刻U-turn.
快得不得了.
見到有機會可以U-turn.
你就U-turn.
一樣.
我們跟隨主都應該是這樣.
我再說.
真理的立場不能改變.
聖經的教導.
不能改變就要辭手.
但做法跟真理無關.
搞甚麼呢?.
有時要U-turn.
很有彈性.

$^{1041}$為甚麼呢?.
因為一錯.
不能再錯.
還要多走十分鐘才U-turn嗎?.
這麼傻嗎?.
當然不會.
騙.
很重要.
所以我們要檢視.
其實每天發生很多事.
很多決定.
很多立場的挑戰在我們人生.
我們要分析.
到底是否跟聖經有關.
所以你和我要讀聖經.
不然你怎知道聖經怎樣教你.
甚麼叫辭手呢?.
對,那是真理.
那怎樣呢?.
就像保羅最後殉道.
也不改變.
但有些不關真理事的.
不要浪費時間.
是嗎?.
我們需要.
在我們人生裡.
當然如果還走錯路.
犯罪那些.
就更加要U-turn.
知道錯還不改.
所以需要立刻.
很有彈性.
不要太固執.
甚麼盲塞.
我都不懂寫.
做人不要太盲塞.
一定要在神的面前.
要立刻變通.
聽清楚,不是真理變通.
是你的做法.

$^{1081}$甚至得罪神的地方.
要變通去改變.
大家有沒有想過.
其實神為何要.
揀選保羅呢?.
有沒有想過?.
保羅被神大用.
聖經十幾卷都寫了.
為何呢?.
有三層天的經歷.
為何神要使用他成為宣教士呢?.
向外邦.
這麼大的責任.
委託給他呢?.
為何要選擇他呢?.
有沒有想過?.
保羅到底有甚麼條件呢?.
神要.
這樣在他人生.
進入他人生.
改變他.
保羅原本是.
猶太人.
生在基里加大宋.
在加瑪尼門下.
受著祖宗嚴謹的.
律法授教.
熱心事奉神.
保羅又是猶太人.
又是羅馬人.
他精通幾國語言.
有極高的教育.
在加瑪尼門下.
受教.
精通出來的.
無論在身份.
語言.
學問.
宗教背景.
都是一流人才.

$^{1121}$當然.
神後來真的使用他.
大用他.
在一切神的國度上.
歷世歷代.
神都使用.
不同的人.
有些是平凡的.
有些是.
出眾的.
神都可以使用.
而保羅而言.
最吸引神的.
相信.
不是他外面擁有的東西.
而是他裡面.
有一個很認真的神.
他裡面有一個.
很認真的心.
是完全堅持.
和很認真.
你都相信.
縱使他過往信錯了.
但他忠於.
祖宗猶太教.
那種絕對忠心.
和堅持.
當他後來回傳.
他知道所信的是誰.
他同樣.
保持.
一種忠心和堅持.
真是至死忠心.
堅忍地為主.
行完.
他的路程.
打完他要打的仗.
你有沒有發現.
神保留.
使用保羅.

$^{1161}$原本有.
一種這麼好的性格.
堅持.
堅持的性格.
而將他這種好的性格.
用在神的國度上.
但神今天.
弟兄姊妹.
神今天.
都保留.
一些我們好的性格.
但是我們很需要.
將這些性格.
放在神的手.
讓神去.
改造它.
成為更加合用.
神會.
間算床.
神看人的.
來生.
並不是保羅外面的條件.
去吸引神.
而是他裡面有一個.
很堅毅.
至死.
縱使過往錯了.
但神看中他.
神都可以容易.
一個這麼好的條件.
就是他.
堅持的性情.
神看我們的來生.
迴轉.
就帶給人生.
有新的方向.
今天的分享.
不一樣的人生.
有新的方向.
為什麼保羅的人生.

$^{1201}$有新的方向呢.
因為在大馬色路上.
主耶穌向他顯現.
他有經歷.
他真實地.
經歷著耶穌.
他有反省.
他真誠地.
安靜地.
反省.
最後.
他有迴轉.
他樂於.
吶喊的迴轉.
神.
單尋找.
使用保羅.
神今天仍然不斷地.
尋找合用的人.
在我們寄居的人生.
實在太短了.
難道就這樣就過了我們一生嗎.
你願意.
你的人生更有.
意義和價值嗎.
用我們今天彼此再鼓勵.
更多時間安靜.
去到神的面前.
多去經歷.
讓我們獨處.
更多獨處.
來到我們神的面前.
留意.
聖靈在我們生命裡.
隨時隨刻的感動.
我相信聖靈現在也在感動你.
多去經歷.
多去反省.
亦多去.
迴轉到神的正路上.

$^{1241}$不枉此生.
我們一起祈禱.
我們都是片刻在神的面前.
自己獨處.
每一個來到我們神的面前.
聖靈此刻都感動我們.
引導我們.
天父多謝你賜給我們舊恩.
多謝你在我們人生裡.
主動的尋找我們.
我們的人生.
遇見主之後.
我們的生命.
不再一樣.
我們有新的方向.
在聖靈引導下.
我們去反省.
和你的關係.
我們有沒有真的.
真實地經歷過你.
能不能夠和你溝通.
能不能夠尋求.
明白你的大靈.
讓我們的人生.
不斷迴轉.
到神的正路上.
不枉此生.
求神的靈不斷.
慰藏你的名.
慰藏你的姓名.
繼續引導我們行走.
永生的路.
幫助.
每一位親愛的弟兄姊妹.
你的百姓.
能夠更多.
因為經歷你.
而更多.
真正的認識你.
改變我們的生命.

$^{1281}$奉耶穌基督的聖名.
Amen.
\newpage



\section{約翰福音 3:16-21}
\label{sec:031789}
\textbf{神愛世人}~~~~駱柱成牧師
\newline
\newline
連結: \href{http://www.church.com.hk/acms/content.asp?site=cdc\&op=show\&type=product\&code=031789\&layout=sermon}{\texttt{031789}} ~~~~ 語音日期: 2026-03-08
\newline
\newline
\hyperref[sec:031788]{\small{< < < PREV SERMON < < <}}
~
\hyperref[sec:index_chronic]{\small{[返順時目]}}
~
\hyperref[sec:031790]{\small{> > > NEXT SERMON > > >}}
\newline
\newline
約翰福音 3:16-21
\newline
\begin{longtable}{cl}
\hline
\hline
章節 & 經文 (和合本修訂版)\\
\hline
3:16 & \begin{tabularx}{0.7\textwidth}{X} 「神愛世人，甚至將他獨一的兒子賜給他們，叫一切信他的人不致滅亡，反得永生。 \end{tabularx} \\ \\ \relax
3:17 & \begin{tabularx}{0.7\textwidth}{X} 因為神差他的兒子到世上來，不是要定世人的罪，而是要使世人因他得救。 \end{tabularx} \\ \\ \relax
3:18 & \begin{tabularx}{0.7\textwidth}{X} 信他的人不被定罪；不信的人已經被定罪了，因為他不信神獨一兒子的名。 \end{tabularx} \\ \\ \relax
3:19 & \begin{tabularx}{0.7\textwidth}{X} 光來到世上，世人因自己的行為是惡的，不愛光，倒愛黑暗，這就定了他們的罪。 \end{tabularx} \\ \\ \relax
3:20 & \begin{tabularx}{0.7\textwidth}{X} 凡作惡的人都恨惡光，不來接近光，恐怕他的行為被暴露。 \end{tabularx} \\ \\ \relax
3:21 & \begin{tabularx}{0.7\textwidth}{X} 但實行真理的人就來接近光，為要顯明他的行為是靠神而行的。」 \end{tabularx} \\ \\
[1ex]
\hline
\hline
\end{longtable}
$^{1}$大家早晨 主內平安.
我們又進入新的一周.
元神繼續帶領我們進入這一周.
讓我們過一個土生土長的生活.
今天我講的題目叫做「神愛世人」.
剛才我們讀了經文.
經文有一節是大家很熟悉的.
就是約翰福音的三章十六節.
我想很多弟兄姊妹都會背.
我相信大家不只是會背這一節.
很多經文都會背.
如果只是會背這一節是不夠的.
但這一節也是很重要的.
因為這一節叫做「聖經中的聖經」.
馬丁路德給這節經文叫做「微型的聖經」.
它裡面已經包括了.
上帝最重要最核心的救恩訊息.
所以這一節是很重要的經文.
我們不只是看這一節.
我們也不只是對這一節要熟悉.
我們要對整本的聖經有整全的認識.
求神幫助我們.
讓我們更清楚祂的心意是怎樣.
這個福音的內容是怎樣.
我們上次講到的是尼哥底姆.
前十五節講祂的對話.
十六節就開始進入宣告.
尼哥底姆是法利賽人.
他是受尊敬的宗教教師.
他是一位年老的德高望重.
有經驗的教師.
他是一個謙卑誠實的人.
他願意來到耶穌基督那裡尋求真理.
耶穌基督在跟他交談當中指出.
你所看到的.
你還未進入得到.
你還未明白很多聖經的真理.
在人看來.
尼哥底姆是宗教最高的知識分子.
在神面前他仍然是很膚淺的.

$^{41}$仍然在地上的層次中停留.
所以耶穌基督就指出.
你看到我走的神跡.
看到我與神同在.
你看到我是從神來的教師.
仍然是不夠的.
這些都不夠.
因為你還未看到我的本質是什麼.
我是神的本質.
神的兒子的本質你還未看到.
於是耶穌就繼續跟他說.
除了從天降下.
仍舊在天的人之.
沒有人升過天.
也就是說.
耶穌基督不是從地而來的.
他是從天那裡來的.
從至高之處那裡來的.
在聖經當中.
有人升過天.
但這些都是從地上被接到天上的.
但耶穌基督是從天那裡降下來的.
以利亞和以洛都被神接去.
但他們是從地上而來的.
只有耶穌基督是從天降下.
又會回到天上那裡去的.
因為他本來就是神最高的寶座下來的.
所以耶穌基督宣告一個真理.
你們是從地上的.
你們的地位就算多高都還是屬地的.
我是從天上而來的.
就算我現在的身份怎麼卑微也好.
我依然的來源是從最高那裡來的.
因為我就是神.
我就是神的兒子.
是從最高處那裡來的.
最後就說到.
耶穌基督將華帶到救恩中心那裡.
就說到摩西舉城這件事.
人子也要照樣被舉起.

$^{81}$耶穌基督是從最高之處那裡下來的.
他願意降卑.
說到馬槽那裡.
最低的一個最不好.
最不衛生的地方可以說是最卑微的.
他在地上的時候.
他來到地上的時候.
也要被舉起.
不是被高舉.
是被舉起.
不是受人稱讚.
受人尊榮.
不是給他一個榮耀的冠冕.
而是要掛在木頭那裡.
要被責罵.
被羞辱.
被釘在十字架那裡.
可以說是吊在半空當中.
還沒回到神那裡.
所以十字架是很羞辱的一件事.
就在這個時候.
如果你看其他福音的書就知道.
耶穌基督說了一句.
向父神禱告.
父啊赦免他們.
因為他們所作的.
他們不曉得.
他們自己做的.
他們不知道後果是什麼.
他們不明白.
所以看到耶穌這個禱告的時候.
你就知道.
他的教導他自己是實踐出來的.
要愛你的仇敵.
為那些逼迫你們去禱告.
世人這麼壞,可以說這麼惡.
這麼敗壞.
這個世界也這麼黑暗.
一點都不可愛.
但是耶穌基督.

$^{121}$他在這個情況之下.
他依然願意為人去寫明.
在歷史當中.
你很少看到比這個更大的愛.
十字架那裡.
就顯明神的愛.
這個愛是我們人很難去明白的.
因為在地上的人.
如果是逼迫你的.
是敵對你的.
你通常很難愛這種人.
所以約翰福音三章十六節.
一開始就說了.
就宣告了.
神愛世人.
世人都犯了罪.
虧欠了神的榮耀.
但是神就是愛這些這麼不可愛的人.
一個墮落敗壞.
充滿罪惡的世界.
一個世代.
羅馬書也說到.
因為我們還軟弱的時候.
基督按著所定的日期.
為罪人死了.
為義人死是少有的.
為仁人死.
或者有感作的.
唯有基督在我們重作罪人的時候.
這麼敗壞的時候.
祂就為我們死了.
神的愛就在耶穌基督.
釘在十字架這裡顯明出來了.
神不是要我們變得好好自愛.
成為一個很優秀的人自愛.
而是當我們做罪人.
這麼敗壞這麼墮落的時候.
已經是愛了.
所以整本聖經你會看到.
使徒約翰說了一句.

$^{161}$也是使徒約翰直接說出來.
神就是愛.
所以神就是愛這幾個字.
不只是說神祂有愛.
有愛和示愛根本上是不同層次的.
示是本體性的.
就是說祂就是了.
你有愛我們人也有愛.
但又不是本體性的.
但神就是愛就說得很清楚了.
神祂的本質就是愛.
所以祂流露的屬性就是憐憫的.
是赦罪的.
是有恩典的.
因為祂本質就是愛.
祂的本體就是愛.
所以你會在祂身上看到祂的特質是什麼.
是有憐憫有恩典的.
但是人的愛呢.
人的愛通常就是受情緒.
或者是環境或者是其他人的影響.
如果是一對熱戀的人.
他會說激情的時候.
他就說遇見你就是我最大的幸福.
但是激情過後呢.
他就說遇見你就是我最大的遺憾.
人通常是隨著感覺而變的.
但是神的愛呢.
神的愛是堅定的.
是真實的也是不改變的.
當神愛世人的時候.
不是因為世人這麼敗壞就變了.
祂甚至將祂的獨生子貼了下來.
讓祂走這條實價的道路.
為要拯救這些罪人.
拯救這些這麼不可愛的人.
所以神的愛是給予的.
是給出去的.
在希臘文當中說到愛這個字的時候.
有家人之愛.

$^{201}$doggy.
有情慾浪漫的那種愛.
eros.
也有朋友之間的友誼的那種愛.
philos.
這些的愛呢.
通常都是出於人的感覺.
或者我要得到什麼.
我要有一個什麼的關係.
這些的愛.
跟神的愛是不一樣的.
因為神的愛是付出的.
是願意給出去的.
不是因為要得到什麼的.
是阿加皮的愛.
是不一樣的.
這種愛跟世人的那種愛.
是很大的區別.
人通常會愛一個人的時候.
你有能力或者你有才.
有那個才華或者有那個錢財.
或者外表很吸引.
或者你讓我很舒服.
我就愛你.
但是神的愛完全不一樣.
我記得我剛剛出去機場工作的時候.
打份工作.
要接觸很多的人.
單位有很多不同的人.
因為很大一個單位.
動一動一個部門就過千人了.
所以.
有些人的個性是很難接受的.
是我之前沒遇過的.
有一個女孩我記得很清楚的.
她是男孩子頭.
說話呢.
動一動就擲你了.
而且你一靠近她.
她又會刺傷你的那種感覺.

$^{241}$所以一般人看到她一眼就覺得.
嘩.
這麼癌瘡的那種感覺.
你很難去愛.
所以我當時就想.
哎呀.
不要搞我這些人.
敬而遠之.
不敢太靠近.
但是.
那段時間我就看過一些屬靈書籍.
一直你要.
其實每一種不同類型的人.
他都要神的愛.
他內心都是希望人肯定他.
去接納他.
去愛他.
所以那時候我就問一個問題.
神的愛.
只是愛那些可愛的人嗎?.
我們人通常會愛那些.
讓我們有安全感的.
讓我們有舒服.
感到舒服的那些人.
但是後來呢.
跟這個女同事.
就多些交流的時候.
你就發覺她都有很多優點.
她很擔當.
因為她自小就跟著.
做師長的那個阿爺.
霸氣到不得了.
但是這個人就很有擔當.
所以她依然有可愛的一面.
所以我們人通常的愛是很狹窄.
我們要學習怎樣去接納人.
在舊約呢.
你會看到.
猶太人很難相信.
神是愛其他民族的.

$^{281}$他不容易接受.
神愛其他的民族.
但是神的愛.
不只是給一個民族.
他要淋到所有的人.
所以說.
神愛世人.
我們的生命是需要不斷更新的.
很多時候我們的想法.
是有偏差的.
跟神的真理.
所以我們需要.
回到神的話語那裡.
去接受祂的調整.
我們不只是.
跟著我們的喜好去思想去愛.
我們的意志.
我們的決定.
是要順從聖靈的.
不是被我們自己那種.
個人的慾望所控制.
我們的情感呢.
也是需要被神的真理去聖化的.
我們的心靈是需要更新的.
需要基督的心去更新.
以基督的心為心.
這樣的時候.
我們只愛那些不可愛的罪人.
當然我們愛罪人.
但是要恨惡罪惡.
神最大的禮物是什麼呢.
就是耶穌基督.
上帝愛世人愛到什麼地步.
就是神愛世人.
甚至將祂的獨生子賜給人.
原文神愛世人.
說得更詳細.
因為神如此愛世人.
不單只是一個程度.
如此愛世人.

$^{321}$而且是一種方式.
就是我要將祂的獨生子賜下來.
祂要藉著這個獨生子.
來愛這個世界.
獨生子是什麼意思呢.
就是說這個兒子是獨一的.
是最寶貴的.
是神最親愛的.
這就是神最大的禮物.
聖經用這個字來形容以撒.
阿伯一號不單止有以撒這個兒子.
其實他還有以實瑪尼.
但是這個兒子是神所應許的.
是最獨特最寶貴.
是不可以替代的.
所以獨生子用在以撒的身上.
我們一般人如果對孩子的那種愛.
在地上可以說是.
人經歷當中.
幾乎沒有什麼比父母的愛更深更大的.
因為父母通常付出了一切.
去為著他的子女.
做父母的時候.
兒子長大了他要為子女找學校.
找最好的環境.
為他們預備最好的將來.
如果一些小朋友.
你會看到在美國有一些寄宿家庭.
你會看到他們的子女稍微有些什麼問題.
媽媽的爸爸就會很緊張.
因為不在身邊.
所以是很緊張的.
因為兒子就是人的心.
心肝可以說.
可以為了兒子將心挖出來的.
所以通常說到子女的時候.
就是心肝寶貝.
所以耶穌基督是神的兒子.
但是耶穌基督是很特別的.
他是父神的兒子.

$^{361}$他是從父而生的.
他不是受造者.
人是受造的.
人是被造的.
但是耶穌基督是從父生的.
父怎樣在自己有生命.
就似給子照樣有生命.
耶穌基督他和父都是生命的源頭.
早期一些教父會解釋《尼西亞信經》.
有一句說到.
《尼西亞信經》說到.
從神而生的神.
從光而生的光.
什麼意思呢.
他要解釋這句.
為什麼神和子是一體的呢.
他們不是分開的呢.
這些胡教家有一個站出來說.
當一個人點起一把火的時候.
一個火把已經著的時候.
那個火引著第二把火.
那裡也有火了.
就變成兩個火把了.
第二個火把.
這個火把.
那個火是不是和原先的火是一致的呢.
到底是誰有先呢.
其實.
你就很難回答了.
因為打一把火的火和第二把火的火.
打一把火的火是不是.
一早就有第二把火的火呢.
所以第二把火和打一把火是不是一致的呢.
所以.
就解釋了這件事.
打一把火的火.
和第二把火的火.
第二把火的火.
一早就傳在第一把火的火那裡了.
所以聖旨不是被做的.

$^{401}$是從父而生的.
聖經用了「生」這個字.
聖旨在永恆當中已經是和父一體的.
他一早就和父是一體的.
但是當有子出來的時候.
父相對子.
子也相對父.
所以父稱子為子.
父稱子為父.
父稱子為子.
子稱父為父.
他們本來就是合一的.
因為耶穌基督說過.
我在你的裡面.
你在我的裡面.
所以耶穌基督是神的獨生子.
他們是同一個本質的.
永恆當中就在一起了.
神就將這個獨生子賜給西人.
亞伯拉罕沒有將他的兒子賜出來.
因為神為他預備了羔羊.
但是天父將獨生子賜下來.
本來這個獨生子在永恆當中就和他一起的.
降下來.
很謙卑地在馬槽那裡.
也被舉起.
被羞辱.
被釘死.
而且他在十字架那裡說了.
父說赦免他們.
因為他們所作的他們不曉得.
所以當約翰看到這些的時候.
約翰在這麼多的門徒當中.
是一直跟著耶穌的.
他一直跟著耶穌.
跟到各國他唯一的一個.
在門徒當中.
所以他看到耶穌基督受了極大的痛苦當中.
他依然將神的愛完全地顯出來.
這種愛不是地上有的.

$^{441}$是從天上來的.
是最高的.
可以說是最完全的.
所以約翰在後面的書信裡會看到.
主為我寫明.
從此.
我就知道何為愛了.
真正的愛是犧牲了自己.
成全了別人.
十字架就是神最清楚.
神的愛最清楚的一個顯明.
神賜下祂的兒子.
為何要賜下祂的兒子.
經文說:叫人得永生.
叫一切信祂的不致滅亡.
反得永生.
正因為耶穌基督道成了肉身.
在十字架上成就了神所定的救恩.
祂將天和地打通了.
讓人可以回到父神那裡.
讓地上的人可以從神生.
可以從聖靈生.
也可以從福音而生.
就是我們一直所說的重生.
不單止讓人有重生.
而且可以得著永遠的生命.
得著永生.
神應許的就是永生.
神的恩賜在基督裡面就是永生.
永生不是佛教所說的獵盤.
不是說人死了之後.
修成了正果之後.
好像將火吹熄一樣.
就煙滅了.
所有的痛苦 所有的慾望都止息了.
都沒有了.
生死輪迴也都停止了.
不是這樣的意思.
永生.
是神的生命進入到我們這個生命.

$^{481}$使人和神永遠的同在.
而且在聖情上.
在神的聖情上有份.
所以耶穌基督說.
我活著 你們也要活著.
我在哪裡 你們也要在哪裡.
與神永遠同在.
永生是耶穌和尼哥弟母一直在討論的.
就是尼哥弟母一直需要的.
人怎樣才能得到永生.
約翰福音三章十六主給了我們答案.
因為神如此愛世人.
因為解釋了前面十五節的經文.
為什麼人需要勝利.
為什麼人需要重生.
為什麼聖靈的工作好像風一樣.
讓人去明白這些事.
這一切的根源只有一個.
就是神愛世人.
這是祂很大的一個原因.
救恩不是兩個源頭的.
可以說是一個源頭.
因為經文說到.
神愛世人.
甚至將祂的獨生子賜給他們.
叫一切信祂的.
是不是就兩個源頭呢.
經文再看上下文的時候.
你就知道不是的.
不是說一個是神賜救恩.
一個源頭是人自己去決定.
人為什麼會信呢.
人為什麼會從不信變成信呢.
為什麼會稱耶穌基督為主耶穌呢.
聖經告訴我們不是隨著人的意思.
是因為聖靈的工作.
風隨著自己的意思吹.
所以是聖靈的工作.
經文說到.
凡從聖靈而生的也是如此.

$^{521}$所以是聖靈的工作將人重生.
在救恩的裡面不是兩個源頭.
是一個源頭.
是這個上帝的源頭影響人.
人才可以去接受這個救恩.
是神先工作.
聖靈先光照人.
聖靈先使人去知罪.
聖靈使人重生.
人才可以去相信耶穌基督.
所以當聖靈使人重生之後.
這個人就明白神的奧秘.
他就明白聖經了.
只可以相信從天上而來這位救主耶穌基督.
人的回應.
人是有回應的.
但是重生仍然需要一件事.
就是人要仰望被舉去的耶穌基督.
凡仰望的人就得著生命.
舊約講到耶穌基督就好像摩西舉銅石一樣.
銅石一望人就得著生命.
所以是聖靈先重生.
人信靠一望祂就得著永生.
所以你們得的是本乎恩.
也因著信.
這個不是出於自己.
乃是神所賜的.
所以十六節這節經文不只是抽出來去理解.
是看回前面十五節.
聖靈的工作是怎樣.
耶穌基督被舉起十字架是怎樣.
你只可以解釋這一節經文.
你只可以明白透這句經文.
是從整章來看的.
耶穌基督自己講過.
不是你們揀選我.
是我揀選你們.
我們能揀選是因為神先揀選了我們.
講到王子施保珍講過.
我們如果可以揀選神的話.

$^{561}$就證明祂已經先揀選了我們.
是神的恩典先臨到我們.
如果不是神給了我們生命能夠揀選祂.
人很難去揀選神.
一個智慧的人.
他不單單看到那些零碎的東西.
一個有智慧的人他可以看到全局.
看到首先是甚麼.
次序是甚麼.
我們認識救恩也需要這樣.
看到聖編所講的是甚麼.
你知道救恩的全貌是怎樣.
次序是怎樣.
在基督的裡面.
神在救恩的裡面.
神永遠是主動者.
神差是他的愛子.
經文講到.
神差是他的愛子降世.
不是要定別人的罪.
為甚麼這樣講.
因為在救約的時候.
你會看到神的工作.
很大的程度就是要顯明人是有罪的.
所以救約的律法.
目的就是讓人知罪.
不是讓人因為有律法而去驕傲.
是使人知罪.
律法是顯明他的聖潔公義和良善.
當人看到律法的時候.
就發現自己.
原來我彌神的標準是參天共地的.
律法就好像一份體欠的報告.
告訴你你的身體狀況是怎樣.
你的身體狀況出了很大問題.
他要指出你哪些指標是不達標的.
指出人的問題是甚麼.
但律法不可以醫治人.
就好像報告表不可以醫治人.
真正醫治人的只能夠醫治人.

$^{601}$有技術的醫生.
所以罪的意思是甚麼.
罪的意思就是告訴我們.
我們達不到神所定的目標.
神的目標是甚麼.
我們沒有達到.
達不到這個目標的時候.
這個就是罪.
不是說我們做了多少不好的事.
是因為我們達不到神的目標.
這個就是罪.
所以人如果沒有活出神他要創造的目標的話.
這個就是罪.
律法使人知罪.
當人達不到的時候.
神就派先知來指責人的罪.
讓人去悔改.
但有律法的猶太人的反應是甚麼呢.
他們的反應就是驕傲,恣意.
和輕看外邦人.
甚至將先知殺死.
所以律法是顯明了人的罪.
先知是責備人的罪.
但都不能夠拯救人.
唯有耶穌基督.
這位大醫生真正能夠醫治人.
真正能夠拯救人.
唯有在基督裡面.
所以新約要仰望耶穌基督.
得著這個生命.
就好像舊約的例子.
凡是仰望同舍的就可以得著存活.
所以神的兒子來到這個世界.
不是要定人的罪.
而是要拯救.
只有耶穌基督.
祂只能夠赦免人的罪.
耶穌基督原來是享受天上的榮耀.
但祂願意來到這個充滿罪惡的世界.
降生在馬槽.

$^{641}$人的本性通常是怎樣呢.
不會說從一個不好的地方.
去到一個更不好的地方.
只會是從一個不好的地方.
去到一個好的地方.
因為人有一個天性.
我們中國人也有一句話.
人往高處走.
水往低處流.
當一個人離開尊貴的地方.
他一定是覺得自己.
去到一個比較差的地方.
他一定是很不習慣.
很不舒服的.
他是很難受的.
甚至覺得自己是很卑微的.
因為他從一個好的地方.
去到一個不好的地方.
他甚至是有委屈的.
耶穌基督祂不是.
祂願意信服神的心.
祂本來是富足的.
但卻為我們的原因成為貧窮.
祂的貧窮是因為要我們成為富足.
我們所支持的一位宣教士.
本來在美國已經是公民了.
在香港也都是公民了.
有自己的住處了.
但是當神呼召來到的時候.
祂就願意放下祂的房屋.
安排好祂的家人.
放下原來這麼舒適的環境.
帶著神的愛去到泰國.
服侍那些父母和社會.
沒辦法照顧得好的小孩.
祂就在那裡教導聖經.
讓他們認識神.
讓人從祂的生命當中.
看到神的愛.
泰國天氣很熱.

$^{681}$祂是很怕熱.
但是祂就是願意克服這種難處.
因為祂要將神的愛帶出去.
所以當環境不好的時候.
一般人的想法就是離開.
很少人說我要留在這裡.
讓這個地方變得更好.
可能祂在那裡改變不了社會.
但是因為祂做出了這些.
讓社會好一點點.
基督就是帶著神的愛來到地上.
我們應該效法耶穌基督.
帶著神的愛.
活出神的愛.
讓這個世界.
讓這個社會變得更溫暖一點.
讓人和人之間可以變得更有尊重.
更和諧一點.
基督來到世上.
不是為了定人的罪.
而是為了拯救.
信他的人不被定罪.
不信他的人罪已經定了.
因為他不信神獨生子的名.
如果上帝不是要定人的罪.
世人是不是就沒有罪.
如果人有罪的話.
這些罪什麼時候定的.
是不是等到以後末日審判的時候.
一次算清.
那個時候才定.
福音這一段經文說到.
不信人罪已經定了.
問題是.
不是信什麼.
什麼時候定罪.
是不是因為不信耶穌.
所以才不被定罪.
是不是因為不信耶穌才被定罪.
不是.

$^{721}$因為人犯罪的時候.
已經定罪了.
因為人一犯罪的時候.
他違背了神的致勝律法.
破壞了神創造人的個有形象.
所以罪早就定了.
而且人早就活在罪當中了.
罪已經定了.
那耶穌基督來做什麼.
是不是耶穌基督來不定罪.
人就沒有罪了.
耶穌基督來是要拯救罪.
正如聖經所說.
他差他的兒子降世.
乃是叫世人因他得救.
是不是所有的人都會得救.
是不是因為耶穌基督來.
全世界的人就自然得救.
不是.
因為聖經所說.
不信的人罪已經定了.
如果是自然沒有罪的話.
就不會說不信的人罪已經定了.
不信的人不是現在才定罪.
而是早已經定了.
但是當耶穌基督來到的時候.
當福音來到的時候.
他拒絕基督的時候.
他就繼續留在罪當中.
繼續留在原來的罪當中.
因為除了神的獨生子以外.
沒有別的路可以靠得救.
他拒絕了唯一的救恩.
所以罪就被確認了.
在我們團契當中.
我們有分享到.
有一個癌症的朋友.
弟兄姊妹關懷他很多年了.
上上為他禱告.
聖靈也感動弟兄姊妹.

$^{761}$一次一次地為他傳福音.
但是他就一次一次地拒絕.
所以直到生命最後一刻的時候.
他都沒有接受耶穌基督.
所以當一個人.
生命最後一刻.
都不接受耶穌基督的時候.
他就帶著自己的罪.
進入到永恆那裡去了.
因為他不接受神所預備的救恩.
這個唯一的救恩.
耶穌基督不是定罪的.
在這麼多的門徒當中.
一早就知道猶大要出賣他.
但是他沒有揭露他.
不是立刻定他的罪.
因為耶穌基督是要世人得救的.
他一直忍辱.
直到最後猶大真的決定賣他的時候.
他講了一句.
你要做的事快去做.
猶大賣主不是耶穌基督逼他去做的.
而是他自己去選擇的.
犯罪是人自己選擇的.
不是神預定他去犯罪的.
因為神是一切美善的源頭.
祂不是惡的源頭.
耶穌基督是將人從罪咎出來的.
不是將人推到罪那裡.
聖靈是高舉基督.
讓人過一個聖潔的生活.
不是引導人去犯罪的.
所以人需要為自己的罪去負責.
但是人的得救不是自己的功勞.
完全是出於神的恩典和主權.
是聖靈的工作.
是耶穌基督成就的救恩.
是出於神的計劃和愛.
從天上來的基督進入這個世界.
要彰顯神的愛.

$^{801}$不是定世人的罪.
而是為了救人.
只要這樣想.
耶穌基督來到應該是很受歡迎的.
祂不定人的罪.
祂是要救人的.
應該是很多人支持祂.
但事實是剛剛相反的.
當耶穌基督來到的時候.
有一些反對的人.
為什麼呢?.
不是因為耶穌基督不夠好.
而是因為耶穌基督成為了他們的威脅.
耶穌基督為什麼被釘在十字架上.
因為這些宗教領袖.
他們看不慣耶穌.
耶穌成為了他們的威脅.
連彼得多也說.
我查不到祂有什麼罪.
但是因為祂成為了他們的威脅.
所以彼得多在群眾的壓力之下.
在人心黑暗的環境當中.
祂只可以釘死耶穌基督.
不是因為光不好.
而是因為人喜歡黑暗.
我們在加州陽光很好.
如果你在外面的話.
在現在這個天氣出去曬太陽的話.
你會覺得很舒服.
當我們曬太陽的時候.
你很自然感到身體會暖一點.
因為光通常會產生熱.
光本來就帶著能量.
所以照到哪裡的時候.
照到我們身上的時候.
照到地的時候.
地就會變熱.
照到我們身上的時候也會變熱.
很多時候曬太陽的時候.
我很喜歡和一位長者握手.

$^{841}$握握手是好朋友.
有一次他就說.
你的手很暖.
我說你的手很冷.
因為我的手比較暖.
他的手很冷.
所以他就感受到我的溫暖.
我的溫暖就傳到他的手上.
但是他笑笑口地說.
你的手很暖.
但是我的心裡火熱.
心裡火熱.
其實是真的.
一個年長的人.
他能夠幾十年回來.
很堅持地回來聚會.
是很難得的.
而且年長了.
身體有很多軟弱.
但是他就很喜樂地回來.
很安靜地坐在我們當中.
他說他心裡火熱.
真的.
從他的臉上可以看到.
從他堅持回來的地方.
你會看到.
如果一個人不回來聚會的話.
是很容易冷淡的.
為什麼.
因為我們就好像一團火一樣.
如果我們不常在一起的時候.
很容易會熄滅.
他說他心裡火熱.
所以當他這樣說的時候.
我的心就感染了.
為什麼.
因為他心裡的火傳給我.
我看我的手的溫度傳給他.
他心裡的溫度傳給我.
屬靈的那種火熱也是.

$^{881}$基督徒在一起.
就是要彼此去感染.
當我們有這種熱的時候.
走出教堂的時候.
我們應該將那種熱帶出去.
將那種光顯出來.
我們要將神的愛傳出去.
將基督的福音帶給更多的人.
而不是受別人那種很冷淡.
很冷漠的態度去影響.
影響我們的信心.
影響我們心靈的火熱.
為什麼這麼多人不喜歡光.
因為光一照出來的時候.
人的行為 那種惡行就顯出來了.
你見過晚上有直升機.
在社區巡邏.
或者在高速公路.
有時直升機執行任務.
一有任務打開燈射下來.
如果是找到那些受傷的.
迷路的.
被困的.
有難處的.
那個人就會向他揮手.
就會向著這裡停留.
被光繼續照射.
因為他知道光來了.
那個救援就來了.
因為那個直升機.
他的救援行動已經要執行了.
射下來.
我就要得到這個拯救.
但如果是一個罪犯呢.
如果在晚上在高速公路開車.
飛得很快.
光一照下來.
巡邏隊查到他了.
一照下來.
他的反應是什麼.

$^{921}$他的反應一定是躲起來.
所以他的反應完全是不一樣的.
一個是接受他的光.
一個是躲避他的光.
逃走 躲起來.
光一照下來.
那個罪犯.
他就被顯露了.
他最怕的是什麼.
他的罪被顯露了.
所以光來到世上.
世人因為自己的行為是惡的.
不愛光.
倒愛黑暗.
同一樣的光.
同一部直升機的光.
人有不同的反應.
不是光不同.
是人的心不一樣.
有人當光.
真光照亮他的時候.
他說 主啊.
你是我腳前的燈.
路上的光.
求你引領我走這條正路.
但有些人呢.
當光一照他的時候.
他說 不要照我.
我就是喜歡這個黑暗.
為什麼.
因為太多壞事要暴露出來.
光一照下來.
那個敗壞要顯出來.
耶穌基督來到這個世上.
一定產生兩種反應.
一種就是接受祂.
另一種就是推開祂 拒絕祂.
當我們傳講神愛世人.
傳講耶穌基督的福音.
不是所有的人都喜歡我們.

$^{961}$可能我們就是被人拿掃把.
扔出去的一個.
因為你跟我說耶穌基督.
我世世代代都是拜祖先的.
拜佛祖的.
我抵擋你.
而且信了耶穌.
太多事我不可以做了.
你不要搞我.
有些人呢.
他會接受耶穌基督.
光來到世界.
就顯出人的心是怎樣.
所以神愛世人.
親愛的朋友.
神已經在十字架上的耶穌基督.
將祂的愛向我們顯明了.
願神賜福給你.
光照你.
幫助你認識耶穌基督.
得著神所賜的永生.
我們弟兄姊妹.
願神幫助我們活出神的愛.
靠著我們的主耶穌基督.
在這個世代.
如光照耀.
照亮這個世代.
或者是帶給這個世代.
多一些的亮光.
願神賜福給我們.
\newpage



\section{創世記 29:9}
\label{sec:031790}
\textbf{聖經人聖經事:正說話的時候}~~~~黃克勤牧師博士
\newline
\newline
連結: \href{http://www.church.com.hk/acms/content.asp?site=cdc\&op=show\&type=product\&code=031790\&layout=sermon}{\texttt{031790}} ~~~~ 語音日期: 2026-03-06
\newline
\newline
\hyperref[sec:031789]{\small{< < < PREV SERMON < < <}}
~
\hyperref[sec:index_chronic]{\small{[返順時目]}}
~
\hyperref[sec:031791]{\small{> > > NEXT SERMON > > >}}
\newline
\newline
創世記 29:9
\newline
\begin{longtable}{cl}
\hline
\hline
章節 & 經文 (和合本修訂版)\\
\hline
29:9 & \begin{tabularx}{0.7\textwidth}{X} 雅各正和他們說話的時候，拉結和她父親的羊來了，因為她是牧羊的。 \end{tabularx} \\ \\
[1ex]
\hline
\hline
\end{longtable}
$^{1}$大家好,我們繼續講聖經人及故事.
由亞伯拉罕以撒雅各,現在講到雅各出場.
這是屬於族長時代三人的重要事跡.
雅各的事,還未進入他人生歷史的直路.
現在開始結婚,還未生小孩.
但很多妓女的雅各承繼了他的母親.
李伯家精密的妓女思考,真是天生的.
但其實背後上帝也在掌管.
這次正說話的時候,就有事發生了.
是怎樣發生的?有上帝的手在做事.
《生命說宣教網絡》很想跟大家一起讀聖經.
多些認識.
看一看,創世記29章,第一節到29章.
創世記只有50章,我們可能很快看完創世記.
第29章.
雅各起行,到了東方人之地.
看見田間有一口井.
東方人之地有一口井.
有三群羊伏在井旁.
因為人飲羊群都是用那井裡的水.
而井口上的石塊是大的.
留意一下這裡.
看見田間有一口井.
我們想的田間和他們想的田間是不同的.
我們想的田就是一幅幅田,千麥,中環,裡面有水,稻米.
他們說的田,其實是田野之間.
田野之間,在feel那裡.
這個feel就指田,是又與不是.
田間可能是可能是.
其實應該是指有一個井的空地,一個大的空地,圓野.
大多數都是谷,即是低一點的井.
不然是掘不到水的.
去山頂掘井是掘不到的.
當然是比較低一點的田野地方,低一點的地方才有水.
在那裡有一個井.
井旁邊竟然有三群羊伏在那裡.
臥在井旁,即是伏在那裡.
在那裡做什麼?.
在等待.
為什麼呢?.

$^{41}$一來就拿水給羊喝,叫做飲羊.
讓羊喝了水就走.
為什麼要伏在那裡等呢?.
三群羊這麼多.
原來大家都是用井裡的水,大家很和諧的.
等的.
原因就是井口上的石頭很大塊.
於是要搬開.
於是時時有羊群在那裡聚集.
羊人就把石頭轉來井口飲羊.
隨之又把石頭放回井口的原處.
即是說,井口很大,石頭又很大.
要把石頭搬開.
搬不開.
把石頭推開,轉,轉,轉,推開.
然後才看到水就拿水了.
井口不會有東西掉下去.
大家就轉一轉,把石頭蓋上井口.
以免掉東西,掉羊,掉人.
這樣也好.
大家就等在一起.
說夠了,我們附近有五群羊.
五群都來了,我們就一起推開石頭.
讓他們喝水.
打完水,推開石頭,就散了.
很有趣,他們不會先倒先喝.
是等在一起.
等在一起變成習慣.
可能於是就說.
不如我們等在一起吧.
你這麼小,我這麼大.
不如我們等在一起吧.
也可以聊聊天.
好啊好啊,沒所謂.
於是就是這樣的情況.
再看看.
他們到了巴蘭.
別士巴.
就是這塊地.
就是以色列尼泊加人居住的地方.

$^{81}$尼泊加人就一直往巴蘭走.
到了巴蘭,伯特利就發了一個夢.
夢見一條梯,天使上去下來.
上帝在梯頂上說話.
然後就取了這個名字,叫做伯特利.
然後他們繼續走,走到近哈蘭.
巴旦哈蘭.
巴旦亞蘭.
巴旦就是延伸.
亞蘭就是一個平地.
平地的延伸.
這附近是一個平地.
現在就是敘利亞地方.
以前就和亞蘭人打仗.
就是這個地方亞蘭人.
亞蘭有個城市很出名叫做哈蘭.
哈蘭是什麼呢?再說一次.
最先就是塔拉帶著兒子亞伯拉罕.
和羅德離開吾爾城.
兩河流最發達在這一帶.
和出口這一帶.
這裡不發達了.
在這一帶.
新宇沃土在這一帶.
於是就帶著他們走.
塔拉走到哈蘭就不再走了.
年紀都大了,沒什麼力量.
可能不舒服了,休息了.
哈蘭是叫做交叉路.
就是每條大道都在那裡經過.
哈蘭都很發達.
他們就住在那裡.
做買賣又變成富戶.
所以亞伯拉罕接續.
他父親的心願.
去應許之地即是迦南.
這條路其實很遠.
不過這裡是一個中轉點.
現在阿國回去找舅父.
又回到哈蘭.

$^{121}$因為哈蘭是尼伯加的家.
也是拿學的兒子.
塔拉在那裡.
繼續看下去.
阿國對牧羊人說.
弟兄們.
好像我們今天教會那樣叫.
弟兄們,你們是哪裡來的?.
他們說:我是哈蘭來的.
弟兄們.
其實這個字是兄弟.
有一點點用譯做弟兄.
可能是當日翻譯.
日本人覺得那些應該不是親戚.
不是家族.
所以就叫弟兄.
其實是兄弟.
你知道社團的人喜歡叫兄弟.
其實又是親戚.
又是四海之內皆兄弟.
兄弟是指親戚和家族.
或者和家庭.
兄弟,你們是哪裡來的?.
你知道牧羊人帶羊出去.
走很遠.
可能過夜幾天才回來.
你們是哪裡來的?.
他們說:我們是在哈蘭那裡來的.
哈蘭.
Koran.
Ch是讀K.
Koran.
我們是在交叉路哈蘭那裡來的.
哈蘭.
我想阿國嘔吐了.
眼睛都大了.
就這樣去到一個井.
看到幾群羊伏在那裡.
三群羊.
他就問:兄弟,你們是哪裡來的?.

$^{161}$哪裡人士?.
哈蘭.
我們家就在哈蘭.
主人在哈蘭.
哈蘭.
就是我現在要去的地方.
他就問:哈蘭.
你認不認識拿學這個人?.
拿學有個孫叫拉班.
你們認不認識他?.
拿學是誰呢?.
Lakot.
Lakot是.
哼.
這樣的意思.
哼.
是亞伯拉罕的爺爺.
亞伯拉罕的爺爺.
就是.
你認不認識老爺爺?.
那個伯伯叫拿學.
他的孫叫拉班.
其實那些人是不會認識拿學的.
是會認識拉班.
因為他們從來沒有見過拿學.
他們只認識拉班.
不過就知道.
拉班的先祖叫拿學.
並不是本地人.
是吳爾那邊的人.
哦.
拿學嘛.
拉班聽過.
認識認識.
他們就說.
我們認識他.
很大戶的.
很多樣的.
有錢人來的.
你認識他嗎?.

$^{201}$這樣.
阿伯就說.
他好不好?.
他平安嗎?.
大家知道這個字.
Shalom.
Shalom就是平安.
指健全.
Shalom好不好?.
Shalom就是四肢在.
沒有死.
最近好不好?.
Shalom.
沒有染到疫症.
Shalom好不好?.
很耐心.
很平安.
境界很好.
這個Shalom.
他們就說.
平安.
這個平安是指什麼呢?.
是指一般的.
他好不好?.
拉班.
他好不好?.
挺好啊.
他們回答他.
平安挺好啊.
看一看.
那些牧羊人就說.
看一看.
他的女兒.
拉傑.
另一群羊.
快要來了.
看到嗎?.
那個人就是.
拉班的女兒.
這樣.

$^{241}$拉班我們都認識.
拉班的女兒.
我們都認識.
帶羊的.
這樣.
他女兒.
看一看.
他女兒拉傑.
拉傑.
就是普通話.
拉傑.
怎樣譯出來呢?.
我也不知道.
拉傑.
怎樣譯出中文呢?.
我也不知道.
因為他叫Rachel.
看一看.
CH讀.
CH讀K.
Rachel.
Rachel怎樣譯出拉傑呢?.
Rachel.
你.
不懂.
不懂譯.
譯了拉傑.
我們都認識拉傑.
Rachel.
Rachel原來是.
癡羊.
癡性的羊.
原來這個名字是這樣.
千萬不要改成拉傑.
一個女性.
癡性的羊.
你看看.
拉傑帶著羊來.
這樣.
拉傑你.

$^{281}$是這個.
拉班的女兒.
嘩.
適合我.
我就是想.
要娶拉班的女兒.
因為媽媽和爸爸都是這樣說的.
我舅父的女兒.
嘩.
馬上眼睛起了心.
嘩.
不要了.
於是.
你看看雅各果然是有計之人.
雅各和其他牧羊人說.
遠遠.
很遠.
對面山群羊是拉傑的.
他們很厲害.
一認就認到.
雅各就說.
日頭還高.
即是未曾日落.
通常.
飲羊是日落的.
即是行完一日.
去到井旁.
不要太曬.
給水羊飲.
然後帶走.
現在日頭高.
沒有人打水.
他們可能去做一些等.
不用走.
日頭還高.
不是羊群聚集的時候.
你們不如.
趕快給水羊飲.
再去放一放.
即是再去放羊.

$^{321}$飲完水再放羊.
為什麼叫人.
趕快給水羊飲.
然後再去放一放呢.
你明白了.
即是洗開那群牧羊人.
就出口術.
嘩.
日頭還很高.
還有半日.
不如趕快.
你們飲好羊.
就走.
就可以再放牧羊.
吃多些草.
羊比較好.
想叫羊群走.
他們就說.
不行.
我們不能走.
我們這裡的習慣.
躲藏的習慣.
是不可以的.
是要等羊群聚齊.
齊.
等羊群聚齊.
人先將石頭轉來井口.
才可以飲羊.
沒辦法.
想用妓女叫他們走.
這是第一點.
亞國正和他們說話的時候.
拉傑領著他父親的羊就來了.
因為那些羊是他放牧的.
即是說.
他的父親拉班.
叫拉傑領一群羊去放牧.
即是說拉傑本身的工作是牧羊女.
很有田野氣勢.
喜歡出去放羊.

$^{361}$即是說.
很捱得.
很做得.
很行得.
很曬得.
跟狼打.
帶著羊和狗.
指著各個牧羊傭工做事.
很厲害的.
正在和他們說話的時候.
拉傑就到了.
第一點我想說的是.
有沒有這麼巧.
去到那裡.
先喝點水.
見到路人.
請問哪裡有井.
我想去喝點水.
當時要找井.
因為井是聚集最多人.
可以問路的地方.
就是井了.
於是就說.
哪裡哪裡有井.
這個方向你繼續去吧.
走半天就有井了.
那就去吧.
去到會議.
為什麼有牧羊.
有三群羊伏在那裡.
不喝水呢.
原來井口有塊大石.
原來這個井的習慣.
是大家等在一起才喝的.
拉班認識嗎?.
認識.
我們很熟.
是老闆來的.
對面很遠.
那個就是拉傑.

$^{401}$拉班的女兒.
她帶羊群來.
來到我們就喝水.
這樣.
正在聊天的時候.
拉傑的羊就來了.
方陽的人.
未必每次都去那個井.
有時這個井.
有時那個井.
有時哪條溪水.
是不一定去得到的.
不過正在聊天的時候.
拉傑就到了.
所以你看到.
這個正在.
有沒有這麼正.
其實這個這麼正.
我自己的解釋.
就是上帝的手.
在這裡工作.
帶領他們相遇.
拉傑出現.
這個就是他要做的事情.
然後.
雅各看見.
母狗拉班的女兒.
拉傑和母狗拉班的羊群.
來了.
看到拉傑和羊群.
都過來了.
去到井旁.
第二點.
第二點是什麼呢?.
雅各就不好了.
雅各真的是這樣的人.
他有一個壞習慣.
就是自己會使橫手.
自己不按規矩做事.
只要有利就行了.

$^{441}$不按規矩.
人家都說等在一起.
才開那個井的石頭等羊.
雅各一看見.
哇.
拉傑來了.
和群羊來了.
就上前.
第二點就是.
雅各是一個壞習慣.
不守規矩.
就上前把石頭轉來井口.
來了.
美女.
來來來.
先喝水.
就開始推那塊石頭.
我已經沒有說.
雅各雖然是小孩.
不過天生神力.
一個人推開.
轉開那塊石頭.
還是說.
其他牧羊人都來幫忙.
不知道.
他馬上上前.
把石頭轉來井口.
看這一句.
是雅各自己去推那塊石頭.
又被他推到.
即是說.
弦韌琴.
美女來了.
來來來.
先喝水.
走得辛苦.
自己就去推.
吸引注意.
自己走去.
都不按規矩.

$^{481}$大家是來齊先的.
等在一起.
一起推.
就這樣.
他就上前.
把石頭轉來井口.
然後.
就喝武寇拉班的羊群.
喝完水.
雅各就與拉傑親嘴.
接著放聲而哭.
這個呢.
我想拉傑真的被他嚇一跳.
拉傑本來正在整羊.
又舀水.
喝完水之後.
喝飽了.
先掃羊.
接著見到雅各走過來.
抱住他.
就與他親嘴.
如果現今.
就會說.
打兩巴仙.
非禮.
這樣.
在當日.
就沒有的.
親嘴的動作.
就.
這樣.
他們就不是很.
不是很什麼的.
就覺得OK的.
當然.
親嘴是表示問安的.
譬如.
《西藥聖經》說親嘴問安.
無要聖潔.
不要想歪念.

$^{521}$親嘴問安.
親嘴.
表示問安.
喝完羊.
這位女士好嗎.
親嘴.
這樣.
面貼著.
接著哭.
哇.
拉傑.
什麼事呢.
這位先生.
沒事吧.
這樣.
雅各就告訴拉傑.
其實我是誰誰.
自己是他父親的外甥.
我是你爸爸的外甥.
你爸爸叫拉班.
對了.
拉班是不是有一個妹妹.
叫利伯加.
是的.
我小時候.
都見到她.
後來就嫁了人.
嫁了二三就生了我.
哦.
你是利伯加的兒子嗎.
是的.
什麼.
你是利伯加的兒子嗎.
哇.
馬上拉傑就跑去.
告訴他父親.
馬上跑回去.
就告訴拉班.
拉班聽見.
我的外甥雅各來了.

$^{561}$要跑出去迎接.
先看看.
抱著他與他親嘴.
親嘴是文安.
接受.
就把他帶回自己家.
雅各就把一切的情由.
告訴拉班.
哎.
舅父.
我告訴你.
我這次出來是著草的.
因為得罪了哥哥.
得罪了爸爸.
這次都不得已走出來.
不過這事是媽媽教我的.
我都聽了她說.
拉班就說.
不要緊.
其實不要緊.
你做了什麼.
第三點.
拉班很顧情.
你實在是我的骨肉.
雅各就和他住了一個月.
我想看看.
你實在是我的骨肉.
是什麼意思呢.
原來他用了兩個字.
骨肉是一樣東西.
我親生骨肉.
不是骨和肉.
但他們說是骨和肉.
是我親生骨肉.
骨肉的意思.
A gem.
A gem是骨.
你是我的骨.
和身體.
肉就是你是我的身體.

$^{601}$你是我的骨和我的身體.
即是說.
你和我是一個淵源.
一個身體.
一棚骨.
一起的.
於是就.
一起住下先.
於是就住了一個月了.
神又安排這個拉班.
沒有看低他.
因為拉班可以看低他嗎.
哎呀.
你伯家的兒子.
一定我會照顧的.
不過心想.
這個臭小子逃跑出來.
採草的.
都不放好兒子.
這樣欺騙爸爸.
欺騙哥哥.
這樣.
幸好我不會聰明.
就叫他來.
否則就死了.
這個兒子都沒什麼用.
逃跑的.
這個外甥沒什麼用.
但是呢.
那你住下就拿些禮物回去吧.
這樣.
沒有啊.
他留在這裡住吧.
你和我是骨肉一樣.
我覺得第三點.
都是上帝祝福.
為雅各預備一個酒瓶佬.
著草的.
這個妻身之所.
就在舅父拉班那裡.

$^{641}$這樣就住下了.
好.
看看今天的金句.
我想特別說的是葉和華做的.
就是雅各正和牧羊人宣話的時候.
拉傑領著羊就來了.
哇.
我覺得這個有沒有這麼巧.
又會去到那個井.
又會遇到一些認識拉班的牧羊人.
然後又遇到拉傑剛剛又來.
又遇到他的舅父拉班.
是歡迎他.
這個正在失敗的失敗者的雅各.
其實完全是上帝的安排.
所以我用這個字正.
正不正就是說剛剛好.
怎麼會有這麼多剛剛好呢.
幾件事一定是上主安排的.
我們一起祈禱吧.
多謝我們的主.
知道你對我們人生是有安排的.
無論我們去到高去到低.
有時候真的去到一條絕路.
有時候真的要逃跑著草離開.
或者什麼情況.
去到這裡.
好像雅各一樣隻身走出來.
衣服都不多.
水都沒得喝.
帶一點點魷魚餅在路上.
要塗一下就出去了.
沿途還要找東西吃.
又沒有帳幕.
他就一直這麼遠路.
就去到哈蘭附近.
就遇到所有要遇到的人.
這麼巧嗎.
完全是上帝的作為.
我們稱讚你.

$^{681}$你在背後一路有掌管.
一路有供應.
一路有帶領.
一路有祝福.
願你的名得在榮耀.
你是我們的主.
帶領我們的主.
是我們的牧者一樣.
讓我們常常仰望你.
等候你的帶領.
你有時候帶我們在青草地.
有時候在溪水旁.
有時候經過死暗幽谷.
不過用你的象羽竿安慰我們.
讓我們更有信心.
我們讚美你.
祈禱奉耶穌名求.
阿們.
\newpage



\section{馬太福音 18:21-35}
\label{sec:031791}
\textbf{未被寬恕?}~~~~曹偉彤牧師博士
\newline
\newline
連結: \href{http://www.church.com.hk/acms/content.asp?site=cdc\&op=show\&type=product\&code=031791\&layout=sermon}{\texttt{031791}} ~~~~ 語音日期: 2026-03-01
\newline
\newline
\hyperref[sec:031790]{\small{< < < PREV SERMON < < <}}
~
\hyperref[sec:index_chronic]{\small{[返順時目]}}
~
\hyperref[sec:031792]{\small{> > > NEXT SERMON > > >}}
\newline
\newline
馬太福音 18:21-35
\newline
\begin{longtable}{cl}
\hline
\hline
章節 & 經文 (和合本修訂版)\\
\hline
18:21 & \begin{tabularx}{0.7\textwidth}{X} 那時，彼得進前來，對耶穌說：「主啊，我弟兄得罪我，我當饒恕他幾次呢？到七次夠嗎？」 \end{tabularx} \\ \\ \relax
18:22 & \begin{tabularx}{0.7\textwidth}{X} 耶穌說：「我告訴你，不是到七次，而是到七十個七次。 \end{tabularx} \\ \\ \relax
18:23 & \begin{tabularx}{0.7\textwidth}{X} 因為天國好像一個王要和他僕人算賬。 \end{tabularx} \\ \\ \relax
18:24 & \begin{tabularx}{0.7\textwidth}{X} 他開始算的時候，有人帶了一個欠一萬他連得的僕人來。 \end{tabularx} \\ \\ \relax
18:25 & \begin{tabularx}{0.7\textwidth}{X} 因為他沒有甚麼償還之物，主人下令把他和他妻子兒女，以及一切所有的都賣了來償還。 \end{tabularx} \\ \\ \relax
18:26 & \begin{tabularx}{0.7\textwidth}{X} 那僕人就俯伏向他叩頭，說：『寬容我吧，我都會還你的。』 \end{tabularx} \\ \\ \relax
18:27 & \begin{tabularx}{0.7\textwidth}{X} 那僕人的主人就動了慈心，把他釋放了，並且免了他的債。 \end{tabularx} \\ \\ \relax
18:28 & \begin{tabularx}{0.7\textwidth}{X} 那僕人出來，遇見一個欠他一百個銀幣的同伴，就揪著他，扼住他的喉嚨，說：『把你所欠的還我！』 \end{tabularx} \\ \\ \relax
18:29 & \begin{tabularx}{0.7\textwidth}{X} 他的同伴就俯伏央求他，說：『寬容我吧，我會還你的。』 \end{tabularx} \\ \\ \relax
18:30 & \begin{tabularx}{0.7\textwidth}{X} 他不肯，卻把他下在監裡，直到他還了所欠的債。 \end{tabularx} \\ \\ \relax
18:31 & \begin{tabularx}{0.7\textwidth}{X} 同伴們看見他所做的事就很悲憤，把這一切的事都告訴了主人。 \end{tabularx} \\ \\ \relax
18:32 & \begin{tabularx}{0.7\textwidth}{X} 於是主人叫了他來，對他說：『你這惡奴才！你央求我，我就把你所欠的都免了； \end{tabularx} \\ \\ \relax
18:33 & \begin{tabularx}{0.7\textwidth}{X} 你不應該憐憫你的同伴，像我憐憫你嗎？』 \end{tabularx} \\ \\ \relax
18:34 & \begin{tabularx}{0.7\textwidth}{X} 主人就大怒，把他交給司刑的，直到他還清了所欠的債。 \end{tabularx} \\ \\ \relax
18:35 & \begin{tabularx}{0.7\textwidth}{X} 你們各人若不從心裡饒恕你的弟兄，我天父也要這樣待你們。」 \end{tabularx} \\ \\
[1ex]
\hline
\hline
\end{longtable}
$^{1}$在一月份的講道當中,我是提及.
我們每個人都是經歷著不同種類的創傷.
我們患病,我們家人患病,離世.
人的心靈患病.
人在傷害人.
人用言語暴力,是彼此的傷害.
當你面對這樣的傷害的時候.
你能不能夠寬恕那些曾經傷害你的人呢?.
我們跟心禱告.
親愛的主,我們感恩因為.
我們眾弟兄姊妹能夠來到聖所當中.
去朝見你,敬拜你.
我們如今,我們是預備.
是拳手主餐.
在主餐之前,我們一同聽你的話語.
我們又知道在這個月,是我們的關顧月.
我們是要去關心.
在我們的家人,承浸家的家人.
不同地方我們淪捨的朋友.
我們又知道當我們要關心其他人的時候.
我們的心靈需要健康.
主,我們今天就求你幫助我們.
用你的話語來醫治我們.
讓我們能夠進入你的話語的奧秘當中.
在你的能力當中.
得蒙醫治,得蒙釋放.
主,我們仰望你,我們愛你.
我們的禱告,是奉主耶穌基督名字求,阿們.
今天我會跟大家讀耶穌基督的比喻.
什麼叫做比喻呢?.
比喻的希臘文的基本字根就是.
Para-Polay.
Para就是旁邊,Polay就是頭窄.
意思就是說,將一樣的東西放在另外一樣的東西旁邊.
我們看到耶穌基督在祂說比喻的當中.
祂將兩件事情放在一起.
好像拋在半空中一樣,讓人看見.
一方面,祂拋出重要的真理和概念.
是關乎神國度的願景.
另一方面,祂拋出一些人物,一些日常的事情.

$^{41}$祂拋出的人物是什麼人物呢?.
男人,農夫,牧羊人,麵包師傅.
父親,兒子,新娘,伴娘,銀行家等等.
這些人都是從事著很平凡的事務.
他們都有家庭的糾紛.
有些人在尋找迷失的羊.
有些人在計算工錢.
這些都是比喻的方式.
為什麼耶穌基督要用比喻呢?.
因為當我們談論神,我們認識神的時候.
我們會發覺有相當的限制,有相當的困難.
保羅告訴我們,人類是有限的.
我們只可以透過模模糊糊的鏡子去看見神.
再者,我們又常常忘記了我們的有限.
我們傾向以為我們能夠把握神.
甚至完全地明白神.
所以當我們開始談論神的時候.
我們會想像一個神明能夠滿足我們的需求.
我們的慾望,我們的喜好.
這些神明的形象就是偶像.
所以比喻是提供我們一些對抗偶像崇拜的解藥.
叫我們不要以為自己這麼厲害,知道神的一切.
另一方面,叫我們能夠進入神的奧秘,神奇妙的恩典當中.
所以不少時候我們看見耶穌基督.
是會透過比喻打開我們有限的智慧.
讓我們進入一個不同的世界.
經歷一個受天國之道所塑造的奇妙世界.
耶穌基督又指出.
這個世界的角度,神的角度,神的角度裡面的真理.
它的規則是和這個世界不同的.
門徒要採取一種別的方式,不同的方式去看世界.
是能夠經歷這種由天國的道理,真理所塑造的奇妙的世界.
要明白神的智慧和奧秘.
我們今天中午看《馬太福音》18章21至35節.
那時彼得進前來對耶穌說:主啊,我弟兄得罪我.
我當饒恕他幾次呢?到七次可以嗎?.
耶穌說:我對你說,不是到七次,乃是到七十個七次.
天國好像一個王要和他的僕人算帳.
才算的時候有人帶了一個欠一千萬銀子的來.
因為他沒有什麼償還之物.

$^{81}$主人吩咐把他和他妻子兒女並一切所有的都賣了償還.
那僕人就服服拜他說:主啊,寬容我,將來我都要還清.
那僕人的主人就動了持心把他釋放了,並且免了他的債.
那僕人出來遇見一個同伴欠他十兩銀子.
便抽著他,合著他的喉嚨說:你把所欠的還我.
他的同伴就服服央求他說:寬容我吧,將來我必還清.
他不肯,竟去把他下在監裡,讓他還了所欠的債.
那同伴看見他所做的就甚憂愁,去把這事都告訴了主人.
於是主人叫了他來,對他說:你這惡奴才.
你央求我,我就把你所欠的都免了.
你不應當連恤你的同伴,像我連恤你嗎?.
主人就大怒把他交給掌刑的,讓他還清了所欠的債.
你們各人若不從心裡饒恕你的弟兄,我天父也要這樣待你們了.
對我來說,人類的歷史一直以來似乎有一個鮮明的主題.
這個鮮明的主題就是復仇.
它在我們的文化當中,在我們的政治當中,在我們的社會當中,在教會當中.
我們每一個人都是在一場復仇的旅程當中.
不少時候我們跟別人辯論爭執,然後生氣,然後憤怒.
在我們周圍我們聽見報復配偶,員工報復上司,學生報復老師.
孩子報復爸爸媽媽,鄰居報復鄰居的故事.
為什麼人會報復呢?.
心理學家就說,當我們受苦的時候,我們當然要找出造成我們痛苦的原因.
所以我們要找出誰搶我們的生意,誰在工作上羞辱我.
誰詆毀我的人格,誰在網上抹黑我.
誰在霸凌我,誰在霸凌我的孩子.
誰讓我受苦,誰讓我受創傷.
我們要找出那些害我們的人,要他們付出代價.
所以我們會紅遍滔滔地說,如果你冒犯我,我怎會不報復呢?.
你可以讓我惡行,我又要還給你,而且是更勝一籌.
你可以叫我受苦,我會讓你受更多的苦.
意思就是說,作為一個人,有人性的.
報復是天經地義,甚至要加碼報復和償還.
要讓對方承受更多的痛苦,以致能夠彌補我的痛苦.
相對世界的方式,耶穌基督就指出.
在神的國度裡,我們以一種與社會不同的方式來生活.
可以經歷一個受天國的真理所塑造的奇妙的世界.
耶穌基督就指出,你們不應該復仇.
神是向那些不配的罪人施恩典和憐憫.
神是不會報復的.
所以今天我們再進入《馬達福音》十八章二十一至三十五節的經文.

$^{121}$這段經文的比喻是講出無限寬恕的真理.
耶穌基督就這樣說.
好象一個王要和他的僕人算帳,僕人的希臘文叫doulos,是解釋甚麼?是奴隸.
這位王要和他的僕人結清所有的債項.
這個僕人,這個奴隸是欠了一千萬的銀子.
一千萬等於國王的身家.
所以意思是不可能償還.
這位奴僕就向王乞求廉恤和寬恕.
俯伏在地就說主啊,寬容我,我一切都還你.
但是面對這麼大的巨債,這條數是不可能還清的.
但是那個僕人的主人就動了持心,將他釋放了,免他一切的罪債.
其實我們突然間看到,巨債,王的廉恤,這個不是關於金錢的問題,這個是關於罪的問題.
這個就是比喻的重點,就是說人剛犯神,人對神的得罪是無法計算的.
所以這個比喻不是單單說一個慷慨的主人能夠寬容他的奴僕.
而是在說神赦免我們莫大的罪惡.
而王廉恤這個僕人的一個要求,一個心願,一個願景是什麼呢?.
就是希望這個僕人饒恕寬容那些虧負了這個僕人的人.
因為王的憐憫和寬恕,這個僕人很開心,邊走邊大聲謝謝.
王啊,謝謝,謝謝,很開心.
不過在路上,他就碰到一位欠他十兩銀子的同事,十兩銀子就等於一個工人一百天的工資.
十兩銀子其實都不是一個小的數目,不過比起一千萬這個天文數字就算不得什麼.
這條數是可以還清的,這個欠債的同事就懇求這個僕人寬容他.
但是這個已經被寬恕的僕人是不願意去寬容.
他抽著這個同伴的喉嚨對他說,你將所欠的還我.
於是就將欠債紙交給他,你快點立即還,他站在權力上訴諸法律的行動.
將這個同伴監在監獄裡面.
看見這一幕,你們應該不會驚訝的,應該都不會驚奇的,是不是?.
儘管我們很多時候都會做錯,但是我們希望別人原諒我們,原諒我吧,原諒我吧.
不過我們不是很想原諒別人的,原諒別人可以嗎?.
要我們原諒我們的敵人可以嗎?.
我們總是抓住不放,抱著我們心中的不平.
我們腦海裡面有一本素寶,記滿了我們面對所經歷的羞辱和痛苦.
我們是要報復.
所以我們聽見人家叫我們去寬恕那個人的時候.
我們不禁心裡自問,不知道他在說什麼.
他有沒有看過郵箱裡面那些欺凌我的惡毒的電郵.
他有沒有在網上被人欺凌過.
他有沒有見過我的老闆,他有沒有見過我的同事.
他有沒有認識我爸爸,他認識我媽媽,他有沒有跟我妹妹相處過.
所以我們面對我們的敵人的方法是什麼呢?.

$^{161}$我們喜歡的方法是平衡.
我們要醫治我們創傷的方法是什麼呢?.
是平衡.
意思就是說,我受了苦,你也要受苦.
你要付上沉重的代價.
但是當我們繼續看耶穌基督的比喻的時候.
我們看見當國王聽見這位僕人領受了這麼大的恩典之後.
他不願意給他同伴一點兒的恩惠的時候.
主人就叫僕人來.
主人就說,你這惡奴才.
你要求我,我就把你所欠的都免了.
你不應當連戳你的同伴,就讓我連戳你媽.
十八章三十三節連戳出現了兩次.
所以一節經文的重點就是說,要施憐憫.
既然第一個僕人已經領受了憐憫.
他應該要效法主人去憐憫其他的人.
現在憐憫是什麼呢?寫著當憐憫.
意思就是應該,必然,不是一個非可有,可無的選擇.
是一個必然的選擇.
所以比喻的基本信息就是說.
當我們經歷到神的連戳的時候.
我們必須向人展現這種連戳.
最終王不單止將第一個的僕人送在監牢裡面.
還將他交給掌刑的人.
要去施刑罰,直至他還清所有的債.
當我們看見一切的時候,你覺得會不會驚訝?.
你會不會說,不公平,真的嗎?.
這個比喻是真的嗎?合理的嗎?.
我們這次不要樹倒,我們就是要進入監牢.
我們不可以有第二次機會嗎?.
難道神都會展開復仇之旅嗎?.
弟兄姊妹們,我們要知道耶穌基督的比喻是真實的.
雖然真是令我們費解.
但祂是想說出天國的奧秘和奇妙.
所以在比喻當中,耶穌基督就是打開我們有限的智慧.
讓我們能夠想像並且進入神奇妙的世界當中.
所以耶穌基督指出的是那些不能夠寬容他人的人.
毫無疑問,他們會經歷靈魂的折磨.
相反,如果經歷了神的連續的人.
又已經被赦免了一千萬銀子的人.

$^{201}$又去寬容那些欠了你十兩銀子的人.
這些人是可以繼續享受神那種滿有愛和憐憫的生活.
是活在一個由天國度理所造造的奇妙的世界.
在神的愛,在神的大能的裡面.
我有一個朋友在四十年前認識了另一位先生.
他很欣賞這位先生的才華.
他們曾經同工過一段時間.
後來這位先生在外國留學和做工.
不過他沒有跟我朋友說他在那裡也是做工的.
所以我朋友知道他在那裡讀書.
他知道之後覺得這位先生不誠實.
他很不喜歡.
四十年來他仍然為此憤恨.
每幾年他會在網上點名提一提這位先生.
在十年前,在網頁那裡.
我見到我朋友又再公開點名批評這位先生.
那時我心裡想,已經三十年了.
三十年這位先生有沒有改變呢?.
有沒有蒙恩呢?有沒有悔悟呢?.
我的朋友有沒有跟這位先生見面,知道他的情況呢?.
他仍然假設這位先生沒有悔悟.
而這位先生他好像沒有說出真相.
這些事情是否仍然值得在網上繼續點名批評呢?.
我見到人的怨恨可以存放幾十年了.
我另一個朋友提到他跟媽媽的關係.
我獲得他的批准可以在這裡討論.
他亦在其他場景提及他跟媽媽的故事.
在2011年5月初.
一個星期五的早上.
他接受一間大學頒發給他榮譽博士的學位.
他買了一件禮服.
回到原生的家庭給父母,兩個妹妹和其他家人看.
他拿這件禮服出來的時候.
媽媽坐在沙發上就爆了.
媽媽說這個被寵壞的人永遠都是「我我我」.
你經常以為自己比別人聰明.
所有人當然停在那裡.
我的朋友就跟過往一樣.
放棄了回應和反擊.
站在媽媽的前面.

$^{241}$他想起十歲的時候.
他拿著A+作品給媽媽看的時候.
媽媽就撕爛了他.
媽媽會扯他的頭髮.
媽媽有些時候會大聲地罵他.
媽媽會用指甲去刺他的皮膚.
他知道媽媽是一個很憤恨的人.
不能夠寬容人的人.
他的怨恨是可以持續很多年的.
一旦你做了一些事情他不喜歡的話.
他就會生氣.
如果他認為你傷害他的話.
他不會原諒你.
在他生命裡面是一個報復的名單.
是說出那些應該受懲罰的壞人.
他們應該要接受甚麼的報復.
後來經過一段日子.
我朋友那時候已經五十多歲了.
他也是愛媽媽的.
雖然和媽媽的關係很艱難.
他也會去探望媽媽.
和媽媽喝咖啡.
有一天媽媽的心情很好.
所以媽媽就和他聊故事.
聊一些惡人的故事.
到了某個時候.
我朋友就大膽地問他媽媽.
那麼多惡人你會不會其中的一個.
都在名單上呢.
媽媽回答說我有啊.
我朋友就說.
媽媽你會不會憐憫那些惡人啊.
媽媽回答說不會.
我不會憐憫他們.
我也不會憐憫自己.
我不會改變我腦海當中的名單.
我不想失去我的記錄.
這就是媽媽的本性.
媽媽的帳簿裡的債務.
不斷地折磨這位母親.

$^{281}$這個故事讓我想起.
另外一位神學家所談及的一幅圖畫.
他叫做Howard Thurman.
他寫了一段文字我想和大家分享.
他說在我們每個人裡都有一片內在的海.
在那片海上有一個島.
在那島上有一座殿堂.
在那殿堂裡有一座祭壇.
在那祭壇上燃燒著一團火焰.
我們每個人都被召喚去踏上旅程.
找到我們內在的海.
找到我們內在的島.
發現那座殿堂.
並在殿堂的祭壇前遇見生命的神.
不過他繼續說.
他說當人去尋找心靈的海島嶼的時候.
發覺這個島嶼已經被很多很多的錯誤去淹沒.
這些錯誤是什麼呢?.
是別人對我們的傷害.
我們對別人的傷害.
這個世界對我們的傷害.
所以在看那島的時候.
Thurman就說.
人發現裡面有很多的雜物.
這些雜物同樣埋藏著遮蓋著那個祭壇.
人於是要找那個祭壇.
就將這些雜物移開.
不過當人要拿走這些雜物的時候.
人又覺得不捨得.
因為這些雜物是我們熟悉的.
這些雜物是告訴我們我是誰.
這些雜物是我自身的現身.
沒有了這些雜物.
我自己都不知道我是誰.
再者就是說這些雜物當中所埋藏著的那種憤恨.
那種仇恨.
是一種力量來幫助我們對抗另外世界的邪惡.
以致我們繼續有力量去保護自己.
以致我們能夠感覺到安全.
不過Thurman就說.

$^{321}$這些雜物當中的憤怒仇恨是沒有幫助我們抗衡這個世界的邪惡.
他們只是殺害我們.
最終毀滅我們.
面對這些生命的現實.
以我們恐懼的情況.
Thurman說我們要去選擇.
選擇會不會尋找那個祭壇.
邁向神的恩典.
他意思是唯有當我們繼續尋找這個祭壇.
移走這些雜物和仇恨的時候.
我們才能夠經歷到神的恩惠.
只有這樣的話我們的生命才能夠被改革.
繼續經歷神那種無限的恩惠.
而不是在憤怒憤恨.
當人能夠移走的時候.
這就是寬恕.
寬恕能夠移走那些雜物.
使我們再有自由去愛人.
能夠接受愛.
所以Thurman就說.
在那祭壇前.
你靈魂最深的意圖赤裸的呈現.
在那祭壇前你聽到神的聲音.
提升你的靈魂.
赦免你的罪.
赦免你的承諾.
是的.
好像Thurman所說.
當我們能夠找到我們生命的祭壇.
沒錯,你和我有缺點.
你和我也有卑劣的念頭.
你和我也有憤恨,有怨恨.
這一切都會被挪走.
我們永遠在永生神的愛當中.
我們能夠去寬恕其他的人.
能夠我們一起去寬容那些欠我們十兩銀錢的人.
這樣的話.
我們就永遠在神的寬恕憐憫當中.
我們繼續的繼續多謝.
神竟然寬恕了我們一千萬銀子的罪債.

$^{361}$親愛的弟兄姊妹們.
沒錯,我們是一同繼續學求神對我們的憐憫.
人對我們的憐憫.
這是我們喜歡的.
但今天這個比喻是很奇特很嚴苛的.
是跟我們說.
不過這是一個美好的訊息.
它提醒我們.
耶穌基督絕不會讓我們忘記寬容他人的重要性.
所以寬容他人是重要.
寬容他人是必須.
寬容他人是應該.
所以耶穌基督說出天國的願景.
你們應當連恤你們的同伴.
像神連恤你們一樣.
你們各人若不從心裡饒恕你們的弟兄.
天父也要這樣待你們了.
有位神學家Louis Smith這樣說.
當我們饒恕他人的時候.
我們釋放了一個囚犯.
而我們發現這個囚犯就是我們,是你們.
所以弟兄姊妹們.
就讓我們今天領受了訊息.
一起進入比喻的奇妙的世界.
我們的心裡,常常從心裡.
我們寬容他人.
饒恕他人的過犯.
饒恕他人欠我們十兩銀錢.
我們一起去做.
再次去做.
做到七十個七次.
阿們.
\newpage



\section{民數記 6:1-27}
\label{sec:031792}
\textbf{上乘之選}~~~~吳周婉貞傳道
\newline
\newline
連結: \href{http://www.church.com.hk/acms/content.asp?site=cdc\&op=show\&type=product\&code=031792\&layout=sermon}{\texttt{031792}} ~~~~ 語音日期: 2026-03-07
\newline
\newline
\hyperref[sec:031791]{\small{< < < PREV SERMON < < <}}
~
\hyperref[sec:index_chronic]{\small{[返順時目]}}
~
\hyperref[sec:031793]{\small{> > > NEXT SERMON > > >}}
\newline
\newline
民數記 6:1-27
\newline
\begin{longtable}{cl}
\hline
\hline
章節 & 經文 (和合本修訂版)\\
\hline
6:1 & \begin{tabularx}{0.7\textwidth}{X} 耶和華吩咐摩西說： \end{tabularx} \\ \\ \relax
6:2 & \begin{tabularx}{0.7\textwidth}{X} 「你要吩咐以色列人，對他們說：無論男女，若許了特別的願，就是拿細耳人的願，願意離俗歸耶和華， \end{tabularx} \\ \\ \relax
6:3 & \begin{tabularx}{0.7\textwidth}{X} 他就要遠離清酒烈酒，也不可喝任何清酒烈酒做的醋；不可喝任何葡萄汁，也不可吃鮮葡萄和乾葡萄。 \end{tabularx} \\ \\ \relax
6:4 & \begin{tabularx}{0.7\textwidth}{X} 在一切離俗的日子，任何葡萄樹上所結的，甚至果核和果皮，都不可吃。 \end{tabularx} \\ \\ \relax
6:5 & \begin{tabularx}{0.7\textwidth}{X} 「在他一切許願離俗的日子，不可用剃刀剃頭。在離俗歸耶和華的日子未滿之前，他要成為聖，要任由頭上的髮綹生長。 \end{tabularx} \\ \\ \relax
6:6 & \begin{tabularx}{0.7\textwidth}{X} 在他一切離俗歸耶和華的日子，不可挨近死屍。 \end{tabularx} \\ \\ \relax
6:7 & \begin{tabularx}{0.7\textwidth}{X} 即使他的父母或兄弟姊妹死了，他也不可因他們使自己不潔淨，因為他頭上有離俗歸神的記號 。 \end{tabularx} \\ \\ \relax
6:8 & \begin{tabularx}{0.7\textwidth}{X} 在他一切離俗的日子，他是歸耶和華為聖的。 \end{tabularx} \\ \\ \relax
6:9 & \begin{tabularx}{0.7\textwidth}{X} 「若在他旁邊忽然有人死了，因而玷污了他離俗的頭，他要在第七日，得潔淨的日子剃頭。 \end{tabularx} \\ \\ \relax
6:10 & \begin{tabularx}{0.7\textwidth}{X} 第八日，他要把兩隻斑鳩或兩隻雛鴿帶到會幕門口，交給祭司。 \end{tabularx} \\ \\ \relax
6:11 & \begin{tabularx}{0.7\textwidth}{X} 祭司要獻一隻作贖罪祭，一隻作燔祭，為他贖因屍體而有的罪，並要在當日使他的頭分別為聖。 \end{tabularx} \\ \\ \relax
6:12 & \begin{tabularx}{0.7\textwidth}{X} 他要另選離俗歸耶和華的日子，牽一隻一歲的小公羊來作贖愆祭。先前的那段日子算為無效，因為他在離俗期間被玷污了。 \end{tabularx} \\ \\ \relax
6:13 & \begin{tabularx}{0.7\textwidth}{X} 「拿細耳人的條例是這樣的：離俗的日子滿了，祭司要領他到會幕門口， \end{tabularx} \\ \\ \relax
6:14 & \begin{tabularx}{0.7\textwidth}{X} 他要將供物獻給耶和華，就是一隻沒有殘疾的一歲小公羊作燔祭，一隻沒有殘疾的一歲小母羊作贖罪祭，和一隻沒有殘疾的公綿羊作平安祭， \end{tabularx} \\ \\ \relax
6:15 & \begin{tabularx}{0.7\textwidth}{X} 一籃用油調和的無酵細麵餅和抹了油的無酵薄餅，以及同獻的素祭和澆酒祭。 \end{tabularx} \\ \\ \relax
6:16 & \begin{tabularx}{0.7\textwidth}{X} 祭司要來到耶和華面前，獻上那人的贖罪祭和燔祭。 \end{tabularx} \\ \\ \relax
6:17 & \begin{tabularx}{0.7\textwidth}{X} 祭司要把公綿羊和那籃無酵餅獻給耶和華作平安祭，又要獻上同獻的素祭和澆酒祭。 \end{tabularx} \\ \\ \relax
6:18 & \begin{tabularx}{0.7\textwidth}{X} 拿細耳人要在會幕門口剃離俗的頭，把離俗頭上的髮放在平安祭下的火上。 \end{tabularx} \\ \\ \relax
6:19 & \begin{tabularx}{0.7\textwidth}{X} 他剃了離俗的頭以後，祭司要取那煮好的公綿羊的一條前腿，連同籃子裡的一塊無酵餅和一塊無酵薄餅，放在他手掌上。 \end{tabularx} \\ \\ \relax
6:20 & \begin{tabularx}{0.7\textwidth}{X} 祭司要拿這些在耶和華面前搖一搖，作為搖祭；這和所搖的胸、所舉的腿一樣是聖物，是歸給祭司的。然後拿細耳人才可以喝酒。 \end{tabularx} \\ \\ \relax
6:21 & \begin{tabularx}{0.7\textwidth}{X} 「這是拿細耳人許願的條例，除了他手頭財力所及之外，他要為離俗獻供物給耶和華。他怎樣許願，就當照離俗的條例做。」 \end{tabularx} \\ \\ \relax
6:22 & \begin{tabularx}{0.7\textwidth}{X} 耶和華吩咐摩西說： \end{tabularx} \\ \\ \relax
6:23 & \begin{tabularx}{0.7\textwidth}{X} 「你要吩咐亞倫和他兒子說：你們要這樣為以色列人祝福，對他們說： \end{tabularx} \\ \\ \relax
6:24 & \begin{tabularx}{0.7\textwidth}{X} 『願耶和華賜福給你，保護你。 \end{tabularx} \\ \\ \relax
6:25 & \begin{tabularx}{0.7\textwidth}{X} 願耶和華使他的臉光照你，賜恩給你。 \end{tabularx} \\ \\ \relax
6:26 & \begin{tabularx}{0.7\textwidth}{X} 願耶和華向你仰臉，賜你平安。』 \end{tabularx} \\ \\ \relax
6:27 & \begin{tabularx}{0.7\textwidth}{X} 「他們要如此奉我的名為以色列人祝福；我也要賜福給他們。」 \end{tabularx} \\ \\
[1ex]
\hline
\hline
\end{longtable}
$^{1}$弟兄姊妹平安.
在過年日期間.
我就去了看埃及文明展.
在西九那裡.
請問有沒有弟兄姊妹有去過看.
或者有打算過去看的.
都有 有打算.
都有看過.
是好看的.
我為什麼想去看呢.
就是因為其實早前.
電視就經常賣那隻法老貓.
我不知道大家有沒有留意.
文化博物館的館長.
就是歷史博物館的館長.
說這隻貓賣到斷市了.
我也是香港人.
有少少這樣的心態.
斷市了 不知道是什麼來的.
放心我不是買來炒的.
我只是好奇它斷市了.
是一隻卡通的貓仔.
好可愛的.
因為我上網都看過那些風格.
我就下定決心.
如果我要去看這個展覽的時候.
我就要買齊一套.
大中小不同的尺寸.
就覺得好可愛.
這樣放在那裡.
不過當我看完那個埃及文明展.
我就決定不買了.
在文明展裡面.
咦 還未開.
在文明展裡面.
我就拍了幾張照片.
它就有一隻貓.
不單止有貓.
還有貓的木乃伊在當中.
整件事是不是很可怕.

$^{41}$有貓的木乃伊.
那個貓的木乃伊.
原來是阿法魯的守護神.
他的守護貓.
那隻貓是代表著一個埃及的女神.
是一個負責保護牠們的生育.
保護家庭的那隻貓.
然後我幅圖裡面.
就有一個法魯的照片.
那個法魯叫做圖特摩斯三世.
我自己都買了一本關於埃及文明的書.
大概在公元前.
一千五百多年前去做法魯.
正正那個年代.
似乎是摩西帶領以色列人.
出埃及差不多的年代.
當然準確的年代.
就無從稽考.
但大約是公元前.
一千五百多或一千六百多年.
那幅圖像裡面.
那個法魯.
其實他的手握著兩塊東西.
一塊是一個權杖.
是代表著.
如果後面的童工處理好.
就示意告訴我.
你開第一幅圖就可以了.
是有一個權杖.
我看的時候.
這個權杖應該是代表著.
一些軍事的能力.
就好像我們讀詩篇.
耶和華是我的牧者.
他的杖.
他的安.
乾.
我們讀聖經的時候.
就安慰我.
但實際那支乾.

$^{81}$那支杖.
其實是有管教的意思.
有一個權力的意思.
埃及那個王的杖.
就手握著一支權杖.
而另一隻手就握著一樣.
類似一個.
我們看古裝片.
鳳天成雲.
他是什麼.
皇帝趙若.
一個皇帝的密令.
又或者好像一把劍.
握著一些東西.
他就坐在那裡.
頭戴著一頂冠冕.
頭頂就有一隻鳥在那裡.
原來.
我翻查一些資料.
原來在古代的埃及.
鳥仔是守護.
又是一種守護神.
他戴著頂帽.
就有一隻鳥在那裡.
加持著這個法老王.
所以.
法老王是被譽為.
或者是他自己譽為.
不單止有王權.
他還有什麼權.
有個神權在那裡.
因為守護神.
是直接在他的頭頂那裡.
去加持他的.
除了之外.
我剛才還有.
我的powerpoint應該還有兩張相.
又有一個人像.
跪在那裡.
有些東西在挑冠他.

$^{121}$另外.
也有一隻很大隻的.
一隻神獸.
我其實拍了非常多的相.
我在看的時候.
我就選了.
剛才其中那幾幅.
想和大家分享.
我在剛才.
我剛才形容.
那些相片裡面.
我就若摸去推敲.
若摸去推敲.
為什麼以色列人.
當神帶領他們出埃及之後.
短短一年左右的時間.
他們就會拜金牛獨呢.
因為原來他們由始至終.
經常見到的.
就是有些神獸在街上.
這純粹是我個人估計.
完全沒有聖經根據.
是個人估計.
我在想.
為什麼他們在曠野那裡.
在那裡吃麻辣吃到悶.
在那裡投訴.
上帝會給一些.
大家記不記得.
給什麼給以色列人吃.
記不記得會飛的.
鵪鶉.
給一些鳥.
給以色列人吃.
純粹個人估計.
因為法老的頭頂的棺.
有隻鳥在那裡.
我不知道是不是.
日後我回到天家.
我會親自問神.

$^{161}$會不會是.
即使真的向法老摑一巴.
你的守護神.
對我來說.
是給一些子民.
日日當鵪鶉那樣吃.
我不知道.
有這樣的可能.
好了.
接著.
這幾幅.
還未可以.
然後.
我又因為對埃及.
都很有興趣.
看完埃及展.
於是我就看一下.
在古代埃及裡面.
我就上網找.
古埃及的社會裡面.
有些什麼圖片.
是和一個.
即是埃及法老.
那些是有關的呢.
於是我就找了幾幅圖.
第一幅圖.
我形容給大家聽.
看看大家估不估到.
有幾個光頭佬.
只是圍著.
下面少少布.
前面有個人躺在那裡.
他們好像.
有些事情正在做.
你估那幅圖.
那些光頭佬是誰.
祭司.
這麼厲害.
我先要讚.
Sunny傳道.

$^{201}$Sunny傳道很厲害.
不過錦麟.
他都跟得很貼.
馬上說是祭司.
是的.
是祭司.
那幾個光頭.
你估他處理.
躺在那裡的那個人.
你估他在做什麼.
是不是在做外科手術.
還是什麼呢.
怎樣.
木乃伊.
是的.
終於到第二幅圖.
有沒有上一幅圖.
看看.
這是第二幅.
是之前那個.
這麼害怕.
難道教會的伺服器壞了.
是的.
我只能夠這樣緊張.
好了.
不要理之前第一第二.
就聽我說.
第一幅圖.
一些光頭是祭司.
如果大家先睜大眼.
那些祭司.
整個人是畫到花朵.
不知道大家看不看得到.
第二幅圖.
大家看不看得到.
正在幫我找.
第二幅圖我又形容了.
現在我做了古代的說書人.
古代的說書人.
沒有圖片.

$^{241}$沒有powerpoint.
不過是有power.
又有point.
好了.
第二幅圖是這樣的.
第二幅圖.
有幾個人跪在那裡.
然後有幾個人形的物體.
站在那裡.
拿著一些槍.
丫叉的東西.
臉是一些野獸.
這幾幅是文物.
這幾幅圖.
自從我見到這隻貓.
是埃及法老的守護貓之後.
我就一隻卡通貓都沒有買了.
因為我不容許.
一些帶著宗教意味的東西.
非上帝的東西在我家裡出現.
下面那個是貓奶衣.
是真的.
是木奶衣.
不過是捲著一隻貓.
是有五隻的.
如果你家裡有五隻貓.
你就看看吧.
中間那個.
我懷疑出埃及的時候.
可能是這個法老.
他的頭冠那裡是有一隻雀.
兩隻手.
一隻拿著拳柄.
一隻拿著類似女王物靈.
或者是一把刀的東西.
然後就坐在那裡.
旁邊就有一個跪在那裡的像.
跪在那裡的像.
然後就有兩塊東西.
是人的身體.

$^{281}$不過是動物的頭像.
是一些神奇的偶像.
或者用中國人的說話.
那些牛鬼蛇神.
其實是在加持祭司.
祭司是要那些東西去加持.
另外有一隻像羊飛羊.
像牛飛牛的東西.
非常大座.
是實物.
在文物展看到的.
當我看到這些文物的時候.
我就不難估計.
為什麼以色列人會製造金牛毒.
見不到上帝的時候會製造一塊東西.
因為他們身處的環境.
他們在埃及遺留了幾百年.
雖然他們的人不是幾百歲.
但代代相傳.
他們已經生活在那個環境.
而且他們成長可以出埃及的時候.
都幾十歲人.
習慣了見那些東西.
所以當上帝帶領他們的時候.
見不到的時候.
很自然就製造那塊東西出來.
說回第二幅圖.
第二幅圖的樓上那幅.
有些人跪在那裡.
有幾個人形的物體站在那裡.
最後面見到有兩個人坐在那裡.
坐在那裡有衣服穿.
見到嗎.
還有頭髮.
那兩個人坐在那裡前面有隻什麼.
有隻貓在那裡.
所以我估計坐在那裡.
那兩個人就是什麼.
法老和他老婆.
因為那隻貓是守護生育.

$^{321}$為什麼我會特別想看這一幅圖呢.
原來剛才我提到.
當說到這個埃及的法老的時候.
他比喻自己擁有神權.
不單有王權.
還有神權.
所以他前面有牛鬼蛇神.
幫他做守護神.
有隻貓做他的生育的守護家庭的神.
跟著有人跪在那裡.
作為他的奴僕婢女.
上下都有東西.
上面有軍隊.
下面有海裡面的魚.
這個我看見原來當時以色列人.
他們在埃及圍爐的時候.
見到的就是有形有象的法老.
原來上帝是有樣子看的.
就好像法老一樣.
跟著下面那幅圖.
我不說了.
如果在網上你見到這幅圖.
你會估計是在說什麼的場景.
第三.
這一幅.
有個人拿著叉.
拿著叉.
想做什麼呢.
打人.
後面有些人在那裡.
兩個人在那裡.
抬東西.
伏力.
跟著有人又被人打.
打到趴在那裡.
有沒有什麼聖經場景你會記起.
關乎古以色列人在埃及的呢.
有沒有什麼場景.
做奴隸.
被人打.

$^{361}$跟著摩西就幫他打死.
打他那個人.
記不記得.
跟著他們就敷冤.
神就去拯救他們.
在這個場景底下.
以色列人做了幾百年的奴隸.
被人打.
各樣的東西.
然後他們敷冤.
上帝去拯救.
而民數記.
就在他們被拯救之後的第二年.
二月初一日.
來到去.
二月初一日.
就寫民數記.
一至大概十章.
都在這一段時間寫的.
就是離開埃及的二月初一日.
他們離開埃及的第二年二月初一日.
他們什麼時候離開埃及呢.
就是對上一年的正月十四.
其實他們只是離開埃及.
大約一年零半個月.
差不多的時間.
有一位猶太裔的牧師.
他曾經這樣說.
他說在民數記的年代.
特別是一至十章左右的年代.
上帝說要數點人數.
然後就要揀選祭司.
揀選利未人.
來服侍他.
這個年代的時候.
以色列人是離開了埃及.
但是埃及並未離開他們.
明不明白我這句話說什麼.
以色列人離開了埃及.
不過埃及未離開他們.

$^{401}$因著這樣的緣故.
我就會想到.
為什麼他們會做金牛獨.
為什麼不是做金魚.
為什麼不是做金蛇.
為什麼會做金牛獨.
因為他們見的就是這些.
有形有體.
為什麼神的十誡.
就是我們沒有形象.
你不可以做任何事情來膜拜我.
因為在埃及.
全地就是牛鬼蛇神.
全部就是這些東西的同在.
我們知道.
過去兩個月裡.
聽了《民數記》.
分別有五章的經文.
有第一二章說到數點人數.
數點人數.
不是上帝要數一下.
究竟我帶了多少人出埃及.
我有什麼功績.
上帝說數那些人用來做什麼.
打仗.
數那些人成為軍隊.
剛才說明.
他們只是剛剛離開了埃及一年.
在這一年裡.
之前又經歷十災.
又過紅海.
又曠野.
很多的事情發生底下.
失驚無神.
原來他們是軍隊.
我想對他們來說.
都是一頭霧水.
都要很大的信心.
才能確認自己的身份.
然後在第三第四章.

$^{441}$上帝就確立了亞倫作為大祭司.
有利未人.
不同的宗族.
有戈核族.
摩西有三個兒子.
不同的兒子就負責.
不同在聖殿的服侍.
有人就抬藥櫃.
有人就負責搬糧柱.
有人就要搬布幕.
各人負責不同的事情.
然後到第五章.
開始就說到聖潔的課題.
說到子民之間.
怎樣聖潔的相處.
怎樣有聖潔的關係.
而今天我們就會說到.
上帝不單止去揀選一些人.
要保持一份聖潔在他當中.
他還容讓人去選擇.
過一個高度聖潔的生活.
而這種人就是拿西爾人.
剛才說到.
他們離開了埃及.
不過埃及並未離開他們.
或者這樣說.
他們對埃及的記憶.
仍然很牢固.
而拿西爾人.
一些願意高度聖潔的人.
即使他們的記憶.
仍然有很多埃及的元素.
但是他們願意為上帝的緣故.
放下這些記憶.
重新 reset 他們的腦.
讓他們充滿對上帝的敬崇.
然後之後第七至十章.
都是說到一些聖潔的事.
而七至十章就說到一些器具.
拿西爾人是一個陰役.

$^{481}$代表著分別為聖歸向神.
而利未人是一個支派.
即是事奉神那些.
包括祭司是一個支派.
是由上帝去揀選.
而拿西爾人就是不論男女.
不論年齡.
不論支派.
你是其餘任何支派都可以.
亦都不論宗族.
你是哪一個家族出來都可以.
只要你願意選定一個日期.
高度的分別為聖歸向神.
我記得何牧師在上個星期.
特別提到分別為聖.
不是代表完全完美.
是代表分出來.
從世俗裡面分別出來.
而這些拿西爾人.
就是願意從世俗裡面分出來.
在飲食方面他們有一些條例.
這裡說到遠離清酒和濃酒.
換句話說就是不可以飲酒.
今天的說話就是遠離茅台.
遠離紅酒和遠離啤酒.
我猜是這樣.
即是濃酒和烈酒.
還有酒造的醋.
我不知道有沒有醋不是酒造.
不是經過發酵.
醋是一種生活日用品.
調味品.
需要品.
葡萄和所有葡萄相關的製品.
無論是新鮮的葡萄.
葡萄汁.
或者一些乾葡萄.
都不可以食.
相對比較拿西爾人和利未人和祭司.
其實他們有很大的分別.

$^{521}$如果按上帝的揀選.
祭司利未人.
由二十開外至五十歲左右.
我們就叫做榮休.
如果我是利未人.
五十歲就榮休也不錯.
我不是尚未五十歲.
我是尚不似利未人.
不是尚未利未人.
因為我這輩子都不會是利未人.
我只能夠是拿西爾人.
而祭司利未人.
只限男性.
而拿西爾人.
男女老幼都可以.
祭司利未人可以食葡萄和其製品.
可以飲酒.
不過只是他公職的時候不可以飲酒.
但拿西爾人.
當你決定高度的聖潔歸向上帝的時候.
你是那段時間不可以飲酒.
你可以選擇一生成為拿西爾人.
你也可以選擇一個年期.
可能是三年.
可能是五年.
總之在這個年期裡.
你就去遵守這些規條.
去培養一個信服上帝的心.
不是因為上帝不喜歡人們飲酒.
或者不喜歡任何葡萄的原因.
當然這也是我回天家再和上帝問問才知道.
而是最重要的是.
培養那些人在一個世俗的環境這麼久.
培養他們不是跟世俗那一套.
培養他們由飲食開始.
信服上帝.
除了飲食之外.
還有一些規條.
就是他們在許願指定的日期離俗歸神.
這個離俗就不是我們印象中的.

$^{561}$剃度的師姑和尚.
就不是遠離俗世.
在深山那裡自己孝敬遠離.
離俗的意思是.
雖然活在俗世.
不過你不跟俗世那一套的方法.
好像剛才說的包括飲食.
除此以外.
他還說到不可以用刀剃頭.
要任由頭髮長長長長.
他要聖潔直到他許願離世歸神的日子滿了.
不可以哀近死屍.
即使父母弟兄姊妹死了.
這麼親的人離世.
他都不可以哀近.
直到離世歸俗的日子滿.
還有為什麼不可以哀近.
因為他是歸耶和華為聖.
因為上帝是生命的源頭.
死亡是有損於上帝的聖潔.
為什麼我剛才形容那幾幅圖.
那些祭司是光頭的.
為什麼我這麼強調光頭.
是因為那世耳人不可以剃頭.
某程度上我相信跟文化背景有些關係.
若果作為歸向上帝的人都剃了頭髮.
你覺得是不是跟他們在埃及的祭司差不多.
可能是差不多.
為免重疊.
可能就是要讓頭髮長長的留下來.
不可以剪.
而我們看參宣.
他在《士師記》裡參宣這位那世耳人.
他的頭髮很長.
一輩子都不剪.
說到能力在他的頭髮上.
我相信能力在頭髮上.
即是法力.
我相信這個不是上帝特別給他的法力.
只不過是一個他象徵歸向上帝的時候.

$^{601}$所有的能力.
那世耳人跟祭司剛才提到的分別.
是他們比祭司更加嚴謹.
那些飲食的條例.
而這裡他們比異教祭司不同的地方是.
他們留頭髮.
不可以剃頭.
不可以跟異教的祭司混亂一樣.
還有另一件事.
大家都很厲害.
我在形容的時候.
幾個光頭的人.
有個人體在前面.
處理木乃伊.
即是處理死屍.
而高度歸向上帝的人.
要撇離近死屍這一回事.
甚至親到父母,家人,兄弟姊妹都不可以.
若果假如他自定的日子.
離世歸耶和華的日子.
失陰無神走著走著.
旁邊有人死了.
那怎麼辦呢?.
他就要有贖罪制和凡制.
然後再另定日期來歸向上帝.
即是他歸向神離族的聖潔的日子.
斬斷了.
因為被玷污了.
於是他要重新來過.
重新來過的時候.
又有一個制.
叫做贖遣制.
然後在他定下來的日子滿了之後.
又有些制.
我讀給大家聽.
凡制,贖罪制,平安制,掃祭,奠祭和瑤祭.
還要將你歸向神的那段日子.
留起的頭髮剪掉.
燒在平安制的上面.
大家是否聽起來覺得很嚴謹.

$^{641}$不知道大家會否覺得.
為何上帝有這麼多要求.
不知道大家有沒有這樣的感覺.
因為大家看這個圖表的時候.
你們就輕鬆了.
我看經文的時候就很不輕鬆.
究竟是甚麼.
這個是甚麼制.
那個是甚麼制.
何時要獻這個制.
何時要獻那個制.
看到我想跟林牧師說.
我不想,我不說.
其實為何要定下這麼多的制.
是因為祂想被玷污的人.
重新有機會回到祂的身邊.
所以有這麼多的制.
而今天我們不需要獻制.
不是因為我們不被玷污.
是因為我們完全聖潔.
只是因為耶穌基督已經成為挽回制.
成為贖罪制.
救了我們.
以致我們不用再搞這麼多東西.
而這些各樣的制物.
背後的意義.
是上帝給我們機會回頭.
跟祂重新親近.
除此以外.
這些制.
贖罪制.
我們不要看這個制字.
我們看背後的意義.
贖罪制的意義是為了潔淨和赦免人.
可以獻斑溝,可以獻白鴿.
凡制是代表完全的降伏.
是獻一隻沒有殘疾.
初生沒多久的羊羔.
接著贖遣制.
是要補償一些損失.

$^{681}$所以價值也挺高的.
要獻一隻公面羊.
接著平安制.
平時贖罪制,凡制,贖遣制是消盡的.
但平安制是一半一半.
消一半.
另一半是獻祭的人可以吃的.
是代表著神和人之間那種和諧,復和.
接著掃祭.
是代表著上帝風調雨順.
給我們很多的供應,很多的環境.
以致我們有不同的收穫.
可以來奉獻,是一個感恩.
而奠祭,就是將酒奠下來.
就不同我們中國人的奠祭.
中國人的奠祭是祭鬼的.
包括我們說奠酒.
在祖先或偶像的面前.
是祭一些靈界.
不過在基督信仰裡.
在古以色列人裡.
奠祭是代表完全傾倒出來.
所以為耶穌傾倒香膏的婦人.
其實她為主有一個奠祭倒出來.
接著瑤祭,大家可以想像.
瑤祭是選擇一隻動物.
最漂亮的鼻的部分.
不知道大家吃羊是否喜歡吃羊鼻.
吃雞是否喜歡吃雞鼻.
我是的,因為特別滑.
大家想像,舉著那塊肉搖來搖去.
就好像大家剛才唱詩歌.
多謝敬拜隊帶得很好.
我也很投入搖來搖去.
是代表甚麼?.
神人的歡慶月臘.
我們不是看上帝有這麼多要求.
有這個制,有那個制.
上帝其實是給我們機會.
玷污了可以回頭.

$^{721}$上帝是要提醒我們.
要完全的奉獻.
上帝提醒我們犯了罪.
可以只有赦免.
提醒我們要降伏在祂的面前.
願意傾倒我們自己.
因為祂就是這樣傾倒在我們身上.
以致如果我們明白的話.
我們願意這樣做的話.
大家就歡慶了.
一起跟神享受這個瑤祭.
這個就是拿西爾人.
他如果不小心去玷污了.
還有當他離族的日子滿了之後.
所獻的祭物.
所以拿西爾人和上帝的關係.
是非常之美好.
好的地方是.
人選擇在世俗裡面.
分別為聖,歸向上帝.
而美好的是.
上帝願意接納人分別為聖,歸向祂.
不同的祭.
是為了給我們機會.
分別為聖,守潔心清.
去到祂的面前.
我記得許道良牧師.
一個播道神學院的前前前院長.
他曾經說過一句話.
我也很認同.
他說神愛我們是我們的權利.
因為神就是愛.
我會演繹多一點.
這是我們與生俱來的福氣.
原裝出廠就有了.
無論你有多好,你有多壞.
你聽祂的話,你不聽祂的話.
上帝都愛你.
甚至我會覺得.
有時你頑皮一點.

$^{761}$上帝都愛你多一點.
就好像我們有兒女就知道.
他頑皮的,你心裡總是想念他.
不知道大家是不是這樣.
我自己就是這樣.
但是,得不得到上帝的喜悅.
祂能不能夠給你一個饒祭.
就是關乎我們願不願意努力付出.
如果我們剛才看回這幅圖.
一個自選分別為性的人.
他要付上的代價.
其實是比祭司利未人更加多.
如果你用金錢物質去計算.
他給的錢是更多.
又要羔羊,又要公山羊.
即是有人在他旁邊死了.
又不關他的事.
不過,總之是玷污了.
又要重新做一次.
上帝愛我們是我們與生俱來的福氣.
因為上帝就是愛.
但是,得不得到上帝的喜悅.
祂願不願納我們.
和我們有一個饒祭的關係.
是視乎我們願不願意努力去建立這段關係.
開始的時候我說到.
民數記說要徵召二十歲以外的人出去打仗.
要建立以色列人成為上帝的軍隊.
而那西爾人.
就是不分性別,年齡,甘心,情願,高度的聖潔.
成為這隊軍隊裡的精兵.
成為軍隊裡的精兵.
他們禁戒清酒,濃酒,酒造的醋.
一切葡萄的產品.
在生活的當中.
不貪求世間的享樂.
甘於平淡.
接著,他們留長頭髮.
不是跟世俗那一套.
他們所眼見的世俗歸向偶像的人是剃光頭的.

$^{801}$不是跟世俗那一套.
而是願意在群眾當中顯露他自己.
他是屬於神的.
有記號的.
接著在關係裡.
不是愛自己,愛父母,愛兒女,愛其他人多於愛主.
而是要愛上帝過於愛其他的一切.
我們願不願意做這些拿西爾人.
做基督的精兵呢?.
做基督的精兵不是一句口號.
不是唱一首歌,信徒精兵歌.
不是唱完你就做了精兵.
而是要在世俗裡分別為性.
無論在我們的享樂裡.
在我們人群當中.
在我們關係裡.
全部上帝是我們的中心.
經文提到神有一個祝福.
這個祝福是我們在崇拜或其他聚會.
都會有的一個祝福語.
就是「願耶和華賜福給你,保護你」.
「願耶和華使祂的念光照你,賜恩給你」.
「願耶和華向你養念,賜你平安」.
神就這樣跟祭司說.
「如果你這樣奉我的名,祝福以色列人」.
「我就會賜福給他們」.
「賜福給你」的意思是上帝很願意.
祂的子民興旺.
願意他們身心平安.
「保護」除了說遇到危難.
神會保護之外.
更有一個意味.
就是上帝很願意高度的關注.
很細心留意我們.
「養念」是一個慈愛的少年.
好像一個慈祥的長輩.
可能是我們的父母.
可能是我們的師長.
可能是一個我們很敬重的人.
用慈祥的念向著我們.

$^{841}$而這個祝福.
不只是給那細爾人.
不只是給那些自選.
高度聖潔歸向上帝的人.
是給所有成為祂子民的人.
這位上帝主動拯救,主動揀選.
聽見人的呼聲.
祂來到人的當中.
帶領以色列人.
讓他們得著自由.
接著樂意接納人自願去親近祂.
讓那細爾人可以來到祂的身邊.
只要你願意,上帝就接納.
不過即使你沒有這樣的感動.
上帝仍然賜福你.
只要你是祂的子民.
祂都願意賜福.
而這位上帝.
值不值得我們高度的聖潔來回報祂呢?.
主耶穌基督主動的來到世界.
從高天降下凡塵.
取了奴僕的形象.
如果我們說.
好像那細爾人那樣.
清酒,濃酒不能喝.
醋不能吃.
葡萄乾不能吃.
我不肯.
但你想想.
耶穌在曠野獸試探四十天.
祂的麵包都沒有吃過一個.
祂的枕頭都沒有.
祂的枕頭的地方也沒有.
當我們還作罪人的時候.
祂願意來到世界拯救我們.
這位上帝.
值不值得我們在世俗當中.
分離,高度的聖潔去歸向祂呢?.
無論我們怎樣想.
祂都願意拯救.

$^{881}$用祂愛子的命.
來保住我們的命.
無論我們成為一個怎樣的人.
祂都想賜福,保護我們.
祂都願意用慈愛的念.
看著我們,對我們笑.
那細爾人是願意分別為聖歸向上帝.
高度自律的精兵.
在我們的生活當中.
我們願不願意追求聖潔.
追求簡樸的生活.
不是一定要事事都享受去信心的呢?.
而我相信這一份的簡樸.
是為了鍛鍊馴服上帝的心智.
為了我們的人生.
不是追求享受.
不只是追求清酒和濃酒.
是追求一份清心.
耶穌以奴僕的形象來到.
祂拒絕了萬國榮華的試探.
我們可不可以以分別為聖.
拒絕一些只是吃,買,玩的試探呢?.
會不會願意效法基督.
傳基督教的福音,使人有自由呢?.
不像一些異教的祭司或異教的人一樣.
特別是農曆新年.
講到異教的人.
最流行在電視或報攤出書或網上.
教你如何趨吉避凶.
教你如何為自己賺多一點.
如何生活好一點.
我們可不可以效法基督一樣.
我們所傳的福音是基督教的福音.
不是為了可以發達,可以享受.
可以玩多一點,吃多一點,健康一點,開心多一點.
是為了可以修補與上帝的關係.
為了神與人之間的一份復和.
為了有永恆的盼望.
而我們又願不願意在一個資訊發達.
你只要打開手機.

$^{921}$甚麼訊息你都可以看到的年代裡.
你甚麼訊息都不看.
先讀經,先默想,先祈禱,聽道行道呢?.
神無條件的拯救.
我們願不願意以上帝為我們的上乘之選.
最好的選擇,唯一的選擇.
為了上帝的緣故.
我們心願自閒.
過高度聖潔的生活.
作耶穌基督的精兵.
以聖潔來尊榮祂呢?.
不知道你想做精兵.
還是想在特選上帝的軍隊裡.
整天側側膊去做逃兵.
我不知道你想做哪一樣.
我想你記住的是上帝無條件的拯救.
無條件的愛.
無時無刻想賜福保護你.
用笑臉對著你.
我們用甚麼來回報祂.
我們一起祈禱.
因為你拯救我們.
在我們還作罪人的時候.
你就來到世界.
親自成為贖駕.
救我們脫離罪的痕跡.
求你幫助我們.
在世俗的當中.
雖然我們活在人群裡.
世界高舉著享受,享樂.
高舉著貪心,高舉著自私.
但求你幫助我們.
一生與你同行.
好像詩歌一樣盡力去愛你.
當我們唱出這句歌詞的時候.
幫助我們.
口唱心和真心誠意.
你已經拯救我們脫離了罪.
求你幫助我們.
不留戀世界的物質,貪婪,榮華.

$^{961}$不享受罪中之樂.
只求一心貼近你的心意.
以聖潔的生命.
作你的精兵.
蒙你的月亮.
聽我們在你面前同心的禱告.
我們這樣祈求.
是奉靠主耶穌基督的聖名.
阿們.
\newpage



\section{民數記 6:1-21}
\label{sec:031793}
\textbf{分別為聖:把自己交給神}~~~~孫靜敏牧師
\newline
\newline
連結: \href{http://www.church.com.hk/acms/content.asp?site=cdc\&op=show\&type=product\&code=031793\&layout=sermon}{\texttt{031793}} ~~~~ 語音日期: 2026-03-08
\newline
\newline
\hyperref[sec:031792]{\small{< < < PREV SERMON < < <}}
~
\hyperref[sec:index_chronic]{\small{[返順時目]}}
~
\hyperref[sec:031794]{\small{> > > NEXT SERMON > > >}}
\newline
\newline
民數記 6:1-21
\newline
\begin{longtable}{cl}
\hline
\hline
章節 & 經文 (和合本修訂版)\\
\hline
6:1 & \begin{tabularx}{0.7\textwidth}{X} 耶和華吩咐摩西說： \end{tabularx} \\ \\ \relax
6:2 & \begin{tabularx}{0.7\textwidth}{X} 「你要吩咐以色列人，對他們說：無論男女，若許了特別的願，就是拿細耳人的願，願意離俗歸耶和華， \end{tabularx} \\ \\ \relax
6:3 & \begin{tabularx}{0.7\textwidth}{X} 他就要遠離清酒烈酒，也不可喝任何清酒烈酒做的醋；不可喝任何葡萄汁，也不可吃鮮葡萄和乾葡萄。 \end{tabularx} \\ \\ \relax
6:4 & \begin{tabularx}{0.7\textwidth}{X} 在一切離俗的日子，任何葡萄樹上所結的，甚至果核和果皮，都不可吃。 \end{tabularx} \\ \\ \relax
6:5 & \begin{tabularx}{0.7\textwidth}{X} 「在他一切許願離俗的日子，不可用剃刀剃頭。在離俗歸耶和華的日子未滿之前，他要成為聖，要任由頭上的髮綹生長。 \end{tabularx} \\ \\ \relax
6:6 & \begin{tabularx}{0.7\textwidth}{X} 在他一切離俗歸耶和華的日子，不可挨近死屍。 \end{tabularx} \\ \\ \relax
6:7 & \begin{tabularx}{0.7\textwidth}{X} 即使他的父母或兄弟姊妹死了，他也不可因他們使自己不潔淨，因為他頭上有離俗歸神的記號 。 \end{tabularx} \\ \\ \relax
6:8 & \begin{tabularx}{0.7\textwidth}{X} 在他一切離俗的日子，他是歸耶和華為聖的。 \end{tabularx} \\ \\ \relax
6:9 & \begin{tabularx}{0.7\textwidth}{X} 「若在他旁邊忽然有人死了，因而玷污了他離俗的頭，他要在第七日，得潔淨的日子剃頭。 \end{tabularx} \\ \\ \relax
6:10 & \begin{tabularx}{0.7\textwidth}{X} 第八日，他要把兩隻斑鳩或兩隻雛鴿帶到會幕門口，交給祭司。 \end{tabularx} \\ \\ \relax
6:11 & \begin{tabularx}{0.7\textwidth}{X} 祭司要獻一隻作贖罪祭，一隻作燔祭，為他贖因屍體而有的罪，並要在當日使他的頭分別為聖。 \end{tabularx} \\ \\ \relax
6:12 & \begin{tabularx}{0.7\textwidth}{X} 他要另選離俗歸耶和華的日子，牽一隻一歲的小公羊來作贖愆祭。先前的那段日子算為無效，因為他在離俗期間被玷污了。 \end{tabularx} \\ \\ \relax
6:13 & \begin{tabularx}{0.7\textwidth}{X} 「拿細耳人的條例是這樣的：離俗的日子滿了，祭司要領他到會幕門口， \end{tabularx} \\ \\ \relax
6:14 & \begin{tabularx}{0.7\textwidth}{X} 他要將供物獻給耶和華，就是一隻沒有殘疾的一歲小公羊作燔祭，一隻沒有殘疾的一歲小母羊作贖罪祭，和一隻沒有殘疾的公綿羊作平安祭， \end{tabularx} \\ \\ \relax
6:15 & \begin{tabularx}{0.7\textwidth}{X} 一籃用油調和的無酵細麵餅和抹了油的無酵薄餅，以及同獻的素祭和澆酒祭。 \end{tabularx} \\ \\ \relax
6:16 & \begin{tabularx}{0.7\textwidth}{X} 祭司要來到耶和華面前，獻上那人的贖罪祭和燔祭。 \end{tabularx} \\ \\ \relax
6:17 & \begin{tabularx}{0.7\textwidth}{X} 祭司要把公綿羊和那籃無酵餅獻給耶和華作平安祭，又要獻上同獻的素祭和澆酒祭。 \end{tabularx} \\ \\ \relax
6:18 & \begin{tabularx}{0.7\textwidth}{X} 拿細耳人要在會幕門口剃離俗的頭，把離俗頭上的髮放在平安祭下的火上。 \end{tabularx} \\ \\ \relax
6:19 & \begin{tabularx}{0.7\textwidth}{X} 他剃了離俗的頭以後，祭司要取那煮好的公綿羊的一條前腿，連同籃子裡的一塊無酵餅和一塊無酵薄餅，放在他手掌上。 \end{tabularx} \\ \\ \relax
6:20 & \begin{tabularx}{0.7\textwidth}{X} 祭司要拿這些在耶和華面前搖一搖，作為搖祭；這和所搖的胸、所舉的腿一樣是聖物，是歸給祭司的。然後拿細耳人才可以喝酒。 \end{tabularx} \\ \\ \relax
6:21 & \begin{tabularx}{0.7\textwidth}{X} 「這是拿細耳人許願的條例，除了他手頭財力所及之外，他要為離俗獻供物給耶和華。他怎樣許願，就當照離俗的條例做。」 \end{tabularx} \\ \\
[1ex]
\hline
\hline
\end{longtable}
$^{1}$今天去到《民數記》的第六章.
我給個題目叫做 分別為聖 將自己教給神.
《民數記》不單止是一個好像數點人數的一卷書.
同樣《民數記》都是記載了一班夢救贖的子民.
他們走到在曠野的路上.
外面有試探 裡面又不停抱怨.
又面對未知的前面.
其實想想都幾像我們自己.
我們在外面有很多試探.
裡面可能在家裡 在職場上.
我們心裡有很多抱怨.
同樣我們一樣 對前路是甚麼都不知道.
但是上帝就在曠野裡面帶領著這班以色列民.
不單止帶他們走到目的地.
同樣將他們在這個走曠野的路上.
讓他們成為屬神的子民.
大家每個星期來到教會.
同樣都是被神去造就.
希望成為一個更健全的基督徒.
今日我們在整個章節第六章的時候.
分別有三個我們要解決的問題.
第一個就是.
甚麼叫做將自己交給神.
分別為性和日常生活有甚麼關係呢.
當我們生命上有破口 有失誤的時候.
我們應該怎樣做呢.
我們怎樣才可以重回上帝的面前呢.
最後到底祝福是做甚麼的呢.
神要透過祭司宣告祝福.
這個對我們現今每一個坐在這裡的弟兄姊妹.
有甚麼意思呢.
我會將它分為三個大的段落.
就是第一個拿西爾人所許的願.
自願分別獻給神.
接著就是當這個誓起完之後.
原來會有些失誤之後要怎樣算呢.
最後是祭司的祝福.
三個重要的點.
就將我們整個第六章從今生奉獻到怎樣去失誤.
然後恢復最後進入神的祝福和差遣.

$^{41}$我們一起開始看看這個第六章.
拿西爾人的奉獻.
成熟的信徒自願去分別為性.
我們一起讀一下第一至第八節.
耶和華對摩西說.
李曉遇以色列人說.
無論男女許了特別的願.
就是拿西爾人的願.
要離族歸耶和華.
他就要遠離清酒弄酒.
也不可喝甚麼清酒弄酒做的醋.
也不可喝甚麼葡萄汁.
也不可吃鮮葡萄或公葡萄.
在一切離族的日子.
凡葡萄樹上結的.
自核自皮所做的物.
都不可吃.
在一切的日子.
不可用剃頭的剃頭.
要由髮留長長了.
他要性結.
直到離族歸耶和華的日子滿了.
在他離族歸耶和華的日子.
不可靠近死屍.
他的父母或弟兄姊妹死了的時候.
他不可因他們使自己不潔淨.
因為他離族歸神的憑據.
是在他頭上.
在他一切離族的日子.
是歸耶和華為聖.
這裡說到一個很特別的族類.
叫做甚麼.
拿西爾人.
神主動的啟示.
無論男女許了特別的緣.
作拿西爾人.
要離族歸耶和華.
大家記不記得.
有一個在聖經裡面.
大家應該很熟悉的人.

$^{81}$拿西爾人.
有沒有記得是誰.
參算.
到底這些拿西爾人是怎樣的呢.
原來拿西爾人.
他不是一個人種.
不是中國人.
他們是一個自願去奉獻.
不是因為律法要求他們這樣做.
而是神將一個更深奉獻的路.
擺在他們面前.
讓他們願意去回應上帝這個願.
對於成熟的基督徒.
給我們一個很大的提醒.
我們熟悉崇拜.
事奉奉獻.
但熟悉之後.
我們經常都會有一個.
我們外在維持是這樣.
但內裡我們已經習慣了.
我們很多時候都會說.
渴望有什麼問題出現.
我們會不會還是很想去親近神.
還是覺得其實我們現在坐在這裡都OK.
都幾好這個狀態.
又或者我們對於毀身的問題.
我們是不是仍然去相信.
神值得我們為祂做出具體的選擇.
而不是停留那些.
我好感動.
這句口號幾好.
拿西爾人亦都不是一個叫我們去做超級基督徒的制度.
反而是給我們一班成熟的基督徒.
作一面鏡子.
照下我們在這個時間.
我們願不願意去盡心專注的.
去在我們的基督身上.
他告訴我們.
作拿西爾人有三個記號.
三個層面的交托.

$^{121}$第一個就是說.
不可吃清酒.
不可喝清酒濃酒.
放下那些可以依靠的安慰.
不剃頭髮.
是要讓人見到我們是屬於神的.
不要捱近死屍.
保守敬拜的時候.
那一種的純潔.
和什麼才是優先次序.
這三個就是告訴我們.
作拿西爾人的時候.
我們將自己交給神.
我們依靠什麼.
我們是屬於誰的.
以誰為我們的生命優先次序.
我們說第一個記號就是.
不要喝清酒濃酒.
放下可安慰的.
清酒濃酒.
我猜就是由茅台到啤酒.
現在這樣說.
由茅台到啤酒都不可以喝.
紅酒白酒當然是不可以喝.
剛才經文裡面都說到.
不可以吃葡萄製的東西.
有一些解經家就說.
為什麼不可以吃葡萄呢.
因為在當時的背景裡.
葡萄所製的東西和吃葡萄.
是代表一個享樂的意思.
大家有沒有看電影或電視.
那些皇帝躺在那裡.
通常妃子都是給什麼去吃.
葡萄.
我覺得好像是挺對的.
所以就叫我們.
不要吃葡萄所做的一切.
叫我們放下這個享樂.
我們說在這個替代性的安慰.

$^{161}$什麼叫替代性的安慰呢.
其實在這個時間.
我們說走.
未必是最大的試探.
但是很多時候.
很多東西給我們的替代安慰.
就是我們現在回家可能很累.
我們要滾一下手機.
現在還有很快的短劇.
一分鐘一集.
一開始就是爬一巴.
我也有看過.
爬一巴然後很快又一集.
滾一下手機.
或者我要很辛苦.
我要購物.
或者我要不開心.
我要大吃一餐.
這些都是一些安慰.
但是上帝告訴我們.
叫我們問一下自己.
當我們遇到很大壓力的時候.
我們第一個要找的是誰.
當我們受傷的時候.
第一個用來麻醉自己的方法.
是什麼呢.
當我們疲倦.
很累.
很辛苦的時候.
我們第一個想依靠的.
又是誰呢.
是不是上帝.
還是一些要讓你感覺刺激的東西呢.
我們說.
在新約裡面.
耶穌跟我們說.
「犯渴者水的還要再渴」.
「人若渴我所賜的水就永遠不渴」.
這裡告訴我們.
我們替代的安慰.

$^{201}$是要回到能夠給活水我們的耶穌.
拿西人的禁戒.
不是告訴我們.
要我們變得枯乾.
什麼不要做.
這個不要做.
那個不要做.
而是引導我們.
回到能夠給活水我們的上帝.
這些禁止.
我們不要用次等的安慰方法.
去填補我們的空虛.
而是叫我們.
回到耶穌基督的裡面.
得享真真正正的安息.
到第二方面.
第二就是叫我們.
不要剃頭.
讓人看見我們是屬主的.
留長頭髮.
當時是一個聖潔的記號.
拿西人的奉獻.
不是就這樣躲起來放在心裡.
而是要在生活上.
留下給我們.
可以辨認的痕跡.
今天我們不會說要留長頭髮.
才是代表我們的敬虔.
而是要在我們的生活上.
怎樣讓我們看到上帝.
每一天我們都會告訴自己.
又或者讓人知道.
我們很忙的.
其實弟兄姊妹.
經常都見到牧師.
第一句就說.
我知道你很忙的.
不用你說.
我們都是很忙的.
但是我們很想讓你看到.

$^{241}$我們忙的背後.
是為誰忙.
我們忙的背後.
是為上帝而忙.
是為弟兄姊妹而忙.
不是為無事忙.
求主繼續讓我們看到.
我們身邊每一個人.
我們的忙是為了什麼而忙.
今天我們在生活上.
我們能不能夠真實地.
去分辨到在日常生活裡面.
我們能夠看得見.
我們和其他人是不一樣的.
我們的金錢使用.
會不會是屬於神的眼光呢.
在我們的人際關係裡面.
我們的反應.
像不像你是一個屬主的人呢.
當我們說話的方式的時候.
你的語氣像不像是屬於上帝的呢.
我們的言行是不是合一的呢.
當我們在隱密處的時候.
沒有人見到我們的時候.
我們的選擇是不是就這樣.
不讓人知道.
我們在家裡和在教會裡.
我們都是同一個我呢.
這些都是我們在生活上.
怎樣去告訴別人.
我是一個基督徒.
其實這個指向.
是要指向基督.
並不是去炫耀自己.
剛才所說.
我們要留長頭髮.
是要讓別人知道.
我們是屬神的.
但是今天我們不再用這個方法.
留長頭髮.

$^{281}$讓我們知道屬神.
你有什麼方法去告訴別人.
你是一個基督徒呢.
可能有時候.
你會戴一個十字架.
告訴別人你是基督徒.
又或者.
你吃飯的時候.
你會不會祈禱呢.
還是你覺得.
在同事面前吃飯祈禱.
有點尷尬.
甚至有些職青告訴我.
他們尷尬到一個地方.
就是就這樣閉一閉眼.
阿們.
這樣呢.
我相信.
這些都不是我們要做的.
我們要做的.
就是顯示出.
我們能夠在日常生活.
都能夠告訴別人.
我們是一個基督徒.
這裡說.
不要讓你的過去的資歷.
取代你今天的模樣.
不要讓昨天的熱心.
替代你今天的馴服.
意思是說.
我們不要讓我們的過去.
限制我們的今天.
我們要馴服在上帝面前.
可能今天.
上帝就叫你.
好像拿西爾人一樣.
去做好這個約.
最後呢.
又去有第三樣東西.
叫我們不要靠近這個死屍.

$^{321}$保守敬拜和純潔.
這個意思是什麼意思呢.
神是最優先的主.
也是生命源頭的神.
拿西爾人不是不愛他們的家人.
而是他要表明.
在這段時間.
特定的時間裡面.
他要把屬神的呼召.
放在最優先.
神教導以色列人.
他們不單止是和萬民的子民.
不一樣的.
他是萬民的子民.
他的不一樣.
不是被強烈的情感牽引.
他們是要承認.
神是他唯一的主.
有很多時候.
我們都是因為忙.
因為有很多的事情.
而將這些事高過於上帝.
但是我們也很感恩.
我們現在生於新約的時代.
在這個福音的年代.
我們可以怎樣呢.
我們今天不需要再靠.
不碰這個死屍來潔淨.
福音甚至告訴我們.
耶穌可以用手去觸摸.
這些麻瘋病人.
叫死人復活.
為他們帶來潔淨.
和生命帶給有需要的人.
一個成熟的基督徒的記號.
是怎樣的呢.
我們是仍能被主呼召.
我們在熟悉的環境裡.
我們仍然能夠聽到.
每一個人的說話.

$^{361}$我們可以放下.
放下一些我們認為.
可以依靠享樂.
和可以依靠的事情.
我們仍能夠將神放在第一.
不是被家庭.
我們的面子.
人際壓力去代替神的位置.
其實拿西爾人的怨.
不單只是民數記那一個時間.
在現在這個年代.
依然會有人.
願意好像拿西爾人一樣.
作這一個怨.
認不認識這位林書豪.
認識是嗎.
當時早前他是在NBA裡面.
一位很出色的華人球員.
今日他回到台灣.
他在台灣裡面.
繼續教籃球打籃球.
同樣他都是一個自由的傳道人.
他是對年青人一輩說.
他願意接受拿西爾人之約.
選擇用減法的生活.
去無全清.
以致讓人知道.
他不單只是屬於神.
他還要將神放在他生命裡面.
讓神成為他的唯一.
我看到這個年青人.
能夠願意.
甘心地去作一個這麼大的怨.
為的就是上帝.
我不知道大家會不會.
就如剛才所說.
我們進入大齋期的時候.
我們會不會願意.
放下一些我們的喜好呢.
會不會放下我們很喜歡飲的.

$^{401}$可樂,咖啡.
一段時間.
讓我們自己去專注於上帝上面呢.
求主繼續幫助我們.
當我們一個成熟的基督徒.
甘心樂意的時候.
是不是就生命裡面.
不會有失誤呢.
或者當我們有失誤的時候.
我們應該怎樣做呢.
在民數記接下來那套這樣說.
「若在他旁邊忽然有人死了.
以致點染了他離族的頭.
他要在第七日得潔淨的時候剃頭.
第八日他要把兩隻斑溝和兩隻初甲.
帶到會幕口交給祭司.
祭司要獻一隻作贖罪祭.
一隻作凡祭.
為他贖那因死屍而有的罪.
並要當日使他的頭成為聖潔.
他要另選離族歸耶和華的日子.
又要獻一隻一歲的公羊羔.
來作贖遣祭.
但先前的日子要歸徒然.
因為他在離族之間被玷污了」.
嘩你看看這裡.
他說有時有很多時候.
我們被玷污不是我們自己想的.
好像他這樣說.
突然之間有一個人死了.
跌在你旁邊.
你都不知道的.
這樣被玷污.
這裡是想說甚麼呢.
在我們生命裡.
我們都不會自己想玷污自己.
但是在生活裡.
我們會遇到很多壓力.
很多你不知道的原因.
可能讓你刻意離開了上帝.

$^{441}$又或者你長期覺得生命很疲倦.
找不到依靠.
你會放縱自己.
甚至有時在生命裡.
你遇到艱難.
你會說為甚麼沒有神在我生命裡.
為甚麼你不幫我.
有很多時候.
這些不知道為甚麼.
就成為了我們一班成熟的基督徒.
甚至可能你剛剛信主的時候.
我們會覺得.
這些玷污不是我們自己想的.
成熟不代表我們不受玷污.
成熟的基督徒.
甚至基督徒.
我們是要知道.
怎樣去悔改.
這裡告訴我們.
我們有很多的制可以獻.
這些制完全是神告訴我們.
我願意給你機會.
你好好的去改過.
當成熟的基督徒上見.
有很多不同的破口.
剛才說我們事奉的時候.
都會有疲倦的時候.
我們會從熱心.
到你說.
今天沒有人來.
我自己就頂上.
又或者我不做就沒有人做.
將事奉的喜樂.
慢慢慢慢降低.
我們甚至會不會對家庭.
已經開始冷淡呢.
我們維持我們自己.
在家庭上的責任.
但是我們已經沒有甚麼.
可以跟家人說.

$^{481}$說的時候都是說.
孩子的學業.
你的工作.
信仰已經不再說了.
又或者一些隱然未見的罪.
外表很正常.
但是當我們自己一個人的時候.
我們能不能夠將我們的情緒.
放在上帝那裡呢.
往往帶著情緒的.
是從來都解決不了問題.
關係裡面我們有沒有苦毒呢.
在我們的心裡面有沒有一些人.
我們認定了他就是這樣的呢.
我們照常來到聚會.
但是我們從來都不會.
為這些人去祝福,祈禱.
在民數記告訴我們.
有一條迴轉的路.
神告訴我們.
我們要承認我們被玷污了.
然後我們能夠去到神面前.
潔淨更新.
然後迴轉再歸向上帝.
在舊約裡面.
當我們被玷污的時候.
我們用不同的劑來獻情.
但是來到這裡.
新約的時候.
我們已經不需要.
為什麼呢.
因為我們有主耶穌.
為我們做了一次性的大祭.
祂將祂獻過給我們.
祂用寶血將你和我的罪.
一次過清洗了.
今天你坐在這裡的時候.
我們怎麼樣去回應祂呢.
曾經有一個牧者提醒過我們.
我們在生命裡面.

$^{521}$你和我在這一刻.
我們都只不過是一個逗號.
或者英文就是comma.
在前面過去的已經過去了.
可能我們有錯誤.
有做得不好.
甚至是得罪神的地方.
但是在逗號之後.
句號之前.
我們依然可以迴轉.
你做得不好.
你做錯了.
今晚我鼓勵你們.
去到上帝面前.
祈禱認罪悔改.
我們一天還沒有走到這個句號的時候.
我們依然能夠改變.
這個就是上帝在這裡告訴我們.
一個很重要的訊息.
我們是有機會的.
上帝是會接納我們的.
還記不記得.
和耶穌釘十字架.
旁邊那兩個人.
其中一個怎麼說.
當他那一刻願意相信的時候.
耶穌跟他說.
你會和我一起在樂園相見.
所以你和我還沒有去到這個句號的那一刻.
我們依然都可以改變.
所以當我們要重新開始的時候.
我們有什麼可以在生活上重新去做.
我們重建我們禱告的節奏.
我們不是有空的時候.
子女要考試的時候才祈禱.
我們是每一個時間都和上帝有親密的對話.
我們去找我們屬靈的土伴.
不是隨便去交一個朋友.
而是讓我們在靈命上可以建立.
可以做就的朋友.

$^{561}$我們去觀賞我們令你跌倒的破口.
可能是一個Apps.
可能是一段亂七八糟的關係.
甚至是一些晚上不好的壞習慣.
今天我們願不願意去放下一個拖延.
而去馴服.
向得罪你的人和好.
歸還 承認和道歉.
這些都是上帝告訴過我們.
我們可以被做就.
最後就是由祭司的祝福.
將神的名放在你和我的身上.
耶和華曉如摩西說 你告訴阿倫和他兒子說.
你們要這樣為以色列人祝福說.
我們一起讀.
願耶和華賜福給你 保護你.
願耶和華使他的臉光照你 賜恩給你.
願耶和華向你養念 賜你平安.
他們要如此奉我的名 為以色列人祝福.
我也要賜福給他們.
這三句很熟悉.
Sorry 原來剛才沒有這麼精文.
難怪讀不到 對不起 太緊張.
這三句是很熟悉的 三句祝福的說話.
在當中 願耶和華賜福給你 保護你.
願耶和華使他的臉光照你 賜恩給你.
願耶和華向你養念 賜你平安.
這是一個祝福.
這個祝福是很重要的.
這個祝福是神將他的名.
放在百姓身上.
一個盟約的宣告.
弟兄姊妹 你願不願意.
讓這個祝福臨到你的身上.
我們說這個祝福是從外而內.
去到我們的生命裡面.
讓我們去守住這個祝福.
當他說使他的臉光照你 賜恩給你.
祝福就是這個中心 是神自己.
不是我們做得多好 多有成就.

$^{601}$祂才來祝福你.
而是出於上帝的恩典.
當我們看到 向你養念 賜你平安.
神要與你和好.
祂要使你一個完整的我.
這裡告訴我們.
當你想想這個圖畫的時候.
祂看著你 告訴你.
我賜你平安.
這個圖畫是多麼美麗.
祂不是遠遠 很遠距離告訴你.
祂看著你 你做得好.
我賜給你平安.
不是 你看看這段經文.
祂是看著你的.
然後對你說 我賜給你平安.
最後 我們說有三個要領受祝福的態度.
首先第一個 這裡沒有打出來.
首先第一個要領受祝福的態度是什麼呢.
想不想得到.
謙卑.
謙卑 怎樣才謙卑呢.
這個很重要.
就是不要早退.
你沒有拿祝福就走了.
不用拿祝福了.
所以除了不要走 不要早退.
拿祝福之外.
我們要用信心去領受這個祝福.
接著我們帶著我們悔改的心去領受.
願意將我們的生命去交賞.
最後我們帶著差遣的意識.
去迎接新的一個星期.
我們今天在民數記第六章那裡.
告訴我們.
我們願意去分別為勝.
好像拿西爾人那樣.
讓我們的心歸向上帝.
我們願意將我們失誤和破口.
去帶到神面前.

$^{641}$同樣我們要相信這個祝福.
是神對我們說話.
將神的名 神的光 和神的平安.
讓我們帶入家庭 帶入職場裡面.
上帝告訴我們.
當我們要對準上帝的時候.
我們就能夠聽到他對我們的說話.
最後我用一個近期.
我很喜歡和弟兄姊妹分享的一個見證.
但我相信這裡很多人都沒有聽過.
所以我會和大家說.
我有一天做牧師的時候.
都會有他軟弱的時候.
都會有很多問為什麼的時候.
我都有問為什麼.
我記得當我決定要去KFCC的時候.
去到現在如今九個月.
我有時候都會問.
上帝啊 其實你要我做些什麼呢.
有一天我靈修的時候.
上帝就給了我一句經文.
他就說耶和華對亞伯蘭說.
你要離開本地 本族 父家.
往我所指示你的地方去.
當我就想.
我剛剛才去完KFCC.
難道上帝又要我去宣教.
還有我都不是一個宣教的人.
我就說上帝.
其實你想要我去哪裡呢.
當我祈園禱張大眼的時候.
在我的書桌裡面.
就有兩隻熊貓和耶穌站在中間.
我就說.
上帝你是否要我去成都呢.
這樣多好啊.
上帝去成都看熊貓.
我就說上帝.
如果你真的要我去成都.
你一定要告訴我.

$^{681}$我要這次很清晰的指示.
否則我不會離開KFCC.
我一直祈禱.
一直祈禱.
當我上班的時候.
我平時都知道.
KFCC那條街叫南寧街.
但是周圍附近的街.
我就不是很清楚.
那天就照樣上班.
下車就走了.
一走的時候.
嘩 見到.
是不是很害怕.
你是我 你害不害怕.
成都街.
你看不看到上面還有兩隻熊貓.
上帝答了我.
是不是.
原來上帝說.
你不要再問來問去.
你是要去那裡.
你是要去KFCC.
成都街都給你.
兩隻熊貓都給你.
這麼清楚.
你都還問.
當我看完的時候.
我心裡面.
第一個感覺就是.
我終於對齊上帝的心意.
我終於聽到上帝.
對我清清楚楚的說話.
丁姊妹.
希望你們繼續為孫牧師祈禱.
孫牧師今次.
不會再硬硬的.
是不是要走.
是不是猜輸了我要去.
不是的.

$^{721}$我今次是很開心的.
跟大家宣告.
成都街.
這兩隻熊貓.
KFCC.
是上帝要我所去的.
亦都是上帝.
所派我們廣福堂.
要去到這個地方.
傳講上帝的福音.
因此.
上帝同樣.
告訴我們一個重要的訊息.
我們會在4月5號.
復活節的當日.
在KFCC開第一次的崇拜.
求主繼續去帶領著我們.
讓我們這班基督徒.
不單止坐在金鐘.
還願意像拿西爾人一樣.
專注的為上帝而活.
在我們的生命裡面.
只有上帝.
亦都同樣.
作為KFCC的總幹事.
都呼籲弟兄姊妹.
走出這4幅牆.
去到社區.
見證上帝.
這個這麼大的福音.
求主繼續幫助我們.
我們一起有個禱告.
差天父上帝.
用你的話語.
去扶持我們.
幫助我們.
讓我們在這個世代裡面.
不會被外界所吸引.
而是用你的話語去牽引我們.
亦都讓我們.

$^{761}$在我們的軟弱裡面.
願意回轉.
求主你繼續去根深我們.
讓我們在失誤的當中.
有勇氣去悔改去承認.
因為我們知道.
我們未去到句號的這一刻.
求主你照你所應許的賜福給我們.
保護我們.
使你的臉光照我們.
賜恩給我們.
向我們養念.
賜我們的平安.
願你的名.
繼續放在我們的身上.
不要離開我們.
繼續差遣我們.
進入我們的家庭.
我們的職場上.
讓我們的生活.
活出像你一樣.
求主你繼續幫助我們.
帶領著我們.
我們同心在你面前仰望.
是奉靠主耶穌得勝的明智祈求.
阿們.
\newpage



\section{以弗所書 5:25-32}
\label{sec:031794}
\textbf{榮耀的基督和祂最愛的教會 (華/粵)}~~~~馬天嵐牧師
\newline
\newline
連結: \href{http://www.church.com.hk/acms/content.asp?site=cdc\&op=show\&type=product\&code=031794\&layout=sermon}{\texttt{031794}} ~~~~ 語音日期: 2026-03-08
\newline
\newline
\hyperref[sec:031793]{\small{< < < PREV SERMON < < <}}
~
\hyperref[sec:index_chronic]{\small{[返順時目]}}
~
\hyperref[sec:031795]{\small{> > > NEXT SERMON > > >}}
\newline
\newline
以弗所書 5:25-32
\newline
\begin{longtable}{cl}
\hline
\hline
章節 & 經文 (和合本修訂版)\\
\hline
5:25 & \begin{tabularx}{0.7\textwidth}{X} 作丈夫的，你們要愛自己的妻子，正如基督愛教會，為教會捨己， \end{tabularx} \\ \\ \relax
5:26 & \begin{tabularx}{0.7\textwidth}{X} 以水藉著道把教會洗淨，使她成為聖潔， \end{tabularx} \\ \\ \relax
5:27 & \begin{tabularx}{0.7\textwidth}{X} 好獻給自己，作榮耀的教會，毫無玷污、皺紋等類的缺陷，而是聖潔沒有瑕疵的。 \end{tabularx} \\ \\ \relax
5:28 & \begin{tabularx}{0.7\textwidth}{X} 丈夫也應當照樣愛妻子，如同愛自己的身體；愛妻子就是愛自己了。 \end{tabularx} \\ \\ \relax
5:29 & \begin{tabularx}{0.7\textwidth}{X} 從來沒有人恨惡自己的身體，總是保養愛惜，正像基督待教會一樣， \end{tabularx} \\ \\ \relax
5:30 & \begin{tabularx}{0.7\textwidth}{X} 因我們是他身體的肢體。 \end{tabularx} \\ \\ \relax
5:31 & \begin{tabularx}{0.7\textwidth}{X} 「為這個緣故，人要離開父母，與妻子結合，二人成為一體。」 \end{tabularx} \\ \\ \relax
5:32 & \begin{tabularx}{0.7\textwidth}{X} 這是極大的奧祕，而我是指基督和教會說的。 \end{tabularx} \\ \\
[1ex]
\hline
\hline
\end{longtable}
$^{1}$在十字架上所成就的事情,那個異形,讓我們能夠來到你面前,一起來敬拜你,.
我們今天來到你面前,我們知道你在你的寶座上,你聽我們的敬拜,悅納我們的敬拜,.
讓我們在地上也能夠經歷到在天上的敬拜,我們在這個時候,我們為粵語部的士工禱告,.
我們為聘木的事情禱告,也為粵語部的領袖禱告,.
讓你自己的教會能夠在這些領導者的引導之下,能夠討你的喜悅,.
我們今天來到你面前,我們定義要聽你自己的話,.
剛才我們已經唸過了今天所要看的經文,今天我們要接著五章二十五節到三十二節,.
我們要來看基督的榮耀和祂的教會,這幾週來教會的牧師和長老們看見牧師長老的需要,.
所以長老們就在各個語言堂分享長老崇高的呼召,長老們在各個不同的組裡面講到長老的呼召,.
是一個很崇高的呼召,而且是一個榮耀的職份,而且是一個相當重的職份,.
今天我們就要用這兩處經文,我們來看,我要從另外一個角度來看,神為什麼要在地上的教會設立長老,.
讓教會跟隨祂的旨意走,教會是隨著神的意思走,我們下面要想一個問題,教會是什麼?.
如果我們給你們機會講的話,如果我們給你們時間能夠分享,我想我們一定會有很多不同的意見,.
我們當中可能有很多不同的意見,今天我們就要從這兩段經文,我們當然整本聖經都是一樣,我們可以看到教會到底是什麼,.
我們可以從這兩段經文來想一想教會是什麼,以弗所書第五章二十五節到三十二節,再加上二十二節到二十四節,對我們來說是很熟悉的,以弗所書第五章,.
以弗所書第五章的經文和第三章的經文我們都很熟悉,因為是跟婚姻有關的,是跟婚姻有關,所以我們一講到婚姻我們大家都很注意,.
但是因為我們太看重了家庭,婚姻,夫妻,我們卻忽略了神的心意,.
但是我們卻忽略了神的心意,我再把三十一節到三十二節讀給各位聽,我重新再從三十一三十二節讀給大家聽,.
他說為這個緣故人要離開父母,這是聖靈講的話,為這個緣故人要離開父母,與妻子聯合,二人成為一體,這是極大的奧秘,.
但我是指著基督和教會說的,弟兄姊妹,這句話把我們帶回到神的心意上面,.
這個經文提到為這個緣故人要離開父母,與妻子聯合,二人成為一體,這是極大的奧秘,但我是指著基督和教會說的,.
因為我們要注意人要離開父母,與妻子聯合,不是說女人要離開父母,不是說女人要離開父母,與丈夫聯合,.
我們因為太注重我們自己的事情,就好像我們看到一棵樹,我們眼睛就盯在樹的上面,我們卻沒有看到整個樹林,.
聖靈要我們看到整個樹林,要明白神的心意,.
我們今天從這幾節經文,再加上前後文,我們可以看到教會到底是什麼,除了這裡的經文之外,前下文和上下文,我們看到神的心意是什麼,.
教會到底是什麼,第一件事情我們看到的,就是教會是神的心意,從起初開始一直到最終結,都是他所設計的,.
教會是一個對墮落的世人早就預定完美的救贖的計劃,就是耶穌基督的福音,我們常常講到預定的時候,我們腦筋總是轉不過來,.
有時候我們想到預定這兩個字的時候,我們的腦筋可能不能轉過來,.
如果我跟你們講說,我今天預定了我要去做什麼,如果我對我自己講,我今天要預定我會做什麼,我不一定能夠做成的,.
為什麼?因為有很多的因素我沒有辦法控制的,我自己的能力也有限,所以我沒有辦法說我預定一定要做什麼事情,.
像我這個年齡,我現在都不敢講這些事情,特別是馬牧師提到,如果我這個年齡的話,我有很多事情我未必能夠做成,.
如果不是神允許,我沒有把握的,如果不是神允許的,我是沒有把握的,教會是神早就預定好了,.
是祂的計劃,要對人類的救贖,就是耶穌基督的福音,第三個,教會是神所揀選的,.
第四個,教會是父神為祂獨生愛子所預備的新娘,最後一個,教會是神所愛的,.
我們先跳過這裡,我們先講教會是神從創立世界以前就已經揀選了,.
我們一起讀以弗所書第一章第四節到第六節,我們一起讀,就如神從創立世界以前在基督裡揀選了我們,.
使我們在祂面前成為聖傑,無有瑕疵,又因愛我們,預定我們,借著耶穌基督的兒子的名分,.
使祂榮耀的恩典得著稱讚,這恩典是祂在愛子裡賜給我們的,看到沒有,神在創世以前,.
祂就已經揀選了祂自己的教會,看到沒有,在創世以前,神已經揀選了祂自己的教會,只是還沒有成就,.
一直到主耶穌出現,主耶穌被釘十字架,從死裡復活,才成就了,直到耶穌出現,釘在十字架上面,這件事才會成就,.

$^{41}$我們再把以弗所書第五章三十一到三十二節,我們再念一遍,我們再將以弗所書五章三十一到三十二節,再重新讀一次,.
為這個緣故,人要離開父母,與妻子聯合,二人成為一體,這是極大的奧秘,但我是指著基督和教會說的..
這一處經文也把我們帶回到創世記裡面,當神設立第一對夫婦,就是亞當和夏娃的時候,.
她已經在她的腦筋裡面,對不起,我不能講說她的腦筋,在她心裡面已經有教會的意念了,.
真的是很奇妙,這個是非常奇妙,我們再回到下面一個,教會是帶有神聖的使命,我們下一個看看教會是帶有神聖的使命,.
我們還是繼續把這一段經文念一下,我們先念第九到十一節,我們一起來,為了藉著教會使天上執政掌權的,現在得知神百般的智慧,.
這是造神從萬世以前在我們主基督耶穌裡所定的旨意,我們從這裡面我們看見,這是神給教會的一個神聖的使命,.
使命是什麼呢?使命是什麼意思呢?就是要讓人知道,我們再回到前面一個,為了藉著教會使天上執政掌權的,現在得知神百般的智慧,.
要讓人和天上執政掌權的知道神的智慧,而也要讓他們知道這個旨意是在主耶穌基督裡所定的,.
這個計劃是在耶穌基督裡面所定的,這些都是創世以前神就已經定了,所以在創世以前神已經在耶穌基督裡面定了,.
要宣告什麼?剛才我們已經講了,就是從萬世以前在主耶穌基督裡所定的旨意,就是神在萬世以前在我們主耶穌基督裡所定的旨意,.
祂要向誰宣告?向世人,向所有的人,教會可以用口傳福音,但是教會的存在就已經向世人宣告,.
還有一個就是向天上執政的掌權的,換句話講就是撒旦的權勢宣告,就是向世人或者向天上執政的就是向撒旦宣告,宣告什麼呢?宣告他的榮耀的救贖計劃,.
就是基督他得勝的福音,就是要向這個專門破壞神的工作的權勢宣告,.
就是要向這個專門破壞神的工作的權勢宣告,要宣告基督他是得勝的,你知道這話是主耶穌自己講的,.
在路加福音那邊他講,他說他看見撒旦從天上落下來,後來是聖靈在啟示錄裡面,.
他向使徒約翰啟示,這件事情在天上已經發生了,撒旦和他的使者在神的國裡是沒有份的,.
他們已經被趕到地上來,你們相信這一句話嗎?這是主耶穌自己講的..
接下去我們就要看教會是什麼,教會是神所愛的,不單是所愛,而且是最愛的,.
你想想看,他從未有創世以前一直到終結,教會都在他心裡面,教會一定是他最愛的,.
神愛基督和愛教會,基督也照樣會愛教會,所以我們把約翰福音第十七章二十六節念一下,.
我們一起來,我已將你的名指示他們,還要指示他們,使你所愛我的愛在他們裡面,我也在他們裡面,.
大家聽不聽得明白呢?神的愛在基督裡面,基督的愛在我們裡面,所以我們會愛神,.
我們也會愛自己的弟兄姊妹,主耶穌當他說把父的名指示給教會,就是對神的認識越來越多,.
父神愛基督,愛他的獨生兒子,他也愛我們,所以這個愛在我們裡面,在教會裡面越來越多,.
我們對神的認識越多,我們的愛就越多,當我們認識神越多的時候,我們對神的愛就越多,.
弟兄姊妹就能夠在基督裡彼此相愛,在這個slide裡面我們看見,正如基督愛教會,為教會捨己,.
你們記得約翰福音三章十六節講什麼嗎?會背的舉手,你們其他沒有舉手的,你們回去的時候一定要把三章十六節背下來,.
神愛世人,甚至將他的獨生子賜給他們,教一切信他的,不致滅亡,罕得永生..
在這裡告訴我們神如此的愛世人,或者是說神如此的愛教會,所以把他的獨生兒子賜給他們,.
猜他的獨生子來到這個世界上,這是神愛世人,愛教會所做出來的行動,.
基督愛教會,他的行動是什麼?他的行動就是為教會捨己,換句話說就是為教會捨命,.
就是為教會受苦,為教會被釘死在十字架上而死,接下來我們就會有一個問題,.
我們常常會問說,基督不是無罪的嗎?基督有沒有罪?說有的舉手,我希望你們沒有人會舉手,.
基督是無罪的,那既然是無罪的,那為什麼要為罪受苦?那麼基督被釘死在十字架上受苦是必須的嗎?.
基督在十字架上所受的苦是必然的嗎?.
好,那麼我們現在來看以賽亞先知怎麼說,我們來看看以賽亞先知是說些什麼?.
我念給你們聽,耶和華就定義,有一個括弧上面寫著喜悅,將他壓傷,使他受痛苦,.
耶和華以他為贖罪祭,又說卻定義在括弧裡面說喜悅,將他壓傷,使他受痛苦,又說以他為贖罪祭,.
換句話講,在父神的心裡面,他的兒子來到這個世界上,一定要受苦,.

$^{81}$所以在這裡講到,神的心裡就說,當耶穌來到這個世界上是必須要受苦,.
這是我們很多時候沒辦法了解的,我們中間有沒有哪一位父親定義要你的兒子受苦,讓他死掉?.
我不相信我們通宗有任何人會定義希望他的兒子受苦受死,.
沒有人啊,不可能有這個人,神卻做了這件事情,.
他使他成為贖罪祭,為世人贖罪,目的是要把罪人領到神的面前,.
讓我們能夠與神和好,主耶穌必須要死嗎?.
主耶穌是不是必須要死呢?當主耶穌要誕生的時候,.
我們常常過聖誕節的時候,我們很高興,哇,Baby Jesus,Baby Jesus,.
當時每次聖誕節我們都很高興,耶穌這個嬰孩,耶穌這個嬰孩,.
當主耶穌誕生為人的時候,神賜給他一個有血有肉的身體,.
跟你跟我是一模一樣的,一樣的身體,.
有很多人稱這個身體叫臭皮囊,我不曉得你們廣東話裡面有沒有這一句話,.
我們這個人的身體是個臭皮囊,我們廣東人都會有一句話,.
有啊,我們這個臭皮囊,為什麼叫臭皮囊,就是很多牧師這樣講,臭皮囊是抵擋神的,.
為什麼講到自己臭皮囊,臭皮囊是抵擋神的,是不義的,.
聖經上告訴我們說,不義的不能承受神的國的,.
不能夠與神聯合的,.
與主耶穌他自己神的生命是沒有辦法合在一起的,.
不能夠與神耶穌這個身體聯合在一起的,.
你們有沒有想過,主耶穌在地上活了33年半,.
他就是穿著這個臭皮囊,.
他自己就穿著這個臭皮囊,.
所以主耶穌在這個世界上33年半時時刻刻都在受試探,.
所以耶穌基督在這個肉體上33多年的話,一直都受到試探,.
可是聖靈告訴我們,他沒有犯過罪,.
要讓人跟神聯合,除了死沒有其他的方法,.
光是死還不夠,還要有一個新的生命,.
所以父神叫他從死裡復活,.
給他一個新的生命,.
新的人的生命,.
因為你們要知道,主耶穌沒有靠自己的能力,就是他本身的神能,.
他完全是順服,.
剛才我們詩歌裡面,我們也聽到剛才我們那個詩歌的時候,對不對,他是順服以至於死,而且死在十字架上,.
所以神使他從死裡復活,.
勝過陰間的權柄,.
把他升為至高,勝過了掌死權的撒旦,.
這就是我們從聖經上看到的,特別是在希伯來書第二章第十四節,.
特別是我們在聖經裡看到,特別是在伯來書那裡,.
他要藉著死敗壞那個掌死權的,就是魔鬼,.
我們從聖經上我們真的是看見,主耶穌他從死裡復活,都要靠神的恩典和聖靈的大能,.

$^{121}$那麼你我我們這些罪人,我們也只能信靠順服,靠著神的恩典,我們才能夠與主聯合,.
這個叫做 Grace of Union,主耶穌就是靠著這個恩典,讓我們能夠與神聯合,.
接下去另外一個問題就出來了,我們看一看教會,我們看到教會的宗派能力,哪裡有教會哪裡就有紛爭,.
我們不要看別人的教會,我們就看看我們自己,越看越不像神能夠靠著我們完成他自己的心意,.
我們怎麼能夠完成神給教會的使命,不要急,我說過了,這是神完美的設計,.
我們再回到耶穌說的第五章二十六節到二十七節,我們就不念了,但是我把它念給你們聽,.
他說要用水界的道把教會洗淨,成為聖潔可以獻給自己,做個榮耀的教會,毫無玷污咒文等類的病,乃是聖潔沒有瑕疵的..
接下去他說,從來沒有人恨惡自己的身子,對不對?你們有沒有哪一個恨惡你的身子的?.
我們總是保養固其,對不對?每天要買菜的時候一定要注意是organic的,都是養品,這個叫做保養品,.
每天要做運動,我們年紀大的人每天都要做運動,對不對?我們就是保養我們自己的身體,基督在教會也是這個樣子,因為我們教會是基督的身體,.
他說你們這因性蒙神能力保守的人,必能得到所預備到末世要顯現的救恩..
要用水藉著道把教會洗淨成為聖潔,這個是在地上的教會才能夠發生的,所以我們在教會裡面我們要傳講神的話,我們有茶經,我們有禱告會,我們有主日學,你們有沒有參加主日學呢?.
要參加主日學,要參加禱告會,要參加讀經會,要聽傳道人所講的信息,為什麼?因為這是主耶穌在預備祂自己的教會,我不知道我們當中有沒有人,你信了主之後,你不太參加教會的?.
這是不對的,因為對你屬靈的生命只有破壞,沒有好處的..
在禮物所書第三章十六節那邊,聖靈感動保羅,他為教會禱告,他為教會禱告什麼事情?.
他是求聖靈剛強我們裡面心照的人,他肯求神更新我們裡面的生命,心照的人使基督因著信住在我們心裡,就是因為基督因著信住在我們裡面..
下面這句話要仔細聽,叫我們能和眾聖徒一同明白基督測不透的愛,一個人信了主之後,只有在教會中與弟兄姊妹碰撞,.
你有沒有和弟兄姊妹碰撞過?.
有沒有和弟兄姊妹碰撞過?.
碰撞的結果是什麼?.
都是他不好..
還是他不好?.
弟兄姊妹我告訴你錯了,當你跟人碰撞的時候,什麼東西被撞出來?.
你碰撞期間有什麼可以撞出來?.
你試試看,下去之後拿一個杯子裝滿熱咖啡在裡面..
如果你可以試一下,回到家找一杯裝滿熱咖啡..
等一下就可以試了..
等一下你就可以試了..
然後你就故意找這個,我都稱他小朋友,跟這個Anissa,跟他撞一下..
你可以試一下,拿一杯熱咖啡去碰一下,特別是Anissa..
什麼東西被撞出來?.
就是你杯子裡面的咖啡..
是不是杯中的咖啡呢?.
你的咖啡是燙的..
咖啡是熱的..
就會燙死你..
熱死你..
你那個咖啡如果是,你那個杯子裡面如果是很黑很黑的水..
如果你的杯子裡面是很黑很黑的水..
撞出來的就是黑水..

$^{161}$撞出來的就是黑色的水..
你那個杯子裡面如果是白的乾淨的水..
如果你的杯子裡面是乾淨的水..
你撞出來的就是乾淨的水..
你撞出來的就是乾淨的東西..
只有當我們在教會裡面跟弟兄姊妹碰撞摩擦..
只有我們在教會裡面跟弟兄姊妹碰撞摩擦..
我們才知道原來我是這個樣子..
我們才明白我們自己是這個樣子..
我才能夠明白,才能夠經歷到基督在我身上的愛..
所以用水接著倒,把教會洗淨..
這就是聖靈在地上教會裡的工作..
所以用水來潔淨我們身體地上的工作..
讓主可以把教會獻給祂自己,做個榮耀的教會..
讓教會能夠獻給自己成為一個榮耀的教會..
接下來我們要看一個非常非常榮耀的關係..
我們講一講一個很榮耀的關係..
祂預備我們做什麼?.
就是做一個新娘子..
要我們做一個新娘..
要準備迎接新郎..
我們去預備去迎接這個新郎..
這是在啟示錄裡面告訴我們講得很清楚..
這是在啟示錄裡面告訴我們講得非常清楚..
這是在啟示錄裡面告訴我們講得很清楚..
我們從創世記在啟示錄的時候,我們看見神為亞當預備了一個妻子..
我們從創世記我們看見神為亞當預備了一個妻子..
祂從亞當身上取下了一根肋骨..
祂在亞當身上取下一條肋骨..
所以夏娃跟亞當是一模一樣的..
所以亞當夏娃是一模一樣的..
因為祂是從祂身上取出來的..
因為祂是從亞當身上取出來的..
聖經上告訴我們,特別是在創世記裡面告訴我們,.
祂說人要離開父母,與妻子聯合..
所以在創世記祂說人要離開父母,與妻子聯合..
主耶穌離開了父家..
基督離開了父家..
與神為祂揀選的妻子聯合..
與祂所揀選的妻子聯合..

$^{201}$那麼教會就是祂的妻子..
教會就是祂的妻子..
所以我們在地上的時候我們受聖靈的裝備..
預備好了,等候祂的丈夫..
等候祂的丈夫..
你們想不想參加這個婚禮啊?.
大家想不想參加這個婚禮呢?.
這個婚禮是在天上的..
這個婚禮是在天上..
是最完美,最榮耀的婚禮..
從創世記一直到啟示錄..
因為祂告訴我們高揚娶妻的時候到了..
新婦也自己預備好了..
新婦已經預備好了..
然後祂又告訴我們,祂說新婦裝扮整齊..
祂說到新婦裝扮整齊..
等候丈夫..
你們有沒有參加一個婚禮,新娘子穿著睡衣,拖鞋出來..
你們有沒有參加一個婚禮,新郎是穿著睡衣,拖鞋的呢?.
我們因為弟兄姊妹,前一陣子才參加誰的婚禮,還記不記得?.
輸一個寶貝女兒,是不是?.
你看她裝扮得整整齊齊的..
然後那個會堂裡面也是打扮得非常,裝扮得非常好看..
而且禮堂也裝扮得非常好看..
將來教會就是要這樣子迎接她的丈夫..
我們講到這裡,我們有兩個問題值得我們思考..
第一個問題,這麼完美的設計,榮耀的結果..
一個這麼完美的計劃,這麼榮耀的目的..
你跟我,或者是教會是不是應該完完全全的信賴神?.
我們教會怎麼能夠完完全全的信靠神?.
甚至我們在患難中也能夠歡歡喜喜的..
我們現在看這個社會上的情形,這個日子會不會好過?.
會越來越難的..
特別是對神的教會..
那麼我們在患難中,我們是不是也看到這個榮耀的盼望,看到這個完美的設計,我們也能夠歡歡喜喜的..
那麼我們在患難中,我們是不是也能夠看到這個榮耀的盼望,看到這個完美的設計,我們也能夠歡歡喜喜的..
那麼這麼完美的設計,這麼榮耀的墓地..
那麼這麼完美的設計,這麼完美的設計..
神要什麼樣的人來引導,牧養祂所愛的地上的教會呢?.
我今天不是來講牧養的人的條件,但是我要用三處經文讓大家看見什麼樣的人神要擷取來領導,牧養祂自己的教會..

$^{241}$但是我會用三處不同的經文來講神是選擇什麼樣的人來牧養在地上的教會..
人若不知道管理自己的家,焉能照管神的教會呢?出入教會不可做監督,恐怕他自高自大,就落在魔鬼所受的刑罰裡.監督也必須來教會裡換有好名聲,恐怕被人毀謗,落在魔鬼的網羅裡..
監督既是神的管家,必須無可指責,不任性,不暴躁,不因酒之事,不打人,不貪無益之財,樂意接待遠人,好善,莊重,公平,聖潔,自持,堅守所教真實的道理,就能把純正的教訓勸化人,又能把爭辯的人駁倒了..
我這做長老,做基督受苦的見證,同享後來所要顯現之榮耀的,勸你們中間與我同做長老的人,勿要牧養在你們中間神的群羊,按著神旨意照管他們,不是出於勉強,乃是出於甘心..
也不是因為貪財,乃是出於樂意,也不是瑕疵所託付你們的,乃是做群羊的榜樣.到了牧長顯現的時候,你們必得找那永不衰殘的榮耀冠冕..
這是一個挑戰,特別是對弟兄來說.我自己看這是聖靈的心意,所以我也盼望在我們當中的弟兄們,你真的是能夠像使徒所寫的,要羨慕善功,.
把自己交給主,我們剛才所念的,我們就是一個印象,我告訴你,很難做到的.你一定要靠著神的恩典,就好像我們長老們他們所分享的,你必須要是一個謙卑順服的人..
弟兄們,你可以跟神說,你說神啊,我願意做一個長老,那麼你如果在神面前講了這一句話,你就要開始操練你自己..
不單是對神的認識上面要操練自己,而且對自己的品格要操練,花時間的,花時間..
我知道有一個人,他在年輕的時候,當他第一次看到提摩太前書第三章,第一節到第七節的時候,聖靈感動他,讓他在神面前做了一個禱告,他就跟神說,神啊,照著這個,我願意操練我自己,將來做一個長老..
經過了很多年,他沒有成為一個長老,但是神呼召他做一個牧師,弟兄們,這是我對你的挑戰,不是我的挑戰,這是聖靈的心意..
姐妹們怎麼辦呢?我告訴你,你可以為這些弟兄們禱告,你可以為你的丈夫祈禱,鼓勵他,操練自己,在神的面前跟神說,我願意..
我們當中也有人,目前沒有丈夫的,你可以為教會裡面的弟兄禱告,特別是那些年輕的弟兄們,就為他們禱告,求神聖靈在他們心裡面插一刀,.
讓他們願意羨慕善功,把自己完完全全獻給神,願意的舉手,願意的舉手,年輕的弟兄們都不敢舉手,.
神這麼完美的設計,這麼榮耀的盼望,難道不能把你的心激動起來嗎?.
\newpage



\section{約書亞記 7:1-26}
\label{sec:031795}
\textbf{暗中的罪,眾人的代價}~~~~李美儀傳道
\newline
\newline
連結: \href{http://www.church.com.hk/acms/content.asp?site=cdc\&op=show\&type=product\&code=031795\&layout=sermon}{\texttt{031795}} ~~~~ 語音日期: 2026-03-08
\newline
\newline
\hyperref[sec:031794]{\small{< < < PREV SERMON < < <}}
~
\hyperref[sec:index_chronic]{\small{[返順時目]}}
~
\hyperref[sec:031796]{\small{> > > NEXT SERMON > > >}}
\newline
\newline
約書亞記 7:1-26
\newline
\begin{longtable}{cl}
\hline
\hline
章節 & 經文 (和合本修訂版)\\
\hline
7:1 & \begin{tabularx}{0.7\textwidth}{X} 以色列人在當滅之物上犯了罪。猶大支派中，謝拉的曾孫，撒底的孫子，迦米的兒子亞干取了當滅之物，耶和華的怒氣就向以色列人發作。 \end{tabularx} \\ \\ \relax
7:2 & \begin{tabularx}{0.7\textwidth}{X} 約書亞從耶利哥派人往伯特利東邊，靠近伯‧亞文的艾城去，對他們說：「你們上去窺探那地。」那些人就上去窺探艾城。 \end{tabularx} \\ \\ \relax
7:3 & \begin{tabularx}{0.7\textwidth}{X} 他們回到約書亞那裡，對他說：「眾百姓不必都上去，只要二、三千人上去就能攻取艾城；不必勞動眾百姓都上去，因為他們人少。」 \end{tabularx} \\ \\ \relax
7:4 & \begin{tabularx}{0.7\textwidth}{X} 於是百姓中約有三千人上那裡去，但他們竟在艾城的人面前逃跑。 \end{tabularx} \\ \\ \relax
7:5 & \begin{tabularx}{0.7\textwidth}{X} 艾城的人擊殺他們約三十六人，從城門前追趕他們，直到示巴琳，在下坡的地方擊敗他們。他們都膽戰心驚，融化如水。 \end{tabularx} \\ \\ \relax
7:6 & \begin{tabularx}{0.7\textwidth}{X} 約書亞和以色列的長老就撕裂衣服，在耶和華的約櫃前臉伏於地，直到晚上。他們把灰撒在頭上。 \end{tabularx} \\ \\ \relax
7:7 & \begin{tabularx}{0.7\textwidth}{X} 約書亞說：「唉！主耶和華啊，你為甚麼領這百姓過約旦河，把我們交在亞摩利人手中，使我們滅亡呢？我們不如住在約旦河的那邊！ \end{tabularx} \\ \\ \relax
7:8 & \begin{tabularx}{0.7\textwidth}{X} 主啊，求求你，以色列人既在仇敵面前轉身逃跑，我還有甚麼可說的呢？ \end{tabularx} \\ \\ \relax
7:9 & \begin{tabularx}{0.7\textwidth}{X} 迦南人和這地所有的居民聽見了就必圍困我們，把我們的名從地上除去。那時，你為你至大的名要怎樣做呢？」 \end{tabularx} \\ \\ \relax
7:10 & \begin{tabularx}{0.7\textwidth}{X} 耶和華對約書亞說：「起來！你的臉為何這樣俯伏呢？ \end{tabularx} \\ \\ \relax
7:11 & \begin{tabularx}{0.7\textwidth}{X} 以色列犯了罪，又違背了我所吩咐他們的約，又取了當滅之物。他們又偷竊，又行詭詐，又把那當滅的物與自己的器皿放在一起。 \end{tabularx} \\ \\ \relax
7:12 & \begin{tabularx}{0.7\textwidth}{X} 因此，以色列人在仇敵面前站立不住。他們在仇敵面前轉身逃跑，因為他們成了當滅的物。你們若不把當滅的物從你們中間除掉，我就不再與你們同在了。 \end{tabularx} \\ \\ \relax
7:13 & \begin{tabularx}{0.7\textwidth}{X} 你起來，去叫百姓分別為聖，說：『你們要為了明天使自己分別為聖，因為耶和華－以色列的神這樣說：以色列啊，在你中間有當滅的物；你們若不把你們中間當滅之物除掉，你在仇敵面前必站立不住！』 \end{tabularx} \\ \\ \relax
7:14 & \begin{tabularx}{0.7\textwidth}{X} 到了早晨，你們要按著支派近前來。耶和華所選的支派，要按著宗族近前來；耶和華所選的宗族，要按著家族近前來；耶和華所選的家族，要按著男丁，一個一個近前來。 \end{tabularx} \\ \\ \relax
7:15 & \begin{tabularx}{0.7\textwidth}{X} 被選的人有當滅之物在他那裡，他和他所有的必被火焚燒，因為他違背了耶和華的約，又因他在以色列中做了愚妄的事。」 \end{tabularx} \\ \\ \relax
7:16 & \begin{tabularx}{0.7\textwidth}{X} 於是，約書亞清早起來，召以色列按著支派近前來。選出來的是猶大支派。 \end{tabularx} \\ \\ \relax
7:17 & \begin{tabularx}{0.7\textwidth}{X} 他召猶大的宗族近前來，選出來的是謝拉宗族。他召謝拉宗族，按著男丁，一個一個近前來，選出來的是撒底。 \end{tabularx} \\ \\ \relax
7:18 & \begin{tabularx}{0.7\textwidth}{X} 他召撒底的家族，按著男丁，一個一個近前來，就選出猶大支派，謝拉的曾孫，撒底的孫子，迦米的兒子亞干。 \end{tabularx} \\ \\ \relax
7:19 & \begin{tabularx}{0.7\textwidth}{X} 約書亞對亞干說：「我兒，我勸你將榮耀歸給耶和華－以色列的神，在他面前認罪，把你所做的事告訴我，不可向我隱瞞。」 \end{tabularx} \\ \\ \relax
7:20 & \begin{tabularx}{0.7\textwidth}{X} 亞干回答約書亞說：「我實在得罪了耶和華－以色列的神。這是我所做的： \end{tabularx} \\ \\ \relax
7:21 & \begin{tabularx}{0.7\textwidth}{X} 我在所奪取的財物中看見一件美好的示拿外袍，二百舍客勒銀子，一條重五十舍客勒的金子。我貪愛這些物件，就拿去了。看哪，這些東西都埋在我帳棚內的地裡，銀子在外袍底下。」 \end{tabularx} \\ \\ \relax
7:22 & \begin{tabularx}{0.7\textwidth}{X} 約書亞就派使者跑到亞干的帳棚裡。看哪，那件外袍藏在他的帳棚裡，銀子在外袍底下。 \end{tabularx} \\ \\ \relax
7:23 & \begin{tabularx}{0.7\textwidth}{X} 他們從帳棚裡把這些東西取出來，拿到約書亞和以色列眾人那裡，倒在耶和華面前。 \end{tabularx} \\ \\ \relax
7:24 & \begin{tabularx}{0.7\textwidth}{X} 約書亞和以色列眾人把謝拉的曾孫亞干和那銀子、那件外袍、那條金子，以及亞干的兒女、牛、驢、羊、帳棚，和他所有的，都帶著上到亞割谷去。 \end{tabularx} \\ \\ \relax
7:25 & \begin{tabularx}{0.7\textwidth}{X} 約書亞說：「你為甚麼給我們招惹災禍呢？今日耶和華必使你遭受災禍。」於是以色列眾人用石頭打死他，用火焚燒他們，把石頭扔在其上。 \end{tabularx} \\ \\ \relax
7:26 & \begin{tabularx}{0.7\textwidth}{X} 眾人在亞干身上堆了一大堆石頭，直存到今日。於是耶和華轉意，不發他的烈怒。因此，那地方名叫亞割谷，直到今日。 \end{tabularx} \\ \\
[1ex]
\hline
\hline
\end{longtable}
$^{1}$我們今天繼續講到約書亞記的內容.
上次曾牧師講到以色列人在耶利哥城的得勝.
指出以色列人對上帝單純的信靠.
即時的行動.
以至整個群體在前行的路上得到上帝的祝福.
最終他們得勝.
不過今天接下來要讀的約書亞記第七章的內容.
雖然同樣是攻城.
但和上次有截然不同的結局.
今次以色列民是出乎意料地大敗.
為何會有這樣的結局呢?.
在七章一節.
聖經的作者給了我們一句主題句.
他說以色列人在當滅的物上犯了罪.
因為由大支派中.
謝拉的曾孫 撒底的孫子.
加米的兒子阿干取了當滅的物.
耶和華的怒氣就向以色列人發作.
以色列人從輝煌的得勝.
到今天慘痛的失敗.
其實這個巨大的反差.
都是在提醒我們信徒一個很重要的屬靈功課.
就是在基督裡面.
個人和群體.
其實是擁有一個不可以分割的關係.
我們說在基督裡面合一.
其實信徒自然而然.
就成為一個榮辱與共的群體.
所以其實透過今天的經文.
再次去嚴正告誡我們信徒.
個人暗中的罪.
其實是需要眾人去負上代價的.
在我們分享之前.
我們先低頭做一個祈禱.
天父 我們求你的話語.
醒察我們的內心.
使到我們立刻為我們內心顯而未現的罪.
未同你求饒恕的罪.
向你去悔改.
以致我們同教會.

$^{41}$再一次去經歷你自己的祝福.
以上的祈禱.
是奉靠著耶穌基督的名求.
阿們.
我們看看究竟在經文裡面.
作者怎樣一步一步帶領著我們.
去揭開這個慘痛的教訓.
所以今天的經文第七章.
可以分為三個序幕和一個中章.
一層一層揭示.
罪怎樣影響整個群體.
不過因著神的恩典.
亦都透露著神.
怎樣帶領著以色列的子民.
返回正路裡面.
在第二至第九節.
記載了以色列人攻陷的失敗.
以及約書亞去尋問上帝的經過.
讓我們一起去讀出二至九節.
他們回到約書亞那裡.
對他說.
眾民不必都想去.
只要二三千人想去.
就能攻取艾城.
不必勞累眾民都去.
因為那裡的人少.
於是民眾若有三千人想那裡去.
竟在艾城人面前逃跑了.
艾城的人擊殺了他們三十六人.
從城門前追趕他們.
直到士巴林再下坡殺敗他們.
眾民的心就消化如水.
約書亞便撕裂衣服.
他和以色列的長老把灰撒在頭上.
在耶和華的約跪前.
俯伏在地直到晚上.
約書亞說:哀哉!主耶和華!.
你為什麼竟領著百姓過約旦河.
將我們交在阿摩利人的手中.
使我們滅亡呢?.

$^{81}$我們不如住在約旦河那邊倒好.
主啊!以色列人既在仇敵面前轉背逃跑.
我還有什麼可說的呢?.
迦南人和這裡一切的居民聽見了.
就必圍困我們.
將我們的命從地上除滅.
那時你為你的大名要怎樣行呢?.
事實上艾城和耶利哥城有不同的地勢.
艾城位於伯特利東邊.
是進入迦南內陸的一個重要據點.
從地點來看.
艾城其實控制著從約旦河谷.
通往山地的一個重要通道.
我們見到伯特利在它旁邊.
其實就是它們兩城相近的時候.
形成一個聯防.
所以其實從地勢來看.
或者從它的規模防備來看.
其實它們都是一個堅固的城邦.
雖然它們較細.
亦都人口較少.
但它們防禦的工作是完備的.
擁有城門.
擁有著一個很堅固的守衛系統.
它們處於高地.
所以屬於易守難攻.
而人口和軍力裡面.
聖經記載.
它們全城男女有一萬二千人.
但是一萬二千人裡面.
亦都擁有訓練有素的軍人.
城中亦有皇帝.
它們的政治體系都是完整.
雖然艾城不比耶利哥城的人多.
但從艾城這個險要的地勢來看.
其實都是一個不容忽視的對手.
但是經文剛剛告訴我們.
以色列人做錯了判斷.
在第三節他們認為.
眾民不必都想去.

$^{121}$只要二三千人想去.
就能攻取艾城.
結果是怎樣?.
未攻進城.
他們就已經落荒而逃.
從對這場戰役的敘述來看.
其實以色列人的失敗.
似乎都是有些實際的原因.
可能是輕視對方的軍力.
認為人少.
所以我們就不用有更嚴謹的軍事策略.
又或者在這個險要的地方.
他們只選擇了一條單一的路線.
認為這樣就已經足夠.
所以從行軍的角度來看.
似乎這個檢討都是很合理的.
不過約書亞很清晰.
這場戰爭對於以色列人來說.
其實不只是一場簡單.
熟悉的土地之爭.
它是一場屬靈的征戰.
就是以色列人逐漸去經歷上帝的應許.
直至他們進入迦南尼地.
也是因為這樣.
他們一直都相信神的應許.
認為神一定是保守他們所向披靡的去殺敵.
但是到如今.
事與願違的時候.
面對功成的失敗.
第一步我們見到.
其實他們不只是停留在檢討軍力和策略.
他們最重要第一步是尋尋上帝.
親愛的弟兄姊妹.
當我們在生活裡面遇到挫敗.
或者一些事情不如預期的時候.
可能特別是在教會的工作裡面.
我們第一個反應是什麼呢?.
我們是否急於分析問題.
我們找一些人為的解決方法.
還是我們習慣學習.

$^{161}$約書亞一樣.
我們先來到神的面前去求問.
檢視自身和教會的屬靈警方.
有時候我們以為犯罪是因為.
恥於人有,敗於人有.
但他是忘記一開始犯罪的機會.
可能是因為我們遠離了上帝.
當人慢慢去遠離神的時候.
其實就醞釀著一個犯罪的心.
回到經文裡面.
其實當以色列人.
去攻艾城的時候遭遇失敗.
約書亞的反應就是撕裂衣服.
殺灰在頭上.
其實就是表達一種很悲傷.
一個疑惑的心.
他向神祈求.
俯伏在神的面前.
求問神為何使我們滅亡.
將我們的命從地上除滅.
那時你為你的大名要怎樣行呢?.
結果聖經說.
耶和華如何回答他?.
結果神回應約書亞.
他說以色列人犯了罪.
讓我們繼續讀十節至十二節.
耶和華吩咐約書亞說.
起來,你為何這樣俯伏在地呢?.
以色列人犯了罪.
違背了我吩咐他們的約.
取了當滅的物.
又偷竊又行詭詐.
又把那當滅的放在他們的家裡.
因此以色列人在仇敵面前站立不住.
他們在仇敵面前準備逃跑.
是答應了被救助的.
你們若不把當滅的物.
從你們中間除掉.
我就不再與你們同在了.
耶和華用了兩個很嚴厲的說話.

$^{201}$詞語「犯了罪」和「違背」.
違背是甚麼意思呢?.
違背就是越過上帝給我們所設的界線.
或者是違反了上帝的命令.
這個就是違背.
這個可能未必指到.
我們世上犯了違法的事.
其實在聖經的場景.
或者在屬靈的層面裡.
更多是指向對信任的背叛.
斯托特穆斯提到違背神.
不單只是外在不信服的行為.
而是人類試圖去散奪神的位置.
拒絕神的權威.
選擇以自己成為生命的中心.
這種博弈其實核心就是身份的錯置.
我們很想成為自己的主宰.
而不是將神成為我們生命的王.
其實在約書亞記開始的時候.
神已經不斷強調他們必須要持守聖潔.
就是分別為聖.
從不忠於上帝的事分別出來.
而神的指紋就是透過.
謹守遵行神的道來去實踐.
以致他們能夠享受神的賜福和順利.
神的應許就是我們得勝的關鍵.
不過這個祝福也不是自動給我們的.
而是建立在一個重要的回應裡.
就是我們要謹守遵行神的道.
這兩個並行的關係不是因果的關係.
而是在屬靈裡面.
神和人的盟約底下的一個關係.
當以色列人去信服神的時候.
人就得著神的應許.
經歷著成功征服迦南的事情.
相反,當我們違背神的命令的時候.
罪就阻隔了上帝的同在.
我們就可能遭遇失敗和懲罰.
不過值得反思的是.
明知如此,為何也犯罪呢?.

$^{241}$可能是來自我們人性的驕傲.
或者有時我們對神命令的輕視.
甚或我們會覺得上帝不會罰我們.
眼前的引誘和利益.
吸引著我們離開上帝.
甚或我們平日很少去深刻.
去經歷上帝與我們的同在.
而最終我們對上帝的命令的理解是有所偏差.
其以上這些原因歸根究底.
反映了一個深層的問題.
就是我們對恩典的誤解.
當信徒以為神的慈愛.
可以無條件無限地包容我們的時候.
卻忽略了一個信復的呼召.
我們的信仰慢慢變成一種廉價的交易.
其實這種現象不只是出現在聖經裡.
歷世歷代的教會信徒.
甚至我們這一刻.
我們每一個可能都在面對這個試探.
在二十世紀德國神學家潘弗華.
正是處於這個處境裡.
他向教會裡發出一個很震撼的提醒.
在他的著作《追隨基督》裡.
他提到廉價恩典和寵價恩典的對比.
他對於1930年代的德國教會.
信徒的信仰和生活他覺得嚴重的脫節.
他看到一個現象.
他說人們以為只要受針參加禮拜.
口頭認順就能夠自動獲得救恩.
卻完全不需要理會.
耶穌基督對門徒的呼召和要求.
這種恩典其實慢慢讓我們教會.
失去了先知性的聲音和改變世界的能力.
反觀我們今天.
我們會不會都不知不覺之間.
將恩典當作理所當然的保障呢?.
而不是我們進來信耶穌.
我們就知道需要為寵價恩典付上代價的道路呢?.
回到經文裡.
這次以色列人的失敗.

$^{281}$其實源自於一個很實質的原因.
就是取了神吩咐不可取的當滅之物.
在攻打耶利哥城的時候.
其實神已經很清楚告訴他們.
祂說:至於你們勿要謹慎.
不可取那當滅的物.
恐怕你們取了那當滅的物.
就連累以色列的全營.
原來在神的裡面.
從來都沒有隱蔽的事.
一個人暗中犯罪.
最終連累了整個的以色列民族.
個人的不信服.
為整個的群體帶來一個災難性的後果.
而接下來的經文也都讓我們看到.
神怎樣一步一步去揭露這件事的真相.
請聽我讀出十六至二十五節的內容.
「於是約書亞清早起來.
使以色列人按著支派近前來.
取出來的是猶大支派.
使猶大支派近前來.
就取了謝拉的宗族.
使謝拉的宗族按著加瑟人丁.
一個一個地近前來.
取出來的是撒底.
使撒底的加瑟按著人丁.
一個一個地近前來.
就取出猶大支派的人.
謝拉的曾孫.
撒底的孫子.
加米的兒子阿剛」.
在這個審問的過程裡.
約書亞採取了一個漸進式的抽籤方法.
去揭露這個罪行.
這個過程單是為了找出罪魁禍首.
同時這個過程是要全以色列民.
去認識這個罪的嚴重性.
從支派去到宗族.
宗族就是家族的意思.
從宗族去到個人.

$^{321}$最終真相完全的揭露.
經過一輪的查問.
最終他們找到謝拉宗族的阿剛所做.
而在接下來的經文.
阿剛都完全去承認了.
自己確實是奪去耶和華吩咐.
不可取的當滅之物.
約書亞對阿剛說.
「我兒.
我勸你將榮耀歸給耶和華.
以色列的神.
在他面前認罪.
將你所做的事告訴我.
不要向我隱瞞」.
阿剛回答約書亞說.
「我實在得罪了耶和華.
以色列的神.
我所做的事如此如此.
我在所奪的財物中.
看見一件美好的士拿衣服.
二百瑟克勒銀紙.
一條金紙重五十瑟克勒.
我就貪愛這些物件.
便拿去了.
現今藏在我帳篷內的地裡.
銀紙在衣服底下」.
我們看看阿剛犯罪的過程.
完全說明了罪在人心裡的機制.
他看見.
之後貪愛.
然後就取去.
看到一件很美好的士拿衣服.
他就貪愛這些貴重的物品.
之後藏在帳篷裡.
試探是具有一種漸進式的隱憂.
將人一步一步推向犯罪.
在這個過程裡.
我們從阿剛的認罪清楚看到.
三個階段.
眼目的試探.

$^{361}$內心的貪戀.
和行動的犯罪.
這三個犯罪的三部曲.
是不是很熟悉呢?.
像不像《創世記》裡.
夏娃犯罪的過程.
「我看見,我貪愛,我就取了」.
在《新約聖基雅國書》裡說.
「但各人被試探,乃是被自己的私慾牽引,被誘惑」.
「私慾既懷了胎,就生出罪來」.
「罪既長成,就生出死來」.
其實在漢奕的階段.
我們就應該要警醒.
我們要為自己的眼目設立界線.
特別是在這個資訊發達的時代.
我們看到什麼.
可能就慢慢塑造我們成為什麼.
網上有一個很有趣的圖片.
它將今天人看手機上癮的程度.
比喻為吸食鴉片.
親愛的大英姐妹.
我相信我們都是機不離手的.
我們每天打開手機.
瀏覽社交網站.
我們慢慢接觸很多資訊的導入.
可能開始的時候.
我們都會警覺有些網站不好看.
甚至有些資訊不合乎神的度.
但可能太多了.
它成為你唯一吸收資訊學習的途徑.
我們已經覺得這些東西無害.
甚至我們不會去過濾.
很有趣的.
這些資訊有時都會以另一種形式出現.
去捉住人的心.
可能是以一個心靈療癒的方式去表現出來.
例如每天的星座運程.
水晶的能量.
命理的運勢.
甚至是求神問卜.

$^{401}$我在神學院裡有一科叫教木學.
裡面有一本書老師要我們看完.
叫解破七宗罪.
這七宗罪有哪七宗呢.
驕傲 嫉妒 憤怒 韜鐵 即是貪食.
情慾 貪婪 懶散.
當說到貪婪的時候.
作者引用了教父阿奎拿對罪的認知.
他說貪婪的罪直接干犯了我們的論社.
因為除非有人損失了他的財富.
否則我們是不會得到任何額外的財富的.
貪婪也跟所有道德上的罪一樣.
是一種干犯上帝的罪.
因為都是為短暫的事去勞碌.
而最後丟棄了上帝永恆的事.
雖然亞干一人犯罪.
但整個的民族需要付上代價.
其實這個不單只是集體的責任.
而是體現了一個很特別的群體生活.
就是盟約之下.
我們原來是緊緊相連.
是有一個緊緊相連的命脈在當中.
就好像聖經裡面所說.
每一個肢體受苦 全身都受苦的群體關係.
當我們說要進入教會.
我們要相信上帝的時候.
我們就成為一個身子.
好像以弗所書所說.
身體只有一個 性靈只有一個.
正如你們蒙照 同有一個子莫.
一主 一信 一浸 一神.
就是眾人的父 超乎眾人之上.
貫乎眾人之中 也住在眾人之內.
同一個身子的關係.
就是一個榮辱與共的關係.
除了個人的信仰.
必須是行在群體的生活裡面之外.
其實我們每一個信徒.
對於整個群體都有一個義務.
就是我們要留意我們個人的得失.

$^{441}$其實是影響整個群體的得失.
也就是說 阿干一樣.
當一個人犯罪.
整個群體其實都是受影響.
但同樣地 其實當一個人得益 得榮耀.
眾之體都一同分享這個喜樂.
所以在第25節裡.
約書亞對阿干說.
你為什麼連累我們呢?.
今日耶和華必叫你受連累.
連累這個字bring trouble.
其實就是說成為全體的絆腳石.
大家記得剛才11節.
當約書亞沉吻上帝的時候.
上帝並不是說阿干犯了罪.
上帝回覆.
祂說 以色列人犯了罪.
的確 從神和子民的文約關係來看.
我們其實是沒有一個純粹的個人信仰.
因為神是呼召整個群體.
以色列人被呼召.
其實就是進入祭司的國度.
聖潔的國民.
我們不是一個人去見證著這件事.
而是一個屬神的群體.
見證進入神的國度當中.
所以當罪存留在一個群體裡的時候.
它可能就好像病毒一樣.
慢慢就會傳染開去.
以至最終整個群體.
去隔絕和上帝的交流.
不能夠體驗上帝的同在.
我們記得剛剛一開始提到的經文.
阿干的罪不單止讓他自己承受代價.
其實連累了一開始攻城的時候.
36個以色列人的士兵陣亡.
其實我們很容易理解.
很多實際的例子.
一個企業的高管如果貪污.
不單止可能股價大跌.

$^{481}$甚至影響了整個公司.
最終或者可能會倒閉.
很多人失業.
一隊的運動員.
如果其中有一個人使用了禁藥.
可能令到整支的隊伍取消資格.
或者一個近一點的見證.
當我們有一個家庭成員犯錯的時候.
可能就會讓整個家庭.
傾盡一切去為他承擔後果.
親愛的弟兄姊妹.
我怎樣看自己和教會的關係呢?.
我相信我們都很重視自己的靈性生活.
不過我們有沒有想到.
這件事都和教會眾人的福祉有關呢?.
約書亞時代其實是在面對一個有形的征戰.
是一場實質的戰爭.
在今天我們其實是在面對一個無形的征戰.
就是一個屬靈的征戰.
我們有沒有被一些罪仍然存留在我們生命裡.
不願意去捨棄呢?.
有什麼罪是容易存在在教會當中呢?.
苦睹,不饒恕,閒言閒語和洩頭的嘴.
又或者我們會不會變得.
對屬靈的事冷淡,不願意毀身.
甚或我們外在的一些背景,學歷.
使得我們變成驕傲和恣意.
神是不會容許我們信徒故意活躍的.
當我們清楚知道某些事是錯的.
我們仍然選擇去做的時候.
這些就叫做明知故犯.
這樣就顯示我們當中沒有悔改.
神也不會讓我們繼續下去.
神會用盡祂的方法去教導我們這些屢勸不改的人.
可能用一個更深的審查.
更能讓我們悔改的事讓我們經歷.
不過同時可能我們的教會也和我們一起必然受管教.
所以到最終章.
到事件的最後.
神要以色列人認真去悔改.

$^{521}$徹底去處理罪的存在.
在第25至26節這樣說.
「於是以色列眾民用石頭打死他.
將石頭拎在旗上,又用火焚燒他所有的.
眾人在阿竿身上堆成一大堆石頭.
直傳到今日.
於是耶和華轉而不發他的烈怒.
因此那地方名叫阿葛谷直到今天」.
阿竿的全家被帶到阿葛谷.
用石頭打死和焚燒.
這麼嚴厲的目的只有一個.
就是將惡和罪從他們民中去除掉.
唯有這樣去潔淨營地.
上帝的同在才會再臨到他們當中.
我在預備這篇信息的時候.
內心浮現了一幅圖畫.
就是一條連接兩個懸崖峭壁的吊橋.
原本這條橋是讓人從一邊安全走到另一邊.
但我們可以想像.
好像這張圖片一樣.
如果當中有一些木條可能是爛了的.
沒了的.
或者是快要斷的.
當人走上去的時候.
不要說要去到對岸.
可能在當中已經有致命的危機.
親愛的弟兄姊妹.
當我們說要活出聖潔的生活的時候.
我們要知道我們要這樣行.
我們才可以見證上帝.
我們才能夠讓人透過我們的見證.
透過教會.
透過我們的努力傳福音的見證.
認識到上帝的信實.
以致他們從這個黑暗的世代裡.
透過教會走向一條永生的道路.
亞干的罪提醒著我們.
在盟約群體裡.
其實我們沒有一個人可以獨善其身.
一個人犯罪.

$^{561}$全體都要負上代價.
不過上帝的恩典夠用.
耶穌基督的代贖.
使得一人悔改.
全體都能夠得到復興.
這就是十字架上的奧秘.
基督一人擔當眾人的罪.
我們卻在他裡面一同得榮耀.
現在是大齋期.
四月就是復活節.
壽苦節的時候.
我們這段時間有沒有特別紀念上帝對我們的恩情呢?.
我們有沒有來到神面前誠心悔改?.
心裡面有沒有覺得.
這些罪是障礙著我們和上帝的交流.
還是在障礙教會承擔上帝的使命?.
我們千萬不要辜負上帝這份寶貴的恩情.
這個重價的恩典.
我們要彼此扶持.
潔淨我們的群體.
以致我們在神的路裡面.
經歷神的同在.
承擔更多的工作.
讓我們一起祈禱.
親愛的天父.
留下寶貴的說話提醒我們.
特別在這個紀念你的日子裡.
我們在你面前認罪悔改.
倘若我們生命裡.
仍然有一些未向你認罪悔改的事.
求你當下提醒著我們.
以致我們願意向你去呼求.
求你的饒恕.
也求你帶領著我們教會.
帶領著我們成為聖潔的群體.
以致我們在這個社區裡面.
真的成為明燈.
照亮人.
能夠去到你的裡面.
奉教主耶穌基督明求.

$^{601}$阿們.
\newpage



\section{箴言 8:12-31}
\label{sec:031796}
\textbf{懇切尋求智慧的必尋見}~~~~蕭壽華牧師
\newline
\newline
連結: \href{http://www.church.com.hk/acms/content.asp?site=cdc\&op=show\&type=product\&code=031796\&layout=sermon}{\texttt{031796}} ~~~~ 語音日期: 2025-08-31
\newline
\newline
\hyperref[sec:031795]{\small{< < < PREV SERMON < < <}}
~
\hyperref[sec:index_chronic]{\small{[返順時目]}}
~
\hyperref[sec:031797]{\small{> > > NEXT SERMON > > >}}
\newline
\newline
箴言 8:12-31
\newline
\begin{longtable}{cl}
\hline
\hline
章節 & 經文 (和合本修訂版)\\
\hline
8:12 & \begin{tabularx}{0.7\textwidth}{X} 我－智慧以靈巧為居所，又尋得知識和智謀。 \end{tabularx} \\ \\ \relax
8:13 & \begin{tabularx}{0.7\textwidth}{X} 敬畏耶和華就是恨惡邪惡；我恨惡驕傲、狂妄、惡道，和乖謬的口。 \end{tabularx} \\ \\ \relax
8:14 & \begin{tabularx}{0.7\textwidth}{X} 我有策略和健全的知識，我聰明，又有能力。 \end{tabularx} \\ \\ \relax
8:15 & \begin{tabularx}{0.7\textwidth}{X} 君王藉我治國，王子藉我定公平， \end{tabularx} \\ \\ \relax
8:16 & \begin{tabularx}{0.7\textwidth}{X} 王公貴族，所有公義的審判官，都藉我掌權。 \end{tabularx} \\ \\ \relax
8:17 & \begin{tabularx}{0.7\textwidth}{X} 愛我的，我也愛他，懇切尋求我的，必尋見。 \end{tabularx} \\ \\ \relax
8:18 & \begin{tabularx}{0.7\textwidth}{X} 財富和尊榮在我，恆久的財寶和繁榮也在我。 \end{tabularx} \\ \\ \relax
8:19 & \begin{tabularx}{0.7\textwidth}{X} 我的果實勝過金子，強如純金，我的出產超乎上選的銀子。 \end{tabularx} \\ \\ \relax
8:20 & \begin{tabularx}{0.7\textwidth}{X} 我在公義的道上走，在公平的路中行， \end{tabularx} \\ \\ \relax
8:21 & \begin{tabularx}{0.7\textwidth}{X} 使愛我的承受財產，充滿他們的庫房。 \end{tabularx} \\ \\ \relax
8:22 & \begin{tabularx}{0.7\textwidth}{X} 「耶和華在造化的起頭，在太初創造萬物之先，就有了我。 \end{tabularx} \\ \\ \relax
8:23 & \begin{tabularx}{0.7\textwidth}{X} 從亙古，從太初，未有大地以前，我已被立。 \end{tabularx} \\ \\ \relax
8:24 & \begin{tabularx}{0.7\textwidth}{X} 沒有深淵，沒有大水的泉源，我已出生。 \end{tabularx} \\ \\ \relax
8:25 & \begin{tabularx}{0.7\textwidth}{X} 大山未曾奠定，小山未有之先，我已出生。 \end{tabularx} \\ \\ \relax
8:26 & \begin{tabularx}{0.7\textwidth}{X} 那時，他還沒有創造大地和田野，並世上頭一撮塵土。 \end{tabularx} \\ \\ \relax
8:27 & \begin{tabularx}{0.7\textwidth}{X} 他立高天，我在那裡，他在淵面的周圍劃出圓圈， \end{tabularx} \\ \\ \relax
8:28 & \begin{tabularx}{0.7\textwidth}{X} 上使穹蒼堅硬，下使淵源穩固， \end{tabularx} \\ \\ \relax
8:29 & \begin{tabularx}{0.7\textwidth}{X} 為滄海定出範圍，使水不越過界限，奠定大地的根基。 \end{tabularx} \\ \\ \relax
8:30 & \begin{tabularx}{0.7\textwidth}{X} 那時，我在他旁邊為工程師，天天充滿喜樂，時時在他面前歡笑， \end{tabularx} \\ \\ \relax
8:31 & \begin{tabularx}{0.7\textwidth}{X} 在他的全地歡笑，喜愛住在人世間。 \end{tabularx} \\ \\
[1ex]
\hline
\hline
\end{longtable}
$^{1}$「錄音或錄影的版權是屬於基督教宣導會不各堂所有..
我們的屬靈傳統之一,是以真理教導為本..
歡迎你將此錄音或錄影的連結傳給親友..
懇請注意,未經授權,請勿轉錄或剪輯本錄音或錄影..
甚願好像《仕途行傳》二十章三十二節下所說:.
「這道能建立你們,叫你們和一切誠信的人同得基業.」.
天地的主,永恆的上帝,.
您一班的子民到您的面前敬拜您,.
主下您的名稱為至高,主下您的能力養館萬有,.
你佩得我們的尊崇,我們的讚美,.
主下求你保守我們,開我們的眼睛,.
能夠見到你的尊榮,見到你的偉大,.
以致我們會輕看我們所面對的困難,.
我們會交託我們所面對的一切挑戰..
如今到您面前來,我們很願意趁我們的聖靈在我們心中工作,.
主下除去很多的憂慮,很多的內心的不安,.
讓我們在您面前知道你的掌權,.
我們就安然地來敬拜你,.
求你賜下你自己聖經的說話,引導我們的心去尋求你,.
引導我們去敬畏你,.
讓我們人生的路行在一條光明無因的路上..
天父我們等候你,求你賜下你的恩賢施恩於我們,.
我們仰望祈禱,奉主耶穌的聖名祈求,阿們..
在這段時間生意很好,也賺了很多錢,.
但後來發覺很多的企業轉移向一些自己同行購買的電腦,.
當他一路查問情況時發覺,.
原來對方很多很多賣出的電腦比他賣出的價錢低很多,.
按他所了解,他覺得沒有可能做的,.
只會虧本,.
但對方就一直出售,所以很多人歡喜向他買電腦,.
後來他慢慢知道原來這幾間公司.
是用一些舊的零件放在新的電腦上,.
以舊換新,以舊聯珍來做生意,.
在這段時間他有一個很大的掙扎,.
他自己也問自己要不要這樣做,.
但他最終決定堅持自己的立場,.
一段很長的時間公司差不多要結業,因為沒有生意,.
但經過一段時間後來發現一件事,.
就是這些大企業發現這些電腦原來有問題,.
然後他們到處去搜尋發現只有弟兄這間公司誠實地賣真電腦,.

$^{41}$所以開始將這些訂單轉移向這位弟兄的公司,.
他後來重新有很好的生意,.
這位弟兄會說,就算神不給我這些生意,.
我仍然要堅持我的選擇,.
在他內心有一個智慧,.
他去選擇行公義公平的事情,.
各位聽眾我們今天看箴言,.
我們常常說到關於智慧的事情,.
箴言第八章,.
智慧用一個第一身,.
好像一個以人式的方式講述自己,.
又呼籲我們去選擇智慧去跟隨它,.
所以你會看到剛才我們讀第八章,.
前面我們沒有讀的,.
是第八章一開始第一節,.
智慧就呼籲我們,.
智慧豈不呼喚,.
聰明豈不揚聲,.
是大聲呼喚我們去聽從它,.
去跟隨它,.
然後說愚無人,你們要學習靈巧,.
愚無人,你們心要明辨,.
不讓你的思想懶惰,.
你們真的要小心去學習,.
你未必馬上會有智慧,.
但是你可以透過一個一個的學習過程,.
在當中得著智慧,.
在過程當中小心明辨,.
分辨什麼是光明,.
什麼是黑暗,.
在當中學懂智慧,.
當你承認你的限制,.
你是愚無人,.
你就會有更大的動機,.
願意去更多的學習去分辨,.
不斷的去領受聰明,.
選擇聰明,.
為什麼要選擇聰明呢?.
因為這個世界上有很多其他的事情,.
是呼喚著你去選擇它,.

$^{81}$在這裡之前的第七章,.
就說到有淫婦,.
我們稍為回顧一下,.
你會發覺在聖經裡面第七章就提到,.
這個邪惡,.
它在不斷的呼喊著,.
浸然用女子淫婦,.
來代表形容邪惡,.
是呼喚人去選擇它,.
所以經文提到,.
一些人被它引誘,.
聖經說,.
好像牛去被宰殺,.
有好像雀鳥急投羅網,.
自投羅網,.
而且急速進入這些羅網裡面,.
但是不知道其實過程當中,.
是賠償自己的生命,.
因為有太多其他世上的事情,.
在不斷呼喚我們跟隨我,.
你得到我,.
你有前途了,.
你得到我,.
你可以出人頭地了,.
有很多不斷的呼喚來到的時候,.
聖經提醒我們去選擇智慧,.
因此來到後面的八章十一節的時候,.
聖經強調一件事情,.
因為智慧比寶石更美,.
一切可喜愛的,.
包括你很歡喜想得到的一些物件,.
一些感情,.
一些成就,.
這些一切一切都不足以與智慧比較,.
智慧是最寶貴的,.
是值得你付上一切代價去尋找去得著的,.
因此來到我們今天讀的一段經文,.
八章十二節就開始說一件事情,.
要能夠有智慧,.
是必須要敬畏神的,.

$^{121}$智慧第一身,.
自己就說,.
我智慧以靈巧為居所,.
又尋得知識和自謀,.
將它的內涵很清楚地說出來,.
智慧是有靈巧,.
是有知識,.
是有自謀,.
另一處聖經的翻譯就說到,.
我有見識,.
見識是說在不同的處境當中,.
衡量很多不同的原則,.
很多不同的價值觀念,.
你知道那個處境,.
明白那個真理,.
才能夠作出合適的決定,.
有知識,.
有健全的判斷力,.
那個心靈是清醒的,.
不會被別人很多的迷惑去影響你,.
你能夠知道是非,.
以致作出有健全健康的判斷,.
以致能夠可以作出光明的一些決定,.
當我們說到這裡的時候,.
大家都會很明白,.
在這個時代我們見到所面對的世界,.
充滿著各種的虛荒,.
有些好像很有道理的,.
小心一聽,.
其實是歪理,.
很多覺得是真實的,.
看清楚,.
假的,.
特別我們今天這個數碼,.
這個AI的年代,.
你有沒有見過,.
有些人物,.
明明是他,.
他說了一句話,.
很真的,.

$^{161}$後來發覺原來是AI製作出來的,.
也是生成的一些物品,.
這個年代我們學會了兩個字,.
叫做心偽,.
虛偽的偽,.
深刻的一些虛假的事情,.
你發覺在你身邊會常常出現,.
這個年代我們更加需要有神給我們的智慧,.
去分辨真假,.
我們不至於受到誘惑,.
就在英國聖經裡面提及到那種智慧的內涵的時候,.
然後接著他立刻在第十三節,.
就去將智慧和敬畏神連繫在一起,.
所以十二節之後,.
十三節聖經就說,.
敬畏耶和華就是恆惡邪惡,.
當然我們知道敬畏神就是聽從神的說話,.
是曾經在其他地方提及的,.
特別在這一年一開始的時候,.
如果大家有印象就是第一章,.
已經開宗明義地說,.
使人領受明智的訓誨,.
這裡說的是智慧,.
就是公義,公平和正直,.
智慧帶來的是一種對神的敬畏,.
或者說敬畏神帶來智慧,.
但是你其中必定包括了公義,公平和正直,.
何時你行在公義,公平當中,.
你見到自己進入了智慧的裡面,.
然後聖經說到,.
敬畏耶和華是知識的開端,.
我們常常讀這句話,.
因此你發覺,.
要能夠得到智慧,.
必須從敬畏神開始,.
如果你沒有敬畏神的心,.
你不會有神的智慧,.
什麼是敬畏神?.
敬畏神就是你會以神為神,.
你知道祂真是神,.

$^{201}$你尊崇祂,.
你聽從祂,.
你會敬畏祂,.
你不會不理祂,.
你不會輕視祂,.
我們廣東話說,.
當祂無道,.
不會,.
因為你真的知道祂是神,.
因此你的心裡面有一個很大的一種動力,.
要去聽從祂的吩咐,.
你也知道祂的屬性,.
祂是慈愛的,.
祂是公義的,.
因此你的生活行為,.
按照祂的屬性來活出來,.
生活,行事,說話,.
按照神的屬性來做,.
就算要付上代價,.
很樂意,.
因為祂是神,.
去聽從祂,.
付上代價去聽從祂,.
因此你知道,.
敬畏神,.
是以神為神,.
是真的敬重祂,.
但是剛才讀的第八章十三節,.
你會發覺聖經從一個比較反面的角度去說敬畏神,.
敬畏也就是,.
恨惡邪惡,.
這個恨惡是一個很嚴重的字眼,.
一種很強烈的抗拒,.
我不要玷污自己,.
我不要邪惡,.
當你讀聖經的時候,.
聖經描述主耶穌的屬性,.
就用了八個字,.
聖經就說,.
我們的主,.

$^{241}$祂喜愛公義,.
恨惡罪惡,.
祂是喜愛公義的,.
祂是恨惡罪惡,.
因此當你想到你是主耶穌的門徒,.
你跟隨祂的時候,.
你同樣需要按照這個方向去過你的生活,.
你就是敬畏神..
我想我們很多時候都會想起舊約聖經裡面的約瑟,.
這位人物他經歷過很多很多的苦難,.
被自己的兄弟出賣,.
被人出賣到一些地方當奴隸,.
面對很多不公不義的對待,.
很困難,.
當他面對著他的主母,.
引誘他,.
要跟他有不正當的關係,.
旁邊的人如果去分析這個處境,.
就覺得你有機會了,.
你有機會可以脫離這種苦難,.
你可以扶搖直上,.
有人看起你,.
你很有前途..
但是當時這位約瑟雖然在苦難當中,.
但是他有一個很清楚內心裡面對神的敬畏,.
你記得他所說,.
「我豈能作者大惡,得罪神?」.
他知道這位神是神,.
他絕對不可以做這些惡事,.
得罪這位天地的主宰..
他一定讓你明白,.
當你的心真正敬畏神的時候,.
你發覺你裡面有一種強大的力量,.
可以救你脫離試探,.
可以幫助你選擇公義,.
選擇良善..
你發覺聖經裡面說,.
如果你不敬畏神,.
你不會得到智慧..
聖經就強調敬畏,.

$^{281}$「我恨惡邪惡,我恨惡驕傲,.
狂妄惡道,和乖謀的口.」.
一路講述出那種敬畏神的生命,.
領受的智慧所產生的那種生命的動力..
接著你會再看的時候,.
發覺後面的經文,.
聖經多一個例子來講述這種智慧..
聖經就講到,.
有一些國家的君王,.
有權勢的人,.
他們能夠可以管治國家,.
是因為他們憑著智慧,.
憑著神給他們的智慧來治國,.
因此帶來了公平..
這些貴族審判官,.
即將智慧能夠可以公義地掌權..
一路講述..
大家都會很明白,.
當我們看回我們今天所身處的世界,.
一些有權位的人,.
如果他不是有神的智慧,.
也都不按照敬畏神的方向來行事和抉擇,.
不懂得去厭惡邪惡,.
你會發覺他們所做的事情,.
帶來了極大的災難..
當你看回很多新聞消息,.
在不同的地方,.
很多的權貴,.
他們不按照這種智慧來處事的時候,.
他們很多的人民都受到很大的災害..
正面的,領受神的智慧,.
你會發覺生命產生一種美善,.
祝福很多身邊的人,.
但相對來說,.
帶來了人很多的災禍..
我們再看,.
在讀落的時候,.
你會發覺聖經裡面特別強調一件事,.
17-21節提到, 愛智慧的得著真實的永恆的豐盛,.
你愛他,你會真正得到豐盛..

$^{321}$聖經17節說, 愛我的,我也愛他..
在這裡我們突然間發覺,.
聖經用一種人際間的關係,.
愛的關係,.
描述到人和智慧的關係..
當然你會明白, 愛不可以單方面的,.
是一種相互的關係,.
你不可以向神求神,.
求你給我智慧, 過我的人生,.
但同時你不愛神,.
或者你不愛得到神給你的智慧,.
那種相互的關係,.
表示著一種互相的毀身,.
互相的接納,.
因此你去愛智慧,.
智慧就是愛你..
但你說, 怎樣才叫做愛智慧呢?.
下面解釋給我們知道,.
懇切尋求我的必尋見,.
懇切去尋求智慧..
這裡所說的懇切, 原來的文字是很強度的字眼,.
是表示著一種很熱切的,.
很努力去尋找智慧..
講述一些人, 當他遺失了某些物品,.
他很早起床,.
一早起床就不斷去尋找,.
直到找到為止..
那種懇切尋求的心,.
在他行動當中表達出來,.
他在當中很想表達出他愛智慧,.
因此他去尋找,.
他去明白, 他去認識,.
他一定知道, 或者我們說到這裡,.
我們都可以問問自己,.
我們有沒有愛神的智慧?.
我們有沒有懇切去尋求智慧?.
如果你沒有, 你不要期望你得到..
然後聖經又說到,.
其實當我們去想愛智慧的時候,.
你放心, 你只要走在公平公義的路上,.

$^{361}$你就尋見智慧..
智慧向我們宣告說,.
我在公義的道上走,.
在公平的路上走,.
亦意味著, 何時你走在公義的路上,.
你就可以尋見智慧..
當你在人生面對著不同的抉擇的時候,.
你要問, 哪一條是公平的路?.
當你決心走在公平的路上,.
你就領受神給你的智慧,.
在這條路上, 走在一個蒙福的路上..
我最近有機會看一位人物的傳記,.
我看的時候很感動,.
這位是南非第一任黑人的民選總統,.
大家可能都認識這位孟德拉先生,.
孟德拉先生本身是一位教徒,.
但是他自己在他的國家裡面, 看到當時的白人統治的社會,.
是一路推展一個叫做種族隔離政策,.
他本身是很反對這個政策,.
因此面對著很多的壓迫,.
他被拉到監牢當中,.
他坐了二十七年牢..
在這段時間,.
他自己公開的承認,.
他說我雖然是教徒,.
但是過去我的信仰好像很遙遠,.
但是在監牢裡面,.
我突然間很深刻的體會禱告的力量,.
在困難當中,.
在逼迫的患難當中,.
我在禱告裡面得到心靈的能力..
很多後來的人寫他傳記的時候講述,.
這種心靈的能力是什麼呢?.
是一種不屈不撓的精神,.
面對著很多的困難,.
很多人對他的傷害,.
但是他仍然堅持他的良善,.
堅持去選擇公義,.
一種不屈不撓的精神,.
很明顯在他生命裡面表達出來..

$^{401}$另外,發覺他生命裡面有一種力量,.
是不被腐毒來控制,.
被人這樣傷害的人,.
很多時候很自然心中出現腐毒,.
但是蒙特拉先生心裡面沒有腐毒,.
從幾件事可以看到,.
當他後來出監之後,.
奇蹟地選上做民選的總統,.
當他進到英國政府,.
見到很多上一代的,.
或者說前任的總統,.
白人總統他們的官員,.
很多人都害怕,.
害怕這個新總統會秋後算帳,.
他們以前曾經怎樣去迫害一些黑人,.
但是這位總統跟他們宣佈,.
你們全部可以留任在新政府裡面..
面對各位的少教,.
這位少教曾經奉命,.
要將蒙特拉所屬的政黨,.
他們的總部炸毀,.
但是蒙特拉知道這件事情,.
但是他跟他說,.
你繼續留任在我的安全部門,.
你繼續做我的主管..
在他後來,.
當他就職典禮的時候,.
他邀請一位人物參與,.
這位人物是監獄的獄長,.
這位獄長曾經在監獄當中,.
是鞭打蒙特拉先生,.
苦待他,.
但是他很願意寬恕他,.
邀請他出席他的就職典禮,.
然後表達你放心,.
我會繼續做你的工作..
這位獄長後來因為深刻感受到,.
蒙特拉先生那種生命的力量,.
他說我很深刻體會,.
他是以德報怨,.

$^{441}$在他生命裡面有一種能力,.
我很想得到..
因此這位獄長後來都信了,.
很多很多事情,.
當我讀的時候,.
我心裡面很深體會,.
其實弟兄他有智慧,.
他有一份從神而來的智慧,.
以至他願意為了敬畏神,.
付上極大的代價,.
也因為他所做的事情,.
他帶給他自己的國民,.
有極大的祝福,.
帶來平安,.
帶來盼望,.
因此我相信,.
今天在這個時代,.
不一定是一些偉大的人物,.
而是每一個在你身旁很多的,.
可能是小人物,.
這些凡夫走卒,.
何時當他們真的敬畏神,.
他們領受神給他們的智慧,.
你發覺在他們的生命當中,.
都可以閃爍出一種光芒,.
將神自己的真實照耀在這個世界上..
我們再看,.
在《夏天的經文》裡面,.
在二十二到三十一節,.
走路講述到神的創造,.
聖經裡面有一個訊息告訴我們知道,.
在神所創造的世界裡面,.
要站立得穩,.
你必須憑藉著神的智慧,.
在神所創造的地方走你的路..
第八章二十七節,.
其中有關於創造的事情,.
我們稍為看一看,.
聖經提到,.
他立高天,.

$^{481}$我在那裡,.
他在冤滅的周圍畫出圓圈,.
上是窮倉堅硬,.
下是冤緣穩固..
我不知道大家有些時候,.
當你去到一些郊野地方,.
或者早上的時候,.
望到很遙遠的一些景色,.
看到地平線,.
看到太陽從遠方慢慢升起,.
你的感覺不知道會怎樣,.
或者你的感覺這麼奇妙,.
為什麼永遠是天有天,.
地有地,.
海有海,.
不會突然間海洋升高了,.
佔了上天,.
我們水浸..
聖經提到,在冤滅水面之上,.
神畫了一個天平線,.
或者地平線,.
畫出一個圓圈,.
上是窮倉堅硬,.
穩定的在天上,.
立定,.
下是冤緣穩固..
一路講述這些情況,.
當你一路讀的時候,.
你發覺聖經裡面在講述,.
為何上帝極其偉大,.
創造奇妙的天地..
但是經文就一直告訴我們一件事情,.
在二十二節裡面提到,.
在這些一切被做以先,.
已經有了智慧,.
我就有了..
然後這些一切被做以先,.
我已被納,.
我已生出,.
英文翻譯更準確的意思是,.

$^{521}$它是被引帶出來,.
帶進我們面前,.
讓我們看到智慧所帶來的美善創造..
當你這樣看的時候,.
你會明白聖經在講一件事情,.
就是智慧是從神而來的,.
真正的智慧是神所賜的,.
智慧是屬於神的,.
是神引帶出來,.
成為了我們生命的祝福..
因此你知道,.
天地萬有,.
是神藉著智慧所創造的,.
而你和我要得到智慧,.
必然是神賜給你的,.
是神而來的,.
你可以從神那裡尋找得著..
然後聖經就提到,.
在三十節裡面講到,.
那時智慧說,.
我在祂旁邊做什麼?.
工程師,.
engineer,.
或者我們中間有做工程的弟兄姊妹..
你留意一下,.
聖經在講一位的智慧,.
祂天天充滿喜樂,.
也意味著祂和神之間那種互相喜悅的關係,.
祂們很開心,很喜悅地一起工作,.
而且祂們緊密地同工,.
聖經特別講到,.
祂在神的旁邊,.
祂成為神的工程師,.
來創造,設計,成全這一切,.
描述著一種很緊密的關係,.
因此你知道,.
我們今天的世界,.
是神藉著祂的智慧所造成的..
再進一步去思想,.
如果這個世界是神用祂的智慧所造成的,.

$^{561}$我們在這個世界生活,.
要知道怎樣行我們的路,.
就必須要憑著創造世界的神的智慧,.
我們才可以行出一條應該有的路,.
得到應該神要給我們的祝福..
或者我有一個簡單的比喻,.
我之前也和大家分享過,.
就是很多年前,.
我自己在家裡買了一個電子牙刷,.
買了回來的時候,.
我從來沒有用過電子牙刷,.
我都沒有怎樣特別用,.
但有一次我自己進去,.
想自己真的試一下用,.
煮一些東西,.
我按了一下按鈕,.
就想著它會一直煮,.
十分鐘之後回來看一看,.
咦,動都不動的,.
沒有反應的,.
不行了,我也有點趕,.
我就真的將牙膽整個抽出來,.
放在火爐上面,.
開火,燒一燒快一點,.
幫一幫它,.
不過真的很快,.
很快它就裂了,.
啪一聲就裂開了,.
嘩,我很害怕,.
為什麼突然間會這樣反應,.
叫太太來看看,.
太太一看到,.
你有沒有看說明書呀,.
我真的沒有看,.
也真的不知道不可以這樣做,.
那次很深刻,.
我突然間發覺,.
如果你要運用一件物品,.
你要去明白它的運作,.
你要看說明書,.

$^{601}$你才知道它的製造過程,.
你才知道它的運作方式是怎樣,.
各位聽眾,你們要明白,.
我們生活在神所創造的天地,.
如果你沒有神給你的智慧,.
你怎樣能夠走一條光明的路,.
你是神所創造的,.
如果你不是按照神的計劃,.
你怎樣能夠活出那種神定義給你的美事,.
但是說明書在哪裡,.
我怎樣知道去看,去明白,.
大家一定不會記得,.
《聖經》裡面詩篇十九篇,.
這篇詩篇我們常常會讀,.
《聖經》裡面詩篇十九篇,.
我們剛開始的時候已經讀過《聖經》文,.
就是講到神的創造,.
但是裡面你發覺在講一件事情,.
這日到那日發出言語,.
這夜到那夜傳出知識,.
你留心整個天地宇宙窮滄,.
你會認識神多一點,.
你會看到那種奇妙的智慧,.
那種美善,.
這個天地當神創造的時候,.
帶著很多訊息向我們講,.
然後《聖經》去到第七節的時候,.
接著就講到神的道,.
神的律法使人蘇醒,.
使如蒙人有智慧,.
神賜下道成肉身的基督,.
也賜下成文字的真道,.
這些真道當你去讀去看的時候,.
你會發覺領受智慧,.
各位弟兄姊妹,.
我們很多時候神像給予鼓勵,.
你們要小心讀《聖經》,.
要好好的靈修,.
不是說讀多一點《聖經》,.
神就喜歡我們多一點,.

$^{641}$我們就蒙福多一點,.
不是,.
是我們真的需要神的話,.
如果沒有說明書,.
我們怎知道在這個世上生活,.
怎知道神的旨意,.
怎能夠分辨是非,.
在過程裡面,.
我們需要神的道,.
神就賜給我們道,.
也如剛才所講,.
在我們身上的時候,.
你知道道成肉身來到我們中間,.
活著的道在我們中間,.
將真理向我們說明,.
也叫我們從祂的生命裡面,.
從祂的救恩裡面得著智慧,.
因此你知道整個自然界,.
神的道和基督永恆的道,.
這份說明書不斷向我們說話,.
你要去尋找去明白,.
因此你知道如何在這個世代,.
能夠活出那種神定義給我們的豐盛的生命,.
當我們讀到這段經文的時候,.
你會發覺到最後,.
去到第三十節的時候,.
突然間講一件事情,.
剛才我們讀到,.
智慧在神的旁邊為工程師,.
他們很緊密地,同心地去做這件事情,.
但是最後一句話,.
聖經就說,.
在他的全地歡笑,喜愛住在人世間,.
這裡不單單講他們那種喜樂的心情,.
更加表達著智慧住在人間,.
也都意味著上帝是讓他的智慧臨到我們,.
他不是很遙遠,在你身旁,.
只要你願意去尋求他,.
你會尋得見,.
他是很喜樂地在我們中間,.

$^{681}$更加是喜樂我們能夠尋求他,.
他喜悅我們愛他,.
以致他可以愛我們,.
你發覺聖經一直講述神自己的智慧的時候,.
讓我們知道,你們尋求我,.
你們放心,歡欣喜氣地尋求我,.
因為我於你們極其重要,.
我都非常渴望可以在你的生命中,.
所以當我們看到的時候,.
就會回想八章一節,.
為什麼聖經就講,.
智慧在呼喊,在揚聲,.
你們來吧,你們尋找我吧,.
各位弟兄姊妹,.
願我們都回應這個呼喚,.
我們去尋找他,.
注著敬畏的心去愛慕神的話,.
把握著你的日子去尋求智慧,.
以致你經歷到神給你的滿足,.
也可以得著人生的智慧,.
知道走當走的路,.
我們一起祈禱吧,.
請你為自己祈禱,.
天地的主,.
人算什麼,你竟然顧念我們,.
主我們都是很多的無知,.
很多的限制,.
我們以為聰明,.
但其實很多事我們都不知道,.
但你竟然樂意將你的智慧賜給我們,.
你竟然將你永恆的道賜給我們,.
主求你感動我們的心靈,.
讓我們能夠愛慕你的道,.
我們渴望能夠得著你的智慧,.
主求你救我們脫離這個世界,.
很多很多的假智慧,.
那些假的說明書,.
讓我們心裡面能夠轉向你,.
去尋求你,.
讓我們在敬畏你的過程中,.

$^{721}$能夠行在智慧的路上,.
主在我們心中的禱告,.
求你垂聽,.
奉耶穌的聖名祈禱,.
阿們..
\newpage



\section{約翰福音 13:34-35}
\label{sec:031797}
\textbf{彼此相愛的標準與實踐}~~~~林誠信牧師
\newline
\newline
連結: \href{http://www.church.com.hk/acms/content.asp?site=cdc\&op=show\&type=product\&code=031797\&layout=sermon}{\texttt{031797}} ~~~~ 語音日期: 2026-03-08
\newline
\newline
\hyperref[sec:031796]{\small{< < < PREV SERMON < < <}}
~
\hyperref[sec:index_chronic]{\small{[返順時目]}}
~
\hyperref[sec:031798]{\small{> > > NEXT SERMON > > >}}
\newline
\newline
約翰福音 13:34-35
\newline
\begin{longtable}{cl}
\hline
\hline
章節 & 經文 (和合本修訂版)\\
\hline
13:34 & \begin{tabularx}{0.7\textwidth}{X} 我賜給你們一條新命令，乃是叫你們彼此相愛；我怎樣愛你們，你們也要怎樣彼此相愛。 \end{tabularx} \\ \\ \relax
13:35 & \begin{tabularx}{0.7\textwidth}{X} 你們若彼此相愛，眾人因此就認出你們是我的門徒了。」 \end{tabularx} \\ \\
[1ex]
\hline
\hline
\end{longtable}
$^{1}$鼎姊妹 午安.
差不多中午了.
今天很開心.
很感恩為窩福堂.
踏入50周年.
因為.
去年朱牧師.
跟我分享.
他將會退休.
亦都.
告訴我今年你們是.
50周年.
每一個月他會請.
窩福堂的仔仔女女.
回來.
跟大家分享.
我們在神的話語彼此.
互相勉勵.
我今天在這裡.
代表我們廣福堂所有的.
鼎姊妹向你們問安.
願上帝賜福你們.
繼續使用你們.
亦都讓.
李牧師.
和整個教育團隊.
執事們和鼎姊妹.
同心合意.
希望福音.
和建立主所喜悅的.
門徒.
今天我.
感受到.
上帝在.
播道會這個大家庭.
特別在窩福堂裡.
做了很多奇妙的工作.
不知道你們知不知道.
你們有11個仔女之外.
有多少個孫子.

$^{41}$你們可能不知道.
不過我也不知道.
我只能夠數.
原來窩福堂.
生了廣福堂在1984年.
而廣福堂.
亦都生了.
同福堂 尖福堂.
同福堂有了.
很多其他的.
仔仔女女 所以那些是你們的.
曾孫.
尖福堂亦都有全福堂.
所以亦都是.
你們的曾孫.
所以我相信.
不止11個仔女.
亦都不止.
11個孫子 可能很多.
在孫子以外.
已經有曾孫.
這個.
屬靈的恩典.
和福氣.
不是叫我們只是.
獨有 祂是要我們.
與人分享.
讓我們的社區 我們的城市.
甚至我們的國家 直到地極.
都能夠充滿著.
上帝的指紋 祂的百姓.
是將救恩.
耶穌基督能夠賜人.
永生真正.
那個盼望的緣.
緣由 與人講得.
清清楚楚 求主繼續.
恩代我們 40年前.
我們看到.
加里遜牧師.

$^{81}$在窩福堂.
當時他就是.
看到一個異象.
他看到一班在海外.
留學回來的弟兄姊妹.
他們在.
中環那邊上班.
他就挑戰他們.
你們有沒有這個負擔.
願意出去 在中環那邊.
開展一個.
崇拜和.
課安的事工.
談何容易呢 如果一開始.
在中環那些地方.
你去哪裡找一個地方來崇拜呢.
那時是1984年.
開了一個.
香港教會歷史的先河.
就是.
我們租用.
酒店的球室.
因為星期日早上.
通常酒店的球室.
是晚上有些宴會.
那麼早.
九點.
十點 十一點.
這些時間 人家不會在球室有什麼功能.
很感恩.
我們結果.
就是84年開展.
港福堂.
由起初幾十人.
到現在有.
接近二千人的一個會眾.
這一切都是神的很多恩典.
我們在十年前.
也開展兩文三語.
的一個服事.

$^{121}$有一個英語的崇拜.
有一個國語的崇拜.
他們的會眾都是各自.
大概二百人左右.
我們也有一個.
清少的崇拜.
再加上我們四堂的語語崇拜.
總共有七堂的崇拜.
在每個週末進行.
神給我們互相鼓勵.
一起在神的角度裡面.
雖然我們身處.
不同的城市的地方.
但我們是有.
同一的使命.
我們都要將福音.
真是很好的與人去分享.
主席也帶過我們.
讀出這一段經文.
《約翰福音》十三章三十四至三十五節.
我相信大家都.
很耳熟能詳.
雖然這兩節聖經.
這麼短.
但裡面所包含的意義.
是很豐富的.
我們從約翰福音.
在聖經的前文裡.
我們會看見.
是主慧門和洗腳.
並預言祂會被人出賣.
然後是由.
十三章到十六章.
主耶穌是親自.
跟我們講最後的一個.
教導.
提醒他們.
他們必須來.
活出這一種.
新的命令.

$^{161}$這個新的命令.
等下我會詳細再和大家講.
但在十三章三十六節到十四章六節.
那裡講到.
真正的生命之道.
耶穌就是我們.
每一個人都需要的.
生命的糧.
我們有了祂的說話.
我們生命才能夠真正的滿足.
我們的喜樂.
平安 盼望的源頭.
都是因為基督耶穌.
祂的說話和應許.
所以當他講到十五章的時候.
十四到十六章就提到.
主耶穌說.
不是丟棄我們.
孤兒寡婦.
祂要透過聖靈住在我們裡面.
雖然祂要升天.
祂要為我們受死.
然後復活升天.
但祂不是丟棄我們.
祂差遣聖靈住在每一個.
神的兒女的心中.
這個關鍵很重要.
這個關鍵正正是.
我們今天能否.
活出神要我們彼此相愛.
這種原動力.
在十五章裡面.
也提到我們要.
住在主的裡面.
這個住這個字.
就不是像旅行一樣.
我在酒店住幾天.
那個地方住幾天.
這個住是.
永遠地.

$^{201}$住在主的裡面.
我們才能夠.
多結果子.
我們是恆常每一天.
你都需要與耶穌結連.
因為祂就是真正的葡萄樹.
而十七章就提到.
主是我們的大祭司.
祂在臨別.
為門徒祈禱的一個重點.
就是要為他們合而為一.
來祈禱.
我們今天可以彼此相愛.
才能夠合一.
才能夠達至上帝賦予.
地上教會的大使命.
這個心意.
所以我們說新命令.
是彼此相愛的「新」.
這個含義是指什麼呢.
首先我們從.
這個新命令.
如果在.
探究原文聖經.
那個希臘文是在說.
Kainos這個字.
不是說時間上的新和舊.
在舊約聖經.
有提醒我們.
我們是有.
這個神給過我們.
這個十誡.
但十誡是刻在石板上.
去到.
這個以西結書.
38章說到.
神要賜給我們.
祂的靈住在我們裡面.
將祂的律法.
刻在我們肉心上.

$^{241}$神的靈要住在.
祂的百姓.
祂的子民的心中.
所以我們今天.
不只是週末才在神的殿.
去敬拜.
我們是七天都敬拜神的.
我們每天無論你是出差.
無論你去到任何地方.
你每天都需要.
去親近神.
因為神的靈 神的電.
就是在我們的內心.
耶穌基督就是那個道成肉身.
滄風腐腐的.
滿有恩典.
和祂的真理.
是住在我們裡面.
這個Tabernacle這個字.
就是一個這樣的意思.
而我們將來去到新天新地.
也是耶穌要.
和我們永遠在一起.
而天堂這個地方.
就是這個致勝所.
是那個真正.
致勝所的實現.
神讓我們今天看到.
這個不是新和舊的分別.
是在性質上.
上帝挑戰.
祂自己的百姓子民.
我們需要.
有這個新的體會.
耶穌說.
我怎樣愛你們.
你們也要怎樣相愛.
我們看到.
舊約有教導我們.
要盡心盡性盡意盡力.

$^{281}$愛我們的神.
而且我們要.
愛人如己.
這個申命記錄章.
或者利未記十九章.
都有的教導.
但這條命令不只是在舊約.
去到新約也一樣.
耶穌回應.
馬克思的提問.
最大的界明是什麼.
就是你要很好的.
全心全意的去愛神.
但不是停在那裡.
因著愛神的緣故.
叫我們能夠去愛別人.
而耶穌就是.
顯色給我們看見.
這個門徒相愛新的要求.
我們可以在.
今天講章餘下的時間.
我希望透過耶穌的生命.
和我們展示.
我們應該如何學習.
彼此相愛.
我們再.
首先來看看.
一位很出名的.
已故.
在上世紀.
是英國的一個大文豪.
就是魯易斯.
他也是一個哲學家.
但他本身原初是一個無神論者.
但後來因為他認識了主耶穌.
他發現.
這是他的生命的.
自保.
他願意透過這本著作.
四種愛.

$^{321}$和我們說什麼是真正的愛.
我嘗試引述.
C.S. Lewis 說的.
頭這三種愛.
都是上帝賦予人類的.
我們可以享用的.
但是我們不要濫用.
這三種.
頭這三種的愛.
都需要規範在.
第四種愛.
就是上帝的.
一個最高最大的.
一個.
美好的.
對我們犧牲的愛.
親情的愛.
首先第一個字.
用Storge這個字.
他是在說.
我們父母和兒女.
或者是我們兄弟姊妹.
這種有血緣關係的.
這種親情.
我們在《西洋聖經》裡面.
看見這個字.
他只是出現.
用母親情.
這個字眼出現.
在《羅馬書》一章三十一節.
或者《提摩太后書》三章三節.
裡面說到.
在末世的時候.
人是專顧自己.
不顧別人.
我們個個都變成很自私.
自我,自大.
在我們的感情.
關係的.
理解裡面.

$^{361}$就算血緣關係.
都變成母親情.
這是一個很嚴重的.
提醒.
我們必須要看到.
今天這個世代正是.
什麼都是說我.
自己的利益.
是先行的.
為什麼這個世界這麼亂.
就是因為當某些國家.
走出來說.
國家的利益是先行的.
我喜歡你國家.
有石油.
我要你給你不給我就打你.
有些無法無天的事.
竟然在我們.
今天這個年代是出現的.
主耶穌回來的日子.
更加近,更加逼近.
頂尖我們要更謹誠.
我們要學習另外第二樣.
就是友誼的愛.
朋友這種.
這種的愛.
Galileo的字是從這個字根裡面.
出來的.
耶穌在約翰15章說.
我從此不再稱呼你們作僕人.
你們是我的.
朋友.
耶穌不將父神的心意隱藏.
祂要向他的門徒.
解說上帝.
的心意.
要他們可以和祂.
傾心吐意來祈禱.
這個都是我們今天.
看到的舊約.

$^{401}$他們和約拿大.
他們那種很純粹的.
同性之間的感情.
那種的關係.
是在這種友誼的愛裡面的.
而第三種.
是情慾的愛.
是說男女之間.
談戀愛的時候.
亦都說到是夫妻之間.
那種的.
情慾的愛.
希臘文就是Eros這個字.
這個Eros其實是.
針對到是一種很親密的.
兩性之間.
那種的關係.
這種關係就只有.
在聖經裡面說.
是一夫一妻男與女.
這種.
交合那種性愛的.
關係而有的.
但他們這種的情操.
不只是說性愛而已.
是說兩者.
他們是有很親密的彼此.
的認識.
這三種都必須.
要臣服在第四種.
就是上帝那種.
Agape Love.
耶穌說我們犧牲在十字架上.
顯明上帝怎樣愛世人.
不只是.
愛某一個民族.
耶穌是愛這個列國萬族.
祂最後的心意.
是很想在每一個民族裡面.
都朝聚有人.

$^{441}$因著信服祂的兒子.
耶穌.
祂在十字架的拯救.
而你也成為神的兒女.
將來我們在天下.
是各族各方各民.
都有擇大眾讚.
到時我們不只是會唱.
廣東話的詩歌.
或者是國語詩歌.
或者是英語詩歌.
我試過被人問.
我們到時唱什麼文呢?.
你說呢?.
我們到時是唱摩西的新歌.
到時是Universal的.
我們是共通的.
在神的靈那種.
愛的交通和朝聚之下.
一切都更新了.
我們今天.
在地上.
因為太多人的罪性和自我.
我們彼此會.
嫉妒,紛爭.
甚至有種族歧視等等.
但到將來我們在天下.
什麼的膚色.
什麼的民族.
我們都在耶穌基督的救贖裡面.
我們有份的話.
我們就是神的兒女.
我們一起去敬拜和事奉神.
的漢奸.
我們朝著Agape Love.
我們去學習.
學習的標準是.
耶穌說我怎樣愛你們.
那耶穌怎樣愛我們呢?.
耶穌愛我們.

$^{481}$是為我們犧牲.
完全.
不介意任何代價.
這一種.
是毫無.
吝嗇.
傾盡祂的生命.
是為我們捨己.
我們今天都需要學習.
從神而來的那種愛.
因為神就是愛.
如果我們沒有這份愛心.
我們沒有了愛的源頭.
第二方面我們有.
愛的源泉.
這個源泉就是.
若有彼此相愛.
是以神的愛作標準的話.
原來我們就真正.
愛父的心.
在我們裡面得以.
源泉.
這個源泉不是說我們變了像上帝一樣.
那種源泉.
這個源泉在希臘文.
Telios那個字.
是在說成熟.
一種真正的成熟的愛.
是用上帝.
為我們犧牲.
而我們的效法主.
對我們這一種的愛.
我們都願意這樣.
來為別人.
為我們的兒女,為我們的朋友.
甚至為我們的仇敵.
你都願意這樣去愛.
當人不了解你.
當人甚至是誤解你.
甚至你明天上班的時候.

$^{521}$可能你的頂頭上司.
或者你的同輩的同事.
是你的仇敵.
你怎樣去化解這些呢?.
你求主給你.
有這一種成熟.
源泉的這一種愛.
去愛他們.
為他們祈禱.
先為自己祈禱.
求主叫你不要生氣.
當別人對你這一種.
誣衊,誣禮.
或者是冒犯.
求主給你有這一種很大的.
勇氣和力量.
你選擇不是報復.
我用基督的愛去愛他.
他未信主的.
你挽回他去信主.
如果他.
是基督徒的話.
你更加要將他挽回.
挽回什麼意思呢?.
因為他跌倒了.
他在軟弱裡面.
神的兒女都會成為別人的絆腳石.
求主幫助我們.
都有這種體現.
第三,這一種來自神的愛.
一定是可以.
結出果子的.
這種果子.
是叫我們.
未信的人看到.
我們這一群神的兒女.
怎樣彼此相愛.
他們認出我們是他們的門徒.
這是一個見證.
但我們要做得到這件事.

$^{561}$你必須.
常住在主的裡面.
主耶穌.
說的這一種住.
不是旅行的那種住.
這是永遠.
住在祂的裡面.
換言之.
每天,每個場景.
每一個時刻.
原來我們都在.
基督裡面被聖靈.
去管理著.
當我們什麼時候遇到.
試探,我們想跌倒.
我們想犯罪.
你馬上留心.
原來聖靈在我們內心.
提醒我們.
不要做這件事.
不要有.
這些不正確的.
你內心的.
不潔的思想或雜質.
如果我們事奉裡面.
都是能夠有一個清潔的良心.
我們和弟兄姊妹就能夠.
更加彼此相愛.
這樣去服侍.
其實一個教會的大家庭.
要彼此相愛,合一去服侍.
有時為何不容易呢?.
因為.
是各人看別人比自己強.
是各人都看自己.
很強.
各人都覺得自己很強.
如果你叫師班.
如果每個師班人.
都覺得自己很強.

$^{601}$為何不找我獨唱?.
每個都要唱到突哂音.
那就很不和諧.
剛才大家聽.
我們敬拜隊唱得好不好聽?.
很多人都點頭.
是不是?.
真的.
一般教會的敬拜隊.
通常是差不多的.
由三個到五個人.
我們的師琴哥哥.
真的彈得很好.
這一班弟兄姊妹.
他們能夠.
這樣去火熱.
去為教會.
忠心地.
每個星期.
一起去敬拜.
透過現代的詩歌.
透過傳統的詩歌.
我們都能夠敬拜神.
所以我們講福堂.
和你們一樣.
我們沒有分傳統.
還是現代.
我們說上帝.
他聽得明奇異恩典.
但他同樣聽得明.
我們要.
剛才那首歌叫做.
活出愛.
很好的.
很有動感的詩歌.
上帝也聽得明.
非洲來的弟兄姊妹.
我們曾經請.
他們一班的兒童師班.
他們上到台上.

$^{641}$是擊鼓.
跳舞.
而且他們唱的時候.
是真的跳起來.
跳起來兩三尺高.
我自己差點擔心台會不會倒.
不會的.
傳統的熱情.
投入.
對於我們比較傳統的教會.
我們說這樣行不行.
其實沒有說不行.
上帝.
他要我們敬拜的原則只是一樣.
你要用心靈.
誠實.
按真理去敬拜親近神.
至於其他的.
只不過是我們那種形式.
音樂有很多不同的風格.
最重要是那些歌詞.
那些歌詞是什麼呢.
大家唱的時候.
唱得明明白白.
用心去唱.
神就閱立了.
祝福大家繼續在前行的日子.
未來的.
如果主還沒有回來.
未來的五十年.
神繼續大大加力使用你們.
彼此相愛標準我們剛剛看了.
這三方面.
都是我們要好好地一起學習.
我們怎樣實踐呢.
耶穌說.
我怎樣愛你們.
你們也要怎樣相愛.
耶穌怎樣愛他的門徒呢.
我們首先在.

$^{681}$最後的晚餐裡面可以看見.
沒有一個門徒懂得.
為別人去服事.
個個都覺得自己.
很高很大.
去到耶路撒冷.
最後這程路.
他們還在爭論什麼.
爭論誰是大佬.
爭論將來.
耶穌你德國降臨的時候.
我要做你的左右.
丞相你是我們的王.
還找了媽媽出來.
我們看到.
這些情景.
我們就知道門徒那種.
事奉的心態.
其實對主耶穌.
一點都不認識.
他們只是在尋求.
自己事奉的那種虛榮.
和未來我要得的.
這些所謂的.
位置.
什麼時候我們的內心不清潔的話.
你的服事是有問題的.
但耶穌.
讓我們看見那幅圖畫.
他自己親自為門徒.
去洗腳.
彼得說.
你是我們的夫子.
為什麼你會做這件事.
但耶穌要彼得.
知道你們要謙卑服事.
你們要被我.
潔淨.
你也要好好的.
用我這種謙卑服事的生命.

$^{721}$作僕人的領袖.
誰願為大.
誰就要去作別人的.
一個服事.
去服事別人.
今天我們需要在神的面前.
也同樣有這種學習.
我在幾年前.
我們教會的一個受難節.
崇拜裡面.
我很感動.
當他們有一個環節.
他們的施班指揮.
和我們聖學部的部幕.
一個傳道人.
他們親自為.
施班員去洗腳.
除了他們的鞋.
除了他們的襪.
去幫他們洗腳.
這件事如果你沒有做過.
你會覺得.
我都沒有.
幫家裡.
自己的太太洗腳.
你要學習.
彼此洗腳是一種經歷.
你的經歷就是說.
原來.
我們好像.
不再是.
我是你的老師.
或者我是你的.
上司.
原來我們都在神的眼中.
我們是弟兄姊妹.
我們一起謙卑地互相服事.
作僕人的領袖.
去為別人.
放下你自己的身段.

$^{761}$第二方面.
要書赦免.
這方面也是很重要.
耶穌說你要愛你的仇敵.
為逼迫你的禱告.
其實我們要愛.
一些對我們好的人.
一點都不難.
但是愛那個.
經常腰心腰肺.
跟你一說.
他就對著你開刀.
衝撞你.
如果你做父母.
你的子女經常這樣.
你會覺得很難受.
但你要為祈禱.
不要放棄.
今天有時.
我妹妹已經三十七年.
我自己看到.
在家庭裡面.
其實是.
父母要用我們的榜樣.
用我們的見證.
愛.
去饒恕.
愛有很多眼淚.
有很多忍耐.
但是你要付出.
你要看到真正的盼望.
就是神的心意.
要我們去包容.
去接納.
我曾經見證過一些青少年.
他在他的成長期.
反叛.
父母很頭痕.
怎麼辦呢?.
就算在教會長大的小孩子.

$^{801}$上了兒童主學.
為什麼到中學的時候.
他就開始不肯回家呢?.
他不肯回家其實一定有原因.
但有一件事.
小孩子很清楚.
哪個叔叔 哪個阿姨.
是愛他 是疼他.
他會知道.
你有愛的時候.
就算在那一兩年間他失去了.
當那些叔叔阿姨.
去找回他們的時候.
為他們祈禱.
跟他們一起同行.
他知道.
他們找到愛的源頭.
是來自神.
為什麼我這麼頑皮.
甚至在學校犯了事.
但是叔叔阿姨.
仍然不離不棄.
要挽回你呢?.
這是用神的愛.
去寬恕.
去接納.
就好像在舊約聖經裡面.
約瑟.
他被他的親兄長出賣.
去到波提佛的家.
本應.
他做了大內總管.
誰知道主母看中了約瑟.
覺得他長得很帥.
臉色紅潤.
天天在勾引他.
但約瑟就死都不就範.
結果有一天他硬來.
主母捉住約瑟.
要他跟他上床.

$^{841}$結果.
把他的外婆一扯.
約瑟寧願.
逃跑.
我不要陷自己.
在這些試探.
和得罪.
不只是上帝.
而是得罪我的主人.
約瑟是一個很敬畏神的人.
所以他能夠在這種.
試探下當機立斷.
他就來逃跑.
結果.
反而被人冤枉.
證據確鑿.
主人一回來.
波提弗的老婆跟他說.
你那個僕人約瑟.
他這樣.
調戲我.
結果波提弗就.
把約瑟下監.
他含冤莫白.
來收監.
但約瑟仍在監獄裡.
得到那些獄長.
和後來.
有法老兩個僕人.
善長和酒正.
都下來了.
但他為他們解夢.
解完夢一個官服完職的酒正.
竟然忘記了約瑟.
約瑟被人忘記.
上帝永遠.
都沒有忘記約瑟.
上帝使用約瑟.
結果後來當法老王.
再發另一個夢.

$^{881}$約瑟為他解夢.
他將約瑟提拔.
他說這個不是普通一個人.
有神的靈住在約瑟的心中.
法老一個不認識神的人.
但在約瑟的見證.
在他生命裡.
看見神.
那種的恩典.
結果約瑟後來成為.
埃及的宰相.
到迦南地.
患大饑荒的時候.
我們看到雅各猜念他的兄長.
來到遞糧.
結果最後.
聖經是這樣說.
約瑟接了他爸爸來.
跟他的兄長相認了之後.
要讓他們知道.
我已經饒恕你們.
不過當雅各一死.
在創世記第五十章.
那群兄長馬上害怕.
他說.
爸爸在這裡的時候.
可能約瑟留下一手.
現在爸爸走了.
會不會你現在就出手了.
殺光我們了.
但約瑟怎麼說.
約瑟說你們不要怕.
從前你們的意思是要害我.
但神的意思.
原是好的.
要保存很多人的性命.
他說你們不要怕.
我要養活你們和你們的婦人孩子.
原來在約瑟這麼苦難.
這麼多坎坷的.

$^{921}$被人出賣.
被自己親人.
凋棄的一種日子裡.
上帝沒有忘記約瑟.
上帝在他最痛苦的人生階段裡.
他繼續造就他.
因為約瑟敬畏神.
而他能夠活出.
上帝這種赦免別人的.
這種愛.
是來自神的那種恩典.
我在這裡和大家分享一個.
也是很感人的見證.
這是發生在二十幾年前.
當我還在加拿大.
牧養教會的時候.
有一個星期五晚上.
我突然收到.
教會一對夫婦打電話給我.
林牧師.
你明天早上可不可以.
陪我們一起.
去接待一個太太.
載她.
去一個婦人.
那裡去認屍.
我說.
我真的沒試過認屍.
處理安息禮拜很多次.
那次是我第一次.
帶人去認屍.
我說這是什麼事.
原來這個太太的.
女兒叫Anna.
Anna是一個姊妹.
是一個信耶穌.
信得很好.
而且讀書很優秀.
她每一科都是State A.
她去到那個大學.

$^{961}$她就是.
準備讀醫科.
將來她要成為.
一個醫療的宣教士.
所以她在大學已經很熱心.
去傳福音.
原來當天發生什麼事呢.
在10月的一個.
很平常的禮拜五.
下午.
他們放學後.
他們約了幾個同學.
一起開車回到.
多倫多.
多倫多是一個大城市.
多倫多那晚.
有一個報道會.
他們請那些未信的人.
來那個地方.
聽報道會.
他們租了一輛車.
而開車的司機.
是一個留學生.
他平常沒有開車.
但這次因為要出來.
所以租了一輛車去開.
可能是因為他開車不是很熟.
平常不太會開車.
特別是在高速公路.
他突然扭tie .
他想cut線.
但扭tie 扭得太大力.
結果整輛車就spin.
就打轉.
轉了出去那個highway.
外國那些.
如果你有住過.
一般那些highway.
都有一些shoulder link.
很闊的.

$^{1001}$讓那些死火的車停在那裡.
等救援.
我第二天去和死者的媽媽.
Anna的媽媽.
和教會一對夫婦.
去到肇事現場.
我觀察過.
看到那些車痕.
那四條輪的車痕.
在地上.
本來應該沒事的.
但很奇怪.
剛剛在那條路有四棵大樹.
就撞了在那個shoulder link.
結果那輛車撞過去.
爛腰撞過去.
是一個高速底下.
很大的一個撞擊力.
不過因為他們四個人.
都是坐在車上.
有安全帶.
所以沒有人拋出車外.
但唯獨.
沒有人慢慢爬出來.
等十字車.
唯獨Anna姊妹在司機位旁.
她垂下了頭.
當刻.
她的脊椎骨.
在那個很大的衝力下.
一撞.
斷裂.
她當場死了.
她回到主耶穌懷抱.
當然這個媽媽.
很傷心.
我接這個媽媽.
去認屍的期間.
她當然會流淚.
但她跟我說.

$^{1041}$她知道將來.
她要在天家.
跟她的女兒Anna再相見.
我幫她處理.
遺體.
和之後的安息禮拜.
幫她打點了.
然後.
她跟我說.
林牧師.
去探望受傷的.
那幾個留學生.
我知道.
他們只是皮外傷.
他們應該沒什麼事.
他們都不在醫院.
仍然想去探望他們.
我就照著她的意思.
去找那三個留學生.
一打開門.
這個肇事的司機.
他走到Anna的.
媽媽的腳前.
抱著她的腳.
然後滔滔痛哭.
大聲喊.
阿姨阿姨.
我對不起你.
是她的一個錯失.
導致Anna姐妹.
離開人世間.
但更加感動.
令我很感動.
就是Anna的媽媽.
對著肇事的司機.
她說.
我知道將來在天家.
會有我的女兒.
但我知道這是神的恩典.
是神拯救了她.

$^{1081}$也拯救了我.
我就是要告訴你.
是因為神的愛.
我同樣要親口告訴你.
我饒恕你.
我選擇饒恕你.
這句說話容易說嗎?.
我相信絕對不容易說.
當她這樣說的時候.
那位留學生就哭得更大聲.
然後我們全家都流淚.
我請他們起身.
我們一起手牽手.
我們一起祈禱.
我們一起將Anna姐妹的安息禮拜.
要怎樣去準備.
我們交託給神.
那一刻我感受到.
聖靈在整個房間.
充滿著.
我們每一個人的心靈.
我們的內心.
充滿著神的愛.
而這位.
馬來西亞的.
男生.
他就在這個安息禮拜.
他決志信主.
他信了主之後.
他學成之後.
他回到馬來西亞.
我回看這件事.
已經隔了二十多年.
但我看到.
原來一個生命.
我們會說這麼好的一個生命.
為什麼神這麼早就接去.
但上帝有時.
所做的事.
不是說我們生命的長短.

$^{1121}$是你今天在你短暫的人生.
你有沒有捕捉.
和把握神給你的機會.
你去服侍神.
還有你把握這些機會.
去帶人信主.
這些全部機會.
不一定是時時有.
不是時時有.
我稍後.
會再和大家分享.
另外一些見證.
恆久忍耐是第三方面.
彼得乘以為.
自己可以為主.
去坐牢.
甚至犧牲.
他都覺得沒有問題.
但耶穌說:西門!西門!撒旦要得著你們.
好西你們.
好像西墨寺一樣.
但我已經為你們祈禱.
失去信心.
你回頭後要堅固.
你的弟兄.
彼得是有.
一個受託.
但他不知道自己.
以為是靠自己.
去服侍.
其實我們沒有一個人可以靠自己.
你是多聰明.
多有才智.
多多恩賜都好.
你絕對絕對不可能靠自己.
去事奉神.
如果你的內心.
往往都是.
以這個做你的出發點.
我很肯定.

$^{1161}$你的事奉遲或早.
不單止燃燒.
而且你一定不能夠榮耀神.
我們今天在神的面前.
神給了我們很多恩典機會.
是要他忍耐.
耶穌怎樣忍耐.
來對付彼得呢?.
耶穌復活之後.
都要找彼得出來.
外面說你愛我比自己的心更深嗎?.
都要問他三次.
問三次.
有些人說是不是因為彼得否認耶穌.
三次所以耶穌問足他三次.
我告訴大家.
耶穌真的不是這麼小器的人.
耶穌的心不知道多寬宏.
和忍耐.
多大.
耶穌是要確保彼得.
你如果要牧養我的羊.
你必須用我的愛.
去愛牠們.
今天我們不是打份工.
去作傳道人.
我經常提醒我的傳道同工.
你今天.
教會每一個弟兄姊妹.
是你的屬靈兒女一樣.
他能夠進來教會.
未信主你也要信主.
他如果已經信主你也要栽培他.
要成為.
神喜悅的門徒.
我們有上班的時間.
但是.
下班後你一定會被淘汰.
我從來都不會.
下班後把手提電話關掉.

$^{1201}$我的手提電話.
24小時都開著.
不過.
教會的弟兄姊妹很疼愛我.
他們不會說睡不著.
打電話給我.
牧師你今晚還沒睡吧.
你跟我一起.
聊聊天好不好.
我當然告訴他們.
你是晚上半夜.
真的有事.
如果你有我的電話.
你必須要我在這.
譬如有你的親人在ICU.
或者有什麼事.
我都願意來.
如果這是真的需要.
我試過.
探一些睡危的病人.
我就留下電話.
醫生說他.
不久人逝.
我說如果真的萬一你需要的話.
而他是半夜離開.
人家.
他沒得選擇.
他不是遷就你的方便.
是你要去.
為了他.
作出服事.
犧牲等等.
其實我昨晚就是這樣.
昨晚我和李牧師.
完了晚課.
我和李牧師和另外.
有位你們的行政.
的傳道.
我們一起吃飯盒.
一起聊天.

$^{1241}$昨晚聊到九點多.
他說李牧師.
第二天早上講到.
我們可以再聊.
但是.
當時我因為.
崇拜和和他們.
吃飯聊天.
我的電話是沉默了.
是靜音.
我一離開了光復堂.
我打開電話.
打開有聲音的時候.
原來有留言.
有聲音訊息.
原來有一個長者.
在廣義醫院.
一個ICU裡面.
等著我.
牧師你可不可以盡快來.
醫生說他現在突然惡化了.
我沒見過.
那個是我們教會的弟兄.
介紹我去見.
原來就是這個弟兄.
很多年工作的.
上司和師父.
這個弟兄.
不斷和他分享福音.
不過他.
很多的知識.
但是他現在面對.
就是生與死.
在癌症的末期.
所以這個弟兄.
很不捨得他這個師父.
將會離開.
他一直繼續用聖經.
和他分享.
我昨天晚上馬上飛的士.

$^{1281}$過去廣義醫院.
然後我就.
見到那個.
老人家的太太.
回家休息.
剩下那個弟兄.
和另外.
他的舅子.
當晚在這個病房裡.
我們四個男人.
我就和這個.
躺在病床裡面插著口.
他開不了眼.
我說如果你是.
林牧師.
我來探望你.
你的徒弟.
Tom.
我說是Tom請我來的.
他很關心你.
很簡單.
我和他說.
如果你聽得到.
你搖搖頭.
他真的搖搖頭.
接著我很簡單.
將福音和他說一次.
如果你也聽見.
你也願意.
去接受耶穌.
我們其他.
幾個也一起.
看到他再搖搖頭.
我再問多一次.
確保他不是.
敷衍我們.
我再問的時候他也是再搖搖頭.
很清楚要信主.
頂尖我們.
是上帝的恩典.

$^{1321}$不是靠我們自己.
我去的時候在路途當中.
我問神的靈.
這個病人是在生與死之間.
我真的沒有把握.
我只有把握.
我知道耶穌你能夠救他.
我知道我講耶穌.
你生命的故事.
是你怎樣愛我.
你也要愛他.
很簡單.
跟他說救恩.
鼓勵他認罪悔改.
接受耶穌.
住在他裡面.
他口一句一句跟我決志.
但是我請他.
內心跟我一起祈禱.
今天下午三點.
我要再去幫他做洗禮.
他的太太.
也會在那裡等我們.
求主也保守.
他的時間也仍然在.
我們不知道.
因為這些全部.
不是我們工作時間表裡定的.
所以有時候你們要愛惜.
你們的傳道人.
愛惜你的牧者.
其實他們的工作.
真的很多時候.
都是為了我們的靈命.
為了我們各方面.
那種付出.
那種代價.
不問酬勞.
只是很願意看到弟兄姊妹生生成成.
你們在基督耶穌裡面.

$^{1361}$你們能夠生根建造.
去多結果子.
能夠多邀請.
未信主的朋友回到教會.
今天地上的教會.
很多地方都會有荒涼.
都會有一個脫節.
慢慢地.
每一間教會都好像在老化.
但是我們.
就不向這個現實低頭.
我們就是要靠著.
耶穌去得勝.
我們要得回我們這個世代裡面.
每一個青少年.
我們要從魔鬼的手中.
將他們奪回.
是歸服在耶穌基督十字架的名下.
這是我們今天這個時代.
每一個地域的教會.
都需要做的.
我們不知道還有多少時間.
過去這幾年.
變化都很大.
由COVID.
到AI的出現.
到現在世界亂七八糟.
最近有人問我.
穆斯基.
是不是快要第三次世界大戰了.
我說.
你知不知道.
第三次世界大戰其實幾年前已經開始了.
他說 是嗎.
現在今天打仗.
不一定是去到你家.
大家平靈白難.
雖然.
有一個總統他的名字叫.
尼塞浦.

$^{1401}$他到處去搶奪人家的東西.
他喜歡的東西.
你不給他他就要搶.
但是今天.
我們在幾年前.
已經有很多國與國之間.
經濟戰.
打經濟戰.
打金融戰.
打這些.
國家安全戰.
大家知不知道.
我們今天已經身處一個.
非常危險.
動盪的一個世代.
耶穌說他回來之前.
好像賊一樣.
要我們警醒.
我們可能聽到一些.
戰爭的風聲.
但隨時耶穌都可以回來.
那四匹馬.
在啟示錄所形容.
已經整裝待發.
已經預備好了.
今天我們說科技.
這些AI 這些晶片.
全部在這裡.
我們今天如果還是沉睡.
我們終不懂得愛上帝.
不懂得愛別人.
我們是抵打的.
你今天反省一下.
你的公司 你的家.
有誰還未信主.
你的朋友圈裡面.
我們剛剛過了年.
還可以繼續拜年.
在.
農曆年裡面.

$^{1441}$我們有沒有派利是之餘.
我們同樣關心他們的靈魂.
他們的得救.
我們請他們.
去報道會.
講你的見證.
這個星期很特別.
除了剛才講廣元醫院.
另外一個在養和醫院.
探了兩次.
都是一個.
很有成就的一個大律師.
但一樣是.
他旁邊有一些朋友.
很關心他.
送了聖經.
原來他有一個拍檔是基督徒.
送了聖經給他.
貼上標籤.
請你看這本書.
由昨天開始.
他的女兒.
從紐約飛回來探望爸爸.
我叫他的太太和女兒.
他精神好的時候.
就在他床邊.
讀這些聖經給他聽.
所以今天.
可能去完廣元之後.
如果他們那邊準備好.
我就會探望另外一個長者.
兩個都是.
都是已經.
八十幾歲的人.
但他們都是癌症.
去到一個最後的階段.
很不容易.
時間我不知道還有多少.
但神給我有多少.
的機會.

$^{1481}$我就準備去做這些事情.
我們在.
恩愛夫婦營裡.
送給我們每一對夫婦的櫃展.
很好.
雖然很衰弱.
但我們用禱告.
來承托.
學習彼此一齊去認罪.
彼此一齊互相饒恕.
為我們的配偶去感恩.
這都是我們要好好學習.
第四方面.
犧牲性命.
耶穌告訴我們.
沒有一種愛.
比人願意為朋友捨命更大.
現在是耶穌.
要我們去死.
不是普通一個人為我們去死.
這種捨命我們就知道.
什麼是愛.
而這個愛就是阿加皮.
是上帝最神聖的那種愛.
這幅照片很有來頭.
這幅照片是一個.
國際.
一個.
攝影比賽.
一個得獎.
原來攝影師跟我們解說.
這幅影片.
有一個背景.
大家看到中間的主角.
其實是一個媽媽.
是一個鹿.
那些鹿的媽媽.
但她旁邊有.
有多少隻獵豹.
看不看到.

$^{1521}$有多少隻.
你們的眼睛很利.
你看到三隻.
因為有一隻.
第三隻在前面那兩隻豹.
遮住了牠的頭.
但你看到牠是新的.
這三隻獵豹已經張開嘴巴.
準備吞下去.
這一隻母鹿.
是必死無疑.
但原來背景.
是一個媽媽.
母鹿在照顧.
兩個小孩的小鹿.
但遠遠看見.
這三隻獵豹.
看到牠們.
牠就來了.
跟著就馬上收起.
另外那兩隻小孩.
叫牠們離開.
去另一邊.
然後牠自己衝去那幾隻野鹿.
那個獵豹那裡.
迎面被牠們追.
分散牠們的注意力.
等牠們去襲擊牠.
當牠們捉到牠.
準備要吃牠的時候.
你看牠的眼神.
是不是很害怕.
牠氣定神閒.
牠是預備了.
壯烈犧牲.
是這個媽媽.
為了自己兩個小孩犧牲.
照片是說這件事.
而我們的上帝.
是犧牲牠的獨身愛子.

$^{1561}$為我們去死.
而我們今天如果不懂得愛神愛人.
我們有愧於上帝.
給我們的拯救.
最後一件事.
老約翰在約翰一書三章說.
我們相愛不要只是在舌頭上.
如果我們見到.
你已經窮乏.
我們不要塞住連戌的心.
否則我們是假的.
如果你說你愛神.
你見到你現在在教會.
甚至教會以外.
一個有需要的陌生人.
你都不要塞住.
你自己連戌別人的心.
相愛不要只是在言語和舌頭上.
要在行為和誠實上.
我們教會講福家.
在去年六月一日.
我們在香港仔.
開始了關愛中心.
這個關愛中心叫.
Care Center.
KFCC.
康復Care Center.
在右手邊在十二月的時候.
我們有聖誕節.
我們邀請那些街坊來聚餐.
我們是針對那裡的基層家庭.
和一些SEN.
有特殊教育需要的孩子.
請他們來.
一起吃飯.
當晚有呼喚訊息.
有見證和他們分享.
帶領那些街坊.
來信主.
這是神給我們的恩典.

$^{1601}$我們樂意和他們分享.
今天我們要走出去.
我們的社區.
也要走出我們的安書區.
我們要謙卑服輸.
饒恕,赦免,恆久忍耐.
犧牲,赦命.
和扶貧相助.
這都是在耶穌的生命裡.
可以看得見.
神是怎樣.
透過他的兒子.
那種完美的示範.
讓我們知道什麼叫愛.
而這個相愛的目的.
是為了給未信的人.
可以認出我們是主的門徒.
作為主的門徒.
這是一個活生生的見證.
我們屬靈的前輩.
騰金輝牧師.
他在2002年.
出版的這本書.
是依乎所書學習.
在聖靈中長進.
我覺得這番說話.
到今時今日.
都是一語中的.
是在說我們時代的需要.
他說這個世界最需要的是愛.
不是仇恨.
不是鬥爭.
而是伸出雙手.
接納基督的愛.
讓祂的愛充滿我們裡面.
這時候一切的隔膜.
一切的難阻.
一切的階級都超越.
如果全人類都在基督的愛當中.
我們是手拉手.

$^{1641}$這就是天國.
人間可以呈現.
天堂的一個美景.
是在乎你和我.
今天撥出神的愛.
如果在祂的孤星裡.
有一句很出名的說話.
那句說話是怎樣說的呢.
他說.
當我們懂得.
去愛別人.
你就.
去幫助.
一個有需要的人的時候.
你就彷彿.
見到神的面.
今天我們要做一個.
清心亦都真心.
去愛神愛人的人.
讓我們行出.
這些行動.
我們有這個行動.
去愛別人.
你就會見到神的面.
在結束的時候我們看到.
耶穌基督的福音.
是改變了我們的生命.
我們今天要合一.
要彼此相愛.
我們要愛神愛人.
愛我們的國家.
愛我們的家庭.
愛罪人但不是愛罪.
我們要團締這個愛.
我們作為神的兒女.
每一個基督徒有這個責任.
去活出基督的愛.
大家都有責任去建立.
國家和我們家庭.
美好的未來.

$^{1681}$願上帝賜福大家.
讓禍福堂在未來的五十年.
如果轉未回來的話.
你們繼續都是忠心.
好不好.
頂子妹你說好不好.
我們手牽手.
一起去傳揚福音.
我們歸榮要給神.
三一的真神.
給我們有這種動力.
我們有這個立志.
我們要得回.
今天仍然有很多在救恩以外.
一些未信主的朋友.
還有我們自小在教會長大的青少年.
讓他們一個都不要流失.
Amen.
謝謝大家.
\newpage



\section{創世記 29:20}
\label{sec:031798}
\textbf{聖經人聖經事:七年如同幾天}~~~~黃克勤牧師博士
\newline
\newline
連結: \href{http://www.church.com.hk/acms/content.asp?site=cdc\&op=show\&type=product\&code=031798\&layout=sermon}{\texttt{031798}} ~~~~ 語音日期: 2026-03-07
\newline
\newline
\hyperref[sec:031797]{\small{< < < PREV SERMON < < <}}
~
\hyperref[sec:index_chronic]{\small{[返順時目]}}
~
\hyperref[sec:031799]{\small{> > > NEXT SERMON > > >}}
\newline
\newline
創世記 29:20
\newline
\begin{longtable}{cl}
\hline
\hline
章節 & 經文 (和合本修訂版)\\
\hline
29:20 & \begin{tabularx}{0.7\textwidth}{X} 雅各就為拉結服事了七年；他因為愛拉結，就看這七年如同幾天。 \end{tabularx} \\ \\
[1ex]
\hline
\hline
\end{longtable}
$^{1}$弟兄姊妹,事情峰迴路轉.
雅各在他舅父拉班居住.
居住了一個月.
然後就出事了.
出什麼事呢?.
結婚.
我們看看聖經人聖經時.
怎麼說這件事情.
上帝怎麼做事.
《生命說宣教網絡》.
很希望你和我們一起去宣教.
每天五元奉獻.
一個月一百五十元.
會幫助到很多人認識福音.
而我們會送一個每天的聖經人士給你.
再說一次.
這個聖經人士本來是每天金句.
每天金句就短一點.
一字金句.
現在聖經人和聖經故事.
你覺得這麼長.
本來以前十分鐘聽完.
現在長了十五分鐘.
現在三十分鐘都有.
同事說.
因為講故事.
故事會長一點.
而且未必有一節很簡單.
精簡金句.
不過哪一節重要一點.
和大家看看故事的重心.
今天是七年如同幾日.
都不是七日.
七年如同幾日.
過了幾日.
最近我媽媽家的姐姐.
就染到疫症病了.
很慌亂很慌忙.
眨眼幾日了.
這麼快幾日.

$^{41}$快一個星期了.
可能病徵快沒了.
很慌.
七年如同幾日.
七年很長.
如同幾日.
拉班對雅各說.
你雖然是我的骨肉.
是我的骨和我的身體.
一樣.
豈可白白地服侍我.
問好.
即是雅各都聰明.
懂得看眉頭眼額.
像少爺一樣.
是主人拉班妹妹的兒子.
像少爺一樣.
拉傑出到牧羊.
他在家裡也像少爺一樣.
他做事.
做什麼都沒有說.
可能是跟著拉傑出牧羊.
我都懂的.
我去幫忙.
還是在家裡.
我跟媽媽梨伯家.
學了兩手好菜.
你們去我煮飯.
我不知道他在做什麼.
我掃掃地.
不過他有做的.
就是服侍.
這點很重要.
這令拉班喜歡他.
豈可白白地服侍我.
即是他有做事.
請你告訴我.
請告訴我.
好好聲氣.
你要什麼為公家呢?.

$^{81}$你要什麼公家呢?.
公家其實是舊的中文.
其實是公知的意思.
Masquerade.
Masquerade就是公知.
你要什麼薪金?.
我想給你一份糧.
人情還人情.
骨肉還骨肉.
數目要分明.
這樣也是.
大家都學到功課.
這件事背後.
公公平平.
你來我這裡.
你又做事.
本來我想招待你.
你真是令我很開心.
因為你去做事.
又服侍.
於是我出糧給你.
我不想你白做.
我應該招待你.
不過你又做這麼多事.
不如出糧.
人情還人情的數目要分明.
拉班都幾聰明.
其實我們做事應該是這樣.
你想要多少薪金?.
這裡拉班有點不聰明.
可能他先放風.
其實拉班都很聰明.
會騙他.
很聰明的.
先放風摸底.
想要什麼薪金?.
想要多少薪金?.
摸底.
十萬九千七.
趕你走才行.

$^{121}$這麼大口.
不用了.
我做事應該的.
我在這裡服侍.
這樣很好.
隨便拿一點.
你說多少就多少.
可能他放風.
大家對答.
談判.
要什麼薪金?.
這個時候.
雅各有沒有想好?.
有.
等他問這句很久了.
其實第十六節.
拉班有兩個女兒.
大女兒叫利雅.
利雅的意思是什麼?.
不要改這個名.
利雅是什麼?.
疲倦.
利雅是疲倦.
可能這個寶寶出生的時候.
利雅出生.
這個寶寶女兒.
經常都睡.
很累.
看著看著又睡.
很累.
不如叫她做很累.
叫她做利雅.
疲倦的.
可能.
生出來之後.
媽媽覺得.
這個女兒很累.
生完之後.
很累.
要休息一下.

$^{161}$為了紀念媽媽.
生得這麼辛苦.
這麼疲勞.
很有可能叫這個女兒做利雅.
就是疲倦的.
都可能.
後者機會大一點.
前者見到寶寶經常睡.
叫她做利雅.
疲倦的.
這是我作的.
極有可能是媽媽生她的時候.
很累.
他們喜歡這樣改.
好了.
小的.
小女兒叫拉傑.
Rachel.
就是雌羊.
剛才說過.
有兩個女兒.
很有趣.
利雅.
是不是不漂亮呢.
又未必是.
不過問題是.
雅國喜歡拉傑.
沒有喜歡利雅這麼多.
是什麼因素呢.
是因為利雅的眼睛沒有神氣.
而拉傑卻生得美貌盡秀.
讀多一次.
利雅的眼睛沒有神氣.
拉傑卻生得美貌盡秀.
我解釋一下.
沒有這兩個字.
就是沒有的.
利雅的眼睛神氣.
拉傑卻生得美貌.
神氣這個字.

$^{201}$不是很有神氣.
這個神氣的意思.
是柔軟溫和.
利雅的眼睛柔軟溫和.
沒有這個字.
有些人看她的眼睛很鮮明.
她的眼睛會吸引你.
一看她的眼睛已經在說話.
沒有說話的眼睛已經在說話.
目光很銳利.
很有吸引力看她的眼睛.
特別你知道以前的人是這樣封住的.
所以她的眼睛美不美.
一看就迫著你.
要吸引你看她的眼睛.
看她的眼睛馬上迷住了.
已經催眠了.
利雅的眼睛不是.
看上去不突顯.
不會吸引你的.
看上去很一般.
很柔和的眼神.
不會刺激你.
一看上去很銳利.
不會的.
其實不是差的.
不是不美的.
據說她沒有神氣.
中文翻譯沒有神氣.
不是那麼好.
利雅的眼睛比較溫柔.
溫和和柔軟.
不會看上去很厲害.
不是這樣的.
這是她的特色.
未必她的樣子不美.
拉結卻生得美貌盡美.
美是美麗.
貌是形象輪廓.
即是說.

$^{241}$拉結是美麗的.
而且樣子有很美的輪廓.
如何美麗的輪廓呢?.
例如我們去到印度.
印度大部分女孩和男孩.
樣子都美麗.
為甚麼呢?.
輪廓很分明.
眼睛很大.
鼻子很高.
合起來.
好看.
例如很多中東女孩.
或者是俄羅斯.
烏克蘭打仗.
很多女孩的輪廓很美麗.
拉結生得很美.
而且有很美的形象輪廓.
盡瘦就是美.
而且外表.
這兩個字.
Yafet Toa 和 Yafet Ma'a.
「美」在哪裡呢?.
「美」在很美的輪廓.
「美」在很美的外表.
整個外表看上去很美.
其實很有趣.
這兩個字比較他的眼睛.
記載的人說.
他的眼睛很溫柔.
很溫和.
不銳利.
雅各不喜歡.
雅各反而喜歡.
很美的輪廓.
和很美的外表.
雅各喜歡拉結.
「我願為你小女兒拉結」.
「服侍你七年」.
Yafet是工作勤勞.

$^{281}$雅各在這裡住.
又在工作.
數目要分明.
也想給你薪金.
不知道你想要什麼呢?.
多少薪金呢?.
雅各說.
「我不用薪金」.
「我幫你打工七年」.
「不過你要把小女兒拉結嫁給我」.
「這就是我幫你做七年」.
「免費,然後娶你女兒」.
「你讓你女兒嫁給我」.
以前是爸爸媽媽作主.
甚至是爸爸作主.
「如果你通過爸爸那關」.
「女兒就嫁給你」.
「我為你工作七年」.
「很高薪金」.
「做七年才讓你女兒嫁給我」.
雅各說.
「我接手養她」.
「我沒問你拿錢」.
「我為你的小女兒服侍你七年」.
拉班說.
「好啊」.
「我將她給你」.
「我好像給別人一樣」.
「As good as」.
「我的女兒也要嫁人」.
雅各說.
「那不如什麼」.
「不要住在帳幕旁邊」.
「住在我們家裡」.
住在帳幕而已.
「與我們一起住」.
於是雅各為拉結.
這代表他答應.
「好,給你」.
其實補充.

$^{321}$「我也是要嫁給別人」.
「不如你開始入來同住」.
「家人」.
雅各為拉結服侍了七年.
因為他心愛拉結.
看著七年如同幾日.
「做七年」.
「牧羊,勞動,煮飯,洗衣服」.
其實洗衣服是不用的.
「牧羊,勞動,煮飯,洗衣服」.
但因為他愛拉結.
所以這七年如同幾日過了.
如果你想時間過得快樂.
又過得快又過得快樂.
不會捱.
你有一件事.
就是你要心愛你的工作.
你要有一個目標.
是你心愛的目標.
於是你做下去.
七年如同幾日.
如果我們以前去打球.
很開心早一晚已經在家.
拿著球拍一拍.
很緊張.
第二天上午完崇拜.
然後吃東西.
下午跑去球場打球.
為什麼打一打已經兩小時.
天黑了嗎?.
這麼快?.
還不夠後.
其實昨晚開始等到現在.
很久.
不過一打球.
打了兩三個小時.
都覺得很快.
為什麼這麼快?.
因為我喜歡.
喜歡就快.

$^{361}$如果不喜歡.
等三分鐘都過得久.
你要有一個心愛的心.
心愛的人.
你會覺得時間很快.
心愛的工作.
你會覺得時間很快.
心愛的事.
你的目標.
你會一直這麼多年.
你都繼續心愛.
繼續覺得時間很短.
都未研究完就死了.
其實最重要.
你有一個真心愛那件事.
愛那個人.
你會過得很有意義的日子.
願主祝福你.
我們一起禱告.
主我們在你面前禱告.
很想我們能夠有人生的目標.
心愛一個人.
心愛一件事.
心愛一個教會.
心愛敬拜.
心愛讀經.
心愛祈禱.
當我們心愛的時候.
那些時間會過得很快.
當我們不愛的時候.
為什麼聖經讀的都還在.
為什麼祈禱期那麼久都未完.
為什麼敬拜那麼長.
因為我不愛.
為什麼要讀書那麼辛苦.
拿個學位.
辛苦到不得了.
不過當你心愛的時候.
時間就會一直過.
在那個漫長的熬夜裡.

$^{401}$是覺得有意義的.
又學到東西的.
主又幫助我們.
時時將自己的精神.
將自己的意志.
真的用一個心愛的態度投入去.
以致無論人與事.
我們都覺得裡面很有滋味.
很有學習.
願主祝福我們.
禱告奉耶穌名求.
阿們.
\newpage



\section{馬可福音 6:34-44}
\label{sec:031799}
\textbf{上帝的邀請}~~~~劉卓聰宣教士
\newline
\newline
連結: \href{http://www.church.com.hk/acms/content.asp?site=cdc\&op=show\&type=product\&code=031799\&layout=sermon}{\texttt{031799}} ~~~~ 語音日期: 2026-03-15
\newline
\newline
\hyperref[sec:031798]{\small{< < < PREV SERMON < < <}}
~
\hyperref[sec:index_chronic]{\small{[返順時目]}}
~
\hyperref[sec:031800]{\small{> > > NEXT SERMON > > >}}
\newline
\newline
馬可福音 6:34-44
\newline
\begin{longtable}{cl}
\hline
\hline
章節 & 經文 (和合本修訂版)\\
\hline
6:34 & \begin{tabularx}{0.7\textwidth}{X} 耶穌出來，見有一大群的人，就憐憫他們，因為他們如同羊沒有牧人一般，於是開始教導他們許多事。 \end{tabularx} \\ \\ \relax
6:35 & \begin{tabularx}{0.7\textwidth}{X} 天已經很晚，門徒進前來，說：「這地方偏僻，而且天已經很晚了， \end{tabularx} \\ \\ \relax
6:36 & \begin{tabularx}{0.7\textwidth}{X} 請叫眾人散去，他們好往四面的鄉鎮村莊去，自己買些東西吃。」 \end{tabularx} \\ \\ \relax
6:37 & \begin{tabularx}{0.7\textwidth}{X} 耶穌回答他們說：「你們給他們吃吧！」門徒對他說：「我們要拿兩百個銀幣去買餅給他們吃嗎？」 \end{tabularx} \\ \\ \relax
6:38 & \begin{tabularx}{0.7\textwidth}{X} 耶穌說：「你們有多少餅？去看看。」他們知道後就說：「有五個，還有兩條魚。」 \end{tabularx} \\ \\ \relax
6:39 & \begin{tabularx}{0.7\textwidth}{X} 耶穌吩咐他們，叫眾人一組一組地坐在青草地上。 \end{tabularx} \\ \\ \relax
6:40 & \begin{tabularx}{0.7\textwidth}{X} 眾人就一群一群地坐下，有一百的，有五十的。 \end{tabularx} \\ \\ \relax
6:41 & \begin{tabularx}{0.7\textwidth}{X} 耶穌拿著這五個餅和兩條魚，望著天祝福，擘開餅，遞給門徒，擺在眾人面前，也把那兩條魚分給眾人。 \end{tabularx} \\ \\ \relax
6:42 & \begin{tabularx}{0.7\textwidth}{X} 他們都吃，並且吃飽了。 \end{tabularx} \\ \\ \relax
6:43 & \begin{tabularx}{0.7\textwidth}{X} 門徒把餅和魚的碎屑收拾起來，裝滿了十二個籃子。 \end{tabularx} \\ \\ \relax
6:44 & \begin{tabularx}{0.7\textwidth}{X} 吃餅的男人共有五千。 \end{tabularx} \\ \\
[1ex]
\hline
\hline
\end{longtable}
$^{1}$各位節目的早晨.
很高興可以在大家當中分享上帝的說話.
今天是母親節 是英國的母親節.
是一個很特別的日子.
在當中都有不少的母親.
願主的恩典在每個母親的身上.
來彰顯恩慈,信實,平安.
讓你的家庭,婚姻和個人的生活上.
都常常經歷上帝的豐盛.
香港或者其他地方的母親節是五月第二個禮拜.
英國就是今天.
我不知道在英國的母親節裡面.
有沒有一些歌曲是大家很喜歡在母親節時點播的.
我不知道.
但是香港裡面有些歌曲在母親節前後特別多人點播.
如果要在當中的你去選一首獻給天下母親的歌曲.
你會選哪一首?.
《真的愛你》 我也是選《真的愛你》.
早前第一天節目就選了另外一些.
我說我比較世俗化 聽一些流行歌.
我就選《真的愛你》.
《真的愛你》第一句歌詞是什麼?.
無法可修飾的一對手.
帶出溫暖永遠在背後.
縱使囉嗦始終關注.
不懂珍惜太內疚.
作者在當中寫歌,作曲,寫詞.
都是那個Beyond的成員.
他就想起媽媽的雙手.
在當中很辛勞的去預備每天的所需用.
無論是煮飯或者做其它的工作.
一切的辛勞其實背後帶出那種的愛,那種真摯.
昔日年輕的他他不懂得珍惜.
他只覺得媽媽很嘮叨,很囉嗦.
但是今天當他長大成人.
或者媽媽身體慢慢開始老去的時候.
他想起昔日的媽媽所付出.
孕育成為今天的他.
當然在母親節的時間.
當我們在當中去紀念在座天上,天下所有母親的同事.

$^{41}$都讓我們同樣來珍惜我們每一個弟兄姊妹.
在教會裡面每一個弟兄姊妹的付出,我們的投入.
其實聖經裡面都是有很多對手的.
今天我選的經文是一個很耳熟能詳的聖經故事.
是《五餅二魚》的故事.
都是記載在福音書裡面.
馬太,馬可,路姦,約翰都有記載《五餅二魚》的故事.
在當中我想有些弟兄姊妹都回教會一段時間.
對這個故事不陌生的.
或者聽過很多次.
我從幼稚園,小學,中學聽過無數次《五餅二魚》的故事.
但是這個故事裡面充滿了的不單單是神大人的手.
好像剛才所說的《真的愛你》裡面的雙手.
那雙手是帶出真,至,愛 透過這雙手去祝福來的五千人.
我們看見神大人的手.
同樣我們看見一雙小手.
將有限的五個餅,兩個魚奉獻給主耶穌.
並且還有甚麼呢?.
還有很多雙手.
很多雙手是一起去參與在這個神蹟裡面.
以至最後這個神蹟.
我們看見不單單為了這五千人的需要.
並且有餘.
還有十二個男子正.
是眾人齊心,努力的成果.
所以先和大家分享這雙大人的手.
帶出真,至,愛.
經文在六章第三十四節一開始是這樣說.
耶穌出來見有一大群人就憐憫他們.
因為他們如同羊 沒有牧人一般.
於是開始教導他們許多事.
馬夫音的記載是這樣.
當讀到「就憐憫他們」.
「因為他們如同羊 沒有牧人一般」的時候.
我就發現這個形容是耳熟能詳.
我突然想起其實在《馬太福音》第九章裡面.
不是說五餅二魚 是另一個故事.
另一個故事裡面記載耶穌初初出來傳道.
整講耶穌走遍各城各鄉.
在他們的會堂裡面教導人.

$^{81}$宣講天國的福音.
又醫治各樣的病症.
他看見一大群人就憐憫他們.
「因為他們困苦無助.
如同羊 沒有牧人一般」.
於是耶穌就跟門徒說.
要收的莊稼多 但是作工人少.
所以你們要求莊稼的主.
差遣做工的人出去收他的莊稼.
兩段獨立的經文 兩段獨立的記載.
但都同樣形容耶穌有一種憐憫的心腸.
看見那個需要 那個需要就好像.
人在當中好像羊一樣.
無依無靠 無帶領 無方向.
而聖經裡面所講的憐憫.
在這兩段經文的憐憫都同一個希臘字.
這個字是說一個從內心深處.
深到一個甚麼地步.
即不是從你的眼 你的頭腦.
在這個上面.
而一路落到甚麼地方.
不單單是落到你的胃.
胃以下 你的肺腑.
你的心腸 你的腸.
腸在這裡 落到腸裡面.
即是那種的觸動.
是一路去到你內心最入最入最入的裡面.
以致產生了一種的行動.
要做一些事情.
所以不是一個表面上的回應.
是一種內心深處很真實的一個回應.
這個行動 這個是憐憫.
憐憫的英文是Compassion.
Compassion這個字的拉丁文是一個組合字.
是有兩個字 Com和Partial這個字.
Com是在講一齊的意思.
一齊 together.
Partial是在講同 一齊去痛的意思.
所以那個意思是感同身受.
一齊去痛 感同身受.

$^{121}$所以那種憐憫是去到你內心深處.
深到你很感受那個需要.
那個痛楚 以致產生一個行動.
因為一切都是源於耶穌看見那個群眾.
是沒有方向 好像沒有目的.
並且在草場上危險的.
因為草場上是充滿了豺狼,野獸.
如果沒有牧羊人去保護的時候.
羊群是很容易面對襲擊.
耶穌在那一刻裡面因為群眾.
遷起一點點來到耶穌基督面前.
被他們的需要所觸動到.
以致他內心深處產生一種的關懷.
要有一個行動 這個就是憐憫.
而在《馬太福音》裡面記載的那個憐憫.
那個行動 聖經裡面說.
在《馬太福音》裡面記載的.
耶穌就對門徒說.
要收的莊稼多 作工人少.
所以你們要求莊稼的主.
差遣工人出去收他的莊稼.
耶穌的行動是產生了一個.
去鼓勵人成為他的門徒.
成為他的跟隨者.
所以在《馬太福音》第十章一開始的時候.
耶穌就招募門徒.
十二個門徒跟隨耶穌.
之後三年多的時候.
耶穌繼續去塗造他們.
建立他們 裝備他們.
縱然他們有軟弱 有跌倒.
門徒的關係裡面不和諧.
耶穌管教他們 提醒他們.
叫他們去切實彼此的相愛.
耶穌基督在十字架上為眾人去寫基.
也為門徒去寫基.
耶穌基督受死 復活.
復活的主在當中跟他的門徒說.
約定他們要在加利利山上離去.
等 門徒就猶豫不決.

$^{161}$去到北邊加利利 在加利利山裡面等候耶穌.
而耶穌在那一刻對著那幫門徒說甚麼呢?.
「所以你們要去使萬民作我的門徒」.
耶穌基督在當中離去呼召人成為門徒.
也都希望門徒成長 接受裝備.
將來繼續延續門徒產生門徒.
建立更多的門徒.
這個是《馬太福音》裡面看見.
那種憐憫觸動到耶穌的心.
是要有更多人參與在上帝天國救贖的工作裡面.
這個是《馬太福音》.
門徒是很重要的.
《馬太福音》裡面 門徒.
但《馬太福音》裡面 耶穌是否建立門徒?.
不是 耶穌亦沒有招募門徒才說是甚麼.
於是就開始教導他們許多事.
是教導.
是說耶穌因為觸動 因為憐憫的緣故.
觸動他的心 所以他就開始做一些事.
《正經》裡面記載《馬太福音》.
就是說耶穌開始行動就做教導工作.
這個亦都很配合在《馬太福音》裡面的整個風格.
《馬太福音》風格裡面記載耶穌.
在十六章裡面記載耶穌很迅速地作出行動回應.
很快的 很快速地去做回應.
而耶穌快速地回應.
就是他自己親自在當中宣講教導上帝的話語.
事實上在整個報道 在整個宣教在上帝救贖工作裡面.
在教導上帝的話語 讓人學習上帝的話語.
是可等的重要.
今天我們看見.
第一本《聖經》中文的《聖經》是何時翻譯成為中文的?.
是二百多年前.
蘇格蘭的宣教士馬拉森來說.
其中最主要的任務就是為中國人翻譯第一本中文《聖經》.
從此讓華人能夠用母語能夠閱讀上帝的話語.
這是一個很大的成就.
他一生不是帶領很多人信耶穌.
但是他一生花上最多的時間精力學習中文.
並且翻譯《聖經》.

$^{201}$以致讓後來的人 其他的傳教士.
除了有《聖經》之外還有很多的工具書去學習中文.
這是打開了近代基督教或者新教的傳教士來華.
一個很重要的一頁.
二百多年前我們已經擁有上帝的話語.
用我們的中文能夠閱讀主的話語.
以致由現在開始有不同的《聖經》的譯本版本出現.
由原先的一個文言文 接著白話文.
接著也有不同的版本 和合版.
當代譯本 新譯本.
或者我們今天正在用的.
正在用的版本是和合版 修訂版 和修版.
有很多不同的版本都幫助我們與時並進.
我們能夠用現代的語言文字來認識上帝的話語.
也都按照不同人的需要 有時候有些《聖經》是細本一點的.
薄身一點的 容易攜帶.
也都有一些大字版的《聖經》幫助長者能夠讀到上帝的話語.
也都有一些參考對照的.
中英文《聖經》的對照版本.
中文加希臘文原文的對照.
也都有很多不同的角色奇色.
不同的大小 不同的口薄 不同的紙質.
不同價錢的《聖經》 我們都擁有.
不知在你家裡有多少本《聖經》呢?.
有些是多過一本 或者我這本是電子版本.
電子版本的《聖經》 在我的iPad裡面.
我想我們擁有《聖經》的一段時間.
但事實上今天在全世界仍然有3333個民族.
這個數字會改變的 會改變的.
但容易些記 3333 3333個民族.
未擁有他們民族第一本屬於他們的《聖經》.
母語的《聖經》.
我們已經是比這三千多個民族早了二百年.
能夠擁有上帝的說話.
因為這個緣故 我們就可以學到主的話.
在主的話裡面去成長.
今天仍然很需要有人從事《聖經》的翻譯.
剛才在休息的時候有些人提問.
還有三千多個民族沒有《聖經》 很可惜.
是啊 還有三千多個 人數不是很多.

$^{241}$佔全球80億人口裡面大約一億多而已.
不是很多 但3333個的民族.
AI的緣故是加速了《聖經》翻譯的進度.
是加速了.
以前可能要三,四十年才能夠翻譯好《新約聖經》.
現在只需要幾年的時間去做.
所以很多的項目已經是開始了.
但開始了之後 需要人做什麼呢?.
就是去做檢查.
確保那個翻譯跟《聖經》裡面的原意.
是沒有很明顯的一種的差距.
那個樽頸位是要有很多很有經驗.
從事《聖經》翻譯的宣教士.
去做很多的教對的工作.
求主席當中的幫助.
我們看見神的話語是重要的.
但我們仍然是需要有人從事《聖經》的翻譯.
縱然有AI 但還需要有很多人.
從事在《聖經》的教對裡面.
在環教會裡面 我們有很多很出色的博士.
《聖經》研究的博士.
但從事《聖經》研究的博士在當中去參與.
不是華人《聖經》的更新版.
是華人參與其它不同種族的《聖經》的製作.
讓他們能夠首先擁有整個民族平生的第一本《聖經》.
這個是很重要的.
教導《聖經》的話語,翻譯《聖經》的工作.
派發《聖經》的工作都是很重要 很重要.
十年前的今天 我有機會帶《短宣隊》.
那時在香港帶《短宣隊》去到希臘裡面.
做福音的工作 宣教的工作.
那次我們最主要是服侍在希臘裡面的敘利亞難民.
十年前敘利亞打仗 有很多難民.
從敘利亞去到土耳其 再坐橡皮艇去到希臘.
從希臘沿著陸路去到歐洲的中部.
特別去到德國,法國的地方.
也有些人去到英國裡面.
所以在服侍裡面也讓我們有很多的機會.
去服侍一些穆斯林敘利亞的朋友.
但同時也讓我們有機會認識希臘這個地方的屬靈的情況.

$^{281}$也有機會去到亞典 亞典是希臘的首都.
希臘大約一千萬人口 大約五百萬人口是住在亞典裡面.
在亞典裡面也去到亞典裡面全亞典最大的教會.
亞典最大的教會我想是和你們在Swingdon的教會裡面.
早堂,午堂的人數合併.
如果你有聯合崇拜.
你們的人數已經超越了希臘亞典最大的教會的總人數.
我們Swingdon很小 我們環教會很小 在英國裡面你覺得你很小.
但你去到希臘亞典的時候.
你現在是在亞典裡面特大的教會.
是全亞典裡面最大的教會 那間教會大約一百人.
你有沒有想像過原來你們的教會的人數.
已經是超越了希臘亞典最大的那間教會.
一般教會大約十幾人.
十幾人? 十幾人的教會?.
不是啊 打開聖經 新的聖經.
一開始的時候保羅去到亞典傳福音.
又去到哥林多 又去到帖撒羅尼加 去到不同地方.
福音在那裡已經有二千年.
難道那裡沒有基督徒?.
那裡有很多人 有很多希臘人.
他們都是一出生就知道自己是一個希臘東正教徒.
作為一個亞典人 他的身份就是希臘東正教.
但是他們有沒有機會去閱讀上帝的話語呢?.
大部份都沒有閱讀聖經的習慣.
因為他們手上的聖經就是我們傳道人.
在神學院讀那本二千年前那本希臘原文的聖經.
所以對他們很多人來說.
他們是懂得讀 但不太掌握裡面的意思.
今天如果弟兄姊妹你看回馬利信版本的聖經.
二百年前的聖經 你懂不懂讀中文字?.
懂 但是你能不能完全明白呢?.
不是很明白 又明白一點 又不明白一點.
是一個這樣的狀態.
所以對於在希臘裡面的基督徒.
他們覺得打開讀聖經這件事是屬於.
以基督教教會來說是屬於傳道人牧師的責任.
我們的責任是坐下來聽而已.
啊?這樣嗎?.
不是啊 我記得五百年前宗教改革的時候.

$^{321}$已經是打破了這個觀念.
是啊 這個在新教裡面開始說的是.
讓每一個人能夠用當時流通的語言.
用他的母語去讀上帝的華語.
這個是宗教改革Reformation.
然後有聖經的翻譯.
而馬丁路德最卓越的工作.
就是將當時的拉丁文翻譯成當時流行的德語.
因為這種德語文字也統一了當時.
在整個羅茵河流域的人.
所以今天在德語裡面叫做Higerman.
就好像中國的普通話裡面.
也有很多不同的方言.
但是希臘沒有.
所以當十年前的時候.
我認識到有一個機構.
他們的工作是做什麼呢?.
就是要將聖經翻譯成現代的希臘話.
並且是派聖經去到每一個家庭裡面.
在雅典以外 剛才說有五百萬人口.
他們期望在十年的時間.
他們已經完成了一個計劃.
但是用了十年的時間.
由2015 希望去到2025年.
就用十年的時間.
派發一百萬本希臘文的新約的聖經.
讓每一個家庭能夠直接閱讀到上帝的說話.
等於教導上帝的說話不是一些很偉大的事.
每個人都需要能夠打開閱讀上帝的說話.
不單單是透過傳道人.
不需要只是教過牧師.
也是每個弟兄姊妹自己.
或者一個非基督徒打開聖經.
直接去閱讀上帝的說話.
因為神的話就是生命.
就能夠改變人的一生.
我們在經文裡面可以看見.
耶穌就是藉著上帝的說話.
去滿足人內心,心靈的需要.
但是此刻眾人在當中聽了耶穌講道很久了.

$^{361}$所以大家都累了.
可能累了.
但是大家都很享受.
耶穌也很投入在教導裡面.
當耶穌很投入在教導的時候.
門徒就發現開始天色已經晚了.
所以門徒就進前來跟耶穌說.
這個地方很偏僻.
已經很快天黑了.
所以要不要叫眾人來散去.
以致他們去到附近的村莊.
去買自己所需要的食物.
天色昏暗.
當時是二千年前.
沒有一個很好的照明系統.
天黑真的入黑不見五紫.
真的看不到.
所以門徒做了一個風險評估之後.
建議耶穌一個可行的方法.
就是不如現在先停止聚會.
休會先.
然後有秩序安排群眾離開.
去到附近的村莊.
按他們需要去買他們的晚餐.
以免造成危險和種種不便.
這班門徒是不是一家很有質素的門徒?.
很有責任感.
很懂得為耶穌打算.
但似乎耶穌有沒有接納他們的意見?.
耶穌沒有.
反而耶穌就反問他們什麼?.
「你們給他們吃飽」.
你們去解決這個問題.
接著門徒聽完之後.
他們又很乖地按著耶穌的建議.
「你們給他們吃飽」.
就開始計算了.
我們這12個門徒.
我們身上整副身家有多少錢?.
有200個銀幣.

$^{401}$200個銀幣.
200個銀幣即是多少錢?.
當時每個人工作一天的工資就是1個銀幣.
200個銀幣就是200個工資.
200天的工資.
200天的工資就大概約莫6個月到7個月的人工.
如果今天在英國裡面.
我們的人工除稅之後.
約莫2000鎊.
也不是太差的.
也不是很多.
約莫2000鎊.
2000鎊6個月的人工.
6個月還是7個月?7個月吧.
7個月即是總共大概有多少錢?.
就是14000鎊.
14000鎊又不是很多錢.
但如果你拿著14000鎊去Tesco.
去買Meal Deal.
買三文治 買薯片 買果汁.
餵飽這麼多人.
14000鎊大概買了多少Meal Deal?.
大概不過3000個Meal Deal.
但是聖經說的人數有多少?.
男人也不算.
所以基本上他們知道他們手上的2000個銀幣.
不能夠解決所有的需要.
即是滿足部分的需要.
或者大家節省一點.
不要吃那麼多 吃少一點.
眾人也可以滿足到那個數字的需要.
但事實上醫不了他的肚腹.
醫不飽眾人的需要.
是不是耶穌刻意將一些解決不了的難題.
加諸於門徒身上.
以至顯出門徒的無知.
門徒的愚昧.
然後顯出耶穌的大能.
耶穌的睿智 耶穌的能力.
耶穌不需要搞這種東西.

$^{441}$耶穌不需要搞辦公室政治.
是不需要的.
耶穌是一個大能的神.
在舊約聖經裡面我們看到大能的神.
按著祂的憐憫 祂的能力.
當以色列人在曠野40年的時候.
神按著祂需要供應嗎哪和鵪鶉40年的時間.
以至以色列人在曠野裡面一無所缺.
當他們缺乏的時候.
祂埋怨上帝就供應了.
今天五千人.
五千人而已.
比對於曠野的時候.
按聖經學者所估計.
約莫有二百萬.
二百萬的群眾在曠野裡面漂流了四十年.
今天五千人 區區五千人.
難道耶穌都不能夠自己去解決.
滿足這個需要嗎?.
似乎耶穌是想門徒明白一個更重要的功課.
所以就拋一個問題出去給門徒.
要門徒去明白.
是呀 手頭上你擁有多少資源?.
二百個銀子.
二百個銀子能不能夠滿足餵飽五千人的需要?.
不可能.
你擁有的資源是這麼多.
可以摸到的.
感受到的重量是這麼多.
是有形的.
英文說是一個tangible.
觸摸到的.
但同時間透過這個醒覺的事件.
這個神蹟.
耶穌要讓門徒明白另一樣東西.
就是有些東西你觸摸不到的.
無形的intangible.
但是存在.
是要藉著信心的呈現.
交託給主.

$^{481}$多少不重要.
按你的能力.
你的信心呈現給主.
然後就好像剛才詩歌裡面所說的.
在主手中一切都奇妙.
就經歷奇妙的滋累.
奇妙的改變.
經歷這個神蹟.
耶穌是要讓門徒很深的去明白.
將你有的.
多和少不重要.
用信心呈現.
在主手中.
這個無形的資源,能力.
就可以去打開.
就可以打開你從來未打開過的一部分.
這個部分.
其實聖經裡面在《馬太福音》的十七章二十節.
耶穌也是這樣跟門徒,群眾說.
祂說:你們若有信心好像一粒芥菜種.
芥菜種的種子是多大?.
就好像一粒芝麻一樣.
是芝麻.
你有信心好像芥菜種.
就是對這座山說.
你從這邊挪到那邊.
它也必挪去.
並且你們沒有一件不能作的事.
沒有一件事不能作.
信心的呈現去打開一個看似你觸摸不到,看不到.
但存在,是真實的.
看似是虛無.
但是真實的.
在主的裡面一切都是奇妙.
在經文六章第三十八節.
耶穌就問門徒.
門徒就說:你們有多少個餅?.
你們去看看.
他們知道後就說有五個餅.
還有兩條魚.

$^{521}$資源就是這麼多.
這個資源是誰提供的?.
門徒?.
經文裡面沒有說.
這個資源不是門徒提供的.
是有一個小男孩.
來獻上一個孩子.
是不是男孩呢?.
不知道.
不過以往的聖經作者都會說是一個男孩.
但可能是一個女孩.
沒有人知道.
總之是一個孩子.
將他自己所擁有的獻情給主.
是五個餅,兩條魚.
份量不是很大.
不是這樣.
五個大麥餅.
記得以前在教會裡面聽過道.
穆斯林強調五個大麥餅.
我們知道大不是形容大和小.
大麥是一種材料.
也有小麥.
這個餅是用大麥來做的.
不是尺寸大.
是一個大麥所做的餅.
兩條魚.
也不是因為這個很神奇的組合.
才產生一個神奇效應.
其他的比例是不行的.
我們煮東西吃.
我們知道要按那個比例.
多少油,多少麵粉,多少蛋才能發起一個麵團.
比例改變了.
次序改變了.
是不可行的.
聖經裡面不是說這個神奇的比例組合.
是要五個餅,兩個魚才產生神奇效應.
一個餅,六條魚不行.
不是這樣的.

$^{561}$正如剛才所說.
好像詩歌.
唱的詩歌.
剛才的分享講道裡面也說.
「在神手裡一切都奇妙」.
只不過在那一刻裡面.
耶穌就按所呈現的.
哪怕是兩個餅,三條魚.
在神手裡.
糟化作奇妙.
不是要五個餅,兩個魚.
假如那個小孩賣獻上的.
舉個例子.
獻上一個石頭給耶穌.
所謂要忍化耶穌.
那麼耶穌會不會在當中將它變成.
餵飽五千人的東西呢?.
我不知道「在主手中一切都奇妙」.
石頭可不可以變成餅呢?.
魔鬼知道的.
在《馬太福音》第四章第三個字.
魔鬼說「你若是上帝的兒子」.
「你就可以將石頭變成餅」.
魔鬼也知道耶穌是有.
轉化,改變,行神蹟的能力.
魔鬼也知道.
但是弟兄姊妹.
我們作為弟兄姊妹.
我們知不知道耶穌有這樣的能力呢?.
耶穌是可以就地取材.
在當中將所有東西轉化.
成為祝福不同人的資源.
所以重點不是魚,不是餅.
是耶穌自己.
是大能的神.
今天祂邀請在座的每一個弟兄姊妹.
也邀請祂的教會,宗教會.
在地上的宗教會.
一起參與主奇妙救贖的旅程.
這個奇妙救贖旅程.

$^{601}$由舊約到主耶穌復活.
耶穌的受死,復活,升天.
一直到今天.
繼續延續這二千年的救贖.
奇妙救贖旅程的工作.
一直繼續邀請祂的門徒.
跟隨耶穌基督人.
地上的弟兄姊妹.
一起參與上帝的救贖使命工作.
參與在整個大使命的行列裡面.
去洗萬民濁耶穌基督的門徒.
不論祂的膚色.
不論祂的文化.
不論祂的語言.
眾人都是何等寶貴.
神看每一個人都何等寶貴.
不是因為祂的種族,祂的膚色.
所以祂不配得到上帝這份恩典的禮物.
是不分彼此.
眾人都是這樣.
何等重要.
在這個奇妙救贖旅程裡面.
我們看見不僅有全能的手.
還有一雙小手.
一個小孩子.
不知名.
這個不知名的小孩.
代表著可能是不顯眼.
可能被忽視.
微不足道.
又或者是覺得.
這樣行不行啊?.
行不行啊?.
連兩百個銀子.
都不能夠滿足這個需要.
區區五個餅,兩個魚.
就能夠滿足這個需要.
你會不會天真一點啊?.
你會不會真的這麼無知啊?.
這麼蠟燭啊?.

$^{641}$但是這個小孩子.
代表的不是入世未深.
不是代表不認識人情世故.
而是代表著一種.
將自己僅有.
用信心獻情給主的心.
這個心在在座的群眾當中.
佔的比例是多少啊?.
五千分之一.
所以說餵飽了多少個男士啊?.
五千個.
所以原則來說.
這個男孩屬不屬於五千個男士.
一分之一.
不屬於.
是五千分之一以外.
不過用這些計算.
我就當是五千分之一.
那個數量.
那個數目是很少的.
零點零幾多個百分比.
我也不懂得計算.
零點零幾個百分比.
是絕無僅有的.
但他卻是將他自己有.
無私的獻上.
那份心.
其實在整個的神救贖.
奇妙救贖的.
工作計劃.
宣教計劃當中.
上帝都興起了很多.
無名小子.
一些無名小子.
不出名的人.
今天我們就叫他們做英雄.
無名英雄或者偉人.
昔日他獻出的時候.
都是一個很小的回應.
就好像何爾斯.

$^{681}$這個小女孩.
她當時六歲.
她將她的心獻給上帝.
她聽見一個故事.
宣教士去到她家裡面去住.
就和她爸爸媽媽聊天.
她就坐在那裡聽.
聽見那個叔叔.
一個美國人講述.
在中國農村.
百多年前中國農村裡面.
中國是一個務農的社會.
所以很多時候農夫都很忙.
他在秋天的收割裡面.
稍為有多一些時間.
他有多一些時間.
他就想起過去的時間.
有些傳教士來到我們村莊裡面.
就講有福音.
要接受福音.
福音能夠改變心靈.
他們又不太明白.
甚麼叫做福音.
所以稍為沒那麼忙的時候.
就去到福音站.
福音站就是宣教士工作的地方.
拖著木頭車.
去到宣教士工作的地方.
宣教士站.
就想著和這班外國人 洋人.
拿一些福音回到村莊裡面.
和人去分享.
農民他們的世界.
就是那麼簡單.
或者是那麼無知.
但是在一個錯配的情況下.
當時那些洋人的傳教士就很忙碌.
就沒有辦法去招呼他.
所以就打發他.
不如你明年再來.

$^{721}$所以那個農民就拖著一個空的木頭車.
回到村莊裡面.
一年之後.
當沒那麼忙的時候.
又想起有些叫做福音.
又拖著木頭車.
又在那段時間裡面去探訪.
福音站的傳教士.
很不幸.
那些傳教士都是很忙.
所以他就說.
如果有多些傳教士去到中國裡面.
就可以幫助這個農村的村民.
讓他們明白福音不是一個有形的物體.
福音是講述耶穌基督的愛.
耶穌基督的犧牲.
耶穌基督所賜下的救恩.
這份禮物.
小小的6歲的荷爾斯聽見這個故事.
就放在心裡面.
當我長大成人.
我希望可以去中國裡面.
讓農民能夠知道什麼叫做福音.
19歲的荷爾斯在城鎮裡面.
不是一帆風順的.
有起有跌.
有些時候也很想在世界裡面.
追逐那個美好.
那種富貴.
那種豐富.
他家裡是開百貨公司的.
所以家裡不是沒有錢的.
有錢的是加州.
加利福尼亞州.
但是在他成長裡面.
他去讀書.
讀完第一年大學.
在史丹福大學讀完書的時候.
他覺得裡面不是很開心.
不是很滿意.

$^{761}$所以100年前的媽媽跟他說.
不如你聽著.
那種前衛.
很難去想像100年前的人是那種前衛.
不如你take a gap year.
去中國裡面做一年短孫.
你想一下現在的短孫好像很偉大.
好像很革命性.
不是100年前在美國裡面.
take a gap year.
去中國裡面.
去感受一下甚麼叫做宣教傳福音.
19歲,20歲何爾斯就離開了美國.
加利福尼亞州.
就去到中國今天的三水一帶的地方裡面.
那些港府人的地方.
就在當中去參與一些服事.
19歲,20歲的小女孩.
很容易就學到一些廣東話.
就在那裡參與.
參與完之後.
他發現他找到他最大的快樂.
於是一年之後回到家裡.
跟家人說.
我要去中國裡面繼續做宣教工作.
就這樣他就在中國裡面有10年,20年.
那段日子是中國最混亂的情況.
軍閥割據之後有很多的衝擊.
接著中日戰爭.
他就在那段時間裡面.
在大陸裡面去為一些婦女提供一些工作.
因為大家沒有收入.
所以安排他們去織一些蕾絲繡花.
然後透過他的脈絡關係.
就做出口這些蕾絲繡花.
去到美國,歐洲裡面賣.
賣了之後就有錢了.
有錢就可以養當地的人了.
婦女沒有社會地位.
婦女也沒有機會來讀書.

$^{801}$所以在所謂的工廠裡面.
晚上就在同一個地方裡面.
就聘用一些人.
教他們寫中文.
讀書,寫字.
為婦女製造工作.
提供就業機會.
亦都給他們學習的機會.
改善他們的生活.
為他們能夠建立一種更優質的生活.
荷爾斯年輕的女子.
就在中國裡面做了這些工作.
加上打仗的緣故.
所以很多家庭.
很多的父母離世.
或者沒有能力照顧子女.
所以很多人都見到荷爾斯這麼有愛心.
所以就將孩童放在.
她的笑工廠的門口.
希望她照顧.
所以她從來不是計劃要開孤兒院的.
但是因為環境的緣故.
所以她就開了一個叫爛童院.
困難的爛.
爛童院 不是孤兒院.
有些是有父母的.
不過家裡沒有能力去照顧.
她就去照顧他們.
為他們提供一個家.
並且亦都是提供一個屬靈的家.
讓他們在當中去認識上帝.
有些就是在爛童院裡面.
一路一路去成長.
成長了之後.
這些兒童.
當時在一個困難的環境下.
一群兒童 一群少年人.
他們人生有什麼夢想嗎?.
有.
他的夢想就是希望在當中來見證神.

$^{841}$傳福音.
就是這樣有些這樣的爛童.
一些的兒童.
青年人的成長.
要在事奉上帝.
所以就安排他們去神學院裡面讀書.
有部分的就去到當時在廣西吳州的.
建道聖經學院裡面讀書.
讀完之後.
當時就差派他們.
要繼續去傳福音的工作.
有些爛童就下來香港.
慢慢去成長 有工作.
就想起以前那個美好.
我們在這個爛童院的生活很好.
有些就做了傳道人.
後來下來香港.
長話短說.
之後何懿詩在45年之後.
完了二次大戰.
回到家鄉美國加利福尼亞州.
應該有十幾二十年沒有回家裡面了.
然後又想起.
會紀念以往的這些兒童.
這些婦女.
所以回到香港.
當時她家裡有些錢.
所以她當時在三水的地方裡面住.
但是她也有些時候會來到香港.
來有一個安靜 退休.
和一個休息的時間.
所以家裡給了她少少錢.
所以她就在長洲那裡買了一個小小的房子.
下來香港休息.
休息 短暫的休息.
然後回到大陸裡面繼續工作.
就是那個背景.
所以她回來香港的時候.
她就在長洲住的地方.
就在當中去迎聚一班青年人.

$^{881}$各方面迎聚他們.
接著這個地方後來.
長話短說因為太豐富的故事裡面.
有一本書叫《誰掌管明天》.
這是中文翻譯.
原文是何懿詩寫的.
瑪麗寫的叫《The Good Hand of Our God》.
上帝那種的好 慈愛 美好的手.
最後她就把這個地方奉獻了出來.
把這個小小的房子奉獻了出來.
這個地方是在長洲山頂道22號的地方.
知道那是甚麼地方嗎?.
見到神的地方.
其中的一個地方.
他們早期擁有的地方.
何懿詩奉獻了出來.
接著她也在建道神學院裡面被邀請.
不如留在這裡來駐校做老師.
教一些神學生.
就是這樣在當中她晚年裡面.
也有十多年的時間在香港作育.
英才建立了不同的傳道人.
也有很多傳道人成為宣教士.
在世界各地裡面來服侍.
從來沒有想過一雙這麼小的手所做的.
能夠成為這麼大這麼大的祝福.
特別是在華人教會.
在華人社會裡面.
在經文裡面我們不單單看見上帝傳人的手.
不單單看見小孩子的手.
更加重要的是我們看見有很多很多很多雙手.
是眾人的手.
是齊心的參與.
所謂齊心事成.
在六章第39節之後.
耶穌吩咐他們.
叫眾人一組一組的坐在青草地上.
吩咐他們他們有門頭.
有其他的群眾共同去參與.
他們就一群一群的坐下.

$^{921}$有一把有五十.
耶穌就拿著這五個餅,兩個魚.
望天祝福.
抹開大餅門頭.
於是照樣去分派了這五個餅,兩個魚.
聖經記載他們都吃.
並且吃飽.
並且將那些零零碎碎收拾起來.
就裝滿了十二個男子.
吃餅的男人共有五千.
所以我們看到整個神蹟的彰顯.
不單有神大人的手.
顯出他的無微不至,真,善,愛.
也有小孩的手.
更多的是很多很多對手.
共同去參與.
參與在上帝的奇妙救贖工作裡面.
這個奇妙救贖工作.
不是單單僅僅滿足需要.
不是僅僅滿足.
剛才說了什麼呢?.
是還有十二個男子來祭祀的.
在屬靈上的意思是超過我們所想所求.
超出你的想像.
從來沒有想過是這麼豐富的.
在這個五餅二肉神蹟裡面.
我們看到充滿了很多的管理.
很多的執行.
還有一些很細緻的安排.
確保流程是順利的.
不會有人去遺留.
這麼多這麼大的聚會.
有遺留其實是很自然的.
但在裡面是經過細心的安排.
確保是沒有人去遺留.
寫記載.
他們都吃並且吃飽了.
是吃飽.
是沒有人被忽略.
每個人的都能夠是充充足足.

$^{961}$並且是有緣.
一切都是所有人共同努力的成果.
所以沒有人是袖手旁觀的.
是眾人齊心努力的成果.
是人人有份.
教會的服侍都是的.
是要眾人齊心的努力.
有人去負責做主席.
敬拜,施事.
也有執事,有管理,有音控.
有詩琴.
有各方面也有會眾我們共同去參與.
如果崇拜裡面沒有弟兄姊妹坐在一起去參與.
這個崇拜不是一個真正的崇拜.
但只有弟子去參與.
沒有主席,沒有敬拜,沒有詩琴.
這個崇拜也不能夠進行的.
是要眾人一起去參與.
我們在《五兵二虞》裡面.
我們只是在當中看見上帝的手.
一開始說上帝的手.
我們看見一雙小手.
這個小手用信心將祂有的見證給主.
更加重要的.
弟兄姊妹更加重要.
今日希望你們可以帶走的最重要一點就是.
你們的角色是很重要的.
是眾人的參與是一個都不能少.
沒有人是旁觀.
每一個人的參與.
才能夠共建上帝整個宣教奇妙救贖的工作.
弟兄姊妹在上帝救贖的工作裡面.
我們可以一起去參與的.
不單單是部分的弟兄姊妹.
為上帝的角度.
為宣教工作來祈禱.
如果我們每人都願意花三分鐘.
每天三分鐘.
我花我三分鐘就微不足道了.
三分鐘.

$^{1001}$如果在教會裡面.
只有陳全道花三分鐘為教會的需要祈禱.
都是微不足道.
但若然我們教會裡面.
舉個例子有一百人.
三乘一百就有三百分鐘.
這個multiply effect是不能少的.
三百分鐘一起為教會的需要祈禱.
剛才所說我們信心少到好似芥菜種子.
產生一個什麼呢?.
一個很巨大的影響.
就算叫一座山從左邊搬到右邊.
都得在神沒有難成的事.
只怕我們不去參與.
同一個道理.
不單單我們為宣教禱告.
可以在網站App Store裡面打宣教日印.
中文的.
有粵語 有普通話.
裡面有人講的 有文字的.
有足夠的資料去幫助你.
去為宣教來祈禱.
如果你說不懂的時候.
用你的手機App Store裡面有很多不同的工具.
或者我們獻上的力所僅有很少.
十磅不是很多錢.
但是對某些頂尖智慧來說.
十磅都是錢來的.
都不是不重要的.
但是十磅可能買到.
可能在中餐館兩籠點心.
可以嗎?.
兩籠點心.
應該都不可以.
倫敦就一定不可以了.
十磅亦都不可以.
我剛才早上.
我還沒有在這裡吃過中餐.
我不知道.
Sweden中餐.

$^{1041}$我只知道十磅在倫敦.
在唐人街是叫不到一個乾炒鵝的.
是不可以的.
我住的地方裡面沒有中餐館.
不過我想吃都沒有.
所以我不知道.
但是倫敦裡面的唐人街.
是十磅吃不到一個乾炒鵝.
但是十磅可以叫大概接近三杯咖啡.
如果是美式咖啡.
齋啡的時候.
不是很多錢.
但是那個Multiply Effect.
就是說如果每人十磅乘以一百.
你想一下.
那個力量 是不容忽略.
如果弟兄姊妹有一個心.
很想去服侍.
但是如果有多幾個人.
我們就齊心就事成.
就可能組織到一個短宣隊.
開始有某個參與去服侍.
上年我有機會和倫敦華人宣導會.
一班弟兄姊妹去到波蘭服侍烏克蘭難民.
這是一個很特別的旅程,一個經歷.
他們參加弟兄姊妹全部都是BNO.
所以當站穩神腳.
在英國裡面生活各方面.
每人都有一個小小的心.
為什麼要去短宣呢?.
時間關係不可以和你們分享.
但是這次的短宣裡面.
正正是滿足他們裡面.
有一個小小向上帝呈現的心.
因為他們這個行動.
我也鼓勵了另外一間教會.
牛津教會.
今年他們在六月到七月的時候.
都會去波蘭服侍烏克蘭難民七天.
一個星期.

$^{1081}$這個就是所說的multiply effect.
那2027年有沒有哪一間教會想延續呢?.
延續工作呢?.
會不會在你們當中呢?.
不知道.
但是大家可以為這件事去思考.
在當中祈禱.
最重要的就是有力就出力.
有錢就出錢.
齊心努力一起去榮耀主.
見證主因迎接耶穌基督的返來.
我們作一個土稿.
天父我們多謝你.
讓我們能夠在教會裡面.
我們一起共同參與你自己天國的工作.
所以我們來自五湖四海.
今天我們就在這個地方裡面聚集.
我們一起去學如何去敬拜主.
如何以心靈和誠實.
來將最大的榮耀獻給我們的救主耶穌基督.
我們仍然在當中可以學習.
如何去以愛還愛.
是因為你的愛的緣故.
讓我們經歷恩典.
又讓我們去學習用我們的手.
用我們的口.
最重要的是用我們的心.
來將我們有獻給你.
無論是多少.
這個都不重要.
最重要的是我們的態度.
我們的心和態度.
願祝福教會.
祝福在當中每一個弟兄姊妹.
特別是一班的媽媽.
讓她在未來的時間.
繼續經歷你的恩典,你的祝福.
紀念不同的人的需要.
無論是身體,心靈各方面.
主你知道.

$^{1121}$你知道我的需要.
就願你的恩典繼續帶領我們.
跨過面前種種的挑戰和困難.
讓我們一起彼此.
作為一間教會.
是同心一起去見證.
耶穌基督的恩賜和信實.
多謝主.
順應禱告奉耶穌名求.
阿們.
\newpage



\section{約書亞記 3:1-6}
\label{sec:031800}
\textbf{你有沒有期待看見神的作為?}~~~~劉少康牧師
\newline
\newline
連結: \href{http://www.church.com.hk/acms/content.asp?site=cdc\&op=show\&type=product\&code=031800\&layout=sermon}{\texttt{031800}} ~~~~ 語音日期: 2010-03-01
\newline
\newline
\hyperref[sec:031799]{\small{< < < PREV SERMON < < <}}
~
\hyperref[sec:index_chronic]{\small{[返順時目]}}
~
\hyperref[sec:031801]{\small{> > > NEXT SERMON > > >}}
\newline
\newline
約書亞記 3:1-6
\newline
\begin{longtable}{cl}
\hline
\hline
章節 & 經文 (和合本修訂版)\\
\hline
3:1 & \begin{tabularx}{0.7\textwidth}{X} 約書亞清早起來，和以色列眾人起行，離開什亭，來到約旦河，過河以前住在那裡。 \end{tabularx} \\ \\ \relax
3:2 & \begin{tabularx}{0.7\textwidth}{X} 過了三天，官長走遍營中， \end{tabularx} \\ \\ \relax
3:3 & \begin{tabularx}{0.7\textwidth}{X} 吩咐百姓說：「當你們看見利未家的祭司抬著耶和華－你們神的約櫃的時候，你們就要起行離開所住的地方，跟著約櫃走， \end{tabularx} \\ \\ \relax
3:4 & \begin{tabularx}{0.7\textwidth}{X} 使你們知道所當走的路，因為這條路是你們從來沒有走過的。只是你們要與約櫃相隔約二千肘，不可太靠近約櫃。」 \end{tabularx} \\ \\ \relax
3:5 & \begin{tabularx}{0.7\textwidth}{X} 約書亞吩咐百姓說：「你們要使自己分別為聖，因為明天耶和華必在你們中間行奇事。」 \end{tabularx} \\ \\ \relax
3:6 & \begin{tabularx}{0.7\textwidth}{X} 約書亞對祭司說：「你們抬起約櫃，在百姓的前面過去。」於是他們抬起約櫃，走在百姓前面。 \end{tabularx} \\ \\
[1ex]
\hline
\hline
\end{longtable}
$^{1}$大家好 我是劉少康牧師.
今天想跟大家思想的題目.
就是你有沒有期待看見神的作為呢?.
我們要讀的經文就是約書亞記第三章1-6.
到底今天還有沒有神跡呢?.
我不知道你說的神跡是指什麼.
是不是跛腳突然間康復.
有病的一聲「沙」的癌症就不見了.
還是我們所指的神跡.
就是神聽穩我們的祈禱.
神對我們的生命有一個很大的扭轉.
我們整個人改變了.
我自己相信神跡仍然今天不斷發生.
因為神是一位有作為的神.
在約書亞記第三章就記載.
他們在等預備過約旦河.
首先約書亞吩咐他們.
你們要自潔.
要潔淨自己.
當你們潔淨的時候.
你就會看見神的作為.
在第一章第五節是這樣說的.
約書亞吩咐百姓說.
你們要自潔.
因為明天耶和華必在你們中間行奇事.
多肯定.
明天你們就會見到神在我們中間.
做一件大奇事.
各位你今天有沒有這樣的期待.
今天或者明天.
或者在很近的將來.
神就要在你的生命當中.
做一些很奇妙的事.
因為神要觸摸我們的生命.
神要改變我們整個人.
第二.
他們要讓約櫃走在前面.
當約櫃起行的時候.
一大群百姓就跟在後頭.
這個象徵神走在你和我的前面.

$^{41}$我們需要緊緊跟著神走.
今天有很多人很心急.
自己做了主意.
今天有很多人自己喜歡怎樣就怎樣.
我們沒有聽到神的帶領.
所以很多時候我們就裝扮.
所以很多時候我們就碰到很多問題.
如果你想見到神.
在我們生命裡面有作為.
我們必須要學習.
自潔.
潔淨自己.
離開罪惡.
第二.
我們認定我們是跟神走的.
不是我在這裡帶頭.
不是我自己喜歡怎樣就怎樣.
不是我自己去作主.
這個功課.
我們必須要學會.
所以我也很想鼓勵在座每一位.
希望你的生命是願意見到.
神在我的人生當中是有工作的.
讓我們願意信服神.
讓我們一起都期待神的作為.
\newpage



\section{歌羅西書 2:1-5}
\label{sec:031801}
\textbf{屬靈上做醜人}~~~~黃紹權牧師
\newline
\newline
連結: \href{http://www.church.com.hk/acms/content.asp?site=cdc\&op=show\&type=product\&code=031801\&layout=sermon}{\texttt{031801}} ~~~~ 語音日期: 2026-03-15
\newline
\newline
\hyperref[sec:031800]{\small{< < < PREV SERMON < < <}}
~
\hyperref[sec:index_chronic]{\small{[返順時目]}}
~
\hyperref[sec:031802]{\small{> > > NEXT SERMON > > >}}
\newline
\newline
歌羅西書 2:1-5
\newline
\begin{longtable}{cl}
\hline
\hline
章節 & 經文 (和合本修訂版)\\
\hline
2:1 & \begin{tabularx}{0.7\textwidth}{X} 我要你們知道，我為你們和老底嘉人，和所有沒有與我見過面的人，是何等地勤奮； \end{tabularx} \\ \\ \relax
2:2 & \begin{tabularx}{0.7\textwidth}{X} 為要使他們的心得安慰，因愛心互相聯絡，以致有從確實了解所產生的豐盛，好深知神的奧祕，就是基督； \end{tabularx} \\ \\ \relax
2:3 & \begin{tabularx}{0.7\textwidth}{X} 在他裡面蘊藏著一切智慧和知識。 \end{tabularx} \\ \\ \relax
2:4 & \begin{tabularx}{0.7\textwidth}{X} 我說這話，免得有人用花言巧語迷惑你們。 \end{tabularx} \\ \\ \relax
2:5 & \begin{tabularx}{0.7\textwidth}{X} 雖然我身體不在你們那裡，心卻與你們同在，很高興見你們循規蹈矩，對基督的信心也堅固。 \end{tabularx} \\ \\
[1ex]
\hline
\hline
\end{longtable}
$^{1}$讓我們在聆聽主的說話裡同心低頭祈禱.
父神多謝你讓我們再次回到你的啟示當中.
去明白你的心思意念.
即將愛之耶穌所成就的救恩.
也讓我們有基礎去明白這個救恩的來源.
並且讓我們能夠有一個非常穩固.
永恆不變的屬靈的生命的基石.
好使我們能夠在其中不斷地生根建造.
求你的帶領.
讓每一個聽的,說的.
都能夠同樣在今天這段訊息當中.
能夠好好地拿捏.
怎樣去作一個合神心意的神的指紋.
活在這個地上能夠真正活出一個有神生命的人.
應有的生命力.
求你的帶領.
讓每一個虛心學習的人都能夠同蒙主你的孕服.
禱告奉靠主耶穌基督得勝的名字而求.
阿們.
在公元1587年的時候.
在天主教教會其中當時作為教皇的司道五世.
Pope Sixtus V.
他在天主教的教廷裡面.
設立了一個史無前例的職位.
這個職位是甚麼呢?.
這個職位有個名.
拉丁文叫 Avocatus Diaboli.
如果正式地讓他去翻譯為英文的話.
就叫做 Devil's Advocate.
如果中文的翻譯.
就應該準確地去描述為.
魔鬼的擁護者.
當你聽到這裡的時候你一定會說.
是不是呢?.
我都說天主教是邪教.
沒理由一個教皇去設立在一個教廷裡面.
或者教會.
不過我們應該正式看天主教是一個政治體制.
在裡面要設立一個魔鬼的擁護者這個角色.
他其實真真正正做些甚麼呢?.

$^{41}$如果當大兄姐妹當你帶著一些已經預先前設的思維去看每一件事.
其實你這種的思想是很缺乏客觀性的.
司道五世他決定要設立這個職位.
最主要一個原因.
就是要達至促成信仰的純正.
你聽下去很矛盾.
他又成為一個撒旦魔鬼的擁護者.
但是又怎樣去保持信仰的純正.
因為在教廷裡面除了說到他是Avocatus Diabolus之外.
他同時有一個功能.
他這個功能是叫做Promoto Fidei.
他的意思就是說去促進一個人的信的過程的推動者.
究竟這個職位是做甚麼的呢?.
這個職位基本上就是每逢教廷提出一個神學的詮釋.
在決議之前一定要經過這個職位的監察或者審核.
他在當中的主要工作就是要將每一個神學.
他當中的意義提出相反的意見.
當他這個提出相反的意見.
就能夠促成在教會當中對這個神學的一些反思.
不過更加難能可貴的這個職位.
就是他要鼓勵這個擔當這個職位的人.
去反對到一個程度是連人家都對於他個人要進行一些反感的攻擊.
這樣就很嚴重了.
換句話來說他要做一個反對的角色.
反對到一個程度要神憎鬼厭.
神憎不憎是另一回事.
至少教廷裡面所有人都憎他的.
很個人地要攻擊他的.
你知道這個世界的人喜歡搞個人攻擊的.
針對人不針對事的.
其實很多時候我經常說我們要針對事不要針對人.
其實人本身就是針對人的.
所以哪個人做什麼那個人總之你不喜歡他.
你就全部他做的事都是錯的.
這種是一種很危險的想法.
但是卻是很普遍這個教廷所設立的角色.
就是要他做一個這樣的作用.
使得所有人憎他憎到一個地步.
不只是他講的言論連他的人都憎了.
教廷為什麼要有一個這樣的設立.

$^{81}$尤其是為什麼這位教皇是這麼特殊要去提出這一點呢.
這個教皇本身都不是一個好的教皇.
如果從歷史的角度.
但是從他做這件事的裡面.
他是想鼓勵教廷裡面.
無論任何一個神學或者甚至一個政治考慮.
都有一個反對的動力在背後.
能夠幫助教廷的立法者.
或者立神學理論的人.
能夠想多一重.
成不成功呢.
根據歷史告訴我們知道.
這位教皇自從他設立了這個職位之後.
少有人肯做這份工作.
甚至就算委任了做的人.
都往往都做不到教皇想做到的目的.
換言之.
有哪個人想做一份工要被黑人憎的.
今天講到尾的時候.
我會講多一點有關於.
擔當這份工作的人的情況.
不過從今天我們所讀的歌羅西書第二章.
第一至第五節裡面.
尤其是去到譚道在第二章第三節當中.
引用了一個字眼叫智慧知識.
尤其是智慧知識這個陳述.
其實就是好像擔當一個反對黨的角色.
或者用廣東話講的叫做做一個醜人.
他的角色在當中成為了教會.
尤其是保羅寫這封書信的目標教會.
哥羅西教會的信徒.
要擔當一個醜人的角色.
去幫助這班的信徒.
尤其是哥羅西教會的信徒.
他們有一個更練正的屬靈生命的成長.
所以今天我們的講題.
定為當一個屬靈上的做的醜人.
屬靈上怎樣去做一個醜人呢?.
Playing the devil's advocate.
我們透過這段聖經去嘗試去看看.

$^{121}$我們作為信徒有沒有一個這樣的醜人.
甚至可能我們做這個醜人.
來幫助我們自己的信能夠更加純全.
好,讓我們一起去打開這段聖經.
我們再次一起去讀一下好嗎?.
歌羅西書第二章第一至到第五節.
我們一起讀.
我願意你們曉得.
我為你們和老底家人.
並一切沒有與我親自見面的人.
是何等地盡心竭力.
要叫他們的心得安慰.
因愛心互相聯絡.
以致豐豐足足在悟性中有充足的信心.
使他們真知神的傲備.
就是基督所積存的一切智慧知識.
都在他裡面藏著.
我說這話.
免得有人用花言巧語迷惑你們.
我甚至須與你們相離.
心卻與你們同在.
見你們循規蹈矩.
基督的心也堅固.
我就歡喜了.
我們提一下前文我們說些什麼.
前文我們說過.
事奉是一件會叫人充滿又關心.
但是又憂心的事情.
但是在這個憂心的過程當中又滿有喜樂的.
因為什麼呢?.
因為你願意所付出的擔憂,憂心.
或者甚至困擾.
其實都是為了一個更美好的結果.
我們用了一個例子.
就好像父母一樣.
為了養育自己的子女.
完全不會否定或者不願意犧牲自己的睡眠時間.
去為了剛剛初生的嬰孩.
半夜三更來餵奶.
半夜三更來換紙尿布.

$^{161}$尿布現在流行尿布.
環保嘛.
不過我見很多不斷說環保的又不願意.
用紙尿片.
但是無論如何你都肯要捱更底夜.
去成就你的兒女.
由嬰孩一步一步成長.
這一份的過程.
而這個過程最終帶來的喜樂.
就是你看到你的下一代.
能夠真真正正成長起來.
成熟起來.
作為一個牧者其實同樣都有這份的心.
作為一個信徒.
也都應該有這份的心志.
今天你是一個初信徒.
你的前一代的前輩.
就會願意希望幫助你成熟一點.
而你作為一個後輩.
你將會成為下一代的信徒的支援的時候.
你亦都同樣是需要付出這份的心思意念.
因為你看到下一代的喜樂.
就是因為他成長起來.
他真真正正能夠在信的上面站得穩.
這個就是我們上文所說到的所謂.
bitter joy.
苦中的喜樂.
那個的意義.
所以更勝過你說教會有幾千萬的建築物.
有成千上萬的人湧進來.
一個成熟的信徒能夠在上帝裡面.
能夠可以真是堅強地站立在上帝的當中.
配得起在上帝面前活著的人.
連天上的天使都同樣滿有快樂.
這個全部是聖經裡面的啟示.
但來到今天我們所看這段的經文.
我們在上個星期已經講過的.
上一段的經文從第一章第24節.
其實和來到我們今天所看的二章一至五節.
本來是一個整體的結構.

$^{201}$上文我們所看到的東西.
其實來到下文裡面卻是要進行一個成就的.
過程的一個另一面的介紹.
就好像一個小銀一樣.
它有字的畫面.
也都有一個黃的肖像的另一面.
即是有公仔有頭.
兩者並存在同一個銀上面.
所以今天我們來到看這個第二個部分.
我們會分開為三部分一起去看的.
第一部分我們去看的是二章第一節.
二章第一節首先保羅講到.
如果要成就這個工作的.
能夠要使到下一代能夠更加成長的.
成熟的.
你必須要有一個很準確的焦點來去事奉.
所以你一定要調整清楚.
你事奉的焦點.
第二章第一節這樣說.
「我願意你們曉得.
我為你們和老底家人.
並一切沒有與我親自見面的人.
是何等地盡心竭力」.
從第二章第一節的開始我們看見.
是承接上文保羅講及到他怎樣為到.
哥羅西教會甚至現在在這段聖經說了.
他附近的另一個雙子城老底家教會.
他怎樣付出他的勞苦.
我們再次重申.
保羅向教會去提出他的勞苦或者他的付出.
主要的目的是希望教會能夠明白到.
保羅付出的種種.
去面對這些困難或者不幸難堪種種的程度有多少.
不是去專注說我保羅為你付出.
焦點不是人.
而是那個人願意付出的程度是多少.
用父母養育兒女的例子.
重點是要兒女明白父母所付出的程度.
不只是專注父母這個人物.
當然你明白程度有多少.

$^{241}$你可能對父母的尊重會更多.
但這不是主旨.
我相信沒有一個父母.
一個好的父母.
為兒女付出不是為回報.
如果你是為回報才為子女付出的話.
我可以告訴你.
你失望的機會很高.
首先其中一個因素是.
他明不明白你付出的程度是怎樣.
明白了之後.
他會不會因為這個程度而感動到他對父母的尊重.
是很複雜的過程.
在今天這個個人主義高舉的狀態.
一般人領受人的祝福是無限的喜悅.
但領了之後就覺得各走各路.
因為英文有一個很好的字眼.
特別在電腦界.
這個叫leisure.
認不認識這個字.
leisure就是吸血蟲.
現在這個世界的人.
你會發覺他的思想就是.
吸了你血之後就離開了.
好像蚊一樣吸了你的血.
他就走了話之裡.
保羅在前文講到.
他專注希望哥羅西教會明白到.
我為你付出這一切的程度有多少.
希望你能夠了解這份程度.
不是來報答保羅.
而是你們曉得怎樣將來.
同樣付出類似或者更多這個程度.
去為下一代的成熟去努力.
這個是前文所提到的.
所以來到這一段聖經第二章第一節.
保羅在這裡最重要的是.
希望哥羅西教會他們明白這一點.
所以保羅在這裡指出他付出的對象.
包括了老弟家教會.

$^{281}$其實相傳在寫歌羅西書的時候.
其實同樣保羅也都寫過一封書信.
是給老弟家教會的.
當然在我們的正典裡面是沒有.
沒有是什麼原因呢.
有些人說是不是聖經不整全呢.
不是.
而這個通常往往聖經沒有正典.
沒有的書卷的目的.
就是因為這卷書卷已經涵蓋了.
另外一些的書卷或者書信.
是遺失了的內容.
所以我們不需要看兩卷一模一樣的書信.
所以上帝的啟示.
都有他自己本身的特色和生命力.
保羅一方面點出他所關心的對象.
不是哥羅西教會一間.
是包括了老弟家教會.
尤其是當我們看見保羅.
根本上和哥羅西教會.
以及老弟家教會都素未謀面的時候.
所以保羅根本就不是要透過這封書信.
希望搭關係然後怎麼樣.
你們將來可以報答我.
為你們的付出.
而是將這個栽培下一代.
使教會成長的棒.
交給你們繼續去承擔.
保羅沒有辦法會再去哥羅西.
保羅再沒有機會可以去到老弟家.
承接栽培這兩個地方的信徒的棒.
是交給你們去做了.
所以當你要去曉得怎樣做的時候.
你就要明白作為一個前輩的保羅.
他怎樣去付出.
程度有多少.
你們就去做至少有保羅這種的心態.
是不是會要求更多呢.
這個就看上帝在你們每個人裡面的心意了.
所以保羅在這裡講到.

$^{321}$再看下一句.
並一切沒有與我親自見過面的人.
這一句話就更加清楚顯示出保羅.
他根本上不是要你們認識我保羅.
然後回饋保羅我.
有些我不認識的我見都沒見過的.
但是這個教導這份心志.
是要你們曉得.
所以為什麼第一節就講到.
我願意你們曉得.
曉得什麼這個就很重要.
不是曉得保羅.
而曉得保羅的事奉的委身.
你們都要有同樣的這份委身.
因為這份委身的最大的特色.
就是盡心竭力.
沒有保留.
弟兄姊妹.
我們很多時候.
作為一個基督徒.
其實怎樣認證你是一個基督徒呢.
或者讓自己知道自己是一個基督徒.
其實很簡單.
當你看到其他信徒的時候.
你會不會盡心竭力去為他的成長.
靈性的好處而去付出呢.
在我信主以及事奉的裡面.
都不斷有一些這樣的弟兄姊妹在我身邊.
他們這些先賢能夠給了一個很好的見證和學習給我.
尤其是我在事奉主的這幾十年的裡面.
老實這麼說.
在我事奉的當中.
這二十幾三十年的當中.
就是全職事奉的時候.
在我身邊不少的弟兄姊妹願意付出.
有弟兄姊妹.
她說牧師.
通常我們叫名字.
因為很熟朋友.
Ellen.

$^{361}$你有什麼需要幫忙.
我可以幫你.
我就跟她說.
可不可以幫我講道.
她說我不懂得講道.
我說你讀過神學.
不懂得講道都不要緊.
那你怎樣幫我.
她說我懂得打字.
我說那行.
我寫好一些講道.
我寫好一些主要學課程.
你幫我打字.
一打打了幾十年.
如果斷編計斷字計.
好像現在幾千一個字.
我記得當年我在香港做出版的時候.
去找人打字.
五元一百字.
你寫篇文章.
那篇講道尤其是上個禮拜.
講了一個半小時.
那篇講章.
你看看你現在找AI拔出來.
那裡成萬字的.
一篇都幾和味.
如果真是你去打字賺錢.
有弟兄他願意為我打字.
打了幾十年.
又有弟兄願意在神學上面.
學到的東西傾盆講給我聽.
然後幫助我.
能夠有很多的討論.
能夠將我的神學.
能夠磨得更加的練靜和清潔.
撇除了很多私人的意見.
你知道在讀神學.
甚至基督徒的成長過程當中.
最大的一個問題是什麼.
就是信徒自高自大.

$^{401}$我在香港的時候.
回教會參加一個主日學班.
講系統神學.
系統神學我本身懂的.
不過我上那個主日學.
其實因為懂不代表我不想.
不懂更加要想.
懂也想的.
為什麼.
因為想聽多些.
使自己更加練靜.
我記得那個牧者一進來.
在主日學班裡面第一句.
大家好我們當中有多少個神學家.
因為那間教會很大.
我以為在身邊是不是真的有些神學家.
他說哪位是神學家.
他說你們就是神學家.
個個都一家的.
個個都認為自己對的那一套.
基督徒最大的問題.
認為自己的信仰就是那一套就對了.
沒有客觀性上帝的說話才對.
我們是可以錯的.
幾十年來.
我事奉的身邊都有一個這樣的弟兄.
陪著我去糾正.
或者跟我討論有關於神學的問題.
一毫子都沒有付出過.
弟兄姊妹.
我作為後輩的.
又怎樣去付出呢學識.
保羅在這裡就在說這件事.
看到嗎.
聖經不是很難的.
在這裡說的.
你有沒有盡心竭力地去接待你的弟兄姊妹.
那些不對的信徒或者邪惡的牧者.
就利用這段聖經說.
所以那些牧者很勞苦的.

$^{441}$扮得自己很辛苦或者將自己去塑造成一個很苦的傳道人.
然後讓那些信徒來騙他們那些憐憫之心.
但保羅沒有這樣做.
保羅在這裡說.
我要你們曉得的是.
我怎樣竭力盡心為你所付出的東西.
而你們要去做的不是回饋保羅.
而是你們多去竭力盡心地去幫助下一代的信徒.
所以事奉這件事.
我們要有一個很準確的目標.
弟兄姊妹.
我們作為信徒.
不是說因為我們去事奉主.
因為耶穌是一個大能的神.
所以有祂撐腰.
然後我們事奉的事情每件事都很高興很開心.
那就真正是事奉了.
不開心的那些就不是事奉了.
我曾經聽過在信仰的路程當中有一個弟兄說.
他說什麼去斷定我所做的事奉是上帝合神心意的呢?.
他說只要那件事我做了之後.
我心裡覺得開心那就是了.
我心想如果你不開心是不是就不是上帝的旨意呢?.
那保羅在這裡怎麼說?.
完全保羅說的是相反的.
保羅說我現在為你們受苦.
不過他開心.
那個喜樂的泉源.
不是因為我個人覺得開心不開心.
那個喜樂的泉源是因為上帝的目標.
怎樣能夠得以成就在宗教會身上.
包括老弟家教會.
包括哥羅西教會.
以及眾多保羅都沒見過的人.
他的成長就成為保羅的喜樂.
讓我講一句話.
現在我們教會全面搞這個所謂網絡的崇拜.
有很多大型節目我不認識的.
其實老實說.
我很想至少知道他們的名字.

$^{481}$如果你們聽到的.
或者你介紹一下你的名字在哪裡.
你告訴我.
讓我知道原來在遠方有一位叫什麼名字的弟兄姊妹.
原來在這裡是努力追求著事奉著.
我不需要一定跟你素未謀面.
我不需要一定要跟你見面.
當然有機會見面是一件開心的事.
過去我剛剛去年回香港已經見了幾個弟兄姊妹.
很開心.
我真的很開心.
原來他們是追求的弟兄姊妹.
他們成熟了.
在上帝的裡面我們是努力往上帝的方向去走.
這個就是我們最大的喜樂.
但是在今天的教會裡面是怎樣.
而今天很多教會的信徒的事奉是什麼.
我開不開心.
你決定就是等於上帝的旨意.
你開心不一定是上帝的旨意.
而且很多時候我們開心不是上帝的旨意.
我可以這麼嚴肅的說.
但是如果你的喜樂.
甚至你的憂慮.
你的困難.
你的面對的很多困擾.
用保羅說我為你受很多的苦.
其實都是為你的成長.
那份真正的喜樂才是真正在上帝裡面.
那份事奉的實質.
另外事奉也不是抱著一種小眾而去做的事.
我說清楚這句話.
事奉不是為了一群小眾才去做的事.
看保羅這裡說.
他事奉的神所涵括的人群有多少.
哥羅西教會.
老弟家教會.
還有很多沒見過的人.
我們教會在這麼多年裡面都做網上轉播或者直播.
我很記得當年我請了一些牧者來.

$^{521}$我請他們的時候就在同工會的聚會當中.
我說我很想請你們來一起.
在我們馬來西亞這個地方一起事奉.
就在我作出這個申述的時候.
曾經有一個我請過來我們教會講道的牧者.
他以後不來了.
你不用擔心他以後不會來了.
他為什麼不會來了.
因為他馬上就和當眾所有牧者這樣說.
如果艾倫和牧師請你去他教會講道.
你們要想過度過才好.
他的講道全部擺上網的.
一字不留擺上網的.
不會剪輯的.
你聽聽我講道沒有剪輯的.
他上面的網上的轉播或者直播.
甚至曾經有信徒聽完之後.
說牧師你這麼直接講這番話.
得罪很多人.
不如我建議你剪了他好不好.
我說不要留在這裡.
那位牧者他說.
因為他講的道是全世界的人都會聽到的.
但他講的道不是我.
他講的道呢.
就只是講給這班信徒的教會的信徒知道.
當他講這番話的時候.
我就要明白.
我就會明白了.
所以我嫁女兒也不請他來喝酒.
就是這個原因.
你明白嗎.
我和他沒有團契.
因為他的講道只是針對李起悅講的東西來講給你聽.
他事奉的對象是一個小眾來的.
只是一小撮的人.
他事奉的對準的目標只是一小撮的人.
本來不是的.
他講的東西.
你拿著歌羅西書.

$^{561}$派給街邊所有的教會的信徒.
他都可以給人聽的.
雖然中文有一句說話.
我不是很想經常用的.
因為這些比較世俗一點.
或者這些叫傳統.
不過有一句說話的確是事實來的.
一個誠實的人.
是事無不可對人言.
我只能夠跟你講.
我不能夠跟他講的這番說話.
那我就會問.
那這番說話為什麼只有我聽得.
其他人不聽得呢.
是不是你想想我在我身上有什麼套取呢.
還是這些只是私語呢.
私人需要知的說話.
一會兒講的.
什麼私人的說話講呢.
我們一會兒看.
不過至少在這一節裡面.
清楚給我們看見.
保羅在事奉的裡面.
不是只是招呼或者對待某一個群體.
希望籠絡他們人心.
不是的.
如果是這樣的話.
不如這樣.
這番說話只是給哥羅西教會聽的.
我寫給老弟家的教會那封信.
另外給老弟家聽的.
或者給那些其他沒見過的受眾.
我寫給另一篇的.
比較普及一點的.
不要有這麼仔細的內容.
有沒有這樣做呢.
沒有.
事奉的焦點很清楚.
是上帝委託的東西.
他盡心去做.

$^{601}$不是見人講人話.
見鬼講鬼話.
保羅作為一個使徒的身份.
他就要做這個見證.
去見證出.
他這個普及性的事奉的心智.
以及他的事奉的行動.
做的事奉.
不是只是討人喜悅.
或者討自己教會會有的喜悅.
而是討上帝的喜悅.
因為這個是要承擔上帝所交託的工作.
第二點.
第二第三節.
在這裡.
讓教會有一個準確的連繫點.
保羅在這裡怎樣說呢.
我們一起看看.
第二節和第三節.
要叫他們的心得安慰.
因愛心互相聯絡.
以致豐豐足足在.
悟性中有充足的信心.
使他們真知神的奧秘.
就是基督.
所積蓄的一切.
智慧,知識.
都在他裡面藏著.
在這段聖經裡.
保羅說到.
教,牧養的其中一個要旨.
就是要讓教會.
能夠連繫在一個點上面.
在上文第一章的二十四節之後的經文.
有已經說過.
我們要連繫在基督的身上.
而來到這一節的經文.
保羅更加進一步地.
去清清楚楚地說到.
那個連繫點就是基督.

$^{641}$我們仔細去看一下.
這裡想說些什麼.
我曾經聽過一個教會.
他們在教會裡面.
常常有很多的爭端.
教會這麼小嗎?.
每一次教會有爭端.
譬如教會建不建停車場.
要不要重建飯堂.
要不要裝修哪一個樓層的廁所.
或者牧者的宿舍要怎麼裝修維修.
每一次每一樣東西.
都有人去有爭端.
其實已經反映到這間教會的信徒的狀態.
不過都不要緊.
每一次教會的長輩都會出來這樣說.
教會要為了合一.
不要再有爭端了.
聽起來很順耳.
聽起來好像很合情合理.
合一嘛.
不要爭端了.
合一在哪裡呢?.
不說.
往往合一的點就是.
縮小一點的規模.
就是錢銀的問題.
第一個節目.
信徒在教會裡面的聯合點.
不是一些什麼什麼金錢的問題.
而是我們對上帝的心思.
是否能夠有清楚的拿捏.
在這樣的過程裡面.
二章第二至第三節裡面.
保羅就清楚在這裡陳明這個點了.
他說.
保羅的牧羊的事奉的結果.
是一個眾多的群體.
有很多的人.
所以要叫他們.

$^{681}$不是某一個.
眾多的群體做什麼呢?.
促成他們連繫在基督的上面.
不過牧師不是這樣.
我讀到這段聖經.
他說.
要叫你們的心得安慰.
因為愛心互相聯絡.
要說聯絡在愛的上面.
沒錯.
愛是其中一個.
他們的聯繫的接著劑.
聽清楚.
愛是其中一個接著劑.
有沒有用過膠水?.
用過膠水是嗎?.
重點不是膠水.
重點是你要用膠水聯合兩樣東西.
譬如寄封信.
有沒有試過寄信?.
那張郵票.
貼錯了信封再撕出來不能貼了.
那怎麼辦?.
找一支膠水出來.
找一支漿糊.
找一支UHU出來.
懂不懂?.
塗在郵票上.
然後找一個合適的信封.
貼在信封上.
重點不是膠水.
重點是那兩樣東西的連接.
是什麼東西?.
好了.
除此之外.
還有什麼?.
是要豐豐富足在悟性當中有充足的信心.
牧師不是基督.
是講信心.
是不是?.

$^{721}$豐豐富足的悟性中.
不是.
豐豐富足的悟性中有充足的信心.
這裡所談論到的是什麼?.
是貼的過程.
你是用信來貼上去的.
貼的時候要有愛心.
要用準確的膠水.
貼的過程中要怎樣貼?.
你要有豐豐富足的悟性.
懂得怎樣去貼.
然後怎樣?.
貼上去.
明白.
上帝的奧秘.
讓他知道上帝的奧秘.
奧秘是什麼?.
很熟悉.
牧師這個字眼好像剛剛聽完你說過.
奧秘這個陳述.
我用了一個例子來解釋.
就好像上帝和你去傾密計.
將一些以往上帝沒有公開講到的東西.
私底下跟你講清楚.
介紹清楚.
那些東西就只是講給我聽.
講了給你聽是做什麼?.
算數嗎?.
只是你知道嗎?.
如果只是你知道.
那就是籠絡你的心.
不是.
原來這裡用了一個句式叫.
就是基督.
原來這個奧秘最終要講的就是.
基督是怎樣.
耶穌是怎樣.
他的工作是什麼?.
請留意.
這裡特別用基督.

$^{761}$不是基督耶穌.
也不是講耶穌.
是基督.
基督就很不同.
有一個很特別的意義.
就是被高納者.
你要和所有的人去介紹.
明白這位被高納者是哪一位.
上帝講了給保羅知道.
就是被召的人.
當你在這個被召裡面.
接著下去你所做的事.
就是要怎樣?.
將所有你所認識到的東西.
所以第三節.
所積蓄的一切智慧知識.
都在他裡面藏著.
不是叫他收起來.
是叫他們能夠得著.
而且很穩固地收藏寶貴.
放在這班人身上.
哪班人身上?.
前面講完.
老弟家教會.
哥羅西教會.
和一切沒見過的.
都要他們學會明白.
得到這個知識智慧.
讓他們好好保護著.
有這個心態.
有這方面的學知.
在《聖經》裡面有一句說話.
《申命記》第29章第29節.
我們打開去看看.
《申命記》第29章第29節.
找到之後我們一起讀好不好?.
好,我們一起讀.
好,叫我們遵行者律法上的一切話.
弟兄姊妹聽不聽到這句話?.
現在上帝將一個隱蔽的事.

$^{801}$告訴了保羅.
而這個保羅就告訴.
所有老弟家.
哥羅西.
以及一切他沒見過的人知道.
原來就是基督.
將一件原本隱藏的事.
將它顯明出來.
其實都不是保羅顯明出來.
基督自己本身顯明了出來.
保羅只不過是將這個奧秘.
和基督連起來.
告訴你們知道.
現在的事情再不是神隱蔽的事.
現在是顯明的事.
基督顯明在你們身上.
基督在做什麼?.
成就救恩.
怎樣成就救恩?.
基督怎樣成就救恩?.
死在十字架上.
什麼叫死在十字架上?.
用回聖經這裡的陳述.
是什麼字眼?.
好像聖經問答比賽.
不是,很簡單的.
那句話剛才我說了很多次.
就是第一節最後那句.
盡心竭力.
死到懸疑.
那你們是不是死到懸疑呢?.
而這件事現在已經顯明了.
顯明的事是屬於誰的?.
是你和我.
將它藏在上帝裡面.
是我們和我們的子孫.
我們要將這個盡心竭力.
使人連繫在基督身上的點.
將它清楚講述給人知道.
我自己也努力去締造.

$^{841}$而且要長傳.
一直去到一代一代傳下去.
在這樣的過程當中.
這個就是我們整個的.
事奉的核心.
還有在這段聖經裡面特別強調了.
積蓄的一切智慧知識.
特別是講知識這個字眼.
保羅當他談論到.
智慧源自於神之後.
智慧用來做什麼?.
大家看節目有沒有想過?.
你有沒有想過.
你用智慧來做什麼?.
我很想有上帝給我智慧.
很多教會.
我經常去教會跟信徒聊天.
我想牧師你為我祈求上帝給我智慧.
你有沒有想過.
你祈求上帝的智慧用來做什麼?.
他們求上帝的智慧.
其實最終的目的就是想掠水.
最好炒股票.
知道哪隻股票會升.
炒哪隻外幣.
那隻外幣會升.
找哪份工作.
那份工作是賺到錢的.
給我智慧.
讓我在那份工作或者做生意.
讓我有智慧掠水.
是不是這樣?.
不是.
智慧和知識連繫.
這裡保羅說的不是我說的.
Have the wisdom with the knowledge.
什麼knowledge?.
Oh, knowledge.
知識.
原來在創世記裡面.

$^{881}$在上帝創造天地的時候.
在伊甸園創世記第二章十九節當中.
記載了有一棵樹.
這棵樹就叫做.
Tree of knowledge of good and evil.
分辨善惡樹.
我們求上帝智慧.
是曉得讓我們分辨善惡的這個知識.
不是求智慧.
所以我們可以去生意.
得到更大的發圍.
工作能夠升官進爵.
財技精湛.
使你賺到盤滿缽滿.
不是.
是要去分辨善惡.
為什麼?.
上個星期我們說過.
伊甸園是一個什麼地方?.
伊甸園是上帝和人的.
私密.
相會的地方.
在裡面是講密集的地方.
所以在伊甸園裡面.
即或人犯了罪.
躲在那裡.
又去弄衣服.
又去怕死.
上帝都仍然在伊甸園裡面.
走到伊甸園裡面.
然後問.
人啊.
你在哪裡啊?.
頂枝梅.
是一個私密討論的空間.
私密討論的空間.
對不起.
我只能用英文.
我不懂中文.
因為在裡面你學會的.

$^{921}$是要分辨善惡.
你現在由一個狀態.
由上帝給你的私密地方.
你學不會分辨善惡.
保羅就要告訴你.
你學會了.
有上帝的智慧.
可以回到原本上帝創造的.
私密交通的地方.
我再說一次.
私密不是穿的那種私密.
是很秘密的.
隱蔽的.
Close the room.
只有你和上帝一起.
然後去和上帝.
將真正學會分辨善惡的能力.
擁有回來.
保羅在這裡所陳述的.
就是告訴我們知道.
怎樣去重歸於這個的狀態.
在這番種種的說話裡面.
讓我們看到.
保羅真是很刻進記心.
怎樣去讓教會回歸.
回到上帝創造的那個樣式.
當我們能夠分辨善惡的時候.
我們就能夠去怎樣知道.
哪些應該做和不應該做的事.
或者有些事.
作為一個弟兄姊妹.
你是先前的前輩.
怎樣栽培後信呢.
你有些說話.
你就一定要說的了.
你會做醜人的.
這個就在接下來第三點上面.
第四至第五節的經文.
就會說給我們知道.
我一會兒會看的.

$^{961}$不過從這段聖經裡面.
讓我們去看到.
人獲得神的智慧最重要的基礎.
就是能夠怎樣和基督再次連繫.
當你和基督連繫的時候.
你就能夠打通你和上帝.
可以溝通.
明白他奧秘的情況.
保羅就在這裡做了一個見證.
當你明白了這個奧秘之後.
不只是自己藏起來.
而是讓裡面眾人一切都知道.
怎樣去和基督連結.
然後你們一同將這個這麼寶貴的東西.
好好擁抱著他.
分享給其他人知道.
所以我們看看上帝做事真的很美麗.
拯救人.
人在哪一個點那裡曾經跌倒.
你就在分辨善惡那裡做不到的時候.
上帝就讓你在那個點重新站起來.
今天的教會.
有沒有真的幫助信徒和基督連結起來呢?.
還是連結在其他的事情上.
好像我剛才所說的那間教會.
什麼都是用錢來衡量的呢?.
最後大家所謂的合一.
都是因為那個預算不用花那麼多錢.
大家就開心一點.
是不是這樣呢?.
今天甚至很多人的關係.
也都是連結在一些利益的身上.
頂替支援.
因為我們根本喪失了.
分辨善惡的能力.
不過很可惜的就是今天.
很多的教會卻是將.
教會變成一個.
好像派屬靈福利的福利所一樣.
或者讓他們來到這個廉價的.

$^{1001}$所謂救恩的裡面去購買.
有時候我真的覺得.
今天的教會就好像.
一個屬靈的義烏.
知不知道什麼叫義烏?.
你知道全世界那些平價產品.
有一個世界舉足輕重的城市叫義烏.
專賣切東西.
你說不知道.
回香港.
一下飛機.
一出了飛機.
第一個大招牌就是這樣.
一定見到的.
歡迎來到義烏市.
買世界產品.
專賣切東西.
教會就是賣平價屬靈的義烏.
教會是不是真的帶你去和基督連結呢?.
還是你只是來到.
只是做一個購買廉價救恩的地方呢?.
好了 最後一點.
第四節到第五節.
我們一起讀一下好不好?.
第四節到第五節.
我說者說.
免得有人用花言巧語迷惑你們.
我生只需與你們相離.
心卻與你們同在.
見你們循規蹈矩.
信基督的心也堅固.
我就歡喜了.
來到第四到第五節.
我給一個副題.
它叫做做一個屬靈的醜人.
或者叫做.
不要怕做這個屬靈的醜人.
我們看回.
天使聖經想說什麼.
他說.

$^{1041}$我說者或者說前文講了一大堆的說話.
目的是免得有人用花言巧語迷惑你們.
在哥羅西教會.
什麼人講的花言巧語.
其實我上一堂已經講過了.
在哥羅西教會裡面.
有很多人去傳揚其他的信仰模式.
而且那些人口才是非常了得的.
他用很多人的傳統.
來引誘你們走下去.
遲些會繼續講的.
使得你們以為你很虔誠.
實質你不虔誠.
這些人很善用於言語.
花言巧語.
來迷惑你們.
以前我在香港.
我只看中醫.
很少看西醫.
看西醫一定很嚴重.
馬上要找藥.
不過我多數看中醫.
我在香港的時候.
有一個中醫.
我媽媽介紹我的.
她說你去看這位中醫.
我還記得他的名字.
叫做冼冠希.
找不到他就死了.
因為我見他的時候已經七十幾八歲.
已經三十幾年前.
除非他是人歲.
我去看他醫生.
一定要等很久.
因為很多人去看.
這個醫生有個特色.
你去看他的時候.
他問得很詳細.
他要你回答很多問題.
我很記得.

$^{1081}$我說冼醫生.
我在這裡這樣這樣.
接著他說.
你是什麼病.
我把脈把到.
但我想知道你為什麼這樣.
接著我就解釋.
我很忙.
我怎樣怎樣.
那時候年輕.
二十幾歲.
接著說了一大輪之後.
你知不知道他怎樣說.
考詞.
我說什麼叫考詞.
我真的這樣這樣.
他再說又打斷我的說話.
花詞.
考詞又斷花詞.
很重要.
我再說多兩次之後.
他就說考詞花詞.
他的聲音不小.
全國.
他的店鋪.
他的店鋪我告訴大家.
他在樓梯底擺攤.
排隊排到出門口.
等他等一個小時.
因為他問很多問題.
全部人都聽.
當他這麼大聲.
全部人都看他.
看過來.
這個人說了這麼多.
原來醫師說他.
你知不知道什麼叫考詞花詞.
即是辯駁.
是藉口.
不要相信.

$^{1121}$每次我去看醫生.
他都是這樣問.
我發覺.
原來不只是我.
個個都這樣問.
個個都這樣回應.
其實他想說什麼.
我和我媽說.
我很討厭這個醫生.
為什麼.
搞到我沒面子.
個個都是這樣.
個個都說花言巧語.
說我們講的是藉口.
我真的.
你知我幾時上班.
幾點開工幾點收工.
我早上七點鐘要回到辦公室.
晚上一點鐘才從辦公室門口走出來.
從銅鑼灣乘小巴.
亡命小巴.
那裡全部街邊都沒有人.
他說我考詞.
我自己不休息睡覺不夠.
我明明是這樣.
我媽就和我這樣說.
個個都這樣說.
其實他想提醒你們.
不要再這樣.
他當然知道.
你們沒辦法逃脫.
這樣的生活.
希望你們自己想想.
是不是很討厭他.
我真的很討厭他.
這些藥真的很好.
師母好像也有份看.
看暗瘡.
不知看什麼.
又可以.

$^{1161}$個個看都好.
不過死了.
不知為什麼可以的死了.
不行的現在健在.
不知為什麼.
不過anyway.
看回這段聖經.
波羅在這裡所講的這番說話.
很明顯地.
去指出.
在你們當中裡面.
有一些人很善用花詞巧語.
去引誘你們.
使到你們.
離開.
離開什麼.
離開你們.
原本所守的.
規矩.
使到你們的心.
不在主的裡面.
離開和基督的關連.
所以我前文所講的.
這一大堆東西是提醒你們.
這班人講的東西.
是想你離開基督.
而保羅在這裡.
就開張明義.
講給.
哥羅西教會知.
不要聽這些人講.
使到你們.
有智慧去辨別.
到他們講的東西.
哪裡講.
前文第三節剛剛講完.
「知識」這個字眼.
「分辨善惡術」.
就是這個字眼.
你們要學懂怎樣去分辨它.

$^{1201}$保羅講到這番說話.
其實是很大膽的.
因為為什麼.
在這個世界上有哪個人.
喜歡聽信義的說話.
阿諛你的說話.
捧你的說話.
保羅你又不在.
又不存在.
遠水不能夠近火.
又不能夠聽著這班人.
你這麼厲害.
你最好在現場跟我辯駁.
你又不在.
但保羅告訴這班.
哥羅西教會的信徒.
不要聽他們這樣講.
你要學懂辨別.
不是哪個人在不在.
不是我保羅在不在.
是你們要懂得怎樣去辨別.
不過保羅.
我雖然真的不在你身邊.
是和你們身體肉身上相離.
但我的心是沒有離開過你們.
當保羅縱使在這樣的情況之下.
他聽到從.
哥羅西教會.
的以巴弗傳來的消息.
原來哥羅西教會的人.
仍然是.
循規蹈矩.
堅信基督.
他就滿心歡喜.
弟兄姊妹.
醜人不怕做.
最重要做了之後.
你的目標.
能夠兼顧你下一代的.
能夠成長.

$^{1241}$那你就心滿意足了.
做父母就是一樣.
你把子女養到那麼大.
捱著班底夜.
有時又要激氣.
叫他做事他又不做.
煮好飯他又不來吃.
吃的時候又談三談四.
不要緊的.
最重要他吃.
吃了之後有力.
他繼續走他人生路.
吃了之後你看到他有智慧.
他經歷過一些事情他自己好像成熟.
你就滿心歡喜了.
保羅就是這樣的情況.
他不怕.
哥羅西教會的人誤會他.
因為當這封信寫給哥羅西教會的人.
連這班花言巧語的人都會聽到.
這些人會不會更加下重藥.
更加花言巧語.
去欺騙哥羅西教會的人.
不奇怪的.
但是保羅最安心的是.
你們仍然守在神的律例典章裡面.
包括你的道德和你的信.
都在基督當中的時候.
他就見到這份成熟的心.
保羅知道自己的受難.
或者自己的受苦.
或者受困厄.
是需要付代價的.
但是他也都看到他的目標很清楚.
能夠幫助到哥羅西教會.
能夠成長.
辨別到這些虛情假意.
或者花言巧語的人.
所以他不怕冒險去做這個熟稜醜人.
來到最後我想說.

$^{1281}$我們作為牧者.
或者作為信徒.
一個前輩的信徒.
你想一下自己信了主.
那一年是哪一年.
我不太記得了.
受浸那一年呢.
都十年八載之前的事了.
比較一下今天你的靈性.
你有沒有幫助到你.
後過你信主的人成長.
或者連你自己都一鑊泡.
在這個世界上.
做牧者.
你如果想知道他是一個好牧者.
或者是一個花言巧語的牧者.
很簡單的.
除了他講道的悉經如何之外.
你看他怎樣對待信徒你就會知道.
如果明明他見到信徒做錯事.
他不出聲.
或者他甚至扭橫折曲地去稱讚他們的話.
那你可以離開那間教會的那個牧者了.
真的.
因為他只會施更多毒害在你身上.
而他從來沒有想糾正你.
我年輕的時候都挺頑皮的.
所以我就經常吃一味菜.
叫藤條炆豬肉.
不過我覺得這些藤條炆豬肉吃完之後.
又挺有意思的.
我媽媽一定這樣說.
或者我爸爸一定這樣說.
他說你知不知道你自己錯了什麼.
你的答案不是知和不知道.
你的答案是我知道我錯了什麼.
因為知和不知道可以敷衍你.
你如果知道自己錯了什麼.
我爸爸媽媽的藤條炆豬肉就有意義了.
保羅對於哥羅西教會其實真的很愛護.

$^{1321}$他和他們都沒有相交過.
沒有認識過.
沒有直接聯絡過.
但是在這封信裡面充分看到.
他只需要在上帝的國裡面屬神的子民.
他從來沒有放過他模樣.
他都想他好.
在研究歌羅西書裡面有一個研究.
保羅的仕途身份究竟是一個universal的身份.
還是一個local church的身份.
就是一個本地教會的仕途身份.
或者是一個普世的仕途的身份呢.
我可以給一個很準確的答案給大家知道.
保羅的仕途身份是一個普世性的身份.
同樣我作為一個牧者.
我牧養的不只是以馬來尼環浸信會的弟兄姊妹.
是 你們可以和我有很多的相交.
更加直接.
但是即使你不是和我直接有相交的弟兄姊妹.
我仍然是會竭力去幫助你.
在食靈生命上面成長起來.
求主幫助我們.
讓這份心智不是我作為牧者有.
不只是保羅有.
我都要承接保羅這份心智.
你有沒有這份心智.
這個是很重要的.
你不要說我不是牧師.
我不是傳道人.
我不用.
不是.
你因為是一個給其他新信的信徒.
早信的.
你要幫助他們.
和你同輩一起信主的.
你和他一起走這條道路.
甚至可能那些給你早信主的.
不過他們有錯誤的.
不怕做屬靈醜人.
告訴他.

$^{1361}$錯在哪裡.
希望大家能夠在我們的信徒上面.
能夠真的做到有回歸創世的伊甸園裡面.
那棵分辨善惡樹原本有.
那個本身的.
屬靈的特質.
讓我一起和心低頭禱告.
父神多謝你給我們一同在裡面仰望.
看見保羅.
不害怕成為一個屬靈上面的醜人.
去成就信徒.
能夠更加成熟地.
曉得辨別豬靈.
有這份重歸創世記裡面.
伊甸園當中的分辨善惡樹的應有的屬性.
求你的帶領幫助.
讓你的靈.
所賦予我們的智慧.
曉得怎樣連接在基督的奧秘的身上.
好使我們能夠實踐.
出上帝你每一個的啟示.
和你所賦予我們真正.
聖徒應有的樣式.
禱告奉靠主耶穌基督.
得勝的名字而求.
阿們.
\newpage



\section{使徒行傳 15:1-21}
\label{sec:031802}
\textbf{當紛爭再起時}~~~~禤凱榮牧師
\newline
\newline
連結: \href{http://www.church.com.hk/acms/content.asp?site=cdc\&op=show\&type=product\&code=031802\&layout=sermon}{\texttt{031802}} ~~~~ 語音日期: 2026-03-15
\newline
\newline
\hyperref[sec:031801]{\small{< < < PREV SERMON < < <}}
~
\hyperref[sec:index_chronic]{\small{[返順時目]}}
~
\hyperref[sec:031803]{\small{> > > NEXT SERMON > > >}}
\newline
\newline
使徒行傳 15:1-21
\newline
\begin{longtable}{cl}
\hline
\hline
章節 & 經文 (和合本修訂版)\\
\hline
15:1 & \begin{tabularx}{0.7\textwidth}{X} 有幾個人從猶太下來，教導弟兄們說：「你們若不按照摩西的規矩受割禮，不能得救。」 \end{tabularx} \\ \\ \relax
15:2 & \begin{tabularx}{0.7\textwidth}{X} 保羅和巴拿巴跟他們發生了激烈的爭執和辯論；大家就決定指派保羅、巴拿巴和本會的幾個人，為所辯論的事上耶路撒冷去見使徒和長老。 \end{tabularx} \\ \\ \relax
15:3 & \begin{tabularx}{0.7\textwidth}{X} 於是教會為他們送行。他們經過腓尼基、撒瑪利亞，沿途敘說外邦人歸主的事，使眾弟兄都非常歡喜。 \end{tabularx} \\ \\ \relax
15:4 & \begin{tabularx}{0.7\textwidth}{X} 他們到了耶路撒冷，教會、使徒和長老都接待他們，他們就述說神同他們所做的一切事。 \end{tabularx} \\ \\ \relax
15:5 & \begin{tabularx}{0.7\textwidth}{X} 惟有幾個法利賽派的信徒起來，說：「必須給外邦人行割禮，吩咐他們遵守摩西的律法。」 \end{tabularx} \\ \\ \relax
15:6 & \begin{tabularx}{0.7\textwidth}{X} 使徒和長老聚集商議這事。 \end{tabularx} \\ \\ \relax
15:7 & \begin{tabularx}{0.7\textwidth}{X} 辯論了許久後，彼得站起來，對他們說：「諸位弟兄，你們知道神早已在你們中間揀選了我，讓外邦人從我口中得聽福音之道，而且相信。 \end{tabularx} \\ \\ \relax
15:8 & \begin{tabularx}{0.7\textwidth}{X} 知道人心的神也為他們作了見證，賜聖靈給他們，正如給我們一樣； \end{tabularx} \\ \\ \relax
15:9 & \begin{tabularx}{0.7\textwidth}{X} 又藉著信潔淨了他們的心，他們和我們之間並沒有甚麼分別。 \end{tabularx} \\ \\ \relax
15:10 & \begin{tabularx}{0.7\textwidth}{X} 現在你們為甚麼試探神，要把我們祖宗和我們所不能負的軛放在門徒的頸項上呢？ \end{tabularx} \\ \\ \relax
15:11 & \begin{tabularx}{0.7\textwidth}{X} 相反地，我們相信，我們得救是因主耶穌的恩典，和他們一樣。」 \end{tabularx} \\ \\ \relax
15:12 & \begin{tabularx}{0.7\textwidth}{X} 眾人都默默無聲，聽巴拿巴和保羅述說神藉著他們在外邦人中所行的神蹟和奇事。 \end{tabularx} \\ \\ \relax
15:13 & \begin{tabularx}{0.7\textwidth}{X} 他們講完了，雅各回答說：「諸位弟兄，請聽我說。 \end{tabularx} \\ \\ \relax
15:14 & \begin{tabularx}{0.7\textwidth}{X} 剛才西門述說神當初怎樣眷顧外邦人，從他們中間選取人民歸於自己的名下； \end{tabularx} \\ \\ \relax
15:15 & \begin{tabularx}{0.7\textwidth}{X} 眾先知的話也與這意思相符合。 \end{tabularx} \\ \\ \relax
15:16 & \begin{tabularx}{0.7\textwidth}{X} 正如經上所寫的：『此後，我要回來，重新修造大衛倒塌了的帳幕，從廢墟中重新修造，把它建立起來， \end{tabularx} \\ \\ \relax
15:17 & \begin{tabularx}{0.7\textwidth}{X} 使剩餘的人，就是凡稱我名的外邦人，都尋求主。 \end{tabularx} \\ \\ \relax
15:18 & \begin{tabularx}{0.7\textwidth}{X} 這話是自古以來顯明這些事的主說的。』 \end{tabularx} \\ \\ \relax
15:19 & \begin{tabularx}{0.7\textwidth}{X} 所以，我的意見是不可難為那歸向神的外邦人； \end{tabularx} \\ \\ \relax
15:20 & \begin{tabularx}{0.7\textwidth}{X} 但是要寫信吩咐他們禁戒偶像所玷污的東西、血和勒死的牲畜，禁戒淫亂。 \end{tabularx} \\ \\ \relax
15:21 & \begin{tabularx}{0.7\textwidth}{X} 因為歷代以來，摩西的書在各城都有人宣講，每逢安息日，也在會堂裡誦讀。」 \end{tabularx} \\ \\
[1ex]
\hline
\hline
\end{longtable}
$^{1}$不知道大家是否喜歡聽演唱會?.
有沒有去搶過演唱會的票?.
我不知道你喜歡甚麼歌星.
張學友,劉德華還是梅艷芳.
我想問一下大家.
如果你有買過演唱會的票.
我相信你都同意的就是.
演唱會的票很貴.
不單止貴.
有時候在未有網購的年代.
我們還要排很長的隊伍.
有沒有這樣的印象?.
如果你不想排隊怎麼辦?.
不想排隊都可以.
你就買黃牛票.
說到黃牛票.
大家知不知道香港有哪個歌星.
他有一個紀錄是最高最高的黃牛票?.
你猜猜是誰?.
我上網查過.
是謝霆鋒.
那你知不知道謝霆鋒.
他最貴的那張黃牛票多少錢?.
原本應該是1200的門票.
在黃牛檔查到20萬.
20萬.
我看到都覺得好像你一樣.
有沒有搞錯?.
是不是有黃票賣給他做生日禮物?.
我想大家想一想.
如果某一些原因.
我不知道是什麼原因.
原來他認識你的.
他謝霆鋒簽了一個靚名.
給你免費的演唱會門票.
開心到不得了.
到了演唱會那天.
拿著張票很開心地走去聽演唱會.
到了門口入閘的時候.
那個護衛一手按住你.

$^{41}$先不要進去.
然後跟你說.
你先交2000元.
不是吧.
張票明明是免費的.
為什麼我要交2000元?.
沒錯.
你的票是免費聽演唱會.
不過你要交2000元的行政費.
那不就是不是免費.
那你給還是不給?.
你給的就是說不是免費.
不給的.
我舔都舔你.
記得要進去聽.
跟他辯論.
今天的經文.
告訴我們.
在第一代的基督徒.
他們所面對的問題.
遠遠嚴重過.
一張免費的演唱會門票.
為什麼這樣說?.
因為他們一直所持守的信仰就是.
信耶穌得永生.
是不是?.
我信耶穌這個伯伯神給我們的恩典.
我們就可以得救進入天國.
但現在不是.
現在有另外一班人來告訴他們.
不是的.
你還要守摩西的律法.
還要受國禮.
這樣才可以.
如果是你,你會怎麼處理?.
讓我們一起看看.
從經文裡面學習.
當日的教會領袖.
如何面對一個這麼嚴重的紛爭.
好了.

$^{81}$第一節.
講到有幾個人從猶太下來教訓.
如果你們不按摩西的規條受國禮的話.
不能得救.
很嚴重.
不過我想大家留意.
這幾個人.
你千萬不要當他們是壞人.
當他們是搞亂教會.
並不是這樣.
他們其實是很敬虔的信徒.
為什麼這樣說?.
因為我們要站在猶太人的歷史來明白.
這班猶太人.
為什麼這樣說?.
因為對他們來說.
幾千年來.
國禮或者是摩西的律法.
是神將他們從萬民中挑選出來的選民.
是從萬民中給他們一個約.
是一個很重要的一環.
這個國禮.
這個摩西律法.
是神賜給他們.
讓他們成為一個神所揀選的群體.
所以在那個角度來說.
大家要明白.
他們不是無中生有的在搞事.
他們其實是忠心於神的吩咐與命令.
他們嘗試保存的是.
他們數千年以來信仰的完整性.
對於他們來說.
外邦人突然間這班人走過來.
又不熟悉摩西律法.
又吃一些我們覺得不潔淨的食物.
他們大概是這樣想的.
這班外邦人.
不可以就這樣大搖大擺走進來.
然後就不遵守戒律.
就說自己得了應許.

$^{121}$這樣還得了嗎.
這班人其實不是否認耶穌基督.
我們先搞清楚這件事.
他們其實就是說.
信耶穌以上.
還要加上某一些東西.
要加上國禮.
加上他們的傳統.
他們的文化.
而這些加上的東西.
並不是原初福音的要求.
從剛才我說的演唱會的例子來說.
並不是這張免費門票無效.
只不過他另外要加行政費而已.
不過無論你叫這個東西是甚麼東西.
總之在原本的東西上再加些東西的時候.
這個就不再是恩典.
耶穌基督在十字架上的壽死就成為謊言.
當你加某一些東西在恩典上的時候.
其實就是耶穌十字架在恩典上拿去的東西.
恩典因此而不再是恩典.
救恩.
如果是耶穌基督再加一些東西的時候.
正如這裡所顯示.
不再是恩典而是報酬.
是我們通過自己的努力賺取回來的報酬.
我們所謂的福音.
就不再是在乎神為我們做了甚麼.
而是在乎我們為神做了甚麼.
這個福音完全搞亂了.
當保羅和巴拿伯聽到這個言論的時候.
聖經怎樣說.
聖經說他們兩個就與他們大大的紛爭辯論.
相當激烈.
為甚麼.
因為很明顯.
這個不單是對某一個事件的不同看法.
而是一個言論直接動搖了福音的根基.
所以他們馬上就與他們辯論.
大家要留意.

$^{161}$辯論的雙方都是基督徒.
他們兩個其實都是接納外邦人進入救恩.
唯一的分別就是.
一群人覺得耶穌基督的信心以上.
還要加上一些條件.
即是說這個紛爭是在教會裡面發生.
當教會出現紛爭的時候.
我們應該怎樣處理呢.
中國人就說.
家家有本難唸的經.
我相信沒有人會否認.
家家有本難唸的經.
不過當我們認定教會是我們屬靈的家的時候.
我們都不會詫異教會一樣有本難唸的經.
不過這本難唸的經並不是指聖經.
這本難唸的經就是說.
每一個地方的教會.
都會有一些不為人知的事情.
我記得我曾經有一次去一間教會講道.
那間教會有一段頗長的時間沒有牧師.
他們自己主動跟我說.
我們因為早前發生了一些事情.
所以牧師就離職了.
我沒有再問.
因為第一我不需要知道.
第二就算他不講.
我大概都知道發生什麼事情.
對我來說教會裡面有不同的問題出現.
其實很正常.
大家想一想.
坐在這裡的全部都是不同個別的人.
我們有不同的成長背景.
不同的習性.
不同的做事方式.
不同的喜好.
不同的嗜好.
所以當我們一齊相處.
一齊來做教會的事工的時候.
當然是會有問題出現.
而且之前那兩次講道.

$^{201}$如果大家還記得的話.
每一次當教會做對的事情的時候.
就會有什麼事情發生.
撒旦就會在背後來做搞擾的工作.
所以從另一個角度去看.
如果教會今天完全沒有問題的時候.
也許教會沒有在做對的事情.
想像一下.
如果我們教會所有的人回來.
我沒有想著什麼.
我只是想著回來認識朋友.
可以陪我去風花雪月.
吃喝玩樂多好.
如果是這樣的關係.
這樣的教會的時候.
我想大概問題不會很多.
這個現象在教會成立的第一天.
開始就已經發生了.
在今天的經文裡面.
當神將人數天天加給教會.
當撒旦感到受威脅的時候.
撒旦一個遠程導彈.
炸下來教會那裡.
為什麼我說遠程導彈.
因為這個言論很有破壞性.
它直接的攻擊整個福音的核心信仰思想.
究竟救恩是因神的恩典.
還是因人的行為呢.
如果你是當日要處理這個紛爭的.
其中一個領袖.
你會怎麼做.
我們再看下去.
當時的眾門徒定規.
叫保羅和巴拿巴和幾個人去哪裡.
去耶路撒冷見使徒和長老.
我想問為什麼要見使徒和長老.
用剛才演唱會的例子.
如果我們不想給他二千元行政費.
又想去看演唱會.
那怎樣.

$^{241}$就跟他理論.
那怎樣理論.
張飛是謝霆鋒給的.
就打電話給謝霆鋒.
我猜你都很難找到他.
除非你是謝賢或者王飛.
不然我猜都很難找到他.
那怎麼辦.
那就找一個經理人.
聖經告訴我們.
使徒與長老是神所授權的一個地上.
他的代表來管理教會.
正如保羅在哥林多前書這樣說.
人應當以我們為基督的執事.
為神奧秘士的管家.
正因如此.
教會是建立在使徒和先知的根基上.
而耶穌基督就是整個教會的防國石.
什麼意思.
意思就是說.
我們所有的真理.
必須出自耶穌基督.
因為耶穌基督是唯一的道路.
真理和生命.
我們看看當日的信徒領袖.
怎樣按著這個真理來處理這個糾紛.
經過一輪的各自雙方的論述之後.
彼得是第一個站出來說的.
彼得不是高言闊論.
他只不過是告訴他們.
他自己親身的經歷.
他的什麼經歷.
就是聖靈感動他來到外邦人哥利樓那裡.
來將福音給他們.
他接著說了一個很重要的經驗.
就是說原來他在沒有做任何儀式.
沒有行國禮.
也沒有在他們遵守了摩西律法之後.
竟然聖靈降臨在他們身上.
就是說他們好像猶太人一樣.

$^{281}$得著救恩.
沒錯.
大部分的猶太人都是喜愛律法.
他們視律法是神給他們一個特別的恩賜.
讓他們從地上萬俗之中來分別出來.
但是當彼得通過救恩的淚鏡.
再重新看這件事情的時候.
他發覺原來你是單靠人的努力.
或者是行律法來得著救恩的時候.
這條路極度痛苦.
而且是近乎不可能.
簡直是不可能.
為什麼?.
因為沒有一個人在神的面前可以做到一百分.
從救恩的角度來看.
不論你是猶太人還是不是猶太人.
都是需要依靠神的恩典.
耶穌基督裡面的信心來得著救恩.
所以彼得就說.
既然是這樣.
為什麼我們要試探神呢?.
好了.
接著就輪到保羅和巴拿巴.
他們分享了神怎樣通過他們兩個在外邦人的裡面.
來行神跡騎士.
他們能夠行神跡騎士.
並不是因為這班人是持守摩西律法.
而純粹因為他們相信福音.
所以他們分享的不是什麼神學的理論.
是一個很簡單的見證.
一個第一手,第一身的經驗.
他看到這班人外邦人.
他們有病的得醫治.
低的可以再行.
被鬼附的可以得到釋放.
通過神的能力來印證.
這兩位使徒他們在外邦人裡面所做的工作.
有時候你想看是否神的心意其實不困難.
你看一看聖靈在做什麼.
你就知道是否合乎神的心意.

$^{321}$當這班人在思考的時候.
雅各就站出來說.
他說剛才這些弟兄所分享的.
其實和我們舊約裡面所講的意思是相合.
以之前的聖經怎樣說的.
之前的聖經說.
將來神要重新修造大衛倒塌了的帳幕.
把破壞的重新修造,建立起來.
叫凡稱為我名下的外邦人都尋求主.
這段經文是來自阿摩斯書的九章十一到十二節.
是講到一個預言將來神會重新修復倒塌了的帳篷.
接著外邦人就會尋求神的命.
這是一個很重要的提點.
就是說外邦人被接納在這個救恩裡面.
不是一個需要討論的題目.
因為這是在應驗的先知的預言.
外邦人被包括在神的救恩裡面.
一直都是神的計劃.
先知的預言.
彼得的經驗.
保羅所經驗到的神蹟奇事.
所有東西都指向同一個結論.
什麼結論?.
在耶穌基督裡面.
如果我們有信心的話.
神的國是屬於所有的人.
屬於萬國萬民.
這就是神的恩典.
不是人的行為.
正如保羅所說.
你們得救是本乎恩.
也是因著信.
不是出於自己免得有人自誇.
在雅各做了一個總結.
簡而言之.
是神的心意.
無論是外邦人還是猶太人.
都是通過耶穌基督的信心.
而得著救恩.
做完這個總結之後.

$^{361}$雅各就停在那裡.
雅各還說了一些話.
他說.
按我的意見.
即是他自己的看法.
提點了幾樣東西.
來作為外邦人一些行為的守則.
讓我們一起看看.
這幾個行為的守則.
第一.
禁戒偶像的污穢.
這裡應該是指到.
當時拜偶像的食物.
雅各叫他們.
拜偶像的食物不要吃.
為甚麼.
因為你給人的感覺就是.
你參與在偶像的敬拜裡面.
所以不要這樣做.
第二.
姦淫.
大家要明白.
當時在外邦人的文化裡面.
去找妓女.
或者去外邦人的廟宇裡面.
亦都有一些廟祭的提供.
不要參加那些.
而外邦人的文化裡面.
他們覺得.
有時一些性經驗.
如果是社交式的時候.
可接受的.
雅各說.
千萬不要這樣.
我們跟隨的不是外邦的文化.
或者不是跟隨世上的文化.
而是跟隨耶穌基督在聖經裡面.
給過我們的教導.
第三個是勒死的生畜.
大家知道在猶太人的律法裡面.

$^{401}$動物要經過很小心的處理.
讓血流乾.
如果是勒死的動物.
體內仍然有血.
所以要禁戒.
第四是血.
大家都知道猶太人避免不吃血.
為甚麼.
因為聖經的教義裡面說到.
血是代表著生命.
所以不要吃.
聽上去大家覺得合理嗎.
我覺得很合理.
不過等一下.
剛才不是說過在耶穌基督的信心上.
不要再加任何東西嗎.
當你加了的時候.
豈不是恩信稱義的恩典.
成為律法主義.
究竟得救是本乎因.
還是因行為呢.
大家在這裡要相當的清晰留意.
這四個建議並不是律法而是愛.
雅各並不是說在救恩上加多這四個條件.
你做不到的時候就不得救.
雅各不是這樣說.
他們這四個建議.
我們要明白的就是.
他們是要維持一個親密團契裡面.
一定要有的互相尊重.
因為對猶太人來說.
這四個東西在他們日常生活裡是完全不可接受的.
如果外邦人要和猶太人相處相交的時候.
猶太人見這些外邦人做這四件事的時候.
搞不定搞不定.
沒錯.
外邦人是不再需要遵守律法.
只需要憑著耶穌基督的信心得著救恩.
但是他們卻不可以隨意的.
想做什麼就做什麼的情況下.

$^{441}$來冒犯他們的猶太兄弟.
所以這是一個很重要的提醒.
這就是在哥林多前書的一些原則.
保羅在哥林多前書這樣說.
凡事都可能但不都有益處.
所以你們要謹慎.
你們一定要謹慎.
不要將你們的自由.
成為其他軟弱人的絆腳石.
救恩完全是神的恩典.
一個教會.
一個健康的教會.
必須要建立在救恩的純正.
和肢體的愛上面.
福音的真理絕對不能夠妥協.
但是當我們信徒要相交的時候.
我們需要互相體諒.
互相尊重的愛在當中.
我想大家去想像一下.
去聽演唱會.
這個交響樂的音樂會.
大家有沒有去過一個交響樂的音樂會.
全部樂器是一樣的.
大家大致上應該沒有.
你去到一個交響樂屋的時候.
你當然期望很多種不同的樂器.
你只有十部和二十部.
如果都是同一個樂器的時候.
你會有什麼感覺.
很單調.
和好像沒有了一種完整的感覺.
交響樂曲應該是由不同的樂器所組合而成.
特別當指揮家拿起指揮棒.
來帶動音樂演奏的時候.
一個美妙的交響樂曲.
就可以充滿了當時的會場.
所以來自不同樂器的音樂.
並不是混亂.
也沒有互相的競爭.
而是只會帶出和諧.

$^{481}$而一個交響樂曲美妙的地方.
就是正正在於有不同的樂器.
正正由於它們的不同.
就能夠演奏出一個美妙的音樂.
其實教會也是一樣.
大家想想看.
我們周圍的人和你都是不同的.
我們都是來自不同的成長背景.
有不同的文化.
不同的習性.
不同的工作方式.
所以大家可以想像一下.
帶著種種的不同.
如果好像一個音樂會裡面.
沒有了指揮的時候.
你有你彈.
你有你拉.
你有你按.
這個演唱會就會搞得亂七八糟.
但是感恩.
我們第一代的教會領袖.
幫教會找到一個和諧.
這個和諧的方式.
並不是要求每一個的會眾.
一模一樣.
倒一個模出來.
每個人一模一樣.
那就不會有紛爭了.
他們不是要求當時的會眾.
他們所做的就是教導信徒.
當我們要和我們自己不同的弟兄相交的時候.
有的時候我們需要限制我們自己的自由.
換言之.
當我們限制自己的自由.
我們不單止可以享受肢體的生活.
也可以幫助教會演奏出一個和諧的樂曲出來.
不過也要同一時間留意.
一個不同樂器的交響樂團.
他要演奏出一個美妙的音樂.
一個一定要有的因素.

$^{521}$是什麼?.
每一個樂器都要聽著指揮做事.
指揮叫你某一個時間彈你就彈.
如果你不聽這個指揮做事的時候.
那首音樂就不會能夠成功地演奏出來.
所以我們要明白.
教會的指揮家是誰.
不需要我講.
我們都明白知道教會的指揮家就是耶穌基督.
祂帶領著信徒通過聖靈.
雖然我們每一個人都不同.
但我們仍然可以成為一個和諧,合一的群體.
在一個健康的教會裡面.
不同的文化.
我們仍然可以一同敬拜.
不同的性情.
我們可以一同事奉.
不同的音賜.
我們仍然可以合一.
唯有這樣教會才能夠為神.
將這個美妙的福音演奏成一個美妙的音樂.
讓我們周圍的地區都聽得到神的音樂.
所以弟兄姊妹.
容讓我們彼此鼓勵.
信服聖靈的帶領.
進行基督的教導.
讓天上的父神得到祂當得的榮耀.
請我們一同低頭祈禱.
我們的教會就是我們屬靈的家.
但我們同一時間都知道.
很多時候因著我們的軟弱.
因著撒旦的攻擊.
教會總是都會出現不同的問題.
但這些問題我們不需要懼怕.
也不需要因此對教會有不好的品評.
我們需要的就是.
在這些面對這些紛爭的時候.
我們需要以基督為我們的首要考慮.
我們需要仍然持守.
福音帶給我們信仰群體的合一.

$^{561}$以致當外人來到進入教會的時候.
看到我們的不同.
卻是合一的時候.
主啊 我們就能夠見證你的榮美.
你的尊貴.
求你恩代我們在這個教會裡面.
持守你給我們的一切.
讓你的榮耀得著高舉.
祈禱奉主明求.
阿們.
\newpage



\section{路加福音 10:38-42}
\label{sec:031803}
\textbf{為馬大殷勤服侍平反?}~~~~鄺偉志牧師
\newline
\newline
連結: \href{http://www.church.com.hk/acms/content.asp?site=cdc\&op=show\&type=product\&code=031803\&layout=sermon}{\texttt{031803}} ~~~~ 語音日期: 2026-03-15
\newline
\newline
\hyperref[sec:031802]{\small{< < < PREV SERMON < < <}}
~
\hyperref[sec:index_chronic]{\small{[返順時目]}}
~
\hyperref[sec:031804]{\small{> > > NEXT SERMON > > >}}
\newline
\newline
路加福音 10:38-42
\newline
\begin{longtable}{cl}
\hline
\hline
章節 & 經文 (和合本修訂版)\\
\hline
10:38 & \begin{tabularx}{0.7\textwidth}{X} 他們繼續前行，耶穌進了一個村莊。有一個女人，名叫馬大，接他到自己家裡。 \end{tabularx} \\ \\ \relax
10:39 & \begin{tabularx}{0.7\textwidth}{X} 她有一個妹妹，名叫馬利亞，在主的腳前坐著聽他的道。 \end{tabularx} \\ \\ \relax
10:40 & \begin{tabularx}{0.7\textwidth}{X} 馬大伺候的事多，心裡忙亂，進前來，說：「主啊，我的妹妹留下我一個人伺候，你不在意嗎？請吩咐她來幫助我。」 \end{tabularx} \\ \\ \relax
10:41 & \begin{tabularx}{0.7\textwidth}{X} 主回答說：「馬大，馬大，你為許多的事操心煩惱， \end{tabularx} \\ \\ \relax
10:42 & \begin{tabularx}{0.7\textwidth}{X} 但是不可少的只有一件。馬利亞已經選擇了那上好的福分，是沒有人能從她奪去的。」 \end{tabularx} \\ \\
[1ex]
\hline
\hline
\end{longtable}
$^{1}$各位弟兄姊妹,早安!.
好,先準備一下,我們一起祈禱吧!.
在崇拜裡面,講道的,是一個你的器皿.
我們講的和聽的,都要聆聽你的話語.
以致我們明白了主你的話語.
更加明白了神,更加反思我們的光景.
以致我們呢,能夠可以更加去敬拜你.
謙卑下來,再一次跟隨你的腳蹤.
我們將以下時間交上,求你幫助孩子.
能夠將你的話語,講得明白.
弟兄姊妹亦都聆聽的時候,不是單是聽到人的聲音.
但是主你透過你的話語的講解.
將你的話語,擺在人的心裡面.
好像兩人離間,幫助我們去看你的話語的時候.
去謹守真行你的道.
多謝主聽我們禱告,奉耶穌基督的名字.
阿們.
有些話還取寵,但其實不是我想出來的題目.
一向都有人來引述這段經文的故事.
就說我們應該像瑪利亞一樣.
好像這個故事一樣.
安靜地在主的腳前聽到.
就要跟神建立關係.
不要像瑪利亞去馬大那樣.
只顧著做,忙到不停.
你沒有聽過這個說法.
前一陣子,準備講道之前一段時間.
就看到一篇文章.
其實不單止是一篇文章.
有時候弟兄姊妹都會這樣說.
其實文章或者那些人的分享大概是這樣.
其實心裡面替馬大不值被耶穌罵.
為什麼呢?你看看馬大那麼勤力服侍.
看到妹妹坐著不做事.
當然是看不過眼.
當然叫她做事.
我覺得馬大就對了.
幾乎是耶穌罵錯了她.
所以我聽別人分享.
有時候說著說著好像說到自己家裡.

$^{41}$他家裡有兩個女兒.
大女兒這樣做,小女兒不這樣做.
全部代入下去.
所以其實馬大要平反.
馬大有它的好處.
這樣說法有沒有道理呢?.
有道理的.
因為有時候在教會界.
我想不到第二個說法.
或者我們在信徒之間.
有時候將這個故事的兩個角色.
講得很盡.
馬大好像很對立.
馬大的形象就是這樣做.
事奉很賣力.
不過太過份了.
但是另外一種形象就是瑪利亞.
應該聽道,追求神.
很安靜的.
甚至有些解經去到一個極端.
其實想來想去都好像有點不對.
你看到一些釋經書.
有些覺得馬大代表了靠行為得救.
聽過沒有?.
瑪利亞是靠信心得救.
為什麼會推到這個推論呢?.
我也不明白.
為什麼會這樣.
可能在解釋裡面.
混合很多東西去解釋.
就變成這樣的情況.
甚至我們很多時候說.
我們最重要注重我們的being.
不要只是doing.
你也聽過了.
這個對不對呢?.
當然是對的.
我們不只是做做做.
我們的內裡.
我們的人.

$^{81}$是怎樣是很重要的.
其實問題在哪裡呢?.
是不是這樣去解釋呢?.
所以下一張.
兩張slide.
第二張就是整段經文.
很長.
我們細心看一下.
究竟問題在哪裡.
是不是這樣對立.
首先我們看馬大的問題.
經文其實也可以再看一次.
如果你看38節.
你看到有個女人叫馬大.
接耶穌到她家裡.
其實她是接待他.
服侍他.
對耶穌是一個好的做法.
為了服侍他才接待他.
接到他家裡就服侍他.
39節.
她有個妹妹叫瑪麗亞.
她在主的腳前坐.
聽到待會會解釋這個.
大家看一下馬大的問題在哪裡.
馬大事後的事多.
心裡忙亂.
進前來.
就說.
主啊.
我的妹妹.
留下我一個人事後.
你不在意嗎.
請吩咐她來幫助我.
其實核心就在.
她說的話.
和她的狀態.
這個狀態.
41節.
主回答說.

$^{121}$馬大馬大.
你為許多事.
操心煩擾.
所以這個情況.
就是馬大的情況.
我們先講馬大的情況.
你看看馬大的問題.
我問一些問題.
去澄清.
她的問題.
是不是工作.
是不是服侍.
她很多服侍.
所以不對.
不是啊.
她的問題在哪裡.
她的問題就是.
她很多服侍的時候.
有兩個形容.
這裡有中文的.
和本修訂版翻譯.
第40節.
這個馬大.
她事後的事情很多.
以致她.
心裡忙亂.
OK.
接著耶穌後面.
這是盧家作者形容她.
接著耶穌正在證實這件事.
再很清楚.
41節.
馬大馬大.
講兩次.
提她.
喂喂喂.
你要留心一點.
提她兩次.
耶穌不是經常這樣叫人兩次.
如果我們去思考這個問題.

$^{161}$我們都可以想想.
有沒有耶穌說.
抗委之抗委之.
有人提你.
誰誰.
提你兩次.
因為她已經很迷失.
你為許多事操心煩擾.
兩個不同的說法.
但是.
問題就是.
她一直做.
她很多事這樣做.
但是.
她的心裡.
忙亂.
有幾個層次.
第一個層次.
就是心的問題.
所以不是.
我自己去看.
一會兒再解釋多一點.
但是我自己去看.
就不是.
馬大有很多事奉.
所以.
就被耶穌說.
不是這個問題.
因為如果推到這樣.
那不用事奉了.
事奉就是不對的.
你只是讀經聽道就算了.
不是這樣.
絕對不是這樣.
但是.
她這個意思就是說.
如果你看她原本那個字的意思.
就是心裡忙亂.
兩個不同的字.
但是她也是在說一樣東西.

$^{201}$分心.
就是.
distract了.
或者是.
有另外一個說法就是.
pull away.
就是被一些東西拉走了.
她不能夠專注.
所以.
問題就是.
她做很多事.
但是她不能夠專注.
專注什麼呢.
其實就是相對馬尼亞.
她在耶穌的腳前.
去聽他的道.
那你說.
馬大在做事.
聽不到道的.
但是其實.
就不是說.
她只是聽不到道.
而是說她的心.
不專一.
她的心.
沒有放.
主.
在首位.
為什麼呢.
你看到.
在.
馬大說的話.
其實這段經文.
有一個重複字眼.
經常出現.
是一個重點.
不知道你有沒有留意到.
就是主字.
馬尼亞就在主的腳前.
一會兒會解釋.

$^{241}$但是你看到馬大.
她也說主.
她說就這樣說.
但是她是不是當耶穌是主.
是不是當耶穌是主.
她的說法.
她吩咐耶穌做事.
你吩咐我妹妹.
拜託你.
現在我就很忙.
現在我就很煩.
我見到我妹妹不做事.
但是請耶穌你.
吩咐耶穌.
你去叫馬尼亞做事.
其實問題就是.
她說的口就說主.
但是她做出來說的話.
就是想自己做主.
甚至想操控耶穌.
你.
拜託你聽我說.
那誰是主呢.
所以馬大的問題.
她就是被分心.
不專注.
以及.
不認識主的主權.
所以.
今天的message不複雜.
不是說我們事奉有問題.
不是說我們很勤力事奉有問題.
如果不是這樣.
那就很大件事.
有時候聽到.
有人說.
很有趣的.
她說見到一個人事奉.
你不要像馬大那樣.
你去馬尼亞.

$^{281}$有時候聽到火都來.
不是馬大.
她忙碌做事.
不是她忙碌做事的問題.
是她忙碌做事.
被分心.
被pull away.
是分心.
以及她沒有專注為大.
她沒有以主的看法.
想法為先.
尤其是.
沒有對神的話語專注.
是那個問題.
不是她做事的問題.
OK.
不要拿那件事繼續說.
在我面前不要說.
起碼.
接著看.
我講講一個例子.
我在很多年前事奉.
有一個機構.
是音樂事奉.
可能大家會聽過.
但如果你在樹梨內.
應該沒有聽過.
大概二十多年前.
有一個機構舞起得很快.
我認識那個負責人.
他很有心.
說起名字.
可能你會認識.
也有一點點名氣.
在基督教音樂界.
他搞的機構.
在幾年裡.
很感恩.
帶動了很多人.
去敬拜.

$^{321}$我記得.
有一次.
他搞了第五年.
他搞了一個機構.
搞了一個很大型的敬拜聚會.
在伊爾沙伯體育館.
灣仔.
呼籲弟兄姊妹一齊去參與師班.
很有趣.
呼籲帶小朋友一齊去敬拜.
所以很虛憾.
我也帶了個兒子.
那時候很小.
穿著T恤.
全個會場穿著T恤練歌.
但很深刻印象.
那一次練歌.
很多人.
他做Worship Leader.
去帶.
但很複雜.
總之長話短說.
他去到一個位置.
一直在練的時候.
時間很緊迫.
我不說那個性別.
那個位置是姊妹弟兄.
他突然之間.
語氣一直很憤怒.
說大家不要吵.
很大聲.
你知道那些音樂人.
運氣起來.
整個會場都靜了.
很大聲.
大家嚇一嚇.
發生什麼事.
後來靜下來.
他再回氣.
再說怎樣去準備.

$^{361}$事實上.
你想想千人大師班.
又有小朋友.
很複雜.
但那一次.
很深刻印象.
因為我認識他.
後來這件事之後.
這個活動過了.
都很成功.
很順利.
但這位姊弟.
他選擇將這個機構.
慢慢地放下.
後來完全消失了.
所以你現在聽不到.
這個叫教音樂的機構.
其中一件事.
我相信不只是那一刻.
但那一刻.
讓我們看到.
他自己後來都會說.
其實他的心靈裡.
是有了問題.
所以問題是.
我們去反思.
馬大的情況.
當然我再說.
不是說不做事.
不服事.
這是另一個要反省的問題.
但就算是服事.
的確是要去反思.
究竟我們的心是怎樣.
是否忙亂.
是否一直事奉.
有沒有去專注神那裡.
有時不只是神的話語.
所以這是馬大問題.
第二個問題.

$^{401}$不是問題.
第二個我們要看.
就是馬利亞.
究竟我們學了甚麼.
是否學她那麼安靜.
即是你個個在聽道.
這就叫做上好的福分.
所以要解釋一下.
首先就是第38節.
第39節很多人都指出.
馬利亞.
在主的腳前.
坐在那裡聽她的道.
即是主的道.
其實是當時的一個說法.
就是你知道.
保羅有個師父.
就是加瑪列.
總之那些師父去教徒弟.
其實就是門徒.
他就用了一個說法.
就是那些門徒.
在那個教師的腳前.
來聽他的教導.
直接說這件事.
所以如果這樣說.
一開始聖經形容馬利亞.
她在甚麼位置.
她在一個門徒的位置.
即是她是聽耶穌的道.
其實後面也很清楚.
尤其是在42節.
耶穌稱讚馬利亞.
祂說不可少的只有一件.
馬利亞已經選擇了上好的福分.
是沒有人能夠從祂那裡奪去的.
那個信息量很大.
雖然是一句說話.
逐句都是.
你要想清楚究竟祂說了甚麼.

$^{441}$以致你明白馬利亞.
被祂讚的是甚麼.
42節那個說.
但是不可少的只有一件.
其實耶穌說話很多時候.
是有很多對比的.
有時說一個道理.
尤其是我們查新生命新生活.
最近正在教.
有時我們會很容易誤會.
在聖經的話語.
譬如說你愛世界.
你就不愛神.
但耶穌有時是用這些對比.
將那件事很極端地說了一個信息.
不是否定了.
因為有個組員就說.
其實我信了主一段時間.
我覺得好像我為家人.
為我下一代預備一些產業.
金錢好像不太對.
很洗俗的思想.
但我說其實在聖經裡面.
也是有愛人如己的.
我們也要來好好地管理我們的產業.
就是我們的兒女.
為我們的兒女去準備一些東西.
不是罪來的.
因為你有愛人.
我們的家人就更加要愛.
你為他準備.
有甚麼好大問題呢.
這樣是否就是愛世界.
你愛金錢呢.
其實兩個真理無敵對.
如果是這樣.
聖經裡面的有錢人就全部不行.
變成仇婦.
我們不是這樣看的.
大家都有錢.

$^{481}$因為個個都不行.
有些錢就不是這個問題.
所以我要說這件事.
所以當耶穌說.
他要去強調一些東西的時候.
馬尼亞只有一件是對比於馬大的甚麼.
你看到耶穌說甚麼嗎.
就是這一節.
為許多的事.
其實是故意說這個對比.
你馬尼亞很多事.
你為那麼多的事忙亂.
對比馬大衛很多事忙亂.
但對比馬尼亞.
只有一件事.
其實是一種強調的手法.
如果我們將這件事無限放大.
馬尼亞坐在耶穌的腳前讀經就行了.
聽到就做了基督徒.
非常好.
他選擇了上海的福分.
我不相信是這樣解釋的.
因為馬尼亞還有.
如果這樣計算.
馬尼亞還有甚麼事跡.
我們可以效法.
又不見我們效法.
甚麼.
馬尼亞獻珍拿達香膏.
一整年的薪水.
其實香水是香膏的意思.
將那麼貴重的香膏.
在耶穌要離開的時候.
來告耶穌.
我經常說的阿膏就是名詞.
告就是動詞.
他將香水去告耶穌.
我們又不會學他.
你要學馬尼亞.
你還要學那件事.

$^{521}$你坐在這裡聽到.
然後你要將你的身家奉獻.
是不是這樣.
我們不是看聖經那麼極端.
拿著.
只有一件.
所以我說.
其實耶穌是在說專注的問題.
在那一刻.
馬尼亞是專注.
去做為耶穌的門徒聽道.
其實也有識經者指出.
這是很不尋常的.
因為如果你說馬大.
作為一個女性在那時候去服事耶穌.
她是選擇在那時候去聽道.
在主耶穌的腳前.
所以這是一個.
耶穌強調.
他的專一.
以及第二個重要的字眼.
就是第二句.
馬尼亞已經選擇了上海的豐分.
是一個選擇.
其實你說馬尼亞在那時候.
可不可以做其他事.
是可以的.
但是她在那時候覺得.
耶穌在那時候的教導.
她是選擇做一個門徒.
好好去聽耶穌的教導.
這是重要的.
而她是這樣去做的時候.
耶穌就讚賞她的選擇.
還有這句裡面.
有一個叫做選擇.
上海的福分.
原本沒有福字.
這是什麼呢.
其實上海.

$^{561}$也是沒有上字的.
不過她翻譯的時候.
說這是上海的福分.
其實你看原本.
或者我看NASB.
這個英文翻譯.
都是說得很清楚.
那個福分其實沒有福字.
就是一個份.
一部分.
portion.
就是一個pizza.
其中一份.
你說這個份是什麼呢.
上面上好就是good.
就是一個好的份.
原來瑪利亞就選擇了一個.
好的portion.
就是pizza選擇大件一點.
或者選擇好一點的.
當然不是說pizza.
是說原來這個份.
有一些參考.
你看回舊約都是用類似的字眼.
雖然原文是不同.
但是譬如說詩篇.
119篇.
那篇詩篇是讚美神的話語.
有176節.
每一節都是.
讚美神的話語.
其中57節說.
耶和華是我的.
福分.
福分其實是沒有福字的.
耶和華是我那份.
OK.
我那份.
其實是什麼意思.
拿你的portion.

$^{601}$你去拿這個地.
而這個地.
是一個應許地.
你去拿你的.
神應許給你的那份.
耶和華是我的份.
就是.
耶和華是屬於我的.
我是很想拿它.
耶和華是屬於我的一份.
所以已經去到一個.
可以說是一個隱喻.
一個比喻.
一個寓意.
你那份我那份.
實質有一份.
而是說你那份.
我那份.
瑪利亞選擇了.
好的那份.
如果這樣看.
瑪利亞.
在當時.
她選擇了在耶穌腳前聽.
但是.
瑪利亞.
在不同的情況下.
她的選擇.
可能有不同.
但是耶穌說.
她選擇了好那份.
我們就問了一個問題.
所以這個就有些處境.
在這個情況.
瑪利亞選擇了這件事.
但整體來說.
耶穌稱讚她.
當瑪利亞表達.
她這樣選擇的時候.
這就正確了.

$^{641}$我會問.
其實在我們的生活.
生命裡面.
祈奉裡面.
我們看不看到神.
給我們一些東西.
或者你祈禱尋求神.
再落地一點.
我現在究竟怎麼辦.
沒有工作.
我祈禱神.
你給我一份工作.
或者我現在有不同的困難.
神你幫助我.
我們相不相信.
神為我們預備了一份東西.
有一個portion.
神是應許的.
給我們得著.
不過.
是有選擇的.
你可以不選擇.
其實基督徒經常都是這樣.
有很多事情.
你可以尋求神.
神告訴你.
你可以做這件事.
你聽我的說話.
你有信心.
愛心.
你這樣得著.
你預備那份.
這是最好的.
但有時候我們真的.
不是.
如果基督徒做了很久就知道.
未必真的選擇最好的.
我們始終沒有信心.
或者沒有愛心.
所以如果這樣看.

$^{681}$我們會問從瑪利亞那邊.
他選擇了聽道.
當然這是好的.
在一個情況.
每個人都不聽神的話.
但你真的很想明白神的話語.
你去培靈會.
你去聽道.
但問題是.
那些份還有很多份.
神為我們預備了很多東西.
我們是否看到.
並且願意去選擇.
再說瑪利亞.
他如果那時候選擇.
和姐姐一起去.
做服事.
是很自然的.
也可能得到別人的讚賞.
但瑪利亞竟然.
迪力派中醫.
做一個女性竟然坐在耶穌面前.
做門徒去聽道.
很奇怪.
這句說話.
要看上文下理.
記不記得上個月.
我說到唱了一首歌.
那首歌.
是什麼故事.
沒有反應不要緊.
大家可能比較累.
或者真的忘記了.
好撒瑪利亞人.
看回經文.
對上一段.
就是那個故事.
其實你看任何一本色經書.
看一些色經書.
比較好的.

$^{721}$他一定是將兩件事拉在一起說.
因為路加福音的編排.
或者所有福音書.
都是這樣.
他不會將故事亂砌.
為什麼.
這個故事會在好撒瑪利亞人的故事.
緊隨著.
這個之後.
其實就是要補充.
整件事發生什麼事.
為什麼耶穌說好撒瑪利亞人的.
比喻.
因為有個律法師走來.
挑戰耶穌.
他說.
我怎樣才可以.
承受永生.
耶穌問他.
聖經怎麼教.
這樣回答.
你要盡心盡意.
愛主你的神.
並且要愛倫社.
然後.
耶穌說.
你這樣就會得到永生.
我再說.
他愛耶穌.
盡心盡意.
愛主你的神.
愛人愛神.
你會信神才愛.
所以耶穌說.
並且這樣活出信仰.
這樣愛人.
你就可以永生.
然後那個人就不服氣.
誰是我的倫社.
然後耶穌就跟他說了一個故事.

$^{761}$說什麼.
就是撒瑪利亞人的故事.
我上次特別說.
那個撒瑪利亞人.
他根本不是猶太人.
他是敵人.
但是他就從良心去回應.
誰是他的倫社.
有需要的人就是他的倫社.
他這樣身體力行.
去愛那個人.
這個就是撒瑪利亞人.
你說這句有什麼關係.
是沒有人奪去的.
這句其實就扯到承受永生的問題.
其實兩個故事.
兩個事蹟.
一個是很撒瑪利亞的.
一個故事.
這個就是一個真事.
馬大和馬利亞.
這個事蹟.
最後其實這是結論.
馬利亞是選擇了.
上好的福分.
就是說原來是優次的問題.
愛人.
你見到撒瑪利亞的故事.
表面上你看到他們是愛人.
但其實.
是愛神愛人.
等於很多人問.
基督徒和非基督徒有什麼分別.
那個都很好的人.
我們的分別.
其實我們不是很好的人.
第二就是.
我們是明白了.
我們是要愛神.
或者因為聖經的真理改變了我們.

$^{801}$以致我們去愛人.
當然我們的本性.
我們的良心.
都有愛人這部份.
但我們真正能夠愛人.
其實是因為我們信了主.
而愛神.
永遠是首先的.
優先就是.
你做多少好事都好.
所以才有哥林多前書一句.
你若捨己身焚燒.
卻沒有愛.
即是你的表那一層.
可以做到好像很愛.
但其實你的內涵是沒有的.
一樣是沒有.
所以就說了.
問題就是我們事奉.
我們聽到.
我們是否真的出於愛神.
愛人呢.
如果我們.
好像瑪利亞那樣.
將神放在主權.
在第一位.
選擇神給我們的墳.
而我們又.
這樣去事奉.
幫助我們去.
反思我們自己.
究竟我們在事奉.
我們做一個人.
或者我們愛人.
我們的動機.
我們是否一個愛神的人.
我們是否很專注.
或者我們在做的時候.
反而會不會像馬大那樣.
被很多事煩擾.

$^{841}$其實就沒有將主放在首位.
整件事是不同的.
主耶穌是要求我們.
愛主愛人.
以致我們這樣才是一個.
真正實踐著神的心意的人.
求你更新我們的內心.
給我們.
要聖靈幫助.
給我們檢視我們生活的每一步.
有時候我們沒有方向.
有時候我們想得到一些東西.
好像得不到.
大主.
你是已經為我們預備了最好的那一份.
求你給我們.
去選擇.
可能我們是知道的.
但我們沒有勇氣.
求你幫助我們.
真正的愛你.
我們會祈禱奉耶穌基督的名字.
阿們.
\newpage



\section{創世記 12:1-9}
\label{sec:031804}
\textbf{召命人生}~~~~任卓賢教授
\newline
\newline
連結: \href{http://www.church.com.hk/acms/content.asp?site=cdc\&op=show\&type=product\&code=031804\&layout=sermon}{\texttt{031804}} ~~~~ 語音日期: 2026-03-15
\newline
\newline
\hyperref[sec:031803]{\small{< < < PREV SERMON < < <}}
~
\hyperref[sec:index_chronic]{\small{[返順時目]}}
~
\hyperref[sec:031805]{\small{> > > NEXT SERMON > > >}}
\newline
\newline
創世記 12:1-9
\newline
\begin{longtable}{cl}
\hline
\hline
章節 & 經文 (和合本修訂版)\\
\hline
12:1 & \begin{tabularx}{0.7\textwidth}{X} 耶和華對亞伯蘭說：「你要離開本地、本族、父家，往我所要指示你的地去。 \end{tabularx} \\ \\ \relax
12:2 & \begin{tabularx}{0.7\textwidth}{X} 我必使你成為大國，我必賜福給你，使你的名為大；你要使別人得福。 \end{tabularx} \\ \\ \relax
12:3 & \begin{tabularx}{0.7\textwidth}{X} 為你祝福的，我必賜福給他；詛咒你的，我必詛咒他。地上的萬族都必因你得福。」 \end{tabularx} \\ \\ \relax
12:4 & \begin{tabularx}{0.7\textwidth}{X} 亞伯蘭就遵照耶和華的吩咐去了；羅得也和他同去。亞伯蘭離開哈蘭的時候年七十五歲。 \end{tabularx} \\ \\ \relax
12:5 & \begin{tabularx}{0.7\textwidth}{X} 亞伯蘭帶著他妻子撒萊和姪兒羅得，以及他們在哈蘭積蓄的財物、獲得的人口，往迦南地去。他們就來到了迦南地。 \end{tabularx} \\ \\ \relax
12:6 & \begin{tabularx}{0.7\textwidth}{X} 亞伯蘭經過那地，直到示劍地方，摩利橡樹那裡；當時迦南人住在那地。 \end{tabularx} \\ \\ \relax
12:7 & \begin{tabularx}{0.7\textwidth}{X} 耶和華向亞伯蘭顯現，說：「我要把這地賜給你的後裔。」亞伯蘭就在那裡為向他顯現的耶和華築了一座壇。 \end{tabularx} \\ \\ \relax
12:8 & \begin{tabularx}{0.7\textwidth}{X} 從那裡他又遷到伯特利東邊的山，支搭帳棚；西邊是伯特利，東邊是艾。他在那裡又為耶和華築了一座壇，求告耶和華的名。 \end{tabularx} \\ \\ \relax
12:9 & \begin{tabularx}{0.7\textwidth}{X} 後來亞伯蘭漸漸遷往尼革夫去。 \end{tabularx} \\ \\
[1ex]
\hline
\hline
\end{longtable}
$^{1}$各位弟兄姊妹平安.
很開心.
首先多謝教會和林牧師的邀請.
讓我有機會在大家當中.
和大家一起分享神學主人.
首先其實我應該.
我覺得我可以這樣說.
我也是自己人.
因為我在這裡聚會的.
但是今天有一點特別.
因為可以說是代表崇基學院神學院.
所以我先戴另一頂帽子.
首先我代表崇基學院神學院.
代表我們的同學.
一眾的師生向大家問安.
多謝大家一直以來對神學教育的支持.
沒有了你們.
其實神學教育是很難走下去的.
因著神學教育.
想和大家一起分享一個題目.
「照命人生」.
當然你會說讀神學一定要講照命.
沒有照命沒有呼召.
怎樣去讀神學呢?.
但其實照命應該是每一個信徒都有的東西.
為什麼這樣說呢?.
就是教會這個字.
教會這個字原來的意思就是.
被呼召出來的一群人.
所以如果你在教會.
你在教會聚會.
你已經是affirm一樣東西.
我是被呼召出來的.
究竟我們被呼召出來做些什麼呢?.
上帝怎樣呼召我.
我又怎樣去回應上帝給我的呼召呢?.
今天剛才和大家讀的一段經文.
正正就是在說呼召一段很重要的經文.
如果我們說創世記的話.
其實這一段可能是.

$^{41}$可以說是創世記裡面的轉捩點.
簡單說一點背景.
創世記有五十章.
創世記不只是說創世的.
如果你看看其實第二章已經說完了.
我們可以說第一章至第十一章.
是在說創世續章的時間.
但是由第十二章至到第二十五章.
是在說一個人.
就是亞伯拉罕.
第二十六章至到第三十六章.
就在說雅各.
然後第三十七章至到五十章.
就在說約瑟.
所以這樣說你會發覺其實創世記.
在很早的時間是說很宏闊的一個圖畫.
但是好像拍攝的話.
他就開始聚焦聚焦聚焦在一個人.
一個民族.
那個就是誰呢?.
那個就是亞伯拉罕的家族.
第十二章正正就是這個開始由很闊的一個環境.
聚焦在一個小的地方.
就是剛才我們剛剛讀的第十二章.
今天在第十二章一至到第九節.
我會分開兩段的時間.
和大家思考兩個方向.
可以這樣說.
我們說第一至到第四章.
我和大家思考一個課題.
什麼叫做照命的人生呢?.
我會歸納為一個旅途的人生.
一個擁抱未知的人生.
一個冒險的人生.
第二點就是一個捉壇的人生.
我想和大家再讀多一次.
創世記第十二章第一至到第四節.
不知道會不會太小.
如果你很喜歡看聖經都可以.
如果是聖經的話就是第十五頁那裡.

$^{81}$如果不是都可以一起看螢光幕.
創世記第十二章一至到第四節.
我們請大家一起讀.
預備起.
耶和華對亞伯蘭說:你要離開本地,本族,父家,和我所要指示你的地去.
我必使你成為大國,我必賜福給你,使你的名為大,你要使別人得福.
我必祝福與他助你的,我必助助他,地上的萬族都要因你得福.
亞伯蘭就遵照耶和華的吩咐去了,羅德也和他同去.
亞伯蘭離開哈蘭的時候年七十五歲.
這裡有一個很特別的地方,如果你剛才相信我說的.
第十二章就是介紹亞伯蘭出場的地方.
有什麼特別呢?.
你對比一下,一個大人物出場的時候很多時候會交代很多背景.
你看摩西,摩西出場的時候交代他BB仔在籃子裡.
但是亞伯蘭呢?亞伯蘭是一個最重要的人物.
他在第十一章從來沒有交代他任何的背景.
一出來就說他是一個老人家.
如果我們覺得我這麼久以來都沒有被上帝呼召,你還是很年輕.
亞伯蘭七十五歲才被上帝呼召.
而且你發覺,如果我們說亞伯蘭這樣的一個經歷是代表著所有人.
因為他在聖經裡可以說是被呼召的一個榜樣.
一個模式,一個象徵的話.
他的呼召很應該都可以應用在我們身上.
其中另一個很特別的就是,其實亞伯蘭沒有說過任何話.
第一至第四章正是耶和華的呼召,耶和華的說話.
很多人在這裡歸納我們叫做五重的福氣.
就是說怎樣賜福於你等等.
然後亞伯蘭就去,就沒有說任何話.
沒有說任何的excuse,沒有計算任何的代價.
然後就去了.
今天我們說面對上帝,當上帝的話臨到我.
或者你面對不同的處境.
可能很多時候我們都想,我們都不需要計算這麼多.
另外這裡其中一個重要的地方,我覺得就是在第一節.
耶和華對亞伯蘭說,你要離開本地.
本族富家,往我所要指示你的地去.
上帝呼召亞伯蘭,亞伯蘭回應.
亞伯蘭去哪裡呢?.
如果你在主日學那裡,你立即回答.
亞伯蘭就去迦南,他去迦南地.

$^{121}$上帝呼召他去迦南地.
但是你看清楚聖經這裡.
其實耶和華沒有叫他去迦南地.
沒有叫他去迦南地.
這個我不夠膽說,絕對是這樣說.
但是我的看法,上帝真的沒有叫他去迦南地.
你如果不相信的話,你看看其實為什麼他去迦南地.
你如果有聖經,我沒有放在這裡.
你有機會打開聖經,不然你就聽我讀.
第十二章,我們剛才說第一至第九節是這樣說.
你看清楚前兩節,他就這樣說.
創世記第十一章,第三十一至第三十二節.
你有機會你看看,不然就請聽我讀.
那裡就交代了再早一點的時間發生了什麼事.
他就說,他拉帶著他兒子亞伯蘭和他孫子.
即是他爸爸叫他拉,他就帶著亞伯蘭.
哈蘭的兒子羅德,以及他的媳婦.
亞伯蘭的妻子撒萊,一同出了迦勒底的梧椅.
要往迦南地去,他們來到哈蘭就住在那裡.
他拉共活了二百零五年,就死在哈蘭.
如果你將這一段經文放進去的話.
我們可以得到一個這樣的結論.
這裡有一幅圖細了一點.
原本他們就在右手邊的位置.
不只是亞伯蘭,是他的父親帶他們上了去哈蘭.
即是最高的位置.
他們的目的地是去迦南.
可能是因為那個地方好.
他就坐在那裡,住在那裡,住了一段時間.
沒有動過.
接著在第十二章那裡.
他說亞伯蘭就叫亞伯蘭.
你要離開本地本族父家.
往我所要指示你的地去.
往我所要指示的地去.
再簡單,再解釋得清楚是什麼呢.
你離開,你去吧.
去哪裡?.
去我遲些告訴你的地方.
在哪裡?.

$^{161}$你遲些就知道了.
你先去吧.
你是亞伯蘭,你去哪裡?.
你可能真的站在那裡.
你是亞伯蘭,你去哪裡?.
你按照你原先想要去的地方.
你繼續去吧.
他們原本根本上就是要去迦南地.
弟兄姊妹.
有時候我們想什麼叫做照命.
我們會有一個很美麗的圖畫.
照命就是突然間上帝和我說.
你轉跑道吧.
我不是否定這件事.
你轉跑道吧,你去觀塘兩法塘吧.
當然有很多是阿們的.
如果不是的話.
你轉工場吧.
你去轉不同的工作崗位.
因為這個是我特定預備了你的.
很多時候我們說.
上帝日如果我等到這個很明確的聲音.
我就去了.
如果我沒有的時間.
我就不去了.
但你看看亞伯蘭.
亞伯蘭他出去的時候.
其實他真的不知道去哪裡.
他真的不知道去哪裡.
在希伯來書的一段經文.
剛才正正說得好.
希伯來書第十一章那裡他怎麼說的呢.
希伯來書第十一章第八節.
他說什麼呢.
他說亞伯拉罕恩爵信蒙召的時候.
就遵命出去.
去哪裡呢.
往將來所要得為業的地方去.
出去的時候.
還不知往哪裡去.

$^{201}$召命不是等於我們完全明確.
我們要去哪裡.
當我們還不知往哪裡去的時候.
可以仍然都是在召命的當中.
大家有沒有經歷過還不知往哪裡去那種情況呢.
我有少少的經歷.
不知道大家可能都會有相似的經歷.
就是怎樣呢.
我之前在德國住了七年.
那時候其實我是一個什麼人呢.
我是一個窮學生.
但在德國七年很悶的.
雖然是一個窮學生.
我們都會想.
我們都有權利去旅行.
只是一個叫做窮遊的時候不流行.
但在德國有一個這樣的計劃.
有一個廉價航空.
German Wings.
當時真的做到一個這樣的價格.
30歐.
30歐你可以做什麼.
吃了餐就沒有了.
30歐我在德國南部.
30歐可以做到什麼.
你可以去漢堡.
你可以去應該丹麥都可以.
你可以去歐洲一些很厲害的城市.
米蘭.
去意大利羅馬.
你也可以去馬德里.
去巴塞隆拿.
去任何這些城市的其中一個.
因為你是不知道去哪裡的.
你可以訂了你的日期.
但是他應該是有三二十幾個這樣的地方.
你就沒得選擇.
他應該去到最後兩天.
他才告訴你.
為了什麼呢.

$^{241}$他就說如果你想選擇.
你可以刪除三個.
每次就加10歐還是15歐.
我加給你我就不是叫窮學生.
所以我是不加的.
其實就簡單來說.
因為我另外一個身份就是我蘇宣的老公.
蘇宣很活潑的.
其實他很想去的.
所以他就不斷說我們去吧.
我記得有一段時間.
他每一天就看著螢幕.
就好像看股票一樣.
我看來幹嘛.
你看完之後他都不會減了.
最後其實我們都去不成.
因為太害怕.
為什麼呢.
這是我的問題.
我最後總結為什麼呢.
因為我已經太老了.
我受不住那個恐懼的情況.
因為他不單止.
你不是去的地方這麼簡單.
因為譬如有些說去波蘭的.
去波蘭就你不是直航的.
你去到一個地方.
你搭17個小時巴士.
你都要照吃的.
那你帶些什麼呢.
人家說波蘭可能會被人打劫的.
最後都去不成.
為什麼呢.
因為我還不知往哪裡去那種恐懼.
籠罩著我的心.
聽姐妹.
其實我們的信仰可貴的地方.
不是我們知道去哪裡.
是當我們還不知往哪裡去的時候.
其實上帝就已經呼召了我們.

$^{281}$你看下亞伯蘭.
你再看看亞伯蘭.
亞伯蘭出了去的時間.
剛才我們讀的.
一會兒我們再讀.
你會發覺他去了很多地方.
如果他真的知道自己去哪裡.
其實他一去到那個地方就停在那裡.
去了士堅.
去了伯特利.
接著再去了希伯倫.
可以這樣說.
他根本上不知道自己去哪裡.
正如我剛才所說.
呼召某程度上不是一個很實在的情況.
很多時候我們會覺得.
上帝啊.
你給一個很實在很明確的地方.
我就確定這個是你的呼召.
反過來說.
如果我現在經歷很不穩定.
我的工作.
我的家庭.
我的人生.
好像都是一些混亂的情況.
我們會問.
上帝為什麼.
你在哪裡呢.
我們會覺得.
這個某程度上應該不是上帝在呼召我.
這個是我和上帝不能夠結連.
但是你看看亞伯蘭.
亞伯蘭.
如果我們說他剛才進入迦南的時間.
沒有時間和你看.
但是如果你看看第十二章下半段.
你會發覺亞伯蘭正正.
當他被呼召進入迦南地的時間.
他經歷他人生很多很艱難的時間.
去到.

$^{321}$剛才讀他漸漸往尼格夫去.
為什麼呢.
後來就因為經歷饑荒.
他去埃及地.
他大財主來的.
如果他留在哈蘭.
某程度上他不用經歷饑荒.
去到埃及地的時間.
他做了一件事.
就是認妻作妹.
人生最大的污點.
就是在他被呼召之後發生.
上帝啊.
我如果不是經歷呼召.
我留在哈蘭.
我人生不會有這個污點.
他唯一的親人是羅德.
在第十三章你就會見到.
他分家產.
他當時沒有小朋友.
他沒有後代.
羅德就是他最親的人.
但是羅德為了財產.
與亞伯蘭他唯一的親人反目成仇.
如果你是亞伯蘭.
你可能說.
上帝這個是你的心意嗎.
你不是呼召了我.
為什麼我還經歷一些.
這麼艱難的時間.
兄弟姐妹.
可能這個也是我們的寫照.
我投靠上帝.
我覺得我是在進行上帝的心意.
為什麼我所經歷的不是一帆風順.
我所經歷的不是像上帝賜福給我.
這麼道路亨通.
如果我們說亞伯蘭何時才叫做守得雲開.
你慢慢去讀.
去到第十七章.

$^{361}$亞伯蘭就上帝.
他真的定居了.
然後上帝就說.
你叫做亞伯拉罕.
因為我要賜福給你.
我甚至乎要賜你有你的後代.
當時聖經加了個註腳.
他多少歲呢 年九十九歲.
即是說他由不穩定到慢慢穩定的時間.
他經歷了二十四年.
弟兄姊妹 你這裡有兩個抉擇.
你可以說.
其實上帝的呼召只是在九十九歲的時候才發生.
如果是的話 你要解釋究竟亞伯蘭這二十四年是發生過什麼事呢.
八卦的.
他根本上不在上帝的心意當中.
這樣說是不合理的.
你可以想像.
我們只可以得出一個結論就是.
其實上帝的呼召 上帝的召命.
不等於我們是一帆風順.
或者上帝的召命不等於.
我可以很清晰 很確定.
不懷疑這樣去走我的人生的路.
或者從正面的這樣說.
當你面對艱難的時間.
其實我們都可以憑信心這樣說.
上帝啊 亞伯蘭在艱難的時間.
他一直都在.
所以現在我面對艱難.
我甚至乎覺得我錯失了一些時機.
無論是在我的事業 在我的家庭 在我的事奉 在我的不同的關係當中.
我覺得有些地方我是否錯失了呢.
我們可以有一種信心的勇氣.
我相信上帝啊 你一直都在.
你仍然都可以化腐朽為神奇.
這是第一點亞伯蘭給我們的鼓勵.
第二點 我們看第五至到第九節.
再想請大家和我讀下半段這段經文.
創帝記第十二章第五至到第九節.

$^{401}$預備起.
亞伯蘭帶著他妻子撒萊和質爾羅德.
以及他們在哈蘭積蓄的財物.
獲得的人口往迦南地區.
他們就來到了迦南地.
亞伯蘭經過那地 直到事件地方.
摩利像樹那裡 當時迦南人住在那地.
耶和華向亞伯蘭顯現說.
我要把這地賜給你的後裔.
亞伯蘭就在那裡向為他顯現的耶和華捉了一座壇.
從那裡他又遷到伯特利東邊的山.
支塔象棚 西邊是伯特利 東邊是艾.
他在那裡又為耶和華捉了一座壇.
求告耶和華的名.
後來亞伯蘭又漸漸遷往格尼夫去.
這裡我們可以看到當亞伯蘭真的離開了哈蘭.
進入迦南地的地方.
你會發覺很混亂的.
因為很多地名.
其實不是你的問題 是亞伯蘭的問題.
亞伯蘭他去了不同的地方.
然後你可以說是一個漂泊的人生.
或者叫做搭象棚的人生.
但我會更加覺得.
如果你留意這裡.
正正就是第七至第九節.
或者是第七至第八節.
都有一個字是出現了兩次.
就是什麼呢.
他為耶和華捉了一座壇.
如果你再看後一點.
你會發覺這個pattern重複出現.
去到第十三章的時間.
希伯倫的時間.
他再到了希伯倫.
他又寫他為耶和華捉了一座壇.
去到第二十二章.
這個很出名的.
就是獻以殺的時間.
他為耶和華捉了一座壇.

$^{441}$短短十四章的經文出現了四次.
而亞伯拉罕在不同的地方.
為耶和華捉了一座壇.
所以我們如果說.
亞伯蘭他的生命是一個照命的人生.
我可以說他帶給我們的啟示是.
如果我的人生.
可以在亞伯蘭的身上看到的話.
亞伯蘭的人生就是捉壇的人生.
我們的人生應該都是一個捉壇的人生.
捉壇的人生.
他這裡很清楚說.
做什麼的呢?是一個求告耶和華的命的人生.
什麼叫做捉壇.
或者我們看回當時候捉壇.
就不少的.
他們的捉壇.
我以前做童軍的.
你不知道有沒有人做童軍.
做童軍最厲害的就是起火.
起到火有少少煙已經很厲害.
他們那個一定是龍煙.
你想想他捉壇是為了什麼.
就是獻祭.
獻祭是整隻牛這樣燒的.
捉壇的人生.
如果你套入在當時候的時間.
其實是一個很牙煙的狀態.
為什麼呢?.
你想想亞伯蘭是一個什麼身份.
亞伯蘭是一個什麼人.
亞伯蘭是一個有錢人.
是真的.
不單止.
亞伯蘭當時候是一個新移民.
他是由哈蘭進入加拿大的一個新移民.
新移民最害怕或者最緊張的是什麼.
你想想.
一個新移民最緊張的是什麼.
就是融入當時候的社會.

$^{481}$我也有少少新移民的經驗.
我住了在德國七年.
我覺得自己很well prepared.
所以當我去到當地的時候.
我做了兩件事.
第一件事就是請當時候的人去我家.
煮東西給他吃.
和和他解釋一件事.
就是我是不吃狗的.
我是狗的朋友.
我是不吃狗的.
因為當時德文裡面.
狗就是Hund.
雞就是Hund.
差很少.
他們經常覺得你們中國人是吃Hund的.
就是吃我們的朋友的.
所以我一去到我就跟他說.
我就是吃雞不吃狗的.
所以我只是煮雞給你吃.
第二件事.
我就扔下自己的身份證.
去亞馬遜買了一雙球鞋.
叫他們侵我去踢球.
你要知道我當時都一把年紀了.
不是很老.
但是正值當時候.
就是德國拿世界盃的時間.
他們是雄心萬丈.
我不知道.
我以前中學生都有踢過一下球.
但是為什麼他們的球場好像大兩倍.
可能我自己幻覺.
每次猜人猜到最後是我的.
不要緊.
我志不在此.
但是我要告訴你.
我知道自己我是新移民.
所以我就跟他們踢球.
每一次他們說Cotton, Snow.

$^{521}$我就去.
總之我記得好幾次.
回來我就要搓鐵打酒.
為什麼.
因為我知道我是一個新移民.
新移民我最重要是有朋友.
我最重要是跟他們同一口徑.
我最重要是不要冒犯到他們.
亞伯拉罕.
他是一個什麼人.
他是一個新移民.
但是他正在做的事.
正正就是格格不入.
他去到那個地方.
他捉一座壇.
慌死人家不知道.
他是求告耶和華的命.
他捉一座壇.
當他燒的時間.
人家只會覺得.
你是不是在挑機.
你為什麼要這麼大的火.
你是不是想打仗.
你帶著一群人.
你是不是想挑戰我們已有的文化.
已有的體制.
你是不是想搞轉那個派對.
兄弟姐妹.
如果我們說.
我們是被上帝呼召的話.
我們都應該要有一個確信.
就是我們被呼召出來.
其實是與世界不同.
我們是與世界有一種格格不入.
那種的狀態.
你看楊福音.
楊福音說.
光來到世間.
世界不接受光.
因為世界是黑暗的.

$^{561}$可能你說.
不是啊.
我們和這個世界.
我們要相適應.
我們不可以走回頭路.
我們不可以故步自封.
因為在古代.
你說我們不是走.
好像古代monasticism.
修院的狀態.
修院就是格格不入.
我們只是求告上帝.
在修院裡面.
因為我們已經和世界結連.
這個是真的.
但是在十六世紀.
宗教改革.
其中一個最大的貢獻是什麼?.
嘉義民.
因為我們是嘉義民傳統.
當然中華基督教會.
有很多不同的傳統.
但是中華基督教會改革中.
嘉義民給我們一個傳統.
一個概念是什麼呢?.
worldly acceptism.
世界就是我修道的場所.
我不是不修道.
只不過我是面對這個世界.
去操練我的格格不入.
我為上帝畫了一個神聖的空間.
被人看到.
這個是上帝.
當我正在做的事情的時候.
別人看到.
上帝就在當中.
這個是我們的心智.
是我們的操練.
是我們的照命的一部分.
神學教育.

$^{601}$同濟學院,神學院.
我們說我們要訓練識見的僕人.
我們不單要做僕人.
我們要有識見.
在多元的世界.
我怎樣探索.
我怎樣去尋問.
我怎樣去見證上帝.
在這個世界.
我怎樣靈巧像蛇.
但是純良像甲子.
我怎樣在多元的世界.
顯示得到.
在公共的空間裡.
上帝啊.
我們為你築起神聖的壇.
這個是很需要智慧.
很需要上帝的恩典.
也很需要教會不同的弟兄姊妹.
我們一起來參與.
因為這個.
不只是給神學院的照命.
這個是給教會的照命.
是你和我共同承擔的照命.
我們一起禱告.
願上帝幫助我們.
亞伯拉罕成為我們的榜樣.
面對很艱難的時間.
給我們有信心的勇氣.
上帝啊.
你一直都在.
你參與在我們的生命.
你可以化腐朽為神奇.
因為你一直帶領我們.
上帝也給我們有見證你.
創造神聖空間的勇氣.
我們的壇.
願上帝你閱立.
我們求你也給我們有智慧.
梁法堂有智慧.

$^{641}$神學院有智慧.
怎樣在這個世代當中.
在很多世界指鹿為馬的時間.
我們堅持耶穌你講說就說事.
當這個世界不斷在紛爭.
在鬥爭的當中.
上帝你教我們去傳播愛.
當這個世界在失望.
絕望的境地.
願我們成為眾人的盼望.
求上帝加力.
以我們成為你照命彰顯的地方.
奉耶穌名求.
阿們.
\newpage



\section{創世記 37:12-36}
\label{sec:031805}
\textbf{聖經人聖經事:上帝的手在帶領}~~~~黃克勤牧師博士
\newline
\newline
連結: \href{http://www.church.com.hk/acms/content.asp?site=cdc\&op=show\&type=product\&code=031805\&layout=sermon}{\texttt{031805}} ~~~~ 語音日期: 2026-03-12
\newline
\newline
\hyperref[sec:031804]{\small{< < < PREV SERMON < < <}}
~
\hyperref[sec:index_chronic]{\small{[返順時目]}}
~
\hyperref[sec:031806]{\small{> > > NEXT SERMON > > >}}
\newline
\newline
創世記 37:12-36
\newline
\begin{longtable}{cl}
\hline
\hline
章節 & 經文 (和合本修訂版)\\
\hline
37:12 & \begin{tabularx}{0.7\textwidth}{X} 約瑟的哥哥們到示劍去放他們父親的羊。 \end{tabularx} \\ \\ \relax
37:13 & \begin{tabularx}{0.7\textwidth}{X} 以色列對約瑟說：「你哥哥們不是在示劍放羊嗎？來，我派你到他們那裡去。」約瑟對他說：「我在這裡。」 \end{tabularx} \\ \\ \relax
37:14 & \begin{tabularx}{0.7\textwidth}{X} 以色列對他說：「你去看看你哥哥們是否平安，羊群是否平安，再回來告訴我。」於是他派約瑟出希伯崙谷，約瑟就往示劍去了。 \end{tabularx} \\ \\ \relax
37:15 & \begin{tabularx}{0.7\textwidth}{X} 有人遇見他，看哪，他在田野走迷了路。那人問他說：「你找甚麼？」 \end{tabularx} \\ \\ \relax
37:16 & \begin{tabularx}{0.7\textwidth}{X} 他說：「我找我的哥哥們，請告訴我，他們在哪裡放羊。」 \end{tabularx} \\ \\ \relax
37:17 & \begin{tabularx}{0.7\textwidth}{X} 那人說：「他們已經離開這裡走了，我聽見他們說：『我們往多坍去。』」約瑟就去追哥哥們，在多坍找到了他們。 \end{tabularx} \\ \\ \relax
37:18 & \begin{tabularx}{0.7\textwidth}{X} 他們遠遠看見他，趁他還沒有走近他們，就圖謀要殺死他。 \end{tabularx} \\ \\ \relax
37:19 & \begin{tabularx}{0.7\textwidth}{X} 他們彼此說：「看哪！那做夢的來了。 \end{tabularx} \\ \\ \relax
37:20 & \begin{tabularx}{0.7\textwidth}{X} 現在，來吧！我們把他殺了，丟在一個坑裡，就說有惡獸把他吃了。我們且看他的夢將來怎麼樣。」 \end{tabularx} \\ \\ \relax
37:21 & \begin{tabularx}{0.7\textwidth}{X} 呂便聽見了，要救約瑟脫離他們的手，說：「我們不可害他的性命。」 \end{tabularx} \\ \\ \relax
37:22 & \begin{tabularx}{0.7\textwidth}{X} 呂便又對他們說：「不可流他的血，可以把他丟在這曠野的坑裡，不可下手害他。」呂便要救他脫離他們的手，把他還給他父親。 \end{tabularx} \\ \\ \relax
37:23 & \begin{tabularx}{0.7\textwidth}{X} 約瑟到了他哥哥們那裡，他們就剝去他的外衣，就是他身上那件長袍。 \end{tabularx} \\ \\ \relax
37:24 & \begin{tabularx}{0.7\textwidth}{X} 他們抓住他，把他丟在坑裡。那坑是空的，裡頭沒有水。 \end{tabularx} \\ \\ \relax
37:25 & \begin{tabularx}{0.7\textwidth}{X} 他們坐下吃飯，舉目觀看，看哪，有一群以實瑪利人從基列來，用駱駝馱著香料、乳香、沒藥，要帶下埃及去。 \end{tabularx} \\ \\ \relax
37:26 & \begin{tabularx}{0.7\textwidth}{X} 猶大對他的兄弟們說：「我們殺我們的弟弟，遮掩他的血有甚麼好處呢？ \end{tabularx} \\ \\ \relax
37:27 & \begin{tabularx}{0.7\textwidth}{X} 來，我們把他賣給以實瑪利人，不要下手害他，因為他是我們的弟弟，我們的骨肉。」他的兄弟們就聽從了他。 \end{tabularx} \\ \\ \relax
37:28 & \begin{tabularx}{0.7\textwidth}{X} 那時，有些米甸的商人從那裡經過，就把約瑟從坑裡拉上來。他們以二十塊銀子把約瑟賣給以實瑪利人，他們就把約瑟帶到埃及去了。 \end{tabularx} \\ \\ \relax
37:29 & \begin{tabularx}{0.7\textwidth}{X} 呂便回到坑旁，看哪，約瑟不在坑裡，就撕裂自己的衣服， \end{tabularx} \\ \\ \relax
37:30 & \begin{tabularx}{0.7\textwidth}{X} 回到他兄弟們那裡，說：「孩子不在了。我往哪裡去才好呢？」 \end{tabularx} \\ \\ \relax
37:31 & \begin{tabularx}{0.7\textwidth}{X} 於是，他們宰了一隻公山羊，拿了約瑟的那件外衣染上了血， \end{tabularx} \\ \\ \relax
37:32 & \begin{tabularx}{0.7\textwidth}{X} 派人把長袍送到他們的父親那裡，說：「我們發現這個， 請認一認，是不是你兒子的外衣？」 \end{tabularx} \\ \\ \relax
37:33 & \begin{tabularx}{0.7\textwidth}{X} 他認出來，就說：「這是我兒子的外衣，惡獸把他吃了，約瑟一定被撕碎了！」 \end{tabularx} \\ \\ \relax
37:34 & \begin{tabularx}{0.7\textwidth}{X} 雅各就撕裂衣服，腰間圍上麻布，為他兒子哀傷了多日。 \end{tabularx} \\ \\ \relax
37:35 & \begin{tabularx}{0.7\textwidth}{X} 他的兒女都起來安慰他，他卻不肯受安慰，說：「我必哀傷著下陰間，到我兒子那裡。」約瑟的父親就為他哀哭。 \end{tabularx} \\ \\ \relax
37:36 & \begin{tabularx}{0.7\textwidth}{X} 米甸人把約瑟賣到埃及，給法老的官員，就是護衛長波提乏。 \end{tabularx} \\ \\
[1ex]
\hline
\hline
\end{longtable}
$^{1}$講這些故事是很開心的.
昨天講了藥室的採醫與發夢.
今天的聖經人聖經士生命說書網絡為大家預備的.
就是上主的手在帶領.
我們看聖經的人或聖經的事最重要是看到.
背後上帝的手在做什麼.
昨天藥室又有採醫又發夢.
上帝的手你看到什麼?.
上帝的手是帶領藥室.
他有一個出生的困境是什麼呢?.
小兒子而且爸爸疼愛有件採醫.
這樣害死他了.
他又不懂做不懂死.
不斷講哥哥的壞話.
還有一點點直和高傲.
向我下拜.
父母都向我下拜.
我的夢是這樣的.
你們要拜我的.
你不是說排第幾.
以兄弟來說你排第十一.
不要以為你很厲害.
不是一是十一.
所以哥哥不喜歡他.
不過上帝的手一直帶著藥室在他身上發生很多事.
我們看看上帝的手在藥室身上做什麼.
創世記37章繼續講創世記12節.
藥室的哥哥往事劍去放他父親的羊.
以色列即是雅各對藥室說.
你哥哥們不是在事劍放羊嗎?.
你來.
我要打發你往他們那裡去.
藥室說我在這裡.
藥室都很乖.
爸爸說藥室你來.
藥室就說我來我在這裡.
挺有趣的.
藥室的哥哥去了事劍放父親的羊.
藥室有沒有去呢?.
沒有.

$^{41}$藥室不用去.
你說哥哥是不是很生氣.
爸爸真的很偏心.
我們出來日曬雨淋.
辛苦到極.
也沒有一口好吃.
有時要和野獸搏鬥.
就說他就好不用出來.
留意這是上帝的工作.
令到雅各的心很不安.
不安就是神的手在後面做事.
怎樣不安呢?.
聽說你哥哥去了事劍放羊.
你和我去看看他們.
藥室都孝順的.
爸爸叫我一定做.
以色列楊雅各說.
留意你去看看你哥哥們.
平安不平安.
群羊平安不平安.
然後回來報信給我.
挺有趣的.
有點無稽的.
你去看看他們.
平不平安.
我掛著他們.
看他們平不平安.
看羊平不平安回來告訴我.
其實是不必的.
我認為沒有什麼需要.
也沒有實效.
因為如果兒子有一個損傷.
有些羊被野獸吃了或被人搶了.
十個兒子十個哥哥.
一定會有人回來報信.
而且不止十個.
有些牧羊人跟著去.
一定有人回來報信.
何須去找他們呢?.
你去找他們.

$^{81}$我掛著他們.
看他們平不平安.
然後回來報信.
我想知道.
其實會不會.
雅各都想約瑟出去闖蕩一下呢?.
應該不會.
因為如果他真的想約瑟出去.
看看羊.
認識一下曠野生活.
跟一些哥哥去照顧.
怎會叫他一個人孤零零在曠野行走.
其實很危險.
雅各一定知道.
不過心裡可能是神令他心不安.
於是他忍不住.
沒有人了.
其實可以叫其他牧羊人.
他沒有叫.
於是找一個心腹.
約瑟你出來找他們.
告訴我他們好不好.
平不平安.
於是打發他出去.
出希伯倫谷去視劍.
約瑟走路不習慣.
走到某程度就走迷了路.
其實都幾難.
又沒有北斗導航.
又沒有衛星導航.
又沒有地圖.
是靠自己方向感.
看看太陽在那邊.
就知道方向感.
向著那邊就去那邊.
見到希伯倫那邊見到視劍那邊.
他就去了.
在田野走迷了路.
好在上帝的工作.
上帝的手在那裡做事.

$^{121}$有人遇見他在田野走迷了路.
遇見他.
如果是好人好者.
大家在曠野遇見都會見個面.
你好嗎?.
哪裡來的?.
有什麼要幫手嗎?.
多好啊.
我有些水.
見你的口好乾.
給你喝吧.
他們很幫手.
大家知道在曠野裡面的生活.
田野間是很難的.
走迷了路.
頭餓又沒有路牌.
又沒有電話.
於是有人遇見他.
喂.
童子小朋友.
為什麼你一個人在這裡.
我去找我哥.
我去視劍.
不過走迷了路.
不知走哪邊.
你教我走哪邊.
視劍.
那些人很好幫手.
你找什麼?.
視劍.
你去視劍找什麼?.
我去找我哥.
求你告訴我們.
他們在何處放羊.
這個問題是不是比較搞笑.
純粹是一個生手.
沒有牧過羊的人.
沒有走過曠野.
若實話.
我去找我哥.

$^{161}$好多人.
你告訴我.
他們在何處放羊.
喂.
在田野曠野中間.
在何處放羊.
你哥在何處放羊.
我怎知誰是你哥.
我又沒有寫你哥.
我怎知誰是你哥.
他不知.
但他在這裡.
我找不到我哥.
即是小朋友走失了.
你找誰?.
我找我媽.
你媽是怎樣的?.
有眼耳口鼻.
有頭髮的.
這樣不行嗎?.
我找我哥.
求你告訴我們.
他們在何處放羊.
這麼沒sense的問題.
怎知上帝的手真厲害.
那些人就說.
哦.
那班我知道.
我剛剛碰到他們.
他們已經走了.
不在視劍了.
而且我們見過面.
我聽到他們說.
幸好我聽到.
我聽到哥說話.
遇到人聊天的人.
剛剛又碰到藥室.
不然真的不知怎辦.
我聽見他們說要往多山.
多山哪邊?.

$^{201}$那邊那邊.
於是藥室就說.
我去追他們.
你追快點.
因為他們走了很久.
都不知是否昨天見到.
藥室就追趕他哥哥.
遇見他們真的在多山.
跟住這個地方追過去.
遇見他們.
有群人在前面.
那麼多人像我哥一樣走過去.
他們遠遠見到他.
為何有一個人在走.
穿彩色衣服.
藥室來的.
看清楚一點.
是呀是呀.
他們很遠見到他.
十個哥哥.
就說趁他還沒有來到庚前.
不如我們想想同謀.
怎樣弄死他在曠野.
這麼狠毒.
所以你見到這些哥哥是否很恨他.
再做夢之後.
越發恨他想他死.
他們彼此說.
那個做夢的人來了.
我們把他殺了.
扔進坑裡.
回去說有惡獸把他吃了.
看他還做什麼夢.
且看他的夢將來怎樣.
弄死他我看你做什麼夢.
其實你見到十個哥哥是很憎恨他.
其中有些特別憎恨他.
就是雅各的妾兒特別憎恨他.
好似你一個半個愛邦人.
特別憎恨他.

$^{241}$裡辯聽見了.
裡辯是誰呢.
裡辯以前中文翻譯為「流辯」.
即是大哥.
雅各的第一個兒子.
流辯沒聽話.
這麼難聽.
於是翻譯為裡辯.
不過裡辯挺有趣.
為何不翻譯為外辯呢.
裡辯.
裡辯聽見.
裡辯大哥流辯很深.
要救他脫離他們的仇.
我們不要害他性命.
不要流他的血.
不如這樣吧.
他不能夠翻譯九個弟弟的意思.
那怎麼好呢.
不要流他的血.
我們將他扔在坑裡.
不要親手害死他.
扔在坑裡即是甚麼.
爬不上來.
扔在坑裡會怎樣.
由他在裡面慢慢餓死.
死也是.
九個弟弟會否喜歡呢.
都喜歡.
不過有甚麼好處呢.
就是我們不用親手動手殺他.
不要親手動手.
由他餓死渴死.
明白嗎.
所以不要親手殺他.
裡面的意思是要救他.
裡面是要救他.
神的手安排了裡面是要救他.
要救他.
不要動手.

$^{281}$其實他想甚麼呢.
要救他脫離他們的仇.
把他歸還他的父親.
裡面是想要救弱勢.
我們先扔他在坑裡.
然後由他餓死我們離開.
不過他可能是屈尾十.
就叫僕人救他上來.
帶他回家.
又或者回到家.
哎呀弱勢呢.
不知在哪裡.
知不知在哪裡.
沒見到我叫他找你們.
沒見到嗎.
沒有.
然後可能靜靜和爸爸說.
這件事是這樣的.
他在哪裡東南西北多山附近.
有個洞快點叫僕人去救他.
快呀不然來不及.
他想歸還父親救他.
弱勢到他高高納累.
他們就抓住他脫下彩衣.
然後就把他衣服拿走.
然後就穿上底衣或帽衫.
推他到坑裡.
幸好那坑是空的.
裡面沒有水.
幸好是空的.
即是說是乾的.
又沒有野獸.
又沒有蛇.
裡面又沒有水.
不然會有水噴進去.
淹死.
或者半缸水.
你都幾慘.
他就是這樣.
於是就扔了進去.

$^{321}$在坑裡面是空的.
又沒有水.
經常都在那裡工作.
第25節.
他們坐下吃飯.
剛剛到那裡吃飯的時候.
在那裡看到有樂陀.
一隊人經過樂陀隊.
看他們穿的衣服似米殿人.
米殿是地方名.
其實他們的種族是以實馬尼.
都是半個親戚.
以實馬尼是亞伯拉罕的另一個兒子.
所以和亞國這班兒子.
都是有少少關係.
都是半個親戚.
「喂喂喂,你們是哪裡來的?」.
「我們是在基列來的」.
「現在用樂陀載著香料」.
「乳香,抹藥」.
「去埃及賣」.
基列是很出名出乳香的.
他們帶著香料去埃及賣.
親戚親戚.
於是大家就講起來了.
尤大就說.
「不如這樣吧」.
「和其他兄弟講」.
「我們殺了我們的兄弟」.
「藏了他的血有什麼益處呢?」.
「如果我們殺了我們的弟弟」.
「流了血」.
「我們本身做了一件壞事」.
「那有什麼好呢?」.
「我們隱藏著這件事」.
「都有血了」.
「我們都變了殺人兇手」.
「不好啊,不如這樣吧」.
「我們賣了他給以實馬尼人」.
「賣給親戚」.

$^{361}$「我們不要下手害他」.
「賣了他又撿回多少」.
「一來大家兄弟骨肉」.
「不要相殘」.
「買了他撿回多少」.
「大家買酒喝下」.
「聊聊天」.
「吃吃東西」.
兄弟們就說.
「這樣也好」.
「其實也好」.
「殺了弟弟也不好」.
「我怕神懲罰我們」.
「那不好了」.
「賣了他也好」.
「好,賣了」.
「米田商人經過」.
「於是哥哥把他從坑裡拉上來」.
「二十四十萬銀子賣給了以實馬尼人」.
「他們把弱勢帶到埃及」.
「你看到神的手在工作」.
「又遇到這個又遇到那個」.
「哎呀,我們哥哥在哪裡」.
「我怎麼知道你哥哥在哪裡」.
「哎呀,在哪裡哪裡」.
「哦,快點去追,去了多山」.
「於是一件事一路下來」.
「每一個紅色字」.
「都是告訴你神的手在工作」.
「於是弱勢就開始去了埃及」.
「裡面可能沒有跟他一起吃飯」.
「走開了」.
「大哥裡面回來後就說」.
「o下,為什麼弱勢不在那個洞裡」.
「去了哪裡」.
「糟了,是不是被東西擔了」.
「怎麼辦」.
「於是就快點走回兄弟那裡」.
「喂,弟弟那個小孩不見了」.
「還不止小孩」.

$^{401}$「那個小孩不見了」.
「怎麼辦」.
「糟了,裡面已經想了一段時間」.
「他以為小孩不見了」.
「其實那些哥哥」.
「又沒有跟他說什麼」.
「他走了出去走一走」.
「回來後,哎呀,找弟弟不見了」.
「可能他也看看有沒有飯吃」.
「照顧他」.
「不見了怎麼辦」.
「他們就把衣服脫了」.
「丟在坑裡」.
「不如這樣」.
「我們找一隻公羊殺了」.
「然後坐在這裡染血」.
「然後回去告訴爸爸」.
「打發人,即是埋伏人」.
「送到父親那裡說」.
「哎呀,我們在田野撿到衣服」.
「你承認是不是弱勢的」.
「瓦割認得」.
「哎呀,真的是我做給弱勢的衣服」.
「有血,只剩下衣服」.
「全爛了」.
「當然有惡獸把他吃了」.
「弱勢被撕碎了」.
「哎呀,弱勢被惡獸吃了」.
「全散了,碎了」.
「哎呀,撕碎了撕碎了」.
「說兩次就說很慘」.
「撕裂衣服,腰間圍上麻布」.
「為他的兒子悲哀多日」.
「他的兒女都來安慰他」.
「不肯受安慰」.
「爸爸不要緊」.
「你還有十一個兒子」.
「還有很多個女兒」.
「不怕的,不要緊的」.
「神都有保守,有帶領的」.

$^{441}$「弱勢說哎呀」.
「雅各說怎麼辦好」.
「不行的,這個兒子這個兒子」.
「你們都是很乖」.
「哎呀,我沒了兒子」.
「我快死了,很老了」.
「悲哀著下陰間」.
「我要去找我兒子了」.
「到我兒子那裡」.
「我要悲哀死,我要去找我兒子」.
「弱勢的父親雅各為他哀哭」.
「米殿人帶了弱勢去埃及」.
「將他賣給法老的內臣」.
厲害了.
這個內臣是一個很親密很信任的大臣.
就是護衛長波提弗.
於是弱勢去了波提弗的家.
開始接近王室.
又有另一番遭遇.
這一天我想和大家看的是.
神的手在工作.
紅色的字都是神的手在工作.
弟子們,如果你發覺.
為何我有這麼多過山車經歷.
人生裡有高有低.
其實是神的手在你的身上工作.
祂一路預備你無助.
令你適合做某一件事.
有祂的心意.
一直都這麼說.
神在每一個人身上有一個計劃.
我自己都不知道.
我自己回望才發覺.
原來神在我身上有一個計劃.
原來我讀了這些書是有這樣的用途.
讀了這些書是有這樣的用途.
讀了國內上海長津大學的書.
令我在國內行走.
別人見到我在中國大學的高級學位.
覺得很厲害,馬上信任了.

$^{481}$於是令我在中國的工作更好.
其實是這樣發生的.
很多事情都是這樣建立.
其實背後都是神的手在動工.
因此將自己交託給主.
遇到什麼事情就求問神.
是不是你要我做的.
如果是,沒有反對,我做.
遲一點就會知道原來神這樣對我有益.
所有事都在祂手裡.
我終生的事都在你手中.
這樣我交託給主一起祈禱.
求主幫助我們.
讓我們更加懂得原來你的手.
到今天仍然在世界上工作.
從來沒有停.
你是色在,金在,以後永在的主.
你雖然在高天之上,但有時間聽每一個小孩的祈禱.
你在高天之上,但你可以管理,控制這個世界.
所有每一件或大或小的事情發生.
甚至我們每一步.
都是主能夠帶領著,而且是帶領著.
多謝我們的主.
讚美你,你的手正在工作,而且大能.
你運籌遺額,真的很厲害.
願你得造榮耀.
我們信靠你,信任你,祈禱奉耶穌名求.
阿們.
\newpage



\section{約翰福音 9:1-41}
\label{sec:031806}
\textbf{瞎眼今得看見}~~~~葉菁華教授
\newline
\newline
連結: \href{http://www.church.com.hk/acms/content.asp?site=cdc\&op=show\&type=product\&code=031806\&layout=sermon}{\texttt{031806}} ~~~~ 語音日期: 2026-03-15
\newline
\newline
\hyperref[sec:031805]{\small{< < < PREV SERMON < < <}}
~
\hyperref[sec:index_chronic]{\small{[返順時目]}}
~
\hyperref[sec:031807]{\small{> > > NEXT SERMON > > >}}
\newline
\newline
約翰福音 9:1-41
\newline
\begin{longtable}{cl}
\hline
\hline
章節 & 經文 (和合本修訂版)\\
\hline
9:1 & \begin{tabularx}{0.7\textwidth}{X} 耶穌往前走的時候，看見一個生來就失明的人。 \end{tabularx} \\ \\ \relax
9:2 & \begin{tabularx}{0.7\textwidth}{X} 門徒問耶穌：「拉比，這人生來失明，是誰犯了罪？是這人還是他的父母呢？」 \end{tabularx} \\ \\ \relax
9:3 & \begin{tabularx}{0.7\textwidth}{X} 耶穌回答：「既不是這人犯了罪，也不是他的父母，而是要在他身上顯出神的作為來。 \end{tabularx} \\ \\ \relax
9:4 & \begin{tabularx}{0.7\textwidth}{X} 趁著白日，我們必須做差我來的那位的工；黑夜來到，就沒有人能做工了。 \end{tabularx} \\ \\ \relax
9:5 & \begin{tabularx}{0.7\textwidth}{X} 我在世上的時候，是世上的光。」 \end{tabularx} \\ \\ \relax
9:6 & \begin{tabularx}{0.7\textwidth}{X} 耶穌說了這些話，就吐唾沫在地上，用唾沫和了泥抹在盲人的眼睛上， \end{tabularx} \\ \\ \relax
9:7 & \begin{tabularx}{0.7\textwidth}{X} 對他說：「你到西羅亞池子裡去洗。」（西羅亞翻出來就是「奉差遣」。）於是他去，洗了，回來就看見了。 \end{tabularx} \\ \\ \relax
9:8 & \begin{tabularx}{0.7\textwidth}{X} 他的鄰舍和素常見他討飯的人，就說：「這不是那從前坐著討飯的人嗎？」 \end{tabularx} \\ \\ \relax
9:9 & \begin{tabularx}{0.7\textwidth}{X} 有的說：「是他」；又有的說：「不是，卻是像他。」他自己說：「是我。」 \end{tabularx} \\ \\ \relax
9:10 & \begin{tabularx}{0.7\textwidth}{X} 於是他們對他說：「你的眼睛是怎麼開的呢？」 \end{tabularx} \\ \\ \relax
9:11 & \begin{tabularx}{0.7\textwidth}{X} 那人回答：「有一個名叫耶穌的，他和了泥抹我的眼睛，對我說：『你到西羅亞池子去洗。』我去一洗，就看見了。」 \end{tabularx} \\ \\ \relax
9:12 & \begin{tabularx}{0.7\textwidth}{X} 他們對他說：「那個人在哪裡？」他說：「我不知道。」 \end{tabularx} \\ \\ \relax
9:13 & \begin{tabularx}{0.7\textwidth}{X} 他們把以前失明的那個人帶到法利賽人那裡。 \end{tabularx} \\ \\ \relax
9:14 & \begin{tabularx}{0.7\textwidth}{X} 耶穌和泥開他眼睛的那一天是安息日。 \end{tabularx} \\ \\ \relax
9:15 & \begin{tabularx}{0.7\textwidth}{X} 法利賽人又問他是怎麼得看見的。他對他們說：「他把泥抹在我的眼睛上，我一洗，就看見了。」 \end{tabularx} \\ \\ \relax
9:16 & \begin{tabularx}{0.7\textwidth}{X} 於是法利賽人中有的說：「這個人不是從神來的，因為他不守安息日。」另有的說：「一個罪人怎能行這樣的神蹟呢？」他們之間就產生了分裂。 \end{tabularx} \\ \\ \relax
9:17 & \begin{tabularx}{0.7\textwidth}{X} 於是他們又對那盲人說：「他開了你的眼睛，你說他是怎樣的人呢？」他說：「他是個先知。」 \end{tabularx} \\ \\ \relax
9:18 & \begin{tabularx}{0.7\textwidth}{X} 猶太人不信他以前是失明，後來能看見的，等到叫了他的父母來， \end{tabularx} \\ \\ \relax
9:19 & \begin{tabularx}{0.7\textwidth}{X} 問他們說：「這是你們的兒子嗎？你們說他生來是失明的，現在怎麼看見了呢？」 \end{tabularx} \\ \\ \relax
9:20 & \begin{tabularx}{0.7\textwidth}{X} 他的父母就回答說：「他是我們的兒子，生來就失明，這是我們知道的。 \end{tabularx} \\ \\ \relax
9:21 & \begin{tabularx}{0.7\textwidth}{X} 至於他現在怎麼能看見，我們卻不知道；是誰開了他的眼睛，我們也不知道。他已經是成人，你們問他吧，他自己會說。」 \end{tabularx} \\ \\ \relax
9:22 & \begin{tabularx}{0.7\textwidth}{X} 他父母說這些話，是怕猶太人，因為猶太人已經商定，若有宣認耶穌是基督的，要把他趕出會堂。 \end{tabularx} \\ \\ \relax
9:23 & \begin{tabularx}{0.7\textwidth}{X} 因此他父母說「他已經是成人，你們問他吧」。 \end{tabularx} \\ \\ \relax
9:24 & \begin{tabularx}{0.7\textwidth}{X} 於是法利賽人第二次叫了那以前失明的人來，對他說：「你要將榮耀歸給神，我們知道這人是個罪人。」 \end{tabularx} \\ \\ \relax
9:25 & \begin{tabularx}{0.7\textwidth}{X} 那人就回答：「他是不是個罪人，我不知道；有一件事我知道，我本來是失明的，現在我看見了。」 \end{tabularx} \\ \\ \relax
9:26 & \begin{tabularx}{0.7\textwidth}{X} 他們就問他：「他給你做了甚麼？是怎麼開了你的眼睛？」 \end{tabularx} \\ \\ \relax
9:27 & \begin{tabularx}{0.7\textwidth}{X} 他回答他們：「我已經告訴你們了，你們不聽，為甚麼又要聽呢？難道你們也要作他的門徒嗎？」 \end{tabularx} \\ \\ \relax
9:28 & \begin{tabularx}{0.7\textwidth}{X} 他們就罵他：「你是他的門徒，而我們是摩西的門徒。 \end{tabularx} \\ \\ \relax
9:29 & \begin{tabularx}{0.7\textwidth}{X} 神對摩西說話是我們知道的，可是這個人，我們不知道他從哪裡來。」 \end{tabularx} \\ \\ \relax
9:30 & \begin{tabularx}{0.7\textwidth}{X} 那人回答，對他們說：「他開了我的眼睛，你們竟不知道他從哪裡來，這真是奇怪！ \end{tabularx} \\ \\ \relax
9:31 & \begin{tabularx}{0.7\textwidth}{X} 我們知道神不聽罪人，惟有敬奉神、遵行他旨意的，神才聽他。 \end{tabularx} \\ \\ \relax
9:32 & \begin{tabularx}{0.7\textwidth}{X} 從創世以來，未曾聽見有人開了生來就失明的人的眼睛。 \end{tabularx} \\ \\ \relax
9:33 & \begin{tabularx}{0.7\textwidth}{X} 這人若不是從神來的，甚麼也不能做。」 \end{tabularx} \\ \\ \relax
9:34 & \begin{tabularx}{0.7\textwidth}{X} 他們回答他說：「你完全是生在罪中的，還要來教訓我們嗎？」於是他們把他趕出去了。 \end{tabularx} \\ \\ \relax
9:35 & \begin{tabularx}{0.7\textwidth}{X} 耶穌聽說他們把他趕出去，就找到他，說：「你信人子嗎？」 \end{tabularx} \\ \\ \relax
9:36 & \begin{tabularx}{0.7\textwidth}{X} 那人回答說：「主啊，人子是誰？告訴我，好讓我信他。」 \end{tabularx} \\ \\ \relax
9:37 & \begin{tabularx}{0.7\textwidth}{X} 耶穌對他說：「你已經看見他，現在和你說話的就是他。」 \end{tabularx} \\ \\ \relax
9:38 & \begin{tabularx}{0.7\textwidth}{X} 他說：「主啊，我信！」他就拜耶穌。 \end{tabularx} \\ \\ \relax
9:39 & \begin{tabularx}{0.7\textwidth}{X} 耶穌說：「我為審判到這世上來，使不能看見的看見，能看見的反而失明。」 \end{tabularx} \\ \\ \relax
9:40 & \begin{tabularx}{0.7\textwidth}{X} 同他在那裡的法利賽人聽見這些話，就對他說：「難道我們也失明了嗎？」 \end{tabularx} \\ \\ \relax
9:41 & \begin{tabularx}{0.7\textwidth}{X} 耶穌對他們說：「你們若是失明的，就沒有罪了；但現在你們說『我們能看見』，你們的罪還在。」 \end{tabularx} \\ \\
[1ex]
\hline
\hline
\end{longtable}
$^{1}$各位弟兄姊妹平安.
這個核眼金德漢健的講題.
是參考這個奇異恩典Amazing Grace.
這首詩歌那句歌詞.
Was blind but now I see.
是走音的.
這句歌詞的靈感.
也都很可能是來自這段經文裡面.
第九章第25節那句說話.
那個核子這樣講.
有一件事我知道.
從前我是核眼的.
今天看得見了.
在這張聖經裡面.
表面上好像只是在講一個盲人.
是得蒙醫治得看見的故事.
但是這張的聖經.
是有一個很深刻的諷刺.
發覺眼睛是生理上面功能正常的人.
反而是看不到上帝的作為.
以下下來我們會用四個部份.
去思考這段經文.
我們首先看看這個聖經裡面.
各種人物的核眼.
對照我們當今的處境.
然後我們看聖經裡面的世界的恩典.
怎樣去回應最後在我們這個世界裡面.
恩典怎樣臨到我們身上.
首先我們看到在這段經文裡面.
有幾種人 有幾堆人.
他們都好像看不見一樣.
首先就是門徒的核眼.
耶穌在路上看見一個生來核眼的人.
那個人其實沒有說過話.
但門徒第一個反應並不是憐憫.
反而是問拉比他犯了罪.
究竟是他犯罪還是他父母犯罪.
使他生來核眼.
這個問題的方向是向後望的.
向以往去望的.

$^{41}$去追溯原因 尋求解釋.
可能這個也是一個心理上的需要.
見到這麼大的一個問題.
一個的苦難.
我們就很想問他是什麼原因.
而最開心的就是知道那些原因不關我們事.
而我們是沒有這個問題的.
耶穌的回應 他的方向是向前望的.
他並沒有接受那些門徒提問的框架.
所以根本就沒有回答他是什麼原因.
反而是立即轉向行動.
他不是去解釋那個問題.
而是去解決那個問題.
他說「我們必須做那差我來的人的工」.
耶穌在這裡不是去解釋黑暗.
乃是成為光 去照在黑暗當中.
值得留意的是第九章第三節.
希臘文文法上應該是視作結果句.
我要留意的是這句不是目的句.
如果字面有時候可以這樣理解.
「那人生來黑暗」.
好像上帝整定他是黑暗.
等到這個時刻.
讓耶穌去展現他的神聖.
然後得到醫治.
這樣是錯誤的理解.
上帝沒有理由這麼殘忍.
要那個人生出來.
我們看到經文是「去到成人」.
即是十幾年的時間.
然後等到今天.
終於讓耶穌去彰顯他的大能.
而是那句說話是「以至於」.
即是說這個不是上帝刻意降臨苦難.
而是上帝在任何的苦難之中.
都可以展現恩典.
而今次在這一刻.
在這個事情裡面.
上帝藉著耶穌基督.
就展現他的醫治的恩典.

$^{81}$「門徒的黑暗」不是他的眼睛有什麼問題.
而是思想方向出了問題.
他們是向後望.
耶穌是向前望.
而那個生來黑暗的人的鄉里.
他們的鄰舍明明是看到東西的.
不過只是看到殘缺.
因為那個人得到醫治的時候.
他的鄉里反而在吵起來.
那個不是平時在吃飯的嗎?.
有人說「這個是的」.
有人說「不是的」.
直到那個人說「我就是了」.
為什麼認不出那個人不是逆了庸呢?.
很可能是因為他們以往看到的.
這個人其實不是一個完整的人.
這個人只不過是一個生來黑暗的人.
他的殘疾成為他唯一的一個身份.
當他已經不再是殘疾的時候.
這個標籤沒有了的時候.
他根本就看不到有一個人在這裡.
而那個生來黑暗的人的父母.
更加很奇怪地.
有一些和那些法理錯誤的人質問他的人的回答.
那個父母是不敢去講.
究竟他的兒子為什麼會有這樣的狀況.
他只是回答說「是的,他生出來就是這樣」.
他不敢去講那個我們說.
英文講Elephant in the room.
即是房間裡的大象.
在這個審訊裡面沒有人敢講.
被人盤問的時候沒有人敢講.
事實上是耶穌開了他的兒子的眼睛.
因為他們很擔心會被人趕出會堂.
因為這樣等於承認耶穌就是基督.
父母被傳召去問的問題的時候.
你看看他怎樣回答.
他們說「至於他怎樣能夠看得見.
我們就不知道是誰開了他的眼睛.
我們都不知道.

$^{121}$你們問他吧,已經成年了」.
這一招很厲害.
一篩就篩給他的兒子.
他是不是真的不知道呢?.
肯定不是.
他回到家裡的.
他不是到處流浪的.
回到家裡的時候一定會問這個問題.
怎麼可能不知道.
但為什麼他們說不敢講.
或者說他們不知道呢?.
是因為可能是時勢真惡.
因為那種恐懼令到他們假裝看不見.
而約翰福音裡面講得很清楚.
就是他們這樣講是因為怕猶太人.
將他們逐出會堂.
這件事對當時來說是非常之嚴厲的事情.
不是說我不回會堂.
就好像我不回教會.
不是這麼簡單.
基本上是切斷他和他的民族.
和所有的人的關係.
而這個關係不是純粹那種感情的關係.
還包括經濟,社會等等.
種種的東西都切斷.
是很慘烈的事情.
所以他很害怕.
他們很害怕的時候.
害怕就不會講真話了.
你問他吧.
他們不是看不見.
不過他們是選擇看不見.
因為他要承認看到什麼的時候.
那個代價太大.
而法理賽人的眼睛.
生理上是健康的.
但他們還是看不見實際的情況.
當那件事他知道.
耶穌在安息日那裡去醫病的時候.
而且不單止是醫病.

$^{161}$只是將泥土弄一下.
這個很明顯在當時的律法.
當時的傳統裡面是不對的.
是違反了律法的.
所以對於他們來說.
這個是大逆不道的事情.
於是他們就在那裡說.
我們知道這個是罪人.
我們知道上帝跟摩西說話.
然後說我又不知道祂在哪裡來.
其實他說那種不知道.
不是真的不知道.
而是不想去知道.
就是在那裡迴避事實.
而最終耶穌在這段經文最後一句.
就說你們說我們.
如果你們真的說我們看得見的話.
你們的罪就存留.
他們那幫法利賽人自以為看得見.
其實他們看不見.
或者不願意去看到.
今天我們的世界.
其實可能都有類似這樣看不見的事情.
雖然生理上我們看得見.
但實際上我們在表現上.
我們好像看不見一樣.
或者我們都會好像那幫門徒一樣.
我們會向後去追究原因.
但忽略了前面的行動.
我們見到受苦的情況.
可能第一個衝動就是問.
為什麼會這樣.
誰搞出來的.
是不是你做錯事.
是不是誰誰做錯事.
而不是問我自己.
我們可以做些什麼.
以致說我們對於那個苦難.
事實上我們看不見.
我們只可能看到苦難.

$^{201}$但看不見承受苦難的人.
另一方面我們可能好像那些鄉里一樣.
我們看見的是一些標籤.
但看不見那個人.
那些鄉里只是看見那個盲眼的人.
那個生理盲眼的人.
但不是真的去認識那個人.
今天或者我們在這個社會裡面.
我們都有一些標籤.
可能那些叫做露宿者,無家者.
那些叫做精神病患者.
那些可能是讀書讀不成的人.
那些叫做邊緣青少年等等.
當那些標籤比起那些人具體的.
活有血有肉的生命更加真實的時候.
我們就好像那幫鄉里一樣.
根本就認不出他.
他有什麼轉變我們都看不見.
而我們也都很可能會好像.
那個盲了眼.
就是生理黑眼的人一樣.
父母一樣.
不願意承認我們房間裡面的大象.
無論在教會,社會.
可能我們都會有一些這樣的大象出現.
有一些情況我們是知道的.
但是沒有人夠膽開口去講.
可能是有人被不公平的對待的.
可能是有一些做法明顯是有問題的.
可能有一段的關係早已經是破裂了.
可能有一些東西.
是已經去到一個很危險的境地.
但是我們可能都會選擇沉默.
因為時勢真惡.
因為代價太大.
在說出真相.
還有我們是不想搞破壞派對.
我們想繼續去維持著一個.
好像很和平的關係之間.
可能我們真的選擇了後者.

$^{241}$我們去嘗試去保護著我們原本有的東西.
我們不願意在一個壓力裡面去說出真話.
又或者我們都會很緊張一些的原則.
那些原則去到一個地步.
我們將那些原則凌駕在那個人的尊嚴.
凌駕在那些人的需要之上.
有時候我們都會.
可能會用一個原則去行事.
用管理的方式行事.
用我們跟隨傳統穩穩陣陣的方式行事.
以致說有個人的需要的時候.
我們沒有很好的理會.
在聖經裡面.
我們看到上帝並不是這樣的.
在約翰福音其實第一章.
就是約翰福音的序言裡面.
都有說光照在黑暗裡.
雖然黑暗卻不接受光.
在耶穌的身上.
祂是將光照亮這個世界.
祂當門徒去問.
究竟為什麼那個人這麼慘的時間.
耶穌並沒有提供任何因果的分析.
並沒有做一個神學的分析.
而是祂立即去行動.
祂彎下身在地上拿著泥土.
然後吐口水下去.
弄一弄祂.
濕濕的.
然後擦在那個瞎了眼的眼睛上.
然後叫祂去洗.
恩典不是拿來去解釋事情.
不是去解釋黑暗.
而是進入黑暗.
光照在黑暗裡.
那位在創世記裡面.
說有光的上帝.
再次在一個混沌面前.
他選擇創造.
苦難不需要拿來去解釋.

$^{281}$更加不是去解釋.
這個就是什麼做錯事的一種懲罰.
而是這個本身就需要去回應.
需要去讓上帝恩典展現出來的一個機會.
然後我們也看到.
耶穌是怎樣首先去回應那個人的需要.
祂真的看到那個人.
我們注意在第一節那裡已經說.
耶穌過去的時候看見一個生來黑暗的人.
在這件神蹟的事件當中.
並不是那個人在大聲喊叫.
大衛的子孫可憐我爸那種.
而是那個人沒有出聲.
而是耶穌停下腳步.
首先看到那個人.
可能祂的香梨,祂的鄰舍.
都看不到那個盲人的情況.
但是耶穌在路過的時候看見他.
看見他不是一個問題.
不是一個標籤.
而是一個有血有肉的具體的人.
這個是恩典.
是一個不計較任何東西.
不去很多的計算.
而是先行的,主動的.
不去看他究竟配還是不配這個問題.
然後耶穌在這裡.
你可以說祂用最平凡,最卑微的材料.
就是祂的口水,祂的泥土去促成這個醫治.
早期的教父愛蔭柳.
曾經這樣說過.
這一件這樣的動作.
其實是在呼應《創世記》第二章.
上帝用塵土那裡做人.
所以令人思想這件事情.
不單止是一個醫治.
更加是一個新的創造.
那位講要有光的創造主.
再一次彎低身.
用雙手去創造一件奇妙的事.

$^{321}$使得那位生來核眼的看得見.
而約翰講明給我們聽.
耶穌叫他去洗那個池.
那個名西羅亞就是鳳釵蜆的意思.
鳳釵蜆沒什麼特別.
特別的就是因為耶穌就是鳳釵蜆那位.
所以物理上.
祂是去了一個池叫西羅亞洗.
但是真正得到醫治的.
就是那位猜祂去洗的.
那位被天父所猜蜆來到世界的.
那位耶穌基督.
耶穌首先聽見然後看到.
而耶穌首先看到.
而那位盲人是首先聽見因為他看不到.
但是他聽見聽到耶穌要他去做的事情的時候.
他就很順服地過去了.
這件事有些人說其實是呼應著約翰福音.
接著那十張羊認識好牧人的聲音.
然後其實在耶穌叫他去西羅亞洗之後.
其實發生很多事情.
他被質問,他的父母被質問等等很多很多事情.
其實很有趣的.
在那個文學的結構裡面.
那個主角如果說是耶穌的話.
那個主角是不見了的.
突然間不出現.
直到什麼時候出現呢?.
就是另外一個主角.
就是那個生來黑眼的人.
他被法利賽人趕出去了.
然後在第九章35節.
耶穌聽說他們把他趕出去就找到他.
當他被那個群體.
被那個主流的社會趕出去的時候.
耶穌就主動地去尋找他.
耶穌不是只是你走過我就依你.
耶穌主動去關心他.
去尋找他.
然後法利賽人就曾經命令那個核子說.

$^{361}$歸榮要給上帝.
其實這句話在古代當時的世界裡面.
其實是逼人去認罪的一個誓言.
而當那個人向耶穌說.
主啊我信和敬拜的時間.
那一刻其實才是真正的歸榮要給上帝.
今天我們在這個世界當中的時間.
耶穌在哪裡?.
我們知道耶穌是在這個世界當中.
他仍然在這個世界裡面.
看見那些被這個社會標籤化.
被這個社會邊緣化的人.
被驅逐的人.
他仍然會問.
他仍然會看得見.
他仍然會主動的去幫助他們.
而恩典往往是在一些最平凡的事情裡面出現.
就好像吐口水.
就好像那些泥土一樣那麼平凡.
恩典未必一定是那些.
可以放到煙花.
可以擺到上IG.
很多人看那些視頻.
那個事情裡面出現.
反而可能是最微小的一些事情裡面.
可能是我們對人的溫柔的問候.
可能是我們願意好好的聆聽.
可能我們是在別人遭遇患難的時候.
我們稍微去伸出援手.
但無論怎樣.
小小的事情.
上帝可以施行祂的拯救.
可以施行祂的醫治.
可以彰顯祂的恩典.
無論那是在教會之內還是教會之外.
而耶穌說.
「趁著白日,我們必須做那差我來者的工」.
耶穌沒有說我.
只是我要做差我來者的工.
我們都要做.

$^{401}$那個「我們」在當時的語境裡面.
就是包含了那些門徒.
即是說耶穌把門徒都納入在祂來到這個世界的使命當中.
耶穌所做的事情.
正正是祂的門徒要做的事情.
事實上在《約翰福音》第20章裡面.
耶穌去向那些很害怕猶太人.
最迫害他們的人.
就是耶穌復活了那天的晚上.
走入已經關上門的那群門徒聚集的地方.
跟他們說「願你們平安」.
然後說「父怎樣差遣了我,我也照樣差遣你們」.
耶穌就是那位被天父差遣進入世界的人.
而復活了的主.
也照樣地去差遣祂的門徒進入世界.
如果說耶穌在這個世界裡面.
就是去宣講和行動的時候.
我們也都應該在這個世界裡面.
既宣講也有行動地.
被上帝所差遣進入這個世界.
與耶穌基督一起去同工.
基督徒的群體是不是願意.
或者明白我們其實是跟耶穌基督一起去同工.
而我們願不願意好像耶穌基督一樣.
願意成為那位被社會趕出去.
被主流的力量趕出去的人的淪陷呢?.
而我們看到那個生危核眼的人.
被人一直盤問很厲害的一些問題.
而他的答覆是很有力的.
因為他講得很清楚.
人家問那個人肯定是罪人.
你知不知道什麼的時候.
他就說那個人是不是罪人我不知道.
有一件事我知道.
從前我是核眼的.
如今能夠看見.
這個是非常真實的情況.
沒有東西可以推翻這件事情.
因為他自己的親身經歷.
剛才講那個Amazing Grace的作者.

$^{441}$叫John Newton.
他晚年記憶力衰退.
有人問他你還記得些什麼?.
他就說我其實很多事情都不記得了.
但有兩件事我永遠記得的.
第一我是一個大罪人.
事實上他以前是去販賣奴隸的.
後來當然是悔改.
所以他很明白.
他是前我失喪今被尋回.
然後第二件事很記得的就是.
基督是一位偉大的救主.
這些不是讀神學那些讀出來.
不是他背誦聖經的經文.
而是他親身遭遇.
經歷上帝恩典.
經歷自己的改變而來的事情.
今天我們在一個大齋期.
或者叫做御苦期.
或者叫做四秦期的時間.
正正是給我們很多機會去審查.
靜下來去審查究竟我們是誰.
我們在上帝面前是怎樣的人.
究竟上帝給我們的醫治和改變是什麼.
究竟上帝在我們的生命當中.
究竟是什麼位置.
這個是我們審查的機會.
羅馬地下那些墓穴的壁畫.
其中一個壁畫是.
將生來核眼的人的遭遇.
成為了一個受洗者的象徵.
為什麼會和洗禮有關.
弄點水抹一下他的眼睛.
又或者根本他就浸在西萊瓦池那裡.
所以點水禮和浸禮都包括在內.
這個我自己想的.
但是他真的是有水.
讓他能夠得到醫治.
得到整個人的改變.
而他的旅程很可能.

$^{481}$是我們每一個跟從耶穌基督的人的旅程.
我們會不會是可以在這個經文裡面.
去看到我們自己.
究竟我們是這位仁兄.
是經歷這麼大的改變.
還是我們繼續成為那些門徒一樣.
不停地問究竟哪個問題.
為什麼是這樣.
是不是他不好.
是不是他父母不好的問題.
或者我們是不是好像那些鄉里鄰舍一樣.
只是看到那個標籤.
那個標誌就是這個人就是生來核眼.
沒有其他 沒有明字.
沒有一個具體的人.
又或者好像他的父母一樣.
我們不夠膽去講一些很明顯.
其實是人都知道的東西.
不過我們都不夠膽去承認.
又或者我們好像法利賽人一樣.
我們只是在想那些原則 規則等等的東西.
我們忘記了究竟我們最需要的是什麼.
又或者我們只是記得那些的律法.
那些的規條.
包括那些安息日的要求.
但我們不記得安息日是為人而設.
而人子是安息日的主.
在這個大齋期的期間.
願我們都承認自己.
其實或多或少.
在不同的情況底下.
我們都是核眼的.
就讓世界的光.
耶穌基督.
祂去塗抹我們的眼睛.
讓我們也都順服的去洗.
洗去我們自己的種種的罪惡.
然後帶著一個經歷上帝恩典的見證.
我們跟其他人講.
好像那位的核子一樣.

$^{521}$有一件事我知道.
從前我是核眼的.
如今能夠看見.
阿們.
\newpage



\section{列王記下 4:8-37}
\label{sec:031807}
\textbf{信心系列:信心奇女子}~~~~曹子光牧師
\newline
\newline
連結: \href{http://www.church.com.hk/acms/content.asp?site=cdc\&op=show\&type=product\&code=031807\&layout=sermon}{\texttt{031807}} ~~~~ 語音日期: 2026-03-08
\newline
\newline
\hyperref[sec:031806]{\small{< < < PREV SERMON < < <}}
~
\hyperref[sec:index_chronic]{\small{[返順時目]}}
~
\hyperref[sec:031808]{\small{> > > NEXT SERMON > > >}}
\newline
\newline
列王記下 4:8-37
\newline
\begin{longtable}{cl}
\hline
\hline
章節 & 經文 (和合本修訂版)\\
\hline
4:8 & \begin{tabularx}{0.7\textwidth}{X} 一日，以利沙經過書念，在那裡有一個富有的婦人強留他吃飯。此後，以利沙每次經過就轉到那裡去吃飯。 \end{tabularx} \\ \\ \relax
4:9 & \begin{tabularx}{0.7\textwidth}{X} 婦人對丈夫說：「看哪，我知道那常從我們這裡經過的是神聖的神人。 \end{tabularx} \\ \\ \relax
4:10 & \begin{tabularx}{0.7\textwidth}{X} 我們可以為他蓋一間有牆的小閣樓，裡面安放床榻、桌子、椅子、燈臺。每當他來到我們這裡，就可以住在那裡。」 \end{tabularx} \\ \\ \relax
4:11 & \begin{tabularx}{0.7\textwidth}{X} 一日，以利沙來到那裡，轉進那閣樓，躺臥在那裡。 \end{tabularx} \\ \\ \relax
4:12 & \begin{tabularx}{0.7\textwidth}{X} 以利沙吩咐僕人基哈西說：「你叫這書念婦人來。」他把婦人叫了來，婦人就站在以利沙面前。 \end{tabularx} \\ \\ \relax
4:13 & \begin{tabularx}{0.7\textwidth}{X} 以利沙吩咐僕人說：「你對她說：『看哪，你為我們費了許多心思，我可以為你做甚麼呢？我可以為你向王或元帥求甚麼呢？』」她說：「我已住在自己百姓之中。」 \end{tabularx} \\ \\ \relax
4:14 & \begin{tabularx}{0.7\textwidth}{X} 以利沙說：「究竟可以為她做甚麼呢？」基哈西說：「她真的沒有兒子，她丈夫也老了。」 \end{tabularx} \\ \\ \relax
4:15 & \begin{tabularx}{0.7\textwidth}{X} 以利沙說：「叫她回來。」於是他叫了她來，她就站在門口。 \end{tabularx} \\ \\ \relax
4:16 & \begin{tabularx}{0.7\textwidth}{X} 以利沙說：「明年這時候，你必抱一個兒子。」她說：「神人，我主啊，不要這樣欺哄婢女。」 \end{tabularx} \\ \\ \relax
4:17 & \begin{tabularx}{0.7\textwidth}{X} 婦人果然懷孕，到了明年那時候，生了一個兒子，正如以利沙向她所說的。 \end{tabularx} \\ \\ \relax
4:18 & \begin{tabularx}{0.7\textwidth}{X} 孩子長大，一日出去到他父親和收割的人那裡。 \end{tabularx} \\ \\ \relax
4:19 & \begin{tabularx}{0.7\textwidth}{X} 他對父親說：「我的頭啊，我的頭啊！」他父親對僕人說：「把他抱到他母親那裡。」 \end{tabularx} \\ \\ \relax
4:20 & \begin{tabularx}{0.7\textwidth}{X} 僕人抱去，交給他母親。孩子坐在母親的膝上，到中午就死了。 \end{tabularx} \\ \\ \relax
4:21 & \begin{tabularx}{0.7\textwidth}{X} 他母親上去，把他放在神人的床上，關了門出來， \end{tabularx} \\ \\ \relax
4:22 & \begin{tabularx}{0.7\textwidth}{X} 呼叫她丈夫說：「你叫一個僕人給我牽一匹驢來，我要趕去見神人，然後回來。」 \end{tabularx} \\ \\ \relax
4:23 & \begin{tabularx}{0.7\textwidth}{X} 丈夫說：「今日不是初一，也不是安息日，你為何要到他那裡去呢？」婦人說：「平安無事。」 \end{tabularx} \\ \\ \relax
4:24 & \begin{tabularx}{0.7\textwidth}{X} 於是她備上驢，對僕人說：「走，趕緊走，除非我吩咐你，不要為了我而慢下來。」 \end{tabularx} \\ \\ \relax
4:25 & \begin{tabularx}{0.7\textwidth}{X} 婦人往迦密山去，到了神人那裡。神人遠遠看見她，對僕人基哈西說：「看哪，書念的婦人來了！ \end{tabularx} \\ \\ \relax
4:26 & \begin{tabularx}{0.7\textwidth}{X} 現在你跑去迎接她，對她說，你平安嗎？你丈夫平安嗎？孩子平安嗎？」她說：「平安。」 \end{tabularx} \\ \\ \relax
4:27 & \begin{tabularx}{0.7\textwidth}{X} 婦人上了山，到神人那裡，就抱住神人的腳。基哈西前來要推開她，神人說：「由她吧！因為她心裡愁苦。但耶和華向我隱瞞這事，沒有告訴我。」 \end{tabularx} \\ \\ \relax
4:28 & \begin{tabularx}{0.7\textwidth}{X} 婦人說：「我何嘗向我主求過兒子呢？我豈不是說過，不要欺哄我嗎？」 \end{tabularx} \\ \\ \relax
4:29 & \begin{tabularx}{0.7\textwidth}{X} 以利沙吩咐基哈西說：「你束上腰，手拿我的杖前去。若遇見人，不要向他問安，人若向你問安，也不要回答。要把我的杖放在孩子臉上。」 \end{tabularx} \\ \\ \relax
4:30 & \begin{tabularx}{0.7\textwidth}{X} 孩子的母親說：「我指著永生的耶和華，又指著你的性命起誓，我必不離開你。」於是以利沙起身，隨著她去了。 \end{tabularx} \\ \\ \relax
4:31 & \begin{tabularx}{0.7\textwidth}{X} 基哈西在他們以先去了，把杖放在孩子臉上，卻沒有聲音，也沒有動靜。基哈西回去，迎見以利沙，告訴他說：「孩子還沒有醒過來。」 \end{tabularx} \\ \\ \relax
4:32 & \begin{tabularx}{0.7\textwidth}{X} 以利沙進了屋子，看哪，孩子死了，放在自己的床上。 \end{tabularx} \\ \\ \relax
4:33 & \begin{tabularx}{0.7\textwidth}{X} 他進去，關上門，只有他們兩個人，他就向耶和華祈禱。 \end{tabularx} \\ \\ \relax
4:34 & \begin{tabularx}{0.7\textwidth}{X} 他上去伏在孩子身上，口對口，眼對眼，手對手。他伏在孩子身上，孩子的身體就漸漸暖和了。 \end{tabularx} \\ \\ \relax
4:35 & \begin{tabularx}{0.7\textwidth}{X} 然後他下來，在屋裡來回走了一趟，又上去伏在孩子身上。孩子打了七個噴嚏，眼睛就睜開了。 \end{tabularx} \\ \\ \relax
4:36 & \begin{tabularx}{0.7\textwidth}{X} 以利沙叫基哈西說：「你叫這書念婦人來。」於是他叫了她來。婦人來到以利沙那裡，以利沙說：「把你兒子抱起來。」 \end{tabularx} \\ \\ \relax
4:37 & \begin{tabularx}{0.7\textwidth}{X} 婦人就進來，在以利沙腳前俯伏於地，向他下拜，然後抱起她兒子出去了。 \end{tabularx} \\ \\
[1ex]
\hline
\hline
\end{longtable}
$^{1}$我們將會繼續看信心的系列.
看到信心.
原來信心是一個可以帶領我們穿越不同時間發生.
在我們身上的屬靈旅程.
今天想和大家分享一個和信心有關的故事.
是記載在《舊約聖經》的列王記下.
第四章第八到第三十七節.
今天想和大家分享三件很簡單的事.
第一件事就是信心的起點.
作為信心的起點其實是為神預留一個位置.
列王記下第四章一開始.
記載了一個貧窮寡婦的故事.
她的丈夫原本是先知的門徒.
即是以利沙先知的門徒.
但是丈夫死了之後就遺留了一大筆的債項.
債主臨門要拿他兩個兒子去當奴隸.
這位婦人就去求先知以利沙幫助她.
結果以利沙就叫她將家中剩餘的一些油.
不停地倒進拿回來的器皿裡.
原來神蹟接著就發生.
油一直在倒 一直在滿.
最後裝滿了所有器皿.
她賣了油 還清了債務.
還應該可以靠一些來養生.
某個程度上 聖經形容這一個.
是好像將她兩個兒子活過來的狀況一樣.
本來要被帶走成為奴隸.
人生應該就好像玩完了一樣.
但是神介入了之後 生命重新開始.
但很有意思的是.
聖經接著記載了另一個更長 更細緻的故事.
而這次不再是一個貧窮的寡婦.
而是一個似乎頗富有的家庭.
不再是兩個差點被帶走的兒子.
而是一個後來真是死去又復活的兒子.
現在這個故事是一個富有的婦人.
一個死而復活的孩子.
換句話說 第一個故事就好像一個序幕一樣.
窮寡婦的故事原來就是接下來這個故事的序幕.
準備我們進入一個更震撼的神蹟.

$^{41}$這次的主角是一位書唸的婦人.
聖經說她住在書唸這個地方.
而且她似乎是一個頗有地位 頗富裕的女人.
當神人以利沙經過這裡的時候.
她就極力挽留她入屋吃飯.
之後每逢以利沙經過都會去她家吃飯.
但事情不是那麼簡單 只是停留在吃飯.
這位婦人和丈夫商量了之後.
決定為以利沙建一個小的房間.
擺上一張床 一張椅 一張書桌 一張硬椅.
讓她每一次經過的時候都可以好好地休息下來.
弟兄姊妹 這件事其實真的一點都不簡單.
婦人不是請一餐飯那麼簡單.
而是為神預備了一個長期的空間.
而為神人預備了一個長期的空間.
聖經形容了這個空間是一個小樓.
這位婦人 這位書唸的婦人所做的.
不單只是金錢的付出 更加是一份心意.
是對神人 對上帝的一份心意.
如果仔細地讀這篇經文 你就會發現一件事.
這位婦人其實沒有什麼特別追求屬靈的事情.
她不是因為先知住在她的家.
所以她就會問很多關於屬靈問題的事的人.
她也不是因為追求屬靈到一個地步.
以利沙每次吃飯的時候就會問她聖經.
因為這些全部都在聖經裡面沒有記載的.
她甚至大概都沒有為自己去祈禱.
叫子妮莎去為她沒有兒子去祈禱.
因為經文其實提醒我們.
她老公已經年紀很大 而且她一直都沒有兒子.
換句話說 如果她真的想求的話.
其實她有很多事情可以求的.
不過 她沒有.
從這個角度去看 這位婦人的信仰.
其實未必真的想得那麼熱切.
她不是真的很主動.
她不是一個拼命追求屬靈事情的那種人.
她只是在這個故事的一開始.
為神人做了一件很簡單的事.
她為神人預備了一個空間.

$^{81}$但她從來都沒有想過.
她為先知所預備的這個小房間.
日後竟然會成為她兒子的急救室.
A and E的房間.
當她兒子死了之後.
聖經說她第一時間抱了她兒子上去神人的房間.
放了這個兒子上去神人的床.
更奇妙的是 之後以莉莎去到那個地方的時候.
關上門 為孩子禱告 令他復活.
當初這個婦人不是為神人預備這個房間.
關上門 是為讓神人休息的嗎?.
去到後期 神竟然透過同一個房間.
同一道門 行出一個死人復活的神蹟.
弟兄姊妹 有時候你為神預備的空間.
你自己都沒有想過將來神會怎樣用.
當你為神預留空間的時候.
神就會在這個空間裡面為你去預備.
你從來沒有想過的事情發生.
我們SOL的教會.
我們原本的理念一直都是.
除了自己的教會之外 也有一些網絡的東西.
音樂事公其實是一個很重要的事公.
是我們一直都由第一天開始就想做.
甚至已經開始了做.
所以我們有這一份心意.
是由第一天開始就已經有.
很感恩我們能夠和南部不同的教會.
我們在開始的時候有一些很好的聯繫.
所以我們就透過這個小小的空間.
一個音樂的 看看有沒有可能合作的.
這樣的一個心意.
於是我們就大家在不同的教會裡面聊.
結果我們沒多久之後.
我們組織了一個微型音樂會.
我相信是一次 當年是一次.
我想都差不多七,八年前了 至少.
是一次華人教會史上.
挺難得的 在英國的一個音樂的小聚會.
我們就將我們自己所創作的.
或者我們崇拜音樂的那一種風格.

$^{121}$大家就在這裡分享出來.
不算很多人 大概七,八十人左右那次來.
但是這是一個很特別的事.
上帝就是預備了我們這個很特別的事公.
事實上當我們為神預留空間的時候.
神就會為我們預備一些從來未必會想到會發生的事.
其實今天有很多人的信仰的開始.
都好像這個婦人一樣.
特別是很多現在從香港來英國生活的人.
他們未必一開始的時候就很有信心.
或者未必根本是很想投入信仰的裡面.
但是他們願意和基督徒一起生活.
接觸 參與教會的活動.
甚至去社區的聚會.
信仰?是的 真的未必那麼有興趣.
崇拜?坦白說.
有時間就去 沒有時間就不去.
小組?更加不用說了.
接觸無妨 但是要他認真.
經常都要參加這個小組的話 就要看看情況了.
其實很多人他們都有穩定的聚會.
但是未必馬上可以做到認真去相信.
甚至對上帝很有信心的那一種追求型的信徒.
那種屬靈的心態.
但是有時候信心的開始不一定是尋求神開始.
而是先為神留一個位置 留一個空間.
就好像這個婦人一樣.
其實以利沙都是一樣的.
以利沙因為這個婦人給出的一個好意.
所以他也回報她兩個好意.
大家如果一會兒就會聽到.
以利沙其實第一個好意就是.
因為知道她是一個頗有錢的人.
所以看看會不會幫她疏通得到.
這是第一個她想做的事情.
第二樣想做的事情就是為她 可不可以幫她有一個遺址.
所以大家的出發點都是想為對方有一份心意.
我聽過一個故事.
有一個香港來的姊妹.
她來英國之後每個星期都帶著兒女去教會.

$^{161}$但問題是她不懂得開車.
所以就要她老公來載她.
有一次她就跟她老公說.
不如你載我之餘也進去教會吧.
老公就說:我都不相信,我進去做什麼?.
老婆就說:那你不是不陪我吧?.
結果她真的進去了.
但她每一個星期都要坐在教會裡面最後一行.
而且詩歌當然不會唱.
祈禱也不會不祈.
聽到就看哪些人說.
說的訊息會不會適合自己聽.
於是她就通常每次傻伯就坐在後面滾手機.
但一年之後她竟然相信了主.
有人就問她:你何時開始相信主?.
她就說了一句很特別的話.
她說:其實我未相信之前.
我相信神已經慢慢走進我的生命裡面.
她說因為第一.
她慢慢習慣了聽道這個環節.
她以前覺得說得很悶.
但慢慢覺得也很有道理.
第二,她開始被基督徒影響到.
她發現教會的人都很關心她.
很樂意去幫人.
第三,她的心慢慢打開了.
以前她覺得信仰只是心理安慰.
但慢慢開始在想.
神是否真的存在呢?.
神是否在這附近呢?.
祂周圍放那麼多基督徒在我身邊.
有沒有甚麼特別的心意在裡面呢?.
結果每一次的崇拜.
每一次的講道.
甚至人們怎樣去關心她的一句說話.
她堆積,堆積,堆積之後.
她就開始慢慢打開了心.
結果在一次荒廢的聚會裡.
她就決志相信了神.
所以弟兄姊妹.

$^{201}$信心很多時候不是突然爆發.
信心往往是從一個很小的空間開始.
當一個人願意為神預留一個空間的時候.
神就慢慢進入了她的生命裡面.
當年我中學的時候.
我班裡面有兩位基督徒.
他們兩位都是臨班的同學.
當年我是中五的學生的時候.
他們兩位就在班裡面傳福音.
很多人都相信了.
他們最後的對象就是我.
跟著就跟我說.
可不可以給他們兩個禮拜的時間.
給他們兩個禮拜的時間跟我一起吃午餐.
我心裡想好我就豪兩個禮拜給你.
結果兩個禮拜裡面.
我就聽著他們講聖經故事.
然後一邊吃飯一邊吃午餐的時候.
在那裡唱福音詩歌給我聽.
我覺得很愉悅.
你知不知道為甚麼.
因為在街裡面唱歌.
物質質的環境裡面那個飯堂.
那個吃東西的地方.
你唱一些歌.
而這些歌又不配合音的詩歌.
令我覺得很尷尬.
但無論如何.
那兩個禮拜.
是我對上帝的一個小小的心意打開了.
一個空間打開了.
當然我不是馬上相信耶穌.
不久之後他們邀請我去一個報道會.
在報道會裡面.
我就舉手決志去相信神.
是的 信心的起點.
是為神預留一個空間在裡面.
然後信心的起點.
其實我們都好心款待神人.
不知不覺他們就為神預留了一個空間.

$^{241}$是的 弟兄姊妹.
很多人的信仰的開始.
不是先尋求神.
而是先為神預留一個位置.
第二件事就是信心的成長.
生命原來可以不停地去經歷神.
當我們再仔細去讀這段聖經的時候.
有沒有發現在整個故事裡面.
聖經大部份時間都稱呼這位女士.
做書唸的婦人.
或者婦人.
甚至人稱代名詞的她.
甚至有時叫她做婢女.
但是到第四章三十節的時候.
經文突然之間用了一個很特別的稱呼.
就叫做孩子的母親.
而且這個稱呼是整個聖經裡面唯一一次出現.
這個改變其實真的很有意思.
因為一開始的時候.
這位婦人似乎對於母孩子其實一點都不在意.
你記不記得我們說.
以利沙想為她做兩件事.
第一件事就是問她.
有什麼可以幫你在人脈裡面開通.
她的婦人都很囂張.
告訴她自己可以在百姓裡面自己搞定.
意思就是她現在樣樣都肯.
不需要別人幫.
第二件事她就說.
她可以幫她有一個小孩.
坦白說.
以利沙說可以做到這件事.
但是這個婦人其實一點都不在意.
其實她老公已經年紀很大.
她當時的文化底下.
女人其實沒有什麼可以做到.
男人可以三妻四妾.
但是女人沒有選擇.
沒有小孩.
但是她只能夠接受命運.

$^{281}$接受事實.
所以當以利沙跟她說.
明年這個時候你要抱一個兒子的時候.
她的反應不是興奮.
而是幾乎可以說是抗拒.
神人啊.
不要欺控婢女.
意思就是說.
不要給我有希望.
她可能也學會了這一套理論.
一套法則.
就是有些事情不期待的時候.
就不會失望.
但奇妙的是一年之後.
以利沙所說的竟然成真.
她真的生了一個小孩.
這個本來沒有想過會出現的小孩.
就是這樣改變了她整個人生.
而經文其實給了我們一個很細節的描述.
暗示了原來這位婦人的信仰開始改變.
當小孩突然出事的時候.
當她的小朋友突然出事的時候.
婦人就急著去見神人.
婦人就急著去見神人以利沙.
但她丈夫就問了她一句話.
他說今天不是初一.
也不是安息日.
那為什麼你要去神人那裡呢?.
這句話其實透露了一件事.
原來這個婦人是有定時定候去見神人的.
是有定時定候去見以利沙的.
換句話說.
用今天的說話來說.
她是定時定候去聚會.
她定了時間去見神人.
聽神的說話.
不過有趣的是.
當她還沒有小朋友的時候.
經文完全沒有提到這樣的習慣.
似乎是自從她有了這個小朋友之後.

$^{321}$她對神開始有興趣.
開始有動力去認識神都未定.
大概我們都不難猜.
因為神在她身上做了一件美事.
就是令她有兒子.
這樣你都應該懂得說.
所以大概可能她就是因為這樣的緣故.
去建立一個定期去見神人.
去聽上帝的說話的心態.
她的信仰開始慢慢成長.
其實這件事是很真實的.
因為人生有很多時候.
是當身份改變了之後.
信仰才會慢慢開始改變.
承認的越來越開始有追求的心態都未定.
以前她只是一位很有錢.
一個普通的婦人.
但現在她是一個母親.
有一個孩子.
而一個人一旦成為了父母.
人生有很多事情都會改變.
現實一點來說.
我不知道大家同不同意.
有了小朋友之後.
我們的脾氣是不是逐漸越來越差了.
我們的脾氣都會改變了.
當然不是每個人都是這樣.
生活習慣都改變了.
以前你可能睡到黃昭白晚.
十點起床都未定.
但有了小朋友的時候.
你可能四五點就會起床.
未來計劃你都可能有變動.
你以前以為一年可以去幾次旅行.
但因為有了小朋友的時候.
你的計劃可能就已經打亂了.
彷彿一切都因為多了一個孩子而完全改變了.
很多人都會經歷了不同家庭的事件.
發生甚至在裡面教養的方法產生衝突.
上一代和我們處理小朋友問題的不同.

$^{361}$會帶來很多張力.
這些都是有了小朋友之後.
會帶來很多家庭的事件發生.
不過當我講到成為父母和有小朋友的時候.
我知道當中未必一定每個家庭都有小朋友.
但我想你知道一件事.
神塑造我們的生命.
不一定只是要透過有小朋友.
有時候神可能要透過家庭.
透過工作.
透過困難.
逆境.
透過失去.
透過你拿不到的東西等等.
每一個人的生命故事都不一樣.
但神都會在不同的經歷裡面塑造我們.
令我們慢慢變成更成熟更像他.
在舒蓮夫人的這個故事裡面.
在她的生命裡面.
神就是透過成為母親這個身份去改變了她.
神有時候不是要先改變人的信仰.
而是要先改變人的身份.
對於這位舒蓮夫人.
當她成為母親之後.
人生完全不同了.
舒蓮夫人以前可能覺得信仰是一種客氣的款待.
但有了小朋友之後.
信仰變成了一個生命的依靠.
但對一個有信仰的人來說.
這個改變其實有另一個意思.
因為神往往透過人生的角色和責任.
來塑造我們的生命.
神給我們的家庭.
給我們有小朋友.
令我們有不同的關係.
其實甚至乎有壓力.
甚至經歷的人生裡面各種的挑戰也好.
不是單單要我們去忙碌這麼簡單.
而是要我們在這些各種的事情裡面.
學習成長.

$^{401}$學習忍耐.
學習愛.
學習付出.
學習信靠.
學習犧牲.
慢慢塑造成為更像天父的一個人.
有時候神不是先改變人的信仰.
而是先改變了人的身份.
先成為母親.
先成為父親.
先承擔起一種責任.
然後在這個過程裡面.
人慢慢開始懂得尋求神.
我認識一位香港家庭的朋友.
幾年前移民來英國.
其實他們在香港的時候.
幾乎沒有接觸過信仰.
很少回教會.
星期日多數是家庭日.
出去吃飯 逛街 看電影 做運動.
教會對他們來說簡直是一種陌生的地方.
但後來來到英國之後.
他們開始為孩子的成長擔心起來.
新的環境 新的學校 新的朋友.
又沒有親戚.
他們很擔心小孩子未必適應到.
所以人們說不如帶他們去教會.
於是他們開始帶小朋友去教會.
再加上很多時候.
很多學校都需要教會來填申請表格.
特別是牧者.
如果你想去基督教學校.
我都已經填過很多張了.
很多時候就是因為這樣進來教會.
做父母當然是很想小朋友找一個好的學校.
但是當小朋友去到一個學校的時候.
他感覺到一個陌生環境.
甚至對這個家庭.
剛才我所說的家庭.
小朋友經常說學校沒有他認識的人.

$^{441}$他非常難過.
所以他後來多了去教會的活動.
去教會的聚會.
令這個小朋友越來越多.
認識了不同的朋友.
而家長也認識了很多不同的教友.
結果他們慢慢就參加了教會的聚會.
他覺得教會的人很真誠.
甚至有些人很主動關心他們.
幫他們陪小朋友一起玩.
而久知他們每個星期都開始去教會.
他們去教會最初只是為了小朋友.
但是慢慢他們開始去聽講道.
有一次講道的內容.
特別是講道「不要為明天憂慮」.
那個爸爸回家之後就想.
其實我一直都很擔心未來.
錢會撕下撕下就沒有了.
工作又未找到.
所以都不知道怎樣算好.
結果他就學習這個講完所說.
要將他們憂慮的事交託給上帝.
慢慢這樣他就開始來祈禱.
雖然不是很正式的祈禱.
但是這樣就打開了他對於信仰的投入.
慢慢一年之後.
這個爸爸和媽媽一起信了耶穌.
甚至洗了泥.
其實這個爸爸就說.
我都不是為了信神而來教會的.
我只是為了小朋友.
但是神就在這段時間慢慢改變了我們.
弟兄姊妹.
有時候神不是先改變人的信仰.
而是改變了人的人生的處境.
在我們教會裡面.
都有一些小朋友很喜歡來教會.
然後家長看到小朋友這麼喜歡來教會.
他就陪小朋友去教會聚會.
是的 沒錯.

$^{481}$上帝很多時候不是先改變人的信仰.
而是改變了人的人生的處境.
可能要先讓我們成為父母.
甚至可以先讓我們的家庭進入一個新的環境.
甚至可以先讓我們承擔一些責任.
甚至面對不同的壓力都未定.
然後在人生的過程裡面.
人開始去尋找神.
其實舒蓮夫人的生命都是這樣.
一開始她只是接待神人.
但當她成為了母親之後.
她就開始更深的去經歷神.
信心就在生命的經歷裡面一步一步地成長.
最後第三件事就是信心的高峰.
就是在困境裡面如何去馬實神.
有一天悲劇突然之間發生.
孩子突然之間不知道什麼原因突發重病.
最後死了在他母親的懷抱裡面.
對任何一個母親來說.
這一刻一定是非常難受 非常痛心.
但是聖經的作者完全沒有描述到母親有什麼特別的情緒.
是不是很恐怖?.
很多時候母親都會有非常大的情緒.
因為突然之間小朋友死了.
甚至是離開了世界.
不單是重病.
對於家庭來說實在是一個很大的打擊.
母親的情緒通常都是非常難過.
又要痛哭崩潰.
但經文完全沒有提到她哭.
沒有提到她崩潰.
甚至沒有說她埋怨上帝.
反而聖經記載的一連串.
她很冷靜 有很迅速的行動.
她第一時間將小朋友抱上去神人房裡的床.
然後關上房門.
然後叫僕人快些預備勞車準備出發.
她甚至沒有跟她老公說這件事.
因為怕他心臟病發 這麼大年紀.
然後她就急急上路.

$^{521}$用了大概兩個小時的車程.
趕到神人以利沙的家裡.
如果這位婦人對神是完全沒有期望的話.
她是沒有理由這樣做的.
如果她覺得一切都完了.
她只需要留在家裡哀傷就已經足夠了.
但她沒有.
她選擇去找神人.
而且她很急.
她叫僕人趕快些.
不要浪費時間.
聖經甚至記載她叫她的僕人不要停.
不要跟人打招呼.
換句話說.
她不想浪費一分一秒.
她第一時間去見到神人.
因為她相信一件事.
如果神要做事.
現在應該還來得及的.
這就是當時她腦海裡面大概的信念.
而她最終見到以利沙的時候.
聖經記載了一個很有意思的動作.
她捉住以利沙的腳.
例如文的動詞.
其實已經是一個我們見過的動詞.
就是當年以利沙去到她家的時候.
這個婦人捉住她.
或者用極造挽留她.
留下她.
要以利沙一起吃頓飯.
用力挽留她.
現在同一隻字.
同一句希伯來文字.
就是捉住她的腳.
不放.
更有意思的就是.
婦人之後和以利沙說了一句話.
她說了這句話.
對不起.
這裡可能沒有打到出來.

$^{561}$也不一定.
沒有打到出來也不一定.
她就說我指著永生的搖滑.
有敢在你面前起誓.
我必不離開你.
這句話其實正正就是當年.
以利沙跟隨伊麗亞的時候.
她說過的那一番話.
在伊麗亞升天之前.
伊麗亞叫她走.
但是以利沙她死都不肯走.
她說我必不離開你.
當年以利沙的師父.
伊麗亞幾次叫她不要留下的時候.
以利沙就是說過這句話.
一年說了三次.
現在同一句話.
鼠年婦人她說了出來.
她好像提醒了以利沙.
你以前也是這樣捉緊神.
現在今天我也是一樣.
弟兄姊妹.
這個其實是信心在考驗裡面的高峰.
一開始的時候.
這個婦人對神其實真的不算太有興趣.
她只不過是為神預留了一個空間.
或者為神人預備一個房間.
但自從成為了母親之後.
她開始慢慢接觸神.
慢慢認識神.
定時性後來聽神人的講道.
而當人生最大的困境出現的時候.
她的兒子死了.
她的信心竟然沒有崩塌.
反而變得更堅定更堅強.
她在捉緊神.
但是要留意一件事.
就是鼠年婦人其實不知道神會怎樣做.
她不一定要相信神會叫她的兒子復活.
甚至在當時的信仰背景裡面.

$^{601}$死人復活都不是一個普遍預期.
但她仍然去找上帝.
為什麼呢?.
因為信心很多時候.
不是我知道神一定會怎樣做.
而是信心就是我知道神仍然存在.
只要神仍然存在的時候.
我就有出路.
我就有希望.
有一個很有名的基督徒.
在十九世紀.
我相信很多人都聽過他的名字.
他叫做Spethford.
原本是一個很幸福的家庭.
他的人生可惜是哲別年生.
受到很多打擊.
生意失敗.
又失去了兒子.
在1873年的時候.
他安排了太太和四個女兒.
坐船前往歐洲休息.
但途中船隻和另一艘船隻相撞.
最後只有他老婆生還.
四個女兒全部都葬身大海.
太太在英國發了一部很短的電報給他.
寫了一句話.
Saved alone, what shall I do?.
只要我一個人活著.
我可以做些甚麼?.
我可以怎麼辦?.
Spethford馬上就坐船去英國.
會合他老婆.
當船走到當年沉船的地點的時候.
船長特別告訴他.
就大概在這個地方.
你四個女兒就在這個地方沉了下去.
就在那一刻.
這個父親站在甲板上.
向朦朦朧朧的大海.
想到四個女兒長眠在海底.

$^{641}$但他沒有向神發出埋怨.
沒有控訴.
反而回到船艙寫下了一首詩歌.
就是今天很多地方仍然在唱的一首詩歌.
我心靈得安寧.
It is well with my soul.
有時享平安.
如江河平幽雲.
有時哀傷來到此浪滾.
不論何環境.
我已無主引領.
我心靈得安寧.
弟兄姊妹.
Spethford並不是因為他知道.
神一定會怎樣解決他的痛苦.
而即使在最深的悲傷裡.
他不知道應該可以怎樣做.
他不知道神會怎樣做的時候.
他仍然抓住馬實上帝.
這個就好像書店婦人一樣.
書店婦人未必知道神會怎樣做.
她未必敢肯定神一定會叫她的孩子復活.
但她仍然趕去找神恩.
因為她相信一件事.
只有神仍然在的時候就有希望.
因為她確信神在的時候仍然有出路.
很多人以為信心是突然出現.
但書店婦人的故事告訴我們.
信心往往是慢慢形成.
先是為神的國預留一個位置.
然後在生命裡經歷上帝.
最後在困境裡學得去馬實上帝.
去抓住上帝.
願主祝福大家.
\newpage



\section{約翰福音 17:20-23}
\label{sec:031808}
\textbf{建立健康的使命教會 (8) 有質素的相交}~~~~江山偉牧師
\newline
\newline
連結: \href{http://www.church.com.hk/acms/content.asp?site=cdc\&op=show\&type=product\&code=031808\&layout=sermon}{\texttt{031808}} ~~~~ 語音日期: 2026-03-15
\newline
\newline
\hyperref[sec:031807]{\small{< < < PREV SERMON < < <}}
~
\hyperref[sec:index_chronic]{\small{[返順時目]}}
~
\hyperref[sec:031809]{\small{> > > NEXT SERMON > > >}}
\newline
\newline
約翰福音 17:20-23
\newline
\begin{longtable}{cl}
\hline
\hline
章節 & 經文 (和合本修訂版)\\
\hline
17:20 & \begin{tabularx}{0.7\textwidth}{X} 「我不但為這些人祈求，也為那些藉著他們的話信我的人祈求， \end{tabularx} \\ \\ \relax
17:21 & \begin{tabularx}{0.7\textwidth}{X} 使他們都合而為一。正如父你在我裡面，我在你裡面，使他們也在我們裡面，好讓世人信是你差我來的。 \end{tabularx} \\ \\ \relax
17:22 & \begin{tabularx}{0.7\textwidth}{X} 你所賜給我的榮耀，我已賜給他們，使他們合而為一，像我們合而為一。 \end{tabularx} \\ \\ \relax
17:23 & \begin{tabularx}{0.7\textwidth}{X} 我在他們裡面，你在我裡面，使他們完完全全合而為一，讓世人知道是你差我來的，也知道你愛他們，如同愛我一樣。 \end{tabularx} \\ \\
[1ex]
\hline
\hline
\end{longtable}
$^{1}$各位弟兄姊妹 早晨.
來到加拿大一段時間.
都很體會 尤其是多倫多.
生活每個人都很忙.
又要工作 又要接家庭.
接送 又要做查柯夫司機.
載小朋友去這裡 去那裡.
很多的事情.
有時候我都在想.
不知不覺教會每個星期我們回來的時候.
其實都成為了一個很重要的加油站.
什麼意思呢.
就是說在星期日回來.
我們就像開車一樣進去油站一樣.
當然是想加滿油.
我們期盼就像剛才那樣.
有美好的詩歌.
一篇神的話語.
都成為了我們的力量.
在接下來的星期前行.
亦都可能臨走的時候.
我們都會跟身邊的弟兄姊妹打招呼.
談談最近近況.
然後我們的油缸就滿了.
繼續一個星期忙碌的生活.
其實沒有什麼錯與對.
我們每個人都是需要補給的.
尤其是屬靈的補給.
不過如果教會只是一個加油站.
其實遠遠是不足夠的.
因為神設立教會.
不是我們自己搞定 自己顧自己.
相聚完之後就各自上路.
神很期盼教會能夠成為一個屬靈的家.
你想想家裡給你一個什麼感覺.
你都期盼你的家人會關心你.
問候你 你最近怎樣.
工作怎樣 有什麼事情發生.
我們不需要裝作一切都很好.
在家裡你的真我就會很流露.

$^{41}$我不知道你回來的時候.
你的真我能不能露出來.
因為你都可能會有一個既定的標準的樣子.
標準的答案.
你都期盼有人可以在你人生的路上.
有人陪你走過此刻所面對的事情.
無論是開心還是不開心.
所以今天在這段時間.
我們說建立一個健康的使命教會.
第二個G.
就是我們怎樣去建立一個有質素的商家.
這個group.
完結之前我們必須要問一個問題.
就是究竟今天我們在教會裡面的生活.
在我們的小組裡面.
怎樣去活出來呢.
剛才那段經文提醒了我們一個很重要的原則.
就是有質素的商家.
我們才要回到耶穌基督裡面.
剛才那段經文是記載在《約翰福音》第十七章.
很多時候都會被稱為大祭司的祈禱.
其實是在耶穌臨升天的時候的一段很重要的禱告.
在這個禱告裡面.
耶穌不單止為了當時的門徒來祈禱.
也為了所有相信祂的人.
所有讀祂的話語的人.
來發出一個祈求.
也為了今天你和我來禱告.
當我們來到細閱這段禱文的時候.
我們會發現這個禱告有一個很核心的主題.
剛才不斷重複的字眼就是「合一」.
「合而為一」.
耶穌說要使他們都合而為一.
但是耶穌沒有停在這裡.
接下來祂再說一句話.
「正如父,你在我裡面」.
「我在你裡面,使他們也在我們裡面」.
其實這一句話是很重要的.
耶穌不是隨便用一個例子來形容教會的關係.
祂沒有說教會的關係就像朋友一樣.

$^{81}$也不是像同事一樣.
甚至也沒有說教會是一個社團.
或者是一個組織.
耶穌用什麼來形容教會的關係呢?.
就是透過三位一體神祂自己的關係.
父在子裡面.
子在父的裡面.
這是一種完全相屬,完全相愛,完全合一的關係.
神學家很多時候都會說.
三位一體的神本來就是一個愛的團契.
有父,子,聖靈.
在這三位一體裡面.
我們看到父,愛,子.
子又純復父.
聖靈又榮耀子.
三個位格之間不存在競爭的關係.
也沒有利益交換的關係.
而是一種彼此榮耀,彼此相愛的關係.
換句話說.
三位一體本身就是一個完全共融的團契.
耶穌在這裡說了一件非常震撼的事.
教會的合一就是要反映三位一體神的關係.
也就是說教會群體的生活.
應該是彼此的結連,彼此的接納.
彼此去分享的一種關係.
我會稱呼它為一個共融的關係.
什麼是共融的關係呢?.
我就想起薯蓉.
大家喜不喜歡吃薯蓉?.
你可能在外面經常叫一個薯蓉.
或者叫一個的時候會奔導過來.
你有沒有想過薯蓉是怎樣做出來的?.
首先你當然要買一袋的薯仔.
我會形容這一袋的薯仔.
每一個獨立的薯仔就好像我們一樣.
一個薯仔,薯女.
起初我們來到當中.
都是很獨立的一個個體.
但是當在整的過程裡面.
這些薯仔就要怎樣做呢?.

$^{121}$先把它煮熟.
接著就要用一些東西把它硬碎.
再要加些什麼呢?.
牛油.
還有呢?.
鮮奶.
喜歡的話再灑一些東西.
每個人有不同的做法.
混在一起的時候.
你還能不能分辨到誰在打誰呢?.
分辨不了.
所以一個共融的關係.
用這個類比去解釋了.
在教會裡面.
其實不再分你,你,我,我.
而是我們只有一個身體.
基督的身體.
當然在這個過程裡面.
是需要時間的.
不會突然間你今天來.
突然間你就去到這樣的狀態.
No.
但是也反映了.
其實我們每一個都需要像薯仔一樣被加熱.
被神的話語繼續去轉變.
扭轉改變我們的生命.
也都上帝透過生命裡頭.
我們一起在過程當中.
一起去經歷,一起去守望.
以致我們能夠逐漸成為一個基督的身體.
所以在教會.
並不是建立在大家的興趣裡面.
不是說你和我的興趣相同.
那我們走在一起.
No.
而是建立在基督.
耶穌基督的裡面.
基督就成為了我們當中的中心.
我們開始進入一個彼此共融的關係.
所以有質素的相交.

$^{161}$第二樣東西很重要的.
就是生命有著這種呈現.
在《新約聖經》裡面.
很多時候我們說相交,團契.
都會去到一個希臘字.
叫做Kononia.
這個字的意思不是一個普通的友誼.
而是說一個共享生命彼此相屬的一種狀態.
這種相交是遠超過.
平時我們所說的一些聚會.
聊聊閒聊.
或者是透過一些活動的層面.
而是說一種在基督裡面.
彼此在當中互相影響.
互相抵禮的一種關係.
在屬靈的道路上.
我們一起彼此的同行.
這個正正就是所謂的有質素的相交.
這種相交不只是帶來一些安慰.
更加在過程中塑造著我們的生命.
轉化著我們的生命.
很多時候我們都會說.
神透過什麼來改變我們的生命呢?.
透過祂的話語.
透過禱告.
不過原來與此同時.
上帝也透過我們身邊的弟兄姊妹.
來塑造我們的生命.
你嘗試看看你前後左右的弟兄姊妹.
我生命的塑造是由你來做的.
今日神就是將我們.
這裡有這麼多個「薯仔」「薯女」放在一起.
我們在當中去經歷.
上帝就是那個廚師.
如何透過在這十八年.
進入第十九年.
無論是在教會的進程.
又或者我們現在有十八個不同的小組.
可能我們每一個小組裡面.
其實我們都在經歷很多不同的事情.

$^{201}$上帝就是透過我們在當中.
如何向我們說話.
所以在帖撒羅尼迦前書五章十一節.
使徒保羅這樣說.
我們一起讀出來.
預備一二三起.
「所以你們該彼此勸勉彼此建立.
正如你們素上所做的」.
這段經文給我們一個很重要的屬靈提醒.
保羅並不是說我們要得著鼓勵.
而是說我們要彼此勸勉.
互相來建立.
在耶穌基督的生命裡面.
鼓勵不單止是我們被動來領受.
更加提醒我們一個重要的事.
就是我們如何主動去給予別人一個行動.
是有action.
把我們放在這個屬靈的家裡面.
讓我們能夠彼此互勉.
彼此來代禱.
使我們在信仰的道路上.
不會只有我們一個.
不會覺得孤單.
而是有身邊的領子妹同路人.
和我們一起去走.
一起去面對.
而在希伯來書十章二十四到二十五節.
亦都這樣提醒我們.
讀出來 預備起.
又要彼此關心.
激發愛心.
勉勵行事.
不可停止聚會.
好像那些停止慣了的人.
都要彼此勸勉.
這裡所說彼此關心的意思是什麼呢?.
我們經常都會說.
彼此關心.
就是說我們有沒有留心.
留意弟兄姊妹此刻.

$^{241}$屬靈的光景.
此刻他所面對的事情.
我們有沒有來到上上去記掛著他們呢?.
有沒有來到去思想.
怎樣去幫助他們去面對.
怎樣去成長呢?.
而激發愛心 勉勵行事.
就好像神透過我們的生命.
好像火一樣重新去挑釁我們.
使我們能夠去經歷那份.
從神而來的信心.
那份的愛.
那份對神的委身.
特別是當我們面對軟弱的時候.
再一次地來到被神去燃點.
在加拿大書保羅.
也有一節這樣的經文提醒我們.
我們一起讀吧 預備起.
你們各人的重擔要互相擔當.
此就會承傳基督的律法.
在這裡所說給我們看到.
真正的相交.
並不是只是單單停留在.
一些基本的關心和問候裡面.
而是我們願意進入對方的生命.
和他一起去擔當.
一起去面對他的重擔.
以致我們在當中能夠彼此扶持和守望.
所以無怪乎.
剛才我們也讀到.
神經常提醒我們.
我們經常掛在口邊.
不可停止聚會.
原因是什麼?.
原因就是因為在屬靈生命的塑造.
往往我們都需要在一個持續.
而且是一個很真實的相交裡面去發生.
在那裡我們一起去敬拜神.
一起去禱告.
一起去分享生命.

$^{281}$也一起去經歷神的作供.
當我們有時候會疲乏.
會有軟弱的時候.
這個屬靈的家.
就成為了我們的扶持和鼓勵.
共融的生命.
其實是不會自動來產生.
而是需要在我們屬靈的生活裡面.
慢慢去建立彼此來學習.
其實你想想.
從你第一天回來教會的時候.
已經有一套共融的生活守則.
灌輸在你的腦子裡面.
不過差在你在不在意.
我嘗試用六個英文字母.
叫做Closer.
Closer的意思大家知道.
就是很靠近.
當我們持續地去實踐的時候.
一方面這六個特質.
不單止讓我們更加親近神.
也幫助我們彼此之間去親近.
我們嘗試近鏡看看.
這六個英文字母所帶出來的特質是什麼.
第一個Closer的C.
Care for one another.
就是一個彼此的關懷.
塑造一個生命的相交.
其實是很需要從一個很真實的關懷開始.
一會兒我們也會互相去問候.
行圈幾句.
我們願不願意可以走多一步.
再問得深入一點.
或者留意到.
弟兄姊妹此刻的生活處境.
願意花時間去聆聽.
和去關心他們.
去發現.
究竟此刻他們背後.
面對著一些什麼事情呢.

$^{321}$可能是一些壓力.
挑戰.
需要等等.
又例如可能在小組裡面.
有弟兄姊妹面對著工作的壓力.
孩子成長的煩惱.
或者家人的健康出現了問題.
我們可不可以主動地去關心他們.
去為他們去禱告.
甚至乎隔幾天你又發一個訊息.
打個電話.
簡單一句的問候.
你最近怎樣?.
我一直都很記掛著你.
可能是一句這麼簡單.
又很直接的一個訊息.
其實就會讓對方.
好感受到.
他自己不是孤單.
而是背後有一群人.
和他一起去同行.
一起去面對.
他被記掛和關心.
所以真正的相交.
我們必須要先要彼此去看到.
看到對方的需要.
第二個Error就是.
Love in action.
即是要以行動去表達愛.
特別是一個屬靈的群體裡面.
愛不只是掛在後面.
而是需要有實際的行動.
有時候愛是.
最近我們當中都有一些頂子妹.
可能她面對軟弱病患.
我們都很好.
有一班頂子妹.
就是每天來支援.
這個家庭膳食的需要.
以至到她都能夠安心地去養病.

$^{361}$同時間我都很記得.
當我在不知道幾年之前.
兩三年之前.
一次咳到快要死.
我在同工的組裡去分享.
吳牧師第一時間就問我.
有什麼需要.
他有一個秘方.
就是用蜜糖.
不知如何混合一些薑.
喝完之後.
真的很不同.
他跟我分享.
他不單說這麼簡單.
隔沒多久.
直接車過來.
買了一支蜜糖.
買了所有的配料.
你煮吧.
在那一刻真的很感動.
原來愛是需要有行動.
所以在那一刻很感動.
我不知道.
其實我們恩泉屬靈的家.
我很感受到大家彼此的關懷和守望.
在當中.
其實我想每一個頂子妹.
來到當中都會不約而同跟我分享.
為什麼這裡這麼吸引你.
其中一件事.
就是這裡她看到愛是有行動的.
所以在塑造生命相交裡面.
愛是會透過這些.
可能對你來說微不足道.
但是它是因著你願意這樣去回應的時候.
在當中產生一個你意想不到的果效.
當這個人經歷了這份愛的時候.
我們又學如何將這份愛繼續傳給下一個.
所以真正的愛提醒我們.
不要只停留在說話.

$^{401}$而是真的能夠有行動.
第三個O就是open heart.
有一個躺開的心.
健康的相交.
其實很需要有一個真誠和開放的心.
試問如果我們每一次回來.
只是分享一些什麼順利的事情.
好像很難啟齒我們所面對的一些掙扎軟弱.
好像都收藏起來的時候.
其實我們這個屬靈的家.
其實很難真正地發生生命之間的真情流露.
當我們如果能夠很坦誠地去分享.
最近我的工作壓力很大.
有點累.
我為著未來為著此刻的戰爭.
都有點擔心.
或者最近我都面對一些事情令我有點不開心.
我們能夠可以很自然地在小組.
或者在教會和弟兄姊妹去分享的時候.
其實都會讓我們慢慢成為一個可以互相去代禱.
互相去守望一個群體.
所以同時之間.
躺開的心不單止是我們願意去分享自己的需要.
亦都包括我們需要有一顆柔軟的心.
去聆聽別人善意的提醒.
當弟兄姊妹可能是出於愛心.
給我們一些鼓勵的說話.
提醒勉勵的時候.
我們都願意以一個謙卑的心靈去聆聽.
看看神透過他的一個真誠的表達.
對自己如何繼續去成長.
第四個就是一個S.
就是Speak Truth in Love.
在愛中去講誠實話.
過去Pastor Alan和Pastor Ada.
都在他們的訊息裡提醒我們.
如果只是有真理沒有恩典.
很容易就給人一種很嚴厲的感覺.
讓人感受到傷害.
如果反過來只是有恩典沒有真理.

$^{441}$可能有時候我們都未必能夠去幫助到.
對方去有一個真誠的成長.
所以在一個屬靈的群體裡我們都求神.
知道其實都不容易的.
尤其是我們中國人.
我們都很怕傷害到別人.
所以有時候我們都是猜一聲不講.
我想我們在上兩次的訊息裡都給我們的提醒.
我們都需要去有那份智慧.
求神給我們有智慧.
我們背後所表達的不是出於論斷和批評.
而是出於愛和關心.
就好像我們當中有些頂尖妹因為忙碌.
而慢慢遠離了教會的生活.
又或者她長期受到工作和生活的壓力.
拖著她以至她和上帝的關係變得遠離的時候.
真正的愛是不會選擇沉默.
亦都不會帶著責備.
而是我們都會懷著一個關心.
溫柔地說最近好像很少見到你.
我很關心你.
有什麼可以為你去禱告.
這樣的提醒其實不是批評.
而是我們用愛的表達.
當我們用智慧溫柔愛心去講出誠實話的時候.
關係不單止不會被破壞.
反而會受到保護.
亦都幫助人重新回到正確的方向.
當然我們都要想一想.
當我們以為我們只是在講誠實話的時候.
其實我們都需要坐在神的面前去反思.
究竟這個所謂的誠實話.
是不是出於神的真理呢?.
抑或只不過是留於個人的看法.
個人的情緒.
或者一些你做事的習慣.
如果是出於自己的時候.
有時候我們真的要停一停.
想一想.
求神都讓我們先去過濾一下.

$^{481}$我們所說的話.
是不是真的能夠造就到弟兄姊妹.
抑或反過來你說出來的時候.
傷害了他.
所以真理如果沒有了愛.
就會傷人.
如果愛沒有了真理的話.
其實就會令到人容易迷失.
而接著那個「E」呢.
就是encourage one another.
彼此鼓勵.
在一個信仰群體裡面.
我不知道.
你自己問一問.
你多不多接收到人對你的肯定和鼓勵.
這是第一.
你問一問自己.
在你這段時間回到隱傳家裡面.
你收得多不多呢?.
第二個問題我問你.
你多不多自己去說.
向人說一些鼓勵和肯定的說話呢?.
除了收之外.
你也要留意.
你自己要說剛才經文說的.
我們要彼此的去造就.
而不是一味的去接收.
這兩方面都是很重要的.
其實我們是一個鼓勵的說話.
箴言也說了.
一句說話其實都可以成為一個人的祝福和勉勵.
金蘋果落在那個網址裡面.
所以我們怎樣去到在.
神亦都今天這樣來看好.
你和我的時候.
我們怎樣能夠去到.
透過上帝給我們一些感動.
一些我們去欣賞對方的事情.
能夠透過我們的表達來出現.
所以鼓勵很多時候不需要一些很長的分享.

$^{521}$有時只是一句簡單的肯定.
一個簡短的祈禱.
特別在人的軟弱的時候.
都成為了他的扶持.
信心都能夠去得到那個堅固.
關係亦都因為這樣而更加緊密.
而最後的R.
就是一個Remain accountable to one another.
彼此的去守望.
彼此的守望不是要控制對方的人生.
而是要在愛裡面關心他屬靈的生命方向.
當我們看到弟兄姊妹迷失的時候.
我們願意來靠近,陪伴.
用禱告,鼓勵的說話幫助她重新回到神的面前.
可能透過我們.
最近你的靈修生活是怎樣.
其實我們都可以豐富彼此之間的問候.
除了問你今天吃飯了沒有.
最近工作忙不忙.
今天早餐吃甚麼之外.
其實你都可以問.
在最近你的生命裡.
神有沒有透過一些事件向你和你說話.
又或者你在生命裡.
這個星期有甚麼掙扎或感恩.
我可以為你去祈禱呢.
當我們整個屬靈的群體.
都是這樣來彼此守望的時候.
我們不再是自己去努力.
去做一個門徒.
而是真的能夠回來.
在耶穌基督的裡面.
神就是將你和我帶領來到這個屬靈的家.
我們一起向前去成長.
\newpage



\section{約書亞記 3:7-17}
\label{sec:031809}
\textbf{神蹟的出現}~~~~劉少康牧師
\newline
\newline
連結: \href{http://www.church.com.hk/acms/content.asp?site=cdc\&op=show\&type=product\&code=031809\&layout=sermon}{\texttt{031809}} ~~~~ 語音日期: 2010-03-02
\newline
\newline
\hyperref[sec:031808]{\small{< < < PREV SERMON < < <}}
~
\hyperref[sec:index_chronic]{\small{[返順時目]}}
~
\hyperref[sec:031810]{\small{> > > NEXT SERMON > > >}}
\newline
\newline
約書亞記 3:7-17
\newline
\begin{longtable}{cl}
\hline
\hline
章節 & 經文 (和合本修訂版)\\
\hline
3:7 & \begin{tabularx}{0.7\textwidth}{X} 耶和華對約書亞說：「從今日起，我必使你在以色列眾人眼前被尊為大，使他們知道我怎樣與摩西同在，也必照樣與你同在。 \end{tabularx} \\ \\ \relax
3:8 & \begin{tabularx}{0.7\textwidth}{X} 你要吩咐抬約櫃的祭司說：『你們到了約旦河的水邊，要在約旦河中站著。』」 \end{tabularx} \\ \\ \relax
3:9 & \begin{tabularx}{0.7\textwidth}{X} 約書亞對以色列人說：「你們近前，到這裡來，聽耶和華－你們神的話。」 \end{tabularx} \\ \\ \relax
3:10 & \begin{tabularx}{0.7\textwidth}{X} 約書亞說：「你們因這事會知道永生的神在你們中間，他必從你們面前趕出迦南人、赫人、希未人、比利洗人、革迦撒人、亞摩利人、耶布斯人。 \end{tabularx} \\ \\ \relax
3:11 & \begin{tabularx}{0.7\textwidth}{X} 看哪！全地之主的約櫃必在你們的前面過去，到約旦河裡。 \end{tabularx} \\ \\ \relax
3:12 & \begin{tabularx}{0.7\textwidth}{X} 現在， 你們要從以色列支派中選出十二個人，每支派一人。 \end{tabularx} \\ \\ \relax
3:13 & \begin{tabularx}{0.7\textwidth}{X} 當抬耶和華全地之主約櫃的祭司，腳掌踏入約旦河水裡的時候，約旦河的水，就是從上往下流的水，必然中斷，豎立成壘。」 \end{tabularx} \\ \\ \relax
3:14 & \begin{tabularx}{0.7\textwidth}{X} 百姓起行離開帳棚過約旦河的時候，抬約櫃的祭司在百姓的前面。 \end{tabularx} \\ \\ \relax
3:15 & \begin{tabularx}{0.7\textwidth}{X} 那時正是收割的日子，約旦河的水漲滿兩岸。抬約櫃的人到了約旦河，抬約櫃的祭司腳一入水邊， \end{tabularx} \\ \\ \relax
3:16 & \begin{tabularx}{0.7\textwidth}{X} 那從上往下流的水就在很遠的地方，在撒拉但旁邊的亞當城那裡停住，豎立成壘；那往亞拉巴海，就是鹽海下流的水全然中斷。於是，百姓在耶利哥的對面過了河。 \end{tabularx} \\ \\ \relax
3:17 & \begin{tabularx}{0.7\textwidth}{X} 抬耶和華約櫃的祭司在約旦河中的乾地上穩穩站著，以色列眾人都從乾地上過去，直到全國都過了約旦河。 \end{tabularx} \\ \\
[1ex]
\hline
\hline
\end{longtable}
$^{1}$今集的題目是「神跡的出現」.
在《約書亞記》第三章七到十七節.
就出現了一個神跡.
原來約書亞吩咐十二個支派.
每個支派派一個人.
就一起抬著神的躍跪.
當神吩咐起行的時候.
這些祭司就抬著躍跪.
就踏著腳進入約旦河.
當祭司的腳一放入水裡面.
約旦河的河水就斷開成為壁壘.
而且那些水是在很淺的地方已經止住了.
結果他們慢慢就走進這個乾了的河床裡面.
到底為什麼會有這樣的事情呢?.
好像不可能.
為什麼會出現呢?.
這些我們就叫做神介入.
在歷史事件當中我們就叫做神跡.
在聖經約書亞記第三章還有一個詳細的描寫.
他說原來約旦河的河水.
在修葛的日子是漲過兩岸的.
大家都知道.
那些河水特別當一些雨季之後.
那些水在上游一直沖下來.
那些河水就會氾濫掩蓋兩岸.
我相信當日當這些祭司踏著腳進入水的時候.
可能剛剛是修葛的日子.
所以河水是大跌漲的.
但是竟然間他們在很遠很遠上游的地方.
那些水已經止住了.
於是聖經有一個很美麗的記載.
就是他們在乾地上過去.
本來是一條很高水很漲的河.
現在竟然間變為乾地.
你生命當中有沒有經歷過神這樣捉摸你呢?.
你有沒有見過一些弟兄姊妹.
他們會講見證.
神就是這麼奇妙.
神就是這麼真實.
神就是這樣在我的生命裡面.

$^{41}$做了又大又奇妙的工作.
讓我們就降伏在神面前.
向祂敬拜.
讚美祂的作為.
原來我們的神就是一位行祈祀的神.
你不要以為只有別人才有的.
我們到處去看一看.
哪個人有什麼神蹟出現.
神都同樣能夠在你.
在我.
在我們的身上顯出祂自己的作為來.
所以約書亞記第三章很感動.
當整大群人就跟著這些祭祀.
一對一對的就過河了.
當過河之後.
他們就進入應許之地.
今天神同樣要帶你進入生命更豐盛的地方.
\newpage



\section{馬太福音 7:7-12}
\label{sec:031810}
\textbf{哪裡會有好東西?}~~~~石偉漢牧師
\newline
\newline
連結: \href{http://www.church.com.hk/acms/content.asp?site=cdc\&op=show\&type=product\&code=031810\&layout=sermon}{\texttt{031810}} ~~~~ 語音日期: 2026-03-15
\newline
\newline
\hyperref[sec:031809]{\small{< < < PREV SERMON < < <}}
~
\hyperref[sec:index_chronic]{\small{[返順時目]}}
~
\hyperref[sec:031811]{\small{> > > NEXT SERMON > > >}}
\newline
\newline
馬太福音 7:7-12
\newline
\begin{longtable}{cl}
\hline
\hline
章節 & 經文 (和合本修訂版)\\
\hline
7:7 & \begin{tabularx}{0.7\textwidth}{X} 「你們祈求，就給你們；尋找，就找到；叩門，就給你們開門。 \end{tabularx} \\ \\ \relax
7:8 & \begin{tabularx}{0.7\textwidth}{X} 因為凡祈求的，就得著；尋找的，就找到；叩門的，就給他開門。 \end{tabularx} \\ \\ \relax
7:9 & \begin{tabularx}{0.7\textwidth}{X} 你們中間誰有兒子求餅，反給他石頭呢？ \end{tabularx} \\ \\ \relax
7:10 & \begin{tabularx}{0.7\textwidth}{X} 求魚，反給他蛇呢？ \end{tabularx} \\ \\ \relax
7:11 & \begin{tabularx}{0.7\textwidth}{X} 你們雖然不好，尚且知道拿好東西給兒女，何況你們在天上的父，他豈不更要把好東西賜給求他的人嗎？ \end{tabularx} \\ \\ \relax
7:12 & \begin{tabularx}{0.7\textwidth}{X} 所以，無論何事，你們想要人怎樣待你們，你們也要怎樣待人，因為這就是律法和先知的道理。」 \end{tabularx} \\ \\
[1ex]
\hline
\hline
\end{longtable}
$^{1}$好,兩姐妹平安,我們一起先有祈禱.
親愛的,多謝你讓我們在這一刻的時刻.
我們要去敬拜你,我們要去讚美你.
我們要聆聽主你的道,藉著你的道去建立我們的生命.
塗造我們,以致我們真的能夠按主你的心意去行.
求主你幫助我們,聽我們祈禱奉耶穌的名字而求,阿們.
今天跟大家去思想的題目叫做「哪裡會有好東西」.
我不知道你有沒有一些好東西要和別人分享.
事實上我們每一個人都很希望找到一些好東西.
最近在中國都有一些好東西,人們爭著去要.
不知道你知不知道,有沒有看這些新聞.
要養一種動物,叫做養龍蝦.
不知道你有沒有看過.
最近在國內很多人爭著的叫做OpenClaw.
就是一個軟件來的,一個養龍蝦的東西.
大家有沒有人知道是什麼東西,先了解一下.
都沒有人知道,只有一個人知道是什麼.
兩個人知道,這些是IT人來的.
什麼來的呢,其實就是一個類似AI的東西.
我們現在用AI,很多資訊科技.
我們平時用的AI就是一個智者.
你有任何問題不明白就問那個智者.
就是平時的check GPT那些東西.
一有不明白就問他,他馬上解答你.
這個就是平時用的AI.
而這個OpenClaw就是一個助理.
明白分別了,智者就是你見到這個就是智者.
你就去問他問題,就是智者.
而這個就是你的好幫手,你的生活助理.
這個軟件這個程式就是安裝在你的電腦裡面.
他就會幫你做你平常要做的事情.
你家裡有什麼要做,平常要做.
你找這個軟件裝在你的電腦裡就幫你解決.
包括什麼呢,譬如你想要入數.
有時候要銀行過錢.
有些東西要寫程式,寫程式都懂的.
你要幫你投資等等的東西.
你只要講給這個軟件聽,他就會幫你做.
而且他會很醒目的,他會知道你的習性.
即是他慢慢知道你的日常生活之後.

$^{41}$他知道什麼時候你就要些什麼.
你講某一些事,你再講回那件事.
他會記得的,然後他就會馬上回應.
你就是這樣想的,你就想這樣做.
他就會幫你按著你的喜好,幫你處理你的事.
你聽到都很重要.
所以在國內裡面很多人爭著要去安裝這個軟件.
成為了在中國裡面的一個養龍蝦的潮流.
為什麼養龍蝦呢,因為他的標誌.
這個圖案的標誌就是龍蝦.
所以他說,幾分鐘就完成了你一個星期的工作.
是不是很棒,有一個助理,幾分鐘而已.
他已經知道你做什麼,幫你處理你工作上需要的東西.
所以我網上看到有一個老闆.
因為安裝了這個軟件之後,第二天就解僱了那個助理.
因為他已經幫你做了所有的工作.
而安裝這個軟件其實都需要一些錢.
所以後來你見到在國內,在深圳騰訊的一個總部.
有很多人排隊,因為好像是免費幫人安裝龍蝦.
很多人去排隊,真是很熱鬧.
你聽到都應該是好東西.
幫你做這麼多東西.
本來你很辛苦做的,現在他都幫你搞定了.
本來你要做幾個小時的,現在幾分鐘就做好了.
但你有沒有想過,當你以為是好東西的時候.
其實是一件很高危的事情.
因為安裝這個軟件的其中一個特色就是.
你要開放你的權限.
就是說你在電腦裡面有些權.
你要給權他,他才可以幫你投資.
你給權他,他才可以幫你做事.
給權力他,他才會幫你把不要的文件刪除.
把你的文件重新分類.
整個過程裡,他差不多掌管你的一切.
當你錯誤運用,你就有難.
所以這個熱潮一方面好像很厲害.
但原來過了沒多久,現在人們都在說.
養這個龍蝦是很嚴重的.
因為當你去運用他的時候.
就算你以為你沒有在用他.

$^{81}$其實他已經在做事了.
而且每一次做事,他都要給token.
要給一點費用給那些公司.
你以為他做一點事,其實他用了你幾千個token.
所以有人說安裝了這個軟件之後.
不知不覺有張賬單是幾萬港幣的賬單.
你都會收到.
等於你請一個工人,請一個助理的錢.
而且他不知不覺在用你的錢.
這是其中一個,當然更重要的是.
他在用你很多的權限和資訊.
甚至令到你很多東西都沒有了你的檔案.
你又回不了頭.
所以這樣說,有些你以為是好東西.
原來可能又不是真的那麼好,要看清楚.
但有時候有些我們以為是普通的東西.
但你有沒有想過這些可能是你的珍寶呢?.
今天我們去看的經文是問一個問題.
你相不相信上帝會將好東西給你呢?.
今天讀的經文是來自馬太福音第七章的段落.
第七節開始到第十二節.
當中有一個問題就是.
你相不相信上帝會將好的東西給你?.
又或者這樣說,你心目中的好東西是什麼呢?.
什麼是你的好東西呢?.
首先當然要有些少回顧.
我們上一次就說到第七章六節.
我們已經進入了登山部份的第七章的部份.
我們上次看到一至六節.
而上次第六節的最後就說.
不要把聖物給狗.
也不要把你們的珍珠丟在豬的面前.
恐怕他們就踐踏了珍珠.
這是一個這樣的結構,踐踏珍珠的就是那些珠.
接著他說轉過來就咬你們.
應該就是在說那些狗.
那些狗會咬人.
那些珠不懂得珍惜珍珠.
而這裡上次我們看到就是說.
原來耶穌特別說不要去評斷別人的時候.

$^{121}$我們更加要懂得去分辨.
有這個智慧去分辨事情.
就好像我們上次看到.
在腓立比書一章十至十一節說.
使你們能分辨是非.
在基督的日子真誠無可指責的人.
更靠著耶穌基督.
結滿仁義的果子.
足歸榮耀稱讚給上帝.
所以我們要懂得去分辨.
特別要懂得分辨到什麼呢.
找回什麼是聖物.
什麼是好的屬上帝的.
什麼是珍珠寶貴的.
原來呢.
接著後面去到第七節.
說到原來有些東西好像聖物一樣.
好像珍珠一樣.
很重要的.
很寶貴的.
我們沒有去珍惜.
我們沒有好好去爭取這些東西.
這些才是我們所要去看重的.
所以今天我們就會看到.
第七至第十二節的經文.
這段經文可能大家都很熟悉.
經常看的經文.
想問一個很簡單的問題.
就是你這段經文.
我讀一次讓大家有些回憶.
我很快地讀一次.
你們祈求就給你們尋找就找到.
叩門就給你們開門.
因為凡祈求的就得著.
尋找的就找到.
叩門的就給他開門.
你們中間誰有兒子求餅.
反給他石頭呢.
求雨反給他蛇呢.
你們雖然不好.

$^{161}$上帝知道拿好東西給兒女.
何況你們在天上的父.
祂豈不要把更好的東西.
給求祂的人嗎.
所以無論何事.
你們想要人怎樣待你們.
你們也要怎樣待人.
因為這是律法和先知的道理.
看完這個段落.
想問一個問題.
這段聖經是說什麼的呢.
回教會這麼久了.
經常都讀這段聖經.
通常你會覺得這段聖經是說什麼的呢.
是說祈禱嗎.
是不是.
同意嗎.
通常你回教會都會說是說祈禱.
但你有沒有想過.
其實這段聖經不是說祈禱.
我們經常讀聖經.
有些習慣就是.
我們慣性回教會這麼久了.
通常讀到這段就是說祈禱.
向天父求.
但我想問.
究竟這段聖經是不是說祈禱呢.
耶穌在登山寶訓裡.
記不記得在哪裡說祈禱.
在第六章.
如果你有印象.
我們一路看登山寶訓的時候.
在第六章耶穌說祈禱.
耶穌在第六章的第五節開始說.
你們禱告的時候不可將假冒為善的人.
你要這樣禱告.
你們不可將愛邦人祈禱.
然後教耶穌教導的土民.
有沒有印象.
所以你記得耶穌教祈禱的時候.

$^{201}$其實是在第六章做的.
耶穌在第六章的第五節.
一路說說祈禱.
甚至說到我們大家熟悉的.
我們所謂的主土民.
就是耶穌教導的土民.
這就是耶穌的祈禱.
但你有沒有想過.
耶穌一路說登山寶訓.
由說八福一路開始說到.
是炎是光一路發展.
第六章已經說完祈禱了.
第七章開頭說.
你們不要評斷別人.
為什麼突然去到第七章七節.
突然又跳回去說祈禱呢.
那你是不是想想.
是很奇怪.
那個思路是不連貫的.
不知道你明不明白我的意思.
所以思路不連貫的時候.
代表了什麼呢.
其實如果耶穌前面說祈禱.
已經教完最好的祈禱.
也教了你怎樣祈禱的話.
其實這個段落.
可能不是說祈禱.
是我們一路誤解了.
是一路我們習以為常.
我們就以為這段聖經.
就是要說祈禱.
所以今天我們也去想一想.
這段聖經是不是說祈禱呢.
可能有機會是.
不過我自己一路看下去.
我就覺得應該不是.
我們看看找一個答案.
我們先看第一個部分.
A的部分.
就是第七節和第八節的部分.

$^{241}$這裡說.
你們祈求就給你們尋找就找到.
叩門就給你們開門.
我們為什麼經常覺得.
這段聖經是說祈禱呢.
因為一開始說你們祈求就給你們.
我們有個習慣.
你們祈求就是祈禱.
所以你們就會自然覺得.
這段聖經就是說祈禱.
因為你們祈求就給你們.
叫我們多點祈禱.
但是你留意.
剛才我們說到.
在馬太福音第六章第五節.
到第八節的時候.
耶穌說你們禱告的時候.
when you pray.
是用了pray的字.
六章第五節用了pray.
之後用了三次pray.
就是祈禱的字.
但這裡不是用祈禱的字.
這裡是用ask.
seek和lock.
你會發現原來這裡說.
耶穌說到這個段落的時候.
其實不是說你們祈禱就給你們.
而是你們祈求.
你去問一些事情的時候.
就給你答案.
你去找的時候就找到.
你去敲門的時候就開門.
所以你看回這三個句式裡面.
其實沒有說.
當你們祈禱的時候就給你們.
就聽你們祈禱.
不是這樣說的.
希望你明白.
我在想這段聖經說什麼.

$^{281}$所以這裡其實用了三個動詞.
ask, seek和lock.
這三個字去表達.
這個段落裡面一個很重要的意念.
而且這三個字是一個命令的字.
耶穌有個命令就是我們要去祈求.
我們要去尋找.
我們要去敲門.
而且這個命令是什麼呢.
是一個我們今天理解為.
一種現在進行式一樣.
就是keep asking.
要繼續不斷去問.
不斷去問.
keep seeking.
繼續去尋找.
keep locking.
繼續去敲門.
敲門敲門.
敲到有人認你為止.
是在說這件事.
所以這裡透過這三個命令.
三個現在進行式的方式的表達.
其實是在表達什麼呢.
一個堅持不斷.
長久不懈的一個過程.
或者一種態度.
去面對著去找到好東西的事情.
其實這裡是在說.
後面我們會說到好東西的時候.
是在說我們對好東西的態度是什麼呢.
原來我們要去堅持不斷.
長久不懈.
不要停.
不要放棄.
當然我們祈禱.
當然我們不要輕易放棄.
但這裡這個段落.
更加指向一個更重要的.
就是一些好東西的時候.

$^{321}$我們更加要好像聖物.
好像珍珠一樣.
努力不斷去追求.
不要去放棄.
有一句成語.
不知道是不是成語.
就是三顧草廬.
大家聽過了.
大家三顧草廬.
你知道在以前的年代.
就是什麼.
沒有獵頭公司.
是不是.
就是你要找人才.
就不是這麼簡單.
今天你登一登報紙.
有些公司幫你找人才.
或者透過網上很快找到.
以前的年代就是.
你要找一個好的人.
合適的人是不容易的.
特別在古代要打江山.
更加需要一些智者.
所以三顧草廬就是來自.
當年去找這個智者的一個故事.
那個智者是誰.
是不是.
諸葛亮.
諸葛孔明.
找諸葛亮的過程裡面.
就是說當時原來在東漢末年.
就是秦朝之後.
秦秦記我們說完了.
漢朝.
漢朝之後西漢之後.
就東漢.
去到東漢末期.
其實皇室已經是慢慢去沒落.
應該是捱不了多久.
所以當時.

$^{361}$其中有一個屬於漢室的後人.
就是誰.
劉什麼.
劉備.
大家很熟悉 劉備還有什麼.
關羽 張飛那些.
大家很熟悉.
三國演義那些人物.
就是劉備很想找一個人去幫助他.
怎樣去復興漢室.
能夠幫他打江山.
以致到他後來.
在人的口裡面.
就聽到有一個智者.
這個智者就是叫諸葛亮.
所以他為了要找到這個智者.
他就去找.
去到他家裡面去找他.
結果第一次找不到.
第二次又找不到.
如果他就這樣放棄了.
那就沒得三國鼎立了.
直到第三次.
終於他見到諸葛亮.
和他一起在談天下形勢.
於是諸葛亮就告訴他.
教他怎樣.
他就說.
現在在天下的形勢裡面.
最好的方法是什麼.
教他先取荊州.
再取益州.
接著就逐鹿中原.
就要拿中原.
於是在過程裡面.
劉備就更加確定到.
諸葛亮就是他要找那個軍師.
邀請他出山.
這個過程裡面.
不是一次就可以了.

$^{401}$三次去敲人家家門.
三次去到他裡面去找他.
聯絡他和他聊天.
以致邀請他.
所以今天我們第一個要想的.
就是我們今天.
其中一種我們.
尋找好東西的態度.
就是我們有沒有.
沒有停下來.
我們要繼續堅持.
當你認定覺得是好東西.
不要輕易放棄.
劉備都三顧草率.
你呢.
特別我們知道我們.
一個很重要的情況.
能不能保持.
我們很大的問題.
就是很多時候我們面對著.
做一些事情的時候.
我們就是.
總是做到一半.
就做不到了.
就是開頭就很慶賀的.
做到一半就開始.
就借皮的了.
搞不定的了.
會不會都是這樣呢.
我們剛剛我們有.
今年有什麼.
有個讀經計劃.
讀經計劃現在已經.
讀到兩個半月的了.
不知你讀這兩個半月怎樣.
我記得初初我們.
讀經計劃頭一個月.
大家很興奮.
會在那裡留言.
說一些話.

$^{441}$或者在分享裡面.
或者在祈禱會裡面.
大家都會說讀經裡面.
有什麼得著.
什麼啟發.
怎樣鼓勵你.
怎樣堅持.
但是呢.
這陣子好像沉寂了很多.
沒什麼聲音了.
為什麼會這樣呢.
因為我們讀得久.
我們就疲累了.
但是我們有沒有想一下.
我們讀經讀得久.
令到我們有讀經的疲累.
還是我們反而因為我們疲累.
我們要更加認真讀聖經呢.
很容易的.
因為我們的身心靈很疲累.
我們就覺得想放棄.
想停止.
明知道讀聖經是好.
不過我們都說.
我還是放下了.
還是正正好因為知道.
你很疲累.
你身心靈都很累.
工作很繁忙.
天天工作到像狗一樣.
這個狀況.
你更加要讀聖經.
在靈裡面給到你滿足和安慰.
很重要的.
耶穌怎麼說呢.
耶穌說.
凡奴婦擔重擔的人.
可以到我這裡來.
我就可以讓你得到安息.
就因為你擔了很多重擔.

$^{481}$你才要來到耶穌面前得到安息.
所以我們正正好.
我們很累.
我們去找神.
這個才是真正的王道.
而在經文第八節.
就很清楚告訴我們.
因為凡祈求的就得著.
尋找的就找到.
叩門的就給他開門.
原來當我們堅持不懈的時候.
耶穌給我們的答案是什麼呢.
我們會得到的.
我們會找到的.
我們會開到門的.
耶穌給我們的答案是肯定的.
不過肯定的背後是什麼呢.
就是凡祈求的就得著.
凡尋找的就尋見.
凡叩門的就給他開門.
弟兄姊妹.
你明不明白這個凡是什麼意思.
有一句話.
有個歇後語.
不知道你們有沒有聽過.
就叫做神仙放屁.
這個歇後語是什麼.
不同凡響.
聽得明白嗎.
神仙放屁的響聲.
是和凡人不一樣的.
所以是不同凡響.
有時候我們都會覺得.
自己好像凡人一樣.
一個很平凡的人.
甚至有時候覺得自己.
是一個很凡的人.
但這裡聖經說.
凡祈求的那個凡人.
是不是代表我們平凡的人.

$^{521}$就沒有資格去尋找好東西呢.
聖經告訴我們.
這個是什麼.
For everyone.
聖經說這個凡祈求.
就是for everyone.
每一個願意去祈求.
尋找叩門的都可以找到.
所以這個更加是.
我們不可以停下來的原因.
因為你一定會找得到.
這些好東西.
重要的就是你去不去找.
有沒有中間放棄呢.
讀聖經能夠幫助你.
你有沒有中間放棄呢.
神有很多好東西等著你.
你有沒有中間放棄呢.
我們對信仰的堅持.
有時候就是要這樣去堅持到底.
不只是說祈禱.
有很多信仰的生活.
我們都需要堅持.
既然有這個時候.
接著後面怎麼說呢.
接著後面去到第九第十節.
他說你們中間誰有兒子求餅.
反給他石頭呢.
求雨反給他蛇呢.
第二個標題.
我給你這個就叫好壞之間.
差之毫厘還是差之千厘呢.
好東西和壞東西之間.
究竟分別是很小還是很大呢.
可能對某些人來說.
是很少的.
但有某些人來說.
就覺得是很大的.
很大的分別.
這裡我們看到.

$^{561}$耶穌特別說了四樣的物品.
包括餅,雨,石頭和蛇.
就是祂所說的四樣的事情.
而這四樣的事情裡面.
很明顯就見到.
有兩個是屬於好的.
就是餅和雨.
不好的就是石頭和蛇.
為什麼這樣說呢.
為什麼說你們中間有兒子求餅.
反而給他石頭呢.
為什麼求餅和求石頭.
是有一個對應的關係呢.
其實這個也是因為.
和當時的文化很相關.
因為原來猶太人本身.
他們其中一種主要的食物就是餅.
而餅很多時候就是圓形的東西.
一些扁平的東西.
這些就是他們餅的形狀.
當然有些大有些小.
但基本上都是一些圓形.
扁平的東西.
原來在巴勒斯坦地.
同樣他們有很多石頭.
在石灰岩或者玄武岩的.
叫做鵝卵石.
也是扁平又是圓的.
尺寸也是差不多大小的.
原來有些石頭都是這樣的.
所以餅和石頭.
在有些時候的角度看.
看得不清楚就會差不多.
特別你想想.
今天我們可能不太明白.
例如你兒女去廚房拿東西吃.
開個燈很光猛.
打開個雪櫃又會亮燈.
看得很清楚你的食物.
但在耶穌的年代.

$^{601}$你知道不是這樣的.
沒有電燈.
所以在家裡找食物的時候.
其實很多時候.
都不是一些很光猛的環境.
光就算很微弱的光打出來.
有時候都不是很明顯.
看得到這是餅還是石頭.
所以對於當事人來說.
很明白有時候不認真看.
都分不清的.
不單是餅和石頭.
有這些相類的狀況.
魚和蛇.
有時候都很近似.
這幅圖的左邊.
就是在巴勒斯坦地區.
都會找到的一些魚.
這個魚就叫這個名字.
我都不懂讀這個名字.
就是這種魚.
有人說這種魚.
是其中一種Catfish.
即是鯰魚.
一些這樣的形狀.
你見得到不是一些.
我們平時想的那些.
鯧魚那樣的菱形形狀.
而是一些比較長身.
會轉彎擺動的魚.
原來當時都有這些魚.
是他們生活中見得到的魚.
而右邊就是一條蛇.
甚至這些就是他們一些.
我又不懂讀這個蛇名.
這些蛇就是.
甚至在水中出現的水蛇.
即是當你在海中.
去撈獲一些魚獲的時候.
一方面固然可以有一些.

$^{641}$正常你見到的魚.
但可能有時候.
這些蛇.
都會在這些魚獲當中.
困在當中.
而且你想想.
這裡是黑色.
我又再考考大家.
在加利利湖裡面.
那些漁夫是何時去捕魚的.
晚上.
你記不記得耶穌有一次.
當那些漁夫經過整夜捕魚之後.
甚麼都沒有收獲.
在洗網的時候.
耶穌叫他們下網.
記不記得.
因為他們在晚上捕魚.
所以如果正常在加利利湖.
晚上捕魚.
撈到一些東西的話.
其實有時候你分不清楚.
那條是魚還是蛇.
如果你把手伸進去.
就這樣拿去.
可能你會被蛇咬.
因為通常在晚上.
在一些比較暗的環境.
又是看得不是很清楚.
有時候分辨不到.
所以我就是說.
有時候分別可能很少.
有時候你真的不經意.
就看不到那個分別.
所以對某些人來說.
有時候餅和石頭.
似是易分.
但很多時候.
有時候是難以分辨.
魚和蛇也一樣.

$^{681}$有時候好像易分.
有時候很難分辨.
但無論如何.
我們就知道.
魚和餅.
是代表著不同的物品.
特別是這個世代裡.
其實有很多東西.
都很難分辨.
就算你是專家.
也可能出事.
也可能中伏.
最近你有沒有看這單新聞.
他說德國有一家.
很出名的散熱名牌.
叫暴力熊.
因為這些原材料.
越來越貴.
於是就去上一個購物平台.
去買原材料.
買了四萬歐羅.
買回來之後.
他說鋼板.
他說鋼板假裝鋁板.
他要鋁板.
然後純銅變成導銅.
結果呢.
這四萬塊的貨.
全部都不能用.
已經是一個.
買開這些材料的公司.
原來去到這些地方.
原來都會被人騙的.
都會被人騙的.
今天也是.
原來有很多東西.
都會令我們搞錯了.
究竟我們今天.
擁有的是餅.
還是石頭呢.

$^{721}$是魚還是蛇呢.
在這裡說.
餅是代表了生命的糧.
代表了一些營養.
供應.
魚是代表一些祝福.
令你很滿足.
身體能夠得到幫助的東西.
但是石頭呢.
如果你吃石頭.
咬崩你的牙.
代表了一些無用和失望.
這些蛇呢.
代表了一些危機.
也代表了不潔.
因為這些石頭是沒有鱗的.
是不潔的.
那麼難以分辨的時候.
這段聖經就告訴我們.
對父親來說.
他一定能夠分辨.
我們做小孩子.
我們有時候隨便.
一看到桌面就拿來吃了.
但是對父親來說.
他為了兒子的原因.
他會分辨好.
什麼是最好的東西.
先給子女的.
所以你相不相信.
天父將最好的東西.
給你我們每一位.
屬靈的兒女呢.
C.S.Lewis.
不知道大家知不知道哪一個.
C.S.Lewis就是.
那裡亞故事的作者.
也是其中有本很出名的書.
叫作Scripted Letter.
魔鬼的信的書.

$^{761}$這個很出名的基督教的神學家.
和作者.
他說他有兩句說話.
我覺得很有意思.
他說如果上帝.
應允我年輕的時候.
所有的祈禱.
我現在可能已經毀了自己.
即是破壞了自己一生.
因為原來有時候.
我們年輕的時候覺得.
很隨意往還.
覺得年輕什麼都可以.
但是原來我們如果年輕的時候.
我們神聽了我們所有的祈禱.
很多時候我們會做了很多錯的決定.
因為我們不懂得分辨.
我們覺得這些是好東西.
我們沒有想過原來那些就是石頭.
那些就是蛇.
是傷害我們的一生.
破壞我們的生命.
所以有時候我們去求上帝一些事.
我們經歷去問神一些事的時候.
有時候有些事神不給我們.
你要想一想.
上帝覺得那些不是好事.
就不給我們.
不是好的不適合我們.
就不給我們.
反過來.
有一些我們覺得不好的事.
我們覺得很差的事.
我們覺得很壞的事.
會不會在天賦眼裡面.
其實是對你有幫助的事.
所以C.S. Lewis.
他另外一句說話是說.
上帝會讓我們跌到人生的谷底.
是因為在這個教訓裡面.

$^{801}$在地別的地方是學不到的.
原來我們人生裡面面對一些.
不想面對的挫折或者困難.
其實可能這些困難.
你學了這個功課.
你一生就會得益.
在我們眼裡面.
可能我們覺得不好.
但原來上帝覺得.
這個才是能夠幫助我們生命的好東西.
這個就是裡面說的那種狀況.
所以弟兄姊妹.
今天我們要想一想.
我們能不能從神的眼裡面去看好壞呢?.
第一我經常從自己的眼裡面去想.
我喜歡的就代表好.
有沒有想過.
從上帝眼裡面看的才是好呢?.
第三點.
惡負還好 天賦更好.
在之後第九第十節之後.
接著去到第十一節.
第十一節的前半節.
是回應惡不好的事情.
經文怎麼說呢?.
它說你們雖然不好.
尚且知道拿好東西給兒女.
你留意在第十一節裡面.
說的那個不好.
不是我們一般人所謂的.
都不是很好的事情.
這麼簡單.
其實在這裡原文裡面用的字.
是用邪惡的字.
你知道不好.
但如果邪惡.
是深一點的.
是那個力量.
那個破壞力是大一點的.
耶穌接著說到.

$^{841}$原來你們是邪惡.
原來你們是邪惡的話.
你還會知道將好東西給兒女.
而這個邪惡的字.
正正好也是馬太福音六章十三節裡面.
救我們脫離兇惡.
那個惡者那個邪惡.
但是我們知道這裡說.
原來我們就算這些人是邪惡也好.
兇惡也好.
原來他們都會將好東西.
給他的兒子.
正如我們所說的.
苦毒不吃兒.
但這個父親他關注的是什麼呢?.
他關注的是.
兒子生活的需要.
一些食物的需要.
一些基本的需要.
但我們知道我們人活著.
不是單靠食物.
我們還有靈性的層面.
這個不好的父親.
有沒有關顧到這個人的靈性需要呢?.
我們知道.
我們的天父一定關心我們靈裡面的需要.
所以在第十一節的後半節.
他怎麼說呢?.
何況你們在天上的父豈不更把.
好東西給求他的人.
所以天父一定會將好東西.
給每一個求他的人.
但是說了這麼久.
究竟什麼是好東西呢?.
當你看回路加福音就有答案了.
路加福音十一章十三節.
差不多基本上是一樣的.
但是這裡有不同.
你們雖然不好.
尚且知道好東西給兒女.

$^{881}$何況天父他豈不更把.
聖靈賜給求他的人嗎?.
原來這裡耶穌說的好東西.
是在說聖靈.
如果我們理解這段聖經.
從第七節到這裡.
如果我們用一個比較低層次的角度去想.
我們祈禱神就給我們.
神就給我們好東西.
我們每一次祈禱神就給我們.
我們這樣理解了這段.
我們只會認為我們每一次向神祈禱.
神一定要給我們好東西.
就是和你祈禱有關.
但這段聖經其實不是說祈禱.
而是說更重要的東西.
就是神賜下了聖靈.
你禱告蒙英運.
你在地上健康是好的.
讓你有食物是好的.
有工作是好的.
但你有一個更深的需要.
就是你靈裡面.
屬靈生命的需要.
就需要有神和你同在.
而這段聖經告訴我們.
我們要找更重要的東西.
那個好東西是聖靈.
弟兄姊妹.
你理解這段聖經.
應該要想一個更重要的層次.
耶穌要給我們更重要的東西.
就是那個聖靈.
聖靈能幫助我們分辨是非.
你記得我們之前說到.
第七章第一至五節說什麼.
不要評斷別人.
如何能夠評斷一些事情.
準確地判斷一些事情.
我們需要聖靈.

$^{921}$聖靈幫助我們的判斷.
不單止是今天的一些事實.
一些Fact check的問題.
更加是我們整個人生方向的問題.
今天我們普遍的教會.
或者很多的教會.
或者傳統的教會.
都有個傾向是什麼呢.
我們就避開去說聖靈的事情.
因為我們很多時候一想到聖靈.
就會想到靈恩.
因為我們很多時候.
一個負面印象就是.
說到靈恩就一定是不好的.
其實說到靈恩.
就是說到聖靈的恩賜.
聖靈送了什麼禮物給我們.
是好東西.
不過今天我們很多時候.
就是強調了聖靈送給我們的禮物.
能不能夠講方言.
能不能夠翻方言.
能不能夠有神蹟.
我們很關心能不能夠有神蹟.
有神蹟才有聖靈.
我們誤解了方向.
耶穌告訴我們.
給我們聖靈.
不是要我們行神蹟這麼簡單.
不是要我們能不能夠講方言.
還是翻方言.
這些都是次要的東西.
當我們強調了次要的東西.
就確實不是很好.
但是聖靈本身是很重要的東西.
很重要的.
我們今天教會都要認真去思考.
聖靈在我們生命裡的意義.
而耶穌講得很清楚.
聖靈強調的不是醫病趕鬼.

$^{961}$不是那些東西.
講什麼呢.
耶穌在約翰福音很清楚.
祂叫我們分辨.
透過聖靈怎麼說呢.
十六章第七至十一字約翰福音.
我把真情告訴你們.
我去對你們是有益的.
我若不去保衛師.
就不會到你們這裡來.
我若去就猜他到你們這裡來.
說的是聖靈.
耶穌說祂要為我們死.
然後升天離開我們.
但是耶穌說.
同時間當耶穌離開的時候.
就會讓聖靈來到.
耶穌不能夠和我們在一起的時候.
耶穌讓聖靈來到.
會被猜來的聖靈來到當中.
與我們在一起.
但是耶穌告訴我們.
聖靈是做什麼的呢.
特別在這個段落.
祂很強調的.
不是說什麼屬靈的恩賜那些問題.
說的是三樣東西.
他說他來的時候.
就是說聖靈保衛師來到的時候.
有三樣東西很重要的.
第一就是畏罪.
第二就是畏義.
第三就是審判.
指證世人.
今天原來耶穌賜給我們好東西.
就是讓我們在這個世代裡.
我們能夠對罪對義對審判.
有更加敏銳.
有更加多的提醒.
第九節說.

$^{1001}$畏罪是因為他們不信我.
特別是我們所犯過的罪.
聖靈幫助我們能夠知罪.
讓我們知道我們犯了什麼罪.
我們怎樣能夠在神的面前.
知道求神赦免我們的罪.
聖靈能夠幫我們處理我們的過去.
我們的過去的生命.
可能都是一塌糊塗.
但是聖靈幫助我們走出這個坎裡.
讓我們知罪.
讓我們知道回到神的面前.
得到無赦的罪.
處理我們的過去.
聖靈也不只是處理了過去.
不理我們今天的我們.
他為什麼是好東西.
因為他也幫我們行義.
甚至說畏義是因我道父乃女去貪.
你們將不再見到我.
當我們不能夠見到他的時候.
我們怎樣能夠行出義的生活呢.
我們怎樣能夠活出義呢.
我們需要聖靈的幫助.
需要聖靈不斷對我們的提醒.
以致我們能夠知道.
怎樣在地上過上帝喜悅的生活.
這種義是上帝要我們怎樣去.
經歷了神的改變.
怎樣去活出今生呢.
需要聖靈很多的提醒.
為審判.
審判是將來的.
末日裡面我們會面對審判.
但唯有聖靈讓我們知道.
哪一個是最終的審判官.
哪一個帶來最後的審判.
而我們怎樣去靠著聖靈的幫助.
我們能夠迎向審判.
每一個人都犯罪.

$^{1041}$每一個人都面對審判.
但基督徒的審判不一樣.
因為我們有神的憐憫寬恕.
我們被神捍衛無罪.
所以我們面對審判的時候.
神都說一筆勾銷.
所以他說.
聖靈幫助我們面對過去的罪.
今日的義和將來的審判.
令我們在屬靈裡面.
得到最大的益處.
所以今天的這段聖經.
我的體會就是.
它不是說祈禱.
它只是比祈禱更大的事情.
就是我們生命裡面怎樣去尋找神.
扣門去認識神.
以致我們能夠在裡面.
去找到祂自己.
找到神才是最重要的.
今日神討個夢應找.
你今日記得神幫助你.
可能你明天就不記得了.
人是很善忘的.
但你捉緊神.
有神和你一起很重要.
所以聖靈能夠引導我們的生命.
問題是我們有沒有去相信.
聖靈對你一生的引領.
以致常常緊貼著.
神給我們的聖靈.
耶穌在登山部分裡面.
很想我們去捉緊神自己.
在上個月.
2月11日.
不知道大家有沒有記得這單新聞.
在佛羅里達州發生的新聞.
在Davy Exit 7附近.
發生一種嚴重的汽車事故.
當中說到那輛混油車.

$^{1081}$撞了前面那輛車.
以致前面那輛車.
整輛混油車反轉了.
而後面的車也收掣不及.
然後再撞上去.
以致這條叫綠線行車的路.
全面封閉.
而在地面裡面有漏油的現象.
所以由早上就開始.
道路封閉交通阻塞.
在星期三的日子.
我們有加勒團的日子.
我從家裡來這裡是沒有問題的.
因為不是那條線.
所以很順利的.
我早上還看到對面塞車.
我還想著.
我們團契完結應該可以了.
誰不知到我們團契完結.
我們離開的時候.
還是沒有清理好所有的油壺.
仍然是綠線全封.
在這個時候我們團契完結.
其實我接我的兒子.
接我的兒子放學的時候.
就要去到.
Flamingo Road 再下一條街.
我忘記怎麼說.
Exit 1.
就是那個出口.
就是那條路裡面.
就是當時就是.
不知道怎麼辦.
因為兒子在等我放學.
甚至乎女兒在等我放學.
而前面基本上就是封路.
封路之後怎麼樣呢.
其實我還有很大段距離.
都沒有去到車子反轉的位置.
我還在這裡叫做595的位置.

$^{1121}$我在595的時候.
還沒有去到95這裡.
車子反轉的位置.
但已經是很嚴重的塞車.
車子能不能走路呢.
很慢很慢地走路.
為什麼呢.
因為綠線要變為一線.
其中去到某一個位置.
唯一一個出口.
就是六條線歸到一條線.
在另一個出口出去.
才可以疏散到車子.
大家知道我吧.
我不懂Ford的路.
我只會走那條路.
結果怎麼樣呢.
唯一的方法是怎麼樣呢.
唯一的方法是什麼呢.
開Google map.
而Google map呢.
我想你都沒有試過.
Google map去到那裡.
教我這樣這樣這樣.
才能走出這條公路.
我不知道你明白不明白.
怎麼可以這樣這樣.
我心想.
不是直去就可以了嗎.
它教我繞來繞去.
甚至繞到一個方向.
就是我要從西面去東面走.
你明白我的意思嗎.
我現在要從東面向西面回家.
但它教我繞來繞去之後.
有一段路是要我從西面走到東面.
接著再有一個急的U位.
再走才能走出去.
當時我看到這條路線.
我就心想.

$^{1161}$我好害怕.
我真的不懂得走.
沒走過.
但唯有怎麼樣呢.
唯有相信吧.
信者才能夠得.
接著走.
接著車子慢慢走.
接著一路走.
一路慢慢繞來繞去.
嘩 驚險.
接著轉來轉去.
轉了很久.
終於走回84.
走回一些本地路.
終於鬆了.
嘩 趕得及.
不是趕得及.
只是遲一點.
遲十分鐘左右.
就接到我兒子放學.
這個經歷裡面.
告訴我.
有很多時候.
我習慣了靠自己.
這樣走這條路.
好像很理所當然.
但原來當有事件.
發生的時候.
原來我是不知所措.
不知怎麼走.
不知怎麼去我的目的地.
不知怎麼去到.
我要去的地方.
我需要Google Map.
在我們人生裡面.
原來有時候.
我們都出了意外.
看節目.
我們人生裡面.

$^{1201}$我們都有時候.
在信仰上面.
或者在人生的.
各樣的經濟.
生活的困難裡面.
我們都會炒車.
我們炒車的時候.
你還以為用自己的方法.
就解決得到嗎.
有時候是不行的.
你要馴服聖靈的帶領.
聽聖靈的聲音.
聖靈就是我們.
最好的Google Map.
或許當它教你走的時候.
好像和你所走的方向.
不一樣.
明明你要去向西.
但它似乎指向你向東.
但你又怎會不明白.
不清楚知道.
在聖靈裡面.
柳暗花明就會怎樣.
又一串.
原來那條路.
是要這樣走的.
唯有你馴服.
聖靈的帶領.
你才走得到.
而耶穌很想我們去知道.
我們生命裡面.
我們都要去找到.
那個好東西.
耶穌告訴我們.
我們天父關注的.
不只是關注物質.
地上的父親.
還會關注一些物質.
但天上的父親.
更加關注.

$^{1241}$我們靈裡的需要.
而誰能滿足.
我們靈裡內心的需要.
就是聖靈.
所以這段聖經.
告訴我們.
我們要找好東西.
不只是禱告蒙英找.
可以是禱告蒙英找.
我們會經歷.
但在過程中.
更重要的是.
我們要找到.
聖靈就是我們的好東西.
我們要馴服.
聖靈的帶領.
而聖靈比餅和魚.
會更好.
而我們要不停地去尋找.
聖靈對我們一生的指引.
今天我們.
習慣了很多.
用我們的科技.
用我們的知識.
用我們的能力.
去處理問題.
但如果你信神的話.
你要知道.
神真正能夠幫你.
處理你人生生命問題的.
真是最好的導師和師傅.
今天的功課.
是什麼呢?.
試試在這個星期裡.
每一天.
當你做一些.
對你來說.
比較重要的決定的時候.
先停下三十秒.
心裡默默地說.

$^{1281}$聖靈啊.
這件事你怎樣看?.
聖靈啊.
這件事你怎樣看?.
聽一下聖靈.
會不會給你一些答覆.
可能聖靈告訴你.
OK.
你怎樣都行.
你自己決定.
但可能聖靈告訴你.
不是喔.
小心點喔.
或者聖靈告訴你.
STOP.
Don't do that.
讓聖靈作主.
先看了Google Map.
才行.
這就是我們今個星期.
要做的功課.
難不難?.
當然我明白.
有些人來說.
家裡沒有.
通常是什麼呢?.
小事就老婆決定.
大事就老公決定.
而家裡是沒有大事的.
所以家裡什麼.
都是老婆決定的.
這個情況就不同了.
這個情況.
不用問你老婆.
不用問你另一半.
總之你覺得有些事.
停一停.
等一等.
跟聖靈說一句.
聖靈告訴我.

$^{1321}$你有沒有什麼心意?.
無論你決定上班.
請人.
轉工.
買車.
買樓.
還是投資.
還是搭飛機去杜拜.
這些事.
你們也祈禱想一想.
求聖靈的聲音幫助你.
我們一起祈禱.
天父幫助我們.
因為我們知道.
我們有時很忽略.
聖靈對我們的提醒.
求主讓我們.
今天這段經文.
幫我們重新思考.
我們找的.
我們要尋找的.
就要尋找最好.
聖靈就是最好的.
因為他們是我們.
隨時的幫助.
隨時的依靠.
也是唯有聖靈能夠提醒我們.
怎樣為罪為義.
為審判去警惕.
幫助我們.
我們能夠靠神.
去渡過我們的今生.
多謝你聽我們祈禱.
奉耶穌的名字而求.
阿們.
好好去運用社會的資源.
包括了.
我們可以讀多一點.
屬靈的書籍.
了解很多屬靈的偉人的掙扎.

$^{1361}$成長.
以致我們更加明白.
我們怎樣去.
藉著他們的見證.
去激勵我們.
過一個屬神的生活.
同時間我們也真的要.
在這個世代裡.
真的要用多點AI.
AI的作用不是我們要被控制.
而是用一些AI.
幫我們Fact Check.
檢視多點事實.
以致我們能夠.
多點資訊去分辨是非.
幫我們搞清楚我們的立場.
搞清楚究竟.
那件事是對與錯.
還是說只是黨派之爭.
還是說是顏色之爭.
還是什麼呢.
求主幫助我們.
最後一個的總結.
就是評斷是為了定罪.
我們有時很多時候.
評斷一個人.
是為了想告訴別人.
你是有罪.
但我們要學習分辨.
就是為了要愛這個人.
幫助這個人能夠成長.
亦都幫助自己.
在愛裡有根有記.
評斷是看重人的過去.
看人家過去有多不對.
但我們今天要懂得去分辨.
就是讓我們建立好的未來.
一個贖神喜悅的未來.
基督徒過去的日子.
我們確實是一個罪人.

$^{1401}$但今天我們認識了主.
我們應該要扭轉過來.
我們給自己努力.
去成為一個上帝喜悅的人.
求主幫助我們.
在這個世代裡.
我們要懂得分辨.
亦都懂得以愛彼此扶持.
而不是讓我們的家.
或者我們的神家.
以論斷作為我們的原則.
求主幫助我們.
我們一起祈禱.
天父我們求你幫助我們.
主在登山裡特別提及這件事.
正好你知道.
今天我們作為門徒.
在這個世代裡.
都會同樣面對著.
自己很想去評斷別人.
自己很想做審判官的那種想法.
求主幫助我們.
讓我們知道.
我們最終要面對的是你的審判.
所以我們今天更加需要警惕.
謹慎我們的言行.
以致我們不需要被主.
同樣用這樣的方式去審判我們.
更加求你給我們每一個.
都能夠有這個智慧.
去清晰我們的思考.
亦都能夠幫助別人.
在神的眼界裡.
能夠懂得去分辨是非.
分辨很大.
求你幫助我們.
聽我們祈禱.
奉耶穌名字而求.
阿們.
\newpage



\section{創世記 39:21}
\label{sec:031811}
\textbf{聖經人聖經事:耶和華在監房與約瑟同在}~~~~黃克勤牧師博士
\newline
\newline
連結: \href{http://www.church.com.hk/acms/content.asp?site=cdc\&op=show\&type=product\&code=031811\&layout=sermon}{\texttt{031811}} ~~~~ 語音日期: 2026-03-13
\newline
\newline
\hyperref[sec:031810]{\small{< < < PREV SERMON < < <}}
~
\hyperref[sec:index_chronic]{\small{[返順時目]}}
~
\hyperref[sec:031812]{\small{> > > NEXT SERMON > > >}}
\newline
\newline
創世記 39:21
\newline
\begin{longtable}{cl}
\hline
\hline
章節 & 經文 (和合本修訂版)\\
\hline
39:21 & \begin{tabularx}{0.7\textwidth}{X} 但耶和華與約瑟同在，向他施恩，使他在監獄長的眼前蒙恩。 \end{tabularx} \\ \\
[1ex]
\hline
\hline
\end{longtable}
$^{1}$這次我們來看《耶和華神在監房與弱室同在,仍然是弱室的故事》.
《弱室的故事》很長的.
要講很久了.
聖經人聖經事新穎說書的網絡為大家提供的.
希望我們當中人就是弱室.
事就是耶和華在監房仍然與弱室同在.
神的手是從來沒有離開過的.
十分之感恩.
好了,繼續下去.
「弱室的主人波提弗聽見他妻子對他所說的話」.
說什麼呢?.
「你的僕人如此如此待我」.
聽到他太太這麼說.
他就生氣.
「把弱室下在監裡,就是王的囚犯被囚的地方」.
「於是弱室在那裡坐監」.
真的很有趣.
這句字原來一百年前也用「坐監」的中文.
「弱室的主人波提弗聽見他妻子對他所說的話」.
留意一下.
主人有名叫波提弗.
但他太太上帝並沒有提起他.
連名字也沒有,到現在也不知道.
「他太太對他所說的話」.
「你的僕人如此如此對我」.
生氣.
這個「生氣」.
第一點要說的就是.
有個人生氣而令弱室坐監.
「生」和「氣」是兩個字.
「生」是kara.
「kara」是引發燃氣.
例如漏油漏氣.
是引發的.
例如砲槍有條心,引發爆炸.
「生」是引發,燃起了氣.
「生氣」是Anger.
但「氣」這個字「af」可以讀作「af」.
是「鼻」的意思.
是引發了鼻.

$^{41}$很生氣.
聖經中很多時會用「鼻」這個字.
是表示怒氣.
他向我鼻子或鼻的氣.
是表示生氣.
引發了生氣.
就將弱室下在監裡.
這個「生氣」導致弱室坐監.
把他押入大牢.
就是他的主人.
這個大牢是甚麼呢?.
是埃及王法老囚禁一些他想囚禁的犯人.
當時這樣可以嗎?不用審嗎?.
皇帝說是審的.
把他押入大牢.
慢慢罰下來.
慢慢處置.
大牢在哪裡呢?.
原來皇帝要把人押入大牢.
大牢的監房在波提弗家.
因為他是護衛長.
這個人交給你看守.
把他押入監房,因為監房在他家.
所以監房的弱室是熄滅的.
而主母也知道監房在哪裡.
不過豈有此理.
我看你這麼帥.
三分顏色上大雄.
叫你也不肯聽.
好!老羞成怒.
自己很沒面子.
老羞成怒過不了自己.
於是害了他.
於是弱室就在那裡坐牢.
不過要留意一件事.
上帝的手有沒有在後面工作呢?.
是有的.
如果你是他主人波提弗.
聽了他太太說.
又信他太太說,你會怎樣?.

$^{81}$你會怎樣做?.
拔劍?拔刀?.
安琪比他用刀,拔刀.
拉住他.
跟他出去花園時,我劈他的頭下來.
或者就在這裡劈他的頭下來.
嘩!調戲我太太.
我這麼信任你,讓你做管家.
除了我吃的飯,什麼銀兩都給你.
整個家都給你.
竟然你欺負我太太.
馬上劈.
如果這個波提弗在皇帝面前.
這麼有地位,劈死一個奴隸有沒有問題呢?.
當然沒有問題.
劈死一個奴隸.
劈死一個埃及公民有問題.
劈死一個奴隸沒有問題.
就算在當日這樣的社會.
劈死一個埃及公民.
都有辦法遮到.
如果在皇帝的護衛長來說.
沒有問題.
只要執個罪名給他就可以.
馬上可以劈死他.
說他造反都可以.
說他收集情報想造反都可以.
碎屍萬段都可以.
割他肉醬都可以.
不過波提弗沒有.
只是將他下在監獄.
生氣.
為什麼會這樣.
神的手在工作.
亦都虐失自己賺取到他的生命延續.
為什麼呢?.
其實如果看完這兩節.
你會發覺他的主人波提弗.
他的善色.
是波提弗不相信他太太.

$^{121}$反而波提弗是相信虐失.
看完之後.
虐失不是這樣的人.
怎麼會呢?.
照我所知他不是這樣的人.
他屬於什麼都不偷吃.
我留意到他.
他應該不是.
扔他入監.
如果不是當場已經斬死.
頭都斬了幾段.
他要坐監.
亦都要.
交少少氣.
那麼多工人看著.
那麼多親戚朋友.
那麼多人在.
亦都是太太看著.
他太太告他.
起碼要交多少氣.
引起被子生氣.
憤氣很生氣.
拉他入監.
拔刀出來.
劈死你.
等候法律.
豈有此理.
起碼要交少少貨.
交少少氣.
我認為波提弗是這樣.
如果不是.
他已經一刀劈死虐失.
照胸口砍下去.
反而坐監.
怎麼會呢?.
一個奴隸.
而且是正規國座.
所以我認為.
波提弗是信虐失.
沒有事的.

$^{161}$下次不要了.
這樣太輕了.
拉他入監.
第二點.
但耶和華與虐失同在.
請你記得.
我曾經講過.
耶和華的同在.
是最為緊要.
最為需要.
無論去到高高低低.
走過死淫休工.
都不怕.
耶和華是我的牧者.
我必不自缺.
牧者不需要什麼.
有牠已經足夠.
牠領我到青草地.
領我到溪水旁.
我雖然行過死淫休工.
也不怕遭害.
因為你與我.
同在.
死淫休工固然可怕.
青草地.
溪水旁固然好.
但最緊要的是.
你與我同在.
就沒有什麼特別需要.
有你同在就最好了.
行過死淫休工.
也不怕遭害.
因為你與我同在.
在監獄與虐失同在.
向他施恩.
很有趣.
令他又再施辱.
在監獄長的面前無恩.
監獄長.
中間所有囚犯.

$^{201}$都交給虐失的手下.
只是奴隸.
那個把全家.
交給他做管家.
這個爹的層次.
皇家監獄長.
虐失這麼厲害.
整整齊齊.
很好.
向他施恩施義.
監獄長竟然說.
我回家休息.
所有監犯的事交給你.
都在虐失手下.
他們在哪裡.
所辦的事.
都是經他手的.
整個皇家監獄.
他搞定了.
犯罪虐失手下的事.
施辱監獄長一概不察.
又是這樣.
除了自己吃的飯出的糧.
什麼都不知道.
因為耶和華與虐失同在.
有耶和華的同在.
耶和華與虐失同在.
是講了兩次.
耶和華與虐失同在.
然後呢.
耶和華與虐失同在.
令到他所做的.
盡都順利.
早兩天才講了.
當神同在.
他會讓你盡都順利.
有時候同在.
未必盡都順利.
主耶穌與他同在嗎?.
有,但主耶穌做事順不順利?.

$^{241}$毫不順利.
要記住.
不是神的心意.
就一定不順利.
有時候是神的心意也不順利.
因為祂要你不順利.
最重要的是.
無論順利與不順利.
主耶和華與你同在.
要記住這件事.
盡都順利.
原來在監獄裡.
在監獄裡.
防守嚴密.
但祂的手在外.
幫不了我.
不,神的手可以在監獄裡.
幫助虐失.
神在監獄中與虐失同在.
就是今天的題目.
在監房與虐失同在.
最重要的字是同在.
因此.
請你特別要記住.
主啊,如果我要祈禱.
神啊.
請你賜我一件東西.
六個號碼.
很多人有六個號碼後.
家散人亡.
很多人有六個號碼.
加上特別號碼後傾家蕩產.
本來好好的一家人.
就散了.
最重要的是.
主啊,如果我要求一件事.
我要求你的.
同在.
這是最重要的.
什麼都假.

$^{281}$有主的同在在監房裡.
都可以,都好.
所以記住,耶和華的同在.
不會因為監房.
來拘禁到人.
有信心.
一起禱告.
求主幫助我們.
雖然有時候.
在一個很困難的情況.
好像在監房裡.
這樣不行,這樣不准.
又缺乏,又有限制.
人又髒.
各樣的遭遇都不好.
又不安.
主啊,不過.
最重要的是有你同在.
心裡面就很平安.
有你同在,事情就會順利.
有你同在.
就沒有偶然發生.
沒有突發.
沒有一件事是主.
為什麼會這樣發生.
我是耶和華神,我也不知道.
你知道所有事.
你管理所有事.
你掌控所有事.
求你在我們的難處當中.
與我們同在.
如果要向你求一件事.
就是求你與我同在.
上天祈禱.
與我們同在,奉耶穌基督命求.
阿們.
\newpage



\section{歌羅西書 3:5}
\label{sec:031812}
\textbf{貪婪,後果}~~~~林天行牧師
\newline
\newline
連結: \href{http://www.church.com.hk/acms/content.asp?site=cdc\&op=show\&type=product\&code=031812\&layout=sermon}{\texttt{031812}} ~~~~ 語音日期: 2026-03-15
\newline
\newline
\hyperref[sec:031811]{\small{< < < PREV SERMON < < <}}
~
\hyperref[sec:index_chronic]{\small{[返順時目]}}
~
\hyperref[sec:031813]{\small{> > > NEXT SERMON > > >}}
\newline
\newline
歌羅西書 3:5
\newline
\begin{longtable}{cl}
\hline
\hline
章節 & 經文 (和合本修訂版)\\
\hline
3:5 & \begin{tabularx}{0.7\textwidth}{X} 所以，要治死你們在地上的肢體；就如淫亂、污穢、邪情、惡慾和貪婪—貪婪就是拜偶像。 \end{tabularx} \\ \\
[1ex]
\hline
\hline
\end{longtable}
$^{1}$大家早晨.
在講道之前請大家一齊擁抱聖的話語.
因為這是毫得無比的.
讓我向祂禱告.
因為藉著你的話.
我們能夠對你有更多的認識.
因為就是我們身邊的弟兄姊妹.
互相扶持,彼此去建立的時候.
我們能夠成為一個更加合乎你心意的器皿.
能夠成為一個更加忠心的群體.
去履行你給我們的使命.
我求你開我們的眼睛.
藉著你的話語.
讓我們看見你的心意.
讓我們看見這個世界.
好叫我們懂得.
怎樣去忠心地去履行你給我們.
各人的召命.
讓我們在自己的生活.
也讓我們在信仰的群體生活.
一同去為你做美麗的見證.
打美好的仗.
為的就是為著這個榮耀.
盼望我們可以預備自己.
好叫我們見主面的時候.
都能夠無愧於你.
我們將以下時間.
恭敬,交賞.
求你親自向我們說話.
觸動我們的心靈.
改變我們的生命.
聽我們禱告.
奉靠耶穌基督.
得勝的名字祈求.
阿們.
有時做錯事情.
沒有即時見到後果.
反而是一個更大的危機.
我在美國讀神學.
在讀神學期間.

$^{41}$認識了一位在美國長大的中國人.
所以他的外形.
他的身體,他的Genes,他的DNA.
全部都是中國人.
但是他的生活,他的習慣.
就像一個美國人.
他曾經在初職的時候很忙.
工作時間很長.
所以吃東西很不定時.
很晚才吃東西.
很晚才吃東西的時候.
怎樣吃呢?.
開車買個快餐就吃.
經常吃漢堡包,薯條,喝汽水.
這成為了他日常生活的一部分.
因為很多美國人都是這樣吃.
因為他可能有中國人的基因.
但他又不是特別胖.
都是差不多正常的尺寸.
但是吃著吃著.
就吃很多很多很不健康的食物.
他去到一個位置.
他覺得不太對勁.
他覺得身體不太正常.
他要去做驗身.
他說到驗身的時候.
他的血脂是脹到抽不到血.
驗不到的.
我不明白從醫學角度可不可以是這樣.
我猜他可能是誇張的.
但就是說到即使他沒有很嚴重的肥胖.
但他吃了這麼多的快餐,這麼多的垃圾.
裡面他驗不到的血脂是脹得很厲害.
是即時必須要去作出一些改變.
否則後果不堪設想.
弟兄姊妹.
當我們的生活習慣.
沒有為我們帶來即時的後果的時候.
我們會否不知不覺地繼續這樣生活.
直至我們看到後果.

$^{81}$才發覺大事不妙.
今日我想和大家探討一下.
在我們的生活習慣.
如何影響大地.
以致帶來一些怎樣的後果.
可能沒有即時的後果.
但我們是否看得到影響有多深遠.
我知道今個月是猜拳月份.
是和猜拳有關的.
不過我想給大家看看.
在我們的生活習慣當中.
我們如何影響大地.
以致影響其他在地上的人.
用一個普世性,跨文化的眼光.
去看我們的生活帶來什麼影響.
可能你聽到很多不同的數據.
很多不同的專家.
說如果我們繼續這樣生活的話.
地球的氣溫就會上升很多.
在幾年之後地球就會消失.
不過這些好像很搖亂.
多一度,半度對我們有什麼影響呢?.
地球過多幾千年之後就會消失.
對我們有什麼影響呢?.
後果好像很遠.
我們就變得不太靠近.
但今日我想藉著第一段聖經.
和大家做一個神學反思.
我不會像一般講道那樣.
很深入地鑽研經文.
但我們藉著這個神學反思.
我們想一想.
我們作為基督徒.
我們怎樣去關愛這個受造的世界.
而不是只看環保.
環保和關愛創造是有一個不同的.
我今日藉著去反思的時候.
我們會先看看我們其中一個問題.
就是因為我們的貪婪.
我們的貪婪正在破壞大地.

$^{121}$以致帶來不同的後果.
從而我們再反思一下.
我們可以做些什麼去改變.
讓我們先看第一樣.
貪婪是令到大地受到破壞.
但在講到我們做了什麼不對的事.
帶來什麼影響之前.
讓我們先看看聖經.
講一下什麼是對的.
我們應該怎樣做.
剛才主席帶我們去看了.
歌羅西書第三章第五節.
但我們先看看頭四節.
請聽我讀出歌羅西書第三章第一節.
「所以你們若真與基督一同復活.
就當求上面的事.
那裡有基督坐在神的右邊.
你們要思念上面的事.
不要思念地上的事.
因為你們已經死了.
你們的生命與基督一同藏在神裡面.
基督是我們的生命.
祂顯現的時候.
你們也要與祂一同顯現在榮耀裡」.
第五節講到一個應用.
就是你不要做些什麼.
作出一個應用.
但在應用之前.
保羅是建立出一個神學的動機.
是因著一個神學的原因.
因著我們與基督一同復活.
因著基督會再回來.
以致我們做的事有所不同.
這裡所用的每一個動詞.
都是一個複數的動詞.
換句話說保羅是與整個教會的群體.
我們一同與基督一起復活.
所以我們一同要求上面的事.
所以我們一同要做些不同的事.
當我們作為一個群體.

$^{161}$我們一起活出一個不一樣的見證時.
那是很美.
那是很震撼.
那是很榮耀的.
而他所說的是.
因著我們已經與主一同復活.
因著主會再回來.
所以我們會一同與主顯在榮耀裡面.
所以我們做的事有所不同.
保羅是按照猶太人天啟文學的觀念.
他將時間和空間合併.
去顯出主的在來.
讓我們看見這個在來.
是在神的計劃裡面一定會發生的.
這個空間是什麼?.
是我們已經與主一同復活.
基督已經坐在神的右邊.
但是時間就是將來主會回來.
我們一同顯現.
當時間和空間合併在一起的時候.
我們知道神幹的事必定會成就.
以致我們緊緊抓住這個盼望.
等待耶穌再回來.
思念天上的事.
以致我們在地上所做的事不一樣.
以致我們可以自死在地上的肢體.
就如淫亂,污穢,邪情,惡念和貪婪.
貪婪就和拜偶像一樣.
今天我們要帶出的問題是.
我們應當因著我們相信耶穌死和復活.
而期盼著耶穌再回來.
以致我們所有群體一同思念天上的事.
以致我們在地上活出不一樣的生命.
以致我們不會被貪婪主導我們.
但是實際上就是因著我們的貪婪.
我們正在破壞大地.
貪婪是其中一個原因.
我們傷害大地.
為什麼是貪婪呢?.
是因為我們想賺到盡.

$^{201}$是因為我們想用最少的付出換來最大的享樂.
是因為我們想享受到盡的時候.
去滿足我們自己.
這一份貪婪.
我們不計後果.
我們罔顧我們的行為.
對大地如何破壞.
我們先從海洋說起.
我們的貪婪如何影響海洋.
首先我們知道海洋是很重要的.
我們的地球大概有80\%是海洋.
20\%是陸地.
而海洋不單是穩定我們的氣溫.
亦是穩定我們在地球上的二氧化碳和氧氣的比例.
亦在海洋中滋生很多生物和物種.
對於整個生態是非常非常重要的.
如果我們弄壞了海洋.
基本上是在破壞大地.
是我們和地球同歸於盡.
講得足一點.
但我們知不知道我們今天.
無論是香港人或其他南亞甚至歐美.
我們吃多少海產.
以致我們如何破壞海洋.
現在的捕魚不會用魚竿來釣魚.
亦不像彼得時用網綁住魚竿然後拉上來.
是用大船和大網抽魚竿來捕捉.
捕捉得多少就捕捉多少.
所以目標雖然是為了捕捉大量的魚.
但一定會有捕捉錯誤.
往往會勾到很多海龜,海獅,鯊魚,海豚,珊瑚.
這些不是捕魚人士想捕捉的東西.
只是捕捉錯誤而已.
他們不會理會捕捉多少就捕捉多少.
但說放生太多是不行的.
因為捕捉得上來時已經死了.
即是在海中一直抽著抽著抽著.
抽一段時間再拉上來.
發現有隻海豚.
你如何救它?到時可能已經死了.

$^{241}$所以首先在捕魚中.
會經常誤中附居,殺錯良民.
在捕魚時有很多不同的海洋生物被捕.
甚至死亡.
但我們主要是吃魚.
所以我們大量吃我們喜歡吃的魚.
但這很影響海洋的生態平衡.
簡單來說.
海洋中可以分開有大魚,中魚,細魚和異母.
很容易理解.
大魚吃中魚,中魚吃細魚,細魚吃異母.
我們吃得最多的是甚麼魚?.
我們吃得最多的是中魚.
大魚,鯊魚我們不是怎樣吃的.
當然有人吃刺.
但我們吃得最多的是大魚.
例如吞拿魚,星斑,老鼠斑都算是大魚.
我們吃得最多的是大魚.
所以你可以想像.
主要是捕魚.
吃著吃著.
海洋中的大魚就少了.
中魚少了.
少了中魚.
大魚吃中魚.
不夠中魚給大魚吃.
大魚就會少一點.
大魚又少了.
中魚又少了.
細魚就多了.
因為沒有人吃它.
多了細魚.
就吃光了蝦母和魚母.
又不夠魚母.
不夠魚母怎麼辦?.
細魚都沒有了.
其實很簡單.
我們破壞了生態的循環.
因為我們吃了中間那一層.
不單止海洋生物.

$^{281}$生物被我們吃得多的時候影響了.
在捕魚的過程裡.
都是很影響海洋生態.
一般來說.
可以想像一艘船.
放一個網在海底撈撈撈.
把深海魚撈回來.
另一幅圖你可以想像到.
在這樣捕魚的過程裡.
它是掛著海底.
一路刷過去掃魚.
可否想像對河床有多大的影響?.
有多少不同的破壞?.
所以在海洋的生物.
在河床裡.
都被我們搞得翻天覆地.
不單止.
還剩下大量的漁網.
漁網多了破爛了.
它就棄置在海洋裡.
幸好可以幫你撿回來.
你知不知道.
根據不同的調查.
海洋塑膠的垃圾裡.
有46\%是漁網.
覺得很沮喪嗎?.
有沒有被騙的感覺?.
當你很努力地說.
慳著這些膠飲管.
慳著這些膠叉膠巾.
不要丟到海裡.
原來你發覺.
有差不多一半原來是漁網.
但沒有人跟我們說過.
我們如何去適用海產.
帶來多少不同的後遺影響海洋?.
除了海,陸地.
有沒有發現.
早幾年經常聽到.
在巴西亞馬遜森林裡.

$^{321}$很多的山火.
無論如何都救不息.
其實我覺得是沒有嘗試去救息.
有時候是刻意去燒森林.
有時候是刻意扮作.
不經意燒了森林.
為什麼呢?.
因為燒了森林後.
你會有更多的空間.
然後可以用這片陸地.
去放牛.
因為養牛.
你賣到牛扒就賺到很多錢.
你不夠地.
我就燒了森林.
以致我有更多地.
以致我可以養牛.
但你明不明白.
一個熱帶森林.
對於整個生態來說.
有多重要?.
一個熱帶森林裡面.
有多少不同的物種.
你破壞了它.
你沒有了那些樹.
又是對於氣候有多少影響.
但因為這樣賺到錢.
因為你養到牛.
你賣到牛.
賺到錢.
這個都不是他們在乎的事.
換句話說.
都是因為我們心底裡的貪婪.
錢要賺到盡.
享受要享到盡.
以致我們罔顧大地的生態安全.
弟兄姊妹.
你有沒有想過.
你想用最低的價錢.
去買到你想要的東西.

$^{361}$但賣東西的人.
又想用最高的價錢.
賺到最大的利益.
那中間誰要去受苦?.
就是很多不同的人和大地.
我想說.
我們去看這個大地的時候.
我們有沒有想到.
這個大地其實不是屬於我們.
是屬於上帝.
我們是管家.
我們去保養這個大地.
如果不是我們的地.
但我們當作自己用.
以致不斷滿足自己的貪婪.
是失卻了神給我們的使命.
我在美國的時候.
我很早拿到車牌.
但我很久都沒有車.
所以我經常開別人的車.
我都開得習慣了.
譬如有一次.
我哥哥要去機場.
我就跟他車去機場.
然後他上飛機.
我就開他的車離開.
經常會這樣去做.
但我經常開別人的車.
譬如我哥哥的車.
我不會太害怕.
因為那輛是十幾年的日本車.
他也撞過了.
我自己再開的時候.
我都不太怕被人欺負.
所以就照樣開.
但話說我讀神學的時候.
當時我參與在清善寮時工.
有一年清善寮去短孫.
我們嘗試去幫他們籌款.
在外國比較常見.

$^{401}$就是我們免費洗車.
隨意按感動的自由奉獻.
這樣籌款給清善寮去短孫.
有些人泊車去教會.
我就幫他們駕駛去後巷.
幫他們洗乾淨.
有些車我駕駛過去.
有些清善寮有車牌.
我就幫忙駕駛去後巷.
有一輛車經過的時候.
導師跟我說.
給你差錢.
這輛是你駕駛的.
不要給年輕人碰.
那輛是一輛很名貴的房車.
很美的Jaguar.
我開了引擎.
我以為沒有開.
因為靜得我聽不到聲音.
駕駛那輛車我就很小心.
千萬不要害怕.
小心點.
不要走那麼近過來.
等我駕駛停車.
等一下.
好 小心點洗.
因為這輛好車不是我的.
我不可以害怕它.
弟兄姊妹 你有沒有看到這個地球.
是神好美麗的創造.
並不屬於我們.
我們應該好好去愛護.
保護.
顧惜這個大地.
而不是任我們的貪婪.
去滿足我們自己.
剛才說到貪婪.
令到這個大地受到破壞.
但同樣我們知道.
貪婪的後果是由別人去面對.

$^{441}$這才是問題.
不單止我們影響這個大地.
我們沒有見到即時的問題.
而且面對的困難是由別人去承擔.
這就是我們要用普世的眼光.
用跨文化的意識去留意.
當這個大地受到破壞的時候.
受苦不是你和我.
而是遠處的人.
剛才說到海產的問題.
為什麼東西可以這麼便宜.
為什麼可以被人賺到盡.
想用最低的價錢.
賣出給享受的人.
可以盡情的享受.
但又想放出去的人.
可以賺到最多.
中間的就是被剝削的人.
這是另外一個很大的問題.
是人口販賣的問題.
今天不能細說.
但你知不知道.
有很多人真的被騙了上船去工作.
舉個例子.
可能在一些窮困的國家.
你每個星期可以賺到一元美金.
接著有人跟你說.
喂,有一份筍工.
我帶你出海工作.
那裡可以讓你一個星期賺到二十元美金.
哇,很好啊.
那我就去吧.
誰知被人上了船.
上了船後.
沒有薪水.
不能下船.
不斷在很惡劣的環境.
全日工作.
一做就做幾個月.
在很惡劣的環境.

$^{481}$強行開張床.
就在那裡睡覺.
很少糧食.
很長的工資.
很長的工時.
而是沒有工資.
被人扣在船上.
你說不能反抗嗎?.
真的不能反抗.
為什麼?.
因為他真的可以一槍打死你.
扔你進海.
接著說有意外.
他不需要負責任.
有很多在泰國的漁船.
就是這樣販賣人口.
去賺取暴利.
而這些人是無人問津的.
換句話說.
我們在吃的海產.
背後可能有很多有血有肉.
別人流出來的血肉.
而是給我們享受的.
我這樣說不是說.
你以後不要吃海產.
但我想問.
你有沒有想過.
後果原來是別人承擔著.
另外.
大家聽過在索馬里亞海邊.
很多海盜.
我絕對不是贊同他們的做法.
他們是不應該去打劫的.
但這些海盜是怎麼來的?.
為什麼會有這麼多海盜?.
索馬里亞那一帶.
其實都是一些窮的國家.
不是很多謀生的方法.
但有一班人本身是捕魚的.
但他捕魚.

$^{521}$就真的在漁船上.
拿魚竿,魚網去捕.
但他捕不到.
為什麼?.
因為一架架大漁船.
一撈就撈了所有魚.
他沒有魚打.
所有魚都被人捕了.
你可以想像.
大財團底下的小市民.
是找不到入息的.
他們沒有辦法生活.
他們才被迫去做海盜的搶劫.
我再強調.
我不贊同他們搶劫.
但我們能否看見.
我們如何逼到他們沒有罰紙謀生.
都是因為我們的貪婪.
我們要賺到盡.
他們找不到.
他們會用其他方法.
另外我想和大家思想一下.
氣候變化對於邊緣化.
被剝削的,貧窮的.
或者是小資源的人.
是一個重大的打擊.
換句話說.
氣候變化是一個危機倍增者.
是令到本身存在的危機倍增出來.
例如糧食的問題.
貧窮的問題.
污染的問題.
戰爭帶來難民的問題.
這些都會因著氣候變化.
令到這些問題倍增.
為什麼我這樣說呢?.
這裡是一個萬家拉小田.
我看過一份報告.
說到萬家拉有一個人.
他勉勉強強可以找到一塊大地.

$^{561}$有足夠的溪水來灌溉.
以至種植一些基本的農產.
他可以這樣生存.
但因著氣候變化的緣故.
雨水驟來驟去.
他很不穩定.
有時他的溪水乾了.
他的田就沒有了.
他就種植不了.
你知道沒有水就出產不了.
出產不了之餘.
也會影響泥土的質素.
以至養分會下降.
再後來有水的時候.
出產都會少了.
當他本身是餐穩餐食餐餐清的時候.
他地上沒有出產.
他賣不了.
他就沒有東西吃.
他不是說我等幾個月.
希望幾個月之後好一點.
他沒有存貨.
對於貧窮人來說.
他承擔不了氣候變化帶來的影響.
所以這是一個危機倍增者.
我們對他的影響.
是遠遠大過我們能夠想像的.
是遠遠大過我們需要承擔的.
環境破壞對貧窮人的影響是倍增的.
有幾個方面.
第一是天然災害.
你可以想像.
如果颳風下雨.
我們有瓦遮頭就沒事了.
如果他住的房子漏水.
吹倒了的.
養分就大很多了.
第二是適應的代價.
搬遷的困難.
和天然資源的減少.

$^{601}$怎樣才算是適應的代價呢?.
剛才我們有說到.
譬如我說今個月旱災.
下個月才有.
對我們來說是沒所謂的.
為甚麼呢?.
因為香港人沒有人試過.
漏水後沒有水的.
就算他製水.
下面也有供水車供應給你.
所以我們不知道甚麼是沒有水的問題.
但如果有一個月旱災.
他們沒有水.
他們沒有佛招灌溉農作物.
會有甚麼影響呢?.
他們是不能夠適應的.
這個代價太大了.
另外搬遷是困難的.
如果這裡不行.
我搬去別的地方.
那塊田就是那塊田.
他搬不到去別的地方種植.
他沒有這個選擇.
天然資源減少.
剛才有說到.
無論是溪水少了.
他也沒有別的選擇去做.
我們是否明白.
環境的污染.
環境的破壞.
這是一個危機倍增者.
對於貧窮人來說.
影響更多.
以一個很簡單的例子.
天氣熱.
我們可以開冷氣.
我開冷氣.
令到天氣更加熱.
不要緊.
我可以再開大一點的冷氣.

$^{641}$但如果沒有冷氣的人.
他就要承擔了.
看不看得到.
當我們去享受的時候.
不是我們去承擔.
這個大地被破壞的影響.
是由必然去面對的.
所以當我們看得到.
貪婪令到大地受破壞.
但貪婪的後果.
是由必然去面對的時候.
我們可以做些什麼呢?.
我鼓勵弟兄姊妹.
放下自己.
讓我們學習去愛護這個大地.
和愛護別人.
我們想想.
不是每樣東西都要去到盡.
錢不是要賺到盡.
享樂不是要去到盡.
有沒有這東西.
是我們可以放下.
以致我們想想.
我們可以怎樣去保護這個大地.
我們怎樣去愛護其他人.
愛護那些在遠處.
我沒有見過他們的人呢?.
是因著一個神學動機.
是因著我相信.
主耶穌必定會回來.
是因著我相信.
這個大地是屬於神的.
不是屬於我的.
我只是管理它.
是因著我相信.
我要忠於基督.
以致我要去面對主.
以致我要思念天上的事.
而不是地上的事.
所以我們作為一個群體.

$^{681}$我們在地上.
活出不一樣的生命.
不是只是照顧自己的享樂.
不是滿足自己的貪婪.
而是要知道.
我願意去愛護這個大地.
我願意去愛護其他人.
我鼓勵大家.
我們去學習更多的了解.
更多的思考.
更多的禱告.
我不是說.
你有什麼不可以做.
什麼不可以做.
從此不再吃海產.
從此不吃牛扒.
我不是這個意思.
不是這樣一刀切的.
但你要想想.
從我自身的角度.
我可以做些什麼.
以致我了解多些.
怎樣去愛護這個環境.
以致我有多些的神學反思.
去想我應該怎樣去做.
以致我可以有更多的禱告.
去慰藥我自己.
怎樣去善用這個身份.
我慰藥在遠處.
我看不到.
但在困苦流離當中的人.
我怎樣為他們禱告.
這些是我們可以做的.
其中一個.
我不知道大家會不會有興趣.
就是你可以去看多些相關的資料.
在Netflix.
我不知道大家有沒有用.
其中一個就是有一套.
叫做Seaspiracy.

$^{721}$在香港Netflix.
據我了解是還可以看的.
我沒有Netflix.
不過我知道.
剛才我說的部分資料.
是來自這篇紀錄片.
但我會說.
我不是完全說話.
我是自擔全收的.
甚至我覺得有些東西有點偏頗.
所以你去看的時候.
你也要思考.
你也不一定要看這個.
但有很多很多不同的資源.
你可以去了解多些.
以致可以想想.
究竟我可以做些什麼.
怎樣去回應這個大地的需要.
更加重要的是.
除了多了解.
去多思考.
多討苦之外.
我也鼓勵你.
去改變生活習慣.
想想我可不可以減弱一些生活.
我可不可以減少一些能源去用.
坊間有很多不同的方法.
說怎樣環保怎樣去做.
你也可以想想怎樣去應用.
不過我想說.
跟環保有什麼不同.
從環保的角度來看.
已經是覺得.
消費者教育是沒用的.
什麼是消費者教育是沒用的.
如果我要跟消費者說.
你們盡量節省餐具.
不要用即棄餐具.
這樣說是沒用的.
是沒見效的.

$^{761}$是會有些改革.
譬如現在.
你要外賣餐具嗎.
加一元吧.
這樣你才有改變.
人們才會減少用.
所以我只是教你.
怎樣不要這樣去做.
是沒改變的.
所以如果你是一個在位的人.
你有辦法.
可以帶來一些改革的話.
這是一件好事.
不過我說的不是環保.
是生態關懷.
是從信仰的角度出發.
不是從環境保育出發.
所以是因著我們每一個.
相信耶穌的人.
我為著主的緣故.
我為著其他人的緣故.
為著大地的緣故.
我願意改變我的生活習慣.
即使你會覺得.
一籮籮的塑膠魚網在海裡.
我把膠巾蓋上.
把飲管蓋上.
有什麼用呢.
有沒有用是一件事.
我做我應該要做的.
這是從信仰的角度.
是一個神學的動機.
所以這並不是消費者的教育.
而是作為一個神學的教育.
我們作為主的門徒.
我們應該怎樣去做.
怎樣去為著其他人去做這件事.
弟兄姊妹.
大家可能知道.
十一年前.

$^{801}$尼泊爾有一個大地震.
在地震之前.
我準備買機票去尼泊爾.
但是地震了.
我就去不了.
但我一直都很心掛把.
很為著他們著緊.
都為他們禱告.
經歷了兩年之後.
我終於可以去到尼泊爾.
我去到尼泊爾的時候.
在山區裡面住了兩晚.
你可以想像一下.
就是在一個泥屋.
踏出來.
一邊有木屋.
在那裡睡了兩晚.
挺有趣的.
跟在山區的人住.
跟他們聊天.
去了解地震對他們有什麼影響等等.
但是我跟他們說.
兩年前地震.
我本來想來.
我來不了.
我一直為你們祈禱.
今日我終於見到你們.
我很開心.
他們很感動.
他們覺得竟然有些素未謀面的人.
一直會為他們禱告.
弟兄姊妹.
可能有一日.
你見到不知在那裡很遠的人.
可能是一個第三國家的人.
你見到他.
你可以跟他說.
我為了你的緣故.
我改變了我的生活習慣.
希望你的生活會更好.

$^{841}$那個其實都很感動.
可能將來在天上面.
你才會見得到.
但是我為了你.
我改變了我自己的生活習慣.
我減少了我的享樂.
我也曾經為你出過一份力.
這也是一份情義.
讓我們用不同的眼光.
用普世的眼光.
去看這個大地.
讓我們懂得.
怎樣為著主的緣故.
盡責.
盡忠.
做回我們當初的事.
讓我們同低頭.
弟兄姊妹想請你在神面前.
有少少安靜的時間.
你會不會覺得.
在你的生活當中.
有因為你要滿足自己的貪婪.
有些過度的享樂.
過度的消費.
會不會買些你不需要用的東西.
會不會有些你不需要的享受.
有些不同的資源.
你浪費了.
你有沒有看輕了.
作為這個大地管家的責任.
你有沒有忽略了.
在地球的另一端.
很多人欣賞我們的行為.
他們承擔著的後果.
你願不願意用一個普世的眼光.
去看這個世界.
去看神給你的使命.
你願不願意.
因著耶穌已經重死復活.
並且會再來.

$^{881}$以致我們思念天上的事.
在地上面.
我們會出不一樣的生命.
天上就是在你面前.
去尋求你的原諒.
相信或多或少.
我們每一個都有滿足自己貪婪的意思.
忽略了你給我們去管理大地的使命.
忽略了我們關心在地球遠處的人.
他們的需要.
讓我們看見很多人.
在貧窮當中.
在飢餓當中.
在疾病當中.
在戰亂當中.
在流離失守當中.
他們面對著困鎖.
面對著氣候變化.
倍增的問題.
讓我們願意的.
為著他們的緣故.
以禱告的心.
去被你改變.
好讓我們在地上面.
我們活出.
我們所想.
所盼望.
天上的事.
好叫我們的行為.
能夠演出你在我們當中.
能夠將榮耀歸給你.
聽我們禱告.
願主基督得生命之祈求 阿們.
\newpage



\section{約書亞記 9:1-27}
\label{sec:031813}
\textbf{汲取教訓}~~~~陳沛偉牧師
\newline
\newline
連結: \href{http://www.church.com.hk/acms/content.asp?site=cdc\&op=show\&type=product\&code=031813\&layout=sermon}{\texttt{031813}} ~~~~ 語音日期: 2026-03-15
\newline
\newline
\hyperref[sec:031812]{\small{< < < PREV SERMON < < <}}
~
\hyperref[sec:index_chronic]{\small{[返順時目]}}
~
\hyperref[sec:031814]{\small{> > > NEXT SERMON > > >}}
\newline
\newline
約書亞記 9:1-27
\newline
\begin{longtable}{cl}
\hline
\hline
章節 & 經文 (和合本修訂版)\\
\hline
9:1 & \begin{tabularx}{0.7\textwidth}{X} 約旦河西，住山區、低地和沿大海一帶直到黎巴嫩的諸王，就是赫人、亞摩利人、迦南人、比利洗人、希未人、耶布斯人的諸王，聽見這事， \end{tabularx} \\ \\ \relax
9:2 & \begin{tabularx}{0.7\textwidth}{X} 就都聚集，同心合意要與約書亞和以色列人作戰。 \end{tabularx} \\ \\ \relax
9:3 & \begin{tabularx}{0.7\textwidth}{X} 基遍的居民聽見約書亞向耶利哥和艾城所做的事， \end{tabularx} \\ \\ \relax
9:4 & \begin{tabularx}{0.7\textwidth}{X} 就設詭計，假扮使者出去。他們拿舊布袋和破裂補過的舊皮酒袋馱在驢上， \end{tabularx} \\ \\ \relax
9:5 & \begin{tabularx}{0.7\textwidth}{X} 將補過的舊鞋穿在腳上，把舊衣服穿在身上，作食物的餅都又乾又長了霉。 \end{tabularx} \\ \\ \relax
9:6 & \begin{tabularx}{0.7\textwidth}{X} 他們到吉甲營中約書亞那裡，對他和以色列人說：「我們是從遠地來的，現在求你與我們立約。」 \end{tabularx} \\ \\ \relax
9:7 & \begin{tabularx}{0.7\textwidth}{X} 以色列人對希未人說：「或許你是住在我附近的。若是這樣，我怎能和你立約呢？」 \end{tabularx} \\ \\ \relax
9:8 & \begin{tabularx}{0.7\textwidth}{X} 他們對約書亞說：「我們是你的僕人。」約書亞對他們說：「你們是甚麼人？是從哪裡來的？」 \end{tabularx} \\ \\ \relax
9:9 & \begin{tabularx}{0.7\textwidth}{X} 他們對他說：「你的僕人是因耶和華－你神的名從極遠之地來的。我們聽見他的名聲，他在埃及所做的一切， \end{tabularx} \\ \\ \relax
9:10 & \begin{tabularx}{0.7\textwidth}{X} 以及他向約旦河東的兩個亞摩利王，希實本王西宏和在亞斯她錄的巴珊王噩所做的一切。 \end{tabularx} \\ \\ \relax
9:11 & \begin{tabularx}{0.7\textwidth}{X} 我們的長老和我們當地所有的居民對我們說：『你們手裡要帶著路上用的乾糧去迎接以色列人，對他們說：我們是你們的僕人。現在求你們與我們立約。』 \end{tabularx} \\ \\ \relax
9:12 & \begin{tabularx}{0.7\textwidth}{X} 我們出來要往你們這裡來的那日，這從我們家裡帶出來的餅是熱的；看哪，現在這餅又乾又長了霉。 \end{tabularx} \\ \\ \relax
9:13 & \begin{tabularx}{0.7\textwidth}{X} 這些皮酒袋，我們盛酒的時候還是新的；看哪，現在已經破裂了。我們這些衣服和鞋，因為路途非常遙遠，也都穿舊了。」 \end{tabularx} \\ \\ \relax
9:14 & \begin{tabularx}{0.7\textwidth}{X} 以色列人收下他們的一些食物，但是沒有求問耶和華的指示。 \end{tabularx} \\ \\ \relax
9:15 & \begin{tabularx}{0.7\textwidth}{X} 於是約書亞與他們建立和好關係，與他們立約，讓他們存活；會眾的領袖也向他們起誓。 \end{tabularx} \\ \\ \relax
9:16 & \begin{tabularx}{0.7\textwidth}{X} 以色列人與他們立約之後，過了三天才聽說他們是近鄰，住在附近。 \end{tabularx} \\ \\ \relax
9:17 & \begin{tabularx}{0.7\textwidth}{X} 以色列人起行，第三天就到了他們的城鎮，他們的城鎮是基遍、基非拉、比錄和基列‧耶琳。 \end{tabularx} \\ \\ \relax
9:18 & \begin{tabularx}{0.7\textwidth}{X} 因為會眾的領袖已經指著耶和華－以色列的神向他們起誓，所以以色列人不擊殺他們。全會眾就向領袖發怨言。 \end{tabularx} \\ \\ \relax
9:19 & \begin{tabularx}{0.7\textwidth}{X} 眾領袖對全會眾說：「我們已經指著耶和華－以色列的神向他們起誓，現在我們不能碰他們。 \end{tabularx} \\ \\ \relax
9:20 & \begin{tabularx}{0.7\textwidth}{X} 我們要這樣對待他們，讓他們存活，免得因我們向他們所起的誓而憤怒臨到我們。」 \end{tabularx} \\ \\ \relax
9:21 & \begin{tabularx}{0.7\textwidth}{X} 領袖對會眾說：「讓他們活著吧。」於是他們照領袖所說的，為全會眾作劈柴挑水的人。 \end{tabularx} \\ \\ \relax
9:22 & \begin{tabularx}{0.7\textwidth}{X} 約書亞召了他們來，對他們說：「你們為甚麼欺騙我們說：『我們離你們很遠』呢？其實你們就住在我們附近。 \end{tabularx} \\ \\ \relax
9:23 & \begin{tabularx}{0.7\textwidth}{X} 現在你們當受詛咒！你們中間必不斷有人作奴僕，為我神的殿作劈柴挑水的人。」 \end{tabularx} \\ \\ \relax
9:24 & \begin{tabularx}{0.7\textwidth}{X} 他們回答約書亞說：「因為確實有人告訴你的僕人，耶和華－你的神曾吩咐他的僕人摩西，把這全地賜給你們，並要在你們面前除滅這地所有的居民。我們因你們的緣故很怕自己喪命，就做了這事。 \end{tabularx} \\ \\ \relax
9:25 & \begin{tabularx}{0.7\textwidth}{X} 現在，看哪，我們在你手中，你看怎樣待我們是好的，是對的，就這樣做吧！」 \end{tabularx} \\ \\ \relax
9:26 & \begin{tabularx}{0.7\textwidth}{X} 於是約書亞就這樣對待他們，他救了他們脫離以色列人的手，以色列人沒有殺他們。 \end{tabularx} \\ \\ \relax
9:27 & \begin{tabularx}{0.7\textwidth}{X} 那日，約書亞分派他們到耶和華選擇的地方，為會眾和耶和華的壇劈柴挑水，直到今日。 \end{tabularx} \\ \\
[1ex]
\hline
\hline
\end{longtable}
$^{1}$鼎姐妹平安.
欺詐這兩個字我相信.
在近一兩年大家聽很多.
也都會很有共鳴.
我們都經常被提醒.
身邊有很多欺詐的陷阱.
我們要小心.
就算一些很有學識.
甚至是專業的人士.
有時稍一不留神.
都會上當 都會中招的.
根據警方的數字.
在過去一年 在2025年的時候.
總共發生了四萬多宗的騙案.
佔了香港整體罪案數字的接近一半.
為了幫助市民去提高警覺.
政府就拍了很多宣傳片.
也都推出了一些手機的程式.
一些APP來幫助市民.
不要那麼容易上當.
提高警覺真的很重要.
不單是生活裡面.
不要那麼容易來受騙.
提高屬靈的警覺.
我們要有清醒的屬靈頭腦.
對我們的生命都非常重要.
在行走天路的當中.
我們要時刻的警醒.
要有這份警覺性.
以致我們不會很容易陷入迷惑.
或者偏行己路.
然後就隨流失去的當中.
今天我們繼續一起思想約書亞記的訊息.
今天我們來到第九章.
我們會從第九章裡面.
一宗欺騙的事件.
去吸取教訓.
去吸取屬靈的教訓.
首先讓我們再一次同心低頭祈禱.
親愛的天父.

$^{41}$我們多謝你.
多謝你賜下你的恩賢.
賜下你的話語.
因為你的話語是我們腳上的燈.
路上的光.
懇求聖靈在此刻.
去好好地預備我們的心靈.
讓我們好好地去準備好.
去領受主日給我們的提醒.
又讓我們定義在生活裡面.
緊緊跟隨你 跟隨你的腳蹤而行.
多謝你聽我們同心的祈禱.
奉耶穌基督的聖名.
阿們.
首先我們重溫少少之前的記載.
當約書亞帶領以色列民.
過了約旦河之後.
他們先後很戲劇性地.
去攻取了耶利哥城和艾城.
他們的鳴聲.
就因此傳遍整個迦南地各個地方.
叫聽到的人都聞風喪膽.
於是當中在迦南地裡面.
很多的王就打算聯合起來.
一齊地去對付 去攻打以色列民.
希望能夠先發制人.
在今日的經文第九章一開始.
頭兩節就記載了有六個的王.
他們預備聯手地去對付.
去攻打以色列人.
跟他們爭戰.
不過與此同時.
卻是有另一群的民族.
他們打算用另一種方式.
去化解以色列民對他們的威脅.
他們就是駐在基片的希美人.
他們打算以詭計去蒙騙以色列人.
希望能夠與他們立約.
以致他們可以存活.
在經文第三至六節那裡這樣記載.

$^{81}$基片的居民聽見約書亞向耶利哥.
和艾聖所行的事.
就設詭計假充事者拿舊口袋.
和破裂鳳寶的舊皮酒袋.
拖在路上.
將補過的鞋穿在腳上.
將舊衣服穿在身上.
他們所戴的餅都是乾的醬料梅.
他們到吉格營中見約書亞.
對他們和以色列人說.
我們是從遠方來的.
現在求你與我們立約.
在《申命記》第20章裡面.
很清楚記載了神給以色列民一個吩咐.
當他們進入迦南地的時候.
除了一些很偏遠的城市之外.
他們不可以和迦南地裡的人去立約.
並且要將他們徹底地除滅.
以免以色列民被迦南地當中的人同化.
一起行在耶和華眼中捍衛真惡的事情.
相信欺騙人他們都知道.
以色列民他們有一個很重要的原則.
因此他們裝扮成為從極遠的地方而來的人.
從經文所見他們確實是花盡心思.
無論他們所穿的衣服.
他們所穿的鞋.
舊的破裂的啤酒袋.
和他們預備買的發霉的食物.
用今日的說法.
這個是心意的詐騙.
是deepfake來的.
做很多的功夫來欺騙眼前的以色列人.
其實以色列民他們都不是沒有防範的意識的.
在第七節那裡這樣記載.
以色列人對這些欺媚人說.
只怕你們是住在我們中間的.
如果是這樣的時候.
怎能與你們立約呢?.
不過原來他們這個防範的意識.
是非常之薄弱的.

$^{121}$很輕易就被欺騙人.
接下來要說一番動聽的話來瓦解.
他們說了什麼呢?.
欺騙人說了什麼呢?.
第八節.
他們對約書亞說.
「我們是你的僕人」.
約書亞問他們說.
「你們是什麼人?從哪裡來?」.
他們回答說.
「僕人從極遠之地而來」.
「是因聽見耶和華你神的名聲」.
「和他在埃及所行的一切事」.
「並他向約旦河東兩個阿摩利王」.
「就是希什本王西弘」.
「和在阿斯他祿的巴山王鄂」.
「一切所行的事」.
「我們的長老和我們那地的一切居民」.
「對我們說」.
「你們手裡要帶著路上用的食物」.
「去迎接以色列人」.
「對他們說」.
「我們是你們的僕人」.
「現在求你與我們立約」.
「我們出來要往你們這裡來的日子」.
「從家裡帶出來的這餅」.
「還是熱的」.
「看吶現在都乾了長了霉了」.
「這些啤酒袋」.
「我們盛走的時候還是新的」.
「看吶現在已經破裂」.
「我們這衣服和鞋因為道路甚遠」.
「也都穿舊了」.
「當我們一路聽這個情節的時候」.
「就好像一個劇本一樣」.
「忌騙人他們用了很多心機去寫這個劇本」.
「經過精心的預備的」.
「什麼要說什麼不說怎樣說法」.
「忌騙人首先就向以色列人」.
「表明他們向他們效忠」.

$^{161}$「自願卑微作他們的僕人」.
「同時藉此去抬舉約書娜」.
「和當中以色列的領袖」.
「忌騙人不單止抬舉約書娜和她的領袖」.
「同時間有抬舉耶和華的明星」.
「在經文裡面說」.
「他們看到耶和華的明星和大能的作為」.
「意思是他們是敬畏神的」.
「他們只是想過什麼要說什麼不說的」.
「他們只是說到神」.
「帶領他們在過約旦河之前的事情」.
「就是他們戰勝了西弘和鄂這兩個王子」.
「至於近來發生的事」.
「過了約旦河之後的事情」.
「他們攻陷了耶利哥城和艾城的事」.
「他們就絕口不提」.
「因為按理如果他們真的從遠方而來的時候」.
「在這段時間其實他們在路上」.
「他們不知道發生什麼事」.
「他們精心想過」.
「應該怎樣去向以色列民說這番話」.
「欺騙人這番動聽的話」.
「使約書亞和以色列民的領袖」.
「他們放下了戒心」.
「當我們一路看的時候」.
「我們見到欺騙的人」.
「他們花盡心思 想盡辦法 委曲求全」.
「或者我們會有一個印象」.
「其實他們是不是一群很弱小的族群呢?」.
「但當我們一路看下去的時候」.
「看到第十章我們就見到了」.
「其實他們是一座大城」.
「好像京城一樣」.
「皇帝都住在當中的」.
「比艾城更大」.
「並且在經文裡面提到」.
「城內的人都是勇士」.
「在第十章第二節就是這樣記載的」.
「可以看見當時以色列民的氣勢」.
「其實頗強盛 頗厲害的」.

$^{201}$「就算是很大的城 很強大的城」.
「他們的居民聽見的時候」.
「聽見他們的鳴聲」.
「他們都心都消化」.
「在這個時候」.
「以色列民正在一個很強盛的氣勢當中」.
「人在成功的時候」.
「可能心裡面都出現了微妙的變化」.
「欺騙人剛才的表現」.
「可以說是毫無瑕疵」.
「一點破綻都沒有的」.
「無論是外表」.
「無論是他們臣服的態度」.
「加上他們敬畏臣的表現」.
「一切看來都似乎是非常合理的」.
「約書亞和以色列的領袖聽見的時候」.
「他們有什麼反應呢?」.
在第十四十五節這樣記載.
我們一起讀.
預備起.
「以色列人受了他們之食物」.
「並沒有問耶和華」.
「於是約書亞與他們講和」.
「與他們立約」.
「用他們活著」.
「會眾的首領也向他們起誓」.
「約書亞和以色列的領袖」.
「果然上當了」.
經文裡面形容.
「約書亞和他們的領袖」.
「受了他們的一些食物」.
意思是拿出他們的食物.
來驗證是否真的發霉.
是否真的從遠方而來.
有些英文譯本譯做.
Examine their food.
看看是否真的.
看看他們的虛實.
看來他們很謹慎.
不過在經文第十四節的下半節裡面.

$^{241}$要強調的不是這一點.
要強調的是他們並沒有求問耶和華.
以色列的領袖們.
他們靠著自己的經驗.
靠著他們眼睛能夠看見的.
手能夠摸到的.
就作出了這個驗證.
作出了這個決定.
很可惜的是.
他們並沒有求問耶和華.
結果怎樣呢?.
真相很快就出現了.
在第十六節至二十一節.
「以色列人與他們立約之後」.
沒多久.
「過了三天才聽見他們是近民.
住在以色列人中間的.
以色列人就起行.
為三天悼了他們的誠邑.
就是欺騙基斐拉,比祿,基列耶林.
因為會眾的首領.
已經指著耶和華以色列的神.
向他們起誓.
所以以色列人不擊殺他們.
全會眾就向首領發怨言.
眾首領對全會眾說.
我們已經指著.
耶和華以色列的神.
向他們起誓.
現在我們不能害他們.
我們要如此待他們.
容他們活著.
免得有憤怒.
因我們所起的誓.
淋到我們身上.
首領又對會眾說.
要容他們活著.
於是他們為全會眾作了.
劈柴挑水的人.
正如首領對他們所說的話.

$^{281}$只是包不住火.
他們立約的消息.
很快就傳開去了.
真相亦都很快的傳回來.
去到以色列民的耳中.
他們就非常之緊張.
他們就立刻去欺騙這個誠邑.
去查過究竟是什麼一回事呢?.
他們發現原來他們不單止是近鄰.
他們更加是一個很有規模的族群.
是由四個誠邑來組成的.
但是這個時候已經為時已晚.
由於他們指著耶和華以色列的神起誓.
指著耶和華的名起誓.
因此他們不能夠爽約.
他們不能夠下巴輕輕.
當沒有發生過.
事實上當我們再看下去歷史的發展.
在大約三百多年之後.
蘇羅就因著去殺害了一些的欺騙人.
沒有守到這個約.
耶和華就懲罰以色列民.
去割下了三個月的饑荒.
這個事件記載在撒慕爾記下二十一章當中.
以色列民指著耶和華的名起誓.
誓已經起了.
他們無奈就要履行這個約.
條件是欺騙人要被詛咒.
永遠的成為以色列民的奴僕.
負責在耶和華所揀選的地方.
就是會莫亦都是將來的聖殿.
為會眾和耶和華的壇作劈柴挑水的人.
在獻祭的地方.
來幫助去提供柴,水,操縱的工作.
是他們終身要負責的.
這個罪事我們暫時看到這裡.
現在這一刻.
讓我們一起去想一下.
究竟這件事件.
在這件事件當中.

$^{321}$給了我們什麼很重要的屬靈教訓呢?.
給了我們什麼提醒呢?.
當我們剛才一邊看經文的時候.
一邊看著整個事情發展的時候.
我相信有一節經文大家會很深刻.
在第十四節裡面.
就記載了一個很致命的致命傷.
就是以色列民的領袖們.
他們在作出這個立約的決定之前.
他們沒有求問耶和華.
當一切表面的東西看來是合理的時候.
以色列的領袖.
他們就很輕率地做了這個決定.
接納了欺騙人他們立約的要求.
他們這麼草率來做這個決定.
確實是令我們很費解的.
其實沒多久之前.
他們才剛剛經歷了兩場的戰爭.
耶利哥城的戰爭.
艾城的戰爭.
都是按照著神的吩咐才能夠得勝的.
特別是艾城的戰爭.
我想大家記得.
他們才是大敗.
之後按著神的吩咐來得勝.
為什麼會這樣呢?.
為什麼他們會這麼快忘記了.
要去尋求神引導的重要性呢?.
會不會之前是因為打仗.
這次是因為敵方來求和.
沒有什麼威脅.
沒有什麼大的威脅.
於是他們就鬆懈了.
他們就忘記了要尋求神.
尋求神的旨意.
弟兄姊妹.
如果按照這個邏輯的話.
只是在面對危難的時候才來尋求神.
在平靜安穩的時候就不需要了.
如果是這樣的話.

$^{361}$我們的信仰會變得很功利主義.
神只是為我們消災解困.
當我們在平順日子的時候.
我們就不要去找神了.
神也不要找我了.
有時候會不會這樣呢?.
有時候我們的信仰會不會這樣.
落入了一個以自己作為主的信仰裡面.
神在人心裡面的作用.
只是為了滿足人自己的需要.
這不是真正信靠神的信仰.
這樣的信仰也很容易迷失.
人按照自己的意思而行.
偏行幾路.
很快就會撞板.
很快就會失去方向.
其實以色列民他們忘記了一個.
很重要很重要的事實.
就是耶和華的軍隊.
耶和華是他們這隊大軍的元帥.
大家記不記得.
在之前第五章裡面說的.
是耶和華率領著他們去征戰的.
哪裡有軍隊是不需要聽元帥的指令的呢?.
是可以就這樣自己來做軍事行動的呢?.
或者自己作一些決定的呢?.
如果有這樣的軍隊的時候.
相信這隊軍隊很快會潰不成軍.
事實上在真實的戰爭裡面.
擾亂敵方的通訊系統.
是一個很重要的戰略.
以致敵方的將領不能夠向前方的士兵.
去發出清晰的指令.
或者這也是人的本性.
人很容易又回到以自己作主的傾向.
做自己生命的元帥.
神的引導只是有需要的時候才用的.
在我們生活的裡面.
有時候會不會都容易落入了這樣的光景裡面.
我們都經歷過.

$^{401}$當我們徬徨無助的時候.
我們都很渴望尋求神的幫助.
神你帶領著我.
你幫助我.
讓我們有智慧,有力量去面對這些困難.
但是當我們在順境的時候.
在平順的時候.
有時候我們不自覺地選擇跟神失聯.
選擇offline.
我們覺得都沒有什麼特別需要去問神.
都沒有什麼重要的事情.
起初可能覺得是一些小的事情不需要問.
接著是一些大的事情都不問了.
起初是一次不問.
第二次又不問.
漸漸地成為了我們生活的習慣.
我們生活裡面各樣大小的事情.
就以自己來作主.
著名的基督徒作家,紐易斯(C.S. Lewis).
他有一本經典的著作.
叫作《School Tape Letters》.
有些中文譯作.
「小心魔鬼很聰明」.
在這本書裡面講到魔鬼用很多的伎倆.
即是你不覺意的伎倆.
來叫人和神疏遠.
遠離神.
其中一個他們很重要的策略.
就是要慢.
慢慢地叫信徒離開跟隨神的軌跡.
慢到一個地步是你不會察覺的.
因為如果太大動作的時候.
你會醒覺的.
你會迴轉的.
你會知道危險.
你會掉頭走的.
所以他們要用盡辦法.
讓信徒覺得其實一小步一小步.
不去緊緊跟隨.
都無傷大雅的.

$^{441}$書裡面講魔鬼會引誘信徒.
就好像一個比喻.
讓他們不相信自己其實已經遠離神.
不相信已經遠離太陽.
取而代之的是.
其實他們不知不覺地.
以很慢很慢的速度.
向著最黑暗最冰冷的地域前進.
在書裡面有一句很有警惕性的說話.
他說通往地獄最穩妥的路是怎樣的呢?.
通往地獄最穩妥的路是漸進的.
是很慢的.
不是很斜的.
坡度平緩.
走起來很舒服的.
軟軟的.
沒有什麼突然要轉彎的需要.
沒有里程碑.
也都沒有路標.
慢慢一小步一小步走在這條路裡面.
就能夠很穩妥地去到地獄的當中.
弟兄姊妹我們真的要很警醒.
我們真的要有這份靈裡面的清醒.
有時候我們一個人會當局者迷.
有時候我們會有自己的盲點.
會有自己的限制.
會有自己的偏見.
所以我們很需要弟兄姊妹作為彼此的守望者.
很重要的.
你看到祂好像慢慢偏離的時候.
你提醒祂.
你拉祂回來.
如果你仍然沒有靈裡的守望者.
主內的守望者.
你可以向神求.
求神給你有好的屬靈的友伴心交.
能夠彼此扶持.
在行動上第一件事.
如果你未做你要做的.
你要回成長班.

$^{481}$有班弟兄姊妹跟你一起.
去行走這條天路.
能夠彼此提醒.
一起走在神的道路裡面.
如果我們沒有人提醒的時候.
如果我們沒有這個警醒的心的時候.
我們會很容易.
一步一步步向一個我們自己建構出來.
我們自己作主的信仰.
結果我們就會垂留失去.
我們就會迷失了.
在另一方面弟兄姊妹.
當我們看第九章的記述的時候.
裡面有一個很特別的地方.
如果你有留意到的話.
裡面是完全沒有提到耶和華的話.
完全沒有.
沒有神的說話.
其實在之前的記載.
之前那幾章的經文裡面.
耶和華的話經常臨到約書亞.
臨到她的指紋當中.
帶領著他們前行.
但是來到第九章.
很沉默.
沒有神的說話.
或者這個就是人不去尋求神的結果.
神的說話實在太重要了.
在我們腳前的燈.
路上的光.
在黑暗的裡邊.
沒有了燈會怎樣呢?.
我們一定會迷路.
並且我們很容易會跌倒.
其實很諷刺的是.
在第九章緊接著的.
第八章最後那幾頁的經文裡面.
正正就是說到約書亞在以巴盧山上.
對全體的百姓.
男女老幼.

$^{521}$來宣讀神的律法.
經文裡面記載.
摩西所吩咐的話.
沒有一句不宣讀的.
剛剛向整個的以色列民族.
由老到亂.
來宣讀摩西的律法.
重溫了.
但是轉眼間來到第九章的時候.
她是沒有了神的說話.
當以色列民的領袖.
來作出一個很重要的決定的時候.
也見不到神的說話.
弟兄姊妹.
當我預備今天講章的時候.
給我自己很大的提醒.
作為牧者.
作為領袖.
很需要神的說話.
在我們生命裡面.
做每一個的決定.
帶領著弟兄姊妹.
在我們當中的導師也一樣.
你也是很重要的領袖.
你怎樣來幫助弟兄姊妹生命成長.
有沒有神的說話在你生命裡面.
我們要謙卑在神面前來審查.
審查我們自己的生命.
有沒有神的話在我們心靈裡面.
其實不單止領袖們.
其實弟兄姊妹也需要.
我們每一個人都需要在神面前來審查.
特別是當我們作決定的時候.
有沒有神的話引導著我們.
引導著我們人生呢?.
大家都很熟悉的.
在詩篇119篇第十一節.
有一個很重要的提醒.
「我將你的話藏在心裡.
免得我得罪你」.

$^{561}$這節的經文很重要.
在《The Message》本裡也很傳神.
我們有沒有將神的話.
好像存款一樣.
傳入銀行的夾萬裡面.
珍而重之.
最安全的地方.
我們有沒有這樣的心呢?.
以致我們不會犯罪.
得罪神.
去到一個一敗塗地的地步.
弟兄姊妹們真的很需要.
神的話在我們生命裡面.
怎樣能夠有神的話在我們生命裡面呢?.
我們要對神的話認真.
我們要看重神的話.
在今天.
無論在教會.
或者在網上.
有很多聖經的資源.
有很多講道的資源.
我們可以有很多機會去聆聽.
去學習神的話.
但是我們要謹記的是.
我們要去實踐.
我們所聽的.
我們所學的.
如果不是.
你所聽.
你所學的.
神的話.
不是你的.
你不要以為.
有時候我們在網上隨便聽哪一套.
或者我們可能很草草地查了經.
草草地去靈修.
神的話就在我這裡.
你要真正去活出來.
去經歷.
神的話.

$^{601}$才是屬於你的.
才是在你生命的裡面.
雅各給我們一個很重要的提醒.
雅各他說.
聽道而不行道的.
就好像怎樣呢?.
就好像人對著鏡子.
看自己本來的面目.
看見走後.
隨即就忘了他的相貌如何.
隨即就忘記.
我們也很喜歡.
同行這個觀念.
很重要的.
有弟兄姊妹一起去.
在天路裡面一起去並肩同行.
我們今年有新開設的.
同行查經證書課程.
我們訓練士工也一直有同行小組.
楊弟兄姊妹能夠在靈裡一起同行.
非常之重要.
我們有屬靈的夥伴.
一起去結伴同行.
是很得無比.
但同時間.
同行這個觀念.
我特別想去強調多一重的意義.
是我們一同去行出神的話.
我們一起去兌現神的話.
一起去經歷神的應許.
我們和弟兄姊妹一起同行的時候.
有沒有這樣的抱負呢?.
一起同行的時候有沒有這樣的使命呢?.
神讓我們學習了.
讓我們一起經歷了.
有友伴了.
我們如何一起去活出神給我們的教導.
以至到神的說話.
真的能夠深深地藏在我們生命裡面.
帶領著我們的人生.

$^{641}$保護著我們的生命.
我們回到這個敘事當中.
在《欺騙人的詭計》裡面我們看到.
一切表面的情理是非常之合情合理的.
無論他們的外表.
剛才我們看過了.
他們的對答的態度.
按照人的常識.
按照人一般的邏輯.
其實沒有什麼懸念.
順理成章的.
不過原來已經偏離了神很重要的原則.
弟兄姊妹.
在我們生命裡面.
我們其實也累積了很多經驗.
我們也累積了一些知識.
甚至是專業的知識.
有時候我們覺得.
我們也有一定的能力.
去處理不同的事情.
於是我們又不知不覺地.
去單單按著自己的意思來行事.
一些好像看來理所當然的決定.
理所當然的做法.
不過我們卻是不知道.
原來我們已經背離了神的原則.
可能越走越遠.
在《箴言》第三章五至六節裡面.
有一個很重要的提醒.
我們也熟悉的.
《箴言》三章五至六節.
你們要專心仰賴耶和華.
不可以依靠自己的聰明.
在你一生所行的事上.
都要認定祂.
祂別指引你的路.
其實都是人性來的.
人很容易會依靠.
很喜歡依靠自己的聰明.
一個人做決定.

$^{681}$可以很快.
可以很有效率.
也必定適合自己的心水.
我不知道你有沒有認識一些人.
或者你自己也是.
很喜歡一個人去旅行的.
什麼時候去哪裡吃東西.
什麼時候去哪裡看什麼.
不需要跟別人商量.
很快就可以決定了.
不過在我們跟隨主的人生路上.
《箴言》這兩節經文.
給我們一個很有智慧的提醒.
弟兄姊妹我們每天睡醒覺.
睜開眼睛.
各樣大大小小的決定.
什麼是我們最穩妥的定準和方向.
經文提醒我們.
是專心仰賴耶和華.
在一切所行的事上認定祂.
不是靠著自己的聰明.
專心仰賴耶和華.
更貼近原文的翻譯.
是完全全心的去信靠祂.
幫我們全心的信靠.
打從心底的去信靠神的時候.
很自然的.
我們就會謙卑自己.
我們要放下自己種種的自恃.
轉移去尋求神.
仰望神.
在一切所行的事上認定祂.
以神的旨意作為依歸.
我們就可以心裡有一份安心.
有一份安穩.
因為帶領著我們的那位.
是那位創造天地.
又極愛我們的神.
他必指引你的路.
更多的翻譯亦叫做.

$^{721}$使你的路徑平坦正直.
英文是Make straight your path.
在我們走在祂的旨意當中.
走在祂的心意當中的時候.
神不單止給我們方向.
祂也會幫我們除去路上種種的障礙.
可能是不容易.
但是神會幫我們去修正我們的道路.
去鋪平我們的道路.
讓我們有祂施恩的手帶領著.
行在寬闊平坦的道路上.
希西加作猶大王的時候.
他行而和華.
眼中看為正的事.
他推行了一連串的宗教改革.
他廢去了一些的優壇.
儲葬,砍下一些的木偶.
不再容許敬拜外邦的偶像.
他決意要扭轉他父親亞哈斯王所作的惡.
相信在當時他推行這些種種的改革.
不是一件簡單的事情.
聖經記載當時他的經歷是怎樣的呢?.
聖經這樣記載.
在列王記下十八章.
「希西加倚靠耶和華以赤烈的神.
在他前後的猶大列王中.
沒有一個及他的.
因為他專靠耶和華.
總不離開.
謹守耶和華所吩咐摩西的誡命.
耶和華與他同在.
他無論往何處去.
盡都亨通」.
希西加專倚靠耶和華.
這個倚靠耶和華.
就是剛才我們看箴言裡面所講的.
「專心仰賴耶和華」.
就是這個字.
希西加倚靠神.
專靠他.

$^{761}$聖經裡面形容.
耶和華與他同在.
使他在智利猶大國.
推動種種很重要而艱鉅的改革的時候.
「盡都亨通」.
弟兄姊妹.
在時代的變幻裡面.
教會也要去轉變.
有很多地方要去更新.
你對回應時代的需要.
回應時代福音的需要.
牧羊的需要.
因此教會正是步向種種的更新.
改變的當中.
無論是軟件的.
硬件的.
是生命的建造.
是建堂的事工.
種種的更新.
建造.
很重要的是我們要敏銳.
神的帶領.
神的心意是怎樣的呢?.
希西加在推動.
剛才所說的.
宗教的改革.
社會的改革裡面.
相信不是一件容易的事.
面對很多的阻力.
面對很多的困難.
但經文裡面提醒我們.
希西加當年的經歷.
是他倚靠神.
他專心去倚靠他的時候.
神就為他開路.
神就與他同在.
叫他所行的道路.
無論往何處去.
盡都亨通.
弟兄姊妹.

$^{801}$在我們將要面對種種的更新的時候.
無論內在的.
外在的更新.
無論是建堂的事工.
是不容易的.
很大的工程.
建造的工程.
生命的工程.
我們怎樣來去承擔呢?.
經文這裡提醒我們.
只要我們專心仰賴耶和華.
不倚靠自己的聰明.
在一切所行的事上.
都認定祂.
祂必指引我們的路.
祂必為我們開平坦正直的路.
弟兄姊妹.
在我們個人的層面都一樣.
在個人的生命裡面.
在我們面對種種不同的困難.
一些可能要做一些不容易的決定.
可能是我們的家庭.
可能是我們的健康等等.
不知什麼的情況.
無論怎樣都好.
經文提醒我們.
不倚靠自己的聰明.
是全心仰賴信靠神.
在一切所行的事上認定祂的時候.
我們就不需要過分的驚慌.
不需要過分的擔心.
這樣的時候.
我們生命裡面種種的懼怕.
種種的掛慮.
就能夠除去.
取而代之的.
是我們靈裡面的平安和喜樂.
因為唯有耶和華是我們的拯救.
是我們的盤石.
是我們的高台.

$^{841}$倚靠祂的人必不動搖.
弟兄姊妹.
倚靠神的生命.
專心仰賴祂的生命.
神必定為我們開路.
祂師兄的手.
必定會一直拖帶著我們的人生.
走在豐盛的路上.
來到結束的時候了.
親愛的弟兄姊妹.
今天我們藉著約書亞記第九章.
神藉著祂的話語去提醒我們.
我們實在很需要有靈裡的警醒.
有這份靈裡的警覺.
我們要有屬靈的清晰的頭腦.
唯有這樣.
我們才有力的走在神的旨意當中.
我們必須要在生活裡面.
向上去尋求神的心意.
去貼近祂的心懷.
去討祂的喜悅.
在我們生命裡面.
很重要要有神的說話.
我們要一起走出神給我們的教導.
以致神的說話是真正屬於你的.
在生命裡面.
我們要專心仰賴神.
凡事認定祂.
不是靠著自己的聰明.
最穩妥的是在一切所行的事上認定祂.
弟子們,如果我們整個群體都是這樣的時候.
神必定會賜福給我們.
神必定會復興更新我們.
神必定會繼續使用我們.
讓我們同心祈禱.
親愛的天父,我們多謝你.
多謝賜下你的話語給我們有很重要的提醒.
求你常常地去幫助我們.
給我們有清醒的屬靈頭腦.
以致我們時刻行在你的心意當中.

$^{881}$讓我們不要依靠自己的聰明.
有時我們有意無意地行在自己的道路裡面.
漸漸遠離你.
主一求你.
師兄的手與我們同在帶領著我們.
兼顧我們的生命.
主一求你去紀念你賜福教會.
讓我們整個的群體都能夠有這份屬靈的警醒.
敏銳你的意思.
我們又願意緊緊地來跟隨你.
在我們一切所行的事上.
我們來求告你.
我們認定你是我們生命的主.
你是我們的元帥.
帶領著我們.
無論個人或是我們整個教會來到前行.
願你幫助我們.
多謝聽我們同心的祈禱.
奉耶穌基督的聖名.
阿們.
\newpage



\section{腓立比書 2:12-16}
\label{sec:031814}
\textbf{成就主的美意}~~~~李適清博士
\newline
\newline
連結: \href{http://www.church.com.hk/acms/content.asp?site=cdc\&op=show\&type=product\&code=031814\&layout=sermon}{\texttt{031814}} ~~~~ 語音日期: 2026-03-15
\newline
\newline
\hyperref[sec:031813]{\small{< < < PREV SERMON < < <}}
~
\hyperref[sec:index_chronic]{\small{[返順時目]}}
~
\hyperref[sec:031815]{\small{> > > NEXT SERMON > > >}}
\newline
\newline
腓立比書 2:12-16
\newline
\begin{longtable}{cl}
\hline
\hline
章節 & 經文 (和合本修訂版)\\
\hline
2:12 & \begin{tabularx}{0.7\textwidth}{X} 我親愛的，這樣看來，你們向來是順服的，不但我在你們那裡，就是我現在不在你們那裡的時候更是順服的，就當恐懼戰兢完成你們自己得救的事； \end{tabularx} \\ \\ \relax
2:13 & \begin{tabularx}{0.7\textwidth}{X} 因為是神在你們心裡運行，使你們又立志又實行，為要成就他的美意。 \end{tabularx} \\ \\ \relax
2:14 & \begin{tabularx}{0.7\textwidth}{X} 你們無論做甚麼事，都不要發怨言起爭論， \end{tabularx} \\ \\ \relax
2:15 & \begin{tabularx}{0.7\textwidth}{X} 好使你們無可指責，誠實無偽，在這彎曲悖謬的世代作神無瑕疵的兒女。你們在這世代中要像明光照耀， \end{tabularx} \\ \\ \relax
2:16 & \begin{tabularx}{0.7\textwidth}{X} 將生命的道顯明出來，使我在基督的日子得以誇耀我沒有白跑，也沒有徒勞。 \end{tabularx} \\ \\
[1ex]
\hline
\hline
\end{longtable}
$^{1}$馮檢基議員.
.
\newpage



\section{馬太福音 26:36-39}
\label{sec:031815}
\textbf{客西馬尼園的爭戰 (華)}~~~~曾興才牧師
\newline
\newline
連結: \href{http://www.church.com.hk/acms/content.asp?site=cdc\&op=show\&type=product\&code=031815\&layout=sermon}{\texttt{031815}} ~~~~ 語音日期: 2026-03-15
\newline
\newline
\hyperref[sec:031814]{\small{< < < PREV SERMON < < <}}
~
\hyperref[sec:index_chronic]{\small{[返順時目]}}
~
\hyperref[sec:031816]{\small{> > > NEXT SERMON > > >}}
\newline
\newline
馬太福音 26:36-39
\newline
\begin{longtable}{cl}
\hline
\hline
章節 & 經文 (和合本修訂版)\\
\hline
26:36 & \begin{tabularx}{0.7\textwidth}{X} 耶穌和門徒來到一個地方，名叫客西馬尼。他對他們說：「你們坐在這裡，我到那邊去禱告。」 \end{tabularx} \\ \\ \relax
26:37 & \begin{tabularx}{0.7\textwidth}{X} 於是他帶著彼得和西庇太的兩個兒子同去。他憂愁起來，極其難過， \end{tabularx} \\ \\ \relax
26:38 & \begin{tabularx}{0.7\textwidth}{X} 就對他們說：「我心裡非常憂傷，幾乎要死；你們留在這裡，和我一同警醒。」 \end{tabularx} \\ \\ \relax
26:39 & \begin{tabularx}{0.7\textwidth}{X} 他就稍往前走，俯伏在地，禱告說：「我父啊，如果可能，求你使這杯離開我。然而，不是照我所願的，而是照你所願的。」 \end{tabularx} \\ \\
[1ex]
\hline
\hline
\end{longtable}
$^{1}$我們每個人都需要睡覺..
有些人遲遲不能入睡,.
有些人隨時隨地都能睡覺..
睡覺對不同的人都有不同的意識..
當門徒們耶穌在禱告的時候睡著了,.
耶穌就責備他們..
在馬可福音第十四章三十七節,.
耶穌就對西門彼得說,.
「西門,你睡覺嗎?不能警醒片詞嗎?」.
但是有時候睡覺卻給不一樣的人帶來不一樣的好處..
有一天,老張去參加了一個有兩千人的語文競賽..
這語文競賽的題目是,.
用一句話來形容什麼是和平,寧靜和快樂的意境..
大家就絞盡腦汁,想盡辦法寫出最好的一句話..
老張就在紙上隨便寫了幾個字..
結果得頭獎的竟然是老張..
大家非常地好奇,到底老張寫了什麼字呢?.
老張只寫了六個字,.
「我太太睡著了.」.
所有的評審頒獎的時候感動得淚流滿面,.
並深深地擁抱著他,.
告訴旁邊的人說,要精心禱告,不要睡覺..
如果你稍微有留意基督教的階級的話,.
再過三個禮拜,就是我們一年力度的復活節還有送難節..
你知道嗎,在這個之前的四十天,.
就是在2月18號,被稱為聖徽星期三,H Wednesday,.
乃是改革中傳統的宗派,.
包括聖公會,路德會,慰禮公會,.
都會把這四十天的時間來鼓勵信徒們要好好地進食禱告,.
好好地在神面前反省,默想,然後悔改,認罪,.
特別來思想耶穌基督為我們的受苦,受難..
因此今天早上為了預備大家參加我們4月3號的送難節,.
我們一起來看一段非常非常寶貴的經文,.
就是耶穌基督在克西瑪尼園的禱告..
我稱這個禱告為克西瑪尼園的征戰,.
因為這是耶穌在地上三十三年半當中最厲害,.
也是最大的一個征戰..
這個征戰記載在馬太福音26章,馬可福音14章,還有路加福音22章..
我們先來看馬太福音的記載26章,我們從36節到39節..
這裡說耶穌同門徒來到一個地方,名叫克西瑪尼,.

$^{41}$就對他們說,你們坐在這裡,等我到那邊去禱告,非常清楚地吩咐..
於是帶著彼得和西比特的兩個兒子,就是雅各和約翰同訓,就憂愁起來,極其難過,38節..
便對他們說,我心裡甚至憂傷,幾乎要死,你們在這裡等候,和我一同警醒..
他就稍微往前走,俯伏在地,禱告說,我父啊,貪惡可心,求你叫這杯離開我,.
然而不要照我的意思,只要照你的意思..
我們先來做幾個非常重要的觀察..
首先這個事情發生的地點,乃是在甘藍山上的克西瑪尼園..
克西瑪尼園的意思就是壓油池,或者是炸油池,.
乃是他們用來壓炸這個甘藍,成為甘藍油的一個地方..
可以說這是耶穌進入禪場去禱告的一個地方..
根據科學家的研究,這個甘藍園的最老的一個樹大概有九百年之久..
所以每年如果你去以色列旅遊的話,.
然後第二,這個事情到底發生在什麼時候呢?.
幾乎是進入到深夜了..
因為前文告訴我們,耶穌基督和門徒們喫了最後一個預約戒的晚餐,也就是最後的晚餐..
在做晚餐的時候,耶穌基督曾經預言有人要出賣他,就是猶大要出賣他..
之後呢,我們看到說他們當中的兩個門徒還在那時候爭論誰要做老大..
不但是如此,耶穌也在那個時候預言彼得在祭教之前要三次不認主..
你可以想象這時候已經進入到深夜的時候了..
然後這裡第三件我們看到的觀察是,耶穌帶了三個得意的門徒,就是彼得,雅各和約翰,把其他八位門徒留在一邊..
這三個門徒彼得,雅各和約翰曾經三次跟耶穌基督在一起,第一次是亞魯的女兒復活的時候,他帶了三個門徒去..
第二次乃是耶穌基督登山變相的時候,耶穌把這三個門徒也帶在身邊..
今天是最後一次在克西滿園禱告的時候,他帶了這三個門徒去,為了讓他們能夠給他心靈上還有情感上的支持,給他打氣..
這我們清楚看見耶穌基督在這裡,雖然他是一個神的兒子,但是他也完全流入人心中裡面的驚恐,他需要帶導夥伴的支持,阿們..
今天早上我們來看耶穌基督在克西滿園的徵葬,可以說是他33年半在地上的日子當中最大的一次徵葬,也是可以說最厲害一次的徵葬..
到底有哪些徵葬呢?三方面的徵葬..
第一個徵葬是肉體和心靈的徵葬,一起來講來.肉體和心靈的徵葬,要得勝的話必須要怎麼樣?要警醒禱告,要警醒禱告..
這是我們最常見的一個屬靈徵葬,肉體和心靈的徵葬..
像保羅所說的,很多時候我們想做的事情偏偏沒辦法去做,我們不想做的事情偏偏要去做,他說這可恥的身體誰能夠幫助我呢?.
這是我們最常見的一個屬靈的徵葬,肉體和心靈的徵葬,如果要得勝的話必須要怎麼樣?必須要警醒禱告,阿們..
這我們看見在另外一個福音書裡面,馬可福音的記載,我們來看看馬可福音的記載,跟馬太福音的記載非常相似..
第十四章我們從三十二節到三十四節,他們來到一個地方,名叫克西瑪尼,耶穌對門徒說,你們坐在這裡,等我禱告..
耶穌帶著彼得,雅各,約翰同去,就驚恐起來,同樣的形容的字眼,極其難過,然後三十四節講的話跟馬太福音講的話差不多是一樣的,對他們說,我心裡甚是憂傷,幾乎要死,你們在這裡等候,警醒..
耶穌就在地上所面臨的第一個屬靈的徵兆,肉體和心靈的徵兆.他肉體知道他將要揹負十字架,為了眾人的罪,釘死在十字架上..
心靈上他是非常的願意,他在想起十字架所帶來的痛苦,所帶來的羞辱,他心中裡面可能會有一些的恐懼,有一些的害怕,所以他說,我心中極其憂傷,甚至要死.我心裡甚是憂傷,幾乎要死..
這是耶穌基督第一次,在四本福音書裡面第一次告訴我們他心靈深處裡面所領受的情感上的傷痛.你可以想象嗎?神的兒子在這裡他也完全表露他在神面前他心靈上裡面的脆弱,他甚是憂傷,幾乎要死..
這也提醒我們在神面前表達我們的脆弱,不是羞恥的事情.不要打腫臉皮,臉皮充胖子.耶穌基督明白我們的軟弱,神也體恤我們的脆弱..
所以耶穌說,我心裡甚是憂傷,幾乎要死.你可以想象憂傷到什麼程度呢?我們看看路加醫生的記載,路加醫生的記載是跟馬太跟馬可福音的記載有所不同,主要有兩方面的不同..
我們看看路加福音第二十二章,四十三節到四十四節,這是其他福音書沒有記載的,有一位天使從天上顯現,加天給他力量..
四十四節,耶穌極其傷痛,禱告更加懇切,汗珠如大雪點滴在地上.路加畢竟是個醫生,觀察入微,他只有天使從天上顯現,加給他的力量..

$^{81}$神的兒子耶穌基督在地上面臨最大的這樣的一個陣仗的時候,他仍然需要天上來的力量幫助他,為他打氣,為他加油..
因此當我們面臨生活上的試探,甚至是失戀的時候,我們可以相信神會採擺天使來幫助我們,阿們..
第二點,我們看見耶穌基督禱告更加懇切,汗如大雪點滴在地上.哇,這個醫生畢竟是醫生..
這個是什麼症狀呢?醫學上來說,這是一個血汗症.這個血汗是非常罕見的一個跡象,乃至當一個人極其傷痛,悲傷,在極大的壓力底下,他身上裡面的微小的血管就爆裂,就摻雜在汗腺裡面,變成血跟汗,摻雜在一起,成為大雪點滴在地上..
你可以想象耶穌這個時候真的是傷痛到極點啊,他的壓力,他所承受的這些痛苦到了極點,所以他說我幾乎要死啊,難怪他的汗像大雪底一樣滴在地上..
你可以想象這個時候,他為什麼要帶三個門徒同去呢?就是希望能夠從他們身上得到一些精神或者是情感上的一個鼓勵,沒想到他們在做什麼呢?睡著了..
我來看看馬可福音第十四章三十七節,耶穌回來見他們,睡著了..
就對比著說,西門啊,你睡覺嗎?不能警醒片時嗎?片時其實原文的意思是警醒一個小時嗎?換句話說,耶穌已經去了禱告的時候,一個小時了..
總要警醒禱告,免得入了迷惑,你們心靈固然怨溺,肉體卻軟弱了..
門徒們真的非常願意來陪伴耶穌,在那裡守望禱告..
所以他們其實離開耶穌的地方,在路加門告訴我們說,丟一塊石頭的距離,大概如果你丟得遠一點的話,三四十尺,近一點的話,二三十尺..
他們可以看到耶穌禱告的地方,但是他們真的非常非常困見,他們就睡著了..
耶穌在這裡非常體恤他們的軟弱,祂說你們心靈固然怨溺,但是肉體卻軟弱了..
像那弟兄姊妹,這豈不是你我常有的關情嗎?心靈固然怨溺,肉體卻軟弱了..
教會很多的士工需要你來服侍,但你想你自己卻沒有時間,沒有精力來服侍主,你說什麼?心靈怨溺,肉體軟弱..
背經行動已經進入到三月份了,你說我年紀大了,記憶不好,背不起來了,讓年輕人去背吧,心靈怨溺,肉體軟弱..
有丈夫的,神靈告訴我們說,你們做丈夫要愛你們的妻子,如同基督愛教會,但是你知道嗎,我的太太總是嘮叨嘮叨,怎麼愛她呢?心靈怨溺,肉體軟弱..
太太說,我知道耶穌告訴我們說,在主裡要凡事順服我們的丈夫,那我的丈夫一點都不體貼,一點都不明慧我的需要,我怎麼能夠順服他呢?心靈怨溺,肉體軟弱..
門徒真的是如此,所以耶穌基督說,你們心靈固然怨溺,肉體卻軟弱,祂非常體恤他們的軟弱,不但他們是一次睡覺,幾次睡覺啊?.
兩次睡覺,三次睡覺?我們看看第十四章,馬可福音,三次九節到四十二節..
耶穌又去禱告,說的話還是先前一樣,如果是刻意的話,請你叫這杯從我面前撤去,然後不要照我的意思,只要照你的意思.又來,見他們睡著了,可能耶穌就一去禱告,一禱告就多少?一個小時..
可能兩個小時快過去了,看他們又睡著了,因為他們的眼睛怎麼樣?甚是睏倦!他們真的是心有餘力不足,也不知道怎麼回答耶穌啊..
四十一節.第二次來,看他們又睡了,對他們說,現在你們忍懶睡覺安歇吧.第三次來了,夠了,時候到了,看啊,人子被賣在罪人手裡了,起來,我們走吧,看啊,那賣我的人進來..
他們不但一次,兩次,三次,很想鼓勵耶穌,和他一起手腕禱告,沒想到一次,兩次,三次,都是睡得一塌糊塗..
在這裡說,他們眼睛非常的睏倦,你知道眼睛睏倦的時候,你真的是頂不住啊,特別是午飯以後,你要上課的話,我記得是在我念神學的時候,哇,我們香港來的一個同學,我們都說,午飯以後就很難撐,就真想用牙籤撐住你的眼皮,教授看到我們睡覺,他說,真不好嘛,對不對?.
這位香港來的同學告訴我們說,他沒問題,我可以睜著眼睛睡覺,到現在不懂他是說真的還是說假的,可以睜著眼睛睡覺,睜著眼睛睡覺,你睏倦到沒辦法,你想不睡都不行啊..
兩個禮拜前,我剛剛從非洲的加納回來,禮拜三早上六點多,我就去機場八點多的飛機,坐了大概二十多小時的飛機,到了非洲的加納,已經是禮拜四早上的九點多,特首穆斯來接我的時候,那時候已經快十點了,我就坐在他的車上,他跟我說要一個多小時,快一個半小時纔到他的教會,那時候他就開始跟我談話,你知道嗎?.
那時候我已經二十幾個小時沒睡了,真的是睏倦得不得了啊,想要跟他談話,沒兩下子我就呼呼大睡了,心靈怨逆,肉體軟弱啊..
這是我們肉體和心靈的一個很大的徵兆,為了要得勝,我們必須要怎麼樣?必須要警醒禱告.那耶穌基督進入到第二徵兆,第二徵兆是什麼呢?就是驕傲和謙卑的徵兆,一起講一遍來,驕傲和謙卑的徵兆,要得勝必須要怎麼樣?要效法耶穌..
這是我們每天所要面對的一個屬靈徵兆,驕傲與謙卑的徵兆.自從亞當氏族犯罪以後,我們人類最的根源就是什麼?驕傲,自大,自我中心,自私自利..
因此呢,驕傲是我們人的本性啊.驕傲就是告訴神說,我不需要你,我靠我自己就可以了,仁義丁天嘛.這是我們人驕傲的本性,所以這是我們要面對的第二個屬靈徵兆,也是耶穌在科西瑪尼園徵兆當中的第二個非常非常重要的一個徵兆,驕傲與謙卑的徵兆..
我們看到最典型的例子,他的門徒當中大家都想爭老大佬,特別是彼得.這位彼得通常說話比他想象的還要快.因此呢,耶穌就在這時候已經告訴他要發生的事情,我們看看在馬太福音第26章34到35節..
耶穌說,我實在地告訴你,今夜雞叫已遷,你要三次不認我.聽不清楚啊?當耶穌說,我實在告訴你的時候,每次我跟你講過很多遍,當耶穌說,我實在告訴你,或我實實在在告訴你的時候,那是非常重要的話,你要放下所有的手機,所有的東西,好好豎起你的耳朵來聽..
耶穌說,我實在告訴你,今夜雞叫已遷,你要三次不認我.彼得不自量力啊,說,我就是必須和你同死,也總不能不認你.哇,他是代表性的人物,是領袖嘛,我怎麼可以讓你失望呢?我就是和你必須和你同死,我也總不能不認你啊..
其他我們都不敢說,但是因為彼得說了嘛,就跟隨彼得了.你可以想象,這時候彼得自大驕傲到一種地步,連自己都不知道自己是怎麼一個人.我就是和你同死,我也不能丟棄你啊.哇,好大的口氣哦..
其實我們都知道,不久以後,耶穌被抓拿到大祭司院子裡面去的時候,彼得就在這位僕人師孃面前三次不認耶穌..
僕人師孃告訴我們說,耶穌轉頭看他的時候,他就跑了出去,在那裡痛哭.他就想起耶穌基督所說的話,計較以先,你要三次不認我.驕傲自大的結果,在那裡痛哭悔改..
後來彼得真的是痛哭悔改,後來成為耶穌基督所重用的大使徒,成為教會當時最棒的一個領袖.所以後來當他寫彼得前後書的時候,我們可以看到彼得完全的改變..
不再是那位驕傲自大,自以為是的,所以彼得前世第五章第六節說什麼呢?所以你們要怎麼樣?自卑,服在神大能的手下,到的時候,他必叫你們升高..
這些話就可以了,但他忘記了.其實在這個之前,耶穌基督已經告訴門徒們說,你們不要效法文士和法利賽人的假冒為善,驕傲,自以為是,自以為義..

$^{121}$所以耶穌在馬太福音23章12節告訴他們說,凡至高的必降為卑,自卑的必升為高.所以凡至高的必降為卑.彼得想爬得很高,靠著他自己的血氣,想要做得最高,沒想到他卻降為最卑微的,跌得最慘..
所以在彼得前世第五章第六節所說的,所以你們要自卑,順服在神大能的手下,到的時候,他必叫你們升高.彼得學會了,他懂得如何自卑,謙卑在神的面前..
所以我們這次要來勝過我們驕傲而謙卑的征戰的時候,我們必須要來效法耶穌基督.反過來我們看看耶穌基督是如何的謙卑呢?.
在馬可福音第十四章,我們繼續看三十五節到三十六節,這裡說,他就稍微往前走,俯伏在地,禱告說,倘若可行,便叫那時候過去.三十六節說,阿爸父啊,在你凡事都能,就你叫這杯撤去,然而不要從我的意思,只要從你的意思..
先看三十五節,這裡說,耶穌稍微往前走,俯伏在地,禱告..
你知道猶太人一般上禱告的姿勢是什麼姿勢?就是這樣做,舉起雙手面向耶路撒冷來禱告.但是耶穌在這裡禱告是什麼禱告啊?俯伏在地,是一種敬拜的姿勢,是一種最卑微,最無助的姿勢.俯伏在父神面前,禱告..
神的兒子如此的謙卑,在天父面前..
下一節告訴我們說,他禱告說什麼?阿爸父啊,阿爸父啊,就你將這杯從我面前撤去,然而不要從我的意思,只要從你的意思.阿爸父啊..
馬可福音裡面跟福音書裡面有時候摻雜了一些亞蘭文,當時耶穌基督的母語是亞蘭文,雖然他講希伯來話,所以亞蘭文是他的母語,在這個時候呢,以最親切的阿爸,阿爸父啊的,這個阿爸乃是亞蘭文,他對他天父非常親密的一個稱呼,阿爸父啊,就你把這個杯撤去,然而不要從我的意思,乃是從你的意思..
可以說耶穌基督自己卑微到這種地步,他就放棄自己可以做的決定,你可以想象耶穌基督是神的兒子啊,來到這個世上完成天父給他所要做的旨意,他可以選擇要做什麼不做什麼,這時候他卻將這個選擇的主權交回給天父,阿爸父啊,就你將這杯撤去,然而不要從我的意思,只要從你的意思..
因此我們要來效法耶穌基督的卑微,效法耶穌基督自己卑微,沒有人逼他的,沒有人要求他的,是他自己謙卑他自己,所以我非常喜歡使徒保羅,當他在描述耶穌基督的謙卑的時候,他以最簡單的字眼來形容描述耶穌基督最奇妙,最奧祕的真理..
在腓立比書第二章,第五節告訴我們說,你們當以基督耶穌的信為信,第六節和第八節說什麼呢?他本有神的形象,不以自己與神同等為強奪的,反倒虛極,取得奴僕的形象,成為人的樣式,就人的樣子,就自己卑微,全心順服,以至於死,且死在十字架上..
簡單的三句話就把耶穌的神性,人性完全形容得非常非常的切切,他本有神的形象,本來就是神嘛,但是卻沒有強詞墮裡,標板自己說我是神的兒子,我要怎麼做怎麼做,他乃是神本身,但是他卻沒有因此而來誇耀他自己,這裡說什麼?反倒虛極,成為人的樣式,取得奴僕的樣式..
為了來到世上神,成為人,他必須要倒空他自己,把天上的榮耀,天上的富貴,榮華都要放下,來到我們當中,住在馬槽裡面的誕生,成為奴僕的樣式..
第三點,既然成為人的樣子,既然有人的樣子,人的樣子就應該怎麼樣?就自己卑微嘛.我常常說,如果你是人的話,一定要懂得謙卑,不懂得謙卑的話就不是人..
他既有人的樣子,就自己卑微啊.他倒是人的話,人我們應該做的本分就是要自己謙卑自己,自己卑微,然後全心順服,以至於死,且死在十字架上..
你看到耶穌謙卑到什麼地步嗎?本來就是神啊,卻沒有因此強詞奪理,取了奴僕的形象,最卑微的,最被人藐視,看不起的僕人的形象,成為人的樣子..
有了人的樣子以後呢,他就自己卑微,全心順服,而且保羅強調說,且死在十字架上,死在最殘酷,最無人道,只有非羅馬人才能夠被釘死的十字架上..
你可以想象耶穌謙卑到什麼地步嗎?下一次當你感覺到自己很驕傲,自大,自誇的時候,麻煩你仰望一下十字架.當你自己覺得自己很棒,自以為是,自以為義的時候,麻煩你做個禱告,耶穌啊,幫助我謙卑吧..
驕傲和謙卑的徵兆,你要來效法耶穌基督,阿們.然後最後第三,也是最大的一個徵兆,是什麼徵兆呢?順服和卑溺的徵兆,一起講一遍來..
順服和卑溺的徵兆,要怎麼樣?要放棄權力.這是一個生與死的徵兆.這是一個抓住權力和放棄權力的徵兆..
最典型,最失敗的例子,莫過於加里爾人猶大,耶穌基督是二門徒當中的一個,他竟然用了三十兩銀子出賣了耶穌基督..
所以在克里瑪利安第三次禱告,看到門徒睡覺以後,跟門徒講話的時候,這裡面告訴我們什麼呢?我們看看路加福音二十二章..
二十二章四十七節到四十八節的記載,說話之間來了許多人,那十二個門徒裡,名叫什麼?猶大的,走在前頭,就見耶穌,要與他親嘴..
耶穌對他說,猶大,你用親嘴的暗號賣人子嗎?.
這個顯得傷盡人心的這位加里爾人猶大,竟然用一個最親密的一個動作,耶穌親嘴作為暗號,來出賣他跟隨了三年半的耶穌基督..
他為什麼如此的叛逆呢?如此的悖逆呢?因為他所期待的和耶穌所要成就的是完全不一樣..
他大恃所望,因為他是一個非常敬虔的猶太人,他期待著這位要來的彌賽啊,乃是要來轟轟烈烈帶領他們反抗革命,推翻羅馬的政府,讓他們猶太人可以自由的..
沒想到耶穌卻謙卑的像僕人一樣,替他們洗腳,服侍眾人,不要受人的服侍,他受不了..
竟然用最親密的動作,出賣了愛他的耶穌基督..
所以悖逆,我們看到這個悖逆的猶大,出賣了耶穌..
悖逆的猶大跟耶穌基督完全的順服,成了一個非常非常強烈的對比..
這也是在三本福音書裡面很清楚讓我們看見的,這個強烈的對比就在耶穌基督完全的順服身上讓我們看見..
我們再看看馬太福音26章,我非常喜歡這節聖經,26章39節,這裡說什麼呢?.
他就稍微往前走,符不在這個馬可福音講得很僵勢,禱告說,我父啊,倘若可行,求你叫這杯離開我,然而不要照我的意思,只要照你的意思..
我們非常熟悉的後半段這句話,然而不要照我的意思,只要照你的意思..
其實我們每一個跟隨耶穌基督的人,我們禱告都要學習這句話,而且兩遍來,然而不要照我的意思,只要照你的意思..
當你很想得到你想要得到的東西,最終你還要說,然而不要照我的意思,只要照你的意思..
當你很想為一件事情懇切地祈求,禱告,甚至進食,最後你要加一句說,然而不要照我的意思,只要照你的意思..
這句話說起來很容易,聽起來也很容易,但是如果你明白上文所講的話呢,這句話要做起來,親愛的弟兄姊妹,真的是不容易..

$^{161}$他的堂若可行,叫這杯從我面前賜去.到底是什麼杯啊?.
這個杯為什麼讓耶穌基督一次,兩次,三次,花了兩三個小時,在科西瑪尼人的征戰當中禱告,迫切,然後汗流大血點一樣滴在地上呢?.
因為這個杯不是美酒加咖啡的杯,這個杯乃是個苦杯啊,這個杯乃是神憤怒的杯,神要審判世人的杯..
在以下的書第51章17節很清楚地告訴我們,耶路撒冷啊,興起興起,站起來..
你從我手中喝了他什麼?憤怒之杯,喝了那使人東倒西歪的絕,以致喝盡.這是神憤怒的杯,什麼憤怒的杯?那是神審判的杯,那是神審判罪惡的杯..
因為耶穌知道,祂所痛苦的並不是實際的十字架上給祂帶來肉體的痛苦和羞辱,乃是祂要承擔世人的罪孽,而讓神的憤怒領到祂身上..
這是耶穌基督說憂傷幾乎要死的重點就在這裡,明白嗎?早在耶穌就來了之前700年,神就接著先知以賽亞告訴我們..
以賽亞書第53章第6節說什麼?我們一起來讀來..
我們都如羊走米,各人偏行幾路,以後華使我們眾人的罪孽都歸在他身上..
你可以想象神的兒子耶穌基督將全世界所有人的罪都壓在他的身上,神的憤怒因此要領到他的身上..
這是為什麼他三次禱告說,若是可以,但若可行,如果可以的話,求你將這杯撤去,求你將這杯除去,求你將這杯從我面前挪去..
然而不要成就我的意思,只要成就你的意思;然而不要從我的意思,只要從你的意思;然而不要照我的意思,只要照你的意思..
耶穌基督最痛苦的時刻並不是十字架之前的鞭打荊棘的冠冕,這只是肉身的痛苦,肉身的羞辱,這彬彬羅馬彬彬對他的羞辱拖磨在他身上..
這些都不是他最痛苦的地方,最痛苦的時候乃是他在十字架上大聲的呼喊..
我們來看看馬太福音第27章46節,約在深珠日中午的時候,耶穌大聲喊著說,你說耶穌這時候血都已經快流乾了,他還鼓起他全身的力量大聲喊著說什麼呢?.
伊利伊利拉巴撒巴格達尼,我的神我的神,為什麼離棄我?.
這乃是耶穌在十字架上第三三次大滅霸裡面最痛苦,最孤單的一刻呼喊.他不是用猶太,希臘文,也不是用希伯來文,他是用什麼?最熟悉的雅蘭文..
伊利伊利拉巴撒巴格達尼,我的神我的神,為什麼離棄我?.
因為就在那個時刻,他要死的那個時刻,眾人的罪孽都壓在他身上的時候,神憤怒的被領導他,神審判領導他的時候,就在那個時刻,天父轉眼不看他的愛子耶穌基督..
這是三十三年半以來第一次天父與他的愛子耶穌基督分開隔離的時候..
從起初轉到天地,直到耶穌基督受洗,聖父聖子聖靈他們三位一體的真神從來沒有分離過,從來沒有分隔過..
就他要死的那一刻,大聖父王說伊利伊利拉巴撒巴格達尼,不是為自己啊,親愛的弟兄姊妹,乃是為你的罪,你的罪,你的罪,我的罪,承受了父神的憤怒,父神的審判在他身上..
當你下次來參加送難節的時候,你明白耶穌基督在科學碼尼言的真章嗎?這是最大的真章,最厲害的真章..
這是為什麼說我起起憂傷,幾乎要死啊.你可以想象,神的兒子從來沒有和天父分離過,但是就是因為眾人的罪孽都壓在他的身上,天父不得不轉眼不看他的愛子..
你可以想象,這是耶穌基督在十字架上最寂寞,最孤單,最傷痛的時刻,父子分離的時刻..
難怪保羅在寫腓立比書的時候,當耶穌基督曾經經歷到人生最低層的時候,但是神卻把他升為世界上最高,最至高的..
我非常喜歡腓立比書第二章第九到第十一節,應該是腓立比書的高潮.他說什麼,所以神將他升為至高,又賜給他那超乎萬名之上的名叫一切在天上的,地上的和地底下的,無不因耶穌的名屈膝,無不口稱耶穌基督為主,使榮耀歸於父神..
Hallelujah!當耶穌基督完全的謙卑,自己卑微,群心,順服,以至於死,且死在十字架上之後,保羅受所謂神把他升為至高,並賜給他那超乎萬名之上的名,因耶穌基督的名,無不屈膝,無不口稱耶穌基督為主..
親愛的弟兄姊妹,克西瑪尼言的真章是耶穌在地上最厲害的真章,最大的真章,也是真情流露的一個真章..
想來想去,不明白耶穌基督為什麼要忍受,要接受這個悲.只有一個答案,因為他太愛你了,因為他太愛你了,因為他太愛你了..
愛到這種地步,他承受了最痛苦的那一時刻.伊利伊利拉巴撒巴赫達尼,我的神,我的神,為什麼離棄我?不為別的,只因為耶穌愛你.阿們,我們一起禱告..
幾乎要死的那種感受.因為你從來不是為你自己,從來不是為你自己,而是為我,為我每一位該死的罪人,承受了神的憤怒,神的審判..
為了拯救我們.謝謝主耶穌..
這麼偉大的愛,無法持度的愛,無法想象的愛,可以讓你承受與天父隔離的痛苦..
主啊,幫助我們..
你為我們付上了最大的代價,最痛苦的經歷,就是因為你愛我們,愛我們到底..
沒有人像你這樣愛我們,真的沒有人像你這樣愛我們..
我們心存感恩,你既然為我們死,我們就當為你活,就當為你活..
因為無人像你這樣愛我們.謝謝耶穌,謝謝耶穌..
你在克西瑪利亞的真相都得神了,我們纔有今天,我們才能夠坐在這裡敬拜你,讚美你,服侍你.謝謝耶穌,耶穌我愛你..

$^{201}$幫助我們更多的愛你,更正的愛你,更深的來愛你..
謝謝耶穌,聽我們的禱告,奉告主耶穌基督,寶貴的民窮,阿們..
\newpage



\section{馬可福音 14:3-9}
\label{sec:031816}
\textbf{為主打破玉瓶的女人}~~~~梁惠玲傳道
\newline
\newline
連結: \href{http://www.church.com.hk/acms/content.asp?site=cdc\&op=show\&type=product\&code=031816\&layout=sermon}{\texttt{031816}} ~~~~ 語音日期: 2026-03-15
\newline
\newline
\hyperref[sec:031815]{\small{< < < PREV SERMON < < <}}
~
\hyperref[sec:index_chronic]{\small{[返順時目]}}
~
\hyperref[sec:031817]{\small{> > > NEXT SERMON > > >}}
\newline
\newline
馬可福音 14:3-9
\newline
\begin{longtable}{cl}
\hline
\hline
章節 & 經文 (和合本修訂版)\\
\hline
14:3 & \begin{tabularx}{0.7\textwidth}{X} 耶穌在伯大尼痲瘋病人西門家裡坐席的時候，有一個女人拿著一玉瓶極貴的純哪噠香膏來，打破玉瓶，把膏澆在耶穌的頭上。 \end{tabularx} \\ \\ \relax
14:4 & \begin{tabularx}{0.7\textwidth}{X} 有幾個人心中很不高興，說：「何必這樣浪費香膏呢？ \end{tabularx} \\ \\ \relax
14:5 & \begin{tabularx}{0.7\textwidth}{X} 這香膏可以賣三百多個銀幣賙濟窮人。」他們就對那女人生氣。 \end{tabularx} \\ \\ \relax
14:6 & \begin{tabularx}{0.7\textwidth}{X} 耶穌說：「由她吧！為甚麼難為她呢？她在我身上做的是一件美事。 \end{tabularx} \\ \\ \relax
14:7 & \begin{tabularx}{0.7\textwidth}{X} 因為常有窮人和你們在一起，要向他們行善，隨時都可以，但是你們不常有我。 \end{tabularx} \\ \\ \relax
14:8 & \begin{tabularx}{0.7\textwidth}{X} 她所做的是盡她所能的；她是為了我的安葬，把香膏預先澆在我身上。 \end{tabularx} \\ \\ \relax
14:9 & \begin{tabularx}{0.7\textwidth}{X} 我實在告訴你們，普天之下，無論在甚麼地方傳這福音，都要述說這女人所做的，來記念她。」 \end{tabularx} \\ \\
[1ex]
\hline
\hline
\end{longtable}
$^{1}$我們開始聆聽神的話語之前.
我們一起低頭我們同心的禱告.
今日.
一班屬神的人.
我們來到.
神你自己的家.
我們可以的來到去敬拜你.
我們可以的來到去.
聆聽你的話語.
這是何等大的恩典和福氣.
神下我們真的為著這個世界.
來禱告我們看到最近.
這一個多星期以來的戰爭.
主我們真的要為著這場戰爭.
來禱告神下我們求你.
掌管整個的過程.
我們真的好盼望.
這個戰爭能夠早日的結束.
主我們知道一切都在你的手裡面.
求你去掌管.
也都求你去幫助.
神下我們身處在美國.
這一個國家.
我們同樣都看到在美國裡面.
都有好多好多的問題.
主這一個是以你的名立國.
的一個國家.
我們求你親自的來保守.
也都親自的來祝福.
讓這個國家再一次.
來高舉你的名.
去榮耀你的名.
主我們來仰望你.
神下我們今天來到主的面前.
我們一同的來等候你.
求你給我們有一個闖開的心.
有一個打開的耳朵.
我們求你向我們每一個人的心靈.
來講說話.
藉著聖經寶貴的話語.

$^{41}$讓我們更加明白.
你對我們的心意是什麼.
我們恭敬仰望等候.
是奉耶穌基督得勝的名求.
阿們.
我今天很想和大家分享.
一段大家很熟悉的經文.
馬可福音十四章三至九節.
如果你信了耶穌兩三年.
你都稍為勤力讀聖經.
你應該看過這段經文.
馬可福音十四章三至九節.
我們一同來讀這段經文.
預備起.
耶穌在伯大尼長大麻風的西門家女坐直的時候.
有一個女人拿著一瓶至貴的珍拿達香膏來.
打破玉瓶.
把香膏掀在耶穌的頭上.
有幾個人心中很不喜悅.
說:何用這樣枉費香膏呢?.
這香膏可以賣三十多兩銀子周濟窮人.
他們就向那女人生氣.
耶穌說:由他吧.
為什麼難為他呢?.
他在我身上做的是一件美事.
因為常有窮人和你們同在.
要向他們行善.
隨時都可以.
只是你們不常有我.
他所做的是盡他所能的.
他是為我安葬的事.
把香膏預先掀在我身上.
我實在告訴你們.
普天之下.
無論在什麼地方傳這福音.
也要述說這女人所做的.
以為紀念.
好.
今天我想和大家分享一個主題.
就是為主打破玉瓶的女人.

$^{81}$我想透過這段經文和大家思想.
有三個重點.
第一個重點是什麼呢?.
就是傳言的奉獻.
第一個重點是什麼呢?.
第一個重點是傳言奉獻.
在經文第三節.
第五節和第八節.
這裡說到耶穌在八大里長大麻風西門的家坐直.
當然不僅僅是耶穌和西門在家裡.
他說有一個女人.
拿著一瓶至貴的真拿達香膏.
在很多人面前打破了玉瓶.
將香膏敲在耶穌頭上.
第五節告訴我們.
這個香膏價值多少?.
可以賣三十多兩銀子.
第八節.
耶穌說他所做的是盡他所能.
他為我安葬的事.
將香膏預先敲在我的頭上.
這裡說到有一個女人.
拿著一瓶至貴的真拿達香膏.
如果我們看約翰福音.
其實有差不多這件事的一個比較詳細的記載.
在約翰福音十二章一至八節裡.
他說瑪利亞拿了一斤至貴的真拿達香膏.
是一斤.
我不知道我們當中有沒有女士們.
你們有沒有香水的習慣?.
有嗎?.
有的說有的.
我也沒有這種習慣.
但如果你平時有買香水.
如果是比較好的牌子.
你會發覺貴不貴?.
都貴的.
小小瓶都可能不便宜.
這裡說的是拿一斤.
不知道.

$^{121}$真是很多.
這裡說拿了這個真拿達香膏有個名字.
叫做真拿達香膏.
這個香膏有個名字叫真拿達香膏.
原來是名牌.
在當時來說是東方一個很著名的香料.
它能夠發出一些很濃郁的香氣.
經文說他從頭到腳.
將一瓶香膏打破了瓶子.
將香膏敲在耶穌的頭上.
其實在舊約的時候.
是皇帝加冕的時候.
才會敲香膏在他的頭上.
我們看到經文第五字告訴我們.
這個香膏可以賣多少錢?.
三十多兩的銀子.
大約是三百多個的銀幣.
約在當時來說.
多少錢的工資呢?.
是一年的工資.
在我們灣區.
我想問一下大家.
如果平均一些一般的收入.
一年的年薪大概一般人是多少錢?.
五六萬差不多吧.
還是再多一點?.
還是少一點?.
多一點嗎?.
一般而已.
一般都更多一點吧.
多少?七萬嗎?.
我來了這裡不是很久.
來了這裡多久呢?.
我搬了來灣區大概八年半.
搬了來這裡.
我發覺到.
我以前去寫是四萬至五萬.
有些人說四萬至五萬.
怎麼行呢?.
在灣區根本上不可以生存的.

$^{161}$那我加多一點吧.
五至六萬吧.
現在你們這樣說.
我又要再加一點.
可能要花六至七萬.
這裡年薪六至七萬.
等一下.
這個女子呢.
是用了她一年的年薪.
去買這個一斤的珍娜達香膏.
在耶穌的面前打破這個玉瓶.
敲在耶穌的身上.
一年的工資.
等於姐妹.
你捨不捨得用你一年的人工.
去奉獻給教會?.
捨不捨得?.
又突然間姐妹說.
不捨得.
是不是?.
不捨得.
其實我們在說.
我們很多時候.
我們唱歌的時候說.
神啊我們將我們自己傳言獻上.
剛才那首歌.
記不記得?.
我們將自己傳言獻上.
到底我們說傳言獻上的時候.
我們是不是只是唱歌的時候.
才會唱呢?.
只有嘴巴去唱.
只有嘴巴去說.
但是我們實際的行動.
我們不會這樣做.
但是你看到.
這個女子.
聖經所記載的.
這一個女人.
她願意.

$^{201}$當時的女人.
都不會出來工作.
我相信可能她是.
可能她家人.
或者各樣原因.
她儲起來.
可能儲了很久很久的那筆錢.
來買一斤的珍拿達香膏.
你覺得難不難得?.
你覺得她是不是傳言獻上?.
等於姐妹.
我們是不是一個傳言獻上的基督徒呢?.
馬太福音26章14-16節說.
猶大衛了三十塊錢出賣耶穌.
在當時耶穌的時代.
買一個奴隸是多少錢呢?.
只不過是十兩的銀子.
你買一個奴隸回來.
你就可以要他生就生.
要他死就死.
買一個奴隸.
三十塊錢可以買三個奴隸回來.
所以在當時來說.
三十塊錢是不是一個小數目呢?.
其實不是一個小數目.
對當時這個女子來說.
是一個相當大的數目.
弟兄姊妹.
你覺不覺得她所做的事.
是盡她所能呢?.
覺不覺得?.
她可不可以成為我們的榜樣?.
可不可以成為我們的激勵呢?.
可不可以提醒我們.
不要再做一個只會說.
只會唱的基督徒呢?.
而是我們願意有實際行動.
去回應上帝的恩典和愛呢?.
弟兄姊妹.
什麼叫做盡他所能呢?.

$^{241}$我覺得其實每一個基督徒的盡都不一樣.
同意嗎?.
你和我的盡都不一樣.
可能當中有些弟兄姊妹.
你要上班.
你的小朋友都很小.
你要照顧家庭.
是不是?.
所以每一個人.
你的盡都不一樣.
我們不需要和我們隔離的那個比較.
他放這麼多時間在教會.
他做這麼多事.
或者他已經退休了.
小朋友都長大了.
他可以拿多一點時間出來.
你不需要和別人比較.
你的盡和他的盡.
有什麼分別.
你不需要.
你只需要問你自己.
你只需要在神面前去醒察.
我已經盡我所能去擺上.
我能夠盡我所能去愛主嗎?.
什麼叫做傳言的奉獻呢?.
弟兄姊妹.
我們可以問問自己.
到底我們是不是一個.
傳言奉獻愛主的基督徒呢?.
傳言的奉獻和愛主的心.
是成正比的.
同意嗎?.
當你越愛主.
你就越願意.
更多的奉獻.
剛才菲烈.
我來的時候.
菲烈就帶我參觀你們教會.
也是第一次走遍整間教會.
菲烈告訴我.

$^{281}$招牧師買這個地方的時候.
很奇妙的.
我相信這些.
你已經聽過很多了.
是不是兩天就可以籌到.
四百萬?.
三百八十萬.
兩天.
有沒有聽錯?.
我相信很多弟兄姊妹.
你是真的傳言的擺上.
是不是?.
為什麼有些弟兄姊妹.
願意奉獻這麼多.
擺上這麼多呢?.
我相信.
是和愛主的心.
是成正比的.
你越愛主.
你越經歷主的恩典.
你就越願意擺上更多.
甚至毫無保留的奉獻.
傳言的奉獻.
弟兄姊妹.
你賺得多.
是不是等於你奉獻得多?.
不一定.
有很多弟兄姊妹.
可能年薪是很多的.
但未必等於他很捨得奉獻.
他收入少.
等不等於他奉獻少?.
都不等於的.
有些他收入很少.
但他很樂意.
來奉獻.
我記得.
每一次回香港.
我有機會遇到一個弟兄.
其實他最近已經告訴我.

$^{321}$他退休.
不過他仍然繼續做兼職.
他在建築業做保安.
你知道那些都不是很高收入的.
但他每一次.
我回去見到他.
他總是塞信封給我.
那裡全部都是美金.
不是港幣.
他說.
梁傳道.
這些錢.
你拿回去.
如果你看到.
神的事工有什麼需要.
你看到有什麼人.
有需要.
你就幫我.
去奉獻出來.
他不是給我.
你不要誤會.
我不敢用他的錢.
我不會放在自己的袋子裡.
我想告訴你.
我會收起.
放在抽屜裡.
寫著是弟兄的奉獻.
我就會看看.
有哪個機構很需要金錢.
或者哪個宣教士有需要.
或者哪一個人有困難.
我就會在那裡抽一些錢.
去拿去奉獻.
代這個弟兄去奉獻.
我為這個弟兄所擺上的.
他說你千萬不要告訴別人.
他覺得勝經說.
左手不可以讓右手知道.
所以我也不會告訴你.
他叫什麼名字.

$^{361}$弟兄姊妹.
求主幫助我們.
我們學習一個樂意奉獻.
做一個樂意奉獻的基督徒.
瑪利亞很清楚耶穌的身份.
和他的工作.
所以他樂意.
耶穌說他是為我安葬的事.
去預備的.
所以他知道耶穌.
很快就要上十字架.
所以他預先為耶穌做了這件事.
你對耶穌的認識有多少呢.
今天很多弟兄姊妹.
信了耶穌十年.
都未看過一次聖經.
有些甚至告訴我.
信了二十年都未看過一次聖經.
我不知道你信了耶穌多久.
我不敢問你們.
有誰已經讀過整本聖經一次.
我怕大家都會尷尬.
是不是.
弟兄姊妹.
求主幫助我們.
好好追求認識神的話語.
對耶穌有更多的認識.
不明白主的心意.
我們是不會全然奉獻的.
明白主的心意.
等不等於我們會有行動的回應呢.
這個留給大家去思想.
你是不是已經盡你所能.
獻上你自己呢.
這個我們需要好好思想.
今天我想在這一點上.
和大家分享一個.
宣教士的見證.
我很多年前.
在一個大型聚會裡.

$^{401}$買了一本書.
叫做《大漠飛沙》.
不知道大家有沒有聽過.
有一個宣教士叫黃剛賜.
有沒有聽過.
黃剛賜宣教士.
這個宣教士.
我記得在2014-2015年.
我們還是住在芝加哥.
住了三年.
記得他們每一年的12月聖誕節.
他們有一個很大型的聚會.
在聖誕節有幾天的時間.
有二三千人參加.
我記得我那次參加的時候.
其中有一個session.
就是黃剛賜宣教士.
他分享他去非洲宣教的經歷.
當時我們很多人聽他的經歷的時候.
我們那時候有二三千人.
每個人聽的時候.
我們都流著眼淚.
我們心裡覺得很感動.
和大家分享黃剛賜宣教士的見證.
在1991年.
馬來西亞籍的宣教士黃剛賜夫婦.
他帶著三個兒女.
分別是十歲,九歲和六歲.
一家飛到非洲的沙漠國家尼日.
這個地方向當地的居民宣教.
不是向華人.
是向當地的人宣教.
為了兒女的教育.
他們沒有辦法.
他們只能將三個兒女放在首都警察會的宿舍裡.
和學校裡讀書.
兩夫婦去到更偏僻的窮鄉僻壤傳福音.
在那一晚.
他們將三個小朋友.
送到宿舍的時候.

$^{441}$六歲的小女孩.
哭得不得了.
又哭又吐.
一筆成聖.
在那個陌生的社監懷裡.
哭得一筆成聖.
平時對離別都看得很開的老兒.
都流下眼淚.
感情很豐富的老大.
一直在哭.
一次又一次.
在跑道裡追出來.
爸爸媽媽已經.
開著車準備離開宿舍.
老大一直追出來.
爸爸看到倒後鏡.
看到老大追出來.
停下車.
問老大什麼事.
是不是有什麼要跟爸爸說.
或者漏了要說的話.
爸爸問他.
什麼事.
老大說.
沒什麼.
爸爸說.
你回去吧.
繼續開車.
開了沒多久.
又看到老大追出來.
爸爸又再問他.
什麼事.
有什麼要說.
老大說.
沒什麼.
回去吧.
為什麼老大一次兩次追出來.
為什麼.
捨不得.
如果你是那對宣教士夫婦.

$^{481}$你捨不捨得.
捨不捨得.
弟兄姊妹.
如果你有三個子女.
十歲九歲六歲.
你捨不捨得帶著他們.
去非洲宣教.
還要把他們放在.
華首都的警察學校.
要這樣分開.
你看到三個子女.
哭得這麼淒涼.
你做父母的心會怎樣.
很難過.
弟兄姊妹.
是很不容易的.
我們說要全然奉獻的時候.
全然擺上的時候.
不單是你一個.
是整個家庭.
容易嗎.
其實一點都不容易.
他寫了這本書.
他後面其實有幾篇.
就是他三個兒女長大之後.
他三個兒女寫了.
我看其中有一個.
他們跟父母去宣教那幾年.
是他們人生裡面.
一個令他們成長.
有很多學習.
一個很寶貴的經歷.
他們不是很負面.
他們是很正面看這件事.
到他們長大了的時候.
他們都是繼續.
好像父母一樣.
一起去服侍.
弟兄姊妹.
求主幫助我們.

$^{521}$做一個全然擺上的基督徒.
第二點想和大家分享.
就是服侍的焦點.
第四節到第六節.
我們一起讀.
第四節到第六節.
一起讀.
(中文翻譯).
這一點我想和大家分享.
就是服侍的焦點.
這裡說到.
其實很多時候.
我們做一些事.
就算做對的事.
做好的事.
都會被人罵.
有沒有這個經歷.
很多人都是.
你在教會裡面服侍.
你很盡心擺上.
都會有人批評你.
你博甚麼.
總之你出於動機是好都好.
都會有人批評你.
這個女子所做的事.
有沒有人批評她.
有的.
在經文裡面說.
有幾個人心中很不起悅.
心裡想.
真的很浪費.
三十兩銀子.
很多錢.
這樣浪費.
三十多兩銀子.
你去周濟窮人.
都不知道多少人受惠.
向這個女人生氣.
其實.
你看回這段經文.

$^{561}$如果你仔細留意.
這段經文由開始到最後.
這個女人都沒有說過一句話.
是耶穌為她說話.
這裡說.
耶穌說.
由他吧.
為甚麼難為他.
他在我身上.
做的是一件美事.
等於姐妹.
其實.
如果我們不清楚.
我們服侍的焦點.
我們很多時候.
因為別人的批評.
別人的聲音.
會影響我們.
事奉的心智.
如果我們是.
清楚我們服侍的對象.
我們就不會那麼容易.
受別人的影響.
或者別人.
閒言閒語.
有時候我們聽到.
有些弟兄姊妹說.
我明年不做了.
做到要這樣受氣.
我不做了.
以前我在香港教會.
去事奉的時候.
我聽到很多弟兄姊妹.
在事奉裡面.
很有挫敗感.
他們要放下他們的事奉.
不想再受這些氣.
弟兄姊妹.
我們會不會.
有時候都會這樣呢.

$^{601}$大家有沒有試過.
有時候去一些聚會.
或者去一些.
會派一些長刊.
很多時候.
一些長刊後面.
有很多人的名字.
有沒有留意.
後面有誰誰負責什麼什麼.
但是很多時候.
這些長刊.
最後有一句.
大家猜不猜到.
最後有一句.
很多人的名字.
下面最後有一句是什麼.
排名不分先後.
為什麼要寫這句話.
大家知不知道.
為什麼要寫這句話.
就是怕被人罵.
等一下.
為什麼我的名字排得這麼後.
為什麼這裡沒有我的名字.
有時候漏了.
不經意漏了.
為防這件事.
就寫排名不分先後.
或者如有錯漏.
請勿見怪.
諸如此類.
這個我們看到.
今天人很在乎.
我的名字在不在.
或者我的名字排第幾.
排到最後.
有沒有搞錯.
弟兄姊妹.
到底今天我們.
事奉的焦點是什麼.

$^{641}$我們是不是很在乎.
人家知不知道.
我們有做事.
或者你在不在乎.
或者你會不會很在乎.
人家怎麼看你的事奉.
今天很多人.
很多弟兄姊妹.
我們的事奉.
可能我們只是為了.
博取什麼.
你會不會的事奉.
沒人見到.
那我不做了.
洗廁所都沒人見到.
我不做了.
弟兄姊妹.
事奉的焦點.
其實很重要.
何用這樣枉費.
其實這是一個價值觀的問題.
你用什麼去衡量.
叫做浪費.
叫做枉費.
你不用什麼去衡量.
你做還是不做.
可能你未信的朋友.
家人會覺得你很傻.
星期天.
為什麼不睡晚一點.
你們教會九點半崇拜.
九點半崇拜.
起碼都要八點半.
要起床.
梳洗.
吃早餐.
開車過來.
為什麼不睡晚一點.
很多人會覺得你很傻.
但你覺得什麼.

$^{681}$你覺得什麼.
你覺得很值.
你覺得值得.
所以你願意這樣做.
是不是.
所以有時什麼叫做.
值得與不值得.
在乎呢.
你覺得那樣東西的價值是怎樣.
如果你覺得服侍主.
就算受氣.
就算其他人不明白.
不要緊.
都仍然願意繼續去擺上.
願意繼續去服侍.
這個就是我們事奉的焦點.
弟兄姊妹.
我們看.
不明白主的心意.
我們當然不會全然奉獻.
明白主的心意.
等不等於我們會有行動的回應.
等不等於.
不一定.
你明白了主的心意.
很多人知道.
我們要奉獻.
很多人知道.
我們要傳福音.
很多人知道.
我們要讀聖經.
等不等於他們會做.
都不等於的.
弟兄姊妹.
是要明白.
加上愛主.
還要加上什麼.
感恩的心.
我們才會全然奉獻.
並且.

$^{721}$不會受周圍環境的影響.
有些人批評你.
有些人用什麼眼光看你.
弟兄姊妹.
不要受其他人的影響.
我們事奉的焦點.
要對準的是神祂自己.
神知道我們的心.
神知道我們擺上的是.
我們的心態.
我們的動機.
已經很足夠了.
我們看看第六節.
這裡說的是耶穌自己.
為瑪利亞.
去辯解.
祂在我身上做的是什麼.
是一件什麼事.
是一件美事.
弟兄姊妹.
我想問你.
你事奉.
你最終的目的.
想得到什麼.
得到別人的稱讚.
招牧師稱讚你.
你真是醒目了.
你真是忠心了.
當然.
招牧師的讚賞和肯定.
都很重要.
但是這個是不是最重要.
是不是.
誰的稱讚才是最重要.
有主的肯定和稱讚.
你覺不覺得心滿意足.
覺得嗎.
如果你還不覺得心滿意足.
請你找招牧師聊聊天.
就是說我們事奉的焦點有問題.

$^{761}$事奉的心態有問題.
同意嗎.
求主幫助我們.
事奉其實.
主肯定和主稱讚.
我們已經很滿足了.
最重要的是.
主耶穌怎樣評價我們的服事.
我剛才好像漏了一點.
沒有說.
我漏了一點.
看到嗎.
服事的焦點.
我後面括了一個什麼.
沒有明字的女人.
有沒有留意.
你心水清.
我們看馬可福音.
馬可記載這段經文的時候.
他沒有記載.
是哪一個女人.
拿這瓶至貴的珍拿達香膏.
去告耶穌.
有沒有留意.
馬可福音剛才我們讀的那段經文.
是沒有記載他的名字.
但我剛才讀約翰福音的時候.
約翰福音是有記載的.
我覺得.
馬可福音是刻意.
作者刻意不寫他的名字.
知不知道為什麼.
弟兄姊妹.
馬可福音.
他不寫他的名字.
其實.
沒有他的名字還不重要嗎.
其實不重要.
最重要的是他做了什麼.
不是他有沒有名字.

$^{801}$在聖經裡面記載.
對不對.
所以有一天.
你發覺到一些長刊.
本來是應該有你的名字.
但是不記得放.
那你不要生氣了.
好不好.
我們不需要生氣.
最重要的是主耶穌.
如何評價我們的服事.
主耶穌明白.
我們的心意.
這是最重要的.
第三點.
很想和大家.
去思想.
把握機會來服事主.
我們一起讀這一節經文.
第七節一起讀.
因為常有窮人和你們同在.
要向他們行善.
隨時都可以.
只是你們不常有我.
主耶穌說.
你去行善.
其實你隨時都可以.
不過.
不是經常都能夠.
像這位姐妹那樣做在我身上.
當耶穌被釘十字架.
當耶穌復活升天之後.
其實瑪利亞已經再沒有這個機會.
所以我們事奉的時候.
弟兄姊妹.
要把握機會.
來服事主.
不要以為.
明年再說.
等我退休.

$^{841}$退休後又要照顧孫子.
等孫子長大一點.
等孫子長大了.
不用你照顧.
然後你又腳痛.
你又不能走路.
你又如何服事.
弟兄姊妹.
其實你仍然可以.
行得走得健健康康.
去服事主的時候.
你要把握這個機會.
我記得在.
2022年的時候.
有一次在Fansgiving.
我和范傳道.
和一個弟兄.
他太太和他的兒子.
我們一起去旅行.
去拉斯維加斯.
去旅行,開車等等.
我記得我們回程.
回到拉斯維加斯的時候.
因為冬天.
很快天黑.
五點左右已經天黑了.
你知道在高速公路.
開車很快.
你開慢一點也會被人發現.
或者很危險.
我們一路開車.
去拉斯維加斯.
回程.
我們在回程的路上.
那個弟兄.
開車.
坐在司機位.
司機位的後面.
一路和弟兄.
聊天.

$^{881}$其他人坐在旁邊.
我發覺一路開車的時候.
突然那個車.
砰一聲.
撞了一些很大的.
我一路看著前面和弟兄聊天.
我也不知道我們的車撞了什麼.
但是很大聲.
你知道撞了什麼嗎.
什麼.
撞了一隻鹿.
但是看不到.
那隻鹿應該是很大隻.
我不知道.
那個弟兄很有經驗.
因為他不是第一次撞鹿.
我說什麼事.
我們撞鹿.
那輛車.
我告訴你那輛車.
你說嚴不嚴重.
很嚴重.
後來.
我們說快點停在一邊.
因為很危險.
那些車一路在後面.
一路衝過來的時候.
其實很危險.
我們勉強將車慢慢移到旁邊.
然後那個弟兄說快點下車.
因為所有的氣袋都出來了.
整輛車很大一股氣味.
我們就下車五個人.
然後那個弟兄問有沒有受傷.
有沒有受傷.
我們過去看看頸有沒有事.
看看腳有沒有在.
我們五個人.
一個人.
一點的受傷都沒有.

$^{921}$一點都沒有.
但是車就.
滔滔不絕.
因為天都很黑.
他還想到處去找.
我們撞的那隻鹿.
找不到.
但是他想找那隻鹿.
但是找不到.
我們懷疑因為衝力很大.
那隻鹿已經拋得很遠.
但是應該是很大隻.
才會令到車撞.
有些人說.
如果撞鹿很危險.
如果鹿在前面玻璃那裡.
他的腳可以插進來.
如果坐在前面.
可能會即時死.
但是我們很感恩.
我們五個人.
撞到頸都沒有.
我記得.
我認識一個姐妹.
她有一次撞車.
車都不至於滔滔.
但她撞到頸.
做物理治療.
醫了幾個月才康復.
但我們.
五個人沒有事.
這是不是上帝的.
保守.
是不是上帝的恩典.
這個意外.
我們隨時都可以死.
但我們五個人.
都沒有事.
在車上有兩個傳道人.
一位是我先生.

$^{961}$另一位.
是我們遠孫中心的訓練導師.
他會講道.
會教訓練.
如果我們.
三個人死了.
真是很可惜.
不過當然.
如果上帝要收回我們的生命.
做夠了.
不用再做了.
這是上帝的心意.
但其實.
我們沒有一個人知道.
我們還有多少日子.
我們還可以事奉.
你不知道.
我們很多時候都開車去旅行.
我發覺我們在這裡.
去高速公路有時都會有牌子寫著.
小心這裡會有鹿出現.
是不是晚上.
還是不要開高速公路.
很危險.
弟兄姊妹.
上帝存留你的生命.
你仍然坐在這裡.
上帝有什麼心意呢.
上帝是否叫你.
繼續坐在這裡.
是不是.
上帝存留你的生命.
叫你繼續在這裡.
不是叫你每個星期.
回來坐在這裡.
有些基督徒坐了十年二十年.
都不動的.
我在教會.
我在香港的教會.
我初初去那間教會事奉時.

$^{1001}$有四百人.
現在我的教會已經去到六百人.
我發覺有些弟兄姊妹.
認識了很多年.
每個星期都見到她.
不過她就回來坐.
不參與服事.
叫她傳福音.
她又說沒有空.
又說不懂得傳.
弟兄姊妹.
我們千萬不要做這些基督徒.
求主幫助我們能夠把握機會.
去做神要我們做的事.
第七節.
多謝你.
多謝你給我們有生命氣息.
多謝你給我們仍然.
可以來到教會.
我們一同的來敬拜你.
神啊你存留我們的生命.
今日我們要.
怎樣將神你給我們的生命.
我們真是.
可以全然的擺上.
去服事主.
我們沒有一個人.
是霸佔地圖.
沒有一個人在這裡.
我們的人生是浪費.
或者枉過.
我們願意你祝福.
每一個弟兄姊妹.
給她們為主所用.
榮耀主的名.
多謝你.
靠賴耶穌基督.
得勝的名頭.
\newpage



\section{約翰福音 9:1-41}
\label{sec:031817}
\textbf{能看見的瞎子}~~~~周曉沂傳道
\newline
\newline
連結: \href{http://www.church.com.hk/acms/content.asp?site=cdc\&op=show\&type=product\&code=031817\&layout=sermon}{\texttt{031817}} ~~~~ 語音日期: 2026-03-15
\newline
\newline
\hyperref[sec:031816]{\small{< < < PREV SERMON < < <}}
~
\hyperref[sec:index_chronic]{\small{[返順時目]}}
~
\hyperref[sec:031818]{\small{> > > NEXT SERMON > > >}}
\newline
\newline
約翰福音 9:1-41
\newline
\begin{longtable}{cl}
\hline
\hline
章節 & 經文 (和合本修訂版)\\
\hline
9:1 & \begin{tabularx}{0.7\textwidth}{X} 耶穌往前走的時候，看見一個生來就失明的人。 \end{tabularx} \\ \\ \relax
9:2 & \begin{tabularx}{0.7\textwidth}{X} 門徒問耶穌：「拉比，這人生來失明，是誰犯了罪？是這人還是他的父母呢？」 \end{tabularx} \\ \\ \relax
9:3 & \begin{tabularx}{0.7\textwidth}{X} 耶穌回答：「既不是這人犯了罪，也不是他的父母，而是要在他身上顯出神的作為來。 \end{tabularx} \\ \\ \relax
9:4 & \begin{tabularx}{0.7\textwidth}{X} 趁著白日，我們必須做差我來的那位的工；黑夜來到，就沒有人能做工了。 \end{tabularx} \\ \\ \relax
9:5 & \begin{tabularx}{0.7\textwidth}{X} 我在世上的時候，是世上的光。」 \end{tabularx} \\ \\ \relax
9:6 & \begin{tabularx}{0.7\textwidth}{X} 耶穌說了這些話，就吐唾沫在地上，用唾沫和了泥抹在盲人的眼睛上， \end{tabularx} \\ \\ \relax
9:7 & \begin{tabularx}{0.7\textwidth}{X} 對他說：「你到西羅亞池子裡去洗。」（西羅亞翻出來就是「奉差遣」。）於是他去，洗了，回來就看見了。 \end{tabularx} \\ \\ \relax
9:8 & \begin{tabularx}{0.7\textwidth}{X} 他的鄰舍和素常見他討飯的人，就說：「這不是那從前坐著討飯的人嗎？」 \end{tabularx} \\ \\ \relax
9:9 & \begin{tabularx}{0.7\textwidth}{X} 有的說：「是他」；又有的說：「不是，卻是像他。」他自己說：「是我。」 \end{tabularx} \\ \\ \relax
9:10 & \begin{tabularx}{0.7\textwidth}{X} 於是他們對他說：「你的眼睛是怎麼開的呢？」 \end{tabularx} \\ \\ \relax
9:11 & \begin{tabularx}{0.7\textwidth}{X} 那人回答：「有一個名叫耶穌的，他和了泥抹我的眼睛，對我說：『你到西羅亞池子去洗。』我去一洗，就看見了。」 \end{tabularx} \\ \\ \relax
9:12 & \begin{tabularx}{0.7\textwidth}{X} 他們對他說：「那個人在哪裡？」他說：「我不知道。」 \end{tabularx} \\ \\ \relax
9:13 & \begin{tabularx}{0.7\textwidth}{X} 他們把以前失明的那個人帶到法利賽人那裡。 \end{tabularx} \\ \\ \relax
9:14 & \begin{tabularx}{0.7\textwidth}{X} 耶穌和泥開他眼睛的那一天是安息日。 \end{tabularx} \\ \\ \relax
9:15 & \begin{tabularx}{0.7\textwidth}{X} 法利賽人又問他是怎麼得看見的。他對他們說：「他把泥抹在我的眼睛上，我一洗，就看見了。」 \end{tabularx} \\ \\ \relax
9:16 & \begin{tabularx}{0.7\textwidth}{X} 於是法利賽人中有的說：「這個人不是從神來的，因為他不守安息日。」另有的說：「一個罪人怎能行這樣的神蹟呢？」他們之間就產生了分裂。 \end{tabularx} \\ \\ \relax
9:17 & \begin{tabularx}{0.7\textwidth}{X} 於是他們又對那盲人說：「他開了你的眼睛，你說他是怎樣的人呢？」他說：「他是個先知。」 \end{tabularx} \\ \\ \relax
9:18 & \begin{tabularx}{0.7\textwidth}{X} 猶太人不信他以前是失明，後來能看見的，等到叫了他的父母來， \end{tabularx} \\ \\ \relax
9:19 & \begin{tabularx}{0.7\textwidth}{X} 問他們說：「這是你們的兒子嗎？你們說他生來是失明的，現在怎麼看見了呢？」 \end{tabularx} \\ \\ \relax
9:20 & \begin{tabularx}{0.7\textwidth}{X} 他的父母就回答說：「他是我們的兒子，生來就失明，這是我們知道的。 \end{tabularx} \\ \\ \relax
9:21 & \begin{tabularx}{0.7\textwidth}{X} 至於他現在怎麼能看見，我們卻不知道；是誰開了他的眼睛，我們也不知道。他已經是成人，你們問他吧，他自己會說。」 \end{tabularx} \\ \\ \relax
9:22 & \begin{tabularx}{0.7\textwidth}{X} 他父母說這些話，是怕猶太人，因為猶太人已經商定，若有宣認耶穌是基督的，要把他趕出會堂。 \end{tabularx} \\ \\ \relax
9:23 & \begin{tabularx}{0.7\textwidth}{X} 因此他父母說「他已經是成人，你們問他吧」。 \end{tabularx} \\ \\ \relax
9:24 & \begin{tabularx}{0.7\textwidth}{X} 於是法利賽人第二次叫了那以前失明的人來，對他說：「你要將榮耀歸給神，我們知道這人是個罪人。」 \end{tabularx} \\ \\ \relax
9:25 & \begin{tabularx}{0.7\textwidth}{X} 那人就回答：「他是不是個罪人，我不知道；有一件事我知道，我本來是失明的，現在我看見了。」 \end{tabularx} \\ \\ \relax
9:26 & \begin{tabularx}{0.7\textwidth}{X} 他們就問他：「他給你做了甚麼？是怎麼開了你的眼睛？」 \end{tabularx} \\ \\ \relax
9:27 & \begin{tabularx}{0.7\textwidth}{X} 他回答他們：「我已經告訴你們了，你們不聽，為甚麼又要聽呢？難道你們也要作他的門徒嗎？」 \end{tabularx} \\ \\ \relax
9:28 & \begin{tabularx}{0.7\textwidth}{X} 他們就罵他：「你是他的門徒，而我們是摩西的門徒。 \end{tabularx} \\ \\ \relax
9:29 & \begin{tabularx}{0.7\textwidth}{X} 神對摩西說話是我們知道的，可是這個人，我們不知道他從哪裡來。」 \end{tabularx} \\ \\ \relax
9:30 & \begin{tabularx}{0.7\textwidth}{X} 那人回答，對他們說：「他開了我的眼睛，你們竟不知道他從哪裡來，這真是奇怪！ \end{tabularx} \\ \\ \relax
9:31 & \begin{tabularx}{0.7\textwidth}{X} 我們知道神不聽罪人，惟有敬奉神、遵行他旨意的，神才聽他。 \end{tabularx} \\ \\ \relax
9:32 & \begin{tabularx}{0.7\textwidth}{X} 從創世以來，未曾聽見有人開了生來就失明的人的眼睛。 \end{tabularx} \\ \\ \relax
9:33 & \begin{tabularx}{0.7\textwidth}{X} 這人若不是從神來的，甚麼也不能做。」 \end{tabularx} \\ \\ \relax
9:34 & \begin{tabularx}{0.7\textwidth}{X} 他們回答他說：「你完全是生在罪中的，還要來教訓我們嗎？」於是他們把他趕出去了。 \end{tabularx} \\ \\ \relax
9:35 & \begin{tabularx}{0.7\textwidth}{X} 耶穌聽說他們把他趕出去，就找到他，說：「你信人子嗎？」 \end{tabularx} \\ \\ \relax
9:36 & \begin{tabularx}{0.7\textwidth}{X} 那人回答說：「主啊，人子是誰？告訴我，好讓我信他。」 \end{tabularx} \\ \\ \relax
9:37 & \begin{tabularx}{0.7\textwidth}{X} 耶穌對他說：「你已經看見他，現在和你說話的就是他。」 \end{tabularx} \\ \\ \relax
9:38 & \begin{tabularx}{0.7\textwidth}{X} 他說：「主啊，我信！」他就拜耶穌。 \end{tabularx} \\ \\ \relax
9:39 & \begin{tabularx}{0.7\textwidth}{X} 耶穌說：「我為審判到這世上來，使不能看見的看見，能看見的反而失明。」 \end{tabularx} \\ \\ \relax
9:40 & \begin{tabularx}{0.7\textwidth}{X} 同他在那裡的法利賽人聽見這些話，就對他說：「難道我們也失明了嗎？」 \end{tabularx} \\ \\ \relax
9:41 & \begin{tabularx}{0.7\textwidth}{X} 耶穌對他們說：「你們若是失明的，就沒有罪了；但現在你們說『我們能看見』，你們的罪還在。」 \end{tabularx} \\ \\
[1ex]
\hline
\hline
\end{longtable}
$^{1}$喜大利兄姊妹 早安.
在開始進入經文之前.
今天想和大家做一個實驗.
因為我之前也提到一個相關的實驗.
那我就想邀請大家.
先問問大家 平時有誰有戴手錶的?.
對 而且今天可能就是.
比較是一些非電子手錶的人.
這個實驗會比較好.
大家首先用沒有戴手錶的手.
就先按著你的手錶.
按一按.
我就想問一下大家.
知不知道你的手錶的錶面.
是用阿拉伯數字 羅馬數字.
還是只有線條的呢?.
先想一想.
想一想自己的手錶是甚麼樣子的呢?.
好 大家都想了 是不是?.
我們一起來看一下你的答案.
看看是否正確.
有沒有人是答對的?.
有沒有人是答錯的?.
也有些人.
好 現在我們又用手按著你的手錶.
好 我想問一下.
大家知不知道你的手錶的錶面.
是否有秒針的呢?.
先想一想.
好 我們又來看一下答案.
好 想問一下有沒有人是答對的?.
有沒有人是答錯的?.
也有些人.
好 我們又按著手錶.
好 最後就是想問大家.
知不知道現在是幾點呢?.
準確地 準確地.
是幾點?.
好 我們來看一下.
又來看一下答案.

$^{41}$誰是答對的?.
誰是答錯的?.
是 也有些人.
其實這個實驗也在週五.
週六也做了.
其實這個實驗是做甚麼的呢?.
它是說一個叫做不注意事慢.
即是甚麼意思呢?.
就是如果有些事你不留意它.
或者你只是看著其他東西.
譬如是甚麼呢?.
譬如剛才我叫大家.
看一下數字是甚麼時候.
你就不會留意究竟有沒有秒針.
然後我問你有沒有秒針的時候.
你就看著秒針.
你就不會在多少點.
實驗結果在這裡.
成不成功也無所謂.
但我最想和大家分享的就是.
原來有時候我們去看東西的時候.
是會很聚焦在一些我們想看到的東西.
又或者是聚焦在.
我們覺得我們現在就是要看的東西.
同時間在信仰上.
其實我們也有時候會.
聚焦在一些我們想聚焦的東西.
而忽略了其他都是很重要的事物.
即使是我們本身的眼在視線範圍之內.
其實應該要看到的東西.
可能我們也會忽略了.
所以求主今日去光照我們.
讓我們能夠看得到我們應該要看到的事.
我們一起祈禱吧.
真是主多謝你設下你寶貴的華語.
主我求你的華語.
真是今日去光照我們生命裡面.
黑暗的地方.
求你讓我們今日單單去學悟神理.
去充滿我們的生命.

$^{81}$主多謝你祈禱 奉主明求 阿們.
今日的經文選擇了《約翰福音》.
《約翰福音》不像其他福音書.
其他福音書可能記錄至少十幾件神蹟.
但《約翰福音》只是精選地記錄了七件神蹟.
而今日我們會看的就是其中一件.
就是醫治身內黑眼的神蹟.
由於今日的經文比較長.
由一到第四十一節整個的第九章.
所以剛才邀請主席讀的只是第三十五到第四十一節.
是整章經文的最後一個結局或結論.
剛才讀的經文如果大家有印象.
就是它提到「他」.
究竟這個「他」是誰呢?.
剛才讀的經文一開始是發生甚麼事呢?.
其實整個第九章大概可以分為三段.
第一到第十二節就是醫治的故事.
是整個罪事的起點.
當耶穌向前行.
祂在聖殿繼續走的時候.
祂見到一個失明的人.
門徒問耶穌.
這個人身內失明是他犯罪還是他的父母呢?.
耶穌回答.
不是這個人犯罪.
也不是他的父母.
而是要在他身上顯出神的作為.
於是耶穌沒有像其他神跡事件那樣.
問這個人要不要痊癒.
也沒有跟他有對話.
反而直接套唾沫在地上.
然後用唾沫和泥抹在盲人的眼睛上.
然後再叫他去西洛亞池子裡洗.
於是盲子就看見了.
這就是第一到第十二節所交代的事情.
是整個大故事的起點.
亦一切是源於耶穌醫治了這個身內黑眼的人.
第二段是第十三到三十四節.
就是交代了因為盲子去開眼的這件事.
而在整個群體和整個社區所引發的紛爭.

$^{121}$就是法利賽人不滿.
原來耶穌去開盒子的眼的那天是安息日.
所以觸動了法利賽人的神經.
於是他們就找這個盒子.
去問清楚發生了什麼事.
問清楚來龍去脈.
法利賽人先後兩次去找這個盒子.
而這兩次找這個盒子去傾談之間.
就是找了盒子的父母去盤問.
而這一段第十三到第三十四節.
會是我們今天最主要看的經文.
這段經文我們會看到不同的人之間的討論.
他們溝通 甚至爭吵.
讓我們透過這段經文去檢視自己的生命.
今天就好像做一個視力測試一樣.
我們去檢視一下.
自己有沒有開始看得有點模糊.
有點不清晰.
求神去光照我們生命不清晰的地方.
那些黑暗的地方.
首先不如我們一起讀這段經文.
一二三起.
你那種善作者 徒難的養家.
又有的是自罵 又有的說是不是.
就算你長大 他們的說是謊.
這幾節的經文就講到.
那些群眾平時跟這個盒子.
經常在一起 或者是他身邊的鄰舍.
就在討論究竟這個盒子是誰.
他們認不出是不是之前那個人.
他們覺得是 又好像不是.
難道那個盒子.
就是去那個池子洗完之後.
他整個樣子就改變了.
似乎其實不是一個外觀的問題.
而是他們對於盒子開眼.
這件這麼奇蹟神蹟的事.
他們很詫異 他們不太相信.
所以既然覺得沒可能發生的話.
就會想會不會其實根本就不是同一個人呢.

$^{161}$這樣才能解釋根本就沒有發生過這個神蹟.
不知道大家遇到一些很神奇.
或者很奇妙的事的反應會是怎樣.
我就想起自己在高小初中的時候.
其實我是有這個脊柱 即是彎症的.
我想問大家有沒有聽過這個病症.
或者聽過 或者大概知道.
簡單來說 就是我的脊柱左右彎了.
或者有少許扭曲.
總之脊柱不是很直.
簡單來說.
當時我大概每三個月都要覆診一次.
而且一次比一次更加嚴重.
我的脊柱就越來越彎.
去到一次就已經去到臨界點.
甚麼叫臨界點呢.
就是醫生說如果我再彎多一點.
度數再增加的話.
我就要去戴一個箍.
讓我可以矯正 不要再彎下去.
又或者是甚至乎要做手術.
因為做手術才可以真正矯正問題.
醫生就說下一次覆診再看看怎樣.
再看看是否真的過了那個度數.
剛好在下一次的覆診之前.
就是Alpha的時候.
即是我們今週都開始Alpha.
而那個星期剛好是講醫治的那一課.
如果大家有上過Alpha.
或者經歷過醫治那一課.
真的會去為你醫治禱告.
於是我真的走到台前.
求一個醫治的禱告.
到下一次覆診的時候.
醫生就幫我量度我的脊柱.
然後就發現不單止沒有惡化.
還是有進步.
我的脊柱竟然變直了.
但是我當時的第一反應.
就是想 啊?真的嗎?.

$^{201}$我覺得每個醫生.
因為當時量度的方法是很人手的.
每個醫生量度的方法都有點不同.
而且我又覺得.
這個醫生是這樣 又不是我上一次覆診的醫生.
每次都不同.
就會想了很多想法.
去想應該不是的 不相信.
但是我會繼續覆診.
之後那幾次覆診.
仍然都是保持在一個進步了的度數.
總之是不同的醫生去量度.
接著去到最後一次覆診的時候.
醫生都說你的病情已經穩定了.
也沒有再彎下去.
也不需要任何其他帶箍 做手術等等.
於是我才真的相信.
原來我真的被治好了.
不知道大家遇到一些很驚奇.
或者神蹟的事的時候.
你的反應是怎樣.
第一下的反應是怎樣.
我們要小心自己.
會不會經常去懷疑.
明明祈禱了 神也聽了.
發生在你面前了.
但是你也去懷疑.
究竟是不是神的工作.
我們要小心.
不要讓自己去到一個地步.
就是事情已經發生在你的眼前.
你還是懷疑.
因為我們要小心.
這是一個初期眼瞎的症狀.
我們要小心這是其中一個病徵.
我們去懷疑的時候.
其實也是一個眼瞎的病徵.
我們要不要讓這個病情繼續惡化下去.
而另一個眼瞎的病徵.
就是對人的冷漠.

$^{241}$今天我們看的經文裡面.
看到因為瞎子去開眼.
不單是對於他自己.
這個瞎子的生命帶來改變.
甚至是為他周遭的環境.
都帶來一些翻天覆地的改變.
可能甚至是帶來紛爭,混亂.
剛才所說這些群眾之間的.
都是一些小紛爭.
只是討論究竟是不是他.
但是接下來由第十三節開始.
就會交代一個更大的紛爭.
就是法理財人之間的不滿.
不如我們一起讀第十三到第十七節.
一二三起.
法理財人男子的大恩典是講識人.
法理財人就問他自求多福放棄你.
他對他本說.
他那來活在我的眼中.
我一世就放棄了.
於是法理財人終究的說.
這個人不是從上到下的.
因為他不受我識人.
你有的說.
一個罪人怎能談得上極的神識呢.
他們之間就產生了分裂.
於是他們又對他談一說.
他可以理你的感性.
你說他是怎樣的人呢.
他自我殲滅.
第十六節說到.
他們之間就產生了分裂.
他們的意見不同的原因.
是源於為何耶穌能夠行神蹟呢?.
到底祂是誰呢?.
矛盾是他們覺得.
耶穌無守到安息日.
所以不是從上到下.
但另一方面.
耶穌又真的行了神蹟.

$^{281}$難道罪人也會行神蹟嗎?.
他們想不通.
為何他們會有這個矛盾呢?.
因為他們跳不出一個框框.
耶穌真的在安息日裡.
行了醫治的工作.
但到底耶穌是否真的無守安息日呢?.
到底安息日的意義是什麼呢?.
在《馬可福音》第二章二十七節說到.
安息日是為人設立的.
人不是為安息日設立的.
對於當時的猶太人來說.
安息日是很重要的.
守安息日是十誡之一.
是摩西的律法.
同時守安息日.
對於他們的歷史上.
他們在被擄的處境裡.
其實都是一個身份的象徵.
他們表達了他們很獨有的身份.
也是在見證神的恩惠和救贖.
所以他們很努力地守安息日.
當然安息日的重要律例是無庸置疑的.
但原來也不是不能改變的.
在猶太人的歷史上.
在兩日之間.
當猶太人起義要抵抗希臘人的時候.
他們的生命受到威脅的時候.
都可以違反安息日不能作供的規定.
他們的經典裡面.
都有提及到臨盆.
即是生孩子或是救火.
其實都不受安息日所限.
於是你可能會想.
「醫好萬人」是否與性命有關呢?.
以下的說話我都不想說出口.
不過容許我直接一點.
當我們有這個想法.
「醫好萬人」是否與性命有關.
出現在我們腦海的時候.

$^{321}$其實會否等同說.
「反正他都是由出世就盲」.
「多等一日都無所謂」.
原來有時候我們一些想法.
對於其他人來說是很冷血或很刻薄的.
安息日是為人所設立的.
不是規矩不重要.
而是我們是否真的明白規矩背後的意思呢?.
安息日本來就是讓人可以與神相遇.
可以很專注地去敬拜神.
甚至不單只是社會上某一些階層.
更加是勞工,瀑婢.
都可以與神相遇.
而回到今日的經文.
就是耶穌靠近這個盒子.
並且醫治他.
為的正正是讓他可以與神相遇.
之前我看到一個新聞.
有一個的士司機接載了一個手脫臼的病人.
去急症室.
其實乘客由上車的時候.
已經痛到臉青了.
所以司機希望很近急症室門口.
可以讓他下車方便一點.
讓乘客可以直接方便一點.
進入急症室.
誰知去到急症室門口的時候.
已經被保安截停了.
司機對保安說.
「我的乘客的手斷了」.
「我放下他 我下完客就走了」.
誰知保安竟然說.
「他的手斷了而已」.
「不是他的腳斷了」.
「他可以走過去的」.
但是有另一個的保安.
看到這個情況之後.
叫司機「你快一點」.
下車就完事了.
今天我們要反思的是.

$^{361}$我們有時候會否放大了一些規矩的形式.
甚至忘記了規矩背後的意思.
以至於我們展現不了鄰舍之間的愛.
甚至乎阻礙了人與神去相遇呢?.
接下來讓我們讀第十八到第二十三節.
一二三起.
.
這幾節的經文作者用「猶太人」這個字.
其實都等同「法理賽人」.
「法理賽人」他們不太滿意核子的回答.
於是他們因為剛才第十八節說到.
他們不相信它以前是生命.
於是他們就去找核子的父母.
希望聽到另一個故事.
其實我們有時候會否都是這樣呢?.
我們去問其他人的意見.
是為了去想聽一些自己想聽的東西.
譬如與好姐妹去逛街 去買衣服.
你就會問「哪件衣服比較漂亮?」.
可能你的朋友就會回答.
「藍色這件吧 藍色很適合你」.
然後你就會說「藍色很抑鬱」.
你的朋友就會說「那你買紅色吧」.
然後你就會說「是的 我也覺得紅色是好的」.
原來有時候我們去問別人意見.
或者去聽別人說話 或者去問別人的事.
其實是想從第三方的口中.
聽到自己想聽的東西.
就好像經文中的猶太人.
他們其實是帶著一個前設去問.
並不是很中立地去查探答案是甚麼.
如果我們太看重自己.
其實這也是一個眼核的病徵.
當我們去看東西完全主觀的時候.
其實我們看的東西.
都只是一些我們想看到的東西.
我們要留意我們有沒有這個眼核的徵狀.
而剛才那幾節的經文.
其實核子的父母的回應都是很奇怪的.
剛才讀的幾節裡面.

$^{401}$他們完全顯示出那種事不關己的冷漠.
就說「我不知道的 他是大人來的」.
「你問他吧」.
他們的情緒 或者他們的回應.
就是有一點不合情理.
我在讀神學之前.
我是一個兒科的職業治療師.
專門做一些SEM 即是特殊學習需要的個案.
例如是自閉症 過度活躍症那些.
經常見到家長 或者和家長聊天.
其實他們只要小朋友有進步.
甚至只是很少的進步.
其實他們都會很開心 很雀躍的.
甚至他們會到處和媽媽群組或朋友說.
這個練習好啊 或者這裡有什麼資源啊.
總之去分享他們能夠幫助到小朋友的一些方法.
但是今天經文裡面的父母.
對著自己的小朋友.
由生命到提得翻譯.
其實是一件很大件的事情.
很好的事情.
但是他們卻是很冷漠.
在第22節其實都已經解釋給我們知道了.
他講到他父母說這些話是怕猶太人.
因為猶太人已經商定了.
若有宣認耶穌是基督的.
要將祂趕出會堂.
這句就已經解釋給我們知道.
為什麼父母會有這樣的反應.
或者我們很難理解.
趕出會堂其實有多大件事呢?.
不過其實在聖經的那個年代.
猶太人的起居飲食都是在會堂裡供應的.
所以你一旦被趕出會堂.
不僅是你的身份沒有了.
更加是你很基本的生活需要都沒有了.
所以沒有人會敢得罪猶太人.
這裡令我想起一些穆斯林的歸居者.
他們一方面.
如果大家可能看網上或者一些警察會講的見證.

$^{441}$其實他們一方面經歷了神的真實.
但另一方面他們都知道.
如果相信耶穌的話.
其實他們就會失去整個的社群.
整個的家庭.
但是感恩的是很多的歸居者.
其實他們仍然很努力地堅持相信耶穌.
或者生活在香港的我們.
今天未必去到一個這麼極端的狀況.
但是我們都要去問自己.
到底我們看見的是困難.
還是我們看到困難背後的上帝呢?.
到底我們是否只是選擇一些容易的事情去做呢?.
接著我們繼續看經文第24到第34節.
一二三起.
於是法輪功學員第二節要表達以前心靈的眼光.
他對答說.
你要將明瞭歸給上帝.
我不知道這人是個罪人.
那人就不回答.
他是不是個罪人.
我不知道.
有一件事我知道.
我本來是神明的.
現在我看見了.
他們就問他.
他到底做了甚麼.
是甚麼改變了你的痕跡.
他回答他們.
我已經倒數你們了.
你們不怕.
為甚麼我又要命呢.
難道你們也要作假的叛徒嗎.
他們就回答他.
你是假的叛徒.
而我們是可信的叛徒.
上帝對過世說話是我們知道的.
可是這個人.
我們不知道是他從哪裡來的.
那人又答對答不說.

$^{481}$他開了我的眼睛.
你們仍不知道他從哪裡來.
這真是奇怪.
我們知道上帝不聽罪人.
唯有辯護上帝.
遵從下次旨令.
上帝才聽話.
從上世以來.
你曾經見有人開了.
親我就失明的人的眼睛.
這人又不是從上帝來的.
甚麼也不能做.
他們回答他說.
你完全是身在罪中的.
還要來教訓我媽媽.
於是媽媽馬上說了一句.
經歷了第一次.
法利財人和核子的討論之後.
剛才他們就說了.
有些紛爭.
然後他就覺得.
不太相信這個核子是否真的盲.
就去找他的父母談.
剛才父母就說.
我不知道的.
你去問他吧.
他是大人來的 懂得說.
於是法利財人又真的聽了父母說.
不可以再去找核子.
於是就是第二次和核子的討論.
在這次的討論中.
我們看到法利財人的情緒更加高漲.
而觸發他們的一句說話就是.
難道你們也要作他的門徒嗎?.
於是法利財人回覆.
你是他的門徒.
我們是摩西的門徒.
用「你」,「我們」去劃清界線.
甚至剛才讀到一段核子很長的回答.
是整個大敘事裡最長的回答.

$^{521}$到最後核子說.
這個人如果不是從上帝來的話.
其實什麼都不能做.
什麼意思呢?.
即是核子對法利財人說.
他是從上帝來的.
聽到這一句的時候.
法利財人就人身攻擊核子.
說你完全是身在最終.
還要來教訓我們嗎?.
法利財人他自以為是.
他藐視他人.
甚至要用人身攻擊.
去指證,批判,批評其他人的錯誤.
為什麼法利財人會這麼激動呢?.
是因為核子所說的話.
觸碰到他們的核心價值.
他們的底線.
他們的身份價值.
身份認同.
所以很被冒犯.
甚至他們從頭到尾.
已經很抗拒耶穌的所言所行.
現在還要有第三方.
甚至是一個核子.
甚至是一個平時討飯吃的核子.
一個低下階層.
去對他人說.
這個人是從上帝來的.
於是法利財人就覺得很被冒犯.
當一些我們一直慣常的做法.
或者自己覺得很正確的事.
去被觸碰,去被挑戰的時候.
大家會有什麼反應呢?.
其實作為基督徒.
我們有時候真的會被神破碎我們的生命.
耶穌跟門徒說.
若有人要跟從我.
就當捨己背起十字架來跟從我.
要捨己就是將我們的老我也丟下.

$^{561}$其實在聖經裡也有不同的人物.
他們也是經過神的破碎再建立的.
例如摩西.
摩西在年輕的時候.
血氣方剛.
想用自己的方法拯救同胞.
於是就將埃及人打死.
然後就要逃亡.
但是故事未完.
神破碎了他之後.
讓他也再建立.
透過不同的歷程.
讓他重新塑造為一個謙卑.
能夠帶領人出埃及的領袖.
又例如普羅.
普羅原名是蘇羅.
原本是一個逼迫基督徒的狂熱分子.
但是神就讓他在大馬士革的路上.
被大光照耀.
硬核了幾天之後.
這一場肉體的破碎.
也是在建立他的生命.
讓他去認清自己生命的一些軟弱的位置.
也讓他去聚焦神.
讓他成為了之後教會的領袖.
成為了使徒保祿.
台北有一個牧師曾經說過.
如果沒有經過破碎.
一個人將會永遠都活在恣意裡面.
唯有經過破碎.
將裡面的老鵝都拿走的時候.
神的恩典才可以臨到我們身上.
所以當這一群法利賽人.
去跟那盒子好像是一個人身攻擊的時候.
去罵他們的時候.
其實他們不只是在拒絕那盒子.
而是在拒絕上帝.
所以弟兄姊妹.
我們要小心我們的眼睛.
有時候會不會只是看到別人的不是.

$^{601}$只是去指責其他人 只是去藐視其他人.
在《箴言》第六章第十六到第十七節.
有講到「耶和華所珍惠的有六樣.
連他心所珍惠的有七樣.
就是高傲的眼 撒謊的舌.
流無辜人血的手.
圖謀惡戒的心 飛跑行惡的腳.
吐謊言的假見證.
並弟兄中暴散紛爭的人」.
剛才我們講了不同眼核的表現.
一開始我都講了這是一個視力的測試.
檢視自己有沒有夢察的地方.
最後我們一起去看一下.
耶穌和核子的對話.
我們一二三起.
你已經看見他 現在和你說話的就是他.
他就是主人我信 參就拜耶穌.
耶穌是可謂神父到這世上來.
使得能看見的看見 能看見的反映神明.
從他在那裡來拜祭的人.
聽見的聲音就是對他說.
大哥我夢察生命了.
耶穌對我看不見 你們的事實明顯.
受到罪孽 但現在你們不輸.
我夢察光天 你們的罪過了.
如果我們比較耶穌和核子的對話.
我們比較兩者的時候.
就會發現他們去看事物.
或者去看耶穌的角度是很不同的.
核子對於耶穌的看法是有所改變的.
由一開始人們問他的時候.
他就說是一個醫治人的人.
到之後就說是先知.
到之後就說是主人我信.
但在弗利賽人的眼中.
其實一直都沒有改變.
耶穌只是一個眼中釘.
是他們很珍惚的人.
核子是帶著一個很開放.
很謙卑的心去聽神的話.

$^{641}$而弗利賽人是一成不變 自以為然.
耶穌在《約翰福音》的前一章.
即是第八章說到.
祂是世界的光.
對於核子來說.
光是好的 祂渴望見到光.
渴望讓祂的眼睛重見光明.
但對於弗利賽人來說.
同樣的光卻是刺眼的.
他們想掩著眼睛不去看.
這個光是刺眼.
甚至是光到令他們都黑眼.
而今天我們要問自己的就是.
對於這個光 大家的反應是怎樣呢?.
我們渴望見到這個光.
還是我們覺得太刺眼.
要掩著眼睛呢?.
今天的講題是能看見的核子.
其實能看見的核子有兩個意思.
第一就是好像本來是失明的人.
他生來就是核子.
但因為遇見耶穌就能看見.
而另一個意思就像弗利賽人.
他們自己以為看得見.
但其實是一個核子.
看不見一些重要的事.
今天我們要問自己的就是.
到底我們是怎樣呢?.
到底我們是很開放很渴望神去改變你的生命.
還是我們覺得很OK.
不需要神的介入呢?.
有時候我們會慢了.
到一個地步是我們都不知道自己是慢了.
但正如一開始.
耶穌去醫治這個核子的時候.
其實不是這個核子主動要求的.
甚至耶穌都沒有問他.
你要不要得痊癒.
直接將泥塗在他的眼上.
甚至當這個核子被人趕出去的時候.

$^{681}$都是耶穌主動去找他.
今天或者我們是眼黑到一個地步.
其實我們真的不知道自己是眼黑.
我們看不清楚究竟我們看得到些什麼.
看不到些什麼.
我們根本就不知道.
但我們就一同謙卑到神的面前.
就是當神用他的手塗泥在你的眼睛的時候.
你就不去推開它.
當神叫你去池子洗的時候.
你就馴服.
求主去光照我們.
現在我們有一點禱告的時間.
今天一開始就說到.
其實今天是一個眼力的測試.
我們去檢視自己的生命.
我們屬靈的眼睛有沒有一些蒙查查的地方.
有沒有一些看不清楚的地方.
這一刻很想邀請弟兄姊妹.
我們真的重新去檢視自己.
我們有沒有一些時候.
可能明明一些神的工作.
神的恩典已經發生在我們面前.
但我們還是懷疑.
有時候我們會不會見到很多困難.
我們就不想做了.
我們就去逃避.
或者我們見不到困難背後的上帝.
或者有時候我們會不會很看重自己.
而看不到其他人的需要.
又或者有時候我們會不會對人都很難忘.
見不到這個世界的需要.
這一刻很想大家在神面前去安靜.
我們去檢視自己的生命.
有哪個地方需要神的光去光照.
我們一同再祈禱.
神我們在你面前.
弟兄妹在你面前.
我們渴望我們能夠看見你.
因為縱然我們有很多軟弱.

$^{721}$縱然我們生命有很多破碎的地方.
但神你沒有離棄過我們.
神你是那個愛我們到底的神.
你從來都不離棄我們.
所以當我們的生命有軟弱的時候.
你不是丟下我們.
你反而是來親近我們.
你渴望我們的生命有改變.
所以這一刻弟兄妹都會去到自己的生命去感恩.
縱然我們有破碎的時候.
縱然有得罪神的時候.
但我們仍然能夠回到教會裡.
我們仍然能夠和神傾訴.
這完全是神的恩典.
請祈禱.
祈主我們這一群人.
我們這個群體在你面前.
謙卑在你面前.
去承認我們真的有很多不足.
去承認我們真的有很多軟弱的時候.
但我們同時間在你面前.
我們都要去讚美你.
因為你是那個愛我們到底的神.
因為你是那個.
看到我們有軟弱去醫治的神.
所以主我們在你面前.
很坦誠地和你說.
我們真的很需要你的光照.
很需要神你的醫治.
很需要神你和我們一起.
去面對我們生命中的軟弱.
或者我們生活中遇到的困難.
所以主求你去幫助我們.
教我們我們真的不懂.
但神求你幫助我們的心.
能夠在你面前堅強起來.
以至於我們能夠去面對.
我們生命中的軟弱.
神求你的光去光照我們.
並且主當你的光去光照我們的時候.

$^{761}$幫助我們不去拒絕你的光.
不去假裝看不到.
所以神求你去建立我們.
今天在你面前的立志.
主多謝你.
祈禱奉主名求.
阿們.
\newpage



\section{腓立比書 1:3-11}
\label{sec:031818}
\textbf{竭力成為傳福音的團契}~~~~王永揚牧師
\newline
\newline
連結: \href{http://www.church.com.hk/acms/content.asp?site=cdc\&op=show\&type=product\&code=031818\&layout=sermon}{\texttt{031818}} ~~~~ 語音日期: 2026-03-15
\newline
\newline
\hyperref[sec:031817]{\small{< < < PREV SERMON < < <}}
~
\hyperref[sec:index_chronic]{\small{[返順時目]}}
~
\hyperref[sec:031819]{\small{> > > NEXT SERMON > > >}}
\newline
\newline
腓立比書 1:3-11
\newline
\begin{longtable}{cl}
\hline
\hline
章節 & 經文 (和合本修訂版)\\
\hline
1:3 & \begin{tabularx}{0.7\textwidth}{X} 我每逢想念你們，就感謝我的神， \end{tabularx} \\ \\ \relax
1:4 & \begin{tabularx}{0.7\textwidth}{X} 每逢為你們眾人祈求的時候，總是歡歡喜喜地祈求， \end{tabularx} \\ \\ \relax
1:5 & \begin{tabularx}{0.7\textwidth}{X} 因為從第一天直到如今，你們都同心合意興旺福音。 \end{tabularx} \\ \\ \relax
1:6 & \begin{tabularx}{0.7\textwidth}{X} 我深信，那在你們心裡動了美好工作的，到了耶穌基督的日子必完成這工作。 \end{tabularx} \\ \\ \relax
1:7 & \begin{tabularx}{0.7\textwidth}{X} 我為你們眾人有這樣的想法原是應當的，因為你們常在我心裡；無論我是在捆鎖中，在辯明並證實福音的時候，你們都與我一同蒙恩。 \end{tabularx} \\ \\ \relax
1:8 & \begin{tabularx}{0.7\textwidth}{X} 我以基督耶穌的心腸切切想念你們眾人，這是神可以為我作證的。 \end{tabularx} \\ \\ \relax
1:9 & \begin{tabularx}{0.7\textwidth}{X} 我所禱告的就是：要你們的愛心，在知識和各樣見識上，不斷增長， \end{tabularx} \\ \\ \relax
1:10 & \begin{tabularx}{0.7\textwidth}{X} 使你們能分辨是非，在基督的日子作真誠無可指責的人， \end{tabularx} \\ \\ \relax
1:11 & \begin{tabularx}{0.7\textwidth}{X} 更靠著耶穌基督結滿仁義的果子，歸榮耀稱讚給神。 \end{tabularx} \\ \\
[1ex]
\hline
\hline
\end{longtable}
$^{1}$在正道之前,我們同心來祈禱.
聖下獨生兒子耶穌基督,求你施恩.
可憐這個罪人.
若以在主你面前,施納宣講主的恩賢.
容讓耶穌基督的寶血.
再一次洗淨這個不配的人的罪.
容讓孩子能夠成為神話語的出口.
能夠將主恩賢解明.
容讓真理聖靈開啟我們心竅.
帶領我們進入真理.
恭敬將以下環節交託.
仰望求主保守恩寧.
祈禱,賴奉耶穌基督聖名祈求.
阿們.
在童話故事裡面.
王子打敗火龍之後.
從城堡裡面救出公主.
然後在群眾和民眾夾道歡迎之下.
揚長而去的婚姻典禮之後.
就是故事的結局.
從我的角度來看.
很少見到童話故事會講婚後的生活.
因為婚後的生活.
必然會遇上一些困難,挑戰.
可能會較較,會吵架.
這個就不是童話生活.
這個就破壞了畫風.
以婚禮作為結束.
因為當中婚後會有很多矛盾,衝突.
「開門七件事,柴米,油鹽,醬醋,茶」.
我們俗語講.
就是去「有情飲水飽」.
即是去一天可以.
一個星期可以嗎?.
一個月?.
一年可以有情飲水飽嗎?.
只飲水?.
不同的成長背景,教育程度.
生活習慣,價值觀等等.
美滿的婚禮.

$^{41}$不會自然而產生美滿的婚姻生活.
是要刻意經營.
彼此遷就.
互相接納,欣賞和體諒.
同樣,歷史中的教會.
經歷不斷的衝突,矛盾,不和,分裂.
失約,輸卷.
特別是「使徒行傳」.
很赤裸躺開地記述了當中的挑戰和困難.
彼得,巴拿巴,西拉,保羅等教會領袖.
是奉神差派.
來開展波瀾壯闊的宣教旅程.
是荊棘滿途,路漫漫,困難重重.
在「使徒行傳」15章.
來記述耶路撒冷會議.
激烈辯論,爭辯.
到底外邦信徒需不需要遵守摩西的律法.
特別是國禮,先靈道得救.
最後,長老阿葛一錘定音.
來到這裡說.
即是不用了.
但是要彼此尊重.
互不隸屬.
即是你有你的陽關路.
我有我的獨木橋.
耶路撒冷教會繼續向猶太人傳福音.
安提阿教會就去到外邦人當中作福音的工作.
大家河水不犯井水.
你有你做.
我有我做.
《使徒行傳》不遺言.
說到教會內部的矛盾.
外部面對猶太人的衝擊.
外邦人的質疑,嘲笑.
使徒保羅是以一個被告人的身份.
來被押解到羅馬.
來等候當時的羅馬皇帝去殺該.
來作最後的審訊.
理想中的教會是十全十美.
每個人都積極事奉靈明長進.

$^{81}$積極傳揚福音.
但現實是風潮中奮進的教會.
套用香港中文大學新亞書院校歌創校校長.
錢穆博士所寫的詞.
是「艱險我奮進」.
「困乏我多情彼此勉勵」.
這樣成為傳福音的團契.
都不是自然而然.
而是要刻意經營.
我們一起努力.
一起打拼.
在這個報道主日.
回應教會今年的主題.
「結連傳承復興」.
從保羅寫給菲律賓教會的文案.
感恩禱告.
我們一起努力.
竭力成為傳福音的團契.
今天講到會有兩個分題.
第一是將福音團契置於.
在我們當中信徒關係的核心.
第二是優先將福音置於.
你我禱告生活中.
將福音的團契置於.
我們信徒關係的核心.
剛才主席讀經文第三和第四節.
他講到「我每逢想念你們.
每逢為你們眾人祈求的時候.
嘗試歡歡喜喜地祈求.
因為從頭一天直到如今.
你們是同心合意.
慶王復任.
參考《使徒行傳》十六章.
保羅和巴拿巴是因為.
可不可以接納上次離隊.
中途脫隊的馬可.
在這個宣教旅程裡面.
馬可擅自離隊.
巴拿巴覺得給多個機會年輕人.
但保羅覺得一次不忠 百次不容.

$^{121}$於是兩個意見不合.
和平分手.
大道朝天 各走一邊.
你有你的楊關儂 我有我的獨木條.
各自傳揚福音.
保羅就帶同西拉提摩塔.
在聖靈的引導下.
轉折來到特羅亞.
在夜間有異象現於保羅.
有一個去馬其頓人站著.
來祈求請求保羅去過渡馬其頓.
我們又來到菲勒比.
見到一個在土國地方.
遇到一個賣紫色布匹的商販.
這個就是呂底亞.
帶領了對方信主.
更住在他的家.
繼續傳揚福音.
建立教會.
初信栽培.
門徒分裂.
但好景不常.
保羅被人誣告.
被群眾毆打.
被官差捉住.
剝了他的衣服.
亂棍毆他.
然後被收在監牢裡面.
在夜晚又餓又凍.
饑寒交迫的情況下.
保羅和西拉夜間歌唱.
神就作出請求.
更加帶領玉卒一家來歸主.
在身心受創的情況下.
保羅完成了福音工作.
獲悉之後回到呂底亞的家.
勉勵一番就走人了.
親眼見到自己屬靈的父親保羅.
這個人帶領我們信耶穌.
竟然因為信仰的緣故.

$^{161}$被人捉拿,被人毆打.
但肥納比教會一點都沒有後退.
他不會因為見到這個屬靈的父親.
這樣的狀況.
去遭難,遭害.
而放棄信仰.
肥納比教會沒有退後.
歸宿不聞世事.
自成一個圍爐取暖.
繼續積極傳揚福音.
茁壯成長.
更加是真金白銀地出錢出力.
支持保羅的福音工作.
在金錢上,禱告上.
和保羅一起同供.
所以保羅的思念是表達.
這份相交相知,肝膽相照.
縱使長江的路遙遙千里.
心仍然是相近的.
雖然是天各一方.
但仍然可以產生共鳴.
這就是保羅和肥納比教會的關係.
肥納比書一章講到的「想念」.
是回應四章十節.
講到「我靠主大大喜樂.
因為你們思念我的心.
如今又發心.
你們向來想念我.
只是沒得機會」.
甚至講到肥納比人.
「你們也知道我初傳福音.
來了馬其頓的時候.
輪到受受的事.
除了你們以外,並沒有別的教會供給我.
就算我再次貼近你們.
你們也一次,兩次打發人.
及供給我的需要」.
肥納比和保羅的關係不是擅自衝動.
不是去買張彩票.
或者抽獎券.

$^{201}$但之後就放兩天.
都不會考慮會不會中.
收工了,買了就做善事.
類似是這樣.
但保羅和肥納比教會的關係.
可以說是同心合意興旺福音.
不僅是肥納比教會和保羅.
是肥納比教會裡面的眾肢體.
都是同心合意的.
來參與保羅的福音工作.
保羅寫信給肥納比.
寫這封肥納比書.
是一封謝函.
從二章二十五到三十節表達.
書信是透過會友以巴弗提.
帶回去肥納比教會的.
這位教友帶來了教會的.
饋贈給獄中的保羅.
當時坐牢就不像今時今日.
可以去食用懲教署.
是政府的.
不是.
你吃的住的生活各樣.
都是自酬的.
沒有人供養.
你就餓死在獄中.
都沒有人理會,沒有人可憐的.
在這樣的狀況下.
肥納比教會來到去差派以巴弗提.
去到監獄裡面帶來他生活的需要.
更加幫忙照顧.
在獄中的保羅起居生活.
現在他病好了.
就打發以巴弗提回去肥納比.
變得教會的會友掛念.
肥納比教會透過不同的認識方法.
一直參與保羅的福音工作.
表達他們兩個的關係.
是從不間斷.
從開始直到如今.

$^{241}$同心合意.
原來的文字有「團契」的意思.
今日的語境.
「團契」一詞逐漸演變成為信徒之間.
親密友誼的關係.
一起飲茶灌水.
教兒相聚.
一起查經.
一起畏罪.
但在保羅的年代.
這個詞通常帶有商業意義.
是這個「Partnership」.
即我們是一起合作的.
用公司法來說.
「Partnership」是無限公司.
是上身的.
在債務來說.
如果是至死方休.
是個人破產的.
不是一間公司破產那麼簡單.
對於「Partnership」來說.
即公司化來說.
就是股東是上身的.
一間公司出事.
股東是有法律責任去完成債務.
不像現在有限公司.
為甚麼有限公司倒閉.
說倒閉又如何.
債權人清盤.
公司沒錢又如何.
與業權人無關.
是一個這樣的關係.
來生死與共.
一榮一雪.
大家都是有份和出事.
就好像去一間食店.
由家人來經營.
三兄弟三哥大嫂二哥二嫂三哥三嫂.
一齊合作經營一間食店.
大哥大嫂天未亮.

$^{281}$就去市場買材料.
預備去地方廚房準備食材等等.
做完早餐.
午市前二哥二嫂出現.
一齊預備午市.
招呼了第一輪午市之後.
大哥大嫂就退場回家.
二哥二嫂接手.
做完午市做完下午茶.
準備晚市的時候三哥三嫂就出場.
一齊經營晚市.
去到第一輪晚市完之後.
二哥二嫂又退場回家.
剩下三哥三嫂做完下面.
去十一點打烊.
清理,執拾,結帳,計算.
十二點多收工回家.
第二天早上大哥大嫂又來.
這樣合作無間去經營這盤生意.
同樣保羅,菲律賓教會.
和菲律賓教會的信徒.
都是這樣參與和合作.
在傳福音的事情上.
大家一齊努力.
真正團體的核心.
是在於有共同的意象.
而甘願犧牲自己.
我們基督徒的團體就應該.
為到福音來甘願捨己.
其中我們應該會流露.
這份溫暖和親密.
應該人走過來會覺得.
我們是大家互相接納,互相欣賞.
但我們的核心是應該超越一切.
我們應該有共同意象.
擁抱同一個願景.
這個意象是呼召我們去毀身在這裡.
就正如這個麻雀同樂會.
這班人是怎樣呢?.
風雨不改就打麻雀.

$^{321}$歌迷會不說名字了.
不知道是哪個歌星或影星.
這班人是無常地支持的.
大家因為想捧這個人,明星或歌星.
付費成為歌迷會的會員.
積極籌劃活動,參與購買宣傳品,打氣等等.
是無常的,但他覺得於有榮焉,參與其中是一份光榮的喜樂.
香港有不同的關注組,燒賣關注組.
你網上會看到的.
有些人會走遍香港各地.
不是各城各鄉.
看看哪裡的燒賣好吃,評分.
怎樣怎樣各樣.
但這些人是義工來的.
不會獲利.
是參與,但他樂此不疲.
是介紹,推銷燒賣的文化.
因為是擁抱共同的意向.
在第十及十一節.
使你們能夠分辨是非,作誠實無過的人.
直到基督的日子.
並靠著耶穌基督結滿人兒的果子.
叫榮耀稱讚貴與臣.
福音團契促進團契成員成長.
結出人兒的果子.
果子可以說是信徒的生命亮麗成長與臣結連.
但同時也是流露日常的待人接物.
耶穌說「好樹就結好果子」.
從「果子」又認出是主的門徒.
耶穌在撒種的比喻裡說到.
耶穌呼籲聽眾結出百倍的修成.
彼得比教會繼續作見證.
我們堅持不懈地為保羅祈禱.
我們捐款支持保羅的方工作.
一起見證他們在福音的重要性.
以及承認福音是優先,是首選.
大家都有這樣的共識.
這也足以成為保羅感恩的理由.
保羅對福音的委托.
以及彼得比教會的見證提醒我們.

$^{361}$福音的團契,福音的夥伴關係.
必須放在我們和其他信徒關係網絡的核心.
即是我們為甚麼走在一起呢?.
崇拜之後或團契.
我們午餐會談些甚麼?.
我們當然會說一些客套的話.
「天氣,時事,體育消息,工作,子女,彼此的傷痛」.
這些都是好的.
我們是「借去不離地」.
是有談及生活的.
但在分享各個層面的時候.
我們是甚麼將我們基督徒連繫在一起?.
這應該是對福音的熱誠.
是福音中的團契.
甚麼是走在一起呢?.
我們不只是圍爐取暖.
我們不只是找在一起就真土.
而是我們的目的和意義.
是為了福音.
第二.
是優先將福音置於我們的禱告生活當中.
在書信的起手.
「為到菲律賓的教會感恩祈禱」.
然後九到十一節這樣說.
都是祈求的內容.
「我所禱告的.
就是你們的愛心在各人的知識和見識上多而由多.
使你們能夠分辨是非.
作誠實無過的人.
直到基督的日子.
並靠著耶穌基督結滿了人兒的果子.
叫榮耀歸於神」.
我們要看回保羅當時的情況,情景.
我們先寫這封書信.
我們先看到那份味道.
保羅是做了甚麼?.
說得難聽一點.
當時是在坐牢的.
保羅是要靠菲律賓的接濟.
他才能繼續在獄中生活.

$^{401}$保羅身在獄中未知結果.
不知道會不會上訴成功.
又不知道會坐多久牢.
又不知道何時獲釋.
等到何時何年何月何日.
這個坐牢的情況.
甚麼都不知道.
但他所關心的.
不是自己.
而是菲律賓教會的愛心.
有沒有振蕩.
這份愛沒有特定的對象.
當然是講到信徒之間彼此相愛.
這是從主耶穌的新命令.
要彼此相愛.
亦包括對未信者的愛.
有人呼音十章.
耶穌說好牧人的比喻.
「愛我的羊」.
但同時耶穌說.
「圈外有羊,他所關」.
耶穌所關注的不只是我們.
也關注其他圈外的羊的需要.
再者.
他指的是理解神的旨意.
菲律賓教會也希望我們能夠洞察認識.
上帝神的話語和道路.
然後知道如何來回應.
按著神的心意來生活.
按著知識增多.
應該是帶來更深的悔悟.
是希望我們的生活能夠去活出.
神在我們生命上的旨意.
祂要求的是我們減少自我中心.
從積極來看.
是要我們有更多的愛心.
去對神的愛和對別人的愛.
這就是保羅所求.
保羅沒有為自己祈求.
反而在第四章.

$^{441}$他打發伊巴菲提回去時說.
「樣樣都有」.
我不是因為缺乏來講這番說話.
明明是一無所有.
但保羅說他自己「樣樣都有」.
因為他很確定一切都是從神而來.
他已經接受了菲律賓教會的饋贈.
他所關注的就是菲律賓教會.
在外面的信徒是否仍然熱心追求傳福音.
他所關注的就是菲律賓教會的不和.
因為今天篇幅所限和時間所限.
每次都被人說我超時.
我今天要收手了.
其實這封信除了答謝之外.
腓立比書是一封勸教的書信.
正正是他們不同心.
教會裡面產生不和.
在第四章保羅指名道姓地說.
去勸喻「有瓦碟」和「秦都記」.
這是新標點胡學本的翻譯.
就是「有瓦碟」和「秦都記」.
在主裡面同心.
保羅希望菲律賓教會繼續傳揚福音.
我們稍加思考就會發現.
保羅所祈求的都是以福音為中心.
這是充滿福音氣息的祈禱.
就是說這個禱告是為了在菲律賓的生命裡面.
能夠推展福音的工作.
透過福音的果子結出在他們的生命裡面.
這個祈求就能夠最終落實.
能夠去榮耀神.
今天我們的禱告生活.
有多少是為自己.
有多少是為他人.
我們上一次為其他人禱告是何時呢?.
我們有沒有仔細地為弟兄姊妹成長來祈禱.
以致有足夠的空間和闊度為主所使用?.
當然我們應該可以為自己需要祈禱.
但是福音的焦點在哪裡?.
通過保羅書信去抄錄保羅的祈禱.

$^{481}$問問自己.
保羅為了甚麼來祈禱?.
保羅從沒有在書信裡面提出.
去要求為他自己來祈禱.
就算有.
都是為了信徒的需要.
而作出這樣的呼籲.
保羅的祈禱教導我們怎樣祈禱.
保羅的示域和關注怎樣提醒我們.
將福音放在第一位.
這意味我們將福音的優先.
放在我們禱告的中心.
教會是為他者如存在.
我們的設立.
不只是因為「我們在這裡真好」.
停留在這裡就完結.
我們應該是在這裡得到牧養.
得到鼓勵.
更新.
然後我們回到我們所在之處.
繼續事奉和榮耀神.
以致我們的生命來到給人看到活潑的見證.
讓我們的家人,親友,同事,鄰舍.
從我們見到耶穌.
然後他們覺得.
嘩!這個神是真實的.
然後對福音產生興趣.
甚至是被邀約回來教會.
然後信主.
然後接受栽培,生命成長.
到他們再邀請他們的親友,鄰舍,朋友,家人,同事.
再回來教會.
這是罪善.
我們是有罪,有善.
但我們不是為了存在而存在.
不是為了繼續而繼續.
我們應該有一個使命來到一起努力.
為福音同心努力.
讓我們來到將福音放置在我們信徒關係網絡的核心.
亦都是放置在我們禱告生活的核心.

$^{521}$福音的必然性無可取替.
我向大家呼籲.
先賣廣告.
參與我們報道部今年的福音聚會.
我是張powerpoint的.
希望大家來到繼續守望.
請你考慮怎樣參與和回應.
這麼多禮物和表達我們很豐富.
總有一件事是你能夠回應和參與的.
我們在4月19日和7月26日有教會的福音主日.
9月6日有破鏡重圓教會的報道會.
我們有101工作坊.
我們在4月12日在街頭派發單張.
去宣傳呼籲人參與福音主日.
我們有啟發課程.
有科務員工作組.
接下來月底就去禮敬懲教所做監獄報道.
5月就去禮敬懲教所做監獄報道.
3月底就去伊利沙伯醫院.
就去床邊關懷報道等等.
其實都是總有一件事你可以參與到.
就算你未能夠親身參與.
都請你來禱告.
為我們守望.
盼望即今日不是聽完就完結.
讓我們來到去在今年的裡面.
我們一齊去立志和回應.
會不會想一下.
我可以向多少人傳福音.
信不信是對方的責任.
神的計劃.
但如果我們不傳.
對方就失卻信主的機會.
神的救恩.
你我有責.
讓我們一同來回應.
我們同心來祈禱.
我們的主我們神.
多謝你.
因為你有數不盡的恩典.

$^{561}$你更加將傳福音的責任.
義務.
來交付給我們.
讓我們今天來聆聽主的說話.
透過菲律賓教會.
他們生命的見證.
去激勵我們.
提醒我們.
有一個真實的回應.
我們如此禱告.
奉主明禱.
阿們.
\newpage



\section{創世記 37:12-36}
\label{sec:031819}
\textbf{聖經人聖經事:上主的手在帶領}~~~~黃克勤牧師博士
\newline
\newline
連結: \href{http://www.church.com.hk/acms/content.asp?site=cdc\&op=show\&type=product\&code=031819\&layout=sermon}{\texttt{031819}} ~~~~ 語音日期: 2026-03-14
\newline
\newline
\hyperref[sec:031818]{\small{< < < PREV SERMON < < <}}
~
\hyperref[sec:index_chronic]{\small{[返順時目]}}
~
\hyperref[sec:031820]{\small{> > > NEXT SERMON > > >}}
\newline
\newline
創世記 37:12-36
\newline
\begin{longtable}{cl}
\hline
\hline
章節 & 經文 (和合本修訂版)\\
\hline
37:12 & \begin{tabularx}{0.7\textwidth}{X} 約瑟的哥哥們到示劍去放他們父親的羊。 \end{tabularx} \\ \\ \relax
37:13 & \begin{tabularx}{0.7\textwidth}{X} 以色列對約瑟說：「你哥哥們不是在示劍放羊嗎？來，我派你到他們那裡去。」約瑟對他說：「我在這裡。」 \end{tabularx} \\ \\ \relax
37:14 & \begin{tabularx}{0.7\textwidth}{X} 以色列對他說：「你去看看你哥哥們是否平安，羊群是否平安，再回來告訴我。」於是他派約瑟出希伯崙谷，約瑟就往示劍去了。 \end{tabularx} \\ \\ \relax
37:15 & \begin{tabularx}{0.7\textwidth}{X} 有人遇見他，看哪，他在田野走迷了路。那人問他說：「你找甚麼？」 \end{tabularx} \\ \\ \relax
37:16 & \begin{tabularx}{0.7\textwidth}{X} 他說：「我找我的哥哥們，請告訴我，他們在哪裡放羊。」 \end{tabularx} \\ \\ \relax
37:17 & \begin{tabularx}{0.7\textwidth}{X} 那人說：「他們已經離開這裡走了，我聽見他們說：『我們往多坍去。』」約瑟就去追哥哥們，在多坍找到了他們。 \end{tabularx} \\ \\ \relax
37:18 & \begin{tabularx}{0.7\textwidth}{X} 他們遠遠看見他，趁他還沒有走近他們，就圖謀要殺死他。 \end{tabularx} \\ \\ \relax
37:19 & \begin{tabularx}{0.7\textwidth}{X} 他們彼此說：「看哪！那做夢的來了。 \end{tabularx} \\ \\ \relax
37:20 & \begin{tabularx}{0.7\textwidth}{X} 現在，來吧！我們把他殺了，丟在一個坑裡，就說有惡獸把他吃了。我們且看他的夢將來怎麼樣。」 \end{tabularx} \\ \\ \relax
37:21 & \begin{tabularx}{0.7\textwidth}{X} 呂便聽見了，要救約瑟脫離他們的手，說：「我們不可害他的性命。」 \end{tabularx} \\ \\ \relax
37:22 & \begin{tabularx}{0.7\textwidth}{X} 呂便又對他們說：「不可流他的血，可以把他丟在這曠野的坑裡，不可下手害他。」呂便要救他脫離他們的手，把他還給他父親。 \end{tabularx} \\ \\ \relax
37:23 & \begin{tabularx}{0.7\textwidth}{X} 約瑟到了他哥哥們那裡，他們就剝去他的外衣，就是他身上那件長袍。 \end{tabularx} \\ \\ \relax
37:24 & \begin{tabularx}{0.7\textwidth}{X} 他們抓住他，把他丟在坑裡。那坑是空的，裡頭沒有水。 \end{tabularx} \\ \\ \relax
37:25 & \begin{tabularx}{0.7\textwidth}{X} 他們坐下吃飯，舉目觀看，看哪，有一群以實瑪利人從基列來，用駱駝馱著香料、乳香、沒藥，要帶下埃及去。 \end{tabularx} \\ \\ \relax
37:26 & \begin{tabularx}{0.7\textwidth}{X} 猶大對他的兄弟們說：「我們殺我們的弟弟，遮掩他的血有甚麼好處呢？ \end{tabularx} \\ \\ \relax
37:27 & \begin{tabularx}{0.7\textwidth}{X} 來，我們把他賣給以實瑪利人，不要下手害他，因為他是我們的弟弟，我們的骨肉。」他的兄弟們就聽從了他。 \end{tabularx} \\ \\ \relax
37:28 & \begin{tabularx}{0.7\textwidth}{X} 那時，有些米甸的商人從那裡經過，就把約瑟從坑裡拉上來。他們以二十塊銀子把約瑟賣給以實瑪利人，他們就把約瑟帶到埃及去了。 \end{tabularx} \\ \\ \relax
37:29 & \begin{tabularx}{0.7\textwidth}{X} 呂便回到坑旁，看哪，約瑟不在坑裡，就撕裂自己的衣服， \end{tabularx} \\ \\ \relax
37:30 & \begin{tabularx}{0.7\textwidth}{X} 回到他兄弟們那裡，說：「孩子不在了。我往哪裡去才好呢？」 \end{tabularx} \\ \\ \relax
37:31 & \begin{tabularx}{0.7\textwidth}{X} 於是，他們宰了一隻公山羊，拿了約瑟的那件外衣染上了血， \end{tabularx} \\ \\ \relax
37:32 & \begin{tabularx}{0.7\textwidth}{X} 派人把長袍送到他們的父親那裡，說：「我們發現這個， 請認一認，是不是你兒子的外衣？」 \end{tabularx} \\ \\ \relax
37:33 & \begin{tabularx}{0.7\textwidth}{X} 他認出來，就說：「這是我兒子的外衣，惡獸把他吃了，約瑟一定被撕碎了！」 \end{tabularx} \\ \\ \relax
37:34 & \begin{tabularx}{0.7\textwidth}{X} 雅各就撕裂衣服，腰間圍上麻布，為他兒子哀傷了多日。 \end{tabularx} \\ \\ \relax
37:35 & \begin{tabularx}{0.7\textwidth}{X} 他的兒女都起來安慰他，他卻不肯受安慰，說：「我必哀傷著下陰間，到我兒子那裡。」約瑟的父親就為他哀哭。 \end{tabularx} \\ \\ \relax
37:36 & \begin{tabularx}{0.7\textwidth}{X} 米甸人把約瑟賣到埃及，給法老的官員，就是護衛長波提乏。 \end{tabularx} \\ \\
[1ex]
\hline
\hline
\end{longtable}
$^{1}$說這些故事是很開心的.
昨天說了藥室的採醫與發夢.
今天的聖經人聖經事生命說書網絡為大家預備的.
就是上主的手在帶領.
我們看聖經的人或聖經的事.
最重要看到背後上帝的手在做什麼.
昨天藥室又有採醫又發夢.
上帝的手你看到什麼.
上帝的手是帶領藥室.
他有一個出生的困境是什麼呢.
小,而且爸爸疼愛,有件採醫,害死他了.
他又不會做,不會死,不斷說哥哥的壞話.
還有一點點直和高傲.
你向我下拜,爸媽都向我下拜.
我的夢是這樣的,你們要拜我的.
你不是說哪位,排第幾.
如果以兄弟來說,你排第十一.
不要以為你很厲害,不是一,是十一.
所以哥哥不喜歡他.
不過上帝的手一直帶著藥室在他身上發生很多事.
我們看到上帝的手在藥室身上做什麼.
創世記三十七章,繼續說創世記十二節.
藥室的哥哥往事件去放他父親的羊.
以色列就是雅各對藥室說.
你哥哥們不是在事件放羊嗎.
你來,我要打發你往他們那裡去.
藥室說我在這裡.
藥室這裡都很乖,爸爸說藥室你來.
藥室說我來,我在這裡.
蠻有趣的.
藥室的哥哥去了事件放父親的羊.
藥室有沒有去呢.
沒有,藥室不用去.
你說哥哥是不是很生氣.
真的,爸爸真的很偏心.
我們出來日曬雨淋,辛苦到極.
也沒有一口好吃.
有時候跟野獸搏鬥.
他就好,不用出來.
留意一下,是上帝的工作.

$^{41}$是令到雅各的心很不安.
不安就是神的手在後面做事.
怎樣不安呢.
聽說你的哥哥去了事件放羊.
你跟我去看看他們.
藥室都孝順的.
在這裡,爸爸叫我一定做的.
以色列雅各就說.
留意一下,你去看看你的哥哥們.
平安不平安.
群羊平安不平安.
然後回來報信給我.
這樣很有趣,有點無稽.
你去看看他們平不平安.
我掛著他們,看他們平不平安.
看看羊平不平安回來告訴我.
其實是不必的.
我認為沒有什麼需要.
也沒有實效的.
因為如果兒子有一個損傷.
有些羊被野獸吃了或者被人搶了.
十個兒子,十個哥哥.
一定會有人回來報信.
而且不只是十個.
有些牧羊人跟著去.
一定有人回來報信.
其實不用去找他們.
他就說你去找他們.
我真的掛著他們.
看他們平安不平安.
然後回來報信.
我想知道.
其實會不會.
雅各也想藥室出去闖蕩呢.
應該不會.
因為如果他真的想出去.
叫藥室出去看羊.
認識一下曠野生活.
跟一些哥哥去照顧.
怎會叫他一個人孤零零在曠野裡走.

$^{81}$其實很危險.
雅各一定知道.
不過他心裡可能是神令他心不安.
於是他忍不住.
沒有人了.
其實可以叫其他人.
他沒有叫.
於是找一個心腹.
藥室你來和我出去找他們.
告訴我他們好不好.
平不平安.
於是打發他出去.
出希伯倫谷去視劍.
藥室走不慣路.
走到差不多就走迷了路.
其實都幾難.
今天又沒有北斗導航.
又沒有衛星導航.
又沒有地圖.
靠自己方向感.
看看太陽在哪邊.
就知道什麼方向感.
向著那邊就去那邊.
向著那邊希伯倫那邊就視劍.
向著那邊就去那邊.
他就去了.
在田野走迷了路.
好在上帝的工作.
上帝的手在那裡工作.
有人遇見他在田野走迷了路.
遇見他.
他們如果是好人好者.
大家在曠野遇見都會見個面.
你好嗎?哪裡來的?.
有沒有什麼要幫忙?.
有些水你的口很乾.
給你喝吧.
他們很幫忙.
大家知道在曠野生活.
田野間是很難的.

$^{121}$走迷了路.
又沒有路牌.
又沒有電話.
於是有人遇見他.
喂!童子!小朋友!.
為什麼你一個人在這裡?.
我去找我哥.
我去視劍.
另一個就走迷了路.
不知哪一邊.
你教我走哪一邊?視劍.
那些人就好好幫忙.
你找什麼?視劍.
你去視劍找什麼?.
哦!我去找我哥.
求你告訴我們.
他們在何處放羊?.
這個問題是不是很搞笑?.
純粹是一個生手.
沒有牧過羊的人.
沒有走過曠野.
若實話我去找我哥.
很多人的.
你告訴我.
他們在何處放羊?.
喂!在這個田野.
曠野中間.
在何處放羊?.
你哥在何處放羊?.
我怎麼知道?.
誰是你哥?.
我又沒有寫著你哥.
我怎麼知道誰是你哥?.
他不知道的.
但他就在那裡.
我哥我找不到.
即是小朋友走失了.
你找誰?.
我找我媽.
你媽媽是怎樣的?.

$^{161}$有眼耳口鼻的.
有頭髮的.
這樣不行嗎?.
我找我哥.
求你告訴我們.
他們在何處放羊?.
這麼沒sense的問題.
誰知道.
上帝的手真厲害.
那些人就說.
哦!那些人我知道.
我剛剛遇到他們.
他們已經走了.
不在市劍了.
而且我們見過面.
我聽到他們說.
幸好聽到.
我聽到那些哥說話.
遇到的人聊天的人.
剛剛又遇到藥室.
如果不是真的不知道怎麼辦.
我聽見他們說要往多山.
多山哪邊?.
那邊那邊.
於是藥室就說.
我去追他們.
你追快點.
因為他們走了很久了.
都不知道是不是昨天見到的.
藥室就追趕他哥哥.
就遇見他們真的在多山.
跟著這個地方追過去.
遇見他們了.
有群人在前面.
有這麼多人就像我哥哥.
他們就走過去了.
他們遠遠見到他.
為什麼一個人在那裡走.
穿彩色衣服.
藥室來的.

$^{201}$看清楚一點.
他們很遠見到他.
十個哥就說.
趁他還沒有來到跟前.
不如我們想想同謀.
怎樣弄死他.
在曠野.
這麼狠毒.
所以你見到這個哥是不是很恨他.
再做夢之後.
越發恨他就想他死.
他們就彼此說.
那個做夢的人來了.
我們把他殺了.
扔進坑裡.
算了.
回去說.
有惡獸把他吃了.
看看他還做什麼夢才行.
且看他的夢將來怎樣.
弄死他.
我看你做什麼夢.
其實你見到十個哥.
是很憎恨他.
其中有些特別憎恨他.
就是雅各那兩個妾的兒子.
特別憎恨他.
就說你好像半個外邦人.
特別憎恨他.
裡辯聽見了.
裡辯是誰呢.
裡辯以前中文翻譯是樓辯.
即是說大哥.
雅各的第一個兒子.
樓辯不要聽.
這麼難聽.
於是譯叫裡辯.
不過裡辯也挺有趣.
為什麼不譯出辯呢.
裡辯.

$^{241}$裡辯聽見.
大哥樓辯裡辯很深.
要救他脫離他們的仇.
我們不要害他性命.
不要流他的血.
不如這樣吧.
他又不能夠譯那九個弟弟的意思.
那怎麼好呢.
不要流他的血.
我們將他扔在坑裡面.
不要親手害死他.
扔在坑裡面即是什麼.
爬不上來.
扔在坑裡面會怎樣.
由他裡面慢慢嘔死.
餓死 乾死.
死也是.
那九個弟弟會不會喜歡呢.
都喜歡.
不過有什麼好處呢.
就是我們不用親手動手殺他.
不要自己親手動手.
由他餓死渴死.
明白嗎.
於是就不要親手下去殺他.
這個裡面的意思是要救他.
裡面是要救他.
神的手安排了裡面是要救他.
這樣就要救他.
不要動手.
其實他想什麼呢.
要救他脫離他們的仇.
把他歸還他的父親.
裡面是想要救藥室.
我們先扔他進坑裡面.
讓他在那裡嘔死.
我們走吧.
不過可能有一個屈尾十.
就叫僕人救他上來.
帶他回家.

$^{281}$這樣也不一定.
又或者是回到家裡.
哎呀 藥室呢.
不知道在哪裡.
知不知道在哪裡.
沒有見到我叫他找你們.
沒有見到嗎.
沒有.
然後可能靜靜跟爸爸說.
爸爸 事情是這樣的.
他現在在哪裡.
東南西北多山附近的地方.
有一個洞.
你快點叫僕人去救他.
不然趕不及.
還給他父親救他.
藥室到他哥哥那裡.
他們就抓住他.
就脫了他的彩衣.
然後就拿著他的衣服.
然後就穿著他的內衣.
把他推到坑裡.
好在那坑是空的.
裡頭沒有水.
好在是空的.
即是說是乾的.
又沒有野獸.
又沒有蛇裡面.
又沒有水.
如果不是有水.
掉下去淹死.
或者弄著半缸水.
你都幾慘.
他就是這樣.
於是就扔了下去.
在那個坑裡面是空的.
又沒有水.
整個弟弟都在那裡工作.
第25節.
他們坐下吃飯.

$^{321}$剛剛到那裡吃飯的時候.
吃著吃著.
見到那邊有駱駝.
一隊人經過.
駱駝隊.
甚麼人呢.
看他們穿的衣服.
像米電人.
是米電地方名.
其實他們的種族是.
以實馬尼.
都是半個親戚.
以實馬尼即是說.
是亞伯拉罕的另一個兒子.
所以和雅各這班兒子.
都是有少少關係.
都是半個親戚.
喂喂喂.
你們是哪裡來的.
我們是在基列來的.
現在用駱駝.
是帶著香料.
乳香 抹藥.
去埃及買.
基列是很出名.
出乳香的.
他們帶著香料.
去埃及那裡買.
親戚親戚.
於是大家就說起來了.
尤大就說.
不如這樣吧.
和其他的兄弟說.
我們殺了我們的兄弟.
藏了他的血有甚麼益處呢.
如果我們殺了我們的弟弟.
留了血.
我們本身做了一件不好的事.
壞事.
有甚麼好呢.

$^{361}$我們隱藏著這件事.
都有血.
我們都變了殺人兇手.
不好啊.
不如這樣吧.
我們賣了他給以實馬尼人.
賣給親戚.
我們不要下手害他.
賣了他.
又撿回多少.
一來大家兄弟骨肉.
不要相殘.
買了他撿回多少.
大家買酒喝.
聊聊天.
吃吃東西.
那些兄弟就說.
這樣也好啊.
其實也好啊.
殺了弟弟也不好啊.
我怕神懲罰我們.
那不好了.
賣了他也好啊.
那就好.
賣了.
那些米地商人經過.
於是這個哥.
就把他從坑裡面拉上來.
把繩子拉上來.
二十四十萬銀子.
賣了給以實馬尼人.
他們就把約瑟帶去埃及.
你見到.
是不是整個手都在工作.
預備好.
又遇到這個.
又遇到那個.
哎呀.
我的哥在哪裡.
我怎知道你的哥在哪裡.

$^{401}$哎呀在哪裡.
哦.
快點去追.
去了多山.
於是整件事一路下來.
我每一個紅色字.
都是告訴你.
神的手在工作.
於是約瑟就開始去了埃及.
裡面可能沒有和他一起吃飯.
走開了.
大哥裡面回來.
o下.
為什麼約瑟不在那個洞裡.
去了哪裡.
糟了.
是不是被人擒了.
怎麼辦.
那就快點走回兄弟那裡.
喂.
弟弟那個小孩不見了.
那個童子的小孩.
那個小孩不見了.
怎麼辦.
這樣.
唉.
糟了.
裡面已經想了一段.
他以為不見了.
其實哥哥沒有跟他說.
他走了出去.
回來後.
哎呀.
找弟弟不見了.
可能他也看看有沒有飯吃.
照顧他.
不見了怎麼辦.
他們就把衣服脫下.
扔在坑裡.
不如這樣吧.

$^{441}$我們就找一隻公羊殺了.
然後.
坐在這裡染血.
然後回去.
告訴爸爸.
打發人.
即是埋伏人.
就送去父親那裡.
哎呀.
我們在田野撿到一件衣服.
你承認.
是不是弱勢的.
這樣.
雅各就承認.
哎呀.
真是我做給弱勢那件衣服.
有血.
撿回一件衣服.
又爛了.
當然有惡獸將它吃了.
弱勢被撕碎了.
撕碎了.
弱勢被野獸吃了.
散了.
碎了.
哎呀.
撕碎了.
講兩次.
很慘.
雅各就撕裂衣服.
腰間圍上麻布.
為他的兒子悲哀多日.
他的兒女都來安慰他.
不肯受安慰.
爸爸不要緊.
你還有十一個兒子.
還有很多個女兒.
不怕的.
不要緊的.
神都是有保守.

$^{481}$有帶領的.
弱勢說.
哎呀.
雅各.
雅各說.
怎麼辦.
不行的這個兒子.
你們都是很乖.
哎呀.
我沒有兒子了.
哎呀.
我快死了.
我很老了.
我悲哀著下.
陰間.
哎呀.
我要去找我兒子了.
到我兒子大了.
哎呀.
我要悲哀死.
我要去找我兒子.
弱勢的父親.
雅各.
為他哀哭.
好了.
還未完.
米殿人帶弱勢去埃及.
將他賣給法老的內臣.
這個內臣.
只是說.
他很親密的一個.
很信任的大臣.
就是護衛長波提弗.
於是.
弱勢就去了波提弗的家.
開始接近了皇室.
又有另外一番遭遇.
這一天.
我今天想和大家看的就是.
神的手在工作.

$^{521}$所以紅色的字.
就是神的手在工作.
弟子會.
如果你發覺.
為什麼我有這麼多.
過山車經歷人生.
裡面有高有低.
其實是神的手在你的身上有工作.
祂一路預備你.
模造你.
令到你適合做某一樣東西.
有祂的心意.
或者一直都這樣說.
神在每一個人身上有個計劃.
我自己都不知道.
我自己回望才發覺.
原來神在我身上有個計劃.
原來我讀了這些書.
是有這樣的用途.
讀了這些書有這樣的用途.
讀了這個.
國內上海財經大學的書.
令到我在國內行走.
別人見到我在中國大學的.
高級的學位就覺得.
嘩 很厲害.
嘩 馬上信任了.
其實這件事令到我在中國的工作是更加好的.
其實是這樣發生的.
很多事情都是大家這樣去建立的.
其實背後都是神的手在動工.
因此將自己交託給主.
遇到什麼事情就求問神.
是不是你要我做的.
如果是的話.
沒有反對.
我做.
其實到遲一點我們就會知道.
原來神這樣是對我有益的.
我們所有事都在祂手裡面.

$^{561}$我終身的事都在你手中.
我交託給主一起祈禱.
求主幫助我們.
讓我們更加懂得原來你的手.
到今天仍然在世界上工作.
從來沒有停.
你是色在 金在 以後永在的主.
你雖然在這高天之上.
但是有時間聽每一個小朋友的祈禱.
你在高天之上.
但是你是可以管理.
控制著這個世界.
所有每一件或大或小的事情發生.
甚至我們的每一步.
都是主能夠帶領著.
而且是帶領著.
多謝我們的主.
讚美你.
你的手成就工作.
而且大能.
你運籌為握.
真的很厲害.
願你得著榮耀.
我們信靠你 信任你.
祈禱奉耶穌名求.
阿們.
\newpage



\section{約翰福音 6:35-59}
\label{sec:031820}
\textbf{榮耀君王:永不飢渴}~~~~林鍵牧師
\newline
\newline
連結: \href{http://www.church.com.hk/acms/content.asp?site=cdc\&op=show\&type=product\&code=031820\&layout=sermon}{\texttt{031820}} ~~~~ 語音日期: 2026-03-14
\newline
\newline
\hyperref[sec:031819]{\small{< < < PREV SERMON < < <}}
~
\hyperref[sec:index_chronic]{\small{[返順時目]}}
~
\hyperref[sec:031821]{\small{> > > NEXT SERMON > > >}}
\newline
\newline
約翰福音 6:35-59
\newline
\begin{longtable}{cl}
\hline
\hline
章節 & 經文 (和合本修訂版)\\
\hline
6:35 & \begin{tabularx}{0.7\textwidth}{X} 耶穌對他們說：「我就是生命的糧。到我這裡來的，絕不飢餓；信我的，永不乾渴。 \end{tabularx} \\ \\ \relax
6:36 & \begin{tabularx}{0.7\textwidth}{X} 可是，我告訴過你們，你們已經看見我，還是不信。 \end{tabularx} \\ \\ \relax
6:37 & \begin{tabularx}{0.7\textwidth}{X} 凡父所賜給我的人，必到我這裡來；到我這裡來的，我總不丟棄他。 \end{tabularx} \\ \\ \relax
6:38 & \begin{tabularx}{0.7\textwidth}{X} 因為我從天上降下來，不是要按自己的意願行，而是要遵行差我來那位的旨意。 \end{tabularx} \\ \\ \relax
6:39 & \begin{tabularx}{0.7\textwidth}{X} 差我來那位的旨意就是：他所賜給我的，要我一個也不失落，並且在末日使他復活。 \end{tabularx} \\ \\ \relax
6:40 & \begin{tabularx}{0.7\textwidth}{X} 因為我父的旨意是要使每一個見了子而信的人得永生，並且在末日我要使他復活。」 \end{tabularx} \\ \\ \relax
6:41 & \begin{tabularx}{0.7\textwidth}{X} 猶太人因為耶穌說「我是從天上降下來的糧」，就私下議論他， \end{tabularx} \\ \\ \relax
6:42 & \begin{tabularx}{0.7\textwidth}{X} 說：「這不是約瑟的兒子耶穌嗎？我們豈不認得他的父母嗎？現在他怎麼說『我是從天上降下來的』呢？」 \end{tabularx} \\ \\ \relax
6:43 & \begin{tabularx}{0.7\textwidth}{X} 耶穌回答，對他們說：「你們不要彼此私下議論。 \end{tabularx} \\ \\ \relax
6:44 & \begin{tabularx}{0.7\textwidth}{X} 若不是差我來的父吸引人，就沒有人能到我這裡來；到我這裡來的，在末日我要使他復活。 \end{tabularx} \\ \\ \relax
6:45 & \begin{tabularx}{0.7\textwidth}{X} 在先知書上寫著：『他們都要蒙神教導。』凡聽了父的教導而學習的，都到我這裡來。 \end{tabularx} \\ \\ \relax
6:46 & \begin{tabularx}{0.7\textwidth}{X} 這不是說有人看見過父，惟獨從神來的，他才看見過父。 \end{tabularx} \\ \\ \relax
6:47 & \begin{tabularx}{0.7\textwidth}{X} 我實實在在地告訴你們，信的人有永生。 \end{tabularx} \\ \\ \relax
6:48 & \begin{tabularx}{0.7\textwidth}{X} 我就是生命的糧。 \end{tabularx} \\ \\ \relax
6:49 & \begin{tabularx}{0.7\textwidth}{X} 你們的祖宗在曠野吃過嗎哪，還是死了。 \end{tabularx} \\ \\ \relax
6:50 & \begin{tabularx}{0.7\textwidth}{X} 這是從天上降下來的糧，使人吃了就不死。 \end{tabularx} \\ \\ \relax
6:51 & \begin{tabularx}{0.7\textwidth}{X} 我就是從天上降下來生命的糧；人若吃這糧，必永遠活著。我為世人的生命所賜下的糧就是我的肉。」 \end{tabularx} \\ \\ \relax
6:52 & \begin{tabularx}{0.7\textwidth}{X} 因此，猶太人彼此爭論說：「這個人怎能把他的肉給我們吃呢？」 \end{tabularx} \\ \\ \relax
6:53 & \begin{tabularx}{0.7\textwidth}{X} 耶穌對他們說：「我實實在在地告訴你們，你們若不吃人子的肉，不喝人子的血，在你們裡面就沒有生命。 \end{tabularx} \\ \\ \relax
6:54 & \begin{tabularx}{0.7\textwidth}{X} 吃我肉、喝我血的人就有永生，並且在末日我要使他復活。 \end{tabularx} \\ \\ \relax
6:55 & \begin{tabularx}{0.7\textwidth}{X} 我的肉是真正可吃的；我的血是真正可喝的。 \end{tabularx} \\ \\ \relax
6:56 & \begin{tabularx}{0.7\textwidth}{X} 吃我肉、喝我血的人常在我裡面，我也常在他裡面。 \end{tabularx} \\ \\ \relax
6:57 & \begin{tabularx}{0.7\textwidth}{X} 永生的父怎樣差我來，我又怎樣因父活著，照樣，吃我肉的人也要因我活著。 \end{tabularx} \\ \\ \relax
6:58 & \begin{tabularx}{0.7\textwidth}{X} 這是從天上降下來的糧，不像你們的祖宗吃過嗎哪還是死了；吃這糧的人將永遠活著。」 \end{tabularx} \\ \\ \relax
6:59 & \begin{tabularx}{0.7\textwidth}{X} 這些話是耶穌在迦百農會堂裡教導人的時候說的。 \end{tabularx} \\ \\
[1ex]
\hline
\hline
\end{longtable}
$^{1}$歡迎大家來到我們的周道崇拜.
上星期我們開始研讀《約翰福音》第六章.
耶穌在加伯隆會堂裡面.
「我是生命糧」的講道.
如果你上星期在這裡.
你會知道我將上星期的經文.
十九節的聖經分為了七個部份.
中心的主題其實就在中間第四那裡.
「我父賜天糧,生命及世界」.
這個主題其實剛剛好是平衡了.
《約翰福音》第三章第十六節.
金句中的金句.
「上帝愛西人,甚至將他的獨生子賜給他們.
叫一切順他的不至滅亡,反得永生」.
我們今天要繼續看耶穌那篇講道.
我們要一氣呵成.
我們由《約翰福音》第六章第三十五節開始.
「我是生命糧」 耶穌第一次講「我是生命糧」.
七個「我是」裡面的第一個.
一直會讀到講道的結尾.
就是《約翰福音》第六章第五十八節.
講道一篇講章一篇講道.
如果你去評價一篇講道的時候.
其實最重要的不是那篇講道.
花不花巧精不精彩.
而是這篇講道清不清晰.
它會不會深入淺出.
讓每一個聽眾都可以聽得明白.
從這個角度來說.
耶穌這篇講道說真的完全不達標.
我可以這樣說.
在一個講道堂裡面.
耶穌這篇講道是應該貼堂的.
但貼堂是做反面教材.
我們叫做「How not to preach a sermon」.
你怎樣講道是會摧毀了一篇道.
耶穌這篇道真的做了很多很多不應該做的事.
如果你在今天的神學院裡叫他重寫.
有趣的是.
耶穌好像就是想他們聽不明白.

$^{41}$無論這班人多用心,多追求,多渴望,多想信耶穌.
其實到最後我覺得如果你用心聽的時候.
你都是非常迷思和困惑.
不知道他在說什麼.
你問「How can it be?」.
耶穌怎可能是這樣講道呢?.
其實答案是因為耶穌所講的道.
都未完全結束.
耶穌在地上講的每一篇道.
其實你最後一句都要加上一句.
「To be continued」.
「待續」.
你看連續劇的時候.
其實故事還未完.
叫你下次再回來.
我們今天講道我都要整體.
我都不知道你下星期回來不回來.
所以我怎樣都要給你一個比較完整的.
但耶穌不是.
耶穌那些講道每一篇到最後都是「待續」的.
雖然他沒有告訴你.
但其實每一篇都是「待續」.
直到有一天耶穌大喊一聲「成了」.
那篇道就完了.
你明不明白?.
所以你之前聽他講的那些道.
真的不知道他在說什麼.
但去到他大喊一聲「成了」那一天.
他一生要講的這個天國福音就大功告成.
去到那一天.
你回頭看出來.
就是「成了」那一天之後.
你會發現.
一天都光了.
原來我現在明白那篇道是什麼了.
因為最後的太陽升起來了.
太陽升起來一照回前面.
完全你懵然不知.
不知道他發生什麼事的道.
是完全明白了.

$^{81}$我希望你明白我講的背後的道理.
我們回到今天.
今天這篇道其實真的不容易明白.
如果你真的要明白耶穌講的這篇道.
「我是生命量」這篇.
剛才我們讀了很多很多的經文.
我們必須要解說明白以下五條問題.
我讀完這麼長的道的時候.
你要解決這五個問題.
你才會明白耶穌究竟在講什麼.
第一條問題是.
「絕不飢餓,永不乾旱」.
飢餓和乾旱究竟代表什麼.
你一定要問這個問題.
第二個問題.
耶穌說你要吃我的肉.
喝我的血.
真的要的嗎.
真的不喝不行嗎.
不吃不行嗎.
第三條問題.
我們的祖宗在曠野死了.
現在需要天糧才可以永活.
為什麼我們的祖宗會死.
為什麼要吃天糧才可以永活.
第四條問題.
「他在我裡面,我也在他裡面」.
這個代表什麼.
最後「末日我要使他復活」.
末日最後的光景會是怎樣的一幕.
如果你能夠解答這五條題目.
你就可以明白耶穌究竟在講什麼.
我們想逐章逐節看完今天的經文.
然後我們要回到這五條問題.
當你讀完這五條問題.
解答了這五條問題.
你就可以明白今次的訊息.
今天的經文是耶穌的講道.
耶穌的講道由35節開始到58節.
其實剛剛好分了兩件三文治.

$^{121}$有兩個部分.
而每個部分都是前後呼應的.
我們先看第一件三文治.
是第六章第35節至到第六章48節.
第六章35節耶穌對他們說.
「我就是生命的糧」.
「到我這裡來的絕不飢餓」.
「信我的永不乾旱」.
第48節就是三文治的結束.
為什麼我說它是三文治.
因為祂說「我就是生命的糧」.
於是第35節和第48節剛剛好是前後呼應的.
中間那句說話就是整篇的主題.
或者這三文治的主題.
中間那句說話可以重整一下.
變成「到我這裡來又信我的人」.
「這人絕不飢餓永不乾旱」.
於是你明白了耶穌不是叫你來了就得到一件事.
然後信我再拿第二件事.
而是你來了信我.
你不再飢餓.
你也永不乾旱.
去到這裡如果你醒目的時候.
你馬上會問一句問題.
你會說「那我好吃好住」.
「我又不餓又不渴」.
「那這段聖經或者這篇道對我有什麼意義呢」.
如果你上網上聽一些牧師講這篇道的時候.
其實有兩個可能的.
大部分牧師都是取第一個可能.
第一個可能就是.
你雖然身體上無憂無慮.
你可以不飢餓不乾旱.
但是人的心靈總是會飢餓總會乾旱的.
於是我們每個人都需要有人生的意義.
有滿足有目標.
你有錢又怎樣.
你心裡面是空虛的.
於是耶穌講的飢餓.
耶穌講的乾渴其實不是在講肉體上的.

$^{161}$是在講精神上的.
大家明白嗎.
你這樣講的時候很多人聽.
「那又是喔」.
於是你覺得這個牧師的講話都可以繼續講下去.
這是第一個可能性.
但這個可能性當然不是耶穌所想的可能性.
第二個可能性.
耶穌從一個屬天永恆的角度去想飢餓和乾渴.
這個從屬天的角度來看的飢餓和乾渴.
是大部分人其實感覺不到的.
但是是最真實的.
最終極的.
但是大家感覺不到的.
我講一個故事出來.
你就明白我在講什麼.
什麼是飢渴.
什麼是飢餓.
什麼是乾渴.
因為故事聖經裡面.
如果你在讀終極問答的時候.
你很快就會去到那個故事.
那個故事叫財主與拉撒盧.
財主在世界的時候有沒有飢餓.
有沒有乾渴.
沒有.
但是財主死了之後.
他就陰間就飢渴和乾渴.
你明白嗎.
因為他在那裡.
紅紅烈火燒到.
他跟阿爸阿媽說.
可不可以派拉撒盧給我兩滴水.
耶穌說了那個.
所以這裡你不要看和我無關.
因為我又不飢又不渴.
我大把銀大把洗.
接著你可以說.
我心裡面狀態都很好的.
我已經安排了我之後退休的日子.

$^{201}$我去哪裡玩.
我心裡面非常健康.
非常旺盛.
我完全沒有這樣的心靈的飢渴需要.
但是有一樣東西.
你感覺不到的.
但是必然會發生的.
你有沒有想過在故事裡面.
那個有錢人根本在他在世的時候.
想都沒有想過他有這個飢渴的可能.
是不是.
這個就是可怕的地方.
因為你現在今天.
你坐在這裡的時候.
我跟你說這件事.
你覺得.
我感覺不到.
我感覺不到地獄的火.
我感覺不到我將來會飢渴.
但是我告訴你就是會.
這件事可怕的地方就是.
你根本想不到.
你根本沒有辦法連繫到.
但是這件事必然會發生.
這是耶穌引導我們去想的.
飢渴的方向.
而信耶穌.
是唯一可以解決這個飢渴乾渴的唯一出路.
信耶穌.
你不但會永不飢渴.
不信耶穌就永遠飢渴.
永遠飢渴不是永遠飢渴.
永遠飢渴是地獄火紅紅.
你開始明白這件事了.
耶穌是唯一的出路.
出路就是信祂.
什麼叫信祂.
你上個星期在這裡就知道了.
Notitia.
你的腦袋裡面.

$^{241}$你要對福音有個認識不夠的.
然後Ascensus.
你要心裡面認同福音裡面.
世人都犯了罪.
虧缺要上帝為榮.
罪的功架乃是死.
最後還有Fiducia.
一個全人的信靠.
全人對耶穌的衛生.
於是我們繼續去看.
耶穌將祂的主旨跟他們說了.
祂說你信我的永遠不乾渴.
耶穌在說永恆的救恩.
可視你們已經看見我還是不信.
耶穌同一班來追求祂.
做皇帝的人跟祂說.
你見到我了.
你仍然不信.
令我想起約翰福音第一章第十四節.
道成了肉身住在我們中間.
充充滿滿的有恩典有真理.
我們也見過祂的榮光.
這班人和十二門徒.
和約翰一樣見過祂的榮光.
只不過他們看見了.
仍然不信.
為什麼不信.
是因為耶穌沒有達標.
耶穌不是他們的類型.
耶穌不是他們想要的君王救主.
以創世第五十三章第二節這麼說.
祂無皆美形象使我們注視.
祂也無美貌使我們.
仰慕祂.
馬太福音第一章第二十一節.
祂的名字叫耶穌.
因為祂要將自己的百姓.
從罪惡裡面拯救出來.
耶穌不是他們的類型.
於是他們不信.

$^{281}$第三十七節三十八節.
凡夫所賜給我的人.
必道我者裡來.
道我者裡來的人.
我總不丟棄他.
因為我從天上降下來.
不是要按自己的旨意行.
而是要遵行差我來那位的旨意.
耶穌在做甚麼.
人們說耶穌在教導預定和揀選.
我的答案是可能吧.
我不知道.
但如果你問我.
我覺得耶穌其實明顯在刻意挑釁他們.
挑釁眼前的一班猶太人.
他們相信自己在負那邊.
但耶穌在挑戰他們.
他們是否真的站在天賦那邊.
你明白嗎.
有沒有排錯隊.
有沒有對自己感覺太良好.
你以為你已經坐定一顆鹿在天賦那邊.
其實你是不是呢.
這是耶穌一直以來對著猶太的領袖也好.
猶太的群眾.
他不斷地警告不斷地勸告.
既是警告又是勸告.
如果你回到前面那一章.
你已經看到.
耶穌和那班猶太人一樣是做同一件事.
你們並沒有他的道存在你們心裡.
因為你們不信他所差來的.
在前面的第五章.
你們查考聖經.
因為你們以為其中有永生.
但這經正是為我作見證.
然而你們不肯到我這裡來得生命.
如果你們記得耶穌所說的一切.
祂不是在定他們的罪.
祂是在警告他們.

$^{321}$祂是在勸告他們.
耶穌在挑戰他們.
去深思到底他們是否真的站在天賦那一邊.
現在有兩隊.
一隊人人都排在那一邊.
耶穌又在另一邊.
又說你排錯隊了.
如果你真的去尋求上帝.
你應該過來我這邊.
人們覺得你瘋了.
而你一個人在那裡.
但耶穌就是在做這件事.
耶穌不斷警告他們.
你回去查一下聖經.
你不信我也信本聖經.
其實這些都是林牧師說的.
我經常叫你不要信我.
你信本聖經.
你想想你究竟信不信.
不要說你因信聖經.
你要想想什麼叫信等等.
這些不是我作的.
我只是鼓勵你想想.
因為我怕你排錯隊.
我怕你以為自己得救.
但其實不是.
耶穌在挑戰他們.
耶穌在勸告他們.
叫他們一旦發現自己排錯隊.
起碼今天你還有機會快快迴轉.
來耶穌那裡信靠跟隨.
上帝所差來的獨一愛子.
回到第六章39節.
差我來那位的子兒.
就是說他所賜給我的.
要我一個也不失落.
並且在末日使他復活.
第四十節.
因為父的子兒.
是要使每個人見了子.

$^{361}$而信的人得永生.
並且在末日我要使他復活.
如果你來到我這裡.
如果你信靠上帝所差來的那一位.
你必定第一不失落.
第二你會得到永生.
所謂永生.
不是說你的身體永遠不死.
不是秦始皇那種.
身體永遠不死.
如果是這樣就不需要有末日復活.
講到尾一早跟你說有末日復活.
就是說你的所謂永生.
就不是你現在的身體.
是存到永遠.
永生是說你一旦信靠了耶穌.
在上帝面前.
我們有句說話.
那個拉丁文叫Coram Deo.
你就在上帝面前.
永遠活著永遠不死.
什麼意思.
在上帝面前.
其他人看不到的時候.
上帝看到.
因為你的身體在上帝面前.
無論是生與死.
上帝都看到你的靈魂.
無論你的身體已經腐化了.
上帝仍然照顧著你的靈魂.
上帝保守著你的靈魂.
因為你的身體無論生與死.
都是上帝的看顧底下.
接著去到末日的時候.
上帝要將你已經壞了的身體.
重新復活.
這個是永生的意思.
不是你永遠不死.
你是在上帝面前不死的意思.
第41節.

$^{401}$猶太人因為耶穌.
「我是從天上降下來的娘」.
就私下議論祂說.
「這不是約瑟的兒子耶穌」.
「我們豈不認識祂的父母嗎」.
「現在祂怎麼說」.
「我是從天上降下來的」.
這裡說猶太人.
這個猶太人跟之前跟耶穌對話的.
那群人又不一樣.
之前那群人如果你記得.
早一天吃了五餅二魚.
接著坐船早上來找耶穌.
那群人耶穌已經打發了他們離開.
因為耶穌指名道姓.
說你們不信.
你夠人見我都不信.
所以你相信那群人已經被耶穌擊退.
這裡的猶太人應該是.
加伯農會堂裡的領袖和群眾.
你記得我們現在在加伯農的會堂.
就是他們的教會.
猶太人的教會.
這群人至少認識耶穌.
認識耶穌和認識祂的家人.
因為聖經學者相信.
耶穌一家本來是在拿撒勒.
後來在他還沒開始去加勒比利時的時候.
他們已經搬去了加伯農好幾年.
於是這群人都看著耶穌長大.
於是他們說我們認識祂的家人.
我們認識祂的.
你橫看店看不見從天上降下來.
他又是另外一個看見了耶穌.
但是看不透耶穌.
道成了肉身住在我們當中.
接著又說我們看見祂的榮耀.
但是這群人和耶穌住了那麼多年.
沒有辦法見到耶穌的榮耀.
這裡說他們私下議論.

$^{441}$如果你們看英文或者原文.
那個翻譯更好.
叫做「發怨言」.
你一聽到發怨言.
Grumble你就知道.
這個發怨言是從舊約聖經出來的.
以色列人一出埃及很快就發怨言.
那個發怨言就代表了他們.
在曠野裡面的叛逆和罪惡.
所以這群人.
很快現在又開始發怨言.
這群人見到耶穌.
但是不信.
第43節耶穌回答他們說.
「你們不要彼此私下議論」.
「你們不要再發怨言了」.
這個是另外一個翻譯.
第44節.
「若不是差我來的父吸引你」.
「就沒有人能到我這裡來」.
「到我這裡來的再末日」.
「我要使他復活」.
「若不是差我來的父吸引人」.
「就沒有人能到我這裡來」.
耶穌在做什麼呢?.
有些人又說.
耶穌很明顯在教導預定和揀選.
我答案也是可能.
那些人一見到這些父差來.
他就是這樣想.
你問我.
我也是那一句.
我覺得他是明顯在挑戰他們.
挑釁他們.
問他們拿個「父」出來打他們.
告訴他們.
其實他們好像不在「父」那邊.
他們排錯隊.
有沒有想過他們排錯隊.
因為猶太人自己最驕傲.

$^{481}$最引以為榮的就是他們有「父」.
有耶和華上帝.
現在耶穌每次都拿個「父」出來打他們.
說如果你真的在「父」那邊.
你應該來找我.
很明顯你不信我.
你不信我的原因.
你根本不是在「父」那邊.
於是耶穌其實有意無意之間.
在挑戰他們.
在挑釁他們.
但是好的那一種.
耶穌對他們說.
你們不要再議論.
什麼意思.
不要發怨言.
耶穌說停.
不要.
你繼續走那條路.
你死定了.
但是你現在可以留態.
你不一定要繼續發怨言.
耶穌在警告他們.
因為你繼續下去.
你就變成了你的祖宗.
但是你可以現在停.
你可以停.
閉嘴.
不要再發怨言.
你再想想我是誰.
你回來.
還有可能耶穌呼召他們.
去迴轉悔改.
第45節.
於是耶穌開始引用舊約.
在先知書上.
所謂先知書.
如果上個禮拜我都有講過.
是以賽亞書第54章.
先知書寫著說.

$^{521}$你們都要蒙上帝教導.
很明顯你們沒有蒙上帝教導.
人們一看到這件事.
那麼耶穌引用以賽亞書.
覺得耶穌拿以賽亞來.
去打他們.
照我想.
就不是的.
我覺得.
耶穌是在用他的恩典廉文.
給他們的提示.
如果我無緣無故.
講到的時候.
打了以賽亞書第54章出來.
你會怎樣呢.
如果你是一個乖兒子乖女.
你回家好好地.
拿回以賽亞書第54章.
再出來看.
你會好奇.
林獄師你不知道在引用什麼東西.
我也不知道他引用對還是錯.
尤其是你覺得.
我都是亂說話.
不要緊我回去看以賽亞書第54章.
既然回去看以賽亞書第54章.
不如再看以賽亞書第53章.
不如再看以賽亞書第55章.
耶穌就是在做這件事.
以賽亞書第53章第4節.
他承認擔當我們的憂患.
背負我們的痛苦.
我們卻以為他受責罰.
被上帝打擊苦待.
他為我們的過犯受害.
為我們的罪孽被壓傷.
因他所受的刑罰我們得平安.
因他所受的鞭傷我們得抑自.
我們都如羊走迷.
各人偏行幾路.

$^{561}$耶和華使我們眾人的罪孽.
都歸在他身上.
你乖仔你回去看.
你跟他頂嘴.
你跟他頂嘴都要跟他arm 嘴.
他說以賽亞書第54章.
你就回去讀.
然後你讀著讀著讀著.
讀到以賽亞書第51章來.
你們所有乾渴的.
都來這水邊.
沒有銀錢的也來.
你們來.
買了吃不用銀錢.
不用代價就可以買酒和奶.
你們為何花錢買那些不是食物的東西.
用勞碌的來買無法使人飽足的東西.
耶穌在說饑渴的東西.
你們要留意聽從我的話.
你就可以吃那美物.
聽我的話.
得享肥美心中的喜樂.
你們當質疑而聽.
第三節第55章.
來到我這裡.
質疑而聽來到我這裡.
很明顯.
要聽就必傳喚.
耶穌玩ye 了.
我要與你們納永約.
就是應許給大衛那可靠的慈愛.
Stephas Love.
耶穌很聰明.
耶穌真是恩典憐憫.
慈愛信實.
旁敲則擊.
引領上帝伴翼的子民踏上回家的路.
很厲害.
於是第46節.
凡聽了父的教導而學習的.

$^{601}$都到我這裡來.
這不是說.
又加了一些東西.
講完他們了.
你小心你是不是在父那邊.
這不是說有人看見過父.
唯獨從上帝來的那才看見過父.
這個其實對應了.
楊浦英第一章第18節.
從來沒有人看見過上帝.
只有傅懷磊獨一的上帝.
或者獨一的兒子.
將上帝表明出來.
一講下他們.
你回去想一想.
亦都講下他自己.
第47節.
即是他在挑戰他們.
他說你要想下你們自己.
是不是真的在父那邊.
跟著告訴你.
你猜下我是誰.
從來沒有人看見過父.
唯有我看見過父.
罪人了.
你們要認真審視.
你現在的肅靈狀況.
這是第一.
第二.
你要轉眼仰望我.
我是誰.
看真一點.
我是你獨一的舊主.
看下你自己.
又看下我.
你又知道.
你現在大件事了.
你亦都知道.
我是你唯一的舊主.
他那天是沒有人聽得明白這篇道.

$^{641}$不要緊的.
這篇道未完的.
但你以後來日方長.
只要你跟著耶穌.
跟著祂去.
你就要不斷思想.
究竟你現在的狀況在哪裡.
跟著耶穌究竟是誰.
第47節.
我實實在在告訴你.
你們信的人有永生.
我就是生命的能.
好 我們完成了.
兩個三文治的第一份.
還有第二份.
我們繼續.
第二份第49節開始.
仰慕音第六章第49節.
你們的祖宗在曠野吃過嗎哪.
還是死了.
這是從天上降下來的糧.
他在說自己.
使人吃了就不死.
我就是從天上降下來的生命的糧.
人若吃這糧吃我.
必永遠活著.
這個到底在說甚麼.
我想你稍後繼續思想.
待會去到那篇道的最後.
我才跟你講解明白.
但我想你留意一樣東西.
這三節聖經是開場白.
他們其實對應了整篇講到的.
最後一節第58節.
58節是這樣說的.
這是從天上降下來的.
不像你們祖宗吃嗎哪還是死了.
吃這糧的人將永遠活著.
你看到前後對應.
你如果看回之前的開場白.

$^{681}$你們祖宗在曠野吃過嗎哪還是死了.
我就是從天上降下來的.
人若吃我便永遠活著.
死了活著.
祖宗祖宗祖宗死了.
你吃我便活著.
於是在這裡你看到前後的對應.
你也看到另外一樣東西.
就是祖宗的死.
對應了耶穌所帶來的新生命.
永遠的生命.
祖宗死了.
耶穌就活著.
耶穌不單止活著.
耶穌有永遠的生命.
究竟耶穌在說甚麼.
我給你一個小小的提示.
耶穌所說的祖宗.
全部都死了在曠野.
當中祖宗裡面少了兩位.
有兩位長者.
就是快手快腳走入了認祂之地.
沒有死在曠野.
那兩位是誰.
那兩位約書和格納.
只有兩個.
其他長者全部死了.
死了在哪裡.
死在曠野裡面.
他們死的原因.
因為他們罪孽.
因為他們叛逆.
因為上帝叫他們進去認祂之地.
他們不相信.
他們說上帝讓我們去送死.
那裡有巨人.
上帝搞甚麼鬼.
不如我們回去埃及.
所有的人都要回去埃及.
唯有兩個長者.

$^{721}$即是當時不長者.
他們要我們進去.
德加勒和約書.
沒有死.
於是我們問這個問題.
祖宗為甚麼會死在曠野.
死的原因.
因為祖宗都是叛逆上帝的罪人.
你開始明白.
為甚麼我們現在需要天降的麵包.
第一要赦罪.
我們的罪需要得到赦免.
第二.
我們要跟隨耶穌.
再不是偏離上帝那條路.
回去埃及.
而是跟隨耶穌.
帶我們進入永遠的安息.
耶穌是天降的麵包.
天降的麵包是生命的麵包.
原因是因為祂是赦罪的麵包.
如果當中有些人聽我說主禱文.
你就知道了.
主禱文.
馬來福音第六章十一二節.
我們日用的飲食.
真是翻得差.
我們每日的麵包.
不是日用的飲食.
日用的飲食.
你立刻說你有甚麼想求.
你就全都求了.
我們每日的麵包.
只是一樣東西.
從天降下來的馬來.
我們每日的麵包.
今日賜給我們.
這個麵包是甚麼.
第十二節.
免我們的債.

$^{761}$中間有個CHI的連接詞.
那個連接詞有很多種不同的翻譯.
其中一個翻譯就是.
就是免我們的債的麵包.
有隻麵包.
不單叫生命麵包.
應該這樣說.
有隻麵包叫生命麵包.
它為甚麼叫生命麵包.
是因為它可以免除你的債.
如果它沒辦法免除你的罪債.
對不起了.
你是不會有生命的.
它之所以是生命麵包.
不是因為它可以保你命.
讓你可以健康些.
沒用的.
是使你可以除罪.
這個就是耶穌帶給我們的東西.
第五十一節的最後一句.
耶穌這樣說.
開始越來越難搞了.
我為世人的生命所賜下的糧.
就是我的肉.
從麵包說到我的肉.
第五十二節.
因此猶太人彼此的爭論說.
這個人怎能把他的肉給我們吃呢.
其實你很難怪當時的猶太人這樣的反應.
所以我這裡特意用了豬肉的圖畫.
猶太人最討厭的就是豬肉.
最不乾淨的就是豬肉.
是非常震撼的一個圖畫.
耶穌說你要吃我的肉.
我覺得耶穌是故意.
故意留難他們的.
是挑釁他們的.
令他們馬上.
正如我剛才所說.
他不是想他們當天聽得明白.

$^{801}$他當天是聽不明白.
但是他想他們記得他的那幅圖畫.
於是說了一些很難受的東西.
挑釁他們.
耶穌所說的道.
那天沒有完.
他只是在預備這班聽眾.
將來聽他將來的道.
看他將來的戲.
我想你看看第五十一節.
有一個邏輯.
第五十一節的邏輯是這樣的.
他說.
我為世人的生命所賜下的糧.
就是我的肉.
耶穌所賜下的麵包是他的肉.
他的肉是為了世人的生命.
他的肉是為了使世人得生命.
他的肉是為了使世上的罪人.
罪得赦免與上帝和好得永遠生命.
你明白這個邏輯吧.
這個邏輯是.
我的肉是為了世人得生命而有的.
如果我的肉沒有獻上.
如果你沒有吃我的肉.
你是沒有辦法得到我們的生命.
原因我們現在知道了.
從十字架另外一邊.
我們知道的原因是因為.
他的肉或者他的犧牲.
他十字架的犧牲.
唯有使罪人罪得赦免與上帝和好.
才可以得到永遠的生命.
所有的一切都是指向耶穌基督的十字架.
繼續看第五十三節.
第五十四節.
耶穌對他們說.
我實實在在的告訴你們.
你們若不吃人子的肉.
不喝人子的血.

$^{841}$在你們裡面就沒有生命.
吃我的肉喝我的血的人.
就有永生.
並且在末日我要使他復活.
你這樣看一定會想.
嘩!真的要吃他的肉.
真的要喝他的血.
耶穌是不是真的逼人太甚呢?.
這次你有沒有發現.
他加了要喝血的.
之前喝食肉.
現在要喝血.
這樣就糟糕了.
因為摩西律法裡面.
明明多次提到不可以吃帶血的肉.
現在還要喝血.
又是肉又是血.
真是瘋的.
造反的.
你講到最後的時候.
我會幫你解決.
這個迷思.
究竟是不是真的要吃他的肉.
喝他的血.
這裡你有沒有留意一件事.
耶穌講了一件事.
他講這個肉和血是誰的血.
他不是單單講我.
他是講人子.
他不是整天都在講人子.
但是人子講到的是.
但耶穌第七章裡面要繼承天國.
就是永恆天國的人子.
嘩!是一個君王.
一個上帝.
已經決定了要將永恆天國給他的君王.
人子竟然要奉上他自己的生命.
但是這個人子重要的地方就是.
對上一次人子.
就是之前有講過第三章.

$^{881}$有出現過的.
就是耶穌和尼哥底母最後講.
重生的時候.
講了一句和人子有關係.
他說當摩西在曠野裡面.
怎樣舉起蛇.
然後他不是我也.
他說人子也照樣被舉起蛇.
好像蛇舉起.
使一切信他的人都得著永生.
所以所有東西都是指向十字架.
當他講人子的血.
人子的肉的時候.
令我們想起摩西舉蛇的那一刻.
所有東西都和十字架有關係.
第五十五節.
第五十六節.
我的肉是真正可吃的.
我的血是真正可喝的.
這裡的中文翻譯又出了一些問題.
如果你看英文就是.
我的肉是true food.
我的血是true drink.
如果我籠統一點翻出來.
就是我的肉是真食物.
我的血是真飲料.
真飲品.
所以大家明白為什麼他將它改了.
但是這個真食物.
true food.
true drink.
真飲料.
是重要的.
為什麼.
因為在養父音裡面.
耶穌經常是真的.
祂是真光.
true light.
祂是true bread.
真糧.

$^{921}$真麵包.
後來還有一樣東西.
祂會成為我們的真葡萄樹.
true vine.
所以沒有了這個東西.
後面你唸不到那些東西.
但是這裡耶穌正在說一些.
祂就是真真真真真真真真真真.
珍珠都沒有這麼真的東西.
強調.
祂的永恆的生命.
第五十六節.
吃我肉喝我血的人.
常在我裡面.
我也常在.
他裡面.
還是那句.
待會我們再回來看看.
什麼叫常在我裡面.
看看常在我裡面.
這件事待會回來最後才看.
第五十七節.
永生的父.
怎樣差我來.
我又怎樣因父而活著.
先講祂和父的關係.
照樣吃我的肉的.
也要因我活著.
在說我們和耶穌的關係.
最後.
那個是麵包的三文治.
底層的麵包.
這是從天降下來的.
不將我們的祖宗吃了.
嗎哪雲是死了.
吃這糧的人將永遠活著.
這個是耶穌的結案陳詞.
這個是耶穌永生.
認祂的高潮.
我想簡單去.

$^{961}$重新講一下第五十七節.
五十七節是這樣說.
祂說我與父.
擁有並分享唯一.
真正的生命.
True life.
祂說得很fancy.
永生的父怎樣差我來.
我又怎樣因父而活著.
不要理中間在說什麼.
祂說的就是.
我與我父親共同擁有.
共同分享.
這個世界唯一真正的生命.
The only true life.
就在我們裡面了.
我在父裡面.
父也在我裡面.
這個跟你沒有關係.
不過現在.
天國賣大包.
如果你願意到我這裡來.
如果你願意全心信靠我.
全人跟隨我.
你就會常在我裡面.
之前那一節.
你就會常在我裡面.
我也常在裡面.
於是你就因我活著.
你會因著我得到一個新的生命.
你會因著我得到永生.
還有.
這個永生.
這個新的生命.
不單只是你和我分享.
不單只是我和你分享的人生.
還是你和我.
我和天父.
天父和我和你三個人共同分享的新生命.
三個人共同分享的永恆生命.

$^{1001}$為什麼這個重要呢.
因為你去到第十七章裡面.
耶穌會為我們定義什麼叫做永生.
這個金句.
第十七章.
耶穌最後的祈禱.
上十字架之年的祈禱.
祂這樣說.
什麼叫永生.
認識你獨一的真神.
並且認識你所差來的耶穌基督.
就是永生.
認識上帝.
認識耶穌.
當然你知道.
你聽我講了很久就知道.
認識呢.
不是頭腦上的認知.
Gnostic.
認識其實是代表愛慕的意思.
亞當認識夏娃就生了一個孩子.
亞當愛夏娃.
最後就是一個親密的愛慕.
認識上帝.
認識耶穌.
愛慕上帝.
愛慕耶穌.
永遠認識上帝.
永遠愛慕上帝.
就是永生.
上帝愛我.
我也愛上帝.
我們在這個永恆的愛的關係裡面.
這個是一個永遠上帝的生命.
在上帝的眼前.
他見到愛他的人.
他會保守他們的一生.
他雖然死了.
末日會復活.
使他加入上帝永恆的天國裡面.

$^{1041}$好 我們吃完兩件三文治.
讓我們回到起初的五條問題.
我們逐一解決.
我們的Bible Study就完結了.
五條的問題.
我們需要解答.
第一 絕不飢餓.
永不乾渴.
是什麼意思.
第二 吃我的肉.
喝我的血.
真的要嗎.
第三 曠野死了.
天良活著.
是什麼意思.
他在我裡面.
我也在他裡面.
是什麼意思.
最後 末日我要使他復活.
復活了之後是什麼一個場景.
是我們需要知道的.
於是我們從第一條開始.
第一條 飢餓乾渴.
絕不飢餓 永不乾渴.
這是耶穌最大最大的應許.
這個究竟代表了什麼.
我想帶你回去看一段舊約的聖經.
你可能還沒想過.
尼希米記第九章第十五節.
尼希米書.
第九章第十五節.
你從天上賜下糧給他們充飢.
使水從盆石流出給他們解渴.
你吩咐他們進去.
得你起誓應許要賜給他們的地.
還有另外一段聖經我沒有放出來.
如果你有機會.
你回去看申命記第二十八章.
申命記的尾章經常說福氣和咒詛.
申命記第二十八章告訴我們.

$^{1081}$如果有人背叛耶和華上帝的盟約.
其中一個眾多咒俗當中一個咒俗.
就是飢餓和乾渴.
你明不明.
所以如果你飢餓你乾渴的話.
這代表了上帝不喜歡你.
當你背叛了上帝的盟約.
反過來說.
如果上帝會供應為你充飢.
上帝會為你解渴.
使你絕不飢餓永不乾渴.
這代表上帝.
代表你對上帝的忠誠.
而因為你對上帝盟約的忠誠裡面.
你所得著的福氣一個是咒詛.
另外一個是福氣.
於是不再飢渴.
代表了上帝喜歡你.
上帝祝福你.
也代表了這裡.
你看到的.
這是出埃及的圖像.
上帝會與你同在.
我想你留意一件事.
這兩個雙重的福氣.
即是Twin Blessing.
是從哪裡來.
是哪裡來到上帝的子民.
答案是第一.
從天上賜下糧食.
從天上賜下的糧食.
充飢嗎哪.
第二是從水.
水是從盤石流出來.
如果你要回到舊約.
你知道摩西要擊打一個盤石.
上帝給了一塊石子裡面.
兩個都是舊約的影子.
沒時間和你看裡面的經文.
但你應該依稀記得.

$^{1121}$無論是嗎哪也好.
或者是上帝給他們的盤石.
是跟著他們的.
在新約裡面也有提出來.
兩個都是指向耶穌的十字.
一個是祂的身體.
一個是祂的寶血.
去到哥林多前書第十一章.
保羅說了一句話.
The Rock is Christ.
那個石頭就是基督.
其實是基督.
所有只是一個影子指向基督.
這是我想你第一.
留意所有東西指向耶穌.
第二.
我想你再留意一件事.
這個雙重福氣.
從天降下來到上帝的子民當中.
是有目的的.
目的不是供應你日用的飲食.
目的是什麼呢.
目的就是.
你吩咐他們進去.
得你.
起誓認許要賜給他們的那地.
是有方向感的.
你不可以說多謝上帝.
又使我衝擊又解放.
我現在想回去埃及了.
不可以的.
你一回去埃及就馬上.
沒有麵包.
石頭扔死你.
也不會有水.
什麼都沒有了.
你又會飢渴到地獄.
所以.
上帝給我們從天而降的麵包.
上帝給我們.

$^{1161}$盤石中的活水.
全部是有方向感的.
方向感.
就是要我們逃離埃及.
方向感要我們奔跑天路.
方向感叫我們見證天國.
叫我們進去得你起誓.
認許要賜給他們的地.
去就去吧.
不是的.
去的時候要打仗.
所以這個我們今天不會說了.
但是你明白.
就是你.
Are you ready to fight?.
就是阿加納去到.
Give me this mountain.
你四十年之前答應過那個地方.
你以為他走過去就可以了.
他要打仗的.
但是他不怕打仗.
他說Give me this mountain.
我等了四十年.
上帝給我的地方.
我預了我粉身碎骨.
上帝認許了我不會粉身碎骨.
我一定拿到的.
所以這個所有的東西.
上帝永不飢渴代表了什麼.
我們擔負了上帝天國的使命.
OK.
還有一點.
還有一點很重要的.
一定要講清楚.
不然待會回去你就死了.
真的不是死.
真的是我.
我對不起你.
就是原來不再飢渴.
不是真的代表不再飢渴.

$^{1201}$你說什麼.
我想帶你去看哥林多前書第四章.
保羅的見證.
死了保羅.
保羅你忘記了耶穌說不再飢渴了嗎.
保羅就是和耶穌不對牙.
保羅這樣說.
第四章第十一節.
我們又飢又渴.
這樣也可以.
我們又飢又渴.
我們赤身露體.
我們挨打到處漂泊.
人還把我們看作世上的污穢.
萬物中的渣滓.
我寫寫寫給你們.
是要警戒你們.
好像我所愛的兒女.
接著我求你們還要什麼.
效法我.
效法我做什麼.
好吃好住.
效法我又飢又渴.
赤身露體.
挨打到處漂泊.
你們知道剛才傳說中發生了什麼事嗎.
就是保羅離開之後.
有一群假教師來了.
妖言惑眾.
跟你說傳福音給你的保羅不對.
去到哪裡被人打到哪裡.
又窮又不行.
你跟我們不一樣.
我們都是相信同一個耶穌的話.
不過我們這樣信.
信得很舒服.
有什麼需要馬上求.
又搞定.
我們不飢不渴.
他又飢又渴.

$^{1241}$接著保羅寫一封信出來.
我被人佔了便宜.
成功神學富貴福音.
神經病的.
令你入地獄的.
於是他寫了這段這樣的東西.
他當他們是他屬靈裡面的兒女.
他說這些人.
那些老師你不要理他.
但我是你屬靈裡面的父親.
你怎麼可以不聽我.
我才是真的.
他們是假的.
於是保羅大聲.
我又飢又渴.
還要叫他們要來效法我.
究竟發生什麼事.
我這樣解釋給你聽.
不再飢渴是一個屬靈永恆世界裡面的實體.
通常這些東西都是感覺不到的.
要憑信心的.
他指的不是地上物質世界裡面你不飢不渴.
而是說在屬靈永恆.
在上帝眼中上面看著保羅.
他從來不再飢渴.
地獄的火碰不到保羅.
永恆的飢餓無法觸碰到.
無法碰到傷害到保羅.
接著保羅說.
為著福音的緣故.
我保羅甘願又飢又渴.
因為在靈裡在永恆世界裡面.
我已經不再飢渴.
你明不明白這個邏輯.
我在這一生裡面我可以又飢又渴.
為什麼.
因為在永恆另外一個世界.
另外最終極最真實的世界裡面.
我已經不再飢渴.
是因為我不再飢渴.

$^{1281}$所以令我今天.
我可以甘願又飢又渴.
我可以這樣再形容.
為了福音的緣故.
保羅甘願放下地上的財寶.
甚至生命.
因為在靈裡在永恆在終極的世界裡面.
他已經得著財寶得著生命.
你明白嗎.
有誰願意放棄今生的財寶和生命.
就是他很肯定他在永恆天國裡面.
已經得著財寶得著生命的東西.
因為我深信我在另外一個世界.
永恆的世界裡面.
我不再飢渴.
所以我願意在今生裡面.
又飢又渴.
得人驚.
你下個禮拜不敢回來.
你真的要應用我們.
你就知道多誇張.
看第二條問題.
我們是不是真的要吃耶穌的肉.
喝耶穌的血.
我吃素的也不行.
原來是可以的.
請告訴我們.
你看看這兩節的聖經.
第六章第四十節.
因為我父的旨意是要使.
每一個人見了子而信的人得永生.
並且要在末日使他復活.
第五十四節.
吃我的肉喝我的血的人就有永生.
並且在末日我要使他得復活.
看到這兩節聖經是彼此平衡的.
然而有一節是車頭.
另外一個是車廂.
兩個一樣是平衡的.
哪一節是車頭哪一節是車廂.

$^{1321}$第四十節是車頭.
第五十四節是車廂.
前面的是一個真實的呼召.
後面的是一個襯托的比喻.
你開始明白了.
前面那個才是真實的.
不要看後面.
當然你遲到.
聽不到前面那句你活該.
但如果你進來.
你第一句林牧師說的時候.
叫你見了子而信.
然後你繼續聽下去.
十句之後又說.
吃他的肉喝他的血的人就有永生.
你知道不用怕不用怕.
前面說過了.
這個他在做一個襯托的比喻.
耶穌不是叫我們吃喝他的血肉.
耶穌要我們其實是上面.
上面那個才是主子.
要他見了子而信.
要他仰望聖子.
如果你記得上個星期.
仰望他.
要他全心信靠他的血肉.
他十字架上血肉模糊的贖罪祭.
信他擁抱他.
那你就有永生了.
奧古斯丁說了一句說話很有意思.
他這個是拉丁文.
credi就是信.
ed就是and.
英文翻出來就是.
Believe and you have eaten.
他說你信你就吃了.
當聖母耶穌吃的時候.
他說你信了就等於你已經吃了.
所以這兩段是故意的.
但為什麼耶穌要故意.

$^{1361}$明明知道他們很受不了要吃他的肉.
喝他的血.
原因就是他要給他一個.
很強烈鮮明的血肉圖畫.
很強烈的鮮明的血肉圖畫.
帶我們去到將來.
又或者帶我們去看《約翰福音》裡面.
三段非常之重要血肉的圖畫.
然後有一天.
你想通《撒尼亞話語》的時候.
你知道那幅音是什麼了.
第一.
他說你要吃我的肉.
sars就是希臘文.
在哪裡呢.
道成肉身.
那個肉.
道成了肉身住在我們當中.
sars.
他的肉.
道成了肉身.
住在我們當中.
充充滿滿的有慈愛有信心.
第二.
這個肉身將要被獻了.
變成血肉模糊的一樣東西.
叫做可能上帝的羔羊要除去世人的罪.
最後血肉模糊的圖畫是.
士兵拿著槍擲他的臘腸.
有血和水流出來.
這些圈套是預備他將來會見到一場戲.
最後的高潮.
於是他明白耶穌所說的一切.
是要我們仰望這個十字架上面的救主.
我們去到第三條問題.
曠野死了天良永活.
以色列人為何在曠野死了.
剛才我們已經說過了.
只要你將那兩位老人家.
約書亞彤和加納拿出來.

$^{1401}$你就知道為何他們死在曠野.
為何那兩個長者沒有死在曠野.
為何其他人死在曠野.
為何耶穌基督的天良才可以使我們得心.
我想帶你去看詩篇第95篇.
這裡說的是出埃及後第一代以色列人.
他們死了在曠野裡面四十年.
上帝讓他們兜兜轉轉四十年.
自然死 但兜你四十年.
即是二十歲以上的出埃及那些.
全部要在曠野死.
使他們不得安息 叫他們失落應許.
這是他們的解剖驗屍報告.
詩篇第95篇 他們的解剖驗屍報告.
為何他們要死在曠野裡面.
因為四十年之久.
我厭煩那世代.
說這是心裡迷糊的百姓.
竟不知道我的道路.
第十一節.
所以我在怒中起誓.
他們斷不可以進入我的岸城.
原來他們的叛逆不是一時失腳.
而是一生的持續.
於是耶穌在這裡跟我們說.
你們不可以像你們祖宗那樣.
你們不要發怨言吃嗎哪.
如果你發誅怨言吃我身體.
你都死的 你明白嗎.
你必須要真心悔改回轉.
你祖宗衰甚麼.
你的祖宗衰不是因為嗎哪不夠堅.
你祖宗衰是因為他心裡.
根本不悔轉 沒有悔改.
但是現在你要吃我的身體.
我的肉我的血的時候.
你一定要真悔轉.
真悔改你才能吃.
你吃不到.
如果你不是真信的話.

$^{1441}$你一樣死.
耶穌說你們要真正的悔改回轉.
當我這類來全心信教.
全人馴服.
你要仰望我的十字架.
領受我天上降下來的能量.
使你們聚得赦免.
使你們與上帝和好.
你又要跟隨我的道路.
記得吃了嗎哪.
喝了活水.
要去哪裡.
要進入上帝的安息.
怎樣進入 要打仗.
給我這座山.
要我們不回頭去喊求.
要我們奮勇向前為上帝做見證.
進入我的安息.
就是永遠的生命.
只有這樣才是永遠的生命.
好 第四條問題.
我們進入最後的狀態.
這裡出了一句.
他在我裡面.
我也在他裡面.
是甚麼意思.
這句說話在《讓步音》裡面出現過兩次.
一次是《讓步音》第六章.
我們剛剛讀過.
另外一段我想大家都很熟悉.
《讓步音》第十五章.
你就知道其實第六章是在預備.
《讓步音》第十五章.
因為會再出來的.
我是葡萄樹.
你們是枝子.
常在我裡面.
我也常在你們裡面.
這人就多結果子.
因為離了我.

$^{1481}$你們就不能做甚麼.
原來常在我裡面.
我也常在你們裡面.
說的是要結果子.
甚麼果子.
生命的果子.
新生命的果子.
不是舊生命的果子.
不是做一些道德上的事.
是一個新生命的果子.
會懂得飛的.
永生的果子.
後來你見到.
尼哥德姆那種就是永生的果子.
神了經那種.
福音的果子.
光明的果子.
要在黑暗的世界裡面做光明的見證.
所以說我是葡萄樹.
你們是枝子.
從前我們是殘花敗柳.
憤怒之子.
殺人都犯了罪.
虧缺了上帝的榮罪.
公開大事.
死死後且有審判.
我們是殘花敗柳.
我們是叛逆之子.
我們是憤怒之子.
但恩爵來到耶穌面前.
因為信靠跟隨耶穌.
與祂聯合.
與祂接枝.
接枝.
我不懂.
我不懂做這些事.
但是神學家是這樣說的.
有一棵樹沒得救了.
那怎麼辦呢.
你可以將它和另一棵很健康的樹接枝.

$^{1521}$我們和耶穌接枝.
於是我們獲得了一個全新的生命.
祂在我裡面.
我也在祂裡面.
重生是一個永生.
這裡說到你接枝是一次.
不用每天接.
接了之後.
你要常在我裡面.
我也常在裡面.
正如我們相信耶穌一次.
但我們要持續的跟隨.
持續的相信.
持續在祂裡面.
祂也在我裡面.
我們吃耶穌的餅是一次.
但我們要持續繼續.
我們日用的飲食.
我們每天的麵包.
今天次給我們.
我們的主禱文是不斷祈禱的.
每天祈禱的.
但是你第一次吃耶穌.
信耶穌的時候.
只是一次.
然而你不斷繼續.
我們喝水也是喝一次.
然後持續的喝水.
常在我裡面.
我也常在你裡面.
是一個持續的意思.
也令我們想起這一節.
很有趣的聖經.
你記得耶穌為他們洗腳那一鑊.
然後耶穌說了一句話.
彼得那時候說.
不如從頭到尾.
再洗一次洗八百.
耶穌說.
凡洗過躁的.

$^{1561}$就不需要再洗.
只需要把腳一洗.
全身就潔淨了.
一次認罪悔改.
持續認罪悔改.
一次認罪悔改是大草的.
下面那裡是細草的.
只要洗腳就可以了.
時間關係.
我們跳到最後的第五個問題.
末日.
我要使他復活.
這是最後的光景.
最後我們要去到的地方.
會是怎樣的一個光景呢.
可能有人問我.
有沒有片看.
Preview有沒有片看.
竟然有.
是有片看的.
片在哪裡.
在啟示錄.
啟示錄第七章第九節.
看見有許多的人.
沒有人能計算.
從各邦各國各支派各民族各語言來的.
站在寶座和羔羊的面前.
身穿白衣手拿棕樹葉去敬拜上帝.
永遠敬拜上帝.
這班人是甚麼料子.
這些人是從大患難中出來的.
曾用羔羊的血把衣裳洗得潔白.
大患難出來的意思是.
他們在大患難裡面忠心了.
如果不忠心就沒有了.
他們出來的原因是.
因為他們忠心至死忠心.
然後他們是甚麼.
他們不再饑渴不再饑渴.
這就是最後末日的壯觀景象.

$^{1601}$寶座中的羔羊必牧養他們.
當時已經牧養了他們.
已經帶領他們去清草地.
溪水旁已經用五餅二魚去餵養他們.
去到這個終極的動勢.
寶座上的羔羊必牧養他們.
令他們到生命水的泉源.
上帝必擦去他們一切的眼淚.
為甚麼.
因為他們在地上又飢又餓.
大家要知道.
我們去到別人的喪禮都要擦去一切眼淚.
對不起真的與你無關.
你想錯了.
這是說一班因為不再飢餓的人.
願意在地上又飢又餓的人.
上帝要擦去他們一切的眼淚.
因為不再飢不再渴.
永永遠遠與上帝同座.
耶穌說.
我就是生命力量到我這裡來的.
絕不飢餓.
信我的人永不飢渴.
因為我們絕不飢餓永不飢渴.
所以我們可以在這個人生裡面又飢又渴.
為上帝盡忠.
活出一個福音的人生.
我們要祈禱.
我們是天父我們多謝你.
我們多謝你的福音.
我們多謝你巨大衛大的福音.
我們太多時候讀這本聖經.
就是想著自己有什麼機有什麼殼.
於是耶穌說.
耶穌說我有什麼機有什麼殼都可以搞定.
雖然我家裡有飯吃.
心靈裡面有很多需要.
要一個男孩要一個女孩要什麼什麼.
我要生一個兒子.
要一隻貓要一隻狗.

$^{1641}$主人你說我不再飢渴.
於是你供應給我.
神經病.
耶穌說的是一個更偉大壯觀的東西.
在永恆天國裡面.
你不再飢不再渴.
地獄的火沒辦法觸碰你.
沒辦法.
地獄撒旦的魔掌沒辦法勝過你.
於是.
因著福音的救贖.
因著永恆的救贖.
我今天可以捨棄我今天的人生.
可以做個瘋子.
浪費我的一生.
這是保羅所活出的人生.
求主你幫助我們有一個這樣的眼光.
明白因為天國是屬於我們的.
因為生命麵包是屬於我們的.
我們今天可以放下地上的麵包.
去追求天上的麵包.
下個禮拜我們回來.
我們會聽耶穌的呼召.
呼召我們活出天上永恆麵包的人生.
我們敢祈禱我們跟著耶穌的名求.
阿們.
\newpage



\section{使徒行傳 1:4-8}
\label{sec:031821}
\textbf{生活就是使命}~~~~李適清博士
\newline
\newline
連結: \href{http://www.church.com.hk/acms/content.asp?site=cdc\&op=show\&type=product\&code=031821\&layout=sermon}{\texttt{031821}} ~~~~ 語音日期: 2026-01-25
\newline
\newline
\hyperref[sec:031820]{\small{< < < PREV SERMON < < <}}
~
\hyperref[sec:index_chronic]{\small{[返順時目]}}
~
\hyperref[sec:031822]{\small{> > > NEXT SERMON > > >}}
\newline
\newline
使徒行傳 1:4-8
\newline
\begin{longtable}{cl}
\hline
\hline
章節 & 經文 (和合本修訂版)\\
\hline
1:4 & \begin{tabularx}{0.7\textwidth}{X} 耶穌和他們聚集的時候，囑咐他們說：「不要離開耶路撒冷，但要等候父的應許，就是你們聽見我說過的。 \end{tabularx} \\ \\ \relax
1:5 & \begin{tabularx}{0.7\textwidth}{X} 約翰是用水施洗，但過了不多幾天，你們要在聖靈裡受洗。」 \end{tabularx} \\ \\ \relax
1:6 & \begin{tabularx}{0.7\textwidth}{X} 他們聚集的時候，問耶穌：「主啊，你就要在這時候復興以色列國嗎？」 \end{tabularx} \\ \\ \relax
1:7 & \begin{tabularx}{0.7\textwidth}{X} 耶穌對他們說：「父憑著自己的權柄所定的時候和日期，不是你們可以知道的。 \end{tabularx} \\ \\ \relax
1:8 & \begin{tabularx}{0.7\textwidth}{X} 但聖靈降臨在你們身上，你們就必得著能力，並要在耶路撒冷、猶太全地和撒瑪利亞，直到地極，作我的見證。」 \end{tabularx} \\ \\
[1ex]
\hline
\hline
\end{longtable}
$^{1}$弟兄姊妹平安 多謝教會邀請讓我可以跟大家分享訊息.
我是加拿大華人神學院多倫多院長李惜清.
2024年9月我從香港回來多倫多出任神學院院長.
我說回來多倫多 因為我年輕的時候就在多倫多成長.
在多倫多工作 然後在這裡信主.
想不到經過三十多年 上帝又將我帶回來成長的城市.
參與這裡的神學教育工作.
多倫多這裡都有不少的教會活動.
尤其是每年秋天都有培靈會 有宣教大會等等.
但是我不知道宣教給你一個什麼感覺.
會不會宣教就好像你一定要離開你所在的地方.
要改變很多 要去到很遠的地方才能夠做到.
你會不會也有這個印象呢?.
其實宣教對你是一件什麼事情呢?.
如果我問你 你是不是一個宣教士 你會怎樣回答呢?.
我想很多人都搖頭.
我記得我讀神學的時候 有一位英國來的老師.
他教一科短期的 密集的課程.
他說他有幾年曾經離開英國 去印度一間神學院教書.
當他去的時候 弟兄姊妹很熱情地為他祈禱 拉他出去.
稱呼他為一位宣教士 因為他要去印度宣教.
後來他回到英國 他繼續在神學院教書.
所以他的工作性質無論在印度還是英國都沒有改變.
不過他說當他回去的時候 沒人會再稱呼他為宣教士了.
好幾年前 我和神學院一位快要畢業的同學聊天.
說起他畢業之後去哪裡事奉.
他會離開香港去到一個所謂創啟的地方.
也就是不能公開傳福音的地方.
去那裡參與一個機構做行政管理的工作.
因為他本身都是做行政管理的.
他說那些人將他歸類為宣教士.
他說原來我這樣就是宣教士 叫做宣教.
他讀神學之前的專業就是行政管理.
所以對他來說 究竟是去海外工作還是去宣教.
是在職還是在機構服務 根本分不清楚.
所以何謂宣教呢? 誰才是宣教士呢?.
我在神學院是專門研究職場神學的.
我也遇到不少在職的同學.
他們經歷上帝 或者上帝使用他們.
將福音帶到不同的人當中.

$^{41}$上帝用他們的方式遠遠超過傳統宣教士的觀念.
我有一個學生 他的工作經常要去東南亞一些不同的地方.
他說他曾經祈禱求神讓他到不同的地方宣教.
去實踐他的宣教使命.
但神很有趣 沒有帶他離開他的工作.
反而他來了神學院讀職場神學.
過了不久 他發現公司要派他出差.
去一個福音不能公開講的地方.
去那裡建立他們的辦公室.
於是他每個月都去那裡出差.
與那裡的同事分享 做了好朋友.
他說感覺就像跨文化一樣.
原來上帝就是用這樣的方式來回應他心裡的負擔和禱告.
我們發現天國真的很大.
上帝做的事是超乎我們想像的.
所以不在乎我們在什麼位置.
而是在乎上帝將我們每一個人擺在什麼位置.
主要是我們的生命的核心.
我們的心智是否回應上帝.
我們很需要學習在此時此地.
在我們的日常生活當中.
我們的生命就是屬於主的.
在我們的工作 在我們的家庭 在我們的社區 在我們的教會裡.
我們要成為一個用生命去宣講真理的人.
因為對於基督徒來說 我們的生活就是使命.
我們實踐使命的場景就在我們平常的生活裡.
無論你去到哪裡都是一樣.
我不知道你每一天的生活狀況是怎樣.
是不是上班 下班 上學 下課.
又或者每個星期去買菜 每天煮飯.
吃完又洗 洗完又吃.
又要想下一餐吃什麼.
你會怎樣形容你的工作和生活呢?.
如果耶穌基督是我們毀身的.
祂是我們的掌權者.
祂帶領著我們.
祂是我們的君王.
那麼耶穌想我們怎樣生活.
怎樣在工作或生活的勞碌當中來活出使命呢?.
我們試試從聖經找答案.

$^{81}$我們今天讀的經文在《使徒行傳》一章四至八節.
可能大家都很熟悉.
經文說 耶穌和他們聚集的時候 祝福他們說 不要離開耶路撒冷.
要等候父所應許的 就是你們聽見我說過的.
若寒池用水施洗 但不多幾天 你們要受聖靈的洗.
他們聚集的時候 問耶穌說 主啊 你復興以色列國就在這時候嗎?.
耶穌對他們說 父憑著自己的權並所定的時候日期.
不是你們可以知道的.
但聖靈降臨在你們身上 你們就必得著能力.
並要在耶路撒冷 猶太傳地和撒瑪利亞直到地極作我的見證.
選《使徒行傳》是因為這裡正正記載了早期的信徒.
在耶穌復活 顯明福音的大能之後 他們如何生活.
想起《使徒行傳》就很自然想起傳福音.
沒錯的 但如果細心讀《使徒行傳》.
路加記載的早期教會如何跟隨上帝的心意呢?.
他們如何跟隨聖靈的帶領呢?.
他們繼續來發展教會.
所以《使徒行傳》寫的不是很狹義的傳福音.
而是一群願意毀身跟隨耶穌的人.
在耶穌升天之後如何生活.
對他們來說 他們日常的生活就是活出信仰.
信仰就是生活 生活就是按照主的吩咐.
聖靈的帶領去做上帝要我們做的事情.
在《使徒行傳》裡 門徒在敵對的環境當中.
可以說是無人無物的.
他們在羅馬的政權底下 被羅馬的政權克制.
後來教會裡面還出現很多紛爭.
在這樣的環境之下 門徒用他們的生命為主作見證.
為主殉道 日光之下真是無生事.
他們面對內憂外患.
就好像今天我們面對一個未信的世界.
一個分裂的社會.
不同的價值觀 我們就在當中生活.
我們其實還沒有走到早期門徒殉道的情況.
當一群人一起向主呈獻我們的心的時候.
願意將基督的愛帶給我們每一天接觸的人.
去服侍他們 按照主耶穌的吩咐.
聖靈的帶領 就做我們每一天很平常 很平凡的事情.
但是上帝就可以通過我們平凡的生活.
去繼續編寫這個天國的故事.

$^{121}$門徒的行傳 教會的行傳.
所以我很肯定地告訴你.
跟當代的信徒一樣.
主耶穌今天同樣期望我們在座的每一位.
我們可以一起去經歷.
聖靈在我們生命裡面每一天的帶領.
不斷地發現上帝超乎我們想像的工作.
就是在主耶穌今天安放你的位置.
你的家庭 你的職場 你的社會.
每一個你接觸的人 每一個生活的處境.
在一個微不足道 好像不見經傳.
沒有什麼所謂的環境裡面.
聖靈可以用我們.
同樣我們遇到每一個人.
用我們的生命來祝福他們.
來見證耶穌就在我們的生命當中.
在使徒行傳裡面.
其實工作或者日常的生活是故事的中心.
所謂的行傳 就是行為活動.
所以敘述當中使徒是做很多不同的事.
他們有宣講 有醫治.
他們有施捨 有決策.
他們會有團契.
他們會管理當中的金錢 分配食物.
製造做不同的事.
有時他們會開會討論.
他們會閱讀寫作.
他們會面對敵對的群眾.
他們有擁有 也有捨棄.
而且他們有更多的時間在一起.
向上帝禱告和讚美.
使徒就是這樣跟隨聖靈的帶領.
來活出使命.
就是很簡單地生活.
神在他們當中.
讓他們經歷各樣的豐富和供應.
他們每個人有不同的崗位和位置.
他們會去不同的地方.
他們的工作是沒有分重要和不重要.
都是按照上帝的心意去做.

$^{161}$可能今天我們的處境和社會環境.
都複雜很多.
我們做的事和崗位都很多元 很不同.
甚至我們有不同的判斷和演繹.
但是我們毀心跟隨耶穌的心智是一樣的.
就像《仕途行傳》裡那樣.
我們要走在一起.
我們就是為上帝的國而努力.
門徒如何開始他們服侍的生活呢?.
很特別的.
耶穌不是叫他們馬上做.
我們會馬上想做些什麼呢?.
如何去做呢?.
但是耶穌吩咐他們.
是要他們等候.
耶穌吩咐他們在耶路撒冷等.
對門徒來說這是很難理解的.
因為對門徒來說.
耶路撒冷是他們受到敵視和迫害的地方.
在最危險的地方.
要他們等.
等什麼呢?.
上帝的計劃就是這樣的.
祂的應許要在耶路撒冷那裡成就.
祂的應許就是聖靈要來到.
要臨到這群門徒身上.
就在聖子被釘十字架的這個城市.
在那裡聖靈降臨.
為聖子作見證.
表明耶穌就是上帝的兒子.
真理的靈要讓人知罪.
要帶來公義和審判.
而這件事的起點就是耶路撒冷.
也是門徒因為耶穌被捉拿.
他們四散彼得不認主.
門徒離棄甚至出賣耶穌的地方.
耶穌就是要在他們跌倒的地方.
重新兼顧他們.
在這裡之前.
耶穌其實一直有提及保衛師聖靈.

$^{201}$要來到的這件事.
而在那裡.
祂再一次提到門徒要領受聖靈.
這一次的洗禮.
再不是約翰外在用水的施洗.
而是一個內在的.
屬靈的.
聖靈要臨到.
讓上帝的子民每一個都有聖靈的同在.
上帝的子民就要走在一起.
成為教會群體的一份子.
也就是說成為基督的身體.
預備好可以扶持著.
其實這一切門徒也不是很明白.
不過他們願意等.
他們聽耶穌的話.
他們就在那裡等.
我的本科.
我第一個學位是讀會計.
我之後就入了會計師樓工作.
就在多倫多這個城市.
我一邊工作一邊考試.
做這行的人就知道.
工作的時間很長.
尤其是報稅的季節.
埋數的季節.
有上班沒有下班.
還要繼續讀書考試.
其實核數師樓.
是有很多不同的老闆.
因為是合夥人.
合夥制度.
當我幫不同的客戶去核數的時候.
就有機會幫不同的老闆去工作.
後來我就回到香港.
在同一間會計師樓工作.
有時候在香港工作的時候.
還會回到國內去公幹或者開會.
每一次去公幹的時候.
我都會把握機會.

$^{241}$從這些合夥人.
這些長輩身上學習.
觀察一下他們.
怎樣跟客戶開會.
怎樣將香港的核數標準.
解釋給國內人去聽.
還有我最想從他們身上學到.
就是他們怎樣去看東西.
在他們豐富的人生經驗裡.
他們怎樣去洞察人生.
未必個個都是基督徒.
但是他們有他們的經歷.
雖然我未必認同他們每一個人的看法.
不過都有很多東西可以學得到.
我記得有一位上司.
他在公司工作了很多年.
當我跟他聊天的時候.
我發覺他有很多的感慨.
他是一位半光頭的地中海男士.
有一次他一面談戲.
他一面跟我說.
我在這裡工作了幾十年.
我的人生,我的時間,我的青春,我的頭髮.
都奉獻給這間公司了.
弟兄姊妹.
你的一生的年日.
你的青春.
甚至你的頭髮.
要奉獻給什麼呢?.
要奉獻給誰呢?.
你願不願意將主權交回給主耶穌.
讓祂做我們生命的主.
讓祂去差遣我們.
有時候我們誤會了.
我們回應神.
祂差遣我們一定有很遠.
或者我們很不願意的地方.
但其實我經驗告訴我.
耶穌差遣我們.
是去到一個最適合我們的地方.

$^{281}$就是我們很自然的.
在我們的日常生活裡.
我們很自然的.
來活出信仰.
就是活出使命去回應祂.
其實聖經裡.
門徒要去等.
要去忍耐.
因為要我們知道.
上帝掌權的世界.
等候主的吩咐.
然後按祂的心意去做.
是信得過祂給我們的.
是最好.
祂差遣我們去的地方.
也是最適合我們的.
耶穌叫門徒等.
不過門徒不知道等什麼.
他們也不太明白.
其實耶穌是要重整我們的焦點.
祂要門徒看到.
不是耶路撒冷的危險.
門徒看到的是.
他們要降伏在主耶穌的吩咐之下.
專注在這位復活的主身上.
單單仰望神.
門徒按照他們的理解.
他們的結論是.
耶穌的國度這一天將要來到了.
所以他們很期待.
在第六節他們問.
主啊,你復興以色列國.
就在這時候嗎?.
他們跟了耶穌這麼久.
他們仍然不明白.
天國其實不是他們想像的.
但在這裡.
耶穌沒有直接說他們不對.
或者更正他們.
因為祂的確為門徒.

$^{321}$留下了祂要再來的盼望.
不過天國何時來到呢?.
卻是他們不能夠知道的.
反而耶穌的答案.
是將門徒的目光帶回現在.
就是此時此地.
他們的使命.
對於世界來說.
他們就是耶穌的見證人.
而他們要在當地耶路撒冷.
四周圍.
猶太傳地和撒瑪利亞.
和所有地方.
直到地極去實踐這個使命.
這個使命是一直延續著.
給整個教會群體的.
所以你和我今天都一樣有份.
不過關鍵不是我們要.
放下甚麼要去哪裡.
以為自己可以為神做些甚麼.
關鍵是主耶穌將我們放在哪裡.
祂讓我們面對甚麼環境.
要我們去做些甚麼.
去回應祂.
可能我們都像門徒一樣.
很期待主耶穌要像我們想像一樣.
帶領我們去更好的環境.
或者可以去哪裡.
為神做些甚麼.
壞消息是主耶穌未必會按照你的想像去做.
又或者可以說.
祂多數都不是我們估計的.
這樣去想像去做.
不過好消息是.
如果你願意放手.
放下自己很執著想要的東西.
我們一起去降伏在耶穌面前.
讓祂重整你的焦點.
你就一定會有全新的經歷.
在上帝安放你的位置.

$^{361}$就是此時此地.
你已經身處的那個位置.
你就會見到上帝給你的機會.
要怎樣來使用它.
記得我年輕的時候.
我很努力讀書.
讀到大學畢業.
就在安大略省這裡.
然後我來到多倫多.
我一邊做事一邊考會計試.
最後上帝的恩典.
我就在加拿大考到會計的牌.
年輕的我就是這樣想.
以後都不用考試了.
不知道你怕不怕考試.
誰知過了幾年.
上帝帶領我回到香港工作.
因為香港有本地的稅務和法例.
所以我又要考試了.
我要考兩科本地的稅法和商業法.
才可以拿牌.
所以我一邊讀一邊鼓勵自己.
讀完這次最後努力了.
以後都不用再考試了.
誰不知上帝有更奇妙的安排.
祂後來呼召我放下我的工作.
去到神學院服侍祂.
祂沒有因為我不馬上去回應.
祂放棄我.
過了好幾年.
我終於都回應祂的呼召.
去到神學院.
最初我以為讀書就回教會事奉.
誰知到了第二年.
神學院的老師又挑戰我.
再去外國讀書.
我當時都有很大的掙扎.
想不到走上帝的路.
會有那麼多令我意想不到的事.
我們的上帝就是如此奇妙.

$^{401}$祂呼召我們每一個人去跟隨祂.
不過祂不會強迫我們.
當我們願意放下我們自己.
去尋求祂的心意的時候.
祂就會帶領著我們.
讓我們看到祂給我們每一個不同的經歷.
在最好的時候.
祂幫我們去開路.
揭開我們意想不到的美好的一頁.
對我們來說其實等什麼呢?.
馬上就可以做了.
對我們來說上帝的方法真是又慢又不理想.
就好像早期的殉道者.
可能我們今天都會覺得他們很蠢.
因為他們的堅持帶來迫害.
又解決不了問題.
我們就是這樣的.
我們解決問題的辦法很多.
我們會想很多不同的方法.
我們會做上帝的那一份.
我們覺得這樣做才是最好.
很有注點.
我們忘記了其實上帝的心意.
未必是我們想像的.
我不知道你是不是也很有想法.
又或者你也在想.
很期待上帝為你去做一些事.
或者開一些路.
或者上帝在我們身上的心意.
祂的帶領不是我們所想像的.
如果是這樣的時候.
你是否仍然信得過祂呢?.
你會不會可以開始觀察一下.
祂給你一些不同的機會.
不同的出路.
或者祂為我們所擺設的.
是超乎我們想像.
高過我們所想所求.
使徒行傳一開始就告訴我們.
要實踐使命.

$^{441}$最關鍵不是我們有能力去做些什麼.
所以不要緊的.
你覺得做不到.
或者不知道怎麼做.
我們是需要領受聖靈.
然後按照聖靈的帶領.
只有祂住在我們裡面.
才可以讓我們的屬靈生命有生命力.
有力量去做耶穌基督的見證人.
有很多時候我們想起日常生活.
職場工作.
都想起要傳福音做好見證.
是很好的.
不過很多時候恐怕我們誤解了.
什麼是做好見證.
以致我們很怕要去做些什麼.
好見證很多時候我們理解.
就是做好人做好事.
或者工作做得好.
是其他人覺得我們是一個好的基督徒.
但是如果做好見證.
只是做好人做好事做好份工.
非基督徒都會做.
可能做得比我們更加好.
有時候當我們按照.
其他人對基督徒的要求去回應的時候.
其實是誰在定義我們是否一個好的基督徒.
其實是這個世界的人在定義.
但是如果你看回《使徒行傳》這裡.
「門徒為主作見證」.
作見證是什麼意思呢?.
就是當主耶穌受死 復活 升天之後.
大家都以為祂已經死了.
早期的使徒要到處去告訴人們.
這位你們見到鼻釘十字架的耶穌.
祂復活了 耶穌活著.
所以早期的《門徒傳》是一個.
耶穌死後已經復活的福音.
於是我們今天說要為主作見證.
不僅是做好人做好事.

$^{481}$而是我們要見證耶穌基督活著.
祂如何活著呢?.
祂活在我們每個人的生命當中.
這樣的話 在職場裡做好見證.
就不是做好人做好事.
而是要讓其他人見到.
你生命裡有活著的耶穌.
讓人見到你的生命是不一樣的.
不在乎一次,兩次能否做到.
我們想做的所謂好事.
而是在長期的相處下.
讓人見到基督徒的生命是如何不同.
是有耶穌的生命.
你的生命要吸引人.
去認識你生命當中的主耶穌.
基督徒共同的使命的確是為主作見證.
不過這不是口說的.
而是我們每一天的生活裡.
生命的演繹.
無論你在家裡,工作裡,社會裡,教會裡.
都是一樣的.
這個共同的使命不容易的.
因為聖經的標準很高.
我們都知道我們達不到.
這個共同使命只有靠聖靈去做得到.
是聖靈改變我們的生命和整個群體.
以致我們可以分享工作的成果.
善用我們的資源,才幹和影響力.
去發揮基督徒群體的炎和光的特質.
影響我們身邊的人.
影響我們整個的文化.
這一切都不是靠我們自己可以做得到.
所以我們一起來發揮這種影響力.
每一個人做的其實很簡單.
我們就是努力去緊貼耶穌基督的說話.
去活出一個上帝兒女的身份和樣式.
有時候我們會成功一些.
會開心,感恩.
但有時候我們會軟弱,失敗,跌倒.
不要緊,我們可以告訴其他人.

$^{521}$我們怎樣去回轉同樣是一個很好的見證.
我們的信仰正正就是一個回轉.
可以悔改,復和的福音.
所以主耶穌已經選擇了你.
無論你是什麼年紀,什麼位置.
就在你現在身處的家庭,日常生活的處境裡.
就是你被主差遣進入世界以生命去演繹福音的場景.
有時候我們以為聖靈臨到.
一定是一些很神奇,很特別的經歷.
當然這些都可能會有的.
但要知道聖靈就是這樣很自然地活在我們當中.
去做感動,帶領,教導的工作.
聖靈不單止個別與我們同在.
它也在我們整個教會群體當中.
來幫助我們學習如何彼此相愛.
聖靈加力,賜力是何時發生的呢?.
不是在我們唱詩歌,舉高雙手.
也不是我們苦苦哀求的時候.
聖靈是很樂意,時刻與我們同在.
是因為我們願意跟隨耶穌,活出神的話語.
當我們願意活出主的道.
按照聖經的吩咐去做的時候.
聖靈就會加力,帶領我們.
就像早期的使徒一樣.
他們靠著聖靈的能力去活出神的話.
我們不一定在說一些人看起來很偉大,很特別的事.
只是很簡單的.
我的生命是屬於神的.
於是我就要活出主的道.
很自然地我們就成為一個有感染力.
彼此相愛的群體.
我們的影響力,生命裡的主.
就會感染我們身邊的人.
在2010年我在英國愛丁堡大學.
拿到神學博士的學位.
回到香港我的神學院教書.
對於我這個新的老師來說.
第一個學期就要教三科全新的科目.
感覺是很不容易的.
加上我小時候.

$^{561}$其實我是一個很害羞,很不喜歡說話的小朋友.
上帝就很特別.
今天讓我每天都說很多話.
還要對著很多人說.
我小時候覺得我一定不行.
所以我小時候就想.
長大後做些什麼好呢?.
我就想我做什麼都好.
千萬不要做老師.
因為我覺得不可能做得到.
上帝的恩典很足夠.
不過當我初初開始去做一個神學院老師的時候.
我就真的很不習慣.
很多人他們教聖經科,教神學科.
用很多參考書,神學書來用.
但是我教的是職場神學.
當時都是一個非常新的學科.
究竟要怎樣教呢?.
職場神學的科目.
也是在晚上來教的.
因為日間同學會上班.
然後下班他們就會上課.
所以每一次我教完書.
就已經是晚上了.
每一次我回家的時候.
除了都會覺得很累之外.
我心裡面就會有一種抖結.
因為我覺得我還沒有教得很好.
這是我當時小小的一個職場掙扎.
我不知道你有這些掙扎的時候.
你會做些什麼.
我只有一招教你.
很簡單的.
當然是祈禱.
因為上帝你帶領我輾轉走到這個位置.
我根本覺得自己做不到.
究竟你的心意是怎樣呢?.
你對於一班在職的弟兄姊妹.
你的心意是什麼呢?.
有一晚.

$^{601}$當我很簡單獨自祈禱的時候.
我心裡面就湧出了一幅圖畫.
我找到最相似的一幅圖畫.
當時我第一個反應.
我心裡面就跟神說.
上帝啊 你的創造真是很美.
我在想這個夜空.
這個星空是非常之美.
但是就好像心裡面有一個感動.
跟我說這樣.
就是這些星星.
就是我安放在不同位置的弟兄姊妹的生命.
這些弟兄姊妹被拆遣到各行各業裡.
後來我就將這個儀仗.
跟我的在職學生來分享.
這幅圖畫感動了很多弟兄姊妹.
因為它是說不單止僕者宣教士.
被呼召和拆遣.
其實主耶穌是呼召拆遣.
每一個願意回應祂的門徒.
去到社會裡不同的行業.
生活的不同處境裡.
很簡單的讓我們的生命為主發光.
每個人都很重要.
因為如果你在你的位置裡不發光的話.
你周圍就黑暗.
因為這是一個黑暗的世界.
上帝都很清楚.
但是每一個人同樣道理.
又不是很重要.
因為一顆星不發光.
還有千千萬萬的星.
所以這幅圖畫是說我們回應上帝的呼召.
被祂拆遣去到祂的心意.
我們所屬的一個位置.
我們那個小小的位置裡.
我們就要將自己看得合乎中度.
讓我們的生命為神來發光.
有時候你的星是光一點.
我的星是暗一點 因為我軟弱.

$^{641}$你可以鼓勵我.
有時候我又鼓勵你.
不知道你有沒有發現.
在這幅圖畫裡.
你都有一個很獨特的位置.
是上帝安放你的位置.
我們每個人都不是最重要.
但是同時每一個人都很重要.
你願不願意回應主耶穌給你的使命.
成為照亮夜空的其中一顆星.
我們就是這樣按照主耶穌的差遣.
進入我們的生活裡.
不同的範疇裡.
生活就成為我們信仰實踐的一個很重要的場景.
而我們的信仰就和我們的生活分不開.
然後我們每個星期聚集在一起.
我們一起去敬拜 一起去迎見上帝.
然後彼此鼓勵分享.
最後被主差遣到不同的地方.
分散到社會不同的位置裡.
這個狀態其實和早期的仕途一樣.
一直都沒有改變.
弟兄姊妹.
你願不願意回應主耶穌的呼召.
做一個真正跟隨耶穌.
回應呼召被祂差遣的門徒.
其實很簡單的.
就是在你現時的位置.
在你的家庭 在你的社區.
在你的工作 在你的教會裡.
我們就是以生命來活出信仰.
讓我們同心圖告.
因為是你親自來呼召我們每一個.
願意我們去回應你的時候.
我們就被你得著.
你差遣我們.
按照你的心意 去到不同的處境.
就在我們身處的家庭 工作 生活.
不同的環境裡.
我們很自然地用我們的生命來回應你.

$^{681}$求你給我們有信心和勇氣.
願意這樣來實踐使命.
因為我們是屬於你的.
我們的生活就是使命.
阿們.
\newpage



\section{以弗所書 1:1-23}
\label{sec:031822}
\textbf{以弗所書 (1) 導論,求啟示的禱告}~~~~劉少康牧師
\newline
\newline
連結: \href{http://www.church.com.hk/acms/content.asp?site=cdc\&op=show\&type=product\&code=031822\&layout=sermon}{\texttt{031822}} ~~~~ 語音日期: 2026-03-18
\newline
\newline
\hyperref[sec:031821]{\small{< < < PREV SERMON < < <}}
~
\hyperref[sec:index_chronic]{\small{[返順時目]}}
~
\hyperref[sec:031823]{\small{> > > NEXT SERMON > > >}}
\newline
\newline
以弗所書 1:1-23
\newline
\begin{longtable}{cl}
\hline
\hline
章節 & 經文 (和合本修訂版)\\
\hline
1:1 & \begin{tabularx}{0.7\textwidth}{X} 奉神旨意作基督耶穌使徒的保羅，寫信給在以弗所的眾聖徒，就是在基督耶穌裡忠心的人。 \end{tabularx} \\ \\ \relax
1:2 & \begin{tabularx}{0.7\textwidth}{X} 願恩惠、平安從我們的父神和主耶穌基督歸給你們！ \end{tabularx} \\ \\ \relax
1:3 & \begin{tabularx}{0.7\textwidth}{X} 願頌讚歸給我們主耶穌基督的父神。他在基督裡曾把天上各樣屬靈的福氣賜給我們。 \end{tabularx} \\ \\ \relax
1:4 & \begin{tabularx}{0.7\textwidth}{X} 因為他從創世以前，在基督裡揀選了我們，使我們在他面前成為聖潔，沒有瑕疵，滿有愛心。 \end{tabularx} \\ \\ \relax
1:5 & \begin{tabularx}{0.7\textwidth}{X} 他按著自己旨意所喜悅的，預定我們藉著耶穌基督得兒子的名分， \end{tabularx} \\ \\ \relax
1:6 & \begin{tabularx}{0.7\textwidth}{X} 使他榮耀的恩典得到稱讚；這恩典是他在愛子裡白白賜給我們的。 \end{tabularx} \\ \\ \relax
1:7 & \begin{tabularx}{0.7\textwidth}{X} 我們藉著這愛子的血得蒙救贖，過犯得以赦免，這是照他豐富的恩典， \end{tabularx} \\ \\ \relax
1:8 & \begin{tabularx}{0.7\textwidth}{X} 充充足足地賞給我們的。他以諸般的智慧聰明， \end{tabularx} \\ \\ \relax
1:9 & \begin{tabularx}{0.7\textwidth}{X} 照自己在基督裡所立定的美意，使我們知道他旨意的奧祕， \end{tabularx} \\ \\ \relax
1:10 & \begin{tabularx}{0.7\textwidth}{X} 要照著所安排的，在時機成熟的時候，使天上、地上、一切所有的，都在基督裡面同歸於一。 \end{tabularx} \\ \\ \relax
1:11 & \begin{tabularx}{0.7\textwidth}{X} 我們也在他裡面得了基業；這原是那位隨己意行萬事的神照著自己的旨意所預定的， \end{tabularx} \\ \\ \relax
1:12 & \begin{tabularx}{0.7\textwidth}{X} 為要使我們，這些首先把希望寄託在基督裡的人，頌讚他的榮耀。 \end{tabularx} \\ \\ \relax
1:13 & \begin{tabularx}{0.7\textwidth}{X} 在基督裡你們聽見真理的道，就是那使你們得救的福音，你們也信了他，就受了所應許的聖靈為印記。 \end{tabularx} \\ \\ \relax
1:14 & \begin{tabularx}{0.7\textwidth}{X} 這聖靈是我們得基業的憑據，直等到神的子民得救贖，使他的榮耀得到稱讚。 \end{tabularx} \\ \\ \relax
1:15 & \begin{tabularx}{0.7\textwidth}{X} 因此，我既然聽見你們對主耶穌有信心，對眾聖徒有愛心， \end{tabularx} \\ \\ \relax
1:16 & \begin{tabularx}{0.7\textwidth}{X} 就不住地為你們感謝神，禱告的時候常常提到你們， \end{tabularx} \\ \\ \relax
1:17 & \begin{tabularx}{0.7\textwidth}{X} 求我們主耶穌基督的神，榮耀的父，把那賜人智慧和啟示的靈賜給你們，使你們真正認識他， \end{tabularx} \\ \\ \relax
1:18 & \begin{tabularx}{0.7\textwidth}{X} 照亮你們心中的眼睛，使你們知道他呼召你們來得的指望是甚麼，他在聖徒中所得榮耀的基業是何等豐盛， \end{tabularx} \\ \\ \relax
1:19 & \begin{tabularx}{0.7\textwidth}{X} 並知道他向我們這些信的人所顯的能力是何等浩大，這是照他的大能大力運行的。 \end{tabularx} \\ \\ \relax
1:20 & \begin{tabularx}{0.7\textwidth}{X} 這大能曾運行在基督身上，使他從死人中復活，又使他在天上坐在自己的右邊， \end{tabularx} \\ \\ \relax
1:21 & \begin{tabularx}{0.7\textwidth}{X} 遠超越一切執政的、掌權的、有權能的、統治的和一切有名號的；不但是今世的，連來世的也都超越了。 \end{tabularx} \\ \\ \relax
1:22 & \begin{tabularx}{0.7\textwidth}{X} 神使萬有服在他的腳下，又使他為了教會作萬有之首； \end{tabularx} \\ \\ \relax
1:23 & \begin{tabularx}{0.7\textwidth}{X} 教會是他的身體，是那充滿萬有者所充滿的。 \end{tabularx} \\ \\
[1ex]
\hline
\hline
\end{longtable}
$^{1}$各位弟兄姊妹平安.
今天要和大家講.
在《新約聖經》當中.
一本很重要的書.
就是保羅他寫信書信的裡面.
一封很神學性.
有人形容是保羅這麼多封信裡面的花冠.
好像保羅寫了很多封信都好像花一樣.
而這封信是最突出的.
其中一個花冠.
好像皇冠那樣.
這封信就是「以忽所恕」.
每次我讀「以忽所恕」的時候.
聖靈都觸動我的心.
我很喜歡這本書.
但是這本書又是很深奧的.
所以我向神祈禱.
我說神先幫助我去明白.
又幫助我怎樣講解.
以致你們每一個都明白了.
「以忽所恕」是講什麼的.
我想邀請大家跟我一起低頭做一個祈禱.
親愛的天父.
我恭敬將這一段時間交託給你.
求你的聖靈在我們每一個人的心中運行.
將你這麼寶貴的話.
好像很活潑的道一樣.
能夠很流暢地進入每一個人的內心.
因為聖靈會在我們每一個人的心中工作.
觸動我們.
感動我們.
求天父您親自工作.
求耶穌賜福給我們.
求聖靈親自動功.
我們這樣恭敬仰望祈禱.
奉耶穌基督的命求.
阿們.
這封信「以忽所恕」被稱為監獄詩信裡面.
其中的一封.
原來保羅的書信裡面有四封.

$^{41}$都稱為監獄書信.
保羅大概是在第三次行程之後.
他就在羅馬坐監了.
其實也可以說是一種拘留.
在保羅的身邊.
就會有起碼四個士兵看守著他的.
有些人說保羅可能是被鎖上鎖鏈的.
所以他是完全沒有自由的.
但當保羅第二次在保羅坐監的時候.
他的自由度就好一些了.
就是所謂在家裡面的監禁.
禁止你出外.
但在家裡可以走來走去.
所以就有不同的說法.
所以保羅在監獄書信裡面.
就是「以忽所恕」.
「肥納比書」.
「歌羅西書」.
和「肥利門書」.
這四卷監獄書信.
而四卷監獄書信裡面.
「以忽所恕」是最出色的.
很特別.
這封信的名字是寫給「以忽所教會」的.
大家要記得.
保羅第一次和第二次行程裡面.
都去過「以忽所」.
特別在第二次行程裡面.
根據《使徒行傳》的記載.
保羅在「以忽所」起碼停了三年這麼久.
所以可想而知.
保羅和這間教會的弟兄姊妹一定很熟.
因為停了三年.
在三年裡他時時說到.
好像成為這間教會的牧師一樣.
因此當保羅要離開「以忽所」的時候.
他特別要找他的真兒子提摩泰來接班.
讓提摩泰去到「以忽所」.
好像接班做牧師一樣.
等於主任牧師退休.

$^{81}$要找一個接班人來做牧師一樣.
所以保羅很重視「以忽所」這間教會.
但是很特別.
這封信好像不是單單寫給「以忽所」這間教會.
所以一般的解經家就說.
這封信是寫給小亞,西亞.
即是Asia Minor的宗教會.
所以這封信很特別.
沒有個人的文案.
通常保羅一開始就會說我保羅向誰問安.
或者書信結束的時候.
我又問候這個又問候那個.
又問候誰又幫我問候誰.
但是這封信完全沒有的.
所以可以說是一封傳閱的信.
由於「以忽所」書是保羅托推基古帶去「以忽所」的.
「以忽所」的.
而保羅有另外一卷書.
就是哥羅世書.
都是推基古帶去的.
所以有人就估計.
同時既然推基古出發.
保羅叫他將兩卷書帶到這個地區那裡.
於是推基古就接受了這個任務.
就去帶這兩卷書.
「以忽所」書是在神學上有一定的深度.
有一次在台灣我去一間教會.
有些弟兄姊妹就跟我說.
劉牧師你今天來聽到真是好.
今天來港道的牧師.
他懂得背很多本聖經的書卷.
我一聽多羨慕.
這個牧師這麼厲害.
可以背完聖經的一些書卷.
那個會友還說.
他不單背.
是由第一章背到最後那章.
那我真是很羨慕.
如果叫我背.
我為什麼背不到.

$^{121}$可能背到一些熟的金句.
原來這位牧師很特色.
他每次港道可以隨口就背.
「以忽所」書三章十四到十七節.
是這樣說的.
然後在歌羅西書十二章七到八節是這樣說.
他真是滾瓜爛熟.
如果我們港道可以這樣.
隨時可以引用一些經文.
那多好.
於是後來神都給我一個感動.
祂說你背「以忽所」書.
於是我就開始學這個牧師.
由「以忽所」書第一章第一節我開始.
我嘗試去背.
可能年紀真是太大了.
每天讀很多次.
又聽很多次「以忽所」書.
但背來背去都未必完全連得上.
透過這樣的聽.
這樣的背誦.
我愛上了這本書.
這本書真是很有能力.
帶著神的話.
所以今天我先說「以忽所」書的悼論.
「以忽所」書是怎樣一本書?.
好像保羅其他書信一樣.
頭半部是比較理論式的.
神學性的探討.
然後後半部.
「以忽所」書有六章.
即是一到三章是理論性佔多.
四至六章是實踐生活性比較多.
而「以忽所」書有一個特色.
就是它神學的論述是包含很闊的.
有關於神論的.
有關於基督論的.
有關於教會論的.
有關於信徒的生活的.
特別教會論是「以忽所」書裡面的精華.

$^{161}$所以如果你想了解今天我們教會應該怎樣來做.
教會應該怎樣組織.
教會裡面有甚麼的積分.
請你讀「以忽所」書.
因為保羅在神裡面領受了一些啟示.
教我們怎樣組織教會.
以下半部實踐的部分.
保羅就談到很多教會裡面弟兄姊妹怎樣同心.
我們在基督裡面要合一.
要同心.
然後又講到夫妻關係.
又講到主僕關係.
又講到要孝敬父母等等.
所以一方面好像很深的神學.
但另一方面卻是平易近人.
所以我希望你都愛上「以忽所」書.
我們就開始先講第一章.
第一章我叫它做「求啟示」的土局.
一開頭保羅就將我們可以帶到宇宙的空間.
從很高空那裡俯瞰下來.
你就會看見神創世以前的計劃是怎樣.
保羅心存感恩.
所以有人又形容「以忽所」書是新約裡面的詩歌書.
大家都知道舊約裡面有一卷很出名的詩歌書.
當然是詩篇.
新約二十七卷哪一卷是一些頌讚.
用詩歌讚美神的呢?.
原來就是這卷「以忽所」書.
大家都知道詩篇是可以唱出來的.
如果有人有神給他音樂的感動.
又有天分的話.
可以將「以忽所」書寫成詩歌.
讓我們唱的時候順便背「以忽所」書.
這樣多好啊.
所以一開頭.
保羅就向這位至高創造宇宙萬有.
三位一體的神.
發出他的頌讚.
「原頌讚歸於我們天上的父神.
祂在創世以前已經揀選了我們」.

$^{201}$所以根據「以忽所」書第一章.
「揀選」這個字是出現了很多次的.
我不開估到底「揀選」這個字出現了多少次.
我想請你每一次.
你看著你的中文聖經.
當出現「揀選」的時候.
你畫著它.
當再出現「揀選」你又畫著它.
你數數看出現了多少次.
保羅說是神在創世以前.
已經揀選了我們.
第一是關於三位一體的神.
裡面的父神.
第二是在基督裡面揀選了我們.
出現父神.
出現聖子基督.
然後接著又藉著聖靈.
在我們裡面工作.
聖父,聖子,聖靈.
保羅的神學思想真的很超越.
原來在創造世界以前.
三位一體的神都存在.
所以很多人覺得.
舊約只講耶和華,只講神.
到新約耶穌基督先來.
然後到耶穌基督臨走要升天.
他就說我猜聖靈來.
好像三位一體分了不同的階段一樣.
不是,聖經裡面很早講.
我們的神就是三位一體的神.
是 Trinity.
Three in one.
One God.
But in three persons.
一個的神,祂有三個位格.
而這三個位格.
父,子,靈是同榮同尊的.
他們同樣位置.
同樣尊崇.
同樣有能力.

$^{241}$我們的神是無所不能的神.
是無所不在的神.
是無所不知的神.
所以這就是我們信仰裡面神學的基礎.
所以保羅由第一章一開始.
就觸及這麼深的神學理論.
因此我們要知道.
保羅很愛這班弟兄姊妹.
將他們帶到一個很高的地方.
叫他們看看原來神就是這麼奇妙的一位神.
這麼愛我們的神.
他老祖在創世之前已經選了我們.
於是基督徒就問.
牧師,是不是神預定了誰得救,誰不得救?.
這就是所謂預定論.
如果預定了我不得救,那我信你也沒用.
如果預定了祂得救.
那祂就每晚周圍犯罪.
到最後祂信耶穌就行了.
全錯了.
神沒有錯.
什麼叫預定?.
我想從另外一個角度說.
大家先不要想有沒有預定.
剛才我說過神是無所不知的神.
祂是無所不能的神.
祂也是無所不在的神.
因為神那種超越.
在神來說是沒有時間之分的.
沒有昨日history.
沒有現在present.
也沒有future.
因為一千年前和現在.
和一千年後其實神起高空望下來.
其實連在一起.
那件事未出現神已經知道.
那件事結局是怎樣神已經知道.
這就是全知.
所以記不記得哥林多前書十三章說.
我們現在好像對著鏡子.

$^{281}$說愛騙.
哥林多前書十三章說.
到那一天就知清.
就全知曉.
有一天主耶穌接我們去到天上.
你在天上望下來.
是沒有了那些所謂空間的感覺.
因為在天上高空一望下來.
什麼都看到了.
原來神就是這樣知曉.
所以不是預定論的問題.
是神的全知.
在神已經超越時間.
所以聖經說.
神看千年好像一天.
一天又好像千年.
所以在神裡面沒有了時間.
因為祂全部知道了.
原來你和我是神認識的.
有時我自己都說.
我未出母胎.
我媽媽未生我.
神已經認識我.
而且神選召我做祂的僕人.
就是這個意思.
神就是一個揀選的神.
祂揀了我.
又揀了你.
又揀了祂.
我們都是被選召的一群.
然後保羅因著神這麼大的愛.
所以他就向神祈禱.
發出一個很偉大的求啟示的祈禱.
我求神讓我來真知道祂.
我不知道你今天是否真認識這位這麼偉大的神.
我同樣發出這個禱告.
求神啟示你.
讓你看見祂的奧秘.
看見神的莫大的奇妙.
因此保羅在第一章的下半部.

$^{321}$就發出這個求啟示的禱告.
他說我向神祈求.
有三樣東西.
你們能夠真知道神.
第一就是神的恩召.
祂的召我們那個救恩.
是有何等的指望.
因著我們得了救恩.
我們有很大的指望.
第二.
祂說因著我們對祂的認識.
我們的能力.
我們有豐盛的榮耀.
我們對神的認識.
是有豐盛的榮耀.
第三.
這位神在信徒身上.
所彰顯的能力.
是何等浩大.
出現三個何等.
第一是何等的指望.
第二是何等的榮耀.
第三是能力何等浩大.
大家看不看到保羅求什麼.
他求神給我們三樣東西.
都是何等的.
神給我們的恩召.
是何等的指望.
祂給我們的威勳.
是何等豐盛的榮耀.
第三.
祂給我們的能力.
是何等的浩大.
我希望各位今天聽完這個悼論.
和第一章的訊息.
你回去很詳細地讀第一章.
我求感動我的靈.
同樣地感動你.
願天父賜福給各位.
\newpage



\section{約翰福音 5:1-18}
\label{sec:031823}
\textbf{屬靈的力量}~~~~鍾智愷傳道
\newline
\newline
連結: \href{http://www.church.com.hk/acms/content.asp?site=cdc\&op=show\&type=product\&code=031823\&layout=sermon}{\texttt{031823}} ~~~~ 語音日期: 2026-03-15
\newline
\newline
\hyperref[sec:031822]{\small{< < < PREV SERMON < < <}}
~
\hyperref[sec:index_chronic]{\small{[返順時目]}}
~
\hyperref[sec:031824]{\small{> > > NEXT SERMON > > >}}
\newline
\newline
約翰福音 5:1-18
\newline
\begin{longtable}{cl}
\hline
\hline
章節 & 經文 (和合本修訂版)\\
\hline
5:1 & \begin{tabularx}{0.7\textwidth}{X} 這些事以後，到了猶太人的一個節期，耶穌上耶路撒冷去。 \end{tabularx} \\ \\ \relax
5:2 & \begin{tabularx}{0.7\textwidth}{X} 在耶路撒冷，靠近羊門有一個池子，希伯來話叫畢士大，旁邊有五個柱廊； \end{tabularx} \\ \\ \relax
5:3 & \begin{tabularx}{0.7\textwidth}{X} 裡面躺著許多病人，有失明的、瘸腿的、癱瘓的。等候水動， \end{tabularx} \\ \\ \relax
5:4 & \begin{tabularx}{0.7\textwidth}{X} 因為有天使按時下池子攪動那水，水動之後，先下水的人無論患甚麼病都能得痊癒。 \end{tabularx} \\ \\ \relax
5:5 & \begin{tabularx}{0.7\textwidth}{X} 在那裡有一個人，病了三十八年。 \end{tabularx} \\ \\ \relax
5:6 & \begin{tabularx}{0.7\textwidth}{X} 耶穌看見他躺著，知道他病了很久，就問他：「你要痊癒嗎？」 \end{tabularx} \\ \\ \relax
5:7 & \begin{tabularx}{0.7\textwidth}{X} 病人回答他：「先生，水動的時候，沒有人把我放在池子裡；我正要去的時候，別人比我先下去了。」 \end{tabularx} \\ \\ \relax
5:8 & \begin{tabularx}{0.7\textwidth}{X} 耶穌對他說：「起來，拿起你的褥子走吧！」 \end{tabularx} \\ \\ \relax
5:9 & \begin{tabularx}{0.7\textwidth}{X} 那人立刻痊癒，就拿起自己的褥子走了。那天是安息日， \end{tabularx} \\ \\ \relax
5:10 & \begin{tabularx}{0.7\textwidth}{X} 所以猶太人對那被治好了的人說：「今天是安息日，你拿褥子是不合法的。」 \end{tabularx} \\ \\ \relax
5:11 & \begin{tabularx}{0.7\textwidth}{X} 他卻回答他們：「那使我痊癒的人對我說：『拿起你的褥子走吧！』」 \end{tabularx} \\ \\ \relax
5:12 & \begin{tabularx}{0.7\textwidth}{X} 他們問他：「對你說『拿起褥子走』的是甚麼人？」 \end{tabularx} \\ \\ \relax
5:13 & \begin{tabularx}{0.7\textwidth}{X} 那治好了的人不知道那人是誰，因為那裡人很多，耶穌已經躲開了。 \end{tabularx} \\ \\ \relax
5:14 & \begin{tabularx}{0.7\textwidth}{X} 後來耶穌在聖殿裡找到他，對他說：「你已經痊癒了，不要再犯罪，免得你的遭遇更壞。」 \end{tabularx} \\ \\ \relax
5:15 & \begin{tabularx}{0.7\textwidth}{X} 那人就去告訴猶太人，使他痊癒的是耶穌。 \end{tabularx} \\ \\ \relax
5:16 & \begin{tabularx}{0.7\textwidth}{X} 所以猶太人迫害耶穌，因為他在安息日做了這些事。 \end{tabularx} \\ \\ \relax
5:17 & \begin{tabularx}{0.7\textwidth}{X} 耶穌就回答他們：「我父做事直到如今，我也做事。」 \end{tabularx} \\ \\ \relax
5:18 & \begin{tabularx}{0.7\textwidth}{X} 為了這緣故，猶太人越發想要殺他，因為他不但犯了安息日，而且稱神為他的父，把自己和神看為同等。 \end{tabularx} \\ \\
[1ex]
\hline
\hline
\end{longtable}
$^{1}$弟兄姊妹 早安.
很高興今天可以和大家一齊去敬拜主.
我們一齊去唱詩 去讚美我們的主.
今天我們的講題是 屬靈的力量.
我們會去看約翰福音的五章一到十八節.
如果你有你的聖經 你都可以去翻開到.
約翰福音五章一到十八節.
不知大家講起屬靈的力量 你會想起甚麼.
我們一聽屬靈的力量.
很多時候會想起一些比較神秘的東西.
例如我們會想起 好像哈利波特那樣.
有些咒術可以去施魔法.
或者可能我們東方人.
我們馬上會想起一些風水的東西.
我們想起一些風水命理.
可能擺個石頭在某個位置裡面.
就可能加你的不知甚麼運.
我們見到西方也有很多這些.
他們會喜歡玩一些塔羅牌.
我想特別在現在這個世代.
有一樣東西叫做「新紀元運動」.
是甚麼來的呢?.
是由印度教裡面出來的.
它說在甚麼東西裡面都有些能量.
你透過攝取那些能量.
你就能夠有些不同的運氣.
它跟你說甚麼水晶.
就不知道跟你帶來甚麼運.
它還要說有甚麼磁場.
這個世界你會發現它都很執著.
在一些屬靈的力量裡面.
你會發覺這個世界的人.
他都很想要一些屬靈的力量.
很想去操縱一些東西.
而當然我們基督徒要很小心.
我們不會去跌入這些陷阱裡面.
我們不會去做這些事情.
但我想我們今天都要有一個小心的點.
就是我們都發現原來屬靈是一樣真實的東西.
原來屬靈的世界是一樣真實的東西.

$^{41}$我們當然不會和他們一樣去濫用一些屬靈的力量.
但是原來我們今天都沒有辦法去否定屬靈的真實.
屬靈的力量是真實的.
我們今天想去問一個問題.
究竟屬靈的力量在哪裡呢?.
我們基督徒既然相信神.
我們相信屬靈世界的存在.
究竟今天我們基督徒.
可以怎樣去經歷這些屬靈的力量呢?.
當我們不是玩風水.
我們不玩這些水晶.
我們不玩塔羅.
我們不信這些儒家冥想.
我們可以怎樣去經歷屬靈的力量呢?.
我們一起有個禱告好嗎?.
我們求上帝帶領我們去進入真理裡面.
主呀!我們在你的面前.
我們求你的恩典.
我們知道唯有靠著你的恩典.
我們才能夠明白你聖經的真理.
所以主呀!我們求你今天帶我們進入約翰福音裡面.
求你讓我們更加明白你.
讓我們知道.
能夠怎樣在你裡面得到力量.
得到那份恩典.
主呀!我們求你去幫助.
開我們的心.
讓我們去領受你自己的話語.
以神禱告.
是奉主耶穌的名祈求.
阿們.
我們今天會先看一段經文.
在我們去做一些反思之前.
我們一起享受一下經文.
我們享受一下故事.
我們一起看看耶穌做了些什麼.
我們今天看的約翰福音第五章.
是耶穌做的第三個神蹟.
如果你心水清的話.
你會知道約翰福音有七個神蹟.

$^{81}$七是完全的意思.
講的是耶穌完全啟示自己.
而今天我們看的是第三個.
如果你有印象.
第一個是水變酒.
是一個叫加拿芬宴的神蹟.
第二個是在約翰福音第四章結尾.
是一個醫治的神蹟.
最奇怪的是.
在約翰福音第五章開頭.
又是一個醫治的故事.
如果你看得勝經多.
你就會在想為什麼聖經那麼重複.
聖經的東西好像來來去去都一樣.
我們今天要看到的就是其實不一樣.
今天這個不單止是一個醫治的故事.
我們就一起進入經文裡.
在第一至五節.
他交代了經文的背景.
他說這些事以後.
即是上一個神蹟之後.
到了猶太人的一個節期.
耶穌就去到耶路撒冷.
在耶路撒冷靠近陽門.
有一個池子.
希伯來話叫畢氏大.
旁邊有五個柱廊.
裡面躺著很多病人.
有失明的,瘸腿的,癱瘓的.
你會看到這個地方是一個很大的地方.
裡面裝著了很多病人.
在這一刻我們還不知道為什麼這麼多病人.
但是第四節他給了一些提示我們.
他說這些病人在等候池裡的水去動.
因為有天使,安士,夏池子搞動那個水.
水動之後先下水的人.
無論有什麼病都可以痊癒.
你會看到第四節是卡住了.
原因是今天很多的聖經學者.
都會說這一段不是原本裡面的.

$^{121}$他會說這一段是後加上去的.
原因可能是他想去解釋一下.
當時有什麼傳說.
在這個池裡的人有什麼信念.
但是無論他怎樣也好.
他解釋了一件事給我們聽.
就是當時的人有一種迷信.
他相信原來這個池是有特別的功效的.
他相信那個池是你浸下去.
如果你是第一個浸下去的話.
你就能夠得到醫治.
接著我們見到經文裡面我們的主角.
在那裡有個人他病了38年.
弟兄姊妹我今年26歲.
如果他病了38年的話.
我要再病多12年.
大家有沒有想過這是多長的時間.
我們有時候說人生可能70歲.
一半的時間都在一個病患裡面.
動不了.
對於這個病人來講.
是一個很絕望的境地.
我們看看耶穌怎樣和他互動.
耶穌見到他躺在那裡.
知道他病了很久.
就問他你要痊癒嗎.
耶穌直接去和他說.
其實我能夠給你一個痊癒.
你要不要痊癒.
我們看看病人怎樣回答.
他說先生在水冷的時候.
沒有人將我放在池裡面.
我要去的時候.
人家就讓我先下去了.
病人好像沒有直接回應耶穌的問題.
耶穌說你要痊癒嗎.
他說先生沒有人將我放下去.
他背後的意思是.
其實他想要痊癒.
但是他覺得醫治的能力在池裡面.

$^{161}$他覺得醫治的能力在水裡面.
所以他只是相信了一句跟耶穌說.
沒有人將我放在水裡面.
所以我到現在還沒有痊癒.
我們看看耶穌說什麼.
耶穌和他說.
起來拿你的玉子走吧.
那個人就立刻痊癒.
拿起自己的玉子走了.
如果你是看一個故事的話.
整件事其實是很快的.
加多一句那個人就已經痊癒了.
那個人等了三十八年.
去等的原來只是耶穌的一句話.
耶穌說一句就立刻痊癒了.
三十八年和立即.
對於那個人來說是幾令人震撼的事.
他等了三十八年的痊癒.
在耶穌的一句說話裡面.
他立刻得到.
我們在這裡看到.
屬靈的力量不是在池裡面.
屬靈的力量是在耶穌的說話裡面.
是在耶穌自己裡面.
如果故事是這樣完就好了.
如果故事就是那個人信耶穌.
成為主的門徒.
一切都簡簡單單平安無事了.
我們很遺憾的事情是沒有這樣完結的.
接著經文說到猶太人出場.
我們每一次看到猶太人.
在日漢福音裡面.
我們都要心感不妙.
因為猶太人每一次都是和耶穌作對.
當你看到猶太人.
你就知道有些不好的事情要發生.
發生什麼事呢?.
在第九節說到.
那一天是安息日.
所以猶太人對被醫好的人說.

$^{201}$今天是安息日.
你不可以拿起你的玉子走.
玉子就是他的箭.
你不可以拿起箭走.
作者在這裡很諷刺.
他把病人叫做被醫好的人.
他故意強調他是被醫好的.
但是猶太人好像完全看不到一樣.
猶太人好像完全是瞎了.
看不到那個人是被醫好的.
他只看到你違反安息日.
你犯了律法.
就只有這麼多.
他只看到這麼多.
那個人怎麼回應呢?.
他說那個使我痊癒的人跟我說.
拿起你的玉子走吧.
那些猶太人問他是誰.
他就不知道了.
因為耶穌已經躲起來了.
這個人也是說事實.
是耶穌叫他走的.
但是你會看到.
這個很明顯是一個推卸責任.
你犯了這件事.
他說不是.
是醫好的那個人叫我做的.
不關我的事.
我們中文有一句話叫「割席」.
馬上跟耶穌割席.
他叫我做的.
不關我的事.
如果事情是這樣結束.
也OK的.
那就是大家相安無事.
耶穌有多特別.
耶穌親自去找他.
耶穌親自去找這個人.
跟他說.
你已經痊癒.

$^{241}$就不要再犯罪了.
免得你的遭遇更壞.
這裡有一樣很特別的事.
在約翰福音的頭四章裡.
從來沒有說過任何一個人.
個別地犯罪.
我們看到約翰福音第一個犯罪的人.
跟他說你不要再犯罪了.
免得你的遭遇更壞.
究竟這個人犯了什麼罪呢.
為什麼說他是一個罪人呢.
我們現在不知道.
看完下一節就知道了.
那個人就去告訴猶太人.
死去痊癒的是耶穌.
等下節目我們留意.
這裡不是一個很中性的形容.
那些猶太人剛剛才說.
那個治他的人犯了安息日.
這個被人治好的人.
他立即跟猶太人去報串.
去打小報告.
說治好他的那個人是耶穌.
這個是一個出賣.
是被治好的人出賣了耶穌.
所以十六節我們立即看到.
那些猶太人就因為這樣迫害耶穌.
因為他在安息日做了這些事.
耶穌的回應都很特別.
耶穌說我父做事直到如今.
我都做事.
這個是一個.
耶穌在引用一個猶太人.
他們很多的討論的一個辯論點.
他們說既然這個世界.
是藉著神的工作去維繫的話.
那不可能上帝是安息日不做事的.
如果不是呢.
世界就在安息日裡面會倒下.
但是除了這個耶穌的引用之外.

$^{281}$你會發現原來耶穌和父神.
是一起去工作的.
我們想一想.
究竟父神是怎樣工作的.
父神是透過說話去工作.
其實耶穌在安息日做了什麼.
去說一句話.
耶穌做的事就是說一句話.
就只有這麼多.
但是耶穌說是父神做事的同時.
我都做事.
原來耶穌和父神是在說話裡面.
一起去做事.
你會發現原來父神設立安息日.
是為了讓土地得到休息去滋養生命.
耶穌所做的都是在安息日裡面.
去醫治去滋養生命.
父神和聖子所做的事是一樣的.
但是跟著十八節我們見到.
一個很諷刺的東西.
當耶穌說父神在安息日都做事.
我都做事的時候.
那些猶太人就更加想要殺他.
因為他不單止犯了安息日.
而且稱神為他的父.
將自己和神看為同等.
為什麼我說這裡很諷刺呢.
是因為那些猶太人說耶穌犯了誡命.
他說耶穌你犯第四條誡命.
你不真的要守安息日.
他說耶穌你犯了第一條誡命.
你不可有別的神.
但你當自己是神.
但是那些猶太人.
他自己都犯了第六條誡命.
第六條誡命說不可殺人.
當他們說耶穌犯誡命的同時.
他自己都犯了所有誡命.
這個是約翰福音給我們看到的諷刺點.
弟兄姊妹我們看完《論經文》之後.

$^{321}$我們集中去看一下.
為什麼病人他經歷了屬靈的力量.
他經歷了醫治之後.
為什麼他沒有應有的信心呢.
為什麼他明明經歷了這麼多.
三十八年的病.
他都能夠痊癒.
為什麼他沒有生出對於上帝應有的敬畏呢.
我想這件事.
我們能夠從對於那個詞的迷信.
我們可以有一點點的痕跡.
我想我們先去想一想.
什麼叫迷信呢.
迷信就是我們覺得屬靈的力量.
背後其實是沒有神的.
屬靈的力量是一些我們可以操控的東西.
例如我們可以透過念咒語.
我們透過擺石頭.
透過帶一些首飾.
迷信就是我們相信屬靈的力量.
背後和神是沒有關係的.
是一些我們可以操控的東西.
那個病人他的問題.
正正就是他覺得那個醫治的力量.
是在那個池裡面.
在那個水裡面.
而不是在上帝裡面.
這樣的話我們會能夠明白.
他怎樣去看耶穌.
對於他來說耶穌就是一個.
能夠操控屬靈力量的人.
不代表他背後是有神的.
如果我們用一個比喻.
我想大家說起一些風水師就很明白.
我們說起蘇文峰 麥玲玲.
如果你今天請他幫你看風水.
我們不要這樣做 是假設.
然後他幫你看完一輪.
佈置完一輪.
你多謝他.

$^{361}$你會不會當他神般拜.
當然不會.
你會不會說我要和你建立一個個人的關係.
我要和你建立一個親密的關係.
你都不會.
這個是交易.
我們覺得他們就是操控屬靈的力量.
就是這麼多.
我們不會當他是神.
要和他建立一個關係.
我想那個病人去看耶穌都是一樣.
他覺得耶穌雖然是有能力.
不代表他背後是有神的.
這個就是他最大的問題.
如果用一個圖我們去看的話.
我們會知道其實父神.
差遣耶穌來到世界上.
耶穌和父神是一起去工作的.
所以耶穌有醫治人的能力.
耶穌都期望人.
因為這樣對於上帝有應有的敬畏.
但對於那個病人來說.
他中間那個連結斷了.
他看不到耶穌背後有父神.
自然都看不到他需要對耶穌有一種敬畏.
他看不到神在耶穌的背後.
如果我們看多一段經文.
你就會更加明白.
究竟為什麼會有這樣的事情發生.
我純粹做一個很快的對比.
在約翰福音9章.
都是一個醫治的故事.
差不多所有事情都一模一樣.
你會見到耶穌主動去醫治一個乞子.
你會見到這個故事裡牽涉一個詞.
都是一個水潭.
你會見到兩個人被耶穌醫治之後.
都不知道耶穌在哪裡.
約翰福音刻意將兩個人.
好像平衡一樣給我們去看.

$^{401}$但不一樣的是什麼呢?.
耶穌講明這個人是沒有罪的.
耶穌講明這個人是一個義人.
他們兩個人的分別在哪裡呢?.
我想透過很快的這段經文來看一看.
在24節他這樣說.
法利賽人就叫那個以前真名的人.
就跟他說.
你要將榮耀歸給神.
我們知道這個人是一個罪人.
在說耶穌.
因為他們覺得耶穌犯了安息日.
所以又覺得他是一個罪人.
我們看看這個人究竟是怎麼回答.
他們互動了一段時間.
我們不詳細去看了.
我們再看看那個瞎了的人的結論是什麼.
從30節那裡這樣說.
那個人回答對他們說.
他開了我的眼睛.
你們竟不知道他從哪裡來.
這真是奇怪.
我們知道神不聽罪人.
唯有敬奉神.
遵行他旨意的神才聽他.
從創世以來.
未曾聽見有人開了生來就失明的人的眼睛.
這個人如果不是從神那裡來.
就什麼都不能做.
我想在這裡我們就明白那個對比.
對於這個核子來說.
去見到耶穌是從父神那裡來的.
去見到耶穌醫治的能力.
不是可以隨便操控的.
而是按著父神的心意一起去走的.
我想如果我們在這段經文裡.
我們要去看到一個屬靈的反思.
就是我們見到真正屬靈的力量.
一定是會按著上帝的心意.
真正屬靈的力量.

$^{441}$不是一個可以我們隨意操控.
好像一些魔法的概念.
真正屬靈的力量.
只能夠在父神的心意裡來.
只能夠在上帝的心意裡來.
當我們明白經文是想傳遞一個這樣的訊息的時候.
我們一起去看看怎樣在我們生命裡.
在我們生活裡去應用.
我想我們今天都面對同一個試探.
我們很容易將屬靈的力量.
放在一些物件裡.
放在一些死物裡.
放在一些方法裡.
而不是將它放在上帝裡.
我們見到這個習慣.
是由古時開始已經在發展.
你會見到猶太人喜歡做什麼呢?.
打仗的時候.
他們怕打不贏.
他們就說上帝的藥櫃是有一個特別的力量.
那個藥櫃有特別的功效.
打仗的時候抬他出去.
就會打勝仗.
最後當然是打輸.
他們完全搞錯了.
上帝的力量從哪裡來.
你見到掃羅也是.
以色列的第一個軍王掃羅.
他在打仗之前.
他說我一定要獻祭.
因為他覺得獻完祭之後.
上帝有特殊的力量會幫他打勝仗.
誰知道薩姆爾遲到.
雖然好像是沒有遲到.
但是他就急著一定要獻祭.
他說死了我不獻祭打不贏.
怎麼辦.
他雖然不是祭司.
他自己匆匆忙忙獻了一個祭.
薩姆爾說你這樣做不按上帝的心意.

$^{481}$你整個國家都沒有了.
你見到後來的猶太人.
他以為屬靈的力量在聖殿裡.
他說有聖殿我們就不會死.
有聖殿就什麼事都沒有.
耶利米就對他們說.
不是的.
你們叛逆上帝的心意.
你們不要以為有聖殿就有用.
上帝要差遣他的審判來到.
我們見到今天基督新教以外.
都有很多不同的儀式.
好像在說屬靈的力量在一些死物裡.
例如你會聽到有些聖水.
有些聖物.
十字架有祝福在裡面.
將屬靈的力量和上帝的心意割開.
說那些東西是有特殊的功效.
我們要去小心這些東西.
我們今天作為基督徒.
我們有沒有這樣的誘惑.
有沒有這樣的試探在裡面.
我想今天我們最大的試探是.
我們祈禱的時候.
我不知道大家私人祈禱會祈什麼禱.
很慚愧地說.
我年輕的時候.
可能十年前.
祈禱最多的禱就是求上帝給我老婆.
十年前不斷去祈禱.
我覺得祈禱多就會快一點.
祈禱多就好像上帝要這樣給我.
可能我上班之後.
我身邊的人最多的祈禱就是.
上帝幫我找一份工作.
上帝給我一份好一點的工作.
是一份高薪厚職的.
上班不要太遠的工作.
可能大家隨著年紀漸長.
我也知道很多人很喜歡求平安.

$^{521}$出入平安.
我希望每一天都無風無雨順順利利.
在我們這些祈禱的背後.
我們是不是覺得我們這樣祈禱.
就好像操弄一些屬靈的力量.
叫上帝祝福給我們.
我們在這樣祈禱的時候.
有沒有考慮過.
究竟上帝的心意是什麼呢?.
我不是說大家不可以這樣祈禱.
我相信我們需要真誠地在上帝面前去祈求.
但是我們要留意的是.
我們不要把這些祈禱當作操縱上帝的工具.
我們在這些祈禱裡面.
我們是去聆聽,去測驗上帝的心意.
我們是跟著上帝的心意去行.
而不是要上帝按著我們的心意去行.
如果我們看著耶穌的主禱文.
我們就會知道在禱告裡面.
最重要的是.
我們在天上的父.
是願人都尊你的名為聖.
是願上帝的國降臨.
是願上帝的旨意行在地上.
如同行在天上.
我們今天如果要祈禱的話.
如果我們真的想祈禱有功效的話.
我們就不能再按著我們的心意而行.
因為上帝的力量不在那裡.
上帝的力量是跟著祂的心意去行事.
剛剛麗仙也說.
我們祈禱的時候.
我們不是好像泥菩薩過江那樣.
求上帝導我們去不知那裡.
我們今天所信的不是一個皇大仙式的信仰.
我們是降伏在主的面前.
跟隨祂的心意.
所以我們今天.
如果我們負面來講.
就不要將屬靈的力量放在錯的地方.

$^{561}$正面來講.
就是真正屬靈的力量.
只可以在執行神的旨意裡面去領受.
弟兄姊妹.
我想我們今天有很多不同的力量在我們中間.
有政治的力量.
文化的力量.
我們甚至都相信其實魔鬼的力量是真實的.
所以有很多人能夠做到很多不同的事.
甚至經文裡面講到.
醫病趕鬼.
都是可能這個世界有人會做到的.
但是我們看看瑪雅提福音怎樣去講.
它說不是每一個稱呼我主呀 主呀的人.
都能夠進天國.
是唯有遵行我父子義的人才能夠進去.
是那天很多人會跟我說.
主呀 主呀.
我們不是奉你的名傳道嗎.
奉你的名趕鬼嗎.
奉你的名行許多義人嗎.
祂要跟他們宣告.
我從來都不認識你們.
你們走開啦.
如果我們今天想要經歷真正的屬靈力量.
我們就必須要去進行上帝的心意.
真正的屬靈力量只可以在上帝的心意裡面找到.
可能在我們悔改裡面.
我們能夠經歷聖靈的大能.
去給我們一個新的生命.
在全福音裡面.
我們經歷到上帝給我們那種大能.
去讓我們去作見證直到地極.
我們可能在教會裡面的事奉.
我們經歷到屬靈的那種恩賜.
我們能夠去事奉上帝.
事奉上帝的子民.
我還很記得我在年幼的時候.
都曾經有追過聖.
我不知道大家有沒有追過聖.

$^{601}$當年已經做了基督徒.
但追聖追得很入迷.
我都知道不對.
我就跟上帝說.
上帝不如我祈多些禱.
我讀多些經.
這樣我都算是一個好的基督徒.
我怎樣追星其實你都不要理我.
我都做盡了我可以做的事.
不要再要求我那麼多.
每當我越祈禱越讀經的時候.
心裡面都很痛.
因為我感覺到上帝離我很遠.
我感覺到上帝不喜悅這件事.
但我都不斷用這種屬靈的操練去麻醉自己.
以為我這樣做多些.
上帝那個屬靈的力量還在我裡面.
直到有一天我真的發現.
和上帝真的很遠.
然後在一個講道裡面.
那個講者呼召.
叫大家要全心全意去敬拜上帝.
不可以再留一絲餘地.
那一刻我真的經歷了屬靈的大能.
當我祈禱那麼久.
讀經那麼久都經歷不到上帝的能力的時候.
在那一刻我發覺.
上帝的大能在推動我.
要我作出一個禱告.
去放下一些不合祂心意的事情.
在那一刻真的很辛苦.
我想差不多十分鐘左右我都沒有舉手.
我都沒有回應那個呼召.
但最後上帝都推動我去做這件事.
在那一刻上帝的大能是很真實的.
我想我們今天如果做一個總結.
就是屬靈的力量是在上帝的心意裡面.
我們不要將屬靈的力量擺錯位.
而是當我們願意去實踐上帝的心意的時候.
上帝的心意就令到我們得力.

$^{641}$上帝願意去給我們力量去進行祂的心意.
我們在經歷這份力量的時候.
我們自然而然能夠感受到上帝.
我想我們今天在這個社會裡面.
我們最大的一個難關就是我們所敬拜的神.
是一個看不見的神.
我們有時候要作見證的時候.
我們不能指著不知道哪裡.
然後跟別人說上帝在那裡.
但今天我們每一個人坐在這裡成為基督徒.
是因為我們個人的生命裡面.
我們都經歷了上帝的大能.
我們每個人都在經歷上帝.
弟兄姊妹我們今天雖然見不到上帝.
但是上帝的聖靈讓我們能夠去經歷祂.
去感受祂.
從而讓我們生出一份信心能夠去敬畏祂.
弟兄姊妹我們今天願不願意去經歷這份屬靈的大能.
當這個世界跟你說很多屬靈的力量.
你可以操縱祂的時候.
我們是否願意降伏在上帝的心意裡面.
去享受那份真正屬靈的力量.
去經歷祂.
去和上帝一起去做工.
從而讓我們在每一天裡面.
我們能夠說今天主都賜我力量.
因為上帝喜悅我去做祂的工作.
弟兄姊妹我們願意嗎?.
我們願意去經歷這份屬靈的力量嗎?.
我們一起去禱告.
親愛的父神.
我們去多謝你.
是因為你很願意在我們生命裡面去做工.
你在我們生命裡面去醫治.
去加力.
你讓我們成為一個新做的人.
主呀我們去求你幫助.
我們求你幫助.
讓我們更加明白你的心意.
因為我們知道.

$^{681}$如果我們今天想去經歷這份屬靈的力量.
我們必須在你的心意裡面去行.
我們必須在你的心意裡面去行走.
主呀我們今天很需要去經歷你.
特別在這個世俗的年代.
在這個好像屬靈力量不被彈劾的年代.
我們很需要真實去經歷你的大能.
從而活出一個不同的見證.
我們求上帝你去幫助我們.
在生命裡面活出你那份力量.
從而讓其他人去看到.
讓其他人去羨慕.
讓其他人願意一起去歸入主耶穌基督裡面.
我們求上帝你這樣去幫助我們.
以上禱告是奉主耶穌的名祈求.
阿們.
\newpage



\section{哥林多後書 3:17}
\label{sec:031824}
\textbf{聖靈使你自由 (華)}~~~~葛國光牧師
\newline
\newline
連結: \href{http://www.church.com.hk/acms/content.asp?site=cdc\&op=show\&type=product\&code=031824\&layout=sermon}{\texttt{031824}} ~~~~ 語音日期: 2026-03-15
\newline
\newline
\hyperref[sec:031823]{\small{< < < PREV SERMON < < <}}
~
\hyperref[sec:index_chronic]{\small{[返順時目]}}
~
\hyperref[sec:031825]{\small{> > > NEXT SERMON > > >}}
\newline
\newline
哥林多後書 3:17
\newline
\begin{longtable}{cl}
\hline
\hline
章節 & 經文 (和合本修訂版)\\
\hline
3:17 & \begin{tabularx}{0.7\textwidth}{X} 主就是那靈；主的靈在哪裡，哪裡就有自由。 \end{tabularx} \\ \\
[1ex]
\hline
\hline
\end{longtable}
$^{1}$我們這一段時間連續地講到.
基督釋放我們.
要教我們得自由.
我們看見最近.
這個中東的局勢越演越烈.
可能你的新聞,可能新聞都被它佔據了.
可能你的焦點都被它佔據了.
你看見這些一下子又轟炸了.
無人機啦,又是這些.
其實劉子梅你看見.
從這些裡面我們真是看見.
主來的日子近了.
是不是,國要攻打民,民要攻打民.
到處都有戰爭,到處都有這些各式各樣的問題.
親愛的劉子梅,今天這個世界.
有什麼能夠使我們有盼望的.
什麼能夠使我們有盼望的呢.
你會發現一次打仗.
曾經是大家覺得.
應該是最安全的,據說.
阿聯酋.
杜拜,他們認為這是全世界最安全的國家.
現在也不安全.
劉子梅,有什麼地方是安全的呢.
有哪個地方,有哪一個人.
有哪一個權勢,有哪一個帝國.
可以帶給我們安全的保證呢.
沒有.
所以.
哥林多厚書第三章第七節,剛纔我們讀的經文說.
主就是那靈,主的靈在哪裡.
哪裡就得以自由.
森林要來使我們得以自由.
森林要來幫助我們.
這個自由不是外在的一個政權的壓制得自由.
這個自由是雖然外面有很多的壓制.
但是卻能夠在靈裡面得到自由.
有很多人可能外面很釋放.
外面很自由,但是靈裡面.
心靈裡面,思想裡面.

$^{41}$卻被壓制著.
這就是主耶穌當他出來傳道的第一篇信息裡面.
在路加福音第四章.
十四節,十九節,剛纔我們所讀的這段經文.
耶穌出來傳道第一篇的信息重要不重要.
重要.
耶穌離開世界所講的最後的一段話重要不重要.
重要.
耶穌出來傳道的第一篇信息.
他打開聖經裡面.
當時候的這些的猶太人的拉比.
就是打開聖經,然後讀出來.
然後耶穌就開始講這一段.
耶穌宣講的這一段是引用在.
以賽亞書第六十一章第一節到第三節.
耶穌提到什麼呢?.
他說.
主的靈在我身上.
因為他用怎麼樣?.
用高高我.
主耶穌用這句話來引用以賽亞書第六十一章.
以賽亞書六十一章就在預言耶穌.
說什麼?.
耶和華的靈在我身上.
因為他用高高我.
所以這邊講到他的靈.
他的靈是什麼呢?就是聖靈.
聖靈在主的身上.
他用高高我.
高是什麼意思?.
在舊約裡面.
先知君王祭司.
都需要經過高麗的儀式.
用什麼油呢?用橄欖油.
所以很多人到耶路撒冷去.
到聖地去旅遊都要帶個什麼油回來?.
橄欖油.
其實你回去自己家拿你們家的橄欖油.
也一樣有效果.
很多人要花一塊錢.

$^{81}$現在可能不止一塊美金.
買那麼小小的一罐橄欖油.
而且說是耶路撒冷來的.
聽著沒?.
這個高.
就在預表著油.
就是預表著聖靈.
聖靈用高高我.
當然這個高在這裡講的.
如果你看耶穌講這個.
他提到的是.
在以賽爾書提到就是說.
他被高利他擔任了這個職份.
但是耶穌在引用這段聖經的時候.
他前面講的說.
他用高高我.
是為了什麼?.
為了什麼?.
叫我傳福音給貧窮的人聽.
拆遣我報告.
被擄得得釋放.
瞎眼的得看見.
叫那受壓制的怎麼樣得自由.
報告神願那人的洗臉.
在這邊提到這些事情.
我把它歸納一下.
首先他第一個提到是什麼?.
聖靈做的一個工作.
第一個瞎眼的得看見.
瞎眼的得看見.
在耶穌引用了這段聖經裡面.
他沒有講.
但是在以賽亞書61章有講.
瞎眼的得看見代表是身體得抑制.
聖靈要來做一件事情.
使那個身體有疾病的得到抑制.
今天美國大概有76.4\%的人.
有一種以上的慢性疾病.
佔了三分之二的人都有這個疾病.
比如說.

$^{121}$第一個叫做肥胖.
大約42.4\%的成年人.
那麼第二個是高血壓.
大概有32\%的人.
還有糖尿病.
12.5\%的人.
還有關節炎.
20\%的成年人.
那麼有的人是同時有兩種以上的慢性疾病.
大概有40\%到51\%的人.
所以你看見.
百業蕭條.
唯有一個業不蕭條的哪一業?.
醫療業.
找工作到處都找不到.
但是醫護人員非常搶手.
是吧.
所以現在大家都要去讀護學.
不僅是女孩子去讀.
男孩子也去讀.
因為做護士的更需要男士.
你要抱200磅重量的病人.
就需要男護士了.
我女兒做護士我知道的.
所以你看見今天這個世界裡面.
到處都是這些生病的醫院裡面都是客滿的.
每一次我到醫院裡面去探訪.
我發現總是都是客滿的.
現在有許多的人有身體.
包括我們自己.
我們生病.
很多的病.
醫生都沒有辦法醫治.
今天越來越多的奇怪的疾病.
最近我們教會里面的肇事母生病.
李李生病.
很多人生病.
然後這個50斤的.
不到50就有50斤.
40斤.

$^{161}$這邊痛那邊痛.
腰椎間.
椎間盤突出.
各式各樣的.
現在我們人生當中真的很苦.
你沒有生病你不知道.
我生過.
真苦啊.
我椎間盤突出痛的時候.
真是痛不欲生.
怎麼做都不是.
開車也不能開.
現在就準備.
何等需要.
我們今天何等需要.
有這樣子的一個醫治的大人在我們中間.
是嗎.
不僅僅是.
身體的醫治.
今天還有許多.
傷心的人.
心靈需要得到醫治.
根據統計.
美國.
全世界.
大概有排行榜前10名的.
大概都是歐美國家.
就是離婚率最高的.
大概都是歐美國家.
那麼我們美國佔第幾位呢.
算了一下.
美國第10名.
美國佔第10名.
那麼中國大陸第幾名呢.
中國大陸緊追美國之後.
所以中國大陸是11名.
從整個數據來看.
亞洲第一名.
No.1是中國.
No.2 臺灣.

$^{201}$這邊還沒念進去.
根據統計.
美國大概有將近40\%的人.
Eventually 最後婚姻會以離婚結束.
第一次結婚的.
會離婚的大概是這個數據.
第二次結婚的.
這個比例升高到63\%.
那麼第三次結婚的.
也就是三次的結婚的人.
這個比例達到了73\%.
真是人生很苦啊.
這些大概都是年輕人少不更事啊.
他們這些人把這個婚姻當兒戲啊.
但是我告訴你.
這些年銀髮族的離婚開始越來越普遍.
什麼叫銀髮族的離婚呢.
就是165歲以上的人群的離婚率.
從90年代以來幾乎翻了三倍.
當我們講這些數據的時候.
你可能沒有什麼特別感覺.
但是請問在婚姻當中還不要談到離婚.
請問失戀痛不痛苦.
失過戀嗎.
失戀懂嗎.
我的中文不標準是不是.
失戀.
聽懂嗎.
懂了.
失戀痛不痛苦.
有沒有失戀過的.
痛不痛苦.
那你想想看.
失婚痛不痛苦.
是吧.
你交往了幾年.
你分手了.
很痛苦.
有的人痛不欲生啊.
但是如果結婚了.

$^{241}$分開了.
那更痛苦.
那還不僅僅是夫妻痛苦.
而且是什麼.
兒女也痛苦.
而且很多時候那是一生的苦.
你懂不懂.
今天有多少的人心是碎的.
在人前.
可能在聚會的時候.
你表現出來是強顏歡笑.
但是心裡面卻是非常的痛.
這個樣子的傷心.
這個世界的金錢沒有辦法滿足你.
這個世界沒有任何的東西可以填滿.
可以抑制.
所以許多的人是傷心的人.
那麼耶穌說.
祂要來使傷心的人得到醫治.
醫好傷心的人.
耶穌說祂要使我們眼睛得看見.
祂要使我們身體得醫治.
上帝要解決我們一切的身心靈的問題.
而且祂要使我們使貝魯德得釋放.
這個貝魯德.
從我們今天.
應用在我們今天來講.
有許多的人被一些的成癮.
被一些藥物成癮所捆綁.
根據統計.
大概有16.8\%的美國人.
12歲以上.
大約就是將近4800萬人.
在過去一年當中.
有藥物濫用的這個問題.
也就是addiction.
那麼這一個藥物的濫用.
其中酒精的.
酒精就是有酗酒的問題.
有9.7\%.

$^{281}$那麼大麻的.
大概有22.3\%.
現在.
如果你有機會去接觸到年輕人.
校園裡面.
在職場裡面.
很多人覺得吸大麻沒什麼關係.
所以在他的分享當中.
在跟朋友的分享當中.
他不會避諱.
哎呀我偶爾也抽抽大麻.
有沒有.
特別在我們家族.
像就是今天這個世界.
你可以看見.
有許多的人被這些的癮所捆綁.
被一些酒癮所捆綁.
被一些煙癮所捆綁.
被這些的鴉片.
被這些的大麻所捆綁.
還有一個更厲害的.
被性的癮所捆綁.
這個叫internet pornography.
這個是更不容易的.
而且是更不為人所知的.
根據統計美國.
大概有87\%的人.
他曾經在過去三個月.
看過這個internet pornography.
也就是網路色情.
當然他不一定是說.
到達一個地步.
就是說他已經被捆綁了.
但是這個比例相當的高.
代表一件事情.
這是今天在美國一個很普遍的問題.
所以你知道嗎.
現在最近發生的那個什麼.
那個叫什麼普啊.
那個.

$^{321}$你們知道那個.
那個淫魔對不對.
其實你知道.
這個背後面裡面.
網路的色情推波助瀾.
帶來的影響.
而且不是一天兩天.
非常非常的厲害.
把整個的心理扭曲了.
我可以告訴你.
這個戒毒啊.
我還看見.
基督教裡面有很多戒毒成功的.
對不對.
比如說叫劉明和牧師.
在臺灣的所辦的這個晨熙會.
但是戒攝影啊.
你怎麼辦.
你怎麼去.
你能不能夠辦一個.
找一個營地.
在這個營地裡面就說好.
你如果有這個有這個.
有這個網路色情引得人.
請到這邊來.
我們來幫助你.
有嗎.
會不會有人來.
不會有人來的.
所以他更可怕.
你懂我意思嗎.
他躲在房間裡面.
誰也不知道.
甚至有人結了婚之後.
知道.
太太知道.
已經沒有辦法了.
親愛的同志們.
今天這個世界有太多.
我們看見沒有辦法解決的問題.

$^{361}$但是親愛的同志們.
耶穌說他要來.
使貝魯特得釋放.
你相信嗎.
雖然這個看起來.
目前人好像沒有辦法解決的問題.
但是你相信耶穌基督可以釋放那些.
被這種色情所捆綁的人嗎.
神一定可以.
魔鬼再怎麼囂張.
他想要用這一些厲害的方法.
但是在我們裡面的.
比在那個世界的更大.
在我們裡面是誰.
就是這位聖靈.
使貝魯特得釋放.
還有呢.
使受壓制的得自由.
被壓制.
這邊我們要講到.
有嚴重一點的是被鬼附.
但是有另外一個跟類似鬼附.
比較接近的.
你知道叫什麼嗎.
叫做精神病.
以前我們在初初出來服事的時候.
還特別請過.
臺灣的龍種的精神科大夫.
魯魯賓大夫.
幫助我們傳道同工.
去知道精神病.
是有些什麼現象.
因為趕鬼我們有經驗.
我們知道趕鬼的情況是怎麼樣.
但是這些專業的精神病.
我們警察來告訴我們說.
精神病會什麼這樣子.
那麼你會發現.
有很多時候精神病跟鬼附怎麼樣.
這兩個很類似.

$^{401}$他們都會幻聽幻覺.
簡單講一個.
但是你會發現.
這個精神病的.
你千萬不能夠把它當鬼附來趕.
瞭解我意思嗎.
那精神病的人.
需要有精神病的方法.
那麼鬼附的呢.
就要趕鬼了.
所以最麻煩的是什麼.
同時有精神病.
又有鬼附.
我第一次碰到趕鬼的.
就是碰到這種case.
好 今天不講詳細了.
但是現在就準備.
耶穌說他來要使受壓制的得自由.
美國今天有很多的人.
有抑鬱的問題.
這個抑鬱的問題.
是一個很嚴重的問題.
你看看.
一二三四五.
每五個當中就有一個.
有這所謂的mental illness.
五分之一.
那二十分之一呢.
這是指的說.
有嚴重的有嚴重的憂鬱症.
現在我們看見.
特別年輕的族群當中.
越來越多的這些有抑鬱的問題.
很多父母親.
以前成長過程中.
從來不知道這一回事.
怎麼碰到自己的孩子.
會有抑鬱的問題呢.
我們這個爸爸媽媽.
也沒有說對虐待你們.

$^{441}$哎呀幫助你們啊.
真是照顧你們的.
無為不致.
怎麼還會抑鬱呢.
不管你喜歡不喜歡.
你會發現這是今天這個社會的現象.
現在這一代的年輕人.
面對的這些各式各樣來的.
你說壓力也好.
或者挑戰也好.
勢態也好.
引誘也好.
造成了今天這一代的年輕人.
很多有抑鬱的問題.
親愛的鄉親們.
怎麼辦.
可以解決嗎.
上帝可以解決.
上帝可以醫治.
耶穌說他來.
他說司令.
搞我要使那些被壓制的德自由.
所以親愛的鄉親們.
今天我要跟大家介紹.
這一位醫治的良方.
這一位能夠解決我們人生當中.
一切這些問題的那一位.
聖靈.
可是講到聖靈呢.
在我們一般的所謂福音派教會.
一談到聖靈.
有的就是把它當作是毒蛇猛獸.
把它當作是.
就把聖靈跟靈恩劃作等號.
聖經都沒有好好讀.
就隨便地貼這個標籤.
很可悲.
使得今天很多的傳統的福音派教會.
你可以看見.
聚會死氣沉沉.

$^{481}$然後呢.
看不見神大人的作為.
生病的看不見醫治的作為.
看不見這些.
剛纔我所說這一切的事情的改變.
所以人來到教會里面要尋找這位神.
因為神如果不能夠做這些事情.
千里之內他去找誰.
他去找誰.
他去找外面的.
他去找那些算命的.
他去找那些行邪術的.
他去找那些可以幫助他的.
你懂我意思嗎.
然後我們在這邊自己標板.
我們自己在那邊自己關起門來自我安慰.
我們是正統的教會.
什麼叫正統的教會.
沒有聖靈的工作的教會叫正統教會嗎.
所以我們今天需要聖靈.
跟你旁邊說你需要聖靈.
聖靈是與基督與上帝.
我們所信的三位一體的神.
所以他說主就是那靈的意思.
主耶穌親自endorse.
主耶穌親自的說聖靈就是我.
我們合一的.
我們兄弟不分家的.
你歡迎我.
你就要歡迎聖靈.
今天福音派有一種說法.
聖靈來是為榮耀基督的.
所以我們不用講聖靈.
我們就講基督.
這件事叫斷章取義.
兄弟姐妹那你這樣講的話.
那我們只講基督.
講不講上帝呢.
耶穌說什麼.
耶穌說他要來榮耀父.

$^{521}$是不是.
他不求自己的意思.
他只求誰的意思.
天父的意思.
那你要算的話.
那應該是老爸最大了.
沒有所以你要整個經文來看.
聖父聖子聖靈.
他們同榮同等同尊.
是不是.
這個叫做三位一體.
這是我們每一個教會的共同的信仰.
所以今天我們需要聖靈的高某.
耶穌不單單是耶穌被聖靈高某.
可以做這個事情.
耶穌講這一段話.
也是在告訴我們歷史歷代.
跟隨他的兒女們.
當我們親眼聖靈的時候.
聖靈在我們身上的時候.
我們一樣也可以如此的去行.
主耶穌在地上的時候.
他這是使徒行載十章三十八節來形容耶穌的.
我們一起來讀好不好.
來預備請.
神怎樣以聖靈和能力高拿下了耶穌.
這都是你們知道的.
他周流四方行善事.
醫好凡被魔鬼壓制的人.
因為神與他同在.
那時候耶穌是在地上.
所以說神與他怎麼樣.
同在.
現在主耶穌在哪裡.
現在當我們說.
神與我們同在.
主與我們同在.
是什麼.
主耶穌說.
他說我去了.

$^{561}$但是呢.
我要差遣另外一位保會師.
他要與你們永遠的同在.
所以現在誰跟我們同在.
聖靈.
你說主同在你說上帝同在.
他們沒有分的.
我再說他們不分的.
是我們在分.
我們分說.
靈魂派聖靈與你們同在.
那是.
那不是真正的.
現在都準備不是的.
聖靈要與我們同在.
主要與我們同在.
神要與我們同在.
聖靈的能力要充滿在我們身上.
聖靈的能力就是神的能力.
但是今天耶穌他離生天了.
那麼聖靈是神應許降在我們中間.
要來幫助我們.
很可惜啊.
每次我們到這時候.
我們這些父母親都要去接孩子.
我等你回來我再帶你們禱告.
趕快快去接快回來.
不僅僅是耶穌.
保羅在羅馬書十五章十八節.
保羅說除了基督教我做的那些事.
我什麼都不敢提.
只提他借我言語作為用神機啟示的能力.
並聖靈的能力使外邊人信服.
保羅在這邊講到言語.
他沒有說不要用原語.
請大家記得啊.
我們今天只有這麼多時間.
我focus在聖靈的能力.
可是這不代表我們不要用言語去做接的.
我們不要有好的行為.

$^{601}$不是保羅說言語行為.
然後說了什麼呢.
神機啟示的能力.
聖靈的能力.
看到沒有.
所以保羅說他傳福音使外邊人信服靠什麼.
靠這一張嘴皮子.
靠他的這樣子口才去辯論.
不單單如此.
他說也靠他的見證他的行為他的生活.
同時他說神機啟示的能力.
還有呢.
聖靈的能力.
特別你知道今天我們在中東啊.
看見這些伊斯蘭教的國家.
我告訴你很多像今天我們禱告的蘇丹阿拉伯人.
這些人我跟你講.
你跟他傳福音傳什麼福音呢.
你還沒有開口就把你殺了.
你懂意思嗎.
你跟他講這些東西沒有用的.
你跟他講這些真理他聽不進去的.
他根本不給你機會去聽啊.
但是這些人當他們遇見神的大人的時候.
所以我們聽見很多伊朗人.
特別是伊朗人.
他們信耶穌怎麼信的.
他們就在這個禱告的中間呢.
就看見天使向他們顯現.
看見耶穌向他們顯現.
你有沒有看過天使向你顯現啊.
沒有.
你有沒有看過耶穌向你顯現啊.
沒有.
為什麼.
因為那是靈派的東西.
我們福音派不講這些.
是不是.
但親愛的兄弟妹.
那為什麼伊朗的人會看見.

$^{641}$在幕後的日子.
神仍然會做這個事情.
是用奇妙的方式.
用神蹟啟示.
用這異象異夢向這些伊朗人顯現.
這些伊朗人看了之後.
他們就相信了.
因為他們是相信這些東西的.
否則他們在這兒祭祀禱告幹嘛.
他們希望看見.
這個所謂的神向他們顯現.
所以當他們看見.
這個顯現的近來是.
他們所相信的那位先知耶穌的時候.
他們突然明白.
耶穌纔是神.
這必須講一篇大道理.
把整個聖經都講完了.
可能都沒有辦法讓他們明白.
親愛的同志們.
所以這就是我們今天所要傳講的.
你知道我有一個學生.
在國內.
他剛纔講的王尹.
他年紀輕輕的.
據他講十七歲以前.
他的青少年時期是被王尹所捆綁的.
每天都要去打.
那個年代還沒有自己家.
沒有上網的條件.
去哪裡呢.
去網吧.
他去網吧.
他說曾經最高記錄.
三天三夜.
我聽過三天三夜打麻將的.
他是三天三夜在網吧裡面上網.
不睡覺.
已經到這種地步了.
爸爸媽媽真是沒有辦法的.

$^{681}$一個女孩長得漂漂亮亮的.
但是就是被這個王尹所捆綁.
親愛的同志們.
有辦法嗎.
很奇妙.
結果他經歷了耶穌基督.
他信了耶穌基督.
聖靈釋放了他.
他徹底的從王尹當中被釋放出來.
親愛的同志們.
這一類的事情我沒有時間跟你講.
很多我自己所看到所經歷的.
有一位師母.
是早年的時候.
神護照她的先生.
給她很清楚的意向.
要到中國去宣教.
但是她心裡面一直不願意.
在那邊掙扎掙扎.
在那邊一直等待.
就是她就扯後腿.
到最後.
結果有一天發現.
她得了乳癌.
哇 這一下子把她.
真的是不知道該怎麼辦.
在神的面前禱告求主醫治.
心差了 光糟啊.
當她看見.
神要她悔改.
神要她學習順服先生.
跟著先生一起去服事.
就在她自己在她的傳奇裡面形容.
她就在一次的聖明的充滿當中.
神完全的醫治了她的乳癌.
沒有靠任何的藥物.
沒有靠任何的這些的醫生.
就在這一次當中.
神靈醫治了她.
親愛的劉子妹.

$^{721}$上帝是可以今天也做這樣的事情的.
因為祂要使貝魯德得釋放.
使瞎眼的德看見.
使受壓制的德自由.
劉子妹.
這位神靈在我們今天.
照樣也要做這樣的工作.
你知道中國80年代的時候.
經歷到很大的復興.
那個復興跟聖靈的工作很有關係.
可是為什麼他們會經歷這個聖靈的大能呢?.
第一個.
他們渴慕.
什麼叫渴慕啊?劉子妹.
我曾經跟你們形容過.
我自己當年渴慕的時候.
牧師在前面講道說.
渴慕的人到前面來.
我是第一個衝到前面來.
跪在前面.
這種光景.
現在在美國.
已經很難看到了.
但是你知道嗎.
我過去在中國服侍的時候.
到哪一個地方去.
黨說道告到前面來.
跪在主的面前.
全部都到前面來了.
渴啊渴慕啊.
路加福音11章13節我們一起來唸好吧.
你們雖然不好.
尚且將好東西給兒女.
何況天父.
豈不更將聖靈給求他的人嗎?.
聖靈要給求他的人.
有人說聖靈充滿.
不需要你去求不要追求.
有沒有聽過這種說法?.
這些聖經怎麼解釋?.

$^{761}$而且原文裡面這個求.
是一個持續不斷的狀態.
表明一件事情.
就是你要一直地求.
你不是五分鐘熱度地求.
你要一直禱告一直禱告一直禱告.
直到聖靈充滿你.
這才叫做渴求.
這個圖片是.
我有一本書別人送我的.
是這個荷蘭方城教會的.
這個負責人叫做張榮亮牧師.
如果你稍微瞭解中國教廷教會你就知道.
他是中國教會五大家的其中一個領袖.
在這本書裡面講到方城教會.
首先他講了自己.
他說1980年臘月底的那一天.
他們舉行了一個七天的一個培訓會.
就在這個培訓會的當中.
他說第五天的時候.
森林如瀑布般大大降臨澆灌.
突然間好像一陣大風吹過.
整個聚會的屋子都被震動起來.
讓屋子裡面充滿了新鄉之氣.
那麼那是他們第一次的這些同工們經歷了.
到了1988年的時候.
那個也是差不多我同時間去到浙江的地方.
他說五旬節的火帶來了河南的復興運動.
教會普遍被震.
五旬節他指的五旬節就是指的說.
1988年河南的教會經歷了復興的工作.
森林大大的降臨好像五旬節一樣.
他說之前10\%是基督徒.
之後1989年到1990年增長到20\%.
過去河南你們國內來的知道.
過去河南人所謂的reputation不好.
他們曾經有舉一個例子.
他說那時候農村不敢餵牛.
一餵牛一定遭盜.
盜賊還留言.

$^{801}$你養牛我發財.
你不養我不來.
你再養我我還來.
我家父是你的合作安排.
望你明年多養牛.
我家急著蓋洋樓.
企業委是謝謝合作明年見.
那麼河南人在當時候的中國被人家這樣形容.
十億人民九億片.
一億河南人做教練.
總部設在駐馬店.
河南的駐馬店.
很多工廠招工.
我廠急招工.
薪高也輕.
願者皆可來.
河南人除外.
但是你知道嗎.
當河南的教經歷的大復興.
復興帶來一個結果.
許多人認罪悔改.
不敢再去偷竊了.
按照聖經的話語面去學習主的道.
按照聖經話去學去做.
河南人給改變了.
誰能夠做這件事情.
當聖靈來工作的時候.
改變這一切.
親愛的九子妹.
今天我們相信當聖靈工作的時候.
比那個原子彈.
比那個核子彈.
比那個美國派的那個轟炸機.
所傳遞的那個炸彈.
更要有能力.
那個能力只能夠摧毀外在的建設.
是不是.
但是聖靈的能力.
是能夠改變人心的能力.
我們今天需要的是這個能力.

$^{841}$聖靈的能力能夠改變人的心.
聖靈的能力能夠改變夫妻關係.
聖靈的能力能夠帶來家庭的和諧.
聖靈的能力能夠醫治我們的傷心.
聖靈的能力能夠醫治我們的depression.
它能夠醫治我們身體一切的疾病.
聖靈的能力可以使我們被擄得得釋放.
受壓制的得自由.
\newpage



\section{撒母耳記上 22:1-2}
\label{sec:031825}
\textbf{我們需要避難所}~~~~李思聰牧師
\newline
\newline
連結: \href{http://www.church.com.hk/acms/content.asp?site=cdc\&op=show\&type=product\&code=031825\&layout=sermon}{\texttt{031825}} ~~~~ 語音日期: 2026-03-15
\newline
\newline
\hyperref[sec:031824]{\small{< < < PREV SERMON < < <}}
~
\hyperref[sec:index_chronic]{\small{[返順時目]}}
~
\hyperref[sec:031826]{\small{> > > NEXT SERMON > > >}}
\newline
\newline
撒母耳記上 22:1-2
\newline
\begin{longtable}{cl}
\hline
\hline
章節 & 經文 (和合本修訂版)\\
\hline
22:1 & \begin{tabularx}{0.7\textwidth}{X} 大衛離開那裡，逃到亞杜蘭洞。他的兄弟和他父親全家聽見了，都下到他那裡去。 \end{tabularx} \\ \\ \relax
22:2 & \begin{tabularx}{0.7\textwidth}{X} 凡生活窘迫的、欠債的、心裡苦惱的都聚集到大衛那裡，他就作他們的領袖，跟隨他的約有四百人。 \end{tabularx} \\ \\
[1ex]
\hline
\hline
\end{longtable}
$^{1}$我們今天來看三文魚記上的22章.
今天跟大家說這兩節的經文.
如果大家記得在上一次.
我想都不會記得 因為隔了兩個禮拜.
一個禮拜可能都不記得 別說兩個禮拜.
兩個禮拜之前是說什麼 有沒有記得.
(笑聲).
啊 厲害 E6厲害.
究竟什麼是我們真正的要求.
我們知道這段期間大衛發覺.
很多事情都不能夠成為他真正的要求.
無論他的妻子 無論他的好朋友 或者良師.
這些全部他的權勢忽然間就沒有了.
他越來越發覺這些不是他真正的要求.
究竟什麼是他真正的要求呢.
在這段時間裡 他很深刻的寫了很多的詩篇.
表達他自己對神的愛慕和敬拜.
他越來越發覺 真是神是我們真正的要求.
所以今天我們來看這22章這兩節的經文.
就可以體會到有大衛的心境.
因為從第22章開始的後幾章.
其實大衛是漂流在曠野的時候.
而很多時候他在睡覺和生活的時候.
躲在哪裡呢?就是躲在山洞裡面.
所以在這幅相大家都可以留意到.
就是在山洞裡面 在洞穴裡面的生活.
以至他能夠逃避蘇羅對他的追殺.
所以在22章開始的時候第一節就這樣說.
我們一起讀這節經文好嗎?請讀.
大衛就離開那裡 逃到阿道蘭洞.
他的弟兄和他父親的全家聽見了.
就都下到他那裡.
所以我們看見大衛逃到阿道蘭洞.
發覺有很多的人跟著他.
聽見的時候都到他那裡去.
除了父家之外 待會讀到第二節.
我們更加見到有差不多有400名的.
這些跟從者都去到他那裡.
阿道蘭洞可算是大衛的避難所.
他就躲在這個山洞裡面.

$^{41}$所以我記得在疫情裡面有人呼籲.
他說我們要躲在洞裡面.
回想起來的時候真的是封鎖了自己.
躲在一起.
但是疫情開始過去的時候就有人呼籲.
我們要出洞了 從洞裡面出來.
所以不要永遠躲在洞裡面.
事實上這段時間是逃難的時間.
所以他真的要在洞裡面.
在詩篇裡面時常提到這個避難所.
我們很熟悉的.
亦都提到這個洞穴.
或者在阿道蘭洞.
特別是指到當時所寫的.
這段時期有好幾篇詩篇.
在一些性侵學者研究.
就發覺起碼有三篇.
雖然他們時間不是按著詩篇的排序先後.
不過應該首先是寫詩篇142篇.
然後詩篇57篇.
然後還有一篇應該都是那段時間寫的.
就是肯定的是34篇.
是不是有人想進來招待.
所以我們先看看這裡的詩篇142篇.
這篇的詩題就叫做《求主眷顧》.
他寫著大衛在哪裡寫.
大衛在洞裡作的憤悔詩.
是不是很難看.
對不起.
因為我們有時在電腦裡面看得很清楚出來的時候.
我忘記了橙色比較難看.
考考你的眼力.
大衛在洞裡作的憤悔詩.
乃是祈禱.
這節應該容易看一點.
我們一起讀好嗎.
請讀下第五節.
耶和華.
我曾向你哀求.
我說你是我的避難所.

$^{81}$在沒人之地.
你是我的福分.
所以詩篇第142篇.
這個詩題就很清楚了.
說到他在洞裡面寫的憤悔詩.
是一篇祈禱.
然後第五節就很清楚.
耶和華說.
我曾向你哀求.
我說你是我的避難所.
在沒人之地.
你是我的福分.
所以他體會到.
其他所有的人可能都離開.
有些死了.
但是唯有在神的裡面.
我們能夠找到真正的福分.
第二篇就是詩篇第57篇.
留意這篇詩題也是橙色的.
分別了.
這篇叫《求主詩援》.
我們看看這篇詩題.
大衛逃避蘇羅.
藏在洞裡那時.
他作者金詩.
交給令長.
調用幽妖毀壞.
所以這篇詩又在哪裡寫.
在洞裡面逃避蘇羅的時候.
大衛有幾篇叫《金詩》.
即是詩中之詩.
其中這篇就是《金詩》.
交給令長即是可以唱的.
是一篇詩來的.
意思就是說.
《求主詩援》.
幽妖毀壞.
第一節是這樣說.
我們一起讀第一節.
請讀.

$^{121}$上帝啊.
求你憐恤我憐憫我.
因為我的心投靠你.
我要投靠在你翅膀的任下.
等到災害過去.
很明顯這裡是指他的祈禱.
他的哀求.
他的詩歌.
他說我的心投靠你.
我要投靠在你翅膀的任下.
所以這裡是指明.
大衛那種投靠神的心.
他知道所有其他東西.
都不是他真正的倚靠.
他來去投靠這位.
去任避他.
保護他的上帝.
最後一篇是34篇.
詩題是在這段階段裡寫的.
這篇是頌讚.
頌讚上帝的美善.
所以學者認為這篇是後過.
剛才那兩篇的.
所以這就是大衛在.
亞比米勒面前裝瘋.
就是指到他後來離開了這個洞.
剛剛離開的時候.
他就去到亞比米勒的王面前.
但當時問題還未解決.
但他已經獻上他的讚美.
因為神一直保守到他.
所以看第八節.
就這樣說.
我們一起讀好嗎.
請讀.
你們要嘗嘗主恩的滋味.
便知道他是美善.
投靠他的人有福了.
所以就指到他經歷過投靠神的.
是什麼話.

$^{161}$是有福的.
所以他投靠神是有福的.
所以他就呼籲.
你們也要嘗嘗主恩的滋味.
便知道上帝的美善.
know that he is good.
知道上帝的美善.
投靠他的人是有福的.
這個是作見證.
即是他投靠神之後出來的見證.
問題還未過去.
心境改變了.
成為一個讚美上帝的心境.
而且呼籲人都要來投靠這位上帝.
因為祂實在是我們人生的避難所.
這個在洞穴裡面經歷.
神是他的避難所.
是一個很真實的感覺.
因為在洞裡面是一個避難的地方.
而他不只是覺得暫時在環境上.
避開這些災難或者攻擊.
他是感受到他的心靈.
和他的生命在人生當中.
上帝就是他的任避.
是他的避難所.
所以他勸導人來投靠神.
投靠神.
好像他一樣經歷一樣.
所以我們回到撒梅記二十二章.
第一和第二節.
我們今次讀第二節.
我們從第一節第二節.
讀完這兩節.
請讀.
大衛就離開那裡.
逃到阿道蘭洞.
他的弟兄和他父親的全家聽見了.
就都下到他那裡.
凡受困逼的.
欠債的.

$^{201}$心裡苦惱的.
都聚集到大衛那裡.
大衛就作他們的頭目.
跟隨他的約有四百人.
所以我們知道不單止他的父家弟兄.
全家.
而且有很多很多這些人.
這些人全部都是什麼人.
全部都是受困逼的.
欠債的 心裡苦惱的.
這個只是代表著我們人世間.
各樣的.
我們所面對的問題.
今天這個世界就是充滿著各樣這些問題.
所以其實我們.
每一個人都好像這些人一樣.
受困逼的 欠債的 心裡苦惱的.
但是他們就好像大衛一樣聚集到他那裡.
大衛就作他們的頭目.
當時有多少人.
有很多人 幾百人那麼多.
就好像梁山英雄一樣.
全部都去了.
那個就是.
這些本來是.
可能是.
流氓 匪類.
流浪.
很痛苦.
但是他們去投靠大衛.
所以我們見到在.
這裡有一種很大的預表性.
詹正義的牧師.
是我的一位良師.
他寫了對大衛很多研究.
其中有一篇文章.
叫做《阿道難動的大衛》.
《阿道難動的大衛》.
我們可以見到.
就是講到這段的事情.

$^{241}$他說大衛在逃難的時候.
到處投靠別人的那些人.
就去到他那裡.
之後有很多這些人就投靠他.
在這裡他實在是學到投靠耶和華.
找到上帝是他的避難所.
所以他就呼籲更多的人.
來到好像他那樣.
來到投靠上帝.
所以這個就是有一種的.
不單止是當時的處境.
避難的處境.
也是大衛的心境.
是他一種的.
靈命的影響力.
好像讓其他這些的人.
都來到他那裡.
來到好像他一樣.
去投靠上帝.
找到真正的避難所.
所以這裡其實可以說是一種預表性.
其實在舊約聖經.
差不多每一本書都有預表性.
到了新約的時候.
當耶穌基督來到的時候.
就逐漸將這些預表揭露出來.
原來亞伯拉罕獻以撒.
是預表神獻耶穌基督.
所以好像類似這些預表.
在這裡預表性可以說.
大衛就是他們的頭目.
就好像今天主耶穌基督.
是我們這些困所流離的迷羊的主.
當我們來投靠這位主.
讓基督成為我們人生的主宰.
教會就是以耶穌基督為首.
所以生命就開始改變.
因為人不能夠找到.
真正在這個混亂世界.
心靈空虛.

$^{281}$我們罪惡污染的人生當中.
找到避難所.
以致我們能夠生命有改變.
所以好像今天我們每一個.
在教會裡面的弟兄姊妹.
或者你還未信耶穌.
在網上如果你參加的.
你明白到我們失散流離的人生.
是充滿著混亂和不安的這個世代.
真正的避難所.
就是主耶穌基督.
在當時大衛成為他們的引導.
可以說是好像牧羊他們一樣.
帶領他們去尋找大衛所投靠的這位主.
所以今天也是我們的需要.
在這個避難所裡面.
我們可以看見有些什麼的轉變.
有些什麼的影響呢.
第一我們看見這個避難所.
就是使關係分離者的重建.
這個就是避難所裡面奇妙的轉變.
因為第一節很明顯說.
他說大衛就離開那裡.
逃到阿杜蘭洞他的弟兄和他父親的全家聽見了.
就都下到他那裡去了.
哎呀這裡很寶貴.
他說他的弟兄和他父親的全家.
請問這些是什麼人來的.
家人.
如果大家還記得.
最初大衛在他的家裡面.
和這些家人的關係是怎樣的.
不是太好是吧.
因為我們知道他爸爸忽視了他.
每個兒子都有份.
大衛沒有份.
叫他做看養.
叫他去做信差.
叫他做那些不太重要的事.
阿哥那些重要的人.

$^{321}$所以就去到參軍.
做將領.
那些阿哥和他的關係好不好.
不好是吧.
他爸爸特意叫他送一些糧食.
送一些問候.
叫他報一下平安.
一見到大衛.
第一句說話.
豈有此理.
你為什麼不看養.
大罵他一輪.
他說你在那裡看熱鬧.
你是不盡忠職守.
無端端給他一些罪名.
痛罵他一輪.
可見在阿哥在爸爸的心目中.
大衛沒有什麼地位.
而且可以說關係不是那麼好.
很多時候被他們鄙視.
被他們貶低.
被他們覺得你不重要.
這個真的很痛苦.
不知道我們有沒有在這樣的家庭裡長大.
一定會後來影響我們自己很多.
想到很多時候沒有家庭的溫暖.
沒有家人的肯定.
沒有家裡的溫暖和愛.
一定會影響日後.
所以很寶貴的就是我們今天基督教.
不單止有一個信仰.
還有一個神.
然後我們有天父接納我們.
我們能夠得到醫治.
我們能夠得到神無條件的愛.
可以幫助我們來重建我們的心靈.
重建我們的人生.
重建我們的關係.
在這裡我們看到很奇妙.
當時這些家人.

$^{361}$他父家竟然投靠大衛.
當然可能是因為大衛後來不同了.
從一個無名小卒變成什麼.
變成戰士長.
當時他打贏了哥利亞.
後來承諾全家都免他們的納稅.
所以無端端真的因為這個兒子.
全家都得益.
所以當場刮目相看.
不過你要留意.
其實當時大衛是逃難的.
所以不是他得好處的時候.
不是因為這個兒子做官所以我們得好處.
逃難的時候都去他那裡.
就證明不只是因為好處.
他的哥哥本來是做蘇俄的戰士.
一些官長.
但當大衛做戰士長的時候.
相信很自然.
這些哥哥可能都做他的戰士之一.
但這個時候也不是因為跟著大衛做他的戰士.
戰士的兵.
所以要跟著他來.
肯定不是.
因為他正在逃難.
如果要自保.
多大都不要跟著他.
如果要找人保障.
我是大衛的家人.
萬一蘇俄殺了我怎麼辦.
我找其他比較安全的地方.
還跟著大衛.
反正殺大衛也殺了你.
所以很明顯這裡不是他們因為好處.
或者已經是跟著大衛.
而是他們真的.
他說他的弟兄和全家聽見了.
這個他們可以選擇的.
就到下渡他那裡去.
所以這個是他們選擇去大衛那裡.

$^{401}$我見到當時他們這段關係已經重建起來.
是一個美好的一種關係.
無論你怎樣也好.
我們都會跟從你.
這個是不簡單的事.
所以這裡我們見到實在是很寶貴的.
大衛又證明是好的.
跟著我們看見第三第四節.
我們讀完它,請讀.
大衛從那裡往摩亞的米斯巴去.
對阿亞王說.
求你容我父母搬來.
住在你們這裡.
等我知道上帝要為我怎用行.
大衛領他父母到摩亞王面前.
大衛住山寨多少日子.
他父母也住摩亞王那裡多少日子.
大家明白這個意思嗎.
所以其實可以投靠誰.
可以投靠摩亞.
父母來到他那裡.
大衛知道不安全的.
你投靠我不好.
所以就將他安置在那裡.
安置在摩亞王那裡.
可見大衛非常掛念父母的安危.
非常孝順.
他沒有因為父母過往對他的鄙視或貶低.
他現在這個時候沒有顧及他們.
相反他仍然表達他的愛心.
表達他的孝心.
所以後來他回到山寨去住的時候.
父母就安穩在摩亞的地方.
以致得享平安.
這種就是見到他們的家庭這種情況.
是見到這種和諧.
彼此相愛.
和互相支持這種情況.
他們都任用他的哥哥.
有一位學者McArthur.

$^{441}$他研究認為在歷代至上.
歷代至上27章.
當那裡提到大衛派一個叫做爾利胡.
爾利胡去管理整個猶太之派.
你知道大衛是從猶太之派而出的.
猶太之派後來代表南國.
所以他派一個很重要的人去管理整個猶太之派.
在歷代至上27章提到這個人叫做爾利胡.
爾利胡的名字這麼熟悉.
原來就是他哥哥的名字.
所以是不是這個哥哥的名字.
這位學者研究McArthur認為.
這個非常大可能就是他的哥哥爾利胡.
所以當時大衛安排.
其實他很多都是親信.
他的內閣後來成員的時候.
都是他的表親.
他的兒子做他的親信.
所以他差遣他的哥哥去做猶太之派的管理人.
是非常合理的.
所以你看到大衛和他哥哥的關係恢復.
沒有因為他哥哥以前看不起他.
現在他任用他哥哥.
所以你看到這種都是一種很寶貴的.
在這個洞穴裡面他們相處的時候.
當這些人來投靠他的時候.
關係就恢復了.
關係就重建了.
所以可以讓我們看見.
今天當我們來投靠這位主的時候.
聖經形容我們.
當我們這些困苦流離失索的人.
關係破裂的人.
當我們來到神的面前.
我們很怕神拒絕我們.
因為我們是一個罪人.
但是神告訴我們不是.
因為我來並不是要健康的人得益.
我是叫那些有病的人得益.
我是叫那些罪人能夠悔改.

$^{481}$耶穌親自來的時候.
就接納所有這些犯投靠的人.
無論我們怎樣都好.
只要我們願意投靠相信主.
主就接納我們.
無條件接納.
使我們關係能夠重建.
聖經裡面說我們逢來與神是隔絕的.
但是因著神.
我們的愛和耶穌基督的救贖.
就得以和好.
《哥羅世書》一章二十二節.
就講明這個道理.
我們一起讀一讀.
請讀.
但如今直像基督的肉身受死.
叫你們與自己和好.
都成了聖結.
沒有瑕疵.
無可責備.
把你們引到自己面前.
聖經裡面說.
但如今.
就是之前本來我們是與神為敵的.
但是如今直像基督的肉身受死.
便叫你們與自己和好.
所以我們都能夠與神和好.
而且不單止他改變我們的生命.
在他眼中.
他稱我們為異.
將聖結的神的生命賜給我們.
叫我們都成為聖結.
沒有瑕疵.
無可責備.
這是很完全的接納.
讓我們能夠引到他自己的面前.
所以.
只要人願意投靠這個靈魂的避難所.
希伯來書講到.
耶穌基督就是我們的避難所.

$^{521}$原來當我們來投靠這位主耶穌基督.
是我們的靈魂的避難所.
祂願意無條件地接納我們.
叫我們與祂和好.
不單止與祂和好.
還改變我們的生命.
叫我們能夠與其他的人和好.
我們很多的經歷.
因為一個人信耶穌.
家庭就開始慢慢和睦.
以前有很多爭吵.
因為主的愛.
主的生命進入了這個家庭.
他們開始重建.
關係開始重建.
因為主耶穌基督的愛.
和祂先接納我們.
以致我們能夠用祂的愛.
去接納我們本來不能夠接納的人.
Amen.
跟旁邊的弟兄姊妹說.
上帝接納了你.
只是你而已.
上帝接納了我.
第二個我們看到.
這個避難所不單止是.
關係者的重建.
更加是受困逼者的轉機.
因為接著說第二十二節.
第二節是這樣說的.
凡受什麼.
凡受困逼的.
所以我們看到這個避難所.
是那些什麼人去.
就是受困逼的人.
得到轉機.
因為在《二切學書》.
二十五章第四節.
《新譯本》這樣說.
我們一起讀一讀這節經文好嗎.

$^{561}$請讀.
你作了困乏人的保障.
貧窮人急難中的保障.
躲佈風雨的避難所.
避嫌疑的陰涼處.
所以《二切學書》說到這位.
梅塞亞救主.
他就是我們困乏人的保障.
我們發覺真是很奇妙的.
這個就是那些困乏人.
本來他們是受逼迫.
受困逼的.
困逼就是可能很露堵.
可能無端端失業.
可能無端端被人誤解.
可能無端端被人陷害.
可能在這個世界裡面.
忽然間好好的遇見一些事.
令到他落在一個困境裡面.
有沒有這樣的事.
人生事實上都會見到.
有這些事發生.
其實當時這些四百人當中.
很多都是這些人.
他們去到大衛那裡面.
很奇妙.
因為這樣投靠大衛.
而更加是投靠大衛所投靠的上帝.
他們的人生就有改變.
我們見到後來這些.
我們剛才提過.
大衛後來任用的就是這些人.
就是這些患難當中跟隨的人.
後來當他去德國的時候.
他就叫這些人成為他的內閣成員.
成為了他那些重用的軍官.
你想想本來那些.
好像梁山108英雄那樣.
很有武功的.
不過很流離失所.

$^{601}$不過去到的時候後來就得到正義.
就是重建他們的人生的意義.
做一些正確的事.
所以生命就改變.
你也看見後來這些困境的人.
他們甚至成為尊貴的人士.
成為了後來在歷代志上稱他們為.
稱他們為耶和華的軍隊.
真是厲害.
耶和華的軍隊.
所以他們當時這班本來是.
這些受困逼的人.
結果就成為了耶和華的軍隊.
這個叫做為主而活.
我們看到本來這些罪人.
後來成為了叫我們做基督耶穌的徵兵.
我們都是基督耶穌的徵兵.
本來我們是受困逼的.
沒什麼人生意義的.
不知道為了什麼而活.
現在我們為主打仗 為主而活.
Amen.
這些幸福小組的功勞很Amen.
我們在做幸福小組打仗.
講一個故事給你聽.
一個見證來的.
在西方很多宣教士來我們中國.
今天我們中國人得到福音.
是因為很多西方的宣教士.
從英國 我們知道大德生先生.
是一個最早期的更正教的宣教士來到中國.
他是倫敦會英國的宣教士.
有些宣教士都是這樣.
去到中國將福音傳給我們中國人.
所以今天我們能夠蒙受福音.
都是一棒傳一棒.
一個傳一個.
所以我們成為受福音的人.
又成為什麼.
成為傳福音的人.

$^{641}$我們領受福音又成為傳福音的人.
當時他們藉著醫療在街上報道.
同時他們印製了聖經.
當時已經有中文的新約福音書的聖經.
所以他們將單行本.
大家有沒有聽過.
單行本即是一本福音書.
整本送給你就是馬太福音.
或者送給你就是約翰福音.
這些單行本的中文聖經.
有馬太有什麼.
馬可有盧迦有約翰.
這四本已經翻譯了.
所以他們拿著到處去送這些聖經.
不過送很多時候是不值錢的.
上課說免費.
上了課就不來了.
所以酌量收一點.
你買回來值錢一點.
買回來就珍惜一點.
不要浪費了.
兩個銅錢很少錢就買回來.
所以他們就去買勸人買.
很多時候他們穿著長衣服.
樣子就不是中國人.
是什麼人.
我們中國人叫他做什麼.
叫他做洋人.
洋人鬼佬.
很有趣.
這些洋人穿著中國衣服.
所以很多人就跟著他.
包括那些心蚊仔.
心蚊仔最喜歡看熱鬧.
這麼有趣.
樣子不同我們.
不過穿著我們的衣服.
所以很多時候就跟著他們一大群人.
所以這兩個宣教士就去到一個大店鋪.
專門賣雜貨的.

$^{681}$請他們要不要買.
好好的.
看看吧.
兩個銅錢很便宜.
結果很多心蚊仔.
很多人跟著他們.
這個老闆就覺得很煩.
又不是來買東西的.
來什麼.
來賣東西.
不但是阻礙他做生意.
所以怕了你.
幫你買了.
就買了他四本科學書.
就把他扔了.
反過來就把他扔了進垃圾桶.
後來他有一位店員.
就好奇.
老闆給錢買了這麼快就扔進垃圾桶.
是什麼來的.
就拿來看.
越看越有味.
越看聖經的話語.
就像兩刃利劍刺透了他的心.
他就開始把那四本都撿起來.
誰知道被其他的同僚看到.
就告發他.
老闆你扔了那些東西.
現在他去撿.
老闆就很生氣.
因為豈有此理.
你不是專門去撿東西.
你是不是平時專門去撿我放下的東西.
偷東西.
是不是你.
上個星期不見了可能又是你.
無端端就扔了一些罪名給他.
而且就罵他一輪.
接著怎麼樣.
走.

$^{721}$哎呀真是慘了.
我好心撿了一些東西來看.
而且這麼好的東西.
結果就被人炒魷魚.
這個人就很慘.
到處走.
心裡很難過.
因為他還要養家.
又失去了工作.
非常的痛苦.
這個時候他就不知道怎麼辦.
怨天尤人又不是.
總之心裡就充滿著群敵.
當他經過一個地方.
忽然間看見什麼什麼什麼塘.
不是玲瓏塘.
不過是教堂.
所以他就直奔進去.
聽見裡面在唱歌.
他就見到賣這些.
單行本聖經的宣教士.
在說道.
所以他就坐下來聽.
越聽就越入心.
越聽就越有聖靈的感動.
他就真的願意去認識信仰.
結果他每個星期就這樣.
接著連續幾個星期都回來.
一邊聽一邊聽.
跟這個宣教士就成為了好朋友.
結果他就信了主.
受了洗.
成為了教會裡面的一個.
事奉的人.
所以他的生命就開始改變.
他從漫無目的.
而且受著困逼的時候.
生命開始改變.
結果美國後來就差一宣教士.
就是美爾美會.

$^{761}$就差一宣教士去中國.
就發覺當地那些人不懂中文.
就很難傳教.
你知不知道當時有很多這些笑話.
我們叫西人做甚麼.
洋人.
洋人見到中國人就叫我們做甚麼呢.
就唐人.
所以那些宣教士就傳道傳道.
我好喜歡唐人.
想認識你們.
你們聽不明白.
那就糟了.
聽到這個就每個人都走了.
先不要走.
我要拿你的命.
那怎麼傳道呢.
所以他就找倫敦會.
倫敦會你們認識了當地的一些同工.
可不可以借兩個給我們.
這個宣教士聽到美爾美會這樣講.
大家都很喜歡.
一個英國來.
一個美國來.
當然合理.
借一個像給他.
借誰呢.
就是這個.
這個叫陳大勇.
他上榜了.
這本書寫叫中國信仰英雄榜.
有機會找找看有沒有找到這本書.
叫讀回書,信回神,做回鬼.
叫陳大勇.
其他有些你們認識的.
比如梁發.
徐光啟那些.
大家如果讀過中國歷史就知道.
這個陳大勇上榜.
榮記的榮.

$^{801}$所以就借了陳大勇給他.
結果呢.
真的奇妙.
他就成為了美爾美會.
第一個.
案立的中國人牧師.
奇妙吧.
我們人生當中有時那些困迫.
可能是要神催迫我們.
去投靠他.
這位陳大勇.
他就是因著這個人生的轉變.
去了一個困迫的人生.
落難的人生當中.
投靠神.
生命改變.
後來被神使用.
結果.
他真的成為了耶穌基督的正病.
所以真是主耶穌今天都這樣呼籲我們.
很熟悉的聖經一起讀嘛.
凡勞苦擔重擔的人.
可以到我這裡來.
我就使你得安息.
受困迫的人.
得到真是生命的改變.
我們第三.
這個避難所是什麼.
避難所是欠債者的拯救.
我們繼續看這節經文.
它說.
凡受困迫的.
欠債的.
所以這裡提到這些欠債的人.
他們得到拯救.
這個避難所是什麼.
避難所就是這些欠債人的拯救.
當時他們都是.
因為他們貧窮.
有很多那些市井之民之徒.

$^{841}$所以他們在這裡借一些那裡借一些.
這樣來生活.
結果他們聽到大衛.
他們投靠.
去到以大衛為他們的頭目.
投靠這位大衛所信靠的上帝.
欠債的.
你們有沒有欠債.
有沒有誰欠債.
你說我沒有什麼欠債.
車是公園.
屋是公園.
有沒有.
很開心.
不過一定欠債.
因為現在我們個個都欠債.
欠誰.
不要這麼快這麼熟悉零紙.
我們欠信用卡.
有什麼不用信用卡的.
全部給現金有沒有.
我相信現在這個時代沒有了.
起碼都欠一個月.
持交那些更欠多一點.
所以不要利搭利利搭利.
欠極都欠不完.
你知不知道在2025年.
就是去年.
平均美國欠信用卡數的人.
每一個人平均.
大約是6300至7300元.
你們差不多吧.
省一點用.
每個月大概欠6000至7000元.
可能都有了.
交什麼費用.
現在用信用卡交那些.
或者日常生活.
平均吧.
有些當然欠多一點.

$^{881}$有些欠少一點.
所以這個數字應該挺準確的.
2025年.
如果利搭利利搭利.
就不止了.
越欠越多.
永遠還不清.
現在回到屬靈那裡.
我們屬靈有沒有欠債.
有啊.
屬靈欠很多債.
你想想如果我們一天犯100件罪.
沒那麼多吧.
10件吧.
都沒那麼多.
3件吧.
有些事無意中得罪了神.
有沒有啊.
3件.
當作差不多吧.
3件一個月就多少.
90至100吧.
一年吧.
當1000吧.
10年吧.
一萬.
你現在幾歲了.
說得.
那你就乘一乘吧.
那條數很金計.
一天3件得罪神的時候.
一生有多少件得罪了神.
如果神要跟你計算.
算一算.
那條債就利搭.
他沒有加利息.
應該加利息.
加上利息就不止了.
還多少還不清.
所以我們每一個人都欠神債.

$^{921}$有沒有啊.
有.
你們現在拿錢.
還有沒有去ATM.
沒有了.
還有的.
ATM即是自動櫃員機.
進去之後有很多按鈕.
有一個按鈕叫deposit.
我記得在疫情之前.
有時支票儲起來.
哪有空逐張進去.
有一次過去自動櫃員機.
有一次過進去.
要按deposit.
deposit即是甚麼.
存款.
還有一個按鈕叫withdraw.
withdraw是甚麼.
提款.
你喜歡按哪個按鈕.
誰喜歡按deposit.
誰喜歡按withdrawal.
這要看性格.
有些人希望存款越來越多.
當然希望存款越來越多.
deposit就好.
有些人覺得存那麼多錢來做甚麼.
這是number來的.
我不用它做甚麼.
他希望用的時候.
將number將這些死的東西.
變成一些他喜歡生的東西.
好的東西.
他喜歡withdrawal.
所以要看性格.
不過還有一個按鈕.
還有一個按鈕叫甚麼.
有很多按鈕.
有一個按鈕叫transfer.

$^{961}$今天我們可能很多人不用ATM.
要用上網.
online.
一樣.
現在我奉獻是用上網奉獻.
不是實體奉獻.
就是withdrawal.
但其實用這個transfer.
transfer去一些更好的地方.
其實上帝.
他就transfer給我們.
在聖經裡面說.
告訴我們.
本來我們每個人都是偏行幾路.
在二作書很熟悉的經文.
他說我們如羊走為.
各人偏行幾路.
耶和華洗我們眾人的罪孽.
都歸在他的身上.
原來我們可以將我們的東西transfer給神.
那些債神幫我們還.
聖經又更加說.
他又將transfer給我們.
一起讀好不好.
請讀.
就是神的義因信耶穌基督加給一切相信的人.
就是將神的義.
神的義是甚麼.
就是神的聖潔公義和他的生命.
將這些因信耶穌基督.
加給一切相信的人.
他將他的transfer給了我們.
將我們的transfer給了他.
就是說我們本來有的罪債.
現在耶穌說我幫你償還.
耶穌在十字架上面代替我們的罪.
為我們一次成為了這個贖罪制.
只是一次就不需要再有多次了.
因為他一次就可以來還清我們一生.
一生所有的罪債.

$^{1001}$Amen.
不單止這樣.
他更加將他的義.
來將耶穌基督的救贖.
這個成就的恩典.
耶穌說成了.
accomplish.
將他轉給我們.
今天我們能夠承受了神的義.
得著不單止罪得赦免.
聖經說.
我們更加成為了上帝的義類.
得著神的生命.
Hallelujah.
所以凡欠債的人.
耶穌當我們投靠這個避難所的時候.
他為我們償還了一切的罪債.
最後我們更加看見這個避難所是什麼.
是俘虜者的安慰.
最後這篇聖經我們一起再讀一次.
凡受困逼的,欠債的,心裡苦惱的都聚集到大衛那裡.
人生真是很多的苦惱.
有一位患重病的弟兄.
他身體極度的衰弱的時候.
他深深有個體會.
他說人生不是像錄影機一樣.
錄影機是怎樣的呢.
錄影機就是我們過往所有的事都是歷史.
如果我們存留下來.
全部都可以成為了我們的記錄.
所以這個痛苦的那一刻.
最好是怎樣呢.
有一個按鈕.
叫做Fast Forward.
Fast Forward就是咻地跳過了.
如果有就很好.
每逢我們面對苦難,苦難,困難.
那怎麼辦呢.
Fast Forward就走了.
他說人生不是像錄影機一樣.

$^{1041}$可以按一個按鈕就越過了不好的部分.
是的.
我們的確要面對人生的生,老,病,死.
這些對我們越來越真實.
但是他說他深深體會什麼呢.
他體會到每一格電影.
即是每一格Digital.
每一格即是每一個片段.
主耶穌基督都在那裡.
只要我們認定,投靠祂.
祂就成為我們的力量和安慰.
這是他在患重病人生階段的時候.
深深體會到的一種靈慧.
真的很寶貴.
真的我們沒有一個人能夠不面對這些.
我們需要面對.
但是我們知道主與我們一起面對.
在25年前911的時候.
很快過了25年了.
2001年的時候.
當時又發生了911之後.
很多人就去到曼哈頓.
有一個長老會.
是一位很出名.
今天已經過世的牧師.
Tim Keller.
長老會的教會.
當時Tim Keller.
看到忽然間來了很多.
比平時一倍以上的人.
有600至800的人.
忽然間.
即是今天我們坐在這裡.
本來只有100人.
忽然間多了300至400人湧進來.
那麼多人進來的時候.
他要臨時改.
好像我昨天在幸福小組的情況.
看到有些人不同了.
我馬上改信息.

$^{1081}$改了第二篇信息.
所以這個就是當時他要改信息.
改什麼講什麼好呢.
聖靈就感動他.
他就講了一篇信息.
結果後來很多很多人得到幫助.
信耶穌,然後教會.
忽然間在那段時間有很大很大增長.
當然當時是一個苦難的時期.
所以他那篇信息就是說.
基督教的信仰是唯一在這麼多宗教裡.
告訴人.
神在不公平.
在苦難.
在迫害.
在無辜的情況下.
神祂自己為我們受苦.
主耶穌基督就是這樣.
在人的無辜的迫害.
唾罵.
攻擊.
毀壞.
不公平.
暴力的底下.
為我們受苦.
所以主耶穌祂能夠明白我們人生各樣的苦痛.
忽然而來的災難.
極大的傷痛.
所以今天.
實在是.
耶穌為我們受過最大的苦難.
祂能夠也成為我們在苦難當中最大的安慰.
祂賜給我們力量.
祂是我們生命當中靈魂的避難所.
當時大衛深深體會到.
所有世間上的無論名利.
無論其他所有的人士.
都無法成為祂最終的依靠.
祂投靠上帝.
在那裡在洞裡邊.

$^{1121}$祂體會到.
就好像山洞一樣.
成為祂人生的避難所.
弟兄姊妹.
我們都需要這種避難所.
如果你還未信靠主耶穌基督.
我們更加求神幫助你.
不要再等待.
在你面對人生各樣的患難.
苦痛.
迷失.
迷惘的時候.
你會見到人生真正的意義.
以致我們能夠由這位主.
親自成為我們的保障.
成為我們人生的避難所.
耶穌基督.
祂忍受了困迫.
擔當了我們的罪債.
經歷了人生的苦難和掙扎.
以致祂能夠成為我們隨時的避難所.
使關係分離者得以重建.
與神和好.
與人和好.
使受困迫的人得到轉機.
以致能夠找到新的人生方向.
使欠罪債的人得到拯救.
罪得赦免.
使苦惱的人得到安慰.
得到新的力量.
能夠生命得以轉化.
成為耶和華的軍隊.
基督的精兵.
願意我們都這樣.
來自私人的呼籲.
這篇詩篇是在911.
很出名的詩篇.
911篇.
第一節就是投靠耶和華的隱蔽.
接著第二節.

$^{1161}$我們一起來讀這一節.
我要論到耶和華說.
祂是我的避難所.
是我的山寨.
是我的神.
是我所要求的.
我要論到耶和華說.
祂是我的避難所.
是我的山寨.
是我的神.
是我所要求的.
頭兩個是真實的環境.
是避難所和山寨.
後面是一個真正的體會.
不只是環境上的避難所.
是我們生命.
人聖靈魂的避難所.
所以祂是我們的神.
是我們一生的依靠.
Amen.
我們低頭祈禱.
請同心禱告.
今天神願意幫助你和我.
成為我們人生所投靠的避難所.
讓我們能夠找緊基督的應許.
明白到主在我們身邊.
讓我們投靠這位.
為我們擔當罪孽.
忍受困逼.
承擔苦難.
但是卻以德性的主.
以致我們能夠靠著基督.
成為新造的人.
找到新的人生方向.
成為主所使用.
耶穌基督的正兵.
奉主耶穌基督的名.
Amen.
\newpage



\section{傳道書 4:9-10}
\label{sec:031826}
\textbf{多代同行 (2) PeerWalk (2) 誰是你的同行者?}~~~~吳燕清傳道
\newline
\newline
連結: \href{http://www.church.com.hk/acms/content.asp?site=cdc\&op=show\&type=product\&code=031826\&layout=sermon}{\texttt{031826}} ~~~~ 語音日期: 2026-03-14
\newline
\newline
\hyperref[sec:031825]{\small{< < < PREV SERMON < < <}}
~
\hyperref[sec:index_chronic]{\small{[返順時目]}}
~
\hyperref[sec:031827]{\small{> > > NEXT SERMON > > >}}
\newline
\newline
傳道書 4:9-10
\newline
\begin{longtable}{cl}
\hline
\hline
章節 & 經文 (和合本修訂版)\\
\hline
4:9 & \begin{tabularx}{0.7\textwidth}{X} 兩個人總比一個人好，他們勞碌同得美好的報償。 \end{tabularx} \\ \\ \relax
4:10 & \begin{tabularx}{0.7\textwidth}{X} 若是跌倒，這人可以扶起他的同伴；倘若孤身跌倒，沒有別人扶起他來，這人就有禍了。 \end{tabularx} \\ \\
[1ex]
\hline
\hline
\end{longtable}
$^{1}$今年整個童福的主題就是多代同行.
我們剛剛過了一個季度.
就開始帶大家去看如何跟神的同行.
接著我們就進入了同輩的同行.
同輩當然就是你們.
話都說少就是大家出來工作三五七年.
沒有人在這裡工作超過十年.
除了我,有些人有.
但是雖然你們工作的時間有差距.
但你的年紀可能你只上了一堂崇拜.
其實你的年紀就應該很接近的.
就是說同輩同行的時候.
我收到這個題目的時候.
我就想起這段聖經.
當然我先講這個題目.
我定了個題目叫瑞士的同行者.
我們經常講的.
我不知道你回來教會這麼久.
這個世界就會這樣說.
你有沒有知己.
如果是年輕人就會說.
你有沒有什麼BFF.
你明白是什麼嗎.
BFF還是什麼BFF.
聽到什麼什麼BFF.
那些代稱我都不太明白.
我當然知道什麼什麼BFF.
但是我想在神的家裡其實更重要的一件事.
在這個群體我們經常講.
同行者是非常重要的.
當你望著這幅圖畫.
幫你代入.
這個AI Gen對你來說.
望下去你都能看到.
遠遠有條很長的路.
你可以說山明水秀一邊.
一邊就簡直是風雨交加的情況.
如果是行山的人就會很明白.
在這些情況之下.
如果讓你可以選擇.

$^{41}$在面對這樣一個環境.
就算你望著前面有多風和日麗.
但有些東西你都看不到.
但是很漫長的一條路的時候.
你會怎麼選擇呢.
你寧願一個人自己走.
還是你想有伴去走.
你自己在心中想想你的答案.
我相信有些人會這樣說.
我自己就夠了.
哪怕就算電視宣傳.
叫你每一個人去行山.
很多人都專門喜歡一個人去行山.
就是那些自己望著的人.
他自己就喜歡去的.
他不用擔心.
他覺得應該很安全的.
不會那麼容易出事.
可能就望著在這種環境底下.
有人都會這樣想.
我想只是自己.
哪怕有些人可能心底裡你自己知道.
怎知是我自己來的.
其實我不想靠別人.
免得被別人覺得我作弊.
甚至乎有些人會覺得.
要找人幫忙.
或者要求幫忙.
某程度上就像是一個弱者的表現.
我自己搞定.
當然也有部分人會說.
如果我看到這種情況.
那不是的.
我怎樣都拉著我老公去.
你怎樣都可能把他拉走.
你覺得有個伴.
你就會覺得安心一點.
無論你的答案是哪一種.
我覺得我們不可以不承認一件事.
如果你望著剛纔那幅圖一樣.

$^{81}$人生走的那條路.
怎止像這張圖那樣的狀況.
可能有高有低.
未必那麼平坦.
很多山泥.
可能山泥傾瀉會出現.
很多沙石又會出現.
不是你所想的那種狀況.
但我覺得上帝給我們一個.
在我們這個信仰群體裡面.
神其實在祂的說話裡面.
給了我們很大的保障.
叫我們知道.
其實祂為了設立教會.
凝聚我們這群人.
其實有一個很重要的目的.
給你和我.
我們怎樣可以跟從上帝這條路里面.
你可以成為別人的同行者一樣.
在我們繼續講下去之前.
不如我們一起去祈禱.
我們再繼續講神的話.
天父我再一次去到禰的面前.
我特別求聖靈.
禰在我們的中間.
禰自由的運行.
神啊願禰的心不單單得到滿足.
主禰叫我們的心能夠與禰的心對齊.
而且今天我們眾人坐在這裡.
無論我們今天所遇見的事情.
有多少東西可能未拆解到.
或者有多少工作我們未完成到.
但主要是願意我們這個心.
再一次去到禰的面前.
在這一刻我們說.
我們單單的去愛禰.
單單的聚焦在禰的裡面.
因為神啊我們說.
我們很想更加明白禰.
認識禰.

$^{121}$主要是祝福著我們.
與我們的同在.
聽我們的祈禱.
奉耶穌基督的名字祈求.
Amen.
第四章的時候必須掌握一件事.
當講傳讀書.
當然這本書在聖經是這樣講的.
一卷的智慧文學.
寫書的人大部分是講所羅門.
是一個很有.
無論是昔日的那種.
你可以說很古舊的那種.
甚至是到了今天.
我們在現代社會里面.
他所提到的智慧的點滴.
是很真實的.
是很.
不是應用不了的.
有時候我們會說.
哇 聖經講的好像很遠古.
很遙遠的.
你會覺得有一種落差.
不是在你的現實狀況.
但我要告訴你們.
不是的.
你看這本智慧文學書.
其實聖經裡面.
有很多在上帝的話語裡面.
隱藏著很多寶藏.
我再次講.
其實真的有很多寶藏.
不過問題是.
你和我還沒有發掘到它的寶藏.
當耶穌說.
祂來要叫我們得生命.
並且得得更豐盛.
我不知道你覺得你自己.
在你認識了耶穌之後.
你的生命能夠經歷多少豐盛.

$^{161}$如果有一刻你會說.
不是的 我真的沒有覺得那種豐盛.
有如你剛纔所講的聖經一樣.
那我就告訴你.
在上帝的話語裡面.
確實埋藏了很多應許.
甚至祂的說話.
今天怎樣能夠成為你和我的祝福.
不過很多時候我們拿不到.
或者我們知道了.
但我們未必能夠.
經驗到上帝的說話的大能.
所羅門寫的傳道書裡面.
其實他講到有幾個重點.
他說人生不單單短暫.
我們很明白這件事.
你的生死你控制不了.
我們在所羅門講的.
其實人生有很多東西.
他經常用一個字叫做虛空.
對我們來講怎樣理解虛空.
當然你可以用這個字.
再用中文去演繹.
可能就是短暫 無意義.
表達一種其實你面對人生.
有很多東西你和我都一樣.
隨著年日.
我們在無論努力追求的裡面.
到最終其實那些東西.
都是可能一觸即逝.
甚至說傳道書裡面.
所羅門更加講一樣東西.
他說日光之下無新事.
什麼意思.
其實好像循環一樣.
人人都處於那種狀況.
當你有一刻你會覺得.
為什麼那個人會這樣.
當你看著你媽媽.
或者像我這樣.

$^{201}$看著你媽媽看著你爸爸.
你會覺得有一些東西.
他為什麼這麼不明白.
為什麼這麼固執.
我想說輪流轉.
到你差不多這個年紀的時候.
可能你都會像你爸爸和媽媽一樣.
你都會在那些位置.
都是一些很固執的位置.
都是一些我們說.
就好像所羅門所講.
真的日光之下無新事.
在這個循環裡面.
你經歷的東西.
甚至你的父母經歷的東西.
我們很多的狀況.
其實是很雷同.
人在整個的歷史.
甚至我們這個生命.
在這個循環裡面.
我們看起來好像沒什麼突破.
人人都有一死.
我們說知道的.
但是面對著整個的.
當你在工作.
你放了工.
然後你睡了覺.
然後你又去工作.
我不知道你在這個循環裡面.
你的感覺是如何.
傳道者在所羅門裡面.
他在這本書.
十二章裡面.
他帶出一個很重要的信息.
他說就是在這種狀況裡面.
人有什麼出路呢.
他就提一個highlight.
就是整個傳道書裡面.
這個是一個重點的highlight.
就是他告訴人聽.

$^{241}$敬畏上帝纔是真正的答案.
就是說當你的人.
你以為好像追求很多.
在物質裡面的需要.
甚至你在生活裡面那種滿足的時候.
你要搞清楚一件事.
這些東西都會過去.
這些東西你以為你捉得緊嗎.
都不是捉得緊.
哪怕是你的錢也是.
你以為你可以怎麼投資.
拿到多少.
將你現在有的東西.
你當是double triple.
你可能會想著再拿多少都好.
你要記住一件事就是.
如果按著所羅門這麼有智慧的這位人.
他在告訴你.
其實有些東西.
不單單你眼前所看為重要的東西.
那個不是你應該視為最重要的東西.
他說敬畏上帝纔是真正的答案.
當然面對著整個世界有很多的矛盾.
對所羅門來說.
敬畏上帝就是進行他的說話.
是人面對這個人的一生.
你唯一應該要走的路.
不過這個我們回來教會很久了.
我們是很清楚的.
或者我們聽了很多.
但對你來說.
敬畏上帝.
甚至我們說是進行他的說話.
我們會聽多少.
我們知道多少.
甚至我們走了多少.
全圖書的四章.
九至十節.
這個是我今天想特別跟大家講的經文.
我們一起讀好嗎.

$^{281}$預備一二三.
兩個人總比一個人好.
因為二人勞碌同得美好的果效.
若是跌倒.
這人可以扶起他的同伴.
若是孤身跌倒.
沒有別人扶起他來.
這人就有禍了.
其實如果你再看全圖書的四章.
九至十節.
其實他十一十二節都繼續講.
形容著一樣東西.
人不能自己一個人.
哪怕你可能會覺得.
如果按照聖經所講.
當第一句表達著.
兩個人總比一個人好.
很多人可能會有個解讀.
所以結婚是最重要.
我要找個盤.
兩個人總比一個人好.
如果你這樣去解讀的時候.
你可能未解得全透.
在經文裡面所講的東西.
所羅門強調什麼呢.
就如我剛纔所講.
他特別要強調一樣東西.
一個人不能孤立.
只能成為一個.
等於有人會覺得.
信仰是我的.
我自己能處理.
你不要搞我.
教會都可以不用回來.
甚至乎什麼團契小組都好.
所以你不要來搞我.
我自己一個人.
我想告訴你.
這是違反聖經所講的.
因為聖經不是這樣講的.

$^{321}$聖經是在所羅門強調.
我們人在上帝的設計裡.
是需要我們彼此的合作.
和你需要有陪伴的.
你纔是最重要的一樣東西.
不只是有個盤這麼簡單.
令你開心.
背後上帝是講一樣東西.
是想我們能夠完成.
他託付給我們的東西.
你們聽清楚.
有時候我們在想一個盤.
就是我們想開心.
沒有這麼沉悶.
那就好一點.
有時候你會覺得.
不是的,我一個人都可以的.
OK的,可以的.
你要聽清楚.
上帝始終在創造裡面.
你已經發現.
為什麼他設計.
當有個夏娃出現.
你見到經文的形容.
也是在表達強調一個人不好.
要多一個人.
好像我剛纔所形容.
如果單憑數字去理解.
一和二的時候.
其實經文不是在表達.
OK,從此以後.
你就由一個變成兩個.
那你就沒事了.
我想告訴大家.
其實他是想強調一樣東西.
不是一變二.
是強調你不要一.
你要get到一樣東西.
你不要一個.
那個不是強調.

$^{361}$當然按照經文.
在傳道書四章裡面所形容.
他在說兩個人比一個人好.
你有沒有想過.
兩個人比一個人好.
其實他是在強調.
比起你一個人.
一個更多的一個人.
不只是指兩個.
可能三個人也好.
四個也好.
你有沒有想過.
多一個人在過程中間.
更能夠讓我們彼此的配搭合作.
完成上帝所託付我們的.
所以我很想強調一樣東西.
其實在我們整個.
我們說的屬靈的生命成長裡面.
如果說我們需要同行者.
你不是在找一個人.
你更加不是在想.
你只是找多一個人.
我就翹起手來.
不是的.
對所羅門來說.
他說同行者.
正正就好像他在傳道書.
整本書裡面所說的.
如何面對這個短暫的人生.
如何面對在這個世界裡面.
你確實有很多不同的事情.
充斥著我們的生活.
我們哪怕都可能.
面對著那種你感到虛空.
你覺得乏力.
你覺得無助.
你覺得在這種循環不息的裡面.
你看不到出路.
但是所羅門強調.
當有同行者在幫助著我們每一個.

$^{401}$面對著這個人生挑戰的裡面.
我們如何可以並肩一起前行.
所以他提到人生需要同行者.
他更強調一樣.
他說我們不要成為一個孤單的獨行者.
不要自己一個.
不僅僅是在你的生活.
包括在你的工作.
包括在你的家庭.
人生難免會遇到很多不同的挫折或者挑戰.
有同行者.
在過程中除了有人扶持你.
在你有限的時候.
感到有絕路的時候.
就好像所羅門所說.
若是跌倒.
這人可以扶起他的同伴.
若是孤身跌倒.
沒有人扶起.
這個人就有禍.
跌倒其實有很多含義.
其實在我們的人生裡面.
無論是身體上的軟弱.
甚至在工作上所面對的挑戰.
屬靈生命裡面你所遇見的試探.
甚至可能正正就是你正在處理的弱點.
保羅提醒.
我們除了要有同行者.
能夠面對我們的人生.
更重要的一件事.
是能夠發揮到同行者在你的屬靈生命前行裡面.
扶持的作用.
今天我們生活在一個.
確實是我們說的科技非常發達.
很多人.
現在的科技發達.
我已經不是第一次聽.
很多人沒有朋友.
聊天到一個地步.
只是和AI聊天.

$^{441}$我不知道你有沒有.
你是不是.
有時候有沒有一些地方.
你找不到之後.
你發現用AI說.
我前兩天.
YA不在.
我在Cosplay見到他.
我就問他.
為什麼你不用上班.
他就說.
我已經下班了.
我說是下班了.
他說我現在正在做些事.
我說是呀.
我剛剛纔和AI聊完天.
我正在問他.
他正在思考.
他是我.
我想我這兩個星期之內.
見的第二個.
又和AI聊天.
AI是他的朋友了.
我覺得在我們整個的.
整個的.
科技這麼發達的裡面.
其實我們怎麼運用.
我們有的東西.
可能對你和我是一個很大的挑戰.
我們擁有的東西.
甚至是哪怕你有IG.
還是有不同的社交媒體.
你們的聯繫都好.
我們在那個的.
我們在那個的溝通.
或者在那個的.
你覺得你有很多的朋友.
其實有多少人.
是能夠和你交心.
我上網看一些研究.

$^{481}$職稱對你們來說.
所以特別為你而設.
看你們覺得有沒有共鳴感.
職稱怎麼看同輩的關係.
原來有幾個特點.
其實在你們這個年紀的人.
其實你們很重視同輩的關係.
就是你坐在這裡.
我有時都知道你們的資訊.
你們有遊戲日.
玩遊戲.
喫的.
今晚又有東西喫.
有Ruby在就總有東西喫.
又玩.
生日會.
我相信大家都很重視朋友的關係.
不過.
這個發言是這麼說的.
雖然大家很重視同輩的關係.
但是關係往往較為表面.
你不需要額頭.
你自己在心中.
想一想.
我們很渴望有被人理解.
我們很渴望有朋友支持.
我們很渴望在面對職場的壓力.
甚至你在生活困難的裡面.
你可以有伴.
不過隨著生活很急速.
時間的有限.
回到來.
其實大家做得這樣就未必做到那樣.
我們會不會停留在一種表面的交流.
最近上班怎麼樣.
挺好的.
我經常都挺好的.
怎麼好呢.
不錯.
不用加班.

$^{521}$挺好的.
準時下班.
還可以上班.
我們可能都是這樣一個表面的交流.
但我們卻缺乏了一個深入的連結.
第二件事他有說.
因為整個職場的氛圍.
充滿競爭.
很多職稱可能內心有很多的害怕.
害怕什麼.
比較.
擔心去顯露自己的弱點.
選擇就是保持距離.
內心就瀰漫著一種很強烈的孤單感.
暗愷你自己知道.
問到你做什麼工.
未實際估到你每個月有幾天休息的時候.
只要你做哪一行.
你就可能會覺得.
他的工資在拿我的三倍.
你自然就會捉起一堵牆.
可能表面上你有很多朋友.
但內心你自己知道.
你感到的孤單.
缺乏真正你可以信任和依靠的人.
你是不是呢.
這個是這樣說的.
認同感.
其實認同感也不只是你們.
但是確實是.
在面對你們這個群體裡.
其實我覺得很難得.
大家志趣相投.
有一些共同的價值觀.
人生的目標.
但是同時.
這個研究這麼說.
職稱的群體.
往往在這種氛圍下.
是很容易有那種圈子出現的.

$^{561}$可能你的同事圈都是三五知己.
那些朋輩混在一起.
甚至在我們這個群體裡.
我們會不會不經不覺.
都有這些圈子出現.
隔壁的人.
我不是想指名道姓.
不過隔壁人.
可能你和那些隔壁友會特別投契.
甚至可能生了B的人.
那些B黨的人.
就會無端端有些圈子出現.
是要小心的.
會不會在不經意之下.
我們是在排擠一些人.
在這個群體裡.
作為職稱群體的你.
你如何形容你在這個群體.
不要想太多.
想想這個群體.
你的關係是怎樣.
會不會停留.
就好像我剛纔所說.
留於表面.
更甚一樣東西.
在這個屬靈的生命裡.
大家有沒有交流,分享.
不是喫不重要.
不是玩不重要.
不是喫蛋糕,生日會那些不重要.
不過如果我們在群體的聚集.
或者群體走在一起.
沒有在屬靈的生命裡.
我們真的有作出那種交流和分享.
我就會說.
我們缺少了一些很重要的部分.
同行者非常之重要.
所羅門所說的重要性.
其實從整本聖經.
都有很多地方提到同行者的重要性.

$^{601}$第一.
我想說.
其實上帝的創造裡.
人就像我剛纔說的.
是需要關係的.
由耶和華所說.
乃人獨居不可.
我為他做一個配偶,幫助他.
人本來就是被作為.
需要彼此的扶持.
不單單只是指向婚姻的關係.
也反映人作為人類.
我們需要有朋友.
我們需要有群體.
這是很重要的.
可能在我們當中.
有很多結了婚的弟兄姊妹.
當然你的配偶.
是否能夠成為你的同行者.
這個要想一想.
理性上夫妻應該是最親密的同行.
但在現實的生活裡.
當你的B開始出世的時候.
當你在想你的B的學習.
入哪間學校.
他究竟如何鋪排他的一生之路的時候.
我不知道你們夫婦.
在面對著B的同事.
你有多關心.
你們彼此的屬靈,心靈的狀況.
我們有沒有溝通.
我們有沒有分享.
我們有沒有繼續彼此的守望.
將大家在你的園約.
你的挑戰裡.
可以坦誠的分享.
我想再重申.
基督徒的生活.
是需要彼此相顧.
好像希伯牛書那樣說.

$^{641}$彼此相顧不是單單聖經的教導.
更加是在我們屬靈成長當中.
不可或缺.
我想提一提.
我想我預備這個信息的時候.
都是上帝提醒我.
我猜有很多特別弟兄.
我特別要跟弟兄說.
如果你心底裡.
常常真的會這樣想.
你有事.
一切很多時候.
你都是想著自己解的.
解決的意思.
你覺得要找人.
是一件羞恥的事.
我為什麼要告訴別人.
我有需要呢.
不想.
當我在預備這個信息的時候.
我就覺得.
當經文提到.
我們彼此相顧激發愛心變來時.
那個是沒有性別的分野的.
就是說.
你不可以在一種.
我們用了一個世界觀.
或者我們在地上.
我們覺得作為一個男人.
應該這樣.
你要搞清楚.
在上帝的設立裡.
祂不是這樣期望.
你這樣活出一個男人.
所謂靠自己.
所以說彼此相顧.
在我們的屬靈群體裡.
其實不單單是顧人.
都被人家顧你.
那個是雙向的.

$^{681}$當然一段的關係.
能夠去建立.
你需要花時間.
花心力.
這個是千真萬確.
沒有不勞而獲的.
不會突然之間的.
但是你要搞清楚.
或者你要記住一件事.
在上帝的創造裡.
每一個人.
你都需要關係.
你不要只是想著.
一個你自己.
第二.
有同行者的不是軟弱.
而是智慧.
傳道者所說的野性.
設計我們彼此相助.
不是因為我們軟弱.
而是它很帶出一件事.
因為人是有限的.
你需要別人的補足.
真正的力量.
是在於你承認你自己的需要.
而不是假裝著無所不能.
今天我們確實在面對著.
整個的社會.
堅強.
獨立.
千萬不要在職場裡.
讓別人看到你的情感.
我聽很多人說.
在職場裡.
不能夠表達內心的掙扎.
害怕也不能讓別人知道.
可能要表露出一種.
很兇的眼神.
硬的眼神.
叫別人不要欺負你.

$^{721}$無論在職場.
甚至在我們整個香港裡.
我們會不會不經不覺.
我們都進入了一種.
我們好像在靠自己.
對於承認自己的需要.
我們變得會覺得.
好像不可以擺在一起.
所羅門說到.
我們面對著兩個人.
總比一個人好.
兩人努力.
同得美好的果效.
確實是.
有同行者.
那不是一個弱者的表現.
卻是帶來真正的智慧.
所以所羅門不斷強調.
在我們這個短暫人生裡.
透過你不止一個人.
來更深去體會.
這個永恆價值的一部分.
你如何能夠在這個信仰群體.
一起經歷那份好處.
第三.
耶穌示範了同行的重要性.
如果仔細想想.
我今天整出來的經文不是很多.
特別是《約翰福音》十五章十五節.
「以後我不再稱你們為僕人.
因僕人不知道主人所做的事.
我乃稱你們為朋友.
因我從我父所聽見的.
已經都告訴你們」.
耶穌在自己選擇.
在走上十字架路上.
在這短短三十多年的人生裡.
確實他都展現了與人同行的榜樣.
他與他的門徒同行.
即使面對即將去丁十字架的路程.

$^{761}$在喀西瑪利園.
他都同樣面對極大的痛苦.
他不是自己走上山.
他都叫著跟隨著他的門徒.
陪他一起祈禱.
如果耶穌都選擇不是獨自.
面對人生最艱難的時刻.
你和我又怎能夠單打獨鬥呢?.
說這麼多.
作為一個同行者.
應該有一些什麼特質.
或者說素質.
我覺得這幾樣東西是重要的.
當你心裡覺得.
如果這一刻你聽到這裡.
你有同行者.
我為你感恩.
但同時可能有些人就覺得.
我未能夠找到這樣的同行者的時候.
你就去看一下.
你究竟怎樣.
不要說你找不找到這個同行者.
你能不能夠成為那個同行者.
有這些特質在你生命裡.
第一.
坦誠相待.
按著剛纔所講.
我們會不會停留於表面的裡面.
真正的同行者.
是幫助我們變得.
你可以更真誠.
坦誠地去展現那個你.
你不需要假裝.
更加重要的一件事.
就是過程中間能夠彼此保衛.
變得可以令你更加成熟.
甚至在當中所講的一些真誠的說話.
是叫你的生命得著更大的益處.
這個是真正.
我說同行者的時候.

$^{801}$我們應該要有那種特質.
好像箴言講一講.
鐵磨鐵磨出印來.
朋友相感也是如此.
能不能夠彼此坦誠地相待.
能不能夠有一個真誠的互信.
這個是應該作為同行者.
我們要有一些這樣的特質.
第二,分擔重擔.
我覺得這個很重要.
不是單單你只是.
我們的happy happy.
經常講一些好的東西.
開心的就經常講一些開心的.
我們能不能夠共度患難呢.
在你很艱難的時候.
他會不會能夠.
就像剛纔所羅門所形容.
可以扶你一把呢.
不只是在順境裡面.
能夠跟你分享那份喜樂.
能夠在你走入低谷的時候.
可以跟你一起哭.
可以跟你一起度過.
互相分擔.
第三,同行者這個素質更重要.
共同奔跑天路.
什麼意思呢.
就是說在這個同行者的關係裡面.
因著大家以上帝為中心.
上帝就是能夠在我們的關係裡面.
不單單教正我們.
互相扶持走這條.
繼續跟從上帝的路.
這個纔是最重要.
所以找的同行者.
不只是開心跟你一起去旅行.
找的同行者.
不只是有人跟你說一些開心的事.
找的同行者.

$^{841}$不只是跟你說一些你喜歡說的東西.
你有沒有這些素質呢.
你能不能夠坦誠相待.
你能不能夠分擔重擔.
你能不能夠在這個信仰的歷程裡面.
互相扶持.
以上帝為中心.
一同繼續堅持跟從上帝的一步.
好了.
我呢.
我跟你說這麼多.
我當然要說一下我自己.
誰是你生命中的同行者.
我在08年神學院畢業.
我其實在2023年左右.
我就認識了這兩個姐妹.
這兩個姐妹其實不是我.
我不是同福大.
這兩個姐妹也不是同福式.
我在外面的一些聚會.
一些祈禱的聚會認識.
穿針引線之下.
我們就有一個祈禱小組.
維持了18年.
由2003年開始.
真的好恐怖.
20幾年了.
2003年但我沒說錯.
2003年我認識她.
這張照片是最近.
2026年的2月份.
我很多年.
為什麼這麼說呢.
因為其實由19年開始.
我們就少了見面.
因為疫情.
少了很多大家聚在一起.
在過去的前部分.
所形容那18年裡面.
我很嚮往那種同行者觀.

$^{881}$當然到今天.
雖然我們很不常見.
不過我們走在一起的時候.
我很感恩一件事.
我們能夠.
我們不只是談笑風生地.
說說大家最近見到的趣事.
我們最喜歡說的就是.
你面對著什麼難處.
我們怎樣可以互相代給.
我這個小組其實是一個代禱小組.
因為我是一個祈禱的聚會認識他們.
那時候最瘋癲的時候.
我們是每晚都祈禱.
每晚約定10點半.
我2003年的時候.
就算在神學院宿舍的時候.
我都是會這樣.
跟他們一起這樣祈禱.
其中一個還去了.
回美國讀書.
我們就直接要.
要不就變成去到很晚.
因為他要起床.
我們等他起床.
要不就是我們要很早起床.
他還沒睡前.
是一種這樣的關係.
我覺得最寶貴的一件事.
當我看著這兩位姐妹的時候.
我想和大家說.
我認識她們的時候.
開始其實是很格格不入.
我們有少少就是.
那就做吧.
鼓勵我們祈禱.
我們就祈禱.
其實起初的時候.
我們完全沒有任何的.
未必是很深交應該說.

$^{921}$祈禱而已.
沒什麼深交.
但隨著日子.
當然慢慢的.
認識更多.
甚至乎就是.
在面對一些很艱難的時候.
我們是拿出來說.
拿出來分擔.
有幾個時刻我很記得.
我其中一個的.
我的祈禱夥伴裡面.
其實她是.
她的手腳是有點殘缺的.
意思就是.
其實她不見全隻手指.
有一年她爸爸進了醫院.
我很記得那個情景.
她就打電話.
叫我們一起去探望她.
去探望她爸爸.
接著我這個姐妹.
去到她爸爸的床邊的時候.
那個姑娘就說.
你們有空.
那麼多人在.
幫她剪剪腳甲吧.
幫她剪剪腳甲.
當時我這個姐妹.
因為她的手有點缺陷.
她是拿不起指甲剪的.
還要幫爸爸剪.
原來老人家的指甲很硬.
是很難剪的.
即是男人的指甲.
不知道你們有沒有試過.
因人而異.
是否很難大力剪下去.
你怕弄丟了.
我姐妹就拿著剪刀.

$^{961}$她就想幫她爸爸.
之後我就做了一個動作.
我說你不要來.
讓我來.
我很記得那個時刻.
其實這個是一個很大的.
讓我們彼此之間的關係.
好再進一步.
即是面對那個艱難的時候.
我不是袖手旁觀.
我不是只是在看.
你爸爸進了醫院.
我們去跟他一起去走.
經歷那種艱難.
其實不只是我幫他.
就算是我媽媽信主的時候.
其中我這兩個姐妹的時候.
是來做現場的見證人.
看著我媽媽決志信耶穌.
我覺得那個位置是很難得.
在生命當中能夠遇到一些.
這樣跟你不單單是交道心.
是真是在你的.
屬靈的生命裡面.
同行的同行者.
為你打氣.
為你扶持.
跟你一起同行.
最近見面的時候.
很難得.
還可以有這樣的關係.
還可以有一個.
那麼難得.
大家是可以.
沒有任何的假裝.
去講自己的.
現在這個可能我們面對的艱難.
我們的挑戰.
所遇見的人.
同事帶下來.

$^{1001}$對他心裡面的扎心的位置.
誰是你生命當中的同行者.
這裡有幾點.
都很想成為大家的鼓勵.
在你的朋輩當中.
我覺得這個是較為容易.
去找到你的志同道合.
不過同一時間.
要記住一件事.
我們每個人都要主動多走一步.
你不要想著等人來接近你.
你有沒有能夠去走前一步.
去建立關係.
學習闖開你自己的心.
我最近有跟一個少年人說.
那個少年人應該比你們小很多.
初初出來工作.
他就很真實地跟我說.
你經常叫我們打開我們的心.
其實我不知道怎麼打開.
可不可以教教我.
怎樣才叫打開我們的心.
我就跟他說.
打開我們的心.
其實我沒有辦法.
我不是跟你做過手術.
直接打開.
不過我有多願意.
讓人進入在我的生命裡面.
你面對一些你最覺得.
你不會跟別人說的事情.
你有沒有想過.
當你願意拿出來.
找到這些志同道合.
建立一定的關係.
能夠分享,分擔的時候.
你知不知道是很強大的.
不需要想著多大的事情.
由小事開始說.
你都能夠成為別人的祝福.

$^{1041}$可不可以請建伯早上來.
\newpage



\section{使徒行傳 16:6-15}
\label{sec:031827}
\textbf{甚麼是你的人生異象?}~~~~關鄧婥寧傳道
\newline
\newline
連結: \href{http://www.church.com.hk/acms/content.asp?site=cdc\&op=show\&type=product\&code=031827\&layout=sermon}{\texttt{031827}} ~~~~ 語音日期: 2026-03-15
\newline
\newline
\hyperref[sec:031826]{\small{< < < PREV SERMON < < <}}
~
\hyperref[sec:index_chronic]{\small{[返順時目]}}
~
\hyperref[sec:031828]{\small{> > > NEXT SERMON > > >}}
\newline
\newline
使徒行傳 16:6-15
\newline
\begin{longtable}{cl}
\hline
\hline
章節 & 經文 (和合本修訂版)\\
\hline
16:6 & \begin{tabularx}{0.7\textwidth}{X} 因為聖靈禁止他們在亞細亞講道，他們就經過弗呂家、加拉太一帶地方。 \end{tabularx} \\ \\ \relax
16:7 & \begin{tabularx}{0.7\textwidth}{X} 到了每西亞的邊界，他們想要往庇推尼去，耶穌的靈卻不許。 \end{tabularx} \\ \\ \relax
16:8 & \begin{tabularx}{0.7\textwidth}{X} 他們就越過每西亞，下特羅亞去。 \end{tabularx} \\ \\ \relax
16:9 & \begin{tabularx}{0.7\textwidth}{X} 夜間，有異象向保羅顯現。有一個馬其頓人站著求他說：「請你過來，到馬其頓來幫助我們！」 \end{tabularx} \\ \\ \relax
16:10 & \begin{tabularx}{0.7\textwidth}{X} 保羅既看見這異象，我們就立即設法往馬其頓去，認為神呼召我們傳福音給那裡的人。 \end{tabularx} \\ \\ \relax
16:11 & \begin{tabularx}{0.7\textwidth}{X} 我們從特羅亞開船，直行駛到撒摩特喇，第二天到了尼亞坡里； \end{tabularx} \\ \\ \relax
16:12 & \begin{tabularx}{0.7\textwidth}{X} 從那裡來到腓立比，就是馬其頓這一帶的一個重要城市，也是羅馬的駐防城。我們在這城裡住了幾天。 \end{tabularx} \\ \\ \relax
16:13 & \begin{tabularx}{0.7\textwidth}{X} 在安息日，我們出城門，到了河邊，知道那裡有一個禱告的地方，我們就坐下來對那些聚會的婦女講道。 \end{tabularx} \\ \\ \relax
16:14 & \begin{tabularx}{0.7\textwidth}{X} 有一個賣紫色布的婦人，名叫呂底亞，是推雅推喇城的人，素來敬拜神。她在聽著，主就開導她的心，使她留心聽保羅所講的話。 \end{tabularx} \\ \\ \relax
16:15 & \begin{tabularx}{0.7\textwidth}{X} 她和她一家都領了洗，就求我們說：「你們若以為我是真心信主的，請到我家裡來住。」於是她堅決請我們留下。 \end{tabularx} \\ \\
[1ex]
\hline
\hline
\end{longtable}
$^{1}$今日我講題叫「什麼是你的人生意象」.
今日的經文是來自《士徒行傳》十六章六至十五節.
是講到保羅如何領受一個馬其頓的意象.
事不宜遲 我們讀讀這段經文的頭五節.
我們一起讀 預備 1 2 3.
因為聖靈禁止他們在亞細亞講道.
他們就經過佛雷加·加拉泰一帶地方.
到了美西亞的邊界.
他們想要往彼推尼去.
耶穌的靈卻不許.
他們就越過美西亞夏特羅亞去.
夜間有意象向保羅顯現.
有一個馬其頓人站著求他說.
請你過來 到馬其頓來幫助我們.
保羅既看見這意象.
我們就立即設法往馬其頓去.
認為上帝呼召我們傳福音給那裡的人.
這裡就講到保羅領受了一個意象.
叫做馬其頓的意象.
有些人叫做馬其頓的呼聲.
保羅是怎樣領受意象的呢?.
其實這段經文講道.
神是關了保羅一道門.
並且又開了一道門給他.
當保羅領受了神給他的意象的時候.
他就馴服回應.
這個就是第一方面我們會講的.
我們看看剛才那段經文.
講到聖靈禁止保羅到亞細亞講道.
當保羅想在亞細亞繼續傳道的時候.
聖靈的靈是不許的.
我們看看地圖.
這是保羅第二次宣教之旅.
如果你看看地圖.
右邊就是亞細亞的地方.
也是今天的土耳其的地方.
愛琴海的左邊就是歐洲的地方.
你看到一些藍色的字是代表一個區域.
而紅色的字就是一個城市.
剛才那段經文講道.

$^{41}$保羅經過忽內加.
中間的地方.
他們又經過加拉泰.
又去到美西亞.
其實在美西亞想去比推尼.
但是聖靈不容許他去那裡.
在美西亞有一個城市叫做特羅亞.
已經去到亞細亞的邊境.
即是邊邊的地方.
在那裡保羅領受了一個馬其頓的意象.
他見到一個馬其頓人.
跟他說:你過來幫我們吧.
保羅領受了馬其頓的呼聲.
其實保羅由亞細亞去到歐洲那邊.
這個舉動是非常重要.
也是舉足輕重.
因為標誌著福音從亞細亞和亞洲那邊.
傳到歐洲那邊.
本來福音是從猶太文化.
影響到當時的希臘和羅馬文化.
也奠定了後來基督教.
大家都知道.
很多宣教士都是從歐洲那邊.
去到世界各地宣教.
所以保羅這個舉動是非常重要.
所以意象是非常重要.
也是神給我們的生命的目標和方向.
在《箴言》也有說.
29章18節說到.
沒有意象民就放肆.
為遵守律法的便為有福.
你有沒有人生意象?.
可以慢慢尋找.
我相信每一間教會.
每一個機構或每一個公司.
都有一個意象.
我們今天說很想開一間公司.
開一間咖啡店.
你也有些人有些目標.
你不會盲頭烏蠅開一間店.

$^{81}$起碼也賺回自己的工資.
我們有一個目標.
不知道大家知不知道.
我們代表柴灣堂.
知不知道為什麼會有柴灣堂.
也知不知道為什麼你今天坐在摩頓台這個地方.
是很有關係的.
和大家說一下歷史.
我想問一下.
誰是2009年已經在同福堂的?.
有呀.
也有些弟兄姊妹.
不多.
要說多一點.
同福堂是1991年成立的.
我自己就是1997年回到同福.
我小時候就回到同福.
轉一轉就快30年了.
當我回到同福的時候.
我們是在天后景香樓聚會.
不知道大家知不知道景香樓在哪裡.
大家知不知道百樂商業大廈在哪裡.
那裡有很多人都知道.
在天后廟道口那裡.
以前有一間衛康樓.
現在改名為Marketplace.
以前我們也有聚會點在百樂商業大廈.
旁邊就是景香樓.
當我回到同福的時候.
是1997年.
其實1997,98年.
聖靈在我們教會裡做了很多更新,醫治的工作.
很多我們的長輩,牧師都經歷了聖靈的更新.
那時候神就復興我們的祈禱會.
我記得我初初回到祈禱會是幾十人.
我們就很溫暖的圍一個圈.
大家有沒有什麼分享就祈禱了.
神開始復興我們的祈禱會.
會無端端有一些未信主的朋友.
會上來我們祈禱會.

$^{121}$然後他就信了耶穌.
是我們想不到的.
我因為那時候.
剛才Liz介紹我是第一代的詩琴.
我就是因為要做詩琴就回來祈禱會.
所以一路回來覺得祈禱會很好.
給大家看一下相片.
這個就是景香樓.
不知道誰有去過.
那個時候我們的台長長的.
這個景象就是我們最後一次在景香樓的祈禱會.
你看到我們都是圍一個圈.
都是很溫馨的.
大家都認識了.
神開始不斷給我們很喜歡祈禱的心.
我們不會錯過任何可以祈禱的機會.
我記得星期日只有兩堂崇拜.
早堂和午堂.
中間我們一班弟兄姊妹.
擠在一個很窄的.
應該是童宮的辦公室.
坐在桌子上就祈禱了.
因為沒有位子坐.
我們很喜歡祈禱.
那個時候神就不斷帶很多人來到我們當中.
很多人信主.
天天都給我們有得救的人數.
開始1998年我們的長輩就想.
會不會我們買一個堂子呢.
因為景香樓都是租的.
還有他們在電器道有一個辦公室都是租的.
他們就在想會不會買一個堂子呢.
一直在祈禱.
在1998年的6月4日.
有一個弟兄就領受了一個異象.
他就和長輩分享.
他就說我在東區聽到.
當他祈禱的時候神就這樣跟他說.
他說我在東區聽到妓女的哭聲.
和失去父親兒童的哭聲.

$^{161}$他說為此我要在東區建立我的教會.
你們要從香港的富豪酒店.
到柴灣地鐵站畫一條線.
在方圓三里之內所有拜偶像的地方.
從偶像面對的地方畫出一條直線.
這些線聚集的地方就是我的電要所建立之處.
這個教會是要幫貧窮人拯救他們.
所以我的子民要禁食祈禱.
你們要分三隊.
每隊七個人.
輪流每餐禁食七乘七乘七日.
因為支離的軟弱在禁食的時候.
可以吃白麵包喝清水或者羊奶.
當你們禁食完成的時候.
事就成了.
我們領受了這個議長.
我們時常叫東區議長.
開始拿地圖出來畫.
由富豪酒店畫到柴灣地鐵站.
剛剛好就是三里.
大家就去找那些偶像面向的地方畫.
畫那些線.
最終發現線聚集在兩個地方.
第一個地方就是維港的海上.
這是不可能的.
第二個地方就是今天大家坐在這裡的地方.
接著我們繼續有祈禱會.
弟兄姊妹有領受四個英文字.
叫做ISIS.
是什麼來的呢?.
當時這個地方是一個戲院.
叫做新都戲院.
原來ISIS這四個英文字就是新都戲院的英文名.
也有弟兄姊妹領受.
看到天橋.
看到投影機.
看到一些南亞裔人.
其實都引證了之前神給我們的異象.
我們就開始去接觸這個新都戲院的業主.
當時他是賣九千萬.

$^{201}$對我們來說是一個天文數字.
經過很多的磋商.
初時說賣半間又賣一間等等.
接著最終我們是以四千二百五十萬賣了整個戲院.
今天你可能會說很划算.
但是當年這個一九九八一九九九.
其實我們整個教會的時候只有五百多人.
這個對我們來說是一個天文數字.
而在我們的獎賞去博取這個地方的時候.
那天就是七乘七乘七的最後一天.
哇 神是不是很神奇.
這個地方就是這樣來的.
所以和東區異象很有關係.
又給大家看一張照片.
這個就是我們買了這個地方.
我們裝修前的時候.
我們來過這裡有一個祈禱會.
誰在這裡.
沒有什麼人在這裡.
那天你想一下這個地方沒有窗.
當時當然沒有冷氣.
還沒有裝修.
只有這些光管.
熱得不得了.
但是我們很雀躍在這裡祈禱.
而且我們拿一些筆.
其實你們現在坐在地上 牆 桌子 各樣東西.
我們寫了很多經文在當中.
真的很開心.
當我們買到這個地方的時候.
也不斷地為摩頓台這個地方去很多的禱告.
我們二零零一年就搬進來這個地方.
終於可以搬進來.
我們搬進來這個地方.
但是我們沒有忘記神給我們的東區遺像.
所以雖然我們在摩頓台.
但是我們時常都去柴灣 小西灣 修基灣等等.
我們做很多大型的嘉年華會.
那時候我們可以去洗樓.
大家知不知道什麼是洗樓.

$^{241}$洗樓就是我們進去公屋.
那時候沒有閘.
現在就全部下閘.
我們真的逐家逐戶拍門.
送禮物給他們 跟他們講福音.
那時候是很開心的.
我們一次過幾百人我們就去洗樓.
那個年代真的很開心.
而且那個年代.
我們每逢十二月廿四日平安夜.
我們就會去會展有一個子夜崇拜.
有一年我們就搞子夜崇拜.
我們說不要浪費場地這麼貴租回來.
不如在下午搞兩場報道會.
我不知道為什麼那時候有這麼多能量.
這個就是我們有一年.
我們那時候是租那個Grand Hall.
就是對著整個維港景.
也是香港回歸的禮堂.
我們下午就做一些嘉年華會.
還有子夜崇拜.
現在做一次這個聖誕崇拜.
好像沒有半條人命.
那時候不知道為什麼有這麼多精力.
可以搞很多.
不過那時候真的這麼開心.
我們就會拿一些巴士.
接很多東區居民去到會展.
讓他們在我們當中去聽福音.
其實我們一直去.
舉辦了很多嘉年華會.
很多人信主.
他們信了耶穌之後.
我們每星期就租一些巴士.
載他們回來摩頓台崇拜.
小朋友就會在現在Crossway那個地方.
以前是一個地鋪.
我們就會在那裡教小朋友聖經.
跟他們玩.
每個星期就有巴士長載他們回來.

$^{281}$而且有些人信了主.
我們有些同工就去他們的家庭.
有些婦女的家庭.
有些茶經班.
那個年代就是這樣.
久而久之.
我們發覺這樣載來載去都不是辦法.
我們都需要在東區有一個據點.
所以在2009年.
我們在6月的時候.
我們就設立了一個資源中心.
在柴灣那裡.
同年12月半年後.
我們也在那裡開始柴灣堂的崇拜.
當時鄭牧師是柴灣堂的堂主任.
2018年因為何牧師要退休.
他就回來做了大家的主任牧師.
所以鄭牧師兩個月內就叫我.
我要回來摩頓台了.
你搞定柴灣堂吧.
臨危受命就要接手柴灣堂.
其實當年神給了我們東區意象的時候.
我們是立即回應的.
也很開心地神給了我們的意象.
也神給了我們摩頓台這個地方.
剛才說到保羅.
保羅收了神給他的馬其頓意象的時候.
他也是立即回應的.
他並且立即起行.
坐船去彭遜起行.
也不斷去求證.
我們看接下來那幾節經文.
10節至12節.
我們一起讀一讀.
1,2,3.
既看見這意象.
我們就立即設法往馬其頓去.
認為上帝呼召我們全呼音給那裡的人.
我們從特羅亞開船.
直行駛到薩摩特拉.

$^{321}$第二天到了阿尼波利.
從那裡來到菲拉比.
就是馬其頓這一帶的一個重要城市.
也是羅馬的駐防城.
我們在這城裡住了幾天.
這裡說到保羅.
領受了他意象就起行了.
坐船起行了.
我們看看地圖.
他從特羅亞.
因為那邊是海.
他就起行了.
去到薩摩特拉.
一個小島.
他住了一晚又起行了.
去到尼亞波利.
已經去到馬其頓那邊了.
他就去到菲拉比這個地方.
菲拉比是一個很重要的城市.
是當時的羅馬殖民地.
也有很多退伍軍人住在那裡.
也是羅馬的駐防城.
其實保羅在這裡宣教.
影響了當時羅馬帝國的文化.
保羅是不是很衝動.
他領受了這個意象就立刻起行了.
你覺得他衝不衝動.
我相信他有考慮過.
因為神已經關了一扇門.
不讓他再在亞細亞宣講福音.
但是神又開了一扇門.
他就立即去.
而且保羅是有團隊的.
他不是自己一個.
他第二次宣教是帶著提摩泰.
路加等等.
我相信他的團隊.
都會一起去印證這件事.
他才開始起行.
我不知道電影節目有沒有試過.

$^{361}$有時候都會在神裡面有些領受.
領受我做什麼.
我們會不會找人一起去印證一下.
你會找些什麼人.
會不會找些跟你同聲同氣的人.
我覺得是這樣走的.
找個朋友聊天.
還是你會找一些.
跟你看法可能有些不同的.
當然是一些屬靈的長輩.
你的組長.
你的區域.
一些同工.
有時候我們都需要找一些.
跟我們有些不同的看法的人.
可能意見會更加中肯.
說回柴灣堂.
其實我接手柴灣堂之後.
都面對很多挑戰.
我2018年做了柴灣堂.
當時是代堂主任.
2018年都挺平穩的.
開始2019.
大家知道發生什麼事.
社運.
2020.
接著就疫情.
其實是非常多挑戰.
因為這段時間.
其實有很多弟兄姊妹.
舉家移民.
你想到移民的家庭.
通常都是比較中產.
有小朋友.
他們很多都是我們的中流砥柱.
是我們的組長.
很多都是我們從銅鑼灣堂下來的弟兄姊妹.
其實令到我們整個人口結構都有點不同.
復常之後.
有些人真的懶得就懶.

$^{401}$就會不回教會.
我有時候在柴灣都會遇到他們.
我說很久沒見了.
你去了哪裡.
為什麼不回來.
有些很直接跟我說.
懶.
可能我身邊也有這樣的人.
真的.
懶得就懶.
怪不得神叫我們不可停止聚會.
停止慣了.
我們就不想再繼續聚會.
所以都有點困難.
很多挑戰.
還有我們在柴灣.
我們其實有三個單位.
我們每幾年就要續租.
這三個單位.
每次續租的時候.
我就會想.
我們是不是應該繼續續租呢.
我們是不是繼續留在柴灣呢.
因為我們有一個資源中心.
在我們柴灣堂.
所以我們是有食物銀行的.
有很多社工會轉移個案給我們.
其實有很多個案.
都是在修基灣,西灣河那邊.
因為那邊有很多tong 房.
他們會轉移這些個案給我們.
還有神又讓我們在修基灣那邊.
可以服侍一個長者日間中心.
我們都會想.
會不會神叫我們移去那邊呢.
會不會搬去那邊服侍那邊的居民呢.
其實我們一直都在祈禱.
都不知道神有沒有其他的心意.
直到有一晚.
是我們祈禱會的晚上.

$^{441}$祈禱會結束之後.
我就進了洗手間.
可能去了久一點.
我出來的時候.
發覺所有人都走了.
連燈都關了.
我就摸黑進了辦公室.
我很少真的把燈關在辦公室裡.
當我進到辦公室的時候.
我看到這個景象.
漂亮嗎.
這是我很感恩的.
我們看到整個柴灣的景象.
當我看到這個景象的時候.
彷彿神跟我說.
你見不見到.
這裡有很多的家庭.
他們未信耶穌.
你願不願意繼續去服侍他們.
我感受到神的心.
當天神叫我們去服侍一群有需要的朋友.
一些無腐的兒童.
其實神的心是沒有變過.
我看到的時候.
我的眼淚就流出來了.
我就拿出我的手機拍了這張照片.
所以你看到有點蒙蒙的.
但我覺得這是神跟我說的一個印證.
也是一個記號.
所以我拍了這張照片.
其實不是太多人會去柴灣.
小西灣這個地方.
除非你住在那裡.
或者你在那裡工作.
你最多都是清明和重陽才會去.
很少人會去這麼多國.
銅鑼灣你就會經常去.
但柴灣真的很少人會去那個地方.
但神沒有忘記.
當天祂給我們的東區遺像.

$^{481}$沒有忘記祂所愛的人.
所以我覺得神要我們繼續留在柴灣這個地方.
可能你聽我分享.
你說終於明白.
原來柴灣堂或者摩頓台這個地方.
是神有一個很清晰的遺像和大靈.
但說到我人生.
我真的不知道神有什麼特別的方向給我.
一個特別的使命給我.
其實如果你說沒有.
你就看聖經.
聖經其實說了.
聖經裡面有一個大誡命和大使命.
不知道弟兄姊妹還記不記得.
大誡命就是.
要盡心盡性盡義去愛我們的主我們的神.
並且要愛倫如己.
大使命就是我們要去洗萬民作的門徒.
奉父子聖靈的名給他們施洗.
還要將神所吩咐我們的都要教導他們遵守.
其實如果你的人生只是做這兩樣東西.
你已經有很長時間做了.
我們有沒有做到呢.
我們有沒有盡心盡義盡性去愛我們的主.
並且愛人如己.
我們有沒有好好將神的福音傳出去.
教導他們遵守神的命令呢.
我經常想,當我們在小事上忠心的時候.
神就會比我們有更大的事.
會啟示我們去更多的心意.
但我們有沒有在小事上忠心於神呢.
如果你說找不到.
或者你今天坐在同福堂.
我覺得不是一個偶然.
如果你今天坐在這個教會.
牧師帶領怎樣去做的時候.
你就跟著去做.
其實這些就是神給我們的意象和使命.
就是這麼簡單.
我們有沒有跟著去走呢.

$^{521}$這是第二方面.
第三方面,保羅在宣教裡面有策略.
他不斷尋求策略和看到突破.
我們繼續讀最後幾節經文.
13節至15節.
我們一起讀,預備,1,2,3.
在安息日,我們出城門,到了河邊.
知道那裡有一個禱告的地方.
我們就坐下來對那些聚會的婦女講道.
有一個賣紫色布的婦人,名叫呂底亞.
是推亞推拉城的人,素來敬拜上帝.
她在聽著,主就開導她的心.
使她留心聽保羅所講的話.
她和她一家都領了洗.
就求我們說.
你們若以為我是真心信主的.
請到我家裡來和住.
於是她堅決請我們留下.
這裡講到保羅去到菲臘比.
他就住了幾天.
他不是立即那天出去傳福音.
他去到安息日.
他就去了河邊.
因為河邊有幾個人祈禱.
其實保羅宣教有一個策略.
他通常去到一個新的城市.
他就去會堂.
在安息日那天,他就會去會堂.
因為會堂裡有一班猶太人會敬拜上帝.
菲臘比沒有會堂.
所以他知道在河邊有一班人聚集禱告.
所以他就去到那裡.
保羅去到河邊的時候.
他遇到一個人.
就是一個婦人叫呂底亞.
這個呂底亞也是一個很重要的人.
為何她那麼重要呢?.
第一,她是一個賣紫色布的商人.
你說:通街都是紫色布的,有甚麼特別呢?.
今時今日當然了.

$^{561}$但當時紫色布是非常之矜貴的.
為甚麼呢?.
因為紫色布要用一種紫灰堆裸的軟體動物.
好像一個裸一樣.
要找到這些裸.
然後拿去裡面紫色的部份去染這些布.
所以這些布是很矜貴,很難得.
而且不是一般人可以買到.
通常是達官貴人,甚至是皇室成員才可以買到.
所以當時紫色布比黃金更加矜貴.
所以呂底亞是一個很有錢的人.
而且她的屋子很大.
她叫保羅住在那裡.
其實她後來開放了她的屋子.
成為早期教會的聚會點.
所以呂底亞是非常之重要.
而且呂底亞是很有錢的.
其實她支持保羅的宣教之旅.
她在經濟上不斷成為一個很大的支持者.
所以這是呂底亞.
第二,她是一個外邦女子.
有什麼特別?.
你想想,保羅第一次到歐洲這個地方宣教.
第一個信主就是一個外邦人.
不是一個猶太人.
並且她是一個女子.
女性比男性的地位當時沒有那麼高.
所以其實也是標誌著.
神的福音是給所有人.
不是什麼階級.
是給不只是猶太人.
是給所有的外邦人.
什麼階級,無論你是有錢也好,貧窮也好.
神的福音就是這樣廣傳.
所以這個對保羅來說也是一個很重要的基石.
福音就是傳給所有的人.
並且她是一個敬拜上帝的人.
呂底亞之前也是敬拜上帝的.
不過她還沒有認識耶穌.
當她聽到保羅分享她的見證.

$^{601}$她的故事的時候.
聖靈就感動她.
讓她相信耶穌.
不只是她自己一個人相信.
她整家人都相信.
還立刻受洗.
你想想,如果今天你去短訓.
一出隊就很多人信主.
並且立刻受洗.
你覺得開心嗎?.
當然開心.
其實對保羅來說是一個很大的鼓舞.
她第一次踏入馬其頓的地方.
歐洲這邊就有人信主.
還不只是她一個人.
整家人都信主.
其實對她來說是一個很大的確據.
去馬其頓這條路是正確的.
所以日後保羅就繼續在歐洲傳道.
也影響了整個世界宣教的歷程.
所以保羅去到馬其頓這一步.
是非常重要的.
所以當保羅去宣教的時候.
他不斷去尋求不同的策略.
他不是盲頭烏蠅.
他有策略去傳福音.
而且神不斷去印證他今天所走的路是對的.
我們又說回柴灣堂.
我們去了柴灣之後.
當初我們跟銅鑼灣堂有很多合作.
我們有舉辦過很多嘉年華會.
可能有些弟兄姊妹也有參加過.
如果你在銅福堂來的話.
我們有舉辦過盆菜宴.
我們經常叫耶穌請吃飯.
其實當中有很多街坊信主.
是很開心的.
因為柴灣堂我們已經到了16年.
進入17年.
我們也在想.

$^{641}$是不是還是用我們第一天的策略.
那個方向去繼續去服侍主呢.
所以我們去年成立了一個核心團隊.
有一個call team.
我們再一次去審視.
去祈禱.
去看一看我們現在在做的事.
是否應該繼續做.
還是應該給我們一些新的發展.
各樣各樣.
所以我們也為柴灣堂.
訂立我們的意象宣言.
我們的牧養策略.
因為我們走了小祖制很多年.
我們也發現很多祖長也很累.
我們是不是繼續這樣走呢.
我們也有想了一些方案.
你為我們祈禱.
也挺新穎的.
希望我們成功.
我們也希望幫助弟兄姊妹成長.
我們一直在祈禱的時候.
神也讓我們看得更加清楚.
我們的合作夥伴.
可能有些是學校.
有些是一些機構.
神一直一直為我們開展.
也是很感恩的.
今年是農曆年.
農曆年我們通常很少出隊.
我們聖誕節就會去布街音.
為什麼農曆年不可以去布街音呢.
我們說不如試一下.
我們就買了一些很漂亮的農曆年掛飾.
我們還選了2月14號情人節那天.
去布街音.
如果你熟悉柴灣.
這個地方就是漁灣村.
我們還對著警察宿舍布街音.
我們還用咪高峰.

$^{681}$會被人罵.
很感恩我們在街坊面前.
眼前無恩.
街坊很開心見到我們.
我們有送禮物.
祈禱祝福他們.
我們感受到那份溫暖和喜悅.
原來我們這樣去做的時候.
神也很喜悅我們.
想起今年復活節.
我們說.
如果你看回那個日曆.
今年復活節是五日紅日.
很長的.
我就想.
其實我們弟兄姊妹服侍都很累.
可不可以讓他們休息一下.
不如我們星期五.
Good Friday那天.
我們就搞一次聯合崇拜.
同工都說好.
我們找一個大一點的地方.
在柴灣找一個大一點的地方.
都不容易的.
只有一個地方叫青年廣場.
我不知道大家有沒有去過.
但現在他們長期被學校.
預約了去做比賽.
如果你有小朋友.
很多時候那些校際音樂節.
就會在那裡進行.
他們長期都預約了那個地方.
我們說找什麼地方好呢.
突然有同工就說.
不如我們回摩頓台.
又不用錢.
我們說好啊.
就打電話給行政部.
問一下這個場那天有沒有人用.
他說沒有.

$^{721}$那好啊.
你先幫我hold住.
我在想.
不如搞一個報道會.
很久沒有搞過報道會.
我還在想.
找哪位弟兄姊妹分享見證好呢.
有一個同工就說.
不要.
找一個知名度高一點的人回來講見證.
弟兄姊妹可以容易一點.
去邀請新朋友.
我說找誰好呢.
大家就說了一大堆名字出來.
後來就想起一位創作歌手.
Kim Man.
不知道你認不認識.
王劍文.
我說不知道他行不行.
他們就大膽去問他.
他的時間又OK.
其實我也不知道他的見證是怎樣的.
過了兩天之後.
那個同工就經常進學校.
看一些中學生.
跟他們建立關係.
希望傳福音給他們.
他那天回到學校.
就說很特別.
老師叫他們看見證.
還要寫反詩.
還給了一條link給我看.
我真的沒聽過他的見證.
我真的很用心去看了一次他的見證.
我覺得他的見證很觸動我.
我就跟同工說.
主若願意就讓這件事成就.
過了幾天之後.
他就說OK.
所以事就成了.

$^{761}$所以我們在壽苦節那天.
4月3日.
我們柴灣堂會在這個地方.
做一個聯合崇拜和報道會.
其實這個位置.
我們應該坐不完的.
我都歡迎弟兄姊妹.
如果你說.
我都有一些未信主的親友.
很想邀請他們來相信耶穌.
很想去聽福音的.
我都鼓勵你.
你都可以拿一些票.
邀請他們回來.
如果你讀這個QR Code.
你就要過去柴灣堂拿票.
但我方便大家.
今天我有同工在下面.
可以給大家一些票.
你真的想一下.
你真的有一些福音對象.
你可以去拿一些票.
每人最多拿4張.
我們位置都有限.
我覺得真的很感恩.
在今年.
神給我們柴灣堂很多新嘗試.
還有東西去得很快.
其實我看見.
建文答應了我們之後.
其實他就要出發.
他現在其實在澳洲.
他要在澳洲短暫一個月.
接著他回來之後.
差不多是我們的報道會.
其實神真的預備了他在我們當中.
我覺得當我們原來.
我們一路去祈禱.
一路去尋求神的時候.
神就會一路一路開路.

$^{801}$給策略我們.
所以如果我們的生命.
你說我今天都不知道.
我前面的意象是怎樣.
我的方向又怎樣.
你祈禱吧.
你求問神.
神是很喜歡給我們啟示.
知道的神.
就算你說我都不清晰的時候.
你先做神叫我們做的大使命和大誡命.
我相信當我們在小事上忠心的時候.
神必定會擴張我們的境界.
讓我們更加明白.
神在我們的心意是怎樣.
讓我們一起祈禱吧.
是的,主.
主,你是一個很喜歡啟示我們的主.
很喜歡和我們一起同工的主.
主,讓你給我們有人生的意象.
有人生的目的.
我不知道在當中.
有沒有一些未信主的朋友.
你說,是的.
我也不知道人生的意義是怎樣.
但是今天我很想邀請你.
將你的生命交給主.
讓主給你一個生命的意義.
給你一個生命的方向.
讓你的人生有著神的盼望.
如果有這些朋友.
你說很想將你的生命今天交給主.
很想邀請你.
你舉起你的右手.
很想為你祈禱.
我不知道有沒有一些朋友.
有沒有這些朋友.
有這些朋友在當中.
我感恩,感恩.
很想邀請你.

$^{841}$待會我們唱回應師兄歌的時候.
來到我的右手邊.
很想為你祈禱.
我不知道有沒有弟兄姊妹.
我也很想發出另一個邀請.
你說,是的.
我也在尋覓神給我的意象.
我的方向.
或者你已經知道我為你感恩.
但是,是的.
我需要神繼續帶領我.
給我能力.
讓我有神的洞察力.
很想神繼續加力給我們.
很想為這些弟兄姊妹去禱告.
讓我們繼續領受神的意象.
神繼續跟我們說話.
繼續在當中神引領我們.
讓我們看到神怎樣一步一步帶領我們.
如果有這些弟兄姊妹.
很想邀請你們將你的手放在心上.
很想為大家祈禱.
求主加力給我們.
是的,主.
我們很需要一個意象.
如果我們生命沒有了一個意象.
我們就.
盲頭烏蠅.
我們不知道要走去哪裡.
主啊,你就是來帶領我們.
讓我們願意去尋找的時候.
神,你就開路.
讓我們找得到.
也是開我們的心.
讓我們知道神你怎樣引領我們.
我們都可以順服去走在你的心意當中.
主啊,你就是這樣.
加力給我們弟兄姊妹.
讓我們有一個安靜的心.
就是領受神你給我們的意象.

$^{881}$給我們敏銳聖靈的聲音.
主,你就是這樣去跟我們同在.
加力給我們.
讓你的力充滿我們.
主啊,多謝贊美你.
多謝主祈禱共處明教.
Amen.
\newpage



\section{雅各書 1:19-22}
\label{sec:031828}
\textbf{有入有出的執道人 (粵/華)}~~~~冼嘉琪傳道
\newline
\newline
連結: \href{http://www.church.com.hk/acms/content.asp?site=cdc\&op=show\&type=product\&code=031828\&layout=sermon}{\texttt{031828}} ~~~~ 語音日期: 2026-03-15
\newline
\newline
\hyperref[sec:031827]{\small{< < < PREV SERMON < < <}}
~
\hyperref[sec:index_chronic]{\small{[返順時目]}}
~
\hyperref[sec:031829]{\small{> > > NEXT SERMON > > >}}
\newline
\newline
雅各書 1:19-22
\newline
\begin{longtable}{cl}
\hline
\hline
章節 & 經文 (和合本修訂版)\\
\hline
1:19 & \begin{tabularx}{0.7\textwidth}{X} 我親愛的弟兄們，你們要明白：你們每一個人要快快地聽，慢慢地說，慢慢地動怒， \end{tabularx} \\ \\ \relax
1:20 & \begin{tabularx}{0.7\textwidth}{X} 因為人的怒氣並不能實現神的義。 \end{tabularx} \\ \\ \relax
1:21 & \begin{tabularx}{0.7\textwidth}{X} 所以，你們要除去一切的污穢和累積的惡毒，要存溫柔的心領受所栽種的道，就是能救你們靈魂的道。 \end{tabularx} \\ \\ \relax
1:22 & \begin{tabularx}{0.7\textwidth}{X} 但是，你們要作行道的人，不要只作聽道的人，自己欺騙自己。 \end{tabularx} \\ \\
[1ex]
\hline
\hline
\end{longtable}
$^{1}$今天我想給大家看看.
大家有沒有用手機app去靈修?.
大家有用手機一些app靈修嗎?.
有上網聽到嗎?.
有用Facebook看一些靈修的message嗎?.
沒有啊,有些有,有些沒有.
我們發現今天的世代是一個很豐富的世代.
現在的世代是一個很豐富的世代.
因為我們打開很多電子用品的時候.
我們很容易聽到神的話語.
因為我們打開一些電子用品的時候.
很容易聽到神的話語.
如果我們去形容的話.
可能這個世代是我們最容易聽到神話語的世代.
如果我們形容一下的話.
這個世代是最容易聽到神話語的世代.
但是聽得多是不是真的靈修得多呢?.
是不是真的吸收得多呢?.
我們的屬靈生命像一個garage.
弟兄姊妹,我們的garage是怎樣的?.
我們的屬靈生命就像一個車庫一樣.
大家的車庫怎麼樣?.
我們的garage有時是塞滿了很多東西在裡面.
有時候我們的車庫裡塞滿了東西.
很多箱,很多新的舊的東西塞進去.
新的舊的東西塞在裡面.
但當我們要找東西的時候呢?.
當我們要找東西的時候.
可能怎麼找都找不到,是不是?.
找都找不到.
找了一整天,找個螺絲起子都找不到.
可能找一個螺絲起子都找不到.
我們很多時我們的生命都是這樣.
有時候我們的生命就是像這樣.
塞滿了很多新的舊的東西.
塞了很多新的舊的東西.
我們生命裡面沒有去整理.
沒有去梳理一些的事情.
我們生命沒有去整理和梳理這些事情.
甚至乎聽到神的話的時候.

$^{41}$我們沒有去實踐.
甚至乎我們聽到神話的時候.
我們都沒有去很好的整理.
就好像Garage那些.
很多很多箱在裡面.
就像這個Garage裡面有很多的箱子.
神的話語在我們裡面,但是用不著.
神的話語在我們裡面,可是用不到.
所以弟兄姊妹.
今天我們一起去聽神的話語.
所以今天我們一起來聽神的話語.
我們一起學要做一個有入有出的基督徒.
我們來學習做一個有入有出的基督徒.
我們今天一起來看雅各書一章.
我們今天來看雅各書.
我們一起來讀一下經文好不好.
四節經文.
1,2,3.
雅各書.
大家可以用國語也行.
用粵語也行.
用英文也行.
雅各書一章.
19至22節.
一起.
我親愛的弟兄們.
你們要明白.
你們每一個人要快快地聽.
慢慢地說.
慢慢地動腦.
因為人的怒氣.
並不能實現神的意.
所以你們要除去一切污穢和累積的惡毒.
要傳溫柔的心領受所栽種的道.
就是能救你們靈魂的道.
但是你們要做行道的人.
不要只做聽道的人.
自己欺騙自己.
很熟很熟的經文.
熟悉的經文是嗎.

$^{81}$雅各說你們要做行道的人.
不要只做聽道的人.
雅各說你們要做行道的人.
不要只做聽道的人.
就是說我們一邊接收的時候.
一邊要走出來.
就是說我們一邊接收的時候.
一邊要醒出來.
行道的起點.
是要懂得好好聽道.
行道的怎麼樣.
就要好好的聽道.
雅各說.
你們每一個人要快快地聽.
雅各說你們每一個人要快快地聽.
快快地聽.
為什麼呢.
原來生命改變的第一步.
是要我們領受神話語的態度正確.
原來生命的第一步.
要我們領受神的話語態度正確.
快快地聽.
雖然聽後面沒有授詞.
快快地聽.
雖然聽後面沒有授詞兵語.
就是說聽什麼.
就是說聽什麼呢.
但如果我們看回上文和下文.
上一節就講到真理的道.
我們看回聖經的上一節.
它就講到真理的道.
下文就是講到.
我們接受了這個真理的道之後.
怎樣在心裡面栽種這個真道.
再看下一節.
我們就講到一個真理最重要的一點.
所以當我們看到今天十九節的時候.
快快地聽.
我們就知道原來講的是.
聽神的道.

$^{121}$我們看到十九節的時候.
就看到快快地聽.
是要聽神的道.
可能看到這裡的時候.
快快地聽.
怎麼聽呢.
快快地聽.
要怎麼聽呢.
因為講就可以快.
因為講可以快.
跑可以快.
聽怎麼快呢.
聽怎麼能快呢.
聽是對方說話快.
我就聽得快.
聽是對方說得快.
我就聽得快.
對方說話慢.
我聽不到快.
對方說得慢.
我就聽不了快.
那怎麼快快地聽.
那怎麼快快地聽呢.
所以雅各這裡講的.
不是一個速度的快.
雅各這裡說的.
不是一個速度的快.
快快地聽是講對神話語的一個反應.
快快地聽是對神話語的一個反應.
預備聽神的聲音.
樂意很心急的.
很渴慕聽神的聲音.
非常渴慕心急的想聽神的聲音.
有沒有試過等電話.
等一個很重要的電話是怎樣.
等一個很重要的電話會怎樣.
擺個電話在旁邊.
把電話擺在那.
望實有沒有調校成音.
看一下有沒有聲音.

$^{161}$電話一響馬上拿起來聽.
一響就馬上拿起來聽.
快快地聽.
神想我們都是這樣.
神希望我們是這樣.
隨時願意有受教的心.
隨時有受教的心.
願意隨時接受神的教導.
隨時on call回應.
隨時on call.
弟兄姊妹.
我們有沒有隨時on call回應神.
當我們收到神的道.
祂的教導的時候.
我們要全敬畏的心去回應.
我們聽到神話的時候要全敬畏的心去回應.
慢慢消化去咀嚼神的教導.
慢慢消化咀嚼神的教導.
所以後面那裡慢慢地說.
都不是說一個速度.
慢慢地說.
不是說一個速度.
不是.
你好嗎.
不是在說你好嗎.
慢慢地說是說.
我們接收到神的道之後.
不要急於回應.
慢慢地說是我們接受神的話語之後.
不是很快地做出回應.
慢慢咀嚼.
不要急於告訴別人.
不要急於教導批評.
慢慢地說.
我很感恩在兒童事工回復kit裡面.
有一班很有愛心,很忠心的弟兄姊妹.
對神忠心,不是對我們忠心.
這班老師有時分享的小動作.
這班老師分享的小動作.
一個簡單的分享.

$^{201}$都很感動我.
有一次有位老師說.
他不能很快地準備好每一次的課程計劃.
有一次有位老師說.
他不能很快地準備好每一次的課程計劃.
為甚麼呢?.
不是因為他效率慢.
不是因為他忙得太忙,沒時間準備.
不是他太忙沒時間準備.
是因為他每次看完課程計劃之後.
他要慢慢消化神的說話.
因為他每次準備好以後,他要慢慢消化神的話語.
他想深入一點神的話才教.
他要先理解吃透神的話語.
因為他知道神的話要消化,要吸收.
他知道神的話要消化,吸收.
在兒童事工裡面,就是一班這麼愛主的基督徒.
他們不只是教神的話語.
他們不是在教神的話語.
他們先去讓神的話語充滿他們的心.
他們先去敬畏神的話語.
他們先敬畏神的話語.
弟兄姊妹,我們聽到神的話語的時候就是要這樣.
懷著敬畏和尊重的心回應神.
弟兄姊妹,當我們聽到神的道的時候.
一般我們會心裡面消化了先.
還是我們會轉化了先.
我們再轉化它.
好句,forward下先.
好啊,教得好啊今天,轉化下先.
轉化先.
弟兄姊妹,我們不要急著share.
我們不要急著去share.
我們回應神先,好嗎?.
我們回應神先.
阿哥特別提到對神的道還要慢慢地動怒.
我們要聽到神的道要慢慢地動怒.
為什麼呢?為什麼聽到神的話會動怒呢?.
為什麼聽到神的話會動怒呢?.
我們有沒有試過聽到神的話動怒啊?.

\newpage



\section{路加福音 15:1-2}
\label{sec:031829}
\textbf{寶貴,保證,賬簿}~~~~梁誠德牧師
\newline
\newline
連結: \href{http://www.church.com.hk/acms/content.asp?site=cdc\&op=show\&type=product\&code=031829\&layout=sermon}{\texttt{031829}} ~~~~ 語音日期: 2026-03-15
\newline
\newline
\hyperref[sec:031828]{\small{< < < PREV SERMON < < <}}
~
\hyperref[sec:index_chronic]{\small{[返順時目]}}
~
\hyperref[sec:031830]{\small{> > > NEXT SERMON > > >}}
\newline
\newline
路加福音 15:1-2
\newline
\begin{longtable}{cl}
\hline
\hline
章節 & 經文 (和合本修訂版)\\
\hline
15:1 & \begin{tabularx}{0.7\textwidth}{X} 許多稅吏和罪人都挨近耶穌，要聽他講道。 \end{tabularx} \\ \\ \relax
15:2 & \begin{tabularx}{0.7\textwidth}{X} 法利賽人和文士私下議論說：「這個人接納罪人，又同他們吃飯。」 \end{tabularx} \\ \\
[1ex]
\hline
\hline
\end{longtable}
$^{1}$弟兄姊妹 早安.
 很高興再次回到神的家.
 跟大家分享神的話語.
 大家今天過得好不好?.
 為什麼這麼好?.
 有沒有過得不太好?.
 去加油的時候 蛤?.
 貴了兩元的.
 會不會這樣?.
 你會不會說 哎呀 貴了兩元.
 財散人安樂 會不會這樣?.
 沒有的.
 人的理性是這樣想的.
 今天是講錢的.
 你知道嗎?.
 剛才有位婦人.
 跌了一元雞.
 我們叫她一塊銅板.
 或者一塊錢.
 你還有九塊錢.
 不用擔心那個人不見了.
 財散人安樂.
 是不是這樣?.
 結果不是.
 找到之後怎麼樣?.
 大排筵席 開感恩會.
 我記得我小時候在香港.
 我的執事做榜樣.
 隔一排就開個感恩會.
 星期六下午.
 大家男女老幼一齊回來教會.
 吃餐飽的.
 我們那時候住在松屋村.
 很窮 沒有錢.
 有東西吃就趕快來.
 什麼豬朋狗友都來一齊吃.
 吃的時候有人問.
 為什麼我要開感恩會?.
 哦 原來什麼什麼事.
 有時候就是病得醫治.

$^{41}$ 生意差點倒了又好了.
 或者是遇到什麼大災難.
 經過祈禱 經過神的帶領.
 將感恩的祭就還給神.
 不單止是獨樂.
 還要重樂.
 全教會都請來.
 哇 很大陣仗.
 你病得醫治.
 那你就要花那麼多花費.
 那時候我印象最深.
 一散了會就搶個綠寶橙汁.
 就這樣想.
 怕用不完.
 所以這三節經文.
 好像很簡單.
 其實裡面跟錢有一點關係.
 但這個錢銀的關係.
 是連接到神的關係.
 所以這次的經文.
 其實我們平時讀聖經就過眼雲煙.
 看到失陽的比喻.
 後面就想看一下浪子的比喻.
 講錢失感情.
 所以就不講錢.
 如果牧師講錢.
 當然是講奉獻.
 所以就一看到講這件事.
 我就說下次不要預我了.
 幸好不是張牧師講.
 是梁牧師講.
 大家可以放心.
 我就不計算了.
 講錢失感情.
 就算感情下個禮拜見不到我.
 就沒有所謂了.
 但之前我們都是同心低頭祈禱.
 讓神的靈幫助我們.
 天父求你.
 多謝你給我們寶貴的生命.

$^{81}$ 並且你是將你自己的獨生外子.
 寶貴的生命賞賜給凡尋求你的人.
 你就必與我們同行.
 與我們同在.
 並且是差遣聖靈.
 保證我們在基督裡面.
 蒙你保守.
 無論我們如福如禍.
 我們都是屬於神的寶貴的子女.
 以致我們真是天天存在感恩.
 來寄在我們的賬簿裡.
 來明白數算主恩是多麼豐富.
 讓我們天天都能夠快樂的迎接新的一天.
 無論信或者逆.
 我們知道我們在神的手中一無牽掛.
 凡事藉著禱告祈求感恩.
 將我們所要的告訴你.
 你就賜出人以外的平安.
 保守我們的心懷意念.
 求你在這個時刻做我們的老師.
 無論說的能夠說得明白.
 讓我們聽的同感一靈.
 聽得清楚.
 並且賜我們力量進行你的話語.
 垂聽我們的禱告.
 奉耶穌基督的聖名祈求.
 阿們.
 傳福音的指引.
 賜給我們.
 使出行動表明我們屬於基督的人.
 浸完就拜拜了.
 拿了證書就走人了.
 是不是這樣.
 要怎樣.
 十一奉獻.
 要講錢.
 因為人的錢在哪裡.
 心就在那裡.
 如果你真的將十一奉獻放在教會.
 你的心會不會在教會呢.

$^{121}$ 如果你每次都被人糾纏的話.
 那麼你回來教會.
 可有可無.
 是不是.
 你說十分之一.
 哇 真是不得了.
 整部家產的十分之一.
 全都給了他.
 唉 那你就要看清楚了.
 哇 今天還有這麼多空位.
 趕快多搬些人來吧.
 因為你的心就在那裡.
 不單止奉獻.
 你說崇拜夠不夠.
 不夠啊.
 理論上我們知道.
 祈禱會的人.
 多過崇拜的人數.
 你想想那間教會會怎樣.
 一定坐滿人.
 但如果調轉的人.
 那你就明白為什麼會這樣.
 禱告的人少.
 你聚會的人數多.
 好不好啊.
 好一點點而已.
 如果你祈禱會的人數.
 多過崇拜的人數.
 你想想.
 哇 這間教會.
 那個聖靈的充滿大不大啊.
 一定大啊.
 一百二十個人在馬開樓上.
 不停地祈禱.
 結果怎樣啊.
 聖靈就降臨了.
 張牧師一講道.
 就三千人信主了.
 是吧.
 講一次就三千人了.

$^{161}$ 第二次怎樣啊.
 第二堂道就五千人信主了.
 你說這是什麼事啊.
 那條道理.
 已經擺得很清楚了.
 但是講就容易.
 信就難.
 做就更加難.
 是不是這樣呢.
 所以在這個比喻裡面.
 我們看到.
 回到經文那裡.
 大家手上有這個秩序表.
 我就不打頭影給大家.
 讓大家可以低頭.
 可以看.
 第一節第二節是講什麼呢.
 就是當時路加寫這封信的時候.
 就看見有一個環境.
 就是稅吏和罪人.
 都要靠近耶穌.
 聽他講道.
 我覺得路加寫得真是好.
 因為我們馬上想.
 聽什麼道啊.
 做什麼事啊.
 為什麼會聽道啊.
 今天我們傳福音.
 不是這樣拉人去信耶穌的.
 不是這樣回來教會的.
 你可能要請他去.
 發發發.
 吃一餐大的.
 吃一餐厲害的.
 吃到飽了.
 喝飽吃醉.
 醉醺醺就被你拉回來教會.
 所以你喝茶很重要的.
 摸摸茶杯底.
 吃過.

$^{201}$ 那時候呢.
 你就.
 將他吸引到.
 哇 為什麼你這麼大牌啊.
 可以請我吃.
 包肥餐啊.
 有什麼目的啊.
 是不是這樣.
 但是你見面三分鐘.
 認識了他.
 請了他吃飯.
 你叫他回來教會.
 就容易了.
 所以你看看第二次的話.
 法利斯法人和文士呢.
 就是批評耶穌了.
 他說這個人接待罪人.
 是不是啊.
 跟他們吃飯.
 哦 原來耶穌先跟他們吃飯.
 回來教會容易很多.
 不是說馬上有.
 你可能一餐 兩餐 三餐 三十餐.
 都不一定來的.
 是不是啊.
 有沒有人試過請人吃飯.
 目的就是要和他回教會的.
 有沒有啊.
 請一餐 不來就算了.
 是不是啊.
 請他三十餐.
 有沒有搞錯啊.
 這裡呢 比喻就講得很清楚.
 主耶穌明白這班法利賽人.
 他們想聽的.
 早應該是黑了眼訓的了.
 聽什麼啊.
 一百隻羊 走了一隻.
 喂 好悶喔.
 哪有這麼傻的啊.

$^{241}$ 九十九隻不理.
 走去找那隻.
 回來 唉 黑眼訓 不想聽了.
 接著就講錢.
 哇 見錢就怎樣啊.
 眼會怎樣啊.
 隻眼開隻眼閉還是怎樣啊.
 見錢就眼開了.
 所以他們就聽了.
 嗯 你想說什麼.
 下面就有三個問題.
 我不知道大家讀聖經讀了多久.
 有沒有人讀過這個.
 失前比元的舉手 我看看.
 啊 這麼少啊.
 從來沒有讀過啊.
 不是吧.
 OK 你讀過 但是讀了一分鐘之後.
 接著就沒有再讀的 舉手我看看.
 坦白一點吧.
 因為講錢嘛.
 是不是.
 會不會這樣啊.
 不知道啊.
 但是我們的注意力.
 往往就落在第三個故事裡.
 看個浪子.
 唉呀 好感人啊.
 浪子回頭 怎樣啊.
 金不換.
 金 前面就是講金錢.
 金換不到啊.
 是不是.
 接著呢 現在呢.
 現在視頻上有很多解經.
 開始集中注意力.
 罵那個長子.
 就罵我們這些基督徒啊.
 唉呀 不生羊啊.
 還在生氣那些浪子回頭啊.

$^{281}$ 但是有沒有人試過啊.
 教會十幾二十人信耶穌.
 大家就不回教會了.
 會不會這樣啊.
 不會啊 很開心啊.
 快樂一 快樂一.
 是不是.
 但是我們就想啊.
 究竟 整年啊.
 是只有那幾丁信耶穌的啊.
 會不會這樣啊.
 我們就要反思了.
 我們可能就是長子.
 浪子來了 好像就.
 沒什麼所謂.
 那這些呢 以後才有機會講.
 現在還要回去.
 這三節聖經呢.
 就有三個問題.
 大家讀的時候.
 不知道有沒有發覺.
 裡面有三個問題.
 第一個問題.
 其實第一個問題都很容易答的.
 他說一個婦人.
 耶穌講了三個比喻.
 這個是第二個.
 所以他有個畫示第八節.
 有一個婦人有十塊錢.
 跌了一塊錢.
 就大掃除.
 連二十八 除垃圾這樣.
 就特意找那個銀子.
 那你想一下.
 這個婦人是誰呢.
 你想一下是誰.
 就是神自己.
 神的旨意愛世人.
 祂是找人.
 不是人找到神的.

$^{321}$ 永遠是神找到人的.
 是不是.
 所以這個婦人呢.
 好多解經家都是講.
 這個就是耶穌.
 耶穌為了要找.
 那個銀子.
 結果就翻箱倒櫃.
 大掃除.
 搵搵搵搵.
 好了 問題就來了.
 那這個錢是誰呢.
 指誰呢.
 就是前面的稅吏和那些罪人.
 是不是.
 咦 又不像.
 否則路加不會這樣寫.
 因為後面就說.
 稅吏和法理賽.
 不是 法理賽人和文士.
 其實這些都是.
 失落的銀子.
 你這樣讀聖經.
 就有意思了 你就會讀下去.
 究竟發生什麼事.
 為什麼要翻箱倒櫃.
 去找這個.
 失落的銀子.
 好像失落的羊一樣.
 大家覺得.
 錢有腳可以走的 舉手給我看看.
 沒有人舉手的.
 你的銀子.
 沒有腳的.
 走不掉的.
 請舉手給我看看.
 客氣一點.
 都沒有人舉手的.
 梁牧師你在說其他東西.
 是呀 我們想想.

$^{361}$ 究竟錢有沒有腳呢.
 銀子會不會走路呢.
 會不會長翅膀會飛走呢.
 道理是不會.
 是不是.
 但是你今天.
 你看看這個世界.
 睜開眼一開.
 油價升了兩元.
 你結果會怎樣.
 你的錢就走了.
 你餓壞的錢.
 昨天給的.
 和今天給的.
 有沒有差別.
 好了 錢究竟會不會飛走呢.
 會不會飛走呢.
 好了 又換一句說話.
 錢會飛走.
 但是.
 錢會不會飛過來呢.
 橫財就手.
 是不是.
 大家都很眉飛色舞.
 說到錢.
 我說一個小小故事.
 希望有時間.
 你知不知道我以前做姓韓的.
 我以前是數銀的.
 禮拜一到禮拜五.
 數銀紙的.
 收工的時候.
 又是數銀紙的.
 嘩 你不是很有錢.
 不是 那些錢是別人的.
 我以前在.
 窩打老道.
 Waterloo Road 香港.
 有間四海保齡球場.
 如果你住過香港.

$^{401}$ 九龍就知道在哪.
 培正的斜對面.
 下面就有一間.
 浸信會書館.
 我中學畢業的時候.
 考不到大學.
 我就在那裡打工.
 一做就做了六年.
 在那裡repeat repeat.
 考到入台大衛止.
 所以我六年.
 就在那裡數銀紙.
 上班就是早上.
 幫老闆開了一扇門.
 那時候也是美南浸信會的.
 孫教士.
 就為了福音的緣故.
 就開了個出版社.
 浸信會出版社.
 後來出版社那些人.
 不會來出版社買書的.
 就索性開書局來賣書.
 那我就走去賣書.
 就賣聖經.
 賣蜀明書籍.
 就賣很多東西.
 所以早上開門.
 就要.
 數銀紙.
 確定有錢.
 可以找續的.
 那才可以開門.
 那收舖的時候.
 那你就算一下.
 結帳有多少人買到東西.
 多少錢.
 那結果發現什麼呢.
 我們就想.
 收工了 七點鐘放工.
 我可以去補習班.

$^{441}$ 準備考大學的補習班.
 七點十五分就要上課.
 十五分鐘之後.
 衝衝衝.
 從四海保齡球場那裡.
 衝到珠海.
 到仙獸班.
 那就.
 想著十五分鐘就搞定了.
 結果一賣書.
 一顆小錢.
 大家怎麼想.
 你收入這麼多.
 你支出又這麼多.
 或者你收來.
 你早上的.
 那些.
 預備的零錢.
 和你收到的.
 和那個對數機.
 少了一塊錢.
 那你怎麼辦.
 馬上去找.
 要數.
 重新數 有沒有數少了.
 數錯了 或者多了.
 一做就遲到.
 五分鐘了.
 那時候的跑步.
 特別快 比別人快.
 因為馬上趕上課.
 所以失落一前.
 對這個婦人.
 就好像我.
 以前.
 開店的時候要找數一樣.
 你少一個"崩"都不行.
 所以這塊錢.
 緊不緊要呢.
 緊要.

$^{481}$ 對她來說非常之緊要.
 一少了這一塊錢.
 都不行.
 她又不是開浸信會書局.
 她只是有十塊錢.
 少了一塊錢.
 是不是這樣算呢.
 原來那個錢.
 不是那個錢本身.
 那個價錢.
 那個價值.
 而是那個婦人.
 和這塊錢的.
 關係.
 是不是.
 有沒有道理.
 我們和神的關係都一樣.
 神.
 敬石你和我.
 都是寶貴的生命.
 是種價買熟回來的.
 所以我講題.
 第一個寶貴.
 就是你有寶.
 和旁邊的人說.
 你很有寶.
 和旁邊的人說.
 你很有寶.
 以前是用來罵人的.
 你很有寶嗎.
 不是啊.
 我們在基督裡面.
 是很寶貴的.
 很矜貴的.
 關係就是這樣.
 所以巴赫萊.
 是英國一個很出名的解經家.
 他特別解釋這段經文的時候.
 他說這個婦人.
 那十塊錢.

$^{521}$ 是可能她的實裝.
 嫁妝來的.
 你少了那顆小銀.
 對她來說.
 就關係重大.
 不是那顆小銀那麼簡單.
 而是那十塊銀.
 配在她的妝架上.
 是不能夠少的.
 那個關係就在那裡.
 那個寶貴的地方就在那裡.
 不是值不值錢的問題.
 今天我們在這個社會裡面.
 時時拿著那個小銀.
 就覺得很有寶一樣.
 可是你想一下.
 是真的.
 你這個小銀掉在地上.
 被人踩.
 踩得髒髒的.
 你會不會撿起來.
 那個小銀就小小的.
 不要了.
 不是的.
 你少了那顆小銀.
 你去Downtown停車.
 你就靠那顆小銀.
 等了15分鐘.
 時間就是金錢.
 所以我們在神的眼中.
 也是.
 神整個的宇宙.
 就將你.
 擺在宇宙.
 天國裡面的一個.
 美好的位置.
 少了你.
 就不像樣了.
 當你想到原來.
 你自己是這麼矜貴的時候.

$^{561}$ 你回到天家的時候.
 神將你擺在那個位置.
 是美麗到不得了.
 這樣的時候.
 我們的生命.
 就會覺得是值得的.
 無論別人怎樣踩我.
 不要緊.
 因為我在神的心目中.
 那顆小銀 雖然髒髒.
 但是一抹乾淨.
 在基督裡面成為聖潔.
 合乎主用.
 在黃金街碧玉城裡面.
 你說好不好?.
 好的舉手.
 勇敢一點.
 不好的就不要舉.
 可能還要聖靈感動你.
 快點去.
 組織一個回應.
 和抉擇.
 這個就是問題.
 錢會不會有腳飛走的?.
 有啊.
 因為這位婦人.
 她是失驚無神.
 跌了這個硬幣.
 她就很緊張.
 要找回它.
 今天我們都是.
 看這個婦人的時候.
 就是看見基督.
 同樣基督是教會的頭.
 教會也是要盡這個本份.
 就去找那些失落的錢.
 我們周圍很多硬幣.
 四處被人劈了.
 是不是?.
 這個社會是.

$^{601}$ 將那些人不當是人.
 在極權的國家裡面.
 死了一個.
 根本就不當是甚麼回事.
 如果一下子死了幾萬人.
 索性在街上燒了就算了.
 人的命根本在極權主義的眼中.
 是完全不值錢的.
 但是在神的眼中.
 每一個錢都值錢.
 你說我是不是很矜貴呢?.
 矜不矜貴?.
 好了,再來.
 錢會不會飛走呢?.
 我剛才說過的比喻.
 錢的確是很容易會飛走.
 因為它的價值每一天都在變.
 尤其是兵荒馬亂的時候.
 以前在南京.
 記得南京大屠殺.
 那些人逃難的時候.
 他們會用很多錢.
 去奢走.
 值不值錢呢?.
 因為不值錢.
 你買一斤米.
 可能要一輛車的鈔票.
 才能買到一斤米.
 今天我們在神的心目中.
 亦都是.
 我們在神的手中.
 可能是零.
 但是有耶穌在前面.
 在支票本的金額上.
 你後面有很多個零.
 前面有個一字.
 就是很值錢.
 如果沒有這個一.
 後面都是零 有沒有用?.
 全部都是零.

$^{641}$ 我們在極權的國家.
 或者一些社會.
 對這個價值的看法.
 是無價值的.
 不值錢的時候.
 這些零就好像.
 就是零.
 但是在神的心目中.
 價值 說得多了.
 時間不夠.
 這是第一個問題 大家知道了.
 現在改正一個觀念.
 錢會不會飛?.
 會.
 那你既然會飛.
 你就知道靠著神.
 神要把錢給你.
 飛到你的口袋裡.
 會不會?.
 你信耶穌.
 耶穌就給錢你.
 是不是?.
 我就發.
 看你要什麼.
 還有看你有什麼.
 你有的五餅如魚.
 你願意奉獻給神.
 你就會怎樣?.
 神就祝福你這個五餅如魚.
 怎樣?.
 吃飽 多少人?.
 讀過聖經很久了.
 多少人?.
 一千?.
 兩千?.
 有人點頭.
 祝福每一位.
 第二個問題.
 找到了.
 站在隔壁.

$^{681}$ 吃一餐.
 厲害的.
 好像開感恩會.
 用不用這麼大陣仗?.
 要不要?.
 道理上是不要.
 怎麼搞的?.
 一元雞.
 花了一百多二百元.
 請人吃飯.
 用不用這麼大陣仗?.
 今天教會也是.
 我們為什麼要.
 弄冷氣 弄音響.
 弄那麼多東西.
 花那麼多錢.
 我們幾個人坐在一起.
 聽牧師講就行了.
 不是的.
 當你這樣鋪排的時候.
 誰也不要.
 他們就會回來.
 為什麼你要請我們吃飯?.
 我掉了一元雞.
 要不要這麼大陣仗?.
 其實因為這樣的緣故.
 他們就發現.
 一個信耶穌的人.
 每一天凡事都謝恩.
 哪怕掉了一元雞.
 他們都謝恩.
 他們就會好奇.
 會認識究竟信耶穌.
 有什麼好處.
 回來悟道園.
 究竟有什麼好處.
 原來我們可以歡喜快樂.
 願意奉獻給神.
 讓神使用我們五餅而言.
 他們就會覺得.

$^{721}$ 我們真的想明白.
 你心中盼望的緣由.
 我們存在著.
 極為溫柔的心.
 來跟他們分享.
 你說好不好.
 那麼這個答案.
 滿不滿意.
 這是第二個問題.
 第三個問題.
 在哪裡呢.
 第三個問題是什麼.
 看看.
 我告訴你們.
 一個罪人悔改在神的使者面前.
 也使者為他歡喜.
 你要圈住他.
 罪人悔改.
 關我什麼事.
 這個就真的.
 非常的重要.
 為什麼我們每次都要回來查經.
 回來祈禱.
 回來崇拜.
 因為你我也要知道.
 罪人悔改.
 不是信主那一刻那麼簡單.
.
 六百年來.
 基督教改教之後.
 給我們一個重新思考.
 悔改這個字.
 是說.
 我看錯了.
 我要修正.
 我錯誤的觀念.
 所以不是.
 決志信主那一刻那麼簡單.
 而是你.
 每一天都要有悔改的心.

$^{761}$.
 這個悔改的意思就是.
 你願意開拓你的心胸.
 讓神的話進入你的心中.
 然後提醒你.
 哦.
 原來是這樣的.
 我們先因為.
 藉著神的話.
 藉著信靠主耶穌.
 聖靈的囂觀.
 感動我們.
 將我們的思維.
 從熟悉的想法.
 回到神那種的想法.
 從我們橫面.
 一二三四這樣的思維.
 我們能夠從上.
 降下來.
 聽到神的聲音.
 回應神那種呼召.
 這不是一朝一夕的事.
 而是每一天的潛移默化.
 不是只靠祈禱會.
 不是靠查經班.
 不是靠崇拜.
 這麼簡單.
 其實崇拜.
 我默會四十年.
 都是發覺越來越覺得.
 很多時候.
 需要每一天謙卑的.
 將神的話吸在我們心裡.
 讓我們.
 靈裡面的運作.
 不斷的去行動.
 所以才落到第三點.
 就是報章.
 其實我跳過了.
 我保證就是.

$^{801}$ 我們做了這麼多之後.
 聖靈不斷的.
 保證我們不會有缺乏.
 保證我們的生命.
 在神裡面是寶貴的.
 神一定會保守.
 但是我們的重點就是.
 要記下神的恩典.
 要數算神.
 每一天給我們做的.
 就好像我小時候.
 教會的執事 開感恩會的時候.
 我們不只是貪那六寶橙汁.
 而是想聽這位.
 屬靈的長者.
 怎樣在生命裡面.
 感恩的心.
 所以到現在.
 六十年了.
 還有這個印象.
 原來 教會.
 是神的.
 兒子.
 從家買熟回來的.
 我們聚在一起.
 本來就應該.
 這樣為那些.
 願意悔改.
 或者思維從熟悉的.
 回到神的思維裡面.
 那種感動.
 那種感恩.
 大家一起開心.
 所以在人的角度裡面.
 第八節到第十節.
 敏茲法拉塞.
 是不會明白的.
 是到聖靈降臨之後.
 他們才覺悟.
 其實中間有人覺悟.

$^{841}$ 可能只是彌哥底母.
 其他都未到家.
 都不一定.
 但是無論怎樣.
 路加寫這段話.
 不單止耶穌想.
 向這些.
 法拉塞人和文士.
 來挽回他們.
 今天這個訊息.
 也是想挽回.
 你我不要只是看錢份上.
 而是看神.
 和你的關係份上.
 以致你能夠.
 將金錢當作糞土.
 以認識耶穌基督.
 為自保.
 所以你很有保.
 是不是.
 所以求主幫助我們.
 好了 說到行動.
 今天和大家有一個挑戰.
 如果神允許的話.
 我希望五月可以回來和大家見面的時候.
 你們可以.
 在神的面前.
 能夠有一個改善.
 怎樣改善呢.
 我們呢.
 去年已經和大家.
 有這些操作.
 希望大家繼續有.
 屬靈的.
 的.
 保障.
 寫寶.
 記事寶的寶.
 障是一筆障.
 一筆障寫下.

$^{881}$ 神怎樣感動你.
 是哪一句說話.
 這幾次聖經.
 讓你認識.
 神是一個怎樣的神.
 我希望你的答案是寫婦人.
 不會寫梁牧師.
 句話提醒我們.
 處境是甚麼處境.
 你究竟是被神.
 找回那個銀子.
 還是你仍然是.
 掉在地上.
 那個銀子.
 需要神去找回你.
 第三.
 你向神的禱告.
 的回應是甚麼.
 我盼望神找回我.
 或者我盼望.
 神更深.
 我回家後.
 我仍然會上視頻.
 看這段經文.
 或者打開聖經.
 看這段經文講甚麼.
 然後你記下.
 你跟神說過甚麼.
 神給了你甚麼答案.
 何時給的答案.
 結果是怎樣的答案.
 你記下.
 整年或整個禮拜.
 神聽你祈禱的.
 結果是甚麼.
 昨天我跟老婆吵架.
 跟師母吵架.
 師母就說.
 你都不遷就我.
 不是這樣說的.

$^{921}$ 說起來.
 我現在是這樣裝的.
 她說我已經遷就你了.
 我都已經.
 取消了約定.
 就要跟你赴約.
 去做這件事.
 她說沒有的.
 沒有這件事.
 我說不是的.
 我本來有一個醫生的約定.
 我就.
 不能夠.
 推卸這個約定.
 我就跟神說.
 神啊 我跟太太有約.
 那怎麼辦.
 我就私底下跟神說.
 神好奇妙.
 給我一個靈感.
 打個電話給.
 護士小姐.
 她說.
 你的約定.
 可不可以提早.
 我說有開場.
 如果有哪個.
 病人請假.
 不能來看的.
 你就放我進去.
 讓我早點去看醫生.
 她說好.
 聽你這麼說.
 但事過未發生.
 再過幾天.
 本來禮拜一.
 我一排你禮拜五下午.
 可不可以.
 我當然可以.
 我就告訴大家.

$^{961}$ 神聽祈禱.
 為了不會爽約.
 結果神怎樣.
 跟我鋪路.
 禮拜一變成禮拜五.
 我現在就跟你親密.
 看看信不信你支笛.
 這些以後再說.
 但起碼.
 在整個過程裡面.
 我看見神.
 有答應我的禱告.
 但問題也要檢討.
 你的禱告.
 是否妄求呢.
 很重要.
 所以你也要記下.
 我不應該去拉斯維加斯.
 賭個大筆.
 輸了或怎樣.
 你要寫下.
 我不應該這樣祈禱.
 要悔改.
 我的女婿還沒信耶穌.
 繼續祈禱.
 將這些筆記下來.
 當你越寫的時候.
 你就發現.
 一年下來.
 神的恩典是數算不盡.
 臨睡前.
 不要擦手機.
 不要做這做那.
 好好思想.
 神你在我今天.
 24小時裡面.
 給我有些什麼.
 我經歷了什麼.
 是神你所喜悅的.
 這樣的時候.

$^{1001}$ 你一覺一定睡到天亮.
 還有一件事要做的.
 不單止你的靈.
 裡面需要神的話.
 需要記下.
 神的恩典.
 但是你要操練自己.
 容許我.
 大家做一個小小的運動.
 請大家.
 站立.
 大家都站得低的.
 站得很深的.
 不站的不用什麼.
 不用勉強.
 認為你站得穩的.
 就站好.
 兩腳.
 跟肩膀齊寬.
 現在.
 將它合起來.
 兩隻腳並在一起.
 你覺得.
 最強壯的那隻腳.
 就留下來.
 比較弱的.
 你就抬起它.
 抬起一點.
 腳尖微微的.
 碰到腳板就可以了.
 如果你覺得站得不穩.
 你就扶著前面.
 是不是.
 然後.
 你自己.
 試試輕輕將手.
 拿開.
 你站得穩嗎.
 站得穩嗎.
 開始搖搖續墜了.

$^{1041}$ 站得很好.
 舉手給我看看.
 不錯.
 都有新的安裝物件.
 好了.
 雙手離開.
 一隻腳.
 輕輕點地.
 然後閉上眼.
 1 2 3.
 好了.
 放下.
 大家感覺如何.
 會不會閉上眼.
 那一刻你會搖搖欲墜.
 不會嗎.
 會嗎.
 為什麼會這樣.
 好了 坐下.
 這個運動呢.
 希望你早晚都做一次.
 你就看著那個鐘.
 一秒.
 你就把它弄下來.
 第二天試試兩秒三秒.
 到五月我回來的時候.
 我希望你可以.
 維持到五秒.
 閉上眼.
 金雞獨立.
 啊.
 聽我講道.
 不是 不是講那麼長.
 就是你完了之後.
 你可以show給 梁牧師.
 我已經做到了 五秒都可以了.
 五十秒都可以了.
 你這樣操練的時候.
 你的腳的腳骨力.
 是不是.

$^{1081}$ 不單止好 還有你的平衡.
 你天天這樣練.
 你就沒那麼容易跌倒.
 我認識一個牧師.
 昨天電話給我 他八十八歲.
 他說去年他跌一跌.
 噢 盤骨就裂了.
 裂了到現在.
 都不能走路了.
 所以呢 馬步要很重要.
 今年是什麼年啊.
 馬年啊.
 我不是叫你紮馬.
 我是當雞獨立.
 當雞獨立.
 你試試 你這樣練.
 還有你閉上眼睛.
 為什麼要閉上眼睛呢.
 因為我們原來單腳站的時候.
 我們靠眼睛維持我們的平衡.
 但是如果你閉上眼睛.
 變成了你潛意識的運作.
 這樣的時候.
 你的潛意識.
 你整個筋脈.
 從頭到尾.
 這樣去運作的時候.
 你能夠站得穩.
 意思就是說.
 你不是靠對眼.
 而是靠你的意志力.
 你去判斷.
 屬靈的生命一樣.
 你將神的話藏在心裡面.
 什麼叫免得我得罪你.
 意思就是說.
 免得我行差踏錯.
 時時要聖靈擔憂.
 要聖靈去提醒你.
 而是每天神的話.

$^{1121}$ 浸淫在我們的心裡面.
 所以我們浸信會.
 是提倡浸.
 不是浸禮那一下.
 而是每天將神的話浸在我們的心裡面.
 我們自然.
 無意識當中.
 仍然可以站得穩.
 今天教會需要復興.
 就是要我們在屬靈的生命.
 要站得穩.
 你站得穩就是要神的話.
 在你裡面.
 你站得穩是要經歷在這個行動裡面.
 如何發放出神給你的能力.
 來運作.
 使得人們四處看看.
 你盼望的緣由原來是這樣的.
 我也想回來悟道園.
 聽一聽張牧師的講道.
 好不好.
 雖然遲了.
 現在三月了 但我仍然立志.
 你說不行了 牧師已經為我祈禱了.
 行行行.
 到了復活節.
 不要到時候仍然.
 原地踏步才行.
 你要在神的面前.
 就是要這樣不斷地操練.
 好不好.
 求神幫助我們 我們一起祈禱.
 天父上帝多謝你賜給我們恩典.
 讓我們認識你.
 成為我們屬靈的基礎.
 更加成為我們的目標.
 就是要榮耀你.
 讓你的生命.
 繼續運行在我們當中.
 用你的靈.

$^{1161}$ 復興我們的生命.
 讓我們不只是.
 練習如何金雞獨立.
 而是天天將神的話.
 藏在心中.
 並且活出裡面盼望的緣由.
 使我們周圍的.
 鄰舍將從我們生命裡面.
 看見神的.
 豐盛是何等的美好.
 祝福你的教會.
 祝福每一位.
 奉你的名.
 祈禱的弟兄姊妹.
 讓未信的.
 能夠相信你.
 相信你的更有信心.
 行出.
 主你在我們心裡面的感動.
 化成為行動.
 奉耶穌基督的名.
 祈禱.
 阿們.
\newpage



\section{創世記 40:8}
\label{sec:031830}
\textbf{聖經人聖經事:夢出於神}~~~~黃克勤牧師博士
\newline
\newline
連結: \href{http://www.church.com.hk/acms/content.asp?site=cdc\&op=show\&type=product\&code=031830\&layout=sermon}{\texttt{031830}} ~~~~ 語音日期: 2026-03-15
\newline
\newline
\hyperref[sec:031829]{\small{< < < PREV SERMON < < <}}
~
\hyperref[sec:index_chronic]{\small{[返順時目]}}
~
\hyperref[sec:031831]{\small{> > > NEXT SERMON > > >}}
\newline
\newline
創世記 40:8
\newline
\begin{longtable}{cl}
\hline
\hline
章節 & 經文 (和合本修訂版)\\
\hline
40:8 & \begin{tabularx}{0.7\textwidth}{X} 他們對他說：「我們各自做了一個夢，卻沒有人能講解。」約瑟對他們說：「解夢不是出於神嗎？請你們把夢告訴我。」 \end{tabularx} \\ \\
[1ex]
\hline
\hline
\end{longtable}
$^{1}$弟兄姊妹平安.
再繼續一起讀聖經.
聖經人,聖經事,生命樹,宣教網絡為大家準備的.
今天繼續講約瑟故事.
我相信我們一直講下去.
講到月尾還是約瑟故事.
因為約瑟故事很長.
很多片段.
在《創世記》裡面.
其實有五個,六個很重要的故事.
其中有亞當故事,羅亞故事,亞伯拉罕故事.
以撒算不算有故事呢?.
如果算的話也是重要的.
有雅各故事,約瑟故事.
就看你算不算了.
很長的就是雅各故事和約瑟故事.
重要性就是亞伯拉罕故事.
也是非常重要的.
「夢出於神」就是今天的主題.
本來是「解夢出於神」.
不過「解夢出於神」是怎樣解呢?.
固然是神告訴我們.
但是夢本身是神給我們的.
「夢出於神」.
這個更加重要.
我們看下去.
《創世記》第四十章第一節.
「這是以後埃及王的酒政和膳長得罪了他們的主」.
就是埃及王.
即是法老.
所有埃及皇帝都叫法老.
法老就擄奴酒政和膳長這兩個大臣.
將他們下在護衛長府內的監獄.
就是約瑟被囚的地方.
大家就會看到這個情況.
「這是以後埃及王的酒政和膳長得罪了他們的主」.
埃及王法老就擄奴酒政和膳長這兩個人.
酒政就是斟酒給皇帝喝的.
負責酒的.
他不只是斟酒.

$^{41}$是負責買什麼酒.
什麼酒皇帝會喜歡喝.
什麼時候皇帝會喜歡喝什麼味道的酒.
有些外國侍臣送來的酒是怎樣的.
有很多種酒.
他就負責酒.
而且最重要的是斟下來.
斟到到處都說「杯子不對」.
這樣是不行的.
膳長是負責膳食的.
斟酒的容易些.
只是負責膳食的容易些.
其實各有難處.
不過膳長是困難少少的.
因為膳食的事情要安排的比較多.
有一次不知什麼事沒有說得罪了皇帝.
其實酒政膳長是皇帝很親密的人.
他們時時和他聊天.
因為是為他煮飯吃.
和他很親密.
有很多心事要說.
有時候說多了 說大了.
說得嬉笑了.
和他很熟悉.
於是生氣了 拍桌子.
就可以坐牢了.
那次兩個人一起都惹皇帝生氣了.
於是皇帝就說.
先坐牢 拘捕他.
馬上叫人拘捕.
護衛長在旁邊馬上過來.
就將他們入下再監禮.
今天夢是出於實.
護衛長將他們交給藥室.
由這一句就知道.
其實護衛長波提弗.
在他太太無告藥室.
他進來戲弄我的時候.
波提弗聽到生氣了.
可能他在想藥室不是這樣的人.

$^{81}$他認識他那麼久.
什麼都不偷吃.
他不是這樣的人.
我認識他.
所以他只是說.
押入大牢等候發落.
沒有將他殺死.
否則一刀刺死他.
一路一路進去之後.
連監獄長都很信任藥室.
因為神與他同在.
上兩次最主要的訊息是.
耶和華與藥室同在.
無論你在哪裡.
不過藥室有神耶和華同在.
他就很信任你.
於是他管理了監獄.
護衛長波提弗知不知道呢?.
知道的.
不過一時間拿不出來.
於是護衛長波提弗就知道了.
原來藥室正在管理監獄裡的工作.
於是就將酒精膳長交給藥室.
雖然是一句這樣的說話.
但你會看到.
其實在監獄裡.
上帝的手仍然在工作.
沒有離開.
上帝的同在沒有離開.
仍然在藥室的身上.
請你記得.
無論你遭遇什麼困難.
多大都好.
最重要是耶和華與你同在.
藥室就侍候他們.
為什麼呢?.
因為他們畢竟都是官.
都是皇帝身邊的人.
藥室只是一個奴隸.
雖然是信任你做管理.

$^{121}$但都是侍候他們.
應該都侍候得很好.
有些日子.
他們在監獄裡坐著.
就能出來.
被囚在監獄的埃及王的酒精膳長二人.
同夜各做一夢.
各夢都有講解.
兩人同在那一晚.
大家都做了一個夢.
各夢都有講解.
即是說兩個夢當然是不同的.
為什麼有這個夢呢?.
在以往的人.
是比今天更加會歡喜解這個夢的.
是覺得天上的神靈.
透過夢跟我們說話.
所以他們做了一個夢之後.
就覺得這個夢那麼清楚.
當然是神跟我們說話.
你有沒有做過夢呢?.
當然有,每晚都做.
但這個夢有時候.
要先記清楚.
不要忘記.
起床就忘記了.
很多時候都是這樣.
除非你一覺醒.
馬上用筆或錄音錄下來.
否則你的夢是忘記了.
之後都會忘記.
如果是我自己讀潛意識治療.
讀催眠治療等等.
其實很有可能.
夢是潛意識跟我們的意識說話.
於是有夢法.
另一個就是神超自然.
想跟我們說話.
於是透過夢跟我們說話.
是有可能的.

$^{161}$如果這裡是夢.
應該是神跟我們說話.
告訴我們一些訊息.
不過有時候夢不是很明白.
到了早晨.
藥室就進去看他們.
兩位早晨.
昨晚睡得好不好?.
見到兩位愁眉苦臉.
有愁悶的樣子.
見到兩位酒精善長都愁悶.
便問法老的二神.
就是與他同囚在他主人府裡的.
你們今天為什麼面帶愁容呢?.
這是我第二點要說的.
第一點.
在上面護衛長交給藥室.
其實護衛長是信藥室.
否則不會交給他.
否則不當面殺他.
都叫人在監獄裡橫手殺他.
戲弄我太太.
在上面護衛長很明顯是信藥室.
第二點.
藥室除了做得好.
很順利做得很賣力之外.
他早晨會看看兩位大人.
昨晚睡得好不好?.
為什麼今天兩位面帶愁容.
是不是早餐不對?.
很聰明,很會問,很會說.
為什麼面帶愁容呢?.
第二點.
藥室鑑貌辨識.
為什麼今天兩位面帶愁容.
這是很重要的.
不要以為自己很忠直,做好所有事情.
其實很多人都很聰明.
自己做事很聰明.
不過是他自己做事很聰明.

$^{201}$但是他不能體會其他人的感覺.
一概不知,只想著自己.
藥室開始學到了.
藥室以前在家裡是怎樣的?.
公子哥兒.
不用出去牧羊,穿彩衣漂亮衣服.
要什麼都有.
經過做奴隸一直到現在.
他已經練成了一個.
很好的人際關係.
做到一個很好的領袖.
很能夠解決問題.
又會看別人眉頭眼額.
已經很厲害.
兩位今天為什麼面帶愁容呢?.
早餐不對.
要怎樣?有什麼事可以幫忙?.
很聰明.
他們說,我們剛剛談起.
我們大家都各人做了一夢.
沒有人能解.
即是說沒有人能解.
兩個人怎樣解呢?.
我不懂得解你那個,你不懂得解我那個.
沒有人能解,即是說.
他們都有機會跟其他犯人聊天.
藥室來到,即是說他們不是一起囚禁的.
這兩個人可能比較重要.
獨立囚禁,分開.
可能這兩個人自己一個房間.
有機會其他人來幫忙工作.
聽到有個夢,沒有人能解.
藥室就說,解夢不是出於神媽.
請你將夢告訴我.
藥室挺厲害的.
為什麼呢?.
藥室是不是懂得解夢的呢?.
有沒有學過呢?.
要不要學的呢?有沒有學的呢?.
藥室很自然地說.

$^{241}$解夢是出於神的.
你們將夢告訴我.
即是說,如果你想解到夢.
是疑夢神給你訊息.
你就一定要問神才行.
你將你的夢告訴我.
看看我能不能幫到你.
這點大家要了解.
解夢不是出於神媽.
在解夢之前,先需要發夢.
這個夢是出於神的.
為什麼我今天的題目是夢出於神呢?.
其實神會告訴我們.
前面是怎樣發展的.
有什麼事會發生.
究竟是怎樣會發生.
發生什麼事,什麼時候發生.
你是不知道的.
如果神沒有說,你是不了解的.
不過這個時間.
藥室經歷了很多.
可能他自己也發過很多夢.
然後神就跟他講解夢.
在他心靈裡.
有神的靈會將夢.
浮現出來.
滔滔不絕地告訴他.
這個夢是什麼意思.
也可能是他基本上不知道.
有一個夢,然後他禱告.
如果你發了一個夢也可以.
主啊,發了一個夢是這樣的.
請你將夢告訴我.
究竟想跟我說什麼訊息.
請說服人聽聽.
我等著你跟我說.
也可以的.
神是否會馬上打雷劈下來.
去講這個夢給我們聽呢?.
可能是,可能不是.

$^{281}$或者是過了一段時間.
你才會發現,原來是這樣解的.
不過這裡顯然藥室.
對於解夢是有把握的.
那它的解夢有把握是什麼呢?.
請你留意幾個字.
也是神歡喜藥室的原因.
「不是出於神嗎?」.
「不是出於神」就是說.
這些東西都是出於神的.
三個字.
「出於神」.
大家一定要學到的.
如果你在後面再看.
藥室的哥哥.
認出藥室之後.
藥室就小小作弄他們.
然後他就認出那些哥哥.
那些哥哥臉都青了.
「哎呀,我都不知道該說什麼」.
「請你放過我們吧」.
他們想了一些辦法.
「爸爸還沒死的時候」.
「就叫你寬待我們」.
「爸爸都死了,怎知道要說什麼」.
藥室說.
「不要緊」.
「我看出這件事是出於神的」.
「所以你們想害我」.
「不過神是出於祂的好意」.
「讓我先來這裡預備一族人」.
「不會饑荒餓死」.
藥室也是說這三個字.
「出於神」.
所以藥室在自己的想法裡.
重新由他小時候的想法.
回想到現在.
他得出一個結論.
三個字.
「出於神」.

$^{321}$如果神要指示你.
可能他用「夢」.
而這個「夢」也是「出於神」.
所以如果你想知道這個「夢」是什麼意思.
就問神.
因為「夢」是「出於神」.
用這個「夢」來解釋一件事.
就是告訴我們.
我們很多時候人生會發生的事.
是「出於神」的.
大家知道是「出於神」的.
你要知道.
了解也是這樣.
這件事一定要記得.
「夢」是「出於神」的.
請你將「夢」告訴我.
你告訴我.
我馬上按下按鈕.
「嘟嘟嘟」就一答了.
其實我也不知道你做了什麼夢.
和「夢」是什麼意思.
我也不知道.
只不過藥室聽了之後.
他心中就會求問.
「神啊,現在我們要解夢了」.
「你將你的心意告訴我」.
「如果你告訴我」.
「我就能夠和別人解釋」.
如果神沒有說.
「我沒有特別的感動,怎樣解釋」.
解釋不了.
「這樣解釋,這樣解釋」.
「神有感動我」.
可能是這樣.
記得嗎?.
最重要的三個字是.
「出於神」.
今天跟大家分享到這裡.
記得那三個字.
「出於神」.

$^{361}$「夢」是神給我一個先兆.
告訴我.
如果有.
「夢」也是「出於神」.
你要知道是很謙卑地去等候.
求神告訴我.
我所有事.
無論人生裡的高低.
艱難,痛苦還是很開心.
都是「出於神」.
即是說.
神是看著,照顧著的.
在神來說.
是沒有意外,沒有意想不到.
出於意料之外.
沒有想到.
神預備好,計劃好.
信靠祂.
我們一起禱告.
親愛的主,我們信靠你.
因為你透過若西的說話告訴我們.
「夢」是「出於神」.
所以解「夢」也是「出於神」.
其實最重要的就是.
知道我們的人生都是「出於你」.
將一切的重擔.
交託給我們的神.
正如聖經說.
「將我們的重擔射給耶和華」.
「因為祂顧念我們」.
我們知道所有事都是「出於神」.
是你自己奇妙的作為與安排.
讓我們靜心觀察.
讓我們存一個敬畏的心去等候.
讓我們學習,有信心.
禱告奉耶穌名求.
阿們.
\newpage



\section{詩篇 127:1-5}
\label{sec:031831}
\textbf{建造,看守,發展}~~~~陳德蓉傳道
\newline
\newline
連結: \href{http://www.church.com.hk/acms/content.asp?site=cdc\&op=show\&type=product\&code=031831\&layout=sermon}{\texttt{031831}} ~~~~ 語音日期: 2026-03-15
\newline
\newline
\hyperref[sec:031830]{\small{< < < PREV SERMON < < <}}
~
\hyperref[sec:index_chronic]{\small{[返順時目]}}
~
\hyperref[sec:031832]{\small{> > > NEXT SERMON > > >}}
\newline
\newline
詩篇 127:1-5
\newline
\begin{longtable}{cl}
\hline
\hline
章節 & 經文 (和合本修訂版)\\
\hline
127:1 & \begin{tabularx}{0.7\textwidth}{X} 若不是耶和華建造房屋，建造的人就枉然勞力；若不是耶和華看守城池，看守的人枉然警醒。 \end{tabularx} \\ \\ \relax
127:2 & \begin{tabularx}{0.7\textwidth}{X} 你們清晨早起，夜晚安歇，吃勞碌得來的飯，本是枉然；惟有耶和華所親愛的，必叫他安然睡覺。 \end{tabularx} \\ \\ \relax
127:3 & \begin{tabularx}{0.7\textwidth}{X} 看哪，兒女是耶和華所賜的產業，所懷的胎是他所給的賞賜。 \end{tabularx} \\ \\ \relax
127:4 & \begin{tabularx}{0.7\textwidth}{X} 人在年輕時生的兒女好像勇士手中的箭。 \end{tabularx} \\ \\ \relax
127:5 & \begin{tabularx}{0.7\textwidth}{X} 箭袋充滿的人有福了！他們在城門口和仇敵爭論時必不蒙羞。 \end{tabularx} \\ \\
[1ex]
\hline
\hline
\end{longtable}
$^{1}$各位弟兄姊妹平安.
作為太古城堂教牧童工的身份.
我在這裡講最後一篇《道》.
又是兒童主日的講道.
我想了很長時間.
想和大家分享什麼題目.
可以有什麼講呢.
我起初想講詩篇第八篇.
因為我在這裡講的第一篇講道.
是講詩篇第八篇的.
講題叫《王牌武器》.
是講到兒童口中的讚美.
就是神的王牌武器.
即是兒童對主耶穌的讚美.
就是神的王牌武器.
所以我們要帶領小朋友親近神.
我希望這樣講.
作為我在這裡第一篇講道的回響.
但在我預備的過程中.
神感動我.
祂說你要講詩篇一百二十七篇.
我當時有點掙扎.
因為我已經開始預備.
亦在搜集材料預備去講道.
但既然神吩咐我要講詩篇一百二十七篇.
我唯有去講.
我又要去預備.
我在預備的過程中一直求問神.
神啊 你要我講什麼呢.
和弟兄姊妹有什麼可以分享呢.
求神賜訊息給我.
之後我就定了這個講題.
建造 看守 發展.
希望和大家一起思想一下.
教會兒童侍工的過去.
現在和將來.
我們一起來祈禱.
請大家一起來祈禱.
神啊 你是大能的主宰.
求主讓我們眾人都在你面前謙卑.

$^{41}$讓講的能謙卑.
忠實地傳達你的訊息.
讓聽的能謙卑.
以順服的心聆聽神你自己的說話.
奉主耶穌基督的聖名祈求.
阿們.
今天講的詩篇第127篇.
大家剛才已經讀了.
簡介一下這篇詩.
這是《上行之詩》.
《上行之詩》在詩篇裡面.
由120篇至134篇.
總共有15篇.
而詩篇第127篇.
排在第8篇.
即是《上行之詩》裡面的第8篇.
剛剛在中間的位置.
之前有7篇 之後有7篇.
我們知道猶太人的文化傳統.
放在中間是最重要的.
是焦點的位置.
所以這篇詩很重要.
作者聖經有提及.
就是《所羅門》.
有人研究說.
可能是大衛寫給所羅門的一篇詩.
但無論如何.
都是與皇室有關.
我們先來看經文.
特別當中的字詞的意思是什麼呢?.
這裡提到房屋.
房屋原文是單數 不是眾數.
到底房屋是指什麼呢?.
是否住在那間屋?.
還是皇宮?.
因為聖殿是皇室的人寫的.
或者是指教會?.
國家?.
還是指人生?.
但肯定不是單單指建築物.

$^{81}$應該是指家室方面.
我在這裡就想用一個.
作為一個聖殿.
或教會來看.
這個教會不是照普世教會.
而是照堂會.
城池.
就是城邑.
城邑就是城市.
包括當中的治理.
秩序 人民的生活.
人際關係等等.
至於建築和看守.
字面解釋都知道.
建築和看守.
勞力和勞碌.
不只是工作.
勞累.
是工作至疲倦.
工作至疲倦.
往然這個字.
我們就不能忽略.
因為他提了三次.
往然是指空虛.
虛妄 空泛.
沒有實質作用的意思.
特別是強調價值方面.
和傳道書所說的虛空不同.
傳道書常說虛空的虛空.
傳道書的虛空.
是強調短暫.
好像一陣煙就過去了.
是指一些事物.
不可以永遠存留的意思.
不是指沒有價值.
而是短暫.
但是詩篇第127篇的.
往然是指沒有價值.
經文說.
若不是耶和華建造房屋.

$^{121}$建造的人就往然勞力.
若不是耶和華看守城池.
看守的人就往然警醒.
你們清晨早起.
夜晚安歇.
吃勞碌得來的飯.
本是往然.
什麼不是往然呢?.
什麼是有價值的呢?.
房屋和城池.
或者可以繼續存留.
但是否有價值呢?.
神的心意是否只是放在.
聖殿或教會的硬件上呢?.
經文又提到.
惟有耶和華.
所親愛的.
必叫他安然睡覺.
這段經文我就很大掙扎了.
有時睡不著.
或者失眠的時候.
我就跟神爭辯.
說你真的不愛我了.
因為聖經說.
惟有耶和華所親愛的.
必叫他安然睡覺.
我不安然睡覺 你就不愛我了.
原來不是這個意思.
安然睡覺的意思是.
必在睡覺的時候給他.
就算你睡覺的時候.
什麼都不要做.
神都會給你的.
因為人去勞苦尋求.
糧食,安全,滿足.
都是耶和華.
在他所愛之人.
睡覺的時候.
神賜給他們.
意思是神賜給他所愛的人.

$^{161}$是他樂意做的.
並不是完全在乎.
他們有多努力.
多熱心.
因為即使在他們睡覺的時候.
耶和華都會將福氣賜給他們.
就好像鎖羅門.
在睡覺的時候.
神允許他有智慧.
讓他可以管治國家.
就像亞當.
在睡覺的時候.
神就從他身上拿一條肋骨.
做一個女人給他.
成為他的配偶.
所以睡覺的時候.
好像什麼都沒有做的時候.
神就祝福你.
少年不是我們說的.
少年十五,二十歲.
猶太人認為少年是四十歲以下.
年輕力壯的時候就叫少年.
因為少年是一種武器.
可以防守也可以進攻.
箭是要磨練的.
你要把樹枝砍下來.
削直它.
然後削尖頭.
再加個箭頭.
再加羽毛.
這樣才會成為箭.
城門口是一些重要的地方.
和城池相呼應.
受敵就是常常來攻擊別人.
不是來不及攻擊你.
而是經常來攻擊你.
不至於羞愧.
贏得來又不輸人.
又不輸陣.
不至於羞愧.

$^{201}$我們略略看看.
這段經文裡.
一些自詞的解釋.
回看我們教會裡.
神的恩典.
在我們教會的兒童部.
或者兒童事工裡.
回首過去.
看到神的恩典很夠用.
人是不能夠自誇.
在2010年.
我們開始Avana事工.
第一次參加基礎導師訓練的人.
已經有28個.
全部一次過考試合格.
成為第一批Avana導師.
神恩供應.
超乎我們所想所求.
當時報名Avana的小朋友.
有100個.
而且隨後的日子.
出席率一直維持在.
90\%以上.
即使遇到學校考試的季節.
都微跌了3\%.
若不是耶和華建造房屋.
建造的人.
就枉然勞力.
在2012年4月開始.
教會加設了.
週六黃昏堂崇拜.
第二年我們Avana.
開設了禮拜六的營隊.
與週六的黃昏堂崇拜.
同步進行.
方便家長放下小朋友.
參加崇拜聚會.
在那個時候.
週六和週日.
Avana的小朋友.

$^{241}$一共超過150位.
另外還有十多位.
在我們的後補名單.
接著的日子.
神繼續賜我們更多的孩子.
和更多的事奉人手.
在2014年.
配合Avana的學生領袖試工.
即學生領袖試工.
青少年部天堂天家.
安排了6位青少年.
同步Avana服侍.
在2016年更加有.
14位學生領袖參與服侍.
其中一位學生領袖.
明天都會成家立室.
有些仍然在當中聚會.
參與不同的服侍.
有些去了外國定居.
繼續參加教會聚會.
和參加教會的事奉.
我們不能自誇.
若不是耶和華建造房屋.
建造的人.
就枉然勞力.
我們看看挑戰是什麼.
隨著隊員的人數增多.
後補名單亦有很多人.
可惜我們人手和地方.
都不足夠.
令我們兒童部的部員.
很頭痛很苦惱.
甚至我們需要勸喻.
那些不積極參加的小朋友離開.
因為我們後補名單的人很多.
如果有不珍惜的.
我們想將位置.
留給那些很想加入的.
我們在4月份的時候.
發出了這個勸喻.

$^{281}$以6月前為期限.
要他們改善.
如果有改善.
下一年就可以繼續參加.
當時因為資源不足.
而要拒絕一些人參加.
實在令人很痛苦.
因為主耶穌說.
讓小孩子到我這裡來.
不要禁止他們.
我們唯有交託給主.
若不是耶和華看守聖祠.
看守的人就枉然警醒.
在2015年.
和接下來的時間.
我們有一些明顯的人事變動.
有導師因為個人健康.
或家人患病需要照顧而退出.
另外有導師.
因為個別的需要.
移民的安排等等.
要離開事奉崗位.
而小朋友的生活又越來越繁忙.
補課補習放外活動.
比賽表演.
令到有不少隊員.
要退出阿灣.
家長當然有他們的考慮.
和決定.
而因為學業問題.
亦有小朋友要退出.
當時我們亦多了一些SEN.
即特殊學習需要的小朋友.
要照顧.
挑戰一浪接一浪.
2016年開始.
很明顯在兒童部.
有移民潮的出現.
2017年有四位小朋友.
因為移民退學.

$^{321}$那是接下來幾年最少移民的一年.
在那時候.
因為香港的政治氛圍.
和家長對教育前景的擔憂.
加強了對移民.
和移居外地的考慮.
接著的幾年.
都有小朋友因為移民要退學.
而高峰期.
就在2020年.
有11位小朋友移民.
另外有12位要離開.
離開的原因是因為他們搬屋.
離開了太古城這區.
因為學業的緣故.
有些是情緒的問題.
有些是需要轉教會聚會.
因為家長裡有些政治立場.
是大家很對立的.
家長就帶著小朋友離開了教會.
有些就回國內生活.
接著陸陸續續.
每一年都有小朋友.
因為移民要退學.
我們也有些導師在這段日子.
移民了或者回到.
他們已經移民的國家生活.
每一位這樣離去的小朋友.
或者導師.
我們都會有祈禱祝福禮.
或去或留.
我們都彼此祝福.
面對移民潮和不可知的大環境.
我都勸喻.
勸勉我們的導師.
要時時警醒.
趁著白日多作工.
珍惜事奉主的機會.
讓所有參加過.
Avana的小朋友.

$^{361}$至少都聽過一次福音.
若不是耶和華.
看守聖詞.
看守的人就枉然警醒.
2020年.
新冠疫情爆發.
我們經歷了停課復課.
停課復課的一年.
在聚會裡又要保持社交距離.
並且實行.
防疫措施和限聚令.
這是充滿不安和挑戰的一年.
2021年開始.
匹高山.
我們兩歲至三歲小朋友的營隊.
因應疫情的嚴重性.
決定暫停收生.
到現在為止.
我們都沒有再復辦匹高山.
我覺得有點遺憾.
根據教育局的指引.
我們不同年級的小朋友.
需要分階段來復課.
當時學校多數將課外活動.
由本來星期一至星期五.
放學後的時間.
改為星期六上午或下午舉行.
所以影響了.
星期六Avana小朋友的出席.
後來連星期日早上的時間.
都有不少比賽.
教會的學生.
在這些活動中.
被爭奪了.
2022年.
是疫情最嚴峻的一年.
新冠病毒第五波疫情爆發.
Omicron.
新冠病毒變異株.
相對於之前的四波疫情.

$^{401}$時間更長.
在2022一整年.
全世界都在它的陰霾下.
渡過.
到4至5月.
政府容許.
宗教場所重新開放.
教會可以恢復.
實體聚會.
但同時要推行.
「安心出行」和「疫苗通行證」.
所以在復課前.
我們要了解.
我們的導師.
是否已經接種了疫苗.
我們的助教是否已經接種了疫苗.
亦要關注小朋友.
接種疫苗的情況.
當時有幾位小朋友.
因為未完成接種疫苗.
而不能回教會聚會.
在11至12月.
有很多導師.
陸陸續續先後確診.
他的家人亦先後確診.
以致需要隔離.
我們人手當時很緊絀.
就算不是當時的導師.
都要回來幫忙照顧小朋友.
2022年.
全年我們流失了11位導師.
包括有移民的.
有搬屋的.
有需要照顧家人.
有需要休息.
亦需要轉換事奉崗位.
這一年有8位小朋友移民.
這個時候.
我們阿灣娜小朋友有79位.
若不是.

$^{441}$耶和華看守城池.
看守的人就往然驚醒.
到了2023年年初.
我們經歷了.
三年多的挑戰.
停課 復課 限聚令.
網上聚會 疫苗通行證.
社交距離等等.
三月份開始.
政府就撤銷了口罩令.
嚴峻的疫情似乎告一段落.
但是因為.
停止了實體上課.
有好幾年的時間.
小朋友再回來的時候.
就更加活躍.
我們的導師要更加吃力去維持紀律.
對隊員亦都需要.
更多的關注.
當時香港的教會.
機構 公司.
各行各業 各個階層.
應該都受到不少衝擊.
阿灣娜香港的.
學生領袖計劃.
這個學生領袖計劃.
亦都因為參加人數不足而取消.
而當時我們教會.
有一個少年人報名.
那一年有三個小朋友.
另一位導師移民.
2024年.
移民的小朋友有五位.
導師移民的有兩位.
移民潮逐漸穩定下來.
參加阿灣娜的小朋友.
有74個.
但出席率就不夠以前穩定.
自從疫情之後.
小朋友的出席率.

$^{481}$大約維持在.
70\%.
而遲到和早退的情況.
比以前嚴峻了.
若不是耶和華.
看守城池.
看守的人就往然警醒.
跟著跟大家分享一下.
我們怎樣在挑戰裡面.
成長和發展.
數算恩典.
經歷挑戰.
我們是在神的帶領之下.
與神同工.
從最初阿灣娜事公開始的第一年.
我們面對人手不足.
士氣的問題.
導師的毀身意願.
隊員的學習態度.
個別的學習差異.
特殊需要等等.
我們只有禱告和同心.
才可以應對.
我記得當時我和兒童部部長.
在辦公室裡面禱告.
我們禱告到流淚.
地板都濕了一大片.
幾乎要用地拖來拖地.
以後的日子.
我們常常藉著禱告來承托事奉.
經歷到主的恩典.
救援.
好像小朋友獻上的.
五餅二魚.
《普通的人》其中一首歌詞.
欠缺裡.
沒有還是夠.
有時要面對家長投訴.
記得有一次.
有家長投訴其他小朋友.

$^{521}$不乖,說他們回來.
說色情話.
所以不讓小朋友回來.
我們導師需要了解.
明白,關心和開導.
兒童事公的導師.
需要很多忍耐.
去承擔我們的事奉.
疫情期間.
令人感到最不安.
最有挑戰性.
我們阿灣娜的導師.
同心已靠主.
有智慧感迅速應變.
有限聚令的時候.
我們就馬上改為網上聚會.
而早到的暑假.
就催促我們.
舉辦全營參加全營投票.
暑假生活相片.
影片分享比賽.
當年暑假三月就開始.
不知你記不記得.
在復課的時候.
我很欣賞導師主動提出.
在聚會前做一個小時回來.
將小朋友有用的座椅.
逐張清潔消毒.
然後在地上畫好.
社交距離措施.
然後才讓小朋友回來聚會.
我為這班導師.
這班夥伴.
我的同工.
我感到十分之感恩.
兒女是耶和華.
所設的產業.
他們好像勇士手中的箭.
箭是要被磨練的.
在2017年暑假.

$^{561}$我們第一次舉辦.
澳門親子訪孫團.
有三日兩夜的時間.
在澳門.
我們帶小朋友去尋找.
孫教士的歷史足跡.
亦都聽聽現職孫教士.
他們的心智和見證.
大家亦都很實地.
經驗到澳門小區的文化.
又嘗試在.
百甲巢公園那裡.
派單張傳福音.
一年我們就讓小朋友.
在澳門訪孫裡承擔事奉.
他們參與出發前的預備.
和回港後的分享聚會.
見到他們成長.
導師都很開心.
亦都很安慰.
再一年的澳門訪孫.
我們就由家長擔任領隊.
小朋友協助.
在2020年.
我們藉著阿灣娜十週年慶祝.
雖然沒有實體的聚會.
但是就凝聚了.
阿灣娜的隊員.
家長 導師 校友.
甚至遠在海外的.
都能夠在網上聚會裡.
聯繫起來.
一起去感恩.
疫情前我們TNT.
有音樂服務小組.
小朋友會定期練習.
用音樂恩賜來服侍神.
在阿灣娜聚會裡帶領.
見拜.
但疫情後我們發現.

$^{601}$小朋友回來的時候.
音樂或其他樂器的.
小朋友大幅減少.
而且有音樂恩賜的.
導師又移民.
因此我們就再沒有.
這個服務小組.
但我們又發覺.
多了小朋友學跳舞.
所以在慶祝阿灣娜十五週年.
我們就成立了兒童舞蹈敬拜隊.
由我們一位有舞蹈恩賜的.
導師帶領他們練習.
在教會唐慶裡.
在唐慶崇拜的聚會裡.
以舞蹈去呈獻.
敬拜上主.
去年即是2025年.
開始的時候.
兒童部也推廣一個.
彼此欣賞的運動.
每一次聚會完結.
每一位導師就會向.
至少一位導師和一位隊員.
表示欣賞或謝意.
以致可以達到.
彼此鼓勵的目的.
成為我們阿灣娜團隊裡.
彼此欣賞的文化氛圍.
很感恩.
一年多以來.
我們多了欣賞 少了批評.
多了感恩 少了埋怨.
我們團隊裡的導師質素.
越來越好.
神的心意.
不會只是放在教會或聖殿.
這些硬件上.
神所重視的.
是要建立和看守一個.

$^{641}$信仰純正 信心堅固的.
屬靈群體.
最重要的是我們需要發展這個群體.
你們清晨早起.
夜晚安歇.
吃勞碌得來的飯.
本是枉然.
惟有耶和華所親愛的.
必叫他安然睡覺.
原來神要賜福給我們.
不是因為我們的熱心.
我們的才能.
而是祂定義要賜福.
那些倚靠祂的人.
人主動做的事情.
如果不是按照神的心意.
就是枉然 毫無價值.
過去一段長時間.
我們兒童部經驗了實體聚會.
可以忽然被取消.
恢復聚會之後.
也有很多限制.
而教會裡相識多年.
共同成長的好朋友.
會因為移民要分開.
而小朋友也會面對這個衝擊.
讓我提議一些具體的行動.
教會要繼續發展兒童事工.
有什麼要注意呢?.
有什麼具體行動呢?.
我會用一些具體的見證.
跟大家分享.
首先是神的話.
我們一定要堅持.
教導小朋友學習聖經.
學習神的話.
有一次我在諾茲故事講訊息.
講到撒旦.
引誘女人吃.
分別善惡樹的果子.

$^{681}$題目很沉重.
要遲早犯罪.
我就用了伊托米子.
來做引言.
大家知道伊托米子是什麼嗎?.
太空油.
我就說那些壞人.
引誘人吸毒.
就找他們做角色扮演.
有人扮壞人.
引誘另一組小朋友.
被引誘要吃太空油.
就看他們會不會說服到對方.
當然他們.
在當中有很多辯論.
反應很激烈.
我很欣賞他們很投入.
但亦有些不合作的情況.
引言部分過了.
我要講聖經故事.
我先讀出經文.
然後再講.
我打開聖經.
讀《創世記》三章一節五節.
很出奇.
我一開始讀經.
他們全部靜了.
本來在爭拗.
說分數.
威力無窮.
我沒有叫他們安靜.
聽我讀經.
一開始讀經.
他們全部靜了.
一粒聲都沒有.
很感動.
很震撼.
出於神的說話.
沒有一句不大能力.
所以我們一定要堅持教小朋友聖經.

$^{721}$第二就是詩歌.
要教導小朋友學習詩歌.
尤其是旋律優美.
歌詞內容豐富的詩歌.
有一次我負責.
Avanna的詩歌敬拜.
我邀請了一位隊員做詩琴.
我選了那首詩歌.
《這是天賦世界》.
其中.
在第三節有歌詞.
黑暗勢力雖然猖狂.
我就想.
那個隊員是小學生.
他知不知道什麼叫黑暗勢力雖然猖狂.
所以練習的時候.
我就問這位小詩琴.
聽到歌詞的說話.
他想都不想.
很直接的說.
我感受到黑暗勢力.
我就問他是什麼.
他說是罪惡的勢力.
因為我自己很容易犯罪.
很直接.
很簡單.
很率真.
如果你用這句說話去問成年人.
用這個問題去問成年人.
他們可能想到的.
就是別人的罪.
是別人的勢力.
而不是自己內心的污穢和黑暗.
尤其是當時.
香港社會的政治氛圍.
對立的情況是非常之嚴重.
第三.
就是要禱告.
禱告可以讓我們.
經歷和經驗到.

$^{761}$神的同在.
在一次TNT的搜查時間.
我負責代課.
因為本身有個導師.
那天來不了.
我負責代課.
我就不會跟著課程去做.
我就很想教他們祈禱.
起初.
其中一個小朋友說.
不用學祈禱.
每晚媽媽都和我祈禱.
我知道怎樣祈禱.
但另一個小朋友就說.
我們聽聽怎樣祈禱.
於是我就和他們.
做了一個.
一個靈穴.
或者禱告.
用神讀的方式去禱告.
我就讀了經文.
然後叫他們投入當中的角色.
去感受自己的感受.
和向神禱告.
我就用浪子的比喻.
這個經文跟他們說.
我就用口語化.
將浪子的比喻的故事.
讀出來.
大家也知道浪子的比喻是說什麼.
小朋友拿了爸爸的身家.
出去花天酒地.
用光了錢.
又遇到饑荒.
又餵豬.
我就叫他們投入當中的角色.
小朋友竟然.
投入到浪子的角色.
他說.
我聞到豬吃的食物的味道.

$^{801}$我看到.
食物上面.
有蟲在爬.
還有很臭的味道.
很骯髒.
另一個小朋友.
他感受到肥牛的美味.
覺得很開心.
我當時很詫異.
因為我未試過帶小學生.
三,四年級.
去做靈穴.
我感受到聖靈的工作.
他們能夠報獲神的愛.
亦感受到自己犯了罪.
但是能夠得到.
神的寬恕.
我帶領他們將這個感受.
告訴神,亦向神感恩.
再和大家分享另一個見證.
2025年1月.
我們兒童部有一個導師退休會.
我帶領導師.
做靈穴.
經文是路加斐恩五章一至十一節.
是講到主耶穌叫門徒.
去深水的地方.
下網打魚.
本來他們一整晚都打不到魚.
但因為他們聽主耶穌的吩咐.
結果打到很多魚.
打到很多魚上來網幾乎裂開.
但又沒有裂開.
那次退休會.
CUPBIS是最多導師參加的.
總共有六位.
我們CUPBIS的導師.
一向很細心,很有愛心.
他們很關心小朋友.
要給予他們最好的照顧.

$^{841}$他們對自己的要求亦很高.
有時我都會擔心.
他們給予自己太大壓力.
所以即使當時的等待清單.
有七個小朋友.
他們寧願保持在質方面做得好.
都不願意在量方面.
因為多收小朋友影響質數.
我曾經費盡唇舌.
說服他們收起來吧.
不要緊的,你們做得到.
不用要求自己這麼高壓力.
但他們都沒有辦法說服他們.
但藉著這次退休.
神親自跟他們說話.
他們不約而同.
一起都有這個領受.
主耶穌叫他們收起小朋友.
收起等待清單的小朋友.
結果我們在三月份開始.
在CUPBIS多開了一組.
後來神都給我們額外的資源.
有些家長願意來幫忙做助教.
我們就可以多開一組.
順利.
分享到這裡.
我真的很希望我們教會.
能夠繼續為主栽培有信仰.
有信心的下一代.
藉著聖經的教導.
詩歌的學習.
禱告的操練.
讓孩子們能夠被磨練.
成為勇士手中的箭.
被主使用.
為主打成仗.
正如剛才小朋友所唱的詩歌.
堅守崗位.
去打成仗.
讓我以聖經的話語.

$^{881}$作為總結.
在哥林多前書三章七節說.
可見栽種的算不得什麼.
挑換的也算不得什麼.
只在那叫他生長的神.
只希望我所分享的.
能夠將榮耀歸給神.
詩篇127篇.
一節二節.
或者你也會背.
可以一起讀.
若不是耶和華建造房屋.
建造的人就枉然勞力.
若不是耶和華看守城池.
看守的人就枉然警醒.
你們清晨早起.
夜晚安歇.
吃勞碌得來的飯.
本是枉然.
唯有耶和華所親愛的.
必叫他安然睡覺.
我請大家一起來禱告.
教會兒童事工的過去.
現在將來.
都是在你的計劃裡.
我們是在神的帶領之下.
與神同工.
懇求主.
讓教會弟兄姊妹.
繼續為主栽培.
有信仰有信心的下一代.
使他們能夠被主使用.
為主打聖仗.
奉獻給教會的弟兄姊妹.
奉主耶穌基督的聖名祈求.
阿們.
\newpage



\section{民數記 7:1-89}
\label{sec:031832}
\textbf{合神心意的奉獻}~~~~林誠信牧師
\newline
\newline
連結: \href{http://www.church.com.hk/acms/content.asp?site=cdc\&op=show\&type=product\&code=031832\&layout=sermon}{\texttt{031832}} ~~~~ 語音日期: 2026-03-15
\newline
\newline
\hyperref[sec:031831]{\small{< < < PREV SERMON < < <}}
~
\hyperref[sec:index_chronic]{\small{[返順時目]}}
~
\hyperref[sec:031833]{\small{> > > NEXT SERMON > > >}}
\newline
\newline
民數記 7:1-89
\newline
\begin{longtable}{cl}
\hline
\hline
章節 & 經文 (和合本修訂版)\\
\hline
7:1 & \begin{tabularx}{0.7\textwidth}{X} 摩西豎立帳幕後，就用膏抹了帳幕，使它分別為聖，又用膏抹其中的一切器具，以及祭壇和壇上的一切器具，使它們分別為聖。 \end{tabularx} \\ \\ \relax
7:2 & \begin{tabularx}{0.7\textwidth}{X} 以色列的領袖，各父家的家長，都前來奉獻。他們是各支派的領袖，管理那些被數的人。 \end{tabularx} \\ \\ \relax
7:3 & \begin{tabularx}{0.7\textwidth}{X} 他們把自己的供物送到耶和華面前，就是六輛篷車和十二頭公牛。每兩個領袖奉獻一輛車，每個領袖奉獻一頭牛。他們把這些都帶到帳幕前。 \end{tabularx} \\ \\ \relax
7:4 & \begin{tabularx}{0.7\textwidth}{X} 耶和華對摩西說： \end{tabularx} \\ \\ \relax
7:5 & \begin{tabularx}{0.7\textwidth}{X} 「你要從他們收下這些，作為會幕事奉的用途，照著利未人所事奉的交給他們各人。」 \end{tabularx} \\ \\ \relax
7:6 & \begin{tabularx}{0.7\textwidth}{X} 於是摩西收了車和牛，交給利未人。 \end{tabularx} \\ \\ \relax
7:7 & \begin{tabularx}{0.7\textwidth}{X} 他把兩輛車和四頭牛，照著革順子孫所事奉的交給他們， \end{tabularx} \\ \\ \relax
7:8 & \begin{tabularx}{0.7\textwidth}{X} 又把四輛車和八頭牛，照著米拉利子孫所事奉的交給他們。他們都在亞倫祭司的兒子以他瑪的手下。 \end{tabularx} \\ \\ \relax
7:9 & \begin{tabularx}{0.7\textwidth}{X} 但沒有交給哥轄子孫任何東西，因為他們所事奉的是聖物，必須抬在肩頭上。 \end{tabularx} \\ \\ \relax
7:10 & \begin{tabularx}{0.7\textwidth}{X} 用膏抹祭壇的那一天，眾領袖前來為獻壇奉獻；眾領袖都在祭壇前獻供物。 \end{tabularx} \\ \\ \relax
7:11 & \begin{tabularx}{0.7\textwidth}{X} 耶和華對摩西說：「眾領袖為獻壇奉獻供物，每天要有一個領袖前來奉獻。」 \end{tabularx} \\ \\ \relax
7:12 & \begin{tabularx}{0.7\textwidth}{X} 第一天獻供物的是猶大支派的亞米拿達的兒子拿順。 \end{tabularx} \\ \\ \relax
7:13 & \begin{tabularx}{0.7\textwidth}{X} 他的供物是：一個重一百三十舍客勒的銀盤，一個重七十舍客勒的銀碗，都是按照聖所的舍客勒，裡面盛滿了調油的細麵作素祭； \end{tabularx} \\ \\ \relax
7:14 & \begin{tabularx}{0.7\textwidth}{X} 一個重十舍客勒的金碟子，盛滿了香； \end{tabularx} \\ \\ \relax
7:15 & \begin{tabularx}{0.7\textwidth}{X} 一頭公牛犢、一隻公綿羊、一隻一歲的小公羊作燔祭； \end{tabularx} \\ \\ \relax
7:16 & \begin{tabularx}{0.7\textwidth}{X} 一隻公山羊作贖罪祭； \end{tabularx} \\ \\ \relax
7:17 & \begin{tabularx}{0.7\textwidth}{X} 兩頭公牛、五隻公綿羊、五隻公山羊、五隻一歲的小公羊作平安祭。這是亞米拿達的兒子拿順的供物。 \end{tabularx} \\ \\ \relax
7:18 & \begin{tabularx}{0.7\textwidth}{X} 第二天來獻的是以薩迦的領袖，蘇押的兒子拿坦業。 \end{tabularx} \\ \\ \relax
7:19 & \begin{tabularx}{0.7\textwidth}{X} 他獻為供物的是：一個重一百三十舍客勒的銀盤，一個重七十舍客勒的銀碗，都是按照聖所的舍客勒，裡面盛滿了調油的細麵作素祭； \end{tabularx} \\ \\ \relax
7:20 & \begin{tabularx}{0.7\textwidth}{X} 一個重十舍客勒的金碟子，盛滿了香； \end{tabularx} \\ \\ \relax
7:21 & \begin{tabularx}{0.7\textwidth}{X} 一頭公牛犢、一隻公綿羊、一隻一歲的小公羊作燔祭； \end{tabularx} \\ \\ \relax
7:22 & \begin{tabularx}{0.7\textwidth}{X} 一隻公山羊作贖罪祭； \end{tabularx} \\ \\ \relax
7:23 & \begin{tabularx}{0.7\textwidth}{X} 兩頭公牛、五隻公綿羊、五隻公山羊、五隻一歲的小公羊作平安祭。這是蘇押的兒子拿坦業的供物。 \end{tabularx} \\ \\ \relax
7:24 & \begin{tabularx}{0.7\textwidth}{X} 第三天是西布倫子孫的領袖，希倫的兒子以利押。 \end{tabularx} \\ \\ \relax
7:25 & \begin{tabularx}{0.7\textwidth}{X} 他的供物是：一個重一百三十舍客勒的銀盤，一個重七十舍客勒的銀碗，都是按照聖所的舍客勒，裡面盛滿了調油的細麵作素祭； \end{tabularx} \\ \\ \relax
7:26 & \begin{tabularx}{0.7\textwidth}{X} 一個重十舍客勒的金碟子，盛滿了香； \end{tabularx} \\ \\ \relax
7:27 & \begin{tabularx}{0.7\textwidth}{X} 一頭公牛犢、一隻公綿羊、一隻一歲的小公羊作燔祭； \end{tabularx} \\ \\ \relax
7:28 & \begin{tabularx}{0.7\textwidth}{X} 一隻公山羊作贖罪祭； \end{tabularx} \\ \\ \relax
7:29 & \begin{tabularx}{0.7\textwidth}{X} 兩頭公牛、五隻公綿羊、五隻公山羊、五隻一歲的小公羊作平安祭。這是希倫的兒子以利押的供物。 \end{tabularx} \\ \\ \relax
7:30 & \begin{tabularx}{0.7\textwidth}{X} 第四天是呂便子孫的領袖，示丟珥的兒子以利蓿。 \end{tabularx} \\ \\ \relax
7:31 & \begin{tabularx}{0.7\textwidth}{X} 他的供物是：一個重一百三十舍客勒的銀盤，一個重七十舍客勒的銀碗，都是按照聖所的舍客勒，裡面盛滿了調油的細麵作素祭； \end{tabularx} \\ \\ \relax
7:32 & \begin{tabularx}{0.7\textwidth}{X} 一個重十舍客勒的金碟子，盛滿了香； \end{tabularx} \\ \\ \relax
7:33 & \begin{tabularx}{0.7\textwidth}{X} 一頭公牛犢、一隻公綿羊、一隻一歲的小公羊作燔祭； \end{tabularx} \\ \\ \relax
7:34 & \begin{tabularx}{0.7\textwidth}{X} 一隻公山羊作贖罪祭； \end{tabularx} \\ \\ \relax
7:35 & \begin{tabularx}{0.7\textwidth}{X} 兩頭公牛、五隻公綿羊、五隻公山羊、五隻一歲的小公羊作平安祭。這是示丟珥的兒子以利蓿的供物。 \end{tabularx} \\ \\ \relax
7:36 & \begin{tabularx}{0.7\textwidth}{X} 第五天是西緬子孫的領袖，蘇利沙代的兒子示路蔑。 \end{tabularx} \\ \\ \relax
7:37 & \begin{tabularx}{0.7\textwidth}{X} 他的供物是：一個重一百三十舍客勒的銀盤，一個重七十舍客勒的銀碗，都是按照聖所的舍客勒，裡面盛滿了調油的細麵作素祭； \end{tabularx} \\ \\ \relax
7:38 & \begin{tabularx}{0.7\textwidth}{X} 一個重十舍客勒的金碟子，盛滿了香； \end{tabularx} \\ \\ \relax
7:39 & \begin{tabularx}{0.7\textwidth}{X} 一頭公牛犢、一隻公綿羊、一隻一歲的小公羊作燔祭； \end{tabularx} \\ \\ \relax
7:40 & \begin{tabularx}{0.7\textwidth}{X} 一隻公山羊作贖罪祭； \end{tabularx} \\ \\ \relax
7:41 & \begin{tabularx}{0.7\textwidth}{X} 兩頭公牛、五隻公綿羊、五隻公山羊、五隻一歲的小公羊作平安祭。這是蘇利沙代的兒子示路蔑的供物。 \end{tabularx} \\ \\ \relax
7:42 & \begin{tabularx}{0.7\textwidth}{X} 第六天是迦得子孫的領袖，丟珥的兒子以利雅薩。 \end{tabularx} \\ \\ \relax
7:43 & \begin{tabularx}{0.7\textwidth}{X} 他的供物是：一個重一百三十舍客勒的銀盤，一個重七十舍客勒的銀碗，都是按照聖所的舍客勒，裡面盛滿了調油的細麵作素祭； \end{tabularx} \\ \\ \relax
7:44 & \begin{tabularx}{0.7\textwidth}{X} 一個重十舍客勒的金碟子，盛滿了香； \end{tabularx} \\ \\ \relax
7:45 & \begin{tabularx}{0.7\textwidth}{X} 一頭公牛犢、一隻公綿羊、一隻一歲的小公羊作燔祭； \end{tabularx} \\ \\ \relax
7:46 & \begin{tabularx}{0.7\textwidth}{X} 一隻公山羊作贖罪祭； \end{tabularx} \\ \\ \relax
7:47 & \begin{tabularx}{0.7\textwidth}{X} 兩頭公牛、五隻公綿羊、五隻公山羊、五隻一歲的小公羊作平安祭。這是丟珥的兒子以利雅薩的供物。 \end{tabularx} \\ \\ \relax
7:48 & \begin{tabularx}{0.7\textwidth}{X} 第七天是以法蓮子孫的領袖，亞米忽的兒子以利沙瑪。 \end{tabularx} \\ \\ \relax
7:49 & \begin{tabularx}{0.7\textwidth}{X} 他的供物是：一個重一百三十舍客勒的銀盤，一個重七十舍客勒的銀碗，都是按照聖所的舍客勒，裡面盛滿了調油的細麵作素祭； \end{tabularx} \\ \\ \relax
7:50 & \begin{tabularx}{0.7\textwidth}{X} 一個重十舍客勒的金碟子，盛滿了香； \end{tabularx} \\ \\ \relax
7:51 & \begin{tabularx}{0.7\textwidth}{X} 一頭公牛犢、一隻公綿羊、一隻一歲的小公羊作燔祭； \end{tabularx} \\ \\ \relax
7:52 & \begin{tabularx}{0.7\textwidth}{X} 一隻公山羊作贖罪祭； \end{tabularx} \\ \\ \relax
7:53 & \begin{tabularx}{0.7\textwidth}{X} 兩頭公牛、五隻公綿羊、五隻公山羊、五隻一歲的小公羊作平安祭。這是亞米忽的兒子以利沙瑪的供物。 \end{tabularx} \\ \\ \relax
7:54 & \begin{tabularx}{0.7\textwidth}{X} 第八天是瑪拿西子孫的領袖，比大蓿的兒子迦瑪列。 \end{tabularx} \\ \\ \relax
7:55 & \begin{tabularx}{0.7\textwidth}{X} 他的供物是：一個重一百三十舍客勒的銀盤，一個重七十舍客勒的銀碗，都是按照聖所的舍客勒，裡面盛滿了調油的細麵作素祭； \end{tabularx} \\ \\ \relax
7:56 & \begin{tabularx}{0.7\textwidth}{X} 一個重十舍客勒的金碟子，盛滿了香； \end{tabularx} \\ \\ \relax
7:57 & \begin{tabularx}{0.7\textwidth}{X} 一頭公牛犢、一隻公綿羊、一隻一歲的小公羊作燔祭； \end{tabularx} \\ \\ \relax
7:58 & \begin{tabularx}{0.7\textwidth}{X} 一隻公山羊作贖罪祭； \end{tabularx} \\ \\ \relax
7:59 & \begin{tabularx}{0.7\textwidth}{X} 兩頭公牛、五隻公綿羊、五隻公山羊、五隻一歲的小公羊作平安祭。這是比大蓿的兒子迦瑪列的供物。 \end{tabularx} \\ \\ \relax
7:60 & \begin{tabularx}{0.7\textwidth}{X} 第九天是便雅憫子孫的領袖，基多尼的兒子亞比但。 \end{tabularx} \\ \\ \relax
7:61 & \begin{tabularx}{0.7\textwidth}{X} 他的供物是：一個重一百三十舍客勒的銀盤，一個重七十舍客勒的銀碗，都是按照聖所的舍客勒，裡面盛滿了調油的細麵作素祭； \end{tabularx} \\ \\ \relax
7:62 & \begin{tabularx}{0.7\textwidth}{X} 一個重十舍客勒的金碟子，盛滿了香； \end{tabularx} \\ \\ \relax
7:63 & \begin{tabularx}{0.7\textwidth}{X} 一頭公牛犢、一隻公綿羊、一隻一歲的小公羊作燔祭； \end{tabularx} \\ \\ \relax
7:64 & \begin{tabularx}{0.7\textwidth}{X} 一隻公山羊作贖罪祭； \end{tabularx} \\ \\ \relax
7:65 & \begin{tabularx}{0.7\textwidth}{X} 兩頭公牛、五隻公綿羊、五隻公山羊、五隻一歲的小公羊作平安祭。這是基多尼的兒子亞比但的供物。 \end{tabularx} \\ \\ \relax
7:66 & \begin{tabularx}{0.7\textwidth}{X} 第十天是但子孫的領袖，亞米沙代的兒子亞希以謝。 \end{tabularx} \\ \\ \relax
7:67 & \begin{tabularx}{0.7\textwidth}{X} 他的供物是：一個重一百三十舍客勒的銀盤，一個重七十舍客勒的銀碗，都是按照聖所的舍客勒，裡面盛滿了調油的細麵作素祭； \end{tabularx} \\ \\ \relax
7:68 & \begin{tabularx}{0.7\textwidth}{X} 一個重十舍客勒的金碟子，盛滿了香； \end{tabularx} \\ \\ \relax
7:69 & \begin{tabularx}{0.7\textwidth}{X} 一頭公牛犢、一隻公綿羊、一隻一歲的小公羊作燔祭； \end{tabularx} \\ \\ \relax
7:70 & \begin{tabularx}{0.7\textwidth}{X} 一隻公山羊作贖罪祭； \end{tabularx} \\ \\ \relax
7:71 & \begin{tabularx}{0.7\textwidth}{X} 兩頭公牛、五隻公綿羊、五隻公山羊、五隻一歲的小公羊作平安祭。這是亞米沙代的兒子亞希以謝的供物。 \end{tabularx} \\ \\ \relax
7:72 & \begin{tabularx}{0.7\textwidth}{X} 第十一天是亞設子孫的領袖，俄蘭的兒子帕結。 \end{tabularx} \\ \\ \relax
7:73 & \begin{tabularx}{0.7\textwidth}{X} 他的供物是：一個重一百三十舍客勒的銀盤，一個重七十舍客勒的銀碗，都是按照聖所的舍客勒，裡面盛滿了調油的細麵作素祭； \end{tabularx} \\ \\ \relax
7:74 & \begin{tabularx}{0.7\textwidth}{X} 一個重十舍客勒的金碟子，盛滿了香； \end{tabularx} \\ \\ \relax
7:75 & \begin{tabularx}{0.7\textwidth}{X} 一頭公牛犢、一隻公綿羊、一隻一歲的小公羊作燔祭； \end{tabularx} \\ \\ \relax
7:76 & \begin{tabularx}{0.7\textwidth}{X} 一隻公山羊作贖罪祭； \end{tabularx} \\ \\ \relax
7:77 & \begin{tabularx}{0.7\textwidth}{X} 兩頭公牛、五隻公綿羊、五隻公山羊、五隻一歲的小公羊作平安祭。這是俄蘭的兒子帕結的供物。 \end{tabularx} \\ \\ \relax
7:78 & \begin{tabularx}{0.7\textwidth}{X} 第十二天是拿弗他利子孫的領袖，以南兒子亞希拉。 \end{tabularx} \\ \\ \relax
7:79 & \begin{tabularx}{0.7\textwidth}{X} 他的供物是：一個重一百三十舍客勒的銀盤，一個重七十舍客勒的銀碗，都是按照聖所的舍客勒，裡面盛滿了調油的細麵作素祭； \end{tabularx} \\ \\ \relax
7:80 & \begin{tabularx}{0.7\textwidth}{X} 一個重十舍客勒的金碟子，盛滿了香； \end{tabularx} \\ \\ \relax
7:81 & \begin{tabularx}{0.7\textwidth}{X} 一頭公牛犢、一隻公綿羊、一隻一歲的小公羊作燔祭； \end{tabularx} \\ \\ \relax
7:82 & \begin{tabularx}{0.7\textwidth}{X} 一隻公山羊作贖罪祭； \end{tabularx} \\ \\ \relax
7:83 & \begin{tabularx}{0.7\textwidth}{X} 兩頭公牛、五隻公綿羊、五隻公山羊、五隻一歲的小公羊作平安祭。這是以南的兒子亞希拉的供物。 \end{tabularx} \\ \\ \relax
7:84 & \begin{tabularx}{0.7\textwidth}{X} 用膏抹祭壇的那一天，以色列的眾領袖為獻壇所獻的是：銀盤十二個、銀碗十二個、金碟子十二個； \end{tabularx} \\ \\ \relax
7:85 & \begin{tabularx}{0.7\textwidth}{X} 一個銀盤重一百三十，一個碗七十。一切器皿的銀子，按照聖所的舍客勒共二千四百舍客勒。 \end{tabularx} \\ \\ \relax
7:86 & \begin{tabularx}{0.7\textwidth}{X} 十二個金碟子盛滿了香，按照聖所的舍客勒，一個碟子重十舍客勒，所有碟子的金子共一百二十舍客勒。 \end{tabularx} \\ \\ \relax
7:87 & \begin{tabularx}{0.7\textwidth}{X} 作燔祭的共有公牛十二頭、公羊十二隻、一歲的小公羊十二隻，和同獻的素祭，以及作贖罪祭的公山羊十二隻； \end{tabularx} \\ \\ \relax
7:88 & \begin{tabularx}{0.7\textwidth}{X} 作平安祭的共有公牛二十四頭、公綿羊六十隻、公山羊六十隻、一歲的小公羊六十隻。這就是用膏抹壇之後，為獻壇的奉獻。 \end{tabularx} \\ \\ \relax
7:89 & \begin{tabularx}{0.7\textwidth}{X} 摩西進會幕要與耶和華說話的時候，聽見法櫃的櫃蓋以上二基路伯中間有對他說話的聲音。耶和華向他說話。 \end{tabularx} \\ \\
[1ex]
\hline
\hline
\end{longtable}
$^{1}$我們看這段聖經的時候.
我們必須要明白.
它是緊接前文的.
是頭六章是以色列起行前的準備.
百姓和安營的位置都有關係.
第五章亦都提到.
他們納潔淨的條例.
是照著神的吩咐.
這個對他們在抗疫路是非常重要.
因為那個和他們的衛生很有關係.
亦都是給他們一個提醒.
神的聖潔人不能冒犯.
我們必須是尊神為聖.
我們都要看到我們的生活.
必須要分別為聖.
第六章其實就是講到.
這個分別為聖歸於神.
其中就是那世爾人的那個條例.
和最後就是大祭司他們的祝福,祈禱.
在跟著我們第七章之後.
我們都會去到第十章都會看到.
神是透過這個會幕.
日間有雲彩.
來到在他們的會幕當中停留.
表示他們不要起行了.
他們就停在這個地方.
夜間亦都在雲彩裡面有火.
所謂火柱.
這些都是很清楚神的同在.
和神的引導.
何時動呢.
就是雲彩要行.
代表神要吩咐他們要行.
他一行的時候.
就整個會幕的基建,組件全部拆出來.
所以工程都是很多事情要做的.
你要明白這個背景.
我們之前其實是有看過.
為什麼利美那三個兒子要分配.
有不同的在會幕.

$^{41}$和這個聖所,至聖所的聖物.
是很清楚的來有吩咐.
他亦都要吩咐他們要守禦節.
紀念他們從埃及離開這個遺老之家.
他們從西乃這個地方.
其實原本是要進入這個.
流乃與物之地就是這個迦南地.
原本很快的.
但是就因為他們的信心不足.
他們接不上神的應許.
所以在第十章到二十五章.
是記載到他們在曠野裡頭.
不斷的信心軟弱,失敗.
亦都經歷上帝多方面的管教.
在二十六章到三十六章.
就是他們在摩訶平原要重新去數點民數.
即是百姓一次.
當時這個新的數點.
就是新的一代起來了.
但是舊的一代.
就亦都很多已經在曠野路裡面.
他們都已經死去.
剩下就是迦納和約書亞.
後來摩西都被神接去.
在這個尼泊爾山上面.
我們看到整件事.
都是要提醒我們今天.
很好的.
我們來到去跟隨神.
第七章你可以這樣來分的.
這幾個段落.
首先就是頭九節.
十二個支派領袖奉獻給回眸的物件.
是什麼呢.
是說回眸.
而跟著那個段落.
是說十二個領袖奉獻給祭壇的物件又是什麼呢.
因為重複十二次.
這個祭壇的奉獻禮連續十二日.
每日一日.

$^{81}$每日就是派一個領袖出來.
這個就給我們看得很清楚.
第三就是十二個支派為祭壇奉獻的總數.
你就將他們每一個支派其實所奉獻在祭壇的物件.
基本上是一樣的.
那些數量你乘十二.
這些數目是不會有錯誤的.
所以是很清楚的.
最後就是一節聖經.
神在他們完成了這個祭壇的奉獻之後.
摩西進入這個聖所.
在至聖所之間那個布簾.
他聽見神對他們說話.
神在我們的敬拜是不單止很渴望和我們相遇.
他也帶有啟示性的說話.
去光照和提醒我們.
所以每一次你都不要輕忽.
不只是周末你回來教會.
港福家這個聖殿裡面.
其實今天神的殿在哪裡.
以新約的聖徒來說.
聖靈就是住在我們裡面.
我們的身體就是聖靈的殿.
所以我們要保養姑息之餘.
我們要分別為聖.
你不可以將這個身體放縱去犯罪.
這個是很重要的一件事.
你每天這個靈修祈禱.
不是一種習慣.
而是每天周而復始性.
你是有一個學務.
學務就是我要與神親近.
我們門女為何這麼興奮.
就是每天早上我們都很珍惜.
這段時間與神獨處.
真是me time.
特別我們牧師傳道人.
真的打開手機.
就已經有很多東西等著你.
所以明天早上我要很小心開電話.

$^{121}$我先開了靈修經文.
就好好祈禱.
親近了上帝.
有神的話語光照我們.
那天多忙也好.
你發覺你有力量.
繼續向前行.
求主幫助我們.
我們一起祈禱.
天父當我們打開你的聖言.
求主的靈教導我們.
讓講的聽的我們同感一靈.
同德造就.
逃避你自己的因言.
光照我們.
叫我們的生命.
不斷不斷要更像你.
能夠在這個地上.
能夠榮神益人.
我們同心禱告.
奉主名求.
阿們.
好第一段經文.
在這個頭兩節裡面.
我請大家跟我一起讀.
好第一節講得很清楚.
摩西立這個帳幕後.
發生的事.
這個其實說用高油去抹那個帳幕.
和祭壇.
以及聖所.
致聖所一切的器具.
全部都是有一個目的.
就是要將這一切的東西.
分別為聖.
這個其實在.
前一卷書.
都是兩卷書.
這個出埃及記.
四十章在結尾的時候.

$^{161}$就講摩西是照著神的吩咐.
在出埃及之後.
第二年的正月.
第一日.
就立起這個帳幕.
就用高油做了這些事情.
同民數記這裡所記載.
是一模一樣.
所以在時間上我們很確定.
這個表示就是他們初初.
建立好這個帳幕.
第一次要啟用祭壇.
要以後一路都要獻祭.
在這個過程裡面.
要做些什麼呢.
就是我們這一章有很詳細.
和我們交代.
在第二節就是說.
是以色列的領袖.
各父家的家長.
代表著十二個支派那些領袖.
他們將那些禮物帶來.
代表他們所管理的.
那支派的人帶來.
但請大家留意第二節這裡所講.
其實沒有明顯說到.
是神吩咐他們.
將這些禮物送來.
在說會幕那些.
沒有說的.
是他們自己想到.
我們這些會幕很多東西的.
只是你看圍著會幕.
這麼多這些板塊.
在搬運的時候.
其實很吃力的.
他們想了一些東西.
他們是自發性的.
來奉獻給上帝.
我們首先要學第一個原則.

$^{201}$我們今天看奉獻.
是甘心樂意的.
沒有人逼你.
不是去勉強.
但你看到一些需要.
你心裡有一個感動.
你自發的.
你願意就是獻上.
而且這個獻上.
我們看著是由領袖開始.
是由這12個支派的領袖族長.
他們代表神的子民.
他們獻上做榜樣.
愛和感恩.
就是這個奉獻的動機.
不是壓力.
不是不正確的動機.
如果有任何雜質的時候.
你是叫這個所謂奉獻.
無論是金錢.
或者是事奉的行動.
你就是帶有.
是一種神喜悅的.
我們需要在這方面謹慎.
我記得是很多年前.
在90年代.
那時候我還在加拿大牧羊教會的時候.
當時我們買了一個大貨倉.
是兩層高的一個建築物.
是很多神蹟騎士的.
其實以我們當時的財政的能力.
我們是買不起的這個地方.
但神竟然為我們.
完全開了路.
連那個賣家.
是全心全意幫我們教會.
去達成那個交易的.
所以我們很感恩.
這裡我不詳述了.
我曾經在其他講章提過這個見證.

$^{241}$但我今天要講的.
就是我們買了之後.
二樓是那些辦公室.
不用怎樣裝修了.
但是樓下是因為大貨倉.
我們要拆掉那些貨架.
要來做敬拜神的主禮堂.
就好像現在這個地方.
不過我們沒有那麼大.
那個可以坐二百多人.
我們後來都坐滿的.
整個教會.
我們願意去敬拜事奉.
願意去傳福音.
願意去建立神的家裡.
和每一個弟兄姊妹.
就好好一起無分彼此.
甘心樂意的擺上.
我記得當年.
我九零年代那時候是三十多歲.
很年輕.
現在有些人都說我四十多歲.
不過我說多謝.
你太過獎了.
可能要換一換老花眼鏡.
但其實我自己回想當日.
我是在教會做傳道服事.
但我兼任這個裝修的工程.
要監工 要施工.
我一個人怎麼做得來.
當然我是請那些判頭.
但是那些negotiation和execution.
其實全部每一個細節.
都其實花很多時間.
但是感恩日以繼夜.
無怨無悔.
沒有over time的.
我告訴大家.
你真的願意做神的工.
雖然我們現在教會有很多.

$^{281}$不好的policy.
就是不想我們神的僕人.
做到全殘.
每天兩節好了.
如果三節的話.
你補回一節.
但我告訴執事們.
其實我們想補.
不是不想補.
不過我們補不到.
我們三月三十一日.
還有一樣東西要埋單.
就是我們每年的年假.
你放不完.
我們要carry forward.
全所有奉獻.
我們不介意這件事.
但我自己覺得最重要的.
就是你是不是事奉有喜樂.
有感恩的心.
是你不去斤斤計較這件事.
當時當我做這個工程的時候.
我祈禱我知道.
你太看得起我.
因為我以前是工程師.
但是我現在做了.
祈禱傳道為事.
這個教會事實上是小的教會.
我兼任這件事.
我求神祈禱.
求神給我.
有弟兄姊妹和我一起同工.
當我一呼籲.
教會有很多年青人.
他們有些是在讀書的.
有些是在做事.
放完工.
他們回來幫忙去看裝修.
後來為了省錢的緣故.
我們見那個判頭.

$^{321}$你處理好那些甲板.
釘好整個禮堂.
那幾幅牆.
我還記得那時那幅牆.
比現在我們往復堂那幅牆還高.
因為那些是貨倉.
很高的.
我們找工人打完那些甲板之後.
我們要油油.
我們為了省錢.
自己來.
我們自己去揀了顏色.
自己買油.
油油容易嗎?.
你油你家的睡房當然容易.
現在整個大禮堂油到你手軟.
怎樣油?.
有那些滾輪,推輪.
但有些高的地方.
我們要爬梯.
當時我自己都很快.
山下已經爬上高梯.
沒有人怕我掉下來.
如果現在往復堂要大裝修.
叫我們這些傳道人去爬梯.
牧師看路.
小心.
可能一個在這裡油油.
可能下面兩三個在這裡.
比你更害怕.
不如你下來.
我們來油油.
這一種大家無分彼此甘心樂意的擺上.
很感恩.
結果我們很快.
我們整個裝修工程比預期.
我們很快完工.
而且在獻堂禮的時候.
很多靈姐妹懷著感恩的淚.
我們看到這一塊是你油的.

$^{361}$那一塊是他油的.
不過不會有色差.
為什麼?.
我們全部都是跟 specification.
你買那些油漆.
不是亂買的.
你說這個蘋果綠.
那個是青綠.
這個是墨綠.
如果你只是說綠色.
你油了不同的綠色出來.
是大件事.
很感恩.
我們是因為有這種.
大家甘心擺上.
所以無分彼此.
第二我們再看.
在接著第三到第九節.
這十二個支派的領袖.
他們將什麼送到耶和華面前.
他送了六部車.
是有縫的.
有縫蓋的.
即是淋雨都不怕的.
有十二頭公牛.
每兩個領袖奉獻一部車.
每個領袖奉獻一頭牛.
那些不是私家車.
這些就是牛車.
每兩隻牛就拖一部車.
所以是要十二頭公牛.
就六部車.
ok 行不行?.
接著拿出來之後.
神就對摩西說.
你要從他們收下這些.
作為會幕事奉的用途.
照著利未人所事奉的.
交給他們各人.
第六節到第九節.

$^{401}$摩西就收了這些車和牛.
他很有智慧.
交給利美那三個兒子.
記得我們民數記之前讀過.
他們分佈在會幕的南邊.
西邊和北邊.
這三個兒子.
而且是各盡其職.
你看摩西如何分配這些資源.
就是這個智慧.
所以領袖不要濫用神家的尊嚴.
所以我們有很嚴格的規管.
我們教會頂尖會收到的.
每一分一毫.
我們有我們的執事會.
我們有我們的財務部.
我們每一個財政預算.
每一個部門目區.
全部都要經過審核.
每一個所謂支出.
哪怕那個支票多小.
我們都要忠心.
而且我們要核數.
所以廣福街有一個很健全的財務.
這麼多年來.
我們都有一班忠心的執事.
他們願意做這些.
你說很tedious.
這些是很細微的事.
但是每個星期都要做.
每個星期我們所收的奉獻.
要即時數點.
然後要入數.
全部不是一個人去做.
是一個團隊去做.
而且彼此是監察的.
我們每一筆數點的奉獻.
是要簽名的 作實的.
廣福街是一個很健全的教會.
大家放心奉獻.

$^{441}$3月31號之前.
你們奉獻就可以有25年的補稅收掉.
我不是賣廣告.
不過這個是提醒.
讓大家知道這個情況.
你看摩西怎麼做.
他是將兩部車和四頭牛.
交給格遜的子孫.
請問格遜是誰.
他是利味的大兒子.
格遜負責做什麼.
他是這個會幕裡.
負責管理和搬運所有的布料.
就是整個帳幕的萬紙.
這個罩盤.
海狗皮做的蓋.
這個會幕的門簾.
這些布匹全部重重的.
如果你找人去搬.
是幾十公斤的.
不是只是從這裡搬去第二個門口.
是在抗爭的路上凹凸不平.
還要不知道什麼時候停.
雲彩停你才可以停.
所以整件事.
他們是為了會幕的需要.
而交給他們的領袖摩西.
摩西照著神的吩咐.
他很有智慧.
分配這些資源.
然後我們又看第八節.
他說四輛車和八頭牛.
交給米拉利的子孫.
米拉利就是利美的第三個兒子.
小兒子.
第三個小兒子負責什麼.
就是整個會幕.
工程結構的組件.
這個是最重的.
最重的包括每一個帳幕的板.

$^{481}$每一個山 每一個柱子.
還有那些大樓的座.
原址四維的.
每一條的柱子 鉗子 繩子.
這些所有器具.
是由他負責去搬運的.
拆也是由他拆.
搬到某個地方.
我們看到雲彩停在這裡.
現在要看好.
每一個板不可以錯.
次序要很清楚.
每一個榫口釘進去的角度.
要很正確.
這些就是工程 施工 執行.
全部是有規有矩.
所以我們看到一點都不簡單.
但摩西就是知道.
他們搬只是那些板.
圍著整個會幕那四周.
那麼多塊板.
所以要四輛車.
不是兩輛車.
是不是很有智慧呢.
然後第九節.
但是那一刻子孫甚麼都沒有.
為甚麼.
沒有牛 沒有車給他.
知不知道為甚麼.
那一刻是誰.
那一刻是利味的第二個兒子.
他負責會幕甚麼.
就是至聖之物.
就是聖所和至聖所.
所有那些.
至聖所的私人寶座.
和在聖所裡面的陳設餅.
桌子 金香爐.
所有這些東西.
是他負責去搬.

$^{521}$但搬呢.
就不是找個袋這樣載著.
他是要找個板.
然後找兩條柱.
兩邊都有人抬起.
一抬起.
就在你肩膀以上.
神的座位.
神的躍跪.
是代表上帝的同在和帶領.
這些全部是要他們的子民百姓.
是遠遠都見到.
前面起行.
就是他們帶著.
他們走到運彩亭的地方.
一起停在那裡.
一起再紮營.
整件事一點都不容易.
所以我們看見.
是有規有矩的.
這個讓我們學到第二個原則.
是一個國盡其職的擺上.
按各人不同的恩賜和能力.
我們去服侍和奉獻.
十二個支派領袖.
他們對會幕是運作有一定的認識.
他們知道在哪裡有需要.
他們有奉獻.
但實際上怎樣去執行去運行.
他們信任他們的領袖摩世.
交給摩世.
那些同樣.
今日地上的教會一樣.
會眾堂是作帶領的.
牧者以及執事們.
我們是團隊一起合作的.
我們是同工的.
我們知道我們需要有這份信任.
信任就是你今日拿起奉獻封.
你會看見.

$^{561}$奉獻封有償費.
償費就是為了教會平日七日的運作.
我們這些水電.
所有這些我們都是有operation cost的.
這些行政的費用.
所有這些都是在這個償費裡面要支出的.
但是有一些特別的款項.
在奉獻封有寫.
宣教奉獻.
慈惠奉獻.
這個就是每年我們會將.
我們收到弟兄姊妹這些專項的奉獻.
我們來到開會審核.
我們要支持對外的一些福音機構.
有些機構.
有些宣教士.
他做的事.
不是我們一般弟兄姊妹可以做得到的.
今日仍然有宣教士在中東.
在戰火底下.
他們在前線仍謹守崗位.
他們是冒死不離開那個地方.
這些我們是知道的.
但是我們為他們祈禱之餘.
就是他們是用他們的生命.
他們擺上他們自己.
無怨無悔無憾地.
這一種是覺盡其蹟的.
一種的奉獻.
我們也都看見.
今日我們就是對這一方面.
還有一些譬如說戒毒.
有些是兒童院.
很多這些除了金錢.
亦都是需要有義工去參與.
我們有去監獄做報道.
有頂區.
姐妹區.
都有弟兄姊妹.
是參與過這些工作裡面.

$^{601}$我們感恩.
因為我們知道.
我們是覺盡其蹟.
我們願意將最好的獻過給上帝.
跟著我們要看的那段經文.
由第十到第八十三節是最長的.
我沒有可能.
今日跟大家說完這八十三節.
但我有可能.
跟大家說第十節.
去到這個十七節.
這個我們就是由他第一個支派.
排頭位的.
就是猶大支派.
他們怎樣來做出這個奉獻呢.
我們一起來倒讀.
第十和十一節.
用高末祭壇的那一天.
眾領袖前來為獻壇奉獻.
眾領袖都在祭壇前獻公物.
耶和華對摩西說.
眾領袖為獻壇奉獻公物.
每天要有一個領袖前來奉獻.
我們在這裡看到.
是當他們完成了高末的會幕.
跟著就高末祭壇.
其實是在發生這個次序.
是會幕.
所以他們有會幕的奉獻.
奉獻完之後.
跟著他們已經做了這個高末.
用高油去抹這個祭壇.
分別為聖.
這個奉獻禮.
其實這個奉獻在原文希伯來文.
是有一個啟用的意思.
其實我們今天那些教會.
那些獻堂奉獻禮.
其實表示這個新的堂子.
現在啟用了.

$^{641}$我們將那個地方分別為聖.
歸神而用.
所以這個我們看見.
在這裡同樣的用意.
眾領袖在他們開始.
我們現在已經立好這個會幕.
也都開始了.
要有我們獻祭.
是為了神的百姓.
也都為了他們之間的和睦.
他們需要獻這些祭物給神.
我們等一下逐一在這些細節位.
我們會來講.
請留意十一節所講.
這個就不是自發性.
這個是上帝的吩咐.
上帝吩咐他們這班領袖.
要將獻壇的奉獻公物.
有規矩的.
每天只是一個領袖來.
你不要十二個領袖一起來.
每天一個.
所以要多少天才完工.
十二天.
是連續十二天.
他們獻什麼呢.
我們看一下.
在十二節到十七節.
第一天獻的是猶大之派.
誰是他們的領袖.
就是亞美娜達的兒子.
那順.
他獻什麼呢.
這裡首先.
十三,十四節是講了三樣東西.
是講他們獻掃祭裡面那些器皿.
銀盤,銀碗和金碟子.
這個銀盤.
你看和合本也好.
和修本也好.

$^{681}$我們這個和修本版本.
都是用錫克勒.
你看錫克勒.
你不會明白究竟是多重.
如果我用公斤.
你就明白.
銀盤是1.5公斤.
大概三磅重.
很有份量.
你舉一下啞鈴.
舉三磅.
有沒有份量.
有的.
你托包米.
你說五磅,十磅.
你知道多重.
是很有份量的.
載著什麼.
載著那些調油的細面.
銀碗.
74克.
就是800克.
大概1.8磅.
都是有些份量.
不過輕一點.
這兩樣東西.
銀盤和銀碗.
是載著這些調油的細面.
金碟子就輕很多.
14克.
錫克勒是110克.
很輕.
110克.
為什麼這麼輕.
因為它要載的東西不是重物.
是什麼.
香.
就是金爐那些香.
香爐那些香.
承載這些東西.

$^{721}$好.
接著.
15至17節.
就是說三種.
另外三種的製.
煩製,贖罪製,平安製.
就是這三種.
每種製都說明.
那些製身的數量是多少.
這裡我等一下再和大家交代.
結尾的時候.
就是阿米拿達的兒子.
那順的公物.
這個格式是重複12次.
我們要問一個問題.
為什麼猶大支派先走.
為什麼不找第二個支派先走.
因為舊約.
一直都是將君王和領袖這種角色.
是結合在猶大這個支派.
因為大衛王都是出於.
這個猶大的支派.
而他的後裔就是.
耶穌基督.
你馬太福音,路加福音.
你看耶穌的家譜.
就是出自大衛的子孫.
是這個猶大的家族裡面.
是在那裡來的.
而耶穌基督就是那個彌賽亞.
他是代表著君王和領袖.
這一種的角色.
都是結合在其中.
所以敬拜帶領來作奉獻.
是由猶大支派來作首的.
而不是由劉變.
那12個兒子.
劉變是大佬.
但他不是在那裡.
由他第四個兒子.

$^{761}$猶大開始的.
所以你看這個chart.
就是憲制的次序.
是照回末四面安營的位置來做的.
回末的最右手邊.
向著那邊就是東邊.
東邊就是猶大,以撒加,西布倫.
行不行?.
接著你再順時針方向.
再去南邊.
就是劉變,西面,迦德.
接著去到西邊.
就是以法連馬,拿西,變南文.
好了,去到北邊.
就是丹,阿徹和拿弗他尼.
12個支派的名字都在那裡.
另外一個.
你看這些通常就是.
剛才那個平面圖.
現在這個就是一個印象圖.
有些畫家畫出來.
是立體很多.
很清楚的.
很清晰的.
這個圖其實很幫到我們.
這裡不單止看到.
如果12個支派怎樣分佈.
你看到利美那三個兒子怎樣分佈.
好,我們也都看.
接著這個總結.
12個支派每天一個前來獻祭.
這個這樣的表述.
是12次.
首先獻公物的是第幾天.
接著那個首領的名字.
和他家族的身份.
而第三就是他所獻的.
這些物件的清單.
掃祭,繁祭,贖罪祭,平安祭.
掃祭是表示那些農作物.

$^{801}$對神獻上感恩.
繁祭是代表要將這個祭身完全燒掉.
祭司沒有份吃.
他們獻祭者都沒有份吃.
這個是歸耶和華.
完全燒給上帝.
贖罪祭也是.
他們來到是為了要贖他們子民的罪.
而平安祭是最多的那些祭身.
是為了什麼?.
為了就是表示他們是彼此有fellowship.
是一個fellowship meal.
我們團契有時都會有一些愛言.
也是這樣的意思.
我們是每一個小組.
回來教會一齊有筵席.
是彼此和大家一起溝通,吃飯.
中國人有這種吃飯的文化是很好的.
飲茶,多飲兩次茶.
你多摸幾次茶杯底又怎樣?.
不知為何好像跟他熟了很多.
因為你飲茶不是叫你每天對著他.
只是飲茶沒有用.
大家一起聊天.
這個是彼此和好的一個表徵.
而最後就是這個首領的公物.
這個格式這四方面.
十二次重複有說出來.
這個是一個清單.
大家再看一次.
銀盤銀碗為了掃祭.
金碟為了香壇.
而三個祭物的祭身.
最多就是平安祭.
義工牛,五工綿羊,五個工山羊,五個工羊羔.
肯定不是只給一個人吃.
是給甚至可能支派裡面.
整個支派不同的家庭一起分享.
就是這樣的意思.
所以今天你有沒有想到.

$^{841}$我們回來神的家.
原來我們除了敬拜上帝.
我們也要與人和好的.
不要只是等到聖餐主日.
在這裡認罪.
才在這裡審查.
你每天都可以做這件事.
所以第三個原則.
我們說我們是全體合一的擺上.
不同崗位,不同恩賜,不同支派.
但是他們是同心歸向一個中心.
就是回眸.
就是代表著神的同在和祭壇.
而這個祭壇的功用.
除了是我們向神認罪.
而祭司代表我們去獻禮物.
去贖罪之外.
還有這些平安祭.
是代表著我們可以彼此分享.
我們一起在神的子民群體裡.
我們是要彼此和好.
可不可以望一望你旁邊.
或者在後面.
有什麼人坐在你附近.
望一望.
望完了吧.
OK.
是不是你的弟兄姊妹.
是呀.
熟口熟面的.
有些為什麼你望他.
他不望你呢.
原來那個是你的仇家.
那快點認罪.
快點和好了.
明白嗎.
我們原來每一次去到神的殿.
是有這個功能的.
我們與神相遇.
神的話提醒我們.

$^{881}$我們不要輕忽錯過這些機會.
這些機會是讓我們看到.
就好像摩西在最後那節聖經.
他獻了這些祭物之後.
他去到聖所裡面.
法櫃裡面的.
耶和華的聲音就出來了.
就向他說話.
這個敬拜產生那種啟示的功能.
其實不是只是發生在週末.
這個是發生在我們每一天.
弟兄姊妹.
每一天我們自己.
與神親近.
我們的靈修.
你有沒有這種操練.
有沒有呀.
我試試問一下.
我大膽的問一下.
有沒有一個禮拜.
至少五天.
你都有靈修的.
舉手給我看看好不好.
哇 至少五天.
都很多啊.
好 非常好.
這個看聖經祈禱.
是你每天當行的.
我在禮拜五那天.
開完一連串的聚會.
還有開會.
因為我要講禮拜四.
那天下午.
我要講禮拜五.
我們門女的功課.
最後的課.
快點 趕快做完.
看到一段聖經很好.
詩篇119篇.
你知道詩篇119篇有多長.

$^{921}$不過我們每次只看一段.
而那天看到.
就是154節.
156節.
159節.
我看到這個段落裡面.
三次都提及.
神要照祂的話.
將我救活.
神要照祂的典章.
將我救活.
神要照祂的慈愛.
將我救活.
我看到這段經文的時候.
同時間我收到一個消息.
就是徐武豪博士.
我們大家都很熟的.
過去都是很多.
每年都來我們廣福堂.
有廣方崇拜.
他雖然現在是退在第一線.
但是他仍在加拿大.
在他們的機構.
是有一些服事.
但是他最近和家人.
弟兄姊妹.
他在郵輪那裡休息.
誰知道上了船沒多久.
感染了壞血病.
壞血病是可以很可怕.
死得人的.
在船上沒有太多的醫療設備.
不過有醫生.
醫生馬上開了抗生素.
希望能夠止住那些病毒.
感染他.
在下一個港口.
他就馬上要下船.
要急送他去醫院.
準備做緊急的手術.

$^{961}$當時就是這麼嚴重.
我將我那天的靈修經文.
不是這段經文.
是出埃及的.
是神怎樣透過東風.
叫那些蝗蟲吹去紅海.
神亦都可以用東風.
叫紅海分為兩邊.
叫他的子民平安.
是走過去.
上帝可以呼風喚雨.
上帝可以行神蹟奇事.
因為上帝是統管萬有.
他是整個大自然的造物主.
我和徐博士分享.
我當日的靈修.
我是看過這段經文.
但我正在做門裡的功課.
我看到詩篇.
這幾節聖經.
我說我為你祈禱.
我將這三節聖經.
將你的名字放進去.
是神要藉著他的話.
他的典章.
他的慈愛.
救活徐武豪.
到目前為止.
我一連串收回的回覆.
都是他現在.
雖然是critical.
但現在在醫院裡.
醫生說暫時要觀察多一些時間.
似乎那些抗生素是有效的.
如果有效就不用做手術.
這個是good news.
但我們看到就是祈禱.
上帝是和我們一起同工的.
Amen.
弟妹你真的要經歷這些東西.

$^{1001}$我們不只是聽到.
你走到.
要take it home.
你今天將這篇道.
你回去.
我在哪些地方.
我沒理由實踐的話.
我就求主幫我去實踐.
而且這個全體合一的擺上.
是一個很美麗的圖畫.
我講最後一個見證.
我就結束了.
今天.
我在過去這兩個禮拜忙到不知得了.
忙到我真的.
我自己都.
我很怕我.
第二天我miss了appointment.
因為我其實基本上不是很記得.
我是要打開那本大寶.
我才知道.
所以我現在真是.
每日.
即是明天早上一起身.
我就祈禱.
看聖經.
神的話.
me time.
我要很需要這些me time.
而神的靈充滿自己.
有神的話語.
盛載.
不用怕.
多忙都好.
你心裡是有那種力量.
和那種的能力.
過去這兩個星期.
我收到兩個都是廣福堂的弟妹.
轉介給我.
就說.

$^{1041}$都是分別有兩個專業人士的老人家.
一個是資深的工程師.
一個是資深的大律師.
兩個都是醫院.
我就去.
其中一個帶吳安琪牧師去.
但神都好有恩典憐憫.
結果都不是吳牧師.
亦都不是我去修閣.
因為當時對方都只是很有禮貌地聽.
但不是我們去修閣.
竟然是這位律師的一個徒弟.
原來他今天不是律師.
是做了傳道人.
所以帶埋聖經去.
跟他一說.
原來他的心已經軟化了.
然後順著那位徒弟幫他洗禮.
Hallelujah.
所以我們今天不一定是.
你向那個人傳福音.
你去修閣.
不要緊.
我們是同工.
我們大家一起.
沒有所謂.
我們是全體合一.
最重要是你肯擺上.
你肯擺上.
那個人真的要吃漢堡包.
吃到第十個他才飽.
這個就是在天家裡的周主裴牧師.
那時幾十年前他親口跟我說.
我們教會請了他來.
跟我們說教會的培靈會.
我那時是神學生.
我做司機.
就載他由多倫多去到埠仔.
所以在途中.
周牧師我們都長途.

$^{1081}$都要吃點東西.
他說不用吃太多東西.
都要吃的.
停在麥當勞.
然後在吃漢堡包的過程.
他就跟我說這個例子.
所以到現在我都記得.
我永遠都不會再灰心.
那個人黑黑面對你.
傳到第二次.
第三次又起了黑面對你.
然後很寬容.
可能到第十次他才順序.
忍耐,愛心.
這就是我們屬靈的操練.
去傳福音.
第二個例子就很奇妙.
那個是我們廣福堂的一個弟兄.
他也是一個工程師.
這個是他的師父.
他邀請我去ICU.
為他這個師父來傳福音.
原來他已經做了很多ground work.
傳福音.
原來他是修閣.
那晚我就在第二個教會講完道.
後來一打開電話.
原來有voice message 有text.
去到的時候.
原來他已經夜深了.
已經十點多.
他太太已經離開了醫院.
但是ICU有他的徒弟.
有他的舅子做見證人.
我跟這位資深的工程師.
我跟他說.
如果你聽到請你點頭.
他又點頭.
把握機會.
跟他很簡單講福音.

$^{1121}$跟著邀請他.
你願意的話.
你都可以接受耶穌.
我說你願意就請你點頭.
他點頭.
哇 點頭.
這麼順利.
我自己都有點擔心.
問多一次.
再一次點頭.
見證人全部做見證.
我為他祈禱.
因為太太女兒不在.
我們說留待第二天.
我再來為你洗禮好不好.
他說好.
第二天我天佑早講完道.
星期日下午.
我就再去幫他洗禮.
Hallelujah.
這次他竟然狀態回擁.
他舉手.
可以確定我們昨天晚上的決志是真的.
現在願不願意洗禮.
是要的.
又點頭又揮手.
他很期待這件事.
我們把赤穎玉頌讚歸給神.
這是全體合一的擺相.
先臨有達至.
所以我們永遠都不是獨行其事.
只是上帝可以獨行其事.
我們是可以團隊一起.
這樣去大人信主.
願上帝幫助我們每一個.
我們是甘心樂意.
各盡其職.
全體合一.
這樣擺上.
奉獻給神.

$^{1161}$建立神的家.
在神的國度.
為神所用.
好不好弟兄姊妹.
好我們一起祈禱.
天父上帝求你幫助我們每一個弟兄姊妹.
都能夠這樣的甘心樂意.
各盡其職.
全體合一.
我們擺上我們自己.
我們奉獻在祭壇上.
被你分別為聖.
為你所用.
得著萬民.
你要歸主.
從我們自己的家庭.
從我們工作的地方.
從我們朋友的圈子.
我們都帶領他們.
你要認識你自己榮耀的救恩.
那種豐富.
願你繼續施予厚恩.
代我們每一個.
使用我們每一個.
能夠成為你自己福音的使者.
成為基督的使者.
我們同心合意的禱告.
奉主名求.
阿們.
\newpage



\section{約翰福音 2:1-12}
\label{sec:031833}
\textbf{慈母的難題:“喜酒飲盡了!”}~~~~許開明牧師
\newline
\newline
連結: \href{http://www.church.com.hk/acms/content.asp?site=cdc\&op=show\&type=product\&code=031833\&layout=sermon}{\texttt{031833}} ~~~~ 語音日期: 2026-03-15
\newline
\newline
\hyperref[sec:031832]{\small{< < < PREV SERMON < < <}}
~
\hyperref[sec:index_chronic]{\small{[返順時目]}}
~
\hyperref[sec:031834]{\small{> > > NEXT SERMON > > >}}
\newline
\newline
約翰福音 2:1-12
\newline
\begin{longtable}{cl}
\hline
\hline
章節 & 經文 (和合本修訂版)\\
\hline
2:1 & \begin{tabularx}{0.7\textwidth}{X} 第三日，在加利利的迦拿有一個婚宴，耶穌的母親在那裡。 \end{tabularx} \\ \\ \relax
2:2 & \begin{tabularx}{0.7\textwidth}{X} 耶穌和他的門徒也被請去赴宴。 \end{tabularx} \\ \\ \relax
2:3 & \begin{tabularx}{0.7\textwidth}{X} 酒用完了，耶穌的母親對他說：「他們沒有酒了。」 \end{tabularx} \\ \\ \relax
2:4 & \begin{tabularx}{0.7\textwidth}{X} 耶穌說：「母親，我與你何干呢？我的時候還沒有到。」 \end{tabularx} \\ \\ \relax
2:5 & \begin{tabularx}{0.7\textwidth}{X} 他母親對用人說：「他告訴你們甚麼，你們就做吧。」 \end{tabularx} \\ \\ \relax
2:6 & \begin{tabularx}{0.7\textwidth}{X} 照猶太人潔淨禮的規矩，有六口石缸擺在那裡，每口可以盛兩三桶水。 \end{tabularx} \\ \\ \relax
2:7 & \begin{tabularx}{0.7\textwidth}{X} 耶穌對用人說：「把缸倒滿水。」他們就倒滿了，直到缸口。 \end{tabularx} \\ \\ \relax
2:8 & \begin{tabularx}{0.7\textwidth}{X} 耶穌又說：「現在舀出來，送給宴會總管。」他們就送了去。 \end{tabularx} \\ \\ \relax
2:9 & \begin{tabularx}{0.7\textwidth}{X} 宴會總管嘗了那水變的酒，並不知道是哪裡來的，只有舀水的用人知道。於是宴會總管叫新郎來， \end{tabularx} \\ \\ \relax
2:10 & \begin{tabularx}{0.7\textwidth}{X} 對他說：「人家都是先擺上好酒，等客人喝夠了才擺上次的，你倒把好酒留到現在！」 \end{tabularx} \\ \\ \relax
2:11 & \begin{tabularx}{0.7\textwidth}{X} 這是耶穌所行的第一個神蹟，是在加利利的迦拿行的，顯出了他的榮耀來，他的門徒就信他了。 \end{tabularx} \\ \\ \relax
2:12 & \begin{tabularx}{0.7\textwidth}{X} 這事以後，耶穌與他的母親、兄弟和門徒都下迦百農去，在那裡住了不多幾天。 \end{tabularx} \\ \\
[1ex]
\hline
\hline
\end{longtable}
$^{1}$各位弟兄姊妹 平安.
今天是母親節 又是福音主日.
所以我綁了一個題目.
特別為你們的.
就是慈母的難題.
喜酒用盡了.
那麼 我們呢.
這麼特別的日子裡面.
還有一個日子.
今天是大齋祠的節日.
而我們的廣道.
也有一個系列的.
這個系列就是七個神蹟.
和七個我事的事成.
這個是屬於神學廣章.
那麼你聽這篇廣道.
聽得神神的.
不是很明白 沒關係的.
因為這個呢.
是這些神學廣章.
是比較深入一些的.
以前我在大學做校務.
是給一些教授聽的.
特別是神學院的教授.
那麼我看起大家.
當中也有一些這麼厲害的人.
那麼.
這一種神學的廣章.
是希望弟兄姊妹.
透過.
對聖經的認識.
能夠更加深入我們的信仰.
七個神蹟.
和七個耶穌說.
我是什麼 我是什麼.
第一個就是.
變水為酒.
看不看到.
不要變了.
看回第二頁.

$^{41}$變水為酒 那麼說耶穌說.
祂是桐桐樹.
第二個醫好大神的兒子.
那麼就當中.
耶穌說祂是道路真理生命.
第三個.
醫好三十八年的病人.
記不記得我上次來講過的.
是新年的.
所以我就先講了這個.
我是羊的毛.
第四個是.
五餅二魚 我是生命的糧.
那麼耶穌餵飽了.
五千人的.
只是計算男人.
計算女的可能是一萬.
五餅二魚.
是不是神蹟.
這個神蹟背後說的是什麼呢.
我是生命的糧.
海上行走 在夜晚黑的海上.
那麼你有聽過.
耶穌平靜風浪.
也是在夜晚.
在海上那些鬼靈石靜.
邪靈當道的時候.
耶穌不僅可以平靜.
祂在那裡行走.
因為耶穌是世界的光.
第六.
醫好生來的兒子.
那麼當中看到耶穌的愛.
祂是好目人.
第七個.
拉撒路復活.
當中耶穌說了一句.
你是復活 我就是生命.
那麼這七個神蹟.
我們今天進入.

$^{81}$第一個神蹟.
那麼我們在聽到之前.
請我們祈禱.
天父上帝.
在今天教會的.
大齋期.
我們願意安靜在主的面前.
聆聽主.
你的話語.
也願你的聖靈.
光照我們.
我們得著你的話.
也願意進行你的道.
奉主名求.
阿們.
各位我們.
第三頁.
由水到走.
主要的.
內容就是說.
要怎樣由水到走呢.
要常在真圖圖.
處的裡面.
那麼耶穌就在一個婚禮.
來行神蹟.
那麼.
有什麼意思呢.
有三個意思.
第一個就在.
生活的預節來說.
婚禮是不是很開心.
有沒有人婚禮哭.
有 感動到哭.
開心到哭.
不是悲哀的.
就算哭也是喜樂的.
第二 兩個人的結合.
人和人之間.
最深的結合.
是什麼呢 夫婦的關係.

$^{121}$夫妻的關係.
說人和人之間.
一個深入關係的建立.
也就是說.
人和神怎樣能夠.
去到一個深入的關係的結合.
你不知道.
怎樣能夠和神.
有深入的結合呢.
這個神蹟.
這個經文.
對你有幫助.
第三 從神學的思想來說.
神學是有三個.
神蹟是有三個意思.
第一個就是一個.
dynamic.
dynamic的漁民.
的意思是.
是很大的能力.
而在人類.
目前認為最大的能力.
是什麼呢 是爆炸的能力.
是不是.
原子彈.
的能力是不是很大.
雖然它破壞力很大.
但是它很厲害.
所以在用一種異能.
就是說令人震驚的神蹟.
好像耶穌平靜.
風和海 是不是很厲害.
很震驚的能力.
第二歧視是miracle.
是指.
卑鄙尋常 令人驚訝的.
遲些我們說到.
耶穌令拉薩路復活.
死人都可以復活.
是不是.

$^{161}$卑鄙尋常.
另外約翰說的是.
記號.
或者用兆頭意思是.
從神而來的記號.
每一個神蹟的目的.
是要指向誰.
是要指向神是怎樣的.
指向神是多麼大能.
多麼對人的愛.
大愛和大恩.
這個兆頭.
是有象徵意義的.
Fine.
古詞.
神的名字.
如果你有聽墨牧師說.
雖然我沒有聽墨牧師說.
但是我猜到.
最早期耶和華的名字.
是甚至不能出聲的.
只有文字.
只有一個符號是YHWH.
有沒有說過.
沒有我補充.
這個原先.
是沒有聲音的.
是一個符號.
神蹟有一個更深的意思.
不是說只是很轟動.
而是.
透過這些神蹟.
讓你知道.
這個神是真的是活的.
所以.
約翰是神學家.
不只是保羅.
約翰的神學更加深入.
他綜合了.
法羅人的思想.

$^{201}$和希臘人當時的思想.
將它混合在一起.
神蹟.
用中文來說.
是最對的.
神的蹤跡.
你透過這麼驚人的事件.
在迦南的婚姻.
不是很驚人.
許牧師.
你以為不驚人.
你很細心去看.
你會知道很厲害的.
是一個符號.
指向神.
很多人信主是因為神蹟.
因為看到神蹟.
或者他體會到.
神在他自己身上的.
神蹟.
第四.
迦南婚姻從水到酒.
是耶穌第一個神蹟.
這個神蹟一點都不顯眼.
不轟動.
不動聲色.
不被察覺.
最多是誰知道.
瑪利亞會知道.
耶穌當然知道.
耶穌門徒也可能知道.
其他就是傭人.
傭人和.
證實那些.
是酒的那些人.
他會知道.
神的工作.
神的神蹟.
不是說很轟動讓人知道.
讓很多人知道.

$^{241}$只是讓很少的人知道.
而且是不動聲色.
你不覺得他很厲害.
但如果你了解的話.
你會知道很厲害.
真的很厲害的東西不一定很轟動.
讓你知道.
真的很大的破壞.
也不一定讓你知道.
你不知道那些東西.
更厲害.
魔鬼看不到.
但魔鬼他做的東西.
破壞力是嚇死人的.
所以我們要了解.
第五頁.
迦拿的婚禮.
經文大家可以看到.
我不會詳細說.
因為一會內容會提.
不重複.
第一節是相當重要.
第三天.
迦尼尼的迦拿有一個婚宴.
耶穌的母親在那裡.
這裡帶出了一個幕.
重點就是.
第三天.
是有時間性的.
因為第二天耶穌是.
招聚了一些門徒.
第一天.
他招聚了彼得.
約翰那幫.
第二天他招聚了.
另外一幫門徒.
第三天.
耶穌出來工作的時候.
去施行的第一個神蹟.
在一個婚禮裡.

$^{281}$耶穌和母親都在那裡.
第六節.
照猶太人結證的規矩.
有六口石.
光.
擺在那裡.
每天可以盛兩三桶水.
六口石的光.
在那裡是在提什麼呢.
就是說那些人.
要來個婚禮.
婚禮之前他們要結證他的手.
結證他的腳.
因為他的腳.
沒有我們有皮鞋擋住的.
有些甚至赤腳的.
很臭的.
手是骯髒的.
你要拿東西來吃的.
沒有筷子的.
所以要用的那些水.
隨便走的水.
不是普通水.
是很乾淨結證的水.
來施行他的神蹟.
接著我們看.
十一節耶穌所走的.
是第一個神蹟.
在加利利的加拿大.
特別是.
加利利.
就是一個地方.
不是只有猶太人.
有很多外邦人在那裡.
所以這個神蹟.
不是只有猶太人知道.
是要給很多外邦人.
加拿大是一個很窮的地方.
為什麼他會不夠酒喝.
因為他們窮.

$^{321}$可能他們只是計算.
有多少人來.
於是準備有多少人走.
窮人當然要看著那條數.
是不是.
你猜好像某個.
很有錢的許什麼.
最近死了那個.
不要說名字.
不是我.
許氏.
不知道什麼.
他那些.
不需要看錢.
但他都壞了.
就是因為不看錢.
他要看著.
想著不是那麼多人.
誰知道很多人來.
為什麼會不夠.
因為他們這個婚禮.
要行足多少天.
行足十天.
七天.
行足這麼久.
好了.
我們要再看.
第七.
將這段經文.
我做了一個撮要.
經文的七幕.
導演是誰.
天父上帝.
主角是誰.
耶穌瑪利亞.
配角新人.
傭人.
以及門徒等.
第一幕說.
娶親的筵席.

$^{361}$一出來就充滿了.
開心快樂.
第二幕快樂.
第二幕開始進入不快樂.
酒用盡了.
沒有了酒沒有了歡樂.
因為最重要有酒.
我們去婚禮都是說.
喝酒對不對.
喝了酒令你歡樂.
猶太人覺得酒很重要.
英國人也是.
外國人那些人要快樂.
要喝酒.
我以前在香港不喝酒.
來英國我都喝.
喝完之後.
我是很快樂的.
不過我不會醉酒.
當然.
特別廣告之前不會.
瑪利亞.
第三幕瑪利亞說給耶穌聽.
酒喝盡了.
第四幕.
耶穌有回他.
瑪利亞.
吩咐用人.
那個耶穌.
他叫你們做什麼.
你們就做什麼.
第五幕.
耶穌叫那些人做三件事.
這三件事本身.
都有意義.
有意思的.
我們一會兒會探討第六幕.
本賢直問新郎哥.
為什麼你不是.
將好的酒放在前面.

$^{401}$將那些次的酒放在後面.
哪裡不對.
因為好的酒.
喝到他們醉了之後.
看著那些酒.
都不是很多.
就給一些差的.
他快樂.
一直都快樂下去.
有道理的.
但是這個也有意思的.
第七幕.
這個是耶穌的第一個神蹟.
顯出上帝的榮耀來.
走神蹟的意義在哪裡.
不只是這個第一次.
接下來的六次.
都是有意義的.
好.
當中我們看到.
有角色在整個婚宴裡.
而這個神蹟.
還有一個很特別的地方.
就不只是.
耶穌一個人.
做獨腳戲.
是什麼呢.
最後讓這個神蹟.
能夠發生.
耶穌就低能一點.
其他的那麼大能.
為什麼這次那麼低能.
耶穌會限制.
自己的能力.
這次祂要低能.
要讓人看到有時候.
神蹟要出現.
特別在教會的事奉裡.
因為.
之前的經文也牽涉到.

$^{441}$教會的事奉.
有不同的角色.
大家彼此合作.
瑪麗亞是搭橋那個人.
久用盡.
人們去找瑪麗亞.
聖經學者研究.
瑪麗亞為什麼找瑪麗亞呢.
原來瑪麗亞應該是.
她的親戚結婚.
所以呢.
她去做一個.
婚禮的統籌.
瑪麗亞.
不簡單.
她就將這個消息告訴誰.
告訴耶穌.
因為她知道耶穌會解決.
她是一個搭橋的人.
你帶些人來教會聚會.
讓他們認識耶穌.
你又在做瑪麗亞角色.
第二個.
用人最重要.
是那些人聽話.
某大黨派.
要那些人順從.
但是.
別人不順氣.
都要逼別人順氣.
順從.
但是這個順從.
是瑪麗亞知道.
耶穌是神.
所以耶穌叫他們做的事.
一定是對的事.
好的事.
那些用人如果一天在問.
為什麼.
你是老闆.

$^{481}$你的.
你的什麼.
你的工人什麼都問你.
老闆為什麼這麼做.
為什麼這麼做 浪費時間.
問到有時會發火.
我叫你做就做.
你都不信我嗎.
所以順從很重要.
本賢值的就是要去.
驗證這些酒.
是不是對的.
對的是酒.
好不好喝.
試食.
那些皇帝試食為什麼.
怕死有人毒死他.
那會不會有人在那裡.
搞事呢.
門徒也是一個.
驗證者.
他驗證.
這件事的發生.
還有誰在這裡.
做事 我這裡沒有寫出來.
請你回答.
有沒有人回答.
我有點失望.
有個人做事.
這四個裡面沒有.
耶穌呀.
最重要的你不要忽略.
我特意不寫下來的.
讓你深刻.
人生.
最尷尬的時刻.
突然間沒有了.
人生有時.
好像瘟疫那樣.
本來應該喜樂熱鬧.

$^{521}$充滿祝福的.
但是突然間.
會沒有了 各位.
我參加過一個婚宴.
永遠我都會記得.
很開心的.
一方一對新人正婚在旺角.
然後正完婚.
就一起去.
喝酒 在旺角.
那麼.
你知道在婚宴.
進行的時間.
那些人.
吃得差不多.
就玩一下新郎哥.
新郎哥.
也很開心被人玩.
有個人就玩他.
喝光整瓶醋.
很大瓶的.
當時.
伴郎怕.
新郎哥喝了有不良效果.
等我來喝.
代你喝.
那些人就讓他喝.
他喝完之後.
開始喝得很開心.
喝著喝著 突然間臉色都變了.
因為很痛.
口痛 喉嚨痛.
一直痛下去.
原來那瓶醋.
是風熱醋.
打人的.
當時有醫生.
一看到 我們有弟兄.
哎呀.
很危急 要馬上報警.

$^{561}$馬上找救護車來.
因為會死人的.
喝了等於毒死他.
很大件事.
這個弟兄.
因為一知道一個青年.
青年團體.
全部大家為他禱告.
後來這個弟兄是康復的.
這也成為了一個神蹟的見證.
很多人這樣喝是死了.
一定死的.
所以我們不知道.
突然間沒有了方向.
突然間沒有心機.
突然間蠢了.
我發覺.
我年紀大了開始.
突然間有時會蠢了.
以前很聰明的.
師母說你現在都不聰明.
會的.
突然間會沒了錢.
是不是.
所以.
人生我們要面對.
這件事會突然失去.
正如大家手機有電.
但是如果你的訊號.
中斷.
你就會有.
這些現象.
我不用多說了.
明明有電的 開到機的.
但是畫面不清楚.
每次打.
WhatsApp的時候.
就在這裡轉圈.
最生氣的就是.
要登入銀行.

$^{601}$看你有多少錢的時候.
OTP.
一次性密碼不給你.
你就對著它哭了.
因為沒有OTP.
你進不去.
試過了沒有 各位.
我們人生就是有.
沒有走.
我們欠缺的時間.
同樣.
如果我們.
和我們天上的神.
我們那條.
屬靈的線是沒有的時候.
你就不能夠.
做到一個好的.
基督徒.
基督徒不是靠自己的能力.
可以做一個好的基督徒.
基督徒是要靠神.
在背後和你一起.
去做.
你和他有好的連線.
你的生命才有成長.
你有見證.
第一段婚宴的缺走.
人生總有.
喜樂用盡的時候.
好像我們.
錢會用盡.
酒會喝盡.
食物都會吃盡.
人手也會用盡.
很多東西我們都是會用盡的.
人生.
很多時候就要經歷到.
用盡.
你才懂得去依靠神.
你樣樣富足.

$^{641}$樣樣都有的.
你不需要神的了.
錢你會用盡.
就算感情都會用盡.
能力都會用盡.
所以我們.
需要依靠神.
但是依靠神.
不只是靠一個婚禮的神蹟.
是一個長遠的.
如何能夠長遠呢.
若翰福音.
看到.
你要連繫著葡萄樹.
耶穌是真的葡萄樹.
你要和祂有好的連繫.
以致你整個生命.
都能夠.
去面對.
有力的面對.
酒用盡了.
耶穌的母親.
這樣對耶穌說.
這個酒用盡.
是表示人會有.
變決.
但是耶穌的回答.
十三.
耶穌的回答.
我第一次.
看都覺得很疑惑.
如何回答呢.
夫人我和你何干.
我的時候.
還沒有到.
我記得我的主教學生.
有一次.
他母親叫他去工作.
叫他去.
將垃圾.

$^{681}$拉出去.
好像我們英國.
不遲都要拉垃圾桶.
我的神聖任務.
兒子對他說.
夫人我與你.
有什麼相干.
母親聽到很生氣.
說你這麼不孝.
兒子就說.
我聽許牧師說.
不是耶穌說.
不是我.
聽耶穌說這段經文.
有樣學樣.
他母親就說.
你快點出去.
倒垃圾.
兒子就說.
我的時候還沒有到.
母親.
聰不聰明.
我們以為這一句.
好像主教學生.
這樣理解其實不是.
夫人.
原文用英文去翻譯.
最好是Lady.
Lady的意思.
是一個尊貴的女士.
耶穌臨死的時候.
釘在十字架.
都曾經這樣對他母親說.
夫人.
就是Lady.
他死了之後.
他就將她交給誰.
他的門徒.
約翰.
寫作聖經的人.

$^{721}$對亞蘭文的意思.
就是說.
這個和我和你.
有什麼相關.
意思就是各人做各人的事.
你做你的事.
我做我的事.
我做我的事我心中有數.
再加一句.
但我時候還沒有到.
他要等候神.
等到他什麼時候.
開始行神蹟.
等候神.
但我時候還沒有到.
很多時候我們有一些事.
我們想做的時候.
向神禱告.
我講過.
神不是像拉丁神燈那些.
一講就馬上來.
神是要我們等候.
等候.
等候神.
對他有信心.
無必照辦.
瑪利亞就和工人說.
他說你們.
做什麼.
你們就照做.
照耶穌所說.
各位.
當我們清楚神的心意.
去做事的時候.
好像.
我很清楚神要我來英國的時候.
我就照辦.
照辦是對他有絕對的信靠.
絕對的信靠.
不再問這個問那個.

$^{761}$在這裡有一個見證.
真事來的.
在第二次世界大戰.
美日戰爭的時候.
有一個叫William的弟兄.
當時他不是弟兄.
他不信的.
他就在.
現場上.
他很害怕.
就問.
那個牧師.
君牧.
牧師說你信有神.
神會和你同在.
你信祂就行了.
他說.
我沒辦法信祂.
我都沒見到祂.
我不覺得祂是存在的.
不過我是不是這樣都可以祈禱.
牧師說.
都可以.
祈祂不要死.
說了很久.
第一次美日戰爭.
在上午當天.
日本人一來.
很快就掃他們.
旁邊那幾個都死了.
是他和那個.
那個牧師沒死.
那個牧師重了傷.
對他說.
你快點拿我身.
我這裡.
因為他動不了.
我這裡有本聖經.
你拿去吧.
因為昨天有個天使說.

$^{801}$他很需要.
這本聖經.
沒有他不行.
這個William有點猶豫.
不過他都快死了.
有本聖經.
防身都挺好的.
拿來.
他不想有什麼事.
到下午的時候.
後來那個牧師死了.
下午的時候.
日本又再一次蜂擊.
結果全部人.
都睡在那裡包括他自己.
他想著這次.
我死定了.
大家都想到他有本聖經.
頂住.
那怎麼頂住呢.
就是因為射他的時候.
中了他胸口.
聖經就將子彈.
不知道怎麼夾住了.
沒有爆出來.
他都受了傷他想著死.
後來醒了.
醒了之後他發覺.
我又完全沒有流血.
我不用死了.
他很開心跟著打開那本聖經.
剛剛那一頁.
就寫了一句這樣的說話.
詩篇91篇7節.
對有千人.
撲倒在你的身邊.
萬人撲倒在你的右邊.
這個災難都不會臨到你.
真實的各位.
這個是不是神的計劃.

$^{841}$這個是神蹟.
所以這個人後來做了什麼.
他做了牧師.
做了傳道.
我們的神就是這麼厲害.
當他肯聽.
如果他不肯聽.
那個牧師.
說的時候.
可能這個神蹟.
發生不了.
好了我們接著看15頁.
這個婚宴.
決走的難題.
我們看到這個瑪麗亞是怎麼處理.
在她的身上.
我們可以看到.
她和耶穌的關係.
他們年齡相距.
耶穌三十歲.
出來做傳道.
瑪麗亞什麼時候生耶穌.
紅貞女是指她.
多少歲 13歲.
所以她43歲.
兩個不是相差很遠.
他們彼此之間.
是有尊重的.
耶穌三十歲是成人.
很成熟的.
彼此尊重地說話.
耶穌叫瑪麗亞.
做什麼.
女士.
還有.
瑪麗亞叫耶穌.
做事也很尊重.
她說他們沒有走.
怎麼都沒有說.
是不是.

$^{881}$只是反映問題.
不會說叫兒子你要走神蹟出來.
給他們看看.
不是這樣的.
她知道耶穌有分寸.
言語間是有智慧有技巧的.
母子之間.
是有默契的.
與你何干.
你做你的事.
我又做我的事.
我做事.
心中有數.
媽媽你都一樣.
心中有數.
有信任對不對.
有尊重.
我們和子女的關係.
都需要這樣.
甚至在教會.
都要彼此尊重.
有默契.
他們的事奉有配搭.
有不同的角色.
他們關係好的時候.
所作的功能夠怎樣.
榮耀神.
又益了人.
這個是我們要學習的.
第二段.
隨便走.
主不只是解決難題.
這個神蹟的重要.
他能夠轉化人的生命.
在這三個動作裡.
耶穌說的三個動作裡.
是有意思的.
第一個.
將缸倒滿了水.
倒滿了水.

$^{921}$然後還要倒到缸口為止.
這是要表示.
這個神蹟的發生.
不是混酒進水裡.
淡一點.
可能這是他們一路以來.
欺騙別人的.
為什麼之後的酒會變差一點.
因為水混酒.
淡一點.
你也會想到.
以前的人也會想到.
不是.
是整個水.
轉變為酒.
第二.
現在可以掏出來.
為什麼不說等一會呢.
現在的意思.
就是現在這個終點.
耶穌已經.
知道這是時候了.
因為時間.
所有的東西.
都是掌握在耶穌.
神的手裡.
於是中斷了.
然後.
他們就掏水.
打了水.
放進去一些器皿.
裝酒的器皿.
大家聽不聽得懂.
好了.
然後還要掃.
掃到管聯席那裡.
用來檢查.
味道.
好了.
他們是.

$^{961}$很順從的.
所以很順利.
事報上也需要的.
我們有了好的計劃.
也知道神的心意如何.
我們就要大家合作去做.
水變酒的酒被驗證.
是上好的酒.
為什麼呢.
由管聯席那裡.
上了水變酒之後.
他們.
就去找.
新郎哥了解.
了解的結果.
是什麼呢.
這些不僅是真正的酒.
葡萄的酒還是好酒.
所以呢.
問新郎哥.
為什麼他們不是先將.
好的酒擺出來.
反而是將不好的酒.
先擺出來呢.
這裡是在說.
殉途你的生命.
可能是.
你會先是很苦的.
你信了耶穌之後.
你有了神跟你同行的話.
你就會由.
不好.
很苦變為有甜.
但是現代的人.
或者當時的人有些人生關係.
今早有酒就今早醉.
你明天有沒有酒.
很多人先洗命來錢.
那個姓許的也是.
不是我.

$^{1001}$今天有酒今早醉.
但是信仰不是這樣的.
我們信仰.
很多時候是由苦.
我們因為信了主.
我們生命有了主.
我們就變得甜.
但這個甜不是說你無災無難.
這個甜是說.
你就算有苦有難.
但是都有主跟你同行.
所以呢.
這個是很有意思的.
水變酒還有三點.
我們很快說.
水象徵了外在的宗教.
形式.
水是用來做什麼的.
水是用來洗人的腳.
洗人的手.
用完之後.
用完之後明天又要繼續.
如果有婚禮的話.
是不是.
是不停不停的.
是外在的.
你洗好了手腳.
可能未必乾淨.
這裡明白.
好.
很多人.
只靠外在的.
權利表現.
去看自己.
我應該是OK.
就可以了.
人洗手是希望.
乾淨自己.
這個不是不好的.
但是不夠的.

$^{1041}$即是要令到人.
能夠由這個水.
變為酒.
能夠生命.
得到不停的改變.
更新.
而這個改變是什麼.
不是劣等的酒.
是上好的酒.
是由不好到好.
好了.
現代人的方式.
十九.
是補救式的人生.
十九.
婚姻出了事補救.
去找輔導.
工作有壓力補救.
情緒崩潰補救.
找人疏導你.
好了.
二十.
由不好變好.
由醜陋變美麗.
這個是耶穌基督.
轉化生命的目的.
而不是補救式的.
或者只是改正式.
就夠了.
好像香港新生會幫助.
那些坐牢的人士.
重新做人.
福音戒毒.
令到那些人藉著福音不戒毒.
成為好的人.
神職也是神人的合作.
我們說了.
接著我們.
就要跳一跳.
去到廿四頁.

$^{1081}$耶穌是真的葡萄樹.
這個真的葡萄樹.
就是.
耶穌說.
我是真的葡萄樹.
凡屬我不結果子.
就要被剪斷.
然後.
你們如果常在我裡面.
第四節.
在你們裡面.
我就是葡萄樹.
你們是枝子.
你們離開我.
就不能作甚.
第七節.
你們常在我裡面.
我的話語也在你們裡面.
你們就會多結果子.
你們多結果子.
神就會得到榮耀.
這個是.
真的葡萄樹.
這個神職是要讓人.
去持續.
你聽這段聖經的說話.
持續和主.
建立一個好的關係.
就是廿六節.
我們接著下去.
要看.
怎樣能夠常在主的裡面.
第三十頁.
常在主的裡面是怎樣.
要接駁生命的連線.
各位.
我們一天不吃飯.
會不會死.
我猜大概都不會死.
幾天不吃都不會死.

$^{1121}$我沒死過.
我沒試過不吃死過.
但是雖然你不會死.
但是你會軟弱.
你沒有神的話.
正如你沒有屬靈的食物.
魔鬼就會來找你.
你有沒有常在禱告裡面.
當你缺酒的時候.
你沒有喜樂.
你沒有平安的時候.
你有沒有去禱告.
我今早很想說一個見證.
不過沒時間.
下午都要說.
各位.
孫中山先生.
他在195年.
在廣州.
起義失敗.
他去了哪裡.
他去了英國倫敦.
在196年的10月.
在倫敦街頭走來走去.
有兩個中國人說.
你不是孫中山.
孫中山很開心.
有人認識他.
英國也有人認識.
以為是好人.
跟了兩個人去.
以為找個賓館住.
誰知去了大使館.
清朝的大使館.
一進去.
孫中山立即被人捉了.
嚴刑拷問他.
那怎麼辦.
死定了.
孫中山唯一.

$^{1161}$只有禱告.
因為準備將他運到.
中國.
然後就殺了他.
不停禱告.
神就恩代他.
有一個外國人.
他找到一個.
講英文的人.
告訴他.
他想找康德黎醫生.
他的老師.
在英國很有權威.
就去幫他.
結果他沒事.
這個見證.
是很真實的.
也和我的教會有關.
這間是中華基督教會.
合一堂.
孫中山當年搞革命.
在我們的教會.
不是這間這麼漂亮的.
道濟會堂.
牧師知道這件事.
知道他被人捉了.
教會就為他.
禱告.
結果他能夠逃出生天.
很奇妙的.
十死.
一進去就沒得出來.
我們看到.
禱告的火效.
常在主的裡面.
如果有好的連結的話.
我們就會有.
看到神的奇妙.
主耶穌改.
水變酒的神蹟.

$^{1201}$這個神蹟.
聖經學者發現.
耶穌是從清水.
四十一頁.
從清水變純酒.
從不好的變上好.
從絕望變希望.
從憂愁變喜樂.
從沒有到有.
沒有走到有走.
從缺乏到豐盛.
從舊制到新制.
從虧損到榮耀.
以後的神蹟.
都是跟著.
包含了以後神蹟.
那種特徵和意義.
都在這個神蹟.
彰顯了出來.
我們接著看四十三頁.
今日是.
教會大佳期.
有三個.
很想大家去默想思想.
你今天.
是不是缺酒的人呢.
你有沒有發覺.
你很多時候都缺酒.
缺喜樂.
或者缺平安.
缺很多東西呢.
我們人生的缺酒.
人生的缺乏.
只是.
人.
在沒有出路的時候.
神是給我們出路的.
人的盡頭.
就是神.
起頭.

$^{1241}$神的起頭.
第二.
你是不是插花式的基督徒呢.
插花.
這裡有沒有花 今天沒有.
很多的花很漂亮.
但它是插花式.
是沒有生命的.
很快就會謝.
有些基督徒表面.
來說.
他好像很漂亮很好.
但他知道自己其實.
裡面生命已經和主.
沒有了關聯.
他沒有根.
他沒有吸水.
他沒有和主有好的關係.
他只不過是表面上沒有東西.
如果你是這樣的基督徒的話.
你要再想想.
你要和主有好的結連.
第三.
你願不願意受天父的.
修繭和潔淨呢.
神修繭我們.
是痛苦的.
正如.
不好意思我也要提.
如果師母說我有什麼不好.
我都會很defense的.
給人修繭.
是辛苦的.
但如果你能夠.
接受祂的修繭.
你就會有成長.
特別是神在你靈修的時候.
在你看聖經.
在你聽許牧師講道的時候.
祂有東西要修繭你的時候.

$^{1281}$你就要去面對.
去改靜.
你的生命永遠都在.
軟弱的當中.
不夠飯吃.
魔鬼最喜歡.
找些軟弱的人.
去打敗他.
我們一起祈禱.
主耶穌說.
我們承認我們自己生命有缺乏.
我們不想再靠.
自己死撐.
求你將我們的生命.
好像隨便走.
更新我們的生命.
也讓我們常常要.
願意在主你的裡面.
你是真的葡萄樹.
求天父.
你修繭我們.
因為聖經說.
我們經過修繭之後.
使我的結果子更多.
好讓我們能夠顯明.
自己是一個有真正.
信仰的基督徒.
我們的生命能夠活出.
永神億人的生命.
奉主名求.
阿們.
\newpage



\section{以弗所書 4:1-7}
\label{sec:031834}
\textbf{基督榮美的身體——歸屬,服侍,成長 (英/粵)}~~~~林穎豪傳道
\newline
\newline
連結: \href{http://www.church.com.hk/acms/content.asp?site=cdc\&op=show\&type=product\&code=031834\&layout=sermon}{\texttt{031834}} ~~~~ 語音日期: 2026-03-15
\newline
\newline
\hyperref[sec:031833]{\small{< < < PREV SERMON < < <}}
~
\hyperref[sec:index_chronic]{\small{[返順時目]}}
~
\hyperref[sec:031835]{\small{> > > NEXT SERMON > > >}}
\newline
\newline
以弗所書 4:1-7
\newline
\begin{longtable}{cl}
\hline
\hline
章節 & 經文 (和合本修訂版)\\
\hline
4:1 & \begin{tabularx}{0.7\textwidth}{X} 我為主作囚徒的勸你們，既然蒙召，行事為人就要與你們所蒙的呼召相稱。 \end{tabularx} \\ \\ \relax
4:2 & \begin{tabularx}{0.7\textwidth}{X} 凡事要謙虛、溫柔、忍耐，用愛心互相寬容， \end{tabularx} \\ \\ \relax
4:3 & \begin{tabularx}{0.7\textwidth}{X} 以和平彼此聯繫，竭力保持聖靈所賜的合一。 \end{tabularx} \\ \\ \relax
4:4 & \begin{tabularx}{0.7\textwidth}{X} 身體只有一個，聖靈只有一位，正如你們蒙召，是為同有一個指望而蒙召， \end{tabularx} \\ \\ \relax
4:5 & \begin{tabularx}{0.7\textwidth}{X} 一主，一信，一洗， \end{tabularx} \\ \\ \relax
4:6 & \begin{tabularx}{0.7\textwidth}{X} 一神－就是萬人之父，超越萬有之上，貫通萬有，在萬有之中。 \end{tabularx} \\ \\ \relax
4:7 & \begin{tabularx}{0.7\textwidth}{X} 我們每個人蒙恩都是照基督所量給每個人的恩賜。 \end{tabularx} \\ \\
[1ex]
\hline
\hline
\end{longtable}
$^{1}$親愛的弟兄姊妹和各位朋友,今天是一個充滿喜樂和很有意義的日子.我們聚集在教會歡迎Jackson加入我們雪孫這個大家庭..
Jackson,我們今天和你一起慶祝,因為你公開地宣揚你不僅是教會的一個過客,並且是我們這個信徒群體裡面的紮根..
I do want you to understand the weight of what is happening today. Meanwhile, I believe today's message is also a great reminder for all of us, like every follower of Jesus..
今天Jackson加入我們教會,成為我們會友當然很重要,但我們深信今天的訊息對每一位在座的信徒都是一個很重要的提升..
We live in a culture that is terrified of commitment..
我們現在生活在一個社會裡面的文化是很怕承諾和毀身的..
We like to keep our options open, and we join clubs, we subscribe to streaming services, and we follow social media accounts, all things that we can cancel with a single click the moment we lose interest..
我們經常會給自己留一條路,加入不同的俱樂部,訂閱不同的音樂和電影串流平台,或者我們去join social media,因為我們沒有興趣,按一個按鈕就可以取消訂閱..
In the modern world, joining something is casual, but joining church is completely different..
It is not a subscription, it is a covenant, it is an initiation into a new family..
我們現在加入教會和加入不同的平台都是很不一樣的..
這個不是一個訂閱的會員制度,是一個盟約,是加入一個屬靈家庭的儀式..
When you become a member, you are like stepping out of the crowd of casual attendants and stepping into the core of the family..
You are saying, I am no longer just a guest in this house, I am taking ownership of it..
當你成為教會成員的時候,你是從一個普通的信徒進入這個信仰家庭的核心,你現在會說,我不再是一個這個家庭的過客,而是成為這個家庭的一份子..
Today we turn to Apostle Paul's letter to the Ephesians, chapter 4..
我們今天來看保羅寫給以弗所教會的書信,以弗所書第四章..
I want you to picture where Paul is when he writes those words..
He is not sitting in a comfortable pastor's office..
He is sitting and chained in a cold, dark Roman prison because of his faith in Jesus..
我想你想像一下保羅當時寫這個書信的情景,他不是坐在很舒適的牧師的房間裡面,而是因為耶穌的緣故,他現在是關在羅馬的監獄裡面..
And from that prison cell, he writes these sixteen verses to paint a brilliant and complete picture of what it means to be a committed member of the Christ Church..
在這個牢房裡面,他寫了第十六節的經文,很生動地和完整地描述了成為一個基督教會忠實信徒的意義..
He shows us that belonging to Jesus cannot be separate from belonging to his people..
保羅教導我們,歸屬於基督和歸屬於他的信徒的群體是不可以分開的..
To be a formal member of this church means three fundamental things..
The first one is deeply engaging with his family, the second is actively committing to his work, and the last is purposefully building up his body..
成為教會的會友意味著三件很重要的事情,第一就是深入參與他的家庭,第二就是積極去衛生他的事工,和很有目的地建立他的身體..
Paul begins in the first one with a powerful and emotional plea, "As a prisoner for the Lord, then I urge you to live a life worthy of the calling you have received.".
保羅在第一節開頭就很有力量地懇求當時的信徒,.
"我為主被囚的勸你們,既然蒙召行事為人,就當與蒙召的恩相稱.".
Notice that he reminds them he is a prisoner. He is saying that, "This faith costs my freedom, so do not treat it lightly. Live a life worthy of it.".
請注意,他提醒他們自己是一個囚犯,他當時在說,.
"我為此信仰付出了自由的代價,所以不要輕視你們享有的自由,更要活出與這信仰相稱的生命.".
But what does a worthy life look like? Is it achieving great worldly success?.
Is it just reading your Bible alone in your room or watching a great sermon online?.
究竟與這信仰相稱的生命是怎樣的呢?.
是不是有很大的成就?.
是不是可以自己在房間裡讀自體聖經,或者上網去看一個精彩的講道?.
No. Look at verse 2 and 3. Be completely humble and gentle. Be patient, bearing with one another in love..

$^{41}$Make every effort to keep the unity of the Spirit through the bond of peace..
這些都不是,請看第二節和第三節..
凡事謙虛,溫柔,忍耐,用愛心互相寬容,用和平彼此聯絡,竭力保守聖靈所賜合而為一的心..
Engaging in church life means embracing the messy yet beautiful reality of living in close proximity with other believers..
參與教會的生活意味著擁抱和每一個信徒親密相處之中,有很複雜又美好的現實..
We live in an age of profound loneliness..
People have thousands of online followers, but no one to call at 2am when their world falls apart..
We desperately need new community, but real community is hard..
我們生活在一個很孤獨的時代,很多人可能在社交媒體上有很多粉絲,.
但當凌晨兩點鐘他們心情很低落的時候,想去找一個人去傾訴,是沒有的..
我們心底裡面真正的團契,但是真正的團契就是很不容易得到..
Paul says, make every effort to keep the unity..
保羅說,要竭力保守合一..
The Greek word for "make every effort" implies a sweating and exhausting kind of work..
竭力在希臘原文的意思是表達著一種很竭盡全力,筋疲力竭的努力..
Unity does not happen by accident. It requires humility, gentleness and patience..
合一並不是隨便就可以得到,它需要謙卑,溫柔和忍耐..
I noticed the phrase "bearing with one another in love", and let's be completely honest with ourselves..
You only need to bear with someone if they occasionally annoy or frustrate you..
留意保羅的說話,用愛心互相寬容..
老實說,我們面對著很善良或者對我們好的人,我們不需要用愛心去寬容..
但當有些人持續去激怒你或者令你很憤怒的時候,我們才需要去寬容他們..
When you become a member of the church, you are joining people you might never choose to be friends with naturally..
當你去加入一個教會成為會友的時候,你可能加入一班原本你不會主動去選擇成為朋友的人..
Maybe there is a successful businessman sitting next to the struggling student and the elderly widow worshipped alongside the noisy or crying toddler..
可能有一個很成功的商人,旁邊就坐著一個很刻苦的學生,一個年紀老邁的寡婦,身邊可能坐著一個一直在吵的小朋友..
We are like rough stones put into a rock tumbler..
As we tumble together, we bump into each other and that is friction..
But that friction is exactly what God uses to smooth our rough edges and make us shine like jewels..
我們就好像一粒粒很粗糙的石頭放入一個滾筒裡面..
我們這班石頭互相去碰撞,產生磨擦..
但是上帝會用這種磨擦去磨平我們的稜角,讓我們好像珍珠和寶石一樣去閃耀光華..
Engaging in church life means that when misunderstandings happen, and they will happen, we don't walk away..
The consumer leaves when they feel offended, but the confident family member, we stay..
參與教會的生活意味著人與人之間的誤會一定會發生,但這些誤會是必然的..
我們不會離開,我們不會好像消費者一樣覺得被冒犯,然後我們就走..
我們會好像盟約的屬靈家庭裡面,我們會留下來..
We sit down, we look each other in the eyes, and we forgive, we bear with one another..
That is true belonging means, when we refuse to let friction destroy our family..
我們會坐低,我們會互相寬恕,彼此去包容..

$^{81}$真正的歸屬感就是說意味著我們去拒絕人與人之間互相磨擦而產生的破壞..
And Jackson, today you are stepping into a community of imperfect people..
We are definitely far from saints, and we might sometimes fail you..
But we are also a family that will love you fiercely, pray with you, and walk beside you through life's darkest valley..
We will not let you walk alone, and this is what true belonging looks like..
Jackson,你今天加入一個不完美的團體,我們每一個人都不是聖人,我們都一定會令你失望..
我們這個大家庭會全心全意地愛你,跟你一起去禱告,陪伴你走過人生最黑暗的低谷..
我們不會讓你孤單地走人生路,這些才是真正的歸屬感..
A belonging is just the first step. True belonging must overflow into serving..
If you belong to a family, you don't just eat the food, you help to wash the dishes..
但歸屬感只是第一步,真正的歸屬感是必須延伸到付出和服侍..
就好像我們是屬於一個家,你不單止只是吃飯,你會幫忙洗碗..
Paul moves on from our shared unity to our unique individual roles..
保羅在這裡跟我們說,我們轉向各自獨特的角色..
In verse 7 he says, but to each one of us, grace has been given as Christ apportioned it..
在第七節他說道,我們各人蒙恩都是照在基督所量級各人的恩賜..
God has given every single one of you sitting here today a specific unique gift. Why?.
上帝賜給我們每一個都有獨特的恩賜,為什麼呢?.
Paul explains in verse 11 and 12 that Christ gave leaders, pastors and teachers to equip His people for works of service, so that the body of Christ may be built up..
保羅在第11,12節繼續解釋,基督賜下領袖,牧師和教師,為要成全聖徒各盡其職,建立基督的身體..
Listen carefully to Paul, this is a paradigm shift. Ministry is not just the job of us, the pastoral staff..
Pastors are the equipers, and each of you, the congregants, are the ministers..
讓我們看看保羅這裡怎麼說,福司不只是牧師的工作,牧師只是裝備我們,而你們每一個會眾都是服事者..
Think of a rescue lifeboat out in the freezing ocean..
想想一個很冷的大海上的救生艇..
The church is not a luxury cruise ship where you pay a ticket so that you can sit back on the deck chair and expect the staff to serve you the gospel on a silver platter..
想想教會不是一艘豪華的郵輪,不是你給了錢買票,你就可以睡在甲板上休息,也不是去等工作人員把福音的聖言放在你的面前..
No, the church is a rescue lifeboat pulling broken people out of the dark waters of sin and despair..
And the moment you are saved, you are pulled into this boat and hand and all to roll..
不是,教會只是一艘救生艇,要將破碎的人在黑暗和絕望的深淵裡面拉出來..
當你得救那一刻,你就上了這個救生艇,還有你會有一支槳,你就要幫忙去撐槳和去救人..
Committing to church ministry means taking a hard look at your life and asking, "Have I become a spectator?".
It means asking, "Where is there a need in this church that I can fill in?".
偉新在教會的意味著你認真去審視自己,你去問自己,我是不是只是一個旁觀者呢?教會裡面有哪些需要我可以幫忙的地方呢?.
I am often reminded of a retired gentleman in his late 60s from a church I previously attended..
我記得我以前回教會有一個七十幾歲退休的老先生..
He humbly supported the community by faithfully attending every service, every prayer meeting, fellowship and every gathering..
他很低調地支持教會,每一次的崇拜,祈禱會,團契和不同的聚會他都會很努力去參加..
And I noticed that every Sunday he arrived at least half an hour early..
To be honest, how many of us can come in early if not on time on Sundays?.

$^{121}$我留意到每一個星期日他都至少早半個鐘回到教會,老實講我們當中有幾多個人是每一個星期日都可以準時甚至提早回來教會的呢?.
Even on the days that he wasn't on the volunteer roster, he simply wanted to make sure everything at church was in order..
He would quietly clean up the mess, arrange the chairs, the tables, and spend intentional times talking with and caring for the young people..
即使他那天不在事奉名單上面,他都會提早來到教會,確保教會一切整整有條..
他會清潔地方,他會擺好些椅子,和會抽時間和人們聊天..
When his actions might have seemed small to some, his quiet and steadfast commitment to serving in the unseen details was deeply profound..
雖然他的行為好像很少,微不足道,但他每一個星期都敢做,和他持之以恆,在做一些人們不留意的服務中..
Your work of service might be welcoming strangers at the front door simply with a warm smile, so they feel safe in the new place..
It might be teaching Sunday school, patiently shaping the hearts of the next generation..
It might be completely unseen, unglamorous work of setting up chairs and tables, sweeping the floors, or fixing the audio-visual equipment..
Your work of service might be welcoming strangers at the front door simply with a warm smile, so they feel safe in the new place..
It might be teaching Sunday school, patiently shaping the hearts of the next generation..
It might be completely unseen, unglamorous work of setting up chairs and tables, sweeping the floors, or fixing the audio-visual equipment..
Yet no task is too small when it is done for Jesus..
When you serve, you are offering a physical sacrifice of praise to Jesus..
任何事情在神的眼中都不是算少,當你服事的時候,你實際就以行動去敬拜神..
I challenge you today, do not bury your talents, do not hide your gifts..
If you do not serve, a part of our church body is paralyzed, and we are literally incomplete without you doing your part..
我今日去挑戰大家,不要埋沒你的才能,不要隱藏你的恩賜..
如果你不服事,教會就有一部分會癱瘓,如果沒有你的參與,我們教會就會變得不完整..
When we engage deeply and serve actively, something miraculous happens. We reach the third and final aspect, building the church community..
We transition from being spiritual infants into mature, resilient followers of Christ..
當我們在教會投入維生和積極服事的時候,奇蹟就會發生,我們就會去到最後一個元素,就是建立教會團體..
令到我們由屬靈的小孩變成成熟,堅韌的基督的跟隨者..
Paul describes why this is so critical in verse 14, when we will no longer be infants, tossed back and forth by the waves, and blown here and there by every wind of teaching..
在第十四節繼續解釋:這樣我們不再作小孩子,被一切異教之風搖動,飄來飄去..
We live in a world of cultural hurricanes, the winds of secularism, the waves of anxiety, the shifting tides of morality. They are tearing people apart..
我們今天生活在一個文化潮流充斥的年代,有世俗主義,有焦慮的浪潮,有道德觀念的不斷變幻,人心的不斷侵蝕..
If you try to stand alone against all these storms, you will be swept away. You need a solid weight of a committed community to anchor you. 假如我們嘗試用自己的能力去對抗這些潮流,我們可能會被捲走.我們需要一個堅定的群體來去堅固我們的根基..
And how does the community do this? Paul gives us the secret in verse 15 and 16. Instead, speaking the truth in love, we will grow to become in every respect the mature body of Him who is the head, that is, Christ..
From Him the whole body, joined and held together by every supporting ligament, grows and builds itself up in love, as each part does its worth..
And how does the community do this? Paul gives us the secret in verse 15 and 16. Instead, speaking the truth in love, we will grow to become in every respect the mature body of Him who is the head, that is, Christ..
Notice the critical balance, speaking the truth in love..
我們要留意其中一個關鍵詞,就是用愛心說誠實話..
Truth without love is harsh, judgmental and crushing. It destroys people..
如果真理是沒有愛的話,就變成很嚴厲的判斷,叫人很痛苦,甚至會摧毀人..
But love without truth is sentimental, weak and cowardly. It leaves people stuck in their sins and addictions..
但如果沒有真理的愛,就很情緒化,很軟弱,很懦弱,它會令到人繼續陷在罪中..
Real biblical community means loving a brother and sister enough to gently look them in the eye and say, "The choice you are making is destroying your life.".
And it means being humble enough to listen when someone warns you..

$^{161}$真正按照真理的團契意味著,如果我們愛的弟兄走到這一步的時候,.
你可以溫柔地跟他說,你所做的選擇正是破壞你的人生,同時也意味著謙虛地聆聽別人的警告..
Think about the human body. For a child to grow into a strong, healthy adult,.
every organ, every bone and every microscopic ligament must do its part..
If your liver stops working, your whole body dies..
就好像我們的身體,一個孩子要成長成為健壯健康的人的時候,.
我們每一個器官,我們每一個骨骼,每一個微小的人大,都是各按其跡..
例如如果我們的肝臟停止工作,我們身體就會死亡..
The church is exactly the same. We build ourselves up in love only when each part does its work..
教會也是這樣,唯有每一個部分都各按其跡,我們才能夠在愛中去建立自己..
I usually, during preaching, us, the pastors, provide some application or examples for the congregation,.
so that you know what you can do to live out the Bible..
我們牧者通常在講道的時候都會提供一些實際的應用或者例子,.
令到你知道怎樣將聖經付諸實行..
But this time, I would like to encourage you yourself to reflect on how you can practice what the scripture says..
這一次我鼓勵你自己去想想怎樣可以將聖經的道理去實踐出來..
I would like to encourage you to think of a person who hasn't been showing up at the service or fellowship,.
a face that you haven't been seeing for a while..
我想鼓勵你去想想有沒有一個人很久沒有來教會參加崇拜或者團體活動..
或者他現在是在人質的谷底,或者他被世界的文化浪潮吹到飄來飄去,.
或者他需要你的陪伴和鼓勵..
我鼓勵你今天下午就拿起電話,發一個訊息,打一個電話,或者邀請他喝咖啡..
你不需要一個神學學位才可以用愛心說誠實話..
你只需要出現,你要讓他知道有人去掛著他們,.
讓他們知道他們很重要,他們是教會不可或缺的一部分..
讓我們立志成為暴風中彼此扶持的群體,不讓任何人獨自去面對這個風浪..
當你和耶穌更接近的時候,你的信仰也會引領你身邊的人更加去接近耶穌..
我們不只是靠自己,我們不能夠就這樣靠自己達到屬靈的成熟,.
我們要一起去共同去成長,不然我們就會停滯不前..
來到訊息的結尾,讓我和傑森再多說幾句..
當你今天毀了這個教會的時候,請記住基督對你那份浩瀚和難以言喻的愛..
這份愛驅使耶穌基督走上月林林的十字架,將你贖回歸入祂的家..
當你回顧四周,看看禮堂裡所有人的面孔,請記住現在我們是你的家人..
今天我誠意邀請你以愛心和我們攜手同行,一起為教會的合一而努力,.
希望你能夠繼續發掘神賜給你的屬靈恩賜,並且與我們一同服事,.
也希望你一同去建立教會,叫我們彰顯耶穌基督的樣式..
而其他的會眾,當傑森今天在我們面前公開宣言,.
我們絕對是有責任,也是上帝給予我們的責任,去歡迎祂,去保護祂,去包容祂的缺點,去裝備祂,去服事,並且與祂同工,直到耶穌基督再來的日子..
願上帝繼續賜給我們救援的恩典,令我們成為一個愛中緊密的團體,在服事中積極行動,不斷去共同成長,效法我主耶穌基督的教會..
讓我們低頭祈禱,.

$^{201}$我們天上的父,我們今天懷著無比感恩的心來到您的面前,.
主,我們今天特別為傑克遜禱告,因為他願意公開宣認加入雪孫這個教會的大家庭,我們祈求你賜給他勇氣,讓他能夠深入參與和委身教會..
讓到我們中間如果有磨擦的時候,求神你賜給我們謙卑,溫柔和耐心,讓我們彼此以愛包容..
求主幫助我們,當我們遇到困難的時候,不退縮,而是努力奮鬥,堅守聖靈所賜的合一..
主,我們祈求你喚醒,你賜給我們每一個的屬靈恩賜,讓我們可以忠心服侍你..
主,我們內心每一個好像想做消費者或旁觀者的念頭,賜給我們能力去承擔和共同努力..
父啊,最後我們求你幫助我們建立這個教會屬靈的身體,求主賜給我們勇氣,教導我們用愛心說誠實話,教導我們不再像屬靈的嬰孩一樣,被世界的風浪搖動..
求神你用每一個的肢體,讓我們緊緊地連結在一起,直到這個教會彰顯主耶穌基督美好的樣式..
我們將每一個人的生命都放在神的手裡..
奉耶穌基督的名頭.
阿們.
\newpage



\section{以弗所書 2:11-22}
\label{sec:031835}
\textbf{福音的本質:在憂傷中帶來的盼望}~~~~陳義生傳道
\newline
\newline
連結: \href{http://www.church.com.hk/acms/content.asp?site=cdc\&op=show\&type=product\&code=031835\&layout=sermon}{\texttt{031835}} ~~~~ 語音日期: 2026-03-15
\newline
\newline
\hyperref[sec:031834]{\small{< < < PREV SERMON < < <}}
~
\hyperref[sec:index_chronic]{\small{[返順時目]}}
~
\hyperref[sec:031836]{\small{> > > NEXT SERMON > > >}}
\newline
\newline
以弗所書 2:11-22
\newline
\begin{longtable}{cl}
\hline
\hline
章節 & 經文 (和合本修訂版)\\
\hline
2:11 & \begin{tabularx}{0.7\textwidth}{X} 所以，你們要記得：從前你們按肉體是外邦人，是「沒受割禮的」；這名字是那些憑人手在肉身上「受割禮的人」所取的。 \end{tabularx} \\ \\ \relax
2:12 & \begin{tabularx}{0.7\textwidth}{X} 要記得那時候，你們與基督無關，與以色列選民團體隔絕，在所應許的約上是局外人，而且在世上沒有指望，沒有神。 \end{tabularx} \\ \\ \relax
2:13 & \begin{tabularx}{0.7\textwidth}{X} 從前你們是遠離神的人，如今卻在基督耶穌裡，靠著他的血，已經得以親近了。 \end{tabularx} \\ \\ \relax
2:14 & \begin{tabularx}{0.7\textwidth}{X} 因為他自己是我們的和平，使雙方合而為一，拆毀了中間隔絕的牆，而且以自己的身體終止了冤仇， \end{tabularx} \\ \\ \relax
2:15 & \begin{tabularx}{0.7\textwidth}{X} 廢掉那記在律法上的規條，為要使兩方藉著自己造成一個新人，促成了和平； \end{tabularx} \\ \\ \relax
2:16 & \begin{tabularx}{0.7\textwidth}{X} 既在十字架上消滅了冤仇，就藉這十字架使雙方歸為一體，與神和好， \end{tabularx} \\ \\ \relax
2:17 & \begin{tabularx}{0.7\textwidth}{X} 並且來傳和平的福音給你們遠處的人，也傳和平給那些近處的人， \end{tabularx} \\ \\ \relax
2:18 & \begin{tabularx}{0.7\textwidth}{X} 因為我們雙方藉著他，在同一位聖靈裡得以進到父面前。 \end{tabularx} \\ \\ \relax
2:19 & \begin{tabularx}{0.7\textwidth}{X} 這樣，你們不再是外人或客旅，是與聖徒同國，是神家裡的人了， \end{tabularx} \\ \\ \relax
2:20 & \begin{tabularx}{0.7\textwidth}{X} 被建造在使徒和先知的根基上，而基督耶穌自己為房角石， \end{tabularx} \\ \\ \relax
2:21 & \begin{tabularx}{0.7\textwidth}{X} 靠著他整座房子連接得緊湊，漸漸成為在主裡的聖殿。 \end{tabularx} \\ \\ \relax
2:22 & \begin{tabularx}{0.7\textwidth}{X} 你們也靠他同被建造，成為神藉著聖靈居住的所在。 \end{tabularx} \\ \\
[1ex]
\hline
\hline
\end{longtable}
$^{1}$各位弟兄姊妹 早晨.
今天很開心又見到大家.
大家睡得好不好?.
吃得飽不飽?.
我們很開心見到現場的觀眾.
和網上在教會禮堂的弟兄姊妹.
或者在其他地方聚會.
他們幫我叫了Uber.
很快到沙田.
很快就回去了.
我自己也做大學生試工.
我知道辦一個營會是很不容易的.
但你們的團隊是很專業.
你們的傳道同工和預備這個營會的弟兄姊妹.
多謝他們的辛勞去預備.
完了之後收到你們傳道人的一個訊息.
給我一個好消息.
他說原來今天可以有多一點時間.
不用像昨晚那樣趕得這麼急.
但放心我們不會整個早上留大家在這裡.
另外就是昨天我也學了一樣東西.
原來可以透過AI去想想怎樣介紹港元.
因為五月份我們學完傳道會有一個活動.
是五月二十四日.
我們請了美國庫勒臣學院的一個老師.
他是講關於一個parenting的講座.
家長的講座給一些小朋友.
如果家長的孩子是小朋友的話.
會想一想怎樣可以訴訟他們信仰的方法.
昨天我的同工一個designer.
就叫我給一些描述關於這個港元.
我就學了你們的方法.
AI幫我去撰寫幾句來介紹.
所以非常好都學了一些東西.
我們的營會是叫做心靈感應.
我們昨天說了心.
是有些時候在一個超速的環境裡面.
我們會有一種恐懼.
恐懼我們和其他人比較的時候.
不如其他人迷惘.

$^{41}$福音的工作不是不重要.
但是只是我們這麼多要做的事情裡面.
可能第三第四或者更加後的.
變成事不關己.
我們要聆聽.
聆聽在上帝面前尋找.
祂給我們的呼召使命是什麼.
祂給我們的感動服侍怎樣的群體.
今天是最後一個部分.
我們要感應去回應.
不是只是知道但是不走的.
而是我們都可以一起勇敢來回應.
昨晚我們用了比較長時間去看.
什麼是一個超速的世代.
我們說了什麼.
我們說了要求一個高學歷.
競爭相當激烈.
我們往往好像被世界潮流.
要推著我們來跑的時候.
每一刻都停不下來.
我們很想透過一些成就.
或者個人的一些地位.
自我實現在當中.
很多很多外面而來的一種期望壓力.
壓在我們身上.
得不到一種滿足.
得不到一種滿足在當中.
這種不斷的自我實現.
有些時候一方面不能滿足我們心靈裡的一些需要.
另一方面令我們沒有一個行動在當中.
福音的事情有時變得事不關己.
於是我們都問了一個問題.
怎樣可以在這個無止境的一個困境裡走出來.
而不是在這裡等呢.
怎樣可以在一個無止境的追求裡得到一種釋放.
聽一聽神對我們的呼召的心意其實是什麼呢.
怎樣可以將我們人生的身份.
不是建基在這些成就.
而是回到紮根於神的愛呢.
這個是我們昨天所看的.

$^{81}$就是原來祂對我們有一個呼召.
祂對我們的呼召是要賦予它的額.
意思是要接受它的王權.
成為它的門徒.
而今天在這一點裡我們再走多一步.
我們再解釋多一點.
如果神呼召我們成為祂的門徒.
去賦予祂的額的話.
其實那個使命是什麼呢.
所以我給了今天的題目.
我們不如回去再思考這個福音的本質是什麼.
祂給我們的使命是什麼.
以致當我們要去感應的時候.
其實我們在回應什麼呢.
所以今天下來的時間有三個部分.
第一部分我們很想回來說.
人類正在面對的問題是什麼.
我們很誠實地去說.
其實我們這個世界正在發生什麼事.
第二我們一起去看看.
教會和福音的本質是什麼.
我們又很誠實地去問一問.
神給我們的使命其實是什麼.
而第三個部分就是.
為什麼上主會將我們放在香港這個地方.
而我們可以怎樣去回應呢.
這就是我們今天的時間.
不如開始之前我們一起祈禱.
天父來到你面前求你幫助我們.
今天的早上再一次回來裡面去思考.
這個世界的本質和福音的本質是怎麼樣.
你都可以明白和聽得到主要是怎麼樣將我們這間教會放在香港這個地方.
在今時今日在這個環境裡面.
你要我們的使命和託付是什麼.
祂來幫助我們直至你只有話語和我們說話.
祈禱去奉耶穌的名求.
大家都可以拿你的聖經出來.
或者如果你有手機裡面的聖經.
都可以搜尋到已不鎖書的哲經文.
我們看回前一點點好嗎.

$^{121}$好像昨晚我們嘗試去明白.
首先進入那段經文.
但是我們再看回前一點點.
究竟保羅在一個怎樣的context裡面去寫這一節的經文.
一開始第一部分說到人類正在面對的困境.
是一個憂傷破碎的世界.
如果我們看回前一點點第二章的時候.
第二章一開始第一節到第三節.
保羅其實是說一個比較宏觀的一種想法.
去說到我們正在面對的問題.
第一到第三節他說了些什麼.
他說我們本來就是和上帝來敵對的.
是未認識上帝之前我們死在過犯的裡面.
是絕望的.
不單只是我們的生活缺少了些什麼.
而是很直接的去說我們是死了的.
是靈性裡面是死了的當中.
第二節接著怎麼說.
他說到那時你們在過犯罪惡裡面生活.
隨從今世裡面的風俗.
順從空中掌權者的一些領袖.
就是現在叛逆的人心裡面運行的邪靈.
他在說一種現實.
這種現實是我們後面彷彿有股力量推著我們.
是一個邪惡的力量.
是一個罪惡的力量.
去推著我們做一些事情.
我們明知是不對的.
但是我們都會不信福神.
是被奴役的一個環境.
放縱我們肉體的一些私慾在當中.
本來這裡說的是一個未信主之前的光景是怎樣.
不是個別人的光景.
是全人類所有人的一個故事在當中.
又是跳一跳去到第十二節.
第十二節接著說到.
要記得那時候.
你們那時候和基督是無關的.
和以色列.
與以色列選民團體來隔絕的.

$^{161}$是所應許的人是一個局外人.
在世界上是沒有指望的.
是沒有神在當中的.
這個是當我們未接受耶穌之前的那個困境.
但是這個不是那個結局.
這個我們如果看回第四節裡面.
這個有一個好消息出現.
他說到賢義神有豐富的憐憫.
因著祂對我們的大愛.
僅是我們因為過犯死的時候.
叫我們和基督一同活過來.
賢義神 英文是But God.
聖經裡面大概出現了四十五次.
裡面的意思就是說原來有一個好消息.
藉著神給我們的恩典原來有一個出路.
這個出路來賦予我們一個新的生命.
重新喚醒我們的心靈.
注入一個新的視野.
給到我們有一個新的可能性在當中.
我以前認識一個大學生.
我一開始服侍的時候就被警察派去理工大學服侍.
認識一個大學生.
他有很嚴重的自我形象的問題.
性格是非常暴戾的.
他小時候不是和他父母一起來生活.
他小時候就是和他的姑媽一起來生活.
小學二年級的時候.
他姑媽生日去酒樓裡面有一個聚餐.
但是那時候發生了一件事情.
他因為不喜歡他的姑媽老老瘦瘦.
於是他就起飛腳.
他的姑媽就差不多整張椅子掉下去.
那時候他的姑媽已經七十八十歲.
之後進了中學之後.
他就天天都很緊張.
他經常都會受到一些親戚來責備.
但是他覺得完全不是一回事.
甚至很生氣.
為什麼你們憑什麼來管教我.
你又不是我的父母.

$^{201}$升了中學之後.
他在學校裡面的生活是很緊張的.
天天不是他撩其他人打架.
就是其他人去撩他打架.
說到為什麼你要經常去打架.
他的原因就是.
原來他發覺裡面給他一個很大的成功感.
給他一個很大的成就感.
當他見到其他人痛苦的時候.
他有一種莫名其妙的開心的裡面.
親戚見了他這樣的情況很長時間.
於是就幫他祈黃大仙.
祈不同的神.
但是都是沒有用.
有一年是他一年級的時候.
我們在PolyU的Canteen裡面接觸了他.
跟他們有一個分享.
原來我們成為神兒女有一個可能性在當中.
我們什麼都沒有想過.
就是在PolyU的Canteen那一天.
他就決志接受了耶穌.
他突然間發覺原來我在神裡面.
原來可以是很安全的.
他突然發覺然而神有一個豐富的憐憫.
給他一個可能性在當中.
他突然間發覺然而神給他一個可能性.
他是可以得到自由的.
跟他之前在過犯罪惡當中的生活.
原來是可以很不同在當中.
這一位的弟兄.
他之後就參加了我們的小組.
慢慢慢慢來到成長.
我為什麼要提這個弟兄的一個故事呢.
我看到一樣東西就是.
我們發覺當我們在罪惡捆綁的勢力的環境裡面.
其實不單止是我們跟神之間的關係疏遠.
其實人跟人之間的關係都是疏遠的.
這個是整個社會裡面的一個現實.
剛剛過了農曆新年.
不知道大家的農曆新年怎樣過呢.

$^{241}$應該農曆新年是一家人團團和氣.
一起來相處的環境.
但是我們發覺現在很多時候.
特別經歷了這個社會的事件之後.
很多的關係裡面是很緊張的.
有很多的一些兄弟姐妹.
或者父母和他們的子女.
很多年都沒有話題.
見了面都是沒有話題再說的.
我們發覺最影響的.
不單止是我們跟神之間的個人關係.
人跟人之間有很多關係上的破碎.
家庭之間.
朋友之間.
國與國之間在打仗.
種族之間.
我們繼續看他的經文.
第十四節.
請大家幫我讀第十四節.
1,2,3.
問一問大家.
他講到使雙方合而為一.
這裡的雙方是在說哪雙方.
這裡的雙方是在說哪兩雙方的人.
他在說的是一群受了割禮的猶太人.
和未受割禮的一些外邦人.
經文怎樣形容他們之間的關係.
他們中間是有冤仇的.
是冤仇的.
發生什麼事了.
用不用這樣這麼嚴重.
第十五節好像告訴我們一些答案.
不如我們再讀第十五節.
1,2,3.
他講到解決這個冤仇的方法是怎樣.
是要廢掉那個既在律法上的規條.
好像告訴我們.
造成這個冤仇的原因.
就是因為那些律法上的規條.
而今天當我們廢掉了這些律法上的規條.

$^{281}$就能夠解決那個冤仇.
這樣就令到我們有更加多的問題出現在我們腦海裡面.
不可思議.
為什麼會這樣.
律法不是應該是聖潔的嗎.
律法不是應該是好的嗎.
為什麼這些律法反而造成了兩者之間的冤仇.
其實發生了什麼事呢.
人類社會究竟發生了什麼事.
以致我們會進入一個破碎的環境呢.
這裡我可能需要給大家一些背景資料.
讓大家明白多一點.
是 以色列人和上主是有一個文約的關係.
而上帝設下了律法來指引他們.
這些律法例如是聖潔不聖潔.
是什麼是潔淨不潔淨等等.
大家要留意.
這些律法的設予的目的.
不單止是期望他們遵守成為一個公義的社會.
還有一個更加大的目的.
以色列人被揀選是要成為上帝祝福世界的救贖的渠道.
意思是透過遵守這些律法.
他們向這個世界的人展示.
一個跟從主的群體究竟是怎樣的.
他們透過聖潔的生活.
去吸引其他愛邦的人來跟從.
如果這樣說的話.
這些律法其實有一個大一點的目的.
是宣教性的 是使命性的.
是要我們去回應的.
即是我們以色列人.
他們被召一開始就是為了世界的緣故.
成為一個展示性的展示在當中.
成為萬國的光.
讓其他人看得到.
有一個例子體驗一下成為神的子民應該的樣子是怎樣的.
所以這些律法本身是好的.
但是我們發覺歷史裡面好像去到某一個位置.
這些本身是好的律法.
反而令到這群以色列人變得更加驕傲.

$^{321}$令到其他人更加多的敵意.
體驗不起身邊沒有律法的一些人.
美好的禮物反而成為了一個仇恨的根源.
因此回到剛才所讀的第十四節.
要拆毀了中間這一種的心態.
這一種間斷的牆.
以致可以終止了一些的仇恨在當中.
但是弟兄姊妹老實說.
其實哪裡只有以色列人和愛邦人的間斷的牆.
其實是一個很普遍的現象.
很多時候上帝都賜給我們一些很好的禮物.
可能是恩賜.
可能是其他的擁有.
人擁有這些本來很好的事情.
但是提升到一個層次.
而體驗不起身邊沒有的一些人.
這些很好的禮物就成為了一種的咒詛和難阻在當中.
而當我們過分去看我們所擁有的時候.
我們可能會覺得自己有更加多的優越感.
環顧四周發覺.
咦?為何其他人和我會不同呢?.
這些禮物成為了一種敵意的間斷的牆.
成為了一種驕傲.
不知道大家有沒有試過被標籤的一種感覺.
我們北美讀神學的時候有一個這樣的經歷.
當時每一年我們都要去一個地方.
有一些華人的神學生的會議.
有一年他在芝加哥.
那間神學院是西部的洛杉磯.
我又沒有去過芝加哥.
我又沒有去過美國東部.
都很興奮的.
想著可以坐飛機過去認識其他在北美讀神學的弟兄姊妹.
去到的時候每人有一張名牌.
名牌裡面寫你叫什麼名字.
你那間神學院.
見到名牌之後我就掛上去.
心裡面很期待很渴望和其他人連結.
於是我就坐下來和一群人聊天.
但他們一見到我的名牌.

$^{361}$都還沒有握手.
還沒有打招呼.
見到我是來自於某一間神學院的時候.
他就開始問很多很尖銳的問題.
他就說你是來自於那間學校的.
最近有一些老師說了一些話.
你怎麼看某一個事件可能是十年前發生的.
你怎麼看另一個老師可能二十年前已經退休了.
你怎麼看那些情況.
弟兄姊妹我們發覺有些時候我們會將人標籤.
久而久之我們失去了一種回應和愛人的心態.
因為我們看人好像一個濾鏡一樣.
將人很豐富的層面.
減弱成為被標籤的那些特質.
我究竟是誰.
陳以珊你是誰.
不要緊.
你就是某某神學院的一個學生飛了過來.
其他事情不重要的.
你就是那間學校的一個學生.
在今天的社會裡我們有些時候.
可能只是留心一些和我們身份很相似的人.
如果那個人和我們的身份是不同的.
有些時候我們要攀附.
才可以和他溝通得到.
中間成為了一個間斷的牆在當中.
人看人不是對對方是一個人.
而是透過一些標籤來衡量那個人.
你做什麼的.
你住哪一頭的.
你是哪一頭上班的.
你有沒有價值的.
你可不可以幫我達到某一些目的在當中.
保羅告訴我們這是第一世紀裡面面對的情況.
亦都是我們今天全人類面對的情況.
人被罪惡克制.
影響人和神和人和人之間的關係.
是沒有出路的.
在第一世紀裡已經很清楚說明.
這個撕裂是不會停止的.

$^{401}$我們的政府可能用了很多不同的東西.
是一些律法 法律 治安 輔導等等.
但是都不會有一個真正的平安在當中.
有什麼方法呢.
有沒有一些解決的方法呢.
我們一起看下去好嗎.
教會和福音的本質是什麼.
我們一起看下去第十六節.
請大家幫我讀第十六節.
一二三.
這裡說到原來只有在基督耶穌裡面的十字架.
我們才可以回應得到這個憂傷破碎的問題.
而教會 你們 我們.
就是這個破碎世界裡面的盼望在當中.
這裡可能我要再仔細去解釋一下.
所以我需要大家的專心.
我們再仔細去解釋一下.
這一種本身是怎樣能夠回應.
處理得到這個世界的問題好嗎.
我們一起跟著我.
我們先看一看另一節經文.
我們去看另一節經文是迦拉太書六章.
裡面怎麼說呢.
它是說到 但我絕不以別的誇口.
只誇我們主耶穌基督的十字架.
自誇這個字有時候會翻譯成一種榮耀.
一種份量.
保羅的意思是說到他成為了基督徒之後.
他停止因為其他的東西給他一個身份的地位.
他不再需要依賴其他的東西.
愛邦人和猶太人對他來說是毫無意義的.
回到第十七節.
他說到他要傳福音給遠處的人和近處的人.
因為在十字架裡面可以消除我們的自誇.
不論是遠處的愛邦人.
或者是近處的猶太人.
他們都需要福音.
他們的身份是不會給他們任何好處的.
只有十字架裡面才可以消除這些自誇.
叫我們在神裡面謙卑下來.

$^{441}$不需要有比較.
不需要玩這個比較的遊戲.
但他繼續說不單止在十字架裡面消除自誇.
他繼續說第十六節.
回到第十六節.
他說到剛才大家幫我讀的.
他說是十字架裡面消除了冤仇.
並且藉著十字架使雙方歸為一體.
同神和好.
什麼意思呢?.
他說到原來是十字架裡面.
是我們消除了自誇.
但十字架裡面我們可以創造一個新的人類.
更加準確的說.
是一個新的人類在當中.
是新的人類的可能性.
他之後的經文就用了一些建築物的比喻.
去說到教會這個群體.
我們聚在一起.
不是純粹因為社會學的原因我們聚在一起.
不是因為你做什麼職業.
你的年紀.
或者這個話題和旁邊那個人很相似.
於是你就回來了.
是同一個小組.
他說不是因為這些社會學的條件.
而是因為一些屬靈的條件.
以至到我們今天一起聚集在這裡.
這個屬靈條件是什麼?.
第20節.
他說到被建造在使徒和先知的根基上.
而耶穌基督他自己就是這個防割石.
什麼是防割石呢?.
防割石是古代裡面一塊.
基本上是一個很大的石頭.
在建樓的時候是很重要的.
是一個很大的石頭.
以至到所有的壓力可以壓在裡面.
例如你們這幅圖畫.
咦?.

$^{481}$為什麼去了?.
對了.
例如你們這幅圖畫.
它是一個差不多有750噸的巨石.
在耶路撒冷的一個西牆那裡.
鋪設這個防割石是很重要的.
尺寸是要很準確的.
以至到將所有的壓力壓下去.
如果它的鋪設是不準確的話.
整個建築,整個房子就會倒塌.
所以保羅說到.
基督就是那塊防割石.
是最關鍵和很重要的.
我們看看以上亞書怎麼說.
以上亞書接著說到.
所以主耶和華這樣說.
看來我要在石案裡面放一塊石頭作為根基.
那個防割石是試驗過的石頭.
是有穩固的根基.
是最寶貴的防割石.
信靠它的人必不著急在當中.
為什麼十字架裡面可以消除了.
回應世界裡面那種憂傷.
破碎人與人之間的問題.
原因是因為在十字架裡面.
它消除了一切的自誇.
並且在十字架裡面.
它出現了一個新的人類的可能性.
並且這個新的人類的可能性.
可以展現出.
一種很不同的款待態度在當中.
這種很不同的款待態度.
在這個世界是給不到的.
英文Hospitality.
款待這個字源來自希臘文.
裡面說的是一種對陌生人的愛.
意思就是說將陌生人都願意帶來.
在自己家裡面.
如果我們回到剛才的經文.
裡面說到我們怎樣用光景.

$^{521}$它說到其實我們本身都是一些愛邦人.
但是神祂自己因為基督的防割石的原因.
祂展現了一個超凡的款待.
將我們納入了神的家裡面.
我們本來都是一群陌生的人.
在世界上沒有神沒有指望.
但是我們都可以成為神家的一份子.
因為這個超凡的款待.
給我們一個新的可能性.
在2020年上.
昨晚跟大家分享了我在2020年的時候.
回到香港.
這幅圖畫是當時我拍攝的.
大家都知道2020年發生什麼事.
就是新冠疫症.
新冠疫情大概在三月份的時候.
來到美國的加州.
我最小的小朋友.
是二月底剛出生.
他一出生的時候.
這個疫症就來到我們的省份那裡.
整個環境是很緊張.
你們看得到超市裡面是沒有東西可以買的.
去超市裡面人和人之間.
要隔開一米的距離.
我當時面對的是一個剛出生的.
剛誕生出來的嬰孩.
另外有兩個大概五歲的小朋友.
還有剛生完小朋友的太太.
在病床那裡.
本身家人希望來幫助我們.
但因為整個疫情他都來不了.
那時候可能是我人生中.
其中一個經歷到最低的低潮.
我基本上不知道可以做些什麼.
而就是那個時候我經歷到.
原來是弟兄姊妹當中.
一個超凡的款待來幫助.
當時弟兄姊妹見到這個情況.
他們問都沒有問.

$^{561}$就輪流來煮飯給我們.
還把飯送到我家裡面.
在門口放下來給我們.
有一天我幫我剛出生的嬰孩去洗澡的時候.
我就把電話放在客廳裡.
我就沒有聽電話.
於是當我聽電話的時候.
我就聽到原來是我當時有教會的牧師.
他叫做鄭牧師.
他就跟我說.
喂 Patrick.
我說我是啊.
他說我是鄭牧師.
我煮了一些雞湯和一碟菜.
我現在在你樓下.
但我等了兩個小時.
你快點下來吧.
那為什麼我等了兩個小時.
因為我幫嬰孩洗澡的時候.
我電話放在外面.
於是我就放在外面.
但鄭牧師為什麼不一早打電話給我呢.
原來因為他很擔心我會拒絕.
於是我就馬上去到他樓下.
見到他駕駛著一輛本田的灰色的七人車.
左手拿著一碟白切雞.
右手拿著一碟雞湯.
他在馬路上為我們祈禱.
我覺得很感動.
當時他已經六十五歲.
過了幾個星期.
我又在幫嬰孩洗澡.
接著我又聽到一個電話.
我拿起電話的時候.
他就說.
喂 Patrick.
我是鄭牧師.
你快點下來吧.
我已經等了兩個小時在下面.
我到他樓下的時候.

$^{601}$我又見到鄭牧師.
他正在駕駛著一輛本田灰色的七人車.
左手他有一個雞湯.
右手有一個白切雞.
但和上次有一樣東西不同.
他當時給了我一盒口罩.
他跟我說.
你和Emily你們快要回香港了.
你們比我更加需要這些口罩.
你知道當時的加州.
如果你要找一些合資格的口罩.
幾乎是沒有可能的.
有很多人在香港寄過來才可以.
為何他有兩盒口罩呢.
是他的女兒從三藩市特意駕駛七小時車下來.
送給他.
他只有兩盒.
但他給了我.
他說因為你要回香港.
你比我更加需要.
我覺得很感動.
於是我收了那兩盒口罩.
於是我六月回來香港.
之後有一次.
又收到他的電話.
原來他在下面等了一段時間.
下去的時候又看到一輛.
本田灰色的七人車.
左邊有一個雞湯.
右邊有一個白切雞.
另外他還給了三封大利是.
說這三封大利是是給你小朋友的.
他就告訴我們.
你回香港要用心去事奉.
在這個過程中.
我不期望弟兄姊妹.
和我的年長六十多歲的牧師.
他是這樣來愛我們.
展現一個超凡的一種款待.
我對他來說.

$^{641}$其實可以說不是很有價值.
因為我是會回香港.
我不會留在他的教會幫他事奉.
但在新冠疫情最嚴重的時候.
他駕車冒著生命危險.
出來將他僅有的口罩留給我們.
什麼是Hospitality.
什麼是一種款待呢.
他說是一種.
即使是不認識的一些陌生的人.
將他帶到自己的家裡.
我們本身都不是屬於神家.
我們本身在這個世界上沒有神.
但我們可以成為天國裡的一個成員.
但這個不是無端端突然間發生的.
這種超凡的款待.
將本身一開始所說的那種仇恨.
標籤.
默化一個人.
只追求自我提升而不理會身邊的人.
間斷的牆壁打碎了.
但不是無端端發生的.
是因為有一個人.
曾經有一個超然的付出和款待.
來成就.
我們看看馬太夫人怎麼說.
馬太夫人說到.
像人所棄的石頭.
就成為了房角石頭的石塊.
當主來到這個世界裡.
他犧牲自己.
才換來我們可以經歷的款待.
我不知道大家有沒有聽過這首歌.
中文翻譯就是房角石.
是澳洲Hillsong Church裡的其中一首歌.
我上網查的時候.
我本身都不認識裡面的故事.
裡面查的時候.
他的作者叫做Ruben Morgan.
他就這樣說.

$^{681}$他說他為什麼會作這首歌呢.
來自於一次在挪威發生的恐怖襲擊.
當時這個作者Ruben Morgan.
飛到挪威的時候.
他發現整個城市都是封閉的.
他和一些朋友來見面.
他突然間感受到整個城市.
因為整個恐怖襲擊.
而有的一種震盪在裡面.
他突然間發現整個城市.
經歷了一種痛苦恐懼.
那一次的恐怖襲擊.
有一個人亂槍掃射.
整條街裡面很多人死了.
而當他走在這條街的時候.
人是恐懼,懷疑,悲傷的環境裡面.
突然間有這個歌詞.
就出現在他的腦海裡面.
他在訪問裡面就這樣來分享.
他說有些時候我們可以很順境去敬拜.
但有些時候我們可以.
在一個風暴的環境裡面去敬拜.
有些時候我們可以在一個懷疑的環境裡面去歌唱.
為什麼我們可以在一個恐懼,懷疑,不確定的環境裡面去歌唱.
因為當我們的人生被推到一個極限.
甚至面臨死亡的時候.
而發現一切的希望和盼望.
只是來自於一個的防角石.
那個是耶穌基督.
所以我們可以在軟弱的裡面去歌唱.
有些時候是更加堅強,更加清澈的當中.
因為知道我們所信的是誰.
而祂就是我們的防角石.
Ruben Morgan就錄了這首詩歌.
提醒我們耶穌是我們的防角石.
所以這首詩歌裡面他怎麼說呢.
他中文的翻譯是這樣.
他一開始就說.
我別無盼望任何一件事情.
唯獨耶穌是的寶血公義.

$^{721}$我絕對不相信這個世界任何的謊言.
只信靠耶穌的聖名.
接著他的副歌就不斷不斷重複.
是耶穌防角石.
耶穌的名字在當中.
什麼是福音呢.
昨晚我們說了.
我們在這個超速的環境裡的一種回應的心態.
而今天我們說我們要回應.
回應什麼呢.
神學家Andrew Root是這樣說的.
全福音.
全福音究竟是什麼呢.
全福音他就說.
Evangelism is about meeting human sorrow.
我覺得他說得很好.
很精簡.
很簡短.
Evangelism is about meeting human sorrow.
全福音就是和人在悲傷裡的一種連結.
而藉著我們的行為.
帶來一種盼望安慰.
承擔其他人的憂傷.
讓我們一起經歷上帝的窮在.
在裡面我們可以展現一種生命的可能性.
我們看了這個世界發生的問題.
我們看了教會和福音的本質是什麼.
最後不如我們看一看.
神為什麼將我們放在香港這個地方.
而我們可以怎樣去回應好嗎.
我們回到剛才的經文.
這次幫我讀21到22字.
1,2,3.
漸漸成為主理的聖殿.
你們也靠他同被建造.
成為神跡著聖靈居住的主宰.
非常之好.
保羅他繼續用建築物的比喻來說.
他說我們這個教會的群體.
不是純粹在十字架裡消除了自誇.

$^{761}$基督是我們的一個防角石.
不單這樣繼續來說.
他告訴我們教會的存在是有一個原因的.
是有一個目的的.
21節他說這個房子是互相緊扣.
我們有不同的小組不同的團體.
互相緊扣.
是漸漸成為主裡面的聖殿.
這種漸漸成為.
我們不是純粹回來聚會就算了.
不是純粹回來迎會就算了.
他說有一個目的是當中.
這個目的是什麼.
如果我們往上看一點.
第十節.
他說我們原是他的工作.
是基督耶穌裡面造成的.
為叫我們去行善.
這是神所預備叫我們所行的.
這裡說得很好.
他說我們是他的工作.
英文說我們是神的工藝.
或者更加準確的說我們是神的工藝產品.
是一個傑作.
而聖經裡面只出現了兩次.
如果我們去看教會裡面的歷史的時候.
你會發覺很多偉大的事工都是由基督徒開始的.
譬如關禁貧窮人.
廢除奴隸的制度.
餵飽沒東西吃的那些人.
我這樣說不是說只有基督徒才會做這些工作.
但是我們要發覺作為主的聖殿的話.
是有一個更加大的原因在後面推使我們.
可不可以在一個憂傷破碎的環境裡面.
做一些事情.
好像Andrew Rutko說的.
Meeting people's sorrows.
在憂傷的環境裡面讓人去見神.
我上年去了柬埔寨.
我想這幾年大家認識柬埔寨應該是因為KK元的事件.

$^{801}$我們就開始對他認識多一點.
全球超過22萬人在柬埔寨或者緬甸裡面.
來做一些詐騙的工作.
數量是很大的.
問一問一個問題.
你知不知道全世界不同形式的奴隸.
奴役制度或者人口販賣.
如果我們將所有這些人數加起來的話.
你覺得有多少個紅館可以裝得起.
如果是以一個紅館的尺寸來量度的話.
你覺得如果將全世界任何形式的奴役制度.
人口販賣等等加起來.
你覺得會有多少個紅館可以填滿.
答案是10個?.
100個?.
1000個?.
答案是4000個.
4000個在當中.
隱藏的暴力正在影響很多不同的人.
5000萬的人在全世界受奴役的影響.
四分一是兒童.
每年的夜遊是很誇張的.
有一些人是抵債的奴役.
即是他幾代的人都在被奴役.
有一些是網上兒童被性剝削的情況.
在2022年每100個菲律賓的小孩子.
就有一個在當中被販賣.
製作一些兒童色情的項目.
曾經有一個人說.
那時我已經失去了逃走的希望.
我覺得你去被囚禁.
也是沒有能力去反抗.
於是他做事來哭訴.
神啊 如果你是真的 求你救我離開這個地方.
我認識一位姊妹.
她是大學畢業沒多久就加入了一個機構.
這個機構的名字是國際行議會.
而剛才那幾個powerpoint是她給我的.
這位姊妹是一個很普通的姊妹.
她都去承認.

$^{841}$有些時候見到這麼多的問題.
我是一個很普通的姊妹我可以做些什麼呢.
她說我見到這個世界.
發生的事情的時候.
她覺得很困難.
但是她突然間發現.
是有些事情我們是沒有辦法去選擇的.
譬如是我們成長的環境.
我是在一個怎樣的環境裡成長.
但是有一些事情.
我們是可以選擇的對我們的人生.
譬如我們的工作.
我們要怎樣去回應一些使命等等.
她就說有一句話去提醒她.
就是一句這樣的說話.
大家可能大眼睛看得到.
是叫做God's Way.
Our Work.
Jesus' Way.
是什麼意思呢.
這位姊妹說原來推動她.
鼓勵她在一個這樣的環境裡服事.
是因為她突然間發現.
耶穌是她生命的防角石.
而她問一個問題是.
如果神真的在我們生命裡有一個.
真實的作為的時候.
我們是否都可以.
為著身邊的人帶來一些的震撼.
一些的影響呢.
如果耶穌在我們生命之中.
有一個強大的作為的時候.
我們是否都可以傳福音給一些遠處的人.
或者給一些近處的人當中呢.
但是我們發覺那個現實是怎樣呢.
現實是很多時候.
當我們環顧四周的時候.
我們都很想成為一個很好的基督徒.
但是現實是很多時候.
我們活著好像跟聖經.

$^{881}$是沒有什麼很大的關係當中.
因為很多時候我們關注的.
可能是我們生活的水平.
我們的目光有些時候.
不是放在神自己的心意裡面.
不是很留意或者環境發生其他的事情.
神對我們的心意究竟是怎樣.
究竟這個世界發生什麼事.
有些時候跟我們是沒有什麼關係.
或者好像昨天所說.
是我們那個次序裡面.
可能排到很後的裡面.
因為我們更加.
學悟和愛我們所擁有的東西更加多.
有些時候我們會說.
只要我有什麼我們就夠了.
只要我有一個好一點的工作我們就夠了.
只要我有多一點錢我們就夠了.
只要我有一個好一點的環境我們就夠了.
但是這個只要我有什麼.
有些時候就正正是因為一個想法.
攔阻著我們尋求終極良善的那一位.
以致我們信仰的回應裡面停了下來.
麻木了我們的熱情.
另一個人不再專注.
是那個可以滿足我們靈魂深處的神.
在他自己的身上.
這個只要我有什麼.
有些時候正在毀滅我們的生命.
復興的Jonathan Edwards.
他有一句說話我覺得很好.
他說Sin turns a heart into a fire.
Sin turns a heart into a fire.
什麼意思呢.
他說我們的心.
好像一個貪婪的火焰一樣.
一直燒一直燒一直燒.
永遠都覺得不夠的一直燒在裡面.
直至整個平原全部都被燒盡為止.
曾經有一個人.

$^{921}$他的名字叫鍾馬田.
他在1899年的時候出生.
他以前是一個醫生.
他獲得倫敦大學的醫學博士.
他慢慢在行醫的過程裡面.
對他的呼叫感應越來越有一個清楚的想法.
他掙扎了很長時間.
他突然間發覺.
在醫學裡面.
他是一個很重要的角色.
但他慢慢慢慢發覺他有一種不滿足在當中.
因為他發覺人的靈性的狀況.
是更加沉重的負擔.
他發覺不是最滿足的裡面.
一方面醫療是一個很重要的.
但不能夠解決人類更加真實根本的問題.
就是我們和罪之間的關係.
我們和神隔絕了.
於是一段時間之後.
鍾馬田離開他的醫療專業.
他全時間奉獻在傳道和牧養的工作裡面.
有一次他接受一個訪問.
訪問的人問鍾馬田.
他說你犧牲了這麼多.
你犧牲了一個很成功的醫學.
成為一個很穩定的傳道人.
成為一個全職的傳道人.
你有沒有後悔這個決定.
鍾馬田當時的回應是什麼.
他說.
It was no sacrifice on my part.
對我來說不是任何一個犧牲.
他覺得他沒有放棄任何東西.
但是他能夠得著全部.
因為他反思過去這麼多年.
他專注在福音的工作的時候.
發覺原來這是一個最大的榮耀在當中.
第一節目我們有兩天的一個營會.
我們昨天和今天去看這個世代的狀況.
我們這個世代充斥著很多不同的憂傷和撕裂.

$^{961}$原本其實是沒有盼望的.
但是因為基督自己的犧牲.
摧毀了人的自誇.
成就了一個新的可能性.
這個世界是給不到我們的.
而我們其實每一個人.
我們本來都是一個局外人.
是不屬於神的家.
但是因為基督自己的工作.
我們有機會進入到神的國度.
我們每一個人都是神的傑出的精品.
我們回來這裡不是純粹醉醺醺就算.
我們是被賦予一個使命.
這個使命是告訴我們.
我們在一個憂衰破碎的裡面.
要幫助人和神相遇.
將一個盼望美好的消息告訴其他人.
而教會不單只是一個建築物.
它是一個有機的生命.
有一個很清楚的目的.
是讓我們可以來見證神的.
而如果我們回到聖經裡面.
我們知道神是地方裡面設立地方教會.
是有目的的.
這個不是一個偶然.
我知道窩夫的出現都不是一個偶然.
是實際上是一個道成肉身的原因.
出現在今天這個地方.
我們出現在這個地方裡面.
在香港是一個更加深層次的意義當中.
或者我們都會想.
可能我們所做的事情都很有限.
但是不論我剛才所說的故事.
可能是那位姊妹.
或者是一個六十多歲的鄭牧師.
只是拿著一些飯來我裡面照顧我們.
而一些地方我們都可以承傳展開.
參與在這個美好的事工裡面.
在這個環境裡面.
很想鼓勵大家.

$^{1001}$我們還在等什麼呢?.
我們不如一起來回應.
最後不如我們一起有些禱告的時間.
我邀請大家.
亦都可以有片刻的時間.
來安靜為自己的祈禱.
我不知道在昨天的晚上.
今天的早上.
或者是一些靈修的時間.
神有沒有在這個過程裡面.
都和你有一些說話呢?.
我邀請你有片刻的時間.
去整理一下自己的思緒.
最後我會有一個結束的禱告.
天父在你們面前.
我們要祈禱.
就是面對這個憂傷破碎的環境裡面.
我們發覺人和人.
人和神之間的關係.
出現很多的破裂在當中.
成為很多間斷的牆.
有時這些牆厚到一個地方.
我們都不知道怎樣和人重新建立關係.
但主我們仍然相信.
今天你將我們.
將窩福放在這個環境.
在今時今日的環境裡面.
是有你自己的心意的.
主我們本來都不是神你自己家裡面的人.
是因為你自己的念文.
給予我們一個新的可能性.
讓我們能夠進入.
同時間求主使用我們.
以致我們都可以將這個和平的福音.
分享給遠處的人.
或者近處的人.
以致我們都可以將這個福音.
在一個人的悲傷的環境裡面.
讓人和主你自己去連結.
做一個新的生命的可能性在當中.

$^{1041}$讓我們有這兩天的營會在當中.
但主你幫助這些感應.
不只是停留在這個營會裡面.
當我們下午離開這個營會出去的時候.
繼續幫助我們勇敢的人去見證神利.
我們禱告去奉耶穌名求.
\newpage



\section{哥林多前書 3:1-15}
\label{sec:031836}
\textbf{你要食飯嗎?}~~~~王正傑牧師
\newline
\newline
連結: \href{http://www.church.com.hk/acms/content.asp?site=cdc\&op=show\&type=product\&code=031836\&layout=sermon}{\texttt{031836}} ~~~~ 語音日期: 2026-03-15
\newline
\newline
\hyperref[sec:031835]{\small{< < < PREV SERMON < < <}}
~
\hyperref[sec:index_chronic]{\small{[返順時目]}}
~
\hyperref[sec:031837]{\small{> > > NEXT SERMON > > >}}
\newline
\newline
哥林多前書 3:1-15
\newline
\begin{longtable}{cl}
\hline
\hline
章節 & 經文 (和合本修訂版)\\
\hline
3:1 & \begin{tabularx}{0.7\textwidth}{X} 弟兄們，我從前對你們說話，還不能把你們當作屬靈的，只能把你們當作屬肉體的，你們在基督裡僅是嬰孩。 \end{tabularx} \\ \\ \relax
3:2 & \begin{tabularx}{0.7\textwidth}{X} 我用奶餵你們，沒有用飯餵你們，因為那時你們不能吃。就是如今還是不能， \end{tabularx} \\ \\ \relax
3:3 & \begin{tabularx}{0.7\textwidth}{X} 因為你們仍是屬肉體的。你們中間有嫉妒、紛爭，這豈不是屬乎肉體，照著世人的樣子生活嗎？ \end{tabularx} \\ \\ \relax
3:4 & \begin{tabularx}{0.7\textwidth}{X} 有人說：「我是屬保羅的」；有人說：「我是屬亞波羅的」；這樣你們豈不是和世人一樣嗎？ \end{tabularx} \\ \\ \relax
3:5 & \begin{tabularx}{0.7\textwidth}{X} 亞波羅算甚麼？保羅算甚麼？我們都是神的執事，藉著我們，你們信了；這不過是照著主給各人的恩賜去做罷了。 \end{tabularx} \\ \\ \relax
3:6 & \begin{tabularx}{0.7\textwidth}{X} 我栽種了，亞波羅澆灌了，惟有神使它生長。 \end{tabularx} \\ \\ \relax
3:7 & \begin{tabularx}{0.7\textwidth}{X} 可見，栽種的算不了甚麼，澆灌的也算不了甚麼；惟有神能使它生長。 \end{tabularx} \\ \\ \relax
3:8 & \begin{tabularx}{0.7\textwidth}{X} 栽種的和澆灌的都是一樣，但將來各人要照自己的勞苦得到自己的報酬。 \end{tabularx} \\ \\ \relax
3:9 & \begin{tabularx}{0.7\textwidth}{X} 因為我們是神的同工，而你們是神的田地、神的房屋。 \end{tabularx} \\ \\ \relax
3:10 & \begin{tabularx}{0.7\textwidth}{X} 我照神所給我的恩典，好像一個聰明的工頭，立好了根基，別人在上面建造；只是各人要謹慎怎樣在上面建造。 \end{tabularx} \\ \\ \relax
3:11 & \begin{tabularx}{0.7\textwidth}{X} 因為，那已經立好的根基就是耶穌基督，此外沒有人能立別的根基。 \end{tabularx} \\ \\ \relax
3:12 & \begin{tabularx}{0.7\textwidth}{X} 若有人用金銀、寶石，草木、禾秸，在這根基上建造， \end{tabularx} \\ \\ \relax
3:13 & \begin{tabularx}{0.7\textwidth}{X} 各人的工程必將顯露，因為那日子要將它顯明，有火把它暴露出來，這火要試煉各人的工程怎樣。 \end{tabularx} \\ \\ \relax
3:14 & \begin{tabularx}{0.7\textwidth}{X} 人在那根基上所建造的工程若能保得住，他將要得賞賜。 \end{tabularx} \\ \\ \relax
3:15 & \begin{tabularx}{0.7\textwidth}{X} 人的工程若被燒了，他將損失，雖然他自己將得救，卻要像從火裡經過一樣。 \end{tabularx} \\ \\
[1ex]
\hline
\hline
\end{longtable}
$^{1}$弟兄姊妹平安.
今日我們要講一個課題是.
你要食飯嗎?.
可否幫我問下你旁邊那個 你要食飯嗎?.
我們一講到要問你食不食飯 當然是一定要食飯.
今日我們要看一下 究竟為甚麼.
哥林多教會的信徒他們不能夠食飯.
仍然要食奶.
當我們看今日這篇經文 哥林多前書.
其實在第三章是發出一個問題.
這個問題除了是講當時哥林多教會那班信徒的狀況.
也同時問我們弟兄姊妹.
其實你現在在基督徒裡是否還是一個嬰孩呢?.
是否只可以食奶 還是你已經可以食飯呢?.
在這裡我們必須要問一個問題.
就是為甚麼這班哥林多的信徒.
他們的屬靈生命是沒有成長的呢?.
同時在這裡也讓我們看見一件事.
原來屬靈生命是要成長的 是要有所成長的.
或者到了我們人長大 成長的時候.
我們慢慢不會想著要吃甚麼來令我們更成長.
因為我們越來越長大 吃東西只是為了維持我們的生命.
除非我們有小朋友.
小朋友出生我們就在想究竟可以給他甚麼吃.
讓他可以健康 讓他去成長.
同樣我們屬靈生命也好像是這樣的.
我們信了主之後 我們信了主一段日子.
我們未必會去想究竟我們生命要怎樣去成長.
究竟我們要追求甚麼來成長.
我們好像只是在維持我們的屬靈生命.
而哥林多的信徒他們更加出現一個狀況.
他們不單止沒有去成長.
他們只是停留在一個外表上 他們看來像一個基督徒.
但是他們的生命一點成長都沒有.
跟我們信主之前是沒有分別的.
這個才是重點.
所以我們更加要問自己.
究竟對於今天我們自己的生命.
是否跟我們以前信主之前是否沒有分別呢.
有沒有改變呢.

$^{41}$以前我們很喜歡罵人.
現在是否仍然很喜歡罵人呢.
這是我們今天要認真去面對的一個問題.
保羅他形容當時哥林多教會這班人是怎樣呢.
保羅說他不能當他們是屬靈的.
只能當他們是屬肉體的.
他們在基督徒裡只是一個嬰孩的.
這班哥林多教會的信徒.
為甚麼呢 因為他們當中有紛爭結黨.
他們又容忍有淫亂的事情在他們中間發生.
他們彼此有訴訟.
你又告我 我又告你.
濫用自由.
即是他們濫用了他們在基督徒裡面的自由.
亦都對婚姻的情況誤解.
聚會裡面的秩序亦都很混亂.
他們甚至有假教師 有假教導.
他們不信復活的事情.
並且保羅更加在第三章的一至十五節裡面.
今天我們看的這段經文更加提到的就是.
哥林多教會裡面的分裂和紛爭.
所以這裡告訴我們的是.
保羅告訴我們 他只是寧願用奶餵哥林多教會的信徒.
我們看經文為甚麼這班哥林多教會的信徒沒有成長.
經文一開始 保羅就說.
他用愛去提醒這班哥林多教會的信徒.
指出從前只能當作他們是屬肉體 在基督裡為嬰孩.
初信的時候他們在基督裡面只是一個嬰孩.
所以當時只能用奶 一個淺意的教導來餵他們.
沒有用飯 即是一些盡心的教導來餵養他們.
但過了一段時間 但如今仍然不能夠.
其實在這段保羅在他們最初信的時候.
去教導他們 過了一段日子 現在再去教導他們.
發現他們仍然不能夠用一些盡心的教導來教導他們.
即是發現他們是甚麼.
他們仍然是一個屬靈的嬰孩.
對於嬰孩來說 這個比喻就是形容他們是甚麼.
即是他們未能夠咀嚼 只能夠吃奶 不能夠吃飯.
這個其實是形容信徒只是停留了一個基本的認識.
不能夠進一步去吸收一些更深的教導.

$^{81}$所以怎樣會成長呢?.
即是等於我們信了主 只知道耶穌拯救我們.
釘十字架 賜我們永生.
我們有祝福 有喜樂平安.
可能我們的信仰就停留在這裡.
我們沒有再去吸收一些更深的教導.
甚至可能接受不到一些更深的教導的時候.
我們又怎樣可以去成長.
怎樣說到我們要去建立教會.
怎樣說到我們要去維生在基督的上面呢?.
而引致我們仍然好像一個嬰孩.
不能夠咀嚼 我們只能夠吃奶 不吃飯.
有幾種的情況.
有時候我們上教會 聖經的教導.
對我們來說可能只是一個知識 一個資料.
我們很喜歡聽一些知識 一些資料.
最好講一些大家未聽過的東西.
對我們來說很吸引.
但這些知識和資料只是停留在我們的理性裡面.
第二樣東西讓我們未能夠去咀嚼.
只能夠吃奶 不能夠吃飯.
就是在說什麼呢?.
就是在說我們不去咀嚼.
即是我們不能夠去反省.
我們可能在聽道 我們在自己的看聖經.
我們甚至可能查經等等.
我們沒有去反省自己的生命.
所以我們看了的只是一堆什麼?.
人生大道理.
我們沒有去反省自己的生命.
第三種的情況就是我們吃不下.
怎樣吃不下?.
一講的時候 我們一聽又覺得.
哎呀 牧師為什麼要講我?.
把我的壞事都講出來.
當然牧師不知道聽到誰講.
所以在講台就說我.
又或者有時候我們吃不下 為什麼呢?.
我們吃不下就是一聽到的時候就說.
對呀 你講得對 就是他.

$^{121}$這個就是他 他就是這麼壞.
對呀 說醒他.
我們都吃不下.
而第四更重要的就是.
我們聽了道.
從來都沒有去行道.
大家記不記得上個禮拜.
永勝執事他講道講什麼?.
約拿大他是怎樣?.
不應call 我們是不是經常都不應call.
聽了道都是不行道呢?.
所以這四種情況都引致我們.
仍然只是好像一個嬰孩.
因為我們不能夠去咀嚼.
我們只能夠去吃奶 不能夠去吃飯.
情況有點像有一套電影有一段名句.
那段名句是這樣講的 它說.
講你又不聽 聽你又不明.
明你又不做 做你又做錯.
錯你又不認 認你又不改.
改你又不服 不服你又不講.
當我們這樣去面對神的話語.
面對教導的時候.
其實我們突然之間是很難去有所成長.
而且對於我們來說.
我們一般認為有成長的基督徒是應該怎樣的?.
可能是他都很熟悉聖經.
經常上聚會 勇於承擔.
他很衛生在事奉裡面.
外表都很敬虔.
你會覺得這些是有所成長的基督徒.
這些是我們的印象.
你幫我看看你旁邊那個是不是一個有成長的基督徒.
我們是這樣認為的.
但是保羅在這裡告訴我們.
信徒沒有成長是怎樣的情況?.
未必只是純粹是講外表.
保羅在這裡告訴我們.
這班信徒 那麼多信徒.
他們沒有成長是因為他們什麼?.

$^{161}$他們仍然是屬肉體.
他們照著世人的樣子去走.
這個就是他們當中會有的情況.
而他們發生什麼事 為什麼保羅這樣說呢?.
就是因為他說.
他們信徒中間有嫉妒 有紛爭.
並且他們分黨分派.
有人說他們是屬保羅的.
有人說他們是屬阿波羅的.
所以在這裡保羅給我們看見.
原來一個人的屬靈生命有所成長.
不是我們單單從他的外表.
我們的角度是否屬靈來看見.
而是在看他的生命.
他的生命是否營救一個像屬肉體的人.
他是否照著世人的樣子去走.
所以我們明白.
弟兄姊妹在說教會健不健康.
其實和我們信徒的屬靈生命的成長.
是很有關係的.
所以我們做信徒的.
我們屬靈生命有成長.
不是因為純粹屬靈的外表.
而是因為在於信徒仍然是屬肉體.
他像世人一樣的時候.
多不多教會沒有成長在他們中間.
因為他們的行事.
他們也是照著世人一樣去做.
弟兄姊妹.
這個是我們很需要去留意在我們的生命.
究竟今天我們是否仍然.
好像屬肉體的生命一樣呢?.
我們過的生活好像和世人一樣.
所以保羅在這裡指出.
信徒仍是屬靈的嬰孩的原因.
就是因為他們屬肉體.
他們照著世人的樣子去走.
為什麼他這樣說呢?.
因為他們照著世人的樣子去走.
是因為當時對於希臘人.

$^{201}$他們是身處希臘文化的時間.
希臘文化是怎樣呢?.
他們很容易就將一些人聖化了.
神聖化.
譬如一些很有智慧的教師.
一些講師在他們中間.
他們就會將他們神聖化.
而多不多教會的信徒.
他們分黨分派.
就正正是因為當時他們有這樣的情況.
他們有人將保羅神聖化.
有人將阿波羅神聖化.
所以有人說我屬保羅.
有人說我屬阿波羅.
以致他們都怎樣?.
所以仍然只能夠吃奶.
對於基督教我們都有追星.
有些人在基督教都很喜歡聽某些講員去講道.
所以有些講員在外面有講座也好.
有工作坊也好.
他的粉絲都會去很擁護.
然後一留意到就會去聽他講道.
聽他的工作坊.
聽他教書等等的情況.
這個在基督教經常都會見到.
但等於我們都必須要知道一件事.
就是其實無論那些名講員也好.
那些牧師也好.
其實他們都是人.
他們都是人.
所以我們不需要去神聖化某些人.
更加不要因為我們神聖化某些人的時候.
然後引致我們中間有嫉妒紛爭.
這個正正是一個屬肉體的行為.
照世人一樣去行.
所以早前有一位暢銷基督教的作家楊肥力.
當他承認曾經有八年的婚外情的時候.
就帶來很大的基督教的震撼.
很多人都會問.
當年他寫了一本書講恩典.

$^{241}$我看過這本書的時候.
他當時有婚外情.
那對我是不是真的可以幫助到他呢.
他講的東西是否正確呢.
弟兄姊妹.
素靈的領袖都只是神的同工.
我們都不要將他們神聖化.
他們都只是人.
更加對保羅的提醒.
在教會裡也不要神聖化某些人.
因為他們都只是神的同工.
不要因為他們而產生任何的嫉妒紛爭.
然後信徒沒有成長.
他們就需要吃飯 需要去追求成長.
所以今天當我們說我們要去追求成長的時候.
我們就正正要面對我們的生命.
我們要過一種順從聖靈的生活.
而不是一種順從肉體的生活.
所以今天弟兄姊妹.
我們要順從的是聖靈帶領.
不是順從我們自己的肉體.
我們自己的喜好.
我們自己的中意.
我們認為對的東西.
我們要過順從聖靈的生活.
並且第二樣我們要分別為聖.
我們不是要和這個世界的人一樣去做.
而是我們要分別為聖.
就是要和我們這個世界的人一樣.
所以保羅進一步要去回應.
那些信徒他們分黨分派.
他就問 阿波羅 保羅算什麼呢.
他說他們都是福音的執事.
他們都是神的僕人.
是照著主的交託去引導信徒相信耶穌.
他們都是受託付來幫助信徒.
保羅更加用了這個農夫來做一個比喻.
去指出保羅他去做栽種.
阿波羅他去囂觀.
栽種和囂觀都算不得什麼.

$^{281}$都是一樣.
那意思是什麼.
說的是他們不是敵對的.
都只是神的同工.
所以在教會我們沒有人是敵對的.
對不對 弟兄姊妹.
做傳福音的他們就做傳福音.
喜歡去教主日學的就教主日學.
喜歡去唱詩的就唱詩.
我們沒有人是敵對的.
我們都是神的同工.
我們都是一樣.
因為保羅說真正叫人生命成長的唯一是誰.
唯一是神.
是那個在乎叫我們信徒生命成長的神.
弟兄姊妹 今天能夠幫助你生命成長的那個.
不是我.
也不是你旁邊那個.
誰能夠幫助你生命成長.
唯有神.
這個是保羅要回應.
當時這班信徒他們將人神聖化的時候.
要告訴他們.
真正能夠叫他們生命成長的是神.
我早年去探訪一間馬尼拉最大的教會.
這間在馬尼拉最大的教會.
主任牧師是一個美國人叫David Sumner.
他原本是做電視台的工作.
不是拍戲的 簡直是管理一家電視台.
他去蒙召菲律賓做主任牧師的時候.
就建立了馬尼拉最大的教會.
那間教會的禮堂可以坐到一萬八千人.
所以他們每個禮拜的崇拜.
就好像在紅磡體育館那樣舉行崇拜.
那麼多人去舉行.
在他們的禮堂旁邊有一座建築物.
那座建築物如果用我們華人來說.
那是一個骨灰瓶.
這個骨灰瓶用來做什麼呢.
就是給他們所有的會友.

$^{321}$如果在他們安息了之後.
他們可以將他們的骨灰安放在教會裡面.
我也跟過他們一起去探訪.
做一些很科幻性的工作.
甚至去到他們的貧民區那裡去探訪一些家庭.
我最深刻印象就是我見到當時一班小朋友.
他們在地下一灘的.
因為那些地很凹凸不平.
剛剛下完雨所以有一些水積聚在那裡.
我見到的是一班小朋友.
很開心很天真地在那些水灘那裡玩.
而那些水灘的水是全部黑色的.
所以我們跟著他去探訪很多的家庭.
然後就很期待去聽David Sumlow分享.
究竟怎樣可以令到一間教會那麼茁壯地成長.
可以發展到那麼厲害.
David Sumlow除了分享他們做了什麼事工之外.
最後他講了一個最重要的要素就是什麼呢.
恩典.
他很清楚地告訴我們.
在教會裡面都需要事工的.
都需要人去努力去事奉的.
但是真正能夠叫教會成長的那個是神.
所以弟姐妹.
今天能夠叫你生命成長的那個是神.
並且保羅之後他更加說.
將來各人要怎樣呢.
他說要照著自己的功夫得自己的賞賜.
他說栽種的阿波羅和囂貫的保羅.
他們各有不同的職責.
他們都是要向神交仗.
都會得到自己的賞賜.
這個是他們有不同的工作.
有不同的神在他們的生命裡面.
差派他們的服侍.
並且將來他們要照自己的功夫得自己的賞賜.
他的形容就是殉道所耕種的田地.
所建造的房屋.
但是都要做一件事是什麼.
就是要為自己的生命去負責.

$^{361}$保羅在這裡很清楚在說什麼.
他是在說無論阿波羅.
無論這個保羅.
他們栽種的囂貫的.
他們做完這些工作都是神差派給他們.
他們都要忠心去服侍.
忠心去服侍的時候.
將來他們都要見主面.
他們要向神交代.
但同時間他們都在告訴我們一件事.
就是說我們殉道.
我們所耕種的田地.
所建造的房屋.
即是形容我們的生命是怎樣.
其實我們都要自己努力去耕種這個田地.
我們都要努力去建造我們的房屋.
我們都要去為我們自己的生命來去負責.
所以我們需要繼續去追求成長.
而不是去維持我們的生命.
所以多倫多的信徒.
他們應當要跟隨.
那個真正能夠叫他們成長的神.
而不是那些被差派去服侍神的人.
同樣地我們都要跟隨那位.
真正能夠叫我們生命成長的神.
在新加坡堅信浸信會.
鄺建宏牧師.
鄺建宏牧師是推動細胞小組的一位先鋒.
他們在新加坡的教會.
大約有一萬人聚會.
有一年在他們教會要預備新一年的計劃的時候.
鄺建宏牧師就分享.
當日他就要為教會去祈禱去求神.
新一年的計劃究竟要做些什麼.
當他在禱告裡面的時候.
他又有一個領袖.
他發現教會很多東西都有問題.
他領袖到的是很大的責備.
他看到教會的敬拜又做得不好.
教會的小組又做得不好.

$^{401}$教會的事工又發展得不好.
所以他感受到很大的責備.
所以他就提出教會要停止小組一個月.
而這個月每一晚都會有認罪祈禱的聚會.
當我們說到小組停止一個月.
大家聽起來就沒什麼感覺.
因為覺得只是停聚會.
這個月不用團契了.
OK 沒什麼所謂 很小事.
我們月會這個月沒什麼事要談.
我們稍後會報告三月的月會會暫停一次.
因為我們沒什麼事要談.
我們要停一些聚會是很小事.
但是對於一間細胞小組教會.
他們要停小組一個月是一件很嚴重的事.
是一件大事.
因為停小組會影響整間教會的推動.
整間教會的發展.
所以對於一間細胞小組教會來講.
要停小組是很嚴重的事.
所以有一位同工就去問鄺健雄牧師.
鄺牧師為什麼你這麼勇敢.
為什麼你這麼敢停小組一個月.
鄺牧師的回應是這樣說.
老闆吩咐你做事.
然後你照做有什麼勇敢.
他說反過來老闆吩咐你做事.
你不去做你才是勇敢.
對不對 弟兄姊妹.
所以今天弟兄姊妹.
我們的生命.
我們不是建立在人的身上.
我們建立在神的吩咐上.
神吩咐我們的.
我們就是照著祂做.
保羅最後說他用建造房屋.
樓房來做一個比喻.
保羅說神給他們的使命恩賜.
就好像一個智慧的公頭.
公頭安排了保羅去撒種.

$^{441}$他立了一個根基.
然後阿波羅去挑冠.
在上面去建造.
他說已立好的根基是什麼.
就是耶穌基督.
是沒有人能夠可以立別的根基.
教會的根基就只有耶穌基督.
今天我們教會的根基.
不是有些什麼人.
而是建立在耶穌基督.
不是因為我在這裡久一點.
又或者我在這裡參與多一點.
又或者我說話大聲一點.
說話小聲一點.
而是在說什麼.
說教會的根基是建立在耶穌基督的身上.
沒有人可以建立別的根基.
而在建造這個根基的上面.
這個建造的工程.
保羅說了有兩種的工程.
一個工程是貴重的.
是有金銀寶石的.
一個是卑賤的.
是草木和街的.
他說無論你究竟是金銀寶石的也好.
草木和街的也好.
這些工程都必然會顯露的.
顯露在什麼時候呢?.
就是在那日子.
就是基督再來的時候.
有火顯現.
審判就會臨到.
所以今天弟兄姊妹.
你的生命怎樣去成長.
究竟你這個生命上面.
你建立的工程.
是用金銀寶石.
還是草木和街去建造.
他朝將來有一日.
都會顯露.

$^{481}$就是主回來的日子.
審判就會臨到.
我們記得.
審判由哪裡開始.
是由神的家開始.
所以貴重的工程是會怎樣.
就會傳得住.
是要得賞賜.
卑賤的工程是會怎樣.
是會被燒.
是會虧損.
就算人得救了.
工程經過火之後怎樣.
什麼都沒有.
所以弟兄姊妹.
寶萊告訴我們.
我們要在耶穌基督的根基上面.
去做好這個工程.
追求成長.
以致我們將來要見主面的時候.
這一切都要顯露.
審判都要臨到.
你到時又打算交什麼賬呢?.
所以重點是.
弟兄姊妹我們都要謹慎.
去建造我們自己的生命.
我們要追求成長.
我們不單要聽道.
更加要行道.
聽道就是我們要用神的話語來裝備自己.
行道就是.
我們真的去實踐.
去經歷神話語的真實.
弟兄姊妹.
神的話語.
從來都是走出來的.
你不走出來的時候.
你從來都經歷不到.
神的話語是多麼真實.
你從來都經歷不到.

$^{521}$神的話語是多麼有能力.
我們必須要走出來.
以致我們經歷神的話語.
有一種植物叫做毛竹.
種下去的時候.
頭四五年.
種的人每天要去淋水施肥.
好像沒有事情發生.
看不到有什麼生長.
但其實這些毛竹在頭四五年.
在地下生根.
還要生很長的根.
當這個根基建立好之後.
第五年開始.
它就會在六個星期內很快.
可以成長到有二十多呎那麼高.
弟兄姊妹.
今天我們必須要建立好我們的根基.
就是耶穌基督.
然後我們做好我們生命的工程.
是一個金銀寶石的工程.
以致我們的生命.
能夠繼續去追求成長.
要追求紮根在基督的成長.
三件事 弟兄姊妹.
我們是由從認識耶穌的話.
認識神的話來開始.
以致我們真的能夠去.
每一天都去反省神的話.
而最終我們真的能夠在真理上面實踐.
我們的生命就能夠可以.
紮根在基督裡面去成長.
今天弟兄姊妹.
當我們回教會.
我們做基督徒.
我們日復日.
我們除了去維持我們基督徒的習慣.
我們更需要去追求我們生命的成長.
以致我們真的能夠可以去榮耀神.
可以去見證神.

$^{561}$所以弟兄姊妹.
我們都要去吃飯.
不要再過那些順從肉體.
好像世人一樣的生活.
然後我們就將我們的生命根基.
建立在耶穌基督的身上.
弟兄姊妹你願不願意呢?.
去更新我們.
讓我們來到教會.
我們參加聚會.
讓我們不是只是追求.
去維持我們的屬靈生命.
而是我們都追求努力去成長.
求主幫助我們.
拿走我們生命裡頭.
我們可能還有一些屬肉體.
我們過著一些像世人一樣的行為.
讓我們真的能夠來到你的面前.
學習怎樣去跟隨你.
順從你更加要分別為聖.
最更加懇求你幫助我們.
讓我們真的在屬靈生命裡面.
我們將我們的根基都建立在你的上面.
讓我們能夠好好建造我們的生命.
我們用金銀寶石去建造這個工程.
以致我們的生命能夠更加像你.
以致我們的生命能夠更加可以去榮耀你.
求主你去幫助我們.
憐憫我們.
以致我們真的.
不單止我們的生命得以被建立.
得以被更新.
更加讓我們的教會.
都能夠成為一間榮耀主理的教會.
主我們多謝贊美你.
我們拭目以待.
是奉教主耶穌基督的生命致祈求.
阿們.
\newpage



\section{創世記 39:6-18}
\label{sec:031837}
\textbf{聖經人聖經事:約瑟一時失慎}~~~~黃克勤牧師博士
\newline
\newline
連結: \href{http://www.church.com.hk/acms/content.asp?site=cdc\&op=show\&type=product\&code=031837\&layout=sermon}{\texttt{031837}} ~~~~ 語音日期: 2026-03-16
\newline
\newline
\hyperref[sec:031836]{\small{< < < PREV SERMON < < <}}
~
\hyperref[sec:index_chronic]{\small{[返順時目]}}
~
\hyperref[sec:031838]{\small{> > > NEXT SERMON > > >}}
\newline
\newline
創世記 39:6-18
\newline
\begin{longtable}{cl}
\hline
\hline
章節 & 經文 (和合本修訂版)\\
\hline
39:6 & \begin{tabularx}{0.7\textwidth}{X} 波提乏把他一切所有的都交在約瑟手中，除了自己所吃的食物，其他的事一概不知。約瑟英俊健美。 \end{tabularx} \\ \\ \relax
39:7 & \begin{tabularx}{0.7\textwidth}{X} 這些事以後，約瑟主人的妻子以目送情給約瑟，說：「你與我同寢吧！」 \end{tabularx} \\ \\ \relax
39:8 & \begin{tabularx}{0.7\textwidth}{X} 約瑟拒絕，對他主人的妻子說：「看哪，一切家務我主人一概不知，他把所有的都交在我手裡。 \end{tabularx} \\ \\ \relax
39:9 & \begin{tabularx}{0.7\textwidth}{X} 在這家裡沒有人比我更大，除你以外，他也沒有留下一樣不交給我，因為你是他的妻子。我怎能行這麼大的惡，得罪神呢？」 \end{tabularx} \\ \\ \relax
39:10 & \begin{tabularx}{0.7\textwidth}{X} 她天天這樣對約瑟說，約瑟卻不聽從她，不與她同寢，也不和她在一起。 \end{tabularx} \\ \\ \relax
39:11 & \begin{tabularx}{0.7\textwidth}{X} 有一天，約瑟進屋裡去辦事，家裡沒有一個人在那屋子裡， \end{tabularx} \\ \\ \relax
39:12 & \begin{tabularx}{0.7\textwidth}{X} 婦人就拉住他的衣服，說：「你與我同寢吧！」約瑟把衣服留在她手裡，逃出外面去了。 \end{tabularx} \\ \\ \relax
39:13 & \begin{tabularx}{0.7\textwidth}{X} 婦人看見約瑟把衣服留在她手裡逃到外面， \end{tabularx} \\ \\ \relax
39:14 & \begin{tabularx}{0.7\textwidth}{X} 就叫了家裡的人來，對他們說：「看，他帶了一個希伯來人到我們這裡戲弄我們。他到我這裡來，要與我同寢，我就大聲喊叫。 \end{tabularx} \\ \\ \relax
39:15 & \begin{tabularx}{0.7\textwidth}{X} 他聽見我放聲大喊，就把他的衣服留在我這裡，逃出外面去了。」 \end{tabularx} \\ \\ \relax
39:16 & \begin{tabularx}{0.7\textwidth}{X} 婦人把約瑟的衣服放在身邊，直到他主人回家， \end{tabularx} \\ \\ \relax
39:17 & \begin{tabularx}{0.7\textwidth}{X} 就用這樣的話對他說：「你帶到我們這裡來的那希伯來僕人進來要調戲我， \end{tabularx} \\ \\ \relax
39:18 & \begin{tabularx}{0.7\textwidth}{X} 我放聲大喊，他就把衣服留在我身邊，逃到外面。」 \end{tabularx} \\ \\
[1ex]
\hline
\hline
\end{longtable}
$^{1}$弟兄姊妹大家好,我們又一齊讀聖經了.
聖經人,聖經事.
我們在聖經的人與事後面.
去認識上帝的手在這裡工作.
認識上帝更加多.
信靠祂更加多.
《生命說書教網絡》很開心跟大家一起讀聖經.
今次繼續講《約瑟的故事》.
聖經的事,其實約瑟的故事,人就是約瑟.
約瑟很長.
現在看到他去到埃及了.
他今次一直都是蒙福的.
有特別的手在後面工作.
保護他,令他萬事亨通.
萬事亨通,好像寫揮春一樣.
真的.
但是沒辦法,他今次一時失神.
看看他怎樣一時失神.
大家如果熟悉他的故事都知道了.
《創世記》三十九章.
他的老闆叫波提弗.
波提弗就是太陽的.
波提弗是將一切所有的都交約瑟手中.
除了自己所吃的飯,別的事一概不知.
這一句昨天是讀到這裡的.
是第六節的上半.
今天讀了下半,原來是這樣.
我都希望我有一點像約瑟.
約瑟原來是瘦瓦盡美.
原來有八個中文字在這裡.
綠色字是解釋.
波提弗將一切所有的交給約瑟,相信他.
你知道如果不是上帝的手在背後工作.
是不會發生的,為什麼呢?.
因為約瑟只不過是一個奴隸.
是買回來的一條牛.
不過很聰明,於是相信他.
但其實他是一個外族人.
本地就算是做僕人,其他親戚都進去工作的時候.
是埃及人來的,自己人.

$^{41}$但他竟然起用約瑟成為他的管家.
擺這個管家是什麼呢?.
將自己一切所有的交給約瑟.
除了自己吃的飯,別的事一概不知.
即是你搞定他.
吃的飯是什麼意思呢?.
今晚想吃豬扒,你給我弄豬扒.
給我弄一隻乳豬,燒的,皮要脆.
除了這些之外,其他東西你給我搞定他.
你不用告訴我,你做吧.
這樣就很相信這個奴隸,希伯來人.
說到這裡,因為要轉到下一件事.
所以在摩西記錄的時候,去到這裡他就這樣說.
約瑟原來瘦瓦盡美,將那些紅色字左.
這個中文的翻譯要解釋一下.
原來瘦瓦盡美.
原來是什麼呢?.
原來是這樣.
由於這個中文原來是我們有些意思的.
原來是這樣,我都不知.
原來是這樣.
約瑟原來瘦瓦盡美.
原來這個字我們會想.
原來是這樣,我都不知.
原來.
原來這個字是「是」.
由於是以往發生的事.
所以「瘦瓦」是「瘦瓦」的意思.
當時約瑟是瘦瓦盡美.
翻譯成「原來」反而是令到我們用慣的中文.
在我們的腦海裡是偏差了一點.
原來約瑟很靚,都不知.
其實不是這件事.
約瑟「瘦瓦」是瘦瓦盡美.
這個瘦瓦原來是兩個字.
就是「瘦瓦」和「瘦瓦」.
「瘦瓦」「瘦瓦」是美麗和外形外觀.
意思是說約瑟是一個很靚的外形.
即是說一看上去約瑟很有型.
外形很靚.

$^{81}$有些人看上去很靚.
或者一看上去是有些品味.
「嘩,這個挺好」.
「盡美」的意思又是美麗.
「瘦瓦」是輪廓,輪廓是指面容.
原來約瑟有兩件事是很靚的.
第一件事就是他有一個很靚的外形.
有型,即是美麗的外形,即是有型.
另外「盡美」是指他的面容.
他的輪廓是很靚的,也是用美麗這個字.
譬如我們去印度工作,去傳福音.
那些印度小朋友,個個印度小朋友樣子都很靚.
輪廓很鮮明,很靚的,男女眼睛又大.
下巴又尖尖,笑起來很好看.
男女小朋友,印度,很靚的.
輪廓是這樣,印度的女生,當然結婚後變大了就沒那麼好.
印度的女士是很靚的,有一個很靚的輪廓.
再加上她們有些衣服,有些高的.
有美麗的有型和面的輪廓,靚.
弱勢就這樣,你看下去有時我們會覺得西人好像靚仔靚女過我們.
那些印度人靚仔靚女過我們.
其實中國是不是不靚呢?靚,見慣了.
於是出事了,這事交給弱勢做管家之後.
弱勢主人的妻,即是主母,耳目送情給弱勢說.
你與我同寢吧,不看綠色的字,看黑色,藍色,紅色的字.
接著主母見到弱勢,那麼靚仔,說話又得體.
我老公那麼信任他,又那麼靚仔又有型又年輕.
自然跟其他埃及的僕人一比,真的不行,粗魯點,真的不行.
見到不好看就很喜歡弱勢.
由於有這件事發生,也就是說她的先生波提弗大部份時間都不在家.
而且可能就算在家裡也不理他,所以主母生活得比較空虛.
這件事交給弱勢做管家之後,弱勢的主母耳目送情給弱勢.
經常跟她打眼色,眨眼眨眼.
然後終於弱勢的主母開口,你與我同寢吧.
你一看到我寫藍色和括號的意思,就是希伯來文是沒有這個字的.
加這個字有什麼意思呢?令我們容易明白一點.
其實最簡單的字就是Yim,Yim就是同,就是with.
同,with,寢,sakab,sakab就是睡覺.
而且還是命令語,主母嘛,命令奴隸弱勢,跟我睡吧.
所以你與我同寢其實是沒有這個字的,只是說跟我睡.

$^{121}$由於中文想表達就加了一個吧字,同寢吧,命令式.
再加上這個感嘆號,你與我同寢吧,其實就是一起睡吧.
有命令的味道,第一,在那個地方,弱勢都空不空虛呢?.
都空虛的,因為一個人那麼遠,被埋到那裡.
做管家有沒有壓力?沒有,因為主人與主母之下,萬人之上.
不是的,你做到這個位置,下面的埃及人撐你的,刺你的,煽你的.
不過他一個人很聰明,做得很叻,大家都佩服他,他搞得定.
不過這個都叫做主人,而且老老實實說,一個年輕人.
嘩,這件事有好處,但是第一點要講的是,弱勢不從.
不從其實是maan,弱勢不從,maan是拒絕,直接,不是不從.
對不起,不可以,不是那些,是不可以,是不從的.
弱勢很直接,是拒絕,maan,弱勢maan,不行.
對他主人的妻子說,看一下,一切家務,我主人都不知道.
他媽所有的都交在我手裡,老闆波提弗那麼信任我.
所有東西都交給我,再者家女沒有比我更加大.
他沒有留下一樣東西不交給我,起碼所有工人他管.
所有保安他作主,所有傢俬儀器怎樣搞,他搞.
還有,要搞這些,沒東西不留下,帳房,即是錢我管,嘩,.
沒有留下一樣不交給我,只有留下你.
這個主人波提弗有沒有跟他說,我家裡的東西全部交給你.
不過我除了我太太不交,其他都交了,沒有講這些.
是常識understood,一般知道是應該這樣.
弱勢自己說,老闆將所有東西都交給我,不過只有你.
因為你是他的妻子,我是不能夠做這件大惡.
第二點要講,弱勢厲害,弱勢有自知之明.
第一,弱勢拒絕,厲害,第二,為什麼拒絕.
因為他知道老闆很看我,我不能夠做這麼大的壞事.
這個大惡用希臘的字,Raw,不是希臘,希伯來的字.
Raw,就是壞事,壞事,我不能夠做這麼壞的事.
去得罪神,咦,他不是得罪老闆,是得罪神.
我做這件事是得罪神,這個罪字譯得好.
Qatar,Qatar,希伯來文如果是譯做英文字母.
如果見到CH是讀K,T有時讀D,有時讀T.
這就Qatar,我不能夠Qatar我的神,Qatar的意思.
直接就是罪,Sin,我不能夠做一件這麼壞的事.
這麼壞的事是什麼呢?是我對於我老闆不尊.
而且這件事在神來講是罪,就是這樣.
後來這個女主人不肯拒絕,天天跟弱勢講.
有機會跟弱勢講,但弱勢都不聽從她.
不與她同寢,也不和她在一處,弱勢厲害了.

$^{161}$跟她一起這麼主動,她連跟她一起她都不去.
誰知有一天,第三點要講的,就是有一天.
弱勢進屋裡去辦事,做管家盡心盡力當然走來走去.
進屋裡辦事的時候,家中人沒有一個在那屋裡.
第三點要講的,就是家,家這句字Baif.
這有兩點要拉長,Baif是房屋,其實這個家字是屋.
House,屋裡沒有其他人在屋裡,重複講每兩次屋字是相同的.
屋裡沒有其他人在屋裡,所以他強調屋裡沒有人.
屋裡沒有人,夫人拉著衣裳說,趁沒有人你與我同寢吧.
這個你與我同寢吧一樣,是命令來的,一起睡.
夫人拉著她的衣裳,沒有捉住她的手,是拔了她的衣服.
一起睡,要一起睡,弱勢把衣裳丟在夫人那裡.
跑到外面去了,金蟬鐵鶴脫了衣服,就把這件衣服脫了出去.
這裡是第三件事,屋裡沒有人,而她進去辦事.
剛剛她才說她不和他在一處,弱勢要醒,怎知道她一事失神.
屋裡沒有人,而她都進去辦事,即是她進去她主人間屋或房裡.
怎知更加被人拉著,本來她的主母只是眼目打眼色和口講.
但今次見到沒有人,她就開始捉住她,又沒有捉住手腳.
就拉著她的衣服,脫了衣服跑了.
夫人看見弱勢把衣裳丟在她手裡跑出去.
去到這樣脫了衣服都不肯,這女人覺得很受羞辱,沒有面子.
於是就叫了家裡的人來,有沒有人?救命呀!救命呀!.
於是其他的家人,傭人就進來了.
然後她說,你們看,他帶了一個希伯來人,他是指我老公他.
這些字會指出這個女主人和她丈夫波提弗的關係很有可能比較遠.
所以她會說他帶了一個希伯來人回來,他帶了一個希伯來人.
進入我們家裡,要戲弄我們,他到我借你來要與我同寢.
我就大聲喊叫,叫大家來,他聽見我放聲喊叫就把衣裳丟在我這裡.
跑到外邊去了,夫人就把弱勢的衣裳放在自己那裡.
等著他主人回家了,就對他如此說,你帶到我這裡來的那個希伯來人.
進來要戲弄我,於是我放聲喊叫,他就把衣裳丟在我這裡,跑出去了.
這次弱勢被人冤枉了,其實他就一時失神,辦事而已.
他都知道不要走那麼近,誰知這次辦事只有他一個人.
竟然進去屋裡做事,可能看些東西,執一些東西,就死了.
真的不知道怎麼辦,我就把衣裳丟在這裡,我拉他,他就脫了衣服走了.
我想問問大家,如果你是弱勢你會怎樣做呢?.
其實弱勢在這件事上最危險失敗是什麼呢?就是他一時失神.
失神是什麼呢?就是屋裡沒有人,他竟然進去,如果他不進去就沒事.
第二,如果他被他拉著衣服,他馬上脫了衣服就走了.
如果弱勢有什麼辦法呢?不脫了衣服,一味地按著,一味地按著.

$^{201}$然後大聲喊有沒有人,有沒有工人在外面.
如果弱勢喊先,工人走進來,見他拉著衣服,可能又保護到自己.
誰知弱勢這一下聰明反被聰明誤,我覺得他馬上脫了衣服,金蟬片殼走了.
於是就被人拿著把柄,如果他死命地按著,有沒有人,有沒有人.
那些人走進來的時候,救命救命,他喊先,可能會好一點,但是他沒有.
可能他都覺得我光明正大,沒什麼,於是在這件事上就出事了.
溫習番,弱勢拒絕,我們都要學習,拒絕.
第二,這件事我得罪人是一件很壞的事,我主人那麼信任我.
第二,另外就是我這樣是罪,得罪神.
第三,就是他那次一時失神,房裡沒有人,他進去.
也不是自己喊先,是被那個女人喊,於是他就被人拘捕了.
之後主人回來聽到,如果你是主人波提佛,你聽到會怎樣.
超級生氣,必然的,因此我今天跟大家看的是,弱勢是一時失神.
有沒有一時失神的時間,當然有,一時失神是會有時沒什麼事.
對不起就算了,但是一時失神有時會搞到是彌天大禍.
就看這件事是什麼,所以什麼時候都要盡量謹慎.
確實是要盡量謹慎,不過話說回來,今天講這番話的人都不是很謹慎.
有時都會一時失神,或者從小到大的人可能個性比較放鬆.
或者防備性比較鬆弛一點,一點都不奇怪.
有時也是他原生家庭影響的人,有時會去到某些位他會放鬆了.
就在乎他的爸爸媽媽,或者他原生的家庭是怎樣模造他出來.
有些是他原生家庭令到他去到那個位很謹慎,特別謹慎.
有些是令到他去到那個位是特別鬆弛了,都不奇怪的.
但是無論怎樣都要自己盡量去小心,因為今天確實是人生人世很險惡.
不要去害人,不過都要提防,因為有些人是很邪惡.
有些人是很卑鄙,有些人是很卑劣的,他們真的會作弄你的.
有些人作弄你是為了獲得某些利益,有些人是真的心地壞.
或者他的原生家庭影響他,或者他過往的歷史影響他.
他去到某些位是變得很邪惡的,不留手的,很邪惡的.
明知做得不對,但是他都會做的,很多人都是這樣的.
平時沒事的,不過去到某些位他就會了.
有時候就是這樣,所以要很小心,小心行得萬年船的.
小心行得萬年路都是要特別,我們求主幫助我們.
在約瑟這件事上,我們有一個學習,有些警惕.
不過很感恩,神是在這件事上,讓他這個一時失神是發生.
而這個一時失神令到他很大的災禍,在前面要坐牢,差點沒命了.
不過在那裡他又模造了,練習了他將來是更加叻.
所以有時候不要怕,有時候真的如果有些什麼跌低了.
又差了,起來,在那裡一味信靠神,將來神是會使用你的.
因為他模造你之後,你才更加硬朗可以承擔多一些工作.

$^{241}$在這件事上我們體察到約瑟一世英明.
他已經擺脫了以前是公子哥兒,一個少爺仔那種情況.
來到沙漠,在曠野綁住來到埃及,脫了衣服給別人看一次.
然後決定買他,然後進去做一個下人,一個奴隸.
而他辦事是賣力,神祝福,有你的祝福最重要,有你的同在最重要.
他開始升上去成為管家,甚至主人那麼信任他.
不過我們看到仍然有邪惡的力量,有魔鬼在當中工作.
要害他的,不過他雖然一路拒絕,不過那一次失神.
就被人捉到痛腳了,以致影響他很慘.
主要幫助我們,無論要些什麼將會擺在我們面前.
我們都繼續信靠你,自己小心.
順天禱告,奉耶穌基督名求,阿們.
\newpage



\section{雅各書 2:14-15}
\label{sec:031838}
\textbf{愛的施予}~~~~陳雅莉傳道
\newline
\newline
連結: \href{http://www.church.com.hk/acms/content.asp?site=cdc\&op=show\&type=product\&code=031838\&layout=sermon}{\texttt{031838}} ~~~~ 語音日期: 2026-03-22
\newline
\newline
\hyperref[sec:031837]{\small{< < < PREV SERMON < < <}}
~
\hyperref[sec:index_chronic]{\small{[返順時目]}}
~
\hyperref[sec:031839]{\small{> > > NEXT SERMON > > >}}
\newline
\newline
雅各書 2:14-15
\newline
\begin{longtable}{cl}
\hline
\hline
章節 & 經文 (和合本修訂版)\\
\hline
2:14 & \begin{tabularx}{0.7\textwidth}{X} 我的弟兄們，若有人說自己有信心，卻沒有行為，有甚麼益處呢？這信心能救他嗎？ \end{tabularx} \\ \\ \relax
2:15 & \begin{tabularx}{0.7\textwidth}{X} 若是弟兄或是姊妹沒有衣服穿，又缺少日用的飲食； \end{tabularx} \\ \\
[1ex]
\hline
\hline
\end{longtable}
$^{1}$靜姐妹 早安 平安.
過去幾個星期 我們從農曆新年開始.
我們就看八福 是不是啊?.
就是在《馬太福音》裡面 八福 耶穌所說的.
知道真正的福氣不是一份只是物質上的富足.
之後梁銓導友帶我們看了怎樣在軟弱裡靠主剛強.
仍然是一種在我們的價值觀裡面很顛覆的.
在軟弱裡面都可以再次剛強.
上星期 劉契聰宣教士又帶我們看了這個世界的需要.
並且開始教導我們 不單止是頭腦上知道一些東西.
原來要落手落腳 原來要落水的 要參與的.
到今天 我都想再進一步和大家去看看.
第一次去思想我們其實真實在每一天裡面.
是怎樣可以實際地參與我們身邊的人的需要.
去看到神所關心的 看到神要我們看到.
哪些有需要的人的 好嗎?.
所以今天想和大家來看《雅各書》.
《雅各書》是一本智慧文學 其實看起來很普通的.
如果照讀 你每一句都明白的 是不是?.
就像剛才主席讀的時候 你不會覺得很難明很難解.
你都不知道他講什麼 你不會這樣的.
如果我們要思想聖經裡面 你選一些你覺得最難明的經文.
可能你會很容易選《啟示錄》那些 很難明很深.
《雅各書》這些一看就明白.
但有時候我們會慢慢再仔細去看聖經的時候.
我們就會發現真正難明的經文.
有時候往往是看起來表面很容易的經文.
《雅各書》看起來是很容易明白的.
每一隻字讀下來你都明白.
但往往就是這些表面很容易明白的經文.
今天我們就值得問 好像我們都知道的.
我們都不是不懂 我們都不是不明白.
對我們來說還有什麼震撼性呢?.
今天對我們的信仰這些經文 你看得明白.
但還為你帶來什麼挑戰呢?.
我相信聖經的作者是不會寫無謂的東西的.
如果是一些人人都很容易做到.
人人都懂 人人都知道的東西.
他沒必要再提醒信徒.
所以既然我們認同 我們不是全然明白.

$^{41}$或者不是全然知道的時候.
我們就值得花時間再去看一些表面上看起來好像很容易的經文.
但還有些什麼是我們未明白 或者未知道的.
盼望神藉著《雅各書》今天再提醒我們.
今天我的題目是《愛的詩語》.
在一開始的時候 二章這裡.
《雅各書》說:我的弟兄們.
你們信奉我們榮耀的主耶穌基督.
就不可以按著外貌待人.
《雅各書》在這裡形容.
當時 如果你一路看下去.
二章 一路二節三節四節 一路看下去.
你就會發現 當時似乎有一些信徒.
他們會看重那些穿衣服很華麗的人.
看起來很高尚的人.
就輕看那些看起來貧窮的 不富足的人.
有時我們會想 這個都很小事.
當然有些偏見 有時我們對人.
但是聖經說 這個不是小事.
甚至《雅各》說 這是一種偏見.
而且這種偏見是帶著惡意的.
是一種對別人惡意的評斷.
嘩 很嚴重 是吧.
《雅各》在這裡說.
你們不可以按著外貌待人.
今天當我們活在豐富或者生活豐盛裡面的時候.
我們都看到世界上有很多的需要.
但是在實際的眼光裡面.
我們是怎樣看待這些生活上貧窮的人呢?.
這些情節 當時有弟兄.
有人進來會堂穿得很漂亮 就讓他坐高位.
看起來貧窮的人 《雅各》形容.
當時就讓他坐在腳椅下面.
這些情節不是當時昔日在猶太會堂裡面發生.
乃是到今天仍然都有機會發生的.
我記得有一次 我有機會來到英國之後.
我就南下 剛好我南下的時候碰到是主日.
我就在那個城市找了一間教會進去聚會.
嘗試參與那一間教會的主日崇拜.
我們崇拜當敬拜隊在唱詩歌的時候.

$^{81}$突然間有一位衣衫藍柳的人從門口走進來.
當然大家左望右望.
相信你憑他的外貌.
你會猜到他可能是一位無家者.
那位衣衫藍柳 頭髮蓬鬆的人.
一路從後面走進來.
然後就坐在這個位置 最前的第一行.
你能夠想像他身上還帶著一陣味道.
如果今天這一刻就是一個這樣的人走進來我們教會.
我不知道大家有什麼反應呢?.
可能我們都不會說出口的 都會給他一個位置坐.
但有時可能我們很直接 心裡面很即時的反應.
你未必有說出來.
可能已經反映了我們心裡面做了一些選擇.
這個似乎是《雅各書2章》正正要挑戰我們的地步.
我們的信仰是帶給我們一種張力.
有時是一種拉扯.
當我們說 雅各說是你們信奉我們榮耀的主耶穌基督.
是怎麼相信的呢?你是帶著什麼信心的呢?.
是帶著一份偏心.
還是真的能夠帶著一份愛心去信奉我們這位榮耀的主基督呢?.
很大的挑戰.
我記得曾經聽過一堂道.
同樣都是那個教會遇著一位無家者走進來.
坐在一個弟兄姊妹的旁邊.
陣陣的異味一路發出.
道完崇拜之後.
講完下去問這位弟兄姊妹.
你剛才OK嗎?.
有一位無家者 有這麼大陣的異味坐在你旁邊.
那位弟兄姊妹回答那個牧者說.
我聞不到異味.
我只是聞到主耶穌馨香的香氣.
有人願意踏入來教會參與崇拜.
我們是怎麼對待每一位呢?.
雅各在這裡說.
我親愛的弟兄姊妹 請停.
神起不是揀選世上的貧窮人.
使他們在信心上付諸.
並且承受祂所應許給那些愛祂的人的國嗎?.

$^{121}$在第一世紀的基督徒.
我相信雅各書的讀者.
在第一世紀大部分都是貧窮的人.
我們會發現保羅在《哥林多前書》都說過.
你們蒙召的按著肉體的智慧不多.
能力的不多 尊貴的都不多.
所以我們能夠想像當時的信徒都落在貧窮.
患難裡的偏向比較多.
在早期的教會歷史裡.
福音的屬靈應許.
對於那些真的活在貧窮裡的人.
其實是很吸引的.
因為當他們落在貧窮裡.
他們被撇棄 被人歧視 被人侮辱的時候.
福音卻鼓勵他們 你是被愛的.
你是被尊重的 你是珍貴的.
確實對當時的貧窮人來說是非常吸引.
所以貧窮人或者窮人.
可以幾乎說是一個術語.
就是指當時經濟上雖然被逼迫.
但是靈性卻是很敏銳的人.
這些就是窮人.
不過我們能夠知道.
雖然雅各說 神豈不是揀選世上的貧窮人士.
他們在信心上富足嗎.
那不是意思說 只是揀選貧窮人.
所以今天也不要誤會.
我不夠貧窮 神今天不揀選我.
我不可以有錢 是不是.
不是 雅各沒有說 只是揀選.
我們能夠相信.
不是說有貧窮人 或者你本身是貧窮.
就一定是信心爆棚 不是這個意思.
但確實 無論雅各還是主耶穌.
都曾經有提醒當時想跟隨祂的人.
如果過分的富足 過分的財富而你不懂得.
將神放在第一位的話.
確實財富是容易成為你的絆腳石.
成為阻礙跟隨主耶穌教導的 是不是.
如果大家記得.

$^{161}$福音書也記載過.
耶穌怎樣提醒財主 是不是.
但祂就悠悠瘦瘦地走了.
沒有辦法跟隨主耶穌.
當然我們不是要刻意將富有人排除在主的國度之外.
但確實主耶穌和雅各.
我相信祂在這裡也提醒我們要留心.
當財富成為你的阻礙的時候.
你要怎樣選擇.
神是揀選世上貧窮的人.
使他們在信心上富足.
昔日可能是很少數的富有人.
地主,商人 他們可能財雄勢大.
然後就壓迫當時窮困的人.
雅各是在提醒這一點.
不要仗著你有的財雄勢大.
就可以看不起去壓迫貧困的人.
昔日是這樣 到今天21世紀也是這樣.
沒有分別.
現在全球財富的分佈是極度不平均的.
在最頂層 全球最富有的1\%的人口.
全球最富有1\%的人口.
是擁有了全球45\%的財富.
而在最底層 全球一半的人口.
50\%的人口.
是擁有全球不夠2\%的財富.
你能夠想像昔日是這樣 今天也是一樣.
雅各要提醒讀者的是.
許多貧窮的人 其實是蒙上帝所揀選的.
目的是要清楚責備那些以自高.
以財富自高的人.
以財富去歧視貧窮人的人.
幾年前 我認識了一位媽媽.
這個媽媽是帶著一個女兒來到英國.
她的女兒和我的小女兒一樣大.
我知道這個媽媽的生活是很艱難的.
我知道她一個人要打三份工去養她的女兒.
她很低調 她盡量不讓人知道她的身世.
因為她曾經在不同的地方遭到歧視 受到傷害.
有不好的對待.

$^{201}$所以她很怕別人再用歧視的眼光去看她.
她很少提起自己的事.
所以我認識她之後 跟她的相處都是很低調.
為了保護她 不要再受那些無謂的傷害.
這個媽媽很堅強.
她不會隨便領受別人的恩惠.
雖然她很有需要.
但她不會隨隨便便去領受別人提供的恩惠.
她是一個愛主的基督徒.
她只是很努力默默地生活.
有一次這位媽媽跟我說.
我有機會跟她一起的時候.
她就跟我說 如果你知道附近哪裡有教會派食物.
你可不可以通知我.
有時候教會派食物會放在Facebook群組,WhatsApp群組.
她說 如果你看到 可以通知我嗎?.
我就說 好啊 如果我看到就通知你吧.
她就說 因為每次我看到訊息 看到那些帖子.
我就馬上趕去那間教會 去到已經派了.
有時候我遲了幾個小時 去到已經沒有了.
她說 如果你收到 你通知我 看我趕不趕得及去拿.
我說 好啊 當然可以 幫你留意一下.
留意到通知你.
接著我就多口問了一句.
我說 是不是很艱難?.
於是她笑著對著我.
她說 沒事的 我知道會有供應.
我很感動.
我看到她 我看著她.
哇 何來你的信心.
真是很大的信心.
之後那天 我剛好當時跟教會團契的職員會開會.
我就一路開會 談一下事務.
那個團契是每個星期都有人煮飯的.
團契頂多就是回來吃飯 吃完就開始團契.
我就留意到每一次團契完了 都是有飯盒剩的.
每次都吃不完的.
我就忍不住開完會之後.
我就跟團契的職員說 不好意思 我想問一件事.
我是不是每次都有飯盒剩的 其實我們這個團契.

$^{241}$可不可以留一個給我?.
我想給一個有需要的家庭.
那些職員就說 當然很好 沒問題 全都沒問題.
你拿一個 如果有剩的我們全都給你都可以.
其實我們每次都是決定 姐妹 拿吧 拿吧.
不要剩的.
有飯盒 可以給這個媽媽.
誰不知幾位團契職員裡面 其中一位姐妹說.
我想起一件事.
她說 其實我弟弟是做蔬菜公司的運輸的.
而這間蔬菜公司的批發.
是每個星期都會給一大箱蔬果.
所有的員工 每人都會留一箱.
而我弟弟是每個星期自己吃完.
派給附近的鄰居 都吃不完的.
每次都吃剩很多 很浪費的.
她說 不如這樣.
我叫我弟弟拿給你.
如果他拿起有需要的.
分一些給那個媽媽.
如果可以的 方便的 就給我吧.
到第二天.
我一打開我的門口.
就有一箱豐富的蔬果 擺在我門口.
她的弟弟通知我.
這位姐妹的弟弟通知我.
那箱蔬果擺在你門口了.
我看著那箱蔬果 真是很讚歎.
上帝啊 你真是既信實又奇妙.
短短三天.
第一天我和姐妹聊.
第二天和職員開會.
第三天 就有一箱蔬果擺在我門口.
從那天開始.
每個星期.
我都會收到一箱蔬果擺在我門口.
去給這位媽媽.
那位媽媽每一次來拿蔬果的時候都會跟我說.
我什麼都不懂.
我又不懂很多聖經.

$^{281}$我只是很單純知道主會供應我.
我只是努力倚靠神.
她告訴我 每一次我來你家拿蔬果和飯盒的時候.
我都會想起聖經裡面一段記載.
就是 她連哪段經文都不能說出來.
她只是說 我記得勝經裡面有一段說.
有一個窮寡婦 她失去了丈夫.
家裡只剩下一瓶油.
神卻利用先知 使用先知去她家.
那瓶油 到極都有.
她說 我只記得這個故事.
大家還記得嗎?.
還記得三天前.
這位姐妹一開始.
她跟我要求的是什麼?.
她只是想要一個訊息.
她求的只是可不可以一個通知的訊息.
神的公認卻是超出了我們所想所求.
她只是想要一個通知.
一個很卑微的呼求.
但上帝的公認.
卻讓我們沒辦法不相信.
上帝到今天仍然是看顧她的兒女.
弟兄姊妹.
今天 或許我們未必人人落在這麼艱難的地步.
或許有也不奇怪.
或許我們中間也有人生活得很不容易.
很缺乏.
甚至看著每天的生活都憂愁.
但憂愁 但這些困難和缺乏.
不是叫我們只懂得憂愁.
而是叫我們盡快轉眼仰望上帝.
我們的上帝是連小麻雀都會看顧.
我們的上帝是連野地的百合花.
都會給她華美的裝飾.
她怎會不看顧我們呢?.
這位姐妹在世上的人眼中她是貧窮的.
但她的信心在上帝的眼中是富足的.
因為我從她身上 我很真實地體會到.
聖經裡所說的 貧窮的人卻在信心上得著富足.

$^{321}$是她讓我真實地經歷了主耶穌的教導.
《雅各書》裡所說的.
耶穌說:不要憂慮.
我們吃什麼喝什麼穿什麼都是外邦人所求的.
你們需用一切的 你們的父都知道了.
你要先求他的國和他的義.
這些東西都要加給你們了.
所以不要為明天憂慮.
明天自有明天的憂慮.
一天的難處 一天當就夠了.
頂智我不知道你今天有什麼令你感覺很缺乏.
但上帝會揀選世上貧窮的人.
使他們在信心上富足.
承受神的國.
那些每天生活在窮困邊緣的人.
他們憑著信心看到神的國.
貧窮人在神那裡經歷富足的信心.
那麼富足的人呢?.
在這裡 雅各也提醒.
或者對富足的人帶出了一個挑戰性的問題.
雅各說:我的弟兄啊.
如果有人說自己有信心沒有行為有什麼益處呢?.
這個信心能夠救他嗎?.
如果在弟兄姊妹中間有沒衣服穿的.
缺少日用飲食的.
你們中間有沒有人跟他說.
平平安安去 願你穿得暖 吃得飽.
他不給他身體所需要的.
那有什麼益處呢?.
信心也是這樣.
沒有行為是死的.
當我們想到要去關心有需要的人.
或者幫助窮困的人.
我們很自然 很快第一時間的想.
就是要跟愛心有關係.
但是呢 你會發現雅各在這裡.
是說跟信心有關係.
雅各為什麼會這樣說呢?.
雅各為什麼會說跟信心有關係?.
為什麼幫助有需要的人.

$^{361}$說愛不就行了嗎?.
為什麼是說信呢?.
似乎雅各.
我們經常說聖經作者不是亂寫的.
似乎神藉著雅各提醒我們.
這一群有信心.
跟隨上帝 跟隨耶穌的人.
我們這一群.
用信心跟從上帝的人.
我們是怎樣顯出我們的信心呢?.
如果有一個有需要的弟兄姊妹.
就站在我們面前.
我們是不是只是說平平安安的去呢?.
願你穿得亂 吃得飽呢?.
是當時猶太人常用的字句.
平平安安的去吧.
翻譯成今天的術語就是.
祝你好運.
雅各說這個是無益處的.
他說有什麼益處就是無益處.
這些空洞的說話當然對人來說是無益處.
對對方無益處.
但我相信雅各不是只是說對對方無益處.
更加是說對你說出這番說話的人.
不會帶來任何屬靈的益處.
所以似乎雅各是在挑戰信徒.
當我們說我們憑信心跟隨神的時候.
如果我們沒有外顯的行為.
那麼我們的信心是不是真信心呢?.
你看到雅各是一直用反問句.
是讓我們自己去思想.
在舊約的時候 以賽亞先知都這樣形容.
當時以色列人天天尋求神.
樂意明白神的道.
還說他們很仗行義的國.
以賽亞書是這樣形容.
說他們未離棄神的典章.
他們向神求問公義的判詞.
但是神卻不回應他們.
他說他們自己刻苦己心還禁食.

$^{401}$但是神都不回應他們 為什麼呢?.
在以賽亞書那裡說.
他們雖然有很多宗教禮儀.
很樂意親近神.
但是他們禁食的時候.
卻仍然追求私利.
他們沒有顧念有需要的人.
於是耶和華神就藉以賽亞先知告訴他們.
你們禁食行很多宗教禮儀.
但是我所要的禁食是什麼呢?.
豈不是你們將食物分給飢餓的人.
將流浪的窮人接到家裡.
見到赤身的就給他衣服遮體.
不隱藏自己 避開自己的骨肉嗎?.
神要他們將真正的意義.
真正跟隨神的意義.
注入宗教禮儀裡.
不是說敬拜神不重要.
不是說禁食祈禱不重要.
不是說如何認真守神的律法不重要.
不過如果你沒有了這些.
那些宗教儀式都不是神所預立的.
神說除非你們敢.
我才會預立英文以色列人的求告.
今天我們如何看待我們身邊有需要的人呢?.
到底幫助有需要的人.
如何顯出信心呢?.
今天可能神都呼召我們要有行動.
雅各一再提醒我們.
信心若沒有行為是死的.
不是說我們要靠行為得救.
很多人都會說雅各這段說話.
是跟保羅所說的恩信稱義違反的.
不是,我們知道聖經的作者沒有矛盾.
兩者都沒有矛盾.
他不是說我們要靠行為得救.
不是這個意思.
恩信稱義仍然是我們信仰的核心.
不過當你信了之後.
當你憑信心跟隨神之後.

$^{441}$是應該有外顯的行為的.
他指出的是真正的信心.
是會流露出愛心的行動的.
在我所住的社區裡.
有一個叫做Soup Run的行動.
送湯行動.
就是在星期日的下午.
每個月一次.
在星期日的下午.
首先完了崇拜之後.
弟兄姊妹會集合去預備三文治.
會預備咖啡茶.
會去煲湯.
平均預備70至80份這樣的食物.
然後到晚上就去到特定的地點.
在那個社區的特定的地點.
去派發食物.
還有一些物資的.
睡袋啊 夏日用品啊 衫褲鞋襪啊.
去給一些無家者.
給一些流浪的人.
給一些貧困的人.
這個機構沒什麼資金.
Soup Run這個機構叫Soup Run Trust.
他沒什麼資金.
全部的資金來源都是靠有心者去捐助.
靠有人去奉獻.
過去幾年.
我很感恩.
每個月.
我都可以和我的女兒一起去參與.
這個派發物資給露宿者的行動.
去關心社區上面有需要的人.
其實從一開始.
是我的女兒帶我參加的.
是她看到街上有很多露宿者.
在城市裡面.
她和我一起去街的時候.
她經常看到有些人躺在街上.
然後她很想用實際的行動去幫助有需要的人.

$^{481}$去分享愛.
然後無意中發現原來有一個這樣的活動.
當時12,13歲的她不知道可以做什麼.
她只知道如果她可以參加.
或者可以有一份去關心這些有需要的人.
於是她和我說:媽咪,你可不可以陪我去參加?.
我很感激神.
每次當我們一起去接觸這些流落街頭的人.
很細緻地聽他們所分享的故事.
他們所分享的經歷.
他們每一個人都有他們自己背後的故事.
為什麼會成為無家者.
讓我們真正地去感受.
每一個生命都是值得被尊重.
哪怕他們已經落到.
沒有人會看到他們.
他們是被遺忘的一群.
他們是被這個社會邊緣化的一群.
沒有人會去關心他們.
沒有人會care.
甚至care他們的生和死.
但是當有人願意付出時間.
付出心力.
讓他們能夠稍微感受到一點溫暖.
一點愛.
他們仍然相信原來他們是被看見的.
他們是有價值的.
當中有一些我們說是常客.
每次都見到他.
他很樂意和我們分享他的難處.
甚至我們已經出隊幾年了.
我們會記得對方的名字.
記得他的家庭狀況.
他的背景.
當然除了派發物資.
我們是和他分享主耶穌的愛.
我們會將主的愛帶給他們.
和他們傳福音.
當中甚至有一些是回到本地教會.
也有一些.

$^{521}$因為我以前服侍的教會.
他們有英文課.
甚至帶了一些回到教會的英文課.
我們又會在職場為他們禱告.
每一次派完離開的時候.
我最不想說的一句話.
就是「See you next time」.
我不會想說這句話.
我包不得我從此以後見不到他.
我包不得每一次我離開之後.
在下一次之前.
他就找到工作了.
找到很合適的地方給他賺錢.
能夠賺生活費.
以致他可以回復正常的生活.
我不需要再見到他.
今天無論是我們.
或者我們的下一代.
已經在很富裕的生活中成長.
就算不是很豐足.
但是要他們體會什麼是貧窮.
切實地體會什麼是貧窮.
我們的下一代已經很困難.
所以很盼望.
透過很親身,很靠近的服侍.
能夠只能夠很稍微地體會.
什麼是貧窮.
也讓我更深的感恩.
去明白聖經裡所說.
當紀念主耶穌的話.
施比受更為有福.
今天,如果你能夠施,你是有福的.
雅各說.
憑著信心去做.
他說經上記著.
要愛倫如己.
你們如果切實遵守這個至尊的律法.
你們就做得很好了.
當我們去施禦的時候.
不是因為我們被人苦逐.

$^{561}$因為他很可憐.
我們要去幫助他們.
讓自己的良心好過一點.
我有份幫助你的.
不是,一點都不是.
我們可以用實際的行動去幫助別人.
是因為我們已經得著神的愛.
是因為我們已經得著從神而來的愛.
才知道我們可以愛倫如己.
這是至尊的律法.
耶穌說第一大誡命是.
你要愛主你的神.
盡心盡性盡力盡力愛主你的神.
其次也相仿效.
就是要愛倫如己.
愛人如己.
這是至尊的律法.
剛才我們有說過在路加福音.
有一個官來問主耶穌.
他怎樣才可以承受永生.
主耶穌說.
他跟耶穌說.
我已經遵守了很多律法.
我又沒有犯法.
還孝順父母.
然而主耶穌說你還欠缺一件事.
你欠缺一件事就是去周濟窮人.
聽到這句話之後.
那個財主就憂愁了.
因為他很富有.
耶穌就下了一個這樣的結論.
有錢的人信神的國.
真是何等的難.
駱駝穿過針的眼.
比財主信神的國還容易.
我相信今天主耶穌不是說財主不能進天國.
祂只是對我們有同樣的呼召.
盼望跟隨耶穌不要讓財富成為我們絆腳石.
相反要根據主耶穌的教導.
跟隨主耶穌的教導.

$^{601}$所以當過去一直思想.
有需要的人在我們附近.
在我們中間.
在我們身邊甚至乎很近的時候.
我就跟教會商量.
跟執事商量.
我們教會是否應該要成立慈慰的奉獻.
其實一直以來的姐妹.
我看到你們為教會不同的需要作出奉獻.
慈慰金是奉獻給神.
雖然好像是給人.
但其實這不是一個普通的捐助.
乃是奉獻.
你仍然是奉獻給神.
只是神藉著你.
藉著你所奉獻的.
去實際上提供實際上的經濟援助.
給有需要的人.
所以慈慰不僅是一種實質的幫助.
更加是一種屬靈的同行.
成立這個慈慰奉獻.
其實最終目的.
都是要將從主義來的.
奉獻給神.
藉著這些獻金.
去幫助社會上有需要的人.
或者我們教會中有需要的人.
甚至再擴大深一層來看.
慈慰不僅是經濟上的援助.
慈慰可以說是一種事工.
其實是一行事工.
因為當我們去關心有需要的人.
陪伴他們同行.
發現他們有需要.
我們不僅是投降實質的需要.
更加是我們陪伴,禱告,同行,紀念.
這都是在作一個好見證.
而且是我們信心的見證.
這個到最終.
當然是對上主愛的回應.

$^{641}$這是主耶穌的教導.
是主耶穌的呼召.
我要總結了.
當教會決定要成立.
慈慰金奉獻的時候.
我就想起了一段經文.
耶穌說.
當人子在他榮耀裡.
同在眾天使來臨的時候.
要坐在他榮耀的寶座上.
萬民都要聚集在他面前.
他要將這些萬民分別出來.
好像牧人分綿羊,山羊一樣.
會將綿羊安置在右邊.
山羊就安置在左邊.
於是王會對右邊的說.
你們這個蒙我父賜福的.
可以來承受創世以來.
為你們預備的國.
因為我餓了.
你們給我吃.
渴了.
你們給我喝.
我流浪在外面的時候.
你們留我住.
我赤身露體的時候.
你們給我穿.
我病了.
你們來看顧我.
我在監獄裡的時候.
你們來看我.
二人就回答.
主啊!我什麼時候見你餓.
給你吃.
渴的時候給你喝.
我又什麼時候見過你流浪.
在外面留你住.
或者赤身露體的時候.
給你穿呢.
又是什麼時候見你病的時候.

$^{681}$或者在監獄裡的時候.
我來看過你.
王就回答.
我實在告訴你們.
這些事你們做在我弟兄中.
一個最小的身上.
就是做在我身上.
我實在告訴你們.
這些事你們做在主的弟兄中.
最小的身上.
就是做在主的身上.
願我們他日都是綿羊.
在王的右邊.
我們同心吐告.
真是主耶穌.
感恩你賜下你寶貴的話語.
你的話語真是很寶貴.
你的話語真是來呼喚我們.
再一次提醒我們.
當今天我們仍然在御苦期裡.
我們說去思想主耶穌的犧牲.
思想主耶穌為我們的捨己.
神今天來都呼召我們.
怎樣去為你捨棄.
求主你幫助我們.
再一次回到你的話語當中.
去聽得見你對我們今天的呼召.
我們是怎樣去用信心跟隨你.
信心若沒有行為是死的.
今天你對我們每一個人.
可能都有特別的呼召.
可能都有給我們特別的體見.
你要我們用你的眼光去看見這個世界的需要.
願主你幫助我們.
讓我們能夠體恤你的心腸.
體會主耶穌你所看見這個世界的需要.
又帶領我們.
讓我們真是甘心樂意.
做在弟兄中一個最小的身上.
願主你閱立我們單純.

$^{721}$單純在你面前甘心信服的呈獻.
奉耶穌基督名字祈求.
阿們.
\newpage



\section{使徒行傳 15:36-41}
\label{sec:031839}
\textbf{路不同,道相同}~~~~駱瑞麟傳道
\newline
\newline
連結: \href{http://www.church.com.hk/acms/content.asp?site=cdc\&op=show\&type=product\&code=031839\&layout=sermon}{\texttt{031839}} ~~~~ 語音日期: 2026-03-22
\newline
\newline
\hyperref[sec:031838]{\small{< < < PREV SERMON < < <}}
~
\hyperref[sec:index_chronic]{\small{[返順時目]}}
~
\hyperref[sec:031840]{\small{> > > NEXT SERMON > > >}}
\newline
\newline
使徒行傳 15:36-41
\newline
\begin{longtable}{cl}
\hline
\hline
章節 & 經文 (和合本修訂版)\\
\hline
15:36 & \begin{tabularx}{0.7\textwidth}{X} 過了些日子，保羅對巴拿巴說：「讓我們回到從前宣揚主道的各城，看看弟兄們的情況如何。」 \end{tabularx} \\ \\ \relax
15:37 & \begin{tabularx}{0.7\textwidth}{X} 巴拿巴有意要帶稱為馬可的約翰同去； \end{tabularx} \\ \\ \relax
15:38 & \begin{tabularx}{0.7\textwidth}{X} 但保羅認為不宜帶他去，因為馬可從前在旁非利亞離開他們，不和他們一起工作。 \end{tabularx} \\ \\ \relax
15:39 & \begin{tabularx}{0.7\textwidth}{X} 於是二人起了爭執，甚至彼此分手。巴拿巴帶著馬可，坐船往塞浦路斯去； \end{tabularx} \\ \\ \relax
15:40 & \begin{tabularx}{0.7\textwidth}{X} 保羅則揀選了西拉，也出發了，蒙弟兄們把他交於主的恩典中。 \end{tabularx} \\ \\ \relax
15:41 & \begin{tabularx}{0.7\textwidth}{X} 他就走遍了敘利亞、基利家，堅固眾教會。 \end{tabularx} \\ \\
[1ex]
\hline
\hline
\end{longtable}
$^{1}$讓我們開始的時候,我們先做一個祈禱.
你每個星期都給我們有利的懷孕.
求你這個時候,你的靈就與我們同在.
親自來教導我們.
特別是,使我們能夠知道.
你將我們聚在一起.
成為弟兄姊妹.
一定是有你的心意.
雖然人的路也許是不同.
但你的道,始終是不變.
求你讓我們真的能夠按著你的心意.
能夠與你同行.
禱告奉主耶穌基督的聖名.
阿們.
大家有沒有看過這張圖畫.
大家應該都很熟悉吧.
猜一猜,想一想也知道它說什麼.
有一個小朋友,他考試作弊.
被他爸爸知道了.
當然是很生氣.
爸爸說,你錯就是要認.
打就是要站穩.
這是原則的問題.
是不能夠妥協的.
如果不是這樣,你將來.
就會一錯再錯.
怎麼可以做得成人呢.
媽媽怎麼說.
她很擔心兒子受不住這些壓力.
內心可能有罪疚感.
會收起自己.
於是陪他.
跟他聊天.
甚至跟他一起來面對.
希望在這個接納之中.
能夠使他重新起來.
其實這兩個人.
這兩位父母.
都是很愛他的小朋友.
但是他們所做的方法.

$^{41}$可能是完全不同的.
他們之間.
大家有沒有想過.
誰是對的,誰是不對的呢.
其實這種情況.
不單在家庭裡出現.
在教會裡.
也可能會發生.
當兩個很愛主的弟兄或姐妹.
他們做事很認真.
他們對同一件事情.
可能也有不同的看法.
那會怎樣呢.
大家有沒有想過.
今天我們所讀的這段經文.
其實都讓我們看到.
在聖經裡.
即使是兩位黃金搭檔.
保羅和巴拿巴.
他們兩個都是神所重用的僕人.
但是到最後.
你看到他們兩個.
可能要分開來走.
大家走不同的路.
路加他非常誠實.
他沒有維護誰是對的.
誰是不對的.
今天我們.
跟大家一起思想一個問題.
到底當我們路不同的時候.
我們還可不可以說是道相同呢.
現在就讓我們回到兩千多年前.
那個非常火爆的.
宣教會議.
來看看這兩位宣教巨人.
他們怎樣吵起來.
我會將三個部分.
分三個部分.
跟大家一起來分享.
經文第39節說到.

$^{81}$保羅和巴拿巴.
是為了需不需要.
帶上馬可.
一起來宣教.
於是兩人就發生了爭執.
甚至說到要彼此分開.
聖經這裡用了一個很強烈的希臘文.
就是爭論這個事.
意思不是摸著酒杯底來討論.
而是.
這個字表達的意思就是說.
是拍台白燈.
這樣是扯到痕的.
大家有沒有想過當時的情況.
說保羅需要去探望.
在上一次宣教的時候.
那裡的弟兄姊妹.
他心裡可能在想.
宣教是一個肅靈的征戰.
宣教團隊是需要.
很嚴肅的紀律和效率的.
所以他可能跟巴拿巴說.
巴拿巴.
上次宣教的時候.
馬可這個小子.
在別家臨陣退縮.
他對這個呼召很不認真.
如果這次我們都帶他去.
可能有很大風險.
雖然在這裡.
聖經沒有說明.
馬可為什麼臨陣退縮.
但是從經文的背景.
大家可以推敲.
想一想原來當時.
他們的旅程.
可能也有些關係.
因為跟別家出發.
他們需要走內陸.
所以要攀山越嶺.

$^{121}$是很辛苦的.
加上大家都知道.
跟耶路撒冷會議之後.
其實宣教的主導權.
去了保羅那裡.
他的主要宣教對象.
當然是那些外邦人.
對一個在耶路撒冷.
跟上來的年輕信徒來說.
不單只是他的意志和體力的關係.
可能在屬靈上.
一時是難以適應的.
所以馬可的離開.
可能對當時的宣教團隊.
帶來了很多的不便.
宣教當然我們都知道.
不是去旅行.
不是請客吃飯.
是要付上代價的.
所以保羅他對這些服侍的態度.
是有所保留.
因此可能他所針對的.
不是馬可他本人.
而是他那種不負責任的態度.
而馬可是誰呢?.
馬可曾經說過.
他是巴拿巴的表弟.
雖然他們兩個人有親戚的關係.
但是巴拿巴.
我們都知道這位勸惠子.
他所看到的.
不是那個有問題的人物.
而是那個還有機會的小朋友.
所以巴拿巴可能就跟保羅說.
他說馬可他上次臨陣退縮.
他真的不對.
如果我們這次不讓他有機會.
他這輩子可能也不能翻身.
耶穌不是要拯救失喪的人嗎?.
各位弟兄姊妹.

$^{161}$他們兩個人都不是壞人.
保羅不一定是脾氣差.
但是他有他的理由.
而聖經更加說到.
巴拿巴他是一個好人.
他被聖靈充滿.
其實他們兩個都是對的.
只不過是他們從不同的角度.
來考慮這個問題.
就好像我們經常說的.
盲人來摸象.
摸到象牙的.
就會覺得真理真是多麼尖銳.
是會刺痛人的.
摸到象耳的那個就說.
恩典是多麼柔軟.
是多麼溫暖.
因此他們可能兩個人是各持己見.
兩不相讓.
其實在教會裡面.
我們經常都會有一個爭論的問題.
到底我們是需要去宣教先.
還是做門徒訓練先呢?.
這個意見.
我在教會裡面都經歷過.
結果引起大家都很不愉快.
因為有人認為.
大使命不單單是說受洗的人數.
是需要萬民來做門徒.
所以當然要先裝備.
其他事慢一點再說.
但有些人認為不是.
所有人都做了門徒嗎?.
要等多久?.
有很多人信了耶穌.
幾十年都沒有看過聖經一次.
那其他人就不用聽福音了.
當然要先去.
然後慢慢再讀聖經.
才去裝備.

$^{201}$大家聽起來.
是否覺得兩邊都有道理呢?.
哪一邊是對的呢?.
其實在這裡看到.
最諷刺的就是.
這兩位被神所重用的僕人.
竟然忘記了自己是如何站起來的.
如果覺得馬可很不配的時候.
有沒有想過保羅在剛進主的時候.
耶路撒冷的門徒都怕了他.
都不敢跟他接近.
是誰來擔保他.
誰來介紹他給使徒認識呢?.
這就是巴拿巴.
巴拿巴是否很有問題呢?.
當然不是.
你反而覺得他是一個很有擔當的人.
願意用他的誠信.
為保羅來擔保.
為他來付上代價.
現在巴拿巴.
他用他事奉的信譽.
來擔保馬可.
那又有什麼不對呢?.
大家覺得怎樣呢?.
但回過頭來想.
所有不同的意見我們都需要去包容的呢?.
大家有沒有想過呢?.
我就覺得不是.
教會今天其實有一種危險的傾向就是.
所謂的無底線的寬容.
我們很多時候不敢說別人不對.
我們對於真理更加不敢捍衛.
因為怕一開口.
就會破壞和諧.
大家有沒有這樣想過呢?.
那真理又是什麼呢?.
真理不只是在神學上.
是否正確.
更加需要在道德上.

$^{241}$是否誠實.
其實這些都是神的屬性.
如果我們只是為了和諧.
不理是非的時候.
我們信耶穌.
我們傳福音又有什麼意思呢?.
我們所說的都是一些不真實.
都是一些虛假的信息的時候.
難道這些東西是可以改變我們的生命的嗎?.
大家都知道在基督宗教裡面.
有些所謂的耶和華見證人.
有些叫做摩門教.
東方閃電.
等等這些異端.
其實他們都是打著信耶穌的口號的.
但在教育上他們是南轅北轍.
可以說是道不同.
但是不是所有道不同的.
都有問題呢?.
可能只不過是路不同.
在十八世紀在英國.
大家都知道有兩位.
這個屬靈的巨人.
一個就是衛斯理.
一個就是懷特斐.
他們都是被神大大使用的僕人.
燃點了當時歐洲大復興的火炬.
但當他們在最火熱的時候.
他們兩個人是一些實際上的分歧.
而分開了.
幾年之後.
懷特斐去世了.
然後又有人特意去問衛斯理.
他說:你覺得在天上.
你會不會見到懷特斐呢?.
這個真是很有挑戰性的問題.
你猜衛斯理怎麼說呢?.
他說:我想我見不到他了.
那個人問:那為什麼呢?.
他說:因為懷特斐他很愛主.

$^{281}$他在天上一定是站在最前面.
靠近主的寶座.
而我站在後面.
只不過是能夠遠遠地看著他的背影.
各位兄弟姐妹.
衛斯理這句話說得很得體.
他很謙卑.
他沒有踩到對方.
更加間接地告訴我們.
在天國那裡.
雖然路不同.
但我們都會在天國裡面.
現在回過頭來.
我們看回保羅和巴拿巴.
他們在爭論什麼呢?.
他們在爭論的只不過是一些人事的問題.
其實很有討論和改善的空間.
正如我們今天這樣.
我們教會裝修.
我們要什麼風格.
要古典還是要現代的.
我們牆上有什麼.
黑色還是白色.
我們崇拜.
要什麼音樂.
彈結他還是彈琴.
不同的人都有不同的意見.
這些其實都沒有關於真理的問題.
只不過是我們人對於效率.
對於美感.
對於習慣的不同的偏好.
但如果我們不分青紅皂白的時候.
我們跟別人吵起來.
所謂的商量就不會很好.
只不過是我們弟兄姊妹之間的內耗.
對恩典的一個揮霍.
結果可能在芝麻綠豆的事上贏了.
卻輸在合一上.
這樣值不值得呢?.
所以其實在這件事上.

$^{321}$會讓我們學會了.
如果是關乎真理的問題的時候.
我們必須要像保羅一樣.
我們要寸步不讓.
正如他在《迦太太書》那裡.
一章那裡曾經說過.
他說:若傳福音給你們.
與我們所傳給你們的不同.
他就當受奏奏.
他在這兩次裡面說了兩次受奏奏.
是很嚴重的.
就是說我們不能夠對於真理方面有所退讓.
但如果說去宣教帶誰去.
去聚餐的時候去哪間餐廳吃.
吃些什麼.
我們可能要學習巴拿巴.
這只不過是路的不同.
我們是需要給人有足夠的空間.
彈性和機會.
盡量是需要與人和睦的.
換句話說.
我們是需要在飢餓的真理上合一.
而在非飢餓的事上.
我們是要有更多的自由.
但是.
當時.
你會看到.
保羅似乎是用了對付真理的氣力.
來處理一個人事上的問題.
他把馬可看成為真理的障礙.
他是忘記了.
馬可其實也是一位弟兄.
他也是需要成長的.
雖然聖經裡面沒有說到他們是誰對不對.
但是路加在這裡讓我們看到.
即使是被神大大使用的僕人.
仍然是會有互不相讓的時候.
打破了我們對.
以為屬靈人一定不會衝突的想法.
其實彼得彼書裡面.

$^{361}$保羅也提到.
與他同宮的有兩位姊妹.
都是出現了分歧的.
那位有阿爹和秦都姬.
在我們日常生活裡面.
我們都可能會碰到.
好像保羅和巴拿巴.
這種兩難的局面出現.
我們又如何去處理呢?.
大家有沒有想過這個問題?.
或者我們可以在耶穌基督的身上.
找到答案.
如果我們翻開《約翰福音》.
我們會看到.
一章十四節那裡.
約翰怎麼跟我們說.
「道成了肉身」.
是住在我們中間.
充充滿滿的有恩典和真理.
道成肉身的耶穌基督.
其實祂是完美地平衡了.
這兩個看起來好像對立.
但是祂是相輔相成的本質.
祂是我們學習的榜樣.
耶穌基督祂沒有降低任何標準.
祂沒有用恩典來掩蓋真理.
相反祂當然是替我們承受了.
真理最嚴重的審判.
因此祂對我們的代價.
為我們寫記的代價是高昂的.
大家有沒有想過.
如果沒有真理的公義.
又會怎樣呢?.
恩典就會變成廉價的縱容.
但是如果沒有了恩典.
人就會在公義的審判裡面滅亡.
但是耶穌基督在十字架上.
祂滿足了真理的要求.
以致恩典可以有堅實立足的機會.
保羅和巴拿巴分開之後.

$^{401}$馬可又會怎樣呢?.
上帝沒有忘記馬可這個小子.
沒有忘記這個在他們兩個人.
夾在中間互相爭論受傷的年輕人.
在巴拿巴那種不離不棄的陪伴下.
馬可的生命經歷了第二次的機會.
他是那位一遇到困難就肅殺的小子.
他後來在彼得的身邊服侍.
他更加成為彼得眼中的兒子.
在彼得的口裡得到大量有關耶穌基督的生平資料.
成為他寫作的一個材料.
我又想起耶穌基督時代.
那個在行淫中被捉到的女人.
當時耶穌怎樣跟她說呢?.
耶穌說:「我也不定你的罪,去吧,從此不要再犯.」.
耶穌沒有審判她.
沒有向她頂石頭.
反而向她施恩.
而恩典使她能夠重新站起來.
重新來審視她自己之所悔改.
這就是馬可經歷的恩典.
使他有的第二次機會.
但一個更加感人的事情.
是發生在幾年之後.
當保羅在羅馬避囚.
他正面對死亡威脅的時候.
他寫給提摩泰最後的遺言裡.
竟然提到馬可這位小子.
這個名字.
在提摩泰後書裡有提到.
他說:「他在全道的事上是於我有益.」.
各位兄弟姐妹.
這句說話可以說是重如泰山.
在保羅的眼裡.
以前那位無益的小子.
現在反而成為有益的人.
是他全道的同工.
大家有沒有看到.
上帝的手不期然地在兩條平行的路上.
不斷地來做工.

$^{441}$換句話說,在這段時間裡,上帝磨練了馬可的韌性.
同時也柔軟了保羅對人的看法和剛硬.
保羅後期在肥納比書裡更加寫出.
他說:「你們當以基督的心為心.」.
大家都知道肥納比書也是他監獄書信的其中一卷.
他當時開始明白到.
其實如果是講真理的時候.
如果你沒有基督的謙卑.
你很容易就會傷到別人.
這裡又使我想起電影裡有一個台詞.
是非常有意思的.
我們這些年紀大的人都會看過鐵金剛的電影裡.
裡面說:「你只能活兩次,而第二次才是你唯一活下去的方式..
You only live twice, and twice is the only way to live.」.
這句說話其實用在基督裡是更加適合.
這次的生命就好像馬可那樣.
憑血氣來事奉.
結果遇到困難的時候就打退堂鼓.
但在這個分道揚鑣的事件裡面.
馬可的生命是被上帝重塑.
以致他能夠寫出那位最早成書的馬可福音.
不單是造就了當時的信徒.
更加造就了所有的信徒直到現在.
所以馬可的生命.
他第二次的生命.
才是他真正復仕的開始.
而保羅呢.
其實保羅的生命都是這樣.
他在大馬式的路上.
他被上帝的綱照.
更加被巴拿巴.
他這位伯樂來賞識.
以致他能夠後來成為外邦人的使徒.
他第二次的生命.
更加比他第一次的.
能夠更加活出更加精彩.
最後我們都看到.
保羅和馬可是和好了.
因為他們在基督裡面.
他們越來越體會到.

$^{481}$神的本性.
有恩典和有真理.
但我想大家要留意一件事就是.
他們的僧侶.
為他們留下了.
難以忘記的那個諾印.
因為這個諾印始終都會在他們的生命裡面.
正如耶穌基督復活的時候.
曾經對多瑪所說.
祂說:你伸出你的指頭來摸我的手.
復活的耶穌.
大家有沒有看到.
祂仍然是帶著.
真理和恩典所留下來的傷痕.
就好像日本的金雞藝術一樣.
上帝沒有擦去我們的傷痕.
祂是用黃金來填充每一道裂痕.
將破碎的過去.
和充滿希望的未來.
奉合在一起,而成為一種獨特的花紋.
我們每一個信徒其實都是神所使用的器皿.
都是有著我們獨特的花紋.
這些花紋都是由真理和恩典.
所交織而成的獨特的記號.
其實這些記號.
正正在時刻地警醒我們.
罪是需要付代價的.
也是我們付不起的.
但基督的愛.
祂是為我們付出了寶貴的生命.
但願基督的這份愛.
時刻地陪伴著我們.
使我們不要小看真理的嚴肅.
也能夠倍加地珍惜那份昂貴的恩典.
這裡我更加想說的是.
耶穌基督的救贖是完全的.
但我們卻仍然活在一個尚未完全更新的世界.
所以我們人與人之間.
我們在教會裡面.
仍然會有誤解.

$^{521}$仍然會有傷痕.
仍然會有分裂.
但福音的和好的能力.
會將這些裂痕轉化成為新的意義.
這些就是和好的真相.
最後我更加需要重申.
上帝沒有製造分裂.
沒有製造衝突.
衝突也不是上帝所喜歡的.
但福音在這些衝突.
在這些分歧裡面.
就像一國那樣越走越遠.
就像保羅和巴拿巴那樣.
他們在安提阿分開了.
但保羅是帶著西拉往西走.
巴拿巴是帶著馬河向東.
福音就是這樣.
隨著他們的腳蹤越走越遠.
不斷地來抗戰.
因此我們事奉的路.
也許是有不同.
但只要我們為了基督.
我們為了福音.
其實我們所走的.
都是同一條福音的道路.
最終一定會在基督裡面.
同歸於一.
總結一下剛才我們所說的.
各位弟兄姊妹.
今天你和我多人都會面對.
一些人事上不同的張力.
我們需要堅持.
需要包容.
我們需要原則.
還是需要恩慈呢?.
但願我們都能夠效忠基督.
是滿有恩典和滿有真理.
同時.
福音也是告訴我們.
即使我們人是不完美.

$^{561}$但神仍然能夠在不完美裡面.
來成就祂的旨意.
正如保羅和巴拿巴.
當年的分開.
福音反而是越走越遠.
大家都可能想知道.
到底保羅和巴拿巴.
後來是怎樣的.
其實聖經沒有說.
但有一件事.
我想告訴大家.
無論是怎樣.
有一天在主一角度裡面.
我們都必會與曾經分開的弟兄姊妹.
再次相逢.
就好像維斯理和懷特肥那樣.
因為在啟示錄裡面.
給我們一個終極的盼望.
就是神要擦去我們一切的眼淚.
到那時候我們必定會明白.
在神的角度裡面.
現在有些人好像保羅那樣.
堅持真理不放.
也有些人像巴拿巴那樣.
很注重這個恩典.
就好像我們的父親和母親那樣.
都會用不同的方式.
來愛祂的小朋友.
但是沒有真理和恩典的神.
祂總是用這兩種.
看上去好像不同.
但是一樣的愛.
來給我們保護和支持.
但無論是怎樣.
但願我們每個人.
在這條個人的路上.
都能夠活出基督的陽色.
直到我們再見主的面.
然後我們一同來禱告.
持必人我的天父.

$^{601}$因為你的話語是充滿著亮光.
主讓我們承認.
我們很多時候.
我們不是太過堅持我們自己.
我們就是太過遷就其他人.
求主你赦免我們的偏執.
主讓你是滿有恩典.
也是滿有真理的主.
求你幫助我們.
在應該堅持的地方我們去堅持.
在應該包容的地方我們去包容.
但最重要的就是.
我們能夠分別這兩者的差別.
主讓我們更加為那些.
因為傷害而產生的疤痕.
而向你呼求.
我們承認傷害是真的.
疤痕也是痛的.
但求你轉化這些疤痕.
成為我們服侍的記號.
便是時常來警惕我們.
罪的代價是重的.
但你的恩典也是夠我們用.
因為你一直帶領我們.
在這條恩典的路上.
禱告奉主耶穌基督的聖名.
阿們.
\newpage



\section{以賽亞書 60:1-18}
\label{sec:031840}
\textbf{先嚐天國的味道}~~~~陳穎康牧師
\newline
\newline
連結: \href{http://www.church.com.hk/acms/content.asp?site=cdc\&op=show\&type=product\&code=031840\&layout=sermon}{\texttt{031840}} ~~~~ 語音日期: 2026-03-22
\newline
\newline
\hyperref[sec:031839]{\small{< < < PREV SERMON < < <}}
~
\hyperref[sec:index_chronic]{\small{[返順時目]}}
~
\hyperref[sec:031841]{\small{> > > NEXT SERMON > > >}}
\newline
\newline
以賽亞書 60:1-18
\newline
\begin{longtable}{cl}
\hline
\hline
章節 & 經文 (和合本修訂版)\\
\hline
60:1 & \begin{tabularx}{0.7\textwidth}{X} 興起，發光！因為你的光已來到！耶和華的榮光發出照耀著你。 \end{tabularx} \\ \\ \relax
60:2 & \begin{tabularx}{0.7\textwidth}{X} 看哪，黑暗籠罩大地，幽暗遮蓋萬民，耶和華卻要升起照耀你，他的榮光要顯在你身上。 \end{tabularx} \\ \\ \relax
60:3 & \begin{tabularx}{0.7\textwidth}{X} 列國要來就你的光，列王要來就你發出的光輝。 \end{tabularx} \\ \\ \relax
60:4 & \begin{tabularx}{0.7\textwidth}{X} 你舉目向四圍觀看，眾人都聚集到你這裡。你的兒子從遠方來，你的女兒也被抱著帶來。 \end{tabularx} \\ \\ \relax
60:5 & \begin{tabularx}{0.7\textwidth}{X} 那時，你看見就有光榮，你的心興奮歡暢；因為大海那邊的財富必歸你，列國的財寶也來歸你。 \end{tabularx} \\ \\ \relax
60:6 & \begin{tabularx}{0.7\textwidth}{X} 成群的駱駝，並米甸和以法的獨峰駝遮滿你；示巴的眾人都必來到，要奉上黃金和乳香，又要傳揚讚美耶和華的話。 \end{tabularx} \\ \\ \relax
60:7 & \begin{tabularx}{0.7\textwidth}{X} 基達的羊群都聚集到你這裡，尼拜約的公羊供你使用，獻在我壇上蒙悅納；我必榮耀我那榮耀的殿。 \end{tabularx} \\ \\ \relax
60:8 & \begin{tabularx}{0.7\textwidth}{X} 那些飛來如雲、又像鴿子飛向窗戶的是誰呢？ \end{tabularx} \\ \\ \relax
60:9 & \begin{tabularx}{0.7\textwidth}{X} 眾海島必等候我，他施的船隻領先，將你的兒女，連同他們的金銀從遠方帶來，這都因以色列的聖者、耶和華－你神的名，因為他已經榮耀了你。 \end{tabularx} \\ \\ \relax
60:10 & \begin{tabularx}{0.7\textwidth}{X} 外邦人要建造你的城牆，他們的君王必服事你。我曾發怒擊打你，如今卻施恩憐憫你。 \end{tabularx} \\ \\ \relax
60:11 & \begin{tabularx}{0.7\textwidth}{X} 你的城門必時常開放，晝夜不關，使人將列國的財物帶來歸你，他們的君王也被牽引而來。 \end{tabularx} \\ \\ \relax
60:12 & \begin{tabularx}{0.7\textwidth}{X} 不事奉你的那邦、那國要滅亡，那些國家必全然荒廢。 \end{tabularx} \\ \\ \relax
60:13 & \begin{tabularx}{0.7\textwidth}{X} 黎巴嫩的榮耀，就是松樹、杉樹、黃楊樹，都必一同歸你，用以裝飾我聖所坐落之處；我也要使我腳所踏之地得榮耀。 \end{tabularx} \\ \\ \relax
60:14 & \begin{tabularx}{0.7\textwidth}{X} 壓制你的，他的子孫必來向你屈身；藐視你的，都要在你腳前下拜。人要稱你為「耶和華的城」，為「以色列聖者的錫安」。 \end{tabularx} \\ \\ \relax
60:15 & \begin{tabularx}{0.7\textwidth}{X} 你雖曾被拋棄，被恨惡，甚至無人經過，我卻使你有永遠的榮華，成為世世代代的喜樂。 \end{tabularx} \\ \\ \relax
60:16 & \begin{tabularx}{0.7\textwidth}{X} 你要吃列國的奶，吃列王的乳。你就知道我－耶和華是你的救主，是你的救贖主，是雅各的大能者。 \end{tabularx} \\ \\ \relax
60:17 & \begin{tabularx}{0.7\textwidth}{X} 我要賞賜金子代替銅，賞賜銀子代替鐵，銅代替木頭，鐵代替石頭。我要以和平為你的官長，以公義為你的監督。 \end{tabularx} \\ \\ \relax
60:18 & \begin{tabularx}{0.7\textwidth}{X} 你的地不再聽聞殘暴的事，境內不再聽見破壞與毀滅。你必稱你的牆為「拯救」，稱你的門為「讚美」。 \end{tabularx} \\ \\
[1ex]
\hline
\hline
\end{longtable}
$^{1}$各位早晨 大家好.
多謝阿健牧師的介紹.
我很記得一首歌的歌詞叫「我仍是我」.
很開心能夠服侍大家.
正如阿健所說.
上帝讓我去完成了一個事工.
現在是擔任在另一個事工的裡面.
可能大部分人都認識我.
當我去探訪其他教會 我都會去介紹自己.
可能都認識我 但有部分人未必認識我.
我有一個太太和兩個小朋友.
我的兒子今年讀一年級的大學.
跟譚傳道的女兒一樣.
我女兒是一年級10.
除了我家裡有一家四口之外.
上年多了四隻貓.
阿健有五隻貓.
一隻是街貓拍我們的門.
我們放牠進來之後就大件事了.
原來牠懷孕 出了四隻.
所以有五隻貓.
本來想展示相片給大家看.
但我忘了放在powerpoint上.
不過有機會有興趣可以來找我給相片看.
今天和大家分享的題目是.
「欣賞天國的味道」.
究竟天國有什麼味道呢.
我們用廣東話來解釋.
天國有什麼味道呢.
譬如你去到一個地方 哇 像一層薯條.
天國有什麼味道呢.
我們又怎樣去品嘗天國的味道呢.
我們一起去學習一件事.
我們有時候在我們的生活裡面.
都有一些情況是我們去試味.
什麼時候我們可以.
大家有沒有試過試味.
大家好像很陌生.
試味 什麼情況我們有機會去試味.
當然是那些超市 超級市場.

$^{41}$有些攤檔可以去試味.
很多人都很喜歡試吃的.
見到攤檔就好像買溜蜜糖.
為什麼那麼受歡迎呢.
因為不用錢的嘛.
不吃也看看吧.
很多時候那些攤檔就說.
買吧 買吧 不吃也看看吧.
試吃的物品通常都是比較新的.
新的意思是可能在市場還沒推出.
他就很想多些人去試.
那就試吃吧.
所以都是很受歡迎 不用錢.
也是一個很特別的物品.
不過我人生第一次的真正試吃.
就是我結婚拜酒.
可能有部分人也有這樣的經驗.
因為你要訂酒席.
先試一下好不好吃.
因為那次是二十多年前.
是我第一次結婚.
也是最後一次.
所以沒什麼經驗.
所以試吃是不是去酒樓試菜是免費的呢.
是不是免費的.
有沒有準備結婚的那些.
不是免費的.
原來要給錢的.
但是也OK 起碼試了.
安心吃了什麼也OK.
但記得一次.
我是要去買一樣東西的.
我是想去買芝士.
所以我就準備去一個W的超市.
去買那樣東西.
但是你知不知道其實男士買和女士買是有不同的.
In general speaking.
男士想買那樣東西就是買那樣東西.
所以想買芝士.
我是習慣去W的超市.

$^{81}$我知道在什麼位置.
所以我心裡已經計劃好去那個位置.
但是女士買東西就比較豐富.
想買A就加ABCDE很多樣東西回來.
是fantastic的.
我是男士.
我打算買芝士就買芝士.
誰不知當我去到那個W字頭的超市.
門口就有一個試吃攤檔.
你覺得他在試什麼呢.
芝士.
我看到.
平常我也很少去那些攤檔.
但是芝士.
我又買芝士.
那我就去試一下.
吃一吃就覺得好吃.
那我做什麼.
我就馬上給錢買.
我連去那條路也省下來了.
買完就馬上回家去吃芝士.
所以其實你發現試這件事.
對推銷一個產品有沒有效.
有效的.
特別在COVID之後多了人在網上購物.
都很多人喜歡在網上購物.
除了便宜之外.
有時候要吸引人去買.
就有個30day guarantee.
就是你買了之後30天你不滿意.
可以完全退貨.
有沒有人因為這件事而中伏.
我真的因為這件事很吸引我買了這件事.
但是就沒有中伏.
OK的這件事.
所以推銷產品去嘗試.
這個策略是很重要的事情.
讓顧客先去品嘗到那件事是什麼.
然後才去買那件事回來.
生活上面我們有試吃有試味.

$^{121}$那在淑齡上面有沒有得試的呢.
有沒有的.
有沒有 試一下吧 無懷有沒有得試的呢.
有句經文大家很熟悉的.
你們要嘗嘗主因的滋味.
便知道他是美善.
投靠他的人有福了.
就是說 你試一下吧.
你不試又怎知道神的福是這麼甘甜的呢.
而當你試的時候.
投靠他的人是有福氣的一件事.
那時候我們和一些朋友去傳福音的時候.
可能在座也有很多經驗.
我們會講一大篇關於福音的東西.
耶穌為什麼要走上十字架.
他為什麼要死.
死了之後又做什麼.
我講了很多的.
大概講了三四個字.
講完之後會怎樣.
願主祝福你.
是不是這樣放他走呢.
不會放他走的.
找著機會去問什麼.
你會不會嘗試去接受這個福音.
最後決戰是很重要的一件事.
譬如當中有些人喜歡釣魚.
釣魚回來.
又怎樣.
Hello 就放他回去.
不是的.
你會抓住那條魚.
回去怎麼煮好呢.
昨天傳福音你都很想去決戰的.
但是在講的過程中.
他有笑.
好像在聊天.
但最後你問他願不願意的時候.
他就開始有些猶豫了.
我要想想家裡.

$^{161}$我要想那件事那件事那件事.
他開始懷疑了.
那一刻我們用一個方法叫三福.
可能當中也認識三福.
我們就舉個例子.
假設前面有個很漂亮的蘋果.
又紅又漂亮又乾紮紮.
不過是蠟而已.
很漂亮很乾紮紮.
接著我就問.
你怎樣才知道蘋果是甜的呢.
怎樣才知道.
咬一口它.
你不咬一口.
你又怎知道蘋果是甜的呢.
我嘗試去挑戰就是.
你試一下.
你試了味道好.
那你會願意繼續去追尋.
而我們加句話.
沒有吃虧的.
你吃了之後.
你不喜歡就可以不要了.
《屬靈生命》上面是有試吃這件事的.
今天我想透過這個機會去思考這件事情.
我看經文是在《二賽亞書》60章的經文.
二賽亞先知是出生在主前七百年的時間.
主前七百年是什麼狀況呢.
當時北角的以色列已經被亞述去擄走了.
因為他們的罪.
所以他們滅亡了.
而剩下的就是猶大南角這個地方.
所以二賽亞當時的信息.
是包括了警告和刑罰的信息.
你的兄弟北角已經因為罪的緣故.
被亞述擄走了.
那你們剩下在南角的猶大.
你們看到這件事就知道.
原來神的審判是真實的.
不是說說那麼簡單.

$^{201}$神除了愛.
也是一個很嚴厲的上帝.
所以在二賽亞書裡面.
其中一個很重要的信息是警告和刑罰.
但是這只是二賽亞書一半的信息.
裡面也說到上帝的預言.
預言神的拯救和指明將來的事情.
這就是整個二賽亞書.
分為兩個很不同的情況.
而今天我們看的二賽亞書60章.
是描述神指明將來的警報.
所謂的警報是上帝借助二賽亞先知.
向一群對象講述將來他們身上會發生的事情.
雖然他們會被擄到巴比利.
但是先知說他們有一天會再次回到裡面.
這群對象就是神的指明.
不過我知道這個二賽亞先知.
他的對象不只是單單是當時的猶太人.
也包括所有的信徒.
就是你和我.
所以今天節目主持人所讀的經文.
不只是讀給當時的猶太人知道.
同樣很想讀給你和我知道.
而當中的預言也不是單單說.
他們將來可以從被擄的地方.
回到猶太人身上那麼簡單.
他所說的是新約時代.
甚至是將來新天新地的警報.
所以整個二賽亞先知書裡面所說的.
是整個上帝的角度.
整個上帝角度的預言.
而這個預言.
即使今天在2026年都未完全實現.
但是即使如此.
今天我們都能夠品嘗到部分天國的滋味.
剛才我留心聽趙牧師的祈禱.
有一句說話很硬的.
就是天國今天已經臨在.
同意嗎?.
天國未完全.

$^{241}$但是它已經臨在我們裡面.
而我們可以品嘗.
可以taste得到天國的滋味.
在哪裡可以taste到天國的滋味?.
即使在W市場品嘗得到.
那麼天國的滋味在哪裡可以品嘗得到呢?.
剛剛搶到差一秒 差半秒.
在教會 在團契 在信徒能夠聚集在一個地方.
就能夠品嘗到天國的滋味.
為什麼?.
因為我們所信的人.
我們都是天國的子民.
而聖經更加保證.
當兩三個屬祂的人聚在一起的時候.
上帝就在我們裡面.
同意嗎?.
所以我強調今天我們聚在一起敬拜.
我們不只是一群人唱歌讀經那麼簡單.
我們一定要記得.
就是上帝在我們當中.
今天我們所要敬拜的.
是那個創造天地萬物的上帝.
所以在教會這個場所.
在小組,團契.
在我們品嘗上帝滋味的一個很重要的地方.
有一次我去一間國際教會去敬拜.
是一間國際教會.
其實我進去那個教會.
我已經感覺到那股氣味了.
國際教會那股氣味 為什麼?.
因為去到的時候我看到有很多不同膚色的人一起敬拜.
有些皮膚比較深色.
不是因為他曬太多.
他根本是很深色的皮膚.
有些頭髮是不同顏色的 不是染髮.
所以你看到裡面是有很多不同國籍的人一起敬拜神.
完了之後我就好奇.
我就問那個主任牧師.
究竟你那間教會有多少不同國籍的人一起敬拜?.
你猜猜有多少不同國籍.

$^{281}$這個都頗有規模的國際教會.
裡面是有45個不同的民族,國籍的人一起敬拜神.
哇~.
然後我想起剛剛的時間 我好像在天家敬拜神一樣.
雖然那時候我還沒死.
但是好像在地如藥如在天一樣的感覺.
我當時是和44個不同的民族,國籍的人一起去敬拜我們的神.
像不像將來我們在運行天國敬拜神一樣.
當你環繞看著有很多不同的膚色的人一起敬拜上帝.
所以在那個國際教會裡面.
我已經品嘗得到將來敬拜神的味道.
在今天的經文裡面.
我們除了去看一下.
以蔡先知去描述天國是什麼樣子之外.
等著每樣我們都去同樣思想.
既然教會是一個可以品嘗天國之味的地方的話.
究竟我們教會有沒有這一股味道呢?.
當人來到教會的時候.
這個地方有沒有那股味道呢?.
這個味道好像外面沒有的.
只有這裡才有.
這個天國味道.
經文裡面看到起碼有三個以蔡亞去講述天國特別的地方.
第一 天國是一個吸引萬人前來的地方.
即是說無論是什麼人都好.
你來到都是歡迎的.
第二是一個呈獻的地方.
就是人是很願意來到.
無分彼此的願意為上帝拜上.
是怯的.
第三個是和平得勝的地方.
這個地方是彼此建立的一個地方.
我們一起去分享一下這三個味道.
讓我們去想一下.
我們教會有沒有這個味道呢?.
吸引萬人前來的地方.
經文開始提到.
這個描述道有很多人被吸引到神的面前.
吸引到神的角度當中.
第三節 萬國要來就你的肝.

$^{321}$第四節 眾人都聚集在你這裡.
我們帶出一點點的想像力.
你會發覺這是一個很壯觀的景象.
很多人從世界各地.
去到被吸引到神的面前去敬拜這個神.
不只是隔壁的宿舍.
或者住在隔壁的鄰居.
是很多人從世界各地去到神的面前.
是一個很壯觀的狀況.
其實曾經有和一些牧者去英國探訪一些教會.
當地的童工安排我們去住一間酒店.
後來發現這間酒店.
是由回教徒去營運的.
那時候我覺得有點奇怪.
我們是一群基督徒的牧者.
我們沒有隱藏身份.
我們就問 為什麼當地的童工.
會安排我們去一間回教的酒店去住呢?.
其實當時我一個迷.
我又不敢問 因為他們的安排好了.
你不要那麼多事幹 你安排我們去住吧.
後來有些經歷 我覺得是神的心意.
為什麼知道那間酒店是回教徒所營運呢?.
你去到他的接待處.
有個年輕人站在前面.
你一見到他 你就知道他是回教徒.
他的衣服 你一見到就是了.
但除此之外 在他的酒店大堂.
就有這幅畫的.
而那個大小比這閻公墓大一倍.
你可以想像到.
在他的酒店大堂有一幅這樣的畫.
那就是說 除非你發雞盲.
你一進去 你一定見到這幅畫的.
你一定知道這是回教的背景的酒店.
我們在當中大概住了3天左右.
有一天 我就把握機會.
和當中那個年輕的侍應去聊天.
我和他聊天之後 我就說.
這個是什麼地方 大家知道是什麼地方嗎?.

$^{361}$我給他提示叫做墨家.
是回教徒每年10月都會從世界不同的地方聚集.
去敬拜他們的神.
這個是那個圖畫.
聊天的時候 這個年輕人就說.
這個是他人生心中一定要去的地方.
是他夢寐以求一定要去到的.
不過當時他很年輕 他說他不夠錢.
但是他不夠錢不要緊 他會努力工作.
他說他的心願就是一定要去到一次.
當我聽到這個年輕人所說的時候.
其實我心裡面都有一些震撼.
這個年輕人 他為了他的神.
有這麼大的一個心願.
去到的地方 就算努力工作都是去的地方.
我心裡想 究竟同樣我們基督徒有多渴望.
去到朝見我們的神呢?.
當然我們知道我們的上帝是無所不在的.
所以我們不是一定要去聖城去敬拜 是吧?.
請問當中有誰去過聖城呢?.
哎 譚傳道有 我還沒有.
而在我接下來的旅行計劃.
都不是在我的Top 5 list裡面.
當然我們知道不是聖城才有上帝的 是吧?.
上帝都在當中.
但是我又沒有那個年輕人這麼渴望.
就是我要去聖城去看回上帝.
耶穌所做的事情.
今天 當然我們不一定要去聖城.
但是同樣我們知道上帝是無所不在的時候.
我們每天有沒有很渴望去到上帝面前.
起碼安靜五分鐘.
那個年輕人的回應對我一個很大的震撼.
我們有多渴望去到上帝面前.
聖經說什麼吸引人去到上帝面前呢?.
聖經說是他的肝.
萬國是因為被一個肝去吸引去到上帝面前.
而這個肝是來自我們上帝耶和華的榮耀.
而當耶和華榮耀.
聖經說他照在他的子民身上.

$^{401}$而又再去反射的時候.
人見到就被吸引.
明不明白邏輯?.
上帝肝照在他的子民身上.
然後又反射出來的時候.
人就見到就被吸引.
所以真正吸引萬國去聚集的.
是耶和華的榮耀.
包括什麼?.
包括神的愛.
祂的公義.
祂的救恩 祂的慈愛.
這一切就是真正要吸引人去到的地方.
藉此我們要問今天.
我們要什麼去吸引人去到教會呢?.
有時候我聽到我父母或我祖父母說.
其實他們小時候在香港.
他們也會去教會.
主要是去天主教.
為什麼他們會被吸引去到教會呢?.
有沒有聽你父母或祖父母說過?.
為什麼會回去呢?.
想認識上帝?.
是什麼?.
因為要去派.
要去派什麼?.
派米 派麵粉 派麵包.
要去派就去吧.
是不是很好?.
我敬佩當時的牧者和天主教的神父.
他們很down to earth.
他知道當時的人有什麼需要的時候.
他就用那些東西去吸引人去到他們的教會.
但是那些我知道.
麵粉 麵包那些只是一層表面的地方.
他來到就不會放過他.
就會和他分享福音的地方.
真正讓人留下吸引的.
不是表面的一些誘惑或利益的地方.
而真正吸引的是神的光.

$^{441}$祂的憐憫 祂的福音的地方.
不過當中聖經剛才說到一個很重要的媒介.
就是反射而榮耀的信徒.
直至我們要明白的是.
原來今天我們信徒的生命.
我們的言行.
能不能夠去反射耶和華的榮耀在人的身上.
如果我們在一個有很乾淨的一塊鏡子.
或者一個反射體的時候.
原來整個榮耀就能夠在人的眼中得到光輝.
那麼我們什麼時候可以發光呢?.
聖經說摩西塊面削發光.
我們什麼時候會發射到神的特質呢?.
很有愛 很有公義.
會不會有時候只是在崇拜這段時間.
我們比較熟齡一點呢?.
當我們離開了教會之後.
我們回到我們工作的環境.
我們的家.
甚至是我們的生活圈子裡面的時候.
我們是否仍然可以幫助上帝去反照他的榮耀呢?.
聖經說人是被吸引到面前.
當他的榮耀投放在屬於他的人身上.
然後反射出來的時候.
人就會被吸引.
所以值得我們不斷去思考.
原來今天我們聚在一起.
我們去聽道.
我們去明白.
然後當我們散去.
離開這個會場.
原來上帝說.
你繼續為我去發光.
透過我們的生命.
人或許會更加想回到上帝面前.
什麼人去到上帝面前呢?.
朱牧師去讀的時候.
我們發現有很多不同的名字.
地方名.
譬如第六節.

$^{481}$成群.
米顛.
或者是士巴.
第七節就是基達.
尼拜約.
這些地方.
可能你們沒有聽過這些地方.
比較熟悉的就是羅撒蘭.
伯利恆.
這些地方是在猶太以外的一些地方.
而這些地方都是由外邦人所居住的.
第四節就是.
你的種子要從遠方而來.
第十節就是外邦人必建築你的城牆.
這是一個很重要的預兆.
就是去到神國的人.
不只是猶太人那麼簡單.
也有很多的外邦人.
都去到上帝面前.
而這個信息.
對當時的猶太人.
即是主前七百年.
他們聽到的時候.
其實如果我們待在他們的環境.
他們是不能夠去接受這個預言.
因為我們知道猶太人.
在過去他們是接受不了外邦人.
外邦人是不潔淨的.
他們的上帝的話.
是屬於自己的.
外邦人是沒有份的.
所以這個預言.
在主前說的時候.
他們這群人一定是很難接受.
為什麼會是這樣的呢.
但我們明白到去到新約.
保羅開始向外邦人傳福音.
甚至去到今天的世代.
這個外邦人得著福音的事實.
已經是完全的發生.

$^{521}$福音已臨到在不同的民族.
膚色和人的語言的身上.
我們知道聖經傳書.
即是整本聖經.
或者是部份特別是新約的聖經.
今天已經用超過二千種.
不同的語言去翻譯出來.
甚至在中文也有很多不同的頁本.
在英文也有很多不同的頁本.
換句話說.
在今天裡面有超過二千多個民族.
都已經信了我們的主耶穌基督.
這件事在以前的時代是不可能發生的.
但今天已經成為一個鐵一般的事實.
但是福音仍然未達到萬民.
在二月份的時候.
我們去參與.
即是趙牧師和Nelson都去參與在聯會裡面.
當中有一對宣教士分享他們宣教的工作.
有一句話我都很震撼.
是今年的二月.
他說在他講的時候.
他說亞洲的地區.
亞洲我們很熟悉.
甚至有時會去旅行泰國,馬來西亞的地方.
他說單單在亞洲的地區.
今天仍有超過一百個民族未聽到福音.
哇 我覺得有這麼多嗎?.
亞洲應該很廣泛的.
他說今天在亞洲仍有一百個民族未聽到福音.
即是說福音未完全傳遍在聖餓的世界裡面.
那什麼時候才實現呢?.
耶穌再次回來.
聖餓說這天國的福音要傳遍天下.
對萬民作見證.
然後末期就能夠來到.
即是說當耶穌再次回來的時候.
原來福音就真是可以傳遍這個世界.
不過我們明白到就是.
今天天國仍然是歡迎所有不同的人來到的.

$^{561}$裡面可能有不同的文化,不同的地位.
但大家都一起的見擺著同一個上帝.
這個是你很明白的一個地方.
而值得去反思的就是.
今天既然教會是一個天國接觸的地方.
而天國又歡迎萬國的人的時候.
我們問 今天教會是否能夠接納不同背景的人.
一起參與在一起見擺.
我上年年尾就完成了我的服事在南宣.
在最後的時候我們也有一個聖誕的聯合聚會.
聯合我們中文和英文一起的見擺神.
當中有年輕人 也有我們的弟兄姊妹.
我那天是講道.
講大家都OK的 氣氛都不錯.
因為是聖誕節嘛.
但是講著講著我發覺後面的.
比如說Uncle Fung那一群人.
開始就 不是他.
後面那些人就開始很緊張了.
那些樣子有點僵硬.
我心想是不是我說錯話了.
他們開始有一個很輕鬆的環境.
就開始有些不太動了.
後來我講完道 下台就回去看看發生什麼事.
發現後面說今天有個新朋友來了.
不過這個新朋友就和我們不同的.
後來這個人是來自巴勒斯坦.
一個年輕人住在我們教會附近.
他看到我們有聯合見擺就來到我們當中.
那當中為什麼這麼緊張呢?.
就是大家見到他就說.
你要看著他.
你看他有沒有帶背囊來.
背囊裡面重不重.
有沒有放低角度呢?.
就走了.
可能大家那時候聽了很多新聞.
就是中東或者一些人去教會破壞.
哇 大家很緊張啊.
大家都盯著這個人.

$^{601}$你幫我一起看著.
希望不會發生什麼事.
這個圖畫給我看到就是.
原來今天我們教會.
我們那時候的教會.
原來對不同背景的人都有一定的看法.
我們不要說得那麼複雜.
今天在一個教會澳洲.
一個不同文化兩文三語的教會的時候.
你發覺在同一個屋簷下去發展粵語,國語.
甚至是英語的時工的時候.
其實都是不容易的.
同意嗎?.
大家齊人在一個地方見拜神.
能夠一起共融一個地方.
要知道教會是讓所有人進入的地方.
所以求神去幫助我們.
我們能夠去學習.
去接納不同背景文化的人一起去見拜上帝.
因為這正正是教會在地上展現天國的一個地方.
以至無論什麼人來到都好.
他一進到這個教會的時候.
他所見到的不是大家的眼光.
而是大家說:Come! Embrace!.
我們歡迎你來到見拜神.
因為這個地方是屬於不同的地方.
第二個天國是一個呈獻的地方.
剛才提到經文裡面有很多不同的民族來到.
被上帝的光吸引.
而且一齊去見拜神.
然後去到第六和第七節.
這兩節經文帶出一幅圖畫.
就是不同的人除了在遠方來到之外.
他們都帶了他們的禮物.
預備去獻給他們的上帝.
當中帶了什麼去獻給神呢?.
有黃金,乳香,有帶羊,有帶公牛.
去到祭壇上獻給神.
而經文更加強調一個很寶貴的地方.
耶和華去閱納他們所獻上的.

$^{641}$即是上帝見到人作一個預備去到祂面前.
神是很開心一件的事情.
而這個景象是帶出一個.
在神國裡面一個很重要的狀況.
就是一個何謂真正敬拜的態度.
即是說人來到的時候是作好準備.
他很準備去見拜神.
並且是很願意去為神擺上一些事情.
這個圖畫當神見到祂會很開心.
人不僅千里迢迢來到.
更加預備好祭物獻給他的上帝.
曾經有機會去參與中國一些短訓.
我們去的地方是雲南.
Uncle Fong也有和我們去參與這個寺宮.
雲南當時有一個超通的地方.
是有一個聖經的學校.
裡面有一些學生.
這些學生是來自雲南周圍很多的農村.
其實當我去到這個聖經學校.
最初覺得只是一間學校而已.
但去到以後都是眼睛睜開.
我發覺裡面的學生是很年幼.
剛才我說我的女兒是10歲.
當時聖經學校的學生比我女兒還小.
他們來這個聖經學校大概讀一年的課程.
然後就回到他們的農村.
回到農村做什麼呢?.
他們就做傳道人.
嘩!當時我也是傳道人.
但我比他們大一倍.
這麼年輕就回去農村做傳道人?.
因為他們當時實在是很缺乏.
他們地方教會沒有傳道人.
所以就送一些很年幼的學生去到這個聖經學校去學習.
學完一年就回去做傳道人.
嘩!很難得.
當時我探訪的時候.
也和一些頂尖的妹妹去關心當中的學生.
記得一個女生.
我和她聊天.

$^{681}$我就問她:你是從農村來的嗎?.
她說:是的.
你農村有沒有教會?.
她說:有.
近不近你的家?.
她說:很近的.
她說:大概我走兩至三個小時就到了.
嘩!走兩至三個小時就去到那間教會?.
那去一次見法不就五個小時?.
那算近的嗎?.
我也很震撼.
然後我追加她說:那你喜歡回教會嗎?.
她說:我很喜歡回教會的.
我每個星期提早起床.
早點走路.
以致我不會遲到的.
那我心想.
你發現有時候你和別人聊天.
你以為是想給予別人東西.
你想關心怎樣怎樣.
但原來你和別人聊天.
有時候是他給你東西.
我聽到的時候.
都是很大的提升.
沒有到慚愧.
這些才是真的去敬拜神.
去教會的那個心意.
那有誰走回教會的?.
有沒有?沒有.
有誰開車超過一小時回教會的?.
都沒有.
有時候我們.
這個提醒也不單是對很多人.
有時候我們對敬拜的心未必是這麼熱切.
甚至有時候我們有些事情都會難了.
我們回到教會去敬拜神.
我們不需要走兩三個小時.
我們開車不需要超過一小時.
但有時候有些事情都會難了我們.
有時候發覺今天天氣不似預期.

$^{721}$今天很大雨.
有時候我聽到一些節目說.
明天我的小孩要去派對.
所以我就去不了教會了.
甚至有時候有些真實的.
我弟兄來教會是負責一些音響的事奉.
第一次他跟我說.
牧師不好意思 我要請一年假.
因為我弟兄是做工程的.
就是幫人建屋.
他預視到在新年會很忙.
很多工程.
他就說我都可以回來教.
不過我不可以肯定我每個星期都回來.
所以他跟我請假.
一年我就做不了這個音響了.
所以值得我們去思考.
我們有多渴望去敬拜我們上帝呢?.
我們敬拜的時候有沒有那份準備的心.
在《舊聖經》裡面有很多地方.
我們看到上帝要求他子民去獻情的要求.
而神很清楚給當時的以色列人知道.
當他們去到聖殿朝見神的時候.
他們應該怎樣去預備他們的祭物.
當然今天我們不是舊約的時代.
我們不需要你們帶羊 你們帶牛 你們帶豬回來.
我們不需要.
也不需要趙牧師和Nelson在門口幫你鋸一些東西.
獻給神.
你們不需要做這個行動.
但是聖經有一個訊息就是 敬拜是一個獻情的行動.
以色列人來到上帝面前是要做好準備.
準備去付出 不是單單去收取.
這個同樣是否在我們教會也有一個味道.
我們很感受一個濃烈的氣氛呢?.
有時候我們發現很多年輕人的聚會.
都是安排在下午的 知道為什麼嗎?.
太早沒有人在聚會裡面.
有時候我們看到很多年輕士工的導師.
都會叫他們的年輕人你們早點睡.

$^{761}$因為明天是敬拜.
因為很搞笑的就是那些年輕人會說.
牧師啊 我們很早睡的.
你們十一點鐘睡 我們兩三點睡的.
我們比你們早.
所以親愛的弟兄姊妹.
今天有沒有想過你帶什麼來獻給你的上帝?.
你發現在程序表上面有很多的程序的環節.
都是獻情的.
唱歌是獻情 祈禱是獻情.
奉獻是一個感恩的獻情.
我們只是給十分之一我們的上帝.
是一個感恩獻情的地方.
但是更加希望我們所獻的是什麼呢?.
我們明白不是單單停留在我們嘴裡唱的詩歌.
也不是單單集中在我們所獻的金錢.
而是神很希望獻的是我們全個人的生命.
不如大家一起讀吧.
我讀了這麼多東西 大家可以讀一下.
我們一起讀 羅馬書12章1至2節.
所以弟兄們 我以神的慈悲勸你們.
將身體聖上當作活祭.
是聖潔的 是神所欺欲的.
你們如此事奉乃是理所當然的.
將我們全人當作活祭.
獻給上帝是保羅A.K.文的訊息.
原來我們本身自己就是獻給神的那份祭物.
讓我們全人積極參與在敬拜的裡面.
所以弟兄妹們 我們一起去學習.
今天我們有時候都會去思考崇拜更生.
程序裡面有什麼詩歌.
以前我都會和我們教會的敬拜部部長去思考.
什麼令到這個教會的敬拜更加活潑呢?.
於是人進來的時候都會迎接到他們的感受.
但我知道最重要更生的地方.
不是詩歌 唱什麼詩歌 或者整個流程.
而最重要更生的地方就是信徒裡面敬拜的那份心意.
在你進入這個會堂之前 你已經知道你在做什麼.
你的對象是什麼的時候 其實你唱什麼詩歌的流程如何.
其實不是很影響我們敬拜神的心.

$^{801}$真正更生的是我們在這裡.
我們回來是要預備我們的心.
所以今天我們為主能夠呈現什麼 弟兄姊妹.
我們怎樣去預備我們自己的生命.
我們可以去敬拜我們的上帝呢?.
最後 天國是一個和平得勝的地方.
今日經文裡面除了看到人去敬拜的時候.
也看到一個很更生的景象.
第十四節 素來輔待你的.
他的子孫都必屈身來就你.
藐視你的 都要在你腳前跪拜.
第十五節 你雖然被撇棄,被厭惡.
甚至無人經過.
我卻使你變成變為永遠的榮耀.
成為累代的喜樂.
這個經文裡面 它提出一個更生的景況.
以前是受苦的 可說是作奴隸的.
現在在神的國裡面 你可以安息了.
之前是被藐視,被瞧不起的.
因為神的國 你就得了榮耀.
就是一個這樣的圖畫.
當然我們同意A.K.文某程度上.
去描述當時猶太人.
回到耶路撒冷的那個事實.
以前他們在巴比倫或者波斯.
他們本來是被擄一幫人.
被抓到一些囚牢.
但是他們一天可以回到自己的地方.
是一個生命很重要的自由釋放的狀況.
而這段經文就是解釋他們生命的重整的地方.
但是我們更加明白.
二次世界大戰的描述,這個預言.
不只是說肉身的自由.
或者不是說他們以前沒有工作.
現在變成有工作.
身體有軟弱 現在得到醫治.
不只是表面的改變.
而是說屬靈,生命,德性的改變.
哪裡可以看到呢?.
我們看二賽書第61章1至2節.

$^{841}$就能夠看到.
究竟二賽書第61章1至2節說些什麼呢?.
這段經文是剛剛在今日經文的後一點.
這段經文可能讀的出處比較陌生.
但是當你讀的時候.
你會知道這段不是陌生的經文.
我們一起再讀吧.
第61章1至2節.
主耀華的靈在我身上.
因為耀華用高高我.
叫我傳好訊息給謙卑的人.
差遣我醫好傷心的人.
報告被擄的得釋放.
被囚的出監牢.
報告耀華的恩憐.
和我們臣報仇的子.
安慰一切悲哀的人.
是不是好像似曾相識呢?.
這段經文誰說過呢?.
除了二賽書之外.
我們知道耶穌說的話.
當耶穌滿了三十歲.
他在會堂讀聖經的時候.
他第一卷的書卷打開.
他就是在讀二賽書61章1至2節.
而他讀完之後說什麼呢?.
他說這是二賽的預言.
今日要應驗在你們的耳中.
換句話說.
耶穌說.
主前七百年二賽所說的彌賽亞.
我就是了.
我來到就是為了福音的緣故.
令你們的屬靈生命得到更新和復興.
耶穌來到不只是醫治我們身體這麼簡單.
他是復興我們的生命.
所以二賽說.
天國是一個得自由得勝的地方.
在天國裡面人再不會被罪了捆綁.
剛才大師也唱了一首詩歌很有意思.

$^{881}$令我們的生命因為福音不再被罪所捆綁.
相反在主的裡面.
心靈得著自由.
生命能夠得著釋放.
天國帶給我們這麼重要的滋味.
不過我們知道.
即使在今天2026年3月的時間.
我們都已經品嘗到這個救恩的滋味.
同意嗎?.
因為藉著福音的緣故.
我們的生命都得著更新.
如果我給大家機會用一張紙.
去寫下究竟救恩對你們有什麼改變.
我相信大家至少寫到一樣.
而這個改變就是上帝怎樣去更新你們的生命.
所以天國是一個和平得勝的地方.
同樣教會既然是一個去品嘗到天國滋味的時候.
這個地方同樣是讓人得勝的.
教會同樣是不斷去建造生命的一個地方.
教會並不是一個彼此攻擊一個的地方.
不知道大家有沒有看過這幅圖畫.
有人點頭.
這幅圖畫是一個畫家.
但他本身是一個墓者.
他對畫畫很有恩賜.
他先畫了左邊的那幅圖畫.
左邊的那幅圖畫大家看到有什麼狀況嗎?.
是真戰 是不是?.
有兩個人在打兩個人.
那誰在打誰呢?.
公仔一個穿著軍裝.
拿著一把劍在打什麼呢?.
戴著頭上有兩隻東西 應該代表魔鬼撒旦.
這把劍代表什麼呢?.
是聖靈的話語 神的說話.
而這幅圖畫表達了神對信徒 對教會.
一個很重要的責任就是.
我們的真正對象是撒旦.
我們同心一起用神的說話去抵擋魔鬼的攻擊.
他的試探 他的罪.

$^{921}$他對教會的破壞.
左邊是合乎聖經的心意.
但那個墓者還沒畫完 他還畫了右邊.
右邊有什麼不同呢?.
都是真戰的.
不過很有趣.
兩個穿著軍裝的人自己打自己.
旁邊有人喝彩 是誰呢?.
就是魔鬼撒旦.
今天在教會裡面.
這個可以體驗到天國地方的地方.
我們屬於哪一個圖畫呢?.
我們是否認定惡者是我們唯一的敵人.
還是有時候我們都會模糊了 我們誤以為其他是我們的敵人.
聖經裡面很多的教會 比如哥林多教會.
其實裡面都是分黨分派.
有人去追捧保羅 有人追捧亞伯羅.
但原來神說我們共同的敵人應該是那個惡者.
所以頂智慧教會應該是一個德性的地方.
甚至我們形容為教會是一個生命建造的工場.
換句話說 每個人進入一個教會.
他的生命都會被建造的.
未認識上帝的.
使這個地方有機會更加認識我們的上帝.
已經認識了的 無論初信也好.
或者已經是一段時間的.
我們的生命都能夠彼此的建立.
教會是一個建立的地方.
是讓生命不斷成長的一個地方.
而這個正正是成為一個試味 一個場所.
頂智慧願意我們一起去參與.
令到教會是一個建立的地方.
不是單單靠節牧師或者Nelson傳道.
而是我們一起同心去建造.
因為工程不是一件輕易的事情.
而是要大家一起去守護.
令到這個地方是建立的.
學習事奉的 鼓勵他們.
誰能夠長出來懂得事奉呢?.
處分不好的時候.

$^{961}$我們除了給一些負面的.
反而說不要緊 繼續一起去學習.
因為我們一起向著這個成熟的地方去邁進.
以下書六十章.
去講述天國有什麼味道.
天國是一個吸引萬人前來的地方.
無論是什麼文化.
什麼階層的人都能夠去參與.
各位知道很多的事情發生.
很多運動出現.
或許會讓人有不同的看法.
但是這些看法會不會影響.
我們一起去敬拜同一個上帝.
不同的文化背景.
不會阻攔我們去敬拜同一個神.
天國是一個呈獻的地方.
是一個地方人願意去參與.
願意去付出 願意去服侍我們的上帝.
最後天國或教會是一個和平得勝的地方.
讓每一個在教會的生命都不斷被建立,被鼓勵.
最終是讓神的名得著榮耀.
願上帝祝福.
讓我們每個教會都能夠在一個地方.
讓人去taste天國的味道.
我衷心的祈禱.
天父 我多謝你的說話.
透過這次才知.
我們如何去預言永恆的角度的味道.
而我們知道這個味道.
雖然今天未完全的實踐.
但是在我們當中.
在這個世界已經存在.
並且透過信徒的聚集.
透過教會,小組.
我們能夠品嘗得到.
天父 我願意我們都能夠成為你一個美好的見證人.
讓你的芬芳能夠透過每個群體.
能夠去呈現得到.
而最重要的每一個生命.
當進到這個教會的時候.

$^{1001}$我們聞到這個味道是外面沒有的.
但是在主裡面是有的時候.
我們會被這個味道所吸引.
我們願意留下.
繼續去認識你.
並且願意繼續去愛你和服侍你.
我求你去祝福我們.
讓我們成為一個媒介.
去反照你的榮利.
求你親自對我們說你的心意.
我們禱告五叩.
救主耶穌基督明智祈求.
阿們.
\newpage



\section{以西結書 37:1-14}
\label{sec:031841}
\textbf{我現在是「活」的嗎?}~~~~林家棋宣教師
\newline
\newline
連結: \href{http://www.church.com.hk/acms/content.asp?site=cdc\&op=show\&type=product\&code=031841\&layout=sermon}{\texttt{031841}} ~~~~ 語音日期: 2026-03-21
\newline
\newline
\hyperref[sec:031840]{\small{< < < PREV SERMON < < <}}
~
\hyperref[sec:index_chronic]{\small{[返順時目]}}
~
\hyperref[sec:031842]{\small{> > > NEXT SERMON > > >}}
\newline
\newline
以西結書 37:1-14
\newline
\begin{longtable}{cl}
\hline
\hline
章節 & 經文 (和合本修訂版)\\
\hline
37:1 & \begin{tabularx}{0.7\textwidth}{X} 耶和華的手按在我身上。耶和華藉著他的靈帶我出去，把我放在平原中，平原遍滿骸骨。 \end{tabularx} \\ \\ \relax
37:2 & \begin{tabularx}{0.7\textwidth}{X} 他使我從骸骨的四圍經過，看哪，平原上面的骸骨甚多，看哪，極其枯乾。 \end{tabularx} \\ \\ \relax
37:3 & \begin{tabularx}{0.7\textwidth}{X} 他對我說：「人子啊，這些骸骨能活過來嗎？」我說：「主耶和華啊，你是知道的。」 \end{tabularx} \\ \\ \relax
37:4 & \begin{tabularx}{0.7\textwidth}{X} 他又對我說：「你要向這些骸骨說預言，對它們說：枯乾的骸骨啊，要聽耶和華的話。 \end{tabularx} \\ \\ \relax
37:5 & \begin{tabularx}{0.7\textwidth}{X} 主耶和華對這些骸骨如此說：『看哪，我必使氣息進入你們裡面，你們就要活過來。 \end{tabularx} \\ \\ \relax
37:6 & \begin{tabularx}{0.7\textwidth}{X} 我要給你們加上筋，長出肉，又給你們包上皮，使氣息進入你們裡面，你們就要活過來；你們就知道我是耶和華。』」 \end{tabularx} \\ \\ \relax
37:7 & \begin{tabularx}{0.7\textwidth}{X} 於是，我遵命說預言。正說預言的時候，有響聲，看哪，有地震；骨與骨彼此接連。 \end{tabularx} \\ \\ \relax
37:8 & \begin{tabularx}{0.7\textwidth}{X} 我觀看，看哪，骸骨上面有筋，長了肉，又包上皮，只是裡面還沒有氣息。 \end{tabularx} \\ \\ \relax
37:9 & \begin{tabularx}{0.7\textwidth}{X} 耶和華對我說：「人子啊，你要說預言，向風說預言。你要說，耶和華如此說：氣息啊，要從四方而來，吹在這些被殺的人身上，使他們活過來。」 \end{tabularx} \\ \\ \relax
37:10 & \begin{tabularx}{0.7\textwidth}{X} 於是我遵命說預言，氣息就進入骸骨，骸骨就活過來，並且用腳站起來，成為極大的軍隊。 \end{tabularx} \\ \\ \relax
37:11 & \begin{tabularx}{0.7\textwidth}{X} 他對我說：「人子啊，這些骸骨就是以色列全家。他們說：『看哪，我們的骨頭枯乾了，我們的指望失去了，我們滅絕淨盡了！』 \end{tabularx} \\ \\ \relax
37:12 & \begin{tabularx}{0.7\textwidth}{X} 所以你要說預言，對他們說，主耶和華如此說：我的子民，看哪，我要打開你們的墳墓，把你們帶出墳墓，領你們進入以色列地。 \end{tabularx} \\ \\ \relax
37:13 & \begin{tabularx}{0.7\textwidth}{X} 我的子民哪，我打開你們的墳墓，把你們帶出墳墓時，你們就知道我是耶和華。 \end{tabularx} \\ \\ \relax
37:14 & \begin{tabularx}{0.7\textwidth}{X} 我必將我的靈放在你們裡面，你們就要活過來。我把你們安置在本地，你們就知道我---耶和華說了這話，就必成就。這是耶和華說的。」 \end{tabularx} \\ \\
[1ex]
\hline
\hline
\end{longtable}
$^{1}$在進入今天的信息之前,我們先有低頭的一個祈禱.
張教父神,我們都將今天的信息都交上給你.
求主你真的讓到今天的信息去打進我們每一個人的心目中.
主耶,我們都求你.
求你真的讓到我們每一個人都能夠將你的寓話語.
真的活在這個世代之中.
我們以上的禱告,奉主耶穌基督成為祈禱.
在進入今天的經文之前,我們都先去了解這次經文的結構.
在整卷的48章的以西結書之中.
33至39章,其實內容都是在說同一個主題.
那個叫做Restoration,或者我們叫做恢復這個主題.
我們簡略去理解這7章的經文大概是說些什麼.
好,我們大概了解一下這7章的經文大概的概念是什麼.
33章其實是在說守望者的警告以及存活這個主題.
然後去到34章的時候,就是關於惡牧人,好牧人的元素.
以及它當中含有一些彌賽亞的盼望.
去到35章至36章的15節.
就開始真的說關於恢復這個主題.
這個位置就是在說土地的恢復.
去到36章的16至38節尾.
就是關於恢復上主聖名的主題.
而去到我們今天的37章1到14節.
就是說關於恢復以色列民作為神子民的身份.
然後打後37章15節至尾就是恢復文曰.
去到38至39就是說關於哥格瑪國.
瑪國的比喻是關於一些宗末的審判.
然後上主的以色列民得以回歸.
並且因為回歸令到上主得以的神聖可以顯出來.
我們今天思想的經文是集中於剛才在說的就是37章1到14節.
關於上主要恢復以色列民作為神子民這個主題.
當我們去看以西結書的時候會發現.
這本書其實和其他先知書有少少不同之處.
就是經文 如果大家有看過以西結書.
大家會發現上主頒佈信息的時候.
先知先知在做什麼 先知好像在做drama.
先知好像在做drama 真的在做一台戲出來.
而每一齣劇 每一齣戲其實都包含著上主的每一個信息.
在今天這一齣劇是關於恢復神子民的契劇.
我們會發現其實今天的經文裡面有很多重複的字眼.
我都不一一讀出來 都放在powerpoint上面.

$^{41}$經文去運用這些重複的字眼.
去帶出這一齣劇的主題.
而這一齣劇發生在一個平原.
讓我們一齊去看看這一齣劇怎樣去發生.
好 要求大家給點時間打開今天的經文.
以西結書37章1到14節 給一點時間大家打開.
一邊看一邊聽.
以西結書37章1到14節.
好 以西結書37章1到14節.
我們看看 我們看一看在第一節.
先知已經被上主的靈帶到一個平原.
這個平原其實我們可以把它翻譯成一個山谷.
這個山谷的景象 大家看剛才讀的時候.
不知道大家有沒有留意 其實這個景象是很嚇人的.
經文形容這個山谷遍滿骸骨.
我們可以理解這個山谷 其實可以想像成一個亂葬崗.
上主將先知帶到這個亂葬崗.
對於先知來說 其實這個場面很震撼 很震驚.
第二節我們會看到先知用了兩個看來表示這種震撼.
如果大家有去過 大家有沒有人有去過柬埔寨.
沒有人去過柬埔寨 承恩有去過 Esther有去過.
柬埔寨有一個很出名的地方.
Killing Field 中文翻譯為殺戮戰場.
這個地方是遊客區.
這個地方有兩個特別之處 第一就是這座塔.
大家看到這座塔很漂亮.
知道裡面放了些什麼嗎.
不是死人 是全部骷髏骨頭 人頭.
在整座塔裡面.
這是其中一個塔裡面的特別之處.
另外這個Killing Field的另一個特別之處.
就是有很多個一個個的水坑.
但這些水坑下面全部都是人骨.
經過的時候你就會聞到一陣味道.
所以其實那時候我去的時候.
我去這個Killing Field 當時我去柬埔寨訪孫.
孫教士帶我去這個地方去看看這裡的歷史.
其實這裡當時的場面的確是很震撼.
因為裡面其實是佈滿著當時赤柬殺害人民的骷髏骨頭.
當然我們回到先知的時候 我們看以西結書.

$^{81}$他同樣都是身處一個充滿著骷髏骨頭的地方.
如果我們看回英文的譯本.
我們見到先知形容這些骷髏骨頭多度是用Very Many.
而且這些骷髏骨頭是極其Very 富剛.
所謂極其富剛即是這些骷髏骨頭的主人.
其實已經死了很久 剩下的只有這些骨.
首先 親愛的姐妹 這是上主比爾西傑經歷的一個異象.
並不是一個真實真真正正的一個亂葬崗.
不過我們很好奇地問.
上主為何要帶先知來到這個亂葬崗.
所以知道亂葬崗背後的意思.
其實就是這些骷髏骨頭的主人.
他們沒有被好好地安葬.
就好像我們俗語所說的死無葬身之地.
我們中國傳統文化裡面.
人死了要入土為安.
但是對於猶太人來說.
他們同樣有這個概念.
人死了要有一個地方好好地安葬.
但是現在這些骷髏骨身處在山谷之中.
沒有被埋葬.
背後的意思是什麼.
容後再揭曉.
當先知在這個亂葬崗行完一轉之後.
上主問了先知一個很奇怪的問題.
大家看一看.
他說他對我說 人子啊 這些骷髏骨能活過來嗎.
這個問題對於我們現代基督徒而言其實不難回答.
但是對於以西結來說不是那麼容易想像.
甚至覺得上主是不是在耍花樣.
在舊的時代對猶太人而言.
他們其實都有死人復活這個概念.
如果我們有印象列王記尚的時候.
以利亞幫過撒勒伐的婦人的兒子.
撒勒伐寡婦的兒子去復活.
對於以西結來說.
其實他都有這種概念.
死人會復活的.
但是我們細心想想.
我們細心看看以西結怎麼回答.

$^{121}$以西結說 主耶和華你是知道的.
其實你見到以西結他的答案是很模棱兩可.
他不敢回答YES.
也不敢回答NO.
其實你理解到這個答案.
他其實是把球交回給上主.
他不敢回答.
因為他在他的認知裡面.
死人當然是可以復活.
不過是當肉身未腐爛的情況之下.
才能夠復活.
現在上主問這一堆骸骨可不可以復活.
一個連肉 連筋 連皮都沒有的骨頭.
請問怎樣可以活過你.
不過上主沒有繼續這個對話.
反而吩咐先知以西結向那些骸骨去發預言.
發預言即是作先知的說話.
而上主的吩咐是採用一個命令語去表達.
我們就出現了第二個疑問.
先知的工作其實應該是要向人說話的.
轉述上主的說話.
但是如今以西結的對象不是活生生的人.
而是非常枯乾的骸骨.
一個沒有生命氣息的死物.
要向他講上主的說話.
如今以西結的對象已經不是一個活生生的人.
而是一個死物.
看到這裡.
一齣劇 一台戲究竟在搞哪一課.
其實我們再看下去的時候.
第十一節答案已經揭曉了.
以西結真正要講先知的說話的對象.
其實是以色列人.
一群剛剛經歷亡國被擄以及聖殿被毀的以色列人.
而上主以極其枯乾的骸骨來比喻他們.
比喻他們整個以色列民族的靈性死亡.
在於當時以色列的群體.
他們其實正處於一個絕望的境況.
在第十一節我們看見.
這群百姓其實發出了三重哀嘆.

$^{161}$在原文裡面.
這三重哀嘆是採用了一些押韻的音節.
就是….
然後是….
還有這個….
即是我的骨頭枯乾.
我的指望失去了.
我們滅絕正盡了.
這三個是押韻的音節.
這三個的押韻音節是層層遞進.
由骨頭枯乾.
去到嚴重到一些處於絕望.
沒有盼望而去到最後.
我們更加見到一個嚴重的情況就是.
真正被滅絕.
同時以色列民這三重哀嘆.
表達其實他們整個民族已經被上主所拋棄.
國家被亞述巴比倫欺壓.
並且入侵國土.
連代表上主同在的象徵聖殿都被毀.
聖殿裡面的神聖器皿也被擄去.
成為敵人的戰利品.
上主已經離開了他們.
所以以色列百姓的絕望.
絕望到一個極點.
就好像一堆枯乾了的骨頭一樣.
沒有肉 沒有筋.
甚至連氣息都沒有.
每日的生活都是充滿悲傷 哀痛之中.
弟兄姊妹 我們會不會好像以色列民一樣.
生命好像一堆骸骨一樣極其枯乾.
生命極其枯乾.
其中一個可能性就是來自我們環境因素.
可能因為來自我們的工作或學業的忙碌.
所帶來的勞累.
導致我們自己身體和靈魂都開始出現疲乏.
很累 很累.
但我們任由這些工作 學業.
去將我們榨乾.
榨乾到一個地步.

$^{201}$其實可能最後影響我們自己的屬靈生活.
我們被現實的生活取締了原本.
原本我們作為基督徒的使命.
這一種枯乾的生命.
亦都可以來自我們過度的事奉.
以服侍來取締親近神.
以為事奉得多就是靈明與好.
事情恰恰相反.
用盆栽來做個比喻.
一棵盆栽裡面其實可以不停地長出很多的枝子.
很多的葉.
但這些枝子和葉.
都是來自原初那棵盆栽裡面.
你買回來的那棵盆栽裡面的泥土和肥料.
若果這棵盆栽我買回來.
我不去用水澆灌.
亦都不再換新泥.
不再加入新的肥料.
很快新生出來的枝子和葉.
都是會枯萎.
原先的養分其實已經完全被吸取.
我們的靈明都是一樣.
我們以為只是沿用舊有.
我們屬靈的根基.
我們就可以以為不停地支出.
但是實情.
我們原有的養分可能已經早被透支.
不足以去支撐我們不停地付出.
又或者面對不同的屬靈挑戰.
以致為何我們好像面對事奉.
又或者面對這些挑戰.
好像沒有生命力.
和看不到出路.
但是弟兄姊妹.
我們生命枯乾.
背後可能是我們自己疏忽了親近上主.
又或者被罪.
或者環境奪去了我們的靈明.
任由我們的屬靈生命慢慢地流走.
我們都知道我們的身體.

$^{241}$是需要營養去維持.
和我們的屬靈生命都是一樣.
我們都要吸收養分.
成為有靈的活人.
而不是枯乾的骸骨.
我們需要攝取的是新的屬靈營養.
我們才能夠保持靈明的健康.
在彼得前書二章二到三節裡面這樣說.
要愛慕那純淨的靈乃像初生的嬰孩.
愛慕那一樣.
好使你們藉著它成長以至得救.
因為你們已經嘗過主因的滋味.
我們再一次回到以西結書的內容.
我們去揭曉為何上主要帶以西結來到這個聯藏崗.
對以色列民來說.
他們面對的這個情況沒錯.
他們是正值面對著國家滅亡.
聖殿被毀.
歸根究底來自他們自己背叛上主.
然後觸發上主和他們之間的文約的詛咒.
如今他們毫無指望.
生命枯乾的情況.
但是上主告訴他們已經扭轉了.
在第五至第十節的經文裡面.
上主要使得這些骸骨生根長肉並且包上外皮.
骸骨的外形終於像一個人.
有這個外形.
但是這個只不過是一個軀殼.
未算是一個人.
上主再要在這個軀殼裡面再加上氣息.
這個才是活過來的人.
當我們去看第五至第十節的經文.
我們有一種似曾相識的感覺.
就是創世記裡面有關上主創造人的記載.
先有形體再有氣息.
以西結書的記載.
上主讓到骸骨成為有靈的活人.
是一種再創造的工程.
目的就是要讓一群失去生命力.
處於絕望的以色列百姓知道我是耶和華.

$^{281}$知道我是耶和華.
既是宣認上主才是獨一的真神之外.
更加重要的是.
恢復被以色列百姓破壞了的上主與子民之間的關係.
知道不是一種頭腦認知.
而是承認並徹底認識的一種關係.
所以上主透過再創造的行動.
目的是要恢復他與上主子民之間的關係.
同時是上主主動去修復這個被破壞的關係.
在第十二至第十三節的經文.
有幾句重複的說話.
就是我要打開你們的墳墓.
把你們帶出墳墓.
同時十二節更加多一個行動.
就是領你們進入以色列地.
而第十三節再出現一次.
知道我是耶和華.
這兩節經文其實是回應第十一節.
以色列百姓的三重哀嘆.
原本他們以為上主已經離棄了他們.
如今上主卻是親自打開.
本來屬於死人所住的墳墓.
並且要將他們帶離開這個死人的住處.
他們成為了有靈的活人.
不再屬於墳墓.
更加重要的是.
他們現在失去的土地.
亦都會被上主拉他們回去.
進入應許之地.
第十四節再一次.
被上主的靈.
將自己的靈放在這些再創造的人身上.
從而令到這班百姓活過來.
有生命靈.
上主親自回應以色列百姓的哀嘆.
從絕望之中.
看見盼望和出路.
當然.
這種拯救不單單只是上主.
去對以色列百姓的拯救.

$^{321}$來到如今這一刻.
我們信靠耶穌基督.
讓到聖靈能夠進入我們心裡.
以致我們有聖靈的提醒和幫助.
能夠成為真正的活人.
弟兄姊妹.
我們每一天.
都會面對生命上不同的挑戰.
但我們又怎樣去回應和面對.
我們會不會好像以色列百姓一樣.
被眼前生活或者困境所限.
令到我們自己感到絕望.
又或者失去了活力.
過著得過且過的生活.
生命就好像骸骨一樣極其富剛.
我們嘗試調轉去想一想.
回顧一下.
我們生命充滿活力的時候是何時.
當時的狀況是怎樣.
當時是否有一股很想為主作供發光發熱的興奮.
我們有這種興奮.
我們有這種熱心.
可能就是我們和上主關係最好的時間.
我們生命富剛不多不少.
我們沒有再一次吸收新的養分.
以致我們屬靈生命慢慢枯萎.
每一次面對不同挑戰.
都看不到出路和盼望.
弟兄姊妹.
上主今日一樣和我們說.
我必使氣息進入你們裡面.
你們就要活過來.
容讓我們去求上主讓聖靈進入我們的生命.
再創造我們的生命.
以致我們的生命不再枯萎.
不再好像骸骨一樣極其富剛.
誠心所願.
\newpage



\section{創世記 40:23}
\label{sec:031842}
\textbf{聖經人聖經事:人會忘記你}~~~~黃克勤牧師博士
\newline
\newline
連結: \href{http://www.church.com.hk/acms/content.asp?site=cdc\&op=show\&type=product\&code=031842\&layout=sermon}{\texttt{031842}} ~~~~ 語音日期: 2026-03-18
\newline
\newline
\hyperref[sec:031841]{\small{< < < PREV SERMON < < <}}
~
\hyperref[sec:index_chronic]{\small{[返順時目]}}
~
\hyperref[sec:031843]{\small{> > > NEXT SERMON > > >}}
\newline
\newline
創世記 40:23
\newline
\begin{longtable}{cl}
\hline
\hline
章節 & 經文 (和合本修訂版)\\
\hline
40:23 & \begin{tabularx}{0.7\textwidth}{X} 然而，司酒長不記得約瑟，竟忘了他。 \end{tabularx} \\ \\
[1ex]
\hline
\hline
\end{longtable}
$^{1}$聖經人聖經事由生命說書的網絡為大家預備.
今天講的題目是「人會忘記你」.
人是會忘記你的.
我們學聖經的人物,學聖經的故事.
是想看看背後上帝是怎樣.
認識他多些,他的手怎樣工作,怎樣能力.
你都會看到,人是會忘記你的.
哎呀,我這麼厲害,怎會忘記?.
個個都說厲害,個個都說記住.
不過好多人丟埋石像都好快被人家忘記了.
是會忘記的,因為人的記憶力都好短暫.
又好多東西分了心,會忘記.
難道真是一世人將你放在心裡面,動都不動嗎?.
人是會忘記你的.
哎呀,這次一定要記得.
好多時叫做忘恩負義.
好了,人會忘記你.
不過其實你要知道,剛剛昨天說.
所有事情都是出於神的三個字.
出於神的.
約瑟故事最重要的就是他後來學到.
所有事都是出於神.
所以人雖然忘記,不過不要緊.
人是會忘記,神不會忘記.
於是酒正就將他的夢告訴約瑟.
我夢見在我前面有一棵葡萄樹.
樹上有三斤枝子.
它好像發了芽,又開了花.
然後上頭的葡萄都成熟了.
那發露的杯仍然在我手裡面.
咦,我還拿著發露喝酒的杯子.
這杯子是他最喜歡的.
那我就拿葡萄,把葡萄摘下來.
就捏葡萄汁在發露的杯子裡面.
我就將杯子遞給他.
這就是酒正的夢.
夢見在我前面有一棵葡萄樹.
葡萄樹不是這樣的.
葡萄樹是爛藤的.
有一個很矮的架.

$^{41}$葡萄是翹著的.
葡萄是一棵一棵的.
這樣摘下來.
話說回來,釀酒的葡萄.
和我們會吃的葡萄籽是不同的.
不同品種的.
那些葡萄籽如果是釀酒的.
是不好吃的.
不是很甜的嗎?不是很甜的.
味道是很有趣的.
它是用來釀酒的.
簡單一點,夢就是摘下來.
捏葡萄汁在皇帝的杯子裡面.
就有葡萄汁了,遞給皇帝喝.
約瑟對他說.
你所做的夢是這樣的.
三根之子就是三篇.
為什麼?不要問了.
約瑟這樣解釋,上帝指示他的.
三根樹枝就是三篇.
三篇之內法老會提你出監.
叫你第一,你官服完職.
這是第一點要講的.
你官服完職.
你仍要遞杯在法老的手中.
和先前所作的酒證一樣.
你會恢復官職做酒證.
這些雖然不是幾多品的官.
但是所有人都尊敬的.
因為你和皇帝那麼近.
遞個杯給他.
每個人都說你和皇帝當然很熟.
每個人都連一品大員都怕你.
因為你是和皇帝親近侍候.
三篇之內你會官服完職.
這個就是第一點所講的.
有畫本很好.
一跌下來又上去.
這是人生很多時候會遇到的遭遇.
跌下來很快一下.

$^{81}$然後三天又上去.
這個V形反彈很多時候都會發生.
其實不是你厲害.
是不是酒證好東西.
不是,只是純粹上帝的手在後面工作.
令到他不歡喜他就坐牢.
叛君如叛婦就是這樣.
一不高興一拍台就要坐牢.
說不定會更改.
第二就是回來繼續斟酒.
現在斟酒不鹹不淡都不會酒.
酒是什麼都不熟.
回來.
於是他就回來.
官服完職.
V形反彈.
第二個.
但你得好處的時候.
求你紀念我.
因為求你施恩我.
紀念我.
因為其實我在法老面前.
跟我說好說話.
救我出來開這監獄.
其實我是希伯來人.
是被人拐來這裡.
拐這個字都很少用.
是被人拐帶過來.
我在這裡都沒有做過壞的事.
叫我應該要坐監.
我沒有做過任何壞事.
那次其實是這樣.
所以我就進來.
禪掌見到.
你解酒精的夢都解得很好.
不如我又講我的夢給你聽.
我的夢是這樣的.
在夢中我頭上頂著三糠的白餅.
頭上很多時候.
他們是將東西放在頭上拿的.

$^{121}$現在很多人都是這樣.
以前特別多.
去到印度更加多.
放.
又放得很穩.
中國很多人是將重物放在頭上.
頂著走的那些女人.
很厲害的.
在頭上有一個竹籃.
竹籃或者用植物.
藤籃,竹籃之類.
裡面有三糠的白餅.
三糠白餅不是三個.
是一個圓形的植物.
或者藤或者竹的器皿.
裡面放了些餅.
有三個這樣的筲箕盤.
頂著在頭上.
極上的糠子裡.
特別是有為法老烤的各種食物.
包括餅.
誰知有飛鳥來吃我頭上糠子裡的食物.
最頂那個沒有遮沒有掩.
很漂亮.
而且是為法老烤食物.
誰知有飛鳥來吃了那些食物.
那藥師就說 哎呀 糟糕了.
你的夢是這樣解的.
三個糠子就是三天一樣.
不過三天之內.
第三了.
法老會斬斷你的頭.
將你掛在木頭上.
有飛鳥來吃你身上的肉.
這樣說 你這個夢就不好.
就說先生.
因為三天內你會死的.
你會被斬頭.
因為三個糠子代表三天.
而且吃最頂那個就是吃了頭.

$^{161}$我說很害怕.
為什麼?.
因為法老不喜歡你.
你那些食物不知道被誰吃了.
我不知道發生過什麼事.
譬如會不會有可能就是.
這個膳長將法老吃不完的食物.
私下送給別人吃.
拿些回家.
被人傳回來.
皇上 皇上 膳長拿了你的食物回家吃.
或者給誰吃.
分給下人吃.
但是不知道多開心 說好吃.
我的食物拿給別人吃.
不是吧 坐牢.
可能不知道發生什麼事.
總之你是將法老吃的食物.
最頂的糠子.
有飛鳥來吃了.
這樣就沒命了.
這個就是第三個.
必斬斷你的頭.
這個是插水式.
死了.
糾正是V型反彈.
人生很多時候都是這樣.
有些就是這樣了.
老老實實有多少個是V型反彈的.
很多都是插水下去就去了下面.
不過無論怎樣都是在神的手裡面安排的.
你要盡力將自己的事做好.
小心做.
害人之心不可有.
防人之心不可無.
這個世界的人都很險惡的.
這個世界都是惡在惡者手下.
我都試過會被人害.
他們想都沒想過.
是你幫過他.

$^{201}$有恩於他的人會害你的.
所以我現在很明白.
本來我心裡面很多不開心.
不過後來學了藥室之後.
這些事是出於神的.
是神差派我.
令我經過學到這些東西.
現在我都長進了.
插下去.
幸好神拉我上來.
每個人都不一樣.
你就會死了.
到了第三天.
法老的生日.
就大排筵席.
眾神設擺筵席的時候.
就將酒政和善長提出監內.
叫那兩個出來.
當然你擺筵席的時候.
記起預備食物和斟酒的那個人.
預備酒的那個人.
你叫那兩個出來.
其實法老還在生那個善長的氣.
做飲食的那個人.
好,你就.
酒政仍然是為我做遞杯.
給我喝酒.
善長呢.
給我斬頭.
掛起來.
就好像藥室那樣講.
所解的話.
即是法室:哎呀,皇上要命呀.
都沒用了.
好,事情這樣發生了.
嘩,藥室解夢很準.
所以當他兩個被叫上來.
見到法老的時候.
特別是酒政官府完職的時候.
現在你馬上換件衣服.

$^{241}$斟酒給我喝.
以後繼續做這個.
侍候我的工作.
嘩,多謝皇上,多謝皇上.
當然在這個多謝期間.
不斷地小心.
馬上洗澡換衣服的期間.
等等.
在裡面.
嘩,好彩席.
你一定會想.
嘩,這個藥室.
這個希伯來少年.
解夢真的挺厲害的.
一定會這樣想起的.
誰知道.
酒政卻不紀念藥室.
四個字.
竟忘了他.
即是說.
應該不會忘記.
但是竟然會忘記.
有沒有這樣的經歷過呢.
有的.
我也有經歷過.
也有做過.
一些很重要的事情.
竟然忘了那個人.
會的.
竟然忘記了.
哎呀,慘了.
他忘記了.
這樣也有.
竟然忘記了.
這個竟然忘記了.
也是在神的手裡.
應該會記得的.
但是竟然會忘記.
你猜神的手會不會在背後做一些事情呢.
當然他可能是老人.

$^{281}$有些失憶.
忘記了.
或者太緊張.
忘記了.
或者他知道要跟藥室說話.
不過在潛意識裡.
還是要小心.
不要隨便說.
因為他說得那麼準.
哎呀,不知道會不會上來搶了我的位子.
哇,他說得那麼準.
可能都被那個擅長的被斬首.
嚇怕了.
很多在心裡抑制了.
竟然忘記了.
無論什麼原因都好.
竟忘了.
就是這樣.
因此你記得.
人是會忘記你的.
雖然你有恩於任何人.
就不要以為他一定會將你的恩記到永遠.
又或者你.
哎呀,漂亮的人.
於是他一定會記得我.
又或者我那麼出眾.
他一定會記得我.
讀書那麼多.
一定會記得我.
那次立了大功.
他一定會記得我.
不是的.
所有事都是三個字.
出於神.
如果是出於神.
在這裡.
不出於神,不在這裡.
神叫那個酒精繼續恢復原跡.
叫那個擅長改矮一個頭.
就發生了.

$^{321}$他令到那個酒精忘記了約瑟.
竟忘了他.
那都是那個酒精不好.
酒精是不記念約瑟.
因為聖經特別這樣說.
酒精不記住約瑟.
不記念是指不記住.
出來了就好了.
過得開就是神仙.
不記念.
酒精是會這樣的.
軟弱的.
脆弱的.
如果你將自己所有的東西.
都想在人的身上.
你就會非常之失望.
一定要記得.
出於神.
神不會忘記.
人會忘記.
交代給主.
我們一起祈禱吧.
主我們見到這個酒精忘記約瑟的時候.
心裡面真的覺得.
很歎息.
真的歎一口氣.
為什麼這樣都有.
他竟然這樣V型反戴官復原跡.
竟然會忘記了和他解夢的約瑟.
三天前.
也見到那個.
真的三天之後.
他就被斬頭了.
一定記得三天之前約瑟解的夢.
主面那麼大.
在監房和回到皇帝皇宮面前.
還是皇帝面前.
不會忘記的.
但是竟不忘記.
有時是會有意無意忘記了.

$^{361}$不過.
你沒有忘記.
我們就將自己交在你的手裡.
人須忘記.
你卻不忘記.
這個就是我們信靠你的地方.
知道你是無所不能.
是不會失憶.
是不會意外的.
求你親自憐憫.
幫助我們更多信靠教導你.
祈禱奉耶穌名求.
阿們.
\newpage



\section{詩篇 130:1-8}
\label{sec:031843}
\textbf{從深處的覺醒}~~~~郭志榮牧師
\newline
\newline
連結: \href{http://www.church.com.hk/acms/content.asp?site=cdc\&op=show\&type=product\&code=031843\&layout=sermon}{\texttt{031843}} ~~~~ 語音日期: 2026-03-22
\newline
\newline
\hyperref[sec:031842]{\small{< < < PREV SERMON < < <}}
~
\hyperref[sec:index_chronic]{\small{[返順時目]}}
~
\hyperref[sec:031844]{\small{> > > NEXT SERMON > > >}}
\newline
\newline
詩篇 130:1-8
\newline
\begin{longtable}{cl}
\hline
\hline
章節 & 經文 (和合本修訂版)\\
\hline
130:1 & \begin{tabularx}{0.7\textwidth}{X} 耶和華啊，我從深處求告你！ \end{tabularx} \\ \\ \relax
130:2 & \begin{tabularx}{0.7\textwidth}{X} 主啊，求你聽我的聲音！求你側耳聽我懇求的聲音！ \end{tabularx} \\ \\ \relax
130:3 & \begin{tabularx}{0.7\textwidth}{X} 耶和華啊，你若究察罪孽，主啊，誰能站得住呢？ \end{tabularx} \\ \\ \relax
130:4 & \begin{tabularx}{0.7\textwidth}{X} 但在你有赦免之恩，要叫人敬畏你。 \end{tabularx} \\ \\ \relax
130:5 & \begin{tabularx}{0.7\textwidth}{X} 我等候耶和華，我的心等候；我也仰望他的話。 \end{tabularx} \\ \\ \relax
130:6 & \begin{tabularx}{0.7\textwidth}{X} 我的心等候主，勝於守夜的等候天亮，勝於守夜的等候天亮。 \end{tabularx} \\ \\ \relax
130:7 & \begin{tabularx}{0.7\textwidth}{X} 以色列啊，你當仰望耶和華，因耶和華有慈愛，有豐盛的救恩。 \end{tabularx} \\ \\ \relax
130:8 & \begin{tabularx}{0.7\textwidth}{X} 他必救贖以色列脫離一切的罪孽。 \end{tabularx} \\ \\
[1ex]
\hline
\hline
\end{longtable}
$^{1}$弟兄姊妹早晨 主內平安.
今天來到今年大齋期的最後兩個禮拜.
今天宣講的經文也是來自.
這個禮拜經課的詩篇.
就是詩篇130篇.
我們剛才已經讀過.
為什麼大齋期會讀這篇詩篇呢?.
這篇詩篇是十五篇上行之詩的第十一篇.
現在的學者普遍都會同意.
就是這十五篇的上行之詩.
其實是朝聖者每年上行到耶路撒冷過節.
在那個朝聖的旅程.
或者在某些特別的禮儀裡面.
去頌唱的詩歌.
有人形容大齋期.
其實本身就是一個.
和主耶穌一起走向十字架的一個.
朝聖的屬靈旅程.
詩句裡面有提到.
剛才大家讀過應該都會留意到.
有提到罪孽,赦免,救恩,救贖.
這些字眼.
在最後四世紀的時候.
奧古士丁已經將它和另外其他六篇詩篇放在一起.
後來稱之為懺悔詩.
而大齋期也是一個特別強調認罪和懺悔的節期.
剛才讀完這篇詩篇之後.
不知道大家會不會覺得.
這篇詩篇好像沒有什麼特別.
都是重複整卷詩篇裡面一些經常出現的主題.
不過宗教改革者馬丁路德.
就很高度評價這篇詩篇.
他認為它是聖經之中真正的大師和教師.
就是將福音的基本真理都勾畫出來.
另外大家都可能會聽過.
秦道慧理會的創始人約翰衛斯理的事蹟.
他在一個靈性陷入低潮的晚上.
他得著了得救確具.
心裡面感到異常溫暖.
這就是大家經常說的艾德門街的經驗.

$^{41}$約翰衛斯理在日記裡面特別提到.
在那天下午他有一個祈禱會.
而在祈禱會裡面他正好還唱了詩篇130篇.
到底這篇詩篇有什麼特別?.
為什麼得到這麼高的評價.
可以感動和激勵人的心靈呢?.
按著說話的對象來劃分的話.
這篇詩篇可以分為一至四節和五至八節.
兩個大段落.
前半部分是私人對神的說話.
而後半部分不是以神作為說話的對象.
而這兩個大段落上下.
各自可以分為兩個小段.
每一個小段包含了兩節經文.
這四個小段形成了這篇詩篇的四個篇章.
以下將會將這四個篇章逐一拆解.
讓大家可以理解這篇詩篇的訊息.
並且可以看到詩人的屬靈覺醒的進程.
在開始之前我們先有個祈禱.
在這個美好的早上.
引領我們弟兄姊妹一同回到你的殿當中.
一同去敬拜你.
願到我們所獻上來的祈禱.
我們的敬拜.
都成為你閱立的一個祭.
又願你的靈在我們當中親自去引領我們.
首先去加力給我.
也潔淨我的口舌.
以致我能夠有力為你按著正義宣講你的話語.
也願你的靈在眾弟兄姊妹心中.
去激發眾人的內心.
讓今天的訊息能夠作為眾人生命中的提醒和引導.
願你與我們同在.
奉主耶穌的名祈禱.
阿們.
我們先看第一個篇章.
開始講到詩人求告神.
詩人這樣說.
「耶和華,我從深處求告你」.
詩人在此特別強調.

$^{81}$他的求告是從深處發出.
也就是說.
詩人意識到自己陷於一個深處的處境之中.
深處即是甚麼呢?.
深處即是海的深處.
被一個人處於一個極之嚴峻的惡劣的處境之中.
好像魚液浸親一樣.
而且不是在水面掙扎.
而是已經沉到深海裡.
海在舊約聖經中.
經常用來代表那股為生命帶來破壞,毀滅和死亡的混沌力量.
而沉沒在海的深處.
就意味著人被大水所包圍和覆蓋.
因此與神隔絕了.
換言之.
詩人意識到自己的活力,動力.
以至他的生命氣息.
都正在逐漸流逝和消失.
他更加發覺.
單靠自己的力量去掙扎.
完全是無補於事.
只會令他有更加大的無力感.
面對著死亡的危機.
人經常都會有否認,憤怒,討價還價,沮喪的反應.
但詩人沒有因此被癱瘓.
他沒有跌入否定一切事情的狀態之中.
拒絕面對現實.
他沒有被憤怒的情緒所困.
終日怨天尤人.
質問為何不幸的事情要淋到他的身上.
他沒有糾纏在討價還價的空想裡.
妄想惡事,禍事會突然之間消失.
他沒有陷溺在沮喪的心境裡.
被焦慮和自憐自憫困住.
詩人就是以懇切求告神.
作為他對自己處境的應對.
他求神聽他的聲音.
則以留心他懇求的聲音.
「聽」這個字包含了回應的意思.
而懇求的聲音裡的「懇求」.

$^{121}$特別是在說祈求神的憐憫和私恩.
所以詩人在這裡.
他不只是希望神成為他的聆聽者.
他更加需要神對他的呼求.
作出法治憐憫之情的回應行動.
詩人這樣應對他陷於深處的困局.
表示他認識到自己無力自救的事實.
他更加知道神的回應和幫助.
才是他唯一的出路.
接著到第二個篇章.
詩人突然將話題轉到神的赦免之恩.
詩人對著神自問自答.
詩人這樣說.
「耶和華,你若救察罪孽主,誰能站得住呢?.
但在你有赦免之恩,要叫人敬畏你」.
「救察」這個字不只是指細心留意.
更加有記錄下來的意思.
詩人用反問的方式去表達.
神當然看到了人所作的一切罪孽.
神不會看不到這個世界的罪惡.
人一切的行為,言語.
甚至最隱密的邪惡思想.
神都看得一清二楚.
但神不將人這些犯罪的記錄.
鉅細無遺地收集起來.
作為審判人的根據.
否則根本沒有人能夠在神的審判台前.
站得住而存活下來.
這就是詩人信心的宣告.
但大家一定會問.
這位原本在深處求告神的詩人.
他心裡究竟想起了什麼.
而因此令他產生思想上的跳躍呢?.
答案很明顯.
詩人意識到他自己陷於深處並不是沒有原因.
他突然覺醒到他自己也要負上相當的責任.
這是詩人一個極之深刻的反省.
和對生命的洞見.
深處的覺醒使他認識到.
自己只能夠尋求神的幫助.

$^{161}$而這樣的覺醒和認識使他甦醒過來.
使他用一個新的角度回顧自己的人生.
重新審視自己生命的真相.
罪孽有扭曲和偏離的意思.
詩人知道這個問題的根源是他自己離開了神.
也就是他對神的不信.
他不信將自己的生命主權交在神的手裡.
會由自己去掌握更加穩妥.
他不相信神的旨意.
也就是聖經的教導.
會比他自己所謂理性的判斷更加正確.
他不信神對他的愛是如此深和徹底.
比他對自己的愛更加多,更加深.
弟兄姊妹.
深處之所以為深處.
重點不是在於外在的環境如何.
而是有沒有神的同在.
受造的人如果離開了次生命的創造主.
自然就會陷於深處而步向滅亡.
這是必然的事.
問題只是人有沒有這個覺醒而已.
詩人提到神有赦免之恩.
愚民直譯就是神有赦免.
莎士比亞曾經說過一句這樣的說話.
他說沒有任何東西比起憐憫.
更加能夠助長罪惡的膽量.
猶太拉比問過一個這樣的問題.
他說假如神總是赦免人的罪.
人會不會變得越來越不畏懼犯罪呢?.
答案當然不是.
有另一個拉比作了一個比喻來解釋赦免的作用.
他說假設一個債仔發覺他所欠的債務.
是他無論如何努力都不能如期償還的.
如果他的債主卻是堅持這個債仔.
一定要按著原本的承諾的時間.
去清還所有的欠債的話.
否則他就會立刻作出法律行動.
這個債仔很可能會立刻放棄繼續還債.
等債主申請他破產.
因為反正結果都是一樣.

$^{201}$但是假如債主願意延長還款的期限.
甚至去減免債仔部分的債務.
債仔反而覺得仍然有希望.
於是就會繼續還款.
這位拉比講得很好.
但是關鍵他仍然未點出.
最重要的是.
神的赦免.
赦免這個字除了是抹去一個人的罪孽紀錄之外.
更加包括了除罪的意思.
也就是說.
赦免是一個改變和轉化罪人生命的過程.
所以詩人在此特別強調.
神的赦免是要使人會因為感恩而敬畏神.
詩人在此特別強調.
這個要敬畏神.
記住.
赦免絕對不是縱容人去放肆.
繼續犯罪.
而是要令人的生命轉變.
活出一個敬畏神的人生.
這個才是赦免的目的.
接著詩人作出了一個堅定的信仰宣言.
這也是他的立志.
詩人這樣說.
「我等候耶和華,我的心等候.
我一樣望他的話.
我的心等候主.
勝於守夜的等候天良」.
在短短的這兩節經文中.
詩人反覆用了三個不同.
卻同樣可以譯作「等候」的希伯來字詞.
有些學者認為.
詩人在此其實是想表達.
詩人雖然仍然看起來好像無盡的黑暗之中.
但他仍然是懷著信心和盼望.
耐心地等候神的介入和拯救.
他相信光明總有一天會來到.
大家暫時等著吧.
不過這樣的理解過於消極.

$^{241}$還未能將三重等候的深層意義表達得清楚.
第一個等候所使用的字.
字根與「承索」有關.
大家一定會好奇.
「承索」和「等候」有什麼關係呢?.
大家可以從兩方面去理解.
首先一條「承索」是怎樣的呢?.
一條「承索」是由很多條較幼的繩繞在一起形成的.
引申出來其實是在說「等候」是怎樣的呢?.
「等候」就是將等候者和被等候的對象.
纏繞起來一起的意思.
另一個理解的方式更加簡單和直接.
「等候」的功能是什麼呢?.
就好像一條拉緊了的承索.
一頭連在等候的對象身上.
另一頭就是由等候者的手緊緊地拉著.
讓等候者能夠拉近他和等候對象彼此之間的距離.
這就是等候.
無論是哪一種的解讀.
第一重等候要強調的.
就是等候者和等候對象彼此有一個緊密的關係.
而且是一個生命結連的關係.
等候者絕對不是無可奈何.
只是聽天由命等運到.
而是他要主動地,努力地.
和等候的對象建立一個越來越緊密的生命結連.
他會積極地做好充分的準備.
去迎接等候對象的臨到.
詩人說「我等候耶和華,我的心等候」.
前面那半句「我等候耶和華」.
就是表明詩人等候的對象是哪一個?.
就是耶和華這位救贖以色列民的.
和他們有生命之約的獨一真神.
「我的心」的「心」字.
原文應該是「生命」.
而不是「心靈」.
所以應該說是「我的生命」.
其實是說詩人的整個生命.
也是指他整個人.
所以詩人想表達的.

$^{281}$就是他整個人包括他的身體.
包括他的思想.
包括他的情感.
包括他的意志.
都要結連在神身上.
都要交在神的手中.
成為神合用的器皿.
從此以後.
詩人生命的主權屬於神.
他人生目標和生活方式的選取.
都不會再離開神的心意.
詩人接著這樣說.
「我也仰望他的華」.
「仰望」這個字也包含了「等候」的意思.
所以有些英文譯本.
將這個字.
包括了NASB或者JPS的譯本.
都會將它譯作「wait for」.
或者譯作「await」.
就是這樣的原因.
第二重的「等候」.
就有信心,信靠和盼望的含義在當中.
而這個字的字根很有趣.
就是與「陷於泥沼」有關.
這個字的字根.
與「陷於泥沼」.
與「陷於困難」有關.
所以第二重的「等候」.
通常是指困難處境的等候.
就好像詩人正處於一個深處的狀態之中.
所以聖經幾個中文譯本.
將這個字譯作「仰望」.
其實是很有意思的 譯得很好.
等候者是抬頭注目在天上的.
那位滿有大能值得信靠的信實之主.
而不是單單低頭.
看著地上擺在他面前的那些難處.
這就是等候,仰望的意思.
詩人在困難之中等候的.
不是其他東西.

$^{321}$而是神的話.
詩人向上仰望,期盼神的話臨到.
神的話包括了神的安慰.
讓到詩人的心得到堅固和激勵.
神的話也包括了神的指引.
就是幫助詩人思考.
在當前的困境中要怎樣做.
神的話也充滿了神的應許.
詩人等待信息的神.
用他的大能去實現神說過的一切應許.
這就是第二重等候所表達的重點.
接著到第六節.
經文是這樣說的.
「我的心等候主.
勝於守夜的等候天亮」.
這節經文如果原文直譯.
就可以直譯成這樣.
「我的生命向主勝於那等候天亮的」.
第三重等候就是將詩人.
和那些等候天亮的人作對比.
這些等候天亮的人.
可以是指那些在晚上看守城門的人.
不過更加可能是指.
在聖殿負責守夜的利微祭司.
這些利微祭司等待神光出現的那一刻.
就獻上當日的每日繁祭.
這兩種等候有一個共通點.
就是不能夠鬆懈.
要專注,要謹慎,要警醒.
「勝於」這個詞.
原文其實是一個介詞.
好像英文的preposition.
它可以解作「勝於」.
也可以用來表達在他們當中的意思.
如果在這裡解作「勝於」.
就表示詩人比守城或守聖殿的人.
更加專注,謹慎和警醒.
但如果採用後者的解釋.
其實就是要強調.
詩人是和其他人一起等候.

$^{361}$救贖黎明的來到.
這也顯明了.
詩人非常有群體意識.
稍後大家會看到.
後者的解釋.
其實和第四個篇章.
即是七至八節.
那裡銜接得會更加好.
總的來說.
三重等候絕對不是被動和消極地挨日子.
不會認為當下是沒有意義.
是毫無作為的.
這樣的三重等候.
只是要呼應在前一個段落.
結尾的時候.
他所強調的重點.
是生命要更新,轉化.
要活出敬畏神的人生.
這篇詩篇.
如果大家讀的時候.
可能會覺得在這裡也可以停止.
但詩人沒有在他做完等候宣告之後完結.
詩人最後向以色列發出呼籲.
詩人這樣說.
「以色列,你當仰望耶和華.
因耶和華有慈愛,有豐盛的救恩.
他必救贖以色列.
脫離一切的罪孽」.
這表明了詩人有一個更加清晰的覺醒.
知道他所持守的是一個群體的信仰.
是一個群體的信仰.
而不是單單追求自己個人的救贖.
詩人知道他不是孤單地在深處受困和受苦.
所以他要和他的同伴一起去求告神.
一起去依靠神的赦免之恩.
以至他們的生命,人生都能夠轉變.
他和他的群體裡的同伴一起等候.
將他們的生命結連在神那裡.
一起在等候之中.
去專注和仰望神的話語.

$^{401}$並且一起過一個謹慎和警醒等候的生活.
當詩人重提神有慈愛,有豐盛的救恩.
以及提到神必救贖以色列.
脫離一切的罪孽的時候.
在這位熟悉舊約聖經的詩人的心裡.
很可能正是聯想到.
《出埃及記》那段的金句.
就是《出埃及記》三十四節六至七節那裡所說.
耶和華,耶和華,有憐憫有慈慰的神.
不輕易發怒,且有豐盛的慈愛和信實.
為千代的人流傳慈愛,赦免罪孽,過犯和罪惡.
萬不以有罪的為無罪,必懲罰人的罪.
至父及子,直到三,四代.
這句話就是神在金牛毒事件之後.
與以色列民重新立約之前.
向以色列民自我啟示的說話.
詩人的呼籲就是提醒以色列的同胞.
他們是與神有生命之約的屬神群體.
而在這個群體裡.
大家會更加容易得到勉勵和慰養.
在這個群體裡.
他們能夠更加好地去實踐.
神要他們去承擔的使命.
就是他們要按著神的話語而生活.
要彼此激勵彼此扶持.
這樣去抗衡罪惡.
向群體以外的人去彰顯神的榮耀.
也就是說要讓那些未認識神的人.
認識到神在人生命之中的分量和重要性.
今天讀的這四個詩篇的篇章.
大家可以看到是一環扣著一環.
正好示範了《屬靈覺醒》進程的四部曲.
第一部曲首先是.
對自身正是陷於深處的覺醒.
也就是說要意識到.
自己的生命正在逐漸流逝.
這是《屬靈覺醒》的第一部.
但是人有這樣的覺醒並不是必然.
陷在罪裡面步向滅亡而不自知的人.
實在多的是.

$^{441}$他們有時候甚至會為了合理化自己的惡行.
讓自己能夠繼續享受那個最終之樂.
於是就將判斷是非對錯.
很壞,善惡的真理和價值觀都扭曲了.
所以作為信徒就要經常去審查.
自己生命的境況究竟是怎樣.
是否在某些方面逐漸遠離神.
而於是乎逐漸走近深處的地方呢?.
甚至已經有跌入深處的危機.
第二部就是要覺醒到.
自己之所以陷於深處的終極原因.
是因為自己的不信.
而與那位創造人.
並且次下生命氣息的萬有主宰隔絕了.
這才是終極的原因.
在深處掙扎的人.
如果他們看不到.
這才是淪落到這個地步的終極原因的話.
就有可能被困在憤怒和自憐的情緒之中.
只會怨天尤人.
甚至去搞破壞.
作為信徒就要經常去審查.
神在自己的生命之中.
到底是遠還是近呢?.
例如.
相對於看手機來說.
自己有多少時間會安靜在神面前.
親近神進入一個深度的禱告.
去聆聽聖靈藉著神話語的提醒呢?.
第三部就是要有三重等候的覺醒.
知道要以積極的等候.
將自己的整個生命全方位與神建立生命的結連.
要注目於神話語的指引和教導.
並且要時刻謹慎和警醒.
以至對神的作為和提示.
可以有更加敏銳的觸覺.
這本身就是一個漫長的等候過程.
實際上也是信徒生命成長的進程.
最後就是群體意識的覺醒.
這也是信徒的屬靈意識.

$^{481}$由個人而走向群體的提升的進程.
一個永遠停留於只有個人意識.
而不願意為群體付出.
而放下自我的信徒.
他就很難真真正正地明白基督大愛的真諦.
在今年的大齋期裡.
讓我們從這四個方向去審查.
我們的屬靈生命的實況.
以上的關卡除了勾畫出.
屬靈覺醒進程的四部曲之外.
它也給我們一個重要的提示.
就是如何向現在的人傳福音.
要從什麼樣的地方入手.
最重要在當中.
最關鍵要去問的一個問題就是.
現在的福音對象.
到底是處於一個怎樣的深處處境之中呢?.
舉一些例子.
究竟他們是不是有家人.
有很多朋友.
但卻找不到可以傾訴的對象.
於是內心感覺到異常的孤單呢?.
或者是不是壓力無處不在.
非常沉重.
令人完全無法喘息呢?.
又或者是工作很忙碌.
營營役役.
停不下來.
但卻又找不到真正的人生意義呢?.
不知道做這麼多事是為了什麼.
又或者是活在焦慮之中.
缺乏安全感.
是因為找錯了生命的保障.
於是心靈始終無法安頓下來呢?.
還是被社會扭曲的價值所牽引.
於是陷在罪的折騰之中呢?.
願我們好好在祈禱之中觀察和思想.
明白福音對象的處境.
這樣大家就可以循著那個方向.
作為分享福音的起始點.

$^{521}$去觸動福音對象的心靈.
認識他們生命的真相.
去引導他們尋求神.
讓基督的福音改變他們的生命.
願道上主賜福給大家.
阿們.
\newpage



\section{詩篇 1:1-3}
\label{sec:031844}
\textbf{建立健康的使命教會 (9) 愛慕遵從神話語}~~~~林耀楠牧師
\newline
\newline
連結: \href{http://www.church.com.hk/acms/content.asp?site=cdc\&op=show\&type=product\&code=031844\&layout=sermon}{\texttt{031844}} ~~~~ 語音日期: 2026-03-22
\newline
\newline
\hyperref[sec:031843]{\small{< < < PREV SERMON < < <}}
~
\hyperref[sec:index_chronic]{\small{[返順時目]}}
~
\hyperref[sec:031845]{\small{> > > NEXT SERMON > > >}}
\newline
\newline
詩篇 1:1-3
\newline
\begin{longtable}{cl}
\hline
\hline
章節 & 經文 (和合本修訂版)\\
\hline
1:1 & \begin{tabularx}{0.7\textwidth}{X} 不從惡人的計謀，不站罪人的道路，不坐傲慢人的座位， \end{tabularx} \\ \\ \relax
1:2 & \begin{tabularx}{0.7\textwidth}{X} 惟喜愛耶和華的律法，晝夜思想他的律法；這人便為有福！ \end{tabularx} \\ \\ \relax
1:3 & \begin{tabularx}{0.7\textwidth}{X} 他要像一棵樹栽在溪水旁，按時候結果子，葉子也不枯乾。凡他所做的盡都順利。 \end{tabularx} \\ \\
[1ex]
\hline
\hline
\end{longtable}
$^{1}$弟兄姊妹朋友早晨.
問一個問題.
「富有」和「有福」的分別是什麼?.
「富有」和「有福」有什麼分別?.
有什麼異同?.
那麼大聲笑的那個先說.
「富有」和「有福」都有一個「有」字.
都有些東西.
除了「有」字之外還有什麼相同的地方?.
都是一個「口錢」.
看到「一口錢」嗎?.
我以前看這個「一口錢」就以為.
「一口錢」就是「一口錢」就是「一幅田」.
古時.
講財產的當然是「一幅田」.
可以耕種可以租給別人就很好了.
「一幅田」就很好了.
我以為這個「富」字和「福」字.
「一口錢」就代表著「富有」.
但是這次我嘗試上網看看.
網上的中文老師有沒有新的教導.
有新的提示.
原來「一口錢」這個字是讀「福」字.
我不知道原來這個字本身是「福」字.
「福」字是什麼意思呢?原來不是「一口錢」.
根據文字上的進展.
推翻轉頭看「甲骨文」裡面.
這幅圖畫代表著什麼呢?.
這幅字原來是一個器皿.
器皿上面畫的是一個蓋子.
中間那兩桶是一條繩子.
下面是一個器皿裝著一些容器.
原來不是「一口錢」,不是「一幅田」,不是「一幅地」.
是一個裝著可以裝酒的器皿.
「福」字原來是裝滿了的意思.
原來「福」和「富」這兩個字.
共同的地方都有些東西.
就是一個器皿裝滿了一些酒在裡面.
是滿日的意思.
是豐滿的意思.

$^{41}$是滿載的意思.
是原來很豐滿,裝滿了東西的意思.
這樣看這兩個字就有些不同的心情.
「富」代表什麼呢?.
財物上的富有.
富有代表著可能很多金錢.
很多地產.
很多財物.
有些有什麼好處呢?自然有富有的人生活就比較舒適.
有很多機會去發展.
但未必帶來幸福.
富有不等於有幸福.
富有也不等於有健康.
富有也不等於有安心過日子.
富有也不等於有教育好分.
那「有福」又是怎麼一回事呢?.
「有福」就不是跟物質上有關的.
「有福」就代表心靈上的滿足.
可能心裡面比較平靜.
遇到困難也有盼望.
人際關係也有意思.
是非物質上的福氣.
我們就用詩篇一章一至三至幾節經文.
給我們一些生活上的智慧.
詩篇和章篇都是類似的.
是智慧的文字.
不是一些神的應許.
你這樣做必然是致富或者自圓得福.
不過給我們一些生活智慧在念人生裡面.
給我們一些啟迪.
今天經文分開三段來看.
如何可以擁有富和福的途徑.
如何辨識富和福的選擇.
如何得到富和福的結果.
第一個擁有富和福的途徑.
有什麼致富之道呢?.
雖然我們這裡不是每個人都有錢人.
但是我們這裡每個人都有一些錢.
有些錢財 有些物資.
致富之道是什麼呢?.

$^{81}$簡單來說就是很多東西.
賺錢.
努力去找工作.
不懂得做 賺得不夠就省錢.
有智慧的省錢做節儉的生活.
都是一個致富之道.
或者做生意.
做買賣.
或者做投資 都是致富之道.
甚至做賭博 作弊.
都是致富之道.
有些人希望一炮就可以.
在賭場裡面賺到多少錢.
贏到多少錢回來.
或者做生意 作弊.
可以能夠得到更多的金錢收入.
都是致富之道之一.
甚至去偷 去搶.
原來都是致富之道.
富本身沒有對錯.
富有並不是一個罪過.
可以是一個祝福.
但是致富之道可以分對錯.
福又怎麼樣呢?.
得福之道又有什麼呢?.
福字有一個字字邊.
跟神有關係.
可能我們聽天由命.
很被動的 很無奈地來到這個日子.
等神 等上天來給我們一點福.
或者我們主動地去沾掛算命.
去尋求自己的將來.
為自己作打算.
或者我們求神拜佛.
求神拜佛當然是.
如果你找到一個真神來求.
那就有反應.
假神 魔鬼 撒旦.
都可能對你的所求有反應.
命運自主.

$^{121}$可能我們什麼都不求.
靠自己 掌管自己的命運.
靠自己可能得到蒙福.
自己給自己有福的途徑.
富和福本身沒有對錯.
是一件好事.
是一件福氣.
帶著神的恩典在裡面.
但是致富之道得福的途徑.
就可以有對錯之分.
今天看的經文是一至三節.
但在第一章最後一節的經文是這樣說.
耶和華知道義人的道路.
也知道惡人的道路.
義人的道路有時候可能很崎嶇.
但是耶和華知道.
神知道 神同樣知道惡人的道路.
我們有時候看身邊的惡人.
可能很順境.
很得意.
很風俗.
但是神知道.
惡人的道路特別滅亡.
富和福.
都是我們期望的.
但是神知道.
什麼是義人的道路.
什麼是惡人的道路.
現在再問你一個問題.
富和福.
你選擇哪一樣.
你會選擇什麼呢.
給你五秒鐘想.
富有和有福.
你要選擇哪一樣呢.
選擇富有的舉手.
難怪見堂時還不夠錢.
選擇有福的舉手.
舉手不要動.
你們心裡其實也想富有.

$^{161}$其實兩樣都想要有.
不是很貪.
最好是兩樣都有.
不過你的手代表一件事.
如果只能夠選擇一件事.
就是希望有富足.
還是有福呢.
我們選擇有福在先.
如果再加上富有就更好.
這是我們辨識富和福的選擇.
我們怎樣去選擇呢.
今日的經文有些像.
八福裡面的.
如述八福有福之人.
什麼是有福之人.
經文裡面的聖經比較突顯.
一開始是說有福的人是怎樣的呢.
其實他是用一個詩篇來描述.
一些有福的人.
他會有什麼的徵兆.
或者有什麼的表達.
有什麼的行徑.
我們一起看看.
這些有福的人有什麼行徑.
可以讓我們一起來學習.
聖經裡面一開始有三個字.
一字就說了.
不從不站不坐.
原來有福之人的表達.
就是他不做這個不做那個.
很多事是不做的.
除了不跟從.
不站立不坐.
甚至不走不跳不看.
不聽不做不說.
很多是不去做.
為什麼要不從.
不站不坐.
不從惡人的計謀.
有福的人.

$^{201}$他有一個特徵.
就是他不跟從那些惡人的計謀.
計謀就有好壞之分.
惡人的計謀就說明是惡人的計謀.
當然不是很好.
沒有人明知是惡人的計謀都會去選擇.
我相信不用問阿貴.
再再叫他舉手都不會有人選擇.
明知是惡人的計謀.
怎會去選擇.
但是就會選擇一些計謀.
人很奇怪.
很喜歡一些.
很袖手.
很短暫.
很小火花就得到.
利益.
很容易做得到.
那些計謀.
我相信現在報稅時期.
三月份快要到.
三四月都是報稅的時候.
報稅的時候.
很多計謀在裡面.
為什麼你報稅.
又合規格又可以省稅.
但是有些是惡人的計謀.
那些又不是叫.
很好的報稅.
那些叫騙稅.
為什麼來到稅項裡面.
騙到政府.
我們說政府已經每年抽了我們這麼多稅.
抽少一點都不要緊.
好像天公地道.
那些好像很大件事.
不知道你們有沒有試過在公司工作的時候.
將公司的物資拿回家用.
那些筆.
那些.

$^{241}$寶.
紙張.
有沒有順手拿回家用.
沒有人會特意去拿回家用.
但是想一想.
再想一想.
拿回去用有什麼問題.
反正公司賺了這麼多錢.
拿回去用有什麼問題.
對我們這麼刻薄.
又沒有什麼福利.
當然其中一個福利.
拿回去用有什麼問題.
原來一些本來中性的計謀.
想一想就變成.
本來是惡人的計謀都變成合理的計謀.
人很喜歡合理化.
我們所做的事.
找藉口找理由去打圓場.
去說服自己.
我這樣做.
也不是很過分.
我旁邊的同事拿更多.
反正大家都拿了.
公司給的福利就算了.
他不問我不說.
也沒有人知道.
一個中性的計謀.
原來是一個惡人的計謀.
可能不是別人.
教我們的.
是我們自己裡面的貪念.
湧出來.
不從惡人的計謀.
我們.
不要將這些計謀合理化.
不站.
罪人的道路.
明知罪人的道路都會走.
沒有人會這樣做.

$^{281}$偷偷拐騙.
去吸毒.
去犯法.
沒有人會故意去做.
在1983年.
移民之後的第二年.
有同事.
請我回家吃飯.
聊天.
吃完飯聊天的時候.
他拿了一支煙出來.
那支當然不是煙.
裡面有.
所謂有鹹的.
吃了會很興奮很開心.
讓我去試一下.
當時.
因為未合法.
也不敢去試.
拒絕了他.
他吃完之後真的很興奮.
現在合法了.
如果再有同事叫我去吃的時候.
我應該有什麼反應呢.
合法的.
同事招呼到.
請你試一下.
不會犯法的.
沒有人抓你.
沒有人告你.
吃一次這麼多.
沒有壞身體.
又好像不會上癮.
站罪人的道路.
我們會不會聽從別人的邀請.
站上他們的道路呢.
有時候我們自己騙自己.
我記得小時候.
有一次同學想去游泳.
爸爸不讓我去游泳.

$^{321}$因為我根本不懂游泳.
那時候是初中的時候.
爸爸不讓我去游泳.
說我一定會去.
同學教我一個方法.
你拿一條泳褲出街.
不要穿出來.
回到學校一起去放學就去游泳.
游泳你就靜靜的回家洗乾淨.
藏起來沒有人知道.
爸爸也不知道.
我說是喔.
靜靜做沒有人知道.
照做 拿泳褲游泳.
游泳回家洗乾淨.
結果怎樣.
香港以前我們沒有洗衣機.
拿條泳褲出來怎會不知道.
不知道他蠢還是我蠢.
我就是一個蠢人.
聽惡人的意見.
跟著惡人的道路.
游泳開心之後就說謊.
告訴爸爸做功課.
結果就自己.
做回自己.
承受自己的惡果.
吃了一餐蚊豬肉.
我們不會主動去參與.
別人那些罪人的道路.
但我們心裡面那種罪性.
會容許我們自己去做.
這樣的事情.
以為可以瞞天過海.
以為天上的神都不知道.
爸爸都知道天上的神怎會不知道.
說一次謊話.
一次半次無所謂.
說多十次八次就有所謂.
習以成常.

$^{361}$就答應有所謂.
不讚惡人的道路.
這是一個傲慢人的座位.
我有一次機會.
跟兩個弟兄聊天.
分開來的.
兩個弟兄都很成功.
工作上都很成功.
很聰明.
我問其中一個.
你是不是一個謙卑的人.
他想了想.
他說.
我想我都算是一個謙卑的人.
我問他為什麼.
他說因為我不喜歡和別人比較.
又不敢和別人比較.
所以他說我都算是一個謙卑的人.
我覺得他也是.
其實我覺得他也是很聰明.
很聰明.
他謙卑地說.
不承認自己是一個聰明人.
他只是說自己是一個謙卑的人.
跟另一個弟兄聊天.
我問他同一個問題.
你是不是一個謙卑的人.
他說.
我是一個十分謙卑的人.
當時我就停一停.
他是一個十分謙卑的人.
我問他.
為什麼你覺得自己是一個謙卑.
甚至是十分謙卑的人.
他說我不在別人面前認聰明.
但身邊實在太多的蠢人.
很多時候我們都以為自己很聰明.
我不知道你自己覺得自己是不是很聰明.
我承認.
我自己覺得自己也很聰明.

$^{401}$會不會聰明到一個地步.
覺得自己成為一個傲慢的人.
自己也不知道.
自己也不知道自己是一個傲慢的人.
我是一個十分謙卑的人.
和說我是一個謙卑的人.
是相當之不同.
我是一個謙卑的人.
可能真是一個謙卑的人.
但是當一個十分謙卑的人.
我相信謙卑極.
也有限.
不從惡人的計謀.
不站罪人的道路.
不坐傲慢人的座位.
好像一些很容易.
去達成的行動.
但其實得來一點也不容易.
和我們自己的內心世界很有關係.
不要怪責身邊的人.
不要怪別人拖你下水.
不要怪別人開一條不好的路.
帶你走在你自己.
我們自己心裡面的.
邪惡感.
犯罪感.
犯罪傾向.
推動我們去跟從.
去學習.
去模仿.
為了自己要得到好處.
原來那些有福的人.
除了工作.
有些事情是他特別去做的.
唯有喜愛耶和華的律法.
就夜思想他的律法的人.
那些有福的人原來對耶和華的律法特別有興趣.
是很喜愛的.
你們有沒有人喜歡律法的?.
有啊.

$^{441}$很難得啊.
劉牧師.
這隻手非常之矜貴.
很多人都不喜歡被人管的.
喜愛耶和華的恩典.
喜愛耶和華的仁愛.
喜愛耶和華的饒恕.
喜愛耶和華的憐憫.
我們全部都喜歡.
不單喜歡耶和華所有的好處.
我們全部人都喜歡耶和華所有的好處.
但有多少個像劉牧師那樣喜愛耶和華的律法呢?.
早幾個月前.
我們都經歷過這個學習和體驗.
幾個月前.
大家還記不記得.
在台上分享過.
收那些超速.
那些告票.
收到幾張.
有多少個收到的?.
有沒有人聽見說.
他以後不再出告票之後很開心.
甚至開心到睡不著.
沒有人喜歡被人管的.
沒有人喜歡整天限速40就40.
每一公里都不行.
一定要40.
每次走過那40個牌.
看到那些攝影機.
我都心裡想有沒有搞錯.
騙錢的.
幾個月之後.
終於下了那些道令.
不再出那些告票之後.
走過那些又牌又攝影機的時候.
你猜一下我走的速度是多少?.
有多少人猜我走40個的時候?.
都有幾個勇敢.
有誰猜我走超過40個的時候?.

$^{481}$這是你,不是我.
我竟然是很努力地.
繼續保持速度在40以下.
剛才超過40.
簡直超過40不多.
為什麼這麼奇怪?.
沒有人管,沒有攝影機.
拍了照都不用出告票.
為什麼仍然是這樣?.
我就開始欣賞那些攝影機.
原來它改了我的習慣.
讓我知道原來小心駕駛是一件好事.
原來是一件好事.
就算明明是律法規限了我去做.
初初不甘心去做.
而變成一個樂意去奉行.
這個律法的時候.
是一種轉變.
是一個成長.
可能我們初初都沒有人喜歡柔和的律法.
這麼嚴厲,這麼多的要求.
好像這麼多的限制.
但是當你認受過.
認識過的時候.
就開始去愛上.
柔和的律法.
我有一個小時候的經歷.
可能之前也分享過.
但是我覺得很有意思.
年幼的時候.
我們現在有攝影機.
年幼的時候我們三兄弟.
十歲八歲.
有一次爸爸說帶我們去看電影.
問我們選什麼電影看.
那個年代很喜歡戰爭片.
我就忘了我喜歡選哪一套.
總之一套就是告白.
上面很多士兵.
拿著槍,炮那些.

$^{521}$戰爭片.
我們說我們選戰爭片去看.
爸爸說不好.
我帶你看另一套.
是哪一套呢.
鮮落飄飄處處民.
我看那個廣告.
小朋友.
那個女人.
在那裡跳舞,唱歌.
我們在馬路上哭了.
爸爸拖了我們幾個馬路.
我們就抓著馬路的欄杆.
我們要看戰爭片.
不要看唱歌跳舞那些.
爸爸說不行.
爸爸說是.
去看電影.
看完鮮落飄飄處處民.
我們都飄啊飄啊.
很開心.
之後還聽不聽爸爸的話.
聽啊.
開始是越來越聽爸爸的話.
爸爸為我們預備的是最好的.
不一定是我們心所想的.
但對我們是最好的.
最有益的.
飄啊飄啊就一生就不飄了.
很多時候.
我們要學習喜愛耶和華的律法.
可能要經歷一些事情.
懂得耶和華的律法有多可愛.
我們就開始.
去喜愛它.
開始去欣賞它.
第二部分更難.
就夜思想它的律法.
我只是問你一句.
看看有沒有手舉.

$^{561}$有多少個就夜思想.
耶和華的律法.
牧師先舉手不用說.
兩位牧師都先舉手.
就夜思想.
耶和華的律法.
剛才說了.
律法的好處.
為什麼詩人說要.
就夜思想.
大家都很明白.
不是就夜思想哪裡記得.
思夜想要記得.
很精彩.
幾天前.
星期四我們有金靈社.
知不知道金靈社.
不知道順便告白.
55歲或以上退休人士.
可以歡迎星期四早上回教會樓下.
旁邊的運動場.
我們有金靈社聚會.
9點半至12點.
非常活躍有遊戲有跳舞有唱歌.
有運動.
有24式太極.
沒打太極一段時間.
兩三年前和太太一起.
跟一個師傅學了太極大半年的時間.
24式都學會了.
記住了.
但一段很長的時間都沒打.
星期四比較空閒.
在健身室.
和弟兄姊妹朋友.
一起打了24式太極.
我竟然不打了幾個月.
都還記得怎麼打.
還覺得自己打得很好.
開始覺得.

$^{601}$心滿意足.
欣賞自己.
原來自己學習了24式之後.
沒打一段時間都還記得怎麼打.
開始覺得自己很厲害.
我都很聰明.
但在思想上.
應該學習走手.
竟然和這段經歷有關係.
你們中間背金句很厲害的表手.
當然沒人敢說.
無恥都不敢.
不然會被人考.
我們從小就開始背金句.
背了很多年金句.
背金句背到一個地步.
別人說上文我們會下理.
別人說上斂我們會下斂.
講半句你會講下半句.
背金句背到一個地步.
就好像我打太極.
我回家再想想.
我打太極雖然有板有眼.
但沒功架.
沒實力.
有姿勢.
沒實際.
想一想就像我們背金句.
背金句背到一個地步.
滾家爛熟.
都會背.
偏偏就不記得不明白.
不懂.
然後突然知道.
正經金句所講的是什麼.
背金句厲害不等於.
我們能行出金句所教導的東西.
愛鄰舍猶如愛自己.
誰不會說.
有幾個能愛得出來.

$^{641}$恐怕我們背金句.
背到厲害到.
自己都以為自己很熟悉.
耶和華的律法.
很熟悉神的教導.
跟自己很熟悉.
二十四式太極一樣.
只有姿勢.
恐不實際.
起碼我們對神的律法.
是晝夜思想.
不單止是幾多年前背的金句.
是背了之後.
還日夜思想.
不只是兩三年前學了太極.
覺得自己懂了厲害了.
還要日日練習才有用.
日日練習.
金雞獨立才能站穩.
我們現在金雞獨立不了.
就因為沒有好好練習.
我們即使以為自己熟悉.
神的話語.
如果不是日夜思想的話.
沒有真正入心入肺的話.
沒有真正可以很自然地.
運用出來.
自然反映出來的時候.
文字對我們都是表面上的認知.
對我們生命.
其實沒有實際的效用.
楊華的律法.
引導我們向善.
向神.
要渡過人生路.
要我們晝夜的思想.
不要輕視.
不要以為自己已經得著了.
仍然需要.
晝夜來思想.

$^{681}$好讓我們能夠自然地反映出來.
不需要背的金句.
對我們的行為.
對我們的存在.
不需要背的金句.
對我們的行為.
在我們生命裡面.
活出神的話語.
這人便為有福.
原來以為有福的人.
就是這樣的描述.
斯人容易記的文字.
他不做這樣不做那樣.
不做那樣.
不讓心裡邪情私慾.
犯罪的意念湧出來.
唯有喜愛耶和華的律法.
晝夜思想的時候.
這人便為有福.
到這個時候.
我都問自己.
我是不是一個有福的人.
可能你都需要去問一問你自己.
你是不是一個有福的人.
在你自己眼中你自己是不是一個有福的人.
同樣地問自己.
我是不是一個有福的人.
有沒有剛才所描述的情況出來.
我不是一個大富大貴的人.
但神有恩典.
讓我不用憂財.
不用憂迷.
我覺得我是一個有福的人.
神沒有讓我很富有.
但神要讓我很足夠.
日子覺得很滿足.
因為滿足日子.
這個就是一個福分.
你問我有沒有一個漂亮的車可以開呢.
十四年前我有一個漂亮的車可以開.

$^{721}$因為我現在的車.
是十四年前已經出車了.
由十四年前一輛漂亮的Toyota.
用了三萬多元買回來的車.
由漂亮的車開到今天.
都還是一輛漂亮的車.
雖然有點生鏽.
但神的恩典在裡面.
讓我十幾年來一輛車都很好.
正常的保養維修.
就已經可以用了十幾年.
最嚴重的一次.
就是最近一次.
終於那個startup壞了.
但在停車場裡面壞.
不是在馬路裡面壞.
又可以有感恩地.
有親戚幫我們把車拖回家.
又可以很快的找到人.
找到弟兄來我家幫我fix.
又可以用多幾年.
十四年後的今天.
現在仍然是一輛漂亮的車.
今天走了三十八萬公里.
我們接下來的目標.
希望要過四十萬.
我最近和太太.
有些商議.
我們改一個目標.
希望我們的目標.
去到四十五萬.
為什麼是這樣呢.
因為可以買少一輛車.
不用那麼快買車.
應該說始終都要買車.
不用那麼快買車.
延遲幾年.
怎麼用有什麼好處.
有什麼好處.
你和我一樣心思.

$^{761}$可以放在肩膀那裡.
如果我們每個人都這樣想.
就真是不錯.
那個提案是說.
我回來延遲你買車買樓.
買股票買遊船行.
什麼都好.
延遲少少可能幫助大家很大.
我回來這裡.
得著負和福的結果.
我們有沒有那些象徵.
我是不是不用病.
不是.
原來我和大家一樣.
都有病有痛.
但病痛裡面仍然有盼望.
仍然有覺好睡.
有膽好食.
這就算是福.
就是說神恩典足夠.
我有沒有一個很快樂的家庭.
我可以說有.
和太太每天在日子裡.
都嘻嘻哈哈傻傻耿耿.
那些開心快樂的日子.
當然我們開心快樂.
那些嘻嘻哈哈傻耿耿.
不可以跟你說那麼多.
但希望我們在平常日子裡面.
都感受到神的恩典在裡面.
讓我們有一個.
很和平的.
很友愛的.
很溫馨的.
神的祝福就在裡面.
有兒有孫.
為什麼會不是福呢.
最大的福是什麼呢.
就是在任何情況裡面.
都凝聚了.

$^{801}$完全是神的恩典.
那你問我像什麼呢.
我覺得像一棵樹.
詩人說.
有福的人就像一棵樹.
栽種在溪水旁的時候結果子.
我所有的傳言都是神所給的.
就好像一棵樹.
栽種在水旁一樣.
完全是水江河.
神的水.
在我們生命裡面流.
讓我們得到滋潤.
得到滿足.
得到快樂.
按時結很多的果子.
我們每一個都可以.
按時結很多的果子.
結福音的果子.
結和平的果子.
結人類的果子.
結恩慈的果子.
結喜樂的果子.
結溫柔的果子.
我們每天都可以按時的來結.
這個全是神的恩典.
月至也不枯乾.
凡他所做的.
有福的人所做的.
盡到順利.
盡到順利其實一點都不容易去理解.
前面上我們覺得我們做的順利.
問題是.
順誰的利呢.
是不是順著我們自己的心意去成就.
為之順利呢.
是不是順著神的心意去成就.
為之順利呢.
盡到順利.
這個字.

$^{841}$prosperity.
在美國的教會很流行.
我們叫他們做成功神學.
有沒有人聽過什麼叫成功神學.
我第一次特意上去網站.
聽他們一個大師講一些.
一堂去讀.
那一堂讀英文叫.
Over and Above.
英文叫.
超過,超越.
你所能夠做的.
聽了一輪之後發覺很有趣.
不停地在重複.
你付出捐獻更多.
神就給更多給你.
如果你不是旁邊那個.
為什麼蒙神賜福給你這麼多.
這麼順利呢.
因為捐獻得,奉獻得比你更多.
所以如果你要得到更多的福分.
更多的財物.
更多的健康.
更多的家庭順利.
你就要捐獻得更加多.
不斷重複這個訊息很多次.
你想多得神的祝福.
你就要多去捐獻.
捐得越多得的越多.
這個是不是聖經所教的呢.
我忽然間在聽的時候.
他講的道理好像什麼呢.
好像我看了一套古裝片.
古裝片裡面.
有個人在街邊.
叫人們嘗試一下.
他這個聚寶盤.
他說這個聚寶盤很厲害的.
你放錢進去會生錢出來.
你試一下.

$^{881}$放一個一文錢進去.
突然間變成兩文錢.
很成功.
第二個說讓我試一下.
放一兩銀進去.
他就放一兩銀進去.
那個人就去佛部養養他.
一兩銀變兩兩銀.
第三個說讓我放十兩銀.
這個是妹妹.
放十兩銀進去.
佛部養一養就變十兩銀.
多了十兩銀出來.
然後財主經過.
放一百兩銀進去.
佛部養一養.
接著怎樣.
不見了.
那個人就說.
祖師爺顯靈.
你的不義之財.
貪心所以你的錢.
祖師爺收了沒有了.
當我聽到那個人說.
捐獻多些就得多得.
神的祝福就有這種感覺.
多獻多捐多得.
都不知道是從哪裡來的.
剛才所有的經文就說.
為愛.
耶和華律法.
晝夜思想.
他的律法.
這人便為有福.
他所做的.
盡到順利.
包括神可以祝福給你.
也都同樣包括.
神可以令你有富足.
像所羅門那樣.

$^{921}$神問他要一樣東西.
你要什麼.
他沒有要求富足.
沒有要求很多財物.
他要求神給他智慧.
神就賞賜給他.
再加添他沒有要求的.
富足.
財物.
真正的盡到順利是什麼意思.
就是在疾病困難裡.
仍然有一份喜樂.
仍然有一份盼望.
仍然有一份好睡.
在金融風暴裡.
在財物.
財務方面有困難的時候.
仍然心裡有份平靜.
有份安穩.
最近我聽一個分享故事.
有一個姐妹說.
她是一個會計師.
曾幾何時在公司做工.
公司要她做一些數字.
可以隱瞞政府一些稅項.
幫公司省下很多錢.
她有兩個選擇.
繼續按照老闆所吩咐的來做.
另一個就是按照她自己的良心.
神所給她的良心.
來拒絕做這件事.
她拒絕做這件事.
結果被老闆解僱.
有一段時間都沒有工作.
還在找工作.
她心裡有沒有平靜安穩.
她有.
她覺得自己在這件事上.
站得穩.
雖然失去工作.

$^{961}$心裡仍然有一份平靜.
和安穩.
信神一定有預備.
這就是她的福.
她就是一個有福之人.
她就在找神的話語.
神的律法 神的要求.
走在神的路上.
不從惡人的計謀.
不聽老闆說.
不站惡人的道路.
不走上那條歪曲的路.
不坐司馬人或奧馬人的座位.
有原則地.
為神站得穩.
這就是有福之人.
這就是我們希望學習的地方.
知乎之道.
蒙福之路.
我們要記得.
剛才說.
第一篇詩篇最後一篇經文.
又說是知道義人的道路.
也知道惡人的道路.
但惡人的道路卻必滅亡.
我們願意走上這條.
登天之路.
走向神的路.
送給大家.
一同學習的.
業餘.
我們一同學習.
愛慕和尊重神的話語.
我們一同成為.
一個蒙福的群體.
一同建立一個健康的教會.
靈裡面健康的教會.
屬靈事實上是一個健康的教會.
心靈上.
每一個都有健康的心靈.

$^{1001}$一同有一個健康的教會.
一同同心祈禱.
多謝你創造我們每一個.
多謝你讓我們有一個.
自由選擇.
這個福分.
這個元素.
你讓我們選擇.
去愛你.
你讓我們選擇去接受你的恩典.
你讓我們選擇去.
跟你有一個美好的關係.
你讓我們選擇去.
聽從你的律法.
你的教導.
你讓我們蒙福.
蒙恩的過一生.
一切都是從你而來.
求你讓我們有一個醒覺的心.
有一個醒覺的靈.
懂識.
何謂善.
何謂惡.
求你的話語傳來在心裡面.
讓我們可以晝夜思想.
能夠很自然地.
很樂意地.
很甘心地.
活出與你與我們活出的樣式.
往常而立.
求神的靈帶給我們.
加力.
加恩加福給我們.
多謝你聽我們祈禱.
和主耶穌基督得勝的命救.
Amen.
\newpage



\section{創世記 1:1-3}
\label{sec:031845}
\textbf{恢復生命的榮美 (華)}~~~~卓爾君牧師
\newline
\newline
連結: \href{http://www.church.com.hk/acms/content.asp?site=cdc\&op=show\&type=product\&code=031845\&layout=sermon}{\texttt{031845}} ~~~~ 語音日期: 2026-03-22
\newline
\newline
\hyperref[sec:031844]{\small{< < < PREV SERMON < < <}}
~
\hyperref[sec:index_chronic]{\small{[返順時目]}}
~
\hyperref[sec:031846]{\small{> > > NEXT SERMON > > >}}
\newline
\newline
創世記 1:1-3
\newline
\begin{longtable}{cl}
\hline
\hline
章節 & 經文 (和合本修訂版)\\
\hline
1:1 & \begin{tabularx}{0.7\textwidth}{X} 起初，神創造天地。 \end{tabularx} \\ \\ \relax
1:2 & \begin{tabularx}{0.7\textwidth}{X} 地是空虛混沌，深淵上面一片黑暗；神的靈運行在水面上。 \end{tabularx} \\ \\ \relax
1:3 & \begin{tabularx}{0.7\textwidth}{X} 神說：「要有光」，就有了光。 \end{tabularx} \\ \\
[1ex]
\hline
\hline
\end{longtable}
$^{1}$弟兄妹平安.
我想這幾週人最關注的就是.
美伊戰爭的消息.
我們看見飛彈亂飛.
油價亂漲.
然後心裡也跟著股市七上八下.
老實說現在的世界.
就像是隨時會炸開的火藥庫.
而許多人的心裡.
也像是剛被轟炸過一樣.
心裡充滿不安.
當然這些不安.
不一定是來自於遠方的戰爭.
而更多是來自於我們生活的內戰.
長期的工作壓力.
財務壓力.
或是那個碎了一地的人際關係.
但我們常常裝著外表沒事.
可是內心啊.
可能早已像家裡那個儲藏室一樣.
門關上.
看起來乾乾淨淨.
門一打開.
哇 裡面是另一個世界.
三年前買的健身器材.
現在變成高級的塞衣架.
還有一堆不知道屬於哪台機器的充電線.
糾纏得像一盤炒麵一樣.
弟兄姊妹我們的內心.
是否也像那一間儲藏室.
亂到自己不敢打開.
如果是的話請不要灰心.
因為其實我們所面對的混亂.
在聖經的第一頁就出現過.
我們常以為創世記第一章.
是關於恐龍化石.
大爆炸理論.
這種科學的辯論.
這些就讓神學家科學家去傷腦筋.
今天我要把焦點從宇宙.

$^{41}$拉回我們的內心.
因為創世記不僅是講世界的起源.
更是當我們在面對失控的生活時候.
創世記要幫助我們恢復生命.
那個榮美的秩序.
那麼今天當我們進入混亂的季節.
創世記第一章教導我們該如何行.
有三點.
第一個用宣告取代抱怨.
跟你旁邊講用宣告取代抱怨.
創世記一章一到三節.
祈出神創造天地 地是空虛混沌 淵面黑暗.
聖經一開始就告訴我們.
世界出現之前.
宇宙的狀態空虛混沌 淵面黑暗.
空虛混沌這四個字在希伯來文裡面.
叫做頭乎肉無乎.
大家可以念念看.
其實聽起來好像豆腐碎成一粒的樣子.
頭乎肉無乎.
這個就是世界的開始.
沒有形狀 沒有秩序.
伸手不見五指的黑暗.
聽起來很遙遠嗎.
其實不是.
這常常就是我們真實生活的寫照不是嗎.
就像你明天早上進了辦公室.
端著咖啡打開郵件的那一瞬間.
未讀郵件三十封.
老闆在群裡面拼命的標記你.
但是你的電腦卻在這關鍵時刻當機了.
那一刻你的處境就是頭乎肉無乎.
又或者你看見健檢的報告.
跟你關心的股票都是滿江紅.
這個也是頭乎肉無乎.
當你在夜深人靜的時候.
突然覺得自己不知為何而忙.
這叫空虛.
當事情多到讓你崩潰.
這是混沌.

$^{81}$當你看不到未來的方向.
這是黑暗.
那麼今天當人面對生活當中的空虛混沌黑暗.
人常有什麼反應呢.
最普遍的反應就是抱怨.
因為抱怨會給人一種虛假的安慰.
只要我把錯歸給環境 推給別人.
推給我的父母 推給運氣.
我心裡就好過一點.
但事實上抱怨不但無法解決混亂.
而且會擴大黑暗.
請留意一條屬靈法則.
那就是話語都帶著能力.
神的話帶有創造的能力.
而抱怨就帶有反創造的能力.
當人不斷的抱怨.
就是不斷在擴大他生命當中的黑洞.
直到那個黑洞吞噬掉所有人的視線.
比如你的團隊正在推動一個很重要的專案.
就有關鍵時刻有人產生失誤.
這個時候如果你選擇抱怨.
為什麼這點事情都做不好.
結果會怎麼樣.
抱怨會像病毒一樣的擴散.
瞬間就變成互相指責的批鬥大會.
士氣瞬間瓦解 溝通全面癱瘓.
原本的失誤會因為抱怨.
而被放大成一場巨大的空虛混沌.
所以創造天地的神.
當他面對宇宙級的空虛混沌.
他沒有抱怨.
他也沒有很焦慮大喊.
怎麼這麼黑啊.
誰把總開康關掉了.
神做了一件最簡單但最具震撼力的事情.
在創世記一章三節說.
神說要有光就有的光.
請注意 神不是在心裡想要有光.
祂是開口宣告要有光.
這個是一個主動的介入.

$^{121}$神拒絕讓混亂變成最終的結局.
神要用真理的光.
在混亂當中切出一條光明的路.
因為光才是帶來秩序的關鍵.
就像舊約的約瑟.
他是父親最愛的小兒子.
還穿著父親給他的限量版彩衣.
但是在一天之內.
他的人生就遭遇史詩級的空虛混沌.
他被他的親哥哥扒光衣服丟到坑裡面.
當作特價品賣到埃及當奴隸.
然後他好不容易在老闆家拼出一點業績.
又被老闆娘誣陷性騷擾.
直接就被打入大牢.
如果你是約瑟有沒有理由崩潰啊.
有沒有理由抱怨上帝不公啊.
當然有啊.
但是聖經記載約瑟在監獄裡做什麼.
聖經記載他在監獄裡管理所有囚犯.
後來法老做惡夢.
埃及陷入國家級的空虛混沌.
約瑟被叫出來.
他沒有抱怨過去的遭遇.
他條理分明宣告未來14年的國家糧食管理計畫.
你們發現了嗎弟兄姊妹.
約瑟走到哪裡.
就讓那裡的混亂變成秩序.
在主人家他把混亂的家務理順.
在監獄裡面他把凶惡的犯人理順.
在埃及的皇宮他把國家糧食理順.
為什麼.
因為他心裡有神的光.
他知道神與他同在.
所以他能夠在黑暗當中創造秩序.
那麼這個對我們今天的生活有什麼意義.
弟兄姊妹.
當你的生活陷入混亂.
比如當你回到家發現家裡氣氛很糟.
孩子吵鬧.
這個配偶說話還帶刺.

$^{161}$這是家裡的空虛混沌.
這個時候你不要跟著大吼大叫.
這個只會讓黑暗加深.
你要成為那個宣告光的人.
當然我不是說你要去講一些虛情假意.
只想逃避問題的場面話.
宣告光也不是鴕鳥般的自我催眠.
宣告光是要引導神的真理.
進入混亂當中.
因為只有真理的光才能夠恢復秩序.
沒有真理的光.
那只不過是一個虛偽的好話而已.
沒有真理的秩序.
只是自欺欺人的粉飾太平.
弟兄姊妹.
我們常常誤解什麼叫做和睦.
我們以為說我們只要不吵架.
大家都裝著笑臉.
用著一種屬靈的術語彼此敷衍.
這個叫光.
錯了這個不是光.
想想看當一段關係有一些隱藏的裂痕.
當一個人的生命有未對付的罪.
當浮世的當中充滿人的意識跟血氣.
你卻走到他們面前說.
平安了平安了沒事了.
這個是宣告光嗎.
不是的.
閻麗美書早有一個嚴厲的警告.
說他們輕輕忽忽的醫治我百姓的損傷.
說平安了平安了.
其實沒有平安啦.
還記得大胃王嗎.
當他奪人妻子.
甚至還借刀殺了對方的丈夫.
他用君王的權力把一切掩蓋得天衣無縫.
在皇宮裡面沒有人敢提這件事情.
是不是很和睦啊非常和睦啊.
國政依然運行祭祀仍然照舊.
但這是虛假的秩序.

$^{201}$所以神就差遣了先知拿丹.
拿丹走到大胃的面前.
他沒有昧著良心說好.
大胃啊我宣告你以後一定會得勝.
拿丹沒有這麼講.
相反的拿丹講了一個有錢人奪走窮人楊高的故事.
哇大胃聽了蹊蹺生煙啊.
說什麼行這事的人該死.
結果拿丹如同閃電一樣.
把真理的光直直刺進了大胃靈魂的深處.
他指著大胃的鼻子說你就是那人啊.
弟兄姊妹啊.
這句話聽了多刺耳啊.
他瞬間擊碎了大胃苦心經營的虛假秩序.
但請問.
拿丹這句話是光嗎?絕對是.
這句話是真理嗎?毫無疑問.
正因為這一句不妥協的真理.
大胃在神的面前徹底的破碎.
然後寫下了流傳千古的詩篇51篇.
說神啊.
求你為我造清潔的心.
使我裡面重新有正直的靈啊.
大胃的生命因為這一道刺痛的真理之光.
從道德的混亂當中被拯救回來.
真正屬靈的次序重新在祂裡面恢復.
當然我不是說等一下你聽完道之後.
你馬上跑去對你討厭的牧師或弟兄姊妹.
指著大胃鼻子說你就是這人啊.
宣告光是誠實.
但有愛心的去宣告神的真理啊.
唯有當虛假的舊秩序被真理的光拆毀.
神聖潔真實的新的次序才能夠被建立起來.
所以面對生命的混亂.
也許你要心平氣和的對你的配偶說.
親愛的.
我們不能再繼續這樣下去了.
我們之間失去了神的愛.
我們必須要一起面對問題啊.
我們必須一起跪下來在神面前悔改啊.

$^{241}$當你面對混亂的財務狀況.
不能成天只是宣告上帝的祝福啊.
你要邀請神的靈運行在你的帳簿上面啊.
禱告說神啊.
這是我的混亂.
但我邀請你的秩序進來啊.
然後用真理的光來審視自己.
我是不是在用消費來填補我內心的空虛啊.
弟兄姊妹啊.
如果我們不承認自己在財務上的貪婪跟無知的話.
再怎麼宣告祝福.
依然會活在財務的咒詛裡面啊.
今天當我們面對同事之間或同工之間的嫉妒紛爭.
宣告光.
絕不是在開會的時候想講場面話.
宣告光是帶著為父為母的心腸.
用聖經的真理來勸告這個團隊說.
同工們.
我們現在的狀態不符合神的喜悅.
我們失去了彼此相愛的心.
先要停止爭吵.
先來彼此認罪互相代酒.
這個才是真宣告光啊.
弟兄姊妹.
我們被呼召不是來世界上做一個粉飾太平的和事佬.
常有人說大事化小事小事化無事.
其實最後都是敷衍了事.
逃避問題妥協真理.
讓錯誤跟罪惡繼續蔓延而已.
弟兄姊妹.
我們被呼召要做真理的註實和根基啊.
所以我們口裡要有神的真理.
當你帶著真理走進辦公室.
那些的說謊作假的黑暗就要退避三色.
當你帶著真理回到家中.
那些自私苦毒的破口就要被堵住啊.
當神在創世記第一章宣告要有光的時候.
那道光是帶著威力的.
撕碎了無盡的黑暗.
而今天這個創造宇宙的光.

$^{281}$這個使死人復活的光就在你們裡面啊.
所以不要再用虛假的安慰來麻痺自己.
不要再用蜀式的哲學來稀釋神的真理.
拿起聖靈的寶劍就是神的道.
在你的家庭當中.
在你的婚姻當中.
在職場在財務在事奉的每一個角落.
大聲勇敢的毫不妥協的宣告神的光.
你會發現當神的光發出來.
混亂要消失.
秩序就要恢復.
第二個 用原創取代比較.
跟你旁邊講 用原創取代比較.
有些時候混亂不是來自於外在的遭遇.
而是來自於我們內在的比較.
C.S. Lewis 說 驕傲的本質就是競爭.
一個人不會因為擁有某樣東西而快樂.
他只有在比別人擁有更多的時候才快樂.
弟兄姊妹 這個是人生很可悲的現實.
但人還是很愛比 對嗎.
但請注意 只要你一比.
你的快樂跟平安.
就開始由別人的生活來決定.
想像一下 你原本對生活還算滿意.
你開著一台有點年紀.
但是冷氣吹了還會涼的車.
帶著家人去吃了一頓溫馨的晚餐.
覺得歲月靜好.
但偏偏你飯後手癢 點開了朋友圈.
剛好滑到那位喜歡炫耀的同學.
換上一款最新的跑車.
繼續一滑 還看到他上週.
在法國的露天餐廳打卡喝香檳.
還對著背後的巴黎鐵塔比個耶.
就在那一秒鐘 你有沒有發現.
很多人的平安就灰飛煙滅了.
突然覺得自己的車就破銅爛鐵了.
晚餐索然無味.
然後開始對自己的配偶發脾氣.
嫁給你怎麼那麼命苦啊.

$^{321}$這就是失控啊.
原本美好的人生.
因為別人發的一張照片就被迫改寫.
弟兄姊妹比較就是對著鏡子裡的自己在開槍.
你以為你在瞄準別人的成功.
但其實卻擊毀自己的形象.
這從來不是神原本的設計 你知道嗎.
創造宇宙的神是怎麼樣看待這個愛比較的世界.
如果你讀《創世記》第一章.
你會發現在前五天的創造裡面.
神造花 草造飛鳥 魚.
聖經重複出現一個詞叫「各從其類」.
這什麼意思.
簡單來說就像買手機.
只要型號一樣.
基本上每一隻都長得差不多.
這個叫「各從其類」.
但是神對你的創造難道也是這樣嗎.
不是的.
在第六天神沒有說「讓人各從其類」.
在《創世記》一章26到27節.
神說我們要照著我們的形象.
按照我們的樣式造人.
使他們管理海底的魚 空中的鳥.
地上的牲畜和泉地.
並地上所爬的一切昆蟲.
神就照著自己的形象造人.
還是照著他的形象造男造女.
在古代近東的文化裡面.
形象這個詞有一個特別的意義.
當時君王會在他們統治的疆界裡面.
樹立自己的雕像.
意思就是說這片土地屬於我.
我的權柄就在這裡.
而這個戴著權柄的君王雕像.
就是形象的意思.
所以當神宣告說.
他按照自己的形象造人.
就是在告訴整個受造界.
今天你我作為人.

$^{361}$就是神放在地球上的一個代表.
不管你在哪裡.
你都帶著創造之神的尊貴跟統治的權柄.
就像C.S. Lewis講的說.
這世界上根本沒有平凡人.
每一個活著的人都帶著神聖永恆的重量.
但問題是.
雖然每個人都是神所創造獨一無二的藝術品.
但很多人卻常常活成別人的仿冒品.
弟兄姊妹.
當你早上洗臉看著鏡子裡的自己.
你是看見神的形象呢?.
還是一堆需要修圖的缺點呢?.
要是我的鼻子再挺一點就好了.
白頭髮再少一點.
要是我能夠像那個網紅一樣.
受歡迎就好了.
很多人拼命想要成為別人的仿冒品.
這就像是一個人.
明明是一個符合價值連城畢卡索的真機.
卻哭著對上帝說.
拜託你把我變成一張韓國歐巴的海報就好了.
這種自我否定就好像是在跟神抱怨說.
你把我造壞了.
但弟兄姊妹.
神把祂的尊貴放在你的裡面.
神讓你看起來不像別人.
那是因為你根本不需要像別人.
就像《釋字記》裡面的機電.
有一天神的使者向機電顯現.
這個老兄在哪裡?.
他在酒駕裡面打麥子.
照道理打麥子是要在開闊的地方.
讓風把那個麥糠吹走.
但這個老兄躲在酒駕裡面打麥子.
沒有效率又滿臉灰塵.
為什麼?.
因為他害怕敵人偷襲.
而就在他灰頭土臉的時候.
天使隊推他說.

$^{401}$大能的勇士耶和華與你同在.
如果你是機電你會怎麼回?.
我大概會看看旁邊.
不好意思你在跟誰說話.
這裡只有膽小鬼一位.
而機電也正是這樣回應.
他說 主啊 我的家族是這個支派當中最窮的.
而我又是我們家最小最廢的那一個.
弟兄姊妹啊.
請問神是跟一個.
在躲在酒駕裡面發抖的農夫開玩笑嗎?.
不是的.
雖然機電自己跟別人一比.
只看到貧窮為小.
但神看機電看到的是大能的勇士啊.
這個世界定義.
這個機電只是一個沒用的懦夫而已.
但神說 不 機電.
你的原廠設定是大能的勇士.
而當機電相信神的定義.
這個原本躲起來的膽小鬼.
帶著三百個人打敗了十多萬的大軍.
這怎麼改變呢?.
機電的肌肉沒有變大.
他的存款沒有變多.
改變的只有一件事情.
就是他找到了原創的自己.
撒布爾記下還記載另外一個人叫米菲波瑟.
他是約拿丹的兒子 是掃羅王的孫子.
他本來應該要在皇宮裡面長大.
但因為一場政治動盪.
他在逃難當中就摔瘸了雙腿.
從此以後他的人生陷入谷底.
並住在一個叫羅底巴的地方.
羅底巴是什麼意思?.
叫做無話可說.
你可以想想看.
如果有人問你.
你家住哪裡?.
無話可說.

$^{441}$這有多淒涼啊.
所以米菲波瑟也真的就在羅底巴.
活成一個無話可說的樣子.
當大衛王因為跟約拿丹的誓約.
要派人把他找回皇宮.
要把掃羅王的一切產業歸還給他.
並且讓他有尊榮的時候.
米菲波瑟的反應是什麼?.
他說:僕人算什麼?.
不過如死狗一般啊.
這一位曾經的王子.
在經歷人生的殘酷.
身體的殘缺跟世界的冷漠之後.
給自己下的定義.
就是如死狗一般.
但大衛怎麼做?.
大衛覆蓋了米菲波瑟的自我否定.
他把掃羅的產業都給了米菲波瑟.
並且讓他坐在大衛王的餐桌.
享用王的宴席.
大衛的舉動就是宣告說.
我不在乎你的腿瘸了沒有.
我不在乎你過去失去了什麼.
但因為你身上留著我摯友約拿丹的血.
所以你是極有價值的.
弟兄姊妹.
大衛所預表的就是耶穌基督.
而今天耶穌基督也在對你發出同樣的呼籲.
你是極有價值的.
因為你身上是留著耶穌基督的血.
你想想看.
這個世代何等的競爭.
何等的殘酷跟激烈.
可能一天之前你還被尊為龜骨的精英.
但一天之後你發現.
自己的門卡連刷進公司都刷不進去.
甚至連被開除的理由都不告訴你.
你看看我們年輕的下一代.
他們更是無情的經歷被成績.
被名校的光環.

$^{481}$被同學的成就狠狠碾壓的痛苦.
有許多的孩子為了申請一個好學校.
找到一份好實習.
焦慮的夜不能眠.
這是一個多麼殘忍的裸底吧.
許多人都成了鼠鱗上雙腿瘸了的那個彌非菩薩.
在那個無話可說的荒涼之地.
好像找不到自己的價值何在.
但是今天聖靈要猛烈的搖醒我們.
要從世界的謊言當中要醒過來.
你的價值如果建立在職位上面.
那麼當你退休或裁員的那一天.
你的世界就崩塌了.
你的價值如果建立在財富上面.
那麼當經濟動盪的那一刻.
你的靈魂也跟著破產了.
你的價值如果建立在孩子的成就上面.
那麼孩子表現如果沒有鄰居的孩子好.
你的自我認同就徹底粉碎了.
不要將你的價值建立在這些會改變的人事物.
你要回到創世記第一章.
回到宇宙開始的源頭.
回到你被造的起點.
在那裡沒有KPI.
沒有SAP.
也沒有訂閱數.
在那裡只有創造天地的主.
用祂莊嚴而慈愛的聲音宣告說.
我們照著我們的形象.
按著我們的樣式造人.
這就是神給你的終極定義.
不是因為你做了什麼偉大的事情.
而是因為你們是他們的兒女.
當你真正明白這一點.
你的生命就會產生荷包級別的自由.
當你看著銀行戶頭裡面的數字.
無論多少你就不再被它綁架了.
因為你知道真正的財富是天上那不可朽壞.
不能玷污也不能衰殘的產業.
當你面對那些在社群媒體上炫耀完美生活的人.

$^{521}$你不用再感到自卑跟焦慮.
因為你清楚的知道那些濾鏡背後的虛榮.
遠比不上你裡面那擁有神形象的榮光.
讓我分享我自己的見證.
當我放下工作要全時間事奉主.
我的父親把這個消息告訴家族裡的一位長輩.
那位長輩怎麼說呢.
他說 佐爾君這個二等兵還學人當牧師.
的確 從人的觀察.
我是屬於那種私底下真的無話可說.
一緊張就口無倫次的那種超級哀人.
我所有的心理測驗全部都是極內向的那一邊.
請問這種人怎麼當牧師怎麼講道.
但今天我站在這邊我要說.
不管人怎麼樣的認定.
神給我的原廠設定就是神的僕人.
Hallelujah.
即使我有軟弱.
但耶穌基督在我裡面不是軟弱的.
耶穌基督在我裡面是有大能的.
所以弟兄姊妹不要再讓那些沒有永恆價值的人事物.
來定義你無價的靈魂.
耶穌基督已經用祂的寶血在十字架上.
為你附上宇宙最高的屬價.
既然神認為.
祂值得用耶穌的命來換.
來換你的生命.
那麼世界上就沒有任何人.
可以資格說你不夠好.
所以第三個.
用聖好對付焦慮.
跟你旁邊講用聖好對付焦慮.
當我們來到創世記第一章的.
尾聲神創造王天地萬物.
祂給了一個評語說.
聖好.
接下來是一個震撼的結束.
創世記第一章31節說.
神看一切所造的都聖好.
有晚上有早晨這是第六日.

$^{561}$天地萬物都造齊了到第七日.
神造物的工已經完畢.
就在第七日歇了祂一切的工.
安息了.
神賜福給第七日定為聖日.
因為在這日.
神歇了祂一切創造的工.
就安息了.
經文說祂就安息了.
為什麼很震撼呢.
因為這個不是我們的習慣.
神為什麼要安息.
祂累了嗎.
是不是祂造完大象.
哎呀 腰痠背痛.
不是.
神安息不是因為祂累了.
而是因為祂完成了.
神懂得欣賞自己的工作.
神看著這一切滿意的說.
夠了 這個很美好.
所以祂選擇停下來享受這一切.
讓我們發揮一下想像力.
如果創造天地的是我們.
那第六天晚上會發生什麼事情.
大概率是將不會有安息日.
我們會焦慮的整夜加班.
我們會對著天使大吼大叫.
加百列.
你看那個長頸鹿的脖子.
是不是太長了.
這個喝水會腦充血.
修一下 修一下.
還有那個斑馬.
那個斑馬誰畫的條紋呢.
這個左邊右邊不對稱.
這個強迫症快發作了.
重畫.
還有那個人.
智障設得有點低.

$^{601}$不行 不行 不能出貨.
如果讓我們來創造世界.
這個世界永遠創造不完.
因為在我們的眼裡.
永遠都不夠好.
我們會改革不停.
然後直到累死為止.
弟兄姊妹.
人之所以會焦慮的原因之一.
就是我們不懂得欣賞現在的狀態.
我們總覺得下一站的風景會更好.
所以我們心裡有一套公式.
等我結婚就好了.
等我升上主管就好了.
等我的孩子考上大學就好了.
等我退休就好了.
結果我們一輩子都在追逐下一站.
但因為我們不知道下一站.
真正會發生什麼.
焦慮就產生了.
弟兄姊妹.
未知永遠是焦慮跟恐懼.
最主要的原因.
因為當人焦慮的時候.
他其實會更焦慮.
他會停下來嗎.
不會的.
就好像迷路的時候.
你會更拼命地去找路.
結果一輩子都在奔跑.
從來沒有停下來.
沒有深吸一口氣.
對現在的生活說一句.
現在這樣甚好.
親愛的弟兄姊妹.
如果你學不會在今天說甚好.
即使你退休到了夏威夷的海灘.
你還是學不會的.
因為焦慮是無底洞填不滿的.
但是神卻讓我們看到.

$^{641}$無論在何種狀態.
他都能夠找到甚好的安息.
想想看耶穌.
當他被釘在十字架上.
那是他肉體最痛苦.
環境最黑暗的時刻.
從人的角度來看.
那是一個災難級別的空虛餛飩.
門徒跑光了.
救主自身肉體掛在那邊.
被世人所羞辱.
但是在他斷氣的前一刻.
耶穌說了一句震撼宇宙的話.
他不是說哎呀痛死我了.
也不是說完了.
計劃失敗.
他是說成了.
成了這個原文的意思就是完工了.
耶穌沒有等到所有人都信主才說成了.
因為他知道.
當他完成父神託付的那份責任.
即使在人眼中看來是最糟糕的狀態.
但因為在神的手中.
就有甚好的安息了.
所以今天我要給各位一個非常實際的挑戰.
學會用甚好來結束每天的生活.
當你下班了.
練習在心裡說今天的工作甚好.
不要繼續讓自我價值.
被那種無止盡的產出所綁架了.
如果你要等到所有事情都完美結束.
才肯休息.
你終將會累死.
當你的孩子拿著成績單回來.
盡力了但不如預期.
你要看著他的眼睛對他說甚好.
不要把孩子當作那種擴充自我成就的工具.
不要孩子成績好.
我們覺得自己是成功的父母.
結果孩子表現差.

$^{681}$甚至有點叛逆.
我們就覺得自己好丟臉.
不要這樣.
如果我們這麼做的話.
就是在用分數.
論斷一個神所創造的靈魂.
相反的我們要深呼吸.
看著孩子沮喪的眼睛對他說甚好.
當你這麼說的時候.
你是在打破一個叫做成就的偶像.
你要對孩子說.
你的價值是耶穌的寶血所買俗的.
不是被考卷所定義的.
你在神的手中.
在父母的愛中.
你本身的存在就甚好了.
今天當你洗臉看著鏡子.
你發現白頭髮又多了幾根.
皺紋又多了幾條.
體力大不如前.
甚至被長久的疾病所折磨.
你可能覺得自己越來越沒有用.
你不要對著鏡子嘆氣.
雖然我們活在一個極度崇拜青春和肉體的世界.
也許這個世界總是在釋放一個信息說.
老去是一個咒詛.
生病是一種失敗.
但是保羅在哥林多後書第四章.
發出了一個震耳欲聾的真理宣告.
他說所以我們不上膽.
外體雖然毀壞.
內心卻一天心是一天.
我們這至善之心的谷主.
要為我們成就那極重無比永遠的榮耀.
所以弟兄姊妹.
你的白頭髮不是衰敗的象徵.
那個是你為了愛家人.
為了服侍神.
熬過無數的夜晚所換來的榮耀光明.
你臉上的皺紋是你生命經歷了無數次.

$^{721}$上帝恩典淬煉的刻痕.
所以不要被這個膚淺的世界騙.
你可以對著鏡子輕輕摸著你臉上的皺紋.
說這個是我為家人.
為神傾倒生命的記號.
甚好.
雖然我們的身體會朽壞.
但是在我們裡面.
那個屬神的生命.
正在發出前有未有的榮光.
我們永活的救主耶穌.
就是我們在度過空虛混沌黑暗的日子.
那永遠的光.
祂也是我們可以說.
甚好的永遠盼望.
親愛的弟兄姊妹.
今天早上神要邀請你進入祂的安息.
放下那些我們認為不完美.
所有的眼淚所有的疲憊.
然後奉主耶穌基督的名.
大聲對你的環境對你的靈魂宣告.
神與我同在祂的恩典供我用.
今天甚好.
阿雷路亞.
所以跟你旁邊的人講.
今天甚好.
再說一次今天甚好.
再說今天真正好.
因為神的恩典供我們用.
願神賜福你們我們一起來禱告.
因為在你裡面.
我們是有何等尊貴的身分.
雖然許多時候.
我們會有一些自卑的心.
覺得我們好像有一些比不上別人.
的一些的焦慮.
或者是這個世界我們周遭的人.
輕看著我們.
但主你的應許是真實的.
我們跟隨耶穌基督的人.

$^{761}$必不致羞愧啊.
我們必不致羞愧啊.
主你已經在我們的生命當中.
放下那榮美的形象.
你賜給我們那尊貴的身分.
今天當我們在耶穌基督裡面的時候.
我們是你所設定.
那個神的兒女.
是那王子公主的位分.
主啊.
你讓我們能夠重新找回.
你在我們生命當中的原創吧.
那個榮耀.
那個尊貴.
那個你所喜悅的樣子.
雖然我們的肉體會衰殘.
但我們的內心要一天心過一天.
看著我們的遭遇.
看著我們的苦楚.
看著我們所經歷的一切.
我們深知道.
在那永恆的國度當中.
這一些都要成為我們榮耀的印記.
主我們謝謝你.
因為你是何等的愛我們.
在我們所行的每一個路上.
你與我們同在.
你真理的光.
復活的光.
創造的光.
就在我們裡面.
只因為我們神聖的道路.
讓我們所走的每一步路.
都帶著平安.
帶著喜樂.
帶著那個甚好的安息.
在我們裡面.
主我們謝謝你.
我們這樣禱告.
奉主耶穌基督的聖名.

$^{801}$Amen.
\newpage



\section{約翰福音 5:17-21}
\label{sec:031846}
\textbf{你為何仍然活著?}~~~~招世超牧師
\newline
\newline
連結: \href{http://www.church.com.hk/acms/content.asp?site=cdc\&op=show\&type=product\&code=031846\&layout=sermon}{\texttt{031846}} ~~~~ 語音日期: 2026-03-22
\newline
\newline
\hyperref[sec:031845]{\small{< < < PREV SERMON < < <}}
~
\hyperref[sec:index_chronic]{\small{[返順時目]}}
~
\hyperref[sec:031847]{\small{> > > NEXT SERMON > > >}}
\newline
\newline
約翰福音 5:17-21
\newline
\begin{longtable}{cl}
\hline
\hline
章節 & 經文 (和合本修訂版)\\
\hline
5:17 & \begin{tabularx}{0.7\textwidth}{X} 耶穌就回答他們：「我父做事直到如今，我也做事。」 \end{tabularx} \\ \\ \relax
5:18 & \begin{tabularx}{0.7\textwidth}{X} 為了這緣故，猶太人越發想要殺他，因為他不但犯了安息日，而且稱神為他的父，把自己和神看為同等。 \end{tabularx} \\ \\ \relax
5:19 & \begin{tabularx}{0.7\textwidth}{X} 於是耶穌回答，對他們說：「我實實在在地告訴你們，子憑著自己不能做甚麼，惟有看見父所做的，他才做；父所做的事，子也照樣做。 \end{tabularx} \\ \\ \relax
5:20 & \begin{tabularx}{0.7\textwidth}{X} 父愛子，將自己所做的一切事指示給他看，還要將比這更大的事給他看，使你們驚訝。 \end{tabularx} \\ \\ \relax
5:21 & \begin{tabularx}{0.7\textwidth}{X} 父怎樣叫死人復活，賜他們生命，子也照樣隨自己的意願賜人生命。 \end{tabularx} \\ \\
[1ex]
\hline
\hline
\end{longtable}
$^{1}$好,今天就跟大家講一個題目,好像很深.
為什麼我們還要活著呢?.
其實大家都知道,我不是第一次講.
2026年開始的時候,世界就變得很亂.
我們覺得亂,世界又打仗了.
你最近有沒有入油?.
現在差不多是六塊錢,有些是七塊錢.
七塊錢,我們有沒有試過七塊錢?.
起碼是六塊錢試過.
給我們帶來一個很大的壓力.
在生活上,似乎現在比以前生活就越來越艱難.
如果我們將人生當作負債表.
我們看看下一張圖畫.
你會發覺其實我們在左手邊.
即是你們的左手邊.
我們覺得生活裡面帶給我們那種壓力,挑戰.
以及那些難處,是比右手邊.
即是平安,快樂,順遂,喜樂.
似乎真的左邊多於右邊.
而且那些問題,也是左邊比右邊大.
既然人生不如意的事,中國人早就已經講了.
是十之八九,不是十之四五.
是十之八九.
生活那麼艱難,為什麼我們還要活下去?.
這些是很深的問題,可能可以說是很深的問題.
所以很多時候我都不想去想這個問題.
但如果要問這個問題的時候.
我相信一個最原始的答案.
你會是說什麼給我聽呢?.
最原始的答案就是,人生出來就有一個求生的本能.
有沒有人生出來求死的?.
baby什麼都不懂,但是也要喫東西.
所以我們發覺人.
從一個最原始的觀念去看的時候.
其實和動物的分別不大.
人有求生的本能,人是為了生存而去生活.
這個大家聽過了.
我們不會隨便求死的.
人和動物不同的地方在哪裡呢?.
大家都有求生的本能,但是不同的地方在哪裡呢?.

$^{41}$不同的地方,其中一樣東西就是.
如果人會發覺,我們只是為了生存而去生活的話.
做人太沒意思了,是不是呢?.
只是為了生存,為了兩餐,為了能夠survive.
繼續生活下去.
做人就沒什麼意義了.
其實美國大家都知道.
雖然打仗,是全世界最富裕的國家.
在美國要生存的話.
其實一點困難都沒有的,是不是?.
我來美國幾十年了,我從來沒聽過美國會餓死人的.
不過呢,我們時不時都會聽到.
在美國有些名人會自殺的消息.
大家同不同意我的說法嗎?.
事實上大家可能不知道,我也不知道.
不過因為講這個題目,我去做了一些research.
我最新找到的資料是2024年.
大家知不知道2024年有幾萬人自殺死亡?.
大家可以看到2024年有48,824個人死於自殺.
大概每11分鐘就有一個人自殺死亡.
其實原來在2024年,想自殺的人不止48,000.
48,000人自殺死了.
其實有290萬人,包括有70萬個青少年人.
是嘗試過自殺,不過不成功.
那是什麼意思?.
大概每11秒鐘,每11秒鐘.
就有一次自殺未遂的個案發生.
不知道大家覺得夠不夠震撼呢?.
如果覺得不夠震撼的話.
在同一年,我告訴大家.
其實是有460萬人曾經制定過.
怎麼自殺好呢?.
想過方法,怎麼去計劃自己自殺.
而大概有1430萬人曾經有過嚴重的自殺念頭.
什麼意思?.
就是說在美國,每24個美國人.
就有一個人在過去一年,曾經很認真地想過要自殺.
而18到25歲這個年齡層的比例.
就更加驚人了,大概每10個年輕人.
是每10個年輕人,有一個認真想過要自殺.

$^{81}$當然除此之外.
還有一個好像我們思考的一個發言.
大家知不知道自殺率最高的年齡層是什麼人呢?.
這條講給你聽,下一章應該是.
自殺率年齡層最高的呢?.
我還沒到,是85歲以上.
不知道我們當中有沒有少數呢?我相信是85歲以上.
大概那些是什麼人呢?.
通常都是隻有自己一個人.
家裡人都走了,或者年老多病那些.
但是知不知道第二個最高自殺的年齡層是什麼人呢?.
是幾多歲?越來越年輕,是10到24歲.
如果大家看到這些數據,這些資料的時候.
我不知道你們有什麼看法,有什麼感覺.
你怎麼去解讀這些數據.
如果你問我的話,最簡單的一個.
你問我為什麼的話,原來在美國.
這個是很多人以為可以發夢成功的地方.
有很多人是生活得不開心.
有很多人生活得不快樂.
甚至是想過要自殺,離開這個世界.
究其原因呢,我很同意Victor Frank.
Victor Frank是澳大利亞一個很出名的.
當然已經走了,一個很出名的精神心理學家.
他寫了一本書,叫做《活出意義來》.
他在這本書裡面說道.
人最痛苦的,不是受苦.
是什麼呢?是找不到受苦的意義.
很有意思,去想一想.
接著他在這本書裡面繼續說.
即使是最極端的苦難,人類在最痛苦,最慘的最極端的苦難當中.
只要人能夠找到一個理由.
這個理由呢,無論是為了某個人,某件事.
或者是某個未完成的心願.
無論什麼都好,只要能夠找到一個理由.
他就能夠繼續活下去.
這是他幾十年做心理治療師的經驗分析.
他也是一個哲學家.
不知道大家有沒有想過呢?.
我們為什麼仍然活著?.

$^{121}$今天我們為什麼仍然在生活呢?.
不知道我們能不能找到一個理由.
可以讓我們繼續活下去.
當耶穌出來傳道時,其實只是三年多的時間.
他一出來傳道,就引起當時社會很多的轟動.
因為聖經形容耶穌講道,教導的時候.
不同於一般猶太人的律法師.
而且耶穌在講道的同時還會醫病,趕鬼,行神蹟等能力.
聖經記載了一個故事.
有一個癱瘓的人走不了路.
讓四個朋友帶他去見耶穌.
希望耶穌能夠醫治他.
如果大家記得那個故事,可以告訴我們.
帶他去耶穌在講道的地方.
整個屋子都塞滿了人,根本不能進去.
於是這四個朋友沒有放棄.
就帶著這個癱瘓的人走到猶太人的平房.
不僅如此,還拆了人家的屋頂.
把這個癱瘓的人吊在耶穌的面前.
聖經告訴我們,耶穌見他這樣的時候.
就醫好了這個人.
叫他起身行走吧.
不單單是起身行走.
請你把你的丹加床也拿走.
之後聖經告訴我們,耶穌再次遇到這個人.
見到這個人的時候,耶穌再次勸他,提醒他.
以後不要再犯罪了.
很明顯這個人之所以癱瘓.
是因為犯罪得罪了上帝.
不知什麼原因,癱瘓了,不懂得行走.
現在我治好了你,好好地生活.
不要再犯罪了.
但是聖經告訴我們.
這個人竟然告訴那些猶太人的宗教領袖.
原來是耶穌醫好了我.
我們看第五章十五,十六節,讓夫音.
聖經說,那個人走了之後.
就告訴猶太人的領袖.
說是耶穌醫好了他.
帶來的結果是什麼.

$^{161}$十六節,從此他們開始迫害耶穌.
因為他在安息日醫病.
我剛纔說過,耶穌出來傳道.
其實是三年左右的時間.
他在世界上活了多久?.
三十三年.
如果從我們今天的角度來看.
哇,三十三年,很年輕而已.
英年早逝.
很短命耶穌,十分短命.
但是耶穌的人生十分精彩.
他怎麼在世界上生活的呢?.
為什麼他活得這麼精彩.
能夠影響世界上這麼多的人.
我們從聖經的記載和故事裡面.
我們可以看到.
耶穌在世上生活的時候.
第一樣我能夠觀察到的是什麼.
他每一天是以天父的事為念.
那些猶太人的領袖.
因為耶穌在安息醫病.
違反了猶太人的戒律.
十戒當守安息日.
你違反了十戒.
於是很不開心.
似乎是不是很有理由呢?.
但是這些人不僅不開心.
還開始想辦法去對付耶穌.
這纔是問題.
那耶穌知不知道呢?.
耶穌當然知道.
耶穌知道,但耶穌的回應是什麼呢?.
為什麼我說耶穌在世上生活的時間裡面.
是經常以天父的事為念呢?.
十七節裡面說.
耶穌對他們說.
我的父親一直在工作.
我也應該照樣做工.
主佑弟兄姊妹.
如果上帝安息會休息的話.

$^{201}$那就弊加火了.
千萬不要安息是病.
千萬不要安息是意外.
千萬不要安息是遇到什麼特別的問題.
因為上帝休息.
休息就什麼都不做了.
這是猶太人的觀念.
上帝是安息的時間.
我們不可以做任何的事.
但是詩篇的作者告訴我們.
我不知道有沒有引出來.
保護以色列也不打盹也不睡覺.
是什麼意思呢?.
原來我們所相信的上帝是不需要休息的.
上帝不需要睡覺.
他休息不睡覺.
做什麼呢?.
原來上帝是不斷地在工作的.
做什麼呢?.
保護他的百姓.
照料他所創造的宇宙萬物.
所以上帝是從來沒有停止過做工的.
作為上帝的兒子.
耶穌知道上帝是他的父親.
那他怎麼可以不工作呢?.
如果大家做工.
大家聽到就算公司不是.
有些公司文化是什麼呢?.
你老闆沒收工.
你為什麼要收工呢?.
你老闆還在辦公室工作.
那你自己先走嗎?.
會不會啊?.
作下狀就留在這裡吧.
做不做就一回事.
所以耶穌說.
我父親也在工作.
我也應該繼續工作的.
我父親一直在做工.
我也應該做工.

$^{241}$不過我要解釋清楚.
耶穌照樣做工.
不是因為出冤.
是怕老闆不開心.
怕父親不開心.
那是因為他尊重上帝.
因為他敬愛他的父親.
所以他時常都在想.
現在我父親想做些什麼呢?.
我要以他的事為念.
他想做些什麼.
我就和他一起做.
這是耶穌的生活態度.
我覺得是值得我們學習的.
神想做些什麼.
我們就跟著去做.
我們每一件事.
都要以我們天父的事為念.
其實這件事.
耶穌不是出來服侍神才知道的.
如果大家記得的話.
耶穌十二順就知道了.
所以為什麼教會很強調.
兒童教育,青少年的教育呢?.
耶穌十二順就知道了.
大家知道我說了什麼嗎?.
耶穌十二順的時候.
發生過一件事.
記不記得是什麼事?.
耶穌跟著他的父母.
育身的父母.
約瑟和瑪利亞去耶路撒冷的聖殿裡面敬拜.
完了之後.
大家大堆走了.
走了一半 為什麼不見到耶穌呢?.
開頭約瑟和瑪利亞就以為.
他在親族小朋友當中周圍去玩.
一路走一路玩.
結果不是.
於是回去找耶穌.

$^{281}$聖經就告訴我們.
在聖殿裡面找到耶穌.
當然很緊張.
兒子 有沒有搞錯?.
為什麼你不跟我們回家?.
知道我們多擔心嗎?.
又少了個宅男.
那時候耶穌說什麼?.
為什麼找我呢?.
好像很無聊.
「豈不知我應該以我父的事為念」.
耶穌十二歲就知道了.
他以後的人生.
他過的每一天.
都要以他天父的事為念.
這是他生活的態度.
耶穌在安息日.
醫好了一個生離.
不是生離.
因為犯罪而癱瘓的人.
其實是好事.
幫助人是好事.
但原來這個世界很現實.
也很殘酷.
有時幫助人會惹麻煩.
甚至要負代價.
耶穌因為幫助這個太子.
醫好這個人.
竟然引來猶太人的宗教領袖的迫害.
其實我們知道當日猶太人的領袖.
是誤會了上帝的意思.
誤會了當守安息日的意思.
如果我們看回聖經.
神創造宇宙萬物的時候.
聖經說.
神用了六日的時候創造天地.
到第七日就揭了一切的功.
這是《創世記》講的.
安息了.
並將第七日定為聖日.

$^{321}$然後.
因為在這一天神就揭了他一切創造的功.
就安息了.
我想大家留意.
剛纔我讀的時候.
不知道大家有沒有聽清楚.
聖經說.
神停止了一切創造的工作.
因為創造的工作完成了.
Done.
The job is done.
結束了.
不需要再創造什麼了.
按照上帝的計劃要創造的事情都做完了.
So the job is done.
就不等於神停止了不工作.
所以神將第七日定為聖日.
要紀念大功告成.
也讓他所創造的萬物有休息的機會.
但猶太人誤會了神的意思.
在安息日.
竟然連醫病,幫助人,做善事都不行.
而當耶穌這麼做的時候.
竟然還要想辦法如何對付耶穌.
中文有一句話.
耶穌做好心.
好像人們反而這樣對待他.
所以我剛纔說.
不知道你們有沒有試過.
有沒有那個經驗.
有時候原來做好人去幫人.
是會惹麻煩的.
我經常聽到大陸要小心.
不要隨便幫人.
隨時你幫那個人會惹麻煩的.
幫人要付代價.
輕的就簡單.
就是時間.
幫人要付時間.
要用精力.

$^{361}$有時候要付錢.
這些是簡單的生外物.
嚴重一點的.
可能會被人說你諸事八卦.
好管閒事.
想不想這樣被人說.
再進一步.
可能會被人排擠.
被人陷害,誣告.
甚至那些幫你的人.
你去幫那個人.
是會恩將仇報.
你幫了他.
他還不知道有意無意.
令你難堪.
令你受責難.
就好像耶穌一樣.
耶穌醫好了的膽子.
我們沒想過他會立刻去捕寸.
將這個消息.
告訴那些猶太人領袖.
以致這些猶太人領袖.
要想辦法去陷害耶穌.
原來這個世界.
有時候是這麼殘酷的.
我希望你沒遇過這種情況.
如果你遇過這種情況.
以後你還去幫不幫人.
要不要繼續幫人.
幫還是不幫.
如果我看耶穌的答案的時候.
耶穌的回應.
其實問題不是要幫還是不要幫.
我們再看看耶穌的回應是什麼.
耶穌說:我的父親一直在工作.
我也該照樣工作.
什麼意思.
耶穌的回應是說.
我們應該以父親所做的事為念.
神做什麼.

$^{401}$神想做些什麼.
祂會怎麼做.
我們就應該怎麼去做.
耶穌要傳達的信息是.
我是和我的父親同工的.
我知道要付代價.
但是當我做父親要我做的事的時候.
是做父親所關心的事的時候.
其實我受到的祝福和恩典.
會比我付出的更多更大.
這件事我覺得我們要彼此提醒和學習.
剛剛星期初的時候.
我和一個同工.
我鼓勵他.
剛剛做牧師沒多久.
我和他分享我做牧師幾十年的經驗.
我和他說.
你終身服侍上帝.
你會經歷到神給我們很多祝福.
我和他分享我自己的經驗是什麼.
很多時候神給我有很多驚喜.
是我沒有想過的.
我和他分享神呼召我做牧師的時候.
當然是幾十年前.
開始的時候我很頑固.
拖了十年才肯做牧師.
但當我清楚知道神要我去牧養教會的時候.
我就求神幫助我.
讓我能夠專心牧養教會.
我學習了做一個好牧師.
到今天我仍然在學.
所以不是說我已經得著了.
正如本來所講.
我仍然是努力向前.
希望繼續做一個神所允許的牧師.
但無論怎樣.
這個呼召這個意念一直放在我心當中.
就成為我人生的目的.
我要去牧養教會.
因為這個神給我的呼召.

$^{441}$有人叫我去做專家.
我說如果要做專家.
我去做一個牧養教會的專家.
這個是神給我的呼召.
這是我人生的目的.
所以為了牧養教會.
我放棄了.
沒有去完成我Columbia University的博士後裔.
只差一份paper 沒時間寫了.
我放棄了我的sabbatical leave.
沒有牛蒡不要緊.
不放假沒問題.
很多時候.
我會將人怎樣看我放在詞位.
什麼對教會是最好的.
那個纔是最重要.
所以因為這樣的緣故.
什麼對教會是最好的.
這個經常放在我心當中.
我就去做.
所以我剛纔提過的東西.
好像對某些人來說.
一份論文都不寫.
你sabbatical都不放假.
對我來說is nothing.
因為原來當我以為.
人家的詞位是犧牲的時候.
上帝給我的祝福是大過我所想.
上帝從來都沒虧待過我.
當我覺得沒能力的時候.
上帝就交給我力量.
我在教會以外.
很多人尊重我.
從來沒人會羞辱我的.
在我開始成年已經做分堂.
2022年的時候.
神就給我一個surprise.
你放棄了你的博士論文不寫.
我給你一個博士.
所以那一年.

$^{481}$上帝out of my thinking.
神學院就給了我一個榮譽博士學位.
我從來沒想過.
但是在這節目裡看到嗎.
給大家一點點見證.
當你學習每一件事.
都將神放為先.
將神給你的使命任務放在前頭.
神從來都不會虧待我們.
所以所有的兄弟姐妹.
不知道現在我們的生活.
I don't know about you.
當然我認識你們很多人.
但是我不知道你們現時生活的情況是怎麼樣.
到現在你們的生活.
緊張的是什麼.
你們市場忌慣的是什麼.
是我們的子女.
是他們的學業.
是他們的工作.
是他們的未來.
是我們何時退休.
退休的時候去玩些什麼.
還是我們忌慣一些其他的事情.
I don't know.
But you know.
You should know.
你應該知道.
不知道我們.
有沒有曾經想過.
都應該以我們在天上的父.
所做的工作為臉.
上帝在做些什麼工作.
我剛纔在粗略地講過.
就是照顧他的百姓.
照料他所創造的萬物.
不過聖經有句說話講得很清楚.
神的心願.
神的心志.
給我們多一件事.

$^{521}$神不願有一人沉淪.
乃願人人都悔改.
上帝沒有停止過做工.
耶穌告訴我們.
所以耶穌說我們應該做工.
做什麼工作.
原來神不願有一人沉淪.
乃願人人都悔改.
所以既然是這樣的時候.
耶穌就做一樣的工作.
我希望今天我們作為基督徒.
上帝的子女.
如果這是上帝的心意.
這個會不會也成為我們的心意.
當耶穌受到這些猶太人.
中國領袖的挑戰的時候.
他要不經意地.
將他生活的態度.
我豈不知我應該以我父的事為念.
這是他生活的態度.
但是他講的時候不經意地說.
我父親做什麼我就做什麼.
誰知道這個不經意的自白.
發自內心的一個表白.
卻令那些猶太人的宗教領袖動殺機.
不單單要害他.
要殺他.
為什麼.
聖經講得很清楚.
因為耶穌竟然稱上帝為他的父親.
所以聖經說.
十八節講到.
更使猶太人的領袖決意要殺害耶穌.
因為他不但破壞了安息日的戒律.
而且說上帝是他的父親.
把自己當作跟上帝平等.
耶穌這個公開的表白.
比不守安息日的罪更加大更加嚴重.
因為聖經講得很清楚.
那些猶太人說.

$^{561}$有沒有搞錯.
你叫上帝創造主做你的父親.
那就是說你是他的兒子.
你是他的兒子.
你的本質和上帝一樣.
那就是你是神.
有沒有搞錯.
這群宗教領袖說.
手段是律法的.
很嚴謹的.
你講一個神的名字都不行.
你竟然說神就是你的父親.
破壞了安息日.
於是動了殺機.
要殺耶穌基督.
如果不是耶穌那時候行完神蹟.
那些人慶賀一下.
我相信馬上就用石頭扔耶穌了.
但是耶穌知不知道.
耶穌都知道的.
耶穌回應是什麼.
我們看看耶穌.
原來他不單單以天父的事為念.
作為他生活的態度.
他還有一個為神而活的心智.
我們看第十九節.
接著耶穌對他們說.
我鄭重.
我很認真地告訴你們.
兒子憑著自己不能做什麼.
看見父親做什麼.
他才做什麼.
父親所做的.
兒子也照樣做.
耶穌說兒子憑著自己不能做什麼.
但如果是神的兒子.
又怎會說憑著自己不能做什麼呢.
似乎有些矛盾.
到底是什麼意思呢.
我們看看耶穌怎樣說.

$^{601}$兒子看見父親做什麼.
他纔會做什麼.
父親所做的.
兒子也會照樣去做.
聽起來好像在重複他以前說的話.
但其實他想告訴當時的猶太人領袖.
也想告訴所有在場的人知道.
他所做的一切.
他的言語,行動和生活.
都是以上帝為依歸.
耶穌的一生.
是為神而去生活.
這是耶穌來到世界的目的.
為神而活.
神差他來他就來 神要他死他就死.
所以雖然耶穌知道.
在安息日醫病和叫神做父親.
會帶來那些猶太人.
特別是宗教領袖的不滿,排擠.
甚至是要決意殺害他.
他也會去做.
而且會繼續這樣做.
因為他知道他所做的.
正是他的天父爸爸在做的事情.
所以 主佑弟兄姊妹.
我們祈禱的時候.
我們有沒有叫上帝做天父.
你想想你自己祈禱的時候.
你們是怎麼祈禱的呢.
你們是怎麼稱呼神的呢.
有沒有試過在祈禱的時候.
叫上帝做天父爸爸呢.
有沒有.
是一個很親密的稱呼.
但是大家明不明白.
我希望今天在這個時間.
你們能夠聽明白.
原來這個稱呼.
又包含了一個很深厚的意思.
既然我們稱神做我們的天父爸爸.

$^{641}$天父爸爸在做什麼.
他想我們做些什麼.
我們有沒有放在心上.
我們有沒有以天父的事為念.
耶穌來到世界.
三十三年是為神而活.
我們今天稱祂為我們的天父爸爸.
我們有沒有學去為神而活呢.
如果我們發覺.
這個是獻給我們的啟示和教導的話.
我們信耶穌之後.
我們就應該學.
從為自己而活.
為他人而活.
變成為神而去生活.
我們應該將我們生活的重心.
從自我轉移到神的身上.
今天很多人.
我相信包括基督徒在內.
我們仍然是為自己而活.
我們想自己生活得開心一點.
快樂一點.
所以你想開心點快樂點.
人是很軟弱的.
很自然我們就會追求感官上的滿足.
譬如喫多點自己喜歡喫的東西.
買多點清新的衣服.
讓自己看起來很presentable.
自己又舒服.
人又看得好.
我們要把握我們的長週末.
其他的假期.
穿著遊玩多點.
為什麼.
因為想生活開心點.
這個世界太艱難.
牧師一開出就說了.
我們希望能夠暫時忘記.
生活上的一些壓力和痛苦.
那我們就可以撐久一點.

$^{681}$我們可以繼續活下去.
不過其實我們都知道.
我們每天所面對的那些壓力.
挑戰和問題.
是不會消失的.
只是暫時舒緩一下.
有時我們為了這些東西.
心裡都很糾結 不知道應該怎麼想.
不過沒辦法.
這就是人生 這是現實.
我們這樣過生活.
你滿足了沒有.
你覺得你的人生.
如果是這樣過的時候.
是不是很有意思.
有時候我知道有些人.
會告訴我們自己.
我不能死.
我要繼續活下去.
為什麼.
如果我們死了 我們不在了.
有些人會很慘.
有些人的世界會崩潰.
他們會很不開心.
甚至生活裡沒有盼望.
比如我們的子女.
如果沒有我們怎麼辦.
我們的父母.
誰來照顧他們.
我們的配偶.
這種被需要的感覺.
有沒有.
你是被需要的.
其實是一種壓力.
對不對.
不過這種壓力很多時候.
就成為我們能夠持續生命.
一個最強的生命的韌帶.
吊著在這裡.
不要走.

$^{721}$繼續撐下去.
讓我們能夠繼續活下去.
但是無論我們是為自己而活也好.
為了照顧我們的父母.
妻子 兒女.
去生活也好.
這個.
都是一種負擔.
你也掛念著.
不知什麼時間.
我們不知不覺.
就會被困在自私和責任的囚籠裡.
唯有當我們像耶穌一樣.
以天父上帝的事為念.
學習為神而活.
於是我們慢慢去看著自己.
向內看 向下看.
準備向上看.
看著耶穌基督.
我們的人生.
就可以有一個新的方向.
有一個新的生活的目的.
生存的價值和意義.
我們要做上帝要做的事.
因為我們以天父上帝的事為念.
我們要做神正在做的事.
因為我們應該為上帝而活.
而上帝.
我們的天父.
他會繼續指引我們.
和我們一起去做一些更大.
更有意義的事情.
為什麼?.
因為神愛我們.
神希望我們活得有意義.
神希望我們活得更加開心快樂.
但不是從滿足我們自己的需要那裡著眼.
是從神和公那裡著眼.
我們看第20節.
五章20節有人否認.

$^{761}$父親愛兒子.
包括女孩子.
把自己所做的只是兒子.
他還要把彼者更重大的事.
只是兒子.
使你們驚奇.
父親怎樣使已經死了的人復活.
賜生命給他們.
同樣兒子也要隨著自己的意思.
賜生命給人.
耶穌說父親愛兒子.
所以將被醫病趕鬼.
甚至行神蹟.
更重大的事.
只是給兒子知道.
令人感到驚奇.
這是什麼事?.
聖經有說.
就是使死人復活.
使死人復活.
賜生命給那些.
不認識上帝的人.
讓他們能夠得救,有機會在天國裡生活.
耶穌說天父伯伯已經把權柄賜給了他的兒子.
這裡很明顯說耶穌.
但我剛纔說過.
你是否神的兒子?.
你是否神的女兒?.
我們都是神的兒子.
聖經說神將權柄賜給他的兒女.
讓他們可以隨意將生命賜給人.
你們有沒有發現我們信耶穌之後.
原來我們有沒有權柄?.
我們有沒有權柄在聖靈的大憐和感動下.
將生命賜給人?.
不是你我現在活在這個肉體的生命.
是將永恆的生命賜給人.
不單單將永恆的生命賜給人.
而且我們能夠幫助人.
在這個世界充滿問題,痛苦的環境裡.

$^{801}$可以活得更加有意思.
有人生的目的.
看到自己存在的那種價值.
親愛的兄弟姐妹.
我已經開始說了.
今天很多人生活得不開心.
很多人生活得不快樂.
因為多個人只是為自己想.
經常為自己的設想而生活的時候.
我們真的很容易會變得越來越自私.
而我們想抓住,有把握,可以擁有的東西.
我們大家都知道.
只是慣份衣人而已.
沒有永恆的價值.
如果我們只是為了責任而去生活的時候.
我們不自覺地慢慢變成一個生產的工具.
我們日挨夜挨.
挨世界.
就成為我們老闆和家人的奴隸.
當然他們不是這麼看的.
不過對我們來說.
人生的意義在哪裡.
不過我相信有更多人.
我一開始也說.
牧師你今天說的東西太難了.
這個問題太深.
其實不是很深.
不過我不想去想這個問題.
人生的意義有沒有搞錯.
很多人不想去想這些東西.
不想去找這些問題的答案.
那又如何.
那就按照生老病死.
這個人生命的週期去過生活.
當你這樣去過生活.
什麼都不問的時候.
不知道你有沒有覺得有時生活都很無奈.
有沒有這樣的感覺.
有沒有曾經有這樣的感覺.
有沒有覺得.

$^{841}$心裡面有時會有一些莫名的空虛.
我們到處都是人.
在你的社交媒體上.
可能有上百甚至過千的朋友.
但是能夠推心置腹.
站在你身邊.
真的明白你說什麼.
體會你的心情.
知道你的感覺的人有多少.
我們這個世界比美國越來越自由.
什麼都可以做.
原來什麼都可以做都是一種壓力.
太多東西可以做了.
做什麼纔好呢.
而當我們找到的時候.
發覺原來很多東西我們做不到.
是的你可以做.
但是你做不到.
我們的壓力更加大.
但是如果神恩代我們.
讓我們能夠幫助人的靈命復甦.
讓他們能夠擁有基督耶穌的生命.
我們的生活就有了一個更崇高的價值和意義.
因此我們看到我們這個工作是多麼的重要.
聖經說.
人有了神的兒子就有生命.
很熟悉的聖經.
後面一句是什麼.
沒有神的兒子就沒有生命.
大家記不記得.
不信者得不著永生.
神的真路尚在他身上.
所以為什麼人生會這麼慘.
弟兄姊妹.
你看到嗎.
神父照你和我作為他的子女.
這樣一個很重要的工作交給了我們.
我們有權柄來自上帝.
只要我們肯隨著聖靈的帶領.
我們可以將人已死的靈性復甦.

$^{881}$將真正的生命賜給人.
哇 這是多麼大的祝福.
你沒想過可以做這些事嗎.
你可以的.
因為神已經將權柄賜給了你.
當我們能夠祝福人 幫助人 認識神.
讓這個人能夠找到人生的價值和意義.
這種滿足感.
就讓我們的生活變得更加的開心和快樂.
有一個同工和我分享.
上個星期他帶了幾個弟兄姊妹去參加商場報道.
大家知不知道商場報道.
"Let's go 1,2,3" 大家聽過嗎.
他說我們教會有幾個人去.
整個早上 到兩點了.
又填了幾張表回來.
就是回應表 和他分享福音.
那些人肯填回來.
但是這個同工告訴我.
他有一個姊妹.
不是我們教會的 是另一個教會的.
這個姊妹我都認識的.
她的身體好不好.
長年見到她都帶著帽子.
但是她一個人竟然有二十幾張表填回來.
她說她不是在這裡等著別人走過.
有興趣來問一下.
然後我們就和她說.
或者你經過我們就聽聽福音.
不是啊 她的姊妹是在商場每家每戶的.
去找人傳福音 講見證 分享生命.
結果.
一個早上.
二十幾人填回應表.
但是這個姊妹.
我覺得這個姊妹.
不是什麼教會的牧師傳道人.
但讓我們看到一個人.
他如何以天父上帝的事為念.
這個人 這個姊妹.

$^{921}$我看到他如何努力.
每天去為神而活.
今天我和你都生活在這個世界上.
這是天父的恩典.
天父的恩典是.
讓我們有機會和祂同供.
但是你願意和祂同供嗎.
你知不知道上帝在做些什麼.
你明天上班回家.
見到很多人.
為什麼神將他們帶到你的身邊.
你有沒有見到神在他們身上做些什麼.
你可以為神在他們身上做些什麼.
我希望弟兄姊妹.
我們一起學習.
如何為神而活.
當你這樣做的時候.
你的生命是會被神翻轉的.
不過這不是我們想要的嗎.
你生活得有意義.
你人生更加有價值.
學習如何為神而活.
我們一起祈禱.
天父上帝 我將時間交在你的手中.
因為這是你的時間.
你在我們當中.
你透過聖經.
你透過牧師的信息.
你透過詩歌.
向我們發出挑戰.
我今天為什麼仍然活著.
我人生的價值和意義在哪裡.
求天父你向我們深發出呼召.
讓我們有勇氣向你再次回應.
未信耶穌的.
要決志信主了.
信了你的.
無論你在哪個年齡層.
學習為主而活.
阿們.

\newpage



\section{以西結書 37:1-14}
\label{sec:031847}
\textbf{復活的生命}~~~~趙承恩宣教師
\newline
\newline
連結: \href{http://www.church.com.hk/acms/content.asp?site=cdc\&op=show\&type=product\&code=031847\&layout=sermon}{\texttt{031847}} ~~~~ 語音日期: 2026-03-22
\newline
\newline
\hyperref[sec:031846]{\small{< < < PREV SERMON < < <}}
~
\hyperref[sec:index_chronic]{\small{[返順時目]}}
~
\hyperref[sec:031848]{\small{> > > NEXT SERMON > > >}}
\newline
\newline
以西結書 37:1-14
\newline
\begin{longtable}{cl}
\hline
\hline
章節 & 經文 (和合本修訂版)\\
\hline
37:1 & \begin{tabularx}{0.7\textwidth}{X} 耶和華的手按在我身上。耶和華藉著他的靈帶我出去，把我放在平原中，平原遍滿骸骨。 \end{tabularx} \\ \\ \relax
37:2 & \begin{tabularx}{0.7\textwidth}{X} 他使我從骸骨的四圍經過，看哪，平原上面的骸骨甚多，看哪，極其枯乾。 \end{tabularx} \\ \\ \relax
37:3 & \begin{tabularx}{0.7\textwidth}{X} 他對我說：「人子啊，這些骸骨能活過來嗎？」我說：「主耶和華啊，你是知道的。」 \end{tabularx} \\ \\ \relax
37:4 & \begin{tabularx}{0.7\textwidth}{X} 他又對我說：「你要向這些骸骨說預言，對它們說：枯乾的骸骨啊，要聽耶和華的話。 \end{tabularx} \\ \\ \relax
37:5 & \begin{tabularx}{0.7\textwidth}{X} 主耶和華對這些骸骨如此說：『看哪，我必使氣息進入你們裡面，你們就要活過來。 \end{tabularx} \\ \\ \relax
37:6 & \begin{tabularx}{0.7\textwidth}{X} 我要給你們加上筋，長出肉，又給你們包上皮，使氣息進入你們裡面，你們就要活過來；你們就知道我是耶和華。』」 \end{tabularx} \\ \\ \relax
37:7 & \begin{tabularx}{0.7\textwidth}{X} 於是，我遵命說預言。正說預言的時候，有響聲，看哪，有地震；骨與骨彼此接連。 \end{tabularx} \\ \\ \relax
37:8 & \begin{tabularx}{0.7\textwidth}{X} 我觀看，看哪，骸骨上面有筋，長了肉，又包上皮，只是裡面還沒有氣息。 \end{tabularx} \\ \\ \relax
37:9 & \begin{tabularx}{0.7\textwidth}{X} 耶和華對我說：「人子啊，你要說預言，向風說預言。你要說，耶和華如此說：氣息啊，要從四方而來，吹在這些被殺的人身上，使他們活過來。」 \end{tabularx} \\ \\ \relax
37:10 & \begin{tabularx}{0.7\textwidth}{X} 於是我遵命說預言，氣息就進入骸骨，骸骨就活過來，並且用腳站起來，成為極大的軍隊。 \end{tabularx} \\ \\ \relax
37:11 & \begin{tabularx}{0.7\textwidth}{X} 他對我說：「人子啊，這些骸骨就是以色列全家。他們說：『看哪，我們的骨頭枯乾了，我們的指望失去了，我們滅絕淨盡了！』 \end{tabularx} \\ \\ \relax
37:12 & \begin{tabularx}{0.7\textwidth}{X} 所以你要說預言，對他們說，主耶和華如此說：我的子民，看哪，我要打開你們的墳墓，把你們帶出墳墓，領你們進入以色列地。 \end{tabularx} \\ \\ \relax
37:13 & \begin{tabularx}{0.7\textwidth}{X} 我的子民哪，我打開你們的墳墓，把你們帶出墳墓時，你們就知道我是耶和華。 \end{tabularx} \\ \\ \relax
37:14 & \begin{tabularx}{0.7\textwidth}{X} 我必將我的靈放在你們裡面，你們就要活過來。我把你們安置在本地，你們就知道我---耶和華說了這話，就必成就。這是耶和華說的。」 \end{tabularx} \\ \\
[1ex]
\hline
\hline
\end{longtable}
$^{1}$靚姊妹早晨.
不知道大家有沒有看過一套電影.
就是《少林足球》.
有沒有看過.
有啊,也有差不多的歷史了.
原來我看一套電影是二十多年前的了.
但是《少林足球》裡面有一句.
也蠻有名的說話.
就是周星馳說了一句.
人沒有夢想.
下一句是什麼.
和鹹魚有什麼分別.
夢想呢.
可能我們每一個人都曾經有過夢想.
可能是.
夢想可以成為我們一個方向.
令到我們可以一步一步地.
向著這個方向去走.
實現這個夢想.
但是對於我們基督徒來說.
夢想也可以是什麼.
可以是上帝給我們的意象.
在這個意象裡面.
讓我們可以為上帝去作供.
讓我們可以努力地.
為上帝給我們的意象而努力.
不過在現實.
就是夢想歸夢想.
意象歸意象.
現實也有現實的考量.
現實其實也有很多事情.
令到我們放棄了夢想.
放棄了上帝給我們的意象.
日本文學作家本鑒九雄.
有一句這樣的說話.
很多人都是三十歲就死了.
不過八十歲才埋葬.
不知道是不是大家的寫照.
因為當中裡面說到.
三十歲的時候.

$^{41}$就會因為很多不同的事情.
而開始放棄了夢想.
放棄了意象.
跟一條鹹魚一樣.
不過到八十歲的時候.
身體才死亡才埋葬.
可能我們在現實裡面.
我們做了很多的事情.
煩惱著很多的問題.
面對著很多的困難.
然後有一天發現.
原來我們在形形異異的生活.
很多的困難當中.
原來我們已經放棄了夢想.
放棄了當初.
上帝給我們的意象.
面對現實.
不再做任何的改變.
不要以為做教目.
做教目一定是在實現上帝的意象.
不一定的.
做教目我們也可以做到像廟祝一樣.
大家明不明白這個意思.
可以做到營運著一個堂會.
營運著每個禮拜有一個崇拜.
就完結了.
也可以是這樣.
今天的經課.
以西結書三十七章一至十四節.
其實正是在告訴我們.
當我們失去了對意象的追求的時候.
上帝就開始工作.
為的是叫我們不要停止.
為他所賜的意象去做功.
以西結書.
如果大家有看過以西結書.
其實是一個什麼樣的時代呢.
它是描寫著猶大國和以色列國已經亡國.
所有的國內精英都被擄到巴比倫.
而以西結先知.

$^{81}$也是被擄到巴比倫之後.
上帝就在巴比倫向以西結顯現.
其實一開始我們看以西結書第一章.
其實已經是一個很突破性的記載.
因為在古代中東世界的信仰觀裡面.
他們所信仰的神明.
是和當時民族的地結連在一起.
很少突破地域.
不過當以西結先知在巴比倫地.
見到上帝的顯現的時候.
其實就是在說什麼呢.
就是在說巴比倫的神明.
都不夠耶和華去到這麼突破.
甚至也突破了當時的以色列人對信仰的看法.
然後當上帝不斷呼召以西結.
去到說出不同的預言的時候.
預言就是先知的說話.
這個預言的方式是充滿著戲劇.
或者我經常這樣說.
就是以西結先知其實是一個戲劇大師.
他就用一個戲劇的方式.
去到表達上帝他的預言.
他的遺像.
而今天說的經文37章1至14節.
也是以西結書裡面很有名的枯骨復甦的遺像.
如果我們看第一節的時候.
我們就見到上帝的手按在以西結上面.
上帝的靈帶以西結去到平原的當中.
為什麼要去到平原呢.
平原呢 昨天林孫說到的時候.
也有說過可以是谷.
也可以翻譯做山谷的谷.
也可以是一個平原.
平原在古代裡面.
是一個軍隊對戰的主戰場.
多數都會在平原當中發生.
譬如我們看啟示錄.
有一場大戰叫做哈米吉多頓大戰.
其實米吉多是一個出名的平原.
專用來打仗的.

$^{121}$很常在米吉多這個平原打仗.
所以就有哈米吉多頓大戰的這個說法.
而平原呢.
很多時候都是軍隊對壘之間的一個地方.
所以當上帝帶以西結去到平原.
其實是帶了他去一個戰場的當中.
而按經文.
這個平原遍滿了骸骨.
而且極其枯乾.
就是說這些人已經死了一段時間.
如果是剛剛死的話.
他應該有血有肉.
但是如果他是死了一段時間.
他才會成為一個枯乾的骸骨.
以西結作為一個骷髏.
他被擄到巴比倫.
他一定知道上帝帶他去這個平原裡面.
是一個什麼意思.
就是這個平原是象徵著.
以色列人和巴比倫人打仗的平原.
而以色列就是戰敗的那一方.
那些骸骨就正正是以色列人.
這一班以色列的軍隊.
他們都戰死沙場.
葬身在這個平原的當中.
上帝帶以西結去到這個平原.
其實是去到一個以色列民不堪回首.
不敢去忘的一個地方.
是一個他們經歷失敗.
恥辱.
甚至是令他們絕望的一個地方.
不知道大家有沒有經歷過失敗.
我都經歷不少的.
我都經歷不少失敗.
也都或多或少.
我們都可能有一些不敢回望的經歷.
又或者我們都曾經感受到過一種絕望的時候.
不知道到現在為止.
大家會怎樣看待這些經歷呢.
可能我們會當是我們自己心的秘密.

$^{161}$收起來.
塵封在我們的心底裡面.
又或者可能是設法忘記了這些經歷.
但我們可能會發現越是想忘記就越是忘記不了.
但忘記就想自己好過一點.
也都有可能是當我們受過這些經歷打擊之後.
我們就被這些經歷影響了我們整個人.
整個人的人生.
我們不敢再做任何的改變.
甚至乎去改變任何實現意象的一些想法.
得過且過地去過活.
現實就是這樣的.
都是放棄.
讓我們可能成為一條鹹魚.
又或者我們看似活著.
但其實只是極其枯乾的骸骨.
但如果我們從經文裡面我們看得見.
上帝並不是想我們被過去的失敗.
受過的恥辱.
去束縛著.
而祂也帶領著我們和祂一起去面對這些經歷.
所以我們可以說.
救恩並不止於死後和上帝一起.
基督教信仰的救恩.
不止於死後和上帝一起過.
上天堂這麼簡單.
救恩也是將我們從過去的種種經歷去拯救出來.
我們回到經文裡面.
上帝問以西結.
這些骸骨能不能夠活過來.
於是以西結就說.
主耶和說:你是知道的.
接著上帝就叫以西結.
向這些骸骨說預言.
骸骨與骨彼此相連.
骸骨上面有根長了肉.
再包上皮.
從骸骨重新成為一個人形.
這個人形就有點像創世記.
創世記裡面上帝用泥土做人.

$^{201}$不過只有一個人形.
他沒有任何的氣息.
所以按著以色列人的信仰.
或者按著我們基督教的信仰.
一個完整的人並不止於只有肉體.
更加有上帝的氣息在人的裡面.
如果我們在希伯來文的原文當中.
氣息可以翻譯成甚麼呢.
可以翻譯成風.
也可以翻譯聖靈.
所以當我們去看這段經文的時候.
我們會發現又有風又有氣息又有靈的時候.
我們就會知道其實這段經文.
是一段充滿著活力.
充滿著生氣的一段經文.
所以說氣息,風和靈.
可以說是生命.
上帝令到生命重新注入去這些骸骨.
變成肉體.
再從這些生命裡.
再為這些肉體注入生命.
令到骸骨重新成為活人.
而且經文裡面還說到用腳站起來.
我們可能年長的弟兄姊妹就會知道.
用腳站起來是需要一定的力氣.
但是如果我們說這段經文.
這段用腳站起來.
其實還要成為極大的軍隊.
是怎樣呢.
其實就好像這樣.
夠力嗎.
聽到吧.
不知道後面的弟兄姊妹聽不聽到.
聽到吧.
應該沒有跌的.
是用腳站起來成為極大的軍隊.
所以上帝的氣息是令到那個人.
充滿著生命力,充滿著力量.
成為一支軍隊.
這支軍隊是要做什麼.

$^{241}$這支軍隊是要以色列的全家.
他們是被上帝帶領離開墳墓.
兼且他們有一個意象.
就是重新回到以色列.
在當時是從被擄的地方.
重新回到以色列.
上帝是令到他們不再說出那個.
看啊,我們的骨頭枯乾了.
我們的指望失去了.
我們滅絕正盡了.
上帝令到他們不再說這句話.
上帝令到他們不再被過去束縛著.
被絕望籠罩著.
而是令到他們成為充滿著生命力.
有盼望,有意象.
重返以色列地的軍隊.
這支軍隊是以色列全家.
上帝藉著祂的話,祂的靈和祂的應許.
去醫治以色列人過去的傷痛.
也賜予了他們將來的盼望.
重新復興他們.
我們在這裡可以看到.
一個復活的生命.
真的不只是將來在天國裡面.
而是在醫治著我們過去.
令到我們不再受束縛.
也令到我們能夠有力.
有一份力量去實踐上帝給我們的意象.
不知道鄧英姐妹.
你有沒有上帝給你的意象.
可能那種意象是.
你怎樣努力去愛你身邊的人.
可能那種意象是.
你很想去改變一些現實的不公義的事情.
可能那個意象是.
你很想去服侍教會,服侍上帝.
又或者可能那個意象是有不同種種的.
我不知道大家有沒有這些意象.
或者可能曾經有過這些意象.
但是你有沒有放棄過這些意象.

$^{281}$或者是不是放棄了這些意象.
可能放棄的原因可以很多.
可以是過去失敗,過去的經歷.
不過我想說.
在這段經文裡面.
正正就是提及到上帝的靈.
是賜給我們生命力.
去克服過去對我們的傷害.
以及讓我們有一個指望.
去為那個意象去努力.
在大齋期的時間.
我們一方面可能在預備我們的新聖靈.
使我們可以好好預備迎接.
耶穌基督的受苦和復活.
我想說耶穌基督的復活.
是作為基督信仰的核心.
所以從來都是這樣說的.
就是復活節是大過聖誕節的.
雖然復活節的假期就比聖誕節的假期短.
不過復活節是大過聖誕節的.
因為沒有復活.
我們整個基督信仰就是沒有的.
但是復活作為基督信仰的核心.
不只指向將來.
復活也涉及到我們的過去.
令到我們能夠有一個從上帝而來的目光.
去看待過去種種的經歷.
去給上帝醫治.
去重新賦予一份生命力給我們.
過去的失敗絕望.
反而能夠成為上帝使用祂施恩的工具.
我們有沒有一些地方傷痛.
是不願意被上帝觸碰和醫治的呢?.
我們不想回望的呢?.
我們不妨在這個大齋期當中.
我們去開放我們的內心.
去給上帝醫治觸碰這些傷痛.
不只觸碰而是給上帝轉化.
令到我們不再受這些經歷束縛.
同一時間.

$^{321}$今天的講題是復活的生命.
復活的生命不再是一條鹹魚.
我會說是一條鮮魚.
也可以這樣說.
我們是一個有力.
特地的一個有力的軍隊.
我們不再只是為生存而活.
而是叫我們能夠擁抱異象.
擁抱上帝給我們的異象.
有力地站起來.
為上帝賜給我們的異象而活.
等著我們一同禱告.
因為你藉著先知的異象.
在你的話語.
在你的聖靈.
在你的應許當中.
使我們能夠有復活的生命.
讓我們能夠勇敢面對過去.
有力地活在現在.
有盼望地面對將來.
父上你是賜生命和異象的主.
求你使我們能夠擁抱你賜予的異象.
成為一個滿有生命力的基督跟隨者.
讓我們為你的異象而活.
努力地跟隨基督.
我們的祈禱是靠那位死而復活.
而且最終得勝的耶穌基督的名.
\newpage



\section{創世記 41:16}
\label{sec:031848}
\textbf{聖經人聖經事:這不在乎我}~~~~黃克勤牧師博士
\newline
\newline
連結: \href{http://www.church.com.hk/acms/content.asp?site=cdc\&op=show\&type=product\&code=031848\&layout=sermon}{\texttt{031848}} ~~~~ 語音日期: 2026-03-19
\newline
\newline
\hyperref[sec:031847]{\small{< < < PREV SERMON < < <}}
~
\hyperref[sec:index_chronic]{\small{[返順時目]}}
~
\hyperref[sec:031849]{\small{> > > NEXT SERMON > > >}}
\newline
\newline
創世記 41:16
\newline
\begin{longtable}{cl}
\hline
\hline
章節 & 經文 (和合本修訂版)\\
\hline
41:16 & \begin{tabularx}{0.7\textwidth}{X} 約瑟回答法老說：「這不在乎我。神必應允法老平安。」 \end{tabularx} \\ \\
[1ex]
\hline
\hline
\end{longtable}
$^{1}$各位姐妹又讀聖經,一起看看在聖經的人物與故事後面上帝是怎樣的.
今日都是生命傳書的網絡為大家準備,去到差不多四月的尾.
借不在乎我就是今日的主題,借不在乎我.
不是我叻,不是在乎我多巴閉,是在乎上帝的.
創世記四十一章,昨日仍然是四十章,現在是四十一章.
過了兩年,第一點,嘩,時間很快的.
又一眨眼,一揭一頁紙就過了兩年了.
過了一章就過了兩年,我還記得兩年前怎樣.
還記得小朋友的時候,幼稚園等家人來接我放學的時候.
心急如焚,是怎樣?眨眼不是兩年,是一生都差不多過完了.
第一,過了兩年都很久了,真的不記得那些事了.
法老做夢,所以上一兩章的事和今日這一章,四十章.
和第四章是有些關係的,是在說夢一直延續下來.
過了兩年,法老做夢,夢見自己站在河邊.
特別說河,甚麼河呢?通常我們都認為是尼羅河.
因為埃及最重要的河流就是尼羅河.
每年尼羅河都是為他們帶來耕種等等.
特別是尼羅河三角洲是非常肥沃的,近出海的部分.
因為那些去到很平地,於是河水就會分散流開.
不是一條河,而是分散滲入海裡,流入海裡.
很多溪流,很多河流,細河.
就這樣令到那裡的尼羅河三角洲變得很肥沃.
而尼羅河有一個特別,就是每一年都會有一次氾濫.
他們很歡迎這個氾濫,就是一氾濫的時候,大家搬走.
一氾濫完,水一路退,他們就一路搬回去.
為甚麼呢?因為被水浸過的泥土很肥沃.
所以他們對這個氾濫很歡迎.
站在河邊,有七隻母牛由河裡面上來.
即是說河是令到孕育了七隻母牛.
是從尼羅河上來的,又美好又肥壯.
沒有母牛,沒有母牛有甚麼好處呢?.
有奶飲,有牛仔肉,還有自己本身,還可能耕田.
於是就說好呀,又美好又肥壯.
就在這個蘆迪中吃草,一片美麗與豐富肥沃景象.
誰知夢裡隨後有七隻母牛從河裡面上來.
都是這條河,估計是尼羅河.
誰知車又醜陋又乾瘦,為甚麼會這樣呢?.
於是這樣看法,就是尼羅河可能缺水了.
於是就滋潤不到大地,於是就會乾又瘦.
其實動物生蓄不會這樣,沒有水,也沒有草.

$^{41}$沒有食物,那就乾又瘦.
就和那七隻母牛一起站在河邊.
這樣也有七隻肥呀.
誰知又醜陋又乾瘦的那七隻母牛.
吃盡了美好又肥壯的那七隻母牛.
嚇到發怒紮醒.
你不要問牠怎樣吃牛,牠的口那麼小.
牠又不吃肉,不要問那些.
發夢吧.
就羞牛醜牛吃了七隻美好的牛.
嚇到發怒紮醒.
嘩,怎麼會這樣呢?.
嚇得我心肝都跳.
幸好,嚇一嚇之後又睡了.
又來惡夢,又睡著了.
第二回做甚麼夢呢?.
就夢見一火麥子長了七個水子.
又肥大又尖尾.
留意一下,是一火麥子,即是一條麥.
一條麥上去有些麥水.
麥水即是一串的麥子有殼.
有重量,所以那些麥子一串有重量.
會這樣凹的.
這樣就有七條水子的.
一支麥子麥乾有七條水子.
很肥啊.
七個水子那麼多.
又肥大又尖尾.
就是這樣.
這個大家知不知道清不清楚.
我都在google剪下來.
這個麥水的麥子.
大麥的麥粒是大粒很多.
很緊密的.
而且大麥有很多很長的鬚.
下一張看得清楚一點.
正好一張是乾的麥子.
很多時候要插進去做裝飾.
這樣就不能用了.
這些就是水子.

$^{81}$左右左右的麥子.
很長的鬚.
多數是一支.
都不會有很多七條麥水.
這一條麥水.
一條麥水.
為什麼這些那麼青.
不是黃的,青即是未熟.
有麥水,不過未熟.
裡面還是軟軟的.
還是翼翼的.
要在秋天乾一點收水.
曬久一點,熟了.
再久一點,一落秋雨又浸爛了.
所以要及時收乾.
不收乾就沒有了.
所以聖經都說.
去有莊稼多,收稼的人少.
一片田地.
莊稼熟了.
你不收得切.
落秋雨浸了之後.
濕透水就會爛了.
不能吃了.
綠色我特地想.
在這張Google找到給大家看看.
剛剛如果是三月份.
或者四月份過完魚節.
魚節是星期四晚.
吃魚節晚餐.
星期五,即是星期四的魚節晚餐.
也是最後晚餐,主耶穌.
然後星期五主耶穌釘十字架.
星期六,星期日.
就是復活主日.
而復活主日很多時候.
會慶祝初熟節.
剛好就是.
多倫多前書十五章第二十字.
主耶穌成為睡鳥之人.

$^{121}$初熟的果子.
而那個復活節主日.
也是初熟節.
這麼應驗的.
初熟是什麼呢?.
初熟的意思就是.
有很多還是綠色的墨水.
綠色的.
即是說它已經出來了.
不過還未熟.
要再過多五十天左右.
五旬.
就是七七節的時候就收割了.
已經熟了.
這個就是濕了的.
初熟節用的.
見到的時候的墨水.
好了,全部是綠色的.
初熟節就是慶祝這件事.
全部是綠色的.
初熟這個字就是綠的.
這個意思.
好了.
隨後長了七個水子.
又細弱又被東風吹招.
東風哪裡來呢?.
你想像一下埃及地圖.
東風即是說在西乃曠野那邊.
吹過來的熱風.
乾燥兼炎熱.
這樣吹,墨紙就會乾了.
燒焦.
燒焦當然了.
五十度攝氏.
我們去聖地的時候.
如果夏天去.
七八月的人去.
即是五十度.
最少.
不是很大汗.

$^{161}$因為很乾.
所以一出汗就沒有了.
很辛苦的.
於是長得又不好.
再加上被東風吹招.
那些細弱的水子.
竟然吞了七個又肥大又飽滿的水子.
又嚇到法老醒了.
哎呀,不了,是個夢.
幸好是個夢.
好了.
到了早晨.
法老就答應說沒有睡.
心裡不安.
馬上派人照了埃及所有術士和博士來.
術士呢.
就是玩一些魔術.
或者叫奇門遁甲那些.
他們是有多少特別.
或者依靠邪靈等等而有的能力.
博士呢.
如果在當時.
譬如就算是耶穌時代的博士.
其實應該可以譯作占星家.
或者是天文學家.
識看星座等等.
這個呢.
就是講一下星座的運程.
這個就是講一下.
一些很奇妙的東西.
我們不在這裡講一句真不真.
哎呀,這些東西真不真實的不知道.
當時的人是相信的最重要.
法老就將所有的夢都告訴術士和博士.
解釋給我聽.
你知道那些王一不懂得解釋.
就叫了所有懂得解釋的人來.
誰知道沒有人能給法老完解.
啊,皇上你這個解釋呢.
那隻瘦牛肥牛是什麼.

$^{201}$那些墨水呢就怎樣怎樣.
就解釋了.
誰知道聽完之後.
法老覺得都不是這樣.
有些不妥.
不是這樣的.
不是這樣的意思.
不是的.
這樣就覺得不是.
就是不能夠圓滿地解釋.
不能圓解.
這個時候.
酒正就對法老說.
哎呀,第二點.
我今天想起我的罪.
是什麼呢.
就是兩年前.
法老呢.
你法老臣服的時候.
將我和善長擺在護衛長府內監獄.
我們二人同夜各做一夢.
各夢都有講解.
就是那裡有一個.
希伯來的少年人.
是護衛長的僕人.
這個僕人的字是指奴隸.
我們就把夢告訴他.
他就將我們的夢圓解.
解得很通.
而真是應驗.
於是就按我們的夢來解.
後來就真是成就了.
我官服完職.
而他善長就被掛起來.
法老就說.
嘩,這麼厲害.
我想你留意這一點.
剛才一直說.
回到第一點.
第一點是過了兩年.

$^{241}$久不久.
很久了.
不是的,很快,兩年.
如果你坐牢.
坐牢就久了.
日日坐牢.
坐兩年很久.
而且那兩年真是.
為什麼還沒有人來救我呢?.
那個酒精就沒有來救我.
善長就應該死了.
三天.
酒精是官服完職的.
我叫他紀念我都沒有.
這樣等了兩年.
於是你以為人是會忘記.
誰知道神沒有忘記.
誰知道就是.
過去一陣子.
到法老.
發夢.
兩年之後法老發夢.
於是酒精就說.
哎呀,我真是糟糕了.
我醉的意思就是.
我真是不好了.
我真的醉過醉過.
我忘記了跟別人說好說話.
什麼好說話,到你說什麼好說話.
這樣的皇上.
就是這樣這樣這樣.
就有個人幫我們解了這個夢.
我今天才想起.
我忘記了他.
兩年了.
法老遂即差人.
去召約室.
馬上在家裡叫他出來.
那些人急忙跑去.
波提佛家裡面.

$^{281}$波提佛應該都知道.
馬上飆了他,第一個.
回去了,幫他剃頭.
刮臉換衣裳.
應該駕沖涼.
剃頭做什麼呢?.
因為其實你見到.
很多埃及人.
他們全部都是剃光了.
好像我一樣光了.
不過會在某個地方留一條邊.
很多東西留一條邊.
一扎就扎一條邊.
馬上.
跟他說.
監裡面沒有剃鬚.
剃乾淨了,很乾淨.
埃及人很乾淨.
好了,就換了衣服.
進入法老面前.
法老對約瑟說,我做了一個夢.
沒有人能夠解,我聽說你能夠解.
好了.
約瑟回答說,第三點.
今日的主題.
這不在乎我.
是神與張平說的話.
回答法老.
真厲害.
後面這句.
就是.
很世故的說話.
你懂得解夢.
我解得很好.
皇上有什麼說話.
是神會將平安的說話.
跟皇上說.
這句是一個.
安慰與奉承的說話.
比較多.

$^{321}$是會說話.
神一定會跟法老皇上說.
安慰的說話,平安的說話.
不過其實.
我想說,這不在乎我.
因為不是我.
能夠解你的夢.
夢是上帝給的.
夢在乎神.
夢出於神.
出於神.
然後.
不在乎我.
所以這幾句要說的.
夢出於神.
人的生命高低.
V形還是直情插水.
在於神.
最重要有神跟我們一起.
第二.
你要知道.
在乎不.
不在乎我.
在乎神.
於是他就開始解夢.
留待明天再說.
這不在乎我.
這五個字又送給大家.
這不在乎我.
這五個字又送給大家.
這不在乎我.
這五個字又送給大家.
\newpage



\section{撒迦利亞書 4:1-14}
\label{sec:031849}
\textbf{創始成終的三一神}~~~~黃梅玲牧師
\newline
\newline
連結: \href{http://www.church.com.hk/acms/content.asp?site=cdc\&op=show\&type=product\&code=031849\&layout=sermon}{\texttt{031849}} ~~~~ 語音日期: 2026-03-22
\newline
\newline
\hyperref[sec:031848]{\small{< < < PREV SERMON < < <}}
~
\hyperref[sec:index_chronic]{\small{[返順時目]}}
~
\hyperref[sec:031850]{\small{> > > NEXT SERMON > > >}}
\newline
\newline
撒迦利亞書 4:1-14
\newline
\begin{longtable}{cl}
\hline
\hline
章節 & 經文 (和合本修訂版)\\
\hline
4:1 & \begin{tabularx}{0.7\textwidth}{X} 那與我說話的天使又來叫醒我，好像人睡覺時被喚醒一樣。 \end{tabularx} \\ \\ \relax
4:2 & \begin{tabularx}{0.7\textwidth}{X} 他問我：「你看見甚麼？」我說：「我看見了，看哪，有一個純金的燈臺，頂上有燈座，其上有七盞燈，每盞燈的上頭有七根管子； \end{tabularx} \\ \\ \relax
4:3 & \begin{tabularx}{0.7\textwidth}{X} 旁邊有兩棵橄欖樹，一棵在燈座的右邊，一棵在燈座的左邊。」 \end{tabularx} \\ \\ \relax
4:4 & \begin{tabularx}{0.7\textwidth}{X} 我問與我說話的天使說：「主啊，這是甚麼意思？」 \end{tabularx} \\ \\ \relax
4:5 & \begin{tabularx}{0.7\textwidth}{X} 與我說話的天使回答，對我說：「你不知道這是甚麼意思嗎？」我說：「主啊，我不知道。」 \end{tabularx} \\ \\ \relax
4:6 & \begin{tabularx}{0.7\textwidth}{X} 他回答我說：「這是耶和華指示所羅巴伯的話。萬軍之耶和華說：不是倚靠勢力，不是倚靠才能，乃是倚靠我的靈方能成事。 \end{tabularx} \\ \\ \relax
4:7 & \begin{tabularx}{0.7\textwidth}{X} 大山哪，你算甚麼呢？在所羅巴伯面前，你必夷為平地。他安放頂上的那塊石頭，人就歡呼：『願恩惠、恩惠歸與這殿！』」 \end{tabularx} \\ \\ \relax
4:8 & \begin{tabularx}{0.7\textwidth}{X} 耶和華的話臨到我，說： \end{tabularx} \\ \\ \relax
4:9 & \begin{tabularx}{0.7\textwidth}{X} 「所羅巴伯的手立了這殿的根基，他的手也必完成這工，你就知道萬軍之耶和華差遣我到你們這裡。 \end{tabularx} \\ \\ \relax
4:10 & \begin{tabularx}{0.7\textwidth}{X} 誰藐視這日的事為小呢？他們見所羅巴伯手拿石垂線就歡喜。這七盞燈是耶和華的眼睛，遍察全地。」 \end{tabularx} \\ \\ \relax
4:11 & \begin{tabularx}{0.7\textwidth}{X} 我問天使說：「那麼在燈臺左右的這兩棵橄欖樹是甚麼意思呢？」 \end{tabularx} \\ \\ \relax
4:12 & \begin{tabularx}{0.7\textwidth}{X} 我再次問他：「這兩根橄欖樹枝在兩根流出金色油的金嘴旁邊，是甚麼意思呢？」 \end{tabularx} \\ \\ \relax
4:13 & \begin{tabularx}{0.7\textwidth}{X} 他對我說：「你不知道這是甚麼意思嗎？」我說：「主啊，我不知道。」 \end{tabularx} \\ \\ \relax
4:14 & \begin{tabularx}{0.7\textwidth}{X} 他說：「這是兩位受膏者，侍立在全地之主的旁邊。」 \end{tabularx} \\ \\
[1ex]
\hline
\hline
\end{longtable}
$^{1}$各位弟兄姊妹平安.
我想先問一個問題.
大家知不知道香港人的鄉下在哪裡?.
日本.
我通常一年都會去一次旅行.
都會選擇日本這個地方.
我都承認是我的鄉下.
上年去旅行的時候我去看煙花.
我不知道大家平時看煙花是不是在香港看的呢?.
如果我在香港看 我通常都是看電視機的.
因為太多人了 在現場看就太擠逼了.
所以都沒有試過在現場去看煙花.
剛好上次去旅行可以有一個機會在現場看煙花.
這些就是那時候我拍的照片.
因為很多人 所以都會預早一點去.
但是預早一點 很多人都已經霸了一些好位置.
要預備去看煙花.
大家都很期待著這場煙花的開始.
本來是漆黑一片的天空.
一開始的時候除了有聲音「磅磅磅」之外.
就是看到很漂亮的.
這些繽紛色彩的煙花出現在漆黑的夜空上面.
真的很漂亮.
我由一開始的「哇」.
一直都是「哇哇哇」 「哇」到它完結的時候.
其實越叫都會越大聲.
因為真的越來越精彩.
大家都知道 通常開始的時候都是簡單一點.
沒有那麼多的色彩.
不過越來越精彩 畫面越來越豐富.
每一個煙花的爆發.
其實我自己心裡就不知道下一個煙花是什麼樣子的.
所以每一次的爆發都是每一次的驚喜.
所以整個過程都是很開心.
好像接了很多很多份的禮物.
一路打開就很開心地看到每一個的驚喜.
我不知道大家有沒有這樣的經驗呢.
今天這段經文就是《選自撒迦利亞書》.
上年我有好幾次都是選了當中的經文來跟大家分享.
不知道大家記不記得.

$^{41}$《撒迦利亞書》裡面一共有多少個異象呢.
記得 好像聽到有八個字.
一共有八個異象.
這八個異象大家記不記得是出現在什麼時間的呢.
根據《一章八節》裡面所講.
第一個異象一開始就是說.
就是說這八個異象一口氣接連地出現.
呈現在先知撒迦利亞的眼前.
其實就是在一個漆黑的夜空裡面出現.
一個又一個為他帶來驚喜 帶來震撼的異象.
就好像我剛纔所講的看煙花.
一個接著一個很驚喜 很震撼的景象.
出現在先知撒迦利亞的眼前.
每一個異象對他來講都是全新的.
他從來都沒有見過.
所以每一個異象都需要由神的使者來向他講解.
解釋每一個異象的內容其實是什麼意思.
相信每一個異象都令到先知感到很驚訝.
因為神是一步一步來向先知去揭示.
神的應許和神的計劃.
而來到今天這段經文第四章裡面.
其實已經是第五個異象.
如果八個異象的話.
中間的部分就是第四和第五個異象.
事實上第四和第五個異象.
就是這八個異象裡面的種中之種.
第四個異象是關於神親自來到去結證大祭司約舒雅.
因為他身上充滿著污穢.
他穿著一件很髒的衣服.
這個形象令到先知撒迦利亞很震撼.
因為沒有想過.
一位本應是聖潔的大祭司竟然是全身的污穢.
這個結證的行動由神親自來到去進行.
代表著神親自去向他所愛的子民.
來到去施行大能的拯救.
除去他們以往所犯的一切的罪.
一切得罪他的地方.
他都日不嘔笑.
這個是極其震撼的景象.
出現在先知的眼前.

$^{81}$如果這樣說.
第四個異象已經那麼精彩.
那麼震撼.
接下來的第五個異象又是怎樣的呢?.
是關於什麼呢?.
讓我們在進入第五個異象之先.
我們一起同心禱告.
親愛的天父上帝.
我們再一次來到禰的根祠.
要去聆聽禰自己的話語.
願禰的靈親自充滿.
我們眾人的心.
讓我們能夠聽得明白.
放在心裡 遊行出來.
以上禱告.
奉主耶穌基督的聖名祈求.
阿們.
其實這段經文第四章.
可以分作三個部分.
第一個部分是一至五節.
講述到金燈臺的異象.
這是上半部分.
然後下半部分出現在十一至十四節.
這是下半部分.
中間加插了六至十節.
就是對這個異象的解釋.
也就是說.
如果我們要看整個金燈臺的異象完整地看.
就要先看一至五節.
然後再跳到十一至十四節.
這就是完整的第五個異象.
為什麼在這個異象當中.
會特別要將它拆開.
中間加入六至十節這個解釋呢.
其實在舊約裡面很多時候都會出現這個情況.
我相信大家有喫三文治的.
應該有人喜歡喫漢堡包的.
你們喫三文治 喫漢堡包.
你們最想喫的是什麼.
中間那個餡.

$^{121}$就是餐蛋.
或者是漢堡包裡面的那塊牛肉.
這個其實就是一個三文治.
重要的部分就是中間的部分.
六至十節那個異象的解釋.
我會先跟大家看第一部分和第三部分.
完整的這個異象的內容.
在第一節裡面講到.
那個跟我說話的天使又來叫醒我.
好像人睡覺被喚醒一樣.
這裡說的是那個天使.
就好像要拍醒仙姿撒迦利亞.
為什麼 因為她好像睡著了.
於是就要叫醒她.
我們剛纔說很奇怪.
明明說接一個異象的.
為什麼看看仙姿會睡著.
其實大家就要去思考到.
其實前面那四個異象.
對仙姿來說是一個很震撼的情景.
在那個漆黑的晚上.
她看到一個又一個她從來沒有看過的圖畫.
對她來說其實很震撼.
她要去消化去思考.
去明白掌握天使所講的解釋的內容的時候.
其實她要花很多很多的心力.
她要很專心去想.
可能就在這一刻.
去到這個異象發放的中間這個時間.
她真的已經心力俱疲.
去到這個情況.
於是她好像睡著了一樣.
半夢半醒的一個情況.
但是不行 不可以停下來.
接著就是第五個異象要出現.
所以天使立刻就要去拍醒她.
要她快點起來去看這個第五個異象.
第五個異象裡面其實出現了一個物品.
就是一個純金的燈臺.
這個就是金燈臺.

$^{161}$究竟什麼是金燈臺呢?.
如果我們看回昌一級的二十五章.
三十一至三十九節裡面.
其實是很清楚描述到.
關於究竟一個金燈臺其實是怎樣來到它被製造.
也很詳細地描述了外形是怎樣的.
大概就是在圖中看到的這個金燈臺.
這個金燈臺其實是出現在會幕裡面.
也是出現在後期的聖殿當中.
是用來作一個照明之用.
是要去照亮聖鎖的.
於是我們就去讀.
《昌一級二十五章》這段經文的時候.
我們會發現和《撒迦利亞書》裡面描述.
那個金燈臺其實是有出入的.
不是一樣的.
當中不同的地方其實是最重要的部分.
因為這樣就能夠讓我們去解釋這個異象.
不同的部分究竟在哪裡呢?.
第一個就是講第二節的話.
頂上有盞燈.
第二個的不同就是講每盞燈有七個管子.
第三個的不同就是講到這個金燈臺的旁邊.
兩旁左右是分別有這個橄欖樹.
還有第四個的不同.
就是在第十二節裡面提及到的.
有兩根橄欖枝.
在兩個流出金色油的金嘴旁邊.
這四個不同之處就是和《出埃及記》裡面.
金燈臺的描述有一個出入的地方.
如果這樣講的話.
講到這裡大家會不會很混亂呢?.
究竟在講什麼呢?.
事實上我們不看圖真的很難明白.
我就嘗試找了一幅圖.
去幫大家明白其實究竟異象裡面的金燈臺是什麼樣子.
其實這個是一個概念圖來理解.
我們發現原來這個金燈臺.
其實是建基於舊約聖殿裡面的金燈臺.
大家想象一下.

$^{201}$先知撒迦利亞見到的這個金燈臺.
和那時候在聖殿裡面見到的金燈臺是不一樣的.
但是又似曾相識.
因為紅色框住的部分其實就是.
和《出埃及記》裡面的金燈臺一樣.
也就是在聖殿裡面的金燈臺是一樣的.
但是除了這個紅色框住的部分之外.
其餘的部分全部都是嶄新的.
全部都是先知撒迦利亞第一次見到的一個景象.
大家想象一下.
先知撒迦利亞.
他以前當然是知道這個聖殿裡面的金燈臺是什麼樣子的.
當然是知道這個聖殿是什麼樣子的.
但是在南國猶大 祕魯 亡國之後.
聖殿被毀了.
當然聖殿裡面這個金燈臺.
都已經是不翼而飛.
不知去了哪裡.
無論是先知撒迦利亞.
還是那一幫祕魯去到巴比倫的猶大人.
他們從那天亡國開始.
他們就沒有再見過昔日那個輝煌的聖殿.
也沒有再見過這個明亮的金燈臺.
但是在這個漆黑的夜晚.
來到第五個異象出現的時候.
神向他呈現的這幅極致震撼的圖畫.
就是一個他似曾相識的金燈臺.
但是規模是大了很多.
而且是和以前那個完全不一樣的.
不一樣的地方包括了那兩棵的橄欖樹.
還有那個金嘴.
也就是圖上面的金道館.
還有那個燈盞 還有那七十九條的管子.
這些全部都是新的.
一個新的圖像出現在先知撒迦利亞的面前.
所以我們就明白到為什麼.
先知要三次地發問.
究竟是什麼意思呢.
他眼前所見到的景象.
他的發問分別出現在第四節.

$^{241}$還有第十一節 還有十二節.
他真的不明白.
因為這些全新的部分是他從來沒有見過的.
究竟是什麼意思呢.
天使反問他兩次.
你不知道是什麼意思嗎.
先知必須要承認他的無知.
他在使者的面前答主:我不知道.
他真的不知道.
他很想明白這個震撼的意象.
這個在他眼前呈現的金燈臺.
到底是什麼意思.
在十四節裡面 使者終於回答.
這是兩個受高者站在普天下主的旁邊.
原來這兩棵橄欖樹就代表著兩個受高者.
受高者的意思在原文裡面是指高郵之子.
這個高郵因為是那棵橄欖樹所流出來的.
所以是一個新鮮的郵 是新郵.
在舊約裡面都會以新郵作為一個祭來呈獻給神.
所以受高者將這個新郵提供給金燈臺.
來點亮金燈臺裡面的七盞燈.
這兩個受高者透過金咀 也就是金道館.
將新郵流入這個燈盞裡面.
也就是戰咀的方向.
然後再透過每盞燈的七條管子流入那七盞燈當中.
也就是說金燈臺的燈.
是被這些新鮮的橄欖油所點亮的.
以至能夠發光.
照亮當時漆黑一片的天空.
在漆黑一片的天空裡面.
竟然出現了一個非常明亮的金燈臺.
這是一個很震撼的場面.
剛纔我們說先知撒迦利亞本來是很累.
半夢半醒.
但是突然被使者拍醒.
看到這個景象.
明亮的金燈臺出現在他的眼前的時候.
是震撼人心的一個情景.
使先知立刻睡了一覺醒過來.
他完全清醒過來.

$^{281}$他很想知道究竟這個異像是什麼意思.
其實這兩個受膏者.
如果根據第四個異像和第五個異像裡面提及的人物.
就是分別為大祭司約書亞和所羅巴伯.
這兩位就是這兩個異像裡面的重要人物.
這兩位人物其實就是在《哈幹書一章》14字裡面所提及的.
是兩位帶領著猶大百姓去重建聖殿的領袖.
兩個很重要的人物.
這兩個人物竟然出現在這個異像當中.
作為橄欖樹提供新油給這個金燈臺.
來備燃點 備點亮.
究竟是什麼意思呢.
我們必須要回到歷史當中.
我們還記得當我們談論到這個第一個異像的時候.
其實都講述過提及過.
這一幫猶大人他們回去耶路撒冷聖城.
他們要重建聖殿的時候其實遇到很多很多的困難.
包括了受敵的難阻.
包括了他們資源的貧乏.
包括了百姓對於這樣的狀況底下要建一個聖殿.
一定不會及得上以前那個輝煌的聖殿.
他們覺得很失望 很灰心.
甚至是輕看了這個重建聖殿的工作.
以至到他們停了16年的時間.
當他們建好了那個聖殿的根基.
他們就放下了這個工作.
一放下就停了16年的時間.
要去到公元前520年.
纔再次在神的呼召底下開始重建聖殿的工程.
就在第一個異像當中.
神這樣來說.
祂說「現今我回到耶路撒冷.
凝思聯吻我的殿 必重建在其中.
準成必拉在耶路撒冷之上.
這是萬君之耶和華說的」.
在第一個異像當中.
神正式宣告祂迴歸聖城耶路撒冷.
祂確認聖殿必定會重建在其中.
因為有聖殿就是代表有神的同在.
那所羅巴伯的角色是什麼呢?.

$^{321}$我們就要回到剛纔說的三文治最重要的「銜」.
6至10節裡面的異像的解釋.
一開始說這是耶和華指示所羅巴伯的說話.
神指示所羅巴伯什麼呢?.
當中還說了三個重點.
第一就是在第七節.
「大山啊 你算什麼呢?.
在所羅巴伯面前你必成為平地」.
大山也就是你.
其實就是指在重建聖殿的過程裡面.
所面對的種種困難.
而這些困難最後大山也會變成平地.
本來是一座一座的高山.
好像一個一個的困難.
是難了這群猶大人重建聖殿.
但是神應許這一座又一座的大山.
一個又一個的困難.
神一定會親自除去.
以至到所羅巴伯什麼都不用做.
他可以繼續專心參與在重建聖殿的工程當中.
第二個重點是.
神說他必搬出一塊石頭安在殿頂上.
這裡的「他」是指所羅巴伯.
神應許所羅巴伯一定會將聖殿最後的一塊石頭.
擺在殿頂的上面.
意思就是說.
我們現在建房子.
平頂就是完成的建築.
這裡其實也是一個這樣的意思.
當這塊石頭安放在殿頂的時候.
就是代表著重建聖殿的工程是完成了.
就是一個這麼辛苦艱難的過程之後.
這個聖殿是會被重建完成.
因此人要大聲的歡呼.
說願因為因為歸於這殿.
或者是歸於這石.
就是這一塊最後平頂的石頭.
當所羅巴伯親手將聖殿最後一塊石頭.
安放在殿頂的時候.
就是代表著聖殿重建完成.

$^{361}$在這個時候所有的猶大百姓都能夠歡呼快樂.
接著在第九節.
神繼續說所羅巴伯的手拿了這殿的根基.
他的手也必完成這功.
這是第三個重點.
一方面要指出.
所羅巴伯是帶領猶大百姓重建聖殿的一個先鋒.
他開始重建聖殿的工程.
但他不只是開始.
他的手也會完成整個的工程.
換言之就是由開始到終結完成的那一刻.
整個的過程都是透過所羅巴伯來完成.
由開始到沒了.
聖殿重建的工程.
都是在所羅巴伯的手裡面.
到了第十節.
神就問了這個問題.
誰夠膽表示這一天的事為少呢?.
為什麼神要作這個反問呢?.
剛纔不是說有很多的艱難嗎?.
原來在哈該書2章3節裡面有這樣的提到.
有誰見過這個殿從前的榮耀呢?.
現在你們看看.
豈不是好像在眼中看如無有嗎?.
在這裡很清楚.
將從前聖殿的榮耀.
和這個重建了的聖殿的樣子作了一個比較.
從前的聖殿和現在的聖殿做比較的時候.
其實真的沒得比.
是一個很榮耀,很華美,很宏偉的聖殿.
和一個又小又不夠漂亮.
什麼都比起第一個差很遠.
描述是"無有",Nothing.
榮耀和無有的落差實在太大了.
大家想不想象到這群猶大的百姓.
其實他們期待著迴歸耶路撒冷聖城.
他們是看著廢墟.
他們很想回想從前那個宏偉的第一聖殿.
他們很想回想起從前那個很繁華的耶路撒冷城.
這個就是這一群人他們的心聲.

$^{401}$所以他們可能會將這個重建聖殿的工程提輕了.
但是神說不可以.
神叫人不可以輕看這個重建聖殿的行動.
這裡說.
"這七阿賴士耶和華的眼睛片刷全地.
見所羅巴伯手拿善陀就歡喜".
是什麼意思呢?.
意思就是在神的眼中.
重建聖殿是非常重要.
他只要見到所羅巴伯手裡拿著的那個善陀.
其實是一個原石.
是在他重建聖殿的過程裡面.
他要用來在當中的工具.
所羅巴伯為什麼手拿著善陀.
神見到就開心.
見到就歡喜.
因為這個代表著所羅巴伯他願意聽從神的吩咐.
去帶領猶大百姓參與重新開始.
那個停了16年時間的重建聖殿的工程.
當神見到他們願意迴轉.
當神見到他們願意去重建聖殿的工程的時候.
神就很開心.
不單止是剛纔所說的人很開心.
會歡呼.
"願因為因為歸於這殿".
原來同時神也會感到非常的歡喜.
神喜悅.
神喜悅重建聖殿這件事.
究竟這段經文想說什麼.
是不是想大讚所羅巴伯.
就是很聽話.
要高舉他呢.
不是.
我想邀請大家一起讀出紅色的字.
一二三.
(紅色的字).
這一字的經文我相信大家都熟悉.
原來是出自第五個異象裡面.
神要用這一句話去提醒.
先知撒迦利亞提醒猶大的百姓.

$^{441}$聖殿之所以能夠被重建起來.
能夠完成能夠成功.
絕對不是依靠所羅巴伯一個人的勢力.
也不是依靠所羅巴伯他一個人的才能.
是依靠誰.
是依靠神的靈.
也就是要全然依靠聖靈.
才能夠將這個那麼困難的工程完成.
如果我們說.
只是依靠聖靈去重建聖殿.
那所羅巴伯的角色.
約書亞的角色是什麼呢.
我們又要回去看這幅圖畫.
在這裡約書亞和所羅巴伯就是那兩個獸糕仔.
也就是那兩棵橄欖樹.
金色的油其實就是金燈臺所需要被點亮的燈油.
而這個燈油就是透過兩棵橄欖樹枝.
去供應給金燈臺.
也就是說所羅巴伯和約書亞.
就是為金燈臺提供的燈油.
他們將燈油倒入金道館裡面.
他們是獸糕仔.
也就是生油之子.
也就是將新鮮的橄欖油倒進去.
那我們要問這個油是從哪裡來的呢.
誰糕他們的.
是神.
神仙子是那個燈油的源頭.
他們的油是從神而來的.
這個燈油其實是指聖靈.
是指神的靈.
所以神仙子說是要依靠神的靈.
方能成事.
這個震撼的場面.
要帶出的就是.
在聖殿裡面會有金燈臺.
在聖殿還沒有重建的時候.
神將這個金燈臺的異像賜給先知撒迦利亞.
讓他看見已經很久沒有看見的金燈臺.
是點亮的在他的面前呈現出來.

$^{481}$代表著什麼.
聖殿必定會被重建起來.
即使眼前有多大的困難.
但是神一樣會透過所羅巴伯的手.
透過約書亞的手來將聖殿重建起來.
他們唯一要做的就是.
要跟隨約書亞和所羅巴伯的吩咐.
這群百姓.
要一起同心合意來重建聖殿.
今天我們知道先知撒迦利亞這個異像.
其實這個預言已經實現了.
成就了.
因為聖殿確實已經在公元前516年重建完成.
這個聖殿被稱為所羅巴伯聖殿.
神的應許確實沒有落空.
但是這個異像不是止於此.
這個異像還有一個更深的層面.
值得我們來思考.
因為我們發現約書亞的身份是大祭司.
他代表的是祭司.
而所羅巴伯不是大祭司.
他是誰?.
他是猶大百姓的一個領袖.
但事實上他是大位的後裔.
所羅巴伯的名字出現在.
馬太福音和路加福音耶穌基督的家譜當中.
所以他是被視為君王的代表.
講到這裡.
有約書亞有所羅巴伯.
也就是有祭司又有君王代表著什麼?.
誰是有這兩個身份的同時?.
是彌賽亞.
是耶穌基督.
主耶穌既是君王既是祭司.
他就是彌賽亞.
而彌賽亞的意思正正就是受高者.
原來這個異像還有一個更深的層次.
是指向將來的.
指向終末的.
指向彌賽亞的國度是要來到這個世界.

$^{521}$意思就是整個世界要被神來更新.
當我們再讀回第九節的時候.
我們看到原來當聖殿完成的時候.
重建完成的時候.
神說你就知道萬君之耶和華.
差遣我到你們這裡來了.
在這裡的我其實就是耶穌基督.
就是彌賽亞.
誰差遣耶穌基督到聖肉身的?.
是天父.
我們發現在舊約的這一段很寶貴很寶貴的經文裡面.
原來出現了三一真神在當中.
天父是差遣聖子來到這個世界.
他要去進行救恩的計劃.
而在耶穌基督被釘死在十字架上.
然後他復活他昇天之後.
耶穌又將聖靈帶到教會當中.
還記得五重節.
其實就是聖靈降臨的日子.
一個很重要的事件.
也就是教會誕生的日子.
是一個很重要的時刻.
是由聖靈降臨開始.
教會的開始就是聖靈的降臨.
於是我們發現三一神他親手開始.
亦都親自去完成每一個的計劃.
當中過程的每一步.
都不是出自人的手.
而是出於這位創始成宗的三一神的手.
今天這個異象對於昔日的猶大群體來說.
帶有的意義是什麼呢?.
有聖殿就是等於有金燈臺.
金燈臺的異象就是預告聖殿必定會被重建.
聖殿會被重建.
金燈臺就會被點亮.
聖殿的重建和金燈臺的點亮.
其實都是在見證這位三一神的榮耀.
確實這群猶大的群體.
他們經歷的是亡國 被腦.
然後要回歸聖城.

$^{561}$重建聖城 重建聖殿 重建聖場.
這一切對於這群亡國被腦的人來說.
其實是非常艱難的事.
什麼都沒有.
人又沒有錢.
怎麼去重建呢?.
開玩笑.
但神說這不是一個玩笑.
透過這位創始成中的上帝.
三一的真神.
在當中.
他有一個超乎我們所想象的計劃.
他親自去成就這件事.
將這個又大又難的事.
沒有可能發生的事.
就是這樣發生在昔日這群猶大的群體當中.
今天呢.
當我們去看到這個金燈臺的異象.
我們會不會說.
其實也不是跟我們很有關係.
我們連舊約的金燈臺都沒見過.
舊約的第一個聖殿我們都沒見過.
不是很關我們事的.
這段經文好像很遠.
但其實.
剛纔已經提到.
這段經文不僅是第一個預言的實現.
其實還有一個未來的嚮導.
就是救恩的計劃.
當這個所羅巴伯的聖殿被重建之後.
在五百多年後.
耶穌基督.
祂就親自道成肉身來到這個世界.
然後祂在十字架上為我們成就這個寶貴的救恩.
當祂昇天之後.
祂差派聖靈來到.
祂要聖靈來到讓信徒明白.
之前祂所教導的每一句話.
每一個信徒當他信耶穌的時候.
他都是接受聖靈的印記.

$^{601}$聖靈進入我們的生命當中.
剛纔提到教會的被建立.
初期教會一開始的建立是從聖靈降臨開始的.
聖靈的降臨.
進到信徒的生命裡面.
因此信徒就是聖靈的電.
我們是聖靈的電.
我們是基督的身體的一部分.
我們是要聽從那位基督.
這個是教會的頭.
教會其實就是現代的金燈臺.
因為我們就是聖靈的電.
我們每一個信徒就是金燈臺的一分子.
今天我們的生命是怎樣的呢?.
我們作為一個基督徒.
我們作為贅聖的弟兄姊妹.
我們生命的流露是怎樣的呢?.
我們有沒有將這個創始成中的三一神.
在我們生命的見證當中流露出來呢?.
聖靈有沒有在我們生命裡面湧流呢?.
有沒有呢?.
這是今天這個金燈臺的異像要給我們的提醒.
要給我們的激勵.
神是期望這個異像放在聖經裡面.
不只是向當時的那群猶大人說話.
是要向我們這群信徒今天仍然這段經文.
要向我們來說話.
當我去預備這個信息的時候.
我很想和大家一起去思考這幾個問題.
如果今天我們說占卜是一個金燈臺的時候.
不知道是不是金呢?.
我們評分吧.
現在很喜歡做一些測評.
又會找些人工智能來評分.
我們試試評估一下.
占卜作為一個金燈臺我們會給多少分.
而當我們這樣去評分的同時.
我們也要去想一下.
我們自己有沒有在當中有貢獻呢?.
而我們又給自己多少分呢?.

$^{641}$另外一個向度.
我們去想.
昔日這群猶大的百姓.
神吩咐他們要重建聖殿.
他們喊著救命.
說很難呀 究竟神在哪裡呀 搞不定.
但最後搞定了.
成就了.
如果今天神將這個異象放在占卜裡面.
要我們去重建占卜的時候.
你會不會拒絕呢?.
會不會好像昔日這群猶大百姓那樣.
不相信會成就呢?.
你看到的是什麼困難呢?.
而我們認為我們的困難.
是否大到跟昔日猶大的困難去比較的時候.
是比他們更難.
還是當時比我們現在更難呢?.
這幾個問題我自己也在思考中.
我覺得今天這段經文.
其實神是很想我們有一個進一步的反思.
我們今天怎麼看我們回來的這個占卜家呢?.
當我們說很希望占卜家可以為我們提供很多很多東西的時候.
我們其實有沒有在當中去付出呢?.
而當今天如果神說祂要重建這個地方.
我們看到的是不是我們沒有錢.
我們沒有人.
我們中間斷了層.
我們真的寸步難行.
我們是否只是將這些困難放大.
而忘記了其實.
我們所相信的這一位神.
是一位創始成中的三一神呢?.
如果我們相信的話.
如果我們再一次去重新仔細讀這段經文的時候.
我們就好像代入先知撒加里亞這個角色.
我們在漆黑的這個世代裡面.
我們看到的是不是神給我們的這個金燈臺的異象呢?.
讓我們一起來為教會.
為我們身邊的弟兄姊妹來祈禱.

$^{681}$我相信如果上帝的心意是要重建這個地方.
祂必定會成就.
問題是我們是否願意在當中參與其中.
讓我們一起低頭同心禱告.
親愛的主耶穌基督.
親愛的天父上帝.
親愛的聖靈.
我們再一次俯伏在禰的根前.
因為禰是那位呼召我們的三一神.
其實我們真的有很多的軟弱.
我們生命裡面有很多的艱難.
我們可能想我們搞不定自己.
我們怎麼來搞定教會呢?.
其實昔日這一幫猶大人都是一樣.
他們將聖殿的工作停了十六年.
今天我們讀這段經文的時候.
我們真的很想聽到神禰自己的心意.
我們將尖鋥前面的路公敬教在禰的手中.
因為禰是尖鋥的主.
禰是尖鋥的王.
無論我們做什麼決定.
是否要重建這個地方.
重建的藍圖是怎樣.
是何時要去進行這件事.
我們將這一切交託在神禰的手中.
神禰求禰去帶領我們眾人.
讓我們真的能夠同心合意看見同一個的異象.
就是神禰所賜給我們的這個金燈臺的異象.
讓我們的生命能夠被聖靈來到充滿.
讓我們上上的來到去.
檢視我們自己和禰的關係.
讓我們能夠來到禰的面前.
去上上的應聚悔改.
去潔淨我們自己.
以致到我們是能夠毫無難阻的.
被神禰來到去使用.
亦讓我們可以毫無難阻的.
來到去聽到神禰自己的聲音.
禰的吩咐.
禰的指引.

$^{721}$以致到我們願意去回應禰給我們的使命.
和呼召和異象.
讓我們去領受的時候.
我們就是這樣帶著信心來到去踏前.
因為我們知道.
神禰是那位創始成中的三一神.
求禰親自去與我們每一個人同在.
當我們軟弱的時候.
神求禰的聖靈來到去幫助我們.
加力給我們.
當我們剛強的時候.
神求禰叫我們不要驕傲.
讓聖靈去成為我們的提醒.
叫我們真的能夠可以.
在神禰所賜給我們的這個詹審家裡面.
我們各人就是按著聖靈所賜的恩賜.
在神禰擺在我們的綱位.
角色身份的裡面.
我們是忠心地來到去服侍禰.
以致到能夠使到詹審這個金燈臺.
是能夠被禰使用.
成為這個黑暗世代光明的見證.
我們這樣禱告,交託.
奉靠我主耶穌基督.
得勝又得勝的名字祈求.
阿們.
\newpage



\section{以弗所書 1:13-19}
\label{sec:031850}
\textbf{受了印記}~~~~俞大鳴牧師
\newline
\newline
連結: \href{http://www.church.com.hk/acms/content.asp?site=cdc\&op=show\&type=product\&code=031850\&layout=sermon}{\texttt{031850}} ~~~~ 語音日期: 2026-03-22
\newline
\newline
\hyperref[sec:031849]{\small{< < < PREV SERMON < < <}}
~
\hyperref[sec:index_chronic]{\small{[返順時目]}}
~
\hyperref[sec:031851]{\small{> > > NEXT SERMON > > >}}
\newline
\newline
以弗所書 1:13-19
\newline
\begin{longtable}{cl}
\hline
\hline
章節 & 經文 (和合本修訂版)\\
\hline
1:13 & \begin{tabularx}{0.7\textwidth}{X} 在基督裡你們聽見真理的道，就是那使你們得救的福音，你們也信了他，就受了所應許的聖靈為印記。 \end{tabularx} \\ \\ \relax
1:14 & \begin{tabularx}{0.7\textwidth}{X} 這聖靈是我們得基業的憑據，直等到神的子民得救贖，使他的榮耀得到稱讚。 \end{tabularx} \\ \\ \relax
1:15 & \begin{tabularx}{0.7\textwidth}{X} 因此，我既然聽見你們對主耶穌有信心，對眾聖徒有愛心， \end{tabularx} \\ \\ \relax
1:16 & \begin{tabularx}{0.7\textwidth}{X} 就不住地為你們感謝神，禱告的時候常常提到你們， \end{tabularx} \\ \\ \relax
1:17 & \begin{tabularx}{0.7\textwidth}{X} 求我們主耶穌基督的神，榮耀的父，把那賜人智慧和啟示的靈賜給你們，使你們真正認識他， \end{tabularx} \\ \\ \relax
1:18 & \begin{tabularx}{0.7\textwidth}{X} 照亮你們心中的眼睛，使你們知道他呼召你們來得的指望是甚麼，他在聖徒中所得榮耀的基業是何等豐盛， \end{tabularx} \\ \\ \relax
1:19 & \begin{tabularx}{0.7\textwidth}{X} 並知道他向我們這些信的人所顯的能力是何等浩大，這是照他的大能大力運行的。 \end{tabularx} \\ \\
[1ex]
\hline
\hline
\end{longtable}
$^{1}$多謝各位平安,我們先一起祈禱.
讓我們可以一同聚集,並拜你.
一同思想你的應許,你的話語.
求附你的聖靈帶領我們.
讓我們從你的話語當中得到安慰,鼓勵.
奉耶穌基督的名祈禱,阿們.
好,今天是「全心全意」系列的最後一講.
這次的系列主要是基於《已忽所書》.
第一章三到十四節提到.
信了主耶穌的多重恩典.
保羅在一章三節提到.
祂,即是指上帝,在基督裡.
曾把天上各樣屬靈的福氣賜給我們.
然後提了五樣特別的福氣.
包括在創世之前已經揀選了我們.
給了我們兒女的名分.
藉著主耶穌基督救贖我們.
又使我們歸入基督的國度裡.
今天我們要思想的是.
上帝給了我們聖靈作為印記.
印記是什麼意思呢?.
印記有什麼功能呢?.
印記有沒有一些象徵的意義呢?.
這些都是我們今天要思考的問題.
對我們這些現代人來說.
我不知道有沒有人聽過.
西方人有時會說.
Signs show and deliver.
一般都是用在比較大的交易上.
要簽合約.
對我們來說.
一般比較大而要簽合約的交易.
大概就是房屋買賣.
假如我們要買房.
看到一間合適的房屋之後.
首先要交一些訂金給中介.
拿一份合同回去.
看了可能自己不懂.
交給律師看有沒有問題.
沒問題就行了.

$^{41}$約定什麼時候交大訂.
即是房價的10\%.
通常買家就簽合約.
交給律師.
賣家都簽合約.
又交給律師.
到了兩方面的買賣家.
雙方都交給律師.
那時候那兩份合同就可以互換.
在合同上其實你看不到.
買賣雙方的簽名.
一方就夠了.
律師可能會告訴你.
合約已經Signs show and deliver.
律師說話很詳盡.
不單止簽了名.
蓋了章已經交換了.
買方也可以對客戶說同樣的話.
合約簽了之後.
如果買方突然改變主意.
我還是不買了.
不買就慘了.
首先你的大訂是不會退回給你的.
然後買方將房子賣了.
如果金額低於原來肯出的價格的10\%.
買方還可以要求你賠錢.
我買不回這樣的價錢.
你的10\%都不夠.
於是仍然可以要求你賠.
所以交大訂合約是非常重要.
今日經文第十三節提到印記.
在羅馬的社會當中.
印章是代表所有權,真實性,保護和安全的象徵.
人們怎樣做印章呢?.
有些是用戒指做的.
有些就像華人一樣.
用另外一塊木石做印章.
他們吸引華人用紅泥.
他們就用蠟.
把蠟或軟質的泥弄軟後.

$^{81}$將印章吸下去印.
一般用在文件,商品,信件,合約,遺囑.
甚至財產的生畜都會印.
一般來說有印記有四方面的作用.
第一是驗證官方文件.
如果有法令,協議.
要證明文件是真的就要印.
第二是標示所有權.
如果把印打在動物身上就有所有權.
如果有東西放在容器.
就要保證品質,品牌差不多.
第三是提供安全保障.
有時在聖經裡會看到.
耶穌的墓穴是要求要吸引的.
以免人們移過墳墓的門.
第四是在法律和商業的情境當中.
用來證明合同上的交易是真實的.
第十四節提到另一個隱喻.
隱喻不是很明顯說明這是比喻.
但我們一看就知道.
就是憑據.
憑據就是剛才說的買屋的例子.
你給的10\%是一個憑據.
了解保羅所用的兩個比喻的意思後.
我們可以思考今天的經文.
在第十三節保羅說.
凡是聽了福音相信耶穌的人.
就會有所應許的聖靈為印記.
這句話的意思就是說.
我們屬於上帝的身份.
是得到聖靈的確認的.
有一個傳福音的方法.
是教基督徒怎樣跟人家傳福音的時候.
怎樣去打開華俠的.
這個方法一般叫做三福.
有沒有人學過三福或聽過三福?.
沒有呀?有呀.
這個方法通常是問兩個問題.
第一個問題他會問.
假如你今天離開世界.

$^{121}$你有沒有把握進入天堂呀?.
無論那個人怎樣回答.
第二個問題就是.
假如你來到天堂門口.
有天使來問你.
我為什麼要讓你進入天堂?.
你怎樣回答呢?.
當然這兩個問題都是假設性的問題.
對一位真正的信徒來說.
我們是有聖靈為印記.
應該有絕對把握可以進入天堂.
但有時當人家用這個方法去問信徒.
有些信徒我不知道是什麼原因.
或者他想表現自己很謙虛.
又或者他受到不正確的教導.
他不敢確定自己是否一定可以進入天堂.
他會回答.
我不知道的.
要到時才知道的.
到時再看看怎樣吧.
今日這一節經文就很重要.
我們是屬於上帝的.
我們是上帝的兒女.
有聖靈作為印記.
打了印的.
即是被上帝確認的.
第二個問題在天堂門口其實是不會出現的.
因為我看不到你身上有沒有聖靈作為印記.
但天使一定會看到的.
天使不可能看不到的.
因此天使是不會攔著我們問這個問題.
我們得救是由於信了耶穌.
而不是由於有神學知識.
然後去天堂門口還要參加考試.
考問你為何要讓你進來.
不會出現這種情況.
簡單來說.
聖靈作為印記.
就是確認我們是上帝子民的身份.
十四節就繼續說「聖靈是我們得基業的憑據」.

$^{161}$這裡是說上帝將聖靈賜給我們.
是一個非常嚴肅的承諾.
就好像我們買房子要交大訂一樣.
非常嚴肅的承諾.
現在我們雪梨挑很貴的.
平均每分鐘百多二十萬.
一交就要交十多二十萬出去.
這個承諾嚴不嚴肅呢?.
這個非常嚴肅的承諾.
上帝先將聖金給了我們.
等到時候我們就會承受所有的應許.
或者我們會問「等到何時呢?」.
這次經文告訴我們.
「等到上帝的子民得救贖」.
什麼時候呢?.
什麼時候上帝的子民得救贖呢?.
我相信時後日期是不重要的.
主耶穌也說什麼時候.
沒有人知道.
只有上帝知道.
但其實上帝說的時候不是一個日期.
而是要等到一個得救人數的數字滿了.
我們看看羅馬書11章25至26節.
那裡這樣說.
「等到外邦人的數目天滿了.
以色列全家都要得救」.
這裡的數字是多少呢?.
我們是不知道的.
不過我們知道一件事.
如果我們想早日得基業.
我們要努力傳福音.
快些傳福音.
越多人信主.
可能就越早達到上帝心目中的數目.
我們繼續思想.
聖靈給了我們做印記.
在現在這個階段.
我們仍然活在世上的時候.
有沒有其他作用呢?.
有啊.

$^{201}$17字說.
在保羅的禱告中提到.
「那賜人智慧和啟示的靈.
賜給你們使你們真認識他」.
原來上帝把聖靈賜給我們.
不是在我們裡面睡覺.
而是要賜我們智慧和啟示.
智慧和啟示有什麼目的呢?.
就是要我們真正認識他.
我相信我們每個人的信仰過程.
都有一些相同的地方.
就是起初是認識主耶穌.
相信主耶穌.
慢慢越來越認識多一些.
慢慢會多的.
主耶穌.
有關主耶穌的事.
有關聖靈.
有關唯一的.
慢慢會多的.
如果我們信了耶穌超過五年.
還自覺對上帝是很陌生的話.
這樣很有可能我們沒有全心全意信耶穌.
如果有人信了二十年.
還覺得上帝很陌生.
我真的有百分之九十五的把握.
認為你沒有全心全意信耶穌.
一個人假如全心全意真心相信耶穌.
又有聖靈住在他裡面.
他一定會越來越認識上帝.
其實好像我們認識朋友一樣.
起初認識的朋友.
當然不是很認識他.
但如果你經常跟他有來往.
相處得久.
又經常跟他吃飯.
我們慢慢就認識那個人的喜好.
皮性.
聖靈除了幫助我們認識上帝之外.
其實也會對我們其他方面有幫助.

$^{241}$尤其是在我們日常的生活方面.
接下來我想跟大家說.
三個聖靈的果子的特質.
來分享聖靈怎樣在日常生活方面.
也會幫助我們.
我們知道聖靈所見的果子是忍耐希臘和平一.
忍耐恩慈良善信十分之絕世.
很多專家都指出.
在這節加拉太書的經文裡面.
果子的字在原文希臘文是單數.
它的涵義不是你結九個果子出來.
你是結一個果子.
但這個果子真的特別.
有九種品質.
也可以將這九種品質分為三個大方向.
忍耐希臘和平一組.
忍耐恩慈良善第二組.
第三組是信十分之絕世.
第一組忍耐希臘和平是講自我修為.
是我們內心要保持的心態.
忍耐恩慈良善有些是被動.
即是對別人怎樣對我們的時候.
有被動的反應.
信十溫柔之際.
就是我們主動和別人相處的時候.
的態度和方法.
由於時間關係.
我就在每一組選一個品質和大家分享.
在第一組我想講的是喜樂.
為什麼我會選喜樂呢.
主要原因是因為我們主日查經.
最近已經查完《仕途行傳》.
剛剛就開始查考《約伯記》.
很多弟兄姊妹都知道.
約伯受苦.
一日之內子女全部死了.
財物又失去.
真是很苦.
但《約伯記》的主題不是討論苦難的來源.
不是討論來源問題.

$^{281}$而是討論什麼是真正的信仰.
在《約伯記》第一到第二章.
介紹他是一位非常敬虔的信徒.
連上帝都為他作見證.
告訴我們.
他是完全正直敬畏上帝.
遠離惡事.
他不反駁上帝的判斷.
他只不過挑戰上帝的眼光.
上帝你的眼光差異.
他就批評約伯信你.
是出於利己的目的.
他是一位功利主義者.
你試試拿走他的東西.
他還信不信你.
還能不能堅持信心.
上帝讓人無故受苦.
還有另一個討論.
上帝讓人無故受苦.
上帝算不算公義?.
上帝是正義.
受苦的原因是否由於人犯了罪.
得罪上帝呢?.
這些其實都是《約伯記》要討論的問題.
主題是何謂純正的信仰.
下面有幾個分題.
這些就是其中一些的分題.
我不是在這裡說答案給大家聽.
你有興趣就來參加《曾經八》.
但是由於講到.
人是會無緣無故受苦的.
我就揀選喜樂.
來跟大家打支預防針.
如果有人現在在受苦.
我勸他要喜樂.
我可能有些白痴沒什麼人性.
人家受苦還叫他喜樂.
怎麼喜樂呢?我也做不到.
但是我想在大家還沒經歷苦之前.
在大家還沒經歷苦之前.

$^{321}$跟大家說你能夠結出喜樂這種特質.
你就以喜樂作為良藥.
就可以面對人生各種困難苦難.
當然我想說.
喜樂不是唯一的良藥.
上帝很恩待我們.
為人預備了眼淚.
幫助我們在痛苦的時候.
可以流淚.
沖洗去我們心中的受苦.
眼淚真是上帝的恩賜.
至於我還有其他方法.
最近有位以前的會友打電話給我.
有一位好朋友的丈夫.
和他的朋友去釣魚.
結果一去無回浸死了.
他很傷心.
傷心的時候又想起了我.
他打電話給我.
為他祈禱.
之後他也覺得舒懷很多.
《羅馬書》第八章提到聖靈會用.
無可言喻的嘆息在我們祈求.
在苦難中我們禱告.
向上帝盡訴心中的苦情.
都是面對苦難的一個好方法.
而人生中苦難是難免的.
聖經有八經五章七次都說.
人生之來必遭遇患難.
如同火花飛騰.
現在我們先講講現代人受苦的經歷.
讓大家知道.
受苦隨時有可能發生.
這個現代人我們都認識的.
就是我們現在的代理主任牧師.
他最近出版了這本書.
《神的苦難與基督徒》.
這本書有些弟兄姊妹都買了.
這本書提到一個他自己經歷的苦難.
一章是講他媽媽的苦難.

$^{361}$後來他提到一個苦難.
在六個月之內發生.
在他住的地方發生了三場暴雨.
第一次暴雨是雷暴.
閃電之後電腦的印刷機和互聯網的modem都壞了.
就收你了.
第二次原來他屋頂有幾塊瓦面漏水.
水漏下漏下.
cutter去水槽又可能塞了.
去不到水水又去了哪裡呢.
不進屋去了兩間房的天花板.
結果承受不住壓力.
兩間房的天花板都跌了下來.
第三次的雷暴又剛巧打爛了一棵樹的樹枝.
砸爛了他的車頂.
又要拿去修車.
夠苦難未呢.
在那段期間他太太和一個兒子都發生交通意外.
沒有受傷.
夠未呢.
未喎.
家裡的冷氣機焗爐打了好幾次.
接二連三一個壞然後到一個.
此外他自己本人去了兩次醫院做手術.
而他一個兒子就病了.
在醫院住了一個月.
我記得我二零二三年做手術都住了一個月.
住一個月醫院不是開玩笑.
然後九十歲的外父爬梯不知做什麼.
大概一米半的爬梯跌了下來.
身體上有四十處骨折.
入醫院住了兩個月.
才讓他出院.
他還有一個朋友和他很幽默地說.
你外父爬梯之前沒有看看他的出世紙嗎.
老人家忘記自己幾十歲了.
面對這麼多人生的麻煩事也好苦難也好.
能不能用喜樂去代替受煩呢.
我相信在人事不能敵.
在上帝就凡事都能.

$^{401}$在新教裡面其實有教我們的.
是要用喜樂去面對苦難的.
在《雅各書》一開播第二節經文.
一章第二節.
在第七門領門遭受各種試煉時.
都要認為是大喜樂.
試煉其實我們知道咸義就是受苦.
雅各勸弟兄姊妹不要在苦難中埋怨.
不要離開上帝.
反而要以喜樂去面對.
有沒有搞錯.
在苦難中還要喜樂?還要大喜樂?.
為什麼?.
在第三節雅各進一步解釋.
因為知道你們的信心經過考驗.
就心忍樂.
在彼得前書一章六至七節有同樣的教導.
雖然你們必須在百般苦難中暫時憂愁.
你們要為此喜樂.
在第七節就解釋為什麼要以喜樂代替憂愁.
使你們的信心既被考驗.
就比那被火試煉仍然能壞的金紙更顯寶貴.
可以在耶穌基督顯現的時候得著稱讚榮耀尊貴.
這裡說到信心被考驗.
是一個被動的語態.
被動的.
我們千萬不要自己去拿苦離身.
不要去找苦難.
涵義是說上帝故意讓我們遇到苦難困難.
是要考驗我們的信心.
在約伯記第一至第二章我們看到撒旦兩次提出.
約伯只是為了想得上帝福而信上帝.
才會乖乖地遠離惡事做好人.
那些信心是出於功利主義的目的.
一點都不單純不純正.
於是上帝就讓撒旦去整約伯.
不過兩次我們看到是誰先開口呢?.
是撒旦開口還是上帝先開口呢?.
兩次都是上帝先開口.
上帝去撩撒旦.

$^{441}$「你在哪裡來?有沒有看到我捕人約伯?」.
是上帝先開口.
表揚約伯的信心在先.
故事也可以說是上帝看得起約伯.
才會讓約伯去受苦.
以至他的信心故事一直存留了幾千年.
直到今時今日.
最後一句也是之前我們遇到苦難時要喜樂的另一個原因.
就是盼望.
在耶穌基督顯現的時候得著稱讚榮耀尊貴.
保羅其實也教導我們以喜樂的態度去面對患難.
我相信很多人都讀過《肥勒比書》.
都知道《肥勒比書》裡面經常講喜樂.
《肥勒比書》只有四章的經文.
一百零四節的聖經.
喜樂出現了十二節.
平均每九節經文就出一次叫你喜樂.
在四章四節還特別在收筆之前.
特別再講一次.
「你們要靠主常常喜樂」.
「我再說你們要喜樂」.
我們是不是為了有苦難而喜樂呢?.
其實不是.
我們不是為了那個令人痛苦的處境而喜樂.
而是為了有上帝的同在而喜樂.
《詩篇37篇1至9》說.
「不要為作惡的心懷不平.
也不要嫉妒那行不義的人」.
當我們看到不義的事情發生的時候.
有正義感的人其實都會心懷不平.
好像美國無端端去打伊朗.
無端端打錯一間小學.
一百七十個小朋友無端端死了.
作為澳洲人.
我都為這種行為覺得很可恥.
還有我們的政府經常看到美國人沒有腰骨.
不但不敢批評美國的違法行為.
反而指責伊朗政權很殘暴.
但是作為一介普通草民.
心懷不平的庸人.

$^{481}$解決不了問題.
沒用的.
所以詩人在《詩篇37篇4》說.
「就當以義和華為樂」.
不要生氣了.
有上帝同在就有喜樂.
我們能夠保持喜樂的心.
就是有上帝同在.
有聖靈的印記.
情姦,罰惡是上帝會做的事.
留給上帝做吧.
等上帝做吧.
希望大家都收到預防針的效果.
如果將來有人遇到苦難的時候.
聖靈又或者感動你想起我.
你都可以打電話來找我.
我會聽聽你的苦情.
接著要講忍耐.
希臘文有一個好處.
因為忍耐有兩個不同的字.
分得很清楚.
一個是講對人的忍耐.
另一個是講對事情的忍耐.
其實對一些自己不能控制的事.
除了忍耐之外.
真的沒有什麼辦法.
我想起太太下班的時候.
通常坐火車去伊斯活火車站.
我去接她.
火車站一出來就經過一個迴旋處.
有時會多車.
遇到好的司機.
他會停在迴旋處.
但不會阻礙路.
我就可以直穿過去.
我直穿是沒有塞車的.
如果不是被迴旋處的車塞住.
但有時遇到的司機.
不知道是不用腦還是專顧自己.
他不理會.

$^{521}$他自己前面不能去.
不理會.
照樣塞住你的路.
等你又不能去.
遇到這些事情.
除了忍耐等候之外.
真的別無他法.
但對人的忍耐有些不同.
如果我們要和別人有好的關係.
即使對方有些事.
有些做事方法和我們不同.
我們都要積極地忍耐.
有機會和配偶同住的人.
就更加要學會忍耐.
因為兩夫妻來自不同的原生家庭.
有很多的習慣.
都是從小養成.
很難改的.
尤其是一些和道德無關的事.
我記得不記得幾年.
我想大概二十年前.
有兩位心理學家.
一起寫了一本書.
《Man is from Mars, Women is from Venus》.
男人來自火星.
女人來自水星.
兩個星球的含義.
就是他們都不是來自地球.
都不是地球人.
來自這兩個星球.
其實這兩個是離地球比較近的行星.
人類最近已經很厲害.
派火箭飛去火星探測.
但據說飛去火星一次單程要九個月.
大約起碼九個月.
所以和配偶相處真的很不容易.
因為不是忍一次.
不是忍一時,要忍一世.
我們基督道婚姻是講一生一世.
不是講何時離婚.

$^{561}$所以要忍一世真的不容易.
因此我們要多紀念配偶.
對我們的好.
多想念對方的好處.
這樣會容易忍耐配偶多一點.
聖靈也會幫助我們.
結出忍耐的特質.
這裡講的忍耐.
就是對人的忍耐.
即使沒有配偶.
我們也有很多不同的人要忍耐.
兒女.
對兒女相處要忍耐.
同同事之間.
兄弟姐妹之間.
朋友之間.
我們也需要忍耐.
總之與人相處.
就是要有忍耐.
一定要有聖靈的幫助.
才可以做得好.
最後講一講溫柔.
溫柔不只是音聲細語說話.
也不是軟弱或膽怯.
不是任人踐踏成為別人的地毯.
其實原文的字就像一匹馬.
馬比人大力很多.
但如果我們將一隻雀環.
將繩子放在牠的口中.
就可以控制牠了.
其實溫柔是一種受控制的力量.
有力量不要緊.
力大不要緊.
受不受得到控制.
或者是信服的勇氣.
這種美德可以說是源於對上帝的信任.
我們願意將自己的能力,權利.
交給上帝去掌管.
對別人溫和.
也尊重別人.

$^{601}$即使自己有能力有權利.
我們也不會霸凌人家.
好像美國人最喜歡欺凌別人.
因此我們有能力有權利.
不去霸凌別人.
其實自己也很難做.
只有靠聖靈幫助才有可能.
所以溫柔也是聖靈果子其中一種特質.
更加是主耶穌的品格.
主耶穌曾經說過.
我心裡柔和謙卑.
你們當乎我的額.
向我學習.
這樣你們的心靈就必得向安息.
柔和謙卑不是有人想學就學得到的.
要有聖靈的幫助才可以.
我今天要做一個總結.
上帝賜給我們聖靈作為一個印記.
認定我們信祂的人是祂的家人.
上帝是我們天上的父.
我們是祂的兒女.
將聖靈賜給我們.
除了保證和保障我們的身份之外.
也是將今生和永生的福氣賜給我們.
在我們今生的生活中.
聖靈會指引我們做榮臣益人的決定.
會幫助我們結出喜樂的果子.
面對人生中必然會有的艱難時刻.
有忍耐和溫柔去和人相處.
也都信服上帝.
上帝就是這樣全心全意愛我們.
將聖靈作為印記給了我們.
我們要怎樣回應上帝的愛呢?.
我相信聖靈會給大家智慧去做選擇.
我們一起祈禱.
你差遣了主耶穌基督來到世上.
更加將聖靈賜給我們作為隨時的幫助.
求聖靈幫助我們能夠結出聖靈的果子.
讓我們每日的生活都充滿了從天上而來的平安喜樂.
並以我們的生命全心全意地事奉你.

$^{641}$奉耶穌基督的名祈禱.
阿們.
\newpage



\section{約翰福音 2:1-12}
\label{sec:031851}
\textbf{行奇事的神}~~~~梁紫健傳道
\newline
\newline
連結: \href{http://www.church.com.hk/acms/content.asp?site=cdc\&op=show\&type=product\&code=031851\&layout=sermon}{\texttt{031851}} ~~~~ 語音日期: 2026-03-22
\newline
\newline
\hyperref[sec:031850]{\small{< < < PREV SERMON < < <}}
~
\hyperref[sec:index_chronic]{\small{[返順時目]}}
~
\hyperref[sec:031852]{\small{> > > NEXT SERMON > > >}}
\newline
\newline
約翰福音 2:1-12
\newline
\begin{longtable}{cl}
\hline
\hline
章節 & 經文 (和合本修訂版)\\
\hline
2:1 & \begin{tabularx}{0.7\textwidth}{X} 第三日，在加利利的迦拿有一個婚宴，耶穌的母親在那裡。 \end{tabularx} \\ \\ \relax
2:2 & \begin{tabularx}{0.7\textwidth}{X} 耶穌和他的門徒也被請去赴宴。 \end{tabularx} \\ \\ \relax
2:3 & \begin{tabularx}{0.7\textwidth}{X} 酒用完了，耶穌的母親對他說：「他們沒有酒了。」 \end{tabularx} \\ \\ \relax
2:4 & \begin{tabularx}{0.7\textwidth}{X} 耶穌說：「母親，我與你何干呢？我的時候還沒有到。」 \end{tabularx} \\ \\ \relax
2:5 & \begin{tabularx}{0.7\textwidth}{X} 他母親對用人說：「他告訴你們甚麼，你們就做吧。」 \end{tabularx} \\ \\ \relax
2:6 & \begin{tabularx}{0.7\textwidth}{X} 照猶太人潔淨禮的規矩，有六口石缸擺在那裡，每口可以盛兩三桶水。 \end{tabularx} \\ \\ \relax
2:7 & \begin{tabularx}{0.7\textwidth}{X} 耶穌對用人說：「把缸倒滿水。」他們就倒滿了，直到缸口。 \end{tabularx} \\ \\ \relax
2:8 & \begin{tabularx}{0.7\textwidth}{X} 耶穌又說：「現在舀出來，送給宴會總管。」他們就送了去。 \end{tabularx} \\ \\ \relax
2:9 & \begin{tabularx}{0.7\textwidth}{X} 宴會總管嘗了那水變的酒，並不知道是哪裡來的，只有舀水的用人知道。於是宴會總管叫新郎來， \end{tabularx} \\ \\ \relax
2:10 & \begin{tabularx}{0.7\textwidth}{X} 對他說：「人家都是先擺上好酒，等客人喝夠了才擺上次的，你倒把好酒留到現在！」 \end{tabularx} \\ \\ \relax
2:11 & \begin{tabularx}{0.7\textwidth}{X} 這是耶穌所行的第一個神蹟，是在加利利的迦拿行的，顯出了他的榮耀來，他的門徒就信他了。 \end{tabularx} \\ \\ \relax
2:12 & \begin{tabularx}{0.7\textwidth}{X} 這事以後，耶穌與他的母親、兄弟和門徒都下迦百農去，在那裡住了不多幾天。 \end{tabularx} \\ \\
[1ex]
\hline
\hline
\end{longtable}
$^{1}$頂枝妹早晨.
終於暖一點了.
經歷了比較寒冷的冬天.
不知道大家有沒有見過這東西呢.
是啊 醒目啊 煲湯的朋友應該會知道這是海底椰.
我也是第一次見它還沒切的樣子.
通常我們見到是切了的 切了右手邊的位置.
大家見到它的時候.
海底椰有什麼功用呢.
紫卡化痰.
紫卡化痰還有呢.
還有呢.
潤肺.
潤肺 沒錯.
其實它有很多功用.
它實際上有清熱潤肺紫卡.
滋陰補腎 山津止寒 美顏潤膚 含維他命BCE 維生素.
有助滋潤皮膚 提升免疫力 強身健體.
適合熬夜和體質偏熱的人士.
特別是在北美 經常開暖氣很燥熱.
在秋冬的時候很適合煲來喝.
所以現在暖化可能你不需要.
但是如果冷和燥熱的時候你就會用它.
海底椰是一種幫助我們轉化身體裡面燥熱的一個上好的藥材.
不知道你們知不知道.
北京都有海底椰.
這位海底椰平時不抬頭.
一抬頭就做了一件奇事.
轉化了人對於榮耀和羞恥的看法.
在藥房福音裡面.
人很急躁地保護自己的榮耀和面子.
而耶穌就是用我們生活中很平常.
最不起眼的事去觸摸我們的心.
甚至改變了我們人的一生.
我們一起來看看.
我小時候在教會聽這段經文的理解是這樣的.
水變酒 大家都聽過.
耶穌是一個可以改變物質的神.
可以將水H2O變成酒C2H5O.
多了很多不同的原子.

$^{41}$有碳也有氫.
耶穌真的很厲害.
成為一個可以改變物質的上帝.
不過今天我不是想用這個角度去看.
我想帶大家從另一個角度去看.
想想耶穌當時在藥房福音裡面的處境是什麼.
作者怎樣去描寫耶穌這個海底椰浦頭.
他做第一個神蹟.
大家就會發現這個上帝不單止是一個改變物質的神.
更加是改變人心的神.
我們一會細緻一點去看這段經文.
第三天在加利利的加拿大有一個婚宴.
耶穌的母親在那裡.
耶穌和祂的門徒都被邀請到宴會裡面.
酒用完了.
然後耶穌的母親對祂說.
祂們沒有走.
很奇怪.
瑪利亞對耶穌說.
祂們沒有走.
這句話好像很平凡.
但是很奇怪的地方就是.
瑪利亞明明知道耶穌不是賣酒.
而她老公就是耶穌的父親.
是一個木匠.
最多是會做木桶.
不會弄酒.
為什麼瑪利亞會對祂說.
祂們沒有走呢.
瑪利亞在期待什麼.
我們接著看下去就有些提示了.
就是耶穌說.
母親,我與你何干?.
我的時候還未到.
耶穌的回應有些奇怪.
祂反應有些大腿.
就是祂在問.
我和你瑪利亞有什麼關係.
而不是在問.
祂們沒有走關我什麼事.

$^{81}$我和你是什麼關係.
這個問題是有些奇怪的.
如果你細心想下去的話.
和瑪利亞的回應說.
祂們沒有走.
耶穌就好像很個人.
很個人的說.
我和你是什麼關係.
這樣.
究竟這裡發生什麼事呢.
這裡看到時候這個字.
我的時候還未到.
時候這個字其實是非常關鍵的.
以致到耶穌是這樣回應母親.
我們一起看下去.
母親對人說.
祂告訴你什麼.
你們就做了.
照猶太人潔淨的規矩.
有六口石缸擺在那裡.
每口可以乘兩三桶水.
這段是經文裡最平鋪直敘的地方.
但其實都是亮點所在.
這裡強調用了六口按照猶太人禮儀所擺放的石缸.
而石缸的大小是什麼.
是專門寫下去的.
可以裝兩三桶水.
證明這是一個頗大的石缸.
而這兩點都是為什麼耶穌會這樣回應瑪利亞的關鍵所在.
我們繼續看下去.
第七.
耶穌對人說.
把缸倒滿水.
他們就倒滿了直到缸口.
耶穌又說.
現在要出來送給燕湖總管.
他們就送了去.
燕湖總管就試了那些水所變的酒.
都不知道是從哪裡來的.
就跟那些要水的人說.

$^{121}$叫新郎來.
跟著就跟他說.
人家都是先放一些好酒.
等人家喝夠了才放一些次的.
為什麼先留這些好酒在現在呢.
這個就是耶穌在耶漢福音裡面所描述的第一個神蹟.
在加利利的加納恆.
是顯出了他的榮耀.
門徒就信他了.
這個事之後.
門徒母親和兄弟就一起到加伯隆住了幾天.
看完這段經文.
我們要收集一下剛剛所見的 所問的一些問題.
第一 南瑪利亞的提問.
他說他們沒有走了.
其實究竟他對耶穌說話.
說了一句話 期待著什麼呢.
第二 耶穌的反應其實有點奇怪.
就反問我和你有什麼關係.
提到時候還沒有到.
是不是在暗示著什麼呢.
第三 潔淨條例是這段經文的關鍵.
是會回答我們上面兩個問題的.
我們等下留意一下.
我們準備好就一起來了.
第一個問題就是.
瑪利亞對耶穌說他們沒有走了.
究竟期待著什麼呢.
在猶太文化裡面.
酒宴沒有酒喝呢.
或是在慶祝的婚宴裡面.
沒有酒喝是一個很大的恥辱.
會讓人覺得新郎的家準備失當失禮.
甚至失去了上帝的祝福的一個表現.
大家有沒有聽過葡萄園呢.
有聽過葡萄園嗎.
在舊約裡面經常提到葡萄園.
是上帝所給人的祝福.
以色列自古就是種葡萄和釀酒的重鎮.
所以葡萄酒對於猶太人來說.

$^{161}$其實是一個好意頭.
代表著上帝祝福這段婚姻.
你不夠酒簡直是一個公關的災難.
代表著這個家庭幾乎告訴別人.
他們可能不是很受到上帝的祝福.
這個是令到家族沒有面子.
受到屈辱的一個情況.
猶太人很注重面子.
那種榮辱的文化影響他們很深.
情況就好像.
大家都有聽過華人在婚宴上說的十個字.
主人家會說什麼.
吃著吃著上台再說話.
會說什麼.
酒微菜薄.
招呼不到.
大家多喝兩杯.
酒微菜薄.
大家經常都聽到這句.
主人家經常用來謙稱酒菜的品質是一半.
份量很少招呼不到.
只是謙稱來的.
試一下真是酒微菜薄.
喝酒可能整桌都只是一人一杯.
喝完就沒有了.
前菜上完之後就完.
主人家就會覺得怎樣.
很失禮.
招呼不夠.
所以大家都理解到.
猶太人的榮辱文化.
沒有酒喝.
很失禮.
而耶穌的母親瑪利亞.
很大可能是親戚或是很熟的朋友.
經文提到他們.
可以見到那些毒人準備走水.
他們應該是可以進入別人家內院的人.
可以看到整個婚宴的運作.
親戚或熟的朋友像現代的婚宴一樣.

$^{201}$是熟的人才可以去後台.
看到他們的運作.
然後看看有沒有什麼可以幫忙.
所以瑪利亞比其他人更在意主人家的面子和榮耀.
很敏銳地觀察到背後的難處.
所以當瑪利亞和耶穌說他們沒有酒.
他期待什麼呢.
他反映了一件事.
就是瑪利亞相信耶穌可以做一些事.
瑪利亞相信耶穌有能力幫主人家挽回面子.
或者避免失禮.
他期待著耶穌介入.
挽回他們的榮耀.
這裡大家理解到吧.
大概是這樣.
但是耶穌就回應了.
母親 我與你何干.
我的時候未到.
意思是我和你是什麼關係.
有些日本人可能會說.
我和這件事有什麼關係.
其實你看到原文.
其實很簡單.
我和你是什麼關係.
很直接.
就是這樣.
表達什麼意思.
你這樣也不明白.
你明知道我是彌賽亞.
而行神蹟會令彌賽亞的身份提前被人知道.
就是說耶穌本身是很低調.
時間到了才彰顯上帝的榮耀.
才做光芒四射的彌賽亞.
而瑪利亞的期望是.
不如你彰顯你的榮耀的時候.
不如幫幫你的家人吧.
在《約翰福音》裡面剛剛說到.
時候這個字是很關鍵的.
每一次說到時候的時候.
通常你都會看到第二個字.

$^{241}$就是榮耀.
時候和榮耀其實很關鍵.
就是因為耶穌知道.
當他一行神蹟的時候.
他就會將上帝的榮耀彰顯出來.
收也收不起來.
所以耶穌的反應不是問他們.
沒有走關我什麼事.
而是問母親.
你明知道我潛水.
我在海底.
在做海底事.
你現在叫我行神蹟.
而你和我的關係.
你還不知道我這樣做.
我要被人知道.
我就是那個彌賽亞.
所以大家明白我.
不過耶穌真是一個很乖的兒子.
之後他照做.
只是他之後做的事.
就打破了瑪利亞的想像.
耶穌不單止令主人家得到榮耀.
是令到他自己的榮耀都顯出來.
我們繼續看下去.
我們知道.
照猶太人的潔淨規矩.
有六口石缸放在那裡.
有些中英文譯本.
就會譯做是給猶太人行潔淨禮的石缸.
或者用於猶太人潔淨禮的石缸.
其實猶太人很注重潔淨禮儀.
因為他們知道上帝在川一級記和利未記裡.
頒佈了一些潔淨的條例.
是想和世人展示.
猶太民族是屬於耶和華.
屬於上帝.
這裡出現了一個理念.
就是猶太人.
透過上帝所頒佈的十誡和舊約的潔淨條例.

$^{281}$去到理解.
去衍生出來一些潔淨的禮儀.
目的是什麼呢.
目的是去到提醒.
用這個外在的舉動去提醒人類.
我是和上帝有關係的.
當人類做這個外在的舉動.
最理想的效果就是.
人知道而且深化自己的行為.
去到提醒自己.
自己屬於上帝.
以致我裡面的內心被轉化.
靠近上帝一點.
而不是單單做外在的行為.
而經文所描述的情況.
就是出現了一個斷層.
當猶太人很努力去做潔淨的禮儀.
還要用一般人.
還要大還要多的石缸.
這裡的石缸其實是很重和貴的.
通常出現在一些有經濟能力.
也很注重潔淨條例的家庭裡.
不是每一個家庭都會有的.
所以這一家人明顯是有頭有臉的有錢人.
為什麼經文要特意提到是兩三桶水.
是因為那個石缸夠大.
容納到的水夠多.
而且貴.
所以他們是非常注重榮辱.
他們覺得潔淨條例做得好的人.
就等於和上帝有好的關係.
得到上帝的祝福.
這個家族就好像在猶太文化裡面.
是一個根正苗紅.
明門正宗.
在乎的是什麼呢.
在乎的是人的榮辱.
而是多於和自己和上帝的關係.
多於自己的心態.
結果是怎樣呢.

$^{321}$結果是這樣的.
這個圖是反轉了的.
只是做外在的禮儀.
只是focus到潔淨禮儀的本身.
很忙碌地很注重外在的表達.
忽略了潔淨禮儀本身想要提醒人.
我屬於上帝.
我們的心要回歸上帝.
所以今天.
如果有機會.
希望大家還有機會.
去耶路撒冷的曲牆.
大家還會見到有些猶太人在守潔淨禮.
在曲牆裡面.
有一個洗手的地方.
洗手盆裡面有一隻銀色的杯.
他們就會用裡面那隻銀色的杯去舀水.
象徵著以前的一些缸.
他們去洗淨自己的污穢.
去到上帝面前才潔淨上帝雙手.
表明自己屬於上帝.
可想而知猶太人有多注重潔淨禮儀.
不過.
有沒有真正提醒自己屬於上帝.
這個是一人而異.
這個我們都要問自己.
所以耶穌.
偏偏用了這個行潔淨禮的石缸.
叫人裝滿水.
再用裡面的水變成酒.
他絕對可以用其他容器去裝水.
亦都不需要用這個石缸去裝水.
這個舉動就打破了猶太人和門徒對於潔淨禮的想法.
因為用來行潔淨禮的石缸其實是充滿著不潔的.
你想想.
月本的水是給人洗手.
洗臉.
甚至是洗腳.
你可以想像到人是從四方八面來.
走了很遠的路來到這個婚宴.

$^{361}$當時的科技亦都沒有那麼先進.
而很多人都是沒有鞋穿的.
穿了一些很粗糙的鞋.
你可以想像下.
那缸水有多麼的污糟.
有多麼的不潔.
裡面有多少泥有多少污漬.
而且客人已經到齊了.
酒都差不多喝光了.
所以那六口石缸.
很可能已經裝了幾輪洗腳水.
已經裝了幾輪.
已經有很多泥.
很多東西在裡面.
而現在怎麼說呢.
重新再裝滿它.
接著你就會見到.
裡面最污糟最不潔的水.
變成最好的酒.
比原本婚宴的酒更加好.
經文是這樣說的.
情況就是這樣的.
就好像參觀這些污水處理廠.
我曾經在新加坡參觀他們的處理廠.
其實他們做得很好.
廠長就帶著我們.
在一排污水處理的缸裡.
就是剛剛在別人家裡的廁所.
抽出來的污水.
接著就是這樣.
還有一股味道.
然後經過一系列的處理.
接著你就會見到.
原本啡黑色的水.
變成了很清澈的水.
接著廠長就拿起杯子.
就在我們前面舀起.
就喝.
接著就問我們.
你想不想試一下.

$^{401}$我們就有一個很大的心理障礙.
我們就覺得.
能不能喝呢.
沒人敢試.
就是這樣.
所以我們心裡都有種感覺.
耶穌不單止改變了物質.
他改變了人對於潔淨的.
一個黑板印象和錯誤理解.
不是那些水乾不乾淨.
不是那些洗完之後.
有多少我們叫做.
道德上的一些不潔可以洗掉.
而是重新將這個潔淨條例的本質.
連繫回人的心裡.
我們做這件事其實是怎樣.
我們回到上帝那裡.
所以面對著這個.
做歧視的上帝.
我們可以怎樣在.
我們屬靈的生命裡有回應呢.
這些石缸本身是.
給人做潔淨禮儀用的.
象徵著外在的一些.
行為,宗教儀式.
內心的一些污穢可以被潔淨.
但是耶穌卻將裡面的水變成酒.
象徵著他用他自己的生命.
他用他自己的血.
去做了人真正潔淨的根源.
凡是一些我們覺得.
最修家,最難開口.
最不潔淨的一些生命的片段.
無論是人生失敗.
貪玩無知而導致的後果.
上癮,破碎的生命等等.
那些都是我們人生裡面.
片段我們眼中都會覺得.
我都不想提.
我想都不想想.

$^{441}$就在我們人生.
很夜闌人靜的時候.
只剩下我和上帝的時候.
那些回憶就好像電影那樣.
影出來.
你想走都走不掉.
或者我們有沒有想過.
或者我們都不敢期望.
自己那些攀回手的一些經歷.
上帝是會做一個神蹟.
會將我們這段最想收埋.
最失敗,最不潔的人生片段.
轉化為我們人生的一些養分.
甚至是榮耀上帝的工具.
在基督恩典裡面我們可以學習.
將我們最不堪的人生片段交給他.
不是靠一些宗教儀式.
或者表面的善行去遮蓋而不息.
而是期待著耶穌基督在我們生命裡面.
做一次真正的轉化.
可能和大家都有講過.
其實我小時候有腦癌症.
上帝就醫治了我.
但是上帝醫治我的過程裡面.
他留下了一些東西給我.
就是一些我對於未來.
或者我在日常生活中的一些恐懼.
就是我很害怕突然間會暈倒.
那種恐懼有時候會吞噬我.
令我覺得原來我說話.
可能突然間我會說不到話.
可能我突然間會完全忘記一些東西.
這種恐懼對我來說是很大的.
醫治了我之後.
我仍然要面對著這個恐懼.
那些片段在當我一個人的時候.
經常都出現.
我其實完全想像不到.
我可以站在這裡和大家說話.
因為就算我上來那一刻.

$^{481}$我前面那一刻還有少少panic attack.
有少少恐懼.
好像我突然間會暈倒.
這個就是我一直都想像不到.
這個行歧視的上帝.
他竟然在我生命裡面.
告訴我一些我小時候最不想面對的東西.
他突然間扭轉說.
不是的 你可以用這個東西.
成為一個榮耀神的工具.
就好像楊憑虎的二章.
他做了一個歧視.
他竟然將一些最不潔淨的水.
一些可以用原汁氣的東西.
代表著我們一些過去.
我們覺得過去了就算的一些生命的片段.
轉化成為一個榮耀神的工具.
但問題是我們願不願意.
我們可以只是做一些表面的東西.
我們願不願意用一些最真實的生命去回應上帝.
當我們說到工作的時候.
我們不知道什麼時候開始.
我們都不是太渴望上帝在我們的職場加入.
每天上班你只是做千篇一律的東西.
有時候你可能要面對著一些巨嬰同事.
不懂工作就要煩你.
心裡變得很煩躁.
不知道什麼時候你要搞定手上的工作.
不知道什麼時候開始我們就會有些想法.
就是上帝應該不會在我們生命裡面.
在工作裡面出現.
我們唯一的盼望就是.
把今天的東西快點完成.
彷彿我們不再期待上帝在我們生命裡面.
我們的工作崗位出現.
不知道什麼時候我們不再期待.
上帝和我們的工作沒有什麼關係.
我只是賺錢.
我只需要這份工作去養家戶口.
除了這些之外.

$^{521}$其實這份工作沒有什麼其他意義.
而耶穌最驚訝的地方就是.
祂竟然把猶太人最嫌棄的洗腳水.
看起來最沒有意義的水.
用完即棄.
那些缸子甚至有錢人會把它打爛.
因為覺得已經沾了不潔.
已經沒有用了.
祂竟然把裡面的水成為了主人家的榮耀.
把自己的榮耀都彰顯出來.
同樣地.
祂都可以把我們生命中本身覺得沒有意義的工作.
成為一個榮耀祂的工具.
問題是你願不願意.
被上帝這樣介入你的生命和工作.
我們只是想.
不如我們今天完成我們的工作就算了.
我們可不可以多一點期望.
多一點想一想.
上帝如何在我們日常生活中.
我們的工作中去介入呢.
最後.
在實踐上我們.
內在生命已變得真實.
那些水本身是照著潔淨禮儀所預備的.
是救藥整套禮儀系統的一部份.
而耶穌其實是故意沒有用新的器皿.
是不是.
祂故意用潔淨的缸裡面的東西去做工.
祂不是想加堆外在的條件和規條.
而是把救藥一切的潔淨.
獻祭 節旗的安排都指向自己.
十字架上的救人.
對於救人來說.
這個神蹟就是一個記號.
我們不可以只停留在外在的儀式.
我的意思是.
我們不是不做外在的儀式.
今天你崇拜.
你會團契.

$^{561}$你會去祈禱會.
每一個我們正在做的事.
我們的宗教儀式 禮儀都在提醒我們.
我們是屬於上帝的.
預備好自己.
成為一個裝在上帝的榮耀器皿.
讓他們將普通的水.
普通的人生.
轉化成一個充滿喜樂和恩典的酒.
所以我猜到這裡.
大家都應該懂得怎樣回答這三條問題.
為何瑪利亞要對耶穌說他們沒有走.
就是因為瑪利亞期望.
耶穌要用神蹟挽回主人家的榮耀.
而耶穌的回應提到事後.
就是因為耶穌用神蹟介入的時候.
就會將上帝的榮耀彰顯出來.
而最後猶太人的潔淨禮儀.
是在提醒人屬於上帝.
耶穌用潔淨禮儀的缸裡面不潔的水.
轉化成最美的酒.
是向我們每一個人發出邀請.
用你最真實的生命去回應上帝.
多一些期待上帝介入我們的工作崗位.
多一些期待上帝在我們內在生命做轉化的神蹟.
讓我們一起祈禱.
天父多謝你.
多謝你親自向我們展示.
一些我們以為最不值.
最不值得提起的.
最羞恥的回憶.
求你幫助我們.
讓你在我們生命裡面.
做到神蹟的啟示.
讓你轉化.
讓我們多些期待.
期待你在我們生命裡面多些互動.
幫助我們.
我們禱告奉耶穌的聖名給予阿們.
\newpage



\section{雅各書 2:1-26}
\label{sec:031852}
\textbf{在信裡沒有TACO}~~~~司徒贊傳道
\newline
\newline
連結: \href{http://www.church.com.hk/acms/content.asp?site=cdc\&op=show\&type=product\&code=031852\&layout=sermon}{\texttt{031852}} ~~~~ 語音日期: 2026-03-22
\newline
\newline
\hyperref[sec:031851]{\small{< < < PREV SERMON < < <}}
~
\hyperref[sec:index_chronic]{\small{[返順時目]}}
~
\hyperref[sec:031853]{\small{> > > NEXT SERMON > > >}}
\newline
\newline
雅各書 2:1-26
\newline
\begin{longtable}{cl}
\hline
\hline
章節 & 經文 (和合本修訂版)\\
\hline
2:1 & \begin{tabularx}{0.7\textwidth}{X} 我的弟兄們，你們信奉我們榮耀的主耶穌基督，就不可按著外貌待人。 \end{tabularx} \\ \\ \relax
2:2 & \begin{tabularx}{0.7\textwidth}{X} 若有一個人戴著金戒指，穿著華麗的衣服，進入你們的會堂，又有一個窮人穿著骯髒的衣服也進去， \end{tabularx} \\ \\ \relax
2:3 & \begin{tabularx}{0.7\textwidth}{X} 而你們只看重那穿華麗衣服的人，說：「請坐在這裡」，又對那窮人說：「你站在那裡」，或「坐在我腳凳旁」； \end{tabularx} \\ \\ \relax
2:4 & \begin{tabularx}{0.7\textwidth}{X} 這豈不是你們偏心待人，用惡意評斷人嗎？ \end{tabularx} \\ \\ \relax
2:5 & \begin{tabularx}{0.7\textwidth}{X} 我親愛的弟兄們，請聽，神豈不是揀選了世上的貧窮人，使他們在信心上富足，並承受他所應許給那些愛他之人的國嗎？ \end{tabularx} \\ \\ \relax
2:6 & \begin{tabularx}{0.7\textwidth}{X} 你們卻羞辱貧窮的人。欺壓你們，拉你們到公堂去的，不就是這些富有的人嗎？ \end{tabularx} \\ \\ \relax
2:7 & \begin{tabularx}{0.7\textwidth}{X} 毀謗為你們求告時所奉的尊名的，不就是他們嗎？ \end{tabularx} \\ \\ \relax
2:8 & \begin{tabularx}{0.7\textwidth}{X} 經上記著：「要愛鄰 如己」，你們若切實守這至尊的律法，你們就做得很好。 \end{tabularx} \\ \\ \relax
2:9 & \begin{tabularx}{0.7\textwidth}{X} 但你們若按外貌待人就是犯罪，是被律法定為犯法的。 \end{tabularx} \\ \\ \relax
2:10 & \begin{tabularx}{0.7\textwidth}{X} 因為凡遵守全部律法的，只違背了一條就是違犯了所有的律法。 \end{tabularx} \\ \\ \relax
2:11 & \begin{tabularx}{0.7\textwidth}{X} 原來那說「不可姦淫」的，也說「不可殺人」。你就是不姦淫，卻殺人，也是成為違犯律法的。 \end{tabularx} \\ \\ \relax
2:12 & \begin{tabularx}{0.7\textwidth}{X} 既然你們要按使人自由的律法受審判，就要照這律法說話行事。 \end{tabularx} \\ \\ \relax
2:13 & \begin{tabularx}{0.7\textwidth}{X} 因為對那不憐憫人的，他們要受沒有憐憫的審判；憐憫勝過審判。 \end{tabularx} \\ \\ \relax
2:14 & \begin{tabularx}{0.7\textwidth}{X} 我的弟兄們，若有人說自己有信心，卻沒有行為，有甚麼益處呢？這信心能救他嗎？ \end{tabularx} \\ \\ \relax
2:15 & \begin{tabularx}{0.7\textwidth}{X} 若是弟兄或是姊妹沒有衣服穿，又缺少日用的飲食； \end{tabularx} \\ \\ \relax
2:16 & \begin{tabularx}{0.7\textwidth}{X} 你們中間有人對他們說：「平平安安地去吧！願你們穿得暖，吃得飽」，卻不給他們身體所需要的，這有甚麼益處呢？ \end{tabularx} \\ \\ \relax
2:17 & \begin{tabularx}{0.7\textwidth}{X} 信心也是這樣，若沒有行為是死的。 \end{tabularx} \\ \\ \relax
2:18 & \begin{tabularx}{0.7\textwidth}{X} 但是有人會說：「你有信心，我有行為。」把你沒有行為的信心給我看，我就藉著我的行為把我的信心給你看。 \end{tabularx} \\ \\ \relax
2:19 & \begin{tabularx}{0.7\textwidth}{X} 你信神只有一位，你信得很好；連鬼魔也信，且怕得發抖。 \end{tabularx} \\ \\ \relax
2:20 & \begin{tabularx}{0.7\textwidth}{X} 你這虛浮的人哪，你願意知道沒有行為的信心是沒有用的嗎？ \end{tabularx} \\ \\ \relax
2:21 & \begin{tabularx}{0.7\textwidth}{X} 我們的祖宗亞伯拉罕把他兒子以撒獻在壇上，豈不是因行為得稱義嗎？ \end{tabularx} \\ \\ \relax
2:22 & \begin{tabularx}{0.7\textwidth}{X} 可見信心是與他的行為相輔並行，而且信心是因著行為才得以成全的。 \end{tabularx} \\ \\ \relax
2:23 & \begin{tabularx}{0.7\textwidth}{X} 這正應驗了經上所說：「亞伯拉罕信了神，這就算他為義」；他又得稱為神的朋友。 \end{tabularx} \\ \\ \relax
2:24 & \begin{tabularx}{0.7\textwidth}{X} 這樣看來，人稱義是因著行為，不是單因著信。 \end{tabularx} \\ \\ \relax
2:25 & \begin{tabularx}{0.7\textwidth}{X} 同樣，妓女喇合接待使者，又放他們從另一條路出去，不也是因行為稱義嗎？ \end{tabularx} \\ \\ \relax
2:26 & \begin{tabularx}{0.7\textwidth}{X} 所以，就如身體沒有靈魂是死的，信心沒有行為也是死的。 \end{tabularx} \\ \\
[1ex]
\hline
\hline
\end{longtable}
$^{1}$今天講到的主題是「在信裡沒有Taco」.
是一個很好笑的題目.
什麼叫Taco呢?.
Taco不是能吃的.
Taco的意思是talk alone, cheating others.
就是說我們講信心,信仰.
不可以只有「講」字.
我們的信仰不可以靠嘴巴.
不可以自欺欺人.
在開始今天的訊息之前.
我們有一個禱告.
在當中可以經歷你的同在,你的帶領.
讓我們今天繼續在你的話語裡.
聆聽你的話語的時候.
求聖靈與我們同在.
開我們的眼睛,開我們的耳朵.
柔軟我們的心靈.
使我們不只是聽到你的話語.
而是被你的話語去感動.
去改變我們,改變我們的人生.
幫助我們去走出去.
走出去榮耀你.
得勝名致祈求.
今天我們來到看到《雅各書》.
因為《雅各書》有一句很出名的.
大家都背得出的.
「信心沒有行為,是死的」.
我弄一下這個,還沒開.
「信心沒有行為,是死的」.
也有人說,那保羅又怎樣呢?.
保羅說我們得救不是「本乎因也因著信」嗎?.
不是出於行為.
如果不是出於行為的意思.
是不是說我坐在這裡.
什麼都不用做就行呢?.
當然不是.
其實保羅和雅各是不是在吵架呢?.
一個說不是出於行為.
一個說沒有行為是死的.
其實他們是不是在吵架呢?.

$^{41}$當然不是.
我知道我已經說過這個大象.
但它好使好用.
我又拿它出來再說一次.
其實保羅在說什麼呢?.
其實保羅不是說.
口裡承認,心裡相信就行了嗎?.
我口裡承認,心裡相信就行了.
但其實重申一次.
在保羅的時代.
口裡承認不是只有一張嘴的.
在保羅寫羅馬書的時候.
口裡承認輕則沒了工作.
重則可以失去生命.
是那個口裡承認.
在那個時代.
我們口裡承認自己是基督徒的時候.
是有代價的.
這是保羅所說的口裡承認.
所以當保羅形容大象的時候.
他會說大象就好像一條蛇.
到雅各摸著大象尾巴的時候.
他會形容大象就像一條繩.
有些人就這樣看聖經.
就會覺得好像在前後矛盾.
究竟大象是像一條蛇還是像一條繩呢?.
是不是他們在吵架呢?.
是不是聖經前後有矛盾呢?.
其實不是.
保羅和雅各只不過是用不同的角度.
去描述同一個真理.
所以他們不是在吵架.
最後第三個人.
主耶穌的大弟子彼得.
中間這個人摸著大象.
好像一塊牆一樣.
我們看看彼得的角度.
由彼得的角度去看.
他怎樣去描述信心呢?.
原來信心只不過是一個開始.

$^{81}$信心是一個基石.
上個星期我們有九個youth浸禮.
我經常跟他們說.
浸禮是一個開始.
它不是一個結束.
浸禮是一個開始.
彼得也是這樣說.
開始之後又要有德行.
德行是什麼呢?.
德行就是character.
即是有少少的果子.
信心之後又要有德行.
德行之後又要有知識.
我們要讀聖經.
才可以明白神的心意.
要有知識.
我們一起要繼續去查經.
有了知識就開始要有節制.
為什麼呢?.
因為那些人懂得多就開始喜歡.
teaching mode.
lecture mode.
有時候我都要節制.
我對著那些youth.
不能夠天天說教.
天天都要教他們.
我的女兒有時候跟我說.
你再說我就.
往反方向去.
所以有了知識就要有節制.
有了節制之後又要忍耐.
要懂得忍耐.
忍耐是一個功課.
什麼是生忍耐?.
患難生忍耐.
所以去到差不多.
基督徒我們要經歷一些患難.
讓我們去練習忍耐.
走到這裡.
原來我們才真正.

$^{121}$去結出這個敬虔的果子.
敬虔原來.
Godliness.
英文叫做Godliness.
敬虔原來的意思就是.
走到這一步.
我們才真正的開始.
懂得去愛神.
要懂得去愛神.
才慢慢開始學習去愛弟兄.
即是回到教會.
我們才懂得去愛我們的弟兄姊妹.
要懂得去愛我們的弟兄姊妹.
我們才有這樣的心智.
走出去去愛眾人.
有這個大愛.
原來彼得將愛神和愛人.
放到結尾.
告訴我們這個大誡命.
不是那麼容易走出來的.
這個愛神愛人的大誡命.
不是那麼容易走出來的.
是一個屬靈成長的結果.
所以這個就是彼得的Perspective.
是彼得的看見.
所以我們總結一下.
主耶穌三個大弟子.
他們在說什麼.
其實他們是互相呼應.
不是互相打架.
原來保羅在說.
行為.
即是我們不可以靠善事.
去產生信心.
因為這個世界很多人都做善事.
做很多善事.
我們不能夠靠做善事.
去產生和神有一個關係.
有一個信心.
而雅各在說什麼呢.

$^{161}$雅各在說.
當我們和主耶穌.
和神有一個關係的時候.
必然會產生一個行為.
一個好的行為.
而彼得在說什麼呢.
彼得在說.
信心是一個開始.
贈禮是一個開始.
贈禮不是我們人生的高光時刻.
不是我們很成熟去贈禮.
贈禮是一個開始.
你有了信心之後.
你才可以繼續去成長.
在主裡面成長.
不是靠我們自己去成長.
所以很多人手做的東西.
是很漂亮的.
靠人手做出來的東西.
輪廓是很分明.
影子是很漂亮的.
而一些真實的影子.
自然流露出來的影子.
你反而一看下去.
都不知道本相是什麼.
你是不知道它是什麼.
自然流露出來的影子.
如果我們對比一下.
人手做出來的影子.
有時候是更加漂亮.
光鮮.
就好像這個世界.
很多慈善機構.
一拿出來就一大張支票.
很漂亮.
很光鮮.
但是.
人手做的東西有一個漏洞.
就是.
一定要有一個聚光燈.

$^{201}$一定要打一盞聚光燈下去.
一定要有一面很漂亮.
很平整的牆.
它才會營造出這個影子的效果.
而一些自然流露出來的影子.
有時候是沒那麼漂亮的.
但是.
自然流露出來的影子.
卻是不需要有聚光燈的.
你晚上哪怕月亮有一點光.
就已經反映出那個影子.
很少好事的光.
就可以反映出.
這個自然流露出來的影子.
而這個影子是更加立體.
更加有力量的.
我感恩活因教會走到現在.
我們出來的意象是一片青草地.
我感恩到今時今日.
我們都是維持得很好.
這一片青草地.
而我更加珍惜的是.
在這個教會裡.
無論我和年輕人也好.
到小組也好.
我都有機會和大家分享我的軟弱.
不單止分享.
不單止分享我的軟弱.
也可以被人接納.
不單止可以接納.
弟兄姊妹還可以有一個角色.
就是可以督察我.
這是很難得的.
你可以在小組裡.
分享你自己的軟弱.
不單止被接納.
還要有一些很成熟的弟兄姊妹.
可以保持你負責任.
這是最寶貴的.
有些人聽完就算了.

$^{241}$為你祈禱.
但是如果你們身邊真的有弟兄姊妹.
她可以督察你.
可以指正你.
而你又覺得很舒服.
這個時候.
我們就真的可以活出這個信仰.
所以我鼓勵.
我們當中還沒有落到小組機會的弟兄姊妹.
要參加這個小組.
教會的其中一個目的是什麼?.
激發愛心.
我們要在小組裡.
坦誠流露自己的軟弱.
被接納.
被保持責任的時候.
我們才可以為神.
作出更加多的工作.
因為其實雅各書說的行為.
不是我們做善事.
不是我們走出去做什麼.
這個行為是我們信心自然.
反映出來的果子.
今天我們有四點.
Tako其實有四個意思.
我整理一下.
Tako比較容易記.
第一個.
在信裡面.
是沒有talk alone.
在信裡面.
不是只有講字.
在信裡面.
是沒有afraid.
在信裡面.
是沒有懼怕.
第三點.
在信裡面.
是沒有calculation.
沒有計算.

$^{281}$第四點.
在信裡面.
是沒有退路的.
沒有opt out.
我們來看看第一點.
不要大隻說.
二章十四節那裡這樣說.
雅各說.
我的弟兄啊.
如果有人說.
自己有信心.
卻沒有行為.
有什麼益處呢?.
這信心能夠他嗎?.
其實是一個諷刺的問法.
雅各其實是說.
這信心.
豈能夠他嗎?.
豈能的意思.
是什麼?.
就是不能了.
這樣的信心.
只有講字的信心.
是不能夠他的.
而這樣的message.
重要的事情.
三次以上.
雅各在十七節到二十六節.
短短幾句經文.
重複了三次.
信心沒有行為是死的.
不是說了一次.
說了三次.
信心沒有行為是死的.
嘩.
他不是輕輕帶過.
他是鄭重的說了三次.
鄭重的說了三次.
我們就要問.
其實什麼叫做行為?.

$^{321}$什麼叫做行為?.
在雅各書裡面.
那個行為.
他不是指明.
你真的要做什麼.
你要做ABCDEFG.
雅各書不是這樣說.
反而他說了很多.
你不要去做什麼.
所以行為不是你去做什麼.
更多的是你不要去做什麼.
我們來看看.
大丈夫有所為有所不為.
雅各書有很多不所為的例子.
第一叫我們不要偏心待人.
不要覺得有些人有錢.
就對他好一點.
不要用口傷人.
我們要勒住我們的舌頭.
不要用口傷人.
最重要在中間夾住的.
就是不要言行不一.
不要說一套做一套.
不要心存妒忌.
不要製造紛爭.
而這個不要言行不一.
這個很重要夾在中間.
很重要很重要.
什麼叫做言行一致呢?.
Integrity.
我經常教導年輕人.
什麼叫做Integrity呢?.
就是當沒有人看到你做事的時候.
你依然做對的事.
沒有人看到你的時候.
依然做對的事.
要言行一致.
很多人說下一代.
Gen Z最大的struggle.
為什麼不回來教會呢?.

$^{361}$就是他看到教會的一些前輩.
說一套做一套.
這個是最令他們不想回教會的.
一個最大的因素.
而我們如果真的想傳承.
將這個信仰傳給我們下一代.
我們最緊要做的事.
就是盡量盡量做到言行一致.
不是說你說得有多漂亮.
你可以說到自己不漂亮.
但你內外要一致.
The power of integrity.
is more powerful than beautiful words.
有時候我們說的話可以很漂亮.
但我們做的事.
是不是跟我們說的話是一致呢?.
這個真的很重要.
很重要很重要.
上星期我們有九個新生命在教會誕生.
榮耀歸給主.
五年前當神父召我去事奉他們的時候.
我真的不覺得自己可以做到什麼.
我一直以來都完全不覺得自己可以做到什麼.
但看到每一年.
每一年都有一批一批的年青人跟隨著.
我們看到中間的核心力量是我們的團隊.
我們有UK,我們有Ricky.
今年我們有Kitman,我們有Kathleen在當中事奉.
去守護他們,去陪伴他們.
我們不是說在中間做了什麼激勵人心的事.
不是說經常都要激勵他們.
但你看到Candace分享.
他們讚美神的是法治內心.
而這個成果不是我們在當中做了什麼.
而是靠每一個大人在裡面事奉的時候.
是言行一致.
他們看到的.
兩年前有兩個年青人走來問我.
為什麼Pastor Jackie在廁所抹桌面.
他不是一個Campus Pastor來的嗎.

$^{401}$真人真事.
我看著他說.
你知道我們說我們是Servant Leader.
不是只有港字.
這些對他們來說是很有力量.
最容易大濟港的是我們.
港就很容易.
但盡量盡量在生活當中守住每一個見證.
不要在他們面前做到大濟港.
還有一次我記得我在港台上.
你知道我有時候說話都不經大腦.
都容易得罪別人.
有一個年青人早上跟我分享了他的故事之後.
我就順口把他的故事說出來.
完了之後.
很好啊.
這個世代的年青人很直接.
完了之後他就跟我說.
Confront我.
他說Pastor Zen, I don't appreciate it.
為什麼你把我的故事這樣分享出去.
我覺得你很不尊重我.
我看著他幾秒.
我立刻有一點想辯護我自己.
但是不行.
我要說I am sorry.
其實我練習I am sorry練習了兩個階段.
第一個階段就是.
I am sorry law.
第一個階段是I am sorry law.
我經歷了一個階段.
要煞住按鈕不要說那個law字出來.
英文就是OK, I am sorry.
煞住按鈕不要說那個OK.
I am sorry就是I am sorry.
不要說那個OK.
我說I am sorry.
接著那兩個字叫No excuse.
No excuse.
我的確是需要問過你得到你的同意.

$^{441}$才去分享你的故事.
其實很難說的.
因為大家都要面子.
煞住按鈕.
I am sorry, no excuse.
這樣就保存不了跟他的關係.
我想我們大人真的要學習.
跟他們有錯就承認.
有錯就承認.
打就站定.
一起將這個信仰.
不要只有一個港字.
保護下去.
但義理也是.
大家看到他是信心偉人.
獅子坑都不怕.
整本但義理書其實比較出現得多的.
其中一個字叫什麼.
叫Stand.
Stand 時立 企硬.
但義理是靠企硬.
去保護他的信仰.
他不是很積極的說要去Lobby.
要爭取一些什麼.
他不是說四代的君王.
兩個政權四代君王.
他依然站在那裡.
歷史裡面有多少個人物.
沒有多少個.
但他靠什麼.
去經歷這四代君王.
他靠什麼.
他靠企硬.
他沒有走出來說.
我要反對同性戀.
我要反對墮胎.
當然你做這些都很好.
你有呼召就很好.
但義理他沒有這樣做.
他只不過是一次一次.

$^{481}$堅守自己的信仰.
有人叫他吃肉.
我不知道你的肉拜祭過什麼.
我不吃.
有人叫他不要禱告.
現在不可以禱告.
他繼續去禱告.
堅守他的信仰.
那個義夢.
皇帝發的那個夢.
叫他解夢的時候.
他站硬.
他將榮耀歸於神.
不是榮耀歸於自己.
而他知道這個夢的細節的時候.
很不吉利.
我要說出來.
分分鐘真的被殺頭.
他依然站硬.
神叫他做的事.
站硬就是靠站硬.
這個就是但義.
所以有時候我們去信主的時候.
我們說要活出我們的信仰的時候.
不是我們要做多少事.
而是在我們生命當中.
能不能夠去保護.
保守我們每一個見證.
靠站硬.
靠言行一致.
言行一致的信仰才是活的.
今年我們說的主題是.
非凡人生.
活出主因.
都是靠保守我們的信仰.
站硬 言行一致.
我們走出這個社區.
才可以將神的榮耀活出來.
可以活出主因.
這是第一點.

$^{521}$第二點我們說到.
在信中沒有懼怕.
沒有afraid.
因為二章十九節.
我們一起讀.
一二三.
你信 神只有一位.
你信的不錯.
魔鬼也信.
卻是賤鯨.
所以我們信神.
不可以只是說我相信神存在.
就夠了.
因為魔鬼都相信神存在.
是不夠的.
不可以只是說我們相信神存在.
不可以停在那裡.
魔鬼都信.
但卻是賤鯨.
我們不是要賤鯨神.
我們不是要afraid of God.
但我們要敬畏神.
我們要敬畏神.
才是智慧的開端.
敬畏耶和華是智慧的開端.
有什麼分別呢?.
為什麼我們不要懼怕.
但要敬畏神呢?.
懼怕就是afraid.
連魔鬼都懼怕.
我們要admire.
我們要敬畏神.
原來敬畏神的意思是in awe of God.
admire.
為什麼我們回來要敬拜呢?.
為什麼我們要唱詩歌呢?.
在詩歌當中.
我們要學會admire.
admire God.
剛才Candace鼓勵大家.

$^{561}$唱得大聲一點.
不是說我們要展露我們的歌喉.
而是我們要express.
我們對神的愛.
詩篇都這樣說.
我們要讚頌耶和華.
使萬國都看到祂的榮耀.
當我們admire神的時候.
世人是會看到的.
有新朋友在我們當中.
他會看到神的榮耀.
因為他看到神的子民.
多麼的admire他.
是不是很抽象?.
最近發生了一個故事.
令我更深刻體會到.
什麼是admire.
話說有一個星期日.
我和我女兒和老公.
去三藩市.
探望我的老爺奶奶.
吃完飯之後.
三藩市.
不要浪費了.
下街吃個糖水.
多麼好呢.
竟然我們提出.
下街吃糖水的時候.
我女兒說我不去.
為什麼?我要留在這裡.
學中文.
你不去吃糖水?.
你留在這裡學中文?.
從小到大.
我們這個年代長大的.
要練這些字的時候.
是什麼時候會練這些字?.
旁邊要有一支東西.
才會學練字.
為什麼.

$^{601}$ASA會有這樣的心.
留在這裡寫中文.
不去吃糖水?.
很快.
他自己就說了.
我要留在這裡陪爺爺.
因為爺爺不可以.
下街吃糖水.
所以我們去糖水.
大概多少時間?.
最快至少一個小時.
我們整整去了兩個小時.
回來看到他還是在練字.
還在練字.
為什麼呢?.
以我.
我是不可以讓他這樣做.
一定是不可以.
我知道沒有辦法.
但他知道他不是經常來探望爺爺.
爺爺不可以下街吃糖水.
爺爺想他學中文.
於是他留下去學中文.
這個不是一個虛構的故事.
是一個十五歲的小朋友.
他甘願做的一件事.
而這個真真正正.
是爺爺很愛他.
爺爺很疼愛他.
而他知道爺爺很疼愛他.
所以他留下去學中文.
是不是他想學中文呢?.
一定不是.
這個女兒我明白.
我知道.
但因為他愛爺爺的緣故.
他會去做一些.
本身不是他想做的事.
這個是冰山一角.
我們跟神的關係都是這樣.

$^{641}$有很多事不是我們想做的.
但當我們去回應神的愛的時候.
我們就會做一些.
平時我們不想做的事.
但願今天我們跟神的關係.
不是一個戰經的狀態.
不是我們回教會.
就說我得不得救.
我有沒有得罪神.
但願我們回來教會.
跟神親近.
或者在日常生活中.
跟神親近的時候.
不是一個害怕的狀態.
不是一個永遠都說.
我得不得救的狀態.
不是這樣.
有時我們都會害怕.
去到天家耶穌說.
我奉你的名趕鬼.
奉你的名趕盜.
參加小組奉你的名.
去大這大那的耶穌.
會不會突然說不認得我呢?.
會不會啊?.
就等於Asa和爺爺一樣.
有一個真正的愛的關係.
Asa不會害怕.
有一天去見到他爺爺的時候.
他爺爺突然說不認得他.
不會的.
所以願意我們跟神的關係.
是一個敬畏的關係.
為什麼有些人在敬拜裡面.
真的會感受到神的同在呢?.
其實很多時間.
好像彈耳鈴.
不是只有講字.
你去俯伏敬拜神的時候.
你用你所有的力氣.

$^{681}$去敬拜神.
你Admire你愛神的緣故.
不是說我喜歡唱歌才唱歌.
而是.
你回來敬拜神.
用你最後那口氣去唱歌敬拜神的時候.
你就在當中經歷.
神的同在.
神的同在經歷了之後.
你跟神有一個關係.
你就不會整天害怕.
我得不得救.
不是得不得救的問題.
而是有多親密的問題.
所以願意大家回來.
我們真的傾出我們的所有.
我們去敬拜.
遇見神 被神去改變我們.
第三點.
在信裡面是沒有Calculation的.
或者有些人說.
你有信心 我有行為.
請把你沒有行為的信心.
指給我看.
我就藉著我的行為.
將我的信心指給你看.
什麼意思呢.
意思是有些人將信心和行為.
分得很開.
信心還信心 行為還行為.
但雅各就說.
當然不是啦.
信心和行為怎麼可以分得開呢.
奶茶走奶.
你見過沒有.
有沒有人這樣Order.
你有沒有見過沒有奶的奶茶.
有沒有見過沒有茶的奶茶.
又不會的嘛.
有些人爭論.

$^{721}$如果在茶裡.
我滴一滴奶.
算不算奶茶.
如果你斤斤計較的話.
你就給人喝一喝.
一喝下去.
茶來的 沒有奶的.
你滴一滴茶在奶裡.
給人一喝就知道.
是茶還是奶.
你強行爭論都可以.
但不可以這樣爭論.
我們做基督徒回來.
不是要沖一杯合格的奶茶.
我們要沖一杯好喝的奶茶.
是不是.
煲湯.
煲湯也是.
可能不懂煲湯的時候.
大家覺得這煲湯會不會好喝呢.
很多料的.
像吃火鍋.
是不是.
是不是煲湯越多料.
就越好喝呢.
真的不是.
我剛剛信主的時候.
我曾經憑我這個.
很火熱的嘴.
嚇怕了很多人.
很多人都怕了我.
接受了我福音的.
可能只有三個.
一開始我追著別人說.
有些人說你再說就和你絕交.
有時候不是你說得多.
做得多.
不是說我聽完這個訊息.
我出去就捉人回來.
不是說這樣就好喝.

$^{761}$我要為Like的姊妹.
做廣告.
煲湯的人.
我們星期五有姊妹組.
煲湯有經驗.
當然我們每一個小組.
都有很好的弟兄姊妹.
很成熟的弟兄姊妹.
煲湯不是說你越多料.
就越好喝.
你也要控制火候.
其實我們.
教會煲湯最厲害的.
可能是一個弟兄.
就是你.
每一步.
都不可以.
skip.
你有心煲湯才好喝.
肉要處理.
我煲湯不厲害.
但我知道你懶.
你做漏一樣.
煲湯不好喝就煮了.
煲湯我講到這裡.
因為我煲湯也不是很厲害.
真的要控制火候.
要有心.
雅各書也是這樣說.
如果有弟兄姊妹.
其實他也很落地.
雅各不是說一個你不認識的人.
是說教會裡的弟兄姊妹.
赤身露體.
又缺了日用的飲食.
你們中間如果有人跟他說.
平平安安地去吧.
願你穿得暖吃得飽.
卻不及他們身體所需的.
有什麼益處呢.

$^{801}$其實雅各說的.
祈禱.
祈禱不是不好的.
祈禱是很好的.
沒有心肝的祈禱.
就不好.
沒有心肝的祈禱.
隨便的祈禱是不好的.
其實祈禱大有功用.
因為以利亞和以利沙.
在《遼王記》也說.
他們是靠祈禱.
使人復活.
但他們是怎樣祈禱的呢.
以利亞是將身體.
伏在男子身上.
什麼意思呢.
以利沙在祈禱的時候.
是眼對眼.
口對口 手對手.
其實是什麼意思呢.
他們在祈禱的時候.
將自己的生命.
完全與男子一致.
他將自己.
代入男子的人生.
他不是隨便為他祈禱.
而以利沙.
他在房間裡來回走來走去.
祈禱走來走去.
很用力的為男子祈禱.
男子才得到神的憐憫.
得以復活.
所以原來我們為人祈禱的時候.
不是那麼輕易的.
願你平平安安地去.
你感冒了.
回去多喝兩杯水.
不是那麼輕易沒有心肝的祈禱.
是你將你的生命.

$^{841}$與這個人一致的時候.
你為他祈禱.
這個祈禱是大有功用.
但說到今天.
沒有人是.
需要吃.
沒有人是.
吃不飽穿不暖.
大家都很有充足的供應.
不需要吃米飯.
我們這個世代.
所以今天怎樣應用呢.
這段經文.
今天的應用.
不是說我們.
要煮多少東西給別人吃.
我們昨天Tracy.
吃到後來有很多東西吃.
我們又浸泥.
叫人吃.
又打包.
大家都皺眉頭.
所以今天關心別人.
不是說我們給別人東西.
不是說我們煮多了米飯.
他感冒.
感冒其實好像說不能喝雞湯.
大家知不知道.
不是說我們煮多了什麼.
不要了吃剩了.
那些我們才拿給關心別人.
不是這樣的.
他們身體所需要的.
其實今天我們教會.
弟兄姊妹.
的身體其實.
不是說有多少需要.
我們今天真正需要的是什麼呢.
是.
Mentally.

$^{881}$今天我們可能.
教會裡面很多弟兄姊妹.
在這裡有需要.
其實這裡都影響身體.
Mentally都影響Physical.
我們怎樣能夠.
將我們的生命.
和弟兄姊妹的生命.
他們的struggle.
變成我們的struggle.
one on one.
去給予弟兄姊妹.
他們所需要的呢.
其實我們作為.
牧者是很欠缺的.
尤其是我自己.
我自己的gift可能.
在關顧這裡.
是有限的.
我經常和自己說.
我最大限度.
one on one 眼對眼.
口對口 手對手的關懷.
可能是同一時間.
不可以超出三個姊妹.
弟兄我不可以關心你.
因為這樣去關心你.
就不合體統.
但我同一時間.
只可以關心.
最多是三個姊妹.
但是.
我們教會有這麼多的.
弟兄姊妹.
那怎麼辦呢.
我建議去挑戰大家.
不是所有人都要關心.
不是所有人都要為他祈禱一次.
你能不能在教會當中.
找到一個.

$^{921}$一個.
其實因為你在家裡.
已經很多人要關顧.
你自己老公上有爸爸媽媽.
又有子女.
很多人要關心.
但你可不可以在教會另外再找多一個.
不要多一個.
你真是去體恤他的需要.
知道他的壓力.
知道他的掙扎.
而去陪伴他.
因為其實他不是需要.
你給他什麼.
他是需要你去明白他.
陪伴他.
因為家替書都是這樣說.
你們各人的重擔要互相擔當.
那可能你就說.
那我自己呢.
我自己都已經很多掙扎了.
我又上班.
上班又有上班的壓力.
爸爸媽媽又病了.
那我自己怎樣.
原來就是說.
如果你願意去關心.
這個人的時候.
那個人又願意關心你.
因為教會是一個互相.
一起生活的地方.
一起去分擔重擔的地方.
有時候你願意踏出那一步.
你去先關心這個人的時候.
別人來到關心你的時候.
才可以填補到.
你那個.
生活不樂的地方.
有時候你以為你要先照顧好自己.
但是有些地方.

$^{961}$是你照顧不到你自己的.
你要靠弟兄姊妹.
來幫你一把.
你可以幫到弟兄姊妹的地方.
你先幫了他.
然後讓他來幫助.
你需要的地方.
這個就是互相去擔當.
如果大家都做到這一點的時候.
其實真的我們可以去.
減輕大家的擔子.
所以在信裡我們不要去計算.
不是說我們付出多少.
而是我們用真心去對待.
最後一點就是說.
在信裡沒有up out.
我們一起讀這個經文.
1,2,3.
因為凡是專守.
所報律法的.
只要在一條上失足.
就歸犯所有的了.
就像那說.
不可姦淫的.
也說不可殺人.
你縱然不姦淫.
卻殺人也是違法的.
犯法的.
就是說我們不可以去選擇.
去守哪一條律法.
不守哪一條律法.
就像我們在Costco買東西.
你看到一條香腸.
很好吃.
你能不能只買一條香腸呢.
不可以的.
你要買完12打.
一大盒.
你不可以選擇.
要不就不買.

$^{1001}$要不就買完.
主人蘇也很有恩典憐憫.
他給我們的.
考試題目很簡單.
已經一早就告訴我們.
預先告訴我們.
他會考甚麼.
課題是甚麼.
而課題就是.
只有兩條題目.
PJ也有說過.
借用PJ的例子.
兩條題目.
要看完兩條題目.
我們真的要去學習.
去愛神.
然後去愛人.
而彼得這樣說.
又要信心又要有空閒.
又要有知識.
又要有節制.
又要有忍耐.
經過所有這麼多的東西.
原來才真真正正.
有這個敬虔.
才懂得去愛神.
才慢慢懂得愛弟兄姊妹.
然後才慢慢愛社區的人.
外面的人.
這個是真正的題目.
而甚麼是假的題目呢.
假的題目呢.
有時我們真的有個假象.
以為回到教會.
倒茶倒水.
做這樣做那樣.
我們以為我們做了這些東西.
就是滿足了.
神給我們的呼召.
其實.

$^{1041}$神真真正正要我們做的.
不是這些東西.
不是說我們回來.
做了多少的東西.
而是透過我們做這些東西的時候.
我們有沒有真正的去.
做到愛神和愛人.
這個功課.
所以我們要.
check both boxes.
不是其中一個box.
所以信心不是一個選擇題.
它是一個你信服.
還是你不信服的題目.
就好像亞伯拉罕一樣.
亞伯拉罕不是因為.
行為稱嚴.
我們的祖宗.
亞伯拉罕將兒子獻上.
這個就是一個行為.
而這個行為.
是一個信心的回應.
不是他做了甚麼.
是一個信心的回應.
我專門將.
奉獻的禱告留在.
這個時刻.
弟兄姊妹我們奉獻.
不是說十一奉獻.
奉獻是說.
不是我們用完用剩.
我們夠了.
將那些東西奉獻給神.
原來奉獻是奉獻你最愛的一件事.
你覺得最珍貴的一件事.
你覺得真的.
原來奉獻給神.
是會有一些令你覺得.
ouch的感覺.
不知道你最近的奉獻給神.

$^{1081}$有沒有令你有少少ouch.
如果你有.
神是會祝福你的.
不是說錢.
是.
你最珍貴的是甚麼.
你將這東西奉獻給神的時候.
神是會multiply.
將你的祝福打開天窗.
去祝福你的.
那時候.
剛剛開始會恩教會的時候.
我們知道.
神是呼召我們.
來建立這間教會的時候.
我的女兒.
Asa.
大概當中有些阿姨.
都知道.
對她來說不是一個最好的改變.
因為如果她來到這間教會.
就意味著.
她平時在教會.
最好的朋友.
她都是會見不到的.
如果來這間教會的時候.
當時.
不是說我一個.
無心無肺的媽媽.
我都有關心過她.
但我知道.
如果我清楚聽到.
這是神呼召我做的事.
我不可以跟著.
女兒的喜好去走.
因為.
爸爸媽媽.
我們的being是很重要.
我們的呼召是很重要.
神對你的呼召.

$^{1121}$是高過.
你子女的喜好.
我們很多時候.
都會將子女的喜好.
先走.
她喜歡哪一間教會.
我就跟著她去哪一間教會.
但其實不是這樣的.
我們父母要保護好.
自己的being.
我們的calling.
我們的呼召是什麼.
當我們願意將神的呼召.
先走的時候.
神一定祝福.
神必定祝福.
雖然這等待是長了一點.
等了半年.
頭半年.
我的女兒.
不是很開心.
神不守舍.
但當我們堅守.
這個呼召的時候.
神蹟是會發生的.
到了退休會之後.
我看到我女兒.
投入了很多.
完全將credit交給神.
我們.
認定自己呼召的時候.
神必定祝福.
我們不要讓子女的喜好.
直接影響我們的呼召.
我們的being很重要.
因為我們帶領他們.
他們在我們身上看到.
言行一致的信仰的時候.
他們必定會被信仰帶來.
最大的影響.

$^{1161}$最後用這句經文挑戰大家.
馬太福音16章25節我們一起讀.
123.
因為凡要救自己生命的.
必喪掉生命.
凡為我喪掉生命的.
必得著生命.
我想挑戰大家.
我們今天不要去守護自己的preference.
不要在神面前說我要這樣我要那樣.
我們先將我們的生命.
交給神.
神必然會打開天窗祝福我們.
我們一起去禱告.
你每天和我們去生活.
你每天帶領我們.
你每天在我們生命裡給我們很多的憐憫.
你一次一次的去原諒我們.
我們做錯事的時候.
你告訴我.
你依然愛我們.
幫助我們去回應你的愛.
幫助我們.
去帶著信心.
在我們的生活當中.
也是有力量的.
去作為一個管子.
被你去使用.
將你的光.
真的可以reflect on this community.
Thank you God.
Amen.
\newpage



\section{以賽亞書 37:1-4}
\label{sec:031853}
\textbf{危難當前}~~~~李應新牧師
\newline
\newline
連結: \href{http://www.church.com.hk/acms/content.asp?site=cdc\&op=show\&type=product\&code=031853\&layout=sermon}{\texttt{031853}} ~~~~ 語音日期: 2026-03-22
\newline
\newline
\hyperref[sec:031852]{\small{< < < PREV SERMON < < <}}
~
\hyperref[sec:index_chronic]{\small{[返順時目]}}
~
\hyperref[sec:031854]{\small{> > > NEXT SERMON > > >}}
\newline
\newline
以賽亞書 37:1-4
\newline
\begin{longtable}{cl}
\hline
\hline
章節 & 經文 (和合本修訂版)\\
\hline
37:1 & \begin{tabularx}{0.7\textwidth}{X} 希西家王聽見了，就撕裂衣服，披上麻衣，進了耶和華的殿。 \end{tabularx} \\ \\ \relax
37:2 & \begin{tabularx}{0.7\textwidth}{X} 他差遣以利亞敬宮廷總管和舍伯那書記，並祭司中年長的，都披上麻衣，到亞摩斯的兒子以賽亞先知那裡去。 \end{tabularx} \\ \\ \relax
37:3 & \begin{tabularx}{0.7\textwidth}{X} 他們對他說：「希西家如此說：『今日是急難、懲罰、凌辱的日子，就如嬰孩快要出生，卻沒有力氣生產。 \end{tabularx} \\ \\ \relax
37:4 & \begin{tabularx}{0.7\textwidth}{X} 或許耶和華－你的神聽見亞述將軍的話，就是他主人亞述王差他來辱罵永生神的話，耶和華－你的神就斥責所聽見的這些話。求你為倖存的餘民揚聲禱告。』」 \end{tabularx} \\ \\
[1ex]
\hline
\hline
\end{longtable}
$^{1}$今天我們看這本書,剛剛讀了37章.
第一章講到希西加王.
這裡講到以賽亞先知服侍了四個王.
其中一個是希西加王.
希西加王聽見司禮伊弗.
聽見什麼呢?.
其實是整個36章.
所以真的要揭出來.
特意拋開沒有的.
有聖經的話就揭聖經.
36章.
他聽見什麼呢?.
是36章所發生的事情.
很多事情發生是36章.
發生了什麼呢?.
我們揭下去看看.
剛才紅色的繩子.
揭出來就是36章.
剛才找人揭了給你.
希西加王十四年.
亞述王西那基納上來攻擊猶大.
一切堅固成章.
成功取原來亞述王的大軍.
來攻擊很多個猶大南國的城市.
攻了很多.
現在圍困耶路撒冷.
亞述王從拉吉差遣拉伯沙基.
拉伯沙基是將軍.
或者說是重要馬仔.
西那基納王的馬仔.
拉撒巴基率領大軍往耶路撒冷.
到希西加王那裡.
他站在上池的水溝旁.
在漂泊地的大路上.
原來耶路撒冷.
我也是看這些片的.
是山城來的.
他需要一些水源.
就是水的資源.
有一個水溝的地方.

$^{41}$用來儲水給耶路撒冷.
他就在漂泊地的大路上.
就要跟希西加王那裡說話.
於是希臘加的兒子加載.
這個是官名.
就是他的內臣.
以利亞敬就是他的臣子希西加王.
並書記謝伯拿和阿薩的兒子史官.
就是歷史的那個叫約瓦.
出來見拉伯沙基.
就是亞述王.
派了一個馬仔叫做拉伯沙基.
他想去見王.
王沒有見他.
就派一些馬仔來見他.
就是剛才那堆馬仔來見他.
在哪裡?.
在城門口那裡.
拉伯沙基對他們說.
你要告訴希西加說.
亞述大王如此說.
你所倚靠的有什麼可倚仗呢?.
有什麼障礙呢?.
你說你有打仗的計謀和能力.
我看不過是虛話.
你到底倚靠誰才背叛我呢?.
看你所倚靠的埃及.
是那壓箱的偽仗.
人若靠這仗.
就必刺透他的手.
原來這個猶大王串通埃及王.
打算叛變或者不聽.
之前他們是向亞述帝國進貢的.
但他打算拍另一個碼頭.
那個叫做埃及.
所以就不聽.
打算不聽亞述王的話.
所以就叛變.
現在這個將軍拉伯沙基說.
你靠誰你夠膽.

$^{81}$竟然靠他其實是一個折箱的偽仗.
折斷了你靠他戳下戳下.
平時有些手掌戳下戳下.
斷了的嘛.
即是鞋手.
靠他玩你不就刺到你的手了.
他的意思是這樣.
你靠一些不靠得住的東西.
你死定了.
大概這樣說.
你竟然背叛我靠埃及.
你還不死.
大概是這樣的意思.
我們跳地看下去.
第八節.
現在你把當頭給我.
主亞述王給我一些抵押品.
我給你二千匹馬.
看你這一邊騎馬的人夠不夠.
就是說他都很驕傲.
我就算送二千匹馬給你.
看你都不夠那麼多人懂騎.
行不行.
如果用翻譯到現在的語言.
我送二千架戰鬥機給你.
看你都不夠那麼多人駕駛.
行不行.
他的意思是這樣.
給你二千匹馬你都不夠人騎.
你還敢背叛我.
怎能打敗我神僕中最小的軍長.
竟然走去倚靠埃及的戰車馬.
你不就死定了.
第八第九節.
我們揭曉後面的說法.
他講了一段.
我們跳到第十一節.
以利亞經.
舍伯拿約亞.
剛才我讀過幾個官長.

$^{121}$猶大那邊的.
對拉伯沙基說.
求你用亞蘭言語和僕人說話.
因為我們懂得.
不要用猶大言語.
和我們說話.
這裡是十一節.
不如用達到我城中百姓的意中.
如果翻譯到現在的我們香港.
不如你講英文我們懂得聽.
不要講廣東話行不行.
他不是很想其他民眾懂得聽.
不過香港人那裡懂得聽英文.
不過那時候的人.
國際語言是亞蘭文.
所以說不如你講英文.
我們懂得聽.
不要講廣東話.
大概是他現在的語言.
我覺得最硬的是這裡.
拉伯沙基的回應.
十二節.
你跳下去看十二節.
拉伯沙基說.
我主差遣我來.
豈是單對你和你的主說這些話.
我派來只是和官員說話.
不也是對這些坐在城上.
要與你們一同吃自己端.
和自己嘟的人說麼.
你看到聖經是這樣寫的.
我不是刻意的.
不過那個字是不讀出來的.
即是說.
你們聽吃嘟和喝嘟吧.
你們想叛變我.
我是要告訴你.
你個個都要什麼.
所以我用廣東話告訴你.
一起去喝嘟的人聽.

$^{161}$即是大字.
你敢不敢受到這個官長.
是幾個....
幾個什麼.
啊?.
串系.
就是他這麼串的.
但是他又串不串得起呢?.
又串得起.
他大軍壓陣.
他又打贏了很多的軍王.
所以他真的挺串的.
於是第十三節.
於是拉伯沙基暫時用猶大言語.
大聲喊著說.
你們當聽亞述大王的話.
王如此說.
你們不要被欺西家欺控.
因他不能拯救你們.
也不要聽欺西家使你們倚靠.
耶和華說.
耶和華必拯救我們.
這城必不交在亞述王的手中.
即是說.
你不要聽欺西家.
救不了你的皇帝.
還有.
不要聽欺西家擄點你說靠耶和華.
行不行?.
都是不行的.
他後面有些證據說的.
他說後面就有提了.
他說其他國家的偶像都救不了他的.
我就不讀了.
這裡在十八節打後.
說其他國家的偶像.
都不能夠救他們脫離亞述王的手.
但我特別想讀一段.
第十六節.
因為亞述王如此說.

$^{201}$你們要與我和好.
出來投降我.
各人就可以吃自己葡萄樹和無花果樹的果子.
和自己井裡的水.
這是勸他投降的宣言.
讓我來領你們到一個地方與你們本地一樣.
就是有五穀和新酒之地.
有糧食和葡萄園之地.
即是說.
如果你們投降.
可以怎樣?.
可以吃自己種的水果.
自己的糧食.
而且我又會帶你們去一個好地方.
跟這裡差不多.
收到嗎?.
我又想問一個問題.
既然亞述大軍如此強勢.
為何要勸你投降?.
打不就行了.
按道理一打就贏.
剛才都說了.
就算給你所有武器.
你都不夠人識用.
為何還要浪費唇舌?.
為何要浪費唇舌?.
只是猜測.
為何要浪費唇舌?.
打不就行了.
為何要說這麼多?.
以前無線電視保鏢會的劇集.
講了一大輪都不打.
即是怎樣?.
為何要這樣做?.
是震動你的軍心.
你真是有計謀.
完全正確.
其實又說回頭.
現實的社會是怎樣呢?.
現實的情況是.

$^{241}$最好就不用打你投降.
因為打你.
就算打得贏的話.
其實他都要付代價.
他打人.
他不用死傷一點嗎?.
他都要死傷一點.
最近的事情.
我發生都知道.
以色列真的炸到你伊朗命.
但其實他打你也要很長時間.
你都很煩.
最好就你投降.
你投降他不用死傷.
行不行?.
不用這麼煩.
但你不投降.
他打贏你.
間中咬你兩口.
你都會很傷.
不過擾亂軍心是真的.
希望你們不要抵抗.
最好你不要抵抗.
馬上投降.
不過我想在這裡.
特別想你留意一件事.
拉伯沙基這個官長.
他在描述.
說他們的大王會領你們.
去一個美麗的新世界.
這個美麗的新世界是怎樣呢?.
我再讀一次.
美麗的新世界是怎樣呢?.
這裡就是讓我領你們去一個地方.
與你們本地一樣.
就是有五穀和新酒有糧食和葡萄園之地.
第17節.
那些惡者很多時候都在描述一個美麗的新世界.
你試著回答這句話.
其實是不是真的呢?.

$^{281}$那裡真的這麼好嗎?.
還有.
為什麼我要去那裡呢?.
可以嗎?.
你說那裡會跟現在這裡差不多.
有這麼多好東西.
那我在這裡不就很好嗎?.
為什麼要去那裡呢?.
可以嗎?.
感不感受到?.
描述一個美麗的新世界.
那個美麗的新世界.
究竟存不存在呢?.
是不是真的你描述這樣呢?.
我想給你一個問題.
其實這個社會現在.
很多人都在描述.
特別是在不信主的世界裡面.
是在描述一個美麗的新世界給我們的.
譬如.
DEI的社會裡面.
自由 博愛 平等 多元.
挺好的.
我們聽的各樣東西都挺好的.
但是那個.
它引領我們去那個地方.
是不是真的一個自由 博愛 平等的社會呢?.
又或者我們的聲音是被聆聽的呢?.
是不是真的呢?.
他描述的.
又或者在政治上.
其實選舉多數那些人都說.
如果選了我就會這樣這樣和這樣.
那個是不是一個真的.
真的能夠達到的美麗新世界呢?.
我想將這件事展示給你看.
事實上.
亞述大軍所經過的地方.
並不是如此.
就會將那些人擄掠到其他地方遷移.

$^{321}$也將其他人擄掠到去薩馬利亞.
移去薩馬利亞.
所以混雜了薩馬利亞人.
Anyway.
這個美麗新世界究竟存不存在.
又或者這麼說.
平時都是.
我們現在聽到很多世界跟我們說的東西.
此是而非.
我們求天賦給我們智慧判斷.
究竟是還是不是呢?.
譬如我的女兒有時候跟我說.
Daddy.
聽說現在讀大學沒什麼用的.
認識人比識字還要好.
我不知道有沒有人聽她說.
我說你聽誰說.
她說Elon Musk和史蒂夫·喬夫.
都是不用讀大學畢業就可以了.
但是多少億人裡面才有一個史蒂夫·喬夫.
和一個Elon Musk呢?.
感不感受到.
很多美麗新世界被我們描述.
究竟是不是真的呢?.
求天賦給我們一個判斷的能力去判斷.
剛才拉薩巴.
這個長官所說的.
拉伯沙基所說的.
究竟是不是一個美麗的地方.
第21節.
百姓就靜默不言.
並不回答一句.
因為王曾吩咐說不要回答他.
在這裡我想說.
我自己感受是挺有智慧的.
有些事情真的不能回答也不好回答.
怎麼說呢?.
沒完沒了.
有些他有誠意地問的問題就要回答.
再說一次.

$^{361}$有誠意問你的問題就要回答.
但是沒有誠意問.
你覺不覺得這個拉伯沙基是有誠意地問問題呢?.
沒有誠意純粹來挑戰你.
其實不需要回答或者不能回答.
感不感受到.
我舉個例子.
我們有時候都會教一些叫做.
初信班的.
教會的初信班.
有些人會亂問一些問題.
不關我的事.
為什麼要打仗.
為什麼特朗普又怎樣怎樣.
然後有誰又怎樣怎樣.
特朗普不關我的事.
我不熟悉他.
我回答這些問題沒完沒了.
為什麼誰打完之後又繼續打.
其實不關我的事.
我又不是美國人.
但是如果我每條都要回答.
回答到沒完沒了.
再說一件事.
我真實經歷的一件事.
當年那些初信班.
基督教和佛教又怎樣.
我嘗試回答.
回答完之後.
他又說.
那和道教又怎樣比.
又回答一輪.
回答完之後.
那時候剛剛讀完神學.
聰明懂得.
現在忘記了.
回答完之後.
那和伊斯蘭教又怎樣比.
我現在不是宗教比較學.
那你信哪一種教.

$^{401}$我想問.
那你信哪一種教.
我解那一種可不可以.
我一種都不信.
那你又問.
可不可以.
有時候如果那個不是有誠意的話.
禮貌回應就可以了.
這裡說到他不用回答.
也沒有回答.
因為王吩咐你不要回答.
所以我們真的回到第37章.
現在真的說37章.
「希西加王聽見神僕告訴他的說話」.
「就撕裂衣服,披上麻布」.
「遵耶和華的殿使」.
「加載以利亞競和書記舍伯拿」.
「並祭司中的長老」.
「都披上麻布去見阿摩斯的兒子」.
「先知以賽亞」.
這裡其實希西加王做了兩件事.
第一件事是什麼?.
披上麻布進耶和華的殿.
是什麼?.
就是禱告.
去耶和華的殿禱告.
第一.
他自己要做的事就是披上麻布.
在一種痛悔或慘慘的時候.
當時披麻,蒙悔.
都是慘慘的時候會做的事.
現在見到有危難.
現在去耶和華的殿.
就是禱告.
第一.
第二.
就派了馬仔去找誰?.
找以賽亞先知.
找以賽亞先知又做什麼?.
我們看下去第四節尾.

$^{441}$第四節尾的結論就是.
「故此求你為剩餘的民讓聲禱告」.
原來找以賽亞是找他去禱告.
兩件事做了.
一個是自己禱告.
一個是找人禱告.
都是禱告.
一個是禱告.
或者叫做代求.
一個是祈求.
一個叫代求.
問題又來了.
他自己祈禱就可以了.
為什麼要找人跟他祈禱呢?.
自己祈禱不是好好的.
為什麼要找以賽亞先知?.
難道以賽亞靈一點?.
不過其實現實生活有時都有點相似.
你見到杜生.
「杜生幫我祈禱,你靈一點」.
但見到牧師.
「牧師你再靈多一點」.
是有的.
有大其大.
有這個想法.
那究竟是否正確呢?.
正不正確?.
那我又要看聖經了.
Book.
「愛主」.
是的.
我們看回聖經如何教導我們.
首先我們看回.
在《馬太福音》.
如果我沒有跟你揭開.
你自己揭開看你能否揭開.
十八章.
《馬太福音》第十八章.
十九二十節.
十八章,你用手機就不要懶惰.

$^{481}$幾快而已.
慢慢找回來是有韻味的.
《馬太福音》.
有些很久沒有找聖經的人.
揭聖經的人.
覺得已經有些難度.
新玉在哪裡呢?.
《馬太福音》在哪裡?.
《馬太》在新玉.
淺淺的那些.
《馬太福音》.
十八章.
考大家一起揭開揭開.
第十八章十九節.
找到沒有?.
我要揭紙的找到.
揭紙版聖經的找到.
十九節.
「我又告訴你們.
若是你們中間有兩個人.
在地上同心合意地求什麼事.
我在天上的父必為他們成全」.
這不就是在說代禱的意義.
我們就是那兩個人同心合意.
求什麼事天父會成全.
所以這是代禱的原因.
還有第二十節.
「因為無論在哪裡.
有兩三個人奉我的名聚會.
哪裡就有我在他們中間.
主同我們同在」.
所以這也是代禱的意義.
行不行?.
第一.
第二.
第二就是《雅各書》.
又考一下.
紙本聖經.
我特意多揭幾頁.
後面.

$^{521}$第五章.
《雅各書》第五章.
第十六節.
《雅各書》第五章第十六節.
找到.
慢慢.
五章十六節.
五章五章十六.
《雅各書》.
好.
「所以你們要彼此認罪互相代求.
使你們可以得意志.
二人禱告所發的力量是大有功效的」.
我們知道為什麼要代禱.
因為二人的禱告大有功效.
而且我們其實有一個condition.
就是要彼此認罪互相代求.
在當中不是無條件的.
有功效.
舉手.
(我都問他 如果不是我祈禱就代我自己不祈禱).
這樣就沒有祈禱.
這是新鮮問題.
找人代禱就行了 自己不祈禱.
太懶惰了 自己也要祈禱一下.
為什麼不祈禱.
我們太古城堂不應該這樣教導.
自己當然要祈禱.
自己也要祈禱.
自己也要祈禱 記住.
不要太懶惰 自己也要祈禱.
我們又看回.
看完為什麼要代求.
我們看回以賽亞書.
重點我們看以賽亞書.
他找人代求.
但還有什麼意義.
剛才說了一些意義.
我還想說37章.
我們要看看這個代求的經歷裡有什麼學習到的地方.

$^{561}$因為第三節.
請你留意第三節.
因為他派馬仔去和以賽亞仙子說的時候.
其實是有一段說話說給他聽.
說給以賽亞仙子聽.
希西加如此說第三節.
今日是急難擇伐靈育的日子.
就像婦人將要生產嬰孩.
卻沒有力量生產.
他不是將36章原原本本說給以賽亞仙子聽.
他是總結了然後說給她聽.
什麼意思呢.
原來我剛才讀的經文是有意義的.
代禱的意義.
但我想今天補充多加一樣東西.
原來我們在叫人代求的過程裡.
聖靈是會有功能在的.
即是有功用法出的.
就是什麼呢.
首先你都要總結一下.
你究竟發生了什麼事情.
你要告訴別人叫人幫你代求.
你要想過的.
你也沒什麼可能說了36章給他聽.
你是總結了再說給他聽的.
他總結的是什麼呢.
急難,擇伐和靈育的日子.
總結了36章整章來說.
他總結了.
那即是怎樣呢牧師.
即是說對於他36章整個經歷.
他回答真一點.
其實總結是這幾樣東西.
好像也很沒力量.
我現在好像生小孩生不出.
沒力量生死梗的情況.
這個全部都不是36章說的.
是他自己感受的.
我相信這是他內心最真誠的聲音.
我現在真的很沒力量.

$^{601}$這是他真誠的.
在這個過程中.
我回教會回了這麼多年.
大家都是這樣.
有時聽完聖靈都會感受.
我們叫做同感一靈.
去感受對方的感受.
你會不會呢.
但有時你不能啞門.
不知道你有沒有試過.
譬如有個弟兄跟我說.
我如今又是急難,責罰,凌辱的日子.
有什麼事呢.
因為我明天要考試.
聖靈會感動我告訴他.
其實明天考試不是急難.
一年前就知道了.
所以照計應該不是急難.
不過可能責罰和凌辱是對的.
但絕對不是急難的日子.
可以嗎.
說到責罰.
其實這是希西加王.
對歷那麼多先知的勸告.
一個領會.
因為先知一直都在說.
為什麼外邦人這麼厲害.
因為你們猶太人不行公義.
有詭詐的天秤.
對孤兒寡婦很壞.
所以上帝透過外邦人管教你們.
你們因為在商鞋買了窮人.
很多先知告訴你算盤.
你體會到這是責罰的日子.
是天父審判我們.
或者是在管教我們.
還是一個凌辱的日子.
究竟這個凌辱的日子.
我剛才讀的經文.
你見到那個官長拉伯沙基.

$^{641}$是多凌辱你們.
在當中我們又有什麼體會呢.
我今天還想和你分享.
在過往上兩個星期.
我也有類似的經歷.
弟兄姊妹找我代討和代求的時候.
我也甚至會感動為你代求的人.
他和你一起去分辨.
天父在我們當中的心意.
讓我們也認識一下.
我們現在面對的是什麼問題.
講兩個故事.
第一個故事.
有一位姐妹很久沒有回教會.
她找我代求.
什麼事情呢.
她也類近說是一個急難.
是一個責罰凌辱的日子.
當然她沒有讀聖經.
沒有這麼文奏奏.
我聽完她的故事也類近.
為什麼呢.
因為她的弟弟借了一大筆錢.
要還債.
現在是急不急難呢.
是挺急難的.
我有沒有和她代求呢.
雖然她很久沒有回教會.
我也和她代求為她祈禱.
祈禱之後我就問她.
你有什麼感受和學習到.
天父有沒有什麼話要跟她說.
她就跟我說.
牧師 我感受到什麼呢.
人情冷暖 世態炎涼.
我問她為什麼會有這個感受呢.
她說我問了很多.
挺有錢的朋友.
向他們借錢還大耳窿.
他們也沒有給錢我幫他們還.

$^{681}$明明那些人很有錢.
幾層樓收租.
也不借錢給我還大耳窿.
我聽完之後.
我憑著.
我相信都是聖靈提醒我.
要告訴你.
你那些朋友其實沒什麼問題.
他不是壞人.
如果你的朋友.
你沒有飯吃.
不借錢給你吃飯.
那些人真的很壞.
世態炎涼.
但不借錢給你還大耳窿.
他有什麼問題呢.
其實應該報警.
那些大耳窿是犯法的.
當然詳細我不能再詳細說.
這是他自己的抉擇.
會不會神都會透過身邊的弟兄姊妹為你代購.
他都感動他.
和你講說話.
神透過人講說話.
在這裡我還想多講一件事.
有些人都和我說.
牧師我們很難明白上帝的說話.
你覺得難明白就不難明白.
一時時吧.
一點點吧.
我在這裡多講一件事.
有時連人的說話都不太聽得明白.
我覺得他都很難聽得明白神的說話.
人和你講說話都聽不到.
不知道你有沒有這個經驗.
彷彿和牆壁講說話.
講完之後他好像沒有聽你.
但他好像在聽你.
連你和他講說話都聽不到.
很難聽到神和他講說話.

$^{721}$我們都練習一下才聽人講說話.
然後再聽天父和我們講說話.
神透過弟兄姊妹.
去為我們說說話.
是神透過弟兄姊妹向你說話.
我還想分享多一個故事.
講到靈獄的時候.
是最近發生的.
有位姊妹她自己很.
整件事是很複雜的.
複雜到她由A公司轉去B公司.
總之跳槽.
跳槽之後.
但去了B公司原來是服的.
因為那間公司.
人工加了一點.
但是很服.
總之很多檻讓你踩.
到最後做了不夠兩年.
快要倒閉了.
要倒閉了.
她在被人家倒閉之前.
當然是很憂心很憂慮.
剛剛她原本舊公司有開幕.
就回到舊公司.
剛剛好無縫交接.
靚仔到不得了.
不過回到去的時候.
就跌了一個Grade.
就是開幕跌了一個Grade.
她本來是在管理團隊.
現在就跌了一個Grade.
總之原本有房坐.
現在坐大房的那一隻.
她說最深感受.
很硬的一件事.
就是很硬是什麼呢.
她說因為大會開完.
譬如五點多開完.
那管理團隊留下我們再聊天.

$^{761}$誰你不用了.
你感受到那個感覺.
我本來是應該要留下的.
不過現在回來是不用留下的.
那個感覺是很羞辱的感覺.
她就跟我說牧師你幫我代肚.
我心理很不平衡.
我出了去之後.
那天我為她代完肚之後.
我相信是聲明提醒我.
你薪水有沒有低到.
不知道為什麼薪水沒有低到.
因為跳了出去之後.
回來又給回原本那個我.
也算是這樣.
還想怎樣.
還有還沒說完.
那一晚你回家做了什麼.
回家之後跟女兒對功課.
那挺好的.
回去跟女兒對功課.
那你覺得她們留到要聊到多晚.
你猜一下.
我問她們聊到多晚.
我不知道她們那晚聊到多晚.
第二天早上.
她們個個都沒有戴隱形眼鏡.
個個都戴上眼鏡.
應該感恩都來不及.
神在你合適的時間.
擺你合適的崗位.
個個都戴上眼鏡.
真是好東西.
求神都幫助我們合適地判斷.
我們現在所面對的困難和處境.
我們看下去結局.
結局是怎樣呢.
多看兩節的結局.
由大家所剩餘的.
這是以賽亞先知跟你說的預告.

$^{801}$在後年由大家所剩餘的.
要往下扎根向上結果.
即是說接下來的日子.
原來往下扎根向上結果在這裡說.
即是將來你們會生機勃勃.
接下來神會萬軍之.
繼續讀下去32節.
必有剩餘的民從耶路撒冷而出.
必有逃脫的人從錫安山而來.
萬軍之耶華熱心必成就者是.
即是接下來神會保守你們.
可以吃回自己所結的果子.
在當中亦往下扎根向上結果.
另外我們看下去第36節.
這就是這段經文的結局.
36節.
耶華的使者出去在亞述營中殺了十八萬五千人.
清早有人起來一看.
到時死事了.
接著亞述王西勒基納拔營回去.
住在尼尼美.
一日他的神尼斯洛廟裡叩拜.
他兒子亞德米勒和莎莉息用刀殺了他.
逃到阿拉納地.
他兒子以撒哈頓接續他作王.
即是最後西勒基納王死了.
他的大軍潰敗.
不知為何是耶華的使者.
神不神奇?.
真的頗神奇.
牧師還想講一個故事給大家聽.
神奇的事是會發生的.
這件事我想了很久是否可以講.
應該可以講,已經過了很多年.
這件事是這樣的.
我在帝興當時也是四處求人禱告.
他是一個engineer.
在做一個工程.
即是建樓.
樓下打裝地盤.

$^{841}$地下有四五層樓.
有停車場.
這邊在走地鐵.
上方又在建天橋.
另一邊地盤又在建樓.
這邊建樓,另一邊地鐵在走.
他負責的那層樓在建.
他是做consultant的角色.
建著建著就發現.
他們那邊每個星期向地鐵移了一厘米.
每個星期移一厘米.
但車是在走的.
這裡移了一厘米.
他做consultant就要想辦法.
怎樣做到那裡停不下來.
因為你在做工程就會令到這裡的泥移了.
這件事真是.
最初他跟我們說要代討的時候.
我就說停下來做堅固就可以了.
已經有辦法做了.
他說每個地盤都有千個工人在開工.
叫不停.
而且那個地盤不是我們的.
有些地盤是別人的.
叫不停.
就算叫停都很難.
因為地盤有千個工人在開工.
怎樣叫停呢.
每天都要花很多錢叫停一天.
每晚都睡不著.
叫了不同的代討.
奇蹟會不會發生呢.
有時是會的.
最後怎樣救呢.
政府改職.
說叫他們全部停下來.
因為上面那條橋說要改設計.
叫全部停下來.
大家沒有聲音.
感恩.

$^{881}$再想想怎樣救那個命.
Anyway.
這些奇蹟有沒有發生呢.
是有發生的.
西拿基納王十八萬兵死了.
是會發生的.
但又是不是經常都發生呢.
但即使不是經常都發生.
但我們都相信我們仍然可以.
可以向下紮根.
往上結果.
繼續見證上帝在我們生命上奇妙的作為.
親愛的天父我求你.
保守我們眾體的姐妹.
我們為自己祈禱.
亦為我們身邊的人代求.
讓我們聽到你聖靈給我們的聲音.
叫我們向世界去見證你的恩典.
總在我們身上.
縱然我們軟弱.
你仍然可以彰顯莫大能力.
叫西人看見.
知道你是耶和華.
祈禱教託奉耶穌名求.
阿們.
\newpage



\section{創世記 41:32}
\label{sec:031854}
\textbf{聖經人聖經事:神命定的,必速速成就}~~~~黃克勤牧師博士
\newline
\newline
連結: \href{http://www.church.com.hk/acms/content.asp?site=cdc\&op=show\&type=product\&code=031854\&layout=sermon}{\texttt{031854}} ~~~~ 語音日期: 2026-03-20
\newline
\newline
\hyperref[sec:031853]{\small{< < < PREV SERMON < < <}}
~
\hyperref[sec:index_chronic]{\small{[返順時目]}}
~
\hyperref[sec:031855]{\small{> > > NEXT SERMON > > >}}
\newline
\newline
創世記 41:32
\newline
\begin{longtable}{cl}
\hline
\hline
章節 & 經文 (和合本修訂版)\\
\hline
41:32 & \begin{tabularx}{0.7\textwidth}{X} 至於法老兩次做夢，是因為神已經確定這事，神必速速成就。 \end{tabularx} \\ \\
[1ex]
\hline
\hline
\end{longtable}
$^{1}$各位姐妹,好開心一齊又讀聖經.
一齊看看在聖經的人和事的背後.
上帝是怎樣工作的.
《生命說書》的網絡是很歡喜跟你一齊讀聖經.
今次仍然讀約瑟的故事.
神命定的,必速速成就.
這是一句信心的說話.
希望大家都可以在當中學到東西.
如果神命定了,事情就會成就.
事情是這樣的.
昨天說到約瑟.
當時法老王發了一個夢.
把約瑟解不到.
那些人馬上叫人把約瑟帶出來.
約瑟就解到了.
於是約瑟就跟法老說.
你講一個夢來聽聽.
不過你要知道解夢是在乎上帝.
不在乎我.
不過神應該會跟你講平安的話.
是不是這樣的.
於是法老就跟他講.
講完之後約瑟就對法老講.
法老的夢乃是一個.
神已經將所有做的事指示你了.
留意這些字.
七隻好母牛是七年.
七個好水子也是七年.
借夢乃是一個.
因為法老明明是發了兩個夢.
第一個夢驚醒了.
七隻肥牛在吃草在河邊.
七隻瘦牛很醜的又瘦的.
吃了那七隻肥牛.
嚇到醒了.
什麼事呢.
可能半夜眼訓又睡了.
又發第二個夢.
七個好肥的墨水.
又有七個好瘦又被東風.

$^{41}$在曠野沙漠吹來的東風.
算是沙漠.
因為如果在埃及東面吹來的.
就是西乃沙漠吹來的熱風.
將那些水子變得很乾.
那七個乾的水子吃了.
那七個肥的水子.
什麼事呢.
很害怕沒人能解釋.
其實法老發的夢.
約瑟說你的夢乃是一個.
無論是七隻好母牛.
七個好水子.
都是指七年.
你的夢其實是一個.
這件事應該令到法老聽完之後.
原來是這樣.
再說再說.
真是好arm 聽.
然後隨後上來的.
七隻又乾又醜陋的母牛.
都是七年.
那七個空的墨水.
吹焦了的都是七年.
不過都是代表方年.
這個就是我對法老所說的.
神已經將所要做的事.
顯明給法老.
你看到這件事很重要.
今天要說的第一點.
就是一個夢.
是上帝指示了你要做的事.
神已將所要做的事.
指示法老.
為什麼沒有法老王呢.
因為法老這個字本身.
是埃及的王的稱號.
不過我們很多時候.
都會加上法老王.
正如波斯阿哈徐老王.

$^{81}$一個夢來的.
神已經將他快要做的事.
預先指示你.
這是一個恩典連文.
然後下去看.
就是這樣.
神已將所要做的事.
顯明給法老.
說了第二次.
所要做的事.
顯明給你了.
是透過.
不是兩個夢.
是一個夢.
這個法老當堂.
厲害啊.
這件事真是厲害啊.
很重要啊.
這樣都被你說到.
然後.
弱勢繼續說.
埃及遍地必來七個大風年.
然後隨後有七個方年.
甚至埃及地.
忘了先前曾經豐收過.
全地會被饑荒所滅.
因為以後來的饑荒.
就是七年.
豐年之後.
七年饑荒.
多到都忘記了.
先前的豐收了.
很多時候人都是這樣的.
好的時候.
就不是很覺悟.
哎呀很好啊.
但是一二都不好.
方年就.
哎呀我真是很慘啊.
都還沒有好過.

$^{121}$好了.
至於法老兩回做夢呢.
留意這裡.
第二點要說的.
是神命定這時.
而且必速速成就.
神說了.
就會速速成就.
神命定的.
所以法老當選.
選一個有聰明有智慧的人.
派他治理埃及地.
特別在這裡提醒大家.
神是直接跟法老說.
說得這麼清楚這麼直接.
是的.
因為他不是很認識這個神.
他們埃及地是有超級這麼多神.
這個時候.
真神要向他顯示自己.
是用一些很特別的方法.
令到他.
嘩.
厲害.
第一沒有人解到.
表示.
若失信的神.
是比其他那些術士.
那些博士解的更厲害.
所以其實顯出.
若失信的神.
是厲害過他們埃及的神.
因此.
雖然他們不是很熟.
不是很認識是誰.
耶和華.
他們不知道的.
不過經過這個夢.
就開始認識.
哇.

$^{161}$若識的神耶和華.
真是厲害.
好了.
而且他說.
神命定這時.
必速速成就.
很希望大家都記得.
如果神有一件事.
是命定在你的身上.
那件事是會成就的.
如果那件事.
是發生在你的身上呢.
反過來看.
你就會知道.
你一定要想的就是.
會不會這件事.
是神安排了.
在我身上發生呢.
神這裡用的字叫命定.
即是說.
祂說會發生的.
祂要做的.
有時候我們有一件事.
發生在我們身上.
為什麼會這樣.
神為什麼不保護我.
神為什麼不安排.
好路給我走.
有沒有想過.
神是命定了.
有一件事.
在你和在我身上發生.
我們做牧師.
這個人.
有沒有一些事.
是不好的.
發生在自己身上呢.
當然有了.
有人陷害我.
我經常都記得.

$^{201}$很慘.
我記得的.
怎樣.
但是我後來一路開始好.
讀完藥室之後更加好.
神命定了一些事.
祂想害我.
但是在神來講.
是祂安排了的.
是好的.
有時候去到.
我以前很多頭髮.
不是這樣的髮型.
去到現在變成這樣.
幾十歲之後變成這樣.
怎樣.
這樣.
找什麼茶.
祂會長回頭髮.
我現在不想了.
現在幾好.
神讓我這樣就這樣.
前兩年病了.
瘦了二十幾三十磅.
那些衣服全部不對.
最後衣服闊了.
好像偷人一樣.
怎樣.
這樣.
神命定了.
速速成就就成就.
這件事我的經歷.
真的速速成就.
在我不覺的.
我入院是沒有做手術.
檢查而已.
不過要開始吃薄血藥.
這樣就睡在那裡.
試藥.
弄了十幾天之後.

$^{241}$一出來.
整個人累了.
整個人瘦了.
我想是.
我到六十五歲.
神給我.
開一個新的一頁的時候.
以前我就是.
求問神.
六十五歲會不會接我回去.
我想回去.
神說你想不要這麼快.
不過你這麼胖不行了.
給你減肥.
我一百九十幾.
沒有想過.
可以跌到一百七十幾.
真的少了三十磅.
現在就這樣了.
沒有氣力.
沒什麼力.
肌肉少了.
胃口小了.
為什麼會這樣.
這樣也不是很好.
神命定了.
這樣就這樣了.
我一路祈求上帝.
讓人來接手我的工作.
但是沒有接手的人出現.
其實可能就是我難了.
有人出現.
沒有出現.
如果我回去.
回去.
神命定了的事.
是速速成就.
所以我說不要緊了.
交託給主.
只要是主命定的就行了.

$^{281}$記得.
神命定了的事.
有什麼病.
生什麼病.
不要緊.
交託給主.
而且是會成就的.
既然神命定了一件事.
就是有七年豐年.
有七年荒年.
既然是這樣.
約瑟就跟法老說.
你要選一個有聰明.
有智慧的人.
派他治理埃及地.
約瑟厲害的是.
他獻策.
他沒有說我自己做.
其實皇上你應該找一個人.
開始去管理好你的糧食.
豐年的時候.
收好.
然後豐年的時候可以吃.
不然就會滅了整個國家.
法老當者讓行.
第三點.
見到這件事就要做些準備.
派官員管理這地.
七個豐年的時候.
徵收埃及地的五分之一.
其實夠不夠呢.
叫他們把將來豐年一切的糧食聚斂起來.
積蓄五穀收存.
在各城裡做食物.
歸於法老的手下.
你看到.
叫他們.
法老皇啊.
皇上你應該這樣行.
當者讓行.

$^{321}$這三點是建議和形勢.
你要做的事.
這幾點我想和大家看看.
第一.
什麼呢.
神指示他透過一個夢.
透過一個連環的夢.
兩個夢其實是一個.
將一些事情.
神所要顯明的事.
說給法老聽.
但他不明白.
於是透過約瑟向他解釋.
然後.
先是七年豐年.
七年荒年.
神命定這事會速速成就.
請你記得.
神這樣指示你.
而且這件事會成就.
快點準備.
而且是速速成就.
一定要小心.
既然知道有這樣的事發生.
那就坐著等吧.
看看怎樣.
很多人有拖延心理.
很大件事.
怎麼辦.
要做些事.
但他從來沒有做.
大家要小心不要拖延.
約瑟給了一個建議.
給法老.
這樣做吧.
這裡我們也學到.
第三點.
當你見到一個形勢.
你要做一些事.
做一些事.

$^{361}$我也想了很久.
可能要回去建築.
我這裡也找人要接手.
幾個機構工作.
那怎麼辦.
我還沒有.
不過我也在想.
可以怎樣做呢.
當怎樣行呢.
如果是這樣.
你要派官員去將那七年豐年的事.
多些徵收.
約瑟厲害.
神指示他.
就是.
在那豐年的時候.
豐年多少年.
七年.
收五分之一.
咦.
收五分之一怎會夠呢.
奇奇怪怪.
收五分之一.
因為你每一年要收稅.
然後呢.
那年的稅收了之後.
你會用了.
那怎麼辦呢.
現在就是.
現在我要儲起的.
就是大家糧食生產出來的.
五分之一.
去留做儲備.
如果五分之一.
七個五分之一呢.
就幾好了.
七個五分之一即是多少呢.
即是有一點四了.
即是說.
是有這七年豐年的一點四倍.

$^{401}$的收成放在這裡.
嘩 幾聰明.
如果一就不是好.
用什麼一點四.
一就是剛剛夠那七年豐年的時候吃.
但如果我多零點四.
即是近一半呢.
那我就可以應付其他突發的需要.
或者可以幫助其他國家.
因為是互來看到呢.
周圍的人都去問埃及.
擲糧.
去那裡去問他們買糧食.
那他多了一點出來.
多了零點四就可以給人了.
那為什麼七個豐年收到五分之一.
你計到就一點四了.
即是說 夠有餘.
還多差不多一半.
所以他們預備好了.
收存在城裡做食物.
收藏在城裡面的倉庫.
城裡面為什麼這樣.
容易管理一點.
這件事就是這樣.
好了 總結.
神指示了的事情.
神要做的事情.
是會速速成就.
是人不可抗力的.
是說我不肯 不行 不行.
神說了就會成就.
另外請你記得.
當你看到 咦 神想這樣做.
那你要怎樣.
你馬上要做的就是.
做一些相應的行動.
不要拖延.
嗯 我看得出是這樣.
於是我要做了.

$^{441}$那譬如呢.
我看得出 咦 神都有工作.
而在宣教上面.
中國是在走一帶一路.
那我也要開始看看一帶一路.
有什麼地方.
是我們做宣教可以做.
追著去做了.
另外 啊 神在哪個地方.
哇 那個科研工作發展得很好.
那你做宣教的人就不要拖延了.
那怎麼辦.
我的主在那裡做事.
那我們這幫人怎麼辦.
當然追著去做了.
啊 那邊發展得多好.
好東西好東西這樣.
這樣不夠了.
你要追著去做.
咦 有時候神指示你.
那邊很困難.
你去那邊吧.
咦 去那邊很好啊.
那當神指示顯明的時候.
你就當這樣走了.
跟著去.
貼著去.
做相應的準備與行動.
那這個是重要的.
在這裡學到.
神的手的工作.
還有祂想我們怎麼樣.
我們同心禱告.
多謝我們的主.
其實我們如果是很多時候.
不是太忙.
靜心去想一想.
去聆聽你的時候.
是發覺其實你是不停地對我們說話的.
在我們生命當中.

$^{481}$每一天裡面.
你不停地跟我們講一些訊息.
有時候我們只是太忙.
沒有聽到.
又或者是沒有記住.
所以比我們有更多安靜的時候.
在主面前聆聽.
聽你說什麼.
既然你向我們顯明你的心意.
你命定之後又會速速實行.
就叫我們有相應的準備.
當這樣走.
讓我們預備好.
可以與你配合.
是跟隨你.
跟隨你的腳步.
讓我們有這個敏銳的.
淑寧的耳朵.
誰聽我祈禱.
奉耶穌基督命求.
阿們.
\newpage



\section{約翰福音 6:59-71}
\label{sec:031855}
\textbf{榮耀君王:永生之道}~~~~林鍵牧師
\newline
\newline
連結: \href{http://www.church.com.hk/acms/content.asp?site=cdc\&op=show\&type=product\&code=031855\&layout=sermon}{\texttt{031855}} ~~~~ 語音日期: 2026-03-21
\newline
\newline
\hyperref[sec:031854]{\small{< < < PREV SERMON < < <}}
~
\hyperref[sec:index_chronic]{\small{[返順時目]}}
~
\hyperref[sec:031856]{\small{> > > NEXT SERMON > > >}}
\newline
\newline
約翰福音 6:59-71
\newline
\begin{longtable}{cl}
\hline
\hline
章節 & 經文 (和合本修訂版)\\
\hline
6:59 & \begin{tabularx}{0.7\textwidth}{X} 這些話是耶穌在迦百農會堂裡教導人的時候說的。 \end{tabularx} \\ \\ \relax
6:60 & \begin{tabularx}{0.7\textwidth}{X} 他的門徒中有好些人聽見了，就說：「這話很難，誰聽得進呢？」 \end{tabularx} \\ \\ \relax
6:61 & \begin{tabularx}{0.7\textwidth}{X} 耶穌心裡知道門徒為這話私下議論，就對他們說：「這話成了你們的絆腳石嗎？ \end{tabularx} \\ \\ \relax
6:62 & \begin{tabularx}{0.7\textwidth}{X} 如果你們看見人子升到他原來所在之處，會怎麼樣呢？ \end{tabularx} \\ \\ \relax
6:63 & \begin{tabularx}{0.7\textwidth}{X} 聖靈賜人生命，肉體毫無用處。我對你們所說的話就是靈，就是生命。 \end{tabularx} \\ \\ \relax
6:64 & \begin{tabularx}{0.7\textwidth}{X} 可是你們中間有些人不信。」耶穌起初就知道哪些人不信他，哪一個要出賣他。 \end{tabularx} \\ \\ \relax
6:65 & \begin{tabularx}{0.7\textwidth}{X} 於是耶穌說：「所以，我對你們說過，若不是蒙我父的恩賜，沒有人能到我這裡來。」 \end{tabularx} \\ \\ \relax
6:66 & \begin{tabularx}{0.7\textwidth}{X} 從此，他門徒中有很多退卻了，不再和他同行。 \end{tabularx} \\ \\ \relax
6:67 & \begin{tabularx}{0.7\textwidth}{X} 耶穌就對那十二使徒說：「你們也要離開嗎？」 \end{tabularx} \\ \\ \relax
6:68 & \begin{tabularx}{0.7\textwidth}{X} 西門‧彼得回答他：「主啊，你有永生之道，我們還跟從誰呢？ \end{tabularx} \\ \\ \relax
6:69 & \begin{tabularx}{0.7\textwidth}{X} 我們已經信了，又知道你是神的聖者。」 \end{tabularx} \\ \\ \relax
6:70 & \begin{tabularx}{0.7\textwidth}{X} 耶穌回答他們：「我不是揀選了你們十二個嗎？但你們中間有一個是魔鬼。」 \end{tabularx} \\ \\ \relax
6:71 & \begin{tabularx}{0.7\textwidth}{X} 耶穌這話是指著要出賣他的加略人西門的兒子猶大說的；他本是十二使徒裡的一個。 \end{tabularx} \\ \\
[1ex]
\hline
\hline
\end{longtable}
$^{1}$歡迎大家來到我們的周六崇拜.
五個星期前我們開始了約翰福音第六章.
是約翰福音裡面最長的一章聖經.
一共有七十一節.
一開始的時候有兩個非常奇妙的天國記號.
五餅二魚 餵飽五千.
然後是黑夜狂風 海上行走.
驟眼看來這兩個天國的記號很不一樣.
好像格格不入一樣.
但當我們仔細去研讀的時候.
我們發現原來五餅二魚和海上行走.
是一對龍鳳胎.
正所謂好事成雙.
五餅二魚的結尾是十二個男子.
而海上行走的起頭是十二個門徒.
海上行走是約翰福音裡面.
唯一一件行在門徒身上的神蹟記號.
是耶穌為門徒度身訂造的神蹟記號.
為要成為他們將來不久之後的.
將來的鼓勵和認許.
因為很快十二個門徒就要被耶穌.
差遣進入一個黑暗的世界.
Into darkness.
將餅碎,剩下的餅就是生命的糧.
耶穌基督的福音.
去傳遍以色列十二支派.
更加將福音帶到地極.
他們需要耶穌的同在.
才可以承受這個黑暗對他們的挑戰.
耶穌也自己認許了他們.
他天天與門徒同在.
從那時開始直到世代的終結.
於是有兩個天國記號.
天國記號之後.
約翰帶我們進入了.
耶穌在加伯隆會堂裡面的一篇講道.
「我是生命糧」.
這是耶穌一篇很長的講道.
在當日現場的聽眾.
我相信大概沒有人明白.

$^{41}$耶穌究竟在講什麼.
聽完感覺上是一塌糊塗.
非常凌亂.
但是只要他們不離開耶穌.
只要他們繼續跟隨耶穌.
在一年之後.
他們就會明白.
耶穌那篇道裡面所講的一切.
因為在那篇道裡面.
所有的一切都是指向他的十字架.
當十字架一上來的時候.
你再聽這篇道的時候.
你會發現你完全明白.
當中耶穌講的所有的東西.
上星期我們看完那篇道.
我將長達三十四節的講道.
今天想跟你很快地混雜一下.
我將它分為了九個部分.
應該是熟口熟面的.
第一部分.
不要為這生操勞.
卻要為永生操勞.
第二部分.
信上帝所差來的就是上帝的工作.
第三部分.
我負從天賜真良.
生命降下給世界.
第四,五,六.
我是生命的良.
信我的永不饑渴.
見子而信的得永生.
末日我要使他復活.
若不是父吸引人.
沒有人能到這裡來.
最後的三組.
我為世人的生命.
賜下我的肉作糧.
第八.
我的肉是真食物.
我的血是真飲料.

$^{81}$最後.
吃我肉在我裡面.
吃這糧的永遠活著.
第六章第五十九節.
這些話是耶穌在加伯隆會堂裡面.
教導人的時候所說的.
如果你留心去閱讀整個故事.
你會發現.
在加伯隆的會堂裡面.
是一共有四組不同的人.
第一組只有一個.
耶穌自己獨自一人.
第二組是耶穌的十二門徒.
第三組是十二門徒以外的其他門徒.
最後還有一組叫做猶太人.
大概包括了加伯隆猶太會堂裡面的領袖.
和其他的群眾.
於是我們在想.
耶穌語重心長的講道.
這篇當作是嘔心瀝血的講道.
會不會打動了他們的心.
使他們在耶穌裡面合二為一.
得蒙救贖.
成為上帝一個大能的軍隊.
如果你聽耶穌講道的時候.
最後都是想這樣.
你會不會奉上他們的餘生.
為天國賣命呢.
聽完之後就嘔心瀝血.
耶穌不單沒有打動他們的心.
反而被他們狠狠的拒絕.
最後只剩下耶穌和祂十二的門徒.
而當中竟然還有一隻魔鬼.
將要出賣耶穌.
真是非常之可悲亦非常之可惡.
今天我想帶你去先看這段經文.
十二節的經文.
就是講道完畢之後.
眾人對耶穌的回應.
這十二節的經文.

$^{121}$我將它分為五部分.
你就拿到大綱了.
那五部分是這樣的.
「這話很難,誰能聽進去呢?」.
第二「靈似生命,我話是靈」.
第三「蒙父因似,道我者呂」.
第四「生命之道,還根從誰?」.
最後「有隻魔鬼出賣耶穌」.
我想帶你去細心閱讀這十二節的聖經.
然後再回來這個大圖畫.
替你做一個總結.
然後還要應用.
你也知道難度很高.
看到最後「有隻魔鬼出賣耶穌」.
那怎樣應用呢?.
但是我們真的想問.
究竟我們看完《讓本音》第六章.
最後給我們的福音的意象.
是怎樣的一個意象呢?.
於是我們先看這十二節的聖經.
《讓本音》第六章第六十節.
「耶穌的門徒中有好些聽見了就說.
這話很難誰能聽盡呢?.
耶穌心裡知道門徒為這話私下的議論.
就對他們說.
這話成了你們半腳石嗎?」.
這裡所說的「門徒」應該是十二門徒以外.
積極進取,跟隨耶穌.
視耶穌為他拉比的跟隨者或追蹤者.
如果你看《上下文》的時候.
可能你會記得.
有一群人經歷了五餅二魚.
然後徹夜難眠.
早上起床坐船過海.
沉覓耶穌.
然後來到耶穌的廣道.
是耶穌的鐵粉絲.
這裡當中的「門徒」也有可能包括了.
耶穌在加拔羅本地的粉絲.
就是他沒有跟隨耶穌.

$^{161}$但他去到本地有本地的粉絲團.
這些就是十二門徒以外的跟隨耶穌的人.
這群人聽完耶穌那場廣道之後的回應是.
「這話很難誰能聽得盡呢?」.
這裡呢.
你這樣看中文.
你是看不到他背後的意思的.
他說「這話很難」.
第一:「這話」的原文是logos.
如果你翻譯出來是.
「這道很難」.
這個道是甚麼道呢?.
「道成了肉身住在我們中間」.
其實他在說的不是「這話很難」.
是「耶穌很難」.
因為耶穌就是上帝的logos.
那道.
而這裡所謂的難.
聖經的原文希臘文.
這個難不是難以理解的難.
是難以接受的難.
就是耶穌很難受.
他們不是聽不明白耶穌說話.
就是聽得明白.
不過聽得明白令他們很反胃.
所以我們再翻譯出來一次.
「這個道成肉身的耶穌.
不難理解,但難以接受,令人反感」.
很不一樣的.
因為很多人一讀「這話很難」.
因為之前他們又要誣衊他.
又要吃他的肉.
又要喝他的血.
那當然是難.
讓你想也想不通.
但這裡不是這樣的意思.
「這個道成肉身的耶穌.
不難理解,但很難以接受,令人反感.
於是他們私下議論」.
如果你上個星期在這裡.

$^{201}$你知道私下議論.
原文的翻譯應該是.
「他們發怨言」.
你看英文是grumbling.
是回到初埃及的以色列人對上帝.
或者說上帝的回應.
「於是他們私下議論.
就如之前在加伯隆裡面的猶太人」.
當耶穌在說一些他們聽不入耳的東西.
他們不斷向耶穌,向上帝發怨言.
發怨言代表什麼?.
有機會回到聖經出埃及記裡面.
以色列人發怨言的理由.
最後是因為不捨得埃及.
因為他們貪愛世界.
所以他們這裡的發怨言.
你一會兒聽下去就明白了.
你就會明白為什麼.
是代表了他們貪愛世界.
然後他覺得耶穌很難忍受.
於是耶穌知道他們發怨言.
就問了他們一句話.
「這話成了你們的絆腳石嗎?」.
我想你繼續看下一節.
看完下一節我再清晰的解釋給你聽.
究竟發生什麼事.
下一節很精彩.
下一節是第62節.
耶穌說.
「如果你們看見人子升到他原來所在之處」.
我猜天上,聖父那裡會怎樣呢?.
我以前讀這節聖經是完全讀不明白的.
第一.
耶穌說「如果你們看見人子」.
這個「人子」不是普通人.
這個「人子」是在說《但以理書》第七章裡面的人.
這個「人子」在《但以理書》裡面的意象.
是升到上帝的面前.
就是耿古常在者.
然後「得了權並榮耀國度」.

$^{241}$他的國度是永遠的,不能廢去的.
他的國度必不會敗壞.
他是萬王之王.
他是萬君之主.
或者你說.
如果我可以看見他一亡.
我看見耶穌升天.
那會怎樣呢?.
我想像如果耶穌說.
「如果我現在馬上在你面前升天」.
那你下一句是什麼呢?.
「我一天都光了」.
支持那些人對耶穌有很多的疑問.
其中一件事就是.
耶穌經常說「我是從天降下來」.
他說「不是的,我們認識你的家」.
現在如果耶穌在他們當面升天.
不會再不相信耶穌了.
很明顯他升了天回到他以前那裡.
如果你記得之前有人求問耶穌.
你可不可以再給我們一個神蹟.
再證明一下你.
如果你更加證明到你多偉大的時候.
我就可以承諾去跟隨你了.
我現在就是跟你三天而已.
因為你做了五餅二魚.
但是如果你做一個更大的時候.
我一生跟隨你.
如果耶穌在他們面前真的升到他原來所處的.
而且他人子那種承受國度的時候.
永遠的國度的時候.
是不是「一天都光了」呢?.
是不是解決了所有的問題呢?.
人們知道他真的是從天降下來.
也都因為他行了一個這麼偉大的神蹟.
於是大家虎目見白不再猶豫.
也都不會再有怨言跟隨他.
是不是這樣呢?.
原來不是.
所以我就告訴你們.

$^{281}$耶穌在這裡.
很有可能是自己在想像.
自言自語.
然後獨自嘆息.
嘆息.
耶穌說.
如果你們見到人子升上天.
升回他以前的地方的時候.
你們會怎樣呢?.
原來在《讓福音》裡面.
人子和上升.
這個遲早的配搭已經出現過了.
人子和上升.
這個配搭在《讓福音》裡面出現過.
在哪裡出現過?.
在第三章第十三節.
耶穌和尼哥德姆的對話裡面.
耶穌說過一句這樣的說話.
他說.
「除了從天降下來的人子.
沒有人升過天」.
OK.
意味就是說.
這個人子現在從天降了下來.
不過他很快又會升回上天.
他說.
「除了從天降下來的人子.
沒有人升過天」.
不過人子很快就會升天了.
第三章第十三節.
說的是人子將要上升回家.
回到他以前的地方.
第十三章第十四節.
會告訴我們人子上升回家的模樣.
OK.
我們現在想一下.
第三章第十四節.
耶穌要上升.
上升要坐飛機還是坐什麼?.
上升的模樣是什麼?.

$^{321}$你記不記得第十四節說什麼?.
第十四節是這樣的.
「摩西在曠野怎樣舉手.
人子也必要照樣被舉起來」.
耶穌要回家.
耶穌要上升回家.
耶穌要上升回到他從前.
未降下來.
起初的地方.
但是他怎樣上升?.
他上升有一個模樣的.
要坐某一種直升機.
某一種電梯.
那個直升機.
那個電梯叫做十字架.
Oh my God.
你現在開始明白.
耶穌在想的是什麼?.
如果你們見到有一天.
我上升回家的樣子.
如果有一天你們見到我.
被釘在十字架上.
你們會怎樣回應?.
耶穌的答案就是.
如果你.
現在我只是跟你說而已.
還沒做而已.
但是有一天.
你真的見到我被釘在十字架上.
你會怎樣?.
你會雙倍地令你反感.
雙倍地令你難以接受.
你更加不會跟隨我.
我們會雞飛狗走.
耶穌在自我想像.
獨自嘆息.
因為他們一眼看下去.
這群人.
全部都是貪愛世界.
追求地上的東西.

$^{361}$而現在他們已經對耶穌極度猶豫.
如果有一天.
耶穌真的掛在十字架上的時候.
耶穌心裡面嘆息.
這群人沒有一個人.
會留在我身邊.
在這裡你終於開始明白.
眾人拒絕耶穌的恩由.
表面上.
人們拒絕耶穌的恩由.
是因為.
耶穌無理強迫他們要吃祂的肉.
喝祂的血.
如果耶穌真的這麼說.
你真的這麼領受.
你都覺得反胃.
也都對於猶太人來說.
是不可能的.
因為祂干犯了摩西的律法.
但是背後真正的原因.
背後的原因.
是因為耶穌整天在說血說肉.
血肉代表了祂要獻上.
作一個犧牲.
如果耶穌來祂的使命.
是獻上血肉做犧牲的話.
那麼耶穌就沒辦法成為.
他們理想當中的榮耀的君王.
祂越是說犧牲.
你心裡面越想.
祂不是我在想的那個君王.
耶穌這種.
沒有回報的犧牲.
只是單犧牲而已.
祂沒有榮耀的.
令他們感到反感.
令他們難以接受.
他們要的是榮耀的君王.
耶穌要成為的是.
以下書第五十三章裡的受苦的僕人.

$^{401}$他們要的是此生地上飛黃騰達.
耶穌要他們跟隨祂.
去追求一個虛無縹緲的天空之城.
耶穌經常說.
將職讚財寶在天上.
大佬 我看不到.
你說有就有.
你跟隨我放棄這個世界.
跟隨我.
你就有財寶在天上.
是的.
但是拜託.
可不可以有個天堂銀行的apps.
讓我看看.
沒有.
你只是說虛無縹緲的天空之城.
你明不明白.
於是這群人覺得.
你越是這樣說話越不實際.
於是這群人聽了一天的道.
終於明白.
他們跟隨耶穌只是一時的衝動.
因誤會而結合.
現在因認識而分開.
就對了.
耶穌根本不是他們的type.
十字架都還沒出.
但他已經知道耶穌不是他們的type.
十字架不是他們的杯茶.
令我想起.
使徒保羅在哥林多前書第一章.
他的那一番說話.
因為十字架的道理第一章第十八節.
在那滅亡的人是愚蠢的.
真是愚蠢.
十字架都還沒出來.
耶穌還有一年才上十字架.
但他們已經依稀見到.
耶穌的生命方向不是他們想要的.
在我們得救的人.

$^{441}$卻是上帝的大能.
猶太人要的是神蹟.
他們不斷問耶穌拿神蹟.
希臘人要求的是智慧.
我們卻要全被釘十字架的.
基督.
這對猶太人是半腳石.
真是半腳石.
耶穌說我現在做的事.
是不是令你半腳呢.
對愛邦人是愚蠢.
對那蒙召的人.
基督總是上帝的大能.
上帝的智慧.
你開始明白了.
讓我再次帶你回到第六章第62節.
耶穌自己自言自語在問.
如果你們有一天見到我的醜樣.
我怎樣上升到原位.
怎樣被掛在十字架上.
你們這群人.
你會怎樣回應.
我掛在十字架上的模樣呢.
耶穌心裡知道.
你們一定會舉腳贊成.
火上加油高聲呼喊.
Crucify him.
這個人真是活該的.
真是活該的.
釘死他.
而竟然再過一年.
所有耶穌的想像.
都會成為事實.
至於我們問.
這群沉醉死在罪中的瘋子狂人.
你不要覺得他們很難明白.
說真的要吃血.
你不要想他們.
他們是一個追隨世界的人.
他們沉醉迷戀這個世界.

$^{481}$活在他們的罪中.
他們是瘋子他們是狂人.
我們問還有沒有得救.
而下一節竟然告訴我們.
是還有得救.
為什麼呢.
因為下一節是這樣的.
聖靈賜人生命.
他們已經瘋了.
這群人無法明白耶穌十字架的榮耀.
但是還有沒有得救.
有啊.
怎麼救.
聖靈賜人生命.
肉體毫無用處.
肉體他們看不到上帝十字架的榮耀.
我對你們所說的話.
耶穌說的.
就是靈就是聖靈就是生命.
這篇聖經很快令我想起.
《讓福音》第三章.
耶穌與以色列的老師德姆的對話.
他說從肉身所生的是肉身的.
從靈生的才是靈.
你們必須要重生.
那個答案.
一群沉醉在這個世界.
最終無法看見上帝榮耀的人.
有沒有得救.
有得救.
怎麼樣都救.
重生.
你在《讓福音》裡面離不開.
走不出重生.
而重生就是要看見十字架的榮耀.
願意背起十字架去跟隨耶穌.
這個叫做重生.
從肉身所生的是肉身的.
從靈所生的才是靈.
你們必須要重生.

$^{521}$如果你們重生的話.
那麼你們就可以有一個新的生命.
你們會脫離肉體現在看不見的東西.
你們怎麼可以進入一個重生的生命.
這個就是值一百萬的.
一千萬一億的問題.
我們如何可以重生.
在《讓福音》裡面告訴我們要重生.
需要幾樣東西.
我們需要聖靈的工作.
聖靈賜人生命.
我們需要有天父上帝的恩兆.
耶穌之前在這裡說.
除非上帝召你來.
其實你來不了我這裡.
是沒錯.
但這裡有另外一樣東西.
除非你聽聖子耶穌的說話.
你會發現.
聖靈的工作和聖父的恩兆.
其實是眼看不見手摸不著的東西.
但不要緊.
第三樣東西是非常清晰在眼前.
有一樣東西你可以做的.
你不可以去捉聖靈.
你不可以去捉天父的恩典.
但有另一樣東西是你和我可以去捉住的.
那樣東西叫做耶穌的說話.
耶穌已經告訴你了.
如果你細心聆聽.
你捉住我的說話.
你就是聖靈.
所以大家不用想.
聖靈我也沒有辦法.
我怎樣去控制聖靈.
我不知道.
不要緊.
你聽上帝的話語.
你聆聽祂.
你將祂藏在心裡.

$^{561}$你捉住上帝的話語.
就是聖靈.
上帝的恩典也在你裡面.
這令我想起《讓福音》第三章第34節.
是講上帝的說話.
上帝所差來的說上帝的話.
耶穌來.
耶穌除了上十字架以外.
祂就是喜歡說話.
喜歡說話.
喜歡說很多很多很多的東西.
而那個東西就是聖靈的彰顯.
當我們擁抱上帝的話.
擁抱耶穌基督的話.
耶穌基督的福音的時候.
通過祂的話語.
聖靈無限量的來到我們當中.
我想再講多一次.
重生是通過天父的恩兆.
重生是通過聖靈的工作.
沒有錯的.
但那些虛無縹緲.
我捉不住.
不要緊.
有第三樣東西我們捉得住.
重生和基督的話語是分不開的.
天父的恩兆.
聖靈的工作和基督的話語分不開.
有一樣東西你可以捉住.
耶穌不斷重複提醒你.
有一樣東西你可以捉住的.
就是我和你說的話語.
當你藏在心裡面.
當你不斷思考的時候.
你就捉到聖靈的工作.
天父的恩兆.
於是第六章第63節下半節.
我對你們所說的話.
就是靈就是生命.
這個應該是金句中的金句.

$^{601}$然而第六章第64節.
「可是你們當中有些人不順」.
你看到這裡耶穌可能有些哀嘆.
我們會說耶穌說可惜.
說了這麼多都沒有人信.
但我想你留意.
這裡耶穌不是在哀嘆沒有人信祂的道.
或者不夠多人信祂的道.
耶穌其實自己很知道.
將來會有千千萬萬人數不盡的人.
包括祂自己在場的.
未聽得懂的十一門徒.
他們都會真正相信這個福音.
他們會真正全人全心的去跟隨耶穌.
但耶穌這裡說的那句話.
「我知道你們當中有人不順」.
是接著下面那句話.
是約翰加插的.
約翰加插這樣說.
他說耶穌起初就知道.
從那時已經知道.
一早就知道.
那些人不順怕.
他一眼看下去就知道誰信誰不信.
哪一個要出賣他.
約翰這句話告訴我們.
耶穌從早開始.
祂是完全掌控著祂的人生.
祂完全掌控著祂的命運.
耶穌的死.
耶穌上集不是一個悲劇.
耶穌的死是祂完全掌管了這一切.
而這一切也是天父上帝的計劃之中.
第65節說到天父上帝.
在當中完全掌管.
於是耶穌說.
「我對你們說若不是蒙我父的.
因此沒有人能到我這裡來」.
說到上帝的恩兆.
不過同時讓我們不要忘記.

$^{641}$上帝的恩兆和耶穌的說話是分不開的.
我們沒有辦法捉住上帝的恩兆.
但我們可以捉住耶穌基督的話語.
這個是我們得救.
是我們得生命的關鍵.
上帝的說話已經來到我們的耳邊.
在我們的教會.
在你手上的聖經裡面不斷運轉.
你現在要怎樣可以得著這個生命.
你要不斷地去思想上帝的話語.
將它藏在心中.
我想帶你回到第六章當日的現場.
這個確實不是一個修葛的日子.
我們對比一下《約翰福音》第四章.
撒瑪利亞那時真是修葛.
撒瑪利亞那些人不要神蹟.
只要聽到他們殷勤款待耶穌.
接納耶穌所說的一切.
但是哪一天.
在《約翰福音》第六章那天.
在加伯隆那天不是修葛的日子.
只有這些猶太的群眾冷漠地拒絕耶穌.
令我們想起第一章第十一次.
他來到自己的地方自己的人卻不接納他.
去到一個什麼程度呢.
第六章六十六節.
從此他們途中有很多人退卻了.
不再和他同行.
「鹿鹿鹿」有人跟我說.
這個真是魔鬼撒旦的章節.
這個不關事我們不信邪的.
但是很明顯這一節告訴我們.
在當日魔鬼撒旦得逞了.
他的門徒這麼近這麼近那麼遠.
在耶穌的面前聽了他現場的講道.
竟然離開.
不過命不該絕.
還未絕望因為還有餘種.
第六十七節告訴我們.
有餘種有些人沒有走.

$^{681}$第六十七節.
耶穌就對那十二使徒.
現在只剩下那個十二使徒說.
你們也要離開嗎.
離開嗎.
接著第六十八節.
彼得回答說.
主啊 你有永生之道.
我們還跟隨誰呢.
你知道我讀這兩節的聖經.
我的問題就是.
究竟第六十七節和第六十八節之間.
隔了多久.
就是當耶穌講完之後.
望著那十二門徒.
你們也要離開嗎.
彼得是秒回.
你有永生之道.
我們還跟隨誰呢.
還是等了兩秒.
五秒 十秒.
說話.
誰說的 我說的.
好 三秒之後彼得說.
主啊 你有永生之道.
我們 你留意他說是我們.
他不是說我還跟隨誰.
他是代表著其他十一個門徒.
一起說我們還跟隨誰呢.
接著他可能望著其他十一個門徒.
喂 說啊.
他說是啊 是不是.
彼得說完之後.
你有永生之道.
我們還跟隨誰呢.
是不是.
動了 是不是.
是 是 是 是.
老師你不用氣餒.
今天我們失去了很多客戶.

$^{721}$不過不要緊.
在未來的日子我們會共同努力.
很快這些客戶就會回來.
真是整場大龍鳳.
但是彼得還沒說完.
彼得下一句.
才更加堅定更加厲害.
第六章第69節.
彼得繼續說話.
他除了說我們不走之外.
我們還跟隨誰呢.
他說我們已經信了.
我們又已經知道.
他在用完成式.
我們已經知道.
看了我已經知道.
你就是上帝的聖者.
這個上帝的聖者.
神學家考量了很久.
究竟他在說什麼.
有可能是以賽亞書裡面.
經常說他是.
又說上帝是以色列的聖者.
也有說在馬國音第一章裡面.
耶穌趕鬼的時候.
那隻鬼叫耶穌是上帝的聖者.
無論如何.
我們相信上帝的聖者.
可能代表了你就是彌賽亞.
你記得彼得也是這樣說的.
耶穌說我是誰.
別人說我是誰.
接著你說我是誰.
接著彼得說你就是上帝的彌賽亞.
你就是基督了.
但是這裡我想你想一下一件事.
就是彼得很強烈的告訴你.
我們已經信了.
我們已經知道你是誰.
認定了你.

$^{761}$彼得的肯定.
彼得給耶穌的鼓勵.
是真還是假.
第一.
第二.
你發覺耶穌對彼得的肯定和鼓勵.
好像沒什麼興趣.
沒聽進去.
因為彼得剛剛代表十二門徒.
一起和我們.
不是一個人.
我們全部都信了你.
我們全部都知道.
你就是上帝揀選的基督.
耶穌的回答是什麼.
你不是揀選了.
我不是揀選你們十二個嗎.
但是你們中間有一個魔鬼.
耶穌的意思是什麼.
就是說.
你代表了十二個人.
將你的餘生奉上.
說你多信我.
但是我一看.
我就知道.
起碼其中有一個不單是不信.
是魔鬼.
耶穌其實是在回答彼得.
彼得你很肯定.
我告訴你.
我也很肯定.
你看錯了.
彼得你根本看不見.
你同伴裡邊.
誰信誰不信.
但是我告訴你.
你看不見的東西.
我全部都看見.
甚至耶穌是在和彼得說.
其實彼得你也不用化這麼濃的妝.

$^{801}$你不用說得這麼大.
你以為我不知道你的狀況嗎.
你以為我不知道你們十一,十二個人的狀況嗎.
你們的心裡面所有的東西.
我都知道.
不用裝模作樣.
這個故事其實令我想起.
另外一個故事.
在《馬可福音》裡面彼得有相似的一些動作.
這個故事是我最喜歡的故事.
有一個年輕的財主來找耶穌.
耶穌說你欠一件的.
耶穌望著他就愛他說.
你變賣你所有來跟從我.
接著那個人就變了面色.
悠悠愁愁的走了.
因為他們他的產業很多.
接著下一節.
耶穌看了周圍.
周圍可能還有很多大排人跟著他.
望一望.
大家都怕耶穌.
接著就跟他的門徒說.
他說有錢人進上帝的國是何等的艱難.
在當時.
我繼續讀下去.
第二十四節.
門徒對這話非常驚訝.
耶穌再唱門徒.
再加幾成重.
孩子們要進上帝的國是何等艱難.
接著.
二十五節.
再加多幾成重.
駱駝穿過.
針眼比財主進入上帝的國更容易.
接著門徒就更加驚訝.
這些是沒有打出來的.
我讀給你聽.
門徒很驚訝彼此的對問.

$^{841}$且讓誰還能不救呢.
自己自言自語了.
耶穌望著他們說.
在人是不能的.
在上帝卻不言.
因為在上帝凡事都能.
接著不知道.
等了多少秒.
聽完耶穌的說話.
在上帝凡事都能.
兩秒.
五秒.
十秒.
三十秒.
有人說話了.
於是.
又是彼得.
彼得走出來.
主啊你看看.
我們已經撇下一切.
跟隨你了.
總要有人說話.
已經太久了.
耶穌說了一句話.
在人是不能.
在上帝凡事都能.
都是在等我們的回應.
無論如何都要演一場戲.
拜託彼得上去.
於是上帝彼得就說.
我們.
他又是代表其他人.
他不是說我.
他說我們.
已經撇下了一切.
去跟隨你.
這個我們已經.
就好像之前在《讓人保險第六》.
我們已經信了.
我們已經知道.

$^{881}$認定你就是那裡.
但你問.
這群門徒真的.
撇下了一切.
信了耶穌.
之前他們都在說.
那誰能救.
要放下所有.
你行嗎.
我行嗎.
當然不是.
我們來都是想賺多一點.
為什麼又要我放下.
我是想要賺多一點.
為什麼又要我放下.
但現在耶穌說成這樣.
你再不回應.
我們就不能救.
快點隨便說一些話.
於是我們已經撇下了一切.
其實他們沒有撇下一切.
很快過了一會兒.
他們會爭誰為大.
有兩個門徒.
要耶穌坐在耶穌的左右.
當他榮耀降臨的時候.
去到耶穌撒冷.
西馬利元彼得拿著一把刀.
割了.
那個人的耳朵.
他以為.
又是上帝天國臨近的時候.
結果他三次不認主.
到耶穌聽十字架的時候.
雞飛狗走.
只剩下約翰在耶穌身邊.
最後猶大出賣.
耶穌全部都是假的東西.
裝模作樣的東西.
回到約翰福音.

$^{921}$耶穌回答他們說.
不用裝.
我什麼都知道.
耶穌回答他們說.
我不是揀選了你們十二個嗎.
但你們中間有一個是魔鬼.
我全部都看透了.
第七十一節.
耶穌者說是.
只著要出賣他的家略人.
西門的兒子.
猶大說的.
他本是十二使徒的一個.
就這樣.
約翰福音第六章就完結了.
約翰福音第六章是最長的.
最豐富的一個單元.
這個單元劇就完結了.
你的感覺如何呢.
我讀完.
我感覺很不爽快.
為什麼這個故事.
因為之後的故事是一個新的故事.
這個故事就完結了.
這樣完結不行.
不僅沒有高潮.
低潮再低潮.
那算什麼福音好消息呢.
我們要解迷了.
我帶你回去看整篇.
講到今天經文的大圖畫.
你看看.
我們將最後的這一幕.
分成五部分.
這話很難.
誰能.
誰聽得盡.
靈似生命.
我說是靈.
蒙父恩賜.

$^{961}$到我這裡.
永生之道.
還跟從誰.
有個魔鬼.
出賣耶穌.
你看到什麼呢.
有沒有發覺其實.
中間二三四部分.
其實很不錯的.
很正的.
靈似生命.
我說是靈.
蒙父恩賜.
到我這裡來.
永生之道.
還跟從誰.
你不要當他亂講.
你當他真的講真的.
有一天他真的悔過的時候.
他真的說.
你有永生之道.
我還跟從誰.
其實我預備這篇讀的時候.
聽了很多不同的講章.
或者看了很多不同的.
有些人真的以為.
它們是真的.
但你想清楚一點.
其實是假的.
完全是搞笑的一幕.
我們知道二三四.
其實是正的.
現在壞就壞什麼呢.
壞起頭.
壞收尾.
是不是.
於是我想你看看.
我們怎樣可以將起頭和收尾.
修一修改.
其實最後.

$^{1001}$那個第五點.
是很容易修改的.
因為有隻魔鬼出賣耶穌.
出賣耶穌.
送耶穌上十字架.
耶穌釘死在十字架之前.
魔鬼大哭一聲.
我幻想而已.
成了.
魔鬼一說成了之後.
耶穌再說成了.
是不是.
所以你明白.
其實魔鬼說成了.
衰了.
它不是成了.
它是衰了.
而耶穌再回他一句.
又是成了的時候.
你會發現.
整個故事完全的反轉.
OK.
於是我將第五點現在改了.
成全那一下變成了.
十架軍王.
救贖成全.
It is finished.
是.
哇.
正了.
你現在由底.
讀上去.
是不是.
現在簡直好像一個火箭升空一樣.
OK.
底就是救贖成全了.
十架軍王.
救贖成全.
接著現在永生之道.
還根從遂.

$^{1041}$是不是.
現在永生之道已經成全了.
接著蒙附恩賜.
到我這對來.
靈賜生命.
我話是靈.
是不是.
整個火箭升空.
現在只差最後那個而已.
你去到上面.
恩典轉化它.
福音改造它.
起初現在還是.
這話很難.
耶穌你很難頂.
因為祂追求世界.
遂聽得盡.
我現在將它變了.
變了做什麼.
榮耀十架.
喜樂.
信服.
你現在有的.
就真是一個完整.
偉大的.
福音圖畫.
你怎樣去應用第六章.
第六章將你反轉的時候.
將所有東西當成全救贖的時候.
將所有東西由下面上去的時候.
最後得到.
那個應用就是榮耀的十架.
我們要喜樂的.
相信.
真相信.
然後順服跟隨.
耶穌呼召我們.
叫我們誓聽祂的說話.
你記得你沒辦法摸索.
天父的恩兆.

$^{1081}$你沒辦法摸索聖人的工作.
但是你可以捉住.
耶穌的說話.
我們細心聆聽耶穌的說話.
接著我們全人全心.
信靠跟隨.
聽祂的說話.
我帶你去看兩段聖經.
羅馬書第十章第十七節.
信心是從聽道來的.
聽道是從基督的話來的.
信心是從哪裡來的.
信心不是真正的信心.
不是從看神蹟而來的.
在我們現在現代.
甚至不是聽見證而來的.
不是聽神蹟見證而來.
信心是從聽真道而來.
聽耶穌的說話.
聽福音而來的.
有了這個信心.
我們就可以做.
楊寶欽第六章第十九節.
叫我們去做.
信上帝所差來的.
就是上帝的工作.
如果你看英文是相反的.
這是上帝的工作.
就是叫我們信.
上帝給了我們一個工作.
上帝要我們的工作就是要信耶穌.
什麼叫信耶穌.
是全人全心去信靠.
跟隨上帝所差來的耶穌.
我們過去兩個星期都在說.
信心.
教會告訴我們有三個角度.
或者三個元素.
第一個是有知識.
耶穌說來跟從我.

$^{1121}$耶穌說為什麼要跟從我.
因為你隨我以我別無請求.
世人都犯了罪.
虧欠你上帝的榮耀.
唯有在基督裡.
你才可以得到那個救贖.
你要永生嗎.
你就要來跟隨我.
這個知識上.
接著你認同了.
之前說這麼多神明.
我怎麼知道誰是真的.
我怎麼知道耶穌一定是真的.
可能佛祖也是真的.
可能什麼都是真的.
但是我現在知道了.
我認同了.
耶穌所說的是真的.
那你跟不跟隨他.
你除了有notitia之外.
你要有assensus.
你要同意耶穌所說的東西.
最後你要信靠.
你要委身佛地主.
完成我們擁抱.
我們跟隨耶穌.
我想你想一下一件事.
就是你記不記得.
當耶穌呼召他們去信他.
跟隨他的時候.
之後那一節說什麼.
之後那一節.
聽到的人.
他所謂的門徒.
就給了他一個挑戰.
哇 他說你玩到這麼大.
叫我們all in.
他問耶穌一個問題.
他說你行什麼神蹟.
好讓我們看見你而信你的.

$^{1161}$之前這裡耶穌挑戰他們.
你的工作就是要信上帝所差來.
就是他自己.
而你的信其實他們聽得明白.
完全聽得明白.
他們頂不住.
我們就聽不明白.
他們聽得明白.
他知道耶穌叫他們all in.
去跟隨他.
接著他們說好.
你現在玩到這麼大.
你行什麼神蹟.
好讓我們看見這個神蹟而信你的.
在約翰福音裡面.
耶穌是沒有直接回應他們這個要求的.
但我想帶你去看馬太福音.
如果你知道在馬太福音裡面.
有一幫民事和法理賽人.
他們來找耶穌.
不懷好意的.
他們問他說請你顯個神蹟給我們看.
挑戰耶穌.
你看看耶穌怎麼回他.
耶穌回答他們說.
你聽我讀沒有印出來的.
邪惡淫亂的世代.
求看見上帝的神蹟.
是吧.
打你兩下.
打你刮你兩巴掌.
但接著他竟然說.
除了先知若拿的神蹟以外.
再沒有神蹟給我們看.
有沒有神蹟看.
有.
打你兩下.
在這些邪惡淫亂的世界裡.
追求世界.
愛慕世界的人.

$^{1201}$我打你兩下.
下一句應該是沒有神蹟看.
走呀回家.
但他竟然說除了若拿的神蹟以外.
沒有其他.
不過除了有一個若拿的神蹟是什麼.
若拿三日三夜在大禹的抱腹中.
同樣人子要三日三夜在地裡面.
這是什麼神蹟來的.
耶穌的受死復活.
我的天啊.
耶穌應該送我們走.
下地獄死我們.
但他說我還有一個神蹟.
我是有一個神蹟.
是我的受死復活.
他們問耶穌給我們看一個神蹟.
你行什麼大神蹟.
讓我們可以全人全心信靠跟隨你.
耶穌的回答是.
我會為你奉上一切.
我的天啊.
我會為你奉上一切.
我會為你受死為你復活.
因為我心愛著你.
因為我愛你們到底.
我行了這個神蹟.
你絕對有所有的東西在手.
你來跟從我.
耶穌正在挑戰他們.
或者邀請他們.
你願不願意.
如果你見到這個神蹟.
你見到人子被高高舉起十字架的時候.
你知道這個神蹟達成的時候.
我求你悔改回轉.
全人全心信靠跟隨我.
因為除我以外再沒有永生.
接著耶穌帶我們回到.
那篇道的對話的前面.

$^{1241}$耶穌回答他們說.
我實實在在告訴你.
不要為那毀壞的食物操勞.
而要為那全道永生的食物操勞.
這食物是人子要賜給你們的.
如果你記得我解經的時候.
食物不是代表食物.
因為最後的食物是人子.
人子是但以理書第七章裡面的人子.
但以理書七章裡面的人子.
不給你食物的.
是給你天國的.
是給你永恆不扭壞的榮耀的.
於是耶穌在說.
你不要追求地上的國度.
不要追求今生的榮耀.
而是要追求天上的財寶.
永恆的榮耀和國度.
要為那東西操勞.
為什麼?.
因為這食物或者這國度.
這榮耀是人子.
但以理書第七章裡面.
我們現在要加進去.
是人子用十字譜交換回來的.
但以理書七章裡面沒有告訴我們.
人子怎樣可以領受國度.
但去到福音書裡面的時候.
我們知道耶穌是需要用他的十字架.
才可以將上帝的國度換回來的.
我們都聽過英國的著名詩人John Milton.
Paradise Lost and Paradise Regained.
那個樂園因為人的罪.
第一個亞當的罪而失落了.
現在第二個亞當來要拯救.
救贖那個樂園.
怎樣?唯有通過十字架.
是十字架讓我們得著.
人子的國度,天國的國度.
所以這個背景圖是不能走的.

$^{1281}$那個十字架永要在那裡.
我們現在慢慢去舊約聖經.
看著但以理書.
但是那個十字架不能掉下來.
因為我們現在看得清楚.
最後人子在新約裡面告訴我們.
是通過十字架將我們換取這個國度.
有一位像人子的.
駕駛著天上的雲而來.
被領導耿古臣在者面前.
他得了權並榮耀國度.
使各方各國各族的人都事奉他.
他的權並是永遠的.
不能敗去他的國度.
必不敗壞.
聖經告訴我們.
耶穌告訴我們.
我們既然從人子耶穌身上.
領受了這個永恆的國度.
我們應該怎樣.
我們應該一無掛慮.
奉上所有單求上帝的國.
我們回到新約聖經裡面.
我們到很熟的經文.
路加福音第十二章.
你們這小信的人.
你們不要求吃什麼喝什麼.
也不要掛慮.
不需要掛慮.
為什麼.
這都是世上愛邦人所求的.
你們所需要的.
你們的天父都知道.
你們只要求他的國.
就是一切東西.
都已經加給你們了.
你們為什麼一無掛慮.
因為你們這小群不要懼怕.
因為上帝的你們的父.
樂意把國賜給你們.

$^{1321}$為什麼我們不需要掛慮.
因為最寶貴的最榮耀的國.
上帝樂意賜給他們.
當我們有上帝的國.
我們就不應該再掛慮.
我們只管將我們的生命.
投注在永恆裡面.
你們要變賣財產.
預備永不壞的錢囊.
和用不盡的財寶在天上.
就是察不能近.
從不能住的地方.
因為你們的財寶在那裡.
你們的心也要在那裡.
上帝叫我們ALL IN.
我要結束了.
結束的時候.
我想跟大家分享.
早幾個禮拜.
我所寫了一個.
對於自己信主得救時候的反思.
我每隔幾年就會再重新再想回去.
我當年是怎樣信主的.
這次是因為.
我們將要開始MAX了.
早幾個禮拜.
我想做一個反思.
來鼓勵我們當中的.
參加MAX的大型節目.
我們當中可能三分之一的人有讀過.
讀過都可能會忘記了.
我想再用我當時所寫的分享.
去結束《約翰福音》第六章.
這個福音的故事.
鼓勵我們將我們的餘生獻給上帝.
這個分享有兩部分.
這個分享的題目叫做.
耶穌成為我救主的那一刻.
我為何每週都珍惜那一刻.
我十四歲時受洗成為基督徒.

$^{1361}$十八歲那年我買了一本新的聖經.
開始每天讀兩章.
並在大學二年級的時候讀完第一次.
大約在那時候.
我開始很天真的向上帝祈禱.
求上帝向我顯明祂的計劃.
假設祂真的對我有個計劃.
我當時並不指望上帝會回答我這個禱告.
於是我繼續按照我自己世俗的智慧.
去規劃我的人生.
我的計劃是先做兩年工.
然後去讀一個有名的MBA課程.
接著我想我的人生就成了.
當時真的這樣想.
與此同時 教會成為了我人生中.
不可缺少的一部份.
亦都是我喜樂的所在.
我喜歡預備查經.
我亦都喜歡和基督徒的朋友.
非基督徒的朋友分享屬靈的事和福音.
儘管我聽道的道的講道都不怎樣.
即是當時人們講的道都很差.
但我仍然熱切地去教會.
即使是最枯燥的講完.
我也不會睡著.
我熱愛做基督徒.
我熱愛聖經.
熱愛服侍.
也熱愛基督教信仰大級我的道德標準.
我也沒有自以為是.
不像我周圍所認識的一些人.
我因十字架而來的罪得赦免而歡欣快樂.
我當時的夢想是有一天我可以成為教會的長老.
我心裡想我會有成功的事業.
然後我將我空餘的時間,天賦,才能.
全部都奉獻給教會的事工.
但問題就在這裡.
我只會把剩餘的獻給上帝.
剩餘的時間,剩餘的精力.
剩餘的恩賜,剩餘的一切.

$^{1401}$老實說,我怎可能將上帝和剩餘相提並論呢?.
儘管我內心和我的外表有著一點的敬虔.
我卻好像該人一樣.
他是一個隨意的敬拜者.
我不是像阿伯那樣將他自己最好的獻給上帝.
這是我必須要翻轉的石頭.
唯有這樣,我才可以成為一個真正的基督徒.
上帝慈悲憐憫地向我啟示了那塊石頭.
上帝引導我,感動我.
溫柔,慈愛地感動.
但卻又非常堅定地感動我.
翻轉那塊石頭,使我可以得救.
使我可以為他和他的國所用.
當時我腦海裡浮現了一幅圖畫.
就是我要放棄一切.
前往一個福音未曾到達的地方宣教.
很明顯,在我自己讀MBA的夢想和宣教的夢之間.
我只可以二選一.
就如耶穌所說.
「一個人不能事奉兩個主」.
最終經過了一年多的掙扎.
那段時間,其實感覺上是自己在不斷地成長,進步.
而不是真的很辛苦地掙扎.
我因而選擇了後者.
舊有的夢想在十字架福音的光照之下.
越慢慢地消失.
而新來的夢想卻在永恆的價值.
榮耀和美好中不斷地增大.
這是一個新愛取代舊愛.
新歡勝舊愛的經驗.
在我迎來二十一歲生日之際.
那塊大石頭被翻轉了.
我也成為了一個新人.
不要放過任何一塊石頭.
你們多人都在思想.
但其意思並非要你尋找每一塊隱藏的石頭.
而是要讓上帝顯明那一塊或兩塊.
必須要被翻轉的石頭.
從而帶來一個真正的屬靈突破.
也許是你人生的規劃.

$^{1441}$也許是你的事業.
也許是某個關係.
你的安全感.
任何阻礙你被上帝使用的東西.
都必須被翻轉拿去.
你才能永遠屬上帝.
為他所用.
但請你不要害怕.
耶穌來不是為了毀滅我們的人生.
祂來是為了給我們唯一的新生命.
祂來要把祂的寶石來取替我們的破石頭.
我們必須相信這一點.
信靠祂.
因為祂會溫柔慈愛.
卻堅定不移地引導我們.
所以我請懇切的禱告.
求上帝顯明祂自己的心意在你身上.
這其實只是第一部分.
還有第二部分.
從另一個角度去看.
就完結了.
Rosario Butterfield曾在訪問中說.
她的罪不是她女同志的身份.
大家知道她的故事.
而是她的不信.
所謂的不信.
不是指她缺少了一個認罪的禱告.
而是她的一生.
全都活在一個沒有上帝的世界中.
因此聖國不信的方法.
不是口頭上重覆又重覆的認罪悔改.
是沒有用的.
而是全心全意轉向上帝.
祈求祂重新進入我們的生命.
掌管我們的一切.
Rosario將這種徹底的人生轉變.
恰當地描述為火車失事.
全被粉身碎骨.
全被摧毀的逆轉.
A train wreck conversion.

$^{1481}$這正是真信心的本質.
將我們人生完全的反轉.
這讓我想起我自己的迴轉故事.
並非十四歲時表面上的信主.
而是二十一歲時.
內心徹底降服於上帝的故事.
有時我會想.
如果我沒有踏出那一步真正信仰的跳躍.
我的人生將會如何.
如果我繼續走我的舊路.
去讀我的書.
建立我的事業.
同時在教會忠心服事.
當然所謂忠心服事.
只是用餘下的時間精神才幹.
因為世界已經奪去了我所有最好的一切.
我會單單失去了一些天上的財寶.
還是我會失去所有天上的財寶.
甚至失去了我永遠的生命呢.
每當我想起那些.
我心裡重新.
那場發生在我二十一歲生日前的人生逆轉.
翻轉我生命中的每塊石頭.
尤其是我曾經刻意隱藏的那塊石頭.
為的是要拯救我.
並且要救我到底.
將你的生命,天賦,計劃全交給耶穌.
這是你所能做到最好的決定,最好的投資.
這比他們據為己有要好得多.
好得無比的多.
人們說:在十字架面前放下你的罪.
耶穌說:把你沾滿罪性的生命,恩賜,計劃.
全部都交給我.
好讓我用我的寶血去救贖他們.
使他們成為天國的武器.
這麼多年過去了.
我真的無法想像.
把我的才能因此用在其他事情.
或者其他人的身上.
每次我打開PowerPoint.

$^{1521}$心裡面都充滿著一個屬天的喜樂.
如果祂沒有來呼召我歸祂所有.
單單讓祂使用.
我的人生就會完全被浪費了.
但因著祂的恩典.
也因著祂堅持要我將我所有的獻給祂.
如今我在祂裡面的人生.
有了目標,有了意義.
還有永恆的獎賞.
在《終極問答》第二章中.
我想像那個有罪的婦人.
人生有三個部分.
第一個部分是她過去在上帝以外的人生.
她犯罪.
沒有上帝在她人生裡面的人生.
第二部分是我們在《路加福音》當中.
所見到她迴轉感恩的場景.
第三部分將在《舍行傳》中出現.
她的人生如此.
我們的人生也應當如此.
我二十一歲時歡喜快樂信主.
那個是我人生的第二部分.
但接下來的與第一部分.
就是我人生的前二十一年的對比.
是更美好的一段.
救主耶穌曾經歡喜快樂為我寫明.
如今我也歡喜快樂為祂寫明.
這是何等奇妙的恩典.
如果我們的第二部分.
是短片的記錄下來的瞬間.
即是一個短片的幾秒鐘.
那麼第三部分就是一部完整的電影.
從第二部分一直延伸到永恆.
在永恆的日天所有電影.
所有基督徒的福音故事.
都會被投射在一個巨型的銀幕上.
同時播放.
到時所有人的目光.
都會清楚看見耶穌基督的十字架.
在人類歷史當中所留下的榮耀.

$^{1561}$罪惡(Guilt).
恩典(Grace).
和感恩(Gratitude).
就是基督徒人生的三部分.
第一部分.
第二部分.
第三部分.
好好珍惜你人生的第二部分.
但更要活出你人生的第三部分.
願一切榮耀歸上帝.
願一切讚美歸於我們的救主.
耶穌 我們一起祈禱.
我們是天父 我們多謝你.
在《約翰福音》第六章裡面.
耶穌講了一篇奇妙的道.
「我是生命量」.
這篇道當時的人沒有辦法聽得明白.
說真的 如果我們在當場.
我們也會聽不明白.
然而再過一年.
再過一年.
在十字架的光照之下.
所有的一切.
就完完全全的彰顯在我們的面前.
耶穌來 說要拯救我們.
而所謂的拯救不是一個廉價.
給一個天堂入場券的拯救.
耶穌真的要我們.
將我們從埃及裡面拯救出來.
從巴別裡面拯救出來.
耶穌要我們的人生.
不再為這個世界而活.
不再為自己而活.
而是單單為祂而活.
而祂這樣做.
不是將祂的喜樂.
建在我們的痛苦之上.
祂這樣做.
是要使我們得到真正的自由.
真正的喜樂.

$^{1601}$多謝你.
因為在整篇講到的起初.
耶穌已經將這個應用.
擺在我們的面前.
你們不要為這個世界操勞.
你們要為我為你們.
承傳的那個國度.
永恆的國度而操勞.
天父樂而將他的國賜給你們.
所以你們今天要擁抱天父對你的呼召.
願我們相信你.
願我們再沒有其他的憂心.
願我們一無掛慮的奔跑前面的道.
將我們的餘生完完全全的獻給你.
我為我們當中每一個年齡層.
每一個人生階段的弟兄姊妹禱告.
願我們都嘗試簡單的.
默默的話求主你.
將我餘生你對我的計劃啟示給我.
要我再不追求自己的計劃.
而是降服在你的計劃當中.
因為你的計劃一定比我們自己的計劃.
好得千千萬萬倍.
你要的不單止是我們的罪.
你要我們整個生命.
唯有這樣我們才可以被你所用.
將我們的人生投注在天國上面.
祝福我們使用我們整間教會.
我們這樣祈禱是奉主耶穌基督的命求.
阿們.
\newpage



\section{以弗所書 2:1-4:32}
\label{sec:031856}
\textbf{以弗所書 (2) 教會——基督的身體}~~~~劉少康牧師
\newline
\newline
連結: \href{http://www.church.com.hk/acms/content.asp?site=cdc\&op=show\&type=product\&code=031856\&layout=sermon}{\texttt{031856}} ~~~~ 語音日期: 2026-03-25
\newline
\newline
\hyperref[sec:031855]{\small{< < < PREV SERMON < < <}}
~
\hyperref[sec:index_chronic]{\small{[返順時目]}}
~
\hyperref[sec:031857]{\small{> > > NEXT SERMON > > >}}
\newline
\newline
以弗所書 2:1-4:32
\newline
\begin{longtable}{cl}
\hline
\hline
章節 & 經文 (和合本修訂版)\\
\hline
2:1 & \begin{tabularx}{0.7\textwidth}{X} 從前，你們因著自己的過犯罪惡而死了。 \end{tabularx} \\ \\ \relax
2:2 & \begin{tabularx}{0.7\textwidth}{X} 那時，你們在過犯罪惡中生活，隨從今世的風俗，順服空中掌權者的領袖，就是現今在悖逆的人心中運行的邪靈。 \end{tabularx} \\ \\ \relax
2:3 & \begin{tabularx}{0.7\textwidth}{X} 我們從前也都生活在他們當中，放縱肉體的私慾，隨著肉體和心中的意念去做，和別人一樣，生來就是該受懲罰的人。 \end{tabularx} \\ \\ \relax
2:4 & \begin{tabularx}{0.7\textwidth}{X} 然而，神有豐富的憐憫，因著他愛我們的大愛， \end{tabularx} \\ \\ \relax
2:5 & \begin{tabularx}{0.7\textwidth}{X} 竟在我們因過犯而死了的時候，使我們與基督一同活過來—可見你們得救是本乎恩— \end{tabularx} \\ \\ \relax
2:6 & \begin{tabularx}{0.7\textwidth}{X} 他又使我們在基督耶穌裡與他一同復活，一同坐在天上， \end{tabularx} \\ \\ \relax
2:7 & \begin{tabularx}{0.7\textwidth}{X} 為要把他極豐富的恩典，就是他在基督耶穌裡向我們所施的恩慈，顯明給後來的世代。 \end{tabularx} \\ \\ \relax
2:8 & \begin{tabularx}{0.7\textwidth}{X} 你們得救是本乎恩，也因著信；這並不是出於自己，而是神所賜的； \end{tabularx} \\ \\ \relax
2:9 & \begin{tabularx}{0.7\textwidth}{X} 也不是出於行為，免得有人自誇。 \end{tabularx} \\ \\ \relax
2:10 & \begin{tabularx}{0.7\textwidth}{X} 我們是他所造之物，在基督耶穌裡創造的，為要使我們行善，就是神早已預備好要我們做的。 \end{tabularx} \\ \\ \relax
2:11 & \begin{tabularx}{0.7\textwidth}{X} 所以，你們要記得：從前你們按肉體是外邦人，是「沒受割禮的」；這名字是那些憑人手在肉身上「受割禮的人」所取的。 \end{tabularx} \\ \\ \relax
2:12 & \begin{tabularx}{0.7\textwidth}{X} 要記得那時候，你們與基督無關，與以色列選民團體隔絕，在所應許的約上是局外人，而且在世上沒有指望，沒有神。 \end{tabularx} \\ \\ \relax
2:13 & \begin{tabularx}{0.7\textwidth}{X} 從前你們是遠離神的人，如今卻在基督耶穌裡，靠著他的血，已經得以親近了。 \end{tabularx} \\ \\ \relax
2:14 & \begin{tabularx}{0.7\textwidth}{X} 因為他自己是我們的和平，使雙方合而為一，拆毀了中間隔絕的牆，而且以自己的身體終止了冤仇， \end{tabularx} \\ \\ \relax
2:15 & \begin{tabularx}{0.7\textwidth}{X} 廢掉那記在律法上的規條，為要使兩方藉著自己造成一個新人，促成了和平； \end{tabularx} \\ \\ \relax
2:16 & \begin{tabularx}{0.7\textwidth}{X} 既在十字架上消滅了冤仇，就藉這十字架使雙方歸為一體，與神和好， \end{tabularx} \\ \\ \relax
2:17 & \begin{tabularx}{0.7\textwidth}{X} 並且來傳和平的福音給你們遠處的人，也傳和平給那些近處的人， \end{tabularx} \\ \\ \relax
2:18 & \begin{tabularx}{0.7\textwidth}{X} 因為我們雙方藉著他，在同一位聖靈裡得以進到父面前。 \end{tabularx} \\ \\ \relax
2:19 & \begin{tabularx}{0.7\textwidth}{X} 這樣，你們不再是外人或客旅，是與聖徒同國，是神家裡的人了， \end{tabularx} \\ \\ \relax
2:20 & \begin{tabularx}{0.7\textwidth}{X} 被建造在使徒和先知的根基上，而基督耶穌自己為房角石， \end{tabularx} \\ \\ \relax
2:21 & \begin{tabularx}{0.7\textwidth}{X} 靠著他整座房子連接得緊湊，漸漸成為在主裡的聖殿。 \end{tabularx} \\ \\ \relax
2:22 & \begin{tabularx}{0.7\textwidth}{X} 你們也靠他同被建造，成為神藉著聖靈居住的所在。 \end{tabularx} \\ \\ \relax
3:1 & \begin{tabularx}{0.7\textwidth}{X} 因此，我—保羅為你們外邦人作了基督耶穌囚徒的，替你們祈禱。 \end{tabularx} \\ \\ \relax
3:2 & \begin{tabularx}{0.7\textwidth}{X} 想必你們曾聽見神賜恩給我，把關切你們的職分託付我， \end{tabularx} \\ \\ \relax
3:3 & \begin{tabularx}{0.7\textwidth}{X} 用啟示讓我知道福音的奧祕，正如我以前略略寫過的。 \end{tabularx} \\ \\ \relax
3:4 & \begin{tabularx}{0.7\textwidth}{X} 你們讀了，就會知道我深深了解基督的奧祕； \end{tabularx} \\ \\ \relax
3:5 & \begin{tabularx}{0.7\textwidth}{X} 這奧祕在以前的世代沒有讓人知道，像如今藉著聖靈向他的聖使徒和先知啟示一樣， \end{tabularx} \\ \\ \relax
3:6 & \begin{tabularx}{0.7\textwidth}{X} 就是外邦人在基督耶穌裡，藉著福音，得以同為後嗣，同為一體，同為蒙應許的人。 \end{tabularx} \\ \\ \relax
3:7 & \begin{tabularx}{0.7\textwidth}{X} 我作了這福音的僕役，是照著神的恩賜，是照他運行的大能賜給我的。 \end{tabularx} \\ \\ \relax
3:8 & \begin{tabularx}{0.7\textwidth}{X} 雖然我比眾聖徒中最小的還小，他還賜我這恩典，讓我把基督那測不透的豐富傳給外邦人， \end{tabularx} \\ \\ \relax
3:9 & \begin{tabularx}{0.7\textwidth}{X} 又使眾人都明白甚麼是歷代以來隱藏在創造萬物之神裡的奧祕， \end{tabularx} \\ \\ \relax
3:10 & \begin{tabularx}{0.7\textwidth}{X} 為要在現今藉著教會使天上執政的、掌權的知道神百般的智慧。 \end{tabularx} \\ \\ \relax
3:11 & \begin{tabularx}{0.7\textwidth}{X} 這是照著神在我們主基督耶穌裡所完成的永恆的計劃。 \end{tabularx} \\ \\ \relax
3:12 & \begin{tabularx}{0.7\textwidth}{X} 我們因信耶穌，就在他裡面放膽無懼，滿有自信地進到神面前。 \end{tabularx} \\ \\ \relax
3:13 & \begin{tabularx}{0.7\textwidth}{X} 所以我求你們，不要因我為你們所受的患難喪膽；這原是你們的光榮。 \end{tabularx} \\ \\ \relax
3:14 & \begin{tabularx}{0.7\textwidth}{X} 因此，我在父面前屈膝— \end{tabularx} \\ \\ \relax
3:15 & \begin{tabularx}{0.7\textwidth}{X} 天上地上的各家都是從他得名的－ \end{tabularx} \\ \\ \relax
3:16 & \begin{tabularx}{0.7\textwidth}{X} 為要他按著他豐盛的榮耀，藉著他的靈，使你們內心的力量剛強起來； \end{tabularx} \\ \\ \relax
3:17 & \begin{tabularx}{0.7\textwidth}{X} 又要他使基督因著你們的信住在你們心裡，使你們既在愛中生根立基， \end{tabularx} \\ \\ \relax
3:18 & \begin{tabularx}{0.7\textwidth}{X} 能夠和眾聖徒一同明白基督的愛是何等的長、闊、高、深， \end{tabularx} \\ \\ \relax
3:19 & \begin{tabularx}{0.7\textwidth}{X} 並知道這愛是超過人的知識所能測度的，為要使你們充滿神一切的豐盛。 \end{tabularx} \\ \\ \relax
3:20 & \begin{tabularx}{0.7\textwidth}{X} 神能照著運行在我們心裡的大能充充足足地成就一切，超過我們所求所想的。 \end{tabularx} \\ \\ \relax
3:21 & \begin{tabularx}{0.7\textwidth}{X} 願他在教會中，並在基督耶穌裡，得著榮耀，直到世世代代，永永遠遠。阿們！ \end{tabularx} \\ \\ \relax
4:1 & \begin{tabularx}{0.7\textwidth}{X} 我為主作囚徒的勸你們，既然蒙召，行事為人就要與你們所蒙的呼召相稱。 \end{tabularx} \\ \\ \relax
4:2 & \begin{tabularx}{0.7\textwidth}{X} 凡事要謙虛、溫柔、忍耐，用愛心互相寬容， \end{tabularx} \\ \\ \relax
4:3 & \begin{tabularx}{0.7\textwidth}{X} 以和平彼此聯繫，竭力保持聖靈所賜的合一。 \end{tabularx} \\ \\ \relax
4:4 & \begin{tabularx}{0.7\textwidth}{X} 身體只有一個，聖靈只有一位，正如你們蒙召，是為同有一個指望而蒙召， \end{tabularx} \\ \\ \relax
4:5 & \begin{tabularx}{0.7\textwidth}{X} 一主，一信，一洗， \end{tabularx} \\ \\ \relax
4:6 & \begin{tabularx}{0.7\textwidth}{X} 一神－就是萬人之父，超越萬有之上，貫通萬有，在萬有之中。 \end{tabularx} \\ \\ \relax
4:7 & \begin{tabularx}{0.7\textwidth}{X} 我們每個人蒙恩都是照基督所量給每個人的恩賜。 \end{tabularx} \\ \\ \relax
4:8 & \begin{tabularx}{0.7\textwidth}{X} 所以有話說：「他升上高天的時候，擄掠了俘虜，將各樣的恩賜賞給人。」 \end{tabularx} \\ \\ \relax
4:9 & \begin{tabularx}{0.7\textwidth}{X} 既說「他升上」，豈不是指他曾降到地底下嗎？ \end{tabularx} \\ \\ \relax
4:10 & \begin{tabularx}{0.7\textwidth}{X} 那降下的，就是高升遠超越諸天之上的，為要充滿萬有。 \end{tabularx} \\ \\ \relax
4:11 & \begin{tabularx}{0.7\textwidth}{X} 他所賜的有使徒，有先知，有傳福音的，有牧者和教師， \end{tabularx} \\ \\ \relax
4:12 & \begin{tabularx}{0.7\textwidth}{X} 為要裝備聖徒，做事奉的工作，建立基督的身體， \end{tabularx} \\ \\ \relax
4:13 & \begin{tabularx}{0.7\textwidth}{X} 直等到我們眾人在信仰上同歸於一，認識神的兒子，得以長大成人，達到基督完全長成的身量。 \end{tabularx} \\ \\ \relax
4:14 & \begin{tabularx}{0.7\textwidth}{X} 這樣，我們不再作小孩子，中了人的詭計和欺騙的法術，被一切邪說之風搖動，飄來飄去。 \end{tabularx} \\ \\ \relax
4:15 & \begin{tabularx}{0.7\textwidth}{X} 我們反而要用愛心說誠實話，各方面向著基督長進，連於元首基督， \end{tabularx} \\ \\ \relax
4:16 & \begin{tabularx}{0.7\textwidth}{X} 靠著他全身都連接得緊湊，百節各按各職，照著各體的功用彼此相助，使身體漸漸增長，在愛中建立自己。 \end{tabularx} \\ \\ \relax
4:17 & \begin{tabularx}{0.7\textwidth}{X} 所以我這樣說，且在主裡鄭重地說，你們行事為人，不要再像外邦人存虛妄的心而活。 \end{tabularx} \\ \\ \relax
4:18 & \begin{tabularx}{0.7\textwidth}{X} 他們心地昏昧，因自己無知，心裡剛硬而與神所賜的生命隔絕了。 \end{tabularx} \\ \\ \relax
4:19 & \begin{tabularx}{0.7\textwidth}{X} 既然他們已經麻木，就放縱情慾，貪婪地行種種污穢的事。 \end{tabularx} \\ \\ \relax
4:20 & \begin{tabularx}{0.7\textwidth}{X} 但你們從基督學的不是這樣。 \end{tabularx} \\ \\ \relax
4:21 & \begin{tabularx}{0.7\textwidth}{X} 如果你們聽過他的道，領了他的教，因為真理就在耶穌裡， \end{tabularx} \\ \\ \relax
4:22 & \begin{tabularx}{0.7\textwidth}{X} 你們要脫去從前的行為，脫去舊我；這舊我是因私慾的迷惑而漸漸敗壞的。 \end{tabularx} \\ \\ \relax
4:23 & \begin{tabularx}{0.7\textwidth}{X} 你們要把自己的心志更新， \end{tabularx} \\ \\ \relax
4:24 & \begin{tabularx}{0.7\textwidth}{X} 並且穿上新我；這新我是照著神的形像造的，有從真理來的公義和聖潔。 \end{tabularx} \\ \\ \relax
4:25 & \begin{tabularx}{0.7\textwidth}{X} 所以，你們要棄絕謊言，每個人要與鄰舍說誠實話，因為我們是互為肢體。 \end{tabularx} \\ \\ \relax
4:26 & \begin{tabularx}{0.7\textwidth}{X} 即使生氣也不要犯罪；不可含怒到日落， \end{tabularx} \\ \\ \relax
4:27 & \begin{tabularx}{0.7\textwidth}{X} 不可給魔鬼留地步。 \end{tabularx} \\ \\ \relax
4:28 & \begin{tabularx}{0.7\textwidth}{X} 偷竊的，不要再偷；總要勤勞，親手做正當的事，這樣才可以把自己有的，分給有缺乏的人。 \end{tabularx} \\ \\ \relax
4:29 & \begin{tabularx}{0.7\textwidth}{X} 一句壞話也不可出口，只要隨著需要說造就人的好話，讓聽見的人得益處。 \end{tabularx} \\ \\ \relax
4:30 & \begin{tabularx}{0.7\textwidth}{X} 不要使神的聖靈擔憂，你們原是受了他的印記，等候得救贖的日子來到。 \end{tabularx} \\ \\ \relax
4:31 & \begin{tabularx}{0.7\textwidth}{X} 一切苦毒、憤怒、惱恨、嚷鬧、毀謗，和一切的惡毒都要從你們中間除掉。 \end{tabularx} \\ \\ \relax
4:32 & \begin{tabularx}{0.7\textwidth}{X} 要仁慈相待，存憐憫的心，彼此饒恕，正如神在基督裡饒恕了你們一樣。 \end{tabularx} \\ \\
[1ex]
\hline
\hline
\end{longtable}
$^{1}$各位弟兄姊妹,平安!.
我們來到《爾忽所書》的第二講.
就會講二,三,四,三章的書卷.
《爾忽所書》第二章一開頭第一節.
保羅就說:你們既然間是夢召.
因此你們就得著這個救恩.
然後他再強調:我們是因信而稱義的.
他說:我們夢召是藉著神的恩典得著這個救恩.
所以他由信徒的經歷開始來說.
就說到我們怎樣去經歷這位神.
所以是說到我們信仰上.
你必須要有體驗.
是我們經歷過.
是我親身試過的.
然後第二章第十節.
我很喜歡.
因為神說:我們是祂的Masterpiece.
我們是神手中的工作.
工作這個字眼可以翻譯做一個傑作.
如果我們有時去博物館參觀.
你會見到很多世界的名畫.
是某某十五世紀,十六世紀的名畫家畫的.
你又會見到一些很出名的雕塑.
是米開蘭基羅雕刻的.
你會眼前一亮.
原來一個這麼出名的大師.
他的雕刻就在我的面前.
有一次有一個人去一個雕刻大師的工場參觀.
見到這個大師對著一個石像.
他就看著看著.
對他來說應該可以收工的.
已經雕得這麼漂亮.
但他見到這個藝術家好像還不滿意一樣.
他把鼻子那裡又削一下.
把手臂那裡又削一下.
有時又用紙來磨滑它.
因為是大理石的雕像.
他要磨到很光滑.
很光滑才滿意.
原來這個就是大師的Masterpiece.

$^{41}$手中的傑作.
所以神原來說.
你,他,我,我們每一個都是神手中的傑作.
所以各位弟兄姊妹.
你千萬不要小看自己.
我沒什麼在幹.
教會的事奉.
找他找他不要找我.
我什麼都不懂的.
我不行的.
我是小薯仔一個.
沒有用的.
牧師不要找我.
那個很厲害的.
那個誰又讀過很多書.
誰有什麼恩賜.
我沒有的.
我什麼恩賜都沒有的.
不要找我.
是不是呢?.
錯了.
聖經裡面是保羅說的.
我們是神手中的工作.
我們夢照就是完成他交託給我們的職責.
所以神很看重你.
神看重我們每一個.
每一個都可以被神使用.
求主今日現在就使用你.
以至你可以為他做很多的工作.
然後第二章下半部.
就保羅說到.
我們門徒就要知道神.
認識這位神.
那裡說如今我們在基督耶穌裡面.
我們已經得清二了.
然後我們就可以將兩下合二為一.
拆去近時那個不和好的牆.
我們在基督裡面成為新的.
我們成為新造的人.
然後我們就與基督和睦.

$^{81}$神就將一個和好的職份託付給我們.
各位.
當你已經與神和好.
你就要勸人與神和好.
這個是神交託給我們的責任.
然後到第三章.
保羅將焦點轉到自己.
保羅在第三章說.
因此神父召我作外邦使徒.
專門向你們這些外邦人傳福音.
叫你們在基督裡面彼此合二為一.
都成為神國度裡面的人.
保羅用自己親身一個很寶貴很難得的經歷.
來講道我們傳福音過的積分.
還有在基督裡面是沒有種族.
沒有各樣不同差別之分.
在耶穌基督裡面.
我們都可以合二為一.
於是他心裡面有一個感觸.
基督高深莫測的愛.
有時在教會裡面敬拜的時候.
敬拜隊會帶我來唱一首聖詩.
How deep is God's love?.
神的愛有多深?.
很多時候唱.
我都不敢流出眼淚.
因為那首歌真的講到.
神的愛那種深度.
神的愛那種偉大.
那種無條件.
我們都是那些浪子.
但是神就好像那位慈父一樣.
祂遠遠見到我們就走過來.
抱緊我們.
可能你身上的衣服髒了一層.
因為你要慰注.
說很髒很一層脆的.
那個父親絕對不介意.
將最好的袍子鋪在你的身上.
將戒指戴在你的手上.

$^{121}$然後他說.
兒子歡迎你回家.
於是他回到家裡.
大排筵席.
宰了一隻肥牛.
要宴請所有的人.
各位,原來今天的父親就是這樣.
用無比的愛來接納我們.
所以保羅在第三章結束.
結束前面理論部分的時候.
他不禁發出一個頌讚的祈禱.
所以我們可以看見保羅的靈性.
他自己的靈性造詣.
他有感而發.
因為保羅同樣.
以前他是迫不及待教會的.
但自從在大馬色路上.
被神光照之後.
他整個人一百八十度的扭轉.
從一個要殺基督徒.
要拉基督徒的人.
變成福音的使者.
是向愛邦人傳這個福音.
好,我們來到第四章.
第四章就是實踐部分的第一章.
講到保羅勸弟兄姊妹要合一.
這個也是今天教會面對一個最大的難題.
你跟身邊的弟兄姊妹認識嗎?.
大家是否很熟呢?.
大家是否合而為一呢?.
有時我們進到教會.
到處打招呼.
「嗨,嗨,早晨,早晨,早晨,早晨」.
但都是表面的.
沒有裡面完全真誠的合一.
保羅在第四章說.
「我為主被囚的勸你們.
既然夢照行事為人.
就要與所夢的因上程.
然後又勸他們.

$^{161}$凡事謙虛,溫柔,忍耐.
用愛心互相寬容.
彼此要合一」.
所以這裡保羅形容了很多的「一」.
一共有七個「一」.
「一主一信一朕一臣」.
七個「一」.
原來我們要向著這個目標前進.
然後說我們在基督裡都是合而為一.
很偉大的信息.
亦都是今日教會裡.
我們很需要聽一篇這樣的道.
這篇道是保羅親自講的.
他說「我勸你們.
我為主被囚,我現在已經坐監.
我都勸你們,我大力地勸你們.
求神幫助我們.
聽到這位已經垂老.
坐緊監坐離開世界的這位使徒.
勸我們合而為一.
不要分黨派.
不要彼此結黨紛爭.
不要你氣我,我氣你.
我又不和你合作,你又不聽我的意見.
彼此走自己的路」.
因此保羅用七個「一」.
我們的神既然是合一.
一個主,一個信仰,一個神,一個真理.
然後我們都要合一.
然後我們教會要滿有基督長成的身.
教會是基督的身體.
所以我們各自作肢體.
但我們卻是合一的.
神給我們有不同的恩賜.
但我們混入不同的恩賜.
是叫這個身體漸漸增長.
在愛中建立自己.
然後我們又要穿上新人.
你要有新人的樣式.
你參加這個婚禮.

$^{201}$見到新娘一定是穿得最漂亮才出場.
她裝扮得最好.
特別早上找人化妝.
化個最漂亮的妝.
又弄好頭髮.
然後又老早選了一件很新款.
很新潮的婚紗.
然後以最漂亮的姿勢出現.
所以我們今天都要成為一個新做的人.
舊的事已經過去了.
你要做新的人.
有新的生活,新的樣式.
大家都知道有很多初信的材料.
就叫做新生活,新生命.
蘇永子牧師作的.
在香港特別好暢銷.
很多人都是讀這套書.
作為初信栽培.
你有一個新的生命.
所以你要有一個新的生活.
我也跟蘇牧師很熟.
有時他就說.
我為什麼寫這套書.
因為信了耶穌之後.
是價值觀念的改變.
所以蘇牧師很著重.
我們信耶穌之後.
我們有新的一套看東西的價值觀.
以往可能賺錢最重要.
以往很貪美.
以往很想出名.
以往很想出人頭地.
但是信了耶穌之後.
掉轉了我們的看法.
我們甘於服侍.
我們甘於作僕人.
施比受更為有福.
求主幫助我們在觀念上改變.
我們有新的生命.
因此我們活出新的生活.

$^{241}$大家將第二,第三,第四章詳細來看.
將主的話藏在心裡.
然後我們去實行出來.
天父別祝福你.
帶領你成為一個新造的人.
有新的樣式.
為主做美好的見證.
願天父賜福給你.
\newpage



\section{詩篇 34:1-7}
\label{sec:031857}
\textbf{當我害怕時 (華)}~~~~高志翔牧師
\newline
\newline
連結: \href{http://www.church.com.hk/acms/content.asp?site=cdc\&op=show\&type=product\&code=031857\&layout=sermon}{\texttt{031857}} ~~~~ 語音日期: 2026-03-22
\newline
\newline
\hyperref[sec:031856]{\small{< < < PREV SERMON < < <}}
~
\hyperref[sec:index_chronic]{\small{[返順時目]}}
~
\hyperref[sec:031858]{\small{> > > NEXT SERMON > > >}}
\newline
\newline
詩篇 34:1-7
\newline
\begin{longtable}{cl}
\hline
\hline
章節 & 經文 (和合本修訂版)\\
\hline
34:1 & \begin{tabularx}{0.7\textwidth}{X} 我要時時稱頌耶和華，讚美他的話常在我口中。 \end{tabularx} \\ \\ \relax
34:2 & \begin{tabularx}{0.7\textwidth}{X} 我的心必因耶和華誇耀，謙卑的人聽見就喜樂。 \end{tabularx} \\ \\ \relax
34:3 & \begin{tabularx}{0.7\textwidth}{X} 你們要和我一同尊耶和華為大，讓我們一同高舉他的名。 \end{tabularx} \\ \\ \relax
34:4 & \begin{tabularx}{0.7\textwidth}{X} 我曾尋求耶和華，他就應允我，救我脫離一切的恐懼。 \end{tabularx} \\ \\ \relax
34:5 & \begin{tabularx}{0.7\textwidth}{X} 仰望他的人，就有光榮；他們的臉必不蒙羞。 \end{tabularx} \\ \\ \relax
34:6 & \begin{tabularx}{0.7\textwidth}{X} 這困苦人呼求，耶和華就垂聽，救他脫離一切的患難。 \end{tabularx} \\ \\ \relax
34:7 & \begin{tabularx}{0.7\textwidth}{X} 耶和華的使者在敬畏他的人四圍安營，要搭救他們。 \end{tabularx} \\ \\
[1ex]
\hline
\hline
\end{longtable}
$^{1}$當我害怕時.
是我今天想要跟大家分享的題目.
弟兄姊妹有沒有害怕的時候呢.
有.
牧師站在這裡也好害怕.
所以當我害怕時.
也要講這個道.
跟大家開個玩笑.
今天早上起來.
有沒有什麼事情讓你感覺到很恐懼的呢.
什麼事情.
讓你感受到很恐懼.
害怕是我們人類的.
共通的語言.
每一個人從小到大.
常常我們都會有恐懼的情緒.
那我們今天呢.
要從大衛的.
什麼聲音.
要從大衛的詩篇裡面.
我們來看見.
當我們落在情緒低落也好.
痛苦也好.
困難也好.
就是不是這麼順暢的逆境的時候.
我們怎麼樣可以在這當中.
依靠上帝.
我們的心.
我們的靈能夠與神連結呢.
大衛的生死的危機.
我們都知道大衛.
他寫的詩歌呢.
很多時候都不是在順境的時候寫的.
都是在逆境的時候寫的.
特別是在他逃亡的這段期間.
聖經並沒有記載說.
大衛他逃亡.
花了多久的時間.
但是大概學者推估是七年到十年.
甚至可能更長一點.

$^{41}$但無論如何.
大衛他在逃亡的這段期間呢.
他經歷的是人生的.
從高山到低谷.
但是怎麼樣在這個當中.
我們可以看見.
神其實是在訓練他.
從一個牧人.
轉變成一個可以帶兵打仗的人.
甚至可以成為一國之君.
大衛的恐懼呢.
我們今天來看.
不是三層而已.
其實是很多層.
只是用三來代表說.
好像一層一層疊加在大委身上的.
大衛如何能夠寫出.
這麼美好的詩篇呢.
我們剛剛讀到的.
我的心.
必應耶和華誇耀.
謙卑人聽見要喜樂.
我的心.
要稱頌耶和華.
你們要和我一同.
稱頌耶和華為大.
我要時時稱頌耶和華.
讚美祂的話必藏在我口中.
大衛寫這個詩歌的時候.
不是在順境.
乃是他在逃亡的過程當中.
的一個處境.
甚至是他在很羞辱.
是很危急的時候.
他寫的.
那我在想 弟兄姊妹.
我們人生有沒有面臨許多的恐懼呢.
大衛面臨大掃羅的追殺.
他很無辜啊.
他可能會心裡想.

$^{81}$為什麼這件事情會臨到我身上.
我這麼忠心的事奉王.
而且上戰場去打仗.
竟然在打敗哥利亞.
這個巨人之後.
竟然掃羅王忌妒我.
要追殺我.
大衛陷入在這種長期的痛苦裡面.
弟兄姊妹.
你有沒有也有類似的.
不是被追殺.
而是一種好像長期的一個痛苦.
家庭的經濟負擔也好.
我們人生.
有沒有做錯過決策.
因為做錯了決策.
選錯了.
做錯了決定 你也感受到.
哎呀 我好後悔做這件事.
哎呀 我當初應該去學.
另外一個專業.
當初應該就不要來美國.
或者當初.
我就應該要及早做一個決定.
有沒有我們說錯話.
信錯人.
有這種懊悔在我們的心中.
另外大衛.
他面臨了生死的危機.
他感受到.
死亡就只差一步.
如果我們去讀薩姆爾記上.
你會看見.
掃羅王三番兩次.
想要拿槍 要刺大衛.
都給大衛躲過去了.
所以大衛就跟掃羅的兒子約拿丹說.
我離死亡.
就差一步.
弟兄姊妹 死亡就在隔壁.

$^{121}$當然我們不會覺得死亡就在隔壁.
但是我們會有.
有沒有痛苦到死的那種感受.
有沒有.
有沒有人被蟲咬痛苦到死.
癢死的有沒有.
有沒有被氣死的.
有沒有 被你小孩氣死的.
有沒有 被你老公氣死的.
是不是被你的爸媽不聽話.
氣死的.
我說爸媽不聽話.
我們父母年老的時候.
有時候好像不太聽話.
需要耐心照顧他.
有沒有.
這種快要死的感覺.
但是不是真的面臨死亡.
那有沒有真的.
我們好像看到了我們的.
我們的健康檢查的.
年終報告.
一年一次的報告.
紅字太多.
有沒有很擔憂.
我們如何在這樣的困境裡面.
我們可以像大衛一樣.
實時讚美神呢.
大衛的讚美.
並不是在順境中.
所以顯然讚美.
是一個意志的選擇.
大家同意嗎.
大衛他並沒有是在.
很順利的時候.
在那裡讚美.
反而是很多的時候.
都是他在逆境當中.
大衛的讚美是實時的.
也就是不是心情好.

$^{161}$或者事情順利的時候.
他才決定讚美.
而是當他的情緒.
害怕的時候.
他用意志決定.
他要來讚美這位上帝.
所以從大衛的身上.
到今天三千年後的.
我們今天會舉一個例子.
他叫克利.
無論是在當時的大衛.
還是三千年之後.
被關在集中營的一個.
淨淺的基督徒女子.
他們心中都有恐懼.
當然我們在讀聖經的時候.
都覺得大衛的讚美.
大衛的得勝.
好像是理所當然的.
你會不會覺得.
跟我有什麼關係.
你會不會覺得說.
是因為神節選大衛.
所以大衛可以這樣得勝呢.
其實大衛的生命裡面.
有很多的破碎.
是我們在詩篇裡面.
你仔細的去讀.
你就會發現.
很多的傷痕.
都是在他的低谷裡面.
那我們今天也會介紹.
克利.
這位姐妹的故事.
克利呢.
他是一個荷蘭的基督徒.
他的家庭裡面.
是一個有淨淺的父母.
那在二次世界大戰期間.
他們因為這個家庭.

$^{201}$他們的房子很特別.
前後有三戶.
本來是一戶.
然後買了後面一戶.
然後呢.
他的父親又想辦法把它連接在一起.
但是這個房子的格局.
不太一樣.
所以你看不出來他到底是二樓還是三樓.
你搞不清楚在房間裡面.
到底是二樓還是三樓.
就是因為這樣的一個房子.
特別呢.
他們家裡面.
二次世界大戰的時候.
藏了許多的猶太人.
保護了這些的猶太人.
他們被納粹.
我們說蓋世太保.
就是德國的秘密警察.
他們所抓去.
但是呢.
這個克利的父親有一個朋友.
知道這件事情就舉發他了.
於是他們全家.
都被捕.
甚至入獄.
那我看了這個克利的故事.
他就自己親自介紹.
他說他們被關.
他跟他的姐姐.
住了三百人的一個空間.
但是卻是住了七百人的.
這樣的一個集中營的房間裡面.
那裡面呢.
長滿了獅子.
和跳蚤.
兄姐妹.
你想想看.
你可以跟獅子.

$^{241}$和跳蚤共舞多久呢.
一天.
兩天.
還是三天呢.
還是一百天呢.
一天一個小時就夠了.
我們家的後院.
有那個跳蚤.
所以呢.
我去院子裡面的時候.
有一天就突然發現.
為什麼這裡癢那裡癢呢.
不是吃這個癢吃那個癢.
而是被跳蚤咬到很癢.
我還到了教會的時候.
弟兄姊妹幫我看.
這是什麼中藥的.
後來我想一想.
是跳蚤.
讓我感受到.
聖經裡面的一個經文.
就是從事不死的.
火是不滅的.
弟兄姊妹在地獄裡面從事不死的.
我們家的這個跳蚤.
讓我感受到.
我絕對不要跟跳蚤永遠的在一起.
那是太可怕的事.
但是這個顆粒.
這個顆粒呢.
他的姐姐叫做.
他的姐姐就跟他說.
我們要為這個跳蚤.
來凡事謝恩.
弟兄姊妹.
在一個很癢的狀況之下.
睡覺都睡不好的情況之下.
你還能夠來凡事謝恩嗎.
很難吧.
主啊.

$^{281}$放在我身邊.
邊講的時候邊搔癢.
弟兄姊妹.
被跳蚤咬得很多的時候.
你可能是很虛弱的.
是很痛苦的.
她睡的.
我是在半夜被養醒.
所以我知道那種痛苦的感受.
這個顆粒的姐姐.
她就是身體非常虛弱.
後來呢.
就死在這個集中營裡面.
但是她.
卻跟她的妹妹講說.
我們要為跳蚤凡事謝恩.
為什麼呢.
原來是神使用了這個跳蚤.
就是.
因為跳蚤很多.
所謂呢.
這個警察呢.
他們不會.
德國的軍人他們不會去檢查.
不會去檢查這個房間.
因為他知道那裡面有跳蚤.
所以他們都不敢接近.
所以在那裡面他們每天有兩次的聚會.
聖經的茶經.
讓許多的.
從未聽過耶穌的人就信了主.
很多跟他們.
關押在一起的人.
最後都被毒氣毒死了.
但是.
有許多的人在離世之前.
他們嘴裡呼喊著耶穌的名字.
他們能為永恆的靈魂.
來拯救這些的靈魂.
他們說.

$^{321}$即便他姐姐也死了.
這都是值得的.
其實.
這都還不算什麼.
顆粒他有一次他提到.
他說他在集中營裡面.
他經歷到最殘酷的.
大家猜猜看.
如果你是姐妹你是女生.
你覺得最殘酷的經歷是什麼.
他說他被脫光衣服.
有七次.
站立在那裡.
受檢就是要檢查.
所以他感受到非常的痛苦.
難堪這種羞辱.
他實在無法忍受.
他突然感受到.
聖經裡面耶穌基督.
也不過是像他.
這樣子.
耶穌基督曾經也是在石架上.
赤身肉體.
而且是掛在石架上.
透過他的痛苦.
他能夠體會到.
耶穌基督的一點點的受苦.
一點點的那種羞辱的感覺.
他說.
這件事情.
讓我非常的感恩.
讚美神.
我能夠承受這些的苦難.
其實克利小的時候.
他就很怕.
有一天他會為耶穌受苦.
他就問爸爸.
他就跟爸爸說.
我很擔心.
有一天我會為耶穌受苦.

$^{361}$他爸爸就跟他講了一個事情.
他就說.
如果我們要去阿姆斯特丹.
你是不是要去買票.
三個禮拜之後.
我們要去阿姆斯特丹.
我們要去那裡.
是不是要提早把票給你呢.
克利就回答說.
不是的.
爸爸是搭車的那一天.
票才會在我的手上.
他爸爸就跟他說.
同樣的.
你今天也不用擔心你為耶穌受苦.
當你有幸為他受苦的那一天.
那一刻.
神就會賜力量給你.
兄姐妹.
我們都怕為主受苦.
但是一個有信心的人.
當你在受苦的那一天.
主會給你讚美的力量.
阿們.
所以我們要.
時時讚美神.
不是在苦難的時候.
我才突然想到說.
我要來讚美神.
大衛他能夠時時讚美神.
是因為他本來就是一個讚美的人.
所以當他在害怕的時候.
他就很自然的.
能夠活出這樣的生命.
他來讚美神.
當讚美神的時候.
焦距.
剛剛我們說.
我們要將焦點轉向你.
我們剛剛唱詩歌的時候.

$^{401}$說焦點轉向你.
我們在這裡讚美神的時候.
就是把焦點轉向神.
不是看台上的敬拜的同工.
乃是看我們的神.
是我們當敬畏的.
當讚美的那一位.
大衛他在.
大衛在看到.
這個戰場上的.
非利士人哥利亞在那裡叫囂的時候.
大衛就問身邊的人.
這個沒有受割禮的.
非利士人是誰呢.
竟敢向永生神的軍隊罵政.
那當大衛呢.
他真的上了戰場之後.
跟這個哥利亞單挑的時候.
這個.
大衛說什麼呢.
大衛跟非利士人哥利亞說.
你來攻擊我.
我們高疆和同己.
但是我來攻擊你靠的是什麼呢.
萬軍之耶和華的名.
就是你所怒罵.
帶領以色列軍隊的神.
妞姐妹.
你看見大衛的讚美.
他模糊了焦點了.
他看見的那個敵人呢.
根本就不算什麼.
他看見的是什麼.
當他讚美神的時候.
他讚美的是那位.
永生神的上帝.
他讚美的是.
那位永生神萬軍的耶和華.
他要帶領我們.
今天就要打勝仗啦.

$^{441}$所以他在這個戰場上面.
他宣告.
得勝的是耶和華神.
今天.
神就要把你交在我的手裡了.
所以讚美.
可以幫助我們.
我們的焦距回到神那裡.
妞姐妹.
有太多的問題.
太多的困難.
我們都可以說.
我有理由不讚美.
可是聖經裡面提醒我們.
我們讚美神呢.
是用我們的心.
信心在那裡讚美.
無論環境如何.
但是我們始終要相信.
神是真實的.
神是大能的神.
神是拯救我們的神.
神是我們隨時的幫助.
所以我們的行動是什麼呢.
我們要建立.
唱詩歌的習慣.
請問一下剛剛在敬拜的時候.
妞姐妹你有大聲唱詩歌的.
我不要說大聲啦.
你有開口唱詩歌的請舉手好嗎.
大部分的人都會對不對.
可能這個歌不熟.
你可能會看一下.
但是以長久來講.
唱詩歌的習慣.
今天我們實在是.
這個場合不合適.
否則我們就來玩一個接龍的遊戲.
就是我們唱歌.
唱到那個字停下來.

$^{481}$然後你就接著.
唱那個最後一個字.
我的意思是說.
當我們裡面有很多的.
這些詩歌的元素.
很多詩歌在我們心中的時候.
我們就可以彼此唱和.
特別是在你困難的時候.
在你心中.
妞姐妹最常想起的那首歌是什麼呢.
你現在想想看.
這個月裡面.
你不在教會.
乃是你一個人的時候.
最常想到的詩歌是什麼呢.
有沒有.
主啊我讚美你.
有很多的詩歌.
都是我們可以朗朗上口的.
特別是短歌.
我們可以讚美神.
我今天特別寫了.
不是寫啦.
把一首歌放在這裡.
這是我在準備獎章的時候.
我心裡面就想到一首歌.
今天早上.
我還特別拿來唱了一下.
有多少妞姐妹會唱這個歌.
這是台灣的華人的.
讚美敬拜.
他們出的一個專輯.
歌的內容是這樣.
夜和花熄了林.
現在來充滿我.
是我臨場出席了讚美的詩歌.
我的口.
若不發出熄了的聲音.
主要心.
起石頭來讚美.

$^{521}$夜和花熄了林.
現在來充滿我.
是我臨場出席了讚美的詩歌.
我的口.
若不發出熄了的聲音.
主要心.
起石頭來讚美.
讓讚美的旋律.
充滿在我的心.
像大衛跳舞彈琴.
撐送住生命.
我的心.
我的林.
你要開口來讚美.
熄了又一刀.
磨在我身上.
我就越唱著.
越喜了弟兄姊妹.
如果你是我的粉絲.
我就越唱著.
越喜了弟兄姊妹.
如果你相信我們的神是真實的.
你也會相信這首詩歌.
當你越唱.
越唱神的林.
就會降在你身上.
阿們.
當掃羅.
嫉妒大衛的時候.
掃羅他感受到一個痛苦.
從神那裡來的惡魔.
聖經這麼記載的.
就是邪靈來繳饒掃羅.
掃羅都很痛苦.
掃羅就去找一個彈琴的.
於是大衛就來到他的身邊.
大衛就彈琴.
所以大衛是常常彈琴.
讚美神的.
當彈琴讚美神的時候.

$^{561}$甚至到一個程度.
神的同在.
可以驅敝邪靈.
大家有這樣的經歷嗎.
當我們讚美神的時候.
聖經告訴我們說.
魔鬼就逃跑.
所以讚美是有力量的.
如果你真的相信.
神是垂聽你禱告.
喜悅你的讚美的話.
我們的讚美.
可以要被提升的.
我們的讚美應該要提升的.
而且我們的讚美要帶出能力的.
我們的讚美能夠幫助我們裡面.
領受到神的靈的同在的.
這裡說到喜樂油已.
高抹在我身上.
神喜歡用他喜樂的油.
高抹我們勝過高他的同伴.
聖經裡面寫的.
其實他指的是指的耶穌基督.
但是今天我們呢.
神也樂意將他的喜樂的靈.
高抹在我們每一位弟兄姊妹身上.
Amen.
我們以前在台灣的時候.
我們聚集在一起常常拿著吉他一起唱詩歌.
就是一遍兩遍三遍四遍五遍.
我們就是喜歡.
這樣來讚美神.
並不是一個特別的聚會.
就是弟兄姊妹聚集在一起的時候.
沒有那麼多的Youtube.
我們就是喜歡.
一起在那裡唱詩歌.
讓我們在那裡寓意見神.
當我們在外面唱詩歌.
在公園裡面一起唱詩歌的時候.

$^{601}$就有很多人拿著相機.
手機.
要幫我們拍照.
甚至呢就會來.
於是我們就在那裡發單張.
傳福音.
當我們在那裡讚美神的時候.
那個氣場.
我們說氣場.
那個環境的氛圍就改變了.
環境的那個.
我們說那個氣氛.
就改變了.
在困難的時候.
要怎麼樣改變呢.
我們要讓我們的神與我們同在.
還記不記得.
我們上次弟兄父親節.
唱的這首歌呢.
師班還記得嗎.
這是我們師班弟兄之前獻的詩歌.
大衛他就是在這種光景裡面.
唱詩歌的.
於你神拉起他.
使他腳步重新站穩.
使他可以唱新歌.
兄弟妹.
唱新歌就代表我們再一次的經歷神.
所以讚美.
讚美的話.
如果你喜歡讚美就代表.
你是一個常常經歷神的人.
第二個我們要分享.
開口呼求神.
不只是讚美.
我們還可以把真實的感受.
帶到神的面前.
大衛其實沒有.
以一個英雄的角色.
來寫詩歌.

$^{641}$大衛他乃是寫說.
我是困苦人.
我這困苦人呼求.
葉和華變怎樣呢.
變怎樣.
垂聽.
大衛沒有說.
我這君王呼求.
神變垂聽.
他乃是說我是困苦人.
當克利.
剛剛說他們在集中營的時候.
他們剛剛被關到.
這個集中營裡面.
你們能不能夠想像.
700個人.
200個人的一個房間.
非常的小.
於是那裡就怎樣呢.
就發生了一件事情.
突然就有一種.
衝突突然爆發.
然後就有人開始.
咒罵毆打.
哭喊.
就像地獄一般.
充滿了仇恨和魔鬼的試探.
就在這個危機時刻.
他的姐姐.
就對他說克利我們來禱告.
於是他們就在那裡禱告.
他們不斷的懇求神.
然後就隨著他們的禱告.
他們說突然發生一件事情.
牢房就奇蹟式的安靜下來.
彷彿好像.
經過了猛烈的暴風雨之後.
的那種平靜.
直到一切完全的歸於停止.
他說.

$^{681}$這正是那個年邁飢餓的婦女.
就是貝琪.
她透過神所賜的禱告.
改變了這一切.
兄姐妹,困苦人的禱告.
神垂聽.
但是我們來看看大衛的困苦.
又是什麼樣的情況呢.
大衛他在逃難的時候.
他就為了躲避掃羅.
於是逃到了.
加特王雅吉那裡.
加特王是誰.
他是菲力士人.
這個哥利亞.
他的故鄉在加特.
所以大衛逃往哪裡.
不逃.
他躲到敵人的陣營裡面去.
而且這個是.
當時的這個哥利亞的家鄉.
雅吉的臣僕就對雅吉說.
這不是以色列國.
大衛嗎.
那裡的婦女.
跳舞唱歌.
掃羅殺死千千.
大衛殺死萬萬嗎.
大衛萬萬沒有想到.
他所逃過去的地方.
他以為沒有人知道.
竟然被認出來.
當時大衛還帶了一把刀.
這把刀是他在逃亡的時候.
逃到雅吉米勒那裡去.
他跟雅吉米勒說了一個謊.
跟他說.
掃羅王猜我有急事.
所以我就一個人來了.
這個植物讓我吃.

$^{721}$所以雅吉米勒就給他吃了橙色餅.
大家還記不記得這個典故.
然後呢.
也給他的那一把.
當時砍下.
哥利亞的那把刀.
就是哥利亞的刀.
大家會想說奇怪.
掃羅不是拿甩石去把哥利亞打敗的嗎.
其實他打敗.
當時他拿甩石打到哥利亞的爐裡面的時候.
是石頭進去.
哥利亞倒下.
然後拔出來砍了他的頭.
大衛拿著哥利亞的頭.
所以這把刀一直在雅吉米勒那裡.
大衛在逃亡的時候.
就逃到那裡去.
他就把這把刀拿走了.
所以呢.
大衛為了要躲避掃羅.
他就把這把刀帶著.
他萬萬都沒有想到.
他竟然被認出來了.
弟兄姊妹.
我在這裡寫了兩樣情.
大衛呢.
他因為太害怕這個雅吉.
太害怕雅吉王.
於是他就在當場.
在這個城門那裡.
胡亂畫字.
你知道最早的這個.
塗鴉派的始祖是誰呢.
哈哈哈.
大衛喔.
所以弟兄姊妹你不要找到教會的一面牆在那裡畫鴉.
我學大衛.
千萬不要.
大衛在那裡畫畫畫什麼不知道.

$^{761}$他可能覺得這還不夠.
要裝瘋賣傻才能夠.
survive下去才能夠活下去.
弟兄姊妹如果是你怎麼樣.
如果在那裡.
你就地取材你要怎麼樣呢.
吃泥土嗎.
吃大便嗎.
大家能不能想像.
你要大衛突然就要想.
我要裝瘋賣傻我才能夠survive.
我才能夠活下去.
所以他就地取材呢.
他就畫隨便畫一畫.
讓自己的口水流出來.
如果他有鼻涕.
我想他鼻涕也一起流.
大衛在那裡裝瘋賣傻.
弟兄姊妹.
我讀到這裡的時候呢.
這次我讀到這裡的時候.
我突然有一個.
心裡面有一個感觸.
弟兄姊妹.
大衛戴的這把刀.
是曾經在戰場上面.
他拔出敵人的刀.
砍下敵人的頭的那一把刀.
是得勝的刀.
當時大衛.
還記不記得我剛剛說.
大衛跟哥利亞叫囂的時候.
大衛說什麼.
你靠的是刀槍和銅戟.
但是我靠的是萬軍耶和華的名.
可是大衛今天卻要了那一把刀.
弟兄姊妹你知道我在說什麼嗎.
他原來是倚靠耶和華的名.
但是他今天.
他的安全感.

$^{801}$是那一把刀了.
這個地方是.
加特的城門.
弟兄姊妹聖經描述我們.
當大衛打敗哥利亞之後.
以色列起來就剿滅.
菲律斯人一直追趕他們.
追到哪裡呢.
聖經就寫到追到了加特的城門口.
今天.
大衛就在這個城門口.
做什麼呢.
在那裡化鴉.
過去的英雄.
過去大家都說.
大衛殺死.
千千大衛殺死萬萬.
這一位英雄.
今天卻在這裡裝白癡.
大衛的手在那裡化鴉.
弟兄姊妹.
而這一首詩篇.
就是大衛離開這個地方.
他活下來了.
他怎麼活下來呢.
他被雅吉王當作瘋子.
他說我們這裡.
豈缺少瘋子呢.
把他趕出去.
他是在這種羞辱的情況之下.
被趕出去的.
而他就用他這一支.
隨便化鴉的那個手.
寫了這個詩篇.
弟兄姊妹.
所以我們在詩篇34篇.
這裡寫到.
大衛在雅比米勒面前裝瘋.
被他趕出去就做這事.
所以弟兄姊妹.

$^{841}$我們的見證.
我們在跟人做見證的時候.
我們在歌頌我們的神的時候.
也許你在很平順的當中.
你覺得是理所當然的.
可是當我們在低谷的時候.
你更能夠感同身受.
主如何在這種痛苦的當中.
他扶持了我們.
弟兄姊妹.
我們信仰中最寶貴的.
就是神在我們.
處在破坑.
在淤泥當中.
拉起我們的那個過程.
是最寶貴的.
我們每個人.
生命中一定都會有.
經歷這樣的困苦.
這樣的痛苦.
然後讓我看見一個.
讓我更感動的一件事情.
是大衛.
他出來他寫詩的時候.
他不是去咒罵那些.
惡人.
他乃是在恐懼當中.
持續地來呼求神.
當然他也來.
省察他自己.
是不是我得罪神.
所以有一些詩篇裡面.
大衛讓我們看見.
他在那裡說.
神啊 求你赦免我.
我幼時所犯的一切的罪和過犯.
求你赦免我 塗抹我的過犯.
大衛在這一種光景中.
在恐懼的光景中.
大衛他求什麼呢.

$^{881}$大衛在那裡承認.
他的軟弱.
同時也讓我看見.
大衛在那裡尋求真理.
大衛在那裡求神.
用他的真理來指教他.
兄姐妹.
這一位大衛.
他的心裡面.
是一直都有正直良善的.
但是偶爾會有走偏路.
偶爾也會有像我們一樣.
會軟弱.
會走投無路.
會好像.
以前是依靠神.
但是突然呢.
好像神沒有垂聽我禱告.
我們裡面也會慌亂.
我們又會再去把這世界的刀.
世界所看重的.
戴在身上.
我們本來依靠的是神的名.
可是我們可能.
在生命的軟弱的時候.
我們會依靠的是人.
甚至我們會說.
我啥都不要靠.
靠自己.
於是我們裡面就形成了一個堅固的音類.
誰都不可靠.
靠我自己.
許多人很難接受耶穌.
他覺得這世界上沒有神.
我這麼一路苦過來了.
我都是靠我自己的.
所以我也不需要這位上帝.
兄姐妹.
神他就是要來拯救我們每一位.
如果我們裡面早知道.

$^{921}$祂是可靠的.
我們早就來尋求祂.
大衛他自己也經歷到.
他的軟弱.
他的痛苦.
所以我說.
他在這樣的一個困苦當中的時候.
他仍然抓住神.
神的道.
神的真理.
才是他唯一要抓住的.
而不是這世界上的誰.
也不是他自己.
乃是要抓住神的道.
即便我說.
他欺騙了祭司.
即便他在那裡裝瘋賣傻.
即便軟弱.
他相信神都垂聽他的禱告.
我剛剛說到.
我讀到大衛寫的詩篇.
最令我感動的.
除了我剛剛說的那一面.
更令我感動的是.
大衛他在困苦當中.
他求神指引.
他當行的道路.
在詩篇25篇裡面寫到.
耶和華求你將你的道指示我.
將你的路教訓我.
求你以你的真理引導我.
教訓我.
因為你是救我的神.
我終日等候你.
還有另外一個詩篇.
詩篇86篇.
也是一個困苦人的禱告.
他說諸神之中沒有人.
沒有可比你的.
你的作為無可比.

$^{961}$你所造的萬民都要來敬拜你.
他說因你為大.
且行奇妙的事.
唯獨你是神.
耶和華求你將你的道路指教我.
我要照你的真理行.
求你使我專心敬畏你的名.
後面講到.
我要榮耀你.
一生榮耀你.
大衛的禱告.
他不是跟神說.
主啊.
幫我解決那些難題.
他也沒有說.
主啊你不聽我禱告.
所以你看我這麼困苦.
他沒有在那裡埋怨神.
他反而是.
經歷了困苦.
他也回去檢視.
我是不是在這過程當中.
其實也得罪了神.
神啊求你赦免我.
神啊求你指引我當行的道路.
兄姐妹.
所以你看.
詩篇第34篇第11節.
今天我們讀的詩篇.
34篇第1節.
34篇11節裡面.
講到的是.
大衛就邀請人來聽.
神的教訓.
我讀給各位聽.
詩篇34篇第11節.
裡面寫到.
對不起第11節.
眾弟子啊.
你們當來聽我的話.

$^{1001}$我要將耶和華的道路.
耶和華的道.
教訓你們.
我再讀一次.
眾弟子啊.
你們當來聽我的話.
我要將敬畏耶和華的道.
教訓你們.
兄姐妹.
在我們的生命裡面.
我們不是只有在順利的時候.
榮耀的時候.
求是神.
你看大衛.
他也會說謊.
他也會軟弱也會保命.
但是他心裡面知道.
他一生.
不管怎麼樣.
他就是要求神來指示他.
當行的道路.
所以弟兄姊妹.
我們每一位.
不管你現在的處境如何.
不管你是成功還是失敗.
不管你所遇到的困難.
感受良好.
有沒有讓你感受到.
你不想服事神的.
這種痛苦呢.
或者我乾脆隱藏起來了.
我這麼失敗.
我過去做了什麼什麼的服事.
好像蠻失敗的.
於是我再也不要這樣服事神了.
有沒有這樣的一個悔恨.
一個痛苦在你的心裡面.
是長久以來.
沒有人發覺的.
或者你過去曾經.

$^{1041}$在一個團契裡面.
你傷害了某個人.
你說了一些不當說的話.
但是到如今.
你還在悔恨連連.
可能大家都忘記了那個時刻.
你說了什麼你做了什麼.
但是那一幕一幕的畫面.
仍然在你的心中裡面.
好像成為一個.
不能說的痛.
一個隱藏的.
有沒有一個懊悔.
是你曾經所犯下的錯.
可能你說了一個謊.
你想盡辦法再來圓那個謊.
以至於最後你身敗名裂.
大衛讓我們在他的詩裡面.
讓我們看見.
他緊緊地抓住神.
唯有神的話語.
是他所要追尋的.
最後我們來看第三點.
這是我這邊的第二點.
最後的一個結論說到.
即使我們在軟弱中.
真理的光.
是讓我們有出路的.
就算魔鬼定罪你.
就算你的心裡面.
給自己的一個評價是不好的.
就算別人怎麼樣說你.
弟兄姊妹.
你要相信神的話語.
會讓你的生命有出路.
阿們.
第三個.
交給神.
大衛說.
耶和華的使者在敬畏他的人.

$^{1081}$視為安寧.
搭救他們.
神使用我們.
神搭救我們.
並不是在.
我們裡面.
你都沒事的時候搭救.
一定是在你生命裡面.
有事的時候.
你看一下.
主的恩典.
通常都是在我們裡面.
需要恩典的時候.
神的恩典給我們用.
所以保羅說.
主的恩典是給我們用的.
我們何時軟弱.
何時就可以靠著神的恩典剛強.
為什麼我們可以放心呼求呢.
大衛他的軟弱是很現實的.
但是大衛有另外一個屬靈的身份.
他是很真實的.
大衛他知道他自己是一個敬畏神的義人.
在詩篇裡面.
他有讓我們看見.
義人的路是越造越明的.
即便你走偏了.
神也不讓你全然跌倒.
就算你軟弱.
神也要扶持你的.
神會保守你.
會引導你走當行的義路.
不管你往左或.
向左或向右.
後面有聲音說.
你要行在其間的.
義人的路是這樣的路.
是被神恩帶賜福的路.
在大衛的詩篇裡面.
他一直常常強調.

$^{1121}$神是看顧義人的.
所以大衛其實他就.
站在義人的一個角度裡面.
各位姐妹.
當你受壓,當你困難的時候.
你的心思意念你是站在.
哪一個角度.
你是不是站在一個義人的位份上.
我們是不是站在.
神兒女的位份上.
信心的眼睛.
是讓我們看見.
所看不見的.
而看不見的其實比看得見的.
更加的真實.
你看見的.
大衛看見的或世人所看見的.
那個哥利亞,大家都以為他很厲害.
對吧.
真實嗎.
世人以為那很真實很厲害.
其實一點也不厲害.
大衛真正看到的更真實的是.
他靠著萬軍之夜耀華的名.
可以打敗這個哥利亞.
所以信心的眼睛.
是需要.
被開啟的.
信心從哪來.
從神的話而來.
所以大衛他能夠看見.
那所看不見的.
在聖經裡面.
耶撒法他也曾經.
經歷到上帝.
派天君天使.
來幫助他打仗.
他什麼都不用做.
在敵人面前只要擺好陣勢.
神就派兵.

$^{1161}$來擊殺仇敵.
耶耀華的使者.
在敬畏他的人.
侍為安寧.
搭救他們.
兄弟姐妹你相信你的神也是搭救你的神嗎.
你相信.
當你在患難中.
神會來搭救你嗎.
大衛他.
在他的詩篇裡面.
他說啊.
凡投靠你的願他們喜樂.
時常歡呼.
因為你護蔽他們.
又願那愛你民的人.
都靠你歡喜.
所以大衛其實.
他就是一個愛慕神民的人.
所以當困難來的時候.
他可以把問題.
交給他.
當問題來的時候.
他的問題.
多小啊.
是一個以色列國.
全國的人都要搜索.
他們的一個危難.
大衛曾經也經歷到了.
這樣的一個仇敵四面.
來包圍他.
躲去哪裡呢.
只能躲到洞裡.
他說呢.
這個.
因為我遭患難的時候.
他必暗暗的保守我在他的亭子裡.
把我藏在.
他帳幕的隱密處.
將我高舉在磐石上.

$^{1201}$他說現在我得以昂首.
高過四面的仇敵.
大衛真實經歷到.
上帝的.
保守.
好像神真的派軍兵.
來幫助他.
所以雖然仇敵.
安營攻打.
他說什麼呢.
我的心怎樣.
我的心也不害怕.
雖然有刀劍.
向我攻擊.
我的心仍然安定.
我們剛剛說到.
那個集中營.
最後一個故事.
克利他剛剛進去的時候.
他就帶了一本聖經.
新約的聖經.
他在做見證的時候.
還拿著他那本新約的聖經.
他說他就是帶著這本聖經.
進集中營的.
英文姊妹.
我剛剛不是說他們很嚴格嗎.
他有一次還被全身脫光了.
七次對不對.
你們覺得他怎麼帶聖經進去的呢.
他就心裡面做了一個.
非正統的禱告.
而且是信心的禱告.
他就心裡想說.
我要求神.
差遣天使來保護他.
但是他又想到.
但是天使都是看不見的.
如果天使保護他的話.
有保護和沒保護是一樣的.

$^{1241}$所以他就改變了一個禱告.
他說主啊.
這次請你讓天使不要成為透明的.
結果他就進去了之後.
他說神蹟真的發生了.
他說守衛呢.
就仔細的檢查他前面的那個人.
還有他後面的姊姊.
結果就沒有檢查到他.
什麼意思呢.
守衛眼瞎了嗎.
好像上帝差遣了一個天使.
把他遮蓋了.
那個守衛沒有看見他.
於是他就把這本聖經呢.
就帶進了.
這個集中營的.
這個牢房裡面.
他們就是用這一本聖經.
在那裡聚會.
查考神的話.
在那裡傳福音.
把福音傳給很多的人.
兄弟姐妹.
我們的上帝是全知全能的神.
祂是把我們隱藏在隱秘處的神.
但是重點是我們要進去啊.
我們要投靠他啊.
克利他最後.
他一生.
有一句很有名的名言.
他說.
Never be afraid to trust.
an unknown future to an unknown God.
意思就是說.
你不要害怕.
永遠不要害怕.
把未知的未來.
交給全知的神.
我們永遠都要記得.

$^{1281}$當我們在.
害怕的時候.
千萬要把你的.
未知的害怕.
未知的恐懼.
甚至對未來的不確定.
通通的帶到.
這位全知全能的神的手中.
安心的交給他.
所以我們最後.
為大家預備的不是.
不是Lunchbox.
也不是地震包.
而是屬靈的.
屬靈的急救包.
平常我們就要.
常常的讚美神.
去唱詩歌的習慣.
去背誦一些的經文.
當我們需要.
遇到困難的時候.
一起禱告開口呼求神.
然後用信心的眼睛.
交託給神.
兄姐妹.
你有屬靈的急救包嗎.
想想看.
有多少詩歌是你記得的.
有多少經文是你可以隨時調用出來的.
那有多少的信心.
是你可以.
把你的各樣的困難.
帶到禱告前的.
\newpage



\section{撒母耳記上 23:1-29}
\label{sec:031858}
\textbf{面對困境有出路}~~~~李思聰牧師
\newline
\newline
連結: \href{http://www.church.com.hk/acms/content.asp?site=cdc\&op=show\&type=product\&code=031858\&layout=sermon}{\texttt{031858}} ~~~~ 語音日期: 2026-03-22
\newline
\newline
\hyperref[sec:031857]{\small{< < < PREV SERMON < < <}}
~
\hyperref[sec:index_chronic]{\small{[返順時目]}}
~
\hyperref[sec:031859]{\small{> > > NEXT SERMON > > >}}
\newline
\newline
撒母耳記上 23:1-29
\newline
\begin{longtable}{cl}
\hline
\hline
章節 & 經文 (和合本修訂版)\\
\hline
23:1 & \begin{tabularx}{0.7\textwidth}{X} 有人告訴大衛說：「看哪，非利士人攻擊基伊拉，搶奪禾場。」 \end{tabularx} \\ \\ \relax
23:2 & \begin{tabularx}{0.7\textwidth}{X} 大衛求問耶和華說：「我可以去嗎？我可以去攻打那些非利士人嗎？」耶和華對大衛說：「你可以去攻打非利士人，拯救基伊拉。」 \end{tabularx} \\ \\ \relax
23:3 & \begin{tabularx}{0.7\textwidth}{X} 大衛的人對他說：「看哪，我們在猶大這裡尚且懼怕，何況到基伊拉去攻打非利士人的軍隊呢？」 \end{tabularx} \\ \\ \relax
23:4 & \begin{tabularx}{0.7\textwidth}{X} 大衛又再求問耶和華，耶和華回答說：「你起身下基伊拉去，我必將非利士人交在你手裡。」 \end{tabularx} \\ \\ \relax
23:5 & \begin{tabularx}{0.7\textwidth}{X} 於是大衛和他的人往基伊拉去，與非利士人打仗，大大擊敗他們，奪取他們的牲畜。這樣，大衛救了基伊拉的居民。 \end{tabularx} \\ \\ \relax
23:6 & \begin{tabularx}{0.7\textwidth}{X} 亞希米勒的兒子亞比亞他逃往基伊拉到大衛那裡的時候，手裡拿著以弗得。 \end{tabularx} \\ \\ \relax
23:7 & \begin{tabularx}{0.7\textwidth}{X} 有人告訴掃羅，大衛到了基伊拉。掃羅說：「神將他交在我手裡了，因為他進了有門有閂的城，把自己關起來了。」 \end{tabularx} \\ \\ \relax
23:8 & \begin{tabularx}{0.7\textwidth}{X} 於是掃羅召集眾百姓，要下去攻打基伊拉，圍困大衛和他的人。 \end{tabularx} \\ \\ \relax
23:9 & \begin{tabularx}{0.7\textwidth}{X} 大衛知道掃羅設計陷害他，就對亞比亞他祭司說：「把以弗得拿過來。」 \end{tabularx} \\ \\ \relax
23:10 & \begin{tabularx}{0.7\textwidth}{X} 大衛說：「耶和華－以色列的神啊，你僕人確實聽見掃羅設法要到基伊拉來，為我的緣故毀滅這城。 \end{tabularx} \\ \\ \relax
23:11 & \begin{tabularx}{0.7\textwidth}{X} 基伊拉人會把我交在掃羅手裡嗎？掃羅會下來，正如你僕人所聽見的嗎？耶和華－以色列的神啊，求你指示僕人！」耶和華說：「他會下來。」 \end{tabularx} \\ \\ \relax
23:12 & \begin{tabularx}{0.7\textwidth}{X} 大衛又說：「基伊拉人會把我和我的人交在掃羅手裡嗎？」耶和華說：「他們會交出來。」 \end{tabularx} \\ \\ \relax
23:13 & \begin{tabularx}{0.7\textwidth}{X} 於是大衛和他的人約有六百名起身離開基伊拉，往他們所能去的地方去。有人告訴掃羅，大衛離開基伊拉逃走了，掃羅就停止出發了。 \end{tabularx} \\ \\ \relax
23:14 & \begin{tabularx}{0.7\textwidth}{X} 大衛住在曠野的山寨裡，在西弗曠野的山區。掃羅天天尋索大衛，神卻不將大衛交在他手裡。 \end{tabularx} \\ \\ \relax
23:15 & \begin{tabularx}{0.7\textwidth}{X} 大衛看到掃羅出來尋索他的命。那時，他住在西弗曠野的樹林裡； \end{tabularx} \\ \\ \relax
23:16 & \begin{tabularx}{0.7\textwidth}{X} 掃羅的兒子約拿單起身，到樹林裡去見大衛，使他的手倚靠神得以堅固， \end{tabularx} \\ \\ \relax
23:17 & \begin{tabularx}{0.7\textwidth}{X} 對他說：「不要懼怕！我父掃羅的手無法害你 ，你必作以色列的王，我必作你的宰相。我父掃羅也知道這事。」 \end{tabularx} \\ \\ \relax
23:18 & \begin{tabularx}{0.7\textwidth}{X} 於是二人在耶和華面前立約。大衛仍住在樹林裡，約拿單就回家去了。 \end{tabularx} \\ \\ \relax
23:19 & \begin{tabularx}{0.7\textwidth}{X} 西弗人上基比亞到掃羅那裡，說：「大衛不是在我們那裡，在樹林裡的山寨中，在荒野南邊的哈基拉山藏著嗎？ \end{tabularx} \\ \\ \relax
23:20 & \begin{tabularx}{0.7\textwidth}{X} 現在，王啊，請隨你的心願要下來，就請下來；至於我們，一定會把他交在王的手裡。」 \end{tabularx} \\ \\ \relax
23:21 & \begin{tabularx}{0.7\textwidth}{X} 掃羅說：「願耶和華賜福給你們，因為你們體恤我。 \end{tabularx} \\ \\ \relax
23:22 & \begin{tabularx}{0.7\textwidth}{X} 請你們回去，再確定一下，調查並看清楚他落腳的地方，是誰看見他在那裡，因為有人告訴我他很狡猾。 \end{tabularx} \\ \\ \relax
23:23 & \begin{tabularx}{0.7\textwidth}{X} 你們要看清楚，調查他藏匿的每一個地方，回來給我確實的報告，我就與你們同去。他若在境內，我必從猶大的千門萬戶中搜出他來。」 \end{tabularx} \\ \\ \relax
23:24 & \begin{tabularx}{0.7\textwidth}{X} 西弗人動身，在掃羅以先往西弗去。大衛和他的人卻在瑪雲曠野，在荒野南邊的亞拉巴。 \end{tabularx} \\ \\ \relax
23:25 & \begin{tabularx}{0.7\textwidth}{X} 掃羅和他的人去尋索大衛。有人告訴大衛，他就下到巖石那裡，留在瑪雲的曠野。掃羅聽見了，就在瑪雲的曠野追趕大衛。 \end{tabularx} \\ \\ \relax
23:26 & \begin{tabularx}{0.7\textwidth}{X} 掃羅在山的這一邊走，大衛和他的人在山的那一邊。大衛急忙躲避掃羅，掃羅和他的人正四面圍住大衛和他的人，要捉拿他們。 \end{tabularx} \\ \\ \relax
23:27 & \begin{tabularx}{0.7\textwidth}{X} 有使者來對掃羅說：「非利士人入境搶掠，請快快回去！」 \end{tabularx} \\ \\ \relax
23:28 & \begin{tabularx}{0.7\textwidth}{X} 於是掃羅不再追趕大衛，回去迎擊非利士人。因此那地方名叫西拉‧哈瑪希羅結。 \end{tabularx} \\ \\ \relax
23:29 & \begin{tabularx}{0.7\textwidth}{X} 大衛從那裡上去，住在隱‧基底的山寨裡。 \end{tabularx} \\ \\
[1ex]
\hline
\hline
\end{longtable}
$^{1}$今天我們來看薩姆爾紀上的第23章.
上次我們看了第22章的經文.
今天我們來看第23章.
第23章我們看這幅畫都知道.
就是說到掃羅在這段期間裡.
大衛被掃羅不斷追殺.
窮追不捨.
追到哪裡呢.
我們看第26節這個經文就說到了.
我們一起讀一讀.
第23章第26節.
請讀.
掃羅在山這邊走.
大衛和跟隨他的人在山那邊走.
大衛急忙躲避掃羅.
因為掃羅和跟隨他的人.
四面圍著大衛和跟隨他的人.
要拿獲他們.
所以見到這是一個很危險的情況.
大衛可以說是身陷險境.
非常危險.
危險到大家同在山那邊.
一個在這邊走.
另外大衛和跟隨他的人就在那邊走.
所以我們知道掃羅不是一個人.
大衛也不是一個人.
上次我們說跟隨他的人有多少人.
有400人.
當然不知道這個時候是否全部一起.
無論如何有很多一大群人.
山也不算很大.
在猶太岡野的山大也有限.
但我們知道危險的地方.
就是掃羅的那群人在山這邊.
大衛也在山另外一邊.
所以大衛急忙躲避掃羅.
因為掃羅和跟隨他的人.
已經四面圍著大衛.
所以可以說不單止一人在那邊.
而且他知道他大概就在這個山那邊.

$^{41}$已經派在山下面的整個山腳圍著.
差不多無路可走了.
所以在一個絕大的困境當中.
可以說是沒有出路.
不知道我們人生人會不會有時候感覺到沒有出路.
即是說無論在生活上有時候有困難.
或者在工作上有時候是疾病或者是經濟問題.
或者是我們家庭的問題.
又或者可能在靈性上.
覺得自己好像面對很多的.
無論是試煉或者是軟弱.
或者是一些瓶頸.
甚至是一些在屬靈上的靈情上.
有時候有些困逼.
可能被人誤解.
可能被人誣告.
或者在時空當中的失落.
很多這些在我們人生當中.
魔鬼藉著各樣的方法.
來叫我們灰心.
叫我們放棄.
叫我們失去我們的信心.
叫我們去退去.
所以可能我們人生都會有不同程度.
或者不同時期.
甚至有時候覺得自己好像沒有什麼出路.
這種叫做困境.
所以今天我們在這個主題裡面.
就是用這個作為題目.
就是面對困境.
不過有出路.
所以這個就是最大的喜訊.
面對困境有出路.
因為上帝使我們有盼望的神.
所以無論如何我們都是能夠有盼望的.
事實上聖經裡面也都應許我們.
就是說我們所遇見的試探.
無非是人所能受的.
神是信實的.
不叫我們受試探過於我們所能受的.

$^{81}$他說在受試探的時候.
總要給你們開一條什麼.
給我們開一條出路.
讓我們能忍受得住.
這個就是聖經給我們的應許.
事實上我們在過往的信徒的經驗.
人生的歷程裡面.
我們都是有多次困境也都會有.
但出路總是有.
所以這個就是神給我們的恩典.
也是他的應許.
我們今天也都看見.
大衛在這個很明顯的困境裡面.
如何能夠有出路呢.
第一我們看見.
經文給我們看到的.
就是他不靠他自己的聰明.
三國義記上23章第9至13節.
我們一起慢慢讀完這幾節.
剛纔是第23章的26節.
我們看一看之前他如何來形容大衛.
23章第9節請讀.
大衛知道掃羅設計謀害他.
就對祭司阿比亞他說.
將以弗德帶過來.
大衛禱告說.
耶穌說:以色列的上帝啊.
你僕人聽真了.
掃羅要往基爾拉來.
為我的原國滅城.
基爾拉人將我交在掃羅手裡不交.
掃羅照著你僕人所聽的話下來不下來.
耶穌說:以色列的上帝啊.
求你指示僕人.
耶穌說:掃羅必下來.
奇妙啊,這就是他來求問神.
求問神會不會這樣發生.
原來大衛有了預算.
預算不是因為他自己的聰明.
和他記性.

$^{121}$人家看不到的我看到的.
他完全沒有掛這些口.
相反他只是來承認.
我們需要求問上帝.
12節,13節請讀.
大衛又說:基爾拉人將我和跟隨我的人.
交在掃羅手裡不交.
耀華說:必交出來.
所以大衛和跟隨他的約有六百人.
就起身出了基爾拉.
往他們所能往的地方去.
有人告訴掃羅大衛離開基爾拉逃走.
於是掃羅不出來了.
所以這個第一個避開了一個危難.
所以這個就是大衛他面對掃羅不斷的威脅他.
他只是依靠自己的才幹聰明技能.
你知道大衛是誰嗎.
大衛是戰士來的.
有很多時候都是面對這些戰場的威脅攪拌.
打仗的時候有很多這些戰術.
所以大衛其實是一個很能夠知道應對.
他做到戰士長,指揮那些戰士的.
怎樣佈陣,怎樣迎敵.
但是我們發覺他卻是來專心求問上帝.
他拿了一個以弗德出來,叫祭司.
以弗德就是當時代表去求神指引的一個代表.
他很鄭重,不是隨便說一下.
然後叫祭司拿一個以弗德出來.
表明一個很官方,一個每個人都看見.
現在我們一起來求問上帝.
他將神擺在他很重要的.
在他人生當中,怎樣面對苦難.
怎樣逃避苦難,怎樣面對威脅.
怎樣能夠有出路的一個很重要的指引.
他不是依靠自己的聰明.
這裡很表明他的心,他依靠神,不依賴自己.
其實除了他之外,當時有多少人跟隨他?.
有600,上一次400,現在還多了一些.
所以這600人其實很多都是精英.
雖然有些未必有社會地位,有些未必是絕頂才幹.

$^{161}$不過這麼多人裡面,其實都有很多精英.
都有很多能幹人,很多戰士,有很多跟過他的.
或者以前跟過掃羅的.
所以我們知道這600人當中,其實都可以大家商量一下.
有什麼策略呢?有沒有人獻計呢?.
有什麼好的主意呢?.
當然我們不是說不要這些.
有時候都很重要的,我們的團隊,有時候有很多意見.
有時候有些很重要的,他這方面的專長.
給我們提供一些很好的參考.
但是大衛卻不只是依靠這些,或者不是專依靠這些.
他很明顯的就是,他不是跟這些人商量對策.
他是來叫人力以忽得出來,求問上帝.
這種表現在危機當中,禱告依靠神.
其實不單止剛纔我們讀過第9節第13節.
你再推前一點,你會發覺他打非理士人都是這樣的.
我們再看前幾節,第2節到第4節.
我們再讀這幾節,請讀.
所以大衛求問耶和華說.
我去攻打那些非理士人,可以不可以?.
耶和華對大衛說,你可以去攻打非理士人.
拯救基爾拉.
跟隨大衛的人對他說.
我們在猶大地這裡尚且懼怕.
何況往基爾拉去攻打非理士人的軍隊呢?.
大衛又求問耶和華.
耶和華回答說.
你起身下基爾拉去.
我必將非理士人交在你手裡.
你看見原來大衛這種情況.
不是他當時面對掃羅,沒有什麼法治.
我們中國人就說林雞.
信佛教的林雞抱佛腳.
有些基督徒也是這樣.
平時不祈禱,也不會求問神.
林雞就抱什麼?.
林雞就抱耶穌的腳.
我們應該時常來到耶穌的腳前.
來聽祂的話.
不過不要林雞才抱耶穌的腳.

$^{201}$大衛一貫以來都是這樣.
一貫以來他無論是打非理士人.
無論是現在面對掃羅的追逼.
你看見他都有一個明顯的表現.
在這兩節經文大家都看見.
他就求問耶和華.
第三節說那些人意見不同.
他說怎麼還去基爾拉那裡.
是更加危險非理士人更多的地方.
你現在說去那裡不就更危險.
所以第三節跟隨大衛的人覺得.
這樣不太妥.
所以大衛再一次求問耶和華.
結果耶和華很清楚指示他.
所以他給人的意見不同.
但他卻知道神的意見最好.
很多時候基督徒求問神.
有時會看神有沒有為我開路.
有沒有開門.
這些都是的.
因為這些是助證.
即是引證神的旨意.
不開門就不是了.
但他不是第一個最重要的關鍵.
第一個關鍵就是跟上帝的關係.
聖靈給他的指引.
在聖經裡面清楚的真理.
這個成為他考慮是否神的旨意.
一個很重要的因素.
不只是看環境.
如果是這樣.
個個不是信耶穌的當然是這樣.
那份工作請我.
我這個是神的旨意.
是不是?不是一定的.
如果是這樣.
個個不信耶穌的那份工作請我.
所以這個就是神的旨意.
就是好了.
現在每個人都這樣想.

$^{241}$而基督徒就不是.
基督徒多了.
我們先跟神有一個關係.
平時有好的操練.
平時都學習如何信服神.
求問祂的旨意.
所以在重要關頭的時候.
我們就知道如何求問神.
如何回覆.
當然求問神有時候都需要等候.
所以大約都求了兩次.
但無論如何我們知道.
這個就是祂的心態.
所以在《箴言》裡面很熟悉的經文.
都這樣講.
我們很熟悉的.
一起.
懂不懂背?.
這個應該是一個金句.
《箴言》三章五至六節.
請讀.
你要專心仰賴耶和華.
不可依靠自己的聰明.
在你一切所行的事上.
都要認定祂.
祂必指引你的路.
我們看到這個就是.
《箴言》裡面一個清楚的教導和應許.
這支經文是很重要的.
所以我們專心仰賴耶和華.
不要依靠自己的聰明.
聰明從神而來的.
經驗從神而來的.
才智都從神而來的.
不過你不要只依靠這些.
因為有些人很聰明.
有時候我們發覺他.
聰明反被什麼?.
反被聰明誤.
有很多人以為自己很聰明.

$^{281}$最後反被聰明誤.
我們最近的春令福音.
退休海上游.
積極安排.
排戲.
今年我們的主題.
我希望大家上春令會海上游這三天.
要全本《儀師帖記》讀一次.
《儀師帖記》有十章.
全本《儀師帖記》讀一次.
《儀師帖記》我們有很多方法一起學習.
有講道.
當然有小組討論.
但我們有些話劇表演.
其中一個人叫哈曼.
哈曼就是這樣.
他自以為聰明.
最後發覺他是聰明反被什麼?.
反被聰明誤.
相反.
當一個人專心依賴耶和華.
他在一切事上認定主.
主指引他的道路.
這是一個謙卑的心.
第二個是聰明的做法.
聰明的做法.
因為我們人有限.
怎麼可以只靠自己的聰明.
所以依靠神認定他.
有一個見證在2003年.
爆發亞洲沙士.
大家還記得這個事.
沙士.
當時這位是臺大醫學院內科主診醫生.
名叫方啟泰.
是一位很好的基督徒.
他後來受訪問時說.
當時他在臺大醫院服務了十年.
但是從來沒有面對過沙士傳染病的恐怖.
也是一個很大的考驗.

$^{321}$令他們束手無策.
當時臺大醫院的第一個反應就是.
立即封鎖.
有一個專門的區域是給沙士病人的.
其他就隔離.
防止他繼續傳染.
經過疫情新冠肺炎就知道那種情況.
可以想象得到.
但是這種情況在當時臺灣也是不斷髮生.
很多病人湧入.
很快特區已經爆滿.
從30張床位.
好幾天後變成70張床位.
超過一倍.
所以我們就見到這種很難處理的情況.
當時因為緊急情況也沒有很多人手.
他們只派一位主治醫生.
和七位住院醫生跟他一起.
全部自願的.
因為當時知道可能有生命危險.
誰願意去做.
所以他面對這個挑戰的時候.
真是心力交瘁.
好像沒有出路.
但是很奇妙.
在他大學三年級信耶穌之後.
有一句經文時常在他人生的座右銘.
他不斷不斷地在學習.
這句經文是什麼呢.
就是剛纔所讀的.
《箴言三章》第五至六節的經文.
他說不要專心依賴耶和華.
不可以依靠自己的聰明.
所以我們發覺當時他面對這個挑戰的時候.
他就來求問神.
神很清楚地告訴他.
這個就是我委派你去見證我的時候.
所以他就憑著這個禱告所得到的指引.
最後來到去挺身而出.
帶領其他七位醫師面對這個嚴峻的考驗.

$^{361}$結果經歷了這個之後.
當人訪問的時候.
他就將這個經歷講出來.
成為一個見證神的機會.
也都是來到看見這位基督徒醫生.
他依靠神所得到一個美好的結果.
所以我們看見這個也是一個非常值得學習.
讓我們一生都繼續操練這個.
不要只是依靠自己的聰明.
要專心依賴耶和華.
第二個我們見到在這個經文裡面.
大衛能夠找到出路.
也是接受有人的援助.
第十五至第十八節.
我們一起讀讀這四節的經文.
大衛知道掃羅出來.
索尋他的命.
那時他住在西忽曠野的樹林裡.
掃羅的遺者約拿丹起身.
往那樹林裡去見大衛.
使他依靠上帝得以堅固.
這個就是他的好朋友約拿丹.
十七到十八節繼續讀.
對他說.
不要懼怕.
我扶掃羅的手.
必不加害於你.
你必作以色列的王.
我也作你的宰相.
這是我扶掃羅之道了.
於是二人在耶和華面前立約.
大衛仍住在樹林裡.
約拿丹回家去了.
在這裡我們也見到大衛有出路.
是因為他接受了這位友人的雙援助.
他援助的寶貴地方就是他得到鼓勵.
在困境當中透過人的鼓勵代替.
很重要的.
在這裡我們見到約拿丹所做的.
就是他不單止讀到大衛的心.

$^{401}$安慰他 給他肯定.
更重要的就是他來幫助他.
倚靠上帝得以堅固.
這是他經歷到神對他的幫助.
透過一位屬靈的知己.
這位基督徒的心理學家.
叫做Paul Turnier.
Paul Turnier是一位知名的基督徒輔導家.
有人問他.
你怎麼去幫助別人.
怎麼做一個基督徒的心理學家輔導家.
他的答案就是.
其實有時候我也不知道怎麼去幫助他們.
這些人來到我面前的時候.
雖然他有很多專業知識.
有很多經驗.
他說其實有時候我覺得自己好像沒有做些什麼.
最重要的就是那個人.
第一 他願不願意找出路.
那個人願意找出路的時候.
他就有救了.
第二 其實有時候我經歷到.
我儘量去明白和支持他.
和表示他歡迎.
他願意說出他心裡面的困難.
所以他說最重要的就是.
我成為他一個好的朋友.
幫他守祕密.
和不要去定他的罪.
所以這個就是他.
一個很有幽默感的基督徒輔導家.
他還說了一堆說話.
這個就是這麼說的.
I am convinced that 9 out of every 10 persons seeing a psychologist do not need one.
對不起 我們當中有些psychologist.
不要得罪了你.
他說10個裡面9個其實都不需要.
They need someone who loves them with God's love.
and they will get well.
所以你們每個人都可以做.

$^{441}$只要你成為一個聆聽別人.
和支持別人.
更加重要的就是用神的愛去愛他.
所以這個是非常.
你看他的樣子是很風趣的.
他面對很多很多人的一些困難.
生命的危險 絕路的時候.
他就只是用上帝的愛.
和一個很願意和他們一起走過這條道路的心情.
來鼓勵他們.
所以神給我們一個很寶貴的身份.
我們都可以使人依靠上帝得著堅固.
所以這個就是哥林多後書保羅的心聲.
我們一起讀經文好嗎.
保羅這麼說.
哥林多後書七章六字.
但那安慰喪棄之人的神.
藉著提多來安慰了我們.
所以保羅是一個大仕徒.
真是身經不知多少的屬靈增賜.
當然去到哥林多後書的時候.
他都沒去到後期.
不過已經面對了他第一次宣教旅程.
和第二次宣教旅程.
和各種的困難.
甚至被人扔頭以為他死了.
但那安慰喪棄之人的神.
藉著提多來安慰了我們.
所以真的有時候我們都可以做那個提多.
原來我們都可以成為別人的安慰.
而事實上最剛強的人.
都有時候需要鼓勵和安慰.
就好像保羅一樣.
Amen.
所以求主幫助我們.
我們也想起這幅圖畫.
當時有在書中記載.
其中比如馬克福第二章.
一開始就記載馬可福音.
就說有幾個人抬著一個毯子.

$^{481}$沒辦法進去.
沒什麼出路.
結果他們在上面拆開了瓦頂.
將這個人吊下去.
原來真的有時候得到神的幫助.
真的可能是透過其他人來扶持我們.
剛纔緩道訓練班的同學們.
除了他們自己的成長.
感受到的是大家彼此激勵.
互相分享.
幸福小組好的地方也在這裡.
它不只是個人報道.
自己一個出去傳福音.
固然這個我們要裝備好的.
已經裝備了.
幸福小組好的地方就是一起做.
所以就是大家彼此鼓勵.
大家一起邀請人.
大家一起為這些未信主的人禱告.
所以這個為什麼我們成為了.
我們希望經常可以透過這個方式來傳福音.
就是我們有這個優點.
當然它也有它的缺點.
但是它的優點很重要在這裡.
就是大家有一種互相鼓勵支持.
以至能夠可以一起去完成神給我們的使命.
今天很可能我們都需要這些.
你成為別人的支持.
別人成為你的支持.
所以如果你還沒有一個很好的小組團契的生活.
就鼓勵你.
如果你從來沒有參加過緩道訓練.
可能在這個時候去考慮.
好好地裝備自己.
第三我們看到在大衛找到出路的時候.
最重要就是今天我們所體會到的.
就是投靠神化險為夷.
哎呀.
化險為夷.
這句話真好.

$^{521}$有沒有危險?.
有.
人生有沒有危險?.
有沒有危機?.
但是能夠化險為夷.
這就是上帝給我們的祝福.
二十五節我們一起讀這幾節.
請讀.
蘇羅和跟隨他的人去尋找大衛.
有人告訴大衛他就下到盤石.
住在馬雲的曠野.
蘇羅聽見便在馬雲的曠野追趕大衛.
蘇羅在這山這邊走.
大衛和跟隨他的人在山那邊走.
大衛急忙躲避蘇羅.
因為蘇羅和跟隨他的人.
四面環圍著大衛和跟隨他的人.
要拿獲他們.
這就沒有出路了.
二十七節.
二十八節請讀.
忽有善者來報告蘇羅.
非理士人犯警掠掠.
請王快快回去.
於是蘇羅不追趕大衛.
回去攻打非理士人.
哪有這麼巧的.
明明就快捉到他.
千軍一髮之間.
沒有路可逃.
真的沒有路可逃.
出路是哪裡來的.
忽然之間而來.
忽然之間.
不知道是誰殺出一個使者.
來告訴蘇羅.
現在非理士人在搶我們的地方.
你不要再在馬雲旁野那裡,快點回去.
所以全部班師.
班師護士一齊回去.

$^{561}$攻打以色列人.
非理士人.
他本來在攻打以色列人.
自己在打自己.
在打大衛.
現在就回去打非理士人.
可圈可點就有一個字.
這是什麼字.
是忽字.
忽然.
哪有忽然的.
原來這就是神奇妙的開路.
有沒有經歷過.
困境的時候忽然間神開路.
在我們人生當中都試過.
不是整天,都試過一兩次.
真的困境當中沒有出路.
忽然間光明就來了.
想起來又害怕又傷心.
但又真的可以用眼淚出來讚美神.
背後的故事在這裡.
詩篇54篇.
他自己講的.
所以這個詩題是.
遭敵迫害求主保護.
他解釋為什麼會寫這篇詩.
西忽人來對掃羅說.
原來這個人是西忽人.
來到掃羅說.
大衛豈不是在我們那裡藏身嗎.
那時大衛作者奮不顧身.
交與令長.
用詩源的樂器.
裡面說什麼,我們一起讀這幾句經文.
上帝呀,求禰以禰的名救我.
憑禰的大能為我伸冤.
上帝呀,求禰聽我的禱告.
留心聽我口中的言語.
因為外人起來攻擊我.
強暴人尋索我的命.

$^{601}$他們眼中沒有上帝.
上帝是幫助我的.
是扶持我命的.
請問細拉是什麼來的.
細拉就是寫詩的時候停一停.
就是說他在那時就默想.
上面那幾句.
再深刻一些思想的這幾句話.
他說因為外人起來攻擊我.
強暴人尋索我的命.
他們眼中沒有上帝.
他說到這裡又生氣又生氣.
又惋惜又覺得很慘.
所以他停下來.
接著第四節.
上帝是幫助我的.
是扶持我命的.
所以他開始讚美上帝.
這就是他的出路.
弟兄姊妹.
人有時也會成為我們的幫助.
不過人的瞭解真的很有限.
因為人只能看到表面的東西.
他看不到事情背後的問題.
所以我們有時真的很灰心.
但是奇妙的是上帝知道一切.
上帝明白.
所以我們就可以得著堅固.
昨天我們幸福小組也和Bess聊天.
我也和她說.
我說我們只不過是一個橋樑.
帶你去認識上帝.
不過你最重要的是你自己來認識上帝.
和讓他成為你的救主.
我們有時只不過幫得有限.
我們只可以做一個橋樑.
但真正幫助的就是上帝.
所以這就是神給我們的應許.
在拿紅書第一章第七節.
我們一起讀這一節.

$^{641}$耶和華本為善,在患難的日子為人的保障,.
並且認得那些投靠他的人..
耶和華本為善,在患難的日子為人的保障,.
並且認得那些投靠他的人..
人未必認得,上帝卻認得,這就是最重要的.
有一次在一間製衣廠裡面.
主管提醒員工每次我們這些機器運作的時候.
都有個牌寫著.
如果在操作的時候發生有些線卡住.
那個機器壞了或者令到線亂七八糟.
請你找工頭求助.
結果有一個新來的員工,他來了不久.
真的有一次做著做著的時候發現.
那些線卡都亂七八糟,滾了一堆.
他就用自己的方法解來解去,解來解去,解了很久都解不開.
所以越解就越亂.
浪費了很多時間,最後他非常失望.
終於看到那個牌說什麼.
請向工頭求救.
所以他就請了工頭來,很快就知道原因幫他搞定.
他就辯護,工頭問他為什麼不早些找他.
他就辯護,他說因為我已經用了最好的方法.
用了最好的方法來解決.
工頭就微微笑.
他說記得最好的方法是什麼.
最好的方法就是來求救,找我幫你.
今天我們真的看見人有時候越搞越亂.
唯有投靠神的人是我們人生的出路.
無論我們生活也好,靈性也好,困境也好.
我們都走著這個應許,以和華本為善.
在患難的日子為人的報償.
並且認得那些投靠他的人.
大衛得到出路,因為他不靠自己的聰明.
因為他接受有人互相的援助.
最後很重要,就是他投靠神,化險為夷.
這也是我們人生在絕路,困境當中真正的出路.
我們低頭祈禱,請同心投降.
我們再一次投靠在主的任旁.
在禰的翅膀底下得著禰的任悲.
因為當我們投靠禰的時候,禰認得我們.

$^{681}$幫助我們能夠學習大衛的智慧.
也能夠學習大衛的謙卑.
讓我們能夠經歷在我們人生當中.
在困境的時候得著出路.
化險為夷,奉耶穌基督的名.
\newpage



\section{約書亞記 14:6-12}
\label{sec:031859}
\textbf{活出使命的跟從者——迦勒 II (華)}~~~~梁明財牧師
\newline
\newline
連結: \href{http://www.church.com.hk/acms/content.asp?site=cdc\&op=show\&type=product\&code=031859\&layout=sermon}{\texttt{031859}} ~~~~ 語音日期: 2026-03-22
\newline
\newline
\hyperref[sec:031858]{\small{< < < PREV SERMON < < <}}
~
\hyperref[sec:index_chronic]{\small{[返順時目]}}
~
\hyperref[sec:031860]{\small{> > > NEXT SERMON > > >}}
\newline
\newline
約書亞記 14:6-12
\newline
\begin{longtable}{cl}
\hline
\hline
章節 & 經文 (和合本修訂版)\\
\hline
14:6 & \begin{tabularx}{0.7\textwidth}{X} 猶大人來到吉甲，約書亞那裡，基尼洗族耶孚尼的兒子迦勒對約書亞說：「耶和華在加低斯‧巴尼亞指著我和你對神人摩西所說的話，你都知道。 \end{tabularx} \\ \\ \relax
14:7 & \begin{tabularx}{0.7\textwidth}{X} 耶和華的僕人摩西從加低斯‧巴尼亞差派我窺探這地的時候，我剛四十歲。我把心裡的話向他報告。 \end{tabularx} \\ \\ \relax
14:8 & \begin{tabularx}{0.7\textwidth}{X} 雖然同我上去的眾弟兄使百姓膽戰心驚，我仍然專心跟從耶和華－我的神。 \end{tabularx} \\ \\ \relax
14:9 & \begin{tabularx}{0.7\textwidth}{X} 那日，摩西起誓說：『你腳所踏之地必要歸你和你的子孫永遠為業，因為你專心跟從耶和華－我的神。』 \end{tabularx} \\ \\ \relax
14:10 & \begin{tabularx}{0.7\textwidth}{X} 現在，看哪，耶和華照他所說的使我活了這四十五年。當以色列人在曠野飄流的時候，耶和華曾對摩西說了這話。現在，看哪，我已經八十五歲了。 \end{tabularx} \\ \\ \relax
14:11 & \begin{tabularx}{0.7\textwidth}{X} 現今我還很健壯，像摩西差派我去的那天一樣；無論是戰爭，是出入，我現在的力量和那時的力量一樣。 \end{tabularx} \\ \\ \relax
14:12 & \begin{tabularx}{0.7\textwidth}{X} 請你將耶和華那日所說的這山區給我。那日你也曾聽說，這裡有亞衲族人，以及寬大堅固的城，或許耶和華會照他所說的與我同在，我就把他們趕出去。」 \end{tabularx} \\ \\
[1ex]
\hline
\hline
\end{longtable}
$^{1}$我今天繼續會講的經文跟那個主題就是我在二月份的時候跟各位有講過的那個活出使命的跟從者扎勒.
我不曉得大概如果你上次有來的話應該是那個二月八號的時候那一場的講到.
我們今天會繼續來看扎勒的生命,他是怎麼樣在敬畏上主的人生當中去活出神給他的使命..
而且他是一個讓我們可以學會怎麼樣去跟從我們的神,跟從耶穌,所以這個是從扎勒身上可以學到的..
如果上次二月份的時候弟兄姊妹還沒來到這個崇拜講的時候,你可以上網去聽那個第一講那個講到活出使命的跟從者的扎勒..
在這個開始的時候,我還是先跟各位來讀一遍這段的經文,其實也不太長,希望我們還多一些印象在這個經文當中..
在約書亞記十四章的六道設計,如果看到的話我邀請弟兄姊妹你們可以一起我們來讀這段的經文,看到的我們來,預備來..
那時猶大人來到吉宅見約書亞,有機立洗主耶穌尼的兒子扎勒對約書亞說:.
"耶和華在扎底斯巴里亞,指著我與你對神人摩西所說的話,都知道了..
耶和華的僕人摩西從扎底斯巴里亞打發我窺探這地,那時我正四十歲,我按著信儀回報他,然而從我上去的重地收拾百姓的心騷化,專心跟從耶和華我的神.".
當日摩西指示說,你走所踏之地,定要歸你和你的子孫永遠為業,因為你專心跟從耶和華我的神..
自從耶和華對摩西說這話的時候,耶和華叫他所應許的使我存活這四十五年,連以色列人在曠野行走,看啊,現今我八十五歲了,我還是強壯,像摩西打發我去的那一天一樣,無論是征戰,是出入伊崗那時如何,現在還是如何..
求你將耶和華那日應許我的這三地給我,那裡有雅納族人,並誇大堅固的城,你也曾聽見了,或者耶和華叫他所應許的與我同走,我就和他們趕出去..
只是他的重點有點不同,如果我們看約書亞記的時候,我們發現約書亞是一個比較我們多去認識他是一個領受神給他召命的那一位,而迦勒就是一個回應召命的那一位,當然他們都同樣領受了上帝的召命,同樣也是回應,只是回應的重點有不一樣,迦勒他怎麼樣回應上帝的召命,回應神的召命呢?.
就是他跟從上主的門,跟從上主的門,這個是什麼意思呢?如果你剛纔有留意你讀的那一段經文裡面,比較常出現的一句話是怎麼樣,就是專心跟從耶和華,在短短的幾節的經文當中都出現了好幾次,這就是怎麼樣去形容迦勒是一個怎麼樣活出神給他使命的一個人..
這個在我們通常用的聖經和文的翻譯是專心跟從耶和華,但在另外一個我跟各位看過的是呂正宗譯文當中講到是滿心滿懷跟從永恆主,這是他一個我覺得是比較更貼近原來聖經希伯來文的意思..
它的英文就是fill up after Jehovah,跟從耶和華神的門.重點是怎麼樣呢?就是他有沒有對準他的跟從的對象,而且因為他能夠對準跟從的對象,我們的耶和華神的時候,他所得到的這種生命當中的門..
這個門是什麼意思呢?這個包括了我們上期有講到神對迦勒的肯定.剛纔我們讀到從四十到八十五歲,我們在座當中有哪一位是八十五歲的?超過八十五歲,我們大概都很年輕,大家還沒到那個年紀,四十歲還沒到了..
聖經的描述就是說,他對神是始終如一,他沒有改變,他對準的是他的神.這也提醒我們,在我們人生當中每一個階段,神都知道我們,無論你是四十歲,也不一定是四十歲,從你出生開始,神都知道,神都認識我們..
在我們人生的每一個階段,神都知道,神也肯定我們每一個人的人生,無論我們要過一個怎麼樣的人生..
這個是迦勒他在經歷到,他繼續去,無論是專心跟從耶和華,或許是他跟從神的滿.迦勒擁有的滿是一種怎麼樣呢?我把它大概這個比較勉強的翻譯,可能是滿心滿懷,可能是很滿意,大概我們能夠抓得比較近一點的滿足..
當迦勒他願意真的是跟從耶和華神的時候,他對準耶和華神的時候,他是為神而活的時候,他是為神去打仗的時候,他就是經歷到那種在生命當中的滿足..
這就是在我們上期有講到的.上期我也講到,剛好準備過年,我們講到一些新年的願望.我不曉得大家有沒有看過新年你會貼春聯和輝春,我們都有..
長門是其中一種我在小時候在家裡面每年都要貼的.你知道貼在哪裡嗎?貼在米缸外面,每年會換一張新的.我們那個時代生活也比較不容易,我們過年前就希望明年的米缸不會缺米,希望米缸都是長門..
所以這個在我們小時候在過年的時候是會常有的,除了福字或者說其他的一個輝春,春聯那些祝福語,長門是我們的..
所以我覺得這個也很有意思.我們新年的願望會不會是,你是有一個跟從神的滿.這個滿就是看我們是聖在什麼.這個米缸裡面是裝滿米..
但如果在我們的人生,在我們的生命當中,你聖在的滿是一種什麼滿?對扎勒來說,在約書亞記所記載的,它是滿有神的音許,滿有神跟人之間的約.這就是扎勒他所充滿的那種生命在他裡面的滿,是神的音許,是神人之約..
另外,我跟各位也看到那個在大衛,大衛他也是用專心順服耶和華,其實專心跟從耶和華的意思是一樣的,但是大衛是另外一個在聖經中去形容他是一個跟從耶和華神的滿,那個充滿的人..
所以我想跟各位從經文來問我們自己一下,到底我們是不是一個很滿的人?這個滿大概是用滿足就好了..
你是不是一個很滿足的人?還是當你想到我很滿足,但你的滿是一種什麼的滿?是家裡面的米缸的滿,還是你銀行的戶頭那幾百萬幾千萬的滿,還是什麼什麼的滿在你的人生當中?你聖在一些什麼樣的滿在你的生命當中,在你的生活裡面?.
扎勒給我們的一個就是他有神的應許,他是滿在神的應許,所以他可以繼續在從四十歲到八十五歲,縱然他跟以色列人在曠野四十年,他仍然滿有神給他的應許,他沒有忘記,他更不會不相信神的應許..
所以,這就是在信靠神的應許那種滿足在扎勒這個人的生命當中..
我不曉得對你來說,你是不是一個已經很滿足的人.當然,講到滿的時候,可以是滿足的那個比較正面一點的形容..
另外一種滿是什麼?我叫它做自滿,或者說心臟不滿.如果你覺得你是不太滿足的人,我不曉得你會覺得有時候你跟人比較,最容易能夠覺得自己很不滿足的話,大概是跟別人比較..
你會覺得跟你差不多年紀,差不多學歷,或者說你讀的書也差不多的人,你看到有些人比你好像他在職場上,工作上甚至是很壞,他賺很多錢,而你跟他比較的話,你就覺得為什麼他可以有那麼豐盛的錢那麼多,我就好像不會能夠有那麼多的滿足..
所以,有時候我們每個人都有一些不滿.尤其當我們年紀越來越年紀比較大的時候,我不曉得你有跟那些比較年紀大,或者說長者他們來相處的時候,你會發現有時候有些長者在他們的人生雖然是過了好像是七八十歲的時候,像加了到八十五歲的時候,.
有可能一些長者,老人家他們對他們的過去的人生也有一些不太滿意的時候,或者不滿的時候,有時候會覺得他們年紀見大的時候,他們可能是也失去的東西越來越多,比如說他失去了,沒有再失去他的青春了,失去他身體的健康了,可能他的戶頭裡面的錢越來越少..
這有可能就是因為失去讓人感到不滿,可能不是隻有長者老人家會有這樣子,有可能你覺得你的人生就好像有些失去的時候你就會覺得不滿..
但當你不滿的時候,一個人的反應是怎麼樣的?如果你失去的時候,你最想是什麼?去抓住,你是很想抓住很多的東西,比如說你越來越像我這把年紀,我在香港已經拿到老人卡,現在已經搭公車,已經兩塊錢就可以,但四月份就開始不行了,我們都已經有那種長者的福利..
但其實我們當我們覺得自己其實是越來越失去一些東西的時候,我們就很想抓住一些東西,我們的不滿就是希望抓住一些東西來填滿我的失去..

$^{41}$所以有時候這就是我們在人生年紀越大的時候,或者說我們好像你覺得你自己失去很多的時候,你就很想抓住一些東西,不管你是想抓什麼,多抓一點錢,多希望身體健康一點,或者多抓這些朋友在你身邊,我們都很想抓住一些..
但我很想跟各位來分享,在我們不滿當中,在我們失去的時候,我們應該抓住什麼.另外,我不曉得你是不是一個很滿的人,有些人他覺得他不會失去,他沒有什麼損失,他覺得很有那種滿足的感覺..
這種我在形容就是一個自滿的人.我不曉得你會不會也是有點像這種自滿的人,或許你是一種不拒絕的滿,你會覺得我已經很夠了,我足夠了,我已經很滿足,很滿意了,我就什麼都要,或許我的信心已經很足夠了..
我們有一句話形容一些人,他們很自滿的人,就是他的信心是爆棚的,他信心很多,很大..
通常這一年多我看新聞,就常覺得有一個國際人物,他是一個信心爆棚的人.每次看到電影,看到他出來講話的,你就覺得他講一句話,他覺得什麼都是沒問題的.這個你知道是誰嗎?.
你看到川普出來的時候,你就看到他是一個非常自滿的人.我不曉得他是不拒絕的滿,還是他本身就是這樣子,覺得我什麼都行的人..
所以,這個是我覺得最害怕的滿,一個自滿的人,不拒絕的滿,他以為自己是信心爆棚,所以他所做的一切都沒問題,一切都是最好的.這種人其實就很危險..
所以,不管你是一些不滿的人,還是你是一個很自滿的人,其實我們都要小心去反省我們的滿是甚在什麼樣的滿.當我們真的覺得很想抓很多的東西,一個很自滿的人他其實也抓很多的東西,就是去希望得到更多的東西..
我想跟各位看,我們跟窮神的滿是怎麼樣的.這個圖畫是這樣子,是看到什麼?你看到.有看到嗎?這個圖畫是一個小孩,他抓住一個大人的手..
常常是一個小孩,上這個了,他抓住他爸的手,這個小女孩.常有,你會抓住你爸的手嗎?會.常常小孩都抓住,特別是去逛街的時候,他一定抓住大人的手,他的父母的手..
所以其實當我們失去的時候,我們好像不滿的時候,或許你是一個自滿的人,也很想抓很多東西.但最重要,我想跟各位說,跟窮神的滿,其中一個也可以這樣子看,就是我們有沒有抓住神?我們跟窮神的滿,就是我有沒有抓住神的音訊?我信靠神的音訊,這個纔是跟窮神的滿..
這個纔是我們需要真正的滿足.這是抓了跟窮神的滿.我上一次也講到,大衛也是同樣有抓住神給他音訊的滿..
在歷王紀上,我們上一期有提到,是大衛跟他的兒子所羅門的一個比較.大衛就是專心順從耶和華,其實也是專心跟從耶和華,跟從神的滿的那一個人,是因為神的音訊,是因為他對神的信心..
所以,跟窮神的滿,其中一個是我們抓住什麼?你拼命在失去的時候,在你覺得不滿的時候,或許在你自滿的時候,你抓的東西是不是最好?是不是真的對你好?我們要抓住的是神的手..
所以,當我們抓住,拖住神的手的時候,那份滿足是最好的.我不曉得在我們小時候,可能我們那個年代,我們很少記憶了,當你小時候,你會抓住你爸媽的手是最滿足的..
其實,更大的感覺是什麼?是最安全.當你抓住你的父母也好,你跟大人出去逛街的時候,你抓住他們的手,這是最滿足,最安全,最放心的時候..
所以,這個就是我們每一個教徒,像迦勒一樣,他就是抓住神的音訊.他是不再懼怕,他是無懼地去跟從,順服,得著神給他的音訊,得著神給他的力量..
這就是迦勒的門.他的自門不是信自己,不是信靠自己,他的門是信靠神,抓住神.所以我上次有提到,當他們在還沒有見到迦南帝之前,剛出埃及的時候,他們就成為十二個探子,出去看迦南音訊帝是一個怎麼樣的地方..
然後十二個探子回來彙報的時候,我們就說迦勒跟約書亞是贊成的.十二個探子都說那個音訊帝是很好的地方,但回來的時候只有兩個,就是迦勒跟約書亞說這個好地方是神音訊給我們的,他們可以上去攻打這個地方,得到這個音訊帝..
但有十個探子說那邊的人太厲害了,太強壯了,每個人都比他們高大,每個人都很會打仗,所以他們說要不要去?不要去..
所以十二個探子有點像我們今天常出現的場景,要投票決定,十個反對,兩個贊成,結果應該怎麼樣?如果從投票的結果就說,那我們要不要去?不要去嘛,十個人反對,兩個人贊成,他們都是十二支派的代表嘛,那就代表有十個支派反對,只有兩個支派贊成去..
那怎麼樣打了?但這個是在這種環境的時候,迦勒他那種滿足,他那種堅持,他無懼地堅持繼續跟從神,他是為了神繼續去願意去領受神給他的音訊,所以這就是迦勒的滿,迦勒的自滿,這是一個很正確的滿..
所以迦勒的回應是怎麼樣的?約束亞記第十四章講到,他說是因為耶和華照他所應許的,他不是靠人的能力,不是靠迦勒他那個支派,可能說迦勒的支派他的人強馬壯,他有很大的軍事力量,他就可以去上去打那個應許地的人..
他不是這樣,他是靠的是耶和華會照他的應許,讓迦勒,讓以色列人可以得到那個應許地,所以這個是迦勒他在神面前願意去跟從神的意思,是對準神..
他在這裡有講到,可能有些看不到,他在最後兩句是怎樣的呢?他說"或者耶和華照他所應許的,與我同在,我就把他們趕出去.".
其實這句話不是說迦勒沒有信心,他用了"或者耶和華照他應許的".其實這個應該是讓我們多明白一點,"或者"的意思不是在中文那個意思,"或者"原來的意思是說,只要神照他應許,讓我跟我同在,我就可以去得到那個地方..
所以他這不是說,好像一個不肯定,不是一個沒有信心的表達,或者那個翻譯在中文應該翻得比較好的,是隻要神幫助就可以,就足夠了..
這不是我自己個人的能力,不是我這個支派,或許迦勒跟耶穌這兩個支派可以做到的,是隻要神幫助我,神按照他的應許給我能力..
所以他這裡有一句是我們在教會當中常用的,就是說"我的日子如何,我的力量也如何.".
其實這個"日子如何"的意思更是要提醒我們說,只要我的神幫助我,只要我得到神的幫助,只要我抓住神的應許就可以了,我抓住神的應許就可以了..
所以這是迦勒他堅持,跟從耶和華,專他為神,這就是迦勒的滿足..
所以我盼望今天我們想到的時候,有時候我們今天這個世代,其實也是在專心去跟從一些東西,我不曉得在這個世代裡面,你是專心跟從一些什麼東西,你以什麼為滿足?.
或許說你跟從就是代表你追求一些什麼為滿足了,在AI世代,或者說在我們用網絡世界的年代的時候,追求或者說跟從就是follow,after,follow..
所以這個世代的時候是很多人很想要更多的followers,這個followers就是網絡世界裡面,我不知道在座當中你的未來人生的願望會不會就是這樣子,你需要很多的followers..
有一次我看電視訪問一些小孩子.我們通常都問小孩子,你的志願是什麼?你的願望是什麼?你知道現在的小孩子,我的志願是什麼?是去做KOL,做網紅..
現在的並不是想去做醫生,律師,做生意人.他們訪問那些小孩子說,我現在已經很想當未來的我的工作,我的志願了.但那些是什麼?念小學的小孩而已..
他們就說,我們現在希望當網紅,他們已經在做了.我告訴你,我相信你們如果喜歡上網的習慣的話,你就知道他們其實就是已經在做,他們希望從小就成為網絡上的一個網紅..
網紅最重要是什麼?followers,很多人的跟從他.最重要是subscribers,很多人能夠訂閱他的網絡.所以這就是現代人,他們需要更多的followers來成為我的滿足..
這有可能就是未來,現在已經是這樣子了.但當我們想到我們以跟從什麼為滿足的時候,我們應該想到我們不是跟什麼的KOL,我們更重要的是跟從我們的神,以神為我們的滿足..
所以,反而是我們可以去問一個問題,我是誰的followers?到底今天你跟從什麼?當我們作為基督徒的時候,今天我們要跟從的不是你跟從哪一個牧師,你跟從哪一個傳道人,你不是跟從你所認識的一些基督徒,我們最重要的是要跟從我們的神..

$^{81}$我們是以跟從神為我們的滿足.當人在追求跟從的時候,我們不是追求什麼的財富,我們也不是追求什麼要很出名,要做一個什麼最紅的網紅,也不是要追求一些權利為我們的滿足..
這就是我們在從扎勒的身上,他不是這樣子.他以跟從,以追求跟從神,抓住神的應許為他的滿足,去見證他跟從神的滿足..
所以當約束亞在約束亞第十四章記載扎勒的講話的時候,他不是說他很厲害,他不是說他從四十歲到八十五歲,他仍然是跟隨神,仍然可以打仗,日子如何,力量也如何..
他不是講他自己很厲害的那一面,他只是跟人家來分享,他是跟隨耶和華神的那個小小的一個人.因為神厲害,他才能夠去打仗,去得到神應許給他的應許地..
所以在這個約束亞經第十四章裡面去講到,扎勒的見證是神的厲害,不是人的厲害.他願意的是因為他追求神,是跟從神,他是勝在神給他的應許,這就是他的滿足..
這個大衛的經歷,我們上一期有提到,大衛曾經他的滿是什麼?在詩篇五十一篇裡面講到,大衛他過去的滿是滿在他的罪惡,他所犯的罪,他得罪神,他殺人,他犯了姦淫的罪..
所以這些大衛曾經有個滿,但後來他願意悔改,他得到的滿是什麼?在詩篇的五十一篇裡面,他勝在的是神的憐恤,神的憐憫,他勝在的是神的歡恕,他能夠新的一個開始,他勝在清潔的心,正直的靈..
所以這個勝在在我們的人生當中,在我們還沒有相信耶穌之前,我們所勝在的滿可能都是一些不太好的..
當我們認識耶穌,我們回到神的面前的時候,我們所勝在的就是神的恩典,神的憐憫,神給我們清潔的心,正直的靈,這就是我們一個新的勝在,是一個跟從耶和華神的滿..
但我想跟各位來看另外一個詩篇,詩人他的經歷,他的追求是什麼?他的跟從神的滿是什麼?我想邀請各位來讀這個詩篇,這是另外一個詩人的經歷,不是大衛的,是另外一個詩人在詩篇裡面寫的..
我覺得這個就是讓我們更深的去明白,或者說我們有可以跟從耶和華的那種滿足是什麼.如果看到的,我們一起來讀這個詩篇,73篇23到26節,看到的我們一起來讀..
預備來,然而我常與你同在,你摻著我的右手,你要以你的訓言引導我,以後一切都到榮耀你.除你以外,在天上我有誰呢?除你以外,在地上我也沒有所愛慕的..
我的肉體,我我的心腸雖殘,但神是我的心的力量,又是我的福分,直到永遠..
這個詩人的經歷,我覺得就是再進一步了,他那種跟隨耶和華的滿,是從加勒的延伸出來,我覺得這個是更讓我們去多明白一點,我們怎麼樣可以做一個活出使命的跟隨者,跟從神的人..
這個詩篇其實很有意思,它特別也提到我的肉體和我的心腸雖殘,就有點像一個老人家講解,有時候我們都覺得,我自己也覺得了,有時候可能身邊的人也會看老人家是這樣,年紀大,好像身體各方面都比較差一點..
但其實這裡提醒我們,年紀大的仍然可以有力量,仍然有那份滿足,不會因為我的身體,我的心腸,我的這個人,真的我要接受我的健康可能比較差一點,但這不代表我就沒有那種滿足..
所以這個詩最滿足是因為他經歷到神與他同在,他被神的右手抓住他,這裡就另外一個圖畫,我給各位先看一下,這個是被神抓住的滿,跟隨神的滿也可以說是一個被神抓住的滿,這個被神抓住的滿跟我們剛纔另外一個圖畫有什麼不同?.
剛纔那個圖畫是小孩子怎麼樣?抓住神的手,那這個圖畫是比較,如果你看的有點不同的是什麼?是那個大人,就像是父母,他是抓住小孩的手,這個兩個圖畫的表達不是那個小孩子去抓住父母的手,是父母,那個大人抓住小孩的手,哪一個比較安全?.
你覺得,如果小孩抓住大人的手安全,還是那個大人抓住小孩的手比較安全?我相信這個應該是這個吧,就父母的經驗應該是這樣子..
當小孩抓住你的手的時候,很容易當他有時候摔開你的時候,他就跑掉了,但當你抓住你的小孩的手的時候,他就很安全,他就不會掉下來,他可以坐在他爸的頭上,所以這個就是我想跟各位說,我們真的是要抓住神的手,抓住神的音訊,特別是在我們很多時候我們都經歷到我們需要幫助的時候,我們就拿掉神的面具去抓住他,求他的幫助..
但我很想跟各位說,從這個世人的經歷當中,我們多一點去經歷是什麼呢?要讓神的手抓住我們的手,不是我們去抓住神,是你要願意被神的手抓住你,這個纔是更好的跟從耶和華神的門,是願意被神的手抓住你的手,是由神來帶領你,被神抓住..
這裡在詩篇裡面講到,是神摻著我的右手,這是神抓住我的手帶領著我,是因為我願意以神遵守神的話語,繼續前行,讓神的話語去引導我,這纔是一個更好的跟從耶和華神的門..
然後這個最後是提醒我們說,我們得的福分不是單單在地上,我們是一個直到永遠,在永恆的福分,所以我們不是追求地上的滿足,我們更多的是要追求永恆的滿,那個永恆的滿足纔是我們應該要追求的..
所以這個是需要很深的經歷,它的滿也不一定是多,不是你要滿足這些越多越好..
我常用一個例子,就是說你跑去喫自助餐,在香港喫自助餐,是不是你喫得越多越好?我告訴你,真的是,你喫得越多越辛苦..
那個滿,如果你喫飽了那個自助餐的福分的時候,你喫到肚子都不撐著的時候,你就覺得回家就會很辛苦,你拼命去喫..
不是多就一定是最好的滿足,所以是滿不一定是多.我也提過就是說在迦南應許地的時候,其實真的是一個形容是流來與蜜的地方,有很多很豐富的土產,但其實那不是整個迦南地都是這樣子的..
迦南地大部分還是沙漠,曠野,所以不一定是所有整片土地都長滿了很豐盛的土產纔是最好的,那個多不一定是最好的,那個多不一定是真正的滿足..
所以滿不一定是你要越多越好,所以我們的滿足是要讓神的手抓住我的手,帶領我,按照神的話語,我們繼續走前面的路..
最近二月份的時候,我不知道你有沒有留意一個新聞,我覺得這是人生當中我們常會被提醒的,就是一個越有錢的人是不是越快樂,越滿足,越有幸福感..
如果你最近二月份有留意一個新聞是講誰,講到馬斯克,你知道馬斯克是誰嗎?馬斯克最近在二月份的時候,很容易在網絡上發現他,他寫了一些東西..
最近在二月五號,他寫了一句話,英文是這樣子的,Money can't buy happiness,就是說錢買不到快樂,買不到幸福感..
你知道馬斯克是一個怎麼樣的人嗎?最近他這幾年都是全球的首富,他的身家你知道有多少錢嗎?給你一點點已經夠用了..
他最近二月份的時候,那個資料是這樣子的,他是有8250個億的美元,不是港幣,也不是人民幣,他是有8250億的美元,這夠不夠多?.
很多,給你一個最後的一點點,已經夠你一生,他很多,但他滿足嗎?他自己講的,金錢不能夠給他買到快樂,幸福..
當然你會說馬斯克是個很特別的人,他真的很特別的人,但就是這樣子,我告訴你,神給我們的滿足不一定是多,最重要是我們有沒有跟從,我們有沒有對準,跟從,順服我們的神,讓神的手抓住我們,帶領我們..
所以,今天我想從扎勒的身上跟從耶和華神的滿,是讓我們更多的明白,滿不一定是多,你覺得你今天好像有點不滿,你覺得你今天可能有失去了一些東西,你想抓住一些東西,希望能夠得到更多,得到滿足..
但我告訴你,不一定是多,也就是你的滿足.我們要跟從神,是相信神給我們,這是最好的,神給我們多少,我們就要帶著感恩的心來領受,是神的允許,更是神的應許,是在神的掌管當中,所以我們放心讓神抓住我們的手,是神的掌管,是神帶領我們..
所以我們要堅持,對準,跟從神,我們的生命就滿在神的應許祝福,經歷最大的滿足..
扎勒的一生是怎麼樣呢?不要以為他是為了自己纔去打仗,也不要以為他是為了以聖列人才去打仗,其實扎勒他是為了神,他的滿足,他的跟從神的意思是,他是為了神而做..

$^{121}$他不是為自己,也不是為以聖列人,他是為了神,他是領受神給他的使命,按照神的吩咐,應許,他就去得那個應許地,他是在神的國度裡面,去跟神一起來配搭,完成神的計劃,這就是跟從神的滿足..
不是為自己,也不單單是要為以聖列人,是為了神,神有他最好的計劃在我們的身上,在教會的身上..
所以我盼望今天能夠從扎勒的身上,我們來學會跟從神的那種真正的滿足,這是從扎勒的身上,生命裡面給我們學習到的..
所以我不知道今天,你以什麼為滿足了?你覺得你的人生滿足嗎?.
我們中國人有句話是這樣子的,有句話是跟滿足很有關係的,叫做"有子就萬事足"..
我不曉得這句話對你來說有沒有什麼意思,可能我們的長輩大概都很看重這句話"有子萬事足".但我們今天,可能我們也可以換一點點,可能你說不是有子萬事足了,是什麼?.
有錢萬事足?有房萬事足?還是有什麼萬事足?我也不曉得.但我想跟各位說,在我們跟從耶和華的滿足當中,最重要的是什麼?.
有主萬事足.今天你能夠對準跟從神,跟從我們的主耶穌,你有主,我相信你就會有真正的滿足,這是從扎勒的身上,可以給我們學會的功課..
盼望我們對其在神面前,經歷到跟從神那份滿足.我們一起來禱告..
跟從神為我們最大的滿足,我們能夠抓住神的音訊,我們能經歷到神,你就加給我們的能力..
今天我們在人生當中,其實我們更多的要對準我們的神,我們是為神而活,我們能夠在神的國度裡面,去被神所帶領的時候,能夠完成神讓我們在你國度裡面的使命..
願主幫助我們,讓我們經歷到我們能夠因為跟從神,我們抓住神,被神抓住的時候,我們得到那份最大的滿足..
讓我們經歷到我們的人生真的是可以有主萬事足,事事足全聽我們的禱告,奉耶穌基督的名求,阿們..
\newpage



\section{約書亞記 10:1-12:24}
\label{sec:031860}
\textbf{有神帥領的爭戰}~~~~方兆雄牧師
\newline
\newline
連結: \href{http://www.church.com.hk/acms/content.asp?site=cdc\&op=show\&type=product\&code=031860\&layout=sermon}{\texttt{031860}} ~~~~ 語音日期: 2026-03-22
\newline
\newline
\hyperref[sec:031859]{\small{< < < PREV SERMON < < <}}
~
\hyperref[sec:index_chronic]{\small{[返順時目]}}
~
\hyperref[sec:031861]{\small{> > > NEXT SERMON > > >}}
\newline
\newline
約書亞記 10:1-12:24
\newline
\begin{longtable}{cl}
\hline
\hline
章節 & 經文 (和合本修訂版)\\
\hline
10:1 & \begin{tabularx}{0.7\textwidth}{X} 耶路撒冷王亞多尼‧洗德聽見約書亞奪了艾城，徹底毀滅，處置艾城和艾城的王像處置耶利哥和耶利哥的王一樣，又聽見基遍的居民與以色列人立了和約，住在他們中間， \end{tabularx} \\ \\ \relax
10:2 & \begin{tabularx}{0.7\textwidth}{X} 耶路撒冷人就很懼怕，因為基遍是一座大城，如京城一樣，比艾城更大，並且城內的人都是勇士。 \end{tabularx} \\ \\ \relax
10:3 & \begin{tabularx}{0.7\textwidth}{X} 耶路撒冷王亞多尼‧洗德派人去見希伯崙王何咸、耶末王毗蘭、拉吉王雅非亞和伊磯倫王底璧，說： \end{tabularx} \\ \\ \relax
10:4 & \begin{tabularx}{0.7\textwidth}{X} 「求你們上來幫助我，我們好攻打基遍，因為它與約書亞和以色列人立了和約。」 \end{tabularx} \\ \\ \relax
10:5 & \begin{tabularx}{0.7\textwidth}{X} 於是五個亞摩利王，就是耶路撒冷王、希伯崙王、耶末王、拉吉王和伊磯倫王，聯合上去，率領他們所有的軍隊，對著基遍安營，要攻打基遍。 \end{tabularx} \\ \\ \relax
10:6 & \begin{tabularx}{0.7\textwidth}{X} 基遍人就派人到吉甲的營中約書亞那裡，說：「不要袖手不顧你的僕人，求你趕快上來拯救我們，幫助我們，因為住山區亞摩利人的諸王已經聯合來攻擊我們。」 \end{tabularx} \\ \\ \relax
10:7 & \begin{tabularx}{0.7\textwidth}{X} 於是約書亞和所有跟他一起作戰的士兵，以及大能的勇士，從吉甲上去。 \end{tabularx} \\ \\ \relax
10:8 & \begin{tabularx}{0.7\textwidth}{X} 耶和華對約書亞說：「不要怕他們， 因為我已將他們交在你手裡，他們沒有一人能在你面前站立得住。」 \end{tabularx} \\ \\ \relax
10:9 & \begin{tabularx}{0.7\textwidth}{X} 約書亞就連夜從吉甲上去，猛然襲擊他們。 \end{tabularx} \\ \\ \relax
10:10 & \begin{tabularx}{0.7\textwidth}{X} 耶和華使他們在以色列人面前潰亂。約書亞在基遍大大擊殺他們，在伯‧和崙的上坡路上追趕他們，擊殺他們，直到亞西加和瑪基大。 \end{tabularx} \\ \\ \relax
10:11 & \begin{tabularx}{0.7\textwidth}{X} 他們在以色列人面前逃跑。正在伯‧和崙下坡的時候，耶和華從天上降下大冰雹在他們身上，直降到亞西加，打死他們。被冰雹打死的，比以色列人用刀殺死的還多。 \end{tabularx} \\ \\ \relax
10:12 & \begin{tabularx}{0.7\textwidth}{X} 當耶和華將亞摩利人交給以色列人的那一日，約書亞向耶和華說話，在以色列人眼前說：「太陽啊，停在基遍；月亮啊，停在亞雅崙谷。」 \end{tabularx} \\ \\ \relax
10:13 & \begin{tabularx}{0.7\textwidth}{X} 太陽就停住，月亮就止住，直到國家向敵人報仇。這事豈不是寫在《雅煞珥書》上嗎？太陽停在天空當中，沒有急速落下，約有一整天。 \end{tabularx} \\ \\ \relax
10:14 & \begin{tabularx}{0.7\textwidth}{X} 在這日以前，這日以後，耶和華聽人的聲音，沒有像這日的，這是因為耶和華為以色列作戰。 \end{tabularx} \\ \\ \relax
10:15 & \begin{tabularx}{0.7\textwidth}{X} 約書亞和跟他一起的以色列眾人回到吉甲的營中。 \end{tabularx} \\ \\ \relax
10:16 & \begin{tabularx}{0.7\textwidth}{X} 那五個王逃跑，躲在瑪基大洞裡。 \end{tabularx} \\ \\ \relax
10:17 & \begin{tabularx}{0.7\textwidth}{X} 有人告訴約書亞說：「那五個王已經找到了，都躲在瑪基大洞裡。」 \end{tabularx} \\ \\ \relax
10:18 & \begin{tabularx}{0.7\textwidth}{X} 約書亞說：「你們把幾塊大石頭滾到洞口，派人在那裡看守他們。 \end{tabularx} \\ \\ \relax
10:19 & \begin{tabularx}{0.7\textwidth}{X} 你們卻不可停留，要追趕你們的仇敵，從後面攻擊他們，不讓他們進到自己的城鎮，因為耶和華－你們的神已經把他們交在你們手裡。」 \end{tabularx} \\ \\ \relax
10:20 & \begin{tabularx}{0.7\textwidth}{X} 約書亞和以色列人徹底擊敗他們，直到把他們滅盡，只剩下少許的人逃進堅固的城。 \end{tabularx} \\ \\ \relax
10:21 & \begin{tabularx}{0.7\textwidth}{X} 眾百姓就安然回到瑪基大營中，到約書亞那裡。沒有人敢向以色列人饒舌。 \end{tabularx} \\ \\ \relax
10:22 & \begin{tabularx}{0.7\textwidth}{X} 約書亞說：「打開洞口，把那五個王從洞裡帶出來，到我這裡。」 \end{tabularx} \\ \\ \relax
10:23 & \begin{tabularx}{0.7\textwidth}{X} 眾人就這樣做，把那五個王，就是耶路撒冷王、希伯崙王、耶末王、拉吉王和伊磯倫王，從洞裡帶出來，到約書亞那裡。 \end{tabularx} \\ \\ \relax
10:24 & \begin{tabularx}{0.7\textwidth}{X} 他們帶出那五個王到約書亞那裡的時候，約書亞就召了以色列眾人來，對和他同去的軍官說：「你們近前來，把腳踏在這些王的頸項上。」他們就近前來，把腳踏在這些王的頸項上。 \end{tabularx} \\ \\ \relax
10:25 & \begin{tabularx}{0.7\textwidth}{X} 約書亞對他們說：「你們不要懼怕，也不要驚惶。當剛強壯膽，因為耶和華必這樣處置你們要攻打的所有仇敵。」 \end{tabularx} \\ \\ \relax
10:26 & \begin{tabularx}{0.7\textwidth}{X} 隨後，約書亞把這五個王殺死，掛在五棵樹上。他們就被掛在樹上，直到晚上。 \end{tabularx} \\ \\ \relax
10:27 & \begin{tabularx}{0.7\textwidth}{X} 日落的時候，約書亞吩咐人把屍首從樹上取下來，丟在他們躲過的洞裡，把幾塊大石頭放在洞口，直存到今日。 \end{tabularx} \\ \\ \relax
10:28 & \begin{tabularx}{0.7\textwidth}{X} 當日，約書亞奪了瑪基大，用刀擊殺城中的人和王，把城中所有人完全滅盡，沒有留下一個倖存者。他處置瑪基大王，像從前處置耶利哥王一樣。 \end{tabularx} \\ \\ \relax
10:29 & \begin{tabularx}{0.7\textwidth}{X} 約書亞和跟他一起的以色列眾人從瑪基大往立拿去，攻打立拿。 \end{tabularx} \\ \\ \relax
10:30 & \begin{tabularx}{0.7\textwidth}{X} 耶和華將立拿和立拿的王也交在以色列人手裡。約書亞攻打這城，用刀擊殺了城中所有的人，沒有留下一個倖存者。他處置立拿王，像從前處置耶利哥王一樣。 \end{tabularx} \\ \\ \relax
10:31 & \begin{tabularx}{0.7\textwidth}{X} 約書亞和跟他一起的以色列眾人從立拿往拉吉去，對著拉吉安營，攻打這城。 \end{tabularx} \\ \\ \relax
10:32 & \begin{tabularx}{0.7\textwidth}{X} 耶和華將拉吉交在以色列人的手裡。第二日約書亞就奪了拉吉，用刀擊殺了城中所有的人，正如他向立拿一切所做的。 \end{tabularx} \\ \\ \relax
10:33 & \begin{tabularx}{0.7\textwidth}{X} 那時基色王何蘭上來幫助拉吉，約書亞就把他和他的百姓都擊殺了，沒有留下一個倖存者。 \end{tabularx} \\ \\ \relax
10:34 & \begin{tabularx}{0.7\textwidth}{X} 約書亞和跟他一起的以色列眾人從拉吉往伊磯倫去，對著伊磯倫安營，攻打這城。 \end{tabularx} \\ \\ \relax
10:35 & \begin{tabularx}{0.7\textwidth}{X} 當日約書亞就奪了城，用刀擊殺了城中的人。那日，約書亞把城中所有的人完全滅盡，正如他向拉吉一切所做的。 \end{tabularx} \\ \\ \relax
10:36 & \begin{tabularx}{0.7\textwidth}{X} 約書亞和跟他一起的以色列眾人從伊磯倫上希伯崙去，攻打這城， \end{tabularx} \\ \\ \relax
10:37 & \begin{tabularx}{0.7\textwidth}{X} 奪了希伯崙，用刀擊敗希伯崙、它的王和屬它的一切城鎮，以及城中所有的人；他沒有留下一個倖存者，正如他向伊磯倫所做的，把城中所有的人完全滅盡。 \end{tabularx} \\ \\ \relax
10:38 & \begin{tabularx}{0.7\textwidth}{X} 約書亞和跟他一起的以色列眾人回到底璧，攻打這城， \end{tabularx} \\ \\ \relax
10:39 & \begin{tabularx}{0.7\textwidth}{X} 奪了底璧和屬它的一切城鎮，又擒獲它的王，用刀把城中所有的人完全滅盡，沒有留下一個倖存者。他處置底璧和它的王，像從前處置希伯崙，處置立拿和它的王一樣。 \end{tabularx} \\ \\ \relax
10:40 & \begin{tabularx}{0.7\textwidth}{X} 這樣，約書亞擊敗全地的人，就是山區、尼革夫、低地、山坡的人，和那裡的眾王，沒有留下一個倖存者。他把凡有氣息的完全滅盡，正如耶和華－以色列的神所吩咐的。 \end{tabularx} \\ \\ \relax
10:41 & \begin{tabularx}{0.7\textwidth}{X} 約書亞從加低斯‧巴尼亞攻到迦薩，又攻打歌珊全地，直到基遍。 \end{tabularx} \\ \\ \relax
10:42 & \begin{tabularx}{0.7\textwidth}{X} 約書亞一舉擊敗了這些王，奪了他們的地，因為耶和華－以色列的神為以色列作戰。 \end{tabularx} \\ \\ \relax
10:43 & \begin{tabularx}{0.7\textwidth}{X} 於是約書亞和跟他一起的以色列眾人回到吉甲的營中。 \end{tabularx} \\ \\ \relax
11:1 & \begin{tabularx}{0.7\textwidth}{X} 夏瑣王耶賓聽見了，就派人到瑪頓王約巴、伸崙王、押煞王， \end{tabularx} \\ \\ \relax
11:2 & \begin{tabularx}{0.7\textwidth}{X} 和北方山區、基尼烈南邊的亞拉巴、低地、西邊多珥山岡的諸王， \end{tabularx} \\ \\ \relax
11:3 & \begin{tabularx}{0.7\textwidth}{X} 以及東方和西方的迦南人、山區的亞摩利人、赫人、比利洗人、耶布斯人，和黑門山下米斯巴地的希未人那裡。 \end{tabularx} \\ \\ \relax
11:4 & \begin{tabularx}{0.7\textwidth}{X} 他們和他們的眾軍都出來，一大隊人馬，多如海邊的沙，並有極多的戰車戰馬。 \end{tabularx} \\ \\ \relax
11:5 & \begin{tabularx}{0.7\textwidth}{X} 眾王組成聯軍，來到米倫水邊一同安營，要與以色列作戰。 \end{tabularx} \\ \\ \relax
11:6 & \begin{tabularx}{0.7\textwidth}{X} 耶和華對約書亞說：「你不要怕他們。明日這時，我必把他們全部交給以色列人殺滅。你要砍斷他們馬的蹄筋，用火焚燒他們的戰車。」 \end{tabularx} \\ \\ \relax
11:7 & \begin{tabularx}{0.7\textwidth}{X} 於是約書亞和所有跟他一起作戰的士兵，來到米倫水邊，突然攻擊他們。 \end{tabularx} \\ \\ \relax
11:8 & \begin{tabularx}{0.7\textwidth}{X} 耶和華將他們交在以色列人手裡，以色列人就擊殺他們，追趕他們到西頓大城，到米斯利弗‧瑪音，直到東邊米斯巴的山谷。以色列人擊殺他們，沒有留下一個倖存者。 \end{tabularx} \\ \\ \relax
11:9 & \begin{tabularx}{0.7\textwidth}{X} 約書亞照耶和華所吩咐他的去做，砍斷他們馬的蹄筋，用火焚燒他們的戰車。 \end{tabularx} \\ \\ \relax
11:10 & \begin{tabularx}{0.7\textwidth}{X} 那時，約書亞轉回，奪了夏瑣，用刀殺了夏瑣王。先前夏瑣在這些王國中是為首的。 \end{tabularx} \\ \\ \relax
11:11 & \begin{tabularx}{0.7\textwidth}{X} 以色列人用刀擊殺城中所有的人，把他們完全滅盡；凡有氣息的，沒有留下一個。約書亞又用火焚燒夏瑣。 \end{tabularx} \\ \\ \relax
11:12 & \begin{tabularx}{0.7\textwidth}{X} 約書亞奪了這些王的一切城鎮，擒獲了這些王，用刀殺了他們，把他們完全滅盡，正如耶和華的僕人摩西所吩咐的。 \end{tabularx} \\ \\ \relax
11:13 & \begin{tabularx}{0.7\textwidth}{X} 至於造在山岡上的城鎮，除了夏瑣以外，以色列人都沒有焚燒。約書亞只焚燒了夏瑣。 \end{tabularx} \\ \\ \relax
11:14 & \begin{tabularx}{0.7\textwidth}{X} 從那些城鎮所奪的財物和牲畜，以色列人都取為自己的掠物。至於所有的人，他們都用刀殺了，直到滅盡；凡有氣息的，沒有留下一個。 \end{tabularx} \\ \\ \relax
11:15 & \begin{tabularx}{0.7\textwidth}{X} 耶和華怎樣吩咐他的僕人摩西，摩西就這樣吩咐約書亞，約書亞也照樣做了。凡耶和華所吩咐摩西的，約書亞沒有一件偏離不做的。 \end{tabularx} \\ \\ \relax
11:16 & \begin{tabularx}{0.7\textwidth}{X} 約書亞奪了那全地，就是山區、整個尼革夫、歌珊全地、低地、亞拉巴、以色列的山區和山下的低地， \end{tabularx} \\ \\ \relax
11:17 & \begin{tabularx}{0.7\textwidth}{X} 從上西珥的哈拉山，直到黑門山下面黎巴嫩平原的巴力‧迦得。他擒獲了那裡的眾王，把他們殺死。 \end{tabularx} \\ \\ \relax
11:18 & \begin{tabularx}{0.7\textwidth}{X} 約書亞和這些王作戰了很長的一段日子。 \end{tabularx} \\ \\ \relax
11:19 & \begin{tabularx}{0.7\textwidth}{X} 除了希未人基遍的居民之外，沒有一城與以色列人講和，都是以色列人作戰奪來的。 \end{tabularx} \\ \\ \relax
11:20 & \begin{tabularx}{0.7\textwidth}{X} 因為耶和華的意思是要使他們的心剛硬，來與以色列人作戰，好使他們全被殺滅，不蒙憐憫，反被除滅，正如耶和華所吩咐摩西的。 \end{tabularx} \\ \\ \relax
11:21 & \begin{tabularx}{0.7\textwidth}{X} 那時約書亞來到，剪除了住山區、希伯崙、底璧、亞拿伯、整個猶大山區和以色列山區的亞衲族人。約書亞把他們和他們的城鎮盡都毀滅。 \end{tabularx} \\ \\ \relax
11:22 & \begin{tabularx}{0.7\textwidth}{X} 以色列人的地中沒有留下一個亞衲族人，只有一些還留在迦薩、迦特和亞實突。 \end{tabularx} \\ \\ \relax
11:23 & \begin{tabularx}{0.7\textwidth}{X} 這樣，約書亞照著耶和華所吩咐摩西的一切話奪了那全地，就按著以色列支派所得的份把地分給他們為業。於是國中太平，沒有戰爭了。 \end{tabularx} \\ \\ \relax
12:1 & \begin{tabularx}{0.7\textwidth}{X} 這些是以色列人在約旦河東，向日出的方向，從亞嫩谷直到黑門山，以及東邊亞拉巴的整個地區所擊殺的王和所得的地： \end{tabularx} \\ \\ \relax
12:2 & \begin{tabularx}{0.7\textwidth}{X} 有住希實本的亞摩利王西宏，他統治的地從亞嫩谷邊的亞羅珥起，包括谷中之城和基列的一半，直到亞捫人邊界的雅博河， \end{tabularx} \\ \\ \relax
12:3 & \begin{tabularx}{0.7\textwidth}{X} 以及從東邊的亞拉巴，直到基尼烈海，又向東通過伯‧耶施末的路，直到亞拉巴的海，就是鹽海，再往南直到毗斯迦山斜坡的山腳。 \end{tabularx} \\ \\ \relax
12:4 & \begin{tabularx}{0.7\textwidth}{X} 又有巴珊王噩，他是利乏音人所剩下的，住在亞斯她錄和以得來。 \end{tabularx} \\ \\ \relax
12:5 & \begin{tabularx}{0.7\textwidth}{X} 他統治的地是黑門山、撒迦、巴珊全地，直到基述人和瑪迦人的邊界，以及基列的一半，直到希實本王西宏的邊界。 \end{tabularx} \\ \\ \relax
12:6 & \begin{tabularx}{0.7\textwidth}{X} 這兩個王是耶和華的僕人摩西和以色列人所擊殺的。耶和華的僕人摩西把他們的地賜給呂便人、迦得人和瑪拿西半支派的人為業。 \end{tabularx} \\ \\ \relax
12:7 & \begin{tabularx}{0.7\textwidth}{X} 這些是約書亞和以色列人在約旦河西所擊殺的諸王，他們的地從黎巴嫩平原的巴力‧迦得，直上到西珥的哈拉山。約書亞按著以色列支派所得的份把這地分給他們為業， \end{tabularx} \\ \\ \relax
12:8 & \begin{tabularx}{0.7\textwidth}{X} 就是赫人、亞摩利人、迦南人、比利洗人、希未人、耶布斯人的地，包括山區、低地、亞拉巴、山坡、曠野和尼革夫。 \end{tabularx} \\ \\ \relax
12:9 & \begin{tabularx}{0.7\textwidth}{X} 這些王是：耶利哥王一人，靠近伯特利的艾城王一人， \end{tabularx} \\ \\ \relax
12:10 & \begin{tabularx}{0.7\textwidth}{X} 耶路撒冷王一人，希伯崙王一人， \end{tabularx} \\ \\ \relax
12:11 & \begin{tabularx}{0.7\textwidth}{X} 耶末王一人，拉吉王一人， \end{tabularx} \\ \\ \relax
12:12 & \begin{tabularx}{0.7\textwidth}{X} 伊磯倫王一人，基色王一人， \end{tabularx} \\ \\ \relax
12:13 & \begin{tabularx}{0.7\textwidth}{X} 底璧王一人，基德王一人， \end{tabularx} \\ \\ \relax
12:14 & \begin{tabularx}{0.7\textwidth}{X} 何珥瑪王一人，亞拉得王一人， \end{tabularx} \\ \\ \relax
12:15 & \begin{tabularx}{0.7\textwidth}{X} 立拿王一人，亞杜蘭王一人， \end{tabularx} \\ \\ \relax
12:16 & \begin{tabularx}{0.7\textwidth}{X} 瑪基大王一人，伯特利王一人， \end{tabularx} \\ \\ \relax
12:17 & \begin{tabularx}{0.7\textwidth}{X} 他普亞王一人，希弗王一人， \end{tabularx} \\ \\ \relax
12:18 & \begin{tabularx}{0.7\textwidth}{X} 亞弗王一人，拉沙崙王一人， \end{tabularx} \\ \\ \relax
12:19 & \begin{tabularx}{0.7\textwidth}{X} 瑪頓王一人，夏瑣王一人， \end{tabularx} \\ \\ \relax
12:20 & \begin{tabularx}{0.7\textwidth}{X} 伸崙‧米崙王一人，押煞王一人， \end{tabularx} \\ \\ \relax
12:21 & \begin{tabularx}{0.7\textwidth}{X} 他納王一人，米吉多王一人， \end{tabularx} \\ \\ \relax
12:22 & \begin{tabularx}{0.7\textwidth}{X} 基低斯王一人，靠近迦密的約念王一人， \end{tabularx} \\ \\ \relax
12:23 & \begin{tabularx}{0.7\textwidth}{X} 多珥山岡的多珥王一人，吉甲的戈印王一人， \end{tabularx} \\ \\ \relax
12:24 & \begin{tabularx}{0.7\textwidth}{X} 得撒王一人，共三十一個王。 \end{tabularx} \\ \\
[1ex]
\hline
\hline
\end{longtable}
$^{1}$主持人:各位弟兄姊妹,平安!.
當我們最近讀到《約書亞記》的時候,.
不知道你心目中有沒有一種複雜的感覺?.
因為《約書亞記》幾章裡面充滿爭戰的畫面,.
城池被攻破,諸王被擊敗,名單一個一個被列出..
但究竟聖經所記載的,所關心的,是不是純粹戰爭的歷史呢?.
尤其是當日的爭戰,比起今時今日,.
我們都已經不是拿刀,拿劍去征戰的了..
但是我們仍然活在各樣的屬靈爭戰當中,.
不論是面對生活上的環境的壓力,.
在不同的際遇裡面所面對的信仰的掙扎,.
世界,文化所給我們的衝擊,.
都會讓人感到無助的..
不單止個人,信仰群體也是,.
不論是你的小組,分隊,成長班,.
你所事奉的部,甚至是城鎮,教會,.
我們怎樣為了主而要贏取生命呢?.
而今天我想跟大家一起看的《約書亞記》十到十二章,.
希望和弟兄姊妹環繞一個主題去思想,.
就是怎樣去經歷有神誓令的爭戰..
當然在認知上,我們都知道有神誓令的爭戰當然好,.
但是在憑信心走在神的旨意當中的時候,.
我們仍然需要作出不少的抉擇..
上次Raymond在第九章提醒我們需要謹慎的,.
不要盡信人所展現出來的外貌,外表,以至落入這些計謀當中,.
而是需要謹慎的去求問神的旨意..
而當我們來到第十章的時候,聖經又給我們看到另外一面了..
當中所講的不單止是神的大能和神蹟,.
也都當中有些很刻意的記載,人怎樣對神回應呢?.
人對人怎樣去回應呢?互動呢?.
我相信是要讓後世包括我們的人,能夠明白一些走在神誓令之下的原則..
而我今天都很高興,很感恩,.
和大家分享一些在經文裡面自律行間,一些精妙之處和心得,.
希望能夠成為我們自己生命,教會群體,作為更好管家的一個準備..
在我繼續講解之前,讓我們先有祈禱..
聖經記載了當時的大能,大力,.
也都記載了一些怎樣來引領人呢?.
怎樣來預納人呢?的一些記載..
天父,讓我們能夠在聖靈光照之下,明白你的真道,.
讓我們生命被餵養,能夠成為一個更合神心意的人..

$^{41}$我們祈禱奉主耶穌基督的聖名祈求..
阿們..
於是五個阿摩利王就聯盟起來要攻打基片..
因為基片和以色列人立約,會更加對他們不利..
在我們讀第九章的時候,都有一點印象,.
這些民族為什麼會這麼精心地去騙約書和那些以色列人的領袖呢?.
以為他們可能都是弱小,大國之下無外交,保住自己的命..
但是不論九章末,當以色列人去接收那些城邑,.
到第十章二節就提到原來基片是一座大城,.
如都城一般比艾城更大,.
並且城內的人都是勇士,原來他們是一個有實力的群族..
所以為什麼以色列人當初發怨言?.
阿摩利王原本想著他的算盤是有基片頂住以色列人,.
但是現在基片竟然背叛他們和以色列人立約,.
必須要嚴厲懲戒,以免更多城邑變節去求和..
於是,螢光幕那裡就是第十章五到第六節,.
我們一起去讀發生什麼事?.
於是五個阿摩利王,.
就是耶路撒冷王,希伯倫王,耶末王,拉吉王,伊基倫王,.
大家聚集,率領他們的重軍上去,對著基片安營,攻打基片..
基片人就打發人往吉甲的營中去見約書亞說,.
「你不要袖手不顧你的僕人,求你速速上來拯救我們,幫助我們,.
因為住在山地的阿摩利人和諸王都聚集攻擊我們.」.
以色列人面對基片人的求救會有什麼反應呢?.
首先背景我們都知道,.
以色列人本身對基片人都已經有所懷疑,.
這班人文化,軍事都差不多,.
騙得我們一次會不會有變節,.
和加納人結盟,讓我們腹背受敵呢?.
加上基片人明顯是用詭計欺騙以色列人立約在先,.
所以當他被人打的時候,去求助的時候,.
這個可能是一個好機會,.
其實對以色列人來說是一個大的考驗..
這個考驗就是約書亞和以色列人應不應該守這個約..
情況不是這麼簡單的,.
假如以我們有所認知的,.
例如中國在戰國的時候在歷史上群雄崛起,.
叢林伐則,弱肉強食的時代思維,.
以色列人就會想,.
他們基片人都是加納人,外邦人來的,.

$^{81}$讓他們鬼打鬼,.
消耗都差不多的時候,.
一來我們又不怕基片人會背叛,.
又不怕那些阿摩利王打完仗之後被削弱,.
我們日後攻取他們更加容易,.
這是一個千次良機..
當然有些人又會想,.
不行,在博弈上不可以不幫,.
畢竟一來我們真的有約,.
但假如反口不幫的話,.
基片人立即翻臉,.
聯同那五王來打以色列人,.
那就讓我們入困境當中..
或者可能有些人會想,.
比較狡猾的,.
拍心口,.
基片人不用怕,一定幫你的,.
但實際上慢慢走,走著走著走著走著走著走著走著走著走著走著走著走著走著走著走著走著走著走著走著走著走著走著走著走著走著走著走著走著走著走著走著走著走著走著走著走著走著走著走著走著走著走著走著走著走著走著走著.
走到戰場的時候,.
雙方打到七七八八了,.
死傷八八九九了,.
才出場,.
還要大大聲地喊,.
基片人我們現在來幫你了,.
衝啊!.
可能直至結果,.
已經疲憊的五王聯盟不堪一擊,.
差點被滅族的那些得救的基片人,.
又對以色列人一千多謝萬多謝這樣,不知道世俗的政治家會不會這樣去盤算,但是約書亞作為神的僕人,.
帶領以色列民的領袖,他怎樣反應呢?.
我們繼續看第七節,.
那裡記載,於是約書亞和他一切的兵丁,並大量的勇士,都從吉甲上去,約書亞他二話不說,親自上去,.
當然帶上全軍,包括大量的勇士,即精銳的部隊,從大本營吉甲那裡上去,要幫助基片人面對五王的攻擊,.
在約書亞的反應裡面,讓我們看到一個屬於人的原則,就是神的百姓必須守約,即使當初的決定有所疏忽,過程不完善,.
納了約,在神面前說了的話,就必須要信實地履行,因為我們所信的神就是守約,詩詞愛的神,.
為什麼以色列人現在可以去到迦南地呢?就是因為神信守他幾百年前和亞伯拉罕信納的約,.
就算以色列人多次犯罪得罪神,在曠野接受神的管教也好,神都沒有反口,仍然紀念他的約,帶領以色列人進到迦南地,.
在這裡看到,神都很期望他的子民是守約的子民,而約書亞的舉動得到神的認可,.
然後第八節,神主動地跟約書亞說,不要怕他們,因為我已將他們交在你手裡,他們沒有一人能在你面前站立得住,.
這些做得正確的事情,人還沒有去求問神,神已經主動地鼓勵了,約書亞不單止守約,他更是全心全意地守約,.

$^{121}$然後第九節我們看到約書亞就中夜從吉格上去,猛然臨到湯門那裡,他率領這個軍隊夜間急行,.
中夜上去,由吉格去到基片那地理上是難行的山路,他們將會面對一場疲憊的戰鬥,不是輕鬆的..
在繼續看這個爭戰之前,我先思想有關於守約,誠信的課題,因為我們身處的時代,欺騙的事情越來越多,.
難免讓我們有所懷疑,繼續做誠實人,做老實人,守約是不是真的值得呢?.
我先在這裡舉另一個聖經的事件記載,也是和欺騙人有關的,大家有聖經,翻開,或者按下去,撒母耳記夏,.
撒母耳記夏,第21章,時空座標就是當時大衛作以色列王第二任的君王,距離約書亞時代已經是四百年之後了,.
21章,開頭經文記載,大衛年間有饑荒,一年,即是連續三年,大衛就求問耶和華,耶和華說,這饑荒就是因為塑羅和他流人血之家殺死羈騙人,.
第二節解釋,原來這羈騙人不是以色列人,是阿摩利人裡面所剩下的,以色列人曾向他們起誓不殺滅他們,.
但塑羅卻為以色列人和猶太人大發熱心,想要殺滅他們,饑荒在古代的時候不時都會發生的,隔兩三四年,失修一次,.
都是天地間正常的循環,好像做生意有賺有蝕,勝敗乃兵家常事,連續兩年饑荒,還可以說是不幸,.
但連續三年,古人的智慧都清楚,當中一定有問題,於是大衛求問神,.
神的回答裡面發現原來是因為之前上一手的君王,塑羅王,他忽然間為了以色列猶太人大發熱心,想要殺滅羈騙人,.
這個在聖經裡沒有記載的,在這裡,上帝篤他出來,塑羅王不知道是不記得,還是無視四百年前以色列人曾經和羈騙人所立的約,.
或者,就好像我們今天的角度所說,現在是什麼時代,還提到明末清初幾百年前所立的約,.
現在怎會有效,但是上帝卻是記得,當日約書人和以色列人怎樣和羈騙人立約,.
四百年後經歷過改朝換代,翻天覆地,是亦都好,神依然紀念,.
仍然會因為塑羅王作為以色列人的代表,破壞這個盟約,未得到平反解決,然後讓以色列人需要承受這個懲罰,.
至於大位怎樣處理,大家可以自己去讀,但是想帶出一個重點,就是我們所信的神,是一位信實守約的神,.
當神的子民,神所選立的領袖,要去跟隨,事奉這位神的時候,亦都需要反映出神的信實,.
在神面前,你可能立過字的,許過願的,不論對神,對人都需要遵守,更加神藉著主耶穌基督和我們所立的新約,.
過了二千年後,今日到將來,仍然有效的,用士徒保羅的字眼,羅馬書十章裡面記載,.
就是你若考慮承認耶穌為主,心裡信神,叫他從死裡復活,就必得救,.
這個確實的應許,因為犯求過主名的,就必得救,各位弟兄姊妹,我想你回想,.
你曾經有沒有向神立過字,許過願,不論你立字,我要過一個聖潔的生活,我會處理這些罪,.
或者曾經夢神呼召的,在神面前回應立字,要全職事奉的,現在去到什麼境況呢?有沒有找數呢?.
在這裡分享我年輕時候的一個經歷,我怎樣由工作裡面轉去全職事奉呢?.
那時候青年的營會,在澳洲雪梨那裡都有的,不同教會的年青人聚集在一起,邀請海外港元來港,.
我在那個情況之下,是夢召全職事奉的,那時候他向神祈禱,希望能夠工作三年,.
當中累積一點點社會經驗,最好儲錢,雖然神有供應,能夠做三年工作就好了,.
接著聽到長者的意見,找工作當然是找一份最高薪金的,於是我膽粗粗說不要繼續在法律方面發展,.
去找IT的工作,那時候發展得很快,搶人搶得很厲害,我在那裡一間國際的ERP,.
Enterprise Resource Planning,企業資源的管理的軟件公司裡面工作,.
不是賺很多錢的銷售員,沒有資格,我做的算是那時候整個產業鏈系列裡面的底層,.
於是我就留在自己公司裡面,兩年裡面已經跳了三次薪金,.
接著還派我去新加坡亞太區的總部開一個新的部門,那時候才做了一年多,薪金都頗高,.
就算那時候一個星期,隔個星期坐飛機回Sydney去家裡吃飯,拍拖,見第一次面都完全不覺得肉痛,.
但時間到了,都差不多三年了,怎樣決定呢?如果你是我的世侄去做,你都會說繼續先做吧,.
那麼難得,大有可為,但我記得在城面前立過字,真的要還的,於是和我的上司遞辭益信,.
那時候副總裁問我,你做得好好的,大展拳腳,搞什麼呢?於是我和他解釋福音,信仰,.
為什麼我遵守在神面前的承諾,好像江湖再見,後會有期,我想表達的就是那種掙扎都很真實的,.

$^{161}$面對各人的利益,各人的好處,不變的那種吸引,但是假如在神面前有立過字的話,就必須要償還,.
神是守約的神,神應許他的子民帶領以色列人進入迦南地,去見證他的大能,人的回應很重要,.
接著我們看第十到十一節,在螢幕上面,或者大陸人手上聖經,約書亞第十章十到十一節,.
耶和華使他們在以色列人面前賄亂,約書亞在欺騙大大地殺敗他們,追趕他們,.
在伯羅倫上坡路擊殺他們,直到阿西加和瑪基大,他們在以色列人面前逃跑,正在伯羅倫下坡的時候,.
耶和華從天上降大冰雹在他們身上,直降到阿西加打死他們,被冰雹打死的比以色列人用刀殺死的還多..
在十到十五節那裡,記載了三個神蹟,首先第一個我們在這裡看到的第十節那裡,.
講到耶和華使他們,就是那幫阿摩利王,在以色列人面前賄亂,.
一個字原文也曾經出現在出埃及的過紅海的記載當中,神使得埃及的軍兵混亂,.
在協助和欺騙的爭戰當中,神再一次讓他的百姓經歷到神的大能,.
在他們的人心裡面來搞擾,以至約書亞和以色列人能夠勝過那五王的聯盟..
不單止這樣,在剛才的經文裡面,我們看到當敵軍逃跑的時候,.
耶和華在天上降下大的冰雹,再打死的人,被以色列人用刀殺的還多..
冰雹的現象在以色列地是有發生的,翻查新聞,2015年10月的時候,.
以色列曾經有一個大雷雨,大冰雹,冰雹大到5厘米的直徑,好像網球一樣大,.
你想像一下在天下這樣突發,這樣破壞了首都特拉維夫和周邊加利利地區一帶..
神安他的大能,運用他所創造的氣象,來為以色列人爭戰..
聖經這樣記載,是要讓人知道,勝負的關鍵,真的不在於人的刀劍,而在於神怎樣去介入,去令他們爭戰..
第三方面更加震撼的,在十二到十四節裡面記載了,約書亞向神一個祈禱,.
就今天的言辭來說,都可以算是很誇張的祈禱,向神這樣求,求神說,日頭你要停在機片那裡,月亮你要停止在阿倫谷,.
意思就是讓他們有整天的時間去當中爭戰,失敗受敵..
當然這個說話在科學上是很難解釋的,怎樣整個天體停頓運作,怎樣行的?.
但是這個卻是有份參與戰爭的人,共同的見證和經歷..
當然聖經的重點並不是在於這個天文的現象,而是在於十四節裡面所說的,.
「耶和華聽人的禱告,沒有杖者日的,是因為耶和華為以色列爭戰.」.
當中再帶出一個關鍵,是神願意為他的百姓爭戰,甚至不單止影響人心,.
不單止運用大自然裡面的天氣的現象,更加影響到各樣的時間..
在第十章裡面,神給我們看到,他怎樣為他的子民爭戰..
看到神的信實,他是一位守約,審判,拯救,掌管著創造物,時間的神..
弟兄姊妹,當我們站在神的約裡面,行在神的旨意當中的時候,.
我們要做的就是忠心於神給我們的呼召,忠心於神給我們納的約..
當中我們所需要的爭戰,神自己會成就..
而我們今天的爭戰,起碼在我們地球這個部分,.
就不是刀劍的戰場,而是屬靈的戰場..
罪孽的權勢對我們影響,世界所給我們的壓力,.
經歷人心的那種詭詐,都真實地影響著我們..
約書17第十章提醒我們,真正的德性不是在於人自己的能力單打獨鬥,.
也不是群體的力量,人多就好辦事,而是在於神是不是在我們當中最令我們爭戰..
而這個爭戰,不單止在第十章裡面看到的,.
到了我們回到第十一章的時候,發覺這個爭戰更加大,更加嚴重,.

$^{201}$發現一個更加明顯的對比,就是敵人比以前更加強大,規模更加多,.
因為這個是整個北面迦南地的聯盟..
十一章四字那裡記載到,他們人數多於海邊的沙,並且有許多的馬匹和戰車..
在古代近東的世界裡面,馬匹和戰車所代表的就是軍事那種優勢,科技的力量..
以色列人在曠野那裡逃難出來,本質上有把刀,有把棍已經算不錯了,.
怎樣跟那些馬匹和戰車來爭戰呢?.
在人的角度來說,這場戰爭根本是不能打的,.
但是它是正正在這種對比之下,更加顯出神作為的清晰..
我們一起讀一下第六節,十一章第六節,.
神怎樣對約書亞說,一起讀..
"耶和華對約書亞說,你不要因他們懼怕,明日這時,.
我必將他們交付以色列人,串燕殺了,.
你要砍斷他們馬的筋,用火焚燒他們的車輛.".
這個說話是神給約書亞一個確據,.
在當中這個命令不單止說到神要將他們交給以色列人,.
更加有一個很特別的吩咐,要砍斷馬的蹄筋,焚燒戰車..
根據學者的解釋,為什麼要砍斷馬的蹄筋呢?.
功用就在於那些馬不能夠再作太過劇烈的運動,.
例如打仗,拉著戰車,周圍橫衝直撞那些,不行了..
簡單的工作,例如耕田,慢慢的運輸,仍然可以的..
意思就是他們將不會再成為征戰的工具,.
也都是間接向約書亞宣告,啟示在這個之後,.
不會再有大型的戰爭了,和平將會來..
這個作用就好像《爾賽書》第二章記載,.
神必在列國當中施行審判,那些國民要將刀打成泥頭,.
將槍,即是戰爭的工具,打成鐮刀,修割的工具..
這國不僅舉刀攻擊那國,他們也不再學習戰事,.
是一個和平的宣告..
況且以色列人假如將來需要征戰的話,.
神就是他們的力量..
正如詩篇二十篇那樣說,有人靠車,有人靠馬,.
但我們要提到的是耶和華,我們神的名..
這個就約書亞記第十一章裡面,被人看到的..
神的確都很明白人的心,假如真的要打仗的時候,.
假如有馬,有戰車,不就很好,拿來用..
或者在困難的工作的時候,譬如我們很容易依靠科技,.
問AI有什麼工具..
神要提醒他的子民,他所更加需要的,.
就是人順靠神為他們征戰..
而在第九節的記載,約書亞在打完仗之後的確,.

$^{241}$根據神所吩咐的一字不漏地將他的戰馬的根提劈斷,.
然後用火焚燒車輛..
這裡展現出約書亞對神吩咐的忠心..
而這份忠心也不單止停在這裡,.
因為在第十一章裡面我們看到有幾個重複的話語,.
分別在第十五,二十,二十三節裡面提到的,.
就是正如耶和華所吩咐摩西的..
這裡所說的,就是約書亞所做的,就是完成著摩西未完成的使命..
有時在這裡的時候,我們都需要被提醒,.
神的工作有時真的會跨越世代的..
摩西他沒有帶以色列人進入迦南地,.
但是神給他的使命,他在約書亞身上繼續..
今日我們所見到的事公,或者想起建堂,.
上一代未必能夠完成的,可能被這一代的領袖當中完成它..
或者我們福岡的傳揚,.
我們尤其是宣教,未必在今世代裡面見到的果效,.
尤其是這些很難傳福岡的地方,.
不論是日本,或者東南亞的地方,.
都可能透過下一代繼續去推進..
教會的歷史,宣教上都是這樣..
但是起碼在約書亞劇裡面,第23節,即第11章尾,.
給我們看到一個重要的總結,.
就是約書亞照著耶和華所吩咐摩西的一切話,.
奪了那全地,於是國中太平,沒有征戰了..
在當中,征戰完成,享受到和平了,.
但是這裡以色列民需要小心,.
今日我們都要小心,安息並不代表躺平,放軟手腳,.
雖然主要的征戰已經完成,.
但是以色列民仍然需要去收取那些地方,.
下個星期,瑞豐,劉牧師就會和我們講解,.
你們擔然不去拿那些地要到何時呢?.
和我們思想..
安息也不代表屬靈上放鬆,.
因為聖經經文記載,.
事司紀之後那本書很快就看到,假如他們對神不忠的話,.
太平也不能夠持久,仍然需要接受神的管教審判..
在整個救恩計劃當中,.
會不會約書亞的記載,.
亦都在預表著主耶穌基督所帶給我們真正的安息呢?.
約書亞他帶領以色列民過約旦河,進入應許之地,.

$^{281}$帶來國中太平呢?.
在新約希伯來書的記載,.
假如那時候他得到應許之地,已經得到安息了,.
神為何在詩篇裡面再提這個安息的日子呢?.
意思就是最終為人帶來安息的是主耶穌基督..
主耶穌基督帶領他的百姓,.
率領他的百姓進入一個永恆的國度,.
在十字架上贏取最終的征戰..
亦都是主耶穌使我們人心裡面得著神所賜的真正的安息..
講了很久了,有個小小的總結,.
第十章,第十一章,經文給我們看到兩大方面..
第一,我們看到神是一位信實守約的神,.
他的應許是不會落空的,.
神亦都很欣賞我們信守承諾的人,.
親自率領他們進行征戰..
當我們回想這個世代對信守承諾放在一個低的位置的時候,.
不論是因為利益,那句話不算數,.
或者因為自己的固執不願意改自己的壞習慣在罪裡面,.
或者放縱私慾,對婚姻的承諾看得很輕,.
以致來來合合帶來很多的傷害,.
這些各種的,都需要讓我們重新回到神的身邊,.
知道我們需要在神的面前守約,.
縱使我們有軟弱,人的能力有限,.
我們都更加需要主人所叫的率領我們去征戰..
第二,我們看到神既然是一位為他子民征戰的神,.
我們看到他不論是影響人的內心,.
運用大自然的力量,控制時間,.
神都很欣賞那些願意樂意信靠他能力的人,.
包括在禱告裡面展示出對他這份信靠的人,.
所以弟兄姊妹說他們今日面對著的外在壓力,.
好像第十一章所形容的那樣,.
好像海邊的沙那樣多,.
或者文化環境對你的挑戰好像戰車馬匹那樣凌厲也好,.
但神仍然對我們說不要懼怕,.
我們需要信靠那位率領我們征戰的神..
最後雖然第十二章我們沒有看到,.
但是拿一拿印的聖經也好,你的手機也好,.
很多全部都是人名,地名那些,.
所記載的是什麼呢?.
就是列舉摩西和約書亞他們在征戰當中所打敗的王,.

$^{321}$摩西打敗兩個王,巴沙王岳,.
然後約書亞也記載他打敗了三十一個王,.
神將這三十三位王趕出神賜給以色列人為業的地,.
這個我們用熟悉的眼光去看,.
會不會象徵著我們在主耶穌基督的率領之下,.
將罪惡從神賜給我們的生命當中逐一趕出呢?.
在準備的時間,.
我會想起在新約的時候,.
史特普羅寫過一些關於罪惡的列表,.
譬如在羅馬書一章,.
裝滿了各種不義,邪惡,貪婪,惡毒,.
滿心是嫉妒,兇殺,增敬,詭詐,毒行,財毀,.
背後說人的怨恨神的,武蠻人的,狂傲自誇的,.
劣造惡事的,違背父母的,無知的,背弱的,.
無親情的,不憐憫人的..
提摩泰後書,.
史特普羅寫給他的徒弟提摩泰,.
那時候人怎樣犯罪法呢?.
人要尊顧自己,貪愛錢財,.
自誇,狂傲,暴毒,違背父母,忘恩負義,.
心不聖潔,無親情,不解怨,好說,殘言,.
不能自虐,性情兇暴,不愛良善,.
賣主賣友,任意妄為,自高自大,.
愛燕樂,不愛神..
當然在聖經裡面,.
形容罪的字眼可能還有很多,.
或者有不同的列表,.
但在準備過程當中,我看這兩段,.
數一數那些不重複的字眼,.
剛剛是33個..
對於我來說,只是下一刻而已,.
這麼巧合..
會不會也是象徵著,.
我們在屬靈生命裡面,.
個人的罪惡,.
就好像在迦南地佔據著.
神所應許賜給以色列人的地的王一樣,.
需要逐一地趕出,.
靠著主耶穌基督的說領之下,.
我們能夠逐一戰勝,.

$^{361}$以致我們得享.
神所賜給我們的平安喜樂,.
未需要去到天堂的時候,.
都要享受到神所給我們的那份安息..
雖然我們所面對的,.
再不是迦南王了,.
再不是拿刀拿槍上陣了,.
但屬靈戰爭,.
卻是在我們日常生活當中,時刻地去打..
不論是面對會激我們的人,.
讓我們發揮出我們的罪性,.
各樣的貪婪,罪惡,不小心,.
或者對人的那種粗暴,.
新約很清楚地告訴我們,.
我們這個其實是跟屬靈的惡魔爭戰,.
好像在耶福所書的六章,.
我們一起讀,.
我們並不是與屬血氣的爭戰,.
乃屬於那些執政的,掌權的,.
管轄者幽暗世界的,.
以及天空屬靈氣的惡魔爭戰..
舊日約書人所面對的,.
就是城池,.
那些不敬畏神的民族,君王,.
今日我們自己面對的,.
是各樣的罪惡,.
或者我們在生命上的影響,.
搞擾,.
或者在教會上面對的,.
不論是文化的壓力,.
或者是價值觀的侵蝕等等,.
這些都是屬靈的爭戰,.
我們需要依靠主耶穌基督,.
率領我們去勝過這一切,.
希望在約書一記裡面,.
讓我們有這些得著,.
我們一起祈禱..
掌管歷史的神,.
你的確是一位守約,.
思前愛的神..

$^{401}$在聖經裡面,.
我們看到你為以色列人民怎樣去爭戰,.
也在歷史裡面,.
看到你怎樣為教會爭戰..
雖然我們所讀到的各樣的神蹟,.
不論是日頭停柱,.
大的冰雹,.
或者人心的混亂,.
希望我們都能夠經歷今時今日,.
神你怎樣率領我們,.
在對抗罪惡上面,.
能夠得勝..
在主耶穌基督裡面,.
得著真正的安息..
求神你讓我們記得,.
主耶穌基督已經在十字架上面,.
成就救恩,.
已經得勝..
讓我們的一生,.
以及成就家,.
都能夠成為紀念你恩典的名單,.
讓後來的人看得見,.
神你為我們爭戰..
我們祈禱,.
奉我主耶穌基督的聖名祈求,.
阿們..
\newpage



\section{創世記 41:38}
\label{sec:031861}
\textbf{聖經人聖經事:神的靈在他裡頭}~~~~黃克勤牧師博士
\newline
\newline
連結: \href{http://www.church.com.hk/acms/content.asp?site=cdc\&op=show\&type=product\&code=031861\&layout=sermon}{\texttt{031861}} ~~~~ 語音日期: 2026-03-21
\newline
\newline
\hyperref[sec:031860]{\small{< < < PREV SERMON < < <}}
~
\hyperref[sec:index_chronic]{\small{[返順時目]}}
~
\hyperref[sec:031862]{\small{> > > NEXT SERMON > > >}}
\newline
\newline
創世記 41:38
\newline
\begin{longtable}{cl}
\hline
\hline
章節 & 經文 (和合本修訂版)\\
\hline
41:38 & \begin{tabularx}{0.7\textwidth}{X} 法老對臣僕說：「像這樣的人，有神的靈在他裡面，我們豈能找得著呢？」 \end{tabularx} \\ \\
[1ex]
\hline
\hline
\end{longtable}
$^{1}$弟兄姊妹,又讀聖經啦.
每日讀一段.
《聖經人》《聖經事》.
看看這個人物是誰.
故事是甚麼.
而看到神的手在背後工作.
心力書書的網絡.
是很開心和大家一起讀聖經的.
今次的《聖經人事》.
就是神的靈在他裡面.
這裡有兩個對象.
有兩個人物.
不是人物來的.
第一個是神的靈.
神的靈呢.
雖然眼看不到.
但是很活的.
很活躍的.
不停地做事的.
神的靈在聖經裡面最先見到是甚麼呢.
神的靈運行在水面上.
創造的時候.
神的靈已經參與了.
父神創造的時候.
找了兩個極佳的工人.
極佳的製作者.
一個是神的靈.
和聖子耶穌.
聖子 他的兒子.
六天就造起一個.
他心目中造的天地.
然後神的靈繼續工作.
從來沒有停止過.
神的靈今天仍然在工作.
神的靈在他裡面.
這個他是誰呢.
若實.
神的靈是不是在每一個人裡面呢.
信了主的就是.
但是當時不是.

$^{41}$如果若實的時代.
舊約的時代.
神的靈未必永遠在一個人裡面.
有時候呢.
如果神感動他就會是.
有時候不是.
特別是未認識耶和華.
神的靈有時候會感動他.
外邦人都會發生一些事.
是跟隨神的.
但是不是像我們今天這樣.
基督徒是有神的靈.
長久永遠在我們裡面.
神的靈和若實.
是怎樣的呢.
41章第37節.
神命定的事是速速成就.
昨天說.
你見到若實由他出生.
一直到他穿起一件彩衣.
被哥哥嬲他.
然後賣他.
然後被主母陷害他.
然後坐了監.
然後現在出來.
解夢.
然後呢 嘩.
這一段說就是.
若實被提升到.
一人之下 萬人之上.
整件事.
都可以見到神的靈.
在背後工作.
神掌管的事.
沒有一件失手的.
沒有一件是會說.
哎呀 一時靈不穩.
跌了 沒有的.
沒有意外的.
因為他都編排了.

$^{81}$而且他有能力.
叫事情按著.
他的想法去做.
神命定的事.
必速速成就.
昨天就說了.
來到這裡承接著昨天.
法老和他一切.
神福都以為.
這事為妙.
哎呀 對呀 說得對.
好呀 真的好.
要找一個人去.
管理 法老就對.
神福說.
不是跟若實說.
跟神福說 大家奉福.
法老對神福說.
像這樣的人.
最重要的是第一.
有神的靈在他裡頭.
我們豈能.
找得著呢.
他們都信很多神.
有很多.
祭司 他們的神.
但是.
法老這個時候知道.
若實的神.
才厲害 才真.
我們那些神.
全部都不夠他來都解不了.
又指示不了.
現在要了 像這樣的人.
要找一個.
有智慧 有聰明的人去找.
去處理這些.
糧食 哎呀.
像這樣的人 有神的靈.
在他裡頭 我們.

$^{121}$豈能找得著呢.
沒有了 只有他.
起碼我現在.
見到只有他.
於是法老就對若實說.
「神既將這事.
都指示你 可見.
沒有人像你這樣.
有聰明 有智慧.
由於神的靈在你裡頭.
你是最聰明 最有智慧的了」.
現在.
我找不到別人了.
就直接.
沒有想 是招聘.
擺一個大龍鳳 招聘一下.
沒有.
是你了 為什麼.
法老都看得出.
因為有神的靈.
在你裡頭.
第一點要和大家思考.
這個.
螢光第一點.
其實.
最重要的就是.
有神的靈在我們心裡.
工作 當時更加是.
今天我們每個信主.
都有聖靈 你要讓.
聖靈在我們的.
心裡工作 在你的.
心裡工作.
讓祂的心意在我身上.
可以運行 指示.
帶領 安排.
這個就叫.
今天叫聖靈充滿.
很多人.
今天將聖靈充滿.

$^{161}$聖靈充滿之後.
馬上變成這樣.
祈禱的時候這樣 唱詩的時候這樣.
很厲害 好像和.
主耶穌很近.
整個靈魂感動到我 多數都閉上眼.
很厲害.
又或者一指.
整個跌倒.
很厲害.
祈禱的時候.
整個人跌倒.
很厲害 激到.
聖經裡面沒有.
聖靈激到一個人的.
聖經例子.
我抄過.
激殺就有 激到就.
沒有 為什麼會這樣發生呢.
為什麼會這樣發生呢 我不知道.
我解釋不了.
不過我想說.
聖靈充滿.
未必一定要講方言.
未必一定要可以禱告激到.
或者可以講一些.
而是.
知道一些沒有人知道的事情.
都是聖靈工作.
可能會令到我們這樣.
不過最簡單的聖靈充滿.
就是.
祂在我生命裡面.
坐在.
司機位.
在我屋裡面.
在主人的座位.
我是一個僕人.
雖然是我的身體.
不過我只是祂的僕人.

$^{201}$祂大力保守.
與我同在.
祂指示我 指揮我.
我有信服.
聖靈充滿.
沒有人有這樣的聰明智慧.
於是.
聖靈會帶領我工作.
認識神是一個最重要的智慧.
認識神是智慧的開端.
不過這個智慧.
未必是世間人想的智慧.
有些人.
沒有其他人那麼聰明.
不過他很敬畏神.
其實神很喜歡他 很疼愛他.
很保護他 給他很多恩典.
給他很多平安喜樂.
但他可能不是讀了很多書.
一個博士學位都沒有.
不過神是很祝福他.
我見有很多.
老人家.
屢遇不改去教堂.
敬拜很好.
我相信將來他們做很多事工.
很多禱告.
很多小小錢捐獻.
好像寡婦兩個小錢那樣獻上.
這些人將來在神面前.
領獎排隊.
一定前過我們很多.
其實神的靈在他心裡.
他一直都尊主為大.
其實他真的認識神.
有聰明有智慧.
回到這裡.
若失有神的靈在他腦裡.
沒有人像他那麼有智慧.
最重要是.

$^{241}$法老都認到有神的靈.
在他腦裡.
這個反而有時我們自己都沒有留意到.
是不停地追尋.
我的學術 我的學歷.
我的聰明智慧 我的人脈.
我多有方法.
與手段.
談判.
其實不是的.
神是欣賞那些聽他話的人.
第二點.
你可以掌管我的家.
唯獨在寶座上.
我比你大.
你可以掌管我的家.
我的民都聽從你的話.
不過.
只有一樣東西比你大.
就是我的寶座.
當我在寶座上.
你仍然比我低.
一人之下萬人之上.
很重要.
法老竟然.
本來在這個希伯來人奴隸.
在這個監獄裡.
提出來的.
竟然說你可以掌管.
我的家.
如果說回.
藥室的故事與經歷.
我想他以前的主人.
波提弗都大起一驚.
原來藥室那麼厲害.
最怕的是誰呢.
最怕的是波提弗的太太.
就是他以前的主母.
陷害他那個.
死了,這次我不行了.

$^{281}$不過看來藥室.
沒有跟他計較.
起碼聖經沒有記載.
你掌管我的家.
這個法老說.
掌管我的家.
其實不是真的掌管我家.
你可以掌管.
我的國家.
我的國家就是我的家.
這裡是沒有分開的.
我的家裡有很多人.
就是我的民.
這個家其實應該是指國家.
我的民是指我的百姓.
都要聽你的.
不過你是最大.
但是比起我來說.
你是次一級.
我在寶座上比你大.
然後法老跟他說.
我現在派你.
治理埃及全地.
管理我的國家.
管理我的民.
埃及全地.
都聽你的.
法老就摘下.
手中打印的戒指.
以前人家的戒指.
很大隻.
因為他以為拆除出來.
點些墨就吸進去那些文件.
或者那些皮.
那些葡草紙等等.
就有了印.
皇上的印.
所以他時時帶著印在手上.
而那些人來到.
拜見皇上的時候.

$^{321}$他們除了親吻他的腳.
很多時候就親吻他的戒指.
這件事.
留到西方後期.
都還是這樣.
親吻教宗的戒指.
皇上的戒指.
表示對他的權力.
是馴服的.
打印的戒指.
你見到有些什麼.
紅色字了.
戒指.
戴在你手上.
給了你戴,我沒有了.
我不戴了,因為給了你戴了.
我信任你了.
你掌管了.
有了我的皇權.
將戒指戴上去.
然後穿上細麻衣.
細麻衣.
當時都很貴重.
因為要用麻.
織出來.
你知道埃及人.
很多都不穿衣服.
只是有布.
但去到重要場合.
都有麻衣穿上.
很漂亮的.
而且將金鏈戴在他的頸上.
那條鏈.
是象徵地位.
鏈上有一個陀.
即是那件東西.
還有那條鏈.
看就知道了.
譬如去到很多場合.
大聚會.

$^{361}$各個宗派一起的.
聯合聚會.
一出場.
全部監督,主教.
個個都戴一條鏈.
鏈上有十字架.
主教,監督.
就會知道.
有條鏈.
人類都是這樣承繼下來的.
戴在他頸上.
然後叫約瑟坐他的附居.
什麼叫附居呢.
一想起附居.
就會想起誤中附居.
為什麼誤中附居呢.
其實就是說.
皇帝是不是只有一架車呢.
不是.
很多架車可能要維修.
可能會爛.
又要多修幾架車.
一出去就不要讓人知道.
一排車.
不知道皇上坐在那裡.
所以想攻擊他.
有時五架車出來.
五條路走過去.
哪架車是皇帝坐呢.
原來是六架才是.
所以皇帝的車.
很多架是一樣.
或者有少少不同.
不知道他坐在那裡.
坐他的附居.
其他的車.
給一架皇帝的車.
而且有戒指.
有細馬衣.
今年有皇帝的車.

$^{401}$還有鶴道的.
是.
皇帝出來的時候.
他的車到.
前面有些馬兵和士兵在叫鶴道.
讓開.
皇帝駕到.
然後鶴道就呼叫.
跪低.
全部跪低.
多厲害.
你出去的時候.
你坐我的車.
士兵在你面前.
就會叫人跪低.
很厲害.
有排場.
這樣的辦法.
戒指,細馬衣.
今年有皇上附居.
還有鶴道的士兵.
叫人跪下.
出來的排場是什麼呢.
排場就是.
皇帝來到.
很厲害.
你看到.
這個藥室.
本來是一個少年人.
穿著彩衣.
一直都迷路去找哥哥.
誰知被人賣了.
陷害到做家奴.
奴隸到陷害.
到現在坐牢.
然後做了皇帝.
給皇帝你做就是在做.
這麼厲害.
治理埃及全地.
好了.

$^{441}$法老對藥室說.
我是法老.
其實這個對藥室說.
其實是眾神一起在聽.
藥室.
你聽著我是法老.
在埃及全地如果沒有你的命令.
不許人擅自辦事.
其實這句.
是一方面說給藥室聽.
另一方面.
是所有群臣都聽到.
做筆記的官.
馬上寫下.
要發文件出去.
其實所有人都.
沒有你的命令.
不許人擅自辦事.
這個擅自辦事.
好笑的.
如果你看皇后本聖經.
很多就是.
動手動腳.
沒有人.
可以擅自.
動手動腳.
因為你知道法老認不認識.
他的官員.
很多說現在要儲糧.
可能有其他低級官員.
說對了儲糧.
先放在我那裡.
要那麼多.
儲多五分之一.
不行.
當然要五分之二.
他收了之後.
加五分之一.
五分之一去自己家.
如果沒有你的命令.

$^{481}$不許人擅自動手動腳.
然後.
法老又.
賜一個新的名給藥室.
你這個希伯來人.
Yosef.
神增加.
Yosef不夠好.
我給你改一個新名.
這個新名很難讀.
不過我都要譯給大家看.
新名叫.
薩法那剔巴萊亞.
那麼長.
薩法那剔巴萊亞.
其實是一個埃及名.
埃及的字.
Jove.
這個Jove.
Jove.
PHR可以讀F.
Jove.
薩法那.
巴萊亞.
那剔.
為什麼會譯成那剔呢.
Jove.
那巴萊亞.
巴萊亞.
Jove.
薩法那剔.
巴萊亞.
巴萊亞就是他活了.
所以整個名字.
薩法那剔就是.
Jove.
兩個音.
怎麼會有四個音呢.
為什麼有四個音我就不知道了.
PH就是一個法字.

$^{521}$那.
然後那剔.
不知道怎麼來的.
PH這個音.
然後巴萊亞.
巴萊亞.
整個名字的意思就是.
薩法然他活了.
就是說神透他說話.
哇.
有生命力.
的意思.
又將安城一個很出名的城市.
安.
的祭司波提拉菲的女兒.
叫阿西納.
就是給她為妻.
挺有趣的.
將別人的女兒給她做太太.
也可以的.
當時的法老真的挺大力的.
這個阿西納呢.
就叫做.
阿善那.
阿善那這個那.
是一個女神的名字.
就是Nav.
阿善.
就是屬於Nav.
女神的.
屬於.
這個女兒.
就是屬於Nav.
這個女神的.
名字就叫做.
整個女兒的名字就叫做.
阿善那.
好了.
有個疑問.
約瑟有沒有問題.

$^{561}$娶一個希伯來人.
娶一個外族的.
埃及女子還是有.
其他信仰的.
有些.
問題的.
就是說其實會令到.
他的信仰容易會走歪.
但是約瑟.
也算是保存得.
很好.
皇上賜一個.
當地的女子給他做太太.
也不能夠.
推搪.
於是就娶了.
很感恩.
以後約瑟並沒有因為.
娶了一個外籍太太.
當地的太太.
不同信仰的太太.
而影響到他的靈性.
也挺好.
也都.
不算是.
很差.
其實.
當時是沒有.
很強烈的.
信仰的根據.
是不可以娶外籍太太.
不過很多人都.
prefer都情願.
自己的子女娶一個.
同地方的.
女士.
例如.
阿布拉罕是要叫他的.
老僕人.
跟我回原地.

$^{601}$娶一個我當時.
老家的.
哈蘭的女子.
來嫁給.
爾薩.
是想不要外面的人.
到爾薩.
生了之後.
爾薩和利伯加頭都赤了.
因為爾素.
娶了兩個哈蘭女子.
黑人來的.
黑族.
於是搞到他們.
真的很慘.
不知道搞什麼.
令他們很苦惱.
不想的.
不過.
沒有說這樣是錯.
所以他娶了一個.
法魯比的.
女士.
叫阿薩.
阿薩那.
就是.
阿薩那屬於.
女神的意思.
名字都很難.
讀.
約瑟就出去.
巡行埃及地.
約瑟見.
埃及王法魯的時候.
他只是三十歲.
上帝真是厲害.
用了以往的日子.
去模造他.
變成一個合適的人.
他就開始坐在.

$^{641}$這個位置.
他和法魯面前出去的時候.
遍行埃及全地.
七個豐年之內地的出產.
是極為豐富.
多得很厲害.
約瑟就收斂了.
這七個豐年的糧食.
將糧食積存在.
各個城裡.
而各個周圍田地的糧食都積存在城裡.
因為在外面.
很容易被農民自己吃了.
沒有了.
放在城市裡.
約瑟第三點.
積蓄五穀甚多.
和海邊的沙無法計算.
因為穀不可勝數.
第三點.
要和大家講的是.
約瑟知道這件事.
他盡忠職守.
盡量去.
收納那些.
五穀.
令到他可以多起來.
是超過了.
他所求的很多.
而神真的祝福.
多到無法計算.
好像海邊的沙那樣多.
如果用海邊的沙.
你會發覺在聖經舊約.
用海邊的沙.
為了天神的星那些.
多到數不完.
一定夠用.
十分豐富.
請大家留意.

$^{681}$在平順的時候.
要儲備好.
不要等到.
要用的時候.
沒有儲備.
儲備好.
準備.
今日和大家分享到這裡.
你會知道.
神在他的手裡工作.
在神的手裡工作.
再從頭看一次.
有神的靈.
在他裡頭.
約瑟有神的靈在裡面.
所以連法老都知道.
是他,非他莫屬.
我們要尋求的是.
讓神的靈在我們生命裡.
坐在守衛.
坐在一個帶領司機.
領袖的位.
然後法老說我信任你.
我將我全個國家.
交託給你.
將我皇帝的東西.
全部給你.
你去做吧.
這些時候很多時候.
有個很出名的人.
嫁給他.
沒有命令.
一人之下萬人之上.
很多時候在這裡的時候.
我們會鬆毛鬆翼.
走錯路.
自己很驕傲很自大.
不過約瑟沒有.
他盡量去做.
將他要做的事情.

$^{721}$七個豐年的事情.
積蓄在五谷十分之多.
盡忠職守.
就是神的靈在他裡面.
我們一起禱告吧.
天護上帝.
我們讚美你.
想不到這個三十歲的年輕人.
沒有讀過甚麼書.
又沒有讀過管理學.
只不過是一個.
天之驕子.
在家裡生活.
沒有做家務.
不用出去無恙.
手真的很有亂.
路又不識.
但是我們的神.
仍然可以無做這個.
天真的小朋友.
無知的小朋友.
讓他經歷了.
做奴隸.
讓他經歷了被誣告.
成為罪犯.
然後將他提出來.
其實你的靈一直在他心裡.
主我們讚美你.
多謝你欣賞你自己無比的.
工作.
有時候我們經歷了.
一些很低沉的時候.
其實神都有心意在我們當中.
最重要的就是.
我們讓你的靈在我們心裡面.
作主.
於是我們就一路跟隨你.
何時你要提醒我們.
主我在這裡求你差遣我.
我等候主的使用.

$^{761}$正如藥室.
不知道怎樣.
不過他一路都是心裡跟隨你.
一叫他解夢.
他解到.
你的心有個寧靜的心.
時時都聆聽聖靈的聲音.
聆聽上帝的聲音.
多謝主.
讓我們開始學習.
很多時候在最苦.
最難的時候.
反而聽到神的聲音.
自己很好的時候.
有時候會聽不到神的聲音.
聽到自己高傲的聲音.
主求你自己幫助我們.
謙虛在你面前跟隨你.
以你做我們人生的主.
做我們的帶領者.
做我們的領袖.
聽我們祈禱.
奉耶穌基督名禱.
\newpage



\section{羅馬書 16:17-27}
\label{sec:031862}
\textbf{在福音中被堅固與合一}~~~~黃浩君傳道
\newline
\newline
連結: \href{http://www.church.com.hk/acms/content.asp?site=cdc\&op=show\&type=product\&code=031862\&layout=sermon}{\texttt{031862}} ~~~~ 語音日期: 2026-03-22
\newline
\newline
\hyperref[sec:031861]{\small{< < < PREV SERMON < < <}}
~
\hyperref[sec:index_chronic]{\small{[返順時目]}}
~
\hyperref[sec:031863]{\small{> > > NEXT SERMON > > >}}
\newline
\newline
羅馬書 16:17-27
\newline
\begin{longtable}{cl}
\hline
\hline
章節 & 經文 (和合本修訂版)\\
\hline
16:17 & \begin{tabularx}{0.7\textwidth}{X} 弟兄們，那些離間你們、使你們跌倒、違背所學之道的人，我勸你們要留意躲避他們。 \end{tabularx} \\ \\ \relax
16:18 & \begin{tabularx}{0.7\textwidth}{X} 因為這樣的人不服侍我們的主基督，只服侍自己的肚腹，用花言巧語誘惑老實人的心。 \end{tabularx} \\ \\ \relax
16:19 & \begin{tabularx}{0.7\textwidth}{X} 你們的順服已經傳於眾人，所以我為你們歡喜；但我願你們在善上聰明，在惡上愚拙。 \end{tabularx} \\ \\ \relax
16:20 & \begin{tabularx}{0.7\textwidth}{X} 那賜平安的神快要把撒但踐踏在你們腳下。願我們主耶穌基督的恩與你們同在！ \end{tabularx} \\ \\ \relax
16:21 & \begin{tabularx}{0.7\textwidth}{X} 我的同工提摩太，和我的親戚路求、耶孫、所西巴德，向你們問安。 \end{tabularx} \\ \\ \relax
16:22 & \begin{tabularx}{0.7\textwidth}{X} 我這代筆寫信的德提，在主裡向你們問安。 \end{tabularx} \\ \\ \relax
16:23 & \begin{tabularx}{0.7\textwidth}{X} 那接待我，也接待全教會的該猶，向你們問安。城裡的財務官以拉都和弟兄括土向你們問安。 \end{tabularx} \\ \\ \relax
16:24 & \begin{tabularx}{0.7\textwidth}{X} 願我們的主耶穌基督賜恩典給你們各位！阿們！ \end{tabularx} \\ \\ \relax
16:25 & \begin{tabularx}{0.7\textwidth}{X} 惟有神能照我所傳的福音和所講的耶穌基督，並照歷代以來隱藏的奧祕的啟示，堅固你們。 \end{tabularx} \\ \\ \relax
16:26 & \begin{tabularx}{0.7\textwidth}{X} 這奧祕如今顯示出來，而且按著永生神的命令，藉眾先知的書指示萬民，使他們因信而順服。 \end{tabularx} \\ \\ \relax
16:27 & \begin{tabularx}{0.7\textwidth}{X} 願榮耀，藉著耶穌基督，歸給獨一全智的神，直到永遠。阿們！ \end{tabularx} \\ \\
[1ex]
\hline
\hline
\end{longtable}
$^{1}$弟兄姊妹平安.
開始之先我們一起禱告.
七位天父我們都多謝你.
今日讓我們一齊回到教會.
一齊看羅馬書最後一段經文.
十六章十七至二十七節.
主要求你開我們的眼睛.
透過保羅的警告.
童宮的名單.
和最後的讚頌.
讓我們看到神是願意兼顧我們的神.
我們有求聖靈引領我們.
不單止聽到.
更是讓你在我們每個人心裡工作.
願榮耀歸給你.
我們禱告奉主明求.
阿們.
今日羅馬書最後一講.
大家記得羅馬書大概講過什麼嗎.
有人記得.
大家都不記得.
好不要緊.
大家回去聽回錄音.
如果我們很快地看一看.
羅馬書一開頭.
其實就是在說福音是什麼.
福音是藉著眾先知.
是舊約聖經的應許.
是要彰顯神的義.
目的是要讓神的名.
在萬國裡使人可以因信而信服.
先是猶太人後是希臘人.
然後我們看到全人類都犯了罪.
但是全人類都可以因順情義.
所以作為信徒.
我們的救恩是完美的.
不單止可以對罪死.
在主裡面活.
我們還可以靠著聖靈有新的生活.
在福音的前提下.

$^{41}$猶太人和愛邦人都會得救.
然後去到十二章.
因著福音的懸固.
我們都會有相應的行為改變.
就是心意更新而變化.
所以教會裡縱然有不同背景文化的人.
我們仍然可以彼此接納.
然後去到上一章.
即是十五章.
我們看到保羅宣教的目的和計劃.
以及他準備很想去造訪羅馬教會.
之後就是Jason上次提到的一大批.
民安的對象.
所以我們看到.
整本羅馬書的目的.
是因為保羅想去西班牙宣教.
所以希望羅馬教會.
成為一個福音的基地.
在這個大前提之下.
保羅必須要讓羅馬教會.
明白福音是給萬民的福音.
不是單單給猶太人.
所以他很清楚解釋福音是甚麼.
而另一方面.
當時羅馬教會的內部.
因為有猶太人和有外邦的信徒.
他們都有不同的文化.
有不同的飲食習慣.
所以充滿了張力.
所以保羅亦都希望鼓勵他們.
在基督恩典裡.
可以彼此接納.
剛才主席已經讀出了今日的經文.
羅馬書十六章十七至二十七字.
其實作為書信的結尾.
保羅大可以說聲拜拜.
然後一個祝福就完.
但我們看到經文很噁心地.
發出了一個嚴厲的警告.
然後列出了一大串的名單.

$^{81}$不過如果我們帶著羅馬書.
那兩個寫作目的來看這段經文的時候.
我們可能更加明白.
如果教會不合一.
如果信徒不懂分辨真偽的時候.
前面一至十五章所說的福音.
其實只是停留在紙上面.
亦都會阻礙到福音在世界的擴展.
所以讓我們看看這段經文.
看看我們有什麼需要學習的地方.
來到今天的經文.
保羅在十七節一開始就叫羅馬教會的弟兄姊妹.
要留意和要躲避兩類人.
留意和躲避都是社交的用語.
即是要在信徒群體裡識別這兩類人出來.
然後主動跟他們割席.
那這兩類人究竟是什麼人呢.
就是一些離間他們.
使他們跌倒違背所學之道的人.
離間即是令人不和爭執.
甚至造成一些敵對的關係.
導致結黨和分裂教會.
使他們跌倒違背所學之道的人.
就是那些誘惑人偏離真道.
進入謬誤和罪惡之道的絆腳石.
如果用聖經的術語.
這些人是在教一些阻礙人成長.
引人走入滅亡之路的假教師.
所以保羅說.
要在教會裡留意有沒有假教師的存在.
要避開他們.
不要和他們交往.
無可否認保羅的用意是要保護教會.
免得弟兄姊妹在真道上被他們玷污.
但是到了這裡我們就會問一個問題.
其實這群人是什麼人呢.
可能是前幾章說到一些堅持猶太傳統的信徒.
特別是之前說到他們不斷評斷其他人.
但都有可能是保羅沒有特定指到一群特定的人.
而是對教會作出一個預防性的警告.

$^{121}$但是無論如何.
作為教會.
我們都要不時用留意和躲避的態度.
防範假教師的滲入.
因為他們會引發紛爭.
但保羅又為什麼不直接.
叫他們和這群假教師討論神學.
辯明真理呢.
原來十八節說到假教師的目的.
根本不是追求真理.
相反他們立心不良.
十八節說到因為這樣的人不服事我們的主基督.
只服事自己的肚腹.
意思是說他們雖然自稱是基督的僕人.
但他們的服事對象只是他們的肚腹.
他們追求的不是弟兄姊妹之間一起生命成長.
而是追求自己的肉體,私慾,名聲甚至是利益.
試問和這些人.
討論真理有什麼用呢.
而且經文說到他們會用花言巧語.
誘惑老實人的心.
這是指這群假教師製造紛爭的人.
會用一些令人覺得很好聽奉承的話.
去騙一些入世未深的信徒.
和容易欺騙的信徒.
弟兄姊妹.
如果回到今天 什麼是服事自己的肚腹呢?.
可能就是將信仰變成一種消費.
帶著自己個人目的去服事的人.
以前聽不少一些前輩提到.
有些教會的弟兄姊妹甚至是教睦同工.
在搶做教會的領袖.
他們的目的未必是服事人.
而是享受那種話事的感覺.
甚至可能是希望在教會裡面.
建立他個人的人際網絡.
所以當我們服事的時候.
可能都要檢視一下自己.
我們是想主服事我們的肚腹.
還是我們真的想服事神.

$^{161}$如果每一個人都只是看自己的利益.
教會一定會被離間 被分裂.
到十九至二十節.
雖然每一個時代都會有假的教師欺騙信徒.
但經文說到信徒其實有兩個保障.
而這兩個保障.
第一個就是教會本身對福音的信服.
第二個則是神自己的應許.
十九節保羅說到羅馬教會對福音的信服.
其實已經眾所周知.
就像羅馬書一章一開始所說.
羅馬弟兄姊妹的信德傳遍了天下.
甚至遠在另一邊的保羅.
在十五章都讚賞他們充滿良善.
有各樣豐富的知識可以彼此勸戒.
羅馬教會已經成為宗教會的榜樣.
這件事也讓保羅歡喜.
所以保羅只是在提醒他們.
要在善上聰明 惡上愚處.
根據十二章打後的脈絡來看.
神眼裡的善不單止是道德上的良善.
還包括了教會裡的合一.
弟兄姊妹之間彼此接納的意味.
而惡就包括了弟兄姊妹之間彼此評斷.
最終使人跌倒的教導.
而最終當教會成為彼此合一.
互相勸戒的群體.
就可以預防假教師的滲入和欺瞞.
而面對假教師的滲入和離間.
保羅在二十節都說出教會的第二個保障.
就是神自己的應許.
二十節說到那次平安的神.
快要把撒旦踐踏在你們腳下.
在這裡我們看到離間和使人跌倒.
幕後的黑手就是撒旦.
撒旦即是敵對者.
祂會派出假教師離間你們.
使你們跌倒.
違背所學之道.
相對的是次平安的神.

$^{201}$因為祂重視教會的合一和和諧.
所以應許必然親自踐踏撒旦的工作.
所以信徒可以靠著神的恩惠得勝.
最後保羅祝福說.
願我們主耶穌基督的恩惠與你們同在.
在這裡想做一個小結.
羅馬書的寫作目的.
是希望猶太的信徒和外邦的信徒.
都可以彼此接納 在真道上站穩.
所以保羅在最後的結尾.
一連串的問安之後.
發出了一個提醒或警告.
要求弟兄姊妹要留意和躲避.
一些離間人.
叫人跌倒違背所學之道的人.
這些人不是服侍基督的.
相反只是服侍自己的肚腹.
他們用花言巧語誘惑老實人的心.
但最終因為教會對福音的信服.
亦因為神的應許.
教會必然被保守 不自分裂.
弟兄姊妹.
想問問現在這個世代.
會不會有人希望離間你們.
使你們跌倒違背所學之道?.
有沒有這些人?.
正常也有.
那怎麼辦?.
我相信羅馬教會.
已經成為了一個榜樣給我們看.
他們對福音信服.
而且充滿善良.
有各樣豐富的知識.
而最重要的是.
他們可以彼此勸戒 互相提醒.
我離開了母會之後.
和以前小組的弟兄姊妹.
仍然每個星期有網上的靈修小組.
我們會用一個我覺得比較紮實的靈修材料.
一起去靈修和代禱.

$^{241}$話說上星期去了營地.
回到家已經很累.
我在靈修 特別是網上的靈修小組.
就渾游障礙.
那晚主題其實是信心和神的工作關係.
去到結尾分享的時間.
有個姐妹突然問.
為什麼靈修資料說到沒有信心的時候.
神好像不會工作.
然後大家彼此分享到.
其實自己信心軟弱的時候.
仍然看到神的工作.
然後將榮耀歸給神.
那一刻我就體會到.
信服福音 彼此勸戒 互相提醒的重要.
縱然我也受過神學的訓練.
但仍然會有疏忽的時間.
本靈修資料其實是扎實的.
但可能也會出現表達上的誤導.
但大兄姐妹追求福音的心.
亦可以彼此勸戒互相提醒的態度.
讓我們更加見到真理.
更加經歷到神的榮耀.
窩福堂其中一個核心價值.
就是不住學習.
今年踏入了第五十週年.
願我們都可以繼續一同不住學習.
一同彼此勸戒 互相提醒.
在這個混亂的世代.
持守所學的道.
然後到二十一至二十四節.
在這裡我們看到八個人的問安.
和之前一至十六節的人物是不同的.
一至十六節是保羅對羅馬教會的問安對象.
而這八個人就是保羅的同工團隊.
一起向羅馬教會問安.
當中我們看到一個混血兒.
三個猶太人 四個外邦人.
八個同工有八個不同的背景.
我們可以看到福音打破了種族和階級的屏障.

$^{281}$讓教會可以有同一的心.
首先我們來看看二十一節.
這裡分別提到提摩泰 路求 耶穌.
和索西巴德的問安.
雖然保羅很簡單.
只是說提摩泰是我的同工.
但是我們知道十三卷的保羅書信裡.
有兩卷書是寫給他的.
另外提摩泰這個名字.
也出現在另外六本保羅書信裡的問安語裡.
可以反映到提摩泰在保羅的宣教工作.
是有重要的地位.
也是保羅的得力助手.
但是其實對一個傳統的猶太人來說.
提摩泰的出身一點都不風光.
在《史陀行傳》裡.
我們知道提摩泰的媽媽是猶太人.
爸爸是外邦人.
所以他是一個混血兒.
提摩泰當初連立約的記號「國禮」都沒有做.
理論上他是被傳統的猶太人輕視的人.
但是作為希伯來人的希伯來人.
即是保羅.
在這裡第一個提起提摩泰.
後面才提到他那三個猶太同胞.
反映了保羅對提摩泰的出身沒有特別的重視.
而這裡也貫徹了羅馬書第二章說到.
真正的猶太人.
即是在立約的群體的人.
是得到福音的人.
不在乎外表和肉體的國禮.
而是在乎內心和聖靈的引證.
然後我們看看下面那三個同胞.
保羅用我的親戚來形容他們.
表示其實他們就是猶太同胞.
在《使徒行傳》裡面.
路求可能就是安提阿的教師.
耶孫應該是在貼薩諾尼加的時候.
接待過保羅一行人的同工.
索西巴德在《使徒行傳》二十章裡面.

$^{321}$和保羅同行的索巴德.
三個人的共通點.
就是散在不同地方的猶太人.
所以我們看到.
在一個猶太人和外邦人.
都不可以同桌吃飯的年代.
保羅的同行團隊.
起碼有兩類不同文化的人.
然後我們看二十二節.
二十二節說到.
我這代筆寫信的得體在主裡.
向你們問安.
《新漢語.日本》這裡加了一個括號.
顯示了這是一個插句.
而這一句明顯就是得體.
在保羅的同意之下.
加插了進去的親自問安.
得體是保羅的書記.
他可能和羅馬教會一些人.
有一定的關係.
所以很想親口跟他們問安.
他也用了他的專業技能.
就是寫字.
參與了保羅對羅馬教會的教導工作.
到這裡可能會有人問.
為什麼保羅又不是不會寫字.
為什麼要有人代筆呢.
羅馬書究竟是不是保羅自己寫的.
首先我們要知道.
在羅馬世界裡面.
代筆是很普遍的.
因為有一個代筆的書記.
可以省去了找寫信的物品的麻煩.
而另一方面寫明是代筆.
反而證明了原作者是保羅.
因為保羅在另一些書卷裡面.
寫明了他的親筆.
這裡列明了代筆.
證明了羅馬書是得體.
記載了保羅對羅馬教會的口述教導.

$^{361}$就好像在以前的年代.
如果我們幫一些不懂得寫字的長輩寫信.
作者仍然是那位長輩.
我們只是代筆而已.
但這裡要留意的是.
我們看到得體可以突然間.
用第一人稱向羅馬教會問安.
顯示了在保羅眼裡.
得體其實和自己一樣.
都是主裡面的同工.
沒有什麼高低地位之分.
然後二十三節.
我們看到該由,爾拉都和豁圖.
該由是一個外邦人.
是保羅第一次去哥林多時為他受洗.
他富有和慷慨.
他不單止接待了保羅.
也接待了全教會的東道主.
爾拉是城裡面的財務官.
財務官在當時應該有一定的社會地位.
而後面的豁圖.
我們很少資料.
我們只知道他是一個弟兄.
還有這個名字通常是奴隸用的.
所以在保羅的同工團隊裡面.
包括了一個富有的外邦人.
一個應該身份地位頗高的財務官.
還有一個可能是奴隸的弟兄.
弟兄姊妹.
如果我們看一看這份名單.
裡面有猶太人,外邦人,混血兒.
有專業的書記,有錢的東道主.
有政府的財務官.
還有一個奴隸的弟兄.
當我們看到社會的高層和一個奴隸.
都可以在同一個團隊裡面.
一起向羅馬教會問安的時候.
我們就明白什麼叫做在福音裡面的合一.
當我們知道階級觀念,種族主義.
再不是福音的障礙的時候.

$^{401}$我們就展現在基督裡面的平等.
而這個也是神堅固教會的方式.
是對抗分裂最強的武器.
但反過來想.
有時候我們會不會無意間.
在教會裡面製造一些隱形階級出來呢?.
話說我以前在舞會.
這個教務同工要我強調的是在舞會.
見過一些信徒領袖在公開的場合裡面.
無理責備保安.
而保安其實當時是在做一些自己的分內事.
也見過一些弟兄姊妹很高姿態的.
指揮行政同工或者清潔姐姐.
當是自己手下這樣洗.
連一句麻煩都沒有.
轉過頭來就對著教務同工客客氣氣抱頭抱頸.
如果福音是令到種族,社會,階層.
可以兩下合二為一.
拆毀中間的牆的時候.
那一刻我就明白為什麼福音是停滯不前.
所以弟兄姊妹.
當我們經過教會的門口進來教會.
面對弟兄姊妹的時候.
可能我們都要檢視一下自己.
有沒有將自己在外面的職位.
我們的身份,我們的地位.
我們的階級觀念放進來.
破壞著基督的身體.
讓無論行政同工也好.
弟兄姊妹也好.
可以感受到自己其實是一個.
備受尊重的福音夥伴.
願神兼顧我們對抗世界的價值觀.
讓福音真真正正呈現在我們的群體當中.
雖然我們知識上知道有假教師.
我們也知道教會的合一.
教會是需要合一的.
但面對假教師的花言巧語.
面對自己的固有思想.
我們憑什麼可以在福音上站立得穩呢?.

$^{441}$保羅在二十五至二十七節說出了最終極的報章.
就是「唯有神可以兼顧我們」.
二十五節說到.
「唯有神能照我所傳的福音和所講的耶穌基督.
並照歷代以來隱藏的奧秘的啟示.
兼顧你們」.
意思就是面對假的教導.
面對教會內部分裂的挑戰.
神仍然有能力.
仍然願意保守我們不受瞞騙.
維持教會的合一.
而這個應許不是一個新的應許.
因為是照著我所傳的福音和所講的耶穌基督.
也照著歷代以來隱藏的奧秘的啟示顯明了出來.
在羅馬書裡.
保羅所傳的福音和耶穌基督.
就是說信徒因為耶穌的死.
可以和基督聯合.
得著一個新的生命.
所以縱然不同背景.
教會內仍然可以彼此接納的福音.
而且這個福音是歷代以來隱藏的奧秘的啟示.
意思就是說.
從前猶太人以為自己才是聖約的群體.
但現在因著基督的死 復活 升天.
舊音不再局限在猶太人身上.
而且延伸到外邦人身上.
信主的猶太人 信主的外邦人.
在神面前完全平等.
福音不再是猶太人的福音.
而是傳給萬民.
使萬民因信而信服的福音.
所以今天我們每一個信徒.
都肩負起這個福音的使命.
而在二十六節.
保羅說到這個隱藏的奧秘.
如今已經顯示出來了.
而且按著神的旨意.
不僅被猶太人看到.
還要被萬國萬民.

$^{481}$在萬國萬民之中啟示了出來.
這個啟示和舊約聖經完全一致.
福音被顯示出來.
目的不是要增加我們的知識.
而是要讓我們產生信服福音的心.
因為就像羅馬書第一章開頭所說.
「這福音是神從前直誦先知.
在聖經上應許的.
他的名(神的名).
在萬國裡使人因信而信服」.
這個就是福音的奧秘.
目的就是要讓萬民知道.
同時認識福音.
讓他們可以因信而信服.
先是猶太人.
後是希臘人.
而最終所有信徒.
都被納入舊恩的群體之中.
弟兄姊妹.
既然今天神已經將所有人.
在福音裡都平等的奧秘顯示了出來.
我們今天又願不願意信服這個真理.
去接納那些跟我們想法不同.
背景不同.
甚至性格都很不同的弟兄姊妹呢?.
而要接納這些不同的人的重點.
焦點就是要放在神那裡.
27節.
保羅將焦點完全轉向神.
以一個讚美作結.
願榮耀藉著耶穌基督.
歸給獨一全自的真神.
直到永遠.
阿們.
唯有神的智慧.
藉著主耶穌基督.
才可以成就一切.
弟兄姊妹.
教會自二千年前建立到現在.
一直都面對著不同的困難.

$^{521}$不同的挑戰.
對外.
二千年前教會會受假教師的影響.
同樣地.
我們今天都面對著不同的異端.
網上也有很多似是而非的資訊.
社會上也有不同的價值觀.
這些都在衝擊我們對信仰的根基.
對內.
教會初期.
面對猶太信徒和愛邦信徒的張力.
今天我們都可能因為背景的差異.
產生了間隔.
造成內耗.
甚至可能因為教會裡的不和.
造成教會的分裂.
但是教會仍然在這裡屹立了二千年.
面對這些如此現實.
如此真實的情景.
我們可能會對教會失望.
可能會對人絕望.
所以保羅在結尾.
都將我們的目光回到神.
唯有神可以兼顧你們.
唯願我們在任何時間.
都可以有被神兼顧的心.
神不是要我們自己死撐.
祂是透過福音去兼顧我們.
而一個被福音兼顧的人.
是不會輕易被假教導動搖.
也不會隨便破壞教會的合一.
我們站穩.
是要為了讓萬民都可以認識.
這位全智獨一的真神.
總結今天的經文.
保羅提醒了教會.
要小心假教師的迷惑.
導致教會的分裂.
教會要在真道上站立得穩.
就要繼續信服福音.

$^{561}$同心抵抗一切破壞教會的攻擊.
同時我邀請了大家.
看看身邊的弟兄姊妹.
打聲Hello都可以.
我們可能有不同的背景.
有不同的地位.
可能有些人已經退休.
有些人還在社會上拼搏.
但我們還是被同一個福音拯救的群體.
這個世界很想我們分裂.
但神就差派了我們在地上.
成為福音的見證人.
我們可能會面對不同的困難.
不同的挑戰.
但我們仍然可以將目光放在那一位.
願意也有能力.
兼顧我們獨一全智的真神上.
所以願我們也可以成為一個.
在善上聰明.
在惡上愚蠢.
合一的群體.
然後將一切榮耀歸給那一位.
讓萬民因信而信服的神.
我們一起禱告.
天父唯有你可以兼顧我們.
多謝你透過經文.
讓我們知道福音這個關乎萬民的奧秘.
今天我們看到縱然有一班.
背景完全不同的同工.
都可以在你裡合一工作.
願這份見證成為我們的力量.
主要我們又求你施恩給我們.
當我們面對假教道.
我們教會的不和的時候.
提醒我們.
你是那位獨一全智的真神.
你是那位有大能.
也願意兼顧我們的神.
求你帶領我們成為你在世上的見證.
然後將榮耀歸還給我們.

$^{601}$直到永遠.
阿們.
\newpage



\section{哥林多前書 13:4-7}
\label{sec:031863}
\textbf{重塑生命系列 (7) 妒忌與愛}~~~~張健庭牧師博士
\newline
\newline
連結: \href{http://www.church.com.hk/acms/content.asp?site=cdc\&op=show\&type=product\&code=031863\&layout=sermon}{\texttt{031863}} ~~~~ 語音日期: 2026-03-22
\newline
\newline
\hyperref[sec:031862]{\small{< < < PREV SERMON < < <}}
~
\hyperref[sec:index_chronic]{\small{[返順時目]}}
~
\hyperref[sec:031864]{\small{> > > NEXT SERMON > > >}}
\newline
\newline
哥林多前書 13:4-7
\newline
\begin{longtable}{cl}
\hline
\hline
章節 & 經文 (和合本修訂版)\\
\hline
13:4 & \begin{tabularx}{0.7\textwidth}{X} 愛是恆久忍耐；又有恩慈；愛是不嫉妒；愛是不自誇，不張狂， \end{tabularx} \\ \\ \relax
13:5 & \begin{tabularx}{0.7\textwidth}{X} 不做害羞的事，不求自己的益處，不輕易發怒，不計算人的惡， \end{tabularx} \\ \\ \relax
13:6 & \begin{tabularx}{0.7\textwidth}{X} 不喜歡不義，只喜歡真理； \end{tabularx} \\ \\ \relax
13:7 & \begin{tabularx}{0.7\textwidth}{X} 凡事包容，凡事相信，凡事盼望，凡事忍耐。 \end{tabularx} \\ \\
[1ex]
\hline
\hline
\end{longtable}
$^{1}$好,請姐妹平安,我們再一次低頭祈禱.
或者對於愛的一些經歷.
其實都是扭曲的經歷.
我們願意在你的愛當中來對照的時候.
我們更加明白你的心意.
也都更加能夠活出.
你創造我們所應該擁有的那種.
好像神這樣去愛人的時候.
我們學習怎樣來愛人和愛神.
我們這樣禱告奉耶穌基督的名祈求.
阿們.
好,今天我們就講窒道和愛.
其實窒道這樣的情感.
或者窒道這種在法治人心的一種慾望.
是很少來當眾講的.
你不知道有沒有發覺.
我們有講憤怒.
遲一些又會講懶惰.
之前有講過貪食.
或者有講其他.
即是驕傲.
那些都有時候都會.
很多時候有時候在公眾的場合.
或者和朋友聊天之間都會講的.
但是我們就在交流之中.
很少去提你是否窒道某個人.
或者某某某.
你知不知道我看那個人很不順眼.
其實我想我是在窒道他.
有沒有這樣講的呢.
只有在我經驗當中沒有聽過.
我自己也沒有這樣講過.
有時候我都會說我都很驕傲.
或者有時候說我都很懶惰.
有時候我想這個人真是.
他這麼貪心真的不行.
我們都會講這些.
因為似乎在公眾的領域當中.
這些是比較可以見得到.
有時候都不道理不認.

$^{41}$那我們就講.
但是窒道為什麼這麼少講呢.
其實就是因為他在公眾的領域當中.
有很多時候都很懂得躲起來.
譬如有時候出於窒道而有.
你對於一個人的一種態度.
你收起來.
別人只是看到你對他有這樣的態度.
其實你可能不知道你是窒道而來的.
甚至你自己都可能.
不會想自己因為窒道而有這樣的反應和表現.
這個就是窒道那個狡猾之處.
窒道很多時候其實是收起自己.
可以用第二個面具出來掩飾自己.
所以今天我們一起來思想窒道的時候.
都讓我們自己好像有些心靈的檢查.
因為坦白講.
我們自己對自己的檢查.
才會知道自己有沒有窒道.
可能你身邊的人都未必會知道.
一開始的時候.
不知道你有沒有想過.
為什麼會有窒道而來呢.
有沒有想過.
為什麼一開始你對於一些人.
當然是人.
會有窒道呢.
有沒有想過.
窒道其實很多時候是從比較而來的.
窒道是從比較而來的.
聖經很出名的一個.
第一個殺人故事.
就是因為窒道而殺人.
那個是窒道殺人案.
是講兩兄弟殺人.
一開始的時候就是戈隱和阿伯.
戈隱是阿伯.
第二節開始.
又申了戈隱的兄弟阿伯.
阿伯是牧羊的.

$^{81}$戈隱是種地的.
有一日戈隱拿地裡的出產.
為貢物獻給耶和華.
阿伯也將他羊群中投生的羊.
和羊的脂油獻上.
耶和華看中了阿伯和他的貢物.
只是看不中戈隱和他的貢物.
戈隱就大大地發怒.
便了臉色.
耶和華對戈隱說.
你為什麼發怒呢.
你為什麼便了臉色.
你若行得好豈不蒙虐難.
你若行得不好.
罪就伏在你門前.
他必戀慕你.
你卻要制伏他.
戈隱和他兄弟阿伯說話.
二人正在田間.
戈隱起來.
打死他兄弟阿伯.
把他殺了.
耶和華對戈隱說.
你兄弟阿伯在那裡.
他說我不知道.
我何時看守我兄弟的呢.
後面刑罰的部分我就不讀了.
大家都知道這個故事.
這個故事其實有幾部分都很特別.
不知道你們有沒有想過.
第一部分其實就是.
他們已經被神趕出了伊甸園.
他們在伊甸園很明顯.
每天都可以見到神.
神每天都與他們同行和他們見面.
這是在伊甸園的日常生活的一部分.
但是他們被趕出伊甸園之後.
很明顯失去了每天見神的部分.
所以絕對不是經常可以見到神.
所以他們在伊甸園沒有提到獻祭.

$^{121}$沒有提到獻祭.
伊甸園沒有什麼蛛絲馬跡.
留下給我們看到他們在伊甸園做的事.
是有獻祭這樣的行動.
但是趕出了伊甸園他們就要獻祭.
我想其中一個很真實的需要.
發覺是離開了伊甸園的需要.
因為我們發覺被趕來伊甸園之後.
人就從神開始沒有辦法經常見到神.
好像藉著獻祭本身與神保持一個交流.
以致可以繼續與神有一種關係.
這是獻祭的主要用意.
獻祭本身是與神有一種交流.
當然獻祭本身後期的發展是與罪有關.
你看到在摩西五經.
特別是後期說到出埃及記.
利未記說到最低限度有兩種祭.
都是與罪有關.
因為要贖罪.
但是這裡所說的該恩亞伯那種獻祭.
就沒有很具體說贖罪祭那些.
似乎是一種對神的尊榮.
我用的例子可能不太貼切.
有時候我們中國人都會去探望一些長輩.
要帶一些手信去.
有時候很久都沒有見過一個長輩.
想與他保持一種關係.
你就去探望他帶一些手信去.
當然不一定是這樣的.
聖經與聖經其實都沒有說得很清楚.
為什麼他們那時候又會獻祭.
但是我覺得從那些蛛絲馬跡.
很直覺地就是這個原因.
因為他們現在已經不是經常見到神.
不是經常見到神的時候.
就與神保持一種關係.
他們也將祭物可以獻給神.
就好像對神的一種尊崇.
一種尊榮.
所以其實是沒有說神只喜歡羊.

$^{161}$沒有說神只喜歡羊.
在後來神與光陰的對話都不是說.
你吾蒙月立,因為你沒有獻羊給我.
因為我是吃肉的,不吃齋的.
他不是這樣說的.
所以根本就沒有說獻什麼是重點.
如果用我剛才說的.
去探望長輩,帶一些手信.
去尊榮你與他保持關係的人的時候.
其實帶什麼是沒有所謂的.
在某個角度來說.
所以那個人就帶他去種田.
阿伯就帶他去養羊.
就帶他的羊.
但是這個天祭的舉動本身就引來了一個殺機.
這個殺機就是說.
那個人就很生氣.
因為神不月立他的祭物.
而月立了他弟弟的祭物.
神不月立他的祭物.
神月立他弟弟的祭物.
他就很生氣.
後來就生氣到一個地步.
後來就出現了殺機.
殺了他弟弟.
這個生氣本身是怎麼來的呢?.
很明顯.
是怎麼來的呢?.
是因為他有比較.
他比較過.
神為什麼要你不要我.
一經這樣比較之後.
他就生氣了那個.
神要他的弟弟.
神月立那個.
純粹是有一個比較而來的.
我也看過有一個心理實驗很出名的.
是這樣的.
一開始就放一個小孩.
其實有兩個小孩.

$^{201}$在兩個不同的房間裡.
各自在房間裡都有幾樣玩具.
給他玩的.
那個小孩其實也很有安全感的.
他一個人在玩.
但是他也玩得很開心的.
在玩他的玩具.
有兩個小孩都玩得很開心的.
但是大家就隔著一個窗.
那個窗有個窗簾遮住了.
大家都不知道.
隔壁房間的小孩在玩什麼.
後來大家都玩得很開心的時候.
那個研究員就拉起了窗簾.
拉起窗簾之後.
大家就看到大家有什麼玩.
其中一個小孩就很不開心了.
為什麼呢?.
因為他的玩具只有一輛很簡單的玩具車在推來推去.
但是他看到旁邊又有遊戲機又有LEGO.
他就很不開心.
他原本是玩得很開心的.
一個人的時候.
但是一比較之下.
他原本玩得很開心的東西.
就再不能夠令他開心了.
於是就反而很不開心了.
這個實驗本身是一個很反映人性的真實.
完全他開心和不開心.
和那個東西好不好玩.
是完全沒有關係的.
一開始的時候.
他根本也是玩得很開心.
但是那個問題就是出於有一種比較之下.
就開始對於自己所擁有的東西就不滿足.
而且不滿足到一個地步.
就是甚至乎是遷怒於你覺得那個比你好的人.
就好像那個人一樣.
他遷怒於他的弟弟.
但是這個故事本身有另外一點也是很特別的.

$^{241}$聖經那裡其實當人離開了伊甸園.
剛才也說過了.
不是經常見到神的.
但是這次就說神主動出現.
去和戈隱說話.
不是和阿伯說話.
如果神和戈隱說.
跟著和阿伯說.
阿伯小心一點.
你大哥可能想殺你.
可能會更有效.
你想一下會不會.
跟戈隱說.
說完之後戈隱仍然要殺他弟弟.
但如果神不跟戈隱說.
跟阿伯說就好了.
跟阿伯說.
不如你小心一點.
你大哥起了殺機.
他想殺你.
那可能阿伯不用死.
不過這是我編出來的.
但是故事本身就不是這樣發生.
但是這個故事發生是很有趣的.
就是說到.
一開始剛才我也說過了.
為什麼要獻祭.
因為他們離開了伊甸園根本見不到神.
他們想和神保持一個關係.
所以獻祭本身也是要討神喜悅.
所謂神悅納和不神悅納.
他們很在意.
現在就可以了.
什麼意思.
現在神這麼大面子.
神親自主動出現和我聊天.
那不就比獻祭還要好.
你說更厲害.
比獻祭還要好.
你想一想.

$^{281}$阿伯雖然說神悅納祭.
阿伯從始至終都沒有見到神.
都沒有.
現在該隱見了.
那該隱見到神的時候.
神跟他說什麼呢.
神就跟該隱說.
你為什麼發怒呢.
為什麼變了臉.
首先神就先問問題.
然後他再問.
你若行得好氣不蒙悅納.
神問了一個問題.
其實神完全知道該隱為什麼會發怒.
神知道.
但是神就要由一個發問的開始.
因為通常在窒道的情況之下.
我們是自己不知道的.
其實發展到後期的時候.
你當然知道.
因為可能你看到這個人的時候.
很不值得他的時候.
你當然知道.
但是在初初發展的時候.
很多時候是不知道的.
初初的時候你都不知道.
慢慢去發展你對於這個人.
有一種窒道的心態.
初初是不發覺的.
但是一開始的時候.
所以神的發問就很有意思.
神想這個窒道本身.
現在還沒有開花結果.
現在就是殺了種子的時候.
神就想.
你想一下為什麼你要這樣.
因為窒道很多時候.
等它開花結果的時候.
很難蒙的.
你有沒有在後院種過薄荷.

$^{321}$你沒有嗎.
我貪得意在後院又種了薄荷.
就後悔了.
因為薄荷你種一個盆子都還可以.
你種在地上.
山到處都是.
蚊子怎麼也有.
蚊子完又有.
其實窒道是這樣.
一開始的時候.
你不發覺它.
你小的時候.
你蚊它容易.
但是如果發展到後來的時候.
你對於人的關係.
對於他整個的看法.
樣樣東西你都覺得.
很不喜歡它.
或者很羨慕窒道它的時候.
是很難.
你的心裡面有很多苦毒.
就已經一切變成果了.
很難搞.
好像有些薄荷掉在地上.
所以神一開始就先用一個問題.
來問他.
問他.
你知不知道為什麼你要這樣.
你想一下你自己.
你用不著窒道人.
你如果行得好.
那你就蒙越納.
其實這個也是一個反面建議.
神的意思就是說.
你用得著蒙越納嗎.
你下一次行好一點.
那我就越納你.
其實神是一個反面建議給他的.
你想一下.
但是他在那一刻就完全聽不到.

$^{361}$完全聽不到.
於是他就.
但是神就再說.
祂說.
罪就伏在你門前.
他必戀慕你.
你卻要制伏他.
罪就伏在你門前.
他必戀慕你.
你卻要制伏他.
這裡說到.
這是一個很特別的一種形容.
這裡說到.
我相信這裡所說的罪.
不是一個普通籠統說的罪.
特別是對應該人.
在他心裡所產生的那種嫉妒.
合理吧.
因為我們現在上文下理都是在說.
他因為嫉妒後來又殺人.
所以這裡所說的罪.
可能其實不是說普通的罪.
是在說嫉妒這個罪.
但是這段聖經形容他嫉妒這個罪.
就很有意思.
首先.
他是在說.
好像躲在你門前想偷襲你.
去殺死你.
吃掉你那些捕獵者.
你就是那些.
如果外面有一隻蠻鷹.
你有一隻小雞.
你小心.
他已經伏在你門前.
他會隨時將你捕獵.
這是第一.
形容到嫉妒本身.
是很有攻擊性的.
他好像一個很有技巧.

$^{401}$或者很有經驗的獵食者.
等待著獵物出現.
他躲在你門前.
當然這裡所說的象徵性.
是說他的心靈的門口.
接著他再說.
他必戀慕你.
你卻要制服他.
這裡說到他必戀慕你.
又是用另一個字去形容.
就是說在真言那裡.
說到那些淫婦.
專門勾引那些.
可能心智不堅定的年輕人.
勾引他們進入妓院.
然後他必戀慕你.
進入妓院.
他跟你發生一些不應該有的關係.
然後將你制服.
令你無法自拔.
這裡說到質度.
有另一個性質就是這樣.
第一個性質.
他好像一個捕獵者.
隨時等候機會.
你一不留神的時候.
他就撲出來.
第二個性質就是.
你如果不去抗拒他.
是他要制服你.
他好像裡面的人.
私慾引誘你心繞繞.
然後你一上了當.
就像釣魚一樣上鉤.
就上了鉤.
然後你就被他釣了.
所以在這段對話當中.
神是很知道人性的.
當然那個人就無一句聽得入口.
後面就真的被質度這個罪.

$^{441}$控制了他的心.
然後就殺人了.
那些情殺案.
很多犯罪的人.
是很聰明.
在這個世界的標準來說.
他不應該犯罪.
有很多事情是很大好前途的.
但是都會發生.
因為質度本身將整個人失去理智.
而推向懸崖.
我看過一單新聞.
那個都是令我覺得很奇怪.
所以我看了headline.
然後就看了內容標題是怎樣.
那個新聞是真事.
因為是新聞.
我幾年前看的.
有些細節我都找不出來.
是說有一單情殺案.
那個罪犯是一個NASA的女太空人.
是一個NASA的女太空人.
為了要殺她的情敵.
我不知道是情敵還是勾引她的丈夫離婚.
我不記得了.
因為幾年前看的.
那個細節令我很驚訝.
或者令我很吸引注意.
就是說開頭那些警察是沒有懷疑是她殺的.
因為好像是一個.
沒有可能的一個殺人的操作.
其中一個沒有可能.
就是在時間上沒有可能.
因為那個人和她住的地方是距離了十幾個小時的車程.
距離了十幾個小時.
就是離開很遠的不同的城市.
如果她坐飛機去殺她.
有那些航空紀錄的.
那她一定不能坐飛機.
如果是她殺的.

$^{481}$只有一個可能性就是開車.
是不是.
如果開車去殺她的時候.
那裡是走十幾個小時的車程.
然後她在案發之前.
有些人有證明她仍然在那個城市.
她沒有離開那個城市.
然後案發之後.
來回了三十幾個小時之後其實一天多一點.
然後她又出現又上班.
又在同一個城市.
所以一開始警察查是沒有證據是她會殺人.
但是後來經過了.
可能有些百問一疏.
可能在中間那裡.
她有入過油的時候.
其實她都盡量不讓人拍到.
但是拍到一點的時候.
就拍到是她.
後來一邊查一邊查.
終於證實了原來是那個女太空人殺了她的情敵.
就發覺原來她犯案是很聰明的.
她是計算了每一個點.
她在哪裡可以避開.
不會有那些被人拍到的.
而且在這十幾個小時.
她是預備了不睡覺.
帶了尿片.
不下車.
不睡覺.
在車上可能只是喝一點水.
整個過程是非常精密.
我想說的是什麼呢.
就是說其實如果是一個有理智的人.
根本是一個.
你可以想像得到.
不是被這個.
好像這裡所說.
好像是疾導將她亂冒而制服了.
根本完全.

$^{521}$那些人訪問回來.
完全都不像是她那樣.
你想想.
做到一個女太空人在美國來說.
各方面是很優秀的.
無論是情傷還是智傷都很優秀的.
你上到.
你對人處事.
怎麼會亂發脾氣.
去做太空人.
找死嗎.
你上到太空那裡打架.
怎麼搞.
不會有這樣的.
一定是情傷智傷各方面.
是很優秀的一個人.
但是就在那一刻.
疾導本身.
就將她完全的力度去控制.
控制到一個地步.
殺人本身.
她都是完全的.
是深思熟慮計劃了一切.
但是只是為了因為疾導.
而將她控制了.
而去殺另外一個人.
這個就是聖經所說的可怕.
疾導一開始的時候不是這樣的.
我相信也沒有人會想像得到.
她會因為疾導而這樣前途盡毀.
去殺另外一個人.
而在另外一個.
聖經那裡說到疾導.
去殺人的例子.
就是大衛和素羅之間那種關係.
一開始的時候.
大衛其實是很效忠素羅的.
其實也不是一開始.
由始至終大衛都是很效忠素羅.
但是你就知道素羅.

$^{561}$是要將她非殺不可.
就是置諸死地而後快的.
是不是.
但是你想想原本來說.
起殺機的原因是什麼呢.
是不是因為大衛要搶王位呢.
大家都知道不是.
是不是因為大衛做了什麼不忠的事.
以致素羅要殺她呢.
大家都知道不是.
一開始的時候.
只是一個很簡單的.
一些人唱歌.
就是說大衛殺了哥利亞很厲害.
於是她就聽到那些婦人唱.
素羅殺死千千.
大衛殺死萬萬.
素羅就.
素羅甚發怒.
不起悅者說.
就說將萬萬歸大衛.
千千歸我.
只剩下王位沒有給她.
是不是.
這個就是那個殺機的那個開始.
很雞毛蒜皮的.
你想一下真的很雞毛蒜皮.
大衛都不止一次.
向素羅表忠過.
就是說有時候那些皇帝多疑.
怕人會搶她的王位.
但是大衛都不止一次表忠過.
說她不會搶她的王位.
連她的兒子約拿丹.
都站在大衛那邊幫她說話.
就是.
所以在客觀的因素.
根本就完全不成立的.
大衛要搶王位這個因素.
完全不成立.

$^{601}$但是那個.
重點是.
為什麼素羅會後來好像鬼上身那樣.
就是要殺大衛.
其實聖經都有說.
神.
從神.
次日.
素羅開始有嫉妒的時候.
就真的好像鬼上身那樣.
他說.
從神那裡來了.
惡魔大大.
乾在素羅身上.
他就在家中胡言亂語.
大衛照常彈琴.
素羅手裡拿著槍.
素羅把槍一輪.
心裡說.
我要將大衛刺透.
釘在牆上.
大衛躲避他兩次.
素羅懼怕大衛.
因為要離開自己.
與大衛同在.
這個.
有些人都說.
他可能真的鬼上身都不奇怪的.
鬼上身要殺人都不奇怪的.
但是.
我不用這個角度去說.
就是從窒道這個角度去說.
就是.
你想想.
為什麼素羅一定要殺大衛呢?.
在這個窒道的背後.
就是他看到.
大衛就是成為他競爭的對手.
是不是?.
在一個正常的關係當中.

$^{641}$人與人之間又可以有很多不同的結連的.
就是.
在大衛和素羅的關係當中.
其中一個正常的結連.
一個是皇帝.
一個是將軍.
是不是?.
一個是.
一個是他的主人.
另外一個就是他家裡的隨僕.
幫他彈琴.
是不是?.
有很多不同的關係的.
在一個家庭當中.
也都有很多不同的關係的.
但是.
當窒道來臨的時候.
那些關係就瓦解了.
瓦解成一個什麼樣的關係呢?.
就是競爭對手的關係.
本來大衛.
如果是一個主.
和一個將軍的關係的時候.
是可以幫助到素羅的.
本來大衛如果是一個岳父.
和一個女婿的關係.
也是一個.
嘩!家裡只有一個女婿.
彼得上一個兒子還要好.
你也聽過了.
一個好女婿彼得上一個好兒子.
是不是?.
這也是一個很好的關係.
但是當窒道去侵入人的心的時候.
一切的關係.
正常的關係都瓦解了.
剩下來的.
只有一個競爭對手的關係.
是不是?.
真的.

$^{681}$有時候我們看到.
那些我們窒道的人.
本來可以好好和他們建立關係的.
本來可以好好合作的.
本來可以好好互相欣賞.
一起慶祝.
一起開心的.
一切關係都瓦解了.
你就看他.
你很厲害嗎?.
我一定要比你厲害.
是不是?.
一切的關係都瓦解了.
變成了競爭的關係.
而在這種競爭的關係之下.
是很痛苦的.
因為.
你知道競爭.
經常都要將自己拉緊的.
好像一條橡筋一樣.
你競爭.
你就要拉緊橡筋.
隨時準備射的.
是不是?.
是很痛苦的.
但是在人際關係當中.
其實很多關係.
你不需要拉得那麼緊的.
在一個愛的關係當中.
一個愛婦和女婿的關係.
很開心的.
在一個皇帝和將軍的關係.
打勝仗也很開心的.
是不是?.
很多的關係.
其實不是很慶祝的.
但是在競爭.
在窒道之下.
一切的關係就變成了競爭.
你想一下.

$^{721}$有時候.
如果教會當中有窒道的時候.
一起事奉的時候.
大家都看對方為競爭的對手的時候.
是很痛苦的.
在你工作的環境也是一樣.
你知道有人窒道你的時候.
有時候你都不.
你本來大家一起談.
大家都想將事情做好.
是很開心的.
但是你根本.
你都不.
你怕做得多的時候.
你就提防影射.
在這個世界.
你做得多.
別人不喜歡你.
你又怕別人窒道.
很多事情就變得很複雜.
因為窒道本身.
將這一切的關係都變了窒.
最終都變成了競爭的對手.
而窒道的根源.
我們想一下.
很多時候.
是出於我們對神的不滿.
是一個人的不滿.
但是在這個故事當中.
主人所說的主人的回答.
是很有意思的.
他說.
朋友我沒有虧待過你.
因為我做好人.
你就眼紅了嗎?.
這句話很有意思.
首先他說.
朋友我沒有虧待過你.
因為我做好人.
你就眼紅了嗎?.

$^{761}$而事實上.
很多時候窒道的根源.
本身去到最後.
我們是不滿神的.
你有沒有想過?.
是不滿神的.
為什麼神給他這麼好.
給我沒有的.
是不是?.
表面上我們是有一個具體的對象.
我們是窒到他.
但是在背後深層.
那個心理的不平衡失衡.
去到一個狀態.
其實最後的手指指向神.
為什麼神不給.
事實上.
去後面.
我都會說一點.
其中一個窒道.
其實是一種心理不平衡的狀況.
因為我們看神.
為什麼會對我不好.
其實神是不是對我們不好.
當然不是.
但是窒道是覺得.
神都是對別人.
對別人好一點.
因為影射之後.
就是對我不好.
那就覺得對你不好.
就是.
這種就是窒道的原因.
事實上.
這是一個很危險的.
大兄姐妹.
窒道之所以為.
Seven deadly sins.
其中一個就是.
除了他將人性當中那些.

$^{801}$令到很多的關係.
令到人失去了理性.
好像做了一些很愚蠢的事.
來破壞.
來傷害別人之外.
亦都是破壞了.
人怎樣去看神.
我都講過一個例子是真實.
那時候我年輕的時候.
那時候還在讀中學.
我有一個弟兄.
就是同學.
那時候他還沒有做基督徒.
有時候我們一起去教會.
他去不去都有去教會.
我有去教會的時候.
有一次我完了之後.
一起坐地鐵回家.
跟他聊聊天.
我就問他.
你覺得信仰怎樣.
他講到一件事.
他說.
我有一件事一直都很不開心.
或者很覺得不公平.
我說什麼不公平.
我們大家一起讀書.
大家上同一間學校.
他說.
我覺得.
為什麼神要給一些人這麼有錢.
為什麼我又沒有這麼有錢.
他不是說我這麼有錢.
他說為什麼我這麼窮.
但中學生來的.
但可能其實不是中學生都有這樣的觀念.
可能成年人更加有.
他說因為這件事.
令我對於信仰很多東西.
我都不知道應不應該接受.

$^{841}$這個問題本身是需要去處理的.
事實上.
但是背後.
我講這個例子.
講到很真實地.
其實在他心目中.
因為有些有錢的朋友或者同學.
令到他心緊於懷.
而緊於懷到一個地步.
他就覺得神是不好的.
他又不是窮到很窮那種.
都是小康家庭那種.
我們一起在學校讀書.
又不是沒有飯吃沒有書讀那種.
但是.
因為對於一些的嫉妒.
在他心目中的時候.
他最終投射到.
神是不是不公平.
神是不是不好.
神是不是像他所說對我那麼好.
這個就是嫉妒其中一個問題.
嫉妒其中一個問題.
如果你不是那麼介懷這些的時候.
你看神本身.
其實你會發覺神有很多恩典.
神是很好的.
但是當嫉妒滲入你的心裡的時候.
剛才講的嫉妒將人的關係破壞.
不僅是人.
人和神的關係都破壞.
你想神都不用另外一種想法去想神.
而當這樣的時候.
嫉妒就暴露出你內心那種不安.
你就沒什麼安全感.
因為你就發覺.
你都不知道神是不是這樣對我的時候.
你最終極對於人生你自我的把握.
其實你是建基於神那裡的.
你最終極.

$^{881}$你最有把握.
你最知道你是被接納的.
從哪裡而來?.
是這個世界而來?.
這個世界當然不能給你.
有時候你做少一點.
也會被人看不起你.
這個世界的標準不斷地變化.
這個世界也用你的成敗得失來衡量你是一個怎樣的人.
你和人的關係也會不斷地變化.
這個世界根本給不了你一個終極的穩妥.
你是完全被接納的.
因為所有的接納.
會有某個程度上的環境因素.
而也會改變.
最終人的感覺都完全被接納.
完全被愛.
無論如何.
我都可以穩妥的.
就是在神那裡而來.
而問題就是.
如果你在一個質度的關係當中去看.
甚至乎這樣去理解神的時候.
你自己的身份就出現了問題.
在浪子比喻其實就是一個這樣的問題.
那個小兒子大家知道他怎麼壞的.
花天酒地就浪費了錢.
然後就回去.
父親就不是太的.
甚至很好.
就接納回去.
那部分我就不說了.
大家都知道.
但是我就想說那個大兒子的部分.
那個大兒子.
就因為那個小兒子回來就很不開心.
他就說.
他就問究竟發生什麼事.
他就知道那個小兒子回來.
父親和他一家人在幫他慶祝.

$^{921}$他就生氣.
大兒子卻生氣.
不肯進去.
生氣到.
我現在不參加這個派對.
現在翻臉了.
他父親就出來勸他.
這個父親出來勸他.
其實對比於那個父親.
在那個村口等小兒子兩個舉動.
其實是一個對稱的.
他對父親說.
我服侍你這麼多年.
從來沒有違背過你的命.
你並沒有給我一隻山羊羔.
叫我和朋友一同快樂.
但你這個兒子和昌記吞盡了你的產業.
他一來了.
你倒為他宰了肥牛族.
父親對他說.
兒啊.
你常和我同在.
我一切所有的都是你的.
只是你這個兄弟是死而復活.
失而又得的.
所以我們未當歡你快樂.
你看到嗎.
所以那個父親去接受那個小兒子.
那個小兒子願意悔改的時候.
那個大兒子是非常憤怒.
憤怒其中一個原因.
他覺得你為何對他這麼好.
他不值得對他這麼好.
你有沒有對我這麼好.
其實他未問這句.
他都在口順叮.
他沒有問出來.
他首先說你為何對他這麼好.
他做了那麼壞事.
花天酒地.

$^{961}$和昌記吞盡了你的錢.
他只不過沒有問.
你有沒有對我這麼好.
他沒有問.
但其實他的意思.
所以我現在不進去.
其實他的意思是這樣.
但其實在整個的關係當中.
反映了他對於自己的位置有一種不安.
你有沒有對我這麼好.
你好像對他好一點.
對比之下你好像沒有對我這麼好.
而就是因為一種這樣的情況之下.
窒道將他們兩兄弟的關係拆毀破壞了.
但其實窒道也將大兒子.
反映了他對於自己.
究竟是不是大兒子.
可以承受家產.
會不會老爸偏心.
再分不是嗎.
他已經分了.
很多事情會不會出來.
因為他對於自己的身份.
顯得毫無把握.
對於神那種.
你不能夠在神那裡.
有一個完全覺得.
神都會將一切最好的給你的時候.
一比較.
你覺得神可能會有些好東西不給你.
給了別人的時候.
你對於自己的身份開始懷疑.
究竟是不是神真的好像.
他所說的這麼好的對我呢.
就好像這個父親.
是不是他好像說得這麼好的對我呢.
是不是.
而我們不需要這樣去看.
窒道就是這樣將人的心去蠶食.
但是福音其實是可以將我們醫治.

$^{1001}$聖經所說.
愛是恆久忍耐又有恩慈.
愛是不窒道不自掛不張狂.
就是這樣的意思.
愛是不會窒道.
其中一個原因.
因為愛是一種.
將自己來去給別人的.
一個這樣的一個付出.
就是一個將自己去給別人的.
一個給出來的.
一個這樣的.
一個這樣來去發生.
而神就是愛.
在神的愛當中.
我們知道神就是將他自己給出來的時候.
我們就在愛當中.
我們已經收到神.
也都收到神的愛的時候.
我們就不需要窒道.
所以在加拉太書那裡.
第五章二十六節.
也都提醒我們.
不要貪圖虛名.
彼此惹氣.
互相窒道.
其中一個窒道的原因.
就是看錯了我們追求的東西.
你看錯了神.
為什麼將一些好東西.
你覺得給了別人不給你.
其中一個原因.
可能是你看錯了你追求的東西.
正如我剛才說的那個例子.
那個朋友.
因為他覺得錢就是這個世界當中.
最最最重要的東西的時候.
他也都沒有別人那麼多的時候.
或者沒有他想要那麼.
他覺得有些人是多一些.

$^{1041}$那些是蒙福一些的時候.
如果他這樣看的時候.
那他就覺得神不好.
是不是.
就因為他所追求的東西.
他得不到的時候.
他就以為.
神沒有將最好的給他.
所以這裡說到貪圖虛名.
和互相窒道是很有關係.
有時候在一個競爭的關係當中.
為什麼會去到競爭呢.
就是我們大家都要想厲害.
大家都要想表現出.
喂.
我比你厲害.
是不是.
在事奉的關係當中.
喂 我有恩賜過你.
在大家一起工作的當中.
喂 我比你聰明.
當一個這樣的關係出現的時候.
你很難窒道的其實.
因為你追求的東西根本是錯了.
所以在.
保羅的《迦太太書》.
在教會生活當中他提到.
有時候我們貪圖虛名的時候.
就是有這種窒道的心態出現.
有時候人與人之間有些窒道是很.
很窒道的.
不是那麼顯著的.
有時候就是爭誰做誰的朋友.
親密一點.
都有窒道.
我不知道你見過沒有.
就是很窒道的.
就是不是那麼顯著的.
所以我們真的要對付自己的內心.
因為這種窒道將所有的關係都變成酸.

$^{1081}$然後最後自己的心都變成苦.
結束的時候.
希望有幾點的屬靈操練.
都要提醒我們幫助我們對付窒道.
就是你想想神的恩典有多好.
數算一下神的恩典.
數算一下有多少你已經是接受了.
連那些不信耶穌的人寫的.
浪子心聲的歌都說.
命裡有時終須有命裡無時莫強求.
為什麼有時候基督徒.
反而可能很多東西我們有的.
反而我們就不懂得去感恩呢.
所以數算主因是對付窒道一個很重要的恩典.
有時候你不需要強求很多東西.
你就已經有的.
你就是數算.
這個一定是神對你最好最適合你的.
第二就是.
如果你有些人你是真的窒道他.
你發覺你自己是窒道他.
你要為他禱告.
你當然為自己禱告.
但只是為自己禱告不夠.
你為他禱告.
你求神給他更多福氣.
你不是為他禱告.
「神啊,拿他的禱告給我」.
不是這樣禱告.
是說「神啊,更大大的加倍給他」.
加倍的祝福他.
祝福他不夠.
真的要這樣.
你說這麼誇張.
不可能做不出這樣的禱告.
就是你做不出這樣的禱告.
你正正就是要這樣禱告.
「神啊,我做不出這樣的禱告」.
但你聽我那個.
即使不願意的禱告.

$^{1121}$你都要聽的.
你加倍的祝福他.
但慢慢慢慢.
你發覺原來他得到的是你.
有份參與是為你.
你為他的祝福.
你發覺不同.
你不是因為他多了.
原來是窒道他.
原來他多了.
原來是你為他禱告.
這個心情是很有趣的.
但慢慢慢慢你又會發覺會改變.
當然了.
欣賞別人是窒道的一個對付.
窒道的良知.
但特別是學習公開欣賞別人.
窒道是藏起來的.
即是你那些氣人.
窒道那些東西.
你都很鬼祟.
因為你很噁心.
但你公開去讚賞別人的時候.
其實都是一個反窒道的藥.
最後其實都是一樣.
你停止比較.
有些東西你真的不需要比較.
各有前因莫善人.
每個人都不同.
所以將來在天間的時候.
神會不會因為你這樣這樣.
因為你跟另外一個人比較.
在分賞次的時候又要這樣比較.
根本就不會的.
而當我們這樣的時候.
我們是更深的紮根於福音的身份.
我們的整個身份.
是建立於神對我們無條件的接納當中.
跟我們得多少.
跟我們做多少.

$^{1161}$跟我們有多少是完全沒有關係的.
\newpage



\section{路加福音 5:33-35}
\label{sec:031864}
\textbf{為主禁食}~~~~黃健羒牧師
\newline
\newline
連結: \href{http://www.church.com.hk/acms/content.asp?site=cdc\&op=show\&type=product\&code=031864\&layout=sermon}{\texttt{031864}} ~~~~ 語音日期: 2026-03-22
\newline
\newline
\hyperref[sec:031863]{\small{< < < PREV SERMON < < <}}
~
\hyperref[sec:index_chronic]{\small{[返順時目]}}
~
\hyperref[sec:031865]{\small{> > > NEXT SERMON > > >}}
\newline
\newline
路加福音 5:33-35
\newline
\begin{longtable}{cl}
\hline
\hline
章節 & 經文 (和合本修訂版)\\
\hline
5:33 & \begin{tabularx}{0.7\textwidth}{X} 他們對耶穌說：「約翰的門徒常常禁食祈禱，法利賽人的門徒也是這樣，惟獨跟你在一起的又吃又喝。」 \end{tabularx} \\ \\ \relax
5:34 & \begin{tabularx}{0.7\textwidth}{X} 耶穌對他們說：「新郎和賓客在一起的時候，你們怎麼能叫賓客禁食呢？ \end{tabularx} \\ \\ \relax
5:35 & \begin{tabularx}{0.7\textwidth}{X} 但日子將到，新郎要被帶走，那些日子他們就要禁食了。」 \end{tabularx} \\ \\
[1ex]
\hline
\hline
\end{longtable}
$^{1}$在傳統的基督教會有七大節期.
例如聖誕節 復活節 可能大家也熟悉.
今天沒有時間一一介紹.
只說一個.
就是復活節之前四十天.
如果不計星期一的話.
就被稱為大齋節.
Lent Season.
我們就在這個季節當中.
復活節之前四十天.
很多基督徒在這四十天當中就效法.
耶穌基督當日在曠野.
就是祂出來公開傳道之前有四十天.
是禁食 祈禱 受魔鬼的引誘.
有些人覺得禁食就是絕食.
我比較喜歡用禁食這個字.
待會我會解釋一下.
禁食有不同的程度.
就是不吃東西.
絕食就是甚麼都不吃.
只有一種.
所以禁食就有不同的程度.
在這幾十天當中.
究竟我們的思想主題在哪方面呢?.
其實就只有一個.
主耶穌曾經為我寫明.
我為祂捨棄甚麼呢?.
主是為我寫明釘十字架.
那我為祂捨棄了甚麼呢?.
禁食其實是一種捨棄食物.
但其他很多信徒在傳統的四十天當中.
就思想這些.
究竟我在心願當中為主捨棄甚麼呢?.
在《馬太》,《馬河》,《路加》.
三本福音當中都有記載.
耶穌被聖靈引領去曠野四十天.
受魔鬼撒旦的試探.
後來就肚餓了.
四十天不吃東西.
我們當然有問問題.

$^{41}$有沒有喝水呢?.
聖經就沒有記載這些.
而且那四十天當中.
不斷地受魔鬼撒旦的試探.
當然在《馬太福音》裡面.
更加詳細記載了有三次的試探.
撒旦說:如果你是神的兒子.
你就將石頭變成麵包.
但除了那三次.
其實在那四十天當中.
都是不斷受撒旦的試探.
耶穌可以四十天不吃.
或者不喝.
我們千萬不要隨隨便便學習.
我可以學耶穌.
我可以有四十天不吃東西.
因為我們知道.
人生存的三個通則.
一般我們要提醒.
你不能夠三分鐘缺氧.
那三分鐘沒有氧氣.
腦袋就會受影響.
所以急救的時候.
那三分鐘要壓心臟.
或者人工呼吸.
三小時如果沒有隱蔽的時候.
如果大熱天或者風雪.
沒有地方可以躲藏的時候.
就不行了.
身體就會失溫.
三天不喝水.
身體就不行了.
嚴重脫水.
去旅行最重要是喝水.
所以坐飛機.
偶爾水也要喝.
因為不喝水是不行的.
如果三個星期沒有食物的話.
你的內臟就開始受寸斷.
所以這個三個通則.

$^{81}$通常我們說人.
普通人不能夠沒有.
不能夠三分鐘缺氧.
不能夠三小時缺氧.
不能夠三日沒有水.
不能夠三星期沒有食物.
所以耶穌是一個特殊的例子.
四十天.
你千萬不要隨便學他.
今天我們要學習的話.
禁食要禁食多久?.
有什麼益處?.
我們今天就談談這個問題.
當年反對耶穌的.
那些是法理塞人和文士.
他們都是猶太人.
剛才我所讀的經文.
就說他們來質問耶穌.
約翰的叛徒時時都禁食.
法理塞人的叛徒也都是禁食禱告.
那為什麼你的叛徒又不禁食禱告呢?.
可見當時的猶太人.
其實是有習慣來禁食祈禱的.
我們先談談禁食的傳統.
為什麼我用傳統這個字.
而不用誡命這個字呢?.
因為在摩西的律法裡面.
從來沒有明顯說什麼時候要禁食.
在律法裡面沒有提到這個.
只有在利未記當中.
第十六章有提到「刻苦己心」.
中文的翻譯「刻苦己心」.
很多解經家就說.
這個就是指禁食.
在利未記當中.
而在大衛王開始.
我們知道他有禁食.
大家還記得.
他犯罪.
生了一個兒子.

$^{121}$但是神讓這個兒子有病.
大衛很傷心.
希望神醫治.
聖經就說.
他表示是懺悔.
認罪.
他就禁食.
期望神醫治這個兒子.
用禁戒身體享受.
表示虔誠.
祈禱.
這個是猶太人的傳統.
後來我們在聖經裡面.
在《以斯帖記》也提到.
猶太人被流放去波斯的時候.
以斯帖王后為了冒險.
主動去見皇帝.
請所有猶太人.
包括她自己的宮女.
為這件事禁食三天.
不喝不吃.
然後去見皇帝.
意思也是一樣.
表示虔誠.
祈禱求神介入.
求神保護.
所以這是猶太人一直以來的傳統.
然而在《爾賽亞書》第58章.
我們讀到先知的責備.
當時百姓以為可以用禁食.
克苦己心.
打動神的心.
應允他們的祈禱.
這是人之常情.
我自己很辛苦很餓.
希望用苦肉計.
神可以受到感動.
但是《爾賽亞》怎樣去警告他們?.
第58章說.
「看!你們禁食的時候.

$^{161}$仍追求私利.
剝削為你們做苦工的人.
看!你們禁食.
卻起紛爭輕重.
以凶惡的拳頭打人.
你們今日這種禁食.
無法使您的聲音.
聽聞與告處」.
即是如果你不能配合.
仍然做很多壞事的時候.
無論你是怎樣禁食.
都無法使你的聲音.
達到神的面前.
《爾賽亞》站在神的立場說.
「我所要的禁食.
豈不是要你們鬆開凶惡的繩.
解開惡上的鎖.
避免欺壓的德自由.
截斷一切的惡嗎?.
豈不是要你們食物分級飢餓的人.
將流浪的窮人接到家中.
見赤身的吸他衣服遮體.
而不是隱藏自己避開念的骨肉嗎?」.
神的意思是.
你們禁食歸禁食.
但是要配合.
不單止不做那些壞事.
還要做一些憐憫的事.
慈悲的事.
這樣才是一種正直的禁食.
所以才宣告神的心意.
禁食不是打動神的心.
答應我們所求的.
禁食是我們要行動配合.
和我們所作所為息息有關.
不但我們彼此之間要和睦相處.
彼此相愛.
更是要爭取社會公義.
照顧那些有需要的人.
憐憫軟弱的人.

$^{201}$所以禁食不是用來支配神的手段.
而是藉著克制我們的慾望.
順服神的旨意.
實行神的誡命.
所以千萬不要以為.
怎樣才能夠讓神聽我的祈禱.
怎樣才能夠讓我的祈禱更加有效.
好叫神就答應我所祈求的事.
千萬不要以為禁食就能夠增加祈禱的力量.
禁食一天不夠.
禁食兩天.
禁食兩天不夠.
禁食三天.
以為越長越好.
我告訴你.
如果這不是神的旨意.
你就算一個星期禁食.
神仍然不會按照你的期望來完成這件事.
還有一種心態不應該有.
就是將禁食作為一種炫耀.
以為被別人屬靈.
在《瑪特福音》.
耶穌自己教導.
你們禁食的時候不要像那些假無為善的人.
臉上帶著羞辱.
因為他們逢頭糾面.
故意讓人看出他們在禁食.
舔得很冷.
我實在告訴你們.
他們已經得了他們的賞賜.
你禁食的時候要梳頭洗面.
不要讓人看出你在禁食.
只要你隱秘中的乎看見.
你負責隱秘中察寒必燃賞賜.
耶穌的教導.
禁食不要告訴別人.
別人不知道.
你要穿著整齊乾淨.
不要讓人知道你在禁食.
假無為善的人就相反.

$^{241}$要讓人知道.
別人就稱讚你很虔誠.
他現在禁食祈禱.
刻苦給自己的心客稱讚.
所以禁食是為了神.
祂知道就可以了.
那些人問耶穌.
為什麼你的門徒禁食.
你的門徒不禁食呢?.
耶穌就說.
新郎和賓客在一起的時候.
怎樣可以叫賓客禁食呢?.
但日子將到.
新郎就要來帶走.
那些日子他們就要禁食.
其實新郎就是指耶穌.
當他還在世上的時候.
門徒應該把握機會多跟他學習.
不需要禁食祈禱.
但他不在的時候.
門徒反而要禁食.
安靜,默想.
在靈裡面與他溝通.
思想耶穌的話.
所以新郎不在的時候.
就談到耶穌後來被釘十字架.
死亡,離開世界.
我們哀傷禁食.
不是因為可憐祂.
雖然今天有很多人說.
等於耶穌釘十字架.
他們不相信耶穌.
但看見這麼慘.
被釘在十字架上.
他們都會覺得很可惜.
都有些憂傷.
我們是為主.
寫命在十字架上.
不是單是為了.
祂受苦被釘在十字架上.

$^{281}$而是因為祂代替了你和我.
本來是我們犯罪.
我們應該被釘十字架.
但現在祂代替我們受罪.
所以哀傷是為了我們自己的罪.
為了我們自己得罪神而哀傷.
我們為自己的軟弱得罪愛我們的神.
憂傷.
所以在禁食期間.
不是單單思念主十字架寫命.
而是為我們自己的罪哀傷.
祂是神.
禁食是為了主.
在一切禁食操練之前.
這是最重要的.
千萬不要為了自己.
為了祈禱要神答應.
所以怎樣怎樣.
主要是為了神的旨意.
為了認識祂.
所以禁食不是為了減肥.
不是為了省錢.
不是為了清除腸胃.
或者有人用排毒.
我不知道是否中藥.
排毒清除腸胃.
或者有病不能吃什麼就不吃.
不是為了這些原因.
而是為了主的原因.
為了思念祂.
為了默想.
這是我們第一要注意的.
接下來我們說禁食的操練.
禁食的目的是為了主.
為了與主更加親密來往.
更加認識主和祂的旨意.
實際上是怎樣的操練呢?.
不知道大家有沒有試過.
其實禁食主要是針對我們對食物的慾望.
作為一種克制.

$^{321}$以至我們放下最基本生存的需要.
在這樣的時候.
我們內心就會有一些感受.
一些想法.
一些體會慢慢浮現出來.
如果你未有過這樣的經驗.
你可能會覺得很玄.
但是有過禁食的.
我與一些弟兄姊妹聊天的時候.
他們也會有這樣的經歷.
就是到了上下的時候.
有些想法就會自然產生出來.
禁食了十二小時,十六小時後.
開始反而腦筋很清醒.
非常警覺.
特別敏感.
即是腎的感動.
其實簡單來說.
禁食是利用我們身體對食物的需要.
作為一種提醒.
你肚餓.
很想吃東西.
看完照片也很想吃東西.
但是我不吃.
告訴自己.
是因為要用這段時間.
我要與神有更多的溝通.
這段時間是屬靈的操練.
我要祈禱.
我要默想神的話.
與祂聊天,與祂見面.
我很想喝果汁.
很想喝咖啡.
很想喝湯.
但是我不喝.
為了主的緣故.
每一次我們身體有這樣的需要.
肚餓,口渴.
這就提醒了.
為什麼我不喝.

$^{361}$因為我是為了主的緣故.
這麼好吃為什麼我不吃.
因為是為了主的緣故.
所以不用用鬧鐘.
你的身體一餓的時候就提醒你.
這段時間為什麼不做.
因為我要與神有更多的溝通.
所以如果你從來沒試過禁食.
不要一開始就說.
我一天就不吃東西.
可能你的身體受不了.
必須從一些最容易做的事情開始.
來學習,來克制慾望.
這個介紹了好幾種不同程度的禁食.
第一種就是節食.
少吃一點.
平時你下午吃東西.
吃八分飽的.
不妨在這一餐吃五分飽.
吃少一點.
下午到晚餐之前可能就會肚餓.
平時你午餐吃到飽.
不會,但是現在吃少一點.
你可能又餓了.
就不吃了.
提醒你忍幾個小時.
很快就過去了.
口渴什麼?要喝東西嗎?.
平時港式奶茶還要加珍珠.
好喝.
現在只喝水.
不享受,只喝水.
原則就是減少食物.
減少享受的飲料.
然後就可以做到了.
節食不是這麼難的.
吃少一點,吃簡單一點.
第二就是避免大魚大肉.
吃最普通的.
最容易買到的.

$^{401}$最基本的.
連Walmart,Target那些都可以買的.
肉類也不講究哪個部分.
什麼T骨那些.
最基本的肉就可以了.
或者是不花錢去買.
以前我們有朋友特別在阿拉斯加.
空雲那些魚.
那些空雲來的,最新鮮的.
不要花錢買那些.
避免大魚大肉.
這個是第二個可以做得到的.
第三就是簡單的烹飪.
我們唐人最懂得煮東西.
煎炒,煮炸.
做很多方法.
又炆,又燉,又焯.
我上網查過AI.
起碼有36種不同的煮東西的方法.
36種.
我都說不出10種.
最普通的.
很多方式來煮東西.
在禁食的時候就只有一種.
只用水煮.
最簡單.
水煮.
也不加調味品.
油鹽,醬醋什麼的.
什麼都不加.
只用水煮.
鹽也不加.
那當然不好吃.
你覺得不好吃.
又鹽又不加鹽.
那就提醒自己.
我這段時間是為了主的緣故.
而不做這些事.
是為了要多讀聖經.
祈禱,默想.

$^{441}$所以不花時間去烹飪.
平時你搞個三個小時才有餐吃的.
現在15分鐘就搞定了.
那你省一點時間.
就可以有時間讀經.
平時你說沒時間讀經.
很忙,沒時間祈禱.
那你現在就有時間.
所以是好的,簡單的烹調.
還有就是吃齋.
吃蔬菜水果.
不吃肉.
其實原則都一樣.
簡簡單單是一切重新.
到最後就只喝水.
那多久才禁食一次?.
在耶穌的時代.
那些法利賽人一般就是.
一個星期有兩天.
每天就是一餐的正餐.
其實一餐的正餐.
都不可以說是整天不吃東西.
可以說是12小時吃一餐.
12小時吃一餐.
那就是他們的禁食.
一個星期有兩天.
有些像現在的回教徒齋戒月.
其實他們齋戒月的Ramadan.
就是日出之前可以吃東西的.
日出至日落就不吃東西.
日落之後他們要吃大餐.
所以基本上就是12小時吃一餐.
就是這個意思.
可以行.
其實新約時代的信徒.
都是跟他們這樣.
不知道大家有沒有讀過一本書.
叫做《十二仕徒遺訓》.
就是遺留下來的教訓 D-1K.
它裡面介紹初期信徒.

$^{481}$初期基督教會的信徒是怎樣禁食的?.
他們都是一星期兩天.
不過他們就註明了.
我們就不像那些猶太人.
星期三,星期五.
我們就是星期一.
我看看有沒有記錯.
猶太人就是星期一和星期四.
他們就說我們不要跟他們同一個時間.
我們就星期三,星期五.
所以是一個星期有兩天.
這是很長的時期.
幾百年.
基督徒是一個星期有兩天禁食的.
同樣的.
禁食就是12小時之內吃一餐.
這就叫禁食.
平時吃三餐.
或者你除了三餐之外.
還有一些零食.
不過就可以學習.
12小時吃一餐.
其實不是那麼難的.
當你的胃部經過訓練之後.
你可以24小時吃一餐.
或者36小時,72小時.
這個時間不是那麼重要.
重點不是說時間越長越好.
不是說越長你的祈禱越有效.
而是說長一點的時間.
你就多一點時間想到要考慮的事情.
神也會讓你看到更加遠,更加長遠的事情.
所以這些都是禁食有關的事情.
有些弟兄姊妹就說.
我有胃病.
有糖尿病.
不吃東西怎麼行.
不行的.
我個人的意見就是禁食.
都是你和神之間的約定.

$^{521}$不是一定的.
如果你身體不受得的話.
不禁食沒有所謂的.
不能夠勉強.
但我又聽到有些人說.
平時真的不行的.
但我這段時間和神有一個約定.
我用未來的12小時.
我只吃一餐.
身體又可以的.
這是神特別的恩典.
有些信徒見證.
雖然身體有毛病.
但特別為這件事情禁食祈禱的時候.
都可以的.
所以只要我們記得.
你是為主禁食.
不是為自己的時候.
應該就沒有問題.
第三,我們就到禁食的益處了.
禁食在這個時代.
特別在美國的文化當中.
好像有些脫節.
美國人如果去遊埠的時候.
就知道哪些人是從美國來的.
當然要吃東西的.
去哪裡買個雪糕.
一停站的時候又去買食物.
當然是從美國來的.
我們不習慣.
但是按照我剛才的解釋.
最起碼禁食一日只吃一餐的話.
其實不是那麼難.
對於我們信主的人來說.
有什麼益處呢?.
我分開幾方面來講.
第一個是個人方面.
是培養我們靈性的品格.
第二就是智慧的方面.
如果你有.

$^{561}$好像耶穌所說的.
他在世界上跟隨的門徒.
不必拘束.
按猶太人的規矩禁食.
他離開世界之後.
反而需要禁食.
因為是為了主的緣故.
意思就是說有禁食的時間.
也有吃大餐的時間.
那什麼時候是做什麼呢?.
這就是智慧.
你要判斷那個時候.
什麼時候是享受食物.
感恩節什麼的.
但是也有一些禁食的時候.
你就是不吃東西.
也是為了主.
怎樣做一個合宜的決定呢?.
這就是智慧.
慢慢有這種操練的時候.
你才學得到.
喜樂的時候就喜樂.
悲傷的時候就悲傷.
還有其他弟兄姊妹.
有一切同心.
有適當的時機.
就是操練我們的智慧.
第二就是操練我們的自由.
你可以吃.
也可以不吃.
我記得在大學的時候.
我對於按時飲食非常重視.
有時說了十二點鐘吃東西.
遲十分鐘好像都.
不習慣一定要找些東西吃.
因為十二點鐘到了要吃東西.
但後來有了禁食的操練之後.
遲一些.
甚至不吃.
無所謂.

$^{601}$可以的.
即是胃受了訓練之後.
就不會抗議.
這就是自由.
你可以吃.
也可以不吃.
你吃多一點也可以.
吃少一點也可以.
這就是一個自由.
這都是我們在禁食的操練之後.
就會在靈性上面.
這方面是能夠享受到.
第三就是節制方面.
如果你有抗拒美食的操練之後.
就算好吃.
但我都不吃.
是為了住的緣故.
吃了五分飽.
你又可以接受七分八分.
又可以.
不因為免費.
有些人吃自助餐.
錢都給了免費.
吃多一點.
自助餐通常會超量.
吃多一點.
有了禁食節制的能力之後.
下一次.
有些引誘來到了.
無論是金錢也好.
權力也好.
美食也好.
你有能力抗拒.
你就慢慢學習.
No means no.
Say no.
習慣了有這樣的習慣的時候.
有些人來到.
我就是不做.
我就是不去.

$^{641}$我就是不看.
如果你有這樣的操練.
沒有這樣的操練的時候.
有時會想想.
少少吧.
但你有這樣的操練的時候.
習慣了.
不吃就是不吃.
不去就是不去.
這樣你就有節制的能力.
免得得罪神.
這是我們個人方面的操練.
是有益處的.
另外一方面就是和世界的關係.
如果你有過飢餓的感覺.
你就能夠體會世界上.
其實很多人.
各種的原因.
是一天只吃一餐.
有些人為宗教的原因.
有些人沒有選擇.
或者貧窮.
有些人打仗.
你因為有過飢餓的經歷的時候.
你就知道.
原來世界很多人都是這樣的感覺.
原來世界很多人.
甚至那些做母親的為了給子女多吃.
一些東西自己捱餓.
都是這樣的.
你就會明白.
為什麼有些人會是這樣.
甘願自己沒得吃.
都要給人有得吃.
你慈悲的心.
憐憫的心.
你得到培養.
你不會單單和人講道理.
或者更加有一個同理心.
明白.

$^{681}$站在對方的角度來哄.
來想事 看事.
當你有得吃.
又可以買得到食物.
你就會傳感恩的心.
嘩 今天可以吃得到.
吃得樂.
可以買得到.
這就是我們和世界之間的關係.
同時你不會浪費食物.
有些還可以吃的.
你就留著.
我知道有些人不習慣.
吃那些隔夜sung .
通常有些家庭.
只有一個人會吃那些隔夜sung .
不知道是媽媽還是爸爸.
其他人是不吃的.
但你有過一個禁食的經歷之後.
你知道食物是很重要.
是寶貴的.
特別你肚餓的時候.
吃這麼多東西都是很享受的.
你就知道省一點.
不要浪費.
這些都是禁食帶來的好處.
第三方面就是和神的關係.
你多一點的禁食.
就是多一點花時間和神來往.
你和神越來越親近.
越來越接近.
經常聊天的時候.
你就會明白神的心意是怎樣的.
其實好像人和人一樣.
你多一點時間在一起聊天.
你經歷同一樣的事情.
彼此之間就心靈相通.
所以我們說那些夫妻.
在一起幾十年的.
說上面那句都知道.

$^{721}$下面那句說什麼.
或者去餐館點什麼菜都知道.
大家一通.
你和神的關係.
如果時時有來往的時候.
你會明白祂的心意.
不需要去到那時候.
哎呀 神的旨意是怎樣.
或者神的心意是怎樣.
你會明白.
因為那時候你和神很熟.
時時都有在聊天.
不需要請牧師為你祈禱.
你自己就可以和神祈禱.
不需要和弟弟姐妹一起的時候.
你和我祈禱吧.
我不懂.
不懂就學.
漸漸你和神多一點聊天的時候.
你就知道說什麼.
有些人不需要說.
有些人明白.
時時都意識到神就在我旁邊.
禁食的時候就達到這個目的.
禁食的時候多一點讀聖經.
你對神的話語.
對神的作為會更加明白.
生活遇到什麼情況.
神的話語自然就會浮現出來.
你會清楚祂的引領.
所以禁食這個操練.
在舊約和新約聖經裡面都強調.
不是單單學怎樣去禁食.
不是用苦肉計去打動神的心.
不是藉著禁食表示我們的虔誠.
而是為了靈裡面能夠與神溝通.
讓我們的心靈更加的潔淨.
可以聆聽神的話語.
所以禁食的焦點.
乃是我們和神的關係.

$^{761}$今年復活節是什麼時候?.
4月5日.
之前這段時間.
我們還說是大齋節的期間.
希望你都可以試一試.
如果你沒有試過的話.
不同程度的禁食.
少吃一點.
或者只是吃蔬菜水果.
無論你是什麼程度.
用多一點時間在神的面前安靜.
默想祈禱.
我相信你會對神有更加多的認識.
亦都盼望你能夠明白神的心意.
體會祂藉著耶穌基督釘十字架受死的愛.
在復活節的時候.
我們會更加興奮地去慶賀.
耶穌帶來的新生命.
請大家低頭同心祈禱.
賜給我們在世界上可以享受到美味的食物.
特別在美國這個地方.
真的想吃什麼都可以買得到.
飲品亦都是很多的選擇.
但求主幫助我們.
為了更加與你親近.
為了更加多花時間和你在一起.
我們禁止自己去享受.
為了操練自己.
為了與世界,與主的關係更加密切.
由這部教導我們在這段時間當中.
更多的思念.
你自己對我們的愛.
我們為你捨棄什麼.
阿們.
\newpage



\section{約翰福音 20:19-29}
\label{sec:031865}
\textbf{多代同行 (2) PeerWalk (3) 你們的……與我何干?}~~~~麥慕思牧師
\newline
\newline
連結: \href{http://www.church.com.hk/acms/content.asp?site=cdc\&op=show\&type=product\&code=031865\&layout=sermon}{\texttt{031865}} ~~~~ 語音日期: 2026-03-21
\newline
\newline
\hyperref[sec:031864]{\small{< < < PREV SERMON < < <}}
~
\hyperref[sec:index_chronic]{\small{[返順時目]}}
~
\hyperref[sec:031866]{\small{> > > NEXT SERMON > > >}}
\newline
\newline
約翰福音 20:19-29
\newline
\begin{longtable}{cl}
\hline
\hline
章節 & 經文 (和合本修訂版)\\
\hline
20:19 & \begin{tabularx}{0.7\textwidth}{X} 那日（就是七日的第一日）晚上，門徒因怕猶太人，所在的地方門都關了。耶穌來，站在當中，對他們說：「願你們平安！」 \end{tabularx} \\ \\ \relax
20:20 & \begin{tabularx}{0.7\textwidth}{X} 說了這話，他把手和肋旁給他們看。門徒一看見主就喜樂了。 \end{tabularx} \\ \\ \relax
20:21 & \begin{tabularx}{0.7\textwidth}{X} 於是耶穌又對他們說：「願你們平安！父怎樣差遣了我，我也照樣差遣你們。」 \end{tabularx} \\ \\ \relax
20:22 & \begin{tabularx}{0.7\textwidth}{X} 說了這話，他向他們吹一口氣，說：「領受聖靈吧！ \end{tabularx} \\ \\ \relax
20:23 & \begin{tabularx}{0.7\textwidth}{X} 你們赦免誰的罪，誰的罪就得赦免；你們不赦免誰的罪，誰的罪就不得赦免。」 \end{tabularx} \\ \\ \relax
20:24 & \begin{tabularx}{0.7\textwidth}{X} 那十二使徒中，有個叫低土馬的多馬，耶穌來的時候，他沒有和他們在一起。 \end{tabularx} \\ \\ \relax
20:25 & \begin{tabularx}{0.7\textwidth}{X} 其他的門徒就對他說：「我們已經看見主了。」多馬卻對他們說：「除非我看見他手上的釘痕，用我的指頭探入那釘痕，用我的手探入他的肋旁，我絕不信。」 \end{tabularx} \\ \\ \relax
20:26 & \begin{tabularx}{0.7\textwidth}{X} 過了八日，門徒又在屋裡，多馬也和他們在一起。門都關了，耶穌來，站在當中，說：「願你們平安！」 \end{tabularx} \\ \\ \relax
20:27 & \begin{tabularx}{0.7\textwidth}{X} 然後他對多馬說：「把你的指頭伸到這裡來，看看我的手；把你的手伸過來，探入我的肋旁。不要疑惑，總要信！」 \end{tabularx} \\ \\ \relax
20:28 & \begin{tabularx}{0.7\textwidth}{X} 多馬回答，對他說：「我的主！我的神！」 \end{tabularx} \\ \\ \relax
20:29 & \begin{tabularx}{0.7\textwidth}{X} 耶穌對他說：「你因為看見了我才信嗎？那沒有看見卻信的有福了。」 \end{tabularx} \\ \\
[1ex]
\hline
\hline
\end{longtable}
$^{1}$今日跟大家分享 大家仍然是跟教會講道的題目.
叫做「多待同行」.
不過你們去到另一個系列叫做「Peer's Walk」.
給了我一個講題 我也給了一個講題.
我的講題 去多了.
「你們的與我何干?」.
大家有沒有猜到我說什麼呢?.
我給這個題目的時候 我也覺得.
我給這個題目的時候 我自己也在想.
大家明白我說什麼嗎? 大家繼續聽下去 可能會明白.
不過大家看到這句話 「你們的與我何干?」.
你會聯想起什麼呢?.
你會聯想起別人在吵架的時候.
可能找你上來做和事佬 你就說.
「關我什麼事啊?你們在吵架 不關我的事」.
或者有人想你靠邊站 你站在他那邊的時候.
你可能也會說 「關我什麼事啊?你們的事」.
在這些廣東話的狀況裡面.
今天和大家看的是《約翰福音》二十章十九到二十九節.
就是約翰耶穌被釘十字架之後.
埋葬第三天復活之後.
第一次與那些門徒一起顯現.
是不是好像「砰砰」的聲音呢?.
相信大家都應該很熟悉這一段經文.
過多一點就是復活節了 這段是耶穌復活了之後的經文.
門徒發生什麼事呢?待會我們會繼續看.
當中似乎有些人不相信 有些人相信.
有些人要看見才相信 所以我們會一起去看.
不如我們先一起祈禱吧 好嗎?.
主我們多謝你 讓我們可以一起去聚集 一起去敬拜你.
主 剛才大敬拜的姐妹都告訴我們.
世界很紛亂 世界有很多事情我們不理解.
但有很多事情與我們都有關係.
特別是在主身上的 都有很多關係的事情.
特別是在聖經身上的事情都與我們很有關係.
所以讓我們知道 聖經不是一本普通的書.
與我們有直接關係.
並且我們周邊的弟兄姊妹.
祂生命的事情都與我們有關係.
我們不是沒有關係的人.

$^{41}$因為我們經常說我們是弟兄姊妹.
我們是一家人.
所以讓我們去體會 主你今天要與我們說話.
向我們說話.
我們的禱告祈求是.
奉主耶穌基督得勝名字祈求 阿們.
先讀一段經文 因為19至29節我會分開幾段.
分開大約三段一起讀 好嗎?.
我們先讀第一段 預備起.
「禮拜日黃昏的時候門徒聚在一起.
因為怕猶太人就把門戶都關上.
耶穌來了站在他們中間說.
願你們平安.
說了這話就把手和肋旁給他們看.
門徒看見了就歡喜了.
耶穌又對他們說.
願你們平安.
父怎樣差遣了我.
我也怎樣差遣你們.
說了這話就向他們吹一口氣.
說 你們領受聖靈吧.
你們赦免誰的罪.
誰的罪就得赦免.
你們不赦免誰的罪.
誰的罪就不得赦免」.
這裡有個背景 剛才我都說過.
主耶穌死後三天.
今天是復活了 第一天.
所以這裡接著的意思就是.
Good Friday 昨天是星期五.
即是明天是星期日.
耶穌就復活了.
在黃昏的時候 門徒聚集在一個樓房裡面.
這個樓房應該不是在大街.
不是熙來攘往的地方.
因為他們提到關閉門.
即是要一個隱藏的地方.
一個很秘密的基地.
這個門是關了的.
為甚麼要關門呢.

$^{81}$關門不是大吉吧.
為甚麼要關門呢.
因為他們怕.
基文說怕 怕甚麼呢.
怕猶太人.
怕猶太人來捉他們.
怕猶太人來迫害他們.
因為耶穌自己的師父才剛被迫害.
所以他們怕猶太人.
怕環境的惡劣.
我想我們現在這個時代很少有師父.
但他們這班門徒沒有師父就沒有甚麼.
沒有目標 沒有方向.
可能更加被周圍反基督教的人去攻擊.
怕環境惡劣.
還有一樣 怕死.
因為羅馬政府.
猶太的文士 祭司 花里塞人.
可能都會捉他們.
要他們受苦 甚至受死.
怕苦難.
因為捉了之後會有甚麼刑.
極刑.
會有想不到的狀況.
如果你看了耶穌傳你都知道.
那些刑具 那些刑罰.
不是我們現在能夠想像.
就算你去到怕苦難 你要逃.
你要離開.
你都要離開別井.
我想我們這個會明白多一點.
因為現在世界的紛亂.
你都可能明白人要逃的時候會是怎樣.
亦都怕前逃的問題.
因為以前有跟隨耶穌.
方向很多都清晰.
但現在似乎光明不在前面.
加上剛才我之前說的種種的怕.
似乎都只有在黑暗裡面.
當人不知的時候.

$^{121}$就有很多恐懼 很多擔心來到.
我不知道在座中的你.
你最近有沒有擔心世界呢?.
擔心香港的前景呢?.
我剛剛去了台灣.
去了八天左右.
我覺得台灣給我的一個狀況很沉.
我以前去一個街市.
全部都很多人.
這次我再去這個街市.
十個店舖只有一間開.
可想而知那個情況是這樣.
你對未來有沒有擔心呢?.
你們年輕一群.
你對未來有沒有擔心呢?.
有沒有懼怕呢?.
門徒是怕的.
但他們怕到一個階段.
就不是四散.
他們反而聚在一起.
聚在一起做什麼呢?.
聚在一起當然是商討一下前面有什麼可以做.
沒有我們現在可以用WhatsApp溝通.
所以一起聚在一起.
一起集思廣益.
一起有分析.
看看大環境是怎樣.
看看形勢是怎樣.
看看前路可以是怎樣.
看到嗎?.
他們原來不是一個人的.
他們不是一個人去面對情況.
不是自己去走這條路.
我知道上星期Michelle都有提醒我們.
我們不是一個人走的.
我們是.
topic都告訴我們.
peer walk.
有你周邊的人.
我不知道你坐在你周邊的人.

$^{161}$是不是和你同行的人.
而在我們人生不同的階段裡面.
我深信.
你都會遇見過不同的人.
當你有東西想和別人分享的時候.
想訴苦的時候.
你會想起誰呢?.
你會想起你現在旁邊的人.
想起一些弟兄姊妹.
想起你教會外邊的一些朋友.
想起你的家人.
想起阿嵐.
這裡是他負責的.
想起他.
想起誰呢?.
為什麼你會想起他呢?.
他是不是幫過你呢?.
他是不是能夠和你同行過呢?.
是不是有些你遇見過的問題.
他有高見給你呢?.
他有高你一線的眼光.
去和你分析呢?.
你才會找他.
你覺得他信得過.
覺得他不會點到黑路你走.
更加可能他給你一個力量.
給你一個幫助.
給你一個鼓勵.
和你一起走.
稍後我會分享我的一些朋友.
回到經文裡面.
這班門徒聚集在一起.
因為什麼?.
怕.
很重要.
因為怕走在一起.
其實怕.
你有沒有想過.
你以前的怕.
或者過去曾經有怕的經驗的時候.

$^{201}$其實會怎樣呢?.
怕會做到些什麼.
或者做不到些什麼呢?.
他們怕就把門關了.
關門會怎樣呢?.
關門就好像我們現在這樣.
困在這裡.
我們關門都可以開門的.
都可以有路走的.
但我想他們的關門是什麼呢?.
只有一條路.
關了就沒有路可走.
就困在裡面.
這個關門還因為他們怕.
怕就會怎樣呢?.
沒有了平安.
沒有了師父.
方寸就大亂了.
這裡有個嬰兒.
但他應該未懂迷路.
兩個嬰兒未懂迷路.
你小時候.
有沒有被人作弄你呢?.
我小時候被人作弄.
就是我媽躲起來了.
你猜我會怎樣呢?.
小朋友只有一個問題.
哭.
原來當一個人沒有平安的時候.
會怎樣呢?.
會哭.
特別是小朋友會失去安全感.
會很害怕.
自然沒有了平安的時候.
就沒有了喜樂.
面帶愁容.
你看看你旁邊那個.
他有沒有平安和喜樂?.
好像強行笑出來一樣.
希望我們每個人都有平安和喜樂.

$^{241}$因為我們裡面有誰呢?.
有耶穌.
但這班門徒.
你看到他們怕.
他更加將他的心門關了.
將他的心門封閉了.
忘記了他同行的耶穌.
同行的耶穌三年跟他同行了.
他們見到甚麼?.
耶穌行神蹟.
耶穌醫治.
耶穌釋放很多人.
但這刻關上門的時候.
心門被關.
就甚麼都不記得了.
耶穌其實說了很多次預言自己會死.
會復活.
但都忘記了.
心門被封鎖的時候.
就算聖靈要進入.
可能都會有些難度.
因為恐懼.
他們就躲起來了.
門徒真的有很多的怕和恐懼.
但主耶穌呢?.
這個主耶穌.
我說這個主耶穌.
我們的主耶穌.
很愛我們.
祂突然之間站在他們當中.
進到一個被關門的牢房裡面.
耶穌要破他們的心門.
耶穌今天也很想破我們一些的心門.
我們每個人都可能有我們的執著.
有我們死都不願意放手的地方.
耶穌要來破我們的心門.
我曾經也有恐懼.
我恐懼到一個地步.
我有很長的日子.
天天以淚洗臉.

$^{281}$接近一年多的日子.
這個恐懼令我不想見人.
連我家人都不敢碰我.
沒有人敢問我做什麼事.
特別是我家人.
我媽媽不會問我.
由得我.
不過我是看醫生的.
所以原來恐懼帶給一個人是怎樣的呢?.
什麼都做不到.
就像被鎖住一樣.
耶穌站在他們當中.
打破了門徒的心門.
打破了門徒對耶穌的一些想法.
打破人對死亡的恐懼.
我想你們沒有想過死吧?.
這裡人人都這麼年輕.
我來了六年多.
我做了接近五十宗的喪禮.
你可想而知.
我對死亡也有另外一個體會.
我未來同福堂的時候.
只做了一宗喪禮.
來到同福堂之後.
差不多兩個月就有了.
我剛才說我剛剛去完台灣.
八天裡面我收了兩個電話.
這兩個電話就是接著要做喪禮.
所以恐懼死亡是帶給很多威脅.
耶穌因為都被折騰了.
門徒眼見這位師父在十字架上.
把他和門徒愛的關係被分隔了.
帶給人絕望和憂愁.
但是在這一刻.
耶穌站在當中的時候.
他們要知道.
耶穌原來勝過死亡.
死亡原來只是開始.
不是終結.
是起點也不是終點.

$^{321}$這顯現真實可以看見耶穌.
我們今天未必可以真實摸到耶穌.
當時的門徒真的很幸福.
可以摸到耶穌.
見到耶穌的實體.
還有打破了人對復活的疑惑.
因為耶穌是超乎物質.
原來不能阻擋復活的耶穌.
不能尋找他的門徒.
穿牆過壁站在他們其中.
也沒有什麼可以限制復活主的行動.
其實我們讀到這一段經文.
原來已經不是第一個單位見到耶穌.
這些門徒已經是第五個單位.
第一個單位是莫大拉的瑪利亞.
第二個單位是一群婦女.
第三個單位是彼得.
第四個單位大家熟悉的.
二瑪吾斯的兩個門徒.
原來這裡已經是第五個單位.
在復活日裡面.
耶穌不斷向人顯現.
耶穌不是鬼魂.
耶穌有血有肉有身體.
祂把手和肋旁給人看.
給祂的門徒看.
復活的主才是門的本身.
大家應該對聖經不陌生的話.
都知道我們信耶穌的時候.
神是扣我們的門.
讓我們可以進入我們的心裡面.
讓我們得著救恩可以得救.
這裡復活的主是一個復活的能力.
穿進那個緊閉的心門.
穿破一個封閉了的心門.
要他們得著生命.
如果你想起《詩篇》23篇的時候.
他們得著飽足了.
有草食了.
可以甦醒了.

$^{361}$不是懵閉閉了.
不在黑暗的裡面.
闊的耶穌來到門徒中間.
向他們顯現.
說了一句話.
「願你們平安」.
可想而知.
耶穌這句話是在回應他們的心.
在安慰他們.
因為封閉的心.
今天耶穌要闖進他們的心裡面.
打破他們的限制.
蒙蔽了的眼睛要開啟了.
要讓光進到他們裡面.
真神這位愛他們的主耶穌真的復活了.
這些門徒一起經歷真神的復活.
這班門徒.
聖經剛才我們讀的時候.
有沒有記他們的名字?.
沒有.
但是相傳.
或者一般的學者都說.
是那十個使徒.
不過就留了兩個.
一個是出賣耶穌的猶太.
另一個就是多瑪.
但是當他們在挑戰和困難的時候.
這十個門徒是走在一起.
不是你有你 我有我.
獨自去面對.
今天.
不知道你面對著林林總總的困難或者環境.
你的同行的人是誰呢?.
你的小組.
是不是你安全分享的地方呢?.
你可不可以赤露躺開自己的狀況讓別人知道呢?.
你很害怕.
你很擔心.
你的情況是怎樣呢?.
我想不同的階段都有不同的擔心.

$^{401}$現在可能我已經不擔心我的職場.
因為我在教會裡面.
我在擔心什麼呢?.
擔心我的身體.
去到另一個階段.
擔心我牧養的對象.
這裡有兩個嬰兒.
即是有新的家庭在建立.
我們的狀況是有不同的.
在我裡面 我有些同行者.
上星期我看到Michelle也有分享.
我的同行者.
我剛才在計算.
我的同行者是跟我一起讀書的.
跟我一起讀神學 在希腹讀書的一班同學.
25年了.
可想而知我已經服侍了多久.
剛才我說 我也有一段時間在恐懼的裡面.
我這班朋友 這班弟兄姊妹.
我這班戰友怎樣對我呢?.
當我在恐懼的裡面.
在我的舊教會的裡面.
當我一要面對恐懼的時候.
有一個說 我馬上放工.
馬上來接你 馬上來接你回家.
你周邊有沒有這些這麼愛你的弟兄姊妹?.
與你的Peer Walk.
我們要建立的.
我建立了25週年.
我們還一起去旅行.
不單是吃飯.
一起去旅行也有難度.
一起去旅行會互相坦誠相對.
我們在旅行裡面吵架.
25年都吵架.
但是吵完架.
其實感情更加深厚.
感情更加好.
聖經裡面.
雖然剛才說流放裡面.

$^{441}$沒有描述那十個門徒是誰.
但是有一個不存在.
我們就讀一下這段經文.
繼續這段經文 我們一起讀.
預備起.
十二個門徒中.
有一個稱為雙生子的多瑪.
耶穌來的時候.
他們有和門徒在一起.
他們的門徒對他說.
我們已經見過主了.
多瑪對他們說.
除非我親眼看見他手上的釘痕.
用我的指頭探入那釘痕.
又用我的手探入他的肋旁.
我決不相信.
哇 你看不看到多瑪的反應.
我不知道你以前也看過這段經文.
你以前看過的時候.
你覺不覺得多瑪真的有些問題.
他有什麼問題.
我很喜歡聯想那些情況.
我很喜歡用廣東話聯想.
是嗎 耶穌真的報了.
耶穌真的來過.
你有沒有和我一樣的反應.
如果現在耶穌真的在街上.
你會怎樣.
沒事發生.
還是你會衝下去下面.
看看耶穌呢.
如果耶穌在.
耶穌不用的.
耶穌不會在街上.
不過我在想.
我們的熾熱是多少呢.
我經常在想.
現在那些粉絲.
他們追那些明星.
很誇張的.

$^{481}$我不知道你中間有沒有人追星.
是不是用追星這個字.
應該不是用追星這個字.
還是追星.
我怕我已經脫節了很久.
我只知道.
我認識一些人是坐飛機去追星的.
嘩 是吧.
你們認識的吧.
OK OK 有人認識.
還有現在那些中粉.
他們有特別的名字.
我也認識這些字.
中粉也有特別的名字.
都不同了我的世代.
但是我想這一群門徒.
其實你看聖經.
剛才我們讀的經文的時候.
好像很淡然.
我們見過主了.
是不是很淡然地說.
說出了事實.
但是你回到剛才.
我們讀的經文的時候.
耶穌是給什麼.
手和肋旁給他們看.
門徒一看見的時候.
有什麼 歡喜了.
你感不感受到.
那個是不是淡淡然呢.
是不是死氣沉沉.
我們已經見過主了.
是不是這麼死氣沉沉呢.
當多瑪和門徒對話的時候.
其實會是怎樣的呢.
如果你是那群門徒.
你見過耶穌.
看著多瑪的時候.
你會說些什麼話呢.
多瑪你真的蝕卻了.

$^{521}$主耶穌來顯現過.
祂賜我們平安.
又給我們吹一口氣.
當時還有一股暖流.
進入我的心裡面.
起碼我都蠟著了.
原來主耶穌.
即日給我們說的話是對的.
是今天我們明白了.
你應該這麼激烈的.
是不是.
但是我們讀聖經的時候.
沒有讀這麼激烈.
我很喜歡圖畫式讀聖經.
因為我覺得才能讀出那個味道.
因為多瑪的反應.
真的很淡淡然.
無怪乎多瑪怎樣形容祂.
多疑的多瑪.
祂的疑惑真是.
我想祂當然是這樣.
我就是要這樣的.
除非我親眼看見.
耶穌的手和釘痕.
用我的指頭探入祂的釘痕裡面.
又用我的手探入祂的肋旁.
我才相信的.
我是不是有些誇張法呢.
我覺得我們讀聖經要讀一下.
為什麼多瑪會這樣的呢.
多瑪為什麼聽完.
十個祂的一起三年的門徒.
這十個人難道說的時候.
真的好像我剛才淡淡然嗎.
多瑪其實心裡面.
又出現了什麼狀況呢.
不過看到這裡.
多瑪的回答是很決絕的.
多瑪堅決地拒絕門徒的轉述.
祂都不相信.

$^{561}$口裡可能沒有說不相信.
但心裡面是什麼.
祂堅決要看到才相信.
祂要具體地自己經驗才相信.
所以如果你再是那十個的門徒.
你還會有什麼反應呢.
對著多瑪.
唉 祂是這樣了.
不要理祂.
還是你會繼續很肉緊.
很有勸勉式地告訴多瑪.
我們見過主.
你會怎樣呢.
如果耶穌在我們生命裡面的時候.
有沒有那種熱和火呢.
我們傳福音的時候.
有沒有那種熱和火呢.
很重要.
門徒 但很有趣.
你看看那些門徒.
他仍然接納著多瑪.
雖然看到這樣的同伴.
但他仍然接納著它.
因為繼續下來.
我們稍後再讀.
第二次顯現的時候.
門徒是和多瑪一起.
門徒在聚集的時候.
不會不告訴多瑪.
我們稍後再聚集.
我們不和你玩了.
你不要來.
因為你這樣對耶穌.
因為你這樣不相信我們.
門徒沒有.
門徒沒有拋棄多瑪.
沒有拒絕過多瑪.
其實聖經有時候叫我們看到一些東西.
可能我們當中有些未信.
或者有些信心不足的人.

$^{601}$和多瑪可能都是一樣.
我看到耶穌.
你沒有看到.
所以你未必相信我所說的.
門徒仍然繼續接納多瑪.
在我們的教會裡面.
我想我們也接納不同的人.
但是我們有沒有雀躍.
或者很緊張他未信.
很緊張他對神的信心呢?.
多瑪似乎真的有很多疑惑在裡面.
十個門徒和他講.
他都沒有打動到他的心.
從另外一個角度看.
弟兄姊妹如果有現實裡面.
如果有弟兄姊妹和我們分享.
我怎樣經歷經歷耶穌.
你通常的反應是什麼呢?.
哦.
你會這麼狠的呢?.
還是你很熱切的聆聽他的見證呢?.
有沒有渴慕他的見證在你生命裡面呢?.
有沒有覺得.
還是耶穌你對他太好了.
你給他有見證.
我沒有.
妒忌.
比較呢?.
我們經驗的耶穌.
我想耶穌不會拋下你的.
我都說我在向日葵服侍.
我向日葵的平均年齡是70.
快要80了.
平均的年齡.
這裡平均年齡多少?.
30而已是不是?.
大家有沒有數過由大堂上去大堂有多少級樓梯?.
沒有了你們.
一下子上去了是不是?.
30級.

$^{641}$不多.
剛剛30級.
一個80歲的老人家走30級容易嗎?.
我們有一些老人家要坐輪椅上來.
大家知不知道我們有個輪椅升降機.
升降機上來.
知不知道一個輪椅一個人一個升降機.
坐一個人和他上去.
然後空升降機下來.
多久?.
應該沒有人知道.
八分鐘.
八分鐘的要.
我今早要用輪椅.
要用輪椅的人有八個.
用多久?.
大家懂得計數.
最高峰是12個.
如果有些老人家不早點回來.
最後那個應該不用上去了.
因為崇拜完畢了.
是不是?.
OK.
就是這樣.
有一個老人家真的要坐輪椅回來.
原來他因為跌了很多次.
在家裡又跌在街上又跌.
所以他自己走路都很緊張.
不夠膽走.
所以一定要坐輪椅回來.
看見30級樓梯他坐升降台上來.
但原來他平日在家附近的公園有走樓梯.
不過是五級樓梯.
有一天真的很多人要排隊下去.
我們的同工都會等到最後一個.
坐完升降台去到地下我們才走.
就聊聊天.
我們說你平時有走樓梯.
不如回來舊的都試一下.
對他來說大不大挑戰?.

$^{681}$很大挑戰.
但他真的被我.
我們不是說一次.
我們說了幾個星期.
有一天他突然在大堂放開他的輪椅.
走了上來.
他的女兒是信耶穌的.
他的女兒每個星期都和他回來.
她說媽媽你是不是傻了?.
走樓梯.
弟兄姊妹.
他的見證可能你聽到.
你有什麼感受呢?.
嘩 真是厲害.
和我無關.
他就是他 當然和我無關.
但你給我.
我覺得這是一個鼓勵的肢體.
是給我一個鼓勵.
原來牧者說話要小心一點.
你們也要小心一點.
不只是牧者.
說話帶有什麼?.
說話帶有能力.
他能夠走上來.
並且我又想.
如果我老的時候.
我會怎樣呢?.
八十歲的人原來都可以突破.
都可以改變.
何況現在的我.
我應該更加可以多突破和改變.
我還有一幅相.
他還有一樣厲害的東西.
給大家看看這幅相.
不對 應該是上了幅.
怎麼不見了?.
這個 他抄的.
這邊是抄歌詞.
這邊是抄聖經.

$^{721}$看到他抄什麼嗎?.
每樣抄三次.
這是《恩典之路》的歌.
八十多歲的老人家.
他激勵了我.
我現在抄聖經.
原來他的見證可以成為什麼?.
我的見證.
可以成為我的激勵.
我們說同行.
他也在和我同行.
我不只是他的牧者和他同行.
他帶給我很多的鼓勵.
再過八天.
多瑪 耶穌就向他顯現了.
再讀前面.
不對 下面.
讀完這個 預備.
過了八天.
門徒又在屋子裡.
多瑪又和他們在一起.
門戶都關上了.
耶穌來了.
站在他們中間說.
「願你們平安」.
然後對多瑪說.
「把你的指頭放在這裡.
看看我的手把.
伸出你的手來.
探探我的肋旁.
不要疑惑 只要信」.
多瑪對他說.
「我的主 我的神」.
耶穌說.
「你因為看見我才信嗎?.
那些沒有看見就信的人.
是有伏敵」.
讀完之後 這麼多事情相識.
場景一樣.
部分對答一樣.

$^{761}$但多瑪不一樣.
因為是兩次的顯現.
給大家看看表.
這個表是對比第一次顯現.
和第二次顯現的時候.
是不是很多類同的地方.
嘩 你這樣看的時候.
都很特別看聖經.
是不是約翰福音的作者.
特別一些寫作技巧.
他精心寫的呢.
我覺得也是神精心去設計.
我大膽說.
耶穌第一次顯現的時候.
多瑪不在場.
聖經沒有交代多瑪不在場的原因.
但神總要觸摸人的心.
愛我們的神.
總會來到我們中間.
甚至在祂的面前.
正正就是在多瑪的面前.
讓祂經歷 讓祂認識.
多瑪這一刻.
經驗完主的時候.
相信祂由懷疑到經驗.
並且復活的耶穌.
成為祂的見證.
親眼見到耶穌.
不單止知道復活的事實.
親身和耶穌同場.
並且聖經沒有記載.
祂有沒有真的碰到耶穌.
耶穌叫祂來摸.
叫祂來碰.
叫祂來觸摸.
祂有沒有碰到呢?.
不知道.
你碰不碰呢?.
你夠不夠膽摸耶穌的釘痕呢?.
但這個觸摸體驗到.

$^{801}$耶穌是一個復活的生命.
不是一個假的.
是一個得性的生命.
並且那個疤痕是乾的.
應該是乾的.
用手去探入耶穌的性痕裡面.
深信多瑪的生命被扭轉.
祂的生命被割伸.
所以多瑪最後一句說什麼?.
祂宣稱,宣認.
我的主,我的神.
完全的降伏.
這次耶穌的顯現.
多瑪有沒有被激勵到呢?.
四十天之後.
耶穌就升天.
接下來應該是仕途行轉.
仕途行轉.
那扇門又聚集了.
後來他們從橄欖山回到耶路撒冷.
那扇靠近耶路撒冷.
有一安息日可就的路程.
他們進了城.
上了一間樓房.
就是彼得約翰雅各.
安德烈肥力多瑪.
巴多羅邁馬太.
阿勒肥的兒子雅各.
激進派的西門.
雅各的兒子猶大等人所住的.
有沒有多瑪?.
有啊.
仕途齊了 為甚麼?.
多瑪被激勵.
多瑪在這個群體裡面.
十一位的仕途齊齊整整.
還有一百二十人在一個地方裡面.
相傳是馬可樓一個大宅裡面.
一百二十人聚集.
再不是剛才.

$^{841}$剛才大家讀的經文.
關上門 隱隱蔽蔽.
這一刻耶穌顯現之後.
這群人有甚麼變化呢?.
這群人膽粗粗 膽大了.
並且再好像沒有怕了.
那麼多人聚集.
我們這裡都沒有百二人.
所以你想不想到.
原來一起見證的時候.
一起遇見神的時候.
有另外一些事情發生.
相傳多瑪在教會歷史裡告訴我們.
他去了印度傳福音.
今天還說在印度.
我去印度沒有發現這關.
我可能去得太少地方.
去了一間叫多瑪教會.
我想說到這裡.
剛才我說我給今天的題目.
是「你們的與我何干」.
我想見證是第一身的.
我們每個人都會有第一身的見證.
第一身的經歷.
但是別人的見證呢.
通常你會怎樣看呢?.
你有沒有真正的渴慕呢?.
有沒有嘗試.
「嘩!他成功了」.
「他經歷了耶穌」.
我有沒有這個渴慕.
還是我不值得呢?.
你的見證.
希望不是路遠的.
不是菜乾,不是鹹魚,不是臘肉.
雖然都很香,很好味,有價值.
但我希望你的見證.
是一個新鮮的見證.
你每天有的見證.
可以跟別人分享.

$^{881}$我最後一個見證跟大家說.
這個是我帶去信耶穌的一個女孩.
我現在榮升了.
我叫淑玲麻麻.
大家明白我說什麼嗎?.
我有多少代.
我其實不太知道.
我叫淑玲麻麻.
後來我見到一些人.
我才知道原來你帶她去信耶穌.
我才知道.
這個女孩在一間很小的教會裡服侍.
這間教會不能給她薪水.
所以她要做兼職.
她要做兼職.
然後去牧養這間教會.
這間教會雖然小.
但起碼有三至四個宣教士可以出來.
他們宣教的地方就是今天打仗的地方.
這幾個宣教士很用心去宣教.
這幾個宣教士有一個告訴我.
「我想明天回去」.
我們現在很想這宣教士回來.
她說「我很想明天回去」.
這個女孩去牧養這間教會.
我認識這女孩在牧養這間教會.
我認識她的時候.
我都不覺得她有什麼.
但我覺得耶穌就是改變她的生命.
她在牧養的宣教士.
她每兩年就去那地方去牧養.
其實去年五月她的經歷.
逃亡.
我為她擔心在她逃亡的裡面.
在這個過程裡.
雖然我被人叫做淑齡奶奶.
但其實我在這個女孩的生命裡.
我更加看到.
原來她服侍的.
和我有很多的距離.

$^{921}$我其實不應該叫淑齡奶奶.
她撇開了她的舒適區.
她過著簡樸的生活.
她在支持她牧養的弟兄姊妹.
她願意去到危險的地方.
去支援他們.
讓我看到的時候.
她的見證.
我也很想.
讓我知道.
其實她也在教我怎樣去牧養.
她也在教我怎樣去關心我的弟兄姊妹.
到位的關心.
還是用她的錢去關心.
一張機票不便宜的.
她每年去兩年就去一次.
她的見證是我有份.
所以弟兄姊妹.
今天來到這裡的時候.
我也很想.
如果當你每逢聽到別人的見證.
和我有關.
你可不可以和自己說.
她的見證和我有關.
因為我們教會是一個親體.
復活的主在我們當中做見證.
我們這個群體.
亦都是彼此相愛的群體.
主是真實的.
主是有復活的大能.
福音的大能.
在見證裡面.
我們學到對方的好東西.
並且我們可以一起感恩.
一起學習.
最後有一段經文.
《禹汗一書》四章十四節.
「父差子作世人的救主.
這是我們所看見且作見證」.
弟兄姊妹.

$^{961}$如果待會.
我不知道你們吃不吃飯.
你們互相認識的時候.
都說你的見證.
新鮮的見證.
讓我們都渴求.
讓我們都繼續在一個群體裡面.
互相去認證.
很重要的.
因為互相激勵.
剛才那個婆婆沒有坐輪椅上來.
其實帶給我們有很多人看見.
向你攜有很多弟兄姊妹提到.
他們都說「嘩!我都想啊!我都要啊!」.
有些是不夠膽走的.
一個是激勵.
一個是鼓勵.
\newpage



\section{詩篇 1:1-6}
\label{sec:031866}
\textbf{畫夜默想之福}~~~~高銘謙牧師博士
\newline
\newline
連結: \href{http://www.church.com.hk/acms/content.asp?site=cdc\&op=show\&type=product\&code=031866\&layout=sermon}{\texttt{031866}} ~~~~ 語音日期: 2026-03-22
\newline
\newline
\hyperref[sec:031865]{\small{< < < PREV SERMON < < <}}
~
\hyperref[sec:index_chronic]{\small{[返順時目]}}
~
\hyperref[sec:031867]{\small{> > > NEXT SERMON > > >}}
\newline
\newline
詩篇 1:1-6
\newline
\begin{longtable}{cl}
\hline
\hline
章節 & 經文 (和合本修訂版)\\
\hline
1:1 & \begin{tabularx}{0.7\textwidth}{X} 不從惡人的計謀，不站罪人的道路，不坐傲慢人的座位， \end{tabularx} \\ \\ \relax
1:2 & \begin{tabularx}{0.7\textwidth}{X} 惟喜愛耶和華的律法，晝夜思想他的律法；這人便為有福！ \end{tabularx} \\ \\ \relax
1:3 & \begin{tabularx}{0.7\textwidth}{X} 他要像一棵樹栽在溪水旁，按時候結果子，葉子也不枯乾。凡他所做的盡都順利。 \end{tabularx} \\ \\ \relax
1:4 & \begin{tabularx}{0.7\textwidth}{X} 惡人並不是這樣，卻像糠秕被風吹散。 \end{tabularx} \\ \\ \relax
1:5 & \begin{tabularx}{0.7\textwidth}{X} 因此，當審判的時候惡人必站立不住，罪人在義人的會眾中也是如此。 \end{tabularx} \\ \\ \relax
1:6 & \begin{tabularx}{0.7\textwidth}{X} 因為耶和華知道義人的道路，惡人的道路卻必滅亡。 \end{tabularx} \\ \\
[1ex]
\hline
\hline
\end{longtable}
$^{1}$我們同心祈禱.
慈悲仁愛的天父.
讓我們透過.
今天所讀到的詩篇.
對我們的人生.
有指引,有方向.
鼓勵我們的生命.
叫我們能夠不自陷.
在罪惡當中.
讓我們能夠死心革面.
是可以悔改.
並且去到擁抱.
耶穌基督的十字架.
主你幫助孩子.
講解你的說話清楚.
和明白我們一同.
去到謙卑領受.
奉耶穌基督的名字祈禱.
Amen.
今天很高興能夠再一次.
來到同福堂.
和大家一起分享上帝的說話.
今天所聽到的.
說話是來自.
詩篇的第一篇.
如果大家聽詩篇的時候.
通常第一篇.
懂得背誦的詩篇.
應該是詩篇二十三篇.
原因是因為它夠短.
只有六節.
但是另一篇可能比較多人.
懂得背誦的詩篇.
應該是詩篇第一篇.
同樣地它也只有六節.
讀給大家聽.
不從惡人的計謀.
不濫罪人的道路.
不助洩萬人的座位.
為喜愛耶和華的律法.

$^{41}$晝夜思想.
者人便為有福.
他要長一棵樹.
栽在溪水旁.
按時候結果指.
葉子也不枯乾.
凡他所做的.
盡道順利.
惡人並不是這樣.
乃將康比.
避風吹散.
所以當審判的日子.
惡人必不站立得住.
罪人在義人的會中也.
是如此.
耶和華認識.
知道義人的道路.
惡人的道路.
卻必定滅亡.
一篇很美麗的詩篇.
當我們聽這篇詩篇的時候.
有兩個背景.
讓大家知道.
第一.
就是你聽到.
「者人便為有福」.
這個字.
在和合本的翻譯裡面.
就放在第二節.
結尾的那裡.
如果你翻查希伯來文.
你就知道「有福的人」.
這個騙語.
應該放在第一節的開頭.
並且第一節的開頭.
的第一粒字.
應該用「福氣」這個字.
作為開始.
所以你會發現這篇詩篇.
就不是那麼負面.

$^{81}$用三個「不」去做開始.
它的主調.
是跟你說「有福的人生」.
是應該怎樣.
話說回來,它是詩篇的第一篇.
所以就存在在詩篇.
150篇裡面的引言的部分.
這個引言的部分.
極度重要,因為它主導了.
你怎樣看這個詩篇.
所以「福氣」這個字.
是第一,所以「福」.
就是我們今天很想談到.
「有福的人生」.
是應該怎樣的,這是第一.
第二,就是你會發現.
這篇詩篇常常對比.
「惡人」和「義人」的分別.
「惡人」怎樣.
「義人」怎樣.
就好像一個叫老口.
叫你作出一個抉擇.
究竟我應該走向「義人」的人生.
還是「惡人」的人生.
很明顯.
「義人」就是「有福」.
「惡人」就必定去到「滅亡」.
所以這一段的經文.
是用「有福」做開始.
用「滅亡」做結束.
就成為我們生命裡面.
的一種重要的抉擇.
這兩個背景.
是讓大家去到知道.
原來它是用三個不同的角度.
去到和我們解說.
究竟「惡人」.
和「義人」有什麼的對比.
第一個的角度.
就是「所在」.

$^{121}$也就是那個介詞「在」這個字.
是點擊率非常高.
所以「義人」.
在哪裡.
「惡人」.
在哪裡.
第二就是他們像什麼.
像什麼.
「義人」像什麼.
「惡人」又.
像什麼.
所以你今天聽到的時候.
就不需要看旁邊像什麼.
你就只需要.
去到看自己.
不要和你旁邊的老公說.
老公你是不是像這些.
是不需要的.
去想自己的生命.
這個是第二.
讓大家知道.
第三就是「惡人」的道路的方向.
是怎麼樣.
「義人」的道路的方向.
卻又是怎麼樣.
恰巧每一點.
都有兩節.
「所在」一到二節.
「所像」三到四節.
「所向」五到六節.
就構成今天講到的三點.
讓大家一起去到思想.
首先就是第一點.
「惡人」和「義人」的對比.
就是在於.
他的「所在」.
原來你會發現.
「在」我用紅色highlight.
總共出現六次.
很公平的.

$^{161}$三個是「惡人」.
三個是「義人」.
這個是第一樣.
讓大家去到知道「所在」.
首先我們會知道.
剛才講過.
原文的一開始.
就是「那人的福」.
又或者「有福的人心」.
去到做開始.
那個「人」字.
如果你翻查希伯來文.
你就知道它是單數的.
單數的「有福的人」.
但你一看那個螢光幕.
那個原文的直譯.
你就會見到.
全部「惡人」都是眾數.
「惡人們」.
「罪人們」.
「犯罪集團」.
整群人.
都是「惡人」.
但「無福的人」就只有.
「一個」.
你能不能回到辦公室.
你就是那個關鍵的「一個」.
然後整群人.
都在你面前「掩掩養養」.
全部都是.
「惡人」.
會不會是這樣呢?.
事實上我們正在面對.
正在身處一個時代.
是一個罪惡滿盈的時代.
稍一不慎.
你就會有人說你「是非」.
稍一不慎.
你就會偷呃拐騙.
殺人金腰帶.

$^{201}$你會發現這個世界是極度的險惡.
但是我們作為基督徒.
很多時候.
就只有「一個」.
面對這樣的情況.
但是不要緊.
只有「一個」都不要緊.
因為你是「蒙福的人」.
你才是關鍵.
我們先看看「惡人」那邊.
「惡人」那邊你會見到有幾個名詞.
是很重要的.
首先第一個名詞.
就是「計謀」.
然後下一個就叫做「道路」.
然後就是「座位」.
很明顯你會見到一個「漸進式」.
通常你犯罪是怎樣呢?.
就有一個「idea」.
即是有一個「意念」.
有一個「計謀」.
你想到然後你走出來.
就是「道路」.
然後你久而久之.
不斷走出來.
那個罪惡.
就成為你的生活方式.
然後你就「坐埋落去」.
「坐埋落去」那個「座位」.
除了可以解作「座位」.
可以解作你的「居所」.
又或者可以解作你的「生活方式」.
有沒有見過一些人.
說謊是真的.
當吃生菜.
成為他的生活方式.
就成為慢慢墮陷.
泥足深陷的一步.
兩步和三步.
這個叫「計謀」.

$^{241}$「道路」和「座位」.
不得止,我們也可以分析他的動詞.
他就說「不從惡人計謀」.
那個「從」字.
其實是解作「行走」.
這樣的意思.
然後不站罪人的「道路」就是「站立」.
然後就是「座位」.
就是「坐埋落去」.
由「行走」然後「站立」.
然後「坐埋落去」.
又是一個漸進式.
就好像你行下街.
見到一條冷巷.
然後那裡有些誘惑.
你就站在那裡.
看一看.
受到誘惑.
然後就「坐埋落去」.
「行路」「站立」.
就「坐埋落去」.
就成為一個漸進式的效果.
慢慢泥足深陷.
墮落萬丈深淵.
就是一個這樣的路程.
我們叫這個「惡人」.
「罪人」和「洩萬者」的路徑.
就是這樣去到「行走」.
反之.
「義人」是怎樣的呢?.
「義人」相對簡單一點.
你看看「惡人」那裡.
都頗複雜的.
又說「惡人」又說「罪人」又說「洩萬者」.
又說「道路」又說「計謀」.
又說「行走」又說「站立」.
搞到你五花八門.
非常之眼亂.
但是「義人」.
就只有一個非常簡約的焦點.

$^{281}$那個焦點.
就是耶和華的律法.
你見到的.
第二節「律法」那個字.
總共出現兩次.
並且用兩個不同的動詞.
去到描述怎樣對待耶和華的律法.
第一個動詞.
就叫做「喜愛」.
為喜愛耶和華的律法.
「喜愛」原來這個字.
是解作「答堂」這樣的意思.
你先看著你的老公.
你覺不覺得「答堂」.
看著你的老婆.
如果你覺得「答堂」.
為何你轉頭.
即是.
「答堂」.
即是告訴我們.
就是.
我們對於耶和華的律法.
甚至要讓你的老婆「答堂」.
原來.
詩篇十九篇已經告訴我們.
耶和華的律法.
全被能甦醒人心.
耶和華的法道.
確定能使漁人有智慧.
然後他用了兩樣東西.
去到進行一個對比.
且比金紙可羨慕.
且比蜂房.
下滴的墨.
去到羨慕.
大家聽到第一個對比就是「金」.
金牌,有沒有買「金」?.
你喜不喜歡「金」?.
喜歡.
我們當然很喜歡.

$^{321}$今天我鼓勵你.
你要比「金」更喜歡的.
就是耶和華的律法.
耶和華的說話.
他的說話.
不只是給你錢那麼簡單.
他的說話是引導你步向.
永恆和永生.
賦予你一個很強勁的生命力.
這是比金紙可羨慕.
且比蜂房.
下滴的墨更加羨慕.
因為在迦南地裡.
被稱為「流奶與墨之地」.
所以最好的物體.
就是蜂蜜.
而蜂蜜是當時.
最甜的一個物質.
所以你要比這兩樣東西.
更加寶貴.
這就是喜愛「答堂」的意思.
第二.
就是晝夜蜜賞.
不說大家不知道.
原來蜜賞這一種的希伯來文.
可以形容兩個動物的叫聲.
第一個動物.
就是甲子.
我住屯門.
經常在屯門河邊跑步.
跑步途中.
會見到一位叔叔.
經常在橋底餵那些白鴿.
他餵的時候.
我經常聽到那些白鴿.
發出一些聲音.
那些聲音就是「咕咕」.
「咕咕」.
那些聲音.
所以如果你看到旁邊的姊妹.

$^{361}$去到晝夜思賞上帝的說話.
她其實是在「咕咕」.
「咕咕」.
這是第一個意思.
第二個意思.
就是獅子的咆哮聲.
如果你看到旁邊的弟兄.
在晝夜思賞上帝的說話.
他其實是在.
咆哮中.
那為什麼經文.
要用這個字呢?.
這是一個很特別的希伯來文.
就是因為.
白鴿自然會.
發出「咕咕」聲.
獅子會自然.
發出咆哮聲.
敬畏耶和華.
義人就一定.
不用審.
是晝夜默想和思賞上帝的說話.
所以這是義人.
一個關鍵的特徵.
那個特徵不是很複雜.
是一個簡約的聚焦.
耶和華的說話.
成為你答堂的東西.
成為你的喜愛.
成為你的實踐.
成為你的價值觀.
成為你.
去到行善的基礎.
所以你會發現耶和華的說話.
成為他的焦點.
這樣的話,就在第三節.
就有另外一個「再」.
就是在上帝的時間下.
他就能夠結出果子.
我講一個假設的故事.

$^{401}$給大家聽.
話說有一個小學生.
他非常喜愛吃朱古力.
喜愛到一個.
極速深陷的地步.
所以每逢放學的時候.
經過七十一.
他就一定例牌.
進去買兩排KitKat去吃.
那他的媽媽呢.
面對一個這樣的情況.
她就不知如何是好.
她就鼓勵她的兒子.
如果你今天放學回來.
你就一條直路不進7-11.
我就獎你一個.
很好的禮物給你.
然後就放學了.
一放學.
就例牌經過7-11.
然後就去到門口.
他就行走.
然後就站立.
然後就跟自己說.
我站著而已.
我沒有想過進去.
然後他就進去了.
我進去而已.
我沒有想過走去那一格.
然後他就走去那一格.
並且拿起兩排KitKat.
還是抹茶味的那一款.
然後就跟自己說.
我沒有想過給錢.
然後他就走去收銀那一格.
還給錢.
就跟自己說.
我給錢而已,我沒有想過拆開.
然後他就拆開了.
然後就說.

$^{441}$我拆開沒有想過放進口.
然後他就放了.
我放而已,我沒有想過吞.
一步一步.
就泥足深陷.
由行走到站立.
然後坐下去.
你知道的.
我們每個基督徒.
如果你要犯一個嚴重的罪.
你一定不會一步到位.
通常.
就是陰一點.
有一位弟兄.
他是一個敬拜隊的成員.
他領場.
放心,這位弟兄不是屬於同福堂的.
他有太太.
有一對很可愛的小朋友.
然後.
每一晚練歌都練到很晚.
練完歌之後.
又不知道為什麼.
最近都跟那個詩琴姐姐.
一起去吃糖水.
我想問.
吃糖水有沒有道德問題.
沒有的,吃糖水而已.
然後有一個弟兄.
見他經常都是這樣.
就提醒他,溫馨的提示.
愛心說誠實話.
就跟他說.
不如你早點回去陪老婆.
不要再跟姊妹吃糖水了.
然後這位弟兄很大反彈.
他就說.
你當我是什麼.
沒什麼.
通常說沒什麼就是有的.

$^{481}$然後.
他們吃下又一晚.
講下公事.
下一晚就講私事.
然後就講心事.
然後就大件事.
你會發現.
慢慢培養.
就成為愛情.
談心事沒有道德問題.
吃朱古力.
沒有道德問題.
你是可以吃.
不過你知道的.
通常你搭地鐵的時候.
你要墮陷在月台之前.
是會有一條黃線的.
穿越黃線.
會不會墮落月台.
其實通常不會的.
我們叫做麻尋位.
有一個距離.
但是如果你走著走著.
站著的時候要小心.
因為那個7-11.
你一定要兩步路快點走.
要保持一個.
安全的距離.
千萬不要讓那個誘惑.
和罪惡成為你的生活方式.
這是第一點.
讓我們知道.
秘訣是什麼.
秘訣就是晝夜思想上帝的說話.
不單止是妄想.
並且成為你生活.
言行舉止的一個實踐.
這是第一點.
讓大家知道.
第二就是像什麼.

$^{521}$你會看到.
經文所講的.
其實是完全不公平的.
異人像什麼.
總共有四句.
惡人像什麼.
就只不過一句.
既然是這樣.
惡人其實是像什麼.
像漢鼻.
是什麼來的.
就是你如果打蜜紙的時候.
那些比較重的蜜紙.
就放在了燒機上面.
然後你就會發現很多水濕濕.
在那個燒機上面.
那你就怎樣分隔.
你就將那些蜜紙.
向天揚一揚.
揚一揚的時候.
你帶風.
所以那些比較輕的漢鼻.
蜜紙就掉到那個燒機裡面.
就從而分隔漢鼻和燒機.
一個非常有效的方法.
這樣.
你就會發現用這個比喻.
去比喻惡人.
他的比喻的焦點是什麼呢.
焦點就是惡人沒有金.
這個是重要的.
第二 惡人.
好像漢鼻在空氣上面揚起.
比喻做非常高調.
奄奄養養.
凸凸戚戚.
很多時候都在你面前.
整天都在你面前出現.
但是不需要羨慕.
不要覺得他們厲害.

$^{561}$因為他們好像一陣風.
就完全.
過多一兩秒.
就化為暴風和去空.
這個就是一個這樣的比喻.
這樣的比喻是要對比異人.
異人是什麼呢.
是一棵樹.
而這棵樹栽植在溪水旁.
和合本也在溪水旁.
你翻查希伯來文.
他不是說在溪水旁.
是在溪水的上面.
你見我原文的直譯.
就是在溪水的上面.
所以就代表他的根部.
是非常直接地.
深深抓住溪水的泥土.
以至他直接地.
去吸取養分和水分.
這是一個比喻.
比喻晝夜思想上帝的說話.
就是你的養分和水分.
去到吸住的時候.
以至你產生兩個後果.
第一.
你就按時候結果治.
第二.
你的葉子也不會枯乾.
大家覺不覺得.
你在家中植那棵植物.
你要搞到你的葉子不枯乾.
是一個極度困難的事情.
又說那些酸鹼度.
又要協調.
又說那些水不能淋得多.
又說淋得少.
是複雜的.
但種種這些的設定.
只要你種植在溪水旁.

$^{601}$全部都是剛剛好.
原來他的比喻是.
根部非常低調.
不會在你面前張揚.
並且他的泥土.
是下面來的.
就不會讓你看到是怎樣.
反之.
他是非常有實質的果效.
因為那個叫果子.
果子你知是一個比喻.
祝靈的比喻.
聖靈所結的果子.
是有人愛喜樂和平.
忍耐恩慈良善.
信信溫柔.
和節制.
你會發現就算很多的沙塵暴.
去到吹起.
他仍然都能夠結果子.
就算很多苦難臨到.
仍然都可以結出.
溫柔喜樂.
愛心的果子.
這一種的生命力.
是由一條河流給你.
話說這條河流有一個典故.
就是來自以西結書的47章.
裡面有提到聖殿石安山上面.
就有一個祭壇.
東邊所流出來的.
生命河.
大家有沒有聽過.
這條河流經的各處.
百物都生長.
然後很多魚類在裡面去到暢泳.
然後有使鹽水.
去到變甜.
這一個就是以西結書的47章.
然後耶穌基督.

$^{641}$就在約翰福音的第七章裡面.
就引用了這個.
舊約的典故.
你聽過的七章三十七節.
就說人若學了.
可以到我這裡來.
凡信我的人就必定.
由他的福宗裡面流昌.
活水.
和江河來.
耶穌用了這個比喻.
告訴你相信耶穌的時候.
就好像將你的生命.
接駁去石安山那條河流一樣.
就好像今天.
詩篇第一篇這個比喻.
你就會看到那個樹木.
就紮根在那條.
生命河裡面.
你就湧入出那個果子.
我認識.
一位姊妹.
她很不幸.
醫生確診.
她是一個第四期的胰臟癌.
末期的了.
所以醫生就跟她說.
這三個月的日子.
當你聽到這個消息.
其實都很坎坷.
很絕望.
但是這個姊妹是基督徒.
她就跟自己說.
既然我三個月之後.
我就會見耶穌.
倒不如好好把握.
這三個月的機會.
就將這個福音.
這個福氣.
拼命介紹給人去認識.

$^{681}$起碼三個月之後.
我都有面見耶穌.
她這樣想.
然後有一次的併發症.
召了八車進了醫院.
整個病房.
全部都是一個癌症的病人.
整晚.
如果你住過醫院.
有時都睡不著.
可能隔壁那個有嘈.
又會跟你量血壓.
很難睡覺.
她反而.
在一個這麼低氣壓的病房裡.
周圍關心.
身邊的病友.
走去關心她.
聆聽她.
又說要為她祈禱.
她的喜樂和愛心.
就好像燈光了整個房子.
以致那晚.
一晚.
她就帶了三個人順著.
一晚帶了三個人順著.
我們就要想一個問題.
那個問題就是.
你患癌症.
為什麼患成這樣.
意思是.
為什麼可以在惡劣的環境下.
仍然能夠結出愛心.
喜樂的過節.
那個秘訣就是那個基礎.
就是那個根部.
能夠抓緊心心的泥土裡.
讓上帝的說話.
成為你的養分和水分.
鼓勵你的人生.

$^{721}$賜給你有盼望.
給你有愛的動力.
讓你棄惡.
行善的.
抉擇和力量的根源.
就是我們主.
所啟示寶貴上帝的說話.
這是第二點.
讓大家知道.
第三點.
第三點就是所向.
惡人們.
和義人們的道路.
他的方向到底是怎樣.
然後你就會看到紫色highlight.
那個義人們.
你就會看到這次.
就不是單數了.
這次是重數.
我們.
每一個孤獨的單數.
聚集一起成為一個聚會.
個個都回到.
彌童福堂.
然後你這次真的可以看一看.
旁邊那位弟兄姊妹.
一看就知道是義人.
他們每一個人聚集在這裡.
彼此同心.
彼此鼓勵.
彼此聚集在一間的教會當中.
能夠彼此鼓勵.
如果是這樣.
你就同心走上.
這條蒙福的道路.
不過他都說了惡人的結局.
就是怎樣.
惡人的結局是怎樣呢.
在審判的過程當中.
他是必不站立得住.

$^{761}$這個是和合本的翻譯.
但如果你看回動詞的時候.
你就會發現.
那個動詞原來不是解作.
站立,是應該解作.
不起來,你看我那個直譯.
他不要說站立.
他連站起來都站不起來.
就毋須說站立.
所以是更加嚴重.
就是在上帝的審判下.
那個惡人站起來.
都站不起來.
這是第一點讓大家留意.
然後道路那個字.
就出現在第六節的.
兩次出現.
分別就是義人的道路.
是被主上帝深深認識的.
義人的道路有時孤單.
耶穌說.
義人的道路是進閘門.
惡人的道路是看莊大道.
卻步向滅亡的路.
進閘門的人是少.
去看莊大道的人是多.
但是那條的道路.
惡人的道路卻必定滅亡.
進閘門的道路.
是要步向永生.
有時是孤單,有時是艱難.
有時你會疑惑.
但是上帝的應許是什麼呢?.
上帝的應許是.
上帝一定知道.
一定明白.
一定看顧.
上帝知道你的委屈.
上帝知道你是堅持到底的那個.
上帝知道你的心地是好的.

$^{801}$上帝知道.
你的決意不會以惡報惡.
這條的道路.
是你會面對很多的敵人.
有人會說你是非.
有人會戳你的背.
有人以惡報惡.
如果你被誘惑.
別人說你是非.
你又說他三個是非的時候.
會不會真的以惡報惡.
而你就成為那個惡人.
會不會我們倒轉頭.
學習保羅的教導.
就是以善去成惡.
而不是以惡報惡.
就要確定你是站在義人那一邊.
你不是成為惡人那一邊.
別人說你是非.
我們可不可以請他吃自助餐.
所以我們.
這是一個比喻.
是一個不同的邏輯和思維.
我們是天國的子民.
我們的基調是愛.
是善.
是寬恕.
一個這樣的基調.
你就會明白這個世界的邏輯.
被打破.
以至教會成為一個光明的見證.
可以照在黑暗的當中.
不用審了.
惡人的道路已經說了.
是必定滅亡.
是不需要心存僥倖.
是一定會滅亡.
我們倒轉頭.
倒轉頭就是康比的比喻.
被封.

$^{841}$去到吹散.
無論他怎樣走.
打橫行也好.
打斜行也好.
都會步向滅亡.
這個終末觀.
就決定你現在的倫理觀.
我再說一次.
你相信一個什麼的未來.
是影響你現在有什麼的抉擇.
如果你相信那個未來.
是神國的想像.
永生的盼望.
那種新天新地.
那種耶穌.
的敬拜.
你就肯定.
不會成為惡人.
因為惡人的道路.
終末是必定滅亡的時候.
因為這個終末的結果.
你現在就會作出這個決策.
道理不是很複雜.
因為我每天住屯門.
我上班的那間神學院.
說實話我是在建道神學院工作的.
接下來九月就要成為院長.
請大家多多為我祈禱.
不過.
anyway.
每天上班.
我都要搭一架巴士叫961.
我每天都去屯門長洲.
屯門長洲去到進發.
每次去九巴apps.
去巴士站.
你一定指定動作打開那個apps.
一定不知道為什麼第一板是廣告.
然後你就看.
剩下兩分鐘.

$^{881}$就來的了.
那架巴士你要不要等.
當然要等.
有時候.
都等得很辛苦.
又不知道為什麼夏天的時候.
那個九巴站的頂.
是用玻璃做的.
有沒有見過.
不知道為什麼那個光線聚焦在你的頭頂.
是會有一點很生氣.
但無論多辛苦.
剩下那五分鐘.
你一定要等.
因為你相信那個巴士一定會來.
所以影響你.
現在就一定會等.
是不是這樣.
多辛苦你都會等.
因為你有盼望.
你是一個終末的盼望.
但如果你不相信那個巴士.
你一定不會等.
因為你知道那個巴士.
一定有冷氣.
所以你等得久都不要緊.
因為想到去就有得歎.
就是一個這樣的情況.
同樣的道理.
你相信一個什麼的未來.
決定你現在選取的方向.
到底是怎麼樣.
有很多時候我們掛在後面.
我們相信我們有永生.
我們知道.
未來終末的盼望是怎麼樣.
你說你相信.
但很多時候我們的生命.
到你選取那個抉擇的時候.
我覺得你不是相信.

$^{921}$就因為你仍然.
用世俗的價值觀.
那種你如我渣.
那種行事路徑.
然後一個大話.
蓋一個大話.
那種罪惡的世俗.
你以為很聰明的那一招.
但誰不知.
這些全部都會化為.
虛空和塵土.
然後其實你是.
不相信有這個終末觀.
如果你相信.
你就會選取二人的路徑.
去學習步向這個.
永恆的盼望.
然後你一定打死.
我都不會成為惡人.
因為那個惡人.
一定要步向這個滅亡.
有機會剛剛過了12月.
我就去了土耳其.
第一個旅程.
一定就去探訪.
那個叫啟示路七教會.
去到一個叫.
別加摩的古城.
如果你有印象.
就在啟示路裡面.
形容別加摩這個城市的時候.
你就會見到有個形容.
叫做撒旦的座位.
就是經文所說.
你一聽到撒旦的座位.
很明顯就是邪惡族心.
那裡.
就是一個鬼王的座位.
簡直.
真的很邪.

$^{961}$但是.
我想想什麼叫撒旦的座位.
原來它是指一個.
愛邦神明的殿.
那個殿就是埃及.
這個希臘.
最高的神明叫做宙斯.
他的宙斯殿.
就坐落在別加摩城裡面.
你見到這個螢光幕.
叫做重構圖.
那個殿.
那個殿的樣子.
很像一個座位.
都很像.
所以聖經絕對是不會.
去到.
提別魂的名號.
所以他不提宙斯這個名.
叫他做什麼.
撒旦就可以了.
這個是新約的立場.
你知道.
這些聖殿一定香火鼎盛.
一定是很多人.
去到供奉.
並且非常多人.
很多利益在那裡輸送.
但是我去到土耳其的時候.
我就很想看看.
這個撒旦的座位.
誰不知就是右邊那幅相.
是我拍了出來.
這個城地.
只有幾級樓梯.
其實是為什麼呢.
因為德國的考古學家.
將整個殿搬到柏林.
德國人很流行這些.
如果你去柏林.

$^{1001}$你見到什麼呢.
柏林是你見到以前利布格尼薩王的古巴比倫城.
整個城就去了柏林.
不過無論如何.
不要提柏林了.
我們提回這個.
我的心很觸動.
很多的感受.
就是因為以前.
顯赫一時的撒旦的座位.
惡人的座位.
到現在.
移為平地.
化為虛空.
這個歷史的教訓.
讓我們能夠看到.
惡人的道路.
一定滅亡.
然後你見不見到這個平地.
剩下兩棵樹.
令我都想起今天的京文那棵樹.
我們同一個地點.
結出的果子.
一定是永恆.
一定是長存.
所以我們千萬.
不要羨慕成為惡人.
一定要跟隨.
耶和華的腳蹤.
成為那個義人.
所以今天三點.
很想送給大家.
有福的人生的三個方向.
第一.
所在到底是什麼.
究竟你成為那個惡人.
慢慢墮陷的路徑.
還是成為一個.
簡約的聚焦.
就是聚焦在上帝的說話裡.

$^{1041}$這個就成為義人.
和惡人的分別.
第二,你葬什麼.
沒有根的康比.
還是有根抓住泥土.
吸取養分的一棵果樹.
這個完全是.
你的抉擇.
第三,他的路徑.
和他的方向.
一個是滅亡,一個是永生.
一個是寬闊的路.
一個是盡責門.
你今天要.
作出這個抉擇.
祈求主的話.
成為我們的提醒.
督責,勉勵.
和成為我們的盼望.
我們一起低頭.
仰望上帝.
當你閉上眼.
的時候.
很想有一兩分鐘的時間.
你自己私下.
和耶穌有教堂.
有禱告,有默想.
回應上帝.
今天對你的生命.
所發出的訊息.
有一兩分鐘的時間.
安靜,隨著琴音的播奏.
我們有一段.
回應上帝的時間.
安靜.
隨著琴音的播奏.
我們有一段.
回應上帝的時間.
安靜.
隨著琴音的播奏.

$^{1081}$我們有一段.
回應上帝的時間.
安靜.
隨著琴音的播奏.
我們有一段.
回應上帝的時間.
安靜.
隨著琴音的播奏.
我們有一段.
回應上帝的時間.
安靜.
隨著琴音的播奏.
我們有一段.
回應上帝的時間.
安靜.
隨著琴音的播奏.
我們有一段.
回應上帝的時間.
安靜.
隨著琴音的播奏.
我們有一段.
回應上帝的時間.
安靜.
隨著琴音的播奏.
我們都每一個闖開我們的心.
來到上帝面前.
今日是一個.
立志去到將你的.
生命奉獻的日子.
如果你聽到這個訊息.
覺得窄心.
你今日真是決意.
決意走上.
敬畏上帝.
晝夜默想上帝的說話.
蒙福的道路.
我真是好想.
我想你輕輕舉起你的手.
好想為你祈禱.
有沒有這樣的弟兄姊妹.

$^{1121}$這個舉手不是向人舉.
是向上帝舉.
主啊我願意立志.
走上蒙福的道路.
請放下手.
我們一起起立.
領受從神而來的賜福.
親愛的天父.
求主你去到大大賜福.
每一位的弟兄姊妹.
特別剛才.
他們有很多都願意.
舉手走上.
這條蒙福的道路的弟兄姊妹.
求主你賜福給他們.
牧養他們.
讓他們意志堅定.
讓他們選上.
上好的福分.
讓他們一生成為.
行善的義人.
晝夜思想上帝的說話.
讓你的話成為他們的提醒.
成為他們的價值觀.
成為他們人生的指引和方向.
以致他們不會走差.
漫漫墮陷.
成為一個惡人.
反之.
要成為一個見證者.
成為見證你福氣的一個門徒.
跟隨耶穌基督的門徒.
主你大大賜福.
給在座.
這裡每一位的弟兄姊妹.
又願我主耶穌基督的恩惠.
天父的慈愛.
聖靈的感動與團契.
常與你們同在.
讓你們一生堅定你們的腳步.

$^{1161}$走上蒙福的道路.
拒絕惡人的道路.
從今時.
直到永永遠遠.
阿們.
\newpage



\section{創世記 45:5}
\label{sec:031867}
\textbf{聖經人聖經事:是上帝差我先來}~~~~黃克勤牧師博士
\newline
\newline
連結: \href{http://www.church.com.hk/acms/content.asp?site=cdc\&op=show\&type=product\&code=031867\&layout=sermon}{\texttt{031867}} ~~~~ 語音日期: 2026-03-22
\newline
\newline
\hyperref[sec:031866]{\small{< < < PREV SERMON < < <}}
~
\hyperref[sec:index_chronic]{\small{[返順時目]}}
~
\hyperref[sec:031868]{\small{> > > NEXT SERMON > > >}}
\newline
\newline
創世記 45:5
\newline
\begin{longtable}{cl}
\hline
\hline
章節 & 經文 (和合本修訂版)\\
\hline
45:5 & \begin{tabularx}{0.7\textwidth}{X} 現在，不要因為把我賣到這裡而憂傷，對自己生氣，因為神差我在你們以先來，為要保全性命。 \end{tabularx} \\ \\
[1ex]
\hline
\hline
\end{longtable}
$^{1}$各位弟兄姊妹大家好.
很快又一天了.
又一齊看聖經了.
今次的聖經人,聖經事仍然是講約瑟的.
現在應該還有一兩次是講約瑟的.
希望我們可以透過聖經人物和他的故事裡面.
看到上帝是怎樣的.
多些認識他和看到他背後.
上帝的手在工作.
這個就是生命說書的網絡為大家提供的.
希望我們多讀聖經.
我們開始了.
今次的題目是.
上帝差我先來.
留意一下.
是誰的手在這裡工作.
是上帝的手.
他做什麼呢?.
是差遣我先來到這裡.
原來這句話出自約瑟的.
我們看看.
創世記四十四章.
這件事是很峰迴路轉的.
幾句背景交代一下.
由於在中東那塊地方.
當日雅各約瑟的時間.
在尼羅河邊.
故事是這樣的.
在尼羅河邊有七隻肥牛在河邊吃草.
然後有七隻瘦又醜的牛.
在那裡吃了那七隻肥牛.
發到第二個夢就是.
有七棵和水.
應該是墨水.
很肥很重地垂下來.
很漂亮.
然後有七棵被東風.
即是西那沙漠吹來的熱風吹乾了.
而且墨殼裡面是沒有墨紙的.
是空殼來的.

$^{41}$很瘦.
於是吃了那七棵肥的墨紙.
「哎呀!什麼事呢?」.
驚醒後沒有人解答.
最後酒證記得約瑟在監獄裡跟他解夢.
就叫了約瑟來.
約瑟來就「嘩!很厲害啊!解到啊!」.
然後就叫了約瑟做全國宰相.
預備七年的食物.
而在雅各.
即是他們住在迦南地.
都一樣饑荒.
沒有糧食了.
於是雅各說「哎呀!餓死了」.
就叫他十個兒子不如去埃及扔糧.
如果舊一點的和合本.
是用我們以前很舊的字叫「扔糧」.
即是去找糧食回來買糧食.
事實上哥哥就去了.
去到的時候見到約瑟.
約瑟就認到他們.
過程峰迴路轉.
大家自己看《聖經》.
創世記最後那幾章.
約瑟就作弄他們.
賣了給他們.
又把錢放回去.
沒有收他們的錢.
嚇到他們.
糧食吃完之後.
又要去扔糧了.
雅各說「你帶回那些錢去告訴別人.
上次他們忘記了收錢」.
誰知那個負責人就說.
「上次的錢我們收了」.
「這次要做什麼?再扔糧?」.
「好吧!請進來」.
「我們宰相想請你吃飯」.
嚇到那十個人.
「十兄弟死了」.

$^{81}$「這次是不是要殺死我們?」.
「要做什麼?」.
終於送了糧食給他們.
賣給他們.
然後又把自己喝酒的杯子.
放在辯雅敏的包袱上.
「上次你帶了最小的弟弟辯雅敏來」.
於是他們就帶來了.
幾經周折之後帶來了.
然後又派兵去追他們.
「你為什麼來扔糧偷我們宰相喝酒的杯子?」.
嚇到他們死.
如果是誰偷的.
那個人就永遠在這裡做奴隸.
抄在辯雅敏的包袱上.
於是一起回到辯雅敏的頭上.
於是就說辯雅敏留下來做奴隸不行.
回到辯雅敏的頭上.
爸爸真的哭死了.
他好像已經沒有藥室了.
如果辯雅敏沒有藥室.
我爸爸一定會死的.
於是這個時候來到第十八節.
尤大就靠近他.
靠近藥室.
靠近一點點也不太近.
「你是誰?」.
「我主求你容僕人說一句話」.
「說一句話給我主聽」.
「不要向僕人發烈怒」.
「因為你如同法老一樣」.
「不要向我發憤怒」.
「我想跟你說一句話」.
「不過你不要憤怒」.
「因為你憤怒我就沒命了」.
「因為你和法老是一樣那麼大」.
「因為僕人曾經向我的父親」.
「為這童子辯雅敏作一個擔保」.
「我若不帶他回來交給你父親」.
即是我的爸爸.

$^{121}$「我便在父親面前永遠擔罪」.
「我的罪過是永遠不會得到赦免」.
「所以既然是這樣」.
「現在求你容僕人註下」.
這個僕人是他自己.
「現在求你宰相大人」.
「讓我在這裡住」.
「第一 替這童子作我主的奴僕」.
「讓童子辯雅敏和她哥哥可以回去」.
「如果童子不和我們同去」.
「我怎能去見我父親呢」.
「恐怕我看見災禍臨到我父親身上」.
即是如果爸爸見不到最後這個兒子.
回去他便很慘他便會死.
我們是擔保他一定回來.
特別猶大說「我擔保他一定回來」.
「不如我留在這裡」.
「替這童子作我主的奴僕」.
這個奴僕不是指傭人有糧出.
是奴隸.
當時的奴隸很慘.
奴隸是沒有鞋穿的.
過了二千年都是沒有鞋穿的.
為甚麼呢?.
因為在主耶穌講過一個比喻.
是浪子比喻.
浪子比喻其實有很闊的內容.
其中一件事是小兒子回頭時.
爸爸說「給他穿衣服」.
「給他戴戒指」.
「給他穿鞋子」.
因為奴隸是不穿鞋子的.
給他穿鞋子即是他是家人不是奴隸.
連鞋子都不能穿.
這個當時的奴僕多慘.
地上又凹凸凸的.
我很欣賞猶大的.
猶大跟約瑟說.
「可不可以讓我講句說話」.
「不如我留下」.

$^{161}$「替謝童子做我主的奴僕」.
「然後希望你放他回去」.
「否則我爸爸就很慘了」.
猶大這個人.
猶大是誰呢?.
其實就是他賣約瑟的那個.
最大的大兒子劉便.
或者有些後來的聖經叫做呂便.
想說「不如我們扔他到坑裡」.
「然後看看怎樣吧」.
「不要殺他」.
扔了到坑裡.
後來猶大說.
「不如有商女經過」.
「我們賣了他」.
「又可以有些錢分」.
「又不用殺他」.
「我們不要動手殺他」.
「賣了他」.
是猶大說賣他的.
到現在.
其實猶大都不知道.
面前應該就是約瑟.
不過這裡他又有些義氣.
其實他都想保住約瑟當日的性命.
於是一定要有東西交換.
「不如我們留下他的性命」.
「賣了他給別人」.
「我們有多少錢就好了」.
「那怎樣交代呢」.
「將他建彩醫治血吧」.
猶大都算是一個.
想救他弟弟約瑟的人.
到這裡猶大都挺好的.
他就說「不如這樣吧」.
「由我來替這個童子做主的奴僕」.
「放他回去」.
「不然我爸爸很慘」.
「一定死的了」.
「哭得哭不死」.

$^{201}$猶大這個人都挺孝順的.
猶大這裡不錯.
不過他曾經賣約瑟.
但聖經是說他想保存他的生命.
所以猶大這個都值得我們尊敬的.
第一就是他替.
我代替.
因此猶大由人類犯罪之後.
我們的神預備救贖.
選了一個人開始.
就是亞伯拉罕.
而亞伯拉罕由他開始.
救贖的線就去到主耶穌.
亞伯拉罕之後有以撒.
以撒就背著柴上山的信服.
代替羊羔做羔羊.
可以獻上自己.
他就把木頭掛在自己的身上.
象徵主耶穌和祭神.
然後到雅各.
雅各開始整個族轉變.
叫以色列.
有沒有叫亞伯拉罕人呢?.
沒有,叫以色列人.
由他開始一個國家一直下來.
由雅各開始下面誰承接他.
關於救贖的線就是猶大.
猶大一直到大衛大衛到耶穌.
原來猶大在救贖的線上.
是一個上帝選擇他的人.
在這裡他也有表現.
他代替.
這個代替.
多多少少象徵到.
我代替我弟弟.
我這個愛的人可以向父親交代.
這也是相當之象徵主耶穌的.
猶大替同子作主的奴僕.
這點跟大家特別先講.
看看下一章,第45章第一節.

$^{241}$後來約瑟跟他們說.
其實約瑟好幾次忍不住.
看到哥哥在這裡.
心裡面很想哭.
又很想大家相認.
但是還不行.
有時走進自己房間哭一會.
擦一擦臉才出來再說話.
這次約瑟撐不住了.
到了這關頭.
45章第一節.
約瑟在左右站著的人面前.
情不自禁想哭想認.
於是吩咐一聲.
人都要離開我出去.
所有人出去.
於是那些服侍的人全部出去了.
約瑟和弟兄們相認的時候.
並沒有一人站在他面前.
所有埃及人都出去了.
他就放聲大哭.
約瑟也織了很久.
放聲大哭大聲到怎樣呢.
真的大哭.
因為埃及人和法老家中的人都聽見.
這句有沒有誇張了呢.
我看得有少少.
埃及人是指什麼人呢.
就是事後的那些人.
所有人都出去.
那些人是空對耳的.
是一班外人.
不知道有沒有害.
叫我們出去.
主人是否真的要出去.
是的出去吧.
所有人都出去.
如果你是那些事後的人和衛兵.
你會怎樣.
當然是出去.

$^{281}$記在門口的門縫裡聽著.
所以他們聽到約瑟放聲大哭一點也不奇怪.
法老家中的人都聽見了.
應該不是哭聲傳得那麼遠.
因為當時地方很大.
和法老王的居所都住很遠.
應該是嚇死了.
於是有些下屬馬上跑去.
報告大王.
約瑟哭了起來.
和十個人在那裡哭了起來.
他們好像有些特別的關係.
聽著吧.
所以我推測這些埃及人是事後的人.
法老家中的人是他們傳給他們的.
都聽見約瑟哭了起來.
約瑟對他的弟兄說.
我是約瑟.
我的爸爸還在嗎.
哇.
他的弟兄不能回答.
為什麼呢.
因為每個人的臉都發青了.
非常驚慌.
你是約瑟.
你知道埃及人的裝束不同.
埃及人多數是沒有頭髮的.
剃得很光.
只剩下一條邊.
穿了很多衣服.
有金鏈.
這個宰相.
法老之下.
萬人之上.
不認得他.
有時候會喜歡塗顏色在眼睛上.
有時候戴一個長嘴在這裡.
他的哥哥不認得他.
約瑟對他的弟兄說.
請你們走近一點來看看我.

$^{321}$大人走近一點?.
是走近一點.
他們就近前來說.
我就是你們的兄弟約瑟.
就是你們賣到埃及的.
哇.
我想那幫哥哥都嚇到暈倒.
變眼敏就不知道什麼事.
還是童子小時候.
好了.
來到這裡.
第二點要講.
就是螢光的紅色字.
第二.
現在不要因為把我賣到這裡.
自憂自恨.
我特別要講一下這個自憂自恨.
約瑟說.
各位哥哥.
十位哥哥.
弟弟.
不要因為當日你賣我.
而自憂自恨.
這裡要學約瑟的地方就在這裡.
不要因為賣我來這裡.
而自憂自恨.
如果約瑟講得出這句說話.
一定是自己先勝過了自己的自憂自恨.
才講得出的.
最會自憂自恨的是誰呢?.
就是約瑟.
被哥哥賣了.
賣了又開頭本來穿著彩衣去找哥哥.
哎呀.
很開心的樣子.
約瑟都不會露面的了.
都去找.
好啊.
出去走走.
然後.

$^{361}$賣了給這個米電人.
賣去埃及.
喂.
賣了做奴隸.
他已經要走路了.
那會怎樣的?.
衣服脫了.
鞋子脫了.
要綁著.
跟著駱駝這樣跑.
這樣走到埃及.
然後.
全身剝光珠.
因為.
好像買一隻雞一樣.
買一隻牛一樣.
是要看全相好不好.
靚不靚.
大不大隻.
白不白.
這樣好不好.
哎呀.
全部在看.
有一個人來到.
大家都讓開.
就是皇上的護衛長波提弗來到.
就是.
這個都挺好的.
行了這個.
看他都挺帥的.
就這個.
於是租了回去.
回去做下人.
不過做得很賣力.
他知道.
弱勢已經離開了以前那些公子哥兒.
他知道如果我再不是這樣.
在這裡沒命了.
於是他很賣力做事.
嘩.

$^{401}$這麼厲害.
這麼賣力做事.
才被提拔到做了管家.
做了管家有沒有壓力?.
最沒壓力就是做最低那個.
做管家最有壓力.
哎呀.
有些人服你.
有些人不服你.
有些人又這樣.
這裡爛的要整.
哎呀.
煩得不得了.
不過弱勢.
很帥.
那裡說他很俊美的.
很帥.
外形輪廓.
Appearance.
那個外形和他的臉的輪廓都很漂亮.
就被主母看中.
記不記得.
在那些猶太.
現在我們讀神學的時候.
就沒什麼讀那些經外的經典.
譬如那些《刺經》.
或者當日其他人寫的著作.
覺得那些不對.
那些《刺經》.
其實是提供了很多.
即是當代同時間有人在那裡做著作.
或者他們流傳下來.
後來有人寫下.
那些拉比的文獻是說什麼呢.
弱勢是很帥.
帥到什麼呢.
是帥到弱勢的主母叫其他的太太.
闊太,官太來一起.
我介紹個人給你認識一下.
很帥的.

$^{441}$是誰.
我管家弱勢.
叫他出來.
他們那些文獻是說.
弱勢一出來.
每個女人都戒手.
即是說可能在削蘋果.
可能看到弱勢.
嘩!這麼帥.
每個都喜歡.
於是主母就更加喜歡他.
就引誘他.
他一時失神.
就一個人進他房.
那一次就出事了.
就被他告他了.
他其實有點失措.
慌忙失措.
主母捉他.
拔掉他的衣服.
他是命令式的.
你與我同寢罷.
跟我一起睡.
這樣的意思.
是命令.
他就馬上脫了衣服跑出來.
如果弱勢可以做的就是.
拔掉衣服.
大聲喊引起引起引起.
可能有些人走進來.
什麼事啊?.
老闆.
沒事沒事.
就跟我一起出去.
可能就沒事了.
他沒有被拔掉衣服.
他沒有叫人.
於是脫了衣服.
當時可能以為很厲害.
誰知被主母起了指令.

$^{481}$惱羞成怒害他.
其實老闆波提弗.
一聽到之後生氣.
下載監獄.
於是知道老闆是相信他的.
如果不是已經拔掉了.
嘩!戲弄我太太.
斬了.
波提弗是很大官.
斬了他.
陰犯全部放在他家.
很厲害.
他放他在監獄.
其實他也相信他.
不過被捉到.
沒辦法.
然後又被酒精忘記了他.
又坐了三年.
接著出來解了夢.
然後做到現在.
嘩!如果弱勢離別了家鄉這麼久.
經歷這麼多事.
最會自憂自恨的是誰呢?.
弱勢比他哥哥先.
由於他在自憂自恨裡.
已經是痛苦不堪.
經過很多年之後.
他終於仰望神.
見到神的工作.
他出來釋放到.
所以見到他哥哥這麼害怕.
他就說不用怕.
不需自憂自恨.
其實是他自己先經歷了.
在這裡跟大家講第二點.
不要自憂自恨.
無論你被人害.
或者你經歷到很低谷.
很慘各樣事.
不過你要記得.

$^{521}$是上帝的手在背後托著你.
支持你.
扶著你.
是他讓那些事發生.
沒有一件事在上帝是偶然.
哎呀!沒想到會這樣.
不好意思了.
害你了.
不會的.
每一件事上上帝都知道.
都預先安排好.
你一定要有這個信心.
我們做基督徒要有這個信心.
所以不要自憂自恨.
仰望主.
頭低就自憂自恨.
越看自己就越慘.
頭臥高.
不會自憂自恨.
所以他才能夠用這句說話.
不要自憂自恨.
去安慰到他哥哥.
就是說他自己先釋放了.
要學習.
第三點.
他解釋到原因是什麼.
怎樣釋放他不再自憂自恨呢?.
就是他知道是.
上帝差我在你們以先來.
為要保存性命.
原來是這樣的.
就是神差我.
這個差字.
你看到我們今天.
在香港基督徒多數都看的差字.
就是差遣禮.
即是做宣教.
差遣.
其實是上帝的旨意.
將我安排差我來這裡.

$^{561}$大家會明白了.
上帝差我.
原來這件事是上帝容讓它發生的.
會不會有些事情是.
一下上帝失手我就死?.
不會.
如果你會有些懷疑的.
請你看多一兩次《約伯記》頭那幾章.
最精彩.
我們很喜歡看的那幾章是什麼呢?.
天上開庭.
上帝的眾子都來到事立.
看看上帝有什麼吩咐.
這個時候.
撒旦魔鬼仍然可以來天上.
可能是上帝給他任務.
讓他可以在世界上.
找一些人.
你可以給他一些試驗.
讓他可以多學一些人生的功課.
上帝就問撒旦.
你去哪裡了?.
我就聽你講在地上走來走去往返而來.
這樣就對了.
你走來走去有沒有見到約伯?.
當然有了.
你以為我不知道你想什麼嗎?.
約伯當然好.
你給他那麼多好東西.
他當然尊敬你.
你試一下.
把他的錢都拿了.
他就不相信你了.
好.
他的身外物全部在你手中.
但是你不可以碰他這個人.
你有沒有看到這裡是上帝給了界限給魔鬼.
你可以攻擊他.
但是不可以攻擊他的人.
然後第二次又開庭的時候.

$^{601}$天上天朝.
喂.
撒旦有沒有見到約伯.
你叫我給你害他.
沒有家產了.
兒女都死了.
他沒有來氣我.
魔鬼就說.
當然了.
錢的東西身外物.
你試一下動他這個人.
他一定當面棄掉你.
上帝就說.
好你可以動他這個人.
但是不可以拿他命.
不可以死.
於是約伯由頭頂到腳掌生瘡.
你知道了.
最慘的連太太都說.
這樣的神你還相信.
你離棄他.
死了他吧.
約伯就說.
你這個愚頑的婦人.
你什麼都不懂.
我們在神道得福都可以受害.
赤身出於母胎赤身歸回.
多厲害.
你看到.
上帝是有旨意.
他有管理.
和他有界線.
若生不可以死.
他可以經歷很多痛苦的事.
搞到自憂自恨.
令他這個本來穿彩衣的公子哥兒.
這個小伙子現在變得很成熟.
於是上帝就說.
我要猜他為所有以色列人保存性命.
現在這裡已經饑荒了兩年.

$^{641}$還有五年不能耕種.
不能收成.
而上帝是猜我在你們以先來.
為要給你們存留餘種在手上.
要大師拯救保存你們的生命.
即是看來猜我來到這裡的不是你們.
你們沒有猜我.
你們賣我.
猜我來的是上帝.
又令到我好像法老的父親.
什麼都問我.
什麼都聽我.
而且我做他全家的主.
並且埃及地傳遞的宰相.
這個是上帝的工作.
不是我.
也不是你們.
不要說你們聰明.
令到我來到這裡.
不是.
是上帝猜我來.
剛才和大家看的第三點.
上帝猜我.
在李沫儀先來到.
為要保存我們族人的性命.
看完這個藥室的故事來到這裡.
心裡十分之感恩和讚美上主的大臨.
他是由小時候都有多少驕傲.
又大言不慚的童子藥室.
藥採醫的.
由他這樣訓練到他承擔起.
埃及的宰相.
而且這個人是很謙和的.
一點都不厲害.
多好啊.
真是一個很合上帝用的人.
很多人看到創世記藥室故事.
就想著亞伯拉罕以撒雅各藥室.
然後再傳下去大衛.
不是.

$^{681}$其實是猶大.
亞伯拉罕以撒雅各猶大.
大衛主耶穌.
藥室最重要的是保存全家的性命.
於是以色列可以延續下去.
重要性在這裡.
是上帝猜測他.
講到這裡.
你會記得.
神的手從來沒有在我們旁邊縮過.
哎呀 糟糕.
或者哎呀 糟糕.
做漏了.
意外.
沒有的.
是祂的手一路扶持.
一路安排.
一路保護.
一路猜險.
最重要懂得的是要看得通.
是上帝給我的.
所以不要自憂自恨.
先在自己的不如意順境與逆境當中.
等等.
不要自憂自恨.
逆境當中出來.
仰望上帝.
看到祂的恩典.
自己先脫離自憂自恨.
於是你才會行得正.
以及愛到其他人.
這是很重要的.
求主幫助我們.
我們一齊禱告.
天賦上帝真的很深的.
很用力的賺給你.
主要是你自己從來沒有誤事.
你從來沒有不對的時間.
從來沒有失手.
從來沒有意外.

$^{721}$每一件事都在你的手裡.
在我們看的時間.
好像過山車一樣.
又高又低.
很可怕的日子.
藥室.
做個奴隸.
綁住.
從來沒有試過的.
衣服又沒有.
穿鞋又沒有.
綁住.
跟著駱駝或者奴仔.
走到埃及.
成為奴隸.
脫光了.
被人檢驗.
然後買了.
做奴隸.
沒有人供的.
然後升高又被主慕陷害.
坐了監.
然後在監裡.
又被他幫助過的酒精忘記了.
然後他又出來.
一直做到現在.
見到那些哥哥.
不發火就假了.
不過他就是這樣.
他看到.
不再自憂自恨.
是因為他們.
藥室舉目仰望主.
知道這是主的差遣.
這個就是我們人生要學的.
每一件事是主差遣.
讓事情發生.
給我們脫離了自憂自恨.
給我們更有力量去服侍你.
垂聽禱告.

$^{761}$奉耶穌基督命求.
阿們.
\newpage



\section{撒母耳記上 8:4-10}
\label{sec:031868}
\textbf{以色列興衰史:合一王朝系列 (1)}~~~~袁惠鈞牧師博士
\newline
\newline
連結: \href{http://www.church.com.hk/acms/content.asp?site=cdc\&op=show\&type=product\&code=031868\&layout=sermon}{\texttt{031868}} ~~~~ 語音日期: 2026-03-15
\newline
\newline
\hyperref[sec:031867]{\small{< < < PREV SERMON < < <}}
~
\hyperref[sec:index_chronic]{\small{[返順時目]}}
~
\hyperref[sec:031869]{\small{> > > NEXT SERMON > > >}}
\newline
\newline
撒母耳記上 8:4-10
\newline
\begin{longtable}{cl}
\hline
\hline
章節 & 經文 (和合本修訂版)\\
\hline
8:4 & \begin{tabularx}{0.7\textwidth}{X} 以色列的長老都聚集在拉瑪，來到撒母耳那裡， \end{tabularx} \\ \\ \relax
8:5 & \begin{tabularx}{0.7\textwidth}{X} 對他說：「看哪，你年紀老了，你的兒子又不行你的道。現在請你為我們立一個王治理我們，像列國一樣。」 \end{tabularx} \\ \\ \relax
8:6 & \begin{tabularx}{0.7\textwidth}{X} 撒母耳不喜悅他們說「立一個王治理我們」，他就向耶和華禱告。 \end{tabularx} \\ \\ \relax
8:7 & \begin{tabularx}{0.7\textwidth}{X} 耶和華對撒母耳說：「你只管聽從百姓向你說的一切話，因為他們不是厭棄你，而是厭棄我，不要我作他們的王。 \end{tabularx} \\ \\ \relax
8:8 & \begin{tabularx}{0.7\textwidth}{X} 自從我領他們出埃及的日子到如今，他們離棄我，事奉別神；正像他們從前所做的一切事，現在他們也照樣向你做了。 \end{tabularx} \\ \\ \relax
8:9 & \begin{tabularx}{0.7\textwidth}{X} 現在你只管聽從他們的話，不過要嚴厲警告他們，告訴他們將來王會用甚麼方式管轄他們。」 \end{tabularx} \\ \\ \relax
8:10 & \begin{tabularx}{0.7\textwidth}{X} 撒母耳將耶和華一切的話轉告求他立王的百姓。 \end{tabularx} \\ \\
[1ex]
\hline
\hline
\end{longtable}
$^{1}$列王記的第一講.
我們剛剛才講完《雜務2記》.
所以今天要開始一個新的系列.
我非常興奮.
因為我很喜歡舊約.
有人問:牧師,為什麼你這麼喜歡舊約?.
現在我們都有新約了.
為什麼還要讀舊約?.
舊約是新約的儒表.
新約是舊約的承傳.
所以新舊約都一樣重要.
千萬不要說有新約而不去讀舊約.
我們今天要開始講《列王記》.
第一講我通常都是介紹.
每一個系列的開始.
那篇道都是為何要讀這本書.
所以今天也不例外.
為何要讀《列王記》.
我在每個系列第一講.
都會講一個介紹的一篇道.
也會講有很多背景,有關這本書的資料.
讓大家來認識,更加了解這本書.
我們剛剛講完《雜務2記》.
《雜務2記》可以說是將《士司記》和《列王記》連在一起.
《雜務2記》就在中間.
在《士司記》的時候.
以色列人沒有王.
但去到《雜務2記》他們希望薩姆爾幫他們立一位王.
而去到《列王記》他們真的有很多的王來帶領他們.
很可惜這些王大部分都不是好的皇帝.
《列王記》是屬於猶太人的歷史.
是近東一帶的歷史.
在《列王記》裡面沒有亞洲的歷史.
沒有美洲的歷史.
這些歷史對我們很關心.
因為我們是亞洲人.
又住在美國.
但是聖經裡面沒有講到這些國家.
其實對神來講最重要的.
就是祂自己指紋的歷史.

$^{41}$就是猶太人的歷史.
在這裡又要講一講.
我們這些讀神學的人.
讀神學的時候都會讀系統神學.
在系統神學裡面.
我們都會讀神論,基督論,聖靈論,聖經論這些科目.
但是有一個科目是完全沒有提到的.
是什麼?.
就是有關以色列人的科目.
其實在聖經裡面.
超過60\%的篇幅都是有關以色列人.
居然在系統神學就沒有有關以色列人的科目.
其實這個是很不應該的.
為什麼會這樣呢?.
就是因為有很多人都憎恨以色列人.
說是他們將耶穌基督釘死在十字架的上方.
所以他們都不想讀有關以色列人的系統科目.
甚至我也講過.
很多人都說沒有千禧年國.
什麼是千禧年國?.
就是在算末世的時候.
耶穌基督會回來.
祂會坐在大衛的王位上.
重建大衛倒塌的帳篷.
這個就是千禧年國.
因為很多人都憎恨以色列人.
所以他們說不會有這件事發生.
我們不妨來看看.
列王記在聖經的位置.
我們想像的聖經可以分為四個部分.
第一個部分就是《摩西吾經》.
創出黎民新.
第二個部分就是歷史的部分.
包括了我們所講過的.
撒慕爾記上下和列王記上下.
還有第三個部分就是先知書.
有《以賽亞書》,《耶利米書》.
《耶利米哀歌》,《以西結書》,.
《但以理書》和《十二先知書》.
我已經曾經講過.

$^{81}$《耶利米哀歌》,《但以理書》和《十二先知書》.
如果大家有興趣可以聽這些系列.
最後第四個部分就是寫作的部分.
《約伯記》,《詩篇》,《箴言》,.
《傳道書》和《雅歌》或者《所羅門志歌》.
我也曾經講過《箴言》和《雅歌》的系列.
大家喜歡的都可以聽.
在猶太人的聖經.
和我們手上的聖經有點不同.
它們分為三個部分.
第一個部分就是律法.
就是和我們的摩西五經一樣.
還有第二個部分.
它們是叫做先知書.
但是它們的先知書和我們的先知書很不同.
等一下我會解釋.
還有第三個部分就是寫作.
我為什麼說它們的先知書和我們的不同.
我們認為是歷史書的.
《約書亞記》,《時詩記》,.
《撒母耳上下》和《列王記上下》.
它們編入了先知書裡面.
叫它們做前先知書.
為什麼它們會這樣?.
明明這些對我們來說是歷史.
但是猶太人會說這些不是歷史.
這些是預言.
神早就預言這些事會發生.
例如在《約書亞記》說到以色列人會進入迦南地.
他們就會說早在《創世記》的十二章和十五章.
神已經預言這些事會發生.
所以不是歷史,是預言.
另外我們今天讀的《列王記》說到以色列人.
因為他們背若最後被神趕出應許之地滅國.
猶太人會說這些事早在《利未記》的二十六章.
和《申命記》的二十八章都已經有預言.
所以他們說這些不是歷史,是預言.
所以他們將這些書本編入了前先知書裡面.
猶太人歷史的開始.
可以說是由亞伯拉罕.

$^{121}$就是希伯來人的老祖宗開始.
亞伯拉罕的事是公元前二千年的事.
也就是說離開我們已經有四千年那麼長.
由亞伯拉罕去到大衛一共是一千年.
去到大衛的時候可以說是以色列人傳承的時期.
可以說是他們最高峰的時期.
但是在亞伯拉罕去到大衛之間.
以色列人是經歷了有兩個靈命倒退,靈命低落的時候.
一個就是他們在埃及的時候.
他們在埃及大約四百年.
但是他們是沒有行國禮.
所以我們見到他們的確靈命很低落.
還有第二個時期就是士詩時期.
在士詩的時候他們是武王來管他們.
各人偏行己路.
當他們犯罪的時候神就興起他們的敵人來攻擊他們.
使得他們成為奴僕.
然後他們就向神來呼求.
神就差派士詩來拯救他們.
但是很可惜拯救了他們之後.
他們又去犯罪,他們忘記了.
所以在士詩記裡面一共有七個惡性的循環.
所以我們見到以色列人在埃及.
和在士詩的時代.
是他們兩個靈命最低落的時期.
我差點想加多一個時期.
就是以色列人出了埃及之後.
在曠野的時候.
他們也是靈命低落.
但我隨即想到.
雖然第一代的以色列人.
除了迦納和約書亞.
他們是可以進入應許之地.
其他的第一代的以色列人.
都因為他們的不信而死在曠野.
但雖然是這樣.
第二代的以色列人他們是服從的.
他們是能夠進入應許之地.
所以我就沒有加上這個時期.
我們說大衛的時候.

$^{161}$可以說是以色列人巔峰的時候.
是鼎盛的時候.
從大衛開始以色列人就一路走下坡.
去到所羅門之後.
以色列國分為南北兩個國家.
北國有20個王.
沒有一個是好的.
南國19個王,1個王后.
但大部分的王也是不好的.
是給他們有兩個復興的時候.
一個是在希西加的時候.
另外一個在約書亞的時候.
北國在公元前的722年.
是被亞述國所滅.
他們是沒有復國.
南國在公元前的586年.
被巴比倫使他們亡國.
他們被擄.
就成為街下囚.
就好像當他們在埃及的時候一樣.
但是神給南國是會有復興.
人民是可以回到他們自己的領土.
應許之地.
雖然他們回到應許之地.
但他們仍然被不同的國家來征服.
剛才已經說過有亞述和巴比倫兩個國家.
接著有馬代波斯,希臘和羅馬.
正如路加醫生在《路加福音》的21章24節所說.
「耶路撒冷將要被外邦人踐踏.
直到外邦人的日期滿了」.
這也正如保羅在羅馬書第九章至十一章告訴我們.
「現在我們這個時候就是外邦人的日期」.
「神是將猶太人暫時放在一邊」.
「現在是外邦人得救的時期」.
「是教會時期」.
「到外邦人得救的人數天滿了」.
「全以色列都要得救」.
在耶穌基督第一次降世.
來到世上為我們釘十字架之前.
全部在世上都是人的國.

$^{201}$我們剛才說到亞述國,巴比倫,馬代波斯,希臘和羅馬.
這些都是人的國.
當耶穌基督來臨為我們釘了十字架之後.
在我們所有信徒的心裡.
我們都有神的國.
而我們也是屬於我們自己的國家.
例如我們現在在美國,我們就是美國的公民.
所以在這一段時間.
在教會的時期.
是有人的國和神的國並存的時候.
這個並存的時候.
會一直到耶穌基督再次來臨.
當耶穌基督再次來臨的時候.
祂就會帶領我們和以色列人進入千禧年國.
這樣就會開始了神的國.
以後就會帶領我們去到新天新地.
再沒有人的國.
有很多學者.
他們會引用約翰福音的十八章三十六節.
耶穌基督說:我的國不是屬於這個世界.
這些人就會說:你看,沒有千禧年國.
因為耶穌基督的國不是屬於這個世界.
所以在地上沒有一個實質的千禧年國.
但是,解錯經了.
耶穌基督是說什麼?.
祂的國不是屬於這個世界.
千禧年國不是人的國.
是神的國.
是由耶穌基督親自來掌權.
所以我們要小心.
千萬不要讓這些人來誣陷我們.
我扯得太遠了.
我要回到列王記.
在《撒母耳記》的時候.
我們說以色列人曾經要求撒母耳納一位王給他們.
其實早在史詩的時代.
以色列人在機殿的時候.
已經希望擁納機殿成為他們的王.
但是機殿很有智慧.
他拒絕了.

$^{241}$他說你們已經有一位王.
這個就是神治基.
在列王記.
終於以色列人他們的願望實現了.
是有很多王來帶領他們.
但是大部份的以色列王都是不好的.
而在列王記裡面.
所有的王都是以一個標準來衡量他們.
是什麼標準?.
就是大位的標準.
如果是有不好的王.
經文就會說.
他不像他的先祖大位專心服從或者是依賴耶和華.
如果是有好的王.
經文就會說.
他行他先祖大位一切所行的道不偏左右.
所以所有的王都是以大位來衡量他們.
其實在列王記裡面.
衡量一位王都是從屬靈的角度.
都是從屬靈的眼光來看他們.
一個很好的例子.
就是德國有一個邪惡的王帝叫做暗尼.
他是一個出色的統治者.
他有強大的王國.
在政治上來說他是一個天才.
在經濟上來說.
他是帶領以色列人過著富庶的生活.
但是在列王記裡面.
他只有七節的經文.
無論他在地上有什麼成就.
神體他只是從屬靈的角度是屬靈的眼光.
另外一個極端.
就是在南國有一個好王帝約西亞.
我剛才說過南國是有兩個復興的時候.
一個就是希西加.
一個就是約西亞.
在約西亞的時候.
南國是有很大的復興.
約西亞在聖殿裡面找到律法書.
所以就幫助他有一個很大的改革.

$^{281}$以至有一個復興的時候.
在列王記裡面.
是有三章的經文.
來描述約西亞的事蹟.
所以我們見到.
這個就和暗尼成為一個強烈的對比.
列王記這本書的書名又是怎樣來的呢?.
就好像很多其他的書本一樣.
在聖經裡面的書本.
都是用希伯來文第一個字成為書名.
例如創世記第一個字就是Bereshit.
就是開始神創造天地.
這個Bereshit就是開始的意思.
所以中文的版本就叫做創世記.
列王記都不例外.
在列王記上的一章一節講到「王大衛」.
「王」這個字是希伯來文的Melech這個字.
就成為了這本書的書名.
而中文的版本我們就叫它做列王記.
其實在希伯來人的聖經裡面.
列王記上下是一本書.
是七十字譯本.
就是第一部將希伯來文的舊約翻譯成希臘文的版本.
首先將列王記分為上下兩冊.
而拉丁文的版本和英文的版本.
都跟著用同樣的手法.
而今時今日.
希伯來文的聖經也都是將列王記分為上下冊.
還有一樣很有趣的事.
就是無論是七十字譯本或者是拉丁文的版本.
都將薩姆爾記和列王記連在一起.
在七十字譯本裡面.
我們叫做薩姆爾記上下的.
被命名為王國的第一卷和王國的第二卷.
而我們的列王記上下.
就被命名為王國的第三卷和王國的第四卷.
在拉丁文的版本也都很相似.
我們叫做薩姆爾記上下的.
被命名為王的第一卷和王的第二卷.
有沒有看到.

$^{321}$拉丁文聖經的注意力是放在王的身上.
而七十字譯本是放在王國的身上.
是有些少不同.
我們認為是列王記上下的.
拉丁文的版本.
叫做王的第三卷和王的第四卷.
所以無論是七十字譯本或者是拉丁文的版本.
所記載的列王記.
都是由素羅一直講到.
由大王西底迦就是南國最後一個王的故事.
也就是說記載了整個以色列王權的歷史.
你說歷代至上下呢?.
歷代至上下就是補充了猶大君王的歷史.
列王記的作者又是誰呢?.
根據猶太人的傳統.
作者就是耶利米先知.
主要的原因就是耶利米這個名字.
在列王記裡面你是找不到的.
在列王記裡面是講到有很多其他的先知.
例如是撒迦利亞先知.
以賽亞先知.
尼加先知.
甚至是西班牙先知.
而耶利米一個這麼出名的先知.
居然在書裡面是沒有他的名字.
甚至在列王記下的二十五章提到耶利米的事跡.
也沒有講到他的名字.
這就是原因有很多學者認為這本書是耶利米所作.
其實在聖經裡面有另外一個例子.
就是約翰福音.
我們知道約翰福音的作者是使徒約翰.
但是在約翰福音裡面.
你找不到約翰的名字.
約翰叫自己做主所愛的門徒.
所以有很多人認為列王記就是耶利米所作的.
原因就是和約翰福音一樣.
但是也有很多人認為列王記的作者不是耶利米.
原因就是在列王記最後發生的一件事.
是發生在巴比倫.
是在公元前的五百六十一年.

$^{361}$照我們所知.
耶利米從來沒有去過巴比倫.
他是去了埃及.
根據耶利米書的四十三章一至七節.
在公元前的五百六十一年.
耶利米已經是大概八十六歲.
他已經是一個老人.
不是很可能能夠再去到巴比倫這個地方.
其實在列王記裡面.
這個作者從來沒有透露他自己的身份.
但是作者就強調了先知的事公.
所以列王記的作者真的可能是一個先知.
而他是和以色列人一起流亡去到巴比倫.
在列王記裡面.
亦都是引用了很多在聖經以外的資料.
例如所羅門記,以色列諸王記和猶大諸王記.
這些書本都沒有被列入聖經的正典.
亦都是沒有神所密使.
但是它們是寶貴的歷史資料.
所講的歷史事實是很可靠.
此外在列王記裡.
亦都是引用了其他聖經裡面的書本.
例如在列王記下18章至20章所講的事跡.
可能就是引用了以賽亞書的36章至39章.
在列王記下的25章的事跡.
可能就是引用了耶利米書的52章.
這個作者的確是一個當事人.
是一個目擊的證人.
所以他所寫出來的事跡絕對是可靠.
此外當然他亦都有聖靈的密使.
所以他說的話絕對是可信靠.
列王記的成書日期又是何時呢?.
我們說作者是一個目擊的證人.
在列王記裡面直到今日這句話.
一共出現了有14次這麼多.
這句話非常之重要.
因為這個作者所描述的這些歷史事物.
在他寫作的時候仍然存在.
一個好的例子就是聖殿.
當這個作者開始寫列王記的時候.

$^{401}$聖殿仍然存在.
當然在成書的時候.
以色列人可能已經是被擄了.
所以聖殿已經被毀.
但是他開始寫作的時候聖殿仍然存在.
列王記我們剛才說最後提到一件事.
在公元前的561年發生.
是甚麼事呢?.
就是南國的王約瓦根在監獄被巴比倫王釋放.
這件事應該是以色列人被擄之後的37年所發生.
也就是說列王記最後成書的日期.
應該在以色列人被巴比倫國被擄的期間.
而最可能成書的日期.
就是公元前的561年至538年之間.
這本書成書的目的又是甚麼呢?.
有兩個原因.
第一個原因就是要記載歷史.
第二個原因就是要教訓以色列人.
我們先看第一個原因.
列王記所記載的事跡.
是由大衛死了之後所羅門登基開始.
這就是公元前的971年.
一直記載到巴比倫人將耶路撒冷淪陷.
將聖殿毀滅.
這就是公元前的586年.
在這些期間.
以色列人是由團結的十二支派.
是合一個王朝.
變成了分為南北兩個國家.
北國是以色列.
南國是猶大.
當以色列國還是一個聯合國的時候.
他們曾經有三個王.
就是蘇羅,大衛和所羅門.
他們每個人大概都作了王四十年.
在聖經裡四十是一個試驗期.
以色列人在曠野四十年到處漂流.
以利亞呢?.
以利亞先知戰勝了巴力的先知.
被耶洗別來追殺他.

$^{441}$他是跑了四十晝夜去到神的山.
去到荷列山.
耶穌基督在曠野被撒旦試探有四十天.
又是四十天.
所以我們見到在聖經裡.
四十真是一個試驗期.
所以無論是蘇羅,大衛和所羅門.
他們都被神試驗了四十年.
他們的開始都是好的.
但很可惜最後他們都失敗了.
就算大衛都不例外.
大衛都是一個罪人.
但是是神的恩典.
是神的憐憫.
稱大衛做最合他心意的人.
而我們剛才都說過.
所有猶大王也好.
以色列王也好.
他們衡量的標準就是大衛.
在所羅門死了之後.
以色列被分為南北兩個國家.
北國有十個支派.
南國有兩個支派.
猶大和變雅文.
北國地大物博人數眾多.
南國比較瘦弱人數稀少.
北國建都於撒瑪利亞.
南國建都於耶路撒冷.
北國第一個王是耶羅伯安.
他因為怕自己的人民以色列人.
會回到南國的聖殿耶路撒冷來敬拜.
所以他設立金牛的偶像.
在雪旦和伯特利兩個地方.
他居然帶領以色列人來拜偶像.
北國有很多壞皇帝.
他們都謀朝篡位.
所以北國經歷了有九個王朝.
能不能想像.
二十個王裡面有九個王朝.
你都別說不可怕.

$^{481}$南國始終是一個王朝.
就是大衛的王朝.
北國二十個王都是不好的皇帝.
南國一共也有二十個王.
一個王后十九個王.
王后就是阿他利雅.
在南國有八個好皇帝.
而有十二個壞皇帝.
真是諷刺.
在南北國一共有四十個王.
居然只有八個王是好的.
也就是說只有五分之一的王是好的.
不過比較少有先知的輔助.
南國有很多有先知的輔助.
因為神和大衛來到立約.
大衛是屬於猶大支派.
而猶大支派是南國兩個支派之一.
所以神特別姑息南國.
差派了很多先知來勸告他們.
北國就比較少.
因為神從來沒有和北國來立約.
在獵王企裡.
對於王的介紹都有很特別的方式.
無論是南北國.
都會說到王登基的日期是何時.
而南國的王.
還會說登基的年齡有多大.
無論是南北國.
都會說到他們統治的期間有多長.
而南國的王還會說統治的地點.
無論是南北國的王.
都會說他們的父親是誰.
而在南國的王的介紹的時候.
還會介紹王的母親.
無論是南北國的王的介紹.
都會說到他們屬靈的評價.
說到他們的死因和他們的承繼人是誰.
我剛才說了獵王企寫作的目的之一.
就是記載歷史.
還有一個目的是要教訓以色列人.

$^{521}$獵王企的作者給我們看見.
就是因為以色列人.
他們對神的話語越來越漠不關心.
他們最後離棄神.
拜偶像以至朝來亡國朝來被擄.
獵王企不單止是準確的歷史.
還是對歷史有準確的解釋.
作者很明顯是一個.
和以色列人一起流亡到巴比倫的流亡者.
他很希望透過以色列人的歷史.
來教訓被擄的人.
告訴他們為甚麼他們會亡國.
為甚麼他們會被擄.
這個作者早就在獵王企告訴讀者.
他們要得到神的祝福.
一定要遵守神的律例典章.
如果他們不遵守的話.
他們就會被懲罰,被審判.
甚至被神趕出應許之地.
很不幸.
我們說大部份南國的王.
和所有的以色列的王.
都是不虔誠的.
都不是好皇帝.
他們都是行了耶和華眼中看為惡的事.
這些君王都是邪惡.
他們不單止沒有禁止以色列人派偶像.
還鼓勵他們引誘百姓來犯罪.
我們剛才都說過.
耶羅波安就是一個好的例子.
神於是差派先知去面對這些君王和人民.
但是都被拒絕.
所以神就透過先知.
來預言以色列和猶大國.
將會滅國,將會被擄.
而最後這些預言都實現了.
其實獵王企是說出了.
為甚麼人民會被擄.
是對這些被擄的讀者來解釋.
原因就是因為他們拜偶像.

$^{561}$所以有件事以色列人真的改變了.
當他們被擄回歸之後.
他們就沒有記載他們有拜偶像.
其實在這獵王企裡.
也讓我們看到神的憐憫.
在這獵王企裡.
有兩個很邪惡的王.
北國的阿哈王和南國的約瓦根.
兩個都是邪惡的王.
但是神就向他們施憐憫.
這也是給被擄的以色列人有很大的鼓勵.
告訴他們如果他們願意回轉的話.
神會寬恕他們.
就好像神憐憫阿哈王和約瓦根王.
這在獵王企裡是一個非常重要的訊息.
獵王企的地理環境又如何呢?.
當然是有關整個以色列地.
就是從丹到別士巴.
由北邊到南邊.
也包括了約旦河東的地方.
在公元前的971年到561年.
當時有四個國家不斷侵入以色列國.
第一個就是以色列人的宿敵是埃及.
在公元前的十世紀.
埃及在所羅門和羅波安的統治期間.
意圖影響以色列人的歷史.
第二個國家就是亞蘭國或者是敘利亞.
在公元前的九世紀.
他們對以色列造成了很大的威脅.
而第三個國家就是滅了北國的亞述國.
在公元前的八百年至七百五十年.
亞述是忙於對抗他們的鄰居敘利亞.
所以他們沒有時間來南侵.
就是因為這樣.
給了以色列和猶大有半個世紀的和平和繁榮.
但是之後他們的王提格拉比列色三世.
開始進攻北國以色列.
在公元前的七百二十二年將北國亡了.
而在公元前的七百零一年.
他們也圍困耶路撒冷.

$^{601}$就是西那基納這個王.
當時是希西加在南國掌權.
幫希西加能夠退敵.
打勝了他們的敵人.
而第四個國家就是滅了南國的巴比倫.
在公元前的六百一十二年至五百三十九年.
巴比倫成為古代的強國.
曾經三次入侵南國猶大.
然後在公元前的五百八十六年.
將耶路撒冷淪陷.
將聖殿二為平地.
列王記的神學觀又如何呢?.
是記載了有三件事都是與神有關.
第一就是神的公義.
在列王記告訴我們.
就算神的子民以色列和猶大.
就算他們不服從神的毀約.
神一樣會審判他們.
而那些邪惡的王絕對是要為了自己所做的事來負責.
神公義的憤怒淋到自己子民的身上.
神是公義的,這是第一點.
第二點是讓我們見到神的大能.
神就是那位全能的神,全知的神.
祂一早已經是差派先知.
被以色列和猶大預言.
他們會滅國,他們會被擄.
最後這些事都應驗了.
而南國呢?.
神亦都預言他們被擄之後.
七十年將會回歸.
而這些事全部都應驗了.
讓我們見到神就是那位全知的神.
是大能的神.
還有第三點就是神的信實.
神和大衛立約.
雖然大衛的後人都毀了約.
但是神沒有將他們毀滅.
就好像神將耶羅伯安一世.
暗利和耶護的家族.
都是將他們毀滅了.

$^{641}$神是沒有毀滅大衛的家族.
當列王記這部書完結的時候.
大衛的王朝已經不再存在.
但是神仍然保存大衛的後裔.
而耶穌基督就是大衛的後裔.
耶穌基督來到世上.
祂是坐在大衛的王位上.
祂將會帶領以色列人再次復興.
帶領以色列人重建大衛的王朝.
在列王記裡面.
亦都有解釋的挑戰.
主要解釋的挑戰有兩個部分.
第一個部分就是年份的矛盾.
第二個部分就是聖約的承傳.
我們先看年份的矛盾.
學者都注意到在列王記裡面.
無論是有關以色列王或者猶大王的年份.
都是有些不一致.
在列王記裡面其實記載了很多歷史的事實.
從在聖經以外的資料.
例如希臘,亞述,巴比倫的考古學和文獻.
我們可以知道有很多這些歷史的事實發生在何時.
而在公元前的892年至566年間.
有很多事情我們都可以非常確實.
知道這些事情是發生在何時.
在亞述國的歷史亦記錄了有關以色列王.
阿哈和耶護的事.
所以我們很肯定知道阿哈去世的時候.
應該是公元前的853年.
而耶護登基的時候應該是公元前的841年.
有了這些固定的日期.
我們就可以知道.
以色列南北國分裂的時候.
應該是公元前的931年.
而北國撒瑪利亞淪陷的日子.
應該是公元前的722年.
而南國耶路撒冷淪陷的時候.
應該是公元前的586年.
我們剛才說在列王記裡面.
那些年份有很多是不一致的.

$^{681}$這只不過是表面的不一致.
其實有兩個部分.
一個部分就是個別王朝的年份有不一致.
另外一個就是總共的年份都有不一致.
我們先看個別年份的不一致.
我們可以用一個例子來舉一反三.
指出問題在哪裡.
例如在列王記上的十六章二十三節.
說以色列王暗利是在屬於猶大王阿薩三十一年開始作王.
而在屬於猶大王阿薩三十八年他去世.
他的兒子阿哈繼承了他成為以色列王.
這樣就有數得計了.
三十八年去世,三十一年登基.
做了王七年.
但是聖經裡面居然說他在位十二年.
這就是不一致的地方.
總共年份的不一致呢?.
如果我們將所有北國王在位的年份加在一起.
我們會得到二百四十一年.
但是我們知道南北國分裂的時候.
是公元前的九百三十一年.
而北國的滅亡是在公元前的七百二十二年.
減一減我們就知道北國的國運應該是二百十一年.
而不是二百四十一年.
那為什麼將所有王的年份加在一起就會多了三十年呢?.
等一下我們會解釋.
同樣猶大國都是一樣.
你將所有南國王的年份加在一起.
你是得到三百九十三年.
但是剛才我們說南北國分裂是公元前的九百三十一年.
而南國的滅亡是公元前的五百八十六年.
你減一減的話就會得到三百四十六年.
而不是三百九十三年.
為什麼將所有王的年份加在一起又會多了四十七年呢?.
有一個很容易的解釋.
就是根據以色列人的傳統.
一個王可以和他的父親父子一起掌權一段很長的日子.
如果是這樣的話.
在聖經裡面的年份就可能有重複有重疊.
我們不妨看回剛才個別王朝的例子.

$^{721}$列王記上的十六章二十三節說.
亞瑟奄利作王十二年.
但是我們剛才也說過.
他是在亞瑟的三十一年登記.
在亞瑟三十八年死亡.
他應該是做了王七年.
為什麼說十二年呢?.
其實他自己做王是七年.
但是他和他的父親森利一起掌權有五年.
所以一共他在位十二年並無矛盾.
那總共的年份為什麼會不一致呢?.
如果父子一起掌權的話.
父的年份你又加進去.
子的年份你又加進去.
那你就會多了.
所以本來北國二百十一年的國運變成了二百四十一年.
本來南國三百四十六年的國運就變成了三百九十三年.
就是這樣的原因.
除了年份計算的不一致.
也都有聖約承傳的不一致.
其實在聖經裡我們知道神和亞伯拉罕立約.
神也和大衛來到立約.
神和亞伯拉罕立約有一樣很重要的東西.
就是地的應許.
神和大衛立約.
很重要的就是大衛有後裔坐在他的王位上直到永遠.
我們先看神和亞伯拉罕立約.
我們先看地的應許.
這個地有多廣闊呢?.
在民數記的三十四章六節.
西邊要到大海為界.
這就是你們的世界.
大海是什麼海?是地中海.
但是在列王記上的四章二十一節.
所羅門統管諸國從大河到非利士地.
直到埃及的邊界.
所羅門在世的日子.
這些國都進貢服侍他.
在所羅門的版圖.
西邊是去到非利士地.

$^{761}$不是去到大海,不是去到地中海.
有很多學者會用列王記上的五章一節來解釋.
經文說:推羅王西南平素愛大衛.
他聽見以色列人高所羅門.
接觸他父親作王.
就差遣臣僕來見他.
他們會說:你有沒有見到.
這個推羅王希蘭臣服於所羅門.
所以他的地屬於以色列.
所以版圖和神應許亞伯拉罕一樣.
但不是的.
這個希蘭王是一個獨立的國王.
他和所羅門是平起平坐.
不能說他的地屬於所羅門屬於以色列.
還有就是不變地的應許.
根據神給亞伯拉罕的應許.
不免是要去到敘利亞.
但我們知道就算大衛鼎盛時期.
都沒有征服敘利亞,只是打敗他們.
所羅門就更加了.
敘利亞的地不是屬於所羅門.
那你說:神的應許是不是落空了?.
沒有落空.
因為剛才我們也說過.
耶穌基督會回來帶領我們進入千禧年國.
在那個時候這些地的應許會被實驗.
還有一個應許就是大衛的應許.
神應許大衛將會有一個後裔坐在他的王位上.
根據列王記上的五章五節和八章二十節.
這個後裔應該是所羅門.
但我們要看神和大衛立約的經文.
在撒母耳夏的第七章.
而在列王記上的六章十二節都有提到.
這個後裔是會謹守神的典章.
遵從神一切的誡命.
哇!所羅門有沒有遵守神一切的律例典章?.
當然沒有,所羅門是一個罪人.
在列王記上的十一章.
告訴我們他娶了很多外邦妃嬪.
這些妃嬪都引誘他離棄神.

$^{801}$所以所羅門是一個罪人.
哇!那你說神給大衛的應許又落空了?.
不是,剛才我們已經解釋.
神給亞伯拉罕的應許應驗在哪個人身上?.
是耶穌基督的身上.
而神給大衛的應許.
也是一樣實驗在耶穌基督的身上.
那你說牧師是你說的.
不是我說的.
在舊約裡面.
以賽亞先知,耶利米先知,以西結先知.
十二先知都有說到.
這些應許都會實驗在彌賽亞的身上.
就是耶穌基督.
他第一次來臨的時候.
讓我們見到他真的沒有罪.
他是聖潔的.
所以就應驗了神和大衛所納的約.
他第二次來臨將會帶領以色列人進入千禧年國.
那個時候.
以色列人的領土.
就會完全實驗了神給亞伯拉罕的應許.
這個也是原因.
為什麼我一開始的時候就說.
有很多人都想告訴我們.
沒有千禧年國這件事.
他們為什麼會有這樣的思想?.
現在有很多神學官都是這樣教的.
好像改革中.
就會告訴我們這些地上的應許都不重要.
最重要的是什麼?.
是將來的應許.
所以神給亞伯拉罕地上的應許不重要.
最重要的是將來應許亞伯拉罕天上的耶路撒冷.
所以他們會將所有地上應許的經文都寓意化.
說這些天國就是屬靈的國.
而不是實質的國.
如果是這樣的話.
我們跟佛教的佛經和回教的可蘭經就沒什麼分別.
佛教說我們現在的身體是臭皮囊.

$^{841}$將來我們會去極樂的世界.
回教他們會說.
你現在服從亞伯拉罕.
將來就有七十二個處女在天堂等著你.
如果每件事都寓意化的話.
我們跟佛教和回教根本沒有分別.
我們跟他們不同的是什麼?.
我們的聖經是實實在在的歷史.
可以用考古學來印證.
所以不要將所有這些東西都寓意化.
神給我們世上的應許也會實現.
證明神是大能的.
如果神給我們的應許.
在地上是不會實現的.
是將來去到天上.
你怎麼知道神應許你將來在天上的事一定會發生?.
我們這位神就是大能的神.
所以祂在地上應許的一切都會發生.
這個我們要很小心.
千萬不要被人矇騙.
我們最後會說列王記上下的大綱.
列王記上下可以分為五個部分.
第一個部分就是王國的傳承.
在列王記上的第一章至第二章.
是說到大魏的衰退,所羅門的興起.
而第二個部分就是王國的興旺.
在列王記上的第三章至第十章.
主要是說所羅門的王國.
第一個部分和第二個部分都是有關合一的王朝.
而第三個部分是王國的分裂.
列王記上的十一章至十四章.
而第四個部分是說到王國的先知.
主要是有兩位先知.
第一位就是以利亞先知.
在列王記上的十七章至二十一章.
第二位就是以利沙先知.
在列王記下的一章至八章.
第五部分就是王國的避諱.
所以列王記可以分為五個部分.
而我這次說列王記這個系列.

$^{881}$會把它分為五個系列.
第一個就是合一的王朝.
主要是說所羅門.
第二是說到北國的王朝.
第三個系列是南國的王朝.
有沒有看到.
我是將南北國的王朝分開了.
在列王記中南北國是混在一起.
很多時候是比較混亂.
你不知道是說南國還是北國.
你不小心的話.
在這個系列我會將兩個王朝分開了.
第四個系列就是以利亞先知的系列.
第五個系列就是以利沙先知的系列.
願神祝福他自己的話語.
願意我們都能夠透過列王記來學習.
千萬不要像以色列人一樣.
願意我們都能夠成為一個夢神喜悅的兒女.
讓我們一起低頭禱告.
讓我們看到你就是那位公義的神.
審判是由你的家開始.
就算你自己的子民猶大以色列.
當他們背若.
當他們不守你的律例典章.
你一樣會審判他們.
懲罰他們.
你也是那位全知大能的神.
你早已預言以色列人會遭遇到這些事.
你也早已預言將會將以色列人放在一邊.
現在就是教會的時期.
願意我們都能夠.
不要學習以色列.
我們要讓你來猜險出去.
做你大使命的事工.
以致我們能夠帶領更多的人信著.
而你也是那位信實的神.
你所應許的都會實現.
無論是地上的無論是將來的事.
一切都會應驗.
因為我們見到耶穌基督第一次的來臨是那麼準確.

$^{921}$有關耶穌基督第二次來臨的預言.
一樣是那麼準確.
願意我們他人能夠繼續.
走在你的旨意.
走在你的路上.
以致將來你回來的時候.
能夠被你稱為一個忠心良善的僕人.
將一切的榮耀頌讚歸給你.
我們禱告.
奉主耶穌基督的名義而後.
阿們.
\newpage



\section{約翰福音 3:22-30}
\label{sec:031869}
\textbf{事奉的榜樣}~~~~駱柱成牧師
\newline
\newline
連結: \href{http://www.church.com.hk/acms/content.asp?site=cdc\&op=show\&type=product\&code=031869\&layout=sermon}{\texttt{031869}} ~~~~ 語音日期: 2026-03-22
\newline
\newline
\hyperref[sec:031868]{\small{< < < PREV SERMON < < <}}
~
\hyperref[sec:index_chronic]{\small{[返順時目]}}
~
\hyperref[sec:031870]{\small{> > > NEXT SERMON > > >}}
\newline
\newline
約翰福音 3:22-30
\newline
\begin{longtable}{cl}
\hline
\hline
章節 & 經文 (和合本修訂版)\\
\hline
3:22 & \begin{tabularx}{0.7\textwidth}{X} 這些事以後，耶穌和門徒到了猶太地區，在那裡他和他們同住，並且施洗。 \end{tabularx} \\ \\ \relax
3:23 & \begin{tabularx}{0.7\textwidth}{X} 約翰也在靠近撒冷的哀嫩施洗，因為那裡水多，眾人都去受洗。 \end{tabularx} \\ \\ \relax
3:24 & \begin{tabularx}{0.7\textwidth}{X} 那時約翰還沒有下在監裡。 \end{tabularx} \\ \\ \relax
3:25 & \begin{tabularx}{0.7\textwidth}{X} 約翰的門徒和一個猶太人辯論潔淨的禮儀， \end{tabularx} \\ \\ \relax
3:26 & \begin{tabularx}{0.7\textwidth}{X} 就來見約翰，對他說：「拉比，從前同你在約旦河的東邊，你所見證的那位，你看，他在施洗，眾人都到他那裡去了。」 \end{tabularx} \\ \\ \relax
3:27 & \begin{tabularx}{0.7\textwidth}{X} 約翰回答說：「若不是從天上賜的，人就不能得到甚麼。 \end{tabularx} \\ \\ \relax
3:28 & \begin{tabularx}{0.7\textwidth}{X} 你們自己可以為我作見證，我曾說，我不是基督，只是奉差遣在他前面開路的。 \end{tabularx} \\ \\ \relax
3:29 & \begin{tabularx}{0.7\textwidth}{X} 娶新娘的是新郎；新郎的朋友站在一旁聽，一聽見新郎的聲音就歡喜快樂。因此，我這喜樂得以滿足了。 \end{tabularx} \\ \\ \relax
3:30 & \begin{tabularx}{0.7\textwidth}{X} 他必興旺；我必衰微。」 \end{tabularx} \\ \\
[1ex]
\hline
\hline
\end{longtable}
$^{1}$弟兄姊妹早晨.
我們又進入了新的一周.
今天我們詩上的經文就是.
約翰福音三章二十二到三十節.
我今天講的.
正道的題目是事奉的榜樣.
上次我們講到是.
前一段是很出名的經文就是.
約翰福音三章十六節.
神愛世人.
聖經清楚告訴我們.
世人犯了罪的時候.
虧缺了神榮耀的時候.
神就愛這個世界.
這個世界很不美好.
也很不可愛.
但是神在這種情況就已經是.
愛我們這些世人.
這個世界是墮落的.
敗壞的.
也是充滿罪惡的.
但是當我們做罪人的時候.
基督就為我們死了.
神的愛就在這裡顯明出來了.
不是等我們變成一個很可愛的人.
神才愛我們.
是當我們做罪人的時候.
神已經是愛我們了.
耶穌基督是按著父神的旨意.
在所定的日期.
道成了肉身.
祂在天上.
很榮耀的天上.
來到這個充滿罪惡的世界.
從至高之處來到這個很卑微的世界.
祂不是來到地上.
去到皇宮那裡.
祂是去到馬槽那裡出生.
我們人都是盼望向高處走的.
但是基督呢.

$^{41}$祂是按著神的旨意.
專行神的旨意.
祂不是升高.
祂反而是降卑.
所以基督來到世界的時候.
祂不是得著很大的榮耀.
好像我們這些世界的榮耀.
祂反而是被掛在木頭那裡.
被羞辱.
被責罵.
被釘在木頭那裡.
可以說是半天吊的.
就在這個痛苦當中.
祂竟然說出一個很偉大的話.
祂說.
父啊.
赦免他們.
因為他們所作的.
他們不曉得.
所以當我們做罪人的時候.
敵對基督的時候.
當時那些猶太人敵對祂的時候.
基督向神發出了一個偉大的禱告.
要愛你的仇敵.
為那些逼迫你們的禱告.
所以祂說要赦免他們.
他們還不曉得.
他們不曉得罪的後果是什麼.
所以耶穌基督的教導.
不單止是說出來的.
祂真正是在最艱難的環境當中.
把他們活出來.
所以基督.
當我們做罪人的時候.
當我們很邪惡的時候.
祂就願意為我們寫命.
這種愛在歷史上很難找到.
也是我們的理性.
我們的頭腦是不明白的.
因為我們.

$^{81}$通常我們自己愛的.
一般都是愛我們所愛的人.
所以你會看到.
這個世界所說的愛是從人出發的.
我們的起點就是從人出發.
是有條件 有界限和有目的的.
儒家一般說的就是.
要按著那個關係去愛.
什麼意思呢.
就是說你要先愛你的父母.
接著愛你的親戚.
再愛你的朋友.
是有個次序 有個親疏的.
你會看到.
墨子是提倡那種謙愛.
要人人都彼此相愛.
彼此帶來的那種利益.
你要愛你自己的父母.
好像愛別人的父母一樣.
你要愛自己的兄弟.
你也要愛其他人的兄弟.
就當他們是兄弟一樣愛.
當然了.
一般來說.
你是不可以愛別人的妻子.
當作是自己的妻子.
這是不可能的.
現代的那種愛呢.
你是愛自己的國家.
愛自己的人民.
去提倡這種思想.
那種愛還是從人出發的.
所以你會看到.
當說到社會的愛的時候.
他說愛社會 愛人民.
接著人民就在後面加多一個字.
叫做愛人民幣.
就是愛人民的錢.
愛 為什麼會這樣說.
因為愛是從人出發的.

$^{121}$他是以自我為中心的時候.
最終回到自己那裡.
就是要利己.
我要得到什麼.
這就是從下而上的愛.
神的那種愛不是.
他是從上而下的那種愛.
是很偉大的.
他不求回報.
他也不是為了得到什麼.
他是無條件的.
當我們這麼敗壞的時候.
基督這麼榮耀.
我們這麼卑微的時候.
他竟然在致命的天上.
來到這麼邪惡的地上.
拯救我們這些受苦的人.
為我們寫明.
我們敵對他.
他竟然能夠愛我們這些做受敵的.
基督的愛是不一樣的.
所以你會看到.
基督的愛在十字架那裡.
擔當我們的罪的時候.
我們看到.
那裡可以說是沒有愛的.
為什麼.
因為基督擔當人的罪.
承受神的審判.
他就被神所利器了.
好像愛那裡就是沒有愛.
沒有憐憫.
因為基督要承受公義的審判.
但是在這個救恩也讓我們看到.
正因為基督願意擔當我們的罪.
他成為了一個愛的源頭.
為什麼.
因為救恩成全了.
他的愛被顯明了.
他願意擔當我們的罪.

$^{161}$反而成為了一個愛的源泉.
成為一個愛的源頭.
所以當我們說今天要講事奉的時候.
你不懂得愛人,愛仇敵的時候.
你很難去事奉.
你也很難明白什麼叫做事奉.
因為事奉.
你是要放下你自己.
你是要專行神的旨意.
你是要信服神.
不是求你自己的榮耀.
今天我們很多的事奉.
很多時候我們說事奉.
我們是計算會.
我們付出多少.
我們希望要有多少的回報.
我們付出了.
我們希望得到別人的尊重.
這種的事奉不是事奉.
因為你有一個東西就是計較.
你還在計較.
你很難明白什麼叫做事奉.
我記得英文部.
英文部有個弟兄.
每逢星期四.
他都在這裡做那些雜工的事情.
我以為他是在這裡受薪而做的.
後來才知道他每個星期都來.
他不是.
他就是因為出於甘心樂意的事奉.
他每個星期按時來.
搞好衛生.
清潔周圍的環境.
他就是每個星期四都堅持.
所以你會看到.
他是不求回報的.
他不是為了什麼利益.
他就明白什麼叫做事奉.
所以當一個人.
計較的時候.

$^{201}$你很難明白.
很多時候我們做多一點的時候.
我們就會說我吃虧了.
但是你應該思考.
神會紀念你.
神紀念你.
很多時候.
人都想利用我們.
我們服侍嘛.
我們服侍的時候.
也不希望別人看低我們.
但是當別人看低我們的時候.
你說你是服侍的.
你會怎麼回應呢.
你應該說我要赦免他.
好像耶穌基督一樣.
赦免他.
因為他不懂得.
交回給神.
耶穌基督.
他就是事奉的一個真正的榜樣.
因為他專行神的旨意.
為人去付出.
也不求回報.
當我們讀今天的經文的時候.
你會看到後面還有一段.
所以你會看到.
上面是說神的神愛世人這一段.
這一段插入了一段歷史.
就是施洗約翰的門徒.
和耶穌的門徒.
有糾紛這件事.
耶穌的門徒在猶太地.
猶太地在那裡居住.
而且跟別人去受洗.
那麼.
而約翰呢.
他是在哀戀受洗.
耶穌是在猶太地受洗.
那麼為什麼要去哀戀受洗呢.

$^{241}$是不是因為一定要去哀戀那裡.
泡在水裡.
才是一個受洗的形式是必須的呢.
但是耶穌基督呢.
他反而是在猶太地受洗.
猶太地.
他不一定是在河裡.
所以我們會看到.
受洗不一定是泡在水裡.
但是.
施洗約翰.
他在哀戀那裡.
因為那裡水多.
哀戀這個地方是什麼意思呢.
就是水多.
很多水的意思.
所以你會看到.
受洗不一定要泡在水裡.
我們也會看到.
經文說的是耶穌基督的門徒.
為人去受洗.
其實耶穌基督不是親自去為人受洗.
為什麼.
因為約翰自己曾經說過.
我是用水去為你們施洗.
那麼在我後面來的位置.
就是耶穌基督了.
他是要用聖靈為你們施洗.
所以你就會看到.
施洗不一定是耶穌基督去施洗.
很可能不是耶穌基督用水去受洗.
因為耶穌基督來.
是要用聖靈去受洗.
是約翰親自說的.
但我在我後面來的.
他是要用聖靈為你們施洗.
所以你會看到.
耶穌基督很可能不是親自用水為人施洗.
是他的門徒為人施洗.
耶穌基督真正的工作是要用聖靈為人施洗.

$^{281}$人是要受聖靈的洗.
那麼約翰的水洗.
他只不過是一個預表.
預表什麼呢.
就是當水從上而流下來的時候.
就預表著聖靈好像從天降下來一樣.
所以這個水洗是一個象徵.
聖靈的洗只是真正的一個實現.
如果不是聖靈去重生我們.
我們沒有辦法得著這個永生.
我們前面的經文已經分享過.
所以你會看到水洗.
約翰的水洗是一個記號.
一個象徵.
表達什麼呢.
表達人要悔改.
表達人要尋求赦罪.
不可以靠著自己去成聖.
他是藉著這個水洗.
表達人要靠聖靈.
靠著聖靈的洗.
你才可以真正去成聖.
所以這個水洗就是表示.
人願意悔改.
人願意進入天國.
這個舉動就表達了.
人的心願是什麼.
他的心智是什麼.
那麼聖靈的洗呢.
聖靈的洗只是真正賜生命的.
是真正可以改變人的生命的.
這個才是真正的實體.
聖靈的洗才是真正的實體.
救恩真正的實現.
是要靠著聖靈才可以的.
所以約翰他只是一個受差的.
他不可能差派聖靈.
而是聖靈是靠著上帝去差遣的.
因為施洗約翰沒有這個資格.
他沒有辦法去做這件事.

$^{321}$他做的只可以是水洗.
所以水洗就是預表.
聖靈的洗才是真正的實現.
人這些外面的水洗.
施洗約翰的水洗.
不可以改變人的生命.
只有耶穌基督.
藉著聖靈.
去將我們潔淨.
我們的生命才可以得著改變.
所以你會看到.
約翰這個洗禮.
是一個象徵 一個預表.
基督才是實體.
舊約的禮儀 獻祭.
也是一個影.
基督才是真正的實體.
你會看到 舊約一直在預備.
新約就盛傳了.
基督來到的時候.
就盛傳了舊約一直所預備的.
所以基督來到.
預表就完全這樣.
完全完成了祂的使命.
舊約我們會看到有很多獻祭的.
獻牛羊.
這些祭物.
只是一個預表.
當基督來到的時候.
祂就成為真正的祭物.
耶穌基督就是代罪的羔羊.
所以你會看到.
我們今天的信徒.
不再獻祭了 為什麼.
因為那是死祭.
現在要獻的是活祭.
因為保羅這樣說.
所以弟兄們.
我以神的慈悲勸你們.
將身體獻上.

$^{361}$當作活祭.
是聖潔的 是神所喜悅的.
你們如此事奉.
你們如此事奉是理所當然的.
所以我們今天要獻的是.
獻上我們自己.
我們要當一個活祭來獻給神.
我們要事奉神.
所以舊約是一個影.
基督就是那個實體.
那個影不可以代替那個真實的.
預備也不可以代替那個承傳的.
同樣原則.
施洗約翰來到.
來到這裡 在這個地上.
他是做預備的工作.
但是基督來到.
他是做承傳救恩的工作.
所以角色是不同的.
施洗約翰他是預備道路.
呼召人悔改.
為基督去做見證.
所以施洗約翰他很清楚自己的角色是什麼.
他不是基督.
他不是耶穌基督.
他也不是那個核心.
那個中心人物.
他也不是那個主角.
他只是那個預備的.
為耶穌基督做預備的.
他很清楚自己的身份.
但是.
當有這個預備和這個承傳.
同時出現的時候.
問題就出現了.
那個問題是什麼.
人開始比較.
約翰還在.
耶穌基督.
來到了.

$^{401}$也在.
大家都在這裡做工.
人就開始比較了.
問題是什麼.
問題在人那裡.
不是在神那裡.
後來你會看到一些情況就是.
耶穌基督出來傳道.
耶穌基督很可能他沒有用水施洗.
因為他用聖靈去施洗.
那麼.
他的門徒開始.
和人施洗了.
那些門徒和約翰.
也都可能在這裡施洗了.
其實兩個人在這裡做工.
其實都沒什麼.
但是.
當人有比較的時候.
他就會有那種競爭.
他就會爭競可異.
就好像我們今天一樣.
如果我們這裡月語堂.
附近一帶都是很少月語堂的.
但是突然之間.
我們搬來這裡沒多久.
旁邊又多一個月語堂出來.
哎呀.
慘了.
這一間.
我們一間就夠了.
為什麼你還多一個在這裡.
所以.
如果你是很熱心傳福音的.
你是我們教會的一些很熱心的人員.
特別是做外展同工的.
那麼你可能就說.
我做了這麼多一下子就被搶走了.
一下子那些人就去旁邊了.
可能你就很生氣了.

$^{441}$所以很多人他學不到.
施洗約翰.
他是做了一個很好的榜樣.
他可以容納很多事情的發生.
他可以容納很多的事情.
容納很多的人.
當耶穌基督的門徒施洗的時候.
約翰繼續去施洗.
這兩樣東西.
他同時去做.
他沒有說他們做了.
我就停止了.
他是同時存在.
其實神的事工是可以同時存在的.
為什麼呢.
因為神的工作是有不同的階段的.
耶穌基督來到了.
施洗約翰還在這裡.
這個階段他們同時存在.
也都在神的工作當中.
我們也都有不同的職份.
不同的搭配.
所以是可以同時出現的.
所以約翰.
不要因為耶穌基督的出現.
而停止他的工作.
只要神給他機會繼續去做.
他就繼續忠心去做.
所以一個原則就是說.
神的工作不是一個人去完成的.
神的國度是要很多的僕人.
一同去參與去完成的.
而且不應該是真正的.
但問題就來了.
當人不明白這個原則的時候.
神的工作是這麼豐富的.
是可以同時用納的.
他們就去比較 去質度.
我們前段時間學了《雅各書》.
人的口是很難去被制服的.

$^{481}$是不止息的惡文.
是滿了害人的毒氣.
人的心是出懷著苦毒的質度.
和紛爭的.
所以約翰說你們不可以自誇.
也不可以說謊話抵擋真理.
因為這些不是從上頭而來的智慧.
是屬地的 是屬情慾的.
所以你會看到教會很容易.
你不難去發現.
有這樣的人.
因為我們的口舌很難去制服.
很多時候我們都有一些無心之失.
為什麼呢.
比如你要稱讚一個人的時候.
你唱詩歌真的很好聽.
比誰還要好聽.
這樣就出問題了.
前面的句子是很好的.
你稱讚他.
真的很詛咒他.
但多了後面的一句.
比誰還要好聽.
你的主學真的說得很好.
但是.
誰又好一點.
這樣的時候那個比較就出來了.
那個紛爭就出來了.
所以我們很容易驕傲.
很容易增進.
如果稍不警醒的話.
我們就落在這種試探當中.
就在這個背景之下.
約翰的門徒來對約翰說.
眾人都往耶穌那裡去了.
耶穌在那裡.
我們本來是一大群人的.
我們本來這裡這麼多水.
應該是多一些人過來這裡洗的.
現在人們都跑到耶穌那裡去了.

$^{521}$他的門徒在那裡私洗.
所以他們說的這句話.
可能是沒錯的.
真是一個事實.
但是.
這個不是單純的事實.
是加了人的眼光進去.
人的眼光覺得不公平.
為什麼說不公平.
從人的角度來說.
當耶穌出來的時候.
是私洗約翰介紹他的.
是把人介紹給耶穌.
這個就是那一位.
我要等候的那一位.
這個就是.
當我看到聖靈降在他身上的時候.
我就知道他是神的兒子.
這個就是那一位.
而且我是把很多的門徒推給他的.
約翰.
這麼偉大的人物.
不是私洗約翰.
是另外一個約翰.
我們這個約翰福音的約翰.
都推給他了.
那麼你耶穌呢.
耶穌你今天.
我已經成為一個不在猶太地那裡.
一個這麼中心的地方.
我已經去到哀怨那裡.
可能是成為一個邊緣人物.
我已經推了下去.
但是你現在還搶走那些人.
這樣就很不公平.
所以他的門徒很可能說.
你知不知道.
你現在有一個對頭.
你有一個對手.
他在那裡輸了.

$^{561}$那些人都跑過去了.
就說 約翰.
你的風頭已經被耶穌基督搶走了.
耶穌基督現在威過你了.
可以說他大粒過你了.
以前你這麼提拔他.
現在他威過你了.
所以這幫人就在挑逗當中.
在比較當中.
在嫉妒當中.
如果約翰的生命不夠成熟的話.
他很可能說.
耶穌基督你真是忘恩負義.
我已經把最好的學生給你了.
介紹給你認識了.
那些人都願意去你那裡了.
你今天 你有今天.
不是傳教我嗎.
你竟然一點面子都不留給我.
如果他靈性不好的話.
他很可能是這樣說.
甚至可能他說.
你耶穌基督現在有今天.
是傳教我的.
如果不是我.
你根本上沒有今天.
這些話其實在我們當中.
每一個人都是很熟悉的.
因為很可能.
我們威風的時候都說過.
你有今天.
是不是當天我介紹一份好工作給你.
不是靠我.
所以這件事就很危險了.
因為神藉著人去幫助人的時候.
人只是參與神的工作.
恩典的源頭依然是神的.
所以當我們這樣說的時候.
你就很容易取代了神.
你就將榮耀歸給你自己了.

$^{601}$其實沒有一個人可以說.
我可以去造就他.
我可以去成就他.
其實真正來說.
那個恩典的源頭一定是從神而來的.
人可以幫人一把.
神可以藉著你去幫他一把.
但是最終的功勞.
最終的恩典是從神而來的.
所以真正的才能.
潛力和才幹.
這些都是從神而來的.
因為保羅也說過.
我栽中了,阿波羅囂貫了.
唯有神才能讓他生長.
所以一個人能夠達到某個地位的時候.
某個能力的時候.
是神祂自己的作為.
所以你會看到.
耶穌基督不是不近人情.
祂在神的歷史當中要顯明一件事.
就是當人要將基督和其他人去比較的時候.
當人和基督捍衛同等的時候.
祂要信服神.
神要在這個地上有祂的作為.
因為耶穌基督是獨一的.
祂只是獨一的真神.
沒有人是可以和祂比較的.
如果有人誤會這件事的時候.
就很危險了.
如果在歷史當中約翰這麼偉大.
耶穌基督也這麼偉大.
如果耶穌基督和約翰是同等的時候.
那麼耶穌基督的獨特地位就失去了.
就很容易誤導人了.
所以神祂有祂的作為.
所以當人將人取代神的時候.
這是很危險的.
因為人本來就是人.
人要做人的事.

$^{641}$神是神 祂是獨一的.
當人說要取代神的時候.
這就很危險了.
很可能這個人.
將這個人提高到神的地位的時候.
這個人在褻瀆神.
因為祂將人取代了神.
所以神按著祂的主權.
當基督慢慢升起的時候.
興起的時候.
約翰祂慢慢就下來了.
因為約翰祂的使命就是做預備的.
當祂完成祂的使命的時候.
祂就完成了神託付祂的.
所以耶穌基督興起.
很多人歸向祂.
約翰祂下來.
這是神的安排.
不是不公平.
而是要顯明.
誰才是真正的忠心.
歷史的忠心是誰.
誰才是真正的獨一的神.
其實我們看經文的時候.
一切都是從神而來的.
施洗約翰自己曾經說過.
我是用水施洗.
以後那個是要用聖靈施洗.
後來這班門徒就跟他說.
你所見證的那一位.
現在在這裡施洗.
而且人都去了那裡.
約翰在這種情況下.
沒有去辯護.
他沒有說什麼.
他只說了一句很偉大的話.
若不是從天上賜的.
人就不能得著什麼.
什麼意思呢.
你們不要為我擔心.

$^{681}$你怕我的地位被人搶了.
你怕為我抱不平.
你們不需要.
因為若不是從天上來的.
人得不著什麼.
人多也好.
人少也好.
這個完全我要交回給神.
這個就是一個真正事奉者的態度.
不是我要去爭.
而是我要交回給神.
將我的工作交回給神.
因為一切都是從神而來的.
所以這個.
當約翰這樣說的時候.
他就站得很穩定了.
因為他不需要爭.
不需要搶.
不需要比較.
只要忠心地完成他的使命.
神給他多的人.
他高興.
神取走那些人.
他已經信服了神.
他也做了他該做的.
他也滿足了.
所以神要他停在哪裡.
他就停在哪裡.
今天我們為什麼會比較.
因為我們忘記了一件事.
就是一切都是從神而來的.
當你忘記這件事的時候.
一切都是從神而來的.
你就很以為很多事情都是我的.
我自己的.
如果你有恩賜.
其實真正來說.
是神賜給你的.
是為著你去服侍神.
神賜給你的.

$^{721}$如果你是一個普通人.
你也不需要覺得自己很差.
因為也是神給你的.
你已經盡了你的能力了.
神同樣有工作給你.
同樣有事奉給你.
沒什麼是可以自誇的.
因為保祿這樣說.
你有什麼不是領受的呢.
既然是領受的.
一切都是從神領受的.
有什麼是值得自誇的呢.
所以真正明白這件事的時候.
我們才不會去驕傲.
不會去嫉妒.
別人做得好.
自己做得稍微.
按現象來說稍微少一點.
因為一切都是從神而來的.
一切都是為了榮耀神.
但是人呢.
通常來說就是要爭一口氣.
人爭一口氣.
佛爭一炷香.
他就要爭來爭去.
這是世界.
就要證明我自己是偉大一點的.
但約翰讓我們看到.
他是不求自己的榮耀的.
他是不急著被人尊重他.
他不需要.
他不在意自己的地位.
因為他知道.
若榮耀是從神來的.
神遲早會給你.
因為神要榮.
抬舉一個人的時候.
到了時候.
神稍微動一動手.
你自然得到榮耀.

$^{761}$不是靠自己的方法去追求榮耀.
如果你是靠自己的.
很可能有一天.
變成你一種羞辱.
因為你是靠自己爭回來的.
我不是說人不需要去盡責.
真正的事奉者.
不是為自己的面去活的.
而是為.
也不是人稱讚你去做些什麼.
也不是說要努力.
得到一些地位.
而是真正的去尋求.
神祂的旨意.
討神祂自己的喜悅.
完成神祂交付給你的.
所以人稱讚你也好.
彈你也好.
你都不要受這些影響.
為什麼.
為什麼這樣說.
因為人.
他的動機不是為了人.
他是為了神.
我們人最終要向.
向一位負責.
就是神自己.
不是向人去負責.
如果你是為了人的話.
人一變.
你就開始搖動了.
如果人一彈.
你就開始跌倒了.
因為你看著的是可變的.
你要定精在神的身上.
你才可以站穩.
所以當我們認清自己的身份的時候.
你才有真正的喜樂.
約翰很清楚地說了一件事.
我不過是人.

$^{801}$基督祂是神.
當基督興旺的時候.
我不是嫉妒.
我而是喜樂.
因為他說了一句很偉大的話.
「他必興旺,我必衰微」.
如果一個人不認識自己的時候.
他很容易越過自己的本份.
做那些不該做的事.
說那些不該說的話.
最後就羞辱他自己了.
蘇格拉底這個哲學家.
他說了一句很偉大的話.
人不應該先去尋求那些外在的.
你要先認識你自己.
如果你要知道其他的話.
你應該先要了解你自己.
他就說到這裡.
但是他沒有繼續說.
他沒有給我們一個答案.
怎樣才可以認識我們自己.
但是聖經就有說.
聖經說得很清楚.
敬畏耶和是自畏的開端.
當一個人如果要認識自己.
你先要認識神.
在神的光照當中.
在神的話語當中.
你認定你自己是誰.
你的位置是甚麼.
我需要做甚麼.
我需要神光照我們.
明白我們的位置在哪裡.
所以當一個人真正認識自己的時候.
他就不會越過那些過界的事.
做那些過界的事.
反而是信服神他自己的安排.
當這樣的時候.
你在自己的位置的時候.
真正做神所交託給你的事的時候.

$^{841}$你才真正有平安和喜樂.
為甚麼很多時候我們沒有平安和喜樂.
就是因為太計較得失.
太在意別人怎麼看我們.
但是約翰說得很清楚.
若不是從天上所賜的.
人得不著甚麼.
所以無論神比我們多也好.
少也好.
都是神的作為.
我們不要拿這些東西去比較.
因為一切都是從神而來的.
我記得我在廣州聚會的時候.
有一個傳道人.
有一個很老的傳道人.
坐了二十年牢.
出來就繼續去事奉.
他沒有因為.
都有點老了.
出來都已經五十多了.
沒有因為年齡而退下去.
而繼續去事奉.
當中有一個姐妹.
經常跟著他.
他也很有恩切.
他講到很有能力.
很有恩慈.
神也使用他.
但是這個年長的弟兄.
更有名望.
神更大地使用他.
可以說他影響了整個中國的人.
所以這個姐妹.
他本來可以多一點機會繼續去講.
但是他看到.
教會的人數還不是很多.
還沒有一個巨型的教會.
所以他就小站講台.
每逢星期三晚.
他就帶一個福音班.

$^{881}$傳福音.
他也講到.
但是少一點.
每個星期三晚.
就在那裡坐定定.
等弟兄姊妹帶人回來.
他自己帶人回來.
他就在那裡分享.
所以一個在講台.
一個在福音那裡.
在那條線那裡.
各按著神給他的託付.
他們沒有爭競.
也沒有比較.
而是彼此去搭配.
所以教會後來.
真的變成了一個巨型教會.
所以事奉不是證明自己.
有什麼能力.
而是真正承傳神的工作.
神的工作是什麼.
約翰繼續講了一句.
我曾說我不是基督.
你們忘記了.
為什麼.
因為你們將我和耶穌基督去比.
我只不過是僕人.
我只不過是被差遣的.
我是為主預備道路的.
你們應該為我做見證.
因為我事奉的心智.
我一直是這樣做的.
我沒有改變過我的立場.
我就是為主去預備的.
但你們今天竟然去講這番話.
我從來都沒有想過.
一個真正謙卑的人.
當一個人知道.
我只不過是神所夢召我.
做該做的時候.

$^{921}$我只不過是服侍神的人.
其實我沒有什麼值得他驕傲的.
所以當一個人認識這件事的時候.
我只不過是僕人的時候.
我只不過是神呼召我出來工作的時候.
我只不過是信奉聖經的教導的時候.
我們的禱告不應該說.
神啊你讓我給多些人去認識我.
而且我要做得比人好.
比人更加之成功.
而是應該這樣說.
神啊無論是生還是死.
我應該是要叫基督在我身上顯大.
這個才是真正的僕人.
不是要我被顯大.
而是要基督被顯大.
當我要比其他人強的時候.
這個不是事奉.
當我所做的是顯出基督的時候.
這個就很合乎聖經所教導的事奉.
因為我們只不過是僕人.
是被差遣的.
我們不是主角.
真正主角是耶穌基督.
所以我們會遇到不少的試探.
如果我們依然想高舉我們自己的時候.
當你有這樣一種思想的時候.
這個念頭出現的時候.
你就開始偏離事奉了.
這個就不是事奉了.
因為你的心態已經錯了.
真正的事奉不是要高舉我自己.
真正的事奉應該是讓人看到.
你所做的不是你有什麼功勞.
而是看到基督.
這個就是真正的事奉.
所以約翰用了一個很美的比喻.
約翰說娶新婦的就是新郎.
新郎來迎接新婦.
那個是屬於他的榮耀.

$^{961}$我只不過是新郎的朋友.
我不是新郎.
所以我的職責應該不是搶新郎的位置.
不是搶基督的位置.
我是要為他預備做很多事情.
我要做那些聯絡的工作.
在這個婚禮當中.
我要成就很多的事工.
目的就是將新婦帶到新郎面前.
這個就是新郎的朋友要做的.
所以當新郎.
特別是在洞房花燭夜的時候.
新郎就在那裡守護房間.
當新郎的朋友在外面聽到.
新郎得著新婦的時候.
他就為新郎高興了.
而不是嫉妒.
如果你嫉妒的時候.
你就出問題了.
為什麼.
因為你不是新郎.
你的位置就很危險了.
所以真正最寶貴的地方是什麼.
我們不是要站在最核心的地方.
我們是要讓主站在最核心的地方.
釋放最大的回報是什麼.
不是成功或者是別人稱讚.
而是服侍主本來就是一件喜樂的事.
他本身就是喜樂的.
不是要有什麼成果.
他本身就是喜樂的.
因為施洗約翰也說過.
我給他解鞋帶都不配.
所以你能夠服侍本來就是一個恩典.
哪怕那些工作是很卑微的.
也是榮耀的.
因為約翰的心態是主必興旺.
我看到耶穌基督興旺的時候.
我就喜樂了.
他來就是為了這個目的.

$^{1001}$他不是要自己興旺.
是要基督來高舉.
所以當我們一邊去服侍的時候.
將人帶到自己身邊的時候.
你就出錯了.
因為人看到的是神的工作.
興旺的時候.
而不喜樂.
這樣這個人就出問題了.
如果當一個人將人帶到自己的面前.
不是基督的面前.
這個人也就出問題了.
我記得我早期一出來事奉的時候.
帶一個團契花了很多時間.
每個星期可能一到兩次查經.
一直要帶.
一帶就帶了兩年.
後來教會說我要調動你去其他地方服侍.
這個團契就交出來了.
所以我當時很不開心.
覺得花了這麼多時間.
後來就想了.
其實弟兄姊妹得到更好的造就.
交給另外一個姊妹.
去讓他們成長.
是更加好的.
因為他們被建立了.
所以人如果心裡是以為.
這個小組是我的.
這個團契是我的.
這個教會是我的.
這樣你就出問題了.
因為這些都是主的.
是耶穌基督的.
我們不是建造自己的王國.
我們是服侍主的.
我們是要承傳主自己的工作.
所以真正的事奉不是只是做事.
不是被動的.
不是像機械一樣.

$^{1041}$而是要明白神的心意.
甘願去馴服.
全人的投入.
聖經形容約翰的時候.
約翰有以利亞的心智和能力.
事奉不是靠一些方法.
技巧或者學問.
最重要的是.
聖經說出來.
靠著心智和能力.
因為心智表明是神呼召你的.
你依然忠於這個心智.
能力是神給你的.
就算有一天你年老了.
環境也好不利了.
但你的心智不變的時候.
神給你的能力依然在的時候.
你依然可以為神去做很多的事.
你依然有影響力.
所以今天事奉最大的問題不是方法.
而是人看到我們不是基督.
如果你在你的分享當中.
人記得你的口才很好.
人記得的不是基督.
你就出問題了.
如果人家說你很厲害.
但不是說基督很偉大.
你就出問題了.
有一個十九世紀.
一個很著名的傳道人.
叫施布珍.
有一群美國的基督徒.
去倫敦做一個短訪.
有人建議他們去兩個地方聽道.
一個是去Joseph那裡.
於是他們早上就去Joseph那裡.
講道很好.
他們聽完就說.
我們敢肯定.
他就是有史以來最偉大的人.

$^{1081}$他講道太好了.
所以他們說.
不如晚上我們繼續聽一場施布珍那裡算了.
但其他朋友就說.
你要去施布珍那裡.
他們在這個勸說之下.
他們就去了施布珍那裡聽.
來到大都會會幕那裡.
去聽施布珍講道.
當聽完的時候.
我們敢肯定.
耶穌基督才是真正偉大的救主.
所以兩個人的服事都很好.
很厲害.
一個人家看到這個講完很有吸引力.
但另外一個人家看到.
基督很偉大很有吸引力.
人家的事奉應該是這樣的.
不是我,是基督.
所以不是要顯出我們有多榮耀.
而是要顯出基督有多榮耀.
不是我們有什麼恩賜.
完全是一切都是出於神的恩典.
我們很多時候以為我們做了很多事.
其實如果你再回想的時候.
你知道很多事都不是我們做的.
為什麼這樣說.
你去探訪就知道了.
可能你的交通的問題.
去做安慰給什麼話語.
你要想.
如果是主日學.
你要預備那個課程.
講道更是了.
有神的靈的帶領.
很多事都是神讓我們經歷.
經歷祂的恩典.
我們才可以完成.
不是我們有什麼功勞.
如果你覺得你做了很多事的時候.

$^{1121}$很高興的時候.
覺得我就是這麼厲害.
那你就出問題了.
為什麼.
因為你還沒死.
你剛才唱的詩歌有死.
你還沒死.
你依然在在乎人的回應.
過於神的榮耀.
你還是在乎我.
還是以我為中心.
還沒死.
所以面對那些稱讚的時候.
你應該將神的榮耀歸給神.
面對那些.
彈的時候.
很多人彈你的時候.
可能是一些不對的話.
不合理的話.
你應該為他禱告.
這些都不是學來的.
為什麼.
因為你要有基督的生命.
你要有基督一個豐盛的生命的時候.
你才可以行出來.
所以當一個人真正.
他看自己.
那個自我死的時候.
他不再為自己活的時候.
他就可以說出一句.
他必興旺.
我必衰微.
你會看在約翰身上.
他很偉大.
但他是一個被聖靈充滿的人.
聖靈充滿的人.
會不會犯錯.
今天可能我們已經做到長老.
已經做到執事.
已經是教會幾十年的人.

$^{1161}$生命已經很豐盛.
但依然會出問題.
約翰在難處的時候.
他坐牢的時候.
他也懷疑過耶穌基督.
他說以後要來的.
那一位是不是你.
還會發出一個問題.
那耶穌基督是怎麼做的呢.
耶穌基督的回答是.
乞子看見.
騎腿行走.
掌大麻風的潔淨.
耶穌基督說這些話.
是要造就約翰.
雖然他有軟弱.
他已經是很豐盛的一個生命.
但耶穌基督依然扶持他.
他沒有拆毀約翰.
他反而稱讚他的同工約翰.
你就是那個點著的明燈.
所以同工和同工之間.
應該是互托互相扶持.
不是應該彼此去攻擊.
所以同工是建立人的.
不是打擊.
因為耶穌基督他自己也是.
壓傷的蘆葦不折斷.
漲殘的燈火他不催滅.
他是扶持他是堅固.
讓人重新得力.
約翰他一生堅持.
他是基督,我是僕人.
他是興旺的,我是衰微的.
他不單止說出來.
他而且活了出來.
我必須衰微.
這個不單止是他的心意.
他是明白神的心意.
當一個人去到一個很頂峰的時候.

$^{1201}$好像約翰他.
他已經是很多的會眾見過他.
聽過他在約旦河的分享.
也有很多人跟從他.
他在這個時候就要退下來.
他依然盼望.
有人比他做得更好.
約翰他的心是這樣的.
他的願望是這樣的.
但是我們中國人的思想很不同.
有一句話叫做.
教息徒弟,餓死師父.
就是反映一個什麼心態.
我怕被取代.
什麼意思呢.
就是如果我傳了一些秘技給他的時候.
我的絕招傳了給他的時候.
他有一天就把我弄倒.
那我就沒有地位了.
所以你會看到很多的技巧.
很多的文化那種技藝.
因為太過保守了.
傳男不傳女.
傳內不傳外.
到最後死的一刻再傳.
就不見了.
所以那種保守和狹窄.
使很多的秘技.
很多的好東西流失了.
如果是武俠小說的話.
如果每一代我都教授一招的話.
降龍十八掌.
一代之後十七掌.
再兩代十六掌.
最後變成降蛇掌.
沒有威力了.
那些精髓都不見了.
世界就是這樣的一個思想.
他是為了保住自己.
反而失去了更多.

$^{1241}$中國華人的文化.
你會不少見.
但是如果這種思想在教會的時候.
就會很危險.
因為我們很多時候的禱告.
主人你可以容許.
但是不要讓他容得比我高.
真正一個有心胸的人應該說.
我盼望下一代比我更強.
更有恩賜的.
因為這樣的時候.
整個國度才可以擴張.
整個教會才可以興旺.
所以我們現在的教會.
更希望更多的年輕人.
更多有能力的年輕人.
有追求的.
如果是多些人的時候.
我們就不擔心教會會退後.
所以一代要比一代強.
整個教會就興旺.
我們要有這樣的心胸.
當保羅剛剛信主的時候.
你會看到巴拿巴就很提攜他.
眾人都怕了他.
但是巴拿巴就願意相信他.
接納他.
還將他帶到那些使徒的面前.
為他做見證.
後來出來宣教.
原先是巴拿巴和保羅.
後來就變成了保羅和巴拿巴.
保羅更威風.
巴拿巴好像落了下來.
雖然依然是泰山.
你會再看巴拿巴.
他後來還有一個馬可.
保羅不接納他.
他接納.
接著.

$^{1281}$當保羅晚年的時候.
他才想起這些.
你來的時候要將馬可帶來.
因為他在全道的事上.
於我有益.
巴拿巴一早就看到了.
他的心胸很廣闊.
所以一個真正偉大的人.
他不只是自己被神所用.
他是讓更多的人.
來到神的面前.
被神所使用.
所以真正的童工不是派人來超過你.
而是甘心成全人.
耶穌基督呢.
耶穌基督更偉大.
祂說你們要做那些比我更大的事.
其實是不可能的.
如果從身份來說.
因為耶穌基督祂就是神.
祂是獨一無二的真神.
祂在歷史當中.
是一個真正有神性的那一位.
所以不可能超越祂.
在身份地位上.
但是如果在服事的事上.
耶穌基督祂來到這個地上.
只是在以色列那裡.
一個很小的地方去服事.
但是我們今天的人.
當時祂的門徒.
是去到更廣闊的地方.
彼得去了羅馬.
多瑪去了印度.
馬太去了非洲.
西門和猶太.
傳統說他們去了中東那一帶.
將福音帶給世界.
所以他們的服事範圍.
是比耶穌基督更廣闊的.

$^{1321}$這個正是耶穌基督.
祂的心所要的.
就是要神的國度不斷地去擴展.
所以今天我們做牧師的.
長老的.
主日學教師的.
帶團契的小組的.
帶禱告會的.
你們都要興起.
你們的同工.
一同去服事.
約翰.
他的事奉很短暫.
可能是一年左右.
但是他花了三十年的時間.
經歷了曠野的生活.
神對他有很多的雕琢.
但是他的服事時間很短.
但是一個人.
他的價值不在乎服事多久.
真正的價值是什麼呢?.
就是有沒有完成神祂的旨意.
完成神祂的託付.
生命的價值不在乎時間.
而是在乎神祂所託付的使命.
有一位很年輕的傳道人.
叫做大衛.
很多年前,可能是幾百年前.
他有機會去耶魯讀書.
耶魯大學.
但是後來因為得罪了一個老師.
因為一句話得罪老師.
斷送了他自己的學業生涯.
最後,按他自己來看.
我沒有前途了.
但是在這個時候.
神依然為他開一條出路.
感動他去印第安人那裡傳福音.
他在那裡傳福音也很艱難.
他很多時候在森林裡禱告.

$^{1361}$為那些靈魂去流淚.
為那些印第安人去流淚.
他身體很弱,很差.
沒多久他就病了.
而且他的事奉只有四年.
離開的時候只有二十九歲.
如果他停在這裡的時候.
你覺得這個人好像擺上了.
但是這麼短,這麼苦,這麼可惜.
也沒有什麼功勞.
但是後來他的日記.
就被一個很偉大的傳道人.
幫他印了出來.
分發出去.
所以他這個日記就影響了很多宣教士.
當中有一個是宣教之父.
威廉·克里.
影響了他.
所以你會看到.
他只是事奉了四年.
但是他影響了幾百年的宣教.
神就藉著他這麼微小的.
這麼短暫的使用他.
讓他成為神手中寶貴的器皿.
約翰也是一樣.
約翰說.
他必興旺,我必衰微.
願神賜福我們每一位.
願神帶領我們,幫助我們.
專行神的旨意.
活出基督.
願我們心裡都說.
他必興旺,我們必衰微.
請低頭領受祝福.
願主耶穌基督的恩惠.
父神的慈愛.
聖靈的感動和交通.
尚與你們同在.
願聖靈感動你們.
讓你們能夠活出基督.

$^{1401}$專行神祂的旨意.
讓你們成為這個世界的光.
光照更多的人.
願你們的心說.
他必興旺,我必衰微.
願主使用你們.
阿們.
\newpage



\section{箴言 18:24}
\label{sec:031870}
\textbf{摯愛的知己好友}~~~~曹子光牧師
\newline
\newline
連結: \href{http://www.church.com.hk/acms/content.asp?site=cdc\&op=show\&type=product\&code=031870\&layout=sermon}{\texttt{031870}} ~~~~ 語音日期: 2026-03-22
\newline
\newline
\hyperref[sec:031869]{\small{< < < PREV SERMON < < <}}
~
\hyperref[sec:index_chronic]{\small{[返順時目]}}
~
\hyperref[sec:031871]{\small{> > > NEXT SERMON > > >}}
\newline
\newline
箴言 18:24
\newline
\begin{longtable}{cl}
\hline
\hline
章節 & 經文 (和合本修訂版)\\
\hline
18:24 & \begin{tabularx}{0.7\textwidth}{X} 朋友太多的人，必受損害；但有一知己比兄弟更親密。 \end{tabularx} \\ \\
[1ex]
\hline
\hline
\end{longtable}
$^{1}$我不知大家有冇朋友架呢.
朋友呢樣ye 其實都幾 點講呢.
你可以話有好多朋友又得.
或者你有幾個朋友亦都得.
因為有知己朋友架嘛係咪.
人生裡面有一首歌我就覺得係好經典ge .
唔知大家有冇聽過呢.
朋友 我當你一世 朋友.
不過呢首歌你唔好以為真係好happy.
因為佢話 非什麼大仇 為何舊知己.
在最後變做 不是老友.
變成唔係老友.
嘩 其實呢個都幾痛苦.
本來係舊知己.
但係跟住最後呢 變唔到老友.
係我ge 年代呢 我好鍾意聽一首歌.
誰不知分開只有數年.
知唔知跟住gwo 句係咩.
共你相逢 先感覺到.
兩家亦在邊.
係咪好sad呀.
經過咁多年相隔jor 一輪之後.
原來發覺兩個都改變jor .
不過都唔係ge .
我諗o下我細個ge 時候都聽過一首好精彩ge 歌.
以真誠我祝福你.
識架喎.
願將一腔熱誠給你.
常為你鼓舞.
呢D 咪老友 呢D 咪知己囉.
咁當然啦 仲有咸校長 唔係咸校長 係譚校長.
仲有咩你為我 我又為你.
都係朋友黎架.
梗呃你手朋友.
當然我亦都記得一首好經典ge 歌.
鄭伊綿.
來忘掉初對 來懷念過去.
曾共度患難日子 總有樂趣.
呢D 你同我同年代架 你唔好話唔識架.
唔知邊個鄭伊健添呀.

$^{41}$喂 誰呀.
我只能夠話無言呀咁樣係呀嘛.
中國人呢成日都講.
中國人成日都講話多個朋友.
好過多個敵人.
又或者係出外就靠朋友.
注重ge 就係gwo 種和諧啦.
gwo 種資源.
朋友ge 概念其實.
如果咁樣理解ge 話.
就係減少D 助力.
就增加多D .
減少D 阻力就增加D 助力.
呢個就係幫助我地.
係複雜ge 人際關係裡面.
搵到一條生存ge 一個ge 希望咋嘛.
而係中國ge 傳統文化裡面.
人脈係非常之重要ge .
所以呢我地追求盡量都係唔同人交惡ge .
我地可以ge 話呢.
中國人ge 智慧呢就係屬於叫做減去.
減少D 敵人就增加我地ge 安全感.
唔知你同唔同意呢.
大概呢呢個都係中國人所講ge 處世之道.
但係曾賢所講ge 似乎係另一番ge 智慧ge 說話.
曾賢所強調ge 就係萬千人ge 裡面.
你要辨認出邊D 先係真正ge 好朋友.
當然啦我地唔一定要否定多個朋友.
好過多個敵人ge 呢一個ge 處事ge 態度.
但我地必須要睇清楚就係.
gwo 一種朋友只能夠陪伴你走平坦ge 道路.
當你要走上坡ge 時候或者上斜路ge 時候.
或者你人生跌jor 落去幽谷黑暗ge 日子ge 時候.
其實唔係咁多朋友呢D 所謂ge .
同你同袍ge 平地裡面行ge 朋友.
可以陪你咁樣走法ge .
頭先有一首歌呢我好有印象.
就係呢關正傑ge 相對無言.
我記得係19年ge 時候.
我識得一個好好ge 朋友.

$^{81}$就係因為一D ge 當時ge 環境ge 改變.
就覺得大家原來係係gwo 個.
即係政治ge 環境裡面係有咁不同.
跟住連Facebook都搏埋黎.
我唔知大家有冇經歷呢D 咁ge 情況.
係呀平坦ge 道路可能好多friend.
但係一走到D 日子裡面.
其實我地就會分得清.
到底邊D 先至係真正係你ge 朋友.
我地到底要過ge 一個活得順利ge 人生.
還是箴言所教導ge 活得有根基ge 人生.
而呢D ge 朋友其實就係係箴言裡面.
係好重要ge 教導之一.
今日我地就嘗試睇兩節ge 箴言.
黎到學習三樣ye .
第一樣ye 就係知己好友.
係要講求ge 係質而唔係講求量.
我地一齊睇一睇第十八章ge 第二十四節.
我地請讀.
朋友的自取敗壞.
但有一朋友比弟兄更親密.
中文聖經返去做濫交朋友呢四隻字.
你聽落都好似有D 好負面性.
你想唔想你ge 仔濫交朋友.
你想唔想你ge 女濫交朋友.
梗係唔想ge 係咪.
有邊一個家長想你讀完書.
或者去完Sunday school之後話.
我鍾意濫交朋友.
無一個家長想見到咁ge 情況出現.
呢個四個中文字係一個非常之誤導ge 字眼.
因為原文根本冇濫交呢兩隻中文字.
其實原文係上半句有三隻中文.
三隻ge 希伯來文字.
意思返過黎就係話.
我有好多朋友ge 人.
佢受到損害.
又或者係佢自取敗壞.
又或者你返去做仇犯悲傷都得.
下半句就有四個希伯來文字.

$^{121}$好得意個箴言好多時候都係七個希伯來文字.
上三就下四或者上四就下三.
下半句就係講.
但有一個朋友ge 朋友呢個字.
其實希伯來文係一個動詞ge 分詞.
意思即係話其實呢個字.
係唔係指一個朋友咁簡單.
係一個愛呢個字.
不過佢係一個動詞ge 分詞.
希伯來文呢個字其實可能好多人都聽過.
因為好多ge 產品都用呢個字做名字.
"I have"呢個字.
呢個字其實意思即係等於知己.
所以有D 翻譯本就索性話.
但有一個知己比弟兄更親密.
所以我就基於呢一個ge 動詞ge 分詞.
我就將佢索性呢句說話翻譯做.
但有一個真正愛你的知心朋友.
遠比一位ge 親屬更親近.
呢一句ge 箴言ge 上半句話有好多朋友ge 人.
證明jor 呢一個人應該社交圈子好廣.
我唔知你同唔同意.
你係屬於一個社交圈子好廣ge 人.
還是唔係呢.
你多唔多朋友ge 呢.
問題係點解呢個咁多朋友ge 人.
竟然會令到自己受到傷害.
跟住甚至經歷jor 悲傷.
甚至好似中文翻譯咁講自取敗壞.
點解會淪落到咁樣ge 情況呢.
更加特別ge 地方就係18章24節.
箴言上半節一開始就話「個人」呢一隻字.
希伯來文呢個字叫做「ish」呢個字.
「ish」呢個字其實好有趣ge 地方就係.
跟住24節ge 下半節.
又有一個字同佢好相似ge .
同「ish」呢個字非常之相似.
希伯來文如果讀出黎叫做「ish」.
聽得聽到係咪好發音好似呀.
「ish」「ish」.

$^{161}$一個係「人」一個係「有」.
「人」就係「ish」.
「有」就係「ish」.
呢兩個字母呢有兩個字母係完全一樣ge .
所以呢個作者係故意咁樣放出黎做一個對比.
上半節ge 主角ge 呢個人「ish」.
強調jor 佢社交ge 行為.
佢好努力呀.
去經營好佢ge 人物.
所以呢佢擁有好多ge 好朋友.
好多ge 社交鋪所謂ge 好朋友.
但結果係點呀.
結果就係佢諗住雙手賺返黎ge 友情.
結果呢去到最終原來係自毀長盛.
破碎啦.
損害啦.
變成jor 惡啦.
甚至毀滅啦.
甚至不幸.
但下半句真正ge 「有」呢個字呢.
呢個「有」「ish」ge 字呢.
就係有一位充滿jor 資深「有」ge 人.
呢個格言好明顯係玩緊D 文字遊戲.
你以為你有一個擁有萬千好友ge 大人物呀.
你覺得你自己係一個「ish」.
但其實你係乜都冇呀.
相反你生命裡面有「ish」.
一個就夠啦.
就係一個知己朋友.
唔知大家有冇咁ge 經歷呢.
有一個知己朋友勝過100個你電話簿裡面ge friend.
你先真正擁有一個真正有價值ge 關係.
有兩樣ye 值得我地留意ge .
第一樣ye 就係呢度好明顯.
係一個質量對於數量ge 對比.
上半節ge 呢個人呀.
「ish」呢個人.
佢經營ge 係「量」.
佢擁有巨大ge 社交人物.
身邊圍繞住無數ge 所謂社交ge 朋友.

$^{201}$呢D 社交ge 朋友ge 裡面.
梗係希望可以搏多D 點讚呀.
拼命擴張多D 人物呀.
以為朋友越多就越安全就越成功.
我唔知大家有冇睇D youtube頻道ge 呢.
佢地梗係通常話比你聽叫訂閱佢ge 頻道架係咪.
咁呢就越多就越好架啦.
越多就越成功.
呢個係社會或者呢個文化.
所比我地ge 一個ge 價值ge 睇法黎ge .
下半節ge 另一個人呢.
佢有「ish」佢有D 乜ye ye 呢.
唔需要多呀.
一個.
一份實質ge 有愛ge 友情.
佢強調ge 就係真正ge 擁有.
一個真正愛你ge 知己好朋友.
就勝過萬千虛浮ge 社交.
咁大家都知道啦.
以前呢我地幾年呢都有做好多gwo D 關於D 社區ge 活動呀.
lor 好多funding呀.
甚至乎呢同一D ge 政府人員呢打交道.
咁呢小弟呢都有機會呢.
去參與好多唔同ge 一D 係倫敦舉行ge 會議.
咁呢有試過就去過一個地方就有國防部ge .
一個官呢走走黎呢同我地講一D 講一D 好多好多ye ge .
咁裡面呢就有一大堆gwo D 香港人ge .
香港人ge 人ge 機構係裡面啦.
咁我見到入去ge 時候呢.
個個lor 住一疊卡片出黎架喎.
見到人呢就係度你叫乜ye 名呀呢張係我ge 卡片.
D 人走埋黎ge 時候呢就問我.
你係邊度呀你叫乜ye 名呀有冇卡片呀.
跟住我翻埋成日身冇.
都唔知可以比D 咩佢話真係冇卡片係手.
我D 卡片係哂我個車.
到今時今日都仲有一個部車.
冇拎係銀包裡面ge 其實我.
咁呢就你會發現都唔知講D 乜ye 好.
有一次仲經典.

$^{241}$有一次去到首相府唐寧街十號.
咁入到去ge 時候呢.
咁呢就裡面有二百個人只得二百個人ge je .
咁入到去裡面呢.
其實我唔係好識D D 人架完全唔識架.
離遠我彷彿見到一個好似牧師ge 形狀.
咁呢就即係我相信呢.
我估ge 就係因為佢著件牧師袍ge 衫.
咁我估佢係牧師咋嘛係咪.
當唔比個角色扮演ge 咩係咪好難講架.
咁我估佢係牧師je .
咁壞人黎架gwo 個.
咁其他全部都唔識架啦.
咁呢就路遠添個個牧師.
咁我冇得費地走埋去打招呼.
咁呢就唔係好識得好.
基本上唔係好識得好多人.
我只係識得兩三個就裡面.
咁但係大家都係傾緊ge 忙於傾計.
但係每一次傾計我見到D 人呢.
又係卡片.
一走埋去就卡片架啦.
就問你係邊度做ye 呀.
有D 咩ge 事工或者機構呀咁然之後呢就可以有D 乜ye 互相幫忙呀.
我心裡面諗我對sol有D 咩可以幫到人.
即係可以幫到人啦.
咁呢就有D 咩可以祝福到你呢.
咁有時都唔知可以講D 乜ye ye .
即係如果要從人脈ge 角度去睇呢.
其實都幾痛苦架喎.
因為真正擁有ge .
就算你lor jor 成疊卡片ge 朋友.
你擁有幾多份真正ge 情義係裡面呢.
所以朋友唔真係需要好多.
甚至希伯來文話一個已經足夠.
第二樣ye ge 就係呢度係講緊擴張ge .
倒塌同埋一個凝聚ge 連結.
上半句提到有好多社交ge 朋友佢地ge 結局.
自取敗壞啦破碎啦受傷啦受傷害啦.
呢D 係向外擴張到極限跟住ge 瓦解.

$^{281}$當你試圖拉住一萬隻手.
咁你隻手會點呀.
可能比佢地扯斷jor 你都未定.
下半句ge 結局呢.
佢話比弟兄更親密.
呢個係一種內向ge 凝聚.
深度ge 連結啦.
係一個生命力掘到deepge 生命力.
佢比血緣關係更加牢固.
因為佢經歷ge 係一個選擇.
用意志去毀身係一個深入ge 關係ge 裡面.
弟兄姊妹朋友到底我地ge 生命裡面追求ge 係乜ye 呢.
社交ge 繁華.
還是gwo 份真實存在ge 愛呢.
真正ge 朋友唔係一個投陷.
而係一個持續愛ge 行動黎ge .
留意朋友呢個字.
朋友呢個字呢.
wear呢個字呢.
其實同毀滅呢個字呀.
wear呢個字呢.
佢有兩個字母係非常之相似ge .
wear呢個字ge 意思唔係指一個抽象ge 敗壞.
佢更加係一種動態ge 行動.
破碎啦.
受傷啦.
一個擁有好多社交ge 朋友.
竟然之間搞到咁樣焦頭爛額喎.
甚至受傷.
點解朋友多會問題更多呢.
但係佢基本上可以營造.
即係佢有三個咁ge 結論.
第一個原因就係因為社交ge 氣息.
好多社交朋友咁你咪要花好多時間應酬囉係咪.
假設下弟兄姊妹如果你D 識得一千個喎.
唔多啊.
一千個香港人.
咁就每人同你食餐飯ge 話你都要用一千日同佢食係咪.
你要應酬呢D 人啦.
其實都係花緊好多時間.

$^{321}$唔好話錢啊.
時間你都花jor 好多.
所以有咁多人ge 話你每一段ge 關係都攤薄到如紙一樣咁薄.
表面就四海皆兄弟.
但呢一種只不過係熱鬧ge 孤獨je .
你根本就冇去投資任何一段關係.
去到可以承托起你整個人生.
生命冇jor 承托力.
你有事ge 時候你就會跌到.
即係招頭爛額架啦.
唔係有幾多個可以分唔係有幾多ge 人係你跌到招頭爛額ge 時候.
企出黎支持你.
第二樣ye 就係自我的粉碎.
當你個個都走埋去建立關係ge 時候.
開始就慢慢就係去討好人啦呢個鍾意講一樣ye 你係討好人啦就迎合人啦.
於是乎你慢慢就冇jor 你自己gwo 種核心ge 自我的價值.
第三樣ye 就係盲目咁信任呀圈子越大雜質就越多.
如果你對邊個都係躺開度大媒ge 話就等於你係將你ge 防線拎走jor .
到一D 不懷好意或者係D 價值觀偏差ge 人混入jor 你ge 生命ge 裡面ge 時候.
你受傷ge 機率就會變大jor .
所以當你設下你ge 防禦就導致jor 你自取敗壞.
其實仍然都係gwo 句問題唔係多.
而係冇界線冇選擇冇深度.
我再一次講.
真言唔係反對多朋友而係講出一個人性ge 真相.
當關係冇界線ge 時候冇深度冇分辨朋友越多風險越大.
圈子越廣人就越空.
最終唔係人傷你就係你親手將自己交比D 錯ge 人敗壞jor 自己.
我地窮盡一生想追求上半句ge 繁華結果會係點呀.
你諗住萬千ge 好友但聖經冷冷咁英雄話比我地知.
咁樣ge 擴張最終只係帶黎破碎失望受傷.
我地每一個人都好想人地點讚我地.
每一個人都好想lor 多D 人係人脈ge 裡面.
我初初做youtuberge 時候呢都有呢個感覺.
大家都知道有時都會掟D 歌上去.
我其中掟jor 一首歌我唔係叫你聽呀對唔住呀我冇叫你點讚架.
亦都冇叫你呢個訂閱啦.
我第一次呢點我就postjor 一首歌大家可能識ge 我年代ge 話.
我ge 年代ge 話.
校長ge 一首歌叫做暴風女神.

$^{361}$你好難估到我會唱一首咁ge 歌ge .
咁我初初掟呢首歌上去ge 時候呢.
嘩都幾多人點讚架喎好奇怪突然之間我見到一個diss我係.
我心裡面好唔開心我話比你聽.
我話我心裡面諗你唔鍾意je 你可以唔比個like我架嘛係咪先.
跟住你整個dislike比我喎你係第一個整個like dislike比我喎.
我唔識佢係邊個黎架坦白講.
但係呢種ge 感覺呢係唔好受架.
當然啦後期麻木jor 啦因為越黎越多dislike嘛係咪.
可能唔同個歌越黎越多dislikege 時候你開始麻木啦呢種ge 現象.
人就係咁呀人好怕比人dislike.
你好想每一個人都認同你.
但係呢個世界並唔係咁樣.
根據呢一節ge 箴言ge 下半句佢話但有一個朋友比弟兄更親密.
正如我剛才所講好明顯同上半句係做一個對比.
量同質ge 對比好多朋友同一個知己朋友ge 對比.
作為知己好友就一定係一位真正愛你ge 人.
唔需要多個弟兄姊妹朋友一個或者就已經足夠.
有咁樣ge 一個好朋友箴言話比血緣關係ge 親屬更加珍貴.
呢一個人佢選擇jor 你成為佢ge 好朋友.
佢扮演jor 就係無論你係咩情況即係佢都不離不棄ge 愛.
就好似上帝ge 愛一樣.
當然啦箴言ge 十八章廿四節佢唔係反對結交朋友.
而係佢警告gwo D 以廣度取代jor 深度ge 危險.
試問人一生裡面真正能夠支撐你渡過黑夜ge 係D 乜ye 人呢.
唔係gwo 五百個不與你為敵ge 人.
係gwo 個比弟兄更親密ge 人 同唔同意.
有一個年輕ge 創業家佢ge 社交能力非常之強.
佢ge 手機裡面超過二千個聯絡人.
朋友圈裡面呢又非常之廣闊.
唔到十分鐘佢就可以接觸到好多人.
佢每一晚都有飯局.
唔係商業ge 應酬就係朋友ge 聚會.
係旁人ge 眼裡面.
佢簡直就係一個多朋友到爆ge 社交圈子非常之唔錯ge 一個人.
有一個晚上佢返屋企ge 路途裡面突然之間胸口劇烈咁痛.
佢就暈低係路邊gwo 度.
佢勉強lor 出手機係螢光幕裡面搵gwo 二千個ge 聯絡電話.
佢指尖一路係咁掃落去一路係咁掃.
佢發現gwo D 平時一齊飲酒gwo 班人.

$^{401}$互相稱兄道弟gwo 班人.
佢一個都唔敢打電話比佢地.
因為怕佢地覺得佢麻煩.
怕呢一份虛弱會毀壞jor 佢苦心經營ge 強者ge 形象.
gwo 個社交平台比佢點讚ge 人佢就驚得罪埋佢地.
呢D 大多都係執得出黎玩ge 人.
根本冇進入佢ge 生命ge 人.
係gwo 一刻佢體會到真面所講.
擁有再多ge 朋友到頭黎.
可能都係係脆弱ge 時間受煩悲傷.
我個好朋友其實就係開創我地SOLge 其中一位ge 創辦之一ge 一位弟兄.
佢係一個醫生黎.
呢位醫生呢佢好特別.
佢人地讀五年醫科呢佢會讀七八年醫科.
因為每兩年就要諗一年班.
佢每一次考試呢都係非常之艱難.
就算考試考到升班呢佢都arm arm 好合格ge je .
咁呢就有一次佢同我傾計.
佢話Benny.
佢話你一次ge 經歷我真係好難忘.
佢話佢gwo 次係非常即係考唔到.
有日考唔到開心.
當時係半夜三點半鐘.
三點半.
佢lor 起部手機.
係佢ge contact list入面係咁掃落黎.
係凌晨三點半.
佢可以搵邊個傾計.
佢當然諗起我啦.
但佢唔敢打比我.
因為Benny應該訓緊覺.
所以佢冇打比我.
佢再掃落去.
原來冇一個佢可以打得落呀.
係冇一個可以搵到.
跟住佢望住ge 天流著淚.
佢只係知道原來唯有上帝係比佢搵到.
係呀其實真係知幾難求呀同唔同意呀.
一個其實或者已經好足夠.
好第二樣ye .

$^{441}$知幾好朋友呢.
其實係任何時間都堅持愛你ge 人.
我地睇一睇第17章第17節.
請讀.
朋友乃時常親愛.
弟兄為患難而生.
呢一節ge 真言原文都係七個字組成.
係結構上面呢一節真言呢.
係將兩類人放係中間.
就係朋友第一類.
第二係弟兄.
上半節ge 真言一開始係講到.
時常親愛.
呢一句ge 表達呢就係一個時間ge 表達黎架.
指出任何ge 時間裡面都愛你ge 人.
而下半句真言結束呢就提到為患難而生.
我想問大家坐係度ge 人.
你坐緊係呢度係咪患難緊呀.
除非你係聽緊我講到覺得係一個患難je .
否則ge 話你冇理由坐係度係經歷緊一個苦難架嘛.
患難唔係日日都有架喎.
唔係分分秒秒都有架喎.
係一個特定ge 時間arm 唔arm .
係一個特定時間先有患難架.
所以兩個有關時間發生ge 事情.
前後呼應咁夾jor 兩個人物係中間.
就係弟兄同埋朋友.
值得留意ge 就係17章第17節呢度親愛ge 呢個字.
其實就係第18章24節出現過.
自己ge 同一個字.
就係a have呢個字.
就係愛呢個字.
17章第17節呢度所講ge 親愛.
又係一個動詞ge 分詞.
代表jor 一個繼續進行ge 一個狀態.
相對於上半句.
下半句一開始有個連接詞中文翻譯做呢可以叫做和.
但係其實呢度更加理想應該話.
一個對比或者叫但係.
咁樣就好特別啦.

$^{481}$其實佢係講緊同上一句ge 對比呀.
即係話當時gwo 個人跌落去一個困難.
警方ge 裡面患難裡面ge 時間.
親屬社群gwo 班弟兄gwo 班親友.
出手救你.
咁樣ge 處境就同上半句呢.
再出一個對比.
患難災難困境.
應該特定ge 時間出現.
而唔係日日出現架.
相對於上半節箴言所講ge 每一個時間呢.
咁就好明顯係對比.
每一個時間同一個特定.
某一個環境出現ge 困難係一個好大ge 對比.
但係呢一個人呀.
你ge 好友都仍然企係你gwo 度.
冇離開過你.
我相信呢句ge 箴言.
唔係話親屬或者弟兄唔重要.
因為經文講得好清楚啦.
就係因為gwo 個時間.
特別出現jor gwo 班你屋企ge 親人.
黎到幫你救你咁所以證明jor 其實緊要ge 呢班人.
只不過呢個箴言.
更加所強調ge 就係要突出ge 就係至愛知己ge 重要性.
呢個朋友.
係任何時間都企係你gwo 邊.
弟兄姊妹諗清楚.
諗清楚你人生gwo 度裡面.
有冇一個咁樣ge 朋友係任何時間.
都支持你.
愛你.
保護你明白你體諒你.
甚至原諒你.
係不論晴天雨天.
你處係成功或者係輝煌甚至係淪落到好似過街老鼠人人喊打ge 時候.
呢一個人.
仍然都係.
企係度支持你.
所以個愛唔係一個名詞.

$^{521}$係一個行動ge .
動詞.
ge 分詞.
當你最唔可愛最失魂落魄ge 時候.
邊一個.
仲係咁愛你.
朋友.
往往都淨係志趣相投ge 連結先可以走返埋一齊.
弟兄呢好唔同啦弟兄就係.
一種.
責任.
甚至係一個血緣ge 關係.
我再一次講呢隻經文ge 核心並唔係要貶低親屬ge 關係.
而係要突出朋友.
兄弟所代表ge 親屬關係.
係呀係法律同埋家族ge 保障.
但係朋友就係靈魂同埋生命ge 甘露.
一個能係一個能夠任何時間裡面都愛你ge 朋友.
已經係將.
好朋友或者好鄰舍呢一樣ye 昇華到一個親屬ge 地位上面.
有一個曾經好紅ge 鋼琴家.
佢因為一場ge 意外導致到佢手指全部受傷曬.
所以佢從此彈出一首好ge 琴.
佢開始飲酒啦.
佢開始黎到去放棄自己啦.
甚至呢暴躁啦破產啦.
呢一個曾經有才華ge 人明星可以話呢.
係非常之高ge 但係依家因為佢咁ge 情況之下.
所有佢D friend全部都離開曬佢.
包括佢ge 家人.
佢最後.
佢住一個非常之破舊ge 公寓連保險費都交唔起.
但係佢身邊裡面.
仲有一個人.
每一個禮拜.
十年如一日.
都走黎睇佢一次.
就係同佢食一餐飯就每一個禮拜六都黎同佢食一餐飯.
有一次呢個鋼琴家又飲醉酒就對住呢個朋友就係咁鬧落去話.
你走啦.

$^{561}$我已經唔再衰你我已經唔係再黎gwo 個天才啦.
我而家只係一個廢物.
你到底想我身上面lor 到D 乜ye ye 呢.
呢個朋友.
無嬲.
即係好安靜咁樣坐係佢身邊.
打掃掃整俾佢整爛曬佢D 玻璃.
然之後就同佢講jor 一句說話佢話朋友.
我黎唔係聽你彈琴架.
我黎係陪我ge 朋友食一餐飯咋嘛.
你彈得好唔好你彈唔彈到.
根本對我黎講.
完全唔重要.
幾年之後呢個鋼琴家死jor .
葬禮gwo 一日.
冇人黎到黎參加佢ge 葬禮呢個朋友就坐係第一行gwo 度.
成個場非常之冷清.
呢個朋友坐係佢ge 長凳gwo 度.
黎到去出席呢一個ge 朋友ge 葬禮.
唔係因為佢係一個成功ge 鋼琴家.
係因為佢係佢ge 好朋友.
所以陪到佢最後ge 一程.
呢D 先至係.
真正ge 知己朋友最後第三.
知己好朋友.
甚至係比親戚更為親密.
如果將第17章第17節.
同埋18章ge 24節拼埋一齊黎睇ge 話.
我地可以見到希伯來文呢兩節經文呢.
都係講到呢一位朋友都係一個人.
即係話其實.
唔需要多架.
根本就證明jor 呢個時常親愛ge 呢個人.
佢係好罕有.
知己根本就係好罕有.
而呢一種朋友係唔需要好多.
因為曾彥講到呢一種知心友ge 時候.
佢話佢ge 感情重甚至比血緣ge 感情.
更加.
黎到佢心口.

$^{601}$唔知大家有冇聽過李伯呢個人物呢.
李伯.
一生豪邁.
確實有好多唔同ge 階段裡面佢有好多重量級ge 好友.
至少有五個人生唔同ge 階段裡面呢.
搵到ge 好朋友唔係好多架咋五個黎講.
第一個係孟浩然.
係李伯年輕ge 時候呢非常崇拜ge .
大佢12年ge 人.
因為呢個人佢係一個非常之.
係一個比較謙卑ge 人佢唔向一D 權貴黎到屈膝ge .
第二個就係何志張.
佢係提拔李伯ge 一個伯樂.
第三個係咪讀道普還是道輔呀.
道咩話.
道輔係對唔住我真係唔識呀.
道輔呢.
佢地係.
被譽為斯坦ge 靈魂ge 交感人物黎ge 兩個都係.
非常之好friend.
咁另一個叫郭子儀係李伯ge 救命恩人.
係患難見真情.
第五個呢叫做吳芝嵐.
係少年ge 時候同阿李伯闖蕩江湖ge 劍客.
原來李伯劍客係一個非常之高劍術ge 一個ge 劍客黎ge .
好似話江湖排第二ge 當時唔知係咪啦.
你可以見到.
李伯咁出名ge 人佢人生裡面都有知己.
雖然唔多.
人生幾個階段裡面都有幾個知己出現過.
我地通常都以為.
知己同朋友梗係知朋友梗係家人係最緊要啦知己同家人ge 話.
係家人係血濃於誰.
但證經提醒我地.
親密唔單止係由血緣決定架.
而係由.
愛ge 質素去決定.
有D 關係係天生ge .
但有D 關係係選擇性.
親屬ge 關係係神安排架.

$^{641}$你一出世就已經存在架啦係咪.
你鍾唔鍾意都好啦.
gwo 個係你家姐gwo 個就係你阿爸gwo 個就係你阿媽.
我細個時候呢.
屋企呢成日都嘈軟巴閉.
阿哥呀我同我家姐成日都鬧架我試過有一次同我家姐炒足兩年冇出個聲.
因為點解呢因為就係聽收音機.
呢個時代冇人聽收音機我都知架啦.
所以我同我ge 時代相差好遠呀.
gwo 陣時我仲係聽呢個商業電台同埋呢一個ge 香港電台.
咁呢就佢揀唔知商業電台我就揀香港電台.
結果呢大家就嘈交.
嘈到兩年都冇出個聲喎.
心裡面諗同埋我成日都問呀.
阿媽你生佢好啦點解生埋我呀.
你一日就生佢一日就生我點解兩個都要生呢.
呢個出生係冇得選擇.
唔到你可以選擇.
但係朋友.
你可以選擇.
朋友ge 路係你人生ge 路上面.
係由你自己一步一步咁建立出黎.
好多時候我地反而係親人ge 面前.
唔敢完全打開我地自己喎唔知你同唔同意呢.
我地好怕佢地擔心.
好怕比佢地鬧.
甚至怕佢地誤解.
可以選擇ge 話都收埋自己.
但係當你遇到一個真正知己ge 時候.
你可能會唔掩飾啦.
可能會將你ge 軟弱呀將你gwo 種唔開心呀.
好坦蕩蕩咁話曬比你D 知己朋友聽.
因為佢接納你.
呢一種.
被看透但係仍然被愛ge 關係.
係非常之重要就係知己ge 好友.
當然啦理想裡面家人本應該就係咁.
係有達到呢一個程度架.
應該最安全ge 地方但係現實裡面我地又知道.
我地知道可能鑑於可能某一場ge 誤會.

$^{681}$傷害.
期望ge 落差.
我地慢慢學識我地收埋自己.
唔再咁講ye .
亦都唔想再講ye .
所以好多時候我地返到屋企呢唔知有冇經歷過呢D 咁ge 情況.
返到屋企食完餐飯.
匿埋個房裡面唔出黎.
到第二朝先出黎叫阿爸阿媽或者哥哥姐姐一面.
屋企係最近.
但係亦都係最遠ge 關係.
ge 一個地方.
問題唔係我地ge 關係ge 其實.
而係我地係呢個關係ge 裡面我地慢慢失去jor 安全感.
反而有一個朋友.
可能唔係日日見.
但係佢明白你.
呢種ge 關係.
唔係因為責任.
而係因為甘心.
唔係因為你.
一定要愛而係你選擇去愛.
箴言十八章二十四節佢話.
比弟兄更親密啦親密呢隻字呢.
其實大家可能都聽過啦.
就係二人成為一體.
黐埋一體ge gwo 隻字.
即係黐到實一實.
講到呢兩個人.
係有一個結連係裡面呀大家係一個.
資深有個結連.
呢一個結連係非常之重要.
我聽過個故事就係一個姊妹分享佢話做手術gwo 一日.
屋企人全部都係外地.
趕唔切返黎.
電話裡面街人就喊住同佢講.
講好愛佢呀唔使擔心呀手術一定會成功呀成功之後呢佢會即刻呢.
希望見到佢呀.
當然佢仍然都感覺到好孤單.
但係.

$^{721}$係gwo 個時候.
有個朋友佢心裡面亦在心裡面其中一個好朋友就講到.
佢冇同佢講好大段ge 道理只係握住佢隻手.
陪佢入到去手術室ge 門口.
手術之後第一個見到ge 就係呢個朋友.
當佢問呢個朋友ge 時候問.
佢話你公司咁忙點解仲會請假黎呀.
呢個朋友講jor 一句說話係話.
因為我知道.
你怕孤單.
所以我作為你ge 朋友.
依家過黎.
黎到陪住你.
我想問下大家一個問題.
你生命ge 裡面.
有冇咁樣ge 朋友呀.
係你順利ge 時候.
佢唔妒忌你.
係你軟弱ge 時候.
佢唔離開你.
係你做錯ye ge 時候.
佢同你講真話.
但佢仍然都係咁愛你.
係你跌到人生最低谷ge 時候.
佢仍然坐係你身邊.
同你一齊走過呢一段ge 日子.
而我地都有冇成為別人ge 咁樣生命裡面ge .
呢一個朋友呢.
係教會ge 裡面.
如果我地每一個人都能夠.
咁樣去同行.
咁樣一齊黎到互相守望ge 話.
呢一個彼此連結彼此承擔彼此守望生命ge 呢個群體.
我相信一定係好夢祝福.
我地因為唔係再成為一個客氣ge 親屬.
而係成為一個真正屬神ge 群體.
彼此屬於對方ge 屬靈ge 家人.
最後.
如果今日.
係你心裡面仍然覺得.

$^{761}$我好渴望有咁樣ge 朋友但我搵唔到.
又或者你覺得自己唔係咁ge 人.
係一個失敗者.
永遠都成為唔到別人ge 朋友.
我想話你知一樣ye .
針對呢度所講比弟兄更親密ge gwo 位.
時常親愛ge 知己.
去終極映射ge .
係一個人.
就係我地ge 主耶穌基督.
佢係罪人.
佢係罪人還jor 罪人ge 時候.
佢就為我地ge 罪釘係十字架上面.
佢係地上面陪罪人行走佢ge 好多ge 道路.
同佢地一齊黎到去生活.
耶穌基督就有咁樣ge 好朋友好知己.
真正ge 改變唔係由我地努力.
去成為一個更好ge 朋友開始.
而係由我地先被主咁樣愛我地作為開始.
當你知道你已經被接納.
被明白.
你先有勇氣放低你ge 面具.
去同人建立真正ge 關係.
你先至有能力.
係別人跌倒ge 時候.
唔係問點解你會搞成咁.
而係講我係度.
我同你一齊.
我陪你.
願我地每一個人.
都能夠係呢個孤獨ge 世界裡面.
靠住主ge 愛成為彼此生命裡面.
一個不離不棄ge 知己好朋友好嗎.
我地一齊禱告.
七位天父我地多謝你.
多謝你.
透過真言ge 訊息.
比我地係人生ge 道路裡面.
去搵真正ge 好朋友.
並且學習成為別人真正ge 知己.

$^{801}$主呀幫助我地係教會裡面.
去學習.
有咁樣ge 一個屬靈ge 家庭.
我地成為對方ge 祝福.
不論佢係咩環境裡面.
我地去支援佢.
去支持佢.
我地知道知己並不需要多.
所以我地.
我地好希望呢個群體裡面.
我地弟兄姊妹都建立到有一D 知己ge 朋友.
我地唔係需要每一個會友都成為我地最深ge 知己.
而係我地會諗.
到底我地會唔會成為到別人ge 知己.
我地可唔可以.
搵到係呢度搵到一個ge 知己.
主我地求你幫我地係呢個地方.
成為jor 屬靈ge 家庭.
搵到屬靈ge 戰友.
屬靈ge 同行者.
屬靈ge 知己好友.
最好ge 屬靈ge 知己.
\newpage



\section{路加福音 18:9-14}
\label{sec:031871}
\textbf{法利賽人與稅利禱告的真偽:生命的反思}~~~~梁成裕牧師
\newline
\newline
連結: \href{http://www.church.com.hk/acms/content.asp?site=cdc\&op=show\&type=product\&code=031871\&layout=sermon}{\texttt{031871}} ~~~~ 語音日期: 2026-03-22
\newline
\newline
\hyperref[sec:031870]{\small{< < < PREV SERMON < < <}}
~
\hyperref[sec:index_chronic]{\small{[返順時目]}}
~
\hyperref[sec:031872]{\small{> > > NEXT SERMON > > >}}
\newline
\newline
路加福音 18:9-14
\newline
\begin{longtable}{cl}
\hline
\hline
章節 & 經文 (和合本修訂版)\\
\hline
18:9 & \begin{tabularx}{0.7\textwidth}{X} 耶穌向那些自以為義而藐視別人的人講了這比喻： \end{tabularx} \\ \\ \relax
18:10 & \begin{tabularx}{0.7\textwidth}{X} 「有兩個人上聖殿去禱告，一個是法利賽人，一個是稅吏。 \end{tabularx} \\ \\ \relax
18:11 & \begin{tabularx}{0.7\textwidth}{X} 法利賽人獨自站著，自言自語地禱告說：『神啊，我感謝你，我不像別人勒索、不義、姦淫，也不像這個稅吏。 \end{tabularx} \\ \\ \relax
18:12 & \begin{tabularx}{0.7\textwidth}{X} 我每週禁食兩次，凡我所得的都獻上十分之一。』 \end{tabularx} \\ \\ \relax
18:13 & \begin{tabularx}{0.7\textwidth}{X} 那稅吏遠遠地站著，連舉目望天也不敢，只捶著胸，說：『神啊，開恩可憐我這個罪人！』 \end{tabularx} \\ \\ \relax
18:14 & \begin{tabularx}{0.7\textwidth}{X} 我告訴你們，這人回家去比那人倒算為義了。因為凡自高的，必降為卑；自甘卑微的，必升為高。」 \end{tabularx} \\ \\
[1ex]
\hline
\hline
\end{longtable}
$^{1}$今天我們分享的題目是.
法利賽人與瑞利禱告的真偽.
而帶來我們生命的反思.
上個星期是劉姑娘帶我們一起看.
十八章一至八節的經文.
裡面教導我們去建立我們生命.
成為一個禱告的生命.
今天當我們去進到九至十四節的經文.
我們會看到耶穌用兩個不同人的禱告.
去作為比喻成為我們的教導.
當我們初初看的時候.
我們可能會覺得耶穌是不是只是繼續教導我們祈禱呢?.
但是當我們很深入去看這個比喻的時候.
其實我們會發現.
除了教我們祈禱.
其實耶穌更加帶我們的生命.
進入一個很重要的反思.
我不知道大家有沒有試過有一些事情.
你覺得自己是很行的.
但是最後發現其實自己又不是真的這麼行的.
我記得我讀書的時候.
就跟我的同學組合了一個叫做弦樂四重奏.
就是說我們四個人一起有一個組合.
就有兩個小提琴.
一個中提琴.
和一個大提琴.
我們就一起參加校際音樂比賽.
我們開始一起練習了.
初初練習的時候很坦白地說.
你不是很想進來的.
因為好像宰雞一樣.
咿呀咿呀咿呀這樣.
但是練著練著.
我們就開始成形了.
那首歌我們越彈就越像樣了.
但是不知道為什麼.
練著練著練著到了一個位置.
總是差一點點的.
總是到了某些位置.
那裡就是有一點不合拍.

$^{41}$於是我就開始生氣了.
我就說停!停!停!.
為什麼到了那些位置我們就會這樣子.
為什麼到了那些位置.
我們就總是不太合拍.
那些同學就不太明白了.
不知道為什麼.
我們繼續練吧.
於是又繼續練.
但是彈啊彈啊彈到某些位置.
到了那裡又差一點點.
又不太合拍.
於是我就說停!停!停!.
為什麼會這樣子.
為什麼總是不合拍.
那些同學就說真的不知道.
不知道為什麼會這樣子.
於是其中一個同學就說.
梁成宇不如這樣吧.
你去找一個音樂老師來聽一下.
除了我們自己練之外.
不如我們找一個音樂老師來聽一下.
聽一下我們彈得怎麼樣.
於是我就怎麼樣.
你叫我我就去吧.
於是我就去了找我們華學音樂老師.
就說我們四個同學有一個弦樂四重奏.
我們有一首歌.
我們參加校際音樂比賽.
但是不知道為什麼練到某些位置.
我們覺得是差一點點的.
可不可以邀請你來跟我們聽一下呢.
音樂老師就說好啊.
我還很大肢的跟他說.
看看哪個同學彈錯了.
於是就來了.
音樂老師就來了.
就聽我們彈.
你彈一次吧.
就彈了.

$^{81}$彈完之後.
我已經急不及待.
放下我的小提琴.
就說老師.
哪個位置出錯了.
我心想肯定是其中一個同學.
不知道哪裡彈錯了.
接著就說梁成宇.
問題在你那裡.
那一刻.
真的逐漸想找一個洞鑽.
開始覺得你的臉很熱.
覺得很懦弱.
其實你知不知道.
你每彈到這些位置.
你就會慢了半拍.
因為你慢了半拍.
所以整個組合就不合拍了.
我說是嗎?.
是我慢了半拍嗎?.
接著那一刻.
當我以為自己很行的時候.
原來不行的那個.
是我.
過了幾十年.
兩年前有一次.
有一位資深同學跟我說.
如果我繼續是這樣.
你跟我說.
他說宇峻你知不知道.
其實你敬拜.
你總是不知為何唱著唱著.
你的情感為一到.
你會自己慢了半拍.
我說是嗎?.
原來我本人有時候會突然間慢半拍.
他說是啊.
你一唱到.
你的感情為一到.
你會不理其他人.

$^{121}$你在那裡就拖慢別人半拍了.
接著我說.
原來有時候我也真的.
也多謝敬拜隊.
有時候他們經常提醒我.
梁牧師.
你越唱越快.
你知不知道?.
所以有時候我也是越唱越快.
別人跟不上我.
所以其實有時候原來問題.
是在我自己那裡的.
所以當我們今天去看這段聖經.
我們每個人都反思我們自己的生命.
好不好?.
最近Pastor Irene跟我說.
原來我們指著人的時候.
原來有三隻手指是指著自己的.
所以今天我們藉著這段聖經.
我們一起去反思自己的生命.
好不好?.
但是你可能會想.
那麼這段聖經.
又帶來什麼反思呢?.
很多時候我們看這段聖經.
我們也有一個錯覺.
我們覺得這段聖經.
只是教我們怎麼祈禱而已.
如果你回到教會很久.
或者你從小回到教會.
你一定有聽過這段聖經的.
但是我以前很記得.
我以前覺得如果人們問.
法理賽人跟瑞利土哥的比喻.
其實在教什麼呢?.
我就會這麼說.
這段聖經應該教我們.
祈禱千萬不要太陽.
祈禱要收起來一點.
好不好?.

$^{161}$不要像法理蠢.
祈禱得太淺.
收起來.
其實這段聖經不是真的說這件事情的.
所以你不要看完這段聖經之後.
下次崇拜你.
我坐在後面好不好?.
千萬不要這樣.
這段聖經不是叫我們祈禱不要太陽.
祈禱要收起來一點.
但是其實如果我們看九至十節.
其實這段聖經早就告訴我們.
這段聖經想教導我們什麼的.
好不好?.
我們一起首先讀九至十節好不好?.
我們一起讀.
一二三.
耶穌向那些像著自己是異人.
藐視別人的設一個比喻.
說有兩個人上殿裡去禱告.
一個是法利塞人.
一個是瑞尼.
原來這段聖經說得很清楚.
這段聖經.
耶穌用這個比喻是要教導我們.
向哪些人教導.
向那些像著自己是異人.
藐視別人.
去說一個比喻.
究竟這段聖經.
又給我們什麼生命的反思呢?.
在進入經文以前.
你可能有一個感覺就是.
啊!那個不是我.
我怎麼會是法利塞人呢?.
有沒有人覺得你是法利塞人?.
是啊,你也不會舉手.
因為你覺得我不是法利塞人.
所以說這段聖經是與我無關的.
但是原來當我們真的去看看法利塞人的心.

$^{201}$我們有時候也反思.
原來其實法利塞人離我們不是很遠.
當我們覺得法利塞人離我們很遠的時候.
原來我們跟法利塞人的距離.
其實是比我們想像中近.
也帶來我們生命.
我覺得有三個很重要的反思.
我們一起看看今天分享的第一點.
今天分享的第一點.
是明白救恩不是靠自己.
我們要提防治疫.
我們首先從十一到十二節.
我們一起讀.
一二三.
法利塞人暫著.
自言自語地禱告說.
我不像別人勒索不義奸淫.
也不像這個稅吏.
我一個禮拜禁食兩次.
凡我所得的都捐上十分之一.
我們首先看這段聖經的時候.
一看就看到這個法利塞人暫著.
很多時候我們就有一個錯覺.
覺得他站著祈禱.
是不是他驕傲?.
很多時候我們祈禱也是站著的.
特別當時原來猶太人通常祈禱也是站著的.
所以他站著也不代表他驕傲.
但是我們稍後會發現他站的位置.
他覺得自己配得夠恩.
而當我們去看他禱告的內容的時候.
他真的覺得自己配得.
當法利塞人站著自言自語地祈禱.
自言自語地祈禱可以有兩個譯法.
一個就是向自己祈禱.
或者為自己的事祈禱.
而我們真的去看法利塞人這個祈禱的內容.
為自己的事祈禱是比較合適的.
我們看看這個法利塞人的祈禱.
相反他很強調自己的作為.

$^{241}$我不像這個.
我不像那個.
我一個星期做這個.
我這樣我做那個.
那麼你說究竟法利塞人的問題是什麼呢?.
我們首先一起看看這個法利塞人祈禱的大方向.
如果我們看這個法利塞人的大方向.
其實這個禱告大致可以分為兩方面的.
但是當我們看完兩方面的時候.
我們又發現他的祈禱.
其實未必跟我們離很遠.
我們看看法利塞人他怎麼祈禱.
他一開始的時候.
他覺得自己不像其他社會上被藐視的人.
我不像別人勒索不義奸淫.
也不像這個瑞利.
當時瑞利的人很討厭他.
為什麼呢?.
因為他們是為羅馬帝國做事的.
他不僅是為羅馬帝國做事.
他還經常收稅的時候收多一點.
塞進自己的口袋.
所以那時候的人覺得瑞利是賣國賊.
並且收多一點.
益自己.
或者這個世道叫做自肥.
所以那些人根本不喜歡瑞利.
所以這個法利塞人說.
我不像別人勒索不義奸淫.
也不像這個瑞利.
勒索和奸淫.
都是反映出舊約十誡的律例.
他說我不像這些人.
我沒有他們這麼壞.
那你可能想.
其實有時候我們也可能跟自己說.
因為我不像那個人這麼壞.
我試試問大家.
有誰覺得你好過殺人犯的?.
你舉手.

$^{281}$是不是?我先舉手.
所以其實法利塞人不是真的差這麼遠.
你不是覺得自己比殺人犯更壞吧?.
所以我們要小心.
我們的心情和法利塞人未必差這麼遠.
好了,然後他另外.
方向是什麼呢?.
他開始去誇耀自己超級虔誠的行為.
因為我一個星期禁食兩次.
凡我所得的都捐上十分之一.
當時舊約的律例.
是要求他們一年禁食一次的.
但是這個法利塞人一個星期都禁食兩次了.
另外他說我所得的都捐上十分之一.
如果你就這麼看.
你就說那只是十一奉獻而已.
其實不是的.
是多過十一奉獻.
為什麼呢?.
如果我們看利未寄這個律例.
利未寄說地上所有的.
無論是地上的種子.
是樹上的果子.
十分之一是耶和華的.
是歸給耶和華為聖的.
我們看這裡.
就是說當時十一奉獻的條例.
是你修成的.
因為當時是農業社會嘛.
就是說你修成的十分之一要做十一奉獻.
但是當這個法利塞人說.
凡我所有的都捐上十分之一.
是多過這個律例的要求的.
為什麼呢?.
他不只是他的修成.
或者他的收入的十分之一去奉獻.
他的收入所買來的貨品.
他也是十一奉獻的.
就是說如果他是一個農業的人.
他的修成十分之一已經是十一奉獻了.

$^{321}$但是那十分之九呢.
是他的收入嘛.
他拿去買一個麵包.
他買了麵包回來.
他會再切十份.
再將其中一塊十一奉獻的.
如果是這樣.
他的十一奉獻就多過當時立法的要求了.
所以其實他是在說神啊.
你看看我多好.
我不像這些人這麼壞.
不僅如此.
我還是一個超級虔誠的人.
那麼.
梁牧師.
我雖然覺得我比殺人犯好.
但是呢.
我沒有覺得我超級虔誠的.
那麼我試一下問大家.
試一下而已.
好.
在我們教會呢.
大家知道如果你參加服事.
我們也有一個出席率的要求的.
就是八成崇拜.
六成小組.
好嗎?.
你誠實.
誰是多過八成崇拜六成小組的.
你舉手.
是不是?.
我舉手才是我多過吧.
是啦.
多謝你的誠實.
是啦.
問題是什麼呢?.
問題不是你出席率好.
問題是.
你覺得你沒有別人這麼壞.
你這個比教會更加高的要求.

$^{361}$就叫你得救了.
所以我們要搞清楚.
是啦.
剛才舉手的那些不需要這麼害怕.
是啦.
你不是因為現在這樣.
是啦.
那你就好了.
從今天開始我就不做八成六成的要求了.
是啦.
我看到我的 helper(助手)開始這麼說.
我很害怕.
喂.
不要啊大哥.
千萬不要搞這些事情.
是啦.
那個法理賽人.
不是他禁食的問題.
不是他奉獻的問題.
而是他覺得靠著這些東西就得救了.
法理賽人的問題是.
當他覺得他靠著這些東西就得救的時候.
他就跟救恩斷絕了.
因為他覺得靠我這個行為.
靠我超級虔誠的行為.
我比社會上很多人好.
我又超級虔誠.
他靠自己而不是靠上帝.
當他靠自己而不是靠上帝.
其實他跟上帝的救恩隔絕了.
因為你不是靠上帝.
所以你發現在他的祈禱裡面.
他只是歌頌著自己的行為.
他沒有求過上帝的恩典.
他更加不覺得自己是一個罪人.
當你覺得你不需要神的時候.
其實你就跟神隔絕了.
你就跟救恩隔絕了.
幾個星期之前我分享過.
我讀中學的時候.

$^{401}$曾經贏過一個歌唱比賽.
如果大家記得.
我贏過一個之後成為了歌星的同學.
因為崇拜有錄音的.
所以我不會說那個是誰.
但是有些人知道.
還有最近他上過電視.
所以有些肢體拍了他.
就問我是不是這個.
因為有錄音我也不會在這裡講.
但是他們私底下問我.
我也很誠實地告訴你是哪一位.
他也很誠實地承認.
還問我弟弟我有沒有唱歌.
但是那一次的比賽為什麼我會贏呢?.
是因為我媽媽教我.
有些肢體就問.
其實為什麼你去到中學才會叫你媽媽教你的?.
你小學找你媽媽教你的時候.
你可能小學也已經贏了.
但是為什麼我去到中學才找我媽媽?.
因為我經常覺得我可以靠自己.
我不需要我媽媽.
我自己會唱.
我有同學教我.
我有老師教我.
我有自己的方法.
我不需要你.
法理賽人的問題就在這裡.
我靠我自己.
我比其他人好.
我靠我自己.
我超級虔誠.
我靠自己就夠了.
神我不需要你.
當我們發現我們不需要神的時候.
我們就跟神的救恩斷絕了.
但問題就是好像我小時候.
媽媽常常在我的身邊.
其實今天我們都是唯有靠著神才能夠得救.

$^{441}$我們不可以靠自己任何的行為.
或者我們任何做的事情得救.
我們是唯有相信耶穌.
好像剛才我們唱的詩歌.
我們是唯有靠著十字架的寶血.
聚德赦免.
我們才能夠恢復跟天父的關係.
成為神的兒女.
但其實神一直在我們身邊.
所以我不知道我們中間有沒有未信耶穌的朋友.
可能你都在思考生命的意義.
但今天我想說.
如果你今天有朋友你未信耶穌.
你想思考生命的意義.
可能你都試過很多不同的方法.
去找生命的意義.
但唯有我們的生命相信耶穌.
將我們的生命交給耶穌.
藉著祂的寶血.
我們得著永恆的生命.
我們跟天父和好.
成為神的兒子和女兒.
所以今天如果有未信耶穌的朋友.
在我們中間我想管理你.
你將你的生命交給耶穌.
但同樣今天我們相信了耶穌的.
我們要問我們自己.
我們有沒有一個真正悔改的心.
真正的悔改是因為知道不是靠自己.
我們不要像這個法利賽人那樣.
覺得我都比很多人好.
我的行為都真是頗虔誠.
然後就看低其他人.
而更加重要.
唯有我們有一個真正悔改的心.
我們才真正得著救恩.
有很多人會問一個問題.
今天我是否需要行為改變才得救呢?.
其實是相反的.
當你有一個真正悔改的心.

$^{481}$你的生命跟隨耶穌.
很自然你靠著耶穌.
你的行為就會改變.
你的生命就會結果子.
所以不是你的生命結果子.
有行為而得救.
而是你生命首先悔改.
生命靠著耶穌改變.
你就會結果子.
而你的生命也會有行為的改變.
所以今天我們都要問自己.
我們有沒有一個真正悔改的心.
當我們有一個真正悔改的心的時候.
我們就去到今天分享的第二個功課.
就是明白救恩是藉著耶穌.
而我們就跟耶穌建立一個正確的關係.
我們一起首先讀第十三節好不好?.
一二三.
那雖你遠遠地站著.
連舉目望天也不敢.
只垂著空說.
神啊!開恩可憐我這個罪人.
今天我們要問一個很重要的問題.
就是今天究竟我們跟耶穌.
在一個什麼樣的關係裡面呢?.
我很喜歡講我以前在少年部.
做中一栽培的時候.
一個很好笑的例子.
大家都聽過就是.
做完中一栽培之後.
我問一個中一的學生.
究竟今天耶穌是你誰啊?.
他就很天真地回覆我.
耶穌是我媽媽的好朋友.
於是我們就笑了.
旁邊的同學就笑他了.
那不就是你叔叔了.
因為我所有我媽媽的好朋友.
她都叫我稱呼他叔叔的.
今天究竟我們跟耶穌.

$^{521}$是在一個什麼樣的關係裡面呢?.
但是今天我們要搞清楚.
原來救恩是藉著耶穌.
我們才能夠跟耶穌.
進入一個正確的關係.
如果我們看這個瑞利的祈禱.
跟法利賽人的祈禱.
最少有五個很重要的分別.
如果我們首先一看就是.
這個瑞利遠遠地站著.
這個瑞利是以色列的男人.
那個法利賽人也是以色列的男人.
昔日在聖殿.
以色列的男人是可以進入以色列人院的.
如果戒低的話.
我幾個星期前的講道.
我講過了.
外族人是不可以進入以色列人院的.
外族人進去之後會被處死的.
但是昔日在聖殿的條例.
婦女也不能進去的.
婦女有婦女區的.
不是婦女木區.
是婦女區.
只有以色列的男人.
才可以進入最前面的以色列人院.
我們相信那個法利賽人.
一定是進入了以色列人院的.
但是當這裡說這個瑞利遠遠地站著.
有些聖經學者認為.
他雖然是一個以色列的男人.
但是他沒有進入以色列人院.
為什麼呢?.
他覺得自己不配.
所以第一件事.
這個瑞利他知道自己不配.
第二件事.
我們發現那個法利賽人.
他是滿有信心向神祈禱的.
神啊!.

$^{561}$我不是像他們那麼壞.
神啊!.
我做了這些事情.
但是這個瑞利他很清楚自己不配.
所以聖經說他舉目望天也不敢.
舉目望天是猶太人慣用的一個祈禱姿勢.
我們有時候也舉目向天.
我們祈禱就是祈禱神.
但是他知道自己不配.
所以聖經說他舉目望天也不敢.
第三件事我們會發現.
這個法利賽人為了自己的成就很驕傲.
我不像這些人那麼壞.
我做了這些事情.
但是這個瑞利他怎樣呢?.
垂著胸.
垂著胸代表什麼呢?.
垂著胸是代表一個極度的慚愧.
以及為了自己最愧疚的表現.
在新約聖經垂著胸這個舉動只出現過兩次.
第一次是在這裡.
第二次是當那些以色列人發現他釘死了耶穌.
而耶穌是一個異人.
他們發現自己釘錯了一個異人的時候.
他們怎樣呢?.
垂著胸回家.
垂著胸是代表對自己的罪.
一個極度的慚愧.
也知道不能靠自己.
所以第四件事我們會看到.
當法利賽人去高舉自己的時候.
這個瑞利形容自己是一個罪人.
我是一個罪人.
而最後法利賽人為自己的事祈禱.
但是他沒有求過神任何東西.
因為他覺得我不需要.
在這裡這個瑞利都為自己的事祈禱.
但是他知道他唯有靠著神.
所以他跟神說.
「神啊,求你開恩可憐我這個罪人」.

$^{601}$所以當這個瑞利祈禱.
令我們想起詩篇這個經文.
我們一起讀好不好?.
1,2,3.
「神啊,求你按你的慈愛憐恤我.
按你豐盛的慈悲塗抹我的過犯.
求你將我的罪孽洗除正信.
並脫除我的罪.
因為我知道我的過犯.
我的罪尚在你面前」.
原來這個瑞利回到神面前.
他很清楚知道自己生命裡面的罪.
他就求神憐憫和寬恕.
他知道自己不配來到神的面前.
所以他在神面前無條件承認自己的罪.
他知道自己不配.
所以他求神憐憫.
但是原來今天唯有我們都要清楚.
我們自己的生命不能夠靠自己.
我們要知道我們的生命唯有靠耶穌.
我們才能夠跟耶穌建立一個真正的關係.
你說為什麼會這樣?.
為什麼我們真的唯有知道.
我們不是靠耶穌.
除了靠耶穌我們不行的時候.
我們才能夠進入一個真正的關係呢?.
今天如果我們有法利賽人的想法的話.
其實我會形容我們跟耶穌可能有兩種的關係.
如果今天我們好像法利賽人那樣.
覺得我都很好.
我做這些事情.
我靠這些我就得救了.
我們可能會跟耶穌有一個.
我叫做給面派對的關係.
什麼叫跟耶穌有一個給面派對的關係呢?.
就是今天我都很好.
我都不像那些殺人犯.
我都很好.
教會說事奉就要八成六成.
不要說一百一百.

$^{641}$我都九成七成了.
於是你開始想.
耶穌.
今天能夠有我這麼好的門徒跟隨你.
真的是你的面子.
耶穌.
你今天有我肥成跟隨你.
你就真的有面子.
真的是耶穌你的榮幸.
我們不會說出來的.
有沒有人願意這麼說出來的?.
你不會的.
但是你裡面是這麼想的.
我經常忍不住要說一說周星馳的.
有一部周星馳的電影叫做《賭俠》.
裡面他想拜賭神為師.
於是他就拍了一段視頻.
那段視頻是這麼說的.
我聽聞有一個人叫做賭神是你.
我聽聞你賭錢都要戴顯影液晶眼鏡.
我不用的.
這裡有一隻煙屎.
只要我一擦.
就變成一隻巢的煙屎.
為什麼?因為我還沒發功.
我發功了.
變成麻將槍也可以.
有我這樣的特異功能人士來拜你為師.
真的是你的榮幸.
他的意思就是今天我這麼好.
我來跟隨你.
真的是你的榮幸.
今天我肥成這麼好.
我跟隨你.
耶穌你真的有面子.
所以今天我跟隨你.
是我給你面子的耶穌.
但是你不會說出口的.
但是可能你裡面就覺得.
今天有我這麼跟隨耶穌.

$^{681}$耶穌你真的有面子.
真的是耶穌你的榮幸.
所以你會跟耶穌進入一個.
叫做給面子派對的關係.
又或者我會覺得.
當我們好像法利賽人那樣想事情的時候.
我們會跟耶穌進入一個.
選擇性跟隨的關係.
就是怎麼樣?.
我喜歡就做.
我不喜歡就不做.
因為為什麼?.
我真的已經很好了.
是不是?.
我已經很好了.
所以我喜歡做就做.
是不是?.
我已經很好了.
所以我不喜歡做的事情.
我一定不會做.
但是今天如果我們不搞清楚.
我們的生命唯有靠耶穌得救.
你可能就跟耶穌進入一個.
給面子派對的關係.
你可能跟耶穌進入一個.
選擇性跟隨的.
但是今天.
如果你知道唯有你靠著耶穌才能夠得救.
你跟耶穌只有一個的關係.
就是師父和門徒的關係.
你的生命都只有一種回應.
就是你的生命會死跟著耶穌.
所以今天如果你的生命不是死跟著耶穌.
你要問問自己.
生命的反思.
今天究竟我們會不會都已經進入了.
法利賽人的想法.
我很喜歡用一個例子.
就是今天如果你現在聽到.
你現在聽到.

$^{721}$火警鐘響了.
放火粥了.
然後在這裡你知道只有一個人知道怎麼走.
我每一堂都找過熟悉的人玩一下.
Pastor King.
我叫他Pastor King.
今天如果這裡火警鐘響.
而只有一個人知道屠心道門在哪裡.
就是Pastor King.
你會跟誰?.
還會不會跟我?.
你不會跟我.
為什麼?.
跟我你會死的.
是不是?.
你的生命只會跟一個人走.
就是Pastor King.
所以我們同工族人都問過這個問題.
我們發現我們很久沒走過火警了.
我們可能都忘記了怎麼走.
所以我們上個星期同工會.
我們說我們要再走一次火警.
所以我們再走一次火警.
你記住跟著我們同工走.
因為我們知道怎麼走的.
行不行?.
但是今天如果這裡火粥.
而只有Pastor King一個人知道怎麼走.
我第一個就跟著Pastor King.
我就死跟他.
為什麼?.
我跟他才有命.
我相信如果今天火粥.
而Pastor King知道怎麼做.
你會推開肥城.
說肥城你借命吧.
不要阻礙我跟著Pastor King.
我要跟著Pastor King才有命的.
同樣今天如果我們知道.
唯有靠著耶穌才得救.

$^{761}$你的生命只有一個回應.
就是死跟耶穌.
因為你知道你唯有跟著耶穌才得救.
但是為什麼我們的生命不會這樣的呢?.
因為我們不覺得.
我們只是靠著耶穌才得救.
我們覺得我們可以靠其他東西得救.
所以今天我們要問自己.
如果我們今天沒有一個死跟耶穌的生命.
我們究竟和耶穌是不是有一個正確的關係?.
究竟耶穌是我們的誰?.
究竟我們真的給面子派對地跟著耶穌?.
我們真的選擇性地跟著耶穌?.
還是如果我們知道唯有靠著耶穌.
我們才得生命.
你的生命就會死跟耶穌?.
而你的生命死跟耶穌的時候.
你一定會結果子的.
你的行為一定會改變的.
但是到了那一刻.
你不要覺得你是靠著你的行為改變.
是因為你知道唯有耶穌.
你才得到真正的生命.
而你的生命就是死跟著耶穌不放.
所以今天原來我們發現.
我們和耶穌的關係只有一種正確的關係.
那個就是你死跟著耶穌的關係.
但是我怎樣我的生命才會死跟著耶穌?.
就是明白除了耶穌.
我們沒有人能夠做其他任何的事情去得著拯救.
當我們死跟耶穌的時候.
我們進入今天第三個的功課.
就是明白怎樣才被稱為異議.
而我們的生命會持續更新改變.
我們一起讀十四節好不好?.
我們一起讀一二三.
我告訴你們.
這人回家去.
被那人倒算為異了.
因為凡自高的必降為卑.

$^{801}$自卑的必升為貢.
我們看聖經的時候.
我們要明白這段聖經在當時耶穌的年代所帶來的震撼.
為什麼呢?.
我們現在看這一段有沒有覺得很震撼?.
其實也是這樣.
這段聖經我讀了幾十年了.
今天你還拿出來講.
因為我們是順著穆家福音去的.
所以一定要拿出來講.
但是每一次聖經一定給我們有教導的.
這一段聖經在耶穌當時的年代是很震撼的.
為什麼呢?.
因為法利賽人是當時社會最尊敬的人.
是當時社會很尊敬的人.
為什麼呢?.
因為他們熟悉律法.
並且謹守律法.
而當時瑞利是社會上最憎恨的人.
因為他們賣國賊.
幫羅馬帝國做事.
還拿著他們這個權柄去自肥.
所以耶穌居然說.
這個人所共憎的瑞利.
他回家.
居然被法利賽人算為義.
你就明白當時的震撼了.
今天如果我問你.
你誠實回答.
究竟梁牧師比殺人犯.
誰倒算為義呢?.
你會選擇誰呢?.
試試你選梁牧師舉手.
糟糕了,只有這麼多人舉手.
是不是?.
你就明白那種震撼了.
是不是?.
耶穌說這個殺人犯比起梁牧師.
這個殺人犯回家倒算為義了.
這就是耶穌說這句話的震撼.

$^{841}$剛才沒有舉手的那些.
我們沒有拍照,沒關係.
但其實.
很重要的是你要明白.
為什麼耶穌會這麼說.
為什麼耶穌會這麼說呢?.
因為如果我們去看回家去.
即是怎樣?.
「下去」即是從聖殿下來回到家.
其實在這裡耶穌想做一個對比.
你覺得自己在聖殿禱告很淑女.
你覺得自己在聖殿禱告.
是怎樣?.
「必蒙月納」.
但是耶穌說這個瑞利的禱告.
當他回到家.
比起你在聖殿禱告倒算為義.
為什麼?.
因為法利賽人的祈禱.
其實沒有求過神.
他沒有覺得自己需要神.
沒有覺得自己需要依靠神得著救恩.
當你知道你不是靠著神得著救恩.
你反而不被接納.
但那個瑞利雖然他做了很多壞事.
但他知道他靠著神.
他倒算為義.
一個在社會上被認同,被尊敬的人.
但因為他的生命不懂得依靠神.
對比一個做了很多壞事.
但知道自己卑微求神憐憫的人.
神說這個人回家倒算為義.
這裡耶穌要顛倒我們的價值觀的想法.
這裡的關鍵是「算為義」.
就是說上帝接納了這個瑞利.
但上帝接納這個瑞利.
他不是說瑞利真的很好.
他很完美的.
他做一個賣國賊.
他這樣自肥很好.

$^{881}$不是.
但為什麼瑞利未稱為義.
因為他知道他不是靠自己.
他懂得回到神面前求神憐憫.
相反法利賽人的問題.
不是他禁食.
不是他奉獻.
這些都是好東西.
法利賽人的問題是.
他以為這些就是他稱義的本錢.
他拿著這個東西.
他來到神的面前.
他就覺得神請你認可我.
因為我拿著這個東西.
我拿著這個東西.
我是應該得到你的認可.
但不是.
是你承認自己什麼都沒有.
你不能夠靠自己.
你回到神面前.
求神憐憫.
你才能夠成為義.
所以在這裡.
凡自高的必降為卑.
自卑的必升為高.
他不是這樣,他好像道德的說話.
他是在說今天如果你覺得你佩德.
覺得你自己有資格.
你要小心.
你沒有再依靠過神.
但今天你知道自己不行.
你唯有靠神.
你反而得到神的接納.
而這段聖經帶給我們生命一個什麼反思呢.
我曾經看過一個很出名的基督教作者.
他寫過有一本有關基督教教育的書.
他裡面有一段說話.
他特別當時是寫給一些青少年的牧者.
但我覺得這段說話不是只寫給青少年的牧者.
給我們的生命也有很多的反思.

$^{921}$他裡面這樣說.
作為青少年的牧者.
他說你千萬不要在青少年人面前.
將自己扮作一個屬靈超人.
因為這個世界是沒有屬靈超人的.
而你在青少年面前將自己扮成一個屬靈超人的時候.
只會讓青少年人覺得你遙不可及.
以及拉遠了他們和你的距離.
但他說相反.
你要讓青少年人知道.
你的生命都有軟弱和限制.
但是你是靠著耶穌跨過勝過和改變.
當青少年人看到你的生命都會有限制.
都會有軟弱.
但是你讓他們看到.
你是靠著耶穌跨過勝過和改變的時候.
他說青少年人會跟隨著你的榜樣.
在他們的軟弱和限制裡面.
靠著耶穌跨過勝過和改變.
所以今天我們不要扮我們是屬靈超人.
因為這個世界沒有屬靈超人的.
我們不要讓別人看到我們是今生不壞.
因為這個世界沒有人是今生不壞的.
但是很重要.
我們要知道原來我們的生命仍然需要持續更新和改變.
因為當我們停在這裡那一天.
其實我們可能慢慢就變成一個法利賽人.
當我們覺得我們的生命不需要繼續更新改變.
我已經很好了.
並且我是九成七成了.
教會說事奉的都是八成六成而已.
而我可以停在這裡.
我不需要再更新改變.
那一刻我們可能已經成為了法利賽人.
但是我們不知道.
所以我們的生命是需要繼續持續更新和改變.
我們可能在教會服侍了很多年.
我們可能覺得我們的聖經都很熟.
我們的事供又很忙碌.
接著我們慢慢覺得自己比旁邊的熟了.

$^{961}$我們開始對其他人的軟弱多了一份批評.
但是對自己的問題多了一份合理化.
我們開始會繼續用口說.
是神的恩典.
但是我們的生活怎麼說.
神其實我都很好.
有我事奉你.
真是你的榮幸.
所以我們要小心.
今天究竟我們仍然帶著神.
我要靠著你繼續更新改變的心.
抑或我們其實回到神面前.
說神我覺得真是挺好的.
請你點頭認可我.
當我們去到那個位置.
我們就要問自己.
我們的生命有沒有去持續更新改變.
其實我們要想.
其實今天能夠事奉是恩典.
就好像每一次當我進入去奉獻的時候.
我說其實奉獻是恩典.
今天誰能夠奉獻.
是唯有神的兒子.
和神的女兒.
才有資格奉獻.
但是今天我們是不是靠我們自己.
能夠成為神的兒子.
神的女兒.
我們不能夠.
所以奉獻是恩典.
同樣能夠事奉也是恩典.
誰能夠事奉.
神家中的人能夠事奉.
但是今天是否靠著我自己.
我自己的行為.
我自己有多好.
我能夠成為神家的一份子.
不能夠.
是靠著耶穌的恩典.
我成為神家的一份子.

$^{1001}$所以如果我是靠著恩典成為神家的一份子.
今天其實我做每一個事奉都是恩典.
而今天我這樣去事奉上帝.
是理所當然的.
沒什麼能夠誇.
所以今天如果我們對自己的事奉開始驕傲.
今天如果我對自己的事奉開始覺得自滿.
我們要小心.
我們可能慢慢進入了法利賽人的心態.
而我們不知道.
所以今天我們要問自己一個很重要的問題.
在社交媒體裡面.
我相信我們會對自己有很多不同的描述.
我們請敬拜隊.
在社交媒體裡面.
會對自己的生命有很多不同的描述.
我不知道你在你的社交媒體裡面.
對你自己又有什麼描述.
如果你上我的社交媒體.
你會列出我有哪幾個家人.
我的家人是誰.
也會說一下我讀過的學校是哪幾間.
也會說一下現在我在哪裡工作.
你也會看到我有什麼群組.
我喜歡什麼.
如果你上我的社交媒體.
你應該看到我很喜歡愛華頓.
所以我加入了很多愛華頓的群組.
昨晚很開心.
剛才有執事會主席來跟我分享.
昨晚的喜悅.
昨晚愛華頓贏了.
如果你看我的群組.
也會有很多我喜歡的音樂描述.
特別我很喜歡Beyond的音樂.
所以你也會看到很多有關Beyond音樂的群組.
最近因為人開始入伍.
我發現我加入了很多歷史片段的群組.
我開始喜歡看80年代的照片.
80年代的廣告.

$^{1041}$最近我加入了一個播80年代廣告的群組.
裡面看80年代的遊戲機廣告.
那些感覺都回來了.
什麼什麼天堂.
什麼什麼家.
很開心.
不斷在看不斷在看人和群組.
很多事情你都會在描述自己.
但是我們有沒有想過.
其實我們生命有一個很重要的描述.
就是今天我們是不是屬於神所接納的人.
我們可以生命很在意我們有不同的描述.
但是你有沒有想過.
今天你的生命是不是神所接納的呢.
但是今天我們的生命.
怎樣才是神所接納呢.
是不是我有很多事奉.
神啊.
因為柏妮可歌只有一個事奉.
但是我有三個事奉.
神啊.
因為這個教會說.
回來事奉要有八成六成的出席率.
我雖然錯過了一次.
但是我都有九成八成.
因為神啊.
你看看現在我經常上新聞看那些人多壞.
我不是像他那麼壞.
我都挺好.
於是神啊.
我應該被你所接納的.
但是不是啊.
今天聖經告訴我們.
是唯有你知道不是靠自己.
是唯有你知道唯有靠耶穌.
而你進入一個死跟耶穌的關係.
並且你常常靠著耶穌你的生命更新改變.
你的生命就為神所接納了.
所以今天當我們去看這個瑞利的生命.
這段聖經沒有說這個瑞利之後是怎樣的.

$^{1081}$但是如果我們進入路加福音第十九章.
我們會發現路加福音第十九章有另外一個瑞利的記載.
那個瑞利叫做撒該.
撒該經歷了耶穌的恩典.
他本來是一個瑞利長.
如果瑞利是人所憎恨的話.
他是那種罪犯眾憎的人.
但是當他經歷了耶穌的恩典.
他的生命有一百八十度的改變.
那麼撒該怎樣改變呢.
留在路加福音第十九章再說吧.
但是今天很重要的就是我們的生命有沒有繼續成長和改變.
但是我們的生命怎樣能夠繼續成長和改變呢.
是唯有我們回到耶穌的十字架裡面.
我們再一次去認定知道.
我們唯有靠著耶穌我們才能夠得著救恩.
耶穌說虛心的人有福.
虛心的人是甚麼人呢.
虛心就是靈女貧窮.
就是說你常常知道不是靠自己.
常常知道唯有靠耶穌.
如果只有我自己我靈女貧窮.
但是靠著耶穌我就生命成長.
唯有這樣我們的生命才會繼續更新和改變.
今天的回應詩寶貴十架.
也是我們很熟悉很喜歡唱的一首詩歌.
但是當我們再唱的時候.
我們有沒有真正感受到這個十架的寶貴.
有沒有感受到十架在我們生命裡面的必然重要性.
並且我們的生命就繼續每一日跟隨耶穌.
\newpage



\section{創世記 49:9}
\label{sec:031872}
\textbf{聖經人聖經事:猶大是個小獅子}~~~~黃克勤牧師博士
\newline
\newline
連結: \href{http://www.church.com.hk/acms/content.asp?site=cdc\&op=show\&type=product\&code=031872\&layout=sermon}{\texttt{031872}} ~~~~ 語音日期: 2026-03-23
\newline
\newline
\hyperref[sec:031871]{\small{< < < PREV SERMON < < <}}
~
\hyperref[sec:index_chronic]{\small{[返順時目]}}
~
\hyperref[sec:031873]{\small{> > > NEXT SERMON > > >}}
\newline
\newline
創世記 49:9
\newline
\begin{longtable}{cl}
\hline
\hline
章節 & 經文 (和合本修訂版)\\
\hline
49:9 & \begin{tabularx}{0.7\textwidth}{X} 猶大是隻小獅子；我兒啊，你捕獲了獵物就上去。他蹲伏，他躺臥，如公獅，又如母獅，誰敢惹他呢？ \end{tabularx} \\ \\
[1ex]
\hline
\hline
\end{longtable}
$^{1}$大家好,今次又一齊讀聖經,看看聖經人和聖經事.
今次應該是約瑟故事裡面的那一段.
不過今次約瑟已經接了爸爸去埃及.
而爸爸已經老了,即是啞國.
已經差不多離開世界,壽終正寢,要歸回列祖那裡.
做了一個遺言或是遺囑或是祝福,臨終之前.
我們一齊去看看,看看神會怎樣做事.
生命樹宣教網絡,鼓勵大家多讀聖經,一齊參與五元宣教.
聖經人聖經時,猶大試過小獅子,讓我們一齊看看聖經.
49章創世記第一節,啞國叫了他的兒子們來.
說:你們都來聚集,我好把你們日後必遇的事告訴你們.
啞國臨死時,叫了他的兒子們,你們都來聚集,我好把你們日後必遇的事告訴你們.
啞國的兒子們,你們要聚集而聽,要聽你父親以色列的話.
你見到啞國叫自己做以色列,亦表示我們以後都叫做以色列人.
無論你是我的兒子,將來發展成怎樣的支派.
叫我們做以色列.
啞國的兒子們,你再讀一次,看這個稱呼.
啞國的兒子們,啞國叫兒子們,做啞國的兒子們.
你們要聚集而聽,聽你父親以色列的話.
開始將自己改名以色列.
第八節,猶大呢?.
原因就是我今次只講猶大.
最重要的十二個支派中,猶大支派最為重要.
在救贖的線上,由阿巴喊,以撒,啞國,即是以色列,到猶大.
由猶大到一個人叫耶西,耶西生大衛.
大衛的苗裔一直到主耶穌和他的母親瑪利亞.
和他肉身養大他的養父約瑟.
約瑟和瑪利亞兩個都是大衛家的人.
是一脈下來的.
不過主耶穌是由瑪利亞同女所生.
救贖工作出來.
先看猶大,主要看猶大.
猶大這個字是Yi Hu Da.
Yi Hu,讀快了就變成Yu.
中文就變成猶大.
猶大的名字是讚美的.
其實是好名.
主耶穌的一個門徒猶大出賣主.
聖經裡面有猶大書,主耶穌的弟弟叫猶大.
其實Yi Hu Da這個名字很多人都會改.

$^{41}$因為讚美主的名字好.
猶大叫讚美的.
今天要講猶大有幾件事.
猶大將來有幾件事.
其實都是在主耶穌身上.
大家覺得主耶穌身上是這樣的.
第一,迎光時,你弟兄們必讚美你.
約瑟曾經經歷過.
十一個和菌向他下拜.
十一顆星向他下拜.
還有太陽月亮向他下拜.
就是約瑟.
其實都是告訴約瑟將來會做埃及的宰相.
但這裡是說猶大將來你的弟兄必讚美你.
讚美,用到讚美的字.
是讚美神的字.
你手必捏住仇敵的頸項.
你父親的兒子們必向你下拜.
即是猶大將來會做王.
因此眾兄弟都向你下拜,都讚美你.
而且你還會戰勝仇敵.
是很有力的王.
捏住仇敵的頸.
讚美你,因為你是王.
第一點是讚美你,因為你是王.
是一個尊貴的王.
而且在眾多兄弟中,只有他會站在這個位置.
第二點是猶大試過小獅子.
猶大是什麼呢?.
除了是王,他是小獅子.
然後是阿爾亞.
這個阿爾亞就是亞國以色列跟猶大說的.
兒子,你找鳥食便上去.
其實是你想去找食.
不會坐在那裡,讓牠來撞到我口.
是你上去撲去拿食物.
你屈身下去.
吃完東西之後你屈低身.
趴在那裡的時候,是臥如弓屍.
你好像小獅子那樣,撲出去找食物.

$^{81}$你吃完東西開始伏在那裡的時候.
就好像弓屍那樣.
當你站在那裡,蹲在那裡的時候.
蹲如舞獅.
你看到這裡特別說了幾樣東西.
就是你像小獅子.
所以第二點.
亞國說猶大是什麼呢?.
是獅子.
獅子象徵什麼呢?.
獅子是萬獸之王.
獅子是象徵王.
在主耶穌的身上.
或者在將來的啟示錄.
或者在以色列的書裡面.
都有出過一些活物.
有四個面或者四個活物.
是獅牛人鷹.
而這個獅象徵的是王的一面.
而在馬可路加約翰四呼音.
很多人都是用獅牛人鷹來表達.
馬太呼音就是說主耶穌是王.
因此有人來送黃金與香茉藥.
送黃金給他.
問他生下了猶太人之王在哪裡呢?.
就是馬太呼音.
如果說獅子.
在以色列是沒有老虎的.
在巴勒斯坦地中東是沒有老虎的.
有獅子有狼等等.
你是小獅子.
是公獅是舞獅.
你小的時候就是小獅.
長大了就是公獅舞獅.
很厲害.
誰敢惹你呢?.
這個「惹」字呢.
中文「惹你」.
Kum.
Kum這個字是直立.

$^{121}$誰敢在你面前站直身呢?.
這樣呢?.
沒有.
因為每個人都見到你都害怕.
好像見到皇帝一樣.
在皇帝面前誰敢站直身呢?.
沒有的話.
全部都是在恭敬.
進去後都是低著頭.
不叫你抬頭.
你根本不敢望.
是不可以望皇帝的.
不是這樣.
所以有套卡通片叫做.
埃及王子.
美國拍的埃及王子.
埃及王子要去.
Let my people go.
跟法老講的時候.
埋到他兩個面幾乎貼在一起.
好像現在那樣.
你讓我的百姓離開.
不會啦.
如果在埃及的時候.
你都不知跪到多遠.
他坐在上面.
以前你在中國見到皇帝.
哪裡見到的.
在很遠就已經跪下了.
很多級樓梯.
就不會埋到上去.
Let my people go.
讓我的百姓走.
是拍戲而已.
是卡通而已.
給小朋友容易明白.
說他是「獅子」.
誰敢在你面前站直身呢?.
你是在說他的威嚴.
他的莊嚴.

$^{161}$好像獅子一樣.
第三樣.
很有趣.
有兩個字.
「龜」.
「龜必不離猶大」.
「象必不離開他兩腳之間」.
「直等細鱸來到萬民都必歸順」.
只是黑色字和紅色字.
是和合本的經文.
第十節.
「龜必不離猶大」.
「象必不離開他兩腳之間」.
「直等細鱸來到萬民都必歸順」.
我解釋一下這段.
因為我很注重解經.
「龜」.
「Sherbet」.
「Sherbet」的意思.
不要讀錯.
不是「Sherbet」.
不是我們現在吃甜品凍的冰.
「Sherbet」是權杖或指揮棒.
拿著它就表示權杖.
現在我們都很好笑.
見到香港小姐競選.
獲選了.
先有一支權杖.
有冠冕.
有權杖.
有斗篷.
拿著那支很長的.
跟著走一個圈.
權杖.
選到做Queen.
當然有權杖.
當然那支杖沒什麼權.
如果你有機會去禮儀教會.
譬如我自己是聖公會中學出身.
有時見到一對聖人.

$^{201}$在聖壇舉行聖祭.
吹前的時候有很多人.
十字架高舉在前面.
然後有屍斑.
一半的屍斑.
有人拿著一本很大的聖經.
「道成肉身來到我們當中」.
後半的屍斑.
然後有一個小朋友拿著一支杖.
不是很長.
下面小一點上面大一點金屬銀色的頭.
後面有一個穿黑衣服的牧師.
今天講道.
講道的有一支權杖在前面.
他拿著.
表示什麼呢?.
他今天是在講上帝的道.
他是有權威的.
最後有一個小朋友拿著一支.
彎的,高的牧羊人的杖.
後面又有一個大人牧師.
這就是那間教會的主任牧師.
他通常是最高位置走到最尾.
當然如果主教來.
他有另一個主教的象徵.
走在最尾.
有一個象徵然後主教走在後面.
所以譬如去聖公會聚會.
去禮儀教會聚會.
你見到遲到跟著隊伍進去.
不行.
越走後面越高級越大.
跟著主教進去真是很好笑.
你還要比他大.
不過你都不用擔心.
因為人們已經攔著你.
你不會進去.
這個龜就是那支權杖.
不過中文叫龜.
但是在中國的龜的應用.

$^{241}$與當日埃及人的應用.
是有點不同的.
龜是權杖不離開猶大.
即是猶大一直都有權.
有權兵.
第二杖是必離開他兩腳之間.
那支杖是他拿著.
他才會在兩腳之間.
那這個杖是甚麼呢.
Kakar.
Kakar這裡是K.
K這裡是K.
Kakar.
Kakar的意思.
Kakar.
這個就是頒令立法的意思.
所以那支杖不是象徵牧養.
是象徵甚麼呢.
我是有拿著法律.
我講的就是法律.
我講的就是御令.
就是頒佈命令.
所以我們是頒佈法例.
頒佈命令.
國家的權與法令的頒佈.
都在猶大那裡.
直到何時呢.
不會離開.
一直即是說不會離開.
直到必不離.
即是說不會離開.
必不離.
一直延伸到何時呢.
延伸到世羅來到.
這件事就難明了.
如果你看一般和合本.
這個世羅.
Shiloh.
就譯了一個音.
Shiloh.

$^{281}$Shiloh的意思呢.
如果我們和合本是譯作.
次平安者.
Shiloh 次平安者.
一看次平安者是誰呢.
耶穌基督嘛.
猶大一直拿著.
以色列皇權.
與立法命令的權.
一直到耶穌基督彌賽亞來到.
Shiloh他們中文和合本.
是譯作次平安者.
但其實Shiloh這個字.
應該是屬於他的.
直到屬於他的來到.
甚麼意思呢.
直到屬於他的來到.
屬於他指誰呢.
屬於猶大的.
屬於他的來到.
這個他的意思呢.
就是屬於他的來到.
這個他呢就是指猶大.
直到甚麼時候呢.
是有一個屬於猶大的來到.
這個通常為甚麼會譯作次平安者呢.
他可能是在暗喻.
耶穌基督彌賽亞.
萬民都必歸順.
到那個時候萬民都會歸順猶大.
歸順呢.
葉克哈順服.
萬民都會順服在他面前.
應該是我們基督徒說.
是主耶穌在來.
大審判.
全世界都順服.
新天新地.
給些圖片看看.
龜和象.

$^{321}$你看這些圖畫.
這些人像.
這些公仔就會知道.
是埃及時代的樣子.
通常他們都很修長.
這個是女的.
這個是男的.
上身多數都不穿衣服的男.
圍一條布.
很漂亮.
你會見到.
這個呢.
這個樣子.
尖嘴的.
這個是食蟻獸的像.
食蟻獸的樣子.
食蟻獸.
伸條舌頭出來食蟻.
食蟻獸是神像.
不過是男性.
食蟻獸都做到神.
埃及什麼都是神.
神奇無敵.
很有特色.
他們會拜他.
希望拜他自己都有這樣的能力.
這個呢.
青蛙.
大家看到.
當然有不同性別.
青蛙是他們的神.
大家記不記得十宅裡有蛙宅.
青蛙自在.
就是這樣.
你見到蛙宅.
剛才說龜.
看回聖經.
說龜與象.
權象與立法命令.
那誰是龜呢.

$^{361}$這個就是龜.
拿著一條鑰匙.
這支就是象.
象的意思就是立法.
龜就是表示王權.
這個象表示立法.
這個表示權.
龜與象.
近一點看.
龜原來是這樣的.
他拿著那支鑰匙.
特別是行禮儀式時.
常常要拿著.
拿著不放.
拿著到處走.
表示他有立法之權.
有控制之權.
龜是這樣的.
我們很容易看到.
有十字架.
是否這麼早已經有十字架呢.
不是的.
只是剛好找出來.
有十字架的影子.
有自己的象徵.
這裡也有一支象.
可能有鑰匙.
這個龜是表示有權力者.
這個就是我們見得最多的.
木乃伊的外殼.
打開裡面有很多層.
才到人的屍體.
看不看得到他交叉手這樣拿著.
這個交叉手拿著.
如果是這一支.
和這一支.
龜是誰呢?.
龜是這個.
這些是水紙.
象就是這個.

$^{401}$也是這樣拿在他身上.
他們可能去到埃及.
見到埃及的王.
龜與象.
其實以色列以前.
是沒有什麼龜與象.
只不過可能去到埃及.
見到埃及的法老.
有這兩樣東西.
是表示掌權與立法.
很厲害.
這個應該歸在猶大的身上.
學了也說不定.
而且第四點要說的.
猶大把小廬山在葡萄樹上.
將廬區山在美好的葡萄樹上.
葡萄樹上,葡萄樹上.
即是用疊句.
小廬山在葡萄樹上.
廬區即是小廬的媽媽.
都山在好的葡萄樹上.
即是有甚麼呢?.
有廬,有牲口.
而且有廬區.
即是有小廬.
其實廬區的意思本身是小廬.
這是一個小小疊句.
主耶穌騎著廬驅子.
即是主耶穌叫門徒帶著廬驅來.
去到的時候見到廬驅.
小廬和母廬.
你這樣帶著廬驅走.
人們便會問你.
你為甚麼帶著廬驅走?.
主要用牠.
人們便會給你.
真是這樣發生.
帶著廬驅來.
要帶著母廬一起跟著.
否則廬驅不安.

$^{441}$牠不聽你的話.
所以有小廬.
小廬即是有大廬.
有母廬.
有母廬即是有公廬.
有小廬,有廬區.
即是有很多牲口.
而且還有很多葡萄樹.
更是好的葡萄樹.
好的葡萄樹.
這個猶大將來會將葡萄酒中洗衣服.
用葡萄酒來洗衣服.
然後用葡萄汁來洗葡寡.
葡寡即是外衣.
外葡,最漂亮的衣服.
你看到第十一節.
第四點就是說.
猶大將會十分富裕.
有廬,有葡萄酒,有葡萄汁,洗衣服.
多得很厲害.
而且是好的葡萄樹.
是很富裕的皇帝.
另外,這個富裕的皇帝是怎樣的呢?.
這個當第五也可以.
眼睛必因酒而紅潤.
牙齒必因奶而白亮.
這個可以當作螢光字第五.
就是甚麼呢?.
眼睛紅潤.
牙齒白亮.
是因為有酒有奶.
很有錢,富有.
另外就是身體很健康.
第五點如果要有的話就是健康.
如果歸納第四點.
第四點就是表示.
他很富裕.
富裕到喝酒喝到眼睛都亮光.
牛奶多到牙齒眼睛都亮白.
很富裕.

$^{481}$這個就是雅各以色列對於猶大的祝福.
就是說他將來是非常富有.
是掌管萬有有權柄的主.
萬民都歸順.
沒有人敢在他面前直立.
大家知道吧.
我們一起禱告.
主要在這件事上.
看到雅各以色列對他的兒子猶大.
有特別大的祝福.
原因就是他是你所揀選的.
在阿巴罕以撒雅各到猶大這個人.
將來救贖的線就是他一直下去.
直到大衛直到主耶穌.
你是透過雅各對猶大的祝福.
表達出你是怎樣祝福得勝的尼賽亞.
讚美你.
我們都一起為尼賽亞感恩.
讚美幫助我們.
更加信靠你.
知道你是掌管萬有.
歸與將不離開你.
你是最有權力的.
你是最會伸張正義立法的.
讚美主.
願你得到榮耀.
祈禱奉耶穌基督名求.
阿們.
\newpage



\section{詩篇 8:1-9}
\label{sec:031873}
\textbf{從美學沉思宇宙與人生}~~~~梁燕城博士
\newline
\newline
連結: \href{http://www.church.com.hk/acms/content.asp?site=cdc\&op=show\&type=product\&code=031873\&layout=sermon}{\texttt{031873}} ~~~~ 語音日期: 2026-03-22
\newline
\newline
\hyperref[sec:031872]{\small{< < < PREV SERMON < < <}}
~
\hyperref[sec:index_chronic]{\small{[返順時目]}}
~
\hyperref[sec:031874]{\small{> > > NEXT SERMON > > >}}
\newline
\newline
詩篇 8:1-9
\newline
\begin{longtable}{cl}
\hline
\hline
章節 & 經文 (和合本修訂版)\\
\hline
8:1 & \begin{tabularx}{0.7\textwidth}{X} 耶和華－我們的主啊，你的名在全地何其美！你將你的榮耀彰顯於天。 \end{tabularx} \\ \\ \relax
8:2 & \begin{tabularx}{0.7\textwidth}{X} 你因敵人的緣故，從孩童和吃奶的口中建立了能力，使仇敵和報仇的閉口無言。 \end{tabularx} \\ \\ \relax
8:3 & \begin{tabularx}{0.7\textwidth}{X} 我觀看你手指所造的天，並你所陳設的月亮星宿。 \end{tabularx} \\ \\ \relax
8:4 & \begin{tabularx}{0.7\textwidth}{X} 人算甚麼，你竟顧念他！世人算甚麼，你竟眷顧他！ \end{tabularx} \\ \\ \relax
8:5 & \begin{tabularx}{0.7\textwidth}{X} 你使他比神微小一點，賜他榮耀尊貴為冠冕。 \end{tabularx} \\ \\ \relax
8:6 & \begin{tabularx}{0.7\textwidth}{X} 你派他管理你手所造的，使萬物， \end{tabularx} \\ \\ \relax
8:7 & \begin{tabularx}{0.7\textwidth}{X} 就是一切的牛羊、田野的牲畜、空中的鳥、海裡的魚，凡游在水裡的， \end{tabularx} \\ \\ \relax
8:8 & \begin{tabularx}{0.7\textwidth}{X} 都服在他的腳下。 \end{tabularx} \\ \\ \relax
8:9 & \begin{tabularx}{0.7\textwidth}{X} 耶和華－我們的主啊，你的名在全地何其美！ \end{tabularx} \\ \\
[1ex]
\hline
\hline
\end{longtable}
$^{1}$我記得我開始思考問題的時候.
當時十三歲.
我就買了一本莊子來看.
看完莊子第一卷叫《逍遙遊》.
什麼叫《逍遙遊》呢?.
我知道我們那個年代有些明星叫呂琦,陳寶珠.
還有謝霆鋒老爸,謝賢.
當時那些戲經常叫《青春逍遙遊》.
為什麼偉大的哲學家莊子也說《逍遙遊》呢?.
他很奇特地說.
他說在北海有一條很大的魚幾千里那麼大.
忽然間一變 變成一隻飛鷹大鵬鳥.
由幾千里那麼大飛上天.
由太墟來到觀賞下面的人間.
小時候也不是很明白.
當時還要看文言文.
是強行要看清楚一些東西.
但後來慢慢明白了他說的道理.
魚可以變成鷹,大鵬鳥.
即是說人生可以轉化.
我可不可以飛到天上.
我要飛到天上 後面有首歌.
我可不可以做個飛到天上來到俯拆人間的.
那種很優美的生命呢?.
這個就成為那個時期.
我的一些很重要的思想.
其實莊子是這樣說的.
他說到天地有大美而不言.
四時有明發而不說.
我看不清楚 因為我的相機壞了.
我先看這裡.
而萬物有誠理而不說.
四時有成發而不議.
很奇怪的就是天地有大美.
不用說也有.
四時的變化有規則.
不用說也有.
但萬物有已定的規則,法則.
不能討論的是有的.
這個宇宙有些很奇妙的東西.

$^{41}$就是世事總會不斷變化.
而不斷中又有規律的存在.
這就是規律讓我開始去明白.
一種很特別的美的存在.
美是什麼呢?.
當你看到一樣東西.
你突然間有一個很深的感動.
感動你感覺很大的美.
而同時它成為一種規則.
你發覺原來美是有共通性的.
在這種美的共通性裡面.
我們開始感覺到宇宙有些無法描寫的特質.
這個叫美.
除了天地的美.
當我第一次見我太太.
未婚前的女孩時.
也覺得很美.
美是什麼呢?.
因為她的單純,溫柔,美善.
後來到現在.
我見到我太太已經是50年前了.
我們認識的時候已經70多歲了.
還是那麼美.
不過當然身體多了很多問題.
到差不多一年多.
這個沒辦法.
我們總會經歷這些種種的困難.
人生還有種種的困難.
但是在這種困難裡面.
我們會去思考.
美是一種很特殊的經歷.
天地間有一種美的意境.
而沒辦法用人的語言去描述.
好像大衛的詩說.
你的名,又說我主你的名.
在全地何其美.
你將你的榮耀彰顯於天.
你讓敵人的緣故從陰霾.
吃禮的口中建立能力.
似仇敵而報仇的閉口無言.

$^{81}$並你所陳設的月亮星宿便說.
人算什麼,理景顧念他.
世人算什麼,理景眷顧他.
你叫他比天使微小一點.
並設給他榮耀尊貴為冠冕.
詩人說道.
上主,你的名是何等的美.
這個何等的美.
讓我們忽然間去感覺到.
就是一種寬宏的榮耀彰顯於天上.
好像我小時候,我記得我望天.
放紙妖,我發覺放紙妖的時候.
看到夕陽西下,那時六歲.
十三歲看莊子,六歲已經感覺到.
就是莊子說的,天地有大美而不言.
那時小時候,什麼叫夕陽西下.
我不懂得想這些名詞.
但覺得美,一路到星空出來.
然後不肯下去吃飯.
媽媽要拿飯菜上來,起天盤吃.
這時候感覺那種美.
是感覺,但它是普遍的價值.
而且是你一種特殊的覺醒.
在這種覺醒裡,我們會明白.
原來聖經提出一個很特別的原則.
就是說從嬰孩和吃奶的後宗建立能力.
嬰孩很純真.
你看這個寶寶,樣子很純.
認不認得她是誰,這個就是我.
還寫著1952年,眼睛小.
那時真的很純,耳朵一樣大.
我現在七十多歲,當然比那時更加奸.
但有一樣要學習.
就是可否維持像我年輕時那麼純真.
只有純真的心才能接近天地的美.
接近上帝的愛.
維持內在的單純.
這個單純使敵人閉口無言.
敵人報仇的就是反對美善的勢力.
這個世界很多人有很多權力的追求.

$^{121}$有情慾的追求,好像愛潑斯坦.
把美善摧毀,建立一些罪惡的共同體.
但人不同的地方在哪裡.
因為他很像嬰孩,他自然流露對人的愛.
對父母的信心,他叫父母抱他.
而這種單純如果維持到七十多歲.
你就可以摧毀世間很多的邪惡.
而這些邪惡令我們感受不到天地的美.
但當我們知道是觀天地的美.
去看生命的奧妙,人心的存在.
人能夠去思考,去反省,去感受.
有愛,有良知,由知到美.
這是何等奇妙的事情.
而這種當我們嘗試觀賞天地.
我們不禁感覺聖經所說的.
人算什麼,你竟顧念他.
這裡聖經嘗試表達人的一個特質.
就是人是弔詭的,他是偉大同時是渺小.
人與上帝與宇宙相比.
其實真是人算什麼,人很渺小.
我小時候看兒童樂園也知道.
原來人在地球很渺小.
而地球在太陽系也很小.
太陽大過宇宙一百萬倍.
但太陽系在宇宙也是很渺小.
在宇宙無數的星河,無數的星宿當中.
我們很渺小,人算什麼.
但我們被顧念是因為上帝與宇宙有一種親情.
他會將我們視為重要的.
所以讓我們有榮耀,尊貴為冠冕.
這就是人的定位.
這個定位就是人既有責任要管理宇宙萬物.
但另一面其實他很渺小.
又渺小又尊貴.
但接著我們開始思考.
我們要沉思宇宙,宇宙人生本來就是一個奧秘.
它浩瀚無盡,我不禁要問.
為什麼有這一切.
我小時候也問過你.
為什麼有這一切,為什麼天地這麼美.

$^{161}$為什麼有我在這裡看天地.
為什麼我會覺得美.
再想下去就是一個最深的哲學問題.
為什麼世間有東西不是沒有東西.
Why are there something, why are there nothing.
如果原本沒有東西就是沒有.
又沒有我,又沒有萬物.
不用想太多,根本不存在.
但為什麼有東西存在.
有我存在,出生在我媽媽.
但見到有父母,又要吃奶奶.
接著長大,有公公,婆婆,爺爺,嬤嬤.
有很多親戚,姨婆.
那時候有很多姨婆.
有九姨婆,三姨婆,七姨婆.
每個人不同,那些人是七口宮.
每個人又不同.
那為什麼有這些東西.
為什麼有太陽,為什麼有星宿.
為什麼有美,為什麼星河.
晚上看見很美.
為什麼有這一切.
這些東西是怎麼開始的.
這個就涉及我們對宇宙的探索.
就是說其實在中世紀以來.
人類對於宇宙的起源.
通常都會提到上帝.
但第一個開始不想提到上帝的.
那個叫做Laplace.
Laplace是一個法國的天文學家.
當時他寫了一本書給拿破崙看.
講天文學.
拿破崙看過序的時候.
為什麼你不提上帝.
因為那時候每個科學家都講上帝.
牛頓寫他的物理學.
他開頭就說上帝寫了兩本書.
一本就是宇宙萬物.
一本就是聖經.
所以歸榮要給上帝.

$^{201}$但這個Laplace都不講上帝的.
Laplace現在我們不需要這個假設了.
因為在我們這個年代.
當時的科學是這樣認為十九世紀.
如果萬有引力定律.
就應該將所有東西合起來做一點這麼小.
那為什麼宇宙張開呢.
除非物質宇宙無限這麼大.
又永恆這麼下去.
永遠地撞來撞去一個運動.
宇宙就是一個無限的物質宇宙.
永恆無限.
所以不是上帝永恆無限.
物質是永恆無限.
所以就將物質視為宇宙的最後根源.
唯物論的根源就是這樣來的.
而由這個唯物論的觀點出來.
就認為宇宙沒有神.
就是無神論那樣來的.
但怎知道二十世紀就有重大的改變了.
就因為二十世紀人開始發現.
在哈勃的望遠鏡開始發現宇宙正在膨脹.
在曼威爾遜山上.
你看這張照片是最厲害的望遠鏡.
發現那些星的光譜是向紅色那邊移過去的.
星的光譜向紅色那邊移過去.
表示什麼呢.
就是說所有的星星都是以高速度離開地球.
星星以高速度離開地球是指什麼呢.
就是說宇宙正在膨脹.
宇宙膨脹不是無限的.
宇宙是正在膨脹的.
是動態的不是靜態的.
如果宇宙是這樣膨脹.
就代表宇宙有一個起點.
如果宇宙有一個起點.
然後就引你很有興趣討論.
就是有一個比利時的神父.
是天文學家叫做勒梅特.
他就提出宇宙是有開始的.

$^{241}$而他認為宇宙有一個大爆炸開始.
就爆出來.
當然那時候很多人反對.
就是因為無神論是反對.
你這樣說就是說宇宙是有上帝的.
有開始的.
你是神父當然是這樣說.
但是誰知道大家就去找.
如果真的宇宙有開始.
應該有些射線被人找到.
我的科學家就去找一些資金.
看看是不是detach到宇宙有那些射線.
叫做背景射線.
Cosmic Background Radiation.
那麼個個都還沒找到的時候.
忽然間有兩個搞電話的科學家找到.
因為原來電話公司那部雷達比較漂亮.
他就能夠接收到宇宙那些射線的資料.
所以當那些射線他接收到的時候.
他們馬上發表論文.
因為他本身都是很厲害的科學家.
不過找不到工作.
就去了電話公司做.
不過應該是大學做教的.
誰知道沒有位置.
誰知道發達.
用電話公司的雷達收到這些射線.
接著他們就拿到貝獎.
就發現了這種宇宙最後的射線的存在.
如果真的寫出宇宙的起源.
那起源呢?.
追回去起源是怎樣的呢?.
那真的很複雜.
如果你追查到宇宙的起源.
你就會發覺宇宙是追到10的.
-36次方那麼多秒.
這個叫做本落黑時間.
那麼宇宙如果在這個本落黑時間.
最開始是-36次方那麼多秒.
即是你明不明白這個叫做-36次方.

$^{281}$即是10之後36個零分之一.
10之後零零零零零零零零零零零零零零零零零零零零零零零零零零零零零零零零零零零零零零零零零零零零零零零零零零零零零零零零零零零零零零零零零零零零零零零零.
之一秒.
嘩!這個秒有多小啊?.
宇宙是這樣開始的嗎?.
嘩!那麼呢?.
那麼再追下去原來做一個叫做-43次方秒就是科學的那些規律就是這個時候開始的了那麼-43秒-10的-43次方秒之前又是怎樣的呢?.
嘩!未有規律的時候又是怎樣的呢?.
是沒有的就是沒有的.
所以最後發覺宇宙是不是由沒有開始的呢?.
那麼當時呢就是因此如果宇宙由沒有開始而為什麼會變成有?為什麼有規律放進去?.
最後是一個科學解不開的問題.
那麼所以接著有一個另外一個大天文學家就是Jasper就是哈布之後他就掌管那個天文台.
他寫了一本書就叫做God and Astronomer.
就是上帝與天文學家.
當他寫這本書的時候他提出的結論就引起了很多人的興趣.
我們研究那麼多那麼多.
嘩!宇宙又大爆炸.
抑或它又起膨脹.
又不知道怎麼開始.
講了那麼多到現在.
講到終於找到最開始的時候-10的-43次方秒.
接著怎樣?.
不是我們知道的了.
所以他說我們好像那些科學家攀上一個無知的高山.
去到高山的山頂征服最高的山峰.
就是我們上到頂最開始是怎樣的?.
怎知道上到那裡已經有人坐在那裡.
怎麼會有人先來?.
不是我們科學家最厲害的嗎?.
那麼你讓誰來?.
那幫是我們神學家.
我們是神學家.
原來我們信神的那幫早就坐在那裡了.
你經常說不承認我們是真的.
你一直找到尾又找到我們在那裡.
嘿嘿嘿.
嘿嘿嘿是我家的.
那麼Yasha就提出這個說法.
形成當時很有趣的一個說法.

$^{321}$接著就是到1992年.
就是要正式去研究這種的科學.
就是當時在NASA的COPY加上研究.
用人造衛星找這些背景射線的資料.
最後弄到這個.
就是電腦弄到這個相片出來.
就是元初宇宙起源那種射線.
今天的分佈狀況.
符合他們預測的.
那麼就是宇宙真是由無變成有.
是規律不知為什麼會出現.
那麼當時的George Wood大科學家說.
It's like looking at God.
紐約時報經常是頭條新聞.
他說我好像看到上帝.
科學家說看到上帝?不是吧?.
結果真是科學家說看到上帝.
這個就是George Wood的樣子.
你看見他傻傻的.
但是他拿著一個叫諾貝爾獎.
那麼他自己人問他.
為什麼你好像看到上帝.
他說我看到上帝的雙標.
雙標是什麼?.
我記得我去巴勒斯坦那一帶.
在伯利恆那一帶.
有些巴勒斯坦人問我.
Hey are you Japanese?.
我最怕人當我是日本人.
所以我最好有件衣服去旅行.
I am not Japanese.
Don't say sayonara to me.
那麼我就說.
No no I am Chinese.
You have to say 謝謝 你好.
那麼他就說.
Oh China.
You are Chinese.
Made in China.
Everything made in China.

$^{361}$他說樣樣都是中國製造.
因為樣樣都有中國的雙標.
那麼你再想想樣樣東西都有.
哪裡製造哪裡製造.
但是Cosmos宇宙.
那個雙標在哪裡?.
Joseph Wood說那個雙標就是.
宇宙這些的射線.
這個射線就是留下來.
有一個非常複雜的精密的調整.
這個就是宇宙的DNA.
就是宇宙有種密碼.
安排了它的規律在那裡出現.
然後再進一步.
不單止宇宙有.
而且生物體都有.
那麼晚的科學家都發現.
生物體裡面的information content.
就是DNA RNA都是密碼來的.
那些密碼的信息內容可以寫.
一萬億個數碼.
一隻單細胞生物.
就是你一個細胞裡面.
都是有一萬億個數碼.
不是吧.
真是.
原來生命體是更加的複雜.
所以結果我們開始發覺.
我們信仰神論.
是一個符合科學的信仰.
因為科學的信仰.
基本上是相信宇宙是由上帝的話.
第一宇宙有起源.
第二宇宙是由母變成有.
第三宇宙的起源.
在第一剎那.
是安排好一切的規律的.
那麼如果有神論.
是符合科學和理性的信念.
因此毛神論有些重大的改變.

$^{401}$有一個叫Hort的堅定.
不相信宇宙有起源的.
但是結果他最後研究了.
就是宇宙的生物體.
如果它能夠合成.
偶然出現的機會率.
只有十的負四萬次方.
就是十之後四萬個零分之一的機會.
才會撞到生命出來.
就等於沒有機會.
所以它經常有些超越的智慧創造天地.
還有一個大哲學家Antony Flew.
我中學的時候已經看過他的東西.
他說他是毛神論者.
但是這個毛神論者.
他幾十年的毛神論.
就是他跟一個叫Habermas的.
基督教的護教士辯論過.
有沒有耶穌有沒有復活.
上帝存在.
之後他很多年後終於寫信給Habermas.
他說我是終於相信有神存在.
另外耶魯大學有個教授也研究.
就是問了六十個那些著名科學家.
寫了一本書叫Cosmobiophiles.
其中有二十四個諾貝爾獎得主.
結果全部都承認有神存在.
而有神存在.
有神創造天地.
宇宙為什麼這麼長.
一百五十億年.
誰知道有個MIT科學家叫Size of God.
根據愛因斯坦的理論.
就是說在速度 光速裡面.
那個時間會縮短.
如果我用光速這樣走.
我回來會年輕一點.
因為時間沒有長命.
如果你不明白就算了.
這些你明白就知道了.

$^{441}$就是說宇宙起源的時候.
其實那個速度很快.
如果用大爆炸的理論解釋.
就是目前的科學的模式.
其實是在宇宙的座標.
二十四小時等於地球.
計算的八十億年.
他有條數計出來.
八十億年.
第二天就四十億年.
第三天二十億年.
到第六天大概是一百五十億年多一點.
如果這樣計算.
跟今天宇宙的年齡差不多.
所以宇宙真的是七天的.
計算法不同的.
原來科學家 聖經也對.
所以相信上帝符合科學.
好了 我們說了宇宙的美之後.
接著還有一種我們美學的思考.
就是人間的苦難.
我們看戰爭裡面.
我們看見那些轟炸.
還有那些在戰亂中的時候.
左邊是烏克蘭.
左邊上面下面那些就是俄羅斯.
烏克蘭其實殺東部的人很多.
殺完司機.
不只是俄羅斯.
但是他們先打.
然後殺東部的人才是俄羅斯出兵的.
然後這個有個記者報導了.
後來西方不讓他講的.
另外一個就是右邊那幅照片.
巴勒斯坦在上方看著天津飛機炸下來.
下面就是他們沒飯吃的時候在哭.
中間下面那個好像是哥哥抱著妹妹.
飛機來了 走了走了.
這個就是一種令人覺得很恐怖的.
情況就是戰亂.

$^{481}$那麼我就想起這幅畫.
就是畢加索在他二十歲的時候畫的那幅畫.
就是一個小哥抱著一隻血的白鴿.
就是白鴿.
原來他後來白鴿成為一個和平象徵.
就是因為有一個血白鴿給他.
他就把他畫了.
因為他的孫子是很喜歡白鴿.
經常用白旗去引白鴿過來.
但是後來德軍打到來了.
當時畢加索住在法國.
因為西班牙本身有戰亂.
所以他住在法國.
住在法國 結果德軍照樣打過來.
打到來的時候.
這個時候的他爸爸保衛巴黎佔中就犧牲了.
那個小哥很痛心.
他就不再用白旗了.
白旗代表投降.
所以他用紅旗引白鴿過來.
德軍操入巴黎之後.
看到挂著紅旗表示抵抗.
就很憤怒的衝上來.
法國的小哥反對他們.
就把小哥扔下樓去死了.
但是白鴿他用刀刺 全部流血.
全部死了 剩下一隻還沒死.
那個伯伯就是那個爺爺.
死了兒子 死了孫子.
他就找畢加索請你畫了這隻白鴿出來.
後來他畫一隻.
到戰後1949年就畫這隻白鴿.
就成為第一次世界和平大會的封面.
但世界和平大會後來是蘇聯和東歐說.
美國和英國就不喜歡.
後來到再舉行世界和平大會.
總共還有兩次.
有兩幅你看見都是他畫的白鴿成為封面.
但是那時候西方是反對他們的.
他們就去華沙開會.

$^{521}$其中有一幅有中文字.
就是中國參加開了第三次世界和平會議.
那次和平會議是1952年.
和平會議就叫朝鮮不要再打仗.
就是韓戰停止.
還有不要到處駐美軍基地.
就引起世界的縱亂.
而那時候畢加索.
之前其實畫了一幅畫很出名的.
就是這個叫格爾尼卡這個地方.
我不知道西班牙話有沒有讀錯.
讀錯沒有人知道.
這個地方被納粹第一次轟炸.
就是世界第一次轟炸人民.
就是納粹在西班牙炸那個佛朗哥的敵人.
那邊是左翼的.
右翼就是佛朗哥.
所以死了很多人.
我們看這些照片就炸爛了那個城市.
所以這個畢加索畫了一幅很出名的畫.
就是格爾尼卡.
知不知道畫什麼.
畢加索通常不知道畫什麼.
他說任你解的你可以.
但是你看這幅畫你會發覺.
中間有一匹馬在那裡亂叫.
就是駕下來那些動物都跑了.
就是雞飛狗走.
那些馬也奔跑.
有一隻牛衝出來.
可以代表很多的暴亂出現.
但是你留意左邊有個媽咪抱著兒子.
仰天好哭.
右邊有個人舉起手.
在火燭的屋裡面慘叫.
地上有個躺在那裡的軍人手拿著.
一把劍斷了.
手還拿著一朵花.
上面有個人.
眼睛瞪來瞪去.

$^{561}$發生什麼事死了那麼多人.
下面有個人是.
瞪來瞪去發生什麼事.
頂有個眼睛好像有個上帝.
在觀察著世間的邪惡和腐爛.
這幅畫是很震撼的.
原來這幅畫是之前還有一幅他畫耶穌的.
耶穌釘十字架的畫.
這裡有幾幅他畫耶穌.
也有很特別的筆觸畫耶穌.
畢加索應該是他背景是天主教徒.
不過後來他是無神論者.
他說他自己是無神論者.
但是他在早期畫的耶穌.
他畫的耶穌這幅畫.
你看一下你明不明白.
他這幅畫他有一種很特別的震撼性.
我看的時候.
因為他左邊.
正中間是耶穌釘十字架.
但是耶穌的頭很小.
看不清楚他的樣子.
但是左邊是有個慘叫的人.
張開嘴巴慘叫.
後面的石頭似乎是復活之後.
開了出來的石頭.
然後中間的耶穌.
他的樣子是不清楚.
但是有一個慘叫的女人在他心口.
有人說這是莫大拉瑪利亞.
或者是聖母瑪利亞.
但是我看不是.
我說這是一個.
人類慘叫的符號在耶穌的心中.
耶穌他釘十字架又將正在慘叫的人類.
放在他心中.
他正在承受他們的苦難.
慘叫的.
我小時候見到我爸爸有精神病.
成為一個很兇惡的人.

$^{601}$經常罵我媽媽.
那時候到中學.
其實我心是慘叫的.
為什麼我爸爸變成這樣.
我本來很聰明又會寫詩.
很學懂的爸爸.
突然變成一個廢人.
媽媽一個人出去打工養大我們.
養到我兩個兒子都讀博士.
媽媽很艱苦.
所以當時我一直都覺得將來我結婚.
一定不可以欺負老婆.
一定要稱讚她.
稱讚到她暈為止.
所以你經常聽到我稱讚老婆的意思.
有個心理背景.
我覺得媽媽受他的傷很深.
人是慘叫的.
我媽媽更加慘叫.
她嫁給我爸爸的時候.
很帥很高大很有才.
但後來因為時代的扭曲.
我們中國人這個年代很悲慘.
所以我爸爸就有精神病.
他沒有去承受.
是我媽媽承受家庭的艱難養大我們.
我們承受這些苦難長大.
我們答之就是.
怎樣使這個世界變得不是那麼悲慘.
中國不要那麼悲慘.
這就是當時開始走上的路.
就是那個慘叫.
所以我看這幅畫.
一個慘叫的女人.
我是很深的震撼.
而下面就有些死了的腳.
死了的人聽十字架.
其他兩個人我說不是.
那個就是人類有死亡.
有痛苦.

$^{641}$慘叫有痛苦有些死了.
但你留意這幅畫的下面.
那個女人可能就是莫德納麗亞.
在洞裡面發現.
有一隻橙紅的甲.
就是白甲 甲刺.
拿了兩樣東西倒出來.
好像是你選生命還是死亡.
是喜樂還是苦難.
是罪惡還是潔淨.
好像基督復活下面.
開了一個新的可能性.
就是如果你選對的話.
上面很奇怪的.
是有一個笑的樣子.
左邊是慘叫.
右邊是笑.
因為復活的基督.
帶來一切新的開始.
所以這幅畫.
你可以離近一點看那個慘叫的樣子.
和對比另外那幅.
格雷尼卡的那幅畫.
是慘叫的媽媽死了兒子.
火中慘叫的人.
以及死在地上的那朵花和劍的軍人.
人類有很多這些悲慘.
接著我再研究回.
耶穌基督的美學表現.
左邊那幅是波蘭的.
一個格雷尼卡的畫家.
是在大教堂裡面畫的.
耶穌釘十字架.
搬十字架的過程.
是經過很多受苦的人.
那些人是.
你看有些是猶太人.
就是被抓了的猶太人.
有些跛腳.
很多個都是病苦的人.

$^{681}$後面那些拐杖.
就是很多病苦的人.
耶穌的受苦是經過這些受苦的人.
而右邊這幅是最新的AI畫.
就是耶穌在戰區.
是德國畫.
在戰亂的地區.
無論是烏克蘭也好.
在加薩也好.
在以色列也好.
在現在正在打的波斯.
就是伊朗.
那些鱷都炸爛了.
他抱著一個Baby.
耶穌抱著他.
在那些受苦的人當中.
而這一幅.
就是波蘭的畫家畫到.
耶穌當時是帶著血淚.
走在女人.
當中女人有個婆婆.
在他十字架面前流淚.
後面有個小哥哥在流淚.
再後面有些目無表情的女人.
就是女性受很多痛苦.
這個世間.
是耶穌知道那些病苦的人.
受苦的女性.
他通過他的呼喊告訴你.
我明白.
這裡耶穌經過那些死了的.
或者被傷害的小孩.
你留意耶穌跌在十字架上.
地上也有些死了的小孩.
有些還沒死的就有德軍在他旁邊.
戴著黑眼鏡.
蓋世太保在那裡.
按著那些小孩.
然後再後面有一條天梯.
那些死了的小孩正在上天堂.

$^{721}$這幾幅畫我看的時候都很震撼.
他在波蘭的.
Black Madonna的大教堂那裡.
我當時進去看.
是看見.
他那個畫家嘗試表達.
耶穌基督是和我們一起.
正在經歷我們經歷的痛苦.
但是.
耶穌基督的畫裡面有另外一個就是.
他復活.
左邊文藝復興畫.
耶穌榮耀他從復活墳墓出來.
那些羅馬兵丁全部掉下來.
就是世間的暴力.
權力是掉下來.
然後.
第二幅就是在Forest Lawn.
在美國洛杉磯.
一個很大巨大的畫.
就是耶穌復活的時候.
看到新天新地.
未來的世界.
基督回來的那個世界.
然後第三幅就是AI畫.
耶穌帶著年輕人.
小朋友.
前面有陽光.
我們可以再復活.
苦難會過去.
一切.
有上帝同在的時候.
我們就以無盡的愛來對這個世界.
使世間的苦難會改變.
減少世間的苦難.
所以這幅就是.
波蘭的畫家最後畫.
耶穌升天.
耶穌升天是帶著一種.
很深的憐憫和苦難的感覺.

$^{761}$這個不是很榮耀的畫法.
是耶穌充滿懷念.
對人間受苦的人.
對於很多.
是在戰亂中.
或者貧窮中.
或者種種無端端的衝突中.
或者婚姻衝突裡面.
很多的人的痛苦.
我今天升天.
但是我會一路紀念你們.
而畢加索有一幅畫.
就是1962年.
就是在那些爛了的.
碎了的武器裡面.
有白鴿飛起.
飛向太陽.
而聖經說到.
當彼得得到聖靈降臨.
跟那些人說話的時候.
那些人覺得很窒心.
他說我們應該怎樣走.
彼得說.
你個人要悔改.
奉耶穌基督的名受洗.
又叫你們聚得赦免.
我就必受聖靈.
去受到上帝所賜下的聖靈.
原來丹瑟士的聖火會帶來的.
就是悔改.
重新聖靈來臨.
我們的潔淨.
我們有新的生命.
我們就可以重新再前去.
我們再去洗更受苦的人.
當中告訴他們.
雖然我們在安定的地方.
我們沒有經歷你們的艱難.
但是我們明白.
我們知道今天我們跟你們一起.

$^{801}$經過這些的痛苦.
其實當中國開放的時期.
當時一個部長去見我們基督徒.
他們叫我先說.
團體叫我先說.
我就說.
本來我是當年有120萬年新的電台.
新的電台請我們後來我沒有接受.
因為我知道我賺這個錢.
我以後就搭business class.
我開靚車有司機.
那些錢在投資.
很多賺錢.
但是我終於沒有要.
因為我覺得中國很艱難.
我當時那些電台可能都要罵中國.
才得到人的歡迎.
所以我寧願不罵中國.
我不要那些錢.
我只是回來跟你.
中國艱難的群眾.
我和你一起走.
我說耶穌基督一生與我同行.
我亦選擇與中國人同行.
當時我見那個部長其實眼有淚水.
後來我們全體向那些中國高級官員.
吃飯的時候唱.
(唱).
他們覺得很奇特.
我們基督徒有一種這麼喜樂的生命.
但是帶著無窮的愛.
這幅圖我們看見是另一種美學表現.
就是當飛機轟炸的時期.
無論是巴勒斯坦也好.
烏克蘭也好.
還是俄羅斯.
還是以色列.
還是伊朗.
那些歌聽到飛機的聲音.
已經很害怕.

$^{841}$怕炸彈掉下來.
那些歌害怕飛機.
我女兒小時候.
我在中神讀書的時候.
我拖她出去公園走.
她經常走啊走啊.
那些歌你就追得連武器都追不到她.
但是忽然間她走回頭很害怕.
然後剛剛用飛機飛過去.
來啟德機場.
那時候啟德機場就馬上跑回頭.
馬上抱著我們.
小朋友怕飛機的.
那些聲音那麼大.
但是那個飛機來的時候會炸.
他們看到的情況是怎樣.
所以這幅圖很震撼.
左上角的那幅圖.
中間那幅圖就已經沒有感覺了.
小朋友在廢墟裡面坐著不知道在做什麼.
右邊就是以色列都被那些黑馬斯殺了的小朋友.
左邊就是巴勒斯坦被殺的小朋友.
下面那些小朋友很害怕.
哥哥兒子都很小.
他不害怕嗎?.
但是他也要安慰小朋友.
所以我突然間好在找到幾幅畫.
耶穌抱著他們.
誰會痛愛去擁抱.
廢墟中的孩子.
耶穌騎著路進去.
耶穌就殺了他.
那幅圖就是耶穌今天騎著路進去的戰亂區.
有一個什麼樣的經驗.
今天我們人類做了什麼.
而有一幅畫是在伯利恆的教堂裡面.
那個牧師他說.
聖誕節如果耶穌今天出生.
祂的生命.
祂會在哪裡出生呢?.

$^{881}$應該是在加沙出生的廢墟那裡.
這幅畫就是另一個美學.
在現代戰爭轟炸下.
飛彈射擊下的廢墟.
以色列也炸到爛了很多.
被伊朗炸回頭.
在廢墟中.
耶穌在那裡出生.
所以在這樣的情況下.
怎麼說普世寬騰呢?.
我們看見這些是苦難.
所以他在教堂陳設了.
在很多的磚頭當中.
拜了耶穌.
以前在馬槽出生.
今天在廢墟中出生的基督.
他這個藝術表達的美學.
就是表達今天耶穌.
仍然是以無窮的慈愛和憐憫.
是跟我們共同經歷苦難.
正如聖經所說.
他們在一切的苦難中.
他是跟他們同受苦難的.
這個是舊約寓言.
未來耶穌是基督跟他們同受苦難.
是他以慈愛和憐憫救贖他們.
所以這幅圖.
這幾幅圖表達一個十字架裡面.
有一個清新世界.
外面是一個殘破戰亂的世界.
結果我們做了什麼呢?.
其實我們看見人間苦難.
都不是因為今天見的.
我在三十年前已經去到中國.
在農村見到留守兒童.
我們在中國做了三十年.
接近三十年是幫助留守兒童讀書.
而他們寫封信給我.
他說如果我是一隻受傷的鳥兒.
你們就是白衣天使.

$^{921}$在你們的幫助之下.
我們的傷口得到癒合.
我正在茁壯成長.
翱翔藍天的一次機會.
你們幫我去開創一番天地.
我好受感動.
因為你們是無私的奉獻不求回報.
很多年後我們再探訪這些小朋友.
我們有北京讀大學的學生.
到他請我吃飯.
那些以前經常哭的女孩.
現在都變成很漂亮的女孩.
每個都讀完碩士出來打工.
每個都很有成就.
右邊那幅圖的男孩.
你覺得好像傻傻的男孩.
其實他已經讀到博士.
在北京大學讀博士.
他博士令文是外國出了.
已經成為出色科學家.
不過越出色那些樣子越傻.
如果那些學生這麼開心.
他還寫一張卡牌.
是多謝我們當年對他的扶持.
讓他有今天的發展.
接著我去清華大學探訪一個.
老人家餘學的很大的思想家.
他是跟馮友蘭.
他臨尾到九十多歲.
終於決志接受耶穌.
另外有一個我自己的老師.
成宗英教授.
是世界首席儒家的專家和大師.
死之前他跟我討論基督教.
討論了三十多年.
不斷跟他解釋基督教.
和中國儒家的關係.
中國文化是可以接受基督教.
以及是有實質的存在.
他最後都同意.

$^{961}$到他死之前.
他說下一次見你.
我再討論罪的問題.
結果是沒有下一次了.
他就死了.
但他死之前.
叫牧師的牧師是我的好朋友.
成教授請我跟他洗禮.
為什麼呢.
儒家多數不肯信耶穌.
他要尊敬耶穌.
但是不信耶穌.
但他都肯信.
因為他想過罪的問題.
因為有神就有罪.
罪的問題.
你要得到永生.
你要潔淨罪.
所以他在這個時候.
終於接受了基督教.
作為儒家同時可以成為基督徒.
以及我們不單止跟學者接觸.
我們也去到農村這些屋.
你看這間屋多爛.
要爬一條梯子上去.
所以爬一條梯子.
每個人都要拍一張照片.
這麼特別的.
爬上一間屋要爬梯子.
由裡面那些屋.
你看見那些的.
那個守江.
我們去探望他.
他那間屋沒有牆.
那面牆是靠樹枝插在那裡釘住.
是沒有水泥的.
所以冬天都很冷.
我們去幫這個守江讀大學.
接著進大學.
以及有些官員一起見到我們.

$^{1001}$他就去幫他蓋一間屋.
然後這個.
守江現在大學畢業了.
所以他們有些寫信給我們.
我們感覺最特別就是.
陌生人可以愛我們.
他經常說我不明白.
為什麼你不認識我們會愛我們呢.
你改變了我.
讓我不再迷惘.
不再坑.
人間有愛.
而我們做這些的事情.
我們都猜不到.
他們從陌生人的愛裡面感覺到.
原來陌生人可以愛他.
然後我們將他們的困難.
也向中央政府反映.
所以我也見過幾個中央領導人.
我也有見過.
我們完結.
最後有一個我們學生.
叫蕭慈柳.
寫了一首詩給我們.
這首詩是十年前的.
她今天已經變成四十歲.
已經是很有成就的女人.
她這樣說 什麼愛是厚重而深沉.
什麼愛是細長而深遠.
什麼愛是如縱水如泰山.
什麼愛如天空如大海.
其實民間.
我們機構民間.
根深蒂固的叔叔阿姨.
就是給我們這個大愛.
就是這樣的.
她說從那時候開始.
他們就撐起一個個年少的夢想.
從當時開始.
我們叫她做叔叔阿姨.

$^{1041}$從那個時候開始.
我們每一個年都在這裡成長.
從那個時候開始.
我們是一個一個.
她永遠繼續愛我們.
最後她說.
直到現在.
我們開始長大.
我們已經朝著夢想前進.
直到現在我們已經被別人叫叔叔阿姨.
直到現在我們可以上大學.
我們成家立業.
我們工作.
直到現在他們仍然這麼愛我們.
陪伴我們.
所以愛可以改變世界.
我們未必叫到中東和平.
但我們每個人付出的一切.
會使世界變得美麗一點.
願意我們每個人都努力.
是經歷從上帝來的愛.
最後我作一個呼喚.
如果在座有哪一位.
你今天覺得想接受到.
這種天地的親情和慈愛.
有哪一位是你希望經歷到.
那種內在的改變.
我們內心或者有很多慘叫.
有很多過去的苦難.
但因為耶穌基督釘十字架.
承擔的是我們.
天下之美.
就是說我生命可以轉化.
好像莊子所說.
我可以轉化.
由一條魚變成一隻太虛之上的飛鷹.
我們今天也是.
我們可以轉化.
從苦難的過去得到安慰和意志.
從過去很多的罪裡面得寬恕.

$^{1081}$而愛對人.
如果你這個時刻有感動.
我作一個祈禱.
我邀請你和我一起祈禱.
如果你是基督徒.
都沒有怎樣經歷神的.
都可以和我一起祈禱.
我們一起祈禱.
嘗試讓我們進入那個美的空間.
從神聖而來的美的空間.
我們祈禱.
親去天父.
在這個時刻.
我來到你面前.
帶著我的過去.
我過去日子有很多辛酸苦楚.
有些放不開的痛楚.
苦難.
都求你.
依持我.
求你安慰我的心.
我知道我也有很多做錯的.
我有很多罪過.
我心裡面也不是完全善良.
我有自私.
有時也會貪婪.
有情慾.
我有很多不同的罪過.
求你潔淨我.
求你寬恕我.
求你使我轉化成一個新的生命.
讓我進入你那個美的空間.
美的世界裡面.
讓我們經歷到你 主.
我現在願意打開我的心.
接受耶穌基督成為我的救主.
讓我過去.
我的苦罪是去除.
又和基督一起復活.
有新的生命.

$^{1121}$讓我對這個充滿艱苦的人間.
走出愛心.
使這個世界能夠因著愛變得美好.
很多人因著愛可以有希望.
主今日我做這個決定.
我願意一生跟從你.
求你賜給我信心.
無論怎樣艱難.
我都有信心.
無論遇到什麼困所.
我都有盼望.
無論遇到什麼人.
我都學會先去愛他.
將衝突化解為愛和寬恕.
耶穌基督.
天父.
你接納了我.
你讓我得救.
禱告奉耶穌基督.
永祈求.
阿們.
\newpage



\section{哥林多前書 1:26-31}
\label{sec:031874}
\textbf{小人物,大使命}~~~~梁錫強醫生
\newline
\newline
連結: \href{http://www.church.com.hk/acms/content.asp?site=cdc\&op=show\&type=product\&code=031874\&layout=sermon}{\texttt{031874}} ~~~~ 語音日期: 2026-03-29
\newline
\newline
\hyperref[sec:031873]{\small{< < < PREV SERMON < < <}}
~
\hyperref[sec:index_chronic]{\small{[返順時目]}}
~
\hyperref[sec:031875]{\small{> > > NEXT SERMON > > >}}
\newline
\newline
哥林多前書 1:26-31
\newline
\begin{longtable}{cl}
\hline
\hline
章節 & 經文 (和合本修訂版)\\
\hline
1:26 & \begin{tabularx}{0.7\textwidth}{X} 弟兄們哪，想一想你們的蒙召，按著人的觀點，有智慧的不多，有能力的不多，有尊貴地位的也不多。 \end{tabularx} \\ \\ \relax
1:27 & \begin{tabularx}{0.7\textwidth}{X} 但是，神揀選了世上愚拙的，為了使有智慧的羞愧；又揀選了世上軟弱的，為了使強壯的羞愧。 \end{tabularx} \\ \\ \relax
1:28 & \begin{tabularx}{0.7\textwidth}{X} 神也揀選了世上卑賤的，被人厭惡的，以及那一無所有的，為要廢掉那樣樣都有的， \end{tabularx} \\ \\ \relax
1:29 & \begin{tabularx}{0.7\textwidth}{X} 使凡血肉之軀的，在神面前，一個也不能自誇。 \end{tabularx} \\ \\ \relax
1:30 & \begin{tabularx}{0.7\textwidth}{X} 但你們得以在基督耶穌裡是本乎神，他使基督成為我們的智慧，成為公義、聖潔、救贖。 \end{tabularx} \\ \\ \relax
1:31 & \begin{tabularx}{0.7\textwidth}{X} 如經上所記：「要誇耀的，該誇耀主。」 \end{tabularx} \\ \\
[1ex]
\hline
\hline
\end{longtable}
$^{1}$弟兄姊妹主內平安.
我今天犯了兩次一個低級的錯誤.
即是講完的低級錯誤.
就是剛才唱歌唱得太大聲.
變了現在就沙聲了.
那就是非常好的敬拜.
那就是激勵了.
那些人說你不要唱太大聲 要講到.
我忍不住的 知不知道.
因為高舉十架 榮耀神.
很難又收回我的聲音.
我的聲音已經很大了.
特別見到大雄好像上了身.
退後兩步就知道他的高音來了.
就知道我跟著他去瘋.
他就忘記了E是我極限.
幾乎叫破喉嚨那個唯獨主寶血那個.
好 那我們就非常之好的敬拜.
將我帶到上帝那裡.
那你就在我們記得應該是.
大約兩年前差不多的時候.
我開始想講舊約的聖經.
其實初初不是這樣想的.
我打算講一下《創世記》.
也可能講一下《創造》.
因為我對這方面也有些興趣.
但不可以說研究的.
因為大家知道我讀理科的.
那我就拋下書包.
就想著可以握一頓飯吃.
跟你講一下讓你不明白就行了.
誰知道完了之後.
姐妹跟我講什麼多元氛圍.
又什麼上帝粒子.
又有很多維度.
我說原來真是他們聽的.
聽得明白的 問回我的.
那我打算一講開舊約.
上個人我覺得.
其實真的我以前從.

$^{41}$以前我的角度我從小回教會回到大.
我看舊約是好像一個故事那樣.
也有很多的故事.
很多年月日 地名 人名.
我覺得似乎是挺有興趣的.
但是好像跟我沒什麼直接的關係.
當我去準備講道的時候.
每次上帝給我亮光.
就是原來不是沒有關係的.
是帶有的關係.
那我一直這樣講下去.
那我就試圖用一個所謂.
宏觀去看舊約.
我們不是說著重每一節.
不是著重每一個年代.
或者每一個人物.
每一個地點.
而是著重上帝究竟是跟我們講什麼話.
還有怎樣跟新約去接軌.
integrate(融入)的.
發覺原來真是很厲害.
上帝真是在舊約的時候.
已經好像希伯來書十章所講的.
舊約的法律 律法是.
我們今天的影子.
是一個預表.
是一個shadow(陰影).
我們原來很多東西.
在舊約那裡已經是預表了.
只不過是今天我們的影子.
因為耶穌基督還沒來之前.
但是 祂所做的事.
跟我們今天所做的.
新約時代或者今次今天所做的事是一樣.
祂所犯的錯.
今天我們仍然在犯錯.
舊約告訴我們.
在人最無助的時候.
就是看到上帝的大能.
今天同樣我們都是一樣.

$^{81}$舊約告訴我們人不斷的犯罪.
但是神有不斷的憐憫和救贖.
今天也是一樣.
特別是今天我們在一個很混亂.
甚至一個很危險的世界那裡.
大家如果有看新聞的.
除非你剛剛才參加完.
這個詞我們會明白.
維斯里旅行團回來.
剛剛從太空回來.
否則你知道現在是很多戰亂.
中東就烽火連天.
可以說是打得很厲害.
之前有俄烏戰爭.
仍然繼續有以巴的衝突.
這個已經是幾千年了.
我們實在是一個很危險的時間.
剛才聽到新聞說.
美國可能會退出NATO.
北大西洋工作組織.
這個組織是第二次世界大戰成立之後成立的.
是維護了歐洲.
包括英國一個安全了幾十年.
這個會有一個很大的轉變.
似乎這個世界是走去極端.
已經沒有中間了.
中間就是沒有立場.
看之前選舉 曼城那裡.
傳統的兩大黨都是大敗.
結果是最高票數的兩黨都是極端派.
極左和極右.
最近在英文班裡面.
我們有時候都會聊聊天.
講一下政治.
我們就說一下.
他說快要投票了.
老師你怎麼看.
即將到5月有個地區的投票.
可能見過老華僑.
就熟悉英國的行情多一些.

$^{121}$我說我也不知道.
因為現在是走去極端.
這個世界.
他說那你選左邊還是右邊呢.
好像很簡單.
不是左就是右.
左邊就是有點仇富.
你有錢他就拿了你.
就是劫富濟貧.
他給你很多福利.
養了很多懶人.
他就給你一個大愛包容.
未必是我們傳統.
就是我們很勤力的香港人喜歡.
右呢.
右極右很歧視外國人的.
他們是民族主義.
很保障自己人.
所以我說現在你是香港人.
你來到這裡.
你又帶了一些錢來.
你就沒有運行了.
選哪一邊都.
是不是.
老實講真是的.
告訴他甚麼.
告訴他這個世界是這樣.
趨勢是這樣.
現在世界的地緣政治是趨向.
可能會遲些真的好像回到三國一樣.
三國鼎立.
又或者將來可能去到啟示錄所講的.
到最後有一個大的極大的權勢.
是可以控制所有的人民的.
這個是聖經所預言.
我不知道.
有一樣東西我知道的就是人是.
雖然現在科技是日新月異.
之前是電腦.
接著又資訊發達.

$^{161}$現在又AI.
人的確在科技上是進步了很多.
但有一樣東西是沒有進步的.
就是我們的思維.
就是我們的道德的觀念.
似乎是這麼多年.
幾千年沒有進化過.
今天我們仍然活在罪裡面.
我選了一段聖經.
剛才主席提醒我.
你知不知道你上個.
去年不知何時何月何日.
他很熟悉.
你已經讀過這一卷經文.
我說我可能自己沒有進化過.
再用回那個經文.
這個經文是給我一個很大的感動.
就是上帝所揀選的是什麼人呢?.
這個也是善施的一個很重要的題材.
之前也說過.
善施是什麼時代呢?.
善施就是他們去到迦南地.
約書亞記記錄到他們已經打贏了幾場仗.
那些地成立了一個地方.
但是仍然不斷受到外族人所.
即迦南不是外族.
是迦南地的人去排斥他們.
本來是沒有問題的.
但是他們開始離開了上帝.
所以《士斯記》一二章開始.
耶和華要興起一些士司.
因為約書亞離開世界了.
所謂士司就是他們的領袖.
當時未有皇帝.
但是聖經說得很清楚.
「士司在世的時候.
耶和華拯救他們脫離受敵的手.
耶和華因為他們受欺壓迫害.
所發出的哀聲就憐憫他們.
但士司一死.

$^{201}$他們又轉去行惡.
比他們祖宗更壞.
去隨從別臣」.
他們去拜其他的臣.
去拜巴力.
事奉叩拜他們.
不放棄他們的惡習和頑硬的行為.
就是每一次上帝興起一個士司.
教導他們.
他們就轉好一點.
當那個士司離開了之後.
他們又變回壞.
結果上帝又將他們交給當地的人.
結果一段時間他們又大聲喊.
大聲喊救命.
上帝憐憫他們.
上帝又興起一個士司.
這樣持續了幾百年.
人的劣性是很難改的.
很劣性很難改.
有人看著你就乖.
沒有人看著你就頑皮.
做老師的都知道.
我小時候就是.
老師看著我們我們很乖.
老師轉身寫黑板的時候.
我們就作反.
就是這樣.
不過今天不知道還有沒有課室.
還有沒有黑板這件事.
可能都沒有了.
現在每個人拿著iPad上課.
小孩子也是.
父母叫你不要.
提醒了你很多次.
OK OK以後乖.
轉眼又忘記了父母的話.
我們人就是這樣.
所以士司記提醒我們.
我們是這樣的人.

$^{241}$我們會不斷犯罪.
但是上帝的拯救.
祂的憐憫是不斷.
祂不斷施恩給我們.
直到今天.
我們都是不斷的犯罪.
上帝不斷的憐憫.
士司記就告訴我們.
上帝所選擇的人是什麼人呢.
之前說了兩個大家熟悉的人.
一個是基奠.
一個是參宣.
大家一說士司記就想起這兩個人物.
的確這兩個人物.
我們想起的.
我們小時候讀上主日學的.
想起他們真的很厲害.
基奠帶著很少人.
就可以打敗那些米典人.
參宣更加厲害.
兩手一推的柱子.
就全部非利士人在廟裡面的人.
全部壓死了.
是他們的英雄.
但是不是的.
聖經不是這樣說的.
大家研究一下士司記.
主日學可能沒有重點.
沒有說這件事.
基奠是一個沒有膽量的鬼.
他是在大後方.
他不敢走去戰場的前方.
是上帝的使者找他出來的.
他不敢的.
他要試探神.
上帝叫他去查巴黎的廟.
他要晚上半夜去查.
他的膽量是不夠的.
信心是不足的.
但是後來成為了以色列的勇士.

$^{281}$參宣更加不堪.
肉體的情慾.
眼目的情慾.
今生的驕傲.
失敗了三樣.
參宣是什麼呢.
參宣是一個嫖客.
見到妓女他就喜歡和她在一起.
不是英雄.
結果最後他衰在女人手上.
大家上次聽到就說了.
但最後上帝仍然繼續使用他.
他悔改.
成為以色列的英雄.
他死的時候.
殺死的菲律賓人多過他生前的時候.
力量從哪裡來的.
從上帝而來.
為什麼會揀選他們.
不是因為他們是英雄.
是一個勇士.
而是上帝認為他們是可以做到英雄.
可以做到勇士.
這正正是我之前所說的.
顛覆了我們今天的思維.
今天我們的講道也是.
小人物居然有大的使命.
我經常都說.
大家可能熟悉小成本大製作.
經常都這樣說.
我經常覺得我們教會是小成本大製作.
每一件事都有神的恩典是充充足足的.
上帝正正就是要揀選一些.
我們不為意的小人物.
《士司記》第四章就很明顯了.
有兩個人物出現.
相信大家如果不是很熟悉《士司記》.
或者不是讀神學.
你是不會很留意這兩個人物的.
第一個就是底波拉.

$^{321}$聽過嗎?.
底波拉我聽過.
底波拉我沒有聽過.
或者很少人聽過.
我又考驗你們.
有沒有人聽過底波拉.
剛才完了之後.
陳全道給我看聖經.
他有些東西教導我.
陳全道非常之好.
他本聖經《士司記》第四章.
他寫滿了筆記.
不是我剛才說到他寫的.
原來他真的看聖經很用心.
底波拉是一個女的先知.
也是當時做了以色列人的士司.
可能有些人不知道.
原來士司有女人可以做.
有.
很奇怪.
上帝揀選她做一個以色列人的領袖.
當時女人的地位有多少呢?.
我真的不知道.
但居然上帝揀選成為以色列人的士司.
而且她是一個先知.
聖經說她做過些什麼呢?.
她說在以法蓮山區.
山區不是城市.
底波拉在一個棕樹之下.
做什麼呢?.
是聽以色列人上來做審判.
她做了一個和事佬.
她做訴訟的.
有些人說誰呢?.
隔壁陳師奶偷了一瓶醬油.
她就去做什麼呢?.
她就做裁判 排解民間的糾紛.
當然上帝給她智慧.
人們是服她的.
但她肯定不是什麼呢?.

$^{361}$肯定不是一個軍事的領袖.
打仗不會有經驗.
肯定不是當時政治的人物.
肯定她家裡不是高官的子女.
肯定不是當時所謂社會的精英.
她只不過是一個很謙虛的.
在一個棕樹之下聽訴訟的婦人.
但居然上帝揀選她成為以色列的侍奴.
她做了什麼呢?.
當時以色列人又走去犯罪.
上帝把她交給了迦南地的一個王爺賓.
有幾十年的時間.
他們當然是在被人統治之下.
肯定不會有好的生活.
就喊救命 上帝憐憫他們.
上帝就和底波拉說.
叫她做些什麼呢?.
不是叫她去打仗.
而是叫她去找人去打仗.
找了一個叫做巴拉去打仗.
那些人名我們不需要太過記住.
大家有時間可以去看《士司紀》四章五章.
結果巴拉帶了一萬大軍.
去和耶賓軍隊的元帥西西拉打仗.
結果是打贏仗.
結果是全軍覆沒.
西西拉也走了佬.
這個元帥全軍覆沒.
能力從哪裡來的呢?.
上帝而來 能力不是從這個士司底波拉.
也不是從以色列的軍隊元帥巴拉.
能力是從上帝而來.
上帝藉著底波拉這個士司.
這個以色列的領袖去和巴拉說.
上帝使用我們有時就是這樣.
他不是直接用我們.
他不是說你去戰場.
不是 順便有他計劃鋪排.
但最重要的是底波拉是願意信服.
願意聽從上帝的說話.

$^{401}$他是一個什麼人呢?.
是一個很普通的婦人.
接著還有一個女人叫做瓦益.
不知道大家有沒有聽過.
這個就不是士司.
但當時她在屋裡面.
當對方的元帥西西拉逃亡的時候.
就逃到這個村那裡.
她就要找地方真的去歇一歇.
喝一口水 休息一下.
去到瓦益這個家裡.
瓦益做什麼的呢?.
她是做那些帳篷.
即是大家的童軍去露營.
就有一些帳幕.
然後有一些釘 釘著那些營.
不是現在這樣.
即是一開 自己彈起.
吹完氣.
以前是真的要逐個釘下去.
她是做那些事情.
這個西西拉就倒楣了.
千挑萬選.
沒有選擇這間屋企.
結果她進去休息的時候.
這個瓦益就用她的一口釘.
釘在她的太陽穴那裡.
就直接殺死了這個西西拉.
上帝居然是使用一個手無搏擊之力的.
一個名不經典的人.
一個婦人.
一個住在村裡面.
做一些很普通手藝的人.
一個很普通的婦人.
沒有名 無位 無權的.
將這個戰役最關鍵的一刻.
就是將這個元帥殺死的一刻.
交了給她.
我們不明白上帝的意思.
搞這麼多事情幹什麼呢?.

$^{441}$上帝有祂的意思.
上帝告訴我們.
這些寂寂無名的人物.
上帝有重擔給他們.
有一個使命給他們要去做.
只要他們聽從上帝的說話.
只要他們忠心事奉.
今天上帝也在跟我們說這些說話.
你看看新約.
真的有很多小人物大使命.
之前我也有說過.
瑪利亞 耶穌的母親.
真的是一個小人物.
上帝揀選她.
成為耶穌的母親.
幾個門徒都是耶穌所揀選的.
頭幾個都是打魚的.
彼得 安德烈 約翰 雅各.
是一個普通的行業.
當時是以漁業為重的地方.
真是漁夫相信有成千上萬.
全部都是小人物.
上帝揀選這四個.
也就是耶穌揀選這四個.
耶穌揀選瑞利人馬太.
上次也說過了.
沒有人喜歡瑞利人馬太.
第一 聽到收稅人已經不開胃了.
沒有人喜歡.
今天我們也是.
第二.
他是幫羅馬人做事收錢的.
欺壓猶太人.
他們是二五仔.
他們宣誓了效忠羅馬政權.
那些人不喜歡.
第三就是當時很腐敗.
很多貪污.
今天也是.
有人的地方就有貪污腐敗.

$^{481}$我們都很清楚.
收10元就代了2元.
所以他們是非常富有.
富戶來的.
大家也見過了.
之後那個蟻仔殺雞.
都是富戶來的.
上帝揀選他們.
這些小人物.
上帝揀選有罪的人.
所以今天不要告訴上帝.
我不行.
上帝說你行.
能力從上帝而來.
不要告訴上帝我有罪.
上帝說不要緊.
你有罪 你悔改.
我會赦免你.
我會使用你.
就是這麼簡單.
上帝很多的故事.
上帝的角度.
都是由小人物開始的.
不是由精英一族的開始.
上帝揀選的不是完美的人.
而是願意為上帝成為完美的人.
願意讓上帝去塗造你的生命.
這就是小人物大使命.
明白了沒有?.
這就是不斷地在那個事實上提醒我們.
以前的人打仗.
不是靠人多人少.
而是靠上帝的大能.
上帝揀選的不是世人所見到的.
所謂有多厲害多威風.
而是上帝認為你是合適的人.
所以為什麼我去年又說.
《廣南多條書》一章.
27次那裡開始.
26次那裡說.

$^{521}$今年又說.
因為這個是我很受感動的.
因為上帝今天正正就是揀選我們.
去成為祂的合用器皿.
因為當我們做牧者.
很多時候去找人在教會幫忙.
很多時候第一句就是我不懂.
第二就是我沒有時間.
其實不是.
時間是上帝給你的.
你不懂.
上帝會教導你懂.
所以這就是小人物大使命的意思.
剛才看了《廣南多條書》那段聖經.
《廣南多條書》是什麼教會呢?.
《廣南多條書》是一個非常繁榮的地方.
就是在外邦人的地方 希臘的地方.
它是一個海港.
其實它是很繁榮的.
雖然雅典是首都.
當時哥林多其實是一個商業的城市.
就好像北京是首都.
但是上海是一個繁榮的地方.
大家明白那個原因.
所以哥林多是很重要.
所以保羅在哥林多建立教會的時候.
知道那裡是很重要的地方.
但是當你一個繁榮的地方.
人多了.
自然怎麼樣?.
人多口雜.
自然有分歧.
自然有不聽話.
我們自己教會也是.
我們教會三十多年.
開頭每個人都很聽話.
我不是說你們不聽話.
為什麼呢?.
因為兩三家人你怎麼會不聽話呢?.
其實很容易.

$^{561}$教會大了.
自然有不同的聲音.
不同的意見.
一定有的.
這個是正常的.
正常.
其實是福氣來的.
有不同意見 有不同聲音.
但是我們千萬不要中了魔鬼的詭計.
哥林多教會中了魔鬼的詭計.
他們開始分黨.
他們抱著名牌.
他們有些人說我熟悉保羅.
保羅很厲害.
保羅這麼遠來.
保羅是我們的元老.
是我們開始的創始人.
是他開始了教會.
我熟悉保羅黨.
我說不是.
那個是埃及亞歷山泰尼的阿波羅神學家.
說話很厲害.
他神學知識.
保羅是很普通的.
我熟悉阿波羅.
有些人說不是 你搞錯了.
我熟悉彼得.
彼得就是教會的基石.
耶穌親自叫他看著教會.
陰間的鎖匙給了他.
我埋堆當然是埋一個最厲害的.
搞得好像黑社會一樣.
找大佬的找一個最厲害的.
他們真的有一點點這樣的情況出現.
結果就是分黨.
其實為什麼分黨呢.
是人的罪.
個人的驕傲.
阿波羅可能不是一個驕傲的人.
應該不是.

$^{601}$但是那些人捧到他很高.
這就是一個對大問題.
我見過很多教會.
有些牧者.
本來他自己沒什麼.
但是被人捧到他很高.
人的軟弱就來到.
這件事是很真的.
人的軟弱就來到.
所以保羅在哥林多前書第一章.
是罵他們的.
是提醒他們.
他說你們蒙召.
就是上帝召你們做事.
不是按著人的觀點.
不是你智慧有多少能力有多少.
不是的.
和世人眼光不同.
上次我說的.
幼稚園選學生.
都要看你的Portfolio.
是用人的眼光去看.
但上帝選人就不是用這些.
保羅說上帝選世上愚蠢的.
是為了令有智慧的人羞愧.
選擇軟弱的是為了強壯的人羞愧.
選擇世上最卑賤被人厭惡的.
就好像那些婦女.
孤兒寡婦.
那些有殘疾的.
那些收稅人的罪人.
為什麼.
令到其他自以為是的人羞愧.
上帝選擇一無所有的.
是要廢掉那些自己以為自己樣樣都有的.
為什麼要這樣做.
保羅提醒我們.
在上帝面前.
一個都不可以自誇.
今天你的能力是從上帝而來.

$^{641}$一定是的.
有些人可能自小.
他們有一個好的環境長大.
這個福氣是上帝給你的.
可能你小時候讀名校.
上帝給你有這樣的福氣.
這個能力不是你自己的.
一定不是的.
上帝選擇是軟弱的令到強壯的羞愧.
所以上帝不是因為你有多厲害才選擇你.
是上帝很厲害.
祂看中你.
你要將思維刀鑽.
是上帝厲害.
祂選了你出來.
不是你自己厲害.
不是說我做到了.
上帝你選我.
不是你做到.
是你靠著上帝你才做到.
明白了嗎.
所以這個希波拉也好.
這個瓦易也好.
這兩個女人.
她在《時事記》四章出現的時候.
她們不是因為她們之前有什麼業績.
不是她們有多厲害.
是上帝的眼光多厲害.
將她們選出來.
因為知道她們是願意信服.
她們聽話.
她們忠心事奉.
瓦易殺死袁帥.
耶賓手下的袁帥.
不是用刀.
不是用槍.
是一個露營紮.
露營的一口釘.
重點不是用什麼殺她.
重點是上帝的能力.

$^{681}$之前說的五餅二魚.
記得嗎?.
早幾個星期.
上帝使用的是什麼?.
是一個小朋友.
使當時幾千個成年人羞愧.
我不相信那些人沒有帶食物來.
但是他們不願意拿出來.
聖經強調是一個小朋友.
一個無名的小朋友.
我們不知道他什麼名字.
他拿了五餅和兩條魚出來.
重點也不是五餅二魚有多厲害.
而是因為上帝耶穌施犬的神蹟.
本來一個很小的成本.
就成為一個大的資源.
可以餵飽幾千人.
明白了嗎?.
這就是逆轉我們的思維.
我們雖然是小人物.
我們的資源不多.
但是憑著信心.
看到上帝的異象.
成就了上帝的大使命.
有一件事我不是誇口.
不值得誇的是.
我們教會有自己的座堂.
真的.
很多比我們大很多倍的華人教會.
在英國沒有自己的地方.
這就是因為之前的牧者看到異象.
我們的資源很少.
但是信心大.
當然是要上帝有供應,有預備.
我們才有自己的地方.
今天我們可以搞很多活動.
我們可以回來練歌.
這樣那樣.
完了之後不用撿起椅子離開.
有自己的地方.

$^{721}$沒有想過兩堂崇拜都坐得差不多滿.
之前買了一百張椅子.
想著很久都不會坐滿.
但是上帝為我們開路.
很多事情回頭看.
上帝有祂的意思.
為什麼2019年發生這麼多事.
上帝有祂最終終極的意思.
當時我們看不到.
現在我們看到了.
我們不要理會過去.
今天我們在教會.
我們要用心去事奉.
上帝給我們的恩典很大.
最重要是我們用忠心去事奉.
底波拉這個女仙姿,女士師.
她沒有上過戰場.
但是她可以做到時事.
她沒有住過皇宮.
她沒有讀過很多書.
她只不過是在榕樹頭下.
我們中國人叫做榕樹頭.
在那裡說話.
在那裡和人聊天.
但是上帝就是使用她.
雅逸也是.
在家裡不用出門.
上帝都可以使用她.
這個就是上帝使用人.
亦都是他們的福氣.
大家如果看到士師的第五章.
這個剛才陳全道提醒我.
這句聖經是非常好的.
就是五章二十四節.
整個五章都是女士師.
底波拉她的歌頌.
頌讚上帝的一篇.
其中她提到五章二十四節.
「願基尼尼人希伯的妻子雅逸.
比眾婦人多得福氣.

$^{761}$比帳篷中的婦人更蒙福氣」.
你為上帝做事.
你更加蒙福氣.
很重要的.
今天我們不需要地上的榮耀.
不需要甚麼金銀珠寶.
不需要別人怎樣稱讚我們.
這些不重要的.
但是雅逸的名字就寫在聖經裡面.
我們今天都記得她.
今天我們是小人物.
上帝要我們做大事命.
我們怎樣做?.
我們怎樣走出去做大事命?.
我們可能會這樣想.
之前劉宣教士來到提醒我們.
要做些事為上帝.
我們去宣教.
在這裡宣教兩個字.
我們就想起甚麼?.
想起不要太忙.
或者我遲些拿假期去短宣.
最好就是去泰國.
陽光海灘玩玩 宣教一舉兩得.
我們就想這些事.
或者我先退休 多些時間.
宣教就是將福音帶到群體那裡.
其實你每一天都可以宣教.
你的宣教工場在哪裡?.
你的宣教工場就是你的家.
家裡有人未信主.
宣教工場就是你上班的地方.
你每天都要上班.
你每天都要上學.
有沒有同學未信耶穌?.
是呀.
宣教工場就在我們教會那裡.
團契小組有些人是初信.
或者未信的.
我們可以查座 福音查座.

$^{801}$我們可以和他們去傳講福音.
宣教工場就是你臉書裡面的幾百個朋友.
最近你幾次和他們講過福音呢?.
其實工場就在這裡.
你不是真的要走去哪裡才做到宣教工作.
上帝的使命不是因為你是個小人物.
不出名.
就不會交給你.
上帝的大使命.
就是將福音普天下.
傳開使萬人作主的門徒.
這個是給我們每一個人.
這個不只是給教會的牧者.
這個很清楚告訴我們.
但是最重要的是.
我們要起來.
真是決心為主工作.
我們要有信心.
很多時候我們不信心.
我們要忠心.
我們要堅持.
很多時候不是一帆風順.
生命不是一帆風順.
大家都知道.
信耶穌又不是一帆風順.
我們知道的.
但是我們要依靠上帝幫助我們.
所以上帝今天不是問你有多少能力.
而上帝說.
你願意.
願不願意給我使用.
這個是最重要的.
我們不需要成為一個大人物.
才去被神用.
我們不需要去讀神學.
學士,碩士,博士.
才上帝使用我們.
上帝使用底波拉,瓦易.
五餅如魚的小孩.
一些漁夫.

$^{841}$他們受過什麼教育呢?.
他們做過什麼偉大的事呢?.
是沒有的.
是願意.
他們願意放下一切.
被上帝使用.
早期有一個安提柯教會.
大家可能讀《士徒行傳》都聽過.
這個教會有什麼人呢?.
沒有什麼人知道.
聖經沒有怎麼提及.
裡面有什麼人.
但是這是一個很重要的教會.
當然耶路撒冷教會是很重要.
因為耶路撒冷是耶穌基督釘十字架的地方.
而聖殿也在那裡.
我們也知道耶路撒教會.
有一些很重量級的人在那裡鎮守.
彼得也經常在那裡.
有耶穌的兄弟雅各.
是他的領袖.
但是安提柯教會很重要.
因為它是和外邦人宣教的.
第一個宣教的教會.
它是猜全宣教士出去傳福音.
保羅和巴拿巴第一個旅程.
都是在那裡開始.
在那裡開始才有所謂基督徒這個名字.
聖經說得很清楚.
安提柯教會就是我們現在打仗打到飛起的地方.
但是當時有一些無名的信徒.
是建立了這個教會.
他們一定是很有信心.
他們如果沒有信心的話.
就不會去搞宣教活動.
是不是?.
自己搞定了.
做多一個崇拜算了.
每個星期都在想宣教.
聖經沒有特別提醒他們.

$^{881}$但是那些無名的英雄.
這群小人物居然可以肩負這個大使命.
成為向西方地方第一個宣教的教會.
這個就是重點.
重點是什麼呢?.
重點是他們一群人.
是願意被上帝使用.
重點就是這群人一定是很團結的.
這個就是重點.
之前所說的無名英雄小人物.
可能都是自己.
但是如果看回一個群體一個教會.
他們能夠肩負大使命.
一定是有一個同心同步.
高舉基督十誡這件事.
這個是最重要的.
大家在事奉上能夠做到同心同步.
我覺得這是一個大的前提.
過往我就好像一個星探一樣.
知道什麼是星探嗎?.
不是精探.
我坐在那裡崇拜.
就聽到後面兩行有個姐妹.
聲音挺好的.
完了之後我就會再問.
聲音挺好的.
有沒有想參加我們敬拜組,敬拜隊.
試過有人回答我.
我只有一把聲.
我只有一把聲就夠了.
不是的.
他說只有一把聲我不懂得看譜.
還有我的拍子感很差.
我就加入了敬拜隊.
我只有一把聲是不夠.
但是看譜不是一定要.
不懂我們教你.
但最重要是什麼呢?.
最重要大家同心事奉.
最重要就是大家一起去.

$^{921}$在靈裡直接事奉.
我們不斷的成長.
這個是重點.
所以敬拜隊很多年了.
我參與現在是牧養敬拜隊的時候.
我強調一件事.
因為敬拜隊不是單單來到.
就是練習那兩個小時.
是 練習是很重要.
我們要做到最好的.
我們不要隨便.
我們要擺最好的給上帝.
這個當然是重要.
大家今天看到我們之前的敬拜.
為什麼會這麼齊心.
氣氛這麼好呢?.
是花了很多的時間.
詩琴妹妹跟我說.
她連電視劇都沒有追這兩個星期.
真的.
大雄跟我們練琴.
她真的有血有淚.
她帶著我們.
這個沒有誇張.
這個是對的.
但是如果你問我.
有些人說我有時候帶得很隨便.
或者這裡不是半拍咩.
我有時候我是忽略這些的.
因為我覺得不是百分之一拍走音.
或者拍子不正確是最重要.
最重要是大家要齊心去事奉.
在主裡靈裡相交.
我們才可以事奉.
任何事奉都要靈裡一插一.
才可以有效的.
所以敬拜隊就是一個敬拜的團契.
我們一起的.
不只是為練詩而言事.
是大家互相去幫助.

$^{961}$你有不足的地方.
我用愛心花很多時間.
大家互相提點去教導.
的確我們教會敬拜隊.
真的有很多很有恩賜的弟兄姊妹音樂商.
也都很願意將時間擺出來.
這就是群體教會.
上帝喜悅的團契的生活.
所以我們敬拜隊堅持有祈禱.
每個月有祈禱的時間.
我們有愛言食飯.
如果你可以的留下.
我們鼓勵.
為什麼呢?.
因為我們要大家在靈裡有合一和相交.
我們才有能力可以高舉基督的十字架.
我看這件事很重.
不是說我練了很多次.
大家入得很齊.
這些是好的.
但是入得齊.
是沒有心那麼齊的重要.
心的合一更加重要.
至於心的合一 上帝一定喜悅.
這個安提柯教會.
聖經不是說它很多.
它像是一個配角.
說是說保羅 說是說巴拿巴.
只不過是有一群人在這裡.
但是想一想.
這個教會不簡單的.
它是早期一個宣教的教會.
和外邦人的傳福音的一個窗口.
這群人一定是上帝喜悅.
一定是很合一.
一定是很聽從上帝的話.
很有意義的教會.
一定是.
我們忽略了安提柯教會.
我們後來都是說哥林多教會.

$^{1001}$說以弗所教會.
說羅馬教會 耶路撒冷教會就一定說了.
我們說這些所謂的大教會.
我看約翰的七封書信.
都是寫給亞細亞的幾間教會.
沒有提安提柯.
安提柯被人遺忘了.
但是它是一個很重要的教會.
一群小人物.
來一個大製作.
一個大的使命.
今天我們要為上帝做一些事情.
首先我們要聽從上帝.
我們願意信服.
放下我們個人的驕傲.
回到上帝那裡.
我們願意忠心.
即是遇到困難我們不會退縮.
很多時候在事奉路程當中.
我們見到有很多人遇到困難就退縮.
我在教會很多年.
在其他教會也有很多年.
我見過很多都是退縮.
意見不合 退縮.
你那些人不中聽.
我說我不和你一起了.
中了魔鬼的詭計.
或者很懶.
不舒服.
一句不舒服.
不喜歡就行了.
就像年輕人那樣不喜歡.
很多時候我們就有這樣的心態.
但是我們願不願意踏出一步呢?.
我們願不願意毀身呢?.
願不願意放下自己的自尊呢?.
我們一起去同心協會 一起去學習.
大家意見不同 我們一起去祈禱.
這些就是最重要的一環.
安提我的教會做得到.

$^{1041}$相信我的教會做得到.
教會人多了.
意見不同一定是有的.
我們教會也有.
這件事我很緊張.
因為我經歷過教會分裂.
不是我們這一間幸運.
之前我講過.
經歷過真的嚴重到.
木仔要下閘放狗趕人走.
經歷過真的要報警抓人.
咳水咳電.
好像黑社會.
不是開玩笑.
經歷過我們去到教會.
發覺原來木仔已經搬了東西走.
我們都不知道.
不在這裡聚會.
這些是很可惜.
也是很失見證的.
很多事情是由小開始.
我們要特別留意.
上帝提醒我們我們是一群小人物.
我們沒有個人的什麼能力.
只不過是上帝肯縱容我們.
我們今天成為一個群體.
是一個很大的福氣.
在孫子的地方有一個群體一起去事奉.
這個是非常之大的福氣.
上帝每一個使命.
都是由小人物開始.
大家看聖經.
從頭看到尾都是.
都是的.
保羅講得很對.
在前書一章.
他起用軟弱的有序的人.
就是因為他要令到自以為自己很強的人.
自以為自己很厲害的人.
是羞愧的.

$^{1081}$明白了嗎?.
是這麼簡單.
整本聖經都告訴我們.
我們是很不厲害.
我們需要上帝的救贖.
的的確確我們沒有進化過.
士師記的時代他們是這樣.
今天我們都是這樣.
好 我們有個低頭.
有個反思的時間先.
頂智慕我們很多時候.
我們會不會覺得自己很不夠資格.
或者不夠能力.
又或者我們不夠聰明.
又或者更加我們說我們自己不夠熟齡.
所以我們就不敢走出來去事奉.
又或者在事奉路途上.
我們遇到困難.
遇到人事的關係.
遇到一些爭執.
我們就退縮.
我們選擇離開.
大夫 我知道今天藉著士師記和我說的就是.
你不是找一個大人物.
你是找一個小人物.
你是找一個願意謙卑自己來到神帝面前.
被上帝塗灶的人.
這個是最重要的.
你讓我們在Swinner地方有這個教會.
我們有不同的團契.
我們弟兄姊妹在一起.
令裡頭有好的相交.
因為父我知道我們有力的事奉.
是要從零裡頭開始.
要零裡頭合一的開始.
父你提醒我們.
昔日安提柯教會一班無名的英雄.
也是因為他們的忠心的事奉.
他們的合一.
他們能夠看到上帝的異象.

$^{1121}$他成為福音傳和外邦人的窗口.
求主繼續去帶領我們.
提醒我們.
今天我們誇的不是誇自己的能力.
我們的日子如何.
我們的力量亦都是如何.
但是你的恩典夠多.
靠著你的幫助.
軟弱可以成為堅壯.
求主繼續帶領我們.
今日求主繼續帶領保守我們未來日子當中.
我們教會都是有主.
與我們同在.
讓我們能夠零裡頭合一.
讓我們不要膽怯.
我們真的要走出去.
我們學習今日所學習的兩位婦女.
她們不出名的.
但是婦女是使用她們.
你將重任交託給她們.
今日婦女也將重任交託給我們.
一切從我們的教會.
從我們的家庭.
我們工作.
我們學校開始.
將福音去講傳.
特別看到今日戰亂的世代當中.
我們更加要加快加倍地.
將福音去講傳.
更加要站得穩.
求主繼續帶領我們未來日子.
我們期待鳳姐的得勝名次祈求.
\newpage



\section{使徒行傳 16:6-15}
\label{sec:031875}
\textbf{神的美意 (華)}~~~~楊芹牧師
\newline
\newline
連結: \href{http://www.church.com.hk/acms/content.asp?site=cdc\&op=show\&type=product\&code=031875\&layout=sermon}{\texttt{031875}} ~~~~ 語音日期: 2026-03-29
\newline
\newline
\hyperref[sec:031874]{\small{< < < PREV SERMON < < <}}
~
\hyperref[sec:index_chronic]{\small{[返順時目]}}
~
\hyperref[sec:031876]{\small{> > > NEXT SERMON > > >}}
\newline
\newline
使徒行傳 16:6-15
\newline
\begin{longtable}{cl}
\hline
\hline
章節 & 經文 (和合本修訂版)\\
\hline
16:6 & \begin{tabularx}{0.7\textwidth}{X} 因為聖靈禁止他們在亞細亞講道，他們就經過弗呂家、加拉太一帶地方。 \end{tabularx} \\ \\ \relax
16:7 & \begin{tabularx}{0.7\textwidth}{X} 到了每西亞的邊界，他們想要往庇推尼去，耶穌的靈卻不許。 \end{tabularx} \\ \\ \relax
16:8 & \begin{tabularx}{0.7\textwidth}{X} 他們就越過每西亞，下特羅亞去。 \end{tabularx} \\ \\ \relax
16:9 & \begin{tabularx}{0.7\textwidth}{X} 夜間，有異象向保羅顯現。有一個馬其頓人站著求他說：「請你過來，到馬其頓來幫助我們！」 \end{tabularx} \\ \\ \relax
16:10 & \begin{tabularx}{0.7\textwidth}{X} 保羅既看見這異象，我們就立即設法往馬其頓去，認為神呼召我們傳福音給那裡的人。 \end{tabularx} \\ \\ \relax
16:11 & \begin{tabularx}{0.7\textwidth}{X} 我們從特羅亞開船，直行駛到撒摩特喇，第二天到了尼亞坡里； \end{tabularx} \\ \\ \relax
16:12 & \begin{tabularx}{0.7\textwidth}{X} 從那裡來到腓立比，就是馬其頓這一帶的一個重要城市，也是羅馬的駐防城。我們在這城裡住了幾天。 \end{tabularx} \\ \\ \relax
16:13 & \begin{tabularx}{0.7\textwidth}{X} 在安息日，我們出城門，到了河邊，知道那裡有一個禱告的地方，我們就坐下來對那些聚會的婦女講道。 \end{tabularx} \\ \\ \relax
16:14 & \begin{tabularx}{0.7\textwidth}{X} 有一個賣紫色布的婦人，名叫呂底亞，是推雅推喇城的人，素來敬拜神。她在聽著，主就開導她的心，使她留心聽保羅所講的話。 \end{tabularx} \\ \\ \relax
16:15 & \begin{tabularx}{0.7\textwidth}{X} 她和她一家都領了洗，就求我們說：「你們若以為我是真心信主的，請到我家裡來住。」於是她堅決請我們留下。 \end{tabularx} \\ \\
[1ex]
\hline
\hline
\end{longtable}
$^{1}$粵語堂的各位弟兄姊妹們 主內平安.
謝謝你們給我這個機會 讓我再次地來到粵語堂講道.
也求聖靈帶領我們講的和聽的都能夠同感異靈.
雖然有一些語言上的障礙.
我們也能夠完全無誤地聽懂神的信息.
我不知道各位有沒有這樣的一種心路歷程.
就是有的時候覺得自己的人生很迷茫.
覺得自己好像站在一個十字路口.
走了很多很多的路 似乎哪兒哪兒都不對.
然後你會覺得生命當中 上帝一直在給你關門.
我不知道這是不是你的生命的現狀.
可能真正相信上帝的人.
一定會在這個時候就問上帝 為什麼.
為什麼你讓我在人生中走了這麼多彎路.
有的時候是很短.
我們每一個人都會在生命的某一個階段.
陷入一種抑鬱也好 焦慮也好 迷茫也好.
有的人這個狀態很長 時間很長.
有的人甚至長達幾十年.
他回首他的一輩子兜兜轉轉.
好像都沒有走到一個正確的方向.
然後我們就會問說 上帝你是什麼意思呢.
我希望今天我們在分享這段經文.
講保羅第二次的宣教旅程.
這段經文的時候可以給我們一些小小的啟示.
我們的使徒行傳現在進入第十六章.
也就是保羅開始了他第二次的宣教旅程.
他第一次的宣教旅程主要是在小亞西亞.
也就是今天依然是在亞洲這一帶.
然後上帝突然給保羅一個帶領.
要讓保羅去到歐洲.
這就是一個第十六章的開頭.
但經文是這樣記載的.
經文說聖靈禁止他們在亞西亞講道.
於是他們就經過 我讓大家來看這幅地圖.
首先聖靈說禁止他們在亞西亞講道.
於是保羅他們沒有辦法.
他們禁止了 他們就開始走.
他們就往東北方向走.
他們經過佛里加.

$^{41}$接著又往北走經過加拉泰.
然後他們又往西繞了一點.
到了美西亞的邊界.
聖靈說他們就想往畢推尼去.
耶穌的靈卻不許.
他們就越過美西亞到了特羅亞.
這就是上帝的一個很奇妙的帶領.
什麼帶領呢.
讓人有點感覺不舒服.
甚至是一個冰冷和殘忍的事實.
禁止 關門.
走了好多的路 然後又不許.
你知道這一個禁止 一個不許.
保羅他們這麼多人 一個團隊.
走了多出八百公里的路.
有可能多費了好幾個月的時間.
為什麼.
我想大家如果要是保羅團隊中的一員.
如果你是希臘的話.
你是他的副手或者你是保羅的話.
你設身處地的想一想.
多麼不容易的旅程.
這讓人多困惑 多沮喪.
只等到保羅他們繞了一大圈.
就這樣白走了朝東北.
白走了一大圈的時候.
他們原本也許可以直接朝西北走.
他們朝東北走了一大圈.
只等到他們來到特洛亞.
神的第三步指示才領導.
說夜間有異象陷於保羅.
有一個馬其頓人站著求他說.
請你過到馬其頓來幫助我們.
保羅即看見這異象就隨即往馬其頓去.
以為神召我們傳福音給那裡的人聽.
終於神的這個清晰的異象臨到了.
然後保羅他們就立刻調整方向.
然後他們就.
我們再看這張圖.
可能大家也不是特別的能看清楚.

$^{81}$於是他們就朝西北.
有幾個非常非常小的城市.
都是海邊的港口城市.
第一天到了薩摩特拉.
然後第二天到了尼泊利亞.
然後從那裡來到了菲利比.
菲利比是一個什麼城市呢.
菲利比是一個羅馬的駐防城.
什麼叫做駐防城.
就是是一個軍事要塞.
然後這裡面.
住了很多很多的羅馬的退伍的軍官.
那麼保羅他們每到一個地方.
最開始雖然他們是去給外邦人傳福音.
但是你想保羅他必須有一個支點.
他的這個支點還是先到安息日.
到猶太人的會堂去給.
在會堂中守安息日的這些猶太人去傳福音.
作為一個切入點.
但是到了菲利比這個城市.
是完全沒有猶太人的.
保羅他們就經過了一番打聽.
知道雖然沒有猶太人.
但是有一群敬畏神的人.
這些人還是外邦人.
他們可能在生活當中.
耳聞或者接觸到了猶太人.
知道了猶太人的這個神.
他們對神就產生了一種由衷的敬畏.
一種渴望.
他們就渴望去敬拜這位神.
他們就自己組織在安息日.
他們還挑了安息日.
他們專門在安息日.
來到城外的一個河邊來聚集禱告.
所以這就是我來說的第二個問題.
叫做打開的門和想不到的人.
所以說為什麼是想不到.
這些人聚集在河邊禱告的這群人.
他們是外邦人.

$^{121}$他們是一群婦女.
因為這個城市的男人在安息日.
外邦人肯定都不會守安息日的.
都是該幹嘛幹嘛.
工作呀 種地呀.
有好多住房城 有很多軍人.
然後來到這兒.
聖經說.
有一位婦人她叫呂帝雅.
她是一個賣紫色布匹的人.
然後神就開啟她的心.
叫她留心聽保羅的講話.
這是一個很讓人感動的畫面.
前面說到了神兩次關門.
終於這回神開門了.
神的開門是開了一個人的心門.
聖經就兩次地強調說.
神讓開導她的心.
叫她留心聽保羅所講的話.
所以你看.
聽和留心聽是兩回事.
很多人真理是聽進去了.
但是呢.
首先可能並沒有用心.
第二可能她的心.
還被很多的糊塗的東西.
錯誤的東西.
以及屬世的驕傲所迷惑.
所以雖然信息講出去了.
但是她聽不懂.
聽不明白.
所以我就經常會.
我覺得很讓人安慰的.
經歷的一些和弟兄姊妹的互動.
就是比如說有人就會說.
我曾經也聽過這個信息幾十年前.
但是我真的沒聽進去.
他說兜兜轉轉.
經過這麼多的人生的苦難.
走了這麼多的錯路之後.

$^{161}$我今天再聽這個信息.
扎人扎住我的心.
所以這就是聽.
聽見.
但是和心完全不一樣.
留心了.
用心了.
這個道理才能真正的被聽進去.
所以當呂迪亞聽.
聽進去了之後.
就很神奇的一件事又發生了.
聖靈立刻讓他做了決定.
讓他相信了主.
而且他和他的一家當下就受了喜.
更重要的是.
當他一受喜.
他就對保羅他們說.
你若以我是真信主的.
就請到我家裡來住.
於是強留我們.
呂迪亞是一個商人.
他做的是布匹生意.
他專門賣紫色的布.
你知道紫色的布.
他們當時是用什麼.
是一種在海里提取的特殊的貝殼.
研磨之後提取的紫色.
所以這個顏色染料造價非常昂貴.
只有羅馬的貴族皇家的人.
非常有錢有地位的人.
纔可以穿紫色.
所以你就想象一下.
你知道呂迪亞的經濟實力.
他是高端客戶.
他的客戶都是高端有錢有地位的人.
所以呂迪亞是不僅有心.
也有經濟實力.
所以他能夠為這個團隊提供住宿.
其實從另外一個側面也是說明.
上帝打開他的心.

$^{201}$讓他成為基督徒之後.
呂迪亞立刻就投入了一個重要的事奉.
這個事奉就是什麼事奉.
Hospitality 接待.
接待的事奉.
Hospitality.
你知道在初代教會.
兩千多年前的人.
人家要傳福音.
可真的是用兩條腿和一張嘴.
實實在在的這幾千公里的路.
是這麼走過來的.
要長途旅行到達一個地方.
才能夠把福音傳給當地的人.
才能夠建立教會.
兼顧信徒.
當年傳福音是一件非常不容易的事情.
所以接待真的不像我們今天.
一起到一宿下一個館子.
或者把你的家敞開.
做兩個菜.
邀請大家.
我覺得在今天這樣的人來說.
做到這一點已經都不容易了.
能夠願意大家一起把家敞開.
接待弟兄姊妹的人.
都慢慢變得越來越少了.
因為我們越來在自媒體的世界當中.
人們越來越來越封閉自我.
非常滿足於沉浸在自己的小圈子裡.
因為所有的廉價的娛樂活動.
唾手可得.
不停地刷手機.
像上癮一樣.
所以讓我們特別的封閉.
但是兩千年前的人.
接待這件事情多麼的重要.
你看這些宣教士.
風塵僕僕地來到了這個地方.
無親無故.

$^{241}$完全對這個地方一無所知.
他們的腳落在的這家人.
那就是他們真正的親人.
有人說.
你的家人在哪裡.
你的家就在哪裡.
但是在初代教會真是這樣.
能夠敞開心接納他們.
為他們提供衣食住行的人.
就是他們的家人.
這個情誼是非常深厚的.
你知道不光是提供喫住這麼簡單.
我給你一個樓上的房間客房.
你去住吧.
不是.
這些接待家庭.
是要為他們提供真正的.
最強烈的情感的支持.
屬靈的支持.
家人還重要.
你要知道宣教士.
面對多麼挑戰性的文化衝突.
和各種各樣的.
甚至是巨大的逼迫.
如果沒有這些人.
他們的施工是做不下去的.
所以在彼得眼裡.
彼得有一段話說.
彼得後書說.
他說所以你們要互相款待.
互相接待.
要不發怨言.
每個人要把神所賜的恩賜活出來.
做百般恩賜的好管家.
所以在彼得眼中.
Hospitality.
接待是一個重要的屬靈恩賜.
保羅也提醒教會說.
聖徒缺乏要幫補.
客要一味的款待.

$^{281}$所以你看.
呂地亞一成為基督徒.
他就開始做了一件很重要的事.
運用自己的接待這個恩賜.
上帝給他錢.
經濟實力.
給他這一份心.
給他這一份熱情的能力.
不是白給的.
用就要用在這個最關鍵的時刻.
你想這個時機也是相當的好.
相當的巧妙.
我剛纔說了保羅之前.
他們繞了一大圈.
800多公里的路.
白走了.
你想想他們有多沮喪.
多疲憊.
然後你知道這件事情之後嗎.
他們在菲律賓這個城市.
之後就遇到了牢獄之災.
因為他們給一個.
被鬼附著的使女傳福音.
結果被那家人給誣陷.
關到監獄裡.
還受了鞭打.
在這中間.
上帝讓他們見到了呂迪亞.
也就是說上帝知道前面有危險.
後面有危險.
中間給你一個好好休息的時間.
給你一個特別能夠愛你.
關懷你的蜀靈的家人來陪伴你.
真的我們感慨上帝的美意.
是不是.
所以這就是我們今天講到的第16章.
非常的短.
我們也把這個故事梳理完了.
那麼下面我們就要進入一個.
特別重要的環節.

$^{321}$就是2000年前古人的生命.
對於我們今天有什麼提醒呢.
跟我們今天有什麼關係呢.
我想要說的第一個重要的關係就是.
關於走彎路的時候.
其實我們人生當中.
不是你特別倒黴.
神不愛你.
所以你經歷了關門.
你繞了路.
你有沮喪的時刻.
其實每一個人.
我們生在這個世界.
我們來到這個世界.
我們必定遲早.
不是今天明天就是後天都有可能.
一定的不是有可能都必然經過苦難.
每一個人都經過.
並不是說我就特別的倒黴.
這是第一點.
第二點.
我就想問大家.
為什麼當我們經歷苦難的時候.
經歷彎路的時候.
我們首先第一個想法就是去埋怨神.
為什麼你不給我清晰的帶領.
為什麼呢.
其實產生這種沮喪情緒的.
第一個特別關鍵的問題就是.
我們都有一種錯覺.
本能的錯覺.
我們都認為我就是這個世界的中心.
說這句話可能很可笑.
因為常識就是顯然了.
你我不過如飛逝去的一個生命.
怎麼可能是這個世界的中心.
但是確實啊.
我們這個主體這個人.
我們都是通過自己的.
作為自己作為一箇中心.

$^{361}$來感知這個周圍的世界的.
所以在我們所感知到的這個世界當中.
我們肯定是中心.
就是這個中心.
所以我們就有一個錯覺.
我們就覺得我們周圍的環境啊.
其他的人啊.
事情都應該圍繞著我轉.
如果別人upside我一點點.
或者說一件事情發生.
沒有按照我的意願.
不管是對人還是對神.
我們都要埋怨.
其實這個埋怨本身.
就是我們的一種錯覺和驕傲.
我們覺得我本應該如此.
這個世界就應該圍繞著我轉.
其實我們的孩子.
哪怕他們去上幼兒園.
他們走入社會的第一天.
可能就是他們去幼兒園.
他們要學的一個重要的功課.
一個五歲的孩子就是要知道.
這個世界不是應該圍繞著我轉.
我不是這個世界的中心.
我可能在家裡是.
但是我來到這個外面的世界.
我就不是了.
我在家裡都不應該是.
如果爺爺奶奶或者爸爸媽媽.
讓孩子以為自己就是這個家的中心的話.
那可能你們的教導是有問題了.
所以當我們成年之後.
我們就有一種理所當然說.
神必須為我開路.
神不能允許我繞路.
神不許給我關門.
神不許給我靜默.
我祈求你就必須給我一個答案.
其實這都是我們身為人的一個錯覺.

$^{401}$我們也不是這個世界的中心.
我們也不是特別的倒黴.
所以我們繞彎路.
每個家庭每一個人都會這樣.
生活本來就如此.
為什麼.
因為我們來到了一個破碎的世界.
我們來到了一個墮落有罪遠離神的世界.
因為我們本身我們自己也是破碎充滿瑕疵的.
所以在這樣一個充滿摩擦和張力的世界.
你覺得有可能生活永遠一帆風順嗎.
就算偉大如保羅.
保羅他可不是遊山玩水去的.
保羅可不是為了自私地滿足自己的人生的願望.
或者是為了功成名就.
青史流芳去的.
保羅是受了神的呼召去傳福音的.
他這一路經歷了多少艱辛.
我們看看保羅自己所說的.
多受苦楚.
多次下奸.
受過鞭刑棍刑.
遭石頭打.
差點打到人家都以為他死了.
三次海難.
旅行遠路.
遭江河的危險.
盜賊的危險.
同族的危險.
外邦的危險.
城裡的危險.
曠野的危險.
海中的危險.
假弟兄的危險.
受勞碌.
受困苦.
不能睡覺.
又飢又渴.
多次沒有飯喫.
受寒冷.

$^{441}$赤身露體.
連穿的衣服都沒有.
除了這外面的事.
還有為宗教會掛心的事.
天天壓在我身上.
為什麼保羅經歷了這麼多的患難.
他沒有埋怨神.
保羅給出了一個理由.
我想想大家也可以問問自己.
為什麼人生當中.
上帝讓我們經歷這麼多的苦難.
我可能剛纔給出了一些理由.
因為我們活在這個破碎的世界.
因為這個世界可以用很多患難.
來磨平我們的棱角.
讓我們謙卑下來.
因為我們人都很驕傲.
然後保羅還給出了我們這個理由.
他說因為你經歷患難.
方顯出神的大能.
讓你更深地經歷神的愛.
其實真的是這樣.
在歲月靜好的時候.
你看看你的鄰居.
不信主的鄰居.
人家照樣很nice.
很有道德感.
在這樣歲月靜好的時候.
信和不信.
你的信仰.
差別不大.
差異不大.
只有真正的患難.
讓人體會到有神的意義.
我覺得作為牧師.
我會經常經歷.
見到很多經歷過患難的人.
經歷過真正大能的人.
在大能當中.
沒有離棄神的人.

$^{481}$當他遇過了這個苦難之後.
真的可以叫做浴火重生.
真的就是脫胎換骨.
真的就是更深地經歷了神的帶領.
他的整個人都不一樣.
他會更加地堅定.
更加地溫柔.
更加地有韌性.
更加地開闊.
所以有一句話說得.
其實很有道理.
說生活中.
凡打不倒你的事情.
都是祝福你的事情.
確實是.
我們像保羅一樣.
保羅所以就告訴我們說.
為什麼要經歷患難.
他說患難有個好處就是.
讓你在患難當中.
經歷到別人經歷不到的神.
很多時候你對神的那種心愛.
和上帝對你特殊的愛和帶領.
那種不離不棄的.
與你同在的那個力量.
你在平順的時候.
你的體會真的不深刻.
也許我們能體會到.
但是我覺得不深刻.
真的在大能的當中.
你就知道神的深刻.
就像保羅所說的.
我們也會四面受敵.
但是靠著神.
我們不會被困住.
我們也會焦慮.
我們甚至可能會得抑鬱症.
我們可能也會有心理的問題.
但是靠著神.
我們最終能振奮起來.

$^{521}$因為我們始終有盼望.
我們會經歷逼迫.
生活當中經濟的逼迫.
身體的逼迫.
家庭關係的逼迫.
但是神承諾.
他不會丟棄我們.
可能人某一些.
有時候有人丟了丟棄了你.
或者說有一段關係丟棄了你.
但是神承諾.
他不會丟棄你.
我們甚至會被打倒.
但是靠著神.
你就想想那個小草.
假以時日.
只要一場春雨過後.
再幹的土地.
又會有草從裡面生出來.
這就是基督徒.
這就是有神的生命.
所以寶經幽幻是非常有意義.
有幽幻因為我們走的是天路歷程.
所以在這個過程當中.
有幽幻只能夠讓我們.
在這個天路歷程走的過程當中.
變得更堅定.
但是沒有神的人.
無論他是順境還是逆境.
他都是在這個世界當中.
巡迴往復.
他出不了世界的這個圈子.
他是破碎的.
是虛空的.
但是我們走天路的人.
我們走的是朝上的路.
所以這就是保羅繞路八百公里.
不是神的錯誤.
不是神殘忍.
而是神有意為之.

$^{561}$其實這也是給保羅.
在進入一個宣教歷程當中的一個訓練.
神讓保羅更深刻地明白一個道理.
我們不總能理解神的用意.
你看你就會覺得.
上帝你就直接明確地指引我去菲律賓.
到歐洲不就得了嗎.
為什麼我走這兒.
你說不行.
那我就只有再走別的地方了.
你又說不行.
所以我們不能總是理解神的用意.
恰恰是因為這樣.
我們更需要堅信神的帶領.
其實保羅也是沒有停.
他邊走邊尋求神的帶領.
邊走邊尋求神的帶領.
堅持著堅持著堅持著.
終於神的帶領來了.
也許此刻的你就正在經歷一些人生的彎路.
可能是婚姻關係緊張.
或者是為孩子而掛心.
或者是你自己的健康遇到了巨大的危機.
或者是你至親的人離開了你.
你可能感覺不到神的帶領.
但是保羅沒有丟下裝備.
保羅重要的一點就是他這個人特別堅持.
他就堅持.
你看他們多繞了多少路.
他們也沒停下來.
他們就一直走.
是這兒嗎.
走到這兒.
禱告.
神也沒開門.
接著走.
是那兒嗎.
走到那兒.
禱告.
沒開門.

$^{601}$但是保羅特別的有毅力.
他沒有停下來.
因為他知道宣教這條路.
人一旦泄氣了.
一旦心冷了.
一旦停下來了.
可能真的你這個勁兒就沒了.
這股心氣就淡了.
就好像.
假如有人說的特別好.
說我們好像就像一塊鐵一樣.
其實我們的出場設置.
我們來到這個世間.
我們就充滿了瑕疵.
我們變成了一個冰冰冷冷.
就像一塊鐵一樣.
我們必須得不斷地回到火裡面.
回到溫暖的源頭.
我們才能夠被烤熱.
我們才能夠把溫度傳遞出去.
我們才能變成一塊熱的鐵.
紅色的鐵.
只要你離開了火的源頭.
假以時日.
你又回到出場設置了.
我們就變成了一種驕傲頑固不化的.
自以為是的.
一塊冷冷的鐵了.
所以越在這樣的時刻.
越是迷茫.
你尋不到神帶領的時刻.
你越不能放棄.
當你心裡煩亂.
你感到無法禱告.
你越應該禱告.
因為只有認真的禱告.
才能讓你的心平靜下來.
才能夠讓你發生本質的改變.
有時候你會覺得生活太忙.
沒有時間讀經嗎.

$^{641}$其實這正是你應該讀經的時候.
忙的人通常都是沒有智慧.
需要指引的人.
忙忙亂亂卻抓不到重點.
你需要的指引正在聖經裡等著你.
你週日想偷懶不來教會嗎.
必須來.
你越是不想來.
你越得拖著自己來.
來教會不是為了人.
來教會是你用心靈和誠實來朝見神.
是你那塊冷的鐵.
不斷地重新回到火的源頭.
這些各種各樣的.
表面上看起來.
好像你覺得是一個宗教的一個行為.
實際上這不叫宗教行為.
這叫屬靈的自制力.
屬靈的操練.
屬靈的功課.
當你認認真真地去做這些功課的時候.
這塊冷的鐵.
不停地回到火的源頭.
讓它把自己燒熱.
因為你知道你離了它.
你立刻便會出場設置.
這就是經文給我們的第一個提醒.
萬路是神的帶領.
第二個提醒就是.
我們從呂地亞來竊入.
叫做打開心門.
真實連接.
神有美意.
當你回到呂地亞.
你看他一信主.
他就把自己的家庭開放.
他就強留.
聖經四個字說強留我們.
這就把呂地亞那種慷慨豪放真摯.
大氣的這種特質.

$^{681}$一下子就刻畫出來了.
你是不是能夠瞬間感受到.
有特別特別溫暖的.
這麼一個很慷慨的.
又心很大.
也很寬大的這種愛的力量.
其實這個魅力.
最重要的是一個什麼根本呢.
其實就是.
呂地亞具有與人連接的能力.
連接的能力.
真的很重要.
很多人這一生.
其實他都不能不會.
也不敢把自己真正的打開.
哪怕是他和他最親的人.
他都沒有真正的連接.
有人說我可能是我性格比較內斂.
我是一個很內向的人.
我不善社交.
但是真實的連接和性格沒有一點關係.
內向的人也需要和人連接.
連接不是說你聊天的時候.
很聊得很嗨.
聊得很好.
這就叫真摯連接.
不是.
真摯的連接是敞開自己.
內向的人也可以勇敢地敞開自己.
不是什麼都跟別人去說.
而是你從心裡有一個敞開的態度.
用真心去擁抱對方.
幾周前我和我們國語堂的一位姐妹聊天.
她說一段話特別地震撼我.
她說對我來說十字架就代表著連接.
向上和神連接.
別忘了還有左右和人連接.
如果一個人.
一個基督徒號稱他真的跟神連上了.
他感受到了神的愛.

$^{721}$他聖靈總是充滿著他.
總有神的愛.
你覺得他可能不跟周圍的人真實地連接嗎.
沒可能.
那他肯定是自然而然地.
他就有一個連接的能力.
就像呂迪亞一樣.
當神把他的心打開.
他接受了福音的那一刻.
他聽到了保羅他們的難處.
他就是自發地.
就是完全本能地他就要跟他連接了.
他本能地就把自己的心敞開.
把自己的家敞開了.
這是一個真正跟神連接的人.
他必然要跟人連接.
你看就像保羅.
保羅其實他這個人的性格特點.
我們讀前面也發現.
比如他和巴拿巴的那些爭執.
你就仔細讀聖經的使徒行傳和保羅書信.
他其實是一個非常非常有個性的人.
他不是一個很.
不一定是一個很容易接觸的人.
早期的保羅.
你想他去大馬士革.
去追捕那些基督徒的時候.
他一個人自願請命大祭司去追殺基督徒.
他要一個人去.
他成為基督徒之後.
他就一個人跑回大樹住了很多年.
直到巴拿巴去把他找出來邀請他.
很有可能保羅的性格.
他真的是他在開始.
他是一個很自我目標性很強.
但是不太願意與別人連接的人.
但是你讀到保羅書信的後面.
使徒行傳的後面.
你就覺得保羅是一個多麼真摯坦誠.
可以用心的人.

$^{761}$你看他在書信後面.
他把提摩泰稱為自己的兒子.
他說到提多說是我心上的人.
他屢次的寫信說.
你一定要把馬可帶來.
他跟馬可有過節.
他說你又要把馬可帶來.
讓他到監獄裡來看我.
他和他團隊的每一個人都成為生死之交.
把他們每一個人都親如家人.
其實你想一想也make sense.
對吧.
保羅經歷了三次宣教旅行.
這一路上遇到了那麼多的風險.
但是他這個人從來沒有抱怨.
從來沒有覺得苦讀.
從來都沒有說疲憊呀或者什麼.
剛纔讀到了原因.
他說因為上帝與我同在.
其實還有一個隱含的原因.
因為他不但和神連接.
他感受到了從神而來的愛.
他也與人連接.
與他的團隊.
他去到的每一個接待家庭.
都給他無限的愛.
他也跟他們完全地敞開心.
成為真正的屬靈的夥伴.
成為真正的屬靈的家人.
比家人還親的家人.
所以說白了.
其實傳福音這件事情.
就是在踐行一個深度的連接.
你覺得我參加一個傳福音培訓班.
我參加一個三福的培訓.
或者我去了教會的短宣.
我就會傳福音了.
我就可以傳福音了.
其實根本不是.
一個傳福音的人.

$^{801}$其實首先你要願意把自己獻出來.
不是說一天24小時都在教會里面.
不是這個意思.
把自己獻出來是把自己的心真的敞開.
願意和還沒有認識神的那些人.
去產生深刻的連接.
只有你真的是把心敞開.
跟人家連上了.
人家對方纔能打開心門.
願意來聽你說的話.
而不是高高在上.
教育你幾句聖經的話.
或者是認給你一個屬靈的四個原則.
一個傳福音的單章.
其實哪怕.
我們都有這樣的感受.
就是哪怕是一個對一個路人.
一個真誠的敞開心的打招呼.
和一個表面敷衍的打招呼.
一秒鐘的兩人相遇的時間.
對方都是有感受的.
更何況.
我們這些號稱帶著福音去拯救靈魂的人.
我們真的一定要像呂迪亞和鮑羅一樣.
把心打開.
所以其實今天的這個整個的這個故事.
我覺得也像我們人生的一個縮影吧.
我們希望我們回去的時候帶回兩個信息.
一個就是.
如果你在彎路在迷茫在挫折當中.
你不要覺得上帝殘忍給你不斷的關門.
這正是上帝帶領你的時間.
他的帶領從來從來都沒有錯誤.
也可能.
你想想也許鮑羅多繞了那些路.
上帝也在預備呂迪亞這邊.
他也在預備他的心.
讓兩個人的時間完美的就是碰在一起.
還有.
你覺得人生的意義是什麼.

$^{841}$人生的意義是掙了很多錢.
或者說功成名就或者怎麼樣嗎.
其實真的都不是.
人生真實的意義給你喜樂的源泉.
其實就是連結.
就是我們所說的關係.
就是我們所說的愛和神連結.
然後也真實的和人連結.
真實的和你的愛人配偶.
你連結了嗎.
你真實的和你的父母連結了嗎.
你真實的和你的弟兄姊妹打開心門連結了嗎.
所以我們都要學習.
在的天父.
謝謝你.
在我們人生當中發生了很多的事情.
我們看不清真相.
暫時看不清真相.
不知道原因.
我們在彎路當中.
我們在黑暗裡.
但是我們堅信你的旨意是美善的.
我們糾結是因為.
我們只站在生命的一個短暫的階段.
而你站在永恆裡.
所以萬事相互效力.
讓愛神的人得益處.
謝謝你啊.
號召我們呼召我們去連結.
可能我們不敞開是因為害怕受傷害.
但是我們不怕.
因為我們有你.
若人對我們有虧復.
我們堅信天父.
你都會加倍的賞賜回來.
所以也求你在我們與人交往的過程當中.
幫助我們更加的真誠.
更加的不計回報.
與人連結.
就像呂地亞把自己的家敞開給保羅那樣.

$^{881}$奉主耶穌基督的名求.
Amen.
\newpage



\section{啟示錄 1:7-8}
\label{sec:031876}
\textbf{耶穌基督的啟示}~~~~譚綺敏傳道
\newline
\newline
連結: \href{http://www.church.com.hk/acms/content.asp?site=cdc\&op=show\&type=product\&code=031876\&layout=sermon}{\texttt{031876}} ~~~~ 語音日期: 2026-03-29
\newline
\newline
\hyperref[sec:031875]{\small{< < < PREV SERMON < < <}}
~
\hyperref[sec:index_chronic]{\small{[返順時目]}}
~
\hyperref[sec:031877]{\small{> > > NEXT SERMON > > >}}
\newline
\newline
啟示錄 1:7-8
\newline
\begin{longtable}{cl}
\hline
\hline
章節 & 經文 (和合本修訂版)\\
\hline
1:7 & \begin{tabularx}{0.7\textwidth}{X} 「看哪，他駕雲降臨；眾目都要看見他，連刺他的人也要看見他；地上的萬族要因他哀哭。」這是真實的。阿們！ \end{tabularx} \\ \\ \relax
1:8 & \begin{tabularx}{0.7\textwidth}{X} 主神說：「我是阿拉法，我是俄梅戛 ，是今在、昔在、以後永在的全能者。」 \end{tabularx} \\ \\
[1ex]
\hline
\hline
\end{longtable}
$^{1}$請姐妹平安.
今天我們先讀經文.
不如我們一起讀吧.
啟示錄的一章七到八節.
我們一起讀吧.
漢啊!他駕雲降臨.
眾目要看見他.
連次他的人也要看見他.
地上的萬族都要因他哀哭.
這話是真實的.
阿們!主神說.
我是阿拉伐.
我是阿梅加.
是昔在今在.
以後永在的全能者.
講一個小故事吧.
有一個人.
發白日夢.
他夢見自己去了旅行.
到了一個很大的教堂.
有一間小的房屋.
他看見有一個.
長長的會議桌.
兩邊坐滿了人.
大家愁眉苦臉.
他問牧師.
為何這麼開心.
牧師說.
因為他們正在開教會研討會.
討論啟示錄的意象和意思.
大家都覺得.
很難明白.
各持己見.
那就不太開心了.
但他又到另一個房間.
裡面充滿了歡樂.
開心的聲音.
他又問牧師.
裡面的人在做什麼.
為何這麼開心.

$^{41}$其實他們沒有讀明啟示錄的意象.
但他們在啟示錄裡.
單單看到耶穌基督.
就開心快樂了.
所以我盼望.
我們一起在啟示錄第一章.
因為單單看見.
耶穌基督就歡喜快樂.
其實我們很容易.
煩惱的.
最近局勢不穩定.
不知道中東戰爭.
要打到什麼時候.
油價漲到什麼時候.
都不知道.
通脹又來了.
大家生活都開始不容易了.
雖然我們不在戰亂的地區.
其實我們很感恩.
但其實我們.
不多不少都會.
不安.
因為世界各樣的事情.
都會影響我們的情緒.
所以我們也看見.
很多天災人禍.
水浸.
落薄.
很多人都會覺得.
和聖經說的世界末日很相似.
在幾百年前.
如果你說世界末日.
很多人都未必覺得有.
但現在看見世界動盪.
看見天災人禍.
教會也有一些壓力.
有些地方.
基督徒都會受逼迫.
所以其實.
心裡都有些憂慮.

$^{81}$其實我們最需要的是什麼呢?.
在憂慮的時候.
我們最需要的是什麼呢?.
我想我們不是要更多的.
預測到世界末來的局勢.
是怎樣走.
資訊是怎樣.
我們不是要這些.
而是我們要重新看見耶穌基督.
所以盼望.
我們今天都可以看見耶穌基督.
其實啟示錄不是一個.
預言的天書.
好像中國人說.
看一看通勝.
看一看現在是怎樣.
有時天氣都看一看通勝.
但其實.
聖經啟示錄.
不是一個通勝.
不過通勝這個字我覺得.
頗有意思.
即是通通勝利.
其實啟示錄第一章.
或者整本啟示錄.
主題都是告訴我們.
主耶穌通通勝利.
最後都是祂勝利.
所以我們最需要看到.
就是勝利.
但我們在神裡.
每樣東西都得勝.
啟示錄第一章.
用了兩個.
很震撼力的詞語.
一開始講到.
啟示這個字.
和漢拿這個字.
第一章一開始第一節.
就講到耶穌基督的啟示.

$^{121}$啟示這個字在.
希臘原文的意思.
就是指揭開.
揭開一些東西.
或者拿開一些遮住的東西.
顯露真相.
或者是拿開.
一個柄子.
所以簡單來說.
原本看不清楚的真相.
就顯明出來了.
好像你去看舞台上.
布幕前面遮住.
觀眾不知道.
布幕後面發生什麼事.
但一揭開布幕.
大家就看見.
舞台上的真相.
所以啟示錄一開始.
講到耶穌基督的啟示.
即是說耶穌基督揭開一些東西.
揭開些什麼呢.
大家之後就會.
更加知道.
我們都有點緊張耶穌揭開些什麼.
牧師在去年講到.
今年我們不如講啟示錄.
其實我馬上就有點緊張.
因為我怕.
我不知道怎樣解釋.
那些經文.
如果大家讀過啟示錄.
不容易理解.
要講得好.
很難.
因為是啟示的文學.
所以其實.
我們一起去學習的時候.
我們一起去學.
我們不單止有一個好奇心.

$^{161}$知道將來的社會是怎樣.
世界是怎樣.
啟示錄不是完全講這些.
但是.
我們要有全天秘的心.
去了解神的心意.
以致我們可以面對世界的苦難.
第二個字就是.
看.
這個字是一個感嘆詞.
即是說你看一下.
不單止叫人用眼睛去看.
還是一個很強調式的宣告.
你看呀你看呀.
是引起讀者的關注.
叫我們去特別.
看到一些東西.
才有回應.
所以既然耶穌基督揭開一些東西.
有啟示.
我們也要看到.
才有意義.
對於我們來說才有生命的改變.
所以啟示錄很特別.
在書中反覆出現.
我看見.
我又看見.
我聽見.
這些都是啟示文學的特點.
在《神比喻》看到很多這樣的意象.
他也描述出來.
他看見漢奸的意象.
也描述出來.
引發我們去注意.
注意什麼呢?.
第七節說到.
看呀!他駕雲降臨.
眾目都要看見他.
連癡他的人都要看見他.
地上萬族都要因他哀哭.

$^{201}$另外.
約翰他看見.
耶穌基督降臨.
這對於他.
當時是很重要的.
他看見耶穌基督降臨.
是以榮耀君王的身份降臨.
對於約翰來說.
那時候是很震撼的.
約翰在哪裡呢?.
他在一個被放逐的.
拔摩海島裡.
周圍都很荒涼.
在監獄裡.
他看見的東西和他現在身處的環境.
是極大的對比.
一個曾經受苦.
好像囚犯一樣.
被殺的.
釘在十字架上的耶穌.
他看見他好像萬王之王.
萬主之主的身份.
降臨.
好像我們剛才唱的那首歌.
這是一個很大的反轉.
所以曾經他看見.
被審判的.
今天他很榮耀地來到當中.
曾經是被羞辱的.
今天很榮耀地.
來到當中做審判者.
所以.
第一章五至六節.
說到耶穌很愛我們.
他愛我們.
用自己的血使我們脫離罪惡.
又使我們成為國民.
做他父神的祭司.
所以主耶穌很想.
我們回到那裡.

$^{241}$使我們脫離罪惡.
成為祂的子民.
用祂的生命拯救了我們.
祂沒有丟棄祂的子民.
所以祂說.
我曾經去拯救你們.
將來也會繼續去拯救你們.
去審判這個世界.
我就來到了.
所以這個對於約翰來說.
是很震撼的.
其實對於當時的.
信徒來說.
都是一個很大的鼓勵.
為什麼這麼說呢?.
因為第一世紀.
你說要信耶穌就要受到逼迫.
那時候羅馬帝國的皇帝.
他也立了自己的像.
我們看.
《舊約聖經》.
也看過.
如果你拒絕去.
拜那個像的時候.
拜羅馬君王的時候.
你就會被.
經濟制裁.
今天也很容易經濟制裁一些國家.
你沒法去生活.
你沒得買賣.
然後你就可能會監禁.
好像約翰那樣放逐.
也是監禁.
甚至死.
所以在教會來說.
他們看不到榮耀.
他們看到的是什麼呢?.
看到羅馬帝國的權勢.
一切榮耀都在羅馬帝國.
皇帝有權有勢.

$^{281}$如果你去反抗.
如果你不去敬拜他.
你是敬拜耶穌.
你就要面臨.
一些監禁.
制裁.
甚至死亡.
所以其實現實.
是很殘酷的.
但我們知道.
主耶穌是一個又真實.
又信實的神.
所以當約翰在拔摩海島.
被囚禁.
被監獄的時候.
他看見主耶穌的榮耀.
這個對於他來說.
是一個很大的力量.
他親身知道神和他同在.
也鼓勵當時的信徒.
所以在這麼艱難的時候.
約翰看見了.
今天你有沒有看見.
你看見的是什麼呢?.
我們很多時候看見.
很現實的環境.
但都未至於被逼迫.
我們看見.
很不公平的事情.
在世界發生.
我們看見很多事情.
很不開心.
很不公平.
但你要知道你所認識的神.
是一個什麼樣的神.
很值得我們思想.
今天我們認識什麼樣的神.
主耶穌是一個什麼樣的神.
對我們來說是很重要的.
其實我們對神的認識.

$^{321}$決定了我們面對苦難的態度.
當我們越認識神的時候.
我們越能夠在困苦中.
站穩腳.
因為知道神和我們同在.
神與我們一起.
祂沒有丟棄我們.
繼續看第七節.
再來的時候.
眾人都會看見祂.
連刺祂的人都會看見祂.
這個很特別.
在什麼年代,什麼時間.
什麼地方的人.
都會看見祂.
曾經被藐視,被羞辱.
被釘在十字架上的耶穌.
無論你信不信祂.
你都看見祂.
將來祂以一個勝利者.
王的姿勢.
回來.
提醒我們信徒.
現在這個世界有很多事情.
是不公平的.
羞辱,困苦是不公平的.
但這些都是暫時的.
將來當耶穌基督.
來到的時候.
榮耀是永恆的.
祂亦都向全世界的人.
顯現.
因為祂要.
向全世界去審判.
為信徒去伸冤.
翻轉所有的不公平.
亦都沒有人能逃避.
所以說到地上萬族.
都要因祂哀哭.
我想到對於我們.

$^{361}$信徒身份的人來說.
我們應該開心快樂.
耶穌來,你會不會開心?.
想一想.
你會不會開心快樂呢?.
應該是很開心快樂的日子.
但祂說到萬族.
都哀哭,為什麼呢?.
因為對於不信耶穌的人.
其他世界.
大部分的人.
是不信耶穌的.
所以對於大部分的人.
他們都會哀哭.
所以萬族都要哀哭.
他們會感到驚惶.
感到恐懼.
因為一個.
神審判的時候.
可能是他們沉淪的時候.
所以他們會哀哭.
但對於我們信耶穌的人.
你記住是一個歡喜快樂.
很幸福的日子.
因為我們永遠可以和主耶穌一起.
所以約翰不單止看到.
這個景象.
他還聽見神對他說.
我是阿拉伯.
我是俄梅加.
是昔在今在.
以後永在的全能者.
這個說話是非常震撼的.
我是阿拉伯.
我是俄梅加.
大家都聽過了.
Alpha和Omega.
是希臘文的第一個字母.
希臘文有二十四個字母.
英文有二十六個.

$^{401}$它是第一個字母.
我們就是A.
它也是A字頭.
翻譯是這樣.
叫Alpha.
最後一個字是Omega.
Omega就是最後一個字母.
表明神.
是歷史的起點.
創造這個世界.
也是歷史的終點.
將來他來到這個世界.
要審判這個世界.
所以人類的歷史.
掌握在神的手中.
然後說到昔在今在.
以後永在.
這個是神的名字.
記得摩西問.
神叫什麼名字.
他的名字是自有永有.
也就是昔在今在.
以後永在.
意思就是無論在什麼時間.
我都在.
過去現在將來.
我都和人類同在.
我是無所不在的神.
時間空間.
都是在我的掌管之下.
然後說到全能者.
這個字.
不只是有能力.
這個字在原文的意思.
就是雙手緊緊地.
握住所有東西.
就是沒有東西.
out of control.
沒有東西離開主宰之下.
在神沒有難成的事.

$^{441}$今天我們很沒有安全感.
因為我們不可以掌控.
我們不可以.
所有東西都在我們掌握之內.
但是在神裡面.
所有東西都是他的掌握之內.
所以當.
約翰看見.
這些異象.
也聽見神這樣對他說的時候.
他就.
安然下去.
這個看見榮耀的主的時候.
給他一個喜樂.
給他一個力量.
我想起《史提反》.
在《史徒行傳》第七章記載.
史提反受逼迫的時候.
你會發覺很特別.
他看著天.
看見了主的榮耀.
為什麼他會看見了.
他沒有描述.
但是他看見神的榮耀.
他自己看到.
他說我看見天開了.
人子站在我的右邊.
被石頭打死了.
很痛被石頭打死.
但是他心裡仰望主.
他臨死的時候很有力量說了一句話.
他說主啊.
不要將這罪歸於他們.
好像主耶穌在十字架上.
第一句說話.
主啊赦免他們.
因為他們所做的他們了得.
所以當我們看見榮耀的主的時候.
不只是給我們喜樂.
是給我們可以去.

$^{481}$整個生命都可以有改變.
我們有力量去面對一些壓迫.
很多時候我們不是不相信主.
但是我們很容易.
被環境淹沒.
當我們身體出現了一些問題.
或者家庭出現了一些狀況.
孩子出現了一些情況.
或者工作有壓力.
經濟有壓力.
甚至在事奉裡面.
有些壓力.
那時候可能我們會問.
主啊我已經禱告了這麼久.
你有沒有聽啊.
你是不是真的在掌權啊.
這些都是我們很容易.
擔心的問題.
因為我們沒有辦法掌控一切.
我們不知道將來如何.
今天一個中東的戰事.
我們估不到.
在二月底突然間發生.
接著呢.
我們都開始受到影響.
以致我們的安全感.
真的不知道建立在什麼地方.
但是主耶穌告訴你.
你的安全感.
要建基在我這裡.
如果我們真的看到主耶穌.
你就會發覺.
這些忍耐等候不是一些重擔.
是真實的盼望.
因為聖經都已經說了.
我們有真實的盼望.
基督徒的盼望.
不是在世界的未來.
不是世界變得好還是不好.
是主耶穌的太平.

$^{521}$而是我們生命的終結.
我們歸宿在神的盼望裡.
在天家的盼望.
我們永遠可以和主耶穌在一起.
所以接下來.
約翰就用自己的經歷.
去描述什麼叫做.
看見榮耀的主.
所以一直看下去.
第九節到尾的時候.
你會發覺約翰就是用他的經歷.
在這裡去說.
他說:我約翰就是你們的弟兄.
和你們在耶穌的患難.
各道忍耐裡.
一同有份為神的道.
並為給耶穌.
作的見證.
曾在那名叫拔摩島的島上.
就是約翰.
寫這個啟示錄的時候.
他說.
曾在那名叫拔摩島的島上.
我們知道.
約翰不是被神死在拔摩島.
如果你去看看.
一些正典.
有些傳統.
是說.
他被釋放.
因為轉了一個王.
那個王離開了.
轉了第二個.
羅馬王的時候.
他就用一些懷柔政策.
釋放多一些受迫害人.
所以約翰後來是被釋放的.
他在.
以弗所教會.
可能他寫下來的時候.

$^{561}$是在以弗所教會.
約翰是因為傳神的道.
被耶穌.
給耶穌作見證.
以致他流放在拔摩島的監獄.
在人的角度來看.
約翰可能是一個失敗者.
今天你為神去作供.
為什麼你會淪落到這樣的地步呢?.
一個老人家.
被社會遺棄.
丟到荒島裡.
你的荒士宮.
做了什麼呢?.
走到最盡頭.
但是神沒有忘記約翰.
也沒有離棄約翰.
在這個絕望的拔摩島裡.
有一件很特別的事情.
發生了.
第十節.
「當主人我被聖靈感動.
聽見在我後面.
有大聲音如吹號」.
所以約翰.
在他最無力,最失望的時候.
在拔摩島的時候.
他就見到這些異象.
他被聖靈感動.
聽見後面有聲音.
其實一個老人家.
他可以怎樣反應呢?.
他馬上轉身來.
第十二節.
他轉身來.
看誰要跟他講話.
其實這個轉身的動作.
是很有意思的.
在希臘文裡意味著一個.
刻意的決定.

$^{601}$刻意地轉身.
想想如果你坐得很舒服.
老人家在看電視.
有人叫你.
你會不會突然轉身?.
可能也不會.
兒子叫你也未必會.
突然轉身.
看誰要跟我講話.
你作這個決定.
轉身也要力氣.
一個這麼老年紀的約翰.
所以你會發覺.
約翰很認真去面對.
他身邊的事情.
很細微的事情.
他也懂得去反應,去欣賞.
看明白這是什麼聲音.
其實.
我們世界裡充斥著不同的聲音.
也有不同的影像.
特別是AI時代.
不用思考了.
一打開電話.
什麼事情都告訴你.
有很多聲音,很多影像.
所以我們要有勇氣.
去跟這個世界說.
我們轉身了,不要再看手機了.
不要再看電視了.
不要再看Youtube.
不要再看ViuTV.
TikTok.
很難.
你叫小朋友這樣做.
也很難,因為已經習慣了.
所以其實這個轉向.
是一個很好的提醒.
我們會不會不跟從這個世界.
去轉向耶穌基督呢?.

$^{641}$也很真實.
要我們去想一想.
當我們轉向的時候.
約翰就會看到.
第十可靠的主.
第十二節.
當他轉向的時候,他看到什麼呢?.
他看到七個金燈台.
金燈台.
然後呢.
一直看下去的時候.
你會發覺他看到.
好像人子一樣.
在中間行走.
金燈台其實是代表教會.
表示耶穌基督.
與教會同在.
在教會裡面行走.
他仍然托著教會.
保守著教會,督責著教會.
然後呢.
約翰就很詳細地描述.
這個人子是什麼形象.
他說身穿長衣.
表示他大祭司的身份.
然後呢.
他腰束金帶.
尊榮.
君王身份.
也都呢.
頭髮是.
白的.
白是代表智慧.
所以在舊約裡面.
也有提到.
一般看到神的形象.
就是白膠.
有智慧.
有永恆的意思.
所以.

$^{681}$眼目如火焰.
明察秋毫,洞悉一切.
人的內心,教會的真實.
隱藏的罪.
他都看到,好像一個審判者.
腳好像光明瞳.
即是說他踐踏仇敵.
即是他很光明的.
拳柄不可以搖動.
他可以踏住仇敵.
聲音如重水.
即是他的聲音.
很雄壯,很威嚴.
令人敬畏.
右手拿著七星.
因為七星其實是.
代表著教會,所以接下來.
你會發覺是講七教會.
即是托著,掌管著教會.
口中套出利劍.
即是神的島.
很銳利.
發出拳柄的審判.
會來到當中.
面貌是放光的.
神的榮耀在當中.
神聖的榮耀在當中.
當約翰很近距離地描述這位主.
認識這位主的時候.
他的反應是怎樣呢?.
第十七節.
他說:「我一看見.
就撲倒在他腳前.
就像死了一樣.
很近距離地撲倒在他腳前.
你想想你自己會是怎樣?」.
我就想.
通常我們說耶穌很愛你.
剛才約翰說.
那麼榮耀.

$^{721}$他應該是一見到耶穌的時候.
給我們的印象.
耶穌那麼溫柔,那麼善良.
好像好朋友一樣.
經常可以去祈禱,向祂吐苦水.
約翰也是.
耶穌的親密朋友.
是親密使徒.
你發覺他親眼看見耶穌復活.
耶穌也向使徒顯現.
他也見過天上的主.
看著祂怎樣去升天.
當他那麼近距離地見到耶穌的時候.
你發覺他的反應不是衝上去.
馬上抱著耶穌.
「主啊,你和我同在,真是好」.
然後馬上去說話.
他沒有這樣做.
他反而突然撲倒.
撲倒在地上.
還在祂的腳前.
為什麼會這樣反應呢?.
你想想,如果你看見這個異象.
你會怎樣反應呢?.
你會開心,還是像約翰那樣?.
其實也可能開心.
但要說到.
可能說到.
約翰很覺得.
見到這個異象.
看到神的臉會像死一樣.
所以當約翰見到的時候.
不知道發生什麼事.
從來沒有見過.
從來沒有見過的東西.
突然在你面前出現.
你可能會嚇到.
所以約翰的即時反應.
突然撲倒在耶穌的腳前.
好像死了一樣.

$^{761}$動也不敢動,因為太震撼.
口中套一把兩刃利劍.
好像飛落我身上.
所以真的很害怕.
人工智能的圖片也不錯.
所以現在是人工智能時代.
你被一個形象.
出來的時候.
你會發覺自己也很害怕.
所以約翰.
他那個時間.
害怕到不敢動.
很震撼.
接下來的經文.
你會發覺.
可能是全段第一章.
我覺得是最感人的部分.
為什麼呢.
耶穌還用右手.
摸著約翰.
跟他說.
不用怕,不要懼怕.
主耶穌很溫柔地.
看到這個老人家.
在他前面很害怕.
可能覺得自己不配.
突然發生什麼事.
不知道耶穌基督來到.
我覺得自己很多罪未洗乾淨.
所以突然很害怕.
耶穌基督.
向他重申.
他說什麼呢.
我是首先的,我是末後的.
又是那存活的.
我曾死過,現在又活了.
直到永永遠遠.
並且拿著死亡陰間的.
鑰匙.
所以耶穌去兼顧他.

$^{801}$你現在看到的我.
是歷史的主宰.
我曾經死過.
但我現在又活了.
直到永永遠遠.
就像剛才.
釋在今在以後.
永在的全能神.
再說給約翰聽.
主耶穌是很榮耀的.
但他也很慈愛.
他是審判的主.
可以很銳利.
但他也是.
安慰人的主.
他掌管歷史的神.
有什麼都在他裡面.
他曾經死過.
是為了救人而死的.
他是救贖主.
死亡和陰間的鑰匙.
都在裡面.
他可以打開門.
讓你再活過來.
所以生死的權柄.
在主耶穌的手裡.
所以當主耶穌這樣跟你說的時候.
你會有什麼感受呢?.
當你患重病的時候.
主敢跟你說.
當你做一個很大的手術.
不知道生死的時候.
主敢跟你說.
當你面對一些.
很煎熬的環境的時候.
主會這樣跟你說.
你有什麼感受呢?.
我相信這是一個很大的安慰.
上星期小祖也有一個感恩的分享.
有一對夫妻.

$^{841}$他們很緊急.
要去一個地方.
去弟弟做手術.
一個很大的手術.
一個很大的心臟手術.
這個手術很有風險.
連醫生都說.
百克有多少都不太清楚.
可能是五十,五十.
可能在手術裡隨時死亡都有風險.
所以他的家人都很多.
走到他身邊陪伴他.
陪他做手術.
這個弟弟是基督徒.
面對這麼大的手術.
他都說自己不怕.
我記得.
那個夫婦分享.
他在醫院裡.
經常拿著一本聖經.
去祈禱.
他連他自己什麼時候要做這個手術都不知道.
因為當時他有肺積水.
要等肺積水清了.
他才可以做手術.
但他心裡有奇妙的平安.
所以因為他信靠神.
他的生命在神的掌管裡.
所以他知道.
無論是什麼情況.
他都願意去交給神.
最後這個手術很成功.
在過程裡很多奇妙的事情都發生.
這對夫婦回到我們當中.
很感恩的見證.
他們都說了很多.
因為他們去的時候是中東戰爭的時候.
所以狂姬.
令他們看到神的安排.
所以這對夫婦回來之後.

$^{881}$姊妹就相信了主.
弟兄也願意去相信主耶穌.
願意去決志相信主耶穌.
所以我們很感恩.
當時我們祈禱的就是這對夫婦.
所以我們聽到見證的時候.
我們心裡都很感恩.
是神自己的帶領.
當我們親身經歷過神的時候.
是神祂自己很奇妙的帶領.
所以最後主耶穌也給了約翰一個使命.
不是單單去看見這些.
祂給了約翰一個使命.
不是去告訴你.
這些你就停在這裡.
祂給了一個使命是怎樣呢.
十九節.
祂說你要把所看見的.
和現在的事.
和將來必成的事.
都要寫出來.
所以當看見神的時候.
不是一個個人的平安.
不是一個神秘的經歷.
是一切很細微的.
在監獄裡面寫道.
一個人不只能夠相信上帝.
我們也說我們相信神.
相信上帝.
能夠因著相信而真實地看見上帝.
不是相信是更加真實地看見神.
怎樣可以真實地看見神呢?.
在一個沒有自由,狹窄的囚室裡.
可能隨時會面對酷刑的囚室裡.
潘福華所看見的不是苦難.
他是看見比圍牆,比一切環境更大的主耶穌.
他對苦難的看法是.
上帝不是要將我們在苦難裡面救出來.
而是要在苦難裡面與我們同行.
盼望我們看見神怎樣在苦難裡面與我們同行.

$^{921}$在1945年4月9日.
在希特拉自殺前21天.
在1945年4月9日.
\newpage



\section{創世記 50:19}
\label{sec:031877}
\textbf{聖經人聖經事:我豈能代替神呢?}~~~~黃克勤牧師博士
\newline
\newline
連結: \href{http://www.church.com.hk/acms/content.asp?site=cdc\&op=show\&type=product\&code=031877\&layout=sermon}{\texttt{031877}} ~~~~ 語音日期: 2026-03-24
\newline
\newline
\hyperref[sec:031876]{\small{< < < PREV SERMON < < <}}
~
\hyperref[sec:index_chronic]{\small{[返順時目]}}
~
\hyperref[sec:031878]{\small{> > > NEXT SERMON > > >}}
\newline
\newline
創世記 50:19
\newline
\begin{longtable}{cl}
\hline
\hline
章節 & 經文 (和合本修訂版)\\
\hline
50:19 & \begin{tabularx}{0.7\textwidth}{X} 約瑟對他們說：「不要怕，我豈能代替神呢？ \end{tabularx} \\ \\
[1ex]
\hline
\hline
\end{longtable}
$^{1}$大家好,今次是我豈能代替神,誰夠膽代替神?.
不過話說回來,我們時時都自己去代替神.
今日生命說書的網絡為大家預備了聖經人和聖經的故事.
讓我們一齊去看看我們的上帝和這個故事.
今天的主人翁仍然是弱失,他說了這句說話,這句說話真是好.
我都要時時記住這句說話,希望你都記得,我怎能夠代替神呢?.
看看經文,創世記去到第五十章是最尾那一章,創世記.
第十五節,這個時間雅各已經死了.
弱失的哥哥們見父親死了,就說去到這個時間危險了.
因為或者,未必,或者弱失懷恨我們.
不是呀,他對我們幾好,沒什麼問題,都沒有什麼心懷怨恨.
叫我們不要自憂自恨,就這樣說這裡.
但是有些就會想,或者他其實裝得很漂亮,可能他是懷恨我們.
會照著我們從前代他一切的惡,足足地報復我們.
你會看到現在很多世界上的事都是.
就是很多時候那個人的行為,說話是由裡面浮出來.
然後投射在別人身上,如果是那些叫什麼呢?.
我想偷襲你,然後就說你為什麼要偷襲我?.
他會先告別人,很有趣的.
又或者是,鬼鬼祟祟,是不是想偷錢?.
其實他自己有個心想偷人錢,他才會看到別人想偷錢.
整天看著我,為什麼這樣?有什麼企圖?其實他是自己先有企圖.
他才會覺得別人有的.
這班弟弟,或者哥哥的弟弟,沒有弟弟,哥哥就會在這裡說.
他會好像從前我們代他那樣,足足地報復我們.
其實就是他們自己這個人心凶宅,他們會.
不過於是由於這樣,他們會看到別人會.
他看不到弱勢那種胸襟的廣闊,他看不到他的靈性.
就是將自己投射在那裡.
或者他說是我們,照著我們從前代他的,足足地報應我們.
那怎麼辦呢?不如這樣,我們打發人去見弱勢,跟他說.
你父親未死以先就吩咐過,你們要對弱勢這樣說.
從前你哥哥們惡待你,求你饒恕他們的過犯和罪惡.
如今求你饒恕你父親臣的僕人的過犯.
他們對弱勢說者說,弱勢就哭了.
其實他們是說叫人去,叫誰去呢?.
其實是自己去,一班兄弟去找弱勢.
弱勢,爸爸死之前跟我們說,叫我們告訴你.
你要饒恕我們,饒恕你的僕人.
現在我們也是跟你,你真的大了,像和菌一樣向你下拜.

$^{41}$好像天上的星十一粒一樣向你下拜.
現在我們做你的僕人,你救回我們,你又饒恕我們.
如今求你饒恕你父親的僕人.
爸爸在他死之前叫我們來告訴你的.
如果你是弱勢你會怎樣?.
不是吧?真的假的?為什麼他不告訴我?.
不是,弱勢沒有.
他們對弱勢說者說,弱勢就哭了.
弱勢也經歷過不少事情,壓力也很大.
弱勢就哭了,經歷過什麼?.
經歷過被賣,圍爐,坐牢,被冤枉.
然後一人之下萬人之上去搞定整個國家的糧食.
七年豐年,七年荒年,不可以亂.
佛老又要信任他.
哇,坐這些位置,真的三十位,很大壓力.
有壓力的人是容易哭的.
而經歷了這麼多的事情.
哎呀,見不到爸爸,想念他的兄弟,想念家鄉.
是會很容易哭,弱勢就哭了.
但這個弱勢哭了,是被他們刺激到的.
這個弱勢為什麼哭了呢?.
哭了就是,其實這個弱勢哭出來就是.
哎呀,哥,為什麼你們這樣都不信我?.
為什麼你們這樣都不信我?.
弱勢這麼聰明,你以為哥走來說這些話.
哎呀,弟,現在我們都做你的僕人了.
爸爸未死之先就叫我們來跟你講求你饒恕我們.
他叫你饒恕我們呀.
唉,在弱勢來講,一早我就饒恕了.
我那天只是玩玩,搞鬼你們而已.
我一早就饒恕了,你們放心吧.
其實呢,如果我是弱勢也會想哭.
哎呀,來到這個階段,爸爸剛剛又死了.
你們就已經這樣了,唉,這樣都不信我.
如果我要害你們,我一早就害了.
根本你們回不了家接爸爸來.
弱勢哭了,其實他是愛他的兄弟.
愛他的兄弟?你們,我這麼愛你們,這麼疼你們.
為什麼你們會這樣想我?哭了.
所以這個哭了,其實我自己覺得是反映出.

$^{81}$弱勢是很愛他們的.
看你哥哥們又來苦苦在他面前說.
我們是你的僕人.
本來我那段就派人去講.
現在這一段就自己去,唉,我們都是你的僕人.
你都聽你的話了,你放心吧.
我們會跟著你,叫我們做什麼都做.
你對我們不要殺害,不要報仇就行了.
弱勢對他們說,不要害怕,不用怕.
不要害怕什麼?是不用怕我.
這個不用害怕,不用怕我.
第二點,我豈能代替神呢?.
我豈能代替神呢?其實是說我不會代替上帝.
這個我也在努力學習.
我豈能代替神呢?什麼意思?.
其實神是安排,神的手在背後.
我們學聖經人,聖經士就是嘗試去看.
神在背後一直工作,帶領著保護我.
又帶領著保護你們,一直到現在.
其實也不是說一定可以的.
起碼我跟我的老闆,跟法老說.
我爸爸,我的哥哥,成族人都來了.
這樣好啊,這樣娛樂啊.
我就想選一塊地給他們.
選最好的那塊地.
其實弱勢很懂的,埃及人是種田的.
講到明,埃及人是很討厭那些畜牧業的人.
不喜歡的,如果你在埃及做畜牧業.
是會看低一點的.
耕田的,看高一點的.
還有要分開吃飯.
埃及人是不喜歡跟外族或畜牧業的人一起吃的.
但是都已經安排了很好的地.
在哥山地給他們.
因此你會知道其實神的手在背後一直工作.
你要信靠神,神會一直默默地在後面工作.
可能你經過死人丘,經過青草地溪水泊.
但是祂的手沒有離開,祂的杖,祂的肝都一直帶領著你.
我弱勢就說,我豈能不代替神呢?.
神安排的這些事,我不會代替祂做什麼呢?.

$^{121}$審判官,其實最多人是代替神做什麼呢?.
自己走出來,這件事我已經查明了.
是你做得不對,你做得錯,你不對,你奸搞.
因此我現在代替神,替天行義.
做了審判官,這個在新約聖經叫做論斷.
論斷就是以審判官自居的意思.
論斷,好像我沒事,你有事,全部問題都在你那裡.
我就說你……論斷別人,審判別人.
我豈能代替神做審判呢?.
如果神要做工作對付你們,祂會做的.
我不能夠代替祂,我不能夠主動地代替祂.
神啊,你多謝我,因為我已經代替你出手了.
不會啊,神啊,我不敢出手.
如果是出手的,你出手吧.
我豈能代替神呢?神說我會的.
神一定會的,審判的主嗎?.
我豈能代替神呢?要學這件事.
約瑟學到了,經歷過那麼多低谷.
他一定有很多怨恨.
不過他已經由自憂自恨裡面出來.
所以他會說到各位兄長,你們無需自憂自恨.
他一定是釋放了他的內心,他才會這樣.
現在再多一步,我是不會代替上帝審判你們的.
追究的,神讓他自己做吧.
第三點,從前你們的意思當時我真的明白.
我都生你的氣,你們是要害我.
不過第三點,神的意思原是好的.
原來在那麼壞的事裡面,其實神在做的事.
神的意思原是好的.
那你要看看,一件很壞的事.
臨到有什麼屬靈的功課可以學到.
我都有教基督徒生命教練.
怎樣用生命教練的方法去幫助人.
可以在他自己的問題遭遇裡面.
可以釋放出來和看到神怎樣帶領.
時時會問的就是,在這件事上.
你學到一個什麼屬靈的功課呢?.
學到什麼屬靈的功課呢?.
很多時候我們會看到,我在這件事上遭遇了金錢的損失.
遭遇了名譽的損失,被人騙了我.

$^{161}$感情被欺騙,被歧視,很痛苦.
在這件事上,你學到什麼?.
屬靈的功課學到什麼?.
正面的屬靈功課學到什麼?.
學習到耶穌,他怎樣經歷那些對他反對的人.
學到屬靈的功課.
第三點,神的意思,原是好的.
他本身有一個美好的旨意.
是什麼呢?.
起碼我看到的是保存許多人的性命.
許多人是指我們全族人的性命.
神這樣安排我,差遣我先來.
是要保存許多人的性命.
所以他覺得自己有任務.
第二就是成就今天的光景.
讓我們不單止保存生命.
今天還發了,厲害了.
我們家族在埃及地,全部人都認識我們.
今天這麼好的光景.
所以不要害怕.
我是會養活你們和你們的婦人孩子.
這裡講得不錯,若是真的有經歷是什麼呢?.
就是我不吃,我的婦人孩子都要.
若是你殺死我們,不要緊.
我的家爺仔那真的不知如何是好.
若是說你不用擔心.
我會養活你們,供應你們.
不單止不殺你們,供應你們.
和你們的婦人孩子.
所以放心吧.
於是若是用親愛的話安慰他們.
若是真的很好.
若是在聖經裡面被認為沒什麼過失的人.
沒什麼過失的人有一次是不得醒目.
或者是留了地步.
以致他一個人進入他主母的房間裡.
沒有其他人.
其實你找他進去,你知道是這樣的.
還這樣?.
你這樣拿著衣服給他脫?.

$^{201}$為什麼要拿著衣服不要讓他脫?.
脫到救命救命,引你引你.
你就沒事了.
反而他有事了.
一下子自負了之類.
他沒有做錯事.
藥室很小心.
不過有疏忽的地方也不奇怪.
但是以罪,大家找不到藥室有什麼罪.
挺好的.
而且他還用親愛的話去安慰他的哥哥.
這個就是藥室故事的結尾.
真的很感恩.
非常好.
在藥室身上學到東西.
再看看.
第一,藥室愛他的兄弟.
哭了聽到.
你們不用說這些.
拿爸爸出來騙我.
其實我不會的.
你這樣也不相信我?.
這個遭遇人家不信任是時時都有的.
第二.
我不會代替上帝做審判者.
等神自己工作.
交託給他.
第三.
神的意思是好的.
是要保住我們性命和成就今天這麼好的光景.
即是神有恩典在背後.
祂的原意是好的.
我們要信任的.
我們在地上生活要怎樣?.
彼此用親愛安慰的說話來生活就好了.
這個也是藥室教我們的.
聖經人,聖經士今天到這裡了.
我們一起禱告.
讚美我們的主.
你真是使用藥室.

$^{241}$讓祂成為我們今天在地上生活的一個榜樣.
祂的人生經歷.
真的可以拍一套戲.
有高有低.
經歷了很多事.
不過.
在這件事的時候看到的.
是原來你的手在祂背後一直工作.
是你的手扶著祂.
是你的手托著祂.
是你保護了祂.
是你為祂開路.
是你給祂智慧.
是你給祂懂得解夢.
是你給祂有管理國家行政的智慧.
有一個和祂的哥哥相處的智慧和愛.
主求你自己憐憫.
幫助我們學習藥室.
讓我們看到.
你才是真正的掌權者.
我們不會搶你的位去做.
我們代替你.
我們豈能代替上帝呢?.
你的旨意完是好的.
願你得到榮耀讚美.
祈禱奉耶穌基督的名.
阿們.
\newpage



\section{腓立比書 2:5-11}
\label{sec:031878}
\textbf{以基督的心為心}~~~~蘇善美宣教師
\newline
\newline
連結: \href{http://www.church.com.hk/acms/content.asp?site=cdc\&op=show\&type=product\&code=031878\&layout=sermon}{\texttt{031878}} ~~~~ 語音日期: 2026-03-29
\newline
\newline
\hyperref[sec:031877]{\small{< < < PREV SERMON < < <}}
~
\hyperref[sec:index_chronic]{\small{[返順時目]}}
~
\hyperref[sec:031879]{\small{> > > NEXT SERMON > > >}}
\newline
\newline
腓立比書 2:5-11
\newline
\begin{longtable}{cl}
\hline
\hline
章節 & 經文 (和合本修訂版)\\
\hline
2:5 & \begin{tabularx}{0.7\textwidth}{X} 你們當以基督耶穌的心為心： \end{tabularx} \\ \\ \relax
2:6 & \begin{tabularx}{0.7\textwidth}{X} 他本有神的形像，卻不堅持自己與神同等； \end{tabularx} \\ \\ \relax
2:7 & \begin{tabularx}{0.7\textwidth}{X} 反倒虛己，取了奴僕的形像，成為人的樣式；既有人的樣子， \end{tabularx} \\ \\ \relax
2:8 & \begin{tabularx}{0.7\textwidth}{X} 就謙卑自己，存心順服，以至於死，且死在十字架上。 \end{tabularx} \\ \\ \relax
2:9 & \begin{tabularx}{0.7\textwidth}{X} 所以神把他升為至高，又賜給他超乎萬名之上的名， \end{tabularx} \\ \\ \relax
2:10 & \begin{tabularx}{0.7\textwidth}{X} 使一切在天上的、地上的和地底下的，因耶穌的名，眾膝都要跪下， \end{tabularx} \\ \\ \relax
2:11 & \begin{tabularx}{0.7\textwidth}{X} 眾口都要宣認：耶穌基督是主，歸榮耀給父神。 \end{tabularx} \\ \\
[1ex]
\hline
\hline
\end{longtable}
$^{1}$丙姐妹早晨.
剛才在一場優美的音樂.
很有意境的詩歌裡面.
開始我們今天的講道.
丙姐妹 不知道大家剛才來到看.
剛才來讀出的一段經文時有甚麼感覺.
我們今天的三內經課.
其中一段就記載在.
菲勒比書第二章五到十一節.
大家方便的話都可以打開大家的新約聖經.
第三百零二頁.
是菲勒比書第二章五到十一節.
在三內經課裡面有一個很特別的地方.
它將菲勒比書裡面一個很中心的思想.
吸了出來 拿了出來.
成為今天的講道主題的經文.
今天的題目是來自菲勒比書第二章五節.
今天的經文在菲勒比書.
題目也是在菲勒比書.
題目是以基督耶穌的心為心.
不知道大家想起基督耶穌的心為心會想起甚麼.
菲勒比書是一本甚麼書呢.
菲勒比書很有意思.
是當時保羅在歐洲建立第一間教會的時候.
來寫給他們的信.
菲勒比教會也是當時歐洲整個布道之旅的時候.
第一間建立的教會.
可以說是心裡面BB的教會.
而菲勒比這個地方也是很特別的.
在羅馬政府的統治之下.
是其中一個馬其頓的省.
這個省在正中間.
可以說是可以出去也可以進去.
地理環境非常之不錯.
當我們想起菲勒比書的時候.
我們經常會想起一段經文.
你要喜樂 你要靠基督喜樂.
但為甚麼菲勒比書保羅寫的時候會這樣說呢.
其實沒甚麼可能的.
這間是當時保羅建立的第一間教會.

$^{41}$也是當時保羅最困難的時候.
他正在去坐牢.
這間教會當時受著的不單止是外面人的攻擊.
異端就來攻擊這一班剛剛信主的信徒.
甚至乎來告訴他們.
他們所信的可能不是真的.
他們跟他們說.
你們要回去你們的國禮.
因為這樣才能夠蒙上帝的旨意.
這樣才能蒙上帝的福.
令到這一班新興的很熱心的小信徒.
心裡有很多的迷惑.
於是保羅就寫了這封信來告訴他們.
這些所信的 這些所傳的.
全部都不是真的.
基督要在你裡面得自由.
而這個自由是喜樂.
因為保羅太深愛這間教會.
當他們遇到這些問題的時候.
一想起就會覺得很喜樂.
為這一班信徒來感恩.
這個就是當時腓立比書寫的時候的背景.
所以書中提到.
他們患難中要喜樂.
他們要靠主來喜樂.
這個並不是口號.
而是當時保羅很希望這一班信徒能夠明白.
最不能夠喜樂的人.
保羅都能夠喜樂.
所以你們沒有什麼值得不喜樂.
喜樂似乎是腓立比書中心的一個思想.
當我來到想的時候.
既然信徒應該要喜樂.
不知道弟兄姊妹.
今天你覺得自己喜不喜樂.
開心嗎.
根據近日的一個調查報告.
不知道你覺得香港人的快樂指數有多高.
前兩年都不錯的.
香港人的快樂指數都有6點多.

$^{81}$在2022年到24年的時候.
今天已經是2026年了.
你覺得你的快樂指數是10分.
裡面你覺得你可以拿多少分呢.
有沒有人是10分的請舉手.
沒有.
掩著頭.
有沒有人覺得是7分的.
有呀 都不少.
似乎是長者為主.
看看這班年輕人的指數有多少.
有沒有人覺得自己不要太高要求.
6分左右.
有沒有6分.
都有.
我們執事都有6分.
似乎大家教育都很開心.
事奉都很喜樂.
有沒有覺得是5分的.
看到了.
可以放下手.
幾個 看到.
有沒有覺得自己低於5分的.
都看到.
不少.
知不知道香港人的.
今年我們得的快樂指數是多少.
都不是太差的.
都有5.49.
全世界裡面我們排行第88.
即是我們在全世界裡面.
香港人是第88位快樂的.
知不知道最快樂是哪個國家.
芬蘭.
他依然穩守7.73.
他依然是成為最快樂的地方.
知不知道我們的調查除了說.
為何香港人不快樂.
哪兩樣東西令我們最不快樂.
會不會都是你令你不快樂的原因.

$^{121}$第一 經濟.
沒錯.
經濟不好.
這幾年的大氣候.
都令到大家都很艱難.
整個香港都很艱難.
第二樣東西.
猜不猜得到是什麼.
猜不到.
是人與人之間的關係.
原來令到我們的心情低落的.
不單止是整個大氣候.
可能我們跟人之間的關係.
會不會變得疏離.
或者人與人之間的關係.
沒有以前那麼緊密.
令我們的快樂.
我們的喜樂指數一點都不高.
在菲勒比書裡面.
告訴我們.
中心的思想是.
我們要靠住喜樂.
我們應該過一個喜樂的生活.
就算在困所中的保羅.
在監獄中的保羅.
都鼓勵我們.
我們要過喜樂的人生.
究竟我們怎樣去重拾.
基督徒喜樂的人生呢.
今天看的經文.
雖然只有幾節.
甚至是將整個的菲勒比書.
最中心的那一段.
我們叫第二章第六節到第十一節.
稱之為「瀑人之歌」.
是幾歲的主耶穌基督.
怎樣來到道成肉身.
為我們拋棄自己.
甚至放棄自己應有的權利.
在「瀑人之歌」裡面.

$^{161}$我想大家留意.
我們今天.
很想跟大家分享一個.
我們如何重得喜樂的主題.
就是第二章五節.
我們一起讀好嗎.
來囉.
「你們當以基督耶穌的心為心」.
這一句聖經可以說是貫穿了.
整個菲勒比書一個很重要.
非常之重要的重心思想.
為什麼這樣說呢.
在二章五節這個.
「當以基督耶穌的心為心」之上.
即是由第一章到第二章四節這裡.
其實是在說.
保羅在說.
我們信徒應該怎樣來生活.
我們怎樣彼此相顧.
不要只顧自己的事.
也要顧別人的事.
我們的行事為人.
怎樣與我們基督徒的身份相稱.
於是在第二章的第五節.
保羅就用了一個小的總結.
來總結菲勒比書的上半部分.
他說「你們當以基督耶穌的心為心」.
這是上文的一個總結.
當我們能夠不只顧別人的事.
也不只顧自己的事.
我們也顧教會的事的時候.
我們很容易明白那個.
「以基督耶穌的心為心」.
到經文下半節的時候.
到二章六節到第四章完結的時候.
這裡是在說什麼呢.
是在說.
當我們要以基督耶穌的心為心的時候.
我們應該要效法基督.
來成為我們的榜樣.

$^{201}$行事為人活出真正的見證.
我們就可以從這些異端的邪術.
來得勝過來.
所以二章第五節.
也是整個菲勒比書下半部.
由二章六節開始.
到第四章的一個開始.
這樣就很有趣了.
第二章五節的.
「你們當以基督耶穌的心為心」.
是上半章的一個總結.
同時也是下半章的一個開始.
這一句經文.
以我的理解.
它是整個菲勒比書裡面.
最重心.
也是最重要的思想.
「當以基督耶穌的心為心」.
怎樣解釋呢這一句.
就是說我們行事為人.
我們就用耶穌基督的想法.
耶穌基督的思想.
來成為我們的想法.
我們的思想.
我們就可以以他的心為心.
是不是這個意思呢.
在二章五節的.
「你們當以基督耶穌的心為心」.
「的心為心」這四個字.
其實在整個菲勒比書裡面.
一共出現了七次.
包括這裡的二章五節.
一共出現了七次.
這個心為心.
不單止譯作「以心比心」這麼簡單.
也可以譯作「以耶穌的意念為我的意念」.
或者是「以耶穌的思想為我的思想」.
如果是譯作「意念」的話.
這裡出現了有幾次.
分別是第一章第七節.

$^{241}$二章第五節.
就是這裡.
四章第二節.
如果以心為心.
不是譯作「意念」.
是譯作「思想」的話.
同樣是出現了一次.
第二章二節出現了一次.
如果以心為心.
譯作「與基督同一個想像」.
「同一個想法」.
這個一共出現了三次.
分別是第三章的十五節.
和第十九節.
和第四章的第十節.
我們會問.
「以基督耶穌的心」.
可以譯作「意念」.
可以譯作「思想」.
可以譯作「想法」.
這些字全部和「心為心」.
是來自同一個字.
對我們有什麼意思?.
或者對我們來解讀這段經文.
有什麼意思呢?.
剛才我弄了一個腓立比書.
「以基督耶穌的心為心」.
包括這裡.
是出現了七次.
這六次和二章五節的這一次.
它們之間有什麼特別之處呢?.
我想弟兄姊妹來留意.
在剛才所說的七處聖經裡.
那六處同樣都是譯作.
「意念」,「想法」,「思想」.
但是唯有第二章五節.
「你們當以基督耶穌的心為心」.
其實它原來的意思是一個特別的譯法.
它不單只是說到耶穌的心.
也不單只是說到耶穌的思想這麼簡單.

$^{281}$當我來看這段經文的時候.
我覺得這一句很能夠觸動我的心.
我們如何能夠以基督的心為我們的心.
但是我想告訴弟兄姊妹.
原來這一句不單只是解作.
用基督的心腸成為我們的心腸這麼簡單.
當我來看的時候.
我嘗試用我很有限的希臘文.
將這一句用它原來的意思來譯出來.
這個意思本來是這樣說的.
「你們當以基督耶穌的心為心」.
但是這句原來的意思.
回到最基本的時候.
其實它是怎樣來說的呢?.
這一句是這樣譯的.
它說「要在你們裡面有這個意念.
就是也在基督裡的」.
完了.
再聽一次.
這個意思原來是這樣譯出來的.
「你們要以基督耶穌的心為心」.
原來意思是.
「要在你們裡面有這個意念.
就是也在基督耶穌裡的」.
其實原來意思是沒有意念的.
不過我覺得加上去就變得完整一點.
「要在你們裡面有這個意念.
就是也在基督裡的意念」.
大家聽不聽得到分別?.
當我們說基督耶穌的心為心.
在這一段經文裡.
在整個腓立比書出現了七次的時候.
唯獨是在二章五節裡.
在這裡出現了一樣東西.
在以上六節經文是沒有出現過的.
大家聽不聽得出來?.
「要在你們裡面有這個意念.
就是在基督耶穌裡的意念」.
這個字就是「基督耶穌裡」.
在原文裡出現了「基督耶穌裡」.

$^{321}$而這個字跟之前那六節並不一樣.
什麼叫做「基督耶穌裡的意念」呢?.
我們說「在主裡我們是一家人.
在主裡我們是互相照顧」.
因為我們明白.
為什麼以基督耶穌的心為心.
要「在基督裡」呢?.
不在基督裡有沒有問題呢?.
在這裡保羅不只是想教導我們.
我們要有耶穌的心.
我們要同心.
我們要有一個意念.
在這裡想說.
如果我們要重拾主裡的喜樂.
我們要學習以基督耶穌的心為心.
而基督耶穌的心為心.
更重要的地方就是要「在主裡」.
「在主裡有基督耶穌的心為心」.
什麼叫做「在主裡」呢?.
我們想了很久.
什麼叫做「在主裡」呢?.
為什麼這個「在主裡」只出現在二章第五節.
沒有在其他經文裡出現.
「在主裡」這個字.
原來的意思.
其實是跟哥林多後書五章十七節有關係.
可能大家聽完這段經文.
就會知道「在基督耶穌裡」是什麼意思.
「在主裡」的意思是這樣的.
在哥林多後書五章十七節.
若有人在基督裡.
這個「在基督裡」正正就是我們剛才所說.
在基督耶穌裡的意念.
那個「在基督耶穌裡」是一樣的.
若有人在基督裡.
他就是新造的人.
如果在主裡 在基督裡的時候.
我們能夠有耶穌基督的心為心.
原來「在基督耶穌裡」一個更重要的提醒就是.
我們有沒有一個新的生命.

$^{361}$去承載這一段跟上帝的關係.
原來你們以基督耶穌的心為心.
你們在基督裡要跟耶穌同心.
不只是解作同心.
原來我們在基督裡跟耶穌同心的時候.
耶穌在問我們.
究竟在基督裡我們已經成為一個新造的人了嗎.
你是否一個新造的人.
若有人在基督裡.
他就是新造的人.
舊事已過都變成新的了.
在基督裡我們能夠意念更新.
我們能夠有思想的改變.
但是我們成為一個新造的人了嗎.
我們脫棄了以前一些我們不應該的想法.
一些不應該的行為或思想了嗎.
究竟我們在裡面已經轉化成為一個新造的人了嗎.
這個很值得我自己去反省.
在基督裡就是解作新造的人.
原來當我們說要跟耶穌同一個心思.
同一個意念的時候.
原來我們拿著一個舊人在裡面的時候.
我們不能夠更新.
我們不能夠體貼著耶穌的心思意念.
因為我們並未成為一個新造的人.
什麼叫新造的人呢.
大家平時在家中喜不喜歡扔垃圾.
可能姊妹比較喜歡.
亂寫來扔垃圾.
我本身也很喜歡扔垃圾.
但是前六七年左右.
經常有一袋垃圾不捨得扔.
那一袋其實也不是垃圾.
但是一些沒用的東西.
就是扔什麼都扔.
那一袋東西總是不捨得扔.
有沒有一些東西是你總是不捨得扔的.
我那袋是什麼呢?就是我的子女的嬰兒衣服.
大家千萬不要覺得噁心.
我會用那袋嬰兒衣服.

$^{401}$偶爾會用密封袋子封住它.
嬰兒衣服洗乾淨.
打開的時候就鎖一鎖.
拉上去.
一個月鎖一次.
為什麼要鎖呢?.
因為總覺得有一股嬰兒味.
他們兩個現在長大了.
沒有嬰兒味了.
但是那股嬰兒味總覺得在衣服裡.
兒子一袋女兒一袋.
鎖著鎖著.
後來我先生疫情開始.
他實在接受不了.
他說你真的很不衛生.
不斷鎖很不衛生.
有細菌的.
我不捨得洗.
洗得多沒有那股味道.
我不喜歡.
總之就是用密封袋.
鎖完再拉上去.
保存味道.
一個月又拿出來.
又鎖一鎖又拉上去.
逐袋鎖.
後來我先生就說.
你扔掉吧.
疫情.
你骯髒了這麼多年.
沒有細菌也有細菌.
你說我捨不捨得扔掉呢?.
當然不捨得扔掉.
這些是寶貝.
後來我在網上找到一個方法.
去解決這兩袋衣服.
不知道大家知不知道.
有一些婦女.
幫人做一些拼布.
有沒有聽過?.

$^{441}$她可以將寶寶的口水肩.
他的被子.
他的衣服.
剪成一塊.
然後用一個像被袋一樣.
將他一幅幅布繡起.
成為一格格的圖畫.
你還可以看看.
幾時幾月這一格是他幾歲的時候.
然後你就不需要扔掉那些衣服.
還可以慢慢抱著他來索.
你覺得這是不是一個挺好的方法?.
是的.
我就拿去做了一張被子.
我就在想.
若有人在基督裡.
他就是新造的人.
在主裡我們成為一個新造的人.
是否當我將這些寶寶的衣服扔掉後.
我就成為一個新造的人呢?.
這個當然有機會是.
但我相信在這裡所說.
在基督裡成為一個新造人.
我相信上帝並不是要扔掉我們原有的東西.
而是將你原有的東西.
重新改造.
轉化.
成為新造的人.
就好像那些寶寶的衣服一樣.
我沒有將它扔掉.
因為在一個方法之下.
他將它變成了一張好像新的被子.
但其實那張被子是不是新的呢?.
其實不是新的.
因為那些一格格的衣服.
全部都是那些小朋友的舊衣服.
就是這樣來轉化.
這個就是若有人在基督裡.
他就是新造的人.
舊事雖然已經過.

$^{481}$變成新.
但上帝是沒有扔掉我們的東西.
他將我們一些.
本來我們都覺得不要或不想要的東西.
來轉化.
我們就能夠有一個新人.
新的樣式.
你們當以耶穌基督的心為心.
我們怎能夠將耶穌的心.
成為我們的心.
他的想法成為我們的想法.
究竟我們內在裡.
我們這個人更新了沒有呢?.
我們這個人成為一個新造的人了沒有呢?.
如果這一刻,弟兄姊妹,你覺得很困難.
有些東西我還未拋得開.
有些東西我還未能夠拿走.
不需要擔心.
我相信上帝有恩典.
我們求天父幫助我們.
將我們的舊有的扭轉過來.
變成新的一樣.
讓我們可以繼續去效法基督.
怎樣以他的心為心.
當我們知道.
我們這個以心為心.
是在基督裡的時候.
是在基督裡證明我們的生命.
需要一個新的改變.
是需要一個新造的人生.
第二樣更重要的是.
當我們知道我們有一個新造的人生的時候.
保羅並不是這樣停止.
他在第二章的第六節和第十一節.
給了我們一個很好的機會去學習.
六節到十一節是「獨人之歌」.
是菲勒比書裡面一段很重要的.
基督論的「獨人之歌」.
將上帝本有的形象.
怎樣來到不堅持.

$^{521}$將上帝來到同等.
反倒虛己.
取掉路報的形象.
成為人的樣式.
既有人的樣式.
就卑微自己.
全心信服.
以至於死.
且死在十字架上.
上帝對人的捨棄.
不是一步一步.
他不以自己為上帝的形象.
不單止.
他放棄.
放棄之餘.
他更加去倒空自己.
成為一個根本不應該是他的形象.
奴僕的身份.
甚至乎成為人的樣式.
正正就是有人的樣式.
他就學習怎樣去卑微.
甚至乎他全心的信服.
以至於死.
且死在最受咒詛的十字架上.
在這裡短短第六至第八節裡.
用了一個層帝式.
來表述主耶穌怎樣去行他的僕人之歌.
一步一步地為人犧牲.
當我們說.
在基督耶穌裡.
我們生命要更新的時候.
或者我們都在想.
在基督耶穌裡.
我們的生命是需要效法基督.
就好像僕人之歌一樣.
效法基督.
怎樣效法呢.
我做什麼來效法.
如果我們看經文的第六節的話.
我認為整個的第六節.

$^{561}$是整個僕人之歌一個很重要的要點.
他本有上帝的形象.
卻不堅持與自己上帝為同等.
反倒虛己.
他不堅持.
不單止不堅持.
而且還是虛己.
什麼叫不堅持.
什麼叫虛己呢.
在我預備這段經文的時候.
我想起我以前的一個中學老師.
他給我的感覺比較冷漠.
比較少說話.
但我有一個很深的印象.
到現在我仍然記得.
我覺得他很愛他太太.
為什麼呢.
他跟我說.
有一次他以前看電影.
不是前後座.
是分高等.
超等.
不是高等.
我真的給他過.
超等貴不貴.
貴的.
可能是兩倍價錢.
最高最好的位置.
他以前和他太太去看電影.
他太太很喜歡看電影.
結果那場戲.
很多人買票.
只剩下兩三張票.
其中兩張是超等座.
一張在走廊旁邊.
另一個位置在武雷宮那麼遠.
你是他.
你是那位老師.
你會怎麼坐呢.
當然是太太坐在走廊附近.

$^{601}$他自己當然坐在另一個位置.
他不是.
他怎麼做呢.
我不知道為什麼.
我小時候聽到.
覺得他真的很愛妻.
他太太坐在走廊的丹丁位.
買了兩張票.
他自己就坐在丹丁位.
旁邊走廊那一級的樓梯.
那是什麼位置呢.
他說不是一個位置.
他就坐在那裡.
太太就坐在旁邊.
陪他太太看電影.
他說最麻煩的是什麼呢.
不是坐在地上那麼簡單.
因為以前那些戲院很窄.
每一次人家上上落落.
買爆谷.
去洗手間又怎樣.
因為他坐下了.
人家不會拍他.
就怎樣.
踢他.
踢他的腰.
踢他的背脊叫他起身.
他就起身.
然後被人家下去.
然後再坐下.
他說整場戲裡.
不記得那場戲做什麼.
只記得被人踢過多少次.
起了多少次身.
我心裡覺得很感動.
雖然現在不會有機會.
我們現在的位置都很舒服.
但為什麼我要這樣做呢.
我相信他這樣做.
他不以自己.

$^{641}$本身買了票的那種消費者的好處.
同等為強奪.
他不坐那個位.
他反倒坐在一個.
他不應該坐的位上.
為什麼.
因為他為了他所愛的.
他甘願放下.
他為了他所愛的.
甘願放下.
弟兄姊妹.
當我們來到看第六節到第十一節的時候.
這個瀑人之歌告訴我們.
上帝不以他是上帝的兒子.
那種高等高高在上的身份.
來跟我們說話.
他是不堅持自己與上帝為同等.
明明可以和上帝一樣.
是同等擁有那種權利和身份.
不單止不同等.
第七字是反倒虛己.
反倒虛己這個字原來意思就是.
將自己倒空.
本來有的變成一無所有.
這個就是反倒虛己.
除了露目的形象.
當我們重拾在主裡面那種.
好像保羅一樣.
這麼喜樂的人生的時候.
我們都會再問.
我們的人生有沒有以基督耶穌的心為心呢?.
而這個基督耶穌的心為心.
有沒有在主耶穌基督裡面.
我們的心思意念已經得著更新呢?.
這個在主耶穌裡面的心思意念更新的時候.
我們有沒有來學習.
我們的主一樣.
願意放下我們應有的權利和身份.
我們不堅持自己所有.
為的是愛我們的主呢?.

$^{681}$我相信這個是我們一生的學習.
當我們能夠生命更新.
當我們能夠學習為自己.
去放棄我們一些權利的時候.
我相信我們對以基督耶穌的心為心.
能夠再走近一步.
弟兄姊妹都是很繁忙.
要做的事都很多.
要你放下你的思緒.
放下你的工作.
每天回到教會裡面我相信大家都很艱難.
但會不會正正在這個御福期裡面.
讓我們都可以看到.
主耶穌為我們的罪釘死十字架上.
祂不以自己和神同等為強奪.
祂反倒虛偈.
正正是這個思想.
幫助我們生命來更新.
幫助我們生命有動力去學習放下自己.
盼望這個放下.
從我們的日常生活開始.
我們放下我們自己可以多睡一會的時間.
我們放下一會可以和我們家人去飲茶的時間.
我們首先回到上帝的殿裡面.
我們敬拜祂.
然後我們就能夠慢慢去體會.
得著以耶穌基督的心為心.
求天父在一個重要的日子裡面.
給我們每一個都同樣有提醒.
我們一起來讀禱告吧.
因為你給我們的恩典.
是大過我們一切的所想所求.
當你在這裡以僕人之歌.
勉勵我們要以基督耶穌的心為心的時候.
主啊 求你讓我們看到.
這個心必須要在基督裡.
當我們在基督裡的時候.
讓我們常常反省.
我們心有沒有更新而變化.
我們有沒有一個新的生命在我們裡面呢.

$^{721}$如果我們仍然覺得很困難.
主啊 求你幫助我們.
當我們要去更新的時候.
我們要效法你的那種捨己.
我們知道是很難.
但是主啊 讓我們學習.
怎樣去放下我們以為我們有的.
我們學習去放下.
我們以為我們可以堅持.
我們為主走上一步.
我們更貼近你的心思.
更貼近你的意念.
多謝你 聽我們禮別時的禱告.
是奉主耶穌基督的聖名祈求.
其他事項.
\newpage



\section{馬太福音 21:1-11}
\label{sec:031879}
\textbf{路不能亂走,車不能亂坐}~~~~林豪恩先生
\newline
\newline
連結: \href{http://www.church.com.hk/acms/content.asp?site=cdc\&op=show\&type=product\&code=031879\&layout=sermon}{\texttt{031879}} ~~~~ 語音日期: 2026-03-29
\newline
\newline
\hyperref[sec:031878]{\small{< < < PREV SERMON < < <}}
~
\hyperref[sec:index_chronic]{\small{[返順時目]}}
~
\hyperref[sec:031880]{\small{> > > NEXT SERMON > > >}}
\newline
\newline
馬太福音 21:1-11
\newline
\begin{longtable}{cl}
\hline
\hline
章節 & 經文 (和合本修訂版)\\
\hline
21:1 & \begin{tabularx}{0.7\textwidth}{X} 耶穌和門徒快到耶路撒冷，進了橄欖山的伯法其時，打發兩個門徒， \end{tabularx} \\ \\ \relax
21:2 & \begin{tabularx}{0.7\textwidth}{X} 對他們說：「你們往對面村子裡去，會立刻看見一匹驢拴在那裡，還有驢駒同在一處，解開牠們，牽到我這裡來。 \end{tabularx} \\ \\ \relax
21:3 & \begin{tabularx}{0.7\textwidth}{X} 若有人對你們說甚麼，你們就說：『主要用牠們。』那人會立刻讓你們牽來。」 \end{tabularx} \\ \\ \relax
21:4 & \begin{tabularx}{0.7\textwidth}{X} 這事發生是要應驗先知所說的話： \end{tabularx} \\ \\ \relax
21:5 & \begin{tabularx}{0.7\textwidth}{X} 「要對錫安的兒女說：看哪，你的王來到你這裡，謙和地騎著驢，騎著小驢—驢的駒子。」 \end{tabularx} \\ \\ \relax
21:6 & \begin{tabularx}{0.7\textwidth}{X} 門徒就照耶穌所吩咐的去做， \end{tabularx} \\ \\ \relax
21:7 & \begin{tabularx}{0.7\textwidth}{X} 牽了驢和驢駒來，把他們的衣服搭在上面，耶穌就騎上。 \end{tabularx} \\ \\ \relax
21:8 & \begin{tabularx}{0.7\textwidth}{X} 許許多多的人把自己的衣服鋪在路上，還有人砍下樹枝來鋪在路上。 \end{tabularx} \\ \\ \relax
21:9 & \begin{tabularx}{0.7\textwidth}{X} 前呼後擁的人群喊著說：「和散那歸於大衛之子！奉主名來的是應當稱頌的！至高無上的，和散那！」 \end{tabularx} \\ \\ \relax
21:10 & \begin{tabularx}{0.7\textwidth}{X} 耶穌進了耶路撒冷，全城都驚動了，說：「這是誰？」 \end{tabularx} \\ \\ \relax
21:11 & \begin{tabularx}{0.7\textwidth}{X} 眾人說：「這是從加利利的拿撒勒來的先知耶穌。」 \end{tabularx} \\ \\
[1ex]
\hline
\hline
\end{longtable}
$^{1}$各位 請問大家有沒有試過搭飛機的時候.
驚喜地被升級上了頭等艙的經驗呢?.
好像間中都有的 有些人都點頭.
有一位網友分享說.
我拿著一張經濟客艙的機票登機的時候.
空姐將我的機票換了另外一張.
就跟她說 女士 你的座位被升級了.
她一時反應不來 就拿著機票向前走.
看著機票座位寫著01A.
就心想可能是經濟客艙裡比較寬大前面那個位.
沒想過竟然被升級成為商務客艙位.
於是她說我蹦蹦跳跳地走進寬敞的商務艙走道.
坐下就馬上開始玩各種座位的按鈕.
這位網友總結說.
好幸運啊 人生第一次商務艙成就達成 超爽.
接下來很快會是節假期了.
不知道你買了機票來呢?.
有沒有期望過意外地被升級到商務艙或者是頭等呢?.
突如其來地被升級到商務艙的確是意外驚喜.
不過常言道 水能載舟 亦能覆舟.
驚喜之後 小心變成驚嚇.
在2026年1月22日的報紙有一段這樣的報導.
香港駐天津聯絡處主任在社交平台分享.
獲噴射飛航高級管理人員安排港澳航程.
坐私人艙室由香港去澳門.
政制及內地事務局獲悉事件後表示.
已決定即時終止該公務員的合約.
經過進一步調查.
認為該公務員早前承傳前澳門市座位獲升級的事件.
有違反防止賄賂條例下.
有關公務員未經許可接受利益的規定之嫌.
已把個案轉交執法部門跟進.
自古至今代步工具都有某種含義.
好像香港初初開埠的時候.
交通工具主要有轎車.
然後就有人力車 有馬車.
但是什麼人才會用這些交通工具呢?.
通常是殖民地的官員和富有的洋行大班.
一般香港市民是用什麼交通工具呢?.
就是自己的雙腳.

$^{41}$我們再想想在耶穌年代 遠古的時代.
在耶穌年代的人是用什麼交通工具呢?.
多數人都是走路的.
在聖經裡面的記載.
我們都見到耶穌和他的門徒經常都是用自己雙腳走路.
但是剛才你讀福音經課的時候.
你又讀到在那個片段.
耶穌好像殺友戒士那樣去安排代步工具.
究竟所為何事呢?.
我們一起看看福音經課.
剛才我們讀過的經文.
看看究竟這麼特別的交通工具的安排又是什麼事情呢?.
我們看看第21章第一節.
耶穌和門徒將近耶路撒冷.
到了白髮旗在橄欖山那邊.
耶穌就打發兩個門徒對他們說.
去對面的村子.
你見到一匹驢 山在那裡 綁在那裡.
還有驢驢一齊 你解開就牽到我這裡來.
如果有人問 你就告訴他主要用他.
這樣就可以了.
於是第六節 門徒就照耶穌吩咐的去做.
就將驢和驢驢驢帶了來.
將自己的衣服鋪在上面.
耶穌就騎上去了.
你見到這個畫面.
這個真是一個安排.
是一個特別的安排.
它不是自動出現的.
其實當時除了耶穌之外.
還有那些門徒.
在往耶路撒冷的路上是絡繹不絕的.
因為那是一個很大的節期.
那個節期是猶太人的逾節.
很多人就去耶路撒冷過節.
一般平民百姓.
其實都是像耶穌和門徒一樣.
去步行前往耶路撒冷.
現在他們已經來到耶路撒冷很近的地方.
他們到耶路撒冷其實只不過是多走一會兒就到了.

$^{81}$為甚麼在這個時候要去找一個助騎呢?.
是不是耶穌那時候已經體力不支?.
是不是耶穌很趕時間要快點去到呢?.
看來也不是.
這一場是特別安排的事情.
它可以是甚麼呢?.
它可以看作一場進城的儀式.
進城儀式是甚麼呢?.
在古代進城儀式.
其實是一個很強烈的象徵意義.
其實不是古代 現在也是.
很強烈的象徵意義.
古羅馬的進城儀式最著名的就是海船式.
是給戰勝的將領和皇帝最高的軍事和宗教榮譽.
他們身穿神聖的服飾.
坐著金色的戰車率領著軍隊.
戰利品和俘虜穿過城門到神廟.
展示他們的權力和慶祝勝利.
和古羅馬的海船式進城儀式相比.
你想像一下 耶穌坐的不是金色戰車.
而是一匹民間運輸用的驢.
耶穌率領的不是軍隊.
只不過是一班門徒.
在路人的眼中.
這些算甚麼進城儀式呢?.
有趣的是.
有些路人都樂意配合去演好這場戲.
所以在二十一章第八節 第九節你會看到.
有些人將衣服鋪在路上.
有些人在樹上砍下樹枝鋪在路上.
又有人高叫「和撒拿歸於大衛的子孫.
奉主名來的應當稱頌」.
「高高在上 和撒拿」.
馬太福音記載這件事到底要表達甚麼意思呢?.
相信大家都留意到.
馬太福音從一開始預告耶穌出生的時候.
就已經用舊約經文來解釋事件.
這裡亦不例外.
二十一章第四節.
作者這樣說.

$^{121}$「這是成就 要應驗先知的話說.
要對錫安的居民說.
看啊 你的王來到你這裡.
是溫柔的 又騎著驢.
就是騎著驢的驅子」.
這位先知是誰呢?.
這位先知就是撒迦利亞.
這段經文引自撒迦利亞書九章九節.
經文這樣說.
「錫安的民應當大大喜樂.
而路撒冷的民應當歡呼.
看啊 你的王來到你這裡.
他是公義的 並且施行拯救.
謙謙和和地騎著驢.
就是騎著驢的驅子」.
馬太福音指出 耶穌的進城儀式.
與古羅馬的海船式進城儀式相距雖然甚遠.
但不要視為兒戲.
因為這是有古時的先知預告的.
錫安的王是騎驢進城.
不是騎戰馬進城.
但是眼前這個狀況.
看到眼前的狀況.
對於那些當時認真思考.
或者有期待的人 實在是非常困惑.
眼前看到的是.
有一個這樣的人騎著驢.
然後有一群門徒就在那裡拍手進城.
與現實上的需要實在相距得太遠了.
現實的需要是什麼呢?.
在羅馬帝國的鐵蹄之下.
在官府和地方勢力的高壓管治之下.
騎驢進城的王有什麼用呢?.
騎驢的人.
無論他是多麼謙和 多麼公義.
他都只是御若冰的秀才.
可以有什麼作為呢?.
根據華人基督徒聖地和平網絡.
在2026年3月23日的報導.
以色列總理萊塔尼亞胡.

$^{161}$在記者會回應問題的時候.
暗諷有些人天真 自欺欺人.
以為光靠道德和正義就能夠打敗敵人.
他援引歷史學家杜蘭特.
在《歷史的教訓》一書裡所寫.
指出歷史證明遺憾且不幸的是.
耶穌基督並未比成吉思汗強.
因為若你夠強大 夠冷酷 夠強勢.
邪惡終將戰勝善良 侵略終將戰勝克制.
2000年後 那個地方的人.
有人是這樣看的.
其實2000年前那個地方的人.
都可能有些人是懷著類似的想法.
當他們知道耶穌只是一位.
從加里尼拿撒勒手無寸鐵的先知.
或許都希望他快點做完這場戲.
就乖乖地做一個離過節的普通人.
免得節外生枝.
如果你沒有這樣的能力.
就不要搞大龍鳳了.
畢竟這是猶太人的大節日.
人們是來慶祝節日的.
不是來受苦的.
是來向耶路撒冷頌讚的.
不是來哀悼的.
現實是.
就算耶穌這位手無寸鐵的先知.
他都遭遇到嚴刑峻罰.
被誘殺錯 沒放過這樣對待.
為何這位騎驢進城的謙和之士.
都被列在被殺之列呢?.
他沒有戰馬 也沒有軍隊.
他怎樣能夠威脅到當權者呢?.
會不會只是因為他在暴力當道的世界裡批評暴力.
呼籲謙和才是拯救之道呢?.
會不會只是因為他在聖殿被異化的宗教狀況之下.
拆除綑綁心靈的律法主義.
帶領人歸向上主呢?.
會不會只是因為他的應許的拯救和自由.
促成了宗教和帝國的合作.

$^{201}$而被判為罪人呢?.
這個人到底犯了甚麼事.
以致被本族出賣.
被羅馬兵釘上十字架呢?.
如果他沒有犯罪.
好人為何會遭遇苦難呢?.
好人遭遇苦難的確是令人傷感.
和不解的千古疑問.
在災難裡.
很多受災的人.
其實根本沒有做過甚麼壞事.
在戰爭裡.
很多平民百姓根本沒有犯過甚麼事.
但他們就是最受苦的人.
為何好人會遭遇痛苦呢?.
耶穌就是其中一個二人受苦的例子.
今天是大齋期最後一個主日.
亦是聖周的開始.
聖周又稱「受難周」.
在基督教的傳統裡.
是復活節之前那一個星期.
用來紀念耶穌受難.
每一年都有這個紀念受難的日子.
原來教會從來沒有輕視好人遭遇苦難.
這些反常的狀況.
亦沒有忘記那些為義受苦的無辜者.
今天的經課裡.
以賽亞書指出.
受苦不一定是犯錯.
受苦不一定是有罪.
受苦可能只是他衷心地活出.
聽從上主 活出作為上主僕人的樣式.
我們看看以賽亞書怎樣說.
第五十章第五節.
主耶和華開通我的耳朵.
我並沒有違背也沒有退後.
人打我的背我任由他打.
人拔我腮下的鬍鬚我任由他拔.
人辱我套我我並不掩面.
主耶和華必幫助我.

$^{241}$所以我不徒愧.
我硬著臉好仗堅石.
我也知道我必不自蒙羞.
稱我為義的與我相近.
誰與我爭論.
可以與我一同站立.
誰與我作對.
可以就近我來.
主耶和華要幫助我.
誰能定我有罪呢?.
怎樣面對這種好人都要受苦的逆境呢?.
怎樣可以面對呢?.
怎樣的人才能夠面對辱罵而不抱愧呢?.
怎樣的人才能夠面對攻擊.
仍然站立得穩呢?.
怎樣的人才能夠面對邪惡.
而自己不變做邪惡呢?.
這些人看似很軟弱.
但其實不可能是軟弱的人.
事實上堅強的人才能夠不被逆境摧毀.
有實力的人才能夠謙和.
行得正,站得正的人才能夠站立得穩.
今天的《書信經》給腓立比書.
保羅指出.
手無寸鐵的耶穌.
其實就是這麼堅實的人.
我們試試看在腓立比書裡面.
保羅怎樣說.
「他本有上帝的形象.
不以自己與上帝同等為強奪的.
反倒虛己.
取悅奴僕的形象成為人的樣式.
既有人的樣子.
就自己卑微,存心,馴服.
以至於死.
且死在十字架上.
所以上帝將他升為至高.
由次給他那超乎萬名之上的名.
叫一切在天上的,地上的和地底下的.
因耶穌的名無不屈膝.

$^{281}$無不口稱耶穌基督為主.
小榮耀歸於上帝」.
耶穌騎驢驅進城而不是騎戰馬.
不是因為他沒有能力.
呼召戰馬來.
只是因為他自己卑微.
耶穌被釘死在十字架上.
每個人都說:這麼厲害.
救救自己吧.
耶穌不是沒有能力救自己.
而是因為他不以自己與上帝同等為強奪.
其實.
一個人.
他有.
他才有得放下.
他有.
他才可以放下一些東西.
沒有.
是沒有東西可以放下的.
耶穌自我降卑.
是顯示他是真正的強者.
他是真正有實力的人.
耶穌為什麼選取騎驢進城呢?.
為什麼以溫和謙卑的態度出現呢?.
甚至被釘十字架而不反擊呢?.
馬太福音引述撒迦利亞書的一句說話.
的下文是這麼說的.
撒迦利亞書第九章第十節.
「我必除滅以法連的戰車和耶路撒冷的戰馬.
征戰的弓也必除滅.
他必向列國講和平.
他的權柄必從這海管到那海.
從大河管到地極」.
這位騎驢的王宣揚的是和平.
不是戰爭.
他認為和平不是來自軍備競賽.
相反是透過裁軍帶來的.
平安不是透過武力得到的.
而是放棄武力才有的.
不過呢.

$^{321}$如果有一位武林高手.
他自廢武功.
他會不會被欺負呢?.
如果上帝的聖子.
和上帝同等的聖子.
取了奴僕的形象.
他會不會被苦待呢?.
在這個世界裡.
這樣的遭遇是非常大的.
而且難免流血.
撒迦利亞先知繼續這樣講九章十一節.
「錫安,我因與你立約的血.
將你中間被擄而囚的人.
從無垂的坑中釋放出來」.
縱使流血也唯有如此.
才能夠帶來真正的釋放.
相反來說.
如果每一個人都武裝起來.
帶來的只是更嚴重的擔憂和懼怕.
距離釋放和自由.
只會越走越遠.
還記得保羅在腓立比書中對信徒的勸勉嗎?.
第二章五節.
我們剛才讀過的最後一節.
「你們當以基督耶穌的心為心」.
然後就開始講耶穌的降臂.
今天是終之主日.
到下個主日 復活主日.
讓我們在這段時間裡面.
默想耶穌基督最後一周的受難.
死亡和復活.
透過為門徒洗腳的降臂.
然後去到赫西瑪利泉的被捕.
被釘上十字架的痛苦.
進入墳墓裡面.
你見到耶穌一步一步地降臂.
對於永遠追求升級的現代人來說.
耶穌的路途是很奇怪的.
但祂是人類的福祉至關重要.
為何降臂之路都會重要.

$^{361}$和能夠造福人群呢?.
讓我們以這個例子來作結.
還記不記得十多年前.
有一套真人騷的電視節目.
叫做《窮富庸大作戰》.
已經是十多年前了.
其中一位參加者叫做田北辰.
他曾經堅信.
如果你有鬥志.
即使是弱者.
你都可以變成強者.
在節目裡面.
他化身為一個普通的清潔工人.
靠每一天工作賺取生活費.
並且住在豪華的籠屋裡面.
徹底體驗普通香港清潔工人的人生.
他做了兩天.
如果你有看的話.
你記得.
他已經受不了.
他說很淒慘.
沒有意思.
有點帳等日子過.
很無奈.
沒有盼望.
最後他接受訪問時.
他的總結是這樣說.
在這兩天裡面.
他在想甚麼呢?.
就是只是在想吃東西.
一心的盼望.
就是希望有下一餐.
他說他終於明白貧窮人的心聲.
怎會想計劃呢?.
怎會想將來呢?.
最重要的是下一餐怎樣解決.
於是他不禁問.
當市場經濟到了盡頭.
到底還有沒有公義呢?.
如果他沒有參加過這個活動.

$^{401}$沒有這個體驗.
他不會問這個問題.
有一位參加者在體驗過清潔工的勞動生活之後.
那時候他每一晚捲曲在不夠二十呎的床位.
然後他經歷了這個體驗之後.
他辭去了那份珠寶公司行政總裁的職位.
跟社福機構合作展開了扶貧的工作.
社會認為升級才是值得追求的.
耶穌的生命告訴我們.
有時自我降級是能夠拯救世界的.
窮富庸大作證的例子也顯示.
有能力的人.
如果願意降級.
願意輸尊降貴地體驗人間疾苦.
就能夠成為其他人的幫助.
也令社會多一點溫情.
願上主幫助我們.
\newpage



\section{馬太福音 27:45-50}
\label{sec:031880}
\textbf{我的上帝!我的上帝!為甚麼離棄我?}~~~~石偉漢牧師
\newline
\newline
連結: \href{http://www.church.com.hk/acms/content.asp?site=cdc\&op=show\&type=product\&code=031880\&layout=sermon}{\texttt{031880}} ~~~~ 語音日期: 2026-03-29
\newline
\newline
\hyperref[sec:031879]{\small{< < < PREV SERMON < < <}}
~
\hyperref[sec:index_chronic]{\small{[返順時目]}}
~
\hyperref[sec:031881]{\small{> > > NEXT SERMON > > >}}
\newline
\newline
馬太福音 27:45-50
\newline
\begin{longtable}{cl}
\hline
\hline
章節 & 經文 (和合本修訂版)\\
\hline
27:45 & \begin{tabularx}{0.7\textwidth}{X} 從正午到下午三點鐘，遍地都黑暗了。 \end{tabularx} \\ \\ \relax
27:46 & \begin{tabularx}{0.7\textwidth}{X} 約在下午三點鐘，耶穌大聲高呼，說：「以利！以利！拉馬撒巴各大尼？」就是說：「我的神！我的神！為甚麼離棄我？」 \end{tabularx} \\ \\ \relax
27:47 & \begin{tabularx}{0.7\textwidth}{X} 站在那裡的人，有的聽見就說：「這個人呼叫以利亞呢！」 \end{tabularx} \\ \\ \relax
27:48 & \begin{tabularx}{0.7\textwidth}{X} 其中有一個人立刻跑去，拿海綿蘸滿了醋，綁在蘆葦稈上，送給他喝。 \end{tabularx} \\ \\ \relax
27:49 & \begin{tabularx}{0.7\textwidth}{X} 其餘的人說：「且等著，看以利亞來不來救他。」 \end{tabularx} \\ \\ \relax
27:50 & \begin{tabularx}{0.7\textwidth}{X} 耶穌又大喊一聲，氣就斷了。 \end{tabularx} \\ \\
[1ex]
\hline
\hline
\end{longtable}
$^{1}$家中的弟兄姊妹平安.
讓我們一起有個祈禱.
多謝你讓我們真的去聽你的說話.
讓我們思考主在這星期為我們受苦的意義.
幫助我們在這星期更加明白主的愛.
明白主對我們的拯救.
願主你幫助我們.
讓我們今天開始在這星期讀多些聖經.
去了解耶穌你自己的受苦和復活.
多謝你聽我們祈禱奉耶穌名字而求.
阿們.
今天和大家去思考一段經文.
來自《馬太福音》27章45至50節.
題目是《我的上帝,我的上帝為什麼離棄我?》.
為什麼選擇這個題目呢?.
因為今天按照教會的年齡.
今天是一個中之的主日.
正好是今天是耶穌光榮進入耶路撒冷的日子.
在當時群眾歡呼.
舉起中樹枝.
鋪在地上.
讓耶穌騎著驢驅進入耶路撒冷的日子.
在這個日子裡面是一個很重要的最後一星期.
特別在聖經裡面有很多重要的描述.
都是在最後這個星期裡面.
事實上在《馬太福音》.
耶穌進入耶路撒冷是在第21章.
但一直說到耶穌受死就去到第27章的段落.
所以你想想是很長的篇幅.
就是說耶穌最後的一個星期.
所以我們今天也去思考一下其中一個的片段.
在這個受難週裡面.
最後一個星期裡面.
其中有些人將祂按照不同的日子.
去說出耶穌一些重要發生的經歷.
包括了剛才我們說的有中樹枝.
歡迎耶穌進入耶路撒冷.
也有潔淨聖殿.
有奉獻銀錢等等.
也有沒有特別技術的安靜的日子.

$^{41}$甚至我們知道很熟悉的.
就是最後的晚餐就是在星期四的晚上.
星期五我們普遍認為受難的日子就是星期五.
然後有一天的安息日.
什麼事情都不能做.
到星期日的清早我們見到耶穌復活了.
所以這個星期我們都會在教會裡面的不同的活動.
包括了今天星期五我們有受苦節的晚會.
在下星期日我們有戶外崇拜.
一個復活節的崇拜.
我們都去紀念我們的主.
而今天《簡約經文》集中在耶穌受難的日子.
就是那個Friday.
就是星期五的日子.
一般我們叫Good Friday.
那個Friday星期五發生什麼事呢?.
特別我們今天讀的經文.
就是在第27章45至50.
就是說耶穌要受死了.
他斷氣之前的一個時刻,一個描述.
讓我再讀一次這一段經文.
讓我們去加深我們的印象.
待會我讀講解的時候.
大家就容易一點記得明白了.
45:從正午到下午三點鐘遍地都黑暗了.
約在下午三點鐘.
耶穌大聲高呼說.
「以利以利拉瑪撒巴各大利」.
就是說我的上帝我的上帝為什麼離棄我.
站在那裡的人有的聽見就說.
這個人呼叫以利亞.
其中有一個人立刻跑去拿海綿招滿了曹.
綁在盧偉鶴谷旱上送給他河.
其餘的人說且等著看以利亞來不來救他.
耶穌又大喊一聲氣就斷了.
這段聖經裡面就說一個時間.
就是我們所認定的一個星期五.
一個Good Friday的日子.
這段經文開始的描述就是說.
從正午就是十二點的時刻.

$^{81}$到下午三點鐘的時刻.
所以你會發現這個段落其實是三個小時的事情.
三個小時是很漫長的.
但是這三個小時是發生什麼事呢?.
其中經文告訴我們.
遍地都黑暗了.
我不知道你有沒有想過.
在日光日白的白天裡面.
整個天都黑了.
會有什麼感覺.
晚上黑漆漆很容易想像.
特別是萬家燈火準備好燃點燈.
預備迎接黑漆漆的晚上.
所以黑夜對你來說也不會覺得很黑.
但是在白天正常的時間.
中午的時間應該是烈日當空的日子.
到了下午三點鐘竟然是遍地黑暗.
氣象氣氛都會很不同.
你會感受到很龐大的壓力.
或者很幽暗的角度.
所以如果你記得你們讀經運動裡面.
讀到出埃及記不記得有這一幕.
記不記得是第幾災.
黑暗籠罩黑暗的災是在第九災.
所以在出埃及記當第九災全地都黑暗的時候.
對於法老來說是很震撼的事.
整個天都黑了什麼都看不到.
特別是在整個正常的日子裡面不能運作.
所以有一種帶來審判好像末日的意味在裡面.
黑暗之災之後下一個是什麼災.
殺長子之災十災裡面.
遍地黑暗之後就殺了他的長子.
現在在這個畫面裡面.
十二點至三點遍地黑暗之後下一幕是什麼.
耶穌這位獨生子就要為我們死.
所以都頗有連貫性有個相關在當中.
就是在這個時刻裡面.
大家沒有遇到帶燈出街的環境.
不會像現在這樣帶著手提電話帶電筒出街的時候.
在這個黑暗的環境裡面.

$^{121}$那些在十字架旁邊的人.
他們就聽到耶穌說話.
聽到耶穌叫.
你可以想像一下.
可能耶穌的面孔都不是看得很清楚.
因為太黑了.
但是他們知道耶穌在哪個位置.
因為祂就釘在十字架上面那個位置上.
而在上面裡面耶穌就發出一種呼聲叫聲.
在這個黑暗裡面耶穌就說.
以利以利拉瑪撒巴葛大利.
聽完之後不知道你有什麼感受.
如果你在當時裡面其中一個在十字架下面的人.
你聽到的時候會想起什麼呢.
以利以利其實以利以利就應該是屬於一些希伯來文.
希伯來他們的語言就是以利以利.
其中就是拉瑪撒巴葛大利.
可能經過一些研究就應該是屬於亞蘭文.
你記得我們說過在聖經裡面.
他們一直以色列人傳統的語言就是希伯來的語言.
但是他們去到他們坊間流行說的話裡面.
他們就會用了亞蘭文.
而另外在國際溝通上就用了希臘文.
幾種不同的語言.
而這裡耶穌就說以利以利用了希伯來文.
然後拉瑪撒巴葛大利就是亞蘭文.
當他呼叫大喊的時候.
你可以想像一下下面的人聽到些什麼呢.
因為你見到在上面喊下來.
特別是耶穌確實為我們受苦流血.
可能身體都疲弱了.
但是當他喊出來的時候.
下面的人突然間聽到聽錯了另一個字.
聽錯了什麼字.
以利亞.
可能就是這樣.
我想像一下耶穌就是這樣.
當他喊完以利之後.
就馬上呼一呼氣.
以利.

$^{161}$可能是這樣.
所以樓下的人就聽到什麼.
以利亞.
所以原來當時我們見到經文的描述.
就告訴我們樓下的人.
就是十字架下面的人.
就似乎他們不是聽到以利以利.
而是聽到以利亞.
以利亞.
他們就聽到耶穌喊以利亞.
以利亞.
喊以利亞做什麼呢.
你想一下.
喊人都有原因的.
什麼事可能就會喊一些什麼人.
特別是一些出名的人.
我什麼時候喊Rainbow.
可能有些銀行的東西不明白.
就喊Rainbow.
有些事我就會喊.
如果某些人在下面聽下去.
他們聽到以利亞的時候.
想像著耶穌喊他做什麼呢.
因為你記得勝經裡面描述到.
以利亞是其中一個.
聖經裡面沒有描述他死亡的人.
甚至乎你見到他升天的時候.
是旋風帶他升天的.
你見到的畫面就是.
以利亞是一個很重要的舊約人物.
甚至乎你記得.
我們說理賽亞再來的時候.
其中哪一個就要先來.
就是以利亞.
馬拉基書裡面第四章.
你見到以利亞.
所以以利亞的來臨.
是有一個很重要的意義.
甚至乎在猶太人的眼裡.
就是這樣的.

$^{201}$猶太人的眼裡.
以利亞是什麼呢.
是末日裡面一個重要人物.
他們認為.
以利亞會在一個異人遭遇困苦的時候.
就會從天而來答救他們.
因為知道以利亞被接上天.
所以他們根據猶太傳統.
就會知道.
以利亞會在一些有需要的人.
特別是一些異人.
他遭遇困難的時候.
當他們呼求以利亞.
以利亞就會救他.
所以你可以想像.
在十字架下面的那群人.
當他們聽到.
以利亞 以利亞.
為什麼離棄我的時候.
就好像.
他們就以為.
耶穌就好像這幅圖畫一樣.
耶穌好像在游水一樣.
就是你現在淹死了.
淹死的時候.
在水裡面就喊.
以利亞 以利亞.
扔個水泡來救我.
類似這樣.
當然在這個畫面裡面.
似乎就是以利亞沒有扔過水泡.
所以就說.
你為什麼還走開.
好像一個這樣的畫面.
原來在下面的人.
似乎就是這樣拿著.
耶穌所說的話的意思.
他們以為耶穌在求救.
向以利亞求救.
以利亞來到救他這個異人.

$^{241}$以致免去他受這個刑罰.
所以他們在47-49節.
我們看到就是.
在十字架下面站著的人裡面.
有一些人聽到.
這個人在呼叫以利亞.
而其中有些人就立刻跑去.
拿海綿去沾醋.
綁在路上遞給耶穌喝.
而還有一些人就.
我們等一等.
就看看以利亞來不來救他.
因為他們就期望.
耶穌似乎在呼叫以利亞.
那耶穌說到自己這麼厲害.
那以利亞是不是會來救他.
這群人就在這裡等待著.
但我們今天第一個分題.
就是在說什麼呢.
旁人看耶穌.
這群旁觀者.
其實站在那裡的那些人.
其實是什麼呢.
他們誤解了耶穌.
他們拿一些醋.
倒給耶穌的時候.
其實是想戲弄耶穌.
甚至其實他們是不信耶穌.
這群人他們對耶穌.
整個的行動帶來一些誤解.
不知道今天你有沒有對耶穌.
這一句話都有誤解呢.
有沒有啊.
我們分開教會這麼久.
可能我們都會.
誤解是什麼呢.
耶穌很慘.
耶穌在十字架上真的很慘.
我們無可否認.
耶穌在十字架上.

$^{281}$我們釘身流血.
帶著荊棘的冠冕.
是很慘的一個畫面.
但耶穌是不是在這個時候.
大家不知道祂很慘.
就在那裡埋慘呢.
在那裡埋慘就去告訴別人.
現在我很慘了.
連神都不要我了.
我們會不會有時候都這樣想呢.
因為有時候我們讀聖經.
我們都會這樣想.
特別很多時候當我們覺得.
神不在我們身邊的時候.
我們覺得好像.
神啊你在哪裡.
神你在哪裡的時候.
我們經常說.
神啊你會不會離我很遠.
神啊你會不會都離棄我.
我們甚至將耶穌說過這句話.
投入在自己身上.
就是耶穌都覺得神不要祂.
所以有時候我們都覺得自己很慘.
神不要我們.
弟兄姊妹.
神會不會不要我們呢.
神會不會離棄我們呢.
如果我們知道神不會離棄我們的話.
那神會不會離棄耶穌呢.
如果你的答案是不會的話.
那耶穌這句話在做什麼呢.
耶穌是不是真的覺得自己很慘.
以致告訴別人.
是神離開了我呢.
今天我們就去想這個問題.
在十誡上面.
耶穌說了這句話.
究竟是什麼意思呢.
我們會不會就好像剛才.

$^{321}$十字架下面那群人誤會了耶穌.
是在求救.
耶穌不行了.
耶穌死定了.
耶穌要呼求 求拯救 求以利亞幫助他.
是不是這樣呢.
求主幫助我們.
在我們的第二部分.
我們去明白多一點.
耶穌在這個十字架上面的情景.
所以第二個的分題.
就是明眼人看耶穌.
明眼人看耶穌.
就是我們怎麼去重新去.
用一個新舊約穿插的角度.
和他們去看看究竟.
耶穌在做什麼呢.
耶穌說這句話是想做什麼呢.
第一個B的分題就是.
禱告交託就一身輕.
當我們懷疑耶穌是不是很慘.
很辛苦很痛苦痛苦到叫救命的時候.
我們都要問.
究竟現在耶穌是不是身負重壓.
重壓到一個地步一定要求救.
這個是不是他的求救聲音.
還是說其實他已經能夠.
昂然踏上這個十字架的路呢.
耶穌上十字架是不是被逼的.
不是的.
耶穌甘願上這個十字架.
但是我們怎麼去看.
他在這個十字架上面的呼求呢.
第一個我剛才說到看的是什麼呢.
那個禱告交託.
禱告交託.
你記得我們一直說馬太福音.
說登山寶訓.
我們說過耶穌祈禱的了.
是不是.

$^{361}$耶穌的祈禱通常叫神做什麼.
符是不是.
你記得嗎.
耶穌教我們怎麼去叫神呢.
他通常叫我們跟神祈禱叫符的.
是不是.
是符來的.
所以我們知道我們說.
當我們說到我們平時說主禱文.
我們說那個主禱文其實不是主禱文來的.
是主教導我們祈禱的禱文.
其實耶穌不會這樣祈禱的.
所以你記得耶穌是怎麼祈禱的呢.
就是在前一晚.
在前一晚在那個我們所謂的.
用一些理解就是星期四晚上.
就是在赫西瑪利園裡面.
耶穌就真的有他的禱文.
我們在赫西瑪利園裡面.
真的有耶穌的主禱文.
耶穌的主禱文裡面他怎麼去祈禱呢.
我們看到耶穌的祈禱就是說.
第一次他說我父啊.
他稱神為什麼.
我父啊.
我的爸爸.
如果可能求你將這個杯離開.
然而不要照我所願的.
是照你所願的.
我們確實知道耶穌在當下那一刻.
完了這個最後晚餐之後.
他在這個赫西瑪利園的禱告裡面.
確實顯出那種禱告的迫切.
和那種的掙扎.
因為他知道關鍵的時刻將要來臨.
猶大將要帶兵丁來到這個園裡面.
要捉拿耶穌.
所以耶穌就在這裡發出一個迫切的禱告.
說我的父親啊.
如果可以就怎樣.

$^{401}$但是我會跟著你的意思.
第二次耶穌的祈禱是什麼.
他也是說我父啊.
這個杯如果不能夠離開我.
我必須要學.
我就願你的旨意成就.
他仍然認定了那位是他的父親.
以至到他願意天父的旨意成就.
第三次他再去這樣禱告.
而經文告訴我們什麼.
說的話和先前一樣.
就是說他仍然呼叫那位的神是什麼.
我的父啊.
原來在那個星期四的時刻裡面.
我們就知道耶穌已經經歷了一個很.
為了祂將要受苦的一個很重要.
迫切的禱告已經作出了.
已經為這件事將要上十字架.
將要受死這件事做了一個很迫切.
很認真的禱告.
你有沒有想過你祈禱之後會有轉變呢.
我們經常說上一節我們經常都會祈禱.
但是我們有沒有想過我們祈禱之後.
是會有轉變呢.
那種轉變有時候不是那件事轉變.
而是你的心呢.
你的心態應該經歷一種轉變.
我們再看多一點.
就是當耶穌之前的時候.
他是什麼狀態呢.
在第26章37-38字告訴我們.
耶穌確實在下西馬聯開始祈禱之先的時候.
他說什麼.
耶穌憂愁起來極其難過.
以至到他對那三個門徒說什麼.
我心裡非常憂傷幾乎要死.
很大的心靈掙扎在裡面.
耶穌真的在經歷.
因為這個是他願意去面對一個很大的事情.
但是當你留意回去到後來.

$^{441}$當耶穌三次祈禱之後.
耶穌說的是什麼呢.
耶穌在後面第42節.
剛才我們讀到的是願你的子兒成全.
那種轉變就是由心靈裡的掙扎.
覺得很辛苦的心靈狀態.
以至到願你的子兒成全.
你明不明白那種狀況.
就是當你經歷了一種認真和交託的禱告的話.
你可以放鬆.
願你的子兒成全.
OK了.
所以耶穌說什麼.
起來我們走了.
即是說預備好了.
耶穌說他已經預備好.
我們可以去面對出賣他的猶大和那些的冰釘.
所以其實我們見到耶穌那個心情的轉化.
其實在這個赫西瑪尼的禱告裡面.
已經能夠克勝了那種.
將要面對那種的痛苦悲傷難過.
這個很重要的.
當你能夠體會到耶穌在那個禱告之後.
能夠經歷了一種轉變.
能夠預備好去上十字架的時候.
你看耶穌所說的話你就明白.
不是賣慘.
不是叫痛苦.
他是承受著痛苦.
不過他不是在不斷呼求.
告訴別人我很慘很痛苦.
耶穌不是這一種的狀態.
耶穌是什麼狀態呢.
耶穌在十字架上說了七句的話.
受苦者通常我們會思考其中一個思考方式.
就是想想耶穌在十字架上所說的七句話.
在四福音裡面.
耶穌記錄了耶穌的十加七言.
按著他的次序就應該是這樣排列.
我們試試看一看這十加七言裡面.

$^{481}$耶穌在十字架上.
如果我們猜想耶穌是很辛苦很痛苦的時候.
只是關心著自己怎樣經歷那種痛的時候.
其實耶穌是不會說這些話的.
我們試試看一看耶穌說了什麼話.
包括了耶穌說父啊赦免他們.
因為他們所做的他們不曉得.
耶穌在說什麼狀況呢.
在說下面要對付他們的人.
耶穌說赦免他們.
其實他們都不知道在做什麼.
你聽不聽到耶穌在十字架上所關心的.
不是關心自己有多辛苦有多痛.
而是下面那群人不知道自己在做什麼.
耶穌說父啊赦免他們.
耶穌還在關注下面要拯救的所有世人.
包括下面這群人.
然後耶穌第二句說的是什麼.
今天你要和我在樂園裡和誰說的.
在旁邊的十字架上其中一個釘十字架的人.
這個人你今天會和我一起在樂園裡.
因為他認定了耶穌是無辜是主.
所以耶穌這句話是關心著那個釘十架的人.
又不是關心自己有多辛苦.
第三句對著下面瑪利亞和約翰說什麼.
母親看你的兒子指著約翰說這是你的兒子.
和約翰說什麼看你的母親.
所以耶穌關心瑪利亞關心約翰.
都不是關心自己.
第四句就是我們今天讀的經文.
以利以利拉瑪撒巴哥大利.
其實是什麼意思.
就是我的神我的神為什麼離棄我.
你發不發現不同了.
有什麼不同.
耶穌的呼求有什麼不同.
怎麼叫.
我的神我的神.
耶穌平時叫什麼.
父.

$^{521}$叫爸爸或廣東話叫爸爸.
是這樣叫的.
但現在是我的神我的神.
不同了.
等下我們再解答這個問題.
所以你發現到第七句說話.
又是怎樣叫.
父呀我將我的靈魂加在你手裡.
整個過程裡面你會發現.
他在後面都不是關心自己有多辛苦.
而且他認定了那個位置是叫父.
你認定了這個原則.
等下你就容易明白我接下來想說的東西.
第五句說話我學了.
等下我會補充一點.
第四和第五句是有一個關聯的.
第六句當然是說成了.
成就了.
所以你一路看下去.
你會發現耶穌十教裡頭幾句說話.
是愛心神的愛.
關心著地上的每一個人.
而後面的說話.
其實是在說他要得勝.
他要完成這個救贖.
他要死在十字架上.
釋放世人.
其實後面說的這件事.
所以我們這樣去整理一下.
第一個我們要想一下.
我們的重點是甚麼.
禱告交託一身輕.
就是耶穌在之前的晚上裡面.
經歷了一個這麼深切的禱告的時候.
就已經能夠將他的掙扎.
完全向天父表白.
以致到神更加被他經歷到.
很明顯就是耶穌已經準備好.
耶穌的心態已經調整好.
然後才去被捉拿.

$^{561}$面對審判.
以致到被釘上十字架.
耶穌也不是說他自己很慘很痛苦.
所以今天我們的挑戰就是.
我們有禱告.
我們有沒有交託呢.
這裡說你一定有禱告的.
我們一定有祈禱的.
但我們有沒有交託呢.
我想起一個畫面.
就是大家都熟悉的東東.
東東就是Tony養的一隻狗.
一隻很可愛的狗.
一個男人和一隻老狗.
他家裡有兩樣東西.
其中這隻狗真的年紀大了.
東東年紀大了.
但是要交給其他人照顧.
交給誰照顧.
交給老婆照顧.
在加拿大照顧.
這個過程是什麼呢.
禱告交託是什麼意思呢.
就是我已經將這隻狗交給你了.
我交了給你之後.
你就不要指指點點.
因為你說了給他在加拿大照顧.
你就不在了.
你就給他全心去照顧這隻狗.
全心照顧東東.
你就讓他去做.
信得過他.
這就是交託.
所以很感恩.
東東之前說半年命都沒有.
如果沒有記錯.
現在東東還生猛跳來跳去.
跳上去Lisa那裡.
照顧得很好.
因為交託是交得對的.

$^{601}$所以禱告就是有好的禱告.
正確的禱告.
就會讓你心態有轉變.
因為你可以放心交給那個人.
而我們禱告的特點是什麼.
聽清楚這句話.
不要誤解.
我們的禱告是什麼.
所託非人.
對不對.
我們的禱告是所託非人.
因為我們的禱告是所託是神.
我們交給人.
有時候有些人都可能不小心.
但我們交給神.
就可以輕輕鬆鬆.
所以耶穌.
他的禱告為什麼要這麼迫切.
因為他確實要將他所有的東西.
交託給神.
神會幫我們去解決.
所以之後的事情.
我們可以輕鬆面對.
耶穌能夠昂然闊步.
踏上十字架.
就是這樣.
其中有一個人物.
大家可能都很熟悉.
這個很快地說一下.
就是喬治穆勒.
就是孤兒之父喬治穆勒.
其中他最出名的就是他的祈禱.
他的祈禱是怎樣呢.
他就是扮孤兒院.
他扮孤兒院裡面.
從來都不籌款.
他只是用祈禱求神.
供應他孤兒院的需要.
他一祈禱交託了之後.
就一切交給神.

$^{641}$所以他從來都不需要籌款.
為他的幾千人的孤兒院裡面.
去籌辦任何需要任何的資金.
甚至乎如果大家都聽過.
其中一個最出名的事件.
就是說他有一次.
在早上的時候.
整個孤兒院已經沒有任何食物.
他仍然帶一群小朋友.
在神的面前去借飯祈禱.
就交給上帝.
借飯祈禱之後.
結果有個麵包車.
因為車的故障的原因.
就剛剛停在門口.
車主就說.
我已經要拿車去維修了.
我的麵包快跟著車了.
就全部給他們孤兒院.
之後又有一些混奶的經過.
不知如何又全部交給他們.
結果那群小朋友又早餐又有奶喝了.
這是遠一點的.
很遠的禱告交託.
近一點的.
因為我們現在是否還可以.
禱告交託一聲輕呢.
近期在2025年12月23,24日.
在香港有一個叫.
普世平安星光夜的報道會.
不知道你們有沒有聽過.
在這個報道會裡面.
其實他們很想在接近聖誕的日子.
搞一個大型的公開的報道聚會.
希望讓很多人能夠聽到福音.
當然這個報道會聚會裡面.
他們主辦的其中一個重要的負責人.
就是王祖藍.
大家可能都知道誰叫王祖藍.
其中他們都想過會不會停了福音的聚會.

$^{681}$因為當時亦都有宏福縣的大火發生.
但當然他們後來談過.
在這個大火發生的時候.
人更加需要福音.
所以他們決定繼續去辦這個報道聚會.
辦這個報道聚會.
其中一個他們大的挑戰是什麼呢.
他們就選擇了一個新的地點.
這個地點其中特點是什麼呢.
是一個露天場合.
在一個露天場合辦一個大型的報道聚會.
其中最大的挑戰是什麼呢.
天氣.
正所謂天氣不似預期.
不一定受你控制.
其中當他們一路搞.
一路賣了很多票出去的時候.
到了那個聚會的五小時之前.
五小時之前.
場地方面就跟他們說你們可能要取消了.
因為天氣太差了.
下大雨.
他說如果你在這樣的環境之下.
辦一個戶外的音樂的報道聚會.
當中你可以想像這個場地.
所有人不單止是突襲.
甚至乎暴露在一個露天環境裡面.
都有很多可能的危險.
還要你播放一些不同的音響.
可能有些電等等的東西.
所以他們面對這件事情.
就是要面臨是否要取消.
而其中王祖林帶領他們一群的同工.
一群的參與者一起去向神祈禱.
其中王祖林說他向神祈禱是什麼呢.
他說和神祈禱.
我求你讓雨能夠停下.
兩點了.
他說兩點鐘你就要停雨了.
他就向神祈禱.

$^{721}$求神兩點鐘去停雨.
結果兩點鐘後真的停了雨.
而且接著可以很順利地繼續各樣的活動.
亦都可以預備好場地.
五點鐘後可以進來拍照.
打卡玩攤位.
之後就參與報道的聚會.
這就是其中報道聚會的其中一些照片.
弟兄姊妹就是這樣.
禱告完之後要交托.
你記住今天要記住的是什麼.
所托非人.
所托是神.
耶穌都是這樣.
所以耶穌在十字架上我們可以理解.
耶穌不是真的有那種痛苦掙扎求救.
我們就問.
耶穌說的話是什麼呢.
耶穌又不是掙扎求救.
耶穌為什麼說我的神我的神.
為什麼離棄我呢.
我們看看為什麼.
第二個重點是什麼呢.
就是萬千應驗在一身.
耶穌是不是大聲呼求求救呢.
還是耶穌在宣告神的救贖應驗呢.
如果耶穌不是求救.
為什麼要大聲喊呢.
大聲喊就是通知所有人.
我宣佈一件很重要的事情.
就是我的神我的神為什麼離棄我.
他要讓我們想起這一句話.
當然我們要在回答這句話講什麼之前.
我預備了一個遊戲.
好了我們就問一下大家.
數學可能很厲害的人.
不知道你懂不懂計算Probability.
即是或然率.
我就想問一下.
如果一個硬幣連續十次集到公.

$^{761}$那個或然率是多少呢.
懂不懂計算.
很簡單一個硬幣有公和字兩面.
你集到公的機會是多少.
二分一.
如果你再下一次集到公的機會.
就要再乘多二分一.
你集十次都是公的機會.
就是二分一乘十次.
因為每一次只有一半機會.
你有十次一半機會.
而十次一半機會每次都是中公.
就是那個或然率.
好所以我想找人出來幫我集這個硬幣.
我希望你能夠在一次集到十個公給我看.
我想找一個人去集.
好誰出來.
Tony.
好Tony你集完之後給大家看.
不然我以為我出狡猾.
你們用誰叫公.
有個人仔頭.
近一點給大家看.
你可以選擇十次公或十次字.
你選哪個.
七次公.
不是七次公.
即是十次全部都是公或十次全部都是字.
你可以選擇公或字.
即是你現在給你十次機會集.
你要十次都是公.
你就符合到這個或然率.
好我們先試試你第一次是否公.
是甚麼.
你先給Marco看一下.
即是不成功.
即是第一次已經不成功.
我們給一次機會好嗎.
現在不是第二次 剛才第一次已經完全失敗了.
你明白嗎.

$^{801}$因為你一斷纜就失敗了就不合格了.
所以現在我們正式再來第二次好嗎.
再給你一次機會.
只要你集到公我再給你機會再集公.
你集到十次都是公你就成功.
那你就可以買那些彩票了.
不是 說笑的.
好你都是選公是不是.
好你現在第一次看你集不集到公.
公 是不是公.
不要 均真一點.
待會到時候說我們出古惑.
好第一次成功了.
好第二次.
又不成功.
收公 都是公.
不過收公.
不成功嗎.
好謝謝你.
因為時間關係你回家自己練習一下.
現在我們只是輕輕集了兩件事.
即是連續兩次公都未得.
你想想要集到連續十次都是公.
你那個機會率是多少呢.
就是一萬分之九點八.
一萬分之九點八.
就是集十次.
如果你集一百次都是公.
那你想像一下應該很難很難很難很難.
是呀原來就是這樣.
我們集公字要集到全部都是這一句.
都已經很難了.
所以你想一下如果聖經裡面.
舊約聖經裡面做了一些預言.
而那些預言全部都應驗在耶穌身上.
有沒有可能呢.
如果純粹靠或然率.
靠那個probability有多大可能成功呢.
因為現在我們剛才集這個公字.
只有正和反兩個可能性.

$^{841}$如果耶穌要做一個預言的話.
而那個預言可能本身就有很多可能性的話.
只是符合一個預言.
已經有很多可能的發生.
令到不能夠符合這個預言.
而聖經裡面舊約有多少個預言.
是講到會應驗在耶穌身上呢.
好了考考大家.
多少個.
超過一千.
三百.
多少.
三百三十三.
OK成交.
奇俠根據一些人的研究做一些統計.
舊約裡面發生在耶穌的預言.
多少個呢.
就是三百個以上.
其中有些學者就嘗試統計了.
三百三十二.
恭喜你了.
你說三百三十當你中了.
這個銀子送給你.
原來舊約裡面.
真的可以歸納在耶穌出生.
事公死亡復活的事裡面.
原來有三百三十二個預言.
如果一個應驗.
我們說這麼巧.
兩個應驗.
這麼幸運.
三個應驗.
什麼狗屎撞馬棍.
撞得對.
三百三十二個.
都撞得中在耶穌身上.
那個可能性比你擲一萬次弓更難.
甚至有些人說.
不如小的預言不要計.
如果再找多一點.

$^{881}$再核心一點的.
原來都有六十至一百個.
不同的人去評估.
六十至一百個預言會發生.
所以如果應驗在舊約裡面的預言.
在一個人身上是極之困難.
不容易的.
而這件事就要應驗在耶穌身上.
而耶穌這一刻就要宣告.
應驗是在祂身上.
耶穌說.
以利以利拉瑪撒巴葛大尼.
《馬克福經》裡面說.
以羅爾拉瑪撒巴葛大尼.
或者簡單一點直接說.
整個說話就是.
我的神我的神為什麼離棄我.
其實當耶穌宣告.
大聲喊叫我的神我的神.
為什麼離棄我的時候.
原來就是引自詩篇二十二篇第一節.
詩篇二十二篇第一節.
所以耶穌這個呼叫.
為什麼祂不是喊父啊.
因為祂不是想說.
和父的關係和禱告.
而是祂要讓我們想起詩篇二十二篇.
你明白我的意思吧.
《十誡七言》第一句就說父啊.
《十誡七言》最後一句也是說父啊.
祂和神禱告的關係.
是用父啊去表達的.
現在這一刻.
祂喊的時候.
祂不是喊父啊父啊.
為什麼離棄我.
不是用祂和祂的關係去表達.
而是用我的神我的神.
為什麼離棄我.
因為原來詩篇二十二篇.

$^{921}$是叫彌賽亞的預言詩.
這首詩歌就是預言著彌賽亞將要降臨.
彌賽亞就是基督.
彌賽亞就是那位救世主.
詩篇二十二篇已經預言了.
最多段落是描述到.
耶穌要來到成為那位彌賽亞.
成為那位拯救者.
其實耶穌在十字架上面.
就是大聲宣告這個詩篇.
這個預言就是由我去宣告.
就好像當天的大衛要宣告這個詩篇一樣.
所以當你從這個角度你就明白.
耶穌並不是很辛苦地大聲喊這句話.
而是耶穌在這裡大聲喊的時候.
就是宣告這個預言.
這個彌賽亞之歌要在我身上去實現.
其中這首詩歌裡面有很多段落.
我們不能夠看完詩篇二十二篇.
不過我們看其中一些經文重點.
你就會看到那些描述.
包括了第二十二篇第一節.
剛才我們讀過很多次了.
就是我的神我的神為什麼離棄我.
就好像耶穌正在宣告.
祂就是彌賽亞.
接著在二十二篇的第七至八節怎麼說呢.
他說凡看見我的都癡笑我.
他爸自己答交託給耶和華.
記不記得我剛才說.
耶穌就禱告交託了.
而這個詩篇裡面正正好就是這個詩人大衛.
他都禱告交託了給耶和華.
他說讓耶和華救他爸.
耶和華既喜愛他可以答救他爸.
有些人就是這樣去攻擊耶穌.
他說有人會救他.
有人會救他.
這個正正好就是在.
在瑪利福音27章41至43節的時候.

$^{961}$那些祭司那些民事路人.
都在十字架上面戲弄耶穌.
他說他靠神神若喜悅他.
現在可以救他.
原來這個詩篇裡面七至八節這個恥笑.
正正應應在耶穌身上.
十字架下面的人.
就是這樣去恥笑耶穌.
十四十五節他說.
我如水被倒出.
我的骨頭都脫了節.
我的精力枯乾如同瓦片.
我的舌頭緊貼上顎.
你將我安置在死灰中.
這裡面正正好就是在說.
那種骨頭脫節的情況.
和口裡面充滿饑渴的一個狀況.
而這個狀況正正好就是描述到.
耶穌被扎他的手和腳.
在十字架上被釘上的時候.
他面對著這種痛苦的狀況.
預演了他那種身體的特徵.
所以在第十到十六節.
都有這樣的情況.
接著說到第十八節.
他們為我的內衣抽籤.
記不記得那一幕.
兵兵將耶穌的衣服抽了出來.
抽籤然後分了衣服.
在詩篇第十八節.
原來也有說這個的地方.
所以原來這個詩篇二十八篇裡面.
一句又一句是預言彌賽亞將要怎樣來臨.
而耶穌正正好宣告著.
這首詩篇裡面的預言.
全部都在耶穌身上能夠得到應驗.
最後一句.
待會我們會回頭再看到這句.
但是我們先讀了.
他們的四極都要賞念耶和華.

$^{1001}$並且歸順他.
列國萬族都要在你面前敬拜.
這首詩篇到了最後的時候.
並不是說有多慘.
而是說有一天之後他要得勝.
他要得勝的時候.
地的四極.
即是五湖四海四面八方的人.
都要紀念神.
而且歸順於神.
列國要在他面前敬拜.
耶穌就要成為那位敬拜的主.
所以到第27章54節的時候.
在下面當耶穌死了之後.
在上面斷氣的時候.
下面的八夫長.
他說和他們一同看守耶穌的人.
看見了地震.
看見了所經歷的事.
就說了一句話.
他真是上帝的兒子.
耶穌宣告他就是神所派來的彌賽亞.
而在下面的人看見耶穌斷氣之後.
看見耶穌所成就的.
他們想起耶穌真是神的兒子.
所以這個應驗表明了.
上帝對我們的拯救.
不是一個偶然發生的事件.
因為有些人不相信耶穌.
就說耶穌的死是不是計劃B.
就是說之前計劃A不行.
接著要用計劃B.
不是的.
在舊約裡面已經講定整個計劃.
整個計劃已經很清楚.
我們有位救贖主要來到我們當中.
耶穌又是一個受難的男徒.
而在舊約裡面已經清楚表明.
而在新約裡面就應驗著這個男徒的出現.
今天神都要為你定下生命的男徒.

$^{1041}$你的生命裡面是交在神的掌管和計劃之中.
神不需要定下計劃A計劃B計劃C.
幾十個計劃.
神對你生命只有一個計劃.
一個生命計劃.
上帝都要在你生命裡面有一個男徒.
當然你的男徒不是一個受難的男徒.
但是他就是一個生命的男徒.
而這個男徒要在上帝的帶領之下.
要一一去實現.
今天我們要認真去問一下.
What is your plan?.
你自己有什麼計劃?.
你可能結婚生子.
存錢買房子.
退休.
很多不同的計劃.
但是今天要問你更加重要的問題.
What is God's plan?.
神在你身上的計劃是什麼?.
你容不容忍你發現上帝在你生命裡面的計劃.
以致你能夠按著上帝在你生命裡面的計劃男徒.
去活出來.
上一次Jenny都說到.
但以理都有個計劃.
不過這個要重溫.
記得嗎?.
上次教我們.
但以理我們都看到神策.
用神的策去想起.
如何在艱難裡面.
在解夢裡面求神去幫助.
包括Pause.
暫停一停.
Link.
聯繫一群人和你去禱告.
Adore.
宣告我們的主我們的神.
Nominate.
向人去介紹這個神.

$^{1081}$是真是我們的主我們的神.
這個值得我們去想想上次的計劃.
不過今天就要想一想.
上帝在你生命裡面有個怎樣的計劃.
求主幫助我們.
在受苦之時想一想.
上帝一早預備一個救贖計劃.
上帝亦都在你未出母胎的時候.
已經為你預備一個人生美好的計劃.
信耶穌是其中一步.
生命成長是另外一步.
不要只是計劃完成了一部分就滿足.
我們應該要完聖所有計劃.
完成上帝的計劃.
上帝要你做什麼呢?.
傳福音?做見證?.
要去服侍人?愛人?愛神?還是什麼呢?.
第一次會值得去想.
因為如果我們找不到我們的計劃.
我們人生活著有時候真是毫無意義.
如果你信了耶穌之後.
你都不只是相信這一步.
然後就不要再按著神給你的計劃.
去走到下一步.
浪費了很多東西.
求主幫助我們.
第三個重點留意就是.
不離不棄神仍在.
當我們讀到這段經文的時候.
我們都很懷疑一個問題.
就是天父是不是真的離棄了愛子耶穌呢?.
還是祂以沉默的同在來支持耶穌呢?.
我們讀這段經文的時候.
無論我們想不想.
耶穌是不是很慘也好.
但耶穌問這個問題都很有意思的.
祂說為什麼神離棄祂呢?.
是不是真的會有離棄呢?.
當然我們知道回應詩篇22篇.
但是神會不會離棄人的問題.

$^{1121}$是常常在我們的腦子裡.
所以這就是我們第三個要問的問題.
神是不是真的會離棄了耶穌.
還是神用一種默默的方式去支持耶穌呢?.
今天我們很多人都會常常去感受到.
神好像不在你身邊.
有一個歌手.
不知道大家認不認識.
叫做陳雷.
這個歌手其中一首在兩年前多的歌.
那首歌的歌名叫做神的不在場證明.
神的不在場證明.
正好想描述有些人在人生裡面.
面對著很困苦很艱難的時候.
好像神就不在了.
神就好像缺席一樣.
總是好像要靠自己去度過自己的生活.
自己的人生.
所以這首歌引起很多的共鳴.
因為很多人都確實好像感受到.
究竟有沒有神在我身邊.
究竟有沒有神陪伴我呢?.
究竟神是不是在呢?.
當我們看回在《馬達福音》.
由耶穌一開始受洗的時候.
發生什麼事?.
耶穌受洗的時候.
「這是天上有聲音說,這是我的愛子,我所喜愛的.」.
天上已經有聲音很清楚說.
耶穌就是我的愛子.
我最疼愛的就是祂了.
是不是?.
不是只有耶穌才叫神做天父.
只有耶穌才當祂是老爸.
是不是老爸不當祂是兒子呢?.
不是的.
天父已經認定了耶穌告訴世人.
祂就是我的愛子.
所以耶穌在登山寶訓的時候.
也告訴我們.

$^{1161}$我們在天上的父.
在《赫西瑪利人的禱告》裡面說.
「我父」.
就是說由一開始傳到最後.
耶穌都認定了祂是天父.
天父也都認定了祂是愛子.
天父和愛子在愛裡面.
是從無分離.
雖然從無分離.
但是我們也都知道.
天父當下確實要保持沉默.
要搞清楚.
沉默不代表祂不在.
有時候祂在.
但是因為某些原因.
祂確實要沉默.
不可以出聲.
每一年我們家裡都會發生一件事.
這兩三年.
這幾年都有了.
就是什麼呢?.
就是我的子女要考鋼琴試.
考鋼琴試.
接下來五月.
他又要考.
特別是我兒子會考第八級的琴試.
每一次我們做父母的.
我們都會帶同子女一起去到那個市場.
而去到那個市場裡面.
他們就在一間教會裡面.
在Ted Street那裡.
那間教會裡面.
就是舉辦這個琴試.
而那個琴就在禮堂的最前面的位置.
我們有沒有陪子女去呢?.
有啊.
但是我們可不可以進去市場呢?.
又不可以.
我們看著兒子和女兒進去市場裡面.
就是那些大的三角琴.

$^{1201}$有時候有一點遮到.
不太能看到他們的樣子.
但是我們做父母.
我們就在門縫裡.
好像後面那些.
有個小門縫.
那些玻璃.
我們父母就在外面.
不停地裝一下.
我們不能進去.
但是我們在那裡.
我們在現場.
我們裝一下.
我們怎麼能夠知道他考得好不好呢?.
我們用耳朵聽.
因為琴聲是會傳到出來的.
你聽到.
有一點刺刺的.
你知道糟糕了.
你聽到.
很順.
今天很有狀態.
我們父母就在外面門內.
我們不可以出聲.
但是我們在那裡.
在那裡默默地支持他.
默默地為他祈禱.
耶穌在這一刻.
他要上十字架上面.
就好像一個考試一樣.
這個大考是耶穌自己要做的.
我們做父母不能代他考試.
耶穌就是要去完成這個考試.
他就要寫在十字架上.
所以「這是我的愛子」.
這位天父就在那裡看著.
他不是不在現場.
不是真的掩面不看.
只不過他要默默無聲.
有時候就是這樣.

$^{1241}$神就在背後.
背後就是他在那裡.
甚至我們面對人生的痛苦掙扎的時候.
我們以為神不在的時候.
其實神在那裡.
甚至神在為我們做一些事情.
甚至神在等待我們考完試那一刻.
考完試那一刻就宣佈.
「Pass 你成功了」.
而這個時刻就是期待著神.
讓你在痛苦的歷練裡面.
能夠見到榮耀.
讀書會不辛苦.
但考完試之後考得好成績.
是不是很開心.
父母都充滿喜悅.
同樣有時候上帝都知道我們很辛苦.
不過上帝就期待著我們能夠跨性的時候.
就讓人見到榮耀.
所以當你正面對困難的時候.
神從來都沒有離開過.
這個反而是神彰顯祂沉默的同在.
藉著你困難的時光.
才能夠發出光芒.
能夠榮耀神.
因為如果連試都不考.
你不可能告訴你成功還是失敗.
你只是逃避.
但上帝期望你考完這個試.
完成了.
你可以很昂然闊步地走出市場.
這個就是榮耀.
神今天都期望我們.
在我們人生的困難裡面.
見到上帝的榮耀.
我們不要失望.
我們不要停止向上帝去呼求.
我們不要停止和上帝的關係.
因為父仍在.
有一個人物.

$^{1281}$大家可能都很熟悉.
不過不知道這裡又熟不熟悉.
大家認不認識李聖林牧師.
有些印象.
李聖林牧師是.
之前是牧職神學院的院長.
很多年前.
我都上過他的課.
可能大家一般人認識他.
其中一個原因反而是他的兒子.
這裡其實就是說.
他的兒子在早幾年.
在一個演唱會裡面.
上面有整個電視機掉下來.
最嚴重受傷的就是他的兒子阿牧.
在這個人生的困難裡面.
我們很常見的.
可能怨天尤人.
甚至乎覺得神沒有到.
可能都會這樣的.
但對於基督教的家庭.
一個牧師的父親.
一個信耶穌的兒子.
他們就要體會到.
上帝沒有離開過.
他們仍然要深信上帝.
真的在他們當中同在同行.
以致他們能夠靠主能夠去.
走出來去告訴人聽.
他們的神仍然是真實的.
經歷這幾年裡面.
一直的醫治陪伴.
最近最新的消息是什麼呢.
這個阿牧不單止可以離開病床.
甚至乎現在他已經去到一個階段.
就可以坐著一個電動輪椅.
可以控制著電動輪椅.
去幫助他能夠出入.
這幾年裡面是很難捱的.
這幾年真的很辛苦.

$^{1321}$我們有時候看一些報道.
都會很心痛.
但是你相信上帝在這裡.
相信上帝在你的家庭裡面.
我們繼續可以靠著上帝.
因為知道神是不離不棄.
所以今天最新的告訴我們.
原來他可以行動多一點.
他可以去遠一點.
他有選擇權去避開人不見.
因為你躺著什麼都不能動.
他可以見他的人.
他可以不見他不見的人.
他有選擇.
所以這裡都說到他最新的代讀.
就是說在當中.
就是他們家裡引用的經文.
就是李成龍牧師引用的經文是什麼呢.
就是羅馬書五章三至五節.
就是在患難中也是歡歡喜喜的.
因為知道患難生忍耐.
忍耐生老年.
老年生盼望.
盼望不至於羞恥.
因為所賜給我們的聖靈.
將神的愛囂貫在我們心裡.
是的.
你相不相信神能夠幫助你.
在你最困苦的時候.
這個操練.
我們可以藉著耶穌生命的見證.
我們去學習這種操練.
所以求主幫助我們.
我們真的能夠認定到.
這句話的真正意味.
所以當你下次還在問.
我的神為什麼離棄我.
你覺得神好像離開你的時候.
不是.
你需要做的就像耶穌一樣.

$^{1361}$禱告交託一生輕.
因為所託是神.
而且我們更加要認定.
上帝在你生命裡有一個藍圖.
我們按照上帝的計劃去行.
而在當中我們要知道.
沉默的同在.
就是因為神等待我們轉化成為光.
求主幫助我們.
我們一起祈禱.
因為我們知道耶穌.
你就是為我們受苦.
但你是勇敢去面對著受苦.
你願意為我們去承擔.
你不是被逼.
你不是勉強.
你不是作難.
你是甘心樂意.
為我們釘在十字架上.
以至成就救贖.
你顯明了你的堅強.
顯明了你的勇敢.
顯明了你的愛.
我們多謝你.
原主你讓我們更加捉緊神你自己.
因為你就是對我們不離不棄.
我們多謝神你.
聽我們祈禱.
奉耶穌名字以求.
阿們.
\newpage



\section{歌羅西書 2:6-15}
\label{sec:031881}
\textbf{超脫人生的塑造基礎 (1)}~~~~黃紹權牧師
\newline
\newline
連結: \href{http://www.church.com.hk/acms/content.asp?site=cdc\&op=show\&type=product\&code=031881\&layout=sermon}{\texttt{031881}} ~~~~ 語音日期: 2026-03-29
\newline
\newline
\hyperref[sec:031880]{\small{< < < PREV SERMON < < <}}
~
\hyperref[sec:index_chronic]{\small{[返順時目]}}
~
\hyperref[sec:031882]{\small{> > > NEXT SERMON > > >}}
\newline
\newline
歌羅西書 2:6-15
\newline
\begin{longtable}{cl}
\hline
\hline
章節 & 經文 (和合本修訂版)\\
\hline
2:6 & \begin{tabularx}{0.7\textwidth}{X} 既然你們接受了主基督耶穌，就要靠著他而生活， \end{tabularx} \\ \\ \relax
2:7 & \begin{tabularx}{0.7\textwidth}{X} 照著你們所領受的教導，在他裡面生根建造，信心堅固，充滿著感謝的心。 \end{tabularx} \\ \\ \relax
2:8 & \begin{tabularx}{0.7\textwidth}{X} 你們要謹慎，免得有人用他的哲學和虛空的廢話，不照著基督，而是照人間的傳統和世上粗淺的學說，把你們擄去。 \end{tabularx} \\ \\ \relax
2:9 & \begin{tabularx}{0.7\textwidth}{X} 因為神本性一切的豐盛都有形有體地居住在基督裡面； \end{tabularx} \\ \\ \relax
2:10 & \begin{tabularx}{0.7\textwidth}{X} 你們在他裡面也已經成為豐盛。他是所有執政掌權者的元首。 \end{tabularx} \\ \\ \relax
2:11 & \begin{tabularx}{0.7\textwidth}{X} 你們也在他裡面受了不是人手所行的割禮，而是使你們脫去肉體情慾的基督的割禮。 \end{tabularx} \\ \\ \relax
2:12 & \begin{tabularx}{0.7\textwidth}{X} 你們既受洗與他一同埋葬，也就在此禮上，因信那使他從死人中復活的神的作為跟他一同復活。 \end{tabularx} \\ \\ \relax
2:13 & \begin{tabularx}{0.7\textwidth}{X} 你們從前在過犯和未受割禮的肉體中死了，神卻赦免了你們一切的過犯，使你們與基督一同活過來， \end{tabularx} \\ \\ \relax
2:14 & \begin{tabularx}{0.7\textwidth}{X} 塗去了在律例上所寫、敵對我們、束縛我們的字據，把它撤去，釘在十字架上。 \end{tabularx} \\ \\ \relax
2:15 & \begin{tabularx}{0.7\textwidth}{X} 基督既將一切執政者、掌權者的權勢解除了，就在凱旋的行列中，將他們公開示眾，仗著十字架誇勝。 \end{tabularx} \\ \\
[1ex]
\hline
\hline
\end{longtable}
$^{1}$鄧小妹好 網上的鄧小妹也好.
在這裡問大家安安.
讓我們在再次聆聽主的說話以先.
讓我們一起同襯低頭禱告.
將我們的新心靈交給我們的主.
多謝神.
慈悲的父神多謝你.
讓我們再次回到你的啟示當中去學習.
並且從當中去認識上主你的性情.
並且看見上主你的救恩計劃.
怎樣從你創造以後一直到今天仍然在生效.
怎樣在我們每一個屬你的真正神的兒女裡面.
能夠生根建造.
好讓我們重歸於創造的原樣.
我們為此獻上感恩.
求你的帶領.
讓每一個無論是聽的,說的.
都能夠在今天這個訊息當中.
讓我們重新去明白到.
我們怎樣走面前的屬靈的道路.
禱告奉靠主耶穌基督.
北聖的名字而行.
阿們.
今天在教會裡面其實很多人.
是處於一種不信主的狀態.
當他們在這種狀態之下.
比不信主的人的表現更加可怕和可憐.
我一直都認為.
所信的東西.
是需要清楚從聖經作為基礎.
但是對於很多這些不信主的人.
錯了.
誤打誤撞.
以為自己信了主的人.
他有一些特色.
他認為他自己所信的東西全部都是對的.
不會去檢查也不會去監察一下.
聖經和他自己信的東西有什麼分別.
他於是乎認為只需要我在教會裡面做到某一些程度.
或者參與某一些活動.

$^{41}$我們就已經信得很無比的正確.
他以為他自己所認信的所謂基督教的教義.
是毫無瑕疵的.
錯的是其他人而已.
不會是我自己.
如果連這樣信的人.
都可以說是真的信的話.
其實聖經就沒有需要存在.
只需要一班人說一些宗教教義.
宗教的術語.
或者一些我們在神學上的理論.
神學的教義.
就已經足夠.
但是當我們是一個認真相信主的人.
屬神的子民.
其實我們為什麼要去讀聖經呢?.
就是因為我們會希望用聖經裡面的啟示.
去糾正我們今天的信究竟有沒有什麼地方扭曲了.
而往往那些不信主的人.
最怕就是拿本聖經出來.
去量度一下自己的信.
因為他接受不了自己的信是錯誤的.
他只是擁抱著一大串扭曲.
並且適合他自己口味的教義.
所以在我信主或者穆會的過程.
我經常和很多的信徒.
或者慕道的人.
都這樣說的.
我說很多的人其實今天活在教會當中.
他不信主比他不信主.
是更加來得有福氣的.
你聽到這句話.
我不代表你不信主是好事.
不過你不信.
比你不信好很多.
雖然兩樣都是很慘的.
為什麼會有這樣的情況呢?.
我相信最主要的一個原因.
是隱藏在教會當中.
有很多能言善辯,花言巧語的人.

$^{81}$而這些人往往他們始著大聲格惡的態度.
用一些手段使很多人.
透過他們的言語去籠絡著他們的心.
而這種人正正就是在上文.
我們剛剛在歌羅西書的第二章裡面.
就介紹了這一類的人.
在第二章第四節說.
免得有人用花言巧語來迷惑你們.
何謂花言巧語呢?.
當然就是講到那些人很懂得去用言語.
運用言語來扭曲事實.
或者用言語來欺控人.
在原文聖經裡面這個字眼很有意思的.
他說他說的話是很中聽的.
很順耳的.
聽上去似是而非的.
不過你又不能夠拒絕他.
因為很適合一個罪人的心態.
所以在上一堂講道的裡面.
我們講到很多時候這類人.
最懂得的就是怎樣去收買人心.
阿波羅在這樣的討論裡面.
帶到去哥羅西教會的當中.
我們可以說他是冒著一個很大的險.
因為他要做一個屬靈上的醜人.
揭開很多人的錯信.
他的目的不是去污衊其他的人.
他的目的是希望人能夠對這些錯誤.
去有辨識的能力.
所以我們前文也講過.
這個所謂分辨善惡的問題.
分辨善惡的問題不是在於那棵樹的存在.
而是在於人拒絕上帝.
希望用自己的標準去分辨善惡.
但上帝給人是一個分辨善惡的標準.
是由上帝去釐定的.
所以今天當我們知道有這樣的背景之下.
保羅就乘勝追擊.
要去更加講得清楚給哥羅西教會.
以及今天我們能夠有幸讀到.

$^{121}$歌羅西書裡面的啟示的每一個信徒.
讓我們去學會怎樣更加進一步去建立一個穩固的屬靈的根基.
使我們的靈性有基礎地去被上帝的說話塑造.
所以今天我給了這一段經文.
從第二章六節至十五節裡面.
稱之為超脫人生的塑造基礎.
我們怎樣能夠得到一份屬靈的生命.
是得到上帝的塑造.
而且超脫了人那種罪性那種的狀態.
從而使得我們去不單止辨別到什麼叫錯信的人.
同時間知道一個正確積極的根基去怎樣去建造.
這個就是今天我們去看這段聖經的主旨.
也都保羅是希望藉此鞏固哥羅西教會.
以及今天我們看到他書信的每一個屬神的子民.
能夠有一個屬靈的洞悉能力.
一方面要抗拒錯誤.
但同時間要把握擁抱成長的動力.
讓我們一起再次回到這段聖經.
歌羅西書的第二章.
我們再次同心一起讀一下.
從第六節到第十五節.
我們一起再看看這段經文講些什麼.
請大家等待節目一起打開之後一起讀.
你們既然接受了主耶穌基督.
就當遵祂而行.
在祂裡面生根建造.
信心堅固正如你們所領的教訓.
你們要謹慎.
恐怕有人用他的理學和虛空的妄言.
不照著基督乃照人間的遺傳和世上的小學.
就把你們擄去.
因為神本性一切的豐盛.
都有形有體地居住在基督裡面.
你們在祂裡面也得了豐盛.
祂是各樣執政掌權者的元首.
你們在祂裡面也受了不是人手所行的割禮.
乃是基督使你們斷去肉體情慾的割禮.
你們既受洗與祂一同埋葬.
也就在此與祂一同復活.
都因信那叫祂從死裡復活神的公羽.

$^{161}$你們從前在過犯和未受割禮的肉體中死了.
神赦免了你們一切過犯.
便叫你們與基督一同活過來.
又圖沒了在批律例上所寫.
攻擊我們.
有礙於我們的自據.
把他撤去.
釘在十字架上.
既將一切執政的掌權的腦內.
明顯吸眾人看.
就像著十字架跨城.
讀完這段聖經.
你得到些什麼呢?.
每個字你讀了都好像知道.
但讀完之後你都不知道他想說什麼.
有沒有感覺?.
我這番說話是對近一班.
信了主或返了教會很一段時間的人.
去說的.
這段聖經讀完之後.
如果是你靈修.
你讀完之後其實就會蓋上聖經.
因為你覺得自己已經.
只不過其實你全部都不知道他想說什麼.
究竟這段聖經想說什麼呢?.
我分開兩大部份來看.
第一個部份.
從第二章的六節至第八節.
在這一段落的裡面.
說了一個比哥羅西教會.
以及我們知道的事實.
越人的事越要我們格外警惕.
我的英文題目叫做.
Extra cautious to those easy delightful advocates.
為什麼要讀英文呢?.
我覺得英文容易理解.
當一些事情你越覺得好聽的時候.
你就要越小心.
要越加留意.
我和師母都很喜歡去旅行.

$^{201}$其實有誰不喜歡去旅行呢?.
不過有些人真的不太喜歡去旅行.
因為有些人喜歡黐家.
不過我就很喜歡去旅行.
為什麼很喜歡去旅行呢?.
因為我去旅行會見到很多我見不到的東西.
不是因為讀萬卷書不如行萬里路.
我是很喜歡去到不同的地方.
跟當地的人和環境有一個連結.
有一個交往.
所以我最近和別人談.
我其中一樣去旅行最喜歡去的.
是什麼呢?就是看那些洞洞洞那些地方.
為什麼呢?.
因為我覺得那個和這個世界有些連結.
我和師母曾經試過去夏威夷租一輛車.
由碼頭開車去到一個海洞.
在海洞開了兩個多小時.
為什麼浪費時間呢?.
我還記得在海洞外面撿了一個木瓜.
後來送給別人吃.
那個海洞在我年輕的時候讀書.
我想大約十三歲十四歲左右的時候.
我看地理書裡面有很多圖片拍過那個洞.
我經常都很想去那個洞.
又好像我很喜歡去旅行.
去到在新加坡有個地方叫Balancing Rock.
有沒有人去過?.
有人在那時候我讀書的時候看過.
不過你們不用去了那塊石斷了.
我去的時候都還看到它在那.
所以有一次我和師母去旅行的時候.
就去到洛杉磯旅行.
就去了一個其實那個區很窄的.
叫做Fashion District.
有沒有人去過?.
Fashion District你去的時候.
你會很失望的.
因為那裡沒有什麼真正的Fashion.
那裡賣的那些衣服.

$^{241}$比利源東西街以及深水Po 那條大馬路.
長沙環道那些衣服還要差.
但是無論如何.
那個很複雜的地方附近就有一個花墟.
我和師母去看花.
師母喜歡花.
我也喜歡花.
我去看.
嘩.
去到那個地方.
有很多各式各樣的盆栽植物.
很多很漂亮的花.
很漂亮的花.
各式各樣顏色都有.
不過我看了很久之後.
那個花墟其實很大.
那個花墟有幾個市場.
分開很大.
我去的時候差不多收工時間.
但是還有很多花賣.
但是那些花.
我突然間就想起一個問題.
那些花這麼漂亮.
顏色這麼鮮艷.
究竟是不是真的呢?.
尤其是那個花墟裡面大部分賣東西的.
都是亞洲和亞洲來的人.
當你發覺在亞洲來的人.
很多時候這些技術很高的.
造假.
那我就.
就和師母說.
那些顏色是不是真的呢?.
有些是紅色.
有些是綠色.
有些是黃色.
尤其是玫瑰.
有些是粉紅色.
粉紅色我也見過.
真的我也見過.

$^{281}$但是為什麼這些花有這麼多種顏色呢?.
我就去找.
師母有看他的花.
我就去看看究竟有沒有人造假.
誰不知當我走到某個地方.
因為接近收工的時間.
造假活動就出來了.
原來那些亞洲人.
他們很厲害.
泡了一盤很多的顏色水.
然後把白色的玫瑰插進去.
然後開晚.
那些叫紅外線燈.
很熱的.
然後那些植物.
那些花.
雖然剪掉了.
因為它有熱力會由根部吸那些顏色上去.
那些花就開始變色了.
各位姐妹.
漂亮的東西.
其實你要格外留神.
讓我們看看.
保羅在這裡.
保羅在第一章第六至第七節裡面.
已經用過三次一個字眼.
其實重複講過的.
叫做柔軟.
Cathars.
保羅作為一個信徒.
你有一些東西要很像.
什麼什麼.
很像什麼什麼.
我不再重複講了.
不然我今天講到有超過一個小時了.
來到現在第二章的裡面.
尤其是去到第二章的第七節.
保羅也都同樣.
再次用同樣這個字眼.
Cathars.

$^{321}$不過和合本的翻譯就改了叫做.
正如.
要告訴你.
你要像一樣東西.
其實這個結構.
跟著下去.
在第二章的第十六節.
就講到一些不正如.
你不要像的東西.
所以你偷偷看看一點點.
好不好.
我們看看.
在第二章的第十六節那裡.
就說.
所以不拘在飲食上或節期上.
越索安息日都不可.
讓人論斷你.
不可.
有些東西你要像的.
但有些東西你不要給人有機會.
有機可乘.
保羅在對於哥羅西教會的模樣上面.
其實要給哥羅西教會清楚知道.
你們的屬靈生命.
要與基督的樣式.
是匹配的.
是類同的.
在這個基礎之下.
因此他說.
你們要因此而似基督.
所以在第二章的第六節就這樣說.
你們既然接受了主基督耶穌.
就當遵祂而行.
你就必須要.
與基督一樣.
按照祂的行為.
思想.
來照著去運作.
那怎樣像祂去運作呢?.
所以第七節就說.

$^{361}$你要怎樣了.
第七節.
在裡面說了一大串的事物.
而這一大串的事物.
都說了有三方面的東西給我們知道.
就要在主基督的裡面.
所以我們要怎樣呢?.
在祂裡面.
這個詞彙是很重要的.
如果那本聖經是你的.
或者現在無所謂了.
電子聖經.
highlight了這個.
在祂裡面.
你要在主耶穌基督裡面.
做三樣的事.
或者像三樣的事.
第一.
你要在祂裡面.
生根建造.
你落點.
你的生命的基礎.
你的生命的養分.
全部是由基督.
來去吸取的.
不是從其他地方吸取的.
這個是談論到整個基礎的因素.
而這個基礎的因素.
只有基督才做得到的.
不是由其他的來源.
其他的人塑造的.
有時不知道大家喜不喜歡種東西.
有些東西是種在某些土壤上面.
你就必須要找回那些泥土.
譬如.
有人喜歡種蘭花嗎?.
我很喜歡蘭花.
不過種到都死的.
不知道為什麼.
對於我來講.

$^{401}$蘭花很難種.
於是我就去問花博.
怎樣才種得好蘭花呢?.
他說你首先要找些適合的泥土給他.
我說我給他一些漂亮的泥土.
那些叫做什麼?.
Pot Soil.
這樣就漂亮死了.
蘭花不可以種在Pot Soil.
是嗎?.
不知道.
蘭花要種在一些芒草裡面.
它要梳起.
但又要保持濕潤.
但又不可以太濕.
濕了它會發霉.
是嗎?.
原來是這樣的.
但回頭.
當我去湯泉道家的時候.
我看到那些蕃茄種得這麼漂亮.
處處蕃茄.
那些無花果樹.
每一棵都長了果.
我說你怎樣種的?.
他說當然是要用黑土.
小肥的泥土給他.
我說是嗎?.
不是用芒草會好一點嗎?.
他說不同的.
種每樣東西有它特定的方式.
用泥土是不同的.
基礎是很重要.
這個基礎.
基督徒的成長.
要在基督耶穌的供應的養分之下去成長.
第二樣東西.
要使你的信心堅固.
堅固的時候做什麼?.
正如你所領受的教訓一樣.

$^{441}$原來是這樣的.
在它裡面原來是要持續地去增長.
所以《原文聖經》就用了一個字眼.
就是一個持續的成長來去介紹.
而介紹這個目的.
就是讓我們知道我們不是停留的.
做基督徒你的生命.
不是一次得救永遠得救.
信了就完結搞定.
不是這樣的.
我經常跟別人說.
婚姻是一樣的.
婚姻不是只有婚禮.
婚姻是比一個婚禮多很多東西.
是持續的.
生長的.
要使一個婚姻.
使一個宿命生命成長.
是有些因素,有些質素.
而這個質素先前可以說是基礎是從基督來的.
但是我們要怎麼樣?.
要持續地生長.
不過在這裡我先介紹多一點給大家知道.
這裡強調了一個叫.
因為這個跟著就是整個的歌羅西書.
催促信徒成長的動力來源.
在歌羅西書裡面是重複又重複地說的.
要不斷地在這裡重複重複去建造.
來向耶穌這個養分吸取之後.
後來會告訴我原來是向上帝.
還要信心堅固.
在什麼上面去堅固?.
原來是在所領受的教訓裡面去堅固.
信徒的生命是需要持續地去重複又重複地再被肯定.
為什麼牧師你講這些跟和合本有些分別的次序?.
因為本身原文的次序跟和合本不同.
所以我按照原文聖經的次序告訴大家.
保羅在這裡向哥羅西教會說.
你往後的屬靈生命要在基督裡面去成長.
要吸它的養分.

$^{481}$以它作為根基來使你成長.
不是從其他地方吸.
你吸其他的東西就吸其他的古靈精怪的東西.
你只是吸基督的養分.
然後怎麼樣?.
你不斷被建造.
建造什麼?建造一個感恩的心.
感恩的生命來在裡面驅動你們去成長.
之後會詳細解釋.
但是在這個成長過程當中.
其實是不斷重複肯定你的信.
確認又確認確認又確認.
信耶穌不是一次過的事.
而是在你的人生當中不斷重複去問.
這條問題.
我是不是真的信主呢?.
如果是的,我信得清不清楚呢?.
如果問我信得不清楚的話.
那不是洩氣.
讓你知道你信得不清楚.
然後下一步我怎樣繼續去得到.
弄得更清楚呢?.
這個就是保羅說給.
哥羅西教會的整個核心.
而這個核心是要讓我們在這個成長過程中學曉感恩.
向上帝感恩.
因為當我們知道原來我們所得著的成長過程.
全部是由上帝的兒子所帶來給我們.
當我們漸漸有耶穌基督的生命樣式的時候.
我們也同樣會在這方面.
但是在這種種的基礎之下.
保羅作出了一個.
似乎從一個負面的角度.
來告訴哥羅西教會知道.
你有些事是需要很謹慎的.
其實上一堂已經說了.
保羅已經要做這個屬靈醜人.
點出哥羅西教會裡面有一些.
其他的人傳其他的信息.
而且他們曉得花言巧語來引誘你們離開.

$^{521}$為什麼呢?.
因為保羅要讓他們知道.
我們不是活在一個完全無菌狀態.
無罪狀態的世界.
我們仍然會在一個有罪的世界裡面.
而罪的影響仍然籠罩著我們身邊.
我們如何謹慎自守.
這就是白節所說的.
白節說什麼呢?.
白節說你們要謹慎.
恐怕有人用他的理學和虛空的妄言.
不照著基督乃照著人間的遺傳和世上的小學.
就把你們擄去.
很嚴重.
說到擄去這樣的情況.
保羅其實從來不介意自己做這個屬靈醜人.
點出那些人的問題.
雖然他沒有點名.
但是保羅說出那些人的問題.
很自然大家都知道那是什麼問題.
有什麼問題?.
保羅昔日點出給哥羅西教會知道.
他們教會裡面有的問題.
今天同樣有.
我們一齊仔細看看.
保羅用了一個語氣.
作為一個很嚴肅的虧權.
所以謹慎這個字眼.
希臘文叫做.
Blappetas.
這個字眼和流心是有同樣的意思.
即是用廣東話.
你夾實呀.
看清楚呀.
這樣的意思.
他提醒原來在哥羅西教會裡面.
有些人和你們不同.
怎樣不同?.
恐怕有人.
不用恐怕了.

$^{561}$不是客氣的說話這麼簡單.
是保羅顯出他自己對哥羅西教會的擔憂.
好像.
我說笑一點.
趁我嫁女兒期間就思想多一點.
其實經常很怕子女.
不知道出去認識了什麼朋友.
身邊什麼人.
認識了什麼人.
就算他多大都好.
你也會擔心他.
會不會被人騙呢.
廣東話叫做被人騙呢.
保羅很擔心的.
擔心你們有些人會騙你.
怎樣騙你呀.
前面說花言巧語.
這裡更加說得很實在.
分開四樣東西.
表達出這類人的表現.
第一.
這類人他們很懂得善用理學和虛空的妄言.
第一.
這裡.
已經說了兩樣.
他們的用理學和虛空的妄言.
另外一件事是什麼呢.
他們就是不照基督.
已經說了三樣了.
還有一樣.
這類人是怎樣的呢.
乃照人間的遺傳和世上的小學就把他們擄去.
中共這類人.
來騙你們的人.
帶你們離開主的人.
擄走你們的人.
他們這四個手段層出不窮.
用理學.
和虛空的語言.
或者把不好的東西一起說.

$^{601}$用人間的遺傳和世上的小學.
告訴你們.
偏偏就不用基督的教導.
你就會問牧師.
究竟這些人的理學.
虛空的妄言.
或者遺傳.
世上的小學是指什麼呢.
可不可以給我一些例子.
有,當然有例子.
在《原文聖經》裡面.
對不起,因為我從來沒有用《原文聖經》.
所以我中文.
看它的次序.
其實有點混亂.
不過方便你們看而已.
畫本就.
不要緊,大概都齊了.
什麼叫做用理學.
來講說話呢.
理學就即是說.
他用一些哲學.
來跟你講說話.
今天的例子最多就是.
什麼神,什麼神.
什麼主義,什麼主義.
在北美洲.
沒有一個比.
我可以這樣說.
這個世界真正的.
社會主義國家.
沒有一個比加拿大.
更社會主義.
什麼都講.
三樣東西.
多元.
是吧.
加拿大最喜歡的,diversity.
各個族裔.
各個人都可以個別獨立.

$^{641}$不同的.
只要每個人獨立處理.
你們中了很多這種毒的.
不要說你們,我們中了很多這種毒的.
第二是怎麼樣呢.
公平.
equality.
每樣東西要公平對待.
雖然他不同,不過你要對他一樣.
其實這個在騙你.
當你看見有些人.
他跟你不同,不過.
他的待遇又給你優厚的時候.
同一間公司而已.
為什麼他可以做得,我不可以做得.
最簡單.
有些公司.
有些.
我不是歧視.
雖然C9法案就快通過了.
通過了洗腦教會更加.
看你是不是真的全主導.
同性戀者.
你看看你的公司.
你好人好者不是同性戀的.
你講.
講男女關係.
被人覺得你是性侵犯.
但是那些同性戀者講性的.
肆無忌憚.
你有沒有覺得這些情況.
你沒有批評他.
同性戀,錯.
馬上C9法案.
看你基督徒還不早見證.
那些時候就是早見證的時候.
我怕,不要說了.
現在教會就是這樣.
鬼殺更吵.
不要通過.

$^{681}$又說敵基督諸如此類.
這個世界越邪惡.
就是越多見證主的機會.
還有一個.
I.
I是什麼.
知不知道.
D就是diversity.
E就是equality.
inclusiveness.
兼修並修.
什麼都要接受.
你不同意都不能不同意.
全部要包攬.
什麼都要接受.
明明你見到那個人.
是有點奇怪.
你不能說他奇怪.
你要接受.
其實你又中了這些毒.
說出來好像.
哎呀,大愛,我們個個都愛.
我們接受所有人.
你還需要辨別朱玲嗎?.
你閉上眼睛就行了.
我愛,我什麼都愛.
教會都是這樣.
而且教會更柔甚.
什麼神,什麼神就是這些.
第二.
這個談論到.
所講及到的有關於.
虛空的妄言.
就是講一些空洞的說話.
空洞的說話.
我們經常都聽的.
你扭開電視.
我們要討論的.
討論,開完討論之後.
討論了什麼.

$^{721}$結果就是討論了.
就是討論了那堆東西.
就是討論了一些東西.
講了等於沒有講.
你有沒有留意.
你在工作環境裡面.
甚至你在家裡面.
你跟很多家人講完.
大家各書幾見.
大家各表現之後.
沒什麼.
但是這樣講了之後.
你就會覺得自己好像有成就.
在參與.
這些人很懂得用這種方法.
使你覺得.
你好像很重要.
我曾經試過在一些.
商業機構裡面工作的時候.
有個同事.
什麼都拿出來講.
什麼都拿出來問.
很多發言.
我自己心想.
你有沒有做功課.
做了功課.
這些問題.
功課已經有答案了.
為什麼你還拿出來講呢.
最可怕的就是.
在這個大環境之下.
虛空謊言的世界裡面.
那個所謂.
負責的教導者.
還要怎麼樣.
Good question.
好問題.
問多點.
還要講.
No questions is stupid.

$^{761}$其實他問的那條就是.
Stupid question.
不過他就說No questions is stupid.
為什麼.
因為不想得罪他.
弟兄姊妹.
是,No questions is stupid.
但是你有沒有做過功課.
有沒有心思細想過.
那條問題.
你想問的問題.
是不是自己真的找不到答案呢.
所以聖經任你問也是.
為什麼我喜.
我欣賞我們教會的弟兄姊妹.
包括網絡上面的所有弟兄姊妹.
參與這個主教學課程.
因為他問的那些問題.
真的想過度過.
找不到答案.
他才會走來問.
雖然有些問題.
是很基礎的問題.
但是真的找個不懂的牧師.
還有什麼牧師又這樣講.
那些大有名堂的.
我不信他又好像不是很好.
不過我告訴他我不信他.
如果被人譴責.
怎麼處理.
所以我們那些弟兄姊妹.
問的問題很有水準.
不是隨便出來.
不斷地騙.
不斷地問.
表達自己好像有存在感.
有重要感.
欺騙人的人就是講這些東西.
塑造這樣的空間.
還沒停.

$^{801}$還有.
我們包括.
就是什麼.
他們很喜歡講那些.
人間遺傳和.
世上的小學.
這個陳述.
就有些.
困難的.
什麼人間小學.
什麼遺傳小學.
講什麼呢.
其實在愚民聖經裡面.
講到這個.
遺傳你不用講了.
阿爺講了什麼.
牧師講了什麼.
長老講了什麼.
他講就好像真理一樣.
不要挑戰.
我曾經聽到一個很著名的牧者說.
他們上一代的神學.
是錯的.
不過不要指證他.
我們要尊敬他.
我想說這些.
巧言善辯.
很多人很喜歡聽他講道.
你尊敬他什麼.
他講錯你就應該告訴他.
他講錯.
點出他錯的地方.
尊敬他就是因為你能夠提出.
讓他知道他的錯.
所以我尊重.
我沒有告訴他我錯.
你這樣才是尊重我.
你不尊重我就引惡揚善.
講小學.
什麼叫小學.

$^{841}$原文用了一個字眼.
叫做.
Stoichkaia.
這個字眼有三個可能解釋.
第一個解釋.
是指在物理世界裡.
最原始的物質.
你們不認識這個字.
圓塔.
元素.
元宇宙.
聽聞已經.
那個元不是原始的元.
是還原的元.
這類人講的東西.
好像很精密.
好像這個世界很微小的東西.
我一會兒給個例子.
另一個解釋.
是世界上最幼稚的教導.
幼稚不是歧視.
是最基礎的.
來顯示.
你不知道嗎.
這麼低能啊.
很簡單的東西.
你不懂嗎.
我告訴你.
第三個可能解釋.
這個字眼.
世上最原始的觀念.
或者.
最原始的靈性.
是什麼東西.
這些就是宗教.
虛無的東西.
牧師.
給個例子.
有沒有聽過傳道人.
他說他就是傳道人.

$^{881}$講DNA.
有沒有聽過.
然後轉頭就講.
複製人.
所以世上沒有一個人不懂.
是我這個專家.
才懂的.
我研究了複製人很多年.
你們不懂.
你們聽我說.
他真的這樣講.
我在現場聽過他這樣講.
而且他不是一次這樣講.
他每次都這樣講.
那時候還要大大聲說.
複製人是不可能的.
後來.
好像有很多複製人.
領子們.
這類人.
藏身在教會裡.
講一些.
你好像不懂的東西.
但好像他才懂的.
專家.
這類人.
為什麼要講這些東西.
講一些.
哲學.
diversity.
equality.
inclusiveness.
多元.
平等.
包容.
大愛.
教會才會講這些東西.
因為講出來很宏偉.
教會多愛大家.
然後還講耶穌很愛你們.

$^{921}$這些所有的人.
所以就算同性戀.
我們人也很愛.
耶穌只是不愛你的行為.
那些罪.
就是耶穌不愛你的行為.
你就脫離這些.
不需要的.
我接受你.
inclusiveness.
講一些你不懂的東西.
複製人.
外星人.
原始人.
你不懂嗎.
不懂.
你不懂就聽我說.
他用這些方法.
去引誘人.
去聽他講.
為什麼他們要.
用這些方法呢.
因為有一個原因.
當你去.
聽這些人講的時候.
這些人其實就已經.
瞄準著你的人性.
弱點.
人性的弱點是什麼.
你問問你自己.
人性最大的其中一個弱點.
就是懶.
又要知道.
信主好,有天堂上.
他們一定讚到.
天上有地下無.
非常之好的.
好正.
一定要來.
一定要接受,不接受你沒有機會.

$^{961}$信主你就死在陰間.
上天堂是非常之好的.
不過要上天堂.
入天國裡面.
又要努力.
這些他不講的.
等他講,他做了給你.
於是乎.
你聽了就當以為自己做了.
他看到你這個.
懶惰的態度.
人性不勞而獲.
但是要最好.
所以你越不需要用腦的東西.
他就講一些越基礎的東西.
講完之後.
是啊是啊,是嗎?.
我不知道,現在知道了.
那我就信主了.
他用這種方法去使到你們.
覺得聽了他們講的東西.
就不需要自己再動腦筋.
不用再努力.
而不用再努力.
去到天國裡面.
聖經怎樣教導的.
你從創世記到啟示錄.
你看看.
真正信主的人.
無一不掙扎在上帝裡面.
但是這些人就用這些言語.
考驗.
來為你包裝了一個套餐.
你聽了就等於做了.
那你就聽他講.
你不聽他講.
你就覺得不信主了.
沒有信好,所以你繼續去聽他.
這種就是這類人的性情.
他針對著人性的.

$^{1001}$弱點.
或者甚至說.
針對人的罪性的性情.
用這種誘惑去聽從他.
而你要很甘心樂意.
隨著你自己的罪性.
去跟著他,不需要處理你自己的罪.
所以他.
什麼都講,他偏偏就不講.
基督,所以在這段聖經.
講得很清楚,就是偏偏.
不照基督.
一講基督就講罪了.
一講罪就要處理你的罪.
一講處理你的罪.
你就要出力了.
要付出代價.
為什麼要付出代價?基督要付出代價.
為什麼你不需要付出代價?.
經常都說我們要仗基督.
有基督的樣式,基督的樣式是怎樣?.
基督就是要付出代價.
為他的信而去付出.
但是.
這類人偏偏就不講這件事.
免得阻礙了.
他們去收買.
人心的大計.
所以.
為什麼我用.
標題說.
越越人的事.
你要越要警惕呢?.
因為那些東西越吸引的時候.
其實就越.
多機會去造假.
因而.
你以為自己得著了.
最正的東西.
但其實你根本一些最正的東西.

$^{1041}$都得不到.
有沒有聽過一些新聞.
現在不是很喜歡整容的.
有些夫婦.
兩個都很靚仔靚女.
高富帥.
另外一個就.
白富美.
然後結成夫婦.
然後開開心心的結婚.
然後結完婚之後.
沒多久就離婚.
當她一生小朋友就離婚.
因為一生出來那個原來.
很醜的.
他們就會指責對方.
為什麼生出來的兒女那麼醜.
原來兩個都很醜.
兒女就遺傳.
最新鮮我有個例子.
越吸引.
你就越小心.
我有條褲.
師傅有條褲.
要找人改.
我們知道有個地方.
華人必到的.
翻版聖地那個地方.
多倫多全世界.
北美洲出名的.
那裡有間改褲很便宜.
我聽過十元改褲.
我去到我都問他.
那條西褲多少錢.
十元.
是不是即日拿.
即日就不行.
即日要加價.
即日要加兩元.
我不即日.

$^{1081}$他說好啊.
你試條褲我釘好.
然後你出來我改.
當我穿了條褲.
釘完之後.
釘好之後他就來了.
要加價.
我說為什麼要加價.
加幾錢.
他說你條褲.
不可以就這樣剪短車翻賣.
我說為什麼不行.
披口.
你不要騙我.
我都識少少縫紉.
我識挑褲腳.
不過不想自己做.
他說我們要怎樣.
拆腳和你拆了.
摺了然後才車翻.
我心想.
和你剪了.
然後摺回去.
摺了入去然後車翻.
不就一模一樣.
況且我條褲不是改好多.
改穿了.
怎樣.
我說加幾錢.
他說兩蚊.
你知不知.
他其實想騙多我兩蚊.
始終他都想.
做十二蚊生意.
好不緊要.
兩蚊我給你.
到師母.
穿條褲出來.
整了之後.
十蚊.

$^{1121}$你條褲不行.
你條褲.
講好多ye .
講完之後.
要加錢.
加幾錢.
五蚊.
牧師.
我奉獻給你.
你不用奉獻給我.
七蚊多給他不是問題.
不過我想用這個例子.
給大家知.
他要了你的生意.
他想吸引你入去.
一個很便宜的價.
十蚊.
多倫多哪裡有十蚊改條褲.
總會有.
不過不是那麼多.
到我幾十蚊多.
十七蚊是最便宜.
要加稅.
另外一間19.99要加稅.
複雜些.
更貴.
十蚊吸引你入去.
踩在頭上.
其實我是否可以不做.
可以的.
不過他看中了你.
你的心已經被我籠絡住.
其實我很控制.
不過我怕麻煩.
越吸引你的東西.
好像我去花.
去花墟那裡賣.
看別人的花.
那麼漂亮的花原來全是假的.
顏色原來是混在一起.

$^{1161}$屬靈生命.
是一樣的.
如果有些人說的東西.
你越喜歡聽的時候.
他可能之後才慢慢.
載覺你.
我多說一點.
因為都很長.
教會很多是這樣.
你有沒有留意一個事實.
你未信主.
未回教會的時候就很熱情.
不斷叫你信主.
對你很有愛心.
甚至煲湯給你喝.
有甚麼帶熱你.
當你信了主.
受了浸.
加入教會之後.
他就開始告訴你很多東西.
你要奉獻.
甚麼叫奉獻.
不過你都受了洗.
入了局.
奉獻多少.
看愛心.
這些最要命.
說些不說些.
接著再說甚麼.
最初就說十分一.
有些就狼胎到.
你全人奉獻.
你入了.
受了浸.
你就會在一種內疚裡面.
活在教會裡面.
你就會不夠.
屬靈.
弟兄姊妹.
你傳道的時候.

$^{1201}$傳道給別人.
我很希望你多問他一些問題.
你說你對信耶穌的人.
有多少理解.
有沒有困難.
你對於他們信主.
我有些朋友.
我那時跟他講傳福音.
我跟他說.
我們做信徒是這樣這樣.
我會這樣想.
我的朋友就說.
做信徒很難的.
很多事不能做.
我說那是呀.
不是我不能做,是我選擇不做.
那你有沒有困難呢.
如果要你信.
有你就不要信.
因為我叫你信.
你信都是騙我的.
不過你騙我不重要,你騙他們自己.
但最嚴重你騙上帝.
不過上帝不讓你騙.
不過撒旦.
是最開心的.
所以弟兄姊妹.
教會裡面.
有很多這樣的人.
有弟兄姊妹跟我說.
牧師我失業了.
我失業了就找工作.
我失業沒錢奉獻.
就不要奉獻.
教會沒錢.
教會沒錢就沒錢.
沒錢教會倒閉.
那就倒閉.
教會倒閉.
收入主名不是.

$^{1241}$你去借錢.
來奉獻之後.
使教會導致.
一個弟兄一個姊妹.
負債累累.
誰最嚴重.
教會倒閉比較嚴重.
還是弟兄姊妹.
負債累累比較嚴重.
我相信在.
上帝眼光裡面.
寧願教會倒閉.
都好過教會.
弟兄姊妹在錯信之下.
好了.
保羅說了這些警告.
其實是人都會知道.
在教會裡面發生什麼事.
接著從第九節.
到十五節裡面.
就告訴我們.
怎樣才能夠.
過一個超脫的人生.
牧師.
你說到說了九個字.
怎麼說呢.
這裡有這麼多字.
所以就不說了.
下課再說.
不過我想告訴大家.
保羅在這段聖經裡面.
將基礎.
以正面.
和負面已經說了.
給哥羅西教會知道.
哥羅西教會的.
信徒以及在他當中.
這些會去欺騙人的人.
會怎樣看待.
保羅呢.

$^{1281}$保羅會不會因此.
而使他寫這封書信.
那個作用是完全被.
英文好說.
Cancel out.
被抹除.
沒有了功效.
很危險.
但是牧養.
其實就是要冒這些危險.
好像養.
教養小孩一樣.
做父母.
什麼的父母才叫做失職的父母呢.
是見到子女.
好的你會去讚.
但見到壞的.
你不去彈.
這樣你就是失職.
不是啊牧師.
這樣會使子女討厭我.
讓我.
借一句說話告訴你.
我寧願你討厭我.
不過你在上帝裡.
被上帝得著你.
我無所謂.
你討厭了我.
但上帝剩下一個.
合他心意的人.
我寧願.
你討厭了我.
或者你將來會去到天國的時候.
見到我的時候.
會說.
真是不好意思.
牧師又見到你.
我討厭了你.
不要緊的.
到時大家在主裡面.

$^{1321}$打招呼.
然後.
回到父神裡面.
一起歌頌上帝.
就可以了.
這樣的屬靈生命.
才有根有機.
讓我們先學會怎樣去警惕.
亦都找到成長的基礎.
然後我們就在下個星期.
對不起下個星期好像是.
上次主日.
即是要等我.
嫁女那天說.
這樣也好.
原本嫁女那天說的堂堂.
是預早說了兩個小時.
不過.
現在可以說短一點.
讓我一起跟上帝講.
父神多謝你.
讓我們在這段聖經裡面.
看見保羅對.
哥羅西教會的關懷.
亦都讓我們作為信徒.
我們怎樣去看見.
上主你怎樣去關懷.
每一個屬你的子民.
你擔心我們.
能夠在這個充滿罪惡.
包圍的世界.
當中被誤導.
被老去.
因此.
你透過你的仕途.
直言不諱的.
去將真相.
揭露在我們面前.
求主你的幫助.
讓我們有悟性.

$^{1361}$去領受這番說話.
讓保羅的啟示.
讓我們更看到.
我們怎樣要活在.
基督的裡面.
曉得.
辨別錯誤.
但同時間.
能夠尋找到.
真正成長的動力.
來走我們面前.
屬靈的道路.
求你的幫助.
求你的閱立.
我們對你的那份情請.
求你.
讓我們有智慧.
來去曉得.
真正走在.
這條屬靈的.
人生路當中.
由基督的樣式.
散發在我們每一天.
的生命生活當中.
同心的禱告.
奉靠主耶穌基督得勝命至教.
阿們.
\newpage



\section{羅馬書 8:28-30}
\label{sec:031882}
\textbf{建立健康的使命教會 (10) 塑造活像基督的生命}~~~~吳磐德牧師
\newline
\newline
連結: \href{http://www.church.com.hk/acms/content.asp?site=cdc\&op=show\&type=product\&code=031882\&layout=sermon}{\texttt{031882}} ~~~~ 語音日期: 2026-03-29
\newline
\newline
\hyperref[sec:031881]{\small{< < < PREV SERMON < < <}}
~
\hyperref[sec:index_chronic]{\small{[返順時目]}}
~
\hyperref[sec:031883]{\small{> > > NEXT SERMON > > >}}
\newline
\newline
羅馬書 8:28-30
\newline
\begin{longtable}{cl}
\hline
\hline
章節 & 經文 (和合本修訂版)\\
\hline
8:28 & \begin{tabularx}{0.7\textwidth}{X} 我們知道，萬事都互相效力，叫愛神的人得益處，就是按他旨意被召的人。 \end{tabularx} \\ \\ \relax
8:29 & \begin{tabularx}{0.7\textwidth}{X} 因為他所預知的人，他也預定他們效法他兒子的榜樣，使他兒子在許多弟兄中作長子。 \end{tabularx} \\ \\ \relax
8:30 & \begin{tabularx}{0.7\textwidth}{X} 他所預定的人，他又召他們來；所召來的人，他又稱他們為義；所稱為義的人，他又叫他們得榮耀。 \end{tabularx} \\ \\
[1ex]
\hline
\hline
\end{longtable}
$^{1}$各位弟兄姊妹,早安!.
在進入救經文之前,我們一起心低頭祈禱.
因為你誓願耶穌基督來到世上.
這個星期特別是一個很珍貴的時刻.
讓我們重新體會經歷.
主你是如何榮耀進城.
你是如何受苦受死.
你是如何為我們復活.
以至我們有永恆的盼望.
主,我們求你憐憫我們.
由我們來到主你的面前.
我們現在打開我們的心靈.
讓我們的生命向主說.
主啊,讓你成為我們生命的主.
去掌權,去改造我們.
去榮耀你,去與你同工.
以至到復活的主的生命.
與我們一起彼此的配搭.
我們願意與主你同行.
我們將以下時間交給天父.
求聖靈在我們當中自由運行.
成就你要作的功.
祈禱感恩,奉主的聰明求,阿們.
不能夠或缺的.
亦都是彼此互相的配搭.
特別當我們有敬拜的人生的時候.
其實我們的敬拜.
不單止在我們現在崇拜當中.
有一個活潑的敬拜.
我們都知道我們就是演員.
神就是我們唯一的觀眾.
我們的生命向祂展現.
敬拜的人生就是我們如何用回應.
用信服去回應祂.
去愛祂.
去更充滿神自己的生命.
我們如何做得到呢?.
就是在第二個元素裡面.
特別是有質素的相交.
我們不單止在小組,在團契.

$^{41}$我們彼此去認識.
成為朋友.
更重要的是.
在有質素的相交裡面.
我們要認定.
如何幫助大家互相.
來活出敬拜人生.
又如何有勤成長的靈命.
如何一起去服侍.
其實是互相緊扣的.
每一個環節都不可或缺的.
如果你回到小組.
你沒有一個敬拜人生的話.
其實只是一個社交的關係.
如果你回到來崇拜.
你不想活像基督的話.
所有的道其實都可以入而出.
所以是你如何和神相交.
你如何去回應.
如何透過這五個.
我們個人在神的面前.
能夠成為活祭.
我們個人能夠是健康的屬靈生命的話.
我們就必定會有一個健康的使命教會.
所以必須要從你和我開始.
所以在第三個元素裡面.
說到勤成長的靈命.
上星期.
Pastor Leo 跟我們說了.
神話語的重要性.
邀請大家愛慕神的話語.
今天我想說的是.
我們愛慕神的話語.
我們來崇拜.
我們去小組.
究竟神的心意.
要我們成長些什麼呢?.
成為怎樣呢?.
這麼大本聖經.
其實今天羅馬書第八章二十八至三十節的經文.

$^{81}$正是告訴我們最重要的.
成長靈命的目標.
就是.
神也預定他們(我和你).
效法祂兒子的榜樣.
我們今天就是去領受清晰的目標.
以至我們的生命是充滿使命的.
就是我們要成為基督的樣式.
又活出基督的生命.
以至我們能夠透過這個健康的.
屬靈的生命.
建立健康的使命的教會.
讓我今天透過這短短三節.
從聖經的教導裡面.
也從我自己過去這半年.
一些我個人的屬靈的經歷.
一些的體會.
很想和大家分享.
以致我們能夠一起建立一個教會.
獻給神願意告訴主.
你塑造我們活像基督的生命.
你願意做這個位身.
主啊 你塑造我們.
去活出像基督的生命.
在預備這個講章的時候.
剛剛我兒子去香港.
他就在三藩市坐飛機回香港.
這張照片就在飛機上拍攝.
這也是我孫子第一次長途飛機.
中國人都是比較節儉一點的.
我兒子為什麼一定要現在飛呢?.
因為他兒子三個星期就到兩歲了.
你們父母是否明白呢?.
很重要的.
如果三個星期之後就貴很多了.
所以他們兩個就劫持我孫子去香港.
你知道你弟弟坐這麼長途飛機.
你做父母也明白.
其實不是那麼容易的.
回到香港.

$^{121}$不要說他回到香港.
我只是看看他孫子.
究竟他像誰呢?.
這張照片我比他大一歲.
你覺得他像不像我?.
一點點吧.
我的反應是說其實也不太像.
但是很多人就說他眼睛像.
但是你說他像不像我兒子呢?.
我兒子跟他差不多年紀的時候是一樣的.
一個個包.
十足.
但是其實眼睛和獅子鼻子是避不開的.
我也是很放心的.
因為當我在這裡預備堂道的時候.
其實神也放心的.
只要我們還在祂的家裡.
只要我們還抓住祂.
因為我兒子小時候.
左邊那張照片就是我兒子.
這張照片就是我.
一般般的.
但是越大越像.
認識我的人就知道他越大越像.
再看看我們大家二十多三十歲的照片.
就越來越像.
為什麼呢?.
因為那些口氣.
那些表達.
那些性情.
那些樣子慢慢就越來越像.
其實就是這樣.
所以你的兒子也會像你.
只要他不要離開家裡.
或者起碼有一段時間跟你在一起.
他也會經歷感受到家庭本身.
對他的影響的因素.
那些基因的重要.
我們基督徒也是一樣的.
我們在找神.

$^{161}$知道神給我們的聖經的話語.
基督的樣式.
就是那個寶貝.
我們要像祂.
在找這個的時候我們就會很不同.
但是當我看這些小時候的照片的時候.
其實有機會去看看我的舊照片.
也有很多感觸的.
這張照片表面上.
就是一個快樂家庭的照片.
為什麼我說表面呢?.
因為是1997年拍的.
1997年當時發生什麼事呢?.
就是我等了17年.
大家記得我的生命故事.
就是做生意.
然後高高低低.
存夠錢就讀神學.
誰知道呢.
1997年大家知道了.
股災.
香港很多人跌負資產.
就是很大劑的.
我也不例外的.
但是決定留在加拿大繼續讀神學的時候.
其實以前在香港可能帶他們出入大酒店.
去拎卡佛買玩具.
現在拎卡佛沒有玩具賣了.
但是當年有的.
但是這張照片給我們看看.
回想起來就是.
沒錢嘛.
就坐Bench.
不是坐Bench.
是坐Bench.
Bench就是在公園.
不用錢的.
就跟子女去玩.
去郊外.
原來錢不是塑造他們的東西.

$^{201}$其實環境很開心.
他們就開心了.
所以其實生命也有很大的改變.
在過程當中.
神塑造了我一個生命的改變.
這張照片給我們看到.
但是我有時候也會不明白.
存夠錢了.
預備好了.
神的呼召又清楚了.
踏上去的時候.
一踏下去不是紅海開水.
是紅海蓋上來.
那怎麼辦呢.
不止這樣.
我在神學院的時候.
我有一個同學.
進了神學院的第二年.
斷了一隻腳.
撞車.
我自己的朋友也問我.
你們這些神學生.
又說神呼召你.
又說神保守.
豐盛生命.
你看看你的同學.
就是斷了一隻腳.
你怎麼解釋這件事.
什麼原因.
其實很多時候我們也會面對這些處境.
在我們的生命裡面可能事與願違.
有很多令人費解的細情.
好像現在打仗經濟的環境.
很多東西都給我們很多的困擾.
是不能夠我們控制的.
或者我們個人的遭遇.
都會面對一些很突如其來的.
很偶然的.
又好像雜亂無章的.
我們也會問.

$^{241}$神啊為什麼呢.
究竟這件事.
怎麼會發生.
我不知道你生命裡面.
你有沒有問過這個問題.
你有沒有事與願違的事情.
有沒有人際關係上的張力.
又或者你自己的計劃.
原先以為神一直帶領你一個方向.
怎麼知道呢.
當中有很多的轉折.
就算我們在建堂的事上.
如果大家看回我們那個.
建堂的網站.
這十年.
哇 多少.
其實十二的神蹟就代表了.
十二個最大的轉折.
是不是.
要不然用什麼神蹟.
但是就是因為轉折.
但是我們更深的經歷.
所以神.
是讓我們去明白一件事.
要清楚明白就是.
有時候我們生命裡面發生的事.
教會裡面發生的事.
你看到一點點發生的事.
好像不合理的.
不理解的.
令人費解的.
為什麼教會會這樣做.
為什麼神會這樣做.
為什麼對方會這樣做.
但是我們的生命經歷往往.
看似是偶然.
好像是雜亂無章.
但是我和你可不可以從當中.
看到神的工作.
那個是生命成長的關鍵.

$^{281}$也是我們一起同行的時候.
一個很重要的關鍵.
一會兒我會將我自己的心路歷程.
和大家分享一下.
但是這段經文.
是我們很熟悉的.
你買東西送給別人.
譬如說壽贈.
買個杯送給別人.
萬事都互相效力.
意思是什麼.
你如果沒了工作.
不用怕.
互相效力.
總之你不用怕.
但是有時候你也覺得.
好像.
好像潤我一樣.
我現在沒了工作.
很慘.
不夠錢.
經濟上很大衝擊.
但是你給我經文.
我一定是對的.
就是無語.
但是我們需要明白清楚.
究竟我們當了這段經文.
我們都知道.
萬事都互相效力.
叫我們神的仁德益處.
就是叫按祂旨意被召的人.
是一個對我們來說的口號.
一個杯上的一句話.
一個紀念品.
還是我們一個經歷和體會.
是深刻的.
是印證在你生命裡面.
其實我們要將這個口號放下.
縮小它.
我們要將這個經文放大.

$^{321}$幫助我們更深明白.
究竟神在這裡講什麼.
保羅跟我們說.
他說.
生命不是偶然.
不是插亂無章.
他說我們知道.
他特別用「知道」這個字.
「知道」即是什麼?.
不是說我們明白.
不是說我們相信.
不是說我們希望.
是共同集體的經歷.
以致我們一齊可以說我們知道.
對於建堂來說.
或者對於你自己生命裡面.
各種季節的事情來說.
你經歷過的事情.
就是你知道的事情.
你未經歷過的事情.
你只可以說希望.
所以讓我們都能夠看得見.
原來如果我們只是仍然將這節經文.
或者我們生命的經歷.
依為只是希望.
其實我們是需要大家.
互相去幫助扶持.
去看到你生命的這些經歷.
對你有什麼益處.
以致你能夠成為知道.
神對這些事情是有祂的意思.
是什麼意思.
不只是隨便說一句金句.
將你任何困難面對的挑戰.
放在一旁.
不去想.
而是要去反省.
要去思想.
因為我們要成為一種共同的經歷.
就像我們教會所說的.

$^{361}$十個建堂的見證.
那些是一個共同的經歷.
是知道的事情.
然後他說萬事都互相效力.
萬事就是生命裡面所有的事情.
好的壞的成功的失敗的.
對你來說不知道哪些是那些的事情.
那些事情會不會.
是一些你期盼的事情未發生.
那些是萬事.
但是那些.
譬如我剛才所說的.
同學抱著腳.
有什麼益處.
怎樣互相效力.
怎樣解釋.
有時候很難解釋.
但是如果這段經文原文告訴我們.
互相效力那個原文的字是在說.
神是主動在做一些編排整合和對齊的.
就是和祂的旨意對齊.
換句話來說就是.
神不是讓萬事隨便發生.
而是祂主動編排整合和對齊.
換句話說發生了一件事.
當它發生之後.
我們就要了解.
究竟神在這些事情當中.
透過聖經的話語.
透過聖靈.
透過我們周圍的群體.
透過環境.
究竟在跟我們說什麼話.
在造就著我們什麼樣的生命.
去成就祂在我們生命當中的旨意和心意.
你一定需要明白.
這段經文就是這個意思.
也讓我們清楚知道.
原來神是看到全局的.
祂就是電影的編劇.

$^{401}$你看到電影.
你只是看到那一刻.
你不知道後面.
就好像你看這些一點點的事情.
生命裡面的事情.
好像煲胎一樣的東西.
在頭頂上.
這是什麼動物呢?.
是人還是鬼還是動物呢?.
那時候我們不知道.
是大象還是烏龜還是鱷魚呢?.
但是.
當神來說.
祂看到全局.
祂知道這隻是.
鯉魚.
是鯨魚.
在我們生命裡面.
祂只看到這一點點.
所以讓我們有忍耐.
讓我們信賴主.
讓我們知道.
看到全局的神.
是能夠看到超越我們眼前所看到的事情.
所以我鼓勵大家.
如果邀請你.
我們保羅說.
既然我們知道是一個共同的經歷.
是肯定的.
既然萬事都是神在當中.
會要成就你和我在神面前的心意.
是叫我們得益的.
如果我們這樣相信.
我們首先就要決定.
我們不要浪費了我們的經歷.
是不是?.
那些經歷可能是很痛苦的.
你不要浪費它.
為什麼有時候會浪費呢?.
就是.

$^{441}$我們在經歷當中埋怨神.
我們沮喪.
我們沒有信心.
我們看到自己的軟弱.
但是就逃避責任.
甚至是苦毒.
怨氣.
很多時候.
我們就會浪費了神給我們的經歷.
因為靠著我們自己.
我們就這樣回應.
其實我們怎樣回應事情.
怎樣與人相處.
怎樣面對挑戰.
就正正反映著.
我們有沒有浪費我們的經歷.
因為神在當中可能是塑造著一些.
像基督的品格.
一會兒我會和大家分享我自己的經歷.
究竟你會問神.
為什麼會有這樣的事情發生呢?.
就像我剛才所分享的兩件事.
不是應該讀神學.
順風順境讀完就事奉.
或者讀神學不應該.
不應該少了一隻腳.
為什麼呢?.
但是如果你改變了問題.
你的答案就會很不同了.
如果你問的問題是.
神啊.
你要怎樣藉著這件事情使我得益呢?.
既然我們深信羅馬書第八章二十八節.
怎樣使我得益呢?.
其實.
特別是我的同學.
他當年.
三十年前.
沒有一隻腳.
他有行動限制.

$^{481}$但是他是一個電腦工程師.
當年其實.
電腦溝通沒有現在那麼方便.
但是他可以在網上製作一些聖經.
以及一些肅靈之源.
給了很多牧者使用.
造福了很多教會.
原來神真是擠在他家裡.
他要做一些更大的事情.
要幫助更多的人.
我們會明白.
但是當你知道原來神在後面.
使他得益.
也讓基督的身體得益.
神仍然掌權.
我們就會問一個合適的問題.
不會問錯了問題.
而知道認定神是掌管的.
我們也會找出.
究竟這些事情的發生對我有什麼益處.
關鍵就來了.
什麼是益處呢?.
究竟這些益處.
廣東話比較精彩.
有什麼好處呢?.
一想起好處就想起什麼益處呢?.
順境.
亨通.
事事亨通.
農曆新年也說了.
萬事亨通.
成功.
不要失敗.
安逸.
平平安安.
安安日日.
總之平平安安過一天.
心想事成.
最好是吧.
健康愉快.

$^{521}$這些都是我們想得益的事情.
但是這裡.
在原文裡面.
究竟聖經的原本意思.
叫愛神的人.
就是你和我被他呼召愛他的.
得益得到什麼益處呢?.
益處是何止呢?.
這是關鍵.
所以第29節就回答了這個問題.
第29節回答了這個問題.
因為祂所預知的人.
在創世之前.
在創造我和你之前.
在我和你出生之前.
神已經在永恆裡面.
看到有一天.
我們信主祂是預知的神.
祂也是永恆的主.
祂預定了.
知道了既然我是會回應祂.
祂也預定了一些回應祂的人.
就是我和你.
就是我們相信耶穌的人.
回應神的人.
就要效法耶穌基督的榜樣.
效法這個詞.
希臘文是在說.
不只是有樣學樣.
是會連基本的DNA.
基因也會改變.
以至你成為了另外一種基因.
這種改變.
徹底由內而外生命的改變.
是怎麼樣的呢?.
就是成為了神的兒女.
耶穌基督就是你的大哥.
你是神的眾兒女當中的長子.
我們在地上.
有父母生我們.

$^{561}$但我們成為神的兒女之後.
我們生命慢慢改變.
我們有兩個身份.
一個是地上我們家人的兒女.
我們父母的兒女.
是天上父親的兒女.
而這個身份不單只是身份的改變.
是本質的改變.
「效法」這個詞是改變我們的本質.
而神命定亦希望.
不單是希望.
祂命定,預定.
亦邀請你和我.
要和祂合作.
所以.
你是需要讓神在你的生命裡去作供.
既然我們清楚知道.
神的心意.
就是萬事都要互相效力.
就是為的.
目的就是要我們成為改變了的基因.
成為真正神的兒女.
成為耶穌基督的朋友.
耶穌基督成為我們的榜樣.
我們能夠活出祂的生命.
既然我們這樣知道的時候.
我們這樣認順的時候.
我們需要一些行動.
是讓神在我們的生命裡作供的.
是不是?.
否則的話.
你說「效法」.
每個人都說「效法」.
但是完了崇拜就忘記了剛才說什麼了.
又或者在小組裡面.
我們只是說一些社交的事情.
去遊船河的事情.
其實那些都是好的.
遊船河不是不好.
我也很尊重大家遊船河.

$^{601}$也有很多牧者都遊船河的.
所以不是遊船河不好.
只不過.
就是說只是遊船.
我的意思是只是.
或者只是社交就不夠了.
因為其實.
在我們的相交當中.
我們是要效仿基督的生命.
亦要互相提醒.
我們要深切誠實地說.
以至我們能夠成就基因的改變.
你有沒有讓神在你的生命作供?.
還是你每一次遇到各種困難的時候.
你只不過依然顧我.
自己的掌權.
其實我們自己要面對的就是.
當我們說得益處仗基督的時候.
並不是說一份更好的工作.
並不是說我們的人際關係.
夫婦關係.
弟兄姊妹的關係.
一定是順風順水.
一定是平安無事.
沒有衝突.
不是說我們的憂慮擔心.
是會一晚之間消失.
亦不是說我們的計劃.
按照我們的時間表去成就.
不是.
得益處是按照神的時間.
讓我們在過程中學習到堅毅.
學習到順從.
學習到謙卑.
學習到勇氣.
學習到純潔,公義.
愛別人的心.
願意寬恕別人的心.
亦學習到對神更深的信心.
對信仰更深的體會.

$^{641}$這一切就是學習到仗基督的過程.
你如何讓神在你生命中作供?.
你在面對不同的處境的時候.
你會變得苦獨.
或者變得越來越離開神.
或者不感受到神的同在.
還是你更加感受到神的同在.
更加經歷著神在你生命中的工作.
更加明白每一天.
你的靈修對你自己生命的鼓勵和提醒.
讓神在你生命中作供.
這幅圖就是天兒.
我最小的女兒Tiffany.
幫我的博士論文.
整個論文中的一個總結.
所以我給大家一個雞精版.
就是這個圖.
這個圖中你會看到.
我們的器皿,瓦器就是我和你.
是被神的手塑造的.
我們大家都明白這個道理.
神的手在塑造我們.
但是塑造我們是有一個平台的.
你看到有一個圓形的搖.
那些瓦器原先是一塊塊的.
就放在神的手上.
其實要透過咖啡色的圓盤來轉.
你是需要那個速度快慢.
要配合那個手的手勢.
我和師母曾經去過一個地方學塑造.
原來很難的.
塑造來塑造去.
啪一聲就扁了.
塑造來塑造去.
真的很難塑造的.
原來神真的很厲害.
要有適當的力度.
適當的速度.
還有你自己也要想著要塑造什麼才行.
不然塑造什麼出來是不知道的.

$^{681}$所以其實.
在我的論文裡面也研究.
見到歷世歷代信徒成長.
其實是需要生命的經歷和遭遇.
那些成為了我們成長的契機.
但是我們不要浪費它.
不要以為那些是經歷.
為什麼神要給我們這些.
問一下神.
透過這些怎麼幫我們.
然後有生命成長小組.
在小組裡面也不要浪費.
因為你的時間.
大家的關係.
就是對彼此的認識.
就可以將經文和你生命的遭遇.
大家互相去看到.
究竟怎樣在你生命裡面實踐.
如果在小組裡面也沒辦法幫到你.
你就找執事.
找教務.
找資深的弟兄姊妹.
你信得過他的.
他會跟你一起去同行的.
對我來說.
當我遇到這些事情的時候.
除了我自己有一些很信賴的.
一起成長的牧者朋友之外.
我也會有一個屬靈的導師.
所謂去process(處理).
將我一些想不通的事情.
經歷裡面.
生命的遭遇.
或者教他現在歷程裡面的一些問題.
我要去清理一下.
process(處理)一下.
究竟神的意思是什麼.
我怎樣去回應神.
究竟在我生命裡面在做什麼.
有時候自己也看不到的.

$^{721}$所以成長需要三個大元素.
第一就是神的話語.
上個星期Pastor Leo講過了.
所以為什麼我們要愛護神的話語.
為什麼我們要靈修.
為什麼我們要上主一學.
因為有神的話語.
這個是其中你成長的關鍵.
看到要成長什麼.
然後有聖靈.
就是將你的生命遭遇.
和你的經文和你的生命遭遇連結.
然後你還不明白的話.
或者你的心不向聖靈打開.
向聖經不夠打開的時候.
就有一些人在你身邊.
是你願意讓他向你的生命說話的.
是與你同行的.
是開放你的生命給他的.
我很開心.
我時時都跟大家說.
這幾十年的人生旅程裡面.
我遇到很多的遭遇.
很多逆境的逆轉.
但是很慶幸.
有一群很資深的屬靈朋友和我一起走.
他不是每天都和我一起走的.
因為有些是很久以前認識的.
但是就算當年不結立的朋友.
但是當有需要的時候.
一找他就立刻通.
立刻談得來.
我想你也有這些朋友.
我們也有這些朋友.
建立這些友情.
幫助我們成長.
就好像這個短片一樣.
其實就是怎樣去塑造的呢?.
就是你會看到.
神的腳.

$^{761}$神的手.
神的腳在將我們生命的遭遇.
生命的環境,處境,困難,挑戰.
就在他的腰上.
速度和手的壓力是需要不同的.
還有他的腰臟知道他現在要做些什麼.
你看到其實一開始是一塊的.
做了一塊之後.
他去了一個杯.
杯中間變成一個.
很多不同的形狀.
就在看到我的生命要變成怎樣的神的手上.
以致那個速度能夠配合.
我們不明白的時候就要問別人.
你不要覺得不明白的時候就放棄.
或者不明白就是因為不明白.
不是的.
其實你就需要憑神的話語.
靠聖靈和身邊的人.
幫助你能夠明白,體會.
以致你的生命是讓神在你的生命裡面作供.
所以你要問自己.
你有沒有一些能夠向你生命說話的朋友.
有沒有一些你遇到處境,困難,想不通的事情.
處理不了的關係的時候.
你可以找到他.
或者他們.
是能夠和你同行.
去看到神的心意.
對你的心意呢?.
或者我可以用幾分鐘時間和大家分享一下.
我自己也問神一個問題.
主啊.
你在這一刻.
在我的生命裡面.
塑造著哪一種基督的樣式呢?.
有時候這個問題.
問一些信得越久的人是越難的.
因為有時候我們有一種屬靈驕傲.
我不是在說你.

$^{801}$我是在說我自己.
那個屬靈驕傲就是.
好像也有點遲了.
起碼我也在做牧師.
有祈禱,有讀經,有事奉.
什麼都做齊了.
但是.
究竟神還在做什麼呢?.
信心好像有一點.
也不是完全.
不過有一點.
每樣都有一點.
那就OK了.
其實在做什麼呢?.
有時候也很模糊的.
其實神是在塑造著我的.
其實我也希望未來幾分鐘.
我和你分享一下.
你不要說牧師你又要講建堂了.
我就是要講建堂.
為什麼呢?.
因為建堂在我的生命裡面產生了.
建立生命的經歷.
那個是真實的.
所以我才要告訴你.
其實真的是.
大家看見那些陶像的手和搖.
其實一直在我來了恩泉十年裡面.
頭幾年就說.
立刻要建了.
開心到Spencer當時做主席.
第二年就要建了.
第二年第二年.
第二到第十年才建.
真的是很久了.
那麼神在我的生命裡面做什麼工作呢?.
我有什麼調教呢?.
我打算做完三年.
建完堂就可以.
不關我的事了.

$^{841}$但是為什麼繼續也要關你的事呢?.
其實是不可以不關我的事的.
那是一刻間的決定而已.
但是如果不信服神.
真的是很大的一件事.
所以你一定要知道.
神究竟在你生命裡面做些什麼呢?.
長話短說.
近期一點.
就是神在我生命裡面.
問祂.
究竟你現在在做什麼處境.
還是在搞什麼呢?.
其實如果你上網站看到12個經歷.
12個神蹟.
就是每一個都在做我的生命的一部分.
不過我想說最近.
也具體一些分享.
就不會這麼模糊了.
近期一點就是.
去年很開心.
終於開動土禮了.
但是動土禮的時候.
其實我們還沒有銀行的正式批核.
如果真的不批核.
那840萬怎麼辦呢?.
是不是?.
又或者.
究竟應該動土還是不動土呢?.
究竟動土之後.
究竟是否應該先等一等呢?.
等多久呢?.
其實有很多元素的.
但是神究竟在我的生命裡面.
或者我們領袖也有不同的意見.
但是我們領袖當中有不同的意見.
我們怎樣去磨合呢?.
怎樣去互相等候呢?.
怎樣去我在當中謙卑事奉.
服善呢?.

$^{881}$以至在神裡面去清楚神的心意呢?.
神究竟在那一刻在我身上做什麼工作呢?.
在我們的教會裡面在做什麼工作呢?.
其實需要常常去思考.
所以呢.
既然動土禮是做了.
8月3日大家一鼓作氣.
於是就畢地了.
因為8月3日大家投了一票.
大家記不記得投了贊成票呢?.
應該都記得的.
你投了贊成票.
有印證就去畢地了.
畢到本能坑.
有很多不同的事情發生的.
但是.
就是一路走.
我很喜歡Pastor Ada說的那句話.
叫做.
行步見步.
你知不知道行步見步的信心.
和見步行步的信心有什麼不同呢?.
見步行步的信心就是看到銀行的批核來了.
看到我們籌款的500萬已經到位了.
看到天氣沒有這麼大雪.
還要在大雪之前已經蓋了一個地基.
什麼都清楚的.
然後我們就走前一步.
那叫做見步行步.
行步見步就是.
神開前面的一步.
你知道肯定是神的開路.
那一步而已.
看不到下一步.
就好像剛才那些圖片那樣一點點的.
你就一起.
透過彼此的辨識.
為什麼我們教育同工.
為什麼執事同工那麼重視那個屬靈辨識.
就是辨識之後.

$^{921}$就踏前一步.
就是走到那一步.
然後才再行步之後.
就看到下一步.
我們過去就是這樣經歷的.
神就是塑造我們這種信心.
於是11月很開心.
因為全年我們一鼓作氣.
你都經歷了.
如果沒有羅華我知道.
你都知道了.
為什麼.
因為我們籌款不錯.
年尾250萬到位.
籌款晚宴又很開心.
700多人來.
很開心.
每樣都好像表面上不錯.
但是到了年尾還是250萬.
就有點沮喪.
你有沒有沮喪.
我不知道你.
不過我也有一點.
為什麼呢.
既然神在前面.
又做了那麼多事情.
你走一步又看到一步.
誰知道走著走著.
忽然間下一步看不清了.
走不走下去呢.
就是神要給我們的考驗.
我自己來說.
11月.
我就覺得大家都很慘.
我自己都覺得很慘.
我都在神面前祈禱.
我說大家最好不要一起那麼慘.
反正我也有些朋友.
不如回香港看看.
可不可以求主憐憫.

$^{961}$在香港幫幫忙.
於是我就回香港了.
兩個星期.
我希望大家為我祈禱.
也希望.
如果250萬.
10個人.
每人25萬.
不是很多.
在香港來說.
你也知道房子是3000萬港幣.
就是10\%的房子而已.
應該如果他有資產的話.
他不會覺得這個很肉質.
那麼.
只要有一個意象.
最重要是神感動他.
OK.
所以我也很屬靈的.
神感動他.
就應該他會給的.
但我也上兩個星期和大家說過了.
回來很沮喪的.
為什麼呢?.
一光熱城.
因為已經看到神的工作那麼偉大.
應該就踢入龍門了.
回去配合.
於是回來和大家說.
大家不用籌款了.
多開心.
你想不想聽這句話.
我想你建堂的知識最喜歡的就是.
從現在開始到建堂.
到25年後.
我們都不需要還錢了.
也不需要再籌款.
想不想聽這句話.
當然想啊.
誰不想呢?.

$^{1001}$我也希望是這樣.
我當然想啊.
我也不想將爛攤子交給下一手了.
但是.
下面是我想的.
其實後來11月.
去了兩個星期.
12月初回來.
我的認罪悔改.
神告訴我.
透過這段經文在靈修裡面.
在靈修裡面神說.
倚靠人.
是很嚴重的.
是被人咒詛的.
神咒詛那些倚靠人的人.
更加是說.
不見福樂來到.
他要住在曠野乾旱之地.
無人居住的炎地.
你聽了也覺得很慘.
你未必明白什麼是炎地.
你沒有住過.
不過總之也覺得很荒涼.
什麼都沒有.
很慘.
這段經文在我靈修的時候.
真的向我心裡有很大的說話.
就是說.
Hubert 究竟.
其實你很屬靈地做這件事.
你很想為弟兄姊妹.
我也知道弟兄姊妹說.
教目已經五年工資捐出來了.
弟兄姊妹都犧牲性的奉獻了.
又是認識的人都叫來晚飯了.
鄭欣存也所有認識的人都釋出給他就去囚福.
捐完了.
但是還欠那個.
但是我們建築成本一千四百萬.

$^{1041}$裡面就是欠那個.
所以我們叫做.
趴越差距就好像一條橋.
欠一個就去到銀行那裡.
神都不可以做一件事嗎?.
但是神就是不做.
神開始讓我看到.
其實我的思想錯了.
當我為大家想的時候.
就變成我好像是神.
我覺得你不用做不如讓我做.
由我做了之後讓你回來開開心心.
最好聽的就想聽.
其實是不是神的心意呢?.
所以.
後來這段經文也說.
「依耶和華為他所依賴的那人就有福了.
他就好像樹栽於水旁在河邊紮根.
炎熱來到毫不察覺.
葉子也必清脆」.
很漂亮吧?.
你看到那個對比嗎?.
「在乾旱之年才厲害.
一無掛慮並且結果不止」.
我聽到這個.
我多開心啊.
但是神的話語.
你相信祂你才開心.
是不是?.
你不相信祂也不會開心.
因為你看到的就是那250萬還沒到位.
就不是看到那250萬到位的那個.
到位的是之前的.
現在後面還有還沒到位.
就看到那些人的反應很不熱烈.
各種原因.
詳細的聽我今期的《牧者心聲》.
明天會上載在那個應用程式裡面.
你會詳細聽到一些比較詳細的見證.
但是.

$^{1081}$原來神看我們的心.
這段經文裡面也說.
於是我心改變.
我就跟神說.
其實不是我一個人的事.
不是我要為你做什麼.
不是我的祈禱好像很屬靈那樣.
天天就問神.
是不是透過這十個弟兄姊妹可以搞定.
神就告訴我.
其實你以為很屬靈.
其實你是搞錯了一點.
差一點.
我現在就告訴你這十個不是你倚靠得到的.
那你還可以倚靠什麼呢?.
就回來倚靠神.
祂禱告的方向不同了.
明白神的方向不同了.
我也問神.
那什麼都做了還有什麼可以做呢?.
這個也是常常要問的問題.
認識的人都認識了.
都問了.
我當時回去就問了.
250萬.
5個人搞定.
6個人也可以搞定.
8個人也可以搞定.
不過太大牛腩了.
大家看到就覺得很害怕.
現在香港的環境不像現在那樣.
但是神改變我的心之後.
當我回來之後.
我就完全不同了.
神一直給我的經文是.
耶和華吩咐摩西說.
你要吩咐以色列人獻禮物給我.
其實神邀請我的就是我要去吩咐.
然後凡是甘心樂意的就給我禮物.
給神禮物.

$^{1121}$你們都可以放下.
都可以收下.
所以最重要的關鍵是甘心樂意.
我不知道在會眾裡面.
還有哪些是甘心樂意的.
但是只有神知道.
但是我的責任就是要吩咐.
其實老實說這也是掙扎.
其實做牧者很多時候都要收集民意的.
民意就是說說得太多了.
民意就是說你沒有錢了.
那就捐完吧.
捐完就沒有了.
我自己也覺得自己捐完就沒有了.
朋友圈都給你了.
但是神就責備.
你這樣就不對了.
神裡面祂有封聖.
你的完結就是神的開始.
祂說你們都可以收下.
所以祂說他們要為我造聖所.
使得我住在他們中間.
是神自己要共衷作.
如果我不讓弟兄姊妹一起知道.
一起同工一起合作去做這件事.
只是回香港搞定回來已經搞定了.
其實是抹殺了你參與和經歷神的機會.
以色列人無論男女.
在經文裡面說.
凡心裡受感動的都帶甘心祭禮獻給耶和華的.
為要耶和華藉著摩西所吩咐的一切.
供他們去做.
因為.
理論上那裡說.
我算什麼?我的百姓算什麼?.
竟能夠如此樂意奉獻.
因為萬物都是從神你而來.
我們將從你手得來的獻給你.
我和你.
我都要再一次想想.

$^{1161}$我們手得來的還有些什麼呢?.
我們有什麼是自己還是擁抱著.
不願意放手給神的?.
還有些什麼是我們的關係網絡裡面.
我們還沒交給神使用呢?.
耶和華我們的神.
我們預備了許多的材料.
為你的聖名建造奠禦.
從你的手而來的.
都是屬於你的.
所以.
親愛的弟兄姊妹.
其實我這個掙扎不是製造出來的.
這四個月沒有機會跟大家分享我自己的生活歷程.
其實是需要克服這些的.
你以為我不想做一個.
其實我們每一個教育都是這樣.
每一個領袖當然想.
最好不用做一些真是討人厭的事情.
伊斯蘭經常說.
又來我再跟他說.
沒有啦.
雖然他很樂意的.
他其實很樂意的.
但是.
究竟我們看神還是看人呢?.
究竟在我們的回眾當中.
有哪些是神已經為我們預備而甘心樂意.
我們不知道的.
我們唯有告訴你.
然後讓你自己去回應神.
對不對?.
所以後來.
這段經文.
就讓我們看到.
我先不詳細讀了.
就是.
心裡有智慧心受感動.
凡有智慧都能做聖所.
所有一切的功都各自離開原本的工作前來.

$^{1201}$你看到.
見堂見聖殿是整個全民皆兵的圖畫.
不是一個人的事.
不是摩西一個人的事.
紅色的字就是.
每天早晨百姓繼續的將甘心祭拿來.
其實經文多處都講甘心祭.
已經超過了分庫建造經驗所需的.
所以是大眾一起的事.
我們大家都閉上眼睛去做的時候.
原來神就學習.
和我們自己.
加上我們所有的網絡.
其他的教會.
就能夠一起成就.
因為.
最後說.
他們所有材料已經足夠整個工程之用.
而且有餘.
這個你和我相不相信呢?.
所以去年.
我們就開始了一個禱告運動.
其實很坦白講.
參加的人是有一部分.
但是我絕對不是去論斷那些沒有參加.
因為你有你的原因.
但是我可以和大家很坦白的分享.
當我看完這幾段經文.
當香港的經歷回來.
我心想.
神就是要我們全間教會.
悔改.
每一個人都有份.
你的份就是禱告.
禱告到有些事情發生.
禱告到神會彰顯到.
祂的心意在你的生命裡面有些什麼可以做.
於是神真的回應了.
有些人突然間會為我祈禱.
我也回到香港.

$^{1241}$3月的時候.
有另一個機會回去.
不是故意回去的.
是有一個機會回去.
於是做了兩個星期.
既然要回去參加一個慶典.
我順便多做了兩個星期.
這次神告訴我.
你媽媽以前的阿姨.
我跟神說.
很久沒見他們了.
小時候.
她說「你什麼歲」.
「都挺行的」.
「你媽媽的兒子知道」.
我姐姐也過了10年了.
我姐姐已經是團友.
他們都是很熟悉的.
因為我以前是一大一小.
我姐姐是最大的.
我是最小的那一個.
拉仔.
所以就照顧我.
所以她的朋友我都認識了.
她經常帶著我出街.
她玩的東西.
我都會跟著她一起玩.
我問神.
熟悉的我都問了.
不熟悉的我都問了.
但是她認為跟我熟悉.
但是我跟她真的不熟悉.
哈哈哈.
跟我家族熟悉是與我無關的.
那我怎麼問呢?.
但是神說.
這段經文我之前也和大家分享過.
最要我心的就是.
原來有一個朋友來而已.
是不是?.

$^{1281}$有一個朋友來家裡而已.
就在隔壁的鄰居敲門.
吵醒別人.
其實很明顯別人不開門.
是對的.
弟弟睡了一覺.
上床了.
大哥.
你這個朋友的朋友.
就好像我媽媽的朋友.
都老到八十歲九十歲了.
你明白嗎?.
然後她的子女.
都已經七十歲了.
然後我姐姐的朋友都六十幾七十歲了.
全部都是不同的生命季節.
但是神說.
就是朋友的朋友.
於是她怎樣呢?.
真的敲門.
那你怎麼可以令到你的朋友的朋友.
真的反應呢?.
這裡就說了.
不要打擾我.
我門已經關了.
孩子們也和我在床上了.
我不能起來給你.
十一月的朋友.
已經告訴我.
真的有大部分都說.
他們的生命季節改變了.
他們的困難了.
香港的經濟改變了.
他們現在有很多問題了.
我就從一個牧羊角度去關心他們.
為了祈禱.
但是.
都是解決不了我的問題.
所以為了祈禱都是要解決問題.
但是這段經文.

$^{1321}$再一次提醒我.
他們會不會再問一次他們呢?.
這次我真的回香港又再問一次他們.
十個裡面有五個肯見.
我心想.
上次已經和你談完.
你夠膽蓋戒我嗎?.
那我繼續為他祈禱吧.
因為我不會藉口弱質告訴他.
見頭有事我再一次和你談談.
他真的見了.
有一個真的.
明天你聽我《牧野心聲》會詳細講那些經歷.
五大見證.
不如不要今天講了.
今天沒時間了.
有一個真的很有反應.
是上一次見.
但是這次才有反應.
我也想不到他會有反應的.
他會有反應的.
就在我《牧野心聲》裡面.
明天大家可以聽.
但是聖經說.
儘管他不理你.
但是.
我告訴你們.
耶穌說.
「雖不因他是朋友起來給他.
也會因他不顧面子地籍求.
起來照他所需要的給他」.
這當然是說到我們和神的關係.
我們需要向神求.
但是同樣地.
告訴我們.
如果神感動我們.
所以我在三月回去那一次.
我每一天都是很長的時間.
有一個小時為我今天會見的一些朋友.
去禱告.

$^{1361}$看看神的意念.
有些說不如今天不要見.
或者見到也不要跟他講.
有些叫我跟他講.
神會告訴你的.
真的.
應該跟他.
這樣就完全不同了.
這次也做了一個要求.
180萬.
如果9個人是20萬.
18個人是10萬.
36個人是5萬.
72個人是25,000.
其實每一個朋友的心裡面都有一個數目.
可能有一部分的朋友是這樣.
但是有些朋友是25,000以下的.
就把連結發給他.
希望他捐.
但是有些朋友你覺得神給你祈禱.
你祈禱之後神給你感動.
我就按照這個感動.
有些我就會問他.
你會不會是那九分之一.
或者十八分之一呢?.
有些我就會問他.
你會不會是那三十六分.
或者十八分呢?.
有些我會問他是否那三十六分.
或者三十二分呢?.
不過最後我會跟他說.
唉.
你也知道.
朋友就是.
三十六個比較容易.
七十二個真的有回應的朋友比較難.
所以你會不會祈禱.
想想神會不會感動你.
就成為我們恩泉那三十六個建堂夥伴的其中一個.
去幫到我們.

$^{1401}$真的拉近我們那個建堂的籌款距離呢?.
其實不是說策略.
我講這些分享就是.
神就給我智慧.
給我體會.
給我引導.
怎樣去面對和處理.
於是有一個姐妹就跟我說.
我就是為你們一尊鄉下.
我願意去配對.
其實她在香港為什麼要做配對呢?.
但是她要鼓勵我們弟兄姊妹.
因為她就看到以上.
我看到的經文.
就是弟兄姊妹要被興起.
我們要同心.
如果別人也看到我們自己在這個恩泉家裡面.
有怨氣的.
有失望的.
有沮喪的.
有不信的.
神更加看到.
求主憐憫我們.
當我們願意被塑造基督的生命的時候.
我們需要清理我們自己.
改變我們自己.
神在這個過程裡面建造了我們什麼呢?.
其實就建造了.
讓我知道.
依賴神是有福的.
而且是結果不止的.
三月的結果上個星期也分享過了.
我不再重複.
是很大的恩典.
但是看到這些經歷的時候.
神要我學什麼呢?.
如果我很直接和大家坦白說.
就是謙卑.
原先的計劃可能不是.
我以為很神聖.

$^{1441}$但原來不是.
以為已經叫了.
原來神又會再有新的可以被邀請的對象給我.
以為那些人拒絕了不會回應.
神又讓我看到.
原來要依靠祂.
有信心.
謙卑.
實踐耶穌所說的不顧顏面.
就是願意去問.
去邀請.
讓人去回應.
不是預判了他們的回應.
所以我們的生命.
其實你遇到任何的處境關係.
和我剛才所經歷的.
不僅僅是建堂的經歷.
就是每一個經歷你不會浪費的時候.
你就回到神的面前.
最後就是講到.
你會說.
如果我真的不明白神的心意.
我又很軟弱.
我又真的很不信.
那我怎麼辦呢?.
我鼓勵你.
在結束的時候.
就是第三十節所說的.
祂所預定的就是你和我.
因為我和你都是預定祂.
神在千古以前已經看到我和你.
是會相信祂.
回應祂.
祂召了我們來.
祂就稱我們為二.
以致我們能夠成聖.
我們成聖的過程.
越像基督.
我們就越能夠榮耀祂.
這是神學裡面所說的「金鎖鏈」.

$^{1481}$神已經允許了.
這五個動詞都是過去式的.
換句話說,神已經拯救了我們每一個.
祂是會保守我們到底的.
祂不會遺棄我們.
祂有一個護航系統的.
只要你相信.
只要你不要放棄.
只要你仍然在神的家.
只要你仍然與神在一起.
神是會保守你.
一直帶領你去到底.
如果你覺得你的信心不足.
我很軟弱.
你就求主憐憫你.
就像一個人掉進了急流裡一樣.
我不知道你是否掉進了急流.
你在游泳.
就像我當日掙扎劍塘一事也是游泳.
但是神要你握緊祂的救生繩.
當一個人掉進急流衝擊的時候.
你多大的體力都無法游上岸.
給你一條繩子你就握住.
你可以放鬆身體.
握住.
最重要是握得緊.
神就拖你上岸.
你就讓祂.
你就回應祂.
你與祂合作.
與祂同工.
所以結束的時候我想鼓勵你.
你抓緊神的應許.
你每一個經歷.
包括我們教會的經歷.
你不要浪費你的經歷.
因為神是掌權.
祂應許.
每一個經歷是祂在當中的編織者.
也要讓神在你的生命裡面作供.

$^{1521}$因為神正在無助你.
你要問合適的問題.
以致你能夠抓緊神.
去到永恆.
所以在你的生命裡面.
如果你真的這麼相信.
你也知道這幾節經文.
你是靠著它.
繼續被神塑造你的生命.
你問的是兩個很重要的問題.
就是在每一個處境裡面.
你問神.
你怎樣藉著這件事使我受益呢?.
受益就是像基督.
所以第二個問題就是.
主啊.
你此刻正在我裡面.
塑造著哪一種基督的樣式呢?.
如果你懂得問這兩個問題.
就算你沒有答案.
你可以找其他人回答.
可以找屬靈導師回答你.
並且你抓住神的救恩繩.
信得過祂帶領你到底.
你就可以一直保守你.
走到底.
所以你的生命.
你的人生.
就算是我們教會的經歷.
也不是偶然的.
你的掙扎也不是枉然的.
你和我都在掙扎當中.
我們都需要成長的.
但是神正在塑造著你.
活像基督.
活出基督.
神也要保守你.
只要你和祂同供.
一起同心同意.
求你憐憫我們.

$^{1561}$求你的靈去激勵我們.
督責我們.
因為神你掌權.
我們願意回應.
願以順服.
祈禱奉靠著耶穌基督的名銜.
阿們.
\newpage



\section{創世記 13:6-18}
\label{sec:031883}
\textbf{世上最「正」的眼鏡}~~~~江惠明牧師
\newline
\newline
連結: \href{http://www.church.com.hk/acms/content.asp?site=cdc\&op=show\&type=product\&code=031883\&layout=sermon}{\texttt{031883}} ~~~~ 語音日期: 2026-03-29
\newline
\newline
\hyperref[sec:031882]{\small{< < < PREV SERMON < < <}}
~
\hyperref[sec:index_chronic]{\small{[返順時目]}}
~
\hyperref[sec:031884]{\small{> > > NEXT SERMON > > >}}
\newline
\newline
創世記 13:6-18
\newline
\begin{longtable}{cl}
\hline
\hline
章節 & 經文 (和合本修訂版)\\
\hline
13:6 & \begin{tabularx}{0.7\textwidth}{X} 那地容不下他們住在一起；因為他們的財物非常多，使他們不能同住一起。 \end{tabularx} \\ \\ \relax
13:7 & \begin{tabularx}{0.7\textwidth}{X} 當時，迦南人與比利洗人在那地居住。亞伯蘭的牧人和羅得的牧人之間起了爭端。 \end{tabularx} \\ \\ \relax
13:8 & \begin{tabularx}{0.7\textwidth}{X} 亞伯蘭就對羅得說：「你我不可以相爭，你的牧人和我的牧人也不可以相爭，因為我們是一家人。 \end{tabularx} \\ \\ \relax
13:9 & \begin{tabularx}{0.7\textwidth}{X} 遍地不都在你眼前嗎？請你離開我吧！你向左，我就向右；你向右，我就向左。」 \end{tabularx} \\ \\ \relax
13:10 & \begin{tabularx}{0.7\textwidth}{X} 羅得舉目，看見約旦河整個平原，直到瑣珥，都是水源充足之地。在耶和華未毀滅所多瑪、蛾摩拉以前，那地好像耶和華的園子，又像埃及地。 \end{tabularx} \\ \\ \relax
13:11 & \begin{tabularx}{0.7\textwidth}{X} 於是羅得選擇了約旦河整個平原。羅得往東遷移，他們就彼此分開了。 \end{tabularx} \\ \\ \relax
13:12 & \begin{tabularx}{0.7\textwidth}{X} 亞伯蘭住在迦南地；羅得住在平原的城鎮，他漸漸遷移帳棚，直到所多瑪。 \end{tabularx} \\ \\ \relax
13:13 & \begin{tabularx}{0.7\textwidth}{X} 所多瑪人在耶和華面前罪大惡極。 \end{tabularx} \\ \\ \relax
13:14 & \begin{tabularx}{0.7\textwidth}{X} 羅得離開亞伯蘭以後，耶和華對亞伯蘭說：「你要從你所在的地方，舉目向東西南北觀看； \end{tabularx} \\ \\ \relax
13:15 & \begin{tabularx}{0.7\textwidth}{X} 你所看見一切的地，我都要把它賜給你和你的後裔，直到永遠。 \end{tabularx} \\ \\ \relax
13:16 & \begin{tabularx}{0.7\textwidth}{X} 我要使你的後裔好像地上的塵沙，人若能數地上的塵沙，才能數你的後裔。 \end{tabularx} \\ \\ \relax
13:17 & \begin{tabularx}{0.7\textwidth}{X} 你起來，縱橫走遍這地，因為我必把這地賜給你。」 \end{tabularx} \\ \\ \relax
13:18 & \begin{tabularx}{0.7\textwidth}{X} 亞伯蘭就遷移帳棚，來到希伯崙，幔利的橡樹那裡居住，在那裡為耶和華築了一座壇。 \end{tabularx} \\ \\
[1ex]
\hline
\hline
\end{longtable}
$^{1}$鄭英姊妹平安.
我一次去眼鏡舖.
看到很多的眼鏡.
今日的眼鏡可以說是多得不得了.
很多的款式.
我見過有一次.
我的朋友戴著眼鏡.
然後我戳一下他.
咦 為什麼沒有鏡片.
原來這些是fashion來的.
我又記得很久以前.
我都是戴隱形眼鏡.
那時候隱形眼鏡都很多顏色.
如果你想自己扮歐洲人.
你可以配褲藍藍地的隱形眼鏡.
不過我想你可能要染頭髮才可以.
否則都扮不像.
我現在戴的這個眼鏡.
有防電腦的藍光.
又防紫外線.
又可以有近視.
又有防老花.
很多功能.
我這個應該是四合一.
我不知道在座有沒有你的眼鏡比較厲害.
五合一的有沒有.
還有什麼散光.
軟視.
還有一些不知道什麼.
有沒有.
超薄.
我這個可能厚過石仔那個.
一會兒量一下他那個多厚.
今日的眼鏡真的很厲害.
但是無論世界上的眼鏡有多厲害.
你那副有多正.
我相信都不是最正.
今日我就要推介一個.
世界上最正的眼鏡給你.
希望從今日起.

$^{41}$你戴著它.
你去到處.
好像剛才那隻蟹.
可以打橫行都可以.
你的人生也都會好.
現在我們就透過今日的經文.
一起去找這副眼鏡出來.
剛才讀的經文的故事.
是講到亞伯蘭和羅德的故事.
他們的關係是叔侄.
話說神父召亞伯蘭離開本地本族父家.
去到神指示他要去的地方的時候.
他就馬上遵從神的吩咐.
聖經說羅德都和他一起去.
即是羅德一直跟著叔叔一起.
跟著他一起生活.
後來亞伯蘭因為迦南地發生饑荒.
他就去到埃及.
在埃及發生一些事.
最後埃及王就把很多牲畜奴婢賜給他.
之後他就回去迦南地.
當他回去迦南地的時候.
其實他是更多的財產.
他的財產很多.
他的侄子羅德都得到不少.
在經文第七節說到.
他們回到那個地方之後.
因為他們太多的牲畜.
很多的財產奴婢.
那個地已經容不下他們兩個一起住在那裡.
因為還有很多牲畜要散佈在那個地方.
所以他們不能夠再住在一起.
牧養亞伯蘭牲畜的和牧養羅德牲畜的.
起了糾紛.
甚至幫他們看牲畜的工人.
都因為這些事經常有紛爭.
阿叔就跟侄子說.
你我之間切不可紛爭.
你的牧人和我的牧人之間都不可以紛爭.
因為我們是至親.

$^{81}$全地不都是在你面前嗎?.
請你離開我吧.
你向左我就向右.
你向右我就向左.
亞伯蘭他是顧念親情.
所以他主動來講出.
不如我們分開吧.
你先選擇.
這位長輩真的不能說.
他看親情比財富更加重要.
他寧願放下長輩的身段和特權.
讓後輩先來作出抉擇.
可見他其實是一個很慷慨大方.
和對神有信心的人.
羅德就老實不客氣.
他忘記了這個叔.
畢竟都是長輩.
請我今日有今時今日的事.
其實都是叔帶協我.
那個地是應許給亞伯蘭的.
他應該是感恩圖報.
不過他沒有.
他忘記了.
總之叔叫我選擇.
我就選擇吧.
接著他們作出的選擇是怎樣呢?.
羅德是先選擇的.
他看到約旦河整個平原.
整個都是平原.
平原當然好.
可以耕田,起屋.
沒有理由選擇山.
慢慢他也住在平原的城市裡.
經文說「直到所多」.
他在城市裡定居了.
以前過的是遊牧的民族.
到這一刻他定居了.
而且那個地方越來越近所多.
從經文看.
羅德什麼都從自己的角度想了.

$^{121}$從自身的利益去想了.
甚至他貪圖一個安人.
相反亞伯蘭是之後才選擇的.
他同樣舉目.
神叫他舉目.
你向東南西北.
你望東南西北的地方.
後來他的生活形態.
仍然是一個遊牧的生活.
經文說他一直來而張鵬.
他來到希伯倫.
希伯倫其實是大衛初登王座時的地方.
他的生活態度就是信靠神敬畏神的態度.
怎樣看得出呢?我們慢慢看下去.
怎樣看得出亞伯蘭過的是信心的生活呢?.
亞伯蘭他一直過著千載不定的象棚生活.
他隨時要起身去第二個地方.
所以即是說其實他是不斷地搬的.
你喜歡搬屋嗎?.
不喜歡.
搬一次真是害怕.
最近聽到一位宣教士說.
他們以前每兩年就要搬一次.
搬了無數次.
而且是國家去國家.
然後城市去城市.
然後真的在本地都不斷地搬.
其實每一次的起步.
再踏出安書區.
都是不容易.
是一個很大的挑戰.
而且每當亞伯蘭去到一個新的地方.
他就在那裡為耶和華築一座祭壇.
可見其實他對神常常有一份敬畏的心.
不過可能我們都或者會想幫羅德平反一下.
他選擇去平原有什麼不妥呢?.
其實我們很正常.
我們正常人都想過得安樂一點.
沒人想拿虎離身的.
羅德的選擇看來是人之常情.

$^{161}$不過在第十三節這樣說.
在耶和華看來.
所多瑪人是罪大惡極.
問題就在於.
他漸漸樂而張鵬去到那個地方.
漸漸靠近了所多瑪.
其實他知道所多瑪是一個嚴重犯罪.
得罪神的地方.
但他仍然沒有看重這件事.
是不斷的靠近他們.
羅德是按他自己的心意去舉目.
他的選擇是按自己個人的喜好.
和世俗以為美好的奴.
約旦河是整個平原.
又有水灌溉.
看來好像伊甸園一樣.
誰不喜歡呢?.
但相反.
亞伯蘭是按照神的意思去舉目.
大家都舉目.
但亞伯蘭是按神的心意去舉目.
他舉目看的是多大的範圍?.
不只是這裡.
羅德看到前面一片的平原.
但神叫他舉目東西南北三百六十度.
因為羅德和亞伯蘭帶了不一樣的眼鏡去看.
去舉目.
亦帶來了不一樣的結果.
羅德選擇了約旦河的平原後.
他漸漸向東移動.
而且住在平原的城市中.
在那裡安頓下來.
一直到所多晚.
本來一個選擇舒適安穩的地方是沒錯的.
但錯就錯在.
他明知神不喜悅那個地方.
他繼續向罪惡之城靠近.
他不理會神的看法.
逐步進入所多瑪這群搏抑神邪惡的人中.
聖經在《創世記》十八章十九章告訴我們.

$^{201}$因為所多瑪這個城罪惡極重.
所以神差派使者去毀滅這個地方.
那時羅德還捨不得離城.
最後是神的憐憫.
天使拉住他和太太和兩個女兒逃出城.
然後他們才可以免了一死.
可惜他的妻子因為沒有聽天使的吩咐.
逃命的時候仍然留戀那個舊地.
當她一回頭就變成一根鹽柱.
而後來羅德都因為他兩個女兒.
想替父親傳宗接代.
結果亂倫而產生了後來的兩個民族.
摩訶和阿曼都是拜偶像搏抑神的人.
而相反亞伯蘭他的巫師好像吃虧了.
但他卻得著了神的賜福.
他看到的不是他自己的眼光.
而是神的眼光要他看一個更廣大無邊界的地方.
神說你看見的地方.
我都要賜給你和你的後裔直到永遠.
這是神的保證.
只要他的信心有多大.
神就賜給他多大的境界.
不單是賜給他一個人 而且是賜給他的子孫後代.
直到永遠.
這裡的子孫後代 你的後裔指的是什麼呢?.
這個我們可以從加拉太書三章六到八節看到.
我請大家和我一起讀 預備起.
正如亞伯拉罕信神 這就算為他的義 所以你們要知道 有信心的人就是亞伯拉罕的子孫.
聖經既然預先看見神要使外族人恩信稱義.
就預先把好消息傳給亞伯拉罕 萬國都必因你得福.
亞伯拉罕即是亞伯蘭 後來他改了名叫亞伯拉罕.
這裡 你的後裔指向他屬天的子孫.
因信成為亞伯拉罕的後代.
今天 你是不是亞伯拉罕的後代?.
有沒有人說不是?.
今天我們都因著信同為亞伯拉罕的後裔.
但我們不是猶太人.
所以這裡的後裔指向我們不是血統血緣上的子孫後代.
是每一個因信耶穌基督而成為亞伯拉罕的後代.
神藉著亞伯拉罕將福氣賜給我們.

$^{241}$使我們這些外邦人都因為亞伯拉罕的子孫耶穌基督.
得以成為他的後代.
要踏出安書區真的很不容易.
我都回想我自己過去的日子.
如果當年二十多年前.
如果當年我不信服神的帶領.
回到中國內地工作.
我都會錯過了很多經歷上帝的信實的機會.
我的生命也不會快速地成長起來.
而且我都會失去了很多可以用來鼓勵其他人.
和見證神的寶貴人生經驗.
今天是全因堂第九屆的猜拳年會.
也是我們搬到這個新的地方的第一次猜拳年會.
我想起過去一年所經歷的一切.
我真是為神奇妙的帶領而感恩.
當我們站在一個搬還是不搬的十字路口的時候.
神保守了我們的心,我們沒有停留在柴灣舒舒服服地過日子.
而是我們願意像亞伯蘭一樣來而將彭.
神叫我們起行,我們就起行.
和你們每一位對我的信任.
而且有這一份勇氣和膽量一起走出來.
當然其實最重要的是.
神每一步其實是祂自己親自開路和帶領著我們.
我們是同心信服在神的心意之下.
作出這個搬塘的舉動.
走出來不是一定事事順利.
甚至乎是困難重重.
但是只要認定是神的心意.
我相信神必負責到底.
正如我當年不貪戀在香港一個穩定舒適的生活.
我一個人回到內地工作.
雖然遇到很多的挑戰和艱難.
寶珠也來過我那裡.
但是那些經歷卻成為了我深刻經歷神和生命成長的契機.
也成為我人生很美好的回憶.
今天神也透過亞伯蘭的事蹟去鼓勵我們.
不要用自己的眼光來看世界.
而是要戴上上帝的眼鏡去看萬事.
神說你要舉目.
你看見的地方我都要賜給你和你的後裔.

$^{281}$弟兄姊妹,今天我們在北角,我們看見什麼地方?.
只要我們看見,神說我都要賜給你.
但是祂說你要起來,縱橫走遍這地.
因為我必把這地賜給你.
神要使用我們成為很多人的祝福.
但是我們必須要起來.
這裡說你要縱橫走遍這地.
就算我給了你這麼多地方.
但是你坐在這裡也無法得著.
我們要起來.
弟兄姊妹,你願不願意為主而起來?.
我們不是單單眼睛要看.
並且我們的腳要走.
我們的信心有多大.
我們就能夠取用神的應許有多大.
信心是要用行動體現出來.
那些才是真信心.
又要縱橫走遍這地.
你有沒有想過東南西北.
東西和南北加起來就構成一個十字架.
一中一橫一豎一橫就是十字架.
剛剛過去的幾天.
我和賓豐寶珠幾個參加福音更新論壇.
我也徵求了講員的同意.
借用她這幅圖畫.
其實我們的眼光很重要.
我們用什麼眼光去看事物呢?.
如果我們覺得這個信仰.
正是我,我是罪人.
耶穌基督在十字架上拯救了我.
我將來得永生,就是這樣.
其實很不完全.
福音不僅是我和上帝之間的事.
福音於我自己,於我的社區.
於我的城市,我的國家.
這個世界都是很重要.
我再一次被提醒.
我們大家都很熟悉約翰福音三章十六節.
「神愛世人,甚至把他的獨生子賜給他們.
叫一切順他的,不致滅亡,反得永生」.

$^{321}$神愛世人,但在英文裡.
「For God so loved the world」.
The world是世界.
整個世界,每一個人.
每一個生物.
每一個受造界.
受造之物,都是神所愛.
今天,我們在這裡.
我們不是要建立一個自己.
一個偉大的教會.
最重要,我們要這個群體.
可以影響我們的城市.
影響世界,帶來改變.
所以今天,神要鼓勵我們.
用神的眼光,帶著耶穌基督的福音.
憑信心走出去.
以致我們所到之處.
都可以成為別人的祝福.
為這個世界帶來影響.
並且將人帶到基督的面前.
讓我們一起同心低頭祈禱.
親愛的天父上帝.
是你先愛我們.
你在十字架上寶血洗淨我們的罪.
以致我們今天能夠成為你的兒女.
來得著永生.
但我們知道.
這並不是這一刻我們可以停下來.
而是神在這個受造界.
每一樣都是你所做.
你要我們都成為這個世界的祝福.
所以求你讓我們有從你而來的眼光.
以致我們真的能夠用你的眼光.
去看這個世界的時候.
我們能夠貼近你的心意.
讓我們都能夠懂得.
如何走入這個世界裡.
成為你自己美好的見證.
幫助我們帶著你的眼光去看這個世界.
以致我們能夠成為這個世界的光.

$^{361}$成為這個世界的炎.
叫榮耀歸給你.
我們同心祈禱奉耶穌基督命求.
阿們.
\newpage



\section{出埃及記 2:10}
\label{sec:031884}
\textbf{聖經人聖經事:從水裡拉出來}~~~~黃克勤牧師博士
\newline
\newline
連結: \href{http://www.church.com.hk/acms/content.asp?site=cdc\&op=show\&type=product\&code=031884\&layout=sermon}{\texttt{031884}} ~~~~ 語音日期: 2026-03-25
\newline
\newline
\hyperref[sec:031883]{\small{< < < PREV SERMON < < <}}
~
\hyperref[sec:index_chronic]{\small{[返順時目]}}
~
\hyperref[sec:031885]{\small{> > > NEXT SERMON > > >}}
\newline
\newline
出埃及記 2:10
\newline
\begin{longtable}{cl}
\hline
\hline
章節 & 經文 (和合本修訂版)\\
\hline
2:10 & \begin{tabularx}{0.7\textwidth}{X} 孩子長大了，婦人把他帶到法老的女兒那裡，就作了她的兒子。她給孩子起名叫摩西，說：「因我把他從水裡拉出來。」 \end{tabularx} \\ \\
[1ex]
\hline
\hline
\end{longtable}
$^{1}$《生命說書》的網絡又跟大家一起看《聖經》的人與故事.
認識背後不停地工作.
在這裡是掌管著上帝.
今次講另一個人的故事.
從水裡拉上來.
從水裡拉上來即是誰呢?.
從水裡拉上來.
摩世.
我們來看看.
《昌合記》第二章第一節.
有一個利未家的人娶了一個利未女子為妻.
即是利未人生的利未人.
如果她生兒子的話.
那女人懷孕生一個兒子.
見她盡美就藏了她三個月.
後來不能再藏就娶了一個蒲草箱.
抹上石漆和石油.
將孩子放在裡頭.
把箱子擱在河邊的蘆笛中.
看看這三節.
有一個純種的利未家的嬰兒出生.
出生後發現孩子很靚仔.
多數嬰兒出生都靚仔.
其實這裡是在說.
見到她很疼她.
就藏了她三個月.
特別留意的是藏了.
藏起來.
不是擺白日宴.
不是拿出來大家見下滿月酒嗎?.
不是.
藏了三個月.
主要原因就是當時.
在《昌合記》第一章一開始.
《創世記》最尾是約瑟做宰相.
到《昌合記》第一章就不是了.
是有不認識約瑟的新王起來.
就是說舊的政權被推翻了.
新的埃及有另一族人上來做埃及的皇帝.
他們是不認識約瑟的.

$^{41}$沒有面子比的.
主要埃及地方都很大.
尼羅河都很長.
在上下尼羅河是分了兩個大部分的埃及人.
他們之間的關係不是很好.
可能以前的政權被打敗了.
反轉.
新的埃及王法老上來.
他就想了.
第一章裡面說.
哎呀.
這麼多以色列人.
他們就叫做希伯來族.
哇.
幾百萬.
現在雖然是很免費在做勞工勞帝.
但是萬一.
如果有外面的國家.
跟他們有什麼密謀要造反呢?.
裡應外合.
哇.
我國家裡面有400萬人反我.
外面又有外國的士兵.
那我就一定沒有了.
那怎麼辦?.
他們又強壯.
因為個個都大隻做事.
做勞帝.
建磚,建城牆,建城.
哇.
那不是了.
不如我們滅他族了.
怎麼滅呢?.
就是叫收生婆.
如果埃及裡面的所有希伯來族的人有嬰兒.
一生出來你去接生.
如果是男丁.
你幫我殺死他.
女的.
任由他存活.

$^{81}$可以斟茶遞水做事.
也可以做粗工.
男的殺死他.
但是那兩個收生婆.
聖經是說.
初埃及的第一章說他是敬畏神的.
所以他不做.
那被埃及王就說.
喂.
法老說為什麼你留下那些男生.
我不是叫你殺了他嗎?.
哎呀.
主上你不知道了.
那些希伯來女人強壯.
她聽到說生了叫我們去.
我們跑了去.
生了出來了.
已經包著抱著了.
那些大人都在了.
所以我沒辦法下手.
這樣?.
唉.
這樣了.
逃過一死.
你看到.
其中一個就是這樣了.
於是下令要滅他們族的.
所以.
這個男子兒子就藏起來了三個月.
後來不能夠再藏了.
開始長大了.
他的父母就將蒲草.
那些草是很高很長的.
就割開它.
剪開它.
研碎.
就用來做紙.
所以埃及的紙就是蒲草紙.
如果將它一支支連在一起疊起來.
就可以做成一個箱.

$^{121}$就用繩子綁住.
加一些好像七和石油.
這個石油.
就不是我們現在說的中東石油那些石油.
是屬於好像鈉青之類的東西.
是可以游下去封住.
令到蒲草箱不入水.
就將小孩子放了進去.
然後將箱子.
就放在河邊的奴隸當中.
敗了.
其實接著怎樣?.
就是說.
唉.
兒子.
就看你命水了.
衝了去哪裡.
有沒有撿回一條命.
誰會撿了你.
不知道.
因為是要殺死男孩子的.
孩子的姐姐遠遠站著.
這個媽媽也很聰明的.
就派了姐姐.
你跟著遠遠地.
在奴隸裡面.
看著看著.
看看弟弟怎樣.
要知道他究竟怎樣.
剛剛好.
第一點.
我想給大家看的就是.
「箱擺在河邊的奴隸中」.
你一放進去.
還能不能自己管?.
還能不能自己控制?.
沒有的了.
可能他媽媽也禱告.
主耶和華.
我們就將這個兒子放在箱子.

$^{161}$放在河當中.
求你自己保護他.
唉呀.
再見了.
這樣就.
放進去了.
放進河裡了.
就是說已經不到你掌控了.
不過.
起碼也想看看發生什麼事.
誰知道第二件事發生.
紅色字.
「法老的女兒來到河邊洗澡」.
先洗個澡.
「他的侍女們也在河邊走來走去」.
那麼他.
咦?.
誰看到呢?.
是法老的女兒.
竟然看到在河邊飄來飄去.
咦?有一塊東西在河裡.
不是鱷魚.
不是,是一個箱子.
什麼東西來呢?這麼有趣.
就叫侍女.
誰誰誰.
「幫我將奴隸裡的箱子拿過來」.
於是他們就馬上硬著頭皮拿箱子過來.
第二.
把箱子放進去.
又要剛好有法老的女兒來到河邊洗澡.
還要看到那個箱子.
也是上帝的手在那裡工作.
「沒有啊,米利暗」.
米利暗,我後來才知道這是摩西.
所以說米利暗說多了.
其實這個侍女知不知道.
法老的女兒來河邊洗澡的呢?.
洗澡的呢?.
不知道.

$^{201}$她什麼時候來洗澡.
你猜不到的.
她可能來這裡,可能去那裡.
未必知道的.
可能她是去那裡.
「哎呀,熱得下去,洗一下身先」.
可能是這樣.
一直摸著摸著.
誰知道法老的女兒來到河邊濕濕的身.
就看到那個箱子.
所以第二點.
又是上帝的手在那裡工作.
打開箱子就看到那個小孩.
可能那個小孩看到她.
又哭起來了.
「哎呀,慘了」.
這個法老的女兒.
即是女性.
看到那個寶寶哭.
「哎呀,寶寶為什麼哭啊」.
「不用怕,不用怕」.
接著她看她的服飾.
各樣東西一看.
「這個是希伯來人的兒子來的」.
「是啊是啊,是啊是啊主人」.
可能那些侍女就會說.
那些宮女就會說.
「是啊是啊,現在主上是要殺光所有希伯來的男孩子」.
「啊?這樣?」.
「哎呀,那麼這個B都要死的了」.
「是啊,如果傳開了,這個B就要死的了」.
「啊?那怎麼行啊」.
「哎呀,哭啊,慘了」.
那麼可能就可憐她了.
那麼這個時候.
她的姐姐很聰明.
就對法老女兒說.
「你知道你們一走出來」.
「可能已經很多人說」.
「喂,什麼事?什麼事?」.

$^{241}$「這個希伯來的女孩子走過來幹什麼?」.
「但是實際讓她走過去」.
「起碼也有幾個衛兵」.
「公主洗澡」.
「或者你也要走遠一點」.
「郡主洗澡」.
「那麼你也不能夠走過去看著的」.
「遠遠地有士兵」.
「近距離」.
「有侍女,宮女」.
「但是讓她走過去」.
「就跟法老的女兒說」.
「哎呀,如果當日這樣」.
「法老女兒就公主了」.
「公主,公主」.
「如果你覺得這個小孩可憐」.
「和我去找一個希伯來婦人的奶媽」.
「來為你奶這個孩子」.
「為你奶這孩子」.
「可不可以呢?」.
那麼很謙虛的問得.
「很合理喎」.
「如果公主喜歡這個小孩」.
「又可以我找一個希伯來奶媽」.
「來養她的」.
「不知道這樣可不可以呢?」.
「好不好呢?」.
「合不合你心意呢?」.
「法老女兒就說可以」.
「這些大人物有權就說可以」.
「就這樣說出來」.
「這樣,這樣,這樣,這樣」.
「她說沒有,可以」.
「童女就去叫了孩子的母親來」.
「請公主等一等」.
「我知道附近有一個這樣的奶媽」.
「你稍等,很快回來」.
「馬上來了」.
「法老的女兒」.
「即是說為什麼這麼快叫到?」.

$^{281}$「即是箱子就在那裡」.
「在那些牆牆的爐位後面」.
「就一路浸在水裡面」.
「姐姐就看著撥開看看撥開看看」.
「然後媽媽在哪呢?」.
「媽媽可能在不遠的岸上」.
「媽媽看著女兒,女兒看著弟弟」.
「我幫你找一個希伯來的奶媽」.
「很快回來的,我認識一個在附近」.
「真的?快點,我走了」.
「好,好,馬上馬上」.
「馬上媽媽入手,馬上跑過來」.
「好了,對她說」.
「你把孩子抱去,為我奶她」.
「我必給你公家」.
「於是婦人,即是她媽媽就抱著孩子去奶她」.
「怎樣拿公家?」.
「怎樣去哪裡拿?」.
「怎樣進到皇宮問哪位?」.
「那些你不知道的」.
「不過這個故事」.
「就只是將最重要的情節講了」.
「這個小孩要不要給他身份證?」.
「如果不是人家要殺這個嬰兒的時候」.
「你都要證明,別動,這個公主的養子來的」.
「可能這樣,我給他一塊肉」.
「或者我給他一塊牌,試試看一亮出來」.
「有沒有見過這東西?」.
「哎呀,皇上,皇上」.
「可能會給他這些東西」.
「所以是保護他的性命的」.
「這裡就沒有特別講了」.
「第三就是,這個法律的女兒竟然講可以」.
「為何這麼出奇?」.
「當然出奇了」.
「法律的女兒可以講什麼?」.
「可以多加一個字」.
「不可以,不可以」.
「為甚麼?」.
「因為我父王是要殺死所有希伯來的男婿男嬰的」.

$^{321}$「所以不可以,不用了」.
「我帶他回去交給我爸爸」.
「交給處事的官就行了」.
「那你沒有了」.
「不要了,不要了」.
「還有不好嗎?」.
「如果他說不可以」.
「這也是上帝在法律的女兒心裡做的事」.
「令她動了愛這個孩子的心」.
「所以他說可以」.
「我還要給你工價幫我養大她」.
「就是這樣」.
「哪個是工作來的?」.
「他說可以」.
「是上帝的工作來的」.
「孩子佔長」.
「婦人就將他帶到法律的女兒那裡了」.
「第四,他就作了他的兒子」.
「真厲害」.
「也是上帝的手在後面工作」.
「見到這個孩子」.
「這麼大了?」.
「多帥啊」.
「聰不聰明伶俐啊?」.
「可能已經在家教他了」.
「下次見到公主的時候」.
「就要說公主平安」.
「如果清代就要說家家吉祥」.
「你要懂得說這些東西」.
「但方方站直一點」.
「像我這樣腰都彎了」.
「站直一點」.
「見到就要跪」.
「不要看」.
「讓她叫你平身」.
「你才可以平身」.
「這麼聰明」.
「真厲害」.
「好吧」.
「這個公主應該還沒結婚」.

$^{361}$「不過有個」.
「就認她做兒子」.
「要這麼大人物」.
「法老的女兒」.
「才能保住這個孩子的命」.
「見到上帝的手在工作嗎?」.
「第四點」.
「就作了她的兒子」.
「這是上帝的手」.
「大大地工作」.
「既然如此」.
「我給你改一個名字」.
「改這個名字」.
「公主改一個名字」.
「一定是埃及名」.
「名字就叫做」.
「好吧」.
「我就叫你做Moshe」.
「Moshe」.
「英文叫Moses」.
「其實翻譯得不好」.
「Moshe」.
「重音在後面」.
「Moshe」.
「可以翻譯成什麼呢?」.
「Mosab」.
「這些名字這麼難聽」.
「Moshe」.
「Moshe」的意思其實很簡單.
「被拉的」.
「不是被警察拉的拉」.
「是被拉上來」.
「拉上來的意思」.
「是我拉他上來的」.
「其實他真的被拉上來」.
「他只不過是一個奴隸的兒子」.
「而且男子是要死的」.
「不過今天他拉他上來」.
「做我兒子」.
「其實應該是這樣」.

$^{401}$「然後就說」.
「如果我們在屬靈就是」.
「這個必死的」.
「希伯來男子」.
「竟然做了」.
「埃及法老的養子」.
「那就是拉上來了」.
「是不是上帝拉他上來」.
「屬靈上這麼看」.
「這個公主就說」.
「好吧」.
「我就拒絕我的兒子」.
「我要跟他改個名字」.
「就叫Moshe」.
「Moshe」.
「意思說」.
「因我把他從水裡拉上來」.
「拉出來」.
「中文聖經和本是意思說」.
「所以會令人覺得」.
「摩西的名字這個意思」.
「Moshe就是因我把他從水裡拉上來」.
「摩西這個名Moshe是什麼意思」.
「是因我把他從水裡拉上來」.
「其實不是」.
「你看到我用藍色寫意思兩個字」.
「括號」.
「如果我用藍色在這些powerpoint」.
「紅色黑色的字」.
「是和合本本身的字來的」.
「藍色都是和合本本身的字來的」.
「不過是原文沒有的」.
「就用藍色括著它」.
「他做我的兒子」.
「我跟他改個名Moshe」.
「跟著就說」.
「改名的時候就說」.
「他叫Moshe」.
「因我把他從水裡拉上來」.
「不是意思」.

$^{441}$「意思兩個字是沒有的」.
「就說我把他從水裡拉上來」.
「其實Moshe這個意思是被拉上來」.
「其實他是不是公主拉上來」.
「是」.
「但最重要的是神把Moshe拉上來」.
「這個人就是摩西了」.
「所有今天跟大家分享的四點」.
「你都會看到」.
「上帝的手在摩西的身上工作」.
「他本來只是一個隨處飄的」.
「一個蒲草箱在河上」.
「但竟然撞到法老女兒」.
「來濕濕個身,沖沖涼」.
「他不一定來這裡的」.
「不是每天都來的」.
「可能是走到哪裡玩的時候」.
「去哪裡的時候」.
「熱了下去河裡沖沖涼」.
「就剛剛遇到」.
「第三件事就是」.
「他又喜歡這個小朋友」.
「跟著姐姐就說」.
「可不可以找一個」.
「希伯來的奶媽給他?」.
「可以」.
「為什麼不說不可以?」.
「不可以,不行」.
「第一,我無端端」.
「我這麼尊貴又未結婚」.
「弄個養子自己不要」.
「而且這樣做又犯了爸爸的指引」.
「不好,不可以」.
「算了,我們走了」.
「可以這樣的」.
「他說可以」.
「跟著就做了他的養子」.
「做我的乾兒子」.
「做我的兒子」.
「嘩,這麼厲害」.

$^{481}$「於是就改名叫Moshe」.
「就拉上來的」.
「其實這個拉上來」.
「是誰拉的?」.
「除了是法老女兒」.
「神透過法老女兒的手」.
「拉了他離開水之外」.
「神是透過整件事」.
「將他在一個奴隸的兒子的身份」.
「拉了出來」.
「有些事情要給他做」.
「這就是今天」.
「我們學到屬靈的功課了」.
「我們一起禱告吧」.
「讚美天父上帝」.
「當你要用人的時候」.
「你就是會由很小的時候」.
「開始模造他」.
「一直讓他適合為他開路」.
「無論有高有低」.
「高的是令他走一條順的路」.
「低的是磨練他」.
「以致他更加適合那個工作」.
「你給他的呼召與差遣」.
「今天我們一開始見到摩西」.
「她是被你」.
「在低處奴隸那裡拉出來」.
「奴隸將要必死的男子那裡拉上來」.
「你是透過了公主」.
「剛剛去到那裡又接受她」.
「於是她就成為公主的兒子」.
「學了埃及的學問」.
「主要是你親手的工作」.
「我們讚美你」.
「知道你的手一直在運作」.
「你創造的大能」.
「你維持的大能」.
「你掌管的大能」.
「一直沒有斷絕」.
「你是繼續工作的」.

$^{521}$「願主得著榮耀」.
「祈禱奉耶穌名求」.
「阿們」.
\newpage



\section{約書亞記 18:1-10}
\label{sec:031885}
\textbf{你們耽延不去得,要到幾時呢?}~~~~劉銳鋒牧師
\newline
\newline
連結: \href{http://www.church.com.hk/acms/content.asp?site=cdc\&op=show\&type=product\&code=031885\&layout=sermon}{\texttt{031885}} ~~~~ 語音日期: 2026-03-29
\newline
\newline
\hyperref[sec:031884]{\small{< < < PREV SERMON < < <}}
~
\hyperref[sec:index_chronic]{\small{[返順時目]}}
~
\hyperref[sec:031886]{\small{> > > NEXT SERMON > > >}}
\newline
\newline
約書亞記 18:1-10
\newline
\begin{longtable}{cl}
\hline
\hline
章節 & 經文 (和合本修訂版)\\
\hline
18:1 & \begin{tabularx}{0.7\textwidth}{X} 以色列全會眾都聚集在示羅，把會幕設立在那裡。那地已經被他們征服了。 \end{tabularx} \\ \\ \relax
18:2 & \begin{tabularx}{0.7\textwidth}{X} 以色列人中剩下七個支派還沒有分得他們的地業。 \end{tabularx} \\ \\ \relax
18:3 & \begin{tabularx}{0.7\textwidth}{X} 約書亞對以色列人說：「耶和華－你們列祖的神所賜給你們的地，你們耽延不去得，要到幾時呢？ \end{tabularx} \\ \\ \relax
18:4 & \begin{tabularx}{0.7\textwidth}{X} 你們每支派要選三個人，我好派他們去，他們要起身走遍那地，按照各支派應得的地業寫明，然後回到我這裡來。 \end{tabularx} \\ \\ \relax
18:5 & \begin{tabularx}{0.7\textwidth}{X} 他們要把地分成七份。猶大在南方，住在他的境內。約瑟家在北方，住在他們的境內。 \end{tabularx} \\ \\ \relax
18:6 & \begin{tabularx}{0.7\textwidth}{X} 你們把地劃成七份之後，就要把所寫的帶到我這裡來。我要在耶和華－我們的神面前，為你們抽籤。 \end{tabularx} \\ \\ \relax
18:7 & \begin{tabularx}{0.7\textwidth}{X} 利未人在你們中間沒有分得地業，因為耶和華祭司的職分就是他們的產業。迦得支派、呂便支派和瑪拿西半支派已經在約旦河東得了地業，是耶和華的僕人摩西給他們的。」 \end{tabularx} \\ \\ \relax
18:8 & \begin{tabularx}{0.7\textwidth}{X} 那些去劃地的人起來正要去的時候，約書亞吩咐他們說：「你們去走遍那地，把地劃分以後，就回到我這裡來。我要在示羅這裡，在耶和華面前為你們抽籤。」 \end{tabularx} \\ \\ \relax
18:9 & \begin{tabularx}{0.7\textwidth}{X} 那些人就去了，走遍那地，按照城鎮把地劃成七份，寫在冊上，回到示羅營中約書亞那裡。 \end{tabularx} \\ \\ \relax
18:10 & \begin{tabularx}{0.7\textwidth}{X} 約書亞就在示羅，在耶和華面前為他們抽籤。約書亞按照以色列人的支派，在那裡把地分給他們。 \end{tabularx} \\ \\
[1ex]
\hline
\hline
\end{longtable}
$^{1}$弟兄姊妹平安.
今日要同大家分享的經文.
記載在《約書亞記》十八章一到十節.
分享的題目是.
「你們丹賢不去得,要到幾時?」.
這個題目是直接引用.
第十八章第三節約書亞所說的話.
其實從第十三章開始到第十九章.
很長篇幅的經文.
都是在說同一件事.
就是以色列民分地.
我們知道前面他們打了很多的仗.
他們很多的爭戰.
到了十三章開始.
其實基本上他們已經獲取了勝利.
那些地已經被征服.
接下來就要分地給各國的支派.
沒辦法在這麼短的時間.
從十三章說到十九章.
說完這麼多的內容.
所以我就截取了當中.
十八章一到十節這一段.
跟大家重點分享.
看看這段經文.
神要對我們說些什麼話.
在分地這件事上.
對我們有什麼提醒和啟示.
我不知道大家在讀聖經的時候.
特別讀到分地這一段的時候.
會不會很快就跳過不讀了.
因為好像都頗長和很枯燥的經文.
很容易就跳過不讀.
但其實當中有很多很深刻的屬靈含義.
我讀的時候就有一個疑問.
為什麼這個地已經是神允許了給他們.
還要他們這麼辛苦去爭戰.
要去打仗.
要去得那個地呢?.
為什麼到現在不是已經征服了嗎?.
不是屬於他們了嗎?.

$^{41}$為什麼還要去得呢?.
我不知道你讀的時候.
會不會有跟我同樣的疑問.
其實在聖經或者在神學角度來看.
神雖然已經賜了給以色列有應許之地.
但這個應許並不是不勞而獲的一個禮物.
而是一場需要與神同工.
以及憑信心奪取的過程.
神賜應許之地給他們.
是一個律法和信仰上的一個約定.
而爭戰就是實際執行和信心的考驗.
這片土地雖然已經賜給了以色列人.
但當時還是被迦南人所佔據的.
神就要以色列人透過爭戰.
來驅逐這些罪惡.
去彰顯祂的公義.
並且在這個過程當中.
讓以色列人能夠去經歷.
怎樣去信靠神和去信服.
所以有些學者就說.
為什麼賜給他們還要他們去得呢?.
其實有幾個方面的原因.
他們歸納至少有以下四個方面的原因.
第一就是信心的考驗和鍛鍊.
進入應許之地不單止只是在地理上的移動.
好像從東移到西.
從西移到南.
不是那個移動.
更多是屬靈上的成熟.
神是要以色列人透過爭戰.
去學習怎樣去依靠神.
而非去依靠他們自己的武力.
在聖經當中提到.
當以色列人因為恐懼而退縮的時候.
他們就沒有辦法進入神的應許.
第二方面就是對迦南惡行的審判.
當時迦南地的居民.
他們真的是罪大惡極.
他們行很多的惡.
包括淫亂.

$^{81}$他們用火將自己的兒女獻祭給假神.
神就藉著以色列人的手.
去對這些邪惡的勢力進行審判.
所以神要他們去爭戰.
爭戰被視為神公義的執行.
另外就是要清除偶像崇拜的誘惑.
神命令以色列人趕出當地的居民.
其實是為了防止他們受到.
這些異教偶像崇拜的誘惑.
歷史已經表明了.
當以色列人沒有徹底去執行神的命令的時候.
他們很快就會轉向這些異教的假神去敬拜他們.
最後一方面就是.
他們要學習與神同工.
躺在那裡是不能夠得到應許的.
神定義要人參與在他的計劃裡面.
這種就是神作主導人去執行.
這個是模式來的.
約書亞他帶領以色列人去爭戰.
其實都是這種模式.
是為了建立人與神之間很親密的立約的關係.
所以總括來說就是.
應許是所有權的賦予.
而爭戰是使用權和治理權的實際取得.
兩樣是有點不同的.
在我們很熟悉的人物.
譬如說大衛其實都是一樣的.
他剛剛很小的時候就被薩姆爾高納為王.
是不是馬上就做了王呢?.
我們知道也不是.
要經歷了很長的時間.
他也要經歷很多的爭戰.
要帶人去打仗.
經歷很多的鍛鍊之後.
最後他才真正登基作王.
所以不是一高納了就真的作了王.
當然高納他的時候.
其實在法律的意義上.
在信仰的層面上.
其實已經是了.

$^{121}$神已經給了他.
不過他還要經歷一些時間去獲取.
這個其實是反映了聖經裡面一個很重要的概念.
就是神的恩典或者神的應許.
往往是伴隨著人的責任和行動的.
這個情況有點類似.
那樣東西是已經屬於你的.
但是你又未得到.
我不知道你有沒有這樣的經驗.
有時候你會收到一些禮券.
我也試過收過.
特別生日的時候.
有些銀行寄來一些禮券.
不是很多,五十元那些.
不過就要你去到.
通常去到沒有雷公那麼遠.
要去兌現的.
還要親自去.
那結果怎樣?.
當然不去那麼遠又浪費時間.
為了那五十元車費時間也虧了.
所以最後還是沒有兌現到.
這五十元禮券是給了我.
是屬於我的.
但是又沒有兌現到.
我又沒有拿到.
這個情況就好像這樣.
又或者有人開了一張很巨額的支票給你.
但是你又放在枕頭底下.
一直都沒有拿去兌現.
最後還是沒有拿到錢.
是屬於你的.
但是又沒有真的拿到.
其實現在以色列民是這樣的情況.
有一些支派已經獲取了土地.
但是還有七個支派還沒有得到土地.
他們雖然征服了迦南.
但是還沒有真正在那個土地上經營和生活.
特別是還有七個支派還沒有得到土地.
所以這段經文是說.

$^{161}$約書亞催促他們不要膽言.
你們要快些去得到那個土地.
要快些在那個土地上經營和生活.
我們看看第十八章的第一到二節.
我們一起讀一讀.
以色列的全會眾都聚集在士羅.
把會幕設立在那裡.
那地已經被他們制服了.
以色列人其餘的七個支派.
還沒有分給他們地業.
其實這裡已經是說到他們的狀態.
雖然是有五個支派已經分了地業.
但是還有七個還沒有分的.
這七個其實都在停滯中.
沒有繼續向前.
他們好像安於現在的現狀.
所以約書亞把他們聚集在一起.
聚集在哪裡呢?.
就在士羅這個地方.
士羅的意思是安息.
這個城其實就在以法連支派所分到的地上.
如果我們看地圖就會比較清楚.
中間黃色那塊地就是以法連他們已經分到的地.
紅圈那裡就是士羅那個城.
約書亞就把他們聚集在那裡.
他也選擇了士羅來設立會幕.
在士羅這個城裡面設立了會幕.
供全民去敬拜神.
當五個支派分得地業之後.
約書亞就把會幕從吉甲搬到士羅這個地方.
其實這個就意味著.
他們的宗教,政治和軍事的中心.
都移到迦南的中心地帶.
會幕其實是以色列百姓很重要的宗教生活的中心.
也象徵著以色列的終極統治者就是神.
因為會幕就是象徵神的同在.
會幕在那裡就是代表神在那裡.
神與他們同在.
所以士羅這個地方.
數來都是一個宗教中心.

$^{201}$如果你看舊約.
特別是在史詩的時代.
基本上那個會幕都設立在士羅的.
但是直到以利作事師的年代.
約櫃就被非理士人擄走了.
所以士羅的地位就開始式微.
在經文裡面也講到.
他們在那裡設立會幕.
然後就說那地已經被他們制伏了.
那地的意思其實是指迦南地的核心地區.
其實就是若澗河西那裡.
若澗河東已經有兩個支派分了地業.
但是若澗河西大部分都還沒有分.
只是分了另外兩個半支派.
還有七個還沒有分到地.
所以他們聚集在那裡.
聚集在那地.
雖然尚未完全消滅所有的迦南人.
但是若澗河西的戰略高地.
其實已經在以色列人的控制之下.
這個其實是意味著整體的勝利.
士羅因此被定為一個中心.
以便後續可以為其餘的七個支派.
去分配剩餘的土地.
當時猶大爾法蓮和嗎哪西已經分了土地.
會幕設立在士羅就是為了.
規劃其餘的七個支派他們的地業.
我們一起來看一看.
那些支派他們分地的情況.
已經得到地業的支派包括裡面.
即是樓面.
迦德還有嗎哪西半支派.
他們全部在若澗河東已經一早分了地.
另外就是猶大爾法蓮和嗎哪西.
另外半支派就在若澗河西分了地.
未分到地的就是.
卞雅敏,西布倫,以撒加,阿哲,南弗塔利和丹.
這七個支派在若澗河西是還未分到地的.
另外還有一個比較特別的支派就是利美.
他們是沒有份去分地的.

$^{241}$因為他們是作為祭司的職責.
耶和華神就是他們的產業.
所以他們就不能分.
他們會散居在其他的支派當中.
這個就是當時的情況.
他們有些已經分了地.
有些還未分.
未分的就停滯在那裡.
沒有繼續向前在地上經營和生活.
所以若舒雅就催促他們不要擔言.
快些去得地.
所以接著若舒雅就去責備他們.
在第三節我們一起讀一讀.
若舒雅對以色列人說.
「耶和華你們列祖的神所賜給你們的地.
你們擔言不去得.
要到幾時呢?」.
「你們擔言不去得要到幾時呢?」.
就是今天的講題.
是若舒雅對他們的責備和催促.
「擔言不去得」其實是指.
以色列人已經征服了大部分的艱難地.
也都設立了會幕.
但是因為他們可能是懶惰.
又或者是恐懼.
又或者是他們安於現狀.
他們沒有積極地將剩餘的土地從原住民的手中驅逐出來.
沒有去佔領.
也都沒有實際分配和居住.
所以這裡不是指他們未征服那個地.
而是指他們未實際去佔領和去經營或者去生活.
這個時候雖然他們已經設立了會幕.
在時羅大局已定.
主要的戰爭都已經獲勝了.
但是他們仍然有很多的土地.
是未去取得的.
所以若舒雅就說他們「擔言不去得」.
「擔言不去得」的意思是什麼呢?.
就是說他們行動上面其實是怠慢的.
或者他們那七個支派已經暫時安頓了這個地方居住下來.

$^{281}$所以對於再去征戰.
再去分配剩餘的土地.
他們會覺得比較遲疑.
又或者他們有軟弱不去做.
所以若舒雅跟他們說「不要擔言」.
「要快點去得地」.
叫他們快點去得地的話.
不是叫他們去攻打.
不是叫他們去佔領那麼簡單.
而是實際上要他們在那個地方上面去經營,去生活.
真的讓那個地方成為他們生活的地方.
所以若舒雅責備他們.
其實也是鼓勵他們要拿出信心.
要履行神吩咐給他們的使命.
不要在已經取得的局部勝利當中.
他們就停滯不前.
而是要完成神所給他們的全部產業的承接.
所以總括來說.
當時的土地已經被他們征服了.
但是還有一些地未被實際佔領.
為什麼他們沒有去做呢?.
有些學者就歸納了有三大方面的原因.
第一方面就是.
他們對神的應許的誤解.
他們未能夠理解神賜土地給他們的終極目的.
他們誤以為應許僅僅是物質上的安逸.
現在有地方住了.
有些財物可以用了就足夠了.
而他們不知道原來他們有屬靈的使命需要去完成.
第二方面就是在征戰之後.
他們有一些屬靈的疲乏.
因為長期征戰.
所以他們的身體都很疲憊.
滿足於已經得到的一些財物.
失去繼續戰爭的一些意願.
第三方面可能因為他們是遊牧民族.
所以習慣了遊牧的慣性.
所以拒絕農耕的生活.
長期的遊牧生活令他們不適應定居和耕種.
他們都是傾向於像以前那樣.

$^{321}$跟迦南人做生意,做貿易.
而獲取一些物品.
而不是在當地自給自足去生活.
這些是他們不能去地的一個原因.
有時我們都會像他們那樣停滯.
當我們在困境裡面的時候.
我們就會很迫切地去祈禱.
求神幫助我們去擺脫這個困境.
但是當這個困境稍稍有轉機.
稍稍有不同的時候.
環境稍稍安穩的時候.
我們很可能就會安於現狀.
我們就會止步不前.
所以神就藉著環境跟他們說話.
特別是在示羅設立這個會幕.
去催促這個信徒.
不可以停留在次好上面.
應該要追求基督裡面的上好的應許.
有時我們經常都會說.
我們在這裡真好.
所以我們就會停滯不向前.
其實可能前面更加好.
所以我們要追求在基督裡面上好的應許.
所以會幕的設立.
不單止是一個敬拜的中心.
亦都是神對那些擔言者提醒.
神的同在就是催促他們的記號.
表示他們不可以再停下來.
要向前走.
接下來我們看到他們的心態.
他們擔言他們不去行.
其實是反映他們一種很深層的心態危機.
首先他們是安於現狀.
他們因為得到部分的恩典.
而遲遲不去爭取.
其次就是他們畏懼困難.
他們不願意親手去清理那些殘餘的迦南人.
寧願去逃避這個雖然困難.
但是必要的一個清理的工作.
最終他們就是缺乏使命.

$^{361}$他們忘記了神賜這個土地給他們的目的.
不是為了他們享樂.
或者他們可以過一個安逸的生活.
而是要在他們當中建立一個聖潔的國度.
要承擔神所託付的神聖使命.
所以這三種心態的交織.
令到這班本來是屬神的子民.
是蒙召得地的子民.
淪落成一個好像屬靈的旁觀者.
所以約書亞要責備他們.
和馬上要命令他們要去.
所以我們看接下來.
在第四到第七節就看到.
約書亞命令他們要怎樣去做呢?.
這裡經文有點長.
所以就不邀請大家一起讀了.
從第四節到第七節.
其實都是約書亞的命令.
在這裡約書亞他責備了以色列民之後.
他就立即頒佈了土地分配的具體的命令.
就是要每一個支派派出三個人.
去走遍尚未佔領的土地.
並且要把土地分成七份.
詳細去繪製這個地界.
然後回來回報給他.
接著約書亞就在耶和華神面前.
為各支派去鑾溝.
公平去分配這個產業.
這個經文裡面都明確指出.
猶大支派仍然在南方.
約瑟加仍然在北方.
他們已經佔據了那個地.
已經分了.
但是利美支派.
但是利美支派他們沒有得到地.
因為他們是祭司承擔這個職份.
他們就不需要獲得這個產業.
因為耶和華神就是他們的產業.
但是迦德和裡面嗎哪西半支派.
早在若間河東就得了地.

$^{401}$只剩下那七個支派未有.
所以這一次約書亞對他們的責備.
其實是對消極態度的責備.
也同時是對神聖的秩序和分工的重申.
大家有沒有留意到.
在讀這個分地的經文裡面.
其實是有些原則的.
其實在舊約裡面.
民數記也有講到這個分地的原則.
如果我們總結出來.
就是兩大原則.
第一就是按人數的多寡.
人數多一點的支派就會分到比較大的地.
人數少一點的就會小一點.
另外就是會「粘溝」即是抽籤.
就會確定那塊地的位置.
究竟在南,在北,在東,在西還是在中間.
這個是確定位置的.
一個就是大小.
另一個就是位置.
但是當你很認真去想的時候.
這兩個的原則會不會有些矛盾的地方.
或者會不會有些衝突的地方.
我剛剛開始讀的時候.
我也覺得好像有些衝突和矛盾.
因為這個經文講到.
他派一些人去勘察那個地.
跟著去劃那個地界,將那些地分成七份.
分了之後就開始去抽籤,即是去「粘溝」.
既然已經劃分好了,分好了.
有大有小,按照人數多寡.
人數多的會大一點,人數少的小一點.
人數多寡,所以劃分那七份就有大有小的.
如果是這樣,再去「粘溝」的話.
會不會有些比較人數少的支派抽到比較大的地.
或者比較人數多的支派抽到比較小的地呢?.
我當時覺得會不會有這個可能呢?.
但是再認真去想的時候,其實是不會的.
因為他們的分地不是一次完成的.
他們分地的行動其實是分了兩次去完成的.

$^{441}$是兩個階段來的.
第一個階段我們看到,有些大的支派其實是先行已經佔了地.
譬如說猶大和約瑟的子孫.
馬來西亞,以法連那些支派已經是佔了地.
正如那七個支派是未分到的.
所以在分七份之前.
其實有些大的支派是基於他們強大的實力.
和對應的人數已經先分好了.
之後到第二階段就分給剩餘的七個支派.
那麼這七個支派怎麼分呢?.
其實那七份不是完全均等的.
我們知道,剛才都說到是按照人數.
所以他們是有大有小的.
這些土地都是有肥沃一點的土地.
有些可能是貧石一點的.
有些可能是平原.
有些可能是山地來的,不是那麼好的.
但是山地又不是完全不好的.
因為山地有很豐富的森林資源.
所以有不同的情況.
那麼這個就是一個不是完全的等分.
這個是一個適配的七份.
甚麼是適配呢?.
即是按照他們具體的情況.
各個支派,他們的人數等等.
那麼這個有些聖經學者他們的分析就是.
在第二階段的時候其實他們要做很多功夫的.
第一就是要派一些的勘察員去做一些資料的搜集.
所以就好像約書亞那樣去派一些人.
每個支派三個人去走遍全地去勘察.
去劃分,劃那個界線.
劃了之後又怎樣做呢?.
他們劃的時候其實不是盲目地去劃的.
這班人是非常清楚剩下那七個支派.
他們的人口的規模的.
所以在劃分土地的時候.
他們就會預先考量到.
這七份土地當中是有不同的情況的.
大小,肥沃,不肥沃等等.
這個有少少像我們之前因為要建堂.

$^{481}$所以我們各個部門就會將我們部門的需要告訴教會.
即是我們需要多大的地方.
用來進行甚麼聚會,需要甚麼的設備.
將這些資料全部搜集了先歸納在一起.
交給一些專業的顧問的公司.
他們就根據我們的資料去劃草圖出來.
這個是要做很多的資料搜集.
要做很多的功夫的.
不是隨便去劃的.
做完這個之後.
劃了之後就進行到第二步.
就是要黏溝.
黏溝就不是為了隨機分配大小.
而是為了消除爭議.
即是說黏溝是代表神賜給你的這個地方.
不是我指派的.
不是那個約翠婭.
不是領袖去指派這個是你那個是她.
而是透過黏溝.
黏溝就是說是神給你的.
你抽到那個就是神給你的.
你就沒有藉口說我不要這個我要那個.
這個就避免了各支派之間的互相的爭奪.
有些可能會爭奪水源好的地方.
比較肥沃的地方.
就避免了這件事.
所以黏溝其實是決定那個位置在那裡.
最後就是按人數多寡.
這個是一個最終的動態的調整.
雖然黏溝是定了一個位置.
但是那個地界已經是死板不可以移動的.
可以向南移一點.
可以向北移一點.
可以是移動的.
所以如果抽完籤之後.
如果發覺某個支派的人數很多.
但是他抽到的地是不夠他們住的.
又或者一個支派的人數很少.
抽到一個地方很大怎麼辦呢?.
他們就會重新再劃分去調整那個地界.

$^{521}$譬如說猶太就是這樣.
在十九章九節我們見到.
猶太因為他的地分得比較多.
但是西面就沒有了.
所以就從猶太當中割了一些地分給西面.
我們一起讀讀第九節.
「西面人的地業是從猶太人地業中得來的.
因為猶太人的分過多.
所以西面在他們的地業中得了地業」.
這裡就是做了一個調整.
猶太多的分了一些給西面.
所以他們最後是會有這個調整的.
剛才說了這麼多.
可能大家都一頭霧水.
都不知道究竟是怎樣分.
如果我用一個比較簡單的日常生活的例子.
可能大家會比較容易明白.
我不知道大家有沒有分過蛋糕.
應該都有.
特別是小朋友生日的時候.
我們需要分蛋糕.
如果有七個小朋友.
你會怎樣切蛋糕呢?.
是不是平均這樣切七份呢?.
通常都不會吧.
因為小朋友可能有一個兩歲.
一個五歲 一個十歲.
都不同的年紀.
他們吃東西的份量可能都不同.
還有你切的時候.
你會發覺可能切那塊有兩粒士多啤梨.
那塊有一個大的朱古力.
那塊奶油又多一點.
每一塊都不一樣.
你不可能很平均地去切.
所以你會怎樣切呢?.
你可能會首先切三塊大塊的.
接著就切兩塊是中間大小的.
中間尺寸的.
兩塊比較小的.

$^{561}$加起來七塊.
三個比較大的小朋友就去抽籤.
抽那三塊比較大的蛋糕.
兩個中間的就抽中間那兩塊.
小的兩個就抽最小的那兩塊.
最後抽完籤之後.
發覺中間年紀的小朋友.
原來那塊蛋糕不夠他吃.
他原來的胃口很大.
食量很大.
但是大的小朋友雖然分了一個大的蛋糕.
但是他吃的東西不是那麼多.
那怎麼辦呢?.
將他多的切一小部分.
分給不夠吃的那個.
這個就是最後那個動態的調整.
這個就是他們分地的原則.
我們所見到.
其實神是很有智慧的.
也給以色列人很有智慧去處理分地的事情.
這個其實都會.
跟我們教會的分配.
地方分配可能我們都可以去參考的.
我們成長班聚會都需要分配地方.
當然我們人數多的是怎樣?.
當然去分配大一點的地方.
人數少的就很自然可能就小一點的地方.
老人家行動不方便怎樣?.
會不會分配去九樓?.
當然不會.
我們都是分配一些方便他們的地方.
也有分配在我們本部的.
也有分配去陪小聚會的.
我們會考慮很多不同的因素去分.
分完之後你都發覺.
原來最後有些不太合適.
那會怎樣?.
會作一個微調.
這些就是原則.
要考慮很多的因素.

$^{601}$當然我們都有限.
我們的智慧有限.
每一個的分配未必很完美.
但是我們分了之後發覺不完美.
可以怎樣?.
我們可以作出最後的微調.
這個就是一個原則.
當約書亞吩咐他們去勘察那個地.
他們之後就很信服了.
馬上就去行動了.
所以在第八節到第十節.
我們見到他們信服和行動.
當約書亞說了之後.
第九節就說他們.
「他們就去了」.
他們就走遍那地.
將城一分成七份.
然後在賊上回到侍羅營中見約書亞.
約書亞就在他們面前.
在耶和華神面前就為他們黏溝.
就分了這個地葉.
在這裡我們見到.
七個支派的代表.
他們立即就執行了.
約書亞所委派的土地勘察和劃分的任務.
他們沒有耽誤.
這次真的馬上去了.
這個行動其實是彰顯著神的公義.
因為他們是依靠神的引導.
去公平分配這個產業.
這個過程是彰顯著.
分配的神聖性和它的程序性.
不是憑人的意思去分.
而是藉著黏溝去體現神的主權.
同時亦都將整個過程記錄在賊上.
以確保分配的透明性和可追溯性.
最終這個土地就能夠.
分配給剩下的七個支派.
完成了神所託付的.
應許之地的分配工程.

$^{641}$所以在整個分地的事情裡面.
我們見到分配產業有深層的意義在當中.
我想至少有四個方面我們值得去留意.
第一就是有獨特的份量.
這個就是按照人數多寡.
來分給我們的啟示.
分地的時候他考慮了人數.
就好像他做人的時候.
考慮了每個人的氣命大小.
和我們的使命其實是不一樣的.
有些人是可以領受五千的.
有些人領受兩千.
有些人領受一千.
我領受五十.
每個人都不同的.
這個不是為了比較優劣.
而是要讓每一個人.
都在他自己的份上.
能夠發揮出最大的功用.
就好像人多的支派.
他們需要大一點的地.
神就像某些人.
他們較大的責任和挑戰.
因為他知道那個人的生命容量.
可以承擔多少.
所以你未必需要很羨慕那些人.
他們很承擔很多的責任.
很重的責任.
因為他們都很辛苦.
要做很多事情.
不需要羨慕.
只是看你那份就可以了.
第二方面就是神聖的定位.
這個就是黏溝所帶給我們的啟示.
黏溝定的是位置.
在哪個位置.
這個是代表神給你的.
不是憑你的努力得來的.
因為是黏溝.
就代表神參與其中.

$^{681}$你抽到的就是神賜給你的.
不是你個人努力得來的.
其實很多事情都是一樣.
我們的出生背景.
我們的外貌.
我們的天賦.
就好像黏溝一樣.
我們沒辦法選擇的.
不是我們選擇的.
是神預定好給我們的.
但當我們去承認.
這是神黏溝給我們的話.
我們就不再抱怨.
為什麼我不是生在比較.
人多又強大的支派當中.
而是我們可以開始說.
神要在這塊土地上.
怎樣去帶領我呢?.
神要在祂賜給我的土地上.
怎樣去使用我呢?.
我們要接受我們所得的.
是神親自給我們的.
是最適配於我們的.
第三方面就是珍視地界.
雖然有地界.
好像有些限制.
但在界限當中其實是有自由的.
我很喜歡詩篇十六篇第六節.
雖然這裡的字有點小.
但如果你看得清楚.
我都邀請你跟我一起讀出.
用醒良給我的地界.
坐落在皆美之處.
我的產業實在美好.
這節經文說到.
神其實是用誠.
量給我們的地界.
祂賜給我們的.
都是祂親自量好的.
是跟我們度身訂造的.

$^{721}$而量給我們的.
是坐落在皆美之處.
我的產業實在美好.
所以神給我們的是最好的.
你不用看著別人的.
如果以法連之派只是看著猶大之派.
為什麼他們分到這麼大的地.
他們的地又平坦.
他們就忘記了.
原來他們分的地是很肥沃.
也有很豐富的森林資源.
所以我們不要看別人的份.
我們要深耕自己的份.
人生的滿足不在於土地的大小.
而是在於你有沒有在那塊土地上.
趕出仇敵和去開墾森林.
第四方面就是動態的豐盛.
調整和恩典.
就好像猶大分地比西面一樣.
神有時也會透過其他人.
透過他們的豐富.
去補足我們的不足.
又或者讓我們的豐富.
去成為其他人的祝福.
這個其實用在教會的場景是很好的.
我們教會作為一個身體.
基督的身體.
我們也要過肢體的生活.
在哥林多前書十二章二十三到二十四節這樣說.
「身上肢體我們看為不體面的.
越發給他加上體面.
不盡美的越發得著盡美.
我們盡美的肢體自然用不著裝飾.
但神能配搭者身子.
把加倍的體面給那有缺欠的肢體」.
這個就是神將我們放在一起.
我們有豐富一點的.
可以分一些給有缺乏的.
我們彼此成為其他人的祝福.
這個就是動態的豐盛.

$^{761}$最後有少許總結和應用.
我們可以從以下三個層次來總結和應用.
第一就是克服未知的恐懼.
雖然神已經透過鏈鉤定下了位置.
但地圖上的地界不等於你實際擁有的.
就好像那七個支派一樣.
應許給了他們.
但他們沒有實際去取得去擁有.
如果神給了我們有恩賜.
就好像神分了蛋糕給我們一樣.
我們是因為害怕.
我們因為怕失敗.
我們不敢去跨出我們的舒適圈.
我們還停留在原地停滯不前,原地踏步的話.
那塊蛋糕或那塊地.
就永遠只是地圖上的虛線.
它不會成為一個實線.
神的應許.
就是神給你的那份需要人的回應.
就是需要你去得地來完成.
第二方面就是停止無謂的攀比.
要專注在你自己的份上.
當人們看著其他的支派.
例如猶大或弱勢已經安頓好了.
他們心裡可能會產生疑問.
「為何我的地這麼難開墾?」.
「為何我的地這麼小?」.
「為何他們先得到地而我還未得到?」.
有這些心態.
會產生很多的懈怠在當中.
我們的服侍有時候也會這樣.
在成長班也好.
在教會其他的服侍也會這樣.
我們會看其他的人說.
「為何他的服侍這麼輕散」.
「他這麼少工作」.
「為何分這麼多工作給我?」.
有時同工也會有這樣的情況.
「為何那個同工工作這麼少」.
「要分這麼多給我?」.

$^{801}$「因為你的能力大」.
「所以就分多一點給你」.
是否應該去感恩才對?.
我告訴你一個秘密.
當我知道我要講這段經文的時候.
我第一個想法就是.
「不如跟其他人交換吧」.
因為我知道我要講十三章至十九章.
這麼長的經文.
而這段經文又是講分地這麼枯燥.
頂尖會讀的時候.
通常都會跳過不讀的.
我心裡第一個想法就是.
「不如找其他人交換吧」.
但後來我馬上又想.
「還是不要吧」.
「因為這個也是神給我的」.
「所以我應該好好去準備,預備」.
很感恩.
當我很認真去預備的時候.
我自己首先被神所提醒.
和被神所公認我有所得著.
原來這段經文雖然看起來很枯燥.
但有很多很深刻的.
屬靈的教導在裡面.
最後就是要轉化被動的等待.
為主動的勘察.
約書亞解決擔言的方法就是.
他們馬上要出去勘察.
要走遍那地.
要劃分七份.
這個應許之地.
不會自動變成你的家園.
如果你只是坐在侍羅的地方.
等神將敵人趕走.
這就是擔言.
就好像以色列人那七個支派一樣.
他們在擔言.
在我們屬靈生命裡面.
我們有時候也會說.

$^{841}$我們在等候神.
對的.
有時候我們真的需要等候神.
但很多時候我們不是真的在等候神.
我們是在躲避神.
得地是一個主動經營的過程.
如果你感覺到你現在的人生在停滯.
你嘗試去勘察一下.
神給你的現狀是怎樣.
我現在是否有些未發揮的恩賜呢?.
我的性格當中是否有些.
還未開墾的森林需要我去開墾呢?.
在神給我的這塊小小的蛋糕上面.
我怎樣可以活得精彩呢?.
擔言不去得的本質.
其實是對神分地原則的不信任.
當我們相信神給我們的糞.
是最適配的.
也是最神聖的.
我們就沒有理由繼續去推諉.
如果我們因此不用的話.
就會枯萎.
我們的產業不去經營.
也會荒廢.
所以我們每天應該感恩地去領受.
神給我的那塊小小的蛋糕.
並且勇敢地走進神給我的地界.
不需要等到環境完美的時候才去做.
因為當你起步去得那一刻.
應該去才真正落實為產業.
所以弟兄姊妹.
今年是甚麼年呢?.
馬年.
所以我們馬上行動.
快馬加鞭.
不要擔言.
我們去得這個地.
神給我們的這塊地.
神給我們的這塊蛋糕.
我們一起同心離道去祈禱.

$^{881}$慈愛天父.
因為你是那位衡量和供應的主.
你對我們的生命有一份神聖的計劃.
就好像當初你為以色列各個支派.
添溝分地一樣.
你也都按照我們生命的容量.
和你測不透的智慧.
量給我們一份獨特的產業.
求你赦免我們的遲疑.
我們今天願意帶著名為感恩的量尺.
重新去勘察你給我們的人生.
無論我們的天父.
我們的家庭.
我們的職位.
還是在我們面前的挑戰.
我們都領受為你最好的恩賜.
正如聖經所講.
你是我的產業.
是我杯中的墳.
我所得的.
你為我持守.
我們不再膽言.
要憑信心踏進你指給我們的地界.
深耕你給我們的恩賜.
成為你榮耀的見證.
奉主耶穌基督的聖名祈求.
阿們.
\newpage



\section{以弗所書 1:15-19}
\label{sec:031886}
\textbf{成為轉化世代的器皿 (華)}~~~~劉彤牧師
\newline
\newline
連結: \href{http://www.church.com.hk/acms/content.asp?site=cdc\&op=show\&type=product\&code=031886\&layout=sermon}{\texttt{031886}} ~~~~ 語音日期: 2026-03-29
\newline
\newline
\hyperref[sec:031885]{\small{< < < PREV SERMON < < <}}
~
\hyperref[sec:index_chronic]{\small{[返順時目]}}
~
\hyperref[sec:031887]{\small{> > > NEXT SERMON > > >}}
\newline
\newline
以弗所書 1:15-19
\newline
\begin{longtable}{cl}
\hline
\hline
章節 & 經文 (和合本修訂版)\\
\hline
1:15 & \begin{tabularx}{0.7\textwidth}{X} 因此，我既然聽見你們對主耶穌有信心，對眾聖徒有愛心， \end{tabularx} \\ \\ \relax
1:16 & \begin{tabularx}{0.7\textwidth}{X} 就不住地為你們感謝神，禱告的時候常常提到你們， \end{tabularx} \\ \\ \relax
1:17 & \begin{tabularx}{0.7\textwidth}{X} 求我們主耶穌基督的神，榮耀的父，把那賜人智慧和啟示的靈賜給你們，使你們真正認識他， \end{tabularx} \\ \\ \relax
1:18 & \begin{tabularx}{0.7\textwidth}{X} 照亮你們心中的眼睛，使你們知道他呼召你們來得的指望是甚麼，他在聖徒中所得榮耀的基業是何等豐盛， \end{tabularx} \\ \\ \relax
1:19 & \begin{tabularx}{0.7\textwidth}{X} 並知道他向我們這些信的人所顯的能力是何等浩大，這是照他的大能大力運行的。 \end{tabularx} \\ \\
[1ex]
\hline
\hline
\end{longtable}
$^{1}$有一段聖經在以弗所書第一章第十五到第十九節由我來唸.
保羅在這裡說因此我既聽見你們信從主耶穌親愛眾聖徒.
禱告的時候常提到你們.
求我們主耶穌基督的神榮耀的父將那賜人智慧和啟示的靈賞給你們.
使你們真知道他 並且照明你們心中的眼睛 使你們知道他的恩召有何等指望.
他在聖徒中得的基業有何等豐盛的榮耀 並知道他向我們這信的人.
所顯的能力是何等的浩大 我們就讀到這裡.
我想今天當你走進教會 你一定會發覺好熱鬧 帳篷很多.
今天有很多教會的轉化處的事工會在這裡呈現出來.
我希望你在散會之後有機會去瀏覽一下.
如果你有機會去看 你就會發現這些年我們教會的轉化處的事工發展的非常的快速.
而且涵蓋的範圍也越來越廣 我們對社區的影響力也越來越加的深遠.
回想這30年來 生命和真是從一個很小的開始.
成長到如今 30年前我們怎麼能夠想得到我們今天的景況.
神的帶領永遠是超過我們所求所想的不是嗎.
當時我們所想的只不過是建造一個堂會而已.
可是神所想的卻是一個影響世代的運動們.
當我們回顧這些年我們所走過的路.
我們真的不能夠不心存敬畏.
因為生命和這意象是遠遠大過於我們當初所曾經想像過的.
事實上這個意象的範圍也是遠大過於傳統的教會所曾經領會過的.
在過去傳統的中國教會的教會觀裡面.
我們所常常想到的還只是我們自己的需要而已.
這就好像保羅在剛才我們所讀的聖經以弗所書第一章15節在問候當地的教會的時候.
保羅說我聽見你們信從主耶穌 親愛眾聖徒.
我請問你 一個教會能夠信從主耶穌 親愛眾聖徒好不好.
好 這就不是我們理想中的理想的教會.
一個教會能夠信從主耶穌 親愛眾聖徒多美好.
這是我們理想中的教會.
可是好像保羅沒有很滿足.
保羅沒有只是停留在這裡而已.
他接下來對他們說 17節裡面說.
我要為你們求神 使你們真知道他.
意思好像說 以前這一切都不過是起步而已.
好像眼鏡還沒被打開.
所以我現在為你們祈求 求神開你們的眼睛.
能夠看見 能夠明白 能夠知道.
知道什麼呢.
知道神在你們身上恩昭是有何等的指望.
神在你們身上所要成就的基業 換句話說 神要在你們身上成就的工作.

$^{41}$是何等的榮耀.
知道神在你們身上要顯明的能力.
是何等的浩大.
換句話說 保羅是在說.
不要只是停留在信從主耶穌 親愛眾聖徒.
不要只是停留在自我需要的滿足上.
要求神開我們的眼睛.
看見神在我們身上要成就的工作 那又大又廣的事情.
生命河不但是有個建造榮耀教會的使命.
生命河不但是有一個流向萬國萬民的使命.
生命河更是一個轉化世代 轉化社區的一個使命.
而這個社區的使命 就要把神那種生命活水的江河.
能夠帶進社區當中 帶來醫治 帶來城市的轉化.
那你說我們要怎麼樣能夠成為這樣一個轉化世代的器皿呢.
今天早上從三方面跟大家來分享.
首先第一個 一定要明白神的心意.
跟你們說 你要明白神的心意.
歷代至上第十二章三十二節這裡講到說.
以薩家族的族長 他們的通達實務 知道以色列所當行的.
他們族弟兄都聽從他的命令.
我們看見這裡講到以薩家族長 是一群有影響力的人.
他怎麼帶出這種影響力 是因為他們通達實務.
換句話說 明白神的工作 也知道自己當行的事.
同樣今天我們若想影響世代 我們也必須明白神對我們的心意是什麼.
首先第一個 神要興起教會成為一支轉化的大軍.
跟你們說 教會要成為轉化的大軍 你跟他說.
我深深的相信 馬太福音五章十四節.
是神要對今天的教會所說的話語.
馬太福音五章第十四節 我們一同來唸.
你們是世上的光 承照在山上 是不能夠隱藏的.
這裡說 你們 請問誰是你們 就是我們.
主耶穌說 你們就是世上的光.
跟你們說 你就是世上的光.
你知道過去 這個觀念很多時候 我們還是沒辦法真的放在身 印在我們心裡.
多少時候過去 教會還一直是把自己定位.
孤立在一個角落裡面的象牙塔.
以至於我們並不覺得我們對這個社區 城市會有什麼影響力.
以至於好像過去的教會對城市 社區也不聞不問.
但我們看見神的心意卻不是如此.
主說 你們就是世上的光.

$^{81}$你們是教會的整體 也是每一位信徒.
主說 你們就是世上的光.
而這個光是必須要照在人前.
不能點在斗底下.
就像主在這裡說 城市要站在山上 是不能夠隱藏的.
這個光不是要我們隱藏起來.
這個光要顯明在人前 叫世人看見 歸榮耀給神的.
你們 換句話說 面對現今黑暗的世代.
弟兄姊妹 教會是有責任的.
教會應當要實際的進入社區 為主發光.
才能夠扭轉這個黑暗的世代.
你們 事實上 越來越看見 神在現今在幕後的世代.
就要為祂自己名的緣故 興起一支極大的軍隊來.
你們 事實上 這也是我們教會今年的主題.
興起一支屬神的大軍.
在以西節注37章裡面 你看見神吩咐以西節.
向風發預言.
我們看見當風降臨在這個枯乾的骸骨上的時候.
結果骸骨就復活 成為有靈的活人.
我們看見這支神的手正在今天.
在信徒生命當中 在教會上所做的工作不是嗎.
但為什麼神要來做這種復興的工作呢.
神為什麼要在我們的生命當中.
叫我們的靈能夠活過來呢.
神為什麼要降臨在神的百姓身上.
能夠在枯骨當中 能夠在教會當中來做更新.
做復興的作為呢 為什麼呢.
以西節注37章第十節 講得很清楚.
我們一同來念 於是我遵命說預言 氣息就進入骸骨.
骸骨變活了 並且站起來怎麼樣.
成為極大的軍隊.
這是正是今天我們所要做的事情.
這些日子來我真的看見神不但在我們教會.
要興起極大的軍隊.
當我在世界各地在宣教的時候.
看見神在世界各地也興起這支極大的軍隊.
是一支轉化的軍隊.
要在這種彎曲背面的世代當中.
顯明著神的榮耀來 Amen.
而你我就在這支大軍當中.

$^{121}$Hallelujah.
給你們說 你就是屬靈的大軍.
而且不要忘記我們今年要背聖經.
所以外面有背經的攤位.
不要忘記去背一下.
八福是非常容易的.
就算你今天忘記了 散會趕快背一下.
也是可以背得好的 OK.
鼓勵你趕快去背一下聖經 Amen.
我們看見神對教會的心意.
不是讓我們彼此相愛而已.
彼此相愛是很重要.
但神的心意是要我們興起成為一個轉化的大軍.
第二 神要我們主動來為主而活.
誰要我們主動為主而活.
弟兄姊妹 今天你不但是要明白.
神對教會整體的心意是什麼.
更重要的 到底什麼是神在你身上的心意呢.
這是更重要的一個問題.
以弗洛索斯第五章十七節.
我非常喜歡保羅在這裡說這句話.
以弗洛索斯第五章十七節.
保羅說不要做糊塗人.
跟你們說不要做糊塗人.
要明白主的旨意如何嗎.
今天我看見多少基督徒就是糊裡糊塗.
信仰糊裡糊塗 生活也糊裡糊塗.
我們看見多少時候我們不知道為什麼而活.
不知道我們為什麼活在 神把我們擺在這裡.
也不知道神為什麼把我們放在.
這種生活的環境當中.
我好像永遠是糊裡糊塗地過日子.
所以保羅在這裡提醒我們.
不要做糊塗人.
也許你過去曾經糊塗.
但求神幫助我們從今天開始.
要主動積極尋求神的心意.
Amen.
你還記得當保羅在往大馬色路上.
忽然被主的大光光照.

$^{161}$他就撲倒在地.
當他撲倒在地的時候.
他就問了在一生當中最重要的兩個問題.
第一個問題他問說主啊你是誰.
他知道這位用大光光照他是他生命的主宰.
但他不知道他是誰.
他說主啊你是誰.
而當耶穌對他說我就是你所逼迫的耶穌的時候.
保羅立刻接下來問了他一生最重要第二個問題.
那主啊我當作什麼.
弟兄姊妹在你一生當中你問過這兩個問題嗎.
也許你已經知道耶穌的名字.
但你可知道耶穌就是你生命的主宰嗎.
你有沒有真正的讓耶穌來主導你的生命.
讓耶穌來指引你生命的道路呢.
你可知道在你一生當中.
最重要的不是你想做什麼.
最重要的是主要你做什麼.
我看見太多人活在現今世代一天到晚只想著我想要做什麼.
我今天要鄭重告訴你.
在你一生當中你想做什麼一點都不重要.
最重要的是主要你做什麼.
如果你確認耶穌是你生命的主.
那你就要認真的來尋求神說主啊那我當作什麼.
弟兄姊我常常覺得今天做一個基督徒.
我們不能只是坐在這裡被動的.
過一個得過且過的基督徒而已.
因為我們若不主動.
我告訴你魔鬼撒旦就會採取主動.
今天我們若不主動.
仇敵就會主動.
A.W. Tozer曾經這樣說.
他說我們如果一天不親近神.
我們就是一天敞開門.
容許罪惡跟錯誤進入我們的生命.
換句話說.
如果我們不主動.
這個世界一定會採取主動的.
如果不主動來影響這個世界.
這個世界一定主動來影響我們.

$^{201}$今天其實若不是我們主動來影響世界.
我們就是被這個世界所影響.
就是這樣子.
這為什麼許多時候你會覺得做一個基督徒.
有時候覺得做了很辛苦的原因.
為什麼.
因為你一天到晚就是被動的被挨打.
你想想看一天到晚被動被挨打不是真的很辛苦嗎.
No 今天做一個基督徒.
我們不是只是被動的被挨打.
你要明白你是基督徒組成的大軍.
我們需要起來主動積極.
尋求主的心意.
為主而活.
只有當我們積極主動.
我們才可能在這樣一個世代.
成為神手中的器皿.
被神所使用來影響這個世代.
Amen.
跟你旁邊再一次說.
不要再被動要主動積極.
跟他說.
Hallelujah.
所以要明白神對你的心意重要.
第二.
要起來為你的環境禱告.
各位要起來為你的環境禱告.
在以利亞的時代.
那個時候以色列人正遭遇到三年半的乾旱.
到處都是饑荒.
神就呼召以利亞上加密山頂去禱告.
列王經上十八章四十一到四十四節.
這段經文也是我們很熟悉.
這裡說以利亞對亞哈說.
你現在可以上去吃喝.
因為有多雨的響聲.
以利亞知道神要降下大雨.
以利亞知道神的心意.
神渴望降下負心來做恢復的工作.
很寶貴.

$^{241}$在四十二節.
亞哈就上去吃喝.
以利亞就上了加密山頂.
屈身在地.
將臉浮在兩棲之中.
這是一種生產式的姿態.
是一種生產式的禱告.
四十三節對僕人說.
沒有什麼.
如此七次.
第四十節.
第七次僕人說.
我看見一小片雲從海裡上來.
不過如人手那樣大.
以利亞說你上去告訴亞哈.
當套車下去.
免得被雨阻擋.
我們看見當時.
以利亞聽見那多雨的響聲.
他知道神將要工作.
他沒有停留在那裡.
他上到加密山頂.
屈身在地.
將臉浮在兩棲中.
他禱告了七次.
神降下大雨.
就解除了乾旱.
同樣我今天感覺到.
神也照樣在呼召我們.
要為我們所處的大環境來禱告.
多好受到面對現今的世代.
充滿動盪的世代.
充滿罪惡跟卑劣的世代.
我常常聽見人在彼此埋怨.
很哀嘆說.
四風日下人心不古.
多少時候我們看見.
我們所處的國家美國.
離開了神的價值觀.
離開神的道路.

$^{281}$越來越遙遠.
很多時候我們心中.
充滿了憂慮跟恐懼.
但是事實上.
地球人要告訴你.
越是在黑暗的時代.
越是神的復興.
要來降臨的時代.
在整個教會歷史當中.
你看見就是如此.
信仰.
若是在長期安逸的環境當中.
就會自然而然的走向死亡.
反而信仰.
是在充滿挑戰的環境當中.
才能夠重新被活潑.
會恢復過來.
所以你不要哀嘆.
現今的世代.
充滿黑暗跟挑戰.
這正是神為我們預備的時機.
神給我們的時機.
在這樣一個時刻.
能夠為主而活.
與神來同工.
因為越是黑暗的世代.
越是在艱難的時刻.
我們看見越是神呼召我們.
要起來為我們環境禱告的時刻.
今天你可能知道禱告的重要.
今天你可能知道.
我們需要慈守神的應許來禱告.
但如果你還只是坐在那裡.
聽聽說說而已.
那還是沒有用的.
我們看見以利亞.
她聽見那多雨的響聲.
她知道這是個時機.
她聽見神的呼召.
但如果說以利亞.

$^{321}$她當時聽見.
看見.
卻不起來.
上加密山頂去禱告的話.
我告訴你.
雨還是不會降下來的.
我們看見乾旱還是不能夠解決的.
所以今天我們要起來.
為我們所處的環境來禱告.
為了家庭來禱告吧.
為了小組來禱告吧.
為教會來禱告吧.
為社區來禱告.
為城市為國家來禱告吧.
求神降下恩雨來.
解除我們屬靈的乾旱.
不但要起來禱告.
更重要要持續的禱告.
跟你們說要持續的禱告.
跟他說.
我們看見當以利亞在那裡.
聽見多雨的響聲.
她知道神要降下大雨來的時候.
我相信.
她能夠上山上禱告的時候.
她怎麼想都應該.
一禱告大雨就降臨下來.
但事實上卻不是如此.
她在山上禱告.
她每禱告一次.
去看有什麼動靜沒有.
扶人回來說沒有.
一次說沒有.
第二次也是沒有.
第三次還是沒有.
換了我們可能說.
哎呀.
大概聽錯了吧.
什麼多雨的響聲.
可能沒有啊.

$^{361}$禱告也沒用.
還是試試其他的方法吧.
但以利亞卻沒有停止.
她持續的禱告.
直到第七次.
她才看到.
像人手這麼大的雲飄過來.
她就知道.
神要準備降下大雨來了.
第二點同樣.
我深深的相信.
今天也是神要呼召.
為我們所處的環境禱告的時刻.
雖然很多時候.
我們看看我們周圍的環境.
好像怎麼禱告都沒什麼轉變.
你說我們的家庭禱告.
我們的兒女禱告.
為我們的社區教會小組.
好像禱告來禱告去沒什麼結果.
連片雲都沒有.
但不要灰心.
神說要持續的禱告.
始終.
神要賜下她偉大奇妙的作為.
我們看列王記上.
十八章四十五節.
非常寶貴的聲音.
我們一同來念.
撒時間天陰風雲黑暗降下大雨.
呀哈就坐車往耶穌那邊去.
好喜歡這句話.
撒時間.
跟我一起來說.
撒時間.
什麼意思.
忽然間.
在你想不到的時候.
撒時間.
想不到的時候.

$^{401}$神就要降下一致的恩告.
神要降下突破的恩告.
神要賜下恢復的恩告.
神要賜下復興的恩告.
在我們中間.
你要繼續為你的家庭來禱告.
為你的兒女來禱告.
為你的小組來禱告.
為社區城市來禱告.
為這個國家來禱告.
不要灰心.
持續的禱告.
直到神工作.
Amen.
Hallelujah.
最後第三.
要效法耶穌倒臣肉身的榜樣.
跟他說.
你要知道耶穌基督.
之所以能夠來到這個世界.
帶來如此翻天覆地的影響力.
我告訴你.
並不是因為他所行的神蹟騎士而已.
為什麼耶穌能夠帶來這麼大的影響力.
最主要的原因.
是因為他倒臣肉身.
親自來到這個世界上.
體嘗過我們人世間一切的憂患.
軟弱跟痛苦.
唯有他能夠明白.
我們人在肉體當中的掙扎跟痛苦.
唯有他能夠在十字架上代替我們而死.
不但死.
而且從死裡復活.
帶給我們得勝的喜樂跟盼望.
倒臣肉身.
是真正能夠帶來我們生命影響的一個關鍵.
這個倒臣肉身.
不單單是指的耶穌基督進入世界的行動.
更是指的一個願意進入世界.

$^{441}$認同我們需要的一種心態.
因為你越來越看見.
這是今天教會最需要活出來的真理.
因為太多時候.
教會覺得這個世界太黑暗了.
想急於離開這個世界.
卻沒想到耶穌正是為這個黑暗世界而來.
耶穌正是因為這個世界的黑暗.
進入到我們中間來.
今天我們若是想影響世代.
我們就必須效法主耶穌倒臣肉身的榜樣.
你說怎麼效法呢.
兩件事情.
第一個多一點了解少一點定罪吧.
跟你朋友說多一點了解少一點定罪.
今天教會太會定罪了.
太會定罪不是嗎.
但你看耶穌基督他卻不是如此.
有一次有一個行淫的婦人被抓了.
到耶穌面前.
按照當時的律法.
一個行淫的婦人要被石頭打死.
眾人期待耶穌要他們把石頭打死這個婦人.
耶穌卻說了一句話.
在我們中間有誰沒有罪的.
先拿石頭打他.
後來眾人一個一個走了.
耶穌就問那個婦人說.
沒有人定你罪嗎.
婦人說沒有.
耶穌說那我也不定你的罪.
要注意喔.
耶穌沒有說她沒有罪喔.
耶穌說我不定你的罪.
從今天開始不要再犯罪.
耶穌給她一個重新開始的機會.
我告訴你這不是今天這個世界最需要的訊息嗎.
多少人渴望一個重新開始的機會.
教會能不能不再定罪.
而是給人重新開始的機會呢.

$^{481}$第二.
多一點行動少一點教條吧.
給你比方說多一點行動少一點教條吧.
教會太會講教條了.
都是教條.
耶穌卻不是如此.
在迦納的婚宴他們沒有酒.
換了我們是耶穌可能會第一件事.
把那個主人叫來.
你知道為什麼你們沒有酒了嗎.
因為你犯罪得罪神了.
悔改.
先悔改恩典就領到.
耶穌沒有.
耶穌先把水變成了酒.
解決了他們的問題.
滿足了他們的需要.
你想想看那天婚宴結束之後.
那個主人會怎麼反應.
那主人這人是誰啊.
帶我到他面前我一定要來跟從他嗎.
不是嗎.
耶穌沒有先叫一個人悔改.
卻先充滿憐憫解決他們的問題.
然後再把他們帶到救恩的裡面.
不是嗎.
多一點行動.
憐憫的行動.
恩慈的行動.
少一點教條.
讓人能夠明白神的愛.
這些年我在教會裡面服事.
我就發現.
你知道很多時候我發現.
教會跟社會之間.
真的很大的鴻溝你知道嗎.
一方面我看見.
這個世界是多麼需要愛啊.
另一方面我看見.
教會裡面也真的有很多愛.

$^{521}$但好像這個愛就是傳遞不出去嘛.
不是說教會裡面弟兄姊妹沒有愛心.
很有愛心啊.
但卻是找不到一個合適的管道.
能把神的愛從我們中間傳遞出去嘛.
我就這些年越來越明白.
愛是需要一個平台的.
什麼是平台啊.
平台就是一個.
使我們可以進入社區.
進入人群的橋樑嘛.
透過這個平台.
一方面使社區的人.
不一定要走到教會裡面.
才能夠領受神的愛嘛.
另一方面幫助教會的弟兄姊妹.
能夠實際的用行動來表達神的愛.
這是為什麼這些年.
我們花這麼多的功夫.
在教會裡面.
不斷打造不同平台的原因啊.
基金會是一種平台.
透過基金會我們可以服事到.
很多在社會當中.
平常我們無法接觸到的人群.
我們食物發放中心現在每個月.
我們可以服事到超過四萬個.
在灣區有需要的家庭.
Hallelujah.
在非洲兒童之家.
我們也收養了.
超過三百多位無家可歸的孩童.
優視頻道也是一種平台啊.
透過優視我們北美華.
在為北美華人提供了一個.
用基督教價值觀潛移默化.
來影響北美華人文化的電視台嘛.
我們透過這樣一種電視台.
其實更重要是所拍出的短視頻.
能夠潛移默化的影響.

$^{561}$美國華人的文化.
職場工作坊也是一種平台.
透過個別職場教練的方式.
我們幫助一個人能夠在職場上.
不但被提升有突破.
而且經歷到生命改變的大能.
非洲職業訓練學校.
也提供的是一種平台.
我們在多哥這個國家.
不但幫助很多的年輕人.
找到合適的工作.
而且開設了生命和診所.
在醫藥缺乏的地區提供了最切身的救助.
除此之外我們看見生命和社區中心.
藉著禱告帳篷的工作.
服侍很多本地有需要的人群.
而生命和教母領袖學院.
也栽培了神國未來的工人.
這些攤位我們今天在外面.
都有很多的攤位.
我希望大家能夠留步去看一下.
關懷一下.
不但關心.
而且希望參與並支持他們的施工.
你光聽我講這些.
你就明白.
為什麼你一天到晚覺得.
我們教會施工沒完沒了.
很多人常常埋怨.
怎麼搞的沒完沒了.
沒錯.
但我告訴你.
我們的施工之所以沒完沒了.
正是因為這個世界的需要沒完沒了.
只有當我們用著愛心的行動.
來深入社區的時候.
我們才可能有效的向世界宣告說.
神就是如此深愛世人.
今天你是怎麼來看我們所處的環境.
今天你是怎麼看我們所處的社區.

$^{601}$我們所處的城市.
你相信神能夠藉著我們帶來改變嗎.
很多時候我們只看見.
城市裡面充滿了問題.
我們只看見社區當中充滿了問題.
我們也看見我們的環境當中.
充滿了各樣的問題.
但事實上我告訴你.
每一個問題也都表明了.
神可以介入的機會.
你我是否願意成為神解答的管道呢.
當年保羅問主說.
主啊我當做什麼.
這個問題帶來他一生的改變.
也影響了當時的時代.
而今天我們照樣需要回到神面前.
認真的面對這個問題.
主啊我當做什麼.
不只是神對教會有一個轉化的心意.
更是對我們神有一個心意.
主我當做什麼.
只有當擺上自己持續地禱告.
並且用愛心的行動深入社區的時候.
我們必然能夠成為社區當中.
帶來影響力神的轉化的奇名.
我們一起禱告.
親愛的我們再次來到你面前.
謝謝主耶穌.
你在我們還做罪人的時候.
你先來到這個世界.
你道成肉身帶來拯救.
主啊再一次.
讓我們效法你道成肉身的榜樣.
在你所給我們的環境.
你所給我們的社區城市當中.
也讓我們照樣把你的愛.
藉著我們能夠湧流出去.
湧流到我們身旁周圍的每一個人.
湧流到我們的家庭.
我們的社區.

$^{641}$我們的城市裡面.
好讓世人看見.
就明白神.
主你真的深愛世人.
主再一次.
開我們的眼睛.
看見的都是你給我們的機會.
讓我們心中不再埋怨.
不再苦澀.
而是學習問保羅的同一個問題.
主在這種困境當中.
主我當做什麼.
在這充滿問題的環境當中.
我當做什麼.
求主親自向我們說話.
賜下勇氣.
賜下智慧.
讓我們成為你手中的器皿.
帶來影響力.
我們謝謝主.
奉耶穌得勝的名祈求.
Amen.
\newpage



\section{馬可福音 10:42-45}
\label{sec:031887}
\textbf{耶穌改變世界的方法}~~~~招世超牧師
\newline
\newline
連結: \href{http://www.church.com.hk/acms/content.asp?site=cdc\&op=show\&type=product\&code=031887\&layout=sermon}{\texttt{031887}} ~~~~ 語音日期: 2026-03-29
\newline
\newline
\hyperref[sec:031886]{\small{< < < PREV SERMON < < <}}
~
\hyperref[sec:index_chronic]{\small{[返順時目]}}
~
\hyperref[sec:031888]{\small{> > > NEXT SERMON > > >}}
\newline
\newline
馬可福音 10:42-45
\newline
\begin{longtable}{cl}
\hline
\hline
章節 & 經文 (和合本修訂版)\\
\hline
10:42 & \begin{tabularx}{0.7\textwidth}{X} 耶穌叫了他們來，對他們說：「你們知道，外邦人有君王作主治理他們，有大臣操權管轄他們。 \end{tabularx} \\ \\ \relax
10:43 & \begin{tabularx}{0.7\textwidth}{X} 但是在你們中間，不可這樣。你們中間誰願為大，就要作你們的用人； \end{tabularx} \\ \\ \relax
10:44 & \begin{tabularx}{0.7\textwidth}{X} 在你們中間誰願為首，就要作眾人的僕人。 \end{tabularx} \\ \\ \relax
10:45 & \begin{tabularx}{0.7\textwidth}{X} 因為人子來，並不是要受人的服事，乃是要服事人，並且要捨命作多人的贖價。」 \end{tabularx} \\ \\
[1ex]
\hline
\hline
\end{longtable}
$^{1}$為了預備今天跟大家分享這個題目.
我就上去問一問AI.
大家知道現在很popular.
我們傳道人都要學一下如何善用人工智能.
問一個問題.
自從人類歷史以來.
對人類社會造成最大影響的是誰?.
問一個問題 對整個世界帶來改變的是誰?.
本來我只是期待一個答案.
誰知他給了我幾個名字.
首先是科學和數學方面的Isaac Newton.
就是牛頓.
然後是技術方面的約翰內斯古·藤堡.
就是Johannes Gettinsburg.
然後是政治和宗教方面的穆罕默德.
在地緣政治方面的就是我們都聽過的成吉思汗.
最後一個.
他說對宗教和社會造成最大影響的就是拿撒勒人耶穌.
所以他給了我幾個名字.
都是影響人類歷史改變整個世界的重要人物.
他收集了很多的資料.
歸納了有這個結論.
我相信大家聽到這些名字的時候都不會奇怪.
因為事實上他們在不同的領域當中.
真的改變了人類的歷史和改變了我們這個世界.
但是我在想.
其中這幾個人當中.
是有意識地想去改變人類社會.
改變人類歷史的.
那個是誰呢.
當然你都知道牧師說的是誰.
對我來說我覺得就是拿撒勒人耶穌.
為什麼我這樣說呢.
因為耶穌在離開世界之前.
很清楚地吩咐門徒一件事.
大家都知道是什麼.
所以你們要去使萬民作我的門徒.
奉使父子聖利明和吉他門施咒.
所以耶穌在離開世界之前.
就把這個使命我們稱之為大使命.

$^{41}$交給所有的基督徒交給教會.
他的意象是什麼呢.
耶穌的意象就是要所有世界的萬民.
萬族萬民都成為他的門徒.
都成為一個君尊的祭司.
神的子女.
天國的子民.
這是耶穌在世界上的任務.
和他見到他需要成就要做好的工作.
我不清楚穆罕默德.
他是一個宗教領袖.
到底穆罕默德有沒有這樣的想法呢.
I don't know.
不過如果有的話.
我很清楚知道他用什麼方法.
他用什麼方法呢.
大家都聽過了.
一手拿著可蘭經.
另外一手拿著一把刀.
所以一手是刀一手是可蘭經.
這個特別是在伊斯蘭教早期擴張的時候.
這個就是很明顯的.
結果因為這樣的緣故.
穆罕默德就成立了一個很強大的阿拉伯帝國.
這個帝國是以神權統治作為藍本.
就像今天的伊朗一樣.
神權統治.
現在怎樣大家自己去判斷.
但是從一個社會的角度來看.
這個做法是一個柔性的信仰.
和硬性統治的結合.
一手拿著可蘭經.
一手拿著一把刀.
其實如果我們看回聖經的時候.
在舊約當中的以色列國.
其實都是一個神權統治的王國.
但是我們發覺人真的太軟弱了.
人都是有罪性的.
一個不小心就會犯罪.
犯罪會怎樣呢.

$^{81}$犯罪就是離開上帝.
所以後來又分了南國北國.
無論怎樣也好.
當那些人做了皇帝有權有位之後.
將神放在一邊.
就將自己變成一個獨裁者.
不單止將神放在一邊.
甚至還會引進了其他很多外國的神明.
所以其實在舊約聖經裡面來說.
王國時代結果就是以色列被滅國.
到了新約時代.
耶穌是怎樣看待這個問題的呢.
大家記得耶穌說過很多話.
不過和這個有關的.
可能大家都熟悉了.
耶穌是怎樣和那些人說的呢.
當他在彼拉多面前的時候.
他說我的國不屬這世界.
我的國如果屬於這世界的時候.
我的神僕就會起來和你們爭戰.
讓我就不用落在猶太人的手裡面.
只是我的國不屬這世界.
這是耶穌在彼拉多面前說的.
事實上耶穌也曾經對那些想陷害他.
想試探他的人.
是曾經說過的.
該殺的物.
該殺的是誰.
當時羅馬帝國的皇帝.
該殺的物就給了該殺.
但是神的物就當歸給神.
所以從耶穌的時代開始.
政教分離的觀念.
這個教義就出來了.
剛才說到耶穌可以改變不單單宗教.
可以改變整個社會文化的一個人.
耶穌不用武力.
不用強權.
不用軍事的力量去統治和改變世界.
那麼耶穌可以用什麼方法去改變這個世界呢.

$^{121}$那麼耶穌有多少個門徒呢.
那時候跟著耶穌.
跟了幾年時間.
他們其實有一個極概念.
這個觀念是跟所有人一樣的.
就是他們相信救主神所應許賜給他們的一個.
是會帶領他們去推翻羅馬帝國的統治.
重建以色列國.
那麼耶穌就很自然的成為王.
當他們看到.
現在耶穌這麼多人擁戴.
氣勢如虹.
竟然就開始想到.
以後如果耶穌得國的時候.
我坐在哪裡呢.
開始想怎麼上位.
怎麼去做頭.
而雅各和約翰的媽媽.
跟耶穌有一些裙帶關係.
大家知不知道.
原來雅各和約翰的爸爸叫西比泰.
他媽媽叫薩羅米.
是瑪利亞的姐姐.
他跟耶穌有一些親戚關係.
西比泰的妻子薩羅米就帶了他們兩個兒子.
就是雅各和約翰去找耶穌.
喂 耶穌啊.
當你得國的時候.
當你成為國王的時候.
可不可以我們兩個兒子.
一個做你的左丞相.
一個做你的右丞相.
來找耶穌.
當然大家如果熟悉聖經的話.
都知道他們這個行動.
當然是引起其他門徒的不滿.
沒有教錯.
你拉關係想這樣上位.
為什麼這些門徒會這樣想.
因為他們同樣有這樣的想法.

$^{161}$如果耶穌真是國王.
誰不想一個在耶穌的左邊.
一個在耶穌的右邊.
這些門徒同時在那個時間.
雖然跟了耶穌幾年的時間.
但是沒辦法的.
受社會的風氣文化影響.
誰不想做一個位高權重.
可以統治人 可以操控人的人.
有權有勢 幾何.
好事來的.
耶穌知道門徒也有這樣的想法.
不單單是雅各約翰.
不單單是他的父母.
於是耶穌就叫了一群門徒來.
就跟他們說.
你們想做阿頭.
你們想做左丞相 右丞相.
你要學做一個謙卑服侍人的人.
你要做一個僕人領袖.
我們看看馬可福音第十章怎麼說.
耶穌就叫他們來.
對他們說.
你們知道外邦人.
有尊為君王的治理他們.
有大臣操權管束他們.
只是在你們中間.
不是這樣.
你們中間誰願為大.
就必作你們的傭人.
就是傭人.
在你們中間.
誰願為首.
就必作眾人的僕人.
因為人子來.
說他自己.
並不是要受人的服侍.
乃是要服侍人.
並且要捨命作多人的屬教.
耶穌跟門徒說.

$^{201}$你們中間誰想做大.
想做阿頭.
那你就先要去做人家的傭人.
你們中間誰要做左丞宣右丞相.
人之下 萬人之上.
就必須要做眾人的僕人.
今日我們姜濤偉教經常聽這些說話.
聽慣了好像沒什麼緊要.
但當時的門徒來說.
有沒有搞錯 耶穌說這樣的話.
因為那個世界是什麼世界.
那個環境是弱肉強食.
以強權軍力去征服世界的時間.
耶穌的教導根本不可思議.
對他們來說是難以理解的一個教導.
做阿頭.
做帶頭的那個.
不講權.
不用武力.
怎樣去炸.
耶穌說你要做一個傭人.
好像一個僕人一樣.
所以今日我的題目是.
要學做一個僕人領袖.
但耶穌沒有用這個名稱.
即是僕人領袖.
如果你在外面做工.
或者你讀書讀領袖學的時候.
你會發覺僕人領袖這個名詞.
是1970年.
一個領袖學家.
叫做Robert Greenleaf.
首次提出僕人領袖這個名稱.
這個名稱的意思是什麼呢.
他說真正的領袖.
不是先追求權力和地位.
而是出於服侍他人的心.
進而自然去帶領其他人.
所以僕人領袖這個名詞.
是Robert Greenleaf提出來的.

$^{241}$Greenleaf.
所提出這個理念.
和耶穌的教導是否不謀而合.
很相似.
相似而已.
不是完全一樣.
因為耶穌的教導門徒.
還強調了什麼.
在字裡的行間.
一會兒你會更加清楚.
耶穌所強調這個領袖.
不單單是一個僕人.
一個傭人.
更加是謙卑去服侍.
這是講到心態的問題.
服侍多於地步.
甚至要為人犧牲.
因為這是出於神的恩賜.
為什麼他這樣說呢.
我們看約翰福音13章.
約翰福音13章第3節說.
耶穌知道.
父已將萬有都交在他的手裡.
而且知道自己是從神那裡來的.
這個認知很重要.
但是又要歸到神那裡去.
第4節.
就離席站起來脫了衣服.
拿一條手巾束腰.
隨後把水倒在盆裡.
就洗門徒的腳.
並用自己所束的手巾去擦乾.
耶穌有沒有權柄.
有沒有.
絕對有.
他自己知不知道.
清楚知道.
因為他也說天父已將萬有.
所有的東西都交在他的手裡.
所有的東西都在耶穌的掌握當中.

$^{281}$所以耶穌知道自己是一個有權勢的人.
正如上個星期.
如果記得我上個星期的目撕講道.
我講了兩點.
第一點是講什麼.
一個基督徒.
或者我從耶穌的身上學習的時候.
凡事以天父的事為念.
凡事以父親的事為念.
還有.
耶穌在這個世界做什麼.
是為神而活.
所以雖然他有權柄能力.
可以用他的力量.
去改變他做很多事情.
但是天父怎麼想的.
以天父的事為念.
如果我要為神而活.
我應該怎麼做呢.
釋經說.
當他知道所有的東西都在他手中.
但是他又知道.
上帝差遣他來世界.
要做那件事.
要完成任務的時間.
就快到了.
那我就快要到了.
要完成任務的時間就快到了.
那我應該做些什麼呢.
天父要我做些什麼呢.
做一件很特別的事.
耶穌就用替門徒洗腳的行動.
向門徒示範.
示範什麼呢.
要謙卑地互相服侍.
並且用自己的生命.
和我們生命中的影響力.
去引領其他人.
跟隨耶穌.
做耶穌基督的門徒.

$^{321}$所以不知道我們看不看到.
耶穌帶領的門徒.
耶穌是領導者.
耶穌帶領的門徒的方法.
是不是很反傳統.
反潮流反傳統.
他沒有用天父賜給他的權柄.
能力去操控去管治他的跟隨者.
有沒有.
沒有.
他沒有用老師的身份.
他是老師.
當時的老師是很被尊敬的.
很有權威的.
特別是那些經學老師.
多厲害.
耶穌沒有用他的身份.
沒有用他的權威.
去指許學生.
做這樣做那樣.
有沒有.
沒有.
相反他竟然和學生洗腳.
這也是不可思議.
因為這只是那些下等的工人.
才會做的事.
但是弟兄姊妹你們有沒有看到.
耶穌這個謙卑服侍的行動.
是影響了門徒的一生.
讓他的門徒都成為一個僕人領袖.
是賺取.
是earn回來的.
賺取了那些跟隨者.
對他們的尊敬信服和跟從.
早安弟兄姊妹.
我真的有很多感恩的事.
所以我鼓勵大家要寫一些感恩的事.
我回來做ISB之後.
我發覺很多弟兄姊妹.
我們當中.

$^{361}$是願意起來服侍教會.
這是我很感恩的事情.
我們有多少牧師.
能夠做多少事.
弟兄姊妹不幫忙.
不參與不服侍.
我們做不到任何的事.
有很多弟兄姊妹.
願意起來服侍教會.
問題是甚麼呢.
很多時候有人服侍教會.
因為教會有需要.
教會有需要.
我有時間.
或者我又願意幫忙.
那我們就去做.
牧師今天說領袖.
我沒想過做領袖的.
不過有需要.
能夠用我.
我又願意.
我有時間.
或者我有甚麼能力.
我就幫忙.
不過不要叫我做領袖.
我發覺在教會裡面.
原來很多人怕做領袖.
我不知道你怕不怕.
怕不怕在教會裡面做領袖.
我寧可做僕人.
去服侍別人容易些.
因為做領袖要負責任.
因為你不是在外面打工.
那麼簡單.
打工是賺錢.
是去服侍上帝.
所以很多時候我們怕負責任.
因為要向上帝負責.
我們又怕自己做得不好.
那就無謂了.

$^{401}$有沒有這樣的想法.
我鼓勵大家.
頂尖我鼓勵大家.
無論你有沒有想過自己做領袖也好.
當你有機會參與服侍的時候.
要學耶穌一樣.
用一個僕人的心態去服侍.
我剛才說過Robert Greenleaf.
他這個僕人領袖的觀念.
跟耶穌好像不謀而合.
但不是完全一樣.
因為他的觀念是來自人本主義.
Humanism.
和靈性的哲學.
Spiritual Theology.
即是Philosophy.
這是建基於人的尊嚴.
和群體的福祉.
如何令整個群體蒙福.
這是基礎上的.
他會否受基督教的影響呢.
很多東西都很接近.
這個我不知道.
可能會是都未定.
因為事實上他說的話.
很多東西跟耶穌的教導很相似.
但他不是一個基督徒.
我找不到什麼資料.
Greenleaf是一個基督徒.
但耶穌的教導.
跟他的教導不一樣.
耶穌的教導是來自.
至高的上帝.
是來自神對人的啟示.
是建基於神的主權.
和神的榮耀上.
為什麼呢.
因為神要做一件事.
用他的做法.
是為了完成上帝拯救世界.

$^{441}$改變人類命運的計劃.
所以我們要學習成為.
神的僕人.
有三個term.
第一個.
傳統的領袖是怎樣的.
第二.
剛才牧師提及過.
僕人領袖.
第三類人是什麼.
神的僕人.
有什麼不同的地方呢.
傳統的領袖是說什麼.
怎樣去帶領人.
怎樣去管理人.
管治人.
僕人領袖是說怎樣去服侍人.
怎樣去建立人.
神的僕人.
是怎樣去服侍神.
榮耀神.
傳統的領袖.
強調什麼.
是強調你的power.
你有多少個power.
你有多少個authority.
你的地位身份.
僕人領袖.
是強調道德的影響力.
怎樣用我的生命去影響你.
而神的僕人.
是強調屬靈的權柄.
我們相信神.
我們相信聖靈.
祂賜給我們權柄能力.
怎樣去服侍和帶領.
傳統的領袖.
為了工作的成果去服務.
最重要是你的outcome.
你的result是怎樣.

$^{481}$總之就是job done.
但是僕人領袖.
是為了人的成長.
很不一樣.
為了人能夠grow up.
能夠成長.
能夠mature.
去服侍.
但是神的僕人.
是為了人生命的更新.
當然說到屬靈生命的更新.
而最後.
傳統的領袖.
是用控制和order.
去帶領他的下屬.
僕人領袖.
是引導群體達到一個共識.
然後一起去做.
神的僕人.
是採用真理為導向.
帶領跟隨者去走神的路.
做神的工作.
所以僕人領袖.
和神的僕人有不同.
和傳統的領袖.
更加有很大的分別.
今天這個時間.
第一節我想講給你知.
耶穌是很想我們成為神的僕人.
我們一個僕人領袖.
神希望你和我一樣.
可以成為神的僕人.
學習怎樣像耶穌一樣.
去服侍上帝.
所以.
我認為教會需要很多人幫忙.
但是當你服侍的時候.
我希望你記住.
我們在教會裡面參與服侍.
不是為了完成一個工作.

$^{521}$做好一件事.
下星期有聖誕節.
下星期五晚有崇拜.
很多人要教很多東西.
又要準備吃東西.
又要聯絡一些人.
聽回來好像要做很多工作.
做好一件事.
等下星期多點人來崇拜.
一起敬拜上帝.
無論如何.
不.
當然這些事我們要做.
但是在我們做的時候.
如果我們明白神的僕人觀念的時候.
我們還要去想一下.
我們在服侍人的時候.
還可以怎樣去幫助.
和我們一起同工的弟兄姊妹成功.
能夠幫助他們.
讓他們做得更加好.
讓他們能夠在事奉當中.
去長進.
去經歷神.
不知道我們在服侍的時間.
除了去做事之外.
如果只是去做事就很慘很辛苦.
平時上班都很辛苦.
回來教會還要去做事.
但如果你改變了這個觀念就不同了.
你不是為了做一件事.
然後就倒閉那麼簡單.
你是希望透過這個一起去經歷上帝.
幫助弟兄姊妹越來越成熟.
生命被神改變.
所以在服侍的過程當中.
我們必須去學習關心.
弟兄姊妹的需要.
他們怎樣想事.
他們怎樣感覺.

$^{561}$他們服侍的時候要付出什麼代價.
要面對什麼挑戰.
我們可以怎樣幫助他們.
讓他們能夠開心快樂地去事奉.
除了服侍之外.
我們有沒有投資時間精力.
跟他們一起.
彼此相教.
互相的代禱.
提醒我們大家一起.
怎樣做一個好的門徒.
所以服侍的時間.
我希望能夠多加這些元素在裡面.
當我們這樣做的時候.
我們教會以後就可以繼續被建立.
教會的事工可以被傳承.
也可以完成神交付給我們教會的使命.
最近我跟不同的同工有分享提醒他們.
因為真的有很多弟兄姊妹願意來服侍.
特別是Mill Peters的弟兄姊妹.
成年已經做過弟兄姊妹.
他們不單單自己的堂會已經有服侍.
當其他堂會有需要.
我真的很感動.
他們take time make effort.
用時間去不同的堂會.
去幫助其他堂會的工作.
所以我很感恩.
不過我提醒同工要小心.
不要養弟兄姊妹.
burn out.
因為有時候我去不同堂會的時候.
為何你又在那裡 去到那裡又在那裡.
好像無所不在.
那些積極熱身服侍的弟兄姊妹.
我說你牧師要小心.
不要讓弟兄姊妹覺得是burn out.
是耗盡.
忙是自然的 但要忙得開心.
要忙得有能力.

$^{601}$要多點關心他們身體心靈的需要.
多點為他們祈禱.
而且讓他們知道.
我們教務同工是關心他們的.
這是我老牧師有時候會提醒年輕的傳道人.
事實上.
前兩天.
聖靈感動說I don't know.
我忽然間想起一個姐妹.
她正在接受化療.
已經有一段時間沒有回教會.
我很久沒見她了.
我知道她受化療辛苦.
所以我都不想去探望她.
怕我們有什麼傳言給她.
我就寫了一個很短的短信去問候她.
原來才知道.
我師母一早已經做了這件事.
我第二次做這件事.
她回覆我的短信令我很感動.
她說牧師我都很掛念你們.
我知道你們從香港回來.
師母不是很舒服.
她說了很多話.
說了她的近況給我聽.
但最令我感動的是什麼呢.
最後在結尾的時候.
她給我的祝福.
她說May our heavenly father.
continue to shower his blessings onto you.
and our church.
Take care.
願我們的天父繼續將祂的祝福.
傾置在你 師母和我們教會身上.
請保重.
你不是我 所以你不知道我的feeling.
我馬上告訴你我的feeling.
我心裡很感動.
I am being touched.
我以為我是關心這個姐妹.

$^{641}$誰知這個是患病在病床上掙扎的姐妹.
仍然沒有忘記關心我和師母.
還有 她的教會.
大家聽到嗎.
教會和外面機構不同的地方就在這裡.
原來我們彼此牧養.
我以為我牧養這個姐妹.
她牧養我回頭.
給我有鼓勵 給我有鼓舞.
聽到嗎.
到底是甚麼推動我們能夠這樣做.
我相信是救主耶穌向你們所展示的.
那種無私 忘我.
還有 捨己的愛.
我們來看看人.
其實我相信這是現實.
一出生 掛包落地.
我們受罪性的影響.
都會有自私的傾向.
大家同意嗎.
我們會先從自己的利益著想.
對自己沒有益處的事.
或者我們不太熟悉的事.
或者對我們沒有好處的人.
通常我們不會理會.
這也是人之常情 人的傾向.
其實耶穌是怎樣的呢.
上星期我跟耶穌安息是治病.
明知那些律法師 法律會攻擊他.
耶穌又知道那些人說.
你醫病趕鬼.
你是有鬼王的能力去趕鬼.
那就誣告耶穌基督.
耶穌因為這樣沒有去趕鬼嗎.
沒有.
安息是治病.
雖然人家這樣攻擊他 誣告他.
他繼續去做神 要他去趕鬼.
而當猶太人去煽動.
以至彼拉多在群眾的壓力下.

$^{681}$讓他們釘耶穌十字架.
耶穌在十字架上怎樣說.
我們說耶穌在十字架上說了七句話.
不過我想提特別其中一句.
耶穌在十字架上看著那些大喊高呼.
要釘他十字架的群眾.
然後耶穌舉目望天.
向父神祈禱說.
說什麼.
父啊 赦免他們.
因為他們所作的 他們不曉得.
為什麼基督教能夠成為.
改變世界文化最厲害.
影響最深遠 最廣的一個宗教.
就是因為耶穌基督.
用自己的生命所表達.
這個偉大無私無條件的愛.
這個愛是來自上帝.
因為愛是從神來的.
耶穌怎樣說 馬達福音5:44.
只是我告訴你們.
你要愛你們的仇敵.
為那逼迫你們的禱告.
哇 耶穌說了很容易.
但他現在就是為那些逼迫他人禱告.
弟兄姊妹 你說我不是耶穌.
耶穌可以這樣做.
神的兒子 我不是耶穌.
這個教導 人怎能做得到.
誰能做得到.
你同意嗎 你可以做到嗎.
耶穌在世上的時候.
曾經說他拿了人的形象.
成為人的樣式.
意思是什麼.
耶穌在世上那三十多年的時間.
是和你和我一樣.
擁有人的肉體.
也擁有人的意識和感情.
耶穌有沒有喜怒哀樂的表現.

$^{721}$Yes, of course.
福音書裡面記載了很多.
耶穌一些情緒的反應.
耶穌都會動怒的.
耶穌會流眼淚的.
所以耶穌其實是明白你的想法.
明白我的感覺.
大家對這個事情的那種觀點.
但是 耶穌在十字架上.
當他看著下面那群人.
丟耶穌十字架的人的時候.
他竟然可以向上帝祈禱.
祈求神的赦免他們.
這告訴我們什麼呢.
原來愛仇敵愛.
被別人禱告這件事.
是可行的.
是做到的.
因為耶穌在那段時間.
都是人來的.
你說是的 但是他是神人異性.
那些老基督徒可能聽過一個人的名字.
叫做Koritenbun.
有沒有聽過.
Koritenbun.
這個Koritenbun是第二次世界大戰的時候.
在納粹集中營當中其中一個倖存者.
1947年的時候.
他有一次跟一個曾經在集中營看守過他的守衛見面.
這個守衛見到他的時候.
就請求Koritenbun.
描述他在集中營對他所做的一切.
在那個時間.
Koritenbun自己說.
他感到他自己的心.
好像冰一樣的冷.
當那個人說你描述我吧.
I am sorry for what I did to you.
就在集中營裡面.
Koritenbun說我的心好像冰一樣的冷.

$^{761}$他說我意識到在感情上面來說.
對這個前納粹的守衛.
我完全沒有愛.
也沒有寬恕.
於是他做了一件事.
於是我心裡面向神默默祈禱.
祈禱什麼呢.
他說主耶穌,我沒有辦法原諒這個人.
但是請你將你的寬恕.
是你的寬恕.
耶穌的寬恕賜給我.
Koritenbun在以後跟我們分享.
他說當我有意識地決定跟這個人握手.
揮出我的手.
跟他握手的時候.
他感到有一種溫暖的熱流.
從他的肩膀開始一直到整個手臂.
是一種超自然的熱愛接管了他的心.
讓他能夠說出那句著名的說話.
是什麼.
我原諒你弟兄.
我全心全意地原諒你.
世俗的愛.
通常是依賴環境.
和有沒有化學反應.
或者我們大家之間相互有沒有好感.
然後就導致我們有愛的行動.
這是世俗的愛.
世界裡面的愛.
但是多少個愛呢.
即使完全沒有愛那種感覺.
但是我們仍然可以選擇是被選擇去採取行動.
Koritenbun他的經驗告訴我們.
基督無私無條件的愛是來自上帝的.
從神那裡而來.
並且是聖靈賜予的.
對他來說他本來不行的.
沒辦法.
他的心很冷.
他覺得不可以饒恕這個人.

$^{801}$但是當他向神禱告.
求神將耶穌的寬恕賜給他的時候.
忽然之間他成功了.
他從神那裡領受了力量.
這種愛是不在乎接受愛的那個人是誰.
或者這個人有沒有值得我們愛的地方.
No, it doesn't matter.
那些不重要了.
原來我們這種愛.
如果是來自上帝的話.
是在乎上帝是誰.
我們認識的上帝是誰.
Corey有一個姐姐叫Bessie.
有少許人知道.
因為她在集中營死了.
捱不過那個時間.
但是她有很多說話值得我們反思.
今天只是其中一句跟你們分享.
她臨死前幾天.
她跟她妹妹說.
There is no pit so deep.
that God's love is not deepest still.
什麼意思.
她說沒有一個深淵是上帝的愛無法觸及的.
這是一個很深奧的說法.
沒有一個最深的深淵是上帝的愛無法去到的.
上帝的愛就是這麼偉大.
所以詩人是怎樣形容上帝的愛的.
詩人形容上帝的愛是長闊高深.
什麼叫做長闊高深.
無可測量.
耶穌為所有世上的罪人而犧牲.
不是為好人而犧牲.
上帝為世上所有的罪人犧牲.
很明顯不是因為人可愛.
No.
耶穌的行動是顯示神那種偉大的愛.
帶給這個破碎世界.
有醫治和有盼望.
所以我們可以清楚看到.

$^{841}$神藉著耶穌所顯示的愛是有意識.
有意志去驅動行為.
而非是一種轉迅即逝的情感.
要不要我再說一次.
這個可能大家也有這個觀念.
不過我想強調的.
聖經所啟示耶穌的愛是先行後感.
記住這四個字.
先行 先做出來.
然後你才有感覺.
意思是我們不要等到有感覺.
要去愛這個人才去做一些事.
這是世俗的愛.
我們先要順從聖經的教導.
神的吩咐去做.
然後那些感覺才會來.
這就是Koritambun的經驗.
親愛的弟子們.
不知道我們有沒有哪個人.
或者有沒有哪些人.
我們覺得很難去愛.
耶穌是怎樣跟我們說的呢.
伯太五章四十六節.
「假如你們只愛那些愛你們的人.
上帝又何必獎賞你們呢?.
就是連瑞麗,瑞君(罪人)都會這樣做」.
親愛的弟子們.
我們來教會這麼久.
學聖經學了這麼多.
我們說參加查經班.
我們也知道應該去饒恕那些得罪我們的人.
我們也知道應該去饒恕那些傷害我們的人.
但很多時候我們做不到.
為什麼我們做不到呢?.
其實彼得也很難做.
彼得也問過耶穌一句話.
彼得進前來耶穌說.
主啊,我弟兄得罪我了.
我應該饒恕他多少次?.
七次夠不夠?.

$^{881}$耶穌的回答是什麼?.
不是七次.
乃是七十個七次.
這個難不難呢?.
彼得問耶穌的問題.
當然覺得耶穌要求敵這件事有難度,怎麼做得到?.
他還有一個背景.
為什麼?.
因為猶太人的老師拉比教那些猶太人說.
那些人得罪你,你要饒恕他.
饒恕多少次?.
三次就夠了.
第四次呢?.
Forget about it.
這是猶太人拉比老師的教導.
彼得提出七次.
他覺得自己很generous.
彼得覺得自己已經很慷慨.
誰知耶穌的回應說.
No,七十個七次.
即是無限量.
無限量對他饒恕.
我告訴你,我的時間不夠.
首先我們要明白.
什麼是饒恕?.
在聖經裡簡單來說.
饒恕是不懷恨.
對得罪你的人不要懷恨.
放下你要報復的權利.
在舊約聖經裡猶太人的老師教他們.
以牙還牙,以眼還眼.
所以現在看以色列打仗.
你真的不要碰他.
你踢他一腳,他踢你兩腳.
不會讓你欺負他.
以牙還牙,以眼還眼.
這是他們的culture.
但饒恕就是.
不要懷恨,放下報復的權利.
你可以這樣做.

$^{921}$舊約聖經裡這樣教.
以牙還牙,以眼還眼.
但No, don't do that.
耶穌教導門徒.
更加要愛仇敵.
愛仇敵的意思.
就是不要讓仇敵的過錯,罪惡.
繼續在你的心裡折磨你.
要不要我再說一遍.
饒恕的意思,愛仇敵的意思.
不要讓仇敵所做的事.
所犯的錯誤,甚至是罪惡.
繼續在你的心裡折磨你.
如果你的仇敵不悔改.
你就不要說他.
這是一般人的反應.
如果是這樣的話.
我們的情緒主導權.
就永遠拿在那些人手裡.
你聽到我的意思嗎.
我們的情緒主導權.
就是被這班人拿著.
聖經告訴我們.
天父見日頭照好人.
也照打人.
記得嗎.
好人,壞人.
日頭都照給他們.
大家都需要陽光.
意思就是說.
就算人在罪裡面不知悔改.
天父上帝仍然向他們施展出.
普遍性的慈愛和寬容.
所以神當然會因為人故意犯罪.
生氣.
不過在施行審判和懲罰的同時.
我們會看到神有恩典.
所以我們沒有必要被那些人的罪.
捆綁和折磨我們的內心.
令我們不開心.

$^{961}$甚至令我們有苦讀.
這是第一點.
第二點我們要明白.
饒恕和修復關係是兩回事.
饒恕一個人不要懷恨.
不要對付他.
不要報復.
修復關係.
我們可以單方面饒恕一個人.
這是我們心裡可以做到的.
無論對方悔改或不悔改.
我們都可以選擇不懷恨.
選擇不去報復.
但這樣不代表我們立刻要信任這個人.
不代表我們要恢復和他原先的關係.
如果要修復關係.
對方必須要悔改.
而且要展現出悔改相應的果子.
聖經有一個很明顯的例子.
耶穌時代的ILS主管.
叫薩戈.
他信耶穌之後.
他怎樣跟耶穌說.
他主啊.
我把我所有的財產一半分給窮人.
還有我訛責過誰.
我還他四倍.
這是悔改的果子.
他真的悔改.
否則誰會這樣做.
分一半財產去周濟窮人.
我騙了誰賠他四倍.
所以聖經的描述.
是更多說到我們跟神之間的關係.
就好像我們如果只相信耶穌是神的兒子.
很多人都相信的.
雖然說得不錯.
但沒有認罪悔改.
上帝會憐憫我們人的軟弱.
我們不悔改因為人的軟弱.

$^{1001}$但我們不可以不能夠跟神重建關係.
成為神的兒女.
但以後當我們認識神更深.
體會到祂對罪人的愛.
包括祂賜給人那種恩典.
於是我們知罪而且悔改.
不同了.
聖靈就會進入我們的生命當中.
而且會賜給我們一個新的生命.
讓我們能夠像耶穌一樣.
去饒恕 去放下仇恨.
所以主要的弟兄姊妹.
為了保護我們的心.
不要繼續被仇恨去腐蝕.
讓我們求我們的天父爸爸幫我們.
加添我們的力量.
讓我們能夠以謙卑 饒恕的心去對待人.
讓我們的生命都像耶穌一樣.
能夠顯示出那種無私.
忘我無條件的愛.
這種愛就是我們能夠改變世界的動力.
耶穌在世界上.
真的出來做工只有三年的時間左右.
但竟然改變了整個世界.
所以我希望今天大家能夠看到.
耶穌沒有用祂做兒子的權柄能力.
沒有用軍事強權去征服世界.
耶穌是用祂謙卑 服侍人的心.
和無私的大愛.
去改變了我們的心和生活的態度.
所以弟兄姊妹.
可能我們沒有想過自己可以改變世界.
因為我們連改變自己 改變家人都不行.
但請大家不要忘記.
我們信耶穌之後.
我們就擁有耶穌基督的愛.
這是聖靈的果子.
如果你沒有 你還未信主.
除了聖靈的果子.
還有其他屬靈的德性.

$^{1041}$讓我們能夠做神的僕人.
當我們願意謙卑 信服耶穌的教導.
我們就能夠發揮愛的威力.
讓我們改變自己 影響他人.
所以弟兄姊妹.
至今 上帝沒有觸動你的心.
你聽明白 你知道了.
有沒有想過接下來的一個星期.
你可以幫助哪一個人.
你可以邀請哪一個人.
來參加我們的主教崇拜.
有沒有想過有哪一個人.
一路你遙述不到.
但今天可不可以好像.
「耶穌 我不行了」.
你把你的遙述賜給我.
我們一起祈禱.
親愛的天主上帝.
不容易的功課.
但這是聖經的教導.
賜給我們一些能力.
去做你要我們去做的事.
因為我們活著.
就是為了你的緣故.
讓我們能夠以天父上帝的事為念.
奉耶穌的名字祈禱.
阿們.
\newpage



\section{傳道書 3:16-4:3}
\label{sec:031888}
\textbf{公義缺席}~~~~譚文倩莊傳道
\newline
\newline
連結: \href{http://www.church.com.hk/acms/content.asp?site=cdc\&op=show\&type=product\&code=031888\&layout=sermon}{\texttt{031888}} ~~~~ 語音日期: 2026-03-29
\newline
\newline
\hyperref[sec:031887]{\small{< < < PREV SERMON < < <}}
~
\hyperref[sec:index_chronic]{\small{[返順時目]}}
~
\hyperref[sec:031889]{\small{> > > NEXT SERMON > > >}}
\newline
\newline
傳道書 3:16-4:3
\newline
\begin{longtable}{cl}
\hline
\hline
章節 & 經文 (和合本修訂版)\\
\hline
3:16 & \begin{tabularx}{0.7\textwidth}{X} 我又見日光之下，應有公平之處有奸惡，應有公義之處也有奸惡。 \end{tabularx} \\ \\ \relax
3:17 & \begin{tabularx}{0.7\textwidth}{X} 我心裡說：「神必審判義人和惡人，因為在那裡，各樣事務，一切工作，都有定時。」 \end{tabularx} \\ \\ \relax
3:18 & \begin{tabularx}{0.7\textwidth}{X} 我心裡說：「為世人的緣故，神考驗他們，讓他們看見自己不過像走獸一樣。」 \end{tabularx} \\ \\ \relax
3:19 & \begin{tabularx}{0.7\textwidth}{X} 因為世人遭遇的，走獸也遭遇，所遭遇的都一樣：這個怎樣死，那個也怎樣死，他們都有一樣的氣息。人不能強於走獸，全是虛空； \end{tabularx} \\ \\ \relax
3:20 & \begin{tabularx}{0.7\textwidth}{X} 都歸一處，都是出於塵土，也都歸於塵土。 \end{tabularx} \\ \\ \relax
3:21 & \begin{tabularx}{0.7\textwidth}{X} 誰知道人的氣息是往上升，走獸的氣息是下入地呢？ \end{tabularx} \\ \\ \relax
3:22 & \begin{tabularx}{0.7\textwidth}{X} 總而言之，人能夠在他經營的事上喜樂，是最好不過了，因為這是他應得的報償。他身後的事誰能領他回來看呢？ \end{tabularx} \\ \\ \relax
4:1 & \begin{tabularx}{0.7\textwidth}{X} 我轉而觀看日光之下所發生的一切欺壓之事。看哪，受欺壓的流淚，無人安慰；欺壓他們的有權勢，也無人安慰。 \end{tabularx} \\ \\ \relax
4:2 & \begin{tabularx}{0.7\textwidth}{X} 因此，我讚歎那已死的死人，勝過那還活著的活人。 \end{tabularx} \\ \\ \relax
4:3 & \begin{tabularx}{0.7\textwidth}{X} 但那尚未出生，就是未曾見過日光之下所發生之惡事的，比這兩種人更幸福。 \end{tabularx} \\ \\
[1ex]
\hline
\hline
\end{longtable}
$^{1}$虛空 勞碌益處.
智慧有善 逆光之下.
最後都是要指向傳道書的核心.
是講敬畏神.
敬畏神就好像我們基督徒做人的一個大原則.
幫我們在很快消散.
好像煙霧的人生裡面.
能夠選擇一條有智慧 有方向的道路.
記不記得上一次.
應該都差不多兩個月前.
在第三章裡面有一首詩.
那裡描述人生裡面一些既定時間會發生的一些狀況.
然後經文裡面亦都有提到.
神由開始去到終結所做的事.
我們人是沒有可能完全明白的.
神這樣做是要我們人在祂面前有那個敬畏的心.
今天來到下一個段落.
三章十六至四章三節.
一開始傳道者就說.
他留意到秉行公義和伸張正義的地方.
都有邪惡.
有沒有看過《Wicked》這套戲.
這套戲裡面的故事圍繞著兩個女巫.
一個叫Elphaba.
你見到尖頂帽子的那個.
她一出生的時候.
她全身的皮膚都是綠色的.
她是一個很聰明.
亦都很有正義感的一個人.
但因為她不願意同流合污的關係.
所以她得罪很多人.
不受歡迎.
另一個女巫就叫Galinda.
或者Glinda.
如果你看過這套戲.
你就知道我為何特意這樣說.
很美麗.
亦都是一個世界女.
她被人捧到上天花板.
在大家的心目當中.

$^{41}$她就是一個好的女巫的象徵.
故事圍繞著他們兩個之間的友情.
還有一些執政掌權的人.
如何去製造恐慌.
和利用權力去扭曲一些事實.
我們先來一起看看今天的經文.
傳道書三章十六四至三節.
我又見日光之下.
應有公平之處有奸惡.
應有公義之處也有奸惡.
我心裡說.
神必審判義人和惡人.
因為在哪裡各樣事務.
一切工作都有定時.
我心裡說.
為世人的緣故神考驗他們.
讓他們看見自己不過像走獸一樣.
因為世人遭遇的.
走獸也遭遇.
所遭遇的都一樣.
這個怎樣死.
那個也怎樣死.
他們都有一樣的氣息.
人不能強於走獸.
全是虛空.
到貴日處.
都是出於塵土.
也都歸於塵土.
誰知道人的氣息是往上升.
走獸的氣息是下入地.
總而言之.
人能夠在他經營的事上喜樂.
是最好不過了.
因為這是他應得的報償.
他身後的事.
誰能領他回來看呢.
所發生的一切欺壓之事.
看哪.
受欺壓的流淚.
無人安慰.

$^{81}$欺壓他們的有權勢.
也無人安慰.
因此.
我讚嘆那已死的死人.
勝過那還活著的活人.
那尚未出生.
就算未曾見過日光之下.
所發生之惡事的.
比這兩種人更幸福.
若果說神在我們每一個人身上.
都有一個美好的計劃.
傳道者在此有些懷疑.
因為他至少看到一個問題.
就是這個世界很明顯地.
存在不公義.
邪惡和痛苦.
傳道者說.
他看到應該最公道的地方.
都沒有公道.
新疆正義的地方.
都變了質.
這裡所說的不公義.
其實是沒有被具體界定.
所以我們不可以單單只在想.
是在說法庭上的那種公義.
因為事實上.
不公義.
是可以在我們的社會結構裡.
去存在的.
或者是在一些應該負責.
維持公義的人.
的身上裡出現的.
很多時我們失望.
不是因為這個世界有壞人.
而是因為我們對一些有權.
有能力的人.
或者制度有期望.
事實上你和我都不可以否認.
人的心一旦牽涉權力.
利益的時候.

$^{121}$就很容易被扭曲.
傳道者似乎一早就看到.
問題不一定在政策上.
而是坐在那些執政掌權的人.
那個心到底是怎樣.
To be fair.
其實我們人.
其實都是很多時會嘗試去.
秉行這個公義.
和做一些正直的事.
只不過結果有時.
是很強差人意.
或者那些準則不是很一致.
結果可以是帶來一些很荒謬的結果.
例如不知道你知不知道.
我們美國在某些州份裡面規定.
若果一個未成年的少女.
要打耳窿的時候.
要有父母的consent.
要有父母的同意.
要簽名的.
但是在同一個州份裡面.
若果一個未成年的少女.
要去墮胎的時候.
不需要父母知道.
父母不需要簽甚麼.
聽下去是不是很荒謬.
是不是會很令人質疑.
究竟我們所定的一些法律.
是不是真的在嘗試保護我們一些小朋友.
因為聽下去完全不make sense.
正如經文那裡說.
應有公平之處有奸惡.
應有公義之處也有奸惡.
這裡的奸惡在希伯來文的原文.
那個意思有歪了 彎了的意思.
即是本來是正的 是直的.
但是當人自己去定那個標準的時候.
就會歪了 彎了.
即是不一致.

$^{161}$甚至可以互相矛盾.
這類的事件.
就正正反映了一個問題.
就是就算我們人最精英的制度.
法庭 法律.
去嘗試實現公義的時候.
結果都可以出現一個這麼大的落差.
當然不是一個單一的事件.
而是一個模式.
因為我們人會自己以為.
我們能夠去建立一個完全公義.
完全正直的系統.
但是實際上都是歪了.
都是彎了.
不完全.
或者公義和正直是混到亂七八糟.
這個正正是要呼應.
傳道者在這裡所說的觀察.
問題不是在制度裡面.
而是我們人本身.
其實是當人嘗試去做神.
去定義那個對錯的時候.
最終出來的公義和正直.
就會變得不是很穩定.
甚至變形.
其實歪了.
彎了的心.
用什麼去將審判之處.
變成監獄之處.
或者你會說.
換另一個人或者另一批人上去.
就OK的了.
但是歷史告訴我們.
換另一個人.
另一批人上去.
只是歷史重演.
不知道是不是傳道者.
自己都站在這個權力之上.
反而他更加清楚看到.
這個制度的腐敗.

$^{201}$和人性的醜陋.
那個扭曲.
就算你是最有話事權的那個.
你都改變不了人心那個扭曲.
三章十七節.
我心裡說.
上帝必審判義人和惡人.
因為在那裡.
各樣的事務.
一切工作.
都有定時.
可能我們會問.
那就是什麼時候?.
如果真的會有審判的時候.
為什麼不馬上現在就去做呢?.
為什麼有些壞人在做一些壞事.
我們看上去他們覺得似乎很好.
They are OK.
為什麼反過來有些人在走一條正路.
在做一些好事.
他們卻要受苦呢?.
我認識一個熱愛工作的弟兄.
他在公司工作很鞠躬盡瘁.
早到遲走.
有一天來了一個年輕人.
老闆就和這個弟兄說.
從今天開始.
你和這個年輕人是天媒.
你們平起平坐.
一起去分擔計劃的工作.
沒多久弟兄留意到.
這個年輕人好像有些特別的待遇.
搶功勞.
卸責.
最擅長擦鞋.
弟兄雖然知道.
但他沒眼睛看.
我繼續做好自己的工作.
繼續埋頭苦幹.
誰知這個受寵的年輕人.

$^{241}$持著老闆看著他.
他便變本加厲.
有一天走到這個弟兄的桌面上說.
我不喜歡有什麼人來找你.
You have to pay me respect.
一個懂得擦鞋.
一個只懂得喊.
越來越多不公平的事情發生.
弟兄決定不行.
我要找老闆談談究竟發生什麼事.
老闆說沒有.
沒有這樣的事.
但情況是越來越差.
如果你是這個弟兄.
你會怎麼做.
你又猜後來會怎樣.
沒多久老闆發現這個年輕人.
其實真是大隻狗.
不單止吞撲.
實際上很多事情他不懂得做.
原來大部份的事情.
是這個不聲不響的弟兄在背後做了.
所以他立刻解僱了這個年輕人.
跟著跟這個弟兄道歉.
為表示誠意的關係.
他就加回這個弟兄的人工.
還給他一份bonus.
你就想.
其實熬了一輪之後.
這個弟兄終於受不了.
他還沒找到新的工作.
他就離開了這一份.
他工作了十年的工作.
聽到這些故事的時候.
我們都很想那個壞人.
這個年輕人被懲戒.
有些事情會不好的發生在他身上.
有權的人去秉行公義.
坐強扶弱.
幫助這個弟兄報仇雪恨.

$^{281}$打到這個年輕人落花流水.
跪在這裡say sorry.
respect.
我們是想這樣的.
但為什麼不是這樣.
壞人很順利.
好那個在這裡要受委屈.
我們基督徒會想.
神啊為什麼你還沒出手.
為什麼你不是立刻去秉行公義.
為什麼不是立刻可以有那個後果.
因為對我們來說.
真正的公義應該是要即時發生的.
要即時有懲罰的.
我們做父母的.
沒理由小朋友做錯了一些事情.
然後你沒事.
一個月之後跟他說.
然後罰他一個月之前的事情.
是要立刻處理的.
神為什麼你不是這樣做.
如果沒有報應.
或者還要我等.
那就是沒有公義.
但在這裡.
正正在這裡.
有一個問題我們要反思一下.
當我們這樣說.
或者當我們這樣想的時候.
其實我們是不是在教神怎樣做事.
我們是不是要求神.
在我們覺得對的時候去做事.
聖經在彼得後書三章九節那裡說得很清楚.
他說神沒有拖延.
而是祂在忍耐.
主沒有拖延他的應許.
就如有人以為他要辭延.
其實他是寬容你們.
不願一人沉淪.
而是人人都來悔改.

$^{321}$意思是神.
如果神是要一早.
或者立刻審判的時候.
你和我.
你和我其實有沒有機會悔改.
你和我有沒有機會悔轉呢.
神之所以延遲審判.
不是因為他不公義.
而是因為我們所信的神.
祂是有憐憫.
有忍耐.
祂選擇等候.
願意更加多人.
能夠得著祂.
願意更加多人得著這個救恩.
所以傳道者說.
都有定時的意思是.
審判是一定會來的.
但不是你和我的時間.
是神的時間.
有趣的是.
這件事會繼續去困擾我們.
因為我們仍然在fight for control.
我們仍然fight for主導權.
傳道書這裡提到.
神容許世界出現混亂.
不公.
其實是要試驗人.
要讓人看到.
我們人很多時候覺得.
我們自己就是神.
我們自己去定規則.
我們自己決定對與錯.
甚至覺得我們自己可以管理這個世界.
做得和神一樣好.
甚至更好.
但是我們人.
其實不適合做神.
為什麼這個世界有這麼多不公平.
正正是因為我們人在扮神.

$^{361}$我們人拒絕創造主.
我們覺得自己可以取代祂.
但最後產生出來的東西.
就是一些歪了的東西.
任何肯認真思考的人.
其實都會看到.
現在執政掌權的人所做的事.
有時不對勁.
有時真的不正直.
某程度上.
其實這些都是神憤怒.
顯明了神的憤怒.
祂沒有即時審判.
而是讓人去承受去扮神帶出來的後果.
你看看周圍.
你就會發現.
當我們人去自己說事.
自己做神的話.
結果就是得出我們現在所見到的東西.
做得一點都不好.
想一想.
其實我們是否應該讓神去做神呢?.
而我們人需要調教一下我們自己.
承認我們不是神.
接著經文突然轉了一個方向.
傳道者說人和動物一樣.
都有氣息最終都會歸於塵土.
想一想.
有血有肉有呼吸.
食和餓病和死.
出於塵土歸於塵土.
想一想.
好像的而且確我們和動物好像沒什麼分別.
但當然是有分別.
三章十一節.
回到之前那裡.
它這樣說.
「又把永恆放在人的心裡」.
意思是唯獨我們人裡面.
有一種對永恆.

$^{401}$神的計劃.
和超越今生意義的渴求.
不是一個感覺.
而是你整個人的Being.
包括我們的思想.
意志.
品格.
和我們的情感.
這個渴求.
叫我們相信.
我們的人生.
不是只有我們在日光之下.
所見到的那些東西那麼簡單.
三章二十二節.
「總而言之.
人能夠在他經營的事上喜樂.
是最好不過了.
因為這是他應得的報償」.
就是說.
你不要再試圖控制所有的事情了.
靠自己去解釋.
判斷.
不要再教神怎樣做.
何時做.
日光之下的視覺.
我們所看到的東西.
所明白的東西.
其實是很有限的.
我們未能夠完全掌握神的啟示.
我們看到.
但是好像模模糊糊的.
被我們自己的有限遮住.
就好像在芝加哥.
你冬天開車一進去的時候.
那些霞氣.
遮住了窗戶.
你看到一點.
但是不看得清楚.
不看得完.
四章一節.

$^{441}$日光之下所發生的一切欺壓之事.
看那受欺壓的流淚無人安慰.
欺壓他們的有權勢也無人安慰.
因此我讚嘆那已死的死人.
勝過那還活著的活人.
但那尚未出生.
就是未曾見過.
日光之下所發生之惡事的.
比這兩種人更幸福」.
四章一節.
這一節裡面.
「欺壓」這個字出現了三次.
在希伯來語的字根裡面.
欺壓之事和受欺壓的是被動的.
而欺壓的人是一個主動的行為.
兩個被動一個主動.
1993年有一套電影叫做《Joy Luck Club》.
其中一幕講述一對夫婦的交易.
講述這對夫婦.
老公Harold和老婆Lina都在工作.
他們沒有小朋友.
為了公平起見.
家裡所有的開支都50/50.
房租,水電,買菜,旅行.
他們在冰箱門口設了清單.
花了甚麼錢就寫在清單上.
花了多少錢.
每個月底的時候就結帳.
確保雙方沒有人吃虧.
聽上去合理嗎.
兩個都在工作.
但其實背後有一個潛藏的問題.
這個隱憂最後因為一桶冰淇淋的關係.
爆發了出來.
老公Harold很喜歡吃冰淇淋.
經常買,經常吃.
而Lina就沒有碰過這桶冰淇淋.
不是因為他白胖.
而是冰淇淋會令他回想起.
他童年的一個很差的遭遇.

$^{481}$有些童年陰影.
但去到月底的時候.
Lina都要負擔冰淇淋的一半費用.
有一天Lina的媽媽就來探望他們.
她經過冰箱見到這張清單.
她就很不客氣地跟這對夫婦說.
她就說這個Harold根本不疼他的老婆.
因為他不單要他老婆負擔一半的冰淇淋費用.
他連他老婆怕吃冰淇淋都不知道.
有些不公義.
表面不激進.
沒有欺壓,沒有暴力.
甚至很斯文,很理性,很有制度.
這種不公義.
其實都很貼近傳道書這裡的訊息.
當公平成為了婚姻的保護機制.
兩個人只是在計算.
確保自己沒有吃虧.
慢慢這個冰淇淋已經不再是冰淇淋了.
因為冰淇淋象徵著傷害.
這種不公義.
其實可能我們經常都會遇到.
問題是你是主動的那個.
還是你是被動的那個.
這個電影情節.
讓我們看到人所謂的公平,公正.
在執行出來的時候.
都可以是歪了,彎了.
不單止不完全,不公義.
亦都帶著很多傷害.
求主幫助我們.
我們不只是看到不公義.
卻是看不到我們自己對人的不義.
或者我們都需要反思一下我們自己.
表面上的公平,公正.
是不是聖經所說的公義.
當然更加值得我們思考的是那個核心的問題.
你會不會不覺,不覺,不留意.
就將自己變成了你生命裡掌管的那個神.
自己去定規則.

$^{521}$自己去決定對錯.
甚至覺得我自己已經可以管理一切.
做得和神一樣那麼好.
甚至比祂做得更加好.
求主的教導.
聖靈的工作.
在我們每一個人心裡面去動工.
我們一起低頭祈禱.
你寶貝的話語.
再一次提醒我們.
其實真正可以秉行公義的只有你一個.
因為你是那個公義的主.
在你裡面公義沒有缺席.
因為真正的公義已經在十字架上去成就.
奉主耶穌基督聖名祈求.
Amen.
\newpage



\section{出埃及記 3:7}
\label{sec:031889}
\textbf{聖經人聖經事:我原知道他們的痛苦}~~~~黃克勤牧師博士
\newline
\newline
連結: \href{http://www.church.com.hk/acms/content.asp?site=cdc\&op=show\&type=product\&code=031889\&layout=sermon}{\texttt{031889}} ~~~~ 語音日期: 2026-03-28
\newline
\newline
\hyperref[sec:031888]{\small{< < < PREV SERMON < < <}}
~
\hyperref[sec:index_chronic]{\small{[返順時目]}}
~
\hyperref[sec:031890]{\small{> > > NEXT SERMON > > >}}
\newline
\newline
出埃及記 3:7
\newline
\begin{longtable}{cl}
\hline
\hline
章節 & 經文 (和合本修訂版)\\
\hline
3:7 & \begin{tabularx}{0.7\textwidth}{X} 耶和華說：「我確實看見了我百姓在埃及所受的困苦，我也聽見了他們因受監工苦待所發的哀聲；我確實知道他們的痛苦。 \end{tabularx} \\ \\
[1ex]
\hline
\hline
\end{longtable}
$^{1}$各位電影節目大家好.
很快又一天了.
又是讀聖經的時候.
生命樹宣教網絡.
為大家預備了《聖經人》和《聖經事》.
希望我們在《聖經》的人物和故事後面.
不單止認識這個人和事情.
而是認識到在背後.
是一隻手在工作的上帝.
他是怎樣的.
認識他更加多一點.
今次這個題目是特別貼切的.
仍然是摩西的故事.
不過是特別貼切的.
就是我們的主.
我們的上帝是說.
「我原知道他們的痛苦」.
這個「我原知道」.
有一點就這樣看中文.
就是我一早就知道了.
我根本就知道了.
其實是不是這個意思呢?.
又有一點不是.
又好像是.
好,先看看《聖經》.
出一集的第三章第七節.
由於我有些解釋.
瓦開了一些位置.
今次又加大了一些字.
讓大家容易清楚一點看.
所以這個7A是表示上半節.
「耶和話說:我的百姓在埃及所受的困苦.
我實在看見了.
他們因受督公的合制所發的哀聲.
我也聽見了」.
先解釋一下.
由於是以前約瑟.
就是雅各的兒子約瑟.
在埃及的時候很厲害.
那個法老是把全個國家都交給他了.

$^{41}$但是其實過了一段時間之後.
就有新的不認識約瑟的皇上.
來做了埃及的法老.
其實是利落河很長的.
利落河有上下兩半的.
兩半是有不同的埃及人.
他們之間時時都爭的.
如果是北埃及.
北埃及即是說.
如果你看地圖.
它是比較出海.
就是地中海那邊是北一點.
如果上埃及反而是南面的.
但上游就沒有出海口.
利落河三角洲那些那麼富裕.
所以北邊那些人就是.
辛苦一點,硬一點,窮一點.
有機會就搶了下面的皇位來做.
他們一上來的時候.
就約瑟邊位.
那麼多外族人在那裡走來走去.
做奴隸的.
就發覺原來這些就是希伯來人.
我們以前的法老就是.
讓他們在那裡住和工作的.
做什麼呢?就是建幾座織貨的城.
於是這個新的法老說.
好,給我建兩座織貨城.
織貨城就是做貨倉擺糧食的城.
以前約瑟都做開的.
就是方年,豐年等等.
大家知道了.
這個新的皇是令到希伯來百姓很慘.
原因就是他想一下.
哇,他們超過百多萬人.
如果萬一我們有外國的人.
聯同他們一起.
外國攻擊.
又聯同他們在裡面造反.
我們不夠他們來的.

$^{81}$因為他們個個的奴隸大大隻.
又很力.
那不是了.
不如我們真的做種族滅絕的時候.
逢親希伯來人的女士生孩子.
如果是兒子.
就叫那個修生婆捏死他.
殺死.
那神就救了他們.
那個修生婆就不肯這樣做.
於是他們就存活了.
後來摩西都差點死.
都是不能夠存活的男孩子.
因此摩西就被他的家人.
找一個蒲草造的箱.
塗點油上去不漏水.
推他在河裡漂浮.
後來被法老的女兒撿到.
養了他做養子.
後來摩西就這樣長大了.
學了埃及的各樣的技術.
學會了.
但是因為一時衝動殺了人.
就跑去米殿曠野逃難了.
終於神就呼召他.
上幾集就說這些了.
「我的百姓在埃及受的困苦我實在看見了」.
困苦這個字.
就是困苦慘情.
即是看到他們很難很慘.
很難很慘.
為什麼呢?因為越來越逼他們.
以前一天要交多少塊磚的貨.
有草供應有沙石供應.
為什麼要草呢?.
因為有草燒乾了.
那些乾了的草就好像有鍵一樣.
拔掉那塊泥.
現在不供應草.
你自己找草.

$^{121}$然後一時之間自己找泥.
但是貨要交這麼多塊磚.
因此他們考古學家是發覺.
這幾座城建到上面的磚.
質地是差了的.
裡面的草是少了的.
因為要找草要找沙石.
找不到.
很慘,不夠的話.
就會被督工鞭他.
喂,你未交到貨.
明天你要交齊.
先打幾槌.
因此被欺壓得很慘.
神就說了.
「我實在看見了」.
他們因為受督工的合制所發的哀聲.
特別是說督工合制他們.
藍色字.
實在這個字是沒有的.
希伯來文是沒有的.
「我看見」.
「raa」就是我看到.
「我實在看見了」.
他們因為受督工的合制所發的哀聲.
這兩個字實在是沒有的.
但是如果希伯來文.
兩次重複「我」「raa」「raa」.
「看見」「看見」.
重複用了兩次「看見」.
「raa」就是看見.
「我」「raa」「raa」.
「我看到」.
所以中文就是「我實在看見了」.
其實神是說.
「我看到」「看到」.
你強調神看到的.
所以當你遇到很多困難.
很多壓力.
各樣東西壓制的時候很辛苦.

$^{161}$其實你要知道.
神是跟你說.
「我」「raa」.
「raa」「raa」.
「我看到」「我看到」.
神看到你那些困難的地方.
還有很有趣的.
祂說「我都聽到」.
阿媽是聽到了.
「我聽到他們的哀聲」.
在這個困難裡面.
哎呀很慘.
要捱到什麼時候呢.
慘了我的子女都要捱.
我上一代都已經捱死了.
「jaa」「jaa」.
「haam」叫「crying」.
這個「crying」不是「哭」.
是「嗌」.
就是辛苦.
哎呀很慘.
「foo」「haam」.
「我看見看見」.
那些督公對他們的壓制.
「我聽到了」.
「了」的.
「我聽見了他們的哀的聲音」.
就是痛苦叫的聲音.
原來雖然我們肉眼看不到神.
但是祂是聽到你的.
祂是知道你的.
祂是看到的.
第七節下半節就是我今天的主題了.
「我原知道他們的痛苦」.
「我原知道Key」.
「Key」是實在的.
「Surely, Key, 我是知道」.
這個字如果讀舊約成都會出的.
「Jaadaa」「Jaadaa」是認識.
認識神.

$^{201}$「Jaadaa」就是這個字了.
認識到你.
「Jaadaa」.
「我是認識知道察覺」.
「我一早就察覺到他們的痛苦」.
「痛苦」這個字.
很痛苦很有趣.
譬如我們今天遇到人聊天.
「這件事怎樣?很痛苦啊」.
「在公司上班很困難很痛苦」.
「在家裡很痛苦」.
「跟我跟誰相處很痛苦」.
這個痛苦多數是心靈的痛苦.
譬如「我原知道他們的痛苦」.
多數是內心痛苦.
剛剛說完他們的哀聲等等.
辛苦到叫.
內心痛苦.
但這個「maqab」.
應該重金在前面.
「maqab」的意思是指兩方面.
是身體與心靈的痛苦.
這個字用得真好.
因為你知道「痛」有很多個字.
譬如我們說很痛很赤.
「maqab」在希伯來文是指兩方面.
就是身體的痛苦與心靈的痛苦.
做奴隸的身體有沒有痛苦呢?.
當然有.
手粗了又弄傷了.
弄傷了也沒有東西可以塗.
又要繼續開工.
要雲和那些泥去執草.
執草一弄傷手就弄傷了等等.
「maqab」是身體有痛苦.
心靈有痛苦.
做奴隸必定是心靈有痛苦.
接著是第八節.
第三節會說.
「我下來要救他們脫離埃及門的手」.

$^{241}$「要令他們出了那地」.
「是到美好寬闊瀏覽與密之地」.
「就是迦蘭人,黑人,阿摩利人」.
「比利時人,希美人,耶布斯人之地」.
「我下來要救他們」.
這個「我下來」.
神說「我下來」.
這是一個很形象化的說法.
如果是「我下來」.
是可以說「沉」.
即是說在上面很好.
現在沉不下來.
或者是「我下來要救」.
就是解救.
「我下來要deliver」.
解救他們脫離埃及門的手.
然後你看到兩個字相對很有趣.
「我要令他們出」.
「出」是mean.
「mean」是從.
「我要令他們出了那地」.
然後「入」.
「入」是美好的地方.
「到」是L.
所以我要帶他們從那地.
入美好之地.
這個「出」是指從.
「到」是指入.
從那裡入美好之地.
那裡是怎樣的呢?.
美好,寬闊,流奶與密之地.
現在是迦南人,黑人,阿摩利人等人住的.
我要帶他們去那裡.
我要帶他們去一個離開埃及.
做奴隸.
越來越辛苦.
做到叫出來那麼慘的時候.
我要帶他們去一個美好的地方.
寬闊.
寬闊是指沒有那麼多壓制.

$^{281}$可以更加自由寬闊.
其實應該照知道迦南地平原不是很多.
是有些山的.
大部分都是山.
不是很多平原.
流奶與密之地.
流奶其實多數我們都是想奶.
密就是蜜糖.
我解釋一下.
如果去過以色列.
我都會告訴你.
流奶的意思是畜牧業.
一個適合畜牧業的地方.
特別適合山羊.
畜牧業.
蜜就不是蜜糖,不是Honey.
而是蜜棗的.
以色列的蜜棗真的多到不得了.
他們吃蜜棗.
而這個營養與熱量很高.
而且結果子超級多.
施洗日漢吃的是蝗蟲野蜜.
極有可能都不是蜜糖.
而是蜜棗的.
很甜.
即是我們這些糖尿不太妥那些.
不太能吃.
很甜.
熱量很高.
營養與蛋白質都很高.
是一個很好的食品.
如果是做運動的話.
會帶他們出去一個很好的地方.
今日和大家分享的是.
我很想大家能夠記得.
神說:我原知道他們的痛苦.
所以你的痛苦不是只有你一個人知道.
很多時候我們都認為.
痛的時候只有自己知道.
是真的.

$^{321}$因為痛這個感覺.
無論是身體或內心的痛.
有一個特點.
只有我自己知道.
你告訴我.
是不能說的.
不能言喻.
我無法告訴你.
所以有時候聽病人.
我聽到會友等等.
那裡痛,那裡痛.
撩下撩下的痛.
怎樣撩下撩下呢?.
他說撩下撩下.
他說得很真.
現在又撩了.
撩下撩下又痛.
什麼叫撩下呢?.
只有他知道.
醫生聽不聽得懂呢?.
只是猜測.
一會兒一會兒.
嘰嘰地痛.
肚子扭著地痛.
那個藥.
是不知道的.
但是神知道.
我實在是認識察覺他們身心的痛苦.
你身心的痛苦.
你那些Marco.
神都是知道的.
認識察覺他知道.
我原知道.
我就是知道.
我確實知道.
你將你的事情交託給主.
祂會照顧你的.
祂了解的.
我們一起祈禱吧.
多謝天父.

$^{361}$你雖然是高.
雖然很遠.
好像不在我們身邊.
肉眼看不到.
不過我們知道.
你原知道我們的痛苦.
你是確實.
這個原知都是確實.
是了解察覺到的.
不是說遠了就是看不到聽不到.
而是你知道的.
因此.
你是行動的主.
於是你說我要下來解救他們.
將他們從這個地方.
帶他們進去那個地方.
你就這樣去釋放了認識的人.
而你就叫摩西去做這件事.
你原知道他們的痛苦.
因為你聽到他們的聲音.
你的眼睛看到.
你用的字是.
我知道知道.
兩次用知道這個字.
來表示你真的知道與關心.
讚美我們的主.
你是這麼愛我們的.
讓我們更多認識你.
更多將自己的難處.
與身心的痛苦向主教導.
順天祈禱奉耶穌名求.
阿們.
\newpage



\section{約翰福音 11:17-32}
\label{sec:031890}
\textbf{俯伏主耶穌腳前}~~~~鄭子禧傳道
\newline
\newline
連結: \href{http://www.church.com.hk/acms/content.asp?site=cdc\&op=show\&type=product\&code=031890\&layout=sermon}{\texttt{031890}} ~~~~ 語音日期: 2026-03-29
\newline
\newline
\hyperref[sec:031889]{\small{< < < PREV SERMON < < <}}
~
\hyperref[sec:index_chronic]{\small{[返順時目]}}
~
\hyperref[sec:031891]{\small{> > > NEXT SERMON > > >}}
\newline
\newline
約翰福音 11:17-32
\newline
\begin{longtable}{cl}
\hline
\hline
章節 & 經文 (和合本修訂版)\\
\hline
11:17 & \begin{tabularx}{0.7\textwidth}{X} 耶穌到了，知道拉撒路在墳墓裡已經四天了。 \end{tabularx} \\ \\ \relax
11:18 & \begin{tabularx}{0.7\textwidth}{X} 伯大尼離耶路撒冷不遠，約有六里路。 \end{tabularx} \\ \\ \relax
11:19 & \begin{tabularx}{0.7\textwidth}{X} 有好些猶太人來看馬大和馬利亞，要為她們弟弟的緣故安慰她們。 \end{tabularx} \\ \\ \relax
11:20 & \begin{tabularx}{0.7\textwidth}{X} 馬大聽見耶穌來了，就出去迎接他；馬利亞卻仍然坐在家裡。 \end{tabularx} \\ \\ \relax
11:21 & \begin{tabularx}{0.7\textwidth}{X} 馬大對耶穌說：「主啊，你若早在這裡，我弟弟就不會死了。 \end{tabularx} \\ \\ \relax
11:22 & \begin{tabularx}{0.7\textwidth}{X} 我也知道，即使現在，你無論向神求甚麼，神也必賜給你。」 \end{tabularx} \\ \\ \relax
11:23 & \begin{tabularx}{0.7\textwidth}{X} 耶穌對她說：「你弟弟會復活的。」 \end{tabularx} \\ \\ \relax
11:24 & \begin{tabularx}{0.7\textwidth}{X} 馬大對他說：「我知道在末日復活的時候，他會復活。」 \end{tabularx} \\ \\ \relax
11:25 & \begin{tabularx}{0.7\textwidth}{X} 耶穌對她說：「復活在我，生命也在我。信我的人雖然死了，也必復活。 \end{tabularx} \\ \\ \relax
11:26 & \begin{tabularx}{0.7\textwidth}{X} 凡活著信我的人必永遠不死。你信這話嗎？」 \end{tabularx} \\ \\ \relax
11:27 & \begin{tabularx}{0.7\textwidth}{X} 馬大對他說：「主啊，是的。我信你是基督，是神的兒子，就是那要臨到世界的。」 \end{tabularx} \\ \\ \relax
11:28 & \begin{tabularx}{0.7\textwidth}{X} 馬大說了這話就回去，叫她妹妹馬利亞，私下說：「老師來了，他在叫你。」 \end{tabularx} \\ \\ \relax
11:29 & \begin{tabularx}{0.7\textwidth}{X} 馬利亞聽見了，急忙起來，到耶穌那裡去。 \end{tabularx} \\ \\ \relax
11:30 & \begin{tabularx}{0.7\textwidth}{X} 那時，耶穌還沒有進村子，仍在馬大迎接他的地方。 \end{tabularx} \\ \\ \relax
11:31 & \begin{tabularx}{0.7\textwidth}{X} 那些同馬利亞在家裡安慰她的猶太人，見她急忙起來，出去，就跟著她，以為她要往墳墓那裡去哭。 \end{tabularx} \\ \\ \relax
11:32 & \begin{tabularx}{0.7\textwidth}{X} 馬利亞到了耶穌那裡，看見他，就俯伏在他腳前，對他說：「主啊，你若早在這裡，我弟弟就不會死了。」 \end{tabularx} \\ \\
[1ex]
\hline
\hline
\end{longtable}
$^{1}$看見姊妹平安.
很感恩,今天早上我們一同聚集在神的殿裡面.
一同敬拜,一同思想神的話語.
在我們開始正道之前,我們一同低頭祈禱.
我們親愛天上的父,多謝你賜給我們你寶貴的話語.
求你今天早上再一次獻給我們眾人的心.
領我們到你自己的私恩寶座前.
藉著約翰所記載的一段經文,向我們每一個人去說話.
不單止使我們能夠明白真道,更加懇求聖靈.
以你的愛再一次激勵我們.
以致我們今天能夠更加信靠我們的主耶穌.
更加愛慕,親近我們的主耶穌.
願你得著一切榮耀和讚美.
禱告奉耶穌德性的明智祈求,阿們.
不知道大家有沒有試過.
可能自己一個,或者是夜晚睡不著.
夜闌人靜,眼光光.
在這裡開始回望過去走過人生的路程.
偶爾會有一些很多的假如.
在我們腦海裡浮現出來.
假如當年那件事不是這樣發生的話.
今天會怎樣呢?.
假如當日那些說話不是這樣說了出口的話.
今天會怎樣呢?.
又或者,有些事情你經歷過突然間變壞.
完全失控,你控制不了.
超乎你想像.
假如我當日早幾分鐘出門.
我就不會遇到那場意外.
如果我當時沒有病到.
我就不會錯失了這麼難得的機會.
假如當年我爸爸沒有被人騙了他的積蓄.
假如我兒子沒有群毒一幫壞人.
令到他不幸沾上了毒品.
這些例子數之不盡.
曾經有一位美國的作家.
他說假如這兩個字.
說它是全世界其中兩個最令人悲哀的字.
我相信我們每個人都經歷過這些.
如果假如的經歷.

$^{41}$有時候我們會被當中的無助感.
苦痛去淹沒,很困惑.
為什麼會這樣呢?.
為什麼神當時好像沒有出手呢?.
為什麼這麼多人偏偏是發生在我身上呢?.
神是否還在掌管我生命的一切呢?.
神是否在乎我是否在關心呢?.
神啊,如果當時你在的話.
就不會是這樣了.
我不知道今天有什麼假如.
深藏在你的心裡面.
但是生命裡面其中一個最沉重的假如.
莫過於站在我們所深愛的人的墳墓面前的一刻.
這個正正是我們今日經文所遇見的一個場景.
主啊,假如.
主啊,假如你早在這裡.
我的兄弟就不會死.
約翰福音十一章這裡記載了一個故事.
可能是約翰記載最多經文最長的一個故事.
記載了一個苦難中的家庭.
和耶穌相遇的一個故事.
這個是約翰福音記載七個神蹟裡面最後的一個.
也都可能是最重要的一個.
我們知道約翰所記錄的神蹟.
容易不是停留在神蹟本身.
拉撒路復活了.
都會再一次死亡.
但是神蹟是指向神的榮耀.
以致我們看到,我們聽到的人能夠相信耶穌基督是神的兒子.
藉著耶穌的名得著生命.
拉撒路的復活指向耶穌即將.
很快祂自己就要進入耶路撒冷.
祂要面對祂自己擲支架的死亡.
和祂要為我們復活.
耶穌的死亡和復活最終使得我們得著這個永遠不死的生命.
弟兄姊妹,如果復活,如果生命今天與我們同在的時候.
縱然我們生活仍有各種的困難,思念,苦痛.
我們仍然可以完全的相信.
依靠耶穌的大能大力.
現在我很想和大家一起看一個畫劇.

$^{81}$若翰寫下來的一個畫劇.
我們分開四幕進入故事裡面.
透過每一幕不同的人和耶穌的相遇,對話.
我們一起去看看若翰這個故事告訴我們的一些真理.
故事的第一幕其實記載在這張頭六節的聖經.
故事一開始我們知道拉薩路病重.
他病到隨時會死.
我們也知道經文告訴我們.
耶穌是這三兄弟姐妹很親密的朋友.
為什麼呢?.
因為當這對姐妹想去找耶穌.
將這個消息告訴耶穌的時候.
他們其實好像很熟悉耶穌.
他們已經知道派人去哪裡找到耶穌告訴他們.
還有第三節留意.
拉薩路這個名字對這裡是沒有出現過的.
只是稱為「你所愛的人」.
說了這個話他們就知道耶穌是在說誰.
到了第五節.
這裡確認了耶穌去愛馬大和他的姐妹和拉薩路.
現在拉薩路隨時不行了.
我們幻想一下如果我們是馬大馬尼亞.
他心裡面當時的感受是怎樣的?.
是很焦急的.
他希望派這個通風報信的人走快一點.
不時地在廚房窗外看一看.
希望耶穌趕得及來.
心裡面想.
我認識耶穌.
耶穌一定會醫治.
一切都會OK.
馬大去張望.
他去等待.
一秒一秒地過.
他期待.
但是越來越渺茫.
最後什麼都沒有.
拉薩路最後死了.
耶穌連影子都沒有看到.
他們心裡一定是很難以置信.

$^{121}$很困惑很失望.
但是我們看到耶穌一聽到.
拉薩路病得不行的時候.
他第一個反應是什麼呢?.
我們再看一次第五第六節.
今天我展示這個新漢語曆本.
給大家在和學本以外多一個參考.
第五第六節說.
耶穌一向愛馬大.
他的姐妹和拉薩路.
所以他聽到拉薩路病了.
仍在原地留了兩天.
我們不要忽略.
約翰為我們寫了好像很不起眼.
第六節第後兩個字.
所以這兩個字.
和本聖經沒有譯到.
但是在原文裡面.
是有這兩個字.
所以這個意思.
剛剛第五節.
耶穌說了他愛馬大.
愛瑪利亞.
愛拉薩路.
但是第六節就說.
耶穌不馬上去幫他們.
反而仍在原地留了兩天.
似乎耶穌是故意的.
耶穌是故意的.
耶穌故意延遲.
而這個延遲.
如果說第五節的話.
就是對馬大.
對瑪利亞.
和對拉薩路的愛.
為什麼會這樣說呢?.
我們留意.
如果我們看回頭幾節經文的時候.
拉薩路死了.
並不是因為耶穌趕不及去.

$^{161}$即使耶穌一聽到這個消息.
馬上走去趕去伯大利.
拉薩路的死.
也不能夠在耶穌來之前避免.
因為我們根據經文記載.
當耶穌抵達伯大利.
拉薩路死了多少天?.
四天.
我們從上一章看到.
耶穌當時在一個地方.
距離伯大利.
是一天的路程.
這個使者大概用了一天去找耶穌.
耶穌說他留多了兩天在原地.
最後用一天走去伯大利.
這是什麼意思呢?.
即是說拉薩路很可能在那個使者一出門沒多久.
就已經不行了.
就已經死了.
當一會兒去到第十七節.
耶穌和他們見面的時候.
拉薩路在墳墓裡已經四天.
為什麼耶穌剛剛要選兩天去延遲呢?.
有些聖經學者將背景告訴我們.
原來當時在猶太人的世界裡.
他們有一個很普遍的觀念.
原來他們覺得人死了之後.
靈魂會在身體附近逗留三天.
希望回到軀殼裡.
但他們覺得當身體腐化開始.
大概三天後靈魂就會離開.
約翰想我們很清楚的知道一件事.
拉薩路在這一刻.
耶穌見到他的一刻.
的而且確已經死了.
也就是說耶穌將要走的一個神蹟.
絕對不能夠被周圍的民眾誤解為.
靈魂自然回到身體的一個說法.
所以耶穌真的因為愛的緣故.
去延遲了兩天.

$^{201}$才啟程走去伯大利.
為的是要顯明神在他身上工作和榮耀.
同時也很想鞏固增加這個家庭對神的信心.
在我們生活裡很多時候很多的焦慮.
很多很混亂的事情要我們去面對.
有時候多到我們不是很看到.
明白到主耶穌的同在到底是什麼一回事.
又或者我們經歷了一些痛苦的時候.
我們很多時候甚至心裡會埋怨.
以為神是不是有一段時間真的離開了我們.
是不是神沒有以前那麼愛我們呢.
我們有沒有試過抱怨神沒有早點出手.
沒有在我們最需要祂危難的時候出現.
祂似乎那時候是不聞不問我們.
神啊你去了哪裡呢.
這個時候我想起詩篇第十三篇.
詩人大衛同樣去問他的神.
叫他耶和華你要忘記我到什麼時候呢.
要到永遠嗎.
你掩面不顧我要到什麼時候呢.
我心裡就算終日受苦要到什麼時候呢.
當時大衛在等神去拯救他.
但是他當刻看著的時候看不到拯救在哪裡.
但是他心裡認識他前愛的神.
信實的神.
他當時仍然選擇去傳言.
倚靠信靠這位神.
那首詩篇最後的結尾.
他這樣說.
但我倚靠你的前愛.
我的心因你的救恩快樂.
我要向耶和華歌唱.
因他用厚恩待我.
我們很多時候可以將我們自己經歷的一些試煉.
就是當時交託給神.
神未必立刻行動.
但是不論我們是否即時在困境裡得著拯救釋放.
或者神是要我們繼續去面對當時的困難.
我們知道一切是按照神的旨意.
我們倚靠耶穌的力量.

$^{241}$我們能夠在當中歸榮耀給神.
同時我們也在加深自己.
甚至是身邊的人對我們的神的信心.
我們每一位正正其實都是耶穌所愛的.
不單只是那三個兄弟姐妹.
耶穌有時的言詞絕對沒有抵觸祂的愛.
沒錯,有時我們會覺得祂比我們希望的遲了一點.
但是在神的眼中,一切都是最適合的時候.
正如一位解經家曾經這樣說.
他說,上帝是最守時的.
祂創造時間,祂從來不會遲到.
故事去到第二幕,我們跳到第十七節.
今天主席開始讀的一段經文.
這一幕是講到耶穌和瑪大去見面.
耶穌終於來到伯達尼,拉撒路在墳墓四天.
兩姐妹在喪禮進行了四天,終於見到耶穌.
這四天裡面,他們有聊天.
他們在想,如果真的見到耶穌,第一句說什麼好呢?.
想來想去,他們最後選擇了這一句,瑪大對耶穌說的第一句話.
他說,主啊,要是你早在這裡,我們兄弟就不會死了.
就是現在,我也知道無論你向神求什麼,神也必賜給你.
我們憑聖經的經文,不能夠很肯定當時他說這句話的語氣是怎樣.
但是我們估計,他多多少少都有些不高興.
抱怨耶穌,為什麼不早點回來?.
為什麼不一聽到消息就回來?.
或者他已經很累了,因為照顧了拉撒路很久.
在喪禮裡面,很多的情緒,幾天了.
他都不太知道可以向神求什麼.
他第二句話有點含糊,他對耶穌是有信心.
或者他只是想,現在耶穌在他朦朧的事奉裡面.
再給他一點安慰,給他一點同情心.
有時我們自己在困難裡面,我們都知道我們要祈禱.
我們要將我們的困難交託給神.
但是有時都說不出口,到底我想求神怎樣呢?.
不知道神怎樣幫我們.
於是耶穌在這裡沒有停留.
他嘗試去加添馬達的信心.
他說你的兄弟必復活.
馬達說,我知道,我知道.
在末日復活的時候,他會復活的.

$^{281}$我們從馬達的反應,我們想像到.
他當時沒有從耶穌這句話有很多的安慰.
因為他和當時出席喪禮的人都知道.
猶太人都很熟悉末日有復活這個信念的概念.
再者,當我們經歷過我們愛的人.
親人離世的弟兄姊妹都知道.
未來的復活對於現正住在哀傷的人.
其實是一個頗冰冷的安慰.
但是馬達似乎忽略了耶穌說的話的重點.
耶穌接著說了一句話,打破了復活這個概念.
復活不只是一個教義,不是一個概念.
復活是一個人.
這個時候耶穌可能特意去看著馬達的眼睛.
他講出了七個「我是」的宣告裡面.
其中一個最重要的.
耶穌說「我是復活和生命」.
相信我的人,雖然死了,必活過來.
又說「所有活著又相信我的人,必永遠不死」.
你相信嗎?.
耶穌這裡去將馬達的焦點暫時離開他即時的難題.
耶穌想指出一個更大的真理.
如果用馬達自己的說話,末日有復活的時候.
耶穌說「這個末日原來已經來臨」.
耶穌本身是起始,耶穌是終點.
復活不再是單單一個教義.
現在活生生有個樣子在你面前.
有個名字叫做耶穌.
耶穌自己就是超越了死亡的生命.
是神賜給我們的禮物.
耶穌要馬達知道什麼呢?.
他的盼望並不是單單在遙遠死了之後能夠上天堂.
而是當下能夠與耶穌基督相遇.
有耶穌基督在他的身邊.
在基督裡面,有一樣東西超越了我們每個人將來審判日有的盼望.
耶穌使我們戰勝死亡.
他的德性也成為了我們當下活著的現實.
也就是說,既然耶穌是復活和生命的時候.
在基督裡面,我們基督徒的生命.
主要不是等死後的生命.
乃是死前的生命.

$^{321}$因為基督徒的生命是與戰勝死亡的主聯合了的一個生命.
死亡無法去觸及我們.
所以有人說教會在本質上又可以稱為一個復活的群體.
我們這是一個復活的群體.
說我們自己已經得救信主是一回事.
但是我們真正理解什麼是在當下我們有生命,我們得著耶穌的生命的意義.
又是另一回事.
我們信耶穌已經與主聯合.
就是今天我們能夠分享這份超越死亡的生命.
當瑪大明白站在他面前的耶穌是復活,是生命的時候.
可能他當時對於將來復活的概念已經被耶穌顛覆了.
第27節,瑪大這樣回答.
他說主啊,我相信,我一直信你是基督,是神的兒子.
是要來到世上的一位.
可能我們會聯想起,彼得曾經也說過類似的.
這裡瑪大說他相信.
但是我們一直在看故事的發展.
我們知道他所信的遠超於他當刻能夠明白.
有時候我們會不會都跟瑪大一樣呢?.
有時候我們都有很多的答案.
有很多基督教正統浸信會信仰的標準答案.
但是我們卻忽視了.
我們看不到我們不是很理會耶穌當刻的同在,祂存在在我們的生命裡面.
耶穌其實問我們每一個人.
你相信這話嗎?.
你會怎樣回答耶穌呢?.
在當時.
很多時候我們去找耶穌祈禱.
很專注在自己心裡面的重擔.
但是我們看不到主耶穌就在我們身邊.
我們怎樣能夠活在.
不單是將來的盼望.
還是今天活著的盼望裡面呢?.
去到第三幕的時候.
進入到瑪利亞.
耶穌叫瑪利亞叫祂出來.
瑪利亞聽到就馬上彈起身.
祂來到耶穌那裡.
第三十二節,瑪利亞急忙起來去找耶穌.
見到耶穌就俯伏在祂的腳前.

$^{361}$在憂傷,祂在哭泣.
在痛苦,困惑,失望的裡面.
瑪利亞當然可能說,我生了耶穌的氣.
我不見祂,不理會祂.
祂來就來吧,這麼遲.
但是她選擇了另一樣東西.
她選擇了馬上起來.
跑去耶穌那裡.
俯伏在耶穌的腳前.
在主耶穌面前去哭訴.
因為這個是她唯一所信靠的主.
也是我們唯一所信靠的主.
之後她跟她姐姐瑪大一樣.
她說了一句話.
主啊,要是你早在這裡.
我的兄弟就不會死了.
好像很面熟.
中文翻譯兩句是一樣的.
姐姐跟妹妹兩句說話是一模一樣.
但是原文有一點點不同.
瑪利亞跟姐姐說的話有什麼不同呢.
原來瑪利亞這句說話.
是比她姐姐說的更加凸顯了.
是我的兄弟啊.
我的兄弟.
不是一個人的兄弟.
是我的兄弟.
耶穌,我知道你愛我的.
為什麼你會容許我的兄弟死啊.
第三十三節,耶穌見到她哭.
其他猶太人在喪禮裡面一起哭.
內心悲憤激動起來.
你將她安放在哪裡呢.
他們回答耶穌.
主啊,請你來看.
然後經文說耶穌哭了.
我們看到在憂傷裡面.
主耶穌與姐妹同在.
他也與我們同在.
他與我們感受我們的傷痛.

$^{401}$當時耶穌是真的哭了.
他切實地顯示他心裡面有憂傷.
很多時候我聽到有些人說.
基督徒,我們是信主.
不應該哭的,我們要堅強.
但是原來我們看到.
在哀傷,哭泣,眼淚.
耶穌說是OK的.
絕對沒有抵觸我們對神的信心.
耶穌沒有阻止這個家庭去悲傷.
相反,他與我們同在.
跟他們同哭.
但是我覺得值得留意的是.
約翰其實用了兩個不同意義的字.
第一個我想留意的是一個哭字.
第三個三節說到猶太人在哭.
第三個五節剛才說耶穌哭了.
兩個中文都是哭.
但是原來不同意思.
第三個三節這個哭這個字是指.
在喪禮裡面經常看到.
嚎哭大哭.
很大聲的,哭出聲的.
大叫的這樣的哭.
三十五節原文用了一個哭字.
耶穌那個哭就沒有這種大叫.
嚎哭的意思.
第二個字是內心悲憤這個字.
在三十三節.
和學本譯是心裡悲嘆,尤甚憂愁.
三十三和三十八節都有這個字.
這裡的悲憤.
其實是包含了一個憤慨的意思.
也就是說這段經文帶出了一個重要的訊息.
耶穌在哭.
不單單是因為他失去了拉薩路這個好朋友.
他的眼淚不單止是情感上的哀傷.
經文說原來他包含了憤慨的意味在當中.
事實上我們看到結局.
耶穌即將要使拉薩路復活.

$^{441}$他應該不會純粹為了拉薩路的死.
很過份的傷悲.
那為什麼.
當耶穌來到墳墓面前.
看到那些人.
他心裡有這麼莫大的一種憤慨呢.
我相信耶穌在這裡的憤怒.
他的對象是罪和死亡本身.
死亡是創世記三章墜入了世界之後的後果.
耶穌的憤怒.
是來自見到人類受罪的捆綁.
從瑪利亞的哀傷.
他感受到我們全部人.
其實都是因為罪而必須面對死亡.
對罪和死亡充滿了憤怒.
耶穌深深明白我們每一個人生命裡的軟弱.
所以當我們在困境裡.
我們有時感到徬徨無助.
甚至很絕望很痛的時候.
我們都可以俯伏在主的腳前.
好像瑪利亞一樣流淚.
將心裡的失望困惑的感情.
都告訴主耶穌.
正如主耶穌和瑪利亞說的話一樣.
其實祂都是滿有慈愛的.
來輕輕叫我們的名字.
祂在聽我們的煩惱.
祂為世上種種苦難不易的事情.
更加是你當時正面對的苦難不易的事情.
祂有憂傷悲憤.
耶穌也在問你.
你將你心中的憂傷埋藏在哪裡呢?.
你願不願意邀請耶穌說.
主啊,請你來看.
耶穌是神,耶穌是聖潔尊貴的神.
但是祂不是一個很遙遠的神.
而是和我們同經生活的風雨大小事情.
甚至和我們同哭的神.
我們遇到任何事情不是自己一個.
耶穌不單在天上很高望著我們.

$^{481}$但是祂在乎的,祂在乎.
祂明白,甚至祂願意進入我們的生命裡.
和我們一起去經歷我們所受的哀傷.
困局的當中.
最後來到第四幕,故事的高潮.
耶穌搞了這麼久,終於走到墳墓裡.
拉薩路已經死了四天.
沒有呼吸,沒有心跳.
開始身體腐化.
第39節,耶穌吩咐把石頭搬開.
馬大跟耶穌說,主啊.
祂的身體發臭了四天.
耶穌說,我不是跟你說過嗎?.
你若信,就必看到神的榮耀.
於是,他們就把石頭搬開.
弟兄姊妹,什麼是信心呢?.
信心不是我們學到很多知識.
我們頭腦覺得是對的.
或者是很正宗的教義.
或者很多宿令的活動去參與.
信心其實有時候可以很簡單.
就是對耶穌的單純的信服和回應.
因為你叫我這樣做,我就去做.
很明顯,馬大擔心屍體腐化的臭味.
他不是很明白耶穌叫他做什麼.
但話雖如此,他選擇了這一刻去信任和服從主.
我不明白耶穌挪開石頭的目的是什麼.
但主既然你這樣說,我就照做.
親愛的弟兄姊妹.
你要挪開的石頭是什麼呢?.
有什麼東西難阻你經歷生命裡.
耶穌要你得著的復活和重生?.
有沒有什麼東西深深埋藏在你內心.
可能在一些暗角裡.
你不太敢或不太想被耶穌碰到的地方.
又是什麼呢?.
可能你覺得那些東西很久了.
真的腐爛了,甚至臭了.
不想打擾耶穌,不想被耶穌看到.
或者可能太久了,放在一旁.

$^{521}$你早就已經不想理會它.
最近我有不同的場合都和不同的人聊天.
教會以外的人.
很多都是到了中年的時候.
當我們聊天,聊得深一點.
他知道我在教會做傳道牧羊工作.
他就開始問教會的事情.
好像有些認識一樣.
原來很多這些人以前都回過教會.
很久沒有回.
以前曾經缺過志,信過主.
但是他們當我去.
「你再回來吧」.
往往都很尷尬,微笑.
好像覺得自己有點不配.
覺得自己離開了神很遠.
「不要了,現在太久了」.
在信與不信之間.
又好像感受到神邀請他回去.
又覺得很多東西已經離開了神很遠.
其實是不容易的.
去面對自己的生命從來都不容易.
你未信也好,信了也好.
信了多久也好.
我們去面對生命的黑暗的時候.
從來都不是就這樣坐在那裡.
一天就過了一天.
耶穌邀請我們將我們心裡的石頭去拿開.
讓祂的光照亮我們的黑暗.
以致祂能夠更新我們.
醫治我們.
祂不是要顯露我們.
令我們慚愧,尷尬.
即時去審判我們.
耶穌是要我們得生命.
得醫治.
面對了這條屍體.
馬達其實什麼都做不到.
你和我面對屍體的時候.
其實什麼都做不到.

$^{561}$但是我們唯一可以做的.
我們可以選擇的.
就是將這個石頭搬開.
讓耶穌的說話.
祂的真光照亮我們的心.
我們可以選擇將這些爛了.
阻隔了神賜生命的.
祂的話語,祂的亮光的石頭拿開.
耶穌是復活和生命.
祂能夠使枯骨重新得著生命.
耶穌問我.
我不是對你說過.
信就能夠得著看到神的榮耀嗎?.
原來我這裡有個好消息.
信心的回應其實可以很簡單.
因為耶穌叫我.
你叫我拿開這個石頭.
我就拿開.
因為耶穌叫我做.
我就去做.
我們看看在墳墓面前.
這個石頭被搬開.
我們聽一下.
耶穌不是很小聲的說.
耶穌是滿有神的權柄.
大聲的呼喊.
拉薩路出來.
拉薩路就在墳墓走了出來.
我們怎樣能夠以這樣的盼望和信心.
面對生命一個又一個的難關.
縱然好像在很黑暗的處境.
但是我們在那裡.
仍然可以看到神的榮耀.
我們不需要每一次都透過我們的困難.
去看主耶穌.
而是可以透過活著的主.
來看待我們的困難.
今天早上.
弟兄姊妹在你的生命裡面.
你感受到神願你的同在嗎.

$^{601}$我們和基督一同復活.
我們和基督的生命去聯合.
不代表我們的困難會盡失.
耶穌說我們在世有苦難.
但是如果我們能夠.
與這位又真又活.
復活的主同行的時候.
我們就能夠謙卑.
憑信心去面對一切的困難.
透過主耶穌給我們的力量.
去承受生命各樣的試煉.
仍然有喜樂的.
有盼望的.
渡過我們每一天.
因為我們每一位信主的人.
現在不再是自己了.
但是屬於基督.
所以或生或死.
我們知道我們都屬於復活的主.
耶穌說我是復活和生命.
信我的人.
雖然死了也必活過來.
所有活著又信我的人.
必永遠不死.
你相信這話嗎.
讓我們一同低頭安靜.
在安靜的裡面.
我們可以與神傾心吐耳.
回應神的邀請.
與今天活著的基督.
去說你想回應祂的事.
主耶穌.
在我們生命裡都遭遇困苦.
有時彷彿有無助的境界.
主耶穌求你幫助我們去看見.
原來在每一個難處裡.
不單止是我們.
主,你在當中,你是復活和生命.
主,我們看見身邊的人.
我們看見新聞.

$^{641}$很多的不公義,很多的苦難.
我們看到覺得絕望.
但是主耶穌,我們知道在你的裡面.
我們得著盼望.
因為你不單止沒有在我們憂傷的時候.
曾經離開我們.
原來你一直感受著我們感受的痛.
甚至曾經為我們的傷痛流淚.
主,你幫助我們.
將我們心裡需要搬開的石頭.
來去看見祂.
可能是纏繞我們很久的罪.
可能是很多年放不下的執著.
可能是那個永遠不能寬恕原諒的弟兄或姊妹.
眼是不順眼的.
主啊,求你給我們勇氣去搬走它.
用讓利去進入我們內心.
即使我們覺得臭.
但是主你不覺.
求你讓我們今天再一次得著完全的醫治潔淨.
願你的榮耀今天在我們每人的生命裡面去彰顯.
我們同心禱告.
是奉靠主耶穌基督得勝的名字祈求.
阿們.
\newpage



\section{腓立比書 3:17-4:7}
\label{sec:031891}
\textbf{腓書回顧 (3) 天民}~~~~盧偉心弟兄
\newline
\newline
連結: \href{http://www.church.com.hk/acms/content.asp?site=cdc\&op=show\&type=product\&code=031891\&layout=sermon}{\texttt{031891}} ~~~~ 語音日期: 2026-03-28
\newline
\newline
\hyperref[sec:031890]{\small{< < < PREV SERMON < < <}}
~
\hyperref[sec:index_chronic]{\small{[返順時目]}}
~
\hyperref[sec:031892]{\small{> > > NEXT SERMON > > >}}
\newline
\newline
腓立比書 3:17-4:7
\newline
\begin{longtable}{cl}
\hline
\hline
章節 & 經文 (和合本修訂版)\\
\hline
3:17 & \begin{tabularx}{0.7\textwidth}{X} 弟兄們，你們要一同效法我，也當留意看那些效法我們榜樣的人。 \end{tabularx} \\ \\ \relax
3:18 & \begin{tabularx}{0.7\textwidth}{X} 因為，我屢次告訴你們，現在又流淚告訴你們：許多人行事是基督十字架的仇敵。 \end{tabularx} \\ \\ \relax
3:19 & \begin{tabularx}{0.7\textwidth}{X} 他們的結局就是滅亡。他們的神明是自己的肚腹；他們以自己的羞辱為光榮，專以地上的事為念。 \end{tabularx} \\ \\ \relax
3:20 & \begin{tabularx}{0.7\textwidth}{X} 我們卻是天上的國民，並且等候救主，就是主耶穌基督從天上降臨。 \end{tabularx} \\ \\ \relax
3:21 & \begin{tabularx}{0.7\textwidth}{X} 他要按著那能使萬有歸服自己的大能，把我們這卑賤的身體改變形狀，和他自己榮耀的身體相似。 \end{tabularx} \\ \\ \relax
4:1 & \begin{tabularx}{0.7\textwidth}{X} 我所親愛、所想念的弟兄們，你們就是我的喜樂，我的冠冕。我親愛的，你們應當靠主站立得穩。 \end{tabularx} \\ \\ \relax
4:2 & \begin{tabularx}{0.7\textwidth}{X} 我勸友阿蝶和循都基要在主裡同心。 \end{tabularx} \\ \\ \relax
4:3 & \begin{tabularx}{0.7\textwidth}{X} 我也求你這真實同負一軛的，要幫助這兩個女人，因為她們在福音上曾與我、革利免和我其餘的同工一同勞苦，他們的名字都在生命冊上。 \end{tabularx} \\ \\ \relax
4:4 & \begin{tabularx}{0.7\textwidth}{X} 你們要靠主常常喜樂。我再說，你們要喜樂。 \end{tabularx} \\ \\ \relax
4:5 & \begin{tabularx}{0.7\textwidth}{X} 要讓眾人知道你們謙讓的心。主已經近了。 \end{tabularx} \\ \\ \relax
4:6 & \begin{tabularx}{0.7\textwidth}{X} 應當一無掛慮，只要凡事藉著禱告、祈求和感謝，將你們所要的告訴神。 \end{tabularx} \\ \\ \relax
4:7 & \begin{tabularx}{0.7\textwidth}{X} 神所賜那超越人所能了解的平安，必在基督耶穌裡，保守你們的心懷意念。 \end{tabularx} \\ \\
[1ex]
\hline
\hline
\end{longtable}
$^{1}$各位兄弟姐妹午安.
歡迎大家來到十佳基督教會.
很感恩我們能夠繼續看「肥辣比書」.
之前我們用了十四篇的講道.
完成了整本「肥辣比書」.
然後上年年底我們開始去回顧.
我們的大題目叫做「肥書回顧」.
一方面我們是回顧「肥辣比書」.
另一方面我們希望透過「肥書回顧」.
讓我們每一個都成為一隻肥豬.
意思是FAT.
Faithful.
我們忠誠委身.
Available我們隨時候命.
以及我們Teachable樂於受教.
我們會用五篇講道來做回顧.
每一次我們都會重溫這個「肥辣比書」.
我弄了一個叫「超濃縮版」.
去幫助大家去記得這本書當中的重要片段.
在一開始的時間.
保羅去稱呼自己叫做基督耶穌的「僕人」.
已經指向這封信的中心「基督之歌」.
因為「僕人」這個字在「基督之歌」會再次出現.
然後保羅送上文案.
他說「願恩惠平安.
從我們的父上帝和主耶穌基督歸給你們」.
「恩惠」這個字是另一個很重要的字眼.
然後「哪日必成者功」是指一章六節.
保羅說他深信上帝會一直在「肥辣比信徒」當中去作功.
到了基督耶穌再來的日子.
會完成這個工作.
「三合一」是指保羅為「肥辣比信徒」所祈求的三樣東西.
「愛心」,「真認識」和「審察力」.
這三種屬靈的能力不斷去增長.
以至榮耀歸給父上帝.
然後「詩」部分.
保羅分享他的近況.
他說他雖然被軟禁.
但是為著福音能夠更興旺.
為著弟兄們能夠更加放膽去傳道.

$^{41}$他非常歡喜.
然後他分享他的內心世界.
他說無論是生是死.
總要讓基督在他身上照常顯大.
因為他活著就是基督.
死了就有益處.
保羅定義他無論生死.
都要為了榮耀上帝.
去到「D」部分.
「天文相稱福音」是指一章二十七節.
保羅說你們作為天國的公民.
行事為人就要與基督的福音相稱.
相稱是彰顯價值的意思.
那麼我們要怎樣去彰顯基督福音的價值呢?.
當中他講到我們要無畏.
要站穩.
要同心.
要競鬥.
所謂競鬥其實是為福音齊心努力的縮寫.
因為希臘文是sun aphilio.
sun是一齊的意思.
aphilio是一群運動員參與對際比賽.
為勝利一同努力的圖像.
然後「合一謙卑他顧」是二章一至四節的內容.
大家可能記得了.
各人看別人比自己重要.
各人不要單顧自己的事.
也要顧別人的事等等.
最後「以基督的心為心」是二章五節.
我們要集中思想.
耶穌基督的心目中.
甚麼才是合一,謙卑和他顧.
然後這部分就是.
菲勒比書最重要的中心.
基督之歌.
耶穌基督經歷八重的下降.
從本如上帝同等.
一直虛己.
一直下降.
直到祂全心信服.

$^{81}$以至於死.
且死在十字架上.
然而基督受苦.
是為要進入榮耀.
祂會得著超乎萬民之名.
萬室都要跪拜.
萬口都要宣認.
耶穌基督是主.
歸榮耀給付上帝.
而因為基督之歌.
我們應當「戰經」作.
不怨言.
恐懼戰經.
是上帝的大能.
祂的恩典.
臨到我們生命當中的記號.
是因為上帝在我們生命中.
又立志.
又實行.
是因為上帝使我們能夠與福音相承.
而所謂不發怨言.
不是自免的意思.
當中是指.
出埃及時期的以色列人.
因為他們念念不忘埃及的偶像.
念念不忘地上的榮耀.
他們興翰上帝的恩典和拯救.
然後顯主道明光照.
是要提醒我們.
我們要高舉 hold forth.
我們要堅守 hold fast.
上帝的道.
讓我們在這個彎曲半謀的世代當中.
好像明光一樣照耀.
然後X部份是兩個與福音相承的生命的例子.
保羅提到提摩泰.
他如何不求自己的事.
真正去關懷肥達比的信徒.
他亦提到以巴弗提.
如何為了完成基督的工作.

$^{121}$賭上生命.
近乎於死.
去到C部份.
保羅再講他自己.
他講到他自己以往是如何靠肉體.
現在他因為基督.
明白地上的事情.
包括能夠讓人自以為是的各樣好行為.
其實都是毫無價值.
都是糞土.
因此他定義繼續去刁棄萬事.
單單為要贏得基督.
但是保羅他不認為自己已經完全了.
所以他應主召.
竭力追.
竭力追求的希臘文是Dioko.
與他之前迫害教會的迫害都是同一個字.
保羅他竭力追求.
是因為基督耶穌已經用救恩私恩買贖了他.
然後去到D2的部份.
保羅勉勵肥達比信徒要效法他.
亦都要留意那些與他們一起去效法保羅的人.
但是這一班有一班人.
他們是口不對心的十家的仇敵.
這班人專以地上的事為念.
相反如果我們是真正天上的國民.
我們會以天上的事為念.
專心等候救主從天降臨.
等候擁有與基督一樣的榮耀身體.
下面主近了同心證.
是在說四章.
在主已經近了的前提之下.
保羅去勉勵肥達比信徒要靠主站立得穩.
他特別提醒兩個姐妹.
友雅蝶和秦都姬在主裡面同心.
要他們回想他們以前.
如何與保羅一起.
Sunna filio是同一個字.
他們如何一無掛慮地為著永恆的事情.
然後去到B部份.

$^{161}$但是保羅不想他們誤會或錯估他們奉獻的作用.
所以他去強調他自己是知足的.
強調他自己不求饋贈.
但另一方面其實這些奉獻.
是他們的天國存款.
是歸在他們帳上的果子.
是能夠將榮耀歸給父上帝.
最後保羅用.
願主耶穌基督的恩惠與你們的靈同在.
作為最後的問安和祝福.
這句呼應著一章二節的問安.
願恩惠平安.
從我們的父上帝和主耶穌基督歸給你們.
肥達比書最重要的就是中間部份.
基督之歌.
我們這五篇的回顧系列.
我們是會由外至內進行的.
之前兩篇我們已經看了.
ABC C' B'和A'.
今天開始我們會進入.
D D'和D'2這個部份.
我們會用兩篇道去回顧這三段.
我們分別會是今天的肥書回顧三天文.
和下一篇.
下一次我們會回來肥書回顧四相情.
今天的內容是一章二十七節和D'2的部份.
我們會分三點.
第一點.
天文福音相稱.
天文後主榮耀.
第二點.
孝榜讓防仇敵.
施天事誇十家.
第三點.
主近了同心正.
無慮祈求永恆.
第一點.
天文相稱福音.
天文後主榮耀.
一章二十七節.

$^{201}$保羅說最重要的是.
你們行事為人.
要與基督的福音相稱.
又或者另一個翻譯.
你們作公民.
總要對得起基督的福音.
在這裡和合本易作最重要的是這個字.
其實在原文希臘文.
是一個很強烈的連接詞.
甚至可以翻譯為.
只要.
有些英文,日本話.
only.
就好像福音書裡面大家可能聽過.
不要怕.
只要信這句經文.
只要就是正正是這個字.
所以如果你記得.
腓立比書一章的內容.
保羅之前在說.
菲律賓的弟兄姊妹.
我很為你們感恩.
我為你們禱告.
我為你們很高興之餘.
我也告訴你們.
我在個人生死榮辱的見證之後.
現在.
我要提醒你們一件最重要的事情.
only.
你最重要.
最不能忘記的這件事.
就是.
你們行事為人.
要與基督的福音相稱.
什麼叫相稱呢?.
原文其實是Asios.
意思是一個天秤.
同等重量,同等價值的意思.
因為在保羅時代的人.
一講到同等重量,同等價值.

$^{241}$就會想起這個工具.
天秤怎麼用呢?.
譬如你要去買一磅米.
他首先在一邊放上一磅的砝碼.
另一邊就倒米上去.
倒到大家平衡.
你就知道那些米是一磅的.
所以這裡是一個很清晰的圖像.
他在說.
你將基督的福音放在裡面.
然後你將你的生命倒到另外一邊.
你試想一下.
天秤一邊放的是基督的福音.
另一方面.
你要將你的生命倒到另外一邊.
這個圖畫會是一個怎樣的圖畫呢?.
我們下一篇導讀會再回來詳細思想.
另外一個很重要的字眼就是行事為人這個字.
原文是Politurme.
意思其實是作功文.
這樣的意思.
因為菲律賓本身是一個羅馬的殖民地.
住在菲律賓的人.
很多都是羅馬的公民.
所以很多人本身是很為他們羅馬公民這個身份去自豪.
他們會穿羅馬的衣服.
他們會用羅馬的名字.
他們會說拉丁文.
他們遵守羅馬的法律.
他們整個的生活方式都在向這個世界說.
我們是羅馬人.
我們屬於羅馬.
腓立比書裡面保羅就活用這個背景.
他去提醒菲律賓信徒.
他說雖然你們地上是羅馬的公民.
但是在被基督的寶血買熟之後.
你們是成為天國的公民.
保羅的意思是說你們明不明白.
你想清楚.
你們其實是天國的公民.

$^{281}$雖然你暫時住在地上.
但其實你的家鄉是天國.
你效忠的對象.
是天國的君王.
所以你的生活方式.
你的人生取向.
你的價值判斷.
都應該是和基督的福音相稱.
你的生活方式.
究竟是向這個世界宣告.
你是屬於天國嗎?.
這個是天國公民的責任.
不過除了責任.
其實天國公民還有另外一個.
和世人自然不同的盼望.
我們看三章二十和二十一節.
我們卻是天上的國民.
並且等候救主.
就是主耶穌基督從天上降臨.
天國公民他們會專心等候救主.
因為救主是主耶穌基督.
他會從天降臨.
等候不是在說法外.
原文那裡其實是在說一個.
專心期待.
迫切仰望.
翹首以盼的意思.
天國公民的生命有一個標誌性的姿態.
他們很迫切地仰望天上.
他們很專心.
很期待耶穌基督的回來.
為什麼迫切地期待呢?.
因為當耶穌再來的時候會做一件大事.
二十一節.
他要按著.
那能使萬有歸服自己的大能.
把我們者卑賤或者謙卑的身體改變形狀.
和他自己榮耀的身體同形質.
在這個短短一節經文.
其實普羅已經三次帶我們回到基督之國.

$^{321}$第一.
那能使萬有歸服自己的大能.
大家會想起基督之國後半部.
萬失跪拜萬口宣認.
耶穌基督是主.
另外卑賤的身體.
和基督之歌裡面.
說耶穌基督謙卑自己.
是同一個字來的.
和他謙卑.
最後.
和基督榮耀的身體相似.
其實這個相似是一個比較弱的翻譯.
更好的翻譯是同形質.
我們上次講過了.
孫莫非.
孫是同樣的意思.
形質是一個叫做內在本質的外在體現的意思.
是讓讀者去返回基督之國.
因為莫非這個字.
整本聖經.
只有基督之歌會出現.
一見到這個字.
就會想起基督之歌.
普羅說.
基督本有上帝的形質.
他為我們成為人.
取了奴僕的形質.
然後他再回來的時候.
他讓我們的身體和他同一個形質.
所以普羅這裡是在說.
他說天上的國民.
你們要期待.
因為那一天就快回來了.
就快到了.
耶穌基督要用他那個.
叫宇宙萬有都歸服他的大能.
將我們這個卑賤的身體.
徹底去反轉.
我們會得到一個全新的榮耀身體.

$^{361}$一個不再受罪惡.
痛苦.
死亡.
克制的身體.
一個和復活基督.
身體一模一樣的身體.
這個是我們的盼望.
而這個盼望之所以存在.
是因為耶穌基督曾經為我們謙卑自己.
卑微到底.
且死在十字架之上.
我們之所以能夠成為天國公民.
我們之所以能夠回到家裡.
是因為主耶穌基督的十字架.
所以普羅這裡不斷帶我們回到回想基督之歌.
其實天國公民這個詞所承載的.
真的不是那麼簡單.
其實是整本聖經的舊因歷史.
上帝從創世以來.
就和他的子民預備的榮耀身份.
當我們要明白天國公民是什麼意思的時候.
我們可能要先回去舊約聖經.
在《出埃及記》十九章.
當上帝剛剛將以色列人.
由埃及維勞之家領出來.
在西奈山的曠野.
上帝就透過摩西對他百姓講他的意象.
上帝說如今你們若實在聽從我的話.
遵守我的約.
就要在萬民中作屬我的子民.
因為全地都是我的.
你們要歸我作祭司的國度.
為聖潔的國民.
上帝帶領以色列人出埃及.
他拯救他們.
不只是要他們脫離奴僕得著自由.
上帝拯救他們.
是要使他們從萬民當中分別出來.
並且.
作一個屬於上帝的國度.

$^{401}$成為一班屬於上帝的國民.
他們存在的目的是作祭司的國度.
不是一個祭司.
是整個國度都是祭司.
他們要在地上代表上帝.
他們要將萬民帶到上帝面前.
這就是天國最初的啟示.
上帝在西乃山設立一個上帝的國度.
這個國度的君王是耶和華自己.
國度的子民.
是以色列人.
國度的憲法是十戒和律法.
國度的使命.
是要向萬國見證這個獨一的真神.
不過問題來了.
以色列人失敗.
他們沒有分別為聖.
他們沒有單單成為上帝的聖潔國民.
他們效法周圍列國.
一起去拜偶像.
行不宜.
他們本來應該去作祭司引領萬民歸向耶和華.
但結果他們自己離耶和華越來越遠.
所以舊約的先知不斷去呼喊.
不斷的去預言.
他說有一天.
上帝要親自來建立一個新的國度.
賜給他的子民一顆新的心.
他要將上帝的律法.
寫在這群子民的心裡面.
然後耶穌基督來了.
他出來傳道的第一句話是什麼.
日期滿了.
上帝的國近了.
你們要悔改信奉音.
耶穌來到是要建立那個從西乃山就預告的上帝的國.
不過他是怎樣建立這個國度呢.
他是不是騎著戰馬.
率領天軍.
武力推翻羅馬政權.

$^{441}$我們知道不是.
他建立國度的方式.
完全顛覆人類的想像.
基督之歌告訴我們.
那位本有上帝形象的.
那位本應該坐在寶座上的君王.
他虛己.
他降卑.
他信服至死.
且死在十字架上.
十字架就是耶穌基督登基坐王的地方.
耶穌就在十字架上面用自己的血.
去買贖一群子民.
他歸給了上帝.
他用自己的身體.
打開那條通往至聖所的路.
他用自己的死亡.
去敗壞那個掌死權的魔鬼.
去到聖經最後一卷書.
使徒約翰就發出這樣的頌讚.
他說他愛我們.
用自己的血.
使我們從罪中得釋放.
又使我們成為國度.
作他赴上帝的祭祀.
大家看到嗎.
《出海及記》十九章的那句話.
就在十字架上應驗.
是耶穌用自己上帝的血.
買贖我們成為他的子民.
上帝在西乃山的應許.
在伽略山上實現.
而我們這些本來不配的人.
就因為耶穌的血.
成為天國的公民.
成為聖傑的祭祀.
這就是為什麼保羅在《肥勒比書》一章二十七節說.
當你們作公民要對得起基督的福音的時候.
其實背後帶著一個很重很重的份量.
因為這個身份不是我們賺回來的.

$^{481}$我們地上移民.
我們有投資移民.
我們有高財通.
我們可以靠錢財.
靠學歷.
靠成就去換一個公民身份回來.
但是天國公民.
是耶穌用他的命買回來的.
當耶穌被掛在十字架上的時候.
那個本應是你我承擔的刑罰.
由他去承擔.
那個本應是你我承受的羞辱.
由他去承受.
這個天國公民的身份真實成本.
我們經常說白白.
白白但是不是免費.
是極度昂貴.
是上帝兒子的血買回來的.
所以如果我們明白這份恩典的寶貴.
如果我們真的明白.
天國公民這個身份的重量.
我們應該怎麼回應呢?.
那就帶我們進入第二部分.
「孝榜樣防仇敵.
斯天是誇十假」.
「弟兄們你們要一同孝法我.
也當留意看那些.
照著我們榜樣行事的人.
因為許多人行事.
是基督十字架的仇敵.
他們的結局就是滅亡.
他們的神明是自己的肚腹.
他們以自己的羞辱為光榮.
尊以地上的事為念.
我們卻是天上的亡民」.
所以如果我們真的明白.
天國公民身份的重量.
我們才能夠明白.
為什麼保羅這裡用這麼激烈的字眼.
他說一群人.

$^{521}$他說他們是基督十字架的仇敵.
去形容那群人.
他們是在菲律賓信徒身邊.
外表上一起去孝法保羅.
不過內心呢.
是以地上的事為念.
為什麼叫他們叫做基督十字架的仇敵?.
因為他們的問題.
不是他們不夠虔誠.
不是他們不夠努力.
不夠屬靈.
他們的問題是他們的生命.
完全沒有反映到十字架的重量.
這些人他們指稱是天國公民.
不過他們的生活形態.
他們價值取向.
他們的人生追求.
其實跟地上的國民毫無分別.
保羅說他們的神明是自己的肚腹.
不是說他們貪食.
喜歡吃東西.
是他們被自己肉體的慾望.
被這個世界的滿足感所支配.
他們人生運轉的中心.
經常只是想我想要什麼.
有什麼對我最好.
什麼能夠讓我快樂.
他們會以自己的羞辱為光榮.
那些本來應該.
我們覺得羞恥的事.
他們會覺得成為光彩.
但是在十字架光照之下.
其實那些是.
要說眼中看為惡的事.
但是他們會將這種生活方式.
當作是正常.
甚至是值得誇耀的事.
以地上的事為念.
這裡我建議你將那個專字刪掉.
因為其實原文當中沒有這個專字.

$^{561}$所以這裡不是說他們真的100\%.
所有時間以地上的事為念.
其實.
說的是他們的思念.
他們的掛慮.
他們的投資.
他們心目中想的事情.
他們的夢想.
是地上的事.
是眼看得見.
手摸得著的事.
而保羅說.
這樣的人.
是基督十字架的仇敵.
為什麼呢?.
因為十字架本來.
是要將我們從這個世界的捆綁當中釋放出來.
是將我們從以自我為中心的生命扭轉過來.
不過這些人他們將十字架當作一張保險單.
他們以為只要口裡稱耶穌為主.
就可以繼續安安穩穩地.
活在地上的國度.
他們從來沒有被十字架震撼過.
他們從來沒有被十字架的愛俘虜過.
從來沒有讓十字架重新定義他們生命的價值.
保羅說我們要留意.
留意看這些十字架的仇敵.
我們要提防他們.
不要像他們一樣.
那我們要成為一個怎樣的人呢?.
不知道你記不記得這個表.
其實十家仇敵的反轉就是天國公民.
十家仇敵思念地上的事.
天國公民就應該思念天上的事.
將來的事.
永恆的事.
什麼是思念天上的事呢?.
簡單來說就是準備移民.
天國公民深信這個世界快將過去.
所以他們心裡所想的.

$^{601}$手裡所作的.
都是為將來移居天國做準備.
對於天國公民來說.
地上的一切都是毫無價值.
因為地上的一切都要被火燒盡.
不知道你記不記得這個故事.
我曾經說過這個虛構的故事.
有一個科學家.
他宣佈地球正面臨末日.
雖然不知道什麼時候會發生.
不過隨時會發生.
而且一定會發生.
所以他就預備一艘太空船要逃離地球了.
他更加邀請所有人一起登上這個太空船.
一同去另一個新的星球生活.
那位科學家就留下一本逃生指南.
接著他就先前往那個新星球去預備地方.
他就說他必定會在地球末日之前再回地球的.
所有願意登上太空船的人.
都會跟著他一起前往新地球.
新星球那裡.
假設我們真的很相信.
確信這個科學家所說的話是真實的.
那我們會做什麼呢?.
我們還會不會繼續現在的生活.
繼續照常上班,上學.
繼續跟世人一樣累積地上的財寶呢?.
我們知道不會的.
理所當然.
我們要想盡辦法找這個太空船.
想登上這個太空船.
理所當然我們會費盡心機.
費盡唇舌.
叫所有我們所愛的人一起去這個太空船.
理所當然的.
我們會不會一直不斷去看那個逃生的指南.
我們不會再重視地上的一切.
因為這些帶不走.
相反我們會想.
將來去到新星球.

$^{641}$我們要有什麼要準備呢?.
我們會累積一些在新星球有用的東西.
這個就是思念天上的事.
我們要登上太空船.
我們要勸人登上太空船.
我們要積累在新星球有用的東西.
不過其實呢.
我一直在想.
其實這個故事跟我們的福音故事.
仍然相差很遠很遠.
因為這個故事只是表達了我們一些.
思念天上的事.
一些功能性,很功利性的東西.
但是完全忽略了我們福音故事當中.
更奇妙,更重要的東西.
因為在這個虛構故事裡面.
地球末日只是一個隨機的災難.
不過在我們真實的福音故事當中.
世界會終結.
是因為世人犯了罪.
因為罪的功架乃是死.
末日和死亡.
原是我們應該受的懲罰.
在這個虛構故事當中.
科學家製造太空船.
是因為他自己也要走.
他只是順便救你而已.
他不需要付太多的代價.
不過在我們真實的福音故事裡面.
主耶穌基督不是需要逃生的一個.
祂是被罪人規附.
理應施行聖潔審判的上帝.
但是因著祂豐富的憐憫和大愛.
祂為我們預備了一個逃生計劃.
用的不是太空船.
是十字架.
祂不是順便去救我們.
祂是捨磚降跪.
倒空自己.
付出重價來拯救我們.

$^{681}$祂本有上帝的形象.
祂謙卑自己取了奴僕的形質.
全心信服.
以至於死.
且死在十字架上.
最後在這個虛構故事裡面.
科學家去了新星球.
他不會再與我們有任何聯繫.
如果我們要去得救.
我們只可以靠自己.
不過在我們真實的福音故事當中.
主耶穌基督升天之後.
祂沒有撇下我們.
祂差派聖靈住在我們心裡面.
祂的福音大能.
祂極豐富的恩典.
要繼續幫助我們逃生.
並且去實行祂留下的使命.
去為祂作見證.
直到地極.
所以為何天國公民是以十字架為榮.
為何他們只誇基督的十字架.
所以如果你問保羅.
為何你只思念天上的事.
為何你會視地上萬事如糞土.
以認識基督為自保.
我相信他不會很功利地告訴你.
我要逃避末日.
我要前往新星球.
不會的.
更大可能他給你的答案.
是十字架的奇妙和偉大.
他會告訴你.
上帝的慈悲是多麼的浩大.
以至人唯一合乎情理的回應.
就是獻上身體當作活祭.
單單思念天上的事情.
他會告訴你.
為上帝,為福音,為天國而癲狂.
是正常的.

$^{721}$因為基督的愛催逼著他.
那份上帝使無罪的.
替我們成為罪的那份愛.
那份一人替眾人死的愛.
催逼著他.
推動著他.
席捲著他.
使他不可能再為自己而活.
他只能夠單單專一地.
為替他死.
並且已經復活的主而活.
他會告訴你.
我不以別的誇口.
只誇我們主耶穌基督的十字架.
因為在這個十字架前面.
世界.
就我而論.
世界已經釘在十字架上.
就世界而論.
我已經釘在十字架上.
所以弟兄姊妹.
那些以地上的事為念的人.
他們為什麼是十字架的首領.
因為他們的生命否定了十字架的大能.
他們的生命,他們的生活.
就是不斷向這個世界宣告.
其實十字架沒什麼大不了.
我繼續可以照常過舊有的日子.
我可以照舊追求我的巴別夢想.
照舊去累積我的地上財富.
但如果我們.
我們真是被基督的血買回來的天國公民.
我們的生命.
應該是向這個世界宣告.
十字架改變了一切.
因為十字架我們不再屬於這個世界.
因為十字架我們屬於天國.
因為十字架我們有榮耀的盼望.
因為十字架.
我們專心等候那位從天而降的救主.

$^{761}$我們要等候得著那個榮耀的身體.
當然等候不是消極的什麼都不做.
等候當中.
我們要同心靜.
我們要求永恆.
我勸有阿諛和秦都基要在主裡同心.
要幫助這兩個女人.
因為她們在福音上曾一同勞苦.
她們的名字都在生命冊上.
你們要靠主常常喜樂.
要讓眾人知道你們謙讓的心.
主已經近了.
應當一無掛慮.
將你們所要的告訴上帝.
上帝必保守你們的心懷意念.
這一大段的經文.
如果我要選一句最關鍵最核心的句子.
我會選主已經近了.
主已經近了的原文其實只有三個字.
原文是Hocrius Angus.
Hocrius就是主的意思.
The Lord.
Angus就是近的意思.
所以你會看到這一句句子.
其實本身是沒有動詞的.
就這樣 The Lord 近.
主 近.
沒有清晰的時態.
但是根據上下文.
其實保羅似乎想同時表達兩個意思.
第一個就是現在式.
現在式的意思就是保羅想表達.
主現在就近在我們信徒身邊.
這句時態是強調空間上主的同在.
是鼓勵你們在受苦當中.
相信主在你身邊.
另外一個就是將來式.
就像翻譯一樣.
主已經近了.
強調時間上基督再回來的臨近性.

$^{801}$他提醒信徒在受苦當中忍耐.
因為主必定會再回來.
施行審判和拯救.
所以保羅用Hocrius Angus.
同時表達現在式空間上的同在.
和將來式時間上基督再回來的雙重意義.
而在這幾段經文當中.
雖然保羅沒有用因為和所以.
但是他的意思其實清楚.
我們可以自己加進去.
因為主近了.
所以你們可以常常喜樂.
在受苦當中.
他的背景是受苦.
在受苦當中靠主常常喜樂.
因為主近了.
所以你們在受苦當中.
可以讓眾人知道你們謙讓的心.
因為主近了.
所以你們在受苦當中.
不過可能更重要的是.
因為主近了.
所以我們要同心靜.
我們必須同心.
專心地見證天國.
傳揚福音.
所以保羅在四章二三節裡面.
他勸有阿諜和秦都基在主裡同心的時候.
留意他其實沒有特別去開庭審訊.
沒有判斷誰是誰非.
他沒有分析誰對誰錯.
保羅在這裡帶他們去一個更高的地方.
他提醒這兩個姐妹.
他說雖然.
你們要想起你們曾經怎樣殉職.
怎樣一起並肩作戰.
一起去競鬥.
接著保羅也帶他們去到永恆的邊緣.
他提醒他們你們兩人的生命.
都在主的生命冊上.

$^{841}$所以親愛的弟兄姊妹.
我相信我們每一個人.
都曾經在教會當中.
經歷過一些人際關係上的衝突和傷害.
在經歷的時候.
其實那份情緒是真的.
那份傷害也是真實的.
在經歷的時候.
我們不會想聽到保羅所說的.
我們聽不入耳.
所以現在我們先聽.
留心去聽.
以致有需要的時間.
我們可以回想起保羅的提醒.
保羅要我們記得主勤了.
主現在在空間上與你同在.
他會完全明白.
你的情緒,你的傷害.
主勤了.
我們會即將進入永恆.
記得在永恆的角度來看.
這一刻的情緒.
這一刻的傷害是暫時的.
在整個天國救贖計劃的角度來看.
在宏觀來看.
英文叫做.
In the grand scheme of things.
其實我們的衝突和傷害.
實實在在在客觀上.
沒有我們主觀感受到的那麼重要.
親愛的電視節目主勤.
正是因為我們作為天國公民.
會專注思念天上的事.
正正因為我們的使命.
是勸人登上太空船的緣故.
正正因為我們相信主勤.
時間非常緊迫的緣故.
我們更加要靠主同心.
一同努力去見證福音.
你想想一班消防員.

$^{881}$在火災現場.
千千萬萬的被困者等待營救之下.
難道我們還會因為衝突傷害.
停下手腳嗎?.
不會的吧?.
所以讓我們一同主勤了.
同心靜.
後半句.
無慮.
祈求.
永恆.
四章六節這句經文.
經常會被斷章取義.
很多人都以為他在說.
我們禱告當中其實什麼都可以求的.
什麼都求得.
上帝祂自有主意的.
祂會決定答不答應你的祈求.
因為這裡說凡事嘛.
凡事就是什麼都可以.
其實不是保羅的意思.
首先我們要留意.
我剛才說過經文的背景.
保羅寫這封信的時候.
菲律賓信徒正在面對威脅.
逼迫.
征戰.
一章尾那裡說過.
他們為福音的緣故.
正在經歷敵人的威脅.
為主去受苦.
而且正在參與屬靈的征戰.
而在這個背景之下.
保羅才說你們要一無掛慮.
將你所要的告訴上帝.
另外.
掛慮這個字其實本身也跟逼迫有關的.
在《馬太福音》十章十九節.
《路加福音》十二章十一節.
當耶穌說.

$^{921}$那些門徒將來會因為福音.
被壓解到官府面前的時候.
耶穌說你們不要擔心怎樣去答辯.
擔心就是掛慮這個字.
本身是跟著逼迫的.
還有保守這個字.
保守這個字不是只是保護這麼簡單.
其實是在說打仗.
在征戰裡面.
你將一隊士兵放在城鎮那裡.
去防守著城市.
就是這個字.
所以這裡其實保羅是在說.
不是說凡事你們所要的.
你們所想的.
你們生活中大大小小一切所要的.
我要升職.
我要買房子.
我要身體健康這些東西.
去祈求.
保羅是在說.
你們在威脅逼迫征戰當中的祈求.
你們在福音裡面受苦.
參與屬靈征戰當中的祈求.
保羅說當你為福音去受苦.
當你身處屬靈征戰的時候.
你不要掛慮.
你只要祈禱.
為著你在征戰當中的需要.
為著我們能夠致死忠心.
為著福音能夠被傳開.
將這些事情告訴上帝.
上帝就必在基督耶穌裡面.
好像一個駐守城池的士兵.
一隊士兵一樣.
保守著我們的心思,怡念.
讓我們不會被恐懼去擊倒.
讓我們在征戰當中.
站立得穩.
正如牧師提醒我們.

$^{961}$禱告是向量不是純量.
純量就是看數字不看方向.
我們禱告多久多頻密多長.
向量看的是數字加方向.
我們的禱告一定要朝對目標.
一定要有一個對的目的.
上帝賜給我們的禱告.
不是要我們拿回家切豬肉切蔬菜.
禱告是天國的屠龍刀與天劍.
用來打仗.
用來實踐天國使命.
所以禱告我們必須方向正確.
我們應該單求上帝的國.
上帝的義.
所以當我們說一無掛慮的時候.
不是說我們沒有掛慮.
而是我們將掛慮帶到上帝面前.
我們不是將所有的需要告訴上帝.
我們將自己在福音征戰這件事上的需要告訴上帝.
當這樣做的時候.
上帝的平安.
就會幫我們抵擋一切的掛慮.
恐懼和不安.
我們結束的時候再重溫三點.
因為我們是天國的公民.
天文雙稱福音.
天文後主榮耀.
我們的身份不是地上.
而是天上.
我們的家鄉不是這個世界.
而是將來的城.
我們的君王不是凱撒.
而是耶穌基督.
這個身份是他用他的血買回來的.
極其昂貴的恩典.
就是因為這個極其昂貴的恩典.
我們效榜樣防受敵.
施天使誇十價.
我們不再以地上的事為念.
不再被這個世界的價值觀捆綁.

$^{1001}$因為十字架的愛催迫著我們.
使我們不可能再為自己而活.
要單單為替我們死而復活的主而活.
也因為這份恩典和十字架.
在我們在地上等候救主在臨的日子.
尤其是在為福音去受苦.
在屬靈增薦的日子裡面.
讓我們能夠一同主近了.
同心靜.
無慮祈求永恆.
因為主近了.
所以我們可以放下掛慮.
將增薦中的需要帶到主面前.
因為主近了.
我們可以放下自己的權利.
謙讓待人.
因為主近了.
我們可以受苦.
爭戰當中.
靠主常常喜樂.
親愛的電視節目.
我們信息結束的時間.
我想問你一個最重要的問題.
你是不是天國的公民.
這個不是一個可以輕忽的問題.
我們其實每一個曾經都是基督十字架的受敵.
我們曾經都以地上的事為念.
以自己的羞辱為光榮.
但是上帝在我們還作罪人的時間.
差他的兒子為我們死.
耶穌基督在十字架上面用自己的血.
買贖一群子民.
歸給父上帝.
祂用自己的身體.
打開那條通往至聖所的路.
祂用自己的死亡.
敗壞掌死權的魔鬼.
今天祂邀請你和我.
再次回轉悔改.
不要再思念地上的事.

$^{1041}$成為天國的公民.
單單思念天上的事.
讓我們一同活一個與基督福音相稱的生命.
在救主再臨的日子裡面.
同心靜無類求.
讓我們的生命向這個世界宣告.
我們是屬於天國.
讓我們一同祈禱.
卻為我們虛己.
降卑.
信服至死.
且死在十字架上面.
你用你自己的血.
將我們從罪當中買贖回來.
使我們成為天國的公民.
主我們本來是十字架的受敵.
我們以地上的事為念.
但是你的愛催迫著我們.
叫我們不再為自己而活.
而是為替我們死而復活的主而活.
主我們已經近了.
求你幫助我們在爭戰當中一無掛慮.
單單為天上的事去祈禱.
幫助我們放下衝突.
同心的去見證福音.
幫助我們單單思念天上的事.
指跨你的十字架.
直到你再回來.
我們得著你榮耀的身體.
得成的冠冕那一天.
服務主耶穌基督的名祈禱.
阿們.
\newpage



\section{羅馬書 14:1-23}
\label{sec:031892}
\textbf{論斷,還是接納}~~~~郭建恩傳道
\newline
\newline
連結: \href{http://www.church.com.hk/acms/content.asp?site=cdc\&op=show\&type=product\&code=031892\&layout=sermon}{\texttt{031892}} ~~~~ 語音日期: 2026-03-29
\newline
\newline
\hyperref[sec:031891]{\small{< < < PREV SERMON < < <}}
~
\hyperref[sec:index_chronic]{\small{[返順時目]}}
~
\hyperref[sec:031893]{\small{> > > NEXT SERMON > > >}}
\newline
\newline
羅馬書 14:1-23
\newline
\begin{longtable}{cl}
\hline
\hline
章節 & 經文 (和合本修訂版)\\
\hline
14:1 & \begin{tabularx}{0.7\textwidth}{X} 信心軟弱的，你們要接納，不同的意見，不要爭論。 \end{tabularx} \\ \\ \relax
14:2 & \begin{tabularx}{0.7\textwidth}{X} 有人信甚麼都可吃；但那軟弱的，只吃蔬菜。 \end{tabularx} \\ \\ \relax
14:3 & \begin{tabularx}{0.7\textwidth}{X} 吃的人不可輕看不吃的人；不吃的人也不可評斷吃的人，因為神已經接納他了。 \end{tabularx} \\ \\ \relax
14:4 & \begin{tabularx}{0.7\textwidth}{X} 你是誰，竟評斷別人的僕人呢？他或站立或跌倒，自有他的主人在，而且他也必會站立，因為主能使他站穩。 \end{tabularx} \\ \\ \relax
14:5 & \begin{tabularx}{0.7\textwidth}{X} 有人看這日比那日強；有人看日日都是一樣。只是各人要在自己的心意上堅定。 \end{tabularx} \\ \\ \relax
14:6 & \begin{tabularx}{0.7\textwidth}{X} 守日子的人是為主守的。吃的人是為主吃的，因他感謝神；不吃的人是為主不吃的，他也感謝神。 \end{tabularx} \\ \\ \relax
14:7 & \begin{tabularx}{0.7\textwidth}{X} 我們沒有一個人為自己而活，也沒有一個人為自己而死。 \end{tabularx} \\ \\ \relax
14:8 & \begin{tabularx}{0.7\textwidth}{X} 我們若活，是為主而活；我們若死，是為主而死。所以，我們或死或活總是主的人。 \end{tabularx} \\ \\ \relax
14:9 & \begin{tabularx}{0.7\textwidth}{X} 為此，基督死了，又活了，為要作死人和活人的主。 \end{tabularx} \\ \\ \relax
14:10 & \begin{tabularx}{0.7\textwidth}{X} 可是你，你為甚麼評斷弟兄呢？你又為甚麼輕看弟兄呢？因我們都要站在神的審判臺前。 \end{tabularx} \\ \\ \relax
14:11 & \begin{tabularx}{0.7\textwidth}{X} 經上寫著：「主說，我指著我的永生起誓：萬膝必向我跪拜；萬口必稱頌神。」 \end{tabularx} \\ \\ \relax
14:12 & \begin{tabularx}{0.7\textwidth}{X} 這樣看來，我們各人一定要把自己的事在神面前交代。 \end{tabularx} \\ \\ \relax
14:13 & \begin{tabularx}{0.7\textwidth}{X} 所以，我們不可再彼此評斷，寧可決意不給弟兄放置障礙或絆腳石。 \end{tabularx} \\ \\ \relax
14:14 & \begin{tabularx}{0.7\textwidth}{X} 我憑著主耶穌確知深信，凡物本來沒有不潔淨的，除非人以為不潔淨的，在他就不潔淨了。 \end{tabularx} \\ \\ \relax
14:15 & \begin{tabularx}{0.7\textwidth}{X} 你若因食物使弟兄憂愁，就不是按著愛心行事。基督已經為他死，你不可因你的食物使他敗壞。 \end{tabularx} \\ \\ \relax
14:16 & \begin{tabularx}{0.7\textwidth}{X} 所以，不可讓你們的善被人毀謗。 \end{tabularx} \\ \\ \relax
14:17 & \begin{tabularx}{0.7\textwidth}{X} 因為神的國不在乎飲食，而在乎公義、和平及聖靈中的喜樂。 \end{tabularx} \\ \\ \relax
14:18 & \begin{tabularx}{0.7\textwidth}{X} 凡這樣服侍基督的，就為神所喜悅，又為人所讚許。 \end{tabularx} \\ \\ \relax
14:19 & \begin{tabularx}{0.7\textwidth}{X} 所以，我們務要追求和平與彼此造就的事。 \end{tabularx} \\ \\ \relax
14:20 & \begin{tabularx}{0.7\textwidth}{X} 不可因食物毀壞神的工作。一切都是潔淨的，但有人因食物使人跌倒，這在他就是惡了。 \end{tabularx} \\ \\ \relax
14:21 & \begin{tabularx}{0.7\textwidth}{X} 無論是吃肉是喝酒，是甚麼別的事，使弟兄跌倒，一概不做，才是善的。 \end{tabularx} \\ \\ \relax
14:22 & \begin{tabularx}{0.7\textwidth}{X} 你有信心，就要在神面前持守。人能在自己以為可行的事上不自責就有福了。 \end{tabularx} \\ \\ \relax
14:23 & \begin{tabularx}{0.7\textwidth}{X} 若有人疑惑而吃的，就被定罪，因為他吃不是出於信心。凡不出於信心的都是罪。 \end{tabularx} \\ \\
[1ex]
\hline
\hline
\end{longtable}
$^{1}$第一節會議,早安!.
今天講的題目是「論斷還是接納」.
我深信這個題目一點都不容易講.
因為論斷其實都是在人群當中時常都發生.
教會最容易發生的事情不是發生天壞等什麼大罪.
教會最容易發生的就是太多論斷的事情發生.
我相信大家都有很多經歷.
可能你有論斷別人的時候.
可能你有被論斷的時候.
對於我來說,我都不敢說我沒有論斷過別人.
可能我都有被論斷的時候.
這個課題都值得我們反思.
因為其實每一刻,每一詞.
每一刻都是提醒自己不要去論斷.
不要進入這個陷阱當中.
在這一年當中,其實我為什麼要講這些題目呢?.
其實我希望這一年都是在處理我們群體之間的關係.
我第一次講的時候是愛出發.
但這次我想講群體之間最容易破壞的.
不單是沒有愛,而是我們很容易彼此的論斷.
我們一起來祈禱吧.
師父我求你來到當中與我們眾人同在.
聖靈求你關照我們的內心.
讓我們能夠看見我們的生命.
經常有很多不滿或者很容易講出來的論斷.
去傷害了人,傷害主人.
主人你幫助我們在這個議題當中.
讓我們能夠去反思.
知道我們如何在教會當中建立一個有愛有真理的教會.
求主來到當中幫助我們.
幫助孩子所講的,幫助姐妹所聽的.
去同樣的力量,我們感到很高興.
我們求主保佑,使我們能夠祈禱.
剛才主席用了很長的篇幅讀這段經文.
事實上都需要長一長.
讓我們知道在這段經文當中.
不是一個故事,而是一個教導在當中.
羅馬書的背景就是在羅馬時代.
即是在羅馬的教會當中.
其實羅馬城當中有很多猶太人.

$^{41}$猶太人當中彼此有些爭議.
有些信了耶穌的猶太人.
有些未信耶穌的猶太人.
彼此爭議究竟耶穌是誰.
引來很大的風波.
所以在歷史上當中.
羅馬皇帝格魯丁覺得太麻煩了.
這班猶太人乾脆將這些猶太人趕出羅馬城.
不要再在羅馬城搞來搞去.
所以不論你信了耶穌的猶太人.
或者未信的猶太人.
你都趕出去.
在主前49年的時候發生了.
所以當猶太人離開了羅馬城的時候.
你試想像羅馬城的教會剩下什麼人呢?.
沒有了猶太人.
信了耶穌的基督徒.
反而是剩下所有的外邦人.
在當中帶領著教會.
直到主後54年.
這班基督徒的猶太人.
能夠回到羅馬城的時候.
發現教會不同了.
以前我們所見到的教會.
五年後發現了.
現在是一班外邦人主導的教會之餘.
發現很多生活的習慣.
文化完全完全不同.
其實可以試想像.
回流的弟兄姊妹.
十年前來悉尼.
十年後回來的時候.
發覺城市很不同.
連教會都很不同.
信了耶穌的猶太人回來之後.
為什麼整個教會不同了呢?.
有什麼不同呢?.
在第二節說到.
有人信什麼都可以吃.
但軟弱的就說只可以吃蔬菜.

$^{81}$有人說不要輕看.
有些人就輕看不吃的人.
吃的人.
不吃的人.
不可以論吃的人.
因為神已經收納他們.
是在說什麼呢?.
是在說有些人覺得.
不要隨便吃東西.
吃東西很容易得罪上帝.
有些人只會吃蔬菜.
另外一個議題就是.
那班猶太人回來的時候.
覺得有些日子是特別的.
有些特質.
每天都是一樣的.
在這個議題上.
很多的衝突.
吃什麼.
日期是什麼.
都是衝突當中.
究竟是什麼衝突呢?.
一班猶太基督徒.
長期守摩西律法.
尤其是注重利未忌.
或者生命忌.
當你看到利未忌和生命忌.
覺得很悶的時候.
哪些可以吃.
哪些不可以吃.
究竟是什麼呢?.
就算我初信主的時候.
你有沒有疑問.
究竟豬籠可以吃嗎?.
其實我很喜歡吃豬籠.
都很多疑問在當中.
一班猶太人說.
摩西律法這些東西.
不可以胡亂吃的.
這些吃了的話.

$^{121}$其實影響我們敬拜上帝的.
所以猶太人的文化和習慣改變了.
看到教會改變了的時候.
他覺得外邦人太過自由了.
太過不敬虔了.
這樣的狀況之下.
來論斷一班外邦的信徒.
但是外邦人的信徒是怎麼樣呢?.
過往他們沒有律法背景的包袱.
所以很容易接受耶穌基督.
在裡面很自由.
對於猶太人的看法就是.
你這班人太過迂腐.
太過守律法在裡面.
形形成成就覺得.
你們信心這麼差.
我們信耶穌應該自由的.
什麼都可以吃的.
不應該這樣的.
所以輕看他們的行為.
保羅怎麼教導他們呢?.
保羅說什麼呢?.
憑著主耶穌基督的心信.
什麼都是沒有不潔的.
唯獨人以為不潔的就是不潔.
說的是什麼呢?.
真正的不潔.
不是食物那裡.
吃吃豬肉都可以的.
而是基督.
你怎麼依靠基督.
基督已經成為那個律法.
所以我們每個人已經在耶穌裡面被接納.
保羅說在耶穌裡面.
沒有東西是不潔的.
問題不是在於食物.
問題不是在於那個節旗.
問題是在於你怎麼看待對方.
所以今天來到當中.
我想用三個問題來高呼.

$^{161}$整個我今天講的講章.
在《馬太福音》第七章有登山寶訓.
《路加福音》第六章有平原補訓.
在這兩個補訓當中.
都不斷講不要論斷.
而《羅馬書》第二章都講不要論斷.
今天來到第十四章的時候.
我們看看如何看待論斷是什麼呢?.
為什麼我們很容易論斷別人呢?.
論斷真正的問題是什麼呢?.
而同時福音怎麼在改變我們呢?.
為什麼我們會論斷呢?.
第一件事是屬靈的優越感.
當時我們經常覺得.
很多時候是跟別人比較.
你這麼不衛生.
你這樣做事的.
你回來教這麼久.
你都不熟悉聖經的.
你這樣服侍.
不是行的.
這個話是.
我們內心裡面.
很多的評價.
對人的評語都在當中.
我們很容易看低別人.
這些評價.
我在這裡深講.
你都不要承諾.
你的靈性很淺.
你信耶穌很淺.
或者你不成熟.
我形容這些人是暗爽自嗨.
為什麼呢?.
通常你說別人的時候.
你覺得你不熟悉靈.
不讀聖經.
說你不行的時候.
其實你內心在說.
我至少不像你.

$^{201}$我比你好.
我比你熟悉靈.
我是不知道的.
你裡面內心的狀況.
告訴自己.
我比你好.
我勁過你.
我怎樣都有價值過你.
所以.
在這個論斷的過程當中.
其實是在說.
在說自嗨的時候.
其實好像聖經故事一樣.
耶穌講的故事就是.
一個法利賽人和一個瑞利的祈禱.
有兩個人在電的祈禱的時候.
法利賽人.
一個是瑞利.
法利賽人站在街口自然說.
神啊.
不像這些人好像勒索一樣.
或者不義.
或者姦淫.
不像那些瑞利一樣.
我每個星期都有兩次.
來禁食.
我還要有十分一個奉獻在當中.
其實當我們不斷去論斷的時候.
其實不是在說.
我們有多麼勁虔.
是直接比較.
你和別人比較.
你不行.
而我覺得我行.
在裡面的心靈當中.
不單純是在說你勁不勁.
你直接比較.
告訴別人.
告訴自己.
我行.

$^{241}$我很厲害.
我很熟練.
我有價值.
透過踩低別人.
來對比.
換來自己覺得自己好一點.
這個就是論斷.
其中一個要處理的問題.
覺得自己聖人一籌.
我不像你.
這些全部都在裡面.
內心當中.
發生在裡面.
沒有人知道.
是你自己知道.
是上帝都知道.
另一種.
就是什麼呢.
就是那種.
自我價值低.
或者白痰齋主.
內心是什麼呢.
雖然他們不會覺得我比你強.
我不會覺得好像比你強.
但是裡面這一類的人.
通常是怎麼樣呢.
就是很喜歡.
先讓我看這裡.
這裡就是.
他們這類就不是說.
剛才說的我比你強.
但是這類型的朋友.
就是怎麼說呢.
他怕自己被否定.
自己覺得自己裡面.
很容易怕被人說出自己的問題.
於是不斷指控別人的問題.
不斷批評.
不斷很容易敏感一些事情.
很容易挑錯處.

$^{281}$有時候覺得不是驕傲的問題.
而是內心有很多不安.
所以你看得到.
有些人不斷左彈右彈的時候.
你不要被他表面壓倒.
他是白痰齋主.
而他裡面也不安穩.
在裡面很害怕.
很害怕被人指出自己的問題.
所以來批判.
而遮蓋了裡面自己的問題.
背後的東西是什麼呢.
很多時候.
他不是關注那個真相.
他不是關注那個問題.
而是很想別人直接批評.
讓自己好過一些.
是別人不好.
人家差 比我還差.
指出對方的時候.
就算專注在對方.
也不在自己身上.
就好像在馬大馬利亞一樣.
馬大馬利亞不斷做事.
我們上次在節目說.
心裡煩擾煩擾什麼呢.
不單是沒有看著上帝.
而是覺得 咦 為什麼我不斷做事的時候.
馬大馬利亞不斷做事的時候.
馬利亞為什麼這麼好.
就是坐在耶穌身上.
而什麼都沒有做.
而我在這裡做這麼多事.
有沒有人看見我呢.
很多時候就是我們內心沒有平安.
很多的衝突在裡面.
為什麼我做這麼多事.
而他沒有做事呢.
為什麼這樣都可以呢.
為什麼這麼不公平呢.

$^{321}$整天不斷左彈右彈.
其實自己裡面很多不安穩.
在裡面.
究竟真正的問題.
是什麼呢.
聖經告訴我們.
不單是內心的問題.
而是一個屬靈上的問題.
在裡面的當中.
你坐了上帝的位.
當你批判別人的時候.
在論斷的時候.
你坐上帝的位.
亞國書說什麼呢.
設立律法和論斷的人只有一位.
這個就是能救人和滅人的位.
你是誰.
究竟論斷別人.
你坐上帝的位.
那個什麼.
能夠救人和滅人的那個.
是在說耶穌自己.
當我們.
論斷別人的頓基的時候.
你想定他的忠告.
他是這樣的.
我明白了.
我認識他這麼久.
他沒有救了.
我認識他這麼久.
已經定了結局.
已經定了營.
彷彿我們正在做上帝正在做的事情.
如果.
耶穌今天.
承擔了.
我們的罪.
來接受審判的時候.
為什麼我們.
依然能繼續.

$^{361}$急著來審判別人呢.
今天耶穌已經承受了審判.
為我們承受審判.
為什麼我們要急著.
來審判別人呢.
這個會不會代表.
你不相信十字架.
那個拯救呢.
這個會不會代表.
你不相信十字架.
已經完成了審判.
他為你承擔了一切呢.
我們一切過犯已經被耶穌.
擔當了一切.
福音說什麼呢.
福音說.
我們已經知道一切.
神是公義的 神是斯聯文的.
所以我們可以放下.
這一切的審判權.
你這個位置是治癒的.
當我們論斷的時候.
我們內心說我比你好.
我比你厲害.
我比你更熟悉.
在一個比較當中.
當我們論斷的時候.
在羅馬書三章說.
世人都犯了再虧缺了神的榮耀.
你要知道.
在世上沒有一個是義人.
包括我和你一樣不是一個義人.
世界每個人都在犯罪.
得罪上帝.
但是我們經常透過踩低別人的時候.
感覺自己很好.
這個是世界給我們的價值觀.
可能是職場好.
家庭的家族關係也好.
在教會關係也好.

$^{401}$我們直接踩低別人的時候.
我比你好.
我比你幸福.
但是我們這樣的心態.
從來沒有經歷過.
福音的能力在當中.
論斷不是在說.
那個人有多壞.
是在說我有多好.
是在證明我有多好.
不是在說那個人有多差.
是在說我有多好.
換句話說.
我們不斷在質疑自己.
在浪子的比喻當中.
在浪子的故事當中.
小兒子.
你知道.
爸爸還沒死.
就問爸爸要家產.
離開了之後.
走這麼遠.
去到一個地方.
花他的錢.
不再影響爸爸.
小兒子回來的時候.
大兒子怎麼想呢.
爸爸.
這個小兒子.
還沒死就問你要家產.
現在花光所有錢回來.
做些不道德的事情.
你這樣對他.
為什麼你和他宰牛.
為什麼你給他戒指.
為什麼你給他衣服.
為什麼你從來沒有為我宰牛.
為什麼你從來沒有給我戒指.
小兒子回來的時候.
你這樣對他.

$^{441}$我做了很多事情.
我為你做了很多事情.
為什麼你對他這麼好.
而你從來沒有為我宰牛.
給我衣服和戒指.
大兒子表面上很馴服.
但心裡很憤怒.
所以他不出席晚宴.
和小兒子的晚宴.
他不太強.
他弟弟這樣也可以回來.
沒理由.
就像大兒子一樣.
我比弟弟好.
我做的事情更多.
我有優越感.
這就是論斷的根.
福音怎麼改變我們呢.
如果我們的生命.
有熟練的驕傲.
有自義.
或者坐上了上帝的位置.
究竟我們信的福音.
怎麼能夠改變我們的生命呢.
很多時候.
我們對於不同意見和文化的人.
我們會容忍.
容忍是什麼意思呢.
我不理你.
你不要惹我.
我們保持距離.
但從來沒有處理過關係.
保羅在.
羅馬書十四章一至三節.
和十五章七節.
不斷說要接納.
信心軟弱的人.
這是新學語.
不同於剛才的版本.
第三節說.

$^{481}$因此你們要彼此接納.
如同基督接納你們一樣.
飽榮耀歸給神.
第十五章七節說.
你們要彼此接納.
如同基督接納你們一樣.
不斷說接納.
接納是什麼意思呢.
接納在原文當中.
Post Lambeno.
意思是分開兩個字.
向前的意思.
而Lambeno.
意思是接受.
向前的接受.
當我們面對不同的意見.
不同的文化.
或者人的時候.
我們怎麼解釋呢.
接納的意思是.
主動地將對方帶到你身邊.
歡迎對方進入.
關係當中.
我們面對很多不同的文化.
面對很多的意見.
很多的事情.
我們不是容忍.
保持距離.
而是接納對方.
將對方放進.
你的關係當中.
當然接納不代表.
認同對方所做的事情.
可能對方有些事情做得不對.
我們接納.
將他帶入你的關係.
但不是要認同他所做的事情.
撒該的故事一樣.
撒該的故事就是.
大家都知道他是瑞麗.

$^{521}$大家都知道他為羅馬.
來到當中.
收取不義的錢財.
在當中.
但耶穌來到他那裡.
撒該 撒該.
今天我要來到你家.
和你吃飯.
耶穌接納罪人.
接納撒該來到他生命當中.
進入他家和他吃飯.
接著撒該說.
主啊 我知道我做錯了.
我願意將所有東西賠償給.
我敲詐我的同胞的人.
耶穌沒有認同.
撒該所做的事情.
不過將這個罪人.
帶到他生命當中.
以至生命當中一步一步的和好.
接納.
不是因為對方直.
也不是因為彼此之間的關係成熟.
而去接納.
而是因為耶穌基督已經接納了我們每一個人.
所以我們有能力去接納.
接納好像耶穌一樣.
將罪人帶到.
他的家一樣.
我們將一個不可愛的人.
當中而進入他的生命.
去接納他的本相.
今天耶穌在我們已經完完全全.
接納了我們每一個人.
保羅一樣.
保羅昔日.
去接納猶太人的教會.
猶太神的教會.
接著當他信了耶穌的時候.
在二萬五時.

$^{561}$他信了耶穌之後.
眼睛開了.
沒有人可以將他.
所有基督徒.
都不敢接觸保羅.
這個人以前迫不得已.
以前這個人殘害教會.
為什麼可以接納他.
唯獨是一個巴拿巴.
他願意來將接觸保羅.
介紹所有領袖.
讓他認識.
重拾和弟兄姊妹的關係.
後來更加成為一個宣教士.
門徒.
能夠再次.
認識耶穌基督.
或者再次將方傳遍.
是因為巴拿巴將他帶到使徒面前.
今天耶穌基督已經.
接納了你每一個人.
你又怎樣接納.
其他身邊的人呢.
第二.
除了學接納.
你還要學放下你的審判權.
我覺得保羅很精彩.
文字上說你是誰.
論斷別人的僕人.
你是他站在.
跌倒自有他主人在.
而且他必要站住.
因為主能使他人站住.
當我讀的時候.
我覺得保羅正在爬開.
你是誰.
你竟然論斷別人的僕人.
第一件事.
就是你問你是誰.
今天.

$^{601}$這裡說了什麼呢.
那個人已經是別人的僕人.
那個是主的僕人.
你是誰在論斷別人的僕人呢.
主已經寬恕他.
主已經為他承擔一切.
為什麼你還要論斷他呢.
但同時你都要問我是誰.
今時今日.
我們都是主的僕人.
主已經為我們承擔一切.
他接納了我們.
我們完全可以接納.
論斷呢.
其實是坐在上帝的位置.
剛才說了我們坐在上帝的位置.
當你論斷的時候你坐在上帝的位置.
你覺得你自己有能力.
你覺得你什麼都知道.
你覺得你公義.
你覺得你是標準.
所以你在論斷別人.
坐在上帝位置當中.
你認清楚你是誰.
你不是上帝.
在這個通姦的故事當中.
那個婦人被人通姦.
把那個婦人拉出去城口.
所有人來到當中準備扔死她.
耶穌在畫一個圈.
一邊這樣說.
你們覺得誰沒有罪的話.
你就拿石頭扔他.
耶穌在圈啊圈啊.
現在人都不見了.
沒有人去拿石頭.
扔死那個奸人的婦人.
因為知道自己也是一個罪人.
最後耶穌對那個婦人說什麼呢?.
他說今天我也不定你的罪.

$^{641}$你要悔改.
弟兄姊妹.
今天我們的弟兄姊妹.
我們的弟兄姊妹.
已經耶穌都不去論斷.
不去審判的時候.
為什麼我們還要來論斷審判我們的弟兄姊妹呢?.
我們要放下我們審判的權利.
在裡面當中.
最後.
保羅在這裡說.
論斷帶來什麼關係呢?.
帶來什麼破壞呢?.
你不可以因為聖物而叫它敗壞.
無論你要吃肉也好.
喝酒也好.
做什麼事也好.
叫弟兄姊妹一概都不可以做.
論斷是在破壞弟兄姊妹的關係.
你有沒有發覺.
教會裡面有很多賤客.
很多賤客.
什麼意思呢?.
你一見到一些不公義的時候.
那些賤客就出刀.
用刀指向對方.
我們自己以為有公義.
我們把公義的刀拿出來.
指向對方的時候.
你做得不好 你有問題.
你驕傲.
一拿刀指向出來的時候.
其實當我們論斷的時候是在破壞弟兄姊妹的關係.
有時候.
當弟兄姊妹見到一個賤客的時候.
他不是見到你在幫忙.
他見到你的情緒.
或者見到你不是很體恤.
他見到你是高高在上.
他感受不到你的破壞.

$^{681}$感受到你的幫助.
反而令到他跌倒.
我曾經勸一個弟兄姊妹回來.
她很久沒有回崇拜了.
我跟她說 你回來教會吧.
你整年都沒有回教會.
我關心你的狀況.
關心你的家庭.
我仍然相信上帝.
我整年都相信上帝.
不過你知道.
教會有很多人說話.
我怕回到教會之後.
我講一些不中聽的話.
很多人會跟我說話.
我受不住.
教會是一個論斷的地方.
我們有真理的時候是要有真理.
但我們怎樣運用我們的真理.
這才是我們要學的地方.
保羅最主要講什麼呢.
你要按愛人的道理.
去跟人相處.
神的話.
神的天國不是在乎吃喝.
是在乎公義和平.
和聖靈的喜樂.
但是你怎樣跟弟兄姊妹.
和好的關係一起同行呢.
你務必追求和睦的事.
和彼此建立德行的事.
問題不是在講你對不對.
問題是你怎樣看弟兄姊妹.
不是你想做什麼.
而是為對方做什麼.
以至不要傷害對方.
要對方有成長.
你不是看自己的公義.
是看對方的生命.
不是硬生生把道理告訴別人.

$^{721}$他要遇見基督.
我呢.
做牧者有時.
有些豬頭骨都找牧者做.
牧者你講吧.
牧者你處理吧.
你跟進一下吧.
事實上我們最需要.
作為牧者.
我怎樣去牧養弟兄姊妹呢.
是否一開始就將所有道理.
你不對你做得不好.
你這樣相處不好.
不是.
很多時候.
對於我來說我都為弟兄姊妹祈禱.
我想怎樣幫助弟兄姊妹.
聽得到我講的話.
當.
如果我這麼attentive.
他跟我說話的時候.
弟兄姊妹聽不到.
弟兄姊妹期望你關心 同行.
以至聽得到上帝的話語.
來叫他成長.
我都為弟兄姊妹祈禱.
到差不多的時候.
就約他吃飯聊天.
吃飯聊天的時候.
不是一開始就說你錯.
你試著明白他背後做的動作的時候.
為何他有一個動作.
為何他講這堆說話.
為何他有這樣的反應.
從中調節.
或者怎樣跟上帝的愛.
或者話語連在一起.
當我這樣跟弟兄姊妹聊的時候.
他反而不是.
傷害到我.

$^{761}$而是知道我跟他同行.
知道真愛在哪裡.
心靈慢慢轉變過來.
弟兄姊妹.
教會不是要是非多.
教會是要有愛在裡面.
我知道我們很多團契.
很多小組.
都有很多的互動在裡面.
這是必然發生的事情.
但當你面對不同意見的時候.
你怎樣實踐真理呢.
怎樣用愛.
Walking in love.
在英文的版本當中.
我們怎樣用愛心來實踐呢.
求主幫助.
我們每一個弟兄姊妹.
希望我們在小組團契當中.
是一個有愛的小組.
一個有真理的小組.
我們一起禱告.
天父求你赦免我們的過犯.
因為我們經常.
有很多的比較.
對人有很多的比較.
或者我們經常.
有很多不安.
以致我們經常對.
身邊的人,我們的家人.
我們的弟兄姊妹,甚至身邊的朋友.
都作出很多的論斷.
主啊.
我們.
來到面前求寬恕.
你先是我們的主,你先是我們的王.
今天你已經.
接納了我們的時候.
讓我們再次.
仰望你.

$^{801}$你是有審判的權利.
而不是我們自己.
讓我們學會放下.
安堂在你的懷裡.
當中.
被你的愛所充滿.
主啊求你給我們.
生命當中有那份.
接納的能力.
縱然很多人.
都是對我們不同的意見.
或者傷害我們.
你怎樣接納我們.
進入恩典當中的時候.
主啊求你給我們那份能力.
去接納身邊的人.
進入一個關係當中.
主啊我們很想.
很想再次.
感受上帝給我們的那份.
愛和恩典.
尤其是和頂尖的那份關愛.
和.
在團契生活當中.
求主幫助我們.
讓我們的生命.
實踐內,見證真理.
主啊多謝你.
再次恢復我們每一個人.
的生命.
我們這樣禱告.
給我們救主寶貴耶穌文字祈禱.
\newpage



\section{馬太福音 21:1-9}
\label{sec:031893}
\textbf{謙和的君王 (華)}~~~~葛國光牧師
\newline
\newline
連結: \href{http://www.church.com.hk/acms/content.asp?site=cdc\&op=show\&type=product\&code=031893\&layout=sermon}{\texttt{031893}} ~~~~ 語音日期: 2026-03-29
\newline
\newline
\hyperref[sec:031892]{\small{< < < PREV SERMON < < <}}
~
\hyperref[sec:index_chronic]{\small{[返順時目]}}
~
\hyperref[sec:031894]{\small{> > > NEXT SERMON > > >}}
\newline
\newline
馬太福音 21:1-9
\newline
\begin{longtable}{cl}
\hline
\hline
章節 & 經文 (和合本修訂版)\\
\hline
21:1 & \begin{tabularx}{0.7\textwidth}{X} 耶穌和門徒快到耶路撒冷，進了橄欖山的伯法其時，打發兩個門徒， \end{tabularx} \\ \\ \relax
21:2 & \begin{tabularx}{0.7\textwidth}{X} 對他們說：「你們往對面村子裡去，會立刻看見一匹驢拴在那裡，還有驢駒同在一處，解開牠們，牽到我這裡來。 \end{tabularx} \\ \\ \relax
21:3 & \begin{tabularx}{0.7\textwidth}{X} 若有人對你們說甚麼，你們就說：『主要用牠們。』那人會立刻讓你們牽來。」 \end{tabularx} \\ \\ \relax
21:4 & \begin{tabularx}{0.7\textwidth}{X} 這事發生是要應驗先知所說的話： \end{tabularx} \\ \\ \relax
21:5 & \begin{tabularx}{0.7\textwidth}{X} 「要對錫安的兒女說：看哪，你的王來到你這裡，謙和地騎著驢，騎著小驢—驢的駒子。」 \end{tabularx} \\ \\ \relax
21:6 & \begin{tabularx}{0.7\textwidth}{X} 門徒就照耶穌所吩咐的去做， \end{tabularx} \\ \\ \relax
21:7 & \begin{tabularx}{0.7\textwidth}{X} 牽了驢和驢駒來，把他們的衣服搭在上面，耶穌就騎上。 \end{tabularx} \\ \\ \relax
21:8 & \begin{tabularx}{0.7\textwidth}{X} 許許多多的人把自己的衣服鋪在路上，還有人砍下樹枝來鋪在路上。 \end{tabularx} \\ \\ \relax
21:9 & \begin{tabularx}{0.7\textwidth}{X} 前呼後擁的人群喊著說：「和散那歸於大衛之子！奉主名來的是應當稱頌的！至高無上的，和散那！」 \end{tabularx} \\ \\
[1ex]
\hline
\hline
\end{longtable}
$^{1}$請問今天是什麼日子?.
啊?復活節啊?.
啊?.
今天是什麼日子?.
今天在教會的上面,我們稱它叫做Palm Sunday.
Palm Sunday就是所謂的忠書之,忠知主日.
那這個日子是怎麼來的呢?.
你可以看見,在馬太夫人剛纔所讀的經文裡面.
還有其他的四福音裡面都有記載.
也就是說,今天是耶穌進耶路撒冷.
然後眾人就是用忠書之來歡迎祂的這個時候.
等一下我們會詳細講.
那麼到了禮拜四的時候.
耶穌為門徒洗腳.
然後跟門徒們一同進最後的晚餐.
到了禮拜五呢,也就是我們知道的.
耶穌為我們訂十字架死了.
那麼我們稱它叫做壽難日.
也就是我們禮拜五的壽難節的聚會.
第三天之後,耶穌從死裡復活.
那個復活呢,叫做七日的頭一日.
也就是下一個禮拜天.
所以下禮拜天就是復活日.
那整個從這個禮拜天一直到下個禮拜天.
稱它叫做Holy Week,就是聖周.
那這一週所發生的事情.
在聖經裡面,四福音裡面佔了非常重要的部分.
馬太福音,馬可福音,路加福音,約翰福音.
都在記載花了很多的篇幅.
馬太用了四分之一,馬可用了三分之一.
路加用了四分之一.
約翰福音用了幾乎二分之一的篇幅.
在記載耶穌受難的最後這一週的事情.
所以可以說這最後的一週.
就是我們整個信仰的核心.
所以今天我們叫做聖中之主日.
有看過中之嗎?.
也許不知道在美國到哪邊去找.
我曾經有一年想要去找.
發現很難找到.

$^{41}$實際地拿出來給大家看.
所以耶穌在今天,就是這個主日.
當他進耶路撒冷的時候.
我們來看看剛纔這段經文.
在這段經文裡面,耶穌進去.
然後你看馬太福音21章第五節怎麼說.
他說:看啊,你的王來到你這裡.
是溫柔的,又騎著驢,就是騎著驢鞠子.
所以耶穌騎著驢鞠子.
這個驢鞠子在路加福音說.
是從來沒有人騎過的.
表示怎麼樣?.
表示這個驢鞠子,它是一個年幼的.
它從來沒有去做過.
從來沒有去被人家騎過的.
它是一個年輕的,一個年幼的.
代表著這是一個什麼?就像聖經上所說的.
這是最好的,出手的,沒有被用過的.
親愛的弟兄妹妹.
耶穌他進城的時候,進耶路撒冷城的時候.
眾人歡呼他,好像是君王一樣的在歡呼他.
可是弟兄妹妹你知道嗎?.
世界上的君王都是騎什麼東西啊?.
你有沒有看過君王騎驢鞠子的,有沒有?.
成吉思汗騎個驢,打遍天下,有沒有?.
沒有,你沒看過.
所以世上的君王他們都是什麼?.
都是騎著戰馬,是不是?.
而且他的戰馬絕對是那個訓練.
是長得最強壯最好的,跑得最快的.
所以世上的君王是騎著戰馬,是騎著馬進城的.
當年羅馬皇帝奧斯都進城的時候.
打勝仗進城的時候.
這個就是在形容他的當時的盛況.
哇!你可以看見眾人在那邊排隊.
在那邊歡呼擁帶他進來.
他騎著那個馬帶著他的眾這個百姓.
帶著他的軍隊一起進城的時候.
把戰利品把他都帶進來.
那是何等光輝榮耀的一個時候,是不是?.

$^{81}$那親愛的同志們,你對照來看.
世界上面的君王都是這樣子威風的.
可是耶穌基督,聖經上告訴我們說.
他是騎著驢雞子.
在撒迦利亞書講這一段聖經的時候.
引用撒迦利亞書的時候,我們來看看.
他說你的王是怎麼樣?.
是公義的,並且施行拯救.
千千和和的,怎麼樣?.
騎著驢.
但是馬太福音纔講到耶穌的時候.
他沒講到千和對不對?他講他是什麼?.
他說看到你的王來到這裡,他是什麼?.
溫柔的.
你有沒有在形容哪一個王是溫柔的?.
哪一個王是千和的?沒有.
你看見世上的君王.
當你想到成吉思汗,我記得我們到希臘去.
我們到了一個新約之旅的時候到希臘.
看到那個成吉思汗.
那個雕像,我們當然沒看到本人.
看到他雕像的時候就是那個很雄赳赳氣昂昂的那樣子.
丟著唄,你看到世上的君王都是這樣子的威風.
他們是那樣子的勇猛,雄壯.
他們甚至可以說是兇暴.
但是聖經裡面讓我們看見耶穌他進耶路撒冷城的時候.
他是千千和和的騎著驢居子.
他是溫柔的.
最近我們.
我想我們大概這一個月都在關注.
不管你喜歡不喜歡.
你都會關注.
為什麼呢?.
因為從猶加的一直上升你就可以感覺到.
怎麼樣?.
你希望這個場戰爭繼續打下去嗎?.
我想我們大家都希望早點結束,是不是?.
現在丟著唄.
你看見整個伊朗.
就讓我們看到整個伊斯蘭教的世界.

$^{121}$你要瞭解到.
伊斯蘭教世界.
整個的發展的歷史.
你就會發現.
伊斯蘭教不能夠小看它.
當年默罕默德.
他興起來的時候.
他說他得到了所謂的真主的啟示.
然後他就起來傳教了.
那麼他傳教.
你要知道當時候的中東的地區.
阿拉伯民族.
其實是拜偶像的.
所以他當然受到這些拜偶像的這些其他的族的反對.
對不對?.
所以就起了衝突.
就他大起張來.
他當時候帶著他的穆斯林.
打遍天下無敵手.
把整個阿拉伯半島統一起來.
632年.
當默罕默德去世之後.
接替他的這些王朝的人.
他們繼續地征服.
一直到了大概 之後的大概77年的時間.
伊斯蘭教把當時候的東羅馬帝國.
也就是說拜佔壇帝國.
還有波斯帝國.
就是今天的伊朗.
全部打敗.
打敗之後.
把他們征服下來之後.
就一統那些地方.
全部都歸信伊斯蘭教.
接著又到北非去.
把這個北非.
所以你現在看.
撒哈拉沙漠以北的地方.
就是北非的這些地方.
全都是被伊斯蘭教征服了.

$^{161}$所以默罕默德.
他是一個教主.
他也是一個君王.
他統一了阿拉伯這個王國.
當年默罕默德.
他帶著他所謂.
他得到的真主啟示的可蘭經.
同時也藉著他手上的這一個刀兵.
把整個當時候甚至基督教所在的地方.
你今天有機會到土耳其去看.
你去看伊斯坦堡.
以前叫做君士坦丁堡.
那曾經是君士坦丁.
遷都到那個地方.
是當時候可以說的首都.
可是你到那邊去看.
聖索菲爾那個大教堂.
你還可以看到.
當年的那個東正教的那個教堂的.
在那個地方.
可是他現在已經變成是清真寺了.
所以當這個整個的土耳其.
這個伊斯坦堡被拿下來之後.
整個的阿拉伯半島就被他統一起來了.
像劉子美,穆罕默德是帶著戰車馬兵.
征服了當時候的世界.
建立了今天可以說是橫跨三洲的.
一個伊斯蘭教的一個帝國.
但是耶穌基督呢.
你知道現在全世界最多的.
以宗教來講最多的是基督教.
當然現在一般講基督教的時候.
把這個天主教,天主教,東正教.
還有基督教統統撐在一起.
反正就是說信耶穌基督的統統加在一起.
現在全世界大概有二十六,七億的人.
是所謂的基督徒.
其次呢.
伊斯蘭教.
伊斯蘭教大概有十七,十八億.

$^{201}$你慢慢會感覺到.
現在你到歐洲去.
你已經漸漸感覺到.
伊斯蘭教的勢力.
穆斯林在那邊越來越多,越來越多.
到有一天很可能.
歐洲變成是伊斯蘭教化.
為什麼?.
因為他們拼命地生.
歐洲人拼命不生.
所以到最後一個結果就是.
從生產率來講.
他們一定就是整個歐洲就變成是穆斯林化的.
親愛的丟子妹.
我在這邊要說什麼?.
默罕默德藉著武力.
來傳播他的信仰.
雖然得罩了這麼多的人.
十七,八億的人今天都信了.
但是親愛的丟子妹.
耶穌基督呢.
他沒有廢一兵一卒.
他靠什麼來征服這世界上二十七億的人?.
請告訴我.
耶穌是進入耶路撒冷之後.
接下去他就要準備要上十字架了.
所以耶穌基督藉著愛.
藉著釘十字架.
藉著犧牲他自己.
他是一個和平的君王.
他不是要來製造戰爭的.
他不是要來打仗的.
因為他說我的國不屬這個世界.
親愛的丟子妹.
我講這個是因為你要讓我們看見.
因為今天整個世界.
包括今天美國的走向.
好像我們不知不覺的感覺.
誰的拳頭大.
誰就聲音大.

$^{241}$是不是?.
誰的兵力比較強.
我們就可以不怕人家來欺負我們.
所以現在全世界大家都在發展軍力.
為什麼?.
日本也在發展它軍隊.
為什麼?.
德國現在也在發展軍力.
為什麼?.
因為好像感覺這個世界一個走向就是.
如果我沒有軍力.
我沒有這個軍隊.
我的軍隊不夠強盛.
我總有一天會像烏克蘭一樣被人家欺負.
在這樣子一個強權文化之下.
丟子妹.
我擔心的是我們基督徒也會不知不覺的被影響.
瞭解我意思嗎?.
耶穌是一位.
謙卑的君王.
所以祂要祂的子民也要謙卑.
Amen?.
在聖經裡面講到謙卑.
很多時候我們來看看.
在所謂的八福天國的信章裡面.
耶穌講到虛心的人有福了.
什麼叫虛心的人?.
虛心的人,原文的意思就是.
淋漓貧窮的人.
他覺得他自己很貧窮.
他很需要.
他裡面好像覺得有一種飢渴.
我很缺乏,我需要.
這叫做虛心的人.
那麼馬太福音第五節又說.
溫柔的人怎麼樣?有福了.
什麼叫做溫柔啊?.
講話嗲裡嗲去的.
動作就是很撒嬌的.
這個叫做溫柔,是不是?.

$^{281}$丟什麼?.
彼得前書第三章第一到第五節.
我這次把這一節把它說過來.
他提到說.
你們做妻子的要順服自己的丈夫的時候.
他提到說不要以外面的這些編法.
這些美裝飾.
而要以內心裡面存著長久溫柔.
安靜的心為裝飾.
所以當正經裡面在講到說.
要長久溫柔的心的時候.
這裡講的是做婦女的.
那請問做弟兄的需不需要溫柔?.
請問巴福是跟誰說的?.
跟姐妹說的.
虛心的人有福了.
溫柔的人有福了.
因為姐妹你們溫柔你們就有福了.
是不是?請告訴我.
請問巴福耶穌對著誰講的?.
是對姐妹講的嗎?.
他對每一個弟兄姊妹都是對我們講的.
換句話說.
巴福裡面講到的就是一個得勝者的品格.
你如果要承受地土.
承受地土就是與主同坐王的掌權的意思.
如果你要做一個得勝者.
將來與主一同坐王.
他說條件是什麼?.
你要做一個溫柔的人.
那好了.
我們今天把謙卑跟溫柔做了一個解釋.
換句話說.
謙卑就是溫柔.
它們兩個是不分家的.
可能謙卑我們比較容易瞭解.
是不是?.
溫柔我們常常覺得.
溫柔那是女性的.
所以我記得多年前.

$^{321}$聽到北加州的江秀琴牧師.
她提到說.
她說在她剛去他們教會的時候.
她要所有的姐妹們.
都來上她教導的順福的這一堂課.
好像是大概是有六講還是七講課程.
但是所有的弟兄.
她要他們來上什麼課?.
謙卑.
那謙卑她講的那個課程.
比這個順福還要多.
這表示是.
當然.
她這樣的意思就表示說.
謙卑比順福還不容易.
是不是?.
丟子沒.
耶穌是一個謙卑的.
一個謙和的一個君王.
所以他要我們也以基督.
保羅說以基督耶穌的心為心.
在腓立比書.
保羅告訴我們說要以基督耶穌的心為心.
以基督他的榜樣為我們的榜樣.
他是一個溫柔的.
他是一個謙卑的君王.
所以他也要我們做一個溫柔謙卑的人.
跟你旁邊說你要做一個溫柔謙卑的人.
平常沒有機會跟你先生講.
這時候趕快講.
有一個姐妹.
當她信耶穌之後呢.
她家裡的人就非常反對她.
特別她先生.
那麼後來牧師就問她了.
她說你先生反對你信耶穌.
你怎麼辦呢?.
這個姐妹回答她說.
他逼迫我.
我就把飯做得更好.

$^{361}$如果他發牢騷.
我就把地掃得更乾淨.
如果他說話魯莽.
我就更溫柔的回答他.
我要做一個賢妻良母.
讓他知道基督徒在各方面的表現.
阿們.
姐妹們說阿們.
弟兄們也說阿們.
換過來一樣.
先留準備.
聖經上讓我們看見一件事情.
我們不是以惡報惡.
我們不是以怒罵還怒罵.
是不是.
聖經上告訴我們說要以善勝惡.
有人打你的右臉怎麼樣.
打回去.
今天你在這世界上看到就是如此.
是不是.
你現在看到就是如此.
打的一個結果是什麼.
我今天才看.
現在這個.
美國的這個炸彈反正已經.
總是炸到一些不該炸的地方去了.
炸到別人的大學去了.
所以現在伊朗說你如果不道歉.
我就開始來炸你在中東地區的你們的大學分校.
這叫做以惡報惡.
是不是.
你會發現伊斯蘭教是如此的.
猶太教也是如此的.
就也就是這樣的觀念懂這意思嗎.
那請問耶穌的教導是什麼呢.
耶穌的教導不是這樣子的.
當然有人說這個關於這個戰爭我們到底應該怎麼樣子.
我不今天不談這個問題這個問題很大的一個問題.
如果你參加禱告會有一次藏雷牧師.
有特別在稍微講了一下.

$^{401}$那麼大家可以瞭解一下.
公義的戰爭.
很多時候打到最後已經沒有公義了.
這個公義不公義呢.
到最後只變成是利益不利益.
是不是.
石油漲價馬上就影響到打仗的方式.
一旦得到了利益之後.
這個戰爭馬上結束.
你會發現這個世界已經走到一個地步就是.
大家分不清楚什麼是標準了.
但是親愛的兄弟妹.
聖經裡面永遠是我們的標準.
神的話永遠是我們的標準.
當然實際在行動實際在運作的時候.
我們不是一天兩天我們就做到是不是.
可是至少我們清楚知道.
我們是天國的子民.
今天伊斯蘭教用武力來征服全世界.
那麼我們基督徒應該怎麼反應呢.
在歷史上面.
十字軍東征就是一個最好的例子.
十字軍東征最後的結果是什麼呢.
還是一樣打敗.
而且帶來一個結果就是.
讓今天的回教徒.
他們對特別是對西方.
有一個很深的不可磨滅的印象.
就是西方的基督教.
就是像十字軍Crusade一樣.
所以有一些的宣教機構.
它的名稱本來叫做Crusade.
現在全都改了.
因為要向回教徒傳福音.
你一講Crusade那好了.
他馬上拒絕你了.
連名稱他們都把它改掉.
親愛的兄弟妹.
我們相信上帝可以不要用一兵一卒.
來使得全世界這十七十八億的穆斯林.

$^{441}$統統要來信號耶穌.
你相信嗎.
伊朗就是一個例子.
這些年來我們聽到很多的報道.
說伊朗因著許多的人.
他們聽到福音的廣播.
或者看見在異夢當中.
在異象當中.
看見耶穌向他顯現.
或者天使向他顯現.
結果他們信了耶穌.
伊朗現在有幾百萬的基督徒.
都是這樣子.
神用這樣的方式.
上帝可不可以用這樣的方式.
來讓這些穆斯林來信主呢.
可不可以讓伊拉克.
可不可以讓科威特.
可不可以讓這些北非的這些土利西亞.
這些很極端的穆斯林.
他們怎麼樣子來信耶穌呢.
你相不相信上帝可以得到他們.
但是絕對不是靠我們的船堅炮利.
主耶穌上十字架.
因著他自己的舌疾.
所以使得全世界全人類都有機會來得救.
所以我們要做一個謙卑的子民.
謙卑的子民.
謙卑很重要.
撒旦墮落就是因為什麼.
就是驕傲.
謙卑的另外一個意思.
在腓立比書第二章第五節到第七節講的.
耶穌基督他本來有神的形象.
不以自己與神同等為搶奪的.
反倒虛極.
取了盧浦的樣式成為人的樣子.
虛極.
耶穌基督虛極什麼意思呢.
虛極就是放下權力.

$^{481}$他本來是神.
他與神同等.
他是神的兒子.
他是三位一體神當中的一位.
是不是 你懂我意思嗎.
從神學的角度來講.
但是他願意虛極.
意思就是說.
他願意把他自己的right.
他的權力放下來.
在美國這是很困難的一件事情.
因為整個的美國的文化教導裡面就是告訴你們.
這是你的right.
有沒有.
你要去爭取你的right.
你不要放棄你的right.
帶來一個結果.
很多時候我們也不知不覺被影響.
但是今天讓我們再次思想.
耶穌他是一個謙卑的君王.
他自己做了最好的示範.
就是他願意放棄自己的權力.
取了盧浦的形象.
定在十字架上面.
謙卑就好像放棄權力.
就好像亞伯拉罕.
在《創世記》裡面提到說.
亞伯拉罕跟羅德.
他們兩個人因為僕人在那邊爭.
那麼兩個僕人在爭.
亞伯拉罕就對羅德說.
你我我們不可以相爭.
你的牧人跟我的牧人也不可以相爭.
我們是骨肉.
所以亞伯拉罕就做了一個選擇.
他說好吧.
遍地都在你的眼前.
你選吧.
你來選.
你要去哪裡.

$^{521}$我都讓給你.
結果羅德一看.
當時候的全約旦河平原.
那個地方園子極美.
水果 蔬菜真是豐盛得不得了.
他就說我要全約旦河的平原.
亞伯拉罕給他了.
請問亞伯拉罕得到什麼.
約旦河平原那個最好的給了他了.
那麼請問亞伯拉罕得到什麼.
你看上帝怎麼跟亞伯拉罕講.
後來上帝跟他說.
凡你所看見的一切地.
我都要喫給你和你的後裔.
直到永遠.
哈利路亞.
他把當時候可能覺得是最好的給了他.
給了羅德.
但是上帝把一切所有的都給了他.
哈利路亞.
對嗎.
今天可能你也在面臨到一些這種問題.
人生當中我們都會有面臨到這種.
夫妻之間也好.
兒女之間也好.
或者是你跟同事跟著朋友之間.
正是在教會服侍的當中.
你都會面臨到這些的情形.
這本來是應該你的權利.
但是別人要把你的權利拿走.
或者別人要搶了你的權利.
或侵犯了你的權利.
我們願不願意把我們自己的權利放下呢.
聖經上說.
神賜恩給謙卑的人.
上帝主黨驕傲的人.
這句話.
我覺得每一個弟兄姊妹.
包括我們教導我們的兒女.
都需要把它背起來.

$^{561}$這個世界是反其道而行.
我們來看看亞哥書第四章第六節.
可以的話你就把它一起練一遍.
就順便把它背起來.
預備 請.
神主黨驕傲的人.
賜恩給謙卑的人.
Amen.
跟你旁邊說.
神賜恩給謙卑的人.
在兩百年前.
在意大利的南部.
有一個偏僻的村莊.
在這個村莊.
有兩個年輕人.
他們兩個人很要好.
這兩個年輕人.
一個叫做瑪利歐.
他這個瑪利歐是地主的兒子.
家裡面很有錢.
另外一個叫做安瑟摩.
安瑟摩的家境貧寒.
他是很貧窮的邪將的兒子.
雖然他們兩個家庭的貧富懸殊.
但是他們卻非常好.
常常玩在一起.
這個瑪利歐他有一個心願.
就是他想要進入到城裡面的修道院去練書.
有一天.
終於他的心願達到了.
他申請要進入這個修道院去練書.
申請成功了.
所以他就要進入城裡面去了.
當他離開這個村莊的時候.
他的好朋友安瑟摩就來跟他送行.
送行的時候他就跟他講一句話.
他說我必定會為你禱告.
結果經過幾年之後.
這個安瑟摩他也來到這個修道院.
做什麼呢?.

$^{601}$不是來練書.
他來做奴僕,做僕人.
好了.
但是他很高興.
他能夠到修道院.
他又再一次跟他的好朋友.
能夠跟他親近了.
這個好朋友瑪利歐讀完書了之後.
就被升為叫做priest.
就是天主教稱他叫做司祭.
或者說我們今天講神父.
他當了神父之後.
第一次講道很緊張.
心裡面在那邊預備,禱告很緊張.
這一個好朋友安瑟摩就跟他講.
他說瑪利歐你放心.
我會為你禱告.
好了.
當他第一次站上講臺的時候.
臺下面這些人會眾眼睛都看著他.
第一次講道兩隻手在那邊發抖.
但是他卻看見他的好朋友安瑟倫.
安瑟摩就在下面.
在那邊在一個圓柱子旁邊.
他就低著頭就在那邊禱告.
圍著他禱告.
那一天他講道之後.
哇.
真是感覺到聖靈同在.
講道非常有能力.
所以當他講完道之後.
臺下的人大家都很高興.
都很說今天你講得真好.
從此之後他一舉成名.
後來又被晉升.
一步一步就變成名講員了.
就被邀請到各個地方去講道了.
結果甚至有一天他被邀請到.
哪裡講道呢.
被邀請到羅馬的聖彼得大教堂講道.

$^{641}$你們有沒有去過聖彼得大教堂.
哇.
那個是不得了啊.
那是一件事情.
好了.
那麼他的好朋友安瑟倫又來了.
他說我會為你禱告.
他每一次當他出去講道的時候.
他就跟著他.
因為後來他就等於像僕人一樣.
事奉他.
跟著他.
為他禱告.
所以那次他要到聖彼得大教堂講道.
那麼這個好朋友也一樣的跟他說.
他會為他禱告.
好.
那一天當講道的時間到了.
那麼他站上講臺.
他已經是講道無數.
已經是世界知名的講員了.
所以當他站在講臺的時候.
他就看著臺下烏鴉鴉的一片黑的人.
他正在講道的時候.
他就感覺自己頭腦一陣混亂.
他就感覺自己那一天好像很沒有力量.
講道好像沒有什麼.
就覺得好像還使不出勁來.
他不知道怎麼搞.
今天怎麼會如此呢.
當他講道結束.
垂頭喪氣的走下講臺.
當天晚上他就去找他的好朋友.
沒有想到.
他的好朋友在那一天的早晨過世了.
原來那個一直在為他禱告的好朋友過世了.
那天的講道沒有他的禱告.
哎呀.
他心裡面非常的難過.
後來過了幾天他就跑到墳墓.

$^{681}$這個安瑟摩的墳墓面前.
他就迫切地跪在那邊.
在那邊禱告.
那麼剛好修道院的院長經過.
就看了他就問他說.
他說你是不是在祈求.
你能夠恢復以前的講道的能力.
這個馬利歐回答院長說.
他說不.
我只求主賜給我一個謙卑禱告的心.
聽著沒.
我們都想要做馬利歐那樣風風光光的人.
院長很訝異.
他以為他在禱告.
要恢復他以前那些講道的能力.
但他經過好朋友的過世他知道.
他今天之所以講道有能力.
被上帝使用.
是因為有一個好朋友在那邊.
晝夜不停地為他禱告.
他更加地羨慕.
有像這一個好朋友那樣謙卑溫柔的心.
聽著沒.
這個世界有很多的榜樣.
特別是這些媒體的發達.
讓我們看見.
你要更多出人頭地.
你要比別人更強.
你要得第一名.
你要做網紅.
你要成功.
你要賺錢.
你要有更多的房子.
別人纔會尊重你.
在這個世界的世俗化的潮流之下.
基督徒很容易迷失了自己.
在我們家庭的生活裡面.
在我們教會的生活裡面.
千萬記得.
不是誰的聲音大.

$^{721}$不是誰的學歷高.
不是誰會講話.
不是誰比他懂.
親愛的同志們.
在神的家中.
慈善看誰比較像耶穌基督.
Amen.
讓我們羨慕的是越來越像耶穌基督.
像他那樣子的謙卑溫柔.
所以同志們.
當我們知道這一位主是這樣子的一位謙卑的主的時候.
他將改變我們在我們自己平常的人生的處事為人.
在我們的家庭.
在我們夫妻關係.
在我們跟兒女的相處當中.
我們開始更多要學習.
不是用世界上的辦法.
而是效法主的心.
謙卑溫柔做一個僕人來服侍.
好 第二個.
當這些眾人他們看見耶穌騎著驢鞠子進城的時候.
眾人在那邊呼喊什麼.
呼喊說 和撒娜.
和撒娜.
這句話.
其實它原來的意思是.
我求你來拯救.
我求你救救我們.
但是在這個地方.
它的意思是什麼.
是讚美神的意思.
是稱頌神的意思.
親愛的同志們.
所以.
這一位主.
他配得我們的讚美.
世上的君王.
帶著自己的威風.
征戰得勝.
這種的威榮進來的時候.

$^{761}$大家去稱服他或者是敬拜他.
是帶著一個什麼.
一個懼怕.
你懂意思嗎.
但是我們的主.
他來的時候.
他是帶著什麼.
他是帶著愛.
他帶著謙卑.
他帶著犧牲.
他來到我們的中間.
有一天我們的主.
他要從天降臨.
你看起書裡面.
還提到說從天降臨.
天使天君跟著他一同駕雲降臨.
丟死被.
耶穌駕雲降臨做什麼.
天使聽見跟著他來做什麼.
來打這一場最後的大戰.
把這些地上的這些壞人.
統統把他消滅掉嗎.
是不是.
聖經沒有這樣說.
聖經講說.
耶穌帶著天使天君.
一同駕雲降臨做什麼.
他來迎接我們.
地上的聖徒們.
死的活的統統要復活起來.
如果我們那時候還在的話.
我們就要在天空中與主相見.
丟死被.
這是我們的盼望.
所以.
當這個終束之主日的時候.
我們想到當年.
這些.
這些眾人他看見耶穌騎著魚進來.
他們當時的心態是因為耶穌.

$^{801}$他們猶太人等待的是一個.
就是要再來的有一位彌賽亞.
聖經上所說的就你們講的彌賽亞.
他們把耶穌就當作是那一個彌賽亞.
耶穌是不是那個彌賽亞.
可是你要知道.
他們所期待的彌賽亞.
是一個政治性的.
要重建大魏倒塌的帳幕.
要把以色列重新的復興起來.
成為一個像大魏王朝那樣的一個復興的王朝.
所以他們歡迎的時候.
其實他們是歡迎一個政治性的君王.
你懂這意思嗎.
但是今天.
我們所歡迎的.
我們所敬拜的.
我們所要歌頌的.
奉主名來的是應當稱頌的.
就是在講這位耶穌基督.
他是奉神的名而來的.
是應當得到我們的歡迎的.
稱頌就是歡迎的意思.
應當得到我們歡迎的.
所以親愛的羅子妹妹.
每一次當我們在這些主日敬拜的時候.
你有沒有想到.
你來到這個地方.
你是來敬拜這位主.
你是來敬拜這位千千和和.
騎著驢精子來到我們中間的.
那位和平的君王.
我們來敬拜他.
我們來讚美他.
我們來歡迎他.
所以當我們在這邊唱詩歌的時候.
我們在這邊歡迎他.
你懂這意思嗎.
羅子妹.
所以.

$^{841}$這個.
我們在教會里面.
羅子妹我常常說.
禮拜天是主的日子.
這是主的日子.
你們在其中要高興歡喜.
為什麼.
為什麼高興歡喜.
為什麼主的日子要高興歡喜.
因為將人所騎的石頭.
如今已成了.
仿假的頭塊石頭了.
那位定死的耶穌基督.
已經復活了.
親愛的羅子妹.
我們所敬拜的這位主.
他是一位又真又活的神.
他不是一個死去的耶穌.
Amen.
羅子妹.
當我們這樣說的時候.
其實對一個不信主的人.
他看到我們怎麼樣子對神的態度.
他會決定.
我到底我們的神是一個什麼樣的神.
求主幫助我們.
不要把我們的神變成死的神.
每一次當我在敬拜的時候.
不管我心情如何.
我看著字.
我就是專心的把這個字裡面的內容.
當我思想去的時候.
我就盡心的從那個字裡面.
從我的口中把它唱出來.
所以它裡面告訴我說.
應該要拍手的時候我就拍手.
應該要舉手的時候我就舉手.
Amen.
定死妹.
不是憑你自己的感覺.

$^{881}$不是你現在的心情來決定.
而是因為耶穌基督.
祂是活著的.
所以我們在敬拜這位是活著的神.
因為以讚美.
Amen.
所以我們要向主來獻上我們的讚美.
每一個禮拜天.
在我們的小組聚會里面.
在我們各樣的聚會里面.
定死妹.
千萬不要把敬拜當成是一個.
好像是一個其次的.
好像是一個就是消磨時間的.
或者是等人的.
不是.
如果你聚會的時候習慣敬拜完才進來.
那你有問題.
瞭解我意思嗎.
敬拜不是隻有聽講道.
敬拜是整個都是一個敬拜.
所以我們來到神的面前的時候.
我們接著唱詩歌.
你說哎呀這個詩歌我不喜歡.
或者怎麼樣.
不是你喜歡不喜歡啊.
這個詩歌我不熟.
不熟的歌我就學習熟嘛.
瞭解我意思嗎.
如果說今天都是你很熟的歌.
你可以唱出那樣子的讚美嗎.
不要給自己找很多理由.
求主幫助我們成為一個讚美的人.
我們主配得你的讚美.
Amen.
我們主配得我們的敬拜.
你看著每一次看到那些穆斯林在敬拜的時候.
他們在那邊跪著那邊五體投地的時候.
就算他們當中有一些人是儀式的.
行事的.

$^{921}$但是我相信他們當中有很多是真心的.
是不是.
那反觀呢.
今天我們基督徒有多少人這樣真心的在那邊.
如果今天要跪.
要拜.
可能我們的已經不太習慣了.
很久沒有拜了.
很久沒有跪了.
當然你可以說我年紀大了.
年紀輕的時候你跪嗎.
年紀輕的時候都不跪.
年老了也不會跪.
我現在應該比較可以進入這個老年階級.
但是我仍然盡我自己所能的.
我可以跳 我就來跳 Amen.
我可以來讚美 我就來讚美.
我可以大聲 我就來大聲.
因為我們的神配得我們的讚美.
求主幫助我們.
聽我準備.
讚美要使我們的生命改變.
讚美不是在你環境好的時候讚美.
讚美是在你環境不好的時候.
仍然要讚美 Amen.
這次我到德州去參加一個研習會.
完了之後有幾天我就到Houston.
去看我的女兒還有外甥.
我女兒跟女婿他們去年.
因為教會有祠堂的負擔.
所以他們的一個聯繫牧師就帶著.
大概有13個15個.
從聖地牙哥他們就全部relocation.
到休斯頓去重新建立異間教會.
他們目標是要針對那些Asian-American.
所以他們本來7月就準備去.
後來因為生孩子所以到了9月生完了.
所以他們今年1月就搬過去了.
我們去看他們.
我要講的是重點是在.

$^{961}$當回來的時候.
我坐飛機從休斯頓回LA.
坐United Airlines.
比如說今天要上飛機.
那我昨天我就先check-in.
check-in之後我就拿到boarding pass.
OK 那我就放心了.
晚上就睡覺了.
第二天早上.
因為我們平常晨早大概6點以前就要起床.
雖然休斯頓慢兩個小時.
但是我還是很早起來.
起來之後我一看.
怎麼United又給我一個信息text說.
他說你的班機已經要起飛了.
我說已經要起飛了.
我還在這裡呢.
然後他說我們會在你起飛之前.
會再有一個last call什麼之類的.
但是我就想不對啊.
我就去看那個text上面還寫著說.
我的check-in的confirmation.
我這樣一點擊進去.
連那個boarding pass都是早上的那班飛機.
早上的班飛機是5點多boarding.
然後大概是6點.
可能6點10分左右就起飛了.
那我接到text的時候.
事實上我是不可能的.
我到機場都要一個小時.
我不可能的坐這一班的.
我也沒有去check-in這一班.
所以好了就打電話去機場.
打了兩次最後的機場的服務中心的人說.
他說我們這裡不能處理.
你必須要到機場去才能夠處理.
好 我們就我女婿就帶著我.
就這樣子我們就全家.
還外甥他施母.
還我們全家陪著我一起送到機場去.

$^{1001}$怕我到那邊萬一被.
所以就把我放就進去了.
我女婿帶著我一去問當地的地勤人.
地勤人也告訴我說.
他說I can't help you.
他說你必須要call這個number.
我說我們call了.
他說你們再call.
哎呀.
Eddie.
好了.
你知道嗎.
那一天早上我們的臨修經文.
我不知道你還記不記得.
就是耶和華必救贖他的僕人.
凡投靠他的怎麼樣.
必不致羞愧.
我就想到.
耶和華必救贖他的僕人.
所以神話一進來.
我就開始.
我就開始禱告.
到後來我就開始讚美.
因為這幾天我們的經文裡面都在講.
就是大衛即使是在很苦的時候.
是不是.
開始讚美.
我就那這個事情來了.
自己要面臨到這個事情了.
我就開始.
我一路上就進了關之後呢.
然後再到那個櫃檯.
到了櫃檯之後呢.
還有兩個小時我的班機.
但是還沒有人在那邊.
我就跑一看有四點多的有一個班機.
我說那我趕緊去坐四點多的.
我去了四點多.
他說他們要快.
但是他說好.

$^{1041}$那我就就剩你一個.
趕快就一查.
說你的機票有問題.
他說你早上已經飛了.
我說I'm here.
他說你已經飛了.
我說我已.
Anyway.
但是呢.
就是從一路上面禱告.
一路上面讚美.
所以當我面對他的時候我就說.
我後來就問他.
我說可不可以請你的supervisor來跟我講.
他找了一個supervisor沒辦法解決問題.
後來找到最高的supervisor來了.
那個把我帶進去.
帶到辦公室裡面去.
幫我弄弄弄.
結果說OK.
搞定了.
他說你還可以去坐同一班飛機走.
但是呢.
你本來的位置沒有了.
現在坐在中間位置.
我本來是IOC.
而且我本來是有extra leg.
都沒有了.
好了.
他說你要到櫃檯去.
看看有沒有辦法給你.
有空位.
我到了櫃檯的時候.
又等了20分鐘.
等到他boarding之前的20分鐘15分鐘.
他開始有人來了.
我又跟他講講講.
講的時候他說.
我說可不可以有IOC.
我已經配了費買這個extra leg.

$^{1081}$我說可不可以有extra leg.
他說沒有.
沒有.
他說你退費他會退給你的.
就在他講的時候.
他突然說.
有兩個位置出來了.
我說那你幫我看看看有沒有.
他說有.
結果就在我本來定的那個位置的.
另外一邊的IOC.
瞭解我意思嗎.
他一弄完就在那一秒鐘之間.
一下子出來兩個seat.
我當時馬上跟他講.
我說.
我說God.
Praise the Lord.
God help me.
他眼睛看著我.
Oh yeah yeah yeah.
God help you.
Hallelujah.
上帝幫助我們.
Amen.
上帝是拯救我們的主.
上帝是救贖我們的.
Amen.
所以我們雖然還沒有看到事情的發生.
但是我們學習用信心來讚美.
Amen.
我講的這個不是隻是一個理論.
這是我們在生活當中.
我們經常會碰到的事情.
但是你要抓住相信神的話.
當你來操練的時候.
你來經驗神的話.
這個聖經.
神的應許就成為你的了.
Amen.

\newpage



\section{馬太福音 21:1-11}
\label{sec:031894}
\textbf{看哪,你的王來到你這裡!}~~~~黃馮潔霞師母
\newline
\newline
連結: \href{http://www.church.com.hk/acms/content.asp?site=cdc\&op=show\&type=product\&code=031894\&layout=sermon}{\texttt{031894}} ~~~~ 語音日期: 2026-03-29
\newline
\newline
\hyperref[sec:031893]{\small{< < < PREV SERMON < < <}}
~
\hyperref[sec:index_chronic]{\small{[返順時目]}}
~
\hyperref[sec:031895]{\small{> > > NEXT SERMON > > >}}
\newline
\newline
馬太福音 21:1-11
\newline
\begin{longtable}{cl}
\hline
\hline
章節 & 經文 (和合本修訂版)\\
\hline
21:1 & \begin{tabularx}{0.7\textwidth}{X} 耶穌和門徒快到耶路撒冷，進了橄欖山的伯法其時，打發兩個門徒， \end{tabularx} \\ \\ \relax
21:2 & \begin{tabularx}{0.7\textwidth}{X} 對他們說：「你們往對面村子裡去，會立刻看見一匹驢拴在那裡，還有驢駒同在一處，解開牠們，牽到我這裡來。 \end{tabularx} \\ \\ \relax
21:3 & \begin{tabularx}{0.7\textwidth}{X} 若有人對你們說甚麼，你們就說：『主要用牠們。』那人會立刻讓你們牽來。」 \end{tabularx} \\ \\ \relax
21:4 & \begin{tabularx}{0.7\textwidth}{X} 這事發生是要應驗先知所說的話： \end{tabularx} \\ \\ \relax
21:5 & \begin{tabularx}{0.7\textwidth}{X} 「要對錫安的兒女說：看哪，你的王來到你這裡，謙和地騎著驢，騎著小驢—驢的駒子。」 \end{tabularx} \\ \\ \relax
21:6 & \begin{tabularx}{0.7\textwidth}{X} 門徒就照耶穌所吩咐的去做， \end{tabularx} \\ \\ \relax
21:7 & \begin{tabularx}{0.7\textwidth}{X} 牽了驢和驢駒來，把他們的衣服搭在上面，耶穌就騎上。 \end{tabularx} \\ \\ \relax
21:8 & \begin{tabularx}{0.7\textwidth}{X} 許許多多的人把自己的衣服鋪在路上，還有人砍下樹枝來鋪在路上。 \end{tabularx} \\ \\ \relax
21:9 & \begin{tabularx}{0.7\textwidth}{X} 前呼後擁的人群喊著說：「和散那歸於大衛之子！奉主名來的是應當稱頌的！至高無上的，和散那！」 \end{tabularx} \\ \\ \relax
21:10 & \begin{tabularx}{0.7\textwidth}{X} 耶穌進了耶路撒冷，全城都驚動了，說：「這是誰？」 \end{tabularx} \\ \\ \relax
21:11 & \begin{tabularx}{0.7\textwidth}{X} 眾人說：「這是從加利利的拿撒勒來的先知耶穌。」 \end{tabularx} \\ \\
[1ex]
\hline
\hline
\end{longtable}
$^{1}$各位弟兄姊妹平安.
今天是中樹節 也是中樹主日.
復活節之前的一個禮拜就是這個中樹節.
從明天開始就是一個聖週.
耶穌基督最後一個禮拜在世界上為我們活著.
雖然只有幾天 但你會發覺有很多事發生.
如果你有時間這個禮拜.
嘗試翻開福音書 讀一下耶穌基督.
最後這個禮拜究竟做了些什麼.
你心裡會很感恩.
也會效法他怎樣去過人生最後的一個禮拜.
說到最後一個禮拜.
每一次來和大家分享神的話語.
都一定有一些我們的肢體病或離開世界.
在人生去到最後一個階段的時候.
有人很快就睡一覺就走了.
也有人要有一段時間在地上.
還要治療或在療養院.
我和大家今早起床祈禱最後準備的時候.
發覺還是要和大家分享一下.
雖然現在要花一點時間.
上個禮拜有兩個安息禮拜.
兩個姊妹.
當然她們都相當好了 都八十多歲了.
有一個就接近九十歲了.
一個就是當她發現自己病了 要去治療了.
她本來就說我不治療了 因為很辛苦.
所以家人就和她談一談.
結果她肯了 怎麼做了兩次呢.
原來她身體是受不了的.
結果就沒有繼續做了.
但是當我和姊妹一起去探望她的時候.
探望過幾次 她都只說了一句.
「我不怕 我已經準備好要去見我的天父.
我的主耶穌 和見我的家人.
讓我先走一步去天上天堂的人」.
她說「你知道為什麼我會說我隨時可以走嗎?.
因為我一生人的心願就是.
將我的家人一個一個帶到神的面前.
包括父母 兄弟姊妹 他的子女 孫子.

$^{41}$我的心願就是要帶他們去信主.
你知道嗎?我本來還欠一個弟弟.
他不肯信主 他的太太也不肯信主.
最近他和我說.
「我們都信耶穌了 我們都信了」.
她說「可以走了 可以向神交障了」.
所以我去探望這些病人的時候.
我們的弟兄姊妹聽到這句話.
他們真的很開心.
因為他們不是每天說.
「為我祈禱 讓我快點康復.
我不想先走 我還有很多事情要享受.
我有很多兒孫滿堂 我不想走.
你快點為我祈禱 讓我康復」.
但是如果你遇到這些姊妹.
或者這些主裡面的弟兄姊妹.
你會發覺很感恩.
他還安慰了我們.
他去探望人很擔心他不知道怎樣.
他說「你不用擔心 我準備好了.
現在最好你祈禱就是讓我快點回去.
馬上回去.
我要去見我的天父 我的主耶穌」.
另外一個他沒有機會這樣.
但是他又很孝順他的子女.
就是將他搬到更近那兩個女兒住的地方.
初初很不習慣.
因為住慣了這裡南灣區.
什麼都在這裡.
甚至他跟我說.
「我的墳地都買在這裡了.
你搬到這麼遠 要開兩個小時車.
我真的覺得很不習慣.
經常都要回來」.
結果經過幾個月之後.
那些女兒就做見證話.
他不再說要回來.
他在那邊很開心.
適應了.
因為那是一個新的老人院.

$^{81}$還有在海邊.
他說每天都可以在海邊走走.
很美的風景.
而且那些女兒很乖.
兩個輪流帶他出去吃東西.
所以他覺得在這裡都差不多.
跟在托倫斯一樣.
所以他跟我們初初很不習慣.
後來就習慣了.
沒有聽到他打電話來說.
「我要回來了 你快點跟我的女兒說.
搬我回來了」.
沒有了.
但是有一天早上.
當那些人去他房間裡看看他.
就發覺他原來已經回到天府那裡了.
他就是睡著睡著 中風就走了.
一直都沒有醒過.
因為那些女兒覺得他都快90歲了.
無謂再搞他了.
要不然要做手術 又要這樣.
那更辛苦了.
就算醒了都好像在受苦.
就決定跟醫生開會說.
他什麼都不做了.
所以頭尾八天.
八天他就很平靜.
很安穩地回到天府那裡.
我們去探望他的時候.
他那時候還是不醒的.
但是你發覺他不是一個病人.
他是一個睡得很熟 很舒服的一個姐妹.
睡得很甜.
又沒有什麼管制 全身都沒有管制.
有兩種過最後一天 最後一週.
最後一個月的姐妹們.
你選哪一種?.
哪有選擇 有選擇就好了.
根本沒有選擇.
我們都沒得求.

$^{121}$神讓我安安樂樂地睡著睡著就走了.
或者說 神啊.
如果你還要我留在世界上.
是不是有些未做完的工要我幫你完成呢?.
那我們就要留下.
就算是受苦 就算是面對疾病.
那種治療 都要留下.
因為神還要用你這個破碎的生命.
一個最後的階段.
主耶穌最後那一週.
祂做了很多事情.
我這裡不會逐一跟你們說.
你們自己回家看一下.
那幾張聖經.
《馬太》,《馬可》,《路加》,《約翰》.
你可以看完它.
你會發覺 原來耶穌做了這麼多事情.
我只是跟大家先來看一節聖經.
這節聖經 我們BSF很熟悉的.
因為我們每個人都讀過這本書.
就是《撒迦利亞書》.
為甚麼不可以呢?.
我們先看看.
《撒迦利亞書》.
九章九節.
大家一起來讀 好不好?.
請.
「錫安的民吶 應當大大喜樂.
耶路撒冷的民吶 應當歡呼.
看吶 你的王來到你這裡.
他是公義的 並且施行拯救.
謙謙和和地騎著路.
就是騎著路的驅使」.
這裡提到撒迦利亞先知.
舊約的一位先知.
他就在講預言.
他說 你知不知道有一天.
錫安的民 就是耶路撒冷的民 居民.
或者再擴大一點講 就是以色列人.
他們要大大的喜樂.

$^{161}$因為有一位王駕到.
他說 你看你的王在這裡.
他為甚麼是公義的.
他是來施行拯救的.
還有他的表現 他的性格.
是謙謙和和的.
不是騎馬 是騎驢.
但是還是驢驅使.
驢驅使就是那隻驢仔.
驢就是媽媽陪著他.
那隻驢仔就在他身邊.
耶穌不是騎馬 是騎著驢仔.
更加謙卑.
你知不知道撒迦利亞先知.
他講了很多關於彌賽亞的預言.
除了以賽亞先知之外.
他是第二個講得最多關於耶穌.
他生平 或者他所遇到的事.
我試一下和大家講.
在第三章那裡 八九節.
他說 妙瑞彌賽亞.
將在一日內 除去這個地的罪孽.
這個預言很大.
就是當他釘十字架的時候.
在十字架那裡 他就是背負了我們的罪.
擔當了我們的罪 除去我們的罪.
在第六章十二十三節說.
妙瑞彌賽亞將建造活石的聖殿.
這是彼得前書二章四至六節講的.
並且在他的寶座上做祭司.
將祭司和君王的職份合而為一.
耶穌基督又是君王 又是祭司.
就在撒迦利亞書那裡預言.
還有就是第九章九和十節.
剛剛我們讀的那個.
他是謙謙和和的站著牢出現在耶路撒冷.
然後十一章十二到十三節.
彌賽亞將以三十塊錢被賣.
接著十二章十節說.
彌賽亞被責 就是被賜.

$^{201}$以色列人將為他悲哀.
耶穌十字架那裡.
那些兵兵不知道他有沒有真的死了.
就用槍刺一下他的肋旁.
看看他是不是真的死了.
結果水和血流出來.
不只是血 水和血已經分開了.
他的心已經破碎了.
還有就是十三章一節說.
彌賽亞將成為潔淨的源頭.
我們每個人都去到他面前.
得蒙潔淨.
撒迦利亞書十三章第七節又怎麼說呢.
他說牧人彌賽亞將被擊打.
羊群就分散了.
這就應驗了在福音書裡面.
當人抓了耶穌的時候.
那些門徒都分散了.
所有跟從他的人都分散了.
十四章四節說彌賽亞將返回.
橄欖山這個指定的地點.
就是說他會再回來.
然後彌賽亞將在全地作王.
就是十四章第九節.
所以你看看這九樣.
就是撒迦利亞書去預言耶穌.
祂要來.
和祂要面對什麼 祂要發生什麼.
所以今天我們只說一段.
就是說耶穌很榮耀地進入耶路撒冷.
我們來看看.
我會用《馬太福音》二十一章.
來先做這個內容.
等一下還會用《路加福音》.
現在第一章這張slide.
可不可以請弟兄們幫我們讀.
一和二節.
請.
耶穌和門徒將近耶路撒冷.
到了伯伐奇在橄欖山那裡.

$^{241}$耶穌就打發兩個門徒.
對他們說 你們往對面村子裡去.
必看見一匹驢山在那裡.
還有驢驢同在一處.
你們解開 牽到我這裡來.
下一章請姐妹們幫我們讀.
若有人對你們說 你們就說.
主要用他 乃人必立時讓你們牽來.
這是成就 是要應驗先知的話說.
要對石安的居民說.
看啊 你的王來到你這裡.
是溫柔的 又騎著驢.
就是騎著驢驅子.
請弟兄們讀這三節.
門徒就照耶穌所吩咐的去行.
牽了驢和驢驅來.
把自己的衣服搭在上面.
耶穌就騎上.
眾人多半把衣服鋪在路上.
還有人撼下樹枝來鋪在路上.
請姐妹們讀.
前行後鎚的眾人喊著說.
「和撒那」原來是有求救的意思.
在這裡是稱頌的話.
「歸於大位的子孫.
奉主明來的是應當稱頌的.
高高在上 和撒那」.
耶穌既進了耶路撒冷.
合乘都驚動了 說「這是誰?」.
眾人說「這是加里尼拿撒勒的先知耶穌」.
這一段是馬太福音記載的.
也講得很清楚.
特別是他講到「和撒那」.
就是大位的子孫.
要稱頌祂 要向祂求救.
因為祂才能夠救我們.
「漢」.
因為我們的題目是「漢」.
你的「黃」.
來到你這裡.

$^{281}$所以「漢」.
其實「漢」在聖經裡出現了很多.
我只是要你來讀幾節.
「漢」就是叫我們睜大眼睛看看.
你要留心地看.
看緊.
當然這個「漢」不只是看就算了.
你已經進入了你的思想.
要經過你的腦裡想一下.
在你的心裡再揣摩一下.
這裡第一節「漢那」.
是約翰福音一章二十九節.
「神的羔羊除去.
或者亦背負世人罪孽的」.
在耶穌剛剛出來做傳道的時候.
祂的表哥施洗約翰.
就是對著祂自己的門徒.
跟著祂的門徒.
來講這句話.
祂不怕別人跟了耶穌.
祂就是要老老實實跟他們講.
我不是那個可以除去你們罪的人.
雖然我幫你洗禮.
所有願意悔改的人都去施洗約翰受洗.
因為祂是負責施悔改的洗.
但是祂很老實跟他們講.
我不過是用水幫你洗.
但是這一位來的.
祂是用聖靈和火力來幫你施洗.
今天我們有沒有把耶穌基督指給人看.
看看「神的羔羊除去.
背負世人」包括你和我的罪孽.
有些宗教是叫人去向人認罪.
向人告解.
是向領袖 屬靈領袖.
去講他的心事.
希望他宣布或宣告.
你的罪赦免了.
人是不可以的.
你得罪他.

$^{321}$他可以說我寬恕你.
但是你所有隱藏了或未現出來的.
他不可以說我赦免你.
所以他可以為你祈禱.
每個人都可以彼此代禱.
為他的軟弱 他的過錯祈禱.
但是你不可以宣告你的罪赦免了.
只有耶穌基督有這個能力宣告.
因為祂在十字架上為我們釘死了.
還了債給我們的天父.
所以今天你知不知道.
一個欠債的人其實是很辛苦的.
現在我們可能信用卡的債.
每個月清的不要緊.
有些人說你還建立你的信用.
日後輪到要買車 買房子.
知道你的信用分數很高.
但是你有沒有覺得.
雖然你每個月賺到夠錢用.
但是如果你還背著房子.
還是有債的.
那個債不是幾千元.
是幾十萬.
你始終覺得自己是一個欠債的人.
但如果有一個人說.
我幫你還了你的債.
我可以宣告你是無債一身輕.
你說多好.
今天我們生命裡面.
很多很多罪惡 虧欠.
如果我們仰望耶穌.
去信靠耶穌.
你所有的罪孽都得到赦免.
因為約翰一書說.
我們要認自己的罪.
神是信實的 是公義的.
必要赦免你的罪.
洗淨你們一切的不義.
提到赦免和洗乾淨.
第二節看到漢娜很深刻.

$^{361}$我從小時候五歲去教會時經常聽.
看看我在門外站在門外敲門.
如果有聽見我的聲音就開門.
我要進到他那裡去.
我與他他與我一同坐直.
這節金句很多時候是用來報道會用的.
誰決志了 舉手吧.
耶穌基督進入你的心裡.
但你知不知道這節聖經是在.
耶穌對七個教會講的裡面.
其中一間教會裡面講的.
即是他不只是對福音朋友.
他是對我們基督徒講的.
對教會講的.
原來有人來了教會.
有人做了基督徒.
耶穌基督還在門外.
他要不斷在敲門.
我們有時聽不到.
有時聽到當沒有聽到.
但他說如果你聽見我的聲音有開門.
我就可以進到你那裡.
我與他 譬如我開門.
他與我一同坐直.
其實是真的 一個信了幾十年的基督徒.
可能他與主耶穌的關係都是門外門內的關係.
即是你有需要求他 你開門.
耶穌基督有事要求你.
麻煩你進來.
我告訴你 那你幫我做吧.
但沒事了 拜拜耶穌.
你可以去做其他事 去找其他人.
我可以了.
但耶穌基督這裡說.
祂經常很希望與我們有一個坐直.
即是有一個生命的交流.
甚至祂讓我們有豐盛的生命.
耶穌說我來是要叫人得生命.
並且得更豐盛的生命.
得豐盛的生命.

$^{401}$得更加豐盛的生命.
今天究竟我在哪一個階段呢.
我是否得了生命.
但我未有豐盛的生命.
更加未有更豐盛的生命.
我不知道是哪一位.
總之我一來到Torrance的教會.
一過Easter.
第一次的復活節.
就有人買了門票.
讓我去Rolling Hills Covenant Church.
在山上的教會.
看Pageant.
那個活畫.
我傻到人家說是Pageant.
Pageant就去看 活畫就活畫.
我一路由頭到尾.
看到尾.
到那個最後晚餐.
我才知道原來是真人.
因為他做到真的.
好像一幅畫一樣.
那些人可以忍住呼吸.
不呼吸的.
釘在十字架上的那位耶穌.
做耶穌的那位.
他真的可以不用呼吸的.
他還要等別人唱完一首歌.
還很長的歌.
我第一次去看的時候.
三十多年前.
我真的很impressed.
很有感動.
這個教會動員多少人.
前面的 幕後的.
音響的 燈光的.
唱歌的 施班的.
播百人.
幾十人 一百人都有.
很多人.

$^{441}$那裡提到一樣東西.
我最深印象.
不是釘十字架復活.
而是那幅圖畫.
耶穌在敲門.
敲門.
我們就敲門.
原來那道門是沒有開關的.
沒有門把的.
外面的人進不來.
是要裡面的人開門給你的.
所以為什麼耶穌在門外敲門呢.
他隨時可以隱形的進來.
他也可以按那道門進來.
耶穌沒有.
他敲門.
因為主權在你那裡.
看看.
耶穌在你門外敲門.
今天我都要每天聽耶穌.
因為有時候我都忘記了祂.
祂在門外.
我在門內.
我在教會的四面牆裡面.
但是耶穌在門外說.
不行的.
你只是在四面牆那裡.
受到保護.
哇 這麼舒服.
外面有很多迷失的羊.
你快點去將牠帶回來.
我都要聽.
第三.
看看.
在小路21章第三節.
看看.
神的帳幕在人間.
祂要與人同住.
他們要作祂的子民.
神要親自與他們同在.

$^{481}$作他們的神.
哇 有歸屬的.
有神 有子民.
那種關係真的很親密.
今天.
我都不想多說.
但是你想一想.
當你去到一個地方.
現在我們在美國.
或者你在香港.
或者在英國.
或者在哪裡都好.
你去到那個地方.
你沒有歸屬感的.
你沒有一位保護你.
供應你.
為你著想的一位王也好.
總統也好.
什麼都好.
你在那裡真的很受罪.
所以很多人在英國回流回香港.
因為他覺得.
原來都不是這麼好的.
回來.
今天我們的天父是最好的.
祂不單單是王.
是主宰.
是創造主.
祂更加是我們的神.
要親自與我們同在.
不是找替身.
我跟你說.
今天我們做父母.
經常找替身的.
你知道嗎?.
你想出去逛街.
吃飯 看戲.
就找個寶寶來幫你看.
沒錢請寶寶.
不要緊.

$^{521}$叫誰?.
阿爺 阿麻 阿公 阿婆.
甚至當時她只有十三歲.
現在是不是十三歲我也不知道.
留在家裡.
我給你很多打機的遊戲.
打一打 看看你贏不贏.
回來告訴我.
就是找人來做替身的.
因為你自己不想在那裡.
但是你要記住.
就算所有替身都走了.
你依然身邊有一位神.
那位神.
祂就像帳幕在人間.
你搬不動的.
那位神經常在那裡.
你和祂在一起.
不過有時我們不讓祂來.
或者我們走了去.
對不起 走了一節.
看看22章第七節.
「我必快來」.
凡是遵守這書上預言的.
有福了.
這裡說耶穌基督快來了.
現在的世界.
所有的情景.
所有的發生的事.
越來越緊張.
但這個其實叫我們不是睡不著覺.
不是吃不下飯.
不是一天閉上眼睛.
錢不知道怎麼投資.
去這裡好還是那裡好.
不是要去做這些事.
最重要的是什麼呢.
提醒我們主必快來.
主回來的日子很近很近.
就在門口.

$^{561}$好了 現在我們要看.
其實我今天主要還要和大家看幾點.
就是為什麼這個王要來到這裡呢.
這個王是怎樣的王呢.
我們應該怎樣去回應這個王呢.
所以第一點就是說.
你的王為什麼要來到你這裡呢.
其實這裡第一件事大家知道.
應驗預言.
在馬太芬那裡就說了.
這件事情就是要應驗先知的話說.
什麼事呢.
就是我們剛才所讀的.
耶穌基督榮耀地進耶路撒冷.
但聖經裡面有很多關於預言應驗的.
我看了一下.
但是太多沒辦法今天在這裡說.
我還在想和魏長老商量.
或者我們主學就是專門讀這件事.
很多啊.
我要印出來都有幾頁紙.
所有聖經裡面說的預言都要應驗.
這個是必然的.
這個應驗是第一次耶穌降臨.
就是撒迦利亞書那裡.
告訴我們耶穌來是要帶給我們盼望.
耶穌來是要做我們的君王.
而且祂要拯救所有的百姓.
哎呀.
怎麼會這樣.
Sorry.
太大力.
第二.
這裡說.
祂來就是要.
這裡說到撒迦利亞書那裡提.
祂是私行拯救.
第一次降臨.
耶穌基督就是在馬槽.
我想了一下.

$^{601}$其實耶穌基督記載祂做王有多少次.
或者說我們對耶穌是王有什麼認識呢.
在祂未出世之前.
天使已經和瑪利亞說.
祂要做王.
祂坐雅各的王位.
當然再推上去就是大衛的王位.
他們都是永生神派祂來的.
祂是君王.
是和平之君.
就是和平的君王.
到祂出世了.
那些博士去哪裡找祂呢.
是去皇宮找祂.
也就是說.
他們都讀聖經.
知道這位彌賽亞如果要降生.
就是在皇族的.
到祂一路長大.
祂五餅二魚.
吃飽了近一萬人的時候.
那些人是怎樣.
來擁戴祂做皇帝.
甚至有一點點強迫性的.
我們擁納你做皇帝了.
耶穌那時候怎樣回應呢.
祂一個人獨自上山.
叫門徒先回去.
你過這條海就先回去.
我要上山祈禱.
祂面對這麼大的試探.
祂獨自在神面前說.
我要馴服你.
不是用這個方法.
就是那些人來擁戴祂做皇帝.
所以祂來是要私行拯救.
祂不是要來做王.
去統治這個世界.
或者統治祂的百姓子民.
其實在舊約聖經裡面.

$^{641}$選納皇帝不是神的心意.
是人的要求.
所以當那些百姓要求.
我要皇啊 我要皇啊.
三母醫跟他說了.
你是不是真的要呢.
這是不是神的心意呢.
而且說明當你要納皇的時候.
你要記得.
祂做了你的皇.
祂就會奴役你.
祂喜歡就拿了你的房子.
田地.
你的子女去服侍祂.
你不可以拒絕.
因為祂是皇.
祂叫他們想清楚.
但是當時聖經說那些以色列人看到.
每個國家都有皇帝.
為什麼我們沒有呢.
不行 我一定要有皇帝.
結果就有索羅王.
有大衛王.
一路就下去了.
所以其實選納皇.
除了耶穌基督之外.
神真的告訴我們.
祂是我們的萬王之王.
萬主之主.
這位王除了來應驗預言.
和施行拯救.
除了祂之外.
沒有別的救法.
那他是怎樣的一位王呢.
他是公義和平的.
看看你的皇來到你這裡.
他是公義的.
並且他要施行拯救.
做一個公義的皇很難的.
所以今天我們一打開報紙.

$^{681}$現在不看報紙的人.
一打開網站.
或者是Facebook什麼地方一看.
到處都有人推翻那個王.
想推翻他們的首領.
為什麼呢.
這一邊是你們選的.
但是沒多久.
你覺得他.
第一年就覺得他不滿意他了.
我要推翻他.
我不要王 我不要王.
想推了他之後.
又選納第二個人來做.
所以其實.
為什麼人們會這樣變呢.
去年是這樣 今年又這樣.
前個月是這樣.
今年這個月又這樣呢.
是因為他們覺得領袖不公義.
不公平.
沒有愛心.
也不是來救民的.
可能是來殘害人的.
他們就要推翻他.
但是你記得.
世界上只有一個君王.
他是你最滿意的.
他是公義和平的.
就是我們的耶穌基督.
另外他是溫柔謙卑的.
他謙謙和和的騎著驢.
就是騎著驢的驅使.
所以如果你看那些宗主主日.
耶穌騎驢子進來京城耶路撒冷.
最準確的圖畫是.
兩匹驢一齊和耶穌進京城.
一匹是媽媽的驢.
因為原文用的字是女性.
一匹是兒子.

$^{721}$耶穌是坐驢驅.
聖經說驢驅是沒有人騎過的.
為什麼耶穌要騎驢驅呢?.
他可以直接騎馬更快.
他是王.
但是他告訴人.
他來不是要用刀兵.
不是要用戰爭來做王.
他是謙謙和和的來.
而且他還選擇騎一隻最小的.
是沒有人騎過.
這個小驢驢驢.
為什麼要有媽媽陪著呢?.
因為他沒有被騎過.
他會害怕.
媽媽就在他身邊陪著他走.
我們去過聖地的人就知道.
最深印象.
當我們要走上修道院.
走很多樓梯.
但是走得快的.
他半小時就能上去.
走得慢的就要45分鐘.
但是我和另一個姐妹.
因為她心臟不太好.
心跳也不太好.
走著走著.
本來就說騎驢驢.
真的有驢驢在背著人上去.
但是騎驢驢又要小心.
那條路很窄.
上山的路很窄.
如果她受怕的話.
她就會瘋了.
要不就迫你去岩石旁邊.
要不就掉下去了.
所以你要騎驢驢.
你都要夠膽才行.
後來上的時候不敢騎.
下來平地.

$^{761}$因為我和姐妹實在太慢了.
最後一個上去.
人家都要下山了.
我們才追車.
我最記得追來追去.
追不到的.
就是追別人走的那些人.
結果真的坐一程雷仔.
被那個人拉我上去坐在那裡.
其實很害怕.
沒有媽媽路在那裡.
結果有個人在旁邊.
拉著她的繩子帶著她走.
都很害怕.
後來去到一段不准雷仔走的路.
就要轉馬了.
我們又騎馬了.
因為要追回巴士.
他們要開的.
結果追到去.
一輛巴士都開了.
另外第二輛在等我們.
所以我們倆就慘了.
令到人家遲到回酒店吃東西.
那時候就有機會坐一下.
原來不簡單.
你要爬上去那一刻是最難.
所以坐驢仔.
那個人捧了我上去.
騎馬要去到一個牆邊.
一個矮的牆邊.
我要站在牆邊再爬到馬上.
但耶穌不是來騎馬.
他是選擇騎雷仔.
而且是一隻未被人騎過的.
耶穌就是這麼謙和.
你知不知道.
當耶穌在世界上做人的時候.
祂不但為門徒洗腳.
那一幕是很謙卑的.

$^{801}$跪在地上.
拿著一塊毛巾.
幫他們逐個洗完.
又抹乾.
很謙卑.
但你知不知道.
還有很多場景裡面.
祂還是很謙卑的.
譬如祂去探訪人.
祂真的是願意用手去扶那個人.
用手去拉那個人.
死了的那個女孩.
艾魯那個女兒死了.
那些人在哭喪.
但耶穌進入了她的家裡.
為這個十二歲的女孩祈禱.
祂就用手去拉她起來.
那個女孩張開眼睛.
耶穌就扶她起來.
彼得的岳母也是.
當耶穌為她祈禱.
祂害了熱病.
為她祈禱.
讓她快點康復.
扶她起來.
那些患大麻瘋病的.
那些盲眼的.
耶穌真的用手去摸他們.
為他們祈禱.
用腋下更有趣.
用口水弄點泥.
然後塗在他們的眼睛上.
那個人叫他去河裡.
去那裡洗一下.
結果他看到了.
所以耶穌是親力親為的.
祂不是說這個骯髒.
那個骯髒.
祂就不做.
耶穌也是這麼謙卑的.

$^{841}$當祂受審.
被捉拿受審的時候.
人家誣告祂很多事情.
打祂.
滴口水在祂的臉上.
但是.
就像聖經所形容的.
祂好像一隻羊.
在展毛的人手下.
是無聲的.
今天神所要的.
神要世人.
要什麼呢?.
要行公義.
好憐憫.
全謙卑的心.
與我們的神同行.
耶穌基督這位王.
是公義的.
是和平的.
是好憐憫的.
是溫柔的.
謙卑的.
與神同行.
三十三年半.
與神同行走到最後一刻.
到祂斷氣之前說.
成了.
我將我的靈魂交給你.
這個是不是.
我們一生要去追求的呢?.
啊.
你的王是怎樣一位王呢?.
榮耀德盛的.
除了公義和平.
除了謙卑溫柔.
祂是榮耀德盛的.
因為那些前行後退的眾人在哭著.
哭著大聲的說.
和珊啊.

$^{881}$求求你救救我們.
當時他們真的需要救的.
為什麼呢?.
他們在羅馬的政府.
那個國家的統治下.
是很辛苦的.
但是呢.
他們以為耶穌來就是幫他們解決政治的問題.
就是可以換一個朝代.
換一個國家.
但是他們不知道耶穌來.
不是這樣的.
但是他們依然看到耶穌是很榮耀的君王.
祂將來是要德盛的.
他們就稱祂說.
奉主名來的是應當稱頌的.
高高在上和珊娜.
祂是將耶穌基督擺在最高的.
好了.
我們最後要想一下就是.
我要怎樣去回應這一位來到我這裡的王呢?.
第一.
應當大大的喜樂歡呼.
這裡說.
錫安的民啊.
應當大大喜樂.
耶路撒冷的民啊.
應當歡呼.
臣.
祂的兒子降臨.
其實這是一個莫大的恩典.
臣的兒子親自降臨.
還有就是.
祂是帶來救贖的.
是給我們有永恆生命的確具的.
祂是叫我們聚得到釋放的.
可以過一個自由的生活.
今天你想想.
你坐在這裡可以自由敬拜.
這個不是每個人都可以的.

$^{921}$有人這一刻可能躲在山洞裡.
祈禱而已.
有人可能真是.
都不知道捐到哪裡去了.
只能夠在一個黑漆漆的地方.
打仗的地方.
隨時被人轟炸的地方.
去呼求神.
大聲的去求神.
救我吧 救我吧.
但今天我們還仍然可以坐在這裡.
是不是應該要更加歡喜快樂.
今天我開車回來教的時候覺得.
應該要喜樂.
因為詩篇都說.
今天是耶和華.
為我們所定的日子.
我們要在其中歡喜快樂.
這是一天.
神所做的一天 我會享受和開心.
你一定要歡喜快樂.
我一路開車來.
其實我開車來是十分鐘之內.
七分鐘.
我就一路重複唱這首歌.
神啊 我不夠喜樂.
你給我喜樂吧.
因為今天是主人.
是你為我定的日子.
你要將一切放開.
交托.
然後去敬拜神.
我們讀一讀《路加福音》.
《路加福音》有另一種說法.
很有趣.
不是很有趣.
是大家互補的.
弟兄們 請你再幫我讀.
他們牽道耶穌那裡.
把自己的衣服搭在上面.

$^{961}$扶著耶穌騎上.
走的時候.
眾人把衣服鋪在路上.
將近耶路撒冷.
正下橄欖山的時候.
眾門徒因所見過的一切.
而能都歡樂起來.
大聲讚美神.
這裡是《路加福音》補充的.
那些人不但把衣服鋪在雷仔道.
地下道.
又找些樹枝.
Palm Tree那些樹枝.
在那裡養一養.
他們還在那裡歡樂起來.
大聲讚美神.
這裡請姐妹讀.
那裡怎麼說.
奉主明來的使益.
在天上有和平.
在至高之處有榮光.
眾人中有幾個法利賽人對耶穌說.
夫子.
擲避你們的門徒吧.
耶穌說.
我告訴你們.
若是他們閉口不說.
這些石頭必要呼叫起來.
這種讚美的稱頌.
在天有和平.
在至高之處.
有榮光.
他們奉主明來的王.
路加福音講明是王.
而且是和平的王.
是榮耀的王.
還有些法利賽人經常妒忌.
妒忌那些人.
那麼擁護耶穌.
夫子.

$^{1001}$耶穌.
你快點擲避你們的門徒吧.
在那裡吵吵鬧鬧.
大聲喊.
耶穌怎麼說.
我告訴你們.
若是他們閉口不說.
這些石頭必要呼叫起來.
今天如果你已經經歷神很多恩典.
但是你從來閉著口不說.
你小心.
當你走路的時候.
那些石頭說.
為什麼你不感恩.
可能突然有隻蟲鼠出來.
吱吱吱吱地喊.
為什麼你不讚美.
即是說.
大自然界的事物.
它們都會知道.
有位創造主.
要將榮耀歸給主.
所以上個星期我們很感恩.
有三位弟兄姊妹在這裡做見證.
嘩!好美呀.
有些人聽完之後告訴我們.
很有感動.
還有其中一個.
講出了他心裡面的說話.
今天我們大家都可以做見證.
不用上來這個臺才做.
當你每天在街道上買菜也好.
或者你接小朋友放學.
孫兒也好 子女也好.
你遇到人你都可以說.
多口就在這裡要多口.
平時就不要那麼多口.
駁人 罵人那些就不要了.
但是你能夠講見證和分享信仰.
多一點聲音都不怕.

$^{1041}$我們這裡坐在這裡的人.
有些已經走了的人.
都是因為在市場 在街道上.
有人跟他們講福音.
請他們回來教會.
他們就回來了.
甚至信了主.
所以我們要大大的歡喜快樂.
第二.
應當要對王.
我是不是又快了.
順服 完全的順服.
這裡提到什麼.
他說耶穌和門徒將近耶路撒冷的時候.
就到了白髮旗.
在橄欖山那裡.
發生了什麼事.
弟兄們 你來幫幫我 我們可以很快完.
耶穌就打發兩個門徒.
對他們說.
你們往對面村子裡去.
必看見一匹驢山在那裡.
還有驢驅童在一處.
你們解開 牽到我這裡來.
若有人對你們說什麼.
你們就說主要用他.
乃人必立時讓你們牽來.
要對王完全的信服.
王怎麼說.
你們要去 往對面的村子去.
指明地點 你要去.
必看見 你要看.
不然就這樣走過就算了.
還要解開.
將驢驢驢驅牽來我這裡.
還要有人問你想偷東西.
你告訴他主要用他.
他就會立時讓你牽來.
每一步你都要跟著這個王的吩咐去做.
好了 姐妹來了.

$^{1081}$門徒就照耶穌所吩咐的去行.
牽了驢和驢驢來.
把自己的衣服搭在上面.
耶穌就騎上.
這裡說什麼.
門徒就照耶穌所吩咐的去行.
今天我們有沒有照耶穌的吩咐去行呢.
我們的兒女 孫子.
有時你請他做一些事.
不要叫他 請他做.
他都不一定全部做完.
做些不做些.
我們自己做別人的兒女.
或者做下屬.
上司叫你做的時候.
我們都甩甩流流.
做不完.
但是你看看這些門徒.
他們對主的回應是.
完全的順服.
完全的聽耶穌基督的吩咐.
不懷疑的就去吧.
所以今天 當你出街之前.
神說我今天要給你一個福音朋友.
你記住了.
你會見到哪一個人.
你是不是在買西瓜.
你拍拍拍.
你過去跟他談.
你說選西瓜是這樣選的.
你就會跟他做朋友了.
你肯不肯呢.
如果神真的這樣叫我.
我都願意去做.
第三.
這裡怎麼說.
應當效法君王的榜樣.
這個都是重複說的.
就是說.
祂是騎著驢來的.

$^{1121}$是溫柔的.
還有耶穌基督自己說了一句話.
祂說我心裡柔和謙卑.
你們當乎我的額.
學我的樣式.
這樣你心裡就得到安息了.
今天如果我們一直都是驕傲自大.
難怪我們裡面是沒有安息的.
就是沒有平伏在裡面.
因為你經常覺得.
還有哪一個人比我厲害.
還有哪一個人比我威風.
還有哪一個人未曾聽我話.
你心裡面就很不服氣.
哪有安息呢.
哪有心安理得呢.
但是如果我們學耶穌基督.
溫柔謙卑 柔和謙卑.
真是願意放下自己.
真是願意不要去計較.
跟跟計較.
這個呢.
就會叫你心裡面很平靜安穩.
反省一下.
在現實的生活中.
究竟哪一個是你生命中的王呢.
你的父母.
有人真是我佩服有一些兒子或者女兒.
他為了要照顧他的父母.
他真是放下了自己.
一路都是陪著父母.
特別是單親那些.
一路陪著他住 一路服侍他.
到最後他要在療養院.
他都每天早上晚晚去服侍他.
幫他刷牙洗臉.
幫他戴上假牙.
幫他做任何他做不到的事.
穿衣服那樣.
我很佩服那些很孝順父母的人.

$^{1161}$你這樣做 你記得你的子女都會這樣對你.
因為他們有眼睇.
你的孫都會這樣對你.
因為他看到這個婆婆媽媽.
是對上一輩那麼好的.
佩偶.
佩偶是不是你生命中的王呢.
她說什麼你就做什麼.
你做慢一點都不行.
因為她是王來的.
你一路都是為了要討好她.
甚至為她犧牲了耶穌.
她不喜歡回教會 我陪她不回.
或者她不喜歡我服侍她 我就不服侍她.
佩偶是不是你心裡面的王呢.
你自己當然是.
其實說到這裡.
不是父母媒佩我 是我們自己經常都做王.
有一首歌 昨天我在這裡找到.
就是汪姑姑 我們以前有一個女傳道.
活到九十多歲的施鳳臣.
幾十年了 她經常唱在台上.
她說我們心中有兩個王.
一個叫我上天堂.
一個叫我落地獄.
上天堂的那個就是我們的耶穌.
我們的王.
叫我落地獄的是誰.
魔鬼 她說就是魔鬼.
她說你要選哪個王呢.
你就會決定你去哪裡.
你可以上網 上Youtube那裡找這首歌.
很有趣的 那些小朋友唱得很好.
好了 還有兒女.
今天其實兒女 或者是兒孫.
都很多時候成為我們的王.
那個甚麼王 小霸王.
真的是小霸王.
他要吃就吃 他不吃就不吃.
他不吃就丟掉.

$^{1201}$現在的小朋友 我不知道怎樣學回來.
可能是打機的時候學.
又可能是看電視學.
沒有大沒小的 你不覺得嗎.
你叫他就可以 你叫他叫你一聲.
他不理你.
今天這個小霸王 小心.
他搶了耶穌的地位.
一些英雄的人物.
有人很崇拜誰.
崇拜偶像.
這個是英雄來的.
但是崇拜了 跟了英雄不久.
這個不行 要換另一個.
他那些偶像是經常換的.
但今天 耶穌是不是你心中的君王呢.
如果是的話 你今天要怎樣降伏在祂那裡.
並且要去敬拜祂呢.
敬拜祂.
你知不知道這班群眾.
將耶穌擁戴為王的那班群眾.
過了幾天.
走了不止 還在大聲叫甚麼.
釘他十字架.
這邊擁護你做王 那邊要釘死你.
所以今天如果有人擁護你做皇帝.
你不要開心.
可能幾天之後.
他又說釘死你.
所以其實你的兒女 孫子.
你都不要把他捧得太高.
他現在就.
到大一點 找都找不到他.
找大學要去最遠的.
然後結婚要搬到最遠的.
你要去探孫子.
make appointment.
不是沒得去看.
嘩 我真的覺得這不知道是甚麼世界.
求神憐憫.

$^{1241}$不要讓你自己做王.
也不要讓你身邊的人.
事物.
做了你的王.
請大帝和我們一起祈禱.
天父我們很感恩.
因為你就是那位最榮耀的君王.
最謙卑的君王.
而且你是那位.
有憐憫 有同情.
是寬恕我們的君王.
不是要奴役我們.
不是要品榜我們的君王.
所以我們願意來到你面前.
將最高的榮耀歸給你.
但願平安 喜樂.
還有盼望.
還有那種熱情.
熱心去愛神愛人.
都歸給我們在地上的每一個人.
奉耶穌基督得勝人名求.
阿們.
\newpage



\section{出埃及記 3:12}
\label{sec:031895}
\textbf{聖經人聖經事:我必與你同在}~~~~黃克勤牧師博士
\newline
\newline
連結: \href{http://www.church.com.hk/acms/content.asp?site=cdc\&op=show\&type=product\&code=031895\&layout=sermon}{\texttt{031895}} ~~~~ 語音日期: 2026-03-29
\newline
\newline
\hyperref[sec:031894]{\small{< < < PREV SERMON < < <}}
~
\hyperref[sec:index_chronic]{\small{[返順時目]}}
~
\hyperref[sec:031896]{\small{> > > NEXT SERMON > > >}}
\newline
\newline
出埃及記 3:12
\newline
\begin{longtable}{cl}
\hline
\hline
章節 & 經文 (和合本修訂版)\\
\hline
3:12 & \begin{tabularx}{0.7\textwidth}{X} 神說：「我必與你同在。這就是我差派你去，給你的憑據：你把百姓從埃及領出來之後，你們必在這山上事奉神。」 \end{tabularx} \\ \\
[1ex]
\hline
\hline
\end{longtable}
$^{1}$各位弟兄姊妹 眨眼又一日又讀聖經了.
今天聖經人 聖經事 是生命說書的網絡為大家預備.
我必與你同在 很安慰的說話.
今天神是跟你說這句話的.
記不記得昨天.
昨天是我原知道他們的痛苦 我知道的.
今天就是說 我必與你同在.
好 我們看下去.
出一句第三章第九節.
現在以色列人的哀聲是達到我耳中.
我又聽見埃及人怎樣欺壓他們.
由於有些綠色字 是我的注釋.
藍色字是希伯來文沒有的.
不過中文有 希望幫助我們更加明白.
現在以色列人的哀聲是達到我耳中.
我又看見埃及人怎樣欺壓他們.
神和摩西這樣說.
現在以色列人的哀聲.
這個哀聲其實是很慘 很痛 很餓.
我們以為哀聲是很辛苦的聲音.
其實不是.
這個哀聲的意思是喊叫.
我都聽過很多次.
有些弟兄姊妹入院.
我都入過醫院住過.
但是我住的時候病房環境相當安靜.
但是我聽過很多在不同的醫院.
特別是公家的醫院.
很多病人很擠迫.
整晚都睡不著 為甚麼.
他說我病房裡每個人都嘔吐聲.
那些病人就大叫.
特別是長者病房有很多大叫.
原因是他們覺得自己孤單沒人理.
所以他們會大叫.
由於中風不能動.
躺在床上會被壓住.
如果把床調斜一點.
病人會慢慢摔下去.
他們就很不舒服.

$^{41}$又或者躺在床上被衣服壓住.
壓得很辛苦 背部又痛.
動不了就開始大叫.
不是嘔吐聲 是大叫.
有兩個幾晚都沒睡.
大聲大叫.
大聲大叫就是這個.
哀聲看中文就是小聲.
哀聲是大聲大叫.
現在以色列人大聲大叫.
是達到我耳中.
達到是進入.
直接翻譯就是.
以色列人的慘叫進入.
進入甚麼.
希伯來文是這樣說.
原文就是.
現在以色列人的慘叫進入.
一句就是這樣.
於是翻譯的人加了少少字.
加了甚麼字呢.
進入我耳中.
第一沒有「耳」字.
也沒有「我」字.
不過現在以色列人的慘叫進入我.
這個一定是.
進入甚麼呢.
進入我那裡.
在哪裡進入呢.
我們自然要想由耳仔.
所以翻譯加了「我耳中」.
現在以色列人的慘叫進入.
即是神是說以色列人的慘叫進入.
因此你們一定要記得.
昨天說的題目是.
我原知道他們的痛苦.
神是知道我們的痛苦.
其實很多時候痛苦有人知道.
真誠地知道和聽.
已經減了一半痛.

$^{81}$有人知道.
不是孤單一個人就很痛.
除了我聽到.
我又看見埃及人怎樣欺壓他們.
看得「raa」.
還有很多次.
我看到看到.
昨天是這樣說的.
我實在看到.
其實我「raa」「raa」.
我看到看到.
說兩次看到.
即是說非常之看到.
我看到埃及人怎樣欺壓.
欺壓呢.
「Lakats」.
「Lakats」呢.
「Lakats」的意思是際語壓.
這個際語壓.
就好像我們壓蔗渣.
就是擠壓.
就真是這個意思.
就不是說我欺壓他.
欺壓呢.
我們看中文字就是欺壓.
「Lakats」是捏他.
好像蔗那樣榨汁.
好像壓橄欖油那樣.
壓橄欖.
好像葡萄那樣踩扁他.
踩爆他.
然後流汁出來.
這個是用在.
擠壓「Lakats」這個字.
是用在攪那些.
水果,食物那裡.
就不是指剝削.
而是指壓他.
我看到埃及人怎樣.
去擠壓我的百姓.

$^{121}$第十節.
留意了.
故此我要打發你去見法老.
使你可以將我的百姓以色列人.
從埃及領出來.
故此這個字.
就譯得不好.
其實這個字是「Eta」.
「Eta」是現在.
所以第十節開始就是.
現在我要打發你.
就不是故此.
因為故此呢.
就是我聽到他們被擠壓得很慘.
所以故此我就.
所以我派你去.
不是.
沒有的所以.
而是現在我聽到.
現在我就派你去.
第一點今天要和大家分享的.
就是我要打發你.
神要打發你.
打發摩西.
我要「Shalak」.
差遣你.
就是這個字.
差遣呢.
「Shalak」.
我差遣你.
什麼叫「Shalak」都可以.
「Shalak」.
差遣你去見法老.
使你可以將我的百姓以色列人.
從埃及領出來.
領呢.
就是大家跟著我.
我帶你出.
「Eta」.
「Eta」是前往.

$^{161}$你可以將我的百姓.
由埃及.
帶他們前往.
不單只是出來.
是領他們前往.
領他們前往.
不是出來了.
出來了就.
哦.
放監.
出來站著.
是前往.
去一個地方.
出來去一個地方.
第二點呢.
摩西就說.
派我.
你不是吧.
這個「派」這個字.
打發.
「Shalak」.
你差遣我.
差派我.
去見法老.
不是吧.
我在皇宮出生.
哦.
皇上.
法老那些排場.
多大啊.
他統領的事情.
多多啊.
我基本連皇宮門口.
都進不了啊.
而且現在.
做法老那個人啊.
我認識他的.
他跟我一起.
長大的.
他那麼大.

$^{201}$我那麼大.
大家以前是.
皇子兄弟.
不過他是.
純種埃及王的皇孫.
我呢.
就是以前法老的女兒.
修養的養子.
但是我們是同樣在皇家學校畢業.
一起在那裡讀書.
一起學武功.
一起學軍事.
一起學算術.
一起學天文.
一起學藝術.
一起學語文.
我們在那裡出來的行政管理.
我們每個人都知道的.
我還是一個.
差點被他想要殺.
因為我殺了埃及人啊.
我現在基本上是有罪在身.
我連.
埃及宮門口.
我都去不了啊.
我現在一回埃及.
已經一入境就被捕了.
連這個入境處.
我都沒過到啊.
所以你叫我去.
摩西就對神說.
我是什麼人.
竟能去見法老.
將以色列人從埃及領出來.
這個問號真的要紅色才行.
我怎樣有可能帶到他啊.
帶到以色列人走啊.
不要說過百萬以色列人.
又沒有手提電話.
怎麼帶啊.

$^{241}$又沒有WhatsApp.
怎麼通訊啊.
怎麼叫他走啊.
再加上我連皇宮都去不了.
去見法老.
他怎麼肯見我啊.
我已經坐牢.
頭都批了.
這裡留意了.
第十二節的上半節神說.
我必與你同在.
今天的主題就在這裡.
我必與你同在.
等一下我要和大家分享第二點.
哎呀忘記了這個螢光.
我必與你同在.
神原知道我的痛苦.
然後他再說.
我必與你同在.
今晚.
今天大家一定要記得這一句.
我必與你同在.
希伯來文就在這裡.
讀多一次.
大家可以試試讀.
因為很容易讀.
都是英文拼音.
因為他譯了出來.
英文拼音我們容易讀.
必定.
我將會.
I will be.
我將會做的.
吾同在.
with.
self.
就是這個.
this.
記號.
of.

$^{281}$a sign.
記得.
我必與你同在.
如果將他翻譯成這一句.
就是.
必定我將會同在.
這個記號.
不確定他說什麼.
其實就是說.
我必定將會與你同在.
而我與你同在這件事.
就是我給你的一個記號.
下面的.
這裡是上半節.
下半節將這個記號譯成證據.
或者證明.
就是一個sign.
你叫我做這件事.
可不可以?.
可以.
你有沒有一些證明證據給我?.
有.
如果我馬上令到你.
由這麼瘦弱的人.
在一秒鐘變成非常灰不健碩.
真的有肌肉.
嘩.
這樣都可以.
如果我和你做這件事.
你就知道我給了你證據.
是我與你同在.
我與你同在.
我必定會與你同在.
其實.
當你要做一件事.
你覺得自己不可能.
什麼叫我去見法老.
不可能啦.
no way.
這樣不行.

$^{321}$不會可以.
不單止不可以.
是不可能.
不過.
神就說.
祂給的答案是什麼?.
一定可以.
祂不是說.
我是要與你同在.
當我與你同在.
就表示一定可以.
這個就是一個記號.
我與你同在就是一個記號.
你記得了.
所以無論什麼時間.
無論多艱難.
見到自己教好不行.
你都要記得.
神是說.
我必與你同在.
我自己都是靠神的同在.
靠祂的應許.
走這個事奉的道路.
走了四十多年.
好像抗夜經歷.
人生抗夜經歷都走完了.
走到現在差不多見主了.
都是神的同在與供應.
養著.
沒有一千也有八百個小朋友.
無論各樣的方式.
怎樣養?.
一個月加起來.
沒有六七十萬港幣.
一個月是不夠的.
另外有一百個傳道人.
另一個機構做宣教的.
沒有四五十萬.
一個月不夠.
一個月一百多萬.

$^{361}$那些小朋友又不夠錢.
那些孤兒還是愛滋病.
我不吃這個給他.
給不了那麼多個.
又有印度.
每一次去到這些時候.
我就會感覺到上帝的同在.
為什麼呢?.
因為祂會猜拳在聖靈跟我說.
你要記得.
聖經說.
耶和華是孤兒的父.
不要以為自己很了不起.
以為自己做了一個父親.
他們是叫我做父親的.
一開始叫我做Daddy.
叫了一下都不叫了.
現在叫我做Da.
Da.
他們是Ta.
Ta是讀D.
Da.
Da即是祖父.
開始叫我阿爺.
阿公阿爺.
叫我已經是Grandfather.
這樣意思.
現在我這樣穿鞋.
手舉起平衡.
他們馬上走過來扶我.
那些小朋友.
很可愛.
但是我想一下.
我那麼少錢.
一個人怎樣都可以.
上帝說你放心.
祂不是你的兒子.
祂是我的兒子.
耶和華說.
我是孤兒的父.

$^{401}$如果你奉獻給主.
憐憫貧窮是借給耶和華.
祂必償還.
你多有錢.
你都借不到錢給神.
不過.
祂說憐憫貧窮是借給耶和華.
這些說話.
是很安慰很鼓勵的.
這個說話也是.
一直鼓勵著我.
我必與你同在.
這就是一個O.
一個記號.
信者說.
我必與你同在.
我讀了.
你將百姓從埃及領出來之後.
你們必.
第三點要和大家說的.
在這山上事奉我.
這山上是何列山.
後面叫西乃山.
你去到西乃山.
在那裡來服侍敬拜我.
你帶他們出來.
見法老帶他們走.
見法老怎麼見都見不到.
我與你同在.
這就是一個O.
證據.
證據這個字.
很可惜很有趣.
在第十個字.
這個證據.
就是我與你同在.
我與你同在.
這一句.
就是.
我必定與你同在.

$^{441}$這就是一個證據.
我與你同在就是.
我要帶你.
我會帶你去見法老.
然後他會讓你走.
這就是證據.
然後.
你會帶他們來這山上事奉我.
這就是我打發你去的證據.
這個證據這個字.
O出現在第一句這裡.
讓大家知道.
這個都很重要.
希望大家能夠記得.
做無論一件大事還是一件小事.
一個大的事奉一個小的事奉.
當你想.
我哪裡行.
我是誰人.
這件工作怎麼做.
你依靠神.
神就說.
我與你同在.
就是一個證據.
告訴你你做到的.
這個你要更有信心.
神歡喜與你同在.
非常感恩.
我有機會和大家預備.
這些聖經人士的時候.
看看亞伯拉罕.
以撒雅各約瑟道摩西.
真的學到很多東西.
這個靈性上.
功力都真的進深了一點.
很感恩的.
大家一起多些學習聖經人士.
看看背後上帝的故事.
他怎麼工作.
願主得著榮耀.

$^{481}$我們一起禱告.
多謝我們的主.
你是說你知道知道.
你聽到聽到.
知道我們的困難與痛苦.
因此現在就打發我們去做一些事.
讓我們勇敢站起來.
這些怎麼做.
好像摩西當日.
他真的嚇一跳.
什麼?.
我這個殺人犯.
埃及都不能回去.
要我回去帶那百多二百萬人.
以色列人出來.
不可能.
不過神就說.
有的.
我給你證據.
證據是什麼?.
上帝.
我必與你同在.
這樣都是證據?.
說而已.
是呀,說而已.
上帝說就帶著能力.
正如在聖經裡面說.
出於上帝的.
沒有一句不帶著能力.
主要就是信你的能力.
信你的同在.
去到哪裡都好.
有什麼情況都好.
有你的同在.
就最重要了.
最好了.
多謝我們的主.
多謝你的同在.
你的同在.
就是你支持幫助.

$^{521}$愛我們.
關心我們的證據.
贊美主.
給我們這個證據.
贊美你.
禱告奉耶穌.
基督的名.
阿們.
\newpage



\section{箴言 27:17}
\label{sec:031896}
\textbf{多代同行 (2) PeerWalk IV Iron Sharpens Iron}~~~~沈炳華傳道
\newline
\newline
連結: \href{http://www.church.com.hk/acms/content.asp?site=cdc\&op=show\&type=product\&code=031896\&layout=sermon}{\texttt{031896}} ~~~~ 語音日期: 2026-03-28
\newline
\newline
\hyperref[sec:031895]{\small{< < < PREV SERMON < < <}}
~
\hyperref[sec:index_chronic]{\small{[返順時目]}}
~
\hyperref[sec:031897]{\small{> > > NEXT SERMON > > >}}
\newline
\newline
箴言 27:17
\newline
\begin{longtable}{cl}
\hline
\hline
章節 & 經文 (和合本修訂版)\\
\hline
27:17 & \begin{tabularx}{0.7\textwidth}{X} 以鐵磨鐵，越磨越利，朋友當面琢磨，也是如此。 \end{tabularx} \\ \\
[1ex]
\hline
\hline
\end{longtable}
$^{1}$今晚分享的信息就是.
你知道今年教會主題是多代同行.
我們是講第二個session.
就是在講藍.
你記得藍說PSWalk我們這次的系列.
PSWalk我問一下你.
PSWalk是喫什麼字的.
你沒有聽到嗎.
藍都暈頭了.
三個星期才起得來.
你看看多體貼這個牧者.
PSWalk其實說喫字.
他想這個系列的時候.
你知道藍想這些.
我經常贊他很厲害想這些東西.
他說喫pillow talk.
就是你的閨蜜.
就是你.
你為什麼叫做Peers.
就是你的閨蜜.
你的閨蜜就是你的好朋友.
我想作為回教會回了很久.
其實我們聖經當然講愛人如己.
是要講的.
但是當我們.
這個系列很小.
我們想去探討.
不是說只是在教會里面.
打好關係好來好去.
怎樣可以不那麼流於表面.
就是由一個可以盡心一點的.
成為一個做人真誠的朋友.
而我很深印象.
聽到藍說到的時候.
他說不要只看別人.
你怎樣可以由自己開始.
成為一個這樣的人呢.
跟著如果.
如果你有.
我可以轉過來.

$^{41}$如果你有印象.
跟著藍說他用了.
用了大衛和約拿丹.
講什麼是Godly Friends.
Godly Wisdom.
就是究竟兩個.
特別是約拿丹.
他本來是王子.
但是他怎樣去Honor大衛.
他在上帝裡面.
他很體重大衛.
Over了他自己的身份.
在那個關係裡面.
跟著再下一個禮拜.
就有Michelle.
Michelle來的時候.
用傳道書去講.
義人兩個人總比一個人好.
我們每一個人.
其實都需要有朋友.
需要有同伴.
需要有人可以將我們.
Offload了我們一些.
心裡面的擔子.
哪一個是你的同行者呢.
我很多年前.
都聽過Michelle去分享.
她這兩個的朋友.
這兩個的土伴.
都是成為她生命當中.
一個很重要的人.
跟著上一個星期.
就是默穆斯.
Gloria來分享多瑪.
這個多疑的多瑪.
在一個信仰的群體當中.
這個群體怎樣接納她.
而接納她的同時.
以致她能夠將這種真實.
即是耶穌怎麼樣啊你.

$^{81}$探她進入你的立堂.
那我纔信呢.
但這個群體怎樣仍然接納她.
以致在這個.
由福音書的尾段.
到這個兆行傳的頭段.
這個多瑪.
有一個生命的轉變.
PS Walk去到最後.
今天最後這個.
賣尾這一講.
一個真誠的朋友.
怎樣可以walk a step further.
又可以講真話.
又可以信得過呢.
通常很多不同的故事.
當然很理想.
搭著肩膀看著日落.
哇 這些就是friendship.
這些就是BBF.
這些這麼美好的環境.
但現實不是這樣的.
現實通常是.
你靠近的那些特別.
如果有些close的.
或者是你的家人.
或者是很熟悉弟兄姊妹妹.
通常是什麼呢.
就困了.
所以今晚很想和大家.
透過箴言一節.
主要一節的聖經.
和大家分享一下.
究竟好朋友.
怎樣互相去塑造生命呢.
箴言呢.
如果你回教會久了.
你知道箴言是整本聖經裡面.
其中一個大的段落.
在智慧文學裡面.

$^{121}$其中有幾卷.
箴言 約伯記 傳道書.
其實還有詩篇和雅歌.
但我抽這三個.
這三個其實都是講智慧的.
在聖經裡面其實都是講智慧的.
但它有少少不同的分別.
箴言的智慧是講一些日常的智慧.
或者是一個general speaking.
所羅門他說作為一個.
他經歷過很多.
在聖經裡面記載.
當時是最有智慧的人.
他寫了很多不同的說話.
但他講的是一些日常的智慧.
應該怎樣去做.
鼓勵你怎樣去做.
去拿到上帝的智慧.
約伯記是講一些特殊的情況.
約伯記是講一些extreme case.
為什麼異人會受苦.
為什麼異人明明在上帝面前的異人.
為什麼突然間會家破人亡.
你面對這些是想不明白.
那是怎樣的.
傳道書的智慧是講.
你將你的人生推得很盡.
喫就喫得盡.
玩就玩得盡.
然後發覺是什麼呢.
去控.
他尋求的是另一種的智慧.
是講在人生的限制裡面怎樣去.
所以你要明白有不同的.
在智慧書裡面有不同的重點.
但是其實這幾本書.
其實中間講的是敬畏上帝.
其實他幾本不同的智慧書.
都用一個perspective.
就是你從敬畏上帝的角度.

$^{161}$進入在我們人生的.
這些不同的處境當中.
來明白怎樣有智慧的選擇.
箴言.
general的.
日常生活的.
約伯紀 special的.
傳道書推到一些人生的界限.
交代了這樣.
今晚是講箴言.
箴言其實是一本結集成書的格言.
格言就是短句.
不同的一些很精練的一些說話.
寫了下來.
你要意會的.
你要反覆地思考的.
你要這樣去想.
這是一些人生的智慧.
happy husband是要怎樣的.
先是happy wife.
結了婚的你就會明白這句話.
有happy wife就有happy life.
你不明白就是不明白.
解釋給你聽也是不太領會的.
你明白的人就會說.
當然啦.
這就是智慧.
是很短的句子的說話.
而剛纔也說.
箴言其實是包括了很多.
日常生活的指引.
是說幫助一個有信仰的人.
在舊約裡面對於上帝耶和華.
神的舊約以色列民族.
他們怎樣去活得好.
他有很多不同的格言.
包括很多不同的層面.
就是家裡,友情,工作,錢,說話,憤怒.
很多這些東西.
很多這些格言裡面.

$^{201}$都有不同的比喻,對比.
智慧人,愚昧人,義人,惡人.
都要有很多這些沉澱,反思.
而《成卷箴言》有不同的部分.
一開始《箴言一章七字》剛纔我說.
敬畏耶和華是智慧的開端.
如果你回去看.
《約伯紀傳道書》其實也類似這樣.
《約伯紀傳道書》也有一句說.
敬畏耶和華是智慧的開端.
《傳道書》去到最後.
十二章 尾二那一節.
傳道者寫了這麼多東西.
總是要你敬畏上帝.
所以其實是一樣的.
在不同的角度.
但是他還是要站在那個位置去看.
敬畏耶和華是智慧的開端.
接著一開始是《所羅門的箴言集》.
之前講了很多從父親的角度.
如何對兒子有不同的教訓.
有十篇不同的教訓.
時間關係我純粹略略地講.
後來再有一些智慧人的不同的箴言.
在後期其實箴言是很闊的.
去到後期.
二十五章我當然是看了不同的經書.
其實是說希西加後期的人.
已經由所羅門時代去到後期.
希西加年代收集了很多.
關於所羅門的箴言.
這些不同的格言.
再加上雅古爾的箴言.
利姆伊勒王的箴言.
或者最後.
如果你回到教會你會聽到.
財德的婦人.
一個尼合斯去到結尾.
都是很多不同人生的智慧.
編寫成一本這樣的書.

$^{241}$為什麼我要講這麼多呢.
因為我們進入去的時候.
我想大家有一個背景.
不是我抽了一句.
有時候我們很容易.
我自己有時候也要提醒自己.
抽了一句.
哇 聖經這樣講.
我最好拿來靈脩.
默想我和他今天是怎樣的.
但如果我們多一點.
能夠從那本書去理解.
箴言其實整本書講了.
很多不同的人際關係.
人際關係這些處境.
其實是一個場景.
在舊約聖經裡面.
其實我們今天比較想.
individual.
但是在舊約聖經.
或者在舊約的設定裡面.
其實很多時候是講群體的.
聖經寫出來的時候.
在舊約的處境.
很多時候是講群體.
以色列人.
這個民族 這個百姓.
他們講的是一個群體的設定.
所以當他們講人際關係的時候.
究竟在這個群體的中間.
這些的互動.
怎樣有一個智慧.
怎樣的選擇.
能夠去祝福整個群體.
所以孤單.
傳道士有講.
單拖那些很容易出事的.
多少你都知道的.
真實的友情.
可以帶來忠告 反省.

$^{281}$幫助我們活出公義.
去到二十七章.
會再收窄一點.
會講多一點人際關係.
不過你知道.
這些短句.
有時都是穿插的.
因為它不是故事.
它也不是一些書信.
一個個段落很清楚.
二十七章其實都有相當多是講朋友.
比如說.
第五節.
今晚不用讀經.
那一句太短了.
當面的責備強於背地的愛情.
朋友家的傷痕出於忠誠.
受敵連連親嘴它是多餘.
膏油和香料使人心情愉悅.
朋友的誠實勸教都是這樣.
它有很多不同的說話.
講朋友其實是很重要 很珍貴.
如果你回去再看不同的針言.
關於朋友的裡面.
十七章十七節都有一句.
朋友乃時常親愛.
弟兄為患難而生.
這些聽起來很好聽的東西.
你有沒有朋友是兩腳插刀的.
有啊 插我嘛.
我們有時候會這樣.
但是聖經描寫的這些其實很重要.
不過聖經針形也很有智慧.
雖然說一些好的友誼.
是很值得我們去珍惜和追求.
但是不要過火.
另外十八章二十四節說.
濫交朋友自取敗壞.
但有一朋友比弟兄更親密.
十八章二十四節.

$^{321}$所以在針言裡面.
其實它講朋友 講友誼.
它有很多不同的角度.
但是我們在這裡收集下來的時候.
其實講我們今天很想看一節.
鐵磨鐵磨出人來.
朋友相感也是如此.
一節而已.
好 我請大家連續讀兩次.
一節的聖經.
好嗎 預備 一二三.
鐵磨鐵磨出人來.
朋友相感也是如此.
你試試閉上眼睛背了它.
鐵磨鐵磨出人來.
朋友相感也是如此.
哇 真是多麼聰明.
鐵磨鐵磨出人來.
朋友相感也是如此.
這一節的經文很有意思.
為什麼呢.
它原來的文字.
什麼是朋友相感呢.
是在說磨你的朋友的臉.
它原來的文字真的是這樣.
英文就說.
鐵磨鐵磨出人來.
YLT是一個原文直譯的聖經譯本.
And a man sharpens the face of his friend.
ESV是近代一點.
它說鐵磨鐵磨出人來.
And one man sharpens another.
就是你對著旁邊的熟悉一點的人.
你用手搓一下 幫你磨你的臉.
我以前看電影的時候.
你幫我搓掉他的臉.
sharpen你的臉.
當你看到鐵磨鐵磨出人來.
朋友相感也是如此.
有什麼圖畫會立刻走出來呢.

$^{361}$是這張嗎.
是否太黑了.
看到這張是什麼嗎.
打鐵的那些.
是這樣的嗎.
是啊 正啊.
Kenneth.
朋友啊.
你不就夠鐵了嗎.
我就是個錘子了.
哇 一樣一樣的衝下去.
我幫你而已.
剛剛Ivan才說完.
不是 不是 不是.
不是這樣的 那個圖畫是怎樣的呢.
我不是找到一份很好.
但大概應該是這樣.
就是那些廚師那些.
磨我的刀.
為什麼呢.
你知不知道鑄鐵是怎樣鑄的.
什麼時候打鐵的.
鑄鐵是 我也是這樣看的.
基本上是融化了鐵.
融化了鐵就去了雜質.
就倒進了容器裡.
就變成一個形狀.
打呢 錘打.
就是敲打呢.
是令到那些雜質或者形狀開始有的.
是沒有印的.
行不行 行不行行不行.
它是打扁.
它開始變成一把劍也好.
刀也好 鋸也好.
它是打到差不多.
你還要磨的.
你知不知道你有沒有磨刀啊.
有沒有人磨刀啊.
啊 Ruby 磨刀你磨這幾寸長啊.

$^{401}$不是 磨刀是怎樣呢.
磨刀其實 磨刀和磨砂紙是一樣的.
你是由粗磨到幼的.
你不是不斷起誓的.
你是磨到它的印要走出來的.
所以當時.
所羅門 你知不知道.
你有沒有印象所羅門是幾時的人.
所羅門是公元前一千年的人.
公元前一千年.
大概是.
公元前一千多年.
到七百多年.
我不知道放在哪裡找了資料.
大概其實那時候.
由後青銅時代到鐵器時代.
即是說當時的文化背景.
就開始有鐵.
以前人們打仗是用青銅.
開始有鐵.
所以對箴言的讀者來說.
你講鐵.
其實好像新約講種田那樣.
其實對他們來說是很日常的事.
鐵磨鐵 磨出印泥.
其實是講什麼呢.
講幾個很重要的東西.
其實鐵是一個很堅硬的東西.
是需要磨的.
需要加工的.
而且原來的文字.
那個sharpen的動詞.
希伯來文的原來的動詞.
不是一個被動的.
是主動去磨.
你磨刀就主動了.
有沒有玩 你擺把刀在那裡.
你不會拿著那個磨刀石.
在桌子上這樣磨刀.
你當然不拿著刀去磨.

$^{441}$你是主動地去磨.
而且那個面有些不同的聲加.
他講的時候.
那些面不是純粹講外表的磨刀.
那把刀本身要利纔行.
你沒見過那些刀厲害嗎.
我沒見過那些最厲害的刀.
我看 講起來好像很誇張.
那些很利的刀.
切蘿白 你有沒有聽過.
切完蘿白吸回去之後.
它是不掉的.
那麼誇張.
我不知道 我沒聽過.
你磨完之後那把刀很厲害.
不是說去樓下的便利店買一把刀.
一百塊錢三把那些.
切片片片的.
跟你在家裡可以.
但是他磨完之後.
那個面是很鮮活的.
什麼屠龍刀爾天劍那些.
不知道是不是這樣.
拿來切蘿白好像浪費一點.
他在講的不是隻是表面.
是你整把刀變成了高級的.
放在一個箴言的處境.
你朋友互相磨.
你整個的生命.
你整個的人格.
你整個的智慧.
你整個的選擇是不同的.
而且他那個磨.
不是一個磨完一次就完的.
原來那個文化上.
我也是看《聖經書》當年.
是講一個不是一次就完的.
是持續是慣性的.
你這樣聽完之後.
你有沒有這些朋友.

$^{481}$有沒有這些可以.
走過來可以令你更加帥的朋友.
而當我.
笑什麼.
而當我想到這裡的時候.
我想到幾樣東西.
第一是.
磨刀.
是要夠近才能磨到.
即是很離身很遠的那些.
是難一點的.
夠近.
認識你.
知道你是什麼.
磨呢 而且磨是一定會不舒服的.
不過還有一件很寶貴的事.
就是說到朋友.
其實他要冒險的.
他不是走過來就說.
喂 不行啊 Ivan 你這樣.
他真是很想為了你的好處.
提醒你的時候.
他可能.
讓你磨回頭.
我覺得《箴言》這套圖畫.
是一個很有意思的圖畫.
當我們講一個真誠的友誼.
講一個好的朋友的時候.
不是純粹講那些好來好去.
也不是講那些.
硬來的那些.
他是很在意你.
但在意你到一個地步.
他是會自己要冒險.
剛剛《箴言》有一段 我回頭看.
讓我翻幾頁.
你翻回來就可以了.
不知道我去哪了.
第十四節 十七順.
清晨起來.

$^{521}$大聲給朋友祝福的.
就算是咒詛他.
你有沒有朋友.
不生性的朋友.
清晨.
你還沒睡醒.
大聲地給你祝福.
對 祝福你.
還沒睡醒 你說什麼.
祝福你 我愛神主利多.
什麼事.
有時《箴言》有這些位置.
剛剛我們講磨.
真是你的朋友.
他會知道什麼是一個時機.
take a risk.
走過來跟你講真心話.
而在《聖經》裡面.
我想一些不同的例子.
當我想的時候.
講起來好像好聽.
我想不同的例子.
其中有兩個.
如果你對舊約有些認識.
摩西和葉剔羅去外赴.
那時候發生了什麼事.
摩西帶領以色人出了埃及.
在埃及那裡.
很多人都要找摩西解決問題.
隊伍排到對面海.
真是很長.
摩西又累 百姓又累.
接著葉剔羅就走過來.
你做什麼.
你在做什麼.
你這樣.
百姓又累 你又累.
這樣不好.
葉剔羅就跟他講.
你這樣不好.

$^{561}$你要在群體中間.
千夫長 百夫長.
即是他要選一些不同的人來幫你.
摩西.
你跟著南出街.
去到紅綠燈的時候.
他舉起像立刻變成綠燈.
你就走過來.
哇 這些神人.
跟得過吧.
一去到紅燈的時候舉起.
就綠光了.
哇 這些.
但葉剔羅不是.
他不是說摩西這麼厲害.
分開紅海.
但他見到.
你這樣做什麼.
你這樣做不好.
摩西聽了他講.
跟著對摩西 對白石.
很大的祝福.
另外.
新約聖經.
阿基拉和阿居拉.
不是 打錯字.
巴基拉和阿居拉.
我沒有 我看錯了.
老花.
巴基拉和阿居拉和阿波羅.
阿波羅在新約聖經說.
在當時已經是一個名講員.
講解聖經很厲害.
但新約聖經有一段記載.
當他講完之後.
有一對夫婦在臺下.
再跟他詳細講解聖經.
如果有一天.
傅牧師講完道.
下來的時候.

$^{601}$你突然說.
不是 阿松.
我看過聖經.
我看過聖經書.
我看過原文.
我看過很多東西.
不是 這裡可以再解釋多一點.
你再跟阿松說.
你敢不敢.
我就問一點.
但你明不明白.
你明不明白.
為什麼新約教會.
初期教會復興.
一個名講員.
在臺上宣講完聖經.
他下去的時候.
在正確的時間.
這對夫婦走過來.
跟他講解聖經.
鐵磨鐵 磨出印.
當我再想聖經不同的故事.
不同的人物的時候.
我再想.
有沒有一些既有成功.
又有失敗或者困難的例子呢.
麻煩吹吹炮棍.
大衛 掃羅.
兩個都是皇帝.
你知道舊約裡面.
有三個皇帝在聯合王國裡面.
掃羅第一個 大衛第二個.
他們都有人去勸戒他.
掃羅他開始的時候.
他像鞍春一樣.
但他是牛高馬大的.
被人選擇做皇帝之後.
就很厲害了.
打仗開頭是很厲害的.
菲律賓人全部打贏了.

$^{641}$後來開始大衛來了.
冒起了.
接著他就有些扭曲.
然後撒母耳走過去問.
在獻祭的時候.
回去看聖經.
我聽到一些牛聲羊聲.
是什麼來的.
當時神叫撒母耳和掃羅打仗.
和阿瑪利人打仗.
打完之後所有東西都要滅掉.
但是掃羅就留下一些.
當撒母耳去對決的時候.
他就說 這些是他們.
百姓要這樣.
但是大衛作為另一個皇帝.
他也有經歷過不同的情況.
除了大衛和約拿丹的故事.
大衛還有一個後來的老婆.
阿比戈.
她的故事是什麼呢.
當時阿比戈.
大衛是落難的.
如果你熟悉大衛的故事.
他落難的.
逃到那裡.
有幾百人.
好像幾百人我記得.
跟著他去找喫的.
到處逃避掃羅的追殺.
去到一個位置.
就有個喇叭.
喇叭就是阿比戈的老公.
就是那裡那些.
說得像的山寨王.
雖然大衛當時已經是.
曾經是掃羅的一個很厲害的.
打贏哥里亞的人.
所以大衛是有名聲的.
去到那裡的時候.

$^{681}$找喇叭就說.
可不可以很有禮貌地.
在那裡借一個地方住.
喇叭就說.
大衛是誰.
這期那麼多人.
背叛主人.
我哪裡都認識.
這樣說話.
跟著呢.
阿比戈一聽.
知道老公這麼說話之後.
立刻亂想.
猜那些人.
派很多的禮物去.
大衛一聽完之後.
火都來了.
就是整個故事很有趣.
你回去聽很有意思.
火都來了.
跟著就帶著那些人.
想著就去.
殺喇叭的了.
跟著阿比戈就去.
跟大衛講他說什麼.
他說.
你今天不要這麼做.
我的老公喇叭.
那些不知原文的名字.
我的老公其實是個網夫.
簡單來說就是老粗.
老粗來的.
他說你今天.
如果你不親手殺他.
將來你在神的面前.
你會覺得你自己做得對.
所以大衛那一刻.
立刻下了道氣.
阿比戈勸退他.
不是說我不想死.

$^{721}$不是的.
用上帝的角度跟他說.
你既然是一個.
在神眼中這麼寶貴的福人.
你不要在這個地方.
因為我老公的緣故.
跟著他老公嚇死了.
如果你看故事.
大衛原來本來想斬我.
嚇死了.
跟著大衛就娶了阿比戈.
大衛有這樣的人勸他.
但同時.
拿丹和大衛的故事是怎樣的.
後來大衛做了皇帝.
跟著拔士巴的故事.
如果你有印象.
就是他用計謀.
先跟拔士巴有姦情.
跟著烏靈雅和拔士巴的老公.
是一個很忠心忠心得不得了的人.
跟著用計謀殺了她老公.
跟著拿丹才用比喻葡圖園的故事.
大衛一開始不知道是真蒙還是假蒙.
為什麼會這樣 誰還要這樣.
用葡圖園的故事.
如果誰這樣.
我在神的面前就要懲罰他.
跟著拿丹就說.
那個就是你.
跟著大衛立刻認錯.
但兩個都是皇帝.
兩個都在被人批評的時候.
有不同的反應.
為什麼一個是黑色一個是紅色呢.
為何上帝最後棄了掃羅.
但又沒有棄大衛呢.
另外新約聖經.
當時耶穌去到新約的時候.
有很多不同的法利賽人.

$^{761}$耶穌也說了很多道.
很多的法利賽人都反對耶穌.
但在這麼多的法利賽人裡面.
有一個.
叫做掃羅.
他在大馬士革和耶穌相遇.
那個聲音跟他說.
你為什麼踢我.
主啊 你是誰.
我就是你所迫害的拿撒蘭耶穌.
這個法利賽人掃羅.
他在這個衝突裡面悔改.
他整個人改變.
但是這個掃羅後來叫保羅.
他仍然跟別人混.
巴拿巴是當初牧養的人.
巴拿巴當時已經是教會領袖.
但他在宣教旅程裡面.
巴拿巴說我要帶馬可.
掃羅說馬可不能帶.
兩個就吵架.
如果你一天回來YA.
看到阿蟲和亞蘭在吵架.
一個就說.
帶上Ivan吧.
衰幾次而已.
繼續不怕啦.
一個就說Ivan不能用的.
吵架 分道揚鑣.
聶麥聖經是這樣的教會領袖.
但是在這種互相的磨擦裡面.
保羅說不能用馬可的人.
去到後期.
他的保羅書信裡面寫.
他說請你將馬可帶來我身邊.
他於我全道的事上有份.
鐵磨鐵磨出印尼.
當我在想.
為什麼有些人在磨的過程.
好一點.

$^{801}$有些就沒那麼好呢.
當我在預備的時候.
在箴言的情況下.
你信了耶穌.
你都會面對著這個世界很多不同的挑戰.
你做好人的時候都會被人誣陷.
你和人講真話的時候.
被人說你為什麼誣陷我.
但是我想有一個很關鍵很重要.
很有意思的地方是.
箴言成卷書.
他在說鐵磨鐵磨出印尼的時候.
不是純粹說你去跟別人挑戰.
不是這樣的.
而是在說一個敬畏上帝的人.
他如何去做他日常生活的選擇.
而去到二十七章的時候.
在你和人與人的關係.
人際關係.
朋友的關係當中的時候.
你如何去選擇一些真朋友.
如何去選擇一些朋友.
你能夠走一步.
能夠和他鐵磨鐵磨出印尼.
是互相的琢磨.
是互相的.
可能有不同.
彼此有不同.
可能在彼此的不同裡面.
都有不同的遷就.
或者都有不同的不舒服.
不僅是你不舒服.
可能對方也有.
可能對方跟你說話的時候.
他需要給你一些勇氣去說出來.
或者反過來.
可能跟你一樣.
可能你要和對方分享一些東西的時候.
其實你也要給他一些勇氣出來.
纔可以和他分享.

$^{841}$但不是純粹為了想和他吵架的緣故.
而是一種關係.
一種信任.
也能夠說出真相的關係裡面.
你很想冒險.
走一步.
不是因為你特別厲害.
而是因為你在上帝的眼中.
你很想在這個關係當中.
做一個智慧的選擇.
因為你敬畏上帝.
去到信息的尾段.
差不多了.
知道左邊的公仔是什麼嗎.
哇 很聰明啊 誰說的.
厲害啊.
我就是在想.
圍爐取暖.
其實我們在一個群體當中.
其實我們兩邊都需要.
你既需要在一個群體當中.
你能夠可以圍爐取暖.
你有很大的安全感.
很大的安穩.
你會有舒適.
但是同時.
不要只是在一個舒適的環境當中.
如何可以在一個關係裡面.
能夠互相有一個塑造.
互相有真話.
互相在上帝的裡面.
能夠彼此的塗鑄.
真言一直在說的是以敬畏上帝.
為一個生命的基礎.
雖然我們在現代社會.
在現代社會很多時候我們出來.
特別是我們做職稱.
我們這些很有自我身份認同.
做職稱的那些.
面對社會工作.

$^{881}$在外面跟別人講真話.
講多少.
你會不會覺得怕被人指責.
不敢講那麼多.
在面對一個職場環境當中的時候.
你如何可以做一個真誠的朋友.
你敢不敢做一個真誠的朋友.
又或者我們先退很多步.
回到教會.
在教會的關係裡面.
你敢不敢冒險.
或者這麼說 為什麼我要冒險.
如果我本身已經是在圍爐取暖.
很舒服 很舒服.
每一個星期回來的時候.
Ruby就弄糖水.
好甜蜜啊.
剛剛崇拜開始之前.
我們纔在後面討論這個.
不要說了.
知道就知道了.
那個食物 什麼時候可以.
好甜蜜啊.
我們也需要.
特別是在這個世代裡面.
社會有很多的動盪.
很多的不確定性.
現在AI多厲害.
厲害到不知道.
未來的工作怎麼辦.
我也不知道怎麼適應這個環境.
我以前讀完那些東西.
都是記的 背的.
現在不背好像沒什麼用.
當你出去的時候.
大家都可能在工作職場裡面.
有很多為了生存的緣故.
他不是專門做那個壞人.
可能我們也習慣在這個模式裡面.
我們回到教會.

$^{921}$我都寧願保持距離.
離身一點.
但是箴言鼓勵我們.
不是說我們做一個漁網人.
不是說叫我們做一個.
沒有智慧的人 亂說話.
不是 你清晨的時候.
去祝福男人.
就是咒詛他.
你的智慧不是這樣.
但是箴言說.
在一些珍貴的關係.
在一些你跟他已經有一定熟悉的關係.
會不會可以考慮.
走一個步驟.
將那個關係.
寄在圍爐取暖的裡面.
同時可以.
困一困.
當我做這個默想的時候.
我想.
我自己也是.
其實我自己也是那些很喜歡和諧的人.
我也很喜歡.
但當我想的時候.
你走一個步驟.
有時是怎樣呢.
其實對我們這些內向的人.
什麼.
當我們這些內向的人.
很怕在別人面前揭露太多.
安全一點.
有什麼祈禱.
我這段時間都挺好的.
感恩的.
你有沒有聽過.
我預備的時候.
我這段時間其實累到瘋了.
剛剛早幾個星期前.
我有一個家人過世了.

$^{961}$這幾個星期.
很多不容易.
然後我今早纔去完國安釋例.
累到瘋了.
我記得阿南提我們港島那些人.
每個星期的小功課.
我今個星期的小功課是什麼呢.
(出一個powerpoint).
今個星期的小功課就是.
當我們鐵磨鐵.
我想說.
當你聽到其他人的意見或者批評.
或者你有些善意的提醒別人的時候.
對方反彈.
你有些心裡面的激氣.
鐵磨鐵.
困一困的時候.
今個星期的小功課如果你是這樣困的話.
你就祈禱.
說 上帝我應該怎麼回應.
在困的過程中.
你讓自己有五秒至十秒停一停.
你去祈禱.
但是.
當我一路走回教會的時候.
我剛纔為什麼會說我自己的東西.
就是.
如果你很少和別人這樣困.
或者已經很有智慧.
很成熟處理這些東西.
沒什麼這些東西.
又或者像我這樣比較內斂的人.
你今個星期的小功課是另一樣的.
就是.
你找一些熟悉的.
或者不太熟悉也不要緊.
你學.
你在一個群體當中.
分享你的近況.
有時我們太過不習慣分享我們的近況.

$^{1001}$又或者我們覺得我們的近況沒什麼.
值得分享.
但當你學.
你去分享的時候.
你習慣在這群朋友裡面分享.
可能是那些雞毛蒜皮的東西.
你講著講著就可以講其他東西.
當你那群朋友聽到.
你又是上班下班.
你今天喫了什麼.
講這些.
講著講著就覺得今天老闆很壞.
講著講著就覺得老闆又來了.
真的發瘋了.
但當你聽得多的時候.
你除了安慰他.
除了圍爐取暖.
不如我們一起祈禱.
可能祈完禱三個月都沒變過.
繼續守望他.
當我去想這些功課的時候.
我覺得真言是一些很實際的東西.
真言是講我們日常生活.
每天面對的一些情況.
需要怎樣做選擇.
很實際.
你繼續我要祈禱.
所以我想請敬拜對象.
當我們講著與同輩同行.
PS Walk.
如果你今天已經有一些.
可以跟你鐵磨鐵的朋友的話.
我為你感恩.
但如果你沒找到這些朋友的時候.
我預備這篇信息的時候.
我一直.
阿南其中講到頭幾句的分享.
我在心裡.
當然你祈禱.
你很想求上帝讓你有這個朋友.

$^{1041}$但整個系列.
一開始的時候阿南說.
為你自己祈禱.
怎樣讓你自己的生命.
可以更加準備成為別人眼中這種朋友.
我們一起祈禱.
真愛耶穌多謝禰.
因為禰很在乎我們的生命.
主一方面.
禰讓你無盡的恩典托住我們.
禰告訴我們知道.
無論我們今天做成怎樣.
我們已經是禰的兒女.
禰的寶血 禰的寵價.
完全的接受我們.
但是禰不滿足在這裡.
主要求的是.
禰讓我們這裡每一個弟兄姊妹.
在我們不同的生命的光景.
不同的人生的階段.
不同的情緒狀況.
禰讓我們的心繼續向禰打開.
渴慕禰 敬畏禰.
教我們在我們日常的生活當中.
可能是一些平淡的事.
在工作 在家庭.
在和你的關係 回教會等等.
逛街 喫飯 看電影.
讓我們有這個準備.
我們準備在不同的境況裡面.
遇到一些困的人.
我們懂得怎樣回應.
讓我們的生命.
經得起這些磨練.
在你的恩典 在你的智慧當中.
我們做討你喜悅的選擇.
因為你的選擇.
在你裡面是一個真理的時候.
讓我們經驗真理那份美好.
那份甘甜.

$^{1081}$祝福我們每一個弟兄姊妹.
祝福我們因為禰的緣故.
我們經歷得到信仰的真實.
我們經歷得到阿爸父對我們的不離不棄.
我們經歷得到這個世界.
給不了我們真實的平安.
求禰親自的祝福.
閱立我們在禰面前的禱告.
敬拜 仰望.
祈禱奉耶穌基督的聖名.
Amen.
\newpage



\section{申命記 6:6-9}
\label{sec:031897}
\textbf{跨代傳承}~~~~何羅乃萱師母
\newline
\newline
連結: \href{http://www.church.com.hk/acms/content.asp?site=cdc\&op=show\&type=product\&code=031897\&layout=sermon}{\texttt{031897}} ~~~~ 語音日期: 2026-03-29
\newline
\newline
\hyperref[sec:031896]{\small{< < < PREV SERMON < < <}}
~
\hyperref[sec:index_chronic]{\small{[返順時目]}}
~
\hyperref[sec:031898]{\small{> > > NEXT SERMON > > >}}
\newline
\newline
申命記 6:6-9
\newline
\begin{longtable}{cl}
\hline
\hline
章節 & 經文 (和合本修訂版)\\
\hline
6:6 & \begin{tabularx}{0.7\textwidth}{X} 我今日吩咐你的這些話都要記在心上， \end{tabularx} \\ \\ \relax
6:7 & \begin{tabularx}{0.7\textwidth}{X} 也要殷勤教導你的兒女。無論你坐在家裡，走在路上，躺下，起來，都要吟誦。 \end{tabularx} \\ \\ \relax
6:8 & \begin{tabularx}{0.7\textwidth}{X} 要繫在手上作記號，戴在額上作經匣； \end{tabularx} \\ \\ \relax
6:9 & \begin{tabularx}{0.7\textwidth}{X} 又要寫在你房屋的門框上和你的城門上。 \end{tabularx} \\ \\
[1ex]
\hline
\hline
\end{longtable}
$^{1}$回到同福就好似回到自己家.
好親切.
既然這麼親切,我又問問你.
我講開道通常有幾點?.
幾點?.
三,就是三,讓你看看.
轉化.
更新.
傳承.
我想過用一節經文兩節經文去講.
但我發覺這麼闊的題目.
很難這樣講.
於是就分開三個.
我覺得在信仰裡面很重要的主題.
跟大家分享.
第一.
我們先講轉化.
就是怎樣信耶穌.
羅馬書十章九節.
我們一起讀一讀.
一二三.
你由口裡認耶穌為主.
心裡信神.
叫他從死裡復活.
先不要讀,讀第一節.
這一節很多信耶穌的人跟你讀.
你由口裡認耶穌為主.
心裡有信.
就叫信了耶穌.
是不是?.
原來還有下一節.
就是跟著那個.
我們一起讀.
天國近了.
你們應當悔改.
信了耶穌.
最重要的是生命的扭轉和改變.
接著下一節是我抄下來的.
這個普及譯本譯得很好.
你們要用行動去證明.

$^{41}$你真的信耶穌.
最終悔改.
而且不要因為.
覺得很安全.
因為我們是亞伯拉罕的子孫.
如果用現在的說話.
我們是某某牧師的女兒.
我們是某某執事的女兒.
我們等等等等.
原來信仰沒有世襲.
信仰不能自動.
我們不如再看什麼叫做悔改.
好嗎?.
悔改.
其中一個意思就是.
如果我們真的有一棵好的樹.
就會結好果子.
壞樹.
就會結壞果子.
悔改有兩個意思.
拉丁文是repent.
意思就是對自己的罪.
感覺悲傷.
感覺懊悔.
有沒有呢?.
我記得我信耶穌的時候.
我和爸爸發生很大的衝突.
我接受不了我有情緒病嚴重的爸爸.
於是我憎他.
我恨他.
愛不到爸爸.
我形容那個階段叫愛無能.
直到遇見耶穌.
告訴我生命可以扭轉.
可以改變.
所以當日的執事問我.
你有沒有罪?.
我說我有.
他說你的罪是什麼?.
我說我恨我爸爸.

$^{81}$就是這個意思.
第二就是Metanoia.
差不多.
你不要理會.
總之在希臘文.
悔改是什麼?.
to turn.
意思是什麼?.
你本來是這樣走的.
但當你信了耶穌之後.
反過來.
方向不同.
我是十八歲信耶穌的.
我本來走的就是.
我要進那間學校.
我要談戀愛.
我要嘟嘟嘟嘟.
但當我遇到耶穌之後.
我方向改了.
主啊.
你是我生命的主.
我還記得.
大學第一年.
就是我信耶穌的那一年.
我將我的宿舍的房間.
每天早上七點鐘打開門.
讓宿舍裡的.
教徒的姐妹.
那裡是女生宿舍.
來我的房間.
為那些未信的祈禱.
逐一祈禱點名.
求主讓他們聽到福音.
當中有一個.
很難搞的.
全級拿第一.
拿獎學金.
我跟她說.
哎呀.
某某某.

$^{121}$你跟我回教會吧.
這些東西留給你吧.
不用客氣的.
Shirley.
每天我見到她.
她都跟我說不用客氣.
我都完全.
不知如何埋首.
但上帝沒有放棄她.
大概二十年後.
我和何牧師.
去多倫多最大的那間教會.
我們要主持一些婚姻講座.
在教會的辦公室.
突然間見到她.
你為什麼在這裡.
她說我在這裡做Amen Assistant.
大家明白嗎.
她是博士.
是全校最厲害的學生.
她做過很多事情.
但當她遇上耶穌之後.
她在教會服事.
上帝的恩典實在奇妙.
對我來說.
我們再下去看看.
除了在學校開放自己的房間.
為不同的人祈禱.
我還將我的故事.
在當時的團契跟別人分享.
當時團契有一個弟兄.
叫Peter Ho.
講了故事給Peter Ho聽.
Peter Ho有一天告訴我.
原來他是導演.
知不知道.
你們不知道吧.
他懂得寫劇本.
他將我的故事.
寫成了一個福音劇.

$^{161}$然後特邀我做女主角.
大家看到上面.
就是當日的劇照.
是超重的我.
和只有一百十磅的.
你們認識的何牧師.
Peter Ho.
認不認得.
在福音劇之後.
我們就開始拍拖了.
I have decided to follow Jesus.
有誰會唱這首歌.
這首歌是我信耶穌之後.
那一年不停在心裡唱的.
我提醒自己.
I have decided to follow Jesus.
唱英文好不好.
唱一次 第一節.
誰會唱 跟我一起唱.
就是說你跟我差不多大.
不要緊 我認了他.
1 2 3.
I have decided to follow Jesus.
No turning back.
好 唱第三節中文.
世界在背後.
十家在前頭.
世界在背後.
十家在前頭.
世界在背後.
十家在前頭.
永不回頭.
當時我參加了很多動靈會.
600人的動靈會.
一呼叫.
500多人走了出來.
我們一起含淚.
唱這首歌.
十家在前頭.
永不回頭.

$^{201}$生命的轉化.
再下去.
生命的更新.
為什麼要講這個呢.
因為信了耶穌之後.
就不是一帆風順.
有高山 有低谷.
有狂風暴雨.
有沒有.
生 老 病 死.
我們經歷很多事.
但我覺得這兩節經文.
是時刻在我心裡.
給我很多提醒.
我們一起讀一讀.
一二三.
所以弟兄們.
神的慈悲.
勸你們.
將身體獻上.
當做活祭.
是聖潔的.
是神所喜悅的.
你們如此事奉.
乃是理所當然的.
不要效法這世代.
要心意更新而變化.
讓自己可以辨別.
上帝善良.
純存.
可喜悅的旨意.
信了耶穌.
是不是可以馬上變呢.
都可能是.
有些人告訴我.
以前講粗口.
一變.
馬上可以不講了.
但另外這個呢.
其實我講道講過的.

$^{241}$這兩節.
有沒有印象.
不要緊.
沒有的.
我都明白的.
其實就說.
有個metaphor是茶包.
放進杯茶裡.
慢慢慢慢滲透.
我們生命的改變就是這樣.
我們拿去看看.
怎麼解釋.
所以.
先不要留意照片.
遲些再說.
羅馬書十一章.
為基礎.
用上帝的同情憐憫.
我們將身體獻上.
當做活祭.
一直我都覺得.
就是將自己整個主權.
交給上帝.
但是上個月.
我就看利未記.
我不知道你怕不怕.
零修看利未記.
那些羊啊.
又怎樣怎樣.
但突然之間.
看到一句重覆的.
就是無瑕疵的生促.
我看我自己.
這裡就痛了.
這裡有些傷痕.
內臟器官就不要說了.
發覺.
我很多瑕疵.
在那天.
不知道為什麼.

$^{281}$好有感動.
一直流淚感恩.
上帝你真好.
我這麼多瑕疵.
你都接納.
以前一隻羊都要無瑕疵.
體察什麼叫上帝喜悅.
悅立.
是理所當然.
是詳細思考的.
出自心靈.
理智.
Your heart and your mind.
再下去.
你可能要知道.
那張照片.
我解釋一下.
那張照片.
是我和何牧師.
信了耶穌之後.
奉獻了做傳道去台灣.
二十多歲的時候的合照.
都像了何.
但從來沒想過.
去到是艱難的.
去到是沒錢買牛肉吃.
所以我極度貧血.
去到是更加沒錢買冷氣機.
所以我生業廢.
那四年艱苦的日子.
但當中.
我很記得.
OMF內地會的宣教士.
每一次他寫信也好.
打電話給我們.
問候我們也好.
在我們最艱難的時間.
我當然是.
Bang Bang Bang Bang.
告訴他.

$^{321}$Reverend Rowe.
你猜他怎樣.
聽到這幾個字嗎.
God is good all the time.
Shirley and Peter.
You have to remember.
God is good all the time.
他沒有說話.
他只是用這句說話.
很溫柔地回應我們.
讓我們慢慢想下去.
叫我們不要效法這個世界.
這個時代的思想.
所想的東西.
我們要去辨別.
要想清楚.
看清楚.
要去選擇.
人家做我們就做.
然後去察覺.
上帝完全成熟的旨意.
最後那句.
察驗什麼.
我們讀一讀.
這句.
1 2 3.
察驗什麼是上帝的旨意.
什麼是那美善的.
可喜歡的.
純全的事.
因為我們每天.
你看手機.
很多資訊.
我們直接就接收.
但是要察驗.
什麼是上帝不喜悅的.
我們看一看.
有幾樣.
我們分享一下.
追隨世俗的價值觀.

$^{361}$你說沒有.
我是做家長工作的.
現在世俗的價值觀.
有幾個字.
要贏在什麼.
氣泡線.
很多人說不好.
但還是要贏.
另外.
享樂自我中心與沉溺.
你說我沒有.
我又沒有沉溺什麼.
我又不會.
要喝很多酒.
但是.
手機呢.
會不會沉溺手機呢.
很多人正在沉溺手機.
機不離手.
一定要忙.
你說不用點.
一會去搭地鐵.
你看看車廂.
有多少人不看手機.
通常都起碼80歲.
肉體的情慾與驕傲.
是說什麼呢.
很多人都說一句.
叫做.
合久別婚.
別人離婚就跟我說.
結了婚久合久別婚.
所以我寫了一本書.
叫做分久必合.
這是我的信念.
結了婚越久就越有默契.
為什麼要說.
合久別婚.
最後與上帝的生活方式.
或者上帝的看法.

$^{401}$背道而馳.
在這裡有誰是50到80歲的.
請舉手.
我舉.
看起來有點不像.
都舉.
不要緊.
我們50到80歲.
50歲的報紙已經叫老翁老婦.
好聽一點叫銀鈴.
金鈴.
還有一個名字叫奇鷹.
有人叫我選奇鷹蝴蝶.
我說不用了.
不用客氣了.
我最近很喜歡.
這個稱呼叫什麼呢.
三個字的有沒有聽過.
壯世代.
Strong Generation.
行得走得.
還有很多東西可以傳承下去.
我們不是在等什麼.
大家知道.
我們不只是坐郵輪.
照顧孫子.
我們還可以為主.
不要奔跑.
也可以的.
我曾經.
被一個腦科醫生.
好像跟大家說過.
看了兩次我的腦素描.
就跟我說.
我覺得你有腦退化的跡象.
我說為什麼.
我的腦退化的跡象.
是比平常人少.
我聽完當然很擔心.
回去問何牧.

$^{441}$那怎麼辦.
他說你現在行得走得.
想得那麼快寫得那麼快.
你的腦退化的跡象比別人少.
如果你正常就很可怕.
前幾天.
我真的證明.
我的腦退化的跡象少沒什麼事.
我看著一個人.
他就是五十年前.
不停叫我回教會的阿姨.
我看了一個小時.
結果我有沒有去問他.
我真的有.
真是他.
但我就說阿姨謝謝你.
你在五十年前.
你不停叫我回教會.
我現在這樣.
是嗎?我忘記了.
但我就發明了一句說話.
White Matter少一點.
Doesn't Matter.
我們很容易被這些.
一路讓我們陷入.
不知道怎麼辦.
檢查完這裡又檢查這裡.
這裡又做那裡又做.
美還美了.
你說那怎麼辦.
我們每天都收到一些壞消息.
我很喜歡旁邊那段.
你上網找找.
有一位牧者就說.
冥想行不行.
他說冥想不是不行.
他自己帶一些東西進來.
撒旦的思想就可以進來.
所以當我們要.
倒空或清空的時候.

$^{481}$最緊要.
專注上帝的話.
對不對.
我們一起讀這幾節.
給你幾節.
123.
這話將我救活了.
我在患難中因此得安慰.
我若不信在活人之地.
得見耶和華的恩惠.
就早已喪膽了.
我心裡多憂多疑.
你安慰我.
就使我歡樂.
我若不是喜愛你的律法.
早就在苦難中滅絕了.
我們不是沒有困難的.
我們不是沒有苦難的.
但是在苦難當中.
上帝的說話會救活我們.
有沒有人經歷過.
舉手.
你看看你的手.
是見證.
當我們覺得別人的說話.
都聽不進去的時候.
原來上帝的說話.
是很有力量的.
我試過身體很差.
差到無法走路.
差到不想回到同福.
因為一定會被人問.
師母你的身體好了嗎.
我不知道怎樣回答.
三個月.
我覺得自己困在一個圈子裡.
直到有一天.
我記得是星期天.
我問上帝.
讓我知道我可以怎樣.

$^{521}$那天上網就聽了一篇英文的道.
那篇道就是一節經文.
希伯來書二章十四節.
它說你們不要再做死的奴僕.
我突然醒了.
原來.
我是死亡的奴僕.
我因為怕病.
其實到最後.
就怕死.
於是就把自己困在那裡.
我記得當天我跪下.
一路哭一路哭.
耶穌我不要再做死的奴僕.
耶穌我要起來.
我起來了.
今天的我.
是可以這樣.
五十下.
有些人知道我打保齡.
打乒乓.
出比賽.
年齡.
放在一邊.
靠著耶穌的能力.
我們可以重新出發.
而對我來說.
更深的一個學習.
就是Let Go.
Let God.
我們有下一代.
我們很想為下一代鋪排很多東西.
我還記得我女兒.
一歲半的時候.
我看見我媽媽很喜歡抱著女兒.
走到家裡的舊樓梯.
她還要一路抱一路唱.
笑什麼笑.
她穿高跟鞋.
這樣抱著寶寶.

$^{561}$你說怕不怕還要歲幾.
我跟她說.
媽媽我給你女兒照顧.
但你要答應我.
不要抱女兒下樓梯.
沒問題.
我在辦公室收到電話.
飛鴻打來.
你媽媽下樓梯.
抱著寶寶.
不僅怕.
我很生氣.
回到家裡.
看見媽媽指著她.
媽媽你又說這樣.
現在我女兒這樣.
你算怎樣.
媽媽托著眼鏡.
女兒我眼鏡這裡裂了.
手這裡瘀了.
你女兒一直在我的懷抱裡.
沒有受傷.
我記得當時我的手這樣.
對不起媽媽.
那晚.
剛巧.
我的靈修製體.
亞伯拉罕獻而殺.
我突然很明白.
我也在那晚.
我跪下.
耶穌.
讓我女兒.
很安全.
但我知道你可以.
Let Go Let Man.
God.
Let Go Let Man.
God.
第三點.

$^{601}$傳承信念的傳遞.
我們用啟卿的方式.
讀一讀.
你和你的子孫要一生敬畏上主.
你的上帝.
執行祂所有的法令和誡命.
就可以得享長壽.
你要盡心盡性.
盡力愛上主.
你的上帝.
最後一個字.
我今天頒給你的誡命.
你要全心全意進行.
並要不斷.
教導你的孩子.
無論在家裡.
還是在路上.
無論躺下還是起來.
你都要談論.
這裡說到.
我們很想很想將信仰.
傳承給下一代.
我現在還要傳承給.
下一代.
我是用沉浸式的.
現在很流行這個詞語.
我見到每一個孫.
我就和他唱.
我要上交山主.
大孫會唱.
第二孫也會唱.
第參孫只有兩歲多.
前幾天.
後母告訴我她也會唱.
兩歲多.
然後叫我出來.
然後播.
我要上交山主.
大孫就唱.
我的媽呀.

$^{641}$幫主.
他會懂得唱.
他不懂得說.
我已經收貨了.
最開心的是他和我一起.
Hallelujah.
拿著咪高峰.
Hallelujah.
深感安慰.
我很記得我女兒.
每一次離開我的背影.
我就和她這樣祈禱.
求主讓她.
行在正路上一生也不偏離.
我們很想信仰一代.
傳一代.
傳下去.
再下去.
另一個信念傳遞的使命.
我們.
又可以用啟認的方式.
讀一讀.
想到你心裡無偽之信.
這信是先在你外祖母羅爾和你母親有離基心裡的.
我深信也在你心裡.
因為神賜給我們.
不是膽怯的心.
乃是剛強.
仁愛緊守的心.
我們一起讀.
最後一節.
從前所交託你的善道.
你要靠著那住在我們裡面的聖靈.
牢牢的守著.
這裡是說.
提摩太的外祖母羅爾.
交給離基.
然後一路傳遞下去.
信仰的傳遞.
除了在自己家外.

$^{681}$在外面也可以.
有沒有找一些年輕人.
跟他們聊聊天.
在過去一年.
我都問我身邊的人.
我有一個未了的心願.
就是我最能夠做的.
就是寫作.
於是我現在是在教.
兩班小朋友寫文.
那種喜樂.
你看到他又不喜歡.
他真的寫.
來到羅老師這班.
我覺得無論他講得那麼生動.
我仍是沒有興趣.
但是.
他後來寫著寫著.
他就發覺到.
他開始對寫作有興趣.
越寫越好.
其中有一個男生.
他本來一行都寫不到.
但有一天他突然寫詩.
你有沒有想過他變了詩人.
我走過去.
同學你是詩人嗎.
他就看著我.
那天寫了四張紙.
是不是很開心.
傳遞的快樂.
就是如此.
這般.
當然還有以下.
最近發生的故事.
想和大家分享.
以下大家看到的這張.
就是我和何牧師.
去了烈火特會.
何特去台北.

$^{721}$張茂松牧師.
給我們有機會參加.
何牧師也在當中.
講了一場.
在會裡.
我聽了一篇很好的道.
是神之心教會的.
牧者講的.
他講.
代宗的戰.
講下一代.
上帝叫我們用刀槍為他們爭戰.
但我們下一代原來是戰.
兒女是耶和華.
所賜的產業.
所懷的胎.
是他所給的賞賜.
所生的兒女好像.
勇士手中的箭.
箭有什麼特色呢.
這裡有寫.
去得更遠比刀劍遠.
對不對.
你喜歡去哪裡都可以.
但為人父母要做那把功.
功就是要將箭.
拉到盡.
有時候大家都覺得痛.
大家都覺得艱難.
但你還是要拉.
或者有時候要放手.
把箭飛過去.
當日他的呼召.
就是鼓勵父母.
成為孩子的菊.
中間那張.
不知道你知不知道.
講話的那個女孩.
是誰.
我的女兒.

$^{761}$上個星期.
有機會跟她同台講話.
可能你覺得.
不是常常都是.
很少機會.
學校請我們講.
和諧家庭.
那個題目已經不容易.
又問我們有沒有吵架.
又問女兒.
成績不好媽媽會不會罵她.
她說不會的.
我們家不會發生這樣的事.
當我在旁邊一路看著她的時候.
其實我很欣慰.
我覺得有些事.
傳承了下去.
雖然我從來沒有想過.
於是.
前幾天.
我打電話徵詢女兒.
今天媽媽說會放你的照片.
她說你放吧.
前幾天.
我去一個.
帶嬰兒.
去嬰兒健身.
有個婆婆走來問我.
你是否何志迪牧師的女兒.
我問她.
她說是的.
我問婆婆.
為何你知道我是何志迪牧師的女兒.
她說你和她的樣子一模一樣.
我深深知道.
我女兒也很驕傲有這樣的爸爸媽媽.
所以今天放了她的照片.
接著的分幾小時.
由於我是什麼世代.
狀世代.

$^{801}$所以我還可以學東西.
我學了剪片配樂.
所以我將.
在烈火特會的領袖.
變成一個.
分幾小時的短片.
當中包括了.
張茂松牧師50週年.
他們拍的片.
旁述是張牧師.
當中包括張牧師上台.
那八千個會眾.
對他的熱烈歡迎.
當中包括.
我覺得他講得最好的.
那兩句說話我錄下了.
也包括他兒子張光偉牧師.
呼籲.
不只是為自己教會增長.
是為全球華人教會復興.
還有.
他兒子為他兒子的祈禱.
我們播放片段.
讓大家看看.
我們開始的時候.
很簡單.
甚至有點簡陋.
但是.
在教會裡頭我們.
彼此相愛彼此扶持.
成為了一個肢體.
我們年輕的歲月.
栽種在這裡.
一晃就是50年.
許多人.
從年輕.
壯年.
到老年.
再至成家.
立業升職.

$^{841}$一路上經歷上帝奇妙的.
帶領.
藉著神的恩典.
我們的生命一直被拓寬.
從原本在.
巷子裡的小教會.
搬到大馬路上.
為了接觸更多.
需要耶穌的人.
有眼淚有歡笑.
這裡已經是8000人的會場.
張牧師.
待會會扶著他上台.
什麼時候.
我們大了耶穌就小了.
什麼時候.
我們小了.
耶穌就大了.
明白.
通通都是遇見耶穌的故事.
對不對.
通通都是遇見耶穌的故事.
當他們一遇見耶穌.
他們的生命就非常奇妙的改變.
就是因為他們遇見耶穌.
於是不但他們的生命改變.
這樣就改變了後面.
2000多年.
全世界不知道多少人的生命.
當他們遇見神的時候.
就有神蹟的故事.
有生命的話.
才能夠帶領華人回向耶穌.
我們願意.
我們繼續.
讓.
所有的華人.
要覺醒過來.
這裡是呼叫年輕人.
為主奔跑.

$^{881}$不同年紀的.
當中應該有3000個年輕人.
一直在這裡.
我都幾乎想下去跑.
這是洪牧師分享的那一場.
他讓我有機會.
我們分享.
同樣要Honor.
撞世代.
然後到最後.
張某松牧師上來.
和我們擁抱.
一個很深的擁抱.
一個很深情的擁抱.
一個接納的擁抱.
一個欣賞的擁抱.
接著.
由於我坐的位置很好.
我看到張光偉牧師.
為他的兒子.
按手禱告.
當日他的兒子.
是擦眼淚.
但我拍不到那一幕.
一個爸爸這麼真誠.
為孩子禱告.
孩子是知道的.
孩子是會感動的.
我馬上想起.
我身邊很多的朋友.
下一代可能沒有回教會.
但我一定.
很想鼓勵他們.
要為自己的下一代禱告.
讓上帝.
帶他們回來.
因為我們深信.
到如今耶和華.
都幫助我們.
所以就讓我們.

$^{921}$切切的.
為我們下一代祈禱.
好不好.
不要忘記他們.
他們走遠了.
但他們會回來.
我相信上帝不會放棄他們.
我們起來吧.
在這個時候.
很想為人父母的.
很想為你孩子.
祈禱的請你舉起你的雙手.
請你舉起你的雙手.
你可以心裡面.
花一點點時間.
你向上帝呼求.
你可以說主人你叫他們回來.
主人你為他們預備未來的配偶.
給大家一點時間.
安靜祈禱.
然後請何牧師上來.
為大家祝福禱告.
當我們今天能夠在這個殿裡.
看到這麼多弟兄姊妹.
願意舉手.
跟神說.
我們願意為我們的下一代.
或許.
有些還有上一代的弟兄姊妹.
來為上一代禱告.
我們深信.
每一間教會.
都是跨代的.
我們有上一代,我們有這一代.
也有我們的下一代.
而我們很深信.
我們的神給我們的應許.
就是當我們願意愛神的時候.
神就會祝福我們一千代.
最遠我們看到.

$^{961}$祂自己的應許.
一定會應驗的.
求主特別帶著我們.
今天看到這麼多的家.
爸爸媽媽.
願意承諾.
為我們的下一代祈禱的時候.
我們求主繼續.
真的按你的應許.
你要祝福我們一千代.
求你與我們同在.
幫助我們,幫助我們的教會.
幫助我們每一個家庭.
願意一生的.
願意跟隨耶穌,走在你的心意當中.
我們也要深信.
這首詩歌.
所提到的,到如今.
耶和華仍然是.
幫助我們,願一切的榮耀.
頌讚,傳歸給.
我們做天主的爸爸,聽我們的祈禱.
奉我救主耶穌基督.
這是一個名校,Amen..
\newpage



\section{雅各書 5:7-12}
\label{sec:031898}
\textbf{百忍成聖}~~~~秦潤祥牧師
\newline
\newline
連結: \href{http://www.church.com.hk/acms/content.asp?site=cdc\&op=show\&type=product\&code=031898\&layout=sermon}{\texttt{031898}} ~~~~ 語音日期: 2026-03-29
\newline
\newline
\hyperref[sec:031897]{\small{< < < PREV SERMON < < <}}
~
\hyperref[sec:index_chronic]{\small{[返順時目]}}
~
\hyperref[sec:031899]{\small{> > > NEXT SERMON > > >}}
\newline
\newline
雅各書 5:7-12
\newline
\begin{longtable}{cl}
\hline
\hline
章節 & 經文 (和合本修訂版)\\
\hline
5:7 & \begin{tabularx}{0.7\textwidth}{X} 所以弟兄們，你們要忍耐，直到主來。看哪，農夫等候著地裡寶貴的出產，耐心地等到它得了秋霖春雨。 \end{tabularx} \\ \\ \relax
5:8 & \begin{tabularx}{0.7\textwidth}{X} 你們也要忍耐，堅固你們的心，因為主來的日子近了。 \end{tabularx} \\ \\ \relax
5:9 & \begin{tabularx}{0.7\textwidth}{X} 弟兄們，你們不要彼此埋怨，免得受審判。看哪，審判的主站在門口了。 \end{tabularx} \\ \\ \relax
5:10 & \begin{tabularx}{0.7\textwidth}{X} 弟兄們，你們要把那先前奉主名說話的眾先知作能受苦、能忍耐的榜樣。 \end{tabularx} \\ \\ \relax
5:11 & \begin{tabularx}{0.7\textwidth}{X} 看哪，那些忍耐的人，我們稱他們是有福的。你們聽見過約伯的忍耐，也看見主給他的結局，知道主是充滿憐憫和慈悲的。 \end{tabularx} \\ \\ \relax
5:12 & \begin{tabularx}{0.7\textwidth}{X} 我的弟兄們，最要緊的是不可起誓；不可指著天起誓，也不可指著地起誓，任何誓都不可起。你們說話，是，就說是；不是，就說不是，免得你們落在審判之下。 \end{tabularx} \\ \\
[1ex]
\hline
\hline
\end{longtable}
$^{1}$好 今天我們會讀《雅各書》.
有沒有人看過《雅各書》.
應該有吧 看了這麼久.
《雅各書》今天我們會講到第五章.
我們一起讀第五章的七字八字.
一二三.
弟兄們啊 你們要忍耐直到主來.
看啊 農夫忍耐等候地裡寶貴的出產.
直到等得了秋雨春雨.
你們也當忍耐堅固你們的心.
因為主來的日子近了.
基督徒說什麼 你們要忍耐.
當我看這段經文的時候.
他講忍耐 我就想肯定是在講背景.
因為《雅各書》背景是門徒四散被逼迫.
很慘的 或者他的身家少了一大截.
或者他被人欺負 被人欺負.
但是我看文法的時候.
他好像不只是講那樣東西.
他不是說你忍耐一下 過了就OK了.
這裡你不要看中文 看他要忍耐這個字.
Marco Fimel 他是用Constative Error Imperative.
很簡單的說 他不是說你忍耐一下.
是不是 好像聽到很悶 是不是.
你忍耐一下 很快就過了.
他不是說忍耐一下 不是傻的.
他一句到尾 這個詞態.
大概我自己meditate回來之後的結論是.
你要有一個commitment.
因為你這一生都是要忍耐的.
剛才他說什麼 忍耐直到主來 是不是.
弟兄姊妹 今天第一件事.
我想大家記得的就是我們這一生人就是要忍耐.
所以我設計了今天我們的主題是一齊讀.
白忍 誠誠.
如果我們可以在我們一生人的裡面.
commit去忍耐 為了主的緣故.
在忍耐的當中 我們可以在神那裡支取力量.
我們靠著祂寶血 得勝的時候.
我們都可以在這個誠誠的道路上有份.

$^{41}$我們都得勝有餘 Amen.
我不知道你現在是不是在忍耐.
你有沒有在忍耐 有沒有什麼難關.
有沒有人最近入過油 是不是.
七國這麼亂 油費很貴.
我昨天入油 肉赤到我七十多元.
七十多元一缸油 OK.
我去吃飯 我計算我現在用了多少錢.
為什麼 因為七國這麼亂.
我去光顧那個打邊爐.
我說最近生意怎麼樣 不就挺好.
是不是 為什麼 信不信.
我馬上信了 因為那個打邊爐那裡.
平時四塊肉 它少了塊肉.
真的生意不好 是不是.
很多時候我們都是少了塊肉.
為什麼 不是你們的錯.
是因為這個世界很亂.
我看這個世界很多新聞.
它說什麼打仗的原因.
我只見到一樣東西.
受苦的都是平民百姓.
其實我們都是平民百姓.
我們這一世很多時候.
我們都需要有很多的忍耐.
我講這篇道 我就在想用什麼例子好呢.
當然網上有很多例子是要我們學習忍耐.
我在想的時候 我表姐就打電話給我.
三四天之前打電話給我.
我相信有人都聽過我表姐的故事.
是不是 有些人未聽過我講多次.
三四年前 我表姐有BB 很健康的.
她生了出來之後 還是很健康的.
後來臉上有些青 是不是.
長話短說 魯莽艷 然後就殘廢了.
怎麼忍呢 是不是.
我表姐在那個時候就好像晴天霹靂一樣.
我自己有個女兒.
如果我見到這件事發生在我身上.
她不是一早已經知道有事.

$^{81}$她本來沒事 因為疏忽而導致.
她魯莽艷毀壞了腦 是完全癱瘓.
很感恩 在那個時候.
我跟她談電話 跟她祈禱 她相信了.
今天她都跟我說 你當時講的那些經文 令我很大改變.
很老實跟你說 我忘記了講什麼經文.
但我只知道那是經文.
陳健華語很有力量 讓我們可以前進.
我表姐很值得感恩.
因為我跟她談話的時候 到今天.
她打電話給我的時候 她有件事改變了.
那時候跟她談話的時候.
她都說不要緊啦 我OK啦 是不是.
我頂到的 我相信了.
我發覺我上一次跟她談話 即是四天前的談話.
她每兩句就一句講耶穌.
她很感恩.
她這個忍耐功課做得很好.
為什麼我講這個例子呢 因為她跟我說了一句話.
她說我去看醫生的時候.
那個護士跟我說 那個姑娘跟我說.
她說你遇到這一輩子都是這樣 你遇到這一輩子都是這麼辛苦.
她說我無所謂啦 我很愛我女兒.
因為呢 雖然她最近才知道 原來她女兒不是很懂得痛.
所以她裡面有發炎 她自己都不知道.
但是她跟她女兒的交流不是用言語.
她看不到東西 她吃不到東西 她走不到路.
但她懂得笑和懂得憤怒樣子 是不是.
她說她很有趣的 她出去見三寶的人她不憤怒.
不憤怒是回家才憤怒的.
我說你知不知道這個正正就是愛.
她說是有享受到和她的愛.
她說我很愛我女兒.
她說我都知道當耶穌再回來的時候.
主再回來的時候 將會得到一個新的身體.
我和我女兒可以抱著一起去聊天.
弟兄姊妹 這個是不是好像火鍊的忍耐一樣.
我不知道你今天在忍什麼.
但是聖經沒有堂醫抱著人生裡面的苦況.
正如耶穌基督也說過.

$^{121}$這個世界是有苦難的.
但是你要記得最後那句是什麼.
我已經戰勝世界.
所以希望我們可以用戰勝的心態.
來面對現在我們百般的磨難.
以致我們忍耐的功課不是被動的.
而是相當的積極.
一起讀highlight那句好不好.
一二三 堅固你們的心.
所以弟兄們 我們要堅固我們的心.
堅固這個字其實都是用剛才那個時態.
都是用conservative errors imperative.
都是說你一輩子都要忍耐.
而你一輩子一句到尾.
我們要一起學習堅固我們的心.
今天你準備好了沒有.
堅不堅固到我們的心.
如果堅固到 你可不可以給我一個amen.
我們來這裡敬拜.
我們不只是聽道 涼冷氣 聊聊天.
我們要讓神改變我們.
以致我們雖然在這個痛苦的世界裡面.
我們都可以得勝有餘.
大家一起讀好不好.
一二三 留守你的信.
在下文裡是這麼說的.
我們一起讀.
弟兄們 你們要把那先前奉主名說話的眾先知.
當作能受苦能忍耐的榜樣.
那先前忍耐的人.
我們稱他們是有福的.
你們聽見過約伯的忍耐.
也知道主給他的結局.
明顯主是滿心憐憫大有慈悲.
剛才說了一種東西.
大概就是這樣一個比喻.
之後他給一個例子給我們.
你看到他的忍耐好像不太一樣.
看和合本的時候就忍耐忍耐.
什麼叫忍耐呢.

$^{161}$原來他在說忍耐的結果就是step fast.
行不行.
他還給一個例子給我們.
就好像看方太與理.
我小時候看方太與理.
他就很快地醃完雞.
丟一點蔥下去就搞定了 但因為時間的關係我已經一早預備了一碟出來.
很美 魔法 我很喜歡看方太與理.
雖然我現在都不會煮東西吃.
但我喜歡吃東西.
我覺得很美 煮得煮再耐就快了.
人生有幾十年 剛才我們查經的時候.
人生有幾十年 煮再耐就快了.
在煮再耐之前.
因為時間的關係.
神擺了很多如同雲彩的忍耐的見證在我們身邊.
讓我們可以堅重我們自己.
就算在這個火煉的世界裡面.
我們都可以得勝有餘.
而step fast這個字在希臘文裡叫做hupomone.
為什麼.
有個加號.
哪兩個字加起來.
under什麼.
回兵.
bye.
為什麼.
我就解作留守.
姐妹我們要在我們的信心裡面留守.
我們很多時候未必留守到.
因為這個世界的痛苦太真實.
我們留守不到的時候.
我們沒有辦法掌握到神的timing.
我們有時候走得快過神.
所以剛才我祈禱的時候.
我們要靠著什麼.
我們要靠著我們的耶穌基督.
我們要跟著這個主去走.
那道經文說什麼.
他說直到得了什麼 秋雨春雨.

$^{201}$我們要等神的timing.
什麼叫秋雨.
秋來也秋去.
所以唱這首歌.
什麼叫春雨.
遲來的春雨.
以前以色列人種東西.
他就是要等兩個季節.
夏天之後地是乾的.
他要下種子.
下種子之前.
他需要秋雨去sprinkle那些泥.
以致可以搓爛它.
放粒種子進去就生長了.
這個大概是9月至10月.
但都未得過.
你長了出來它會曬焦的.
所以他要等到3月至4月.
他要等個春雨.
因為他在秋成之前.
他要多一陣雨季之後.
他才可以發芽.
才可以生長.
剛才說的precious的一個出產.
一個豐富的出產.
所以農夫在當時候.
他不單止有秋雨.
他還要等春雨.
但很多時等待的過程.
令我們沒有了忍耐.
你知不知道當時的外邦人.
他在等春雨的時候.
他們去做些什麼?.
拜偶像.
為什麼?.
覺得漂亮.
等於姊妹.
我們很多時信了主之後.
我們都好像得到了秋雨一樣.
很多神蹟令我們信.

$^{241}$本來心靈好像很乾.
然後春滿雲雜.
然後我們得以改變.
但不要忘記我們要等什麼?.
春雨.
我們等春雨的當中.
我們有很多試探.
世界會試我們.
我們有時會忍耐不到.
今日的經文就正正告訴我們.
我們需要在主裡面學習.
忍耐的功夫.
Amen.
其實申命記那裡.
他講得很清楚.
看到嗎? 秋雨,春雨.
我很快讀給你們聽.
在申命記11章.
「是也,你神所眷顧的從歲首到年終」.
「歲首到年終也,你的神的眼目時常看管那地」.
「你們若留意聽從我今日所吩咐的誡明」.
「愛也,你們的神盡心盡性事奉他」.
「他必按時幹什麼」.
「秋雨,春雨在你們的地上使你們可以收藏五穀」.
「生,酒,和,尤」.
其實當時神已經告訴以色列人.
你不用怕.
從歲首到年終.
祂都在.
你只是要跟著我.
你學習去愛我.
你看清楚舊約.
好像很早Demanding.
但是你要知道.
神已經一早愛了以色列人.
已經預備了道路.
在格欖地得地為業之後.
神告訴你.
你學習愛我.
當我們學習愛神的時候.

$^{281}$之前.
其實神一早已經施愛我們.
所以我們剛才要義愛.
還愛.
弟兄姊妹.
在等待春雨的時候.
我們要做些什麼?.
我們要學習怎樣愛神.
我們要學習聽從神的吩咐.
在聖經裡面.
弟兄姊妹.
我們要留守我們的信.
我們要充滿忍耐的功夫.
很多時候我們留守我們的信心.
我們的忍耐.
不是依靠神蹟.
不是依靠這個環境.
當你看聖經的時候.
你看五餅二魚很開心.
Napa的酒很貴.
有個我忘記什麼名字.
兩百多元一瓶.
可以試六十元一杯.
我試過一次之後很後悔.
很貴.
跟我叔叔的事.
但耶穌不用.
祂沒有酒.
隨便酒.
一講就行了.
很棒.
五餅二魚.
大家都聽過五餅二魚.
耶穌憐憫他們.
覺得他們沒有東西吃很慘.
拿五餅二魚餵飽五千人.
肚子餓就有東西吃.
很棒.
你信不信.
都信.

$^{321}$但之後那些人走了.
耶穌說這些人是來吃東西的.
弟兄姊妹我們不要做那些人.
我們要留守我們的信.
我們要像約伯一樣.
我們一起讀.
一二三.
他的妻子對他說.
你仍然持守你的純淨嗎.
你欺騙了神.
死了吧.
約伯卻對他說.
你說話像愚頑的婦人一樣.
唉.
難道我們從神手裡得福.
不也受禍嗎.
在這一切事上.
約伯並不以口犯罪.
弟兄姊妹.
二人第一次罵老婆就在這裡.
所以你不要罵老婆.
你不是二人.
但你要知道約伯那時很慘.
但他又沒有錢.
又死了兒女.
又滿身毒瘡.
總之很慘.
在當時他都沒有放棄神.
他留守他的信.
他很忍耐.
但他都不敢嘆氣.
一起讀好嗎.
唉.
一起讀.
唉.
你可以相信神的面子.
原文沒有這個字.
但他翻譯了感嘆的意思.
有時真是令我們.
唉.

$^{361}$為什麼又是這樣的時候.
弟兄姊妹.
我們會祈禱.
我們祈禱的時候.
希望神立即出手.
有時神是立即出手的.
我們開這個教會.
神是立即出手的.
立即有地方.
一個月就搞定了.
這是神蹟.
但弟兄姊妹.
我們要等這個春雨.
很多時我們會有很多的挑戰.
很多的挑戰.
裡面有很多的阻瀾.
神又未在裡面.
神由歲守到年終.
都在這裡.
我們面對挑戰的時候.
都在這裡.
或者我們突然有一天.
喇叭不響的時候.
都在這裡.
我們要怎樣.
忍耐.
我們要留守我們的信.
我們可以.
一齊讀一次.
人生就是這樣.
氣虛嘆息.
但當你嘆息的時候.
你要知道.
耶穌已經成了.
耶穌知道各樣的事已經成了.
為要使經上的話.
應驗就說我可了.
然後耶穌償了那一套.
就說.
成了.

$^{401}$便低下頭.
將靈魂交付給神了.
弟兄姊妹.
耶穌基督的死.
是不簡單的事.
耶穌基督.
他應該在世上的時候.
受很多苦難的時候.
他都.
但他留守他的信.
他忍耐到最終.
耶穌基督的死.
告訴我們.
我們的春雨未必在今生裡可以得到.
我們需要用.
永生的眼光.
去看這個春雨.
有很多時候.
你會見到弟兄姊妹病了.
醫治不了.
有時候你會見到自己很糟糕.
不知為何這麼倒楣.
我想告訴你.
我們要知道.
留守的那個.
你.
即是你.
即是我.
留守得到.
在痛苦艱苦的歲月裡.
留守得到的就是神蹟.
我們很多時候去祈禱.
就是求神.
將環境改變.
我們就好了.
但很多時候.
如果你面對這樣的環境.
你祈禱很久.
神都沒有出現.
沒有改變環境的時候.

$^{441}$如果你真的信神.
知道神很愛你的話.
即是說神要用環境.
來改變我們.
同意嗎.
我再說一次.
我們很多時候在忍耐的時候.
我們求神改變環境.
但神不改變環境的時候.
如果你真的信他.
是愛你的.
由歲首到年終.
即是說神要用環境來改變我們.
讓我們得以成聖.
讓我們得以更加像他.
Amen.
我不知道你現在面對什麼環境.
或者你.
婚姻關係的環境很差.
一會兒我們有意外.
不是意外.
是越拿越愛.
1點鐘.
或者你現在.
經濟上很大壓力.
祈禱很久都沒有聲音.
或者你健康出現問題.
或者你現在很不開心.
由歲首到年終.
神都在這裡.
Amen.
有時候我們看不到.
我們怎樣才可以知道.
我就是神.
曬著那個靜脈.
會燒焦的時候.
我們需要很多養分.
我們需要神的愛.
我們第二點是什麼.
一起讀.

$^{481}$神的愛.
孕育在愛裡面.
在經文裡面.
他講那個忍耐.
第一個忍耐.
他那個字是兩個字加起來.
哪兩個字加起來.
Long和Anger.
Passion.
你裡面要有火.
你裡面怎樣才可以有火有能力.
有力量.
你需要神很多的愛.
我將他自作主張.
恆久忍耐.
聖靈在這裡.
剛才好像Jack講.
愛是什麼.
恆久忍耐.
我們需要清楚神的愛.
我們要孕育在神的愛裡面.
你見到.
他說你們若留意.
我今日所吩咐的誡命.
愛耶華你們的神.
盡心盡性事奉他的時候.
你會得到那個愛.
我們要看聖經.
弟兄姊妹.
你有沒有試過.
用Tucker.
來比喻.
你不和太太聊天.
你怎樣可以更加愛對方.
你要溝通.
你信神的話.
你當然要知道神想怎樣.
然後你祈禱.
透過經文.
聖靈會回答你的時候.

$^{521}$那個感受.
是你獨有的.
Fiona經常說和神聊天.
你要看聖經.
我們要看聖經.
我們要學習怎樣看聖經.
除了我們待會有.
越難越愛.
我們早上都會一起看聖經.
就在ASB room.
現在越來越多姊妹來看.
你儘管走進來看.
帶一本聖經就行了.
做一個基督徒我們必須要學習.
看聖經.
你不看聖經你不知道.
原來神的愛的深度.
是這麼廣的.
以致我們信仰有一個深度.
以致我們可以留守.
以致我們留守的過程當中.
我們可以得到更加多的愛.
剛才我說的表姐.
她怎樣可以.
頂多這麼久.
我給了她很多的愛.
女兒出事之後.
她相信了主.
我說你一定要去找教會.
她去幾間.
回到那間.
她說神很好.
我的人改變.
我很平安的時候.
護士不知為何幫我寫一封信.
去申請什麼屋.
政府給了她一間比較大的屋.
因為她有很多儀器掛在那裡.
她要叫女兒去.
我的表姨身去洗澡.

$^{561}$叫她出來.
不單止這樣.
神很愛她.
她找到一間教會.
她填了申請表.
教會是找一個姊妹.
一對一照顧她.
還要幫她做治療.
她有很多治療.
要唱大聲唱歌.
當她做了治療的時候.
我表姐就可以去敬拜.
我表姐每一個禮拜.
敬拜的時候.
她都見到神的信實.
她說雖然我女兒沒有改變.
雖然這一世都是這樣.
但她見到神很愛她.
她也參加很多茶經.
她一天睡三四個小時.
每次茶經的時候.
她得到很多的滋養.
弟兄姊妹.
神沒有將她的環境改變.
但他將我表姐改變.
我表姐雖然不懂很多聖經.
但我見到她是一個神跡.
這個神跡你也是可以的.
你只要盡力.
盡心盡意去事奉.
去學習愛主我們的神.
學習聖經的時候.
你都可以.
基本上耶穌基督的救恩.
可以使我們每一個人.
都是一個神跡.
你們相不相信.
相信給我大聲的Amen.
你見到.
James在《雅各書》裡面.

$^{601}$他很著重.
又講全律法.
又講至尊的律法.
原來至尊的律法是在講什麼呢.
愛人如己.
他是在講愛.
你留守的時候.
你是需要這個愛的.
當時的以色列人.
門徒四散的時候被逼迫.
他很需要這個愛.
James說.
你們要守的是至尊的律法.
不是全律法.
因為全律法你犯了一條.
你都是不行的.
我們都是不行的.
但是至尊的律法是很難守的.
我們要學習愛你的仇敵.
有沒有仇敵.
都沒有.
有沒有不喜歡的人.
都有的.
以前很多仇敵.
就算看見街上的人.
都會罵他.
人很多仇怨在裡面.
這是很難的.
但在至尊的律法裡面.
神愛我們.
永遠給我們一個出口.
有祂的恩典.
所以當我們在這個.
苦難當中.
我們忍耐當中.
我們頂不住的時候.
我們會犯罪.
雖然我們犯錯.
但神給我們一個出口.
只要我們盡心盡意去事奉祂.

$^{641}$去學習祂的話語.
這就是一個成聖的過程.
這就是白人成聖.
不要做金.
不要成為一條金條.
我們要愛神.
在留守的時候我們要學習愛神.
耶穌問門徒彼得.
祂說.
你愛這些.
比愛我更深.
我跳一跳.
這些是什麼.
對你今天.
這些是什麼.
在你等待的過程.
的時候.
或者你的工作有些挑戰.
你會不會愛你的工作.
比神更深.
或者你在.
身體健康.
很大的挑戰的時候.
你會不會愛自己的健康.
比神更深.
有時神或許不會改變環境.
但或者.
你的轉變.
你的平安成為一個神蹟.
一個見證可以影響到.
你身邊的人.
我們每一個禮拜都有.
祈禱會.
我們都為很多不同的事情祈禱.
有個弟兄很年輕.
患了第四期.
癌症.
我們是為他祈禱的.
他痊癒了嗎.
他還沒有.

$^{681}$但他是一個神蹟.
他的平安是一個神蹟.
他的憧憬是一個神蹟.
他反而安慰了身邊的人.
這是一個神蹟.
這個神蹟這個見證.
到我們回到主耶穌基督的時候.
這句說話.
他可以在天國說.
那時候很厲害.
我都希望可以.
到時有話說.
人只有一世.
人生沒有幾個十年.
你想在天國說多一點.
現在我們就要留守我們的信.
領略神的愛.
枯乾的時候.
我們得到滋養.
你看到嗎.
我們愛到天涯.
甚麼是台國.
我想說.
很多時候.
我們都是不同的人.
我們去接收愛的禮物.
都是不同的.
有時候我們在忍耐當中.
我們很想得到愛的感覺的時候.
我們要懂得欣賞.
和弟兄姊妹.
給予愛心的禮物.
每個人都不一樣.
就像在台國.
每個人都怕被撞傷.
有些弟兄姊妹.
會寫很多東西.
你去到就不要再走了.
要斜望45度.
否則會被撞傷.

$^{721}$撞傷之後很痛.
急救時就要1,2,3,4,5.
你認識這些弟兄姊妹嗎.
或者牧師.
我不說是誰.
寫了很多東西.
這是愛.
我和他談了這麼久.
他真的愛.
他不想有甚麼事出錯.
寫了很多政策.
有些人說不要走過去.
很容易撞傷頭.
流血亂撞.
撞下去.
真的很痛.
有些人喜歡這樣.
要懂得欣賞.
有些人不是.
愛是包著角.
包著.
很多台國都包著.
很噁心.
但這是他的愛心.
有些人好像沒有愛心.
你見過有人上班沒有愛心嗎.
我太太經常說.
你好像沒有愛心.
你又不笑又不哭.
好像要流血.
有些人不是這樣.
你走過去台國時.
他這樣.
包著.
你上班回來就縮回.
他保護你.
弟兄姊妹你明不明白我想說甚麼.
在這個教會裡.
大家給予愛的禮物是不同的.
我們要懂得欣賞.

$^{761}$不同的愛語言.
我們自己都要練習.
不同的愛語言.
我正在學習.
你和我聊天.
微信.
你見到.
有些.
他寫了很多.
幫你祈禱.
整個圖文寫了出來.
寫這麼多.
剛剛祈禱完又要寫.
我不寫的.
我練習.
有時我們要懂得表達愛心.
你就知道神.
有時好像不出現.
你說歲首到年終.
神是在的.
來表達祂的愛.
這就是我們得到的愛心.
這是我們期盼主在來的.
傳言.
你想不想主在回來.
其實.
如果你看回這段經文.
第九節.
他說審判的主在門口.
為甚麼審判的主在門口.
究竟你期待的.
耶穌基督回來的耶穌基督.
是你的主.
還是審判的主.
如果你在你的人生裡.
就算很多忍耐的時候.
你得到愛滋養的時候.
你會知道.
你很想祂回來.
因為你們也當忍耐.

$^{801}$堅固了你的心.
主的日子來了.
撐多一會.
我們可以一起.
唉.
唉的時候.
我們都可以成為寶貴的出產.
成為一個見證.
白銀成聖的一個見證.
最後.
耶穌基督給我們的命令.
是甚麼.
大使命.
你懂不懂背.
你們要去.
使萬民在我的門徒.
奉父子聖靈的名.
給他們施洗.
凡我所吩咐你們的.
都教訓他們遵守.
我就是常與你們同在.
以免你這樣做.
我才與你同在.
我試試給一個不同的.
譯本給大家讀.
呂鎮忠譯本.
愚民夕夕一起讀.
所以你們要去.
使一切外國人都做門徒.
給他們施洗.
歸於父子聖靈的名.
教訓他們遵守.
我所吩咐你們的一切事.
看把一切日子.
我都和你們同在.
直到今世的完結.
清楚很多了.
不是和你說條件.
看吧.
唉的時候.

$^{841}$看一看.
我一直都在 由歲首到年終.
所以我們要學習.
怎樣留守.
怎樣等主再來.
我們要留守我們的信.
當你要.
面對很多風雨的時候.
我們要留守.
我們要尋找神的愛.
透過他的話語.
教會的團契.
以致我們可以一起期盼主再來.
或者一起祈禱.
是由歲首到年終.
我們學習忍耐功課的時候.
主啊與我們同在.
讓我們知道你是在這裡.
因為你答應了.
由歲首到年終.
從未試過離開我們.
主啊.
轉化我們.
讓我們成為更加像你的僕人.
讓這個世界看到我們這群基督徒.
就知道你是一個又真又活的神.
耶穌基督名其求.
\newpage



\section{馬太福音 24:15-24}
\label{sec:031899}
\textbf{末世三個伊斯蘭教的徵兆與真假救主}~~~~袁惠鈞牧師博士
\newline
\newline
連結: \href{http://www.church.com.hk/acms/content.asp?site=cdc\&op=show\&type=product\&code=031899\&layout=sermon}{\texttt{031899}} ~~~~ 語音日期: 2026-03-29
\newline
\newline
\hyperref[sec:031898]{\small{< < < PREV SERMON < < <}}
~
\hyperref[sec:index_chronic]{\small{[返順時目]}}
~
\hyperref[sec:031900]{\small{> > > NEXT SERMON > > >}}
\newline
\newline
馬太福音 24:15-24
\newline
\begin{longtable}{cl}
\hline
\hline
章節 & 經文 (和合本修訂版)\\
\hline
24:15 & \begin{tabularx}{0.7\textwidth}{X} 「當你們看見先知但以理所說的那『施行毀滅的褻瀆者』站在聖地（讀這經的人要會意）， \end{tabularx} \\ \\ \relax
24:16 & \begin{tabularx}{0.7\textwidth}{X} 那時，在猶太的，應當逃到山上； \end{tabularx} \\ \\ \relax
24:17 & \begin{tabularx}{0.7\textwidth}{X} 在屋頂上的，不要下來拿家裡的東西； \end{tabularx} \\ \\ \relax
24:18 & \begin{tabularx}{0.7\textwidth}{X} 在田裡的，不要回去取衣裳。 \end{tabularx} \\ \\ \relax
24:19 & \begin{tabularx}{0.7\textwidth}{X} 在那些日子，懷孕的和奶孩子的就苦了。 \end{tabularx} \\ \\ \relax
24:20 & \begin{tabularx}{0.7\textwidth}{X} 你們要祈求，好讓你們逃走的時候，不遇見冬天或安息日。 \end{tabularx} \\ \\ \relax
24:21 & \begin{tabularx}{0.7\textwidth}{X} 因為那時必有大災難，自從世界的起頭直到如今，從沒有這樣的災難，將來也不會有。 \end{tabularx} \\ \\ \relax
24:22 & \begin{tabularx}{0.7\textwidth}{X} 若不減少那些日子，凡血肉之軀的，就沒有一個能得救；可是為了選民，那些日子將減少。 \end{tabularx} \\ \\ \relax
24:23 & \begin{tabularx}{0.7\textwidth}{X} 那時，若有人對你們說：『看哪，基督在這裡！』或『在那裡！』你們不要信。 \end{tabularx} \\ \\ \relax
24:24 & \begin{tabularx}{0.7\textwidth}{X} 因為假基督和假先知將要起來，顯大神蹟、大奇事，如果可能，要把選民也迷惑了。 \end{tabularx} \\ \\
[1ex]
\hline
\hline
\end{longtable}
$^{1}$弟兄姊妹平安.
很高興又有機會來到千象城基督教會.
跟大家分享神話語.
今天要講的題目.
也挺應時的.
因為是要講伊斯蘭教.
我們知道美國現在在打伊朗.
伊朗成為一個伊斯蘭教的國家.
是1979年出事.
不是那麼久的.
今天就是要講伊斯蘭教.
就是三個徵兆和真假救主.
世上有很多宗教都是和基督教無關.
例如佛教,印度教.
甚至是日本的神道,先祖.
都是和基督教無關.
但有很多宗教又好像和基督教很有淵源,很有關係.
例如就是伊斯蘭教.
為什麼世上有那麼多好像和基督教有關的宗教呢?.
因為撒旦很想用假宗教來迷惑人.
很多人說.
伊斯蘭教的神就是我們聖經舊約的神.
你們有沒有聽過?.
他們說伊斯蘭教的阿拉神就是基督教舊約的耶和華.
也有很多人說.
其實伊斯蘭教都是相信耶穌基督,接受耶穌基督.
所以葛培里甚至說.
我認為所有愛基督或者認識基督的人.
無論他們知不知道.
他們都是耶穌基督的身體.
而他接著就說包括莫斯林人.
是不是真的這樣呢?.
是不是你只要認識耶穌基督.
只要你自己宗教來認識耶穌基督.
都是耶穌基督的肢體呢?.
我們來看看究竟莫斯林人認識的耶穌基督是怎樣.
只是一個偉大的先知.
不是神的兒子.
他們絕對不相信三位一體的神.
所以伊斯蘭教所相信的耶穌是不是我們的耶穌?.

$^{41}$絕對不是我們的耶穌基督.
我們和伊斯蘭教根本沒有共同點.
《西遊記》裡面的孫悟空都有真假美后王.
電影裡面的包公都有真假包公.
雙胞胎鬧出.
在《末世》亦都有真假的救主.
在《馬太福音》的二十四章二十四節.
因為假基督假先知將要起來顯大神跡大騎士.
倘若能行連選民也就被迷惑了.
一直在世上都有很多人冒認自己是耶穌基督.
說自己是先知.
他們能夠迷惑世上很多人.
但經文說不能迷惑我們這些神的子民或選民.
到《末世》剛才《馬太福音》的二十四章二十四節.
就會具體地被實現.
真是有假基督來到世上.
我們看假基督之前.
我們首先來看看伊斯蘭教.
他們相信的耶穌基督是怎樣.
其實伊斯蘭教和基督教是一樣.
都有我們的末世觀.
他們的末世觀都是很有趣.
根據伊斯蘭教他們的末世觀.
耶穌在末世是一個很重要的人物.
在看耶穌之前.
我們首先要知道伊斯蘭教的教義.
他們的末世觀從何而來.
就是有兩個部分.
一個是他們的古蘭經或可蘭經.
古蘭經或可蘭經是毛罕默德所寫的.
是公元大約610至632年.
是一個人寫成.
我們手上的聖經.
四十多位作者.
六十六部書.
多過幾千年.
三四千年寫成.
全部是一致的.
你要四十多位作者寫出來的東西都是一致.
後面一定是有聖靈密使.

$^{81}$你要一個人寫出來的東西是一致的話.
不是那麼艱難.
除了古蘭經或可蘭經.
他們還有他們的Hadith.
就是他們的聖訓.
或者叫做Sunnah.
是他們的常道.
又是誰寫的?.
都是毛罕默德寫的.
所以他們的教義是有兩個部分.
亞拉臣的話語就是可蘭經.
另外還有他們的傳統.
他們的常訓就是毛罕默德的話語.
其實這個和天主教很相似.
天主教有他們的聖經.
和我們的聖經都很接近.
但他們多了次經.
但是除了聖經.
他們還有很多的禮儀傳統.
這個就和伊斯蘭教一樣.
可蘭經是亞拉臣的話語.
他們認為.
但是他們仍然都有聖訓常道.
就是毛罕默德的話語.
這個和猶太教也很相似.
猶太教也有他們的神的話語.
就是我們的舊約.
但他們仍然有拉比的傳統.
拉比的話語.
所有世上的假宗教.
都有兩個權威.
就是神的權威和人的權威.
或者傳統的權威.
只有真宗教唯一的權威是什麼?.
是神的話語.
我們不妨來看看伊斯蘭教.
所說的耶穌基督是怎樣.
祂只不過是一個人.
不是神.
祂沒有死亡.

$^{121}$也沒有復活.
祂就好像以利亞先知一樣.
被神接了去.
祂因為沒有死在十字架上.
為我們來贖罪.
所以我們也不能夠被祂救贖.
在這篇我做了很多參考的資料.
我有很多參考資料.
我不想好像有很多牧師傳道人.
他們說到伊斯蘭教會怎樣.
有很多時候我發覺是三安八造.
有很多時候是無中生有.
我們不需要的.
我們只需要說出真真正正.
伊斯蘭教的教導是怎樣.
所以我做了很多的研究.
很多的參考資料.
在我這篇導的後面.
我會留在PowerPoint.
如果大家有興趣可以下載來看.
根據伊斯蘭教.
沒有人能夠替另外一個人贖罪.
所以耶穌基督只是先知.
祂完成了祂的任務之後.
就被神接了去.
祂現在就在亞拉的旁邊.
但是根據他們的教導.
在末世亞拉神會將耶穌基督再次.
猜派來到這個世上.
伊斯蘭教有很多先知.
為什麼一定要猜派耶穌基督回來.
因為根據伊斯蘭教的教導.
解靈還須是靈人.
是耶穌基督第一次來臨.
造成了很多的誤解.
所以一定要猜派耶穌基督本人回來.
糾正這一切的誤解.
尤其是對那一本書的人.
說出他們是信錯了.
那一本書的人是哪些人.

$^{161}$猶太人,基督徒.
我們相信的那一本書就是聖經.
所以耶穌基督回來是要糾正這一切.
祂沒有為我們死.
祂沒有救贖我們.
祂只不過是一個人.
祂沒有三天復活.
我們的永生盼望.
我們被救贖的希望全部都化為烏有.
這個就是所謂錯誤的教導.
耶穌基督回來是要糾正這一切.
糾正了之後.
祂就會結婚生子.
祂就會死.
然後祂就會被埋葬在無罕默不的旁邊.
這個就是伊斯蘭教所相信的耶穌基督.
你想想.
如果你認識這個耶穌基督.
能夠成為耶穌基督的身體.
絕對不是.
根據伊斯蘭教的末世觀.
在末世有三個徵兆.
環繞著三個很重要的人物.
第一個就是馬克迪.
馬克迪是他們十二個伊瑪穆當中的第十二位.
伊瑪穆的意思就是領袖的意思.
在伊斯蘭教的宗教一共有十二個伊瑪穆.
而這個馬克迪就是最後的一個.
馬克迪這個字是從一個阿拉伯文的字根HDY所出來.
這個字根在可蘭經裡你能夠找到很多處都有寫這個字根.
但是馬克迪這個名字你在古蘭經裡是找不到的.
根據伊斯蘭教的正訓或者是常道.
這個馬克迪是無罕默德的後人.
他在末世來到人間.
他要帶領世人進入公平公義.
他要戰勝所有其他的宗教.
尤其是那部書的人就是猶太人和基督徒.
他要建立世界性的伊斯蘭國.
所有不接受伊斯蘭教的人都要被處死.
他要在世上統治七年.

$^{201}$有沒有注意到這個人好像似曾相識.
他是要帶領人保證有公平公義.
他是會勝了尤勝.
他是會統治七年.
像誰呀?.
聖經裡的迪基督.
但以理書第九章第十七節說.
一七之內他必與許多人堅定文約.
聖經告訴我們.
在人間最後是有七年的真戰.
七年的大災難.
有個迪基督會走出來.
會和很多國家堅定文約.
尤其是和以色列.
他會保證公平公義.
似圖約翰.
在啟示錄的六章二節說.
並由冠冕賜給他.
他便出來勝了尤勝.
有沒有看到.
就是勝過那些那部書的人.
以前的伊拉克總統沙丹赫森.
記不記得這個人.
他曾經叫很多人幫他畫畫像.
他自己騎在一匹白馬的上去攻擊耶路撒冷.
他認為自己就是馬克迪的一個御像.
在聖誕以理書的八章二十四節.
經文又說.
他的權柄必大.
卻不是因自己的能力.
他必行非常的毀滅.
事情順利任意而行.
又必毀滅有能力的和聖文.
又是那部書的人.
說神給權柄讓他勝了尤勝.
聖文必交在他的手中.
一載二載半載.
加起來就是三年半.
基督會和以色列來立約.
三年半之後.

$^{241}$他便會毀約逼迫以色列人.
這三年半是非常難捱的一段時間.
在聖經裡說得很清楚.
所以很多人說.
根本沒有這件事.
沒有七年的災難.
你看看.
用三個不同的方法.
告訴你這三年半.
三年半四十二個月.
一千二百六十日.
全部都是三年半.
在聖經裡說得很清楚.
根據伊斯蘭教的教導.
穆斯林人將會和猶太人.
作最後的真戰.
穆斯林人會追趕猶太人.
使猶太人沒有地方躲.
躲在石頭和樹的後面.
然後根據他們傳統的故事.
連這些石頭和樹都會出聲說.
真主的僕人真主的僕人.
是誰啊.
穆斯林人啊.
我身後就有猶太人.
你們快點來殺死他們吧.
只有一種樹不會這樣說.
這是伊斯蘭人自己說的.
這些樹叫做葛爾喀德樹.
因為是猶太人的樹.
所以他們不會說猶太人躲在他們後面.
根據穆斯林的學者.
有大量這些葛爾喀德樹.
在樹以色列.
其實這些樹就是我們所謂的黃羊樹.
真的在樹以色列有很多很多.
那麼這個重點是什麼.
伊斯蘭人和猶太人將會有一場很大的真戰.
很大的屠殺會發生在以色列.
我都在網上查一查.

$^{281}$是不是伊斯蘭教真的要將以色列人趕盡殺絕.
但是在網上找不到.
伊斯蘭教不會承認.
但如果你們查一件事就很有趣.
你們查一下有什麼國家曾經誓言要將以色列人趕盡殺絕.
要將他們滅絕.
你能夠發覺維基百科列出這些國家.
所有很多阿拉伯國家.
加上埃及,伊朗都要將以色列人趕盡殺絕.
這個就讓我們看到.
其實真的是伊斯蘭人的教導.
是要這樣做.
將以色列人趕盡殺絕.
這個馬哈迪的軍隊將會由一個叫做庫拉斯坦的地方出發.
庫拉斯坦在哪裡呢?.
就是現在的伊朗.
所以我今天這篇導讀真的很應時.
你知不知道現在伊朗的軍隊都是拿著一面黑旗.
根據伊斯蘭教的教導.
將來馬哈迪的軍隊都會拿著黑旗.
現在伊朗的軍隊就是拿著黑旗.
上面寫著懲罰這個字.
這個馬哈迪最終會到達伊拉.
就是耶路撒冷原頂的清真寺.
在那裡要樹立他們的軍旗.
很有趣.
聖經裡面也告訴我們.
這個迪基督將會聯合很多的列國來攻打以色列.
我們今天剛才讀過的經文.
撒迦利亞書的十二章第二節.
我必使耶路撒冷被圍困的時候.
向四圍列國的民成為令人昏醉的杯.
這物事也論道猶大.
這個是舊約.
新約又怎麼說呢?.
更加的具體.
在馬太福音的二十四章十五至二十三節.
經文說:你們看見先知但以你所說的話.
還有一節經文在舊約.
就是但以你所說的第九章.

$^{321}$也有說同樣的話.
那毀壞可憎的就是迪基督.
他站在聖地.
那時在猶太的應當逃到山上.
在房上的不要下來拿加里的東西.
在田裡的也不要回去取衣裳.
為什麼呢?很迫切.
當迪基督來攻擊以色列人的時候.
你沒有時間做這些事了.
你要逃命.
當那些日子懷孕的和奶孩子的又和了.
你想想如果你懷孕.
你拖著小孩逃難更加辛苦.
你們應當祈求叫你們逃走的時候.
不遇見冬天或安息日.
冬天山路難行.
安息日有很多事都不能做.
千萬不要在那個時候迪基督來攻擊你們.
因為那時必有大災難.
從世界的起頭直到如今.
沒有這樣的災難.
後來也必沒有.
若不減少那日子.
凡有血氣的沒有一個能得救.
只是為選民那日子必減少了.
那時若有人對你說.
基督在這裡或說基督在那裡.
你們不要信.
馬太福音的二十四章二十四節.
就是我開始的時候和你們讀的經文.
末世又怎樣?.
有假基督有假先知.
所以我們講完伊斯蘭教的馬赫迪.
就是聖經裡面的假救主.
這裡說還有假先知.
所以帶領我們去到伊斯蘭教.
第二個重要的人物是誰?.
就是他們的耶穌伊薩.
他們的耶穌不是我們的耶穌.
千萬不要搞錯.

$^{361}$他們的耶穌也不是馬赫迪.
根據他們的信仰.
馬赫迪是大過耶穌.
因為他大過耶穌.
才可以指出耶穌的錯.
第一次來臨.
令到這麼多人誤解你.
根據伊斯蘭教的信仰.
我剛才也講過.
耶穌只不過是一個偉大的先知.
不是神.
他沒有死,沒有復活.
也沒有救贖人.
他再次回來.
是要帶領所有的基督徒.
去放棄他們本來對耶穌基督的觀點.
對耶穌基督的看法.
要他們接受所謂真正的耶穌.
就是這個只不過是一個先知.
不是神的耶穌.
他回來的時候.
將會成為一個莫斯林激進的分子.
他會搭著兩個天使的翅膀回來.
去到敘利亞大馬士革.
他會在一座白色尖塔附近降落.
降落之後他會怎樣?.
他會向比他更大的馬赫迪禱告.
他會承認馬赫迪是他的主.
他會前往麥加利亞朝聖.
來拜亞拉神.
他將會成為莫斯林最偉大的宣教士.
在這裡末世還會怎樣?.
還要審判所有非莫斯林的人.
那些不相信伊斯蘭教的人.
要指出他們的罪行.
這個耶穌回來唯一的目的是什麼?.
就是要成為馬赫迪得力的助手.
大家記不記得在聖經裡.
迪基督最得力的助手是誰?.
是賈先知.

$^{401}$所以其實他們的耶穌.
就是聖經裡的賈先知.
他回來有三件事要做.
第一件事就是要摧毀所有十字架.
非常諷刺.
你想想.
耶穌基督要摧毀十字架.
什麼意思?.
就是他回來會迫不及待教會.
這個就是這個意思.
根據伊斯蘭教.
他們的耶穌會效忠迪基督.
他就是迪基督的先知.
聖經裡所說的賈先知.
啟示錄裡第十三章第十五章告訴我們.
聖經裡的賈先知會叫人屈服.
向他們的獸來敬拜.
獸是誰?.
獸就是迪基督.
所以這裡是一模一樣.
我們聖經裡的賈先知.
就是他們的耶穌.
除了要摧毀所有十字架.
這個耶穌還有神跡啟示的能力.
聖經也告訴我們.
賈先知也有神跡啟示的能力來迷惑人.
不僅如此,他的口還非常了得.
他還會花言巧語來哄人.
來敬拜獸,敬拜迪基督.
除了他會摧毀所有十字架.
除了他有神跡啟示的能力.
還有一樣很有趣.
就是他會除去一種叫Jesua的歲修.
根據伊斯蘭教的教義.
這個耶穌會建立全球性的律法.
以致他可以除去這個歲修.
很多學者對這個歲修都有不同的意見.
但有一個比較合理的解釋.
就是當時迪基督的政府.
他們的背後都有宗教信仰的動機.

$^{441}$所以他們會勸人加入伊斯蘭教.
接受這個授印.
接受之後就會減你的稅.
你不需要繳稅.
如果你是進入伊斯蘭教的話.
我們美國人就最認識繳稅.
美國人最怕繳稅.
因為繳得這麼重的稅收.
所以你想想,到了末世.
真的有很多人可能因為不想繳稅.
而加入伊斯蘭教.
你加入了伊斯蘭教.
受了授印,你才能夠做買賣.
其實我們現在已經有了二維碼.
你拿手機出來「嘟一嘟」就給錢給我.
我拿手機出來「嘟嘟」又給錢給你.
我最近聽見有些人去完大陸回來.
那些黑衣都不收現金了.
他拿手機出來「嘟一嘟」.
又過了他的銀行.
將來你沒有了這個授印.
你怎樣做生意?.
你連生存都生存不了.
你要加入伊斯蘭教.
受了這個授印.
你才能夠做生意,可以生存.
這個迪基督的授印是什麼?.
我們都知道是666,三條六.
聖經沒有告訴我們.
三條六是什麼意思.
但是有說得很清楚.
六是人的數目.
七是一個完美的數字.
所以六永遠都是差一.
去不到七.
不管你的人有多大的權勢.
有多大的智慧,你有多厲害.
你靠自己永遠達不到神的標準.
這個就是原因我們需要耶穌基督.
根據他們的教導.

$^{481}$這個耶穌最後就會殺死他們的假救主.
這個就帶領我們去到他們.
第三個人物就是他們的假救主.
叫做達加爾.
達加爾這個名字的意思是混合的意思.
或者混亂,迷惑人的意思.
這個人就是伊斯蘭教的賈彌賽亞假救主.
根據他們的描述.
這個人要比馬赫迪和他們的耶穌更加曲折離奇.
這裡我只能說出他們最普遍的傳統.
對這個人有什麼認識.
有人說他的右眼是盲的.
有人說他的左眼是盲的.
有人說他根本不是盲的.
只不過他的眼珠是怎樣.
好像浮動的葡萄一樣.
很生鬼.
這個人會做什麼.
他是要阻止伊斯蘭教的馬赫迪和伊斯蘭教的耶穌.
所以為什麼伊斯蘭人認為這個是假救主.
因為他是敵對他們的馬赫迪和耶穌.
其實根據伊斯蘭教的教導.
這個人和我們相信的耶穌是很相似的.
首先是有四點.
首先他是一個猶太人.
我們知道耶穌基督是一個猶太人.
出生在伯利恆.
第二點.
這個人將會自認是彌賽亞.
這個伯加爾他會自稱是真主.
真主的意思是什麼.
就是彌賽亞的意思.
根據他們的教導.
這個伯加爾將會騎著一隻驢驅子行遍全地.
我們知道聖經裡面的耶穌基督.
第一次來臨的時候.
騎著驢驅子以王的身份進入耶路撒冷城.
在撒迦利亞書的九章九節.
他這樣就證明他是誰.
證明他自己就是聖經裡面的彌賽亞.

$^{521}$所以他有承認自己是彌賽亞.
而且他真的親口有承認.
在約翰福音第四章是向撒瑪利亞婦人.
在馬可福音十四章是向彼拉多.
他也有承認自己是彌賽亞.
所以這個就像他們伊斯蘭教裡面所說.
他會承認自己是彌賽亞.
還有就是他有神跡騎士的能力.
根據伊斯蘭教的教導.
這個伯加爾是有醫治的能力.
他只要手摸一摸任何人.
那個人就會痊癒.
甚至可以使人起死回生.
但是根據伊斯蘭的教導.
這個只不過是現眼法,不是真的.
但是我們的耶穌真真正正是一個醫治者.
真的能夠第一次來臨的時候.
有任何人請他醫治的話.
他都來謝不拒.
而且他也能夠使人起死回生.
在約翰福音十一章.
使拉撒路起死回生.
還有第四點.
他是有統治的能力.
根據伊斯蘭教的教導.
伯加爾是有權柄來統治這個世界.
聖經也告訴我們.
耶穌基督是有權柄.
他第二次來臨.
他是會帶領我們進入千禧年國.
他是有統治的力量.
萬失都要敬拜,萬口都要承認.
就是他第二次來臨.
最後的結果是怎樣呢?.
最後的結果就是.
伊斯蘭教他們的耶穌.
會殺死伯加爾.
就是假的救主.
但是根據我們的聖經.
剛剛相反.

$^{561}$耶穌基督是會戰勝並殺死迪基督和賈先知.
在哪裡見到?.
在啟示錄十九章第二十節.
「那獸被擒拿,那在獸面前曾行騎士,迷惑獸獸印記.
和拜獸像之人的賈先知也與獸同被擒拿.
他們兩個就活活的被冰在燒著的硫磺火壺裡」.
最後迪基督和賈先知都會被耶穌基督制服.
所以綜合以上我們見到.
伊斯蘭教他們的救主.
所謂的救主.
就是我們聖經裡的迪基督.
他們的耶穌就是聖經裡的賈先知.
還有第三個人物.
他們的假救主伯加爾.
他就是我們真正的耶穌基督.
我們其實已經說完了.
在末世伊斯蘭教三個重要的人物.
最重要就是迪基督.
但不妨我們來看看.
究竟迪基督的國又是一個怎樣的國呢?.
我們知道迪基督可能是怎樣.
但是迪基督的國.
最好入手的經文就是啟示錄十三章.
告訴我們當獸出來的時候又怎樣.
獸就是迪基督.
有七頭十國.
七頭是什麼?.
是七個帝國.
由埃及,亞述,巴比倫,馬代波斯,希臘和羅馬.
一共已經是六個國.
還有第七個國.
其實是十個聯合國.
是迪基督的十個聯合國.
就是那十個國來代表.
如果我們記得的話.
但以里斯裡面的第二章.
記不記得.
尼布格尼撒他夢見有一個人像.
有一個巨人的人像.
他的頭是金造的.

$^{601}$是巴比倫.
他的胸和臂是銀造的.
是馬代波斯.
你說臂是兩邊的.
馬代和波斯很有意思.
還有他的肚和腰是銅造的.
就是希臘.
還有兩條腿是鐵造的.
是羅馬帝國.
有人說:牧師,羅馬帝國只有一個.
怎麼會有兩條腿呢?.
等一下我會解釋.
兩條腿之後是怎樣呢?.
有十隻半泥半鐵的腳趾.
這十個就是十個聯合國.
就是我剛才所說的十隻國.
你有沒有看到聖經完全一致.
所以迪基督來到的時候.
是有十個聯合國會支持他.
這就是啟示錄裡說的十隻國和十個冠冕.
這十個聯合國將會是羅馬帝國演變而成.
因為是兩條腿接著十隻腳趾.
很多人都留意到歐洲的局勢.
就是看這十個國家.
他們尤其注意的是什麼呢?.
共同市場,歐洲的共同市場.
曾經共同市場有十個成員.
那些基督徒開心得不得了.
你看看這個還不是那十隻腳趾.
現在已經超過十個成員.
但最重要的是耶穌基督回來的時候有多少成員.
很多學者注意歐洲.
但其實我覺得同樣應該注意的就是中東.
我剛才說羅馬帝國是兩條腿來代表.
聖經真的很有預言.
因為我們知道羅馬帝國是分開了西羅馬帝國和東羅馬帝國.
西羅馬帝國是在歐洲慢慢瓦解了.
東羅馬帝國是去到中東.
西羅馬帝國瓦解了之後.
東羅馬帝國仍然存在超過一千年.

$^{641}$在聖經裡面所說的新約聖經所說的羅馬帝國.
百分之六十以上的領土已經被伊斯蘭教控制.
而且迅速地向西方發展.
我剛才說的迪基督可能是出於伊斯蘭教.
所以令到這個更加有意思.
究竟將來迪基督的國是十個歐洲的國家.
還是十個中東的國家.
還是好像兩隻腳一樣.
五個是歐洲的國家.
五個是中東的國家.
沒有人知道.
但是我相信迪基督絕對和歐洲和中東很有淵源.
極可能是中東和歐洲出.
也可能有雙重的身份.
他可能是一個歐洲人在中東長大.
他可能是一個中東人在歐洲長大.
也有可能他的父母一個是歐洲人一個是中東人.
有一節的經文很有趣.
就是但以理書的七章八節.
見其中又長起一隻小國.
先前的國中有三國在這國前.
連根被他拔出來.
這國有眼像人的眼.
有口說誇大的話.
我剛才說迪基督的聯合國有十個聯合國.
但以理書的七章八節告訴我們.
其中三個國會被小國.
就是迪基督自己的國家連根拔起.
除了他們.
十個國家除了三個還剩下多少個.
剩下七個.
有很多學者都會問.
這個小國就是迪基督自己的國家.
究竟是十個國家中其中一個國.
還是第十一個國.
其實答案應該就在約書亞記裡.
這麼奇怪.
約書亞記裡.
我曾經說過約書亞記.
其實約書亞記和啟示錄很相似.

$^{681}$約書亞記其實是啟示錄的一個預表.
一個預象.
真是約書亞記和啟示錄有很多相似的地方.
你看見我這個表在powerpoint你就知道.
密密麻麻寫到密密麻麻.
我今天沒時間和大家說所有相似的地方.
我只需要說兩點就行.
第一點就是約書亞進入迦南地.
是要戰勝迦南十族.
耶穌基督在末世.
要戰勝迪基督的十個國家.
有沒有見到相似的地方.
還有一件事很有趣.
約書亞的名字是希伯來文的耶和華氏救主.
耶穌基督是希臘文的耶和華氏救主.
有沒有見到相似的地方.
我們回頭看一看.
為什麼我說答案在約書亞記裡.
根據約書亞記.
他要進入迦南.
戰勝迦南十族.
在創世記的十五章神已經說.
你們將來會進入迦南地.
迦南有十個民族.
但是當約書亞去到的時候.
他發覺已經變了七個民族.
不再是十個民族.
有三個民族不見了.
但是也少了一種人叫利弗欽人.
你說牧師不見了三個民族.
只剩下七個.
就好像《但義理書》裡所說.
但如果你說又不見了利弗欽人.
就不見了四個民族.
不是三個.
最好的解釋是什麼.
利弗欽人變成了希美人.
這樣又怎樣.
是真真正正不見了三個民族.
但是剩下七個民族.

$^{721}$這就是照著《但義理書》第七章第八節所說.
所以其實約書亞記裡面給我們一個兆頭.
就是這個敵基督的國.
是十個國裡面其中一個國.
而這個國是另外一個國變了他.
不是他本身就是那個國家.
所以不是第十一個國.
但是就是另外一個國家變了他.
其實這件事也很有趣.
因為如果你記得的話.
希美人早在《創世記》的三十四章已經被滅族.
在《創世記》的三十四章.
雅各有一個女兒叫做底那.
有一種人叫做希美人.
他族長的兒子叫做時劍.
看中了底那就強姦了他.
然後以色列人很憤怒.
尤其是利美和西勉.
他說你們想娶我們以色列的女子嗎.
你們要行國禮.
於是就騙他們行國禮.
當他們行國禮痛得不得了很弱的時候.
利美和西勉就將他們全部殺了.
所以在《創世記》三十四章已經滅了族的希美人.
你說約書亞的時候.
去到迦南地為什麼又有希美人.
就是死灰復燃.
我相信就是利弗欣人.
用了希美人的名字來重新建立這個民族.
你又會問為什麼利弗欣人.
為什麼不做利弗欣人要做希美人.
其實在聖經裡面利弗欣人不是那麼好聲譽.
記不記得在《創世記》的六章四節.
和《民數記》的十三章三節.
都說當時有一種人叫做偉人.
這種偉人有很多解經家都說.
是墮落的天使和當時的女子所生.
因為經文說是神的兒子和人的女子所生出來.
成為偉人.
記不記得當時神要做什麼.

$^{761}$要用洪水來毀滅人類.
其中一個原因我相信就是有這些事發生.
所以神是要毀滅人類.
所以這些利弗欣人就是偉人的親戚.
不是有好的聲譽.
有可能就是這樣他們改變了.
用希美人的名字來重新建立這個民族.
所以希美人就是死灰復燃的民族.
還有一樣東西很有趣.
希美人代表了新約翟基督的國.
你回看在述但以理書的第九章預言.
將來迪基督的國會和以色列立約.
是一個詭詐的約.
我們剛才也說過.
三年半他們就會毀約.
所以是用欺騙的方式來和以色列人立約.
希美人在約書亞記第九章.
記不記得他們又是欺騙以色列人.
當時他們知道神要毀滅所有的迦南人.
他們就扮到自己好像很遠來的.
不是迦南人.
很遠很遠水路來到和以色列人立約.
又是一個欺騙的約.
有沒有看到.
所以我說這些希美人其實就是迪基督國的一個儒匠.
牧師你說了這麼久都沒有告訴我.
迪基督的國在新約究竟是哪個國家.
其實一個很好的經文就是啟示錄十七章第十節.
「又是七位王五位已經傾倒了.
一位還在一位還沒有來到.
他來的時候必須暫時存留」.
這節經文很有趣.
五個國家已經來到.
歷史告訴我們.
就是埃及,亞述,巴比倫,馬代,波斯,希臘.
一個還在.
當使徒約翰寫啟示錄的時候.
羅馬帝國還存不存在.
仍然存在.
而約翰告訴我們.

$^{801}$還有一個未來就是羅馬帝國所變的十隻腳趾.
十個聯合國還沒有來到.
為什麼經文說暫時存留.
因為他們只有三年半.
其他所有的國家.
他們的國運了幾百年.
這個只有三年半.
所以你說是否暫時存在.
聖經的預言真的非常精彩.
你看完啟示錄第十七章第十一節.
就告訴我們迪基督的國是誰.
就是小國但以理書第七章裡.
十七章十一節.
「那先前有如今沒有的獸」.
獸就是迪基督的國.
「就是第八位.
他也和那七位同列並歸於沉淪」.
前面那七位和他們同列是誰.
你不要說我沉氣.
又是埃及 亞述 巴比倫 馬代波斯 希臘.
羅馬一就是羅馬帝國.
羅馬二就是十個聯合國.
迪基督的國是誰.
第八個就是獸的國.
「那先前有如今沒有」.
你看看前面那些國家.
有誰沒有做.
現在有沒有埃及.
有埃及.
有沒有巴比倫.
你說沒有.
但現在的伊拉克其實就是巴比倫.
你說有沒有馬代波斯.
當然有.
現在和美國打仗打得乒鈴bang lang .
那個就是伊朗 馬代波斯.
有沒有希臘.
有希臘.
有沒有羅馬.
有 現在是意大利.

$^{841}$還有那十個半天半泥的腳趾未出現.
所以我們不知道.
但在清單上有哪個國家沒有了.
亞述國.
很可能亞述國是迪基督的國.
將會死灰復燃.
如果真是這樣.
那時候像不像希美人.
是死灰復燃的國家.
所以聖經真是精彩到不得了.
也讓我們看到.
真是很可能末世迪基督的國就是亞述國.
還有一件事很有趣.
在聖經裡亞述國真的有比喻迪基督的國.
你說真的?.
是的 以賽亞書的七章至十四章.
還有尼加書的五章五節至十五節.
你看看尼加書的五章五節.
當亞述人進入我們的地境.
踐踏宮殿的時候.
我們就立起七個墓者.
八個首領攻擊他.
亞述國何時曾經踐踏耶路撒冷的宮殿?.
有沒有發生過?.
在歷史當中從來沒有發生過.
最近的就是在希西加的時候.
記不記得亞述王西拿基納.
走來圍困耶路撒冷?.
但是希西加走去告狀.
告給神聽去聖殿.
神啊 你看這些亞述人真是目無神.
他們是要踐踏我們.
於是一晚之間.
神就將十八萬五千個亞述士兵全部殺了.
所以根本沒有發生過.
這一節的經文是何時發生的?.
末世迪基督的時候.
來攻打以色列人才會發生.
以賽亞書的十章十二節說.
主在石安山和耶路撒冷擊敗亞述國.

$^{881}$在以賽亞書的十四章二十五節說.
在神的地上山上擊敗亞述國.
這兩節的經文.
你可以說是希西加的時候已經實驗了.
但是我相信是有雙重的預言.
將來會再發生.
就是在迪基督的國的身上.
而在說新約也預言了.
耶穌基督會戰勝迪基督的國.
就是在《書啟示錄》的第十九章.
另外有一件事很有趣.
我們剛才讀了尼加書的五章五節.
讀了尼加書的五章六節.
叫亞述國做什麼?.
叫亞述國的關口做「寧祿」的關口.
「寧祿」是誰?.
是世上第一個獨裁者.
迪基督是世上最後一個獨裁者.
有沒有看到世上第一個獨裁者和最後一個獨裁者.
都和兩個地方有很大的淵源.
一個就是巴比倫.
一個就是亞述國.
《創世記》的十章第十至十一節.
寧祿國的起頭是巴別.
巴別在哪裡?就是巴比倫.
所以他和巴比倫很有淵源.
末世的迪基督將會以巴比倫成為他政治和經濟的中心.
在《啟示錄》的十四章,十六章,十七章和十八章都有告訴我們.
還有亞述國呢?.
寧祿的起頭是巴比倫.
但他去了亞述.
建立利尼尼成為亞述國的首都.
有沒有看到他和亞述國也是有淵源.
而末世的迪基督,我們剛才說過.
他會令亞述國死灰復燃.
所以無論第一個獨裁者和最後一個獨裁者.
都和這兩個地方有很大的關係.
在處於末世的時候,我們要注意時事.
我剛才說過.
不單單要注意歐洲,我們還要注意中東.

$^{921}$現在讓我們看到.
不單止要注意亞述國.
即昔日的亞述國.
也要注意伊拉克和敘利亞,為什麼呢?.
因為在歷史告訴我們.
在巴比倫形成之前.
伊拉克和敘利亞都是叫做亞述.
所以如果我們注意這個部分的地勢.
注意有沒有國家出來.
不單止伊拉克和敘利亞,還要注意亞述國.
在這裡結束的時候,我想和大家看看.
我今天和大家說了.
伊斯蘭教的耶穌是否等於我們的耶穌?.
絕對不是.
伊斯蘭教的神是否我們的神?.
也不是.
所以不要以為伊斯蘭人就是認識神,認識耶穌基督.
不用和他們傳福音,絕對要和他們傳福音.
還有在《馬太福音》24章42節說什麼?.
所以你們要警醒.
因為不知道你們的主事那一天來到.
我們都要警醒.
我以前時時在想,耶穌說要很久才回來.
為什麼呢?.
剛才說以色列有列國來攻擊他們.
現在有誰保護他們呢?.
美國,是吧?.
是一哥來的,保護他們.
哪個列國有膽量來攻擊他們呢?.
有,但規模不大.
你看他們現在打伊朗都沒有人敢出手.
但我隨即想到.
一聲浩響.
當所有基督徒被提去之後.
美國還會不會幫助以色列呢?.
當然不會.
反過來可能攻擊以色列就差不多.
所以真的隨時可以發生.
一聲浩響.
基督徒走了七年的災難就開始.

$^{961}$極基督就來到世上.
剛才我們都看到.
極基督來到世上是多麼恐怖.
《馬太福音》24章第15至23節.
所以我們真的要趕緊傳福音給我們身邊的人.
如果我們身邊有我們很愛的朋友,親戚,微信主.
我們真的要快些傳福音.
以致他們能夠免去這一場災難.
我們自己也要警醒.
要做好自己.
否則一會兒別人得救.
反而我們自己不得救.
好像保羅所說.
願神祝福他自己的話語.
讓我們低頭禱告.
讓我們看到.
這個末世很快就會來到這個世上.
現在已經是蠢蠢欲動.
願意我們也能夠警醒.
能夠將福音傳給我們身邊還未信主的人.
願意我們也能夠做好自己.
以致你回來的時候.
我們也能夠被你稱為一個忠心良善的僕人.
將榮耀頌讚歸給你.
我們現在禱告.
奉主耶穌基督的名而求.
阿們.
\newpage



\section{箴言 6:6-11}
\label{sec:031900}
\textbf{重塑生命系列 (8) 懶惰與安息}~~~~張健庭牧師博士
\newline
\newline
連結: \href{http://www.church.com.hk/acms/content.asp?site=cdc\&op=show\&type=product\&code=031900\&layout=sermon}{\texttt{031900}} ~~~~ 語音日期: 2026-03-29
\newline
\newline
\hyperref[sec:031899]{\small{< < < PREV SERMON < < <}}
~
\hyperref[sec:index_chronic]{\small{[返順時目]}}
~
\hyperref[sec:031901]{\small{> > > NEXT SERMON > > >}}
\newline
\newline
箴言 6:6-11
\newline
\begin{longtable}{cl}
\hline
\hline
章節 & 經文 (和合本修訂版)\\
\hline
6:6 & \begin{tabularx}{0.7\textwidth}{X} 懶惰人哪，你去察看螞蟻的動作，就可得智慧。 \end{tabularx} \\ \\ \relax
6:7 & \begin{tabularx}{0.7\textwidth}{X} 螞蟻沒有領袖，沒有官長，沒有君王， \end{tabularx} \\ \\ \relax
6:8 & \begin{tabularx}{0.7\textwidth}{X} 尚且在夏天預備食物，在收割時儲存糧食。 \end{tabularx} \\ \\ \relax
6:9 & \begin{tabularx}{0.7\textwidth}{X} 懶惰人哪，你要睡到幾時呢？你甚麼時候才睡醒呢？ \end{tabularx} \\ \\ \relax
6:10 & \begin{tabularx}{0.7\textwidth}{X} 再睡片時，打盹片時，抱著雙臂躺臥片時， \end{tabularx} \\ \\ \relax
6:11 & \begin{tabularx}{0.7\textwidth}{X} 你的貧窮就如盜賊來到，你的貧乏彷彿拿盾牌的人來臨。 \end{tabularx} \\ \\
[1ex]
\hline
\hline
\end{longtable}
$^{1}$弟兄姊妹平安.
那麼我們再一次低頭祈禱.
好 天父我們願你帶領我們.
我們願你今天的話語再一次.
提醒我們 教導我們.
願你在我們心中說話.
不單止是聽到人的聲音.
也都聽到神的呼召.
奉耶穌基督的名祈求.
阿們.
好了 那麼今天就說.
七宗罪的最後一宗.
那麼今天就說到懶惰這一宗罪了.
那麼就.
也是經過了一個系列的時候.
其實都讓我們去反映到人性.
其實是在神面前.
我們是需要怎樣去自省和對付自己.
那麼最後這個.
說到懶惰的罪.
我相信就是.
我們大家都很不陌生的 是不是?.
就是.
無論你多勤力的時候.
都想懶一下的.
那麼我想這個也是人性的一部分.
就是都不是.
不是說特別是.
某些人才有懶惰而已.
其實都不會的.
那麼特別.
你看到現在.
在之前中國.
我不知道你有沒有聽過一個字叫做老鼠人.
有沒有聽過?.
他們很流行聽.
之前已經說躺平了一段時間了.
是不是?.
大家都聽過躺平了.
那麼.

$^{41}$但是躺平也是躺在那裡不動.
都發覺.
都不夠低.
那麼.
後來有些人製作那些短視頻的.
就去推銷一個字叫做老鼠人.
他就說老鼠人的一天是怎樣過的呢?.
老鼠是怎樣的呢?.
老鼠就是.
你捐了進去老鼠洞那裡.
平時白天是不出來的.
接著呢.
到了夜闌人靜沒有人的時候.
就走出來.
看看有什麼東西吃.
然後找到一個東西吃.
接著就擔回到老鼠洞那裡.
接著就在那裡繼續吃.
又不再出來.
接著呢.
第二天也是這樣.
第三天也是這樣.
一天就是這樣.
老鼠人的一生就是.
有一點點東西吃.
拿回到洞裡.
夠吃就行了.
接著就不會走出來.
所以老鼠人這個字就是這樣.
那些自媒體的人就說.
我老鼠人的一天是怎樣過的呢?.
他就用最少的錢.
就好像躺平一樣.
什麼都不做.
就好像老鼠一樣在洞裡.
什麼都不出來.
但是可能用最少的錢.
可能譬如今天用了兩塊買了一個杯麵.
今天吃了一個杯麵.
接著就繼續整天坐在那裡打遊戲機.

$^{81}$看電視.
就過了一輩子.
他說今天老鼠人的一天.
就只是用了兩塊就可以過了一天了.
這樣就很火熱起來.
在網上很多人有同感.
你就看到其實這是一個挺令人心裡面有同感.
才會火熱起來的.
你想一下那些短視頻.
當然是那些人看到之後.
不僅覺得很有趣.
而是覺得.
以真正講正我的心聲.
我也想做一個老鼠人.
不過不行了.
很累了要上班了.
接著就一傳十,十傳八.
那些短視頻就火熱起來.
就有一個老鼠人這個名字了.
但是這個現象本身跟我們今天說的懶惰似乎是剛剛相反的.
你說老鼠人就是懶惰的極致.
但是我相信.
今天我想從另一個角度去看.
就是今天所看的這個題目不僅僅是懶惰.
也是講到安息的.
我相信本身如果從這個懶惰這個事情本身.
是跟安息是很有關係的.
等一下我會解釋為什麼會有關係.
一開始的時候.
剛剛讀過那一段聖經.
就講到螞蟻.
講到螞蟻是很厲害的.
沒有人指著牠.
指揮牠怎麼去找東西吃.
但是牠自己就懂得找東西吃.
又會儲一些食物過冬.
這個是聖經的觀察.
這個也是很真實的.
就是說如果你懶的時候.
你就學一下螞蟻.

$^{121}$你學一下螞蟻你就不會這麼貧窮了.
這麼貧乏了.
但是問題就是.
首先我要問一個問題.
就是究竟.
怎樣才算是懶惰?.
是不是不做事就代表懶惰呢?.
我問大家.
是不是不做事就代表懶惰呢?.
我相信聖經對於懶惰的定義.
是沒有這麼狹窄的.
不做事有時候也不是代表懶惰.
但是聖經對於懶惰的定義.
第一個定義它就說.
因勤不可懶惰,要心裡火熱.
常常服侍主.
懶惰是對召命的一種拒絕.
這個就超越了.
你做不做事這麼簡單.
其中在聖經的一個定義就是.
我們每個人其實是有一個召命的.
人活在世上.
是有想神擺在我們的位置.
要完成某一些工作的.
這個就是聖經的召命的觀念.
每個人.
在我們人生當中.
我們都不是一個.
好像浮萍一樣.
任由水沖來沖去.
好像落葉一樣.
任由風吹來吹去.
我們每個人存活在世上.
在神的眼中.
神看我們是有一個使命的.
我們每個人.
都是有我們的使命活在世上的.
你做媽媽的時候.
你的使命可能就是照顧你的小朋友.
照顧孩子.

$^{161}$帶領他們一直成長的時候.
做一個正直善良的人.
這個是媽媽的使命.
是不是?.
可能你在工作的時候.
你做一個員工的時候.
你是很努力地為這個公司.
去完成這個公司所交託給你的職務.
不是好像做破壞的.
或者只是白鬥人工的.
這個是你在工作當中的一個使命.
是不是?.
聖經就很清楚在羅馬書第十二章所說.
他說:因勤不可懶惰.
要心裡火熱常常服侍主.
接著後面還有一連串的說法.
他說如果你有恩賜去款待的話.
就要熱心地去款待人.
你有恩賜去教導人的話.
就要努力地去教導.
你有恩賜去勸勉人的話.
就要勸勉別人.
你看回羅馬書那段經文就是這麼說的.
意思就是說.
神在我們每個人的生命當中.
其實給了你獨特的有不同的恩賜.
有不同的品格.
有不同的技能.
有不同的經歷.
有不同的熱忱.
這個就是神做你是很特別.
這個也是你在活在世上的使命和召命.
你要發揮它的時候.
就是發揮你的召命.
所以聖經所說.
如果每個人都有一個這樣的召命的時候.
我們就不要懶惰.
所以懶惰的意思就不只是說.
今天我什麼都不做.
或者今天我做的事情很少.

$^{201}$一件事兩件事.
做不做.
這樣去看懶惰.
懶惰首先來說就是.
你看回你人生的召命.
有沒有火熱的去完成.
這個就是今天說到懶惰的第一點.
這一點是很重要的.
事實上有些人.
我們廣東話有句話叫做無事忙.
有沒有聽過.
無事忙是什麼意思.
就是說你看到他.
是走來走去不斷地不停地做很多事情.
但是你不知道他在做什麼.
意思是做完之後.
其實好像做了和不做是沒什麼大分別.
對於身邊的人.
對於整件事情是沒有任何的影響的.
你明白我說什麼嗎.
這個叫做無事忙.
走來走去做做了又好像沒做一樣.
見到人很多事情.
每件事都好像想幫忙.
在幫你幫你.
但其實跟著都好像沒什麼分別.
所以懶惰本身不是說你做不做.
做多少事情.
而是你那個召命有沒有去實踐.
譬如我作一個例子.
一個學生.
在學生讀書的時期.
他的召命當然就是努力讀好書.
在學習的階段好好地去學習.
以致學到他所應該要學的那些技能.
但如果在那個學生的階段.
他是很努力做很多事情.
做很多可能課外活動.
很多其他的幫同學手.
或者學校有什麼慶典的時候.

$^{241}$很努力地佈置.
你說他是不是一個勤力的人呢.
某個程度起碼他沒有懶惰過.
他做很多事情.
但是在那一刻.
他的字有沒有完成他的召命呢.
其實可能就沒有了.
他搞了很多課外活動.
原來讀書的時間沒有讀書.
他沒有完成他的召命.
這個例子本身你就明白到究竟為什麼.
聖經看我們一個人.
是不是懶惰呢.
其實不是單說.
有時候聖經好像要我們不斷地做.
不斷地做,做到不停.
為之不懶惰,其實不是這樣看的.
是看我們每一天.
我們明白到神所交託給你的.
你有沒有積極地去回應去做呢.
還是反而是將神所給你的.
你放下不做.
即使你做了其他的事情.
這樣也是沒有意思的.
這個就是第一點.
神給你的召命.
你有沒有積極地去回應呢.
在你人生的崗位當中.
我相信其實你說.
我不知道自己的召命.
所以你要想一下.
其實多或少你都知道的.
是不是.
即使你說不是啊.
我現在退休了.
我的召命就是享受世界.
那你說這是不是你的召命呢.
當然你要自己在神面前去反省.
但是我相信這就是第一點.
說到懶惰的時候.

$^{281}$你要明白自己每個人人生.
其實你活在世上都有一個召命的.
在你現在這個階段的召命.
可能跟你二三十年前的階段是很不同的.
是不是.
你不能夠用你二三十年.
我已經完成了二三十年的召命.
跟我現在就沒有召命.
這個完全是說不過去的.
是不是.
所以.
這裡就是第一點.
就是懶惰其實就是對於召命那種拒絕.
我現在不做神所交給我的那些.
第二呢.
就是對於屬靈的一種的冷漠.
對於屬靈的事情的一種的冷漠.
在啟示錄那裡有個老弟家教會就很出名的了.
這句話.
我知道你的行為.
你也不冷也不熱.
我不不得你滑冷滑熱.
你既如溫水也不冷也不熱.
所以我必從我口中把你吐出去.
這裡就說到.
老弟家教會處於一種是對於屬靈的事情.
對於神的事情.
一種不冷不熱的一個狀態.
對於屬靈的事情一個不冷不熱的狀態.
說你是不是熱心呢.
你自己也說不出熱心.
說你是不是懶惰呢.
那你又覺得我應該也不懶吧.
我每個禮拜起碼也有回來.
有些人也不回來.
是不是.
主一學我也有想過上的.
不過我不是每次都能上.
有時候我們就是用一個這樣的狀況.
當中來到很滿足.

$^{321}$但是聖經那裡主耶穌就很不滿意.
以我們來說我們也覺得應該滿意了.
不冷就收貨.
但是主耶穌不滿意.
主耶穌說.
我巴不得你或冷或熱.
你既如溫水.
所以我必從我口中把你吐出去.
意思是什麼.
意思就是說.
主耶穌看到你想作嘔.
是不是.
是這樣的意思.
你想想你看到一個人.
你兜口兜面跟他說.
我看到你也是想作嘔.
這個是很不客氣的.
如果你跟你的朋友.
跟你熟人說了這樣的話.
可能他氣你一輩子.
然後朋友也沒得做.
但是這個就是主耶穌用來形容.
他怎樣去看老弟加教會的情況.
所以第二種的冷度就是.
說到屬靈的一種的冷漠.
這種的屬靈這種的冷漠.
本身就是.
我想是發自內心的.
是發自內心的.
很多事情已經算了.
不用這麼熱心.
求求其其.
或者做一些吧.
不用做這麼多的.
那種心態你明白我說什麼了.
做一些吧.
不用做這麼多的.
都有做過.
有得交差就算了.
就是這樣.

$^{361}$這個是第二種的冷度.
第三種的冷度.
就是對於善的一種厭煩.
所謂聖加拉太書所說.
你們行善不可喪志.
若不灰心.
到時候就必要修成.
後面在貼壽羅尼加後書所說.
弟兄們你們行善不可喪志.
這裡兩句都說到行善不可喪志.
意思就是說.
你做一些對的事情.
做一些認為正確的事情.
做一些好的事情的時候.
到了一個階段.
你就覺得.
唉呀.
都不想再做了.
覺得做夠了.
這種喪志就是這樣的意思.
就是覺得夠了.
做了這麼多.
都應該夠交數了.
就不想再做了.
不想再堅持了.
這個就是.
在聖經當中.
今天我們有思想.
這個冷度.
這個題目的三個方向.
我相信.
所以你搞清楚.
究竟在說什麼.
是很重要的.
因為不只是說.
你做多少事情.
不做多少事情.
我想每個人都需要.
有作息.
有平衡.

$^{401}$其中一點就是.
有些人覺得很累.
因為可能做得太多.
做得太多都會.
物極必反的時候.
就變成什麼都不想做了.
是不是.
所以根本就不是那個問題.
而是你要反省究竟.
聖經那裡.
神怎麼去提醒我們.
怎麼去看冷不冷度.
第一就是照命.
是不是.
第二就是.
你那個.
是對於屬靈的事情的冷漠.
第三就是.
你對於善的事情有沒有厭煩.
所以如果是這樣的時候.
其實你看回.
聖經那裡.
就有那個圖畫.
就有幾幅圖畫.
說到那些懶惰人是怎麼樣.
第一幅圖畫.
剛才已經讀過了.
就在箴言那裡說.
那些的.
他說有些人經常都想睡覺.
經常都想睡.
要睡到什麼時候呢.
你抱著手躺的片時.
打論片時.
讓我多休息一會兒.
這就是那幅圖畫.
叫醒的時候.
多睡一會兒.
你的生命就會很貧乏.
他就這麼說.

$^{441}$這就是在箴言所說的.
那個懶惰人的圖畫.
所以在第一個圖畫當中.
它反映了一個.
很有意思的一個心態.
就是.
他不是說不做.
他是怎麼樣.
他說讓我多休息一會兒.
是不是.
他說多休息一會兒.
所以箴言就諷刺說.
你要休息多久.
要睡到什麼時候.
要欺負到什麼時候.
坐著翹著手翹到什麼時候.
其實就是這樣的意思.
這是一個很圖畫式的表達.
而這種表達.
他的意思就是說.
描述一個在懶惰的狀況之中.
就是一種想拖延的.
就是一種很想拖延.
很想不是即時發生就好了.
或者不需要即時做就好了.
讓我拖延一下.
其實這種拖延的心態.
我們都會明白的.
有些事情你不想做的時候.
你就想拖延.
是不是.
主要為什麼我們想拖延.
因為那些事情你不想做.
那些事情你很想做的.
什麼人叫.
你急不及待等著做.
就是因為你不想做.
不想做就想拖延.
所以第一個箴言描述了.
這種懶惰的心態.

$^{481}$就是如果你對應剛才那三個.
懶惰的定義的時候.
其實會不會是你對於神的召命的一種抗拒.
對於神的召命的一種抗拒.
因為你就是不想做神所要你交託給你做的.
譬如舉例來說.
有個人現在上班.
他很不喜歡這份工作的.
但是這個是唯一的.
找到肯.
他可以錄用他的工作.
那天天就睡到最後.
上班就說到差不多遲到.
因為他根本是很不喜歡這份工作的.
是不是.
所以第一個就是那個.
對於神的召命.
我們心裡面是.
用什麼態度來回應的呢.
有沒有反映過.
其實有時候你的懶惰.
其實出於你心裡面的不滿.
譬如我不甘心做這件事.
有些人是這樣的.
就是.
他在公司的態度很差.
很懶,什麼都不做的.
那你跟他聊聊天.
你發覺其實他也很有能力.
很能夠做得到這份工作.
但是為什麼會有這樣的態度呢.
你跟他聊聊天的時候.
你發覺其中有一種人是這樣的.
他覺得是怎樣的.
他覺得我這麼聰明.
為什麼要做這麼差的職位呢.
浪費了我的時間.
我就隨便做.
我明明應該是高兩級的.
現在就要我做這麼低莊的事.

$^{521}$你明白嗎.
這是其中一種.
但是如果將他轉換.
去到屬靈那裡.
也會有的.
會不會是呢.
剛才所說的其中一個懶惰.
就是那個召命.
他說我這麼有恩賜.
我這麼有才華.
神可不可以用我厲害一點呢.
讓我去一個更大的事奉崗位.
讓我更有更多的影響力.
可以去做更偉大的事情.
如果是這樣的話.
你很難做得好前面的事情.
我們中國人有一句話.
這句話是很貼切的.
很應軍的.
可能你也聽過.
第一個字就是好字.
這是四個字.
你知不知道.
好高無遠.
大家都聽過了.
好高無遠.
其實就是這樣.
你不斷地想.
你自己有多厲害.
做到多厲害的事情.
你想像.
但其實就像這裡所說的.
拖著拖著.
面前的事情拖著拖著.
然後又不做好.
你的貧窮就會來到.
其實就是這個.
這是在《箴言》的第一個圖畫.
就是這樣.
就說到一個懶惰的人的情況.

$^{561}$第二種的情況.
就是馬太福音說到天國的比喻.
天國說.
主啊 你交給我的五千銀子.
請看 我又賺了五千.
主人說 好.
你這又良善又忠心的僕人.
你在不多的事上又忠心.
我要把許多事派你管理.
請看 我又賺了五千.
主人說 好.
你這又良善又忠心的僕人.
你在不多的事上又忠心.
我要把許多事派你管理.
請看 我又賺了五千.
主人說 好.
你這又良善又忠心的僕人.
我要把許多事派你管理.
請看 我又賺了五千.
主人說 好.
你這又良善又忠心的僕人.
你在不多的事上又忠心.
我要把許多事派你管理.
請看 我又賺了五千.
主人說 好.
你這又良善又忠心的僕人.
你在不多的事上又忠心.
我要把許多事派你管理.
請看 我又賺了五千.
主人說 好.
你這又良善又忠心的僕人.
我要把許多事派你管理.
請看 我又賺了五千.
主人說 好.
你這又良善又忠心的僕人.
我要把許多事派你管理.
請看 我又賺了五千.
主人說 好.
你這又良善又忠心的僕人.
我要把許多事派你管理.

$^{601}$請看 我又賺了五千.
主人說 好.
你這又良善又忠心的僕人.
我要把許多事派你管理.
請看 我又賺了五千.
主人說 好.
你這又良善又忠心的僕人.
我要把許多事派你管理.
請看 我又賺了五千.
主人說 好.
你這又良善又忠心的僕人.
我要把許多事派你管理.
請看 我又賺了五千.
主人說 好.
你這又良善又忠心的僕人.
我要把許多事派你管理.
請看 我又賺了五千.
主人說 好.
你這又良善又忠心的僕人.
我要把許多事派你管理.
請看 我又賺了五千.
主人說 好.
你這又良善又忠心的僕人.
我要把許多事派你管理.
請看 我又賺了五千.
主人說 好.
你這又良善又忠心的僕人.
我要把許多事派你管理.
請看 我又賺了五千.
主人說 好.
你這又良善又忠心的僕人.
我要把許多事派你管理.
請看 我又賺了五千.
主人說 好.
你這又良善又忠心的僕人.
我要把許多事派你管理.
請看 我又賺了五千.
主人說 好.
你這又良善又忠心的僕人.
我要把許多事派你管理.

$^{641}$請看 我又賺了五千.
主人說 好.
你這又良善又忠心的僕人.
我要把許多事派你管理.
請看 我又賺了五千.
主人說 好.
你這又良善又忠心的僕人.
我要把許多事派你管理.
請看 我又賺了五千.
主人說 好.
你這又良善又忠心的僕人.
我要把許多事派你管理.
請看 我又賺了五千.
主人說 好.
你這又良善又忠心的僕人.
我要把許多事派你管理.
請看 我又賺了五千.
主人說 好.
你這又良善又忠心的僕人.
我要把許多事派你管理.
請看 我又賺了五千.
主人說 好.
你這又良善又忠心的僕人.
主人說 好.
你這又良善又忠心的僕人.
我要把許多事派你管理.
請看 我又賺了五千.
主人說 好.
你這又良善又忠心的僕人.
我要把許多事派你管理.
請看 我又賺了五千.
主人說 好.
你這又良善又忠心的僕人.
我要把許多事派你管理.
請看 我又賺了五千.
主人說 好.
你這又良善又忠心的僕人.
我要把許多事派你管理.
請看 我又賺了五千.
主人說 好.

$^{681}$你這又良善又忠心的僕人.
我要把許多事派你管理.
請看 我又賺了五千.
主人說 好.
你這又良善又忠心的僕人.
我要把許多事派你管理.
請看 我又賺了五千.
主人說 好.
你這又良善又忠心的僕人.
我要把許多事派你管理.
請看 我又賺了五千.
主人是神.
因為你代替了故事.
主人就是神.
那些僕人就是我們.
那些錢是什麼呢?神給了我們什麼呢?.
恩賜.
神給我們的財,能力,恩賜.
就是我們的本錢.
就是每個人.
天生我財必有用.
就是這個意思.
就是神.
都給了一些本錢.
有些人.
神給多一些.
是要他投資多一些.
你不要以為.
我真的厲害.
天生異品.
天才一個.
如果你真的這樣想.
我都恭喜你.
因為你天才一個.
神將來要你多一些.
你真的做天才這麼多事.
否則就很大件事.
這個就不說了.
一千那個.
前面那兩個人.

$^{721}$剛才說的那個是襯托.
但在這個襯托當中.
都反映了一件事.
他們就拿他們的人生.
每個人都有一個人生.
然後將他的人生.
來投資.
你這樣來想像.
你的人生.
神給了你一些東西.
給了你的能力.
然後就投資.
投資.
最終我們都會回到神面前交賬.
你的人生的投資.
投資成怎樣呢?.
前面那兩個都投資得很好.
神都很滿意.
他們的投資.
最後的人生是怎樣呢?.
就是沒有投資.
他的人生沒有投資.
就.
然後神就問他.
你為什麼不投資?.
是不是?.
你什麼都不做.
將那些錢埋在地上.
你起碼都放在銀行修飾.
意思就是這樣.
問他不投資.
然後.
他那個不投資的人.
怎樣回答?.
他就.
想了一些東西.
令人更加氣憤.
他說.
主啊,我知道你是忍心的人.
沒有中的地方要收割.

$^{761}$沒有散的地方要聚斂.
我就害怕.
去把你的一千銀子埋藏在地裡.
你看你的元銀子在這裡.
這裡你看到.
多麼令人氣憤.
他的答案.
簡單一句來說.
就是.
我也不是不想投資.
不過.
有什麼問題?.
問題不在我這裡.
是在你這裡.
你這麼兇.
我怕投資失敗.
就慘了.
所以不關我的事.
是因為你的問題.
所以我不投資.
這個就是.
這個就是他的原話.
不是我作出來的.
是不是?.
如果對應.
剛才我們所說.
這個比喻本身.
所說的話.
會不會有發生呢?.
你的人生有沒有投資.
在神想要你投資的地方.
以致神將來要去收帳.
和你結帳的時候.
神就覺得可以.
神就稱譽你.
又良善又忠心的僕人.
你在你所做的事上.
雖然不多.
但是你有忠心.
是不是?.

$^{801}$反過來說.
我們有沒有呢?.
我們的人生有沒有投資呢?.
你們的人生有沒有投資呢?.
投資在神想要你投資的地方.
有些人.
耶穌有說過另外一個比喻.
有人拿了錢,然後怎樣呢?.
花天酒地.
那個是什麼比喻呢?.
就是浪子比喻.
其實都是一樣的.
大家都是要拿了老爸的錢.
不是,一個拿老爸的錢.
一個拿主人的錢.
其實都是一樣的.
花天酒地那個就是第二個比喻.
我就不在這裡說了.
這個就不是花天酒地.
這個就說你的人生.
有沒有投資在神想要你投資的地方呢?.
可能.
神想你投資.
你的人生.
不是你現在在做的事.
你有沒有想過?.
因為我們今天講懶惰.
懶惰其中一個就是回應神的召命.
懶惰其中一個就是對屬靈的事不冷不熱.
懶惰其中一個就是對善事已經不堅持去做.
那會不會我們的人生.
就是剛才講懶惰三方面的東西.
譬如神的召命.
神明明要你在另外一個崗位再做.
我聽過有一個.
這個是一個挺有名的.
不是聽的.
實際上是一個挺有名的傳道人.
他就是.
在年輕的時候.

$^{841}$他有兩個朋友.
一起都有機會.
當時在中國文革之前.
他就出來了.
然後就來到美國讀書.
讀了書之後.
之後就.
兩個都覺得有呼召.
覺得神既然給了他們這麼好的機會的時候.
想去給神用.
但是後來有一個.
當然後來就做了傳道人.
後來也是在神學教育那裡.
其實也做了很多事情.
但是另外一個.
跟他一起出來讀書的朋友.
後來就沒有做到傳道.
就在美國那裡.
他就讀數學.
後來在大學那裡教書.
在那個朋友的見證會那裡.
他就說.
我們兩個一起在中國那裡出來.
兩個都有機會.
兩個都覺得看到.
神都給我們這樣的機會的時候.
都想繼續去.
想想怎樣去為神.
為天國做點事情.
但是呢.
他後來就退休了.
後來就退休了.
後來就做了一個大學的教授.
就沒有讀到神學.
他在那個見證會那裡說.
相反他那個朋友.
怎樣給神.
在他一生裡不斷地給人使用.
給神使用.
這個是真事.

$^{881}$不過我不說名字的意思.
就是在洛杉磯.
後來那個傳道人死了也是在洛杉磯.
但是.
你就看到.
有時候.
神是給了你一些的召命的.
你有沒有去滿足呢.
有沒有去做好呢.
還是你選擇了一些其他的原因.
你就放下了你的召命.
而你放下的原因.
其中一方面.
好像這個僕人.
其中有一個說法.
都很出人意表.
他不是想想他自己.
他是想想神.
他出人意表.
說神是一個忍心的神.
如果投資失敗了.
就會很嚴重.
廣東話說得很白的.
就是我這樣說.
但是意思其實是一樣.
神是一個很嚴苛的神.
如果我做錯了決定.
我的人生就會很糟糕.
聽得明白了嗎.
很多時候.
很多人對於.
他不是不想做神的事.
而是整天想著一個方向.
就是說如果我做錯了.
神會怎樣罰我.
有沒有人這樣想呢.
可能有可能沒有.
你自己才知道.
有沒有想過.
如果神要你做一件事的時候.

$^{921}$你放膽可以這麼說.
或者你盡力去做的時候.
你不需要想神怎樣罰你.
或者你不需要想.
你的結果會很糟糕.
你有沒有這樣想.
反過來其實真的很多人是這樣.
我都聽過有一個弟兄說.
我都覺得如果可以讀神學.
然後可以事奉神就很好了.
但是我覺得還是不行.
為什麼呢.
因為如果我現在放下.
我現在手上的東西.
奮身去跟從神的時候.
如果萬一我又.
做得不是那麼好的時候.
我兩頭不到 會很大件事.
然後收入又不夠.
可能生活又有問題.
那就很慘.
所以到頭來.
從始至終都是在原地.
沒有走出來跟從神事奉神.
那個想法其實都是很類似的.
是不是.
我不敢奮身去跟神.
因為如果奮身出去.
如果失敗了的時候.
我就很大件事.
他不做其實也很大件事.
在這個故事裡面說.
只不過這部份的人沒有這樣想.
你沒有想過你不回應神.
也是很大件事.
但是人就是這樣想.
我回應神.
然後有什麼出錯呢.
你回應神.
做了什麼大全道.

$^{961}$做了什麼.
引領豐功偉績.
那當然好.
但這個世界不是那麼多.
有很多窮全道.
生活又很困難.
子女又叛逆.
那就很大件事.
不如找一份固定工作算了.
就是從這個角度來想.
沒有想到另一方面不跟神.
就更大件事.
在這個故事的其中一個難度.
第二方面.
就是你不想自己的難度.
而是想神的問題.
第三個就是不冷不熱的.
就是《老弟家教會》.
剛才也讀過.
主耶穌寫信給老弟家教會.
"那為雅門的,為誠信真實見證的,在神創造物物之上為元首的書.我知道你的行為,你不冷不熱,我不不得你或冷或熱,你寄予溫水,我就將你吐出去".
這一段.
這裡其實就是另一個圖畫.
就是對於屬靈的事情的冷淡.
有一種冷淡.
有些人呢.
很有趣的.
他做很多事情都很熱心.
見到朋友有問題.
又很熱心幫忙.
見到哪裡有事情.
又很熱心去找資料.
每件事都很熱心.
整個人都是很有熱心的人.
但是呢.
對於屬靈的事情呢.
就很冷淡.
就是.
不知道你有沒有遇過.
我也有遇過.

$^{1001}$你見到每件事.
有些需要幫助.
又很熱心幫助.
見到有些興趣的東西.
又很努力去找資料.
好像.
每件事都是很有動力的.
但是對於屬靈的事情呢.
就很沒動力.
讀聖經,讀兩讀.
然後覺得很悶,就放下不讀.
接著我就問.
為什麼你又不讀.
接著他說,讀兩讀很悶.
就沒興趣了,就放下.
接著我就說.
他就說.
讀兩讀,我不明白.
開始就覺得悶.
就放下不讀.
接著我就問他.
咦,這個不像你.
你平時工作.
你很熱心的.
你不明白的時候,你到處去找資料.
為什麼會做好.
讀聖經的時候,讀兩讀.
就馬上放下,不會繼續下一步.
你明白我說什麼嗎.
完全是兩個標準.
對於屬靈的事情.
和對於屬世的事情.
他叫你.
找一些資料.
搜集.
怎樣炒股票.
那些人很熱心,很厲害.
但是叫你資料搜集一下.
給你第一次.
主一學.

$^{1041}$叫你資料搜集.
去帶領主一學.
帶一季都不行.
哇,不行.
這個要做很多搜集功夫.
不行.
這樣.
你看到很明顯.
有兩套不同的標準.
這個就是不冷不熱.
那個.
反映的現象.
而在今天來說.
所謂的懶惰.
是有很多重的表現.
第一個就是.
之前有一堂課也說過.
就是分心的那種懶惰.
現在根本不能夠專心.
對於屬靈的事.
對於回應神的事.
對於神所要交託你.
老老實實去做一些事情.
就好像八足那樣.
太分心.
尤其是一個電話就已經足夠你分心.
第二就是情感上的懶惰.
這種就最直接.
就是很累.
什麼都不想動.
不想開新項目.
不想有什麼.
不想又有新的東西.
來了.
第三種就是屬靈的懶惰.
屬靈的懶惰.
你也聽過了.
剛才也說了一些.
對於屬靈.
又要去跟從神.

$^{1081}$又要事奉神.
又要去.
有些屬靈的.
活動要參加.
有些事奉要參與.
有些屬靈的操練.
現在要練習.
很麻煩.
照明的剛才也說了很多.
不說了.
道德上的.
哎呀.
這麼麻煩.
舉例來說.
我做個例子.
說扔垃圾.
有些人為何明明.
前面街口有個垃圾桶.
又不扔就扔在街上.
有沒有見過.
因為很麻煩.
你做對的事.
又要麻煩.
你做一些.
隨便就算了.
因為你做對的事.
是會麻煩的.
所以這是道德上的懶惰.
我.
做個例子.
舉例來說.
你去投票.
如果你認為.
投票是一個公民責任.
你也不可以亂投.
是不是.
你也要研究.
什麼對.
什麼不對才投.
但你說很麻煩.

$^{1121}$所以最後.
我也不去投票.
這是道德上的懶惰.
你有沒有想過.
投票是公民責任.
你要參與社會的時候.
很麻煩.
就不投.
這是道德上的懶惰.
說了這麼多原因.
其實.
最終.
我們.
想想自己的內心.
第一個.
就是失序的愛.
失序的愛.
其實.
懶惰不是.
一個做不做的問題.
而是你愛什麼.
你愛什麼.
你愛的東西.
與神創造你的時候.
神給你的愛.
是不同的.
就是這個原因.
比如.
你是.
在那個情況下.
神要你.
很.
很努力的.
學習的.
舉例來說.
有些階段的學生.
在那個時候.
他應該要愛慕.
學習愛慕知識.
但如果反過來說.

$^{1161}$在那個時候.
他愛了別的東西.
比如.
我舉個例子.
以前我們讀書的時候.
父母跟我們說.
不要拍拖.
他愛了別的東西.
愛了別的女孩.
這樣就分心了.
其實就是.
你說他是不是懶惰呢.
其實這個不單止是懶惰的原因.
讀書,然後他又不讀書.
又走課,不是懶惰的原因.
而是他愛錯了別的東西.
在那個時候他應該愛慕學習.
愛慕知識,然後愛慕了.
可能拍拖那些東西.
其實很多時候.
有時候我們的人生是這樣.
我們愛錯了另外一些東西.
以致到.
表面上你說.
我只是不想動,但其實不是不想動.
特別淑齡.
為什麼淑齡.
叫你去旅遊很開心.
叫你去退休會不開心呢.
就是說.
叫你去.
吃自助餐很開心.
但是叫你去.
參加培靈會不開心呢.
因為我們愛的是.
物質的東西.
而淑齡的東西就不愛.
所以人生就是.
那個愛慕的秩序.
第一個是搞錯了.

$^{1201}$第二就是恐懼.
剛才那個第二個.
比喻那個booking的恐懼.
就是怕失敗的恐懼.
或者怕.
如果我不做就不錯.
不做就不錯.
少做少錯,不做就不錯.
你如果以這樣的態度.
一做就會錯.
一做就會.
如果被人知道.
我失敗了.
很沒面子的時候.
那你就什麼都不做.
有個姐妹說.
她讀神學.
有朋友跟她說.
你真的很大膽,你去讀神學.
其實事實上是的.
你想想.
你高了差不多一年多.
很久沒有回學校.
又要去看書.
又要做功課又要寫文.
隨時都.
搞到自己.
很累.
很不適應.
你想到各種.
令你恐懼的原因.
令你恐懼不讀的原因.
是會有的.
但其實第二種就是恐懼.
你怕自己失敗.
你怕自己.
搞不定.
你怕自己不適應.
但你知道其實要這樣做.
第三種.

$^{1241}$就是將安逸變成偶像.
總之就是很舒服.
不要動就最好.
是不是.
很舒服,不要動就最好.
之前.
可能你都聽人.
都有說過.
我們每個人.
都有一些.
安逸圈.
comfort zone.
很舒適的環境.
你走出去就會不舒服.
可能你都聽過.
但是.
我最近又反過來想.
其實我們是不是有.
舒適圈.
什麼是舒適圈.
什麼是安逸圈.
有沒有想過.
作為一個人.
怎樣才算是安逸.
那是不是.
你現在.
心理.
和物質方面.
很滿足就算是.
舒適圈.
很舒適的圈.
那個是不是真正的舒適.
有沒有想過這個問題.
那個是不是真正的舒適.
其實可能不舒適.
你心裡面.
有些不安.
覺得很不滿足.
可能你有一種.
覺得.

$^{1281}$對神有一種虧欠.
其實那些都不舒適.
是不是安逸圈就真的那麼安逸.
你反過來想.
其實只不過.
你盡量不去想那些不安逸的東西.
但反過來說.
你走了出來的時候.
是不是又是那麼不安逸.
又未必是的.
你走了出來跟從神的時候.
我自己的經歷來說.
我都不覺得我很不安逸.
就是.
神也不會說你跟從他的時候.
一定要.
天長江大任於.
必先傲其筋骨.
傲其體膚.
這樣對我的.
你明不明.
就是.
哪樣是安逸哪樣是不安逸.
根本不是這樣的.
只不過我們是神化了.
不動就是安逸.
但其實你不動之下.
你那種懶惰.
那種懶惰不是不做事.
那種召命那種回應.
那種有時候的內疚.
那種人生.
看人生那種虛空.
你一樣是面對的.
就是這樣我們也失去了永恆的眼光.
就是其中一個懶惰的原因.
不是說你不做事.
而是你做的事是沒有永恆的眼光的.
是不是.
你沒有永恆的眼光.

$^{1321}$不是代表你不做事.
你很努力.
對於世界物質的東西很努力.
但問題是.
失去了永恆的眼光.
那.
怎樣去面對呢.
首先我們要誠實的.
從一點點的馴服開始.
你要克服.
譬如你說我現在對於屬靈的懶惰.
我的懶惰.
譬如就由靈修開始.
那你就由.
每天很少的決定.
有一個時間.
來到靈修開始.
你現在一定是從很少的開始.
未必是需要想得很厲害.
你說.
你覺得神現在給我的照明.
譬如你做媽媽的.
你現在神給你的照明.
是一個教好你的小朋友的.
那你教好你的小朋友的照明.
其中一方面.
你小朋友怎樣會在.
成長的過程當中.
做一個good person.
一個好的人.
不單止是.
他的物質方面豐富.
他在品格方面.
他在屬靈方面.
他在信仰方面.
也是紮實的.
這樣的時候.
你要想怎樣在這個小的.
小小的馴服來幫助你.
完成你的照明.

$^{1361}$是不是.
第二種.
就是用敬拜來.
取代了.
內在的自我消耗.
一開始說.
那個躺平.
和老鼠人的例子.
現代人.
其實有兩種.
有一種.
是非常努力積極.
要去.
好像要去.
好像這個世界要去競爭.
但有另外一種.
就覺得競爭.
都沒有用.
躺平了.
因為.
大家都很內耗.
兩種其實都是同一個原因.
有兩個不同的反應.
同一個原因.
就是失去了.
內在真正的.
那種安息.
一個內在很不安息的人.
可能唯一他能夠感覺到.
有一點點踏實的.
就是賺錢.
他很努力的賺錢.
或者有一個人心裡面.
也很失去了安息的時候.
他覺得努力都沒有用的時候.
就躺平.
其實都是一樣.
失去了內在真正的安息.
而這種安息就是.
要在神那裡而來.

$^{1401}$要是敬拜.
而當我們.
是用安息.
這個內耗的時候.
我們就會感受到.
我們是被神接納的.
一個懶惰的人.
是覺得.
算了.
做也好,不做也好.
好像已經放棄了自己.
但是你感受到在神面前有安息的時候.
你的人生.
你是不覺得需要算.
因為.
神是接納我.
我會照命.
而.
還有.
還有兩點.
所以我們有一種.
生活可以定一些秩序.
怎樣去.
對付你的懶惰.
你定一些秩序.
譬如哪個時間.
是做什麼.
這個時間是給神的.
這個時間是給自己休息的.
這個時間是.
要專心不要分心.
去有一些.
生活上的調整.
懶惰本身.
不是只是做事不做事.
而是你生活有沒有一些界線.
因為你有這樣的界線.
你才懂得回應神.
而.
也是.

$^{1441}$在這樣之下.
我們就.
懂得去分.
我們是.
有些休息是.
逃避性的休息.
有些休息是更新性的休息.
什麼叫逃避性的休息.
就是根本有些.
你什麼都不想動.
什麼都不想做.
只是想要逃避.
但是有些逃避性的休息.
休息完之後.
你仍然是很累的.
你只是躺著也是很累的.
然後再躺也不夠.
這種就是逃避性的休息.
因為你根本不知道自己想怎樣.
你根本不知道要去哪裡.
但是有另外一種.
就是你需要休息.
七天有一天休息.
在神面前.
重新定你的.
把你的生命好像有一個.
有一個更新之後.
你又知道繼續想.
繼續上路.
你.
去.
就是你去.
做運動的時候或者你行山的時候.
你有時候走到累.
你要休息.
但是那種休息完之後.
你精神.
或者你體力又回來了的時候.
你就會繼續.
繼續這樣走.

$^{1481}$因為你知道有一個目的地.
你不會在那個休息.
本身就停在那裡.
但是前面那種逃避性的休息.
是沒有目的的.
就是什麼都不想做,什麼都不想動.
所以求神.
幫助我們,如果我們是這樣的時候.
常常都是.
能夠記得神的託付的時候.
我們的人生.
就不是懶惰的人生.
正如一開始都講過.
整堂都講了很多.
那種懶惰不懶惰.
不是說做事多少.
而是我們對於.
神的託付.
對於我們人生當中.
我們怎樣來回應.
求神幫助我們.
\newpage



\section{出埃及記 3:15}
\label{sec:031901}
\textbf{聖經人聖經事:耶和華是我的名}~~~~黃克勤牧師博士
\newline
\newline
連結: \href{http://www.church.com.hk/acms/content.asp?site=cdc\&op=show\&type=product\&code=031901\&layout=sermon}{\texttt{031901}} ~~~~ 語音日期: 2026-03-30
\newline
\newline
\hyperref[sec:031900]{\small{< < < PREV SERMON < < <}}
~
\hyperref[sec:index_chronic]{\small{[返順時目]}}
~
\hyperref[sec:031902]{\small{> > > NEXT SERMON > > >}}
\newline
\newline
出埃及記 3:15
\newline
\begin{longtable}{cl}
\hline
\hline
章節 & 經文 (和合本修訂版)\\
\hline
3:15 & \begin{tabularx}{0.7\textwidth}{X} 神又對摩西說：「你要對以色列人這樣說：『耶和華－你們祖宗的神，就是亞伯拉罕的神，以撒的神，雅各的神差派我到你們這裡來。』這是我的名，直到永遠；這也是我的稱號，直到萬代。 \end{tabularx} \\ \\
[1ex]
\hline
\hline
\end{longtable}
$^{1}$各位弟兄姊妹.
又一天了.
眨眼就二十四小時了.
又來讀聖經了.
生命樹宣教網絡.
為大家提供聖經人,聖經事.
看看在聖經的人物和事情的背後.
上帝做些什麼.
認識祂更加多.
特別今天這一集.
真的認識神很多了.
因為簡直是耶和華是我的名.
這個話是誰說的呢?.
讀聖經多些也會知道耶和華是我的名.
當然是耶和華說的.
是神介紹祂自己.
這裡有很多不同的解釋.
甚至有神學爭議.
今天就照我知道的字義.
跟大家嘗試去解釋一下.
仍然是在說摩西的故事.
大家記得摩西在抗疫那裡.
看到荊棘被火燒著.
其實在那裡荊棘被火燒著是經常發生的.
因為在抗疫.
熱到什麼程度.
我都說我們去到耶利哥.
熱到50度C.
即是50度攝氏.
可能不止.
只不過我拿它的溫度計已經去到盡了.
很熱.
又很乾燥.
所以如果有些植物,荊棘.
乾了之後自己著火.
是一點也不奇怪的.
摩西覺得奇怪.
只不過是著火.
有火.
為什麼燒這麼久還在那裡.

$^{41}$為什麼不會熄滅.
不會有煙.
還是火.
他覺得很奇怪.
看了很久.
真的這樣.
他才過去看看發生什麼事.
摩西也很聰明.
他當然是一個異象.
他過去看.
於是神在那裡.
在火裡面呼召他.
過去的時候.
叫他跟我帶以色列人出去.
他們的哀聲我聽到.
很慘啊.
在那裡做奴隸.
來到這裡.
《昭合紀》第三章第十四節.
如果你是摩西.
摩西就說.
叫我帶以色列人出去.
怎麼帶.
摩西說不行.
沒有可能.
原因就是.
第一.
現在的法老.
其實是我以前一輩的兄弟.
跟我們一起上學.
我們也是上王子學院.
我認識他.
而他是想殺我.
因為我殺了埃及人.
所以我不會回去叫他.
來放我們以色列人走.
而且上帝你知不知道.
這是我作的.
上帝你知不知道.
過百萬的人在那裡.

$^{81}$我們的同胞是免費奴工.
奴隸.
怎麼會放我們走.
如果我回去.
我是不會去見到他.
這個也對.
原因就是你去到宮廷門口.
進入王宮門口.
你進去已經被衛兵攔住.
不會進去.
而且怎麼會去到他面前.
沒有可能.
去到他面前.
我就一定沒命.
因為現在回來投案.
所以他知道不行.
上帝說不怕.
我叫他去就行.
第十三節是這樣的.
摩西對神說.
我到以色列人那裡對他們說.
你們祖宗的神打發我到你們這裡來.
如果他們問我.
他叫什麼名字.
我要對他們說什麼呢.
我加了一個綠色的字.
原因就是這個字.
Hene.
在聖經裡面有.
但是中文沒有翻譯出來.
Hene這個字在聖經出現過很多次.
最出名的一節.
就是以賜阿書第九章.
我可以差遣誰.
誰肯為我不去呢.
以賜阿說.
我在這裡請差遣我.
我在這裡這四個字.
其實希伯來文是一個字.
就是Hene.

$^{121}$Hene不是我在這裡.
Hene說過很多次了.
就是「看看」或「嘿」的意思.
很多東西很重,誰幫我拿上車.
很多椅子要收完,崇拜誰幫手.
「嘿」這樣的意思.
舉手「我啊」.
看看這裡.
看看.
「我啊」這樣的意思.
這個字連個「我」字都沒有.
不過中文就是.
「嘿」表示我在這裡.
於是我就在這裡.
在這裡摩西對神說.
「嘿」的意思不是「開」.
是「看看」.
叫我去看看這樣的意思.
看看.
我去到那裡.
他們問我.
你貴姓,我怎樣回答.
如果他們問我.
打發你來的神叫甚麼名.
我怎樣回答.
很有趣.
其實.
這個摩西人是希伯來人.
他們一直下來.
由阿巴拉罕以撒雅各一直下來.
怎會不認識呢.
還有這個問題.
祂叫甚麼名字.
以色列人應該不會這樣問.
不會問甚麼名.
就是阿巴拉罕以撒雅各的神.
祂叫甚麼名.
沒有以色列人會這樣問.
他們全部都認識祂.
所以這句我個人認為.

$^{161}$摩西在太極推手.
不行了.
萬一祂問我叫甚麼名.
你叫甚麼名.
其實就是叫你來的那個貴姓.
何時叫你.
他哪位來的.
真不真.
那時我怎樣回答呢.
要特別介紹這個字.
「看看」.
如果我去.
一定死無生.
不行的.
以色列人都不會跟我.
這個很出名.
《出一級記》三章第十四節.
我A是因為電視幕太小.
寫不到.
太長.
我用大一點的字.
讓大家容易看.
所以分了上半節.
神對摩西說.
竟然問我甚麼名.
我就說我的名字給你聽.
神對摩西說.
「我是自有永有的」.
這個「一」這個螢光.
螢光的字就是我要說的第一點.
「我是自有永有的」.
這個很出名的英文叫.
「I am who I am」.
有些會譯作.
「I am that I am」.
都可以的.
主要是來自這幾個.
希伯來文的字.
「Hiyah Ashiyah Hiyah」.
即是甚麼來的呢.

$^{201}$其實如果將字逐個字拆解翻譯.
「Hiyah」就是「是」.
「是」或者叫「I am」那個「am」.
如果說英文.
我經常去柬埔寨.
柬埔寨的小朋友.
學幾句英文.
他就會問.
「What is your name?」.
問我.
問完我又會問他.
「What is your name?」.
他的答案是.
「My name is …」.
譬如你問我叫甚麼名.
我答「My name is Jerry」.
其實明白的.
英文都可以的.
不過不好的.
為甚麼不好呢.
如果說.
「What is your name?」.
「What is your name?」都不好.
即是很直的.
「May I have your name please?」.
「我可不可以知道你的名字?」.
我不會說「My name is Jerry」.
我會說「I am Jerry」.
或者我直接說「Jerry」.
「I am Jerry」.
這樣就是「I am Jerry」.
於是第一個「am」.
就很有可能是「I am」.
如果在這裡看.
就是一個名字.
以下就是一個名字.
有些人認為這裡不是.
「神呀,你叫甚麼名?」.
很有可能神沒有回答他的名字.
所以這個不是名字.

$^{241}$祂只是回答了.
祂是怎樣的性質.
祂沒有回答祂的名字.
所以這個「am」.
很有可能是名字.
可能不是名字.
譬如中文翻譯.
「自有永有」.
是不是名字呢.
「你姓甚麼?」.
「我姓黃」.
你的名字是「自有永有」.
「黃自有永有」.
是名字.
但其實你見「自有永有」.
不是名字.
是一個性質.
是一個特質.
或者是一種特別的狀態.
於是有些人就說.
這裡有講神的名字.
原因是摩西問祂甚麼名字.
答案當然是應祂的名字.
其實可能不是名字.
有可能是名字.
如果以色列傳統是名字.
真的有點難教.
「Haya」.
第一個「Haya」是「am」.
就是「I am」.
有可能是這樣.
「Shia」是一個連接詞.
連接著上面.
「你是甚麼名字?」.
「I am Jerry」.
「我 Jerry」.
有連接性.
但沒有連接詞.
「Shia」是連接詞.
連接這兩個字.

$^{281}$這兩個字中間連接著.
「Shia」是「那個」.
「我就是那個」.
「那個」.
於是英文叫做「who」.
或者「that」都可以.
「誰昨晚拿了這麼多禮物來?」.
「我就是那個」.
「那個」就是表達.
「who」或者「that」都可以.
於是就是「I am」.
「Shia」就是「I am」.
「I am who」.
這個字又來「Haya」.
「Haya」就是「am」.
或者「be」.
於是這個就是.
「我就是叫做am」.
「你叫甚麼名?」.
「我叫做我是」.
或者「我叫做」.
認不出來.
如果用「be」都可以.
因為這個是verb to be.
「am」是verb to be.
用「be」就表示「存在」.
於是就是「存在」.
用了中文翻譯做「有」.
「存在」就是「有」.
於是「我就是那個我」.
或者「我就是那個存在」.
那麼「自由」「永有」這些概念.
是從哪裡來的呢?.
就是由「am」來的.
「am」是一個verb to be.
而且它是一個簡單的現在式.
一個simple present tense.
這個simple present tense.
就是表示上帝任何時間.
它都是現在式.

$^{321}$而這個「am」是發展出我們.
就是發展出現在我們經常說.
主要是說「你昔在今在以後永在」.
為什麼呢?.
就是說任何時間.
你都是現在式.
那麼「和」了就是過去了.
過去了.
過去式了.
將來式就是未到.
那麼過去了就不再存在.
那一天pass away.
過去了.
將來未出世.
所以兩個都不存在.
只有我現在式.
但是在上帝來說.
我什麼時候都是現在式.
所以世界之初.
我已經在現場了.
我已經在那裡.
將來的世界.
又是我在那裡.
在這個世界之前的.
從耿古到永遠.
我都是現在式.
昔在今在以後永在.
那麼就是由這個「am」發展出來的.
是現在.
發展出來自由榮遊.
沒有人去做了出來.
它都不會消滅的.
那麼這件事很重要的.
很重要的.
大家知道.
近期香港很多人移民.
在歷史上香港1997年.
7月1日就是由英國殖民地統治.
就交回給中國管理.
很多人擔心有事發生就走了.

$^{361}$當時的中國領導人.
鄧小平伯伯就說.
不用走.
因為由97年7月1日開始.
50年不變.
馬照跑舞照跳.
沒問題.
就是50年.
這個資本主義社會制度不變.
就說了.
不用走了.
為什麼還要走.
因為很多人想.
50年不變.
看來當時的鄧伯伯.
50年後都不在了.
說這句話50年後不變的人都不在了.
那就是會變了.
所以每個人都害怕.
哎呀 逃跑了.
當然後來又回來了.
這個.
這個什麼時候都是現在式很重要.
就是說他是不死的神.
所以他說的話是可以的.
不會說完了之後我都死了.
你知道現在都叫做一朝天子一朝臣.
以前是這樣說.
現在不是了.
那就糟糕了.
再看下去.
又說你要對以色列人這樣說.
自由永有打發我到你這裡來.
又再用自由永有.
是那個自由永有的.
你看這個寫法.
這個自由永有M就不是名字了.
是說一個他存在的性質.
那個自由永有的.
就不是那個名叫自由永有的.

$^{401}$就不是那個名叫M的叫你來.
是那個自由永有的.
所以有人認為.
在這裡其實不是一個名字.
不是一個名字.
不過以色列人.
猶太人是將它發展成為一個名字.
再看下去.
好了.
來到第二點.
神又對摩西說你要對以色列人這樣說.
開名了.
耶和華.
第二點我要跟大家說.
你們祖先的神.
就是阿巴拉罕以色列的神.
打發我到你們這裡來.
帶你們走的.
他的名字叫做耶和華.
他剛才又沒有說.
阿Sir你叫什麼名字.
耶和華.
他沒有這樣說.
我是自由永有.
我是M.
我是我是.
我是是.
這樣就說了一個性質多一點.
不過在那裡後來就發展了神的名字.
而耶和華這個字是由那裡發展出來的.
但是這裡才說.
這個耶和華這個字其實.
阿巴拉罕以色列的神.
那時候阿巴拉罕的時候.
會不會已經知道神的名字叫做耶和華呢?.
都知道的.
因為你看到在創世記阿巴拉罕的時候.
就已經寫了神的名字叫做耶和華了.
不奇怪的.
因為你知道.

$^{441}$摩西五經是摩西寫的.
他早就知道耶和華.
他後來才寫回摩西五經.
所以他已經叫了神的名字叫做耶和華了.
所以有這些耶和華的表達在這裡.
就好像新約很多人認為.
去到馬太福音十六章.
才開始耶穌說「你是彼得」.
可能在那裡才主耶穌為彼得叫了一個新的系列名.
以前祂叫做西門.
所以叫西門彼得.
西門是希伯來名.
彼得是系列名.
但是為什麼馬太一寫就寫西門彼得.
很早就寫西門彼得了.
因為祂在以後回憶那些行述的時候.
祂就記得彼得一直被叫做西門彼得.
祂在回憶的時候就寫回這個舊名了.
耶和華很有可能是後來發展出這個由自有永有.
「耶」「耶」出來的名字.
其實我第二點想和大家說.
你要記得我們的神是誰.
是我們祖宗的神.
以色列人我不是.
亞伯拉罕的神.
以色列的神.
都是我們的神.
我們是因為信心而承接成為亞伯拉罕的子孫.
所以你要記得.
在我們來說你都可以有一個信心.
就是我的祖宗.
屬靈的祖宗是誰.
信心的偉人亞伯拉罕.
信服的表達者以撒亞國.
是上帝特別選擇他.
要成就他的國度.
所以沒有亞伯拉罕國.
沒有以撒國.
是神用亞國改名以色列.
就變成以色列國和以色列人了.

$^{481}$你見到三個都有不同的重要性.
而我們都是他們屬靈的後裔.
所以你要記得當你有很多困難.
有很多事情怎麼辦呢?.
你可以記起.
亞伯拉罕以撒亞國祖宗的神.
耶和華會一直帶領你和我的.
因為我們都是屬於他的百姓.
在困苦圍爐之地.
他會帶領我們出來.
他會打發摩西象徵教主.
帶領我們出來.
記得不要忘記自己祖宗的神.
我們的神是耶和華.
亞伯拉罕以撒亞國的神.
我們看看第三點.
耶和華這個字是什麼來的呢?.
就是第十五節.
耶和華是我的名.
直到永遠.
這也是我的紀念.
直到萬代.
第十五節的下半節就結束了.
這句說話很有意思.
上帝說.
耶和華.
這個字就是我的名字.
耶和華是我的名字.
耶和華才是他的名字.
而「自有永有」呢?.
很有可能不是他的名字.
是他表達我的色在今在.
以後永在.
我的永遠不變.
其實如果將這個「am」.
在英文是「m」.
「be」或者「work to be」的「be」.
翻譯成「自有永有」.
表示他的「always」.
「simple present tense」.

$^{521}$「幾時都可以簡單的現在式」.
是對的.
但還有一件事.
由於他「幾時都可以簡單的現在式」.
所以「自有永有」可以.
我歡迎你多加一句.
「自有永有」和「永不改變」.
「有」在這裡不死.
不過變了.
我現在六十多歲.
我未死.
會變了嗎?.
我變了很多.
我年輕的時候.
頭髮不知有多少.
但我變了.
現在沒有了.
當然現在還有的.
我自己看的.
以前見到大太陽.
「打球吧」.
現在見到猛太陽.
「比比吧」.
都是我.
為何會變了?.
那麼多樓梯級.
「不怕,鬥快」.
現在樓梯級.
「有升降機嗎?」.
都是那個人不同了.
我變了.
差了.
所以「自有永有」是指存在.
「永不改變」是指他的能力.
所以只譯成「自有永有」是不夠的.
我個人認為.
應該神學上的「自有永有」.
是「永不改變」.
沒有變.
《雅各書》說.

$^{561}$「他沒有轉動的影移」.
如果不是.
變也死的.
以前我喜歡公義.
現在變了.
我喜歡邪惡.
以前「牧師,你那麼喜歡吃牛扒」.
「買牛扒你吃」.
現在不是.
喜歡吃豬扒.
買光了.
不高興.
買了豬扒來.
我喜歡吃菜.
你服侍不了.
你要是一個.
交往不了.
你要是一個不變.
存在且不變.
這個很重要.
留意這裡.
「耶和華」第三段跟大家解釋一下.
「耶和華」是來自.
通常我們今天看神學的書.
「耶和華」是譯成「YHVH」.
沒有響音怎麼讀.
讀不了.
但是猶太人.
以色列人是讀成「Yahweh」.
或者我們都是這樣.
有些是「Yahweh」.
中文很多時候譯成「阿威」.
「阿威」「阿威」之類.
有些叫「阿危」.
「危」是山字.
下面一個「鬼」字.
很高的山.
「Yahweh」.
讀不了.
「Yahweh」讀得到.

$^{601}$不過要加些響音.
「Yahweh」這個字很有可能是來自.
我們剛才讀過的那個字根.
是「Yah」.
這個字是它的root.
是它的字根.
都會見到「H」和「HY」這些字.
通常猶太人和以色列人.
他們常用兩個字.
其中一個是「Yah」.
「Y-A-H-H」.
這個就是「Y-A-H-Yah」.
「Y-A-H-Yah」.
這個是重一點的「H」.
「Yah」.
大家有沒有.
不是「有沒有」.
一定聽過.
「Hallelujah」.
「Hallelujah」是甚麼來的呢.
其實「Hala」.
「Hala Hala」就是讚美.
名詞「Hala」.
變了做命令式.
讚美是「你要讚美」.
變了命令式.
「你要讚美」.
變了「Hallelujah」.
「Hallelujah」.
「你要讚美」.
中文加了一個「你」字.
比較明顯.
「要讚美」變了命令式.
「讚美」變了「Hala」.
「要讚美」變了「Hallelujah」.
然後「Yah」.
最後的「Yah」是甚麼呢.
「Yah」.
其實就是「耶和華」.
所以「Hallelujah」的意思就是.

$^{641}$「你要讚美耶和華」.
中文的「你要」.
會中文的 imperative 即是命令式.
會用了「當」.
「當讚美耶和華」.
就是命令式.
明白嗎.
「當」不是指應不應該.
應當.
「當讚美耶和華」.
其實意思是 imperative.
命令式.
「你要讚美」.
是誰呢.
讚美「Yah」.
所以「Hallelujah」.
就是這裡了.
所以他們很多時候都說「Yah」.
那這個「YHBHB」是甚麼來的呢.
加了兩個響音.
你才能讀到.
不然沒有響音怎麼讀呢.
譬如.
我們有個同事姓吳.
姓吳英文是 NG.
那些西人.
怎麼讀呢.
因為你沒有響音.
那個字 NG 你發不到音.
很多美國人會將姓吳讀成 ing.
他加一個響音下去.
變成 NG 就讀成 ing.
如果不加響音讀不讀.
那「YHBHB」也讀不到.
所以他們加了個 o 和 i 在這裡.
就變成「YHOVIH」.
這個很多時候會讀到英文的.
Jehovah.
You will.
那就是 Jehovah.

$^{681}$其實這個應該是讀「Y」的.
You.
一讀了「Y」.
你會讀到.
You.
Jehovah.
讀快一點就變成.
Jehovah.
You are.
那就是 Jehovah.
除了說「耶和華」.
就是應對這個「耶和華」.
而英文「耶和華」.
其實如果你用這個來讀.
你讀不到「耶和華」.
不過呢.
讀中文「耶和華」的時候.
是跟著這個底子.
又跟英文底子.
你才讀到「耶和華」.
其實呢.
J 字如果讀「Y」就好很多了.
另外如果是猶太人或者以色列人.
他們會翻起.
Lord 翻譯成.
將「見到耶和華」這個字.
他們不會讀「耶和華」.
不會讀「Yahweh」.
他們不會.
「Yah」也可以.
不過很少.
不敢.
因為他們.
上帝偉大的名.
立約的名.
立約的神的名.
他們不敢讀.
那又怎樣.
他們不會讀「YH」.
讚美「YH」.

$^{721}$「VH」.
他們不會的.
「Yahweh」他們很不敢讀.
那他們就會讀.
The Lord.
「耶和華是我的牧者」.
他們會很多時候.
讀「The Lord is my shepherd」.
他們不會這樣讀「耶和華」.
「The Lord」其實是來自希伯來的「Adonai」.
「Adonai」的意思是主.
他們一見到「耶和華」這個名.
他們就不讀了.
自然就用「Adonai」代替了「The Lord」.
所以他們很多時候.
如果他們讀「The Lord」.
就是在說「耶和華」.
那麼你會知道.
大概是這樣.
在這裡神是說「Yahweh是我的名」.
這裡就是名了.
「Yahweh」是來自什麼呢?.
就是來自正話「Haya」.
「Haya」是什麼呢?.
「Am」.
「I am」.
「Am」.
「自有永有」.
其實「自有永有」是名呢?.
如果這樣看又是名.
又是名.
你叫什麼名?.
我叫做「Am」.
那是什麼意思呢?.
我時時都現在式.
我永遠都不改變.
所以你就用這個來叫我的名.
意思是這樣.
你用這個來叫我的名.
而猶太人,以色列人.

$^{761}$就將它發展成為「耶和華」.
「YHVH」這個字.
這個是我的名.
也是我的紀念.
紀念是什麼呢?.
如果你想紀念我.
你就紀念這個東西.
如果你想記得「Memory」.
譬如我們離開世界了.
無論你什麼狀況.
都找一個石碑寫著你的名字.
弄一張照片.
那個碑文.
那個石碑是用來做什麼呢?.
紀念嘛.
當你想記起它的時候.
就記起那個東西.
所以你記起神的時候.
就記起什麼呢?.
記起「耶和華」.
這是我的名.
是我要來給你做紀念的.
明白嗎?.
意思就是當你記念神的時候.
就記起祂的名字.
用祂的名字來做一個紀念.
也對的.
很多人用了一個地方.
這裡就是何時何時.
我和我爸爸來過這裡.
何時何時我和我爸爸來過這裡吃飯.
何時何時來過這裡玩.
地方也是一個紀念.
不過直接表達一點.
神就說用我的名字紀念我.
用我的名字來紀念我.
不要用地方.
用我的名字就行了.
這樣也好.
因為你不在那個地方就離開了.

$^{801}$哎呀,回不去了.
但是名字永遠都在你腦海裡.
跟隨著你.
這個也很重要的.
好了,今天跟大家分享到這裡.
剛剛來到.
七一及紀念禮.
三章十四十五節.
是幫助我們認識耶和華多一點的.
是祂說.
這是我的名字.
你要來紀念我.
願主祝福你.
多一點紀念主耶和華.
我們一起祈禱吧.
我們的神.
你告訴我們.
你的名字叫做M.
是自有永有.
永不改變.
你又直接用以色列人.
將你的名字發展出來.
成為耶和華這個名字來紀念你.
我們就用這個名字來紀念你.
雖然當中的演變.
是有一點複雜.
或者有一點年代.
不過不要緊.
我們最重要知道.
你是這個名字.
每記起你的時候.
我們就記得.
你是我們屬靈的祖宗.
是我們歷代屬靈祖宗的神.
你是跟我們納入的上帝.
你是那個永不改變的上帝.
讚美主.
願你得到榮耀.
我們更多紀念和信靠你.
祈禱奉耶穌名求.

$^{841}$阿們.
\newpage



\section{民數記 8:1-26}
\label{sec:031902}
\textbf{分別為聖,轉眼為耳}~~~~甄盧鳳雯傳道
\newline
\newline
連結: \href{http://www.church.com.hk/acms/content.asp?site=cdc\&op=show\&type=product\&code=031902\&layout=sermon}{\texttt{031902}} ~~~~ 語音日期: 2026-03-28
\newline
\newline
\hyperref[sec:031901]{\small{< < < PREV SERMON < < <}}
~
\hyperref[sec:index_chronic]{\small{[返順時目]}}
~
\hyperref[sec:031903]{\small{> > > NEXT SERMON > > >}}
\newline
\newline
民數記 8:1-26
\newline
\begin{longtable}{cl}
\hline
\hline
章節 & 經文 (和合本修訂版)\\
\hline
8:1 & \begin{tabularx}{0.7\textwidth}{X} 耶和華吩咐摩西說： \end{tabularx} \\ \\ \relax
8:2 & \begin{tabularx}{0.7\textwidth}{X} 「你要吩咐亞倫，對他說：點燈的時候，七盞燈都要照亮燈臺前面。」 \end{tabularx} \\ \\ \relax
8:3 & \begin{tabularx}{0.7\textwidth}{X} 亞倫就照樣做了；他點燈，照亮了燈臺前面，正如耶和華所吩咐摩西的。 \end{tabularx} \\ \\ \relax
8:4 & \begin{tabularx}{0.7\textwidth}{X} 燈臺是這樣造的：燈臺是用金子錘出來的，連座帶花都是錘出來的。摩西照著耶和華所指示的樣式造了燈臺。 \end{tabularx} \\ \\ \relax
8:5 & \begin{tabularx}{0.7\textwidth}{X} 耶和華吩咐摩西說： \end{tabularx} \\ \\ \relax
8:6 & \begin{tabularx}{0.7\textwidth}{X} 「你要從以色列人中選出利未人來，潔淨他們。 \end{tabularx} \\ \\ \relax
8:7 & \begin{tabularx}{0.7\textwidth}{X} 你要這樣做來潔淨他們：要用潔淨的水彈在他們身上，又叫他們用剃刀剃刮全身，洗淨衣服，潔淨自己。 \end{tabularx} \\ \\ \relax
8:8 & \begin{tabularx}{0.7\textwidth}{X} 然後他們要取一頭公牛犢，以及同獻的素祭，就是調油的細麵。你要另取一頭公牛犢作贖罪祭。 \end{tabularx} \\ \\ \relax
8:9 & \begin{tabularx}{0.7\textwidth}{X} 你要帶利未人到會幕前，並且要召集以色列全會眾。 \end{tabularx} \\ \\ \relax
8:10 & \begin{tabularx}{0.7\textwidth}{X} 你要把利未人帶到耶和華面前，以色列人要為利未人按手。 \end{tabularx} \\ \\ \relax
8:11 & \begin{tabularx}{0.7\textwidth}{X} 亞倫要從以色列人中將利未人奉獻在耶和華面前，作為搖祭，使他們事奉耶和華。 \end{tabularx} \\ \\ \relax
8:12 & \begin{tabularx}{0.7\textwidth}{X} 利未人要按手在那兩頭牛的頭上；你要將一頭作贖罪祭，一頭作燔祭，獻給耶和華，為利未人贖罪。 \end{tabularx} \\ \\ \relax
8:13 & \begin{tabularx}{0.7\textwidth}{X} 你也要使利未人站在亞倫和他兒子面前，將他們奉獻給耶和華，作為搖祭。 \end{tabularx} \\ \\ \relax
8:14 & \begin{tabularx}{0.7\textwidth}{X} 「你從以色列人中將利未人分別出來，利未人就歸我了。 \end{tabularx} \\ \\ \relax
8:15 & \begin{tabularx}{0.7\textwidth}{X} 你潔淨了利未人，奉獻他們作為搖祭之後，他們就可以進會幕事奉。 \end{tabularx} \\ \\ \relax
8:16 & \begin{tabularx}{0.7\textwidth}{X} 因為他們是從以色列人中全然獻給我的；我選他們歸我，代替以色列人中所有頭胎的長子。 \end{tabularx} \\ \\ \relax
8:17 & \begin{tabularx}{0.7\textwidth}{X} 因為以色列人中凡頭生的，無論是人或牲畜，都是我的。我在埃及地擊殺所有頭生的那日，已將他們分別為聖歸我。 \end{tabularx} \\ \\ \relax
8:18 & \begin{tabularx}{0.7\textwidth}{X} 我選利未人代替以色列人中所有頭生的。 \end{tabularx} \\ \\ \relax
8:19 & \begin{tabularx}{0.7\textwidth}{X} 我從以色列人中將利未人給亞倫和他的兒子作為賞賜，在會幕中為以色列人事奉，又為以色列人贖罪，免得以色列人因挨近聖所而遭受災禍。」 \end{tabularx} \\ \\ \relax
8:20 & \begin{tabularx}{0.7\textwidth}{X} 摩西、亞倫和以色列全會眾就向利未人這樣做。關於利未人，凡耶和華怎樣吩咐摩西，以色列人就向他們照樣做了。 \end{tabularx} \\ \\ \relax
8:21 & \begin{tabularx}{0.7\textwidth}{X} 於是利未人從罪中潔淨自己，洗淨衣服。亞倫將他們奉獻在耶和華面前，作為搖祭，又為他們贖罪，潔淨他們。 \end{tabularx} \\ \\ \relax
8:22 & \begin{tabularx}{0.7\textwidth}{X} 然後利未人進去，在亞倫和他兒子面前，在會幕中事奉。關於利未人，耶和華怎樣吩咐摩西，他們就向利未人照樣做了。 \end{tabularx} \\ \\ \relax
8:23 & \begin{tabularx}{0.7\textwidth}{X} 耶和華吩咐摩西說： \end{tabularx} \\ \\ \relax
8:24 & \begin{tabularx}{0.7\textwidth}{X} 「這是有關利未人的：二十五歲以上的人都要前來任職，在會幕裡事奉。 \end{tabularx} \\ \\ \relax
8:25 & \begin{tabularx}{0.7\textwidth}{X} 到了五十歲，他們就要從事奉的工作中退休，不再事奉， \end{tabularx} \\ \\ \relax
8:26 & \begin{tabularx}{0.7\textwidth}{X} 只可在會幕裡輔助他們的弟兄盡責，他們自己不再事奉了。關於利未人的職責，你要向他們這樣做。」 \end{tabularx} \\ \\
[1ex]
\hline
\hline
\end{longtable}
$^{1}$各位主內的弟兄姊妹.
你們信不信.
感不感受到耶穌很愛你們?.
這麼小聲.
這本是聖經.
是記載上帝的啟示.
和祂作為的歷史事實.
要我們信耶穌基督是神的兒子.
要我們信因為耶穌的名得生命.
大家都知道.
聖經是上帝賜下的話語.
是不變的真道.
是我們每日生命的糧.
日夜要思想.
大家有沒有想像過.
外面黑色這個皮.
就好像上帝兒子耶穌基督.
穿著的外套.
背起所有人的嘴.
裡面白色發光的.
就是上帝光照我們內心的說話.
好像一條道路.
道成肉身.
在我們中間帶領著我們去走.
祝福著我們.
我們想深一層.
聖經上帝的話.
是聽先還是看先?.
亞伯拉罕 以撒 雅各 摩西.
是聽還是看上帝的話先?.
兩樣都有.
但是先聽後看.
信心的古人是用謙卑敬虔的態度.
去聽上帝的說話.
然後照著上帝的吩咐.
重新做人.
今天我們是不是這樣?.
問問你身邊的弟兄姊妹.
是不是謙卑敬畏上帝.
聽道 行道?.

$^{41}$一時時吧.
所以我們今天要回來崇拜.
要回到團契主日學.
一起持續給主塑造我們成長.
民數記記載當時敬畏上帝的以色列人.
聽上帝吩咐摩西 雅倫的方法.
去做準備.
每個人都知道自己的職責和界限.
向著迦南應許地去進發.
開始的時候.
以色列人敬畏上帝.
很直接 馴服上帝就照做.
但是時間一久了.
就出現了人性的軟弱.
由百姓領袖都得罪上帝.
他們由聽上帝的吩咐.
變成聽自己的需要.
見到眼前對自己沒有利益的東西.
他們就會埋怨 不滿 不服.
比較計較 甚至投訴.
面對這班人慈愛公然的上帝.
很有耐性.
對他們有提醒 有懲罰.
也都透過律例.
幫他們知道自己是犯罪.
不要放棄過他們.
上帝已經跟他們納藥.
給了認許他們.
問題是他們生活的目標優先次序.
怎樣選擇回應上帝的帶領.
進入應許地.
舊約以色列人的延伸.
就是新約的教會.
就是你跟我一班信耶穌基督的人.
我們信主走出黑暗之後.
現在努力走在光明救恩.
就好像曠野的旅程當中.
向著新天新地去進發.
等耶穌再來跟我們一起分享大賢直.
我們要從文素機學習.

$^{81}$以此為鑑.
不要步他們後塵.
要持續敬畏上帝.
單單聽上帝的話.
做好準備.
知道自己基督徒的職責和界限.
努力地向前行.
即是有沒有在做今天第八章的提醒.
分別為聖.
轉眼為耳.
現實告訴我們.
我們關聖是先看了再聽.
不是先聽後看.
例如聽崇拜講到.
我們很容易用自己的思想框架.
去看誰在講.
判斷聽不聽.
用不用時間去聽.
時間很寶貴.
即是我們會選擇講完去聽.
不是為了聽上帝的話.
去聽和去學.
母領受會說講完.
而不是看自己有沒有預備好見上帝的心態.
我以前也是這樣.
不出奇.
因為我們是有限的人.
成長學的是批判思維.
習慣理性分析.
多過用心聽無限的神跟我們講話.
無可厚非.
但是想深一層.
創造世界掌管萬有的上帝.
先是講給人聽和看人.
我們是受造物.
是應該謙卑敬畏地.
去聽上帝跟我們講什麼.
感恩祂給我們自由.
去選擇回應做不做.
但當然自己要負責任.

$^{121}$先後次序.
是上帝先講給人聽.
人聽了上帝說話之後.
回應上帝.
和上帝互動.
其實舊約摩西五經.
就是上帝跟人講的啟示.
吩咐,命令,引導,責備.
人是要聽.
民數記也不例外.
他的希伯來文第一個字是「耶和華說」.
即是聖潔的上帝跟我們講話.
我們基本的責任是自潔.
才可以站在聖潔的上帝面前.
「人非聖潔不能見主面」.
耶穌親自提醒過我們.
有耳可聽的就應當聽.
聽上帝之道而遵守的人.
是有福的.
第八章二十六節.
經文正正是上帝跟以色列人講的三個吩咐.
摩西,祭司,利未人,傳會眾都要聽.
一至四節吩咐祭司點燈照光回望.
五至二十二節吩咐利未人要潔淨進入回望.
二十三至二十六節吩咐回望事奉的要有界限.
即是事奉的人無論你是哪一個.
都要敬畏上帝謙卑自己.
先聽上帝的吩咐.
潔淨自己進入上帝同在的回望裡事奉.
祭司要跟從上帝的要求.
點燈照光回望的內部.
利未人要潔淨自己.
分別為聖設界限.
走進回望事奉.
上一章其實已經交代了.
摩西清楚知道上帝在回望跟他們一起.
各個族長奉獻了所有的祭物.
摩西主動走進回望.
站在聖所至聖所的前面.
打算開口跟神說話.

$^{161}$但是她聽到法櫃上的私人座.
前面有聲音.
就是上帝跟摩西說話.
摩西只有一件事可以做.
是什麼?.
聽.
我們一起謙卑自己.
敬畏地聽上帝的話.
耶和華曉遇摩西說.
你告訴亞倫說.
點燈的時候七盞燈都要向燈台前面發光.
亞倫便這樣行.
他點燈台上的燈是燈向前發光.
是照耶和華所吩咐摩西的.
這燈台的做法是用金紙除出來的.
連座大花都是除出來的.
這裡我們聽到上帝吩咐的第一點.
祭司要照耶和華指示點燈的方法.
令燈台發光.
上帝清楚吩咐可進入聖所的祭司點燈的方法.
要向前發光.
摩西的回應就是聽耶和華的命令.
告訴亞倫怎樣點燈.
沒有加自己的意思.
其實上帝一早已經吩咐摩西怎樣去製造這個燈台.
摩西聽上帝所指示的樣式.
用金突出連座大花的燈台.
摩西要很用心聽清楚上帝的指示.
用錘在金上面做出一個很細緻的燈台.
而亞倫的回應就是用心聽摩西的吩咐.
照耶和華的樣式.
教導點著聖所金燈台上的七盞燈.
清楚目的是要向前發光.
照光會幕內部.
亞倫的回應是什麼.
就是第三節.
辯者讓行.
強調跟隨上帝的帶領.
就是用心聽清楚上帝指示的樣式照做.
不是用自己有限的習慣或者想法去做.

$^{201}$大家記不記得我們上年啟示錄一張.
講過七個燈台就是七個教會.
教會是要按上帝指示祭司點燈的方法.
向前發光.
這個是分別為勝很重要的起點.
舊約上帝揀選祭司.
新約揀選我們.
不是祭司習慣怎樣做.
做多少.
是祭司有沒有用上帝吩咐的方法.
照光會幕的內部.
一到四節是一個很好的提醒.
上帝不是問我們為祂做多少.
做到多少.
是問我們肯不肯用上帝吩咐的方法.
給上帝的話照光聖殿.
照下我們現在的生命實況.
教會真實的情況.
因為我們每個人就是神的殿.
就好像手機鏡頭.
要抹乾淨給光線才可以照進去.
我們個個都用手機拍照.
但鏡頭常常被手指印弄髒.
拍出來不夠清.
問題是什麼.
是手機不夠好.
還是鏡頭太髒.
有時我們會覺得人生很模糊.
方向很亂.
很多時不是上帝的光不夠.
而是我們的鏡頭太多髒東西.
例如犯錯不應.
自以為是.
貪心.
懶惰.
說謊.
驕傲.
拖延.
我們肯不肯給上帝抹一抹我們的眼睛.
給自己反映出上帝的樣式.

$^{241}$給上帝的光照清楚自己內心的灰暗.
謙卑地認罪悔改.
承認自己真的有罪.
很需要上帝.
弟兄姊妹.
你的生命鏡頭.
即是眼睛.
肉眼和心眼.
有沒有些東西你知道上帝不喜悅.
但一直不肯給上帝抹走.
於是不肯說對不起.
不肯寬恕人.
錯的通常是人.
對方才有責任.
自己就沒有反省的責任.
你的眼睛就是身體的燈.
當你的眼睛明亮.
全身就會光明.
當眼睛昏花.
全身就會黑暗.
所以你要注意.
免得你裡面的光暗了.
第一個小結.
我們基督徒生活的態度.
是要全敬畏上帝謙卑的心.
去聽清楚上帝吩咐的樣式.
向前發光.
不是聽自己的方式.
到二十二節.
我想邀請大家聽金色的字.
然後藍色的字就一起讀.
因為都幾長.
耶和華曉遇摩西說.
你從以色列人中選出利未人來.
潔淨他們.
潔淨他們當這樣行.
用垂墜水彈在他們身上.
又叫他們用剃刀刮全身.
洗衣服潔淨自己.
然後叫他們取一隻公牛毒.

$^{281}$並同獻的掃祭.
就是調油的細面.
你要另取一隻公牛毒.
作贖罪祭.
請會眾讀.
將利未人奉到會幕前.
招聚以色列全會眾.
將利未人奉到耶和華面前.
以色列人要按守在他們頭上.
阿倫也將他們奉到耶和華面前.
以色列人當作遙祭.
視他們好辦耶和華的事.
利未人要按守在那兩隻牛的頭上.
你要將一隻作贖罪祭.
一隻作凡祭.
獻給耶和華.
為利未人贖罪.
你也要使利未人站在阿倫和他兒子面前.
將他們當作遙祭.
奉給耶和華.
這樣你從以色列人中將利未人分別出來.
利未人便要歸我.
此後利未人要進去辦會幕的事.
你要潔淨他們.
將他們當作遙祭奉上.
因為他們是從以色列人中傳言給我的.
我揀選他們歸我.
是代替以色列人中一切投身的.
以色列人中一切投身的.
連人大生畜都是我的.
我在埃及地擊殺一切投身的雅天.
將他們分別為聖歸我.
我揀選利未人代替以色列人中一切投身的.
我從以色列人中將利未人當作賞賜給阿倫和他的兒子.
在會幕中辦以色列人的事.
又為以色列人贖罪.
免得他們埃近聖所.
有災殃臨到他們中間.
摩西,阿倫並以色列傳會中便向利未人如此行.
凡耶和華子著利未人所吩咐摩西的.

$^{321}$以色列人就向他們這樣行.
於是利未人潔淨自己.
除了罪,洗了衣服.
阿倫將他們當作遙祭奉到耶和華面前.
又為他們贖罪,潔淨他們.
然後利未人進去.
在阿倫和他的兒子面前.
在會幕中辦事.
耶和華子著利未人怎樣吩咐摩西.
以色列人就怎樣向他們行了.
第五至二十二節.
我們聽到上帝的吩咐第二點.
被揀選的利未人要潔淨.
上帝清楚吩咐摩西揀選利未人.
他們先要經過一個完整的潔淨禮.
第一,先由摩西用除罪水潔淨利未人.
潔淨他們禮儀上的不潔.
第二,利未人自己要全身剃毛,洗衣服.
潔淨自己,栽獻贖罪祭和犯祭.
潔淨他們生活上嚴重的不潔.
第三,以色列人按手在利未人的頭上.
將利未人獻給耶和華作遙祭.
好像活的祭.
替代全以色列人中一切頭三.
專心獻給上帝服侍.
時間關係.
有關各種的祭.
你們可以自由看利未記.
潔淨利未人被分別出來.
不是升他們做特別階級.
而是全個群體同心的獻給聖潔的上帝.
利未人不再聽自己和世界的聲音去計劃人生.
而是聽上帝的聲音帶領分別出來.
自願的走入光宗去服侍.
和服侍百姓.
這個是很重要的提醒.
是整體分別為性的潔淨.
包括個人和群體.
合一的同心.
不單止是自己和上帝有親密的關係.

$^{361}$同時都需要群體合一的分辨.
這個就是上帝的聲音.
避免自以為是.
自以為而.
引起衝突和不和.
今天你和我都是新約的祭司.
上帝不是只選擇牧者傳道人.
是選擇所有屬於祂的人分別為性.
人受聖靈感動.
在內心信耶穌.
在群體中公開的受洗.
同主同死同埋葬同復活成為新造的人.
重新我們每日過基督教道的樣式生活.
好像穿了手術袍的醫生和醫護團隊.
進手術室之前他們一定洗手.
消毒.
穿手術袍.
戴口罩.
手套.
一旦進了手術室.
他們不會再做其他沒有關係的事.
因為他們意識到他們在一個特別神聖的環境.
手術袍代表那一刻他們被分別出來.
只有一個目的.
就是救人.
同樣上帝在舊約的尼微人.
延伸到新約就是我們基督徒.
被分別出來.
讓我們穿上一個新的身份.
就是上帝兒女的身份.
不再只是為自己打算.
而是為上帝的家.
為人的靈魂要付出.
弟兄姊妹.
你們何時會穿這件袍.
不要只是回教會才穿.
其實無時無刻.
我們都穿著這件.
分別為聖的身袍.
言行合一的跟隨基督.

$^{401}$小心隨低.
變回一個不信上帝.
沒有上帝的人.
第二個小結.
我們要潔淨.
分別為聖是全天候.
個人和群體.
我們再聽第23到26節.
耶和華曉遇摩西說.
尼微人是這樣.
從二十五歲以外.
他們要前來任職.
辦會幕的事.
到了五十歲要停工退任.
不再辦事.
只要在會幕裡.
和他們的弟兄一同侍候.
謹守所吩咐的.
不再辦事了.
至於所吩咐尼微人的.
你要這樣向他們行.
這裡我們聽到.
上帝吩咐的第三點.
會幕要有界限.
一同侍候.
謹守上帝所吩咐.
上帝清楚定了尼微人服侍的年齡.
二十五歲開始學習預備.
大約三十歲全面服侍.
五十歲就退下操重的工作.
改為幫手輔助.
不是收山不理事.
而是角色上的轉變.
由出手做.
變為在旁邊扶持.
守護指導.
這一幅有界限的圖畫提醒.
事奉是有分季節.
原因是上帝重視.
工人的健康和持久.

$^{441}$不是要人做到殘.
做到burn out 才是聖潔.
分別為聖.
不單是手做.
更多是耳聽.
聽得懂上帝的聲音.
也聽到別人的需要.
好像樂團中資深的樂手.
一個樂團有年輕的演奏家.
很有技術 很火熱.
不過整隊樂隊整體可不可以配合得到.
通常都是要靠資深的樂手.
有些樂手退休了.
不再上台表演.
會做導師.
教年輕人怎樣指揮.
怎樣互相配合.
怎樣不搶鏡.
他們現在最重要的能力.
不是有多專業.
而是耳朵有多敏銳.
聽總指揮.
聽樂隊整體的聲音.
同樣 屬靈生命的成熟.
不是用多少自己超凡的方法.
更加不是我自己可以做到多專業.
而是對上帝的聲音有多敏銳.
多肯為整體教會的建立.
去聽和去順從.
由樂意聽人 做就人.
年長或資深的信徒.
不是教會的邊緣人.
而是上帝安排給我們的屬靈導師.
上帝就是要用你們內化了的經文.
實踐過 錯過 經歷過 改過的經驗.
加上你們真誠的同理心.
去聽 去關心 去扶持.
去傳承給下一代.
第三個小結就是.
教會全體要合一同心.

$^{481}$同感一靈.
不分你我 互助互補.
這一幅大致是剛才形容舊約的會幕.
和新約後的教會的情況.
我用AI整給大家看看 參考一下.
但無論如何我們都是按聖經教導.
帶人回到上帝面前.
聆聽上帝的吩咐.
和引領教會群體發光 榮耀上帝.
總結 今天三點.
第一 個人發光.
基督徒生活的態度是要盡自己的本份.
敬畏上帝 謙卑地聽清楚上帝的吩咐.
向前發光.
第二 潔淨.
分別為性的潔淨是包括全天候.
個人和群體性.
第三 要全體發光.
教會全體要同心合一 同感一靈.
毋分利我 互助互補.
敏銳聽上帝的大靈.
我們和聖誕樹都很相似.
被揀選出來盛載燈光.
聖誕節之前家人會去揀一棵普通的樹回家.
一開始這棵聖誕樹不再是一樣.
我們會幫它抹塵 擺好.
然後將燈掛上去 插上電.
一開燈 全屋最吸睛的就是它.
同樣 上帝在人群中揀選我們做祂的指紋 祂的教會.
不是因為我們特別聰明 特別聰明.
而是上帝想將祂的光掛在我們身上.
叫人一望就想起祂.
如果你是一棵聖誕樹.
上帝現在想在你身上加多一串燈.
你會不會說 夠了 太光了 很累了.
還是你願意再光一點.
照亮更多黑暗的角落呢.
我們有沒有每一天走入光裡面.
即是有沒有安靜時間.
讀經 祈禱 靈修 親近上帝.

$^{521}$讓上帝照照自己.
有沒有一些習慣 態度 關係.
一照就知道不對.
但一直不肯讓上帝動我們.
聖誕樹每一件裝飾都有意思.
星星 百里行之星.
是引導人到耶穌面前.
燈泡象徵基督是世界的光.
紅色的飾物象徵基督的寶血 愛的記號.
但這些裝飾從來都不是為樹自己.
是為經過的人.
叫他們記得勝誕的主角是基督.
同樣 上帝在我們生命加添了恩慈 經歷 祝福.
其實不是為了我們自己好看.
而是叫人因為我們生命的改變.
想起耶穌的大能和恩典.
上帝擺在你們生命的裝飾是什麼.
你的才剛 走過的傷痛 得著的醫治.
這些會不會只是成為你自豪的資本.
還是讓人想起上帝無條件犧牲的愛.
不是每一個人都要放下自己所做的.
去全職做傳道人 牧者 宣教士.
但我們每一個人可以在自己的崗位上.
清楚地承認.
我是主所用的利未人.
時間是最難奉獻的祭.
但它是最反映人的心.
我們的說話 處事方式.
反映著我們屬於誰.
聖誕節完了 聖誕樹要拆走.
修好燈飾 要搬走樹.
客廳突然空了.
但見過的人會記得.
那段時間家裡很光 很溫暖.
很多人聚集.
這些記憶留在心裡.
聖誕樹退場.
但燈光曾照過的地方 仍然存在.
下一年 或者再出場.
再被裝飾.

$^{561}$或者換了另一棵摺盤.
同樣利未人到了下一個階段.
要退下前線.
不等於他們不再重要.
反而成為見證.
提醒後來的.
如何在上帝面前忠心地服侍一生.
分別衛星不只是開頭的熱情.
是一生的節奏.
有起步 有奔跑 有傳承.
可能你現在的人生階段很忙.
鼓勵你們每天抽點時間.
轉眼離開娛樂 手機 遊戲 社交媒體 電視 旅行.
學習轉移去聽上帝的說話.
聽一聽上一代的屬靈忠告.
建議年輕的找一位屬靈的前輩.
固定的請教 請他為你祈禱.
無論你現在多麼繁忙.
都是上帝特別使用你的時間.
要抽時間安靜 分別衛星.
用零碎的時間去祈禱.
聽聽上帝跟你說什麼.
在決定之前停一停 問下主.
或許你是長輩之心的信徒.
鼓勵你們用最寶貴的商議.
和真誠的心.
去關懷 代討 做陪談 服侍.
成為年輕一代的屬靈長輩.
因為你們更加懂得上帝的旨意.
更加懂得安靜聽上帝.
更加懂得明白人.
願我們每一個學務給上帝光照.
從內心承認自己需要被潔淨.
全教會分別衛星.
發更大的光 榮耀上帝.
好像抹乾淨的手機鏡頭.
被光線清楚照入.
好像被裝飾的聖誕樹.
不是再為自己發光.
是為主發光.

$^{601}$在不同的人生季節.
用不同的方式.
為教會你們屬靈的家服侍.
學習由「出聲」到「用耳聽」.
「幫手」或「停手」.
成為大家的祝福.
本身我今天預備講章差不多.
早兩天退休.
聖靈提醒我要多講一點.
就是今時今日.
教會很需要認真面對.
整體的被潔淨 被光照 向前行.
這本書《教會復興的障礙》.
出現在我面前.
書是這樣說的.
今日的教會已經不知不覺.
被世俗化侵蝕得相當嚴重.
這種侵蝕並不是罪大惡極.
大錯特錯的事.
而是一些普通的事.
就是忽視聖經所講的教會紀律.
Discipleship need discipline.
門徒訓練是需要紀律.
我們很容易忽視蒙恩之道.
忽視公開的誓言.
忽視洗禮.
不忠實地參加崇拜.
會傳一些這個好那個不好.
不恰當 未經證實的說話.
基督徒更常常用相對的態度.
去討論教會紀律的問題.
所關心的不是絕對是非問題.
而是自己所提的意見有沒有達到效果.
所關心的不是上帝藉著聖經所要求的是什麼.
而是他們所提的見解能不能夠贏得人心.
上帝的教會紀律最核心的問題.
並不是明不明智 審不審身.
也不是我們中不中意.
而是不是上帝在聖經所吩咐教會的做法.
昨天我們牧者們裡.

$^{641}$很多牧者分享 很多長輩分享.
我們最深刻一句.
就是教會要做教會要做的事.
這一間是不是分別為聖的教會.
有沒有找到打開的聖經在哪裡.
還是這一間.
AI的講源.
還是上帝吩咐是這一間.
我們今天學到用心去聽清楚上帝的吩咐.
按上帝的樣式.
令上帝的教會整體發光.
並不是用世俗的樣式.
自己想的樣式.
多人同意的樣式.
現在邀請大家在心裡面問主這三個問題.
主啊 在哪一個範圍.
你要我走出黑暗 給你光照.
主啊 你在我生命加上了什麼裝飾.
是你想用來祝福人.
主啊 你要我怎樣和弟兄姊妹相處.
讓你的教會發光榮耀你.
我們用半分鐘的時間.
聽聽上帝怎樣回覆你.
九曰以賽亞先知提醒.
列國啊 要近前來聽.
眾民啊 要質疑而聽.
地和其上所充滿的.
世界和其中一切所出的.
都應當聽.
新約亞國書提醒.
親愛的弟兄姊妹們.
這是你們所知道的.
要脫去一切的污穢和盈餘的邪惡.
全溫柔的心.
領受能救你們靈魂的道.
只是你們要行道.
不要單單聽道自己 欺控自己.
唯有詳細察看那傳遍使人自由之律法.
實在行出來.
就在他所行的事上必然得福.

$^{681}$若有人自以為虔誠.
卻不勒住他的舌頭.
反欺控自己的心.
這人的虔誠是虛.
在上帝我們的父面前.
那清潔沒有玷污的虔誠.
就是看顧在患難中的孤兒寡婦.
並且保守自己不顛掩世俗.
如果我們現在看一看身處以色列的人.
當地警暴一響.
他們九十秒就會走進去放空洞.
今天我們在香港聽到可能會.
看看真還是假.
先看清楚一點.
這就是我們未發現基督徒.
是處身屬靈真賤的現實之中.
盼望我們明白為何要警醒分別為聖.
轉眼為異.
立即行動.
聽道行道.
不要賭主還未回來.
或者猜他何時回來.
因為沒有人知道主何時回來.
與他更加不要拖延自己靈明的成長.
最後聖靈有話說.
來自希伯來書三章七至十三節.
你們今日若聽上帝的話.
就不可眼著心.
好像在曠野感野.
上帝發怒試探上帝的時候一樣.
在那裡你們的祖宗試探上帝.
所以上帝厭煩那世代的人.
說他們心裡常常迷糊.
竟不曉得上帝的作為.
上帝就在路中起誓說.
「眼深的人斷不可進入我的安息」.
弟兄姊妹們.
你們要謹慎.
免得你們中間或有人存著不信的惡心.
把永生上帝離棄了.

$^{721}$總要趁著還有今天.
天天彼此相勸.
免得你們中間有人被罪迷惑.
心裡就乾硬了.
我們同心土共.
慈愛公義的天父上帝.
在預府期用民數記八章提醒.
我們已經恩信耶穌是神的兒子.
被分別為聖.
耶穌為我們的罪.
釘身十字架分別為聖.
令我們被稱為義.
被稱為神的兒女.
不是我們有幾好.
我們知道我們沒有耐性.
確實有犯罪得罪你.
更加在生活上.
離開了你在聖經吩咐我們.
要和教會一起潔淨發光.
求你光照我們內心.
潔淨我們心思意念.
運用你所賜的恩賜和經歷.
成為教會和祝福人的見證.
引導我們在人生不同的季節.
持續警醒每日的優先次序.
真正懂得分別為聖.
轉眼為義.
聽著你的聲音帶領我們向前行.
單單定睜仰望你.
不再是其他東西.
立志擺上我們的心.
給你察驗.
逐漸改變我們生活的態度和方式.
為的是要榮耀你.
奉主耶穌基督聖潔名字祈求.
阿們.
\newpage



\section{民數記 8:1-26}
\label{sec:031903}
\textbf{分別為聖的服侍}~~~~吳安強牧師
\newline
\newline
連結: \href{http://www.church.com.hk/acms/content.asp?site=cdc\&op=show\&type=product\&code=031903\&layout=sermon}{\texttt{031903}} ~~~~ 語音日期: 2026-03-29
\newline
\newline
\hyperref[sec:031902]{\small{< < < PREV SERMON < < <}}
~
\hyperref[sec:index_chronic]{\small{[返順時目]}}
~
\hyperref[sec:031904]{\small{> > > NEXT SERMON > > >}}
\newline
\newline
民數記 8:1-26
\newline
\begin{longtable}{cl}
\hline
\hline
章節 & 經文 (和合本修訂版)\\
\hline
8:1 & \begin{tabularx}{0.7\textwidth}{X} 耶和華吩咐摩西說： \end{tabularx} \\ \\ \relax
8:2 & \begin{tabularx}{0.7\textwidth}{X} 「你要吩咐亞倫，對他說：點燈的時候，七盞燈都要照亮燈臺前面。」 \end{tabularx} \\ \\ \relax
8:3 & \begin{tabularx}{0.7\textwidth}{X} 亞倫就照樣做了；他點燈，照亮了燈臺前面，正如耶和華所吩咐摩西的。 \end{tabularx} \\ \\ \relax
8:4 & \begin{tabularx}{0.7\textwidth}{X} 燈臺是這樣造的：燈臺是用金子錘出來的，連座帶花都是錘出來的。摩西照著耶和華所指示的樣式造了燈臺。 \end{tabularx} \\ \\ \relax
8:5 & \begin{tabularx}{0.7\textwidth}{X} 耶和華吩咐摩西說： \end{tabularx} \\ \\ \relax
8:6 & \begin{tabularx}{0.7\textwidth}{X} 「你要從以色列人中選出利未人來，潔淨他們。 \end{tabularx} \\ \\ \relax
8:7 & \begin{tabularx}{0.7\textwidth}{X} 你要這樣做來潔淨他們：要用潔淨的水彈在他們身上，又叫他們用剃刀剃刮全身，洗淨衣服，潔淨自己。 \end{tabularx} \\ \\ \relax
8:8 & \begin{tabularx}{0.7\textwidth}{X} 然後他們要取一頭公牛犢，以及同獻的素祭，就是調油的細麵。你要另取一頭公牛犢作贖罪祭。 \end{tabularx} \\ \\ \relax
8:9 & \begin{tabularx}{0.7\textwidth}{X} 你要帶利未人到會幕前，並且要召集以色列全會眾。 \end{tabularx} \\ \\ \relax
8:10 & \begin{tabularx}{0.7\textwidth}{X} 你要把利未人帶到耶和華面前，以色列人要為利未人按手。 \end{tabularx} \\ \\ \relax
8:11 & \begin{tabularx}{0.7\textwidth}{X} 亞倫要從以色列人中將利未人奉獻在耶和華面前，作為搖祭，使他們事奉耶和華。 \end{tabularx} \\ \\ \relax
8:12 & \begin{tabularx}{0.7\textwidth}{X} 利未人要按手在那兩頭牛的頭上；你要將一頭作贖罪祭，一頭作燔祭，獻給耶和華，為利未人贖罪。 \end{tabularx} \\ \\ \relax
8:13 & \begin{tabularx}{0.7\textwidth}{X} 你也要使利未人站在亞倫和他兒子面前，將他們奉獻給耶和華，作為搖祭。 \end{tabularx} \\ \\ \relax
8:14 & \begin{tabularx}{0.7\textwidth}{X} 「你從以色列人中將利未人分別出來，利未人就歸我了。 \end{tabularx} \\ \\ \relax
8:15 & \begin{tabularx}{0.7\textwidth}{X} 你潔淨了利未人，奉獻他們作為搖祭之後，他們就可以進會幕事奉。 \end{tabularx} \\ \\ \relax
8:16 & \begin{tabularx}{0.7\textwidth}{X} 因為他們是從以色列人中全然獻給我的；我選他們歸我，代替以色列人中所有頭胎的長子。 \end{tabularx} \\ \\ \relax
8:17 & \begin{tabularx}{0.7\textwidth}{X} 因為以色列人中凡頭生的，無論是人或牲畜，都是我的。我在埃及地擊殺所有頭生的那日，已將他們分別為聖歸我。 \end{tabularx} \\ \\ \relax
8:18 & \begin{tabularx}{0.7\textwidth}{X} 我選利未人代替以色列人中所有頭生的。 \end{tabularx} \\ \\ \relax
8:19 & \begin{tabularx}{0.7\textwidth}{X} 我從以色列人中將利未人給亞倫和他的兒子作為賞賜，在會幕中為以色列人事奉，又為以色列人贖罪，免得以色列人因挨近聖所而遭受災禍。」 \end{tabularx} \\ \\ \relax
8:20 & \begin{tabularx}{0.7\textwidth}{X} 摩西、亞倫和以色列全會眾就向利未人這樣做。關於利未人，凡耶和華怎樣吩咐摩西，以色列人就向他們照樣做了。 \end{tabularx} \\ \\ \relax
8:21 & \begin{tabularx}{0.7\textwidth}{X} 於是利未人從罪中潔淨自己，洗淨衣服。亞倫將他們奉獻在耶和華面前，作為搖祭，又為他們贖罪，潔淨他們。 \end{tabularx} \\ \\ \relax
8:22 & \begin{tabularx}{0.7\textwidth}{X} 然後利未人進去，在亞倫和他兒子面前，在會幕中事奉。關於利未人，耶和華怎樣吩咐摩西，他們就向利未人照樣做了。 \end{tabularx} \\ \\ \relax
8:23 & \begin{tabularx}{0.7\textwidth}{X} 耶和華吩咐摩西說： \end{tabularx} \\ \\ \relax
8:24 & \begin{tabularx}{0.7\textwidth}{X} 「這是有關利未人的：二十五歲以上的人都要前來任職，在會幕裡事奉。 \end{tabularx} \\ \\ \relax
8:25 & \begin{tabularx}{0.7\textwidth}{X} 到了五十歲，他們就要從事奉的工作中退休，不再事奉， \end{tabularx} \\ \\ \relax
8:26 & \begin{tabularx}{0.7\textwidth}{X} 只可在會幕裡輔助他們的弟兄盡責，他們自己不再事奉了。關於利未人的職責，你要向他們這樣做。」 \end{tabularx} \\ \\
[1ex]
\hline
\hline
\end{longtable}
$^{1}$主持人:各位聽眾早晨.
主持人:《新穎記》第八章第一到第四節是這樣記載的.
主持人:耶和華吩咐摩西說你要吩咐亞倫對她說.
主持人:點燈的時候七盞燈都要照亮燈台前面.
主持人:亞倫就照樣做了.
主持人:她點燈照亮了燈台的前面.
主持人:正如耶和華所吩咐摩西的.
主持人:燈台是這樣做的.
主持人:燈台是用金打出來的.
主持人:連座帶花都是打出來的.
主持人:整個沒有任何配件加進去.
主持人:摩西照著耶和華所吩咐的樣式做了燈台.
主持人:《文俗記》第七章詳細記載了回望奉獻禮的詳情.
主持人:有時候我們看這章經文.
主持人:我們通常都是這樣的.
主持人:看一看下面的都差不多.
主持人:讀完一段就停了.
主持人:因為覺得下面說的都是差不多.
主持人:但上次林牧師告訴我們.
主持人:為什麼《聖經》要這樣記載十二次.
主持人:回望的奉獻代表了聖所器具的啟用.
主持人:也是祭司事奉的開始.
主持人:再按《中埃及記》三十七章的記載.
主持人:聖所裡面安神的命令.
主持人:有祭壇,香壇,燈台和公桌.
主持人:公桌就是我們擺陳設餅的那張桌子.
主持人:祭壇和香壇在前一章.
主持人:當以色列十二支派的領袖來到回望作出奉獻的時候.
主持人:如果我們細心留意已經啟用了.
主持人:但在祭司進入回望事奉之前.
主持人:燈台和桌子也都要被啟用.
主持人:即是說所有的器物要先被啟用.
主持人:然後祭司先進入回望.
主持人:燈台的火被點著之後.
主持人:《聖經》告訴我們要永遠都點著.
主持人:在港澳堂崇拜完結之後.
主持人:這裡的燈就會熄滅.
主持人:但在回望裡面不是.
主持人:不是在祭司進入之前才點著.
主持人:不能省的.

$^{41}$主持人:點著就要永遠點著.
主持人:不能夠讓它熄滅.
主持人:《利未記》二十四章第二第三節告訴我們.
主持人:這麼重要的職責是由大祭司負責.
主持人:他每天都要看油夠不夠.
主持人:讓這盞燈不會熄滅.
主持人:而桌子當中的陳設餅.
主持人:也要由祭司定期更換.
主持人:相信是一個星期換一次.
主持人:而這層設餅只有祭司可以吃.
主持人:我想七天之後好不好吃呢?.
主持人:不知道.
主持人:不過《聖經》裡面告訴我們.
主持人:還有人吃過的.
主持人:記得是誰嗎?.
主持人:大衛.
主持人:但按《聖經》這裡說.
主持人:在《聖經》的教導裡面.
主持人:只是祭司吃的.
主持人:當回望一切的器皿完成啟用禮之後.
主持人:回望就可以用了.
主持人:不過硬件預備好了.
主持人:軟件也是不能怠慢.
主持人:所以接著就是利未人進入.
主持人:預備承擔服事之責的潔淨儀式.
主持人:利未人是十二支派之一.
主持人:大臣將他們分別出來.
主持人:代替以色列人當中的長子.
主持人:在《民俗記》第三章十一到十三節是這樣記載.
主持人:耶和華吩咐摩西雪漢.
主持人:我從以色列人中選了利未人.
主持人:代替以色列人中所有投胎的長子.
主持人:利未人要歸我.
主持人:因為凡投生的是我的.
主持人:我在埃及地擊殺所有投生的那日.
主持人:就把以色列中所有投生的.
主持人:無論私人或牲畜.
主持人:都分別遺姓歸我.
主持人:他們定要屬我.
主持人:我是耶和華.

$^{81}$主持人:我想大家記得.
主持人:所有以色列的家庭門尾門框抹了血之後.
主持人:奪命的天使就越過去.
主持人:以至這些以色列房屋的長子無被殺.
主持人:但神說以後這些長子是屬於我的.
主持人:對於這一班被分別出來服侍祂的人.
主持人:在進入會幕之先.
主持人:神說他們先要進行一個潔淨的禮儀.
主持人:這個禮儀的目的.
主持人:是要將利未人從世界當中分別出來獻給神.
主持人:所以神說你要從以色列人當中選出利未人來潔淨他們.
主持人:我們知道進入迦南地.
主持人:神說利未人沒有地沒有田.
主持人:沒有屬於他們的財產.
主持人:因為他們是被選召出來服侍神的人.
主持人:其他的人就會將他們的奉獻.
主持人:奉獻到會幕或者聖殿的時候.
主持人:這班利未人就從當中.
主持人:從他們的事奉裡面.
主持人:來得著他們的公認.
主持人:我們知道利美原本是以色列十二個支派之一.
主持人:但神吩咐他們將這班人分別出來.
主持人:幫助祭司亞倫和他的兒子在會幕面前執行任務.
主持人:辦理帳幕的事.
主持人:神很清楚地說出他選利未人的原因.
主持人:這裡是重複說的.
主持人:正如第三章我們看到的.
主持人:但在第八章這裡.
主持人:神再說一次.
主持人:在十六到十九節.
主持人:祂說因為他們是從以色列人當中傳言獻給我的.
主持人:我選他們歸我代替以色列人中所有投胎的長子.
主持人:因為以色列人中凡投生的無論是人或牲畜都是我的.
主持人:我在埃及地擊殺所有投生的那一天.
主持人:已經將他們分別為姓歸我.
主持人:我選利未人代替以色列人中所有投生的.
主持人:我從以色列人中將利未人給亞倫和他的兒子作為賞賜.
主持人:在會幕中為以色列人事奉.
主持人:又為以色列人贖罪.
主持人:免得以色列人因埃及聖所而遭受災禍.

$^{121}$主持人:這段經文比較長.
主持人:也講多了一點.
主持人:第三章就說.
主持人:神說我選擇以色列人出來.
主持人:是代替所有以色列人當中的長子.
主持人:然後這裡祂就說.
主持人:我將這一班人賞賜給亞倫和他的兒子.
主持人:在會幕當中來服侍.
主持人:並且為以色列人贖罪.
主持人:免得以色列人因埃及聖所而遭受災禍.
主持人:利未人蒙神所召.
主持人:代替以色列的長子在會幕當中做服侍的工作.
主持人:不過他們受職的儀式和祭司不同.
主持人:祭司要承受聖職.
主持人:就要接受高立.
主持人:聖經說他們要分別為聖來服侍神.
主持人:但是利未人只需要經過一個潔淨的儀式.
主持人:將自己從以色列民當中分別出來歸給神就可以了.
主持人:再代替以色列人當中所有的長子.
主持人:聖話要用他們的事奉來為全民贖罪.
主持人:民俗在第八章十九時這樣說.
主持人:要用他們的事奉來為全民贖罪.
主持人:儀式之後摩西就將他們交給亞倫和其他的祭司.
主持人:使他們一生都在會幕當中來服侍.
主持人:直至年日滿足.
主持人:去到五十歲就可以退休.
主持人:可能我們看到這裡就很開心了.
主持人:咦,這樣說,我們先計算一下.
主持人:我到五十歲應該可以退休,什麼都不用做了.
主持人:繞著手等其他的弟兄姊妹來服侍我.
主持人:一會兒我們看看是不是這樣.
主持人:可能有人會問.
主持人:既然祭司和利未人都是在會幕當中服侍的群體.
主持人:他們之間有什麼不同呢?.
主持人:如果我們仔細看.
主持人:只有祭司可以接近祭壇.
主持人:能夠獻祭,利未人不可以.
主持人:他們只可以在會幕當中做輔助.
主持人:祭司的角色,辦理裡面各種非祭祀性的事務.
主持人:所以文訴記三章第八節說.

$^{161}$主持人:他們要看守會幕一切的器具.
主持人:為以色列人執行任務辦理帳幕的事項.
主持人:如果我們一路看回之前的經文.
主持人:我們知道所有利未人做什麼的.
主持人:拆台拆燈,都很像我們的托姆.
主持人:搬東西,抬東西,整理好東西.
主持人:這個是他們要做的.
主持人:但從文訴記八章五到十五節.
主持人:我們見到利未人要經歷這個潔淨儀式有幾部分.
主持人:第一個是潔淨之禮.
主持人:然後說吩咐摩西說.
主持人:你要從以色列中選出利未人來潔淨他們.
主持人:你要這樣做來潔淨他們.
主持人:要用潔淨的水澆在他們身上.
主持人:然後叫他們用剃刀剃刮全身.
主持人:洗淨衣服,潔淨自己.
主持人:幸好現在按林牧師不是這樣的.
主持人:因為剃光頭髮了.
主持人:但當時是有這樣的要求.
主持人:這個潔淨之禮有兩部分.
主持人:第一就是潔淨的水澆在身上.
主持人:然後自己用剃刀剃刮全身.
主持人:又洗乾淨衣服.
主持人:這裡的潔淨之水.
主持人:在新漢語譯本翻譯為贖罪的水.
主持人:也是以色列人.
主持人:當他們要行潔淨祭的時候所用的.
主持人:裡面加了紅武牛灰的除污水.
主持人:在文述第十九章有九節會提到.
主持人:就是水裡面加了紅武牛燒出來的灰.
主持人:成為除污水.
主持人:按理說的律法.
主持人:原來人觸摸到死屍之後.
主持人:就要成為不潔之民一共七天.
主持人:如果在這七天之內.
主持人:人進入了回膜.
主持人:就會令到回膜成為不潔.
主持人:所以在這一班預備進入回膜服侍的美人當中.
主持人:不知有沒有.
主持人:可能有.

$^{201}$主持人:曾經在這七天之內接觸過死屍.
主持人:這種嚴重的不潔.
主持人:與神所要求聖潔的生命完全相反.
主持人:不能進入回膜.
主持人:以致污染了回膜.
主持人:所以利未人進入回膜服侍之前.
主持人:所有的利未人都要先進行潔淨之禮.
主持人:免得在不知道的情況底下.
主持人:觸犯了律法當中潔淨的條例.
主持人:中國人常常說什麼.
主持人:不知者不罪.
主持人:神說不行.
主持人:你不知道也好.
主持人:先做這個禮儀.
主持人:免得你進入回膜.
主持人:就添污了這個回膜.
主持人:然後他們要用剃刀剃刮全身.
主持人:洗淨衣服潔淨自己.
主持人:這個是進一步的.
主持人:叫利未人來潔淨他們自己.
主持人:就是他們這個行動.
主持人:是要他們潔除一切可能的皮膚病.
主持人:或者其他的污穢.
主持人:因為摩西利未記提到.
主持人:患了大麻瘋的人.
主持人:在利未記十四章第八第九節是這樣記載.
主持人:求潔淨的人要洗衣服.
主持人:剃去所有的毛髮.
主持人:用水洗燥.
主持人:他就潔淨了.
主持人:然後他可以進營.
主持人:不過仍要在自己的帳外居住七天.
主持人:到了第七天.
主持人:他要剃所有的毛髮.
主持人:頭髮 鬍鬚 眼睛的眉毛.
主持人:他全身的毛都剃了.
主持人:然後他要洗衣服.
主持人:用水洗身才潔淨了.
主持人:這個是當時處理皮膚病的方法.
主持人:所以按潔除一切可能的皮膚病的禮儀都是其中一個潔淨之禮當中的部分.

$^{241}$主持人:做齊了.
主持人:才是為了參與會幕的服侍而完成了潔淨的禮儀.
主持人:我想一下也挺有趣的.
主持人:在以色列回眾裡面.
主持人:突然間有一班人利美族的突然間剃光全身所有的毛髮.
主持人:眉毛也沒有 頭髮也沒有.
主持人:現在做到了.
主持人:我們用AI就可以看看自己的樣子是怎樣.
主持人:如果什麼都沒有的時候 以前就不知道.
主持人:第二 除了潔淨之禮之外.
主持人:我們還要做的就是贖罪之禮.
主持人:文宿的八章八到十節十二和二十節這樣說.
主持人:然後他們要取一個公牛毒以及銅獻的掃祭.
主持人:就是調油的細面.
主持人:你要另取一頭公牛毒作贖罪祭.
主持人:你要帶利未人到會幕前.
主持人:並且要召集以色列全會眾.
主持人:你要把利未人帶到耶和華面前.
主持人:以色列人要為利未人按手.
主持人:利未人要按手在那兩頭牛的頭上.
主持人:留意的是以色列人要為利未人按手.
主持人:然後利未人又按手在那兩頭牛的上面.
主持人:你要將其中一頭作為贖罪祭.
主持人:一頭作繁祭獻給耶和華.
主持人:為利未人贖罪.
主持人:摩西亞倫和以色列全會眾就向利未人這樣做.
主持人:關於利未人反而說怎樣吩咐摩西.
主持人:以色列就向他們這樣做了.
主持人:我們不需要問.
主持人:是否每個以色列人都摸到利未人為他按手.
主持人:我相信那是禮儀式的做法.
主持人:有時我們都會為一個人祈禱.
主持人:不是每個人都碰到那個人.
主持人:可能都是在外面伸起他的手.
主持人:為這些利未人來按手.
主持人:但我們會問究竟按手為甚麼.
主持人:完成了決定之禮之後.
主持人:以色列的領袖作為全以色列人的代表.
主持人:他們要來到利未人面前為他們按手.
主持人:我們會問究竟以色列人為利未人按手是甚麼意思.

$^{281}$主持人:聖經當中按手有兩種不同的意義.
主持人:第一是傳遞祝福和能力.
主持人:第二是傳遞罪.
主持人:即是為那個人祝福.
主持人:將能力賜給他.
主持人:或者將你的罪給了他.
主持人:從這章經文來看.
主持人:以色列人按手在利未人的頭上.
主持人:學者告訴我們.
主持人:應該是使他們的罪轉移到利未人的身上.
主持人:因為聖經記不記得說甚麼.
主持人:利未人的服事是為以色列人贖罪.
主持人:所以聖經學者告訴我們.
主持人:這個很大機會是說.
主持人:這班以色列人將他們的罪轉移到利未人的身上.
主持人:然後利未人將自己好像祭身一樣的獻給神.
主持人:為以色列人贖罪.
主持人:利未人又將他們的手按在贖罪祭的牛身上.
主持人:我們記得勝經告訴我們.
主持人:就是代表他將他的罪轉移到這頭牛的身上.
主持人:所以以色列人按手在利未人的身上.
主持人:是代表利未人要代替全體的以色列人.
主持人:讓他們的罪得到代贖.
主持人:因為聖經說的.
主持人:免得他們靠近回眸而被懲罰.
主持人:當然我們見到利未人的獻祭不同於動物的獻祭.
主持人:他們不是一個死的祭.
主持人:而是一個活的祭.
主持人:他們將他們的生命獻在神的面前.
主持人:當然我們知道利未人本身不是無罪之身.
主持人:今天我們相信耶穌基督.
主持人:當耶穌基督受死為我們贖罪的時候.
主持人:祂不需要再找其他的東西來做.
主持人:因為祂本身就是一個無罪之身.
主持人:但是利未人不是.
主持人:所以為什麼我們見到利未人.
主持人:當以色列人將手按在他們身上.
主持人:將罪轉戒在他們身上的時候.
主持人:他們同樣需要找待贖的牛.
主持人:按照聖經所說.

$^{321}$主持人:一頭作贖罪祭.
主持人:一頭作凡祭獻給耶和華.
主持人:為他們自己來贖罪.
主持人:這經文裡面提到.
主持人:當他們講贖罪的時候.
主持人:提到三種的祭.
主持人:第一是凡祭.
主持人:凡祭在希伯來文有上升的意思.
主持人:就是說透過火去焚燒祭物.
主持人:而產生的煙上升.
主持人:代表了一種輕香的氣.
主持人:上達神的面前讓祂感到喜悅.
主持人:而在獻凡祭的時候.
主持人:獻祭者要將沒有殘疾的公牛.
主持人:公綿羊 公山羊 或者班溝.
主持人:甚至初甲帶到祭壇面前.
主持人:就是說你有錢的話就獻牛.
主持人:沒有的話你可以獻初甲.
主持人:我經常看這經文.
主持人:經常來到這裡.
主持人:幸好新約不用.
主持人:不然這裡就要弄一個祭壇.
主持人:大家不知道帶些什麼來.
主持人:一邊講一邊燒.
主持人:當然我們現在不需要這樣做.
主持人:當時他們就將.
主持人:繁祭就會將這些不同的動物.
主持人:帶到祭壇的面前.
主持人:然後他們按手在繁祭生的頭上.
主持人:代表他們將他們的罪.
主持人:轉介在這個祭生當中.
主持人:祭生就代表獻祭者被殺.
主持人:且被燒在祭壇上為人贖罪.
主持人:但是繁祭不是為了指定的罪.
主持人:和過犯所獻.
主持人:就是說我不是知道我自己做錯了什麼.
主持人:所以來到獻這個繁祭.
主持人:繁祭的目的是潔淨人天良的虧欠.
主持人:就是說我都不知道.
主持人:究竟我有沒有犯罪.

$^{361}$主持人:但是人的罪性裡面.
主持人:會有犯罪的傾向.
主持人:可能真的做了一些事情.
主持人:神不起悅的.
主持人:所以我要透過繁祭表達.
主持人:我對神那種的認罪.
主持人:而繁祭最大的特點就是.
主持人:祭生全獻在壇上.
主持人:毫無保留.
主持人:不會留下的.
主持人:而獻之前要將內臟和腿洗得乾乾淨淨.
主持人:整個獻祭過程.
主持人:象徵人被除去污穢.
主持人:將自己完全奉獻給神.
主持人:一點都不保留.
主持人:繁祭的祭生.
主持人:連祭司都不能吃.
主持人:是要全部燒給神.
主持人:代表我將我自己所有的不潔.
主持人:在神面前完全的來到去獻上.
主持人:經文另外提到跟繁祭一同獻上的有素祭.
主持人:素祭的原意是禮物或者是供物.
主持人:是利未忌我們常常見的五祭當中.
主持人:唯一不需要流血.
主持人:沒有肉類的祭物.
主持人:通常是由細面油和乳香製成.
主持人:代表對神的感恩.
主持人:和將生命獻情.
主持人:可以這樣說.
主持人:類似一個感恩祭.
主持人:然後最後一個就是贖罪祭.
主持人:贖罪祭也是利未忌當中規定的祭祀之一.
主持人:目的在於處理因無心之失或不潔而導致的罪.
主持人:你知道的,但是你不刻意去犯的.
主持人:所以你要在神面前來認罪.
主持人:就是透過濟生的血.
主持人:潔淨你的污穢.
主持人:平息神的怒氣.
主持人:達成人和神關係的修復.
主持人:利未人.

$^{401}$主持人:當他要承接聖職.
主持人:他就要做這三件事.
主持人:為他自己的罪性.
主持人:要感恩.
主持人:為他知道的,但不是刻意犯的罪.
主持人:來到神的面前.
主持人:用禮儀的方式來潔淨自己.
主持人:以致能夠來服侍神.
主持人:完成了潔淨之禮.
主持人:利未人就要歸神所有.
主持人:所以聖經是這樣說的.
主持人:阿倫要從以色列中將利未人奉獻到耶和華面前作為遙祭.
主持人:使他們事奉耶和華.
主持人:使利未人從罪中潔淨自己.
主持人:洗淨衣服.
主持人:阿倫將他們奉獻在耶和華面前作為遙祭.
主持人:又為他們贖罪.
主持人:潔淨他們.
主持人:然後利未人進去在阿倫和他兒子面前.
主持人:在會幕中事奉.
主持人:關於利未人耶和華如何吩咐摩西.
主持人:他們就向利未人照樣做了.
主持人:當利未人完成了潔淨之禮之後.
主持人:又完成了贖罪之禮.
主持人:然後祭司阿倫就根據耶和華的命令.
主持人:用遙祭的禮儀.
主持人:將利未人奉獻給神.
主持人:遙祭原文是特別的貢獻.
主持人:即是一份特別的禮物.
主持人:一份特別的奉賞.
主持人:這個祭的貢物.
主持人:飯祭是燒了的.
主持人:完全不可以留下的.
主持人:但是遙祭這個禮儀.
主持人:搖一搖.
主持人:剩下來的祭物.
主持人:是給祭司吃的.
主持人:不是燒了的.
主持人:所以在聖經利未記第十章十五節.
主持人:所舉的腿.

$^{441}$主持人:所搖的胸.
主持人:他們要與火祭的支柔.
主持人:一同帶來當搖祭.
主持人:在耶和華面前搖一搖.
主持人:只要歸你和你的兒子.
主持人:當作永得的份.
主持人:就是照耶和華所吩咐的.
主持人:即是說當人們拿祭物來這裡.
主持人:做一個搖祭的時候.
主持人:可能是胸部.
主持人:可能是腿.
主持人:可以舉一舉.
主持人:在神面前搖一搖.
主持人:然後就可以吃了.
主持人:當然不是馬上吃的.
主持人:可以留下來不用燒的.
主持人:作為神給這些祭司.
主持人:是他們的養生之物.
主持人:所以搖祭在律法上.
主持人:是屬於平安祭的一部份.
主持人:是祭司應得的份.
主持人:根據聖經所說.
主持人:先要在神面前搖一搖.
主持人:然後祭司就代表耶和華.
主持人:吃那個祭物.
主持人:不是耶和華自己搖的.
主持人:是給祭司吃的.
主持人:所以祭司阿倫在這個時候.
主持人:當他們完成了潔淨之禮.
主持人:就把利未人奉獻在耶和華面前.
主持人:作為搖祭.
主持人:是代表神將利未人賜給祭司.
主持人:成為他們的輔助.
主持人:成為他們的幫助.
主持人:在會務當中參與服事.
主持人:當然我看到這個經文的時候.
主持人:難道阿倫那些祭司.
主持人:逐個利未人抬起他.
主持人:多大隻都不行.
主持人:所以我相信可能是一個儀式.

$^{481}$主持人:就.
主持人:我不知道他用什麼搖一搖.
主持人:可能搖一搖那個利未人的代表.
主持人:這個我猜的.
主持人:總之當時可能用一個禮儀.
主持人:是代表這個搖祭.
主持人:用一個象徵式的方式來進行.
主持人:然後這班利未人.
主持人:就成為了祭司的輔助者.
主持人:一起在會務當中來服事.
主持人:好了.
主持人:去到服事的年日.
主持人:他們都很有興趣.
主持人:又說分布摩西.
主持人:這個是有關利未人的25歲以上的人.
主持人:都要前來任職在會幕事奉.
主持人:到了50歲.
主持人:他們就要從事奉的工作中退休.
主持人:不再事奉.
主持人:不留意後面的句子.
主持人:只可在會幕裡.
主持人:輔助他們的弟兄盡責.
主持人:但又再重複.
主持人:他們自己不再事奉了.
主持人:關於利未人的職責.
主持人:你要向他們這樣做.
主持人:到底說什麼呢.
主持人:我很怕說完之後.
主持人:我們當中50歲以上的弟兄姊妹.
主持人:就開始跟傳道同工辭職.
主持人:就說不關我的事了.
主持人:經文指出.
主持人:利未人潔淨歸神之後.
主持人:他們從25歲到50歲.
主持人:在會幕當中服事.
主持人:到了50歲就退休.
主持人:只做輔助的工作.
主持人:不過如果我們對文述的四章.
主持人:記憶猶新的話.
主持人:你就會發現當中有關神明令利未人.

$^{521}$主持人:的歲數.
主持人:跟這一章有些差距.
主持人:在文述的四章.
主持人:有關利未人工作的年日.
主持人:是由30歲.
主持人:去到50歲.
主持人:我們問當中的五年是什麼呢.
主持人:為什麼第四章就說.
主持人:30歲到50歲.
主持人:這裡就25歲到50歲.
主持人:是不是因為人數少.
主持人:就降低年齡多一些人.
主持人:在聖經裡面.
主持人:任何一個地方告訴我們.
主持人:這個差異的原因.
主持人:但是按照猶大拉比.
主持人:他們做一個估計.
主持人:他們猜測.
主持人:25歲到30歲.
主持人:是一個訓練期.
主持人:就是他們要來到.
主持人:跟著已經在參與.
主持人:事奉的那些利美的同胞.
主持人:看一下他們怎樣做.
主持人:你想一下拆那些圍板.
主持人:怎樣弄好它.
主持人:都不是那麼容易.
主持人:又拆又開.
主持人:原來不是每個人都懂拆懂開.
主持人:你都要知道怎樣拆怎樣開.
主持人:不會弄傷自己.
主持人:又能夠讓那些板可以用久一點.
主持人:不會那麼快壞.
主持人:所以拉比就說可能很大機會.
主持人:就是他們25歲就要開始做學師.
主持人:30歲就正式.
主持人:參與立起拆卸.
主持人:降台搬運的工作.
主持人:都有些道理的我相信.
主持人:都很大機會是這樣.

$^{561}$主持人:但是利未人.
主持人:世世代代都做到50歲嗎.
主持人:從民數記三到四章.
主持人:我們明白利未人在抗疫時代.
主持人:主要的工作.
主持人:剛才用了四個字眼.
主持人:就是負責立起拆卸帷幄.
主持人:和在行走的路途上.
主持人:就是神要他們起行他們就要起行.
主持人:於是他們就要拆卸那個帷幄.
主持人:在行走的路途當中.
主持人:負責降台搬運帷幄所有的器物.
主持人:當然後來我們知道.
主持人:因為他們奉獻了有些車有些馬.
主持人:有些是可以用這樣的方式去車載.
主持人:但是帷幄裡面的.
主持人:藥櫃聖壇那些全部是要人抬的.
主持人:所以可以這樣說.
主持人:利未人所做的.
主持人:特別是在抗疫時代.
主持人:是屬於勞力的工作.
主持人:所以神為他們定下了工作的時限.
主持人:就是去到50歲.
主持人:可能體力下降了.
主持人:因此他們就退下來.
主持人:不過他們退下來.
主持人:就是去到幫忙指點他們.
主持人:或者需要一些協助的時候.
主持人:都會做一些功夫.
主持人:不過他不是扮演最主要的角色.
主持人:但是後來帷幄.
主持人:當大衛王將帷幄永久設立在耶路撒冷.
主持人:和後來聖殿永久建立之後.
主持人:利未人的職責就有了改變.
主持人:因為他們不再需要做那麼多艱苦的工作.
主持人:不需要經常拆台搬運帷幄.
主持人:因此在歷代至上23章.
主持人:我們就見到大衛王就指示利未人從20歲開始.
主持人:就在聖殿當中服侍.
主持人:哎呀 又早了.

$^{601}$主持人:但是就沒有寫退休年齡.
主持人:是不是漏了?沒理由的.
主持人:因為真的沒有記載.
主持人:好了 對於利未人在預備進入帷幄服侍之前.
主持人:先以潔淨和贖罪的禮儀.
主持人:來處理可能犯了律法中不潔的條例.
主持人:又或者犯了不是刻意違反律法的罪行.
主持人:其實是在提醒我們每一個服侍主的人.
主持人:在我們每天的生活當中.
主持人:我們有沒有分別為性的心態來服侍神.
主持人:又或者將我們再擴闊一點點.
主持人:今天你和我都是神所呼召的祭司.
主持人:今天在我們每天的生活當中.
主持人:我們是不是常常都刻意留意.
主持人:我自己是不是分別為性來過我每天的生活.
主持人:是的 在新約時代.
主持人:神沒有刻意要求我們在服侍祂之前.
主持人:先做這些潔淨之禮或贖罪之禮.
主持人:不過使徒保羅在一節我們很熟悉的聖經裡面.
主持人:卻提醒我們.
主持人:「弟兄們,我要神的慈悲勸你們.
主持人:將身體獻上當作活祭.
主持人:是聖潔的,是神所喜悅的.
主持人:你們如此事奉乃是理所當然的」.
主持人:已故的一位基督徒作家叫狄賢.
主持人:我不知道我們當中有沒有人讀過他的書.
主持人:他曾經這樣來講這段經文.
主持人:他說既然是活祭的意思就是.
主持人:我可以選擇在那個祭壇裡面下來.
主持人:就是說我要用意志告訴我自己.
主持人:我是真的甘心情願一生將我自己擺在這個祭壇上.
主持人:以致我真的能夠成為神呼召我做的那個祭師.
主持人:我知道的.
主持人:我可以難處的時候可以下來.
主持人:開心的時候就留在那個壇上.
主持人:但這次聖經和你未人將自己當作活祭獻上.
主持人:在會幕當中服侍耶和華.
主持人:有異曲同工的意義.
主持人:所以我們常常問我們自己.
主持人:在我日常生活當中.

$^{641}$主持人:我有沒有在意自己的生活是否合神的心意.
主持人:是否配得上服侍神的身份.
主持人:你有沒有依靠聖靈的幫助.
主持人:你的思想,眼目,行為不偏離神的教導.
主持人:神說我們每一個都是祭師.
主持人:意思就是說我們每一個都是應該服侍祂的門徒.
主持人:不知道我們上主教多久.
主持人:你在教會裡面有沒有服侍的崗位.
主持人:正如你未人一樣.
主持人:他們服侍的重心不是在獻祭方面.
主持人:而是在會幕其他的需要上.
主持人:同樣地在教會當中.
主持人:你的服侍可能和我們的全套童宮有點不同.
主持人:不過在教會裡面仍然有各種各樣的服侍崗位等著你.
主持人:你進來的時候有師士會將程序表給你.
主持人:歡迎你.
主持人:有弟兄在敬拜的崗位當中.
主持人:用詩歌帶領我們來敬拜神.
主持人:或者用詩班的形式來代表我們將.
主持人:咀唇的祭向神獻上.
主持人:有影音幫助我們看得清楚.
主持人:以致我們能夠了解.
主持人:有組長來牧養我們.
主持人:幫助教務童宮.
主持人:帶你查經.
主持人:關心你.
主持人:在你病的時候為你禱告.
主持人:在你有需要的時候來關心你.
主持人:茶水那裡給你倒咖啡.
主持人:給你一件蛋糕.
主持人:讓一切有序地來進行.
主持人:有多個堂務.
主持人:抹乾淨所有的東西.
主持人:擺好所有的桌子.
主持人:擺好所有的椅子.
主持人:讓你來到的時候.
主持人:你很享受在裡面服事聚會的過程.
主持人:如果你現在沒有服事的崗位.
主持人:我鼓勵你好像舊約的利未人一樣.
主持人:在教會裡面找一個你服事的崗位.

$^{681}$主持人:問一下神 究竟祂想你在祂的角度裡面.
主持人:在哪個位置去事奉.
主持人:我常常這樣說.
主持人:教會裡面所有的服事都是珍貴.
主持人:都是寶貴.
主持人:都是一樣的重要.
主持人:有時候我們很容易看.
主持人:這個服事很重要.
主持人:這個崗位很尊貴.
主持人:這個位置高一點.
主持人:這個低一點.
主持人:對不起 這是人的錯誤.
主持人:只要我們願意善用神給我們的恩賜.
主持人:在神給我們的位置.
主持人:那個就是最美好的服事的崗位.
主持人:另外我們有機會聽到弟兄姊妹說.
主持人:我一開始都問了這個問題.
主持人:我在工作上面退休了.
主持人:所以我在教會上面的服事都退休了.
主持人:有一次我真的跟一個人.
主持人:不要問我是誰.
主持人:有一個人聊天的時候.
主持人:他真的跟我說.
主持人:剛才我退休了.
主持人:所以教會事奉我也給你了.
主持人:你給我 我沒用的.
主持人:他真的在引據點.
主持人:他真的在引文數據.
主持人:八章二十五節.
主持人:他說到你五十歲.
主持人:他們就要從事奉的工作中退休.
主持人:不再事奉.
主持人:去印證他的退休是合聖經的.
主持人:不過正如剛才上文所提.
主持人:我們知道當時利未人所做的是.
主持人:搬運鋼台的勞動工作.
主持人:到了會幕固定.
主持人:設立耶路撒冷.
主持人:和在聖殿時期.
主持人:利未人就沒有退休的年齡.

$^{721}$主持人:是的 我們可能不再年輕.
主持人:我們不再做搬運鋼台的工作.
主持人:以免弄傷腰部或影響身體.
主持人:但我相信在教會當中.
主持人:有一個事奉我們都做得到.
主持人:有一個代禱的事奉.
主持人:無論你到什麼年紀都做得到.
主持人:教會裡大大小小的事.
主持人:在教會裡服侍的傳道童,公,牧者.
主持人:文職,堂務都需要大家的禱告.
主持人:因為所有在教會裡.
主持人:你的組長都需要來祈禱.
主持人:因為所有在教會裡服侍的人.
主持人:比你更近地影響撒旦的攻擊.
主持人:他需要你在後面用禱告.
主持人:托住這班事奉者.
主持人:所以我相信.
主持人:聖經裡面說.
主持人:是的 可能去到某個年紀.
主持人:因為身體能力或體力的緣故.
主持人:我們需要在某些崗位裡退下來.
主持人:但聖經說我們可以做輔助他的弟兄.
主持人:聖經就說在會幕裡.
主持人:輔助他們的弟兄來盡他們的責任.
主持人:所以民俗記八章二十五節的退休定義.
主持人:不是說他們什麼都不做.
主持人:而是要從一個帶領者的角色.
主持人:或者由有名字的崗位的角色.
主持人:退下來.
主持人:轉而成為一個輔助者的角色.
主持人:不代表輔助者不重要.
主持人:不代表輔助者只是一個閒角.
主持人:你想想 港福堂這麼大.
主持人:如果只是靠.
主持人:大約有四十個教務同工,文職,堂務.
主持人:整間教會這麼大.
主持人:只靠這四十個人.
主持人:我相信很多事情都解決不了.
主持人:就算你覺得你是一個輔助者的角色.
主持人:你也扮演著一個非常重要的角色.

$^{761}$主持人:和我們一起在神教國裡努力.
主持人:各位弟兄姊妹.
主持人:讓我們以你美容來榜樣.
主持人:注目我們自己是否過了一個聖業的生活.
主持人:在教會裡尋找一個服事的崗位.
主持人:甚至在我們年齡漸長的情況下.
主持人:在服事上面不言休.
主持人:而是轉換我們的角色.
主持人:繼續用神賜給你的金靈歲月.
主持人:在教會裡尋找適合的崗位.
主持人:願意我們一生都服事我們的神.
主持人:我們一起禱告.
主持人:傳領上帝我們來到你的面前.
主持人:因為是你讓我們可以服事你.
主持人:最後當我們看你美人.
主持人:他們進入會幕之前.
主持人:他們先潔淨他們自己.
主持人:最後就提醒我們.
主持人:今天我們是不是每天都小心謹慎.
主持人:敏銳於我們自己.
主持人:是不是過一個聖潔的生活.
主持人:我們看到利未人.
主持人:縱然他們只是在會幕當中.
主持人:不是負責獻祭.
主持人:而是負責很多其他的事務.
主持人:但主要我們知道會幕.
主持人:沒有他們會幕都不成會幕.
主持人:因為他們要立起來.
主持人:要搬運.
主持人:要做好一切的事.
主持人:以至祭司才能夠來獻祭.
主持人:做好他們應當做的工作.
主持人:我們就明白.
主持人:無論是祭司還是利未人.
主持人:其實都在會幕當中不可或缺.
主持人:將你的安排是美好的.
主持人:因為我們每一個人原來都可以.
主持人:按著你給我們的呼召.
主持人:用你給我們的恩賜.
主持人:無論我們在什麼年齡當中.

$^{801}$主持人:我們都可以服侍你.
主持人:幫助我們.
主持人:讓我們每一個在你面前的人.
主持人:都能夠尋找到你給我們服侍的崗位.
主持人:以至我們帶著聖潔的生命.
主持人:彼此的服侍.
主持人:互相的輔助.
主持人:去成就你自己角度裡面要成就的工作.
主持人:我們這樣禱告奉耶穌基督的命求.
主持人:阿們.
\newpage



\section{詩篇 95:1-11}
\label{sec:031904}
\textbf{洪水中的約——安睡在浮動的動物園 (英/粵)}~~~~歐世達牧師
\newline
\newline
連結: \href{http://www.church.com.hk/acms/content.asp?site=cdc\&op=show\&type=product\&code=031904\&layout=sermon}{\texttt{031904}} ~~~~ 語音日期: 2026-03-29
\newline
\newline
\hyperref[sec:031903]{\small{< < < PREV SERMON < < <}}
~
\hyperref[sec:index_chronic]{\small{[返順時目]}}
~
\hyperref[sec:031905]{\small{> > > NEXT SERMON > > >}}
\newline
\newline
詩篇 95:1-11
\newline
\begin{longtable}{cl}
\hline
\hline
章節 & 經文 (和合本修訂版)\\
\hline
95:1 & \begin{tabularx}{0.7\textwidth}{X} 來啊，我們要向耶和華歌唱，向拯救我們的磐石歡呼！ \end{tabularx} \\ \\ \relax
95:2 & \begin{tabularx}{0.7\textwidth}{X} 我們要以感謝來到他面前，用詩歌向他歡呼！ \end{tabularx} \\ \\ \relax
95:3 & \begin{tabularx}{0.7\textwidth}{X} 因耶和華是偉大的神，是超越萬神的大君王。 \end{tabularx} \\ \\ \relax
95:4 & \begin{tabularx}{0.7\textwidth}{X} 地的深處在他手中；山的高峰也屬他。 \end{tabularx} \\ \\ \relax
95:5 & \begin{tabularx}{0.7\textwidth}{X} 海洋屬他，是他造的；旱地也是他手造成的。 \end{tabularx} \\ \\ \relax
95:6 & \begin{tabularx}{0.7\textwidth}{X} 來啊，我們要俯伏敬拜，在造我們的耶和華面前跪拜。 \end{tabularx} \\ \\ \relax
95:7 & \begin{tabularx}{0.7\textwidth}{X} 因為他是我們的神；我們是他草場的百姓，是他手中的羊。惟願你們今天聽他的話！ \end{tabularx} \\ \\ \relax
95:8 & \begin{tabularx}{0.7\textwidth}{X} 你們不可硬著心，像在米利巴，就是在曠野瑪撒的日子。 \end{tabularx} \\ \\ \relax
95:9 & \begin{tabularx}{0.7\textwidth}{X} 那時，你們的祖宗試我，探我，並且觀看我的作為。 \end{tabularx} \\ \\ \relax
95:10 & \begin{tabularx}{0.7\textwidth}{X} 四十年之久，我厭煩那世代，說：「這是心裡迷糊的百姓，竟不知道我的道路！」 \end{tabularx} \\ \\ \relax
95:11 & \begin{tabularx}{0.7\textwidth}{X} 所以，我在怒中起誓：「他們斷不可進入我的安息！」 \end{tabularx} \\ \\
[1ex]
\hline
\hline
\end{longtable}
$^{1}$大家早安,再一次有個機會來和大家分享.
每一個月來到的時候都會提醒我季節的轉變.
上一次我來的時候是一個很嚴寒的冬天,現在是春天已經來了.
我今天都會略略提到神和季節的一些事情.
無論那個季節是怎樣,神對我們的信實是不改變.
讓我們一同來祈求.
我們親愛的父神,多謝你對我們的信實.
求你幫助我們明白怎樣向你有忠誠.
所以無論這個世界在什麼季節裡面.
或者你是怎樣呼召我們去做的事情.
我們都是用公義來回應你.
引導我們進入你自己的華語.
奉主耶穌的聖名祈求.
我們首先溫習一下約是什麼意思.
約是有兩方面的意思.
一件事就是在約裡面有不同的立約的人.
第二件事就是這兩者之間所需要履行的義務.
所以我們看到有人有兩方面和他們需要的義務.
有一個很舊的電影叫做回家.
這個是關於一個家庭他們有一些寵物.
他們有兩隻狗一隻貓.
他們需要帶動這些寵物去一個很遠的沙漠.
但是那些寵物似乎對他們的主人有一些義務.
所以這兩隻狗和貓就離開了那座農場.
這是一個很戲劇性的故事.
是關於那些寵物要走多遠的路去履行他們的義務.
雖然這跟約沒有關係.
但是這兩者之間的關係還是相似的.
似乎不是完全講到約的意思.
但是我們看到在這個關係裡面有兩者和他們需要的義務.
在詩篇95篇的時候是講到我們神的關係.
最基本的本質就是祂是我們的創造主.
我們是祂的創造物.
所以在這個關係裡面神是居於高的地位.
我們首要的回應就是需要去敬拜祂.
我們覺得在這些創造和造物的應有固有的關係.
甚至一些動物也是按著祂的本性來敬拜神.
如果是順著本性去做的.
本身就是在榮耀神.
神和亞當和夏娃所立的約.

$^{41}$所以特別是因為這個關係有進一步的界定.
在第七節講到我是神的子民也是祂的羊群.
這就是神和亞當和夏娃在伊甸園所講到的東西.
亞當和夏娃能夠一直保守自己的神的子民,自己的羊群的話.
神就會看顧他們.
祂會給他們安排一些好的草場和會保佑他們.
可能在我們創造的本性中.
有些東西是不在約的範疇裡面.
因為我們是按著神的形象而被造.
所以我們和神有很多地方的相似.
當我們工作的時候很像神的工作.
當我們被神設計成為有那個關係.
神也是和人有那個關係.
這個是我形容祂是一個本性.
我們自己人的本性.
就算有些不認識神的人.
如果他按著神做他的本性來生活的話.
他也是在榮耀神.
但在我這個約中的關係是有些特別的事情我們需要做.
重點就是我們怎樣去做和怎樣去建立這些關係.
我們會盡力來按照神所盼望我們所要做的事去做.
但我們知道亞當和夏娃最終是失敗了.
他們沒有停留在成為神的羊群的一部分.
你可以這樣想.
如果亞當和夏娃在神自己的草場上的時候.
魔鬼引誘了夏娃.
嘗試欺騙她離開神的草場.
跟她說是不是在別的地方的草場更加好呢.
為什麼神會不容許你去那些草場呢.
所以夏娃就吃了那個禁果.
亞當看見這個事情發生的時候.
他看見一些不同的觀點.
亞當知道自己在做什麼.
當他吃了那個禁果的時候.
他知道他在不依賴神.
他在選擇耶穌而不是選擇神.
這就是為什麼世界上的罪.
落在亞當而不是在耶穌身上.
所以換句話說.
我們得到的罪孽是從我們的父親而來.

$^{81}$所以在天主教裡面.
因為我們的罪從亞當而來.
所以在瑪利亞的身上.
亞當的罪是從耶穌身上而來的.
所以在天主教裡面.
亞當的罪是從耶穌身上而來的.
所以在天主教裡面.
亞當的罪是從耶穌身上而來的.
所以在天主教裡面.
亞當的罪是從耶穌身上而來的.
所以在瑪利亞的身上.
和我們的關係是不同的.
有一些事情是必須要做.
來將這些罪惡除去.
因為那個罪孽.
再不是我們本身.
成為我們本身所做的一部分.
當我們根據那個故事.
從創世到羅亞洪水來的時候.
有些事情會發生.
我們可以根據一個公義的時代.
直到羅亞來.
還有在這段歷史裡面.
我們也看到罪惡的澎湃和擴張.
所以直到一個地暴神.
覺得他必須要做一些事情.
他是通過洪水來毀滅這個世界.
所以這個是帶動我們今天所要看的經文.
我們不需要用太多的言語.
就可以描述到人類在洪水前所面臨的處境.
在創世記第六章第五節那裡.
他說看到地上的罪惡很大.
終日所思想的盡都是惡.
就是說終生所想的都是惡的.
這是一個很嚴峻強烈的描述.
你要記住這個事情.
當我們繼續看以下四節的經文.
我們回到剛才所看的第五節.
現在我們回到第一節.
在這段經文裡面.

$^{121}$是比較難處理令人費解的一段經文.
就是當人世上多起來的時候.
有生兒女的時候.
神的兒子看見人的女子美貌.
就隨意挑選娶來為妻.
這個提到是神的兒子.
娶了人的女兒為妻.
到第四節那裡說到.
他們結婚生子之後.
生了的後代成為上古英武有名的人.
可能這幾句話挑起我們一些想法.
這些事情發生在洪水之前.
只有八個人是沒有被洪水淹沒.
我們很相信他們都會常常說.
在洪水未來之前所發生的一些事情.
我們無法知道這些故事.
只是有一些提引.
有時候我覺得可能.
這個洪水可能在中國的文化.
中國的字裡面表達了出來.
有八個人在一艘船上面.
無論如何.
可能這些是說到以前.
在洪水那段時候.
有些什麼事情發生.
你可能會提到那些人是.
上古有名的人是什麼人呢?.
他們說到兩個明星.
Michael Jordan可能比Lebron James.
打籃球打得更加好.
在這段時間那些人叫做拿飛利人.
是不是他們是從上下來的一些天使呢?.
是不是他們是一些王子貴族.
是吃得好穿得好.
所以他們更加威武壯大呢?.
或者有些人說.
有一個可能就是他們是一些邪教.
是行一些淫亂的事情.
可能這些都是有可能的一些解釋.
這個可能是繪出一個圖畫.

$^{161}$就是那時候是有些事情.
是不好是最邪惡的事情發生.
好像第二節那裡說到.
所有的人都會隨意按著.
他們所想的事情來做事.
如果你看到一些你喜愛的東西.
你就會隨意來奪取.
無論不理人的.
不顧忌人的想法.
你要直接去做.
他們所做的事情.
是和神的心意去關懷.
去培養那種概念完全沒有相反.
這個去到一個程度.
神說我必須要來動手做一些事情.
除了是關於納菲他利人.
還有其他的話語來描述當時的事.
他描述神的時候.
是講到神有很深的情感.
就如在第七節的尾那裡.
他後悔他所做的事.
有時候我們會想.
我們這個神是否反覆無定呢.
其實這個其中一個部分.
就是講故事的一些技巧.
但是神本身是知道所有的事情.
無論是以往的,現在的,將來的.
但是他願意在這個時間的流裡面.
來與我們同在.
他可以在現今的時代.
在這個時間.
看見我們所犯的罪.
而感到悲哀.
有時候作為一個父親.
我好像一個水煲一樣.
有時候小孩子開始做一些.
我認為他們不應該做的事情.
有時候你放手讓他們做.
看他們有什麼事情發生.
但是他們會很頑固地不斷地做.

$^{201}$直到你忍不住.
有可能是神的描述.
就好像作為一個父母一樣.
但是他自己的忍耐是何等的長.
直到一個地步.
事情是壞到一個地步.
神說我需要來作供.
在這個時候神做了一件事.
就是限制人可以活多久.
他將他們在在世能夠活的年日.
用120年為限.
有可能會令我們想.
究竟以前的人是否真的活了這麼久.
我自己個人認為有這樣的可能性.
因為洪水影響到地和所有大氣的環境.
可能會影響到人的性命.
比如說可能影響到當時的養的供應.
所以究竟全世界是否真的有這麼大的洪水發生.
也是一個很擇手的問題.
有一件事就是神發現.
當人的年齡增長的時候.
他自己是越來越壞.
可能對我們年紀大的人是一種警告.
我不知道大家當中有多少人已經是年長.
有時候覺得沒所謂.
但是當我們越是年紀大的時候.
越覺得有這樣的警告.
神在這個情況下.
將人的壽命縮短.
可能是對我們的一個警告.
在第八次的時候.
那個故事開始轉變.
要看到羅雅在神的眼中蒙恩.
我們發覺當我們根據那個代的話.
根據那個義的後裔的話.
有時候故事是比較難處理.
因為它是用口來傳.
因為講到羅雅是一個義人.
與神同行.
羅雅是對神忠心.

$^{241}$與神同行.
然後也是描述了人類的罪惡.
這個講到當時大地也充滿了強暴.
讀回創世記以前.
其中有一個人物叫拉麥.
這個人比其他舊的聖經人物更加強暴.
所以到了第十三節的時候.
神決定要毀滅地上一切的生命.
然後用羅雅作為一個新的開始.
在這個洪水裡面.
神做了兩件事.
第一件事就是要懲罰那個罪惡.
這個是因為亞當和夏娃所犯的罪.
雖然神有能力.
當亞當和夏娃犯罪的時候.
可以令到他們死亡.
但是神表達了他自己的內心.
而且是不輕易地發怒.
但是直至罪惡去到一個地步.
那個殘暴這麼壞的時候.
神決定來施行他的公義.
他最終來實行.
對於邪惡的實行他的公義.
但這些事情不是因為洪水來到就停止.
所以一次又一次的.
在整本聖經裡我們看到.
神是用同樣的方法來處理.
讓罪惡來到存在.
直到一個地步.
他需要施行他的懲罰.
在這裡我是用季節來形容.
神有不同的季節.
有容許的季節.
也有施行審判的季節.
或者我們回到詩篇第95篇.
我們可能看到我們所形容的一些東西.
這個是形容當時以色列人.
未進入他應許之地之前.
在這裡也是說到.
神對人的忍耐.

$^{281}$也是對人所發出的怒氣.
我們也可以從新約裡面.
看到類似的經文.
在彼得後書第二章裡.
他們提到一串烈神懲罰的事情.
甚至是懲罰的天使.
我們在這裡的經文看到.
洪水和索多瑪俄羅摩拉的故事.
他不單止是說到.
他如何來懲罰以前的世代.
他說到這些事情提到.
主要是作為人不敬虔的人的一種尷尬.
神雖然是有忍耐.
但他仍然會懲罰罪惡.
如果你不太相信彼得所寫的說話.
你也可以去到耶穌所講的說話.
讓我讀一讀.
馬太福音講到.
在那日子那時辰.
沒有人知道.
連天上的使者也不知道.
子也不知道.
唯獨父知道.
羅亞的日子怎樣.
人子降臨也要怎樣.
當洪水以前的日子.
人朝上吃喝嫁娶.
直到羅亞進方舟的日子.
不知不覺.
洪水來了.
把他們全都沖去.
人子降臨也是這樣.
我曾經和大家提到.
關於啟示錄的東西.
當我們去到末世的時代的時候.
神是會完完全全的.
徹底一次的來懲罰作惡的人.
我們提到神是這麼忍耐我們.
但祂讓所有的事情.
一路來惡化.

$^{321}$直到一個地步.
去彰顯祂自己的公義.
很清澈明瞭地來實行祂的公義.
甚至在今天我們這個時代.
我們觀察這個政局的時候.
我們都會說.
這個已經到了一個極惡的地步.
需要有些事情要改正來發生.
有時候那些人.
其他人可能會叫政府來做一些事.
但是作為一個信徒.
我們要呼喚神來加入這些事情.
第一件事就是.
神在羅亞這個故事裡面.
就是要懲罰那個惡人.
他告訴我們一個例子.
這件事情就是這樣來發生.
第二件事他做的就是.
他重新按照一個公義的家庭.
來重新開始.
神會是懲罰罪惡.
但是祂會保存忠信的人.
我們覺得在以西結書所指出的事情.
是很有意思的.
因為我們看到.
在當時的當中.
如果當中有異人的時候.
神會因著這些異人來保存其他人.
就如耶穌說到的一些信徒.
就是世上的異人.
世上的異人有機會來保存當時的世代.
但是在以西結書一段的經文裡面.
他提到.
羅亞,但以理和約伯.
他說當時的情況是.
最大惡極.
甚至他們不會因為這些異人來保存他們的性命.
他們知道那些信他們的人得到保存.
但是不包括他們的兒女.
當我們說到羅亞的兒女的時候.

$^{361}$是有一點意思的.
他不一定是因為這些兒女的公義來拯救他們.
因為他們是羅亞的兒女.
而是說他們是藉著這些留下的人.
讓這個世界重新開始.
我們進入一個整體結構的故事.
整經整體結構的故事裡面的一部分.
那個大題目就是.
神怎樣來處理世上的邪惡.
所以在洪水裡面說到其中一個可能的答案.
如果神在滅世上所有的邪惡重新開始的時候.
會有什麼後果.
是不是因為重新開始就能夠所有的東西都好呢.
答案是不是.
我們可能是可以利用那個Pixar.
皮克斯的一種講故事的方法.
來想一想關於這個聖經的故事.
皮克斯這個公司很懂得怎樣去講故事.
他會教你怎樣好好來寫一些故事.
可能我們所講的就是一個講故事的模板.
譬如說我們講一個故事有一個人物叫Sam.
Sam做同樣的事.
有一天他所做的事不妥.
有一個問題.
Sam會用一個方法來嘗試去解決那件事.
但是沒有用.
所以Sam又用另外一個方法來處理.
又不妥.
再來一次.
再次嘗試.
直到有一天.
當他有一些改變.
他終於能夠改正那個問題.
直到他每件事都能夠順利地做出來.
我們可以將聖經的故事擺入一個類似的模式.
我們看到東夏灣每日與神在伊甸園裡面同行.
但是有一天他們兩個吃了禁果.
事情就改變了.
怎樣去修改這個問題呢?.
其中一個可能性就是洪水.

$^{401}$如果神重新開始利用一個義人羅亞的話.
重新開始會不會有些不同呢?.
但是答案是不能夠.
我們知道神知道這樣的做法是不能夠改變事情.
但他仍然要藉著洪水來告訴我們.
這個不能夠.
我們知道所有的洪水.
沒有辦法將這些人重新帶到神創造.
起初創造他們的形態.
所以提醒了我們知道.
我們做某一件事就能夠改變事實.
我們看到Pixar其中一個故事.
其中一個人物在故事內容裡面有不斷的改造.
我們不單止看到神在曹擱.
我們法國人在整個過程裡面有所改變.
我們不單止看到神做一個重啟.
而是利用羅亞和他的家人做一個重新的開始.
但是這個不同的地方.
和亞都夏娃的約有所不同.
神是有一個供應讓羅亞和他的家庭能夠生存下來.
就算是保存所有的動物也是神的約的其中一部分.
你可不可以想到當神要羅亞建造方舟的時候.
的意思.
不單止是羅亞要救他自己.
而是羅亞參與神建造方舟這件事.
而且更加學習到如何維護這些神所創造的其他的萬物.
所以不單止是他要和子女建立的關係.
當羅亞建造方舟的時候.
建立不同的單位的時候.
他們會考慮到動物的需要.
動物園可能不太適合動物居住.
但是他們在管理的其中我們可以學到很多東西.
在建立建造的期間.
羅亞學了如何去愛惜這些動物.
而在建造的過程中整個家庭也有參與.
我還記得有一次我提到.
如何去運作這間教會才能最有效地做.
如果是每一個團體自己來管理的話.
效果會比較好嗎?.
非傳道提醒我一件事.

$^{441}$是不是說每一個團體或每一個部份.
互相合作是否更加好呢?.
我們需要跨越不同的格局來建立一個教會.
在某個程度來說.
羅亞的每一個家庭分子都要跨越這個格局.
來互相合作做這件事.
愛護這個世界來建立一些東西.
我的女兒快要離開家庭進入學校.
可能方舟會成為羅亞子女的一所學校.
讓他們有機會去學習.
方舟是神在暴風雨中對羅亞的一些供應.
而羅亞是在神的工作中得以安息.
在聖經裡我們看到救贖是神的一種安息.
而是我們最終能夠在神的裡面得到安息.
在詩篇95篇裡講到神的救恩被描繪為我們的安息.
就是我們最終的信靠.
我們想一想羅亞如何在方舟裡面睡覺.
睡在一個浮動的動物園.
動物很臭.
外面的暴風雨是狂風暴浪.
這個巨大的方舟在波濤上上下顛覆.
但羅亞和他的家人在這個情況下安然入睡.
他們是安息在神自己的供應裡面.
中暑節是講到耶穌進入耶路撒冷.
我們不知道當時的群眾看到耶穌進入耶路撒冷時在想什麼.
但如果他們大聲呼籲紮理神進入耶路撒冷是應當稱頌的.
所以他們知道耶穌本身就是一個供應.
但耶穌是否會帶領他們擊敗羅馬人呢?.
是否耶穌會帶領他們進入一個更繁榮的景象呢?.
但無論他們怎樣想他們都認為耶穌是對當時的人的一種供應.
在彼得所寫的書信裡再次講到挪亞方舟是神的作為是一個很重要的工具.
彼得講到羅馬是藉著水流得來的.
在洪水裡面有它的特定功用.
它是要將世上的罪惡洗滌乾淨.
讓這個世界重新有機會從公義重新開始.
他講到洪水和方舟是和我們的浸禮有相似的地方.
我們的浸禮是需要經過水.
是很描述我們藉著耶穌基督能夠將我們的罪惡完全洗滌乾淨.
它清理了我們心裡面的良心讓我們能夠重新開始.
講到我兒子他現在十四歲有些身體的味道.

$^{481}$所以我女兒不在我講一下也無所謂.
很奇怪為何有些小朋友不喜歡洗澡.
因為洗澡的話將你其他的污穢味道完全洗掉.
神講到我們每一個人都有那種醜陋的地方.
唯一能夠洗淨這些就是在耶穌基督裡面將這些東西完全洗到乾淨.
我們以為我們像諾亞一樣能夠被拯救.
但我們卻覺得自己就像周圍的人一樣.
被雨淋淋而被洗淨.
我們必須想像我們需要洗澡.
另一方面,有時醫生幾乎沒有洗乾淨自己的手.
我們都聽過.
有些醫生可能洗自己的手洗得不太乾淨.
因為很多的細菌是我們眼睛不能夠看到的.
雖然神講到祂將我們洗得乾淨.
我們仍然需要每天來洗我們的手,將我們剩下的污穢洗乾淨.
這些所有的圖畫都是描述神如何來潔淨我們.
他提到當我們受洗的時候,神洗去我們一切的罪孽.
當我們從水裡出來的時候,就描述我們已經被潔淨,可以重新開始.
就好像諾亞從方舟出來的時候有一個重新開始的機會.
在中時節,我們不知道耶穌是怎樣的.
祂的樣子是怎樣的.
有機會放假的時候,他去巴西那裡有一個很大的十字架牆.
我們可以看到耶穌的雙手張開來歡迎我們.
但這讓我經常在想,耶穌的樣子是怎樣的?.
讓我給你們看一個奇怪的圖像,看看耶穌在聖誕節是怎樣的.
看看耶穌究竟是怎樣的,特別是在中時節的日期.
當你看到耶穌進入耶路撒冷的時候.
你可以形容,耶穌進來的時候,就好像一艘方舟那麼大的船.
耶穌本身就是我們的供應,能夠得到拯救.
祂是我們唯一能夠倚靠的神.
我希望你們每個星期日有休息的時間.
希望羅亞在方舟裡面的那種安息,而我們本身是在基督耶穌裡面的那種安息.
讓我們每次休息的時候,我們能夠在主裡面有安息.
就如我們睡覺,讀書,或者看一些東西.
我們一個屬靈的休息就是安息在神自己的救贖裡面.
讓我們一起祈禱.
我們親愛的父神.
多謝你愛我們,對我們有很多的供應.
我們多謝你,耶穌基督就是你給我們最好的供應.
我們原本應該是受到你的審判.

$^{521}$我們多謝耶穌基督,將這種審判來洗去.
幫助我們學習如何在基督裡面安息.
讓我們能夠從其中重新開始.
藉著聖靈來做我們當做的事.
和主耶穌鳴求.
\newpage



\section{約翰福音 12:12-33}
\label{sec:031905}
\textbf{當你高唱迎接榮耀君王}~~~~周曉暉牧師
\newline
\newline
連結: \href{http://www.church.com.hk/acms/content.asp?site=cdc\&op=show\&type=product\&code=031905\&layout=sermon}{\texttt{031905}} ~~~~ 語音日期: 2026-03-29
\newline
\newline
\hyperref[sec:031904]{\small{< < < PREV SERMON < < <}}
~
\hyperref[sec:index_chronic]{\small{[返順時目]}}
~
\hyperref[sec:031906]{\small{> > > NEXT SERMON > > >}}
\newline
\newline
約翰福音 12:12-33
\newline
\begin{longtable}{cl}
\hline
\hline
章節 & 經文 (和合本修訂版)\\
\hline
12:12 & \begin{tabularx}{0.7\textwidth}{X} 第二天，有一大群上來過節的人聽見耶穌要來耶路撒冷， \end{tabularx} \\ \\ \relax
12:13 & \begin{tabularx}{0.7\textwidth}{X} 就拿著棕樹枝出去迎接他，喊著：「和散那，以色列的王！奉主名來的是應當稱頌的！」 \end{tabularx} \\ \\ \relax
12:14 & \begin{tabularx}{0.7\textwidth}{X} 耶穌找到了一匹驢駒，就騎上，如經上所記： \end{tabularx} \\ \\ \relax
12:15 & \begin{tabularx}{0.7\textwidth}{X} 「錫安的兒女啊，不要懼怕！看哪，你的王來了；他騎在驢駒上。」 \end{tabularx} \\ \\ \relax
12:16 & \begin{tabularx}{0.7\textwidth}{X} 門徒當初不明白這些事，等到耶穌得了榮耀後才想起這些話是指他寫的，並且人們果然對他做了這些事。 \end{tabularx} \\ \\ \relax
12:17 & \begin{tabularx}{0.7\textwidth}{X} 當耶穌呼喚拉撒路，使他從死人中復活出墳墓的時候，同耶穌在那裡的眾人就作見證。 \end{tabularx} \\ \\ \relax
12:18 & \begin{tabularx}{0.7\textwidth}{X} 眾人因聽見耶穌行了這神蹟，就去迎接他。 \end{tabularx} \\ \\ \relax
12:19 & \begin{tabularx}{0.7\textwidth}{X} 法利賽人彼此說：「你們看，你們一事無成，世人都隨著他去了。」 \end{tabularx} \\ \\ \relax
12:20 & \begin{tabularx}{0.7\textwidth}{X} 那時，上來過節禮拜的人中，有幾個希臘人。 \end{tabularx} \\ \\ \relax
12:21 & \begin{tabularx}{0.7\textwidth}{X} 他們來見加利利的伯賽大人腓力，請求他說：「先生，我們想見耶穌。」 \end{tabularx} \\ \\ \relax
12:22 & \begin{tabularx}{0.7\textwidth}{X} 腓力去告訴安得烈，然後安得烈同腓力去告訴耶穌。 \end{tabularx} \\ \\ \relax
12:23 & \begin{tabularx}{0.7\textwidth}{X} 耶穌回答他們說：「人子得榮耀的時候到了。 \end{tabularx} \\ \\ \relax
12:24 & \begin{tabularx}{0.7\textwidth}{X} 我實實在在地告訴你們，一粒麥子不落在地裡死了，仍舊是一粒；若是死了，就結出許多子粒來。 \end{tabularx} \\ \\ \relax
12:25 & \begin{tabularx}{0.7\textwidth}{X} 愛惜自己性命的，就喪失性命；那恨惡自己在這世上的性命的，要保全性命到永生。 \end{tabularx} \\ \\ \relax
12:26 & \begin{tabularx}{0.7\textwidth}{X} 若有人服事我，就當跟從我；我在哪裡，服事我的人也要在哪裡；若有人服事我，我父必尊重他。」 \end{tabularx} \\ \\ \relax
12:27 & \begin{tabularx}{0.7\textwidth}{X} 「我現在心裡憂愁，我說甚麼才好呢？說『父啊，救我脫離這時候』嗎？但我正是為這時候來的。 \end{tabularx} \\ \\ \relax
12:28 & \begin{tabularx}{0.7\textwidth}{X} 父啊，願你榮耀你的名！」於是有聲音從天上來，說：「我已經榮耀了我的名，還要再榮耀。」 \end{tabularx} \\ \\ \relax
12:29 & \begin{tabularx}{0.7\textwidth}{X} 站在旁邊的眾人聽見，就說：「打雷了。」另有的說：「有天使對他說話。」 \end{tabularx} \\ \\ \relax
12:30 & \begin{tabularx}{0.7\textwidth}{X} 耶穌回答說：「這聲音不是為我，而是為你們來的。 \end{tabularx} \\ \\ \relax
12:31 & \begin{tabularx}{0.7\textwidth}{X} 現在正是這世界受審判的時候；現在這世界的統治者要被趕出去。 \end{tabularx} \\ \\ \relax
12:32 & \begin{tabularx}{0.7\textwidth}{X} 我從地上被舉起來的時候，我要吸引萬人來歸我。」 \end{tabularx} \\ \\ \relax
12:33 & \begin{tabularx}{0.7\textwidth}{X} 耶穌這話是指自己將要怎樣死說的。 \end{tabularx} \\ \\
[1ex]
\hline
\hline
\end{longtable}
$^{1}$跟大家分享神的話語.
每一次來到都感受到大家很熱情的歡迎.
大家是一個很有活力的群體.
特別我們歡迎主耶穌基督進入我們的生命.
今天是一個特別的日子.
其實全世界有不同的地方的教會.
都在今天稱為「中樹主日」.
因為是進入受難周.
耶穌騎驢進城的日子.
然後就接連在耶路撒冷發生各樣的事情.
去到祂的受苦,受難.
並且進入墳墓.
然後下個主日就是復活主日.
所以今天讀這段經文是很貼心的.
是的 剛才梁牧師都帶我們去唱了.
榮耀君王的詩歌.
我們不可以只停在唱.
今天這個經文跟大家去想.
在「中樹主日」或者其他教會稱為「Palm Sunday」.
是耶穌一個很特別的時刻.
祂得榮耀的時刻.
這也是這個經文不斷問的問題.
對大家來說到底什麼是榮耀?.
會不會是一些很光彩,很光榮的事情?.
你家人很欣賞你.
終於對得起父母,光宗耀祖.
萬人擁戴.
現在是不用擁戴的 是百萬人follow.
然後「流芳百世」.
這些詞語你會想起這些就是榮耀.
這是世上的一個定義.
「用來荒第十二章」的經文.
似乎也在描述耶穌在祂世上最光輝,最受人擁戴,最被人follow的一刻.
剛才你讀的十二節那裡說.
「一大群上來過之的人.
聽見耶穌要來耶路撒冷」.
說到第二天就是有第一天了.
第一天是什麼?前一天發生什麼事?.
其實這一章的第九節就告訴了我們背景.
那時候有一大群猶太人知道耶穌在那裡就來了.

$^{41}$他們不單是來找耶穌.
還想看一個人.
後面你看到黃色的字那裡說.
「死裡復活的拉薩路」.
你讀.
有人說你很熟悉這一段記載的.
有一大群人找耶穌親眼想見識一下那個已經死了的拉薩路.
竟然復活這件事這麼奇怪.
死人有什麼好看的呢?.
當然真的很厲害.
十幾年前我記得有一宗新聞.
當時也引起熱話.
有人在青馬大橋.
跳橋不死.
青馬大橋.
去機場那條路你一定經過的.
那條橋有多高.
這樣跳下去不死.
於是大家就去追查.
這個人的故事.
引起一個很熱的話題.
何況.
死人復活.
所以你明白為什麼那些人.
這麼熱烈要去找耶穌.
看拉薩路.
因為耶穌叫拉薩路復活是在他死了之後的.
四天.
今天就.
涼了一點.
其實這幾天也挺熱的.
開始在冬天的時候.
忘了放東西進冰箱也可以.
現在就不行了.
要很小心.
不要把食物放在冰箱.
因為會爛會臭.
在中東的天氣.
一個死屍.
四天應該都腐爛不得.

$^{81}$很嚴重.
但是這個拉薩路.
在墳墓出來.
他不是殭屍不是喪屍不是鬼魂.
第十二章一開頭還說.
他和耶穌吃飯.
見證著.
大家知道他從死裡復活.
很轟動.
很威風.
正因為這樣耶穌來到耶路撒冷的時候.
人們要出來迎接他.
耶穌進入了這個.
聖殿的地方.
是猶太人的政治,宗教,經濟.
文化的中心.
耶穌來的時候.
很多很多人來.
所以這裡的點點是代表人頭的.
這裡也不夠.
我不知道你前天.
有沒有人在新界住.
你知道.
火車出事.
我沒有經歷過.
但是我太太.
剛好在那條線坐車.
後來找到一張.
圖是說大衛站.
大衛站那張圖.
是在扶手電梯.
下去大堂拍過去的.
我看那張圖.
那些人物.
電梯都擠滿了人.
幾條電梯都擠滿了人.
整個大堂都是人.
基本上動不了.
堆滿了在那裡的情況.
這就是耶穌進城的時候.

$^{121}$那種情形.
因為聖經學者去推論.
本身耶路撒冷.
大約有二萬五千人住.
二萬五千人.
在過節的時候.
因為禦節類似我們新年.
每個人都想湧入去過節.
在那個時間.
學者推論.
耶穌那時大約有.
八萬五千人在其他地方.
加起來就是十一萬人.
在這個城.
如果你去過聖地.
這個耶路撒冷城不是很大.
就是城牆範圍.
那你想像.
在我們前面.
就是園州村.
這麼大的範圍.
圍著有十一萬人.
那是怎樣呢?.
很多啊!.
如果有十分之一的人出來看耶穌.
都是一萬一千人.
如果是百分之一.
都是一千一百人.
所以你可以想像.
那種威風.
那麼多人出來跟隨他.
是很厲害的.
而第十三節就說.
他們不單止出來看耶穌.
還做了一件事.
拿著中棗樹枝.
或者棕樹枝.
去迎接耶穌.
並且喊口號高唱.
「和珊瓦奉主名而來的那一位.

$^{161}$以色列的王配授稱頌」.
他們拿著棕樹枝.
去迎接耶穌.
還歡呼.
你可以看到圖.
這是在進城的路上.
有人在撿樹枝.
要去歡迎耶穌.
揮舞著大家去唱詩.
那為什麼要用棕樹枝呢?.
我問你.
菊花是什麼時候用的?.
快到了,清明了.
數目你會用.
你不會在情人節送菊花吧?.
那個就是你的死期.
同樣的,你情人節是不會送康乃馨的.
康乃馨什麼時候送?.
大家知道了,母親節.
那個植物是有意思的.
有一個象徵意義的.
那棕樹枝是象徵什麼呢?.
有一段故事.
原來在耶穌降生的150年前.
要計算很前的日子.
當時你知道.
以色列是亡國的.
並且不斷被不同的外族統治.
他們已經去到希臘的皇朝.
後來再分散了四個國家.
然後一直在統治.
在150年前.
猶太有個人叫西門就起義.
並且在公元前142年.
戰勝了希臘皇朝的軍隊.
收復了耶路撒冷城.
那個時候的群眾.
就揮舞棕樹枝去歡呼慶祝.
這個象徵.
也印在這個短命的.

$^{201}$瑪加比皇朝的銀幣上.
就是棕樹枝的徽號.
還刻有「解放以色列」這幾個字.
不過你明白耶穌的時候已經不是了.
已經是羅馬人.
因為羅馬興起了之後.
他們又不斷去擴張.
也再一次將耶路撒冷.
將猶太地變為植物地.
他們拿起棕樹枝揮舞.
又喊口號.
喊什麼口號呢?.
他們喊的口號.
是象徵著一件事.
「和撒拿」.
其實是來自詩篇的.
下半部分就是詩篇.
188篇25至26節.
那裡的詩篇是這樣唱的.
「求你施行拯救」.
原本在希伯來的文字.
「求你施行拯救」的讀音就是和撒拿.
所以他們去到新約時期.
他們喊的和撒拿.
其實是一句希伯來文.
和撒拿.
像什麼呢?.
小朋友可能要理解一點.
街口有個士多.
為什麼叫做士多?.
store.
英文嘛.
為什麼你搭巴士?.
bus.
翻譯的讀音.
所以和撒拿是一個翻譯的讀音.
意思就是「救命啊」.
「求你救我」的意思.
所以你叫詩篇「求你施行拯救」對著神說.
這裡說「救命啊」或者「拯救啊」.

$^{241}$「奉主名而來的那一位.
陪受稱頌」.
不止,他們還加了幾個字.
對比一下.
兩個文字.
白色的字.
「以色列的王」.
「揮中樹枝的徽號」.
「嗌著和撒拿」.
還要說「以色列的王」.
那你記得.
是哪些人將耶穌釘上十字架?.
羅馬兵丁嘛.
耶路撒冷當時因為試過幾次動亂.
所以彼拉多有派軍隊.
駐紮在聖殿附近.
怕再出事.
他們甚至建了一個塔.
在聖殿旁邊.
那個塔就是士兵監視聖殿裡的行動.
那裡有很多羅馬兵.
你當著羅馬兵叫.
「以色列的王啊」.
這樣很明顯.
你聽到他們的期望嗎?.
聽到他們心裡期待.
這個榮耀的君王來.
你竟然可以令死人復活.
更大的事我們可以期待了吧?.
期待他會不會像當時的馬加比王朝一樣.
從外族手上收回耶路撒冷.
甚至復興猶太國以色列國呢?.
他們是這樣有個期待.
這個榮耀的任務.
那就是說要打仗了.
打仗的話.
應該是怎樣的狀態呢?.
騎戰馬.
所以剛才那張照片是給你這樣看的.
應該是作戰的狀態.

$^{281}$就好像你今天打仗.
你應該是派.
當然現在流行無人機.
如果是地面部隊的話.
裝甲車.
坦克車.
都會是了.
導彈車.
你不會騎單車去的.
但耶穌不是騎戰馬.
祂騎一頭鹿.
即是騎單車去.
所以有一個很大的矛盾對差.
耶穌的行動像是向這班人潑了一盆冷水.
但他們仍然都很熱烈.
因為已經在歡迎.
即使是單車.
死人復活.
都很行啊.
我們怎樣去做都應該勝利了吧.
他們歡迎耶穌進來.
是耶穌似乎最光輝.
最多人跟祂的時候.
但當人人都這樣狂呼萬歲.
歡迎耶穌去做這個榮耀君王.
沒有人明白到底耶穌要怎樣做.
耶穌的榮耀是怎樣來的.
沒有人明白.
就算跟了耶穌三年的門徒都不明白.
所以十六節不是這樣說.
門徒起初不明白.
人人都以為自己知道.
人人都覺得上帝在這一刻就要做這件事.
上帝你要做這件事.
我們派了一個人來.
他肯定要做這件事.
值得讓我們想想我們今天怎樣看.
我們自己跟神的關係.
看到世界這麼亂.
我們會祈禱.

$^{321}$神啊.
你幫那些人.
你指著戰爭.
同時我們期待上帝做的事.
會不會是我們自己想出來的.
會不會我們好像門徒一樣.
我們起初不明白這些事.
歷史的走勢很多時候都是這樣.
我們人以為自己的期待.
對著神祈禱都可以.
但往往有些跡象發生.
我們馬上對號入座.
這個是怎樣.
往往都是錯的.
因為我們根本不明白.
耶穌的光榮榮譽那種萬人擁戴.
流芳百世.
不是奪取政權.
人不明白.
很多人都誤解.
有一個年輕人跟他媽媽相依為命.
他媽媽很辛苦很節儉.
才可以儲到一些錢送他去外國留學.
後來他畢業回來.
入了一間很大的跨國公司工作.
他入到一個團隊裡面.
他發現有幾個同事很古怪的.
每一次做整單生意.
這班人就會聚在一起.
所以他作為新手.
入了這個團隊.
他就學了.
終於他做成他第一單生意.
其實那些同事還沒有跟他傳父的呼應.
不是很對.
應該要帶他去信耶穌.
帶他回來阿根拓.
其實好像一個咒語.
或者一個儀式.
他沒有說錯.

$^{361}$萬物都是由神而來.
是神所賜的.
怎樣為他生命帶來福氣.
他以為是他期待他所想要的東西.
那些生命的成功.
那些光榮.
那些出人頭地的情況.
對得起媽媽那種處境.
頂尖的我們可能.
我們都很多人信了耶穌一段時間.
我們都會期待神會為我們做一些事.
甚至神會幫我們生命裡面的一些情況.
解決我們的困局.
當然我們可以這樣祈禱.
但是我們是不是只是用了世界給我們的那種成功方程式.
繼續求神給我們的東西呢?.
或者根據世界的虛榮方法.
以為上帝給我們看見那種好.
就算是教會.
人多就是好.
當然人多就是好.
多些人信耶穌.
但是我們會不會也會跌入一種世上的虛榮.
以為那是我們要達至榮耀的境地呢?.
這段經文給我們看到.
我們要去明白主的心.
好像門徒一樣.
後來他們明白.
這段聖經給我們今天的人看到.
我們已經知道耶穌之後會做的事.
我們應該更加比門徒容易去明白.
耶穌要我們走的光榮之路.
榮耀之路.
不是在這個進城萬人擁戴的情況.
而是接著我們會繼續去分享.
如果騎路進城萬人空巷.
不是一個榮耀.
什麼時候才是呢?.
中間就加插了一段.
繼續講榮耀的插曲.

$^{401}$法利塞人說.
那些人全跟隨他.
全世界都跟隨他.
哈利路亞他們不知道自己說對了.
他們預言了.
之後全世界都要跟隨耶穌.
這裡二十節.
有一個插曲就說.
有人來找耶穌.
那些是希臘人.
外族人來的.
不是以色列人.
他們來找一個叫費力.
其實他是門徒.
想見耶穌.
然後費力就找另一個門徒.
有人想找耶穌.
我找德烈.
兩個人就去找耶穌.
耶穌就跟他講這個說話.
耶穌沒有說見不見他.
耶穌就講了一句.
最底的說話.
人子得什麼.
榮耀的時刻到來.
又是那個主題.
榮耀.
不是進城那一刻榮耀.
那就是說.
接下來要做一些很厲害的事.
那很厲害的事是什麼呢?.
是不是大家所期待的那種.
我們終於可以修復我們的失地.
復國了.
耶穌說.
人子榮耀的時刻到來是什麼呢?.
他就說.
我老實跟你講.
我所講的榮耀.
是這樣理解的.

$^{441}$耶穌講了一個很短的比喻.
一粒墨子如果不落在地裡.
死了亦夠是一粒.
很容易明白.
我不知道現在學校還不能一人一花.
不用自己種的.
如果我們小時候就會派種子.
甚至現在流行一些回禮.
就是一些環保的卡紙.
有種子在裡面.
你整張卡紙.
即是那個包裝.
丟到泥土就可以種到東西.
如果那些種子.
繼續在你的袋裡面.
它仍然是種子.
它不會發生任何事.
這裡是這樣的意思.
如果你繼續放在那裡.
不是放在泥土裡面.
它不會有任何事發生.
但如果放在泥土裡面.
埋了它.
好像沒有了.
但你知道它會生長的.
它會出來的.
這裡表達一個生命.
如果我們不放下.
它不會成長.
不放下我們很執著的事情.
我們不會突破的.
而耶穌說.
不是喔.
死亡喔.
為什麼你說那一刻.
榮耀是說死亡的.
你連死人復活都可以做到.
為什麼你會覺得.
這個榮耀和死亡有關呢.
死亡和榮耀王有什麼關係呢.

$^{481}$當時門徒不明白.
旁邊的人都不明白.
但聖經告訴我們.
當耶穌得榮耀的時候.
不是在城門口.
也不是在他潔淨聖殿.
是在他接受十字架的酷刑.
被釘死的時候.
沒有人想過用這個方法.
沒有人明白.
耶穌說.
當我從地上被舉起來.
就是他上十字架.
就更多人跟他.
真正跟隨他的.
是看到他在十字架受死.
這是他得榮耀的時候.
這個至高的榮耀.
也說明他會透過死在十字架.
戰勝死亡和撒旦的權勢.
所以回到耶穌講完這個比喻.
講的解釋.
他說那個墨子不落在地上死了.
仍舊是一粒.
若果死了.
結出許多枝粒來.
然後他繼續講生和死的對比.
他說愛惜自己生命的人.
就喪失生命.
嘩.
小心聽.
這裡不是說我們可以隨意對待自己的生命.
但是他說你要恨自己的生命.
怎麼理解.
有誰不愛惜自己的生命.
小心保護自己的生命.
我試過有一段時間.
頭經常痛.
連續幾天都是這樣.
開始覺得很辛苦.

$^{521}$有一天文明修的時候.
有點頭痛.
持續的頭痛.
我都覺得麻煩.
我應該都很怕死.
我們很容易去解決頭痛問題.
不賣廣告了.
賣個什麼平板電腦.
馬上解決它.
這裡不是說養生或者對付自己疾病.
不要理那些.
不是這個意思.
耶穌是在表達.
會不會我們捧著我們的生命的目標.
到一個位置.
我們只是看著它.
甚至有些人.
你身邊還有一兩個這樣的朋友.
見到他就介紹你吃一些保健食品.
他不一定做過全消.
但是他自己吃很多.
他不斷去鑽研.
他是一個專家.
接著就說168不行.
又怎麼樣.
很多的方式去養生.
去令自己身體越來越健康.
保持青春.
但是你會發現.
往往他的一生就在做這件事.
他所有時間都是研究這件事.
他一生就捧著自己的生命.
在愛惜到一個位置.
他好像沒有其他生活.
等於說沒有其他的事.
令到你變成這樣.
你很愛惜你的生命.
你的活動.
那樣東西成為你唯一.
你一生都圍繞著這件活動去過.

$^{561}$回頭看很多東西.
你都沒有.
很多東西你都喪失了.
耶穌去指出我們人一個盲點.
有時候我們害怕死亡到一個位置.
我們不斷去找方法避開.
所以耶穌的說法.
我們依然不是太明白.
我看見一個台灣的見證.
有一個女孩子叫王成玲.
她父母對她都很嚴格訓練.
因為父母親都讀很多書.
要求很高.
所以她兩姊弟在學校的成績都很好.
出來工作.
也都進了一些好的公司.
公司工作裡面她很努力.
公司嘗試欣賞她.
派她去海外開新的分公司.
後來只是因為自己的弟弟急病.
同時自己年老的母親又跌傷.
唯有放棄在海外的很好的工作.
回流.
然後去照顧弟弟.
但一直照顧弟弟的時候.
她媽媽都對她仍要求高.
你要這樣照顧她那樣.
她覺得很大壓力.
我一生都在追我的指標.
要成功 要優秀.
要做得好.
我很累.
於是有一天她突然間覺得.
不行了.
我要找我的自由.
我要找我真正願意去接納我.
無論我是怎樣的人.
我得到那份的愛.
於是她就走去.
不是回教會.

$^{601}$她去印度學一些禪修的方法.
學一些神秘宗教訓練的儀式.
操練.
學咒語.
她祈求她能夠心想事成.
隨心所欲.
得到她想要的東西.
她有多厲害.
學到那些方法.
她又可以用來猜中彩票多少錢.
她又賺到一些錢.
十幾年她的人生就圍繞這兩條路.
想過自己想要的生活.
盡量令自己有那份感覺自由.
可以控制的人生.
不再被人指責.
結果她學了很多催眠.
各樣的靈性療法.
探過不同的大師.
跟他們學習.
占卜.
風水.
什麼都好.
她都是這樣.
一直將她的生命投放在這裡.
這是她愛惜她生命的方法.
但是你知道的.
我們今天在教會裡面.
我們明白.
這樣.
是不是真的愛惜生命呢?.
《約翰福音》三章十六節.
是很多人都懂的.
一起讀吧.
神愛世人.
甚至將她的獨生子賜給他們.
要一起信她的不致滅亡.
凡得永生.
這裡所講.
三章十六節的真理.

$^{641}$同一本書.
今天講了十二章.
《約翰福音》.
他在這裡用的不致滅亡.
滅亡這兩個字.
就是耶穌說.
你愛惜生命.
反而會喪失.
生命的喪失這兩個字.
你會滅亡.
你一直捧著你的人生.
你以為你可以掌握到.
然後去追尋.
你會滅亡.
所以不是說.
耶穌說你不要愛惜自己的生命.
不是說不要注意健康.
不要注意衛生.
不是這些.
而是我們太愛太愛太愛.
這個正是敗亡的生命.
正如這個女孩子一樣.
她太想趨吉避凶.
她到處求神問福算命.
自己學到一個位置.
她很厲害.
找到她要的東西.
還自己成為了一個大師.
她在當時不同的電台.
電視台來訪問她.
問她接下來一年會怎樣.
有不同的明星去找她幫助.
要去得治療.
或者去通靈.
她很開心.
我達到大師的地位.
人人景仰.
我想怎樣就怎樣.
當時有一個少少名氣的女歌手.
晚上來找她做一個催眠治療.

$^{681}$催眠到一半.
突然間這個女歌手.
臉色都變青了.
變了臉色.
然後起身.
開始自己傷害自己.
還用對很凶狠的眼.
看著這個聖王的大師.
你為什麼要這樣幫他.
一把男人聲.
在這個女歌手裡面出來.
基督徒就明白.
是什麼原因.
嚇到這個大師.
馬上把她這十幾年所學過的.
所有咒語.
所有的神明都搬了出來.
依然沒有改變.
那個變了臉的女歌手.
還一步一步走近她.
那時候她很害怕.
這個王大師.
什麼都試過了.
沒有用.
她看到旁邊有一個窗.
窗外對面有個教會.
教會那裡有個十字架.
還有一個.
上帝救我.
耶穌救我.
接著那個女歌手就.
整個人躺在地上.
過了一會兒她醒了.
這個女歌手一路走.
她說其實我已經.
有這些東西搞我.
搞了十幾年.
王大師這個時候就明白.
原來她一生所追尋的.
都不過一下就沒有了.

$^{721}$人是敗亡的.
她又想起叫耶穌的名字.
求拯救.
耶穌來幫她.
和撒那.
她經歷耶穌的幫助.
不是令她名成利就.
而是得到一種真正的解脫.
後來她真的進了教會.
有人和她祈禱.
她才感受到.
她從來沒有感受過的平安.
後來她信了主.
也將這個見證講了出來.
一個沒有基督的生命.
結局只是喪失.
結局只是滅亡.
所以恨惡生命到一個地步.
你是要知道.
你不是要再捧著這個敗壞生命.
不斷去刁鑽它.
當然我們人要負責任.
要努力的.
我們生活很多的事.
我們要做的.
但不是你捧著它.
不斷求,不斷盼望上帝去祝福你.
要成為一個怎樣的人.
我們反而是要將人生的話事權.
交給耶穌.
將這個舵交給祂掌管.
讓祂去握著這艘快要沉的船.
這個時候就會有出路.
弟兄姊妹,我們今天的生命是怎樣呢?.
耶穌說一粒種子落在地裡面死了.
是能夠復活的.
我們的生命是怎樣呢?.
我們的生命是不是.
我們交給主又拿回來.
我們的舵交給神.

$^{761}$是的,你都回來敬拜了.
很好,還要九點三.
做這一堂.
但是我們會不會出了教會門口.
我們又馬上搶回大盤來握呢?.
或者這個世界又會搶回來.
做你的司機.
帶你去一個失喪敗亡的路徑.
那種追尋.
弟兄姊妹,我們在《宗廚主席記得》.
耶穌和你說的榮耀.
祂所示範的.
是祂怎樣順服神的心意.
去成就這一切.
放下祂的生命.
不再執著.
能夠活出更精彩的人生.
當然生命主權這個說法.
幫我們記得.
我們可以放手和交托.
Let go and let go.
就是停止我們對我們自我生命.
那種絕對掌握.
我們自己要怎樣就怎樣.
而是交給神告訴你.
怎樣繼續向前走.
簡單來說.
各位,你信了耶穌.
你的命是耶穌的.
如果是這樣的話.
就不是決志那一刻.
我們就信了耶穌.
接著就OK了.
那就等永生.
也不只是祈禱讀經.
回教會的時候.
我們就跟隨耶穌.
而是每一時每一刻.
都好像你的壁畫上這樣寫.
你是要背著十字架天天跟從主.

$^{801}$是天天.
是不是?.
緊緊跟隨這個榮耀的王去服侍祂.
耶穌說你要跟我的.
跟我的話.
是怎樣跟法呢?.
我請大家再讀二十六節.
預備起.
「人如果服侍我,就必須跟隨我,我在那裡,服侍我的人間要在那裡,人如果服侍我,我父就會給他尊榮.」.
騎老進城的榮耀.
不是耶穌要的.
耶穌說我榮耀的時刻.
是指著十字架的.
這一段繼續說的榮耀.
是耶穌說父神都會將那份尊榮給你.
而這份尊榮不是你在世上那種追求的高峰.
那些是很脆弱的.
這個尊榮是當我們去順服神的心意.
我們跟隨耶穌的時候.
走這條順服的道路.
我們會經歷到的.
所以這裡耶穌所說的有兩方面的意思.
一方面是回頭看.
回頭看耶穌是怎樣由祂順服父神的心意一直走到十字架.
這條路.
祂怎樣示範.
祂怎樣走過.
我們可以跟.
而第二方面就是講將來.
我們這樣跟隨耶穌.
我們本來由生命敗亡.
喪失生命.
走向滅亡.
耶穌的死換來我們的永生.
這個榮耀生命的結局.
神說我會尊榮這些一直跟從耶穌的人.
而那個結局是榮耀的.
會分享神的榮耀.
聖經很多處都是這樣說.
所以我們不單歸榮耀給神.

$^{841}$我們知道有一天這個榮耀.
上帝都會跟我們分享.
那是什麼榮耀呢?.
五字.
榮耀在聖經的字的意思是代表重量.
份量.
重要的意思.
神體重你.
神記得你.
所以為什麼耶穌說.
祂升天之後我去為你們預備地方.
記得你的.
你一生走這個跟從我一樣的十格道路.
我記得你.
你有一個位置.
我會保證你在上面.
父神知道你重視你.
體重你.
記住你.
尊榮你.
不會丟下你.
而我們的結局是確定的.
就好像看恐怖片.
恐怖片最好看的過程就是.
當中被人嚇.
如果一早劇透.
主角嚇完一輪都不會死.
那就沒什麼好看了.
或者不怕.
現在已經劇透了.
我們走到底就是.
榮耀的天家.
有永生.
如果是這樣的話.
這個劇透了的劇情.
會不會令到我們今天.
看我們的生命完全不同.
那就不怕了.
為什麼我們今天仍然為很多事去怕呢.
就是因為我們忘記了.

$^{881}$天父會記得我們.
體重我們.
並且我們會分享他榮耀.
天家裡面.
我們是有份的.
劇透了.
你們沒什麼好怕的.
什麼時候打完仗.
猶太會怎樣.
這一生我們有很多困難要面對.
但是我們知道.
無論我們這一生遇到高或低.
困難還是順利.
我們最後的結果是一樣的.
這樣會改變我們的人生.
這個就是耶穌復活的生命.
在今天已經對你和我有效.
有個弟兄他從事專業的工作.
他小時候已經知道.
他信了主之後.
他說我要.
有三個字我要事奉主.
不是梁牧師約定你.
每人有三個字.
同時他有另外三個字.
他一直記住.
是他人生的目標.
他有三個字.
賺大錢.
即是專業.
應該透過專業可以賺大錢.
這是基本香港人的default.
他很喜歡跑步.
所以他每次跑步都是在思考的時候.
未來見這個客人我怎樣做成這生意可以賺大錢.
或者未來的市場怎樣處理投資.
希望可以賺大錢.
一直在頭腦裡想這些.
有一天他又去跑步.
跑一跑.

$^{921}$不小心踢到.
啪啪.
跌在地上.
撞到頭破血流.
他坐在地上.
他突然間好像醒了.
難道我一生就是這樣過.
又說要事奉神.
在那一刻上帝呼召他.
你真是來事奉我.
他後來成為了一位傳道人.
當然離開了他的專業.
去讀神學等等.
十幾年之後.
他重遇當時他的舊同事.
有些舊下屬.
通通都已經升職.
做了高層.
住豪宅.
駕靚車.
很多很風光的成就.
但這些一群同事.
每次見他都有同一個表達.
為什麼你還是這麼年輕.
你知道有時候做投資界.
或者投行那些.
都是爆肝地做.
一夜夜衝那些業績.
當然他的生命有改變.
使他事奉主的人有不同.
重點不是外貌.
而是他很清楚他一生在走什麼路.
很穩定有個永生的結局.
所以無論這一生是富貴貧賤.
他都沒有問題.
這種輕省頂折.
會不會我們信了住一段日子.
我們放下了.
我們把重擔都拿在背上.
把擔憂都堆在心裡.

$^{961}$但耶穌明明是說.
你背起十字架.
但我的額是容易負的.
凡路苦擔重擔來到這裡.
我會給你安息的.
等於復活的生命.
在今天怎樣來到.
就是你記得榮耀的王.
怎樣走十字架榮耀的道路.
並且他會跟你分享這份榮耀.
所以當我們遇到今天各樣的困難.
記得我們已經有劇透了的結局.
不會有什麼好怕的.
你都知道會怎樣.
我的年代.
我成長的年代看《大長今》.
不知道你們有沒有看《大長今》.
我最初看《大長今》是很慘的.
一個韓劇的人物.
他又千辛萬苦.
終於才可以捱過宮廷的生活.
但我很想不想去找一下《大長今》的歷史.
知道後來皇帝又怎樣去對待他.
或者韓國播了先.
報紙都會有劇透.
我看了劇透.
之後這個人又會對他好.
那個壞人又會被懲罰.
那就沒那麼傷心了.
今天你遇到的事是會令你傷心的.
令你驚嚇的.
但回想我們有榮耀的將來.
確定的.
沒甚麼好怕的.
這個是復活生命在今天對我們的作用.
所以我們跟從去順服著這條路.
我們是會得到真正的榮耀.
神記得你.
重視你.
落在地裡面死了.

$^{1001}$表達我們願意去放開我們的生命.
不要再爬這個世界設定給你的階梯.
捧著那個隨時會吹熄的蠟燭.
好緊張.
好焦慮.
好怕跌低.
望向我們的十字架.
記得神真正給我們那種光照.
將生命擺在他面前.
化為生命的種子.
我們的生命被他栽培使用.
結出真正的果子.
如果是這樣的話.
在中樹主人我們看到耶穌給我們的榮耀的概念是不同的.
願我們都能夠立志去過這個順服的人生.
是的,其實頂智無人.
原來真正的榮耀是來自信靠天父.
天父怎樣叫我們過人生.
我們就跟著做.
當我們這樣做的時候.
我們不知不覺一次一次經歷將生命放下.
然後埋在地裡卻長出果子的豐盛.
頂智無人你可以得到.
待會我們就完了崇拜.
可能有些事我們就回家.
你仍然面對家裡的困難.
孩子的掙扎.
工作,明天上班.
很不容易.
現在的世界是這樣的情況.
但你記得.
聖經劇透了.
不要擔心.
每一天我們就放手.
交託給主.
經歷每一天祂給你.
復活生命所帶來的豐盛.
我們一起祈禱.
父神啊,願你繼續引領我們.
學像主基督那種順服你的心意.

$^{1041}$卻登著真正榮耀的經歷.
在今天我們歸榮耀給你.
同樣我們知道你極寒蟲我們.
讓我們都得著那種永生的福分.
並且你復活生命.
在今天可以改變我們.
幫助我們每一個.
多謝你.
我們祈禱奉耶穌基督的聖名祈求.
Amen.
\newpage



\section{出埃及記 4:2}
\label{sec:031906}
\textbf{聖經人聖經事:你手裡是甚麼}~~~~黃克勤牧師博士
\newline
\newline
連結: \href{http://www.church.com.hk/acms/content.asp?site=cdc\&op=show\&type=product\&code=031906\&layout=sermon}{\texttt{031906}} ~~~~ 語音日期: 2026-03-31
\newline
\newline
\hyperref[sec:031905]{\small{< < < PREV SERMON < < <}}
~
\hyperref[sec:index_chronic]{\small{[返順時目]}}
~
\hyperref[sec:031907]{\small{> > > NEXT SERMON > > >}}
\newline
\newline
出埃及記 4:2
\newline
\begin{longtable}{cl}
\hline
\hline
章節 & 經文 (和合本修訂版)\\
\hline
4:2 & \begin{tabularx}{0.7\textwidth}{X} 耶和華對摩西說：「你手裡的是甚麼？」他說：「是杖。」 \end{tabularx} \\ \\
[1ex]
\hline
\hline
\end{longtable}
$^{1}$各位弟兄姊妹大家好.
一齊讀聖經.
生命樹孫鴿網絡.
跟大家看聖經人和聖經的事.
在聖經的人物和故事的背後.
去認識上帝多些.
是這個內容的主題.
今天就來到這個題目.
「你手裡是什麼?」.
這句說話是誰說的呢?.
你手裡有些什麼?.
我手裡有這支指揮棒.
「你手裡是什麼?」.
這個是上帝說的.
我們的主文摩西.
繼續是摩西的故事.
他這個人.
《昭希及記》第四章.
仍然是神在荊棘中和摩西說話.
你和我.
我聽見以色列人.
我的百姓的哀聲.
我紀念亞伯拉罕以撒雅各.
我和他們納的藥.
我聽到.
我知道.
於是現在我要叫你.
去帶他們出來.
跟法老說讓他們走.
摩西第一個.
千萬不要.
不要叫我.
不要等我來.
摩西第一個回答.
如果他們問我.
你叫什麼名字.
我怎麼回答?.
其實我都不太知道.
摩西為什麼會問這個問題.
因為其實他們都知道.

$^{41}$如果他一說.
大家希伯來人.
亞伯拉罕以撒雅各的神.
叫我來帶大家走.
沒有人會夠膽問.
神叫什麼名字.
他們連名字都不敢說.
Yahweh.
他們不敢說.
這句字不會.
是啊是啊.
他們不敢說.
為什麼他們不會問.
都是說 Elohim.
或者 Elohim.
是Yah.
這樣就好.
現在摩西就繼續.
摩西就說.
他們必不信我.
又不聽我說.
他們一定會說什麼.
他們不是問什麼名字.
是會問第二個問題.
耶和華並沒有向你顯現.
他們不會聽我說.
這個比較合理.
各位同胞.
大家留意聽我說.
過來過來.
亞伯拉罕以撒雅各的上帝.
叫做M那個上帝.
是Yah.
叫我來帶大家離開.
那些人就會說.
你不是玩了嗎.
摩西先生.
想我們死嗎.
現在我們已經做事的那塊磚.
已經被打了.

$^{81}$離開出去.
還沒有走到門口.
沒命了.
不要妨礙我開工.
我很忙.
一定會這樣.
摩西就說.
他們不會信我.
不會聽我的.
他們一定會說.
耶和華沒有向你顯現.
你不要玩了.
不要作弄了.
神什麼時候向你顯現.
耶和華哪有向你顯現.
這個比較合理.
但是耶和華對摩西說.
這是第一點.
螢光紅色第一點.
神問摩西.
你手裡是什麼.
他就說是一支杖.
其實他們時常都要拿著一支杖.
因為經常走路.
有時動來動去.
動來動去就像現在旅行行山.
好人好者對腳健康.
行山的人都要拿一支長高的杖.
行山杖有時幫助力量.
有時用來撩動東西.
都要的.
而且牧羊人一定有一支長的杖.
有時用來分開羊.
有時用來隔開野獸.
打牠之類.
通常這支杖頂有個圈.
如果有隻羊滑到坑裡.
牠就拉牠的羊上來.
那隻羊是這樣動.
救那隻羊.

$^{121}$所以你手裡是什麼?是杖.
為什麼我問這個問題是第一點.
其實當神要用你的時候.
神會看看你手裡有什麼.
你手裡有什麼.
我就用你手裡有的東西.
你手裡有什麼.
有一支杖.
我就用這支杖.
你手裡有什麼.
有一個鶴位.
我就用這個鶴位.
你手裡有什麼.
我就沒什麼我有些朋友.
我就用你的朋友.
是你有些什麼.
你手裡有什麼.
有一隻鑊鏟.
好啊.
你手裡有什麼.
我有支杖.
好,我就用你那支杖.
耶和華就說.
將你支杖雕在地上.
他一雕下去.
就變成蛇了.
摩西便跑開.
這裡有故事說.
你手裡有什麼.
我有一支牧羊人的那支杖.
拿著你也看到.
說實話也看到.
如果我看到.
我會不會問你手裡有什麼.
其實是想他集中看.
我手裡有一支杖.
神看到其實不用問.
不過就說你有什麼.
我有支杖.
好,丟在地上.

$^{161}$丟在地上.
一丟就變成一條蛇.
哇,摩西便跑開.
那條蛇是不是一條.
好像玩具塑膠的蛇.
看起來也不是真的呢.
不是,應該是一丟在地上.
那支杖變成一條蛇.
是很生猛的.
哇,嚇到他.
哇.
在沙漠地.
是有蛇的.
有些蛇是很毒的.
你知道他們.
將埃及進入迦南.
走在曠野的時候.
試過一段時間有火蛇咬他們.
是一些很毒的蛇.
咬到很痛.
而且在曠野地.
如果被蛇咬到.
一個人沒有人救到你.
摩西看到那條蛇.
哇,搞不定.
馬上跑開.
這也是一個反應.
也表示那條蛇是很真.
真是一條蛇.
摩西對摩西說.
如果是我.
多加一句.
我自己作的.
喂,不要走,回來.
蛇啊,上帝.
我的杖被你變成一條蛇.
你真是厲害.
伸出手來.
拿住他的尾巴.
他就會在你手中.

$^{201}$變回一支杖.
嘩,不是吧.
為什麼要變成蛇呢.
其實也象徵著埃及.
埃及的法老.
哇,就是一條這麼兇猛的蛇.
不怕的,你聽我說.
如果你走過去.
捉住他的尾巴.
就可以了.
是喔,給他一些聯想.
好了,捉住他的尾巴.
他就會在你手中.
變回一支杖.
於是.
就叫他們可以相信.
是耶和華他們祖宗的神.
就是亞伯拉罕的神,以撒雅各的神.
以撒的神,雅各的神.
向你顯現.
這個.
蛇變回杖的神跡.
會跟著你的.
你去到.
耶和華的神.
亞伯拉罕,以撒雅各的神.
叫我帶你們走.
不要玩我們了.
我們還不夠慘嗎.
神必要向你顯現.
看著,變了一條蛇.
捉住他,變回一支杖.
哇.
他們就會相信你.
我給你這個神跡.
跟著.
他們就知道.
是我向你顯現過的.
耶和華又對他說.
第二點.

$^{241}$把手放在懷裡.
把你的手放在衣服裡.
摸到肚子.
摸到身體.
然後他就把手放在懷裡.
及至抽出來.
不料一拉出來.
哇.
手上長了大麻蜂.
又雪那樣白.
這些地方.
曬得多.
污糟,比較深色.
一拉出來.
哇.
大麻蜂是什麼病.
在聖經裡不算很清楚.
大家都認為.
是麻蜂.
但其實大麻蜂很有可能.
是一種很嚴重.
當日有的.
好幾種皮膚病.
醫不了.
很嚴重,皮膚會白.
會爛.
會辛苦.
大麻蜂.
整隻手都死了.
如果這樣.
很難醫治.
然後又說.
再把手放在懷裡.
他就把手放在懷裡.
及至從懷裡.
抽出來.
不料手已經復原.
與周身的肉一樣.
把你的手.
放在你的身體裡.

$^{281}$看看.
哇.
死了,大麻蜂啊.
上帝.
會不會弄傷整個身體.
有沒有問題.
然後又放回懷裡.
再拉出來.
哇,復原了.
真是眼前一亮.
給了兩個神蹟.
他握手.
你去到.
你示範給他們看.
倘若他們不聽你的話.
也不信頭一個神蹟.
必信第二個神蹟.
你會見到.
這個神蹟向誰.
我給你兩個神蹟.
頭一個,第二個.
我給你兩個神蹟握手.
你去跟以色列人說.
他們不信,不會信我猜險你.
做給他們看.
他們就會信.
你見到.
第二點.
第一點,你手裡有什麼.
我手裡沒什麼.
很普通,一支爛鬼木匠.
木匠人的杖.
是差不多的.
瓷就沒有了,瓷杖都在我手裡.
很普通.
丟在地上.
我可以用你的杖.
顯出我的能力.
你手裡有什麼,獻給主.
第二點.

$^{321}$你手裡有什麼.
沒什麼.
我沒什麼,連支杖都沒有.
行.
你有沒有手,有.
肯就行.
我就用你沒有的東西.
你的手.
用你的手,bare hand.
什麼都沒有.
都行,你有什麼就獻什麼.
沒什麼,不要緊.
你的手都可以獻.
上帝都可以用它.
第三點.
接下來.
倘若他們不聽.
都不聽.
竟然頭一個神蹟又不信.
第三點要講.
如果他們不信.
第一個又不信,第二個神蹟.
都會信.
如果竟然.
借兩個神蹟.
若都不信.
又不聽你的話.
你就從河裡面取些水.
倒在旱地上.
你在河裡面.
拿上來的水.
就會在旱地上變做血.
做一個水變血.
不覺得十宅有水變血.
但其實.
這個水變血的災.
在這裡不是向.
埃及人發出.
是向以色列人發出.
給摩西的第三個.

$^{361}$拿手神蹟.
第一個.
掌變蛇,蛇又變回掌.
第二個,手裡面長大蠻蜂.
然後又醫治.
第三個就是.
將水倒在地上.
眼眶乾乾的.
用水舀上來.
倒在地上就變成成地血.
嘩,頭兩個.
如果他們不信.
他們都會信.
第三個,我給你三個神蹟.
拿去握手.
讓我的子民.
相信.
第三點想講的.
就是你一定要.
對上主有信心.
對上主有沒有信心.
有沒有信心.
摩西就說.
我.
我不是對你沒信心.
我對那些人沒信心.
我對自己都沒什麼信心.
對那班以色列人沒信心.
這就是摩西.
好,上主就說.
你想他們有信心,行.
我透過你行三個神蹟.
第一個,第二個,第三個.
第一個他們會信,不信有第二個.
第二個都不信,有第三個.
水變血.
所以你看到.
我想講的第三點就是上帝.
有沒有計劃.
原來上帝.

$^{401}$是很有計劃的.
想了先.
你看他創造就知道.
第一日,第二日.
第三日,第四日,第五日.
第六日,第七日休息.
做好了,然後.
一條救數的線.
簡直是人亞伯拉罕一直下來.
到將來.
重新恢復.
與人同在比伊甸園更好.
住在新天新地裡面.
你看到.
上帝是有計劃的.
有第一個神蹟.
第二個神蹟,第三個神蹟.
你不.
信.
不信自己.
不信.
那班人會接受.
上帝是有工作的.
千萬不要不信神.
其實這位.
是有計劃的神.
他是預備了計劃.
計劃A,計劃B,計劃C.
救贖沒有計劃A,計劃B,計劃C.
救贖是必定成功.
是一個計劃.
所以救贖上千萬不要說.
舊約不成功,於是神計劃B.
新約,千萬不要這樣講.
這是錯的,教義上.
信上信錯的.
做這些實際上的事.
神是有第一件,第二件,第三件.
十個災.
一至十,愈來愈重.

$^{441}$第十災,滅掌子.
是有計劃的.
上帝的計劃是會成功的.
成不成功.
在乎我,當然不是.
沒有你都可以.
沒有你,沒有我都可以.
有祂自己就行了.
今天跟大家分享到這裡.
第一件事就是.
你手上有什麼.
好,我就用那支杖.
杖邊射,射邊杖.
我什麼都沒有.
現在只有一隻手.
就用你一隻手吧.
我都做到事的.
戴上風,康復.
如果這樣都不行呢.
好,再超自然一點.
在祂面前.
舀些河水,叫祂舀都還可以.
然後倒在地上就會變成血.
成地血.
嘩,這樣他們會相信.
其實後來以色列人是相信的.
真的撿東西走.
不然怎麼叫祂走.
你看到上帝是有工作.
有計劃的.
有時候覺得很混亂,不知怎樣發生.
那些人又這樣又反對.
不知怎樣,放心.
交託給主,祂有祂的時間.
祂有祂的工作.
祂會做到的.
你將你手上的東西給祂就行了.
手上沒有東西都可以.
祂都做到的,一起祈禱吧.
讚美我們的上帝.

$^{481}$你是一個滿有計劃的上帝.
因為你是一個.
很睿智.
很有智慧的上帝.
你是一個有先見之能.
看到很遠.
無限遠的上帝.
你知道我們那些軟弱.
我們是會推搪的.
你一早就認識知道我們.
你就為我們預備了.
很多方法.
甚至在這裡都幾少有.
幾罕有的聖經裡面.
是神給三個神職.
給摩西握手.
叫他帶意識的人出來.
其實是看到神是.
你是多麼體重.
要帶你的百姓.
在這個圍爐之家出來.
在這個被逼迫的地方.
出來.
所以你很罕有地.
將三個神職交給一個人.
給他握手.
我們見到真的少有.
你自己又滿有意志與計劃.
你願意你的心意成就.
願你的旨意.
是行在地上.
願你的國降臨.
主我們就好像摩西.
手上有支杖交給你.
什麼都沒有.
只有一隻手就這樣一隻手交給你.
或者將我周圍的環境.
河水交給你.
你都可以做事.
願你的名得在榮耀.

$^{521}$禱告奉耶穌名求.
\newpage



\section{以弗所書 5:1-6:24}
\label{sec:031907}
\textbf{以弗所書 (3) 仇敵與全副軍裝}~~~~劉少康牧師
\newline
\newline
連結: \href{http://www.church.com.hk/acms/content.asp?site=cdc\&op=show\&type=product\&code=031907\&layout=sermon}{\texttt{031907}} ~~~~ 語音日期: 2026-04-01
\newline
\newline
\hyperref[sec:031906]{\small{< < < PREV SERMON < < <}}
~
\hyperref[sec:index_chronic]{\small{[返順時目]}}
~
\hyperref[sec:031908]{\small{> > > NEXT SERMON > > >}}
\newline
\newline
以弗所書 5:1-6:24
\newline
\begin{longtable}{cl}
\hline
\hline
章節 & 經文 (和合本修訂版)\\
\hline
5:1 & \begin{tabularx}{0.7\textwidth}{X} 所以，作為蒙慈愛的兒女，你們該效法神。 \end{tabularx} \\ \\ \relax
5:2 & \begin{tabularx}{0.7\textwidth}{X} 要憑愛心行事，正如基督愛我們，為我們捨了自己，當作馨香的供物和祭物獻給神。 \end{tabularx} \\ \\ \relax
5:3 & \begin{tabularx}{0.7\textwidth}{X} 至於淫亂和一切污穢，或是貪婪，在你們中間連提都不可，這才合乎聖徒的體統。 \end{tabularx} \\ \\ \relax
5:4 & \begin{tabularx}{0.7\textwidth}{X} 淫詞、妄語和粗俗的俏皮話都不合宜；總要說感謝的話。 \end{tabularx} \\ \\ \relax
5:5 & \begin{tabularx}{0.7\textwidth}{X} 要確實知道，無論是淫亂的，是污穢的，是貪心的（貪心的就是拜偶像的），在基督和神的國裡都得不到基業。 \end{tabularx} \\ \\ \relax
5:6 & \begin{tabularx}{0.7\textwidth}{X} 不要被人虛浮的話欺騙了，因這些事，神的憤怒必臨到那些悖逆的人。 \end{tabularx} \\ \\ \relax
5:7 & \begin{tabularx}{0.7\textwidth}{X} 所以，不要與他們同夥。 \end{tabularx} \\ \\ \relax
5:8 & \begin{tabularx}{0.7\textwidth}{X} 從前你們是暗昧的，但如今在主裡面是光明的，行事為人要像光明的子女— \end{tabularx} \\ \\ \relax
5:9 & \begin{tabularx}{0.7\textwidth}{X} 光明所結的果子就是一切的良善、公義、誠實。 \end{tabularx} \\ \\ \relax
5:10 & \begin{tabularx}{0.7\textwidth}{X} 總要察驗甚麼是主所喜悅的事。 \end{tabularx} \\ \\ \relax
5:11 & \begin{tabularx}{0.7\textwidth}{X} 那暗昧無益的事，不可參與，倒要把這種事揭發出來。 \end{tabularx} \\ \\ \relax
5:12 & \begin{tabularx}{0.7\textwidth}{X} 因為，他們暗中所做的，就是連提起來都是可恥的。 \end{tabularx} \\ \\ \relax
5:13 & \begin{tabularx}{0.7\textwidth}{X} 凡被光所照明的都顯露出來， \end{tabularx} \\ \\ \relax
5:14 & \begin{tabularx}{0.7\textwidth}{X} 因為使一切顯露出來的就是光。所以有話說：「你這睡著的人醒過來吧！要從死人中復活，基督要光照你了。」 \end{tabularx} \\ \\ \relax
5:15 & \begin{tabularx}{0.7\textwidth}{X} 你們要謹慎行事，不要像無知的人，要像智慧的人。 \end{tabularx} \\ \\ \relax
5:16 & \begin{tabularx}{0.7\textwidth}{X} 要把握時機，因為現今的世代邪惡。 \end{tabularx} \\ \\ \relax
5:17 & \begin{tabularx}{0.7\textwidth}{X} 不要作糊塗人，要明白主的旨意如何。 \end{tabularx} \\ \\ \relax
5:18 & \begin{tabularx}{0.7\textwidth}{X} 不要醉酒，酒能使人放蕩；要被聖靈充滿。 \end{tabularx} \\ \\ \relax
5:19 & \begin{tabularx}{0.7\textwidth}{X} 要用詩篇、讚美詩、靈歌彼此對說，口唱心和地讚美主。 \end{tabularx} \\ \\ \relax
5:20 & \begin{tabularx}{0.7\textwidth}{X} 凡事要奉我們主耶穌基督的名常常感謝父神。 \end{tabularx} \\ \\ \relax
5:21 & \begin{tabularx}{0.7\textwidth}{X} 要存敬畏基督的心彼此順服。 \end{tabularx} \\ \\ \relax
5:22 & \begin{tabularx}{0.7\textwidth}{X} 作妻子的，你們要順服自己的丈夫，如同順服主。 \end{tabularx} \\ \\ \relax
5:23 & \begin{tabularx}{0.7\textwidth}{X} 因為丈夫是妻子的頭，如同基督是教會的頭；他又是這身體的救主。 \end{tabularx} \\ \\ \relax
5:24 & \begin{tabularx}{0.7\textwidth}{X} 教會怎樣順服基督，妻子也要怎樣凡事順服丈夫。 \end{tabularx} \\ \\ \relax
5:25 & \begin{tabularx}{0.7\textwidth}{X} 作丈夫的，你們要愛自己的妻子，正如基督愛教會，為教會捨己， \end{tabularx} \\ \\ \relax
5:26 & \begin{tabularx}{0.7\textwidth}{X} 以水藉著道把教會洗淨，使她成為聖潔， \end{tabularx} \\ \\ \relax
5:27 & \begin{tabularx}{0.7\textwidth}{X} 好獻給自己，作榮耀的教會，毫無玷污、皺紋等類的缺陷，而是聖潔沒有瑕疵的。 \end{tabularx} \\ \\ \relax
5:28 & \begin{tabularx}{0.7\textwidth}{X} 丈夫也應當照樣愛妻子，如同愛自己的身體；愛妻子就是愛自己了。 \end{tabularx} \\ \\ \relax
5:29 & \begin{tabularx}{0.7\textwidth}{X} 從來沒有人恨惡自己的身體，總是保養愛惜，正像基督待教會一樣， \end{tabularx} \\ \\ \relax
5:30 & \begin{tabularx}{0.7\textwidth}{X} 因我們是他身體的肢體。 \end{tabularx} \\ \\ \relax
5:31 & \begin{tabularx}{0.7\textwidth}{X} 「為這個緣故，人要離開父母，與妻子結合，二人成為一體。」 \end{tabularx} \\ \\ \relax
5:32 & \begin{tabularx}{0.7\textwidth}{X} 這是極大的奧祕，而我是指基督和教會說的。 \end{tabularx} \\ \\ \relax
5:33 & \begin{tabularx}{0.7\textwidth}{X} 然而，你們每個人都要愛妻子，如同愛自己一樣；妻子也要敬重她的丈夫。 \end{tabularx} \\ \\ \relax
6:1 & \begin{tabularx}{0.7\textwidth}{X} 作兒女的，你們要在主裡聽從父母，這是理所當然的。 \end{tabularx} \\ \\ \relax
6:2 & \begin{tabularx}{0.7\textwidth}{X} 當孝敬父母，使你得福，在世長壽。這是第一條帶應許的誡命。 \end{tabularx} \\ \\ \relax
6:3 & \begin{tabularx}{0.7\textwidth}{X} 見上節 \end{tabularx} \\ \\ \relax
6:4 & \begin{tabularx}{0.7\textwidth}{X} 作父親的，你們不要激怒兒女，但要照著主的教導和勸戒養育他們。 \end{tabularx} \\ \\ \relax
6:5 & \begin{tabularx}{0.7\textwidth}{X} 作僕人的，你們要懼怕戰兢，用誠實的心聽從你們肉身的主人，好像聽從基督一般； \end{tabularx} \\ \\ \relax
6:6 & \begin{tabularx}{0.7\textwidth}{X} 不要只在人的眼前這樣做，像僅是討人的喜歡，而是作基督的僕人，從心裡遵行神的旨意， \end{tabularx} \\ \\ \relax
6:7 & \begin{tabularx}{0.7\textwidth}{X} 甘心服侍，好像服侍主，不像服侍人， \end{tabularx} \\ \\ \relax
6:8 & \begin{tabularx}{0.7\textwidth}{X} 因為知道每個人所做的善事，不論是為奴的或是自主的，都必按所做的從主得到賞賜。 \end{tabularx} \\ \\ \relax
6:9 & \begin{tabularx}{0.7\textwidth}{X} 作主人的，你們待僕人也是一樣，不要威嚇他們，因為知道他們和你們在天上同有一位主，他並不偏待人。 \end{tabularx} \\ \\ \relax
6:10 & \begin{tabularx}{0.7\textwidth}{X} 最後，你們要靠著主，依賴他的大能大力作剛強的人。 \end{tabularx} \\ \\ \relax
6:11 & \begin{tabularx}{0.7\textwidth}{X} 要穿戴神所賜的全副軍裝，好抵擋魔鬼的詭計。 \end{tabularx} \\ \\ \relax
6:12 & \begin{tabularx}{0.7\textwidth}{X} 因為我們的爭戰並不是對抗有血有肉的人，而是對抗那些執政的、掌權的、管轄這幽暗世界的，以及天空靈界的惡魔。 \end{tabularx} \\ \\ \relax
6:13 & \begin{tabularx}{0.7\textwidth}{X} 所以，要拿起神所賜的全副軍裝，好在邪惡的日子能抵擋仇敵，並且完成了一切後還能站立得住。 \end{tabularx} \\ \\ \relax
6:14 & \begin{tabularx}{0.7\textwidth}{X} 所以，要站穩了，用真理當作帶子束腰，用公義當作護心鏡遮胸， \end{tabularx} \\ \\ \relax
6:15 & \begin{tabularx}{0.7\textwidth}{X} 又用和平的福音當作預備走路的鞋穿在腳上。 \end{tabularx} \\ \\ \relax
6:16 & \begin{tabularx}{0.7\textwidth}{X} 此外，要拿信德當作盾牌，用來撲滅那惡者一切燒著的箭。 \end{tabularx} \\ \\ \relax
6:17 & \begin{tabularx}{0.7\textwidth}{X} 要戴上救恩的頭盔，拿著聖靈的寶劍—就是神的道。 \end{tabularx} \\ \\ \relax
6:18 & \begin{tabularx}{0.7\textwidth}{X} 要靠著聖靈，隨時多方禱告祈求，並要為此警醒不倦，為眾聖徒祈求。 \end{tabularx} \\ \\ \relax
6:19 & \begin{tabularx}{0.7\textwidth}{X} 也要為我祈求，讓我有口才，能放膽開口講明福音的奧祕， \end{tabularx} \\ \\ \relax
6:20 & \begin{tabularx}{0.7\textwidth}{X} 我為這福音的奧祕作了帶鐵鏈的使者，讓我能照著當盡的本分放膽宣講。 \end{tabularx} \\ \\ \relax
6:21 & \begin{tabularx}{0.7\textwidth}{X} 今有親愛、忠心服事主的弟兄推基古，為了你們也明白我的事情和我的景況，他會讓你們知道一切的事。 \end{tabularx} \\ \\ \relax
6:22 & \begin{tabularx}{0.7\textwidth}{X} 我特意打發他到你們那裡去，好讓你們知道我們的情況，又讓他安慰你們的心。 \end{tabularx} \\ \\ \relax
6:23 & \begin{tabularx}{0.7\textwidth}{X} 願平安、慈愛、信心從父神和主耶穌基督歸給弟兄們。 \end{tabularx} \\ \\ \relax
6:24 & \begin{tabularx}{0.7\textwidth}{X} 願所有恆心愛我們主耶穌基督的人都蒙恩惠。 \end{tabularx} \\ \\
[1ex]
\hline
\hline
\end{longtable}
$^{1}$明朝末平安.
我們來到《以弗所書》的第三講.
就是《以弗所書》最後的第五和第六章.
我給一個題目給他.
叫做「受敵和屬靈的爭戰」.
你知不知道.
你有一個敵人.
你問牧師是嗎?.
我有敵人嗎?.
這個敵人很厲害的.
其實他時時對準著你.
你自己都不知道.
這麼厲害的?.
那誰是我的敵人?.
「受敵」就是誰?.
保羅來到最後一章第六章的時候.
就突然之間揭出一件很重要的事.
我相信有很多的教徒.
可能你仍然不知道的.
就是原來這個世界.
是有兩個世界觀的.
一個我們叫做「靈界」.
一個叫做「屬地的世界」.
「世俗的世界」.
所以原來我們需要知道.
有「靈界」和有「屬地的世界」.
保羅有一次曾經說.
我好像被提到三層天上面.
我們若摸估計.
天上面分為起碼有三層.
我想問你第一層天代表什麼?.
就是你現在在外面望上天.
或者晚上.
今晚的月光很光.
這個就是第一層天.
第二層天是什麼?.
就是有些國家派太空人.
飛上月球.
又說去火星.
又說要探測什麼.

$^{41}$「外太空」「太空以外」.
第三層天呢?.
保羅說他去過的.
就是屬於「靈界」.
屬天上的.
記不記得《約伯記》一開頭那裡.
就是說撒旦來到神的面前.
他說我發覺地上有一個人.
很蒙你祝福.
那個就是約伯.
他說「神啊如果你放手給我.
可以攻擊這個約伯.
我擔保他就放棄你了」.
神看到撒旦這樣來挑戰.
他說「好,我准你攻擊他.
但你有某一些限制」.
這個就是屬於「靈界」.
天上面「靈界」的事.
有神.
也有撒旦.
他們是互相為仇的.
所以我們知道.
惡者專門攻擊我們.
好像《約翰福音》裡面說.
耶穌說我是牧羊人.
我是羊的門.
都測利是要偷竊,殺害,毀壞.
這個就是仇敵今天要做的策略.
他要毀壞,攻擊.
他是一個謊言之父.
所以撒旦是一個撒謊的人.
而且是謊言之父.
他最擅長裝模作樣.
所以撒旦可以化作光明的天使來欺騙你.
所以今天我們知道我們有一個仇敵.
仇敵就是要打敗這些信耶穌的人.
所以要向我們背後放箭.
於是保羅說.
我們要知道.
我們今天的仇敵是屬靈氣的惡魔.

$^{81}$所以你打仗要知道和誰打才行.
你是否他的對手.
好像大衛.
爸爸叫他去探望哥哥.
大衛就去.
誰知有個巨人來挑戰.
大衛很生氣.
他說你怎可以辱罵我的永生神.
於是他就想出場打仗.
嚇得哥哥馬上制止他.
你這個小子夠膽和這個巨人打嗎?.
後來就聽到有人肯出場.
記得讓他出場.
但是大衛知道他勝敗不是在乎自己.
勝敗在乎耶和華.
是神讓他得勝的.
所以各位弟兄姊妹.
今天我們有個惡者.
就是撒旦.
撒旦要打敗我們.
叫你跌倒.
叫你犯罪.
叫你離開神.
但是我們基督徒必須要站立得穩妥.
於是保羅在第五章二弗所書裡面說.
你們當效法我.
好像我效法基督那樣.
哇!保羅真是夠膽.
你們學我吧.
因為我正在學基督.
不是說你學我.
因為我學基督.
所以你們可以學我這樣去一起學基督.
然後保羅就勸他們.
你們要憑愛心行事.
你們的言語要小心.
一切惡毒,淫辭,妄語都不可以出口.
各位弟兄姊妹.
今天原來我們要很小心我們的嘴唇.
我聽過有一個課程叫「苗圃百福」.

$^{121}$有八種福氣.
其中第一課就是「口出恩熱」.
你知不知道你的口說的話.
有時是帶著能力的.
雅各書曾經說.
舌頭是百體裡面最細的.
但卻能說最大的謊話.
可以叫一座樹林都燒著.
你小小的火種.
還可以點燃整座樹林.
你說的話很有能力.
我們特別廣東人.
罵人有時真是抵死你.
其實很惡毒.
抵你死.
這樣都說得出.
同時有時我們小時候.
聽長輩罵我們.
特別有些長輩或父母.
我只是老式的那些.
「你這麼蠢」.
有時還加一句「蠢過XXX」.
於是你就覺得.
原來我很蠢的.
我沒用的.
因為那些長輩都罵我.
「你沒用的」.
「又不合格」.
「你沒鬼用」.
還加了「鬼」.
所以這些就是惡毒害人.
你想想真的很毒.
抵你死.
「你蠢到XXX」.
「你沒用啊」.
「你完全沒用」.
於是這些就進入我腦裡.
「我沒用」.
好像一直打針一樣.
於是你就覺得你沒有自信心.

$^{161}$今天有很多人在他成長的過程裡.
就受到這些仇敵的攻擊.
但如果你稱讚一個人.
這樣就不同了.
「你今天很漂亮」.
「你真是聰明」.
「我怎麼搞都不行」.
「你一下子就幫我解決了」.
「你真是很厲害」.
「我很欣賞你」.
「你是很醒目」.
一稱讚別人.
當場整個人都醒了.
真的不同了.
原來口出祝福的說話.
叫人成長.
鼓勵人.
激勵人.
聖經裡面說.
我們要多講欣賞別人.
鼓勵人的說話.
口裡面一句惡言.
都不可以出口.
特別是舊約的《箴言》.
講到很多關於我們怎樣控制自己的舌頭.
言語如金蘋果落在銀網子裡面.
良言是良藥.
所以求神幫助我們.
時時口裡面講這些祝福.
鼓勵人.
帶領人的說話.
不是去攻擊惡毒性的說話.
這就是保羅給我們的教導.
所以原來聖經是教很多細節的.
我們多講祝福的說話.
我聽過一位牧師說.
有個人他在家中種了很多盆栽.
他日日有時去幫盆栽除草.
適當時候就淋水.
每一次為盆栽淋水的時候.

$^{201}$祝福這盆盆栽.
盆栽盆栽你長得很靚.
你帶給我很多喜樂.
我見到你就開心.
你越來越長.
真的很奇妙.
這盆盆栽就越生越靚.
有一次出街就見到門口突然有野草.
很醜的樣子.
於是他對野草說.
野草你快點萎縮.
不要在我面前阻礙我.
讓我覺得很不開心.
那些野草很快就枯乾枯死了.
他說原來我的祝福很有能力.
叫這盆植物長得越來越靚.
我一咒詛這盆野草就死了.
固然這是一個有少許說笑形式.
但如果你真的這樣對一位家人.
你試試鼓勵他.
你試試讚他.
特別我們做父母的.
對子女.
子女拿成績表回來.
他苦瓜乾的面.
你也猜到一定不方好成績.
但你說似乎你都盡了力.
你都拿到59分.
雖然不合格.
但你用讚賞.
拿到59分.
總比拿40多分好.
你也好.
下次再努力一點.
下次可以拿70分.
他馬上很痛.
因為原來你沒有罵他.
這麼蠢的.
這次又不合格.
因為你是用一種批評的語言.

$^{241}$不是一個鼓勵的角度.
求神讓我們是口出人言.
所以是有靈界.
原來魔鬼很多時候.
用我們這些口中一句話.
打低一個人.
但如果我們講一些祝福.
從神而來的話.
就可以激昂一個人.
所以以弗所書到了第六章.
保羅來到結束的時候.
他就很鄭重勸告.
你們要認識.
你們有空中掌權的惡者.
這個惡者是在空中來的.
他真的很厲害.
所以你要小心.
於是你要穿上全副的軍裝.
就能抵擋那惡者一切的攻擊.
所以今天我要問你.
請問你有沒有穿上全副的軍裝?.
全副軍裝是什麼?.
會不會背?.
不是穿一件,兩件,三件.
是七件軍裝.
全部都要穿上.
而且你要很詳細了解.
到底這些軍裝講什麼?.
我發覺教會裡有些弟兄姊妹.
知道一點,不知道一點.
是啊,你穿上全副軍裝.
我已經穿上了.
什麼叫那些軍裝?.
保羅的講解是非常詳盡的.
第一.
你要束上帶子.
以真理當作帶子.
古代人的箭步是鬆身的.
所以走動很不方便.
唯一的方法是用腰帶一紮紮實.

$^{281}$整個人當作醒目了.
衣服又會繞著你.
你走路也不行.
束起.
不要全部放下來阻礙你走路才行.
所以束腰是很重要的.
我們中國人說要識紮馬步.
你學功夫就要找你的紮馬.
所以用真理束腰.
什麼是真理?.
神的道在你心中.
你有真理.
你站在真理那邊.
詩篇第一句就說.
不聽惡人的說話.
不站罪人的道路.
然後不助洩萬人的作為.
不聽先聽.
然後第二不站.
不要站在那裡.
不要坐在那裡.
你是祂的同黨.
求神幫助我們.
有真理在我們裡面.
我們以真理束腰.
第二.
護心鏡遮空.
心是我們最重要的器官.
被人一槍打過來.
鋸成武槍.
射箭.
射中紅心.
射中你的心.
中箭而身亡.
所以你有一個護心鏡遮空.
什麼是我們的護心鏡呢?.
就是公義.
當作護心鏡.
你有正義在你裡面.
我站在正義那邊.

$^{321}$我為正義說話.
我不給那些假道理騙了我.
因為仇敵經常找些假東西來.
我公義識分辨.
什麼是真.
什麼是假.
然後第三.
你穿什麼鞋.
那雙軍靴很重要.
因為你的腳要站得很穩.
因為走很多路.
打仗.
一滑腳.
整個人就滑下去了.
所以你的鞋很重要.
不要穿那些會脫爪的鞋.
全部是很穩陣的鞋.
你的鞋讓你走得很穩妥.
讓你走得很淡定.
以平安的福音.
當作走路的鞋.
福音追隨著你.
你去到哪裡就傳福音.
去到哪裡就講神的平安.
去到哪裡就講神的大愛.
這個就是你的鞋.
你的腳踵是全陽福音.
而福音跟隨著你的腳踵.
第四.
你有一個盾牌.
古代的盾牌.
就不像今天那些矮矮的.
那些是長盾牌.
你如果矮一點.
縮了下去.
整個人都遮住了.
不知道大家有沒有看電影.
拍古裝片.
那些盾牌震.
前面一辣.

$^{361}$大家拍著.
就好像整面牆一樣.
所有的兵就跪下.
簡直看不到.
那排箭射過來.
全部擋走了.
這個就是聖經裡面所講的.
盾牌.
以什麼?.
以信德.
就是你的德行.
就是你的招牌.
擋住了.
我很有信用的.
我說這樣就這樣了.
我照做的.
我不會騙人的.
我不會說謊的.
我不會沒有信用的.
我答應你什麼時候交貨.
我就交貨.
哇!好守信用.
信德做你的盾牌.
可以抵擋敵人一切的火箭.
所有的箭射過來.
全部擋了.
這邊又擋.
那邊又擋.
上面又擋.
下面又擋.
全部擋了敵人的箭.
但是有一位神的輩人真的講得好.
沒錯!盾牌擋了敵人的箭.
但是如果射向你的背部.
是的.
後面沒有東西擋.
所以敵人很厲害.
暗箭難防.
在你背後射過來.
但是如果你有一群的箭友.

$^{401}$全部用一個擋.
用一個盾牌震.
就全部抵擋了.
古代有些人作戰是很出名的.
其中一個就是斯巴達人.
哇!斯巴達的軍隊.
真的用盾牌震.
上下前後全部擋住.
好像一隻烏龜一樣.
原來這個盾牌震中間有很多士兵.
它一下子全部擋住.
就全部得寶.
當敵人一退的時候.
就開一個盾牌震.
就衝出去打敗敵人.
所以各位我們需要靠弟兄姊妹互相的守望.
彼此的連結.
所以就組成一個盾牌震.
前後左右全部包圍.
我們都得到保障.
然後還有頭盔.
頭很重要.
「救恩造你的頭盔」.
我已經得救.
我已經得到神的拯救.
我是被稱為義的.
我是神的兒女.
所以我就頂住一切.
今天魔鬼怎麼攻擊.
吩咐他退去.
因為我是神的兒女.
我是神的兒子.
你是神的女兒.
所以神愛他的兒女.
所以魔鬼要走.
記不記得耶穌在四十天曠野禁食的時候.
魔鬼三次來試探他.
他最後是吩咐撒旦退去吧.
因為他有勸兵.
吩咐魔鬼要走.

$^{441}$然後第六是聖靈做我們的劍.
聖靈的寶劍.
這裡真的要解釋一下.
原來「劍」這個字.
在希臘文裡是指短的劍.
短的劍有什麼用呢?.
很多時候我們中國的劍是長劍.
看武俠片.
一拔劍出來.
特別是日本的武士.
那把劍很鋒利.
一削下來.
什麼都削得爛.
但是聖經裡說的「寶劍」是短劍.
即是在你的腰上插著一把像刀子一樣.
長的劍是打遠方的敵人.
比如你對我,我對你.
但到靠近肉搏的時候.
就要靠你的腰上拔一把短的劍.
一把一把插在他的心上.
什麼是我們的短劍呢?.
聖經的說話.
就是我們聖靈的寶劍.
神的話就是我們的寶劍.
突然之間.
用一句金句.
隨口而出.
我靠著那家,我力量之.
凡事都能作.
這句就給你很多勇氣了.
當敵人攻擊你的時候.
不怕.
我靠著神.
凡事都能作.
今天神是我的保障.
祂是我的高台.
祂是我的避難所.
我就宣告神的話.
於是能力就從上而來給你.
有很多人都說.

$^{481}$當我起困難的時候.
我背聖經.
所以各位.
你必須要背聖經.
背到滾瓜爛熟.
因此當你在不同的敵人攻擊下.
用聖經的金句打敗他.
你隨口而出.
耶和華必保護我.
耶和華是我的目的.
我必不自缺乏.
我雖然走過死刃的幽谷.
我不怕遭害.
就是這樣背出來.
記得.
這些金句.
要藏在心中.
到時找聖經.
那句聖經好像記得一點.
又不記得一點.
那是沒有用的.
所以背經是很重要.
我在一個牧師的家庭長大.
所以自小已經參加主學.
就懂得背很多聖經.
原來背的聖經.
弟弟哥哥背的.
最記得清楚.
神的話是我腳前的燈.
路上的光.
你或向左或向右.
神會指示你.
這是正路.
行在其間.
二人的路.
好像黎明的光.
越照越明.
教養孩童使他走.
當行的路.
到老他都不會偏離.

$^{521}$記得很清楚.
所以求神幫助我們.
要用我們的短劍.
一拔出來就刺死了敵人.
記不記得耶穌在曠野接受撒旦的試探的時候.
他是接受三個試探.
三個試探.
都是引用聖經.
打退受敵.
第一個叫他用石頭變餅.
耶穌就說.
人活著.
不是單靠食物.
乃是靠神口裡所出的話.
第二句.
第二個試探.
他說.
你如果拜我.
我將天下榮華給你.
耶穌又引用.
當拜主你的神.
然後又說.
你聖殿殿頂.
跳下來.
神會猜天使都接住你.
耶穌又再引用.
不可試探主你的神.
大家知不知道這三句聖經.
全部引自哪裡.
同一卷書.
可能你不知道.
原來耶穌引用的三句聖經.
全部出自申命記.
摩世五經裡面的.
創出黎民生.
全部在申命記裡面.
請問你申命記有沒有從頭到尾讀過一次呢?.
你懂得申命記懂得多少呢?.
雖然你沒有劍可用.
那就被敵人殺死.

$^{561}$各位聽眾我們一定要很熟悉神的話.
神的道在我們心裡.
你隨時可以拔劍.
靠著聖靈的寶劍.
打敗仇敵.
所有兵器裡面.
唯有劍是你的武器.
然後最後第七.
靠著聖靈.
隨時多方的禱告祈求.
亦都為眾聖徒祈求.
禱告是我們一個最大的武器.
所以七樣東西.
前面五樣是你身上的裝備.
第六樣是你的劍.
然後第七樣藉著禱告多方祈求.
然後保祿說你們都要為我祈禱.
就完結了這卷這麼寶貴的以弗所書.
我希望我講完之後.
你都愛上以弗所書.
好好地讀.
我們讀聖經真的要很仔細.
很用心地讀.
將神的話藏在心裡面.
神的話就在你的生命中產生功效.
願天父賜福給各位.
\newpage



\section{以賽亞書 50:4-9}
\label{sec:031908}
\textbf{誰能定我有罪呢?}~~~~林佩華姊妹
\newline
\newline
連結: \href{http://www.church.com.hk/acms/content.asp?site=cdc\&op=show\&type=product\&code=031908\&layout=sermon}{\texttt{031908}} ~~~~ 語音日期: 2026-03-29
\newline
\newline
\hyperref[sec:031907]{\small{< < < PREV SERMON < < <}}
~
\hyperref[sec:index_chronic]{\small{[返順時目]}}
~
\hyperref[sec:031909]{\small{> > > NEXT SERMON > > >}}
\newline
\newline
以賽亞書 50:4-9
\newline
\begin{longtable}{cl}
\hline
\hline
章節 & 經文 (和合本修訂版)\\
\hline
50:4 & \begin{tabularx}{0.7\textwidth}{X} 主耶和華賜我受教者的舌頭，使我知道怎樣用言語扶助疲乏的人。主每天早晨喚醒，喚醒我的耳朵，使我能聽，像受教者一樣。 \end{tabularx} \\ \\ \relax
50:5 & \begin{tabularx}{0.7\textwidth}{X} 主耶和華開啟我的耳朵，我並未違背，也未退後。 \end{tabularx} \\ \\ \relax
50:6 & \begin{tabularx}{0.7\textwidth}{X} 人打我的背，我任他打；人拔我兩頰的鬍鬚，我由他拔；人侮辱我，向我吐唾沫，我並不掩面。 \end{tabularx} \\ \\ \relax
50:7 & \begin{tabularx}{0.7\textwidth}{X} 主耶和華必幫助我，所以我不抱愧。我硬著臉面好像堅石，也知道我必不致蒙羞。 \end{tabularx} \\ \\ \relax
50:8 & \begin{tabularx}{0.7\textwidth}{X} 稱我為義的與我相近；誰與我爭論，讓我們來對質；誰與我作對，讓他近前來吧！ \end{tabularx} \\ \\ \relax
50:9 & \begin{tabularx}{0.7\textwidth}{X} 看哪，主耶和華必幫助我，誰能定我有罪呢？看哪，他們都要像衣服漸漸破舊，被蛀蟲蛀光。 \end{tabularx} \\ \\
[1ex]
\hline
\hline
\end{longtable}
$^{1}$各位弟兄姊妹午安.
在分享整個話語之前.
讓我們一同的同心低頭禱告.
我們恭敬的將以下的時間交上.
求你開開我們的耳.
讓我們可以聽得到你的話.
由於你直接的話去教導我們.
更加讓我們有悟性去明白真理.
以致我們可以欣著你的話.
生命被激勵.
得著更生改變.
求你的令引導我們.
聽我們的禱告.
是奉耶穌基督的名士求.
阿們.
多謝站員組剛才帶來的詩歌.
其中一首都是我預備經文時一直聽著的.
亦多謝Julia很有力地讀出這段經文.
將神的話語活然在我們的眼前.
活然在我們的耳裡.
今天是教會年曆的第六主日.
亦是中之主日.
紀念耶穌光榮進入耶路撒冷城的事蹟.
當時耶穌騎著一隻驢仔進入耶路撒冷城.
遠號的群眾揮舞著中樹枝.
叫喊著「和撒拿!和撒拿!」.
一個榮耀的君王要進度.
這個星期亦是聖周的開始.
是耶穌基督在世最後一個星期.
在螢幕上.
你見到耶穌是很浩浩蕩蕩的君王進入耶路撒冷城.
但當祂進入耶路撒冷城之後.
接著的一個星期是有很大的反差.
因為祂將會以最謙卑的形象.
最信服的形象.
為世人的罪釘死在十字架上.
我們看到神的揀選.
祂所命定的僕人.
原來不是高高在上.
不是用權力來操控眾生.

$^{41}$而是苦就卑微謙卑信服.
就像耶穌那樣.
以致受苦為我們寫命.
今天我們一起去看的這段經文.
所說的僕人.
亦是像耶穌基督那樣.
是一個信服的僕人.
我們今天會看《以塞瓦書》.
裡面說到以塞瓦先知.
以塞瓦是主前八世紀的先知.
是南國猶大人.
他的名字的意思是耶和華拯救.
先知的工作是傳遞神的話語.
包括對人的警告,提醒.
但亦同時會傳達神要拯救人的一份應許.
整本《以塞瓦書》有66章.
分開了兩部分.
第一部分是1至39章.
當時是主前八世紀.
以色列和猶大受到亞述國統治的時期.
這一部分是講述神公義的審判和刑罰.
第二部分是40至66章.
當時已經是主前六世紀.
當時猶大亡國.
聖殿被毀.
百姓被擄至巴比倫.
他們的生活不容易.
亦失去了一份對神的信心.
雖然後來神的慈愛公義.
讓他們可以回歸去重建他們的聖城.
重建他們的聖殿.
但不是很多人願意去聽.
願意去跟隨.
第二部分是講述神的恩典,拯救和安慰.
在第二部分當中.
有四首「獨人詩歌」.
分別是42章1至9節.
講述獨人的身份和使命.
第二首是記載在49章1至7節.
講述獨人的普世使命.

$^{81}$而第三首和第四首「獨人詩歌」.
是講述這位獨人怎樣的馴服.
甚至受苦.
為了要承傳神所交付的使命.
而今日我們所看的.
就是第三首「獨人詩歌」.
第50章1至3節講述什麼呢?.
50章好像一個法庭的場景.
第1至3節講述以色列民.
因為不聽神的教導,犯罪,得罪神.
所以遭到神的審判和懲罰.
最後亡國.
被露到巴比倫.
這群百姓被露到巴比倫.
在絕望,信心的動搖中.
因著他們的信仰警報.
神施慈愛.
透過他們的獨人.
亦是我們今天所講的第4至9節.
給這群以色列民鼓勵和馴服的教導.
我們先看第4至5節.
主耶和華賜我受教的舌頭.
使我知道怎樣用言語輔助被罰的人.
主每天早晨喚醒我的耳朵.
使我能聽像受教者一樣.
主耶和華開啟我的耳朵.
我並未違背,並未退後.
經文是以主耶和華開始.
在第4至9節裡.
總共出現了4次主耶和華.
分別是第4,第5,第7和第9節.
而剛剛我們所讀的兩節經文裡.
提到我們不同的器官和他們的功用.
第4節提到主耶華賜僕人有受教的舌頭.
以致知道怎樣用言語輔助被罰人.
主又喚醒我的耳朵.
使我能聽像受教者一樣.
第5節提到主耶和華開啟我的耳朵.
在第4節開始說主耶華賜受教者的舌頭.
受教者的英文翻譯為Disciple.

$^{121}$是門徒的意思.
意思是門徒受到裝備領受耶和華的話.
受教者的舌頭是透過門徒的口傳遞神的話語.
這位僕人先學習神的話.
用他的耳朵聆聽神的吩咐和教導.
目的就是要輔助軟弱疲乏的人.
之前提過先知是代言人.
要傳達神的話.
所以對於僕人來說.
舌頭和耳朵都非常重要.
第4節下半節提到主每天早晨喚醒我的耳朵.
使我能聽像受教者一樣.
主每天喚醒僕人顯示僕人和神的關係十分密切.
另外是神主動喚醒他的僕人.
每天讓他和神相交.
是神讓他的僕人能夠聽見他的話.
能聽是一個馴服.
是願意領受神的旨意.
作為神的僕人必須能夠被神教導.
由神賜給他有受教者的耳朵.
讓他能夠聽見神的聲音.
亦給他受教者的舌頭.
以致他可以用神的話去輔助靈裡軟弱疲乏的人.
第5節提到主說:開啟我的耳朵.
開啟我的耳朵亦象徵著馴服.
這位神的僕人必須要聽得見.
亦願意馴服進行.
後半節說:我並沒有違背也未退後.
這種語調你聽得出僕人的決心.
僕人明知道自己作為先知的代言的即是.
他絕對不會違背,不會退後.
違背的意思是反叛.
但這位僕人絕對不會反叛耶和華.
無論聽人的反應是好還是壞.
這位僕人都知道他要忠誠去宣講神的話.
當面對人的攻擊的時候.
他都不會退後.
他願意一生代言和輔助有需要軟弱的人.
這就是他的生命使命.
不知道我們是否也像這位受教者一樣.

$^{161}$我們願意每天與神親近.
願意去聽神的話.
謙卑去領受祂的教導.
更加願意遵行.
用受教者的舌頭去幫助身邊的人.
第六節說到:人打我的背,我任他打.
人拔我兩下的鬍鬚,我由他拔.
人侮辱我,向我吐唾沫,我並不掩面.
這位信奉的僕人去聽神話.
也是遵行神的旨意.
但結果他需要受苦.
第六節說到這位忠誠的僕人.
如何受到一些痛苦.
第六節說到這些受苦包括了肉體被人打.
也是心靈上的受苦,被侮辱,吐唾沫.
在古代的近東,男人的鬍鬚代表著他的尊嚴.
將人的鬍鬚拔下來或剃掉.
也是一個很大的侮辱.
人打我的背,我任他打.
人拔我兩下的鬍鬚,我由他拔.
人侮辱我,向我吐唾沫,我並不掩面.
這位僕人受到這一連串的羞辱.
但他沒有反擊.
反而容讓一群敵人將他淫而凌辱.
有色經學者說,打和羞辱這位僕人的人.
可能是被擄的以色列群體.
他們不滿意僕人呼籲他們離開安書的巴比倫.
返回危險的曠野路,返回耶路撒冷.
也有學者說,打僕人的人可能是巴比倫人.
因為在第四十六章和第四十七章裡.
僕人不停說巴比倫帝國將會滅亡的信息.
所以巴比倫人不喜歡他,要打他.
在耶利米書三十七章十四至十五節裡說.
以利亞不聽耶利米的話就逮捕他.
把他帶到官場那裡.
官長們擄擄耶利米打了他.
把他囚在約拿丹民事的房屋中.
因為貪玩者屋子,當作監牢.
在聖經記載當先知傳講神的話的時候.
百姓通常都不喜歡聽.

$^{201}$耶利米所說的話.
好像這兩節的聖經所說.
曾經都被百姓拒絕.
在四呼音裡也說到耶穌出來傳道的時候.
都是受盡藐視,辱罵和敵視.
但耶穌忍受了這些侮辱.
被賣,被審,交到外邦人的手中.
被鞭打,被奚落,被羞辱.
這六節的經文所說的一切.
都曾經施行在耶穌基督的身上.
但耶穌基督就是為了罪人的緣故.
包括你,包括我.
忍受了諸般的羞辱.
這一份是耶穌基督要彰顯的.
長闊高深的愛.
第六節的經文我們看到.
僕人定要信服.
願意為主受苦.
無怨言,無反悔.
因為他深知是上帝所吩咐的.
所以無論遇到任何困難.
他都不會後退.
這首是第三首的「僕人之歌」.
當去到第四首「僕人之歌」的時候.
我們會看到進一步描述到.
僕人將要受到什麼樣的痛苦.
什麼樣的凌辱.
以致這位僕人可以完成神的救贖.
完全信服成為神代罪的羔羊.
做贖罪祭成為世人的一份救恩.
第七至九節裡這樣說.
主耶和華必幫助我.
所以我不抱愧.
我硬著臉面好上堅尺.
也知道我不自蒙羞.
請我為義的與我相近.
誰與我爭論讓我們來對質.
誰與我作對讓他近前來吧.
看那主耶和華必幫助我.
誰能定我有罪呢.

$^{241}$看那他們要像衣服漸漸破舊.
被重曙光.
在這三節經文裡.
三節提到主的同在.
第七節主耶和華必幫助我.
第八節稱我為義的要與我相近.
第九節主耶和華必幫助我.
誰能定我的罪呢.
作為神的僕人.
力量的源頭是源自神的同在.
這位僕人很清楚知道.
神和他的關係.
在這份心力的關係裡.
他能夠忍受.
他能夠領受.
有心智堅定有勇敢去面對一切一切.
剛才第六節提到受苦的目的.
是要叫神的僕人蒙羞.
但這位僕人受苦.
但是不蒙羞.
因為第七節提到.
這位僕人知道耶和華一定會幫助他.
所以僕人不抱愧.
因著耶和華的幫助.
他的表現是硬著臉.
好像堅石.
一點都不怕.
一點都不退縮.
硬著臉好像堅石.
表示僕人願意肩擔這個重任.
一個不容易的任務.
他有毅力.
有堅忍.
可以繼續下去.
僕人知道他不自羞愧.
第八節提到稱僕人為義.
是主耶和華.
因為耶和華判斷他是公義的.
而且與他相近.
第八節裡面.

$^{281}$爭論對質.
作對.
這些都是法庭的用語.
這好像是一個法庭的場景.
控告僕人的.
是對抗神的勢力的一班人.
這些人雖然說話厲害.
但在神面前其實不堪一擊.
為受苦僕人做辯護的就是神.
裡面提到因為神與這位僕人相近.
然後第九節.
有一連串的發問.
誰可以與他爭論.
誰可以定我有罪.
誰可以判定他有理.
在之前42章1至4節說.
看那我的僕人.
我所父慈所揀選.
心裡喜悅的.
我已將我的靈賜給他.
他必將公理傳給萬邦.
他不灰心不喪膽.
直到他在地上設立公理.
眾海島都等候他的訓誨.
主耶和華給僕人的使命.
就是要他在地上設立公理.
更加要將公理傳給萬邦.
所以在第九節.
那一連串的發問裡面.
僕人再次表明了.
他對於自己的身份和任務.
都有很明確的認定.
所以雖然身體上受到摧殘.
受到痛苦.
但是他反而與這一群反對者.
有正面的交鋒.
提出質問.
誰能定我有罪呢?.
意思是沒有人可以指控他.
因為耶和華判斷他是公義的.

$^{321}$他是進行耶和華交給他的使命.
第九節下半節這樣說.
他們都要將衣服漸漸破舊.
被蛛蟲所咬.
是比喻僕人所面對的挑戰.
反對 控告他的.
都不會長久.
只是短暫 最終都會滅亡.
當僕人受到這一連串.
不公平的對待的時候.
他可以讓神為他伸冤.
他認定神一定會幫他.
因為沒有人能夠定他的罪.
沒有人可以在法庭的爭論裡.
勝過這位神.
有時他可以像堅石一樣.
去面對敵人的羞辱.
面對敵人的擊打.
僕人受苦的忍耐的力量是源於神.
他深信眼前這一切的不公平.
終有一天都會畫上句號.
句號之後是通往神恩典的審判.
最終擊打他 惡待他的人.
都會被神審判.
神必定會為他伸冤.
第九節再次提到主耶和華要幫助我.
這些描述都讓我們看到.
神是公義.
當我們面對不了解的痛苦.
不能夠明白的經歷的時候.
神就用祂的大能為我們主持公道.
我們知道神就是掌權者.
在任何情況下我們都可以仰賴祂.
無論我們有多少問號.
多大的困難.
我們都可以仰賴這位神.
以致我們可以勝過我們的仇敵.
這幾年世界的情況.
香港的情況都不是我們所想像.
無論經濟 工作 生活的環境.

$^{361}$甚至教會都有很多大大小小的困難.
我不知道在你的日常生活裡面.
你面對著什麼挑戰.
我曾經聽過弟兄姊妹說.
他們在工作上會被上司批評.
因為做不完的事.
又不知道自己的位置.
是否可以繼續保留.
因為經濟不好會減少預算.
但面對著的是永遠做不完的工作量.
人一路咳 頭一路咳.
他們就是面對這些的情況.
我不知道你面對著的是什麼情況.
在《以塞瓦書》裡面.
僕人盼望和相信神.
讓他能夠繼續面對眼前的苦難.
並且可以堅持忍耐.
等候到上主的公義 平安和更新來臨.
這首僕人之歌.
講到儒表著受苦的彌賽亞.
到了第四首的僕人之歌.
我們更加會看到很立體的.
這位僕人將會受人侮辱 唾罵和鞭打.
但他甘願為我們付出生命的代價.
2016年有一套電影叫《鋼腳靚》.
我打錯字 應該是《Hacksaw Ridge》.
是由《真人真事》改編的.
講述第二次世界大戰期間.
有個叫Desmond Dawes的人.
因著信仰的緣故去參軍.
但他堅持不拿槍 不殺人.
Dawes這份堅持.
結果受盡其他人的欺凌 鄙視.
甚至還未上戰場已經要入獄.
但最後因為他的堅持.
終於可以以軍醫的身份.
與美軍一起出戰.
在一場重要戰事中.
他一晚救了75個傷兵.
每救一個他都會對神說.

$^{401}$請你給我救多一個.
於是他一個一個的傷兵救了回來.
起初他有時也不明白.
到了危急的時候 他會問神.
你想要我做些什麼.
但儘管他有疑問 但他沒有退縮.
就是靠著神這份恩典.
去渴望進行神給他的使命.
去盡力去救每一個傷兵.
用行動去表達對神 對身邊人的愛.
直到最後一刻 他就因著這份堅持去做到.
因為他很清楚自己要做什麼.
也是為了什麼原因要這樣做.
所以周圍人的反應.
周圍人對他的欺凌都變得不重要.
電影的一開始有這句經文.
是《選治爾賽亞書》40章30至31節.
「就算年輕人也要疲乏困倦.
強壯的也必全然跌倒.
但那等候耶和華的必重新得力.
他們必如鷹展翅上騰.
他們奔跑卻不困倦.
行走卻不疲乏」.
當我們要去學習謙卑馴服的時候.
有時候我們真的會經歷到當中的黑暗.
經歷到當中的苦難 當中的痛苦.
甚至好像這位僕人一樣被擊打.
很正常我們都會有困惑 有疑問.
甚至灰心 沮喪.
但只要我們堅定不移的信靠這位大能的神.
最後神會將亮光按著他的心意.
按著他的方式賜給尋求他的人.
面對艱難的時候.
儘管我們不明白 會疲乏 會困惑.
而我們還是能夠持守一份盼望.
相信主耶和華必幫助我們.
因為我們所經歷的 他一早已經知道.
也一直與我們同在.
他會幫助 施下力量和安慰.
在此很想和大家分享一首詩歌.

$^{441}$叫「我心靈得安寧,it is well with my soul」.
是詩人Spetford施比福所寫的.
這位詩人原本是一位成功的律師.
也是熱心教會的服侍.
在1871年一場大火中.
他所有投資的房地產都已經燒了.
損失慘重.
之後沒多久 他四歲的兒子也因病離世.
兩年後 他原本與太太和女兒一起去歐洲度假.
同時參與一個報道大會的事奉.
但臨出發前 他因為有其他事務要處理.
於是請太太和女兒按照原先的計劃.
搭船先起程.
誰知遊船去到法國威爾斯近海的時候.
被一艘英輪撞沉了.
只有太太生存.
詩人接到這個消息.
隨即去到英國坐船.
去到遐渴這個出事的地點.
去到這個地點他不是埋怨.
而是感受到神安慰他心中的傷痛.
他所領受到的就是這一切都是神安排.
這份安慰讓他的心靈得安寧.
於是他寫了這首詩歌.
詩歌裡有四節.
每一節都有不同的心情轉折.
第一節講到不論何環境都可以知足.
第二節講到撒旦雖然來侵.
眾試煉雖來臨.
但在試煉裡堅定的相信.
第三節回想我中罪.
全釘在十字架上.
他想到自己是不配.
但他的罪因著耶穌基督被釘十架的緣故得釋放.
讓他的心寬暢.
讓他不自主的從心底對神發出讚美.
第四節私人仰望主在你.
表達無論世上有多少悲傷和痛苦.
但因著一份堅信.
他在患難裡仍有平安和永恆的盼望.

$^{481}$在這裡很想和大家分享一個小小的見證.
在一月初開始.
我反覆感冒咳嗽失聲.
又咳嗽又失聲.
到二月初有一天早上睡醒.
兩隻耳朵好像被吸了一隻耳機.
當有人和我說話的時候.
我聽到有回音.
如果一起吃飯的時候.
是一個災難.
因為很辛苦.
聲音很嗡.
所以那段時間真的寢食不安.
但神為我預備了一個很好的耳鼻喉科醫生.
他幫我找到那個原因.
我講漏了除了嗡之外.
早上起床是嗡.
但到晚上我聽到有電流聲.
第二天我的耳朵聽到有兩條聲音.
如果你是懂音樂的話.
你知道半音和全音.
但我聽到四分之一.
我特意去記錄那些頻率究竟是多少.
有些是四分之一.
有些是八度.
有些是沒有聲音.
當時我在想.
糟了,怎麼辦?.
我應該找誰?.
我當時第一樣想的就是.
如果我沒有耳朵.
我可以怎樣事奉?.
是耳咽管濕了.
他給了一些藥給我.
也幫我中了細菌.
但他告訴我這些細菌有抗藥性.
他預計這些情況有機會再發生.
結果真的發生了.
在發生之前.
這段時間我曾經去搭飛機.

$^{521}$我也猜測會不會有問題.
但我忘記了問醫生.
所以就憑著信心上了機.
結果真的很辛苦.
我從來沒有試過搭飛機耳朵這麼辛苦.
當我預備這篇講章的最後兩天.
我的耳朵又開始有點矇.
我又開始聽到這篇電流聲.
我當時不是害怕.
而是覺得又來?.
而是去祈禱.
我同樣這樣想.
假如我真的要再次經歷這件事.
甚至我最忠告耳朵聽不到話的時候.
我知道我仍然可以事奉.
我聽不到外界的聲音.
但我一定會聽到神的聲音.
我知道這些經歷是神給我的考驗.
讓我更加去依靠祂.
更加去跟隨祂.
就好像詩歌剛剛第一次這樣說.
我們有時會在平安.
但有時會有困難.
會有浪途.
但無論是什麼境況.
我們的心靈都因著神得安寧.
讓我們同心一誓.
有個圖告.
主我們知道前面的路可能很崎嶇.
可能會波濤洶湧.
可能有很多困難.
但我們深知你會引導我們.
你會幫助我們.
祈求你繼續帶領我們.
讓我們看到你要我們為你做什麼工作.
不單讓你成為我們的依靠.
讓我們在前面的道路上帶著一份盼望.
更加讓我們有這份信心.
這份使命.
去將這份心恩厚愛廣傳.

$^{561}$特別我們知道在這段時間.
有很多不同的人.
都有經歷著困難.
經歷著傷心,哀痛.
祈求你一天去安慰.
我們深願你的愛.
有一天臨到他們當中.
讓他們得著這份盼望.
一份永恆的生命價值.
求你繼續使用我們.
聽我們圖告.
在奉耶穌基督的名次求.
阿們.
\newpage



\section{創世記 24:50-67}
\label{sec:031909}
\textbf{活於應許的故事裡}~~~~翁曉怡傳道
\newline
\newline
連結: \href{http://www.church.com.hk/acms/content.asp?site=cdc\&op=show\&type=product\&code=031909\&layout=sermon}{\texttt{031909}} ~~~~ 語音日期: 2026-03-28
\newline
\newline
\hyperref[sec:031908]{\small{< < < PREV SERMON < < <}}
~
\hyperref[sec:index_chronic]{\small{[返順時目]}}
~
\hyperref[sec:031910]{\small{> > > NEXT SERMON > > >}}
\newline
\newline
創世記 24:50-67
\newline
\begin{longtable}{cl}
\hline
\hline
章節 & 經文 (和合本修訂版)\\
\hline
24:50 & \begin{tabularx}{0.7\textwidth}{X} 拉班和彼土利回答說：「這事既然出於耶和華，我們不能向你說好說歹。 \end{tabularx} \\ \\ \relax
24:51 & \begin{tabularx}{0.7\textwidth}{X} 看哪，利百加就在你面前，可以將她帶去，遵照耶和華所說的，給你主人的兒子為妻。」 \end{tabularx} \\ \\ \relax
24:52 & \begin{tabularx}{0.7\textwidth}{X} 亞伯拉罕的僕人聽見他們這些話，就向耶和華俯伏在地。 \end{tabularx} \\ \\ \relax
24:53 & \begin{tabularx}{0.7\textwidth}{X} 僕人拿出金器、銀器和衣服送給利百加，又將貴重的物品送給她哥哥和她母親。 \end{tabularx} \\ \\ \relax
24:54 & \begin{tabularx}{0.7\textwidth}{X} 然後，僕人和隨從的人才吃喝，並且住了一夜。早晨起來，僕人說：「請讓我回我主人那裡去吧。」 \end{tabularx} \\ \\ \relax
24:55 & \begin{tabularx}{0.7\textwidth}{X} 利百加的哥哥和母親說：「讓她同我們再住幾天，也許十天，然後她可以去。」 \end{tabularx} \\ \\ \relax
24:56 & \begin{tabularx}{0.7\textwidth}{X} 僕人對他們說：「耶和華既然使我道路亨通，你們就不要耽誤我，請讓我走，回我主人那裡去吧！」 \end{tabularx} \\ \\ \relax
24:57 & \begin{tabularx}{0.7\textwidth}{X} 他們說：「我們把她叫來問問她。」 \end{tabularx} \\ \\ \relax
24:58 & \begin{tabularx}{0.7\textwidth}{X} 他們就叫了利百加來，對她說：「你和這人同去嗎？」她說：「我去。」 \end{tabularx} \\ \\ \relax
24:59 & \begin{tabularx}{0.7\textwidth}{X} 於是他們送他們的妹妹利百加和她的奶媽，同亞伯拉罕的僕人，以及隨從他的人走了。 \end{tabularx} \\ \\ \relax
24:60 & \begin{tabularx}{0.7\textwidth}{X} 他們就為利百加祝福，對她說：「我們的妹妹啊，願你作千萬人的母親！願你的後裔得著仇敵的城門！」 \end{tabularx} \\ \\ \relax
24:61 & \begin{tabularx}{0.7\textwidth}{X} 利百加和她的女僕們起來，騎上駱駝，跟著那人去。僕人就帶著利百加走了。 \end{tabularx} \\ \\ \relax
24:62 & \begin{tabularx}{0.7\textwidth}{X} 那時，以撒住在尼革夫。他剛從庇耳‧拉海‧萊回來。 \end{tabularx} \\ \\ \relax
24:63 & \begin{tabularx}{0.7\textwidth}{X} 傍晚時，以撒出來，到田間默想。他舉目一看，看哪，來了一隊駱駝。 \end{tabularx} \\ \\ \relax
24:64 & \begin{tabularx}{0.7\textwidth}{X} 利百加舉目看見以撒，就急忙下了駱駝， \end{tabularx} \\ \\ \relax
24:65 & \begin{tabularx}{0.7\textwidth}{X} 對那僕人說：「這從田間走來迎接我們的人是誰？」僕人說：「他是我的主人。」利百加就拿面紗蓋住自己。 \end{tabularx} \\ \\ \relax
24:66 & \begin{tabularx}{0.7\textwidth}{X} 僕人把他所做的一切事都告訴以撒。 \end{tabularx} \\ \\ \relax
24:67 & \begin{tabularx}{0.7\textwidth}{X} 以撒就領利百加進了母親撒拉的帳棚，娶了她為妻，並且愛她。以撒自從母親離世以後，這才得了安慰。 \end{tabularx} \\ \\
[1ex]
\hline
\hline
\end{longtable}
$^{1}$大英哲會 晚安.
我們先一起祈禱.
掌管著時間和歷史的神.
多謝你願意將我們納入你的故事裡面.
你不但願意稱我們為兒女.
更願意用我們這群不完美的人.
在你的故事裡面.
在你的救贖工作裡面.
願你在今天的訊息裡面.
讓我們一起去領受你給我們的訊息.
讓我們更加去認識自己 認識神.
以致我們都可以憑信.
活在你的故事裡面.
祈禱.
我前陣子去看了一套電影.
那套電影是所有人進去都會哭出來.
不知道大家有沒有看過一套這樣的電影.
因為看完之後覺得很感動.
我在裡面哭死了.
出來之後後勁未過.
所以不斷在社交媒體上看他們上載的花絮.
裡面有一些導演的訪問片段.
他講到有很多大家最終很深刻.
一說起就很想哭的一些畫面.
其實都是劇本沒有寫.
而是一班演員去揣摩和投入角色去發揮出來的.
那些演員會說如果是什麼什麼角色的話.
在這裡他們應該會這樣反應.
或者他們去描述那時怎樣出現這個畫面.
就是他們演戲的時候.
因為很投入角色當時的心境.
一些劇本沒有的對白他們就脫口而出.
看完之後.
看到一個好的劇本.
會在故事的發展裡面塑造一個立體的角色.
但往往他們不會將每一句說話和每一個表情寫上去.
他們的花絮上都有對比著那個劇本.
和他們真的做出來講的對白.
好的演員就會在這個故事發展裡面去投入角色.
並且完整對這個角色的想像.

$^{41}$而一個好的導演就會和演員一起去討論.
究竟怎樣演戲會比較合理.
去引導演員去思考.
去推進故事的合理發展.
我覺得這是一個很好的比喻.
去討論人怎樣去活出一個有神的生命.
神按著我們的本性.
幫我們每一個人都度身訂做了一個角色.
有我們獨特的設定.
可以在神的故事裡面佔一席位.
亦都參與在這個故事裡面和其他角色互動.
發揮一點點推進這個故事的作用.
如果我們願意用心去揣摩我們的角色.
和這個故事的發展路向.
又經常和作為導演和編劇的神去交流想法.
慢慢我們在一些沒有明確指令的位置.
都可以交出合理和豐富的故事.
甚至可以令導演都覺得很欣賞的一些表現.
如果「以殺取妻」是整個創世的一個很長的劇集的其中一個部分.
今天我們所說的就是「以殺取妻」的最終回.
很想和大家整個故事,整個廿四章再看一次.
我們一起看看不同的人物如何帶著他們自己的獨特.
在大概理解神的心意之後.
順應著向前推進故事的發展.
亦都嘗試學習他們如何以信心去見到神信實的引領.
所以如果大家有實體或電子聖經的話.
我們都邀請你們打開創世的廿四章.
方便我們待會來回在故事的不同情節之間.
整個「以殺取妻」的部分.
由亞伯拉罕(以殺的爸爸)老邁開始.
當亞伯拉罕想起神的應許如果要延續的話.
停在「以殺」是不足夠的.
「以殺」也要有後代.
所以他將「以殺」要娶老婆這件事交託給一個老僕人.
亞伯拉罕比起他以前.
他沒有再因為急切要完成神的應許.
而選擇妥協在現實的限制裡.
這次他堅持「以殺」的妻子必須是他本族的「人」.
而不可以是迦南地的女子.
亦堅持「以殺」一定要留在這個應許之地.

$^{81}$不可以回到他們本族的家鄉.
而老僕人的任務有兩個部分.
第一要找到合適的人選.
第二要將這個女子帶回來.
之後老僕人就帶著他的策略去找尋他少主的妻子人選.
他很有可能在出發前已經有對一個理想妻子特質的想法.
以致他選擇帶十匹駱駝去.
除了承載他的聘禮之外.
亦成為他考驗人選的一個工具.
由老僕人到拿學的城 可能是哈蘭開始.
他完成任務的兩部分結構都很相似.
都是他先到達一個場景.
然後實行他的策略.
兩次他都可以得到頗明顯順利初步的成功.
然後有一個懸念.
最後得到答案任務完成.
任務的第一個部分是他要找到一個適合的妻子.
當老僕人到達哈蘭.
決定以一個祈禱作為他的策略.
他定下的條件其實不容易達成.
因為他需要那個女子主動提出一些東西.
但他的禱還未祈完 李百嘉已經出現了.
又真的按他所禱告.
不單打水給他 還自請餵駱駝打水.
一個如此神奇的順利令老僕人安靜下來.
然後他想起還有一個懸空的問題.
李百嘉究竟是不是他主人本族的人呢?.
當他問過之後發現是.
老僕人立即以敬拜讚美神的信實手揚.
肯定這一切都是因為有著神的引領和護佑.
才可以如此順利成就.
而當老僕人在拉班的帶領下進入了他們的家.
就進入了這個任務的第二個部分.
這次他要爭取將李百嘉帶回迦南地.
而他這次的策略是見證分享.
他將亞伯拉罕如何蒙神賜福.
以撒如何作為他會承繼一切的兒子.
亞伯拉罕對他的要求.
他如何來到迦南還未祈完禱就已經遇上李百嘉.
以及如何應驗了他的祈禱.

$^{121}$他連自己的敬拜讚美也包含在故事中敘述一次給拉班看.
以致拉班和比土尼都看得出這件事是出於耶和華.
而他們就同意並完成了送聘禮的步驟.
不過這次的懸念到了第二天早上的時候.
亞伯拉罕提出可否讓李百嘉和他們多住十天才離開.
亞伯拉罕堅持要離開.
他們拉扯之間就叫了李百嘉來問.
「你現在走不走?」.
李百嘉就離開了並祝福了李百嘉.
到了李百嘉和以撒見面的時候.
老套人應該有將整個故事再說一次.
然後以撒除了李百嘉並愛他.
也標誌著以撒和李百嘉在創世的大故事中.
承繼了亞伯拉罕和撒拉開展新一代的故事.
延續神的應許.
雖然這個故事往往被稱為以撒娶妻.
但以敘事而言.
最主要的角色是老僕人和李百嘉.
他們的性格特點比較鮮明.
老僕人在這個故事中沒有記下他的名字.
但他在這個故事中存在感最強.
也說了最多的話.
所以我們可以看到他是一個相當忠心的人.
因為他在得到任務時有認真對待.
他有認真想過一次這個任務中可能會遇到什麼事.
所以他提出了一個難題.
第五節說到如果那個女子不肯來怎麼辦.
他要知道亞伯拉罕的優次原則.
究竟是不帶那個女子來還是我帶你兒子去.
他也有想過如何完成這個任務.
所以他帶了駱駝和一些平禮.
他有想過用禱告和見證去完成任務.
而他在他的禱告和見證中.
也用了「我的主人」來稱呼亞伯拉罕.
他也很堅持盡快完成任務.
減低一切變數.
以致他選擇先分享他的見證說明來意.
先去接受拉班對他的款待.
第二天就馬上想走.
這些都顯出了他對亞伯拉罕的忠心.

$^{161}$對完成他主人給他的任務的堅定.
可能當中也包括了.
他已經確定當中有神的引領.
他知道自己選擇的人是對的.
正在走的方向是對的.
而希望謹慎快速地去完成.
不過其實每次看到這個故事中.
這個僕人堅持要說出他的長長的見證才吃飯.
又堅持第二天馬上就要走.
我也會想.
這個做法是不是唯一正確的做法呢?.
他先吃了飯再說見證.
選擇李伯嘉會不會帶來什麼改變呢?.
如果他讓李伯嘉和他的家人多留十天.
對事情的發展又會不會有影響呢?.
他的決定究竟是不是唯一正確的道路呢?.
當然這些都是假設性的問題.
我們也不會知道.
但其實亞伯拉罕的吩咐中.
沒有強調過這件事最好快點搞定.
神令老僕人的道路亨通.
事情順利.
也不必然代表神期望老僕人要馬上回程.
我們作為知道後續故事.
有上帝視覺的人回頭看.
可能也會覺得拉班那麼狡猾.
當然要又快又準快點走去王道.
但當時老僕人不知道拉班那麼狡猾.
更加有可能老僕人本身是一個謹慎得來.
希望快點做完事情的一個人.
所以他會在祈禱的面前求神.
讓我今天遇見那個娶妻的人選.
他也會很有策略性地.
還不知道李伯嘉是誰的女兒.
就已經送了金環金手續給他.
又在向拉班介紹亞伯拉罕的時候.
特意提及亞伯拉罕的財物和老僕.
他既具體化了神的賜福.
其實也實際地吸引拉班去同意這頭婚事.
他留宿一晚後就馬上決定回程.

$^{201}$比起判斷他急速的節奏是否恰當.
是否唯一正確的做法.
我相信更加重要的是.
我們見到老僕人是一個沒有上帝視覺的人.
他決定這樣做可能和他的性格有關.
就像我們一樣.
都是一個真實獨特.
會因著我們的經驗性格而去決定如何行事的人.
經文也沒有對他的做法下很多評論.
但我們可以見到神有使用他所做的事.
成就了這一個應許延續的故事.
當中的重點應該是老僕人有敏銳於神信實的供應和帶領.
也會為見到神的信實而敬拜讚美神.
他雖然會以我主人亞伯拉罕的神去稱呼耶和華.
但他也會直接向耶和華禱告尋求指引.
也會在看得出有神的成就而敬拜讚美.
他不單止作為一個旁觀者去見證神對亞伯拉罕當年的看顧和私恩.
他也肯定自己可以預見到一個合適的女子是有神的帶領在裡面.
所以他會在27字裡說.
「耶和華,我主人亞伯拉罕的神是應當稱頌的,因他不斷以慈愛信實待我主人」.
然後還加了一句「至於我,耶和華一路引領我,直到我主人的兄弟嫁女」.
他在48字裡又重覆一次「我主亞伯拉罕的神,因為他引導我走合適的道路」.
「使我得著我主人兄弟的孫女,及我主人的兒子為妻」.
他在56節說要帶尼白加走的時候也有說.
「耶和華,既然使我道路亨通,你們就不要耽誤我」.
所以老僕人有自己去經歷神.
他對神的信靠不單止幫助他成功完成任務.
也讓他在當中很深刻地經歷到神的信實.
在這章講了很多東西的老僕人.
可能最記得的是當日他連禱都未祈願.
尼白加已經出現在他眼前,應驗了他的祈求.
令他沒什麼好說,看得出是耶和華令自己道路亨通.
跟很多說話的老僕人相比,尼白加說的不是很多.
但他一出現,觀眾也會猜到他是故事的女主角.
因為他的名字在22章23節已經得以提及過.
跟老僕人有點不同,尼白加是一個行動派.
而且是一個大方勇敢的人.
他在打水遇到一個完全陌生的男人老僕人.
請求給他一些水.
他作為一個未婚的女性,沒有任何害羞或不確定.

$^{241}$他二話不說就拿下瓶子托在手上給他喝水.
也主動提出為駱駝刀打水.
然後他就將瓶子裡的水倒在槽裡.
又跑到井旁打水,為所有駱駝打了水,直到喝足了.
他本身是一個熱心款待人,也會照顧生畜需要的人.
這份特質剛好跟老僕人在他的禱告中所看中的相符.
而他在23節回答了老僕人有關自己的身份和家可否給他住的問題後.
他馬上跑回家.
在這一刻,他能否判斷到金環和金手鐲是聘禮的一部分.
還是他只覺得是一份謝禮,我們不太知道.
但以他這個熱心款待的特質,可能他更在意的是預備招待老僕人的一行人.
同時也想將自己遇到的事告訴家人.
老僕人在34至49節的見證,利伯嘉很大可能是跟他的家人一起聽著.
因為51節拉班說:漢娜,利伯嘉就在你面前.
所以他應該是有一起聽.
也跟拉班和彼套利一樣,看得出這是出於耶和華.
而他也願意順服神的引領.
以至於他在58節很簡潔有力地回答:我去.
在希伯來文裡只有一個字.
他就是這樣簡潔無條件地帶著奶媽和呂卜離開自己的家,面對很多未知.
他不知道以薩是一個怎樣的人.
不知道以薩的外貌是怎樣,性格是怎樣.
不知道跟以薩一起生活會是怎樣.
他這個決定也很難逆轉.
但僅僅是因為看得出這件事出於耶和華.
利伯嘉就願意以信心毅然前行,不拖延.
顯出他非一般的決心和勇氣.
而在他終於見到以薩的時候,他就急忙下了駱駝.
而知道以薩就是他的結婚對象的時候,他就拿面紗蓋著自己.
這兩個動作都是表示他作為一個妻子,願意謙卑在丈夫面前.
也是側面地替他再剔多一行在聖經裡的好妻子特質.
利伯嘉就這樣決定.
對於利伯嘉的哥哥拉班,不知道大家對他又有什麼印象呢?.
一般我們應該比較記得他在後期的絞結.
因為令本身欺騙人的雅各在這裡跌了一跤,上了很久的一課.
在二十四章裡,他的絞結沒那麼突出.
雖然學者都頗一致地認同.
拉班很可能在三十節見到金環和戴在他妹妹手上的金額的時候.
就已經覺得這是一頭好勳士.
而在三十四至三十六節.

$^{281}$僕人有一點刻意提到亞伯拉罕的財產.
和這一切財產都會給了二世的時候.
都可能真的提到進入拉班的心裡.
不過他也未必很想自己的妹妹願嫁.
五十至五十一節,拉班和彼土利對魯僕人說完見證後的回應是.
這事既然出於耶和華,我們不能向你說好說歹.
漢娜尼伯加就在你面前,可以將她帶去.
遵照耶和華所說的,給你主人的兒子為妻.
他看得出耶和華是這件事裡的人靈.
以至無論他覺得這件事是好還是不好.
他都沒有什麼可以說.
不知道大家有沒有試過.
既然是這樣,既然很明顯是神的意思.
那就唯有這樣吧.
不過他睡了一晚後.
他表示希望多留尼伯加十天才出發.
可能他想想一些對策,令尼伯加不用嫁那麼遠.
又或者他真的很疼妹妹.
希望在告別之前可以祝福她多一點.
可以多相處一陣子.
可能這兩個可能性都有.
不過相信他真的疼尼伯加這個妹妹.
因為他們當時往往都是父母長輩去處理決定婚姻大事.
但他們願意將最終的決定權交給尼伯加.
而當他這個很有想法,很堅定的妹妹說他去.
他也從善如流,為她祝福告別.
將決定權交給將駕的女子.
在整個世界不是一個尋常的做法.
他也呼應了二十四章一開始的僕人去問亞伯拉罕.
如果難女子不幸,而不是她的家人不幸.
都埋下了尼伯加做最終決定的伏筆.
而尼伯加這麼堅定決定嫁的爾薩.
在整個故事的存在感相對都很低.
他更多的時候是作為亞伯拉罕的兒子.
我主人的兒子,或者我的主人去被提及.
而他以名爾薩被提及的次數很少.
我們也不知道為什麼當時已經是成人的爾薩不自己去尋找妻子.
而要他父親的盧浦仁幫他尋找妻子.
不過怎樣也好,在這張有限的描述裡.
爾薩是一個安靜,順從的人.

$^{321}$當他聽了盧浦仁將整個故事複述一次的時候.
他也看出這件事出於耶和華.
所以他就領尼伯加進了母親撒拉的帳篷.
娶了她為妻並且愛她.
67節寫道爾薩自從母親離世以後.
才得了安慰.
可能是指爾薩在撒拉離世後受傷.
剛剛好在尼伯加來的時候結束受傷期.
也有可能是指他與尼伯加建立自己的家庭後.
有尼伯加的陪伴就走過了母親離世最難過的時候.
我們不知道.
但尼伯加與爾薩的個性很不同,很互補就很明顯.
即使爾薩自己沒有很大參與尋找妻子的過程.
但盧埔人帶回來的是一個可以與他互補.
而他也會愛的女子.
也獻出了神在當中又信實地去護佑.
尼伯加本身不是在亞伯拉罕這個文藥的家庭裡的一個人.
但他在神的揀選和護佑下.
得以進入這個藥,這個應許當中.
在故事裡我們也看到尼伯加與亞伯拉罕或撒拉.
也有一些相似或互相呼應的敘述.
尼伯加在第十八二十節裡急忙拿下瓶子托在手上.
急忙把瓶子裡的水倒在槽裡.
這個表現與亞伯拉罕在十八章接待耶和華和使者時.
急忙進帳棚叫撒拉做餅.
又跑到牛群裡親自揀選又嫩又好的牛交給僕人去預備的敘述很相似.
所以尼伯加與亞伯拉罕是很相似的熱心款待客類特質.
而二十四章的四節,七節,三十八節和四十節.
都是相關於十二章一節.
耶和華去呼召亞伯拉罕離開父家和本族的地的呼召相當有關.
而去到二十四章要離開父家和本族的地.
就由亞伯拉罕變成尼伯加.
而尼伯加也與亞伯拉罕一樣離開後沒有再回去.
而尼伯加在離開前得到的祝福.
「我們的妹妹,願你作千萬人的母親,願你的後裔得著受敵的城門」.
與撒拉在之前十七章神要改名的時候所給他的祝福.
「他要興起多國,必有百姓的君王從他而出」.
亦是很類近的方向.
他顯示出尼伯加承受的應許與撒拉承受的一樣.
甚至他之後也與撒拉一樣經歷了一段不辱的艱苦.

$^{361}$才在求告神後懷上雙胞胎.
尼伯加的信心和堅定.
令他在神的護佑當中進入了這個應許的故事裡.
甚至比起以撒可能更像承繼了亞伯拉罕與撒拉的某些表現.
甚至他之後神對他兩個兒子的預言都是直接對尼伯加說的.
以撒娶妻這個故事可以順利進行.
有神信實的引領固然是一個必要的條件.
但當中每個角色或多或少的信心和信服亦都很重要.
這個故事也讓我們看到帶有信心的每個人.
按著自己的性格,特質,經驗都會有不同的實踐方式.
但只要願意憑信信服 總會見到神的信實引領.
可能是在別人的見證裡去印證.
可能是自己的禱告被應運.
可能是其他方式去見到神的信實引領.
神很奇妙地將不同特質的人放在同一條故事線裡.
尼伯加剛好是一個有心有力款待人又照顧生畜的人.
符合了僕人覺得理想的條件.
拉班剛好是一個會被僕人提及財產吸引.
但又因為看得出是神的引導而不敢強硬留住尼伯加.
或是尼伯加和以撒的性格剛好互補等等.
這些好像剛好的安排.
他們之間有的細微互動.
都是獨特於這群有這樣的特質的人才會有這樣的互動.
以致他們一起成就了延續神應許的故事.
不知道大家會不會有時我們看聖經的人物.
都覺得他們有信心又看到神的信實帶領.
都好像很理所當然.
看回當然覺得他相信當然應該相信.
但仔細看亞伯拉罕和神納了約.
知道神會由撒拉得到一個兒子繼承他.
他會成為萬國之父.
但神沒有特別指示他何時生仔.
他後來才知道.
神亦沒有特別指示他何時為以撒娶妻.
亞伯拉罕的僕人知道自己的任務.
是在神的私恩下找到主人本族的女子.
但神亦沒有指示他如何找到.
或者告訴他他中間會經歷甚麼.
李伯嘉在僕人的見證中.
聽到自己進入婚姻是有神的心意在其中.

$^{401}$但神亦沒有指示他如何與以撒相處.
也沒有一來應許他很快生子.
其實這些經歷可能與我們很相近.
今天我們有整本聖經告訴我們.
神對世界的心意是一個大的方向.
有我們對神,對人互動的大原則.
但實際淺行上.
可能神也有給我們一些發揮的空間.
按照我們自己的性格特質.
獨特的處境和經驗.
一些經歷.
在與神和其他人的互動中.
我們可能有發揮的空間.
如果我們可以像路僕人一樣.
時常回望自己的經歷.
而看得出上帝的引領.
見到上帝的信實.
並為此敬拜讚美神.
或者我們像尼伯加一樣.
勇敢地看得出神的心意.
憑順去信服.
我們也可以在這個故事中繼續推進.
我們未必在每一件事上.
都可以得到神很確實地告訴我們.
這一下應該怎麼做.
但如果我們認識自己.
在我們的經歷.
在神的話語中.
看到的方向.
我們可以做到一些適合的決定.
最令我們不在神的故事中.
是耶穌到城獄身.
一次過附上贖駕.
成就救恩.
使我們得以進入應許的故事.
就像尼伯加不在這個故事中.
是服人去找他.
讓他聽得出.
自己有神的心意.
令他進入這個故事中.

$^{441}$加上他的信心和堅定.
他就進入了應許的故事.
我們也跟尼伯加一樣.
並不是停在進入故事中.
尼伯加沒有停在那裡.
他繼續和以撒去相處.
有他們的婚姻.
生了他們的兒子.
在繼續知道神的心意後.
用他的方法令應許延續.
當然他之後做的方法是否正確.
之後的道應該會提到.
但故事一定是繼續延續下去.
我們進入了應許.
進入了神的約之後.
也是繼續延續下去.
我們也活在這個繼續向前的故事中.
我們也要繼續和神互動.
去確認他的心意.
確認我們的方向和做法正確.
也要和身邊的人彼此印證.
見到大家都在走正確的方向.
見到大家的實踐.
也在神給我們的框架裡.
當我們盡量嘗試成為一個好角色.
在神的故事裡按著大方向去發揮.
我們就可以一起活出一個有感染力.
也吸引到人想加入這個劇組的故事.
願我們在神信實的護佑下.
都可以平順生活.
也在這個動態一直向前.
我們不能停在某些經歷.
在生命的真實裡.
我們也要和信仰群體彼此印證.
彼此支持.
一起平順活出一個可以感染到人.
吸引到人進入舊印的故事.
\newpage



\section{馬可福音 15:42-16:8}
\label{sec:031910}
\textbf{從恐懼中走出來}~~~~禤凱榮牧師
\newline
\newline
連結: \href{http://www.church.com.hk/acms/content.asp?site=cdc\&op=show\&type=product\&code=031910\&layout=sermon}{\texttt{031910}} ~~~~ 語音日期: 2026-04-05
\newline
\newline
\hyperref[sec:031909]{\small{< < < PREV SERMON < < <}}
~
\hyperref[sec:index_chronic]{\small{[返順時目]}}
~
\hyperref[sec:031911]{\small{> > > NEXT SERMON > > >}}
\newline
\newline
馬可福音 15:42-16:8
\newline
\begin{longtable}{cl}
\hline
\hline
章節 & 經文 (和合本修訂版)\\
\hline
15:42 & \begin{tabularx}{0.7\textwidth}{X} 到了晚上，因為這是預備日，就是安息日的前一日， \end{tabularx} \\ \\ \relax
15:43 & \begin{tabularx}{0.7\textwidth}{X} 有亞利馬太的約瑟前來，他是尊貴的議員，也是盼望著神國的，他放膽進去見彼拉多，請求要耶穌的身體。 \end{tabularx} \\ \\ \relax
15:44 & \begin{tabularx}{0.7\textwidth}{X} 彼拉多詫異耶穌已經死了，就叫百夫長來，問他耶穌是不是死了很久； \end{tabularx} \\ \\ \relax
15:45 & \begin{tabularx}{0.7\textwidth}{X} 既從百夫長得知實情，就把耶穌的身體賜給約瑟。 \end{tabularx} \\ \\ \relax
15:46 & \begin{tabularx}{0.7\textwidth}{X} 約瑟買了細麻布，把耶穌取下來，用細麻布裹好，安放在巖石中鑿出來的墓穴裡，又滾來一塊石頭擋住墓門。 \end{tabularx} \\ \\ \relax
15:47 & \begin{tabularx}{0.7\textwidth}{X} 抹大拉的馬利亞和約西的母親馬利亞都看見安放他的地方。 \end{tabularx} \\ \\ \relax
16:1 & \begin{tabularx}{0.7\textwidth}{X} 過了安息日，抹大拉的馬利亞、雅各的母親馬利亞，和撒羅米，買了香料，要去膏耶穌的身體。 \end{tabularx} \\ \\ \relax
16:2 & \begin{tabularx}{0.7\textwidth}{X} 七日的第一日清早，太陽出來後，她們來到墳墓那裡， \end{tabularx} \\ \\ \relax
16:3 & \begin{tabularx}{0.7\textwidth}{X} 彼此說：「誰要替我們把石頭從墓門滾開呢？」 \end{tabularx} \\ \\ \relax
16:4 & \begin{tabularx}{0.7\textwidth}{X} 她們抬頭一看，看見石頭已經滾開了，原來那石頭很大。 \end{tabularx} \\ \\ \relax
16:5 & \begin{tabularx}{0.7\textwidth}{X} 她們進了墳墓，看見一個年輕人坐在右邊，穿著白袍，就很驚奇。 \end{tabularx} \\ \\ \relax
16:6 & \begin{tabularx}{0.7\textwidth}{X} 那年輕人對她們說：「不要驚慌！你們尋找那釘十字架的拿撒勒人耶穌，他已經復活了，不在這裡。來看安放他的地方。 \end{tabularx} \\ \\ \relax
16:7 & \begin{tabularx}{0.7\textwidth}{X} 你們去，對他的門徒和彼得說：『他要比你們先到加利利去，在那裡你們會看見他，正如他從前所告訴你們的。』」 \end{tabularx} \\ \\ \relax
16:8 & \begin{tabularx}{0.7\textwidth}{X} 於是她們出來，從墳墓那裡逃走，又發抖又驚訝，甚麼也沒有告訴人，因為她們害怕。 \end{tabularx} \\ \\
[1ex]
\hline
\hline
\end{longtable}
$^{1}$今天是復活節.
大家都知道復活節是紀念,慶祝耶穌的復活.
我相信如果大家是基督徒的話.
每年最少聽一次耶穌復活的故事.
我們有時候做講員.
真的不知道從什麼新的角度來說.
同一段經文,同一個故事.
不過我想問大家一個問題.
究竟我們都熟悉耶穌基督復活的故事.
但是究竟這個復活的故事.
如何影響我們今天的生活呢?.
上個星期大家都知道.
我去了中國度假,在中國旅遊.
在我旅遊的時候.
我聽到新聞的風聲.
美伊打算襲擊伊朗的基礎設施.
大家有沒有聽到?.
我馬上跟師母說.
週末一定開火.
為什麼呢?如果大家有留意的話.
特朗普最喜歡在週末開始軍事行動.
大概是想週末不會太過影響股市或者燃油市場.
正如我所說.
趁在週末的時候.
大家看新聞都應該知道.
在週末的時候.
伊朗兩個的大鋼鐵廠.
一座的發電廠.
民用的核設施等等.
就受到攻擊.
亦如我們所期望.
伊朗亦採取報復的行動.
一連串的軍事行動.
你打我,我打你.
我們發覺.
我們今天不單是活在第三次世界大戰的可能性裡面.
更加是活在一個隨時一觸即發的核子戰爭的危機裡面.
用聖經的語言來說.
我們今天全世界都是活在一個末世的時代.
正如聖經所提醒我們.

$^{41}$在末世的時候會有什麼情況發生.
國家與國家之間有戰爭.
民族與民族之間有仇恨.
天災人禍多到不得了.
如果我們想這些的時候.
真的情緒立刻低落.
我們要問的就是.
究竟耶穌基督如果真是復活的話.
祂所復活的能力如何.
讓我們可以在一個沒有平安的世界裡面.
可以得到平安呢.
而耶穌基督復活的能力如何.
讓我們在一個應該沒有恐懼的世界裡面.
如何在恐懼裡面走出來呢.
今天我們從一個比較實踐式的角度.
來看一看復活的故事.
在十字架上為我們犧牲.
聖經告訴我們.
在星期五.
就是他斷氣的那一天的晚上.
聖經說有一個阿彌馬太的浴室.
去到彼拉多當時的羅馬巡撫那裡.
來要求彼拉多去領取耶穌基督的身體.
我們要問這個浴室是誰來的.
馬海就告訴我們.
他是一個尊貴的義士.
這個尊貴的義士一直等候神的國的來臨.
而另外一個福音講座就告訴我們.
這個不單止是這個公會裡面的義士.
也是一個耶穌的門徒.
不過他家具是暗暗地做門徒.
為什麼要偷偷摸摸地做門徒呢.
聖經就說因為他懼怕猶太人.
大家有沒有覺得奇怪.
這樣就怪了.
是不是這個浴室所隸屬的公會.
當日去審判耶穌.
定耶穌死罪的呢.
原來在一個這樣的公會裡面.
竟然有一個耶穌的跟隨者.

$^{81}$我們就問.
你既然在公會裡面.
你既然是這個義士.
為什麼當日判耶穌死罪的時候.
這個浴室不出聲.
或者他有.
不過聖經沒有說.
又或者他當時嚇到不得了.
已經完全沒有了.
打算去為耶穌說話的可能了.
我們不知道為什麼.
因為聖經沒有說.
不過有一樣東西我們可以很肯定的.
就是當門徒因為耶穌的死.
而恐怕到不得了.
而四圍逃走的時候.
這個浴室是唯一一個.
可以夠勇氣走到彼拉多面前.
要求取回耶穌的身體.
無論是什麼原因.
這個舉動是需要很大的勇氣的.
同意嗎?.
為什麼需要勇氣呢?.
因為耶穌基督是怎麼被釘在十字架上的.
他是被釘在十字架上.
是以一個罪犯的身份.
所以當這個浴室要回耶穌基督的身體的時候.
就是說認同這個罪犯.
大家不要忘記.
他是一個尊貴的義士.
當他這樣做的時候.
他是將他的名聲,地位.
放在一個很大的風險裡面.
隨時可能被人踢出公會.
可以看到這個冒險是相當大的.
既然這麼大的冒險.
他得到什麼好處?.
我想了很久.
我都想不到這個浴室.
在這個舉動裡面.

$^{121}$有什麼實質的好處,利益呢?.
沒有.
除了一樣.
無愧的良心.
所以他在這個恐懼與勇氣的面前.
他選擇勇氣.
我們不知道在這件事之後.
他是不是真的被公會踢了出去.
但是有一樣可以很肯定的.
在這件事情發生之前.
他是一個暗暗地做耶穌的門徒.
但是這件事之後.
不能再暗暗地了.
所有人都馬上知道.
原來你就是耶穌的跟隨者.
大家記不記得耶穌曾經有一番說話.
耶穌說.
「你們」.
這個「你們」指的是跟隨耶穌的人.
所有跟隨耶穌的人.
就是世界的光.
什麼意思?.
他說世界的光應該照在人的面前.
叫世人看到門徒的好行為.
就將榮耀歸給他們在天上的父.
中國有一句說話.
紙是包不住火的.
大家有沒有試過用紙包住火?.
包不住的.
如果我們有聖靈的火.
燃點在我們生命裡面的時候.
世界上沒有一樣的東西.
可以將這個聖靈的火.
或者聖靈可以引發我們人生的工作果效.
沒有辦法可以隱藏得了.
頂主目想問大家一個問題.
你在你的工作地方.
你在你的家庭.
你在社交的圈子裡面.
你的朋友.

$^{161}$你的親人.
知不知道你是基督徒?.
有些基督徒就很不想公開他們的基督徒身份.
為什麼?.
公開了之後.
我隨時會被人取笑的.
我隨時被人用一些很怪的眼光來看我.
對我不不不友好的對待的.
拿我著數的.
曾經有一個見證就說.
那時的十多歲的信徒.
特別是在這些年紀的信徒.
很多時候我祈禱.
叫了一碟燒鵝髀飯.
來到的時候祈禱張開眼睛.
不見了雞腿的.
大家試過沒有?.
但是我們要明白一件事.
真實的門徒.
真實耶穌的跟隨者.
是需要付代價的.
當藥室要求.
拿回耶穌基督的身體的時候.
馬上公會裡面的每一個人都知道.
這個藥室是耶穌的門徒.
所以縱然.
他能夠留在這個公會.
也許從當天開始.
所有公會裡面的人.
都用一種奇怪的眼光.
望這個藥室.
他可能因此而失去.
他在公會裡面的聲譽和地位.
不單止是這樣.
如果我們再看聖經的時候.
馬太福音是這樣記載的.
馬太福音說.
這個藥室拿了耶穌的身體之後.
他就將這個耶穌的身體.
安放在自己的新的墳墓裡面.

$^{201}$新的墳墓裡面.
這是一個甚麼墳墓.
這是一個在一個大的石頭裡面.
慢慢造出來的一個墳墓.
大家可以想像一下.
怎樣造.
人手造.
所以這個墳墓一定貴到難以想像.
要這麼辛苦才造出一個這樣的墳墓.
所以即是說甚麼.
即是說普通人買不起這個墳墓.
是一定要很有錢的人.
才可以買得到這樣的墳墓.
這樣就應驗了.
在耶穌幾百年之前的一個先知的預言.
這個先知叫做以塞瓦.
以塞瓦怎樣說.
他說這個耶穌.
他雖然未行強暴.
口中亦都沒有詭詐的說話.
但是人使他與惡人同埋.
接著是甚麼.
誰知死的時候與財主同葬.
就是這麼巧妙.
幾百年前對耶穌的預言.
就是這樣應驗了.
即是說甚麼.
即是說不單止約瑟真是願意.
他的聲譽地位.
亦都是受到很潛.
很有可能的損失.
亦都是真金白銀付出代價的.
因為一個很昂貴的墳墓.
頂主夢我要大家知道一件事.
救恩是免費.
但絕不是廉價的.
昨天.
昨天我去一個老人院.
我和一個老人家去探望他.
順便和他再一次做了一個決志祈禱.

$^{241}$做完祈禱之後.
他就看著我.
這樣就可以了嗎.
這樣就可以信耶穌了嗎.
我說是呀.
你想怎樣.
你不用付入會費.
不用填申請表.
不用做這樣不用做那樣.
你只要信.
但我和他提醒一件事.
是真心的信.
就這麼簡單信耶穌.
不過.
我們的救恩雖然是免費.
但卻不是廉價.
甚麼意思.
因為耶穌基督是為我們的緣故.
付上一個很大的代價.
由於耶穌基督為我們付了一個.
這麼沉重的代價.
在十字架上為我們犧牲了他的生命.
所以今天我和你.
能夠不需要付分毫的代價.
就能夠承受這個救恩.
但我們還是需要留意一件事.
在我們信起耶穌之後.
在我們領受了這個永恆的福分之後.
當我們行走整個人生道路的時候.
那個時候我們需要付代價.
也許我們付的代價是身體上面的.
精神上面的.
物質上面的.
或者是關係上面的.
但是整條的人生路上面.
我們是需要付代價.
弟兄姊妹.
約瑟在這裡告訴我們.
當我們在恐懼.
疑惑.

$^{281}$憂慮的時候.
我們需要好像約瑟那樣.
拿出勇氣.
正如這個路加福音所說.
耶穌說.
如果你真的有人想跟隨我的時候.
你要捨己.
捨己.
天天背起他的十字架.
來跟從我.
意思就是說.
我們需要提出我們每一天生活的勇氣.
從我們的恐懼裡面.
背起我們的十字架.
從恐懼走出來.
因為知道我們人生真正的價值.
人生真正的價值在哪裡.
人生真正的價值.
不在於我們擁有世界上有多少東西.
有多少間屋.
銀行裡面有多少錢.
不在乎這些.
也不在乎我們人生究竟有多長命.
真正在乎的是什麼.
真正在乎的是.
我們的生命是不是放在天父上帝的手裡面.
如果我們有這樣的認知.
我們有這樣的勇氣.
我告訴大家.
我們在一個充滿恐懼的世界裡面.
我們就能夠可以在恐懼裡面走出來.
我們再看.
在約瑟拿到耶穌的身體之後.
聖經說他就用細麻布裹好.
然後安放在盆石中.
作出來的墳墓裡面.
也都大概是因為.
可能是羅馬政府的命令.
他們將一個很大的石頭.
就滾到墳墓的門口那裡.

$^{321}$擋住墓門.
整個的過程.
都是被莫大拉的瑪麗亞和約西的母親瑪麗亞.
這兩個婦女看得清清楚楚.
他們知道.
耶穌放在哪裡.
他們都知道墓門是被一個大石擋住.
大概是因為整個的過程.
都是相當的很緊拙.
很短時間裡面就做了所有的事情.
所以聖經就說.
沒錯 耶穌的身體被包裹了之後.
但是耶穌的身體並沒有用香料來糕末.
所以也是這個原因.
聖經說在安息日之後.
大家知道安息日是星期六.
在安息日之後.
接著那天是星期日的早上.
三個的婦女.
莫大拉的瑪麗亞.
雅各的瑪麗亞.
和薩羅米.
三個婦女就帶著香料.
天未亮就出發.
去耶穌的墓那裡.
打算去糕末耶穌.
我想大家留意.
這三個婦女.
她們出發之前.
很清楚 因為她們看到的.
她們出發之前.
已經知道一個很大的石頭.
是放在墓門那裡.
大家知不知道這個石頭有多大.
或者有多重.
猜一猜.
根據第一世紀的考古學發言.
一個猶太人的墳墓.
一個這樣的石頭.
通常都是圓的.

$^{361}$好像這個圖上面所顯示.
是圓形的.
直徑大概通常是一兩米.
而厚度是有二十到三十厘米.
重量呢.
這樣的石頭大概是一到兩噸左右.
一噸是什麼觀念.
一噸呢.
大概是有十到十五個人的重量.
十到十五個人的重量.
不知道有多少人.
如果你有參加小組的話.
我猜大概每一個小組都有十幾人.
想像一下.
有一個早上.
你整個小組的人聚在一起.
然後有三個婦女.
天未亮就走來你的小組那裡.
大概還沒吃早餐.
天未亮就出發了.
還沒吃早餐的三位婦女.
走到你們十幾個人的小組裡面.
嘗試去推動你這個小組的重量.
大概就是這樣.
不單止這樣.
根據考古的發現.
剛才所看到的那個圓的石頭.
通常是在一個坑那裡.
然後那個坑位.
是慢慢向墳墓門傾斜下去的.
所以大家想像一下.
推那個石頭到墓門去封閉墓.
就比較容易.
因為由上向下推.
但如果你要把那扇門開開.
是由下向上推的時候.
更加困難.
我相信我們知道.
這三位女士都知道.
他們憑三位婦女的力量.

$^{401}$是沒有辦法可以推開那個墓門的.
既然是這樣的時候.
她們都知道推不開.
問題就是既然推不開.
她們還去做什麼呢.
為什麼她們還去.
你會不會去.
你明知道去都進不去.
還去做什麼.
大家知道我上個星期去武漢.
我都跟大家說過.
我去武漢看櫻花.
嘩 很漂亮.
不過在武漢其實有兩個地方看櫻花.
第一個他們叫做東湖的櫻花公園.
第二個就是武漢大學.
現在大家所看到的.
很漂亮.
不過不是我拍的.
這個就是東湖的櫻花公園.
兩個看櫻花的地方.
都需要在網上預約的.
我們很幸運.
去到的時候.
因為我們在澳洲的時候預約不到.
去到的時候就找到一個當地人.
幫我們預約.
去到這個東湖公園那裡.
很開心.
不過就預約不到這個武漢大學.
不要理他.
先去東湖公園.
嘩 看完之後很漂亮.
拍到很多很漂亮的照片.
接著我們就考慮了.
咦 我們預約不到這個武漢大學的.
那去還是不去.
去還是不去.
那我決定說去.
沒有預約到的.

$^{441}$那都去.
去.
為什麼.
因為我們很想看櫻花.
所以就去.
好了.
幾經波折.
去到這個武漢大學.
你以為我們人人都去嗎.
進不去.
為什麼進不去.
其實我們去到那個門口的時候.
有一個大學生.
那些義工來的我相信.
那個年輕的帥哥.
他好心幫我們.
我幫你們馬上預約.
但發覺.
除非那個預約的人有國內的電話.
香港 台灣 或者澳門的電話.
如果不是.
我們拿著澳洲電話那些是預約不到的.
所以沒辦法.
進不去.
是有失望的.
不過我沒有後悔.
我沒有後悔.
為什麼.
我想去看櫻花.
那我當然要努力嘗試.
行不行.
我控制不了的.
我相信大概是同樣的心情.
這三個婦女.
就是這樣的心情.
明知不行.
但我想告耶穌.
照去.
去到如何.
控制不了.

$^{481}$就交給神.
所以他們就這樣去.
因為他們愛耶穌基督.
因為他們想高沒耶穌.
所以縱然他們知道.
進不去的.
照去.
不過我們不禁要問一個問題.
就是說.
他知不知道那個是一個很敏感的時間.
耶穌幾天之前才死在十字架上.
是用一個罪犯的原因.
來死在這個十字架上.
不單止他們知道.
他們推不開那塊石頭.
他們更加不知道.
那時是有羅馬兵兵看著門口.
他們去到的時候.
他們不知道那些羅馬兵兵.
會怎樣對付他們.
他們也不知道這樣做的時候.
他們會面對什麼政治或宗教風險.
很多不知道的.
不可知的事情.
等著他們.
好了,約翰.
使徒約翰在《約翰一書》這樣說.
他說:愛裡沒有懼怕.
愛既完全.
就把懼怕除去.
因為懼怕裡面含有刑罰.
懼怕的人在愛裡未得完全.
雖然這次的經文所說的.
不是所有的恐懼.
這次的經文所說的是.
當我們在耶穌基督或神的愛裡面的時候.
我們不再需要懼怕神的審判和刑罰.
這裡所說的懼怕就是對神的審判懲罰的懼怕.
不過我們不得不承認.
在我們生活經驗上告訴我們.

$^{521}$愛其實是可以幫助我們克服很多的恐懼.
舉個很簡單的例子.
想像一下.
你在一個陌生黑暗的環境裡面.
而你是特別他黑的.
突然間半夜漆黑一片的時候.
隔壁房間的兒女和小孩在尖叫.
你會怎樣?.
黑漆漆的.
等天亮再過去看看發生什麼事.
又或者不要不要.
讓我先歇一歇.
是不是這樣呢?.
我相信如果你真的愛你的兒女的話.
立刻進去房間看看發生什麼事.
為什麼?.
因為父母愛兒女.
所以我同樣相信.
這三個婦女明知道.
他們未必能夠進入那個墳墓裡面.
但是仍然去.
明知道有很多不能夠預知的風險在裡面.
但是這三個婦女仍然去.
為什麼?.
因為他們對耶穌的愛.
超越了他們的恐懼和焦慮.
好了我們再看.
當這幾個婦女去到這個墳墓的時候.
很稀奇很詫異.
為什麼?.
因為那個石頭竟然已經被移開了.
所以讓我們看到.
人不能夠的東西.
神確是可以.
好了我想問大家一個問題.
為什麼這個石頭會被移開?.
很簡單.
牧師你這樣也問?.
因為這幾個婦女不夠力.
所以神幫他們搬開.

$^{561}$是不是這個原因?.
是不是因為我們做任何事情.
當我們不夠力的時候.
神就馬上會幫我們.
神就會派天使下來.
幫這三個婦女推開這塊石頭.
是不是這個原因?.
或者我用另一種方式問你.
如果這三個婦女.
因為都推不開這塊石頭.
我們決定不去.
都進不去.
去做什麼?.
你猜如果這三個婦女沒有去墳墓的時候.
你猜神會不會派天使去推開這塊石頭呢?.
牧師當然會.
因為耶穌要復活在那段時間.
真的?.
耶穌的復活是不是需要把石頭搬開.
然後才能夠復活呢?.
如果是的話.
我們一定所說的復活的大能.
都比不上那兩噸的大石頭.
不是.
一個兩噸的石頭.
不能夠阻止耶穌的復活.
如果是這樣的時候.
為什麼神要派天使去拿開那塊石頭呢?.
讓我告訴大家.
因為正當耶穌被害之後.
祂的所有門徒.
嚇得不得了,害怕得不得了.
個個就四散.
個個就躲在關閉門窗的後面.
只有這三個婦女.
只有這三個婦女有足夠的信心.
足夠的勇氣.
天未亮.
走一段很長的路.
去到這個墳墓裡面.

$^{601}$去想著去告耶穌.
這個石頭.
神將它拿開的原因.
是因為當這三個婦女仍然不知道.
如何拿開這個石頭的時候.
他們就踏出信心的一步.
願意走一段長的路.
去告耶穌.
全因.
他對耶穌的信心.
信得過這個就是神的兒子.
全因他對耶穌的愛.
信得過這個就是他們生命的主.
這個石頭.
移開.
以至這三個婦女.
能夠進到這個空的墳墓裡面.
目擊.
看得到.
耶穌的復活.
也許愛不能夠除去所有的恐懼.
但是告訴大家.
愛能夠讓我們仍然在恐懼裡面.
靠著耶穌復活的大樓.
繼續向前行.
這個就是那個分別.
就好像這三個婦女.
通過他們的行動.
告訴我們.
沒錯這三個婦女.
並不是無畏無懼的.
他們有恐懼的.
不過.
他們對主的愛.
讓他們可以克服.
他們面對的困難.
而繼續前行.
這個告訴我們.
縱然.
沒錯在整個世界裡面.

$^{641}$這個恐懼是客觀地存在的.
不過愛能夠克服很多的恐懼.
弟兄姊妹.
除了.
我們能夠在耶穌基督復活的大樓.
我們得著勇氣向前行.
我們對神的愛.
也是另外一個質素.
讓我們可以幫助我們.
在恐懼當中走出來.
我們再看.
在這三個婦女見到這個石頭被滾開之後.
他們進去.
進去墳墓那裡.
聖經告訴我們.
他們見到一個少年人.
穿著白袍坐在裡面.
然後就指給他們看.
這個就是安放耶穌身體的地方了.
接著.
這個少年人就跟他們說.
你們現在去告訴.
耶穌的門徒.
和彼得.
大家看這個話說不說很奇怪.
耶不彼得不是門徒的其中一個嗎?.
為什麼說完門徒之後.
又要另外再特別刻意的.
提名彼得的呢?.
為什麼呢?.
大家記不記得.
在耶穌被釘之前的時候.
當耶穌受審問的時候.
發生什麼事了?.
三次彼得.
我不認識這個人.
三次彼得否認.
耶穌.
我想大家進入彼得這個人.
當時的心情.

$^{681}$我想大家去揣摩.
感受一下他當時的情緒是怎麼樣的.
一個跟了這麼多年的耶穌.
跟了他很多年.
貼身的跟著他.
到耶穌被抓去審問的時候.
有人認出.
彼得你不是認識這個耶穌嗎?.
我不認識.
三次不認住.
最後耶穌.
被釘在十字架上.
如果是你的心情會是怎麼樣呢?.
我相信我是很羞恥.
很內疚.
覺得自己是一個失敗者.
覺得自己.
自己都看不起自己.
會不會有這樣的感受?.
但是這裡卻是告訴我們.
這裡所展示的是一個神.
恩典一個很漂亮的圖畫.
這個漂亮的圖畫告訴我們.
無論人跌到多深.
無論人做錯的事情有多嚴重.
救恩遠遠大過我們所有犯的過錯.
救恩永遠都能夠完全的.
可以遮蔽我們最深的罪孽.
仍然可以在神面前得著寬恕.
仍然可以在神面前得著救恩.
只要我們帶著信心.
只要我們悔改.
只要我們能夠將我們的生命.
重新降服在耶穌基督的面前.
我相信每一個人都經過萬萬長夜.
在床上點來點去睡不著的日子.
特別我最近發覺年紀越大.
就越容易睡不著.
不過我們要留意一件事.
無論我們如何輾轉反側.

$^{721}$無論那一晚是多麼長的時候.
總會有神希的出現.
總會有明天的出現.
我告訴大家.
耶穌基督的復活.
就是在救恩裡面的黑暗盡頭的一個神光.
這個憤怒已經是一個空的憤怒.
耶穌基督已經復活.
正因耶穌復活.
我們才有希望來活在這個世上.
帶著平安與喜樂.
剛才正如我一開始的時候所說.
這個世界是一個充滿暴力,戰爭,仇恨.
我們所面對的很多都是這樣.
雖然我們說牙牧師無所謂.
我們活在澳大利亞.
我們不是怎樣感覺到的.
我們不是在戰爭的前線.
又不用死人.
不過其實世界是互相影響的.
大家不知道最近有沒有入油.
現在兩個半一公升.
這個世界.
就算我們不是在戰爭的前線.
但我們仍然在一個世界的地球圈裡面.
是互相受到影響.
今天我們的物價飆升.
今天我們看到這個世界很多不公義的事情發生.
如果大家有股票的時候.
我們的股市好像過山車,高高低低.
除非你認識特朗普,是他的親戚朋友.
可以在股市動盪的時候賺他一筆.
如果不是,你和我現在不是賺錢,是虧錢.
我們都是普通的市民.
我們面對這些種種世界發生的事情.
我們感到很無助.
我們感到很無奈.
我們就好像一個大海裡面,大風大浪.
你的船沒有下到撈,隨著風一直被打一樣.
這個世界有千百個原因.

$^{761}$我們是可以感到恐懼和焦慮.
但是好的消息就是.
當我們信耶穌的時候.
我們知道耶穌能夠從死裡復活.
祂就彰顯了為我們所有的人告訴我們.
祂有勝過死亡的能力.
以致我們可以依靠這個耶穌基督.
可以從世界的各種恐懼裡面.
靠著祂的恩典,靠著祂的能力走出來.
因為告訴大家.
這個世界真正的平安,真正的喜樂.
不是來自這個世界.
是來自我們的耶穌基督.
所以弟兄姊妹.
遠遠我們好像這個故事裡面所學到的.
我們有勇氣來選擇我們站在耶穌基督的一方.
也有足夠神對耶穌的愛跟隨祂,服侍祂.
也有承擔神給我們的使命.
去告訴我們周圍的人.
這個我們所信的耶穌已經復活.
祂復活的大樓不是一個抽象的觀念.
而是每一天給過我們力量.
去活出祂的平安和喜樂.
請我們低頭祈禱.
耶穌基督我們多謝你.
這個傳頌了二千多年的復活故事.
每一次我們講起我們.
每一次反省的時候都讓我們重新的知道.
是這個大樓讓這個救恩能夠成為一個可能.
是這個大樓也都讓我們軟弱的人.
能夠在這個不安動盪的世界裡面.
仍然找著平安.
也都因為這個大樓.
我們在各種不同類型的恐懼裡面.
我們可以找到出路.
我們多謝你.
求你恩代,求你大靈.
祈禱,奉主名求.
阿們.
\newpage



\section{路加福音 24:36-53}
\label{sec:031911}
\textbf{復活主帶來翻轉}~~~~譚綺敏傳道
\newline
\newline
連結: \href{http://www.church.com.hk/acms/content.asp?site=cdc\&op=show\&type=product\&code=031911\&layout=sermon}{\texttt{031911}} ~~~~ 語音日期: 2026-04-05
\newline
\newline
\hyperref[sec:031910]{\small{< < < PREV SERMON < < <}}
~
\hyperref[sec:index_chronic]{\small{[返順時目]}}
~
\hyperref[sec:031912]{\small{> > > NEXT SERMON > > >}}
\newline
\newline
路加福音 24:36-53
\newline
\begin{longtable}{cl}
\hline
\hline
章節 & 經文 (和合本修訂版)\\
\hline
24:36 & \begin{tabularx}{0.7\textwidth}{X} 正說這些話的時候，耶穌親自站在他們當中，說：「願你們平安！」 \end{tabularx} \\ \\ \relax
24:37 & \begin{tabularx}{0.7\textwidth}{X} 他們卻驚慌害怕，以為所看見的是魂。 \end{tabularx} \\ \\ \relax
24:38 & \begin{tabularx}{0.7\textwidth}{X} 耶穌對他們說：「你們為甚麼驚恐不安？為甚麼心裡起疑惑呢？ \end{tabularx} \\ \\ \relax
24:39 & \begin{tabularx}{0.7\textwidth}{X} 你們看我的手和我的腳，就知道實在是我了。摸摸我，看，因為魂無骨無肉，你們看，我是有的。」 \end{tabularx} \\ \\ \relax
24:40 & \begin{tabularx}{0.7\textwidth}{X} 說了這話，他就把手和腳給他們看。 \end{tabularx} \\ \\ \relax
24:41 & \begin{tabularx}{0.7\textwidth}{X} 他們還在又驚又喜、不敢相信的時候，耶穌對他們說：「你們這裡有甚麼吃的沒有？」 \end{tabularx} \\ \\ \relax
24:42 & \begin{tabularx}{0.7\textwidth}{X} 他們給了他一片烤魚， \end{tabularx} \\ \\ \relax
24:43 & \begin{tabularx}{0.7\textwidth}{X} 他接過來，在他們面前吃了。 \end{tabularx} \\ \\ \relax
24:44 & \begin{tabularx}{0.7\textwidth}{X} 耶穌對他們說：「這就是我從前和你們同在時所告訴你們的話：摩西的律法、先知的書，和《詩篇》上所記一切指著我的話都必須應驗。」 \end{tabularx} \\ \\ \relax
24:45 & \begin{tabularx}{0.7\textwidth}{X} 於是耶穌開他們的心竅，使他們能明白聖經， \end{tabularx} \\ \\ \relax
24:46 & \begin{tabularx}{0.7\textwidth}{X} 又對他們說：「照經上所寫的，基督必受害，第三天從死人中復活， \end{tabularx} \\ \\ \relax
24:47 & \begin{tabularx}{0.7\textwidth}{X} 並且人們要奉他的名傳悔改、使罪得赦的道，從耶路撒冷起直傳到萬邦。 \end{tabularx} \\ \\ \relax
24:48 & \begin{tabularx}{0.7\textwidth}{X} 你們就是這些事的見證。 \end{tabularx} \\ \\ \relax
24:49 & \begin{tabularx}{0.7\textwidth}{X} 我要將我父所應許的降在你們身上，你們要在城裡等候，直到你們領受從上面來的能力。」 \end{tabularx} \\ \\ \relax
24:50 & \begin{tabularx}{0.7\textwidth}{X} 耶穌領他們出來，直到伯大尼附近，就舉手給他們祝福。 \end{tabularx} \\ \\ \relax
24:51 & \begin{tabularx}{0.7\textwidth}{X} 正祝福的時候，他離開他們，被帶到天上去了。 \end{tabularx} \\ \\ \relax
24:52 & \begin{tabularx}{0.7\textwidth}{X} 他們就拜他，帶著極大的喜樂回耶路撒冷去， \end{tabularx} \\ \\ \relax
24:53 & \begin{tabularx}{0.7\textwidth}{X} 常在聖殿裡稱頌神。 \end{tabularx} \\ \\
[1ex]
\hline
\hline
\end{longtable}
$^{1}$讓我們一齊先禱告.
親愛的阿爸父.
今日是耶穌基督復活的日子.
讓我們再次反思你在我們生命奇妙的作為.
主啊.
希望我們今天都一齊去轉向你.
願你來塗灶我們的生命.
因你活著.
我們充滿喜樂.
因你活著.
我們生命有盼望.
讓我們經歷你更多的恩典.
起來傳揚你的名.
我們將以下時間交在你手裡.
求你恭活你僕人的口.
使你的話語可以進到我們每一個人的心裡.
以致我們可以活出你.
多謝你聽我們祈禱.
奉主耶穌基督得勝一命祈求.
阿們.
有一個教會的會友.
因為生活太忙碌.
很久沒有回來教會崇拜了.
有一天他突然接到牧師公的信.
裡面說道教會要關門了.
星期日是最後一次崇拜.
希望你回來見一面.
他見到是最後一次崇拜.
他也很珍惜.
於是就去參加崇拜了.
回到教會他也看到以前的弟兄姊妹.
人數也不少.
坐得滿滿的.
最後一次就像識別會一樣.
大家聊聊天.
互相問問安.
但很奇怪.
他見到教會教堂前面有個棺木.
他說今天是哪位教友的安息禮拜呢?.
哪位教友回到天家呢?.

$^{41}$大家都不知道.
崇拜開始.
牧師就歡迎弟兄姊妹回來.
說很久沒有見到這麼多弟兄姊妹.
回到神的家去敬拜神.
平時崇拜就是十個八個坐在那裡.
今天是最後一次崇拜.
下星期教會就關門了.
所以你們不用擔心了.
我不會再寫信給你們了.
今天見了面之後.
我也不知道什麼時候再見.
我也為大家祝福.
接著就說道.
他說最後有個項目.
就是希望大家可以去瞻仰為榮.
大家看到前面有個棺木.
大家來瞻仰為榮.
大家都很奇怪.
誰回到天家呢?.
大家都不知道.
很好奇.
排隊去瞻仰為榮.
大家去看完之後.
都流眼淚了.
你們猜猜棺木裡放了些什麼呢?.
空墳墓我們聽得多.
沒有空墳墓的東西.
我們是否真的看到空墳墓.
會流淚,會開心呢?.
我有想過.
會不會是耶穌的畫像放在裡面.
或者是一個像放在裡面.
復活之後動高的像.
會不會是這樣呢?.
想想這樣.
大家是否會難過,會開心呢?.
原來牧師在棺材裡面.
放了一個鏡子.
當大家看下去的時候.

$^{81}$看到誰呢?.
看到自己.
自己躺在棺材裡.
弟兄姊妹你們有沒有流淚呢?.
人生是很短暫的.
所以當時他們看到自己的時候很難過.
他們明白耶穌說.
人若賺得全世界.
他賠償自己的生命有什麼益處呢?.
人還能拿什麼來換自己的生命呢?.
於是大家都完了崇拜之後.
大家都說.
我們決定要回來教會了.
請牧師不要再關門了.
下個星期繼續回來教會.
盼望我們教會一直不會關門.
其實我們本來死在過犯罪惡裡.
主耶穌的救贖.
以致我們可以活過來.
我們經常在爾弗所書都讀到這段經文.
我們相信主之後.
其實主耶穌也沒有離開我們.
祂很想我們經常回到祂的懷抱裡.
特別是今天是復活節的時候.
我們很盼望再次思想.
主耶穌復活對我們的意義.
剛才讀到的經文.
在路加福音24章36-53節.
我們看到門徒本身是很驚慌害怕的.
後來他們還讚神.
所以復活的主會帶來一個生命的翻轉.
你們相信嗎?.
希望我們都經歷過.
希望我們在座的每一個人都可以享受這個福分.
根據《新約聖經》說到.
其實明確記載了主耶穌在死裡復活之後.
向門徒和一些人顯現有十到十一次.
我們一起來看看.
第一次是向莫大拉瑪利亞顯現.
那時候莫大拉瑪利亞走到耶穌的墳墓裡.

$^{121}$很傷心的樣子.
耶穌親自顯現給她看.
安慰她 更新她.
路上一直有其他婦女來的時候.
耶穌亦向她們顯現 請她們報信.
由彼得顯現過.
在二馬五司路上有兩個門徒.
在路加方面有記載.
他們不明白發生了什麼事情.
耶穌在路上解明聖經給他們聽.
接著來到我們剛才讀的經文.
第五次的顯現.
十個使徒 當時多瑪不在當中.
在耶路撒冷的地方.
接著耶穌再向十一個門徒顯現.
因為多瑪不在.
所以特意再次向他們顯現.
接著在加利利的海邊.
七個使徒回去打魚的時候.
耶穌特別向他們顯現.
特別問彼得「你愛我比這些更深嗎?」.
接著向十一個使徒再次顯現.
有大使命的頒佈 叫他們傳福音.
接著在加利利一帶向五百多人顯現.
不同的人.
亦向耶穌的弟弟雅各.
最後在撒冷山上升遷.
很多人去看著.
所以耶穌復活之後.
有四十千支狗.
多次多方向不同的人顯現.
為的是要去兼顧他們的信心.
給他們安慰.
也是託付主耶穌的大使命.
要他們去作主耶穌的見證人.
所以今天我們要講的.
就是耶穌第五次的顯現.
你會發覺有一半左右.
都是單獨向一些使徒.
即是耶穌揀選的十二個門徒去顯現.

$^{161}$我們知道猶大出賣主 所以只剩下十一個.
所以今天我們要講的.
就是耶穌在祂復活的第一天晚上.
在祂死後第三天復活.
第一天晚上.
向一群聚在一起的十個使徒顯現.
耶穌向他們顯現的時候.
有什麼情況呢?.
今天我們要詳細去看看.
耶穌受難.
被釘死.
那時候門徒的心情.
他們的光景是怎樣呢?.
在《讓封音》二十章十九節.
提到那天就是七天的第一天晚上.
門徒所住的地方.
因怕猶太人門都關了.
這裡特別提到.
他們很害怕.
很害怕猶太人.
很害怕其他人.
很害怕被人拘捕.
很害怕發生什麼事.
所以把門關得緊閉.
人都不可以進來.
倒在房間裡.
心裡充滿恐懼.
也不知道未來的方向是怎樣.
因為突然發生的事情.
大家都很恐懼.
剛才我們讀過那段經文.
在《路加福音》二十四章.
提到.
就在這時.
耶穌站在他們當中.
說:願你們平安.
耶穌其實很愛這班門徒.
他一復活的時候.
特別要向門徒顯現.
因為很明白.

$^{201}$如果有懼怕.
有什麼情緒的時候.
如果不處理好.
其實什麼事都做不到.
在情緒裡.
很難去做事情.
就算很容易的事情.
都會變得很難做.
因為人有情緒.
所以有情緒的時候.
跟小朋友說話.
他們都聽不進去.
要處理好人的情緒.
所以復活的第一天晚上.
當門徒聚集在一起.
很害怕的時候.
耶穌立刻向他們顯現.
就是說:願你們平安.
祂很想他們得到.
從神而來的平安.
祂很想處理好他們心裡恐懼的問題.
但是門徒對於耶穌的復活.
不是很明白.
也沒有心理準備.
也可能太理性.
門都關了.
耶穌怎麼進來呢?.
如果真是耶穌的時候.
怎麼進來呢?.
所以突然間害怕起來.
難道我們見鬼了?.
所以耶穌特別再跟他們說.
你們為什麼去受煩呢?.
為什麼起疑念呢?.
你看看我的手腳.
就知道真是我了.
你摸摸我的雲.
是沒有骨沒有肉的.
我是有的.
有骨有肉.

$^{241}$你來看看我的手腳.
然後就把自己的手腳.
給門徒看.
他幫助門徒.
真是相信復活.
所以門徒.
在四十一至四十三節.
門徒開心到不敢相信.
覺得很稀奇.
為什麼叫開心到不敢相信呢?.
又開心又不敢相信.
在四十一至四十三節.
開心到不敢相信.
這個翻譯翻譯得很好.
原文也是這個意思.
很開心很開心.
耶穌真的復活.
但又不敢相信.
是不是真的呢?.
為什麼會這樣呢?.
前幾天他們才在十字架.
耶穌釘死的下面.
驚慌失措.
又不知道怎麼走.
他們親眼看著.
跟隨了三年的耶穌.
以為可以拯救.
以色列人的老師.
是他們的夫子.
突然被害.
還在酷刑下處死.
被釘在十字架上.
這麼辛苦的處死方法.
對於他們來說.
一切所發生的事.
都來不及消化.
所以突然.
耶穌又活生生地.
站在他們面前.
情緒好像一下子.

$^{281}$在低谷裡.
又害怕又辛苦.
又不知道怎麼做.
突然又衝上去.
「耶穌你還在?你真的在?」.
又很開心.
所以他們都不太確定.
這是否真實的事情.
這是一個心理保護.
自己保護自己的心理反應.
不想再受傷害.
所以耶穌.
再繼續去問他們.
「你們有沒有東西吃?」.
門徒給了他一片燒魚.
他接過來.
吃了它.
這個動作是告訴門徒.
「我可以吃東西」.
「鬼魂是不能吃東西的」.
「我是身體的復活」.
「你是否明白身體的復活是可以吃東西的?」.
其實這是指.
耶穌已經戰勝了死亡.
戰勝了罪惡.
身體復活.
也是耶穌.
常常預言給門徒.
身體復活的.
一個真正的真實情況.
其實耶穌也應許.
我們信徒也有身體的復活.
我們將來審判的時候.
身體是可以復活的.
所以我們看到第一點.
在懼怕裡面.
主親自站在我們當中.
願我們平安.
神很想我們平安.
即使耶穌的顯現.

$^{321}$門徒也摸過耶穌.
但他們是否這樣.
定下來.
信心得堅定呢?.
不是吧?.
因為他們可能還有很多疑問.
所以我們的信仰.
如果要堅固.
除了經歷神之外.
我們還要去明白真理.
我們要知道.
什麼是會回事呢?.
我看到很多弟兄姊妹.
初初信主的時候.
聽我們分享.
禱告很容易去蒙英穩.
很特別去經歷一些恩典.
也很容易感受到.
耶穌真的同在,我祈禱.
神就答應我了.
很多初信的弟兄姊妹.
也有這些經歷.
我自己也有.
想想自己初信的時候.
可能也有.
因為神真的讓我們初信的時候.
特別經歷.
這是聖靈的工作.
但可惜當困難一來到.
工作一忙,世界的吸引力太大.
又想去玩,又沒有時間.
也沒有在聖經裡.
真正去紮根.
知道聖經裡說什麼.
所以很容易離開神.
所以其實明白真理.
明白聖經是很重要的.
所以耶穌繼續.
跟他們說,在44-45節.
耶穌對他們說.

$^{361}$摩西的律法.
先知的書,和詩篇上所記.
凡指著我的話.
也必須應驗.
於是耶穌開他們的心竅.
使他們能明白聖經.
一直說到.
一切在舊約裡.
因為那時候還沒有新約.
一切摩西的律法,先知的書.
其實很多都指向耶穌.
會來的,這個世界.
也會復活的.
舊約一直在說這些的時候.
其實門徒也很熟悉.
所以耶穌知道.
其實門徒還沒有明白.
所以他就說.
耶穌會開他們的心竅.
使他們能明白聖經.
這裡說到.
人是可以明白神的話語的.
神的話語其實不是那麼難明白.
我們每一個人.
只要我們有.
知道耶穌.
跟我們在一起的時候.
我們是可以去明白的.
不是靠自己的聰明.
不是要怎樣研究.
例如說聖書很難明白.
不需要,也可以去明白.
因為你只要親自去讀聖經.
你是會明白的.
所以當我們有追求的心.
我們不是當聖經.
越讀越沒有味道.
很枯乾.
或者是.
當一些知識.

$^{401}$很深奧.
講到好像很深奧.
有些人會說.
好像很深奧.
要了解文化背景.
要了解才明白.
所以我不明白.
有些人可以開始當.
一個知識的累積.
去讀聖經.
這個也很難去明白聖經.
有時候有些人自以為很明白.
好像文是化理賽人.
以為很明白神的教訓.
就用聖經去教訓人.
讀聖經也不可以去明白.
但當耶穌.
去幫我們明白的時候.
我們其實知道.
聖經其實不難明白.
只不過很多時候.
我們有自己的框框.
有自己的想法.
當一些很難明白的東西.
我們也很難去接受.
要明白神的話語.
也真的需要我們去謙卑.
去到神面前去祈求神.
去開啟我們的心竅.
以致我們更加明白.
祂的受苦和受死的意義.
更加知道祂復活的大能.
去明白神的話語.
所以繼續在46節.
耶穌也和我們說.
照經上所寫的.
耶穌必受害.
第三天從死裡復活.
然後耶穌再去解釋給門徒聽.
基督的受難,受死,復活.

$^{441}$不是偶然發生的.
不是突然間發生.
好像你們現在突然間看到.
耶穌預定了一個救贖計劃.
神很愛我們.
很愛祂做的人.
祂差派了祂的獨生兒子.
耶穌來到這個充滿罪惡的世界.
其實目的就是.
要將自己的百姓.
從罪惡裡救出來.
唯有神的救贖.
我們才有真正平安的出路.
所以耶穌基督.
為什麼在那個時間.
在羅馬政權之下.
祂要來到這個世界呢.
難道祂要自己被釘死嗎?.
這個那麼辛苦的酷刑.
為什麼祂要選擇這些時間.
要來到彰顯給我們看呢?.
是因為祂愛我們.
門徒跟隨了耶穌三年.
也看過耶穌走過不少神蹟.
醫病,趕鬼.
盲的可以看見.
甚至連拉撒路復活.
祂也可以看見.
只是在祂釘死之前.
沒有多久發生.
耶穌也在祂.
跟隨的.
和門徒的生活中.
祂多次說到.
自己會受苦.
被殺.
第三天復活.
跟門徒也預言過.
但是當這些.
預言的真實.

$^{481}$發生在耶穌身上的時候.
門徒其實仍然是.
不理解什麼叫死人復活呢.
所以他們仍然是.
很害怕.
關大門.
不敢去.
怕那些人進來.
不敢去面對這件事.
所以其實.
看見神蹟.
或者在身上遇過一些特別的經歷.
不等於你是真的明白.
真理的.
真的要耶穌去開我們的心.
以致我們更加去明白.
耶穌在我們身上的旨意.
所以我們看見在困惑裡面.
其實主可以開我們的心竅.
是要我們去.
明白真理.
可能你也會問.
在日常生活裡面.
主耶穌怎樣開我們的心竅呢.
我想可以在很多方面.
其實都可以的.
首先就是你有沒有去讀經祈禱呢.
有沒有去.
自己去親近神呢.
也有沒有跟其他人.
去相交呢.
跟弟妹有團契啊.
有小組啊.
一些營會.
我自己都是在營會裡面.
去感動以致我去奉獻給神.
或者是透過一些事奉.
可能是.
你很忙的時候去事奉神.
但為什麼這麼有力量呢.

$^{521}$神去感動你.
或者是透過一些短宣裡面.
看見神的作為.
的表現.
這些都可以有神的感動.
有神開你的心竅.
不知道你有沒有.
試過一些這樣的經驗.
之前你靈修讀經.
對一些經文沒有什麼感覺.
讀過就算了.
但某些情況之下.
你再讀的時候.
突然間好像明白了.
突然間好像神跟你講說話.
那一刻就很受感動.
我在靈修讀經的時候.
經常都有神的話語去感動我.
有時候在預備講章的時候.
我真的覺得.
要怎樣去講.
讓姐妹們明白.
神又藉著一些話語.
再次去感動我.
所以困難中.
我也想起一些經文.
我也是被安慰的.
所以我很喜歡去親近神.
其實我們真的要去親近神.
我們才可以去明白.
神的話語.
神藉著他的話語.
去感動我們.
開我們的竅.
所以當門徒的心被打開的時候.
真正明白.
耶穌的復活.
救贖的計劃之後.
耶穌就繼續跟他們說.
去託付耶穌自己的使命.

$^{561}$在47-48節.
引要奉他的名.
全悔改赦罪之道.
從耶路撒冷.
直傳到萬邦.
你們就是這些的見證.
你們就是這些事的見證.
今天我們都是.
主耶穌的見證人.
所以當.
耶穌打開了門徒的心怯的時候.
不單止要我們.
單單去明白真理.
更加要我們去改變自己.
要活出.
主耶穌要我們活出的樣式.
就是做主的見證人.
其實要改變並不容易.
試試在家庭裡.
你以為你可以改變你的配偶嗎?.
你結婚之後.
大家都很相愛.
我們都愛對方才結婚.
但結婚之後.
你才發覺.
你要改變你的另一半.
是很困難的.
你不要這樣想.
這樣想你會很失望.
你會很不開心.
為何他不改變呢?.
為何你自己不改變呢?.
要改變自己.
改變自己才可以.
活出神的樣式.
所以唯有.
遇見主耶穌.
祂的愛.
祂怎樣去愛我們呢?.
我們才會心裡感動.

$^{601}$心裡有火熱.
其實這是一個很喜樂的事情.
如果你試過讀經文.
讀到很感動.
或者你在營會裡.
除了氣氛令你很開心之外.
你被神的話語觸摸.
其實你會很開心.
可能那一刻.
你會流淚.
好像剛才說.
在棺材裡.
看到自己.
會流淚.
但另一方面.
你會悔改.
流淚之後.
你會很有力量.
覺得教會不可以關閉.
你有一個責心.
會覺得我們要去活出真理.
我們知道.
我們要去.
愛我們的鄰舍.
彼此弟兄姊妹要有愛.
我們才可以彰顯神的愛.
你就很自然.
會喜歡回來教會.
喜歡去煮飯.
去服侍弟兄姊妹.
喜歡互相關心.
覺得.
神的家實在太美好了.
因為神太愛我們.
我記得我在大學的時候.
當時我信主兩年.
我亦都經歷了.
神打開我的心.
很奇妙的作為.
我當時找到.

$^{641}$大學的第一份工作.
很開心.
這是一個大廈的會計師樓.
做一個核數的工作.
但上班的時間.
剛好我答應了.
要去門徒訓練營的營會.
我當時在大學.
信主.
我亦都.
在大學參加.
大學的門徒訓練營.
亦都負責大組.
亦都有人帶著我.
我有我的組長.
我亦都是組長.
帶著其他人.
我信主兩年多.
剛好上班時間.
就是營會的時間.
上班的頭三個星期.
就是訓練.
一班大學生.
一起去接受.
訓練課程.
所以很難去請假.
我和我組長說.
這次營會我不參加.
因為我剛好要上班.
我組長.
很嚴厲的跟我說.
不行.
怎可以.
這份工作不是神給你的嗎.
你又說是神給你的.
你要做組長.
你不可以不去.
因為你的組員誰理你.
很嚴厲的跟我說.
你帶著也可以.

$^{681}$我不熟悉你的組員.
你要去帶.
因為神既然.
給你在這個位置.
你既然做組長.
你就要去負責你的羊群.
我當時心裡很難過.
我覺得我做組長很嚴厲.
你怎做組長.
我心裡真的這樣想.
組長你跟我說這些.
我第一份工作.
為何你都不了解我.
我心裡是這樣想.
然後我繼續靈修禱告.
我跟神說.
神啊.
我可不可以不去.
雖然羊是我帶順耶穌的.
也是我一直跟著他們.
但我不去.
有其他人都可以幫我帶.
但我在.
在我靈修裡面.
我看到神責備我.
神也鼓勵我.
我組長也跟我說.
你試試另一份工作.
如果神不讓你去.
我也覺得.
你盡了你的責任.
你也不會不開心.
因為公司不讓你請假.
我覺得也是.
如果我連.
另一份工作也不去做.
到時我可能會.
一直去訓練.
一直不開心.
我就覺得.

$^{721}$如果神給我另一份工作.
那就去吧.
其實我之前很不願意.
因為我發覺.
公司知道我訓練也請假.
我另一份工作要寫理由.
因為訓練不能請假.
我還寫請假.
不是我有甚麼事.
因為我也要有.
真正的理由.
不是我病了.
因為我要參加大學的門徒訓練營.
做組長事奉.
我就這樣寫.
然後很奇妙.
一個星期之後.
把表格遞給我.
我的請假是獲批.
我其實心裡很感恩.
神給我的體驗就是.
你走回自己要做的事.
你去經歷神的恩典.
這個已經是.
三十多年前.
一個很深刻的恩典.
我還記得.
三十多年前.
因著這個恩典.
我知道.
我走這條路.
是神帶領我的.
無論我遇到甚麼情況.
神會負責任.
讓我們有出人意外的平安.
所以我願意.
將我的生命交給神.
讓我們有這樣的經歷.
所以當你不去試過.
去走一些.

$^{761}$神要你走可能難成的事.
的時候.
你會發覺你很難去經歷神.
因為你只是在享受神的恩典.
就好像一個.
一個.
嬰兒.
只會被人餵奶.
不懂得去服侍神.
去服侍人.
所以耶穌.
開了門徒的心協.
之後他知道.
門徒親眼體驗.
耶穌真正復活的時候.
這個時候.
是要去頒佈大使命的時候.
所以他也告訴門徒.
你要做我的見證人.
但他也知道.
在這個時候.
要門徒去做這些事.
是不容易的.
所以在49節.
他說:我要將我父所應許的降臨.
降在你們身上.
你們要在城裡等候.
直到你們領受.
從上頭來的能力.
不是要你馬上去做.
而是要你去.
回到耶路撒冷城.
在城裡等候.
在耶路撒冷城等候.
直到你等候到.
從上頭來的能力.
這個能力是什麼呢?.
就是聖靈.
所以耶穌去應許.
聖靈降臨到他們身上.

$^{801}$在約翰福音裡.
當耶穌去最後晚餐之後.
耶穌用了很長的篇幅.
去跟門徒說道.
聖靈會降臨在他們身上.
所以當耶穌知道.
人是很容易軟弱的.
大家也在擔心.
會不會我再去做神的見證人.
的時候會被抓呢?.
不單只是被.
羅馬兵.
去抓.
更加可能被猶太人去迫害.
連自己人也會被迫害.
所以人其實是很軟弱的.
耶穌是知道的.
做他的見證人是不容易的.
今天我們也知道.
很多宣教士在他一入宮場.
沒多久也會身亡.
我們也知道.
在過去也發生過不少.
但是.
從上頭來的能力.
就是聖靈的能力.
幫助我們完成神所託付的大使命.
不是靠自己的勇氣.
不是靠自己的才幹.
而是靠著.
聖靈的大能.
因為神去負責一切.
聖靈催迫你做的時候.
只要我們去忠心.
去信服去行的時候.
其實神是負責一切的.
所以他叫門徒先去等候.
等候從神而來的能力.
你就會發現.
你心裡火熱去做.

$^{841}$所以第三樣我們看到.
在躲藏裡面.
神似下聖靈.
叫我們有能力去做他的見證.
今天我們也是.
復活主的見證人.
所以我們有能力去做.
神的見證.
如果我們教會失去了.
傳福音的負擔.
其實會變成一個封閉的社團.
我們很容易會凋零.
很容易會關門.
所以我們看到.
很多教會.
特別是西人教會.
一代一代去關門.
一個一個的關門.
有些教會.
賣了去.
回教徒.
賣了.
所以其實是很難過的.
所以我們.
不單止要自己去享受神的愛.
也更加要出去.
去傳揚神的道.
所以最後來到經文最後的.
50到53節.
接著耶穌.
向他們祝福.
然後升天.
你會發覺在路加福音裡.
這一段是最後一段的經文.
是以耶穌去祝福和升天.
去結束的.
但很特別地也記載了門徒的反應.
他們的反應.
不像當初耶穌被釘死.
那麼絕望的反應.

$^{881}$而是.
耶穌再次離開的時候.
他們沒有難過.
反而是去.
敬拜神.
去拜主耶穌.
表示他們信服.
也大大地歡喜.
不單止歡喜.
大大地歡喜.
很開心很開心耶穌去升天.
然後回到耶路撒冷.
去上載.
殿裡送讚神.
他們不是再躲在.
屋裡.
出去在殿裡.
開著門大大聲聲.
去敬拜神.
鼓勵人們去信耶穌.
因為他們在憂愁裡.
經歷神的恩典.
所以我們看見.
在憂愁裡主帶我們進入喜樂.
上上去.
送讚神.
這是一個很大的轉變.
從驚慌害怕.
到大大的歡喜.
從關起門.
去躲藏.
直到上在殿裡.
送讚神.
是什麼讓這群門徒.
生命可以這麼大的翻轉呢?.
就是遇見了.
復活的主耶穌.
主耶穌去升天的離開.
不是真的離開.
是另一種的形式.

$^{921}$與門徒同在.
聖經記載.
祂經常在父神的右邊.
為我們去代求.
使聖靈與我們同在.
祂舉手的祝福.
成為了門徒信心的印記.
所以.
因為復活升天的主耶穌.
有一天.
祂會再來到我們中間.
去拯救審判這個世界.
看見主耶穌的復活.
門徒的生命被翻轉.
也不再去驚慌害怕.
因為祂帶著.
主復活的確據.
也有聖經的真理.
也帶著.
一個福音的使命.
也知道.
祂有從上頭來.
聖靈的能力.
所以你會發覺.
接著在仕途行轉.
會發覺他們很勇敢去傳福音.
也成為主耶穌.
為主殉道的.
一群見證人.
所以今天.
我們也要.
復活的主翻轉我們的生命.
我想到我們有三方面.
可以應用.
首先我們知道.
主耶穌賜下平安給我們.
所以我們要將一切的.
驚慌憂慮帶到神的面前.
其實我們常常.
都被各樣的驚慌.

$^{961}$嚇到.
突然間.
為什麼會打仗呢?.
打仗以為與我們無關.
是不是真的與我們無關呢?.
你會發覺不是.
整個世界都動盪起來.
又價又高.
物價又高.
又怕糧食運送不到.
偏遠的地方.
人民更加.
焦慮.
所以其實.
突然來的事情都令我們害怕.
有憂慮.
未來也有憂慮.
突然失業.
突然畢業找不到工作.
怎麼辦呢?.
一年也找不到.
怎麼辦呢?.
真的很難過.
亦都可能你檢查身體.
突然說你.
有糖尿病.
也好糖尿可以控制.
癌症.
真是晴天霹靂.
對死亡.
對疾病.
對身體越來越衰殘.
我們都有恐懼.
或者是對人際關係.
夫妻之間.
突然吵架.
小孩很難教.
家庭很緊張.
很多張力.
這些都是.

$^{1001}$我們日常遇到的困難.
問題.
說錯一句說話.
那個人突然在你面前.
刻念.
自己整晚都睡不著.
很多這些小事情.
都會令我們很焦慮.
求主幫助我們.
將這些焦慮.
將這些不開心.
憂慮的事情.
當驚慌的事情.
帶到主耶穌的面前.
藉著祈禱.
藉著讀聖經.
藉著和主耶穌的親密關係.
我們將一切交託給神.
主耶穌很想.
賜平安給我們.
很想解決我們心裡的.
罪疚感或者一些.
懼怕的人生的問題.
很想給我們喜樂.
很想和我們同行.
很想加力讓我們.
向前面繼續走.
但我們願不願意將我們的人生.
交給祂呢?.
願主自己親自去看顧.
主耶穌也很想.
開我們的心竅.
使我們明白聖經.
所以我們要立志.
每天親近神.
回到神的話語裡.
我們需要去明白.
神在我們身上的旨意.
所以我們要.
自己和神獨處.

$^{1041}$意思是自己有自己靈修.
禱告,讀經的時間.
自己去知道神.
怎樣向我們說話.
但同樣的.
我們也要一起去.
明白神的話語.
一起去彼此.
分擔我們的難處.
剛才門徒也有十個人.
十個聚集在一起.
十一個聚集在一起.
還有其他人.
所以其實是要在.
教會裡去眾讚神.
一個人的力量是孤單的.
所以聚集.
也是需要的.
為什麼我們要有聚會呢?.
有小早茶經.
有崇拜,營會.
有短孫.
有不同形式的聚會.
以致我們在不同形式裡親近神.
不要回來崇拜.
形式化地覺得.
我基督徒就回來了.
你要帶著一個謙卑的心.
回來真正去.
遇見耶穌.
真正去遇見復活的主.
所以求主去開我們的心竅.
每天都可以.
在生活裡遇見主.
特別是.
在你親近神的時間.
在你回來崇拜的時間.
盼望今天我們都遇見主耶穌基督.
和你說話.
以致你心裡可以有能力.

$^{1081}$最後讓我們用信心.
踏出一步.
主耶穌要賜下你能力.
我們起來回應.
神的愛和大使命.
因為聖靈賜下能力.
也讓我們心心可以經歷.
其實我們做見證可以很簡單.
以前我們經常說.
街頭報道.
很害怕和陌生人聊天.
其實你想想.
很多時間可以去見證神.
例如吃飯.
不會所有朋友都信了耶穌.
吃飯的時候.
和一些未信耶穌的朋友.
吃飯的時候.
聊幾句.
將這些訊息帶給他們.
或者在很多群組裡.
很多群組.
小學同學.
中學同學.
我最近在中學同學的群組.
發一些小故事.
一些復活的訊息.
有一個佛教徒.
他馬上回覆.
我現在在抄佛經.
在抄佛經.
他說你想和我傳教嗎.
我說不要緊.
你在抄佛經.
我這些都是希望大家開心.
希望大家都有喜樂.
所以.
給你們一些小訊息.
你們不要覺得很冒犯.
突然間.

$^{1121}$他說在抄佛經.
你會覺得好像.
好像不知道怎樣.
你可以很歡樂.
很喜樂地告訴他們.
我希望神賜的平安你有份.
所以其實可以在一些很短的訊息裡.
都可以問候一下.
或者你見到教會裡面.
最近誰沒有回來.
小組裡面誰沒有來.
你就關心一下.
很少的分享很少的關心.
其實已經可以很感動弟兄姊妹.
也可以很體驗.
神在你身上的.
作為.
所以盼望我們都.
被神的愛去感動.
知道復活的主可以更新.
改變我們的生命.
我們也要相信神有這個能力.
去改變我們的生命.
我們一起祈禱.
親愛的主耶穌.
從死裡復活.
勝過死亡和罪惡的權勢.
求主你開我們的心穴.
讓我們真正去明白聖經.
認識你是.
那位獨一的真神.
是唯一的救主.
你來到這個世界上.
目的是要救我們.
脫離罪的綑綁.
回去主耶穌基督.
你的懷抱裡.
主求你.
去翻轉我們的生命.
帶領我們在黑暗裡.

$^{1161}$進入光明的道路.
求主你也貼下聖靈的能力.
讓我們的生活不再是.
憂慮和恐懼裡.
讓我們可以勇敢地.
踏出一步.
作主你的見證人.
將福音傳遍.
傳開.
讓我們的親朋好友.
都可以去認識你.
相信你.
得著平安.
主呀!這是我們心底裡的渴求.
我們真的很想我們的朋友.
很想我們的愛的家人.
去認識主.
求主你為我們去承受.
你奇妙的作為.
多謝主聽我們祈禱.
奉主耶穌基督得勝的名祈求.
\newpage



\section{以賽亞書 52:13-53:12}
\label{sec:031912}
\textbf{擔當我們憂患的主}~~~~高銘謙牧師博士
\newline
\newline
連結: \href{http://www.church.com.hk/acms/content.asp?site=cdc\&op=show\&type=product\&code=031912\&layout=sermon}{\texttt{031912}} ~~~~ 語音日期: 2026-04-05
\newline
\newline
\hyperref[sec:031911]{\small{< < < PREV SERMON < < <}}
~
\hyperref[sec:index_chronic]{\small{[返順時目]}}
~
\hyperref[sec:031913]{\small{> > > NEXT SERMON > > >}}
\newline
\newline
以賽亞書 52:13-53:12
\newline
\begin{longtable}{cl}
\hline
\hline
章節 & 經文 (和合本修訂版)\\
\hline
52:13 & \begin{tabularx}{0.7\textwidth}{X} 看哪，我的僕人行事必有智慧，他必被高升，高舉，升到至高之處。 \end{tabularx} \\ \\ \relax
52:14 & \begin{tabularx}{0.7\textwidth}{X} 許多人因他驚奇---他的面貌比別人憔悴，他的外表比世人枯槁--- \end{tabularx} \\ \\ \relax
52:15 & \begin{tabularx}{0.7\textwidth}{X} 同樣，他也必使許多國家驚奇，君王要向他閉口。未曾傳給他們的，他們必看見；未曾聽見過的事，他們要明白。 \end{tabularx} \\ \\ \relax
53:1 & \begin{tabularx}{0.7\textwidth}{X} 我們所傳的有誰信呢？耶和華的膀臂向誰顯露呢？ \end{tabularx} \\ \\ \relax
53:2 & \begin{tabularx}{0.7\textwidth}{X} 他在耶和華面前生長如嫩芽，像根出於乾地。他無佳形美容使我們注視他，也無美貌使我們仰慕他。 \end{tabularx} \\ \\ \relax
53:3 & \begin{tabularx}{0.7\textwidth}{X} 他被藐視，被人厭棄；多受痛苦，常經憂患。他被藐視，好像被人掩面不看的一樣，我們也不尊重他。 \end{tabularx} \\ \\ \relax
53:4 & \begin{tabularx}{0.7\textwidth}{X} 他誠然擔當我們的憂患，背負我們的痛苦；我們卻以為他受責罰，是被神擊打苦待。 \end{tabularx} \\ \\ \relax
53:5 & \begin{tabularx}{0.7\textwidth}{X} 他為我們的過犯受害，為我們的罪孽被壓傷。因他受的懲罰，我們得平安；因他受的鞭傷，我們得醫治。 \end{tabularx} \\ \\ \relax
53:6 & \begin{tabularx}{0.7\textwidth}{X} 我們都如羊走迷，各人偏行己路；耶和華使我們眾人的罪孽都歸在他身上。 \end{tabularx} \\ \\ \relax
53:7 & \begin{tabularx}{0.7\textwidth}{X} 他被欺壓受苦，卻不開口；他像羔羊被牽去宰殺，又像羊在剪毛的人手下無聲，他也是這樣不開口。 \end{tabularx} \\ \\ \relax
53:8 & \begin{tabularx}{0.7\textwidth}{X} 因受欺壓和審判，他被奪去，誰能想到他的世代呢？因為他從活人之地被剪除，為我百姓的罪過他被帶到死裡。 \end{tabularx} \\ \\ \relax
53:9 & \begin{tabularx}{0.7\textwidth}{X} 他雖然未行殘暴，口中也沒有詭詐，人還使他與惡人同穴，與財主同墓。 \end{tabularx} \\ \\ \relax
53:10 & \begin{tabularx}{0.7\textwidth}{X} 耶和華的旨意要壓傷他，使他受苦。當他的生命作為贖罪祭時，他必看見後裔，他的年日必然長久。耶和華所喜悅的事，必在他手中亨通。 \end{tabularx} \\ \\ \relax
53:11 & \begin{tabularx}{0.7\textwidth}{X} 因自己的勞苦，他必看見光就心滿意足。因自己的認識，我的義僕使許多人得稱為義，他要擔當他們的罪孽。 \end{tabularx} \\ \\ \relax
53:12 & \begin{tabularx}{0.7\textwidth}{X} 因此，我要使他與位大的同份，與強盛的均分擄物。因為他傾倒自己的生命，以致於死，也列在罪犯之中。他卻擔當多人的罪，為他們的過犯代求。 \end{tabularx} \\ \\
[1ex]
\hline
\hline
\end{longtable}
$^{1}$我們同心祈禱.
慈悲人愛的恩主.
因為你死裡今日復活.
我們看到你怎樣戰勝死亡.
給我們有重生的生命.
讓我們不再在罪惡當中打滾.
步向滅亡的道路.
反而是得著永生.
求主你藉著第四首「獨人之歌」.
再次與我們的生命說話.
幫助孩子講解你的話語清楚和明白.
我們一同去到謙卑領受.
奉耶穌基督的名字祈禱.
阿們.
一連四次的「獨人之歌」.
今天已經去到第四首.
也就是最長的一首.
這一首歌.
我們通常會覺得是由第五十三章的第一節開始.
但誰不知如果我們根據分段.
應該是第五十二章的十三節開始.
直到第五十三章的最後一節作為結束.
這段經文.
首先略略和大家分享.
大概我們應該怎樣對待和看待它.
首先第一.
你會發現那個「獨人」的身份.
正如我昨天和大家分享.
那個「獨人」的身份.
如果我們說第一首的「獨人之歌」.
其實就是以色列文的這個「獨人」.
但是第二,第三,第四首這三首.
全部那個「獨人」.
就不是指以色列文.
而是指先知的本人.
又或者所預表的耶穌基督.
為什麼呢.
就因為我們特別看第四首這一段.
我們就會發現那個「獨人」.
是用「我」進行稱呼.

$^{41}$然後「擔當我們的憂患」的「我們」.
就好像和這個「獨人」是兩樣東西.
就是說這個「獨人」是要擔當我們的憂患.
如果是這樣.
「我們」就是指以色列文.
又或者是泛指全球的人類.
變成和這個「獨人」是兩個不同的個體.
這是第一個讓大家知道.
「我們」的確在以賽亞書裡面.
是真的有得罪上帝的描述.
特別在以賽亞書的四十二章的二十四節.
也提到.
誰將雅各交作出來.
去作老物.
將以色列交給搶奪的.
豈不是耶和華嗎.
然後下一句大家看到.
就是我們所得罪的那一位.
那個「我們」原來是得罪上帝的.
很清晰.
我們會發現在四十二章.
我們能夠看得到.
所以你會發現「我們」是犯罪的.
「僕人」是沒有犯罪的.
「僕人」是要擔當我們的憂患.
這是第一個背景.
讓大家留意.
第二個我們會發現.
你頭那三節.
他首先說了「僕人」的狀態.
他將會去到一個什麼境地.
為優先的.
你會看到裡面很多的描述.
都是描述這個「僕人」的淒慘.
他怎樣被人鞭打.
又或者怎樣被上帝激害.
又好像一個羊.
不開口將要去到宰殺之地.
類似這些.
是一個很卑微的描述.

$^{81}$但同一時間他一開始.
你會發現他一開始的描述.
是一個非常高升的描述.
特別第一句.
「我的僕人行事必有智慧.
必被高舉上升.
且成為至高」.
一看你就會發現有兩個「高」.
第一個「高」就是被高舉上升.
我想問這個是主動還是被動?.
被動.
第二個「高」是「且成為至高」.
這個是主動還是被動?.
主動.
即是告訴我們.
無論是被動還是主動.
「高」這個位置.
是專屬給這個「僕人」.
那為什麼真的要用兩個「高」呢?.
如果你真的寫這段經文.
一個「高」就可以了.
他很高的.
那就很高.
但是為什麼真的要有兩個「高」呢?.
還是要一個主動一個被動呢?.
因為這個在《以賽亞書》裡面.
在一個專屬的兩個「高」的「打孖」出現.
在整本《以賽亞書》裡面.
總共出現三次這麼多.
第一次大家是很熟悉的.
就是來自《以賽亞書》第六章的第一節.
我們知道那一段就是.
以賽亞先知本人他蒙召的一開頭.
他這樣說.
「當烏西亞黃奔的那一年.
我看見主耶和華坐在高高的寶座上」.
大家有印象.
你和合本譯作兩個「高」.
他不是譯錯的.
不過如果你重看原文.

$^{121}$他真的這樣說.
他讀作「Ram Vanisha」這樣的讀音.
「Ram」就是自高.
是主動的高.
「Vanisha」就是被動的高.
就是說無論是主動還是被動.
他坐在高高的寶座上.
其實就是在說上帝.
在《以賽亞書》裡面這樣的記載.
「高」這個位置是不可以形容任何人.
「高」這個位置是專屬給耶和華的.
因為這個是以賽亞的意象.
他看到耶和華是高高在寶座上.
所以如果有任何一件事.
又或者人物膽敢成為「高」這個位置的話.
在《以賽亞書》的第二章.
就會將他降低.
我想活現耶穌基督在福音書的教導.
凡自高就必被降低.
自卑的話就必被升高.
變成無論是主動還是被動.
耶和華就是專屬有「高」這個位置給他.
另外一段如果你會查考.
就是57章的第十五節.
有一句這樣的說話.
我讀給大家聽.
57章它就說.
那自高自上永遠長存.
名為聖者的如此說.
我住在自高自聖的所在.
也與心靈痛悔謙卑的人同居.
要使謙卑的人的靈去到甦醒.
要使痛悔的人的心去到甦醒.
同樣地在57章的第十五節.
它是這樣形容上帝.
就說祂是自高自上永遠長存的上帝.
同樣地同樣的原文都出現在這裡.
就是說一個是主動高一個是被動高.
所以貫穿整個以賽亞書起碼有三處.
你會發現是主動高和被動高.

$^{161}$都是形容這個上帝.
不過我們回頭去這裡的時候.
52章的第十三節.
竟然是用主動的高和被動的高.
去形容一個僕人.
你會覺得很有趣.
就是說根據這樣的邏輯.
這個僕人不簡單.
他不只是人這麼簡單.
他甚至不只是一個僕人.
其實他是上帝.
因為唯獨高這個位置.
才是專屬給上帝.
所以這位的僕人是有神聖的身份.
所以難怪指向耶穌基督.
新約引用特別是彼得前書.
直接整個僕人之歌.
引用在彼得前書那裡.
是指向耶穌基督這件事.
不過當他說要成為高.
我們想像的高.
和這裡提到這個僕人的高.
根本是一個不同的概念.
我們在想那些高.
就應該是很厲害.
光芒四射.
一定是遠離凡塵.
非常之高昇.
就是那種很厲害.
君王君臨天下的那種高.
但是你看這一段的時候.
就顛覆了我們這一種的想法.
但是竟然這一位的僕人.
是他的面貌比別人憔悴.
他的形容比別人呼高.
這一種這樣的面貌.
就和高這個位置形成一個強烈的對比.
僕人的呈現.
這個僕人也會洗淨許多的國民.
特別留意洗淨這個字.

$^{201}$是一個祭司的用詞.
英文就叫做sprinkle.
sprinkle是什麼呢.
就是一個祭司.
他為了要潔淨聖所的污垢.
所以他就宰了那隻牛或者一隻羊.
拿了那些牛血或者羊血.
進到這個會幕的聖所裡面.
就將那個紙頭站於血中.
然後向著聖所的曼寺.
彈血七次.
大家有印象.
這個是利未記第四章的描述.
那個彈血七次的彈字.
其實就是這個字.
sprinkle 潔淨這樣的意思.
為什麼要彈血七次.
為什麼是七字.
就是象徵完整完全這樣的意思.
血用來做什麼的.
血是用來洗乾淨東西的.
你不會用血來抹地的.
你不是用來抹地的.
這些是用來洗乾淨人的血.
人的罪.
他這個僕人有一個這樣的特徵.
他是面目憔悴.
但卻將會成為至高.
是一個顛覆我們的思維.
這個背景底下.
我們繼續看下去.
這個僕人到底是一個怎樣的情況.
第一節.
他就說我們所傳的有誰信.
耶和華的蚌壁向誰顯露.
第一句.
其實就是說.
這個僕人將會用一個.
在人類裡面.
沒有辦法能夠預期的一個方式.

$^{241}$去進行他的拯救.
接著我們首先看到.
他的經文.
是非常多去描述這個僕人的特徵.
這個特徵很有趣.
他就說.
他在耶和華面前伸掌如亂牙.
將根出於乾地.
也無皆營美容.
我們看見他的時候也無美貌.
使我們羨慕他.
這一句我們不是很明白.
什麼叫做根出於乾地.
什麼叫做亂牙.
為什麼要用亂牙.
第二.
原來他一點都不帥.
不但如此.
原來他比一般人更加醜陋.
無皆營美容的一個這樣的狀態.
這樣的狀態在以賽亞書的場景底下.
去到怎樣理解呢.
其中一段我想大家去到記得的.
就在以賽亞書的十一章的第一節.
有一段我想大家熟悉的經文.
那段的經文就是在說.
耶西的殘鋼必發一個枝條.
大家有沒有印象.
有一段這樣的經文.
大家知不知道什麼叫殘鋼.
殘鋼即是你斬了一棵樹.
就剩下一個樹等.
我們叫這個叫殘鋼.
你一看到那個殘鋼.
即是被人砍伐了.
即是被人沒有了.
然後這個殘鋼沒有了生命.
因為死了.
然後這個殘鋼上面.
看下去是一個沒有盼望的情況.

$^{281}$他就告訴我們.
耶西的殘鋼.
耶西是誰.
耶西就是大衛的爸爸.
大衛的爸爸就說.
耶西的殘鋼就必發一個枝條.
在耶西這個殘鋼裡面.
就要生一條枝條.
很明顯是一個象徵.
那個象徵就是耶西這個爸爸.
就會生一個兒子.
那個枝條就很明顯是指耶西的兒子.
那這條枝條為什麼是在說枝條.
就是枝條的特徵就是若不襟風.
沒有人羨慕.
非常之坎坷的一段經歷.
我想大家都知道.
耶西見阿先知撒母耳來到.
我們第一課都說了.
他七個大佬都送上去.
好像佳麗那樣.
送去撒母耳那裡.
看看有沒有耶和華指使要高納他.
但是耶和華都說了.
提醒撒母耳千萬不要看他的外貌.
因為耶和華看人不像人看人.
人是看外貌.
耶和華是看內心.
所以就要選一個連七個人都不存在的.
這個爸爸覺得他不應該在這裡.
就是大衛這個孩子.
沒有人羨慕.
沒有人覺得重要.
這個大衛是帥哥.
不過在這樣的設定下.
他的美貌又或者他的外形.
用現在的術語.
就是他的profile.
他的profile簡直是不太好.
類似這樣的記憶.

$^{321}$這樣的特徵.
就叫做耶西的殘鋼.
就是說上帝拯救的盼望.
寄望竟然是一個殘鋼.
竟然是一個在當時的profile來看.
是一個非常underestimate.
是一個非常不好的情況.
同樣這樣的類比.
就放在這個僕人裡面.
再次提這個殘鋼.
他提得更加清晰.
因為他說這個殘鋼.
是新一國的枝條.
而這個枝條是伸長如嫩芽.
去描述這個僕人.
如果是這樣.
就想起我們在以賽亞書第十一章第一節的記憶.
那個記憶就是原來這個僕人.
像誰呢?.
像大衛.
大衛是耶西的殘鋼.
這個僕人伸長如嫩芽.
這是第一個記憶.
讓我們看到像大衛.
這是第一.
第二,將根出於乾地.
乾地我們再次想起殘鋼.
因為殘鋼是乾的.
是完全沒有期望.
能夠可以生長一個嫩枝出來.
誰不知那個根.
就正正出於乾地.
為什麼要特別提到乾地?.
乾地通常給人的感覺.
是不會寄望在這裡.
有任何好東西可以出來.
如果我說成這樣.
大家都知道.
很明顯就是自福音書.
在說那撒勒這個地方.

$^{361}$這個那丹葉.
曾經說過.
這個彌賽亞出於那撒勒.
那丹葉也說過.
那撒勒還能出什麼好東西.
原來他忘記了有以賽亞書.
因為以賽亞書.
就說這個殘鋼的盼望的枝條.
不是出於肥沃的土地.
而是出於乾地.
而乾地就是有一個枝條走出來.
所以那個僕人有一個大衛的形象在這裡.
一說到大衛.
你一定想起大衛的後裔.
這件事情.
大衛的後裔很明顯整件事情.
是指向耶穌基督.
是一個這樣的情況.
如果是這樣.
我們就明白下一句.
下一句就是他無家營美容.
不是真的指他很醜陋.
而是說他low profile.
就是說他的設定不是很好.
就好像什麼.
又是好像大衛.
也好像誰呢.
就是耶穌基督他自己.
不僅如此.
我們常常藐視他.
就好像他爸爸藐視大衛一樣.
所以被人嫌棄.
多受痛苦.
常經憂患.
他被藐視.
好像被人掩面不看一樣.
我們是完全不會尊重他.
變成這樣.
我們看到這個僕人的特徵.
一看下去.

$^{401}$是非常低微的一個僕人.
但種種他裡面的描述.
他是很像大衛.
同樣地.
他也很像我們指向的耶穌基督.
這是很弔詭性.
考驗我們每一個人的生命.
究竟我們真的看到耶穌基督本人的時候.
我們真的知不知道這是上帝.
或者去到這一點.
我也在這裡引用一位新約學者.
他的名字叫F.F.Bruce.
他解釋肥納比書第二章.
大家很熟悉的一段經文.
當以基督耶穌的心為心.
大家有印象嗎.
他本有神的形象.
不以自己與神同等為強奪.
反倒虛己取了奴僕的形象.
成為人的樣式.
既有人的樣子.
就自己卑微全心信服.
以至於死.
且死在十字架上.
所以上帝就稱他為至高.
大家聽著.
至高.
此給他那超乎萬名之上的名.
叫一切在天上的.
地上的.
和地底下的.
無不口稱耶穌為主使榮耀歸於父神.
大家記得.
肥納比書二章第五到第十一節.
是最早期教會的信經.
既然是信經.
大家每一個都要懂得背.
這是信經.
你會發現.
F.F.Bruce解釋這一段.

$^{441}$他說要成為人的樣式.
要成為奴隸的形象.
以至於死且死在十字架上.
我們以為.
耶穌基督.
是拿走了神的形象.
放了奴隸的形象在他身上.
即是脫去神的形象.
穿上奴隸的形象.
F.F.Bruce就認為不是這樣.
他不是說拿走神的形象.
取了奴僕的形象.
而是奴僕的形象當中.
顯出神的形象.
大家分不分得到分別.
他不是脫了一件衣服就穿一件.
他仍然是穿那件衣服.
不過那件衣服是你的想法.
是完全不能夠想像.
即是說神的形象.
是怎樣顯出來的呢.
神的形象是用奴僕的形象.
去顯出來的.
所以這一句.
就在學術界裡.
我們叫做Frequently Quoted.
即是說很多學者.
都是引用F.F.Bruce的這一句.
去進行這個解釋.
我又引用這一句.
去解釋這個以賽亞書的53章.
即是說.
越似奴隸 越似上帝.
是不是很顛覆你的想法.
看看旁邊的姐妹淺似奴隸.
是一個越卑微.
你看上去越謙卑.
越沒有佳形美容.
在那裡就顯出神的形象.
就是一個這樣的設定.

$^{481}$我們就看完這53章這裡.
我們看看這個奴隸.
似上帝形象的奴隸.
他是怎樣去對待.
和怎樣去成就救恩.
他就這樣說.
他誠然擔當我們的憂患.
背負我們的痛苦.
我們卻以為他受責罰.
被神擊打腐敗了.
哪知他為我們的過犯受害.
為我們的罪孽壓傷.
因他受的刑罰.
我們得平安.
因他受的鞭傷.
我們得醫治.
如果我們看回希伯來文.
你會發現第四節.
他就說他誠然.
誠然這個字是和學本的翻譯.
如果你看回原文.
應該翻譯成surely.
如果你用中文翻譯.
他肯定擔當我們的憂患.
你可以這樣翻譯.
第二原來那個他.
擺在和學本的最頭.
在原文裡面真的這樣.
即是說他真的強調這個他.
是他實在地肯定地.
擔當我們的憂患.
而擔當就是carry這個字.
carry即是說擔當我們的憂患.
我們昨天看僕人之歌的時候.
大家記不記得.
上帝就好像一個摺衫婆.
生了那個僕人出來.
從出胎.
上帝就抱抱他.
這個抱抱這個字.

$^{521}$正正就是擔當這個字.
即是carry著我們的憂患.
背負著我們的痛苦.
我們略略查一下希伯來文.
我們就知道這裡提到的憂患.
就不只是說你工作很辛苦.
又或者被人迫害這些憂患.
原來他原文其實是應該譯作sickness.
即是什麼意思呢.
不是憂患這麼簡單.
是病患.
又或者有些人sick.
是代表他受傷.
其實都有這個字這樣的意思.
亦都是說全球人類的狀態.
因為我們其實是受傷的.
我們是病入膏肺的.
我們因為罪的緣故.
所以是受苦的.
但是這個擔當就成為這個格式.
這個格式就是這一句.
擔當我們的憂患.
背負我們的痛苦.
大家默想一下.
大家覺得怎麼樣.
人類充滿痛苦.
人類充滿憂患.
人類非常多的擔心.
日日很多的病痛.
有戰爭有瘟疫.
有彼此的廝殺.
彼此的殘害.
彼此的陰謀彼此的毒害.
人類只是急速直線.
步向滅亡的狀態去到進化.
這個世界科技越進步.
人類的道德沒有進步過.
即是越來越退步.
在一個這樣的道德下滑的光影.
我們以為我們正在經歷的憂患和痛苦.

$^{561}$是我們獨個兒一個人來承受的痛苦.
但是這段經文給我們一種的震撼是什麼呢.
原來你經歷的憂患.
不是一個孤獨一個單一的狀態下去到存在.
彷彿就是你自己面對.
你自己去承受.
你自己日日背負著的這件事.
原來不是.
原來我們的憂患.
只是處於一種有這個僕人承受著的狀態.
這個承受.
也就是他正在擔當的一個狀態.
成為全球人類的一個很大的訊息.
那個訊息就是原來你拿著很重的憂患.
就好像一個鉛一樣.
拿得很辛苦的時候.
這個拿著.
冥冥當中.
原來耶穌基督的十字架.
是幫你一起承受著這件事.
一起承受著這件事.
就成為你的憂患的狀態.
憂患從此不再孤獨.
你面對的痛苦從此不再孤單.
原來是由耶穌基督一起去到背負的狀態下.
去呈現這個憂患出來.
我還記得很久之前.
我做傳道的時候.
認識一位弟兄.
這位弟兄很年輕.
三十多歲.
突然之間身體某個地方比較痛.
他去看醫生.
一檢查.
嚇到因為有第四期的胰臟癌.
胰臟這件事很tricky.
因為一檢查.
一痛就是第四期.
醫生就叫他太太做好心理準備.
因為他可能剩下三個月的日子.

$^{601}$晴天霹靂.
還有一個在讀幼稚園的兒子.
是一個很坎坷的一段經歷.
最後這位弟兄.
其實是一年之後才去世.
那段時間.
有時候我都會去探望他.
我還記得最後一次見面.
去探望他.
去到屯門醫院的時候.
他就和我分享.
有時候你覺不覺得.
做傳道或者做牧師.
其實都幾無能的.
你不知道可以為他做什麼.
你只可以為他祈禱.
和為他祝福.
聽他的傾訴.
醫生都可以給他一粒藥吃.
但是你什麼都做不到.
只可以陪伴他.
去到的時候.
我聽他去到分享.
他就和我說.
高傳道.
你知不知道.
我打橫訓又痛.
打直又訓又痛.
趴在這裡又痛.
點點他都痛.
現在我喝水.
喝下一杯清水.
那個痛好像剖開喉嚨一樣痛.
很辛苦.
非常辛苦.
他的面容非常憔悴.
很不舒服.
我都不知道說什麼是好的.
還是要他為他祈禱.
不過我都有一個好奇.

$^{641}$好奇問他.
弟兄.
你怎樣可以面對得到.
和捱得到這個痛苦.
其實我都好想知道.
因為我望見他.
會不會我們人類.
很多人的未來都是這樣呢.
甚至我都想起我自己.
真的.
這個不是在講咒詛.
是在講statistics.
即是說你死於癌症的人.
是三分一.
你有三分一的機會.
是這樣去到去世.
我都真是好想看一看這位弟兄.
可不可以給我一些備權.
去到怎樣面對.
這個這樣的痛苦.
他就和我說了一番說話.
我到現在都記得.
那番說話就是.
高全道.
我跟你說.
我近來默想耶穌基督的十字架.
和他的復活.
他說他所經歷的痛苦.
對比十字架的痛苦.
就可能真是九牛一毛.
但是今天我所經歷的痛苦.
能夠可以拉車邊.
可以經歷到耶穌的痛苦.
他最後說了一句.
我覺得我配成為主的門徒.
他說了這些.
令到我會覺得.
好像耶穌2.0.
就呈現在我眼前.
即是說耶穌背負他的痛苦.

$^{681}$和他一起經歷這件事情.
他知道耶穌知道他痛.
他知道耶穌是這樣走上這條路.
他就視他的痛苦.
是一個跟隨耶穌基督的步伐.
的一個必經的過程.
而這個過程.
就是十字架的經歷.
怎樣幫助這位弟兄.
去carry最痛苦的人生.
隨後他去世.
在他的安息禮拜裡.
我很感動.
他是一位跟隨主耶穌基督.
上帝的僕人.
在一個這樣的形象裡.
活現了這段經文.
上帝是擔當我們的憂患.
背負我們的痛苦.
不過去到現在這一刻.
我都心想.
其實我看著他.
我想起他.
之後我能不能夠做得到呢?.
有時信仰是可以很純真.
有一位弟兄.
因為我在友愛堂負責弟兄組.
這個弟兄組是不親姊妹的.
全部都是弟兄.
但是他們年紀都大.
這十幾年開這個組.
都送走了五六位弟兄.
我們合起來.
每一位弟兄的安息禮拜.
我們組織男人師班.
到底在你的安息禮拜.
能夠點一首歌.
確定你還沒死之前.
我們就在安息禮拜裡.
去唱頌.

$^{721}$其中一位弟兄患了血癌.
他上天去屯門醫院輸血.
然後到最後他坐輪椅.
非要回友愛堂的祖堂.
他上上笑容滿面.
見到我高牧師.
很感恩.
整天都說感恩.
接著我問他太太.
他是不是裝模作樣.
在牧師面前經常說感恩.
太太說不是啊.
他家裡也是這樣的.
所以我很佩服.
如果我將來有血癌.
會不會好像他那樣呢.
所以他點了一首歌.
就叫做《每一天》.
大家有沒有聽過《每一天》.
每一天 所度過的每一刻.
我得著能力勝過試煉.
我依靠天父奏長的公認.
我不論再恐慌與掛念.
很美麗的詩歌.
如果你認識的話.
都是我那個年代.
不過你會明白.
這首歌怎樣carry住.
這位弟兄每天.
差不多每個星期.
都要去屯門醫院輸血的經歷.
他痛的.
他不是不痛的.
但他每天都懷著喜樂和感恩.
上帝carry住這位弟兄.
上帝carry住那位弟兄.
直到白髮.
上帝都carry住.
就是因為我主耶穌基督.
步上十字架.

$^{761}$他成為全球人類的拯救.
他與我們同在.
他carry住我們的生命.
所以上帝升他.
升為至高.
超乎萬名之上的民.
所以他今天復活.
並且升天.
我們今天去到歌頌站尾.
原來死亡.
不是我們人生的句號.
死亡只不過是逗號.
是要經歷上帝使你復活.
榮耀的.
再來的一個句號.
而這個句號.
是要宣告他戰勝死亡.
他能夠身體復活.
他能夠戰勝病患.
他能夠可以紫色刀劍.
帶來真正的和平.
我們今天慶祝這個永恆的盼望.
所以這一首詩歌.
都記在我們的心中.
盼望僕人之歌.
再次成為我們的鼓勵.
我們一起祈禱.
慈悲仁愛的天父.
多謝你透過今天的經文.
我們知道你常經憂患.
背負痛苦.
你是至高至上的上帝.
卻與心靈謙卑的人去到同住.
叫我們謙卑的人去到甦醒.
提醒我們的耳朵.
讓我們不自失去盼望.
因為我們今天相信.
並且承認.
你已經去到復活.
你戰勝死亡.

$^{801}$賜下生命.
讓我們得著永恆的確具.
奉耶穌基督的名祈禱.
阿們.
\newpage



\section{出埃及記 4:11}
\label{sec:031913}
\textbf{聖經人聖經事:誰造人的口呢?}~~~~黃克勤牧師博士
\newline
\newline
連結: \href{http://www.church.com.hk/acms/content.asp?site=cdc\&op=show\&type=product\&code=031913\&layout=sermon}{\texttt{031913}} ~~~~ 語音日期: 2026-04-01
\newline
\newline
\hyperref[sec:031912]{\small{< < < PREV SERMON < < <}}
~
\hyperref[sec:index_chronic]{\small{[返順時目]}}
~
\hyperref[sec:031914]{\small{> > > NEXT SERMON > > >}}
\newline
\newline
出埃及記 4:11
\newline
\begin{longtable}{cl}
\hline
\hline
章節 & 經文 (和合本修訂版)\\
\hline
4:11 & \begin{tabularx}{0.7\textwidth}{X} 耶和華對他說：「誰造人的口呢？誰使人口啞、耳聾、目明、眼瞎呢？豈不是我－耶和華嗎？ \end{tabularx} \\ \\
[1ex]
\hline
\hline
\end{longtable}
$^{1}$各位弟兄姊妹 大家好.
今天就是繼續跟大家講聖經人,聖經事.
是生命樹宣教網絡為大家提供的.
今天就是因為各樣的關係.
這兩天香港是一個比較長的長週末假期.
各樣的事 也是很開心的母親節.
活動各樣比較多.
趕不及做一個在我們平時用大電視的錄影.
因此這次就改了用Zoom來做一個錄製.
因此大家在Zoom見不見得到我一個很微型的樣子.
好像張喬喬仔那樣 不要緊的.
最重要是看那些聖經 不是很需要看我的.
聖經人物和聖經的故事.
最重要就是看到它背後上帝在這裡做了些什麼.
希望聖經人和聖經事是可以幫你認識上主多些.
這個是最重要的 也是我們一直在禱告的事.
今天的題目仍然是在講摩西.
這句話是神和摩西說的.
誰做人的口呢?.
神很多時候都很有趣的.
他是很厲害的 如果我們學人生教練的技術.
即是生命教練 做教練時常都會問人問題的.
上帝經常問問題的.
問了問題之後就令到你回答問題.
就算你在口裡不回答 心裡也會自己回答.
這是人來的 人的潛意識會這樣的.
誰做人的口呢?主咯 跟著看下去.
好 繼續看摩西故事.
我們已經開始由阿巴拉罕伊薩雅閣藥室.
看完《創世記》就來到《出一及記》看摩西故事.
誰做人的口呢?.
《出一及記》四章第十節.
摩西對耶和華說.
主啊 我素日不是能言的人.
就是從你對僕人說話以後也是這樣.
我本是住口奮舌的.
摩西對神另外一個推搪的辦法就是.
嘩 你這麼厲害.
又預備了三個神蹟給我.
去見就一定行 但摩西也不想去.

$^{41}$先問一個問題就是.
如果以色列人問我你叫什麼名字.
是你吩咐我帶他們來的.
帶他們走的 那我怎麼說呢?.
其實以色列人是不會問這個問題的.
因為他們由阿巴拉罕時代下來.
就已經知道耶和華神了.
都不會問的 隨便找一個.
真是住口奮舌的摩西問過這樣的問題.
然後神就說 行了 我叫做阿M.
然後就一直下來.
他不會讓我見他.
我給你三個神蹟握手.
然後他就一定放的.
嘩 這麼厲害.
摩西又來了 下一招是什麼呢?.
我本是住口奮舌.
素日不是能言的人.
我平時是不會說話的.
說話咯咯聲的.
我住口奮舌的.
就是現在你跟我說話.
我都沒有進步到.
耶和華就對他說 你想一想.
誰做人的口呢?.
可能前面如果我是神.
我就會加一句.
我早就知道你不會說話了.
不過你要想一想 誰做人的口呢?.
神啊 耶和華神啊.
其實是想說.
有些人是你說自己住口奮舌不行的.
有些人是啞的.
有些人是聾的.
有些人是失明的.
有些人看不到東西.
那些其實都是我做的.
那這個誰做人的口呢?.
這個是一個回應與挑戰來的.
第一 你是不是說我做得不好?.

$^{81}$是不是說我做你做得住口奮舌不夠好?.
不是 不是 不是.
我不是這個意思 這樣就行了.
其實有些啞 有些失明 有些乞眼.
其實都是我在背後有所安排.
有我的心意的.
所以如果你說你住口奮舌.
其實都是我讓你住口奮舌的.
但是你的口是我做的.
好 現在不怕.
由於這個口是我做的.
所以四章十二節.
現在你去吧 儘管去見法老.
我必賜你口才.
剩下一枝螢光筆.
黃色.
然後現在你去.
我必定會賜給你.
你不用怕.
我給你口才.
還有教你要當說的話.
這樣行了吧.
這樣就是我說的了.
這樣就是你一定有口才.
而且知道說什麼.
第十三節.
這一節很難明白.
第十三節很難明白.
現在摩西就說.
主啊 你願意打發誰就打發誰去吧.
他是什麼意思呢?.
是不是現在乖了很順暢.
這樣去吧.
你既然打發我 我就去吧.
但是有一點不是.
原因就是現在這裡是第十三節.
接著下去的第十四節.
耶和華是向摩西法老生氣.
因為他說這句話.
所以這句話要另外再想.

$^{121}$主啊 你願意打發誰就打發誰.
現在這個「你願意」其實就是.
Now I beseech.
就是Now I ask.
就是現在請你打發誰就打發誰.
這是什麼意思呢?.
其實會不會這句來到下面第十四節.
說生氣呢?.
主啊 既然是這樣.
那你打發誰就打發誰.
其實就是說你可以打發一些口才厲害一點的.
那些口才全部是你做的.
你打發他厲害一點的口才去就行了.
那不用叫我啊.
就喜歡叫誰就叫誰吧.
其實我不拘的 我不介意的.
雖然這件工作實在是超級偉大.
還有超級好.
帶領我的百姓離開圍爐.
做奴隸的地方 非常好.
不過你叫誰就叫誰吧.
其實暗喻呢 就不用叫我啊.
我推測是這樣.
所以你看到第三點就來了.
就是耶和華向摩西發怒說.
第十四節了.
耶和華向摩西發怒 生氣啊.
因為它上面第十三節說了一句話.
不是有你哥哥利未人阿倫嗎?.
如果說其他人 有.
就是你哥哥阿倫.
很有趣的.
你哥哥利未人.
其實摩西的利未人不用說都可以了.
不過他們喜歡說利未人.
利未人特別要提 利未人是什麼意思呢?.
就是一幫將來我要用他的支派的人.
就是要來事奉我的.
現在沒有說 他們都不知道的.
因為還沒有到再合紀去到西海山.

$^{161}$其實你哥哥就跟你同宗同族的那個利未人阿倫.
記不記得?你肯定記得他的.
我知道他是能言的人.
為什麼會你知道呢?.
因為我做的 誰做的?我做的.
現在不僅如此 他已經等著出來迎接你了.
其實現在摩西自己一個人在米殿巷惹禍.
他如果下一章你會見到聖經.
他就回去跟他的岳父說.
我想回去埃及地看看我的兄弟.
他的岳父葉蒂羅就說好 你回去吧.
就帶著他的女兒回去.
他現在阿倫在埃及那裡都在做奴隸的.
他一見到你 他心裡就很歡喜.
其實他出來迎接你.
就是說他將會出來迎接你.
他很掛念你.
一見到你 他心裡就很開心 為什麼?.
哎呀 原來弟弟摩西還沒死啊.
這麼多年沒見.
而且他回來帶著使命 第三 一.
你要將當時的話傳給他.
我就要賜你和他厚財.
又要指教你所當行的事.
你要將當時的話傳給他.
先將我跟你說過的東西.
不要加鹽加醋 不要說那麼多.
不要說那麼多閒話 沒有那麼多無謂的東西.
將當時的東西先傳給他.
於是你們兩個可以一起同心.
然後我要賜傳給他.
於是就有一個你和他出現了.
你將他 於是你和他是擁有著相同的.
在神那裡聽到的說話.
這樣就好了 因為這樣他就會說對話了.
不然這樣說是很厲害的.
嘴巴很厲害 不過沒用 說不中.
我要賜你和他有厚財.
而且我還要教你 當行的事.
所以這裡是有厚 又有時.

$^{201}$又說要懂得說 還有懂得做對的事情.
哎呀 那我怎麼知道 見法老王啊.
真是隨時一個眼色就殺頭改低幾寸.
那麼害怕 那怎麼辦?.
不怕 放心 我給你厚財.
我教你怎麼做.
好了 他要替你對百姓說話.
我們這裡是有特別的解釋的.
「他」是指亞倫.
亞倫要代替你對百姓說話.
於是你以他做「厚」.
「他」當作「厚」.
而「你」就當作「上帝」.
不是吧?那個人當作上帝?.
其實是這樣的意思.
意思就是 你說自己厚財不好.
那你說不行 你喜歡派別人就派別人吧.
派誰就派誰吧 其實心裡想.
不要派我就行了 只要派誰就派誰.
那你可以把那些厚財給那個人好一點.
那麼上帝就發火了 第十四節 是不是?.
你記不記得你那個哥哥亞倫啊?.
他很能說話 很能承諾的.
他現在等著你 他見到你就會很開心.
你回去找他一起做事吧.
他這麼厲害 我要教你兩兄弟怎樣做.
你將我的意象告訴他.
於是「他」亞倫要代替你對百姓說話.
「他」當作你的「厚」.
而「他」是跟著你說的.
他會以你為上帝的帶領.
那就不是說 以你真的做了上帝.
而是說 以你做上帝帶領的表示.
你說甚麼 啊 神的聲音這樣.
所以他就會很忠誠地聽你說的.
於是這裡分開兩件事了.
你手裡要拿著像 一扔到地上就變成蛇那支像.
你手裡要拿著這支像 這個也是一個憑證來的.
這支像 然後你要拿穩它.
不然就掉了 哎呀 糟了 這支像不見了.

$^{241}$這支樹枝行不行?不行.
是這支像 「好行神蹟」.
這樣講得很清楚的.
在這裡我自己覺得 哎呀.
你見到上面說.
他要說當時的話 你們兩個都有口才.
現在他做「厚」 你做「帶領」.
好像上帝那樣.
如果這樣看 摩西就因為那句話.
因為那句推搪.
「啊 神啊 你喜歡派誰就派誰吧」.
「其實你可以派別人啊」.
有些他沒有說出來.
不過心裡面可能說了.
神就發怒 第十四節.
結果就是甚麼呢?.
亞倫出現了.
亞倫出現是令到摩西的光環.
是分了一半給亞倫的.
本來神只是叫摩西 沒有叫亞倫.
是因為摩西在推搪 太極推手很厲害.
結果他的工作 或者說我這樣講.
他的榮耀有一半給他的兄長亞倫拿了.
如果摩西自己說.
「好吧 既然你賜我口才」.
「我就聽你講 死就死 大膽去吧 好啊」.
其實自己又出口 又做帶領者.
不需要找亞倫的.
而亞倫後來由於他代摩西出口.
也都你知道在抗日出了事.
就是摩西上了山 西奈山很久.
等候律法沒下來.
那些人作亂.
「哎呀 摩西上山這麼久」.
「都不知道他遭遇甚麼事」.
「不知道有沒有死了在上面」.
「那怎麼辦啊」.
然後說逼得他很厲害.
然後亞倫走出了一步.
「好吧 大家把你們家裡太太的金耳環」.

$^{281}$「金首飾和家裡的金」.
「我們一起用它做一隻金牛」.
「明日起來 明日我們要向耶和華守祭」.
其實那隻金牛他們想不到.
沒有了摩西 當作上帝的那支象不見了.
「糟了 那怎麼辦啊」.
「那我們找一個神來自己帶我們吧」.
「找誰啊?」.
你知道亞倫一沒有了摩西做上帝的帶領.
代表他就失手了.
我們想起埃及帶我們的就是金牛.
「那你就找金牛吧」.
於是就這樣.
「來啊 明日我們要向耶和華守祭」.
那這個呢 他是向耶和華守祭.
他用金牛代替耶和華.
所以其實亞倫就是給了他進來.
分享了摩西的榮耀.
另一方面也是給了他一個隱憂.
神就說「你手裡拿著這支象」.
「好行神蹟 拿穩了」.
今天跟大家分享到這裡.
有三點的.
第一點就是.
誰做我絕口笨舌.
其實知道自己絕口笨舌是好的.
因為你知道自己絕口笨舌.
你會依靠神.
第二神就問「誰做我絕口笨舌?」.
「我啊 神自己說的 必是我啊」.
因此你依靠我就可以了.
我會教你怎麼說話.
我會教你怎麼做事.
「找別人行不行?」.
「又來?好」.
其實我可以不叫你.
我就叫你哥哥也可以.
不過我叫亞倫大哥帶領你.
跟你一起去.
於是就有神明說.

$^{321}$就叫亞倫出來了.
「你啊哥」.
這樣叫他出來.
神向他發怒.
第一次我們不要惹到神向我們發怒.
如果神要你有工作做.
你就做吧.
「哇 不夠力啊 不夠力啊」.
「用你力的嗎?」.
「上帝說我有力啊」.
「誰最大力啊?」.
「哎呀 我沒本事」.
「本事?」.
「啊 真是」.
「上帝最有本事」.
你依教他就可以了.
不要像摩西那樣.
令到神向他發怒.
將你自己的工作交給別人.
或者不用你.
那就很慘了.
我們一起禱告吧.
基本上上帝我們求你幫助.
我們真的很需要你.
有時候我們面對一個很大的挑戰的時候.
或者覺得自己.
七六八十都不想做了.
這件事留給年輕人.
這件事不要緊.
我在這裡已經過慣了四十年.
在曠野牧羊的生活.
又回到政治裡面.
帶二百萬人出埃及.
又和埃及王對撼.
哎呀 這件事不要了.
不是我.
年輕一點先吧.
找年輕一點的做吧.
簡選了這個人就差異他.
幫助我們對你更有信靠.

$^{361}$令我們有一個聆聽的心.
聽到你對我們的心意.
我們就全心信服.
而且勇敢的去執行.
幫助我們.
知道你是做我們嘴裡的主.
你是值得讚美的.
願你得到榮耀.
祈禱奉耶穌名求.
阿們.
\newpage



\section{使徒行傳 10:34-43}
\label{sec:031914}
\textbf{跨越邊界的愛}~~~~林玉英牧師
\newline
\newline
連結: \href{http://www.church.com.hk/acms/content.asp?site=cdc\&op=show\&type=product\&code=031914\&layout=sermon}{\texttt{031914}} ~~~~ 語音日期: 2026-04-05
\newline
\newline
\hyperref[sec:031913]{\small{< < < PREV SERMON < < <}}
~
\hyperref[sec:index_chronic]{\small{[返順時目]}}
~
\hyperref[sec:031915]{\small{> > > NEXT SERMON > > >}}
\newline
\newline
使徒行傳 10:34-43
\newline
\begin{longtable}{cl}
\hline
\hline
章節 & 經文 (和合本修訂版)\\
\hline
10:34 & \begin{tabularx}{0.7\textwidth}{X} 彼得開口說：「我真的看出神是不偏待人的。 \end{tabularx} \\ \\ \relax
10:35 & \begin{tabularx}{0.7\textwidth}{X} 不但如此，在各國中那敬畏他而行義的人都為他所悅納。 \end{tabularx} \\ \\ \relax
10:36 & \begin{tabularx}{0.7\textwidth}{X} 神藉著耶穌基督—他是萬有的主—傳和平的福音，把這道傳給以色列人。 \end{tabularx} \\ \\ \relax
10:37 & \begin{tabularx}{0.7\textwidth}{X} 這話在約翰傳揚洗禮以後，從加利利起，傳遍了猶太。 \end{tabularx} \\ \\ \relax
10:38 & \begin{tabularx}{0.7\textwidth}{X} 神怎樣以聖靈和能力膏了拿撒勒人耶穌，這都是你們知道的。他到處奔波，行善事，醫好凡被魔鬼壓制的人，因為神與他同在。 \end{tabularx} \\ \\ \relax
10:39 & \begin{tabularx}{0.7\textwidth}{X} 他在猶太人之地和耶路撒冷所行的一切事，有我們作見證人。他們竟把他掛在木頭上殺了。 \end{tabularx} \\ \\ \relax
10:40 & \begin{tabularx}{0.7\textwidth}{X} 第三天，神使他復活，使他顯現出來； \end{tabularx} \\ \\ \relax
10:41 & \begin{tabularx}{0.7\textwidth}{X} 不是顯現給所有的人看，而是顯現給神預先所揀選為他作見證的人看，就是我們這些在他從死人中復活以後和他同吃同喝的人。 \end{tabularx} \\ \\ \relax
10:42 & \begin{tabularx}{0.7\textwidth}{X} 他吩咐我們傳道給眾人，證明他是神所立定，要作審判活人、死人的審判者。 \end{tabularx} \\ \\ \relax
10:43 & \begin{tabularx}{0.7\textwidth}{X} 眾先知也為這人作見證：凡信他的人，必藉著他的名得蒙赦罪。」 \end{tabularx} \\ \\
[1ex]
\hline
\hline
\end{longtable}
$^{1}$各位弟子們平安.
邀請大家再次低頭祈禱.
用我口中的言語.
頌弟兄姊妹的心思意念.
同夢上主的願.
祈禱奉主耶穌 基督的名求.
阿們.
剛才的講題是講「跨越邊界的愛」.
那邊界不知道大家腦海裡浮現的是什麼呢?.
深圳河?.
是需要簽證過的海關?.
還是那種道不同不相為謀的距離感呢?.
我們想像一下.
如果今天澳門特區政府突然宣布.
因為港澳兩地情同手足.
所以今年派錢.
香港人都有份.
只要你拿著香港的永久居民身份證.
跨過港珠澳大橋去到澳門.
不用排隊.
立即領一萬元.
你覺得怎樣呢?.
我相信大家不單很開心.
簡直會瘋狂.
你會驚嘆的是.
澳門特區政府竟然那麼大方.
去到這個地步.
原來那條心粉.
那條福利的邊界.
是可以消失的.
不過現實之中.
邊界往往是很冷冰冰的.
在二千多年前.
這條邊界更加是一道不可逾越的鴻溝.
今日經文的主角是彼得.
他認為很嚴謹.
嚴謹到近乎死硬派的猶太領袖.
對他來說.
生命中有一條不能夠逾越的紅線.
甚麼紅線呢?.

$^{41}$猶太人是上帝的選民.
他們是聖潔的.
所以那些不是猶太族裡的外邦人.
你和我都是外邦人.
是世俗的.
是不潔淨的.
當時的規矩寫得很清楚.
猶太人是不能夠進入外邦人的家.
進入後自己都變得不潔淨,骯髒,骯髒.
就算你說不屋企不要緊.
去外面的餐廳吃飯.
對不起,也不行.
你是不可以和外邦人同桌吃飯的.
不過今日這段經文最震撼我們的地方.
就是這位彼得.
竟然進入羅馬軍官哥尼劉的家講道.
就好像兩界一原本不同國籍.
甚至是迎面交鋒的戰機.
應該會打架的.
竟然在一個交互點停下來.
握手成為朋友.
這個是打破千多年的傳統.
粉碎了族群之間隔下的一個歷史時刻.
其實你有沒有留意.
我們每個人心裡都有條邊界.
可能是種族.
可能是立場.
可能是身份.
甚至可能是我們心底裡的一份自卑.
好像我這樣滿身傷痕的人.
上帝還會愛我嗎?.
今天是基督復活主日.
我們慶祝的不單止是耶穌復活了.
更加慶祝上帝主動去拆了道牆.
彼得在這段經文裡一開口.
第一句說話就擲地有聲地說.
我真的看到上帝是不偏袒人的.
這句說話對彼得來說.
是打破了他幾十年以來所持守的價值觀.
為何他會有這麼大的轉變呢?.

$^{81}$史蘭主第十章很詳細告訴我們原因.
因為上帝對彼得進行了一場很震撼性的教育.
上帝讓彼得看到一個異象.
有一塊大布從天降下來.
裡面裝滿了猶太律法裡面.
這位不潔淨的飛禽走獸.
上帝吩咐彼得.
喂 彼得 宰了它來吃.
彼得反應很快.
他說 主人 不行的.
我從小到大都很自律.
那些俗世的不潔淨的食物.
我從來都沒有吃過.
不過上帝卻連續三次對他說.
上帝所潔淨的 你不可以當做俗物.
這三次對話更加擴張了彼得的心胸和眼界.
上帝告訴他 彼得.
你的上帝太小了.
你的愛也很窄.
正當彼得正在疑惑和掙扎的時候.
聖靈的聲音就到了.
聖靈吩咐彼得要去見哥尼劉.
哥尼劉是誰?.
哥尼劉是羅馬的軍官.
他是外邦人.
在猶太人的眼裡.
他更加是侵略者 壓迫者.
是幫海撒大帝收稅.
去管治他們的敵人.
不過聖經對哥尼劉的評價.
卻令人感到很意外.
他說哥尼劉虔誠.
是一個敬畏上帝的人.
經常周濟窮人.
還幫助了猶太人.
而且時時向上帝祈禱.
最奇妙的地方是.
原來上帝也透過聖靈.
在哥尼劉禱告的時候.
去吩咐他 派人去請彼得來.

$^{121}$當彼得跨越了他的地理.
和心理邊界.
去到哥尼劉家裡的時候.
發現哥尼劉已經預備好了.
請所有家人坐在那裡.
去等究竟聖靈要透過彼得.
向他們講什麼話.
兩個原本水火不相容邊界的人.
同樣一回事.
因為他們願意聽聖靈的話.
而相遇了.
這個事蹟給我們什麼啟示呢.
各位靈姐妹.
我們會不會在心裡將人分裂了.
這個搞政治的.
他很難信耶穌.
那個人生活很放蕩.
教會也不太疼他.
這個人是我的死對頭.
我不想和他坐在一起.
你不要跟我回教會.
很多很多這些成為了我們心裡的邊界.
不過上帝看人是怎麼看的呢.
他不是看你的國籍.
他不是看你的過去.
他甚至不看你現在的身份.
他看的是什麼呢.
他重視的是你裡面的火.
尋求真理的心.
耶穌在世界上最後那三年.
他是一個最典型的擦腸鬆綁的人.
他醫治了大麻瘋的病人.
這群人代表了一群被社會所遺棄的人.
耶穌和瑞利一起同桌吃飯.
這群是代表當時的漢奸.
代表了一群叛亂國家的人.
耶穌更加安慰了撒瑪尼亞人.
代表了一群被歧視的族類.
耶穌用行動告訴我們.
上帝的恩典是沒有界限的.

$^{161}$只要你願意尋求祂.
這條邊界就會消失.
不過這位充滿自我的耶穌.
他是被當時自以為最聖潔.
最守規矩的宗教領袖送上絕路.
他們生安白造.
捏造罪名.
說他顛覆政權.
連審判他的彼拉多也說.
「我查不到他有什麼罪」.
不過那些被仇恨遮蓋了眼睛的群眾.
寧願釋放一個殺人放火的重犯巴拉巴.
都要立志大聲說.
要將耶穌釘在十字架上.
十字架是人類歷史裡最殘酷.
最羞辱的邊界.
十字架標誌著一個人被上帝唾棄.
被人類咒詛.
要被社會完全毀滅.
當耶穌斷氣那一刻.
聖經說「聖靈裡的萬子.
從上到下裂開兩半」.
這是一個很強大的象徵.
告訴我們人類和上帝之間.
由罪所築成的牆崩塌了.
耶穌死了這三天.
去了哪裡?發生了什麼事呢?.
相信不少香港人都聽說過KK園.
那裡是一個黑暗被邪惡勢力操控.
進去後你沒有想過可以出來的地方.
耶穌的復活就好像.
祂親自殺入黑暗的地方.
直搗黃龍.
將死亡的門撬開.
將罪惡的鎖鏈斬斷.
耶穌是否真的復活了?.
比特系社一節特別強調.
他們一群門徒是和耶穌同吃同喝的人.
復活的耶穌不是鬼魂.
祂不是門徒.

$^{201}$因為傷心過度產生的幻覺.
耶穌叫門徒去摸他們手上的釘痕.
叫他們和祂一起吃飯.
這是一個肉身的復活.
代表著耶穌的能力.
可以觸及我們現實生活中.
每一寸的痛苦和軟弱.
今天剛剛好是聖餐主日.
等下當我們拿起這個聖餐餅和聖餐杯的時候.
我們不是在演戲.
我們是和這位活生生.
跨越了死亡的耶穌基督同台吃飯.
我想對在座還未信耶穌的朋友.
又或者你還在邊緣觀察中的朋友說.
雖然一會兒你暫時不能去領用聖餐.
不過這並不代表上帝將你拒絕在門外.
相反耶穌已經為你預留了位置.
祂渴望鑽入你的生命當中.
祂不是要審判你.
祂是要分擔你的重擔.
讓你經歷 讓你體會祂對你的愛.
祂不是住在天上遙不可及的一個神像.
祂是一個可以陪你走過高山低谷.
一起吃飯 一起面對生活難關的朋友.
最後彼得在43節做了一個終極的邀請.
凡相信祂的人必蒙赦罪.
罪是什麼呢?.
罪不單止是我們犯了法.
罪更加是一道很厚的牆 很厚的邊界.
祂讓我們與上帝隔絕.
讓我們與人之間充滿了猜忌.
充滿了仇恨 充滿了比較和自私.
不過復活的意義就是要修復關係.
上帝主動跨過這條邊界找回我們.
為我們跨出自己的邊界 跨出我們的框框.
去接觸那些與我們有隔閡的人.
去接觸那些被我們封在邊界以外的人士.
在今天的聖餐禮安靜的時刻.
我有兩個邀請.
第一 對於已經相信耶穌的弟兄姊妹.

$^{241}$我邀請你 當你接受餅和杯的時候.
請你安靜地去問上帝.
主啊 我心裡還有哪一道牆沒拆掉呢?.
是不是我對家人還有怨恨?.
是不是我對某位同事 同學還有偏見?.
是不是我對某位教友還有意見呢?.
我願不願意好像彼得一樣.
聽從聖靈的吩咐.
去跨出我的舒適圈.
去接觸 去認識.
嘗試去愛那個我原本很不喜歡的人.
對於未信耶穌在尋求祂的朋友.
如果你今天覺得心靈很疲累.
覺得這個世界很冷漠.
或者你覺得自己離上帝很遠.
我想邀請你在心裡做一個很簡單的祈禱.
請你說 主耶穌.
如果你真是那位萬人的主.
宇宙的創造者.
請你跨越我心裡這條邊界.
去進入我的生命.
我不想再被罪和絕望困住.
請你帶給我真正的平安和自由.
復活至的愛是跨越邊界的愛.
這份愛不需要你有什麼優越的背景.
又或者你要成為你心目中很理想的人.
祂只需要你有一顆願意接受耶穌的心.
願意這份愛今天就拆毀我們中間所有的牆.
邀請大家跟我一起去祈禱.
親愛的主耶穌.
我們每一個人都有一條邊界.
這條邊界給了我們很多很多的安全感.
同樣.
將我們不喜歡,很害怕的人隔絕在外面.
上主求你寬恕.
我們知道你死在十字架上.
為了要將我們和上帝所隔開的牆去拆毀.
讓我們可以和你和好成為你的兒女.
同樣你也期望我們願意擴張我們所設定的邊界.
向那些我們拒絕,覺得有隔閡.

$^{281}$不願意接納或者寬恕的人.
一些排除在我們邊界外的人去修復關係.
你的復活.
標示著你有能力去幫助我們.
上主.
其實我們真的不知道怎麼做.
不過只要我們願意.
並且來到你的面前持續地去禱告.
你一定會指教我們.
並且我們愛的力量一小步一小步地.
去跨越這條邊界.
上主我們身邊還有很多未信的人.
求你挑和我們的心.
賜給我們決心和行動.
用愛去擴張這條邊界.
引領更多人來到你的面前.
去領受你給我們的救恩.
讓我們的生命可以更趨圓滿.
我們這樣交託,祈禱,賴藉.
救主耶穌基督的名而求.
阿們.
\newpage



\section{哥林多前書 15:1-11}
\label{sec:031915}
\textbf{基督復活的大能}~~~~王正傑牧師
\newline
\newline
連結: \href{http://www.church.com.hk/acms/content.asp?site=cdc\&op=show\&type=product\&code=031915\&layout=sermon}{\texttt{031915}} ~~~~ 語音日期: 2026-04-05
\newline
\newline
\hyperref[sec:031914]{\small{< < < PREV SERMON < < <}}
~
\hyperref[sec:index_chronic]{\small{[返順時目]}}
~
\hyperref[sec:031916]{\small{> > > NEXT SERMON > > >}}
\newline
\newline
哥林多前書 15:1-11
\newline
\begin{longtable}{cl}
\hline
\hline
章節 & 經文 (和合本修訂版)\\
\hline
15:1 & \begin{tabularx}{0.7\textwidth}{X} 弟兄們，我要你們認清我先前傳給你們的福音；這福音你們領受了，又靠著它站立得住， \end{tabularx} \\ \\ \relax
15:2 & \begin{tabularx}{0.7\textwidth}{X} 你們若能夠持守我傳給你們的信息，就必因這福音得救，否則你們是徒然相信。 \end{tabularx} \\ \\ \relax
15:3 & \begin{tabularx}{0.7\textwidth}{X} 我當日所領受又傳給你們的，最重要的就是：照聖經所說，基督為我們的罪死了， \end{tabularx} \\ \\ \relax
15:4 & \begin{tabularx}{0.7\textwidth}{X} 而且埋葬了；又照聖經所說，第三天復活了， \end{tabularx} \\ \\ \relax
15:5 & \begin{tabularx}{0.7\textwidth}{X} 還顯給磯法看，又顯給十二使徒看， \end{tabularx} \\ \\ \relax
15:6 & \begin{tabularx}{0.7\textwidth}{X} 後來一次顯給五百多弟兄看，其中一大半到現在還在，卻也有已經睡了的。 \end{tabularx} \\ \\ \relax
15:7 & \begin{tabularx}{0.7\textwidth}{X} 以後他顯給雅各看，再顯給眾使徒看， \end{tabularx} \\ \\ \relax
15:8 & \begin{tabularx}{0.7\textwidth}{X} 最後也顯給我看；我如同未到產期而生的人一般。 \end{tabularx} \\ \\ \relax
15:9 & \begin{tabularx}{0.7\textwidth}{X} 我原是使徒中最小的，不配稱為使徒，因為我曾迫害過神的教會。 \end{tabularx} \\ \\ \relax
15:10 & \begin{tabularx}{0.7\textwidth}{X} 然而，由於神的恩典，我才成了今日的我，並且他所賜給我的恩典不是徒然的。我比眾使徒格外勞苦；其實不是我，而是神的恩典與我同在。 \end{tabularx} \\ \\ \relax
15:11 & \begin{tabularx}{0.7\textwidth}{X} 無論是我或是其他使徒，我們都如此傳，你們也都如此信了。 \end{tabularx} \\ \\
[1ex]
\hline
\hline
\end{longtable}
$^{1}$弟兄姊妹平安.
今日是復活節的主日.
可否跟你隔壁的一個人說一聲.
復活節快樂.
有一對夫婦.
他們到以色列去旅遊.
旅遊的期間.
太太突然去世.
先生需要幫她處理後事.
他跟處理後事的人說.
問他怎樣處理.
處理後事的人跟他說.
你可以用一千五百元美金.
可以將你太太送回你們的國家.
但如果你在以色列幫她處理後事.
只需要一百五十元美金.
先生一聽到之後就說.
馬上就說他選一千五百元.
要將他太太送回她的國家去處理後事.
去埋葬.
幫他處理後事的工人就跟他說.
你真的很愛你太太.
願意花這麼多錢將你太太的遺體送回國家.
先生就跟他說.
因為我聽到有一個人.
曾經在以色列這個地方埋葬之後.
三日會復活.
所以我要將我太太送回我國家.
今日弟兄姊妹.
我們回來慶賀耶穌基督的復活.
我們要去思想.
要去問的是.
究竟耶穌基督的復活.
對我們的生命有甚麼意義.
對很多人來說.
耶穌基督的復活是三日的公眾假期.
相信我們當中除了基督徒之外.
其實很多未信主的朋友.
就算不信主.
都應該要多謝耶穌.

$^{41}$因為耶穌的出生和復活.
都帶給他們公眾假期.
並且每一年都會有請假攻略.
如果想再上班的人.
你就很明白很需要請假攻略.
弟兄姊妹.
今日你知不知道復活節我們可以怎樣有請假攻略.
復活節我們可以請四日就放九日.
可能今日我們有弟兄姊妹不在我們當中的時候.
他們都已經去了請假攻略.
但對於我們基督徒來說.
究竟復活的意義是甚麼.
主耶穌已經完成了九因.
我們都得了天堂的入場券.
每年我們都是紀念耶穌基督.
祂為我們的受苦.
所以我們有受苦節的崇拜.
在主日就會有復活節的崇拜.
但我要問的是弟兄姊妹.
是否就是到此為止.
耶穌基督的復活對我們的意義是否就是這樣.
所以今日我們要去看這段經文.
我們一起去明白.
究竟復活對我們的意義究竟是甚麼.
當我們看這段經文的時候.
特別是在哥林多前書的十五章.
我們會看到復活是甚麼呢.
是我們基督信仰的一個核心.
亦都是改變我們生命的一個根本的力量.
一個基督復活的大能.
保羅在這裡用了很清楚而且有力的方式.
去提醒我們福音的本質.
福音的本質其實保羅他表達的時候.
不是一種很抽象的理念.
即不是我們理解不到的.
是我們理解得到.
同時間亦都在講的.
不是我們很努力要去怎樣.
因為這個宗教而幫我們導我們向善.
即我們總之由一個不行善的人.

$^{81}$而又變成一個行善的人.
而是在講一個很真實又發生.
是帶著能力是一種改變的能力.
一種更新的能力.
這份能力是在講耶穌基督從死裡復活.
所賜給我們每一個.
所以在我們生命裡頭.
我們都會有一份基督復活的大能.
但更重要的是我要問.
弟兄姊妹 你有沒有經歷過基督復活的大能.
在你生命裡頭去彰顯呢.
復活不單止是一個歷史的事實.
更加是一個當下運行在我們生命裡頭.
我們信仰當中的一個能力.
保羅在三節四節裡這樣說.
「我當日所領受又傳給你們的 第一就是基督照聖經所說.
為我們的罪死了 而且埋葬了.
又照聖經所說 第三天復活了」.
保羅在這裡是用了一種傳統的信經的方式來表達.
是代表著這一段說話不是他自己的一些看法.
一些的創見.
而是教會在早期的時候.
是一個最核心 最不能夠被改動的一個信仰內容.
所以他強調的是什麼.
這個信息他傳給他們是他領受的.
也是他傳給他們的.
表明這個福音是有傳承 是有根基 是有歷史的.
保羅說耶穌死了.
是說明耶穌為我們的罪照聖經所說真實地去承受了死亡的審判.
這個不是一個象徵 而是一個事實.
並且埋葬了.
埋葬是確認耶穌的死亡是真實性的.
是排除了他是一種假死的可能性.
他確實是進入了死亡.
並且耶穌他更加在第三天之後復活.
照聖經所說他復活了.
這個是歷史的轉捩點.
也是我們信仰的根基.
所以如果耶穌沒有復活.
十字架只是一次的失敗.

$^{121}$但正正因為耶穌基督他已經復活了.
十字架就成為我們一個永恆的德性.
所以復活不單止是證明耶穌的身份.
更加是彰顯著基督他復活的大能.
而在這裡保羅並不是要求人.
只是一種麻木的相信.
有時候我們回到教會.
可能我們都是純粹在理性上.
我們知道 我們相信.
保羅他更加給了一些可被查核的歷史證據.
他列出了一連串的真實的見證人.
這些人是曾經親眼見過復活的主.
這份名單不是一些文學的修辭.
而是他用了猶太人一些法庭式的陳述.
就是說耶穌的復活是可以被見證的.
是可以確認的一些歷史的事實.
而那個見證人的名單是這樣的.
裡面有基法 即是彼得.
他是最早的一個個人的見證人.
主親自向他顯現.
還有十二使徒 他們是一起集體見證的.
除非他們是一個集體的幻覺.
所以他們集體都見證得到.
好像耶穌他真的復活了.
還有就是五百多位的弟兄.
其中當時保羅說.
在他寫這本書的時候.
這些人仍然在世.
即是可以怎樣.
找回他們信文可以去核實.
並且當中有雅各.
雅各是誰 雅各是耶穌的弟弟.
原本是不信的.
但後來他都成為了耶路撒冷教會的領袖.
還有就是眾使徒和保羅他本人.
保羅他更加是怎樣.
他從前是迫不及待.
現在變成一個傳福音的人.
成為一個使徒.
這個本身就是一個奇蹟.

$^{161}$但是保羅為什麼在這裡要.
寫出這份名單呢.
保羅在這裡列出這些名的目的.
不單只是說一個信仰的辯護.
更加是要建立一個信心的基礎.
我們的信仰不是建立在一些神話.
或者一些故事裡面.
而是我們建立在一個真實的人.
他所經歷到的一些真真實實.
真是發生過的一些事情上面.
而特別要值得注意的是.
在這份名單裡面有一個人是雅各.
雅各他雖然是耶穌的弟弟.
但是他原本對耶穌的使命是感到很懷疑的.
大家如果有時間的話.
可以回去看《約翰福音》的七章五節.
那裡也有表達.
但是在復活顯現之後.
雅各他徹底被改變.
甚至他願意為信仰去殉道.
這種轉變的本身.
就是一種復活的真實性.
是一個很有力的證據.
所以姐妹,我們怎樣可以告訴別人.
我們的信仰是真的呢?.
我們信的耶穌是真的呢?.
是你生命的改變.
而保羅在這裡的重點.
他不是要去辯論.
耶穌是不是真的已經復活.
他要的是.
我們再往前想一想.
就是一個更加重要.
更加深的問題就是.
究竟耶穌復活對我們每一個人的生命.
究竟有什麼影響.
弟兄姊妹.
耶穌的復活.
究竟對你的生命有什麼影響呢?.
弟兄姊妹,我們都很明白一件事.

$^{201}$人性是怎樣呢?.
我們普遍都是不想改變.
也都不願意改變.
但是很多時候當我們見到一些人.
他真的突然之間性情大變.
可能他在他人生裡面.
可能發生了一些很重大的事.
或者可能他受極大的創傷.
所以他整個性情大變.
又或者有些人.
可能他為什麼生活習慣被迫改變呢?.
很多時候都是健康響起警號的時候.
你看到一些男士都是.
我見到很多男士都是.
只要老婆不在他身邊和他一起吃飯.
他都多數是會叫男人的浪漫.
豆腐火腩飯.
但什麼時候對於男士來說.
他會改變他的生活習慣呢?.
就是當他身體出問題的時候.
由吃東西很不檢點變得很檢點.
由從來不做運動變成要做運動.
我認識一個中學同學.
他本身上學是從來都是睡覺的.
我由中二開始認識他.
中三,中四他都是睡覺.
他睡覺睡到什麼程度呢?.
就是睡到有一次在英文課.
我們嘗試拆了他兩隻鞋的鞋帶.
然後把他綁在一起.
他到醒了才發現.
然後就說髒話罵我們.
但突然之間有一天.
我們開始發現他不同了.
他每天都很專心去溫習.
每天都很專心上課.
整個人上學不是很整齊.
而是真的很整齊.
這種改變還要不是一天兩天.
而是一段很長的時間.

$^{241}$我們見證著這樣的改變.
當時我們很震撼.
為什麼一個每天睡覺的同學.
可以突然之間變得這麼用功呢?.
後來我們發現了一件事.
原來他暗戀我們班裡的一個女同學.
她的生命遇到一件很大的事情.
改變了她的生命.
改變了她的習慣.
所以我們明白.
當時一班的信徒.
甚至是門徒.
他們突然之間.
他們的生命有這麼大的改變.
由耶穌給他們十字架.
他們四散.
他們的生命去到一個.
這麼低谷的環境.
但竟然他們的生命有這麼大的轉變.
他們不單止努力地去傳承福音.
更加為福音殉道.
其實這是確實肯定的一件事.
就是在他們的生命裡頭.
他經歷了很重大的事.
就是他們真是見證了.
耶穌基督他已經已然的復活.
然後保羅他更加說的是.
他不是說一些神學的大道理.
而是他告訴我們.
他自己的生命如何改變.
保羅他說:我原是使徒中最少的.
不被稱為使徒.
因為我從前迫不得神的教會.
然而今日我成了何等人.
是蒙神的恩才成的.
這段經文最震撼.
也最個人的部分.
保羅他不是在談論一些很抽象的神學.
而是他用自己的生命來做一個見證.
保羅他曾經是迫不得基督徒的人.

$^{281}$是反對福音的人.
是與神對立的人.
他更加參與了斯提反的殉道.
他四處去拘捕那些信徒.
但保羅他就是這樣一個人.
在他人生裡經歷了很大的轉變.
他在大馬士革的路上.
遇見了那位復活的主.
一切都改變了.
如果大家還記得菲勒比書說.
保羅他整個人的價值觀都改變了.
他從前以為有益的.
現在都因為信耶穌基督而當作有損的.
因為他以認識耶穌基督為自保.
所以保羅說他今日成了何等的人.
這句話是包含了復活的大能.
去改變他整個的生命.
是最真實的宣告.
保羅他過去的轉化.
他過去在碧北教會.
他形容自己是使徒中最少的.
不配稱為使徒.
但復活的主向他顯現之後.
這個經歷並沒有激怒他.
反而他被神的恩典所轉化.
因為你可以想像得到.
當時保羅他本身是一個碧北基督徒的人.
突然之間轉態.
變成信奉基督的人.
其實當時他面對的環境.
面對的人.
其實可能有很大壓力.
因為很多人都會質疑他的身份.
他所做的.
但保羅並沒有因為這些經歷而被激怒.
反而他因為神的恩典所轉化.
這裡給了我們一個很極大的盼望.
就是無論你過去是怎樣的樣子.
復活的主可以重新定義你.
我們不需要假裝過去的事從未發生過.

$^{321}$但過去的這些事情.
不再是你的身份最終的定義.
聖經裡有很多這樣的見證.
大衛他犯罪和悔改.
讓他不認主到他心靈的復興.
甚至相傳他進入了殉道.
是他倒釘十字架去殉道.
活大拉的瑪利亞.
他心靈的破碎和被主的接納.
所以因為復活的主不是要來評判.
去批判我們每一個破碎的人.
而是要來更新我們和改變我們的生命.
在我們的過去可能我們會充滿失敗或羞愧.
可能有些是你不敢抬頭的記憶.
可能你犯過罪 虧欠了神.
可能你對神也好 對人也好.
有做過一些錯誤的事情.
甚至有過一些遠離神的歲月.
也曾經冷淡過 懷疑過 甚至背叛過.
可能我們也被罪一直捆綁.
我們反覆去掙扎.
但也無法自拔的一些習慣或傷痛.
但今天耶穌告訴我們.
他承受的這份救恩.
是告訴我們我們的過去.
是不能夠限制了神在我們身上的工作.
復活的主可以進入我們的過去.
帶來我們真實的更新.
我認識一個人 一個弟兄.
我記得當年我在一個課程裡認識他.
當時這個弟兄很掙扎.
很掙扎因為他蒙召.
他想要成為一個傳道人.
但他一直為了他的過去放不下.
所以他一直未能夠決定到.
要去神學院讀書 去受裝備.
後來有一次機會我聽他分享.
他跟我說為何他一直不敢去做傳道人.
要去入神學院受裝備.
他就跟我分享.

$^{361}$原來他過去信仰有七年的時間.
他沒有回教會 他離開了教會.
為何呢 因為他說他最初回教會的時候.
其實他都很愛主.
回到教會甚麼事都參與 甚麼都做.
人叫他做的事他都會做.
但當他越參與教會 越多承擔的時候.
開始有些事可能做得不夠好.
有時甚至可能令到人覺得你做得很差等等的情況.
開始聽到很多批評 開始聽到很多埋怨.
他很受傷 他很不開心.
有一天他說他完了崇拜之後.
他甚麼人都不跟他們說話.
然後自己去了海邊大哭.
還罵上帝 他說上帝這麼忠心地去服侍你.
為何要受這麼多傷害 受這麼多批評呢.
他就決定從此不回教會.
所以他說他那七年的時間.
他是過著嫖賭飲當吹甚麼都做的生活.
這樣的生活維持到第七年的時候.
有一天 他說那天是主日.
他下午三點多才起床.
當時他起床的時候.
因為之前那晚喝了很多酒.
還有些酒氣未過.
所以他坐在桌子旁邊喝水.
當時他媽媽走出來跟他說了一番話.
兒子 其實你發生甚麼事.
以前不是好好的回教會.
很開心的 為何現在這樣過日子.
當時這個弟兄醒一醒.
他想起我以前是回教會的.
我以前是信耶穌的.
所以他醒悟過來.
他決定要回教會.
當他回教會的時候.
他都不敢去接觸任何人.
亦都不敢參與任何的事奉.
他很怕自己又跌入事奉的陷阱裡.
所以他選擇坐在一旁.

$^{401}$但當他慢慢回到教會的時候.
他有一次崇拜當中.
他聽到神呼召他成為一個傳道人.
他講完這個見證經歷的時候.
他說他到現在都仍然是.
接受不到自己過去有七年背叛神.
所以他還未決定是否去報讀神學院.
我當時就跟他說了一句話.
我說其實弟兄如果不是有過去的七年.
又怎會有今天的你呢.
所以他之後就立刻去報讀神學院.
並且今天已經成為一位牧師.
在教會中心去服侍.
弟兄姊妹.
耶穌基督復活的大能.
是能夠改變我們的過去.
然後保羅更加在講的是.
復活的大能不單止改變我們過去.
更加改變我們現在.
保羅在講的是.
不單止是在講我們的生命.
我們的罪會被赦免.
他自己也在經歷.
他更加是怎樣呢.
他更加是徹底的被轉化.
他說並且他所賜我的恩不是徒然的.
我被眾使徒格外勞苦.
這個恩不是在講一種很被動的慰藉.
而是在講的是帶出保羅的生命.
是有行動是有改變的力量.
這裡的重點是恩典不是被動的.
而是帶出我們行動和改變的能力.
所以同樣的弟兄姊妹.
我們也經歷了耶穌基督這份救恩.
去改變我們現在的生命.
是帶出行動和改變的力量.
保羅他是怎樣的.
他說他從前是逼迫教會.
他曾經是最大的敵對者.
他四處去拘拿信徒.

$^{441}$當他遇見那位復活的主.
在大馬士革路上面.
他生命的轉捩點.
他改變到他去建立教會.
他成為外邦人的使徒去廣傳福音.
並且他更加最後說什麼.
他說他是為主而活.
這個不是他.
乃是神的恩與他同在.
他是在經歷神的恩典.
他整個改變是神的恩典.
這個轉變不是靠意志力.
也不是靠任何宗教的訓練.
是因為有復活的生命在他裡面去運行.
所以弟兄姊妹我們有時候搞錯了.
我們信耶穌都可能誤會了.
我們看見聖經的吩咐.
然後我們就用我們的意志.
和因為我們認為信仰是導我們向善.
我們要做一些好的事出來.
所以我們靠我們的意志去行道.
我們靠一些所謂的宗教訓練.
去養成一些習慣.
弟兄姊妹我不是說這是錯.
不過當我們看完聖經的教導.
然後我們靠我們的意志力.
靠一些所謂的宗教的規條去訓練我們.
能夠做一個基督徒的話.
其實我們只是回到什麼.
回到靠律法得救.
靠行為得救.
但是靠律法得救.
靠行為得救.
這件事已經告訴我們是怎樣.
是行通.
所以今天我們必須要的是.
當我們聽完道.
頂尖妹妹我們一定會有很多想法.
哎呀很難.
彼此相愛.

$^{481}$很不容易.
你幫我看看你旁邊那個.
你要和他彼此相愛容易嗎.
如果那個是你旁邊的你老公.
可能都不容易的.
是不是.
可能我在想要彼此相愛.
你已經覺得很不容易.
然後我們要十一奉獻.
我們又覺得很困難.
我們要為福音去受苦.
我們更加哇.
要付這麼大的代價.
我們怎麼行.
因為我們都停留在什麼地步.
我們以為靠意志.
靠宗教訓練.
頂尖妹妹.
今天保羅在這裡就要告訴我們.
我們必須要的是什麼.
我們要經歷復活的大能.
我們的生命才能夠去改變.
所以在這裡我們都需要很誠實的去問我們自己.
當我們信了耶穌.
我們是不是能夠更加去愛人呢.
我們對我們身邊的人.
是不是真的能夠有更多的耐性呢.
有更多的憐憫和接納呢.
我們信了主之後.
我們是不是更加有盼望呢.
面對困難的時候.
是不是仍然能夠仰望主的信實呢.
我們信了主之後.
我們是不是更加有力量呢.
在每一次的試探.
我們在受人的真摯當中.
我們是不是都感受到聖靈的同在和拯救呢.
頂尖妹妹.
如果.
當我們發現.

$^{521}$改變在我們生命中.
是不明顯的話.
很可能.
很可能我們過去的信仰.
只是相信一種觀念.
一個觀念.
而不是真正經歷了.
基督復活的大能.
頂尖妹妹.
信仰不是單單說一些.
知識的累積.
而是說.
生命的更新.
而復活的大能.
是能夠幫助到我們.
使我們活出一個新的身份.
保羅說.
然而這不是我.
乃是神的恩與我同在.
這句話就帶出整段的經文.
最深刻的真理是什麼呢.
保羅他做了很多事.
他付出了很多勞苦.
但是他的結論是什麼呢.
這不是我.
乃是神的恩與我同在.
真正的改變不是.
我有變得多好.
而是基督在我裡面活著.
頂尖妹妹.
真正的改變.
不是因為我們回來教會.
我們做得多好看.
我們做了多少事.
我們多有能力.
甚至可能被人稱讚.
真正的改變.
是什麼呢.
是基督在我們裡面活著.
是改變了你整個的生命.

$^{561}$你面對順境的時候.
你處於一個享盡歡呼聲的時候.
你的生命你仍然記得.
是因為基督在你裡面活著.
甚至反過來.
在你逆境的時候.
在你被受攻擊.
被受傷害的時候.
你仍然記得.
基督在你裡面活著.
所以復活的大能.
使我們可以活出一個新的身份.
不是靠宗教的努力.
所以復活的生命.
不是靠遵守一些規條.
也不是靠更加努力祈禱.
讀經所累積的功德.
我們很容易就這樣去教人.
你沒信心嗎.
祈多些禱吧.
你沒力量嗎.
看多些聖經吧.
這些有沒有錯.
沒有錯.
但是我們更加需要的是.
我們生命經歷這位復活的主.
他復活的大能.
去被他改變.
另外更加不是一個道德的提升.
福音不是讓人成為一個更好的人.
而是讓人成為一個新造的人.
是一個本質上的更新.
而非行為上的修正.
並且是我們要生命的更新.
就是復活主要進入一個人的生命.
帶來的是由內而外的一種轉化.
我們有一顆新的眼光.
新的愛.
新的盼望.
新的力量.

$^{601}$弟兄姊妹.
我們的生命最大的困難.
我們最不能改變的是因為什麼.
是因為我們帶著舊的眼光.
舊的愛.
舊的盼望.
舊的力量.
回到教會內的人.
我們傳統形容的那兩個字是什麼.
就是老餓.
我們仍然在用舊的眼光.
舊的愛.
舊的盼望.
舊的力量.
所以保羅在加拉太書2章20節那裡說.
我已經與基督同釘十字架.
現在活著的不再是我.
乃是基督在我裡面活著.
這裡和哥林多前書15章的主題相呼應.
是在說一個新的身份.
不是自我的實現.
而是在基督裡面我們能夠得到重生.
一個新的生命.
所以弟兄姊妹.
我們要讓復活的大能進入你的生命.
這段經文不單止要讓我們知道一件歷史.
而是邀請我們進入復活的生命.
所以我們要接受這個福音.
不只是純粹在腦裡面理性的理解.
保羅說.
你們若不是徒然相信能以持守所傳給你們的.
就必因這福音得救.
信仰不只是「知」字.
不只是在說「知道」.
而是我們要持守.
知道耶穌復活是知識.
而我們更重要的是什麼.
將我們的生命交給主.
這才是信仰.
所以今天弟兄姊妹.

$^{641}$你是不是真的將你的過去.
你的現在和你的未來.
都交在主的手裡.
還是停留在一個頭腦的認同.
只是心裡面仍然由自己去掌管自己的生命呢.
按我們的喜好.
按我們的經驗.
按我們自己認為最好的去過我們自己的人生.
你要讓復活的大能去改變你的內在.
復活的大能不是先改變你的環境.
而是改變你這個人.
復活的大能帶來的是從絕望到有盼望.
即使在黑暗的低谷裡.
我們仍然能夠看見主的信實.
從懼怕到勇敢.
不再被恐懼和焦慮所限制.
從罪的捆綁到得到自由.
那些我們經常在反反覆覆.
在掙扎徘徊的習慣.
我們能夠得到釋放.
所以我們不是只求環境改變.
而是求神先去改變我們自己.
最後我們要活出一個被改變的生命.
保羅說他格外的勞苦.
不是為了證明自己.
也不是靠他的意志力去撐著.
而是因為他被神的恩典去激勵.
一個真正經歷了復活的大能的人.
他的生命必然不一樣.
這種不一樣不是表現給人看的.
而是從內而外的.
是自然流露的改變.
所以我們在自己的家庭裡.
在職場裡 在教會裡.
我們身邊的人.
是否能夠從我們身上.
看到這份復活的見證呢.
你身邊的人知道你是基督徒.
會說原來你是基督徒.
還是會跟你說你是基督徒.

$^{681}$這就是我們能夠活出一個被改變的生命.
和你分享一個見證.
這個人的生命見證.
這次和你分享是有少許不同.
這是我過去三年多四年的時間.
是見證著這位弟兄的生命如何改變.
我們叫他做高佬.
高佬本身是住在屯門大.
他很小的時候就進入了黑社會.
他形容自己前半生只有拳頭和義氣.
那個不是忠義.
我的語音輸入是忠義.
中學二年級的時候.
他遇到一個好老師.
雖然遇到好老師之外.
但他中三的時候就輟學.
就算有多好的老師也沒有用.
之後他就進入了黑社會跟大佬.
首先去到九龍.
在舞廳裡工作.
接著有Disco.
現在成為了油尖旺.
深水Po 的地頭蟲.
那個時候他的生活就是.
無日無夜的喝酒.
尋歡作樂.
手臂上有刀疤.
因為黑社會當然是要劈友.
他形容光回險學.
財政也有很多債務在他前身.
也反映著他的心靈是很空虛.
很空虛的一個寫照.
在起伏的人生裡.
他有一個很愛錫他的爸爸.
他形容父親的愛.
就好像一座沉默的高山.
為什麼這樣說呢.
因為他爸爸通常不會說太多話.
但他爸爸的愛.
全部都盡在一起喝酒.

$^{721}$在樓下吃個早餐.
而最深情的說話就是.
每一次高玉出事的時候.
爸爸跟他說.
家裡的大門隨時為你而開.
而在他人生裡.
他最大的轉捩點就是.
他從一位弟兄帶他回教會.
起初他仍然有很大的掙扎.
他經常都怕自己保持不住.
但他聽神的教導.
他亦感受到弟兄姊妹的愛.
慢慢這份福音的大能.
都紮根在他心裡.
之後他上了傳描班.
認識了很多弟兄姊妹.
他生活變得很充實.
他還要很勇敢的.
去了很多福音座會去分享見證.
而且他的生命也開始改變.
他由以前經常流連風月的生活.
現在已經變得不再喜歡那些生活.
他又開始顧及他的家人的感受.
堅決要去活得像主.
他經常都跟我分享一句說話.
就是我不會讓這些人.
讓這些事令到他的老我又走出來.
他漸漸地在教會成為了他的像.
他的剛扶持了他.
雖然他成長的過程裡.
他經歷了他的好朋友離世.
我記得當晚我和弟兄一起.
去志佑的家門爆門.
進去之後就躺在床上離世.
之後他爸爸的安息.
我更加記得在他爸爸安息那天.
我趕去大埔的醫院裡.
因為我知道醫院的醫生.
一直處理得很差.
我本來打算去看著弟兄.

$^{761}$讓他不要說髒話罵他或打他.
結果看完醫生的時候.
我的感覺是.
我很想打爆這個醫生.
但這個弟兄竟然一句髒話都沒有說.
甚至沒有任何的行動.
之後他更加經歷了和他的兒子有決裂的時候.
但是他心裡仍然很感受到的神養.
他知道的是.
他今天活著的全部都是恩典.
他每一天都活在這樣的裡面.
所以這個弟兄他今天在教會裡.
很努力去事奉.
很積極的去傳福音 去派單張.
他更加繼續去帶他身邊的朋友.
特別是以前在黑社會裡的朋友.
回教會.
當這些朋友回教會的時候.
回得不好.
他去督促他們.
甚至載他們回教會.
弟兄姊妹.
這個是復活的大能.
讓一個人生命的改變.
所以普羅他會這樣形容.
我今天成了何等人.
是蒙神的恩才成.
弟兄姊妹.
我們今天成為何等人.
都是蒙神的什麼.
恩才成.
基督的復活不只是2000年前的一個神蹟.
基督復活的大能在今天仍然運行著.
這種能力不會因為時間而消逝.
也不會因為時代而改變.
基督的大能在今天直到永遠.
一直神要運行在我們中間.
讓我們可以去經歷.
讓我們可以去活出一個新的身份.
可以改變我們的過去.

$^{801}$更新你的現在.
重塑你的未來.
我心願這一句說話.
都成為我們每一個弟兄姊妹.
心裡面.
謹記的一句說話.
一個見證.
就是今天我們成為何等人.
都是蒙神的恩才成.
當我們經歷.
基督復活的大禮.
我們的生命可以不一樣.
弟兄姊妹.
你願不願意.
你的生命.
因為我們回來復活節的主日.
我們不只是慶賀主耶穌的復活.
完成了一個聚會.
而是我們真的能夠經歷了.
基督復活的大能.
我們的生命去改變.
去更新.
以致我們身邊的人.
能夠從我們的身上.
能夠知道.
耶穌基督的復活是真實的.
是發生過的.
是能夠改變人生命的.
弟兄姊妹.
你願不願意.
由我們在同低頭有一段禱告的時間.
差主耶穌.
為我們的罪.
釘身在十字架上面.
你更加在三日之後復活.
今天我們.
因著你的恩典我們能夠回來教會.
我們可以去敬拜你.
可以去讚美你.
是我們都經歷了何等大的恩.

$^{841}$主我們要獻上感恩和讚美.
主我們今天來到你的面前.
我們懇求你.
懇求你讓你基督復活的大能.
都充滿我們生命裡頭.
求你去根深去改變我們的生命.
拿走我們的老我.
拿走很多纏繞我們的罪.
拿走很多我們不願意改變的壞習慣.
主以至我們的生命來到你的面前.
我們單單要學著你.
讓我們的生命真的能夠可以.
讓人看得見耶穌基督.
你真的復活.
你復活的大能在我們生命.
你改變了我們的人生.
改變了我們的價值觀.
主就願意你親自去祝福我們.
幫助我們.
主說在我們過去的日子.
我們在教會裡面.
我們帶著很多我們的老我.
我們和弟兄姊妹去相處.
我們有很多自己的執著.
很多的堅持.
有時候我們可能都傷害了對方.
主我們懇求你.
基督的寶血充滿在我們中間.
讓你基督的寶血.
去化解一切人的冤仇.
讓你基督的寶血.
去洗淨我們一切的罪.
以致我們來到你的面前.
我們能夠可以坦然無懼.
來到你的私人寶座前.
能夠蒙恩惠得拯救.
主我們就懇求你親自去幫助.
祝福我們.
不單祝福我們每一個弟兄姊妹.
更加祝福我們整個的教會.

$^{881}$讓我們這間教會.
能夠可以彰顯你基督復活的大能.
是因為我們每一個的生命.
都在你們面前.
能夠蒙恩得著改變.
主我們就多謝讚美你.
求你繼續祝福我們.
我們感激老交託.
是奉靠主耶穌基督的聖名自祈求.
阿們.
\newpage



\section{哥林多前書 15:50-58}
\label{sec:031916}
\textbf{叫人全然改變的復活主}~~~~黃紹權牧師
\newline
\newline
連結: \href{http://www.church.com.hk/acms/content.asp?site=cdc\&op=show\&type=product\&code=031916\&layout=sermon}{\texttt{031916}} ~~~~ 語音日期: 2026-04-05
\newline
\newline
\hyperref[sec:031915]{\small{< < < PREV SERMON < < <}}
~
\hyperref[sec:index_chronic]{\small{[返順時目]}}
~
\hyperref[sec:031917]{\small{> > > NEXT SERMON > > >}}
\newline
\newline
哥林多前書 15:50-58
\newline
\begin{longtable}{cl}
\hline
\hline
章節 & 經文 (和合本修訂版)\\
\hline
15:50 & \begin{tabularx}{0.7\textwidth}{X} 弟兄們，我要告訴你們的是：血肉之軀不能承受神的國，必朽壞的也不能承受不朽壞的。 \end{tabularx} \\ \\ \relax
15:51 & \begin{tabularx}{0.7\textwidth}{X} 我如今把一件奧祕的事告訴你們：我們不是都要睡覺，而是都要改變， \end{tabularx} \\ \\ \relax
15:52 & \begin{tabularx}{0.7\textwidth}{X} 就在一剎那，眨眼之間，號筒末次吹響的時候。因號筒要吹響，死人要復活成為不朽壞的，我們也要改變。 \end{tabularx} \\ \\ \relax
15:53 & \begin{tabularx}{0.7\textwidth}{X} 這會朽壞的必須變成不朽壞的；這會死的總要變成不會死的。 \end{tabularx} \\ \\ \relax
15:54 & \begin{tabularx}{0.7\textwidth}{X} 當這會朽壞的變成不朽壞的，這會死的變成不會死的，那時經上所記「死亡已被勝利吞滅了」的話就應驗了。 \end{tabularx} \\ \\ \relax
15:55 & \begin{tabularx}{0.7\textwidth}{X} 「死亡啊！你得勝的權勢在哪裡？死亡啊！你的毒刺在哪裡？」 \end{tabularx} \\ \\ \relax
15:56 & \begin{tabularx}{0.7\textwidth}{X} 死亡的毒刺就是罪，罪的權勢就是律法。 \end{tabularx} \\ \\ \relax
15:57 & \begin{tabularx}{0.7\textwidth}{X} 感謝神，他使我們藉著我們的主耶穌基督得勝。 \end{tabularx} \\ \\ \relax
15:58 & \begin{tabularx}{0.7\textwidth}{X} 所以，我親愛的弟兄們，你們務要堅固，不可動搖，常常竭力多做主工，因為你們知道，你們在主裡的勞苦不是徒然的。 \end{tabularx} \\ \\
[1ex]
\hline
\hline
\end{longtable}
$^{1}$菩提神保佑我們在聆聽主的說話以外.
一起同心低頭祈禱.
父神多謝你讓我們能夠在這個復活節的日子.
一同去思想.
並且歌頌讚美我們在這個德性的位份上.
能夠站穩著腳.
因著你的愛子耶穌基督所奠定的基礎.
讓我們在今天活在這個地上.
從來不擔憂我們會死在罪惡禍犯的當中.
但卻是讓我們向罪惡獻上我們的力量去掙扎下去.
讓我們不會陷在罪惡的當中.
藉此見證我們屬神你的子民.
求你的幫助帶領.
無論是先上還是先下的每一個.
我們都祈求主你的靈與我們同在.
誘導我們去明白你的啟示.
和上主你拯救的那份意願.
觸動我們每一個來更毀身在你當中.
禱告奉靠主耶穌基督德性的名字.
阿們.
天主教的教徒在復活節或者受難節這段時間過節.
很多基督教教徒又在這個受難節和復活節的時候又過節.
甚至連猶太人都在這個時間.
同時過一個他們稱之為愉悅節的節期.
個個都在慶祝.
個個都有一些宗教活動.
但是似乎他們每一個的主子卻有點不同.
譬如天主教會在這段時間很多時候都放在耶穌基督受苦的現象.
所以你會見到不少的國家以天主教為背景所建立.
他們都會在這段時間做出一些譬如抬十字架在街.
或者找一些信徒扮耶穌然後被人真的去鞭傷.
天主教很主動去做這些事.
為何他要做這些事呢?.
因為畫面夠震撼性.
越震撼你就覺得你的宗教感就會同時被挑旺起來.
所以宗教在這些層面上其實是在操弄人的感情和感性.
當你以為你的感情被挑動,靈性好像已經更加貼近上帝的時候.
其實你反而是更遠離開上帝原本在這個節期本身復活節所做的目的.
猶太人又如何呢?.
猶太人就將整個的焦點放在從迫害以至於解放這個焦點的上面.

$^{41}$所以猶太人很多時候在這個逾越節的日子.
同時都去介紹他們如何脫離捆綁而得到上帝的應許.
這種只是尋求應許解脫的思維.
也很多時候流傳在基督教會之間.
我講到這裡不妨我們也講一下基督教會.
基督教會就更加變本加厲.
進入一種叫武枉管,胡作非為的狀態.
環顧在過去這幾十年的裡面.
我看見教會絕多數在復活節或者壽苦節,節期這段時間.
就成為他們大搞福音演唱會的機會了.
很多詩歌分享,默想會,或者有些演唱會做出來的話劇.
名義上好像說是一種默想或者甚至叫做報道的機會.
但是其實只是包裝著教會根本也拿捏不了復活節本身的核心.
講到這裡究竟聖經裡面在這個每年三月與四月之間的時段.
究竟應該要向一個什麼焦點而去進發或者去再重新提醒呢?.
這個就是今天我想和大家一起去看這段.
在聖經當中極難解釋的其中一段經文.
也是今天我們這個復活節的主日崇拜裡面的講道.
就是哥林多前書第十五章第五十至五十八節.
我給了一個題目它叫做「驕人傳言改變的復活主」.
我們的焦點其實要學習認識這位的復活主.
怎樣改變我們然後下一步我們怎樣去做.
讓我們再次回到這段聖經的內文.
更加清楚經文想講些什麼.
然後重拾復活節的主旨是些什麼.
請大家一起打開哥林多前書第十五章.
從第五十至到第五十八節.
請大家打開之後我們一起讀.
好,我們一起讀.
弟兄們我告訴你們說.
血肉之體不能承受神的國.
必扭壞的不能承受不扭壞的.
我如今把一件奧秘的事告訴你們.
我們不是都要睡覺.
乃是都要改變.
就在一剎時.
斬眼之間號筒末次吹響的時候.
因號筒要響.
死人要復活成為不扭壞的.
我們也要改變.

$^{81}$這必扭壞的總要變成不扭壞的.
這必死的總要變成不死的.
這必扭壞的既變成不扭壞的.
這必死的既變成不死的.
那時經上所記.
死被德性吞滅的話就應驗了.
死呀,你德性的權勢在哪裡.
死呀,你的毒鉤在哪裡.
死的毒鉤就是罪.
罪的權勢就是律法.
所以我親愛的弟兄們.
你們務要堅固.
不可搖動.
常常竭力多做主公.
因為知道你們的奴婦在主裡面不是徒然的.
在過去我們看過很多很多的經文.
談論到有關於生與死或者復活等等的問題.
包括《但義理書》我們會去看的.
但是來到哥林多前書.
保羅去談論到有關於生與死與及復活.
似乎是將所有整幅零零碎碎的圖畫.
合而為一讓我們看出那個意義.
這份聖經講什麼呢?前文.
前文是談論到有關於復活之後.
究竟有一個什麼樣的身體.
為什麼會有一個這樣的討論呢?.
其實涉及於教會.
在特別哥林多教會當時.
有一些人.
可能一派或者很多派的人士.
他們將耶穌基督復活.
全息.
成為好像.
猶如今天所說的一種罪得釋放的狀態.
但是這個罪的釋放.
他就更進一步去指到其實是從肉體脫離.
人可以透著信了耶穌之後.
再不受這個罪的肉體所捆綁著.
因為他認為犯罪是因為肉體.
因為人的情慾.

$^{121}$人身體有需要.
而信了耶穌之後.
就能夠讓人脫離這個身體.
靈與身體分開.
構成了一種很早期.
我們所看的類似.
這個諾斯底主義的哲學思想.
流傳在教會裡面.
保羅收到哥林多教會的書信之後.
發覺他們有這樣的思想.
而不合耶穌基督這個死與復活的主題.
因此來到這個部份.
他強調有關於復活的意義.
耶穌的死是承擔罪.
但耶穌的復活卻是一場聖像的宣言.
當耶穌的復活構成了德性的結果之後.
就能夠涵括所有同樣相信耶穌基督的人.
可以進入到這個德性的位份.
所以當我們在這段聖經裡面.
看到保羅的介紹的時候.
他就要讓我們重新回到.
耶穌基督的生命帶來的改變.
尤其是祂為我們每個人帶來的改變.
雖然身體死了.
但卻可以連身體都同時復活了.
在這個我們稱之為一個雙重的條件之下.
信徒同樣相信耶穌基督的.
同樣好像耶穌基督.
不單止靈與上帝再次結連起來.
並且祂的身體同樣會復活.
而這個復活卻是讓我們有一個新的意義.
帶來活在這個地上.
所以今天我們看這段聖經.
其實基本上保羅是要宣告給我們知道.
每一個相信耶穌基督的人.
他們的人的生命應該有同樣.
好像基督的形象一樣.
基督的形象是包括由塵土所構成的人的身體.
以及屬天的人應有的靈性的生命的合體.
怎樣可以做到呢?.

$^{161}$所以這段聖經是很難解的.
這段是不容易解的聖經.
大部份的大英姐妹讀完以為自己明白.
但其實這段聖經裡面有一個很重要的訊息.
我們逐步一步去看.
第一部份.
從第五十節至第五十三節.
在這裡我們看見.
原來復活這件事是為所有人帶來人生的轉折點.
首先我要大家明白一點.
復活不是只有基督徒才會復活.
聽起來可能會很詫異.
但如果我們再次重新認真去看聖經裡面的介紹.
復活並非只有基督徒才有的專利.
我們一起打開但爾里書第十二章第一節看看好不好.
剛剛我們才說完但爾里書沒多久.
但爾里書就說了這一點.
我們一起打開但爾里書.
十二章.
找到沒有?十二章第一和第二節.
我們一起讀.
這裡怎樣說?.
那時保佑你本國之民的天使長米加勒必站起來.
並且有大艱難從上國以來直到此時沒有這樣的.
你本國的民眾凡名錄在冊上的必得拯救.
誰在塵埃中的必有多人復蘇.
其中有永生的有受羞辱永遠被憎惡的.
看不看到?.
不是只有基督徒信上帝.
我寫了我的名字在生命冊上.
我於是乎才有機會復活.
不是.
所有人都復活.
不過基督徒相信主的人的復活.
他們多一重的事情會發生在他上面.
從這段聖經當中.
原來他們會進行一連串的改變.
而這些改變是在他的身體上進行了改變.
我們看看第五十至第五十三節說什麼.
「弟兄們,我告訴你們說,血肉之體不能承受神的國,必扭壞的不能承受不扭壞的.」.

$^{201}$這是道出了一個很基本的承受條件.
要承受一樣東西你必須要有一個條件.
如果你要承受神的國你就必須要有神的國的條件.
你才可以承受神國子民的身份.
甚至神國子民之後的生活.
聖經不停在第五十節說.
「我如今擺一件奧秘的事告訴你.」.
保羅將一件奧秘的事.
奧秘是指一些什麼東西.
平日不會很張揚傳開的.
同樣的字眼用在《達爾利書》裡面.
也都說到一些東西是包裝的.
封嵌了的.
要在某個時間才會告訴你的.
這個奧秘.
要有條件才可以拿到這個奧秘.
這個奧秘是什麼?.
他說「我們不是都要睡覺,乃是要改變.」.
第一個聖經的難題.
在解釋這段聖經裡面.
單單這一句話就有四種不同的翻譯.
為了免卻大家的混亂.
所以我沒有將那四種的翻譯.
去將它詳細層列出來.
但是如果你們有興趣去看的話.
我會鼓勵你們去看一本聖經的譯本.
叫《New English Translation》.
是特別為了這個地方作了一個註腳.
將四種不同的解釋.
都同樣講了給我們知道.
說什麼呢?.
「你們不是都要睡覺,乃是要改變.」.
「睡覺」首先這個字眼在聖經裡面.
是指死亡.
保羅在這裡所說的.
「我們不是都要全部死亡,乃是要改變.」.
前文我們才講到那個所謂承受的條件.
我們本來的狀態是什麼?.
我們是一個血肉肢體來的.
我們是一個會扭壞的身體來的.

$^{241}$為什麼?在《創世記》已經講了.
因為人的犯罪.
所以人的身體已經經歷死亡.
所以保羅在這裡針對死亡.
「你不是要死」.
《創世記》才講你要死.
因為人始早犯罪.
影響所有人都是要在死亡的裡面.
但是為什麼這裡說你不需要死?.
是不是真的死?.
其實這裡所講到的意思是什麼呢?.
就是講給我們知道.
我們作為信徒的.
我們的死.
肉身的死.
有一個改變的過程.
是其他不相信主的人沒有辦法去經歷的.
所以他說「乃是要改變」.
如果你的身體是死了.
好像《創世記》那種死.
現在透過相信耶穌的這個死亡.
會使他的身體改變過來.
這個是唯獨信主的人才會經歷的事情.
如果你沒有這個改變的話.
你是根本上沒有辦法.
進入到去承受神的國當中.
你說是不是?牧師.
是不是不變就不承受得了?.
我給你一個例子.
有沒有試過煮東西吃?.
煮東西吃你通常有些東西要放油的.
那你是不是燒紅了鍋?.
然後才攤上微微沒那麼熱.
就灑些熱油,灑些冷油下去,是不是?.
但是如果在這個時候.
你倒些水下去會怎樣?.
嘩!彈到七彩.
如果我在家裡我一定叫救命.
因為要做很久,要擦很久.
水和油混在一起.

$^{281}$同樣道理你不用用熱油.
用冷油倒進冷水裡.
我無論你怎樣攪都好.
它都不會混在一起.
情況就很類似.
當一個人如果你和神的生命沒有條件的時候.
你是融在一起的.
用廣東人一句話叫「話不投機半句多」.
如果你和上帝的心態,屬靈的狀態.
沒有同樣的條件.
我就算讓你進到天國裡面.
你和上帝都沒話說.
所以保羅在這裡強調一件事.
我們每一個人雖然肉身好像會死.
但是原來這個死不是帶去.
最終和上帝脫離的那種死亡.
而是會有改變.
改變成為配合,符合上帝天國的狀態.
這個稱之為信主.
但是這個信不是可以拖延的.
再看第52節.
「就在一煞時,斬眼之間,號筒末次吹響的時候,因號筒要響,死人要復活,成為不朽為的,我們也要改變.」.
還沒停,怎樣改變?.
「這必朽為的,總要變成不朽為的,這必死的,總變成不死的.」.
原來這裡說給我們知道.
有一連串的改變.
因著你相信即使你死了那一刻.
然後有一個突然的改變.
這個改變是「一煞時,斬眼之間」.
什麼意思呢?.
就是你等都沒得等.
今天很多人當你傳福音.
甚至基督徒也好.
或者教徒.
基督教的教徒.
我要強調.
你跟他說信主的事情.
或者追求認識真理的事情.
他們不是很理睬你的.
你就會帶著愛心的跟他說.

$^{321}$弟兄或者姐妹.
你都信了主了.
你自己都信了主,受了洗了.
為什麼不見你回來查經.
又不回來祈禱會.
弟兄姊妹去傾談聖經的時候.
學習,分享.
在屬靈生命上經歷.
你完全不想參加,為什麼呢?.
他通常都會這樣回答你.
他就會說.
等下先啦.
到時先算啦.
現在很忙啊.
你如果叫人信主也是的.
你會跟他說你信主啦.
很多人會這樣推你的.
我都知道你的主好的.
我都知道我是罪人的.
等下先啦,等多一會兒先.
到時看下怎樣先算啦.
我告訴你.
不是我告訴你,應該是聖經告訴你.
你不用等.
沒得你等.
你是信就信,你不信就不要信.
你不要在這裡等,不要在這裡拖.
你是信了你就認真去信,你不要在這裡等.
到時看下怎樣先啦.
因為我告訴你,到時沒機會讓你等,三時之間.
眨眼之時.
就出現了,出現什麼呢?.
出現號統末次吹響的時候.
這句聖經是什麼來的?.
圈住它,是你的就.
第一個介紹到這段聖經.
講到德性的問題.
整段聖經裡面有三次提及.
德性的問題.
是嗎?所以我不見這個字.

$^{361}$號統末次吹響的意思.
每逢打一場仗在古近東的文化.
即或在今時今日都是.
打勝仗的時候.
勝利一方就會吹響.
勝利的號角.
場仗未打完,不過你贏定了.
勝負已奠定.
你就會吹響號角.
譬如.
你將他的領袖.
將他已經.
殲滅.
就會吹起.
末次的響號.
原來德性那一刻是.
突然間來到的.
耶穌基督.
所造成的這個生命的改變.
讓你可以有個改變的身體符合上帝的.
是突然間就來了.
不能等,到時再決定.
不是這樣.
因為這個響號.
看清楚這句說話.
因這個號.
響起的時候.
死人就復活成為不留懷的.
這裡所指的是哪些死人.
是指的是哥林多前書裡面那些受眾,即基督徒.
你們會明白到,在這個情況之下,你們要復活就會進入一個.
改變的過程.
這個改變的過程是由一個.
不留懷的身體.
去承擔.
所以保羅說.
我們也要改變.
不單止不留懷,而且是不會再承受死亡所帶來的捆綁.
為什麼信了主和不信主的人都同時要去復活呢?.
因為還有一個末後的審判.

$^{401}$審判是沒有negotiation的.
只不過是決定而已.
將你的結果陳明.
然後指出證據.
以前我做.
核數的時候.
有一樣東西.
我老師教的.
他說核數核到最後的時候,我們要寫一個報告.
在這個報告裡面.
你只是告訴他知道結果.
還有原因就可以了.
如果你在銀行工作,或者在這個社會工作,你們都會經歷過一個字眼.
Findings.
找到一些東西.
我基於這些東西就走來給你一個決定.
到時不會再有任何的.
辯駁的機會.
你說不是啊,我是不信啊,我是要辯駁,你可以辯駁.
不過在屬靈的世界裡面.
耶穌的審判就是最終審判.
不過從這段聖經裡面給我們看到.
原來當一個信了主的人.
看起來好像和其他人都同樣死在他自己的年日.
因為衰老的原因.
但是原來他們的死.
卻引發到下一個改變.
因此.
基督徒.
是不怕死的.
頂多你說牧師,我怕死了之後有不捨得我家人的狀態.
我告訴你.
我也會有.
但是我們不會擔心死亡.
因為當我們死亡的時候,就是身體改變的時候.
問題就是你是不是真的信了.
信了就可以改變,不信了你就仍然存在同樣的身體裡面.
你都會復活,不過你復活之後.
是沒有改變的.
最後.

$^{441}$上帝給你一個舊的身體.
讓你在懲罰當中去承受.
所以當我們看回這個描述,其實保羅是告訴我們知道.
原來我們可以被納入到一個勝利的局面裡面.
回到我們今次這個復活節主日的講題.
可以來到每年這個三月和四月之間的時間.
教會不是來搞音樂演唱會,或者什麼欣賞會,默想會.
不是搞這些東西.
而是因為我們可以有機會納入到耶穌基督這個.
勝局的裡面.
他使我們能夠有復活的機會.
賦予我們一個新的身體.
我們應該是在這個復活節的時候,我們帶著一個能參與在這個勝局.
裡面一份的感恩.
這個之後第57節裡面會更詳細去講的.
但是,如果我們不是因為這樣而去參與在這個復活節的子期.
而是只是想自己覺得有宗教感強一點.
覺得自己今年好像屬靈多一點.
甚至搞得好像天主教會那樣搬十字架給邊廂,甚至拿十字架出來打釘.
各樣都好.
對不起,大家知道嗎?你完全off the focus,你完全是脫離了復活所帶來進入勝局的主旨.
我們是進入到勝局的裡面.
不是還在糾纏在釘十字架有多痛苦這麼簡單.
因為祂為我承擔了一些我們受不了的事情.
但是,我們已經進入了去到這個勝局當中.
我們不用再受這些釘傷了.
你再糾纏在,糾結在受釘傷的狀態之下.
換句話說,你沒有脫離到進入復活的階段.
今天很多教會因為沒有了這個焦點.
所以將教會包裝成為一個只是用來使你有宗教感的一個場所.
頂籠想演多兩齣宗教的舞台劇給你看.
有沒有帶你進入一個復活的狀態,怎樣向這個世界一起去同主一起面對呢?.
第二個部分,第54節至到第55節.
因為我們在上文知道我們已經進入了這個勝局當中.
所以當我們是勝過罪惡的時候.
我們發出歡呼之聲.
我們看看這幾節講些什麼,第54和55節.
我讀給大家聽,你慢慢去欣賞了,好不好?.
第54節這樣說,「這必留壞的既變成不留壞的,這必死的既變成不死的.
那是經上所記,死被得誠吞滅的話就應驗了」.

$^{481}$55節,死呀!你得星的權柄在哪裡?.
死呀!你的獨勾在哪裡?.
如果你是用的聖經是一些比較新的版本.
你就會留意到,尤其是第55節.
是特意用了一個詩歌題材這樣分了一句出來去記載的.
你說為什麼翻譯的人要這樣做?.
不覺意呀?.
不是.
因為這一句是一個歡呼聲.
歡呼聲.
有沒有聽過以前的人唱山歌?.
唱山歌是怎樣唱的?.
我就是看電視,我沒有現場聽過.
原來以前的農村有兩個山頭.
唱山歌,一邊有些女士,一邊有些男士.
遠遠的男士看到遠遠有美女,追求她.
然後唱幾句,唱幾句好聽的最好的男高音給她聽.
接著,當遠處的女士聽見這麼漂亮的男聲.
她又回應,有一個美麗的女聲回應那邊的男士.
聖經,當勝利的時候就是做這個歡呼.
怎樣呀?.
當你和我都變成不扭懷不死之後.
我們在這個勝局裡面贏了.
接著怎樣呀?.
就被上在聖經上記載.
死被得勝吞滅,第二次談論到有關於勝利的問題.
而且這次很直接講到.
死對於我們來說已經沒有任何影響.
於是乎有個回應.
回應就是死呀,你得勝的權勢在哪裡?.
這個說話是什麼?聽起來就好像很文雅.
其實這是很囂張的說話.
反過來跟死亡說,死呀,死先生,死代表魔鬼,代表罪.
死呀,罪呀,魔鬼呀,炸彈呀.
你贏我在哪裡,憑什麼權並贏呀?.
再說一句,死呀,你的毒鉤在哪裡呀?.
這兩節的經文是從哪裡出來的呢?.
分別是在以賽亞書第25章第8節.
那裡記載到死是被得勝吞滅.
不信我們看看好嗎?.

$^{521}$以賽亞書第25章第8節,我們有片去看看好嗎?.
我們現在,說真的,很多時候回教會,很多人聖經都不帶的了.
現在只帶電話,不過電話這麼方便,他都不知道為什麼,都不想聽的.
他都不想讀一下它,讀劇集他就會了.
好,我們看看,以賽亞書第25章第8節,我們一起讀好嗎?.
他已經吞滅死亡,直到永遠,主耶華必擦去各人臉上的眼淚.
又除掉普天下他百姓的羞辱,因為這是耶和華說的.
然後耶華在這裡說得很清楚,告訴他們死被永遠吞沒.
在上帝裡面,他不用怕怕死,因為死被上帝降服了.
而所有相信祂的人,就不再在這個受死的裡面得到悲傷.
上帝安慰他,所以弟兄姊妹,如果你有家人,他信了主死.
沒錯,你會為到他離開而有一些依依不捨.
但你應該同時為此而喜樂,因為他在上帝裡面得到最大的喜樂和安慰.
不過問題是你將來是不是也會和他同在呢?.
不單止,在這裡,以賽亞將這個信息說了出來.
山歌裡面的男士高音,馬上給我們看看另一邊的回應.
就在荷西亞書第十三章十四節.
我們打開荷西亞書看看,小先知書,荷西亞,荷爾摩.
十二卷的第一個部分,第十三章第十四節.
我們找到之後一起讀,好不好?.
好,我們一起讀.
我必救贖他們脫離陰間,救贖他們脫離死亡.
死呀,你的災害在哪裡呢?陰間呀,你的毀滅在哪裡呢?.
在我眼前,絕無後悔之事.
看到嗎?聖經自己也在唱頌.
保羅在這裡寫給哥林多教會的信徒知道.
這一件是大快人心,令人振奮的事情.
每一個哥林多的信徒,每一個看到哥林多前書來到第十五章的弟兄姊妹.
都應該為此而歡呼讚頌.
就好像詩篇裡面很多時候那些「Hallelujah」.
大家拜帝威,知不知道這個插句口的.
「Hallelujah」和「Allelujah」有什麼分別?.
有什麼分別?牧師,有呀.
「Hallelujah」和「Hallelujah」是差了一個前置的字標.
差了一個希伯來文字母.
「Hei」什麼意思?.
每逢在一個動詞上,如果那個發音用「Hei」.
就是變成呼籲的意思.
所以我們叫「Hallelujah」,就是「來吧,讚頌主」.
「來吧,當我們在一個德性的局面的時候.

$^{561}$聖經自己都和應,為什麼我們不是因此而為此而和應呢?.
我們可以向死亡大膽的說,你動得了我嗎?」.
可以講得俗一點,「你殺了我嗎?」.
你沒有機會,因為聖經說了「我們贏了」.
說了這麼久,為什麼會贏了呢?.
在下面第56至58節裡講給我聽.
不過當我們進入最後的部分的時候.
這裡講了些什麼給我們知道?.
復活節的主旨,復活節的主旨帶我們進入了一個改變的裡面.
讓我們有符合上帝天國子民的身份.
使我們能夠知道我們在一個聖局之下,怎樣走面前的路.
而在這裡就告訴我們,進行一個戰爭的最後的呼喚.
一起去面對死亡所帶來的後果.
「我贏了,不用擔心,結局已定,不過還有些手尾要執」.
弟兄姊妹,你有沒有試過,有時候工作.
當你做不定的時候,有人幫你搞定了主要的問題之後.
剩下的就交給你去處理.
你那種釋放,那種的重擔,完全是沒有需要擔憂.
用多個例子,你喜不喜歡看劇集?.
喜歡吧?.
喜不喜歡看日劇,韓劇那些專門做醫院的劇集呢?.
很多很囂張的醫生做手術的,是不是?.
有沒有?看過了.
有沒有留意,很多時候那些很囂張的主角.
做擔當醫生的,去做一些其他人都做不了的手術之後.
做完之後那個手術,接著他做什麼?.
他叫身邊的那些助手,搞定了他.
譬如做心臟手術.
哇,不停地流血,劇情很擠迫,很緊湊.
不停地流血,為什麼我找不到流血的地方呢?.
最後那個醫生很高超的技巧,找到流血的地方,幫他止血.
把心臟治好了.
然後最後怎麼樣?.
他就走了.
然後叫身邊那些人,助手,搞定了他.
那些助手要做的事情,不是很高難度的事情.
可能縫在傷口上,加多少針,現在不用針了.
現在好像用釘書機,啪啪啪.
最近我也聽說過,牙齒也是這樣.
我聽說牙醫不同以前了.

$^{601}$有些膠水黏的,那麼厲害.
那些輕輕的東西留給那些助手做.
上帝在整個拯救的計劃裡面.
最難的已經找了耶穌做了,釘在十字架上.
Complete,現在剩餘的那些.
讓我們這些去收拾最後的手尾.
而那些東西根本是很輕的.
我們有沒有跟主在這個時候繼續走這條路呢?.
還是捉主耶穌回來,好像有些事情還沒搞定.
你有沒有見過男主角,女主角醫生回來.
「讓我先連接兩針吧」.
不是的,有些事情剩下留給你去做.
為什麼?為什麼?.
是不是那個醫生不會做?不是.
讓你分享到這份成功.
我很喜歡講這個例子.
我在神學院有個教授,有個兒子很小的,三歲.
有一次下大雪,要出去後院.
不是,要去車路那裡鏟雪.
兒子說「我都幫」.
很小的,三歲,拿著鏟子,變救力鏟.
鏟了兩鏟就已經在那裡玩雪了.
他很開心,車路還有很多雪.
接著他爸爸沒有出聲.
接著鏟鏟鏟,鏟清那條雪,雪路之後.
接著就跟兒子說.
「兒子,你那裡還有一束雪,可不可以幫我撥走它」.
接著那個小朋友的兒子就去鏟那一束雪.
然後就跟爸爸說.
「爸爸,我多聰明,幫到你鏟完那些雪」.
鏟清妹,上帝是這樣的.
祂不是,就算那些事情是祂做的,祂都不領功.
讓你和我有份去領功.
其實我做爸爸都一樣.
有的時候做了很多事情.
叫那些小朋友,叫那些子女來幫忙.
其實那些事情,其實那些事情.
其實我做完,比叫你來做還要快.
為什麼要叫你呢?.
讓你分享那份成功感.

$^{641}$如果你不想的.
我告訴你我做爸爸的思想是怎樣.
這個壞孩子,壞女兒.
想和你一起分享份喜樂,你都不想.
還以為我煩你.
今天很多信徒就是這樣.
上帝在德性的時候.
讓你的復活身體可以有符合進入天國的狀態.
你仍然不想做.
上帝和你分享,你不想.
那你過什麼復活節呢?.
我多看兩場那些《Passion of Christ》.
我不知道《基督的受難記》.
或者回教會多搞兩堂崇拜.
搞幾齣話劇給你看.
你就覺得自己很虔誠.
你可以在這個地上很虔誠.
但原來你的生命沒有改變.
我恐防你的信不能夠帶你去到這裡所說的.
由扭壞變成不扭壞.
由血肉肢體仍然不能夠領受神國的血肉肢體.
最後你要死的仍然藏在死亡當中.
因此去到最後.
從第54節至第55節.
當我們經過了這個所謂聖歌序的歡呼之後.
56至58節就告訴你.
這份感恩能夠促使我們在神這個聖局之下.
更竭盡所能去投入.
我們一起讀一下這幾節經文.
56至58.
一起讀.
死的毒毆就是罪.
罪的權勢就是律法.
我們的主耶穌基督得勝.
所以我親愛的弟兄們.
你們務要堅固.
不可搖動.
常常竭力多做主公.
因為知道你們的奴婦在主裡面.
不是徒然的.

$^{681}$當我們去看看這段聖經第56節.
重新講到有關於罪和死的運作的時候.
我們發現到撒旦其實是很狡猾的.
怎樣狡猾?.
他說死的毒毆就是罪.
沒錯.
死的毒毆就是罪.
但這個罪的毒毆.
是依仿什麼來得到力量去毆你呢?.
或者毆我們呢?.
看清楚.
罪的權勢就是律法.
律法.
誰頒佈律法?.
上帝頒佈律法.
撒旦的毒招.
是借上帝力量來殺我們.
因為他知道上帝的律法頒佈之後.
是揭露人的罪.
這樣就對了.
撒旦就想到.
是啊.
還不好?.
用上帝的條件來挾持上帝所做的人.
指控他們有罪.
毒吧這招.
他不用出什麼東西.
但是也告訴我們知道.
其實撒旦的罪的毆是怎麼來的.
他的權柄.
不是由他自己而來.
是由上帝的律法而來.
他想著這樣就贏了.
上帝我還不將你軍.
明白什麼意思嗎?.
用上帝的力量來將軍你上帝所喜悅的人.
其實是轉過頭來.
用上帝的力量來指控上帝.
是你弄死人的.
因為你頒佈律法.

$^{721}$明白嗎?.
魔鬼撒旦多可怕.
他借上帝的刀來拆上帝的人.
轉過頭來拆上帝.
毒啊這招.
不過保羅說.
上帝卻是很精明地將他的詭計打破.
怎麼打破?.
使我們借著主耶穌的特性.
第三次在這段聖經講到特性的問題.
我們不單止有吹響幕後的豪閣.
奠定聖局.
我們不單止可以向罪的毒毆.
去跟他宣告.
你是失敗的.
你不會在勝利的裡面.
因為我們就在勝利的當中.
最後這裡告訴我們.
原來我們的源頭.
特性不是因為弟兄姊妹.
你和我有什麼巴閉.
而是因為耶穌基督所造成的事.
耶穌做了什麼事?.
耶穌在十字架的上面.
被釘死.
耶穌是按照羅馬人.
以及猶太人的律法.
被釘死在十字架上面.
但是卻因為他的死.
與及復活扭轉了.
原本你以為用律法去指控上帝的人.
以至間接指向上帝.
這種的手段被打破了.
弟兄姊妹.
沒有耶穌的復活.
你不斷釘死他們是沒有用的.
沒有耶穌基督的復活.
我們是沒有盼望的.
但正正因為耶穌有復活.
所以你和我的身體會改變.

$^{761}$我們不需要再在死亡的當中.
而且我們可以有一個配得上帝的身份.
是耶穌所做的工作.
打破了所有罪 透過律法向我們指控.
所以我再不需要擔心被人指控我們.
有些事情做得不好.
尤其是在屬靈世界上面.
你如果真心相信上帝.
所差派的獨生兒子耶穌基督.
祂的復活.
就能夠帶你和我同樣進入復活的身份.
當你在這個復活的身份裡面的時候.
是怎樣表現呢?.
我們看回58節.
當我們是在一個復活的狀態之下.
我親愛的弟兄們.
你們務要堅固 不可搖動.
常常竭力多做主公.
這裡連續說了兩件事.
第一件事 在心智的上面.
你要堅固不動搖.
其實這些陳述和昔日約書亞記.
當約書亞接受委託.
同樣都是這樣.
不要動搖.
只管堅固在哪裡.
在耶穌基督的教導上面.
而我們的表現是怎樣呢?.
行動上是怎樣呢?.
常常竭力多做主公.
首先我要強調「主公」這個說話.
是上帝所喜悅的事.
不只是教會的事奉或者教會的事工.
何謂竭力呢?.
在原文聖經最好的翻譯就是「出類拔萃」.
英文用了一個字眼叫「outstanding」.
不只是站在那裡那麼簡單.
是要非常之突出.
意思是我們所做的事是竭盡所能去做的.
不是隨隨便便做了.

$^{801}$或者做過了.
不是這樣.
是要做到加零一的.
我用一個簡單的說話.
用我做港獨的例子.
不是說了.
而是講到聖經裡面上帝的心思是甚麼.
不是讀了.
讀了就說完了.
你留意一下現在的港獨.
你留意一下現在你們上網有很多人聽.
很多人讀了然後就講他的故事.
那段聖經講甚麼呢?.
上帝的心是講甚麼呢?.
想講些甚麼給我們知道.
沒有講的.
同樣道理.
如果你做大師的.
唱了.
是不是唱了?.
我以前做師班的.
我唱師班的.
我的聲音也不差.
是這幾年老了的聲音.
但是你唱詩是不是唱了?.
不是.
那首詩歌裡面想講些甚麼.
你帶不帶出來給人知道?.
即使你出來領詩也是一樣.
唱三節,四節.
唱完了.
唱了,事奉了.
是不是這樣?.
不是.
要做到加零一.
帶給人看到這首詩歌裡面想講些甚麼.
不是叫你說話.
不是講很多東西.
但是至少你將那首詩歌的核心講了出來.
讓大家唱的時候.

$^{841}$將那個核心講了出來.
如果我們不是這樣做的.
沒有這份生命態度的.
做甚麼基督徒呢?.
隨隨便便.
就是做些事奉.
隨隨便便搞個報道會.
隨隨便便有些人進來.
他真信假信你不要理.
他要受浸.
隨隨便便浸吧.
他要浸就讓他浸吧.
不要問那麼多了.
他要信主就讓他信吧.
是不是這樣?.
弟兄姊妹.
這個求其事奉.
我最不想的一件事是甚麼呢?.
你帶了很多這樣不是真的信主的人.
將來全部都復活.
全部去到耶穌基督面前.
你還要搞到耶穌.
跟耶穌說 喂喂喂 這些是不信主的.
你又說排錯隊了.
到時那些你以為他一路信主的人.
what a big surprise.
就趕他回去排一條對的隊.
那條死在過犯當中的隊.
是不是這樣呢?.
我最不想麻煩耶穌就是這樣.
今天很多的教會是這樣做.
不過不要說我們做甚麼.
你自己是怎樣信主的?.
是不是在復活上面有盼望.
跟耶穌參與其中一起呢?.
並且願意跟耶穌一起去到.
一起完成這個德性的結局呢?.
很多人說事奉主是有很多的捨棄.
很多世人甚至認為信了主的人.
有很多事情都不能做 有很多損失.

$^{881}$阻礙了我人生的希望和盼望.
但是我要多說一句.
沒錯 基督徒有很多事情是不會做的.
因為我們選擇不去做.
因為那些事情不合上帝的心意去做.
但我們卻會盡心盡力地去做上主交給我們的事.
復活節應該說的是復活的能力.
怎樣讓我們在往後的日子.
靠著一個可以復活的生命.
跟主一起毀身走前面的路.
完成上帝交代給我們.
在我們時日生命空間裡面所能夠做的事.
吸引更多人 真正信主的人.
回到上主的當中.
求主幫助我們.
讓我們一齊同心低頭度過.
父神多謝你讓我們在今天這段聖經裡面的學習.
在保羅的教導裡面.
將你的啟示講得清清楚楚.
今天我們很多時候.
活在地上.
我們失去了那份基督徒真正可以有復活盼望.
復活的生命的改變的這個基礎.
來去活在地上.
民心可能是我們根本信主信得不清楚.
但今天讓我們明白的時候.
求恩主再一次接納我們的認罪悔改.
做一個真正匹配由主耶穌基督所帶來的復活的生命.
而不是一個掛名的基督徒.
求你的帶領.
讓徹底的生命改變.
在每一個聽到今天這堂訊息.
以及讀到聖經裡面啟示的人的生命當中.
讓我們是一個不折不扣的基督徒.
全然靠這位賦予我們復活的生命主.
耶穌基督.
禱告逢靠主耶穌去到,得勝名自求..
阿們..
\newpage



\section{路加福音 24:13-35}
\label{sec:031917}
\textbf{遇見復活的主 (粵/英)}~~~~李適清博士
\newline
\newline
連結: \href{http://www.church.com.hk/acms/content.asp?site=cdc\&op=show\&type=product\&code=031917\&layout=sermon}{\texttt{031917}} ~~~~ 語音日期: 2026-04-05
\newline
\newline
\hyperref[sec:031916]{\small{< < < PREV SERMON < < <}}
~
\hyperref[sec:index_chronic]{\small{[返順時目]}}
~
\hyperref[sec:031918]{\small{> > > NEXT SERMON > > >}}
\newline
\newline
路加福音 24:13-35
\newline
\begin{longtable}{cl}
\hline
\hline
章節 & 經文 (和合本修訂版)\\
\hline
24:13 & \begin{tabularx}{0.7\textwidth}{X} 同一天，門徒中有兩個人往一個村子去；這村子名叫以馬忤斯，離耶路撒冷約有二十五里。 \end{tabularx} \\ \\ \relax
24:14 & \begin{tabularx}{0.7\textwidth}{X} 他們彼此談論所發生的這一切事。 \end{tabularx} \\ \\ \relax
24:15 & \begin{tabularx}{0.7\textwidth}{X} 正交談議論的時候，耶穌親自走近他們，和他們同行， \end{tabularx} \\ \\ \relax
24:16 & \begin{tabularx}{0.7\textwidth}{X} 可是他們的眼睛模糊了，沒認出他。 \end{tabularx} \\ \\ \relax
24:17 & \begin{tabularx}{0.7\textwidth}{X} 耶穌對他們說：「你們一邊走一邊談，彼此談論的是甚麼事呢？」他們就站住，臉上帶著愁容。 \end{tabularx} \\ \\ \relax
24:18 & \begin{tabularx}{0.7\textwidth}{X} 兩人中有一個名叫革流巴的回答：「你是在耶路撒冷的旅客中，惟一還不知道這幾天在那裡發生了甚麼事的人嗎？」 \end{tabularx} \\ \\ \relax
24:19 & \begin{tabularx}{0.7\textwidth}{X} 耶穌對他們說：「甚麼事呢？」他們對他說：「就是拿撒勒人耶穌的事。他是個先知，在神和眾百姓面前，說話行事都大有能力。 \end{tabularx} \\ \\ \relax
24:20 & \begin{tabularx}{0.7\textwidth}{X} 祭司長們和我們的官長竟把他解去，定了死罪，釘在十字架上。 \end{tabularx} \\ \\ \relax
24:21 & \begin{tabularx}{0.7\textwidth}{X} 但我們素來所盼望要救贖以色列民的就是他。不但如此，這些事發生到現在已經三天了。 \end{tabularx} \\ \\ \relax
24:22 & \begin{tabularx}{0.7\textwidth}{X} 還有，我們中間的幾個婦女使我們驚奇：她們清早去了墳墓， \end{tabularx} \\ \\ \relax
24:23 & \begin{tabularx}{0.7\textwidth}{X} 不見他的身體，就回來告訴我們，說她們看見了天使顯現，說他活了。 \end{tabularx} \\ \\ \relax
24:24 & \begin{tabularx}{0.7\textwidth}{X} 又有我們的幾個人往墳墓那裡去，所發現的正如婦女們所說的，只是沒有看見他。」 \end{tabularx} \\ \\ \relax
24:25 & \begin{tabularx}{0.7\textwidth}{X} 耶穌對他們說：「無知的人哪，先知所說的一切話，你們的心信得太遲鈍了。 \end{tabularx} \\ \\ \relax
24:26 & \begin{tabularx}{0.7\textwidth}{X} 基督不是必須受這些苦難，然後進入他的榮耀嗎？」 \end{tabularx} \\ \\ \relax
24:27 & \begin{tabularx}{0.7\textwidth}{X} 於是，他從摩西和眾先知起，凡經上所指著自己的話都給他們作了解釋。 \end{tabularx} \\ \\ \relax
24:28 & \begin{tabularx}{0.7\textwidth}{X} 他們走近所要去的村子，耶穌好像還要往前走， \end{tabularx} \\ \\ \relax
24:29 & \begin{tabularx}{0.7\textwidth}{X} 他們卻強留他說：「時候晚了，天快黑了，請你同我們住下吧。」耶穌就進去，要同他們住下。 \end{tabularx} \\ \\ \relax
24:30 & \begin{tabularx}{0.7\textwidth}{X} 坐下來和他們用餐的時候，耶穌拿起餅來，祝福了，擘開，遞給他們。 \end{tabularx} \\ \\ \relax
24:31 & \begin{tabularx}{0.7\textwidth}{X} 他們的眼睛開了，這才認出他來。耶穌卻從他們眼前消失了。 \end{tabularx} \\ \\ \relax
24:32 & \begin{tabularx}{0.7\textwidth}{X} 他們彼此說：「在路上他和我們說話，給我們講解聖經的時候，我們的心在我們裡面豈不是火熱的嗎？」 \end{tabularx} \\ \\ \relax
24:33 & \begin{tabularx}{0.7\textwidth}{X} 於是他們立刻起身，回耶路撒冷去，看見十一個使徒和與他們正在一起的人聚集在一處， \end{tabularx} \\ \\ \relax
24:34 & \begin{tabularx}{0.7\textwidth}{X} 說：「主果然復活了，已經顯現給西門看了。」 \end{tabularx} \\ \\ \relax
24:35 & \begin{tabularx}{0.7\textwidth}{X} 於是，兩個人把路上所遇到，和耶穌擘餅的時候怎麼被他們認出來的事，都述說了一遍。 \end{tabularx} \\ \\
[1ex]
\hline
\hline
\end{longtable}
$^{1}$(英語) 各位姐妹早安!.
(英語) 我不知道復活節對你來說有什麼意義?.
(英語) 很多年前,我在多倫多這裡.
(英語) 就在復活節.
(英語) 我就在這個城市受浸,加入教會.
(英語) 所以對復活節,尤其是浸禮,復活節的浸禮,真的很有感覺.
(英語) 我相信我們每個人都有跟耶穌相遇的故事.
(英語) 我們今天選擇的經文,就是發生在主耶穌復活之後.
(英語) 門徒他們遇到復活的主耶穌.
(英語) 經文記載在路加呼音二十四章十三至三十五節.
(英語) 我們一起來看看,先看中文,然後看英文.
(英語) 耶穌對他們說.
(英語) 你們一邊走一邊談,談什麼呢?.
(英語) 他們站在那裡面帶愁容.
(英語) 兩個人中間有一個叫格劉巴的回答.
(英語) 你在耶路撒冷的旅客當中.
(英語) 是唯一不知道這幾天在那裡發生了什麼事的人嗎?.
(英語) 耶穌對他們說.
(英語) 什麼事呢?.
(英語) 他們對他說就是拿撒勒人耶穌的事.
(英語) 他是個先知,在上帝和眾百姓面前.
(英語) 說話行事都大有能力.
(英語) 祭司長門和我們的官長竟把他戒去.
(英語) 定了死罪釘在十字架上.
(英語) 但我們素來所盼望要救贖以色列民的就是他.
(英語) 不但如此,這些事發生到現在.
(英語) 已經三天了,還有我們中間的幾個婦女使我們驚奇.
(英語) 她們清早去了墳墓.
(英語) 不見了她的身體,就回來告訴我們.
(英語) 說她們看見了天使顯現,說她活了.
(英語) 又有我們的幾個人往墳墓那裡去.
(英語) 所發現的正如婦女們所說的.
(英語) 只是沒有看見她.
(英語) 耶穌對他們說.
(英語) 無知的人啊,先知所說的一切話.
(英語) 你們的心信得太遲鈍了.
(英語) 基督不是必須受這些苦難,然後進入他的榮耀嗎?.
於是,他從摩西和眾仙子起.
由經上所指自己的話都幫他們解釋一次.
(英語) 他們走近所去的村的時候,耶穌好像還要繼續向前走.

$^{41}$(英語) 他們就強留著他.
(英語) 時間晚了,天快黑了,請你和我們住一下.
(英語) 耶穌就進去要與她們住一下.
(英語) 坐下來和她們用餐的時候.
(英語) 耶穌拿起餅來,祝福了,張開,遞給她們.
(英語) 她們的眼睛開了,這才認出他來.
(英語) 耶穌卻從她們眼前消失了.
(英語) 他們彼此就說.
(英語) 在路上他和我們說話,給我們講解聖經的時候.
(英語) 我們的心在我們裡面,豈不是火熱的嗎?.
於是他們就立刻起來,回去耶路撒冷.
(英語) 見到十二個使徒和他們聚集在一起.
(英語) 主果然復活了,已經顯現給西門漢了.
(英語) 於是兩個人將所遇到的事情.
(英語) 和耶穌擦餅的時候.
(英語) 如何讓他們認出的事情都說了一次.
我們現在要聽聽英文的故事.
那一天,兩人去了一座村子叫MAS.
(英語) 約七公里離耶路撒冷.
(英語) 他們在談論這些事情.
(英語) 當他們在談論的時候.
(英語) 耶穌自己靠近了他們.
(英語) 但他們的眼睛仍然在認出耶穌.
(英語) 耶穌對他們說.
(英語) 你們為什麼在走著走著.
(英語) 站著站著看著悲傷呢?.
(英語) 於是一位叫Cleopas的使徒回答.
(英語) 你是唯一一個來耶路撒冷的使徒嗎?.
(英語) 耶穌並不知道這些事情在那裡發生了.
(英語) 耶穌對他們說.
(英語) 什麼事情?.
(英語) 他們對耶穌說.
(英語) 關於耶穌基督.
(英語) 耶穌是一個神父.
(英語) 祂在神和所有人面前.
(英語) 真正的勇氣和道德.
(英語) 祂如何使主教徒和領導人.
(英語) 使祂被懲罰至死.
(英語) 並將祂釋放.
(英語) 但我們曾經希望.

$^{81}$(英語) 祂能夠救救耶穌.
(英語) 是的.
(英語) 除此之外.
(英語) 現在已經是第三天.
(英語) 這些事情發生了.
(英語) 更何況.
(英語) 有些我們公司的女性.
(英語) 令我們驚訝.
(英語) 她們早上就在寺廟裡.
(英語) 當她們找不到祂的身體時.
(英語) 她們還說.
(英語) 她們甚至看到了.
(英語) 天使的視線.
(英語) 說祂是活著的.
(英語) 那些和我們在一起的人.
(英語) 到寺廟裡找到了.
(英語) 正如那些女性所說的.
(英語) 她們看不到祂.
(英語) 祂對她們說.
(英語) 噢 愚蠢的人.
(英語) 心懦弱的.
(英語) 相信所有的聖徒所說的.
(英語) 難道不需要.
(英語) 讓耶穌要忍受這些事情.
(英語) 而進入祂的榮耀嗎.
(英語) 開始的時候.
(英語) 耶穌與所有的聖徒.
(英語) 祂對她們解釋了.
(英語) 所有的經書.
(英語) 關於祂自己的事情.
(英語) 所以她們走到寺廟裡.
(英語) 到那裡去.
(英語) 祂做的就是.
(英語) 祂要走得更遠.
(英語) 但是她們強烈地呼籲祂.
(英語) 祂說 跟我們一起.
(英語) 因為現在是晚上.
(英語) 日子已經過得很遠了.
(英語) 所以祂進去跟她們一起.
(英語) 但是祂在桌上.

$^{121}$(英語) 跟她們一起.
(英語) 祂拿了麵包.
(英語) 祂祝福了她們.
(英語) 砸了它.
(英語) 還給了她們.
(英語) 她們的眼睛睜開了.
(英語) 她們認出了祂.
(英語) 祂消失了.
(英語) 她們對彼此說.
(英語) 祂說 我們的心.
(英語) 祂說 祂在路上.
(英語) 祂在路上說 祂在路上.
(英語).
我相信我們每一個人.
我們經歷耶穌.
故事都不短.
都有這麼豐富的情節在裡面.
(英語).
復活的主耶穌出現在門徒身邊.
(英語).
跟他們走了好一段路.
(英語).
到最後他們才認得耶穌.
(英語).
其實想真一點.
(英語).
如果耶穌一出現.
門徒就認得祂了.
(英語).
你會覺得怎樣呢?.
(英語).
門徒很害怕見鬼.
(英語).
因為他們從來都沒想過.
耶穌會復活的.
(英語).
但是耶穌用了一個很溫柔的方法.
(英語).
他們一路走一路聊.
(英語).

$^{161}$門徒覺得心裡面火熱.
(英語).
最後耶穌用了一個他們很熟悉的動作.
(英語).
在吃飯的時候擦餅.
(英語).
他們認得耶穌了.
(英語).
可能我們都有不認得耶穌的日子.
(英語).
曾幾何時我們以為.
祂不聽禱告,祂不在.
(英語).
但是其實耶穌一直都和我們一起.
(英語).
只不過有時我們被自己想法和想要的東西.
阻礙了我們.
(英語).
有一個很古老的傳說.
講到印度一個部族.
(英語).
當年輕人成長的時候.
(英語).
他們到一定年紀有一個成年禮.
(英語).
當少年人準備要做成年人的時候.
(英語).
他們的爸爸就要帶他們去一個森林裡面.
(英語).
然後他們蒙著眼就單獨過一夜.
(英語).
他們不可以打開眼睛到第二天早上.
(英語).
他們不可以喊救命.
(英語).
因為他們要過到這一晚.
仍然他生存就變成成年人.
(英語).
然後他成年之後就不可以將這個經歷和別人分享.
(英語).

$^{201}$這一晚有一個年輕人.
(英語).
他就蒙著雙眼坐在這個漆黑的森林裡面.
(英語).
四周圍肯定有很多猛獸.
(英語).
他不單止聽到那些風吹草叢的聲音.
(英語).
而且他坐在那個樹幹好像有點動.
(英語).
雖然他很害怕.
(英語).
全身都繃緊.
(英語).
不過他都沒有打開蒙著眼的那個東西.
(英語).
他只是動都不動的坐在那裡.
(英語).
因為他知道這樣他才會成為成年人.
(英語).
終於他經過了這個恐怖的夜晚.
(英語).
他感覺到陽光照在他的臉上.
(英語).
他知道他可以打開蒙著眼的那個面罩.
(英語).
這個時候他發現原來他自己的父親.
(英語).
整晚他爸爸都守護著他.
(英語).
我們的主耶穌就在我們不知道的時候.
(英語).
他一直都在這裡.
(英語).
是什麼阻住了門徒他們認得耶穌呢?.
(英語).
或者門徒他們太過沉醉於自己的想法.
(英語).
他們不覺得他們真的會遇見耶穌.
(英語).

$^{241}$可能我們都試過.
我們很想用自己的方法去找神.
(英語).
我們去決定信不信耶穌.
(英語).
但其實我們連怎樣去找他.
去認識他都沒有能力.
(英語).
是耶穌親自去找我們.
(英語).
他就在我們的身邊.
就在我們的經歷裡面.
(英語).
他經常都在等著我們去迴轉去歸向他.
(英語).
所以我們做的很簡單.
(英語).
我們向我們的主開放.
(英語).
我們專心一意地去認識他跟從他.
(英語).
在路加夫音整個故事的敘述裡面.
(英語).
核心在第23節那裡說.
(英語).
復活的耶穌用另外一個.
好像隱藏的形式那樣跟他們相遇.
(英語).
在聖經當中就記載了.
他怎樣向他們說.
這條受苦然後走向榮耀的路.
(英語).
在路加夫音24章25節那裡.
耶穌對他們說.
(英語).
可能你都信了主了.
(英語).
可能一段日子.
(英語).
甚至好像那些門徒那樣跟耶穌生活了幾年.

$^{281}$(英語).
但有時候我們居然不認得耶穌.
(英語).
因為我們根本沒有認真地想過.
他會介入在我們的生活當中.
(英語).
在我們日常的小事裡跟我們會有互動.
(英語).
在門徒的生命裡.
(英語).
耶穌的死令到他們陷入一個很深沉的悲哀.
(英語).
他們的希望好像粉碎了一樣.
(英語).
因為他們希望耶穌就在地上做王.
(英語).
但耶穌的國是超乎我們想像的.
(英語).
耶穌怎樣回應他們這個想法呢?.
(英語).
他說門徒信得太遲鈍了.
(英語).
其實門徒都知道舊有聖經所說的彌賽亞會來到的.
(英語).
但是他們未能夠用一個信心去回應.
(英語).
今天我們不單止知道耶穌為我們.
(英語).
道成肉身,死而復活.
(英語).
我們知道,而且我們更加有聖經的話.
(英語).
我們有歷史的基礎.
(英語).
我們有歷世歷代信徒的見證.
(英語).
讓我們可以很肯定.
(英語).
然後在當中我們一起去學習和成長.
(英語).

$^{321}$你有沒有為自己的生命去設計成長的時間和空間呢?.
(英語).
我想當時的門徒就是他們口講跟耶穌.
(英語).
他們期望耶穌去做一些事.
(英語).
但是耶穌說他們的信心太遲鈍了.
(英語).
你有多久沒有把握機會去繼續認識神呢?.
(英語).
去找一個空間給自己去遇見耶穌.
(英語).
你記不記得當初你收針的時候那種的經歷呢?.
(英語).
耶穌祂是繼續和我們一起與我們相遇.
(英語).
有一個年輕人他在一間木廠裡面工作.
(英語).
上班第一天他斬了十棵樹.
(英語).
實在很勁自己都很開心.
(英語).
老闆都很滿意.
(英語).
第二天他同樣努力斬了八棵樹成績都不錯.
(英語).
一個星期之後他同樣努力工作.
但是生產力就越來越少了.
(英語).
老闆都覺得可能出現了一些問題.
(英語).
這一天他只是斬了一棵樹.
(英語).
老闆覺得很難接受.
就叫他來辦公室談一談.
(英語).
年輕人對老闆說.
(英語).
我每一天都很努力地斬樹.
我已經盡了力了.

$^{361}$(英語).
老闆說.
(英語).
老闆對年輕人說.
(英語).
年輕人說.
(英語).
我實在太忙了.
(英語).
年輕人說.
(英語).
有時候我們都很沉醉.
在自己做的事情上.
(英語).
每一天都很忙.
我們努力地去生活.
(英語).
過了一段時間.
我們好像還是原地踏步.
(英語).
我們很容易失去了真正的焦點.
(英語).
或者我們生命的重心.
(英語).
你要怎樣去用你餘生的時間.
以至不失去這些最重要的東西呢.
(英語).
上帝就是這樣將時間,能力,恩賜.
放在我們的手上.
(英語).
然後祂和我們一起.
(英語).
所以我們要好好地去把握神給我們的時間.
(英語).
發現耶穌在我們生命裡面的工作.
(英語).
其實門徒他們沒有什麼能力去認出耶穌.
(英語).
是到了耶穌親自去行動了.
(英語).

$^{401}$為他們擦餅.
(英語).
然後讓他們認得出.
(英語).
所以我們是要有一個願意親近神的心.
(英語).
讓神親自在我們的生命裡面工作.
(英語).
弟兄姊妹,你有沒有試過像門徒般絕望的日子.
(英語).
好像信了耶穌,但是祂在哪裡呢?.
(英語).
有時我們身邊的一段關係結束了.
(英語).
有一個機會我們錯過了.
(英語).
醫生說的一句話讓我們的心沉到谷底.
(英語).
這些感覺好像故事已經寫到最後一頁.
(英語).
但是耶穌的復活是提醒我們.
(英語).
人以為不可能的,上帝可以令到它可能.
(英語).
在人以為結束的地方.
往往就是上帝開始去工作的時候.
(英語).
原來在人生的路上,無論在什麼時刻.
(英語).
主耶穌一直與我們同在.
(英語).
祂親自找回我們.
(英語).
每一天就像在二馬五司路上和我們同行.
(英語).
讓我們不斷經歷上帝奇妙的恩典和帶領.
(英語).
耶穌一直與你同在,你認不認得祂呢?.
(英語).
我們同心禱告.

$^{441}$(英語).
因為是你親自尋找我們.
(英語).
在我們迷惘,失落,沒有盼望的時候.
你尋找我們.
(英語).
你用很溫柔的方法帶領著我們.
(英語).
我們求你讓我們認得你.
(英語).
開我們眼睛,讓我們看到你的心意和帶領.
(英語).
以致我們可以每一天都經歷你的恩典.
(英語).
我們多謝你.
(英語).
我們禱告奉耶穌基督得勝的命求.
(英語).
Amen.
\newpage



\section{哥林多前書 15:1-8}
\label{sec:031918}
\textbf{復活的主}~~~~鄭兆奇牧師
\newline
\newline
連結: \href{http://www.church.com.hk/acms/content.asp?site=cdc\&op=show\&type=product\&code=031918\&layout=sermon}{\texttt{031918}} ~~~~ 語音日期: 2026-04-05
\newline
\newline
\hyperref[sec:031917]{\small{< < < PREV SERMON < < <}}
~
\hyperref[sec:index_chronic]{\small{[返順時目]}}
~
\hyperref[sec:031919]{\small{> > > NEXT SERMON > > >}}
\newline
\newline
哥林多前書 15:1-8
\newline
\begin{longtable}{cl}
\hline
\hline
章節 & 經文 (和合本修訂版)\\
\hline
15:1 & \begin{tabularx}{0.7\textwidth}{X} 弟兄們，我要你們認清我先前傳給你們的福音；這福音你們領受了，又靠著它站立得住， \end{tabularx} \\ \\ \relax
15:2 & \begin{tabularx}{0.7\textwidth}{X} 你們若能夠持守我傳給你們的信息，就必因這福音得救，否則你們是徒然相信。 \end{tabularx} \\ \\ \relax
15:3 & \begin{tabularx}{0.7\textwidth}{X} 我當日所領受又傳給你們的，最重要的就是：照聖經所說，基督為我們的罪死了， \end{tabularx} \\ \\ \relax
15:4 & \begin{tabularx}{0.7\textwidth}{X} 而且埋葬了；又照聖經所說，第三天復活了， \end{tabularx} \\ \\ \relax
15:5 & \begin{tabularx}{0.7\textwidth}{X} 還顯給磯法看，又顯給十二使徒看， \end{tabularx} \\ \\ \relax
15:6 & \begin{tabularx}{0.7\textwidth}{X} 後來一次顯給五百多弟兄看，其中一大半到現在還在，卻也有已經睡了的。 \end{tabularx} \\ \\ \relax
15:7 & \begin{tabularx}{0.7\textwidth}{X} 以後他顯給雅各看，再顯給眾使徒看， \end{tabularx} \\ \\ \relax
15:8 & \begin{tabularx}{0.7\textwidth}{X} 最後也顯給我看；我如同未到產期而生的人一般。 \end{tabularx} \\ \\
[1ex]
\hline
\hline
\end{longtable}
$^{1}$今天是復活節.
是一個充滿盼望和榮耀的日子.
因為耶穌基督復活了.
我們也會復活.
記得勝經說.
復活在我 生命也在我.
信我的人 雖然死了 也必復活.
凡活著信我的人 就永遠不死.
你信這話嗎?.
我相信這是上帝給我們的一句話.
是一個很大的應許.
耶穌死了之後 三天就復活了.
在我們香港的生命當中.
很多時候我們會經歷到.
不同大大小小的復活經驗.
我們LG有一個單車隊叫LG Cycling.
現在我們LG系列比較多.
我們有LG Cafe,LG Choir,LG Cycling.
即是一路來.
甚麼Smart TV 很多系列.
昨天他們一起去騎單車.
我其實也騎過一兩次.
其實昨天我很大掙扎.
我可不可以加入他們.
因為他們是開始升級了.
因為無限的升級.
15個mile,18個mile,20個mile,24個mile.
昨天看到他們24個mile.
我說不要了.
今天星期天又要分享訊息.
我怕我出不來.
但有一個意念就是.
耶穌基督死了之後.
三天復活了.
我又想如果我昨天騎單車.
騎到死了.
我要三天後才復活.
我真的今天來不了.
這些都是弟兄姊妹一起做的時候.
你感覺到上帝在我們當中.

$^{41}$我們這個Cycling不是真的就這樣去騎單車.
我們希望藉著這個Cycling.
我們能夠接觸一些新朋友.
我們一起將這個恩典和福氣.
給眾多大能傳給他們.
在聖經很清楚的跟我們說.
他說我們從前領受.
又交給你們最要緊的是什麼事呢.
他就說基督照聖經所記載.
為我們的罪死埋葬.
接著又照聖經所說.
第三天復活.
並且怎樣呢.
並且你會開始發覺聖經說了很多顯現.
並且向誰顯現.
向羈髮顯現.
然後又向十二個使徒顯現.
以後又有一次向五百個弟兄顯現.
他們中間大多數現在還在.
即是當時在那時的現在.
有一些已經睡了.
睡了即是我們在相信主的人.
我們知道我們不是就這樣死了.
我們只是睡了.
以後又向雅各顯現.
再後來又向使徒顯現.
最後向我顯現.
我就是寫這段經文的仕途保羅.
他說我好像一個未到產期而生的人.
你會看到耶穌基督死了之後復活.
他復活之後他不是馬上升天的.
他向很多人顯現.
最後仕途保羅說他也向我顯現.
而我好像一個未到產期而生的人.
知不知道什麼叫未到產期而生.
早產.
為什麼他無端說他只是一個早產仔.
他為什麼說自己早產了.
即是說他未相信主之前.
上帝已經向他顯現.

$^{81}$我不知道你有沒有經歷過.
即是上帝那種顯現是怎樣.
是很難理解的.
但其實整個聖經你看回.
很多.
他向過很多人顯現.
瑪利亞 幾個瑪利亞.
一群婦女 彼得.
二馬五師路上的一些門徒.
十一個門徒 多瑪.
加利利顯現有群門徒.
有五百個人 雅各.
升天之前又顯現.
即是耶穌基督復活了之後.
他沒有將自己收買.
有一個姐妹跟我分享.
她說耶穌向她顯現.
我很驚訝 看到耶穌.
我問她 她跟你說些什麼.
她說 耶穌沒有跟我說什麼.
但她感受到耶穌那份愛.
因為是一份十字架上的愛.
很濃的一種愛 很平安的一種愛.
我再想想 其實有沒有可能呢.
你覺得有沒有可能呢.
其實也有可能.
因為我最初想 耶穌在十字架上說了這麼多.
我們回來的時候.
也說耶穌基督在十字架上說了很多.
不是很多 說幾句 至少七句.
說哪七句 記不記得.
十字架上所說的.
十字架上所說的七句話.
第一句.
「赦免他們 因為他們所做的 他們不曉得」.
「我實在再告訴你們 告訴被釘的那個人」.
「今天我就要跟你在樂園裡」.
「母親 看你的兒子」.
對他的門徒說 「看你的母親」.
「我的神 我的神 為何你棄我」.

$^{121}$「我學了 成了 父啊 我將我的靈魂交給你」.
當耶穌在說這些話的時候.
我們發覺 其實祂也在十字架上.
但是當祂復活了之後.
我們剛才就說 祂不斷向不同的人顯現.
剛才也有讀過這段經文.
祂向傷心失意的穆達拉瑪利亞顯現.
我們一起讀經文 好嗎 再一次.
祂說 耶穌對她說.
「婦人 你為甚麼哭 你找誰呢」.
瑪利亞以為耶穌是個園丁 就對她說.
「先生 如果是你把祂拿去的」.
「請告訴我 你把祂放在甚麼地方」.
「我好去搬回來」.
耶穌對她說 「瑪利亞 祂轉過身來」.
用希伯來話對她說 「拉波尼」.
就是老師的意思.
絕望傷心的瑪利亞.
以為耶穌的身體被人拿走了.
所以她不斷地在哭 很傷心.
很愛的人已經離世了 連屍體也沒有了.
有時候是的 當你傷心絕望的時候.
你一定會哭的.
我相信當時的瑪利亞會哭得很大聲.
但有時候哭這件事 我覺得很特別.
你一哭的時候 你聽得出它其實是真哭還是假哭.
還是傷心哭 還是不是傷心哭呢.
有沒有試過這些.
小朋友有些哭 很明顯是假的.
至少要拿奧斯卡金像獎.
我參加過一個 大家都參加過的.
微信主的喪禮.
我那次去參加的時候 哭得很厲害.
每個人都在哭.
但我聽他們哭的時候.
我不知道是真哭還是假哭.
但我也哭 為什麼呢.
因為那些香煙攻我的眼睛.
那我也哭.
大家也有可能 也不是可能.

$^{161}$也會出席過很多基督徒的喪禮.
也會哭的 弟兄姊妹.
但那種哭不是一種很絕望.
或者是很控制不了的那種.
是一種盼望 還有一份愛的哭.
是很不一樣的 雖然大家都是哭.
耶穌去問瑪利亞 你找誰.
其實祂是不是明知故問呢.
我相信祂不是 祂當然知道.
但祂也這樣問祂.
因為上帝想醫治和安慰祂.
希望將祂在那個悲傷絕望的泥沼裡拯救出來.
有時候我們也找很多東西.
就好像瑪利亞那樣去找一些東西.
在你人生裡找很多東西.
在不同的階段找很多東西.
我也找過很多東西.
但對我來說.
我覺得我最傷心的一段時間就是.
我找不到上帝.
不是耶穌的屍體.
在我心裡我找不到上帝.
因為我自己做了很多事情.
上帝不是很喜悅.
但是同一時間我又很渴望有上帝的描述和念文.
我在這裡找祂.
我嘗試在我的工作裡找祂.
我想在很多不同的地方.
甚至在教會裡找祂.
但也找不到.
但剛才經文有一個很關鍵的時刻.
耶穌做了什麼.
耶穌只說了一句話.
祂說瑪利亞.
當然你這樣看這個瑪利亞.
這裡沒有人叫瑪利亞的.
瑪利亞 沒有.
我想我們這裡沒有人叫瑪利亞.
叫瑪利亞 不算吧.
叫瑪利亞.

$^{201}$我當然這樣想.
叫瑪利亞.
經文說.
瑪利亞怎樣.
她馬上認得到耶穌.
然後她不是叫耶穌.
她叫拉玻尼.
拉玻尼.
就是那些.
你聽到的一些很.
很我們不明白.
是她們兩個之間的一些對話.
她叫瑪利亞.
我相信不是一般的叫法.
我想你也有吧.
有沒有人叫你的名字.
是那個人才會叫你的.
我記得我奶奶怎樣叫我的.
我奶奶叫我奇仔.
奇仔.
這個是我奶奶才會這樣叫我.
沒有人會這樣叫我.
因為我是中山人.
我叫鄭紹騏.
鄭紹騏在中山話來講就是鄭奇.
鄭奇.
奇仔.
她這樣叫我.
是她這樣叫我.
我奶奶不會這樣叫我.
我奶奶叫我臭小子.
我奶奶叫我奇仔.
一叫我奇仔我就知道.
我就認了.
我就會了.
我當然不會說拉玻尼.
我說奶奶.
很煩.
諸如此類.
你和上帝有一個關係的時候.

$^{241}$上帝會叫你的名字.
有時候會覺得.
可能祂沒有真的叫你的名字.
是你和祂之間有一種很特別的關係.
耶穌這個顯現.
也告訴我們有時候好像瑪利亞一樣.
很痛苦.
很困惑.
很失望.
以為神不在.
但其實.
今天上帝告訴你.
祂仍然知道你的名字.
祂分分鐘有時候可能叫你的名字.
祂甚至可能不叫你的名字.
祂性靈感動你的時候.
你就知道祂呼喚你.
復活的主仍然今天在我們心靈中可以顯現.
但這個顯現未必是你眼見得到.
但祂沒有離開過我們.
除了傷心失意的瑪利亞.
我們剛才有說.
祂有顯現過給這個失意的門徒.
這個失意門徒有兩個門徒.
在二萬五司路上一直走.
這次耶穌說了一些話.
我們一起讀.
一二三.
耶穌說.
無知的人啊.
先知所說的一切話.
你們心裡信得太遲鈍了.
基督正在受害.
然後進入他的榮耀.
不是應當的嗎.
於是他從摩西和眾先知.
把所有關於自己的經文.
都給他解釋明白了.
這是路加福音的第24章.
伍馬爾斯是一條村.

$^{281}$這條村距離耶路撒冷25公里.
當這兩個門徒一直在二萬五司的路上.
離開了耶路撒冷.
走到伍馬爾斯的路上.
他們很灰心.
為什麼.
為什麼會灰心.
因為他們無法接受和相信.
他們那位彌賽亞.
耶穌基督被釘上十字架.
甚至很殘酷.
我們都經歷過.
回於Dare.
他們很灰心失望.
我相信我們在我們的宿命當中.
或多或少會經歷一些灰心失意.
無知的人.
耶穌這次有些教訓性.
他說你無知的人.
先知所說的一切.
你心裡仍然不明白嗎.
因為這兩個信徒.
其實是懂得勝經的.
耶穌為什麼要這樣說他們無知的人.
說他們遲鈍.
信德遲鈍.
因為這兩個門徒不是沒有聽過.
不是沒有聽過先知所說的.
但他們不明白.
他們忽略了耶穌基督.
原來是要先受苦.
後得榮耀的真理.
有時候我們不是不懂聖經.
或不是沒有聽過.
但有時候我們真的不明白.
未體會到.
你今天大家活得好不好.
有沒有人覺得自己活得不好.
很苦.
那是應該的.

$^{321}$先苦後榮耀.
這就是耶穌基督在地上的生活.
但如果耶穌基督和那些門徒.
用聖經去告訴他們.
到底是什麼意思的時候.
他說從摩西和先知開始解釋給他們聽.
跟他們查經.
你會用什麼聖經跟他們講呢.
用舊約聖經你會用什麼聖經跟他們講呢.
很多.
舊約裡面講了很多耶穌基督.
由《創世經》已經開始講了.
三章十五節.
那個女人的後裔.
記不記得.
用蛇.
怎樣我踏傷你的頭.
弄傷你的腳.
他在講耶穌基督.
逾越節的羔羊.
出埃及記.
詩篇也有.
受苦的僕人.
就是耶穌基督不斷告訴我們.
其實基督來創下是要受苦的.
剛才經文也有講.
《易創雅書》第五十三章.
最出名的受苦的僕人.
耶穌受很多苦.
今天我們受多少苦呢.
而當我們覺得.
耶穌基督受苦是一個失敗的時候.
這絕對不是神的心意.
在《易創雅書》提醒我們.
很多耶穌基督受苦的事情.
所以怎樣可以讓一個人.
由灰心變成火熱一點呢.
聖經這樣講.
他們彼此說.
在路上.

$^{361}$耶穌基督跟我們講.
將聖經解釋給我們聽的時候.
我們的心裡就火熱起來了.
我曾經也很失意過.
我信主沒多久的時候.
我爺爺就過世了.
我爺爺那時候還沒信耶穌.
我信了耶穌.
他還沒信耶穌.
我很擔心很失意.
但是那時候.
上帝也用聖經的話去提醒我.
他告訴我上帝不願一人沉淪.
其實那時候我還有我祖母.
叫我乾兒子的祖母.
我說祂還在.
雖然我爺爺不在.
因為我跟我爺爺的關係是挺好的.
因為他從小到大.
我從小.
不是到大.
我還沒到大的時候.
他從小就已經開始作弄我了.
他很喜歡作弄我.
他在澳門住.
我爺爺我祖母.
他每次來香港的時候.
他很厲害的.
他在中環碼頭.
因為我住在西環.
他在中環碼頭下船.
從澳門坐船過來.
然後走上來我們那裡.
走來西環.
我住在高街的時候.
有沒有人知道我在說什麼.
沒有人知道我在說什麼.
就是港大那邊.
我住在高街.
他說要走一二三街來到高街.

$^{401}$很厲害的.
他腳步力很強.
所以我對我爺爺有很多的記憶.
他怎樣作弄我.
作弄得我很煩.
我向我爸爸投訴.
我說你爸爸搞我.
真的很多事情.
但是那時候離世了.
我很傷心.
因為我已經過來美國.
但是上帝提醒我.
我還有我媽媽.
我跟她傳福音.
有一年我真的回到澳門.
我就幫一個當地的宣教士.
幫他們辦一個青年營.
當時我也是年輕的.
當然是青年營.
去了五六天左右.
就跟我媽媽去傳福音.
雖然那次我媽媽是沒有信主的.
我臨走之前.
但是我的堂弟就信了耶穌.
就是我堂弟和她一起住的時候.
我堂弟信了耶穌.
長話短說.
後來當地的宣教士.
不斷去探望我這個媽媽.
最後我媽媽是信了主的.
她打電話來美國.
特意打來長途電話的時候.
打來美國說.
「阿奇,我信了耶穌」.
我不知道該怎麼說.
總之她是這樣的意思.
「我信了耶穌」.
「你要跟我祈禱」.
她叫我現在在電話裡跟她祈禱.
那時候很高興.

$^{441}$有一種很灰心.
很失意.
突然間那個人很火熱.
原來上帝是工作.
上帝用祂的話語.
用祂的大能去燃亮我們的生命.
以致後來我真的.
很多時候想嘗試擺上.
會在教會.
在很多地方跟人傳福音.
有這份熱忱在當中.
祂不單是向一些灰心失意的人.
祂更加向一些惶恐的門徒去顯現.
在星期日黃昏的時候.
當門徒聚集在一起的時候.
他們很害怕.
因為他們害怕像耶穌那樣.
被釘被抓.
他們怕那些猶太人.
他們關門.
但是耶穌來到的時候.
在他們中間說.
「願你們平安」.
今天我希望你們大家都有平安.
但我不知道你們每一天有沒有平安.
但耶穌基督說「願你們平安」.
在當中不單止有人徬徨.
其中有一個人叫多瑪.
你記得嗎.
這個多瑪是一個.
我們叫他多疑的多瑪.
他有很多疑心.
他怎麼說.
他說.
「把你的指頭放在這裡」.
就是耶穌見到他之後.
很多疑心的時候.
耶穌說「把你的指頭放在這裡」.
「看看我的手吧」.
「伸出你的手來」.

$^{481}$「探探我的肋旁」.
就是叫你摸摸我.
很多事情希望他能夠明白.
然後還跟他說一句話.
我們一起讀好嗎.
「不要懷疑 只要信」.
然後多瑪就跟他說.
「我主我神」.
耶穌說「你因為漢奸才信嗎」.
「那些沒有漢奸的」.
「有信的更有福」.
「是有福的」.
在這裡當中我強調了.
不要懷疑 只要信.
有時候我們也會像多瑪一樣.
總是信不下去.
但是今天上帝也知道.
告訴我們.
祂沒有放棄過我們.
祂知道我們是小信.
剛才我說我的堂弟信了耶穌.
他信了耶穌之後.
他一直長大.
當時我太信主的時候.
他應該是中學生.
小學生,中學生 很小的時候.
後來慢慢長大.
在澳門的時候.
他在賭場工作.
在澳門很多人都會在賭場工作.
他開始對他的信仰慢慢離開.
他以前也有回教會跟同學.
後來長大了就沒有了.
後來他也覺得做賭場.
也是比較混濁一點.
他轉了做警察.
接著還要做G4.
G4那些特工隊.
保護何厚華那些.
我不知道.

$^{521}$保護特區政府那些人員,官員.
就算他改行歸正.
他也覺得跟他所認識的上帝距離很遠.
他有很多懷疑.
他開始長大了.
看到人生有很多白領.
當你見到很多東西的時候.
要不就總括.
耶穌是真的.
你的世界就是這樣的.
耶穌還不是真的.
好嘛.
你可能悲觀一點.
你看到那麼多東西.
哪有上帝?.
不可能的.
你不相信.
他後來就認識了一個女孩.
接著就回到教會.
我相信那個女孩.
現在是他的太太.
那時候在信仰上幫了他很多.
一起在教會事奉很熱心.
一直都和我有時候聯絡一下.
我回去的時候也探望他.
我看到上帝在他生命當中做過很多工作.
他兩夫婦後來在教會開始做傳道人.
幫助很多當時吸毒的人.
青少年.
婦孺.
很多工作在那裡.
上年我這個堂弟.
做了傳道人之後.
上年就安納了牧師.
整件事我看到是很難以置信的.
由我爺爺過世.
我祖母信耶穌.
接著我堂弟他自己做了傳道人.
現在做了牧師.
很多懷疑.

$^{561}$我記得他那時候很多時候跟我說.
很多懷疑他不信.
但是耶穌復活這件事.
是叫我們不可以不相信.
所以我們看到.
其實耶穌基督向誰顯現?.
一直我們看到.
祂不是向那些宗教領袖.
祂不是向那些有權有勢的.
祂是向那些傷心,軟弱,懷疑的人.
今天我們是不是都是這一種的人?.
如果你是這一種人的時候.
你不要只覺得很對不起主.
其實主是知道的.
不要懷疑只要信.
我相信這句話耶穌說的時候是很磁場的.
祂不是要責備多瑪.
當你很多懷疑很多不信的時候.
你聽到耶穌基督跟你說不要懷疑只要信.
這個這麼磁場的說話的時候.
我們都會降伏在心裡面.
我們可不可以好像耶穌基督那樣.
向那些傷心的人,失意的人,軟弱的人.
都顯現一下我們還沒死.
但我們可以繼續完成耶穌基督的工作.
我不知道今天耶穌是不是要顯現給今天的人.
但今天有很多人都需要我們去.
將耶穌基督這一份復活的力量跟我們分享.
大家都記得這些卡了嗎?.
教完了嗎?交了功課了.
挑戰大不大?.
我覺得這裡有很多顯現.
有幾個我對我自己都很深刻.
他說問一個朋友最近最大的挑戰是什麼.
那我問我老婆.
聽到嗎?你最近最大的挑戰是什麼?.
答案很簡單.
不要發脾氣.
後來還有一個挑戰是.
你只要說我支持你.

$^{601}$不要給建議.
我說我支持你.
那就是我要做好一點.
不要讓他發脾氣.
我覺得挺好的.
我們每一天都為主做一些事.
有一個就是先停十秒.
三次,每天有三次先停十秒.
煮一下你覺得怎麼樣.
那是不是整個早餐都挺好吃,煮一下你覺得怎麼樣.
很有趣,你每天都跟上帝給個機會大家去建立關係.
請人喝咖啡.
留下語音給別人,告訴他們.
我想為你禱告.
有些經文叫你信給別人.
這些任務卡我覺得我們不要就這樣停了.
我們會不會繼續有第二部分呢?我們要再談一下.
但是我們不要停,因為耶穌基督的復活是沒有停過的.
祂的顯現是沒有停過的.
祂的顯現是透過我們每一個人一起去完成這件事情的.
好嗎?.
最後一個問題.
如果今天耶穌基督真的向你顯現的時候.
你覺得祂的樣子是怎麼樣的?.
我覺得祂應該戴著一個冠冕.
這個叫十誡的冠冕.
這個十誡的冠冕是祂的冠冕.
但是這個是有詞的冠冕.
我們一會用一首律詩歌回應就叫做十誡的冠冕.
今天無論我們有多難這條路.
我們知道上帝這十誡的冠冕已經為我們承受了一切.
我們一起祈禱好嗎?.
因為當我們很多時候去到人生傷心失意軟弱的時候.
你總不離開我們.
主人我們甚至會懷疑你.
即使你在十字架做了這麼多的事情.
我們有時候回想我們平日的生活.
我們好像真的糟蹋了主你所做的一切.
你不斷向世人告訴我們.
你沒有離開過我們.

$^{641}$願意主這個十誡的冠冕不只是你帶.
讓我們每個人都帶著這個冠冕.
去接觸上帝你所愛的每一位有需要的人.
祈禱奉主耶穌明教.
阿們.
\newpage



\section{約翰福音 19:38-42}
\label{sec:031919}
\textbf{明唔明,唔明就黎明}~~~~林家棋宣教師
\newline
\newline
連結: \href{http://www.church.com.hk/acms/content.asp?site=cdc\&op=show\&type=product\&code=031919\&layout=sermon}{\texttt{031919}} ~~~~ 語音日期: 2026-03-28
\newline
\newline
\hyperref[sec:031918]{\small{< < < PREV SERMON < < <}}
~
\hyperref[sec:index_chronic]{\small{[返順時目]}}
~
\hyperref[sec:031920]{\small{> > > NEXT SERMON > > >}}
\newline
\newline
約翰福音 19:38-42
\newline
\begin{longtable}{cl}
\hline
\hline
章節 & 經文 (和合本修訂版)\\
\hline
19:38 & \begin{tabularx}{0.7\textwidth}{X} 這些事以後，亞利馬太的約瑟來求彼拉多，要把耶穌的身體領去。他是耶穌的門徒，只因怕猶太人，就暗地裡作門徒。彼拉多准許了，他就把耶穌的身體領走。 \end{tabularx} \\ \\ \relax
19:39 & \begin{tabularx}{0.7\textwidth}{X} 尼哥德慕也來了，就是先前夜裡去見耶穌的那位，他帶著約一百斤的沒藥和沉香。 \end{tabularx} \\ \\ \relax
19:40 & \begin{tabularx}{0.7\textwidth}{X} 他們照猶太人喪葬的規矩，用細麻布加上香料，把耶穌的身體裹好了。 \end{tabularx} \\ \\ \relax
19:41 & \begin{tabularx}{0.7\textwidth}{X} 在耶穌釘十字架的地方有一個園子，園子裡有一座新墓穴，是從來沒有葬過人的。 \end{tabularx} \\ \\ \relax
19:42 & \begin{tabularx}{0.7\textwidth}{X} 因為那天是猶太人的預備日，而那墳墓又在附近，他們就把耶穌安放在那裡。 \end{tabularx} \\ \\
[1ex]
\hline
\hline
\end{longtable}
$^{1}$我不會解釋今天的講題,好像一個講題一樣,大家明就明.
在今天的經文裡面,我們剛才趙孫為我們讀出這段經文的時候.
我們發現有兩個人物.
第一個人物就是阿彌馬太的約瑟.
第二個人物就是之前已經和耶穌討論過人重生的問題.
那位尼哥德姆.
尼哥德姆我們相對比較熟悉.
反而阿彌馬太的約瑟我們較為陌生.
今天我們只集中去看阿彌馬太的約瑟這個人物.
當中他的經歷作為我們今天要效法的對象.
以及當中他的經歷如何成為我們一些鼓勵.
好,邀請大家去打開今天的經文.
約翰福音的十九章三十八至四十二節.
邀請大家打開,也邀請大家再一次去讀一讀這段經文.
去代入一下這件事情.
你當是親身去見證這個敘事.
你是第一身人的時候.
你是有些什麼,看到些什麼和有些什麼感受.
約翰福音十九章三十八至四十二節.
好,打開的第一節目我想邀請你們.
你們一路讀,你們嘗試去代入.
你就是其中的一個見證人.
好,預備,一二三.
喪葬的規矩.
用細麻布加上香料.
喪葬的規矩.
用細麻布加上香料.
喪葬的規矩.
關於亞利馬泰的約瑟這個人物.
其實聖經裡面所記載的其實不多.
但是在四卷的福音書裡面.
其實竟然同時間都有記載這個人物.
而且所記載的事件都是關於.
他主動向彼拉多去求耶穌的身體.
這位亞利馬泰的約瑟究竟是一位怎樣的人物呢.
綜合四卷福音書的記載.
他是一個財主.
是一個尊貴的議員.
一個盼望神國的人.
而且在這個路加福音.

$^{41}$是更加形容描述他是一個善良正直.
不會攀附他們的珍貴的一個很重要的官員.
最重要最重要的是.
他是耶穌的門徒.
我們從而可以掌握到其實這位的約瑟.
他是一個有錢.
有身份.
有地位.
而且正直的一個猶太人.
不過我們去比較四卷福音書的描述.
唯獨馬泰和約翰這兩卷福音書.
提及他是耶穌的門徒.
而今天我們去看約翰福音.
更加只是將他的焦點.
集中他作為耶穌這個門徒的身份之中.
我們先去看一看十九章.
三十八節整個經文裡的結構.
我們看到豹邦已經打了給大家看.
我亦都highlight了兩個很重要的動詞.
經文其實用了兩個分詞來描述.
這位阿里馬泰的約瑟是一個怎樣的人.
第一個分詞就是他是.
他是誰呢.
他是耶穌的門徒.
第二個分詞描述再多一點.
再補充這個門徒是怎樣的門徒.
而是一個一直隱密地去做門徒的一個人.
然後才再正常地去直接說.
為什麼這個約瑟他去做一個隱密的門徒.
一直藏起自己的門徒.
原因其實就是害怕猶太人.
當然約瑟不是唯一一個作為隱藏了的門徒.
在約翰福音的十二章四十二節.
亦都提到耶穌傳道期間.
其實有不少的官員都信了耶穌.
不過礙於法利賽人的緣故.
他們怕被這班人去說出猶太的會堂.
而不敢承認自己是信耶穌.
約瑟就是其中一個隱藏了的門徒.
來到這裡,我們不要急著定論.

$^{81}$這個約瑟很軟弱.
為什麼這個約瑟會怕猶太人.
而不敢承認自己是信耶穌.
但是當我們切身處理去想.
當時承認耶穌是我的老師.
是什麼一回事.
當我們去了解切身處理去想的時候.
去親身經歷的時候.
我們未必會這麼輕易去下這個約瑟.
去說這個時間他做一個隱密的門徒.
是一個軟弱的門徒.
我們不會這麼快下這個結論.
剛才提到約瑟就是其中一個身份.
就是一個尊貴的猶太的議員.
而約翰福音十二章四十二節.
其實也提到只要承認你是跟隨耶穌的.
你就會被人趕出猶太會堂.
一個身份尊貴的猶太人.
被人趕出猶太會堂.
不是吧,很害怕.
再者,耶穌被釘十字架的罪名是什麼.
耶穌被釘十字架的罪名.
就是叛亂整個羅馬政府.
鼓動一群民眾反抗羅馬帝國.
對於身為猶太人的約瑟.
嘩,我跟隨的老師被羅馬政府釘了這個罪名.
我還不害怕.
害怕也很正常.
客觀來看.
約瑟因為外圍環境因素.
而被迫隱藏自己的門徒身份.
他看似軟弱.
但我們想想.
跟隨耶穌的十一個門徒.
其實也不遑多讓.
那十二個門徒包含猶太人.
也是不遑多讓.
猶太人出賣耶穌.
給了三次不認主.
當耶穌被捕.

$^{121}$受審被釘十字架的門徒去了哪裡?.
門徒四散.
我們不能只斷言.
哎呀,約瑟做隱藏門徒.
他怕猶太人斷定他是軟弱.
反而我們更加看到的一件事.
他沒有像那十一個門徒一樣.
四散.
他沒有離開過耶穌.
我們看見當耶穌斷氣死了之後.
他是唯一一個.
走去彼拉多面前.
去求耶穌的身體.
將耶穌的遺體拿走去安葬.
對比其他十一個門徒.
約瑟其實這個行動.
是由那一刻開始已經由暗轉為明.
真正地展現當一刻.
他自己是耶穌真正的門徒.
我們嘗試從約瑟去求耶穌身體這個記載.
來理解約瑟這個轉變.
是什麼一回事呢?.
十二章三十八節提到約瑟來求彼拉多.
要將耶穌的身體領去.
領去這個字本身用了一個假設的語氣.
即是耶穌前去進見彼拉多.
去求耶穌身體的時候.
其實他不知道自己行不行.
彼拉多給不給他.
根據當時羅馬的律法.
任何的死囚.
任何的死刑犯他死了之後.
他們的親人可以向羅馬政府領取他們的屍體.
但是唯獨被判處叛國罪的死刑犯.
任何人都不可以拿.
包括他的親人都不可以拿回他的屍首.
當我們都知道耶穌.
他被羅馬政府冠以反抗羅馬帝國的罪名.
被釘上十字架.
按照羅馬的法律.

$^{161}$任何人都不可以拿走耶穌的屍體.
更莫說約瑟和耶穌母親無故.
沒有任何血緣關係.
更加不符合可以拿走死刑犯的屍體的資格.
所以當約瑟前去見彼拉多的時候.
他都是盡他的能力嘗試去爭取這個機會.
讓他盡自己最大的力量.
能夠好好安葬這位老師.
加上他是一個猶太人的身份.
公然去拿走耶穌的身體.
明顯他將自己門徒的身份公開.
在馬可福音的記載.
他描述約瑟是放膽地進見彼拉多.
對於這麼公然去執行這件事情.
其實會為約瑟帶來危險.
不單單猶太同胞.
很大機會會把他逐出猶太會堂.
甚至他去進見彼拉多的時候.
他去求耶穌的身體.
彼拉多可能會將他視為耶穌的一黨.
將他一樣釘上十字架.
我們不清楚約瑟為何會由暗轉為明.
經文沒有交代過約瑟那種轉變的原因.
又或者在約瑟眼中.
他真的視耶穌是真真正正猶太人的王.
所以他放膽地公然去領耶穌的身體.
無論什麼原因也好.
約瑟甘願冒上這個性命的風險.
從一個害怕猶太人而隱藏自己門徒身份的人.
直到耶穌死後.
放膽公然進見彼拉多去求耶穌的身體.
對比其餘十一個門徒.
約瑟他在那一刻展現什麼叫做跟隨耶穌.
弟兄姊妹我們要明白約瑟所身處的環境.
要公開承認自己是耶穌門徒其實一點都不容易.
因為隨時都會受到同胞的排擠.
甚至會為他自己帶來性命的危險.
所以他將自己作為門徒身份隱藏.
我們不多不少其實都可以理解.
其實時至2026年的今天.

$^{201}$很多我們叫做創啟地區的地方.
都不容許基督徒存在.
很多宣教士都在當地都要隱姓埋名.
隱藏他們是基督徒的身份.
更加更加不可以讓人知道.
更加不可以不明目張膽地.
讓人知道他們在傳福音.
公然傳福音在當地來說.
隨時都會為自己帶來危險.
輕則當地政府會將你驅逐出境.
重則真的會帶來性命的危險.
很多宣教士都會用其他的身份去履行他的福音工作.
另一方面在創啟地區的福音對象都是一樣.
當他們願意離開自己所信的宗教.
歸信基督的時候.
他們都不能夠很公開地跟我身邊的人.
跟我家人說我是基督徒我已經信了耶穌.
因為他們所身處的地區.
是處於一個基督信仰封閉的地方.
不容許他們改變他們原本.
一出生以來的一些宗教.
難道面對這一班人又或者宣教士.
我們能夠斷言去指斥這班貴主者.
很軟弱這麼沒有信心嗎.
難道我們真的可以這樣指斥他們嗎.
在這些環境其實是迫使他們要暗地裡去作門徒.
當然列舉以上的例子不是要為弱勢去護航.
而是很多時候外圍環境因素.
迫使一些信徒要去隱藏自己的身份.
但隱藏不代表放棄信仰.
丟下信仰.
對弱勢而言.
他的確是因為害怕猶太人.
而將他作為耶穌門徒的身份去隱藏.
但在耶穌被處決之後.
正值就是一個最黑暗和充滿危險的時間.
一班明明跟隨基督的門徒.
四散了.
但弱勢這個隱藏的門徒.
仍然放膽走去彼拉多面前.

$^{241}$我可不可以求你給我耶穌的身體.
我幫祂安葬.
他不再刻意隱藏自己的身份.
最後他沒有丟下自己是耶穌門徒.
他擁有的身份.
弟兄姊妹.
我們呢.
即使面對現在的環境可能很艱難.
但我們仍未至於像弱勢那段時間.
是個嚴苛的時間.
至少我們仍然可以如常去敬拜.
如常去跟我身邊的人討論信仰.
我們不是弱勢.
也不是身處在創啟地區的宣教士.
我們不需要刻意去隱藏.
我們自己是基督徒的身份.
又或者對大家來說.
我都沒有去隱藏這個身份.
不過.
弟兄姊妹.
我們有沒有試過.
因為周遭的環境.
所帶來的壓力.
令到我們這個身份.
好像被隱藏了.
黃明道一生為主堅守.
但在極大的政治壓力底下.
他都試過跌倒.
在1945年左右的時間.
對於一個剛剛成立的新政府.
所推行的宗教政策.
他其實有些保留.
他認為.
那一刻他認為.
上主才是世間至高無上的權威.
政治和宗教應該要分立.
信徒應當要信服神的緣故.
才去信服地上的政府.
前提.
是當時的掌權者.

$^{281}$必須所發出的指令.
不違反,不違背神的旨意.
以及兩者不可以衝突.
他強調信徒應該是信服神的.
但到了三自運動展開.
黃明道因為拒絕加入.
連同他的妻子一同被捕.
起初黃明道.
他仍然為神堅持真理.
不願意屈服.
但是當他被捕之後.
在獄中受到不同程度的壓迫.
以致他對於可能被槍斃.
他最終選擇屈服.
並且簽紙願意表示加入三自教會.
他簽紙之後.
他獲得釋放.
但對於黃明道來說.
獲釋是一種折磨.
獄身雖然自由.
但他無時無刻都受到良心的責備.
有時他會說自己三次不認主.
有時他會說自己是猶大.
出賣了耶穌.
在他獲釋之後.
他選擇馬上脫離三自教會.
隱姓埋名.
但卻是埋下第二次被捕的種子.
當他第二次坐牢的時候.
他覺得自己很大機會被判處無期徒刑.
當時他為了讓自己的刑罰能夠減輕.
判得輕.
他在招供的時候.
他說了很多的謊話.
但結果仍然判他無期徒刑.
黃明道在那一刻他很憤怒.
但又感到很悲觀失望.
在那一刻他覺得.
難道我真的要在監獄等死嗎.
他問神.

$^{321}$神啊 你為什麼這麼殘忍.
你叫我遭遇這些打擊 被判無期徒刑.
他說完這番話.
那一刻他突然出現了尼加書七章七到九節的經文.
裡面說到.
至於我 我要仰望耶和華.
等候救我的上帝.
我的上帝必應運我.
我的仇敵不要向我誇耀.
我雖跌倒 仍要站起來.
雖坐在黑暗裡.
耶和華卻作我的光.
我要承受耶和華的惱怒.
直到祂為我辯護.
為我伸冤.
因我得罪了祂.
祂要令我進入光明.
我必得見祂的公義.
因為神這番話令他醒悟.
他說謊源於他害怕.
他害怕死亡.
將原本沒有做過的事.
全部都招認自己做了.
也因為害怕.
他選擇向人屈服.
而招致上帝的憤怒.
得罪了神.
這次的提醒.
上帝這番話.
突然在他腦海出現.
令他對神的不滿煙消雲散.
他再次向神禱告認罪.
禱告過後.
他心裡充滿喜樂.
以及從神而來赦罪的平安.
這次的經歷.
令黃明道裡面的靈性.
得以更新和復甦.
對於黃明道來說.
被判無期徒刑.

$^{361}$是他最黑暗的時間.
但他和約瑟一樣.
在最黑暗的一刻.
他選擇仰望神.
讓他能夠在黑暗之中.
去看見光.
約瑟以及一眾跟隨耶穌的人.
耶穌的死亡是正值一種充滿悲傷.
惶恐不安.
不知所措.
以及黑暗的時間.
那一刻耶穌的死亡.
對於他們來說.
連唯一的盼望都失去.
但約翰福音隨即記載.
在這個黑暗的時間.
好像罪惡超張.
看不見光.
但光卻是在耶穌復活的時候出現.
我們看見.
約瑟以及尼哥德穆.
他將耶穌安葬在一個新的墓穴裡面.
墓穴本身就是為了死人而預備的地方.
但耶穌卻是將這個新的墓穴.
從這個新的墓穴裡面復活.
這個代表死亡的象徵的墓穴.
卻是成為復活盼望的起點.
將黑暗變為光明.
約翰福音一章四到五節裡面講到.
在他裡面有生命.
這生命就是人的光.
光照在黑暗裡.
黑暗卻沒有成過光.
弟子們.
耶穌基督回來之前.
這個世界仍然處於黑暗.
因為這個世界仍然被罪所污染.
而且只會越催敗壞.
對我們來說.
我們可能還沒看見真正的光.

$^{401}$但我們怎樣自處.
特別當我們還可以這麼公開地.
放膽地敬拜我們的上主.
我們所身處的地方.
又不是一些像創啟地區一樣.
又或者弱勢時代的環境.
我們仍然沒有受到很大的壓迫.
我們明明可以作門徒.
但我們卻是選擇暗暗地作門徒.
將我們基督徒的身份從明轉為暗.
我們會不會因為一些壓力.
將我們基督徒的身份隱藏.
又或者其實壓力還沒來到.
我們自己先行收起基督徒的身份.
今日神呼召我們.
就要像弱勢一樣由暗轉為明.
這個轉變不是我處處都要大聲去報道.
大聲去說我是基督徒.
而是在這個黑暗的環境.
我們敢不敢去為真理發聲.
敢不敢去愛那些被社會排斥.
被欺壓的人.
特別面對一些社會不公不義的事情.
我們又會不會敢於去發聲.
敢於去關心受欺壓的人.
弟兄姊妹.
弱勢的行動其實是鼓勵我們.
就算面對黑暗看似看不到盼望的環境.
弱勢卻是從暗轉為明.
公開為耶穌站台.
祂知道站台的後果是什麼.
但祂不再隱藏.
而是跟隨耶穌到底.
在這一刻對於我們來說.
耶穌雖然暫時離開這個世界.
但我們也要記得.
耶穌仍然在掌管著.
就算罪惡.
惡人好像獅子一樣.
張牙舞爪.

$^{441}$但是到了終末耶穌回來的時間.
一切的罪惡必然會被消滅.
黑暗不再存在.
最後很想引用我很喜歡的一套電影對白.
去鼓勵可能正處於黑暗.
看不見盼望的你們.
黎明之前的夜晚是最為黑暗.
但我保證黎明一定會.
這句不單單只是一句電影對白.
而是一個應許.
因為這個黎明就是基督.
今日如果你們正處於一個黑暗.
害怕軟弱跌倒的時間.
請記住黎明一定會來.
因為復活的主已經勝過黑暗.
希望我們能向主認罪.
求主賜給我們像虐室一樣的膽量.
從暗轉為明.
讓我們同心同意.
參考上帝.
在今日的聖周六裡.
正正就是耶穌基督被釘十字架死亡.
對於當時的門徒來說.
可能是一個很黑暗很盼望的時間.
但轉眼間耶穌卻復活.
為世人帶來光明盼望.
對我們來說.
可能我們處於不同的黑暗.
想躲藏自己門徒身份的時間.
但主卻告訴我們.
透過虐室的故事告訴我們.
他如何為耶穌站台.
他如何勇敢放膽地在彼拉多面前冒著危險.
去奪取耶穌的身體.
他願意公開自己.
在這個黑暗帶來危險的時代.
他願意公開自己跟隨基督.
當我們真的軟弱.
當我們真的不能發聲.
我們很想去隱藏自己的時候.

$^{481}$主要是求你讓我們有這個膽量.
求你讓我們有這個膽量.
就好像虐室一樣.
放膽地去關心這個世界.
放膽地去關心我們身邊的人.
放膽地去承認我們是基督的門徒.
當我們軟弱的時候.
求主幫助我們.
讓我們在這個黑暗世代真真正正宣告.
耶穌基督就是世界的光.
耶穌基督就是我們的神.
求你讓我們有這個膽量.
以上禱告奉祖耶穌基督聖名祈求.
\newpage



\section{出埃及記 8:19}
\label{sec:031920}
\textbf{聖經人聖經事:這是神的手段}~~~~黃克勤牧師博士
\newline
\newline
連結: \href{http://www.church.com.hk/acms/content.asp?site=cdc\&op=show\&type=product\&code=031920\&layout=sermon}{\texttt{031920}} ~~~~ 語音日期: 2026-04-02
\newline
\newline
\hyperref[sec:031919]{\small{< < < PREV SERMON < < <}}
~
\hyperref[sec:index_chronic]{\small{[返順時目]}}
~
\hyperref[sec:031921]{\small{> > > NEXT SERMON > > >}}
\newline
\newline
出埃及記 8:19
\newline
\begin{longtable}{cl}
\hline
\hline
章節 & 經文 (和合本修訂版)\\
\hline
8:19 & \begin{tabularx}{0.7\textwidth}{X} 術士對法老說：「這是神的手指。」法老心裡剛硬，不聽摩西和亞倫，正如耶和華所說的。 \end{tabularx} \\ \\
[1ex]
\hline
\hline
\end{longtable}
$^{1}$各位聽聽節目,大家平安!.
《生命傳書》的網絡.
每天為大家預備讀聖經的一段內容.
繼續現在都是用聖經人物和聖經故事.
和大家去認識背後的上帝.
哇!祂對奇妙的手很厲害.
這次這個題目真是直接了.
「這是神的手段」.
大家要知道.
我們習慣了我們到今天的中文.
就是「手段」.
有點貶義.
「祂很會使用手段」.
覺得貶義.
「哇!使用手段很奸」.
其實你要知道.
「手段」這個詞應該可以說是.
犯法或者神的處理力量.
叫做「手段」.
不要馬上一見到手段就覺得很差.
其實祂的智異手段本身不是這樣的.
「手段」即是「懂得怎樣做事」.
「用手怎樣做事」.
而且有些人說「亂做啊」.
不是的.
有段的一段一段一段.
甚麼意思?是一步一步一步.
這就是神的手段.
我們生命樹做宣教.
不妨多說兩句才進入聖經.
我們生命樹做宣教.
是栽培當地的宣教工場.
當地的人.
去完成他們的宣教使命.
即是說柬埔寨人,印度人,斯里蘭卡人.
緬甸人等等世界各地的人.
他們信耶穌.
是誰的第一個責任呢?.
當然是當地的人.
我們在香港.

$^{41}$柬埔寨人信耶穌,印度人信耶穌.
我香港信徒的責任是不會的.
第一個責任當然是當地的信徒.
他們要完成這個宣教使命.
我們香港信徒做甚麼?.
我們能夠做的就是幫助他們.
完成他們的宣教使命.
於是無端端我會在袋裡拿錢出來.
支持我們做事.
無端端這麼辛苦.
我就會研究神學.
然後教書.
所以我們是用教學.
來提升他們的能力.
能夠做更好的本地宣教.
栽培更多本地有遠見的領袖.
有能力的領袖.
正如聖經所說.
要將你所學的教導那些中心能教導別人的人.
這就是寶萊的宣教戰略.
教導那些中心能教導別人的人.
這就是工作的手段.
所以生命樹的宣教手段是甚麼呢?.
就是栽培當地中心能教導別人.
訓練別人,帶領教會的人.
所以這個手段是好東西.
看看神的手段.
「這是神的手段」.
嘩!很開心.
大家要知道.
《春一及記》第七章十四至十一章第八節.
是講十災的.
信主久一點就懂了.
十災.
這十個災是甚麼呢?.
大家一看就知道了.
先找一支螢光筆畫一下.
其實這麼多個災.
有些災的名字比較長.
有些人喜歡變成四個字.

$^{81}$有些變成兩個字.
有些變成三個字.
我覺得不要緊.
長一點就短一點.
第一就是「水變做血」.
就是聖哥利諾河的水變做血.
幸好我自己有個缸裝起水.
都變了.
所以在埃及地.
那段時間是沒有水喝的.
而且臭得不得了.
水變做血.
聖經就講埃及人在河附近挖地.
希望河水滲進來.
經過泥土變成乾淨的水可以喝.
聖經沒有講到他們挖出來的是水還是血.
不過照說在家裡的缸裡面.
瓶裡面的水都變成血.
應該挖出來的水都是血.
那段時間臭得不得了.
而且熱得不得了.
真是死人的.
不吃東西一段長時間都撐得住.
沒有水喝就死了.
除非你活生生喝那些魚臭了的血.
這個災很厲害很大.
全部的災是做什麼呢?.
都是打擊埃及人拜的神.
農業社會一定要拜河.
尼羅河.
沒有了它一定不能耕種.
好,我把它變成血.
挖出來的青蛙他們也會拜.
青蛙有什麼好處呢?.
青蛙厲害的地方是出水能跳入水能游.
而且很有彈跳力.
因此身體表現很強壯.
他們就說不如拜牛.
拜牛不會下水.
沒有青蛙這麼厲害.

$^{121}$又沒有繁殖這麼多.
兩妻一起去拜青蛙.
那次蛙災弄得全屋都是.
法老王連睡的床都是青蛙.
他們也拜青蛙.
弄得很臭.
後來神就說算了.
停這個災.
青蛙屍體堆積如山.
臭得不得了.
第三就是塞災.
這個字我們通常讀「早」.
不過原來是讀「塞」.
塞就是上面這一部分.
讀「早」也可以.
這個真的要留意.
塞災就是神令那些塵土.
那麼多的塵土變成塞.
我想腳真的鋪滿塞塞.
我也不知道當年埃及人怎樣過.
又沒有毛髀膏.
怎麼辦?.
痛苦又癢又痛又咬一咬.
大家說是說.
我們經常出門去那些艱難地方.
如果有蚊子咬.
打死就打死.
趕走牠之後很癢.
癢是因為蚊子會放牠的口水.
到你的身體裡面,令到血液會停.
然後它又知道很幼的針在那裡吸血出來.
不停就塞了一支針.
原來蚊子的針是有一組針的.
有些是放牠的口水.
有些是吸血的.
於是你把牠拔走了.
就有個很微的傷口.
哎呀,好癢呀.
由於牠放了口水下來.
身體覺得有點異物.

$^{161}$於是各方支援.
一方有難,八方支援.
很多液體就會衝去那裡.
嘗試緩解災情和侵襲.
於是很多液體就令到傷口腫了一塊.
而且癢是一個敏感過敏的反應.
於是你就會咬.
指甲有很多菌.
而且會咬傷了更大的傷口.
於是就會發炎.
所以蚊子咬完之後.
蚊爛的地方就會一點紅色.
而且紅了很久.
因為它發炎進入了裡面.
因此如果被蚊子咬了.
咬了就算了.
忍,不要咬.
那要做什麼呢?.
塗藥膏.
通常有很多藥膏.
都會有少少不傷害人的類固醇.
或者什麼.
幫你去止癢.
或者有些自然物質幫你止癢.
你就不要咬.
忍,忍忍忍.
否則整個腳都會發散.
好了,營災.
又拜蒼蠅.
就是因為蒼蠅.
農業社會很需要人去耕種.
越多人越好.
多人做工.
南南去去都不驅.
所以就看到蒼蠅.
嘩,生產力真是厲害.
這麼多蒼蠅.
烏蠅蟲多得不得了.
很厲害.
於是他們就拜蒼蠅.

$^{201}$蓄液之災就是將瘟疫降在生畜.
很多生畜很開心.
誰知有些生畜.
脫毛病就很心痛.
原來有一種菌會感染羊.
間中都會發覺.
令到羊毛很長,脫毛很難看.
慢慢就站不起來.
協調不了身體.
走一步會跌,很有趣.
那就死了.
有些菌會攻擊牛.
我們俗稱的蜂牛症.
牛會肌瘤受損.
跌在地上會死.
病毒或菌會跑進人體裡.
第六就是生瘡.
嘩,渾身生瘡.
站不起來,拉不起來,蹲不起來.
又不能蹲起來,哎呀,很慘.
全部都是你自己說不出話來.
藥膏都買不錢.
薄災開始是天氣變了.
不是下雪.
是雨水下的時候.
那些水滴聚在一起.
或者還沒聚在一起.
但已經凍了.
凍一凍在天空裡結成冰.
於是嘩,掉下來.
有些薄是很大的.
好像一個乒乓球.
一個高爾夫球那麼大的.
可以大到.
那些是打爛的.
車的玻璃.
打爛了車頂,凹了.
打到人都死了.
原因就是高空雜物掉下去.
黃災呢.

$^{241}$黃春一過.
唉,他們食糧很厲害.
所有綠色的東西都被他們吃光了.
田沒有了,沒有收成,黃災.
那如果你看看第八.
去這個蒼埃及記看到黃災.
結果過了時間.
好了,就收回那個災.
是叫風.
就是吹過來.
西風一吹.
全部整個埃及裡面的蝗蟲.
是沒有留下一隻.
全部吹進去西拉抗疫.
那就嘩.
那這個埃及人見到.
埃及法老見到會不會.
嘩.
真是厲害.
如果你是埃及法老.
有這樣的災發生.
如果我在政治上說.
唉,我們不夠他鬥.
他的神真是厲害.
我們讓他走.
你一點都沒有失威.
因為他的神實在太厲害了.
我們就我不想這個國家.
我們受損太多.
所以我們放他們走了.
可能這樣都可以了.
那不會有人上來反你的了.
那如果今天不行.
今天要轉移視線.
那就要殺那些以色列人.
好了,王在.
特別厲害的第九.
和第十就不得了了.
這兩個真是很厲害.
要留意到.

$^{281}$這兩個發揮.
黑暗之災.
和殺長子之災.
這個以能力的彰顯來說.
這兩個是最大的.
因為王,博.
那些東西全部都是在地球表現.
就算這個博,災,樂,博.
都在地球的大氣層,天空.
水面血就是河.
那些蟲,那些濕,凝那些.
在山莊就是身體.
但是黑暗之災.
是天空的太陽.
去了另一個層級.
太陽.
埃及人拜不拜太陽的呢?.
拜到極了.
原因就是.
太陽是他們生命之源.
沒有太陽怎麼種植?.
必須要有太陽.
所以他們很需要太陽.
拜太陽.
其實神在這麼多個災.
除了第十之外.
都是攻擊.
大部份都是攻擊他們的.
假神的信仰.
太陽是神.
不說我是不當我是東西.
好,我認識你登三日.
有三日黑暗之災.
很厲害.
這個黑暗之災裡面.
神是特別做一些事情的.
就殺長子之災.
一家人只是長子死.
這樣都挺有趣的.
這個很難做到的.

$^{321}$好了.
十個災.
真是厲害.
那今天呢.
今天的題目是什麼呢?.
「謝是神的手段」.
看看神做些什麼.
十個災.
好了,再下去了.
請你留意.
在第三個災.
失災.
在這個失災裡面.
就有一句說話.
第七章開始.
然後就去到第十一章.
七八九十十一.
五章聖經.
如果五章聖經分十個災.
大概兩至三個災一章聖經.
有些描述的.
失災是去到第八章.
就是這裡了.
第八章十九節.
竟然出了一句說話.
這句說話就是我今天要跟大家說的.
「謝是神的手段」.
是出埃及記八章十九節.
「謝是神的手段」.
嘩.
上帝真是滿有大能.
真的不要得罪祂.
或者祂要做的事祂一定做到.
如果不是.
不會發生的.
哎呀,為什麼這麼多惡人在這個世界?.
神讓祂在這裡.
或者是激勵你.
要你知道什麼是善惡.
或者是.

$^{361}$祂不好.
但是神仍然存留祂一段小時間.
給祂有悔改的機會.
哎呀,那些惡人怎會悔改?.
已經壞到極了,怎會悔改?.
不要等了,殺死祂.
一見到祂,一側祂的腳.
用指頭一按就沒有了.
如果一按,按死這個惡人.
第一個先按死我,先按死你.
好了.
「謝是神的手段」.
「手段」這個字.
是叫做S-Bar.
這個很難讀,因為有個T.
S-Bar.
S-Bar的意思就是指頭.
Finger.
所以你可以譯作「謝是神的手指」.
「謝是神的手指」.
「神的手段」.
這個「神的指頭」.
Finger.
如果大家有宇哲晚餐的習慣.
你會發覺其中有一個環節就是.
講耶和華的「指頭」.
是會點一些葡萄汁.
不是葡萄汁,是葡萄酒.
點一些葡萄酒,一點紅色.
在碟子上,最好是白碟.
點一點,然後讀.
第一個災,水變血之災,點一點.
第二個災,點一些紅酒.
第三個災,塞災.
讀完十個災之後.
就見到有十點在碟子上.
紅色的點.
是「謝是神的指頭」.
猖狂的八章十九節就說「謝是神的指頭」.
好了,留意這裡.

$^{401}$結束了.
今天每天一月三點.
主要就是一點.
「謝是神的指頭」.
「At Bar」.
這個詩篇八篇第三節.
「和你所陳設的月亮星宿」.
就說,哎呀,人算甚麼.
不要覺得自己這麼厲害.
「你竟顧念他世人算甚麼」.
「你竟眷顧他」.
「你指頭所陳設的」.
就好像神拿著一個燈泡.
擺上去.
啊,月亮,小小的.
拿個大一點的.
哦,太陽燈一盞.
還有滿天星.
其實這個陳設佈置是甚麼呢?.
我再說,應該很多時候我們都要留意.
聖經說「創世記何時創造」.
「創造的日子是何時」.
「為甚麼會這樣,為甚麼會這樣」.
「這個同進化論有沒有問題」.
「地質學又怎樣」.
為甚麼聖經不說?.
聖經為甚麼要說?.
其實神是告訴你.
「起初我裝修我的房子」.
「被上手弄得亂了,爛了,滿了,髒了,頭上有水」.
「哎呀,亂了,甚麼都沒有,光都沒有」.
「我就開始裝修我的房子」.
那是誰裝修呢?.
是神的靈開始裝修.
其實如果你看其他經文,希伯來書等等.
是聖子耶穌在聖父的旨意下.
聖子耶穌和聖靈開始創造諸世界.
其實如果你把它簡單一點想.
就是神告訴你.
「起初我裝修好我這間被弄亂的房子」.

$^{441}$「因為我想再有人住在裡面」.
「我先裝修好它,然後才做人」.
那怎麼裝修呢?.
「我先做了空氣和天上的水,地下的水分開」.
「然後有海洋,有陸地」.
「然後有第三天有青草那些植物在地裡生出來」.
這有很多講的字眼,我講錯了.
「然後我就做小鳥,做魚,做貓,做動物」.
「然後最後做人」.
「那做了多久啊?神啊」.
「做了六天」.
「六天?是不是六個二十四小時?」.
「是」 「啊?你有沒有可能啊?」.
「講出來有沒有人信啊?」.
「我需要有人信嗎?」.
「神說我是我做的」.
「我不信」.
「不信就不要信」.
「啊?這樣也行?」.
「那你不用顧一下我的感受了」.
「不用,因為我是我做的」.
好了.
「那你裝修這間房子做了多久啊?」.
「都說六天了」.
「那第七天呢?」.
「第七天我已經不用做了」.
「坐下來慶祝,欣賞我做完這間房子」.
「登基為王了,做了六天了」.
「所有人知道自己會登基的皇帝」.
「會不會一爸爸一死」.
「父皇一死,馬上自己馬上做啊?」.
「那些謬團三位就會了」.
「但是其實都距離登基是一段時間」.
「是預備好了那些新的國員,內閣成員」.
「各樣東西委任了」.
「然後才開始登基做王啊」.
「起碼都有時間做好自己的皇冠啊」.
「自己住的地方都做好了」.
「然後呢?做了多久啊?」.
「做了六十天」.

$^{481}$「成長了多久啊?」.
「六天」.
「那這間房子大不大啊?」.
「都不是很大,地球那麼大」.
「嘩!我裝修這間房子幾百呎」.
「已經搞了三個月了」.
「那即是什麼啊?」.
「即是你的工人不行了」.
「我不是說他不行」.
「即是說沒有上帝的工人那麼厲害」.
「那上帝的工人做得那麼快」.
「在六天內做好整個宇宙萬物?」.
「當然啦,誰來做?」.
「聖子耶穌和聖靈啦」.
「是不是比你厲害的那些呢?」.
「哦,是的」.
「所以其實創世記者說」.
「神用祂的手段」.
「用祂的指頭」.
「就做好這個世界了」.
「所以對著埃及法老不讓走啊」.
「讓我的指頭動一動吧」.
「我動了十次」.
「或者不動啊」.
「十次就搞定了」.
「哇,這是神的手段」.
「這是神的指頭」.
「我們要做什麼呢?」.
「我們要做的是觀寒你的指頭」.
「哇,厲害啊」.
「人算什麼?」.
「你警告其他人」.
「人算什麼?」.
「謙卑」.
「觀寒」.
「感受到神的眷顧」.
「哇,太好了」.
「這個就是我們要做的了」.
「如果實際上學到什麼呢?」.
「觀寒上帝的指頭」.

$^{521}$「然後,哎呀」.
「我算得什麼?」.
「你竟然會顧念我們」.
「顧念你的百姓」.
「非常感恩」.
「今天和大家分享到這裡」.
「希望你多些觀寒上帝」.
「識觀寒上帝的指頭」.
「自己謙卑」.
「哎呀,我算得什麼?」.
「然後依靠祂的眷顧」.
「同心祈禱吧」.
「多謝我們的主」.
「讓我們在這段經文裡」.
「再次學到怎樣去跟隨你」.
「就是你自己叫我們」.
「觀寒你所做的事」.
「觀寒神怎樣叫一個」.
「無口財路的摩西」.
「去做這個工作」.
「怎樣觀寒你」.
「以色列人是被動的」.
「但是卻能夠做出這麼多事情」.
「他們是受恩惠眷顧的」.
「算不得什麼?」.
「主,我們都學習觀寒」.
「我們學習謙卑」.
「我們學習依靠你的眷顧」.
「多謝你,讚美你」.
「口供耶穌明求」.
「阿們」.
\newpage



\section{馬太福音 28:1-10}
\label{sec:031921}
\textbf{空墳墓的訊息}~~~~陳沛偉牧師
\newline
\newline
連結: \href{http://www.church.com.hk/acms/content.asp?site=cdc\&op=show\&type=product\&code=031921\&layout=sermon}{\texttt{031921}} ~~~~ 語音日期: 2026-04-05
\newline
\newline
\hyperref[sec:031920]{\small{< < < PREV SERMON < < <}}
~
\hyperref[sec:index_chronic]{\small{[返順時目]}}
~
\hyperref[sec:031922]{\small{> > > NEXT SERMON > > >}}
\newline
\newline
馬太福音 28:1-10
\newline
\begin{longtable}{cl}
\hline
\hline
章節 & 經文 (和合本修訂版)\\
\hline
28:1 & \begin{tabularx}{0.7\textwidth}{X} 安息日過後，七日的第一日，天快亮的時候，抹大拉的馬利亞和另一個馬利亞來看墳墓。 \end{tabularx} \\ \\ \relax
28:2 & \begin{tabularx}{0.7\textwidth}{X} 忽然，地大震動；因為有主的一個使者從天上下來，把石頭滾開，坐在上面。 \end{tabularx} \\ \\ \relax
28:3 & \begin{tabularx}{0.7\textwidth}{X} 他的相貌如同閃電，衣服潔白如雪。 \end{tabularx} \\ \\ \relax
28:4 & \begin{tabularx}{0.7\textwidth}{X} 看守的人嚇得渾身顫抖，甚至和死人一樣。 \end{tabularx} \\ \\ \relax
28:5 & \begin{tabularx}{0.7\textwidth}{X} 天使回應婦女說：「不要害怕！我知道你們是尋找那釘十字架的耶穌。 \end{tabularx} \\ \\ \relax
28:6 & \begin{tabularx}{0.7\textwidth}{X} 他不在這裡，照他所說的，他已經復活了。你們來！看看安放他的地方。 \end{tabularx} \\ \\ \relax
28:7 & \begin{tabularx}{0.7\textwidth}{X} 快去告訴他的門徒，說他已從死人中復活了，並且要比你們先到加利利去，在那裡你們會看見他。看哪！我已經告訴你們了。」 \end{tabularx} \\ \\ \relax
28:8 & \begin{tabularx}{0.7\textwidth}{X} 婦女們急忙離開墳墓，又害怕，又大為歡喜，跑去告訴他的門徒。 \end{tabularx} \\ \\ \relax
28:9 & \begin{tabularx}{0.7\textwidth}{X} 忽然，耶穌迎上她們，說：「平安！」她們就上前抱住他的腳拜他。 \end{tabularx} \\ \\ \relax
28:10 & \begin{tabularx}{0.7\textwidth}{X} 耶穌對她們說：「不要害怕！你們去告訴我的弟兄，叫他們往加利利去，在那裡會見到我。」 \end{tabularx} \\ \\
[1ex]
\hline
\hline
\end{longtable}
$^{1}$頂子妹平安!復活節快樂!.
耶穌基督的復活是我們信仰裡的一個很重要的核心.
耶穌基督勝過死亡,是永活的神.
因著耶穌基督的復活.
叫我們一群信祂的門徒.
我們一群屬神的子民.
能夠在生命裡有活潑的盼望.
因此,復活節是一個非常快樂的時刻.
所以我很想在這個時候邀請大家.
跟旁邊的弟兄姊妹說一聲.
復活節快樂!好嗎?.
復活節快樂!.
我記得在多年前.
當我去英國旅行的時候.
我去大英博物館參觀.
其中有一個展館是特別多人的.
是一個木乃伊的展館.
去看的人絡繹不絕,不斷有人龍.
因為裡面擺放了很多具木乃伊.
有些是有很長的歷史.
上千年的歷史.
保存得很好.
有些甚至是還有頭髮的.
令人感覺到很驚訝的是.
為什麼一個血肉之軀.
能夠可以好好保存這麼久呢?.
古時的防腐技術.
竟然是這麼精湛.
確實是令人嘆為觀止.
當我參觀的時候.
除了感覺到真的大開眼界之外.
我心裡面同時有一份感嘆.
我在想.
這些木乃伊.
在他們生前的時候.
無論多麼風光.
在這一刻擺在眼前的.
只有死亡.
只有死亡的終局.
只有工人鑑賞的死亡.

$^{41}$今天我想和大家一起去重溫.
聖經裡面一段非常重要的敘事的記載.
這段經文是記載了耶穌基督復活的情景.
耶穌基督復活是一個鐵一般的事實.
不過並不是讓我們來研究.
讓我們去鑑賞.
而是要向我們.
在今天仍然向我們發出.
很重要生命的訊息.
我們怎樣來回應呢?.
所以我盼望藉著這段經文.
能夠給我們有亮光.
給我們有啟發.
首先讓我們一起去讀出這段經文.
在《馬太福音》28章1至10節.
我們一起讀.
預備起.
「安息日將盡.
七日的頭一日.
天快亮的時候.
沒大拉丁瑪利亞和那哥瑪利亞來看墳墓.
忽然地帶震動.
因為有主的使者從天上下來.
把石頭滾開坐在上面.
他的像貌如同閃電.
衣服潔白如雪.
看守的人就因他嚇得渾身亂顫.
甚至和死人一樣.
天使對婦女說.
不要害怕.
我知道你們是尋找那釘十字架的耶穌.
他不在這裡.
照他所說的已經復活了.
你們來看安放主的地方.
快去告訴他的門徒說.
他從死裡復活了.
並且在你們已先往加利利去.
在那裡你們要見他.
看啊.
我已經告訴你們了.

$^{81}$婦女們就急忙離開墳墓.
又害怕又大大的歡喜.
跑去要報給他的門徒.
忽然.
耶穌遇見他們說.
願你們平安.
他們就上前抱著他的腳拜他.
耶穌對他們說.
不要害怕.
你們去告訴我的弟兄.
叫他們往加利利去.
在那裡必見我.
我們再一次同心祈禱.
親愛的恩主.
在復活節的早晨.
我們來到你面前.
去再一次多謝你.
多謝你十字架的恩情.
多謝你的復活.
以致我們能夠恩信得生.
在這一刻我們一起去思想神的話.
盼望聖靈去預備好我們的心靈.
啟迪我們的心.
讓我們能夠聽得明白.
得著激勵.
讓我們能夠在每天生活裡面.
去榮耀你.
多謝聽我們的祈禱.
奉耶穌基督的名.
阿們.
在我們剛才所讀的經文的上文.
第二十七章的後半段.
記載了祭司長和花里賽人.
他們因為恐怕耶穌基督的門徒.
來偷耶穌的遺體.
接著就說耶穌復活了.
因此他們就去找彼拉多來協助.
希望他能夠派一些人去看守墳墓.
經文記載.
彼拉多就派了一隊兵丁.

$^{121}$和他們一起去.
並且吩咐他們.
要盡你們的所能把守妥當.
於是他們就一起去.
去到墳墓用封條封好墳墓.
嚴嚴地去把守.
這個高度設防的景象.
和我們剛才二十八章所讀的經文的開始.
是一個非常大的對比.
兩個柔弱的婦女.
莫大拉的瑪麗亞和另一個的瑪麗亞.
她們在天未亮的時候.
就是耶穌基督被釘的第三天的清晨.
仍然天朦朧的時候.
她們來到墳墓前面.
希望能夠為耶穌補上一些香膏.
這個強弱懸殊的對比.
似乎正要帶出將要發生.
一個非常震撼的事情.
非常震撼的一幕.
經文記載.
忽然地帶震動.
主的使者從天上下來.
把石頭滾開.
坐在上面.
看守的兵兵嚇得不停地顫抖.
甚至跟死人一樣.
可能已經暈了.
一個天使對比一隊的兵兵.
顯出神的能力.
其實在天使移開石頭之前.
主耶穌基督已經復活了.
祂已經不在墳墓裡面了.
墳墓已經是空的了.
面對前所未見的景象.
婦女們非常害怕.
天使於是就安慰她們.
並且宣告在第五至第六節.
天使就對婦女們說.
不要害怕.

$^{161}$我知道你們是尋找那釘十字架的耶穌.
祂不在這裡.
照祂所說已經復活了.
天使這個極震撼的宣告.
對於婦女來說.
是一個最好,最佳的安慰.
沒有任何一句說話比這句說話更加安慰.
因為她們所愛的耶穌已經不在了.
已經復活了.
耶穌曾經三次對她們的門徒預言.
祂將要受難.
並且在第三天復活.
耶穌基督的復活是按著祂的預言.
按著聖經的應許來應驗的.
是神究竟計劃的一部分.
是必然成就的.
耶穌基督復活的大能.
遠超一切其他任何的權勢.
死不能夠將祂拘禁.
婦女們又驚又喜.
她們急不及待地.
要將這個大好的消息向門徒報告.
同時當風聲傳到祭司長和長老那裡的時候.
祭司長和長老就很擔心.
他們設法扭曲這件事實.
經文記載他們就藏起了冰釘.
要他們告訴別人.
是在他們睡覺的時候.
耶穌的門徒偷了他們的身體.
不過.
祭司和長老的詭計.
並不能夠改變耶穌基督復活的事實.
在約翰福音裡面的記載很有意思.
約翰特別記載.
包裹起耶穌身體的西馬布的細節.
和包裹著耶穌的頭的果頭根.
是怎樣捲在一起.
放在另外一個地方.
有些聖經學者甚至認為.
那個果頭根應該還有人的形象.

$^{201}$有人的形在那裡.
我們大家試想一下.
如果耶穌的身體真的被偷了的話.
偷的人.
哪裡能夠慢條斯理地.
放好這些西馬布和果頭根呢?.
事實上耶穌基督復活的事實.
除了有空墳墓的證明之外.
還有很多的見證人.
就好像今天經文裡面所記載的婦女.
婦女在當時的社會裡面是不受重視的.
她們的見證是不能夠成為正式的呈堂證供的.
去到法庭那裡.
是不會當是一個正式的見證的.
所以如果要有一個更加有力的見證.
應該找一些有名望的人.
更有說服力的人.
來做這個第一見證人的.
不過經文卻是記載.
正正就是這兩位的婦女.
她們就做了這個第一的見證人.
原因很簡單.
就是她們真的在那一刻.
在現場去看著整個的情景的發生.
她們並不是刻意找回來的見證人.
除此之外.
在幾本福音書當中都記載了.
耶穌基督復活之後.
向門徒和向不同的人來顯現.
保羅就在《哥倫波前書》十五章裡面.
這樣總括去記述.
在十五章的三至八節.
「我當日所領受又傳給你們的.
第一就是基督照聖經所說.
為我們的罪死了.
而且埋葬了.
又照聖經所說.
第三日復活了.
並且顯給基法漢.
然後顯給十二使徒漢.

$^{241}$後來一時顯給五百多弟兄漢.
其中一大半到如今還在.
卻也有已經睡了的.
以後顯給雅各漢.
再顯給眾使徒漢.
沒了也顯給我漢」.
保羅在寫《哥倫波前書》的時候.
相信是耶穌基督復活之後.
二十多年的時間.
保羅說.
還有幾百位的目擊證人.
他們仍然在生.
如果有需要的話.
他們可以隨時站出來做見證.
有些不相信耶穌基督復活的人.
或者會這樣說.
可能是因為門徒.
或者群眾這班婦女太傷心了.
驚慌過度.
所以產生了幻覺.
他們誤會了耶穌真的復活了.
這個說法其實是站不住腳的.
要眾使徒甚至更多的門徒.
一起有這個幻覺.
其實是不可能的.
相反.
使徒他們在耶穌基督復活之後.
他們生命的改變.
卻是有目共睹的.
他們以往是膽怯退縮的.
如今他們卻是勇敢地為主作見證.
聖經裡面記載.
當耶穌基督被捉拿的時候.
門徒都離他而去.
當耶穌基督復活之後.
到《使徒行傳》裡面的記載.
我們見到這班門徒.
他們是勇敢和堅定.
去為主作見證.
雖然為福音的緣故被毆打.

$^{281}$但他們卻是心裡歡喜.
因為他們覺得.
他們深信.
他們算是配授被主命受辱的.
門徒生命的改變.
是耶穌基督復活.
一個很重要有力的證據.
弟兄姊妹.
耶穌基督的復活是千真萬確的事實.
這個這麼重要的事實.
除了是一件歷史事件.
除了是我們信仰裡面.
我們教義裡面一個極重要的根基之外.
在今天對我們要發出什麼訊息呢?.
我們又如何以生命來回應呢?.
在經文裡面描述婦女的反應.
相信給我們有很重要的啟發.
當耶穌基督斷氣的一刻.
圍觀的人都垂著胸.
很傷心地離去.
但是有一些婦女.
相信包括.
穆迪拉瑪利亞和另一位瑪利亞在當中.
她們卻是遠遠地站著.
她們不捨得走.
她們不想走.
她們很傷心.
她們想看守著耶穌基督的身體.
來到今天經文.
在清晨的時候.
她們來到墳墓的前面.
相信她們仍然是帶著這份悲傷的心情.
不過當她們見證了空墳墓的這一幕.
聽見天使的宣告之後.
她們的心靈裡面有了很大的轉變.
在經文第八節裡面這樣說.
「婦女們就急忙離開墳墓.
又害怕又大大的歡喜.
跑去報給他的門徒」.
雖然婦女們因為從來沒見過這樣的情景.

$^{321}$心裡面感覺到害怕.
也感覺到敬畏.
但同時間.
她們的心卻大大的歡喜.
這是極大的轉變.
她們原先非常沉痛的心.
如今卻是豁然開朗.
因為她們知道.
主耶穌基督的死並不是終局.
主耶穌已經勝過了死亡.
為人帶來了活潑的盼望.
這份盼望是超過人所能夠測度的.
這份盼望是能夠轉化哀傷無望的心靈.
是能夠讓人找到出路的.
因此她們急不及待.
要將這個好消息去傳開.
耶穌基督的死亡.
在她們心目中 在她們眼中.
原先是一個終局.
沒有了 完結了.
不過這一刻.
呈現在她們眼前的景象.
空焚墓的景象.
天使的宣告.
卻是遠遠超過了她們原先自己的意念.
她們所能夠想到的.
丙姐們有時可能在我們的生命裡面.
也有類似的情況.
當我們面對一些似乎.
沒有出路的困局.
當我們面對一些人生極大的低谷的時候.
有時我們的眼睛會變得迷糊.
特別是我們屬靈的目光.
會變得模糊.
或者這些困局.
長時間的在我們生命裡面.
可能我們因此一段很長的時間.
都沒有試過大大的歡喜.
但是這次是新的.
復活節 主耶穌基督的復活.

$^{361}$是要喚醒我們.
主已經勝過了死亡和一切的權勢.
黑暗已經過去了.
我們已經進入了復活的清晨.
盼望的亮光.
從高天淋到我們的身上.
叫我們能夠憑著信.
可以跨過這些種種的人生的困局.
生命的更新,改變是有可能的.
因為耶穌基督已經復活了.
若有人在基督裡.
他就是新造的人.
舊事已經過去了.
都變成新的.
在前幾個星期.
公佈了全球幸福指數的報告.
每一年都公佈的.
今年有147個國家排名.
香港排第幾呢?.
香港排第90.
我無意將香港和其他各個國家來對比.
會不會我們比起一些落後的國家.
可能都比不上.
我覺得這樣意思不是很大.
因為在計算分數的時候.
總是有不同的著重點.
每一個國家都總是有他自己的處境.
不過總的來說.
我們排第90.
我相信都不是一個值得自豪的事情.
放進去我們的人生.
弟兄姊妹.
我們思想一下.
我們的幸福指數又是建基於什麼呢?.
有什麼是能夠叫我們大大的歡喜呢?.
耶穌基督的復活.
能夠怎樣去更新.
去改變我們的心靈.
去轉化我們的心呢?.
保羅在《多倫多前書》15章裡面.

$^{401}$他這樣說.
基督若沒有復活.
你們的信便是徒然.
你們仍在罪裡.
如果基督沒有復活.
我們的信就徒然了.
我們仍然在罪裡面.
被罪困了.
又或者我們反過來說.
既然耶穌基督已經復活.
我們的信就必不徒然.
罪是很可怕的.
甚至是摧毀性的.
不過如今耶穌基督已經勝過罪的權勢.
勝過黑暗的權勢.
因此我們憑信.
罪就不能夠再克制我們的生命.
也不能夠再蒙蔽我們的眼睛.
我們已經從死亡的黑暗裡面.
步向了充滿生機的神希的當中.
耶穌基督的復活.
也為我們的生命帶來全新.
賦予全新的定義.
全新的意義.
死的獨勾.
罪的權勢已經大敗而回了.
我們靠主的生命.
我們對主這一份的信靠.
是必不徒然的.
保羅對耶穌基督復活的大能.
作了以下的總結.
所以.
我親愛的弟兄們.
你們務要堅固不可搖動.
常常竭力多做主公.
因為知道你們的奴婦.
在主裡面是不會徒然的.
務要堅固不可搖動.
不是因為我們怎樣可以.
勉勵去支撐下去.

$^{441}$而是因為耶穌基督已經復活了.
耶穌基督復活的大能.
能夠轉化我們的生命.
能夠釋放我們.
弟兄姊妹.
或者我們可能都經歷過一些不容易的情況.
可能有一些.
看來好像沒有辦法解決的困難.
大到沒有辦法處理的事情.
一直壓在我們的生命裡面.
又或者有一些.
不能言喻的痛楚.
在我們心裡面一直揮之不去.
我想說的是.
弟兄姊妹們請記得.
耶穌基督復活的大能.
是我們可以倚仗的.
絕望不是終局.
絕望在信靠耶穌的人生命裡面.
是不應該有位置的.
生命的更新.
生命的突破.
是可以的.
主耶穌基督的能力.
釋放了我們.
罪不能夠再核制我們.
不能夠再壓倒性地.
來核制,困鎖我們.
主耶穌的恩慈.
祂的能力.
祂會引導我們.
走在豐盛喜樂,平安的路上.
在今天經文裡面.
重複提到一個地方.
加利利.
首先是天使去吩咐婦女.
你要告訴門徒.
去加利利這個地方.
你將要見到我.
接著主耶穌親自去吩咐.

$^{481}$兩位瑪利亞.
叫他們告訴門徒.
你們要去加利利這個地方.
去到那裡你就會與我相遇.
事實上.
耶穌基督在和門徒吃最後晚餐的時候.
已經約定了他們.
主耶穌說.
在我復活以後.
要在你們之前.
去到加利利這個地方.
為什麼又是加利利呢?.
為什麼不是其他的地方?.
有聖經學者.
他們這樣說.
或者是因為加利利.
在比較北的地方.
比較靠近外邦的地方.
在這個地方.
來頒佈大使命.
是有很特別的意思的.
是有特別的意義的.
因為福音將要傳到萬邦.
而在聖經裡面.
如果曾經形容加利利是外邦人的加利利.
這個也有一定的理據.
不過.
我相信在當中有另一個含義.
對於門徒來說.
是別具意義的.
加利利是主耶穌最初呼召門徒的地方.
有很多.
大部分的門徒.
都是在加利利當中來去夢召的.
在加利利地區的一帶.
去跟隨耶穌基督.
走遍各城各鄉.
將福音去傳開.
聆聽耶穌基督的宣講.
見證祂施展大能.

$^{521}$去憐憫,去醫治.
一點一滴的在當中去認識主.
加利利承載著門徒和耶穌之間.
很多很珍貴的回憶.
包括了他們撇下餘忘的跟從主.
那份初心.
包括了他們在加利利海上遇到風浪.
耶穌基督平靜風浪的時候.
他們那份小信.
包括了他們見證.
耶穌基督五餅二魚的能力.
等等的經歷.
在兩個月前.
我有機會回到我的中學母校的開放日.
我很多年沒有回去.
真的很久很久沒有回去了.
當我回去的時候.
很多成長的回憶.
湧現在心頭.
經過課室.
經過走廊.
想起當時我讀男校.
跟一班男同學很頑皮的一些事情.
我們做一些很愚蠢的事情.
被老師罰等等.
很多成長的階段.
浮現在心頭裡面.
當我再回望.
一直走過的路途.
是神的恩典.
一直在人生不同的階段.
去帶領著.
去保守著.
門徒跟隨主的路途都是一樣.
跌跌碰碰.
有時在高漲當中.
有時在愚昧失腳的裡面.
弟兄姊妹.
有時都是我們人生的寫照.
有時可能我們都是這樣.

$^{561}$主耶穌呼召門徒.
回到他們夢照的地方.
回到加利利的地方.
做一次去思想.
去回想他們的本相.
他們的初心.
以及主耶穌一直伴隨著他們的恩典.
我相信主耶穌為的.
是要再一次建立起他們的生命.
堅固的生命.
在《約翰福音》二十一章裡面.
有一段很感人的記載.
耶穌在復活之後.
在一個早上.
天未亮,天朦亮的時候.
在提比利亞海的海邊.
提比利亞海又叫做加利利海.
向門徒顯現.
當時天未亮.
門徒他們打魚打了一整晚.
一點收穫都沒有.
耶穌施行神蹟.
讓他們大有收穫.
後來當門徒上岸之後.
他們感覺到驚訝的是.
耶穌竟然已經為他們預備了早餐.
有炭火,有餅,有魚.
為他們預備.
請他們過來.
「小子,我們一起吃」.
在吃完早餐之後.
耶穌三次去問彼得.
「你愛我嗎?」.
三次問他.
同時再次把羊托付給彼得.
整個過程對於門徒來說.
特別是對於彼得來說.
是刻骨銘心.
是很深的再一次的接納.
再一次的信任.

$^{601}$再一次的委派.
之後.
耶穌和他們在加利利山上再聚.
向他們頒佈了大使命.
這段經文我們很熟悉.
在《馬達福音》28章16-20節.
「十一個門徒往加利利去.
到了耶穌約定的山上.
他們見了耶穌就拜祂.
然而還有人疑惑.
耶穌進前來對他們說.
天上地下所有的權柄都賜給我了.
所以你們要去使萬民作我的門徒.
奉父子聖靈的名給他們施浸.
凡我所吩咐你們的.
都教訓他們遵守.
我就常與你們同在.
直到世界的末了」.
在熟悉的地方.
在加利利的山上.
復活的主向門徒再一次發出呼召.
再一次的將更大的使命.
來交付給他們.
主耶穌宣告.
天上地下所有的權柄都賜給我了.
耶穌基督的復活.
就是這些所有的權柄的明證.
主耶穌要門徒明白.
你們如今去領先授.
承擔這麼大這麼重要的使命.
不是靠著你們自己的能力.
其實回到加利利.
你們就知道了.
你們就知道自己的本相.
其實沒有什麼能力.
你們能夠靠的.
是復活的主的權能.
祂所賜的力量.
親愛的弟兄姊妹.
在我們跟從主的歷程裡面.

$^{641}$有時可能都跟門徒差不多.
有時我們都會有跌跌撞撞的時候.
甚至好像彼得一樣.
曾經跌倒.
甚至灰心.
我們嘗試一下回想.
我們回想我們夢照.
跟從主的時候的初心.
我們有沒有已經失掉了.
起初的愛心.
我們有沒有已經失掉了.
起初那份熱誠.
或許我們真的需要花點時間.
安靜在主的面前.
我們去回想我們的本相.
如果我們對自己誠實.
有謙卑在主前的時候.
我們就會察覺到.
其實主的恩典一直都沒有離開過我們.
無論我們在什麼光景裡面.
託付的主.
祂很希望能夠重新建立起門徒的生命.
祂也很希望重新能夠建立起我們的生命.
賦予我們剛強.
賦予我們力量.
讓我們能夠重新振作起來.
去承擔主給我們.
要託付給我們更大的使命.
每年教會的主題是「活出使命」.
我不知道你在主的面前.
你的領袖是什麼.
神給你看見什麼需要.
你的獨特是你看見的.
給你的託付.
當你去思想的時候.
你心裡面會不會有一份猶豫.
或者有一份膽怯.
覺得自己可能做不來.
承擔不來.
弟兄姊妹,真的,我們靠的是主耶穌基督賦予的大能.

$^{681}$是祂所賜給我們的能力.
當然我們都要盡我們所能.
我們都要盡心竭力的來擺上.
但是我們知道.
我們能夠承擔的都是出於神自己.
弟兄姊妹,空墳墓是一件千真萬確的歷史事實.
不過不單如此.
空墳墓也是要向我們今天發出生命的訊息.
並且是要我們來回應的.
弟兄姊妹,我們願意怎樣來領受.
怎樣來回應呢?.
耶穌基督的復活為我們帶來確切的盼望.
活潑的盼望.
耶穌基督的復活勝過死亡.
讓我們有永恆的保障.
無論下一刻發生什麼事都好.
我們心靈裡面有一份絕對的安穩.
我們有力量每天踏實地跟從主.
空墳墓也提醒我們生命的突破.
生命的更新是可以的.
困難不是壓倒性的.
沒什麼大得過主的恩典.
因著基督復活的大能.
我們因順能夠脫離黑暗.
能夠化憂愁成為大大的歡喜.
能夠放下重擔.
或許我們會有生命裡面仍然有軟弱的時候.
但請記得我們能夠去依靠主的恩典.
復活主的恩典和能力.
叫我們能夠剛強起來.
叫我們能夠有力量去承擔主給我們的託付.
叫有力量走在主所應許給我們豐盛的道路上.
讓我們同心抬頭祈禱.
親愛的恩主.
在復活的主日的時候.
我們再一次向你獻上感恩.
多謝你十價的恩情.
多謝你讓我們透過你的復活.
因著信有這份活潑的盼望.
主也求你加力量給我們.

$^{721}$或許我們有膽怯有軟弱的時候.
求你將我們的眼目轉向你.
讓我們定精在你自己的能力.
在你自己的同在上.
你是那位耳膜內來的神.
你是那位活著的神.
你在我們生命裡面帶領著我們人生的神.
願你去幫助我們.
讓我們每一天能夠帶著這份力量去榮耀你.
帶著這份力量去見證你.
讓人能夠在生命裡面看見基督.
讓人能夠看見這份盼望的珍貴.
多謝聽我們同心的祈禱.
奉耶穌基督的聖名.
阿們.
\newpage



\section{馬太福音 28:5-6}
\label{sec:031922}
\textbf{主復活帶來七大恩典 (華)}~~~~劉彤牧師
\newline
\newline
連結: \href{http://www.church.com.hk/acms/content.asp?site=cdc\&op=show\&type=product\&code=031922\&layout=sermon}{\texttt{031922}} ~~~~ 語音日期: 2026-04-05
\newline
\newline
\hyperref[sec:031921]{\small{< < < PREV SERMON < < <}}
~
\hyperref[sec:index_chronic]{\small{[返順時目]}}
~
\hyperref[sec:031923]{\small{> > > NEXT SERMON > > >}}
\newline
\newline
馬太福音 28:5-6
\newline
\begin{longtable}{cl}
\hline
\hline
章節 & 經文 (和合本修訂版)\\
\hline
28:5 & \begin{tabularx}{0.7\textwidth}{X} 天使回應婦女說：「不要害怕！我知道你們是尋找那釘十字架的耶穌。 \end{tabularx} \\ \\ \relax
28:6 & \begin{tabularx}{0.7\textwidth}{X} 他不在這裡，照他所說的，他已經復活了。你們來！看看安放他的地方。 \end{tabularx} \\ \\
[1ex]
\hline
\hline
\end{longtable}
$^{1}$第二節聖經在馬太福音第28章第五到第六節.
這裡天使對婦女說不要害怕.
我知道你們是尋找那定十字架的耶穌.
他不在這裡.
照他所說的已經復活了.
你們來看安放主的地方.
第二節今天是復活節.
首先在這裡恭賀大家.
復活節非常非常快樂.
Amen.
Hallelujah.
兩千多年前天使向那些到墳墓來尋找耶穌的婦女說.
他不在這裡.
照他所說的已經復活了.
就從這一句話開始展開了基督教歷史.
嶄新的一頁.
因為主的復活是我們整個基督教信仰不可少的支柱.
主的復活證明了我們的主耶穌.
祂果然是神的兒子.
主的復活也證明了耶穌基督能夠有勝過罪惡死亡的詮釋.
主的復活更證明了耶穌基督已經為我們人類打開一條又新又活的道路.
這樣我們可以不再活在罪惡的匣子當中.
在基督裡面我們能得到從神而來的新生命.
Amen.
主的復活是我們信仰的核心.
也是我們基督徒得勝生活的關鍵.
難怪保羅在哥林多前書十五章十四節他這樣說.
他說基督沒有復活.
我們所傳的便是枉然.
你們所信的也是枉然.
我們看見很多人今天特別在信仰的道路上軟弱無力.
最主要的原因就是不明白主復活所帶給我們的恩典.
重要的是我們只把復活的真理當作一個屬靈的知識.
說說聽聽而已.
並不覺得主的復活與我們每天的生活有什麼密切的關係.
但事實上主的復活是我們基督徒得蒙神恩典祝福的關鍵.
Amen.
所以今天早上要跟大家分享主復活帶給我們的七大恩典.
平常只講三點所以今天復活節多講一點.
第一我們看見主的復活帶給我們得救的確句.

$^{41}$跟你們說主復活帶給我們得救的確句.
跟他說.
羅馬蘇帝五章第十節保羅說.
因我們做仇敵的時候且藉著神兒子的死得與神和好.
既已和好更要因他的生得救.
我們開始做一個基督徒.
你需要一個確信.
在基督裡面你是得救的人.
Amen.
別說你說我這麼軟弱.
怎麼而且有時候常常在罪惡過犯當中.
我怎麼是可能是得救呢.
這些聖經很確切的告訴我們兩件事情.
第一藉著神兒子的死.
我們得與神和好很重要.
因為本來我們是死在罪惡過犯當中.
我們與神是隔絕的.
但神的兒子主耶穌基督來到這個世界上.
為了罪死在十字架上.
代替我們承受神對罪惡一切的刑罰.
這為什麼每次在受難節的時候.
我們一定會念以塞爾書五十三章的經文.
以塞爾書五十三章第四節到第五節這裡說.
他誠然擔當了我們的憂患.
背負我們的痛苦.
我們卻以為他受責罰被神擊打苦帶了.
第五節哪知他為我們的過犯受害.
為我們的罪孽壓傷.
因他受的刑罰怎麼樣.
我們得平安.
因他受的鞭傷我們得一致.
這裡講到耶穌是為我們的罪來到這個世界上.
請問誰是我們呢.
就是你跟我嗎.
第二點你需要憑信心看見.
耶穌在十字架上所背負的.
不但是別人的罪.
更是你的罪我的罪.
耶穌在十字架上所承擔的.
不但是別人的憂患.

$^{81}$更是你的憂患我的憂患.
而且耶穌在十字架上已經承擔了神對罪惡一切的刑罰.
而神所審判過的地方神不會再審判.
你們這為什麼當主耶穌在十字架上斷氣那一剎那.
這已經讓聖殿的脈子從上到下裂為兩半了.
表明我們與神之間本來的罪的隔絕現在已經被逐劇了.
我們的罪已經得到赦免了.
所以藉著主耶穌基督我們能夠重新與神和好.
不但藉著耶穌的死我們與神和好.
第二藉著耶穌的生藉著祂的復活我們也能夠得救.
對不對主耶穌基督復活太重要了.
因為今天如果主耶穌只是為我們死在十字架上就結束了.
那對我們來講只不過是淚灑黃花悲劇一場不是嗎.
而且第三天從死裡復活了.
主的復活證明祂有勝過一切罪惡死亡的權兵.
在基督裡面一切罪惡的鎖鏈已經被斷開了你們.
這為什麼哥林多後書第五章十七節.
我們非常熟悉這一同來年來.
為了人在基督裡他就是新造的人.
舊事已過都變成新的.
我很喜歡在這裡說是新造的人.
觀眾說今天無論你過去在多少的罪惡過犯當中.
無論你過去有多少的傷害跟軟弱.
只要你今天選擇在耶穌基督裡面.
選擇相信並接受耶穌做你個人的救助.
你在基督裡就是個新造的人.
這個新造的人是有那種起初神創造的意識在裡面.
是按照神起初的創造是神眼中看為美好的.
今天就算你說你做了基督徒有時候也在罪惡過犯當中.
但只要你每一天願意迴轉到基督的裡面.
每一天選擇在基督裡面.
你每天都可以做新造的人.
在基督裡面你要確信你是得救的.
神用重駕把你買贖回來.
你的生命是屬於神的.
沒有人能夠將你從神手中奪去的.
所以第一個主的復活帶給我們得救的確訣.
第二個主的復活帶給我們得勝的生命.
跟你們說主復活帶給我們得勝的生命.
哥林多後書第二章十四節這裡說.

$^{121}$並藉著我們在各處顯揚那因認識基督而有的香氣.
你知道保羅在寫哥林多後書的時候.
其實他是面對各樣的困境跟難處.
但是他卻說在各樣的困境當中.
他仍然能夠靠基督而跨盛.
你知道什麼叫跨盛.
跨盛在原文裡面就是走在得勝的行列當中.
今天是耶穌基督得勝的仇敵得勝的罪惡的權柄.
我們走在耶穌基督得勝的行列當中.
我們也可以在得勝的行列當中一同來跨盛.
你說這怎麼可能我們這麼軟弱.
如果你來讀哥林多後書第四章第七節.
保羅講得很清楚.
保羅說我們有這什麼寶貝放在哪裡.
瓦契裡.
保羅形容我們了像瓦契.
瓦契是什麼.
瓦契是最平凡的契名不是嗎.
最脆弱非常容易破碎的.
今天我們人常常也是這樣不是嗎.
我們覺得自己何等的軟弱.
我們感覺自己何等的平凡.
何等的容易破碎.
但是保羅說雖然我們看自己像瓦契很軟弱.
但是在我們這個瓦契裡面卻是個寶貝.
哈利路亞.
跟你們說嘿在裡面有個寶貝跟他說.
這個什麼寶貝.
保羅繼續說這個寶貝是一個出於神的能力.
如果你仔細讀哥林多後書第四章.
你就明白這裡所說的能力.
就是那叫耶穌基督從死裡復活的大能.
這個能力是那叫耶穌基督從死裡復活的能力.
是一個勝過一切黑暗詮釋的能力.
是一個勝過一切仇敵攻擊的能力.
而今天這個能力也內註在你我的裡面.
你們這為什麼保羅接下來在第四章第八第九節.
裡面說這個因著這個能力.
在四面受敵卻不被困住.
心裡作難卻不知失望.

$^{161}$遭逼迫卻不被丟棄打倒了卻不知死亡.
所以今天無論仇敵在我們身上的攻擊有多麼厲害.
無論我們所遇到的環境多麼艱難.
在基督裡面我們是可以天天向仇敵誇聖的.
你們所以你們永遠不要埋怨你現在的環境.
今天我看見很多人活在這個世界上.
常常一種受害者的心態.
什麼叫受害者的心態.
就永遠責怪別人嘛.
永遠責怪環境嘛.
都是因為這個人我才遇到這樣的光景.
都是因為這個環境我才會如此受苦.
今天在基督裡面你絕對不會是一個受害者.
在基督裡面你是個得勝者.
你們因為你的主已經從死裡復活.
他正坐在寶座上.
他的手正在工作.
你們羅馬書第八章二十八節怎麼說呢.
保羅說我們曉得萬事都怎麼樣互相效力.
叫愛神的人得益處.
就是按他旨意被造的人.
你知道在原文裡面.
其實是保羅說我們曉得神使萬事互相.
不但是萬事互相.
而是神使萬事互相效力.
換句話說這一切都是神所容許.
這一切都在神的掌握當中嘛.
神的手正在工作.
神的手正在為你開道路.
神的手正在帶領你走上得勝的道路當中.
哈利路亞.
所以今天無論你遇到生活的困境難處什麼.
你要每一天來宣告.
在基督裡面你是個得勝者.
哈利路亞.
第三主復活帶給我們神同在的應許.
跟你們說主復活帶給我們同在的應許跟他說.
約翰福音十四章十六節太寶貴了.
這一主說我要求復復.
就另外賜給你一位寶慧師.

$^{201}$叫他永遠與你們同在.
當主義是從復活升天之後.
他採選聖靈住在每一位信他的人心中.
這一節經文講到聖靈的工作兩個很重要的特點.
第一個聖靈是我們的寶慧師.
什麼叫寶慧師.
寶慧師有兩種含義.
一個就是安慰者.
一個就是輔導者.
換句話說聖靈在我們身上的工作.
就是在我們每天的生活當中帶給我們安慰跟指引跟引導.
哇 第二點這兩件事情難道不是我們在人生的過程當中.
最需要的兩件事情嗎.
多少時候當我們處在人生的困境當中.
不知道何去何從的時候.
我們多麼希望有人用他的寶貝來安慰我們.
扶持我們告訴我們下一步當如何走不是嗎.
而聖靈就是這位住在我們裡面的神.
他是我們的安慰師是我們的輔導者.
他是天天在我們身邊要賜給我們安慰引導的神.
那你說為什麼我們從來沒有經歷過呢.
那是因為從來沒有尋求過他.
多少基督徒接受了主耶穌領受了聖靈之後.
好像就把聖靈冷凍在我們生命的一個角落裡面.
我們從來不去尋求他.
也不去花時間聆聽他的聲音.
我們總喜歡依靠自己的力量.
依靠自己的經驗來生活不是嗎.
很多時候特別是中國人不到盡頭.
我們總是不會來到神的面前.
但事實上你不需要走到盡頭了.
跟你們說你不需要走到盡頭了跟他說.
今天聖靈就住在你裡面.
他是你的安慰師他是你的輔導者.
只要向他來祈求.
今天他就能成為你的安慰跟力量跟幫助.
第二我們看見聖靈永遠與我們同在.
換句話說他永遠不會離開我們的.
這是何等寶貴的應許跟安慰.
你知道我們人常常就是活在這種憂慮跟恐懼當中.

$^{241}$缺乏安全感不是嗎.
我們常常害怕萬一.
萬一我做錯了怎麼辦.
萬一我走錯了怎麼辦.
我們會失去所有的一切.
但今天做一個基督徒.
無論我們的環境是如何.
有一件事你可以確信的.
那就是無論我們人生是成功或失敗.
無論我們現在光景是剛強跟軟弱.
今天永遠不會離開我們.
他永遠與我們同在.
我常說就好像我們昨天做父母的人一樣.
我請問你今天我們做父母的人.
我請問你你兒女在被逆在軟弱.
你會因為你兒女的被逆軟弱.
你就不再愛他嗎.
不會.
因為你永遠 兒女永遠愛他嘛.
我常覺得像我們這樣的軟弱的父母.
竟然我們都會如此做.
那何況愛我們的天父呢.
他豈不更加愛我們嗎.
聖經上說我們的神.
天父他是永遠的慈愛來愛我們.
什麼叫永遠的慈愛.
那就叫永永遠遠嘛.
哈利路亞.
聖經上還宣告主的同在.
永遠不離開.
好叫你心裡的力量.
能夠得以剛強起來的.
第四.
我們看見主復活帶給我們什麼.
醫治的恩典.
給我們主復活帶給我們醫治的恩典.
你還記得在十字宣第四章當中.
彼得跟約翰他們來到聖殿要禱告.
在門口遇到一個生來是瘸腿的人.
那個瘸腿的人指望他們得到周祭.

$^{281}$彼得就對他說金和銀我都沒有.
我只把我所有的給你.
我奉拿撒那兒耶穌的名字.
叫你起來行走.
聖經上那個人就起來.
跳著走著讚美神.
哈利路亞.
我們看見彼得跟約翰.
接下來就對眾人這樣說.
十字宣第四章第十節.
你們眾人和以色列百姓都要知道.
站在你們面前這人得權益.
是因為什麼.
你們所定十字架.
神將從死裡復活.
拿上了耶穌基督的名字的緣故.
耶穌的名字太寶貴了.
在十二節裡面.
彼得跟約翰既說.
除他以外別無拯救.
因為天下人間怎麼樣.
沒有賜下別的名.
我們可以靠著得救.
這個得救在原文是得權益的意思.
我們看見今天在基督裡面.
我們有得醫治的恩典於一.
耶穌基督從死裡復活.
神讓他身上高天坐在寶座.
祂的名字超乎萬名之上的名.
祂的名字有全病.
勝過一切的罪惡與疾病.
你們.
事實上今天如果你讀四福音.
你會發現醫治.
一直是在我們所傳的.
天國的福音的內容當中.
當耶穌在地上的時候.
祂不但是傳天國的道理.
祂更是醫治病人.
你看見耶穌在地上.

$^{321}$總是千方百計用各樣的方法.
來拯救釋放.
醫治那些有需要的人群.
我非常喜歡馬太福音.
四章二十四節這裡所說的.
就是祂的名聲就傳遍了敘利亞.
那裡的人把一切害病的.
就害各樣疾病.
各樣疼痛的.
被鬼附的顛尖癱瘓都帶來了.
耶穌就怎麼樣.
就治好他們.
這裡說一切害病的人.
不是只有少數的人.
一切害病的人.
這也包括你我在其中.
這裡說是各樣的疾病.
不是只有少數的疾病.
也包括你的病我的病.
我們的主耶穌.
祂不但是靈魂的救主.
祂更是我們醫治的救主.
今天無論在你的身體上.
有任何的疾病.
你都可以帶到主耶穌面前來祈求.
你要來相信祂.
因著主耶穌基督的復活的名字.
你的疾病要得到痊癒的.
Amen.
Hallelujah.
第二點你要知道.
信心是一個選擇.
信心不是一種感覺.
信心是一種選擇.
有一次有一位父親.
他的兒子被鬼附了.
他把他兒子帶到門徒面前.
門徒趕不出這個鬼.
他也把他兒子帶到主耶穌面前.
他對耶穌說什麼呢.

$^{361}$他對耶穌說.
夫子啊.
你若能做什麼.
請你幫助.
很有意思.
耶穌沒有回答說.
我能或我不能.
耶穌反問他.
你若能信.
再信的人凡事都能.
很有意思.
這位父親問的是.
你若能做.
但耶穌說.
No 你問錯問題了.
問題不是我是否能做.
而是你若能信.
再信的人凡事都能.
這個父親就說.
我信 我信.
但我信不足.
他知道他裡面沒有信心.
但他因為耶穌的話.
他應證的話.
我選擇相信.
我信 我信.
雖然我信不足.
我選擇相信.
而當他選擇相信的時候.
耶穌就醫治了他的兒子.
Amen.
同樣今天你要來相信.
復活的主耶穌基督的名字.
因為天下人間.
沒有賜下別的名.
我們可以靠著得痊癒的.
Hallelujah Amen.
第五.
主的復活帶給我們傳福音的權病.
跟你們說.

$^{401}$主的復活帶給我們傳福音的權病.
當主耶穌從死裡復活之後.
他把一個大使命賜給我們.
在馬太福音.
二十八章十九到二十節裡面.
非常熟悉的大使命的經文.
我們一同來念來.
所以你們要去.
使萬民做我的門徒.
奉父子聖名的名.
給他們施洗.
凡我所吩咐你們的.
都教訓他們遵守.
我就常與你們同在.
直到世界的末日.
這是我們非常熟悉的.
大使命的經文.
但很多時候.
我們讀這個大使命的時候.
我們常常比較專注在.
傳福音這件事情上.
卻忽略了其實.
主耶穌更是給我們.
傳福音的權病.
事實上大使命.
不是從馬太福音.
二十八章十九節開始.
是從二十八章第十八節就開始.
十八節怎麼說呢.
耶穌進前來對他說什麼.
天上地下.
所有的權病都賜給我們.
哈利路亞.
換句話說.
今天為什麼我們能夠去.
是因為主已經得勝了.
一切罪惡黑暗的權病嗎.
今天為什麼我們可以去.
因為主已將得勝的權利.
賞賜給我們了.

$^{441}$今天我看見很多的基督徒.
常常不大敢向人傳福音.
因為他們只看見.
很多的難處你知道嗎.
向有錢的人傳福音.
覺得很困難.
因為很有錢.
向貧窮的人傳福音.
也覺得很難.
因為他們太沒錢了.
向有知識的人傳福音.
覺得很困難.
因為他頭腦很大.
向沒有知識的人傳福音.
也覺得很難.
因為他什麼都不懂.
你只看見難處很難.
真的很難.
但你不要只看見人的難處.
你要明白在基督裡面.
是有權病的.
哥林多後書第十章.
第四到第五節這裡說.
我們征戰的兵器.
本不是屬血氣.
乃在神面前有能力.
可以攻破堅固的營雷.
第五節說.
將各樣的計謀.
各樣難處人認識神.
那些至高之一蓋攻破.
又將人所有的心意奪回.
使他順服耶穌基督.
Amen.
第五節這些年.
神讓我有機會.
在全世界各地傳主的福音.
我有機會向有知識的人傳福音.
我有機會向沒有知識的人傳福音.
我有機會向有錢的人傳福音.

$^{481}$我有機會向沒有錢的人傳福音.
但我越來越明白.
福音本是神的大能.
要求一切相信的.
今天做一個基督徒.
我們活在地上.
一定要有一個國度性的眼光.
不只是看見.
我們自己的需要跟環境.
你需要看見神的國度.
你需要看見.
無論這個環境.
這個世界怎麼樣的改變.
神的國是永遠長存的.
神的國度必然降臨.
神的旨意必定成就.
Amen.
神的心意仍然是叫萬人得救.
不願一人成仁.
而神今天也呼召我們.
讓我們與神同工.
一同來擴展神的國度.
Amen.
第六.
主復活帶給我們施捨的祝福.
跟著我說.
主復活帶給我們施捨的祝福.
此論現在第二十章.
第三十五節非常重要的經文.
保羅這裡說.
凡是給你們做榜樣.
教你們知道.
應當這樣勞苦扶助軟弱的人.
又一章紀念主耶穌的話.
怎麼說.
獅比獸 更為有福.
我們看見主耶穌的生命.
就是一個施捨的生命.
祂就像一粒麥子.
埋在地裡死了.

$^{521}$事情是說一粒麥子.
落在地裡死了.
就結出許多果粒來.
耶穌基督就是這樣.
為我們死了.
而且從死裡復活.
我們就看見.
祂生命結出許多的子粒.
包括我們在其中.
耶穌基督.
就是用祂自己的生命.
來做比喻.
教導我們.
獅比獸是更為有福的.
你說為什麼.
獅比獸更為有福呢.
按照人的想法.
得著才是有福不是嗎.
領受才是有福不是嗎.
但事實上一個人.
若是只有領受.
他雖然可能擁有很多.
但他永遠不會明白.
自己存在的價值.
也不會明白.
自己為什麼而活.
只有當一個人.
在付出的時候.
他才開始明白.
自己存在的意義跟價值.
所以我常說.
這是為什麼.
一個偉大的企業家.
明星運動員.
自終會成為一個.
偉大的慈善家的原因.
因為當他在得著的.
一切的時候.
這一切並不能夠帶給他.
真實的滿足跟意義.

$^{561}$只有當他在付出的時候.
他才找到.
他生命的意義跟價值.
事實上今天當你看見.
你所付出的一切.
所能夠帶給別人.
生命的改變.
甚至社區的改變的時候.
我告訴你那種喜樂.
絕對是言語無法形容的.
因為你看見.
你撒下的種子所帶來的收成.
而什麼時候神的恩典.
越多從你生命當中湧流出去.
你我就越多可以經歷.
神豐盛的恩典.
這就好像.
在五秉二虞的神蹟裡面.
主所教導我們了.
你要知道五秉二虞的神蹟.
我們非常熟悉神蹟.
那個神蹟發生的關鍵是什麼.
關鍵就在於給.
不只是那個孩子.
把他的五秉二虞給了出去.
而且是在傳遞的過程當中.
沒有一個人是把秉跟虞.
保留在自己身上.
他們每一個人.
也是繼續傳遞.
繼續給出去.
整個神蹟都是在不斷的傳遞.
不斷的給予當中發生的.
所以獅筆獸的確是更為有福.
你們.
這為什麼今天我們在崇拜中.
特別是為了國度發展基金.
收特別奉獻的原因.
國度發展基金.
是我們特別為我們跨文化的.

$^{601}$就是分堂.
他們各地建堂的需要.
所設的基金.
特別是當我們去年.
我們能夠還完建堂貸款.
這是何恩德的恩典.
也一方面一同紀念到.
世界各地的分堂.
他們建堂的需要.
讓我們慷慨的付出.
好讓他們也能夠享受.
一個屬於自己的會堂.
你們.
我們越多給神.
神必越多給我們的.
你們.
最後第七.
主復活帶給我們永生的盼望.
跟你們說.
主復活帶給我們永生的盼望.
約翰福音十一章二十五節.
非常寶貴的經文.
我們一同來念.
來 耶穌對他說.
復活在我 生命也在我.
信我人雖然死了也必復活.
對我們基督徒來講.
這是我們生命最大的盼望.
因為雖然我們信了主.
我們生命屬於神.
但是我們肉體.
仍然活在這個敗壞的世界.
我們仍然要經過.
生老病死的過程.
我們也遭遇到.
罪惡世界所帶來的痛苦.
死亡的確是我們人類.
在地上所不可避免的傷痛.
但對基督徒來講.
我們一點不需要害怕死亡.

$^{641}$因為死亡.
只不過我們是在地上.
勞苦的結束.
在基督耶穌裡面.
死的人生.
硬許我們有永遠的生命.
你們.
丁鳳姐.
這些年我在教會服事.
我主持過多少追思禮拜.
有些是不信主的人.
有些是信主的.
我就發現這兩種人的.
追思禮拜是非常不一樣的.
對那不信主的人而言.
死亡就是個結束.
你看見他們的追思禮拜.
充滿了哀傷跟無助.
但對信主人來講.
雖然在追思禮拜裡面.
有一種不捨的傷痛.
但是你可以感受到.
是一種確切的盼望.
因為知道有一天.
我們在天上.
不但要與愛我們.
那個復活的救主相見.
我們更是要.
與我們所愛的親人相見.
你們.
我還記得有一次.
我跟一位美國的牧師.
去醫院.
探望一位病重的童工.
這位童工.
事實上再過兩天.
他就過去了.
我永遠無法忘記那一天.
那位牧師.
怎麼握著那個童工的手.

$^{681}$對他說這句話.
他說.
I'll see you tomorrow morning.
我明天早上.
並會再看見你.
我們必定再相見.
何等寶貴的應許.
多麼確切的盼望.
這都是因為.
在主耶穌基督裡面.
我們得到的恩典.
耶穌已經從死裡復活.
祂的復活.
帶給我們永生的盼望.
你們.
今天我們所信的主耶穌.
祂已經從死裡復活了.
今天是個宣告得勝的日子.
我深深相信今年復活節.
也是神要帶領我們.
新的工作季節的開始.
無論.
無論你的過去的.
生命的光景是如何.
復活節就是要告訴你.
在基督裡面.
你給我一個嶄新的開始.
憑著信心來領受.
這神七大恩典吧.
主的復活帶給我們.
得救的雀距.
得勝的生命.
神同在的應許.
醫治的恩典.
福音的權柄.
施捨的祝福.
以及永生的盼望.
領受神的恩典.
你們.
你就必然能夠在.

$^{721}$每天生活當中.
靠主來得勝.
你們.
一起來禱告.
親愛的我們再一次.
來到你面前線上.
因為雖然我們在.
軟弱罪惡當中.
但耶穌你先愛了我們.
來到這個世界上.
為我們的罪死.
在死之家上.
你不但為我們死第三天.
從死裡復活了.
應著你的復活.
帶給我們生命的盼望.
應著你的復活.
帶給我們得勝的雀距.
主啊無論我們環境.
我們的生命的光景.
是如何.
今天在基督裡面.
我們可以靠著你.
天天向仇敵來跨生.
向環境來跨生.
主啊今天早上.
我們懇求你復活的大能.
復活的喜樂.
復活的生命.
再次永留在我們中間.
主啊帶領我們.
在各樣的環境當中.
像保羅一樣.
能夠向我們的環境來跨生.
謝謝主.
將一切榮耀都歸給你.
奉耶穌得勝的祈求.
阿們.
\newpage



\section{馬太福音 16:25}
\label{sec:031923}
\textbf{天國戰略}~~~~李業基牧師
\newline
\newline
連結: \href{http://www.church.com.hk/acms/content.asp?site=cdc\&op=show\&type=product\&code=031923\&layout=sermon}{\texttt{031923}} ~~~~ 語音日期: 2026-04-08
\newline
\newline
\hyperref[sec:031922]{\small{< < < PREV SERMON < < <}}
~
\hyperref[sec:index_chronic]{\small{[返順時目]}}
~
\hyperref[sec:031924]{\small{> > > NEXT SERMON > > >}}
\newline
\newline
馬太福音 16:25
\newline
\begin{longtable}{cl}
\hline
\hline
章節 & 經文 (和合本修訂版)\\
\hline
16:25 & \begin{tabularx}{0.7\textwidth}{X} 因為凡要救自己生命的，要喪失生命；凡為我喪失生命的，要得著生命。 \end{tabularx} \\ \\
[1ex]
\hline
\hline
\end{longtable}
$^{1}$今日講到「在地約天 On Earth as it is in Heaven」系列的一講.
這個系列講的主要思想是.
天不是一個遙遠和未來的地方.
而是今天在地上可以體驗到的事情.
如果基督徒和天父建立了關係的時候.
在地上就可以過好像天上一樣的生活.
這個「在地約天」的概念似乎今天在基督教界頗流行.
有提倡成功神學的人主張.
不需要等到來生了.
你今天就可以有你最好的人生.
You can have your best life now.
有推動流行音樂敬拜的工作者推銷口號.
今天你來到教堂就可以做天堂敬拜.
Hugh Song 是一個頗流行的教會體系 也有他們音樂的自奉.
有一首流行的詩歌「What a beautiful name」.
他甚至假設幻想耶穌的心境是.
你不想天堂裡面沒有我們.
所以你突然將天堂帶到地上.
You didn't want Heaven without us.
So Jesus, you brought Heaven down.
但聖經是否真的講過我們今天能夠將天堂搬到地上呢.
首先講天堂這個字.
如果大家看看通用的中文聖經《逆本》.
根本就沒有天堂這個字.
天堂其實中文是借用佛教的術語.
解釋一個基督教類似的神學概念.
在《地藏經》裡面就說.
「一念天堂,一念地獄」.
但是機械借用的時候 我們就有時候要小心一點.
不要將聖經沒有的概念.
譬如說西方極樂,人間淨土都借過來.
聖經其實是清楚界別天和地的.
耶穌復活之後做了什麼.
升天嘛.
所以祂現在在天上.
但是祂直到預定的日子.
天使長的聲音和神的號角當中.
主就要親自從天降臨.
說今天可以在地上體驗到好像天上的生活.
其實是容易有時候不小心的時候產生誤解的.

$^{41}$飛機翼徒直覺的反應可能是.
喂 你們最近有沒有看新聞.
這個世界越來越像人間地獄就真了 是不是.
福音之所以對我們來生有盼望.
是因為對我們來生有盼望.
就是因為來生和今世是不同的.
因為我們得救在於這盼望.
然而看得見的盼望就不是盼望了.
誰會盼望自己所看見的呢.
如果永生只是重播.
今天我們已經在地上過過的生活.
永生就真是永不超生了.
佛家所說的那些亞避地獄無間道.
你有沒有看過那部電影的第一場呢.
第一幕就解釋給你聽什麼叫做無間道.
所以我們需要準確的去解釋這個在地約天的概念.
才不會產生一些誤會一些的混淆.
首先要由天這個Heaven這個詞說起.
創世記其實一開始就跟我們解釋.
起初神創造天地.
但是其實不是很多談天.
只是要說地 當古代人聽到地的時候.
我想他們想像的不是一個.
好像我們打個Google Earth那樣藍色的水晶球.
地是人居住的地方.
是居間 是人間.
因為神做人要在地上管理大地作皇.
用創世記的說法.
就是人要生養眾多 遍滿全地.
用後來耶穌基督吩咐門徒的說法.
大使命的說法.
人要將福音傳到地極.
所以如果用我們今天的宇宙觀.
其實地是可以涵蓋宇宙任何人類可以踏足的地方的.
如果人在火星上面建立這個殖民地.
這個都仍然是地.
希伯來聖經裡面說.
尤其是創世記一章一節裡面說.
天通常是眾數的.
因為古代近東的宇宙觀.

$^{81}$天是有很多層的.
其中一個的理解就是.
第一層就是天空.
我們叫做Sky.
有些飛鳥 有些飛機.
第二層就是外太空.
就是說例如太陽 月亮 星星居住運行的地方.
第三層以上是看不到的.
另一個說法可能我們今天就叫做屬靈的領域.
這個或者靈界.
似乎就是保羅提到第三層天.
第三層天似乎也都在他的《光環多後書》裡面描述.
有一個人上過第三層天.
他也都似乎將這個東西帶到樂園這個觀念上面.
但是參考聖經.
我們就會得到一個很明顯的結論.
天上不是完全是樂園.
啟示12章第7到第10節提到天上發生了戰爭.
米加勒和他的天使們與龍交戰.
龍和他的使者們也應戰.
可是龍不能得勝.
天上再也沒有他們的地方了.
於是那巨龍被摔了下來.
他就是那古蛇被稱為魔鬼和撒旦.
那迷惑普天下的.
他被摔在地上 他的使者們也與他一起被摔了下來.
對啟示6描寫天上的戰爭.
解經家是有不同的理解的.
有些人將天上戰爭看為伊甸園之前的事.
有些人認為這個是對應十字架和復活天上的一個情景.
也有解釋這個天上的戰爭仍然是今天正在發生的事情.
以後才有機會跟大家詳細研究啟示錄.
但至少可以肯定.
天上似乎都不是那麼平靜.
不是全部是樂園.
雖然撒旦再沒有立足掌權的地方.
但約伯記告訴我們.
撒旦仍然可以站在神的面前去指控人.
而且我們知道天和地.
今天的天和地有一天都會被毀滅.

$^{121}$希伯來主義引用詩天說.
天地都會消滅 理卻長傳.
這個理是神.
天地都會將外衣漸漸舊了.
你要將天地與內衣更換.
天地都改變了.
取而代之的是啟示錄21章中描述的新天新地.
我又看見了一個新天新地.
因為先前的天和先前的地已經過去了.
海也不再有了.
所以我們基督徒說.
信耶穌上天堂.
小心一點.
不只是因為天堂這字.
其實中文聖經找不到.
更加要清楚.
不要上錯天堂投錯胎.
將會捲起舊天舊地.
要去的是新天新地.
所以我們明白.
天上其實不只是指一個與世無爭的安樂窩.
可能是指一個戰場的時候.
在地若天 在地上如同天上.
這句話其實我們就可能有不同的理解.
這句話哪裡來的呢.
其實最直接的來源是主土民.
主土民怎麼說.
他說 所以你們要這樣禱告.
我們在天上的父.
願人都尊你的名為聖.
願你的國降臨.
願你的旨意行在地上如同行在天上.
我們日用的飲食今日賜給我們.
免我們的債如同我們免了人的債.
不叫我們陷入試探.
救我們脫離那惡者.
最後一句 通常古卷我們認為是沒有的.
天父的名在地上.
是不是已經被尊為聖呢.
天父的國是不是已經降臨呢.

$^{161}$天父的旨意是不是已經行在地上如同天上呢.
顯然不是.
這就是為什麼我們需要繼續為他禱告.
直到他成為現實.
在神學裡面我們有一個理論.
叫做已啟動的末世論.
Inaugurated Ascatology.
末世其實已經來了.
基督第一次來的時候這個末世已經開始了.
但是第二次和第一次來之間.
就是我們所處的末世了.
地上開始變成好像天上的光景.
但是還沒有完成.
所以主禱民在地上如同在天上.
所求的不是今天.
就是天下太平.
因為今天天上都還沒有太平 還沒有和平.
耶穌說 你們不要以為我來小事大給地上和平.
我來並不是帶來和平 而是刀劍.
主禱民在地上如同在天上.
所求的是天國在地上如同在天上一樣.
戰爭得到勝利.
在這場原先在天上爆發的戰爭當中.
撒旦其實在天上已經無權了.
但是撒旦仍然在地上垂死爭戰.
而且撒旦在地上仍然掌權.
仍然能夠用他剩下的權柄勒索 威逼 利誘這個世界.
所以基督徒的大使命是擴展天國直到地極.
給人人都有進天國的機會.
什麼是進天國呢.
簡而言之就是棄暗投明.
辯節 改投最終要勝利天賦的一方.
倒轉槍頭去對付魔鬼.
但是你轉到另一邊的時候是否不用打.
不是的 仍然要打的.
保羅提醒我們.
我們對付的仍然是天上的惡魔.
因為我們不是與血氣肢體搏鬥.
而是與那些統治的 掌權的 管轄者黑暗世界的.
以及天上邪惡的屬靈勢力搏鬥.

$^{201}$所以耶穌說 我的國不屬於這個世界.
天國其實不是一個地域.
不是一個地上的政權.
也當然不是一個教堂 是一個建築物.
天國是天賦掌權的社會 社群.
你有這個社群會員的身份 你就屬於天國.
所以其實是每一個屬於天國的精兵.
走上戰場之前的禱告.
主土民祈求天國的勝利.
但如果我們連打仗也不知道的時候 又怎會可以勝利呢.
所以我們對戰是很重要的.
今天似乎有很多基督徒誤解了.
以為進入天國就是進入天堂.
一個和平 安樂 與世無爭 無試煉 無磨煉的樂園.
以為我們的工作就是把教堂變成天堂.
但其實進天國只是表示我們接受投降.
便捷 加入神的軍隊.
爭戰是否會繼續 當然會.
星期五我們有 Via Dolorosa.
十價之路 一個展覽 也是一個禱告行.
來紀念耶穌受苦.
Dolorosa 拉丁文的意思其實是悲哀或者痛苦.
所以我們有時中文也叫苦路.
成功走過這個十價之路是必須經歷一些痛楚的.
但神造人的目的不是為了要折磨人.
我們的目的不是要痛楚.
其實也不是為了要追求快樂.
神造人的目的就是要我們在這個萬世的戰爭當中.
有份在這個勝利當中有份.
到最終可以配得得到我們的冠冕.
和基督系同作和.
這個和其他宗教的看法很不同.
佛教認為痛苦是人生的最終極問題.
基督教不同 痛苦和快樂只不過是必經之路.
有多少痛苦 有多少快樂.
如果得不到最終勝利的時候.
都是白費心機.
所以我們不會將注意力只放在眼前的快樂和痛苦.
今世的安樂.
如果我們只看著前面的事情.

$^{241}$眼前的事情就是本末倒置.
所以基督教因為有這個不同的思想.
我們向遠的望 我們向中的望.
很多時候有些事情用今天的世界來看.
就是有些矛盾的.
有些甚至人覺得是似非而是.
基督透過表面的失敗得到最終的勝利.
他死在十字架.
但是十字架又將死亡本身置諸死地.
這個就是一個似非而是的道理.
置諸死地而後生 舍身反得生.
這個也是天國徵給我們每一位.
需要學習第一個的天國戰略.
耶穌說 凡想救自己生命的必喪失生命.
凡為我喪失生命的必得著生命.
你怎樣在天國的戰爭當中找到你的增戰崗位.
你怎樣找到你的真正的人生目的.
你覺得你的崗位會不會是.
你未成為基督徒之前一直在做的事情.
還是神奇要你去冒險做一些以前不願意不敢做的事情.
今天我覺得有很多基督徒.
認為找自己的崗位就是發掘自己的天才.
看看自己有什麼天賦 所謂恩賜.
找到自己的崗位.
曾經有一段時間很流行的.
做一些十六型性格的測驗.
找到自己的傾向.
原來我的恩賜是什麼 就找事奉的崗位.
我問一下大家.
要寫在十字架上需要什麼恩賜.
孔孝唯一需要的恩賜就是明白.
你的恩賜其實不是很有用的智慧.
天國徵兵第一樣要明白的天國策略就是.
捨身辦得生.
最重要的不是你有什麼才幹.
最重要不是天賦給你什麼恩賜.
最重要的是你願不願意犧牲.
尤其是做教會的事情.
你要明白到地上的鬥爭.
其實是反映到天上的戰爭的時候.

$^{281}$就不能夠只依靠你自己肉身的聰明和智慧.
用一個唯物主義.
人本主義的想法去打這場仗.
打開你的屬靈眼睛.
預期魔鬼有攻擊.
但同時亦都會相信神的恩典是夠你用的.
有些教會只會向裡面看.
不願意發展新的侍功.
更加不願意派人出去建立新的新堂,新的侍功.
多多少少我覺得都是基於有限公司的假設.
我在這裡自身難保.
怎樣去幫助其他人呢?.
這就是未學識捨身反得身的天國戰略.
好像你去到戰場前面前線的時候.
神不能給你支援.
如果你不肯出人手的時候.
就永遠不會得到人手.
如果會怎樣教會是這樣想的時候.
我想就沒有聖地安祖了.
我記得蕭世超牧師籌備建立聖地安祖的時候.
有一個特別的異象分享會.
有一位來自苗彼得國語的頂就站起來.
我們不能夠一直派人出去.
我們剩下的那些人怎麼辦?.
我還很清楚地記得.
他有些酸溜溜的感覺.
因為那時候苗彼得國語.
隔開十幾年都是二十丁人左右.
我都有些同情.
好像搞壞了又走 搞壞了又走.
但我記得那天我站起來說.
會怎樣的傳統.
我們永遠是將我們最優秀的派出去的.
We always send the best and the brightest.
尤其是很多卡普組.
我那時候也沒有想過.
加入這個best and brightest.
今天我們開始這個英文崇拜.
可能大家都會質疑.
喂 我們的設施.

$^{321}$我們的教務同工.
我們的義工的人手夠不夠.
何來有力再去谷一個新的侍工出來呢.
你看看都是有點困難的.
我們現在下會點每個週末.
有十個崇拜.
平均五個是二十人左右.
確實是有點缺的.
不過保羅在那個時候後書提醒我們.
我什麼時候軟弱.
什麼時候就剛強了.
寫生犯得生.
這個就是天國精兵.
第一個要學習的天國戰略.
我今天還有兩個.
似非而是的天國戰略.
想和大家分享.
第二個天國的戰略.
偉人是僕人.
馬方第十章四十二到四十三節.
耶穌說.
你們知道外邦人有君王作主治理他們.
有大臣操權管轄他們.
但是在你們中間不可這樣.
你們中間誰願為大就要作你們的用人.
在你們中間誰願為首就要作眾人的僕人.
世界常規的方法總是把權柄給做大的領袖.
所以即使是耶穌的門徒.
都是以贏得領袖的位置.
明爭暗鬥.
有很多教會都是.
那些衝突其實都是權力鬥爭.
很多時候我認識一些新的信徒.
進來教會看到教會衝突就很奇怪.
有些新信徒就問我.
教會你們賺多少錢.
我就說如果大家都十一奉獻的.
教會的收入應該可以媲美一些街坊生意.
塘水舖之類.
那些新信徒就會很奇怪.

$^{361}$他說為了個麥當勞營利都不夠的生意.
搞到臉紅耳熱.
值不值.
難怪有時候外面的人都以為教會的人發神經.
我當然不認為教會的價值.
可以用我們的收入來衡量.
一個靈魂的永生價值.
沒有辦法用金錢來量度.
所以如果我們真的竭力爭辯.
是為了真理.
就算有時候有些爭吵.
我覺得衝突是值得的.
但當日門徒不是為真理而爭辯.
他們自己是為自己爭取地位.
他們以為有位置就有話事權.
你說話人就要聽.
你就能夠有力量做大事.
耶穌怎樣回答.
說到位置這件事完全是神安排.
神沒有選你.
爭取都沒有用.
更加重要的是明白什麼.
要做天國偉人.
位置上未必一定是領袖.
但一定有僕人的心腸.
世界的想法總是.
越多人聽從你的意願.
你就越大力 你就越有效能.
你就越能夠偉大.
但耶穌說最偉大最有效的.
其實是屬於照顧到.
最多人意願的公僕.
有一次我和一群堂會領袖開會.
那個堂會正在掙扎當中.
又是那些廿人座那些.
我問其中一位領袖.
你需要什麼.
我可以做什麼來幫你.
他說最好你派兩個組的人.
來長期幫忙 長住在這裡.

$^{401}$他的意思是.
派兩組人聽他指揮.
幫他忙.
我問派他來做什麼.
我聽了就知道.
他要那些人做一些.
他們一直都在做.
但是效果又不是很好.
而且很多人因為這件事.
正在做一些離開的事情.
我就提出一個橫架.
不如我找一群專家.
他在這裡的.
都會長期.
不過不是你叫他做事.
是相反.
是他告訴你.
你需要改變什麼.
你需要做什麼好不好.
我想你也猜到他有什麼反應.
習思廣益 僕人領袖這個策略.
其實很多時候不是這個世界的想法.
有些社群就不太合適.
那天陳子明老師來的時候.
我想一般人也在這裡.
你也聽到.
如果來過那些.
如果你做黑社會社團大佬的時候.
習思廣益 僕人領袖.
當然不靈驗.
是不是.
所以教會靈不靈驗.
就看你現在是不是真的教會.
如果真教會就靈驗.
不是真教會的時候就不靈驗了.
但如果這個是一個屬天的社群.
有聖靈保守.
帶領眾人的心的時候.
這個僕人領袖的心態.
就是事功要搞得成功.

$^{441}$事功要搞得大.
最重要的秘訣.
每個人的聲音都能夠聽到.
那個領袖也都是願意聽人的聲音.
基督教第三個.
似非而是的天國戰略是.
以愛還仇恨.
耶穌說.
你們聽見有人說.
要愛你們的淪附.
恨你們的仇敵.
但我告訴你們.
要愛你們的仇敵.
為逼迫你們的禱告.
叫你們作天父的兒女.
這個系列.
我想我不知道.
因為其實我們很多嘉賓講完.
我不知道有沒有.
已經講了這個好撒瑪利亞這個故事.
我認為這個好撒瑪利亞人的故事.
最能夠說明這一點.
但是.
好撒瑪利亞人的故事.
我想大家都聽過.
有一些地方我覺得忽略了.
有時候我們不看清楚.
如果找一個好的例子.
鼓勵人們參加這個十步色經.
剛才賣廣告.
這個好撒瑪利亞人的故事.
是很好的一個材料.
因為你看完.
上聞下理.
耶穌說這個撒瑪利亞人故事的場合.
其實處理的問題是.
人怎樣可以得到永生.
有一個律法師來試探耶穌說.
老師.
我該做什麼才可以承受永生.

$^{481}$就是說.
我要做什麼才能上天堂呢.
其實.
不跟大家讀.
因為大家我想都聽過這個故事.
耶穌和這位律法師都同意.
要得到永生.
人要愛神.
也要愛鄰舍如同愛自己.
唯一爭議的地方就是.
誰是我的鄰舍呢.
所以耶穌說完這個故事之後.
耶穌就問這個律法師.
這三個人當中.
大家知道這個故事.
誰是落在強盜手中的人的鄰舍呢.
律法師就說是憐憫他那個.
我們聽聽這個回應.
好像根本連撒瑪利亞人這個名字都不肯提.
我們都知道為什麼.
猶太人鄙視撒瑪利亞人.
耶穌就對他說.
你去照樣做吧.
照樣做什麼.
耶穌是否叫這個律法師去做一個好撒瑪利亞人呢.
還是耶穌是叫他去愛那個撒瑪利亞人呢.
因為那個撒瑪利亞人就是他的鄰舍.
他得不得到永生.
就是在乎他去不去愛那個撒瑪利亞人.
他的鄰舍如同愛他自己.
耶穌說這個故事.
明顯是有些想刺激一下猶太的律法師.
明白這一點是解釋聖經的一個關鍵.
很多時候可能我們教會大的從兒童主流開始.
我們說到這個撒瑪利亞人故事的時候.
已經不期然地將自己代入了.
這個撒瑪利亞人.
很撒瑪利亞人的角色裡面.
以為我們可以慷慨幫助有需要的人.
為自己贏取永生.

$^{521}$以為這個是故事的教教訓.
我們其實需要的是將自己代入那個.
被搶劫 貧窮 受迫害.
被人打到半死的人.
他被人救回來之後應該做些什麼.
就是愛自己一樣.
去愛那位救回他的好撒瑪利亞人.
這個人可能曾經幫過你.
甚至可能是你的救命恩人.
但因為做了一些事情.
你覺得他迫害你.
你反而當他是他的仇人一樣.
可能他真的心裡面不太喜歡你.
覺得你很討厭.
但也可能他從來只是你心裡面的假想敵.
但無論如何.
當你能夠將你的仇敵看為你的鄰舍的時候.
天國裡面你其實就消滅了一個仇敵.
這個好撒瑪利亞人.
當然可能就是叫耶穌的加利利人.
但其實也可能是在你身邊.
一個神預定得救.
還沒有聽到福音.
是等待你去存方的人.
有一首歌說.
「以愛還愛」這個奇妙的道理.
終生都銘記.
其實天國裡面有一個比「以愛還愛」更加奇妙的道理.
就是「以愛還仇恨」.
讓我總結一下.
今天其實跟大家不是說太多.
因為都很簡單的道理.
開始的時候跟大家解釋.
在地上如同天上.
這句說話在聖經當中最直接的來源就是主土文.
所以這個主土文裡面提醒我們.
其實這個天上面不是一個安樂窩.
天上面不是一個樂園那麼簡單.
天上面仍然有一個增戰.
不過神已經特成.

$^{561}$地上面不同的地方是殺難魔鬼.
仍然有試探人.
有使我們跌倒.
這個惡者仍然是在做工.
所以我們仍然在這個地上面有這個天國的增戰.
主土文不是求地上面或者天下有太平安樂的日子.
今天就來到.
不是求教堂變成天堂.
因為其實天上仍然有戰爭.
我們所求的其實是天國在地上面能夠擴展.
我們所求的是地上這個增戰和天上一樣得到勝利.
我們要勝利就要把握一些天國的戰略.
今天和大家分享了三點.
寫生 犯特生.
委人 是僕人.
以愛 還仇恨.
這三個就是耶穌基督走上這個十價之路.
並且得到勝利帶著的天國戰略.
如果我們願意效法我們的主耶穌基督.
用同樣的天國戰略的時候.
我相信我們都能夠過一個特性的人生.
我們一起低頭禱告.
幫助我們在這個地上.
我們剩餘的日子裡面.
真的發揮你所安排我們的崗位.
要做的事情.
鼓起勇氣.
帶著信心去走這條難走的路.
就好像主耶穌基督背負這個十字架.
一直是完成這個他的使命.
我們將弟兄姊妹.
如果今天當中的.
仍然是不明白這個福音的真理.
求神你也給他們有勇氣繼續學習.
當中如果已經清楚自己有信仰的時候.
但是未願意再走一步的時候.
也鼓勵我們.
真的拿起我們的全副軍裝.
為天父來作戰.
禱告是奉主耶穌基督聖名義求.

\newpage



\section{約翰福音 21:1-19}
\label{sec:031924}
\textbf{全人得痊癒}~~~~俞大鳴牧師
\newline
\newline
連結: \href{http://www.church.com.hk/acms/content.asp?site=cdc\&op=show\&type=product\&code=031924\&layout=sermon}{\texttt{031924}} ~~~~ 語音日期: 2026-04-05
\newline
\newline
\hyperref[sec:031923]{\small{< < < PREV SERMON < < <}}
~
\hyperref[sec:index_chronic]{\small{[返順時目]}}
~
\hyperref[sec:031925]{\small{> > > NEXT SERMON > > >}}
\newline
\newline
約翰福音 21:1-19
\newline
\begin{longtable}{cl}
\hline
\hline
章節 & 經文 (和合本修訂版)\\
\hline
21:1 & \begin{tabularx}{0.7\textwidth}{X} 這些事以後，耶穌在提比哩亞海邊又向門徒顯現。他怎樣顯現記在下面。 \end{tabularx} \\ \\ \relax
21:2 & \begin{tabularx}{0.7\textwidth}{X} 西門‧彼得、叫低土馬的多馬、加利利的迦拿人拿但業、西庇太的兩個兒子，和另外兩個門徒，都在一起。 \end{tabularx} \\ \\ \relax
21:3 & \begin{tabularx}{0.7\textwidth}{X} 西門‧彼得對他們說：「我打魚去。」他們對他說：「我們也和你一起去。」他們就出去，上了船；那一夜並沒有打著甚麼。 \end{tabularx} \\ \\ \relax
21:4 & \begin{tabularx}{0.7\textwidth}{X} 天剛亮的時候，耶穌站在岸上，門徒卻不知道他是耶穌。 \end{tabularx} \\ \\ \relax
21:5 & \begin{tabularx}{0.7\textwidth}{X} 耶穌就對他們說：「孩子們！你們有吃的沒有？」他們回答他：「沒有。」 \end{tabularx} \\ \\ \relax
21:6 & \begin{tabularx}{0.7\textwidth}{X} 耶穌對他們說：「你們把網撒在船的右邊，就會得到。」於是他們撒下網去，竟拉不上來了，因為魚很多。 \end{tabularx} \\ \\ \relax
21:7 & \begin{tabularx}{0.7\textwidth}{X} 耶穌所愛的那門徒對彼得說：「是主！」那時西門‧彼得赤著身子，一聽見是主，就束上外衣，跳進海裡。 \end{tabularx} \\ \\ \relax
21:8 & \begin{tabularx}{0.7\textwidth}{X} 其餘的門徒因離岸不遠，約有二百肘，就坐著小船把那網魚拉過來。 \end{tabularx} \\ \\ \relax
21:9 & \begin{tabularx}{0.7\textwidth}{X} 他們上了岸，看見那裡有炭火，上面有魚和餅。 \end{tabularx} \\ \\ \relax
21:10 & \begin{tabularx}{0.7\textwidth}{X} 耶穌對他們說：「把剛才打的魚拿幾條來。」 \end{tabularx} \\ \\ \relax
21:11 & \begin{tabularx}{0.7\textwidth}{X} 西門‧彼得就上船，把網拉到岸上，網裡滿了大魚，共一百五十三條；雖然魚這樣多，網卻沒有破。 \end{tabularx} \\ \\ \relax
21:12 & \begin{tabularx}{0.7\textwidth}{X} 耶穌對他們說：「你們來吃早飯。」門徒中沒有一個敢問他：「你是誰？」因為他們知道他是主。 \end{tabularx} \\ \\ \relax
21:13 & \begin{tabularx}{0.7\textwidth}{X} 耶穌走過來，拿餅給他們，也照樣拿魚給他們。 \end{tabularx} \\ \\ \relax
21:14 & \begin{tabularx}{0.7\textwidth}{X} 耶穌從死人中復活後向門徒顯現，這是第三次。 \end{tabularx} \\ \\ \relax
21:15 & \begin{tabularx}{0.7\textwidth}{X} 他們吃完了早飯，耶穌對西門‧彼得說：「約翰的兒子西門，你愛我比這些更深嗎？」彼得對他說：「主啊，是的，你知道我愛你。」耶穌說：「你餵養我的小羊。」 \end{tabularx} \\ \\ \relax
21:16 & \begin{tabularx}{0.7\textwidth}{X} 耶穌第二次又對他說：「約翰的兒子西門，你愛我嗎？」彼得對他說：「主啊，是的，你知道我愛你。」耶穌說：「你牧養我的羊。」 \end{tabularx} \\ \\ \relax
21:17 & \begin{tabularx}{0.7\textwidth}{X} 耶穌第三次對他說：「約翰的兒子西門，你愛我嗎？」彼得因為耶穌第三次對他說「你愛我嗎」，就憂愁，對耶穌說：「主啊，你無所不知，你知道我愛你。」耶穌說：「你餵養我的羊。 \end{tabularx} \\ \\ \relax
21:18 & \begin{tabularx}{0.7\textwidth}{X} 我實實在在地告訴你，你年輕的時候，自己束上帶子，隨意往來；但年老的時候，你要伸出手來，別人要把你束上，帶你到不願意去的地方。」 \end{tabularx} \\ \\ \relax
21:19 & \begin{tabularx}{0.7\textwidth}{X} 耶穌說這話，是指彼得會怎樣死來榮耀神。說了這話，耶穌對他說：「你跟從我吧！」 \end{tabularx} \\ \\
[1ex]
\hline
\hline
\end{longtable}
$^{1}$祈禱兩姐妹們平安.
我們先一起祈禱.
讓我們在復活節的時候.
我們都可以一同來聚集.
來讚美你,紀念主你的復活.
扶你一聲令在我當中幫助我們.
明白使徒約翰寫下的訊息.
讓我們都得著鼓勵.
奉耶穌基督的名祈禱.
阿們.
教會今年的主題是「全心全意」.
其實這個主題有一個含意.
就是說很多信徒會友.
信耶穌都不是很全心全意.
所以教會的牧師們.
去年開會的時候就決定今年的主題.
我們要鼓勵大家.
我們要全心全意去跟隨主.
今年的復活節.
我們就是從使徒彼得的生命去思想.
主耶穌怎樣幫助他全心全意的.
進行耶穌給他的使命.
為天國努力.
去得人如魚.
去牧養教會.
如果你兩天之前有來參加.
英文堂的「受難節崇拜」的話.
就會知道.
那個講道是講到彼得三次不認主.
彼得在主耶穌未曾被釘十字架之前.
曾經對主耶穌說.
我就是必須與你同死.
也絕不會不認你.
當主耶穌在赫西瑪利園.
被人逮捕的時候.
也是彼得帶了一把刀.
砍了那個大祭司的捕人的一隻耳朵下來.
彼得我耶穌的心可以說是可見一斑.
但是在大祭司府裡面的時候.
他由於害怕.

$^{41}$由於怕被人抓.
他就三次否認耶穌.
堅持自己是不認得耶穌的.
和耶穌沒有關係的.
不是耶穌的門徒.
但是當雞叫的時候.
耶穌看一看彼得.
他就出去同哭.
彼得很可能是由於自己軟弱來同哭.
他本來以為自己是全心全意跟隨耶穌.
甚至為了耶穌可以捨棄自己的生命.
但是在危險的處境出現的時候.
他發覺原來他的心是破碎的.
如果你星期五有來.
這個心是灰色的.
他並沒有全心全意去為耶穌而死.
他的心是不完全的.
因為他的心靈曾經軟弱.
但是肉體是軟弱的.
我相信彼得走出去同哭.
表明他的跌倒.
其實是令到他自己都對自己感到很失望.
他對自己失望.
對前途失去了信心.
耶穌呼召他去得人如願.
他曾經跟著耶穌到處傳福音.
看過主耶穌所行的神蹟.
在未去耶路撒冷之前.
他與其他使徒跟著主耶穌.
去到凱撒利亞菲勒比的地方.
當被問到他們對主耶穌的認識的時候.
是彼得承認耶穌是基督.
是永生上帝的兒子.
當時他還被主耶穌稱讚.
主耶穌還對他說.
我還告訴你.
你是彼得.
我要把我的教會建造在這盆石上.
陰間的權柄不能勝過他.
我要把天國的弱勢給你.

$^{81}$凡你在地上所捆綁的.
在天上也要捆綁.
凡你在地上所釋放的.
在天上也要釋放.
我們可以想想彼得.
當時的心情是怎樣的呢.
會不會很興奮呢.
會不會那時候他覺得.
哇 真是自己人生的高光時刻.
主耶穌似乎.
將一個非常重大的責任委託了給他.
給他有極大的權柄.
耶穌將天國的弱勢給了他.
還對他說.
凡你在地上釋放的.
在地上捆綁.
在天上也要捆綁.
你在地上釋放.
在天上也要釋放.
聽起來好像有大權.
掌管別人的生死.
彼得當時真的可能飄飄然的.
但是好戲還在後頭.
過了六天之後.
主耶穌還帶著彼得 雅各和約翰.
領他們悄悄地上了高山.
然後在他們面前變了形狀.
耶穌的臉明亮如太陽.
衣裳潔白如光.
彼得真的願意為耶穌死.
他說那句話的時候是真心的.
只可惜肉體軟弱了.
面對強權的時候.
他不敢認耶穌.
最後耶穌被釘十字架.
我相信彼得和其他的使徒.
都是陷入了一個迷惘困惑之中.
事情發生得太突然了.
所以耶穌曾經向他們的門徒.
預先講了三次.

$^{121}$我這次去耶路撒冷.
我跟你們說.
是會被人辱罵 被人陷害.
被人釘十字架.
最後還會復活.
但是門徒對這些說話.
其實聽就聽了.
不是很明白.
聽到卻不明白.
所以事情真的發生的時候.
彼得就跌倒了.
後來耶穌復活了.
我們顯現了之後.
顯現了之後又消失了.
又不知道去了哪裡.
這次門徒仍然未完全清醒過來.
結果在楊寶欽二十一章三節.
我們看到彼得可能覺得前路茫茫.
經常在家裡坐也沒用.
沒什麼意思.
不如吃飯 不如重操僱業.
去打魚 去賺錢.
當時有六個門徒跟他在一起.
覺得與其無所事事.
那就一起去.
一起跟了彼得去打魚.
這個時候門徒在船上.
可能是百感交集.
他們在這個湖上.
跟耶穌有很多不同的經歷.
有一次當他們打魚打了一整晚.
都沒有收穫的時候.
耶穌就吩咐他們.
把網撒到船的另一邊.
結果他們就打了很多魚.
另外一次他們在船上.
遇到突然的風浪.
幾乎沒有命 他們就喊耶穌救命.
耶穌就起來 斥責風和浪.
現在他們幾個門徒在船上.

$^{161}$又在晚上工作.
又是一無所獲.
下一餐都沒有著落.
當他們還在船上沉思的時候.
天色開始變.
太陽的曙光開始出來了.
突然間他們聽到從岸上有聲音來問他們.
孩子們 你們沒有魚吧.
這個是楊卜二十一章五節.
這個就是環聖和新漢語的翻譯.
就不是像和合本那樣.
和合本就是問他們有沒有魚.
但是根據希臘文的文法翻譯.
其實耶穌知道他們沒有魚.
所以就會對他們說.
孩子們 你們沒有魚吧.
這個問題是問門徒有沒有魚獲.
不是問門徒有沒有魚獲.
其實他知道門徒沒有魚獲.
因此當門徒確認他們沒有魚獲之後.
那個聲音告訴他們.
將網撒向船的右邊.
那邊有魚.
他們就撒網下去.
這次竟然撒了很多魚.
連網都撞到拉不上來.
因為太多魚了.
後來我們知道有153條大魚.
約翰真的是一個比較聰明的門徒.
他可能想.
一個站在岸邊的人.
怎麼會知道船的右邊有魚呢.
魚是在水裡面到處游的.
怎麼可能被一隻船分開.
這邊沒有魚 這邊有魚呢.
他就馬上聯想到.
在岸上講話的人一定是主耶穌.
就告訴彼得 是主.
彼得真的是一個行動型的人.
他一聽到是主.

$^{201}$就快快的拿起衣服跳進水裡.
他不是想走.
而是想快快的游上岸.
迫不及待.
游上岸去見主耶穌.
主耶穌都知道門徒的需要.
對他們說 你們來吃早飯吧.
主耶穌不單止為他們煮好了早餐.
還親自將餅和魚分給他們.
主耶穌是首先為門徒.
提供了他們肉體的需要.
讓他們吃飽.
人是上帝所做的.
上帝做人是需要吃東西.
吃有營養的食物.
去維持生命和健康.
主耶穌是先為門徒.
預備了一個營養豐富的早餐.
十五節才開始跟彼得聊天.
吃早餐在猶太人的習俗裡面.
其實是非常重要的.
我們每個月都會守主餐.
守一次主餐.
有沒有人想過為什麼要守主餐.
主耶穌為什麼要設立主餐呢.
有什麼意義呢.
我們看幾段舊約經文去了解一下.
猶太人吃早餐的習俗.
在創世記二十六章二十六節.
和三十節三十一節.
阿比米勒和他的顧問阿胡薩.
和他的軍隊沿水飛角.
從基拉以來到以撒那裡.
以撒為他們擺設筵席.
他們就一起吃喝.
他們清早起來彼此起誓.
以撒送他們走.
他們就平平安安地離開他去了.
阿比米勒是菲律師人的王.
他本來是很妒忌以撒的.

$^{241}$不喜歡以撒的.
並且把以撒那些僕人挖的井.
全部用沙子封了.
後來直接要求以撒.
不如你離開我們去別的地方.
不要和我們住得這麼近.
當以撒走了之後.
他的軍長就有個不同的建議.
就和阿比米勒說.
這個人可能真的上帝很祝福他的.
不知道將來他強大起來.
會不會來打我們.
不如這樣吧.
我們去找他.
和他立一個和約.
以免他日後來攻打我們.
以撒就答應了和他們立和約.
答應了立和約.
之前做什麼呢.
一起吃飯.
吃完飯之後再立約.
我們看另一段經文.
這次是有關雅各的事蹟.
雅各的事蹟我們很多人都知道.
雅各為他岳父打工.
打了二十年工.
他覺得岳父和兒子.
開始對他的態度有變了.
對他開始不太友善了.
於是帶了他的妻子兒女.
和所有的家產.
不辭而別離開拉班.
走了幾天之後.
拉班知道了.
有人告訴拉班.
他走了.
拉班就找人一起去追雅各.
但是追到之前.
上帝警告了拉班.
所以最後拉班見到雅各之後.

$^{281}$就要求和雅各立約.
要雅各好好對我的女兒.
並且還要答應互不攻擊.
他同意立約之後做什麼呢.
31章46節 創世記31章46節.
他們納石頭堆成一堆.
於是在那裡在石頭旁邊吃河.
答應了立約之後.
又是吃一頓飯的.
好了 看完創世記.
我們看一段出埃及的經文.
當上帝拯救以色列人出埃及之後.
就去到荷列山.
荷列山上帝和人民立了約.
宣佈了十戒.
去到24章8到11節.
又有這樣的記載.
他說上帝吩咐摩西要亞倫加.
祭司和以色列的70位長老上山.
上山做什麼呢.
將法律再頒佈給他們.
頒佈給他們是不是.
不是的.
我們看看8到11節說什麼.
摩西把血灑在百姓身上.
說 汗呀 這是立約的血.
師爺和按照這一切的命令和你們.
立約的憑據.
摩西亞倫 拿達 亞比胡.
以及以色列長老中的70人都上去.
看見了以色列的上帝.
在他的腳下彷彿有藍寶石鋪到.
名正如天.
他不把手伸在以色列領袖的身上.
他們瞻仰上帝.
接著做什麼呢.
又吃又喝.
在上帝面前吃喝.
就是和上帝一起吃飯.
上帝和以色列人立約之後.

$^{321}$他按照人間的習俗.
讓人民的領袖.
他們就是代表人民.
在他面前吃飯.
從這些記載.
我們可以看到.
立和約 吃飯.
是猶太人的習俗.
因此主切立聖餐.
是要我們紀念的.
不只是他為我們所做的事.
更加重要的是.
是要我們紀念.
我們是藉著主耶穌和上帝.
立了和約.
所以每當我主持聖餐的時候.
都會強調主耶穌的說話.
這杯是用我的血所立的新約.
主餐是一個立約的餐.
我們以時而時紀念.
我們是跟上帝立了約的人.
也是曾經立誌要跟隨耶穌的人.
主耶穌和門奴吃完早餐之後.
也有重要的話對彼得說.
主耶穌和彼得這一段話.
其實很多牧師都用來講道.
有很多牧師.
包括中外牧師.
我也聽過.
都很小心研究主耶穌.
問了彼得三次.
你愛我嗎?.
並且有些人研究希臘文.
從希臘文裡就知道.
希臘文的「愛」字.
原來有四個字都可以翻譯為「愛」.
他們發現在這段對話當中.
前兩次主耶穌的問題.
用了希臘文的一個字.
叫做「阿加皮」的「愛」字.

$^{361}$很多人理解這個字.
「阿加皮」這個字是指上帝那種.
無私,純正,毫無雜質的愛.
而當彼得回答主耶穌的時候.
「我愛你」的時候.
他就用了另外一個希臘文的「愛」字.
叫做「Fileo」.
這個「愛」字是指人與人之間的「愛」.
如果大家知道美國有一個州.
其實是由兩個希臘文組成的.
前面的「Fila」就是「Fileo」那個字.
後面「Delphia」就是兄弟的意思.
就是兄弟相愛的州.
我不知道他是否真的相愛.
但是很多牧師發現.
耶穌問彼得.
「你是不是阿加皮」.
彼得就說「我Fileo你啦」.
他們就認為.
耶穌是不滿意彼得頭兩次的答覆.
所以他就用人間的「愛」.
因為他用人間的「愛」.
去回應上帝那種純正的愛.
所以在第三次的時候.
耶穌用了「Fileo」問彼得.
好像鋼隔問彼得一樣.
所以彼得就傷心難過.
其實這個說法是對希臘文的誤解.
我首先解釋一下.
「阿加皮」和「Fileo」這兩個字的確有分別.
但是分別就不是那麼大.
有一位世界知名的神學教授叫Don Carson.
其實他已經在1984年寫了一本書.
書名叫做《Ecstatic Fallacies》.
我翻譯為《色經的謬誤》.
在這本書裡面他就指出.
將「阿加皮」和「Fileo」.
說成是兩種完全不同等級的愛.
是一個色經的謬誤.
其實這兩個字是可以互換使用的.

$^{401}$這兩個字是同義詞.
尤其是在約翰福音裡面.
我們留意到約翰寫書的風格.
他喜歡用同義詞.
而不喜歡將一些字重複使用.
那為甚麼主耶穌要問彼得三次.
「你愛我嗎?」.
一個很簡單的原因.
就是由於彼得曾經三次否認著.
而主耶穌就為了這位教會的.
未來領袖的聲譽.
要給他三次機會.
去表明他是真心愛主的人.
我留意到其實使徒約翰.
在第十四節已經有個暗示.
要我們注意三次的問題.
十四節就很簡單.
他說耶穌從死人中復活向門徒顯現.
「者是第三次」.
其實這句話我發覺.
說不說是沒有所謂的.
這本書快要結束了.
我們知道二十一章之後沒有了.
約翰福音寫完了.
突然間要提第三次有甚麼意義呢.
事實上從二十章開始.
耶穌復活之後向人顯現.
其實不只三次 是有四次.
但是第一次是向瑪利亞顯現.
瑪利亞算不算是門徒啊.
我想約翰是不算那次的顯現.
所以他就不算那次.
其後的三次約翰都強調.
耶穌是向門徒顯現.
在這裡的門徒不單止是普通的門徒.
是指特別的十二個仕徒.
如果我們不讀第十四節.
只讀第十三節.
然後跳到第十五節.
我相信沒有人會覺得上文不接下理的.

$^{441}$我試試讀給大家聽.
耶穌走過來拿餅給他們 也照樣拿魚給他們.
他們吃完早飯 耶穌對西門彼得雪.
有沒有覺得少了東西.
有沒有覺得不太通順.
不會的 我又不覺得不通順.
所以第十四節最重要的關鍵是提醒我們.
我們要留意第三次的問題.
我們繼續思想.
為什麼耶穌要問比德三次你愛我嗎.
最近有人曾經問過我.
他說耶穌曾經說過一句話.
他說 凡在人面前不認我的.
我在我天上的父面前也不認他.
哇 比德三次不認耶穌.
他會不會得救.
答案是肯定的.
沒錯 人都是會有軟弱的.
在特殊的情況下.
都有可能否認主.
比德否認主的時候.
是主耶穌被捕之後.
比德是在大祭司的住宅的花園.
當時他的情緒又不穩定.
又有被逮捕的危險.
以致他軟弱跌倒.
當然 從縮靈的角度去看.
這次是上帝和撒旦之間的角力比拼.
所以聖經中記載到主耶穌.
預告比德會三次不認主之前.
還說了一些話給他聽.
耶穌怎麼說呢.
「西門西門 撒旦已得到准許來試探你們.
要像農夫篩墨子一樣來篩你們.
但是我已經為你們祈求.
使你不至於失掉信心.
你再迴轉歸我以後.
我就要激勵你的弟兄們」.
比德什麼時候迴轉呢.
我相信就是他吃完早餐之後.

$^{481}$和主耶穌交心 傾談之後.
他迴轉了.
主耶穌給他三次機會.
去表達他對主耶穌的愛.
來贖回他三次不認主的過錯.
就好像他犯了三次罪.
主耶穌給他三次認罪的機會.
在第一次主耶穌問比德.
約翰的兒子西門.
「你愛我比這些更深嗎」.
這些是指什麼.
從當時的場景來說.
這些應該是指一大堆魚.
他剛剛撈了一大堆魚上來.
一百五十三條大魚.
那些魚大概都可以賣不少錢的.
主耶穌問比德.
「你愛我比我這些物質更多嗎」.
「你愛捉魚的職業多過愛我嗎」.
當比德回答說愛主之後.
主給他一個新的專業.
新的人生方向.
「你餵養我的小羊」.
「你牧養我的羊」.
「你餵養我的羊」.
重複了三次.
「餵養」和「牧養」.
都是兩個不同的希臘文字.
但都是同義字.
所以主耶穌要求比德.
轉變他的職業.
不要再做漁夫了.
你要全心全意投身天國.
做好牧養的工作.
接著主耶穌都將比德的生命.
會怎樣終結告訴他.
不過最後補充一句.
「你跟從我吧」.
似乎是鼓勵比德.
無論你人生會發生甚麼事.

$^{521}$人生的目標一定要對準耶穌.
跟隨耶穌.
今天是復活節.
主耶穌是從一個黑暗的墳墓中.
走出來進入一個有光的世界.
今天的經文.
發生的場景都是很有象徵意義的.
門徒在黑暗的大海中.
一無所得.
當他們筋疲力竭.
肚子餓的時候.
岸上有炭火的光出現.
黎明的曙光.
亦都在天上出現.
主耶穌跟他們一起吃早餐.
再跟門徒立約.
讓比德有三次的機會去認罪.
然後委任比德去牧養教會.
今天我不知道有沒有人在我們當中.
正在經歷人生的低谷.
前路沒有方向 沒有量過.
我要告訴你.
耶穌要使你全人得全遇.
要滿足你肉體的需要.
亦都會給你一個新的人生方向.
新的人生使命.
使你活出一個豐盛有意義的人生.
只要你願意全心全意愛主.
跟隨這位復活的主.
我們一起祈禱.
不單止為我們的罪婦賞贈駕.
亦都從死裡復活.
為所有的信徒愛她的人.
都能夠得到存餘.
求您幫助我們全心全意愛主.
並且以我們的生命來事奉主.
奉耶穌基督的名祈禱.
阿們.
\newpage



\section{出埃及記 10:21}
\label{sec:031925}
\textbf{聖經人聖經事:摩西向天伸杖}~~~~黃克勤牧師博士
\newline
\newline
連結: \href{http://www.church.com.hk/acms/content.asp?site=cdc\&op=show\&type=product\&code=031925\&layout=sermon}{\texttt{031925}} ~~~~ 語音日期: 2026-04-03
\newline
\newline
\hyperref[sec:031924]{\small{< < < PREV SERMON < < <}}
~
\hyperref[sec:index_chronic]{\small{[返順時目]}}
~
\hyperref[sec:031926]{\small{> > > NEXT SERMON > > >}}
\newline
\newline
出埃及記 10:21
\newline
\begin{longtable}{cl}
\hline
\hline
章節 & 經文 (和合本修訂版)\\
\hline
10:21 & \begin{tabularx}{0.7\textwidth}{X} 耶和華對摩西說：「你向天伸出你的手，使黑暗籠罩埃及地；這黑暗甚至可以摸得到。」 \end{tabularx} \\ \\
[1ex]
\hline
\hline
\end{longtable}
$^{1}$各位頂子妹,大家好!.
又係時間同大家一齊讀聖經啦!.
生命樹宣教網絡,做宣教工作都十幾年.
好想透過傳福音,宣教讓人認識福音.
亦都可以有更美好的靈修.
今次就係用jor 聖經人物,同埋聖經故事.
做我地ge 題目都一段時間架啦.
今次ge 聖經人,聖經事呢,就一路講緊摩世.
你知由創世記講起,而家出埋埃及記啦.
阿拜罕,以撒,雅各,約瑟,而家出埋埃及記摩西.
摩西就由佢出世一直到而家就十災.
叫大意識的人離開埃及.
今次ge 題目係摩西向天生長.
嘩!呢個故事真係犀利啦.
不過呢其實話事話呢就好多人都冇留意呢件事架喎.
咁你就留意十災呀,好緊要呀,點呀點呀.
D 細節呢我地好多時候讀聖經呢就冇留意到.
不過好有意思ge .
今次呢,就摩西向天生長,有D 象呀.
你又去睇下先.
大家都記得呢一章呢,係尋日用過ge .
係十災之中就記載係摩西及記七至十一章啦.
咁係第三災呢,係個失災.
呢D 字我地讀早ge ,原來呢佢應該讀失ge .
好!失災!.
咁就借試係摩西及記ge 第八章十九節.
去到失災ge 時候呢,嘩!遍地都係失.
咁我唔知你又冇經歷過,又係比失咬.
我呢就經歷過.
咁就特別呢,係以前細個Dge 時候啦.
都唔係好細個,十幾歲,中學生.
就係爸爸呢開間鋪頭ge .
有間五金鋪,有時有D 野貓捐jor 入去.
嘩!留低jor 好多跳蝨.
跳蝨最難搞啦.
咁當你發覺隻腳好痕ge 時候,特別夏天.
嘩!好熱啦!咁gwo 陣時冇冷氣架嘛.
著住短褲落去鋪頭幫手.
嘩!隻腳好多時都瀑jor 三五七隻蝨係度.
好痕呀!咁一郁佢已經跳jor 啦.

$^{41}$咁呢個失災係令人地好困擾.
而且搞到隻腳真係花曬一點點望.
好啦!失災!.
咁係gwo 度《出口入記》八章十九節就講緊一件事喎.
大家留意下.
咁我都開返D 螢光筆線.
用返黃色啦.
咁呢度呢,係八章十九節出口入記呢.
gwo 度係咁樣寫ge .
《這是神的手段》.
嘩!好寶貴!神ge 手段.
特別要留意,係神ge 手段黎架.
神要做ye .
咁神ge 手段.
S bar.
呢個字好難讀ge .
咁S bar.
S bar係指頭呀.
係finger.
咁gwo 度講呢.
《這是神的手指頭》.
咁呢就係佢ge 手指頭就做呢件事啦.
嘩!好厲害!.
咁就神係度工作喎.
另一件事呢,我地睇下呢.
同大家重溫je .
咁今次呢,我地係講黑暗之災ge ,今日.
另一節聖經我地好少留意.
就係出口入記四章十七節.
係gwo 個時候呢.
摩西就你知道啦同神傾緊計.
太極咁推手呀.
唔肯做唔肯做唔肯去.
咁樣喎.
一味呢就話唔去.
咁又話自己口才唔好.
咁神就話.
你有個國公叫阿倫.
好口才架,同你一齊去.
咁呢.

$^{81}$派埋阿倫去,實得啦.
而且我會教你地兩個點樣講.
咁係呢度我上次就話呢.
摩西呢個推堂呢.
就得jor 個機會.
本來就自己成日叫佢.
一個人帶曬,而家多jor 個領袖.
阿倫同佢分jor .
呢D 榮耀啦.
咁又於是呢就.
你地繼續講啦.
阿倫會幫你做口,不過.
就有第十七節.
神就話你手裡要拿.
這像,記得.
拎住呢支像.
好行神,呢支像係咩呢?就係佢.
神就話我去到佢地唔信.
你手裡面拎住咩?.
拎住jor 支像囉,你擋落地啦.
一擋就變jor 一條蛇.
gwo 條蛇呢係好活,好真.
嚇到摩西係跑開.
因為佢係曠野見過D 蛇.
嘩好毒呀,好犀利呀.
咁一見到條蛇,佢已經跑開啦.
咁即係話呢,係變jor 一條真蛇.
唔係一個幻象.
咁你叫佢執返,執住佢條尾.
就得啦,咁.
執住佢個弱點,捉住佢條尾.
咁摩西梗係震騰騰.
真係咁樣,又變返支像.
咁神就話啦.
你手裡要lor 隻像.
就吼呢,我叫.
阿倫幫你吼.
你自己要拎住呢支像.
好行神蹟.
因此呢.

$^{121}$呢支像呀,係摩西.
黎講呢,係非常.
之重要呀.
大家留意到喎,好.
咁今日要講ge 呢,出一集記.
第十章,黑暗之災.
就係一節夜話話.
對摩西說.
咁呀,呢個.
螢光紅色.
第一點,通常我都講.
三點,甚至同佢講.
你向天生.
象,使埃及.
地黑暗,呢個黑暗.
會黑到點呀.
黑到真係黑到不得了.
黑到似乎摸.
得著,呢個係一個形容.
黎,黑到.
嘩,真係.
黑得.
實在到好似摸到.
咁黑,咁即係咩呢.
超級黑,好黑.
唔係少少暗下,千萬唔好.
哎呀,有雲遮住.
陰天,唔係,係直直.
黑到伸手不見.
唔止,即係黑暗.
似乎摸得著,即係話.
黑到壓埋黎咁就係.
咁呢,就一個咁ge .
睇法,咁有時如果你真係去到.
好黑ge 地方啦,咁呢就.
即係搵D 棉花蓋住自己.
嘩,黑到不得了ge 時候.
你覺得個黑呢,係去到你前面.
伸手不見,唔止.
咁神叫你向天生.

$^{161}$象,記得.
咁呢個象呢,係.
就係呢度啦,就係上帝.
曾經叫摩西.
你gwo 支象,擋落地下就會變.
做蛇,撿返佢會變做象.
咁呢一個神,即係你記住.
呢支象就係我要.
要黎幫助你,你拎.
係手裡面,好得意ge 喎.
有支象拎住呢,好似.
係你有感覺D .
實在D ,有D 渣拿.
你見好多人呢,特別天主教導呢.
佢都好歡喜拎住個十字架.
咁呢就.
訓覺都拎住佢ge 喎,擺.
旁邊ge 喎,拎住個十字架呢.
好似係覺得,唉呀.
舒服D ,有信心D ge 喎.
咁gwo 支象呢,其實係一個.
上帝眷顧同在.
與行神職ge .
一個表達黎ge .
你向天伸張.
咪向地喎.
向天呀,黑暗呀嘛.
我又唔係向住你呢個人.
又係向地,我又唔係.
係向天.
向天即係乜ye 呢,係你.
向天伸張,係你伸出.
依靠ge 象,其實呢.
係代替你手黎ge .
你向天,向神伸出.
依靠ge 手.
依靠ge 象.
咁呢個時候,神就會做ye 架啦.
你.
記得,你呢,需要ge 時候.

$^{201}$依靠自己,諗好多.
計仔,唉,點樣.
可以嚇窒呢個.
法老呢,點樣可以.
做D 乜ye 呢,唔需要架.
你只需要係對抗你.
仇敵ge 時候,或者好艱難ge .
情況ge 時候,你係舉你ge 手.
舉你ge 象,向天.
伸張,向天伸手.
呢個時候.
神就會做ge 工作.
會令到埃及地.
本來係一個好猛太陽ge 地方.
黑到幾乎.
摸到咁黑,嘩,真係.
呢個表達真係好細.
係神自己講ge .
向天伸張,唔好忘記向天伸張.
係向主ge .
依靠,而且呢.
個象係曾經.
係經歷過神蹟ge ye 黎架.
咁你諗返以往上.
係你身上做過ge ye .
你繼續拎住呢個依靠ge 心.
向上,向神.
伸出你ge 手.
神係會幫助你ge .
第二樣ye 呢,摩西就好.
照樣向天伸張,冇懷疑.
真係假阿神得唔得架.
係咪咁架.
咁你會唔會真係.
試一次先,要唔要做個實驗.
先呀,冇架.
摩西就好有信心,佢就話.
向天伸張,將D .
象呢,舉樣先.
你知佢都慣jor ge .

$^{241}$你將D 象指向紅海水就分開.
咁佢D 象好緊要.
因此呢支象亦都係.
佢係一個摩西ge 象徵黎架.
係呢,係信靠ge 象徵.
第二點同大家分享呢.
就埃及地遍地.
就污黑了.
咁呢.
呢個應該叫做污黑之災.
就arm D 架.
咁我上面都有.
使埃及地黑暗啦.
咁因此都叫黑暗之災啦.
其實黑暗唔夠,因為個黑暗黑到呢.
係污黑.
污黑了三天.
嘩,咁呢個都幾得意喎.
大家.
有冇帶支電筒.
早晨起身出街返工架呢.
做咩帶支電筒呀.
咁唔使打呀,因為而家有手機.
有燈架嘛.
咁所以你唔使帶啦.
帶支電筒出街係咩呢.
係邊一個話比你聽.
聽朝實會天光架呢.
咁咪以往歷史.
都係架啦,從來都冇一日.
個天係唔光架,咁你到.
天光天黑,天光天黑,實有架啦.
所以呢個必然啦.
點解必然呀,其實唔一定.
必然架喎,係可以冇架喎.
咁你拎住支電筒.
明日返工就易D 喎,係可以.
冇架喎,不過我地認為.
有,因為上帝呢係咁.
規律地行架.

$^{281}$呢個世界混行緊ge .
不過到第二日,而家起身.
啊,乜個天咁耐都未.
光架呢今日.
咁當然當時我一頭就冇鐘.
冇錶啦,不過佢地梗有D .
辦法呢,去睇下時間.
已經打晒.
間咁耐喎,幾多點.
架喇喎,點解仲未.
光架呢,跟住.
嘩,照睇個時間下午.
架喇喎,仲未光,夜晚.
黑呀黑,今日冇jor .
光,真係奇,我心.
好驚,如果你發覺.
有一日,全日係烏黑.
我諗你驚到不得了.
咁以後,唉,算啦,果一晚.
希望明日天光,又黑.
又黎到一日.
死啦,個田點呢.
死啦,個樣事情點呢.
嘩,你驚到不得了.
唔知我親戚gwo 邊係點呢.
有好多ye 驚,到第三日.
光架掛,點知又黑.
嘩,真係我諗.
暈低,真係嚇到.
個心驚到不得了.
埃及遍地,烏黑三天.
係上帝做ge 神蹟黎架.
超級神蹟.
烏黑你三天.
呢個神蹟真係好呀.
好令人吃驚呀.
真係伏到唔睇頭地.
點知.
第三點.
先令人驚.

$^{321}$第23次話,三天之久.
人不能相見.
因為冇燈架嘛.
冇手機有燈架嘛.
出唔到街.
點火把點唔到咁耐.
冇咁多油,冇咁多燃燒ge 物料.
人不能相見.
誰也不敢起來離開本處.
唯有最安全.
亦都冇一個安全.
更安全ge 地方就ok啦.
誰也不敢起來離開本處.
三點紅色字要同大家睇.
唯有以色列家中.
唯有以色列人ge 家中.
都有亮光.
真係奇啦.
係以色列人ge 屋企裡面呢.
就有亮光.
咁呢個神蹟真係更加可怕啦.
真係可怕到極.
我地埃及人你從埃及ge 角度睇.
我都黑晒.
係以色列住gwo 幾個.
住係大部分住係gwo 度.
gwo 幾個係貨倉城.
要佢地起緊ge 地方呢.
就有光喎.
佢地繼續係日頭有光.
夜晚黑日太好光.
嘩,呢個佢地ge 上帝真係犀利喎.
點會咁呢.
係能夠特別咁分開.
佢地光我地黑喎.
三日喎.
我諗係呢個時間D 埃及人.
嘩,呢個以色列ge 神.
先係真ge 神.
叫咩名啊,以色列個神叫咩名.

$^{361}$叫耶和華喎.
嘩,呢個耶和華.
先係真神.
其實你會知道呢,上帝做好多方法.
想人地認識到.
邊一個先係真神.
呢個就係啦.
第二個唔係呀.
你會知道呢,就係呢個景況.
唯有以色列家中.
都有亮光係神.
一步去表達.
就係話比你聽你gwo 個太陽神.
唔係神.
佢係拜太陽架.
原因就係耕種呀嘛.
農業社會嘛,佢需要好猛烈ge 陽光架.
咁所以佢要拜太陽.
上帝就好.
拜太陽啦.
我熄燈三日,睇你點.
你地會知道.
係邊個可以將太陽熄滅jor .
三日.
不過仍然係以色列家.
都有亮光.
其實之前我相信.
係gwo D 災裡面呢.
好有可能冇特別講.
應該埃及地裡面.
以色列人聚居ge 地方.
都受災.
不過黎到黑暗之災.
係以色列.
所在之地有所分別.
點解咁講.
因為之前冇講到.
有以色列人係唔受災.
係黎到呢個.
黑暗之災先講.

$^{401}$唯有以色列家中都有亮光.
咁大家留意下.
咁即係摩西gwo D 像.
今日同大家講到呢度啦.
咁我要再.
留意,請大家留意返.
遮視神的手段.
手段呢個字.
今日我地就覺得.
好貶義,灑手段.
手段其實係.
做ye ,隻手分段.
分開黎.
做ye ge 方法.
係神做ye ge 方法.
係佢ge 指頭係度工作.
好.
最緊要你有一支.
內心ge 像.
係依靠上帝.
憑信心而行事ge .
神會用gwo 支像.
黎為你施展神蹟.
第一點就係你要向天對象.
唔好向地,唔好向人.
冇用ge ,靠自己,靠人都唔得.
向天依靠神.
伸出信心ge 像.
或者信心ge 手.
神就會工作,令到埃及地.
烏黑到摸得著.
三天,神就會做ye ge .
而且呢,係做得好徹底ge .
第三點就係.
上帝可以有咁大ge 能力.
係度彰顯,以色列.
家中有良國.
哇,好欣賞,好讚美.
上帝咁奇妙ge 作為.
求主呢比我地更加有.

$^{441}$信心去依靠佢.
我地一齊報告啦.
係呢度唔能夠.
唔獻上讚美.
我地要歌頌我地ge 神.
因為你實在係大能ge 神.
你係能夠叫.
日頭冇光.
係真係烏黑到三日.
好似壓埋你咁黑.
D 人係根本上.
唔敢出街.
係埃及經歷呢個日子.
而竟然間.
係係以色列地.
係有亮光.
咁ge 分別.
呢個光根本上人完全.
絕對唔能控制.
冇任何ge 人為ge .
能力與邪術,巫術.
可以做到ge .
冇人能夠挑戰呢一個.
咁偉大奇妙ge 神蹟.
所以我地就見到.
睇到呢個紀錄.
主求主係加我地信心.
比我地.
zar 緊gwo 支.
信心ge 象.
或者我地事事.
舉起我地信心ge 手.
向我地ge 上主.
黎到禱告.
求幫助.
你係能夠.
比我地一個幫助.
你係能夠使.
日頭黑暗三天ge 主.
比我地記得.

$^{481}$我地手中.
拎住信心ge 象.
我地信心ge 手係有作用ge .
比我地能夠向你.
向天舉起.
我地就得jor 幫助.
依靠耶和華ge .
並不蒙羞.
主我地多謝你自己應許.
願你得jor 榮耀讚美.
禱告奉耶穌命求.
\newpage



\section{以賽亞書 53:4-5}
\label{sec:031926}
\textbf{苦難中的出路}~~~~李健明牧師
\newline
\newline
連結: \href{http://www.church.com.hk/acms/content.asp?site=cdc\&op=show\&type=product\&code=031926\&layout=sermon}{\texttt{031926}} ~~~~ 語音日期: 2026-04-05
\newline
\newline
\hyperref[sec:031925]{\small{< < < PREV SERMON < < <}}
~
\hyperref[sec:index_chronic]{\small{[返順時目]}}
~
\hyperref[sec:031927]{\small{> > > NEXT SERMON > > >}}
\newline
\newline
以賽亞書 53:4-5
\newline
\begin{longtable}{cl}
\hline
\hline
章節 & 經文 (和合本修訂版)\\
\hline
53:4 & \begin{tabularx}{0.7\textwidth}{X} 他誠然擔當我們的憂患，背負我們的痛苦；我們卻以為他受責罰，是被神擊打苦待。 \end{tabularx} \\ \\ \relax
53:5 & \begin{tabularx}{0.7\textwidth}{X} 他為我們的過犯受害，為我們的罪孽被壓傷。因他受的懲罰，我們得平安；因他受的鞭傷，我們得醫治。 \end{tabularx} \\ \\
[1ex]
\hline
\hline
\end{longtable}
$^{1}$讓我們一起祈禱.
基督的復活.
是得勝.
但是祂釘十字架之前.
祂的受苦.
也對我們世界上每個人都有好處.
求天父你帶領我們.
讓我們在聖經裡面.
能夠深入地去感受.
基督耶穌的愛.
讓講道聽得到幫助.
奉耶穌名祈禱.
今天我講的是苦難中的出路.
想問問大家.
苦難的來源.
痛苦的來源.
究竟苦難和痛苦.
兩樣有什麼分別.
有什麼相同呢.
其實我們找答案.
很容易的.
現在在哪找資料.
找得最快的.
AI.
你走上AI就找到的了.
總括來說.
苦難和痛苦.
是有分別的.
但是都有相同.
痛苦就好講究.
多數都是講一些人.
親身或者個人性的痛苦.
是不牽涉到很多人的.
可能是單獨一個人.
牽涉少部分的人.
痛苦就是這樣解釋.
但苦難就不是.
苦難比較大一點的.
牽涉的人數會多很多的.
比如說天災人禍.

$^{41}$天災人禍.
很痛苦的.
現在看中東打仗.
死了多少人.
那兩個苦難.
苦難是講很大範圍的痛苦.
究竟苦難和痛苦.
有什麼相同.
就是講到人的痛苦.
你想想.
人的痛苦.
人在這個世界上.
無論他怎樣過.
或者不好.
無論他的家庭底細有多少.
他都相對上.
他都會受過痛苦.
個人的痛苦.
在困苦的裡面.
在困難的裡面.
他可能情緒上的痛苦.
可能家人帶來給他個人的痛苦.
但這個世界上.
有一個人的痛苦.
對全世界的人.
都有影響的.
所以我希望今天透過.
《以賽亞書》五十三章四節五節.
我們去看看.
究竟這個人.
為什麼他的痛苦.
對里圖.
對我.
對世界上所有人.
都有影響.
《以賽亞書》五十三章第四節.
他承認擔當我們的憂患.
背負我們的痛苦.
我們卻以為他受責罰.
被神擊打苦待了.

$^{81}$這個擔當的意思就是說.
這個人.
那個他.
這個他是說誰.
這個他就說耶穌基督.
現在另外一個問題就出現了.
《以賽亞書》.
說這個耶穌基督.
將來會受苦受難.
釘在十字架上面.
《以賽亞書》的議員.
說了耶穌基督.
到耶穌基督出生.
這段時間分了多少年呢.
我都不知道.
我在哪裡找資料.
因為有些東西是很特別.
突然間問你.
沒有人知道的.
他原來經歷了七八至八百年.
就是說《以賽亞書》才知道.
預言這個他.
這個耶穌基督出生.
到耶穌基督真的出生的時候.
相隔了七八年至八百年.
這個耶穌基督.
他親身擔當的意思就是說.
他一個人全都抱在身上.
憂患痛苦.
他全都抱在身上.
到他死的時候.
釘十字架.
人家個個都以為.
神不要他了.
但是不是.
五十三章第五節.
那知他為我不的過犯受害.
原來這個他.
這個耶穌基督在歷史上.
真的出現了這個人.

$^{121}$而這個人他的受苦.
是為了什麼呢.
為了你和我的過犯.
為世界上每一個人的過犯.
擔當他們的犯而受苦.
過犯是什麼.
過犯意思就是罪惡犯罪.
罪怎樣解釋.
我經常都講到都會提醒人.
你要告訴別人罪是什麼.
很多時候我們傳福音.
告訴別人你有罪.
人家是很反感的.
罪是什麼意思.
記住我下次來港島.
又會再問你.
我每次港島講到罪.
我就告訴你很簡單.
應該做的不去做.
不應該做的.
偏偏去做.
阿當阿媽不是這樣嗎.
神叫她不做她偏偏做.
神叫她做.
你要守這條誡命.
你要守我的吩咐.
她偏偏不守.
那就是罪.
意思就是說.
耶穌基督來到這個世界上.
去受苦.
不是很多人知道的.
但當我們信了耶穌.
我就知道祂的受苦.
就是為了你和我.
去擔當你和我的罪.
以致我們能夠懂得.
應該做去做.
不應該做不做.
喜伯羅書二章十八節.

$^{161}$祂自己竟然被試探而受苦.
就能搭救被試探的人.
意思是什麼.
試探等同受苦.
試探等同痛苦.
是不是.
原來耶穌基督來到這個世界上.
聖經講到明.
祂是為了受苦而受苦.
受苦的裡面.
是為了你和我而受苦.
我想問一些有子女那些.
你有沒有留意.
小孩子幾多歲.
起手講一個單字.
你猜一下.
沒什麼錯對的.
幾多歲起手講一個單字.
祂就這樣講的.
有一個弟兄.
有個牧師廣見證.
他查過.
平均來講.
普遍來講.
是一歲至兩歲你講得對.
主席講得對.
一歲至兩歲就懂得講單字.
但是廣見證這個牧師.
告訴我.
他到三歲都不懂得講單字.
不單是這樣.
到他讀小學的時候.
他告訴我.
他說很痛苦.
為什麼.
他不能一次過講一句句子出來.
意思就是說.
我愛爸爸.
整句句子是這樣.
我愛爸爸.

$^{201}$但他不能講整句.
我愛爸.
我愛講不了爸爸.
他到小學的時候更慘.
為什麼慘呢.
因為他讀不了書.
他不是不想讀.
他讀了之後.
他表達不了.
我們那時候讀書要背書的.
背書的.
你背不了書.
同學們恥笑你.
你每次都讀不了書.
念背不了書.
老師就說你很懶.
為什麼你這麼懶.
所以這個人.
這個人從小到中學的階段都很痛苦.
他覺得這是苦難.
他一生背的.
他想著一生一世都背的.
誰知道到中三的時候.
他有個師父.
不知道是師父還是老師.
他帶他去信耶穌.
信耶穌之後他生命改變.
開始懂得讀書.
讀中學.
讀儀科.
讀大學.
讀碩士.
他完全康復了.
為什麼康復了.
因為耶穌基督說.
我為師貪而受苦.
是要搭救受苦的人.
他真的靠著耶穌基督.
生命改變好不好.
不好.

$^{241}$有什麼事情發生.
他結了婚之後.
有三個小孩.
三個小孩都是什麼.
和他一樣.
一個你已經很痛苦.
三個你說痛不痛苦.
做爸爸媽媽看著.
痛不痛苦.
不痛苦.
為什麼不痛苦.
因為這個牧師說.
他結了婚之後.
他很清楚留意到.
他自己的情況.
他怕遺傳給小孩.
所以當小孩出世的時候.
他很留心.
究竟這個小孩.
是不是像我這樣.
第一個小孩.
我說三個都是.
第一個小孩真的這樣.
他怎樣應付.
因為他受過苦難.
所以他能夠幫助他.
受苦難的小孩.
就這麼簡單.
所以他一早就發現.
他第一個小孩是這樣.
所以他馬上找專家.
馬上找醫生.
找一些輔導.
找一些兒童科的人.
幫助這個小孩.
你想想.
原來耶穌基督說.
我親身去體驗.
經驗過這個苦難.
我就能夠搭救.

$^{281}$拯救有苦難的人.
現在這個牧師.
他就親身經歷過.
他本身的苦難.
所以他能夠.
將他經歷.
親身經歷.
由耶穌基督.
信了耶穌基督之後.
他生命改變.
他就將這個信耶穌.
生命改變.
從這個苦難走出來.
交給他的小孩.
站在我旁邊這個牧師.
是見到神學院的.
講師來的.
你知不知道.
他得到什麼病.
有沒有人聽過SEN.
SEN是什麼.
Special Education Need.
他就寫了這本書.
所以各位感恩姊妹.
我相信這兩個.
我估計.
個個都信了耶穌.
你知不知道我們很有福.
耶穌基督說.
我經歷過苦難.
所有能夠幫助受苦難的人.
我們信了耶穌.
我們不會沒有苦難的.
因為主禱文說什麼.
不叫我們遇見試探.
接著說什麼.
救我們出來凶惡.
不是叫我們不遇見試探.
意思是我們不要在試探裡面.
被試探或者苦難裡面打垮我們.

$^{321}$而是我們靠著耶穌基督.
有條出路.
這個就是這樣.
所以各位感恩姊妹.
你知不知道我今天講苦難.
我有什麼特別的表情.
開心 為什麼開心.
因為耶穌基督說我受苦難.
拯救了你的生命.
你不要這麼苦.
他說我已經苦了.
我告訴你有條出路.
這條出路是什麼.
當你經歷過之後.
你就可以幫助人.
你好像我這樣.
現在我在天上.
你在地上.
我在地上就可以幫助人.
在十字架上死了復活.
現在到我和你.
所以我們很開心.
我們有個使命.
就是當我們經歷苦難.
你不要擔心.
不要傷痛.
人性當然會痛痛的.
但是你不要痛痛一輩子.
你告訴自己.
我在苦難走出來.
因為耶穌基督說.
經歷過苦難.
才可以幫助受痛苦的人.
第二.
苦難中得平安.
因他受的鞭傷.
因他受的刑罰.
我們得平安.
得平安的意思就是.
身體上得到完全的復元.

$^{361}$康復的意思.
原來耶穌基督說.
他受的刑罰.
我們得平安.
很重要的.
耶穌基督說.
我親身經歷過.
然後我又受過刑罰.
在刑罰之中.
我忍受了.
我戰勝所想.
我帶了平安給你們所有人.
有一個仁兄.
在香港.
很小的時候.
他學壞了.
加入黑社會.
差不多吸毒了.
他的父母就覺得.
不好啊.
希望兒子可以改.
剛剛移民.
他家人帶他去美國.
移民去美國紐約.
希望兒子可以改.
他對紐約有沒有改?.
沒有啊.
他去了唐人街.
加入了唐人街黑社會.
他告訴我.
唐人街很有趣的.
街頭一幫人.
街尾一幫人.
大家互不侵犯.
河水不犯井水.
一打架有事的時候.
他們不像香港那樣拿刀.
他們拿什麼?.
擲牆.
所以在他的環境裡.

$^{401}$他是沒有平安的.
他都不知道哪一天.
他說死在窗底下.
其中一個人.
所以他一直沒有平安.
有一天.
他在唐人街走.
走著走著.
有什麼事發生?.
他突然間被人箍住.
有什麼事情發生?.
被人箍?.
我是黑社會的.
你一搭我的肩膀.
你猜我有什麼反應?.
你搭我那邊.
我不會搭這邊的.
一搭我這邊.
我會搭這邊的.
通常他在那邊搭這邊的.
你一轉過去.
警察是這樣的.
被人一箍的時候.
我們這樣被人一箍.
被警察一箍著.
我們會怎樣?.
一伸手下去.
一轉頭.
那隻手就甩了.
他不懂.
他那時候小孩子.
他插在他旁邊.
被人箍著.
原來箍著的那個人.
是什麼人?.
什麼人?.
是警察.
你現在被人襲警.
把你拉進去.
美國很有趣的.

$^{441}$小孩子和老人家.
他們是很優待的.
他小孩子.
寫了名字.
記錄了.
放了他.
他是害怕的.
他回到家裡.
跟家人說我做了這件事.
他的對策是怎樣?.
他有什麼辦法?.
他跟媽媽說.
我不在美國了.
你在美國就要坐牢.
他說我回香港.
那個時候.
我已經拿了綠卡.
他十幾歲就回香港.
回來香港.
有沒有做好人?.
沒有.
比以前更差.
起手吸毒.
吸得很深.
還要坐四次牢.
到沒辦法的時候.
他家人有什麼提議?.
他爸爸媽媽說.
兒子.
這樣吧.
你回去美國吧.
他跟媽媽爸爸說.
回什麼?.
我拿了綠卡的時候.
我十幾歲拿了綠卡.
回到美國我三十二歲回去.
二十九還是三十回去.
他爸爸媽媽說.
你試試吧.
那個人講見證告訴我.

$^{481}$他說他那張綠卡.
是小孩上來的.
他回去的時候三十歲.
樣子跟不一樣?.
一看樣子就知道.
不讓你進去.
你的樣子也不一樣.
但是沒辦法.
他硬著頭皮.
因為家人要他過去.
他回去美國.
過境的時候.
讓不讓他進去?.
他讓他進去.
所以你答沒有問題.
你不要怕答對還是錯.
沒有對還是錯.
我希望你們留心聽.
你留心你叫個頭.
你就會有精神.
又讓他過去.
過去了回了美國.
有沒有做好人?.
這次你很大聲主席.
沒有.
他繼續回去唐街.
但心情更加不平靜.
更加沒有平安.
因為他有舖毒品引過美國.
他在美國要找毒品吃.
他整個人都在情緒.
在家裡面.
他得不到平安.
坐也不是.
站也不是.
走到唐街又害怕.
一天他的哥哥就跟他說.
親哥哥.
不如這樣吧.
弟弟.

$^{521}$你已經十幾二十年了.
吸毒坐牢.
你回來又是吸毒.
我帶你去其他州戒毒.
他去不去?.
他去.
他很煩.
因為他沒有平安.
希望不要再吸毒.
因為不吸毒就不用去唐街.
不用去唐街的時候.
他很平靜很平安.
他去另一個州戒毒.
他告訴我戒了三個多.
差不多四個星期.
他戒掉了.
戒掉了之後怎樣?.
他仍然沒有平安.
他仍然心很虛.
就是這個人.
他經常問自己.
我什麼時候會吃飯?.
我什麼時候走到唐人街裡面.
又逃過那幫人一起.
我什麼時候會被人打?.
我什麼時候會受傷?.
所以他完全沒有平安.
但是有一天.
他走過一間教會.
他看一看教會.
有人在那裡派單張.
邀請人進去聽福音.
美國是這樣的.
在教堂門口就這樣派.
他拿著單張進了教堂.
你知道新人進來教堂.
我們真的好好招呼.
我們有牧師跟你聊天.
有執事跟你聊天.
去到散會的時候.

$^{561}$牧師單獨跟他聊天.
叫他信耶穌.
他說你信耶穌就有平安了.
他說我怎麼信耶穌?.
牧師說你祈禱吧.
你祈禱求耶穌幫你.
他聽了之後.
我相信聖靈感動他.
他回到家.
他晚上真的在祈禱.
你知道他在祈禱什麼嗎?.
祈禱很有趣.
他跟我說見證的時候.
他的祈禱就是這樣.
他說耶穌.
如果你真的是耶穌是神.
你就幫我明天起床.
不用吸煙.
我們知道戒毒就容易過.
戒煙就很難過.
很難戒的.
我們吸毒的時候.
或者吸煙.
我的包應該在床頭.
起床就拿著吸煙.
第一次吃完飯就吸煙.
那你在說.
耶穌啊.
拜託你明天幫我.
我不會吸煙.
行不行?.
他第二天我起床看著爆煙.
我平時起床就吸煙.
但我起床後.
對爆煙完全沒有感覺.
他知道這是耶穌幫了他.
他突然內心平靜.
安穩又平安.
之後他做什麼?.
他回到教會.

$^{601}$找了一個牧師.
順著耶穌洗禮.
他開始很輕鬆.
他不覺得有威脅.
他走出街上的時候.
他內心平靜得不得了.
不過我想問.
你猜他在美國做什麼工作?.
72行都沒有.
他的工作不是72行裡面的.
他當兵.
當什麼兵?.
水兵.
他退休了.
剛剛退休三年.
你想想.
原來耶穌基督說.
我身心經歷了苦難.
在苦難裡面我受刑罰.
在刑罰受過之後.
我就帶出一種平安給你們每一個人.
他叫Ken.
他得到平安.
他在痛苦裡面.
他不知道怎樣做.
然後他突然間信了耶穌.
在耶穌裡面.
他親身經歷過.
從一個不平安裡面.
走到平安裡面.
人全身輕鬆了.
所以耶穌基督說.
我受的苦.
是要來幫你和我.
幫世界上所有的人.
耶穌基督親身經歷過苦難.
然後他受的刑罰.
使我們有平安.
第三.
苦難裡面得醫治.

$^{641}$因他受的鞭傷.
我們得醫治.
得醫治的意思.
是屬靈生命裡面的醫治.
耶穌說.
我會給你平安.
我亦都會醫治你.
不單止給你平安.
肉身裡面有平安.
你的靈魂裡面你都得到醫治.
請大家留意這兩個.
這對夫婦是我太太.
去年應該是六七月的時候.
就去加拿大的西北角那裡.
加加里.
她那裡分享見證.
她分享之後.
不知道怎樣.
溫哥華的人.
好像有人聽到.
就傳了到溫哥華.
就是有人請她.
馬上到溫哥華廣場.
她認識這對夫婦.
這對夫婦很好.
她去年應該是.
十一月左右來了香港.
十一月左右來到香港.
她請我喝茶.
不知道吃飯什麼.
我們見面之後.
第二次見面.
最後的時候.
她開口問我.
她說天真牧師.
她都知道我叫天真牧師.
我黑社會叫天真仔.
所以在監房裡面.
沒有人會叫我.
叫我做李建明牧師的.

$^{681}$多數叫我做天真牧師.
她說天真牧師.
我有件事想請求你.
我有什麼事.
她說我有個親戚.
有個表弟進了我的監房.
你可不可以去探望他.
我當然可以.
我有個特別者.
可以隨時出隨時入.
我進監房.
可以凌晨兩點鐘都可以進去.
沒有人不讓我進去.
但是兩點鐘你進去幹什麼.
人家都睡著了.
你幹什麼進去.
我說可以沒問題.
我去探望你的表弟.
我記得我在唐福那裡.
在大嶼山唐福懲教所.
我進去看到這個人.
我第一樣.
我的做法就是.
我一進去.
大約我的開場白就是這樣.
我是黑手的.
什麼意思.
我坐了很多次監.
我是慈雲山十三太保茅頂華的小弟.
一聽到華哥的名字.
人人都很出奇.
你是他弟弟嗎.
因為華哥很出名.
現在也很出名.
第三我就會告訴你.
我是吸毒的.
有八次案底.
坐三次監.
監視了八分之十年.
我通常開場白是這樣.

$^{721}$但是我知道這個人很特別.
所以我沒有開場白.
我問他.
原來這個人很特別.
這個人是怎麼樣.
他原來在加拿大讀完大學.
然後去英國讀碩士.
這對夫婦的表弟.
就是這個男方.
Joseph的表弟.
他讀完碩士回來香港.
他的專業水準很厲害.
他是電腦工程師.
專門應付病毒.
他在英國讀完碩士回來香港.
就做很專業的事.
接著在香港洗禮受浸加入教會.
我就問他說.
你這麼好 為什麼你進來.
他告訴我原因是怎麼樣.
我真的相信他.
他說原來他有憂鬱症.
憂鬱症他去買藥.
這個憂鬱藥很貴.
所以我告訴他.
病毒可以代替.
他真的試試病毒行不行.
真的行.
變相病毒便宜的時候.
就省了很多錢.
所以有一天.
我代替了二萬多元.
所以我就全部買了病毒.
打算儲糧.
儲糧吸毒打仗.
那次一拿了二萬多元的毒品.
就這樣抓了他.
很悲慘 他很頹喪地告訴我.
那時他很悲慘.
一個很專業的IT人才.

$^{761}$那時很興旺的病毒.
我就告訴他.
你不用擔心.
我有八次運動 坐三次監獄.
我信耶穌解毒.
我叫天真牧師.
你知道嗎.
他聽了我的名字之後.
他有什麼反應嗎.
然後呢.
我說我是天真牧師.
他兩隻手直接捉住我的手.
不停地搖.
哎呀 我很榮幸.
我很高興認識你.
我很想認識你.
很想見你很久.
我就很出奇.
我說麻煩你 先鬆手.
我說你認識我的.
我不認識你的.
為什麼我一說了我的名字.
你這麼興奮.
他告訴我一件事.
當我被抓了之後.
他在荔枝角收押所.
一進去 他第一次.
他整個人沒有心思了.
他想想.
為什麼我會被抓進來.
因為他藏的毒品很厲害.
要判很重的.
所以整個人沒有精神.
你猜他想什麼呢.
他想自殘.
他想得很灰.
很不想留在這個地上.
裡面的囚友有很多很好的.
其中有一個是在囚人仕途一起.
應該是同房還是什麼.

$^{801}$因為同一個地方.
看到他這樣的時候安慰他.
他說你不用擔心.
你不用想一邊.
有個牧師叫天真牧師.
他比你還壞.
他是黑手的.
八次判他坐三次監.
吃了百分之十年.
他信耶穌戒毒.
現在是牧師.
經常進來的.
所以他聽了之後.
因為他信耶穌.
他覺得耶穌可以幫到他.
他回到耶穌那裡.
他得到醫治.
他不會再想那些負面的東西.
所以整個人.
很完全的承擔.
他犯的罪的責任.
所以進去見到他的時候.
他很開心.
各位等姐妹.
所以你受苦.
困難.
人生必然的.
爭在心還是錢.
大痛就是小痛就是小痛.
但你不要覺得痛苦苦難.
或者處境裡面.
困局裡面.
你覺得很慘.
不要這樣想.
耶穌說我開一條出路給你.
那條出路是什麼呢.
要你經過苦難.
你經過苦難.
你才可以幫人.
好像那個牧師.

$^{841}$姜牧師.
現在見到神.
因為是講師.
他經歷了幫人.
我呢.
我又經歷了苦難.
很痛苦.
那你知不知道吸毒.
家裡不要你.
一來到當你是廢人.
透明人.
那種痛苦.
很慘的.
但是.
信耶穌.
耶穌說我的痛苦比你更痛苦.
我們跟耶穌對比.
耶穌的痛苦.
我們是濕濕的.
但有一個好處就是.
耶穌給了我們一條出路.
我用一段聖經做作結出.
多傳書十章十三節.
你們所遇見的試探.
無非是人所能受的.
神是信實的.
必不叫你們受試探.
過於所能受的.
在受試探的時候.
總要給你們開一條出路.
叫你們能忍受得住.
剛才希伯羅書二章十八節.
在說什麼.
在說我受試探.
我是過受試探而受苦.
我遇見試探而受苦.
就能夠達到受試探.
試探就等同受苦.
現在我們說.
你們所遇見的受苦.

$^{881}$無非是人所能受的.
神是信實的.
必不叫你們受苦.
過於所能受的.
在受苦的時候.
必給我們開一條出路.
叫我們能忍受得住.
所以各位弟兄姊妹.
今天最好了.
今天復活節.
今天復活節的時候.
我們看到耶穌基督.
在受苦 釘在十字架.
兩個種的痛苦.
學到復活裡面得勝.
得勝 以致我們信耶穌.
你知道嗎 耶穌是神醫.
醫好我們的生命.
醫好我們的性情.
醫好我們的性格.
我經歷過痛苦.
你知道我有多少個博士學位嗎.
有沒有留意.
兩個.
一個瘋的.
榮譽博士.
一個自己讀的.
一個吸毒.
坐牢有問題的人.
怎樣去到這個地方呢.
我七十二歲才拿過博士學位.
所以我的痛苦.
我的受苦.
我的傷痛.
我能夠站起來.
我怎樣信耶穌.
去到能夠拿博士.
我幫了很多人.
因為我是親身經歷.
我又叫了很多人讀神學.

$^{921}$現在已經有一個讀完神學做牧師.
是殺人犯.
現在另一個在中國.
中文大學崇天神學院.
正在讀碩士學位.
第二個.
第三個可能出來了.
他們個個都受過痛苦.
所以他們知道怎樣去幫助有問題的人.
各位弟兄姊妹.
今天很開心.
為甚麼開心.
知不知道為甚麼開心.
今天是復活節.
是得勝的日子.
另外還有一件事.
耶穌說.
我由痛苦在十字架上.
最傷痛的時候.
能夠經歷.
感受.
為你.
為我開一條出路.
以致我們在這個地上.
不會怕的了.
我們一起祈禱.
天父我們多謝你恩待我們.
求你祝福我們.
讓我們都能夠在經歷.
苦難.
痛苦.
困境的裡面.
能夠.
靠著耶穌開的出路.
我們戰勝.
我們能夠見證神的恩.
求你祝福我們.
奉耶穌名祈禱.
阿妹.
\newpage



\section{以賽亞書 52:13-53:6}
\label{sec:031927}
\textbf{受苦義僕}~~~~李應新牧師
\newline
\newline
連結: \href{http://www.church.com.hk/acms/content.asp?site=cdc\&op=show\&type=product\&code=031927\&layout=sermon}{\texttt{031927}} ~~~~ 語音日期: 2026-04-05
\newline
\newline
\hyperref[sec:031926]{\small{< < < PREV SERMON < < <}}
~
\hyperref[sec:index_chronic]{\small{[返順時目]}}
~
\hyperref[sec:031928]{\small{> > > NEXT SERMON > > >}}
\newline
\newline
以賽亞書 52:13-53:6
\newline
\begin{longtable}{cl}
\hline
\hline
章節 & 經文 (和合本修訂版)\\
\hline
52:13 & \begin{tabularx}{0.7\textwidth}{X} 看哪，我的僕人行事必有智慧，他必被高升，高舉，升到至高之處。 \end{tabularx} \\ \\ \relax
52:14 & \begin{tabularx}{0.7\textwidth}{X} 許多人因他驚奇---他的面貌比別人憔悴，他的外表比世人枯槁--- \end{tabularx} \\ \\ \relax
52:15 & \begin{tabularx}{0.7\textwidth}{X} 同樣，他也必使許多國家驚奇，君王要向他閉口。未曾傳給他們的，他們必看見；未曾聽見過的事，他們要明白。 \end{tabularx} \\ \\ \relax
53:1 & \begin{tabularx}{0.7\textwidth}{X} 我們所傳的有誰信呢？耶和華的膀臂向誰顯露呢？ \end{tabularx} \\ \\ \relax
53:2 & \begin{tabularx}{0.7\textwidth}{X} 他在耶和華面前生長如嫩芽，像根出於乾地。他無佳形美容使我們注視他，也無美貌使我們仰慕他。 \end{tabularx} \\ \\ \relax
53:3 & \begin{tabularx}{0.7\textwidth}{X} 他被藐視，被人厭棄；多受痛苦，常經憂患。他被藐視，好像被人掩面不看的一樣，我們也不尊重他。 \end{tabularx} \\ \\ \relax
53:4 & \begin{tabularx}{0.7\textwidth}{X} 他誠然擔當我們的憂患，背負我們的痛苦；我們卻以為他受責罰，是被神擊打苦待。 \end{tabularx} \\ \\ \relax
53:5 & \begin{tabularx}{0.7\textwidth}{X} 他為我們的過犯受害，為我們的罪孽被壓傷。因他受的懲罰，我們得平安；因他受的鞭傷，我們得醫治。 \end{tabularx} \\ \\ \relax
53:6 & \begin{tabularx}{0.7\textwidth}{X} 我們都如羊走迷，各人偏行己路；耶和華使我們眾人的罪孽都歸在他身上。 \end{tabularx} \\ \\
[1ex]
\hline
\hline
\end{longtable}
$^{1}$各位平安.
今日是復活節慶祝的主日.
所以看一看旁邊.
如果我告訴你,主復活啦!.
有點氣勢,主復活啦!.
我回來是很開心的.
剛才你說沒甚麼神氣.
主復活啦!主復活啦!.
興奮一點.
我們今天又看一段經文.
首先說一下,是否每個人都拿了一本聖經回來?.
有誰未有聖經?.
監禁你啊!.
有些流亡之餘.
流亡之餘,流亡之餘.
各位有本聖經.
有些流亡之餘.
我有手機,手機大了.
拿多幾本不怕.
流亡之餘.
一人一本,拿多些.
手機不算.
因為我和你一起.
戒手機.
戒一個小時都好.
前面,這位.
這位也是.
這位也要.
一人要一本.
每人拿一本聖經.
既然給了你,不斷戒.
還有誰未有聖經?.
有了嗎?.
後面還有些人未有.
我們一起戒聖經.
很久,上次我.
無論如何.
我們今天戒手機.
戒一個小時.
我們看聖經.

$^{41}$我們今天看這段經文.
叫做《僕人之歌》.
《受苦的二僕》.
在《爾塞亞書》.
《爾塞亞書》有多少首《僕人之歌》?.
總共有四首.
這是第四首.
我和你看完它.
看看是不解的.
知道原來是這樣.
學者研究完才看到.
徒手看經都看不見.
原來有四首.
所以要上主日學高高高.
不要緊.
下星期上主日學.
我們現在看第一首《僕人之歌》.
第四十二章.
我們看看第四十二章.
我們只是看.
但是.
我們也感受一下.
第四十二章.
第一到第四節.
其實不止,到第九節.
不過我們只是看四節.
我們品嚐一下.
為什麼叫做《僕人之歌》.
《爾塞亞書》有四首.
第一首.
如果旁邊不習慣揭聖經就幫幫他.
看看.
唐容聖經九百六十八頁.
九六八.
揭一下,慢慢,慢慢.
知道你很久沒揭過.
慢慢,等你,等你.
好,我們看四十二章.
第一節.
看啊我的.

$^{81}$僕人,對,我所扶持.
所揀選心裡所喜悅的.
我要將我的靈賜給他.
他必將功利傳給外邦.
他不圈洋也不揚聲.
也不使街上的人聽見他的聲音.
壓傷的路位他不折斷.
將殘的燈火他不吹滅.
他憑真實.
將功利傳開.
他不灰心不喪膽.
直到他在地上設立功利.
海島都等候他的訓誨.
所以說到.
這位僕人是怎樣.
他是神所揀選的.
將靈賜給他.
特別留意,特別的地方是甚麼.
他必將功利傳給外邦.
即是在以色列人的心目中.
他們是神的選民.
所以好東西就是給他們.
不是啊.
神的真理其實一早與期望.
是給外邦的.
萬國要因為神的子民得福氣.
所以是給外邦的.
而且這裡說到.
這位僕人壓傷的路位.
不折斷,將殘,燈火不吹滅.
即是怎樣.
很溫柔的.
有和平的君王.
不是那些用權力逼你的.
是很好的.
不灰心不喪膽.
神給他的任務是將功利.
設立海島等候他的訓誨.
這首是第一首僕人之歌.
我們看第二首.

$^{121}$我們揭下去第四十九章.
四十九章,揭揭後面.
四十九章.
多揭幾頁就到了.
我們這次看五節六節.
也是不止的.
不過是少少的.
四十九章,五六兩節.
我們讀出來,五六兩節.
四十九章,五六兩節.
好了,我們讀了.
耶和華從我出胎.
奏奏我作他的.
僕人.
要使雅各歸向他.
使以色列到他那裡聚集.
原來耶和華看我為尊貴.
我的神也成為我的力量.
現在他說.
你作我的僕人.
使雅各眾之派復興.
使以色列人中.
得保存歸會.
尚為小事.
我還要使你作外邦人的綱.
叫你施行我的救恩.
直到地極.
這裡同樣重複出現是什麼呢?.
原來神揀選的,出母胎的時候.
神已經委派他.
是上帝所揀選成為僕人.
復興誰?復興.
神的子民,雅各和以色列.
復興子民.
但不是就這樣停手的.
復興他們,神作他們的力量.
使他復興之後.
其實只是小事.
還要使他作.
外邦人的綱.

$^{161}$所以將救恩要帶給外邦.
也將救恩要帶到哪裡?.
地極,所以這裡也明顯.
重複呼應著什麼?.
這位僕人.
帶到外邦,帶到地極.
但先復興神的子民.
神的子民復興的原因.
目標是什麼?將救恩.
將功利帶到萬邦.
到地極,所以行不行?.
不只是自己有福氣.
而是將福氣要帶給.
未認識主的人.
神的僕人,第二首僕人之歌.
我們繼續看下去.
第三首,五十章.
四至五節一揭後面就是了.
第五十章.
我們只是讀兩節.
四至五節.
這次沒有那麼明顯.
但後後前前我們都看到.
應該都是僕人之歌.
我讀給大家聽.
主耶和賜我受教的舌頭.
使我知道怎樣用言語輔助.
疲乏的人,主每早晨提醒.
提醒我的耳朵.
使我能聽,像受教者一樣.
主耶和開通我的耳朵.
我並沒有違背.
也沒有退後.
人打我的背,我任他打.
人拔我腮甲的鬍鬚.
我由他拔.
人辱我,吐我.
我並不掩念.
我們跳一跳,看少少後面那裡.
講到十節.

$^{201}$你們中間誰是敬畏耶和華.
聽從他僕人之話的.
這人幸在暗中,沒有亮光.
當依靠耶和華的名.
仰賴自己的神.
我這裡想提醒一下.
我們學者都認為.
這首都是僕人之歌.
使他明白神的話語.
又用神的話語.
輔助疲乏的人.
而且神的話語每天都提醒.
另外一個特別的地方.
人辱我,吐我.
並不掩面.
你又會感受到誰.
主耶穌,被吐湯沒.
也沒有掩面.
我們看到神的僕人.
這裡所指的.
耶穌基督也顯明出來.
我們再看第四首僕人之歌.
我們剛才讀過.
五十二章.
再重覆這個主題.
五十二章.
十三節開始.
五十二章.
十三節.
我的僕人行事必有智慧.
必被高舉.
上升且為至高.
先停一停.
先講背景資料.
我們看完四首歌.
我們都感覺到.
誰最正確.
主耶穌最正確.
但有一個問題出現.
沒理由.

$^{241}$以賽亞與主耶穌爭幾百年.
沒理由幾百年後.
人們才明白.
以為先正確.
幾百年後.
聽以賽亞的說話.
完全與他無關.
沒理由.
這個僕人是誰.
大家明白我的意思嗎.
這是論文級題目.
不知道是誰.
可以論文級討論.
有幾位參選人.
都不太合身.
包括古烈王.
古烈王是誰.
是波斯的皇帝.
波斯皇帝有什麼特別.
神的.
這個參選人.
以賽亞書有提名.
說他是神的僕人.
為什麼.
因為他釋放以色列人.
避老回歸.
但他是不是信耶和華.
有些問號.
敬畏耶和華.
他周圍都做這些事情.
其他國家都是.
讓以色列人回鄉.
他不僅是.
給以色列人.
但這也成就了神的心意.
神說了七十年後.
那時候回歸島.
這是上帝的僕人.
這是其中一個參選人.
不過要多說少少.

$^{281}$我們很難.
那個人不是完全適合神的心意.
但神也可以使用.
所以有些人來陰我.
穆斯這個政要.
國家領導人做了這件事.
他是不是神的僕人.
這些問題陰我.
我當然不回答你.
有些對有些不對.
先說少少.
這是其中一個參選人.
但有一件事他沒辦法中的.
因為他沒怎樣受苦.
所以我們其他的都中不了.
這是其中一個.
還有哪些參選人呢.
以賽亞先知自己.
有些人覺得是的.
但有些人覺得耶利米比較像.
耶利米受苦多了.
所以覺得耶利米才像.
這些都是其中一個參選人.
還有一個.
其中一個參選人.
也有點像.
那個叫索羅巴伯.
索羅巴伯跟我們不是很朋友.
不太熟悉他.
跟他同一代的人.
我們比較熟悉.
尼希米和以斯拉.
我們比較熟悉.
因為他有尼希米記和以斯拉記.
所以我們熟悉他的名字.
索羅巴伯是跟他同一代人.
我跟羅姑娘文姑娘.
是同一代人.
跟哪些不是同一代人.
包包會員牧師.

$^{321}$我仰慕.
但我不是同一代人.
那些長一點.
陳宗道.
危託星.
那些當然更長一點.
不是同一代.
剛才說了三個同一代.
我小小感受一下.
因為他們覺得索羅巴伯有點像.
索羅巴伯做了什麼.
尼希米你比較熟悉.
尼希米做什麼.
重建城牆.
一手拿兵器一手作功.
所以他都很重要.
不過他做這件事的時候.
對以色列人都有自己的好處.
不用被人攻擊.
幫助自己.
還有一個人很有益處.
就是以斯拉.
以斯拉做什麼.
對以色列人有好處.
是宣讀律法書.
帶律法書去耶路撒冷.
是否重要.
也是重要.
那索羅巴伯又做了什麼.
最重要的.
重建聖殿.
其實他是同一代人.
最早期做省長派過去做.
但不是代表很多人很歡迎.
因為其他民族不滿意.
所以他受很多威脅.
很多攻擊.
重建聖殿.
先知都有說.
有幫他.

$^{361}$他住的屋有屋頂.
但耶和華的殿荒涼.
於是乎他是其中.
最重要的推手.
推動重建聖殿.
索羅巴伯.
都有點相似.
這幾個合適.
但沒有一個中了.
不過總之都是其中的候選人.
當然再說一次.
我們中的是主耶穌.
但這是幾百年後.
我們遙遙看見.
受苦的僕人.
主耶穌.
行事有智慧.
這不是人間的智慧.
其實在保羅書信都有提到.
神的智慧.
神給的智慧.
能參透萬事.
沒有人能明白他.
我們平時的智慧.
跟神的智慧不同.
神的智慧.
是讓受苦的僕人.
被高舉.
許多人因他驚奇.
想都想不到.
他的面貌被別人潮水.
他的形容被世人Google.
只讓他必洗淨.
許多國民軍王要向他閉口.
因所未曾傳予他們的.
他們必看見.
未曾聽見的他們要明白.
他不是很聰明.
後面更明顯.
我們所傳的有誰信.

$^{401}$說你都不信.
耶和華的蚌臂向誰顯露.
他在耶和華面前.
生長如嫩芽.
即是很嫩的.
象根出於乾地.
嫩芽在乾地生長.
即是生得很艱難.
但他無佳形.
美容.
我們看見的時候.
也無美貌.
使我們羨慕他.
即是看起來不是很風光.
很靚仔.
所以特別要提多.
近年家會.
特別我們的.
電視節目.
被人批評.
叫做成功神學.
即是那些Prince.
是甚麼的.
說他成功神學.
經常說你做多些生意.
管高些.
令多些人受到祝福.
神會賜福給你.
做生意做大些.
請多些人.
又不是說完全錯了.
但神的心意.
或者受福的僕人.
無佳形美容.
沒有說你不准做君王.
神也有些合神心意的君王.
那個大衛.
你做君王合神心意.
不過你要合神心意才行.
最高位最繁榮的時候.

$^{441}$犯罪.
受神審判.
重點不在那裡.
我們今天都有些學習.
成功是可以.
不成功.
無佳形美容.
神也可以使用.
重點是甚麼.
忠心良善先正是.
所以我們要做忠心良善.
要效法主受苦的僕人.
受苦的義僕.
所以我不能不想起.
我自己成長教會的一位叔叔.
那位叔叔是一位.
在教會做了幾十年.
這個工作我們教會是沒有的.
知道甚麼是唐役嗎.
是古老教會.
大大間才有的.
唐役是負責.
打理這間唐的事務.
即是說唐役.
掃地.
執拾.
教會爛了換燈打他.
掃地什麼都是他.
他和牧師一樣.
都是住在教會.
其實我也做過這些教會.
他就住頂樓.
他就住在門口.
不是說我好一點.
我住頂樓.
那些樓多數都沒有升降機.
住頂樓走到氣都咳了.
不過我還年輕還可以.
唐役這位叔叔.
我認識他的時候.

$^{481}$他已經在這裡工作.
他一直做到退休.
他就是一位.
東南亞.
我忘記他從哪裡來的.
他在香港生活.
所以他的廣東話我沒辦法聽得懂.
我不知道他在說什麼.
但我知道他在說廣東話.
不過聽不懂.
但無所謂.
他負責清理教會.
做教會的事務.
他不懂廣東話.
但他做了幾十年.
做到退休.
直至他退休的時候.
教會也有為他開一個歡送會.
所以我也是校法.
陳姑娘歡送會記得來.
還有亭康的歡送.
我們也來.
我也是學學別人.
他歡送會的時候.
我已經出來事奉了.
我有針在這裡.
很多年輕人很捨不得.
這位叔叔退休.
這麼厲害.
牧師退休都沒有年輕人眼濕濕.
這麼嚴重.
這麼厲害.
我有收封的.
過人之處我也想學校法.
我是學不到的.
不過我也告訴你.
他做什麼.
是負責什麼.
還有一個工作.
打掃了關門.

$^{521}$教會的關門.
他是key man.
在門口關門.
但他經常都不依照指示.
指示是十點鐘要關門趕人走.
但他不趕.
年輕人喜歡留到幾點.
所以不要太早睡.
留久一點.
所以他不聽指示.
學生很喜歡他.
教會可以讀書.
打籃球.
有人說讀書.
但打籃球為主.
但他不趕你走.
再說一次.
我聽說他收過很多投訴.
浪費教會的電.
這個後話.
他殺食.
他家裡是怎樣.
他在家裡住.
他也會煲糖水給自己的兒子吃.
他煲的很大一鍋.
真的大一鍋.
吃了誰.
在那裡的年輕人.
讀書的年輕人.
來吃糖水.
很有趣的.
自己媽媽煲的.
兒子不吃.
人家煲的比較好吃.
很多年輕人都領受過他的恩惠.
是不是很好.
這個不是最重要.
殺食的地方.
後來針對他.
我們很多年輕人.

$^{561}$喜歡跟這位叔叔講心事和秘密.
不講給牧師聽.
都講給他聽.
很多心事.
喜歡哪個女孩子.
講給他聽.
只是講給他聽.
重點在哪裡.
不是.
聽不明白.
就算到處講都聽不懂.
所以好厚密.
真的絕.
上帝好使用.
重點是什麼.
忠心良善 上帝會如立.
不一定要做高位.
做基本的工作.
忠心良善 上帝都可以使用如立.
我們孝法的是誰.
孝法主耶穌 受苦的義僕.
我們繼續看下去.
我們看第三節.
厭氣 多受痛苦.
常經憂患 他被藐視.
好像人掩面不看一樣.
我們也不尊重他.
第一件事 今天我重要.
希望你記得.
我們也不尊重他.
因為他不是很威風.
不是很帥 是不厲害.
不厲害到一個地步.
我們也不尊重他.
我特別想你記住這句金句.
你今天回去的時候.
慢慢回答.
說那些以色列人.
不尊重他 同一代的人.
但今天這句金文.

$^{601}$突然憋在我心上.
很重要 想和你說.
我們信神的子民.
我們是神的子民.
我們有時看慣看熟.
聽慣聽熟 我們也不尊重他.
你說是嗎.
牧師 哪裡說的.
以賽亞先知也說.
二十九章說 你們用嘴唇尊敬我.
心卻遠離我.
馬太福音也是這樣說.
主耶穌也是這樣說.
你們這些子民.
用嘴唇尊敬我.
心卻遠離我.
我們也不尊敬他.
這裡我今天最重要.
希望你今天回去的時候.
慢慢回答.
其實我們每個星期回教會.
是不是真的尊敬.
尤其是經文說的是敬拜神的人.
或者做神事奉的人.
馬太福音明顯說的是事奉神的人.
教人律法等等.
但其實心沒有尊敬神.
這個我自己要提醒.
我自己要提防.
把事情習慣了.
習慣了就覺得沒什麼.
不行就這樣.
看慣了就麻木.
小心提防.
我們看慣看熟.
上帝話語教訓的話.
我們會覺得是這樣.
特別的就是 都是不行的.
這是我們沒有尊敬神的話語.
那我真的要跟進.

$^{641}$我不行是我現在不行.
不知道什麼時候可以.
這裡我還想說.
我們的價值標準.
有時候分開了.
你怎樣來判斷.
怎樣來評估.
舉一個例子.
我們教會有一對夫婦.
做家庭的服侍.
他們專門搞claim.
叫夫婦建立婚姻關係.
很有意義.
但我有時候都會鼓勵人.
去吧 去這個camp.
是好東西.
有一對夫婦我叫了她很多年.
她一開始.
不知43年.
這麼長怎麼去.
現在3日2夜短了.
3日2夜去到了吧.
她有另一個樣子.
牧師 這麼貴 幾千元.
她經常覺得教會的東西.
不用錢.
她去酒店住3日2夜都要錢.
怎麼這麼貴.
她一直咬著我.
這麼貴都不去.
我剛剛才知道.
現在這個復活節.
去郵輪一家四口.
還要搭船去日本上船.
那又不貴.
去這些就貴.
你什麼標準.
小心提防.
不過我還想說多一件事情.
不是完全對的.

$^{681}$我們用什麼心情去敬拜主.
我們難明.
為什麼我們這麼多人的罪可以歸到他身上.
我上星期查經.
我跟其他教會的同工去查經.
查羅馬書的.
查經的時候.
講到這件事的時候.
有一個同工.
因為我不是原創商.
是人家說的.
我覺得挺合理的.
為什麼眾人的罪可以歸到他身上.
那是信用卡.
你刷卡刷了很多.
但是一次過幫你還了.
那叫什麼.
你們知道的.
一筆清.
就沒有了.
信用卡的原理.
這樣想好像很明顯.
不過我再答真一點.
其實主題不是信用卡.
其實是貸款卡.
香港是不流行的.
貸款卡.
外國比較流行.
貸款卡.
知道是什麼嗎.
貸款卡.
付了錢.
主題是付了代價.
不過問題是.
這麼多人.
這麼多罪孽.
看一看這張信用卡.
是誰給你附屬的.
現在我們附屬信用卡.
看一看.

$^{721}$如果那張信用卡.
是成哥給我附屬.
那我真的不擔心.
行不行.
成哥附屬.
我的信用卡.
是成哥附屬信用卡.
讀吧.
因為主為我們成就救恩.
普天下的人.
不但為我們的罪.
也為普天下人的罪.
其實是遇上普天下的人.
不過我們願不願意領受.
繼續說下去.
他被欺壓.
受苦的時候卻不開口.
他像陽谷被遷怒.
宰殺之地.
像陽在展無的人.
手下無聲.
他也是這樣不開口.
因受的欺壓和審判.
他被奪去.
至於他同世的人.
誰想他受鞭打.
從活人之地被剪除.
是因我百姓的罪過呢.
第九戰.
他雖然未行強暴.
口中也沒有鬼宅.
人還是他.
與惡人同埋.
誰知死的時候.
與財主同葬.
又何時應驗.
他受埋葬的時候.
與那個叫做若薪.
他是一個財主.
向彼拉多領主的身體.

$^{761}$安葬在他的身份無裡面.
所以與財主同葬.
耶和華決定.
以將他壓傷.
使他受苦.
耶和華以他為贖罪制.
他必看見後罪.
並且連長連日.
耶和華所起的事.
必在他手中亨通.
他必看見.
自己老虎的功效.
便心滿意足.
有許多人.
因認識我義僕.
得稱為義並且.
他要擔當他們的罪孽.
所以我要使他.
與偉大的同憤.
與強盛的均分老民.
因他將命傾倒.
以致於死.
他也被列在罪犯之中.
他卻擔當多人的罪.
又為罪犯代求.
路加福音又中埋.
路加福音講到GSO.
為了那個罪犯.
願意悔改.
今天你要與我同在樂園了.
不過我特別想.
highlight第二戰.
剛才說.
我們小心.
不要不尊重他.
第二戰.
要使主心滿意足.
第十一節.
我們特別highlight.
他必看見.

$^{801}$自己老虎的功效.
便心滿意足.
今天我第二件事.
希望你默想.
想清楚.
回到離開禮拜堂的時候.
都要想著.
我都要使主心滿意足.
主怎會心滿意足.
就見到他老虎的功效.
就會心滿意足.
那老虎的功效又是什麼.
因為許多人.
因為他得救.
罪孽在他身上.
先解釋一下.
什麼叫心滿意足.
心滿意足.
我接觸很多長輩.
特別是上一代的長輩.
他們自己都窮窮地.
又沒有讀書.
有很多環境因素.
但是他們的子女.
讀到書.
有飯吃.
他們就心滿意足.
捱得很辛苦都OK.
你敢不敢.
見過這些長輩.
其實我們上一輩的人.
很多都是這樣.
自己心滿意足.
我也說說我自己.
我自己心滿意足.
其實我有點擔心.
我講出來.
我怕你說我曬命.
不過.
大家弟兄姊妹都不要這樣.

$^{841}$就讓我曬陣.
其實養我兒子.
我有時會心滿意足.
我兒子幫我做一件事.
令我心滿意足是什麼呢.
其實我有追姦盤特出病.
這個病是需要.
物理治療師說.
每天要拉筋半小時.
我其中有一個動作.
一個動作而已.
我是做不到.
做得很艱難.
就是用一條繩子.
吊起自己的腳.
然後就拉他.
就是拉起他的腳.
那一刻就是這樣.
躺在那裡.
拉他的腳.
我是拉過自己.
治療師教過我.
但是我拉到頸都痛.
所以應該是做錯了.
不夠.
總之不對.
後來我就說.
兒子.
你這麼健碩.
不如幫我壓下去.
他不用拉.
壓下去就可以了.
我把他的腳壓下去.
他每天幫我壓.
拉十分鐘.
這十分鐘是很甜蜜的.
我心滿意足.
這兒子真的值得養.
幫我拉十分鐘.
但是不要誤會.

$^{881}$他不是說.
爸爸什麼時候跟你叫.
不是這樣的.
每天都三催四請.
又要叫兒子.
先吃飯.
先看電視.
先吃飯.
先做功課.
不過他每天都總會做.
夠不夠.
夠了.
你還想怎樣.
我心滿意足.
既然說得兒子.
還要說女兒.
我也是.
知道我是職業病.
因為我經常都探望長輩.
探病.
長輩病到差不多.
又看不到.
又聽不到.
起碼他眼睛不大.
看不到.
又聽不到.
我已經給了一個記號.
我的女兒.
將來.
爸爸去到那個狀態的時候.
你記住你來探望我.
摸一摸我.
頭裡面一個部分.
摸三下.
我就知道你來了.
我不告訴你.
免得你亂摸.
扮我的女兒.
我只是告訴你.
你摸一摸.

$^{921}$摸三下.
我就知道是你來了.
我就會.
我就會什麼.
心滿意足.
你在這裡.
真乖.
來了.
明白嗎.
我想你明白.
心滿意足.
但我想說重點.
心滿意足.
是要令主心滿意足.
心滿意足.
要令主.
因為.
因為為什麼呢.
全福音.
使人未信主人.
可以領受到救恩.
主心滿意足.
其實很多人都未認識福音.
其實今個復活節都提醒我們.
其實許多人要因為受苦的義僕.
能夠得盟救恩的時候.
就可以使自己心滿意足.
親愛的天父.
我求您就賜福給我們.
每一個弟兄姊妹.
讓我們都把握機會.
將您寶貴的救恩.
帶給我們身邊失散的人.
滿足您的心腸.
奉耶穌名求.
奉耶穌名求 阿們.
\newpage



\section{羅馬書 1:1-4}
\label{sec:031928}
\textbf{基督復活的大能 (華)}~~~~葛國光牧師
\newline
\newline
連結: \href{http://www.church.com.hk/acms/content.asp?site=cdc\&op=show\&type=product\&code=031928\&layout=sermon}{\texttt{031928}} ~~~~ 語音日期: 2026-04-05
\newline
\newline
\hyperref[sec:031927]{\small{< < < PREV SERMON < < <}}
~
\hyperref[sec:index_chronic]{\small{[返順時目]}}
~
\hyperref[sec:031929]{\small{> > > NEXT SERMON > > >}}
\newline
\newline
羅馬書 1:1-4
\newline
\begin{longtable}{cl}
\hline
\hline
章節 & 經文 (和合本修訂版)\\
\hline
1:1 & \begin{tabularx}{0.7\textwidth}{X} 基督耶穌的僕人保羅，蒙召為使徒，奉派傳神的福音。 \end{tabularx} \\ \\ \relax
1:2 & \begin{tabularx}{0.7\textwidth}{X} 這福音是神從前藉眾先知，在聖經上所應許的。 \end{tabularx} \\ \\ \relax
1:3 & \begin{tabularx}{0.7\textwidth}{X} 論到他兒子－我主耶穌基督，按肉體說，是從大衛後裔生的； \end{tabularx} \\ \\ \relax
1:4 & \begin{tabularx}{0.7\textwidth}{X} 按神聖的靈說，因從死人中復活，用大能顯明他是神的兒子。 \end{tabularx} \\ \\
[1ex]
\hline
\hline
\end{longtable}
$^{1}$我們起來讀今天的經文.
羅馬書第一章第三第四節.
我想有時候我們在等這個東西.
所以大家要習慣自己拿自己的聖經.
羅馬書第三第四節.
我們起來讀這段聖經.
來 預備 請.
論到他兒子我主耶穌基督 按肉體說是從大衛後一生的.
按聖善的靈說應從死裡復活 以大人顯明是神.
羅馬書這一段聖經裡面他講的基督的人性.
他說他兒子我主耶穌基督.
按肉體說就講的耶穌基督的人性.
從肉身的角度來講.
耶穌是大衛的後一生的.
大衛是以色列裡面他們心目當中的一個最偉大的君王.
所以他是君王的後裔.
預表著什麼呢.
預表著耶穌基督是君王.
但是這位君王他曾經在地上面行走過.
他成為人的樣子.
像我們剛受洗的時候他道成肉身 對不對.
incarnation.
道成肉身就是表明說神.
但是他卻成為人的樣子.
這是一個神蹟.
他不透過人所謂的這種男女的交合.
他是童貞女瑪利亞所懷孕生出來的.
因此他沒有罪.
他是聖靈感孕的.
這些我所要講的這些.
如果不是神是沒有辦法做到的.
今天沒有任何一個其他的人.
可以是這樣子出生的.
那麼我現在要講的重點是.
他按聖善的靈說.
應從死裡復活.
以大人寫明是神的兒子.
這個地方講到耶穌基督的神性.
所以保留在這地讓我們看見一件事情.
就是耶穌基督從死裡復活.

$^{41}$這一件事情的重要性.
我們慶祝復活節 對不對.
剛過去是受難節 是不是.
基督教裡面的這個節慶裡面.
你最喜歡哪個節日.
復活節 真的嗎.
我覺得大部分人比較喜歡什麼節.
聖誕節 為什麼.
因為有大餐可吃.
因為聖誕節有很多好像很熱鬧.
有沒有燈啦這些東西.
反而復活節.
在今天的基督教會有點忽略它了.
可是我們信仰的核心.
不是只是耶穌基督為我們出生.
為我們受死.
耶穌基督也為我們復活.
這是非常重要的.
所以保留在菲利比書第三章.
第十節 第十一節他說.
使我認識基督.
曉得他復活的大人.
並且曉得和他一同受苦.
效法他的死.
或者我也得以從死裡怎麼樣復活.
剛才我們在受洗的時候.
我們為什麼要把他浸到水裡面.
這個聖經上保留在羅馬書裡面提到說.
因為我們受洗是因為歸入他的死.
所以藉著這樣子的一個動作.
這樣子的一個儀式.
象徵著我們的救人老我已經死了.
跟你旁邊說你已經死了.
怎麼到教會來老師叫我死不死的.
就是因為你的救人沒有死.
所以問題就出來了.
夫妻吵架 夫妻這個很多的問題.
兒女的問題 跟別人的問題.
一切的問題都是我們這個救人對不對.
所以這個救人該死不該死.

$^{81}$該死 跟你旁邊說你該死.
這個救人要死了之後.
然後怎麼樣.
所以受洗怎麼樣.
起來是復活.
所以保羅說他要曉得.
我們跟他一同受苦並且從死裡復活.
受難節是在祭列耶穌基督為我們受死.
天主教比較.
我不敢因為他們官方的網站裡面.
他們當然不會這樣說.
可是整個天主教你會發現.
你走進天主教堂裡面去看.
耶穌是怎麼樣.
定在十字架上面.
你看不管是天主教 東正教.
他們的十字架.
你看就是他們掛一個十字架對不對.
請問這個十字架上面有誰 耶穌.
因為他們很看重這件事情.
他們常常想到的耶穌為我們受死.
很好 這非常好.
但親愛的同志們.
我們的信仰不能夠停留在這裡.
我們的信仰必須看見.
這一位耶穌基督.
這位定十字架上的耶穌.
現在基督教的基督徒掛的十字架.
耶穌在哪裡.
沒有了對不對.
這個十字架上面沒有看見耶穌 為什麼.
因為表明一件事情.
耶穌已經復活了.
這個世界.
這個佛教也好.
伊斯蘭教也好.
印度教也好.
很多的教都有他的所謂的教主.
對不對.
但是這些教主都死了.

$^{121}$所以我們信仰的核心是很重要的.
表明說我們是受死而且復活.
第二個.
為什麼復活這麼重要.
保羅在這個地方說.
以大人寫明是神的兒子.
神的兒子是什麼意思.
這個概念不能用我們中文的想法來講.
我是我爸爸的兒子.
所以我不是我爸爸對不對.
懂不懂.
懂嗎.
但是耶穌他是神的兒子.
表明他就是神.
你聽得我在講什麼.
這是不一樣的邏輯.
因為在猶太人的觀念裡面.
當講到神的兒子的時候.
就是在講到他是神.
所以神的兒子就有神性對不對.
OK 好了.
那這一個跟我們一般所提到的.
我們基督徒都是神的兒女.
不太一樣.
了解我意思嗎.
我們都是神的兒女.
但是我們有沒有神性.
沒有.
我們是不是神.
不是.
所以耶穌基督從死裡復活這件事情.
在整個世界.
從以前到現在.
請問有沒有哪一個人死去了之後.
然後又活起來了.
然後就再也沒有死了.
有沒有.
你可能聽過有人死了復活了.
去過殿堂.
但請問他後來.

$^{161}$他後來又到哪去了.
他還要死了是不是.
你有沒有看過一個人死了之後復活.
然後他就再也沒有死亡了.
耶穌基督.
所以復活這件事情就證明.
這裡說以大人寫明是神的兒子.
寫明這個字在希臘文裡面.
就是地平線的意思.
我們英文的horizontal.
horizont.
這一個地平線.
什麼叫地平線.
你看.
當你從這個圖片裡面.
你這個海這樣看過去的時候.
你看到什麼.
你有沒有看到一個.
一個平線.
那個就叫做地平線.
是不是.
所以這個寫明原文的意思.
它就是把它標示出來.
把它劃定界限.
這個意思.
換句話說.
耶穌基督他從死裡復活這件事情.
神接著這個死裡復活.
大人死了死裡復活這件事情.
要來把它標明出來.
把它顯示出來.
顯示什麼.
顯示耶穌基督與所有一切的人.
所有一切其他地上的這些教主.
不一樣的就是因為.
耶穌基督他是神.
Amen.
跟你旁邊說.
耶穌基督是神.
當我們講到耶穌基督是神的時候.

$^{201}$我要再強調.
耶穌基督不是神.
裡面的一位神.
因為你如果碰到伊斯蘭教的.
你跟他講說.
I believe in God.
他會說什麼.
I believe in two.
我也相信.
但是他的God是什麼.
Allah.
你懂我意思嗎.
所以你跟伊斯蘭教的人.
跟他談話的時候.
你不要講God.
為什麼.
你也不要跟猶太教的人講God.
為什麼.
納坦雅胡也跟你說.
I believe in God.
但是他不believe in Jesus.
我們所相信的這位耶穌.
他是God.
他是三位一體的神的一位.
懂這意思嗎.
聖父聖子聖靈.
所以這位神他是.
獨一的真神有沒有.
認識你獨一的真神.
獨一的真神.
他又是獨一的.
但是我們所信的.
又是三位一體的神.
我今天沒有時間仔細講.
需要了解的.
請你參加門訓.
那第二位神的復活.
耶穌基督的復活.
把他跟世上所有一切其他的人.
劃清界限.

$^{241}$有多少人去過這些地方.
這個釋迦摩尼他出生的地方.
去過在盧比尼.
就是尼泊爾的盧比尼.
但是據說他在那邊.
佛家叫做什麼.
在那邊原籍的地方.
那個是在印度一個北部.
可能我想我們大概都沒有去過.
那麼請問現在釋迦摩尼在哪裡.
你到他的.
你到這個地方.
他這個名字是翻譯過來很奇怪.
我們中文叫做居什麼.
居斯納杰羅.
就這個地方.
你上去Google可以找得到.
這個地方其實據說.
這個釋迦摩尼他死了之後.
人把它火化之後.
把他的舍利子.
舍利子知道吧.
把他分成了八個.
所以送往了八個地方去.
所以這八個地方.
都把它弄了一個塔.
一個廟之類的.
來供奉他.
所以你要去找他的墳墓.
你找不到他這一個單獨的.
你只找到他的舍利子懂嗎.
但是我請問.
就算是他舍利子了.
請問莫罕默德耶穌.
那個叫.
釋迦摩尼有沒有再復活.
沒有.
莫罕默德.
釋迦摩尼死了.
他沒有再復活.

$^{281}$莫罕默德呢.
莫罕默德是伊斯蘭教的創辦人.
在他的手下.
南征北討.
把伊斯蘭教整個阿拉伯半島.
通通把它統一了.
後來他的街坊人.
陸陸續續的.
把當時候的北非.
甚至歐洲一些的地方.
通通把它拿下來.
就形成了今天的伊斯蘭教世界.
所以當初他們就是用武力.
擴展了他們的宗教信仰.
但是這位偉大的莫罕默德.
這是他的.
據說是他的靈犬.
請問他現在在哪裡.
在哪裡.
他還活著嗎.
沒有.
世界上面所有偉大的教主.
都死了.
世界上面所有偉大的君王都死了.
成吉思汗死了.
打遍歐亞非山中.
是不是.
跟耶穌差不多聯名.
他死了.
當他死的時候.
據說他叫他的手下.
把在棺木裡面挖兩個洞出來.
讓他的手伸出來.
因為他說我要向世界上的人.
告訴世界上的人.
我空空的來.
我也空空的走.
我什麼都沒有帶走.
一個曾經差池風雲.
把歐亞非山中都統一的.

$^{321}$打下來的一個成吉思汗.
最後他人生走到終點的時候.
他什麼都帶不走.
你今天有多少棟房子.
你今天有多少的存款.
黃金多少噸.
多少兩.
多少盎司.
兄弟們這一切當你走的時候可以帶走嗎.
我們中國人聰明啊.
帶不走怎麼辦呢.
燒紙錢.
是不是.
把汽車把洋房通通燒.
希望將來到天.
到了這個天.
什麼極樂世界的時候.
他還可以繼續開車.
你到天堂還去開車嗎.
真是很可笑.
你有沒有發現我們.
我們自以為聰明.
替自己造神.
但是耶穌基督他從死裡復活.
當你去耶路撒冷.
到聖地去看的時候.
你看見了一個.
耶穌所復活的那個地方.
那個石洞裡面.
現在在那個地方裡面.
裡面沒有人了.
那邊寫著耶穌復活了.
哈利路亞.
耶穌復活了.
所以你不要到死人當中去找死人.
你不要在死人當中去找活人.
耶穌基督已經復活了.
所以耶穌基督的復活.
證明他是神.
他不只是基督教的創辦人.

$^{361}$是一個教主.
他就是神.
所以今天因著耶穌基督的復活.
我們是有盼望的人.
所以你要來相信這位耶穌基督.
跟你旁邊說你要相信他.
不要像重金一樣等待.
等待了十幾年了.
是不是.
我們中國人講說.
晚上睡覺脫了鞋.
第二天起來能不能穿鞋.
你都不知道.
是不是.
黃泉路上有老有少.
你到歐洲去.
你看到有一些的教堂裡面.
有那個小棺木.
有沒有.
中國人覺得很忌諱.
怎麼到教堂去看到這個棺木.
西方人他們受基督教的影響.
他們不覺得沒問題的.
反正他們覺得從教堂裡面.
直接到天上去比較快一點.
請問你有沒有看到.
那個小小的孩子的棺木.
那些可能是君王的孩子.
或者是什麼皇帝的.
這些的親戚.
小孩子.
是不是我們的年齡.
就代表我們的人生多長呢.
人生不是你能夠掌握的.
趁著今天抓住耶穌基督.
趁著有機會來認識這位耶穌基督.
約翰福音一章十二節.
這裡告訴我們什麼叫做信.
我們起來念好不好.
凡接待他的就是明的人.

$^{401}$他就賜他們全品做神的兒女.
如果你不會背的話.
你要把這節背起來.
很多人基督徒信耶穌信了很久.
你問他說你怎麼知道你是基督徒.
你怎麼確認你得救了.
不知道他怎麼回答.
把這些神經背起來.
這節神經是你得救的一個把握確定.
這裡說凡接待他的就是.
那個瓜胡是我多加的.
就是信他明的人.
接待是什麼.
接待就是接受.
接待他就是耶穌.
接待他到哪裡呢.
接待到我們的心裡面.
他說你凡是把耶穌接到你心裡面來.
你就是信他明的人.
他就賜給他們全品做神的兒女.
他就給你全品做神的兒女.
你就是神的兒女.
但是請你記得.
這個原文裡面其實它這個相信.
它是一個英文把它翻譯.
它的介系是叫做believe in two.
什麼叫做believe in two.
這個意思就是說你要信進去了.
我們中國人講說你信得太迷了有沒有.
我以前爸爸媽媽常常.
他們還沒有信主.
常常就說你信迷了.
你有沒有被人家這樣講過.
有沒有被人家這樣講過.
沒有啊.
那你信得還不夠.
人家說我信迷了.
我不是信迷了.
我是信路了.
我們信這位耶穌不是受洗的那一天.

$^{441}$那叫做信耶穌.
信耶穌是一個ungoing是一個持續的.
從我受洗的那一天開始.
我就要一直的信靠他持續的信靠他.
這個信心要越來越加增越來越加增的.
這個是這裡的信的意思.
而且他原文裡面是一個持續進行式.
表明就是一個持續進行式.
不是一個一次的one time thing.
不是一次的一次受洗.
好了從此等著上天堂.
不這種信仰是一個非常廉價的信仰.
我們所信的這位耶穌.
不是只是受洗的那一天的耶穌.
他是我們一生要相信並且要跟隨.
並且要追事奉他的這位耶穌.
你要有這樣子的信心.
你才能夠經歷到他這裡所說的.
保羅說在信的人身上所顯的能力.
是什麼 二等好大.
丘子梅你又經歷過神的能力嗎.
我不知道你有沒有經歷過.
神的能力你又經歷多少.
你在預備這信息我回想我自己的一生.
我怎麼會知道.
我怎麼會有把握確信這位神呢.
我想到神怎麼樣醫治我多年的疾病.
所以這位神他的人力是有大人的.
他能夠醫治人的疾病.
我想到我被聖靈充滿的時候.
我親身的經歷的感受到.
那種overwhelming的一種的喜樂.
那不是我天生的抑鬱的性格能夠有的.
那是一個supernatural.
從上面澆灌下來.
澆灌到我的裡面的那一種的喜樂.
我知道那就是神.
那就是神的大人.
使得剛才像我們完美姐妹一講的.
我們本來是活在懼怕恐懼憂慮的這一個世界.

$^{481}$我們本來是好像覺得不知道人生方向的.
但是因著信了耶穌之後.
耶穌基督那個復活的大人.
進入我們的生命裡面.
而且你要不斷的去經歷他.
你要不斷的相信他.
他就會在你的身上.
顯出這個能力是何等的浩大.
Amen.
我很喜歡中文這個何等.
何等就表示什麼.
哎呀形容不出來.
形容不出來的浩大.
Amen.
我們的神他是這麼樣子偉大的神.
他有這樣子浩大的能力.
要彰顯在我們這些信的人身上.
有一次在.
這個有一個著名的指揮家.
叫吳定一.
我們第一弟兄也是著名的指揮家.
沒有話講.
這個上海音樂學院.
但我現在講的是有一個.
在二十世紀一個很著名的瑞士.
瑞士的一個指揮家.
也是作曲家.
那麼他有一次.
也是在復活節的時候.
他帶領師班還有交響樂團.
他們在彩排韓德爾.
你們知道韓德爾.
他所作的那個彌賽亞這一曲這個聖曲.
那在彌賽亞.
其中在這個第三部分的時候.
那麼講到耶穌基督的復活.
講到這些的時候.
那麼他有一個女高音的獨唱副歌.
叫做我知我的救贖主活著.
這時候在彩排的時候.

$^{521}$這個女高音就在那邊唱.
那個真的是聲音非常的美妙.
呼吸順暢.
發音清晰.
真是無懈可擊.
他唱完之後.
大家所有的這些交響樂團員.
還有這些班員都在等待.
指揮會給他大大的讚賞.
沒有想到這個指揮.
看著這個女高音.
然後他就跟她說.
他說女兒.
你相信你的救贖主活著嗎.
大家本來期待這個指揮要給他.
大大的獎賞一番對不對.
沒想到他竟然問了這一句.
而且帶著有點憂傷的表情在問她.
這個女高音一聽了之後.
就有點很尷尬就說.
他說I think I know.
我想我是知道的.
然後這個指揮就跟她講.
他說如果你真的知道你的救贖主活著.
好.
那我們再來練習.
所以他就帶她們再一次的練習.
他說你要把你自己所認識的.
這位耶穌基督唱出來.
就帶她們練習.
再一次的在練習的時候.
後來又輪到這個女高音在那邊唱的時候.
他唱到我知我的救贖主活著.
I know.
I know that my redeemer lives.
他唱到這裡的時候.
他用他自己從他生命裡面.
所經歷的那種感受.
唱出了一個.
整個大家聽了之後都非常的感動.

$^{561}$這個姐妹自己唱的時候.
也是眼淚盈眶的流下來.
她用她的生命.
來唱出了這一個救贖活著.
Amen.
最後這個指揮.
也眼淚盈眶的對她說.
她說現在我知道.
你真的是知道.
你的救贖活著了.
童子妹不管你信主多少年.
今天我要挑戰你.
你知道你的救贖活著嗎.
你相信這位救贖活著嗎.
當你相信知道這位救贖活著的時候.
在你面臨人生當中不管什麼難處.
夫妻的問題.
婚姻的問題.
兒女的問題.
或者是各式各樣的人生的挑戰.
你面對到人生當中.
這一切的挑戰和苦難的時候.
你有把握.
你有信心.
因為在新的人身上所顯的能力.
是何等的浩大.
而這浩大的能力.
要帶領我們乘風破浪.
這個浩大的能力.
要幫助我們能夠勝過一切萬難.
Amen.
\newpage



\section{彼得前書 1:3-9}
\label{sec:031929}
\textbf{基督復活的大能}~~~~李思聰牧師
\newline
\newline
連結: \href{http://www.church.com.hk/acms/content.asp?site=cdc\&op=show\&type=product\&code=031929\&layout=sermon}{\texttt{031929}} ~~~~ 語音日期: 2026-04-05
\newline
\newline
\hyperref[sec:031928]{\small{< < < PREV SERMON < < <}}
~
\hyperref[sec:index_chronic]{\small{[返順時目]}}
~
\hyperref[sec:031930]{\small{> > > NEXT SERMON > > >}}
\newline
\newline
彼得前書 1:3-9
\newline
\begin{longtable}{cl}
\hline
\hline
章節 & 經文 (和合本修訂版)\\
\hline
1:3 & \begin{tabularx}{0.7\textwidth}{X} 願頌讚歸於我們主耶穌基督的父神！他曾照自己的大憐憫，藉著耶穌基督從死人中復活，重生了我們，使我們有活的盼望， \end{tabularx} \\ \\ \relax
1:4 & \begin{tabularx}{0.7\textwidth}{X} 好得到不朽壞、不玷污、不衰殘、為你們存留在天上的基業， \end{tabularx} \\ \\ \relax
1:5 & \begin{tabularx}{0.7\textwidth}{X} 就是為你們這些藉著信、蒙神大能保守的人，能獲得他所預備、到末世要顯現的救恩。 \end{tabularx} \\ \\ \relax
1:6 & \begin{tabularx}{0.7\textwidth}{X} 雖然你們必須在百般試煉中暫時憂愁，你們要為此喜樂 ， \end{tabularx} \\ \\ \relax
1:7 & \begin{tabularx}{0.7\textwidth}{X} 使你們的信心既被考驗，就比那被火試煉仍然能壞的金子更顯寶貴，可以在耶穌基督顯現的時候得著稱讚、榮耀、尊貴。 \end{tabularx} \\ \\ \relax
1:8 & \begin{tabularx}{0.7\textwidth}{X} 雖然你們沒有見過他，卻是愛他；如今雖看不見，你們卻因信他而有說不出來、滿有榮光的喜樂， \end{tabularx} \\ \\ \relax
1:9 & \begin{tabularx}{0.7\textwidth}{X} 因為你們得到信心的效果，就是靈魂的得救。 \end{tabularx} \\ \\
[1ex]
\hline
\hline
\end{longtable}
$^{1}$最近發生這幾個禮拜.
忽然之間.
沒有計劃,沒有想到.
但是有時死亡的陰影就會籠罩我們.
我們個個都知道.
沒有一個人能夠可以.
預計明天,是吧?.
但是正如今天使母帶領我們所唱的詩歌.
「恩主活著」.
我們就知道是誰掌管.
我們的未來.
更重要的不單只是明天.
是關乎到我們永生,永死.
究竟我們離開這個世界會去哪裡.
這些問題.
今天我們也再一次透過.
復活節的主人.
我們能夠可以彼此提醒.
這位.
勝過死亡,復活的救主.
對我們人生是多麼重要.
我們再一次做一個祈禱.
我們就開始今天的訊息.
因為你掌管我們的生命.
你更加掌管我們的明天.
掌管我們前面人生.
甚至去到永恆的道路.
主,但是你給我們選擇.
讓我們能夠可以選擇跟從你.
就好像今天.
我們有兩位弟兄姊妹接受洗禮一樣.
他們選擇接受你,跟從你,委身於你.
這真是我們每一位基督徒.
最寶貴,最有價值人生的抉擇.
不是因為我們選擇.
是因為你揀選我們.
是因為你的恩典和憐憫.
臨到我們身上.
好讓我們今天不拒絕你的恩典.
願意接受你給我們跟從.

$^{41}$願意接受你給我們的恩典和簡說.
讓我們能夠可以藉著基督復活的大能的應許.
能夠過一過新的人生.
充滿盼望的人生.
反映你的能力在我們身上的人生.
主,多謝你聽我們禱告.
奉耶穌基督的名.
下門.
今天我們因為復活節.
所以我們看一看另外一段經文.
就不在舊約了.
這是彼得前書一章三到九節的經文.
彼得前書是誰寫的?.
使徒彼得寫的.
如果我們說.
基督教的信仰最重要的元素是什麼呢?.
我相信我們每一位.
或者我們是基督徒.
或者甚至不是基督徒都聽過.
其實這個就是一個很重要的信息.
耶穌基督是一位復活的主.
今天我們所信的不是一個宗教的教主.
我們所信的是一位.
聖過死亡復活的真神.
這位的救主.
聖經這樣說.
它說:若基督沒有復活.
我們的信便是枉然.
你們仍在罪裡.
這個很清楚.
在彼得前書提醒.
我們今天的信仰.
我們如果信一位沒有復活的基督.
我們的信便是好像其他宗教一樣.
所有其他宗教都是一樣.
他們可以信的是人.
或者是人製造出來的一個理念.
或者一個假的神.
但是耶穌基督復活.
我們的信便不是枉然.

$^{81}$如果耶穌基督沒有復活.
我們的罪仍然沒有辦法去清除.
我們仍然與神隔絕.
這就是何等重要.
不知道大家有沒有看過.
一部電影叫做.
Passion of the Christ.
好多年前了.
誰看過?.
這麼少人看過.
這部由Mel Gibson所導演的.
一位名導演.
基督受難.
沒有看過?.
真的?.
誰看過?現在沒有印象.
誰看過?.
差不多20年前了.
Passion of the Christ.
如果大家還記得.
它最後一個鏡頭.
Passion of the Christ是指耶穌受難.
所以耶穌受死.
那應該就完結了.
但是這部電影最後一個鏡頭.
不是停在耶穌受死的鏡頭.
它的鏡頭是停在一個空的墳墓裡.
證明耶穌是一個復活的主.
有人問Mel Gibson.
為什麼這個製作要停在這個鏡頭.
他說如果沒有的話.
故事就未能夠完整.
所以它開始拍第二部.
大家有沒有聽過?.
它現在拍的一部叫做.
The Resurrection of the Christ.
即是基督復活.
其實已經在去年八月開始拍攝.
預計會在明年2027年上映.
總共有兩部份的上下集.

$^{121}$The Resurrection of the Christ.
預計會在2027年3月上映上集.
然後2027年5月上映下集.
所以期待到明年的時候.
我們可以觀賞這部電影.
所以在他心目中很清楚.
如果只是耶穌受難是未完整的.
耶穌復活才是故事的結束.
在整本新約的聖經裡.
關乎基督耶穌復活的事實.
或者說關乎復活的真理.
總共有人統計過有150次之多.
單單新約27本的書信裡.
提到耶穌復活的事實.
或者關乎復活的真理.
總共有150次之多.
可見這是基督教信仰的核心.
非常重要.
所以今天我們和彼得.
就是他其中一位的十二門徒之一.
後來所寫的一段回想.
和聖靈啟示他成為神的話語的一段真理.
讓我們思想.
我們看看彼得前書一章.
第三到第九節這段經文.
我們同聲讀出來好不好.
請大家一起讀.
願眾贊歸於我們主.
耶穌基督的父上帝.
他曾照自己的大憐憫.
藉耶穌基督從死裡復活.
重生了我們.
叫我們有活潑的盼望.
可以得著不能扭壞.
不能玷污.
不能衰殘.
為你們存留在天上的基業.
你們者因信蒙上帝能力保守的人.
必能得著所預備到末世要顯現的救恩.
因此你們社大有喜樂.

$^{161}$但如今在百般的試煉中.
暫時憂愁.
叫你們的信心既被試驗.
就比那被火試驗仍然能壞的金子.
更顯寶貴.
可以在耶穌基督顯現的時候.
得著稱讚榮耀尊貴.
留意第五節.
你們者因信蒙上帝能力保守的人.
今天的題目就是基督復活的大能.
這大能不單是耶穌透過復活獻出來.
也是因為祂的復活賜給所有蒙具保守的人.
就像剛才做的見證一樣.
上帝能力的保守.
也要顯現出來.
最後兩節.
第八節請讀.
你們雖然沒有見過他.
卻是愛他.
如今雖不得看見.
卻因信他.
就有說不出來.
滿有榮光的大喜樂.
並且得著你們信心的果效.
就是靈魂的救恩.
基督復活是一個事實.
是一個禮物.
是一個神給我們救恩的禮物.
問題是我們願不願意用信心去接受它.
我們看到.
第一在這裡.
第九節.
他說你們就得著信心的果效.
得著信心的果效是什麼.
就是靈魂的救恩.
所以第一方面.
我們看到基督復活為什麼這麼重要.
祂的能力顯現出來.
在哪裡.
因為祂帶給我們.

$^{201}$帶來一個靈魂的救恩.
其實一開始.
剛才我們所讀的第三節.
已經說出來了.
這個靈魂的救恩是什麼.
就是藉著耶穌基督的復活.
重生了我們.
叫我們有一個新的生命.
然後第五節.
他又再一次說.
他說你們這因信蒙神能力保守的人.
必要得著所預備.
到末世要顯現的救恩.
然後最後剛才我們所讀的.
就是指到這個靈魂的救恩.
耶穌基督從死裡復活.
是因為祂顯明.
祂可以賜一個新的生命給我們.
這個新的生命就是重生的生命.
當耶穌復活的時候.
不只是祂的肉身.
回到祂還沒死之前.
而是祂的肉身是一個改變了的身體.
祂的肉身改變了的身體.
是跟舊的肉身的生命是不同的.
對於我們來說是更加清楚的.
當我們接受主耶穌基督的時候.
這個就是一個靈魂的救恩.
祂是在我們的生命當中.
起了一個新的創造.
所以來過帝母這位律法師問耶穌.
我們怎樣可以見到神的國.
耶穌很清楚地告訴他.
我們一定要重生.
就是一個新的開始.
怎樣能夠進入神的國.
必定需要一個好像耶穌一樣.
一個復活的新的生命.
才能夠進入耶穌永恆的國.
所以我們看到.

$^{241}$這個就是指我們靈魂的救恩.
影響我們將來會得到一個新的復活的身體.
如果耶穌沒有從死裡復活.
我們就沒有犯法得到這個.
靈魂得救的新生命.
正如聖經很清楚地說.
血肉之體不能承受上帝的國.
必扭壞的不能承受不扭壞的.
而在約翰福音三章第六節.
就是關乎尼哥帝母和耶穌的交談的時候.
耶穌說從肉身生的就是肉身.
從靈生的就是靈.
所以這裡指的是靈魂的救恩.
這是一個奇妙的救恩.
是出於神的憐憫.
出於神的預備.
藉著耶穌基督的復活.
讓我們能夠得到這個.
重生的靈魂的救恩.
我們所有弟兄姊妹.
當我們經歷過主.
在我們身上生命的改變.
重生這個奇妙的創作.
當我們相信主的時候.
感覺到一種新的生命的開始.
我們一定會經歷過.
就像剛才師母帶我們唱的時候.
說耶穌是活著的.
因為我們經歷過這位活著的主.
我們過往都不少的迷信偶像.
因為我們中國人有很多民間宗教信仰.
經歷過偶像崇拜的迷惘.
偶像崇拜背後的邪靈的真實的影響.
以及攻擊和誘惑.
但是因著信靠耶穌.
生命得到改變.
就能夠脫離偶像.
信靠真神.
重生一個光明的生命.
有很多弟兄姊妹.

$^{281}$不少的.
我們經歷過這種生命的改變.
這種靈魂的救恩.
我們剛剛開完幸福小組.
有好幾位弟兄姊妹講他們的見證.
講他們從以前沒認識主之後.
到認識主生命的改變.
有些講到他們生活放縱敗壞.
是非常的迷惘.
但是經歷過上帝的救恩.
拯救了他和他整家人.
他感覺到神奇妙的作為.
讓神從一個敗壞的生命人生當中扭轉.
以至他們能夠今天坐在這裡.
敬拜上帝.
太太 兒女 媽媽.
一同的來敬拜主.
生活改變.
人變成行為端正.
家庭蒙恩.
有一個事奉的人生.
這個實在不是人的作為.
是上帝重生.
叫我們經歷靈魂救恩的實在.
有些弟兄姊妹分享.
他們過往人生怎樣的憂慮.
不安 靠自己.
剛剛我們聽完這些見證.
但是我們看見.
上帝扭轉我們的生命的時候.
我們懂得怎樣依靠神.
有一個交托的人生.
活出真正的喜樂.
這個都是生命重生.
靈魂救恩的改變.
有些人我們也經歷過.
人本來很自卑.
覺得自己的教育水準很低.
自己的背景很差.
覺得自己沒有什麼用.

$^{321}$但是很奇妙.
上帝救恩令到他.
明白到上帝愛他.
接納他.
無條件地來施恩給他.
他生命覺得很寶貴.
知道我在神的眼中很重要.
所以他的人生開始改變.
他的生命變得有價值.
他不單止不自卑.
還幫助其他人認識耶穌.
很大膽地來分享福音.
很大膽地來告訴別人.
我是多麼的蒙印.
這些都是生命的改變.
是靈魂的救恩.
重生的後果.
有些人以前說謊.
當作吃生菜.
說得出就說.
沒有經過大腦.
騙得了別人就騙.
騙得過就過關.
但是信了耶穌之後.
有聖靈感受.
他的責備提醒.
說話認真很多.
說話謹慎很多.
說錯了.
其實差一點點而已.
都馬上澄清.
其實我不是這個意思.
不要以為別人當他是說謊.
為什麼會有這樣的生命的改變.
因為這種是重生的生命.
是靈魂的救恩所達至的.
不代表我們沒有問題.
正如剛才我們的姊妹.
在世代的時候分享.
不代表我們的生命.

$^{361}$已經完整了.
再沒有瑕疵.
再不會犯罪.
但是得著這個靈魂的救恩.
就代表有一種新的生命.
與我們這個舊的罪惡的生命對衡.
讓我們能夠在很多時候裡面.
當我們去明白真理.
禱告尋求神的時候.
聖靈特別的.
來幫助我們.
對於我們生命有一種提醒.
對於我們憂愁的時候.
能夠得到屬天的喜樂和平安.
有一種聖潔的渴慕.
希望自己的生命更加的聖潔.
更加離開罪惡.
能夠渴慕有神的平安.
和同神和好的關係.
來過我們的生活.
這種就代表一個重生的生命.
靈魂的救恩.
有真理的原則.
有內心的愛.
去愛神.
還有去學習去愛別人.
昨天我和師母有時間抽空.
去看.
因為這個禮拜裡面.
有一套電影在電影院上演的.
我也給了一些弟兄姊妹知道.
昨天我去電影院看的時候.
碰到我們教會的弟兄姊妹.
這套電影叫做.
Big Awakening.
是一個大的復興.
是講到我們美國.
美國今年建國250年.
最初建國的時候1776年.
後來到第十年.

$^{401}$開始有一種很大的分裂.
還有因為英國一直以來.
那種假信仰.
一種表面的信仰.
是基督教國家.
但是其實沒有真正重生的生命.
所以就影響了當時的英國社會教會.
也都影響了當時.
剛剛開始建國的這13個新的.
美國最初型的一個國家.
當時他們的信仰.
只不過是一種表面的信仰.
說是基督徒.
但是真沒有基督徒的生命.
根本沒有重生得救的.
這種生命的表現.
其實這種情況.
今天在我們現代的社會.
甚至美國說是基督教很多.
基督徒很多.
說用基督教的原則立國的國家.
都是屁屁皆是.
不單止社會.
甚至有些教會都是.
自己說自己是基督徒.
剛才我們的弟兄講的見證就很清楚.
但是卻沒有真實重生的生命.
這部電影就是講到當時.
這位被神呼召非常大用的.
George Whitfield.
大家有沒有聽過他的名字.
就是懷德肥.
喬治·懷德肥.
當時他被神揀選.
成為一個復興家.
來到全港重生的真理.
他就是自己.
因為在監獄裡跟著一群基督徒.
就是約翰·韋斯理和查理·韋斯理.
自己讀聖經.

$^{441}$大聲讀出來.
一邊讀就是約翰·方第三章.
一邊讀這個重生的真理.
他就開始明白到重生的真義.
然後他的生命改變.
然後他就一生傳揚.
全港這個重生的重要性.
這個重生的真理.
叫很多很多那些掛名基督徒.
很多很多當時那些.
是回教會.
或者還沒有回教會.
是奴隸的.
是社會低下階層的人.
因著這種聖經福音的信息.
福音 耶穌基督復活.
賜給人重生的生命.
靈魂的救恩的信息.
有一個很大的復興.
讓今天美國還能夠.
可以繼續經過這二百多年.
沒有在最初那十年就散了.
或者分裂了.
是一個很重要的關鍵.
所以當時他有一位好朋友.
今晚我們都認識的.
Benjamin Franklin.
在銀紙上面還有他的樣子.
就是因為這樣.
和他一起在印刷上面.
傳揚了這些信息出去.
叫人去聽這位.
是大的復興家.
George Westfield他的信息.
以致引來很大的影響力.
而且也因為這樣影響了他.
當時他就叫國會.
當時準備分裂的時候.
要回到神的面前.
來禱告這個重要性.

$^{481}$所以如果大家有時間.
還有幾天.
我建議大家可以看看這部電影.
很便宜.
現在在電影院有放.
所以我們想到這個重要性就是.
原來耶穌基督賜給我們靈魂.
為什麼祂要復活呢.
因為祂賜給我們一種真的生命.
這個真的生命就是靈魂復生.
靈魂重生得著救恩的生命.
不是一個外表的信條.
不是一個宗教形式.
不是一個宗教的信仰.
是一種新的生命.
是一種救恩.
靈魂的救恩.
賜給我們重生的生命.
所以祂從死裡面復活.
證明是一種新生命的開始.
大使們.
今天這種耶穌基督所帶給我們的能力.
一個大能的果效.
就是叫我們從死裡面復活的新的生命.
究竟我們是否每一個都有這個把握.
都有這個經驗.
都有這個清楚靈魂救恩重生的一個確據.
第二我們在這裡看到.
就是帶來復活的大能.
帶給我們復活的盼望.
就是這種活潑的盼望.
在第三節和第四節.
我們再讀一次好不好.
請大家同聲讀出這兩節經文.
願頌讚歸於我們主耶穌基督的父上帝.
他曾照自己的大憐憫.
藉耶穌基督從死裡復活.
重生了我們.
叫我們有活潑的盼望.
可以得著.

$^{521}$不能扭壞.
不能玷污.
不能衰殘.
為你們存留在天上這個永恆的基業.
這種生命是延續去直到永恆.
這種是永恆不扭壞的生命.
他經歷了這種復活.
永恆不扭壞的生命.
他經歷了這種永恆不扭壞的生命.
他叫我們經過死亡.
肉身的死亡.
好像耶穌一樣.
經歷過肉身的死亡.
但是和耶穌一樣.
我們能夠經過死亡之後.
有身體的復活.
這個就是幫助我們明白到.
兩個不同的終極的結果.
一個結果就是永遠的滅亡.
一個結果就是永生.
帶來給我們延續到永恆的一個應許.
這個就是我們今天藉著基督所得復活的新生命.
我們再說這本書是誰寫的.
是彼得寫的.
從前的彼得.
他未經歷過耶穌的復活.
所以當死亡的陰影籠罩的時候.
他非常害怕.
他與主同行三年那麼多.
朝夕相處.
彼得的性格是很剛強的.
是很勇敢的.
可以說有時候是拍胸口就上去的.
但是他面對死亡的威脅的時候.
死亡的危險的時候.
他拜在一個小小的侍女.
講一講幾句話的手上.
因為他懼怕要和耶穌一齊.
來面對死亡.
當人問他這個是不是你的夫子.

$^{561}$或者你是否跟過他的門徒.
這個侍女.
如果以前的彼得.
怕胸口.
什麼都夠膽.
其實他之前也是這樣.
死都不怕.
不過真的面對的時候.
卻是非常膽怯.
當主耶穌死了的時候.
被釘死的時候.
所有門徒都分散.
包括彼得.
所以我們見到當時所有門徒都懼怕.
因為死亡是真實的.
勝過死亡是他們沒有經驗過的.
是沒有體會過的.
所以無論你性格多麼剛強.
無論你回了教會.
或者聽過基督教的信仰多少年.
當我們面對死亡那一刻.
我們都會束手無策.
今天法蘭西斯大兄講的.
就是一個真實的見證.
當我們面對死亡的時候.
我們真的束手無策.
不要說他有五分鐘.
或者十五分鐘.
有些連一分鐘都沒有機會.
所以當我們今天面對死亡的事實.
我們是沒有想.
因為想逃避.
但是如果我們清楚去想一想.
當那一天來到的時候.
我們都會懼怕.
因為如果你不知道自己去哪裡的時候.
死亡就吞噬我們.
在一位很出名的人名.
叫做Harvard Hughes.
我們有些弟兄姊妹.

$^{601}$有沒有誰以前在Hills做過.
有吧.
他是一位美國的富豪.
他是飛行員.
他是航空的工程師.
同時也是電影製片人和投資者.
他創立並且經營自己的電影公司.
但是對航空事業產生非常深遠.
非常濃厚的興趣和深遠的影響.
因為他自己有電影製片廠.
又有自己的飛機公司.
所以他是名成利就的.
不過一位如此出色的人物.
一位如此名成利就的富商.
以及如此有創造能力.
以及如此出色的工程師.
卻是很怕死.
因為這個不是錢買得到.
不是你的地位可以賺來.
更加不可以用你的方法去逃避.
所以他就想了一個方法.
在他年老的時候.
他為自己安置了一個無菌的.
封閉的空間在那裡住.
無菌的.
即是進去那裡消毒.
所以沒有細菌感染的.
那就進去.
凡要進去幫他做事.
或者去探訪他的人.
一定要戴手套和消毒.
才可以進去那裡.
所以他是很怕病.
很怕死.
在他年老的時候.
他是很明顯非常脆弱.
心靈很空虛的一個階段.
但是當死亡要找到他的時候.
他根本沒有辦法知道.
當死亡敲門的時候.

$^{641}$也不會通知.
所以最後他只是這樣靜悄悄地.
被死亡帶走.
今天我們知道.
死亡背後就是魔鬼的勢力.
將來他要和所有還沒得救的人.
都進入永恆的火壺裡.
這就是最終的結局.
所以如果這樣想一想.
死亡不只是死那一刻可怕.
是死後更加可怕.
但是彼得不同的地方.
就是他後來認識了.
主耶穌復活的大能.
當他第一個門徒和約翰.
去到這個空墳墓的時候.
他的心靈就開始有新的認識.
後來當主顯現.
復活的主在他顯現的時候.
他就越來越體會這種復活的大能.
之後到基督升天.
他就成為一個最重要的人物.
去宣講證明耶穌是復活的主.
所以在彼得前書這段經文裡.
他講出這種復活的盼望.
永恆的盼望.
極大的保障.
不朽的基業為他存留.
這就成為我們每一個基督徒.
活潑盼望最終的去向.
Amen.
這個人大家認不認識.
就是邱吉爾爵士.
大家知道邱吉爾是20世紀.
世界聞名的英國政治家和文學家.
他曾經兩次擔任過英國的首相.
在英國受到德國攻擊的時候.
帶領英國人.
過渡第二次世界大戰的危機.
一個非常重要的人物.

$^{681}$同時在1953年得到諾貝爾的文學獎.
所以他是一個出色的政治家和文學家.
這位才氣滿溢的偉人.
當他離開世界的時候.
因為他是一個基督徒.
所以他親手安排自己的追思禮拜.
要唱的詩歌.
以及整個安息禮拜的程序.
他選擇了他的喪禮在聖保羅大教堂舉行.
因為他知道有成千上萬的人.
來追思他的生平.
所以在當時報道的時候.
在1965年.
英國首相的國葬的葬禮.
有數十萬的民眾.
數十萬民眾.
來到道別.
他既然知道有這麼多人來.
他又預備了這個程序.
所以到了最後的安息禮拜結束.
牧師最後祝福之後.
就有一個特別的安排.
在聖保羅教堂的一邊.
有一位喇叭手.
吹出一首號角的一首樂曲.
叫做「熄燈號」.
就是熄燈了.
所以整個會場就開始熄燈了.
就知道這個安息禮拜結束了.
聖保羅教堂的一邊.
有數十萬的民眾.
來到道別.
他既然知道有這麼多人來.
他又預備了這個程序.
牧師最後祝福之後.
就有一個特別的安排.
在聖保羅教堂的一邊.
有一位喇叭手.
吹出一首號角的一首樂曲.
叫做「熄床號」.

$^{721}$就是熄燈號.
當時那些群眾參加的時候都很暗笑.
而且非常覺得是為之一震.
因為邱吉爾一生都是非常幽默的人.
所以連他的追思禮拜.
自己設計的時候.
他都要宣佈用一個很幽默的方式.
來講出一個重要的訊息.
一個大好的訊息.
就是人不是熄燈.
離開歷史就完了.
乃是起床.
進入新的天地.
活潑的盼望.
為什麼我們可以這麼肯定.
為什麼邱吉爾能夠知道自己.
再一次起床.
原因是因為.
耶穌基督已經復活了.
第三我們看到.
不但帶給我們靈魂的救恩.
活潑的盼望.
最後在這段經文也告訴我們.
帶給我們超越的喜樂.
超越的喜樂.
就是本來人生都會有憂傷.
有難過.
但是可以超越這些一切.
第六節到第八節.
我們再來讀一次.
請讀.
因此你們是大有喜樂.
但如今在百般的試煉中.
暫時憂愁.
叫你們的信心既被試驗.
就比那被火試煉.
仍然能壞的金子.
更顯寶貴.
可以在耶穌基督顯現的時候.
得著稱讚.

$^{761}$榮耀尊貴.
你們雖然沒有見過他.
卻是愛他.
如今雖不得看見.
卻因信.
他就有說不出來.
滿有榮光的大喜樂.
最好明顯這兩節.
第六節第八節經文.
開頭和結束.
我們看到彼得所傳的信息.
就是這種超越的喜樂.
這個喜樂不代表我們沒有困難.
他說在百般的試煉中.
暫時憂愁.
而且這個試驗有時.
是不容易的.
好像火燒金一樣.
但是目的是一樣的.
就是要令這些金更加的純.
以致將來能夠得著.
那種榮耀和美善.
在神為我們將來.
預備永恆的時候.
見他面的時候.
要得到他的稱讚.
得著他給我們的尊貴.
今天我們見不到他.
何來這種能力呢.
就是因為他復活給我們的證據.
耶穌基督已經復活.
是一種證據.
是一種發生的事實.
不是單單一種信念.
或者一種哲學思想.
是一種證據.
所以我們因著這個證據.
我們可以將信心.
放上去的時候.
不是迷信.

$^{801}$是一種有根有據的.
真實的信心.
所以他說.
我們因信就說不出來.
滿有榮光的大喜樂.
當時當耶穌被釘十架.
被埋在墳墓的時候.
門徒是憂愁的.
預先耶穌也說了.
他說你們不要憂愁.
知道他們會憂愁.
他說因為你當信我.
就是因為信.
所以後來的門徒.
結果看見復活的.
明證這個事實的時候.
他們就產生了.
很大的信心.
但是耶穌向那個多瑪.
懷疑的多瑪說.
如果那些人.
未見的時候就信的.
這個是更加有福的.
所以我們看到今天.
我們知道.
距離耶穌基督復活的時候.
已經有二千多年了.
但是這個事實是不可以推翻.
這個歷史是被記載下來的.
是可以考究的.
是可以考古的.
是可以考證的.
所以因著這個事實.
我們雖然看不到.
耶穌基督復活的身體.
但是我們這個信心.
幫助我們看見.
耶穌復活的事實.
所以如果一個宗教.
它不能勝過死亡之後.

$^{841}$所得的應許.
那這個只不過是.
一種哲學的信念和思想.
但是能夠勝過死亡之後的事實.
有了確據.
這個就是一個.
真實可信的一種信仰.
正如Rick Warren牧師所說.
今天Francis說.
有幾多個百分比.
失去了那條血管.
一百個百分比.
一百個百分比就是死了.
不過幸好他只失去了一條.
還有幾條撐住.
但是人會死.
這件事實Rick Warren說.
是一百個百分比.
無一倖免.
就是說我們都要面對.
所以有很多宗教的思想.
很多哲學的思想.
幫助我們去解決.
死亡的問題和迷思.
甚至死亡之後的一些出路.
我們很多中國人.
都是信佛教.
當然佛教有分很多種.
大乘佛教.
小乘佛教.
而在佛教裡面.
其實他們信輪迴.
大家可能都聽過.
其實輪迴不是一次.
是六道輪迴.
六道輪迴.
包括了天道.
修羅道.
人道.
惡鬼道.

$^{881}$禽獸道.
最後一個叫地獄道.
只可以在人道的時候修煉.
其他就不能修煉.
所以修煉得好.
就有機會去修羅道.
或者天道.
修煉得不好.
就會去惡鬼道.
禽獸道.
最後是什麼.
地獄道.
但是這個修煉.
就是浮浮沉沉.
有時候行 有時候不行.
所以有時候要下去.
有時候要上去.
其實可以說不是輪迴.
可以叫做浮沉.
但是人就是落在這種浮沉浮沉.
其實是沒有出路的困境.
其實也不合邏輯的困境.
最後如果小乘佛教.
他們是進入劣盤.
就是說所有都幻滅一切.
所以是沒有解決死亡之後.
基督徒不是阿Q精神.
我們不是說上了耶穌.
無災無難無憂無慮.
耶穌也這樣說.
祂說你們在今世有苦難.
也明白門徒會有憂慮.
他自己面對十字架的苦難的時候.
經過苦難.
同樣經歷憂慮.
憂傷至死.
在赫西瑪利園的時候.
是極度的憂傷.
但是不同的地方就是.
基督教給我們這個復活的應許.

$^{921}$耶穌勝過了這個苦難.
勝過了死亡.
所以能夠給我們跳躍了.
這些人世間的苦難和憂傷.
用信心明白到有一個永恆的盼望.
永生的祝福.
以致我們可以忽然間.
有一種很大的喜樂產生.
雖然死亡是懼怕的.
執掟是痛苦的.
憂傷是不可以逃避的.
但是永生的割據.
耶穌復活的真實.
就令我們超越這些.
得到一種平安,一種喜樂.
一種可靠.
耶穌基督復活.
讓我們的心靈能夠親近祂.
因為祂的愛.
我們有說不出來的大喜樂.
就是因為這種愛.
讓我們知道一種安全感.
祂不只是一位高高在上的神.
祂是為我們受苦.
來到這個世界.
與我們同經苦難.
經過死亡的主.
所以當我們經過苦難.
經過死亡的時候.
就好像跟主一起.
經過苦難 經過死亡一樣.
因著祂的愛.
我們就可以有平安.
因著這個信心.
讓我們看到.
我們靈魂救恩的果效的時候.
我們有一種超越的喜樂.
知道我們生命在神手中.
可以很安穩.
在1986年5月.

$^{961}$曾經在中國陝西省.
一個叫做扶風鎮.
離西安118里的.
一個法門寺的寺裡面.
有一個很大的發言.
一個考古的發言.
當時歷史學家和佛教徒.
整個世界的佛教徒.
都有一個非常.
鼓舞人心的消息.
因為發現了.
原來我們所崇拜.
所信仰的釋迦牟尼.
是真實的.
為什麼呢.
因為在那個法門裡面.
一個暗暗的地牢裡面.
裡面有三個箱.
一個解圓 一個解圓.
最後那個打開.
發現了一個.
法門的完整的手指骨頭.
經過考驗 驗證之後.
證明是釋迦牟尼的手指骨頭.
所以當時整個全世界的佛教界.
聽到這個消息的時候.
非常高興.
但是對於我們基督徒來說.
我們帶給我們一個什麼樣的信息呢.
令到我們更加喜樂.
因為釋迦牟尼真是死了.
骨頭都在那裡.
釋迦牟尼死了.
骨頭埋葬了直到如今.
耶穌基督卻勝過死亡.
他的屍體再不見.
骨頭再不會在那裡.
因為他已經什麼.
他已經復活了.
同樣是考古學家.

$^{1001}$歷史學家考古的證據證實.
找不到耶穌身體.
更加找不到他發霉的骨頭.
釋迦牟尼是一位偉人.
也是一個很虔誠的宗教信仰者.
但是可惜他不是.
可以帶我們勝過死亡.
得著永生百合的真正的救主.
因為他自己都沒有辦法勝過死亡.
基督徒敬拜的是一位復活的主.
永活的真神.
神的兒子.
是我們靈魂的救主.
賜給我們極大的信心喜樂.
但願主復活的大能.
繼續幫助我們今天.
我們人生會面對困難 試煉.
但是主的愛.
祂與我們同行的愛伴隨著我們.
主給我們的信心.
聖靈幫助我們的信心.
讓我們能夠超越這些試煉.
知道有一天.
這些最後要得著神給我們的榮耀.
尊貴和稱讚.
叫我們今天仍然緊緊跟隨著主.
靠祂復活的大能.
去榮耀上帝 顯出復活真生命的信徒.
同樣你還未相信接受耶穌.
還未得救.
或者未有重生的經驗.
快快來到主的面前.
懇求這位復活的主.
為你死赦免你的罪.
為你復活賜給永生的主.
賜給你靈魂的救恩.
每一個人我們都需要有這種復活活潑的盼望.
更加需要在今生面對各樣的患難危險的時候.
有這種超越的喜樂.
這種唯有復活的主能夠賜給我們.

$^{1041}$我們低頭同心祈禱.
我不知道在座當中.
或者在網上參與的.
有沒有哪一位還未相信主耶穌基督.
或者你未有重生的經驗.
你要得救這個新生命.
因為復活的主就是要賜新生命給我們.
如果你願意的.
今天你可以來到神的面前.
用信心接受祂.
打開你的心.
你說主呀 這位復活的主.
愛我的主.
求你進入我的心中.
賜給我重生的生命.
永生的把握.
在座當中.
不知道有沒有未信主的朋友或者慕道友.
如果有的請你輕輕舉起你的手.
我為你祈禱.
我為你祈禱 求神幫助你.
讓你接受主耶穌基督.
成為你的救主.
得著永生的生命.
不知道有沒有在座當中.
有沒有這些慕道的朋友.
或者在網上的.
參與的.
如果你還未信主.
希望你藉著復活節.
今天這個崇拜.
接受主耶穌.
做你新生命.
賜生生命給你的救主.
有沒有 我最後問一次.
有沒有的話.
你輕舉起你的手.
我特別為你祈禱.
網上的我看不到你.
不過我都會為你禱告.

$^{1081}$如果你願意的.
不是最重要我看到.
是上帝看到你.
上帝看到你的心.
是最重要的.
我們一起祈禱.
天父啊.
賜給我你的獨生兒子.
耶穌基督.
為我死在十字架.
擔當我的罪.
更加從死裡面復活.
證實你能夠賜永生給我的主.
我願意打開我的心門.
誠心接受你成為我的救主.
做我人生的主宰.
賜給我新的生命.
改變我的人生.
有個新的生活.
使我能夠得著.
你靈魂的救恩.
給我有活潑永生的盼望.
得著你與我在人生當中的.
木樣大靈和賜福.
求你聽我的祈禱.
我們天父.
我們再一次的為了.
我們為了所有.
我們已經信主的基督徒.
讓我們知道你復活的大樓.
幫助我們.
雖然有試煉.
但靠著你能夠得勝的人生.
讓我們的心事上因為靠近你.
得著這種愛裡面的甘甜.
得著超越環境的喜樂.
多謝讚美我們的主.
榮耀歸給我們的神.
歸給這位復活的基督.
奉主耶穌基督的名.

$^{1121}$阿們.
\newpage



\section{哥林多前書 8:6}
\label{sec:031930}
\textbf{不信三位一體的神能得救嗎?}~~~~袁惠鈞牧師博士
\newline
\newline
連結: \href{http://www.church.com.hk/acms/content.asp?site=cdc\&op=show\&type=product\&code=031930\&layout=sermon}{\texttt{031930}} ~~~~ 語音日期: 2026-04-05
\newline
\newline
\hyperref[sec:031929]{\small{< < < PREV SERMON < < <}}
~
\hyperref[sec:index_chronic]{\small{[返順時目]}}
~
\hyperref[sec:031931]{\small{> > > NEXT SERMON > > >}}
\newline
\newline
哥林多前書 8:6
\newline
\begin{longtable}{cl}
\hline
\hline
章節 & 經文 (和合本修訂版)\\
\hline
8:6 & \begin{tabularx}{0.7\textwidth}{X} 但是我們只有一位神，就是父，萬物都出於他，我們也歸於他；並只有一位主，就是耶穌基督，萬物都是藉著他而有，我們也是藉著他而有。 \end{tabularx} \\ \\
[1ex]
\hline
\hline
\end{longtable}
$^{1}$今天在復活節,應該跟大家講一個有關復活節的訊息,.
但我又隨即想到去年英文堂退休會,他們的主題是什麼?.
是「三位一體的神」..
所以我想想,咦,都挺好啊!.
在這個復活節,跟大家講一個有關「三位一體」的訊息,.
幫助我們能夠更加認識這一位為我們的救贖而死而復活的耶穌基督..
我給他的題目是什麼?.
「不相信三位一體的神能得救嗎?」.
我相信這是一個很多人都希望知道的答案..
在結束的時候,我會回來解答,究竟我們不相信三位一體的神能不能夠得救呢?.
首先,三位一體這個神學的教義,是非常之重要的教義..
但是很可惜,很多人都不明白,甚至連基督徒都不例外..
有很多非基督徒,例如回教徒,耶和華,尖精人,就更加不用說..
他們根本不相信三位一體的神..
他們也會指出,在聖經裡面根本沒有三位一體的字眼..
但是我們要知道,在聖經裡面有很多很重要的神學思想,都是沒有那個字眼的..
我只需要舉一個很簡單的例子,在聖經裡面有沒有主土文?.
絕對有,在《馬太福音》第六章,《路加福音》第十一章,耶穌基督教我們怎樣來禱告..
但是你在聖經裡面,你永遠找不到主土文這三個字..
這個就是一個例子.三位一體的神都是一樣..
聖經裡面滿有三位一體的教導,但是沒有三位一體這幾個字..
那在哪裡能夠見到有三位一體的神呢?.
我們很熟悉的大使命,在《馬太福音》的二十八章,十九至二十節:.
"所以你們要去使萬民作我的門徒,奉父子聖靈的名給他們施洗.凡我所吩咐你們的,都教訓他們遵守,我就常與你們同在,直到世界的末了.".
有沒有看到?是奉父子聖靈的名為他們施洗,這個就是三位一體的神..
還有呢,在高林多前書的十二章四至六節:"因是,緣有分別.聖靈卻是一位,職是也有分別.主卻是一位,公用也有分別.神卻是一位,在眾人裡面運行一切的事.".
有沒有看到?一位的神,一位的主,還有一位的聖靈,就是三位一體的神..
在二筆所書的四章四至六節:"身體只有一個,聖靈只有一個.正如你們蒙照同有一個指望,一主一信一洗一神,就是眾人的父,超乎眾人之上,貫乎眾人之中,也住在眾人之內.".
又是講到聖靈只有一個,一主一神,還講得很清楚這位神是誰,就是眾人之父,是天父.所以,給我們看到三位一體的神..
還有一節的經文是我時時都會用的,就是在雪徑拜之後,我都會用這節經文來祝福大家,就是高林多後書的十三章十四節:.
"願主耶穌基督的恩惠,神的慈愛,聖靈的感動,常與你們眾人同在,直到永遠.".
有沒有看到這裡已經講了三位一體的神,是有不同信徒的關係..
耶穌基督的恩惠,神的慈愛,聖靈的感動..
稍後我們會深入地來看這個經文.有人會說:"牧師,你讀的這些經文全部都是新約的.".
這個三位一體的玩意,是不是新約的使徒作出來的?舊約都沒有講到..
不是的,舊約都有的.我們可以看看舊約三位一體的教導..
首先我們從神的名字來看,我們可以先從一體來著手..
三位一體,一體的部分,相信所有的回教徒,耶和華見證人都會同意的..
一體的神,一位的神.這個亦是在《申命記》第六章第四節所講的:.
"以色列,你要聽!耶和華,我們的神是獨一的主.".

$^{41}$那你說這就是了.是獨一的主,你的基督徒搞出三位一體..
你知不知道這個"一"字,獨一的"一"字,可以解釋為和諧,可以解釋為合一..
在《以賽亞書》第四十三章,四十四章及四十五章,.
《以賽亞先知》亦有提醒我們:"除他以外,別無真神.".
聖經清楚告訴我們,有一位神,他的名字就是耶和華..
耶和華這個名字在希伯來文,是由四個子音組成,就是 YHWH..
這個字在聖經裡面,一共出現超過五千五百次..
猶太人是很怕讀這個字的,因為這個字很神聖..
所以他們害怕,如果他們說出來讀的時候,就會觸犯了十誡裡面哪一條呢?.
第三誡不可妄清神的命..
所以他們不敢讀..
但是猶太人寫希伯來文的時候,他們只寫子音,不寫母音..
所以你不讀,不讀,慢慢你就不會讀..
看到你的 YHWH 怎麼讀?後來就有拉比,是將希伯來文裡面的主這個字 "Adonai",.
裡面的母音加上 YHWH 裡面的子音,就變成了什麼?.
耶和華, Jehovah 或者是 Jehovah..
所以就有了耶和華這個讀音..
但是其實我相信真正的讀音應該是 "Yahweh",而不是耶和華..
我們講了一體的部分..
現在我又要講講三位的部分..
又是從神的名字來看..
神的名字除了這個很神聖的 YHWH 耶和華之外,.
還有一個就是 Elohim..
Elohim 這個字是在聖經裡面出現超過 2200 次,是一個複數,是一個眾數來的..
這個字的單數是 L,或者是 Eloah,是單數..
在聖經裡面只是出現超過 300 次,但是複數,眾數是超過 2200 次..
我們怎麼知道這個字是眾數,複數?.
在聖經裡面有很多例子..
例如天使有不同的天使..
有一種天使叫做撒拉佛,是 Seraph..
Seraph 這個字就是單數..
眾數或者複數是什麼? Seraphim..
這一種天使 Seraph 或者 Seraphim,在《爾賽亞書》第六章有講到..
Seraphim 就是和 Elohim 一樣,是眾數或者是複數..
另外一種天使就是基魯伯..
Seraph 這種天使,Seraph 是單數..
這種天使很可怕,有四個面..
有獅子的面,人的面,牛的面和鷹的面..
是很可怕的..
在哪裡出現呢?在《創世記》的第三章,以及《爾西傑書》的第一章..

$^{81}$單數就是 Seraph,眾數或者複數就是 Seraphim..
所以大家能夠見到 Elohim 這個字,的的確確是複數或者是眾數..
這樣就很有意思了..
在《聖經》裡面一開始,《創世記》第一章第一節講什麼呢?.
「起初神創造天地」..
「神」這個字就是用 Elohim 是複數或者眾數..
但是「創造」這個動詞呢,是單數的動詞..
有沒有見到,就是三位一體的「神」,就在《術語》這一句話裡面..
其實在《聖經》裡面還有很多其他的代名詞,.
講到「神」是用眾數或者複數,而不是單數..
在哪裡呢?《創世記》的一章二十六節,三章二十二節,十一章七節到八節,《爾西傑書》的六章八節..
我們首先來看看《創世記》的一章二十六節和二十七節,很有意思..
《創世記》的一章二十六節,神說:.
「我們要照著我們的形象,按照我們的樣式做人,.
使他們管理海裡的魚,空中的鳥,地上的生畜和全地,並地上所爬的一切昆蟲.」.
有沒有看到這裡照著我們的形象,照著我們的樣式,是眾數是複數..
但是到了下一節的經文,就是二十七節,.
神就照著自己的形象做人,乃是照著他的形象做男做女..
一節的經文是用到眾數,複數的代名詞,.
另外一節又用到單數,三位一體的神,非常清楚..
有人會問:.
「一位神就好了,好像回教徒,耶和華見證人相信一位神.」.
為何你們基督徒要搞出三位一體?.
其實三位一體的確是聖經的教導..
一體的神,為甚麼?.
因為三位位格的神,都有一樣的神性,也有一樣的權柄,也有一樣的旨意,所以他們是一體的..
但是為何是三位?.
因為他們有不同的特徵,不同的任務,也有跟信徒不同的關係,所以是三位位格..
我稍後會再詳細解釋..
其實三位一體的神,用圖解來代表,就好像是三個圓圈,代表三個位格..
但是因為是三位一體的神,就變了一體..
所以三位一體,就好像我剛才的圖表一樣,.
三個位格的神,分開三個圓圈,但是三個都是合為一體,這個就是三位一體..
也因為這樣的時候,有很多錯誤的教導..
例如有人說根本是三位神,父是一位神,子是一位神,聖靈也是一位神,多神論..
這個是錯誤的教導,是一位神,三個位格在一起..
還有另外一個錯誤的教導,是叫做形態論..
你說形態論是什麼?始終都是一位神,但是用不同的形態來出現,就好像一杯水一樣..
一杯水是水來的,你把它放進冰格裡會變成冰..
如果你把它拿出來,放進蒸爐裡蒸一下,又會變成蒸汽..

$^{121}$但是來來去去都是那杯水,一會兒就變成冰,一會兒就變成蒸汽,一會兒就變成水,來來去去就是那杯水..
就好像有人說,形態論的神是什麼?.
是的,最初就是天父在雪天堂,後來天父在天堂就說我要救贖人類..
於是這位天父就道成肉身,變成了耶穌基督,來到世上為我們釘十字架,然後祂復活..
復活了之後,變成了聖靈,然後就在我們每個信徒的身上..
有沒有看到?來來去去都是這位神,是以不同的形態出現,就好像冰水和蒸汽一樣..
這個其實是異端來的,不是正確的教導..
最好否定形態論的經文在哪裡?就是耶穌基督受洗的時候..
我們知道,耶穌基督受洗是很重要的經文,在四福音裡面,四部福音書都有講到..
當耶穌基督被施洗之後,祂是在水裡面上來,天忽然為祂開了..
然後經文說,神的靈紡不像鴿子降落在祂的身上..
有沒有看到聖靈和耶穌基督同時出現?.
還不止,還有聲音在說天上說:「者是我的愛子,是我所喜悅的.」.
三位一體的神是同時出現的..
不是說來來去去都是那位神,不是說來來去去都是那杯水,.
而是水,精氣和冰一起出現,而是三位一體的神能夠一起出現..
我剛才說冰,水和精氣不是一個三位一體很好的比喻,.
就是因為普通人都沒有見過水,精氣和冰同時出現,你見過沒有?.
沒有見過,但是真的會有發生的..
在樹壓力是大氣層的千分之六,.
當氣溫是百分之一的攝氏的時候,有一個地方叫做「Triple Point」..
這個「Triple Point」就是水,精氣和冰同時出現..
我覺得這個就是神留在樹世間,給我們見到一個非常好代表三位一體的象徵..
我以前在UCLA和USC教書的時候,我時時要解釋這個「Face Diagram」..
我和學生們是不可以傳福音的,我是教科學的,我不是教神學..
但是我很喜歡指出一件事,就是你有沒有見到水和冰當中的那一條界線?.
是紅色那一條,是由右下角去到左上角的..
這個就是它的傾斜度..
普通的物質是怎樣的?.
這個傾斜度是由左下角去到右上角的..
你說這個分別在哪裡?這個分別就大了..
因為水的傾斜度是這樣,水和冰的交界的傾斜度..
以致冰的密度是比水小的..
但是其他的物質,冰永遠是重過水的,因為它是固體的..
固體的密度當然是重過液體的,但是水就剛剛相反..
有什麼那麼重要?.
如果你說有一個湖,表面是積滿了冰,因為天氣很冷..
那它就浮在水面上,你去買一杯波霸奶茶,.
那些冰浮在上面,不會浮在杯底上..
那下面的魚就可以生存了..

$^{161}$但是如果你的傾斜度,是好像其他物質一樣,.
即是說冰是重過水的時候,那就大件事了..
結了冰之後,冰就沉下去湖底,.
第二層的冰又結冰,又沉下去湖底,到最後會怎樣?.
整個湖結滿冰..
有沒有生物?沒有生物..
這個就是我時時想指出給學生看,.
但是在這個宇宙當中,是有一位造物者,.
他是很愛惜他的造物,愛惜生物..
他做到水是這樣,.
以致那些生物在湖底,.
當很冷很冷的時候,仍然都能夠在湖底,在海底有生物..
其實,一個很好的圖解來解釋三位一體的神,.
就是這個三角形..
大家看看這個螢光幕..
一個三角形有多少隻角?.
三隻角,是吧?.
剛剛好就是,一隻角是代表天父,.
一隻角是代表聖子,一隻角是代表聖靈..
如果你注意到的話,你會見到三角形有三條邊..
每條邊都寫著「不是」..
即是說甚麼?.
聖父不是聖子,聖子不是聖靈,.
聖靈亦都不是聖父..
但是在這個三角形最中間的地方,有一個「神」字..
而在每隻角都有一個箭咀指著這個「神」字,又代表甚麼?.
天父是神,耶穌基督是神,聖靈都是神,.
但是「父」不是「子」,「子」不是聖靈,是吧?.
所以這個就是一個非常好的圖解,.
讓我們見到三位一體的神是怎樣的..
我剛才也說過,為甚麼我們需要有三位一體的神?.
因為一體的神,就是因為三個位格都有同一樣的神性,.
有同一樣的權柄,同一樣的旨意..
為甚麼需要三個位格?.
因為三個位格的神都有不同的特徵,.
和信徒都有不同的關係..
我們先看一體的部分,為甚麼是一體?.
就是因為三個位格的神都有一樣的神性,.
有甚麼神性呢?.
他們都是無所不知,都是無所不能,都是無處不在,都是自有榮有..

$^{201}$我們先來看看第一點,就是無所不知..
看著歷代至上的二十八章第九節,.
經文說:「他監察眾人的心,知道一切的心思意念.」.
我們這位神是怎樣的?.
他監察人心,他知道人想甚麼,他是無所不知..
你說我們的天父呢?.
在馬太福音的六章第八節,耶穌基督說:「因為你們沒有祈求以先,.
你們所需用的,你們的父祖已知道了.」.
是否在說你地上的父親?.
不是,是在說天父..
天父知道你需要甚麼,他是無所不知..
那你說耶穌基督呢?.
當耶穌基督在世的時候,有人抬了一個毯子來給祂依治,.
祂就對毯子說:「你的罪被赦免了.」.
當時有一群民事,他們的心裡面都在想,他們沒有發聲,.
只在心裡面在想:「嘩!這個人在說潛亡的話,怎麼可以這樣說呢?」.
只有神能夠赦免人的罪..
耶穌基督知道他們的心在想甚麼,.
馬太福音九章四節和馬可福音二章八節都有說,.
所以耶穌基督都知道他們的心在想甚麼,他是無所不知..
在約翰福音第二章為耶穌做了很多的神蹟,.
有很多人當時說自己信了,.
但是在約翰福音二章二十四節經文說甚麼?.
耶穌知道他們不是真正相信,沒有將自己交給他們..
所以耶穌絕對是能知一切..
有人說:「牧師,你不要騙我的膝蓋不吃辣椒醬!」.
我讀聖經,我都知道聖經說甚麼,.
馬太福音二十四章,馬可十三章都說耶穌不知道自己何時會再次回到這個世上..
你說祂知道一切?.
耶穌是百分之百的神,祂是能夠知道一切,.
但是祂也是百分之百的人,當祂在世上的時候,.
祂是甘願有些事祂不去知道,.
而祂何時回來就是其中之一樣..
祂是百分之百的神,祂也是百分之百的人,.
有時祂是有人的限制..
你說聖靈又怎樣?.
在爾塞亞書四十章十三節,.
說:「誰曾測度耶和華的心,或作他的謀事,指教他呢?」.
我們是人,我們真的不能夠知道一切..
但是在哥林多前書二章十至十一節說甚麼呢?.

$^{241}$聖靈能參透萬事,.
所以無論是父子聖靈,都是無所不知..
還有第二點,就是無所不能..
在約百多的四十二章第二節,約百話神萬事都能作,.
祂的旨意不能被攔阻..
在創世記十八章第四節,.
為耶和華喜有難成的事,.
我們見到天父真是無所不能..
耶穌基督呢?.
記不記得福音書裡面告訴我們,.
耶穌基督的能力有多大?.
祂能夠醫治趕鬼,祂能夠平靜風浪,.
連大自然都要聽祂的話..
祂還能使人死裡復活..
還有呢,稍後我們都會見到在約翰福音一章三節,.
說萬有都是耶穌基督所做的..
還有呢,在喜百來書的一章三節,說祂托住萬有..
在哥勞西書的一章十七節,說萬有都靠祂而立..
但是,亦都有人說,.
喂!牧師,不妥喲!.
喂!你說耶穌基督是無所不能?.
但是在約翰福音第四章,.
當耶穌基督去到撒瑪利亞婦人在亞國的井旁,.
想跟撒瑪利亞婦人傳福音的時候,.
祂已經筋疲力倦,.
祂又很口渴..
哇!你說祂是無所不能,為什麼祂會這樣?.
喂!不奇怪!在馬太福音的第八章,.
馬可的四章,以及路加福音的第八章,.
說到耶穌基督平靜風浪之前,祂在做什麼?.
祂在睡覺.哇!你說,.
無所不能的神都要睡覺嗎?.
是啊!祂是百分之百的神,.
但是祂亦都是百分之百的人,祂是有人的限制,有人的需要,.
所以祂亦都是會疲倦,亦都需要睡覺..
那你說聖靈呢?.
撒瑪利亞書的四章六至七節,.
說萬君之耶和華說,.
「不是依靠勢力,不是依靠才能,乃是依靠」什麼話?.
「我的靈光能成事,大山啊,你算什麼呢?」.

$^{281}$「你必成為平地.」.
聖靈亦都是無所不能..
而保羅在羅馬書的十五章十八節,.
說到祂是有神蹟其事的能力,.
祂是歸功於誰?是聖靈..
是聖靈幫祂有神蹟其事的能力..
所以我們見到聖靈亦都是無所不能..
三位一體的神都是無所不能..
還有呢?三位一體的神都是無處不在..
在耶利米書的二十三章二十四節,.
說到人不能夠在隱密處藏身..
為什麼人不能夠躲起來?.
因為神是充滿天地..
如果神是充滿天地,我們能不能躲起來?.
我們不能躲起來..
而在賽亞書的四十三章二節亦說到,.
當我們在水中經過,在火中經過,.
神都與我們同在..
這一節的經文給予我們很大的安慰..
無論我們在做什麼,.
雖然經歷苦難,但是神都與我們同在..
這就是我們的天賦..
那麼你說:耶穌基督呢?.
你不要告訴我祂都是無處不在..
首先我們要說的是耶穌基督的神聖..
在二十三章二十四節說到祂的神聖是怎樣呢?.
「是那充滿萬有的.」.
這裡給我們看到,原本耶穌基督的神聖是與天賦一模一樣,.
是無處不在..
但是因為祂是造成肉身來到世上,.
祂在世上是有祂的限制,是有祂的限制..
祂在世上的時候有一件事我們就能夠看到,.
就是祂的好朋友拉撒路死了,.
耶穌基督要去見祂..
門徒就跟耶穌基督說:.
「拉比啊?你夠膽去見祂?」.
「現在的猶太人可能會用石頭扔死你,你夠膽去?」.
耶穌基督在約翰福音十一章九至十節回答說:.
「白日不是有十二小時嗎?人在白日走路,就不至跌倒,因為看見者世上的光.」.
「若在黑夜走路,就必跌倒,因為他沒有光.」.

$^{321}$這節經文最好的解釋就是,.
耶穌基督知道祂距離釘十字架的時候不遠了,.
祂很快要釘十字架,祂還有很多事情要做..
就是因為祂在世上不單止是百分之百的神,祂是百分之百的人..
祂是有時空的限制,祂在世上還有很多事情要做..
這裡就讓我們看到..
但祂的神聖亦都好像天父一樣,是充滿萬有..
那你說聖靈呢?.
這一句的詩歌,是大衛的詩歌..
在詩篇的一百三十九篇,.
也是我很愛讀的一個詩歌,就是:.
「我往那裡去逃避你的靈,我往那裡逃避你的面..
我若升到天上,你在那裡,我若在陰間下塔,你也在那裡..
我若展開清晨的翅膀,飛到海極居住,.
就是在那裡,你的手必引導我,你的右手也必扶持我.」.
大衛的詩歌,「我往那裡逃避你的靈,父子聖靈都是無處不在.」.
還有最後,關於神聖的一點,就是父子聖靈都是自有永有的神..
其實自有永有的神的經文,你可以在《二賽亞書》的四十章,五十七章,.
以及《申命記》的三十三章,你都能夠見到..
但是相信我們最記得的,就是出埃及記的第三章,第十四節..
記不記得摩西去見神,神在《樹經的火焰》裡面,.
神要摩西帶領以色列人出埃及,摩西就說:.
「如果以色列人問我,你的名字叫做什麼,我怎麼回答?」.
神就跟他說:「我是自有永有的神.」.
但是我相信當年在《樹經的火焰》裡面的那位,就是耶穌基督..
為什麼我會這樣說呢?.
你看看出埃及記的三章二節:.
「耶和華的使者從經濟裡火焰中,向摩西顯然.」.
有沒有看到是誰向摩西顯然?耶和華的使者..
在三章四節:「耶和華神見他過去要看,就從經濟裡呼叫說:摩西!摩西!他說:我在這裡.」.
一節的經文說是耶和華的使者,一節的經文說在經濟裡面的是耶和華神自己..
也就是說耶和華的使者和神是可以相替換..
耶和華的使者就是神,神就是耶和華的使者..
為什麼要說耶和華的使者呢?.
這裡就是在舊約給我們見到三位一體的神..
神學家都一致認為,在舊約的耶和華使者就是未道成肉身的耶穌基督..
因為神是零,沒有人能夠見到神..
所以在舊約給人見到的耶和華的使者都是未道成肉身的耶穌基督..
那你說:「嗯?未道成肉身,那祂在哪裡找一個身體呢?」.
這個就是一個奧秘,我們不知道..

$^{361}$但是相信祂是要附在人的身上,給人見到..
亦好像跟瓦割蜀交那樣,你要跟瓦割蜀交,你一定要有一個身體..
所以在舊約的耶和華的使者都是未道成肉身的耶穌基督..
另外有一節的經文很有趣,約翰福音的八章五十八節..
耶穌說:「我實實在在地告訴你們,還沒有亞伯拉罕就有了我.」.
這節的經文,你讀中文的版本好像沒有什麼奇怪..
但是你讀英文就很有趣了,因為你能夠見到文法..
英文是什麼?「Before Abraham was, I Am」..
有沒有見到這裡很奇怪?.
用在亞伯拉罕身上是什麼?.
是「was」,是過去式..
但是用在耶穌基督身上呢?「I Am」..
不單止在亞伯拉罕的時候,祂是「I Am」..
祂說這句話的時候,亦都是「I Am」..
將來祂亦都是「I Am」..
祂是誰?祂就是自有永有那位..
所以我說祂就是在出埃及記三章十四節和摩西見面那一位..
在啟示錄的二十二章十三節,.
耶穌說:「我是阿拉伯,我是俄媒家.我是首先的,我是門後的.我是初,我是終.」.
有沒有見到?祂就是自有永有那位..
你說聖靈呢?.
在希伯來書九章十四節,叫聖靈做永遠的靈..
所以三位一體的神,都是自有永有..
都是有一樣的神聖,所以是一體的神..
還有,除了祂們是有一樣的神聖,是一體,祂們都有一樣的權柄..
有什麼權柄?.
有創造,有謝罪,有使人清義,使人成聖,使人復活..
我們可以一樣一樣來看..
創造先..
耶穌基督在約翰福音的五章十七節,祂說:「我父自創造以來,作事直到如今.」.
給我們見到什麼?.
是天父創造一切,祂說:「我父創造以來.」.
但是祂創造之後,祂是否立即停守休息?.
不是啊!耶穌說祂創造以來,祂是一直來掌管一切的運行..
耶穌基督也有參與一切..
在約翰福音的第一章第三節經文中說:「萬事都是藉著祂造的.」.
這個「他」就是耶穌..
「凡地造的,沒有一樣,不是藉著祂造的.」.
有沒有看到?.
你說是天父創造,現在又說是耶穌基督創造?.

$^{401}$是一體的神..
聖靈呢?聖靈也有參與創造..
在《創世記》的一章二節說:「神的靈運行在水面上.」.
天父當神是創造宇宙一切的時候,神的靈運行在水面上..
你說:「牧師,這裡不是創造,這裡是運行在水面上.」.
天父當神是創造宇宙一切的時候,神的靈運行在水面上..
三位一體的神也有創造,所以祂們是一體..
還有呢?卸罪的權柄..
三位一體的神同樣有卸罪的權柄..
只有誰能夠卸罪?.
只有神能夠卸罪..
《基督文》耶穌基督教我們:求天父赦免我們的罪..
在《馬太福音》第六章..
你說:天父赦免罪,耶穌基督呢?.
我剛才也說過,有人抬著毯子給耶穌基督依治,祂說:你的罪被赦免..
還有,《馬太福音》九章二節,《馬可》二章五節,.
《路加福音》五章二十節,耶穌基督都有赦免人的罪..
所以父子也一樣..
那你說聖靈呢?.
在《約翰福音》二十章二十一至二十三節,耶穌又對他們說:.
願你們平安,父怎樣差遣了我,我也照樣差遣你們..
說了話就向他們吹一口氣說:你們受聖靈..
你們赦免誰的罪,誰的罪就赦免了..
你們留下誰的罪,誰的罪就留下了..
這裡是不是說門徒可以赦免人的罪?.
當然不是啦,人不立過赦免人的罪..
這就是為什麼耶穌基督要先吹一口氣,將聖靈賜給祂,是一個象徵式..
當耶穌基督真的為我們將十字架復活之後,聖靈就降臨,.
感動人去接受耶穌基督..
當那些使徒有聖靈在他們身上,走向傳福音的時候,對方如果接受耶穌基督,.
他們的罪被赦免,是誰赦免他們的罪?.
很明顯是聖靈..
所以父子聖靈都有赦免人的罪的權柄..
還有就是使人清義..
在羅馬書裡我都和大家講得很清楚,什麼叫做清義..
就是當一個罪人接受耶穌基督的時候,他的罪被赦免,.
耶穌基督的義披在他身上,.
以致神看到他,不是看到他的罪,是看到耶穌基督的義..
以致我們能夠與神和好..
然後慢慢我們才越來越成聖,越來越像耶穌基督..

$^{441}$所以清義就是當一個人接受耶穌基督的時候,被神稱為一個義人..
禮爾賽亞書的一章十八節,.
說你們的罪須像朱紅必成雪白,須紅如丹顏必白如羊毛..
天父怎樣?使我們朱紅的罪都能夠變成雪白,透過耶穌基督的救贖..
那你說耶穌基督呢?.
羅馬書的十章四節,律法的總結就是基督,.
使凡信他的都得著義..
迦拉契書的二章十節也說,.
既知道人稱義,不是因行律法,是因信基督稱義..
所以我們見到,你是接受耶穌基督的話,你就可以被耶穌基督稱為義人..
無論聖父和聖子,都能夠使人稱義..
還有聖靈呢?.
在哥林多前書的六章十一節,.
如今你們奉主耶穌基督的名,.
並藉著我們神的靈,已經洗淨,成聖稱義了..
這節的經文很有意思,.
我們見到聖靈不單止使人稱義,還怎樣說?.
使人成聖,兩樣東西都加在一起..
所以這就帶領我去第四點,.
三位一體的神,都能夠使人成聖..
成聖是什麼?就是我們在基督徒的生命當中,.
我們靠著聖經潔淨我們,越來越聖潔,.
越來越像耶穌基督,這個就是成聖..
通常都是聖靈的工作,我剛才也說過,.
哥林多前書的六章十一節,說聖靈幫助我們成聖..
在彼得前書的一章二節,也有同樣的經文..
那你說天父和耶穌基督呢?.
在鐵撒羅尼加前書的五章二十三節,.
說願賜平安的神,親自使你們傳言成聖..
有沒有見到?天父使我們傳言成聖..
耶穌基督呢?希伯來書的二章十一節,.
因那使人成聖和那些得以成聖的都市出於一,.
所以他稱他們為弟兄也不以為此..
誰稱我們為弟兄呢?耶穌基督..
所以父子聖靈都能夠幫助人成聖..
還有最後一點,與我們今天的復活節息息相關,.
三位一體的神都能使人復活..
在約翰福音的五章二十一節,.
父怎樣叫死人起來,使他們活著,.
子也照樣隨自己的意思使人活著..

$^{481}$一節的經文就說出,.
天父和耶穌基督都有一樣的權柄,能夠使人復活..
聖靈呢?在羅馬書的八章十一節,.
願意叫耶穌從死裡復活者的靈,.
若住在你們的心裡,.
那叫基督耶穌從死裡復活的,也必藉著住在你們心裡的聖靈,.
使你們必死的身體又活過來..
這節的經文對我們每一個基督徒都很重要..
在復活節我們說,我們都有永生的盼望,我們都會復活..
是誰使我們復活呢?這節經文說得很清楚..
如果我們真的接受耶穌基督,我們身上是有聖靈,.
當我們死了之後會怎樣?.
我們如果真正相信的話,亦都會復活..
就好像聖靈使到耶穌基督復活一樣..
我們說了一體的部分,為什麼是一體的神?.
因為三位一體的神都有一樣的神性,都有一樣的權柄,所以是一體..
那為什麼需要是三位?.
我們現在要說三位..
為什麼呢?因為父子聖靈都有不同的特徵,就好像冰水和蒸氣一樣..
父的特徵是什麼?父是基督的頭,這個就是父的特徵..
但是在說這一點之前,我首先要說的就是,三位一體的神都是平等的..
不是說父是高過子,或者是高過聖靈..
三位一體的神都是平等的,只是一個次序的問題..
在這裡就和《以不所書》五章二十二節和二十三節所說的很相似..
說到丈夫和妻子,男女是不是平等的?絕對是平等的..
那為什麼經文說丈夫是妻子的頭?是一個次序來的..
神是首先創造當,再創造下,是有一個次序的..
不是說丈夫就是大過妻子,是平等的..
所以我們看到,父就是基督的頭,基督遵行父的旨意,約翰福音的五章十九節..
耶穌基督不會做任何照祂自己意思的事,祂絕對是服從天父,就是因為父是基督的頭..
而且耶穌基督現在已經回到天父的右手邊,為信徒祈求,向天父祈求..
這一節的經文很多人都不明白,我有一個朋友,他就不明白這一節的華語,他不能夠接受..
羅馬書的八章三十四節,希伯來書的七章二十五節,和約翰一書的二章一節都說得很清楚,耶穌是向天父祈求..
我的朋友說,耶穌也是神,天父也是神,怎麼可能向天父祈求呢?.
這個就是因為父是基督的頭,是有一個次序的問題,而耶穌基督是服從天父..
為什麼祂要為信徒祈求呢?因為撒旦不斷地在攻擊基督徒,就好像祂當年攻擊約伯一樣,所以耶穌基督是為我們祈求..
很多經文都是這樣說的,聖靈也是為信徒來向天父祈求..
為什麼呢?羅馬書的八章二十六節告訴我們,因為我們不懂得如何禱告,所以聖靈是幫助我們..
說了聖父的特徵,現在第二,我要說聖子耶穌基督的特徵..
耶穌基督的特徵是什麼呢?道成肉身來到世上.是不是天父道成肉身來到世上呢?不是..

$^{521}$是不是聖靈道成肉身來到世上呢?不是,是耶穌基督..
為我們釘十字架,今也復活節,為我們贖罪復活,讓我們有永生的盼望..
還有呢,祂的特徵,祂要彰顯出天父..
我剛才也說過,天父和聖靈都是靈,從來沒有人見過..
所有世上見過神的是誰呢?是道成肉身的耶穌基督..
我們怎知道,約翰福音的一章十八節說,從來沒有人看見神,只有在乎懷裡的獨生子將祂顯明出來..
所以,耶穌的特徵是什麼呢?祂是道成肉身來到世上,祂是彰顯出父的榮耀..
那麼你說,聖靈的特徵呢?聖靈是附在舊約人的身上,我說附在處,不是住在處..
當舊約的人屬靈的時候,聖靈就會臨到他們的身上,就像參宣和素羅一樣..
但是當他們不屬靈的時候,聖靈是會離開的..
但是在新約,聖靈是住在人的身上,是住在每個信徒的身上,是不會離開的..
如果你是一個真正的基督徒,聖靈是不會離開你的..
而且聖靈從來不會說關於自己的事,祂永遠是將人指向耶穌基督..
剛才我們也說耶穌基督是服從天父,聖靈也是不說自己的事,將人指向耶穌基督..
還有祂是感動人進入父那裡,以忽所書的二章十八..
稍後我們會回來看這一節的經文..
三位一體的神,我們看到,再重複一次,為什麼是一體?有一樣的神聖,有一樣的權柄,一樣的旨意..
為什麼是三位?因為有不同的特徵,他們和信徒的關係都不同..
你說,父子聖靈和信徒有什麼關係?.
天父是在宇宙之先已經簡選了我們,是照著祂的旨意所預定的,以忽所書的第一章..
這就是天父和我們的關係,是祂簡選我們..
耶穌基督是祂渡成肉身來到世上為我們釘十字架..
現在祂成為我們的宗保,在天父面前為我們祈求,為我們辯護,這就是耶穌基督..
你說聖靈和我們的關係呢?.
首先,聖靈感動一個人能夠接受耶穌基督..
接受了之後,祂是住在這個人的裡面..
然後祂是幫助帶領我們成聖..
然後祂是幫助我們能夠結出聖靈的果子..
加拿大書的五章二十二節,有沒有看到三位一體的神和我們的關係都不同?.
所以是需要三個位格,這就是原因..
我剛才是指出了一個三角形,我說是很好的三位一體的一個圖解..
現在又有一個三角形..
這個三角形是本於我們今天的經文..
剛才我們開始的時候,已經讀了兩節的經文..
一節的經文是哥林多前書的八章六節,另外一節的經文是以弗所書的二章十八節..
我們先看哥林多前書的八章六節..
然然我們只有一位神,就是父,萬物都是本於他,我們也歸於他,定一位主,就是耶穌基督,萬物都是藉著他遊的,我們也是藉著他遊的..
這裡講到三位一體其中兩位位格的關係..
一切都是本於父,是祂揀選我們..
一切都是藉著耶穌基督在十字架上所成就的..

$^{561}$這個部份是講出三位一體其中兩個位格..
以弗所書二章十八節是講另外一位位格,就是聖靈..
以弗所書的二章十八節:.
因為我們兩下藉著他對一個聖靈所感,不以進到父的面前,我們是被聖靈的感動,以致我們是能夠進到父的面前..
我們看到三位一體的神的運作,一切都是本於父..
是祂揀選的,是藉著耶穌基督在十字架上所成就的..
接著是被聖靈感動,以致我們是能夠接受耶穌基督,是去到神的面前..
有沒有看到開始和結束都是父,是藉著耶穌基督為我們做,被聖靈感動..
這個就是我們得救的一個圖解,三位一體的神的運作,我們怎樣能夠得救..
在結束的時候,我們再稍為溫習一下,為什麼是一體的神?.
因為有一樣的神聖,有一樣的權柄,有一樣的子義,三位一體的神..
為什麼是三位?.
因為父子聖靈都有不同的特徵..
起碼一樣,耶穌基督導誠肉身,是能夠見得到,父和聖靈都見不到..
除了特徵之外,他們和每一個信徒的關係都不同,所以需要三位..
我現在又回到一開始的時候,問一個問題,不相信三位一體的神能不能夠得救?.
我可能可以用一個比喻,你突然間走過去,見到前面有一間房子,你打開了門進去,.
見到桌上有一隻美味的佳餚,你馬上就狼吞虎嚥吃完這頓飯..
有人會問你,你想不想知道誰為你預備這一餐美味的佳餚?.
你說:不用了,我吃了就夠了..
有很多人都是這樣的,說自己接受了耶穌基督,我得救了..
他們會說什麼?聖經說的,我是真正口裡承認,心裡相信..
那我就得救了,我還需要知道這麼多嗎?.
喂!你讀漏了經文,我講《羅馬書》第十章的時候,我都講過了..
這節經文全部是講什麼?你要口裡承認,心裡相信,神叫耶穌從死裡復活..
就是我們今天的主題,就是復活節..
你要認識這位神,祂是為你死而復活的那一位..
不是單單說你口裡承認,心裡相信,你就能夠得救..
所以如果你是一個真正得救的人,你會怎樣?.
你會去認識這位為你死而復活的神,你會去認識這位三位一體的神..
你說:我接受了耶穌基督,我都不知道三位一體是什麼..
Never too late!你仍然可以去認識這位神..
而一個真正得救的人,很清楚,大家聽著..
一個真正得救的人,不會說得救就趴屁股離開..
一個真正得救的人,是會繼續去認識這位主耶穌基督,這位三位一體的神..
而一個不是真正得救的人,完全沒有興趣,不會去追求,不會去認識這位三位一體的神..
願神祝福他自己的話語,讓我們一起低頭禱告..
願我們不單單這樣,自己得救就算了..
而是我們真真正正有這個心,來認識你這位三位一體的神..
因為一個真正得救的人,就一定會有這個心,來更加的認識你,更加的閱讀你的華語,明白你的心意,明白你的旨意..

$^{601}$願當你回來的時候,我們每一個基督徒都能夠被你稱為一個忠心良善的牧人,將那榮耀頌讚歸給你..
我們這樣禱告,奉主耶穌基督的名義求.阿們..
\newpage



\section{出埃及記 12:30}
\label{sec:031931}
\textbf{聖經人聖經事:無一家不死一個人}~~~~黃克勤牧師博士
\newline
\newline
連結: \href{http://www.church.com.hk/acms/content.asp?site=cdc\&op=show\&type=product\&code=031931\&layout=sermon}{\texttt{031931}} ~~~~ 語音日期: 2026-04-04
\newline
\newline
\hyperref[sec:031930]{\small{< < < PREV SERMON < < <}}
~
\hyperref[sec:index_chronic]{\small{[返順時目]}}
~
\hyperref[sec:031932]{\small{> > > NEXT SERMON > > >}}
\newline
\newline
出埃及記 12:30
\newline
\begin{longtable}{cl}
\hline
\hline
章節 & 經文 (和合本修訂版)\\
\hline
12:30 & \begin{tabularx}{0.7\textwidth}{X} 法老和他眾臣僕，以及所有的埃及人，都在夜間起來了。在埃及有大大的哀號，因為沒有一家不死人的。 \end{tabularx} \\ \\
[1ex]
\hline
\hline
\end{longtable}
$^{1}$各位弟兄姊妹,大家平安!.
生命樹宣教網絡又邀請大家一起讀聖經.
聖經人物,聖經故事.
一起看看在人與事的背後.
上帝做些什麼,想些什麼,認識神更多.
今天的題目仍然是摩西大以色列人出埃及的故事.
今天的題目是很慘的.
無一家不死一個人.
你記不記得是什麼事呢?.
無一家不死一個人.
當然,那一家又不是長子.
又沒有生小孩.
那就沒有死了.
所以概括地說,這是殺長子之災.
無一家不死一個人.
看看聖經.
出一及記十二章二十九節.
第一點想和大家看的,是《到了半夜》.
凡是用黃色螢光,有一個紅色的數目字.
我要講每一段經文,和大家看聖經人士.
我都會用三點,我講到多數都是三點.
最好的演講都是三點.
第一,到了半夜.
耶和華把埃及地所有的長子.
就是最高的作寶座的法老.
直到最低的是被擄囚在監獄之人的長子.
以及一世投生的生畜盡都殺了.
第一點要和大家看的,到了半夜.
這個到了半夜是什麼意思呢?.
就是沒有人想到的時間.
沒有人預算到.
沒有人會做好準備.
例如主耶穌都說,盜賊來偷東西什麼時候?.
是半夜你不知道的時候挖通房屋.
當然我們都很難明白.
如果現在偷東西,怎會挖通房屋?.
因為都是石屎來的.
但是主耶穌挖通房屋.
就是因為當時他們住的是泥屋.
泥和石砌的屋.

$^{41}$是有機會用水濕一下,就會挖通房屋.
沒有水都挖得通.
用水就是溶那些泥.
半夜的時候你不知道,到了就來了.
所以要警醒,就是這樣講.
這裡講到了半夜.
上帝的工作,上帝的手什麼時候做事.
我們不知道.
說不定是現在,說不定是半夜.
你沒有想到,這樣就發生了.
耶和華是將埃及地所有的長子.
我剛才讀的時候都提了.
最高的坐寶座的法老.
家裡都被牢囚在監獄裡的人.
被牢就是埃及出去打仗.
就將很多人,當時都是流行.
將別人的人民百姓,看到就抓回來.
做什麼呢?回來做奴隸嘛.
男的就可以做操工.
女的就可以服侍.
大隻一點的,跟著出去幫忙耕田.
出去建城牆,建什麼都可以.
被牢的已經最低微了.
完全沒有身份,沒有地位.
好像一隻牛一樣.
由法老最高,到被囚在監獄裡的人的長子.
以及一切投身的生畜.
總之他是長子身份第一隻.
生出來第一隻的生畜.
盡刀殺了.
法老和一切神僕被埃及眾人夜間都起來了.
長埃及有大埃號.
無一家不死一個人的.
這裡有些竅妙.
大家要想想.
那個被殺死的長子.
是不是正在睡覺.
死了.
不是,應該他睡覺的時候.
熟睡的時候.

$^{81}$應該是有一個痛苦的喊叫.
如果不是,家人不會醒來的.
不會知道他死了.
可能睡著睡著.
那個大兒子的長子忽然之間.
啊,什麼事呢.
當然最有可能第一個醒來就是媽媽.
最聽到小孩子一叫.
就立刻醒來了.
因為有那個小孩子的喊叫聲.
或者啊一聲,啊一聲,哎呀.
於是才吵醒了其他人.
才會夜間都起來的.
哎呀,死了,大兒子.
於是大埃號.
是無一家不死一個人.
所以這次第一點就是.
都半夜是無一家不死一個人.
當然這個是概括地說.
那正如是不是所有家都有人死的呢.
那不一定的.
不過極少數沒有人死了.
譬如舉個例子.
是兩夫婦結婚.
那個兩個都不是長子.
都不是頭生.
又沒有小孩子.
又沒有生口.
就不死囉.
但是這樣是概括地說.
無一家不死一個人.
嘩,這個災大不大啊.
殺長子之災,好大.
不過你要知道是.
法老和一切神僕並愛及眾人.
都有這樣的事情遭遇.
那包不包括以色列人呢.
是不包括的.
因為以色列人就有第一個逾越節.
就是將血擦在門眉門框上面.

$^{121}$於是就說滅命天使經過的時候.
其實這個滅命天使.
是不是一種天使呢.
還是說有時候舊約是將天使.
神耶和華是用天使來表達自己的.
就是神將他們殺死了.
總之這個滅命的呢.
就經過見到門眉門框上.
有一個血就不進去殺那一家的長子.
如月就是過世,跳過月過這間屋.
不進去.
那就概括地無一個人.
是家裡面不死一個人的.
好了,然後再看下去.
三十一節夜間就這樣.
嘩,忽然之間睡一覺.
可能就照顧著這個.
法老長子的婢女大聲喊.
咦,為何公子大聲喊呢?就去看看.
啊,然後死了.
嘩,而且死狀都可能有一點恐怖.
於是馬上大聲喊.
哎呀,救命呀,救命呀.
就喊了那些御醫,警衛.
接著就趕快通報皇上.
不得了啦,不得了啦,皇上.
於是就吵醒了.
法老就馬上見到這件事.
這件事非同小可.
他就明白.
啊,這件事是摩西的上帝做的.
於是法老召了摩西亞倫來說.
起來,連你們大意識的人.
從我門中出去.
就依照你們所說的.
虛視鳳耶和華罷.
所以第二點給大家看的.
就是虛視鳳耶和華罷.
在什麼時候發生呢?.
夜間其實是在說就是那一晚.

$^{161}$就是那一晚.
就半夜的時候.
哎呀,那些長子都死了.
就是那一晚,法老馬上急召摩西亞倫來.
跟我馬上叫摩西亞倫來.
他們在哪裡.
然後就來了.
一來到就說.
你好東西,這句我是作的.
你好東西,你的神厲害.
竟然殺了我所有的長子.
你好東西.
那他呢?.
你發覺到在法律來講.
如果你是法律,你會怎樣?.
叫了摩西亞倫來.
因你.
跟我馬上拿他出去.
剁了,剁成肉醬.
蒸了他.
豈有此理.
讓我兒子也死了.
我們以色列人死了這麼多長子.
這兩個罪無可赦.
哇,殺死他.
千刀萬割.
他不是,他是.
起身.
帶你的以色列人走.
依照你們所說的去事奉耶和華吧.
就是說呢.
其實法律是知道這件事是出於耶和華.
這兩個人只不過是他的使者而已.
那他在這裡呢.
他是知道.
啊,這次鬥不過耶和華了.
算了,起身,走啦.
那我要講一下這個起來.
come,come這個字呢.
是起身站立.

$^{201}$意思就是說.
不要等了.
就好像人家在睡覺的時候.
起身.
起身,起身,要馬上走.
那這個起來呢.
就是馬上起身.
叫齊你的以色列人起身.
馬上跟我走.
帶著你的以色列人.
睡覺都要叫醒他.
馬上跟我以色列人.
在我民宗出去.
就照你所說的.
去拜你那個耶和華事奉他啦.
那他講的那些字都挺好的.
帶著以色列人.
在我民宗出去.
照你所說的.
事奉耶和華.
如果你發怒應該怎樣.
你兩個害得我們慘了.
我真是很想殺你.
不過算了.
你這個神這麼厲害.
帶著那些以色列人走啦.
快點走.
你出去拜什麼拜什麼.
要拜你的.
你快點去吧.
不要陀衰我們.
這樣講的嘛.
他是說你留意一下.
起身.
不要再睡了.
帶著那些以色列人.
半夜你跟我馬上走.
你去你們事奉.
去事奉你們所說的.
耶和華啦.

$^{241}$這樣喔.
他提這個耶和華的名字.
厲害.
這個名字算厲害.
他都知道厲害了.
那其實呢.
你要知道在這裡.
發怒是開始九個災.
甚至去到前一個.
第九個災.
黑暗之災.
而全部黑到.
那些孝子黑到迫過來.
這麼黑髮.
伸手不見五指那種黑髮.
三天不能夠出去.
但是以色列人住的地方.
是有光的.
這麼厲害.
哇.
那他都開始.
已經有點猶疑了.
不過呢.
他仍然心很硬.
都是不讓他們走.
於是就說.
好啊.
第十個災要來了.
那第十個災呢.
他是知道的了.
哇.
竟然這樣.
好吧.
你馬上去事奉耶和華.
就是說.
他開始連這個外邦的法老王.
都認識耶和華了.
這次.
哇.
都有些尊敬的心了.

$^{281}$好.
為什麼這麼說呢.
就跟大家去看看吧.
再看下去.
第二點.
第二點就是去事奉耶和華.
所以其實.
如果你有機會去事奉耶和華.
不要推遲.
不要像埃及人那樣.
牢籠住.
以色列人四百三十年.
想自由競拜.
沒有.
現在你有機會.
就去事奉耶和華.
神是會扭轉局勢.
令到這個法老王.
讓他們去競拜耶和華.
而且都多少敬畏.
哇.
這個神很厲害.
懂得耶和華厲害.
如果我們這些認識神的.
都不懂得神的厲害.
又有自由.
都不去敬拜祂.
不去事奉祂.
那你說.
多差呢.
連法律都說.
你去事奉耶和華吧.
這個神真的厲害了.
他都懂嗎.
懂的.
再看下去.
第三十二節.
也依你們所說的.
羊群牛群帶走.
要多少牛羊都帶走呢.

$^{321}$留意.
第三點.
並要為我祝福.
黑色字.
或者藍色字.
或者紅色字.
才是經文和合本的本身.
綠色字是我的解釋.
小一點點.
你們走吧.
羊群牛群帶走.
現在不要等了.
立刻去事奉耶和華吧.
這樣吧.
立刻去拜那個吧.
這樣就可以了.
但他現在說.
你去事奉耶和華.
懂的.
好了.
留意.
依照你所說的.
牛群牛群帶走.
並要為我祝福.
很多時候.
我就聽到很多解釋.
並要為我祝福.
就是.
你帶走.
不過記得為我祈禱.
你帶走.
不過為我祈禱.
其實不是.
並要.
加上.
是also.
相同地.
為我.
F.
F這個在文法上叫.

$^{361}$accusative.
就是如果中文叫賓格.
或者叫受格.
通常不翻譯的.
所以.
如果應該是F.
是沒有的.
並要祝福.
那你走吧.
帶走那些牛群牛群.
走吧.
並要祝福.
但是由於這個F.
其實只是不翻譯.
是並要祝福.
我是受格.
我是受那個人.
所以沒有這個F.
通常就不翻譯了.
中文就翻譯翻譯了.
並要祝福.
那你不明白.
並要為我祝福.
就叫摩西你帶他們走.
但是同樣.
並要.
同樣要為我祝福.
不是很明白.
其實要明白的地方在這裡.
在祝福這個字那裡.
要大家留意一下.
barack.
barack呢.
它的字源是跪下.
to kneel.
或者是bless.
現在翻譯了.
祝福就是翻譯了.
它的bless.
將這個barack.

$^{401}$就翻譯了做bless.
祝福.
其實你會見到英文.
叫中文這裡就不准了.
我們有時候praise the Lord.
有時候我們bless the Lord.
bless 祝福.
那祝福.
你祝福主.
你是哪位.
你比祂大.
說祂祝福.
不是.
這個bless.
是我們是向主.
頌讚敬拜的意思.
如果在對主來說.
主對我們bless.
may God bless you.
那就願上主祝福你.
oh our Lord we bless you.
we bless your holy name.
我們是讚美你的名.
我們是尊崇你的名.
這個是敬拜的意思.
而且barack.
其實它的字源是跪下.
所以是kneel.
主要我們kneel down.
在你面前.
還有我們bless you.
我們在敬拜讚美你.
所以如果我將它譯.
就是帶你的牛群羊群走.
走的同時.
請你見到耶和華的時候.
為我跪下.
即是什麼意思.
帶你們走.
馬上走.

$^{441}$去事奉耶和華.
牛羊就帶走.
不過同樣地.
走的時候.
如果當你見到耶和華的時候.
請你代替我跪下.
跪下即是什麼意思.
就好像現在說.
當你見到耶和華.
你的神的時候.
事奉祂的時候.
你代我擺一擺.
明白的意思就是.
你都給我擺一擺祂.
我都真是求求祂.
保護我們.
擺一擺你.
你帶你帶.
其實在這裡.
法祿是開始這一節和上一節的語法用語.
是表示法祿去到第十章.
十章子之章.
他開始明白.
嘩.
這個耶和華是真的.
廣東話.
香港俗語叫做.
嘩.
真的.
有料到.
和我們拜開那些神.
完全沒得揮.
嘩.
我們那些全部都是沒用的.
連我們擺的太陽.
它都可以關我們燈三日.
而在以色列那裡有光.
而它是會在沒有血.
在門眉門腔的.
所有長子都死.

$^{481}$連生畜都死.
這麼厲害.
誰知道誰是長子.
一群羊.
你怎麼知道誰是長子.
嘩.
很厲害.
孿生.
你怎麼知道誰大一點.
大幾分鐘.
大幾小時.
你怎麼知道.
不過.
嘩.
這個神真的厲害.
你們走吧.
去敬拜你的耶和華.
去事奉祂.
帶著那些牛羊走.
見到耶和華的時候.
同樣.
和我拜一拜.
和我對祂下跪.
我就不跟你去了.
你見到的時候.
你當我向祂下跪.
不過中文就譯做.
並要為我祝福.
這個祝福.
因為是巴拉克.
那個字源.
是可以是跪下.
這個跪下.
是New或者Bless.
這個Bless有時候.
回頭就會譯為祝福.
但其實.
Bless the Lord.
和Praise the Lord一樣.
頌讚主.

$^{521}$讚美祂.
不是祝福主.
怎麼祝福主.
是讚美主.
Bless the Lord.
所以.
是和我敬拜祂.
和我拜一拜.
法律在這裡.
是去到這樣才認識.
結束了.
不要去到.
有人死這麼交關.
不要去到這麼大件事.
才回轉.
不要去到這麼大件事.
這麼悲慘才回轉.
早一點就做一個.
信命的兒女.
早一點就聽從.
上帝的話.
你沒有這麼慘.
我自己人生裡面.
都經歷過不少.
都經歷過一些.
很慘的慘況.
如果早一點知道.
早一點回轉.
早一點懂得敬畏上帝就好了.
過去了.
現在幸好法律都會說.
跪一跪.
你和我拜拜.
都認識他.
不過不要去到這樣.
好了.
結束今天.
和大家重溫一次.
你會發覺到.
到了半夜.

$^{561}$第一個你不知道的時間.
上帝的手就會工作.
將頭身的刀殺了.
無一家不死一個人.
第二.
你有機會事奉耶和華.
去吧.
去事奉耶和華.
珍惜每一個自由.
服侍主的機會.
第三.
定要為我拜拜.
法律都開始認識主.
而我們不要去到這麼慘情.
才認識他.
才跪一跪.
拜拜.
一早就開心.
跟隨他就好了.
我們一起祈禱.
主幫助我們.
給我們歡喜樂意.
跟隨你.
不要去到一個很艱難.
不要去到一個.
事情都壞了的時候.
才認識你.
其實你是又真又活.
你一早就已經存活著.
所有事都是你掌管.
而你會在你的時間.
不是我的時間.
是你的時間.
半夜.
可能你就工作.
白天你就工作.
做你的事.
主我們在這裡讚美你.
要跟隨你.
給我們一個順服的心.

$^{601}$不要像法律一樣.
要去到最後.
遭遇很大的災難.
才懂得回轉.
求主祝福.
誰聽祈禱.
奉耶穌名求.
下門.
不過這裡要補充一點.
就是你發覺到.
後來法律.
是回心轉意.
想想想.
死這麼多人.
氣不順.
給我追他.
追這幫百姓.
帶他們回來.
帶不到.
就全部死一項東西.
給我殺光.
嘩.
很厲害.
不過上帝又救他.
終於法律這個第二次.
自己跪下了.
我們下次再看.
今天跟大家回到這裡.
明天我們再見.
拜拜.
\newpage



\section{使徒行傳 12:1-25}
\label{sec:031932}
\textbf{越獄風雲 (1)}~~~~林鍵牧師
\newline
\newline
連結: \href{http://www.church.com.hk/acms/content.asp?site=cdc\&op=show\&type=product\&code=031932\&layout=sermon}{\texttt{031932}} ~~~~ 語音日期: 2026-04-04
\newline
\newline
\hyperref[sec:031931]{\small{< < < PREV SERMON < < <}}
~
\hyperref[sec:index_chronic]{\small{[返順時目]}}
~
\hyperref[sec:031933]{\small{> > > NEXT SERMON > > >}}
\newline
\newline
使徒行傳 12:1-25
\newline
\begin{longtable}{cl}
\hline
\hline
章節 & 經文 (和合本修訂版)\\
\hline
12:1 & \begin{tabularx}{0.7\textwidth}{X} 約在那時候，希律王下手苦待教會中的一些人， \end{tabularx} \\ \\ \relax
12:2 & \begin{tabularx}{0.7\textwidth}{X} 用刀殺了約翰的哥哥雅各。 \end{tabularx} \\ \\ \relax
12:3 & \begin{tabularx}{0.7\textwidth}{X} 他見猶太人喜歡這事，也去拿住彼得。那時候正是除酵節期間。 \end{tabularx} \\ \\ \relax
12:4 & \begin{tabularx}{0.7\textwidth}{X} 希律捉了彼得，押在監裡，交給四班士兵看守，每班四個人，企圖要在逾越節後把他提出來，當著百姓辦他。 \end{tabularx} \\ \\ \relax
12:5 & \begin{tabularx}{0.7\textwidth}{X} 於是彼得被囚在監裡，教會卻為他切切禱告神。 \end{tabularx} \\ \\ \relax
12:6 & \begin{tabularx}{0.7\textwidth}{X} 希律將要提他出來的前一夜，彼得被兩條鐵鏈鎖著，睡在兩個士兵當中；門前還有警衛看守。 \end{tabularx} \\ \\ \relax
12:7 & \begin{tabularx}{0.7\textwidth}{X} 忽然，有主的一個使者顯現，牢房裡有光照耀；天使拍彼得的肋旁，叫醒了他，說：「快起來！」鐵鏈就從他手上脫落下來。 \end{tabularx} \\ \\ \relax
12:8 & \begin{tabularx}{0.7\textwidth}{X} 天使對他說：「束上腰帶，穿上鞋子。」他就照著做了。天使又對他說：「披上外衣，跟我來。」 \end{tabularx} \\ \\ \relax
12:9 & \begin{tabularx}{0.7\textwidth}{X} 彼得就出來跟著他走，不知道天使所做是真的，以為見了異象。 \end{tabularx} \\ \\ \relax
12:10 & \begin{tabularx}{0.7\textwidth}{X} 他們經過了第一層和第二層監牢，就來到往城內的鐵門，那門就自動給他們開了。他們出來，走過一條街，忽然天使離開他去了。 \end{tabularx} \\ \\ \relax
12:11 & \begin{tabularx}{0.7\textwidth}{X} 彼得清醒過來，說：「現在我真知道主差遣他的使者，救我脫離希律的手，和猶太人所期待的一切。」 \end{tabularx} \\ \\ \relax
12:12 & \begin{tabularx}{0.7\textwidth}{X} 他明白了，就到那稱為馬可的約翰的母親馬利亞家去，在那裡已有好些人聚集禱告。 \end{tabularx} \\ \\ \relax
12:13 & \begin{tabularx}{0.7\textwidth}{X} 彼得敲外門時，有一個使女，名叫羅大，出來應門， \end{tabularx} \\ \\ \relax
12:14 & \begin{tabularx}{0.7\textwidth}{X} 認出是彼得的聲音，歡喜得顧不了開門，就跑進去報信，說彼得站在門外。 \end{tabularx} \\ \\ \relax
12:15 & \begin{tabularx}{0.7\textwidth}{X} 他們對她說：「你瘋了！」使女堅持真有其事。他們說：「那是他的天使。」 \end{tabularx} \\ \\ \relax
12:16 & \begin{tabularx}{0.7\textwidth}{X} 彼得不停地敲門；他們開了門，一見是他，就很驚奇。 \end{tabularx} \\ \\ \relax
12:17 & \begin{tabularx}{0.7\textwidth}{X} 彼得做個手勢，要他們不作聲，就告訴他們主怎樣領他出監；又說：「你們要把這些事告訴雅各和眾弟兄。」然後，他離開往別處去了。 \end{tabularx} \\ \\ \relax
12:18 & \begin{tabularx}{0.7\textwidth}{X} 到了天亮，士兵中起了不少騷動，不知道彼得到哪裡去了。 \end{tabularx} \\ \\ \relax
12:19 & \begin{tabularx}{0.7\textwidth}{X} 希律找他，找不著，就審問警衛，下令帶走他們處死。後來希律離開猶太，下凱撒利亞去，住在那裡。 \end{tabularx} \\ \\ \relax
12:20 & \begin{tabularx}{0.7\textwidth}{X} 希律向推羅和西頓的人發怒。他們那一帶地方是從王的土地供應糧食的，因此就託了王的內侍大臣伯拉斯都的情，一心來求和。 \end{tabularx} \\ \\ \relax
12:21 & \begin{tabularx}{0.7\textwidth}{X} 希律在所定的日子，穿上朝服，坐在位上，對他們演講。 \end{tabularx} \\ \\ \relax
12:22 & \begin{tabularx}{0.7\textwidth}{X} 民眾一直喊著：「這是神明的聲音，不是人的聲音。」 \end{tabularx} \\ \\ \relax
12:23 & \begin{tabularx}{0.7\textwidth}{X} 希律不歸榮耀給神，所以主的使者立刻擊打他，他被蟲咬，就斷了氣。 \end{tabularx} \\ \\ \relax
12:24 & \begin{tabularx}{0.7\textwidth}{X} 神的道日見興旺，越發廣傳。 \end{tabularx} \\ \\ \relax
12:25 & \begin{tabularx}{0.7\textwidth}{X} 巴拿巴和掃羅完成了供給的事，就回到耶路撒冷，帶著稱為馬可的約翰同去。 \end{tabularx} \\ \\
[1ex]
\hline
\hline
\end{longtable}
$^{1}$歡迎大家來到我們的「週六崇拜」.
為我們上演了一場非常精彩的月玉風雲.
這個故事是出於《笑逸行轉》的第十二章.
我們問通過彼得月玉救贖的這個故事.
上帝到底要我們收到什麼訊息呢?.
如果你有機會上YouTube.
然後尋找有關《笑逸行轉》第十二章的.
《查經》或者《港道》的時候.
你會發現十之八九其實都是跟懇切禱告有關係的.
其中有一個視頻我看完之後特別歡喜的.
它的主題是.
《齊心禱告就能打開哪座監牢》.
它說親愛的朋友.
眾聖徒的禱告上帝聽見了.
祂以超自然的方式將彼得從監牢裡釋放出來.
接著祂問.
祂說現在在你身邊.
也有家人朋友正在困在艱難的處境裡面嗎?.
鼓勵你可以請教會的弟兄姊妹們.
就不要來Crossway了.
找其他吧.
陪你一起去為他這個很艱難的處境一起祈禱.
接著我們相信.
當你們同聲開口同心信靠.
神就必成全祂的拯救.
接著最後跟你們說.
親愛的.
按著祈禱就可以了.
親愛的主耶穌.
求你幫助我們的家人朋友.
一定要指名的.
接著求你憐憫他.
拯救他脫離艱難的困境.
禱告奉耶穌基督的聖名祈求.
阿們.
靈修小品將很多人的心聲都說出來了.
每個人所想所求.
都表達了在這個故事,這個禱告裡面.
然而我們必須要問.
這樣的理解,這樣的應用.

$^{41}$是不是可以的呢?.
是不是正確的呢?.
答案就是當然不可以.
當然一點都不正確.
這樣的演繹是完全忘記了整段聖經的開始.
有一個人叫雅各.
而他被人割了.
一開始的時候第一件事發生就是雅各的殉道.
雅各是約翰的哥哥.
是十二門徒當中第一個殉道的.
當時他大概都是三十多歲.
很可能已經成家立室.
有自己的妻兒.
死得很慘.
你怎麼可能用這樣的理解.
只要我們同心祈禱就搞定了.
如果你這樣去理解.
這樣去應用這段聖經的話.
那雅各就是死得很冤枉.
死得死不瞑目.
我們問我們究竟應該怎樣去明白.
《笑傲行轉》第十二章的故事.
死亡者告訴我們.
《路加福音》和《笑傲行轉》.
都是出於路加醫生所寫的.
他們其實是一部分為兩冊的合一著作.
如果你有機會去看《路加福音》的開始.
以及《笑傲行轉》的開始.
你會發現兩本書都是寫給同一個人的.
《路加福音》第一章第一節這麼說.
提亞非羅大人 這些事我詳細考察了.
也想按住次序寫給你.
讓你知道你所學的道都是確實的.
接著《笑傲行轉》第一章第一節.
提亞非羅大人.
沒有大人.
提亞非羅.
我在第一本書中.
這是第二冊.
議論到耶穌從開頭所做所教導的一切事.

$^{81}$直到耶穌被接上升的日子.
你會發現兩本書都是寫給提亞非羅.
提亞非羅的希臘文名字的意思.
是愛上帝的人.
你看英文叫Theophilus.
Theo就是上帝.
Philus就是愛上帝.
一個愛上帝的人.
我們相信他是一個愛邦的信徒.
也很有可能是一個羅馬的高官.
因為叫他做大人 提亞非羅大人.
於是當你這樣去理解《路加福音》的時候.
《路加福音》是第一冊.
《路加福音》所說的是.
耶穌開始所做所教導的一切.
直到祂上升的日子.
接著《笑傲行轉》是第二冊.
是什麼呢.
《笑傲行轉》就是耶穌繼續所做.
繼續所教導的一切事.
從祂被接上升的日子開始.
於是當你這樣去理解這本書.
其實是上下冊的話.
所說的一切都是耶穌的行轉.
因為耶穌不單止受死.
祂也復活.
所以到了第二冊的時候.
耶穌沒有死.
祂活在祂的仕途當中.
仕途行轉不是一個正確的題目.
是耶穌的行轉.
耶穌通過祂仕途所做的一切.
對於老家來說.
《笑傲行轉》的主題.
就是第八節.
《笑傲行轉》最主要的主旨.
是「但聖靈降臨在你們.
必得著力量.
在耶路撒冷.
由帶全地撒瑪利亞.

$^{121}$直到地極.
作我的見證」.
這裡所說的聖靈.
其實是耶穌自己的靈.
所以不是耶穌的表弟表哥.
就是耶穌的自己.
這個聖靈是通過.
復活升天的耶穌所差遣下來的.
祂自己升了天.
坐在上帝的右邊.
接著祂將自己的靈.
當然我們沒有辦法明白.
又是祂又不是祂.
祂將自己的聖靈降下.
於是通過祂的聖靈.
祂與祂的使徒一起同在.
從這個角度來說.
我們可以說.
《路加福音》的下冊《笑傲行轉》.
更好的題目應該叫做《夜靈行轉》.
《耶穌行轉》.
但是耶穌不是身體與祂的使徒一起.
祂是通過祂的靈.
於是變成《夜靈行轉》.
而不單止是所謂的使徒行傳.
所以今天所說的題目.
如果你將它張開一點來說.
就是《夜靈行轉之月玉風雲》.
上冊.
下星期回來就是下集.
《夜靈行轉之月玉風雲》.
綜合以上的重點.
我們可以將使徒行傳.
以及《路加福音》.
或者《路加福音》和《笑傲行轉》.
分為三大部份.
剛才我說完的東西.
你綜合它.
第一.
《路加福音》和《笑傲行轉》.

$^{161}$是關於耶穌基督的.
耶穌基督的事.
耶穌基督是上帝所應許的薇塞亞.
所謂基督就是希臘文.
薇塞亞就是希伯來文.
其實就是上帝在舊約聖經裡.
已經預言.
祂將來會有一個君王救主來到.
耶穌就是基督.
耶穌就是薇塞亞.
另外一個部份.
B部份是說到.
耶穌會將祂自己的聖靈.
祂自己的靈差遣到地上.
接著.
才成了第三部份.
《仕途行轉》.
你會發現這裡.
每一部份底下都有兩個空格.
如果你按照你所認識的.
《仕途行轉》和《路加福音》.
去填補這六個格的時候.
你會更加清楚明白.
《路加福音》《仕途行轉》的大綱和主題.
於是我們開始了.
第一部份A部份.
基督耶穌 基督耶穌做了什麼.
基督耶穌受死復活.
基督耶穌升天作皇.
接著.
由A到B.
就是因為祂受死復活.
因為祂升天作皇.
現在祂可以.
把祂的靈差遣到地上.
接著我們要問.
差遣到地上做什麼.
這個不是那麼簡單.
不是很多人能回答.
差遣聖靈到地上.

$^{201}$但是我們的弟兄姊妹應該很熟悉.
以至地上的死在最終的罪人.
可以經歷一個叫做以馬五師的經歷.
以馬五師就是耶穌復活的那一天.
有兩個傻瓜門徒.
竟然離開耶路撒冷.
不記得耶穌說了祂自己會復活.
而復活了的耶穌.
跟他們一起走以馬五師的路上.
跟著跟他們一起查經.
查了兩三個小時的經.
在那裡敲醒他們.
指出他們的罪.
跟著告訴他們自己.
整本的聖經告訴他們.
他一定要受苦進入榮耀.
於是在查經的過程裡面.
這兩個經歷了以馬五師經歷的門徒.
他們眼睛明亮.
心裡火熱.
這樣才會有仕途行轉.
如果他們沒有眼睛明亮.
心裡火熱.
他們是不會有仕途行轉的.
仕途行轉不是一個殭屍片.
喪屍片.
仕途行轉是一個復活的聯盟.
這些人眼睛明亮.
心裡火熱.
於是仕途行轉還有最後兩個空格.
他們要地的見證.
他們要宣揚上帝的天國.
從耶路撒冷開始.
跟著由大撒瑪利亞直到地極.
這是整本仕途行轉的方向.
朝向地極.
從哪裡開始.
從耶路撒冷開始的.
但是他最終的方向.
最終的高潮是在地極.

$^{241}$所以我最後告訴他們.
仕途行轉可以分為兩半.
頭一半是第一章至第十二章.
第二半是十三章至二十八章.
二十八章是最後一章.
頭一半的主角.
所以什麼是主角.
所有主角都是耶穌.
所有主角都是上帝.
頭一半的主角就是主要的配角.
頭一半的主要配角是誰.
是彼得.
頭一半裡面福音侍公是圍繞著耶路撒冷的教會.
圍繞著猶太的信徒.
到了後半章十三章至二十八章.
轉了主角.
轉了主要的配角.
主要的配角不是彼得而是蘇羅.
或者是到了第十三章.
蘇羅的名字改為保羅.
我們所知的保羅.
於是保羅接了棒.
成為主要的配角.
而福音侍公十三章開始是圍繞著.
耶路撒冷,猶太,撒馬尼亞以外.
外邦人的信徒.
可以看到這是明顯的一個分別.
於是在這裡你會發現.
小行轉第十二章尤其是非常重要的.
因為它的轉接點.
如果你有機會.
我們下個星期再回來看的時候.
你會發現小行轉第十二章是.
蘇羅信主以後第一次與彼得同框拍照.
我不知道你有沒有留意到.
那時候的蘇羅其實都在耶路撒冷.
整件事發生的時候.
彼得被人救出來的時候.
蘇羅在哪裡?.
蘇羅不是在安提厄.

$^{281}$他是在耶路撒冷.
於是被人拍了.
被記者拍到他們同框.
因為他在告訴我們.
整個故事開始進入一個膠棒的時刻.
接下來.
保羅就會將這幅音傳到地極.
信不信由你.
我甚至找到一個聖經學者.
他很仔細分析整本試圖行轉的二十八章聖經.
將二十八章聖經分為四十九部份.
Can you believe it?.
已經長到我沒辦法.
我一定要橫向給你看.
接著他將四十九部份.
編成為一個前後呼應的文學架構.
有二十四部份的經文是前後呼應的.
接著一定要單數的.
你明白嗎?.
一定要單數的.
他一定要四十九.
因為你有二十四部份.
接著中間.
中間就是整本書的中心.
整本書的中心是什麼?.
居然就是試圖行轉第十二章.
哇!小豪.
彼得的《月玉風雲》.
竟然是整本書裡面最重要.
最具重量的改變.
整本書的方向的一章.
就是我們這次要讀的一章.
你開始明白其實.
為什麼一篇到其實搞不定.
因為背後有太多的東西需要去考量.
於是我們再問我們怎樣去明白.
試圖行轉第十二章.
這個這麼重要的一章聖經.
我們先看看它的文學結構.
聖學者告訴我們.

$^{321}$試圖行轉第十二章可以被編排成為一個.
ABCDCBA前後呼應的架構.
其實你讀下去你就會明白的.
上面的A一開始的時候是雅各的死.
接著你發覺整個故事結尾的時候是希律的死.
那個AA都是跟友人死.
接著第二.
B是士兵守住彼得.
接著下面的B是士兵失去了彼得.
都是跟士兵有關的.
接著CC 彼得跟著天使.
接著另外一個C是彼得的天使.
即是人家以為彼得沒有回來.
是彼得的天使.
彼得的天使是彼得的天使.
兩個C都是有天使的出現.
接著 所有都不重要.
重要就是中間那塊東西.
中間的D.
主句在哪裡.
第十二章 第十一節.
「主救我脫離」.
讓我讀第十二章 第十一節給你聽.
「彼得清醒過來」.
記得弟弟其實演得很好.
他要抹眼.
「他清醒過來」.
弟弟清醒過來 彼得清醒過來說.
「現在我真知道主差遣他的使者」.
「救我脫離希律的手」.
「和猶太人所期待的一切」.
這個就是整個故事最主句.
就是主差遣他的使者.
「拯救我脫離希律的手」.
和猶太人所期待.
猶太人期待當然就是.
讓他死無全屍.
然後永不消生.
你明白我的意思吧.
不過你讀完這個大綱其實都不夠.

$^{361}$如果我們真的要明白.
《笑逐行傳》第十二章.
在整本《笑逐行傳》當中的位置和作用.
我們需要更加深入地去思考這個故事.
我想大家幫我用這個.
五個W 一個H 的分析工具.
這個又是幫到我們的.
「When」.
「Who」.
「Where」.
「What」.
「How」.
「Why」.
「何時」.
這個故事是何時發生的.
何人這個故事裡面有什麼人.
「Where」這個故事在什麼地方發生.
然後「What」這個故事發生了什麼事.
然後「How」.
這件事是怎樣發生的.
最後為何會有這個故事.
所以我們首先問一個很重要的問題.
「When」這件事是何時發生的.
如果你看那個影片.
你見到他影片上面.
這裡他打了出來是 Passover 的.
這個很重要.
顯然他講的時候沒有清晰講.
這個時候這件事是.
彼得被營救的這件事.
是發生在除敲節.
發生在逾越節.
發生在他將要被人削斷的前一夜.
這個故事裡面有什麼人.
我說有很多人.
我真的沒有辦法盡寫.
這個故事裡面有彼得.
有教會.
有上帝.
有天使.

$^{401}$有希律.
有士兵.
其實還有另外一幫人.
猶太人.
因為希律要做的是要討好猶太人.
這個另外一個.
不過沒有地方寫了.
這個故事在哪裡發生.
如果你當這個是一個舞台劇的時候.
你發現這個舞台劇需要有三個佈景的.
第一.
是一個死亡的監獄.
第二.
那裡鐵門很重要.
你出到那鐵門你就生.
你在鐵門之內就死.
生與死之間.
有一個鐵門接著出去了.
就去街道了.
接著還要繼續走.
就去到瑪利亞的家.
我們問何氏.
究竟整個故事裡面發生的是什麼事.
我用這句話總結整個故事.
主拯救彼得脫離希律的手.
以及猶太人的期望.
行不行.
這個你可以背過去.
就是剛剛第十二章.
第十一節.
接著How.
神其實是怎麼做這件事的呢.
有重要的.
主差遣天使.
天使很重要在故事裡面.
要叫醒彼得.
彼得睡了覺.
當然你可以想為何彼得睡了覺.
有人說他有平安.
即是翌日死了.

$^{441}$你也可以說他被人打得很累.
那沒辦法.
當然是睡覺.
天使做了一個動作叫醒彼得.
接著.
天使是將那個鐵門自動打開.
是這樣去救彼得的.
記住這些東西.
我想你冷靜一下.
停一停.
想一想.
你想想這五點所寫的東西.
我們將它連貫起來.
我想問問你.
你看到甚麼.
我再重複一些主句.
除敲節逾越節.
主拯救彼得脫離.
猶太人的期待.
主差遣天使.
叫醒彼得.
鐵門為他自動打開.
你看到甚麼.
你想起甚麼.
如果不是彼得.
如果是主拯救某人.
從猶太人的期待裡面.
拯救出來.
主差遣他的天使.
叫醒某人.
為他將那道門打開.
你會想起誰呢.
我想起的是耶穌.
整個故事橫看斷看.
就是耶穌受死 埋葬.
復活的影子.
天使幫他打開那道門.
當然那不是鐵門.
是石門.
你會發現很多東西.

$^{481}$和耶穌很有關係.
耶穌何時被埋葬.
受死 是逾越節 是除敲節.
於是你開始明白.
其實你這樣想一點都沒有錯.
因為使徒行傳不是使徒的行轉.
使徒行傳是耶穌的行轉.
一會有機會我想和你讀完經文.
讀完之後你會更加清楚.
更加肯定整個故事不是講彼得.
整個故事是講耶穌.
我們回到使徒行傳第十二章.
這個故事不單單是彼得個人的穴獄救贖.
其實你會發現故事更加廣大.
有外面的一層.
那一層是前面有雅各的死.
後面有希律的死.
所以我們讀經文的時候.
你不要只想著這個故事是關於彼得的拯救.
這個故事還要關雅各的死.
接著結束的時候有希律的死.
一會再回來.
第六我們要回答這個問題.
這個故事何故.
或者我可以這樣問.
整個故事是因為什麼事情而發生的.
接著最後結果又是什麼.
我的答案是這樣的.
整個故事是因為有人徹徹禱告去發生的.
是對的.
雖然靈修小品最後的結論是很差的.
但是他說得對是和祈禱有些關係的.
而最後的結局是什麼呢.
主導興旺.
你不要說最後的結局是彼得的救贖.
最後的結局是主導興旺.
越佛廣傳.
所以我想你記住.
這個是由禱告開始.
結果是主導興旺越佛廣傳作為結束.

$^{521}$我想用一幅圖畫給大家做一個小結.
這幅圖畫表面上我們看到的.
是彼得個人的穴獄救贖.
但是其實後面還有另外一個人.
他其實不是說彼得.
他是說耶穌基督受死埋葬.
復活出墳的影子.
其實他在說一個更偉大的東西.
就是你記不記得當年耶穌如何被埋葬.
去了一個死亡的監獄裡面.
又有冰釘守住.
又有塊石頭擋住.
接著上帝派遣了祂的天使.
將祂拯救了出來.
所有的一切都是聖靈的工作.
聖靈的工作最後有一個目標.
那個目標是什麼?.
使主導興旺越佛廣傳.
地極見證宣揚天國.
所以為什麼呢?.
這個故事裡面最後保羅就會出鏡.
因為結果你到了第十三章的時候.
保羅就要地極見證宣揚天國.
不如下個星期回來.
我們才有機會再看多一點.
好.
你大概有一個想法了.
我想跟你一起去看看今天的經文.
接著看完那二十多節.
我們最後再做一些總結.
我們先從《市場行轉》第十二章一至四節開始.
你會發現左邊我給你的是.
彼得,月玉救贖的故事情節.
但是右邊右邊右邊是.
耶穌受死復活的平衡故事.
這完全是看你多熟悉你的聖經.
你很熟悉的時候你會發現.
這個故事左邊右邊很平衡.
我們先看左邊.
若在那時候希律王下手虎袋教會中的一些人.

$^{561}$用刀殺了約翰的哥哥雅各.
希律王下手虎袋教會.
我們有句英文說話叫.
Like father like son 有其父必有其子.
這個所謂的希律王不是他的兒子.
是他的孫子.
他有一個爺爺很出名.
他的爺爺叫做大希律Herod the Great.
他的爺爺最出名是追殺耶穌.
耶穌一歲多兩歲的時候.
你記得在百里行的時候.
馬太保欽第二章.
最後希律被人欺騙了.
被那些智者欺騙完之後.
他殺了所有兩歲以下的嬰兒.
耶穌就走了.
這個就是大希律.
現在這個希律是他的孫子.
他的孫子繼承了祖業.
繼續輔導耶穌在地上的身體.
那位爺爺迫發.
耶穌剛出生沒多久.
現在孫子迫發耶穌在地上的身體.
他當中做了什麼呢.
他斬首他用刀殺了約翰的哥哥瓦割.
所以什麼叫用刀呢.
我自己一用刀就這樣刺過去.
原來不是.
所謂用刀就是斬首.
這個又很重要的.
如果我問你有哪個希律.
就斬首了另外一個人.
你有沒有印象呢.
那個就不是同一個希律.
是另外一個希律.
不過有另外一個希律.
你看右邊的謎底.
有一個希律.
那個希律不是叔叔.
那個希律是約翰.

$^{601}$這個約翰是誰的約翰.
施洗約翰.
我已經斬了.
就是用刀殺的.
你開始明白了.
瓦割等同施洗約翰.
然後彼得就等同耶穌基督.
你明白嗎.
你的影子是漂亮到這樣出汁的.
於是在這裡.
瓦割 約翰的哥哥.
被人割了.
其實盛傳了.
耶穌給他的一個應許.
或者預言.
你記得嗎.
在《馬可普欽》第十章裡面.
耶穌快要去到耶穌撒冷的時候.
瓦割和約翰來找他.
他說在你榮耀裡.
賜我們一個坐在你的左邊.
一個坐在你的右邊.
當然他們完全不知道.
耶穌要上十字架.
兩個傻子.
接著耶穌說了一句話.
耶穌說我所喝的杯.
你都要喝.
我所受的洗.
你都要受.
耶穌不是處罰他.
耶穌知道他後來會成為一個真正的基督徒.
他們願意背起十字架去跟隨耶穌.
耶穌告訴他.
你現在不明白你在說什麼.
你不知道你求的是什麼.
但是有一天你會明白.
還有我預了你會為我的緣故.
為福音的緣故受苦進入榮耀.
我答應你一定有你的份.

$^{641}$你要榮耀你一定有你的份.
不過你的榮耀不是你現在要的那個榮耀.
你的榮耀是十字架的榮耀.
我們知道最後.
瓦割和約翰這兩兄弟.
兩個都是為耶穌為福音的緣故.
受苦進入榮耀.
一個是十二門徒中的第一個被人殺.
另外一個是最長命的那個.
沒有被人殺.
不過九十多歲.
死在拔磨島上.
兩個人前後包抄.
中間有十個門徒.
全部都是被人相信的.
真棒.
瓦割是第一個死.
約翰是最後一個死.
但是沒有分別.
兩個都是為了耶穌的榮耀.
祂的天國受苦進入祂的榮耀.
所以鋸到瓦割第二節.
他看見猶太人歡喜者時.
也去拿著捉了彼得.
那時正是徐孝節期間.
希律王捉了彼得.
押在監裡交給四班士兵看守.
每班四人四四一十六.
但其實他們是每三個鐘頭轉一次監.
守夜.
每次都是四個人而已.
企圖要在逾越節時把他提出來.
當著百姓扮他.
扮他就是要殺死他.
他們要做的是猶太人歡喜的事.
就在徐孝節在逾越節的時候.
將他殺死.
這又是耶穌受難的影子.
因為耶穌在什麼時候受死.
就是在徐孝節在逾越節的時候.

$^{681}$於是你安靜下來.
你開始明白整個故事.
左邊講的瓦割和彼得.
右邊講的是施洗約翰.
和後來的耶穌.
我們開始讀右邊的經文.
希律王說施洗約翰已經斬了.
接著徐孝節又叫逾越節.
斬了.
祭司長和文士再想辦法怎樣殺害耶穌.
他們大聲催逼彼拉多.
彼拉多不想殺耶穌.
最後的原因是因為猶太人歡喜.
要求他將耶穌釘十字架.
他們的聖蔭終於得勝.
於是彼拉多照他們的要求去定案.
不是因為彼拉多想他.
是因為彼拉多要討好猶太人.
耶穌死了.
第23章 第53節 盧迦布音.
他們將耶穌的身體取下來.
放在一個作孽而成的墳墓裡.
這就是耶穌的監牢.
大家明白嗎.
你不可以殺了彼得.
殺了彼得是不能復活的.
只可以讓他坐牢.
耶穌可以復活 沒事的.
現在耶穌死了之後.
他們將耶穌放在.
耶穌一個境界深厭的死亡監獄裡.
就是祂的墳墓.
我們繼續看.
第12章 第6節.
提撻出來的前一夜.
彼得被兩條鐵鏈鎖住.
睡在兩個士兵當中.
門前還有警衛看守.
我們記得耶穌的墳墓.
一樣有警衛看守.

$^{721}$忽然有主的一個使者顯現.
牢房裡有光照樣.
天使拍彼得 叫醒了他.
說「快起來」.
鐵鏈就從他手裡脫下.
天使對他說「束上腰帶」.
「穿上鞋子」.
他就照樣做了.
又對他說「披上外衣」.
「跟我來」.
我再說.
左邊是彼得穴獄救贖的故事情節.
右邊是耶穌受死復活的平衡故事.
他又說彼得是怎樣被人鎖住.
鎖住鐵鏈.
有兩個人在他身上.
門口還有其他人.
其實是在對比耶穌當時.
被人困在他的墳墓裡.
然後這裡很有趣.
他說天使來到.
天使說「快起來」.
「快起來」是甚麼.
「快點啦」.
但他發現你拍下去的時候.
天使有萬條絲鯉.
究竟他急還是不急.
「急啊 別人快點就醒來」.
「還是不急 慢慢來」.
「那些人不會醒來」.
「我下藥下得很重」.
「那些士兵不會醒來」.
天使好像很急.
但又叫他「束上腰帶」.
又要穿好鞋子.
然後還要穿上外衣.
如果我拍肚子戲.
我就穿上西裝.
穿得漂亮一點.
幫你打領個胎.

$^{761}$穿好頭髮才走.
究竟天使是急還是不急.
其實我覺得他一點都不急.
這是左邊的故事.
接著我們讀右邊的故事.
婦女帶著所預備的香料.
來到墳墓.
發現石頭已經被滾開了.
正如鐵閘被打開了.
他們就進去.
只是不見耶穌的身體.
正在這時困惑的時候.
發現有兩個人站著衣服發光.
所以他們去救彼得的時候.
爐房裡有光照耀.
因為這兩個就是光明使者.
衣服發光.
婦女就苦伏在他們面前.
兩個人對他們說.
為什麼在死人中找活人.
他不在這裡已經復活了.
好像彼得已經被救走了.
我們繼續看第十二章第九節.
彼得跟著天使走.
不知道是真的.
以為是見了異象.
直到第十節.
他們經過了第一層第二層的監牢.
來到往城裡的鐵門第三道.
過了就是自由了.
那門就自動為他們打開.
他們出來.
天使離開他.
彼得清醒過來說.
現在我知道主差遣他的使者.
救我脫離希律的手.
和猶太人所期待的一切.
我想和你看試圖行轉第二章.
試圖行轉裡面第二章.
是彼得自己如何去描述.

$^{801}$耶穌基督的壽死復活.
左邊 彼得說.
是上帝差遣他的使者.
救他脫離希律的手.
和猶太人所期待的一切.
右邊 彼得會告訴他們.
上帝是如何拯救耶穌的.
他既按照上帝確定的旨意.
和預知被交於人.
他的死都是上帝安排的.
你們就藉著不法之人的手.
將他釘在十字架上.
殺了 第二十三節.
第二十四節 第二章.
上帝卻將死的痛苦解除.
使他復活.
因為他原不能被死所扣禁.
於是第三十二節.
這耶穌 上帝已經使他復活了.
我們都是見證人.
上帝以立他為基督.
上帝拯救彼得.
其實是一個影子.
是因為他拯救了耶穌.
死亡無法拘禁耶穌.
正如這個死亡的監獄.
無法拘禁彼得一樣.
我們繼續看.
第十二章 第十二節.
他明白了.
就到那稱為.
瑪可的約翰的母親瑪利亞加許.
在那裡有好多人聚集的禱告.
第十三節.
彼得敲門的時候.
有一個使女名叫羅大.
出來應門.
你記得在右邊耶穌復活之後.
是有幾個女士們.
這裡只有一個羅大.

$^{841}$一個羅大是頂舉幾個女士們.
出來應門.
認出彼得的聲音.
歡喜 不顧 不開門.
跑進去報信.
那些人有什麼反應呢.
他們對她說你瘋了 不信.
使女堅持真有其事.
他們說那是她的天使.
這裡說有些人聚集禱告.
我們最後回來了.
我們要問這個問題.
究竟他們禱告了什麼.
其實經文沒有告訴我們.
他們禱告了什麼.
很多人都是猜測.
我就告訴你.
我猜他們禱告了什麼.
就不是大家想的那些東西.
當他們不相信羅大的報告.
你發現在右邊.
他們一樣不相信婦女的報告.
我們讀一讀.
路加福音第24章第6節.
他不在這裡已經復活了.
第9節.
於是他們從墳墓那裡.
回去把這一些事.
告訴十一的使徒.
其餘的人.
他們這些話使徒以為是胡言.
就不相信.
同一件事.
一個從監牢裡面出來.
另一個從死裡復活.
兩個都不相信.
第36節.
正說這話的時候.
耶穌親自站在他們的當中說.
願你們平安.

$^{881}$第37節.
他們卻驚慌害怕.
以為所見的是一個魂.
如果你看他英文.
It's a spirit.
他是一個靈魂.
你發現左邊的時候.
當羅大說彼得回來了.
他們說你瘋了.
然後說可能是他的天使.
有兩個可能.
有些人說.
猶太人相信每個人都有一個守護天使.
就不是彼得回來了.
是他守護天使來.
可能是哭喪.
回來說好慘啊彼得.
報告一下而已.
另一方面.
他們見到耶穌.
他們不相信耶穌.
他說以為是一個靈魂.
可能是耶穌的靈魂.
或者其他的靈魂.
這是什麼可能呢.
無論你左邊或者右邊.
你會發現.
他們都不相信是真實的身體.
聖經學者說.
你會發現這兩個故事裡面.
又多了一重的連貫.
一邊相信這是一個靈.
右邊相信耶穌都是一個靈.
他的身體沒有真的發現.
當然最後兩個人都被發現.
他們真的活生生的站在眾人面前.
無論如何.
這一切都是對應著福音書裡面.
耶穌復活的故事.
你看不看到.

$^{921}$全部都是和福音書有關係的.
左邊的天使對應右邊的靈魂.
因為兩個都是沒有身體.
我們繼續看第十六節.
彼得不停地敲門.
他門開了門.
一件事他就很驚奇.
彼得做過手勢.
要他們不作聲.
就告訴他們主使怎樣.
領他出監.
又說你們要把這些事.
告訴那國和眾弟兄.
到了天亮.
兵士中起了不少騷動.
不知道彼得在哪裡.
希律找他找不著.
就審問這些警衛.
下令將他們處死.
右邊你看看耶穌復活的情況.
耶穌說不要害怕.
你們要去告訴我的弟兄.
左邊又是你要去告訴某些人.
右邊又是你要告訴某些人.
叫他們往家裡來去.
在哪裡會見到我.
祭司長和長老聚集商議.
就拿許多的錢給士兵.
說你們要這樣說.
夜間他們睡覺的時候.
他的門徒來把他拿走了.
士兵收了錢就照所吩咐的去做.
左邊殺人滅口.
等他們滅聲.
右邊用錢去滅聲.
左邊告訴你.
不要讓人知道彼得逃獄.
右邊說他們.
不要讓他們知道.
耶穌復活.

$^{961}$最後第21至25節.
希律在鎖定的日子.
穿上他的潮服.
坐在位上.
對他們演講.
民眾一直喊著說.
這是神明的聲音.
不是人的聲音.
就是奉承一下他.
希律不歸榮耀給上帝.
所以主的使者立刻擊打他.
他被蟲咬就斷尿期.
我不跟你說背景.
究竟希律發生了什麼事.
有人說他是病了五天.
很慘才斃命.
不是當場死了.
但我想你看到一件事.
就是希律彷彿帶著一個.
撒旦魔鬼的影子.
他那種高傲.
你跟他說.
這不是人的聲音.
是上帝的聲音.
他是沾沾自喜.
他要取替上帝的位置.
而他跟撒旦魔鬼的平衡就是.
魔鬼送了耶穌上十字架.
但最後因為耶穌救贖成全.
又被上帝擊打受審判.
而希律想要殺彼得.
但殺不到彼得.
結果彼得復活了.
結果希律被上帝審判.
又死了.
我想你聽我讀歌歐西書第二章.
比較複雜一點.
我由第十三節開始讀.
一直讀到第十五節.
歌歐西書講到.

$^{1001}$撒旦或者天空掌權者的失敗.
因為耶穌的十字架.
你們從前在過犯和未受割禮的肉體中死了.
上帝卻赦免了你一切的過犯.
使你們與基督一同過來.
還沒講原因.
原因是十字架.
逃去你在律例上所寫的.
敵對我們 束縛我們的自據.
把他撤去釘在十字架上.
十字架拯救了我們.
十字架勝過了撒旦的詭計.
撒旦的詭計就是你和我一起和他一鍋熟.
一起落地獄.
但上帝通過十字架拯救了願意悔改的人.
於是撒旦完全的失敗.
完全的崩潰.
第十五節.
歌歐西書第二章.
基督既將一切的執政者.
掌權者.
這個講的是撒旦的魔鬼力量.
即是其他魔鬼的權勢解除了.
他們的權勢就是要我們和他們一起落地獄.
上帝通過十字架拯救了我們.
於是他將他們的權勢解除了.
就在凱旋的行列中.
將他們公開示眾.
將十字架跨城.
耶穌贏了.
贏了什麼? 贏了撒旦魔鬼的東西.
所以.
希律好像得勝了.
將彼得放在墳墓裡.
但是.
彼得竟然從死裡復活.
而彼得的復活最後.
是引發了希律被上帝審判.
希律完全的失敗.
最後那兩節.

$^{1041}$上帝的道日見興旺.
越發廣傳.
於是有巴拿巴和掃羅完成了他們的供給就從耶路撒冷回去.
這個最後的翻譯有一點點不同.
幾乎有聖經.
有人說他們回到耶路撒冷.
其實應該是從耶路撒冷回去的.
下個星期回來我們再看這件事.
你會發現.
上帝的道興旺越發廣傳.
對應了整個小行轉裡面的大主題.
就是第一章第八節.
但這聖靈會降臨在你們身上.
你們就必得著大能大力.
並在耶路撒冷.
由帶傳遞撒瑪利亞.
直到地極作我的見證.
而十三章開始.
保羅小小的出現.
結果會引發保羅被差遣出去.
將福音傳到地極.
我想用一個圖表.
太複雜了.
列舉了仕途彼得 月肉救贖.
和耶穌受死復活的多項平衡.
我們列舉了九項.
今次我們由耶穌的受死復活開始.
因為誰先的.
是耶穌受死復活先的.
耶穌的受死復活.
引發了仕途彼得的月肉救贖.
你發現太多多項的平衡.
耶穌在徐考節 逾越節的時候受死埋葬.
在徐考節 逾越節的時候.
彼得被押在監獄.
天使滾開墳墓的石頭.
天使打開監獄的鐵門.
婦人見證耶穌的復活.
羅大聽見彼得的聲音.
仕途不相信耶穌復活.

$^{1081}$眾人不信彼得的逃生.
門徒以為見到的是雲.
眾人以為是祂的天使.
不是真實的身體.
耶穌說要把這事告訴仕途.
彼得說要把這事告訴瓦割和其他的弟兄.
太多平衡了.
第七 士兵收錢復活受禁.
右邊 士兵被殺徒生受禁.
第八 神使耶穌復活作亡.
右邊 神使彼得脫離犀律.
或者跨性犀律.
最後 十字架跨性權勢解除.
十字架跨性了撒旦.
使祂的權勢全部被解除.
犀律王斷氣 主的道興亡.
哇 因為耶穌活著.
所以我們才有這套戲.
才有《邵行遠》第十二章.
因為耶穌活著.
地上的教會經歷祂的同在.
不單止經歷祂的同在.
地上的教會更是彰顯祂受死復活的樣式.
這個很重要.
祂的樣子是很相似的.
那個樣式就是耶穌的樣式.
我們回到耶穌.
被地上教會的認許.
祂必與祂們同在.
《邵行遠》第一章第八節這樣說.
「但聖靈降臨在你們,必得著力量,災耶路撒冷」.
「由大全地撒瑪利亞,直到地極作我的見證」.
《馬太福音》大使命中耶穌說.
「你們要去使萬民作我的門徒」.
「為他們施洗教導,看我必天天與你們同在」.
那種同在不是有空找我,召喚我.
我就會出現的那種.
那些同在是.
我怎麼說呢.
是祂沒有一刻與教會分開的.

$^{1121}$教會在地上就是與祂同在.
教會在地上就是活出祂的壽死復活的樣式.
我就必與你們同在.
因為教會領受了.
祂宣教全楊福音的使命.
在這個使命裡面.
耶穌是跟祂,教會分不開的.
我天天與你們同在.
直到世代的終結.
我的時候就差不多了.
因為下個星期才回來看下集.
但是我要結束上集的時候.
我想和你分享三藻的聖經.
用三藻的聖經來回答三條問題.
這三條問題就是.
一,上帝是否一定會拯救信徒脫離仇敵.
第二,遭受逼迫的教會應當如何禱告上帝.
第三,我們若徹徹禱告上帝.
祂將如何回應我們.
第一條問題.
上帝是否一定會拯救信徒脫離仇敵.
我的答案是.
祂一定會.
你要小心.
你聽聽我的解釋.
《少行俊》第十二章第一節.
希律王下手虎袋教會 用刀殺了約翰的哥哥瓦割.
主第十一節.
差遣他的使徒 他的天使.
救我彼得脫離希律和猶太人.
兩個不同的結局.
瓦割和彼得非常不同的遭遇.
令我想起什麼呢.
令我想起丹尼里三個朋友.
面對烈火搖的時候對尼布加利薩王的回應.
我不知道你記不記得那篇道.
可能一年多前的事.
《丹尼里書》第三章第十六節.
尼布加利薩.
因為地方有些我沒有打出來.

$^{1161}$我讀第三章第十六節給你聽.
尼布加利薩王啊 者是我們不必回答.
什麼叫不必回答.
那時尼布加利薩給他再多一次機會.
他很愛惜這些有才幹的年輕人.
他說我再來多一次 音樂起.
只要你這次苦服就沒事.
前面沒事.
他們說尼布加利薩啊.
者是你不必再回答.
我們不會再回答你.
意思是我們是不會敬拜你那個金像的.
接著第十七節.
他們說即使如此.
即使如此是什麼意思.
如果我們不敬拜.
我們就被鷹在烈火堯當中.
即使你要將我們鷹在烈火堯當中.
接著你看他下一句.
他說我們所事奉的上帝.
他能夠能將我們從烈火堯中救出來.
我不知道他會不會.
但我知道他能救.
接著最後還有一句.
最後那句才是重點.
不過王他必救我們 脫離你的手.
他會不會將我們從烈火堯中救出來.
我知道他可以.
他會不會我不知道.
但是無論如何.
他一定絕對必拯救我們 脫離你的手.
如果你回顧那篇道.
我告訴你.
他們背後想的兩種拯救方法.
第一個拯救方法就是.
通過烈火堯 直回天國.
這個是拯救.
他們不會落在尼布格雷薩手上.
因為是烈火戰車.
《舊約聖經》裡面也有.

$^{1201}$以利亞上天.
他們想著下去.
誰知道上帝做了的士.
Uber我回去.
通過烈火堯 直回天家.
另外一個方法是.
將他們從烈火堯中救出來.
繼續作見證.
直到回家的時候.
所以你會發現.
雅各和彼得.
正是這兩種拯救的現代版.
所謂現代版就是新約版.
雅各走運了.
抽中獎了.
我知道大家不是這樣想.
你覺得雅各抽了下下籤.
彼得抽了上下.
其實不是.
雅各中了頭獎.
彼得出事了.
彼得遇獎了.
雅各走運中了獎.
立刻回到天國.
彼得要繼續努力.
等待終極救贖.
還要等.
就是這樣的分別.
所以這是第一問.
上帝是否必然拯救.
答案是必然拯救.
第二個問題.
遭受逼迫的教會.
應該如何禱告上帝.
我想告訴你.
這就是第五節的問題.
於是彼得被囚在監獄.
教會卻為他徹徹禱告上帝.
徹徹禱告上帝.
問題來了.

$^{1241}$禱告什麼.
我告訴你.
絕大部分的牧師和學者.
其實我聽了這麼久.
聽不到有一個人.
不是說他們在祈禱的.
是彼得的釋放.
所有的牧師和學者都告訴你.
他們的徹徹禱告.
就是祈禱彼得的釋放.
說真的我真的不太同意.
我覺得他們的想法是.
將他們自己.
讀入故事裡面.
自己害怕死.
於是他們將自己讀入.
你說教會的祈禱什麼.
就好像我被捉了.
你覺得教會的祈禱.
當然是祈禱我得釋放.
但我認為他們應該.
將彼得放入故事裡面.
應該將初期教會.
放入故事裡面.
而不是將自己放入故事裡面.
你想想如果彼得.
有機會跟你說.
今天你要怎樣祈禱的時候.
彼得會教他們祈禱什麼.
在小行主第四章裡面.
彼得和約翰被人捉了.
接著受到.
猶太長官的威嚇.
不讓他們再傳耶穌.
接著去到第四章.
第23節.
兩人被人釋放之後.
他們就回去將祭司和長老.
對他們講的說話.
和弟兄姊妹分享.

$^{1281}$弟兄姊妹聽見了就同心合意.
可不可以叫徹徹祈禱.
可以的.
同心合意徹徹.
高聲向上帝說.
說些什麼.
下次早告訴我們.
不要被捉.
不是.
說主現在求你監察.
他們的威嚇使你僕人放膽.
講你的道.
大哥.
很不一樣的祈禱.
下次.
他們再威嚇我們的時候.
求你讓你僕人放膽.
講你的道.
你這樣祈禱上帝.
你猜上帝會不會回答你的祈禱.
當然會回答.
所以第31節就是.
他們討個亂之後.
聚會的地方震動.
他們被聖靈充滿.
放膽傳我立刻回你.
上帝的道.
他要什麼.
放膽傳講你的道.
他得到什麼.
兩節之後.
放膽傳講你的道.
所以我想問問你們.
你們猜.
他們徹徹禱告.
禱告什麼.
不要死.
還是放膽.
阿彼,阿彼得也好.
雅各也好.

$^{1321}$放膽宣講你的道.
這不只是第四章.
第五章如果你有機會.
回去看看的時候.
彼得和使徒又被人抓.
五章第四十節.
那些議會的人叫他們來.
議會的人叫使徒來.
最後就是.
打完之後放他們出去.
把他們打了.
又吩咐他們不可.
奉耶穌的名講道.
然後把他們釋放了.
都是同一樣.
不可以再講耶穌.
打完之後就最後一次機會.
第四十一節告訴他們.
他們歡歡喜喜地離開議會.
因為他們算配為者名受辱.
你看看這個故事.
你現在再問.
教會徹徹禱告禱告什麼.
他們是不是.
麻煩我還有妻兒老少.
麻煩多放我出來.
還是剛才我所講的主啊.
現在求你監察他們的威嚇.
使你的僕人可以放膽全港來讀.
我橫看斷看.
我看不到他們會求上帝.
饒他一命再出來.
還有一樣東西.
又有聖經學者留意到一樣東西.
他這個徹徹祈禱這個字眼.
希臘文的字眼.
在路加福音裡面其實已經出現過了.
出現過在哪裡.
在耶穌的赫西瑪利園祈禱裡面.
耶穌這樣說.

$^{1361}$父啊求你將這杯斥退.
然而不照我的意願.
而是要成全你的意願.
耶穌痛苦焦慮.
禱告更加懇切.
耶穌的懇切禱告.
對應了教會的懇切禱告.
汗如大血點滴在地上.
耶穌赫西瑪利園的祈禱不是一個拯救命.
救命釋放的祈禱.
而是一個馴服至死的禱告.
所以我再問你.
當耶穌的身體在地上徹徹禱告的時候.
祂們禱告些什麼.
我想祂們禱告的是.
主啊求你監察.
他們的威嚇.
使你僕人放膽.
講你的道.
最後.
第三條問題.
我們如果徹徹禱告上帝.
祂將如何回應我們.
祂的回應是這樣.
巴拿巴語訴羅完成了事.
就從耶路撒冷回去安提柯教會.
在安提柯的教會裡面.
有事奉主和禁食的時候.
事奉主就是在敬拜的.
就是我們現在在做這些事情.
當然祂也有禁食.
祂在禱告的時候聖靈說.
要為我分派巴拿巴和訴羅.
去做我照祂們所做的事.
於是祂們禁食禱告後.
按手派遣祂們出去.
這個福音必定要傳到地極.
下個禮拜我們回來.
一起繼續思想月肉風雲的下集大結局.
我就是祈禱.

$^{1401}$我們天父我們多謝你.
這是你的話語.
絕對不是我們自己在想的很多事情.
求祝你幫助我們.
真的好好消化這一段.
很奇妙的故事.
跟我們有什麼關係.
我們應該怎樣去形容.
雖然我們仍然有一點點模糊.
雖然我們還沒到下集大結局.
但是我們已經開始感受到.
這段聖經那種大能大力.
第一祂告訴我們.
上帝必然拯救我們脫離我們的敵人.
第二祂告訴我們.
受逼迫的教會.
要祈求的是.
我們大能大力的宣廣你的話.
第三.
當我們切切禱告的時候.
你必踩遣.
這三點是不可思議的東西.
求祝你幫助我們.
緊握著你的真道.
去領受你的踩遣.
保守我們的心.
讓我們繼續思想你的話語.
直到下個星期我們回來看下集.
這樣祈禱奉主耶穌基督名求.
阿們.
\newpage



\section{約翰福音 20:1-18}
\label{sec:031933}
\textbf{基督今復活}~~~~高國雄牧師
\newline
\newline
連結: \href{http://www.church.com.hk/acms/content.asp?site=cdc\&op=show\&type=product\&code=031933\&layout=sermon}{\texttt{031933}} ~~~~ 語音日期: 2026-04-05
\newline
\newline
\hyperref[sec:031932]{\small{< < < PREV SERMON < < <}}
~
\hyperref[sec:index_chronic]{\small{[返順時目]}}
~
\hyperref[sec:031934]{\small{> > > NEXT SERMON > > >}}
\newline
\newline
約翰福音 20:1-18
\newline
\begin{longtable}{cl}
\hline
\hline
章節 & 經文 (和合本修訂版)\\
\hline
20:1 & \begin{tabularx}{0.7\textwidth}{X} 七日的第一日清早，天還黑的時候，抹大拉的馬利亞來到墳墓，看見石頭已從墳墓挪開了， \end{tabularx} \\ \\ \relax
20:2 & \begin{tabularx}{0.7\textwidth}{X} 就跑來見西門‧彼得和耶穌所愛的那個門徒，對他們說：「有人從墳墓裡把主移走了，我們不知道他們把他放在哪裡。」 \end{tabularx} \\ \\ \relax
20:3 & \begin{tabularx}{0.7\textwidth}{X} 彼得和那門徒就出來，往墳墓去。 \end{tabularx} \\ \\ \relax
20:4 & \begin{tabularx}{0.7\textwidth}{X} 兩個人同跑，那門徒比彼得跑得快，先到了墳墓， \end{tabularx} \\ \\ \relax
20:5 & \begin{tabularx}{0.7\textwidth}{X} 低頭往裡看，看見細麻布還放在那裡，只是沒有進去。 \end{tabularx} \\ \\ \relax
20:6 & \begin{tabularx}{0.7\textwidth}{X} 西門‧彼得隨後也到了，進了墳墓，看見細麻布放在那裡， \end{tabularx} \\ \\ \relax
20:7 & \begin{tabularx}{0.7\textwidth}{X} 又看見耶穌的裹頭巾沒有和細麻布放在一起，是另在一處捲著。 \end{tabularx} \\ \\ \relax
20:8 & \begin{tabularx}{0.7\textwidth}{X} 然後先到墳墓的那門徒也進去，他看見就信了。 \end{tabularx} \\ \\ \relax
20:9 & \begin{tabularx}{0.7\textwidth}{X} 他們還不明白聖經所說耶穌必須從死人中復活的意思。 \end{tabularx} \\ \\ \relax
20:10 & \begin{tabularx}{0.7\textwidth}{X} 於是兩個門徒回自己的住處去了。 \end{tabularx} \\ \\ \relax
20:11 & \begin{tabularx}{0.7\textwidth}{X} 馬利亞卻站在墳墓外面哭。她哭的時候，低頭往墳墓裡看， \end{tabularx} \\ \\ \relax
20:12 & \begin{tabularx}{0.7\textwidth}{X} 看見兩個天使穿著白衣，在安放耶穌身體的地方坐著，一個在頭，一個在腳。 \end{tabularx} \\ \\ \relax
20:13 & \begin{tabularx}{0.7\textwidth}{X} 天使對她說：「婦人，你為甚麼哭？」她對他們說：「因為有人把我主移走了，我不知道他們把他放在哪裡。」 \end{tabularx} \\ \\ \relax
20:14 & \begin{tabularx}{0.7\textwidth}{X} 說了這些話，她轉過身來，看見耶穌站在那裡，卻不知道他是耶穌。 \end{tabularx} \\ \\ \relax
20:15 & \begin{tabularx}{0.7\textwidth}{X} 耶穌問她：「婦人，你為甚麼哭？你找誰？」馬利亞以為他是看園子的，就對他說：「先生，若是你把他移了去，請告訴我，你把他放在哪裡，我去把他移回來。」 \end{tabularx} \\ \\ \relax
20:16 & \begin{tabularx}{0.7\textwidth}{X} 耶穌對她說：「馬利亞。」馬利亞轉過身來，用希伯來話對他說：「拉波尼！」（「拉波尼」就是老師的意思。） \end{tabularx} \\ \\ \relax
20:17 & \begin{tabularx}{0.7\textwidth}{X} 耶穌對她說：「不要拉住我，因為我還沒有升上去見我的父。你到我弟兄那裡去告訴他們，我要升上去見我的父，也是你們的父，見我的神，也是你們的神。」 \end{tabularx} \\ \\ \relax
20:18 & \begin{tabularx}{0.7\textwidth}{X} 抹大拉的馬利亞就向門徒報信：「我已經看見了主。」她又把主對她說的話告訴他們。 \end{tabularx} \\ \\
[1ex]
\hline
\hline
\end{longtable}
$^{1}$請姐妹平安.
請再一次同心祈禱.
Hallelujah 主已經復活了.
讓我們能夠去到你面前.
經歷你的豐富.
主的愛多謝你.
求你開我們的耳朵.
讓我們細心的聆聽.
開我們的眼睛.
讓我們可以見證到你的榮美.
更以我們的手腳行動.
去回應你的吩咐和大愛.
有緣無人口中的言語.
弟兄姊妹之間是心思意念.
蒙聖靈的啟迪蒙主的月立.
禱告奉耶穌基督的名教.
各位弟兄姊妹.
今天是復活的日子.
今天是基督復活的日子.
當大家剛才聽到這篇書.
一段在復活節裡面.
我們聽到的很熟悉的經文.
禮行坊裡面記載到耶穌就在那一天復活.
請你們一開始說七日的第一天.
清晨天還黑的時候.
胡達拉的瑪利亞就去到墳墓那裡.
一開始描述就是一個天陰陰.
甚至是天黑未光的日子.
可能也有點像今天.
陰陰天.
有人說今年比較特別.
因為復活節和清明是同一天.
所以可能是陰陰天.
但是無阻我們去經歷上帝光照我們的日子.
這個黑暗陰陰天的時候.
聖經裡面說.
這位胡達拉的瑪利亞就去到墳墓那裡.
首先我們要想想為什麼他要去墳墓那裡.
那天不是他們的清明.
應該我們今天才是清明.

$^{41}$才會去墳墓那裡去拜山.
為什麼胡達拉瑪利亞那天要去墳墓.
其實也不是無緣無故的.
根據聖經裡面描述耶穌在受難的日子是什麼時候.
是星期五.
所以剛剛我們的星期五是公眾假期.
弟兄姊妹有沒有留意.
星期五是公眾假期.
今年比較特別.
今年公眾假期有些重疊.
所以我相信很多弟兄姊妹可能出了外遊.
剛剛星期五就是受難的星期五.
我們相信當天就是耶穌基督受難的日子.
是星期五由早上九點開始一直釘十字架.
到了下午大概三點鐘 耶穌就斷氣了.
當耶穌斷氣之後 接著出現的問題.
我用這個字叫做問題.
因為當時猶太人他們的傳統日曆的習慣.
在星期五的黃昏六點鐘開始就是他們的安息日.
所以要數日子.
今天先大家數一數日子.
讓你們清醒一下.
猶太人的傳統日曆法跟我們有點不同.
我們是每天早上開始到晚上的.
中國人傳統日出就移作 日入就移息.
我們都認識.
猶太人的傳統日曆就是從晚上開始.
每天大概黃昏就是他們那天的開始.
一直到第二天的黃昏.
一個整個日子.
所以從創世記你都會發現.
上帝創造 晚上有早晨.
所以從晚上開始.
每一個星期到星期五的黃昏六點鐘就是安息日.
所以當耶穌釘十字架的時候.
就是安息日的前一天.
星期五的早上到下午.
所以當時他們很緊張.
究竟釘十字架怎麼處理呢.
因為如果耶穌在十字架上面未選氣.

$^{81}$又經過去到安息日怎麼辦呢.
因為整個安息日就沒有人理會了.
所以他們要在安息日之前.
要確認這位釘十字架的耶穌死了.
並且拿了屍體下來.
沒理由把屍體掛在那裡掛了這麼長時間.
所以聖經裡面就說.
當時有兵丁拿茅去刺耶穌的肋膀.
有些人不太明白為什麼要做這個動作.
是因為要證實耶穌真的死了.
因為大概看到耶穌斷氣之後.
他們就拿一根茅去刺耶穌的肋膀.
看到血和水分開.
有說血和水分開.
人死了之後.
身體結構的狀況.
血和水就會分開出來.
所以就看清楚耶穌是死了.
如果當時耶穌未斷氣未死怎麼辦呢.
他們就會敲斷.
就是當時掛在十字架上的罪犯的腳骨.
讓他沒有力氣去支撐身體.
就讓他快點斷氣而死.
這是當時的做法.
所以他們要檢查清楚耶穌是死了.
因為接著他們要很快將耶穌的屍體收下來.
因為不想留他在安息日的時候.
仍然掛在十字架上.
所以整個故事就是.
由星期五的黃昏之前.
耶穌斷氣.
接著黃昏的時候開始安息日.
屍首被取下來.
接著也開始做埋葬.
過了一個安息日的日子.
就是星期五的黃昏到星期六的整天.
到星期六的黃昏就是晚上了.
所以到星期日的早上.
差不多開始有陽光之前.
婦女.

$^{121}$特別是《聖經》裡面說到.
莫大拉的立馬利亞.
就來到墳墓那裡.
是很想為耶穌繼續做祂的身後事.
因為當時他們文化習慣.
傳統就會圍著屍體去果果屍布.
並且有香膏.
抹身體等等.
去為了做安葬的準備.
就是說很明顯.
當耶穌受死在十字架之後.
到祂下葬很短時間.
太倉促了.
還沒做完所有安葬的事情.
所以他們很想繼續去處理.
並且很多哀悼的事情還沒做完.
所以婦女們在主日的清晨.
就到了這個墳墓當中.
因此我們就知道.
耶穌在第三天復活.
大家記住.
不是耶穌三天後復活.
先聽清楚.
不是耶穌三天後.
就是不是二十四乘三.
七十二小時復活.
耶穌不是死後七十二小時復活.
耶穌是死後再第三天.
第一天就是安息日前的一天.
接著就是安息日.
接著就是安息日後的那一天.
這就叫做第一,第二,第三天.
如果計算時數.
四十個小時肚子.
可能是可以計算的.
星期五黃昏.
就是下午.
到星期日早上已經復活了.
所以其實以時數來說.
可能是四十四時附近的最多.

$^{161}$但是我們就叫做第三天.
這個很重要.
並且因為這樣.
你記住.
所以我們復活節一定是星期日.
試過我弟兄會問.
牧師今年復活節是星期幾呢?.
復活節一定是星期日.
為什麼我們今天在星期日崇拜.
猶太人傳統是守安息日的.
是不是?.
舊約十戒都說.
當守安息日.
其實我們不是很守安息日.
剛剛說安息日是星期五黃昏.
到星期六黃昏.
這是他們的傳統.
但是當耶穌復活之後.
我們將主日變成聖日.
所以我們在星期日的時候.
回來崇拜.
所以復活節一定是星期日.
還有有些弟兄會問.
為什麼年年復活節不同呢?.
今年是清明這個日子.
有時候是早一點.
有時候是遲一點.
因為我們計算的習慣.
傳統復活節是根據之前.
耶穌設立聖餐.
聖餐裡面說.
是跟月月節晚餐有關的.
而月月節猶太人計算的方法.
是用月圓有關係的.
所以又要退回星期日.
又要退回月圓.
就是月圓你也知道.
我們跟那些類似舊曆.
就是初一十五那些有關的.
所以又要跟他們猶太人.

$^{201}$傳統的復活節之前春分有關.
所以復活節比較複雜的.
每一年就是春分之後.
第一個月完結之後的星期日.
比較複雜的.
總之告訴你.
大概每年這個時候就是復活節.
一定是星期日.
也是在這段日子裡面.
過了猶太人的復活節.
整件事是這樣計算出來的.
有些背景有些少知識.
另外就是.
所以你就知道為什麼馬利亞.
當天要趕去墳墓.
因為很多事情都很想做的.
這個耶穌我們的老師.
過了新疆十字架.
德國很短時間要趕快安葬.
可能包了果絲布尼.
沒有情流就走了.
還有很多事情要做的.
因為對他們來說.
這個殯葬禮儀也很複雜.
所以七天的第一天.
就是主日.
就是今天這個復活的日子.
一清早天還沒亮.
還在黑.
快要亮的時候.
胡大拉的馬利亞就來到墳墓裡面.
來到墳墓.
來到墳墓他看到什麼呢?.
聖經裡面說來到墳墓.
他第一時間看到的就是.
看到不見了.
看到沒有東西.
看到一個absent.
就好像點名一樣.
老師來報到同學.

$^{241}$接著就看到什麼呢?.
看到absent.
看到不見了.
看到石頭挪開了.
看到身體不見了.
聖經裡面說原來在復活的時候.
第一個發現.
胡大拉的馬利亞發現就是.
沒有了.
不見了.
看到失去.
頂只會其實我們人生裡面.
有時候也是這樣.
我們人生裡面常常都會經歷這一種.
好像黑暗裡面的空洞一樣.
我們很想找很想找.
甚至黑黑暗暗的時候.
每一刻都要去找.
結果看到的就是失去.
無論面前是什麼狀況.
眼裡面看到的.
就是不見了的東西.
可能我們生命裡面.
有時候會不見了工作.
不見了身邊的人.
或者對社會現狀感到絕望的時候.
就可能好像看到一個空墳墓一樣.
有時候我們就是這樣.
在當中看到很多失去失落.
當我們看到失落的時候.
我們有什麼行動.
有什麼反應呢?.
聖經裡面說.
胡大拉的馬利亞.
第一時間的反應.
很典型的一個邏輯的想法.
一定是有人.
將主的屍體從墳墓裡面拿走了.
所以他就立刻走回去.
跟門徒說.

$^{281}$啊,不見了.
屍體被人拿走了.
我們很多時候用我們自己的理解.
去理解那些不見了的.
是不是?.
人生裡面有時候很多東西都不見了.
但是我們就會.
哦,是的.
那些不見了就是不好.
那些不見了就是失去.
失落.
我們很喜歡我們的理解.
甚至我們以為的邏輯.
是的.
因為馬利亞來到墳墓.
想做什麼呢?.
他就是想找耶穌基督的屍首.
一來到就是不見了的屍首.
我們往往.
將我們人生局限在這個空的墳墓裡面.
我們往往甚至以為.
上帝就存在在過去.
存在在過去的儀式生活.
一旦有很多震撼的時候.
我們覺得甚至上帝都好像被拎走了.
不見了.
我們常常都發現.
為什麼很多東西好像很失落.
連信仰.
連我們的上帝.
可能我都覺得.
上帝你去了哪裡?.
為什麼我不見你作為?.
我見不到.
我本來打算找這些東西.
我在這個地球上找公義公平的.
為什麼上帝你好像不見了?.
我在我生命裡面找成功成就的.
為什麼上帝可以讓這些東西不見了?.
當馬利亞看到這些不見了的時候.

$^{321}$他很緊張.
他覺得有人拿走了主的身體.
他就立刻跑回去找門徒.
結果聖經裡面很特別.
說有兩個門徒在跑步競賽奔跑.
很仔細地說.
這兩個門徒.
一個就是彼得.
另一個就是聖經裡面沒有說什麼名字.
那門徒.
傳說我們相信約翰寫的.
這個是約翰寫的自己.
所以他用這個角度去表達這件事.
當時就深信彼得和約翰在跑步競賽.
跑到這個墳墓裡面.
很有趣的.
裡面說得很仔細描述的.
大家去到誰先.
但是沒有進去.
又看看.
後來那個又進去了.
看完之後又出來.
只看到那些頭巾和細麻布.
就擺放在一處捲著.
就不像真的被人搬走了墓地.
突然間就好像耶穌穿過了那些細麻布和頭巾.
那些細麻布和頭巾仍然在那個位置.
這就是當時那兩個門徒所看到的事情.
接著看到之後.
跟著聖經裡面說.
他們看到.
那兩個門徒看到就相信了.
但是他們就不明白.
聖經裡面是這麼解釋的.
接著他們就走了.
這個描述有什麼意思呢?.
這個描述當然正在證明.
看到胸墳墓的不只是德馬利亞.
有門徒.
更有趣的地方就是告訴我們.

$^{361}$這兩個門徒很鉅細無遺.
說他們怎麼去到那個地方.
那件事情當然很實在.
另外其實可能也在反映我們的人生.
我們的生命的追求.
甚至我們的信仰的追求.
有時候也好像是這兩個門徒一樣.
就是跑來跑去.
跌跌碰碰.
有不同的節奏.
有不同的狀況.
有早一點去到.
但是又沒有走進去.
有晚一點去到.
但是又很緊張地走進去.
這個就是在反映我們的信仰的節奏和狀況.
更有趣的地方就是.
當他進去之後發現.
細麻布.
那頭巾.
放在那個位置.
不是很混亂.
不像盜墓.
也不是他們想像的屍體也沒有了.
但是也是很有秩序地.
繼續放好那個頭巾和那個屍布.
很有趣.
所以這邊更描述這兩個門徒.
看見就相信.
但是他不明白.
頂智媒信仰不是我們人生的競賽.
每一個人信仰的起步節奏都不同的.
有些人可能很早認識信仰.
走得好像很快.
有些人走了幾步之後可能停滯.
有些人信仰可能一生裡面很慢很慢地追求.
信仰有不同的節奏.
信仰不是要競爭.
信仰是要你的信心和你的行動去投入.
信仰是你要將信先行.

$^{401}$儘管你有事仍然不明白.
這個很值得我們心思.
有時候我們常常覺得.
我們要搞清楚,釐清所有事情.
我們才可以信.
但是信仰就是正正告訴我們.
你因為有些事情不明白.
你才要相信.
完全明白的事情就不是你投放的信心.
而你投放信心的時候可能你不明白.
不過這個投放的信心.
令你可以繼續有行動.
所以這兩個門徒就是投放了信心.
雖然他們仍然不太明白.
但是他們也有行動.
去生命裡面繼續奔跑.
上帝有祂的節奏.
上帝給每一個人成長都有祂的節奏.
所以頂智門你不要心急.
你有你的節奏.
你只有在你的節奏當中.
很堅定上帝的恩典就夠了.
我常常都有一個比喻.
可能這裡也說過.
我每一個人生命.
很多時候好像是扭一下螺絲一樣.
扭螺絲.
有人扭得快一點.
有人扭得慢一點.
有些螺絲帽子深一點有些淺一點.
最重要是你對方向.
繼續可以動.
你就會進去了.
有些人走得快一點.
有些人走得慢一點.
有些人好像穩固一點.
有些人好像緊張一點.
不過最重要是你對方向.
繼續向前.
無論有時候可能繞圈轉圈.

$^{441}$感覺上.
但是你繼續向前就可以了.
更重要聖經裡面說的.
這個果頭根告訴我們一樣很重要的東西.
上帝除了給我們每一個人有自己的節奏.
同樣上帝都有祂的秩序.
復活原來就是正正告訴我們.
只要回到一個合理的秩序.
這個果頭根擺到.
這個果絲帽擺到.
就告訴我們上帝有祂的計劃和祂的秩序.
好了.
很有趣的地方就是.
你會看到這些弟兄.
這兩個的門徒.
就走來走去.
走走走.
看看看.
(講解中).
這個都是不同人有不同的反應.
兩位的弟兄的門徒.
他的反應就是走來走去.
很快走來走去.
其實可能走來走去沒什麼用.
不過都是走來走去.
好像在做一些事情.
但是這位的姐妹.
她就繼續在那裡.
繼續在那裡.
她哭.
她很想用情緒感受.
在那裡表達她的傷痛.
跟著聖經裡面說的很特別.
有天使顯現.
有耶穌顯現.
她還誤會了她是圓丁.
她還跟耶穌說.
先生.
如果是你帶走她.
告訴我她在哪裡吧.

$^{481}$瑪麗亞認不出耶穌.
這是一個很嚴重的事情.
她這麼愛耶穌.
跟著耶穌幾年.
她很想為耶穌做最偉大的事.
甚至耶穌釘身十字架.
那些門徒都害怕得跑了.
不敢來.
她仍然冒著很多的危險.
在天還黑的時候.
就馬上到墳墓裡面去找.
為了為她所愛的耶穌做最後的事情.
但是她竟然認不出耶穌.
她竟然認不出.
這個認不出.
令我們好心的體會.
為什麼她認不出呢?.
當然了.
復活的耶穌可能有些樣子.
上面不同了就認不出.
這是其中一個可能.
但是突然間一個人出現.
而她是耶穌的話.
她認不出背後反映著什麼事情呢?.
反映著就是很明顯.
瑪麗亞當時在想的是不可能.
她沒有想像過這是耶穌.
她只是在想.
我要找耶穌的屍體.
耶穌會出現的在這裡.
不過是用屍體的形象出現的.
耶穌會出現的.
她應該就是這樣.
接著我就應該繼續為她做身後的事情.
我們有的時候用我們的邏輯.
用我們的思考.
用我們的理解.
而框了我們在信仰當中.
就算復活的主在你身邊.
顯然你都認不出.

$^{521}$有時候有些朋友跟我說.
穆斯我也想相信耶穌.
你告訴我.
如果耶穌突然在我身邊出現.
我就相信了.
我經常說.
如果真的出現就嚇到你暈倒.
你不用相信.
立刻直看見祂.
但是更重要的是.
其實耶穌經常在身邊顯現的.
很多天使.
很多神蹟騎士.
很多上帝的愛環繞保護你.
不過你不知道.
你不察覺.
你不認為是.
甚至你不相信.
上帝的神蹟騎士.
經常都在我們身上發生.
上帝的天使常常都在我們左右.
但是我們不察覺.
因為我們的期望是這樣的.
上帝你是這樣的.
我應該是這樣的.
上帝你是這樣的.
這個世界應該是這樣的.
我有我的期望.
我有我的想法.
我當然看不到.
我估計不到以外的事情.
復活就是告訴你.
一些你估計不到的事情.
因為沒有人經歷過.
沒有人明白.
而這個復活就是要戰勝.
一切我們所想到的框框.
所以瑪利亞無法認出耶穌.
但是耶穌沒有因為這樣.
去斥責祂.

$^{561}$甚至說你用不了我走吧.
耶穌是跟祂很親切地.
用一個祂最熟悉的方法.
讓祂認得到.
耶穌是呼喚祂.
瑪利亞.
《聖經》裡面說的特別描述.
瑪利亞就用一個回應.
這個叫做拉玻璃.
其實是當時的語言.
就是老師.
這個對話很溫馨.
就是耶穌喊.
好像平常一樣.
可能過去幾年就是.
偶爾耶穌就會喊瑪利亞.
瑪利亞就回應.
老師是一個很自然.
也是很親密的關係.
耶穌用一個祂最熟悉的方法.
去呼喚祂.
讓祂那一刻就體會.
是的.
這位就是我的老師.
這位就是耶穌.
不再是那位原主.
原丁先生.
管轅的先生.
而是我的主.
我的老師.
一句這麼溫馨的說話.
打破了那種.
當時瑪利亞仍然很心痛.
很蒙蔽看不到的眼睛.
那份悲傷.
原來我們的主認識我們每一個.
我記得楊方十章三節這麼說的.
我們的主就好像牧羊人一樣.
他按著名叫自己的羊.
將羊領出來.

$^{601}$你有沒有想過.
上帝會叫你的名字.
不是你做夢突然聽到有人叫你的名字.
不是那種.
而是上帝叫你的名字.
意思就是他認識你.
他知道你.
他和你有親密的關係.
他叫你的名字.
這個很重要.
耶穌在這裡用這個方法.
你想想整個復活的主.
復活的主第一句的說話是.
不是說很多很偉大的道理.
而是跟一個他愛的.
跟隨他幾年的瑪利亞.
跟他講一聲.
瑪利亞.
是一個呼喚.
是一個親密的呼喚.
是一個關係.
復活不是單單一件歷史的事件.
一個歷史的事實.
而是一個人關係的復活.
是人和上帝關係的復活.
人和自己關係的復活.
並且讓你有能力和身邊人關係的復活.
當你聽到主呼喚你的名字的時候.
你不再是那個很慘的.
好像瑪利亞一樣.
好像受害者一樣.
不見了耶穌身體很可憐.
而你就是見證者.
你就是見證著基督復活的那一位.
接著往後聖經裡面說.
耶穌和瑪利亞說.
不要拉著我.
祂要升天.
告訴我們.
當耶穌離開了這個地上的肉身.

$^{641}$當祂要復活的時候.
祂不再和往時的方法和門頭繼續相處.
往後祂也猜性靈在我們當中.
和我們一起.
在這段經文裡面.
我們看到一個很重要的事實.
這位瑪利亞.
這位穆迪拉的瑪利亞.
她由一位去哭泣的婦女轉化.
變成一位我們相信是描述為使徒中的使徒.
因為她在當中就真的看見.
她不僅看見.
她看見之後.
她還回應.
聖經裡面用很簡單的字眼說.
她看見耶穌.
她認到耶穌之後.
她就快快跑回去.
不是好像之前說的兩個門徒走來走去.
不知道做什麼.
而是她快快跑回去告訴其他門徒.
我已經看見了.
復活節的行動.
不是只要去討論空墳墓.
是怎樣的.
研究一下.
理解一下信仰.
不是想想你的說法.
是我的想法.
基督復活這件事件.
是讓我們生命裡面看到的.
是我們願意跟別人分享基督的大能.
是我們在限制裡面.
在難處裡面.
在破碎裡面得到醫治.
在絕望黑暗當中能夠得到光明.
各位朋友.
各位弟兄姊妹.
今天雖然這一刻好像仍然有點暗.
但是我們知道天已經亮了.

$^{681}$瑪麗亞所面對的黑暗已經不再黑暗.
我們可以怎麼回應呢.
第一樣我們可以回應的.
弟兄姊妹.
希望你可以回轉.
不要只看著你生命裡面的空墳墓.
看著那些遺憾看著那些不見了.
而是你轉身來.
去聽一下上帝要呼喚你的聲音.
回轉回到主的面前.
我們也可以整理一下我們自己.
甚至就正如聖經裡面看到的果頭根.
果絲寶一樣你可以整理一下.
復活代表著秩序性的混亂.
代表著我們可以找回我們跟上帝關係的節奏.
請你在你混亂的生活當中.
去重新建立跟上帝的關係.
重新找回你應有的秩序.
可能你跟身邊的人際.
可能你跟你的家人.
可能你的情緒都出現很多混亂難處.
放在正確的位置.
找回應有的秩序.
更找回你跟上帝應有的節奏.
以致能夠知道怎麼去繼續理解.
認識跟從相信這位主.
還有各位弟兄姊妹.
今天是一個很重要的日子.
瑪利亞成為一個很重要的榜樣.
她大膽地去宣告.
她一開始很害怕,很擔心,很傷痛.
但是當她經歷了主之後.
她有能力馬上走回去宣告.
當時環境沒有轉變.
仍然有很多人可能很害怕.
仍然有很多當時的兵丁.
可能在捉拿耶穌身邊的門徒,餘黨等等.
但是瑪利亞不懼怕.
她走回去大聲宣告.
我已經看見主了.

$^{721}$弟兄姊妹,我們是不是都願意去宣告呢?.
願意去分享呢?.
願意在這個黑暗的時候.
將光明帶來這個世界呢?.
基督已經復活了.
祂不在這裡了.
祂不再在這個空墳墓當中.
但是祂在我們的心裡面.
祂在我們的群體裡面.
祂在今天的地上,教會當中.
仍然作出呼籲.
以致我們能夠經歷祂復活的大能.
復活是我們每一個人都能夠經歷到的.
頂致祂會讓我們一起經歷祂的大能.
以致祂的差遣.
讓更多人和我們一起成為見證者.
體會生命的豐富.
復活力量的精彩.
請同心祈禱.
主耶穌,我們多謝你.
是你的恩典讓我們看見不一樣.
是你的恩典你親自向我們顯明.
讓我們知道如何去回應.
主耶穌給我們更多的信.
在我們未明白以前.
就讓我們投放更多的信.
在主你所賜的恩典當中.
在今天復活節當中.
就讓我們一起體會復活的大能.
去細心聆聽你向我們每一個人的呼籲.
你會按我們的名字去呼喊我們.
去提醒我們.
去跟我們表示你已跟我們建立親密的關係.
幫助我們脫離混亂黑暗.
幫助我們離開.
以致我們不留戀的失去.
而找到真實主你的實在.
求你成為我們生命的力量.
願你復活的大能在我們身上顯大.
以致我們生命活得不再一樣.

$^{761}$禱告奉耶穌基督的命求.
Amen.
\newpage



\section{}
\label{sec:031934}
\textbf{當黑暗籠罩世界:神在十字架上的回應 (粵/華)}~~~~張健庭牧師博士
\newline
\newline
連結: \href{http://www.church.com.hk/acms/content.asp?site=cdc\&op=show\&type=product\&code=031934\&layout=sermon}{\texttt{031934}} ~~~~ 語音日期: 2026-04-03
\newline
\newline
\hyperref[sec:031933]{\small{< < < PREV SERMON < < <}}
~
\hyperref[sec:index_chronic]{\small{[返順時目]}}
~
\hyperref[sec:031935]{\small{> > > NEXT SERMON > > >}}
\newline
\newline
$^{1}$大英之門平安.
今天我們一起來思想主義的基督為我們釘十字架.
今天我們一起思想主耶穌基督為我們釘十字架.
開始之前我們先低頭祈禱.
開始之前我們先低頭禱告.
天父我們多謝你.
因為你愛我們.
你差遣你的兒子為我們死在十字架上.
並差遣你的兒子死在十字架上.
因你的死我們可以得到生命.
我們可以得到生命.
今天我們一起來思想十字架的時候.
今天當我們一起思考十字架的時候.
讓我們不單止只是回顧一些以前發生了的事.
我們並不是只是一般的回想以前發生的事.
而是讓我們思想十字架今天怎樣承擔我們的重擔.
叫我們亦都是將我們的生命連在十字架上.
如果不是十字架的救贖.
我們都不是現在這樣的人.
如果不是你在十字架上的救贖.
我們今天都不是這個樣子的.
願你自己的十字架繼續冒著我們.
願你的十字架時刻幫助我們.
奉耶穌基督的名祈求.
Amen.
今天大家都知道.
我們紀念耶穌基督為我們釘十字架.
當然今天我們都知道.
是紀念耶穌為我們釘十字架上.
有些時候作為一個牧師.
有一次我和其他牧者聊天的時候.
都發覺一件事.
有些時候作為一個牧師.
我跟其他人在談論的時候.
我都發現一件事.
每年的受苦節,復活節的講道很難講.
就是覺得每年的受難節和復活節的講道.
都是很難的.
因為很多時候最精華的內容.
其實非言語可以表達得到.

$^{41}$因為其實很多精華的內容.
不是言語可以表達的.
而是需要你去感受.
所以唱詩歌就很好.
所以唱詩真的是很好的.
唱詩歌你聽了幾十次,唱了幾十次.
你不會覺得又是炒爛飯.
就詩歌就算你唱了幾十次.
你也不覺得是在炒冷飯.
但是講道不一樣.
但是講道就不一樣了.
但是相信這個都是一個很重要的一個時間.
讓我們思想十字架對我們的意義.
當然我們也還是要知道.
這是一個很重要的時刻.
讓我們思想十字架的意義.
當我們思想十字架的時候.
你不知道有沒有想過.
其實我們今天怎麼看這個世界.
是需要用十字架的眼光來看的呢.
當我們思想十字架的時候.
不知道你有沒有想過.
我們看待今天的世界是.
需要用十字架的眼光來看的.
如果你看這個世界現在是一片美好的時候.
我保證你平時一定不會想十字架.
如果你看待這個世界是一片美好.
我估計你都不會想到十字架這件事的.
最多就是你做基督徒.
現在過受撫節的時候.
就可能參加一下一些有關受撫節的內容.
最多就是作為基督徒.
你就是在受難節的時候.
參加一些關於受難節的一些活動,內容.
因為受撫十字架是.
在黑暗的地方給我們思想到.
神為我們背我們的重擔.
因為十字架是在黑暗的地方.
讓我們來思考神為我們背十字架.
如果你根本是不感覺到.

$^{81}$這個世界很黑暗的時候.
我們也不覺得需要有什麼十字架.
如果你不覺得這個世界是黑暗的.
其實你也不覺得我們需要十字架.
最近不知道大家有沒有聽新聞.
最近不知道大家有沒有聽見這些新聞呢.
最多的當然就是現在美國和伊朗在打仗.
最多的就是美國和伊朗正在打仗是嗎.
其實不經不覺俄羅斯和烏克蘭也打了仗四年.
其實不知不覺俄羅斯和烏克蘭打仗也已經四年了.
可能其實已經見怪不怪.
其實好像已經見怪不怪了.
我前幾天我看新聞的時候就發覺.
都很少報導.
我在看新聞的時候都發現好像很少報導這件事了.
我要特地去找的時候就找到一兩條.
當我專門去找的時候呢.
才找到了一兩條而已.
最近就說俄羅斯有些無人機.
炸了烏克蘭的一個醫院.
連那些嬰兒都要逃難.
就是講俄羅斯的無人機.
炸了烏克蘭的一間兒童醫院.
很多的小朋友都要逃難.
因為都已經見怪不怪了.
根本都不報導了.
因為已經見怪不怪了.
所以根本就不報導了.
接著又看看除了這兩場打仗.
還有什麼新聞呢.
緊接著想一想.
除了這兩場戰爭以外還有什麼新聞呢.
原來其實這個世界很多地方都沒有人報導.
是很黑暗的.
其實這個世界有很多的地方沒有人報導.
其實都是很黑暗的.
蘇丹內戰你不知道有沒有聽過.
蘇丹內戰不知道你有沒有聽過.
已經內戰了幾年了.
已經內戰了好幾年了.

$^{121}$而很多小朋友是被人抓去做奴隸的.
很多的孩子被人抓去做奴隸.
他們抓那些小朋友之前.
有些是先殺死他們的父母.
在很多小孩子被抓之前.
是要把他們的父母殺死的.
整個世界其實都是很多黑暗的地方.
整個世界其實都有很多黑暗的地方.
但可能因為我們活在一個很舒服的環境.
我們好像都不想去留意任何的黑暗.
但好像我們因為生活在一個很舒適的環境當中.
所以我們根本就不想去思考這些黑暗的事情.
但是聖經給我們另外一個角度去看這個世界.
但是聖經給我們另外一個角度去看待這個世界.
沒錯這個世界一開始的時候.
神創造的一切都是好的.
其實一開始的時候神創造這個世界都是好的.
但是有一種黑暗好像就把神創造的.
好的一種光明去遮蔽.
但是有一種黑暗好像是把神創造的這些好的光明給遮蓋了.
整個聖經的故事根本就是在講這件事.
整個聖經的故事其實就是在講這件事.
就是一去到創世記第三章之後.
根本就是講到神要怎樣去把黑暗的世界恢復光明.
其實在創世記第三章之後.
就是一直在講神怎樣把這個黑暗的世界恢復到光明.
神做了很多事.
之前就試過用羅亞來重新開始.
之前就試過使用諾亞來恢復.
跟著又找阿巴拉罕一家人又重新開始試過.
跟著找以色列人的時候又讓他們頒布律法.
讓他們又重新開始試過.
跟著又對以色列人又頒布律法又重新再試過.
跟著頒布完律法根本是很戲劇化的.
跟著又弄了金牛族出來大家都知道.
那麼頒布律法其實有戲劇化的大家都知道.
跟著又找一個神所喜悅的人做皇帝.
大衛重新開始又試過.
緊接著又簡選了一個人大衛又來重新試過.
跟著又甚至乎連葛珈芒亡了.

$^{161}$亡了葛之後.
讓他們重新開始在耶穌時代神自己來到.
又重新開始又再一次嘗試.
甚至國破家亡之後又派了耶穌又重新來試過一次.
你看到整個聖經的故事就是神不斷不斷地來嘗試.
其實整個聖經故事我們看到神就是不斷不斷地來嘗試.
我不知道你小時候有沒有玩煲蠟的中秋節.
煲蠟.
就是點蠟燭玩蠟燭玩火.
不知道你小時候有沒有玩過蠟燭.
中秋節很多時候在香港是下雨的.
不知道你記不記得夜晚很多時候下雨.
中秋節的時候在香港一般都是在下雨的.
不知道你記不記得.
通常有三個中秋節起碼有兩個中秋節下雨.
起碼有兩三個小時都是要下雨的.
但是我們又很喜歡玩火.
但是我們很喜歡玩火的.
落著雨在外面點蠟燭.
儘管外面下雨我們也去點蠟燭.
我們發現其實是很狼狽的.
其實你發現很狼狽的了.
一邊點了很久終於點了.
點了很長時間是濕的終於燃起來了.
燃了一會有風一吹又熄了.
燃了一會風一吹又滅了.
然後又要重新再點過.
緊接著又要重新再點.
整個過程是不斷地重重複複.
整個過程就不斷地這樣重複.
有時候我看聖經的故事的時候.
其實某個角度神都好像是在做相同的事情.
有的時候我看聖經感覺到神好像是在做相同的事情.
每一次人對於神的背弱.
而神又要重新開始了.
就好像那些火光點著了又熄了.
整個人類對於神的那種拒絕.
不是一種單是無知那麼簡單.
不是無知那麼簡單.
是很主動的.

$^{201}$是主動的.
是有意識地去抗拒的.
是有意識地來抗拒光.
如果是單是無知就很容易搞.
其實來講如果真的是無知就很容易解決了.
無知你就教導他.
有教育那就搞定了.
如果真的是無知的話.
那你不是通過教育就可以解決了.
但是你知道教不定的.
有時候有些情況.
但是你知道有的時候教育是搞不定的.
因為是人刻意地去拒絕.
所以在聖經上.
神差遣先知多次多方地不斷地去跟人說.
所以神差遣先知不斷地多次地跟人說.
但是是被拒絕的.
但是都是被拒絕的.
其實在舊約聖經的故事.
大部分的先知都沒有什麼好下場.
其實舊約聖經大部分的先知都沒有什麼好下場的.
當然也不是每個人都被人追殺.
但大部分其實是沒有人聽的.
但大部分是沒有人聽的.
而黑暗也不代表人一時軟弱選錯了.
黑暗也不代表人一時軟弱選錯了.
而是你看到在聖經的故事當中.
整個的所有宇宙萬物都被一種黑暗勢力侵蝕.
而是你看到在聖經的故事當中.
整個的宇宙萬物都被一種黑暗勢力侵蝕.
今天靈修的時候我看到《利未記》.
《利未記》你知道有很多人不喜歡看.
其中因為有很多很繁繁複複的細節.
其中因為有很多很繁繁複的細節.
發霉的時候,洗不掉又要怎樣處理?.
衣服如果發霉了,要怎樣洗掉?去掉它?.
有很多人說,這些究竟講什麼?.
有的人說,你究竟在講什麼?.
聖經為什麼講不完發霉的東西都講了這麼久?.
怎麼聖經講這些發霉的事情都講了這麼久?.

$^{241}$但是背後想到,就好像這裡所說.
發霉是一個什麼樣的現象?.
其實它是在講發霉是一個什麼現象呢?.
其實是一種污穢的東西去侵蝕.
不斷地來擴散一個現象.
它其實是在講污穢的東西不斷地侵蝕發散一些東西.
人是需要處理.
人是需要處理這些嗎?.
應該說,神吩咐人是需要處理它.
而是神吩咐人必須要處理這些事.
在那件事情當中,其實是神要教以色列人.
其實現在我們就是活在一個這樣的宇宙當中.
其實呢,我仔細想一下,我們就在這樣的宇宙當中.
所以當每一個以色列人他們要進勝所有的朝見神的時候.
他要先處理那些污.
也就是說,以色列人親近所有去朝見神之前.
他一定要處理這些發霉的事.
我們現在就是活在一個發霉不斷地污穢擴散的世界.
其實我們現在就是處在一個這個霉菌不斷擴散的世界.
罪惡就是這一種這樣的污穢.
罪惡就是這樣的一個霉菌.
而你知道那些污不是在表面的.
而且你知道這個霉呀並不只是在表面的.
當它擴散的時候.
它會吃進去裡面的.
當它擴散的時候,它會入侵到裡面去的.
這個我相信大家很清楚.
這個我相信大家都很清楚.
有一次我不知道.
家裡的車庫裡面.
擺了一些可樂的汽水.
有一次我不知道我的車庫裡面擺了一些可樂的汽水.
罐裝的.
放在那裡我不知道.
放在那裡我也不知道.
放了一段長時間的時候.
你知道那些罐就開始爆了.
你知道過了一段時間,這個罐子就爆了.
然後那些可樂就漏到整個櫃裡面都是.
這個可樂就流到櫃子裡面到處都是.

$^{281}$但我都不知道.
我也不知道.
直到我發現的時候.
整個櫃都黑了,發霉了.
等我發現的時候,整個櫃都變黑了,都發霉了.
最下兩層全部都是黑色.
後面那兩層全部都黑色.
我不斷地用東西洗.
我不斷地用東西去洗.
又用漂白水洗.
又用漂白水去洗.
又用醋去洗.
用你想得到的東西,在家裡找得到的東西都拿來用去洗.
把家裡面能想到的東西都拿去洗.
不斷地擦.
但是擦不擦得掉呢.
擦不擦得掉呢.
都是擦不掉.
都擦不掉的.
整個櫃下面的東西都變成了黑色.
整個櫃子都變成了黑色.
因為那些霉已經進入了裡面.
最後其實就要拆了整個櫃,扔掉了.
最後沒有辦法,只能拆了,扔了.
我們就是活在一個這樣的世界.
那我們就活在這樣的一個世界.
而去到那個地步,根本是無處不在.
去到這樣一個地步,就無處不在.
黑暗好像籠罩著某個地方.
在人際的關係之中有黑暗.
黑暗籠罩在人際關係當中有黑暗.
在現在我們這個時代是.
有人類歷史以來最有效率可以聯繫到人的一個時代.
其實可以說我們現在是.
自從有人類世界以來.
最有效率能夠聯繫到人的關係的這樣一個時代.
你想想你以前要打個長途電話.
已經幾好困難了,是不是?.
你想一下以前打個長途電話都很困難了,是嗎?.
已經不說你發電報那些了.

$^{321}$都不說那些,發電報就更困難了.
但現在呢?.
你撥下手指就已經可以聯繫到.
一些香港周圍世界各地的朋友.
你只要動一下手指就能聯繫到香港.
世界各地的那些朋友了.
但現在的人的關係是最孤單的.
但現在人的關係是最孤單的.
現代人的關係最孤單.
已經知道現在現代人的抑鬱症是最嚴重.
也就是說現在人的抑鬱症意義是最嚴重的.
因為人越多聯繫的媒介就越少與人聯繫.
其實越多的媒介,媒體就越少與人的聯繫.
不是的,我一天發幾百個信息的.
但你說不是的,我一天要發幾百個信息.
我又聽過有個朋友說.
他說他來了美國兩年而已.
說他在WeChat群組有二萬多個人.
他說來了美國幾年以後.
他在微信群組的朋友有幾億個?.
二百多個.
二百多個?.
不,是二萬多個.
二萬多個,兩萬多個.
但是你想想在這些聯繫當中.
有多少是真正的交心呢?.
但是你想想這麼多人當中.
有幾個是真正可以交心的呢?.
就是因為沒有了真正的交心的聯繫.
人就陷入很多的孤獨當中.
就是因為沒有這些真正的交心.
所以人就陷入了孤獨當中.
而整個制度都是一樣的.
整個制度都是一樣的.
而現在在我們的制度裡面.
就好像人與人之間越來越不能夠信任.
在我們的制度裡面就好像人與人之間.
越來越不能夠信任.
你看看現在有多少新聞你是信的.
你想一下現在有多少新聞你是相信的?.

$^{361}$如果你還是看什麼新聞信什麼新聞.
你就真的還是活在二十年前的思想模式.
如果你現在還是看什麼新聞就信什麼新聞的話.
那你真是還活在二十年前的那樣一個情況之下.
現在根本不可以這樣信的.
現在根本你就不能夠這樣去信的.
現在你根本不知道一個什麼是真什麼是假的時代.
很多東西你信不過.
現在你根本不知道什麼是真什麼是假的.
很多東西你信不了的.
而在人的心靈的深處.
更加有一種黑暗是好像驅不散.
甚至在人的心靈深處也有一種黑暗是驅不散的.
在人的心靈的深處不斷地去尋求.
究竟我現在活在世上有什麼意義.
在人的心靈深處人不斷地在尋求.
我究竟在這個世界上有什麼意義.
現在好像活在一個世界.
好像走馬燈一樣不斷地切出這麼轉.
在這個世界好像走馬燈一樣不斷地這樣兜兜轉轉.
但是你不知道自己活在這裡有什麼意義.
你不知道自己在這當中有什麼意義.
我不知道為什麼要活在這裡.
我活著的目的是什麼呢.
我為什麼而活呢.
我為自己而活.
但是我自己都不知道究竟有什麼值得.
但是我自己都搞不清楚有什麼好值得的.
就好像黑暗籠罩著整個世界.
就好像黑暗籠罩著這樣的世界.
我在幾個星期前剛剛學了一個新的名詞.
我在幾個星期前剛剛學了一個新名詞.
叫做New Dark Age.
你聽過沒有.
叫做New Dark Age.
聽過沒有.
就是Dark Age你知道聽過了.
黑暗時代就是中古世紀.
叫做Dark Age.
就是中古時代黑暗時代.

$^{401}$最近我們現在是慢慢慢慢.
其實已經開始走向一個New Dark Age的時代.
慢慢的我們已經走到了一個新的黑暗時代.
在這個New Dark Age的時代人就不能夠分真假.
在這個新的中古黑暗時代.
我們就沒有辦法分真假.
人也都很難建立到一個真實的關係.
也很難建立真實的關係.
在這個New Dark Age的時代.
人也都越來越不能夠進入一種深入的思想模式.
在這個時代人也越來越難進入一個新的思想模式.
為什麼不能夠進入深入的思想模式呢.
為什麼不能夠進入新的思想模式呢.
因為你的專注力沒有了.
現在很多小孩其實已經發展到.
其實他們已經沒有了專注力.
很難完成看完一本書.
現在的孩子的專注力已經喪失到沒有辦法去好好的看完一本書.
也沒有了深入的思想模式之外.
也都不能夠建立到智慧.
而且也沒有辦法建立智慧.
因為智慧是怎樣來的呢.
因為智慧是怎麼來的呢.
智慧是你學會怎樣做一些事情.
而且人與人之間的交流.
慢慢慢慢累積這些經驗而得到智慧的嘛.
智慧是怎麼來的呢.
是你慢慢去學習一些東西.
然後人與人交流得到的智慧.
而現在很多東西我們交了給AI幫我們決定的時候.
根本你那些東西你都不知道怎麼決定.
你都不知道怎麼做.
你拿來智慧.
那現在我們把什麼東西都交給AI.
你都不知道AI是他怎麼得出這件事情的結論.
那你怎麼得到智慧呢.
現在整個社會就.
那本書其實是一些分析科技研究書.
現在我們慢慢慢慢進入了一個New Dark Age的時代.
那這本書實際上就是在研究.

$^{441}$我們慢慢的怎麼樣走入了這個新的中古黑暗時代.
但是那本書沒有講到一樣東西.
但是這本書並沒有講到一件事.
就是在基督徒我們知道.
除了那些科技好像帶領人進入黑暗當中.
人的心因為有罪的緣故.
很多時候是主動的進入黑暗的.
我們作為基督徒我們知道.
就是說雖然說沒有其他的東西把我們帶入.
但是因為人的內心深處是有罪的.
所以會把我們帶入這些黑暗時代.
而正正是這樣的緣故.
主耶穌來為我們釘十字架.
也就是這樣的緣故.
主耶穌來為我們釘了十字架.
其實整個世界的黑暗多方面.
其實根本你留心去想.
留心去看的時候.
我們是沒有辦法走出來的.
這個世界的黑暗.
如果你留心去想.
留心去觀察.
你是知道是沒有辦法走出來的.
但是聖經那裡就講到.
主耶穌親自進入人類的黑暗當中.
但是聖經就講到主耶穌親自進入了最深的黑暗當中.
十字架就是黑暗的深淵.
所以在《萬學福音》那裡說到.
當主耶穌釘十字架的時候.
遍地都是黑暗.
所以馬可福音講到.
當主耶穌釘十字架的時候.
遍地都是黑暗的.
這種黑暗我想是象徵性的.
這種黑暗其實是象徵性的.
當然那個時候的天文現象是黑暗.
也可能是當時的天文現象是黑暗.
但是講到整個世界根本都是.
好像因為耶穌釘十字架.
就是人類最黑暗的時刻.

$^{481}$也就是說耶穌釘十字架的.
也是整個人類最黑暗的時候.
在那裡主耶穌經歷到.
當然就是被人氣絕.
被人譏笑.
被人離棄.
被人羞辱.
在那個時候主耶穌經歷了.
被人氣絕.
被人譏笑.
被人離棄.
被人羞辱.
而在祂最後呼喊的時候祂說.
我的神啊我的神啊為什麼離棄我.
而在那最後的時刻祂呼喊.
我的神我的神為什麼離棄我.
在每一個黑暗的時刻最黑暗的.
不是單是那種的受苦.
在那個最黑暗的時刻.
最黑暗的並不是那個受苦.
而是神的離棄.
而聖經講到主耶穌就是這樣.
進入了人類的黑暗當中.
聖經也就是講主耶穌就這樣.
進入了人的黑暗當中.
而祂進入黑暗是有一個目的的.
祂進入黑暗是有一個目的的.
聖經說祂被藐視.
被人厭棄.
多受痛苦.
常經憂患.
祂誠然擔當我們的憂患.
背負我們的痛苦.
耶和華使我們眾人的罪孽都歸在祂身上.
主耶穌來到不單單是表達一種愛.
祂是要進入去人的每一個黑暗的深處.
特別最黑暗的.
就是罪的那個格局那裡.
而拯救我們.
進入到最因為罪被隔絕的黑暗的深處.

$^{521}$而所以在十字架最黑暗的時刻.
也是救恩最能夠顯現光明成就的時刻.
也就是十字架是黑暗最猖狂的時候.
也就是救恩光能顯現的最明顯的時刻.
在十字架那裡.
讓我們看到人的虛偽.
在十字架那裡.
讓我們看到人的虛偽.
我們平時不覺得自己虛偽.
我們平時並不覺得自己虛偽.
但在十字架的面前的時候.
當神光照我們的時候.
我們就發覺.
其實我們對神的心.
不是我們平時表現的那樣.
在十字架面前.
當神的光照耀我們的時候.
我們就發現.
我們好像並不像平時所顯現的那樣.
十字架也都看到政治那種殘酷.
十字架那裡我們都看到了政治的殘酷.
你不要期待這個世界有任何的.
無論是多好的政府.
多好的制度.
能夠將人去拯救.
你不要想像這個世界上有多好的政府.
多好的制度.
能夠把人得到拯救.
最終的權力其實是將人壓榨.
最終權力是把人壓榨.
十字架也都看到人心那種盲目.
在十字架那裡.
我們也看到人生的盲目.
你看到那些群眾在中旅主日的時候.
還是喊著和沙拿和沙拿.
奉主命來的.
事應當清俗的.
後來隔了都不夠一個禮拜.
接著就說釘他十字架.
釘他十字架.

$^{561}$我們看到在中旅主日的時候.
這些群眾搖著中旅主都在那裡說.
歡迎奉主來的名是值得稱頌的.
但是一個星期不夠.
他們就要把耶穌釘十字架了.
有時候我們就是這樣多麼的盲目.
很容易被人推動我們.
去鼓動我們的情緒.
有的時候我們就是這麼盲目.
很容易被人推動.
來鼓動我們的情緒.
在十字架裡也看到人心那種懶惰.
在十字架那裡我們都看到了人心的冷酷.
很多時候我們有很多事事不關己.
我們看到一個殘酷的世界.
我們不覺得關自己事.
在很多的時候我們覺得事不關己.
看到這個殘酷的世界.
我們覺得跟我們沒啥關係.
十字架當然.
也都看到主耶穌被人出賣.
十字架我們都看到了主耶穌被人出賣.
人際關係那種脆弱.
即使是最好的門徒.
都顯得是很脆弱.
人際關係的脆弱.
即使最好的朋友也是離棄了.
而那種制度的暴力.
也看到了制度的暴力.
十字架那裡表露無遺.
明明是貝拉多知道耶穌沒有罪.
但是就將他釘十字架.
十字架都表露無遺.
彼多拉明明知道主耶穌是無罪的.
但是還是施予暴力.
而當然十字架的最後就是死亡.
當然十字架最後就是死亡.
這一切其實都是看到.
主耶穌好像經過十字架的時候.
顯現出人世間那種黑暗的狀況.

$^{601}$主耶穌在十字架上面.
都匯聚了這些黑暗的狀況.
而這一切的黑暗.
其實好像當我釘十字架那一刻.
所有他不只是說出來.
而是承擔在他身上.
這種黑暗並不單單只是講出來.
而是承擔在他的身上.
當我們去思想.
主耶穌這樣釘十字架的時候.
其實我們不單止看到.
主耶穌承擔這種黑暗.
當我們想到十字架的時候.
我們也會想到.
主耶穌不單單只是承擔了這種黑暗.
弟兄姊妹你有沒有看到.
這些黑暗好像鏡子一樣.
其實都是照在我們的生命當中.
弟兄姊妹你們有沒有想到.
這種黑暗也好像一個鏡子一樣.
照到我們身上.
我們每個人都是活在.
一個這樣黑暗的世界當中.
我們每個人都活在這樣的.
我們每個人都活在.
一個這樣的情況當中的時候.
有什麼出路呢.
我們每個人活在這樣的情況當中.
有什麼出路呢.
十字架就是那個唯一的解答.
十字架就是這樣唯一的解答.
神不是用一種.
在旁觀者的方法.
來解決黑暗的問題的.
神並不是以一種旁觀者的.
這種態度來解決這個問題.
黑暗的問題.
神是自己進入了黑暗當中.
而是神自己進入了黑暗當中.
他自己承受了黑暗.

$^{641}$而為我們帶來了光明.
所以因為基督進入了.
我們人世的黑暗當中的時候.
我們就不再需要懼怕.
因為耶穌進入了這個黑暗當中.
所以我們就不再需要懼怕.
在我們生命當中.
某一種黑暗是基督沒有走過.
在我們的生命當中.
沒有一種黑暗是基督沒有走過的.
有時候你感覺到孤單.
有時候你感覺孤單.
有時候你感覺到身體的軟弱.
可能是身體的軟弱.
有時候你為你的家人擔心.
也可能是為家人擔心.
有時候為你在這個世界的局勢.
感覺到心靈的不安.
也是因為這個世界感到心靈的不安.
這一切基督都好像是為我們走過.
這一切基督都好像為我們走過.
因為祂是承擔我們每一個重擔.
因為祂承擔了我們的重擔.
祂受鞭傷為我們得到醫治.
祂受鞭傷是為了我們得到醫治.
所以我們不需要自己去擔這個重擔.
所以我們不需要自己去擔這樣的重擔.
我們仍然是活在一個黑暗的世界.
我們仍然是活在這個黑暗的世界.
沒有改變的.
直到有一天我們知道神會回來.
但在這一天還沒回來之前.
我們仍然活在黑暗的世界.
那你說怎麼辦呢.
在祂還沒有再來之前.
我們還是活在這個黑暗的世界.
那我們怎麼辦呢.
我們就知道跟耶穌同行的時候.
祂會一直背著這個重擔.
那我們知道以耶穌同行的時候.

$^{681}$祂一直在承擔這個重擔.
十字架就告訴我們.
這個世界的黑暗不是我們自己走.
是神和我們一起走.
這個十字架的黑暗就讓我們知道.
我們在這個世界上的黑暗.
不是我們自己在走.
而是神在與我們同行.
所以沒有一個在黑暗當中的深淵.
是神的恩典觸摸不到.
所以沒有一個黑暗當中的深淵.
是神的恩典所不能觸及的.
接下來我們會一起領受主餐.
接下來我們一起領受主餐.
主餐其中一個意義就是.
我們是分享主耶穌的生命.
主餐的其中一個意義就是.
我們分享主耶穌的生命.
當我們一起領受的時候.
讓我們一起來思想.
當我們一起領受的時候.
我們一起思想.
我們雖然活在一個黑暗的世界.
但是主耶穌與我們同行.
祂的生命也在我們裡面.
我們是與祂一起走這段在世界的路.
我們是與祂一同走這個世界的路.
我們一起低頭祈禱.
我們一起低頭禱告.
天父我們多謝你.
因為你為我們寫了十字架.
因為你為我們死在十字架上.
想到我們的主走過人生每一段黑暗的路.
想到主與我們一起走過人生的某一段路.
每一個重擔你為我們擔.
每一個痛苦你與我們一同經歷.
讓我們不單單只是在頭腦上知道.
當我們領受主餐的時候.
你的聖靈親自與我們同在.
進入我們的內心.

$^{721}$親自安慰我們.
照亮我們.
我們這樣禱告.
奉耶穌基督的名祈求.
我們禱告奉耶穌基督的名.
阿們.
\newpage



\section{哥林多前書 15:1-4}
\label{sec:031935}
\textbf{復活和基督徒的盼望 (華/粵)}~~~~姚臣泰長老
\newline
\newline
連結: \href{http://www.church.com.hk/acms/content.asp?site=cdc\&op=show\&type=product\&code=031935\&layout=sermon}{\texttt{031935}} ~~~~ 語音日期: 2026-04-05
\newline
\newline
\hyperref[sec:031934]{\small{< < < PREV SERMON < < <}}
~
\hyperref[sec:index_chronic]{\small{[返順時目]}}
~
\hyperref[sec:031936]{\small{> > > NEXT SERMON > > >}}
\newline
\newline
哥林多前書 15:1-4
\newline
\begin{longtable}{cl}
\hline
\hline
章節 & 經文 (和合本修訂版)\\
\hline
15:1 & \begin{tabularx}{0.7\textwidth}{X} 弟兄們，我要你們認清我先前傳給你們的福音；這福音你們領受了，又靠著它站立得住， \end{tabularx} \\ \\ \relax
15:2 & \begin{tabularx}{0.7\textwidth}{X} 你們若能夠持守我傳給你們的信息，就必因這福音得救，否則你們是徒然相信。 \end{tabularx} \\ \\ \relax
15:3 & \begin{tabularx}{0.7\textwidth}{X} 我當日所領受又傳給你們的，最重要的就是：照聖經所說，基督為我們的罪死了， \end{tabularx} \\ \\ \relax
15:4 & \begin{tabularx}{0.7\textwidth}{X} 而且埋葬了；又照聖經所說，第三天復活了， \end{tabularx} \\ \\
[1ex]
\hline
\hline
\end{longtable}
$^{1}$丁鴻誠:丁鴻誠媒體平安.
丁鴻誠:復活節快樂.
丁鴻誠:我這個禮拜同事跟我講說.
丁鴻誠:復活節是我最不喜歡的一個節日.
丁鴻誠:我有同事跟我說復活節是我最不喜歡的日子.
丁鴻誠:我明白他為什麼這樣子說.
丁鴻誠:我會明白他為什麼這樣說.
丁鴻誠:因為他不信.
丁鴻誠:因為他不信耶穌.
然而對我們這些已經信 已經接受 也信靠基督的人.
對於我們已經是信 已經是信靠主耶穌基督的人.
這是極大的福氣.
這個是極大的福氣.
那我們今天一起看短短的一段經文.
出自哥林多前書的第十五章.
我們今天要看一段短短的經文.
就是出自在哥林多前書第十五章.
我們謝謝施斑剛才給我們帶來的信息.
我們多謝施斑給我們帶來一個很好的訊息.
因為在西布賴書第一章第一到第二節.
第一到第三節有這樣美好的一個神的話.
在希伯倫書第一章一到四節.
所以祂早已立祂為承受萬有的.
也就是祂創造諸世界.
祂已經是叫祂成就萬有.
成就祂創造諸世界.
祂是神榮耀所發的光輝.
是神本體的真相.
常用祂的全能的命令拖住萬有.
祂藉著祂的全能拖住萬有.
祂洗盡了人的罪.
就坐在高天自大者的右邊.
第四節祂所承受的名.
祂在第四節說祂所承受的名.
既比天使的名更珍貴.
比天使的名更加尊貴.
就遠超過天使.
就準備敬拜.
丁姐妹敬拜.
誰是忠心.

$^{41}$誰是敬拜的忠心.
如果今天你我都崇拜.
我們所謂的祖日崇拜.
你的心可能還在飄來飄去的話.
如果你來到這個敬拜.
但你自己的心仍然是飄來飄去的時候.
我希望我們能夠聚焦.
我希望你能夠聚焦.
因為如果祂不是忠心的話.
我們的敬拜.
我們就要打一個問號.
如果祂神不是你的敬拜.
忠心的時候.
我們的敬拜就要打一個問號.
所以首先並不是你好我好為主.
這個重要.
但不是主要.
在這個離對道.
你好我好不是很重要.
但是如果你跟神沒有和好.
如果你沒有跟神來和好.
那我們盼望神.
剛才我們所唱的.
所祈求的這個美好的福音.
能夠臨到你.
我希望光從我們所看.
所唱到所讀到的.
能夠引導你們.
我們今天看這段聖經是.
開始弟兄.
那個負責的弟兄有點緊張.
因為58節.
而且在今天短短的20分鐘.
怎麼可能講完.
其實這段經文長長二十幾節的話.
怎麼能夠在短短的時間講完呢.
我今天只會注重一到四節.
然後我會取其中.
其他十五章的其他經節來一起看.
我今天會集中在一到四節.

$^{81}$但仍然都會取其他的經文作為我的講道.
那你發現說在第一節.
他說弟兄們我託弟兄姊妹們.
我要你們認清.
先前所傳給你們的福音.
在這裡講第一節講到弟兄.
當然包括姊妹.
我要你們認清.
我先前傳給你們的福音.
所以我們發現說在這裡.
有幾次福音出現.
我們看見這一節裡面.
有幾次福音的出現.
它不但是我們曾經聽過的.
接受的這個福音.
它不是單單是.
我們曾經聽過接受的一個福音.
所以你發現曾經.
我們問人家說你信了主沒有.
通常他們第一句話是.
我曾經對吧.
當你問一個人的時候.
你順順主媒他說我曾經.
我在某某步道會我走到前面來.
在某某的報道會.
我是信了耶穌.
或者是有某某人向我.
分享了福音的福音.
我接受了.
或者某某人向我傳福音.
我接受了.
但我們留意一下說.
使徒保羅一開場.
你知道哥林多前面說很奇妙.
我看見保羅一開場.
說那麼多前書很奇妙.
他前面講了哥林多很多很多的問題.
直到最後.
幾乎是到十五章那邊.
才進入所謂的這個真理性的一個.

$^{121}$一個非常清楚的一個.
十五章以前.
保羅在提到一間教會.
有很多很多不同的問題.
直到第十五章.
他才講到關於真理的問題.
所以他希望你認清.
他曾經跟他們傳的這個福音.
他提醒他們要認清.
保羅以前跟他們傳的福音.
所以認清我們領受了.
所以我們接受了.
我們接受了.
對吧 這是過去式.
認清已經是知道.
也是接受.
然後我們發現說福音.
我們仍然要靠著天天來生活.
但是我們發現.
這個福音是需要我們每天天天來生活.
為什麼呢.
因為在這裡我們發現說.
你領受了.
又靠它站立的住.
因為在第一節那裡講到.
他們領受了.
但仍然需要靠它來站立得住.
剛才希伯來書給我們看見說.
神接著他的命令拖住萬有.
但是在希伯來書我們看見.
他提醒我們.
神是靠著基督來托著萬有.
但這裡看見說.
我們靠著福音.
能夠今天站在這裡.
站在這個恩典上.
能夠向神敬拜.
是神為我們所做的事情.
在這裡我們看見.
我們是靠著福音領受了.

$^{161}$然後靠著它站立得住.
然後靠著它來敬拜神.
所以過去 現在.
所以講到是過去和現在.
那第二節請留意.
他講到說你能夠持守.
在第二節那裡講到.
他能夠持守.
而且他講到這個福音很奇妙.
他說就必因這福音得救.
就是他很奇怪用一句就說.
必因這福音得救.
所以如果英文講說.
他發現不單是我們現在進行是.
We are now being saved.
那也講到第二節.
講到有將來的一件事情.
他不單止就是一個連續式的.
而且是講到將來.
弟兄姊妹 福音要持守啊.
弟兄姊妹 福音是需要持守的.
他們自己給我們很多的.
一些的所謂的這些提醒.
如果可以用這個詞的話.
我們可以講是.
在很多地方他給我們一個提醒.
改革中就是.
我們在post reformation之後.
給我們很多的一些的提醒.
在改革中的裡面.
他們裡面是給我們一些很多的提醒.
你發現說我們得救的卻據.
我們發覺我們得救的卻據.
是因為我們能夠繼續不斷的持守.
所以我們並不是說.
我知道這個多有把握.
這個是真實的.
福音已經成就了.
我們不是說.
我們看見福音已經是成就了.

$^{201}$但是對於我們怎麼樣子.
所在所信的福音上.
我們能夠站穩繼續的持守.
這就是我們福音對於我們的信的這個卻據.
就是說當我們得到一個福音的時候.
我們能不能夠繼續持守.
這個是很重要的事.
所以長話短說.
福音如果今天或對於你的將來.
還是說你的第一個反應.
如果是我都聽過了.
我請丁公子妹重新再想一想.
如果福音對你來說.
只是說我曾經聽過話.
我希望你再一次重新思考.
福音如果我可以簡單的來描述的話.
就像一個在進行式的.
就是時間在進行這樣.
如果我來形容福音的時候.
我們會用一個進行式的福音來形容.
你發現過去我們曾經領受.
這個圖給你看見說.
我們過去曾經領受.
那我現在我仍然繼續靠著福音站立.
在今天我們仍然是需要靠著福音來繼續站立.
我們要忍耐到底.
我們需要忍耐到底.
就必因福音得救.
就必因借福音得救.
所以你發現福音在第三個福音.
在這個圖片上面.
第三個福音.
當我們看到福音第三次的時候.
它就告訴我們將來要成就這件事情.
是告訴我們將來會成就的事.
所以並不是說救恩.
你現在得到救恩沒有完成.
不是說你今天所得到的福音沒有完成.
那耶穌在十字架上說成了的時候.
祂不但是成了之後.

$^{241}$祂埋葬三天後從死裡復活.
現在坐在天父的右邊.
不是說耶穌基督是受死之後.
三天之後復活.
現在是升天坐在神的右邊.
這個已經是成了一個事實.
這個已經是一個成就了的事實.
但落實在你我的身上.
還是一個進行式.
但是將來的落實在我們身上.
落實在我們身上的時候.
仍然是一個繼續的形式.
所以丁永哲美.
你能夠持守福音.
不是因為我曾經擁有.
那後面是什麼.
我們不是說只是講到.
我們曾經擁有這個福音.
我們擁有的是一個生命.
會繼續在我們裡面動它的美善工作.
我們所擁有的生命.
就是在我們裡面繼續它的善工.
剛才我們老師會也讀到.
剛才師班也有略略的提到這一節.
就是在彼得賢書的一章三到五節.
剛才師班也略略提到.
這節的經文是彼得賢書.
那我們發現什麼是活的盼望.
什麼叫做活的盼望呢.
A living hope.
我們發現說在這裡.
我們剛才講到說我們要沉浸性.
要靠著福音站立.
以後要繼續忍耐到底.
這好像是我們.
那在彼得賢書這一段.
我要你們看.
神為你我做了什麼.
在剛才我們看到的經文是講到.
過去我們曾經領受.

$^{281}$而是今天繼續來到站立.
然後將來忍耐到底.
但是那個是對我們將來有活的盼望.
所以在這段神經.
我們要看到說.
可以從另外一面看到說.
神為你所成就的是什麼.
我們可以看到神在你的生命裡面成就了什麼.
所以曼廷弟兄.
可不可以麻煩你幫我念一下三到五節.
請聽我讀第三到第五節.
願頌讚歸於我們主耶穌基督的父神.
他曾照自己的大憐憫.
藉著耶穌基督從死人中復活.
重生了我們.
使我們有活的盼望.
好得道不扭壞.
不沾污.
不須殘.
為你們存留在天上的記憶.
就是為你們這些藉著信蒙神大能保守的人.
能獲得他所預備到末世要顯現的救恩.
當末世要顯現的救恩.
因為提到在末世要顯現的救恩.
所以丁主耶不要看錯了.
得救不只是靈魂得救而已.
我們不要看錯了.
我們的得救不是單單靈魂得救.
神預備的救恩是靈魂體.
或者如果你是二元論的話.
靈魂和體.
神所立定的救恩是靈魂.
包括靈魂體.
那我們看見說神重生了我們.
復活重生了我們.
使我們有活的盼望.
那第二部分.
你剛才可以看在上面說.
我們得到的積業是不能夠衰殘.
不能夠學壞.

$^{321}$所以你現在有多少錢拿沒有關係.
因為第一次在第三節說到神重生了我們.
令我們能夠得到不衰殘.
在天上的積業.
這個不會影響到我們想要得到的一切.
要多少.
所以我們發現說第五節它講說.
你們藉著信蒙神大能保守的人.
第五節說到你要藉著信蒙神大能保守的人.
有個翻譯說.
我用英文先講.
他說Capt for Jesus Christ.
用英文是Capt for Jesus Christ.
就是神的能力保守的.
所以你信穩不穩妥.
非常非常穩妥.
所以問大家說你所信的穩不穩妥.
非常之穩妥.
不單是因為我拼命忍耐.
堅持到底要抓住基督.
不是說單單只是我迫切地去忍耐.
緊緊地去抓住基督.
你不要忘記神的手比你大.
但是你不要忘記神的手是比我們大.
沒有人 沒有任何權勢.
能夠把我們從祂的手中奪去.
沒有任何權勢能夠把我們從神的手中奪去.
我們接下來看第三到第四節.
我們看一下第三和第四節.
所以保羅在這裡講到.
我當日傳給你們的最重要的就是.
所以第三節說到.
我當日所領受全級裡面最重要的就是.
你如果發現你有合合本的話.
剛才師父會幫我們唸的.
我們一起誦讀的這段聖經.
它是說我當日所傳給你們的是第一.
是 剛才主席所讀的是.
我全級你們的第一.
所以有沒有第二.

$^{361}$有沒有第二件事.
如果你是engineer的話.
曾經在想說第一聽完了.
結果還是沒有第二.
如果你是做工程的講完第一.
必須要有第二 是不是.
豈有此理 講這麼多 五十八節完沒有聽到第二.
但是我們從整篇的聖經裡面讀到五十幾節.
我們都找不到第二這個字.
但是應該我這裡是用新標點的合合本.
我是用新譯的和合本.
它這裡告訴我們最重要的是.
它用的詞就是說最重要的.
三點 交三點.
所以第一個 基督為你 我的罪死了.
就是說基督為你和我的罪死了.
埋葬了 又照聖經所說第三天復活.
埋葬了又在三天之後復活.
這很重要.
所以這件事是重要的.
所以我們不能夠說Easter.
我們教會沒有做那種彩蛋的活動.
我們沒有彩蛋.
Easter Egg.
我們沒有找那些復活蛋的活動.
我這次很怕.
哇 我們那組環球我忘記了誰是主角.
有時候有些人說 為什麼這件事忘記了.
我們的符音 傳的符音很簡單.
但是重要.
我們所傳的福音是很簡單.
但是很重要.
我沒有放在上面 所以這裡告訴我們說.
如果我們所信的基督他沒有復活的話.
我們這裡沒有講到.
但是這裡講到 如果基督沒有復活的話.
死人沒有復活的話.
如果死人沒有復活的話.
那我們人就是死在最終.
我們的人就是死在最終.

$^{401}$你發現說 有人問說 你怎麼知道他死 他復活了.
但是有人會問 你怎麼知道他已經是死了 也是復活了.
照著聖經所說.
保羅是 多謝保羅.
他說是照著聖經所說.
我不知道你的信是建立在你多少年信了在教會.
我不知道你的信是建立在什麼.
是不是因為你多年在教會裡面.
但是我們一定要相信聖經所說的話.
但是我們必須要相信聖經裡面所講的說話.
三天從死裡復活.
三天之後從死裡復活.
是應驗 之前所預言的.
是實現了英雲的.
所以丁鴻準備 沒有一樣事情是偶然的.
如果你對聖經不熟悉或有疑惑的話.
我鼓勵我們繼續的讀 繼續的學習.
我鼓勵大家繼續去讀 繼續去學習.
所以我很快的到第二段.
所以這邊我們不需要讀.
但是我想提醒兄弟妹妹.
既然我們以後要復活.
所以不要看錯了 有次序.
我們將來 提醒兄弟妹妹.
我們將來會復活.
我們復活是有一個次序.
當然這裡講到出生的果子是基督.
然後在祂來的時候.
那些屬於基督的也要復活.
也有提到說基督要首先復活.
還有其他的信耶穌的人會跟著復活.
那我們給你們很快的提煉一下.
所以或許你們可能次序上有一點不同意.
不過這個是我自己最末世的一個認知.
這個是我自己在墨西.
墨西 對不起.
所以知道墨西的信息.
所以基督要先復活已經成就了.
基督已經復活 已經成就了.
那以後你我也要復活.

$^{441}$以後你和我都會復活.
那可能我們還活著的話.
如果基督現在來的話.
我們也要復活.
我們不用復活了.
其實我們就直接提升了.
我們沒有機會復活.
如果我們仍然活著的時候.
基督來的時候.
我們是與祂們一起.
不需要再復活.
但是如果等到以後將來我們死了.
那以後我們也要復活.
如果我們將來死.
然後跟著基督來的話.
我們將來也都會復活.
那舊約的信徒怎麼辦.
你說舊約的信徒又怎麼搞呢.
也同樣的 他們也要復活.
他們都會同一樣東西復活.
所以你可以不同意先後次序.
但是重要的是在族裡面的人.
他們都要復活.
可能你對於那個次序有些不認同.
但是我們凡是信耶穌的.
都會復活.
那不信的.
你不知道他們也要復活.
你知不知道不信的人都會復活.
所以人人都要復活.
所有的人都要復活.
只是時間次序不一樣.
但是那個次序是不一樣.
但是如果是到那時候再復活的話.
是面對神大寶座.
白色大寶座的一個選判.
對於不信的人.
我們都會面對神白色大寶座的審判.
所以人人都要復活.
每個人都要復活.

$^{481}$你什麼時候復活.
你在什麼時候復活呢.
救恩的成就 包括我們的身體.
神的救恩 包括我們的身體.
所以最後我給你看五十一節.
因為時間關係.
我很快提醒說.
保羅告訴我們一個奧秘.
因為時間緣故.
保羅在這裡提到一件奧秘的事.
我們不是都要睡覺.
而是都要改變.
我們不是要睡覺.
而是都要改變.
所以以後你的身體.
我們復活體就像基督一樣.
我們這個身體.
以後我們復活的時候.
都好像基督一樣.
不要想像matrix那樣.
能夠穿越那個牆.
這東西我不知道.
不是穿越牆.
但是我相信是不一樣的.
一個體質因為是神給我們這個復活體.
是要承受那個永遠的.
我相信是我們將來復活的形體是不一樣的.
是不會朽壞的.
不會扭壞的.
我們基督徒有這個盼望.
基督徒是有個盼望.
靈魂的得救.
不是說以後的永永遠遠.
我就浮在天空空氣這樣子繼續的.
靈魂而已.
我們得救不是說.
我們飄浮在空中這個狀態.
而是我們身體也要復活.
來承受這個永遠的光.
而是說我們的身體也都會復活.

$^{521}$來承受這個光.
基督復活了.
所以我們敬拜有一個轉移.
所以我們敬拜有一個不同的意義.
一個轉移.
我們的敬拜有一個轉移.
所以不單是祖日.
你現在過去我們都是什麼.
都是敬拜自我的.
有時候我們的敬拜是很自我性的.
我們要轉移到什麼.
唯獨基督.
我們要轉移到基督那裡.
第二部分.
我們是神給我們的救恩.
是全人的救恩.
神給我們的救恩是一個全人的救恩.
通常我發現我們傳福音給人的話.
我們少了這部分.
很多時候我們傳福音.
可能我們都會缺少了這部分.
你們有多少人跟他們講說.
我們將來要復活.
有沒有人我們傳福音的時候.
跟你說你將來會復活.
我發現很少人.
比較少.
基督會復活了.
那就好.
我得救了.
靈魂OK了.
那就OK.
我們會說基督復活了.
我們信耶穌之後.
我們的靈魂復活了.
身體也要復活.
身體也都會復活.
那我們有一個得救的確訣.
我有一個得救的確訣.
所以我鼓勵你們.

$^{561}$為這福音.
為這你所信的.
你要持守.
你要堅持到底.
我們鼓勵大家.
對你們所信的.
要持守和堅持.
總有一些困難的時候.
當然是會有一些困難的時候.
總有一些受試煉的時候.
總也會有一些試煉的時候.
但是我們要靠主站穩.
我們要靠主站穩.
第四我們有一個極大的產業.
換句話說.
你不用看別人眼紅.
我們不需要看別人眼紅.
所以有時候發現來教會的話.
看到這個人開著什麼車子.
在停車場繞來繞去.
有時候我們會去停車的時候.
人家在開什麼車呢.
丁主持人這些帶不走.
這些物質我們帶不走.
但是神應許給你我的.
將來有一個極大的應許.
神將來應許我們有一個極大的經驗.
我們在這邊很多退休的.
我們能關心我們退休金還有多少嗎.
很多我們當中已經是退休的.
我們都很關心我們的退休金有多少.
丁主持人在基督裡.
您這個已經是我不曉得他的天文數字.
但是對基督來講.
我們對基督來講是一個天文的數字.
不愁吃不愁穿.
而且我們不用擔心.
吃的是上好的東西.
我們都相信我們食後也是上好的東西.
我們能夠享受魚主同在.

$^{601}$能夠享受魚主永遠的同在.
我們能夠魚主永遠的同在.
最後要勇敢地活著.
我們要勇敢地活著.
不要怕死.
我們在軍隊當中.
我們最怕的敵人是那種不怕死的敵人.
我們在軍隊裡面最怕的是什麼呢.
是不怕死的敵人.
你們看過這種人嗎.
不知道大家有沒有聽過.
我碰過幾次了.
我真是對他們很多的尊重.
我碰過很多次.
但是我對他們有很大的尊重.
基督徒可以不怕死.
因為我們的主為我們死而復活.
因為我們不怕死.
因為我們的主已經為我們死.
為了我們復活.
祂也釋放我們往.
跟其他人不一樣.
已經釋放我們不再怕死的那種預備.
我只有到今天.
我只有到六十五歲.
七十五歲.
八十五歲的這個期限.
所以祂釋放我們.
是叫我們不會怕死.
不會想到我們什麼時候會死亡.
英文有個字叫zeal.
一個熱情.
是這樣翻譯嗎.
是吧 是一個熱情.
這個跟年紀沒有關係的.
這個跟年紀也有關係.
你更年期也好.
還是什麼期也好.
退休期也好.
沒有關係的.

$^{641}$跟那個什麼期是沒有關係的.
是不是退休期 什麼期.
但如果你對主的熱情應該是什麼.
越來越大的.
但對神的熱情應該越來越大.
這個是跟外面的人不一樣.
可以嗎.
所以我知道每天早上起來很困難.
有時候我覺得自己早上起來很困難.
但是對主的這個zeal應該是什麼.
與日俱增.
我們希望我們對神那種熱情是與日俱增.
願主祝福我們 我們一起禱告.
願神祝福我們神的懷緣.
因為你的死 你的復活.
教你的救恩成就.
因為你的死 你的復活.
將救恩成就.
也給我們這些陰性的.
有永遠的保護.
你給我們上的人有永遠的盼望.
但願你自己的救恩.
親自領到在座的每一位弟兄姊妹.
我希望你自己的救恩.
臨在在座的頂智媒.
讓他們清楚知道所信的是什麼.
讓他們清楚知道他們所信的是什麼.
以致他們每日的生活的重心.
以致他們每日生活的重心.
是因著你 為著你.
Amen.
謝謝.
\newpage



\section{腓立比書 3:10-11}
\label{sec:031936}
\textbf{基督復活的大能}~~~~杜崇基牧師
\newline
\newline
連結: \href{http://www.church.com.hk/acms/content.asp?site=cdc\&op=show\&type=product\&code=031936\&layout=sermon}{\texttt{031936}} ~~~~ 語音日期: 2026-04-05
\newline
\newline
\hyperref[sec:031935]{\small{< < < PREV SERMON < < <}}
~
\hyperref[sec:index_chronic]{\small{[返順時目]}}
~
\hyperref[sec:031937]{\small{> > > NEXT SERMON > > >}}
\newline
\newline
腓立比書 3:10-11
\newline
\begin{longtable}{cl}
\hline
\hline
章節 & 經文 (和合本修訂版)\\
\hline
3:10 & \begin{tabularx}{0.7\textwidth}{X} 使我認識基督，知道他復活的大能，並且知道和他一同受苦，效法他的死， \end{tabularx} \\ \\ \relax
3:11 & \begin{tabularx}{0.7\textwidth}{X} 或許我也得以從死人中復活。 \end{tabularx} \\ \\
[1ex]
\hline
\hline
\end{longtable}
$^{1}$各位弟兄姊妹平安.
復活節的主日是.
一年裡一個最重要的主日.
其實我們每一個主日.
其實都是紀念耶穌基督的復活.
但是我們在教會年裡.
這個復活期不僅是今天的.
由今天一直以後有八天.
在年曆裡都是一個復活期.
我們每天都是紀念主復活的大能.
在我們身上.
不過今年很特別.
因為今年我們的清明節.
和我們的復活節是同一天進行的.
我覺得很有意思的.
清明節就是神宗追源.
我們想起我們的先人.
其實說起祭祖.
在我們中國人來講.
因為信仰傳了進來.
基督信仰傳進來之後.
曾經都有過一個很大的爭議.
我們基督教曾經有一個想法.
就是我們不應該祭拜或者祭祀.
以致讓我們的祖先得尊榮.
所以來到近代.
其實我們一般都已經有一個.
相當一致的立場.
就是今天我們仍然都是選取了.
這個紀念先人的正義.
正義就是.
我們要紀念先人在以往的人生中.
他有什麼榜樣.
他有什麼地方值得我們去追念.
所以我們去到的時候.
不是拿祭玉去拜祭.
不是拿香燭去.
因為我自己要去聯合道基督教墳場.
牌子都寫得很清楚.
我們不可以帶這些祭玉祭物去.

$^{41}$因為重點就是我們要追念逝去的人.
帶給我們的鼓勵和回憶.
以致我們能夠同樣尊榮.
神在他生命中的帶領.
所以每次我去掃墓的時候.
我都會帶領家人一起有一個祈禱.
這個祈禱就是我們去追念.
先人在我們生命中.
給我們留下的足跡.
留下的榜樣.
亦都為了今天我們能夠認識主耶穌.
我們真的為此獻上很大的感恩.
願意我們的生命都繼承.
我們真的在先人留下來的美好榜樣.
我們繼續前進.
同樣今天也是一個很特別的日子.
因為我們在復活節的下午.
我們會舉行浸禮.
復活節和浸禮.
放在同一天正好提醒我們.
復活節不僅是2000年前發生的一件事.
而是今天仍然發生能力的.
浸禮是超越一個禮儀.
而是我們一個見證.
我們與基督同死同埋葬同復活.
開始一個新的生活.
謝暗生牧師在《復活在當下》這篇文章裡.
提及紀念耶穌基督復活.
有三種時式的信仰意義.
第一,我們是紀念歷史上耶穌基督的受死和復活.
這件事在過去已經發生了.
但是我們也要在當下.
我們的信仰建立在祂的受死和復活之上.
這是現在的信仰.
謝暗生也提醒我們.
我們將來的永生.
是建立在當下對耶穌基督的受難和復活的認信之上.
復活的重點不僅是人死後上天堂.
而是人能否當下常與耶穌基督同在.
我們如果要活得一個有意義的人生.

$^{81}$就必須明白信耶穌.
不僅是說我們有永生.
信耶穌就是我們是永生.
我們就是永生.
而且這個永生是活在當下的.
有生命力的信仰.
不是說我今生不好.
來生才好.
絕非是這樣的.
有生命力的信仰.
不僅是說我們這一刻就是多麼美好的.
前大律師貝勒奇嘉仁為他寫下了一句庫文.
There is no time what is memory.
有人譯為光陰飛逝追憶何求.
這個正正是說明人生當下的重要.
也正好度盡了今天.
很多人不知道怎樣去活在當下.
造成生活中很多無謂的人和事.
浪費了人生.
所以如果我們能夠天天將耶穌基督的復活.
活在當下的生命裡面.
並以此信念來審視我們的人.
判斷是非和價值.
這個就是復活節當下的意義.
其實我們香港是得天獨厚.
但是我們復活節和聖誕節都是一個很長的假期.
所以這幾天我想不少弟兄姊妹都選擇出外去旅遊.
不知道大家覺得人有沒有特別少了呢.
暫時看就差不多了.
但是無論如何這兩個節期.
也是基督徒向外宣告信仰的應訊.
不單只是對未信的人一個很容易分享福音的機會.
同樣也是一種堅定我們在神裡面信心的慶典.
所以教會於受難節同時舉行崇拜.
表示我們對耶穌基督的受死是有一個紀念.
想問問大家.
有誰在星期五參加了苦路十四站.
舉舉手看看.
哇 都不錯.
有些人這次沒有機會來.

$^{121}$大家知道這次的重點不在聚會上.
重點是我們親身和主走十架道路.
在這個過程裡.
我們都紀念到主在這條十架道路上為我們跌倒.
亦都有婦女為他去抹他的血.
還有我也很體會到大英姊妹是全程投入在這個過程當中.
我們用五感.
我們用釘嘗試紀念主來展示釘痕的手.
你聽不聽到釘的聲音.
叮叮叮叮.
仰望十字架的時候.
到最後那個空墳墓.
即是那個有很多蠟燭的地方.
都是給我們有好心的體驗.
但是今天這個空墳墓.
這個墳墓的石已經滾開了.
這個墳墓裡面的主耶穌基督已經復活了.
所以今天我們藉著兩段經文.
我們去思想到復活到底是甚麼呢.
復活為甚麼會是那麼重要呢.
和復活的大能如何成為我們新生活的動力和盼望.
讓我們一起禱告.
今天是復活節.
今天是紀念主耶穌復活的日子.
今天是整年裡面一個復活的最高峰.
今天願意我們的心都聚焦在主耶穌的身上.
我們更加願意保羅所說復活的大能.
都是在我們的生命裡面不斷產生的影響力.
繼續的激勵我們去前進.
繼續激勵我們去跟隨你的腳中而行.
願天父月立帶領我們以下這段時間.
專一的去思想你復活的大能.
我們這樣在你面前祈禱感恩.
奉基督得勝名至祈求.
阿們.
哥林多前書15章3到11節.
這段經文就告訴我們.
基督照聖經所說為我們的罪死了.
亦都埋葬了.
又照聖經所說第三天復活.

$^{161}$並且向很多的見證人顯現.
最後亦都向保羅顯現.
使他從逼迫教會的人.
變成一個傳揚基督的人.
弟兄姊妹請留意.
這不是我們個人覺得很感動的事件.
這是一個曾經發生過.
有人見證.
而且帶來一個永恆改變的真實事件.
所以復活是福音的核心事實.
保羅在哥林多前書15章說.
我當日所領受又傳給你們的.
第一就是什麼.
第一的意思就是最重要的.
最不能夠少的是什麼.
就是基督照聖經所說.
為我們的罪死了.
他埋葬了.
他照聖經所說.
第三天復活.
他顯然而且有很多的見證人.
弟兄姊妹.
復活不是附加在十字架之後的彩蛋.
復活是宣告十字架的救恩已被天父接納.
罪與死的權勢已被打破.
耶穌不是一位偉大的殉道者.
而是又真又活的主.
而且保羅特別強調顯然.
向基法就是彼得.
向十二個使徒.
向五百多個弟兄.
向雅各.
向眾使徒.
最後亦都跟保羅顯然.
換句話說基督的復活.
不是一個私人的幻覺.
而是一個眾人可以見證.
可以查驗.
且帶來生命改變的事實.
如果耶穌沒有復活.

$^{201}$基督信仰就只剩下道德的勸勉.
和宗教的情懷.
但是因為耶穌真的復活了.
基督信仰就成為救恩.
能力和盼望.
你或者會問.
復活跟我今天的生活有什麼關係呢.
保羅接下來就用自己的生命去回答.
他原本是一個逼迫教會的人.
卻因為遇見復活的主.
整個人生方向翻轉.
復活不單只是改變了一個人的情感.
而且是能改變一個人的身份,價值,道路和使命.
在我曾經服侍的教會裡面.
有一個很令我印象深刻的見證.
一個見證的名字就叫做從馬會到教會.
當他未認識主之前.
他常去的地方是馬會.
但是他有一次很特別.
走來我們的祈禱會當中.
他感覺到一股平安在他的心中.
而聖靈就不斷在他的心裡面做感動.
他後來去慕道.
他後來去受針.
他今天雖然不是什麼教會的領袖.
但是他在教會裡面.
很忠心的服侍神.
從馬會到教會.
我又見到一位的姐妹.
最初她是很怕死亡的.
因為她身患了一個癌症.
這個癌症最初是很令她擔心的.
因為這個病可能已經到了第三期.
不知道醫好的機會是怎樣.
她本來很害怕死亡.
但是當她認識了主耶穌之後.
她的生命有改變.
她知道她沒有辦法去計算她在世的日子.
她只可以珍惜每一天她的生活.
她不單止克服了死亡的恐懼.

$^{241}$她更加主動關心一些靈養的病人.
以致她很有力.
因為她跟人分享見證的時候.
她說我自己這個病.
我也不知道自己還有多少年日.
不過我有了主耶穌之後.
我心裡面有這份平安.
即使我今天晚上離開這個世界.
我也能夠有把握.
可以去到天堂與神同在.
很多靈養的病人聽到她的分享的時候.
心裡面都很受安慰.
甚至有些都能夠被邀請去決治相信耶穌.
又有一位老闆.
他未信主之前他的生命是很苛刻的.
即是說他很少去讚人的.
他只是去找出對方那些彈的地方.
他就在那裡說.
其實可能已經做得很好了.
但是他仍然都是看那些做得不夠好的地方.
但是當他認識了主耶穌.
他明白了這個復活大能的時候.
他開始改變.
他懂得去讚賞他的員工.
以致他給員工一種感受.
就是說我們是一組人.
我們這組人一起努力去工作.
達成我們工作上的目標.
他的生命從一個去挑剔人的生命.
變成一個去懂得欣賞人的生命.
又或者我們都經歷到很多不同的改變.
我不知道剛才所講的見證.
會不會你都有經歷過呢.
我們有時經常都會很羨慕有些人說.
好像保羅一樣.
他本來是逼迫教會的.
但是他現在欣賞著有一個一百八十的改變.
是一個很劇烈的改變.
我不知道有時你會不會有點羨慕.
那些人生命有一個這麼大的改變的時候.

$^{281}$自己好像很平凡.
沒有什麼特別.
覺得自己都是一個普通人.
我的改變是什麼呢.
但是我想再講.
其實我們每一個人都因為耶穌基督的復活.
因為信靠耶穌.
我們都有改變.
可能不是人體來劇烈的改變.
但是一定有改變.
如果沒有改變.
你就真的要想一想.
你是不是真的信得清楚.
是不是真的信得真.
最重要是有變.
所以復活的大能.
使我們認識基督.
並且進入一個新的生命.
接著我們看腓立比書三章十到十一節.
保羅說使我認識基督.
知道祂復活的大能.
這裡要留意.
保羅不是說使我知道一些關於基督的道理.
他講的是使我認識基督.
這個認識不只是一個頭腦的資訊.
而是生命的相交和關係.
我不知道你明不明白.
那個認識和那個經歷是有很大的差異.
當我有機會和一些未信的朋友聊天的時候.
其實你有沒有聽過耶穌.
他說有.
我在中學的時候聽過耶穌.
不單止聽過.
我考試也考過.
那時候會考.
我也考過聖經科.
我也有獎勵的.
但是你今天有沒有認識耶穌.
你今天有沒有跟隨耶穌.
他就說沒有.

$^{321}$為什麼.
這個只是他頭腦的認識.
不是祂生命的經歷.
所以當我們說復活的大能是什麼意思.
至少有三個層面.
復活的大能就是使死變成活.
弟兄姊妹.
復活大能最直接的意思.
就是神能叫死人復活.
但是在我們日常生活裡面.
死不是沒了口氣那麼簡單.
很多人活著.
心很多時候是麻木了.
有時候他的盼望是枯乾了.
有時候他的良心是沉重的.
關係是破裂的.
人生是失去方向的.
所以復活的大能臨到的時候.
就是神能夠將我們從屬靈的死亡.
心靈的枯槁.
罪的捆綁中帶出來.
賜下一個新的生命.
所以基督徒不是要更努力的去做個好人.
基督徒是領受了一個新的生命.
第二就是復活的大能使不可能的成為可能.
保羅提到復活不只是講能力.
也馬上講到.
並且知道和他一同受苦.
效法他的死.
這句話很真實.
基督徒的道路不是逃避痛苦.
而是在痛苦中仍有盼望.
仍然能夠忍耐.
仍然能夠去愛.
仍然能夠站立.
為什麼.
因為復活的大能不是讓我們免於風浪.
而是讓我們在風浪中不被吞沒.
你記不記得.
耶穌和祂的門徒坐船渡海的時候.

$^{361}$突然間狂風大作.
耶穌在船上休息睡覺.
因為祂太累了.
祂的門徒走來說.
「夫子啊,我們喪命了,你還不顧嗎?」.
接著你知道耶穌起來說了一句話.
「靜了罷,住了罷,風浪就大大平息」.
但事蹟背後還有另一個令人震驚的地方.
就是當時祂的門徒再一次懷疑.
究竟耶穌的身份是誰.
原來耶穌不是一個律法師.
只是耶穌不是一個老師.
原來耶穌是掌管大自然的主.
其實很多時候.
當我們遇到困難的時候.
我們最需要的不一定是答案.
而是需要力量.
當你面對失去.
面對病痛.
面對壓力.
面對關係的破碎.
你可能會說.
「我真的撐不下去了」.
但復活的大能不是對你說.
你要堅強一點.
復活的大能是對你說.
「是我與你同在」.
「我能夠扶持你」.
「我能夠帶你走過去」.
因為主曾經進入最深的黑暗.
死亡又從死亡走出來.
復活的大能使我們有終極的盼望.
不讓死亡去定義.
本來最後說.
或許我又得以從死人中復活.
這不是不確定神會不會保守他.
而是一種很謙卑的說法.
我願意一路跟隨基督到底.
直到那一天.
完全的復活和得榮耀.

$^{401}$復活給我們一個終極的盼望.
死亡不是句號.
憤怒不是終點.
今生不是全部.
我們的勞苦不是徒然的.
我們的眼淚不是白流的.
我們在主裡面的盼望.
不是一個自我的安慰.
而是建立在基督已經復活的事實上.
復活的大能.
怎樣能夠成為新生活的動力和盼望呢.
將哥林多前書十五章.
和菲勒比書三章放在一起.
我們看見了一個完整的圖畫.
哥林多前書十五章告訴我們.
復活是福音的核心.
是已經發生的事實.
是有見證的人有改變的.
菲勒比書三章告訴我們.
復活不單止是一個外在的事件.
更是要成為我裡面生命的實際.
我認識著我經歷能力.
在苦難中跟隨.
走向終活的盼望.
接著說那些很神學的.
不過希望大家都吸收到.
耶穌復活的重要教義.
是一個重要的教義.
祂的復活是有初期的猶太人.
對復活的觀念為一個背景.
是有一個歷史的證據.
包括空墳墓和耶穌多方的顯然.
祂的復活是身體更新的復活.
所以我們說耶穌的復活.
是靈和身體的復活.
特別我們強調耶穌的復活.
是身體的復活.
不是一個精神的復活.
耶穌的復活是一個身體的復活.
這個重要性就是說.

$^{441}$將來我們離開這個世界.
你有沒有想過香港火化很多.
我們多數火化.
所以我們化了身體就沒有了.
你又想過怎麼辦.
其實你不用怕.
耶穌的身體復活告訴我們有一個復活的身體.
所以我們那個復活的身體很重要.
復活的身體就是說.
將來是不是用我們那些化了灰的.
然後組織回那個身體.
不是 你有沒有想這個問題.
當然不是 一定不是.
一定是神重新再做一個復活的身體給我們.
所以這個復活的身體就好像耶穌復活的身體一樣.
我們可以經歷那個完全的改變.
所以今天我們這個復活的身體特別提醒.
復活就是叫信徒在末日有復活的盼望.
又開展了教會的新時代.
將救贖的歷史推向活世.
耶穌復活證明祂是以色列所等候的彌賽亞神的兒子.
這個表明了祂的教導是神聖的.
是神的啟示.
十字架和復活雖然是耶穌不同的工作.
但是我們既要將兩者要分別出來.
卻同時不可分割.
並且應該將兩者作合一的討論.
將十字架的復活和合一.
將令救贖清二成聖和末世更生的教義.
得著更清晰的解說.
所以今天的信息要很清楚地落在這一句上面.
「基督復活的大能帶給我們新生活的動力和盼望」.
不知道大家是否收到.
為什麼我這麼聰明呢?.
這是我的老師李逸唐牧師系統神學基督論那一段.
但是我想說什麼呢?.
其實我想說.
我們經常說耶穌基督受死.
耶穌基督復活.
其實是兩件不同的事情.

$^{481}$但是兩件合起來是很重要的.
兩件合起來.
讓我們將整個救恩很清楚地.
將整個圖畫表達出來.
有一隻蝴蝶從泳中爭奪出來的過程中.
牠是很奮力地掙扎.
但是正是這份努力.
讓牠的翅膀變得更堅韌強健.
所以你記得有一個說法.
如果我經常覺得爭奪太辛苦了.
不如不要爭奪.
不如我們剪掉牠 讓牠出來.
你知道結果是怎樣嗎?.
很快死.
因為牠沒有經歷過爭奪.
牠的疫味發展的時候.
基本上牠出來.
牠自己不能飛.
牠自己的變態過程不能變出來.
所以這份堅韌的力量.
這份翅膀的掙扎.
就好像耶穌復活的力量一樣.
教我們在困難中成長.
從而賦予我們再次振翅高飛的勇氣.
重生的時刻讓我們體會到.
復活的大能是我們生命中真正的動力.
所以當我們說信徒的清義和復活.
有一個密不可分的關係.
幫助信徒投進耶穌的生平.
死亡和復活的經歷裡面.
正正就是聖靈.
約翰福音14章26節說.
但保衛斯就是赴恩我的名.
所要差來的聖靈.
祂要把一切的事教導你們.
並且要使你們想起.
我對你們所說的一切話.
今天我們能夠清義成為主的門徒.
正正因為有復活的生命在我們當中.
聖靈就教導我們明白.

$^{521}$耶穌和耶穌所說的一切話.
我們今天成長的時候會是怎樣的呢.
有三方面我們去思想的.
第一方面我再不是為自己而活了.
我是為復活的主而活.
基督復活的宣告就是說.
耶穌是主 既然祂是主.
我的人生就不再是由我想要什麼來主導.
而是由主喜歡什麼來主導.
這個會帶出一個新的價值.
新的選擇 新的優先次序.
在金錢 情慾 權力 面子 成就之上.
我開始學習將主放在第一.
第二個動力就是真實的更生.
我可以重新開始.
因為赦免與更生是真的.
復活節最深的安慰就是失敗不是定局.
很多人以為自己不配了.
回不去了 趕不到了.
但是復活不是叫我們將過去洗白.
而是將我們帶到十字架前承認罪.
領受赦免.
並且靠復活的大能活出更生.
今天如果你心裡面有罪疚.
有捆綁.
有你不敢說出來的黑暗.
你要知道主不是要來將你趕走.
而是要將你救回來.
第三就是盼望.
就讓我們在苦難裡面仍有方向.
因為結局在主的手中.
當你在痛苦裡面.
最可怕的就是看不到意義.
以及看不到出口.
復活的主不一定立刻挪走每一個難處.
但是他將我們的未來握在手裡面.
他曾經勝過死亡.
他能夠將破碎帶向更生.
他應許最終的復活和新天新地.
因此我們可以流淚.

$^{561}$但是絕不會絕望.
我們可以哀傷.
但是不被吞噬.
弟兄姊妹.
浸禮為什麼常常都在復活節前後舉行呢.
因為浸禮是在眾人面前.
用一個很清楚的記號宣告.
我已經與基督同死同埋葬.
重復活.
從今以後我不再只屬於自己.
我是屬於主的.
我要走新生命的路.
對今天下午要受浸的弟兄姊妹.
我想對你說.
浸禮不是表示我已經完美了.
浸禮是表示我已經是屬於主的.
羅馬書第六章三到五節這樣說.
難道你們不知道.
我們這受浸歸入基督耶穌的人.
就是受浸歸入他的死嗎.
所以我們藉著浸禮歸入死.
和他一同埋葬.
是要我們行事為人.
都有新生的樣子.
像基督藉著父的榮耀.
從死人中復活一樣.
我們若與他合一.
經歷與他一樣的死.
也將經歷與他一樣的復活.
你不是靠自己的力量去做基督徒.
你是靠復活的主將生命賜給你.
並且每日供應你所需.
這個洪牧師指出.
復活節浸禮之所以特別感人.
正因為他將整個福音信息.
很具體地呈現在教會面前.
當受浸者進入水中.
象徵與基督同死同葬.
從水里上來.
象徵與主一同復活.

$^{601}$開始新的生命.
這不僅是受浸者個人的見證.
也再次提醒會眾.
我們所信的主已經復活.
我們領受的救恩是真實的大能.
我們所蒙照活出的.
也是一種更新了的生命.
因此每一次復活節的浸禮.
不應該只叫我們為受浸的人數感恩.
也應該叫我們重新思想教會.
有沒有忠心地培育門徒.
有沒有認真地陪伴人.
走過信主和受浸之間的道路.
對那些還未受浸.
或者還是觀望的朋友.
我也誠懇地說.
今天你聽到的.
不是怎樣加入一間教會.
你所聽到的是一位復活的主.
只是呼喚你.
祂知道你的故事.
祂也能夠改寫你的結局.
你可以從今天開始.
用信心去回應祂.
有我們一起的禱告.
弟兄姊妹.
在這段安靜的時間.
我很想發出三個的邀請.
你可以在心裡很安靜地對主說.
如果你今天已經信主.
求主讓你更加認識基督.
更加經歷祂復活的能力.
求主幫助我們.
不只是停留在聚會和知識.
而是我們能夠更深地在禱告.
在話語.
在信復裡面經歷祂.
在你的禱告裡面.
你可以對主說.
主啊.

$^{641}$使我認識你.
知道你復活的大能.
\newpage



\section{出埃及記 4:23}
\label{sec:031937}
\textbf{聖經人聖經事:為何要殺長子}~~~~黃克勤牧師博士
\newline
\newline
連結: \href{http://www.church.com.hk/acms/content.asp?site=cdc\&op=show\&type=product\&code=031937\&layout=sermon}{\texttt{031937}} ~~~~ 語音日期: 2026-04-05
\newline
\newline
\hyperref[sec:031936]{\small{< < < PREV SERMON < < <}}
~
\hyperref[sec:index_chronic]{\small{[返順時目]}}
~
\hyperref[sec:031938]{\small{> > > NEXT SERMON > > >}}
\newline
\newline
出埃及記 4:23
\newline
\begin{longtable}{cl}
\hline
\hline
章節 & 經文 (和合本修訂版)\\
\hline
4:23 & \begin{tabularx}{0.7\textwidth}{X} 我對你說過：放我的兒子走，好事奉我。你還是不肯放他走。看哪，我要殺你頭生的兒子。』」 \end{tabularx} \\ \\
[1ex]
\hline
\hline
\end{longtable}
$^{1}$各位弟兄姊妹一起讀聖經.
生命樹宣教網絡為大家繼續提供聖經人,聖經事.
看看人事之後上帝的手是怎樣.
今天繼續是摩西的故事.
來到為何要殺長子問號.
為何要殺長子呢.
大家知道在十災裡面.
去到第十災是殺長子之災.
一至九災都沒有容許意識人離開.
於是神就用第十災.
其實如果你留意一下第九災.
我上兩次都講了在一至九災裡面.
其實最後這一災都很重要.
第九災就是黑暗之災.
遍地都黑暗三日黑到.
那些黑好像壓過來一樣.
好像摸得著.
即是很黑伸手不見五指這麼黑.
但是以色列地竟然有陽光.
這麼厲害的神蹟.
但是法祿的心都很硬.
仍然是信自己的能力.
信我會捱得過的.
既然這樣.
去到第十災殺長子之災.
他知道這件事真的不能開玩笑了.
這個神真的應付不了.
這個神不能夠壓倒.
反而被上帝壓倒.
耶和華.
他開始認識了.
為何要殺長子呢.
原因就是這裡的.
可能都不是很多人會留意到.
跟大家解釋一下.
回到柴契記第四章.
最初在米甸抗爺第十九節.
耶和華在米甸抗爺對摩西說.
我第一點要說的紅色螢光.
你要回埃及去.

$^{41}$在這個地方神就差遣摩西回到埃及.
其實一早在荊棘被火燒著的時候.
已經叫他回去.
但是未正式出發.
這個時候四章第十九節.
叫他出發.
叫摩西你要回埃及去.
回去時機到了.
因為尋索你命的人都死了.
尋索你命的人是哪些人呢.
就是當時的法老王.
就是說摩西殺了人逃跑.
當時的法老就想殺摩西.
法老想殺摩西是因為他殺了人.
所以已經通令全國.
哪些哪些人給我留意.
不要讓摩西回來.
一回來就抓他.
所以他有些手下的大臣.
某些人都知道.
那些人都全部是尋索你命的人.
都死了.
你知道當時摩西多少歲嗎.
已經八十歲了.
那時候已經是一個相當老的人.
而尋索他命的人都過世了.
耶和華神就說你回埃及去.
現在到時候回去.
摩西就順命了.
帶著妻子和兩個兒子.
叫他們騎上路.
一起回埃及地去了.
摩西手裡拿著上帝的像.
他回埃及的時候拿著上帝的像.
帶錢還重要嗎.
重要.
行李重要.
最重要摩西要帶著上帝的像.
其實那些像不是上帝的像.
不是上帝給他一支.

$^{81}$天空掉下來.
是他用那一支在米田曠野.
用來行路動動下.
又用來牧羊的那支像.
其實是他自己的.
不過這支像現在叫做上帝的像.
因為上帝是透過這支像行神蹟.
掉到地下變成蛇.
然後拿起尾部抓起.
又變回一支像.
上帝的像.
神蹟同在的證據.
所以第一點要跟大家看.
其實你會看到.
很有趣.
神有沒有叫我們獨善其身.
信主好.
不如去找一個沙漠.
在那裡長住隱居.
或者去深山野嶺住.
在那裡離開這個世界.
一味敬拜沙漠教父.
沙漠教父不是不對.
他們在曠野專心敬拜上帝禱告.
好的.
不過似乎不是叫摩西這樣做.
叫摩西回去埃及.
你回到那群受苦.
在做奴隸的人當中.
體諒他們.
帶他們出來.
這就是我耶和華神的心意.
是叫他回去埃及.
所以今天.
我們都有一個相同的心智.
就是要回到受苦.
多苦的世界當中.
要帶那些人出來認識上帝.
去到神為他們有應許的地方.
因此我們仍然要.

$^{121}$我們可以說.
不在這個世界當中.
不屬於這個世界.
但是我們要進入這個世界.
回埃及去.
回埃及去不是吃黃瓜韭菜.
不是吃那裡的那些好東西.
不是在那裡享受.
而是帶他們出來認識上帝.
手裡面.
如果你要回埃及.
必須要手裡面拿著上帝的像.
否則你拿著自己的像.
很快就會被埃及.
被這個世界同化了.
融化了你.
你要帶著上帝的權能回去.
才能夠回去埃及做點事.
好了.
第二十一節的上半節.
耶和華就跟摩西說.
你回去吧.
你回到埃及的時候.
第二點.
要留意.
將我指示你的一切歧視.
行在法老面前.
讀多一次.
你回到埃及的時候.
要留意.
將我指示你的一切歧視.
行在法老面前.
要留意.
回去不要說.
哎呀 好久不見.
老友 老友.
哎呀.
先吃飯.
哎呀 我家人阿倫在這裡.
做奴隸這麼多年.

$^{161}$現在怎樣.
還很辛苦.
不是那些.
是回去你要留意.
留意這個字.
Raar.
Raar就是汗.
你回到埃及的時候.
要看我指示的一切歧視.
這個「將」.
你要留意.
將我指示你的一切歧視.
這個「將」字.
藍色的「你」字.
其實原文是沒有的.
所以你回到埃及的時候.
就要看我指示的一切歧視.
其中有一個字是沒有的.
這個指示就是Zoom放在.
或者交給.
然後Yard.
Zoom Yard.
Zoom Yard是手.
Zoom是放在.
其實如果簡單這麼說.
原文意思就是.
你要看我放在你手的一切歧視.
原來神是將一些能力.
放在我們的手裡面.
放在摩西的手裡面.
其實那支「將」.
都是普通的木棍而已.
不過有神將他的能力.
放在那裡.
那就厲害了.
放在摩西的手裡面.
那摩西的手.
那支「將」才有能力.
不是.
那支「將」叫做.

$^{201}$神上帝的「將」.
其實有能力.
其實是神將能力.
給了摩西的手.
所以摩西拿支「將」.
丟在地上.
就可以變成「蛇」.
拿起來變成「棍」.
如果他交給亞倫.
亞倫丟在地上.
變成「蛇」.
拿起來變成「棍」.
但是如果再交給第三個人.
丟在地上.
也是一支「棍」.
拿起來也是一支「棍」.
其實那個能力.
不在那支「將」那裡.
而是在摩西的手那裡.
或者他的能力.
要交給哪一個人.
上帝透過他.
傳遞給亞倫.
這樣也可以.
但是隨便找一個.
二打六過來.
拿著他的「棍」.
丟在地上.
不會變成「蛇」的.
因為其實能力.
不在那支「將」那裡.
上帝的「將」是一個象徵.
是在摩西的手那裡.
好了.
你翻開它的時候.
要留意的.
要留意的是「看」.
要看我放在.
你手的一切其事.
要留意.

$^{241}$要看放在手的一切其事.
甚至沒有這個「理」字.
Zoom up.
就是放在手.
留意一下.
將我放在你手的.
一切能力其事.
行在法律面前.
就是「手」.
於是我上次說的就是.
哎呀 它向天舉象.
哇 就黑暗了.
其實這個向天伸出的象.
那支「將」是象徵於神的能力.
但是其實能力在「手」那裡.
大家要留意.
其實叫我回去埃及.
重返這個世界.
哎呀 這個世界這麼多.
人有這麼多邪惡的力量.
我哪裡能夠勝過.
不用怕.
最重要是有神的能力.
和你同在.
象徵物件就是那支「將」.
不過最重要的是你的「手」.
你記得要看.
要看我放在你手的那些歧視.
走出來.
所以其實你要知道.
神叫我們回去為祂做一些事.
叫我們去為祂做一些事.
祂是有能力放在我們手上的.
當祂差遣一個人的時候.
祂是會給能力.
給相應的能力給祂的.
給祂的同在給祂的.
如果沒有這個能力.
可以說是神的同在.
如果沒有神的同在.

$^{281}$是沒有那些能力的.
因此你要知道.
神有差遣.
同時祂有能力.
而這個能力最重要的是祂的同在.
第二點要看的就是.
你看看.
神將很多能力放在你的手上.
然後呢.
但是我要.
但是21節的下半節.
祂就補充.
雖然有這個能力.
不過我要使祂的心肝硬.
必不容百姓去.
你的手不會有能力.
不過我會令到.
法律的心肝硬.
從這裡就要留意了.
那就是說了.
你令到祂的心肝硬.
然後你又給能力摩西行.
神即其事.
然後殺祂長子.
又釋祂的燈黑暗之災三天.
然後祂不肯放你.
你又說為什麼你不肯放.
其實呢.
原來是你背後使祂的心肝硬.
我要解釋.
這個使呢.
其實意思是容許.
讓.
容讓祂的心肝硬.
任憑或者說.
這個「讓」或者是「任憑」.
在聖經出過很多次.
神是會任憑一個人.
令到他的心離開神.
讓一個人的心肝硬.

$^{321}$其實神可以一變.
祂不肝硬就可以了.
但其實祂本身肝硬.
如果神有你肝硬.
有你背道.
其實都是你自己的責任.
不過如果神有很大的憐憫.
才會拉你回頭.
例如舉個例子.
主耶穌對著出賣祂的猶大.
說今天今晚有個人要賣我.
是不是我是不是我.
甚至猶大都問是不是我.
主耶穌就說.
跟我一起站守在盤子裡.
就是坐我旁邊那個.
而且是坐高位的.
坐在他右手邊.
是守直客人的位置.
那才能夠跟他一起伸手到近.
到可以拿盤裡面的餅.
就是了.
猶大就說是不是我.
那主耶穌就說.
你要做的你就儘管做吧.
那祂就站起來出去賣耶穌了.
那祂的心是肝硬的.
其實這裡主耶穌都給一個機會猶大.
哎呀 那你都知道.
其實主耶穌我真的很壞.
就為了三十兩銀紙.
一年的供價我竟然叫人賣你.
那現在怎麼辦呢.
沒有啊 好啊.
站起來就出去賣耶穌.
是祂的心肝硬了.
就是讓祂的心肝硬.
那責任不在神那裡.
而是祂的心自己肝硬.
那就讓你吧.

$^{361}$譬如說那個爸爸.
有一個兒子 浪子.
爸爸 你其實還沒死.
不過我知道你有一天始終會死.
死了你就分家產.
我和哥哥兩兄弟.
那你當然是分一半給我的.
或者分某部分給我.
或者哥哥佔兩份.
我佔一份.
他是長子嘛.
將我應得那份給我.
那然後拿了就走了.
兒子不如你不要走了.
跟我一起住.
出去很艱難的.
在家千日好啊.
老爺 我管你你傻.
我現在就是要走.
所以你快點給錢.
就當我預先拿你的遺產.
哇 我真的給你一個五.
你應得的只有一個五.
不過好吧 給他吧.
任憑讓他一個很肝硬的心走.
不過後來他不肝硬.
他就回轉了.
是我要讓他的心肝硬.
他不會讓我的百姓走的.
因為我要他知道我.
好了 第二節.
你要對法老說.
耶和華這樣說.
以色列是我的兒子.
留意一下.
這個第四章講了.
以色列是我的兒子.
一去到的時候.
就已經跟法老說了.
以色列是我的兒子.

$^{401}$是我的長子.
第三個紅色字要留意.
為什麼要殺他的長子呢.
是因為以色列是我的兒子.
而且是我的長子.
長子 第一個.
我都說過了.
用我的兒子去好事奉我.
你還是不肯用他去看.
我就要殺你的長子.
很少人會留意.
看到第四章第23節.
為什麼要殺長子之災.
這麼厲害呢.
原因就是.
以色列是我的長子.
跟你說了.
你要讓他走.
如果你不讓他走.
你還是不讓他走.
這樣你記住.
因為你扣留著我的長子.
我就要殺你的長子.
那我保持著他.
不用殺我的長子.
這麼厲害.
其實在你那裡做奴隸.
做了430年奴隸.
痛不痛苦.
極度痛苦.
為什麼不停去糟蹋他.
在那裡做奴隸.
沒有希望 沒有人生.
其實比死更難受.
所以如果你還不放.
我就殺你的長子.
其實我那些430年.
我的長子.
一個一個都死在你那裡.
死在你的威權.

$^{441}$死在你的壓力底下.
在埃及430年遺老.
這樣就很慘.
做奴隸.
現在你要放他.
如果不放他.
我就要殺你的長子.
我會對付你.
這個就是殺長子的原因.
去到最後竟然九個災.
第九個黑暗之災.
這麼厲害.
你都不接受.
都不肯讓的話.
好吧.
於是我就要執行.
我跟你說的4章23節.
看啊.
我要殺你的長子.
第一次.
很多人沒有留意到.
第十災厲害.
殺長子厲害.
為什麼呢.
是因為.
神在第四章叫摩西.
你去到就跟法老說.
跟法老說什麼.
以色列是我的兒子.
是我的長子.
你要放他走.
如果你430年都不肯放.
我就會殺你的長子.
他也是心肝硬.
不肯.
到最後.
神就做這個工作了.
於是長子被殺.
今次我跟大家分享.
為什麼會殺長子.

$^{481}$不過有兩點.
令我們特別要留意的.
就是原來神是差我們.
回到這個世界當中.
雖然我們不屬於這個世界.
但是我們是入世的.
我們的心是出世的.
但是我們的人是入世的.
在這個世界裡面.
我們拿著上帝的能力.
我們手滿有上帝的能力.
我們在那裡.
是要看一下.
神指示加些什麼.
在我的手裡面.
去幫助有需要的人.
神是有公在有同在的.
而到最後.
神是有祂自己審判的作為.
就是會做祂的工作的.
我們一起祈禱吧.
天父讚美你.
你自己是一位大能的主.
去到這裡.
講到殺長子的時候.
令我們記起.
頭的九個災.
滿有大能.
第十個災是更大的能力.
而當中.
你是開了出路.
為人預備了愉悅的血.
只不過法老的心剛硬.
主要幫助我們.
讓我們跟隨你的時候.
我們小心.
免得我們的心.
都好像法老一樣剛硬.
任憑我們剛硬的時候.
我們就做錯了.

$^{521}$我們就會去到很慘的境地.
主要幫助我們.
你妻我們去到這個世界.
幫助有需要的人.
被困的人.
在做罪惡中奴隸的人.
我們帶他出來.
憑著你給我們的恩典.
你給我們的同在.
以及你的象徵.
那支有能力的象.
將人帶到你面前.
幫助我們不會學法老.
剛硬到底.
讓我們有一個.
對你是更柔軟.
信服又聽從的心.
隨天禱告.
奉耶穌命求.
阿姆!.
\newpage



\section{帖撒羅尼迦前書 5:16-19}
\label{sec:031938}
\textbf{復活:恢復感恩的心}~~~~謝又生牧師
\newline
\newline
連結: \href{http://www.church.com.hk/acms/content.asp?site=cdc\&op=show\&type=product\&code=031938\&layout=sermon}{\texttt{031938}} ~~~~ 語音日期: 2026-04-05
\newline
\newline
\hyperref[sec:031937]{\small{< < < PREV SERMON < < <}}
~
\hyperref[sec:index_chronic]{\small{[返順時目]}}
~
\hyperref[sec:031939]{\small{> > > NEXT SERMON > > >}}
\newline
\newline
帖撒羅尼迦前書 5:16-19
\newline
\begin{longtable}{cl}
\hline
\hline
章節 & 經文 (和合本修訂版)\\
\hline
5:16 & \begin{tabularx}{0.7\textwidth}{X} 要常常喜樂， \end{tabularx} \\ \\ \relax
5:17 & \begin{tabularx}{0.7\textwidth}{X} 不住地禱告， \end{tabularx} \\ \\ \relax
5:18 & \begin{tabularx}{0.7\textwidth}{X} 凡事謝恩，因為這是神在基督耶穌裡向你們所定的旨意。 \end{tabularx} \\ \\ \relax
5:19 & \begin{tabularx}{0.7\textwidth}{X} 不要熄滅聖靈； \end{tabularx} \\ \\
[1ex]
\hline
\hline
\end{longtable}
$^{1}$著名的科學家霍金.
有一次在一個學術研討會上的演講.
當他演講完畢的時候.
有一個年輕的記者走到霍金面前問他.
霍金先生.
你的疾病將你永遠固定在輪椅上.
難道你不會因為失去了這麼多東西而悲哀嗎.
問了這些問題.
大家都不知道怎麼辦.
所以會場一片的靜寂.
但是霍金面上掛著笑容.
用他不是太靈活的手指.
一直按著他前面的鍵盤的鍵.
然後螢光幕上就出現以下的字.
霍金說我的手指還能夠活動.
我的大腦還能夠思考.
我有終身追求的理想.
我也有我所愛和愛我的親人和朋友.
最重要的就是我還有一顆感人的心.
就在這個寂靜的環境裡面.
這個年輕的記者心靈被震盪.
他恍恍學到一個人生的很重要的課題.
是什麼呢.
就是感恩.
以下的照片是在我的網上找到的.
它裡面的意思是.
請用一個詞語去描述.
現今世界需要些什麼.
當然問不同的人有不同的答案.
但是我覺得.
現今的世界最需要的是什麼呢.
感恩的心.
有人說.
感恩的心是基督徒一切美好品格的開始.
我就想一想.
其實是真的.
如果沒有這個感恩的心.
我們今天所說的基督徒的品格.
我們稱之為「聖靈果子」.
例如人愛,喜樂,和平.

$^{41}$又是怎樣來的呢.
其實我很相信.
我們的神很在意祂的指紋.
就是你和我.
我們有沒有感恩的心.
如果大家記得.
《路加福音》有這樣的故事.
耶穌治療了十個麻瘋病的病人.
但是結果.
這十個病人去了哪裡呢.
只有一個撒瑪利亞人.
走回來耶穌的面前大聲地讚美神.
並且苦伏在耶穌的腳前.
向祂道謝.
接著耶穌就不禁問.
大哥.
我治療了十個人.
其他九個人去了哪裡呢.
除了這個撒瑪利亞人之外.
就沒有人回來道謝.
所以神的心是很想我們有一個感恩的心.
神的心意是很堅定和一貫的.
你看看《舊約聖經》,《新約聖經》.
耶穌的信息是很清楚的.
在舊約時代.
神已經開始要祂的指紋去操練這樣的心.
所以耶穌說.
你們不要只是口裡說多謝.
不要只是說.
你在行動上也要參與.
配合一些很清楚的儀式去獻祭.
所以舊約的人.
是會請祭司為他們獻上平安祭.
那些是感恩祭.
這些獻祭是做什麼呢?.
向神表達感恩.
舊約除了要獻祭表達感恩之外.
另一種表達多謝神,稱謝神的方法.
就是歌頌讚美神.
就好像我們開聲讚美神一樣.

$^{81}$如果你看舊約的詩篇.
是充滿感恩的詩篇.
感恩的詩歌.
其中一個代表作就是我剛才讀給大家聽的.
詩篇一百零三篇.
大衛經常都說.
我的心你要稱眾耶和華.
所以舊約是充滿這些感恩的.
好,去到新約.
你說沒有的.
沒有獻祭的.
在新約時代.
我們是不需要再依循這些儀式去獻祭.
但是在新約聖經裡面.
有一個很重要的作者就是仕途保羅.
他不斷地說.
舉個例子.
在《歌羅西書》和《以忽所書》.
保羅用幾乎同一樣的字眼說什麼呢?.
他用詩章,聖詩,靈歌,懷著.
感恩的心.
歌頌神.
接著不僅是這樣.
凡是奉我們主耶穌基督的命.
說完又說.
好了.
不僅是這兩個書卷.
他在《帖撒諾黎加前書》.
還說得更清楚.
《帖撒諾黎加前書》五章.
大兄姐妹.
《帖撒諾黎加前書》五章.
十六至十九節是今天的主要經文.
讓我們一起讀這個主要經文.
一二三.
「要常常喜樂,不住的禱告,凡是謝恩」.
「因為這是神在基督耶穌裡向你們所定的旨意」.
「不要消滅聖靈的感動」.
這段聖經可能大家都很熟悉.
有些弟兄姊妹識貝.

$^{121}$通常祝福人.
要敲這段聖經.
又寫給他.
很熟悉的.
但是有一句.
可能跟大家都消化不下的.
蓋不下的.
是什麼呢?.
「凡是謝恩」.
「凡是謝恩」通常大家會聯想.
「凡是謝恩」就是.
什麼都謝恩.
你想想.
好東西當然是謝恩.
不好東西怎麼謝恩呢?.
今天這麼多戰爭.
一旦有戰爭就很多東西被摧毀.
很多人命失去.
這些都謝恩?.
或者經濟不好.
有人失業.
不會這樣的.
可能在同學那裡.
同學考試肥佬.
因為你考試肥佬.
我謝恩.
不會這樣的.
所以這些你不能強行謝恩.
明明是不好的東西.
明明是一些不幸的東西.
你不可能謝恩.
你是謝恩不下的.
但是弟兄姊妹.
神不是要我們這樣.
神的心腸不是.
你強行承受了那些謝恩.
硬生生的.
不是.
其實「凡是謝恩」的英文.
可能大家知道了.

$^{161}$Give thanks in all circumstances.
其實「凡是謝恩」的原文的意思是什麼呢?.
它是說.
無論任何的境況.
人生有些逆境.
有些順境.
但是無論什麼境況.
我們都要有這種心態.
我們學習去帶著感恩的心.
會說.
這個不是你強行為一些壞的東西來感恩.
而是你帶著感恩的心.
去見到更加多的東西.
所以這種是一種人生觀.
是一種生活的態度.
是要學習的.
既然是人生觀.
既然是生活的態度.
是要學習.
是要操練.
是要實踐的.
那究竟神為什麼要我們操練感恩的心呢?.
究竟操練感恩的心.
有什麼重要性呢?.
神在這一方面的終極心意是什麼呢?.
弟兄姊妹.
我很相信.
神要我們做的事.
不會無聊的.
神要我們做的事.
一定對我們有益處.
所以在以下時間.
我很想跟大家說.
操練感恩的心.
這種的生活態度.
究竟對我們有什麼益處?.
是有兩個益處的.
第一個益處.
就是你操練感恩的心.
是會幫助你去管理好你的思想.

$^{201}$我剛才貼上來的前書裡.
它說「常常喜樂,不住的禱告,凡事謝恩」.
其實為什麼神會將這三句放在一起呢?.
原文是並列的.
很密一密的.
為什麼呢?.
有它的原因.
為什麼呢?.
因為這三樣是有密切的關係的.
剛才「凡事謝恩」我解釋過了.
「不住的禱告」不是像字面那樣解釋的.
在字面上我們就說.
你經常逛街,嘴巴張開祈禱.
或者你躲在密室裡七二四祈禱.
是重要的.
不過這不是最重要的地方.
重點在哪裡呢?.
重點就是.
你常常意識到神的臨在.
常常敏銳聲明的臨在.
常常無論什麼事都會去找回神.
其實三言兩語說.
就是你那種關係.
和神的關係,常常敏銳到神在.
喜樂.
喜樂不是強行扮個樣子出來.
我明明不開心的.
扮個喜樂出來.
喜樂是你心靈裡面那種快樂感.
也是我們說的聖靈的果子.
未必是扮出來的.
你是心裡面那種感受,快樂感.
其實一個人.
一個懂得感恩的人.
不是所有事情都順利的.
當所有事情順利.
所有發生的事情沒有困難.
那你自然的想法當然要感恩.
一個懂得感恩的人.
不只是你順利才感恩.

$^{241}$一個懂得感恩的人.
他不是沒有困難的.
而是他有那種能力.
去管理好他的思想和情緒.
等於如果你不管理好你的思想情緒.
你就會慢慢跌入那種負面的循環.
等於那種負面的循環,負面的情緒.
你不管理它,就由它管理你.
所以一個懂得感恩的人.
是管理好自己的思想和情緒.
這是聖經教我們的.
他說心言說你要保守你的心.
勝過保守一切.
因為一生的果效是由心發出.
保守你的心就是管理好你的思想和情緒.
如果不是,它就管理你.
如果負面的情緒管理你.
你一生的果效就是這樣.
如果你管理得好.
你一生的果效就是喜樂.
這是神教我們的.
怎樣管理呢?.
等於可能大家都有機會遇到逆境.
遇到逆境的時候.
大家其實是有選擇的.
有兩個選擇.
第一個選擇是你不斷地埋怨.
不斷地埋怨的結果是怎樣呢?.
不斷地埋怨的意思就是說.
這個壞,那個又壞.
那個環境又不好.
為什麼那些環境對他好,對我不好.
為什麼神代對他好.
祝福他,又不祝福我.
可能慢慢地你會發覺.
當你當成很多的self talk的時候.
你就會覺得神代薄了你,欠了你.
慢慢地這個人的人生.
就變成了負能量的人.
弟子們有沒有見過負能量的人?.

$^{281}$你害怕嗎?.
其實你想想.
其實你很害怕負能量的人.
一旦靠近你,你就變成負能量.
不斷地說.
你看到很多埋怨是有點害怕的.
負能量的人不單止影響別人.
其實他自己也很辛苦.
是不是?.
你不斷地被那些說話.
自己類類的self talk.
不斷地convince你.
我多麼差,多麼慘.
結果他一生就是被這些苦毒.
這些埋怨去折磨.
其實負能量的人是很辛苦的.
你想過這一個生活.
要麼是你選擇另一個途徑.
就是相信真理.
真理是甚麼?.
神說我們有得救的.
我會拯救你的.
我會永遠與你同在的.
你這個困難是有終結的.
是會有盡頭的.
你會提升的.
我會為你開新路的.
這些就是真理.
你選擇這些的盼望.
你選擇神的真理的思想.
用來做甚麼?.
去代替一些東西.
代替一些絕望的感受.
代替撒旦的謊言.
大家有沒有聽過撒旦的謊言?.
或多或少都有的.
你腦裡如果出現的是甚麼呢?.
是你沒得救的.
是你最差的.
每個人都比你好.

$^{321}$你不會有出路的.
你永遠都沒有盼望.
這些是撒旦的謊言.
神的心意是甚麼?.
讓這些神的真理代替.
踢走這些撒旦的謊言.
慢慢你就會有輕省的心.
頂多我現在說的這些真理.
並不是德國說的.
是有些實際上.
有些大學機構.
實際上有研究證明.
這個真是對的.
美國哈佛大學.
曾經做過一個這樣的研究.
他們請了一組人.
然後請他們在兩分鐘之內.
輪流這樣說感應的事.
這幫人就照做.
在這個過程裡面.
那些研究員.
記錄他們的反應.
研究員有四個發言.
第一個發言.
就是原來感恩這個行動.
可以迅速活化你的大腦.
第二個原來不斷的感恩.
可以減少人的焦慮.
第三個可能大家想不到.
可能有些睡得不好的大師妹.
失眠的大師妹.
原來不斷的感恩.
可以改善你的睡眠.
還有最重要的是甚麼?.
第四個.
原來不斷的感恩.
是在內心裡面產生一種快樂感.
這個不只是哈佛大學的實驗證明.
這個發言和當時另外一間大學.
加州大學做的差不多的實驗.

$^{361}$都一樣證明到.
所以有專家這樣說.
感恩.
操練感恩的心.
原來帶給我們很多好處.
有三大好處.
第一個好處.
原來不斷操練感恩的心.
是幫助你對抗這些負面.
以及消極的情緒.
如果你有一天發覺自己的情緒.
很負面很消極的時候.
你嘗試感恩.
列出感恩的事.
操練感恩的事.
第二件事.
就是你操練感恩的心.
慢慢培養一種習慣.
這習慣是甚麼呢?.
是發現你身邊原來有這麼多美好的東西.
因為你平時看不到的.
有習慣可以幫助你發現更多.
第三.
那個好處就是當你操練感恩的心的時候.
你會發覺.
你慢慢進步了.
你不會讓那些你現在手頭上擁有的東西.
是你所當的.
我們經常都說.
take it for granted.
你不會.
這些你現在手上的東西當沒有了.
原來你發現有美好的東西在你身邊.
大英姊妹.
這樣的人.
你會慢慢產生一種快樂感.
這種是沒有辦法形容的.
當你經歷到.
當你真的參與行動經歷的時候.
你就會慢慢發現這種快樂感.

$^{401}$所以今天的結論我鼓勵大家.
大英姊妹讓我們一起操練感恩的心.
在聖經裡面有些暗示.
我們剛才說上上起落.
不住地禱告.
凡事的謝恩.
其實它是在暗示一些東西.
它在說什麼呢.
上上不住地凡事.
在說什麼呢.
不斷地重複.
不斷地繼續一直下去.
其實它在說什麼呢.
操練.
所以神的心是鼓勵我們操練感恩的心.
操練的重點是什麼呢.
不斷地行切 不要放棄 繼續.
最近我們教會有一位弟兄.
想出自傳.
想我幫他寫個罪.
你知道幫人寫罪.
你不能夠老作的.
所以你要看完他的自傳.
看完他的書.
但是當我一邊看這位弟兄的成長.
他的故事的時候.
我發現一個特點.
就是他那種歇而不捨的操練.
他那種不屈不撓的精神.
等於我們今天這個世代.
我們如果要操練一些好的品格.
我們必須要有這東西.
行切 不屈不撓.
所以這位弟兄.
他有這樣的生命的特點.
就讓他無論在事業.
或者在興趣上面.
都有卓越的成就.
我想如果我們操練感恩的心.
我們是很需要這個行節.

$^{441}$很需要不斷.
所以今天我鼓勵大家.
和我一起有以下的練習.
是包括什麼呢.
以下的練習.
是包括我們試試每一天.
麻煩另外一張.
我們每一天試試記錄三件值得感恩的事.
可能有一些弟兄姊妹試過了.
可能你們小組也試過玩這些.
但是頂多我們試試一個行節的練習.
三件感恩的事未必會發現大事.
你知不知道有些細微的事.
都可以感恩的.
譬如你回崇拜.
你打算天黑了.
如果下雨我就連鞋子都濕了.
就算戴雨傘.
褲子都濕了.
昨天就是這樣.
昨天三點鐘崇拜.
到了一點多就傾門賴雨.
我寫在心裡整個天都黑了.
我心想主啊求你不要下大雨.
我都信心少.
不要下大雨.
讓頂多姐妹回來崇拜順利.
誰知兩點鐘下雨了.
一滴都沒有.
所以頂多姐妹你回來的時候有沒有下大雨.
沒有.
頂多姐妹這些不是一些很厲害的事.
是一些日常生活的事.
你可以值得感恩的.
幾微小的事都值得感恩.
另外呢第二樣東西.
向你的同事.
你回到學校向你的同學.
向你的家人.
向你的朋友.

$^{481}$你試試多謝他.
頂多姐妹多謝別人是否需要原因的.
不用的你試試.
今天回去.
你走去老公.
走去你老婆.
老婆說我多謝你.
接著你老婆說你為什麼要多謝我.
老公說不用原因多謝你.
因為你是我老婆.
所以我多謝你行不行.
你試試給人多謝.
給人多謝的人是否開心.
多謝別人的人是否開心.
頂多姐妹這是一個美好的互動.
你試試這樣.
當然你要找找原因.
你要發現原因.
一定會有東西令你多謝他.
第三樣東西.
大家都會做的.
我相信大家都會顧全你的後腔清潔.
就是你晚上臨睡前.
你刷牙後對著自己的鏡子.
你想一想那天.
有什麼我做得好的呢.
又或者那天我有些表現.
我是欣賞自己的.
我今天明明是很生氣的那個人.
但是我可以克服自己.
我今天很慷慨給了我一點點貼士.
這些都是值得欣賞自己的.
頂多姐妹我試試.
我們試試這樣.
你可以每天做完三樣東西.
又或者做一兩樣東西.
但是頂多姐妹.
記住操練的重點是行節.
剛才有些頂多姐妹會拍照.
其實你不用拍照.

$^{521}$因為待會大家結束崇拜.
會有兩張硬紙派給大家.
雖然設計不同.
但是都是一樣的.
所以待會結束的時候.
大家離開的時候.
我們歡迎隊的頂多姐妹會派一張給大家.
那你大家可以拿一張.
記住不要懷著一個.
今天正是牧師叫我拿的.
那就拿吧.
接著走到街口就舉手.
扔在垃圾桶.
待會離開的時候.
順便在垃圾桶裡裝一裝.
開玩笑的.
頂多姐妹其實這個操練很重要的.
是幫助我們有一個感恩心.
其實神的心意很清楚的.
祂很希望我們有一個感恩心.
當我們操練感恩心的時候.
會有喜樂.
其實復活節是什麼.
復活節是 restoration.
恢復的日子.
恢復什麼呢.
恢復我們跟神的關係.
以及我們未犯罪之前.
我們應有的平安和喜樂.
神的心是很想恢復這種的喜樂.
頂多姐妹我再問大家一次.
你想不想有一個喜樂的人生?.
想.
你記住這個是神的心意.
不只是我們自己想.
是神也想.
神很想我們有一個喜樂的人生.
所以祂鼓勵我們去操練感恩的心.
這是第一個重點.
第二個好處是有關屬靈的問題.

$^{561}$神不只是叫我們自己操練感恩的心.
神是會幫我們的.
會幫助我們有感恩的心.
這是屬靈上神要幫助我們的.
我們今天.
你環顧這個世界.
雖然我們很想有感恩的心.
但是你環顧這個世界.
這個世界喜不喜樂?.
不喜樂的.
是不是?.
周圍的人.
你見到很多人在街上走.
愁眉苦臉的.
這個世界不是喜樂的.
為什麼呢?.
因為殺彈偷了我們的東西.
在約翰福音他講了一件事.
他說「盜賊來」.
無非是要偷竊,殺害和毀壞.
那個盜賊是誰?.
就是殺彈.
殺彈偷了我們的什麼?.
不是偷雞.
偷了我們的喜樂.
頂子妹.
如果有車的話.
頂子妹要不把窗戶弄好.
就殺彈伸手來偷了你的喜樂.
這是心靈的窗.
所以你想想.
可能有些人會問自己喜不喜樂?.
不是很喜樂的.
這個世界環顧我們身邊.
都不是很喜樂的.
因為殺彈偷竊了我們的東西.
但是頂子妹.
復活節是怎樣的?.
Vestoration.
神要恢復我們本來應該有的喜樂.

$^{601}$這是神的心意.
所以我們講的貼薩諾黎加前書.
他有一句.
他說什麼?.
「這是神在基督耶穌裡向你們所定的旨意」.
這是神的意思.
不只是我們想有喜樂的人生.
是神也想你有喜樂的人生.
所以我們要跟神配合.
我剛才說我們主動去操練感恩的心.
是我們要做的.
但是我們要靠著神的能力.
為什麼我們要靠著神?.
有兩個原因.
第一個原因是.
我們本來都不懂得感恩的.
你同意嗎?.
你覺得人是埋怨底還是感恩底?.
埋怨底嘛.
這個大家都明白.
我也承認.
人在埋怨底的人本來是不懂得感恩的.
我有一次去一間大百貨公司的超級市場.
本來是去買東西的.
但是那個超級市場.
剛剛有一個推廣員推廣一種新的食品.
是朱古力蛋糕.
很好吃.
大家有沒有聞過新鮮朱古力的香味?.
那個推廣員的這個產品很好.
我忘記了那個牌子.
但是那些新鮮朱古力的香味.
充滿了大部份的超級市場.
搞得我們為之心動.
還有試吃.
我很想過去試吃一件.
但是那個小妹妹比我更快.
她一聞到朱古力香味.
就眼前一亮.
接著一個漸步.

$^{641}$不知道為什麼這麼快.
飛過去拿了一件正在吃的.
她旁邊的媽媽就跟她說.
妹妹 應該說什麼?.
大家 你覺得她會說什麼?.
謝謝 是不是?.
我也想著 謝謝.
就好像剛才小妹妹那樣.
謝謝.
不是的.
妹妹想了一會.
她跟那個推廣員說.
可不可以多給我一件?.
我們本來都不懂得感恩的 是不是?.
我們的底子是這樣.
可能人的麻冤底是多過感恩底.
再加上什麼呢?.
我們身處這個世界.
我們身處這個環境.
我們這個社會.
好 我想問大家兩個問題.
你打開手機 你看資訊.
你覺得這些資訊和新聞.
是負面新聞多還是正面新聞多?.
負面嘛 是不是?.
你記住這個世界是市場主導的.
發佈這些資訊.
報告這些新聞的人.
覺得大多數人都是喜歡聽負面的新聞.
負面的資訊的.
那他就會這樣出來了.
今天沒有了《水果日報》.
當年《水果日報》很多負面新聞.
很多負面的報導.
那我問你.
你明知道有些是假的.
有些是負面的.
為什麼你會看呢?.
那他就說.
我要看看它怎麼負面.

$^{681}$就是這樣.
這個世界是這樣的.
那所以呢.
這個是市場主導的.
所以呢.
再加上什麼呢?.
今天大家知道AI多厲害 是不是?.
如果 AI 知道 Emily 喜歡看負面新聞.
你就大劑了.
那些負面新聞.
就會鋪天蓋地地轟炸你.
一個人你不斷地看負面資訊.
負面新聞.
你那個人會越看就越負面.
越看就越 Upset.
越看就越心悶.
越看就越覺得這個世界沒有希望.
所有的人都這麼壞.
所有人都對我不好.
所有這個世界沒有出路.
丁姐妹.
你記住.
這個是撒旦的目的.
你記住.
你記不記得.
偷竊了你的喜樂.
你沒有了喜樂.
沒有了盼望的時候.
你就會覺得神不真實.
神對你不好.
你就會離開神.
這個就是撒旦的目的.
所以不要中撒旦的詭計.
這是第一個問題.
第二個問題.
你覺得這個社會.
多數喜歡稱讚別人.
還是喜歡暴露別人的呢?.
咦? 大家怎麼不回答?.
暴露別人是吧?.

$^{721}$這裡也挺一致的.
因為我們大家都在這個社會生活.
這個社會喜歡暴露別人.
多於稱讚別人的.
有一個詞語.
是我以往的體會.
跟現在的體會是不同的.
就是說這個詞語.
以往我知道的解釋.
跟現在的新解釋是不同的.
是什麼呢?.
反思.
可能大家跟我一樣.
我以前覺得反思.
是反省我自己.
發覺有些什麼我需要改變的.
我就改變.
這是反思.
大家都知道.
但是現在的反思是不同的.
現在反思有失的定義.
就是我幫你反思.
就是什麼?.
暴露別人.
等於我們身處這樣的世界.
我們身處這樣的客觀環境.
其實或多或少.
我們都會被影響的.
再加上我們的底本來不懂得感恩.
再加上這個世界是這樣.
所以為什麼我們需要神呢?.
不只是靠我們自己.
我們自己搞不定.
我們需要神的能力.
所以在聖經裡面說.
除了我們要操練感恩的心之外.
我們需要聖靈的能力幫助我們.
所以我們需要渴望神的聖靈充滿.
剛才我們說的貼紙羅里加前書.
以及《以忽所書》.

$^{761}$都很相似.
他們所做的事.
神的依靠神.
以及結果都很重要.
在《以忽所書》.
他說我們要讓聖靈充滿我們.
而貼紙羅里加前書說.
不要消滅聖靈的感動.
其實都是在說同一件事.
不斷地被聖靈充滿.
而帶來的表現是什麼呢?.
是一樣的.
就是常常感恩.
凡是謝恩.
也是常常感恩.
其實他在說什麼呢?.
他在說我們不斷被聖靈充滿.
讓我們操練感恩的心事半功倍.
當我們不斷操練感恩的心的時候.
我們會有那種喜樂.
這就是神想你和我過的人生.
所以我們是需要渴望聖靈的.
我們需要渴望不斷地被聖靈充滿的.
我們不單只靠自己.
我們必須有這種做法.
所以大家如果說我渴望聖靈.
我很想被聖靈充滿.
稍後我會呼召大家.
到我的左手邊.
代討者為你禱告.
讓聖靈繼續充滿你.
我們要如何去操練感恩的心呢?.
我們必須要行動的.
我們必須要參與的.
我們必須要投入的.
我很想說兩個處境.
第一是在群體裡面.
你參與見到聖靈的工作.
你有感恩的心跟別人分享.
於是大家都會強調大家有感恩的心.

$^{801}$大家會有那種喜樂的心出來.
第二種是個人的經歷.
我先說群體的經驗.
我最近聽到越來越多的大英姊妹在分享.
她在農曆年期間兩場外宴.
她們不約而同在她們的分享.
無論是WhatsApp,口頭的分享.
跟我說.
她們出現了一些共同的字眼.
那是什麼呢?.
感恩,非常感恩,超級感恩,很感恩.
喜樂,很開心,很歡慶,很喜慶,很溫暖.
這些是很開心的字.
為什麼我們參加了農曆年期間兩場外宴?.
同福堂的外宴跟一般教會有一點分別.
一般的教會我們覺得外宴就是.
自己家庭裡面的大英姊妹圓了銅牌.
大家一起吃一頓飯的外宴.
但是同福特色的外宴就不同了.
我們請了一些微信的朋友一起參加.
所以圈子大很多.
這是同福堂特色的外宴.
農曆年的時候有兩場外宴.
我們招待了不同的人.
第一場外宴是年二十三.
在年二十三裡面.
我們招待了一些獨居和商居的長者來吃團圓飯.
那一天.
那些哥哥姐姐穿得很漂亮.
他們覺得我們回去家裡吃團圓飯.
所以很尊重那個場合.
大英姊妹形容這個場合.
她又出現很多字眼.
就是那些溫馨,溫暖,開心.
賓至如歸.
好像一家人那樣的詞語.
那天有七十八位哥哥姐姐.
獨居和商居的長者參加我們的聚會.
當天竟然有一半的哥哥姐姐.
當場決志信耶穌.

$^{841}$這班哥哥姐姐是很笑的.
信耶穌.
當他們吃完之後.
她說「哎呀,又吃又拿,真是不好意思」.
就一直笑著這樣說.
就好像在家裡那樣.
跟家人分享.
第二個聚會的外宴就是春明.
是大年十二.
那個春明招待了一班低收入的家庭.
有四十一個家庭.
一百多位的來賓.
在當中其實一直明白.
在當中的過程裡有很多的感動.
包括什麼呢?.
等於今天香港物質豐盛.
可能大家外出吃飯.
有時候都不是.
好吃就是大家不要介意那些錢.
什麼地方都不要緊.
就算是很好的地方.
可能大家都去得到.
但是你有沒有想過.
有些家庭是從來都沒有去過.
好像黃室堡那些大商場吃飯的.
你有沒有想過.
我們譬如去喝十幾道菜.
十道菜.
有些家庭是從來都沒有吃過這樣的晚飯.
有的我們抽獎.
其實大家都知道.
我們教會不會有一些名貴的抽獎禮品.
但是在當中有些年輕人.
一些家庭的年輕人.
我們抽到一些我們的中獎.
抽到獎.
嘩!開心到彈起來.
我覺得至少這樣的地方是什麼呢?.
就是那班參與的頂智媒在當中.
他親眼看到.

$^{881}$無論那些人.
無論你是什麼背景.
什麼出身.
你的處境景況是怎樣.
他都是一樣被尊重.
他一樣被歡迎.
一樣被接納.
一樣被愛.
頂智媒這一個是具體的愛.
頂智媒一看到的時候.
他們發出感恩的心.
因為原來神的愛具體到臨到人間.
所以他們看到這個情境.
他們會感恩的.
會開心的.
因為他們參與其中.
當中.
大家記不記得.
同福堂有另外一個特點.
我們有祝福閣.
祝福閣就是一些代老者.
有些來賓在酒樓弄了一個祝福閣出來.
他請他們來為他們服侍祈禱.
很多的眼淚.
神就藉著為人禱告.
將祂的愛帶到人的心裡面.
那一天有二十多位新來賓去接受耶穌.
我們想想.
可能是哥哥姐姐.
不是.
很多年輕人.
有些甚至全家.
在台上全家一起走出來.
決志信耶穌.
頂智媒.
這個你不要以為只是兩個聚會.
其實這個是一個過程.
神父照同福堂.
好像我們剛才奉獻的那樣說.
神父照同福堂是要服侍弱勢社群的.

$^{921}$我們這麼多年.
其實有一些弟兄姊妹.
不斷去探訪這些獨居和商居的長者.
認識他們.
慢慢有一種關係和信任.
讓他們經驗到那種愛.
大家知不知道香港還有木屋區.
在瓦宮岩那個地方.
還有一個很大的木屋區.
我在做門徒訓練.
我帶了一班門徒一起去探訪那些木屋區.
有些弟兄姊妹從來都沒有去過.
我們就一起去.
怕老鼠.
大家一起怕.
結果是沒有老鼠的.
我們探訪的時候.
我由頭到尾站著.
由頭探訪站到尾.
因為沒有椅子坐.
那就好了.
他覺得你很愛他.
所以一直探訪他們的時候.
一直那個關係信任就建立了.
所以那個春銘.
那個盤菜宴.
那個團年飯.
其實是神的愛具體的表現.
我很感恩.
因為我覺得蒙福的.
不只是讓我們和福堂有福份去服侍他們.
其實我們有份服侍人是蒙福的.
我覺得最得著祝福的.
未必完全是那班被服侍的朋友.
弟兄姊妹.
而是我們有份服侍的弟兄姊妹.
每一位參與到的.
他都見到神的具體的愛.
見到聖靈的工作.
見到這份愛是具體的.

$^{961}$來到人的生命裡面.
而那些人是很開心的回應.
所以這班弟兄姊妹.
真的經歷到神的具體的工作.
他們因為這樣而感恩.
因為一直這樣感恩他們有喜樂.
人家聽到也要感恩回應.
他們也有喜樂.
弟兄姊妹這班參與的弟兄姊妹是有福的.
弟兄姊妹我很盼望將來.
會不會是你也是其中一個有福呢.
是不是.
這個是我們群體.
所以我們怎樣去操練感恩的心.
是你要參與.
是你要行動.
你投入其中.
你就明白.
這種感恩心這種喜樂是怎樣來的.
第二種就是我們被聖靈更新改變素做.
弟兄姊妹今天我是牧師.
我跟大家說到就好像很有熱忱.
很正面的樣子.
其實我的前半生.
是一個負能量的人.
我以前看東西.
就是先看到不好的東西.
比如說阿宗.
你做一二三.
你就做到那樣東西.
我會問什麼.
做不到又怎樣呢.
你想想那些人會這樣的角度去看東西.
先看到不好的東西.
先看到不好的事情.
如果不行又怎樣呢.
當然你要為他講話.
這個人未有囚謀.
但是你想想底子是負能量.
所以那時候很不開心.

$^{1001}$但是聖靈的工作是很奇妙的.
它會光照你.
它會打開你的心.
它會打開你的眼.
它會讓你看到突然間發現一些東西.
你是以前看不到的.
我就是.
我突然發現什麼呢.
原來神擺在我身邊有這麼多.
關心我愛我的人.
這是一個很大的發現.
我們試一切東西理所當然是很自然的.
但是神給你發現這東西.
不會讓你羞愧.
好像我這樣.
我發現這東西不會讓我羞愧.
自稱那時候懷疑了這麼多東西.
不是的.
反而很開心的.
從那時候開始.
我就每天學這個感恩.
一開始不習慣.
但是一直發現感恩.
因為我有一個好老婆.
我有師母.
她接納我.
聽我說話.
不離不棄.
我有好的教會.
讓我們一起去敬拜神.
我在教會裡面成長.
我有組長.
我那時候不是事奉的.
我不是牧者.
我也有好的組長.
有好的弟兄姊妹同行.
我有福份去祈禱弟兄姊妹.
你們知不知道.
祈禱都是恩典來的.
我們能夠祈禱都是恩典來的.

$^{1041}$我有份參與禱告祭壇.
祈禱會.
可以和弟兄姊妹一起去祈禱.
一起敬禮神.
因為聖靈在我心裡面工作.
我有機會去拾一奉獻.
這個也是恩典.
我有機會.
我有能力去奉獻.
因為神應許.
因為我經歷到那個應許.
我有信心.
當我受到奉獻的時候.
神會祝福我.
這個是恩典來的.
有很多人看不到的.
因為我有你們.
我有弟兄姊妹一起.
即使人生有很多困難.
但是神也幫助我勝過.
在人生裡面有很多困難.
我看到原來神一直在帶領著我.
一直感恩我.
原來有很多感恩.
我也發現一件事.
當你越感恩.
你就發覺越多事情可以感恩.
我再說一次.
當你越感恩.
你就發覺.
你越來越多事情可以感恩.
如果你不去感恩.
你就會覺得一切都是理所當然.
於是我的心就越來越開心.
我想問弟兄姊妹多一次.
你想不想你有個開心的人生.
Amen.
所以頂智慧.
我鼓勵大家.
這個不只是我們的想.

$^{1081}$我們的想是神的心意.
所以這個很重要.
大家要明白這件事.
好 最後一個比喻.
大家都喝過牛奶吧.
如果你不喝牛奶.
放在家裡的桌子裡.
放四五天會怎樣.
變壞了.
臭了 是不是.
就好像圖一樣.
但是牛奶可以變為乳酪.
你知不知道我有沒有吃過乳酪.
很好吃 是不是.
很好吃的.
但是關鍵在哪裡.
關鍵是牛奶發酵的時候.
究竟有什麼菌在那裡.
如果是一些雜菌.
或者幹菌.
或者其他的奇怪的菌.
牛奶就會變壞 發臭.
但是如果有乳酸菌.
加進去的一個程序裡面.
就可以變為大家都很喜歡吃的乳酪.
大一點的頂子妹.
葡萄也一樣.
葡萄可以變壞 發臭.
那你說怎樣走呢.
關鍵也是.
它可以變為好喝的紅酒.
但是關鍵在發酵的過程.
究竟它是雜菌.
還是酵母菌.
如果是雜菌.
就變成發臭.
但是如果是酵母菌.
經過這個程序.
就會變成大家都很喜歡喝的紅酒.
頂子妹就是我們的人生.

$^{1121}$究竟我們人生的酵母菌.
和乳酸菌是什麼呢.
這個是重要的.
是決定性的.
我很相信就是大家願意.
跟聖靈配合去操練這樣的心.
以及聖靈的梟貫.
不斷的充滿我們.
但是最重要是什麼.
這一切的來源.
這個心意.
這個能力.
都是我們要信的神加給我們的.
所以頂子妹.
你想過一個什麼樣的人生呢.
剛才大家都回答我了.
大家都很願意過一個喜樂的人生.
頂子妹這個也是神的心意.
你選什麼呢.
我們一起祈禱.
主啊讓你的靈充滿我們的時候.
把你的真理放在我們的心裡面.
讓我們很明白你的心意.
今天我們心裡面的盼望.
我們會感恩.
因為這個也是.
你對我們最終極的心意.
所以帶著感恩的心去讚美你.
頂子妹當我們一起禱告的時候.
聖靈對大家有兩個邀請.
第一個邀請是.
已經信於耶穌的頂子妹.
可能你新信.
或者已經信了一段長時間.
但是今天聖靈要再一次去印證.
祂的心在你的靈裡面.
所以神要呼召你.
無論你信的日子如何.
今天你願意在神的面前.
去選擇那個喜樂的人生.

$^{1161}$願意去操練這個感恩.
以致聖靈一直這樣充滿你.
讓你心靈的重擔去離開你.
並且那個喜樂.
要在你生命裡面完全的永留.
如果願意這樣做的頂子妹.
我邀請你把你的手放在你的心上.
七個主.
當我們的頂子妹.
很願意過一個喜樂的人生.
我們很感恩.
因為是神你的心意.
你很想我們帶著喜樂.
帶著感恩.
並且在我們生命裡面去榮耀你.
所以我奉耶穌基督名.
讓你的復活的大能.
充滿我們頂子妹.
復甦我們的生命.
復甦我們的讚美.
復甦我們的敬拜.
讓我們操練這個感恩的心的時候.
主你繼續加能施力給我們.
聖靈在我們生命裡面工作.
讓我們看見比以往更多的東西.
更豐富的東西.
更喜樂的東西.
更有盼望的東西.
讓我們生命有一個突破.
讓我們生命裡面有那種喜樂.
喜樂的心擁有出來.
如串聯一樣.
主祂讚美你.
我們多謝你.
當我們繼續閉上眼睛禱告的時候.
我不知道在座當中.
有沒有一些朋友.
今天我們很清楚.
主祂的心意是要讓你有喜樂的人生.
所以祂很想你跟祂建立關係.

$^{1201}$我們有喜樂的人生.
我們可以操練感恩.
但是我們最重要的決定.
就是聖靈在我們當中的.
僥倦能力和工作.
我們怎樣才有聖靈的工作呢.
就是我們跟耶穌要建立這個關係.
我們怎樣跟耶穌建立關係呢.
很簡單.
只是打開你的心去接受耶穌.
你要馬上跟耶穌建立關係.
聖靈要在你的生命裡面工作.
所以在座當中.
有些朋友.
今天你願意跟耶穌建立關係.
你很想打開你的心.
接受耶穌的.
我邀請你.
如果你願意接受耶穌的.
我邀請你在你的座位舉手.
有沒有你的朋友在我們當中呢.
如果有的.
我邀請你在你的座位舉手.
有沒有呢.
有沒有你的朋友呢.
還有沒有呢.
還有沒有你的朋友呢.
今天你願意打開你的心去接受耶穌.
讓喜樂的靈去充滿你.
將你的人生改變.
有沒有呢.
有沒有你的朋友在我們當中.
如果你相信耶穌的.
我邀請你在你的座位舉手.
有沒有你的朋友呢.
好.
看到了.
好 我們一起祈禱.
是女人.
主 我們讚美你.

$^{1241}$因為今天是你的日子.
是你復甦的日子.
是恢復一切你自己的祝福的日子.
所以今天是一個感恩.
是一個讚美你.
是榮耀你的日子.
所以當我們繼續用這首回應詩去敬拜你的時候.
就讓你的榮耀繼續充滿我們每一位.
聽我們禱告.
奉耶穌基督得逞人民.
阿咩.
\newpage



\section{約翰福音 20:19-23}
\label{sec:031939}
\textbf{門關了,主在中間}~~~~葉慧敏傳道
\newline
\newline
連結: \href{http://www.church.com.hk/acms/content.asp?site=cdc\&op=show\&type=product\&code=031939\&layout=sermon}{\texttt{031939}} ~~~~ 語音日期: 2026-04-05
\newline
\newline
\hyperref[sec:031938]{\small{< < < PREV SERMON < < <}}
~
\hyperref[sec:index_chronic]{\small{[返順時目]}}
~
\hyperref[sec:031940]{\small{> > > NEXT SERMON > > >}}
\newline
\newline
約翰福音 20:19-23
\newline
\begin{longtable}{cl}
\hline
\hline
章節 & 經文 (和合本修訂版)\\
\hline
20:19 & \begin{tabularx}{0.7\textwidth}{X} 那日（就是七日的第一日）晚上，門徒因怕猶太人，所在的地方門都關了。耶穌來，站在當中，對他們說：「願你們平安！」 \end{tabularx} \\ \\ \relax
20:20 & \begin{tabularx}{0.7\textwidth}{X} 說了這話，他把手和肋旁給他們看。門徒一看見主就喜樂了。 \end{tabularx} \\ \\ \relax
20:21 & \begin{tabularx}{0.7\textwidth}{X} 於是耶穌又對他們說：「願你們平安！父怎樣差遣了我，我也照樣差遣你們。」 \end{tabularx} \\ \\ \relax
20:22 & \begin{tabularx}{0.7\textwidth}{X} 說了這話，他向他們吹一口氣，說：「領受聖靈吧！ \end{tabularx} \\ \\ \relax
20:23 & \begin{tabularx}{0.7\textwidth}{X} 你們赦免誰的罪，誰的罪就得赦免；你們不赦免誰的罪，誰的罪就不得赦免。」 \end{tabularx} \\ \\
[1ex]
\hline
\hline
\end{longtable}
$^{1}$今天我們講到的主題叫「門關了,主在中間」.
剛才Stanley大兄領的詩歌.
其實我一邊唱的時候就想.
其實我不用說了.
他已經透過詩歌幫我表達了我整個信息.
我們沒有溝通過.
純粹是他想他的詩歌.
我想我的講道.
但是很配合我想說的.
我們很感恩.
丁姐妹,當我們講到復活的時候.
你第一時間想起什麼?.
想起有幾天假期.
想起超市的復活蛋.
想起一段很熟悉的經文.
還是想起我們每年都紀念的一個信仰事件呢?.
但是如果我們問一個更加誠實的問題就是.
復活是真的.
但是對你來說.
是歷史還是這一刻呢?.
因為我們生活裡面仍然有壓力.
仍然有焦慮.
仍然有不安.
今天很想從約翰福音二十章十九至二十三節裡面.
我們看到一群在恐懼裡面的門徒.
如何遇見這位復活主.
也讓我們問自己.
今天我們願不願意打開我們的生命.
讓復活主再次站在我們中間.
讓我們一起祈禱.
耶穌你是那位在死裡面復活的主.
主我求你幫助今天是我們每一位.
無論是.
站立在你跟前的講員.
還是一群坐在這裡的會眾.
神大求你在裡面與我們每一個人說話.
讓我們更深的去看見.
是啊,你的復活對我們生命的意義.
主求你幫助我們,帶領我們.
我們將這套禱告交託.

$^{41}$奉主名求,阿們.
我們不如一起讀一下這幾節經文.
約翰福音二十章十九至二十三節.
我們一起讀.
萬上門徒所在的地方.
因怕猶太人門都關了.
耶穌來站在當中對他們說.
願你們平安.
說了這話.
就把手和肋膀指給他們看.
門徒看見主就喜樂了.
耶穌又對他們說.
願你們平安.
父怎樣差遣了我.
我也照樣差遣你們.
說了這話.
就向他們吹一口氣說.
你們受聖靈.
你們赦免誰的罪.
誰的罪就赦免了.
你們留下誰的罪.
誰的罪就留下了.
好,謝謝.
在第十九節他提到.
二十章.
其實在第十九節他提到.
那日晚上.
這個「那日」是指什麼時候.
其實我們可以回到.
二十章第一節.
就是.
門徒已經聽過耶穌復活的消息.
那個「那日」晚上.
其實是指耶穌死後三天.
即是星期日.
即是祂復活那一天.
到了十九節.
約翰特別說「那日」晚上.
其實是指同一天.
復活已經發生了.

$^{81}$而來到晚上.
門徒仍然是關著門.
其實在一至十八節裡面.
有門徒知道耶穌復活.
在墳墓裡面不見了.
也有婦女知道耶穌復活了.
但是來到第十九節.
門徒仍然是把門關上.
經文說.
因為怕猶太人門都關了.
這個說話很短.
但是充滿著張力.
我們可以看看門徒的生命狀態.
門徒當時不是一個普通的休息.
他們內心是一個封閉的狀態.
外在的環境有真實的威脅.
耶穌剛剛被釘完十字架.
三天前才被釘完十字架.
宗教領袖敵視耶穌.
門徒自然會想.
下一個會不會就是我呢.
這種恐懼不是幻想.
而是很現實.
裡面他們有深層的恐懼.
他們不單止是怕人.
更加是怕未來.
三年來跟隨的主死了.
方向沒有了.
盼望好像斷線一樣.
他們不知道下一步怎樣去行.
他們心裡面是混亂.
他們聽到婦女說耶穌復活.
但是又未親眼見到.
信又不是完全不信.
不信.
信又不是完全信.
不信又不是完全不信.
信仰變成一個拉扯更加令人不安.
所以.
他們關門.

$^{121}$不單止是關實了一道門.
而是關注了他們自己的心.
而其實今天.
我們都可能是活在這種.
關門的生命裡面.
這幾年.
我們大部分.
我想百分之九十九個百分比的會眾.
都是由香港來到英國裡面.
因為BNO Visa我們來到.
表面是一個新開始.
但是內心其實一點都不容易.
可能你正在做一份不太熟悉的工作.
開始懷疑自己的價值.
有可能你的家人仍然在香港.
你心裡面很多的牽掛.
很孤單.
很多事情要重新開始.
很不習慣.
表面上好像很OK.
其實裡面不是很OK.
慢慢我們就.
因為跟身邊的人又不是很熟.
慢慢我們未必會說很多.
把自己的感受收起來.
連自己都不再觸碰.
門就是這樣慢慢的關上.
這種狀態其實跟門徒當日一樣.
生命被困住.
前面一片的模糊.
但正正就在這種門關了的狀態裡面.
聖經就說.
它說:耶穌來,站在其中.
這句經文是一個很震撼的地方.
就是不是門徒的恐懼.
而是耶穌的行動.
耶穌來,站在當中.
圓聞其實是有一種突然出現的意思.
不是慢慢走進來.
而是門仍然關著.

$^{161}$但耶穌已經在中間.
是,門是關著了.
但是耶穌仍然來到.
門徒沒有開門.
但是主主動的進入.
一點很重要.
耶穌不是等門徒準備好才出現.
祂沒有等門徒信心恢復.
祂沒有等門徒冷靜下來.
祂也沒有等門徒.
祈禱,祈禱到一個救贖靈的狀態.
耶穌就這樣來了.
而且耶穌來到之後.
祂沒有擇地門徒.
祂沒有說.
為什麼你們這麼沒有信心.
祂也沒有說.
我之前不是跟你們說我會復活的嗎.
祂沒有質問,沒有壓力,沒有要求.
耶穌只是站在當中.
這個站在當中很有意思.
祂不是在門口觀察著你們這些人後做什麼.
祂不是在遠處等候.
而是進入他們的空間.
進入他們的處境.
成為他們中間的一份子.
弟兄姊妹.
這就是復活主的心.
祂不怕我們恐懼.
不怕我們混亂.
甚至不怕我們封閉.
我們以為自己的狀態會阻礙神.
其實不會.
因為復活主第一個動作.
不是去改變環境.
而是臨在.
今天也一樣.
我們可能還沒準備好.
我們可能也有很多事情.
還是很混亂.

$^{201}$甚至我們不知道怎麼祈禱.
但是主已經在我們中間.
經文接著說.
耶穌來到他們中間.
第一句說話就是.
願你們平安.
這句說話我們可能聽過很多次.
但是在當時的處境裡面.
是一個極其有力量的.
門徒的環境沒有變動.
外面仍然是危險.
局勢仍然是緊張.
未來仍然是不確定.
但是耶穌沒有先處理環境.
祂先將平安帶到這班門徒中間.
弟兄姊妹聖經說的平安.
不是問題解決了的平安.
而是主在的平安.
不是因為外面風浪停了.
而是因為主站在船上.
門徒以為開門才有平安.
但是耶穌說.
其實我在這裡.
你們就有平安.
所以平安不是一種環境.
而是一種的同在.
當耶穌在門徒中間的時候.
恐懼就開始失去了力量.
當主的聲音進入了他們的耳裡面.
他們的混亂就開始被整理.
復活帶來第一份的禮物.
不是復活彈.
不是使命.
不是能力.
而是平安.
一種因為主而來的平安.
今天主同樣跟我們說.
願你平安.
不是因為我們已經沒有問題.
而是因為祂在我們的生命裡面.

$^{241}$所以真正的問題.
不是我們還有沒有恐懼.
而是我們有沒有留意到.
主已經在我們中間.
當我們在恐懼裡面.
我們記住我們不是一個人.
是我們有耶穌跟我們一起.
或者今天我們要做的.
不是努力地撐下去.
不是找很多方法去解決我們的問題.
而是嘗試打開我們的心門給主耶穌.
向主誠實地講出我們真正的恐懼.
不需要很完美.
不需要整理得很好.
只需要承認.
耶穌我們真的有點害怕.
有點擔心.
門可以關.
但是主仍然會來.
問題不是門開了沒有.
而是我們有沒有留意到.
耶穌就在我們當中.
十九節講到主怎樣進入門徒的恐懼.
敘事者在二十節帶我們看見.
當人真正遇見復活主的時候.
關係是會被更新.
生命也開始有轉變.
這種因為復活而有的新關係.
經文在二十節說.
說了這話.
就是講完願你平安之後.
耶穌做了一個行動.
祂就說.
就把手和肋膀指給他們看.
耶穌給他們看手.
給他們看肋膀.
這個動作很深.
耶穌沒有只是說一句.
是我啊.
你認不認得我啊.

$^{281}$祂是主動展示他的傷痕.
手上的釘痕.
肋膀被刺穿的傷口.
為什麼復活之後還要保留這些傷痕.
因為耶穌要門徒知道.
站在你們面前那個.
仍然是那位為你受苦.
為你捨命的主.
耶穌的復活不是要抹走過去.
而是將傷痕轉化為恩典的記號.
復活不是否定十字架.
而是宣告十字架並無白費.
今天我們的生命都一樣.
神不一定會抹走我們經歷過的傷痛.
但是祂會重新定義我們的傷痛.
可能在你心裡面.
都有一個位置還未康復.
但是神可以將這些傷口.
由痛苦變成見證.
由頗衰變為恩典.
門徒看見主之後.
經文說祂就起落.
留意轉變發生得很快.
上一刻他們仍然在恐懼裡面.
而這一刻他們已經充滿著起落.
但是環境沒有改變.
門仍然是關著的.
外面仍然是危險的.
改變的不是外在.
而是因為他們看見主.
起落不是環境變好.
而是我們重新的看見主.
這班門徒重新的看見那位.
愛他們愛到底愛他們.
愛到為他們死在十字架的主.
重新的相信耶穌沒有離開他們.
重新的知道耶穌仍然在他們生命中間.
所以起落不是開心了那麼簡單.
而是他們和耶穌的關係重新的連結起來.
門徒的轉變不單單是情緒的變化.

$^{321}$而是和主的關係重新的連結.
之前十字架.
他們看見耶穌在十字架上的離開.
他們可能感覺到很失望.
死亡令他們困惑.
和主之間好像斷了一條線.
但是當再次看見這位復活主的時候.
一切就重新的連結起來.
他們知道耶穌沒有離開過.
耶穌仍然是那位愛他們帶領他們的主.
弟兄姊妹今天我們都一樣.
有時候我們的信仰不是消失了.
而是我們和耶穌的關係好像變淡了.
我們仍然回教會仍然回來崇拜聽道.
但是我們內心好像少了一份真實的喜樂.
我們回來未必聽到什麼.
我們查經都未必有經文進入我們裡面.
我們和主好像的關係冷淡了.
但是復活主今天仍然向我們顯明他自己.
當我們再次看見耶穌的時候.
我們的心就會再次的活過來.
如果說剛才說的第一點是主進入恐懼.
第二點就是講到關係被恢復.
而第三點我給他的主題就是.
復活不單止是去安慰人.
而是差遣人.
復活主的臨在從來不叫人停留在一個安全感裡面.
而是帶人走向使命.
二十一節耶穌這樣說.
祂說:耶穌有對他們說:願你們平安.
父怎樣差遣我,我也照樣差遣你們.
耶穌第二次說:願你們平安.
不是重複而是一個推進.
從經文觀察出來.
願你們平安在十九節和二十一節那個屬靈的意義.
我們可以看為.
第一句,十九節的平安是對門徒內心那種的恐懼去醫治他們.
而二十一節這裡說的平安是預備差遣.
耶穌說了一句說話.
願你們平安之後,祂說了一句震撼的說話就是.

$^{361}$父怎樣差遣我,我也照樣差遣你們.
父差遣耶穌是進入世界,接觸破碎,與人同行,活出愛.
現在耶穌說:你們都要是這樣.
門徒再是那一群關起門的人.
他們是一群要被差出去的人.
剛才說,這幾年來英國的朋友,可能都是內心有不同的挑戰,不同的掙扎.
對我們一群信徒來說,其實都可能不多不少,都是這樣.
我們來到英國,如果有我們原本認識的朋友家人,都更好一點.
因為起碼有人會明白你.
但是其實很多人來到之後,其實是由零開始.
工作由零開始,教會由零開始,住的地方由零開始.
在教會裡面,大家不同背景,不同性格,不同經歷.
但我們又要很迅速地一起去相處,甚至一起事奉.
裡面有很多的摩擦,有很多的誤會,甚至有很多的不舒服,令到心裡不開心.
可能只是一次的會議,我們覺得別人很難合作.
又可能是一個小組的分享,有人說了一些話,你覺得令你很不中聽.
以前我們可能會選擇退後,或者心裡就是落閘.
我們少一點接觸.
但是耶穌在這次經文提醒我們,祂差遣我們.
不是叫我們用回自己習慣的那一套去面對我們身邊的人.
面對我們的社會,面對我們的世界.
而是學習好像主耶穌這樣進入人的生命,與人同在.
去學習彼此的相顧,去學習凡事包容,凡事相信,凡事盼望.
去體會愛是永不止息這個功課.
不是因為對方容易阻礙相處,而是因為主先是這樣對我們.
以致我們有那份力量去愛我們身邊的人.
二十二節祂這樣,耶穌給了這班門徒一個新的能力.
祂說:「說了這話,就向他們吹一口氣,說你們受聖靈」.
因為畫面令人想起,創世記神向人吹氣,人就成為有生命.
就是說復活的主不單止是差遣我們.
祂是將新生命放在我們的當中.
門徒之前很害怕,很膽怯,很退縮.
但現在主將祂的氣息吹入他們的生命裡面.
將他們原本有很恐懼的人變為有生命,有能力,有使命的人.
聖靈不是抽象的力量,而是神親自住在我們裡面的一份支持.
當我們不知道怎樣去愛的時候,我們就求耶穌,求聖靈去教我們.
當我們沒有能力再堅持的時候,我們就求主去托著我們.
當我們不知道前面怎樣去行的時候,我們就求主去引導我們.
所以不是靠我們自己做好一點,而是靠聖靈會出一個被差遣的生命.
丁姐妹,如果我們將約翰福音十九至二二節連在一起看的時候.

$^{401}$我們可以看到一條屬靈的盡情.
復活主進入門徒的恐懼,帶來平安,恢復跟耶穌的關係.
亦都賦予使命,並且將聖靈的能力賜下.
這個次序很重要,神單單只是在安慰這班門徒.
祂是一步一步帶這班門徒走進一個新生命的方向.
耶穌不忽略我們的恐懼,不忽略我們的擔心.
但祂不讓我們停留在恐懼裡面,祂讓我們有平安.
祂不單止要我們經歷祂的同在,祂更呼召我們.
要帶著這份同在進入這個世界裡面.
所以基督徒的生命不單止是與神親近,而是被神差遣.
不單單是領受主的恩典,而是要成為恩典流通的器命.
被差遣不是一定要做大事,而是在你現在的位置裡面.
忠心地活出主的同在.
當我們回望這段經文的時候,我們看見上主那個奇妙的計劃.
一班原本被恐懼困住,關上門的人遇見主動臨在的主.
他們由驚慌變成喜樂,由封閉成為被差遣.
由無力變成領受聖靈的氣息.
這段經文不單止是門徒的經文,也是我們每一個人的寫照.
今日有些人,可能有些人仍然是在一個關門裡面.
我們心裡面有壓力,有恐懼,有迷失.
不是不相信,但是覺得再走下去很沒力.
但是復活主今日仍然是進入我們的處境裡面.
他不需要等我們先調整好自己.
他就站在這裡說願你平安.
亦有些人,你和主的關係未必是斷開,你未必是不相信.
但是你覺得總是有點疏離.
你仍然參加聚會,仍然來敬拜.
但是裡面少了一份親近.
今日主再次向你譴明,他自己不是一個抽象的神.
而是一位帶著釘痕,仍然深愛著你,愛到底的主.
當我們再次轉向主的時候,我們的心被更新.
我們的生命重新被延讓.
亦有些人,我們已經經歷主的同在.
但是我們停留在一個位置,習慣了被安慰.
未有力量去踏出回應和差遣那一步.
今日復活主就對我們說.
父怎樣差遣我,我也照樣差遣你們.
他不單止是差遣你,他也將聖靈放在我們每一個人的生命裡面.
以致我們可以在我們的家庭,在我們的工作,在我們的社區.
在我們的教會,活出被差遣的生命.

$^{441}$所以弟兄姊妹,我也很想藉著復活節這個主日.
邀請大家有一個簡單而真誠的回應.
如果你仍然活在關閉和恐懼之中.
你可以心裡面去祈禱.
主,我求你進入我的生命裡面.
如果你覺得和主有距離.
你可以去禱告,主,讓我再次見到你.
如果你願意回應差遣.
你可以去禱告,主,我願意被你差遣.
我們現在也想邀請大家安靜.
在你的座位裡面安靜一會兒.
問自己,你現在是關著門,還是願意打開呢.
如果你心裡面真的有我這個表裡面所說的狀態.
你有沒有勇氣去禱告主,去求主在裡面幫助你呢.
\newpage



\section{腓立比書 1:1-4:23}
\label{sec:031940}
\textbf{腓立比書:喜樂的秘訣 (1) 全書概論}~~~~劉少康牧師
\newline
\newline
連結: \href{http://www.church.com.hk/acms/content.asp?site=cdc\&op=show\&type=product\&code=031940\&layout=sermon}{\texttt{031940}} ~~~~ 語音日期: 2026-04-08
\newline
\newline
\hyperref[sec:031939]{\small{< < < PREV SERMON < < <}}
~
\hyperref[sec:index_chronic]{\small{[返順時目]}}
~
\hyperref[sec:031941]{\small{> > > NEXT SERMON > > >}}
\newline
\newline
腓立比書 1:1-4:23
\newline
\begin{longtable}{cl}
\hline
\hline
章節 & 經文 (和合本修訂版)\\
\hline
1:1 & \begin{tabularx}{0.7\textwidth}{X} 基督耶穌的僕人保羅和提摩太寫信給住腓立比、在基督耶穌裡的眾聖徒，以及諸位監督和執事。 \end{tabularx} \\ \\ \relax
1:2 & \begin{tabularx}{0.7\textwidth}{X} 願恩惠、平安從我們的父神和主耶穌基督歸給你們！ \end{tabularx} \\ \\ \relax
1:3 & \begin{tabularx}{0.7\textwidth}{X} 我每逢想念你們，就感謝我的神， \end{tabularx} \\ \\ \relax
1:4 & \begin{tabularx}{0.7\textwidth}{X} 每逢為你們眾人祈求的時候，總是歡歡喜喜地祈求， \end{tabularx} \\ \\ \relax
1:5 & \begin{tabularx}{0.7\textwidth}{X} 因為從第一天直到如今，你們都同心合意興旺福音。 \end{tabularx} \\ \\ \relax
1:6 & \begin{tabularx}{0.7\textwidth}{X} 我深信，那在你們心裡動了美好工作的，到了耶穌基督的日子必完成這工作。 \end{tabularx} \\ \\ \relax
1:7 & \begin{tabularx}{0.7\textwidth}{X} 我為你們眾人有這樣的想法原是應當的，因為你們常在我心裡；無論我是在捆鎖中，在辯明並證實福音的時候，你們都與我一同蒙恩。 \end{tabularx} \\ \\ \relax
1:8 & \begin{tabularx}{0.7\textwidth}{X} 我以基督耶穌的心腸切切想念你們眾人，這是神可以為我作證的。 \end{tabularx} \\ \\ \relax
1:9 & \begin{tabularx}{0.7\textwidth}{X} 我所禱告的就是：要你們的愛心，在知識和各樣見識上，不斷增長， \end{tabularx} \\ \\ \relax
1:10 & \begin{tabularx}{0.7\textwidth}{X} 使你們能分辨是非，在基督的日子作真誠無可指責的人， \end{tabularx} \\ \\ \relax
1:11 & \begin{tabularx}{0.7\textwidth}{X} 更靠著耶穌基督結滿仁義的果子，歸榮耀稱讚給神。 \end{tabularx} \\ \\ \relax
1:12 & \begin{tabularx}{0.7\textwidth}{X} 弟兄們，我要你們知道，我所遭遇的事反而使福音更興旺， \end{tabularx} \\ \\ \relax
1:13 & \begin{tabularx}{0.7\textwidth}{X} 以致御營全軍和其餘的人都知道我是為基督的緣故受捆鎖的； \end{tabularx} \\ \\ \relax
1:14 & \begin{tabularx}{0.7\textwidth}{X} 而且那在主裡的弟兄，多半都因我受的捆鎖而篤信不疑，越發放膽無所懼怕地傳道。 \end{tabularx} \\ \\ \relax
1:15 & \begin{tabularx}{0.7\textwidth}{X} 有些人傳基督是出於嫉妒紛爭；有些人是出於好意。 \end{tabularx} \\ \\ \relax
1:16 & \begin{tabularx}{0.7\textwidth}{X} 後者是出於愛心，知道我奉差遣是為福音辯護的。 \end{tabularx} \\ \\ \relax
1:17 & \begin{tabularx}{0.7\textwidth}{X} 前者傳基督是出於自私，動機不純，企圖要加增我捆鎖的苦楚。 \end{tabularx} \\ \\ \relax
1:18 & \begin{tabularx}{0.7\textwidth}{X} 這又何妨呢？或是假意或是真心，無論如何，只要基督被傳開了，為此我就歡喜。我還要歡喜， \end{tabularx} \\ \\ \relax
1:19 & \begin{tabularx}{0.7\textwidth}{X} 因為我知道，這事藉著你們的祈禱和耶穌基督的靈的幫助，終必使我得到釋放。 \end{tabularx} \\ \\ \relax
1:20 & \begin{tabularx}{0.7\textwidth}{X} 這就是我所切慕、所盼望的：沒有一事能使我羞愧；反倒凡事坦然無懼，無論是生是死，總要讓基督在我身上照常顯大。 \end{tabularx} \\ \\ \relax
1:21 & \begin{tabularx}{0.7\textwidth}{X} 因為我活著就是基督，死了就有益處。 \end{tabularx} \\ \\ \relax
1:22 & \begin{tabularx}{0.7\textwidth}{X} 但是，我在肉身活著，若能有工作的成果，我就不知道該挑選甚麼。 \end{tabularx} \\ \\ \relax
1:23 & \begin{tabularx}{0.7\textwidth}{X} 我處在兩難之間：我情願離世與基督同在，因為這是好得無比的； \end{tabularx} \\ \\ \relax
1:24 & \begin{tabularx}{0.7\textwidth}{X} 然而，我為你們肉身活著更加要緊。 \end{tabularx} \\ \\ \relax
1:25 & \begin{tabularx}{0.7\textwidth}{X} 既然我這樣深信，就知道仍要留在世間，且與你們眾人一起存留，使你們在所信的道上又長進又喜樂， \end{tabularx} \\ \\ \relax
1:26 & \begin{tabularx}{0.7\textwidth}{X} 為了我再到你們那裡時，你們在基督耶穌裡的誇耀越發加增。 \end{tabularx} \\ \\ \relax
1:27 & \begin{tabularx}{0.7\textwidth}{X} 最重要的是：你們行事為人要與基督的福音相稱，這樣，無論我來見你們，或不在你們那裡，都可以聽到你們的景況，知道你們同有一個心志，站立得穩，為福音的信仰齊心努力， \end{tabularx} \\ \\ \relax
1:28 & \begin{tabularx}{0.7\textwidth}{X} 絲毫不怕敵人的威脅；以此證明他們會沉淪，你們會得救，這是出於神。 \end{tabularx} \\ \\ \relax
1:29 & \begin{tabularx}{0.7\textwidth}{X} 因為你們蒙恩，不但得以信服基督，而且要為他受苦。 \end{tabularx} \\ \\ \relax
1:30 & \begin{tabularx}{0.7\textwidth}{X} 你們的爭戰，就與你們曾在我身上見過、現在所聽到的是一樣的。 \end{tabularx} \\ \\ \relax
2:1 & \begin{tabularx}{0.7\textwidth}{X} 所以，在基督裡若有任何勸勉，若有任何愛心的安慰，若有任何聖靈的團契，若有任何慈悲憐憫， \end{tabularx} \\ \\ \relax
2:2 & \begin{tabularx}{0.7\textwidth}{X} 你們就要意志相同，愛心相同，有一致的心思，一致的想法，使我的喜樂得以滿足。 \end{tabularx} \\ \\ \relax
2:3 & \begin{tabularx}{0.7\textwidth}{X} 凡事不可自私自利，不可貪圖虛榮；只要心存謙卑，各人看別人比自己強。 \end{tabularx} \\ \\ \relax
2:4 & \begin{tabularx}{0.7\textwidth}{X} 各人不要單顧自己的事，也要顧別人的事。 \end{tabularx} \\ \\ \relax
2:5 & \begin{tabularx}{0.7\textwidth}{X} 你們當以基督耶穌的心為心： \end{tabularx} \\ \\ \relax
2:6 & \begin{tabularx}{0.7\textwidth}{X} 他本有神的形像，卻不堅持自己與神同等； \end{tabularx} \\ \\ \relax
2:7 & \begin{tabularx}{0.7\textwidth}{X} 反倒虛己，取了奴僕的形像，成為人的樣式；既有人的樣子， \end{tabularx} \\ \\ \relax
2:8 & \begin{tabularx}{0.7\textwidth}{X} 就謙卑自己，存心順服，以至於死，且死在十字架上。 \end{tabularx} \\ \\ \relax
2:9 & \begin{tabularx}{0.7\textwidth}{X} 所以神把他升為至高，又賜給他超乎萬名之上的名， \end{tabularx} \\ \\ \relax
2:10 & \begin{tabularx}{0.7\textwidth}{X} 使一切在天上的、地上的和地底下的，因耶穌的名，眾膝都要跪下， \end{tabularx} \\ \\ \relax
2:11 & \begin{tabularx}{0.7\textwidth}{X} 眾口都要宣認：耶穌基督是主，歸榮耀給父神。 \end{tabularx} \\ \\ \relax
2:12 & \begin{tabularx}{0.7\textwidth}{X} 我親愛的，這樣看來，你們向來是順服的，不但我在你們那裡，就是我現在不在你們那裡的時候更是順服的，就當恐懼戰兢完成你們自己得救的事； \end{tabularx} \\ \\ \relax
2:13 & \begin{tabularx}{0.7\textwidth}{X} 因為是神在你們心裡運行，使你們又立志又實行，為要成就他的美意。 \end{tabularx} \\ \\ \relax
2:14 & \begin{tabularx}{0.7\textwidth}{X} 你們無論做甚麼事，都不要發怨言起爭論， \end{tabularx} \\ \\ \relax
2:15 & \begin{tabularx}{0.7\textwidth}{X} 好使你們無可指責，誠實無偽，在這彎曲悖謬的世代作神無瑕疵的兒女。你們在這世代中要像明光照耀， \end{tabularx} \\ \\ \relax
2:16 & \begin{tabularx}{0.7\textwidth}{X} 將生命的道顯明出來，使我在基督的日子得以誇耀我沒有白跑，也沒有徒勞。 \end{tabularx} \\ \\ \relax
2:17 & \begin{tabularx}{0.7\textwidth}{X} 我以你們的信心為供獻的祭物，我若被澆獻在其上也是喜樂，並且與你們眾人一同喜樂。 \end{tabularx} \\ \\ \relax
2:18 & \begin{tabularx}{0.7\textwidth}{X} 你們也要照樣喜樂，並且與我一同喜樂。 \end{tabularx} \\ \\ \relax
2:19 & \begin{tabularx}{0.7\textwidth}{X} 我靠主耶穌希望很快能差提摩太去見你們，好讓我知道你們的事而心裡得著安慰。 \end{tabularx} \\ \\ \relax
2:20 & \begin{tabularx}{0.7\textwidth}{X} 因為我沒有別人與我同心，真正關懷你們的事。 \end{tabularx} \\ \\ \relax
2:21 & \begin{tabularx}{0.7\textwidth}{X} 其他的人都求自己的事，並不求耶穌基督的事。 \end{tabularx} \\ \\ \relax
2:22 & \begin{tabularx}{0.7\textwidth}{X} 但你們知道提摩太是經得起考驗的，他與我為了福音一同服侍，待我像兒子待父親一樣。 \end{tabularx} \\ \\ \relax
2:23 & \begin{tabularx}{0.7\textwidth}{X} 所以，我一看出我的事怎樣了結，我希望立刻差他去， \end{tabularx} \\ \\ \relax
2:24 & \begin{tabularx}{0.7\textwidth}{X} 但我靠著主自信我不久也會去。 \end{tabularx} \\ \\ \relax
2:25 & \begin{tabularx}{0.7\textwidth}{X} 然而，我想必須差以巴弗提到你們那裡去。他是我的弟兄、同工和戰友，是你們差遣來供應我需要的。 \end{tabularx} \\ \\ \relax
2:26 & \begin{tabularx}{0.7\textwidth}{X} 他很想念你們眾人，並且極其難過，因為你們聽見他病了。 \end{tabularx} \\ \\ \relax
2:27 & \begin{tabularx}{0.7\textwidth}{X} 他真的生病了，幾乎要死。然而神憐憫他，不但憐憫他，也憐憫我，免得我憂上加憂。 \end{tabularx} \\ \\ \relax
2:28 & \begin{tabularx}{0.7\textwidth}{X} 所以，我更要盡快送他回去，好讓你們再見到他而喜樂，我也可以減少憂愁。 \end{tabularx} \\ \\ \relax
2:29 & \begin{tabularx}{0.7\textwidth}{X} 故此，你們要在主裡歡歡喜喜地接待他，而且要尊重這樣的人， \end{tabularx} \\ \\ \relax
2:30 & \begin{tabularx}{0.7\textwidth}{X} 因他為做基督的工作不顧性命，幾乎至死，為要補足你們供應我不夠的地方。 \end{tabularx} \\ \\ \relax
3:1 & \begin{tabularx}{0.7\textwidth}{X} 末了，我的弟兄們，你們要靠主喜樂。我把這些話再寫給你們，對我並不困難，對你們卻是妥當的。 \end{tabularx} \\ \\ \relax
3:2 & \begin{tabularx}{0.7\textwidth}{X} 應當防備犬類，防備作惡的，防備妄自行割的。 \end{tabularx} \\ \\ \relax
3:3 & \begin{tabularx}{0.7\textwidth}{X} 因為真受割禮的，就是我們這藉著神的靈敬拜、以基督耶穌為誇耀、不依靠肉體的。 \end{tabularx} \\ \\ \relax
3:4 & \begin{tabularx}{0.7\textwidth}{X} 其實，我也可以靠肉體；若是別人以為他可以依靠肉體，我更可以。 \end{tabularx} \\ \\ \relax
3:5 & \begin{tabularx}{0.7\textwidth}{X} 我出生後第八天受割禮；我是以色列族、便雅憫支派的人，是希伯來人所生的希伯來人。就律法說，我是法利賽人； \end{tabularx} \\ \\ \relax
3:6 & \begin{tabularx}{0.7\textwidth}{X} 就熱心說，我是迫害教會的；就律法上的義說，我是無可指責的。 \end{tabularx} \\ \\ \relax
3:7 & \begin{tabularx}{0.7\textwidth}{X} 只是我先前以為對我是有益的，我現在因基督的緣故而當作是有損的。 \end{tabularx} \\ \\ \relax
3:8 & \begin{tabularx}{0.7\textwidth}{X} 不但如此，我已把萬事當作是有損的，因我以認識我主基督耶穌為至寶。我為他已經丟棄萬事，看作糞土，為要贏得基督， \end{tabularx} \\ \\ \relax
3:9 & \begin{tabularx}{0.7\textwidth}{X} 並且得以在他裡面，不是有自己因律法而得的義，而是有信基督的義，就是基於信，從神而來的義， \end{tabularx} \\ \\ \relax
3:10 & \begin{tabularx}{0.7\textwidth}{X} 使我認識基督，知道他復活的大能，並且知道和他一同受苦，效法他的死， \end{tabularx} \\ \\ \relax
3:11 & \begin{tabularx}{0.7\textwidth}{X} 或許我也得以從死人中復活。 \end{tabularx} \\ \\ \relax
3:12 & \begin{tabularx}{0.7\textwidth}{X} 這不是說我已經得著了，已經完全了；而是竭力追求，或許可以得著基督耶穌所要我得著的。 \end{tabularx} \\ \\ \relax
3:13 & \begin{tabularx}{0.7\textwidth}{X} 弟兄們，我不是以為自己已經得著了；我只有一件事，就是忘記背後，努力面前的， \end{tabularx} \\ \\ \relax
3:14 & \begin{tabularx}{0.7\textwidth}{X} 向著標竿直跑，要得神在基督耶穌裡從上面召我來得的獎賞。 \end{tabularx} \\ \\ \relax
3:15 & \begin{tabularx}{0.7\textwidth}{X} 所以，我們中間凡是成熟的人，總要存這樣的心；若在甚麼事上存別樣的心，神也會把這些事指示你們。 \end{tabularx} \\ \\ \relax
3:16 & \begin{tabularx}{0.7\textwidth}{X} 然而，我們達到甚麼地步，就當照這個地步行。 \end{tabularx} \\ \\ \relax
3:17 & \begin{tabularx}{0.7\textwidth}{X} 弟兄們，你們要一同效法我，也當留意看那些效法我們榜樣的人。 \end{tabularx} \\ \\ \relax
3:18 & \begin{tabularx}{0.7\textwidth}{X} 因為，我屢次告訴你們，現在又流淚告訴你們：許多人行事是基督十字架的仇敵。 \end{tabularx} \\ \\ \relax
3:19 & \begin{tabularx}{0.7\textwidth}{X} 他們的結局就是滅亡。他們的神明是自己的肚腹；他們以自己的羞辱為光榮，專以地上的事為念。 \end{tabularx} \\ \\ \relax
3:20 & \begin{tabularx}{0.7\textwidth}{X} 我們卻是天上的國民，並且等候救主，就是主耶穌基督從天上降臨。 \end{tabularx} \\ \\ \relax
3:21 & \begin{tabularx}{0.7\textwidth}{X} 他要按著那能使萬有歸服自己的大能，把我們這卑賤的身體改變形狀，和他自己榮耀的身體相似。 \end{tabularx} \\ \\ \relax
4:1 & \begin{tabularx}{0.7\textwidth}{X} 我所親愛、所想念的弟兄們，你們就是我的喜樂，我的冠冕。我親愛的，你們應當靠主站立得穩。 \end{tabularx} \\ \\ \relax
4:2 & \begin{tabularx}{0.7\textwidth}{X} 我勸友阿蝶和循都基要在主裡同心。 \end{tabularx} \\ \\ \relax
4:3 & \begin{tabularx}{0.7\textwidth}{X} 我也求你這真實同負一軛的，要幫助這兩個女人，因為她們在福音上曾與我、革利免和我其餘的同工一同勞苦，他們的名字都在生命冊上。 \end{tabularx} \\ \\ \relax
4:4 & \begin{tabularx}{0.7\textwidth}{X} 你們要靠主常常喜樂。我再說，你們要喜樂。 \end{tabularx} \\ \\ \relax
4:5 & \begin{tabularx}{0.7\textwidth}{X} 要讓眾人知道你們謙讓的心。主已經近了。 \end{tabularx} \\ \\ \relax
4:6 & \begin{tabularx}{0.7\textwidth}{X} 應當一無掛慮，只要凡事藉著禱告、祈求和感謝，將你們所要的告訴神。 \end{tabularx} \\ \\ \relax
4:7 & \begin{tabularx}{0.7\textwidth}{X} 神所賜那超越人所能了解的平安，必在基督耶穌裡，保守你們的心懷意念。 \end{tabularx} \\ \\ \relax
4:8 & \begin{tabularx}{0.7\textwidth}{X} 末了，弟兄們，凡是真實的、凡是可敬的、凡是公義的、凡是清潔的、凡是可愛的、凡是有美名的，若有甚麼德行，若有甚麼稱讚，你們都要留意。 \end{tabularx} \\ \\ \relax
4:9 & \begin{tabularx}{0.7\textwidth}{X} 你們從我所學習的，所領受的，所聽見的，所看見的事，你們都要繼續去做，賜平安的神就必與你們同在。 \end{tabularx} \\ \\ \relax
4:10 & \begin{tabularx}{0.7\textwidth}{X} 我靠主大大喜樂，因為你們關懷我的心如今又表現了出來；其實你們一直都關懷我，只是沒有機會罷了。 \end{tabularx} \\ \\ \relax
4:11 & \begin{tabularx}{0.7\textwidth}{X} 我並不是因缺乏而說這話，因為我已經學會無論在甚麼景況都可以知足。 \end{tabularx} \\ \\ \relax
4:12 & \begin{tabularx}{0.7\textwidth}{X} 我知道怎樣處卑賤，也知道怎樣處豐富；或飽足或飢餓，或有餘或缺乏，任何事情，任何景況，我都得了祕訣。 \end{tabularx} \\ \\ \relax
4:13 & \begin{tabularx}{0.7\textwidth}{X} 我靠著那加給我力量的，凡事都能做。 \end{tabularx} \\ \\ \relax
4:14 & \begin{tabularx}{0.7\textwidth}{X} 然而，你們能和我分擔憂患是一件好事。 \end{tabularx} \\ \\ \relax
4:15 & \begin{tabularx}{0.7\textwidth}{X} 腓立比人哪，你們也知道我開始傳福音、離開馬其頓的時候，在收支的事上，除了你們以外，並沒有別的教會和我分擔。 \end{tabularx} \\ \\ \relax
4:16 & \begin{tabularx}{0.7\textwidth}{X} 就是我在帖撒羅尼迦，你們也一再差人來供給我的需用。 \end{tabularx} \\ \\ \relax
4:17 & \begin{tabularx}{0.7\textwidth}{X} 我並不求甚麼饋贈，只求你們的果子不斷增多，歸在你們的賬上。 \end{tabularx} \\ \\ \relax
4:18 & \begin{tabularx}{0.7\textwidth}{X} 但我已經如數收到，並且有餘；我已經充足，因我從以巴弗提受了你們的饋贈，當作極美的香氣，為神所接納、所喜悅的祭物。 \end{tabularx} \\ \\ \relax
4:19 & \begin{tabularx}{0.7\textwidth}{X} 我的神必照他榮耀的豐富，在基督耶穌裡，使你們一切所需用的都充足。 \end{tabularx} \\ \\ \relax
4:20 & \begin{tabularx}{0.7\textwidth}{X} 願榮耀歸給我們的父神，直到永永遠遠。阿們！ \end{tabularx} \\ \\ \relax
4:21 & \begin{tabularx}{0.7\textwidth}{X} 請問候在基督耶穌裡的各位聖徒。跟我一起的眾弟兄都問候你們。 \end{tabularx} \\ \\ \relax
4:22 & \begin{tabularx}{0.7\textwidth}{X} 眾聖徒都問候你們，特別在凱撒家裡的人問候你們。 \end{tabularx} \\ \\ \relax
4:23 & \begin{tabularx}{0.7\textwidth}{X} 願主耶穌基督的恩與你們的靈同在！ \end{tabularx} \\ \\
[1ex]
\hline
\hline
\end{longtable}
$^{1}$各位體驗主持人平安.
我會用一次的時間.
來介紹保羅書信當中.
很重要的一卷書.
這卷書往往被稱為.
喜樂的祕訣.
如果你想知道.
如何才能夠喜樂呢.
喜樂有什麼方法.
有什麼祕訣呢.
那你一定要讀菲勒比書.
在我們人生裡面.
我們時時都會遇到很多的困難.
挫折.
或者是一個失敗感.
沮喪的裡面.
覺得壓力很大.
我們沒辦法勝過.
有時候我們的仇敵就很厲害.
想用各樣的方法.
將我們打低.
因此他叫我們覺得自己.
沒力氣.
抬不起頭.
掙扎不起來.
但是保羅很特別.
他竟然在坐牢的時候.
竟然鼓勵弟兄姊妹.
要靠主喜樂.
而且是常常喜樂.
有時候我們會問.
為什麼保羅這麼厲害呢.
坐牢都會喜樂的.
而且他那種喜樂.
是由內心裡面湧流出來的.
不是環境帶給他的.
你坐牢.
又沒有好東西.
又沒有自由.
各樣東西都不行.

$^{41}$又受鎖鏈綁住.
你真的會覺得.
是很大的打擊.
但是保羅他是由心裡面.
發出一份喜樂.
因此我們要問.
為什麼保羅可以能夠靠主喜樂.
你想不想要同樣可以靠主喜樂.
菲律賓這家教會.
是保羅第二次宣教行程的時候.
他第一站就去的地方.
記不記得有一晚.
保羅見到一個馬其頓的異象.
因為神禁止他.
不要在小亞細亞繼續傳道.
於是他那一晚就做了個夢.
見到一個馬其頓人.
站在他的地方.
向保羅招手.
他說你過來馬其頓.
幫助我們.
所以我們通常都叫這個異象.
馬其頓的異象.
保羅聽見歐洲的呼聲.
於是他轉變他宣教的方向.
不再在小亞細亞那裡走.
就去了歐洲.
馬其頓那方的頭一個大城.
就是菲律賓.
菲律賓的城位處一個很關鍵的位置.
那是一條東西的通道.
一條很大的通商要道.
所以菲律賓的城市.
是在一個通商的關鍵的位置.
所以菲律賓是非常發達的一個城市.
當保羅去到馬其頓的菲律賓的時候.
他就在安息日.
去找有沒有猶太人的會堂.
結果真的被他找到了.
原來有一幫猶太的婦女在那裡聚集.

$^{81}$他們其中一位的領袖.
就是賣紫色布的一位女士.
就是Lidia女帝.
大家都知道紫色是一種很高貴的顏色.
而這種紫色布要由一種特別的染料.
把它染成紫色.
由於這種是在貝殼裡面.
提煉出來的染料.
所以是特別名貴的.
很貴纔買得到這種染料.
所以如果賣紫色布.
通常都是一些尊貴的婦女.
有錢人家.
好像今天我們說.
你買這麼漂亮的布料來做衣服.
所以是高價貨.
而這位女帝雅.
就從事一個這麼高級的商業行業.
因此女帝雅就向保羅發出一個邀請.
她說你和你的幫手可以來我家住.
你不要客氣.
你來我家住.
於是保羅就住在這位女帝雅姐妹的家.
接受她的款待.
因此保羅和菲納比教會之間.
有很緊密的關係.
保羅其實只是在菲納比留了很短一段時間.
她就離開去其他地方報道了.
因此這封菲納比這封書信.
菲納比教會知道保羅坐牢.
於是就派了一個代表.
這個代表就是尹巴忽提.
尹巴忽提有雙重的責任.
第一就是帶教會的消息去告訴保羅.
你放心吧.
我們教會很好.
第二就是帶著這家教會的經濟支持.
就是帶一筆錢特別去送給保羅.
所以保羅真的很感動.
這家教會不僅在精神上支持她.

$^{121}$更加在經濟上很實質的來支持她.
所以保羅在菲納比書裡面.
就贊這家教會對她有很多的幫助.
她說你們真的在經濟上幫助了我.
但接著保羅說.
我無論在什麼情況下.
我都學會知足.
在負足的時候我知足.
在缺乏的時候我同樣知足.
哇!保羅的人生哲學.
可以教我們很多功課.
各位弟兄姊妹.
你有沒有知足呢?.
無論處於富裕,困難.
很穩妥的時候.
或者很艱難的時候.
我們都要學會知足.
所以保羅在菲納比書.
一卷這麼重要的書裡面.
不單止說起落的祕訣.
更加她強調我們和基督之間的關係.
菲納比書是一卷高舉基督的書.
菲納比書一共有四章.
第一章保羅一開始就說明.
他是要向他們傳福音的.
他說有些人傳福音.
是出於一些自私的理由.
或者是為了自己的緣故.
所以很多是分黨結派.
我是屬於什麼的.
我是屬於誰的.
我和誰是一批人的.
保羅說無論你傳福音.
你是為了自己.
為了什麼的名利.
無論如何.
究竟基督被傳開了.
因此我就要歡喜.
沒錯有些人在傳道工作裡面.
有一點點的私心.

$^{161}$或者他為了自己個人的利益的緣故.
為了個人的名字的緣故.
但無論如何保羅說.
到底基督就被傳開了.
然後保羅又說到他對死的看法.
所以很多人估計.
這卷的《肥臘庇書》.
是保羅比較晚年才作的.
從第一章19,20節開始.
保羅就說到如何面對死亡.
他說或生或死.
總叫基督在我身上照相顯大.
我死了是有益處的.
因為我死了是毫得無比的.
因為我與基督一起.
保羅他不怕死.
他對死亡有一個很清晰的看法.
他知道他走完這條世上的路.
他去哪裡.
各位弟兄姊妹.
這就是我們的把握.
這就是信仰帶給我們的盼望.
很堅定的盼望.
我們知道我們將來去哪裡.
有幾樣東西我們是沒辦法知道的.
你什麼時候來到這個世界.
你爸爸媽媽什麼時候生你.
你什麼時候出世.
你不知道的.
同樣我們走.
我們離開這個世界也是不知道的.
有些人說.
什麼這麼沒人情的.
應該告訴我.
我什麼時候就走.
我老實說.
如果我告訴你.
下個月你就走了.
你這個月不用睡覺了.
亂亂的.

$^{201}$我就走了.
我還有很多事沒做.
糟糕了.
我不知道我太太的子女怎麼樣.
我就走了.
不要走了.
我走了.
所以其實神很大的恩典.
祂沒告訴我們.
我們什麼時候來到這個世界.
同樣祂都沒告訴我們.
你什麼時候會離開這個世界.
完全在神的掌握裡面.
有些人離開這個世界是很突然的.
譬如說遇到一些交通意外.
你想不到.
你怎麼知道有些車會撞壞你.
你就突然之間就回到天家.
但是在神的時間.
神的預備裡面.
還是最好的.
保羅對死亡一點都不畏懼.
他帶著一個很坦然的心.
來面對死亡.
所以他說.
我活著是對你們有益.
我死了是有好處.
因為我回到天上.
與主同在.
是好得無比的.
到了第二章.
保羅就提出.
我們應該要以基督為我們的榜樣.
耶穌基督就是我們的榜樣.
我們要效法祂.
我們要以耶穌基督的心為心.
所以我們要學基督.
要像基督.
這個就是我們教徒的人生.
各位你像不像耶穌呢?.

$^{241}$你有沒有想著我要效法耶穌.
我要好像耶穌那樣呢?.
人家見到我.
這個人怎麼這麼溫柔.
這個人怎麼這麼磁場.
這個人怎麼這麼仁愛.
這個人怎麼帶人節物這麼好的.
讓我們效法基督.
讓我們以耶穌的心為心.
就好像明光照耀在地上一樣.
我們都是世上的光.
光是要顯起黑暗的地方.
越黑暗的地方.
光就更加顯明.
到了第三章.
保羅又提出.
我以耶穌基督為自保.
我以耶穌基督為我的標杆.
為我的目標.
我一心奔跑.
我要得著基督.
好像基督得著我的一樣.
我向著標杆直跑.
我是忘記背後.
努力面前.
向著標杆直跑.
保羅那種奮鬥.
保羅那種前進.
保羅那種劫力.
為基督而活.
因此我都勉勵各位弟兄姊妹.
你的人生同樣要有目標.
我們人活著是為了什麼呢?.
在一本《教理》的文章裡.
第一條就是問.
人活著是為了什麼?.
你活著的目標是什麼?.
答案是什麼?.
我們活著是為了榮耀神而活.
因此我的生命.

$^{281}$以榮耀神為目標.
所以保羅說.
我只求什麼?.
我活著就是基督.
我以基督耶穌的心為心.
而最後第四章.
保羅在結束的時候就提到.
要上上起落.
要靠主上上起落.
原來基督就是我們的能力.
我們因為有基督在我們裡面.
我們充滿能力.
我們可以靠主起落.
事事起落.
而且我們要期待耶穌快點回來.
今天是一個末世的時代.
到處都有戰爭.
而且這些戰爭越來越複雜.
所以我們更加需要警醒.
等候主的再回來.
保羅說應當一無掛慮.
將你所要的告訴神.
神所賜出人意外的平安.
必在基督耶穌裡.
保守你們的心懷意念.
沒錯這個世界.
很多東西叫我們掛慮.
我們掛這一樣掛那一樣掛那一樣.
但是聖經說.
交託給神.
原來我們凡事就無憂慮.
神會將一種出人意外的平安.
放棄我們裡面.
叫我們靠主可以跨盛.
求主將費納比書.
一卷金特別的書.
一卷很重要的書.
放在各位的心裡.
讓你常常的反覆思想.
讓你將神的話藏在心裡.

$^{321}$免得你得罪神.
求神使用你.
祝福你 願神賜福給各位.
\newpage



\section{約翰福音 20:24-29}
\label{sec:031941}
\textbf{復活主榮耀的記號}~~~~甘雪芬傳道
\newline
\newline
連結: \href{http://www.church.com.hk/acms/content.asp?site=cdc\&op=show\&type=product\&code=031941\&layout=sermon}{\texttt{031941}} ~~~~ 語音日期: 2026-04-05
\newline
\newline
\hyperref[sec:031940]{\small{< < < PREV SERMON < < <}}
~
\hyperref[sec:index_chronic]{\small{[返順時目]}}
~
\hyperref[sec:031942]{\small{> > > NEXT SERMON > > >}}
\newline
\newline
約翰福音 20:24-29
\newline
\begin{longtable}{cl}
\hline
\hline
章節 & 經文 (和合本修訂版)\\
\hline
20:24 & \begin{tabularx}{0.7\textwidth}{X} 那十二使徒中，有個叫低土馬的多馬，耶穌來的時候，他沒有和他們在一起。 \end{tabularx} \\ \\ \relax
20:25 & \begin{tabularx}{0.7\textwidth}{X} 其他的門徒就對他說：「我們已經看見主了。」多馬卻對他們說：「除非我看見他手上的釘痕，用我的指頭探入那釘痕，用我的手探入他的肋旁，我絕不信。」 \end{tabularx} \\ \\ \relax
20:26 & \begin{tabularx}{0.7\textwidth}{X} 過了八日，門徒又在屋裡，多馬也和他們在一起。門都關了，耶穌來，站在當中，說：「願你們平安！」 \end{tabularx} \\ \\ \relax
20:27 & \begin{tabularx}{0.7\textwidth}{X} 然後他對多馬說：「把你的指頭伸到這裡來，看看我的手；把你的手伸過來，探入我的肋旁。不要疑惑，總要信！」 \end{tabularx} \\ \\ \relax
20:28 & \begin{tabularx}{0.7\textwidth}{X} 多馬回答，對他說：「我的主！我的神！」 \end{tabularx} \\ \\ \relax
20:29 & \begin{tabularx}{0.7\textwidth}{X} 耶穌對他說：「你因為看見了我才信嗎？那沒有看見卻信的有福了。」 \end{tabularx} \\ \\
[1ex]
\hline
\hline
\end{longtable}
$^{1}$在復活節的主日 向各位問安 祝你平安.
願意我們帶著一份感恩 謝喜樂的心.
讓上帝的說話成為我們的鼓勵 我們同心的禱告.
在復活節的主日 我們去回望我們走過的信仰路.
在上帝你的恩典 白白的恩典 讓我們今天有這份福氣.
帶著一份謝喜樂 讚頌的心.
願上帝你的話語 成為我們的激勵 以致無論講的.
聽的也同為蒙福 聽我們的禱告 奉教的救主 耶穌基督.
寶貴聖名去求 阿們.
大家身體上有沒有留下一些記號呢?.
我二十多歲的時候 做了一個手術.
所以這個肚腹 現在仍然有一個包含在這裡.
記得那時候 醫生帶著一些很大的歉意 很鄭重地跟我說.
他說 因為有實際的需要 他要打直幫我開這個傷口 不可以打橫.
一方面 癒合需要比較多的時間.
他更加著意的是 見我當時也是二十多歲的年輕人.
如果我會穿一些露出肚腹的服飾 可能就遮不到包含.
他怕我介意 好像身體留下了一點點瑕疵.
同樣 在肚子上的包含 對一個開刀 生孩子的媽媽來說.
相信有著完全不一樣的意義.
她絕對不會看來這是一個瑕疵 一個羞辱.
是一個美麗 偉大 愛的印記.
今天是復活節主日 主耶穌復活升天之前.
曾經向很多人去顯現.
問問大家 如果你當時有機會見到復活主.
你會預期是一個怎樣的形象在你面前出現呢?.
既然升過了死亡 是不是應該散發著光芒.
有一個威榮的形態去顯現呢?.
你看他剛才讀過的經文.
復活主留下了手上的釘痕 立旁的傷痕作為記號.
展現在門徒的面前 為何要留下這些傷痕?.
這是一個羞辱的記號 還是榮耀的記號?.
讓我們著今天的經文思想.
我們重點會看約翰福音20章24至29節.
耶穌向多瑪顯現的場景.
我們先了解一下上文.
約翰福音來到20章 開始描述耶穌復活之後的場景.
作者刻意描述了當中不同人物 他們看見不同的景況.
我們握腰來概念一下.
耶穌復活的當天的清早 穆大拉瑪利亞來到墳墓.

$^{41}$他看見墳墓的石頭被人挪開了.
他跑去找西門彼得和耶穌所愛的那個門徒.
兩個門徒走到墳墓裡.
在看見墳墓裡 耶穌的屍體不見了.
但果頭巾和細麻布仍然捲著放在那裡.
聖經記載 那個門徒就信了.
瑪利亞在墳墓外哭 最初看見兩個天使.
然後看見能夠說出祂名字的男士.
最後認出祂就是耶穌.
於是穆大拉瑪利亞就向門徒報信.
我已經看見主了 有福術主耶穌跟他說的一番話.
我要升上去見我的父 也是你們的父.
見我的神 也是你們的神.
如果你的門徒聽到這番信息.
你的心情會是怎樣呢.
是不是應該說 哇 謝 好喜樂 主復活了.
聖經描述 就是剛才的輔助經文.
在穆大拉瑪利亞告訴門徒看見主.
同一個晚上 當時的門徒只是聚聾在一起.
但不是為了這個信息而興奮 而是什麼呢.
他們仍然害怕 害怕什麼 害怕猶太人.
可能是當時一些宗教領袖或官員.
所以身處的地方的門全部都關上了.
不過中原們是觀賞耶穌可以出現在門徒的當中.
但耶穌不是一個幻影 而是帶著復活的身體跟他們說.
願你們平安.
還向他們展示他的手和肋膀 見證著已經戰勝了死亡.
聖經描述 門徒一看見主就喜樂了.
接著門徒領受聖靈 被差遣領受使命.
接著我們進入今天的經文24-29節 耶穌向多瑪顯現的場景.
第24節 那十二使徒中 有個叫底土瑪的多瑪.
耶穌來的時候 他們有和他們在一起.
作者刻意去提出身為十二使徒之一的多瑪.
當耶穌向他顯現的時候 沒有和門徒在一起.
雖然沒有明確的表現 在暗示著 身處一個這樣的非常時期.
他不應該不和其他門徒在一起.
25節 門徒告訴多瑪這個大好的消息 我們已經看見主了.
並且復述主如何向他顯示他的手和肋膀.
但多瑪並未說服到他不相信 多瑪這樣回應.
除非你讓我看見他手上的釘痕 用我的指頭探入那個釘痕.

$^{81}$用我的手探入他的肋膀 我絕對不會相信.
多瑪要求自己親身去驗證 他不但要看 還要甚麼呢.
還要磨這個物理學上的驗證.
意味著他覺得耶穌復活的這個信息 通不過去理性的思維.
可能他覺得這群門徒會不會其實是個錯覺.
又或者 耶穌只是一個幻影 一個靈魂.
不可能帶著身體復活在人面前.
26至27節 過了八天 門徒同樣聚集在房間.
門同樣地關上 耶穌再次的顯現 同樣向門徒說.
願你們平安 一切重復著早前.
耶穌向門徒顯現的場景 唯一不同的是.
這次多瑪在場了.
耶穌第二次的顯現是特別為了多瑪.
耶穌沒有責備他 而是撫咒他的軟弱.
邀請他做 說他要做的事.
於是耶穌和多瑪說.
將你的指頭伸到我這裡來 看看我的手.
將你的手伸過來 看我的肋膀.
一次中的話 不要疑惑 你總要相信.
聖經記載到 多瑪是否真的用手去摸.
但他哪怕在耶穌面前做了一個很真切的宣告.
28節 多瑪回答跟他說什麼.
我的主 我的神.
很奇怪 我們仿佛看見多瑪從一個信仰路上.
一個低谷 突然之間驟然升上來的一個轉捩點.
多瑪如何從失落 懷疑裡面.
轉化了這個信心的宣告.
說得出 我的主 我的神.
我們嘗試揣摩一下多瑪走過的心路歷程.
約翰夫音的上文另外有兩處記載了多瑪.
在兩個不同的場景 凸顯到他的個性.
第一處記載在十一章.
當時有人告訴耶穌 他的好朋友拉薩路病重.
兩天後傳來拉薩路的死信 耶穌決意要去白大利.
當時猶太宗教領袖正在密謀要害耶穌.
在門徒一片的反對聲裡 多瑪這樣說.
我們就去跟他一起死吧.
表面看來 多瑪不滿意耶穌的決定 有些很晦氣的說話.
但畢竟他用行動 連用行動跟著耶穌叫做忠心跟隨.
盡上門徒的責任 有些晦氣而已.

$^{121}$另一處記載在約翰夫音的十四章.
十家路越走越近 耶穌知道跟門徒一起的時日無多.
就透露他很快要回到天父那裡去.
他跟門徒說 我要去的地方 那條路你們是知道的.
多瑪怎麼回應呢 主啊 我都不知道你去那裡.
怎麼會知道那條路啊 好像很小孩子似的.
耶穌作出一個我們很熟悉的回應.
我就是道路 真理 生命 若不藉著我.
沒有人能夠到父那裡去.
多瑪的反應表達出一個很率直的本性.
不過又確實缺乏一些屬靈的智慧和動見.
好了 來到我們今天第三段經文.
就是今天看的一段經文.
按當時的背景 我相信 多瑪跟很多門徒一樣.
本來期盼著耶穌在政治上可以去復興以色列國.
但是竟然眼前的這位耶穌被釘十架的現實.
粉碎了他們的夢想 陷入一個很絕望的低谷.
當復活主帶著有釘痕的肉身.
竟然顯現在他們面前的時候.
我們可以想像多麼震撼.
多瑪早前向其他門徒理直氣壯地說.
要這樣那樣 這樣的條件我才會相信 是不是.
她說這番話的時候 耶穌明明不在場.
但是如今眼前的耶穌瞭如指掌.
就去滿足你說要這樣那樣的要求.
展示了他一份超自然的力量.
多瑪最初要求一個很理性的驗證.
但體現了耶穌 超越人理解的神聖.
多瑪看到自己 一個自以為是.
這份高傲 無知 愚昧.
內心浮現出一份錯愕 羞愧.
顯示一份敬畏 帶著信心的跳躍.
去宣告 我的主 我的神.
這個宣稱給予多瑪激勵了個人.
與主耶穌自己的相遇.
以確信主耶穌對方的神聖.
願意將自己的生命投放在神的主權帶領下.
耶穌按著我們每一個人很獨特的景況.
介入我們的生命 帶來生命的改革.
復活主母按著這班門徒的心意.

$^{161}$成就他們的願望 帶來地上的政治改革.
卻是揭示了不一樣的屬靈視野.
我相信多瑪 她慢慢去回憶.
去整理耶穌以往的教訓.
慢慢就謊言大悟.
耶穌所說的德性.
不是推翻羅馬政權 是戰勝死亡.
多瑪的例子讓我們看見.
可能有人會對信仰提出各樣.
通不過去理解的疑問.
回應這些疑問不是沒有意義.
不過等於你回應了疑問 他就會相信.
唯有人願意去開放 讓神自己的工作.
在神的光照底下.
人親自去遇見 去經驗神.
才能夠被說服 真正去體會.
所以老是一個歷程.
你往往可能土象多瑪.
又要求各樣的驗證.
有時沒有朋友也會這樣說.
如果我家人的病得到醫治.
如果我的工作可以順利一點 我就相信了.
帶著這樣的期盼.
我也是人之常情.
上帝確實也願意去體貼我們的需要.
我們的心意去成就.
就算我們信了一段日子.
我們何嘗不是常常帶著我們的限制.
將復活主的作為.
框在我們自己想要的那個框框裡.
但我們總不是原地踏步.
我們追求屬靈生命不斷的蛻變.
不斷的成長.
回望我自己走過的信仰歷程.
可能我也分享過.
在2003年SARS的時候.
因為患了情緒病.
希望得到醫治 所以我步入教會.
當時心底裡.
何嘗在想.

$^{201}$如果這個神醫到我的病就好了.
問我的信心 我的信是什麼.
我根本不知道 是不是.
我只是想病得醫治而已.
好了 過了一段日子.
初信的時候.
剛好工作遇到一些不太穩定的事.
我向上帝這樣祈禱.
你什麼都可以搞我的.
你不要碰我的工作 拜託你.
你說我信什麼.
耶穌是我的救主.
寫給我最有永恆的盼望.
我當然要信.
作為生命的救主.
仍然要在這裡拉鋸一下.
又過了一段日子.
直至我很真誠地體會上帝的愛.
心甘命抵自己說要辭了工作.
要進神學院去接受裝備.
這個時候我想我稍稍體會到.
什麼是我的主 我的神.
而這份體會是.
「現在進行式」.
直至很近很近.
跟大家分享一個很深切的體會.
很多個月之前.
藉著一些事件我經歷到我的生命.
被上帝的拆毀和建立.
當時很不容易.
當我已經感到很沉重的時候.
執事會的主席竟然走來.
向我提出邀請.
你會考慮承擔作為其中一位.
管理委員我們叫MC的角色.
我自覺我不是一個管理的材料.
我不明白上帝你為什麼要搞.
這麼大的工程又拆毀又建立.
現成的不就好了嗎.
和上帝有很多拉鋸.

$^{241}$甚至有很多淚土.
花了一段時間.
懇切的禱告祈求.
引證上帝將我推向去接受.
這個挑戰.
推向其中一項埋門的那一腳.
我的下午.
執事會主席又跟我說.
說回你祈禱了很久.
何時可以有一個最終的決定.
就在那天晚上.
我在地鐵站的門外.
我遇見一位我沒見很久的.
新學院的舊同學.
他見面的時候竟然什麼都不說.
第一句就問.
你還在窩福事房嗎.
他竟然記得.
就知道你的堂主任說要退休.
找到新的堂主任沒有.
什麼都不問.
只是問這件事.
我說沒有.
有些叫MC在那裡暫時管理.
你服侍了這麼久.
你一定是其中一位MC.
我說不是.
是的 你一定是.
就這樣寒暄幾句.
就走了.
很奇怪.
當天的晚上.
我就回應了我的上帝.
我沒想過上帝為我預備的出路.
不是去減輕我的生命.
被他修剪過程.
那種其實有點疼痛.
而是將更大的目標.
擺在我面前.
叫我為了這個目標去學習.

$^{281}$我對我的主.
我的神.
有一個很新的相遇.
而我起碼.
在這個時刻.
現在我的主.
我的神在我生命裡面.
大過了.
他佔據了更加大的位置.
當然我不能夠擔保.
是不是經常都會將盡.
信仰是一個歷程.
今天的經文描述到.
當門徒知道.
耶穌已經復活勝過死亡.
理應興奮.
但何嘗卻害怕.
躲在門要關上的房間裡.
多瑪仍然疑惑.
不過我們在落馬書之後.
我們要看《仕途行傳》.
在《仕途行傳》.
我們看到這班門徒是怎樣的.
剛強勇敢地去擔當.
建政福音的工作.
我們當中可能有著.
我們現在不同的屬靈光景.
可能有些仍然你看他不敢.
仍然好像有點軟弱.
看自己.
對自己都可能.
我們仍然有懷疑.
仍然有害怕的時候.
但一個真誠願意尋求神的心.
我相信神會不停在你生命裡面工作.
多瑪表現出疑惑.
那復活主怎樣去回應呢.
二十九節.
耶穌跟他說.
你因為看見了我才信嗎.

$^{321}$那沒有看見卻信的有福了.
大家認為耶穌這個回應.
是不是在責備多瑪呢.
二十九節.
是不是在罵他呢.
傳統上我們很喜歡給多瑪一個.
負面的綽號叫做什麼.
多兒的多瑪.
或者耶穌是不是在比較信心的根基.
有高低之分.
或者有不同的詮釋.
我按我的理解.
針對多瑪不相信其他門徒的見證.
要自己親身去驗證.
耶穌這個回應.
確實帶著一些的提點.
不過我覺得頂多是一個溫和的訓練.
不是真正的責備.
其他的門徒.
一班的婦女.
當中有些最初何嘗米.
帶著失望.
傷心.
誤解.
疑惑.
大悟.
停滯不前.
也是一步一步邁向.
更全面更成熟的相信.
耶穌說.
那沒有看見卻信的有福了.
我相信是藉著多瑪的場景.
引出當時的背景一個.
即將發生很大的轉變.
因為復活主.
耶穌很快要升天離開門徒.
跟著祂的人.
不可能再憑肉眼看到耶穌.
看到耶穌所行的神蹟奇事.
這班門徒的信心.

$^{361}$要面對很大的考驗.
耶穌要門徒步入一個.
更加成熟的信仰歷程.
不要只想著看得見.
那沒有看見卻信的有福了.
這番說話.
也是對當時約翰福音的讀者.
而說的.
如果我們看下文二十章的.
三十節三十一節.
耶穌在他門徒面前.
另外行了許多神蹟.
沒有記錄在這書上.
但記載這些事.
是要使你們信.
耶穌基督是神的兒子.
並且使你們信他.
好欣著他的名得生命.
耶穌復活升天之後.
人不可能再見到他.
但人如果能夠因為.
門徒所傳的道.
神的話語試圖見證.
而他信福音的人.
他們就有福氣了.
這正正是約翰要寫.
約翰福音的目的.
門徒確實有特權.
曾親身見過耶穌的神蹟.
其事聽過很多教導.
不過同時代的.
猶太宗教領袖.
何嘗不是親眼見到.
耶穌的所言所行.
沒有因為這樣而信.
反而怎樣.
要密謀殺害他.
信心不是單單肉眼的看見.
而是屬靈的看見.
透過聖靈的工作.

$^{401}$信心去跟上帝的關係.
一段段屬靈的朝聖之旅.
隱形的時候我問大家.
你會預期復活主.
怎樣顯現在人前.
其實想真一層.
復活主大可以用.
很輝煌,尊榮,大能的形象.
在這裡顯現.
即是夾著雲這樣降下來.
然後有些光,耶穌光.
但是選擇向門徒顯現的.
是幫門徒害怕的時候.
而且他手上的釘痕.
肋膀的傷痕.
仍然留在這裡.
神大可以挪走的.
為什麼復活主要選擇.
保留這些作為一個記號.
我想有多重的意義.
首先復活主.
不是一些沒有身體的靈魂.
好像有些人所說的.
一個幻影.
不是人的想像.
復活主是會說話.
人能夠看見,摸得到.
他人類歷史裡.
是道成肉身的神.
真真實實來到歷史當中.
曾經為我們罪.
被釘死在十架上的耶穌.
帶著同樣的身體.
留下一些記號.
作為復活的證據.
其次.
耶穌保留這些傷痕.
標誌著他願意為他者而活.
為救贖世人.
甘願犧牲受苦.

$^{441}$耶穌復活.
又戰勝了死亡和罪的權勢.
釘痕.
不是一個羞辱的記號.
成為復活主展現.
祂犧牲的愛的一個印記.
什麼是榮耀?.
在約翰福音.
耶穌被賣的那一夜.
被賣的那一夜.
祂向父神這樣禱告.
父啊!時候到了.
願你榮耀你的兒子.
使兒子也榮耀你.
是被賣的那一夜.
說這番話.
耶穌的被賣.
榮耀父神.
父神也榮耀祂.
耶穌也行過很多的.
神蹟騎士.
被人歡呼鴿道進入城.
但這不是祂最榮耀的復事.
祂最榮耀的復事.
祂指認的榮耀不是能力.
不是地位.
而在愛和信服父神的旨意裡.
是受苦僕人的形態.
去完成使命.
耶穌要門徒去體會.
祂手上的釘痕.
是復活住一個.
很榮耀的人.
當耶穌升天.
人不可能再親眼目睹耶穌的時候.
那怎麼辦?.
門徒成為了第一代的.
最重要的見證人.
用祂的身教言教.
用祂的生命呈現出.

$^{481}$基督信仰的要義.
讓我和你繼續成為復活主的見證人.
記得有一次.
有位姊妹才剛剛信主沒多久.
她的丈夫有天跟她來教會.
她說:看看教會是什麼一回事.
我怕她生氣.
原來不是.
發現太太信了主後.
很開心整個人都不同了.
她想看看你們教會教什麼.
很快她的丈夫都信了.
弟兄姊妹,我們的生命.
怎樣呈現到復活主的榮耀.
看到神在我們生命裡的工作的痕跡.
今天我們不可以.
用我們的五官摸到.
看到耶穌.
其實是一個挑戰.
我們可以生活得當神不存在.
特別當我覺得.
順境的日子.
靠自己可以自給自足的.
和神走得遠一點.
其實我覺得我自由一點.
不要搞我這麼多.
看我的Netflix.
讓我去激勵,為什麼我要去擺相.
我們需要不停.
有很多熟練的操練.
去敏感神的同在.
去實踐經歷神的愛.
真實的信心.
我相信是建基於聖言.
一些合符聖經的教導.
作為合神心意的基礎.
聖靈.
聖靈在我們心裡的工作.
建立和神的關係.
作為我們得力的源頭.

$^{521}$一班的聖民.
信徒的生命見證.
彼此的守望同行成為.
鼓勵和提醒.
在生活的各樣困難.
挑戰裡面.
我們當然可以向上帝祈求.
得幫助,得安慰.
我深信上帝也不論是.
我們的心意,供應我們的需要.
體恤我們的軟弱.
不過,如果你只是停留了.
求上帝去承受我的心意.
要怎樣怎樣我才信呢?.
我們的屬靈視野就跳不出.
我們自己這個框框.
我認識一位姿體.
他成長歷程裡面.
他經歷過家暴.
所以他覺得上帝對他不公平.
他想了很久很久.
他今天仍然不願意接受瑞麗.
也認識另外一位姿體.
同樣有著一個.
不愉快的成長背景.
因為他經歷到.
神的愛的依治.
漸漸走出傷痛.
帶著使命去服侍.
有同樣經歷的人.
我要強調我不是要去比較.
我不是要去比較兩個人靈魂的高低.
每個人要面對的.
掙扎都是很獨特.
外人不容易去理解.
我只是想提出.
神邀請我們用新的眼光.
去看我們的治癒.
去醫治我們的傷痕.
盼望可以轉化成為.

$^{561}$服侍其他人的資源.
弟兄姊妹.
當我們遇上這樣的挑戰.
可能你在教會.
要承擔一個新的崗位.
可能客觀環境.
很多的變動.
可能家人自己的健康.
沒有以前那麼壯健.
各樣的需要等等.
當我們被安放在一個.
靠我自己.
不能夠成就的光景.
我們被迫唯有信靠神.
成為我最大的出路.
也成為一個契機.
在這個過程中.
我們跟我們的主.
我們的神有一個新的相遇.
可能過程不容易.
在這個相遇中.
你的靈命就經歷到.
不一樣的成長或者突破.
耶穌被釘十家.
本來不折不扣.
是一個很大的羞辱.
怎麼會不是羞辱.
但因為承受了人類的救恩.
轉化成為了一個榮耀的記號.
耶穌沒有說過.
這個本來不是一個羞辱.
沒有說過不會痛.
耶穌邀請人.
「若有人要跟從我.
就當社己天天背起.
自己的十字架內跟從我」.
你去捨棄和承擔.
一定是並存的.
你都看不到有什麼.
要去承擔.

$^{601}$找不到原因為什麼要捨棄.
我簡單一個小例子.
星期五完了壽夫節崇拜.
肚子餓買點東西吃.
進了超市.
其實第一樣最吸引我的.
是什麼呢.
是杯麵.
然後是一些忌廉蛋糕.
做了一些掙扎之後.
我最後還是選了一個.
雞胸肉的雜菜沙律.
但惡化叫我吃一杯雪糕.
為什麼呢.
因為有更大的意義.
放在我面前.
更大的意義是什麼.
要關注我的健康.
所以我就放得下.
唯有.
我們看到一個.
有更大意義的使命.
等待我們去承擔.
我們可以用新的眼光.
去看我們的光景.
我們忍受到當下的難處.
我們放下當下的安舒.
我看不到.
更大的意義的使命.
我放下來幹什麼呢.
耶穌.
兩次向門徒顯現.
第一句說話.
都是什麼呢.
願你們平安.
當然這是一個當時猶太人.
很恆常的吻口語.
但其實有著很深層的意義.
平安源自希伯來文的.
Shalom.

$^{641}$因為人和神有美好的關係.
以至感到很圓滿.
生命有價值.
耶穌要門徒跳出他們自己.
期盼的框框.
發現神賦予給他們的人生意義.
耶穌.
同樣激勵我們每一個人.
我們找到人生更大的意義.
以至我們心甘命抵.
我們可以.
捨棄.
逼不下來.
逼就很辛苦.
那林多後書這樣說.
那林多後書五章十四十五節.
原來基督的愛激勵我們.
因我們這樣斷定.
一人既替眾人死.
眾人就都死了.
並且他替眾人死.
是叫那些活著的.
不再為自己活.
乃為替他們死而復活的主活.
我們背上十格的復活主.
他留下手上的釘痕.
用他犧牲的愛.
去激勵我們.
我們有個更大.
更有意義的使命.
擺在我們面前.
我們也去背負.
我們自己的十字架.
學後耶穌基督的樣式.
活出一個能夠.
祝福人榮耀神.
豐盛的生命.
我不知道復活主.
對他們有什麼話.
在他們心裡面有沒有.

$^{681}$上帝安放了一個更大.
更有意義的使命.
在他們心中裡面.
你願不願意為著這個使命.
因為它有意義的.
你心甘命抵.
或者至少你去學.
你去放下.
你忍受得到主.
在我面前有個新的使命.
我也在學著要去放下.
上帝每一個人有不同的心意.
不同的說話.
一起為著更大.
更有意義的使命.
我們學習去背上.
我們每個人不同的十字架.
一起去跟從耶穌.
這個復活主對我們的期盼.
的土局.
參的恩主.
你在十架上面.
去彰顯你的大愛.
你的愛成為我們的激勵.
我們能夠愛因為你神.
才愛我們.
願意我們在我們的生活.
的每一個層面裡面.
讓我們更加多去經歷.
上帝你的同在.
以致我們宣告.
我的主.
我的神的時候.
讓我們有更加盡心的體會.
我們去聽到上帝你的呼召.
在我們生命裡面.
擺放一個更加有意義.
更大的使命的時候.
我們就有能力.
去放下一些.

$^{721}$我們學習去捨棄的事.
背負我們每個人的十字架.
去跟隨你.
主我求你.
祝福我們的生命.
也讓我們的生命成為其他人的祝福.
聽我們的禱告.
耶穌基督的聖名去求.
阿們.
\newpage



\section{約翰福音 21:15-19}
\label{sec:031942}
\textbf{愛得更深}~~~~李炎強牧師
\newline
\newline
連結: \href{http://www.church.com.hk/acms/content.asp?site=cdc\&op=show\&type=product\&code=031942\&layout=sermon}{\texttt{031942}} ~~~~ 語音日期: 2026-04-04
\newline
\newline
\hyperref[sec:031941]{\small{< < < PREV SERMON < < <}}
~
\hyperref[sec:index_chronic]{\small{[返順時目]}}
~
\hyperref[sec:031943]{\small{> > > NEXT SERMON > > >}}
\newline
\newline
約翰福音 21:15-19
\newline
\begin{longtable}{cl}
\hline
\hline
章節 & 經文 (和合本修訂版)\\
\hline
21:15 & \begin{tabularx}{0.7\textwidth}{X} 他們吃完了早飯，耶穌對西門‧彼得說：「約翰的兒子西門，你愛我比這些更深嗎？」彼得對他說：「主啊，是的，你知道我愛你。」耶穌說：「你餵養我的小羊。」 \end{tabularx} \\ \\ \relax
21:16 & \begin{tabularx}{0.7\textwidth}{X} 耶穌第二次又對他說：「約翰的兒子西門，你愛我嗎？」彼得對他說：「主啊，是的，你知道我愛你。」耶穌說：「你牧養我的羊。」 \end{tabularx} \\ \\ \relax
21:17 & \begin{tabularx}{0.7\textwidth}{X} 耶穌第三次對他說：「約翰的兒子西門，你愛我嗎？」彼得因為耶穌第三次對他說「你愛我嗎」，就憂愁，對耶穌說：「主啊，你無所不知，你知道我愛你。」耶穌說：「你餵養我的羊。 \end{tabularx} \\ \\ \relax
21:18 & \begin{tabularx}{0.7\textwidth}{X} 我實實在在地告訴你，你年輕的時候，自己束上帶子，隨意往來；但年老的時候，你要伸出手來，別人要把你束上，帶你到不願意去的地方。」 \end{tabularx} \\ \\ \relax
21:19 & \begin{tabularx}{0.7\textwidth}{X} 耶穌說這話，是指彼得會怎樣死來榮耀神。說了這話，耶穌對他說：「你跟從我吧！」 \end{tabularx} \\ \\
[1ex]
\hline
\hline
\end{longtable}
$^{1}$到底愛是甚麼一回事呢?.
愛主耶穌又是怎樣的一回事呢?.
我想問大家.
如果有人問你.
你愛我媽媽.
特別是你旁邊那位老伴.
特別是弟兄.
被太太問的時候.
弟兄有甚麼感受呢?.
你不用告訴我.
你自己想一想.
我會找幾個弟兄躲在一起說.
為甚麼躲在一起說呢?.
你自己再去揣測一下.
的確有時候.
原來我們與人之間的那種情感.
特別對於我們中國人來說.
愛要說得出口.
是一點都不容易.
然而中國人對愛的描述.
又很有豐富又有深厚的感情.
有一句說話可能你也聽過.
「問世間情為何物,直教生死相許」.
有多少人聽過?.
也不錯,你們在哪裡聽過呢?.
金融當然是,對不對?.
所以其實金融也很有影響力.
不過我更加想你知道.
這首詞這首詞出自哪裡嗎?.
有沒有人認識?.
認識的話我請你喝茶.
不會吧?.
是不是?.
不會吧?.
這句是出自袁浩文的「硬憂詞」當中.
他問世間情為何物,直教生死相許.
這個描述是他看到一對大硬.
一隻死了.
而另一隻竟然投地來殉情.
他為了這個場景.

$^{41}$他很震撼於愛的力量.
原來竟然可以讓人生死相害.
是愛驅使這種決定.
今天我跟大家看了經文.
耶穌同樣都問了一個關於愛的問題.
不過不是問情是何物.
而是一個很直接地問彼得.
你愛我嗎?.
前面那個情是何物.
是一個哲學的構問.
一個問題.
一個思考.
而後者耶穌問彼得的.
他是一個很直視的.
很面對面的.
他要你回答的一個質問.
而耶穌的問題指向的不是情.
而是普遍性.
而是一個很深層很靠近的.
是你和我對耶穌.
對這位復活了的主耶穌.
最真實的回應.
或者我邀請大家在進入經文之前.
嘗試打開你的聖經.
或者請....
那裡是什麼來的?.
在音響那邊.
弄一個經文出來.
我嘗試用耶穌的對話來讀出.
你們就用彼得的說話讀出行不行?.
只是讀對話的內容.
我們嘗試一下感受那種張力.
好,準備好了.
約翰的兒子西門.
你愛我比這些更深嗎?.
你餵養我的小羊.
約翰的兒子西門.
你愛我嗎?.
你牧養我的羊.
約翰的兒子西門.

$^{81}$你愛我嗎?.
你們都很優秀.
不知道怎麼回答.
好,到此為止.
看不看到那種壓迫力在這裡?.
如果做丈夫的你.
你太太三次這樣問你.
到了第三次又有什麼感受呢?.
我們在同工會裡面討論這個問題的時候.
各有所看法.
不過這個不重要.
我想大家回到整個經文裡面.
給我們的啟發.
在這段說話裡面.
當我預備這段經文的時候.
你都給我回憶起.
在剛剛3月21日的時候.
報導會裡面.
長者報道裡面.
幾位弟兄姊妹分享了.
他們怎麼在艱難之中仍然堅持愛家人.
愛妻子.
愛他們的兒子.
在他們身上的愛.
我看到的不是很容易過得了.
是有疲倦的.
有時候有艱難的.
然而就是因為他們有耶穌基督的愛.
在他們心裡面不斷地去更新他們.
讓他們展現出這種愛的活力出來.
那麼愛是一個怎樣的力量呢?.
而這個印象.
也讓我去想起今天經文裡面.
耶穌問彼得的時候.
「你愛我媽」這句話.
彼得不是一個完美的門徒.
甚至他三次來到他不認耶穌.
但是這位復活的主耶穌.
沒有放棄他.
而是用一個很直率的問題.

$^{121}$將他心裡面對主耶穌的愛重新點起.
轉化他生命裡面的力量.
轉化他迎向上帝主耶穌給他的照命的力量.
今天就讓我們嘗試用一種回顧.
回憶的角度.
進入歷史的場景.
回到這個海邊.
再次細心看看這位復活的主耶穌.
怎樣去影響彼得的人生.
同時透過這個體會怎樣影響你和我.
讓我們就是先來同心吐槽.
參拜主我們去仰望你的恩典臨在.
願聖靈在我們心中帶給我們感動.
讓我們透過聖經的記載.
主耶穌你和彼得的相遇.
給我們提醒.
給我們啟發.
給我們挑戰.
重尋我們對你的愛.
求主你帶領我們以下時間.
也求你藉著僕人有限的言語.
來去成為引導.
讓我們一同進入經文當中.
看到上帝你心碎的意思.
求你與我們同在.
禱告奉耶穌基督聖名祈禱.
阿們.
十五至十九節.
是整個約翰福音.
二十一章的重點記載.
那是從使徒約翰的視覺.
來回憶耶穌在海邊.
和門徒重聚的一個重要片段.
如果我們細心做一個比較.
其他福音書裡面我們看不到這個記載.
頂多是在《馬太福》裡面提到.
耶穌在復活之後和他們相約.
相約會語在《加利利》這個記載.
那樣.
約翰為什麼要詳細地來記載這個片段呢?.

$^{161}$在約翰的眼中.
耶穌是一個怎樣的道成肉身的上帝.
約翰又怎樣將這番話留在心上呢?.
這個正正就是我去探索.
耶穌對話裡面那種精景的地方的起點.
或者當我們去了解一下.
當中詳細的過程的時候.
我們很容易發現在當中有三組的對話.
而在當中有三對的致意的對比.
我們先來留心看十五至十七節當中.
裡面的內容.
在之前我們也想要回顧一下當中的過程.
二十一章的開始記載了門徒.
從一個日常的生活作息.
他們去到打魚了.
打魚不了準備休息了.
到認出主耶穌的經過.
他們的心情,他們的情緒產生了很大的變化.
由認出主之前的一個很疏落的人生對話.
收工了,沒有了,簡單的對話.
而變為最後只有主耶穌很簡單的吩咐.
你拿東西給我吃.
那種由一個相對熱鬧有聲的地方.
進到一個比較寧靜,安靜的一刻.
一漢就用一個不敢.
在第二十一章裡面說到.
他們認出了耶穌,但是不敢問他是誰.
這個不敢的字眼來反映出.
他們被主耶穌的出現震撼了的心情.
他們不敢.
他們不敢開聲問這位神聖的主.
而在這個時刻.
主耶穌首先打破沉默.
他叫他們拿東西吃.
同時之後就問了彼得這三個直接的問題.
我們可以綜合十五至十七字的對話記載.
來去了解耶穌三問彼得的用意.
在當中,最為我們熟悉的是一個神聖的愛.
Agapiel這個字.
和朋友之愛希臘文Phileo這兩個字的表達.

$^{201}$那是一個兩種層次的表達.
神聖的愛和朋友之間的愛.
除此之外.
我們還有其他的對話裡面不同的字的對比.
「慰養」對比「牧養」,「小養」和「羊」.
這些同義詞的互換.
再加上最後我們留心到.
彼得在回應主耶穌的質問的時候.
他由「是的」,一個很肯定的語句.
到變為「你是無所不知的」.
一個這樣的轉變.
這些修辭手法.
也將我們帶入了兩人對話裡面的精彩之處.
這也是主耶穌.
他為了要重新點起彼得心裡面對他的愛的同時.
也要他正視.
要他正視彼得你自己的問題.
正視彼得的問題.
正視你自己對上帝那份豪情壯語.
背後裡面有一些的細節.
還記不記得當初他在眾人面前就說.
「就是你死,我們也要跟隨你」.
剛剛同樣也是在約翰福音十三章裡面.
在門徒在那一晚.
跟耶穌御節晚餐的時候的豪情壯語.
我們先來看看阿卡普對菲利奧這個交換的作用.
耶穌一開始就問.
「你可不可以用一個無條件的愛來愛我呢?」.
而且是兩次連續追問.
而彼得兩次回答.
都是用「朋友之愛」來回應.
你想想,如果在兩個不同的期望.
層次感來對話的時候.
聽的人有什麼感受?.
「我對你這麼好」.
「原來你是這樣的」.
「我今生付出這麼多」.
「原來你也是這樣的」.
這就是兩個字帶出來的張力.
聖經學者對這個字的運用.

$^{241}$在語文上有不同的見解.
有人認為在這裡沒有影響彼得的心情.
耶穌的追問「你愛不愛我?」.
這個「愛」這個字.
並沒有影響後來彼得的第三次的憂愁.
的確.
或者我們可以這樣看.
彼得不是因為這個「愛」這個字眼而憂愁.
而是主耶穌三次的追問令他著急.
很直線地來看.
這個符合他的性格.
他是一個率直,沒有掩飾的人.
就正如他三次不認主的場面.
問得急了.
就說「我不認識他」.
然而我想重提在旁觀察的那位約翰.
他是怎樣看主耶穌和彼得的對話.
按照道理.
耶穌和彼得的對話是用亞蘭文來講的.
而不是用希臘文來講的.
而這個約翰福音也是約翰在很多年之後.
在爾弗所寫成的.
所以在約翰福音裡面所描述的這個場景.
是使圖約翰相隔多年以後的回憶.
他沉澱過之後.
回想過之後的回憶.
在這一刻.
在這裡使圖約翰告訴我們.
是主耶穌這樣一步一步在用這種手法.
這種方式來重建彼得的人生.
所以我個人會傾向這兩個「愛」的詞語的交換作用.
有它的獨特用意.
具有引發深層思考的作用.
這個「愛」字很有壓迫力.
而「愛」的深度如何量度.
的確一點也不容易.
有時候很難看到表面.
我對於這個「愛」的思考.
我覺得我們中國人一點也不容易去表達.
那種精彩之處.

$^{281}$但是也有它自己心碎含蓄的地方.
我記得很多年前.
我媽媽突然離世.
我爸爸突然的反應.
隨後的反應就是失落.
不知道怎麼辦.
我有時候覺得很奇怪.
為什麼他有這樣的反應.
在我認識我爸爸媽媽這麼多年來.
我從來沒有聽過我爸爸說「我愛你」.
你記住,我們這些老人家.
我已經是老人家,我爸爸媽媽更加是老人家.
我媽媽也不會說「我愛你」.
但是從我爸爸的眼神.
他往後的那種情緒.
他那種失落感,他那種無奈感.
甚至乎他在靈堂的那種感覺.
就呈現出這種思念,珍惜,愛護.
我想這是一種愛的對稱.
愛有時候不止只是在言語.
它是透過生命的見證,行為的表現.
來呈現出對方的珍惜.
說沒用的,看你怎麼做.
這就是愛有意思的地方.
我們再來看另一個字.
Bosikou 放養.
對照 Panmanou 牧羊的字.
這兩個希臘文其實代表著兩個意思.
「餵養」這個字是表示提供食物.
讓牠滿足.
我們從福音書,從牧者的職事來看.
牧羊這個字就是用神的說話供養.
和滋養信徒的屬靈生命.
而這裡強調的是提供靈養一種這樣的行動.
我餵養牠.
作為牧者,作為教會的導師.
我怎樣餵養我的群羊?.
神的說話是最好不過.
甚至乎是最重要,不可或缺的.
而另一個「牧養」這個字.

$^{321}$就是一個很闊的範圍.
是全面一點的.
包括了引導,看管.
將羊群趕回羊圈的一個闊一點的層面.
我們可以看到.
不單止是為牠有提供食物.
甚至乎有管理,保護的.
如果我形象一點,具體一點來說.
我們做兒童工作的導師們.
就有這個這樣的表現了.
要讓他們有聖經教導.
又要和他們吃茶點.
還要帶著他們去廁所.
有沒有類似的情況?.
「餵養,牧養」.
而在這個對話當中.
這兩個字出現的次序.
如果你留心看是「餵養,牧養」和「餵養」.
在這幾個經文裡面.
我們可以看到.
這個次序很有意思.
按照這個次序的結論來看.
主耶穌很想帶回使徒彼得.
要知道回到上帝的話語裡面.
才是最重要的.
回到神的話語裡面.
才是對群養最重要.
特別我們再次去想想.
當彼得還未能夠明白到.
主耶穌以神聖的愛來回應的時候.
牧養的全面性,全方位性的基礎.
正正是襯托出.
不可或缺的就是神的話語.
就是說你有很多事情是做不到的.
你未必能夠為你的群養犧牲你的生命的.
你有限制的.
不過你不要錯過了聖經的話語.
我相信這個正正也是我們.
「窩福堂頂」姐妹所堅持的.
如果我們經常都記住.

$^{361}$蘇sir所說的三個「不住」.
是三個什麼?.
「不住」是什麼?.
不住....
不住禱告.
不住讀聖經.
還要不住....
我都忘記了第三句.
有人提醒我?.
三個「不住」.
要打我又是下次.
如果不行的話不要加薪水.
三個「不住」.
「不住」好像是「相愛」.
「關懷」「不住」「關懷」.
三個「不住」.
提醒我們.
牧羊群羊就是要有三個元素.
但是最基礎的.
最根本的是什麼?.
是聖經的話語.
我們看到.
牧羊裡面.
耶穌如何來引導彼得.
回到一個很重要的原則上.
然後第三組字就是.
「羊」對「羊」的意思.
希臘文的「胚怒」.
對「帕普帕托」這個字.
總之你記住.
「羊」和「小羊」是兩個希臘文.
在這個理解裡面.
十一世紀的一位神學家.
他對經文的註釋有一個這樣的說法.
他為這兩個字做一個很重要的區別.
他說「小羊」這個字代表著信仰中的初學者.
初信的.
而「羊」就代表那些成熟的.
已經知心的信徒.
對應著前面的「餵羊」和「牧羊」.

$^{401}$這兩個職責的字來說.
我們就明白到.
為什麼耶穌先要提醒彼得去餵養小羊.
然後用「牧羊羊」.
最後再用「餵羊羊」這個字.
針對性地去教導他.
就是說你要留心.
留心那些需要缺乏,欠缺神話語的「羊」.
你要一步一步引導他.
然後一步一步地全面去照顧他.
在這裡我們看到是一個遞進式的.
一層一層地來「牧羊」.
好好「牧羊」我們的「羊群」.
在「牧羊」的實務上面.
耶穌清楚地為彼得指示了正確的方向.
提醒他要以神的話語為重.
面向著整存一個成長生命的即是.
同樣.
這也是我們今天本人.
自己心裡面認定的一個使命.
講好聖經.
講好神的話語.
成為弟兄姊妹終身受用的一個內容.
三祖字的「同義」或是「近義」字.
為我們去拆開了對話中的層次感.
而這個背後有什麼用意呢?.
在這裡給我們有一個很多的反省.
這些對話的內容.
都是為了看到耶穌對彼得直接的衝擊.
而這個衝擊就是一個很直接的問題問.
你愛我比這些更深的思考.
換言之.
你要愛得深其實就是要以行動來回應.
在愛的實踐上面.
耶穌就發出了這個題目.
教彼得你怎樣去愛耶穌所愛.
耶穌怎麼說?.
我的小羊.
我的羊.
這些對象.

$^{441}$這些耶穌所關心的.
耶穌所愛的.
是否都能夠成為你心之所愛呢?.
成為你事奉生命見證的目標呢?.
這是耶穌發出的挑戰.
最後我想和大家看的是彼得在三個回應裡面.
一個很有意思的回答.
剛才你們回答到第三次的時候.
都有一點點失誤.
跳過了.
其實正正就是這種感覺.
第三次的時候.
我們看到前兩次由「是的」.
轉變為「你是無所不知」的一個前後的分別.
「是的」這個希臘文的字的意思是肯定的.
是的.
沒錯的.
對的.
很有把握的.
明確表達.
然而第三次的回答.
答多了.
起碼如果你從中文來看.
已經是字數不同了.
內容也不同了.
他沒有那麼肯定地說「是的」.
他反而說「你無所不知」.
「你知道我愛你」.
把整個回答的焦點不是放在自己身上.
放在耶穌身上.
「你知道的,你還要我說嗎?」.
他的憂愁就在這裡去急.
為什麼還要這樣問我呢?.
整個問題當中.
讓我們看到.
彼得和自己的自信轉向了.
放下鎖手了.
彼得在說.
「我怎樣對你,你應該比我更加清楚」.
可能在我們的生活當中.

$^{481}$有時候我們也會這樣回應.
「我怎樣對你,你應該知道的」.
不過如果對方不是像耶穌那樣.
沒有直接回答這個問題.
而是說.
「我怎樣知道你怎樣愛我?」.
「你會怎樣選擇呢?」.
一個不容易的回答.
一個不容易的回應.
彼得似乎在這一刻的表現就是.
將那個主動權交回給耶穌.
我都不知道怎麼說.
你說吧.
由從一個自信一下子.
變成一個放輕了.
不敢再有豪情壯志的表達.
而正正在這個轉變當中.
也是主耶穌很想彼得去認清一件事.
是認清什麼?.
彼得.
你的愛有限.
你的愛裡面有一點點東西.
還需要處理得好.
或者這個「這些」.
我們不容易來定義的內容.
這個「這些」也未必是我們需要刻意細心去考究.
當中代表著什麼意思.
然而「這些」這個字.
也是同樣給彼得在心裡面.
一個很有意思的提醒.
「是呀!我有些東西.
有些東西愛」.
「還多過愛你耶穌」.
可能這個「這些」在你們每個人心裡面.
你的生命裡面都有不同的內容.
但是總的而言.
這些「這些」正正是阻礙了.
和耶穌的親近.
阻礙了事奉耶穌那份的決心.
阻礙了我們愛人的決心.

$^{521}$又或者甚至乎.
那些「這些」.
是阻礙了我們承認自己的不足.
耶穌在這一刻.
就要彼得去坦然承認.
你有些東西是不行的.
不可以靠你自己的.
這就是耶穌想彼得要學的功夫.
你往後的事奉.
你往後的人生.
不是靠你自己一腔熱誠.
你懂得多少東西.
有多有魄力.
你記住.
你要靠我.
主耶穌.
耶穌讓彼得看到.
他自己能夠承認錯誤的勇氣.
這個很重要.
承認錯誤.
對於牧者來說.
的確是一個很重要的內容.
特別對我來說.
是一個很重要的提醒.
做牧師有時候也有很多做得不好的地方.
要坦然接受自己的限制.
坦然面對著弟兄姊妹的提醒.
不容易.
但是需要.
我很期望我們弟兄姊妹.
對我多些的督促.
指教.
因為我也有限.
我也有時候好像彼得那樣自以為是.
我很坦率.
我很直接.
不過.
有沒有細心面對自己的黑暗.
面對自己的不足呢.
這是整段對話當中.

$^{561}$給我的提醒.
耶穌一步一步地要引導彼得走出他生命裡的低谷.
在這一刻.
預備他迎向挑戰.
迎向使命.
可能你和我都有這個時刻.
又或者無論是過去,現在,將來.
你都將會有這個時刻.
在這裡我想提醒大家.
記住.
主耶穌說這句話.
「我願意愛你」.
你願不願意呢?.
耶穌最後在第三次回應彼得的愛的時候.
他將「Felio」.
從「Capel」這個字轉為「Felio」.
他認同了彼得的愛的限制.
彼得不敢衝向這個神聖的愛的懷抱.
但是主耶穌反而願意去擁抱這位有限制的彼得.
用愛來對應他.
最後這兩節很簡單,很直接.
第十八節這樣說.
「我實實在在地告訴你.
你年輕的時候自己束上帶子.
隨意往來.
但年老的時候你要伸出手來.
別人要把你束上.
帶你到不願意去的地方」.
耶穌說這話,是指彼得會怎樣死來榮耀神.
說了這話,耶穌對他說.
「你跟從我吧」.
在這裡耶穌重申了這個呼召的付出.
「你跟我吧,不過你要預備.
你要預備將要被排斥,被否定.
你將要面對著一個身不由己的處境.
甚至乎最終你要在一個人的眼光來看.
不光彩,不榮耀的結局.
不過,不過,上帝看這一切都是榮耀的.
因為你是跟隨耶穌的腳蹤」.
而這個即是正正就是擺在彼得的面前.

$^{601}$「你餵養我的羊」.
很具體的,很清楚的,有目的性的,有目標意義的.
而做好這個即是的力量.
耶穌就說「你跟從我吧」.
這個跟從我很有意思,很清晰,很有力量.
這個跟從背後是馴服,降服.
「放手,放下你自以為是的事.
放下你引以為傲的事」.
靠耶穌.
「有些事你不行的,你會做錯的」.
這個就是跟從我最重要的地方.
或者這段說話.
這個內容我很多時提及的,是我們作牧者的即是反省.
但我更加提醒大家.
每一位願意跟從主耶穌的弟兄姊妹.
你們的生命見證很有影響力.
無論你在職場,在家,在學校.
有你們的位置.
怎樣做好上帝擺在你們的使命.
是要靠耶穌.
多禱告,多讀經,多實踐.
耶穌問:「你愛我,比這些更深嗎?」.
這個問題不單止是問彼得.
也是問我們在座每一位.
我想給大家一點時間沉澱一下.
剛才我跟大家分享的內容.
特別是「這些」.
是甚麼阻礙你們去跟從主耶穌.
好,讓我們一起同心吐槽.
你激勵我們去跟隨你.
激勵我們去放手,去愛你更多.
主,為願你的聖靈繼續提醒我們.
給我們有能力.
讓我們在軟弱當中,剛強起來.
愛你更多,愛你更深.
就像一個小孩子一樣,放手跟隨主.
多謝主,誰替我們祈禱.
奉耶穌基督聖名祈求,阿們.
\newpage



\section{出埃及記 12:42}
\label{sec:031943}
\textbf{聖經人聖經事:這夜是耶和華的夜}~~~~黃克勤牧師博士
\newline
\newline
連結: \href{http://www.church.com.hk/acms/content.asp?site=cdc\&op=show\&type=product\&code=031943\&layout=sermon}{\texttt{031943}} ~~~~ 語音日期: 2026-04-06
\newline
\newline
\hyperref[sec:031942]{\small{< < < PREV SERMON < < <}}
~
\hyperref[sec:index_chronic]{\small{[返順時目]}}
~
\hyperref[sec:031944]{\small{> > > NEXT SERMON > > >}}
\newline
\newline
出埃及記 12:42
\newline
\begin{longtable}{cl}
\hline
\hline
章節 & 經文 (和合本修訂版)\\
\hline
12:42 & \begin{tabularx}{0.7\textwidth}{X} 這是向耶和華守的夜，他領他們出埃及地；這是以色列眾人世世代代要向耶和華守的夜。 \end{tabularx} \\ \\
[1ex]
\hline
\hline
\end{longtable}
$^{1}$第二節我們又一齊讀聖經.
生命樹為大家預備的聖經人士.
今天看背後上帝的工作.
今天講的「借夜視耶和華的夜」.
這個是上帝講的 耶和華講的.
我們的神講的.
這一天是耶和華的夜.
那就是說什麼?.
這個是我生日.
那就是說今天是我生日.
就是我說了算.
「借夜視耶和華的夜」.
「借夜視我們的國慶」.
那你知道那個是屬於誰的.
「借夜視耶和華的夜」.
是指這個晚上特別屬於耶和華.
每一個晚上都屬於耶和華的.
是 不過這個晚上特別一點的.
好 追及的第十二章.
以色列人經過了第十災.
於是就發怒起來.
全部埃及人都哭了.
因為每一家都死一個人.
於是就叫他們走啊走啊.
馬上走 半夜就趕他們走.
哇 他們就沒有鞋子玩劇.
就拿著很多東西走了.
在哪裡走呢?.
以色列人就從蘭塞起行往疏角去.
除了婦人孩子暴行的男人約有六十萬.
這個就是十二章第三十七節.
以色列人從蘭塞.
蘭塞是一個城.
這個城是誰建的呢?.
是以色列人奴隸建的.
他們做什麼呢?.
是用來建得很漂亮很大.
用來擺貨的.
擺什麼貨呢?.
各樣的財物 食物 穀等等.

$^{41}$墨紙這些擺在裡面.
就是由以色列人建的.
以色列人就建了兩個積貨城.
城裡面叫做積貨城.
其中一個就叫做蘭塞.
以色列人從蘭塞起行.
他們就不只是住在蘭塞一個地方.
他們很多是住在戈山.
一個很大很肥沃.
是成為埃及最漂亮的地方.
他們就集合了.
於是他們就一起去到蘭塞.
然後出發了.
蘭塞是Romesus.
Romesus大家會知道.
在埃及有一個王很出名的.
就叫做蘭塞二世Romesus第二.
可能就用他這個名字.
因為他爸爸就是第一.
Romesus.
如果你在瀏覽器那裡看.
蘭塞都會看到他的樣子.
甚至有人在木乃伊那裡找到.
是看到他的骨頭.
因此由他的骨頭.
可以做回他的圓形的樣子.
大概是怎樣.
如果你在谷歌的圖都找到的.
Romesus.
Romesus是太陽之子的意思.
他們很尊敬太陽的.
不過沒用的.
黑暗之災三日熄你燈.
看你太陽之子是什麼.
唯獨以色列地方有亮光.
真厲害.
他們由蘭塞起行往疏閣.
疏閣在聖經裡面有兩個疏閣.
起碼兩個.
一個就在埃及.

$^{81}$另一個就在以色列地的北端.
就是瓦割跟他的哥哥以素相認的時候.
就在那附近.
瓦剝渡口.
然後以素先行自行回以東.
然後瓦割說我們這些年紀輕.
正在乳養嬰兒.
還有那些牛羊都小.
那些女人又多.
不走得這麼快的.
你們先走.
我們跟著來跟著來.
他一看見哥哥走了.
他馬上調頭去了疏閣.
那裡有一個疏閣.
疏閣的意思就是貨攤.
或者叫帳篷.
就好像我們今天那些叫做嘉年華.
一個Booth這樣.
可能當地有很多貨攤.
很多人去那裡.
他們由蘭塞起行去疏閣.
疏閣是C字頭.
其實C是讀S的.
Sukkot.
就是疏閣.
就去那裡.
那裡有多少人起行.
除了婦人孩子.
暴行的男人約有六十萬.
哇.
今天要看的第一點就是.
厲不厲害.
男人有六十萬.
於是你說以色列人出一級有多少人.
男人有六十萬.
一個男人當他有太太.
一兩個子女.
那就是說.
有一個男人背後可能是.

$^{121}$總數是四個或者五個人.
如果是母親.
明白嗎.
爸爸就當成是男人.
所以譬如一個男人.
有他太太.
有兩個子女.
這樣.
所以有四五個人.
你見到如果有六十萬個男人.
那就是說大概浩浩蕩蕩.
那裡可能有近二百萬人.
這個真是大數目了.
二百萬人一起離開.
所以你可以明白.
為什麼法律上是不想讓他們走的.
因為這二百萬人是免費的勞工.
是不需要給人工的.
只是給他們飯吃就可以了.
還要交多少磚頭.
要建多少個城牆.
哇.
二百萬人.
怎麼會讓他們走呢.
又怕他會背叛自己作亂.
但又不捨得這二百萬的勞力.
勞動力.
那怎麼辦.
真是不知道怎麼辦.
所以不讓他們走.
但到最後還是讓他們走.
快點走 快點走 快點走.
二百萬人.
第一 看到上帝的能力多麼豪大.
那你們也知道為什麼摩西會太推手不肯帶.
二百萬人怎麼帶啊.
不要說二百萬.
兩百個人在今天帶他去旅行.
我們教會二百人去旅行.
旅遊巴四輛.

$^{161}$在哪裡集合.
哇 城裡人誰上第一號車誰上四號車.
這個又說還沒上完廁所.
那個又說婆婆還沒到.
哎呀 你真是暈不下來.
二百萬人去到中間吃飯 什麼酒樓.
哇.
頭都赤了.
有些又不知道怎麼辦.
又不知道找哪一輛車.
哎呀 這個中暑 那個見暈.
總之呢.
哎 二百人都搞不定.
還有今天有手機.
二百萬人出來.
沒有手機 沒有WhatsApp.
沒有通訊 怎樣通知啊.
大聲 公都不夠大聲.
所以呢 很艱難的.
如果我是摩西 我都不敢帶.
好了 第一 看到上帝的大能.
就要帶這幫人走.
看到摩西的痛苦.
好了 由太陽之子的城開始起身行了.
好 先看看地圖.
這個地圖我都是在谷歌剪出來的.
以免說我侵權了.
我就貼了一些重要的地方上去.
大家可以看到.
這個地方 埃及.
這個地中海來的 埃及.
這個地方呢.
就是西賴曠野了.
北的叫舒爾曠野.
好了 這個叫西賴曠野.
或者叫西賴半島.
紅海呢 下低上來.
紅海去到北呢.
是分成兩條的.
好了 然後呢.

$^{201}$就過來這個加底斯的曠野.
再進去就是希伯倫耶路撒冷.
其實如果由這裡蘭塞起行.
走到這裡呢.
很近而已.
但是他們卻在曠野那裡.
漂流四十年.
這個是下次有機會再說的.
就是說聖堅是說他們由蘭塞起行.
這個是以色列人出埃及到迦南的旅程.
第一步就是在蘭塞.
Mremesis 起行.
另一個蘭塞是一個.
埃及人逼他們要他們起的積貨城.
第一個.
另一個叫比東.
比東.
比東也是另一個積貨城.
他們兩個都是由以色列人.
做奴隸起的.
這裡是尼羅河.
由南向北流.
流入地中海.
通常一去到平原.
河就會分開變成尼羅河三角洲.
這裡.
這個三角洲是Delta.
這個尼羅河三角洲.
很肥沃.
而且每一年會有水浸.
一有水浸就好了.
水浸泛濫就好了.
浸過所有地.
流走水乾了之後.
很肥沃.
這塊地就叫做戈山.
戈山為什麼以色列人的奴隸.
聚集在戈山最美的地呢?.
原因就是很久之前.
就是約瑟的時代.

$^{241}$以色列人一家人去到.
當時的法老就很器重約瑟.
就說我家人來了.
你隨便選一塊地.
選我們埃及最美的地給他們住.
所以他們是住在戈山的.
但是現在後來有不認識約瑟的新王起來.
就開始逼迫他們了.
就要他們做奴隸了.
所以建的城都是他們住的地方附近.
蘭塞和比東.
由蘭塞起行往疏閣.
就到這裡了.
然後再走比哈希陸.
就在這個狹窄的位置過紅海.
就進入到這裡了.
因為這個派兵追他們.
追到這裡.
比哈希陸他們就過海了.
大概是這樣的.
一路繼續走.
就在輸野抗議一路走的時間.
去到西冷半島附近.
就經過一個很苦的泉水.
又叫做馬拉的苦水.
然後下一個就是以林.
哇! 這些水很甜啊.
我想給大家知道.
蘭塞和疏閣.
第一個站,第二個站,第三個站.
就是比哈希陸過海了.
大家知道了.
然後這條路這樣走了很久.
繞一圈.
然後才進入到迦南地.
好了.
第38節挺有趣的.
又有許多閒雜人.
並有羊群,牛群和他們一同上去.
他們60萬男人.

$^{281}$於是我就推測他們近200萬人.
連同婦女和小孩.
還有無數牛羊.
因為他們就拿了別人的牛羊,衣服.
那些金器等等都拿走了.
我就說「給你,給你,給你」.
這都是神的工作.
「讓他們快點走,快點走」.
「快點走,快點走」.
有許多閒雜人一同.
一起上去.
上去哪裡呢?.
上去英去之地.
哇! 他們聽到.
他們走啊?.
就說「皇上,放他們走啊」.
「法律王,放他們走啊」.
然後看到十個災.
這幫人跟著耶和華神真厲害.
我們也跟著他了.
埃及那些神全部都不行.
我們也跟著他.
跟著他就有上帝保護.
有耶和華保護.
有許多閒雜人就跟著他了.
那麼「閒雜」這個中文翻譯呢.
Arab, Arab.
Arab這個字其實是「混雜」.
「混雜」和「違善」的意思.
甚麼是「混雜」呢?.
其實Arab這個字是「違善」的意思.
「違善」就是說.
那些布不是有打直的線.
有打橫的線的嗎?.
那好像我們叫「橫維直經」.
譬如世界上我們定這個地點.
經緯度是多少?.
「經」就是指直的.
然後「緯度」是橫的.
啊! 加上那一點.

$^{321}$那麼「違善」是打橫的.
「經」線是打直的.
於是「違善」是指.
又有許多違善和牛羊和我們一起.
甚麼意思呢?.
即是本來我們看東西就是這樣看的.
那麼有很多人.
甚麼人? 甚麼人?.
甚麼人? 甚麼人? 甚麼人?.
啊! 怎麼知道有些其他的人呢?.
所以呢.
就是打橫的.
應該是誰? 誰? 誰?.
我們是以色列人.
以色列人 以色列人.
來了一個不是以色列人.
是他的朋友.
哎! 這個是他的朋友.
啊! 一路數下來.
啊! 這個是他的朋友.
啊! 這個是他的隔壁屋.
所以呢.
就是這樣數下來.
然後打橫.
所以是這個數數的概念.
或者計算的概念.
很多和以色列人有關係的人.
隔壁屋啦.
喂! 喂! 今夜我們要離開了.
我們正在吃雨節.
我們就會有些血這樣擦在門眉門腔上.
你會不會都查?.
什麼事? 我們是埃及人.
哎呀! 你查吧.
要不然這個滅長子之災.
你的長子就沒命了.
這樣嗎?.
那寧可信其有.
不可信其無.
我也查吧.

$^{361}$哎! 查了真的沒事.
所以有些埃及人都會相信.
哇! 真的厲害.
我看了九個災.
現在經過第十個殺長子之災.
而這些血就保護了我.
過了餘月.
哎呀! 接著他說.
好啊! 那我們走吧.
你走不走啊?.
我們雖然隔壁屋.
我不是以色列人.
我是埃及人.
不過都在一起.
我都跟著你了.
可能是這樣.
很多懷疑線.
就是這樣送我來.
啊! 這個是朋友.
這個是同事.
啊! 這個是他隔壁屋.
這樣就跟著來了.
有很多牛群羊在一起.
上去 神鷹去的地方.
上去 Allah 是上去.
為什麼叫上去呢?.
為什麼不叫下去呢?.
其實是向南走的.
為什麼不叫上去呢?.
因為好像我們香港人這樣.
很有趣的.
中環.
香港島的中環.
是一個很重要的地方.
所以說你去哪裡啊?.
我上中環這樣的.
你去哪裡啊?.
我下筲箕灣.
我下西環.
上中環.

$^{401}$為什麼不是下中環.
上筲箕灣的呢?.
因為覺得中環比較重要.
其他地方沒有那麼重要.
那麼我們上去神鷹去之地.
哇! 所以是上去的 Allah.
好了.
他們用埃及帶出來的生麵.
烤成無餃餅.
為什麼用生麵烤無餃餅呢?.
生麵就是說.
那些麵粉搓了.
就一團未落發粉.
這樣啊.
就烤成無餃餅.
為什麼呢?.
因為生麵原沒有發起.
因為他們被催逼離開埃及.
不能擔言.
也沒有為自己預備什麼食物.
所以他們一出埃及之後肚子餓.
原來早些經文是說什麼呢?.
他們就說走啊走啊走啊.
半夜就趕他們走了.
個人拿著東西.
他們將他們的麵本來是晚上吃.
可能臨睡前就搓好了.
有芫荽落了去或者之類的.
或者未落.
走了.
就馬上用盤子將那些麵粉扔下去.
那一疊麵扔下去.
就用衣服包著.
然後就揹在身上或者車上就走了.
但是還沒有發起.
所以那塊麵是生的.
還沒有發起.
到了白天早上半夜三更.
累得不得了出門.
到了白天早上肚子餓.

$^{441}$那怎麼辦呢?.
吃點麵包吃點餅吃.
還沒有發起.
無教啊.
不要緊.
肚子餓不如先吃一吃.
無教是硬很多的.
這次不理了.
照樣吃吧.
因為都沒有預備其他食物.
所以他們開始第一天是吃無教餅的.
他們是不是紀念什麼呢?.
不是.
是趕不及敲而已.
趕不及下教而已.
那就捱第一天而已.
好了.
以色列人住在埃及共有430年.
第二.
紅色字.
哎呀 忘記了螢光天.
正滿了430年那一天.
神的日子計算得很準的.
430年.
滿了那一天.
耶和華的軍隊就從埃及出來.
這裡有一個很特別的字來形容以色列人.
大家留意.
叫Jabba.
Jabba是host.
就是軍隊一群軍隊.
或者一群軍隊一群人.
一群host.
神是將耶和華的軍隊.
就是以色列人.
用耶和華的軍隊來形容以色列人.
哇 他們就浩浩蕩蕩.
好像軍隊似的出來了.
其實說他們是什麼像軍隊呢?.
不是.

$^{481}$只是一盤散沙而已.
不過你要知道.
上帝是看以色列是祂的軍隊.
上帝今天看你和我都是祂的軍隊.
軍隊有很多種.
有些是吃失米的.
有些是很精英的.
有些是很精銳的.
有些是很肥的.
有些是跑不動的.
有些是打仗就走了.
一個都不同.
不過神是將我們看為祂的軍隊.
而我們要更加精煉一點.
以致我們可以為神去打仗.
是稱呼我們做軍隊.
神的能力是計算得很準的.
正滿四百三十年的那一天.
他們就從埃及地出來.
第三點,這夜是耶和華的夜.
因為耶和華領頭出於埃及地.
所以當向耶和華.
謹守以色列眾人世世代代.
應該謹守的.
守什麼?守這夜是耶和華的夜.
剛才說了.
這個晚上.
本來都是一個普通晚上.
不過是神帶他們出埃及.
所以這夜是耶和華值得紀念的夜.
這麼說的話.
就是說這一晚是.
上帝自己都看得很重.
那一天是我帶他們出埃及的.
那一晚你們要記住.
記住什麼?記住神的大能.
四百三十年滿.
那一天二百萬人離開埃及.
經過十災特別地.
十災殺蔣子之災.

$^{521}$然後將他們帶出來.
還將埃及的財物拿出來.
牛羊,金器,衣服.
一齊拿出來.
做幾百年這個奴隸.
不用撿回多少.
那要給多少?神說好.
我叫他們給衣服,給金器.
就一起走了.
就好像早一晚那樣.
吃無教餅吃多一天.
那一天就紀念了.
所以你見到是.
這一夜是耶和華的夜.
是應該謹守的.
我希望大家都能夠記得.
神很多時候在我們生命裡面.
都有一天都有哪一個夜.
是祂做了很特別的工作.
你要記得.
這一夜是耶和華的夜.
是祂的首親自工作.
施行神職其事在你的身上.
不要忘記.
世代謹守.
你要記得.
紀念神的作為.
我們一起祈禱.
求主幫助我們紀念你的作為.
四百三十年就那一天.
滿了.
神你又帶領以色列.
二百萬人離開埃及.
走到抗日當中.
主要是你自己幫助了他們.
帶領了他們施展你的大能.
他們沒有錢.
因為他們都是奴隸.
但是那一夜.
埃及人就償還他們的錢.

$^{561}$把衣服,金器,牛裙,羊裙都給了他們.
他們就出去了.
一次拿回.
主要是我們很感恩.
看到你自己奇妙的作為.
幫助我們常常記得.
這一夜是耶和華的夜.
這一日是耶和華的日.
讓我們謹守.
紀念你奇妙的作為.
常常感恩.
祈禱奉靠耶穌基督的命.
下門.
\newpage



\section{詩篇 90:1-17}
\label{sec:031944}
\textbf{原來我們一直錯判了人生}~~~~廖諾文牧師
\newline
\newline
連結: \href{http://www.church.com.hk/acms/content.asp?site=cdc\&op=show\&type=product\&code=031944\&layout=sermon}{\texttt{031944}} ~~~~ 語音日期: 2026-04-05
\newline
\newline
\hyperref[sec:031943]{\small{< < < PREV SERMON < < <}}
~
\hyperref[sec:index_chronic]{\small{[返順時目]}}
~
\hyperref[sec:031945]{\small{> > > NEXT SERMON > > >}}
\newline
\newline
詩篇 90:1-17
\newline
\begin{longtable}{cl}
\hline
\hline
章節 & 經文 (和合本修訂版)\\
\hline
90:1 & \begin{tabularx}{0.7\textwidth}{X} 主啊，你世世代代作我們的居所。 \end{tabularx} \\ \\ \relax
90:2 & \begin{tabularx}{0.7\textwidth}{X} 諸山未曾生出，地與世界你未曾造成，從亙古到永遠，你是神。 \end{tabularx} \\ \\ \relax
90:3 & \begin{tabularx}{0.7\textwidth}{X} 你使人歸於塵土，說：「世人哪，你們要歸回。」 \end{tabularx} \\ \\ \relax
90:4 & \begin{tabularx}{0.7\textwidth}{X} 在你看來，千年如已過的昨日，又如夜間的一更。 \end{tabularx} \\ \\ \relax
90:5 & \begin{tabularx}{0.7\textwidth}{X} 你叫他們如水沖去，他們如睡一覺。早晨，他們如生長的草； \end{tabularx} \\ \\ \relax
90:6 & \begin{tabularx}{0.7\textwidth}{X} 早晨發芽生長，晚上割下枯乾。 \end{tabularx} \\ \\ \relax
90:7 & \begin{tabularx}{0.7\textwidth}{X} 我們因你的怒氣而消滅，因你的憤怒而驚惶。 \end{tabularx} \\ \\ \relax
90:8 & \begin{tabularx}{0.7\textwidth}{X} 你將我們的罪孽擺在你面前，將我們的隱惡擺在你面光之中。 \end{tabularx} \\ \\ \relax
90:9 & \begin{tabularx}{0.7\textwidth}{X} 我們經過的日子，都在你震怒之下，我們度盡的年歲，好像一聲嘆息。 \end{tabularx} \\ \\ \relax
90:10 & \begin{tabularx}{0.7\textwidth}{X} 我們一生的年日是七十歲，若是強壯可到八十歲；但其中所矜誇的不過是勞苦愁煩，轉眼即逝，我們便如飛而去。 \end{tabularx} \\ \\ \relax
90:11 & \begin{tabularx}{0.7\textwidth}{X} 誰曉得你怒氣的權勢？誰因著敬畏你而曉得你的憤怒呢？ \end{tabularx} \\ \\ \relax
90:12 & \begin{tabularx}{0.7\textwidth}{X} 求你指教我們怎樣數算自己的日子，好叫我們得著智慧的心。 \end{tabularx} \\ \\ \relax
90:13 & \begin{tabularx}{0.7\textwidth}{X} 耶和華啊，我們要等到幾時呢？求你轉回，憐憫你的僕人們。 \end{tabularx} \\ \\ \relax
90:14 & \begin{tabularx}{0.7\textwidth}{X} 求你使我們早早飽得你的慈愛，好叫我們一生一世歡呼喜樂。 \end{tabularx} \\ \\ \relax
90:15 & \begin{tabularx}{0.7\textwidth}{X} 求你照著你使我們受苦的日子，和我們遭難的年歲，使我們喜樂。 \end{tabularx} \\ \\ \relax
90:16 & \begin{tabularx}{0.7\textwidth}{X} 願你的作為向你僕人們顯現，願你的榮耀向他們子孫顯明。 \end{tabularx} \\ \\ \relax
90:17 & \begin{tabularx}{0.7\textwidth}{X} 願主－我們神的恩寵歸於我們身上。願你堅立我們手所做的工，我們手所做的工，願你堅立。 \end{tabularx} \\ \\
[1ex]
\hline
\hline
\end{longtable}
$^{1}$今天是復活節.
我想我們基督徒也會相信.
耶穌已經復活了.
但我們也值得去想一個問題.
耶穌復活怎樣影響我們今天的生活.
耶穌復活怎樣影響我們的行事為人.
耶穌復活怎樣影響我們的價值觀.
或者說耶穌復活的大能.
今天是不是我們每一天都能夠經歷到.
耶穌就在我們裡面.
然後成為我們的盼望.
成為我們的力量.
今天的題目叫做.
原來我們一直都錯判了人生.
今天我會跟大家分享.
詩篇第90篇.
有時候我們對生活的一些事情.
都錯判 看錯了.
跑步之後很多人說流很多汗.
我要補充電解質.
於是要喝什麼.
是啊 你也會說補X力.
但其實是錯的 知道嗎.
不好意思 這是直播.
其實應該要喝水的.
因為你喝運動的飲料.
其實你跑步消耗的卡路里.
你喝一瓶.
其實等於你沒有跑過 知道嗎.
因為你喝一瓶的卡路里.
是差不多比你跑一小時高的.
所以其實不需要喝補X力.
你錯判了.
其實喝清水就夠了.
另一個錯判.
是我最近才發現的.
原來減肥不是吃少一點.
有一陣子我要減肥就吃極少.
瘋狂做運動.
然後弄到自己快要虛脫.

$^{41}$我打算一個星期後上榜.
打算會減肥 誰知道.
減不了肥不但更重了.
於是我好生氣.
原來我去AI問一下.
原來減肥不是要吃少.
是要吃得對.
問題是如果你吃得太少.
你的身體有個自動機制.
脂肪會儲存.
然後流失肌肉.
因為你每天都要有差不多.
可能1200-1300.
你睡在那裡不動.
都要消耗1200-1300.
你吃低過1200-1300.
身體就會消耗肌肉.
我的肌肉消耗了.
脂肪又自動儲存.
難怪我一個星期後.
磅重沒有跌.
還要升了.
後來我問了AI.
原來你要減肥 要增肌.
吃多點東西.
不過要吃得對.
多點菜 多點蛋白質.
然後就是吃這些.
就是旁邊這個.
不要吃炸雞.
不要吃煎炸的東西.
原來不是要吃得少.
是要吃得對.
還要做負重運動.
還要去跑步.
原來人生有時我們看錯了.
我們以為是這樣.
但原來不是這樣的.
不過這些是無傷大雅的.
判錯了 無傷大雅.

$^{81}$不過人生有時有些事.
你錯判了.
是死得人的.
有一個這樣的事例.
有一個老友記在護老院.
應該他要吃軟餐的.
知不知道什麼叫軟餐.
就是吃糊仔.
他吸煙不了東西.
但是那天.
應該是可能新入職.
還是不知道是誰的姐姐.
就誤會了.
餵了他吃固體食物.
吃餵了正餐.
餵餵餵的時候.
一個老友記就噁噁噁.
就窒息送去醫院就死了.
然後發現他在喉嚨那裡.
有些食物的殘渣.
是因為這位老友記.
是不可以吃固體食物.
他吸煙不了.
應該要吃糊仔.
不過看錯了.
在美國有一個護士.
他餵一個五歲的小朋友.
兩歲的.
餵一個兩歲的小朋友吃藥.
應該是要餵1.5個mMOL.
他餵了15.
因為看少了少數點.
應該是餵1.5.
他餵了15.
最後的結果.
那個小朋友.
是吃了10倍的藥而身亡.
有時有些東西看錯了.
是無傷大雅.
但是有些東西看錯了.

$^{121}$是死得人.
餅姐妹.
究竟你和我的人生.
如果錯判了.
可能會影響很大.
今天我們看詩篇第90篇.
大家很熟悉的一字經文.
我們一起讀.
好啊 預備 起.
求你賜教我們咸藥.
數算自己的日子.
保佑我們得著智慧的心.
很多人都很喜歡這句禁句.
相得一求你教我們.
怎樣數算自己的日子.
我小時候有日曆要撕的.
一年365日逐日撕日曆.
撕到十月越來越薄.
然後撕到十月月三十一日.
只剩下一塊紙皮.
又要換另一新的一年.
逐日撕日曆也挺好的.
因為我和你的生命正在倒數.
這個經文說相得一求你教我們.
數算撕日曆自己的日子.
正在倒數.
叫我們得著智慧的心.
其實我很喜歡當代日本.
它譯得更加全神.
我們一起讀.
好啊 預備 起.
求你教我們明白人生有限.
使我們早有先會的一天.
原來甚麼叫數算自己的日子.
就是我們要知道你和我是有限期的.
食物罐頭會有 過期的.
有些買來日本的手信.
有賞味的日期的.
過了賞味日期.
你都不吃就不好吃了.

$^{161}$因為吃壞你的.
其實你和我有一天會過期的.
你不和你旁邊的人說.
你知不知道你會過期的.
不過不知道甚麼時候.
有沒有鑿住額頭.
哪年哪月哪日就過期.
但其實你和我的人生.
是在逐天撕裂的.
逐天撕裂.
還有真正有智慧的人.
是知道人生有限.
你就要在人生這麼有限的日子.
做一個有智慧的人.
今天這篇講道會講三個人生的錯判.
透過這個詩篇告訴我們.
第一個很多人都看錯的就是.
經常以為自己還有時間.
我還有時間未到Deadline.
工作未到Deadline.
功課未到死期.
所以我慢慢做.
拖一下先.
玩一下先.
很多人都以為自己還有很多的時間.
然後大安子義.
我一直都有幾張超市現金券的.
通常現金券期限是兩年.
我大安子義有段時間都未到.
但我印象中是二月到期的.
我在剛好二月的時候.
心血來潮.
拿了幾張來看看我的超市現金券.
想在限期之前買東西.
誰知一拿出來原來過了期一天.
我有一個衝動想下去.
跟那個超市的姐姐說.
通融一下吧.
過了一天而已.
我特意很有誠意的.

$^{201}$讓我通融一下.
你猜得到嗎?.
(猜得到).
是啊 我都沒有做這個傻事了.
原先我是打算去求她的.
但是一天而已 是不是很心痛?.
是不是很不值錢?.
一天而已 你知不知道?.
一天而已.
原本那裡有幾百元的.
接著我真的吐血了.
但是過了期就是過了期.
沒有什麼輝子的 知道嗎?.
你和我的生命有一天都會過期的.
不要以為你還有時間.
有時候說總有機會的.
下次吃飯.
最喜歡的在街上約.
下次一起吃飯.
下次一起去玩.
下次去旅行.
下次見面.
遲些再說吧.
很多人都以為不要搞了.
我遲些再探望父母.
遲些再探望祖父母.
遲些再回應呼召.
我最近在分堂.
最多頂子妹回答我就是.
遲些再讀《門女》.
我們有第二屆《門女》了.
6月開始.
有些頂子妹說我暫時不讀.
我說沒關係.
《門女3》兩年後等著你.
但是坦白說.
兩年後都不知道主是不是回來了.
或者你和我兩年後.
還有沒有命讀《門女》.
又沒有人知道.

$^{241}$很多時候我和你都以為.
下次再說吧.
有時間的.
遲些再說吧.
台灣有一個30歲的男士.
被確診有糖尿病.
但是他恃著自己年輕力壯.
又不控制飲食.
他有吃糖尿病藥.
不過不準時吃藥.
10年後40歲.
他末期腎病.
要永久洗腎.
他很後悔過去那10年.
沒有認真面對他這個身體檢查.
如果我們在座的頂子妹.
是畏疾忌醫.
經常都說.
遲些再幫你檢查.
我不想面對.
遲些再身體檢查.
遲些再見我的家庭醫生.
你一直拖拖拖拖.
小心.
真正有智慧的不是我.
已經.
這裡有人說.
時間可控.
我控制到我的時間.
計劃下半年做什麼.
下半年旅行.
我要去進修.
我要買樓.
什麼時候結婚.
很多時候.
人以為事在人為.
我的健康我的家庭.
我的人際.
我全部在掌控之中.
人的錯判很多時候都是.

$^{281}$以為.
你有沒有試過很多事情.
你以為這樣.
但原來是這樣.
有沒有經歷過.
你以為是這樣.
但原來是這樣.
你以為是這樣.
原來是這樣.
我想我們做人這麼多年.
我們經常都被以為害死了.
有些以為是無傷大雅的.
但有些以為.
是會害死你和我.
所以我們一起看.
摩西怎麼說.
你以為還有時間.
我們一起錄.
預備起.
主啊.
你世世代代.
往我們的居所 諸山未臣身上.
被你這世間的基層造成.
從今古老永遠的主神.
摩西一開始這個詩篇.
她還讚上帝的永恆和安穩.
她說主啊.
你是世世代代.
是我們的居所.
諸山都未臣造成.
世界都未有.
你就是神.
你從永遠到永遠.
你是永遠的上帝.
你甚至比那些大的山更加穩固.
然後摩西繼續說.
人的短暫和生命的脆弱.
我們一起讀.
1 2 3.
上帝是永恆.

$^{321}$但是你和我是短暫.
上帝好像一座山那麼安穩.
但是人的生命很脆弱.
很多人以為是這樣.
但是原來是這樣.
摩西就說.
真正的智慧是什麼.
智慧第一樣就是承認人生是脆弱的.
我們是上帝用泥土做.
然後當我們死的時候.
就塵歸塵 土歸土.
作為牧師.
我們久不久就會幫人做安葬禮.
有時候現在一些安葬禮.
殺灰是殺在一些花園那裡.
就真的變成一袋灰.
是骨灰.
就殺下去泥土那裡.
其實就是代表人很卑微.
塵土值多少錢.
塵土很卑微.
也一拿上手越拿緊就越流得快.
塵土是脆弱.
代表你和我的生命很脆弱.
真正的智慧是承認人生的短促.
在中東的草早上是很脆弱的.
但是因為中東的天氣很乾燥.
有時候刮起熱風的時候.
真的傍晚就枯乾了.
草的壽命是一天的.
代表人生短促到你根本沒辦法預計.
承認人生是無常.
好像人生睡一覺.
不過睡一覺發夢呢.
人生如夢.
發夢好像很真實的.
你有沒有試過發夢急尿.
然後有個廁所在.
你真的覺得很真實.
然後你就吐了.

$^{361}$然後你吐完你就醒了.
我想說人生如夢.
發夢是每樣都真實的.
你有沒有試過發夢被怪獸追.
被大猩猩追.
全部你都覺得很害怕.
全部發夢是驚心動魄的.
而且發夢的時候.
你真的很緊張很害怕.
發夢每樣都好像是真的.
但是夢醒的時候.
夢醒時分.
原來只是發夢而已.
原來摩西形容人的一生.
發夢每樣都很真實.
但是當夢醒的時候.
哦.
那就完了一輩子.
真正的智慧也承認.
你和我是不能掌控人生的.
他說人生如洪水沖去.
在原文是指洪水.
你知道洪水可以沖走一輛車.
沖破一間屋.
是幾秒而已.
海嘯有沒有看過那些影片.
海嘯一來.
那輛車幾秒就從這裡衝去不知道哪裡了.
你和我的人生快到一個地步.
根本你不能留住人生.
時間流逝是不能留住.
我們要承認我們不能掌控人生.
最後真正的智慧.
是承認神有絕對的主權.
究竟你和我什麼時候過期.
什麼時候到期.
是上帝說了算.
你留意這套經文說.
是你使人歸於塵土.
是你叫他們如水沖去.

$^{401}$那個的主題是上帝.
所以今天我們學第一個功課.
不要以為還有時間.
真正的智慧是承認人生短促.
不能掌控.
在弘福院其中一個生患者.
他後來在社交平台寫了一千字.
回憶當天所發生的事.
他說當時我老婆打給我.
我才知道火災.
我馬上穿衣服去逃生.
當我打開門那剎那眼前一刻.
濃煙將我吞噬.
他說濃到一個地步.
我開電話的電筒都不能穿透.
我感覺到呼吸很困難.
我馬上回到家裡.
然後我知道我家大廈的大堂.
已經變成火海.
我不能夠走下去.
我逃生最後那條路都沒有了.
然後他說我知道自己將會囚禁.
在這座名為家家的煉獄裡面.
好像等死 等救援.
他說那些烈焰和濃煙.
讓我明白人很渺小.
當有火 有濃煙.
我連呼吸都是奢侈.
我感慨人生無常.
我們從來都不是生命的主人.
然後他在房子裡等消防員救他.
他望著他房子的所有東西.
他說有什麼我可以帶走的呢.
是不是廢寢忘餐.
捱耕底野噴油完成的模型.
滿心歡喜收藏的限量超合金.
各種的名牌奢侈品.
小朋友喜歡的玩具.
我老婆喜歡的包包.
我全部都很想拿走.

$^{441}$不過什麼都帶不走.
原來人生就是這麼無常.
做夢都沒有想過.
一場火災一鋪清袋.
這是我們接觸很多橫幅園的家庭.
當我們派支票在代表報道會的時候.
很多家庭告訴我們.
住了40年了.
沒有想過一鋪清袋.
做夢沒有想過人生這麼無常.
在網上有一篇貼文.
他說每一個人大概有三萬天.
原來一輩子有三萬天.
當你是83歲命.
當然現在香港比日本更厲害.
全世界最長命.
男人是84左右.
女人是88.
我不知道是不是因為香港人要供樓.
所以上來就給我們這麼長命.
我不知道是不是因為我一定要供30年.
所以可能上帝叫我長命也要繼續供.
就是說如果當你是83歲命.
你有數得計.
每個人很公平.
大概三萬天.
那個貼文是這麼說的.
很刺激的.
有沒有人20歲在這裡.
20歲的你剩下二萬多天.
30歲的你剩下一萬九千天.
是撕日的.
現在跟你一起撕吧.
40歲你過了一半.
就是一萬五千多天.
我就入了五了.
我就剩下一萬一千八百天.
60歲已經變成四位數字了.
八千天.
70歲剩下四千五百天.

$^{481}$80歲剩下九百天.
當然現在很有機會會是九十幾歲.
甚至一百歲.
我想說什麼呢.
就是我們的人生一直在撕日.
在撕啊在撕啊.
其實最重要就是我們今天追的東西.
是不是我們真的想要的呢.
有一個安寧病房的護士.
他常常照顧在臨終的病人.
他訪問了超過一百個臨終的病人.
就發現原來很多臨死的人.
去到人生最後階段.
他就回望一輩子追這個追那個.
但是臨死之前他發現追錯了.
原來我追錢追很多東西不是我想要.
可不可以換回我想要的東西.
原來沒有了.
然後這個護士寫了一本書.
叫做《The Top Five Regrets of the Dying》.
他搜羅了五大後悔.
就是很多臨終的病人覺得自己錯判了人生.
第一.
他很後悔有些話沒有親口說.
原來很多人臨死之前不是想重活一次.
是很想回到過去某個時間.
跟某個人說對不起.
跟某個人說我原諒你.
又或者跟某個人說其實我喜歡你.
不過一直都沒有說.
第二很後悔的就是我沒有好好陪我家人.
很後悔為什麼我不回家吃飯.
為什麼我在家吃飯只看著手機.
然後人在心不在.
很後悔沒有做我想做的事.
我活在父母的期望.
活在社會的期望.
但我心裡面有些事很想做.
不過我不敢做.
我怕失敗.

$^{521}$我又怕我做的事賺不到錢.
所以我正在做的事不是我喜歡的事.
他去到臨終很後悔沒有做自己.
第四.
我很後悔太多時間上班.
不停的加班.
為了工作而工作.
我錯過了我子女小時候.
怎樣陪他學走學說話.
照顧他上學照顧他放學.
現在子女都長大了.
不理我們了.
他們有自己的朋友.
很後悔沒有好好跟自己的配偶聊天.
很後悔我沒有好好建立不同的友誼和關係.
他發現原來事業有成.
彌補不了親情的缺陷.
最後很後悔沒有維繫友誼.
沒有建立真正的友誼.
他說來到最後在病床上.
原來那些地位和成功帶不了人有安慰.
很多臨終病人.
最後最安慰到他的就是誰來探病.
是我的家人.
我的親友.
我花很多年建立的好朋友.
我的閨蜜.
他還會來探病.
人臨終的時候不會再看DSE自己考多少分.
不會再看我銀行存摺了多少錢.
不會再看我房子供滿了沒有.
更重要是誰在床邊陪我走人生最後的階段.
很多臨終的病人後悔錯判了人生.
人一世有多少個主人呢.
原來4000個而已.
這也是社交平台給我的.
我看完馬上彈起來.
不是的 你自己有數可以算的.
一年只有52個星期.
給你80歲命.

$^{561}$其實只有4000個星期可以回主人.
當我這樣想的時候.
我就覺得我要很珍惜每一個星期回來崇拜和弟兄姊妹.
不要那麼輕易說我這星期不回來了.
不要緊 我家裡看直播崇拜.
下面可以穿麻煙筒.
又可以喝咖啡.
沒什麼所謂的.
其實你跟我如果從零歲開始.
每個星期回來崇拜.
你只有4000個主人回來而已.
所以當我看到這個貼文.
我都不知道多麼感恩.
我還有機會講道.
我還有機會來跟大家一起崇拜.
當然我現在暫時一年是一次左右.
或者兩次回來主堂崇拜.
但我很珍惜每一次跟大家一起.
分享神的話也好.
我們敬拜也好.
其實我們能夠回崇拜的日子.
都是倒數中的.
有沒有人知道今天會不會是最後一次.
可能是我最後一次講道.
如果有人說每一次.
你都以最後一次的心情來抹聖餐.
來唱聖詩 來聽道.
我相信都會很不同.
原來一生人只有4000次而已.
那你就不要說又是我當主播.
又是我大敬拜.
又是我講道.
又是我彈琴.
不是的.
可以彈琴 又可以講道是恩典.
逐次減的.
所以今天第一點.
有些什麼你覺得不緊急但重要.
不過你一直在拖延.
可能是跟某個人說對不起.

$^{601}$可能是向某個人說謝謝.
向家人傳福音.
可能是身體健康 天天靈修.
又或者讀門雷.
時間不等人.
你以為可以有下一屆門雷讀.
其實沒有人知道.
所以鼓勵你願意把握今天.
你踏出第一步.
當然這個詩篇也講到第二個人生的錯判.
我們一起讀這節經文 好嗎.
1 2 3.
誰曉得你怒氣的詮釋.
誰妄想你愛神的敬畏.
曉得你的憤怒呢.
摩西正在反問.
摩西帶領120萬以色列人出埃及.
他見到那些以色列人.
然後他就很慨嘆.
究竟有誰懂得驚神.
有誰曉得神的怒氣.
按照上帝該受的敬畏去敬畏他.
其實摩西這樣問.
即是有沒有以色列人懂得驚神.
沒有.
不然他不會這樣求的.
第二個錯判.
很多人以為自己的罪神看不到.
我可以躲起來.
沒事的 神罰我的.
以為可以隱藏.
以為瞞天過海.
你看不到我的.
或者這裡這麼多人.
神沒有空逐個看我的.
祂不知道我在想什麼.
祂不知道我秘密在做什麼.
甚至神知道又怎樣.
神不追究.
追究了耶穌.

$^{641}$剛才吃了醒餐.
神不罰我的.
罰了耶穌.
這些全是以為.
你以為是這樣.
但原來未必是這樣.
以為神不會發怒.
救弱的神才會惡一點.
新約的神很慈愛.
耶穌為我們死.
神不會生氣的.
OK的.
以為自己沒有罪.
我又不是作奸犯科.
我都不是很壞而已.
我沒有他旁邊那麼壞.
我沒事的.
很多人以為沒事.
但我們留意聖經怎麼說.
我們一起和我一起說.
一二三.
我們因你的怒氣而消滅.
因你的憤怒而經過.
你將我們的罪名擺在你面前.
將我們的隱惡擺在你面光之中.
這裡摩西說原來我們會有罪孽.
是指外在的行為犯罪.
隱惡是隱然未現的罪.
可能指到態度.
可能指到心態.
可能指到心思意念.
沒錯我們有罪孽有隱惡.
不過神會將它擺在哪裡.
將我和你的罪擺在哪裡.
垃圾桶.
擺在了耶穌身上.
還有擺在哪裡.
你留意.
很恐怖.
原來上帝將你和我的罪孽擺在他面前.

$^{681}$然後將你和我的隱惡擺在上帝面光之中.
原來上帝不會看不到.
還要放大來看.
將你和我的罪擺在他面前.
然後神會有什麼反應呢.
不要緊.
耶穌頂罪了.
OK的.
沒問題的.
這個經文多次出現「嬲」.
神發嬲.
神憤怒.
神震怒.
神的怒氣.
其實神的怒氣是表示祂的愛.
是表示祂忌邪的愛.
神發嬲是警告罪人要悔改,迴轉.
我們做父母的.
會不會對鄰居的小朋友發怒呢.
他頑皮不關我的事.
但是我兩個女兒不乖我就會生氣.
我生氣就是因為我愛他.
有時候真的想打下來.
打啊!快點啊!.
當然不會出手打的.
但是是自己的子女才會氣自己.
也有人這樣說.
當你的父母不再罵你.
不再生你的氣.
不再向你發怒.
就是放棄你了.
是不是.
就是不想再理你了.
隨便吧.
自生自滅.
我沒有感覺了.
整個舊約上帝常常對以色列人發怒.
其實神的憤怒是一個銀的另一面.
就是神的慈愛.
我們常常會錯判神的慈愛等於不發怒.

$^{721}$你看錯了.
如果今天你的父母還有父母向你發怒.
向你囉嗦.
這個是愛你啊.
向你辱謹.
或者說你有你的老婆還向你發怒.
是愛你啊.
如果你的老婆說.
我放棄你.
我不想再罵你.
我不想再訛你.
這個是她對你的放棄.
所以上帝在舊約.
他真的生氣.
以色列人拜金牛毒.
摩西下來就弄爛兩塊法板.
然後可拉黨.
因為他挑戰摩西的權柄.
然後地都裂開.
吞了可拉黨.
你說神是不是很生氣.
生氣到地都裂開.
然後有以色列人拜巴力.
然後上帝用瘟疫擊殺了以色列人.
然後以色列人拜偶像埋怨上帝.
上帝叫火蛇咬死很多以色列人.
然後他預備了銅蛇的救恩.
在摩西帶領以色列人出埃及上.
上都見到神發怒.
但是摩西看到.
神發怒是代表他的愛.
神透過發怒是希望迫使以色列人知道.
神的怒氣會叫我們消滅.
叫我們要息驚.
否則我們不敬畏神的人生就是一聲嘆息.
然後轉眼成空.
便如飛而去.
所以今天我們要學第二個功課.
我們以為我們的罪可以隱藏.
可以躲起來.

$^{761}$但是摩西說真正的智慧.
是活在神的光中清除隱惡.
大家有沒有留意到.
最近這兩年多了很多天眼看著我們.
警隊有一個銳眼的計劃.
這個銳眼計劃是幫助破案的.
直至到今年3月.
銳眼計劃幫助破案700多宗.
然後捉了1200人.
還要在24小時內通緝到嫌疑犯.
九成的罪案.
例如兇殺,行劫,爆竊.
兩天內就能夠破案.
現在是不是很厲害.
因為布滿天羅地眼.
然後還要說.
2028年會增加6萬支CCTV.
還要聯合香港鐵路.
還要加上公屋.
加上政府的建築物.
都會有很多眼睛看著你.
這麼多眼睛看著你.
你真的無法躲避.
但是你知不知道有一個眼鏡更重要.
我們一起跟它玩.
預備 1,2,3.
「耶和華的眼睦」.
(耶和華的眼睦偏差全地).
(要天大幫助監察宋傳誠實的人).
(這就是反正圓滿).
(此後也必有新戰的事).
「耶和華的眼睦」.
(偏差港九新界).
(偏差整個的世界).
(周圍).
(如果你說周圍都有天眼,銳眼).
(你不敢犯法).
(但我跟你更加知道).
(耶和華的眼睦偏差全地).
它看到什麼?.

$^{801}$它看到哪個人的心存誠實.
也看到哪個心是鬼炸.
哪個的心偏於邪.
它看到人內心的動機和思想.
它也看到行事愚昧.
就是外在的行為.
神有看到.
原來我們走不出神的法眼.
所以網上我們也看到.
很值得我們思考的就是.
基督徒常見隱然未然的罪.
有五樣.
第一樣是屬靈的驕傲.
常常用聖經去答人,糾正人.
只看到人的不好.
看不到自己的問題.
看到人眼中有刺.
看不到自己有良慕.
這就是屬靈的驕傲.
第二就是心中的偶像.
當你喜歡的東西沒有了.
破了.
然後你就崩潰.
甚至向神發怒.
懷疑神的愛.
那些東西就成為了你的偶像.
第三.
虛假的平安.
當你參加宗教活動.
然後你自我感覺良好.
幸好我回來採摘了八達通.
OK啦.
這三個月我不怕了.
究竟回來崇拜是為了採摘八達通.
有一個服務紀錄.
還是我來潮見耶和華上帝.
我要在樹上把榮耀頌讚歸給祂.
有時候有些基督徒.
他是想著來了參加宗教活動.
然後自我感覺良好.

$^{841}$然後繼續對罪沒有什麼敏感.
然後遊走在罪的邊緣.
灰色地帶.
然後相信神會赦免我.
這個叫做虛假的平安.
第四.
未能覺察自己有苦毒.
什麼叫做未能覺察自己有苦毒呢.
當你有不喜歡的人.
你聽到他不好的消息.
例如他有病.
例如他撞車.
例如他被人刷牛乳.
然後你心中會暗爽的.
這個就叫做未覺察的苦毒.
那你留意一下.
你如果有些眼中釘.
你知道他有不好的消息.
然後你爽呀.
正呀.
抵呀.
YES呀.
那你小心了.
又或者你對著你不太喜歡的人.
你言談間有骨.
冷言冷語.
那你也要小心.
最後就是虛榮.
以及討人歡喜.
當你去事奉沒有人讚你.
沒有人欣賞你.
那就算了不做了.
還要被人批評.
不要搞了.
不戴勁扮了.
好不要來港島了.
或者港島的時候.
說一些人喜歡聽的.
最好就不要說雜地的說話.
如果是這樣.

$^{881}$這個是隱然未然的罪.
我答應了一個弟兄.
今天要說一下以下這個見證.
因為之前那三堂課就沒有了.
那我有一個隱幕.
其實我最初來雅根堂.
港島啦.
那我也很想.
把我的港島重溫那幾點.
放幾張Powerpoint在我的Facebook.
其實我的初心都真的是.
想一些弟兄姊妹可以重溫港島的內容.
還有我也很有心機做Powerpoint.
我初心真的是放進去的.
然後慢慢試過三五六次.
當然有很多人讚.
有很多人不同的批評.
後來發現.
我常常按進去看.
有多少人Like呢.
然後有哪個有讚的呢.
有哪個又怎樣.
我發現我追緊Like.
我追緊有人認同.
我最初的初心是想總之給大家重溫.
然後告訴大家什麼時候又會上架.
可以再聽了.
最初初心是沒問題.
但是慢慢開始.
神提我.
你的動機.
你做這件事不是錯.
但是你追著人.
經常要按進去看.
有人Like你很開心.
但是似乎一兩天沒人Like過你會失落.
我就知道我動機出了問題.
我是想藉著這件事.
好像是想褒獎自己.
或者想去騙人Like.

$^{921}$然後呢.
我都做掙扎過好幾次.
我跟神說不是啦.
總之我都是想弟兄姊妹.
祝福弟兄姊妹可以給他們有Like.
還有有時候弟兄姊妹會PM我.
問我拿回泡泡.
那不如我放上Facebook啦.
但是後來我都受不了.
就是聖靈一直提醒我不是啦.
總之你常常都追著.
有Like你很Like.
沒Like你會失落.
我追著了這件事.
後來我就決定我不會再放上Facebook.
就是我講到的內容.
不會再.
因為我知道我追緊什麼.
這個是隱然未然的隱惡.
我也是第一次講的公開.
但是當我預備這篇講到.
我覺得.
我不想我自己變成一個試探.
於是我就這幾年無聊.
我分堂無學.
回來港島.
我都不會將自己的港島.
去放上去Facebook.
因為我知道這個是我的隱惡.
所以我要對付.
這個其中都是另一個隱藏的問題.
就是有一架飛機.
北歐航空今年2月.
由瑞典飛去西班牙.
飛了一個小時之後.
竟然在機艙裡面發現一隻老鼠.
然後一定要捉牠.
但是更加慘的是.
你看到老鼠.
你不知道牠咬了哪些電線.

$^{961}$於是飛機要繞彎回瑞典.
盤旋三個小時.
然後下機.
所有工作人員要全面搜查.
老鼠有沒有咬斷電線.
因為如果是這樣.
其實一隻老鼠可以令飛機墜毀.
如果線路錯了.
或者被咬了.
一隻老鼠而已.
令到整架飛機要回頭.
今天在你生命裡面.
有沒有要捉老鼠.
有沒有一些隱然未現的惡.
有時候不知道.
有時候是當我靜下來.
我才發現神跟我說.
不是啊 你最近的讚好數.
沒有人讚好你.
你似乎失落.
你不只是想弟兄姊妹重溫那麼簡單.
這個是聖靈說的.
你有你要捉的老鼠.
當你和我要捉了那些隱惡.
清除了我們要活在神的光中.
不要自己騙自己.
看不到而已.
沒事的.
OK的.
第三.
這個詩篇到最後.
其實是很開心的結束.
中間就說.
如果不敬畏神的人生是八隻字.
怒虎受犯.
轉眼成空.
但是這個詩篇的結尾.
是可以叫我們一生一世什麼.
歡呼.
喜樂就是YES.

$^{1001}$那怎麼能夠由怒虎受犯.
變成歡呼喜樂呢.
不過很多人就追錯了快樂.
他以為怎麼才能夠是歡呼喜樂呢.
就是及時行樂.
有錢.
有包包 有名牌.
就代表我的身份價值.
愛情至上.
我有另一半就最滿足到我.
但是最慘的就是.
可能我跟你一輩子在追錢.
追緊名牌 追緊包包.
原來發現當你得到的時候.
其實都不外如是.
完全滿足不了自己的心.
原來真正的快樂.
是怎麼能夠拿到的呢.
香港的幸福排名指數.
現在跌到多少位.
九十.
九十啊.
還低過中國.
中國也有六十五.
我們就九十.
所以難怪我們經常都北上.
好像是.
北上.
可能又不關事.
可能北上不對.
我遲些也會去.
所以再開心一下.
留在香港我們九十.
北上又到十五.
看起來好像很有反應.
我們一起讀最後這幾節.
預備.
一 二 三.
求你指教我們昨日所算是你的一日子.
好叫我們一將志為心.

$^{1041}$以後啊.
我們要等到幾時呢.
求你轉門.
為你僕人後悔.
求你使我們做做把得你的善惡.
好叫我們一生一世歡呼喜樂.
求你照著你使我們受苦的日子.
和我們遭難的年歲.
叫我們喜樂.
願你的作為將你僕人展現.
願你的榮耀向下門子孫甜蜜.
願你使我們神的榮耀歸於我們身上.
願你天下我們手所做的功.
這個詩篇的結尾是充滿盼望.
而黃色的字是我們最喜歡的.
我們希望能夠有智慧.
有愛.
有歡呼.
有喜樂.
有榮耀.
有榮美.
不過追錯了地方.
我們以為這個世界比我們有歡呼.
能夠有喜樂.
能夠這個世界比我們有愛.
比我們有榮耀.
追錯了.
這個詩篇的結尾告訴我們.
真正的智慧.
慈愛.
歡呼.
喜樂.
榮耀.
全部在神裡面.
不是在這個世界.
所以這個詩篇最後用六個祈禱結束.
是.
摩世用這六個禱告告訴我們.
如何能夠真正的一生.
YES.

$^{1081}$歡呼喜樂呢.
第一.
祈求神指教我們活出智慧的人生.
第二.
耶和華.
整個詩篇第一次出現耶和華這個字.
耶和華就是上帝與以色列人立約的名字.
沒錯.
以色列人犯罪.
但是摩世重提耶和華的名字.
神是守約施慈愛.
祈求神轉回憐憫罪人.
然後就求早早保得神的慈愛.
慈愛這個字也是上帝立約的一個.
最重要的忠心.
就是神堅定地愛以色列人.
希望每早晨都飽享神的慈愛.
你是不是每早起床靈修的.
每早就向神大大的張口.
然後讓祂的愛充滿我們在全日.
都能夠歡呼喜樂.
唯有神的愛會叫我們YES.
我不知道你有沒有這個反應.
還是我說你明天去日本.
你說YES.
或者升職加薪YES.
領受神的慈愛YES.
原來摩世祈求上帝讓我們早早.
意思是每個清晨.
又或者說在人生最年輕的時候.
小朋友的時候.
青年人的時候.
來信耶穌保得神的慈愛.
第四.
我們都會面對苦難.
在受苦的日子遭難的年歲.
能夠以神為喜樂.
第五.
我很希望看到的不是我的榮耀.
我升職加薪.

$^{1121}$是希望看到神在我家庭那裡表演.
在我工作上表演.
在我的生命上表演.
我要見證的是神的大能和神的榮耀.
我追的不是追自己出名.
是要追神的榮耀.
你追的是要看到神的大能在你的生命上.
然後最後.
就是要祈求今世的工作有永恆的價值.
我所做的求你兼顧.
我所作的功求你兼顧.
當人真的為上帝而活.
為上帝做工作.
他就有永恆的價值.
這個詩篇是誰寫的.
摩世寫的.
但是當我昨晚回去再默想的時候.
摩世自己進不了迦南地.
他都不能帶領以色列人進迦南地.
但是他成功交棒給誰.
約書亞.
楊約書亞成功帶領以色列人進迦南地.
原來摩世完成神給他的照命.
他是放下自己的得失.
他沒有看自己進迦南地是成功的.
他只是看他帶領了以色列人四十幾年在曠野.
縱然自己進不了迦南地.
但是他成功交棒.
成功知道約書亞帶領以色列人進迦南地.
原來摩世所做的有永恆的價值.
縱然他自己進不了迦南地.
他放下自己的得失.
最重要是神的旨意成就.
這個叫做有永恆的價值.
所以最後我們要學的就是.
你以為追到快樂.
其實最重要的人生就是飽享神的慈愛.
進入永恆.
在分堂有一個姐妹.
她2020年辭去了做了28年的金融機構.

$^{1161}$她就打算開一個YouTube Channel.
來宣揚詩歌伸展慕道.
大半年之後.
她患上卵巢癌第三期.
她問神一個問題.
為什麼?你有什麼心意?.
神就跟她說.
我不是叫你不停的拍片.
我是要你用你的癌病.
有同理心的去明白其他癌症病人的需要.
這個姐妹就加入了基督教癌關做義工.
過去那五年做了兩次大手術.
不間斷化療,標靶治療,免疫治療.
中藥治療,十次電療.
她在2024年癌病復發.
醫生說只剩下一年半命.
但這個姐妹跟我說.
我早兩個星期跟師母一起吃飯.
她說你知不知道我做這麼多治療最有效是什麼?.
喜樂的心乃是良藥.
她說最能幫到我的就是音樂治療.
縱然我的癌令我很痛.
但我會向我的癌細胞嘗試過.
我一邊跳舞讚美神.
我就覺得喜樂.
我就開心.
我不怕回天家.
不過我最擔心的是我未信主的媽媽.
於是她不斷帶領媽媽回教會.
鼓勵她回喜樂祖.
在過去那兩年.
她很努力帶媽媽信耶穌.
然後在我沒記錯是去年.
她媽媽就在這個舞台接受灑水禮.
當然這個姐妹.
她今天一有生命.
一有健康她就會去講建聖.
我從來不會覺得她是仇的.
她比我更喜樂.
比很多弟兄姊妹都更喜樂.

$^{1201}$所以由她的口講出來.
喜樂的心乃是良藥.
她最重要是早早飽享神的慈愛.
弟兄姊妹人生有三萬日很短.
但這三萬日是為了我們進入永恆.
這三萬日你和我怎樣用.
是讓我們預備好進入永恆.
我們每個弟兄姊妹只有四千個主日可以返.
但這四千個星期日其實都是準備好.
我們預演將來永恆在天家.
在那裡來敬拜主.
所以這四千個主日是我們預演的.
是一個儀式來的.
所以我們好像很短.
但四千個主日.
每個人只有三萬日.
但這三萬日是為我們預備永恆的.
所以你看看你旁邊那個.
我們不止三萬日.
我們有永恆 是嗎.
但你怎樣用這三萬日.
預備你進入永恆.
是非常重要.
今日講三個的錯判.
很多人以為自己還有很多時間.
但不一定.
很多人以為自己的惡是可以隱藏.
一定不可能.
很多人以為自己追世界很多事.
但追完發覺不外如是.
真正的智慧在於承認人生短促.
不能掌控.
真正的智慧在於活在神的光中.
清除隱惡.
真正的智慧在於飽享神的慈愛.
進入永恆.
在這個星期.
我聽到一個童工分享令我覺得很棒.
他說在IG POST有人分享.
基督徒一輩子就是要進入神國.

$^{1241}$進入天國.
不過好像很遙遠.
很抽象.
一點都不興奮.
然後他就用了一個比喻.
他說我們通常買機票.
當你在網上買機票.
信用卡付了錢.
你就開始興奮.
其實你是買明年才去旅行.
或者下半年才去旅行.
但你一按了信用卡付錢.
Yes 訂了機票.
又訂了酒店.
你就開始未出發.
才興奮.
接著你會看YOUTUBE.
YOUTUBER在哪裡介紹.
哪裡有什麼好吃的.
然後做一些計劃.
然後去哪裡玩.
然後我要去哪裡買東西.
你未出發就才興奮.
我覺得這個比喻很振奮.
弟兄姊妹我跟你都未進入天國.
但今天我們在預備的時候.
你有沒有這個Yes.
已經有入場券.
你已經有永生.
你知道你會回到天家.
但當你在回到天家的時候.
你有沒有未出發才興奮.
似乎大家好像沒有興奮.
好好珍惜四千個主日.
很開心今天跟大家一起敬拜.
我們一起.
讓我們透過一種思覺.
我們將榮耀雞給主.
也將我們的心歸給主.
我相信我們唱一唱.

$^{1281}$四千個主日很快.
不知道誰先呢.
我們就會回去進入永恆.
但我們現在好好地練習自己.
好好用這三萬日.
預備我們見上帝的一刻.
\newpage



\section{詩篇 16:1-2}
\label{sec:031945}
\textbf{投靠神的人有福了}~~~~呂君望牧師
\newline
\newline
連結: \href{http://www.church.com.hk/acms/content.asp?site=cdc\&op=show\&type=product\&code=031945\&layout=sermon}{\texttt{031945}} ~~~~ 語音日期: 2026-04-05
\newline
\newline
\hyperref[sec:031944]{\small{< < < PREV SERMON < < <}}
~
\hyperref[sec:index_chronic]{\small{[返順時目]}}
~
\hyperref[sec:031946]{\small{> > > NEXT SERMON > > >}}
\newline
\newline
詩篇 16:1-2
\newline
\begin{longtable}{cl}
\hline
\hline
章節 & 經文 (和合本修訂版)\\
\hline
16:1 & \begin{tabularx}{0.7\textwidth}{X} 神啊，求你保佑我，因為我投靠你。 \end{tabularx} \\ \\ \relax
16:2 & \begin{tabularx}{0.7\textwidth}{X} 我曾對耶和華說：「你是我的主，我的福氣惟獨從你而來。」 \end{tabularx} \\ \\
[1ex]
\hline
\hline
\end{longtable}
$^{1}$弟兄姊妹平安.
今天很特別是清明節.
而也是復活節.
兩個節日都是跟一個已死的人有關.
不過很不同.
清明節是讓我們去緬懷過去.
腎中追遠.
無論是餓在土地裡面.
或者骨枯萎裡面的仙人.
他們的鞋骨仍然存在.
他們的骨灰也都存在.
但是復活節的主留下的是一個空的墳墓.
所以兩個節日雖然同一天好像這麼巧.
但帶來的是截然不同的意義.
清明節是讓我們去紀念我們的仙人.
復活節是紀念我們為我們死而復活的主.
同樣是紀念.
不過在聖經裡面的紀念卻有不同.
聖經的紀念不只是懷念過去.
頭腦上的知道.
這個就是你的仙人.
他就是你的誰誰.
他曾經怎樣怎樣對我們.
以致我們要帶給我們後一代.
不要忘記我們的祖先.
但是聖經的紀念卻提醒我們.
因著我們復活的主.
我們不但要感恩.
因著我們復活的主.
我們知道怎樣去面對現在所處身的環境.
所經歷的遭遇.
更加因為我們的主是復活.
我有信心去面對將來.
因為他說過他會榮耀在內.
就算我今天有多困難.
有多少事情解決不了.
有很多事情不明白也好.
他會告訴我.
他是與我們立約的三一的真神.
我們主是三一道成肉身的主.

$^{41}$曾經死在十字架上.
但是已經復活了.
他會再來.
他不再是一個單單已死的偉人.
而是活著與我們同在的真神.
當聖經說到神紀念什麼什麼的時候.
其實不只是說神記得什麼.
聖經神紀念其實是表示.
這位的神他會施行拯救.
成就他所說過的應許.
正如出合記二章二十四節.
神聽見他們的哀聲.
就紀念他與阿巴拉罕.
以撒亞國所立的約.
神紀念他與以色列先祖所立的約.
所以他會拯救以色列人出埃及.
同樣聖經常常說.
你們要紀念窮人.
要紀念受迫迫的人.
意思不只是說窮人很可憐.
受迫迫的人真是困難.
我們不只是在這裡可憐.
不只是很記掛.
而是我們對於窮人.
以及面對著迫迫的人.
我們有什麼行動去作回應呢.
所以很多人很喜歡在復活節受洗.
因為復活節的意義.
不只是紀念已死的耶穌.
我們透過這個復活節.
接受聖禮.
是與基督同死.
同埋葬.
同復活.
所以弟兄姊妹.
今天是復活節的主日.
你用什麼態度.
你帶著什麼心境.
去紀念我們復活的主呢.
是否只是如常.

$^{81}$因為今天是星期日.
所以就回來.
還是你帶著一個感恩的心.
剛才我想問大家.
既然我們的主是復活了.
那你現在.
你覺得你自己是否有一個復活的心境.
去敬拜上帝.
還是好像清明節那樣.
很苦啊.
死了死了.
還是怎樣.
用什麼心境呢.
我希望我們每個主日.
其實我們基督教.
以主日.
是要表明我們的信仰.
是一個復活的信仰.
不是一個已死的信仰.
所以我們是要帶給.
未知未信的人.
來到教會的時候.
去找到盼望.
而不是找到絕望.
今天的講題叫做.
投靠神的人有福了.
是肯定的宣告.
講到投靠.
究竟世人一般有什麼依靠呢.
先不說世人.
說我吧.
我記得我在這裡事奉做導師的時候.
因為太多事情做.
那時候人手不夠.
很多時候買東西的東西.
我都希望找一些弟兄幫忙.
不是找他們賺錢.
找他們笨.
找他們幫我做事.
而是大家分擔一下.

$^{121}$是不是啊.
也都等他們能夠有些承擔.
團結你有份的.
這些東西都跟你有關係的.
我記得弟兄你幫我買一個.
錄到音就行了.
便宜一點.
D字系列的錄音帶就可以了.
D字大家明白的.
最便宜的那些.
但我收到回來的.
就是S,X那些.
就是那些錄歌那些.
啊.
我沒有罵他.
我說多謝弟兄.
不過不是我想你買.
不要緊.
你買了照用.
原來印是有限制的.
我不知道他聽不明白.
還是怎樣.
因為印.
不是可以成為我們.
最終的可靠.
不是他一定有心.
是可能人的限制.
可能領悟.
領悟.
或者其他的限制.
或者不明白我問什麼.
要求什麼.
有時候我都會.
跟那些比較.
特別是在舞會.
應該是舞會.
有一些弟兄.
大家都同一個團體.
很熟的.
有一次我跟他說了一個秘密.

$^{161}$那不是秘密.
誰知道那個秘密.
在團體裡面.
就講為人知.
幸好可能那個秘密.
在他的心目中.
什麼秘密來的.
開心的東西.
秘密來的.
那就傳開了.
所以頂子妹.
當我們說.
我們依靠一些人的時候.
原來人.
他們有自身的限制.
他們.
甚至可以說.
不多去控制.
他想幫你的.
但是他幫不了.
因為他自己都自身難保.
今天世界人.
可能是依靠綜合的國力.
依靠強勢的領袖.
依靠人際的網絡.
同盟的關係.
甚至依靠自己的小聰明.
特別.
一些小聰明.
以為賺到全世界.
一場大火.
這些小聰明.
就成為他自己的囚禁.
就成為他的網絡.
對於香港人來說.
最大的依靠是什麼.
當然是報東張.
很現實的.
但是當你報東張的時候.
你要有心理準備.

$^{201}$可能要將你所有東西都要挖出來.
多數人因為報東張.
不需要拿來摔.
今天我們很依賴科技.
是科技幫助我們很快得到資訊.
幫助我們很多事情可以解決.
但是科技也幫助我們人.
去做一些不法的事.
製造戰爭.
製造矛盾.
這個就是世界的人的依靠.
究竟世人追求什麼福氣呢.
我講的是世人.
特別那些未信的人.
可能很多都是期望身體健康.
一切順利.
生活無憂.
兒孫滿堂 世界和平.
可以上樓.
年年旅行.
提早退休.
吃得 屙得 睡得.
長命百歲.
我相信這些.
大家不會完全認同.
我相信我們不會完全拒絕.
難道你不想睡得著嗎.
難道你不想身體健康嗎.
我們都想的.
這些是世人所期望的福氣.
有個特徵是什麼.
短暫.
甚至不一定是福氣.
你得到都好哭.
哎呀 終於上樓了.
終於買到我的安樂窩了.
從此就是你惡夢的開始了.
你不要想著年年旅行.
不要想著吃自助餐.
就因為這層樓.

$^{241}$你不能去旅行.
就因為這層樓.
你可能要捱麵包.
所以頂子妹.
得到不一定是真正的福氣.
但是私人告訴你.
什麼是我們真正的福氣.
一起讀吧 預備起.
陳先生 你保佑我.
因為你幫我打鼻.
我贊助你做法書.
你給我錢住.
給我福氣.
給我尊嚴.
這些聖經是我好喜歡.
是我人生的座右銘.
如果大家記得.
2002年是我們教會每月的金句.
我希望你背得到.
特別有些頂子妹背了.
這篇被稱為大衛的金詩.
我們沒有辦法知道.
真正確實的背景.
但是說到金.
金是什麼.
很堅實的.
金是可以經得起火的熬煉.
越燒就越結實.
真金是不怕熊爐火的.
但是為人生很多的經歷.
被人追殺.
甚至連命都差點沒有.
於是他作了這些詩.
就是在他最苦難.
最危險的時候.
他認定上帝是去依靠.
認定上帝是他唯一的產業.
所以被稱為金詩.
原來人生.
不一定是一帆風順就是好的.

$^{281}$太平順我們可能會躺平了.
習慣了.
你要我再起來做運動.
我做不到.
所以人生一些的熬煉.
走越大衛經歷這些苦難.
他體會到上帝的恩典夠他用.
他體會到.
他仍然在上帝的保護當中.
所以頂子妹.
我們願意在大衛的金詩當中.
因為時間關係.
以及這篇詩的分段都不容易.
時間亦都不容許我全篇去講.
我只是集中講第一節第二節.
裡面已經足以令我們今天.
掌握一個很重要的訊息.
第一個訊息很想告訴頂子妹.
而是說神是我們真正的依靠.
再讀一次好嗎.
這節聖經預備起.
神啊對你無厭惡.
因為你不怕氣.
我神對你無惡處.
對於大衛來說.
這個「你」是指誰.
是神啊.
因為他對神的就是神啊.
求你保佑.
我投靠你.
對象就是神啊.
投靠你其實就是說.
找到一個避難所.
找到避難所.
就好像我們以前.
幾十年前經濟不好的時候.
很多國內人來到香港.
投靠一些親屬.
我們在喪禮也好.
在哪裡也好.

$^{321}$我來到什麼都沒有的.
就投靠在哪裡就這樣.
投靠就是這個意思.
就在別人的保護底下.
可以繼續生存下去.
就好像今天.
仍然世界上有很多難民.
他們要投靠在什麼.
一些可以接納難民的國家裡面.
成為他們的避難所.
對於今天住在以色列的.
以色列人來說.
防空洞就是他們的避難所.
是不是.
一聽到警報.
他們就要跑去防空洞.
等消息他們出來.
什麼都做不到.
投靠在大衛的表達裡面.
是全程傳言的信靠.
就好像一個嬰兒.
全程靠著抱著他的父母的.
一樣的意思.
用廣東話來說.
很傳神的.
就是你對著別人說.
我全靠你了.
一會去給你吃飯.
我全靠你了.
是這樣的意思.
就是不再想其他東西了.
全靠你.
但是對於私人來說.
誰是唯一的依靠.
是耶和華.
我們耶和華大家都認識.
記住他不是姓耶的.
這個一定要說.
不是笑話.
耶和華是一個拼音.

$^{361}$神的名字其實就是那四個字.
YHWH是希伯來的字母.
我不讀了.
但是發不到音.
在中世紀的時候.
於是就變了一個叫做.
英文的耶和華.
因為拼音.
中文聖經就變了耶和華.
其實在希伯來的聖經裡面.
耶和華真正的名字.
是叫做雅威.
雅威是一個至聖的上帝.
這個才是真正神的名字.
在希伯來的聖經裡面.
它是說雅威的.
不說耶和華的.
所以在舊的聖經裡面.
希伯來的聖經裡面.
有很多用雅威的.
幾千個的.
無論耶和華也好.
雅威也好.
它背後的意義是什麼.
我們所信的上帝耶和華.
雅威是自有永有的.
是我就是我.
我就是存在的那位.
我是息在今在以後.
存在的全能者.
而且耶和華這個字很有意思.
是那位與我們立約的神.
特別在舊約裡面.
耶和華神與以色列立約.
今天上帝在基督裡面.
都與我們立約.
我們在基督裡面.
我們可以聚得赦免.
我們在基督裡面.
我們能夠成為神的兒女.

$^{401}$是不是.
不過在西方英語的聖經裡面.
大部分都不敢直譯耶和華.
雅威也不敢直譯.
他直譯為什麼.
主.
你看NRV譯本.
他用大家LORD.
所以英文聖經裡.
你看到那個字.
其實他表達的是耶和華.
或者雅威.
好了.
第二個很重要的字.
在詩人裡面說.
神是什麼.
神啊.
我投靠你.
在和合本裡面.
有我的心.
不過在原文裡面沒有.
所以任由他.
神啊.
這個神是用一個簡稱.
Elie.
Elie是一個神的意思.
通常是用在詩歌體裁的.
所以在舊約裡面.
詩歌體裁以外.
是用以來因的.
但是因為今天這個.
是詩篇.
他用一個簡稱.
這個是字本身的基礎.
這個字.
無論如何.
這個字是表示.
這位上帝.
不只是眾神之一.
他是獨一的神.

$^{441}$他是創造天地萬物的主.
《創世記》一章一節說.
起初神創造天地.
起初以來因創造天地.
是一位創造的主.
好了.
我的好處不在你.
你是我的主.
這個主在猶太人裡面.
他不敢叫耶和華.
不敢叫雅威.
他就說主啊主啊.
不是貪容易說的.
主在猶太人的稱呼裡面.
就是神啊神啊.
就是耶和華啊耶和華.
雅威啊雅威啊.
原來神的名字是有能力的.
神的名字是表達了.
他是一位怎樣的神.
短短一節聖經裡面.
大衛告訴你.
我們的神是創造宇宙萬物的主在.
我們的上帝是.
自有永有自存的上帝.
我們的上帝既是神.
也是全人類的主宰.
也是大衛的主.
弟兄姊妹.
今天你所信的神.
在你心目中.
是一位怎樣的神呢.
你會將祂當做主呢.
你會真的叫祂做阿道拉呢.
就是祂是主呢.
這個是值得我們思想的.
我們的上帝不只這麼有大能.
創造的主.
先知以卓啊告訴我們.
我們一起讀.

$^{481}$預備起.
我的意念.
奧秘的事屬於上帝.
上帝從來不給我們插頭.
不是上帝這麼多秘密.
第一我們的限制.
第二有些事我們不需要知道.
是不是.
很多時候我們很喜歡尋根究底.
但有些話你知道都無謂.
沒什麼需要.
上帝是給我們摸不透的全能者.
我們未認識神的時候.
我們的上帝.
我們的所謂的神.
是由我操控的.
我記得我以前都扮偶像的.
我說你不靈呢.
以後我不會再裝的了.
原來祂是我的.
我使用的僕人而已.
我是去米飯班主而已.
但今天我認識了神之後.
我知道我在神面前.
沒有任何可誇.
我只有信服.
因為祂是全能的真神.
而且先知更加告訴你.
上帝是.
很多事我們未必明白完.
我也不明白.
為什麼耶穌不直接救我們.
為什麼要這麼辛苦釘十字架.
是不是.
正所謂搞這麼多事幹什麼.
發發脆脆.
一來到全世界全部得救.
那就簡單了.
省水省力.
上帝不是用這個方法.

$^{521}$我們甚至今天都不明白.
世界為什麼這麼亂.
為什麼神的子民要受苦.
為什麼一兩個人的錯誤的抉擇.
或者私心也好什麼都好.
就要這麼多人民受苦.
我們不明白.
我們不明白上帝你在哪裡.
是不是.
但是先知告訴你.
神的路你是看不透的.
神的心你是摸不著的.
不是上帝亂來.
今天有一些所謂偉大的領袖.
他說話是亂來的.
因為他是人.
他也是私心.
但上帝不會私心.
上帝仍然掌握在你裡面.
上帝仍然坐著為王.
不要太擔心.
當我們很擔心在以色列的弟兄姊妹的時候.
剛剛過去3月18日以色列的講座裡面.
薩托基穆斯跟我們分享.
沒錯.
都不少的以色列人離開了以色列國.
因為害怕.
死人的.
但是教會仍然謹守崗位.
在他分享之前.
他說前一晚.
教會仍然有40至50個頂尺的妹妹.
在這裡祈禱,讀經,唱詩.
為國家守望.
不是四個五個.
是40至50人.
在牧師所涉及的教會宣道會.
平均年齡是25歲.
什麼意思.
不是那些可以走的都走光了.

$^{561}$不是.
留下老人家在這裡走不了.
真是愚民.
不是.
是這幫.
是他要告訴我們.
在這個戰火似乎很絕望的境地裡面.
教會頂尺仍然謹守崗位.
為整個國家守望.
薩穆斯是百里行聖經學院的院長.
他說在學院裡面的.
是一幫年輕的學生.
他們選擇留下.
服侍跟從主.
他們不是走不了所以不走.
他們很多是碩士生.
碩士學位.
很多都是本身做生意的.
去進修神學.
他們絕對可以移民.
可以走的絕對可以走.
但是他們選擇留下.
因為他們帶著使命.
留在自己的國土裡面.
甚至他被一個牧師提醒.
他說信仰不是奢侈品.
可有可無.
信仰是關乎生存的問題.
所以弟兄姊妹.
我經常叫朋友你信耶穌.
不要當信仰信耶穌是.
可有可無.
寄託.
寄託不需要信耶穌.
寄託很多事情可以做.
但是信耶穌是關乎生與死.
禍與福.
不可以不信.
你選擇信不信是另一個問題.
那位牧師提醒了薩穆斯.

$^{601}$信仰是關乎生存.
而薩穆斯語重心長的.
你提醒我.
他絕對有資格.
他處身在一個戰火裡面.
說著說著的時候去暫停.
不知道暫停了幾分鐘.
因為聽到警報.
接著他需要收線.
然後轉換另一層樓.
繼續轉播.
我們的影音幹事聽到.
牧師原來傳來砰砰聲的.
真的戰火的.
不是開玩笑的.
不是電視劇集來的.
牧師鼓勵我們香港信徒.
不要因為一些困難.
不滿就離開香港.
除非上帝要你來.
如果你隨便因為這些事情.
就離開香港.
誰為主人所作見證.
這些說話不是在香港牧師說的.
是在以色列的牧師說的.
是以一個遠方.
深深經歷到.
什麼在戰亂裡面.
仍然謹守崗位.
持守使命的態度去說這些話.
勉勵我們香港的頂子妹.
如果我們一隨便走.
如果我們全部都走了.
誰為主作見證.
誰去傳呼音.
這些朋友可能還在困苦無望裡面.
他不知道自己有永生的把握.
他沒有什麼依靠.
他連神都沒有依靠.
因為不知道.

$^{641}$如果我們全部走.
有誰能夠去關心他們.
為他們祈禱.
將真正的盼望帶給他們.
所以薩牧師勉勵我們.
他有資格勉勵我們.
頂子妹.
不是我們不可以離開香港.
但是不要隨便.
因為真正的依靠.
不是你離開這個地方.
去到另一個地方.
就是你的烏托邦.
世上有神的地.
只有我們跟從上帝.
我們才有真正的平安.
否則沒有一個地方是和平.
帶給你真正的平安.
我們是神的子女.
我們要扎根在基督的裡面.
我們要全心全意.
好像我剛才所說.
神啊我全靠你了.
我必無依靠.
我們所依靠的.
是大能永存的真神.
耶穌說我就是生命.
我就是復活.
信我的人雖然死了.
也必復活.
神是我們唯一真的依靠.
因為祂是自有永有的.
祂是全能至高度的真神.
祂是復活的主.
正如主耶穌說.
我父做事直到如今.
我也做事.
上帝沒有離開.
沒有缺席.
上帝在做事.

$^{681}$耶穌也在做事.
聖靈更加和我們一起做事.
第二方面.
耶和華所賜福.
才是真正的福.
我們一起讀.
預備起.
好.
一說到福氣.
我就想起八福.
你看看八福.
馬太福第五章.
八福不是身體健康.
一帆風順.
心想事成.
不是這八樣東西.
頭尾告訴我們.
我們的福氣是什麼.
頭尾.
最後那一句.
是什麼.
下半句是什麼.
第三節的下半句.
第十節的下半句.
是什麼.
耶穌怎麼說.
天國是他們的.
原來頂子妹.
我們最大的福氣.
不是上樓.
我們最大的福氣.
不是找到一個.
沒有戰爭的國家.
原來你和我最大的福氣.
神已經給了我們.
就是天國.
是屬於我們信耶穌.
神的兒女.
八福重點不是給什麼.
給什麼給什麼給我們.

$^{721}$你看看整個八福.
有那種天國子民品格的流露.
虛心人有福料.
哀動人有福料.
飢渴無意人有福料等等.
溫柔人有福料.
原來真正的福氣.
在於我們有沒有流露.
我們從神而來.
那種新的生命.
當你流露這些新的生命的時候.
你就一個真正有福的人.
世界的人拿福氣.
世界的人對福氣是什麼.
要拿的 拿的 拿的 拿的.
但聖經給我們的福氣是什麼.
是你從基督裡面.
流露一種不再一樣的生命.
和品格.
這裡說 我曾對耶穌說.
你是我的主 我的福氣.
唯獨從你而來.
和本就說什麼.
我的好處不在你以外.
我曾經解釋過.
我自己的解釋.
容易明白的是什麼呢.
一切的東西.
如果出於上帝.
都是對我有益的.
就好像我小時候.
媽媽逼我喝.
絲葛脫敏如乾油.
白色那種.
哇 痛苦到暈.
最後我捏著鼻子才喝完.
是的 我不想的.
不過如果我爸媽.
你要我喝.
我相信你不會毒死我 是不是.

$^{761}$他一定不會毒死我.
因為我是他唯一的兒子.
我肯定他不會毒死我.
所以就算我不喜歡.
我不明白.
我覺得很痛苦.
我就很明白.
你是我的主 我的好處不在你.
簡單來說.
我們人生一切的遭遇.
如果不是上帝容許.
是不會發生的頂智媒.
如果上帝不讓你去旅行.
一定上帝有他的意思.
上帝逼你要面對這些東西.
就算你不想的.
都是為你的好處.
和收本直接簡單來說.
我的福氣唯獨從你而來.
很簡單 是不是.
只有你 是你的路.
其他沒有了.
簡單來說.
神的福氣是真福氣.
其他的福氣.
無牌貨 假貨 山寨貨.
個個都以為是好東西.
用一用就沒有了.
啊 多好啊 你的心機.
聽一聽.
哦 壞了.
神的福氣 真福.
是有永存的價值.
其實和合本.
不在你意外.
是真正的意思.
作者大衛克義勇.
不在你意外.
是要凸顯.
那福氣真是從神那裡而來的.

$^{801}$不是在其他地方.
所以這個不字不是多餘的.
只不過我們覺得.
不字是不要消極.
凡是用一種消極的語氣.
其實是什麼.
是要凸顯.
更積極的意思.
更正面的意思.
我希望你明白大衛的心意.
天主教的版譯本.
甚至說.
我對天主說.
就是我對耶和說.
唯有你是我的上主.
唯有你是我的幸福.
在這次的翻譯裡面.
天主教的版譯本.
甚至告訴你.
耶穌說.
不是 上帝啊.
你就是我的主.
上帝啊.
你的福氣.
才是我真正的福氣.
這個只是我今天想表達的.
我們今天依靠誰.
根據私人告訴你.
我們唯一的依靠就是上帝.
我們要的福氣是什麼.
不是別人給的.
那些是山大貨.
那些是假貨.
我們真正的福氣就是.
正牌的原廠貨.
就是從神而來.
甚至好像這裡所說.
第五節所說.
是我的產業.
是我杯中的份.

$^{841}$我所得的.
你為我持守.
真正的福氣.
在私人第五節已經很清楚告訴你.
原來我們最大的福氣是什麼.
神就是我們的產業.
父母.
你給什麼都無所謂.
你已經是我父母了 是不是.
你已經是我最大的福氣.
有你做我父母.
我還擔心什麼.
明白我的意思嗎.
以前都聽過一個遇到故事.
叫一個小孩.
有個有錢人對別人說.
你們隨便拿.
你拿什麼都可以.
拿了是你的.
有個小孩.
就抱著那個有錢人.
我要你.
那你不是父母.
今天很多人就以為那些是.
那些鮑魚.
那些花膠最好 是不是.
原來神才是你唯一要拿的.
你擁有這個至高的上帝.
一切的福氣的上帝.
那你是不是最有福氣.
那剛才那個遇到故事而已.
我相信.
但是.
詩人說.
地上的福氣.
有也好.
是會過去會失去的.
你以為放在保險箱是最穩當嗎.
原來會水浸的.
無端端說.

$^{881}$那個保險箱不見了都可以.
銀行來的.
或者水浸都可以.
地上沒有任何一個絕對的保障.
真的.
但是詩人說.
神為我持守.
我們的上帝是.
無所不知.
無所不能.
無所不在.
空間都會哭哭眼睡.
是不是.
就是那一秒你哭哭眼睡就沒了東西了.
但是上帝是不打盾的.
詩人告訴你.
所以你怕什麼.
你抱著一些你保障的東西.
但是你不去交代給上帝.
這是不是很愚蠢的做法.
你擁有的不淺.
你怕不知道會不會被人偷了.
你擁有一個新的生命.
不知道會怎樣.
你有很多擔心.
但是如果是屬於上帝的.
他必定為你持守.
你有沒有這份信心.
如果上帝要保守的.
沒有人可以在神的手裡面.
奪取屬於上帝.
或者上帝給我們的東西.
這是我們需要的信心.
另一次聖經.
也是來自詩篇第十六篇第八節.
我讓耶和尚在我面前.
一他在我右邊.
我就不自動搖.
很明顯從以色列的頂子妹當中.
我們看到這句話.

$^{921}$不過有一個條件.
詩人說我讓耶和尚在我面前.
你要上帝給你福氣.
你要自己不動搖.
你有沒有讓上帝在你面前呢.
上帝會不會.
你將他放在保險箱裡面.
有需要才拿出來.
拿不完的了.
你的事將上帝放在你面前.
讓上帝引導你的人生.
讓上帝成為你的保護.
那你就能夠不自動搖.
這個就是以色列的頂子妹.
他們真的以神為首.
他們能夠清楚自己的使命.
就算在戰火這麼混亂的當中.
他們仍然不自動搖.
正如以下先知所說.
他是你這世代安定的力量.
豐盛的救恩.
智慧和知識.
敬畏耶和華.
是錫安的自保.
我們能夠穩定.
我們能夠得到真正的安定的力量.
不是環境.
不是哪個人的保證.
而是敬畏耶和華.
好了.
神繼續透過詩人給我們很多福氣.
第十一節 你必將生命的道路指示我.
在你面前有滿足的喜樂.
在你右手中有永遠的福樂.
很多時候我們都會.
為我們擁有一些東西.
我們很滿足.
很喜樂 是不是.
我每次吃完自助餐.
我都立誓說以後再不吃了.

$^{961}$但我發覺我經常都吃宴.
因為吃完自助餐.
不會永遠都不餓.
你吃完之後你會覺得很飽.
不想再提食物的東西.
但原來人面對食物.
是有不滿足的.
這是很多人生的道理.
很深的經歷.
地上的東西沒有辦法滿足你的.
如果能夠滿足.
就不會有那麼多有錢的人.
繼續用不法的手段.
繼續賺錢.
有時候看新聞你會覺得匪夷所思.
哇 你這麼有錢.
你還不是做這些傻事嗎.
就是因為人的心得不到滿足.
如果有滿足就不會做這些事.
當然有些是病態.
有些是慣性.
有些是沉溺.
但更重要的是.
人裡面沒有真正的滿足.
經濟學裡面有一種叫做無窮的欲念.
你找到別人的無窮欲念.
你就發達了.
經濟學來說.
你觸摸到人哪些東西是無盡的.
你繼續供應.
你就可以賺很多錢.
是的.
只有上帝能夠填滿我們內心那種無底心疼.
丁智梅.
願你讓上帝.
成為你的滿足喜樂.
願上帝和華遠成為你永遠的福樂.
這是真實可以經歷.
是私人可以體會的.
就算大衛面對危險苦難的當中.

$^{1001}$他找到上帝是他的滿足.
哪個人被人追殺.
被人在危機裡面笑得出來.
只有在神裡面的人.
所以丁智梅.
或許你現在也有些東西很心虛.
今早才聽到一位弟兄分享.
聽起來我們也覺得替他有點擔心.
但上帝成為他的依靠.
成為他的滿足.
是的.
正牌的福樂.
是因為他的富足並不加上憂慮.
有人的富足換來的是什麼.
是孤單.
換來的是很多憂慮.
這麼多錢怎麼辦.
會不會被人侯住.
被人綁票呢.
是不是.
我自然有了錢之後.
很多朋友不是黏著我.
反而遠離我.
是的.
會孤單的.
大家不同了.
一起讀這段吧.
預備起.
都有機會去探望林中的丁智梅的時候.
有些丁智梅擁有十分的平安.
十分就是真的平安.
很完全的平安.
不是隨便演.
不是演戲.
有一次有個弟兄.
我忘了是誰.
是弟兄.
肯定的.
他說.
我昨晚做夢.

$^{1041}$我夢見我在天堂.
我相信上帝.
要讓他有十分的平安.
沒多久當然他就死了.
神我恩待他.
讓他知道.
他的唯一歸宿就在上帝那裡.
什麼時候我們去堅心仰賴.
神我全靠你.
我將來的定局.
我離開世界.
我的歸宿.
都在乎你的時候.
你就能夠有十分的平安.
否則.
就算你一百多歲也好.
你仍然面對死亡.
你有恐懼.
恐懼不是年輕人的事情.
是那些沒有上帝在他們內心的丁智梅.
同樣出現的.
丁智梅.
有一句話.
人到盡頭就是神的開始.
你認不認同.
我就不認同.
不是人到盡頭你才是神的開始.
哎呀我拜的那些不行.
求的那些不行.
那我就唯有求耶穌了.
那些只有未信耶穌的人才會這樣做.
我們是神的兒女.
我們基礎在上帝和他的應許上面.
我們的信仰是信這位復活的主.
所以不是等到盡頭我們才依靠他.
是在任何事情開始的時候.
你都是專心的投靠神.
我把我這件事交到你手中.
求你引導.
所以我最後這首詩歌鼓勵大家.

$^{1081}$不是唱了.
不過這首詩歌是要依靠主.
是真有福的.
因為他必隨時蒙扶助.
又要承受所應許的賞賜.
必緊緊跟從主的腳步.
弟兄姊妹.
願你今天透過大衛的詩篇給我們鼓勵.
透過以色列弟兄姊妹的見證給我們鼓勵.
是有些事情我們不明白.
但是我知道上帝知道.
如果我知道神和我同在我不會擔心.
我知道這件事是上帝容許的.
是神給我的.
我就去接受.
今天你出街.
最重要是帶什麼.
手機.
一機在手.
閒通廣九.
是真的.
手機你可以吃東西.
可以搭車.
如果你綁上信用卡.
你做什麼都可以.
搭車,身分證.
你把身分證都放進去都可以.
同樣.
今天你人生要什麼在你身邊.
不要像十方姐那樣.
全身有袋子,桶子,棍子.
什麼全身綁著自己.
毫無自由.
我們拿著耶穌基督.
就是我唯一的依靠.
無論什麼環境.
有什麼事.
耶穌幫我.
耶穌很開心.
求主幫助大家.

$^{1121}$一起祈禱.
希望上帝願你透過大衛的詩篇給我們提醒.
很多時候我們嘗試用自己的方法.
去找很多我們認為的保障.
找很多我們認為的福氣.
殊不知這一切都不是真正我們的依靠.
不是我們真正的福氣.
求主讓我們今天明白.
你是我們唯一的依靠.
真實的依靠.
也是我們真實的平安和安慰.
願你的道存在我們心裡面.
奉耶穌基督的聖名.
阿們.
\newpage



\section{列王記上 8:37-45}
\label{sec:031946}
\textbf{禱告讓主名被高舉}~~~~潘建堂牧師
\newline
\newline
連結: \href{http://www.church.com.hk/acms/content.asp?site=cdc\&op=show\&type=product\&code=031946\&layout=sermon}{\texttt{031946}} ~~~~ 語音日期: 2025-09-07
\newline
\newline
\hyperref[sec:031945]{\small{< < < PREV SERMON < < <}}
~
\hyperref[sec:index_chronic]{\small{[返順時目]}}
~
\hyperref[sec:031947]{\small{> > > NEXT SERMON > > >}}
\newline
\newline
列王記上 8:37-45
\newline
\begin{longtable}{cl}
\hline
\hline
章節 & 經文 (和合本修訂版)\\
\hline
8:37 & \begin{tabularx}{0.7\textwidth}{X} 「這地若有饑荒、瘟疫、焚風、霉爛、蝗蟲、螞蚱，或有仇敵圍困這地的城門，無論遭遇甚麼災禍疾病， \end{tabularx} \\ \\ \relax
8:38 & \begin{tabularx}{0.7\textwidth}{X} 你的百姓以色列，或眾人或一人，內心知道有禍，向這殿舉手，無論祈求甚麼，禱告甚麼， \end{tabularx} \\ \\ \relax
8:39 & \begin{tabularx}{0.7\textwidth}{X} 求你在天上你的居所垂聽、赦免、處理。因為你知道人心，惟有你知道世人的心，求你照各人所行的一切待他們， \end{tabularx} \\ \\ \relax
8:40 & \begin{tabularx}{0.7\textwidth}{X} 使他們在你賜給我們列祖的土地上一生一世敬畏你。 \end{tabularx} \\ \\ \relax
8:41 & \begin{tabularx}{0.7\textwidth}{X} 「論到不屬你百姓以色列的外邦人，若為你的名從遠方而來， \end{tabularx} \\ \\ \relax
8:42 & \begin{tabularx}{0.7\textwidth}{X} 他們因聽見你的大名和大能的手，以及伸出來的膀臂，來向這殿禱告， \end{tabularx} \\ \\ \relax
8:43 & \begin{tabularx}{0.7\textwidth}{X} 求你在天上你的居所垂聽，照著外邦人向你所求的一切而行，使地上萬民都認識你的名，敬畏你，像你的百姓以色列一樣，又使他們知道我所建造的是稱為你名下的殿。 \end{tabularx} \\ \\ \relax
8:44 & \begin{tabularx}{0.7\textwidth}{X} 「你的百姓若奉你的派遣出去，無論往何處與仇敵爭戰，他們若向耶和華所選擇的城，以及我為你名所建造的這殿禱告， \end{tabularx} \\ \\ \relax
8:45 & \begin{tabularx}{0.7\textwidth}{X} 求你在天上垂聽他們的禱告祈求，為他們伸張正義。 \end{tabularx} \\ \\
[1ex]
\hline
\hline
\end{longtable}
$^{1}$此錄音或錄影的版權是屬於基督教宣導會北角堂所有.
我們的屬靈傳統之一是以真理教導為本.
歡迎你將此錄音或錄影的連結傳給親友.
懇請注意未經授權請勿轉錄或剪輯本錄音或錄影.
心願如《仕途行傳》二十章三十二節下所說,這道能建立你們,叫你們和一切承聖的人同得基業..
今天當我們打開你的話語的時候,求主引導我們的心靈.
實在你也都知道我們各人的心,你知道我們需要聆聽你聖靈給我們有怎樣的提醒.
你都很清楚我們整個群體,我們在宣教的學習上面,我們需要怎樣成長.
求主就藉著你的說話,特別在這幾個禮拜裡面,我們一同學習的時刻,你更多的去到成為我們的幫助.
我們同心仰望的禱告奉耶穌基督的名字,阿們.
踏入九月就是差傳月,所以接下來這幾個禮拜,三個禮拜都會說差傳的訊息.
然後到最後一個禮拜應該就會再回到《曾言書》的總結.
今天我們的主題叫做《禱告讓主明高舉》.
我們會特別藉著剛才所讀的一段聖經思想.
我們在差傳的學習上面,特別是在禱告上面.
我們怎樣去到重視上帝給我們的使命.
我們說禱告的時間就是很重要,禱告跟我們基督徒的生活可以說是分不開.
很劃缺的,我猜大家每天都會禱告,是吧?.
或多或少,是吧?教會來說,每一個聚會幾乎都一定會有禱告.
所以那個關係跟我們的信仰是非常密切.
而教會本身,當教會成立的時候,就是那些門徒一起聚集.
在當中他們去到祈禱,教會就這樣開始.
教會的擴展後來就藉著他們得著聖靈的恩賜,可以傳揚福音.
就像神蹟其事一樣,令到很多很多的人信主.
這個復興其實之前是什麼?起點就是禱告.
禱告跟教會是分不開的.
而當耶穌去到潔淨聖殿的時候.
祂就引述了《宇創也書》裡面說:.
「經上記著,我的殿要稱為禱告的殿」.
就是所羅門王興建的聖殿.
這個聖殿除了代替了古代會幕的憲制之餘.
其實還有特別強調祂是一個禱告的殿.
這個功能其實在舊約會幕的時候是比較少強調.
祂特別強調這個殿是讓人去禱告的.
我們剛才所讀的經文就是.
在這個獻殿的儀式裡面.
所羅門王帶領著全會眾去祈禱.
在剛才所讀的大家大概會感覺到.
是很多次去求饒或去垂聽禱告.
垂聽凡是向這個聖殿所發出的禱告.

$^{41}$因為這個方向表示人是向饒華禱告.
是要專心依靠饒華.
然後結果就是讓主的名被高舉.
所以我會說這個禱告的殿.
就是榮耀上主的基地.
甚至我會說其實禱告本身就是榮耀上主的基礎.
甚至在這裡說也是一個基地.
說到祈禱的時間.
可能不同的文化是有不同的.
當時所羅門王帶領全會眾的祈禱.
就不像是我今天這樣主席和我這樣.
有個抗音器所以仍然可以很輕聲.
大家聽得很清楚.
所羅門王要很大聲.
然後才可以去到讓會眾可以聽得到.
實際上不同的文化對祈禱的感覺是不同的.
我們華人教會弟兄姊妹會都是比較含蓄的.
如果主席請大家.
這一節請大家用祈禱的聲音唱.
大家通常會有什麼反應.
小聲地唱.
或者我們廣東話說用喉嚨定的聲音唱.
因為這就是我們平時和上帝祈禱的習慣.
祈禱和天父說的.
如果一起祈禱不要吵著隔壁就更小聲.
所以用喉嚨定.
不過根據我的經驗.
去過韓國的教會參觀.
有時候在短宣泰北去經歷.
如果主席叫大家用祈禱的聲音唱.
會是高聲唱.
因為他們平時的祈禱都是高聲的祈禱.
很多時候在贊美的祈禱裡.
是高聲一起去到眾贊主.
這是他們的性格他們的氣質.
這只是我們信徒不同文化裡.
不同性格的表然而已.
沒錯大聲祈禱.
眾贊的祈禱.
可以很清楚地去到.

$^{81}$明顯地說我就是.
這個祈禱是榮耀神.
不過如果反過來.
榮耀神是否一定要用這種方式.
是否一定要很大聲.
或者說高舉主明.
是否一定要用很宏偉的禱告內容.
或者氣氛.
我就會說其實高舉主明.
不是人的工作.
榮耀神不是人的工作.
因為榮耀神是當祂施恩的時候.
顯明祂的工作的時候.
祂的榮耀就彰顯了.
所以我說禱告.
本身是榮耀上主的基地.
不是一定要透過.
不一定.
不一定透過高聲的禱告.
這裡所說的.
何時上帝的榮耀是彰顯的.
就是當禱告之後.
蒙映雲.
上帝的工作顯明的時候.
上帝的榮耀就彰顯.
在這一段裡面.
是要帶出這個訊息.
剛才所讀的.
朵樓門王祂的禱告裡面.
它是多次.
或者首先就是.
向耀華去到呼求.
求上帝去回應祂的指紋.
在當中這樣說.
在第37節裡面提到.
無論遭遇任何災禍疾病.
你的百姓以至烈.
或眾人或一人.
內心之道有和.
向這殿舉手.

$^{121}$保頓祈求什麼.
禱告什麼.
求你在天上.
你的居所誰聽.
赦免處理.
四十節說.
使他們在你賜給我.
門列祖的土地上.
一生一世敬畏你.
在舊約的時代.
因為神的指紋.
得到上帝很好的眷顧.
保守.
常常享受平安.
所以如果有災難來了.
即是有事發生.
為什麼會發生呢.
他們就會想.
是不是我得罪了神呢.
所以他們希望透過祈禱.
神就赦免他們.
憐憫他們.
那些災禍就會停止了.
古代裡面常常都會這樣想.
如果這樣發生的時候.
當神的指紋.
經歷到上帝他的作為.
他的恩惠.
他們就會將榮耀.
畏給上帝.
多些經歷的時候.
他們的心靈.
就得到鼓勵.
在生活裡面敬畏主.
主的命.
因此就被高舉.
所有門王就是這樣去到.
為這些人去求.
不但只為神的指紋去求.
他下一步就是為.

$^{161}$外邦人去求.
第四十一節裡面就這樣說.
輪到不屬你百姓以色列的外邦人.
若為你的名從遠方而來.
他們因聽見你的大名.
和大能的手.
以及伸出來的膀臂.
來向這殿禱告.
求你在天上你的居所垂聽.
照著外邦人.
向你所求的一切而行.
使地上萬民都認識你的名.
敬畏你.
像你的百姓以色列一樣.
由此他們知道我所建造的.
是稱為你名下的殿.
這些來到上帝的殿.
去到祈求的外邦人.
應該就不是說那些.
在以色列人生社會裡面.
做勞僕的那些外邦人.
這裡說的是從遠方而來的.
應該是特登來的.
所以我們就祈求.
求上帝垂聽他們的祈求.
甚至這裡說的是.
垂聽他們一切所求.
好像不用篩選那樣.
好像特別優惠他們那樣.
其實他們在說什麼呢.
就是一切就是說什麼呢.
就是後面所說的.
目的就是叫他們.
都能夠認識主的名.
敬畏主.
好像神的百姓那樣.
所以這一切是說.
所羅門王為神的百姓.
所祈求的恩典.
希望神都垂聽他們的禱告.

$^{201}$以致同樣的臨到.
這些外邦人身上.
就是說這些外邦人.
可以享受.
以致人所享受神的恩典.
大家在右上角.
看到的是什麼呢.
筷子就好像中國人吃飯.
說什麼呢.
多傷人多傷筷子.
上帝的恩典無窮無盡.
給他有什麼所謂呢.
他是無影響的.
沒問題.
所以希望外邦人.
他們都得到.
其實這也是上帝的心意.
又好像在羅馬書裡面所說.
那本來不是我子民的.
我要稱為我的子民.
本來不是蒙我的.
我要稱為蒙我的.
嚴格來說.
從民族來說.
我們不是以色列民族.
今天我們都成為了神的子民.
這就是上帝的心意.
其實在當時來說.
興建聖殿的時間.
就不是以色列人自己.
只是為內做的事.
因為他們所選用的一些材料.
就在這一段《葉王記》上.
在左勞門王寫信給推羅王的一封書信裡面就提到.
那些木是什麼.
香柏木.
是來自黎巴嫩的.
砍伐樹木的人是西頓人.
不是以色列人.
所以上帝在以色列人.

$^{241}$他們興建聖殿的過程當中.
也都是讓周圍的民族世界認識他.
以致他們有機會.
好像以色列人那樣去敬畏神.
要認識他.
外邦人怎麼認識他呢.
剛才在所讀的第四十二節.
就說他們因聽見你的大名和大能的手.
以及伸出來的膀臂.
來向這殿禱告.
這裡說的是.
會來聖殿禱告的外邦人.
不是說好像我們今天說的外勞.
路人甲也不是這些.
他們是慕名而來的.
是說從遠方已經聽到.
尤華的大名.
大概可能好像.
《約書阿記》裡面記載.
他們知道以色列人過紅海的經歷.
那些神跡歧視很厲害.
他們的神很厲害.
一早已經聽過了.
又或者真的有些人.
外邦人他們經過祈禱.
蒙英韻.
就感覺到上帝的大能.
如果這樣的話.
我們合理的估計是什麼呢.
這些外邦人如果透過禱告.
經歷到神的恩典.
那麼開心.
他當然會跟周圍的親友說.
於是其他人聽到.
說是呀那麼厲害.
我又去了.
有點像我們今天什麼.
有人告訴你.
某個地方旅行.
很好玩的.

$^{281}$很好吃的.
你又心動又想去.
又或者說機票特別便宜.
這陣子你又心動.
我又想買了.
好處嘛.
我們真信這種好的效果.
就是所羅門王.
他會預見得到的.
好了那麼今天《萬國問》問.
怎樣能夠在禱告當中.
去經歷神.
以致去到高舉.
上帝的名得到高舉.
在我們的身上.
在羅馬書裡面十章就說.
猶他人和希臘人並沒有分別.
因為人人都有同一位主.
他也厚待求告他的每一個人.
因為犯求告主明的就必得救.
然而人未曾信他.
怎能求告他呢.
未曾聽見他.
怎能信他呢.
沒有傳道的.
怎能聽見呢.
常會有奉差遣的.
怎能傳道呢.
如經上所記.
報福音 傳喜訊的人.
他們的腳蹤何等皆明.
這一段經文大家都很熟.
我們經常都說.
所以畢宣家對差遣方面.
都有很熟悉裡面所說的.
是讓畢宣家領受了.
這個差遣的意象.
不單止只是有一個意象.
實際上帝只會很積極地.
積極地去參與.

$^{321}$很多宣教的活動.
除了很多的短宣.
還有很多的小組報名.
根據我最近看的.
2024年的年報就告訴我們.
畢宣有多少的差關小組呢.
有十八個那麼多.
應該就是很多的小組.
長期的這些小組.
很多的團契小組.
都有叫做宣教使者.
去到在團契裡面.
推動差遣的關注.
根據這個年報.
去年就已經有超過一百人.
即是有一百個大單位那麼多.
所以帝姊妹是很熱心參與的.
不只是這些.
無論其實可能是.
平時的參與短宣.
去亞洲區 歐洲區 非洲等等.
或者有些更長的參與.
信徒的宣教.
又或者是我們猜測的祈禱.
或者奉獻.
全部都是實踐這個大使命.
都是在回應神給我們的這個意象.
在裡面我們很清楚知道.
我們做很多這些事當中.
如果當我們猜出宣教的事的時候.
是不是就是我們做完了.
好像鼓勵了他出去.
平平安安地去爬這樣.
不會吧.
我們知道很重要.
禱告對出去的宣教使很重要.
聖經都有說.
保羅是一個很偉大的宣教使.
去了很多很多的地方.
他是不是覺得靠自己 不是.

$^{361}$他是跟帖撒羅尼加教會說.
「弟兄們 請也為我們禱告」.
這個「我們」是說.
寫這封書信的那幾個人.
保羅 西拉 提摩泰.
這個宣教的團隊.
需要他們的禱告.
保羅在這裡.
要說出一個很重要的事實.
雖然保羅 西拉 提摩泰.
都是教導他們.
養他們的宣教使.
但是他們需要這些被牧羊的人.
為他們去禱告.
在後書第三章裡面就說.
「無聊弟兄們 請你們為我們禱告.
好讓主的道快快傳開.
得著榮耀.
正如在你們中間一樣」.
貼撒羅尼加教會是明白的.
神的恩典 神的救恩.
怎樣臨到他們身上.
應該都臨到其他的地方.
不僅僅是宣教士要去.
那些祈禱都要跟過去.
同樣期望保羅在那羅西書裡面.
也都說.
「也要為我們禱告求神.
給我們開全道的門.
能宣講基督的奧秘.
而不所書又說.
也為我祈求」.
這個單數 保羅自己.
「讓我有口才」.
我們看完保羅所寫的東西.
他仍然覺得.
需要有人祈禱.
幫助他能夠有口才.
放膽去講明這個福音.
所以我會用一個字眼.

$^{401}$去形容祈禱.
在差泉裡面那種重要性.
就是「根」這個字.
差泉和禱告的關係.
我說禱告是差泉的根.
根就是根源.
差泉的根源就是禱告.
我們看回《史道行傳》.
第十三章.
差泉運動開始的時候.
就是保羅,西拉.
他們禁食禱告的時候.
領受差泉的使命.
事實上我們知道.
聽專教士的分享.
當他們講述他們的夢召經歷的時候.
都是經過很多的禱告.
尋求神的心意.
我是否要踏上呢?.
還有呢?要去哪裡?.
這都是經過很多很多的禱告.
當他們被差演的時候.
是需要我們祈禱.
在差泉裡面.
在那個時間裡面.
到了工場的時候.
是更加需要我們的禱告.
我說禱告是根.
因為禱告是屬靈能力的根源.
差泉很需要是聖靈的工作.
禱告就是我們順從聖靈的帶領.
在差泉的工作裡面.
禱告和行動是分不開的.
如果只有行動.
這個行動或許只是人的工作.
如果只是禱告.
但是沒有行動.
可能只是停在等候的階段.
或者還未完全信服聖靈的帶領.
在這個方向上.

$^{441}$一個差泉培訓的機構.
All Nations裡面的領袖就說.
禱告是宣教的燃料.
很重要.
在講宣教的時候.
我們一定要重視祈禱.
因為聖靈會燃點.
我們每一個的禱告.
發揮屬靈的果核.
宣教是一個屬靈的工作.
必須有很多人的禱告去支持.
更加是宣教的目的是.
要更多更多的人.
甚至是萬國萬民.
要得著救恩.
是一個很大的工程.
需要更多更多的人的祈禱.
差泉的事工.
需要大家一起去祈禱.
並不是等人去祈禱.
等人支持.
好像有點被動.
其實差泉宣教的禱告.
應該是我們主動的一種參與.
所以換個角度來講.
我們平時都會祈禱.
而我們的祈禱.
是需要一份差泉的心懷意念.
為什麼呢.
我會想起.
我們幾乎每個禮拜.
在崇拜裡面都會說.
願人人都尊你的名為聖.
這句我們經常要學習的主土文.
首先就是求什麼呢.
有三個願.
第一個願就是.
願人人都尊你的名為聖.
怎樣去實踐呢.
怎樣去實現呢.

$^{481}$我會說就是差泉.
因為主土文不是說.
願多些人尊你的名為聖就可以了.
是說願人人所有人都要尊主名為聖.
這個是我們的方向目標.
所以只要世界上還有一個人未信耶穌.
這個願都是成為我們的禱告.
我們講禱告的時間.
經常都會講有五種.
這字很小.
由手指公數下去就是.
讚美,感恩,祈求,代求,認罪.
這幾種.
差不多都包了很多範圍.
這五種範圍都和差泉是相關的.
因為當我們為差.
尊教士去祈禱的時候.
往往會有感恩.
也都會知道有很具體的需要.
在感恩的當中.
我們會讚美主.
讚美祂的慈愛.
祂的大能.
祂豐盛的榮耀.
在祈禱裡面.
我們都會坦然的承認.
其實我們自己都有很多的限制.
很多的缺陷.
不足.
求神幫助我們去成長.
代求也都很重要.
我們要去為人的需要.
為尊教士.
為工場裡面福音的對象去祈禱.
在這幅書裡面這樣說.
祂見無人拯救.
無人代求.
甚為詫異.
就用自己的膀臂施行拯救.
以公義賦予自己.

$^{521}$所以當神看到人有苦難.
需要拯救的時間.
祂是期望有人為他們祈禱.
祂在等我們的代禱.
所以我們如果經常都說.
願人人都尊主的名為聖的時候.
我們就要去進一步想.
怎樣去做呢.
就是我們需要去到一個機構.
上主幫助.
尊教士幫助我們.
去實踐主的使命.
讓人聽到福音.
信耶穌.
去祈禱.
我們要為前線的尊教士去祈禱.
為福音對象的祈禱.
那就是更加完整.
完完整整地去表達.
願人人都尊你的名為聖這個心願.
另一方面.
其實這種的代求.
這種屬靈的行動.
是最能夠將我們和尊教士.
在主的裡面連在一起.
我們喜歡說connect.
在禱告裡面.
我們和尊教士是一起.
會經歷到上帝國度的實現.
正如主通文後面那兩個願是什麼.
願你的國降臨.
願你的旨意行在地上.
如同行在天上.
所以當地上有人歸向主的時候.
一定是上帝的心意成就了.
神的國度臨到多一個人身上.
神的國度在一個地方更加被彰顯.
所以其實當我們有份去這個禱告的時候.
其實我們為了一個地區去祈禱.
那個地區的裡面彷彿就是.

$^{561}$在這個國度彰顯的裡面.
我們都有份.
所以聖子梅就是.
你可以透過祈禱.
為世界不同的地方去祈禱.
在那些地區裡面.
彷彿上帝國度擴張的時候.
在神的眼中.
在那個地區裡面.
你有份.
有你的名字.
你成為了那個地區裡面的代禱的勇士.
說到宣教裡面的祈禱.
在二十世紀的時候.
就是有一個關心.
一個中國福音志工的組織.
叫做內地會.
專門在中國的土壤裡面去傳行福音.
在一段時間裡面.
就是內地會發覺.
我們感覺到神的工作發生.
福音很興旺.
但是他們沒有什麼很偉大的策略.
又不是推動一些很偉大的運動.
但是神很奇妙的工作.
創辦人戴德生先生.
都在後知上.
上帝的工作為什麼這麼奇妙呢.
背後還有什麼呢.
有一年戴德生先生去英國聚會.
當初去到港島.
港島之後.
有一個年輕人去跟他聊天.
聊著聊著就發覺這個年輕人.
很清楚在中國內地裡面的福音的情況.
甚至他講得出.
誰某某名去信主.
我為他祈禱.
原來就是這位年輕人.
他有一位朋友.

$^{601}$他去到中國內地傳福音.
凡是他遇到福音對象的時候.
他就將名字寄過給.
當時是寄的.
就是給這個年輕人.
這個年輕人就不斷地為他們去祈禱.
原來是這樣發生的.
戴德生先生就明白.
原來禱告對宣教這麼重要.
不單止在當地的人要祈禱.
遠方的祈禱.
神同樣的去到誰聽.
所以禱告是給我們在.
一個主律的裡面.
在神的國度裡面.
我們彼此的去到結念.
今天我們很享受互聯網.
很多年前.
剛剛用WhatsApp.
不久的時間.
我記得我的外甥.
他去了美國公幹.
但我記得那天是他生日.
於是我早上就打一個.
生日快樂的字給他.
我就期望.
可能是一個小時.
還是一天之後才看到.
誰知五秒之後.
去了學會聽.
多謝舅父.
很厲害.
以當時來說就覺得很開心.
這個WhatsApp真的很好.
這麼遠的.
五秒已經答到我了.
當然現在我們不足為怪.
很經常經驗.
互聯網是會幫助我們.
這一種資訊.

$^{641}$我們可以很connect到.
屬靈上面.
就是靠著我們的禱告.
在禱告裡面.
上帝幫我們給支的結念.
一起去經歷到上帝的作為.
剛才我提過.
禱告和行動.
其實兩者是很有密切關係.
禱告會帶來行動.
但其實行動.
亦都會挑黃禱告.
有一間教會.
他們派了一對夫婦.
去到太北宣教.
有一對短宣隊.
去到探望這對宣教士.
去到的時間.
才發覺原來這對宣教士.
平時要到兩個小鎮.
去傳福音.
要駕車去的.
駕車去是一些山路.
鋪滿了很多的山石.
已經是崎嶇.
而且右邊行的時間.
那個右邊就是懸崖.
萬丈懸崖.
那些短宣隊經過.
都很害怕.
但他們反省的就是.
當他們回去的時候.
發覺他們只是這麼久.
才去一次.
宣教士不斷地去.
不斷地去到.
無這種危險.
去到宣教.
所以當他們回到教會.
去祈禱的時間.

$^{681}$他們更加的迫切.
更加的迫切.
不要說這麼危險.
聽起來都害怕.
其實我們平時去探訪宣教士.
我們都會.
更加容易去體會.
他們工場的實況.
和他們生活上的需要.
有一次我去探訪.
一個宣教士的朋友.
都相識多年.
好像我去到探訪的時間.
當然是說.
一場來到.
當然是一起去到傳福音.
第二天就擺佳堂.
這個是我一個.
很重要的學習.
因為那次.
我也是第一次.
第一次直接.
跟這個穆斯林.
回教徒.
去到有福音的接觸.
平時只是祈禱而已.
就派單張給他們.
我既然這麼沒有經驗.
我就只是講職派單張.
就是什麼.
講一句祝福的說話.
就給他們單張.
他們就收了.
我發覺我的經驗.
所以第一輪的時間.
我發覺只有一半人.
肯接他們.
我就反省.
換句說話講得好一點.
他們接受多一點.

$^{721}$就多一些人去接.
接著之後就開始有些男士.
就來我身邊.
就跟我聊天.
他們聊天的意思就是.
他們站在穆斯林的立場.
去看完這個單張.
他們就說.
這個這樣就不對.
就跟我開始辯論.
當然我就很努力地.
去跟他們分享等等.
效果是每個人都有不同的.
我就會體會到.
原來他們對自己的信仰.
是很清楚的.
他們都涉獵過.
知道我們講什麼.
所以就去看見.
這個單張裡面的內容之後.
他會不同意.
他會挑戰我們.
所以其實就是.
我會感覺到就是.
原來他們.
是這麼清楚自己的信仰.
他們這麼有準備.
那我們呢.
我們都要一起這樣學習.
我們都要有準備.
我自己在當中就說.
怎麼學習.
怎麼去跟他們分享.
跟他們解釋呢.
後來孫教士就問我.
你見過.
有見過多少個.
然後我就告訴他.
他就記錄在案.
這個是平時的習慣.

$^{761}$就是豐富他的資料庫.
我是更加豐富.
因為我講完之後.
他就明白.
又給我看一些短片等等.
原來他們在公園裡這樣辯論的.
這裡我就清楚一點.
所以當我回來的時候.
在為孫教士祈禱的時候.
有更多的真實的圖畫.
去幫助我.
我會覺得更加的熱切.
更加知道怎麼.
為這個孫教的工場.
為孫教士所面對的挑戰.
去祈禱.
所以我們看見的就是.
禱告帶來行動.
行動也都會挑旺我們的禱告.
剛才我們講到.
這種的結連就是.
我們在主的裡面.
透過禱告.
我們一起經歷神的工作.
但是這個結連.
這個connect.
更加重要的是什麼呢.
最重要的結連.
是跟基督的心去結連.
在前書裡面這樣說.
誰曾知道主的心.
誰會教導他.
至於我們.
我們有基督的心.
這個經文的意思.
本來是說.
神的智慧和計劃.
是我們沒有辦法參考.
是在講舊約的時候.
耶穌基督說.

$^{801}$不知道什麼.
沒有人知道.
上帝有這麼奇妙的計劃.
但是耶穌基督來到的時候.
我們知道.
上帝的計劃是怎麼有的.
呈現了.
是一種隱藏的奧秘.
現在是一個公開了.
已經實現了的計劃進行當中.
這個就是神的心意.
就是基督的心.
所以舊恩臨到萬國萬民.
是主的心意.
當我們為宣教的發展.
和宣教士的需要.
去到祈禱的時間.
我們不單止只是跟他們.
在靈裡面去到結連.
跟那個地區結連.
在神的心目中.
我們又可以有我們的名字.
更加是跟神的心去到結連.
當我思想到這裡.
我就記得.
我們所認識熟悉的.
尊敬的.
藤甘畢牧師他講過.
猜傳是神的心.
當我們參與猜傳的時間.
我們觸摸著上帝的心.
與此我們都是一起學習.
透過禱告和行動.
去更體貼上帝祂的心意.
讓我們同心祈禱.
因為你將救恩賜給我們.
更加是你將禱告的權利.
賜給我們.
讓我們在禱告裡面.
能夠經歷你的作為.

$^{841}$經歷你的恩典.
求主幫助我們.
在宣教的意象上面.
讓我們不斷學習.
在我們每一年.
有猜傳月的時候.
求主幫助我們.
能夠不斷的進步.
過去怎樣進步.
讓我們今年可以更進一步.
求主幫助我們.
不單止我們在每一天.
我們讀經祈禱的裡面.
我們的心轉向你.
不單止在我們每一個星期.
我們回到教會崇拜的時間.
我們的心轉向你.
求主幫助我們.
在我們的禱告裡面.
我們的心轉向你.
去體會主你的心懷意念.
體會你見到.
世界裡面這麼多的需要.
我們的心要貼近你的心.
求主就幫助我們.
這樣去轉向.
求主向我們施恩.
帶領我們能夠不斷的學習.
不斷的成長.
禱告奉耶穌基督的名字.
\newpage



\section{歌羅西書 2:6-15}
\label{sec:031947}
\textbf{超脫人生的塑造基礎 (2)}~~~~黃紹權牧師
\newline
\newline
連結: \href{http://www.church.com.hk/acms/content.asp?site=cdc\&op=show\&type=product\&code=031947\&layout=sermon}{\texttt{031947}} ~~~~ 語音日期: 2026-04-12
\newline
\newline
\hyperref[sec:031946]{\small{< < < PREV SERMON < < <}}
~
\hyperref[sec:index_chronic]{\small{[返順時目]}}
~
\hyperref[sec:031948]{\small{> > > NEXT SERMON > > >}}
\newline
\newline
歌羅西書 2:6-15
\newline
\begin{longtable}{cl}
\hline
\hline
章節 & 經文 (和合本修訂版)\\
\hline
2:6 & \begin{tabularx}{0.7\textwidth}{X} 既然你們接受了主基督耶穌，就要靠著他而生活， \end{tabularx} \\ \\ \relax
2:7 & \begin{tabularx}{0.7\textwidth}{X} 照著你們所領受的教導，在他裡面生根建造，信心堅固，充滿著感謝的心。 \end{tabularx} \\ \\ \relax
2:8 & \begin{tabularx}{0.7\textwidth}{X} 你們要謹慎，免得有人用他的哲學和虛空的廢話，不照著基督，而是照人間的傳統和世上粗淺的學說，把你們擄去。 \end{tabularx} \\ \\ \relax
2:9 & \begin{tabularx}{0.7\textwidth}{X} 因為神本性一切的豐盛都有形有體地居住在基督裡面； \end{tabularx} \\ \\ \relax
2:10 & \begin{tabularx}{0.7\textwidth}{X} 你們在他裡面也已經成為豐盛。他是所有執政掌權者的元首。 \end{tabularx} \\ \\ \relax
2:11 & \begin{tabularx}{0.7\textwidth}{X} 你們也在他裡面受了不是人手所行的割禮，而是使你們脫去肉體情慾的基督的割禮。 \end{tabularx} \\ \\ \relax
2:12 & \begin{tabularx}{0.7\textwidth}{X} 你們既受洗與他一同埋葬，也就在此禮上，因信那使他從死人中復活的神的作為跟他一同復活。 \end{tabularx} \\ \\ \relax
2:13 & \begin{tabularx}{0.7\textwidth}{X} 你們從前在過犯和未受割禮的肉體中死了，神卻赦免了你們一切的過犯，使你們與基督一同活過來， \end{tabularx} \\ \\ \relax
2:14 & \begin{tabularx}{0.7\textwidth}{X} 塗去了在律例上所寫、敵對我們、束縛我們的字據，把它撤去，釘在十字架上。 \end{tabularx} \\ \\ \relax
2:15 & \begin{tabularx}{0.7\textwidth}{X} 基督既將一切執政者、掌權者的權勢解除了，就在凱旋的行列中，將他們公開示眾，仗著十字架誇勝。 \end{tabularx} \\ \\
[1ex]
\hline
\hline
\end{longtable}
$^{1}$先讓我們在聆聽主的說話.
讓我們一起同心低頭禱告.
創造天地萬物的主.
我們的父神.
我們再次多謝你.
給了我們有生命氣息.
來到你的跟前.
好讓我們再次在你的啟示當中.
明白創造的意義.
並且明白今天我們活在地上為人.
最重要的生命的目的是什麼.
我們為此再次將我們的心靈.
再次求主你的聖靈帶領我們.
進入到上帝你啟示的當中.
明白主你真正讓人脫罪的那份善意.
好使每一個能夠認罪悔改的人.
都能夠在主你的真道上面.
能夠得著永生的盼望.
我們為此再次獻上禱告.
讓每一個無論是善上善下.
實體崇拜的每一個弟兄姊妹.
我們都能夠同樣在上帝你的啟示當中.
蒙最大的福氣.
禱告奉靠主耶穌基督.
得勝的名字而求.
阿們.
在這個世界的上面.
信主的人和不信主的人.
最大的分別就是.
信主的人會去參透其他.
在聖靈引導之下所見到的世事.
而不信主的人往往在這個世界的裡面.
卻是只看自己的看法為絕對的真理.
換言之說得比較細俗的說話就是.
一個人只認為自己正確的人.
那些人就很肯定是一個不信主的人.
而不信主的人其實在這個世界的上面.
其實是最大損失的人.
他們以為自己是已經樣樣東西都得著.
但是卻是原來活在一個.

$^{41}$他將會失去一切的狀態.
今天我們再次回到歌羅西書第二章的裡面.
在前面我們從第六節到第八節的裡面.
我們看到一個現象.
就是在這個世界當中.
很多人很懂得說很多的說話來討好你.
你喜歡聽的東西.
他就會提供給你聽.
甚至你不想投入的事情.
他也都可以被重奏輕地.
說一些你不需要有多少投入的事情.
這些現象其實.
我們現在這個資訊世界特別.
我們所謂的社交網絡.
social media.
流行的世界就最典型明白顯露出來.
你留意一下看無論在今天所看的所謂YouTube也好.
你看的TikTok也好.
你看的X也好.
你看的Flash也好.
你看所有的Instagram也好.
點出名的東西.
你都會看到.
只要你開著電話.
開著電腦.
他就會灌輸一些你喜歡聽的東西.
他不需要是真實的.
他只需要你喜歡.
你就會很自然被吸引下去.
在這樣的狀態之下.
他就會在模型當中.
藉著你的收看賺到很多的金錢.
所以當我們看回.
這段聖經歌羅西書第二章第六至第八節的時候.
保羅很清楚提醒哥羅西教會.
你們必須要非常之小心.
我們活在這個世界的上面.
縱使你是信了主之後.
你不會是免疫於這個世界對你的吸引.
所以在特別第八節.

$^{81}$第八節在和合本聖經裡面說.
我們要謹慎.
恐怕有人用他的理學和虛空的妄言.
不照著基督.
乃照著人間的遺傳和世上的小學.
就把你們老去.
我們強調了老去這個字眼.
是你很不由自主地被吸引拿開了.
他用什麼手法呢.
和合本和原文希臘文版本最大的分別.
就在次序上有些不同.
在和合本上面.
他會將有一句說話擺前了很多.
就是不照著基督.
你會留意到.
他緊接恐怕有人用理學和虛空的妄言.
馬上就說不照著基督.
在原文聖經裡面的次序.
不照著基督這個說話.
是擺到最後.
乃照著人間的遺傳和世上的小學的最後.
換句話說.
保羅當他寫這封信的時候.
提醒教會每一個信徒.
在這個世界裡面.
很多人用很多花言巧語.
去擄掠你的生命.
他什麼方法都會做.
只要你適合聽的說話他就會說.
但最不適合聽的就擺在最後.
什麼是最不適合聽的.
基督就是最不適合聽的.
在這個世界很多人經常說.
你說耶穌.
其實一聽這句說話.
廣東話.
你說耶穌其實是很強烈的反感.
就是說什麼意思呢.
耶穌的事情你不要告訴我.
為什麼耶穌的事情不要說.

$^{121}$因為耶穌說的事情是針對人的罪性來講.
人們是不喜歡聽的.
但人們喜歡聽什麼.
信義的東西.
信義的是什麼.
說一些你喜歡聽的古靈精怪的東西.
皮毛的東西.
好像他很懂得世界原始的事情.
所以說一些小學.
小學這個字眼在原文不是小學.
小學 primary school 的意思.
是說一些很飄渺.
好像很原始的東西.
用英文說.
他說 matter.
matter 你明白什麼嗎.
最原始的東西.
你說穆斯林我也不知道.
你有沒有具體的例子說.
今天什麼人很喜歡說一些微小的原始的東西.
DNA.
幹細胞.
這些很耳熟.
為什麼呢.
這些東西說了之後.
說了之後你聽完以為自己懂.
其實對你沒有影響.
但是如果說基督就不同了.
基督每樣東西都是針對你的罪性.
回教會很多人.
尤其是有些牧者.
他明白人性是這樣的時候.
他不忠於上帝.
於是乎他就說一些你喜歡聽的東西.
信義的東西.
讚頌你的東西.
他就說盡.
但是如果一說到罪的問題.
說到死的問題.
他就不說.

$^{161}$因為這些是難堪的人.
是趕客的.
不過真正信主的人一定要說到這些題目.
因為信主的人一定要面對自己的罪.
而這個正正是創世記開始.
已經說.
人因為犯了罪.
在罪的當中.
上帝不忍見人在這樣的狀態之下.
需要處理他.
在這些種種的背景之下.
保羅提出了這個警醒.
這個就是我們上文所看的.
但是知道這一切.
這個世界種種的古古惑惑的做法.
那你應該怎樣真的能夠回到上帝的裡面.
可以成為一個新造的人.
可以成為一個.
可以和基督耶穌一同死.
同一同埋葬.
以至於可以再次復活.
有生命的盼望呢.
保羅就在今天這段聖經.
從第九節開始解釋.
所以在和合本裡面.
就介紹了一個很重要的字眼.
人們認為.
因為什麼東西.
我們才和這個世界的做法不同.
沒錯.
你可能在這個世界裡面.
你和人可以打交道打得非常好.
你可以在這個短短的人生裡面.
和人打交道.
討人喜悅.
然後可以在當中賺取到很多利益.
沒錯.
這個真的是事實.
我不怕說我以前在香港做會計的時候.
我的老闆就是我的上司.

$^{201}$都和我說.
我在做數的時候.
我發覺有些數是不清楚的時候.
我和他說.
我說這些是什麼數來的老闆.
因為我做下去我簽名的.
他只是說了一句.
你不知道這麼多東西.
你照做就行了.
我說你如果要我照做.
不如你簽名吧.
他馬上翻臉.
在辦公室那裡罵到全層樓.
全個辦公室都知道.
說我不合作.
你如果和他合作.
那條數可以有百幾二百萬可以拿回來.
自己同事自己慢慢用.
所以你看看.
在這個世界上面.
基督徒處事和這個世界的人處事.
其實是截然不同的.
怎樣不同.
就是這段.
從第九節開始的「因為」要解釋給我們知道.
好了 不如我們一起回到這段聖經.
我們讀讀好嗎.
Joe剛才字正腔圓.
我們讀了一次.
我們再次自己讀讀熟悉這段經文.
說些什麼內容.
好 請大家一起打開.
《歌羅西書》第二章.
從第六節到第十五節.
我們一起再讀一次.
找到一張我一起讀.
你們既然接受了主基督耶穌.
就當遵行而行.
在他裡面生根建造.
信心堅固.

$^{241}$正如你們所領的教訓.
你們要謹慎.
恐怕有人用他的理學和虛空的妄言.
不照著基督乃照人間的遺傳.
和世上的小學.
就把你們擄去.
因為神本性一切的封城.
都有形 有體地居住在基督裡面.
你們在他裡面也得了封城.
他是各樣執政掌權者的元首.
你們在他裡面也受了.
不是人手所行的割禮.
乃是基督使你們脫去肉體情慾的割禮.
你們既受洗與他一同埋葬.
也就在此與他一同復活.
都因信那叫他從死裡復活神的功用.
你們從前在過犯和未受割禮的肉體中.
死了.
神赦免了你們一切過犯.
便叫你們與基督一同活過來.
猶圖沒了在律例上所寫.
攻擊我們 有礙於我們的自具.
把他撤去 釘在十字架上.
既將一切執政的掌權的老來.
明顯給眾人看 就像著十字架跨勝.
上週我們過了復活節的節期.
讓我們一起思想基督耶穌的復活.
在這個世界有人說復活是真的嗎?.
吹水而已.
騙人而已.
死了怎麼會復活呢?.
沒錯 死了的人根本是不會復活的.
是不是呢?.
這種說法就是我們先前所說的.
這些世上的小學.
人從創造的時候.
很清楚地告訴我們一件事.
人不是只有軀殼.
不是只有短短幾十年的時間.
因為神造人的時候.

$^{281}$有將神的靈擺在人的裡面.
人的靈和神的靈是同樣的靈.
所以神的靈是永恆的話.
人的靈也是永恆的.
困惑點就在這裡.
因為人犯了罪之後.
這個身體被敗壞.
而這個靈也同樣開始被侵蝕.
所以在聖經裡面.
包括誕耳理書裡面講得很清楚.
原來人無論你信主好不信主好.
你最終都會復活的.
不過你的復活有不同之處.
當人如果相信上帝罪得赦免.
我們稍後會說.
上帝就會給他一個新的身體.
一個不扭壞的身體.
是可以活在永恆裡面的.
如一個不信主的人.
人死如燈滅.
No, this is not the right thing.
這個不是事實.
因為人間講出來.
讓你容易聽.
一了百了.
原來聖經告訴他.
人縱使你不信主也好.
同樣會復活.
上帝會給你一個新的身體.
那個是原本會扭壞的身體.
你說是嗎.
穆斯林我不擔心.
我火化的.
我燒了.
或者我埋葬了.
化了水.
你不用擔心.
上帝會給你那個原本的你.
你說怎樣.
這個上帝的工作不到你說.

$^{321}$因為他是你的創造主.
他有辦法將舊的身體給你.
但當你用舊的身體去復活的時候.
你就會在復活過後被審判.
然後去承受你所擔當的罪.
是甚麼懲罰.
我們基督徒為甚麼要傳福音告訴別人.
最主要的原因就是.
我們不想看到人在這樣的光景.
而且是一個永恆的事情.
更重要的一件事是.
因為我們看到.
上帝的心意是讓人活在他裡面有永恆的享受.
在上帝的創造裡面.
能夠反映出神的創造的榮美.
那些不信的.
或者傳一些想你不信主的人.
他因為容易.
你也覺得很容易.
所以我不信主.
所以我就可以很自由.
做甚麼都可以.
是 你可以做甚麼都可以.
上帝也可以做甚麼都可以.
但是到上帝做甚麼都可以的時候.
就不到你去選擇.
保羅在這裡提醒了歌羅西書的教會.
讓哥羅西教會的人.
他明白到有這麼大的吸引力.
那不如回來.
我們看一下.
上帝怎樣扭轉這個人會死亡的局面.
而怎樣帶來一個新的生命盼望給我們看.
第九節至第十五節.
這裡上帝透過保羅這個過來人.
去介紹基督耶穌這位主.
是怎樣為你和我這些相信上帝的人.
肯認罪悔改的人.
進入一個超脫的人生.
是和我們一般世上所有其他人所傳的信息.

$^{361}$是截然不同的.
是脫離開他們所看的人生.
進入一個上帝創造人的人生.
說甚麼呢.
所以一開始在這裡說.
因為神的本性.
一切的豐盛.
都有形有體地居住在基督裡面.
因為基督是上帝的兒子.
祂來到這個世界的上面.
祂就將上帝的豐盛.
以一個實體形式展現出來.
放在你和我的面前.
請留意這句說話.
這個實體是上帝豐盛.
豐盛這個字眼可以解釋為.
完完整整上帝的神聖.
上帝是看不見的靈.
看不見我就不用相信.
看不見即是沒有.
是不是.
不是.
很多東西你都看不見.
空氣你都看不見.
氧氣你都看不見.
就算大氣電波.
大興子媒網上的.
現在你在看直播.
我們的直播都是用網上看不到的東西.
叫電波.
在傳遞給你.
所以有一點邏輯的人.
你都會明白.
看不到不代表不存在.
但是上帝將一個看不到.
用人看不到的神.
以一個實體的耶穌展現在這個世界裡面.
二千多年前祂來到這個世界裡面.
展現給我們看.
上帝是怎麼想的.

$^{401}$我的女兒剛剛結婚.
今晚擺酒.
大興子媒請準時到.
聊聊天也好.
在很多的事情裡面.
上帝享受的是你跟祂的交心.
一個好的婚禮.
不是搞得多麼堂皇.
而是你跟賓客是不是有交往.
所以昨天祂叫我去借飯禱告.
我強調的是一班.
所有家人和朋友的關係的建立交心.
在藉著一個婚姻的成立.
而建立一個更穩固的關係.
我用了一個字眼.
Bondage.
大興子媒.
上帝是想我們跟祂建立一個關係.
所以我們看不到祂.
祂就展現給我們出來去看.
就用了耶穌基督.
顯然是將上帝原原本本的屬性.
就是這裡所說的.
上帝的本性一切的豐盛.
一個原備的上帝.
將祂顯露在所有人的面前.
因此當哥羅西教會.
他們能夠真的相信耶穌基督之後.
他們理應他的生命.
同樣充充足足地.
好像有基督耶穌.
那份顯現出的神聖的實體現象.
所以基督徒會在這個世界上.
其實有時候你不需要告訴別人.
別人會看得到 聞得到.
大家都知道我都在稅務局做事.
稅務局是一個講錢的地方.
報了稅沒有 報了是吧.
在做講錢的地方裡面.
我最初在稅務局是做聽電話的.

$^{441}$搭一些普羅大眾.
好像你和我這樣的人.
他打電話進來.
他說你為什麼打我稅.
我就跟他說.
我說在你說我們要打你稅之前.
你有沒有看一下你報的稅對不對.
曾經有個女孩子打來中學老師.
她收到一張單.
不是罰 要她還回退稅錢.
七千多元加幣.
她哭在電話裡.
她說怎麼搞.
我說你報錯了稅.
她說我的朋友幫我報稅.
我說那你的朋友幫你報錯.
她說我現在沒錢.
我說為什麼你沒錢.
你剛剛退了一個月.
你就把錢還出來.
她說我拿了那七千元去玩了.
我剛剛去了歐洲旅行.
七千元挺富有.
我說你就還回那七千元.
我沒錢.
在電話裡她哭了很久.
我跟她說所有的方法.
當然是正途.
是稅務局可以幫她慢慢還的.
稅務局是一個講人情的地方.
很多人說是嗎.
她說算追稅嗎.
不是 講人情的地方.
為什麼 我知道.
因為你牧師在裡面做的.
我都跟你講人情.
你不是故意錯 我不會罰你錢的.
頂尖門.
在對話裡面.
我的老闆坐在我旁邊.

$^{481}$聽了我講了.
我的老闆跟我說.
你是不是一個基督徒.
你可以不用講那麼多東西給她聽.
我說但是.
無論是作為一個基督徒也好.
作為一個稅務局的公務員也好.
我們有責任將最好的知識告訴她.
因為我們是要講一個真相給人知道.
你問那條我是不是一個基督徒.
是 我是一個基督徒.
我作為一個基督徒.
我是一個講真理的人.
講道理的人.
我不會跟你爭辯的.
所以 頂尖門.
我從來沒有叫你爭辯她信主的.
你會不會留意.
我不爭辯的.
上帝的說話講給你聽.
你有靈性就會聽得到.
沒有靈性的人.
你跟他爭辯是沒有意思的.
回來看這段聖經.
上帝要求的就是.
將每一個基督徒都有一個.
上帝的原原本本的生命本性顯露出來.
就是基督顯露出來.
所以跟著第十節說什麼.
你們在他裡面也得了封聖.
他是各樣執政掌權者的元首.
這句話就告訴你.
當我們同樣在上帝裡面的時候.
你跟我的生命反映出來.
就像基督耶穌反映出來的神的屬性.
什麼屬性.
我們是想人回到上帝裡面.
而不是想人入教.
今天的教會最錯的地方是什麼.
就是想抓人入教.

$^{521}$為什麼要抓人入教.
因為我的工作要靠你們教徒奉獻先生存.
就像一些苗族一樣.
這些想法就是世上的小學.
世上的傳統宗教觀.
他扮宗教的.
所以他要靠吃教法.
我不否定在這個世界有很多教會是這樣做.
不過你來到以馬來尼還聖信會.
我就不是這樣的做法.
也不是這樣的想法.
我是想你回到上帝得到永恆的生命.
因為我們的上帝就想有這樣的生命.
當我們能夠像我們的上帝.
活在這個地上的時候.
這個就是上帝的心思.
也是為什麼他要把我們在這個世界上面.
擔當不同的職位.
甚至在這裡說.
是各樣掌權者的元首.
因為每一個在這個世界上面的基督徒.
應該反映出來.
是上帝怎麼管治這個世界.
而不是和這個世界同流合污.
你說牧師很難的.
在這個世界不同流合污.
這個世界不懂幹活.
這句話也是這個世界的人說的.
你懂幹活一點吧.
曾經有個信徒跟我說.
他說黃牧師.
你整天都在這麼小的教會裡面事奉.
你說得又這麼好.
你不如去大教會詛咒多一些人吧.
我不需要去大教會都可以詛咒到很多人的.
問題是你是不是一個上帝想詛咒的人.
而你是不是一個可被上帝詛咒的人.
只要你能夠明白上帝的心.
你作出所有的決定.
和在這個世界的運作.

$^{561}$就上帝喜悅你在其中.
這個也是保羅告訴哥羅西教會的信徒知道.
那怎樣可以做到這件事.
是不是靠我們幾力.
是 我們要出力的.
信專是要出力的.
要有投入的.
不是說.
說句我信就算了.
很多宗教都是這樣.
你信不信.
信.
信就行了.
接著就進入一種交易信的信仰方式.
交易.
大家明白嗎.
你來酬神.
我就給你一些好處.
基督教不是這樣想的.
上帝不需要這些東西的.
如果你看回舊約聖經.
上帝根本不稀罕那些祭物.
所以聖經舊約說.
上帝說我喜愛祭祀.
不是的.
祂要的是你能夠和祂同心.
怎樣同心.
要進行一個手術.
因為你現在不同心.
上帝提供了一個手術給你.
十一節.
祂說你們在他裡面也受了.
不是人手所行的割禮.
乃是基督裡面.
乃是基督使你們脫去肉體情慾的割禮.
你們既受洗與他一同埋葬.
也就在此與他一同復活.
都因信那叫他從死裡復活的神的功用.
你們從前在過犯和未受割禮的肉體中死了.
神赦免了你們一切的過犯.

$^{601}$便叫你們與基督一同活過來.
在這裡連續幾節經文有一個字眼.
你全部字你都知道它是甚麼.
不過你完全可能不知道它想說甚麼.
這個字眼你一定要解釋清楚.
你就會明白是甚麼意思.
就是一個字割禮.
circumcision.
首先祂提出這個割禮不是人手的割禮.
因為猶太人有一個傳統.
只要你們身為猶太人的第七天.
你就會進行割禮這個儀式.
割禮在哪裡記載過.
在創世記第十七章.
從第九至第十四節.
我們一起打開好嗎.
我們一起看看這段聖經說甚麼.
創世記第十七章九至十四節.
創世記是整本聖經最開頭的第一卷書.
亦都是所有整個基督教信仰上的基礎.
第十七章第九至十四節.
找到之後我們一起讀好嗎.
我們一起讀吧.
神又對亞伯拉罕說.
你和你的後裔必世世代代遵守我的約.
你們所有的男子都要受割禮.
這就是我與你並你的後裔所立的約.
是你們所當遵守的.
你們都要受割禮.
這是我與你們立約的憑據.
你們世世代代的男子.
無論是家裡生的.
是在你後裔之外用銀紙從外人買的.
生下來第八天都要受割禮.
你家裡生的和你用銀紙買的.
都必須受割禮.
這樣我的約就立在你們肉體上作永遠的約.
但不受割禮的男子必從民中剪除.
因他背了我的約.
割禮根據這段聖經的內容.

$^{641}$是一個上帝與人立約的標記.
這個標記是很特別的.
它是立在男性的身上.
是將男性陽具的外層的皮切除.
在這個行為上.
很多人會覺得做了.
尤其是猶太人.
覺得做了就有一個真正屬神的身份.
不過當我們去看整段聖經.
其實它強調的不是在上帝的約身份那麼簡單.
它強調了一個字眼.
重複出現過很多次.
就是世世代代.
你有沒有留意?.
很多大興之輩讀到這段聖經.
就只是專注於遵守神的約.
沒錯.
遵守神的約是事實.
但不是一次的事情.
而是長期恆常的一件事.
而且是傳代地去到繼續代代相傳.
換句話來說.
這是一個要有教導和指摘.
或者督摘的工作.
為甚麼上帝要和以色列人.
立約納得那麼.
即是用廣東話.
納得那麼奇怪呢?.
要割了羊皮.
還要走在男性身上嗎?.
因為有幾個重要原因.
其中一個最重要的原因.
如果你要保持世世代代的話.
你就要長期被提醒.
男性當時在家庭裡是作主的角色.
我再重複說多次.
作主不是大事.
不是專權的世界.
有些大興之輩特別今天這個世界.
很主張男女平等.

$^{681}$其實男女平等這個觀念.
根本是邪惡的基本上.
要完全平等是不可能完全平等.
你試試男士走到女廁.
看看裡面的女人叫不叫.
叫你非禮強姦.
你說為甚麼女士衝進男廁.
又不見有人這樣告她呢?.
歧視.
所以其實男女是有分別.
有分別不是歧視.
一個家庭.
男為作為最終的決策者.
我經常說我家裡最高決策人是誰.
你看他現在坐在這裡.
他遠遠都在遙控家裡的生意.
為甚麼?我媽媽不行嗎?.
我媽媽做生意都很厲害.
你問問她厲害嗎?.
不敢回答你.
因為她怕說錯話.
當然厲害.
但他們要請示過我爸爸.
為甚麼?歧視嗎?.
不是.
立約為甚麼要把國禮行在男士身上.
因為男士經常去洗手間也好.
他又見到.
他行房的時候跟妻子行房.
又見到.
這些只是正途事.
那些不正途的時候是做甚麼?.
照忌的話.
糟糕了.
對 跟我老婆有婚姻關係.
多麼忌憚.
有些人提醒也不照做.
因為她根本生命沒有改變.
但大姐姐妹.
聖經用了一種這麼特別的做法.

$^{721}$是要這麼私隱地不斷提醒男性.
你在這個家庭擔當的角色.
你要常常幫助你的家庭.
提醒他們在上帝的約之中.
你要照著行.
今天教會不需要行國禮 為甚麼?.
因為今天這段聖經.
講了一個事實給我們知道.
這個提醒我們要常常守約的做法.
不是用人手刻在人的身體上.
而是透過主耶穌基督.
所以在第十一節說.
你們現在所受的國禮.
和你們身上的基督徒.
不過其實很多都是非猶太人.
他們可能沒有行過國禮.
但你們也同樣有一個國禮在你身上.
因為這個國禮不是人手刻在身上的國禮.
即是再不需要男士常常去廁所提醒他.
你要守約.
你去行房又要提醒.
做壞事的時候你去召妓也好.
你也要提醒.
不需要了.
因為這個國禮不是用實體來刻在人的身體上.
乃是基督使你們脫去肉體情慾的國禮.
這個是由基督耶穌所施行的.
在屬靈生命層面的上面.
將你的情慾割了出來.
所以基督徒現在行上帝的說話.
再不是因為你身上有沒有行過人手的國禮.
而是上帝主耶穌基督的說話.
來管束著我們每一個.
管束很嚴重的牧師.
你可以說是很嚴重.
不如說提醒好不好.
都可以.
不過提醒是有一個權柄在後面.
在這個過程之下.
人的生命就截然不同了.

$^{761}$怎麼不同.
復活節的信息來了.
我們看看十二節.
他說你既受洗.
即是用古今東的文化.
受洗其實是一種投入得到潔淨的形容.
與他一同埋葬.
與耶穌一起死.
死這個經驗.
是將你的情慾埋葬了.
然後與他一同復活.
復活就是我們上個星期所說的復活的信息.
上帝無論你信主好不信主好.
都將你復活.
不過復活的時候信了耶穌的人.
改變了上帝的生命的身體給你.
可以存留在永恆的裡面.
如果你說我不信.
我就是不信.
人人說我也不信.
我全對.
OK.
你既然這樣說.
上帝說沒所謂的.
他不會強迫你的.
我把舊的身體還給你.
不過我把你放進罪的火焰裡面去受.
但如果與上帝的耶穌基督有同樣生命的人.
他就可以同樣得到一個復活的身體.
我們因此是因著信.
叫他從死裡復活的神的功用.
同樣活在我們身上.
所以從這段聖經的記載裡面.
強調給我們看的.
那個人的生命的改變.
不是純粹是一個肉體的行為.
不是一個只是一時三刻的一種感覺.
或者衝動.
信主全部不衝動.
衝動的信仰叫迷信.

$^{801}$真正的信仰是恆常地提醒自己.
走在上帝的道理裡面.
所以為什麼我們做傳道人.
是不斷傳道.
不斷提醒大家聖經說什麼.
對我沒有獎金沒有獎品.
但是有個很好的身份去到上帝裡面.
如果有這個身份.
有這個生命之後.
會怎樣呢.
十三節開始告訴我一個很特別的改變.
他說你們從前在過犯和未受國禮的肉體中死了.
神赦免你們一切的過犯.
便叫你們與基督一同活過來.
又逃沒了在律例上所寫.
攻擊我們有礙於我們的字句.
把它刪去釘在十字架上面.
這一連串的說話在說什麼.
這一連串的說話是告訴我們.
當如果一個人從未信主之前.
進入到信主之後的改變.
有什麼是拆解了.
當我們去看第十三節裡面說.
你過去的現象是怎樣的.
你在犯罪裡面.
未受國禮的時候.
你是死在這個罪當中的.
但是在這樣的狀態之下.
上帝不是要你否定這個事實.
的而且確我們是曾幾何時.
在一個罪的狀態之下.
我們要承認的.
不要去遮著對眼以為.
叫做什麼.
去掩耳盜鈴.
沒有人聽到我聽不到我看不到.
就等同於世界什麼人都看不到.
你知不知道現在很多基督徒都是這樣.
很多人都是這樣.
我認為不是就不是.

$^{841}$全世界都不是.
這個世界的對與錯真實.
不是到你去否定.
我看不到就沒有事實發生過.
不是.
你認為是什麼事.
你認為也不一定是這個世界的事實.
在這樣的種種狀態之下.
你只要承認自己的昔日的狀態.
所以寫德究見證一定要寫你過去.
你其實怎麼知道你自己的罪是什麼.
但是當你知罪之後.
你能夠悔改的話.
上帝就赦免你的罪.
怎麼赦免.
基督耶穌入手.
用了一幅很適合哥羅西這個城市的描述.
便叫你們與基督同活過來.
又用陶末了在律例上所寫.
攻擊我們有礙於我們的自居.
把它刪去釘在十字架上.
我說為什麼很適合哥羅西教會聽這句話.
哥羅西是一個什麼地方.
是一個重要商業交匯點.
在今天的土耳其仍然是一個重要的交通.
陸路交通的要塞.
是一個商業城市.
在一個商業的城市裡面.
做商業就有很多票據.
我家裡做生意的.
所以我們很明白的.
你知道那張單是什麼意思嗎.
加拿大也是一樣的.
你買賣一定開單.
譬如你去餐館買東西不開單.
你就知道他肯定有些狡猾的東西.
逢開單或者憑據的意思.
就是說你跟他有個交易.
你欠了他的貨銀.
這個意思簡單一點就是債.

$^{881}$你欠了別人的債.
當你欠了別人的債.
你就要還債.
這說什麼.
人原本犯罪充滿罪債.
不過因為你肯認罪.
有個人肯出面幫你還了他.
他叫耶穌.
罪債沒有人可以幫你還.
只有耶穌一個可以.
因為他來這個世界.
就是要去幫人還罪債.
不過你要肯信他他才行.
你不信他.
他就算肯幫你還.
也沒有意義的.
對你沒有效果的.
他還了幫你還了.
你說我不信他.
他不是幫我還債的.
那就自己負債.
不過從這段聖經裡面.
當你肯信主.
能夠認罪悔改之後.
因為他幫你還債.
於是乎你就可以活過來.
不用受這個罪債的結果.
因為罪的結果就是死.
肉體的死亡.
於是乎怎麼樣.
你肯認的時候.
這條罪債就會被刪除.
被刪除.
所以這是整本聖經.
唯一一次用這樣的陳述.
就是有一張字句.
來記錄了你的罪債的紀錄.
就好像一張單.
罪有一張單.
記好你的罪.

$^{921}$是甚麼.
要結帳.
不過這次結帳.
就不用你自己結帳.
是誰.
耶穌幫你結帳.
因為這個債你沒有辦法還.
當你願意被耶穌還了這條債之後.
弟兄姊妹請留意.
十四節說是圖抹在.
律例上所寫.
攻擊我們的債.
為甚麼耶穌幫我們還罪債.
是這麼重要呢.
因為他幫你還的這條債.
沒有人可以回來追你.
他還得清.
你有沒有經常收到電話.
有些人說你是不是欠錢.
要不要借些錢幫你還債.
報稅期間最多是這些.
電視YouTube經常跳出來.
誰需要還稅的債.
銀行有些低息的債務.
可以讓你借來還債.
甚至有沒有那些.
RSP早幾年前.
大家記不記得.
退休保險金.
逢報稅期間.
銀行就會問你.
公了沒有RSP.
銀行可以低息借給你.
收音機經常都賣.
報紙都賣.
鼓勵你多借一些.
那些債你一輩子都還不了.
罪債更加難還.
因為你沒有辦法還.
不過耶穌幫你還清了.

$^{961}$沒有人追到你債.
沒有人來追你.
他怎樣來追擊你.
你都沒有效.
兄弟姐妹.
當這一切的罪債.
上帝將它.
透過耶穌.
將它撇除之後.
不是我說的.
是聖經說的.
把它撤去.
撤去是什麼意思.
我們讀會計就清楚了.
Write off.
撇帳.
當你沒有還債了.
還完了.
沒得追.
Write off.
那怎樣.
就沒有人再追到你.
追責你.
追趕你.
一個人認罪悔改.
是一個最輕省的人生.
一個人不是真正認罪悔改的人.
就會將很多重擔掛在自己身上.
他的心就長期被壓制了.
所以在這段聖經裡面.
最後15節說.
既將一切執政的掌權的老來.
明顯給眾人看.
就像著十字架跨盛.
這是說什麼.
這句話如果我們看回英文聖經.
Disarming the rulers and authorities.
He has made a public disgrace of them.
Triumphaning over them by the cross.
首先在圓文聖經連續用了.

$^{1001}$兩個持續的現在進行式的participle.
持續地做的.
做哪兩樣東西.
就是在公眾的地方裡面.
首先解除他們的威脅.
這個世界上所有的人.
什麼權力.
什麼追我們的事情.
都被上帝解脫了.
撇除.
然後怎樣.
這些原本想來擄掠我們的空真.
讓我們全部無權再去追究我們.
然後聖過的十字架.
持續在我們的生命上有影響.
只要你一天耶穌擔保著你.
你都不用怕有人來追你.
你犯過錯的.
梁子梅你知不知道這個世界上.
有一件事最怕的就是.
結婚的時候播片.
那條片有很多你小時候的出獄照.
那些做壞事的照片.
播出來的.
不過這些是小事.
最大的事是什麼呢.
就是如果你曾經做過一些壞事.
有沒有被人揭發的 你想著.
做壞事這個世界一定被人揭發的.
你不用擔心.
一定不過看用什麼方式.
但其中有一個人知道.
是你的死黨知道.
你的夥伴知道的時候.
那你一生就被他控制了.
怎樣控制你.
他會無時無刻提醒你的.
喂 上次那件事你知不知道.
你做過些什麼.
你突然打個噠.

$^{1041}$是嗎 你都知道.
你夥伴怎麼會不知道.
現在我有別的事你幫忙.
你幫不幫.
英文有個字眼.
這叫 Blackmail.
你一生被他擄掠著.
他無時無刻就來提醒你.
特別他有事的時候.
要你做他的噗的時候犧牲你.
他就用這些方法.
你犯的罪越大.
他綁你的技術就更加有效.
但這說什麼.
耶穌勝過了這一切.
人最大的問題是罪的問題.
罪最多時候來譴責我們.
就是用律法來顯露出我們是不合法.
誰最喜歡用律法來威脅我們.
就是撒旦.
他用上帝的律例來威脅上帝所創造的人.
這招絕了.
但耶穌說不用怕.
你肯信我 我保你.
不過你要肯認罪悔改.
不單止給你身份.
我保住你.
可以在上帝裡面活著.
這個世界有哪一個肯做這件事.
一分一毫都不收你錢的.
經常我們聽見這個世界.
Follow the money 跟著錢走.
沒錯這個世界真的跟著錢走.
但基督徒活在這個世界上面.
我們不會因為錢而出賣.
我們在上帝裡面這麼重要的身份.
因為我們知道.
我們這個身份是活在永恆的裡面.
不只是這幾十年之間的事.
你計算一下.

$^{1081}$如果你是這麼厲害的.
你是會計師.
你是做商家的.
你是自己會計算的.
你以為損失了幾十年.
因為你信主而決定.
在這個世界所損失的.
這個世人所謂的利益.
同時相比在永恆的裡面.
在上帝的永恆的創造裡面.
這條數是怎麼計算的.
如果你繼續要選擇.
只是在這個世界活這幾十年裡面享受.
而喪失了這個永恆的生命.
我想你才是世界上最大的錯誤選擇.
哥羅西教會在這段聖經裡面.
保羅充分流露出他對教會的關心.
他希望哥羅西教會的人.
回憶他們在一個如此珍貴的身份裡面.
而帶著感恩的心而去活.
好像第2章第7節最後一章所說.
來一路增長在神國裡面.
為什麼要有帶著感恩的心.
因為感恩就是人生能夠推動我們.
活得更加像上帝的原料.
我曾經認識了一個年輕人.
這個年輕人很小的時候.
他爸爸已經去世了.
他只剩下母親.
他母親也沒有受過什麼教育.
經常都覺得自己很沒用.
但是他又不懂得怎麼幫兒子.
於是這個年輕人希望找到他的親朋戚友.
可不可以幫助一下他們.
但是每一個他見到的親戚朋友.
每一個見到他這樣的狀況.
都是回答.
我們都很不行的 怎麼幫你呢.
我相信最主要的原因在他的分享裡面.
或者我觀察到的就是.

$^{1121}$非常麻煩.
好像中國人說的那句.
「國家自掃門前說」.
他無可奈何 到處去找.
後來他找了一個他的親戚.
他不介意幫他.
他見到他沒有書讀.
於是乎他到處去找學校.
看看有沒有得收.
終於找到一間學校.
雖然不是很好的學校.
不是很出名的學校.
甚至可以說是很多人不喜歡的學校.
但是至少有書讀.
他幫了他進去那間學校讀書.
後來當他小學畢業了.
那間學校因為沒有什麼學校.
承認他的學歷.
升不了中學.
他這個親戚又去幫他找一些中學.
終於找到一間中學.
跟校長說.
「你讓他試試吧 他很願意讀書的」.
「你幫幫他吧」.
因為我有份去.
校長就這樣試試.
這個年輕人又成功.
又讀到書.
然後還很成功地繼續讀上去.
雖然這個年輕人後來在這個世界.
不是什麼很出眾的人物.
不是什麼名流.
沒有大富大貴.
但是有一件事.
他卻帶著一份感恩的心.
在他的生命上面繼續活下去.
什麼人他都願意幫助.
無論大或小的事他都願意幫助.
透過這件事我可以看到.
一個人帶著感恩的心.

$^{1161}$所發生的功效.
那種的影響能力是多麼大.
丁子梅.
耶穌基督在這個世界上面為我們死.
我們肯相信祂.
你其實對祂有多少感恩的心.
還是你只是捍衛一個宗教.
我回到祂就相信祂就回到祂.
回到祂就做一些儀式.
做了就有履行了我的責任.
還是你真的帶著感恩的心.
活在上帝的裡面.
人能夠有感恩的心.
都有如此的功效.
你和我作為基督徒.
是不是這份感恩的心更加強.
因為我們得到的拯救.
是一個永恆上主包底的拯救.
既然如此.
為什麼我們的信會不是全程投入呢.
求主幫助我們.
哥羅西教會能夠有一位.
跟祂素未謀面的保羅.
去寫一篇這樣的書信.
來這裡告訴他們.
你的人生是上帝欣賞的.
所以你更加要在這個信上面.
曉得用心的回應.
沒錯 上帝未必一定需要你.
捧回兩隻金豬.
斬回兩隻油雞去祭祀祂.
上帝只期望一件事.
就是你有沒有和祂同心.
讓我一起同心低頭禱告.
父臣多謝你給了我們.
在這段聖經裡面的學習.
我們作為一個基督徒.
在這個世界上面.
是不是真的成為一個.
上帝你所喜悅的人呢.

$^{1201}$是不是一個有感恩過活的人呢.
還是徒然只是認為.
自己永遠是對的呢.
求你的幫助.
讓我們每一個無論是在座的.
或者在網上現在收聽的.
或者是之後去聽到這個信息的人.
但願上主你都給我們.
真真實實地去成為一個.
應追悔改的人.
全情投入在與上帝交心這個層面.
求你的幫助帶領.
禱告奉靠主耶穌基督得勝明治教.
阿們.
\newpage



\section{約翰福音 21:1-14}
\label{sec:031948}
\textbf{主為你做早餐}~~~~禤凱榮牧師
\newline
\newline
連結: \href{http://www.church.com.hk/acms/content.asp?site=cdc\&op=show\&type=product\&code=031948\&layout=sermon}{\texttt{031948}} ~~~~ 語音日期: 2026-04-12
\newline
\newline
\hyperref[sec:031947]{\small{< < < PREV SERMON < < <}}
~
\hyperref[sec:index_chronic]{\small{[返順時目]}}
~
\hyperref[sec:031949]{\small{> > > NEXT SERMON > > >}}
\newline
\newline
約翰福音 21:1-14
\newline
\begin{longtable}{cl}
\hline
\hline
章節 & 經文 (和合本修訂版)\\
\hline
21:1 & \begin{tabularx}{0.7\textwidth}{X} 這些事以後，耶穌在提比哩亞海邊又向門徒顯現。他怎樣顯現記在下面。 \end{tabularx} \\ \\ \relax
21:2 & \begin{tabularx}{0.7\textwidth}{X} 西門‧彼得、叫低土馬的多馬、加利利的迦拿人拿但業、西庇太的兩個兒子，和另外兩個門徒，都在一起。 \end{tabularx} \\ \\ \relax
21:3 & \begin{tabularx}{0.7\textwidth}{X} 西門‧彼得對他們說：「我打魚去。」他們對他說：「我們也和你一起去。」他們就出去，上了船；那一夜並沒有打著甚麼。 \end{tabularx} \\ \\ \relax
21:4 & \begin{tabularx}{0.7\textwidth}{X} 天剛亮的時候，耶穌站在岸上，門徒卻不知道他是耶穌。 \end{tabularx} \\ \\ \relax
21:5 & \begin{tabularx}{0.7\textwidth}{X} 耶穌就對他們說：「孩子們！你們有吃的沒有？」他們回答他：「沒有。」 \end{tabularx} \\ \\ \relax
21:6 & \begin{tabularx}{0.7\textwidth}{X} 耶穌對他們說：「你們把網撒在船的右邊，就會得到。」於是他們撒下網去，竟拉不上來了，因為魚很多。 \end{tabularx} \\ \\ \relax
21:7 & \begin{tabularx}{0.7\textwidth}{X} 耶穌所愛的那門徒對彼得說：「是主！」那時西門‧彼得赤著身子，一聽見是主，就束上外衣，跳進海裡。 \end{tabularx} \\ \\ \relax
21:8 & \begin{tabularx}{0.7\textwidth}{X} 其餘的門徒因離岸不遠，約有二百肘，就坐著小船把那網魚拉過來。 \end{tabularx} \\ \\ \relax
21:9 & \begin{tabularx}{0.7\textwidth}{X} 他們上了岸，看見那裡有炭火，上面有魚和餅。 \end{tabularx} \\ \\ \relax
21:10 & \begin{tabularx}{0.7\textwidth}{X} 耶穌對他們說：「把剛才打的魚拿幾條來。」 \end{tabularx} \\ \\ \relax
21:11 & \begin{tabularx}{0.7\textwidth}{X} 西門‧彼得就上船，把網拉到岸上，網裡滿了大魚，共一百五十三條；雖然魚這樣多，網卻沒有破。 \end{tabularx} \\ \\ \relax
21:12 & \begin{tabularx}{0.7\textwidth}{X} 耶穌對他們說：「你們來吃早飯。」門徒中沒有一個敢問他：「你是誰？」因為他們知道他是主。 \end{tabularx} \\ \\ \relax
21:13 & \begin{tabularx}{0.7\textwidth}{X} 耶穌走過來，拿餅給他們，也照樣拿魚給他們。 \end{tabularx} \\ \\ \relax
21:14 & \begin{tabularx}{0.7\textwidth}{X} 耶穌從死人中復活後向門徒顯現，這是第三次。 \end{tabularx} \\ \\
[1ex]
\hline
\hline
\end{longtable}
$^{1}$主播我想大家去想一想一個畫面.
有一個人.
站在江邊.
夜色深沉.
四面楚歌,你猜這是什麼人呢.
一聽到四面楚歌,人人都知道.
項羽.
大家都知道,項羽是中國歷史上一個非常有名的將領,如果大家不熟悉歷史的話.
你看電視劇都看不少.
這個項羽就是活在.
秦朝末年楚漢相爭的時代.
他曾經在戰場上擊敗無數的.
極為強悍,所以他有一個稱號叫西楚霸王,大家都知道在這個.
秦朝滅亡之後.
項羽和劉邦.
大家知道兩個爭天下.
雖然一開始的時候,這個項羽佔優勢.
不過久而久之,後來局勢慢慢的逆轉.
在一個叫做.
蓋夏之戰那裡,項羽就被劉邦的軍隊包圍.
突然之間有一天晚上.
全來楚歌.
他們的士兵就以為我們的家鄉.
已經淪陷了.
因此而軍心.
潰敗.
這就是我們所說的四面楚歌這個成語的由來.
項羽.
後來突圍.
然後就去到一個烏江邊.
他的部下就勸項羽,不如你回江東.
重整.
旗鼓,我們再來過.
但是項羽就說.
我沒有顏面.
再見江東父老.
就是這樣,大家都知道結果是怎樣.
他就在這個烏江.
置敏而死.
就是這樣結束了他的一生.

$^{41}$等一下,我想你把畫面又轉到另外一個畫面.
來到一個的.
也是水邊.
也是一個人.
在清晨的海邊.
大家還可以聞到一些.
漁港的氣味.
這個人是誰.
就是西門彼得.
西門彼得.
好像項羽那樣.
帶著.
很大的失敗.
失敗.
失落的感覺.
不過這個不是一個普通的失敗.
因為聖經說得很清楚.
彼得在耶穌將要.
面對十字架的時候,他說,主啊,我願意為你死.
好幾個小時之後.
他竟然三次不認主.
聖經告訴我們,當.
雞叫的時候.
彼得突然想起他曾經.
跟他的主說過.
接著聖經說.
他就出去.
同酷.
兩個人.
彼得和項羽.
都是來到水邊,都帶著無比的失落感.
不過他們的結局.
相當的不一樣.
項羽在那個時候決定完結他的生命.
在他的生命畫上一個句號.
不過彼得呢.
彼得確實在那個時候決定.
在他人生.
揭去下一頁.
兩個很不同的中國.

$^{81}$究竟有什麼的不同,令到這個結局這麼不同呢.
如果今天我們所讀的經文告訴我們.
分別.
在於一個早餐.
我們現在看一看.
這個究竟是一個什麼的早餐.
竟然帶來人生這麼重大的分別.
經文開始的時候就說.
這些事以後.
是什麼事情呢.
事情就是.
耶穌復活.
以及兩次.
顯現.
給他的門徒看到.
這個復活.
耶穌的顯現.
激情過後.
你都要回復平靜的生活.
回復平靜的生活之後,門徒突然之間覺得.
我下一步應該怎麼走呢.
我下一件事要做的是什麼呢.
突然間在一個迷失的情況當中,這個反應相當的.
正常的.
因為人總是在人生高潮過後.
或者嚴重的毀滅性損失之後.
總是覺得迷失的.
為什麼迷失呢.
我未必可以再重複我之前的成功.
我也沒有能力去補償我所失去的東西.
突然之間.
何去何從呢.
我有什麼可以做呢.
我應該向哪個方向走呢.
突然之間.
迷失的當中.
這個都很常發生的就是.
當人在一個迷失的狀況裡面的時候會怎樣.
很容易.
走回以前.

$^{121}$舊有的模式.
熟悉的生活當中.
回復回以前,那就最好了,什麼都不用理會.
對彼得來說,他不單只是,大家要明白,他不單只是一個迷失的當中.
剛才他講過.
他當時的情緒,心情相當的複雜.
他夾雜著的是什麼.
羞辱.
後悔.
內疚.
來到這一刻的時候.
雖然他已經和.
耶穌已經和門徒顯現了兩次,但是聖經.
沒有提到在這兩次的顯現裡面,耶穌和彼得有直接的對話.
沒有.
帶著一個這麼複雜的情緒.
不知道要去哪裡.
又不知道要做些什麼.
回頭看.
回頭看周圍的人的時候.
彼得就說.
我去打魚.
為什麼那個時候彼得要去打魚.
沒有魚吃.
沒有錢花.
是不是這個原因呢.
我相信不是這樣的原因.
我相信原因是因為.
彼得去到那一刻的時候,突然間發覺生命停頓了.
他回頭看.
一大堆痛苦的回憶.
向前看的時候.
完全都是不可知.
他需要一些事情.
去.
分散他的注意力.
他也需要一些事情去.
消耗他的體力.
對於一個.
整生都是打魚的彼得來說.

$^{161}$當然.
最.
自然.
最容易.
的選擇.
就是回去打魚.
而他的同伴也說,好,我們跟你一起去打魚.
不過.
整天晚上.
勞碌了一天的晚上.
一條魚都打不到.
我想大家去.
進入.
當時的彼得.
去感受一下.
你猜當時的心情是怎麼樣的.
當然是失望了,弄了一整晚都沒有魚釣到.
還有呢.
很沮喪.
很挫敗.
原本是一個,我想大家去這樣想一想.
彼得原本是一個打魚的人.
但是.
在三年之前.
這個耶穌來到.
認識彼得之後,就叫彼得說,以後你不要再打魚.
你以後跟隨我.
以後得人.
如得魚.
好了.
彼得接受了這個呼喚.
跟隨著這個耶穌,跟了三年.
在整整的三年裡面.
他目擊看得到耶穌所行的神蹟.
他聆聽聽得到耶穌所講的真理要教訓.
不過.
在最後的時候.
他都好.
豪言壯氣的講了一番說話,就是.
他說我.

$^{201}$就是.
必須與你同死,也總不能.
不認你.
幾豪情壯志.
不過講完這句話沒多久之後,就三次不認主.
好了,現在主耶穌已經離世.
雖然他已經復活.
但是.
彼得再不能夠好像以前那樣跟隨著一個物質的耶穌.
不知道還有什麼可以做.
彼得選擇.
回到以前的生活.
去打魚.
因為這個.
一輩子都是做打魚的.
我們想像一下.
當時去打魚的地方,那個水域.
很熟悉.
那艘船是他自己的船.
下的網是他自己的漁網.
跟他一起出海打魚的,都是他慣常跟他一起出海打魚的漁伴.
但是卻是.
打了一萬魚都.
完全.
打不到魚.
我相信那時候的感情是怎樣的.
失望.
挫敗.
是不是啊.
但是這個情緒為什麼會這樣呢.
當然了,打了一萬魚.
我們叫做什麼.
白做.
White do.
什麼都打不到.
但是我相信的原因不是因為他們那晚打不到魚.
我相信.
彼得當時那種挫敗.
失落感.
並不是來自打不到魚,而是來自什麼.

$^{241}$而是來自.
因為彼得想回到以前那種的生活.
不過.
他卻是.
發覺.
以前的生活.
再不能夠像以前那樣提供他充實.
安全感.
自信.
似乎.
過往的生活.
在現在來說已經再行不通.
似乎他心靈是一個空了的位置.
他想用以前集以為常的生活來填補那個空間.
發覺.
不能夠再填補.
就在這個時候.
聖經告訴我們.
主耶穌.
站在岸邊.
也許他們來這個門徒當時有一段距離.
所以門徒當時認不出這個就是耶穌.
耶穌就說.
問他.
小子.
你們有.
吃的沒有.
有沒有打到魚啊.
我想問大家為什麼耶穌會問這個問題.
為什麼耶穌會問.
他們打不打到魚.
很簡單.
因為耶穌想知道他們打不打到魚.
是不是這麼簡單.
耶穌不是不知道.
耶穌真的不知道他們打不到魚嗎.
耶穌既是全知以全能.
耶穌必定知道他們那一晚是打不到魚.
既然知道為什麼還要問.
我相信原因很簡單.

$^{281}$耶穌想他們.
自己去.
檢視他們當時的生命實況.
整個.
彼得出去打魚的故事.
其實不是在說打魚.
也都不是在說.
要.
謀生.
而是一個更深層次的比喻.
去形容一個.
離開了耶穌基督的生命.
我們可以很努力工作,我們可以應用我們的經驗,我們的技巧,我們的知識.
但是沒有了神的帶領和賜福的時候.
一切的勞苦.
都可以完全沒有果效.
弟兄姊妹.
這是一個永遠不變的真理.
如果我們想得到神的寬恕的時候,我們首先要認罪.
如果我們想接受領受神的供應的時候.
我們首先要.
承認.
自己的缺乏.
如果我們需要神的帶領的時候.
我們必須要.
將我們生命.
降服在耶穌基督的手裡.
弟兄姊妹.
其實耶穌當日問的問題.
你有.
吃的會有.
其實都是問我們每一個人.
弟兄姊妹有沒有想過,我們如果有一個魚網的時候,我們這個魚網裡面有沒有魚.
你說.
人很多時候都要顧及面子,如果我們死撐.
免傷我們的自尊,有,我的魚網當然有魚,已經足夠我們用了.
如果你今天回答的時候.
你和耶穌基督的故事.
我想就停在那裡.
耶穌基督說得很清楚.

$^{321}$沒有病的人.
容不著醫生.
如果我們來到耶穌基督的面前,我們不承認有一個生命的需要的時候.
耶穌基督和我們.
不再需要有任何的關聯.
就是因為這個答案,所以第一代的.
教會就是這一班的信徒.
來建立.
我再看下去.
當我們.
很老實的回答.
沒有,我們的網是沒有魚的.
耶穌就給了一個很簡單的.
命令.
他說.
那你們就將網.
撒在船的右邊.
記不記得.
這樣的故事其實都發生過的.
如果我們看這個《路加福音》五章三到六節,那裡有一個很類似的故事.
那個故事就是說當時耶穌上了西門的船.
然後就出去了,就在船上教訓眾人.
說完之後,耶穌就叫他們出海打魚.
然後彼得又這樣跟他說.
主啊,我們整晚打魚都打不到,現在還打什麼魚啊,不過既然你這樣叫到的時候,好,我們就下網.
就這樣下網的時候.
圈著.
許多的魚.
兩個故事.
其實都是.
呼召著他的門徒信靠和信服,剛才我們所唱的那首歌.
其實那幫人是一個很有經驗的漁夫.
所以當耶穌這樣叫他們的時候,大家想一想.
他們大可以說.
主啊,你講道是聰明的.
打魚都是我們比較熟悉的,我們已經打了一整晚,沒有魚就是沒有魚了,而且我們也累得不得了.
不過他們沒有駐當退堂.
這一幫信徒就立刻跟著耶穌的.
指示來做,就這樣信服.
這一個的反應.

$^{361}$其實告訴我們.
關於.
信服.
一個很重要的事實.
我們可以在我們的事工.
我們的工作,我們的關係,我們屬靈的操練上面.
打了一整晚魚.
有沒有什麼感覺.
我們很努力在我們不同的生活範疇裡面.
打了一晚魚.
但仍然.
毫無收穫.
跟著耶穌說.
沒有你以前的成語拜.
現在下網.
這個這樣的命令.
大家覺不覺得.
荒謬呢,我做了一整晚,沒有就沒有了.
耶穌仍然是在那個.
特定的時刻裡面,叫他的門徒來信復的.
告訴我們什麼.
告訴我們其實信復是需要.
將我們的經驗.
我們的知識,我們之前的成功與失敗,完全的.
放在一旁.
信復就是這個意思.
大家.
再看下去.
有沒有想過耶穌為什麼一定要那幫門徒將網放在.
右邊.
左邊可以嗎.
有沒有分別呢.
為什麼一定要右邊呢.
我記得我很多很多年之前,不知道大家有沒有去過深海釣魚.
Deep Sea Fishing.
我去過一次.
我不知道現在是不是.
二三十年之前.
我們出海釣魚的時候,十幾個人.
好了,到了我們.

$^{401}$決定落網.
魚鉤的時候.
在船的一邊.
一邊落鉤,另外一邊完全沒有人在那裡釣,不准在那裡釣.
那邊要做什麼呢.
被人嘔.
我第一次去,在那邊嘔六個小時.
船家不會因為有一個人要嘔,就要乘船回去.
不會的,你自己在那邊嘔,不要阻礙那些人釣魚.
但我相信.
當時耶穌叫.
門徒在右邊落網,不是因為右邊要你嘔,多餘一點.
不是這個原因.
既然是這樣,那有什麼原因呢.
我相信沒有什麼特別的原因.
我相信.
正正是沒有什麼特別的原因.
就是信復的意思.
大家想一想.
真正的信復.
不需要我們完全理解.
不需要我們完全覺得合理.
不需要我們覺得.
是一個好的決定.
信復的意思就是信復發明令那個人,我們是否信得過他.
這叫信復.
是不是.
門徒當時沒有爭拗.
主紀人叫我落網在右邊.
不會有人走去說為什麼不在左邊.
沒有,他們立刻就在右邊.
來落網,為什麼.
因為耶穌基督叫他們這樣做.
大家落網.
接著.
就有這個神蹟出現.
大家看到這個神蹟了,門徒,聖經第七節所說,耶穌所愛的那個門徒.
突然之間認到.
是主.
即是說什麼.

$^{441}$即是說在整個的過程.
這個過程,那班門徒並不認得出這個人就是耶穌基督.
直至神蹟發生了之後.
那就奇怪了.
為什麼他們都會.
為什麼都會.
都會信復.
也告訴我們另外一個很重要的信復的.
功課.
很多時候我們.
去執行一些.
指令的時候,我就說.
特別是神的.
給過我們的使命.
很多時候我們就說.
你給我看得清清楚楚.
我就信服了,我覺得.
夠了.
安全了,感覺好了,我才信服.
做你要我做的事.
聖經不是這樣說的.
聖經就是說.
當我們信服.
當我們不明白.
當我們覺得不對的時候,我們仍然按照神的意思去做.
做了之後.
我們就看到.
神的心意.
我們就看到.
賜命令的神,是一個如何的神.
這個告訴我們,很多時候我們.
不知道,認不出是耶穌基督.
但當我們做了之後,當我們信服了之後.
我們就能夠.
認得到耶穌基督,信服其實某程度上就是幫我們打開屬靈的眼睛,看得到這個耶穌是哪一個.
也看得到我們作為耶穌門徒,是應該如何的做門徒.
所以信服可以說是,讓我們跟耶穌基督可以建立一個更深的.
經驗.
對耶穌基督更深的理解,以至我們跟耶穌基督可以建立一個.
更親,更強的關係.

$^{481}$頂子妹.
不知道大家有沒有.
接收到耶穌基督給過我們的命令.
將網.
下在你的船的右邊.
沒有,我都不是打魚的,我不是說這件事.
耶穌基督給過我們的命令.
可能是.
寬恕一些我們不想寬恕的人.
參與一些我們不想參與的事奉.
或者.
某些我們不想跟他們有任何關係的人來傳福音.
又或者面對一些我們覺得毫無能為力的生命的.
短板.
可能神要我們處理.
我們的反應會是怎樣.
哎呀,神啊,這個不是很.
不對的,沒什麼分別的,不用,我已經試過了,不行的.
會不會是這樣的反應?頂子妹.
希望我們大家都能夠從這個故事學到.
一個沒有魚的網.
和一個滿有魚的網.
的分別.
不是在於.
是否幸運.
不是在於我們有沒有經驗,有沒有知識.
而是在於我們有沒有信服.
屬靈的果子並不是一定.
我們.
更努力.
就可以有更好的屬靈果子.
我們要有屬靈果子.
我們要的是.
更大的信服.
我們可以整天晚上勞碌.
更加的努力.
或者是我們可以更加多的會有更好的策略.
不過.
我們的網仍然可以是空的.
當彼得和他的同伴.

$^{521}$聽耶穌的吩咐落網之後呢.
正如剛才路加福音的那個故事.
不過奇怪的就是在於.
很多個魚,聖經告訴我們.
很多個魚.
不過什麼.
不過那個網沒有難度.
拿不回來.
我們要明白這個神蹟並不是耶穌基督.
特地來炫耀祂捕魚的能力.
這個神蹟.
是讓我們提醒門徒.
祂是誰,耶穌基督是誰.
亦都提醒.
如果.
我們一心跟隨耶穌基督.
我們總是能夠分享到.
耶穌的封聖.
祂的供應.
以及當我們面對.
種種人生困難的時候.
耶穌基督總是有出路.
讓我們可以繼續前行.
好了,當門徒認到這個就是耶穌基督的時候.
聖經說他們就上岸了.
好了,當他們上岸之後.
就見到.
那裡有炭火,上面有魚,有餅.
簡單來說,耶穌基督在岸上已經為這幫門徒預備了早餐.
大家想像一下.
做了一整晚打魚的.
門徒.
你見到一個已經是.
現成準備好的早餐,你會怎樣?.
開心到不得了,馬上可以吃東西.
但想一想.
我相信當時的彼得.
他的心情.
不是單是開心.
我相信彼得當時的心情.

$^{561}$很複雜.
因為大家記得.
不是太久之前.
在.
一個大祭司的院子裡面.
當時耶穌被捉拿.
然後在這個大祭司的院子裡面.
來審問耶穌.
當時彼得暗暗跟著他們.
去到這個大祭司的院子裡面.
當時那天晚上很冷.
聖經說一幫人.
起了炭火,然後整幫人圍著炭火取暖.
大家有沒有留意這兩段經文有一個共通的地方.
都是炭火.
第一個.
耶穌在炭火裡面為門徒預備早餐.
另外一個炭火就是.
彼得在那個炭火裡面失腳三次.
三次不認主.
大家想像一下.
可能在.
彼得的腦海裡面.
每一次見到炭火就想起他的三次不認主.
在他的腦海裡面.
炭火.
就等同了他失敗的落印.
也許彼得在那個時候就想.
耶穌記得會不會就是用這個機會來跟我算帳呢.
會不會問我彼得你當日做什麼三次不認我呢.
也許在無數的晚上.
彼得一想起這個失敗的事情的時候.
他就跟自己說.
主啊.
你問我這個問題的時候,我會怎麼回答你呢.
我相信在他腦海裡面轉了很多次.
我應該怎麼回答.
我所愛的主.
為什麼在他最需要.
的時候.

$^{601}$我是.
不認他呢.
不過.
很有趣的就是.
在這個事情發生的時候.
耶穌基督.
沒有問到.
這個問題.
沒有問到彼得.
為什麼你不認我.
因為.
耶穌.
這次來到門徒的當中.
並不是為了拆毀.
而是為了重建.
主耶穌說.
將你剛才打的魚拿幾條來.
然後我們來吃早餐.
弟兄姊妹.
我想大家很清楚明白.
這個就是福音.
什麼叫福音.
原本人應該是.
面對神的懲罰.
但是這個例子告訴我們什麼.
神不單沒有懲罰.
神卻是為人預備了早餐.
邀請這幫人.
來一起參加這個早餐.
在那一天.
大祭司院子裡面.
耶穌彼得跌倒三次不認主.
就在這個炭火邊.
你可以說.
從那天開始.
這個彼得.
一直被.
羞辱.
愧疚.
後悔.

$^{641}$來侵蝕他的生命.
不過現在呢.
現在.
同樣是炭火.
耶穌來到.
炭火為他預備早餐.
讓他經驗到神的恩典.
耶穌基督的寬恕.
可以這樣說.
炭火不再是彼得.
提醒他.
這個失敗.
而是.
讓他.
經驗得到神的恩典和賜福.
記不記得耶穌曾經說過.
記不記得耶穌這樣說.
將幾條.
拿幾條你們打過的魚來.
大家怎麼理解耶穌基督這個要求.
為什麼耶穌基督要他們將.
那一次的生活拿幾條魚來呢.
哦.
也許耶穌.
不知道原來有這麼多人來吃早餐的.
預備得不夠.
所以叫門徒多拿幾條魚來才夠吃.
是不是這個原因呢.
也許如果是我預備那個早餐的就有可能.
犯了這麼基本的錯誤.
不過如果耶穌基督是一個全知全能的神.
耶穌基督.
必定不會犯一個.
這麼簡單的錯誤.
那究竟為什麼呢.
我相信原因很簡單.
耶穌.
邀請.
每一個跟隨祂的人.
來參與.

$^{681}$祂所.
開始的工作.
祂所為人所設立的事工.
祂要人.
積極的參與.
一同的.
來建立.
我們也許.
能夠付出的不多.
但當我們願意.
服侍的時候.
我們就能夠參與神.
豐盛的供應裡面.
以致我們可以在參與的過程裡面.
領受神.
更多的恩典.
彼得也跟著這樣做.
馬上就在船上拖著網魚.
留意.
這裡說.
那個網沒有因此而爛.
雖然有很多的魚.
很明顯.
根據如果是這樣說的時候.
當時的彼得和他們的同伴都知道.
魚網應該是破的.
不過沒有破.
很明顯.
拿到那麼多魚和網沒有破.
兩個都是神跡.
讓神.
提供這個供應.
這兩個神跡有沒有更深層的意義呢.
我相信除了耶穌基督是.
呼喚.
祂的門徒.
來建立教會以外.
用回耶穌基督自己的說話.
呼召這些人的時候是.
得人愚得魚.

$^{721}$不是打魚.
而是為了得著人的靈魂.
很明顯.
你叫一班打開魚的人.
現在來建立教會.
無論經驗.
無論知識能力.
似乎都超越了這班漁夫的能力.
但是.
聖經告訴我們.
這班雖然是做漁夫的人.
但是他們建立教會的事工上.
他們的收穫卡是滿滿的.
網絡沒有爛.
什麼意思.
網絡沒有爛.
大家記得在耶穌基督上十字架之前.
耶穌基督曾經跟彼得這樣說.
西門西門真是彼得.
撒旦要得著你們.
好西你們.
像西墨子一樣.
但我已經為你祈求.
叫你不至於失了信心.
這是很重要的一句話.
大家.
聽不聽到耶穌基督所說.
撒旦.
想要得著.
你們.
我想問大家這個你們是指什麼人.
彼得,不是.
你們,中蘇.
這裡的你們是指所有真心跟隨耶穌基督的人.
這麼大件事,是的.
所有真心跟隨耶穌基督,又願意事奉去服侍祂的人.
撒旦.
都想篩走我們.
撒旦.
都想處理處置我們.

$^{761}$不過.
耶穌基督.
跟所有的門徒說.
我們不需要懼怕.
因為耶穌基督會為我們祈求.
叫我們不至於失了信心.
而使徒保羅.
也都告訴我們.
現在耶穌基督.
已經從死裡復活.
現今.
在神的右邊.
替我們祈求.
這個我們就是每一個跟隨.
耶穌基督的人.
所以弟兄姊妹.
假若.
我們今日都是一個.
絕望.
或者是困難.
的境況的時候.
我們感覺好像一隻船沒有了舊.
在大風大浪裡面飄來飄去這樣的感覺.
我想大家.
去理解.
從這個故事告訴我們.
主耶穌.
在岸上面已經為我們預備好早餐.
記不記得.
一開始的時候我講項羽的故事.
項羽和彼得.
有一個這麼大的分別.
就在於這個早餐.
這個早餐不只是吃的早餐.
這個早餐代表著神的恩典和憐憫.
所以我們每一個人如果帶著困難.
帶著重擔的時候我們都可以來到耶穌基督的面前.
將我們的重擔交給祂.
信服.
聽祂的命令.

$^{801}$我們就會得著神的.
祝福.
我們就會得著神的.
報酬.
因為我們忠心.
和我們的努力.
而神的應許就是說.
我們的網.
必定不破.
因為神賜福.
我們每一個的事奉.
每一個服侍祂的心.
請我們同低頭祈禱.
耶穌基督我們多謝你.
這個故事再次告訴我們主耶穌.
你的救恩是充滿你的愛.
你的救恩是充滿你的憐憫.
你的救恩是充滿你的關愛.
\newpage



\section{出埃及記 12:11}
\label{sec:031949}
\textbf{聖經人聖經事:這是耶和華的逾越節}~~~~黃克勤牧師博士
\newline
\newline
連結: \href{http://www.church.com.hk/acms/content.asp?site=cdc\&op=show\&type=product\&code=031949\&layout=sermon}{\texttt{031949}} ~~~~ 語音日期: 2026-04-07
\newline
\newline
\hyperref[sec:031948]{\small{< < < PREV SERMON < < <}}
~
\hyperref[sec:index_chronic]{\small{[返順時目]}}
~
\hyperref[sec:031950]{\small{> > > NEXT SERMON > > >}}
\newline
\newline
出埃及記 12:11
\newline
\begin{longtable}{cl}
\hline
\hline
章節 & 經文 (和合本修訂版)\\
\hline
12:11 & \begin{tabularx}{0.7\textwidth}{X} 你們要這樣吃羔羊：腰間束帶，腳上穿鞋，手中拿杖，快快地吃。這是耶和華的逾越。 \end{tabularx} \\ \\
[1ex]
\hline
\hline
\end{longtable}
$^{1}$鄧兄姊妹大家好 繼續讀聖經.
聖經人 聖經事 生命樹 宣教網絡.
為你提供的.
今天的題目是.
《這是耶和華的逾越節》.
講逾越節都講了三四天了.
因為這個節期實在在我們基督徒的信仰裡面.
或者是聖經一直下來節期的傳統裡面.
是超級極度重要的.
希望大家留意多一點.
不過話是話 逾越節大家聽過很多次.
但是聖經讀過很多.
但是那些細節和十災等等那樣東西.
那些細節很多時候我們都沒有留意到的.
這次跟大家又去看看 講一下.
好了.
《這是耶和華的逾越節》.
這句話是誰講的.
只有兩個.
要麼是主耶和華講的.
要麼就是人講的.
這個就是耶和華的逾越節了.
要麼就是耶和華自己講.
我們的神自己講.
這個就是我的逾越節.
他會講.
這是耶和華的逾越節的.
《出一及記》十二章第一節.
耶和華在埃及地.
曉遇摩西亞倫說.
上次提了.
如果摩西沒有了.
他又說不行 我又不懂得說話.
那你有一個哥哥亞倫會說話的.
我就叫他來幫你.
於是就好了.
其實如果摩西很開心.
很爽的接了這個任務.
好吧 你叫我去我就去.
那亞倫就可能沒有份了.

$^{41}$不是沒有份都會出場.
不過就沒有變成摩西和亞倫.
經常都一樣.
他自己推得多.
反而被他的哥哥亞倫出來.
跟他一起.
如果這樣講.
講得不好就是.
在他的侍宮的承擔與榮耀上面.
是分了一些給亞倫的.
第二節你們要以本月為正月.
為一年之首.
綠色字就是我的解釋.
又說在埃及地和摩西亞倫說.
以色列人要以本月做正月.
那以往正月是什麼時候呢.
不是很特別說.
不過在這裡說.
現在要以這個月份做一月.
在這裡知道正月或叫做正月.
就是一月的意思.
是一年之首.
那這個月份是什麼月份呢.
在以色列人來說.
這個就叫做離散月.
離散月.
Li-San.
C是希伯來文的譯音.
C是讀S的.
San.
San都可以.
San就更加像San了.
Li-San.
離散月.
那Li-San是什麼意思呢.
就是他門飛行的意思.
他門飛行.
那為什麼會這樣呢.
不知道.
總之Li-San這個字.

$^{81}$就是他門飛行的意思.
就是用正月.
離散月做正月.
有原因的.
現在他們初一級了.
是新一年的開始.
新的計算的開始.
新的開始.
那這個圖可能不是很清楚.
不過就是網上都找到了.
你打以色列人的月曆月份.
就會有這個圖出來了.
正月就叫做離散月.
外面這個白色就是.
猶太人的或以色列人的月份.
離散月就是正月.
然後就是.
一路下去不同的月份.
那就有些名字.
有時正月都會出的.
比如十二月.
就是阿達月.
阿達月就有普爾節.
那西灣月就有七七節.
那接著就不是很多節了.
到催角節呢.
就是去到這個.
那些字都很模糊.
那你們看到就是離散月.
離散月有什麼呢.
是有愚月節和隨交節的.
那有這兩個節日都在這裡.
那愚月節完呢.
馬上就是隨交節的了.
而愚月節那一個星期日.
因為愚節是星期末的.
愚節晚餐就是星期四晚做的.
那在馬上的星期日就叫做初熟節了.
那也是隨交節的開始了.
那就不多說了.

$^{121}$就讓大家知道了.
那離散月.
他們的正月就是我們三月中到四月中.
就是現在西曆了.
好.
你們吩咐以色列傳媒眾說.
本月正月初十日.
各人要按著負加取羊羔.
取什麼?取羊羔.
就是羊仔小一點.
不是很大隻的那些.
那我想大家留意一下.
留意一下先.
第一點要說的呢.
就是用這個月.
做新的年份正月的開始.
那就是離散月.
是一年之首.
那有意思的.
就是現在出埃及了.
出埃及之後是一個全新的開始了.
大家明白了嗎?.
那就是我們西曆的三月中到四月中.
你們吩咐以色列傳媒眾說.
各人要按著負加取羊羔.
是按負加.
譬如一家人有五兄弟.
都未曾結婚的.
那就按負加.
那就是有媽媽又有妹.
於是一個家人有多少人?.
有父母有五兄弟.
那就是七個人了.
然後有三姊妹.
那就是十個人了.
還有兩個僕人.
那就是十二個人了.
那就按著負加需要的數目.
取一隻足夠大的羊羔.
一隻.

$^{161}$那就是不夠吃的.
那你就要拿一隻大一點的.
就不用太大隻.
要適合吃的.
適合尺寸.
好了,第二.
第二點呢.
第二點其實是說.
逾月節要求的有些什麼.
介紹一下逾節的東西.
大家讀下去就要小心一點了.
那第一,取一隻羊羔.
按照負加.
是一家一隻.
那家裡人不多,那怎麼辦?.
第一,可能吃少一點.
第二呢,其實都夠吃的了.
因為一隻羊很大隻.
很多肉的.
又有毛餃餅吃的.
那一家一隻羊.
不好,我有錢弄兩隻.
打媽不行.
一家一隻.
按著負加的人口的意思.
一家一隻是什麼意思呢?.
是希望全家都得救.
那這個呢.
綠色字是我解釋來的.
那這個綠色字呢.
是對這些《聖經》裡面.
和合本原來的文字的字呢.
是想做一個我的解釋.
或者有些是寓意的解釋.
就不是原文的解釋.
這個是寓意的解釋.
一家一隻羊.
一隻隻就夠了.
那是全家人吃的.
全家人都有份分吃這隻羊的.

$^{201}$那希望全家得救.
若是一家的人太少.
吃不了一隻羊羔.
那怎麼辦呢?.
本人就要和他隔壁的鄰舍共取一隻.
就是說很珍貴的.
那隻羊不要浪費.
盡量吃得完.
你們預備羊羔.
要按著人數和飯量來計算.
不要眼闊肚窄.
弄隻羊大一點.
誰知道吃不完.
浪費了.
不好的.
一家一隻.
我的贖領的意義就是希望全家得救.
那如果這個恩典有證.
那怎麼辦呢?.
我和隔壁來分享.
就按照我們的要求和飯量.
第五節.
要無殘疾一歲公羊羔.
或者你們是在綿羊裡面拿.
或者山羊裡面拿羊羔都可以的.
羊羔即是小的.
都可以.
那選一隻.
選好.
要留到本月十四日.
在黃昏的時候.
以色列全會眾擺羊羔宰了.
那你看到.
他要求都很多.
紅色字就是要求.
第一.
看上面的要求.
回去.
一家一隻.
按著附加.

$^{241}$希望是全家都得救.
一起享用.
如果多了就找隔壁一起.
按照飯量.
好了.
那一點了.
一家一隻.
第二.
無殘疾的.
即是說.
選一隻羊手御節.
御節十四天就要吃.
要做.
那怎麼辦呢?.
選一隻脫毛的.
選一隻很鈍的.
很水的.
病病的.
糟糕的.
這隻適合.
那不行.
要是無殘疾的.
無殘疾是象徵無罪.
無罪的.
無污穢的.
無罪.
而且是一歲的公羊羔.
一歲可以說是.
算做.
即不是隔遠遠.
走上來都有些力.
都算是成長了.
或者是在綿羊或者山羊都可以.
不要緊.
留到十四日御節的時候.
定了日子的.
御節什麼時候?.
我沒有空.
早一點守御節.
是不可以的.

$^{281}$一定是那一天.
是上帝指定的御節的日子.
十四天.
在黃昏.
黃昏.
早上起床見餓就宰了.
不行.
黃昏.
是一天的開始.
你看到.
除了是一年的開始.
都是一天的開始.
意識就是全會眾把羊羔宰了.
是.
各人都是在黃昏的時候.
十四天.
在黃昏的時候.
全部一齊宰了那隻羊.
一家一隻.
你看到.
是全會眾一齊.
是注重大家一起都守節.
第七.
國家要取點血.
在羊羔取點血.
塗在吃羊羔的房屋左右的門窗.
這個是讀「康」的.
不要讀「門框」.
門窗上和門眉上.
當夜就要吃羊羔的肉.
用火烤料.
與無餃餅和苦菜同吃.
國家要取點血.
拿羊羔的血.
塗在吃羊羔的房屋左右的.
門窗上和門眉上.
拿點血.
要血.
為甚麼要血呢.
若不流血罪不得赦免.

$^{321}$就塗在門窗上和門眉上.
要塗在那裡的意思.
就是血是遮蓋了神的憤怒.
遮蓋了.
意思是神就越過了.
遮蓋.
舊約裡面滯生的血.
是暫時遮蓋神的憤怒.
而不是清除了.
去到新約的時候.
才比較多用的字就是.
我們若認自己的罪.
神是信實的.
是公義的.
必要赦免我們的罪.
洗淨我們一切的不義.
洗了.
舊約比較上是遮蓋神.
暫時看到那些血.
不發露著.
塗在門窗上和門眉上.
屋的外面.
我覺得是象徵神的憤怒被遮蓋了.
被血遮蓋了.
當然就要吃羊羔的肉.
當你吃的那些肉.
是和你建立了生命的關係.
耶穌說.
這個是我的身體.
你吃吧.
如果你不吃.
就與我無分無關的了.
還要喝我的血.
有人這樣說.
要吃耶穌的肉.
喝耶穌的血.
這麼肉酸.
於是很多人就有誤會.
這個是不是食人教.
不是的.

$^{361}$這個是表示.
我吃耶穌的肉.
主耶穌說.
你吃我的肉.
喝我的血.
是建立了生命的聯繫與關係.
用火烤.
我覺得比較象徵十字架受苦.
火烤痛不痛.
哇.
痛到不得了.
簡直是縮手.
好熱.
火烤.
不是燙一燙.
火烤.
超級的痛苦.
與無教餅和苦菜同吃.
無教餅是象徵無罪.
同一個無殘疾.
是一樣意思的.
那隻羊要無殘疾.
就是說.
我們的主耶穌是無罪的.
不是我說.
是主耶穌是無罪的.
我們要吃無教餅.
就是我們要帶著一個清潔的心.
我都無罪的.
去到神面前.
那你起碼要維持自己的聖潔.
苦菜就是在艱難中生活.
在做奴隸.
但我都要守余節.
我多忙都要守.
辛苦艱難都要的.
不可吃生的斷.
不可吃水煮.
是一定要用火烤的.
是怎樣.

$^{401}$整隻帶著頭腿五臟.
用火烤去吃.
不可以生吃.
不可以水煮煮熟.
要用火烤.
是全隻頭腿五臟一起烤.
這個是象徵主的釘十字架.
是全個人受苦.
不是少少受苦.
是全個人受苦.
是全個人被火烤.
痛不痛.
當然痛.
不得了.
好苦.
好痛.
表示痛苦.
亦不可剩下一點留到早晨.
留到早晨就要用火燒了.
臨去烤.
是燒了.
不可以靜.
要盡量吃光.
一留到早晨就要燒了.
不拿走.
不會說先冷藏.
不會吊起它.
慢慢吃.
腳的羊腿在身上.
隨時帶著吃羊肉乾.
不行.
是要靜下來要燒光.
所以盡量適合份量.
不要太多.
只有一次機會.
這個魚字羔羊會不會每晚等你.
不會.
只有那一晚.
那一隻才是魚.
魚字的羔羊.

$^{441}$只有一次機會.
就是神的滅命天使走過.
或者神自己走過.
見到血.
遮蓋了你的罪.
魚羊不進去殺死你的長子.
只有一次機會.
不可以留到早晨.
明天就沒有了.
第11節吃羊羔的時候要怎樣.
腰間束帶.
腳上穿鞋.
手中拿著趕緊地吃.
吃的時候要怎樣呢.
第一.
腰間束帶.
即是說警醒等候.
衣服很闊.
外衣很闊.
內衣也闊.
長.
於是就要用毛巾.
皮帶.
繩子綁著.
束腰的意思.
就是做事.
或者準備行動.
隨時可以做.
先綁腰.
隨時可以動.
就不用等.
我又要先拿條帶.
拿條帶給我.
然後慢慢束起來.
不會束腰.
隨時上著準備.
快速準備好.
腳上穿鞋是表示.
隨時出發.
神什麼時候帶我們服侍他.

$^{481}$叫我們離開去事奉他.
我隨時出發穿鞋.
我說出發吧.
說吧.
哎呀.
我不見了一雙鞋.
只有一雙.
又要找.
不是吧.
隨時出發穿著鞋.
手中拿杖.
為何要拿杖呢.
是預備行遠路.
行很遠.
那你去隔壁屋借些鹽.
要不要拿支杖動過去呢.
當然不用.
近嘛.
在村頭村尾近嘛.
但如果你要去曠野很久.
那你就要拿支杖行.
你就借到力.
很多時候就有點用了.
所以就拿杖是預備行遠路.
趕緊地吃.
這個象徵著抓緊時間.
就要慢慢歎先.
開少少紅酒.
慢慢歎藥.
不是.
是趕緊地吃飽他.
這個表示抓緊時間.
因為警政等候.
隨時出發.
預備行遠路.
要吃飽.
那你趕緊地吃.
吃吃吃他.
這是耶和華.
第三點.

$^{521}$這是耶和華的逾越節.
接著神就說.
這是耶和華的逾越節.
這是我的逾越節.
我很看重的.
好了.
回去你看看.
第二點.
第二點是講到逾越節要怎樣的.
第一點.
一家一隻.
接著飯糧.
第二點.
沒有殘疾.
第三.
一歲的.
雄性的.
要留到十四日.
在黃昏的時候.
和全會眾一起吃.
已經有八樣東西了.
然後拿些血.
第九樣.
要圖在門榮門康上.
遮蓋我們的罪.
第十樣.
就要吃羊羔的肉.
建立關係.
十一.
就用火烤.
第十二.
十二家受苦.
無餃餅.
我們要無罪的.
十三.
苦菜.
用苦菜表現我們生活的艱難.
做奴隸.
就要苦菜紀念.
十四.

$^{561}$然後就帶著全隻吃.
十五.
第十六.
就是不可以剩下這裡的早晨.
十七點.
然後腰間束帶十八.
穿鞋十九.
臘醬二十.
港銀利吃廿一.
你數可以再數細些多些.
就數得沒那麼細.
二十或者廿一件事.
在這個餘節.
過節與羔羊.
都有二十一樣東西有關係的.
這是耶和華的餘月節.
如果要守餘節.
是要這樣的.
大家記得嗎.
你說我們的上帝主耶和華.
看不看重餘月節.
超級好看重.
就說了.
這是耶和華的餘節.
即是說.
餘月節是我的.
緊不緊要.
如果我們的主都看餘節這麼緊要.
我們要不要看餘節緊要些.
當然要啦.
以色列人說今天猶太人看餘節緊不緊要.
非常緊要.
緊要過年.
所以千萬不要想著三月中.
四月中去以色列.
哎呀我們受虧之復活節.
去看看先聖地.
逼人到不得了.
猶太人世界各地.
會回去以色列守節.

$^{601}$很多在以色列國的人.
會去揚撒冷守節.
水洩不通.
塞車到不肯.
過守餘節.
好看重的.
所以我們都要留意.
這是耶和華的餘節.
因為是祂救恩設立的象徵.
是一隻羊羔.
沒有罪沒有殘疾的羊羔.
被殺.
祂的血遮蓋我們的罪.
而祂是傳言獻上.
而且是受苦受死被敲.
而我們是全個以色列會眾.
在指定的時間一起吃的.
一起過節的.
而且過的時間.
是要記得警醒等候.
隨時預備著帶我們出發.
走新的路.
預備走遠路.
抓緊時間.
而且這一個月份.
就是正月.
新的正月.
就是說.
若有人在基督裡.
他就是新做的人.
舊事已過變成新的了.
這是一個新一頁的打開.
月節緊不緊要.
緊要.
這是耶和華的逾月節.
讓我們一起祈禱.
我們是感恩.
我們看重這個逾月節.
在我們還作罪人的時候.
在遠古上古的時代.

$^{641}$你就預備了逾月節.
讓我們開始看到.
這隻羊流的血.
為我們受苦.
遮蓋了我們的罪.
多謝你.
因為你為我們打開新的一頁.
這個是正月.
新的開始.
我們一起守節.
這是耶和華的逾月節.
讚美你.
多謝你的救贖.
讓我們都去堅心.
我們去守住這個救恩.
守住這個節期.
誰聽禱告.
奉教耶穌基督的名.
阿們.
\newpage



\section{使徒行傳 2:22-32}
\label{sec:031950}
\textbf{我們都是這事的見證人}~~~~趙承恩宣教師
\newline
\newline
連結: \href{http://www.church.com.hk/acms/content.asp?site=cdc\&op=show\&type=product\&code=031950\&layout=sermon}{\texttt{031950}} ~~~~ 語音日期: 2026-04-12
\newline
\newline
\hyperref[sec:031949]{\small{< < < PREV SERMON < < <}}
~
\hyperref[sec:index_chronic]{\small{[返順時目]}}
~
\hyperref[sec:031951]{\small{> > > NEXT SERMON > > >}}
\newline
\newline
使徒行傳 2:22-32
\newline
\begin{longtable}{cl}
\hline
\hline
章節 & 經文 (和合本修訂版)\\
\hline
2:22 & \begin{tabularx}{0.7\textwidth}{X} 「以色列人哪，你們要聽我這些話：拿撒勒人耶穌就是神以異能、奇事、神蹟向你們證明出來的人，這些事是神藉著他在你們中間施行，正如你們自己知道的。 \end{tabularx} \\ \\ \relax
2:23 & \begin{tabularx}{0.7\textwidth}{X} 他既按著神確定的旨意和預知被交與人，你們就藉著不法之人的手把他釘在十字架上，殺了。 \end{tabularx} \\ \\ \relax
2:24 & \begin{tabularx}{0.7\textwidth}{X} 神卻將死的痛苦解除，使他復活了，因為他原不能被死拘禁。 \end{tabularx} \\ \\ \relax
2:25 & \begin{tabularx}{0.7\textwidth}{X} 大衛指著他說：『我看見主常在我眼前，他在我右邊，使我不至於動搖。 \end{tabularx} \\ \\ \relax
2:26 & \begin{tabularx}{0.7\textwidth}{X} 所以我心裡歡喜，我的舌頭快樂，而且我的肉身要安居在指望中。 \end{tabularx} \\ \\ \relax
2:27 & \begin{tabularx}{0.7\textwidth}{X} 因你必不將我的靈魂撇在陰間，也不讓你的聖者見朽壞。 \end{tabularx} \\ \\ \relax
2:28 & \begin{tabularx}{0.7\textwidth}{X} 你已將生命的道路指示我，必使我在你面前充滿快樂。』 \end{tabularx} \\ \\ \relax
2:29 & \begin{tabularx}{0.7\textwidth}{X} 「諸位弟兄，先祖大衛的事，我可以坦然地對你們說：他死了，也埋葬了，而且他的墳墓直到今日還在我們這裡。 \end{tabularx} \\ \\ \relax
2:30 & \begin{tabularx}{0.7\textwidth}{X} 既然大衛是先知，他知道神曾向他起誓，要從他的後裔中立一位坐在他的寶座上。 \end{tabularx} \\ \\ \relax
2:31 & \begin{tabularx}{0.7\textwidth}{X} 他預先看見了，就講論基督的復活，說：『他不被撇在陰間；他的肉身也不見朽壞。』 \end{tabularx} \\ \\ \relax
2:32 & \begin{tabularx}{0.7\textwidth}{X} 這耶穌，神已經使他復活了，我們都是這事的見證人。 \end{tabularx} \\ \\
[1ex]
\hline
\hline
\end{longtable}
$^{1}$大英姐妹早晨.
我在講到開始的時候.
我就想和大家進行一個互動的遊戲.
這個遊戲是這樣的.
我會講一個詞語出來.
你就用一個手勢去表達那個詞語.
然後你看看身邊的人.
和你做的那個手勢是否一樣.
沒輸贏的放心.
明白玩法了吧.
我講一個詞語,你就做一個手勢出來.
C班都要做的.
後面主禮,雙禮都要做的.
首先第一個詞語是甚麼呢.
就是「打電話」.
你看看大家打電話的手勢是否一樣.
大家都是六字.
有些不是,有些是C的.
比較少的是C.
大家都是六字.
那收線呢.
收線就是這樣.
吸進去吧.
有沒有人按鍵?.
有的,有人按鍵.
有人按鍵收電話.
給心心.
有些人就是這樣給心心.
有些人就是這樣給.
甚麼都有.
幫人拍照.
有個框的.
有沒有人拍照是這樣拍?.
這樣拍嗎?.
明白我甚麼意思嗎?.
這樣拍照嗎?.
或者是這樣拍?.
有沒有這樣拍?.
有的,都可能有的.
好,「打字」.

$^{41}$「打字」.
全部都是用鍵盤的.
有沒有人是這樣打的?.
有人是這樣打的.
「聽音樂」.
就這樣按,對吧?.
「聽音樂」.
有人是這樣,我看到.
有人是戴耳機這樣聽的.
即是頭戴式耳機這樣聽的.
有人是這樣戴的.
好,我不知道兩位.
我看不到後面.
所以不知道.
好了.
大部分我們的「大拳姐妹」.
做的動作一樣.
不過個別是有點不同的.
其實為甚麼我會跟大家玩這個遊戲呢?.
這個遊戲本身我是上網看到的.
其實它最主要是甚麼呢?.
就是主要想講出.
就是「Y世代」和「Z世代」之間的分別.
那我大概那個年齡.
甚麼是「Y世代」甚麼是「Z世代」.
假設「Z世代」是甚麼呢?.
大概是2000年之後出世.
即是大概最大都是26歲.
這樣的年輕人.
然後「Y世代」.
那「Y世代」就是你當2000年前.
當然還有「X世代」.
還有「嬰兒潮」那些.
那就講講在做手勢之間的那個分別.
那我昨天就聽到一位姊妹去分享.
就說她.
即是姊妹大概.
即是大概某個年紀去到.
那接著她就說.
她聽到就會一個.

$^{81}$和年輕人談一對話.
那她就說.
她和那個年輕人說.
有沒有「OK」的手勢.
大家可不可以做「OK」的手勢.
「OK」是吧.
那原來呢.
她就說那個年輕人呢.
是覺得這個不是「OK」.
是「三」.
「三」是一二三個三.
是「三」.
那接著她就問.
那你們有沒有「OK」的手勢.
她說沒有的.
沒有「OK」的手勢.
那我們就.
嘩! 開心.
好像一個大發現那樣.
那我想問一下大家.
在打Message.
即是打這個信息的時候.
WhatsApp信息的時候.
大家有沒有用「句號」.
有沒有人用「句號」的.
有用「句號」的.
大英節目可不可以舉一舉手.
有用「句號」的.
對了.
那我就聽說過一個說法.
就是現在譬如Geno C.
即是Z世代的年輕人.
他們就說.
他們現在就不用「句號」的.
原因是什麼.
通常他見到一個「句號」的時候.
他就會覺得那個.
那個信息呢.
是一個很惡.
即是好像結束了對話的一個意思.

$^{121}$就覺得有個「句號」就很惡.
所以他們就多數不用「句號」.
那我就聽到.
o下! 我是用「句號」的.
我是有用「句號」的.
那怎麼辦呢.
好像會不會令到他有些誤解呢.
那剛剛就說了一些不同的例子.
就譬如用手提電話.
未普及的時候.
我們都是用固網電話打電話的.
那我們就會用錄音.
因為有個聽筒.
有個話機這樣.
那收線就有個錄音.
就連接起來這樣.
那但是因為手提電話普及了.
那我就有問.
我就有問過昨天.
周末競拜的兩個.
即是比較年輕一點的同學.
我就問他.
有沒有見過一個.
「攪撥式」的電話.
即是「嘟嘟嘟」.
即是一攪就「嘟嘟嘟」這樣.
我想我小時候都見過的.
都還見過的.
那他們說有見過.
那我就反問了他.
是不是在博物館那裡見.
那因為其實他們出生的年代.
這個已經很少用固網電話了.
或者再沒有見過這種的電話了.
而是用可能是iPhone.
即是用一個聽觸式屏幕的電話.
那所以譬如收線都未必是再按.
而是按鍵.
那不同時代出生的人.
他們都有不同的經歷.

$^{161}$那我們也都會形成了.
可能各自有不同的想法.
因為生活模式.
社會的狀況都有不同.
雖然我們會說.
不同年代的人可能就會有爭拗.
即是譬如世代之爭.
可能前幾年我們都還說世代之爭.
或者不同年代的人.
好像很難溝通這樣.
因為各自有不同.
就會有爭.
就會容易有爭議.
不過.
如果我們去看聖經的時候.
我們可能會發現一件事.
其實我們各有不同的時候.
可以是承傳了.
上帝的美意.
我們各有不同的時候.
可以是承傳上帝的美意.
今天的經課是.
《使徒行傳》二章十四節至到.
二十二至三十二節.
不過我想大家打開聖經.
因為今天我們會講得闊一點.
我們會講一講.
《使徒行傳》的二章一至三十二節.
今天就是我們唐用聖經的.
生日第一百八十日.
至到一百八十二日.
我相信大家都差不多了.
其實我們這段經文.
那個時間是一個在五秦節的時候.
在五秦節的時候.
五秦節是門徒領受聖靈之後.
彼得就向著不同地方的猶太人.
去講的道.
今天的經文是這樣.
其實是一篇講道.

$^{201}$我們可以說.
這個就是教會歷史裡面.
教會誕生之後的第一篇講道.
是教會的第一篇講道.
我們先了解一下一些前文背景.
在《使徒行傳》的二章開始的時候.
聖靈就臨到了一幫使徒裡面.
使徒就被聖靈充滿.
就講起別國的話題.
這個別國的話題.
就是向著當時來自天下各國.
不同虔誠的猶太人.
因為那些猶太人都聚集在耶路撒冷去過節.
其實天下各國的人.
我們可以想一下.
他本身是猶太人.
但是他出生的地方.
可能是一個埃及.
或者可能在隔離鄉隔離國家.
其實他們從小到大聽的.
未必是希伯來話.
也未必是亞蘭文.
就是當時猶太人最流行的語句.
可能他聽的是什麼呢.
可能他講的是埃及話.
或者是黎巴嫩話.
就有點像什麼呢.
譬如我們呢.
假設我們移民.
移民之後.
我們的第二代第三代.
可能他都未必會廣東話.
可能他懂的就是什麼呢.
英文.
但是他的樣子.
你都知道他的樣子其實是一個華人的樣子.
其實是一樣的情境.
這些人他們都聚集在耶路撒冷.
因為隔離國家相對比較近.
但是沒有飛機.

$^{241}$走路過去都是遠的.
去到耶路撒冷.
這些人聽見門徒就用.
他們各自的香談去講話.
而當時的門徒.
就被認到是伽利利人.
但是門徒就講起不同的話語來.
按照經文.
我有數過大概有多少個地方.
差不多有十三至十四個地方.
門徒就講的是十三至十四種的語言.
他們的舉動就令到其他人都不明白.
這是什麼意思.
所以有些人就說門徒是新酒灌滿了.
即是什麼呢.
醉了.
醉了的人會突然間講起外語是不出奇的.
因為我都見過有些這樣的人.
我曾經見過有醉了的人.
他本身是講廣東話的.
他突然間講大量的英文.
我們就知道他真的醉了.
不過呢.
我就未見過有人醉了可以講十四種語言.
沒有見過.
所以有人就覺得門徒的舉動是不可思議.
也有人覺得他們這麼奇怪.
為什麼會做這些事情.
但是有一樣東西就是.
他們的舉動是吸引了一些人去注意他們.
可能我們看這段經文的時候.
我們會把重點放在.
為什麼他們可以講到這麼多種語言的上面.
所以有些人就會認為方言.
即是聖靈賜的語言.
方言其實都是講的是不同的語言.
這是其中一個解法.
也有人認為這段經文是講的是.
反巴別塔的一種神學的理念.
我們可能會注意在這個位置上面.

$^{281}$但是如果我們去看這段經文的時候.
我們會發現一樣東西就是.
真正的重點是.
這班門徒是有一份吸引人去注意他的勇氣.
當其時耶穌就死了.
當然門徒也都經歷了他們復活.
和經歷了耶穌的升天.
但是其實門徒的身份.
仍然是一個死刑犯的跟隨者.
耶穌釘在十字架上面.
他們 耶穌的罪名是什麼.
猶太人的王.
即是一個政治犯.
所以其實他們說跟隨門徒.
即是跟隨耶穌的門徒.
都有可能是什麼.
就是反叛的餘黨.
他們可能是一個這樣的身份.
如果照正常的狀態底下.
他們就應該是躲起來.
保持低調.
不要被人發現到他們.
最起碼再過一段時間.
最起碼再過三五七年的時候.
等風聲沒那麼緊才出來.
照正常的情況是這樣.
不過聖靈一來到他們當中的時候.
門徒的行為就是吸引了很多人的注意.
所以聖靈一來到.
聖靈賜給仕途的第一個恩賜.
其實並不是那種講出不同語言的能力.
而是聖靈讓門徒能夠站出來.
去為耶穌基督作見證的勇氣.
問一問大家.
傳福音要不要勇氣?.
就算我們不是一個信耶穌會受逼迫的地方生活.
我們不是創啟地區.
我們要隱姓埋名的地方生活.
傳福音仍然需要很大的勇氣.
無論我們在學校.

$^{321}$無論我們在社會.
在公司.
在家庭.
或者在一個我們自己身處的地方崗位裡面.
當我們去和其他人分享耶穌是主.
或者我們很想和大家講一個福音的時候.
其實都很需要勇氣.
因為我們有時候都會知道.
可能會被人請食檸檬.
被人拒絕.
甚至乎可能會被人辱罵.
被人取笑.
等等等等.
而聖靈第一個給信耶穌的人的恩賜就是勇氣.
好了.
我想問問大家.
有見過這張單張嗎?.
收到這張單張嗎?.
先給後面.
給後面這張單張.
這張單張.
有吧?有吧?.
好了.
我想問問大家.
有沒有拿過這張單張給其他未信的人?.
有沒有拿過這張單張給未信的人?.
我這張單張.
我仍然放在家裡.
先講講我自己掙扎.
我仍然放在家裡.
我想邀請我媽媽來.
接著我太太不斷推我.
說「是時候要叫你父母信耶穌了」.
因為我家裡大部分人都未信耶穌.
所以她都叫我緊張一點.
加快腳步叫人信耶穌.
但是我很想叫我媽媽來.
5月9日和5月10日的.
福音主日聚會.
不過我仍然有掙扎.

$^{361}$我就放了這張單張在家裡.
我仍然在想一個方法.
怎樣邀請她.
怎樣開口去邀請她.
其實當我去看回這段經文的時候.
原來聖靈反而提醒了我.
聖靈祂自己就會給我們勇氣.
去面對和去開口.
為耶穌基督去做見證.
所以可能我們到現在.
我們都不知道怎樣開口.
或者我們不知道這張單張怎樣給人.
或者怎樣傳福音給其他人.
不知道怎樣做.
不過我覺得在聖經裡面.
正正提醒了我們.
原來聖靈給的第一個恩賜.
不是給個別的人的恩賜.
是給我們每一個基督徒的恩賜.
第一件事就是為耶穌基督.
作見證的勇氣.
所以其實我們可能缺乏勇氣.
可能我們會害怕面目.
可能我們會怕吃檸檬.
可能會怕被人笑.
不過原來聖靈所給的那份勇氣.
是足以令到當其時.
生命受到威脅的門徒願意站出來.
接著去到講出教會史裡面的第一篇講讀.
原來我們所缺乏的勇氣.
聖靈是會給的.
讓我們可以祈求聖靈給我們這份勇氣.
為了耶穌作見證.
剛剛我們繼續進入經文裡面.
彼得高聲地講出教會史裡面的第一篇講讀.
我們分析這篇講章的時候.
我們能夠分辨到.
其實聖靈都會繼續去做怎樣的工作.
首先我們看一看彼得的這篇講道.
我們看回這篇經文.

$^{401}$史壇行傳的二章十四節.
至到三十四節.
首先彼得在這篇講章裡面.
他就說叫猶太人和所有住在耶路撒冷的人.
都要注意要聽他說的話.
接著他就引用了.
約珥書的二章二十八至三十二節的經文.
去指出使徒之所以能夠講到不同的語言.
是因為好像約珥書所說.
在末後的日子.
上帝的靈是要囂貫在凡血肉之軀.
不論男女老少.
都會有上帝的靈在當中.
去解釋到他們為什麼會做一個.
這麼奇怪的舉動.
是因為聖靈的充滿.
令到他們可以講到不同的語言.
是因為聖靈的充滿.
他們可以站起來.
去宣講上帝的話.
接著彼得就講到.
耶穌就是上帝的異能歧視神蹟.
並且耶穌就住在人的中間.
亦都死在不法之人的手上.
但是上帝卻是死去復活了.
然後彼得就再引用.
詩篇的十六篇和一百一十篇.
去到陳明耶穌就是基督.
亦都是大衛的主.
彼得之所以引用這麼多的.
舊日的經文主義的原因.
是因為他宣講的聽眾.
都是猶太人.
以及歸依了猶太教的人.
然後就說這些熟悉希伯來聖經.
即是我們舊日聖經的人.
去到帶來一個新的詮釋.
讓他們能夠在熟悉的經卷裡面.
認識到耶穌就是那位基督.
而聖靈亦都使用了彼得.

$^{441}$對猶太教以及希伯來聖經的認識.
讓他可以在眾多的猶太教人的面前.
去到講述耶穌就是那位彌賽亞就是那位主.
不過我們試想像一下.
如果彼得用回同一番的說話.
同一番的講道.
去到二十一世紀的香港人.
去到講的時候.
你覺得那些聽眾會不會明白.
首先問一問大家回教會一定年日的大興智媒.
大家有沒有看過《約珥書》.
有沒有看過《約珥書》.
我相信大家知道《約珥書》是聖經的其中一卷.
我相信這件事.
好了 再問大家.
大家會不會讀到第16篇和第110篇.
詩篇第16篇和第110篇.
先講清楚.
我看過《約珥書》的.
不過我都背不到詩篇第16篇和第110篇.
背不到.
我們都要看回經文.
我昨天都有問過牧師.
牧師都說背不到第16篇.
我們就會說.
如果他本身是一個缺乏猶太教文化背景的人.
然後我們就將這一段的講章.
去到和21世紀的香港人去到宣講的時候.
我相信十居其九的香港人.
都不明白他在講什麼.
在這個時代.
在這個地方的裡面.
我們怎樣去傳福音呢.
我們怎樣去講述.
耶穌就是基督呢.
我們再看一件事情.
其實聖靈是使用了彼得自己對猶太教的熟悉.
猶太教的文化.
他自己作為一個猶太人他熟悉的程度.
然後他使用彼得去宣講這一篇的講道.

$^{481}$而上帝安排我們每一個弟兄姊妹.
在這個時代的裡面.
去到生活.
其實就是要讓我們在這個時代的裡面.
去到見證耶穌基督.
聖靈其實是使用了我們每一個人的年齡.
經驗.
時代背景.
社會地位.
能力.
興趣等等.
去到為耶穌做見證.
所以每一個時代其實都有不同做見證的方式.
可能在聖經的時代.
他們就正在經驗一種聖靈的充滿.
然後講起別國的話語的狀態.
而今時今日.
可能我們正在用的那種媒介.
去到為耶穌做見證.
一定是什麼呢.
一定是互聯網.
一定是社交媒體.
可能我們在Facebook.
可能我們在Instagram.
可能是Tray.
現在是Tray.
是不是Tray.
我都忘記了.
First.
忘記了.
我都沒有用.
真的沒有用到那個apps.
或者其實可能在六七十年代的時候.
其實教會也做了很多事情.
譬如什麼呢.
譬如以前是會派奶粉.
也有一些康樂的活動.
以前比較難去做camp的時候.
教會就去統籌做camp.
然後也都有很多人.

$^{521}$就在營會裡面去到順住.
也都有很多不同的人.
用不同的恩賜方法.
他的技能.
去到將福音帶給身邊的人.
譬如可能是拍片.
康樂的活動.
營會.
音樂報導等等.
但是另一邊廂.
聖靈也都讓到我們每一個信徒.
讓到我們可以有一個生命力.
去到見證耶穌就是基督.
譬如可能是我們待人接物.
我們是不是在見證耶穌基督.
其實我們在遇到困境患難的時候.
我們有沒有表達到我們.
其實我們不懼怕任何的死亡.
或者是威脅.
而是能夠靠著聖靈.
去有能力去面對.
可能其他人看到我們的時候.
看到我們的生命力.
看到我們信耶穌的人.
原來我們生命有不同的時候.
聖靈就在用我們的生命去表達.
原來耶穌是主.
那份力量是多麼大.
可能我們世代之間.
每個人都有各有不同.
我們的社會地位有不同.
我們的興趣有不同.
我們的能力有不同.
我們的背景有不同.
我們打電話的手勢可能有不同.
但是有一些事情.
有一些經驗.
其實正正因為我們的不同之處.
所以我們可以.
讓我們可以被聖靈使用.

$^{561}$去到接觸到不同的人.
讓我們可以和那些人.
能夠好像同聲同氣一樣.
去到表達得出我們明白他的需要.
而這一個需要的答案.
就是在耶穌身上.
聖靈幫助了我們.
讓我們能夠接觸到更加多不同的人.
讓我們可以用不同的方式.
去見證耶穌就是那位死而復活的主.
我們各有不同.
但是我們都是這件事的見證人.
這件事就是耶穌已經死而復活.
而跟隨他的人都會有那份復活的盼望.
讓我們一起祈禱.
因為在復活期的當中我們看得見.
由第一世紀的門徒.
他們就傳說耶穌就是那位復活的主.
而這份復活的盼望也都讓我們.
可以得著自由.
讓我們可以不懼怕死亡.
父神我們求你的聖靈都幫助我們.
讓我們都可以有那份.
傳說福音的勇氣.
讓我們可以有那份力量.
去向不同的人.
不同的群體.
去宣講耶穌是主.
也都求你幫助我們.
可能我們缺乏勇氣的時候.
求你加勇氣給我們.
當我們覺得自己好像沒有準備的時候.
讓我們可以有份信心去依靠你.
父神對我們又為到.
你賜給我們有不同的背景.
不同的社會地位.
不同的身份.
不同的興趣.
不同的年齡.
出生在不同的時代而感恩.

$^{601}$因為這些都是有你安排的美意.
就求你的聖靈使用我們.
讓我們各有不同.
大興姐妹之間雖然各有不同.
不過我們在各有不同的裡面.
都去見證耶穌就是那位我們的主.
就是基督.
就是那位死在十字架上.
但也是戰勝死亡的主.
讓更多的人可以因著我們.
去跟隨耶穌基督.
去得著這份永遠的盼望.
我們的禱告教託的一幕.
是靠我們救主耶穌基督的名.
好嗎?好嗎?.
\newpage



\section{約翰福音 20:19-31}
\label{sec:031951}
\textbf{我信身體復活}~~~~李均熊牧師博士
\newline
\newline
連結: \href{http://www.church.com.hk/acms/content.asp?site=cdc\&op=show\&type=product\&code=031951\&layout=sermon}{\texttt{031951}} ~~~~ 語音日期: 2026-04-12
\newline
\newline
\hyperref[sec:031950]{\small{< < < PREV SERMON < < <}}
~
\hyperref[sec:index_chronic]{\small{[返順時目]}}
~
\hyperref[sec:031952]{\small{> > > NEXT SERMON > > >}}
\newline
\newline
約翰福音 20:19-31
\newline
\begin{longtable}{cl}
\hline
\hline
章節 & 經文 (和合本修訂版)\\
\hline
20:19 & \begin{tabularx}{0.7\textwidth}{X} 那日（就是七日的第一日）晚上，門徒因怕猶太人，所在的地方門都關了。耶穌來，站在當中，對他們說：「願你們平安！」 \end{tabularx} \\ \\ \relax
20:20 & \begin{tabularx}{0.7\textwidth}{X} 說了這話，他把手和肋旁給他們看。門徒一看見主就喜樂了。 \end{tabularx} \\ \\ \relax
20:21 & \begin{tabularx}{0.7\textwidth}{X} 於是耶穌又對他們說：「願你們平安！父怎樣差遣了我，我也照樣差遣你們。」 \end{tabularx} \\ \\ \relax
20:22 & \begin{tabularx}{0.7\textwidth}{X} 說了這話，他向他們吹一口氣，說：「領受聖靈吧！ \end{tabularx} \\ \\ \relax
20:23 & \begin{tabularx}{0.7\textwidth}{X} 你們赦免誰的罪，誰的罪就得赦免；你們不赦免誰的罪，誰的罪就不得赦免。」 \end{tabularx} \\ \\ \relax
20:24 & \begin{tabularx}{0.7\textwidth}{X} 那十二使徒中，有個叫低土馬的多馬，耶穌來的時候，他沒有和他們在一起。 \end{tabularx} \\ \\ \relax
20:25 & \begin{tabularx}{0.7\textwidth}{X} 其他的門徒就對他說：「我們已經看見主了。」多馬卻對他們說：「除非我看見他手上的釘痕，用我的指頭探入那釘痕，用我的手探入他的肋旁，我絕不信。」 \end{tabularx} \\ \\ \relax
20:26 & \begin{tabularx}{0.7\textwidth}{X} 過了八日，門徒又在屋裡，多馬也和他們在一起。門都關了，耶穌來，站在當中，說：「願你們平安！」 \end{tabularx} \\ \\ \relax
20:27 & \begin{tabularx}{0.7\textwidth}{X} 然後他對多馬說：「把你的指頭伸到這裡來，看看我的手；把你的手伸過來，探入我的肋旁。不要疑惑，總要信！」 \end{tabularx} \\ \\ \relax
20:28 & \begin{tabularx}{0.7\textwidth}{X} 多馬回答，對他說：「我的主！我的神！」 \end{tabularx} \\ \\ \relax
20:29 & \begin{tabularx}{0.7\textwidth}{X} 耶穌對他說：「你因為看見了我才信嗎？那沒有看見卻信的有福了。」 \end{tabularx} \\ \\ \relax
20:30 & \begin{tabularx}{0.7\textwidth}{X} 耶穌在他門徒面前另外行了許多神蹟，沒有記錄在這書上。 \end{tabularx} \\ \\ \relax
20:31 & \begin{tabularx}{0.7\textwidth}{X} 但記載這些事是要使你們信耶穌是基督，是神的兒子，並且使你們信他，好因著他的名得生命。 \end{tabularx} \\ \\
[1ex]
\hline
\hline
\end{longtable}
$^{1}$丹瑩姊妹平安.
又可以在這裡跟大家分享神的說話.
或者我們首先低頭祈禱.
祈求上帝的幫助.
我們在天的父,求主你親臨在我們中間幫助我們.
叫我們心思所想的都被你祝福.
用我口中所說的言語,心裡的意念.
說出來的時候就成為你的道.
去幫助丹瑩姊妹.
我們這樣禱告,加上是奉主耶穌明子祈求.
究竟基督徒是信什麼呢?.
我們每個星期都會念《信經》.
《信經》裡面有一句話說.
我信身體復活.
但其實大家如果誠實一點去問自己的話.
心裡是不是真的相信這句話呢?.
其實我懷疑很多基督徒不相信身體復活.
其實通常基督徒相信什麼呢?.
我懷疑一般的基督徒其實是相信另一樣的東西.
那樣東西叫做靈魂不滅,死後就上天堂.
因為很多時候我們以為基督教的盼望就是這樣.
人死了,身體就留在地上.
最後腐化了,什麼都沒有了.
而真正的我,那個靈魂就去到天堂.
離開這個會老,會痛,會壞的身體.
這套想法聽起來是不是很熟悉呢?.
我想不單止基督徒或者很多非基督徒都這樣想.
基督教其實是相信這些東西的.
你給我傳福音就叫我靈魂上天堂.
但是這個想法或者這套語言不是《聖經》的語言.
《聖經》不是這樣說的.
這個其實是一個希臘的思想.
遠自荷馬的史詩到後來柏拉圖很詳細的解釋.
身體是次等的,暫時的.
靈魂才是不老的,不滅的.
靈魂才值得去拯救.
所以其實,弟兄姊妹,這套將靈魂和身體徹底分開的想法.
不是從《聖經》而來的.
今天我們詩篇經課裡面,詩篇第十六篇,第十節是這麼說的.
大家可以打開再看一次.

$^{41}$因為你必不將我的靈魂撇在陰間.
也不叫你的聖者見扭壞.
我們讀靈魂的這個字,在希伯來文叫nevish.
我們很容易以為一個可以離開身體.
單獨存在的那個部分,我們叫做靈魂.
但在舊約裡面,nevish這個字,靈魂這個字.
不應該稱它為靈魂.
nevish這個字是指一個不會脫離肉身的東西.
它是和身體緊緊結合的東西.
甚至可能指整個活著的人.
nevish可以是很多東西.
但一定不是一些外在於身體的,叫做靈魂的東西.
所以詩篇第十六章,第十節那裡.
因為你必不將我的靈魂撇在陰間.
也不叫你的聖者見扭壞.
身體的扭壞和靈魂其實是同一樣東西.
不是兩樣東西.
在舊約的信仰裡面,人不是有一個靈魂.
然後有另一個東西叫做身體.
人本身就是一個有身體的生命.
死亡不是靈魂離開身體.
而是整個人落入死亡裡面.
而神的拯救也不只是拯救一部分.
那部分叫做靈魂.
而整個人是不被死亡拘禁.
是整個人.
後來這個nevish翻譯成為希臘文的時候.
就用了suke這個字,希臘文suke就是psychology的字根.
psychie.
慢慢這種希臘的意願的思想就潛入了我們基督教的思想裡面.
將靈魂想成為逃離肉身不老的那個本體.
你看看彼得前書,彼得在《使徒行傳》.
今天我們敘事經課那裡.
他引用了詩篇的十六篇.
他要證明的不是靈魂不死.
而是神是要洗埋葬的耶穌從死裡面復活.
從一開始就不是一個靈魂上天堂的信仰.
基督教信仰我們所盼望的是一個身體從死亡中被拯救的盼望.
但不知不覺連教會都受這種西方的希臘的思想影響.
我們口裡面就說身體復活.

$^{81}$心裡面就期望靈魂上天堂.
最後我們可以看看今天的福音經課.
我們聽到或者我們看到我們熟悉的多瑪.
她要求見到的是什麼?.
她要求的是身體復活的證據.
除非我看見她手上的釘痕.
用指頭探入那個釘痕.
用手探入她的肋膀我總不信.
或者大家可以真心想一想.
不用舉手心裡面想一想.
如果我今天問你死了之後你認為會發生什麼事呢?.
可能你心裡面都會答我希望靈魂上天堂.
但我們唸信經的時候唸什麼?.
我信身體復活.
所以我們所認識的多瑪.
我們經常覺得她是一個多疑.
要看到要摸到才相信的多瑪.
但我們沒有留意到多瑪的信心.
比我們大部分的信徒更加準確.
因為她真的相信身體復活這件事.
所以你會見到後來耶穌不單沒有怪責她.
沒有說她多疑.
反而更加滿足她的要求.
約翰福音接著說什麼?.
耶穌再次留在門徒中間.
沒有對多瑪說一句雜誌說話.
反而說伸過你的指頭來摸我的手.
伸出你的手來探入我的肋膀.
大家留意耶穌沒有說.
你為什麼要這樣要求?.
你為什麼這麼質疑這件事?.
你需要贖靈一點?.
沒有,耶穌叫多瑪的.
是一個帶著釘痕,帶著傷口的身體.
向她顯現,叫她摸.
換句話說,耶穌沒有否定多瑪對身體復活的期待.
反而用自己的身體去肯定.
這就是復活的內容.
那一刻,多瑪所看到的不是靈魂的顯現.
不是天堂的異象.

$^{121}$是一個釘了十字架的身體.
站在她面前.
弟兄姊妹.
這就是我們教會一直在說的.
所謂那句很奇怪的信經.
我信身體復活.
這句話不是多餘的.
也不是象徵性的.
這其實是我們對世界的一個宣告.
神沒有放棄身體.
祂沒有放棄這個受傷的世界.
復活的耶穌仍然有傷痕.
這意味著什麼?.
受苦的經歷.
你在世界身體所受的苦楚.
那些傷痕.
那些經歷.
不是沒有存在過.
不是像粉筆字一樣抹去.
而是復活的時候賦予新的意義.
因為死亡是沒有最後的發言權.
死亡不能告訴我們身體所經歷的結論是什麼.
所以基督教的盼望是什麼?.
不是靈魂逃離身體.
而是神親自進入死亡那裡.
將死亡裡面的身體.
用新生命的東西帶回來.
身體有一個新的生命.
所以我們今天念信經的時候.
我們不只是念一句傳統的說話.
而是宣告一個很大的盼望.
不是說我將來要去哪裡.
而是神要將我跟這個世界一起去更新.
一起去創造.
這樣弟兄姊妹.
我信身體復活.
在今天又有什麼意思呢?.
我信身體復活這句說話.
背後其實還有很多神學的思考和討論.
就算我剛才說舊約那些什麼流程.

$^{161}$也要講一段時間.
不過通常我都會叫大家進神學院讀課程.
不過如果你有興趣.
繼續去想深思這個課題的話.
我可以建議你讀一本書.
就是賴特Anti-Right所寫的.
Surprised by Hope 或者已經在台灣的校園翻譯了.
成為天堂有什麼好期待.
再施復活忠活與教會大使命.
這本小書很值得大家去看看.
大家更加明白什麼叫我信身體復活這個意思.
但是我信身體復活對我們今天信徒有什麼意思呢?.
首先這句話提醒我們.
我們今天的身體不是多餘的.
不是用原積氣.
等它老了,壞了,病了,死了.
就丟掉它不管的.
這個會疲倦,會衰老.
會生病的身體.
就是那個神所愛.
要拯救將來去更生的身體.
你不要問我神怎樣更生這個身體.
因為我真的不知道.
我有時候想過.
我可不可以有些東西留下.
有些東西是不愛的呢?.
譬如有一顆痣我不太喜歡.
可不可以不愛的?.
但是這個就是上帝要拯救,要更生的身體.
它怎樣做我不知道.
有很多神學家想了很多的理由.
或者很多的想像.
但是其實我們身體不斷地在更生自己.
我今早就Google過問.
我們每隔多少天身體就會更生呢?.
人體細胞會不斷地自我更生.
一般在六個月.
就是一百八十天內.
我們的身體就會換掉九十八個百分比的細胞組織.
而不同的器官的細胞更生.

$^{201}$週期差異很大.
例如表皮兩個星期就會更新一次.
胃黏膜五天就會更新一次.
紅血球的壽命大約一百二十天.
心臟細胞就可以長達二十年.
我們其實不斷地在更生.
上帝的創造不是停留在一個狀態.
上帝其實不斷地再做,再創造.
我們相不相信上帝的再創造.
在我們死亡之後.
會使我們有一個新的生命.
帶著一個新的身體呢?.
所以身體不是多餘的.
當我們面對病痛,軟弱,年老.
面對死亡的時候.
我們不可以說不要緊.
反正這個身體也沒什麼用.
但是我們敢不敢去相信.
神仍然在乎,神仍然記得.
神仍然會拯救.
就是這個會衰敗的身體.
我們的傷痕.
我們的疾病.
我們的痛楚.
上帝是care的.
這件事情甚至會carry下去的.
甚至你可以說有些人認為.
他的痛苦經歷就是造就了他這個人.
死了之後那些東西就煙消雲散.
不是的.
我相信身體復活.
我們的生命是會延續下去的.
所以今天的生命是重要的.
是要care的.
上帝是care的.
it matters.
第二.
剛才說了第一.
不知道大家記不記得.
第二.

$^{241}$就是我相信身體復活.
就改變了我們今天對人的身體的看法.
既然神沒有放棄人的身體.
所以今天所有的貧窮,受傷.
被忽略,被壓迫的身體.
並不是無關重要的.
我們的信仰不是關乎.
靈魂將來去哪裡.
而是關乎今天我們怎樣看.
今天我們眼前的生命.
所以那些靈魂身體分開.
將肉體的需要和靈魂的需要分割.
那些教導,那些思想.
大家要小心去看.
這些是屬世俗的.
這些是屬肉體的.
這些是屬靈魂的.
沒有這樣的事情.
我們是一個人.
上帝是關心我們所有的需要.
無論是所謂身體的.
還是那些不屬於身體的.
最後.
我相信身體復活.
是讓我們這個仍然充滿著.
死亡和破碎的世界裡面.
我們可以活得不一樣.
因為我們知道.
這個世界不是注定要被丟棄.
其實它在等待將來神的更生.
而且這個將來.
不是一個模糊的,不確定的將來.
因為今天在我們的書信經課裡面.
彼得前書的說話就說.
藉著耶穌基督從死裡復活.
神已經重生了我們.
給我們一個活潑的盼望.
可以得著不能扭壞.
不能玷污.
不能衰殘.

$^{281}$為我們存留在天上的基業.
大家留意.
存留在天上的基業.
不是讓我們離開地上的肉身.
才可以到天上拿到那些基業.
存留在天上的基業.
好像媽媽說我留了一些湯在雪櫃裡.
你不需要住在雪櫃裡.
才能喝湯.
只是將最好的.
不會扭壞,不會玷污的.
為我們存留著.
今天的我.
種種和我身體有關的事情.
無論是我世上的學習.
我肉身的經歷.
我五官所領受的經驗.
我的成長.
甚至因為肉身的緣故.
所有的關係.
多謝剛才白師兄的清明頌.
我們的父母.
我們的妻子朋友.
是因為我們有肉身.
才會建立這種關係的.
如果沒有肉身.
我們只有天父.
我們弟兄姊妹.
只是屬靈的雙交.
不是,我們是肉身的交往.
我請你吃飯.
你請我吃飯.
是不是?.
是因為有肉身.
因為在社會裡.
我們有肉身.
我們要經營很多不同的關係.
這些關係不會因為我們死了而消失.
是不是?.
而正正就是.

$^{321}$這些和肉身種種關連的事情.
現在就在神的手中被保守.
擺在天上.
讓我們將來去領受.
不能扭壞.
不能玷污.
不能摧殘.
上主的基業.
它在神的手裡面被保守.
不是靠我們去捉緊的.
所以彼得前書那裡繼續說.
我們那些因蒙神能力保守的人.
必能得著那道末世要顯現的救恩.
然後彼得前書.
就再鼓勵當時的弟兄姊妹.
為什麼我們今天在試煉當中.
仍然可以站立.
在破碎裡面仍然可以盼望.
因為我們有個將來.
這個將來不是一個問號.
而是一個因為主基督的復活.
他作為一個初熟的果子.
來確立的一個應許.
我們望見耶穌身體的復活.
我們就知道.
將來我們的應許.
就寄存在上帝的手中.
是靠上帝的能力托住.
不會被死亡奪走.
不會被時間消亡.
甚至更新了以後.
就不能夠毀壞.
不能夠玷污.
不會再衰殘.
所以今天.
當我們一會兒再讀《使徒信經》.
那一句很傳統的文字.
傳統的說話.
我們不只是簡單的讀一段文字.
而是我們用信心去宣告.

$^{361}$我信身體復活.
我信神沒有放棄我.
我信神沒有放棄這個肉體物質的世界.
我信.
今天一切會毀壞的不是結局.
一切破碎中的不是終結.
因為在復活的主裡面.
一切都要被復原.
被更新都成為新的了.
我們一起祈禱吧.
最要求你去聖靈在我們心裡面.
使我們有確懼.
當我們再遇到復活的主的時候.
叫我們明白到.
祂的身體就是身體復活的見證.
就是身體復活的證據.
以致我們可以繼續去盼望.
我們今天肉身活著.
仍然是可以被得救.
被更新.
所有我們今天所經驗的神你都會重視.
甚至乎在我們將來死亡之後.
仍然都可以在你裡面得到更新.
所以叫我們有這份信心幫助我們.
我們這樣禱告是奉諸耶穌明志祈求.
阿們.
\newpage



\section{馬可福音 16:1-8}
\label{sec:031952}
\textbf{石頭滾開的奇蹟}~~~~黃順成牧師
\newline
\newline
連結: \href{http://www.church.com.hk/acms/content.asp?site=cdc\&op=show\&type=product\&code=031952\&layout=sermon}{\texttt{031952}} ~~~~ 語音日期: 2026-04-12
\newline
\newline
\hyperref[sec:031951]{\small{< < < PREV SERMON < < <}}
~
\hyperref[sec:index_chronic]{\small{[返順時目]}}
~
\hyperref[sec:031953]{\small{> > > NEXT SERMON > > >}}
\newline
\newline
馬可福音 16:1-8
\newline
\begin{longtable}{cl}
\hline
\hline
章節 & 經文 (和合本修訂版)\\
\hline
16:1 & \begin{tabularx}{0.7\textwidth}{X} 過了安息日，抹大拉的馬利亞、雅各的母親馬利亞，和撒羅米，買了香料，要去膏耶穌的身體。 \end{tabularx} \\ \\ \relax
16:2 & \begin{tabularx}{0.7\textwidth}{X} 七日的第一日清早，太陽出來後，她們來到墳墓那裡， \end{tabularx} \\ \\ \relax
16:3 & \begin{tabularx}{0.7\textwidth}{X} 彼此說：「誰要替我們把石頭從墓門滾開呢？」 \end{tabularx} \\ \\ \relax
16:4 & \begin{tabularx}{0.7\textwidth}{X} 她們抬頭一看，看見石頭已經滾開了，原來那石頭很大。 \end{tabularx} \\ \\ \relax
16:5 & \begin{tabularx}{0.7\textwidth}{X} 她們進了墳墓，看見一個年輕人坐在右邊，穿著白袍，就很驚奇。 \end{tabularx} \\ \\ \relax
16:6 & \begin{tabularx}{0.7\textwidth}{X} 那年輕人對她們說：「不要驚慌！你們尋找那釘十字架的拿撒勒人耶穌，他已經復活了，不在這裡。來看安放他的地方。 \end{tabularx} \\ \\ \relax
16:7 & \begin{tabularx}{0.7\textwidth}{X} 你們去，對他的門徒和彼得說：『他要比你們先到加利利去，在那裡你們會看見他，正如他從前所告訴你們的。』」 \end{tabularx} \\ \\ \relax
16:8 & \begin{tabularx}{0.7\textwidth}{X} 於是她們出來，從墳墓那裡逃走，又發抖又驚訝，甚麼也沒有告訴人，因為她們害怕。 \end{tabularx} \\ \\
[1ex]
\hline
\hline
\end{longtable}
$^{1}$弟兄姊妹 主佬平安.
有幸我們和阿哥牧師認識.
作為左鄰右里有個機會過來一同分享見證上帝的道.
這是一個美好的事情.
所以今天和大家分享.
我們先開始祈禱.
讓我們每一位弟兄姊妹一同聚集在這個地方裡面.
讓我們一同思想耶穌基督你所成就的事情.
求主帶領幫助 聖靈讓我們開啟.
讓我們見到你的工作.
祈禱奉主耶穌基督聖名祈求 阿彌陀佛.
弟兄姊妹其實上個禮拜復活節.
復活節基本上講的內容.
其實你大致上你都應該是想到.
其實它是我們基督教一個很重要的事件.
它是整個基督教的確立.
最對基督教的發展.
以致我們所有經歷的事情來講.
都是必然是重要的.
所以我亦歸另一個感覺.
如果你今天回教會久了.
你會留意到.
而這段經文其實多數都在復活節前後.
或者甚至復活節中毒的.
為什麼今天又會再講呢.
很自然地因為我覺得這件事很重要.
重要到我們不厭其煩再詳細的.
我們思想這個都是必須的.
否則你就會成為一個無根的信徒.
耶穌的死 耶穌的復活.
對我們整個的基督教來說.
是一個必須不可迴避.
甚至我們常常思考的事情.
所以我決定再在這段經文裡面游歷.
而不遵從比如說很多時候教會年曆.
我們開始時會講耶穌的升天.
門徒的聚集 差遣 五巡節等等.
我反而希望從弟兄姊妹.
再次從另一個角度.
再次經歷復活裡面帶來那種特別的意義.

$^{41}$不過我的焦點放在那個石頭上.
放在那個石頭上.
弟兄姊妹 如果是這樣講的時候.
我再次去講復活這件事.
對整個人類的歷史來說.
其實是一個很決定性的轉折點.
它將人類的歷史分了兩部分.
分為之前和之後.
而從我們個人的生命層面來說.
它更加是我們每個人.
從死亡走向光明.
從絕望走向盼望.
從黑暗走向光明的關鍵時刻.
這一場天國的地震.
一場在天堂的深處發動.
一路徹底地影響地上的秩序.
直到今天都是.
今天我們聚集一起的時候.
我們剛才就講了.
焦點放在那個巨大的石頭.
被滾開的事件.
這塊原本塞住 堵住.
耶穌霧門入口的石頭.
它不只是一塊物理上的.
巨大的石頭.
它其實更加是深刻地象徵著.
我們每一個人的生命裡面.
那些似乎不可以動搖的困境.
那些我們憑著自己的力量.
沒有辦法移開的阻攔.
以及上帝怎樣用祂復活的大能.
將這一切徹底地翻轉.
為我們帶來一個全新的生命.
一個全新的開始.
我們有兩段經文.
《馬可夫音》很清楚記載到.
在七日第一的早晨.
有幾個婦女 莫大拉的瑪利亞.
阿哥母親瑪利亞 薩羅米.
這幾位婦女他們買了香膏.

$^{81}$要去高抹耶穌的身體.
在他們前往墳墓的道路上.
心裡面充滿了憂慮.
彼此對問著.
誰會為我們推開那塊石頭.
經文很明顯地告訴我們.
那塊石頭很大.
這塊塞著部門入口的石頭.
你想想你去做這件事情的時候.
其實它是一個很巨大的阻攔.
但這塊石頭擺在這裡.
其實是承載著三種不同的視角.
也代表著我們三種不同的人生的處境.
這三位愛主的婦女.
她的視角看到的.
這塊石頭是他們和耶穌之間最後的阻隔.
是他們內心深處過去的.
傷痛絕望的具體的化身.
她曾經跟隨耶穌.
聽從過她的教導.
看過她所號痕的神蹟.
他們曾經在耶穌身上.
寄託著以色列德屬的盼望.
不過他們親眼看到.
耶穌被釘死在十字架的整個過程.
場面是殘酷的 是絕望的.
以至他們所有的盼望都似乎隨著事件.
十字架上高的吶喊.
去徹底的破裂.
現在他們是帶著最後對耶穌表達的愛意和敬意.
這是一個高沒 已經死去.
而且埋葬了三天的身體.
這是帶著深深的傷痛的告別.
是一種無法接受.
又要強迫自己去接受的現實.
在他們心裡這塊巨大的石頭.
不只是擋在霧門入口的障礙.
更加似乎是一個樹立在這裡.
成為了他們過去和美好的回憶裡面.
一個很強大的無法跨越的高牆.

$^{121}$是他們信仰和希望最終歸於虛無的終點站.
他們在路上高比次的對問.
哪個人能夠為我們滾開這塊石頭.
這個問題其實是每一個.
經歷過生命重大失落.
承受過深刻痛苦.
陷入過徹底絕望的人.
在內心的深處的共同呼叩.
他們帶著香料去到墳墓裡面.
想要用這種方式來保存一段.
他們無法真正理解的記憶.
因為其實是他們對逝去的愛.
和破碎的生命所做的最後一些致敬.
在他們眼裡面這塊石頭.
代表了一切美好事物的終結.
代表了無能為力的現實.
代表了過去被過去的傷痛所阻擋.
無法向前行的生命的狀態.
或者我們從當時的當權者.
去看這塊石頭的話.
這個是猶太工會的領袖.
從普彼拉多他們的角度去看的話.
這塊石頭代表了另一回事.
它代表了對現實的控制和權力.
是一種想要藉著控制我們看到的東西.
來確保自己的地位位置穩固的企圖.
他們記得耶穌曾經說過.
三天後我要復活.
這樣的說話是令他們很心裡不安的.
他擔心耶穌的門徒可能會偷走屍體.
然後就會說謊.
於是引起了動盪.
他們要求彼拉多.
將這個冰丁在墓門裡面看守.
甚至我們用蒙羅馬的官方的封印.
即是在石頭上封住.
對他們來說這塊石頭.
代表著他們權力意志的具體的展現.
是要企圖控制著歷史的發展方向.
壓制他們所認為真理聲音的象徵.

$^{161}$他們要用這塊沉重的石頭.
將耶穌這個曾經搞亂天下的人.
威脅他們既得利益.
亦亦宗教權威的危險人物.
永遠封存在歷史的塵埃裡面.
讓它從此消失.
這塊石頭代表著這個墮落的世界.
不斷地想要說服我們接受的謊言.
這個就是現實就是這樣.
你怎樣都改變不了.
它代表著人類內心深處的恐懼.
代表著壓迫性的制度.
代表著一切想要阻擋.
上帝的真理和生命的道.
進入世界的黑暗的權勢.
他們天真的以為.
只要封住這塊石頭.
就能封住真理.
他們以為只要派兵丁駐守.
就能確保自己的權力.
和地位永遠的穩固.
不過他們每晚都沒有想到一件事情.
他們要面對的根本不是一個.
需要他們去封鎖的死人.
而是那位掌管宇宙萬有.
勝過死亡權勢的生命的主.
如果我們從婦女從當權者.
或者從耶穌自己的眼睛看到這件事.
從耶穌基督的視角來看.
這塊石頭又是代表什麼呢.
它代表著死亡.
象徵著阻隔.
象徵著人類因為罪.
而和聖潔的上帝之間.
那種深厚的鴻溝.
死亡是人類最終的敵人.
所有恐懼的終極源頭.
當耶穌在十字架上.
大聲地喊愛神愛神.
為什麼離棄我的時候.

$^{201}$他經歷的就是那種.
和父神完全的隔絕的極致的痛苦.
是一種比身體的死亡.
更難以承受的靈性的黑暗.
但復活的故事.
向我們揭示一個很重要的真理.
就是這塊代表著死亡.
和隔絕的石頭.
這塊沒有辦法阻隔著上帝的愛.
更加沒有辦法禁錮.
那位本身就是生命的主宰.
就像創世記所說的.
亞伯蘭的事情一樣.
上帝和亞伯蘭立約.
今天裡面.
我們再次思想復活.
耶穌的時候.
這塊死亡的石頭.
這樣嚴嚴實實地.
封住了耶穌的墳墓的時候.
是三人一體的真神.
他權力親自介入人類歷史.
這個關鍵的時刻.
耶穌的復活向全宇宙證明了.
沒有任何力量.
無論是罪惡的權勢.
是死亡的毒鉤.
是陰間的門戶.
能夠阻隔著神.
和祂所愛的人之間的聯繫.
這塊在世人眼中沉重無比.
沒有辦法動的石頭.
在神的眼中.
不過是一個注定要被挪開.
被擊碎的障礙.
如果我們轉去馬太福音裡面.
看到另一件事情的時候.
他就會說.
突然地是大震動.
因為有主的使者.

$^{241}$從下來就把石頭滾開.
採上高.
這個不是一個普通的地震.
我們說這是一個天國的地震.
你大概猜到地震是什麼.
是不是.
天國的地震就是.
由上帝所在的地方裡面.
撼動人間的一切.
這是一個很重要的事情.
我們要想想.
這個天使為什麼要滾開石頭.
其實他不是為了讓耶穌.
能夠在裡面走出來.
復活的主裡面.
我們相信是一個榮耀的.
超越時空限制的身體.
如果我們記載在.
其他聖經的經文裡面.
耶穌是突然出現.
突然離開的.
所以其實耶穌要走出來.
不需要推開石頭.
那麼天使為什麼要做.
這些多餘的行動呢.
其實答案是非常簡單和清楚.
石頭被滾開.
是為了讓那些婦女能夠走進去.
親眼看到那個空墳墓.
是為了讓後來的門徒.
好像是彼得約翰能夠進去.
親眼看到那些細麻布.
和那頭巾放的位置樣式.
也都為了今天的我們.
能夠透過聖經的記載.
在心裡面看到那個空墳墓.
然後真實的相信耶穌復活了.
神之所以要挪開那塊石頭.
不是因為祂自己需要這樣做.
而是因為我們需要.

$^{281}$不是因為其他的原因.
而是因為祂知道我們的軟弱.
知道我們的信心不足夠.
祂知道我們需要一個.
可以見到的見證.
一個能夠觸摸歷史的事實.
來幫助我們在心底的深處裡面.
相信死亡已經被擊敗.
生命已得到最終的勝利.
這種復活的思想.
這種我們思想復活的大能.
今天仍然在運行中.
它很具體的.
或者調轉問.
它怎樣具體地滾開.
我們每個人生命裡面.
那些不同的石頭呢.
首先是那塊稱之為.
過去的石頭.
我們每個人心裡面.
都有一個屬於過去的墳墓.
這個墳墓裡面.
裝滿了我們各種各樣.
很後悔的事情.
無法釋懷的怨恨.
曾經破碎的夢想.
和那些無法彌補的傷痕.
這些東西就好像一塊沉重的石頭.
一幅一天這樣壓在我們的心裡面.
壓到我們呼吸不了.
讓我們一直活在自責.
孤獨的陰影之下.
沒有辦法真正地面對今天.
更加沒有辦法勇敢地走向明天.
但是耶穌的復活.
在今天裡面.
應該帶給我們的訊息就是.
那塊過去的石頭.
主已經為我們滾開.
神不要我們永遠停留.

$^{321}$在那個充滿死亡氣息的墳墓.
面前哭泣.
就好像天使對那些婦女說.
為什麼要在死人裡面找活人呢.
她不在這裡已經復活了.
我們為什麼還要一直留在那個.
充滿了腐朽氣味的過去裡面.
去尋找那個已經復活.
已經勝過一切的主呢.
她要帶領我們走出過去的捆綁.
帶領我們進入一個全新的開始.
若有人在基督裡面.
他就是新造的人.
舊事已過都變成新的了.
這個是復活的生命所帶來的確懼.
第二塊.
我們稱之為現實的石頭.
在我們每一天所面對的日常生活裡面.
我們會遇到各種惡樣的.
各個層面的巨石壓在我們的生活裡面.
可能是經濟的巨大壓力.
可能是你和人際關係那種難以化解的衝突.
可能是身體的疾病所帶來的痛苦和恐懼.
可能是對未來前途的迷惘和焦慮.
甚至可能是我們無法解除的惡習.
和反復的纏繞著我們的軟弱.
這些現實的問題.
就好像一塊一塊一塊的石頭.
鎖住在我們人生的道路上.
讓我們覺得無路可走 無計可施.
我們心裡也像那些婦女一樣發出哀嘆.
誰能夠為我滾開那塊石頭呢.
復活的主告訴我們.
那塊名叫困境的石頭.
我們的主同樣有能力將它滾開.
他不是挪開了石頭後才告訴我們要怎麼做.
他總是要我們帶著信心.
帶著我們手上僅有的香料.
亦都是我們雖然微小.
但真誠的 真心的付出和擺上.

$^{361}$走向那個看來毫無希望的墳墓的道路上高.
就已經準備好他的天使在那裡等候做工.
他那位為我們在曠野開路.
在沙漠開江河的神.
我們的責任不是要等到石頭滾開才願意出發.
而是像那些婦女一樣.
即使心裡充滿了憂慮和疑問.
仍然憑著對主那份撲掃的愛.
邁開腳步走向主所指示的方向.
當我們憑著信心行動的時候.
就會在行動裡看到神大能的彰顯.
過去今天的現實.
甚至有第三塊捆綁著我們的石頭.
就是對未來的石頭.
對大部分來說.
未來本身就是一個巨大的未知的石頭.
我們害怕死亡這個終極的未知.
害怕死後將要面對的審判.
害怕年老的時候沒有人照顧.
害怕最終孤獨一個人走完人生的路程.
這種對未來的深刻的恐懼.
總是不知不覺癱瘓我們現在的行動能力.
讓我們不敢冒險 不敢付出 不敢真正的去愛.
但是復活的主挪開了那塊石頭.
他叫死亡的終極石頭.
因為耶穌基督已經從死裡面復活.
死亡對於基督裡的人不再是生命的終結.
而是通往更豐盛生命的門戶.
就好像仕途保羅所發出的勝利的挑戰.
死了你的德性的權勢在哪裡.
死了你的毒鉤在哪裡.
因為耶穌已經復活.
死亡的毒鉤已經不掉.
他對我們已經沒有任何終極的剋制能力.
最後 假如我們回來教會很久 上了很久的福音班.
或許有一件事 你會留意到.
就是某福音的結局是比較特別的.
它有很多時候是好像突然的終結了.
其實我去到第八節的時候.
我們有很多的抄本都認為.

$^{401}$其實原本來說應該是作者刻意在這裡停止下來.
不再寫下去 為什麼呢.
無論你怎樣都好.
如果我們從第八節停下來去想的時候.
於是我們發覺到.
其實整個婦女的反應是惶恐 混亂和害怕.
天使只是宣告了一件事.
她復活了 然後吩咐你去告訴門徒.
耶穌要在家裡等他們.
就是這樣的話.
整件事突然間停下來.
其實是挺窄小的.
你看劇集就等於沒有了一個快樂的結局一樣.
你會停下來 那怎麼辦呢.
但是如果我們在相信或者思想中.
如果復活的耶穌基督沒有出現在馬可的記載裡面的話.
是沒有一個人能夠親眼看見復活的耶穌.
我們沒有辦法直接遇見祂.
只能夠聽見一個陌生的少年人的應許.
相比之下馬太的記載.
就是耶穌和門徒 上婦女和門徒顯現.
路加甚至描述在二馬五師路上的同行.
餐桌上的抹餅.
約翰記載了耶穌向門徒.
瑪利亞多瑪多次的顯現.
你看到馬可是一個很特別的說法.
這種說法太特別了.
特別到我們可能有很多的聯想.
為什麼聖靈要啟示馬可在這個位置這樣寫呢.
因為這個似乎不完美.
充滿了人性軟弱的結局.
其實只是聖靈留給我們每一代信徒的邀請信.
邀請我們每一個親自走入這個故事裡面.
用我們自己的生命 用我們跟隨主義的經歷.
來將這個故事繼續寫下去.
將它完成 馬可福音.
或許讓人感覺很驚慌就收尾.
沒有一個開心的結局.
但它要提醒我們.
信仰的道路從來就不是一帆風順的康莊大道.

$^{441}$它充滿了掙扎的掙扎.
內心的困惑 甚至一時的膽怯.
但是真正的信心.
不是裡面沒有恐懼.
而是即使帶著恐懼.
我們仍然選擇跟隨主義.
仍然選擇繼續往前走.
復慰的故事不是一個完美的.
與我們無關的童話故事.
而是一個震撼人心的.
需要我們用一生去回應的開始.
它要求我們每一個人聽見這個消息的.
都必須做出自己的選擇.
當你看見空墳墓.
你是繼續逃跑躲避.
還是轉身去.
將你看到的.
所經歷的告訴其他人.
這婦女她害怕是真實的.
面對超自然的事情.
面對天使的顯現.
面對空中的墳墓.
有誰不害怕.
但是他們沉默只是暫時的.
因為馬太告訴我們.
他們後來是跑去.
將這個大喜的信息告訴門徒.
他特意留下這個局面.
是因為正視我們的恐懼.
同時也超越我們的恐懼.
去完成這個故事.
我們也看到.
門徒起初的反應.
同樣是不順聽而橫.
馬可福音特別強調.
當這些人聽到耶穌的復活的時候.
他們是不相信的.
隨後耶穌向二萬五師路上的門徒顯現.
他們回去告訴其他門徒.
他們也是不相信的.

$^{481}$最後復活的耶穌.
親自向十一個門徒顯現的時候.
耶穌是擇避他們不順.
心裡是硬的.
這個記載是一個重要的事情.
不是因為門徒的信心大.
所以創造了耶穌復活的事實.
相反是因為耶穌確實已經復活了.
這個客觀的事實.
逐漸兼顧了門徒.
本來軟弱的和動搖的信心.
我們常常以為.
我們必須有足夠大的信心.
神才在我們心裡做事.
但是復活的耶穌.
祂告訴我們.
即使我們不順.
在我們的恐懼.
在我們的軟弱裡.
神仍然向我們主動的顯現.
祂的愛和祂的權能.
遠遠超過我們信心的軟弱.
天使在吩咐婦女.
去告訴門徒的事實.
特別提到一個名字.
就是彼得.
各位曾經三次公開否認主的彼得.
這個可能只是躲在某個角落裡.
自責 慚愧.
覺得自己沒有資格跟隨主的彼得.
復活的主仍然沒有忘記祂.
仍然在尋找祂.
仍然要接納祂.
這個表明無論你曾經跌倒的多深.
多恨.
無論你覺得自己多不配.
復活的主都在尋找你.
尋找你.
祂要親自移開你心裡那塊.
因為失敗 自責而產生的石頭.

$^{521}$《馬福音》這個特意留下來的.
未完成的事情.
其實是寫給我們每一個時代的教會.
亦是寫給我們.
這個故事真正的據點.
是要我們每一個人來畫上.
耶穌復活之後.
向門徒頒布了那個偉大的使命.
祂說你們往普天下去.
傳福音給萬民聽.
我們不再是站在旁邊的旁觀者.
我們是這個偉大故事的參與者.
是被復活的大能所差遣的人.
我們要去到處宣揚這個福音.
而主亦都同在用神蹟隨著.
證明我們所傳的道.
親愛的弟兄姊妹.
在整個講道完結的時候.
今天我要再問大家一個.
最直接最重要的問題.
這個就是你此刻的生命裡面.
你所面對最大的那塊石頭是甚麼.
是否好像那幾個父女一樣.
你正是被過去某些的傷痛失敗.
壓到喘不過氣.
始終沒有辦法釋懷.
沒有辦法真正地往前走.
是好像那些當權者一樣.
你被現實的恐懼和控制欲所捆綁.
害怕放手.
給上帝掌管你的人心.
害怕面對那些你無法預測的改變嗎.
還是像那些逃跑的父女.
和躲藏的門徒一樣.
你正是被未來的恐懼.
甚至可能是對死亡的恐懼.
深深的克制.
以至你無法自由地活在這一刻裡面.
無論你的生命裡面.
那塊石頭是叫甚麼名字.

$^{561}$是不信嗎.
是罪惡感嗎.
是恐懼.
是絕望.
還是死亡本身.
復活這件事所要傳遞給你的信息就是.
你一切都沒有辦法改變復活的主.
在面前站立得住.
他已經從死裡復活.
已大能夠證明他是上帝的兒子.
他擁有完全的權柄和能力.
要將你生命裡面一切的阻攔.
移開.
讓你去認識他 經歷他.
跟隨他.
頂子妹.
阻攔你的石頭.
通通都已經在他的應許之下.
被滾開.
今天這位的主.
這位仍然活著的主.
向你發出邀請.
他真是邀請你往前來.
到他的胸墳墓裡面親自看一下.
透過神的話語.
透過你心靈誠實的禱告.
真實地去看見那位復活的主.
已經勝過一切仇敵.
他也邀請你去告訴其他人.
告訴其他人知道.
將你手中緊緊抓住的那些香料放下.
將你心裡面一直以來的哀哭.
變成讚美的舞蹈.
去告訴你身邊的家人 朋友 同事.
知道這位復活的主.
在你生命裡面成就了何等的大事.
他更加要你勇敢地站出來.
不要讓自己繼續捲縮在自己的恐懼.
和自我的限制裡面.
復活的主已經呼召我們.

$^{601}$挺身昂首.
撥出亭夢照的恩.
相稱的復活生命.
石頭已經滾開了.
這不是一個口號.
這是已經在歷史裡面成就的事情.
並且今天仍然在我們的生命裡面.
運行的偉大的奇蹟.
天國的地震已經發生.
並且仍然繼續.
直到今天仍然在震動著每一個人的心靈.
耶穌基督已經戰勝了死亡.
他那個胸塘塘的墳墓.
就是我們得救的確懼.
就是我們永遠的盼望.
願我們不要再在死人裡面.
尋找那個永活的主.
不再在絕望裡面尋找盼望.
願我們每個人都躺開我們的心扉.
迎接著那位復活的主.
進入我們生命的每一個層面.
願我們不再逃跑躲藏.
而是勇敢地站出來.
成為他復活的見證人.
因為他活著.
我們也都能夠面對.
我們明日一起的挑戰.
因為他活著.
我們生命裡面.
從此擁有一個全新的意義.
和永恆的使命.
讓我們帶著復活而來的真正的起落.
勇敢地走出去.
讓這場從天而降的地震.
透過我們每一個人的生命.
持續的震動.
在這個仍然活在死亡陰影的世代裡面.
將那些不能夠震動的永恆角度.
帶給每一個仍然在尋找盼望.
仍然是需要聆聽這個好消息的人.

$^{641}$願復活的主.
賜福給我們每一個.
從今直到永遠.
阿們.
讓我們祈禱.
親愛的主願意你幫助我們.
讓我們一再去思考復活帶來的意義和力量.
求主幫助.
讓我們生命裡面被緊握的石頭.
主你為我們挪開.
讓我們看到上帝的工作.
上帝的恩典.
讓我們在當中被聖靈幫助.
使我們能夠釋放走出去.
將福音傳開.
求主帶領幫助.
以帶領幫助祈禱風主耶穌.
祈禱聖明祈禱.
阿门.
\newpage



\section{路加福音 24:13-32}
\label{sec:031953}
\textbf{愁容不再,火熱常在}~~~~石偉漢牧師
\newline
\newline
連結: \href{http://www.church.com.hk/acms/content.asp?site=cdc\&op=show\&type=product\&code=031953\&layout=sermon}{\texttt{031953}} ~~~~ 語音日期: 2026-04-12
\newline
\newline
\hyperref[sec:031952]{\small{< < < PREV SERMON < < <}}
~
\hyperref[sec:index_chronic]{\small{[返順時目]}}
~
\hyperref[sec:031954]{\small{> > > NEXT SERMON > > >}}
\newline
\newline
路加福音 24:13-32
\newline
\begin{longtable}{cl}
\hline
\hline
章節 & 經文 (和合本修訂版)\\
\hline
24:13 & \begin{tabularx}{0.7\textwidth}{X} 同一天，門徒中有兩個人往一個村子去；這村子名叫以馬忤斯，離耶路撒冷約有二十五里。 \end{tabularx} \\ \\ \relax
24:14 & \begin{tabularx}{0.7\textwidth}{X} 他們彼此談論所發生的這一切事。 \end{tabularx} \\ \\ \relax
24:15 & \begin{tabularx}{0.7\textwidth}{X} 正交談議論的時候，耶穌親自走近他們，和他們同行， \end{tabularx} \\ \\ \relax
24:16 & \begin{tabularx}{0.7\textwidth}{X} 可是他們的眼睛模糊了，沒認出他。 \end{tabularx} \\ \\ \relax
24:17 & \begin{tabularx}{0.7\textwidth}{X} 耶穌對他們說：「你們一邊走一邊談，彼此談論的是甚麼事呢？」他們就站住，臉上帶著愁容。 \end{tabularx} \\ \\ \relax
24:18 & \begin{tabularx}{0.7\textwidth}{X} 兩人中有一個名叫革流巴的回答：「你是在耶路撒冷的旅客中，惟一還不知道這幾天在那裡發生了甚麼事的人嗎？」 \end{tabularx} \\ \\ \relax
24:19 & \begin{tabularx}{0.7\textwidth}{X} 耶穌對他們說：「甚麼事呢？」他們對他說：「就是拿撒勒人耶穌的事。他是個先知，在神和眾百姓面前，說話行事都大有能力。 \end{tabularx} \\ \\ \relax
24:20 & \begin{tabularx}{0.7\textwidth}{X} 祭司長們和我們的官長竟把他解去，定了死罪，釘在十字架上。 \end{tabularx} \\ \\ \relax
24:21 & \begin{tabularx}{0.7\textwidth}{X} 但我們素來所盼望要救贖以色列民的就是他。不但如此，這些事發生到現在已經三天了。 \end{tabularx} \\ \\ \relax
24:22 & \begin{tabularx}{0.7\textwidth}{X} 還有，我們中間的幾個婦女使我們驚奇：她們清早去了墳墓， \end{tabularx} \\ \\ \relax
24:23 & \begin{tabularx}{0.7\textwidth}{X} 不見他的身體，就回來告訴我們，說她們看見了天使顯現，說他活了。 \end{tabularx} \\ \\ \relax
24:24 & \begin{tabularx}{0.7\textwidth}{X} 又有我們的幾個人往墳墓那裡去，所發現的正如婦女們所說的，只是沒有看見他。」 \end{tabularx} \\ \\ \relax
24:25 & \begin{tabularx}{0.7\textwidth}{X} 耶穌對他們說：「無知的人哪，先知所說的一切話，你們的心信得太遲鈍了。 \end{tabularx} \\ \\ \relax
24:26 & \begin{tabularx}{0.7\textwidth}{X} 基督不是必須受這些苦難，然後進入他的榮耀嗎？」 \end{tabularx} \\ \\ \relax
24:27 & \begin{tabularx}{0.7\textwidth}{X} 於是，他從摩西和眾先知起，凡經上所指著自己的話都給他們作了解釋。 \end{tabularx} \\ \\ \relax
24:28 & \begin{tabularx}{0.7\textwidth}{X} 他們走近所要去的村子，耶穌好像還要往前走， \end{tabularx} \\ \\ \relax
24:29 & \begin{tabularx}{0.7\textwidth}{X} 他們卻強留他說：「時候晚了，天快黑了，請你同我們住下吧。」耶穌就進去，要同他們住下。 \end{tabularx} \\ \\ \relax
24:30 & \begin{tabularx}{0.7\textwidth}{X} 坐下來和他們用餐的時候，耶穌拿起餅來，祝福了，擘開，遞給他們。 \end{tabularx} \\ \\ \relax
24:31 & \begin{tabularx}{0.7\textwidth}{X} 他們的眼睛開了，這才認出他來。耶穌卻從他們眼前消失了。 \end{tabularx} \\ \\ \relax
24:32 & \begin{tabularx}{0.7\textwidth}{X} 他們彼此說：「在路上他和我們說話，給我們講解聖經的時候，我們的心在我們裡面豈不是火熱的嗎？」 \end{tabularx} \\ \\
[1ex]
\hline
\hline
\end{longtable}
$^{1}$好,弟兄姊妹,主內平安.
讓我們一起有祈禱.
主要是幫助我們在你面前去謙卑,去學習.
幫助我們去闖開我們的心門,去聆聽.
幫助我們更認識你.
以致我們生命裡找回我們的那團火.
以致我們能夠好樣為主你去打美好的仗.
圓著你幫助我們眾人.
聽我們祈禱,奉耶穌名字而求.
阿們.
本來這幅圖不是這幅.
原來本來做出來的這幅圖.
AI幫我做出來的這幅圖.
其中左邊那個人.
有三隻手.
所以我看完之後.
不太妥當.
AI也不太信得過.
所以後來我再AI中的AI再修正.
才出了今天這幅圖.
事實上今天我們去看一個題目.
「仇容不再,火熱常在」.
所以有時候我們的念容.
也很反映出我們的心理狀態.
或者我們的表現.
所以今天跟大家先玩一個遊戲.
就是我給大家看到有四個Emoji.
四個Emoji的圖像.
考考大家能不能夠.
按著這幅Emoji的圖像.
其實能不能夠猜出它是什麼感受.
或者什麼情緒呢.
好來了.
我們試試猜猜最左上角那個.
左上角你覺得這是什麼情緒.
或者什麼感受呢.
這個Emoji是什麼.
開心.
這個比較容易.
開心.

$^{41}$Happy.
很Happy.
接著右上角.
右上角.
你覺得這是什麼呢.
戴著太陽眼鏡.
我忍不住說很Chill.
但其實這個答案不是Chill.
是Relax.
證明大家想的.
跟原本的圖像不一樣.
好那是什麼呢.
就是Smart.
他覺得很Smart.
很醒神.
很聰明.
所以就考考下一個.
左下角.
左下角你猜猜這是什麼感受呢.
什麼Emoji呢.
張開嘴巴.
看到牙齒.
憤怒.
憤怒通常有些火在旁邊.
你看那些Emoji都是這樣.
好猜猜這是什麼.
A字頭的.
Angry.
Angry都有些火.
好厲害.
Amy覺得Awesome.
你都覺得.
這個就是Awesome.
好最後一個.
右下角是什麼感受呢.
不順.
上班通常都會有這種感受.
那些夥計都知道.
上班有什麼感受.
就是Stressed.

$^{81}$會有壓力.
上班.
如果上班都沒有壓力.
這個應該是有點問題.
所以.
迫著不用擔心.
你上班有壓力.
是正常的.
所以.
我們看到不同的Emoji.
其實都反映不同的情緒.
感受.
事實上我們的面容.
真的可以讓人看到.
我們今天是什麼狀況.
所以如果問問你.
今天你的Emoji是什麼呢.
如果今天我要.
你用Emoji去表達的話.
Emoji會是哪一個呢.
有心心還是有淚水.
還是嘔白抱.
還是臉都紅了.
的畫面呢.
還是有些星星.
還是不知怎樣的畫面.
事實上真的是這樣的.
今天可能你都有你的Emoji.
所以大家.
第一次見面的時候.
留意一下對方的Emoji.
他的面容是怎樣.
懂得說一些合意的話.
很開心的.
什麼事這麼開心.
讓我們能夠彼此去關心.
而今天其中的兩位主角.
他的Emoji是什麼呢.
他的Emoji.
你剛才看經文.

$^{121}$或者看我們今天的題目.
都猜到了.
就是什麼呢.
就是受容.
因為在我們剛才讀的經文裡.
我們看到24章17字說.
他們就站著臉上戴著受容.
所以這兩個是什麼人呢.
就是仇人.
有仇容的就是仇人.
仇人見面.
就是怎樣.
憤愛.
是的.
所以今天這兩個仇人.
滿有仇念的仇人.
見到主的面之後.
他將要經歷什麼.
憤愛眼紅.
因為他會火熱起來.
那麼盼望今天的訊息都能夠幫我們.
我們作為一個仇人也好.
或者一個不同情緒的人裡.
我們能夠見到主面.
能夠將我們化為火熱.
的一個身體.
或者心理的狀況.
今天跟大家去讀的經文.
就是在路加福音24章13.
到32節的部分.
其實稍後都會有少許提及.
稍後的經文.
不過重點就是我們今天.
去到第32節的部分.
今天最主要是.
分開做兩個的分題.
去思考今天的經文的段落.
第一個的分題就是.
從個人的期望.
到屬神的盼望.

$^{161}$從個人的期望.
就是你自己的期望和期盼.
以至到去到.
從神而來的那份盼望.
到屬神的盼望.
從神而來的那份盼望.
那種的轉變.
我們讀到的經文.
一開始的時候是.
第13節說.
同一天門途中有兩個人.
住在一個村子去.
往一個村子去.
這村子名叫二萬五師.
離耶路撒冷約有二十五里.
讓我們先去了解.
這個經文的背景.
因為這個經文一開始的時候.
其實它特別去講到.
就是在同一天裡面.
同一天是哪一天呢.
當我們要理解.
這班門徒去二萬五師的路.
究竟同一天是指哪一天呢.
當我們去.
試一下回想.
其實這個同一天.
其實就是七天的第一天.
因為這段經文當你對應.
在其他福音書裡面的時候.
就會知道這一天就是.
那個七天的第一天.
就是聖經裡面所講耶穌復活的.
那一天.
耶穌復活的那一天.
在上個星期我們的復活節的.
戶外崇拜的主日裡面.
我們都會看到約翰福音第二十章.
其中我和大家去到.
看到第二十節的部分.

$^{201}$當中其實你看.
其實有幾個段落.
當中其實也有兩個時間的概念.
因為在約翰福音.
第二十章第一節的時候.
它說七天的第一天.
就是耶穌復活的那一天的.
天環黑的時候.
穆大拉的瑪利亞.
來到墳墓裡面.
發現墳墓是空著.
而當中後來.
他們就跟門徒說.
而當中彼得和.
耶穌所深愛的那個門徒.
我們相信就是約翰.
也去到墳墓裡面.
結果他們就看到墳墓真的空著.
而是回到他們的住處.
住的地方裡面.
而後來瑪利亞.
也在墳墓外面哭.
以致後來.
遇上了天使和他對話.
後來他也把.
天使和他報的訊息.
帶回門徒那裡.
但是當你們繼續看到.
去到第十九節的時候.
就把那個.
穆大拉的瑪利亞.
在墳墓裡面遇上.
天使遇上耶穌之後.
就已經跳到那一天的晚上.
那一天的晚上.
就是七天的第一天的晚上.
門徒也因為.
害怕一些猶太人的迫害.
或者是爛了.
所以他們就在.

$^{241}$所在的地方裡面關了門.
後來耶穌就來到.
站在他們當中就說願你們平安.
於是就把他的手.
和臘袍給他們看.
門徒就帶來了喜樂.
事實上你會發現.
在約翰方裡面由一至十八節.
其實就是在說那一天的早上.
而去到.
十九至二十節就是那一天.
復活節那一天的晚上.
其實就是.
你會看到其實由早上.
到晚上之間其實還有一個.
時間的空檔.
還有一個時間的空檔.
就在同一天發生.
而今天的經文正正好.
就是在這一個復活節當日.
門徒已經知道.
穆大拉的瑪利亞.
穆大拉的瑪利亞報來的消息之後.
就會發生了.
接著中間這件事情.
所以我們承接著上個禮拜.
關於復活節的訊息.
我們繼續去思考多一點.
復活節的訊息.
我們完了才回去看.
我們的登山部分的段落.
這裡說就是這樣.
就在同一天.
雖然我們沒有明確知道.
究竟是很早還是很晚的時間.
但是我們知道發生的中段.
就是這些.
可能上午的時間.
或者中午甚至黃昏的時間.
這裡說的就是.

$^{281}$門徒中有兩個人.
去一個村裡面.
那個就是二萬五師.
離那裡就很遠.
但是我們要想想為什麼這兩個門徒.
會離開了.
耶路撒冷而去到二萬五師呢.
我們又要回想.
那一個早上.
大家剛才看回耶穌復活的時候.
其實.
那些門徒的.
那些人的表情是怎樣的.
其實很複雜的.
不是單一個面容.
就能夠描述得到.
那幫人那種的.
情緒或者那種面容.
因為如果你看回.
約翰福音 馬赫福音 盧家福音.
你會見得到就是.
原來在那天的清早裡面的時候.
其中剛才說到.
去到墳墓.
其實如果你再看回其他聖經的段落.
其實就不只是一個婦女.
去到耶穌的墳墓.
帶著一些香料去到.
耶穌的墳前.
其實還有其他的.
不過 莫大拉的瑪利亞.
似乎是當中這個見到耶穌的時候.
一個很重要的人物 主角.
所以就會凸顯出.
在約翰福音裡面凸顯了他.
所以當你看約翰福音的時候.
就會見到這個莫大拉的瑪利亞是怎樣的呢.
他的情緒就是一個哭的情緒.
所以呢.
天使問他為什麼哭.

$^{321}$耶穌見到他都問他為什麼哭.
當你又看到馬赫福音的段落.
說到這一天復活的日子的時候.
他們這幫婦女.
他說見到什麼呢.
見到一個穿白袍的少年人的時候.
他們很驚奇.
甚至乎他們也有.
恐慌 又驚慌.
去到後來第八節.
甚至乎見得到就是說什麼呢.
當他聽到天使跟他們說的話的時候.
他們出來的時候.
從墳墓那裡就逃走.
又發抖又驚訝.
因為他們害怕.
盧教父還未給我們說.
二十四章四字五字說.
原來他們都為這件事去困惑.
接著後面說.
他非常害怕.
所以你可以想像一下.
原來耶穌復活這件事.
對於莫大拉的瑪麗亞.
甚至乎當時一起有份參與.
去到的那幫婦女來說.
其實經歷了很多複雜的情緒.
因為這件對他們來說.
是一件震驚的事.
震驚不單止耶穌的遺體不見了.
更加是他遇上天使跟他們對話.
甚至乎也都見到主耶穌.
整個的心路歷程.
也都發現出他們.
很多裡面交織著的情緒感受.
以至到的人.
可能都會有很多不同的表現出來.
其實除了他們.
還有比德.
甚至乎可以包括約翰.

$^{361}$但是當然我們看回.
路加福金第二十四章十二節的時候.
特別提到比德的時候怎樣呢.
比德起來的時候.
就跑到墳墓前.
就俯身往裡看.
看見細麻布就回去.
因所發生的是心裡驚訝.
原來比德.
見到不見了耶穌的遺體.
不見了的時候.
他裡面都是很多驚訝的.
情感在裡面.
那其他門徒呢.
其他門徒的時候.
路加福金又說到.
他們這些.
就是那些婦女回來匯報.
回到他們聚集的地方.
匯報的時候.
那些使徒以為是胡言.
就不相信.
以為這些婦女亂說.
接著就不相信他們說的話.
因為他們沒有親眼見過.
也不相信.
真的會發生死而復活的事.
所以你可以想像一下.
這幫人呢.
就面對著這種不同.
因為耶穌受死復活.
之後所帶來的.
不同的狀況呢.
帶來很多很複雜的情緒.
而這幫人.
在不斷地知道.
遭到這些事的時候.
其中有兩個人.
就有一個不一樣的行動.
做出來.

$^{401}$所以除了我們知道.
時間是那一天.
今天我們讀到的人物.
都會有兩個的人物.
就是在說這幫聚集的門徒裡面.
其中有兩個人.
一個叫夾劉巴.
就是.
另外一個就在整個經文段落裡面.
就沒有寫到名字的.
但是我自己相信.
他的名字就應該不是叫.
渣劉灘.
因為他真的.
應該都很認真地跟隨著.
但是夾劉巴和這個朋友.
他們似乎就是.
在這件事上午聽到.
耶穌不見了.
甚至有人說他復活的時候.
他們的不信令到他們.
似乎有這兩個人.
就選擇了什麼呢.
離開了耶路撒冷.
首先我們要問.
他們在耶路撒冷哪個地方呢.
你有沒有想想.
以前的社會不是今天的社會.
今天的社會很簡單.
我們應該在哪裡等呢.
我們一大家發短訊.
大家在不同的地方.
都知道在哪裡等.
但是在耶穌的年代.
沒有WhatsApp 沒有手機.
怎樣能夠.
門徒能夠聚集呢.
就一定他們有一個.
很重要的聚集點.
聽得明白嗎.

$^{441}$不然他們不會知道.
去哪裡集合.
而他們集合的聚集點.
極大可能就在.
耶路撒冷的哪裡呢.
我們相傳極大可能.
就是在馬可.
因為我們知道.
馬可的母親接待耶穌.
甚至我們相信.
最後晚餐的地點.
都是在馬可的母親的.
upper house 樓上的地方裡.
他們進行最後晚餐.
所以似乎.
在當時的時間裡.
這個地點馬可樓.
就是在耶穌.
受死前和受死後.
門徒一個聚腳點.
一個會合點.
一個情報收集中心.
其實就是這樣.
所以你會知道.
為什麼當這些婦女.
見不到耶穌的遺體的時候.
收到天使告訴她.
耶穌復活的時候.
她們能夠準確地去到.
告訴門徒呢.
因為她們就回到馬可樓.
立即可以通知她們.
給她們東西看完之後.
又再回去的時候.
又回到馬可樓.
所以到了復活節晚上的時候.
耶穌顯現的時候.
又在他們聚集的地方.
在那裡顯現.
這就是他們行常出現的聚會點.

$^{481}$而這幅圖就是今天.
其中一個.
在耶路撒冷裡面.
被認定是其中一個.
就是馬可樓的地方.
這就是在耶路撒冷的地方.
一群門徒聚集的地方.
但是現在就是.
當聽到耶穌的消息之後.
其中這兩個門徒.
似乎在一個懷疑不信.
失望的狀態裡面.
就離開了他們的團隊.
你聽得明白嗎.
他們聚集 可能還要商量.
接著怎樣發展這件事情.
但是他們在耶穌復活的那一天.
就已經決定離開了耶路撒冷.
離開了馬可樓.
離開了的大隊.
而去了另外一個地方.
去哪裡呢.
就是去二馬五師.
二馬五師是一個什麼樣的位置呢.
這裡你們地圖可以看到.
小小的藍色圈.
就是耶路撒冷.
而左邊的藍色圈.
就可能是.
二馬五師的地方.
因為如果你看回.
剛才經文告訴我們.
這個是一個村子.
而且這個村子.
其實按照馬可福音的描述.
他說這個是一個鄉下.
鄉下.
就是說.
不是一個那麼熱鬧繁華的城市.
未必是那麼多人熟悉的地方.

$^{521}$甚至乎.
原來這個二馬五師的地方.
在學者.
一些歷史學家的研究裡面.
都是某程度上是.
不見經傳的.
不見經傳的意思就是說.
其實隔了很多百年.
他們都未能夠確定到.
哪一個地方就是正式的.
二馬五師的地方.
所以你看得到這裡有四個紅點.
這四個紅點都是.
相傳說有可能.
是二馬五師的地方.
不過重點都不是.
準確地哪一個紅點.
就是二馬五師.
而是他們離開了耶路撒冷這個中心.
而是去到二馬五師的時候.
就是願意走了一個很長的路.
去了一個村裡面.
去了一個鄉下的地方裡面.
去.
當然我們要了解一下.
這個是一個什麼樣的地方.
或者他想為什麼要去這個地方.
而當中的距離有多遠呢.
我們看看你的和修本.
你的翻譯叫做.
二十五里.
其實這個二十五里.
用現在的.
計算方法你就簡單一點可以理解.
大約就是七個馬路.
即是七個現在我們的.
英里.
又或者我們用公里去計算.
就是十一個公里.
我們跑.

$^{561}$有個全馬.
半馬和一個不知道什麼馬.
十公里的馬.
多一個就是十公里.
跑十公里要跑多久.
二汶你跑開.
一個小時可以嗎.
跑不可以.
如果真的跑就應該要一個小時.
但你要知道這兩個門徒是什麼呢.
心情低落.
即是悠悠愁愁地離開耶路撒冷的時候.
總不能想像一下.
一個很憂愁的兩個人.
很輕快地在跑步.
所以我們可以想像一下.
大約他們一路慢慢走.
離開了二馬五絲的時候.
可能就不是說.
跑步一小時.
甚至乎不是平路.
不是今天跑平的馬路.
是山路有上有落的時候.
可能兩個小時.
甚至乎再多一點的時間.
他們就在這個時間裡面.
走這麼遠的路程.
去到.
二馬五絲的地方.
而在這個.
整個路段裡面.
他們有什麼做呢.
走這麼長的路有什麼好做呢.
當然不是說是非.
就是什麼呢.
聖經告訴我他們彼此議論.
他們交談議論.
說他們一路在走的過程裡面.
就是因為.
他們有一件事困擾.

$^{601}$他們不斷去講及這件事.
講完又講.
所以是一種反覆討論.
不斷交換意見.
去講一件事.
究竟他們一路走.
一路走一路走.
究竟在講什麼事呢.
所發生的一切事.
是一件什麼事情呢.
我們見得到.
在第十五節告訴我們.
在他們交談的時候.
耶穌都親自走近和他們同行.
耶穌親自.
想知道發生什麼事.
很有興趣知道.
究竟是什麼事發生了呢.
於是耶穌就對他們說.
你們一路走一路走.
一路談的.
彼此談論的是什麼事呢.
耶穌就問.
耶穌就這樣悄悄地走過來.
問這兩個.
心情一般般的門徒.
就來問他.
發生了什麼事.
那兩個人有什麼反應呢.
其中有一個人.
就是叫夾劉巴這個人.
他的反應是什麼呢.
他說你在耶路撒冷的旅客中.
唯一還不知道.
這幾天在哪裡.
發生什麼事的人嗎.
因為對於.
在當時耶路撒冷.
離開的人.
就是那幾天在那裡的人.

$^{641}$自然就知道.
他們當時的熱搜.
是什麼事.
如果用今天的角度.
照Google一搜尋.
打Jesus那個字.
那兩個英文就自然出現了什麼呢.
例如.
Jesus加CR就自然出現.
Cross, crucification等等.
就是歸耶穌的事情.
當你打.
Jesus DEA.
就自然說是他的死亡.
事實上你可以想像一下.
就是這樣.
在耶路撒冷那個地方.
一件如此震撼的事件.
成為當時耶路撒冷所有.
在那裡出現過的人.
的英文討論.
成為了熱搜.
但是現在竟然耶穌.
不知道而去問.
這兩個的門徒.
所以這兩個人.
當中的格劉巴就很驚訝.
就說你不知道嗎.
於是他就嘗試去解答.
告訴耶穌聽.
究竟發生了什麼事.
究竟發生了什麼事呢.
按照接下來的經文告訴我們.
關於有三個重點.
不過.
接著其中一個一開始的時候.
他特別說到這件事.
是什麼呢.
他說.
就是撒撒冷人耶穌的事.

$^{681}$其實這挺有趣的.
因為這兩個人.
就對著耶穌.
說耶穌.
我們經常都要說耶穌.
但是這兩個人就對著耶穌.
說撒撒冷人耶穌的事.
其中有三個重點.
第一.
他說耶穌原來在他們眼裡.
一直跟著耶穌的時候.
因為這是一個先知.
在上帝和眾百姓面前說話.
行事都大有能力.
這個人.
耶穌這個人是滿有能力的.
是有很高的地位.
甚至好像先知一樣.
甚至可能會帶來一個.
國度的復興.
一個國家的.
再一次興起.
所以說他們第一件.
說到耶穌生前的時候.
他們看到耶穌的形象.
是很厲害的.
你想像一下.
萬千寵愛在一身.
所有焦點.
都似乎在耶穌身上.
甚至乎你記得.
在之前一個星期.
那些人歡迎耶穌進入耶路撒冷.
和撒冷和撒冷.
所以是一個這樣的焦點人物.
充滿奇葩.
接著.
這個甲巴劉和他的朋友.
和門徒一起說.
第二件事是什麼.

$^{721}$但是.
但是有件事發生了.
就是祭司長和他們的官長.
把他抓了.
而且釘死了他.
在十字架上.
其實釘這個死罪.
不是由祭司和.
民視等等去做的.
而是由羅馬政府.
我們看到.
那些人.
去做這個決定.
但是現在似乎在門徒眼裡.
就是因為這些祭司長.
和祭司民視等等.
他們發動這個問題.
令到耶穌被釘死.
於是就是這樣.
他們知道耶穌被釘死了.
而且這件事已經發生了三天.
代表了什麼呢.
希望還滅.
剛才第一點裡面.
講的那種希望.
到了第二點的時候.
因為他被釘死罪.
以致整個救贖.
都帶來一個幻滅.
還有.
在幻滅之外.
還有第三點就是.
那幾個婦女還發現了.
耶穌他們這個.
老師的遺體不見了.
即使.
這些婦女告訴他.
天使顯然.
有人告訴他天使活.
但是他們因為什麼.

$^{761}$只是沒有看見祂.
未能親眼看見耶穌復活的時候.
他們心裡仍然.
不太相信這件事是真實.
就是說.
耶穌既然死了.
還要有人說祂復活.
還有一些謠言.
有些傳言好像關於耶穌發生了.
究竟是怎樣呢.
究竟什麼是真什麼是假呢.
你又記得我們之前.
在上一次都說到.
當中祭司長.
他們發一些什麼假消息.
發假消息就是說.
有人偷了耶穌的屍體.
所以在這個.
門徒之間似乎.
很多人不能夠肯定究竟發生什麼事.
以致連.
老師的屍體都沒有了.
不見了.
第一姐妹你有沒有想想為什麼他們.
帶著仇用呢.
因為門徒面對著他們.
很高的期望.
和一個很大的失望.
因為他很清楚說.
我們素來所盼望要救贖.
以色列的民族.
就是他了.
想著早幾天就要發生.
這個高峰事件.
但是現在跌落一個深谷裡面.
一切希望幻滅.
甚至進入一個迷惘的當中.
第一姐妹.
原來有時候我們都會這樣.
有時候我們發現.

$^{801}$我們生命裡面很多不開心.
很多的仇用滿面.
正正很多是因為.
我們本來帶著希望.
但是我們卻面對失望.
我們一次又一次.
為自己製造很多的希望.
很多的期望.
但是有時候又一次又一次.
我們面對著我們的失望.
所以呢.
真的值得我們去想一想.
究竟我們有沒有經歷過.
對神的期望和失望呢.
我們的期望.
會不會也讓我們.
不知不覺帶上仇用.
你的計劃是不是神的計劃呢.
你是不是明白神的計劃呢.
真的.
我們人生裡面.
總是面對著很多.
希望越大.
失望越大.
是吧 我們興合合地.
用一筆錢去投資.
去搞生意 我們希望.
做好去發財來做.
但是有時候.
在Florida做生意.
真的不容易.
可能一間又一間倒閉.
可能你那間也會是其中一間.
那個時候.
就是最令你最痛.
最迷惘 最受苦的時間.
有時候我們又會想著.
啊 我看到這個人.
這個人不錯.
會不會成為你的另一半呢.

$^{841}$開始希望.
有好結果 拍拖.
很大期望.
但是一邊拍一邊又發現.
不是你想要的東西.
那個失望就覺得.
唉 我浪費了這麼多時間.
究竟怎麼辦呢.
有時候好像我們這樣.
我們又會期望著.
我們這間教會.
要找好牧師.
現在經歷了三年.
你就知道你有多失望了.
因為呢.
你期望到你真的面對的時候.
那個狀態有時候.
是不一樣的.
你期望的牧師都不是萬能.
有很多事情以前你見過.
某些牧師做得到的.
不代表我做得到.
有時候都會在這些位置.
都有失望的.
怎麼辦呢.
第一次會想想.
希望與失望之間的問題呢.
在最近.
我們加勒團都看了一個.
《恩怨之聲》的影片.
當中那個影片裡面.
訪問到一個人.
他就叫鄧英奇.
他就講到這個人.
其中都跟我們.
有些路途可能都很相似.
就是他是一個移民.
他是一個新移民.
他是一個哪裡的新移民呢.
他就是去了.

$^{881}$加拿大卡加里的一個新移民.
事實上他自己.
一路他都.
覺得自己能夠有機會.
很快就會順利去.
去到加拿大開展一個.
全新的家庭生活.
因為他自己在他.
讀書的時期裡面.
他都是去了澳洲留學.
他經歷過在一些外國人的.
地方裡面生活.
很明白外國人的生活處境.
也有一些溝通的能力.
也都.
很期望他自己.
還是年輕有很多的適應.
於是他在2021年.
他就帶著一家四口.
就移民去到加拿大.
的卡加里裡面.
很大的期望應該.
就是.
他那時候已經信了耶穌.
他應該都是上帝有祝福.
上帝一定會帶我.
有一個好的新的地方.
在那一年裡面.
他經歷了一個什麼狀況.
就是他.
初初去到他都申請一些特快.
那些移民的方法.
結果在.
一路申請的過程裡面.
經歷了一年的時間裡面.
一路都搞不到.
身份的轉變那個簽證.
就是那個轉變.
以致他面對一個身份.
的危機究竟.

$^{921}$怎樣呢他移不移得成.
這個去到加拿大裡面呢.
因為當他.
不斷在這裡做的時候.
他的工作簽證.
留在這裡的時候初初都有工作簽證.
但是後來他.
都面對著.
被老闆炒魷魚.
失去工作.
就是跟著.
他的積蓄不斷減少.
他的移民申請.
後來發現還是被拒絕.
他生命裡面.
出現了一個自殺的念頭.
他覺得自殺可能是一個方法.
因為他說.
其中一個他講到他的.
想法就是如果我自殺.
我死了之後政府可能會可憐.
我兩個子女.
被他們批到簽證.
是不是很傻這種想法.
是不是很傻但是.
去到絕路的人.
什麼傻的想法都會想到.
如果不是都不會想自殺.
我想在這個過程裡面.
他似乎帶著很大的期望.
去到但是.
失望一次又一次.
一直搞了很久都未能夠.
順利搞到他的移民.
甚至面對身份危機.
面對著人生的曠野.
他怎樣.
走出來呢.
過程裡面他有回教會.
他在加拿大裡面.

$^{961}$回到教會裡面.
似乎參與了一些小組.
怎樣去認真去讀聖經.
他唯一的信念.
就是相信神如果當他是.
子女的話這個神.
這個天父爸爸是不會.
不理他的因為他都在做爸爸.
就是憑著這個信念.
於是他每一天去堅持.
他要順著由創世.
一路去讀聖經讀聖經就好像.
我們參與我們的讀經運動一樣.
一路堅持去讀堅持去讀.
他就在讀聖經的過程裡面.
明白多些神得到神的安慰.
特別在約書亞記裡面.
他得到很大的幫助.
所以在這個過程裡面.
他慢慢他經歷一種生命的.
轉變就是什麼呢.
他開始禱告不同了.
他說他.
最初的時候很想求神.
給他解決方法.
現在呢他就只好想.
在神的裡面找到那份平安.
他一開始的時候.
那個祈禱的要求就是.
神啊你要批出那個PR.
一個永久居留權.
但是後來他的禱告就是.
神啊如果這個不合你心意.
就求你.
求按住你的意思.
因為你給我的才是最好的.
就是從一個.
自己人生的計劃裡面.
一路在想自己可以怎樣怎樣.
以至到去到那一刻.

$^{1001}$他知道他要回到上帝裡面.
上帝你的計劃是什麼.
上帝你的計劃是最好的話.
我無論去或者留.
我都按住你的計劃.
是最好去走下一步.
弟兄姊妹原來我們的生命裡面.
就是怎樣能夠由.
那種個人的期望.
能夠轉到神的盼望呢.
就是我們能夠透過一步一步.
明白上帝的心意.
我們就能夠轉眼仰望耶穌.
定睛在神身上.
這個弟兄能夠這樣去.
轉變他.
以至到他今天雖然有些地方.
未能夠完全解決到他的.
一些簽證的問題.
一些PR的問題.
但他在神裡面有安慰有盼望.
事實上我們見得到.
這裡耶穌怎樣能夠幫助.
那兩個的門徒.
能夠解決他生命的問題呢.
這裡說在25節開始說.
無知的人先知所說的一切.
說你們的信心.
你們的心信得太遲鈍.
基督不是必須受這些苦難.
然後進入他的榮耀嗎.
於是他從摩西.
和眾先知起.
凡經上指著自己的話.
都給他們作了解釋.
我們見得到這班活生生的門徒.
在一個沮喪低沉裡面.
他們面對一個.
面對著死去的耶穌的時候.
他們不知道怎麼辦.

$^{1041}$他們有生命.
卻活得似生不如死.
當他們面對著.
活著的耶穌.
站在他們面前的時候.
他們仍然不知道.
他們原來已經見到.
他們的主.
耶穌說他們的心信得太遲鈍.
如何能夠幫助我們.
處理我們信心太遲鈍.
我們的信心不足.
我們對神那種信.
那種盼望好像很飄忽.
這裡說神的話.
是真實的.
當你的信心不足的時候.
應該要回到神的話裡面.
明白上帝的大圖畫.
所以你才能夠明白.
耶穌面對的苦難.
正正好要帶來的是.
將來的榮耀.
耶穌就是這樣說.
祂必須要受這些苦難.
以致到將來的榮耀.
這幅是更全備的圖畫.
人面對苦難.
面對死亡.
往往就是希望還滅.
但是耶穌.
面對苦難 面對死亡.
就是一個神聖的救贖.
為我們燃點希望.
所以我們要懂得用神的眼光.
回看那個圖畫.
我們就知道.
個人的期望有時候真的不太重要.
如果你能夠配合到.
神的盼望 神的大圖畫.

$^{1081}$神的大計劃的話.
你可以活得更好.
所以經歷這個.
信心太遲鈍的時候.
就有這些耶穌親自.
去解釋眾先知.
舊約裡面所說的一切.
怎樣在耶穌身上.
能夠應驗.
當你如果能夠同樣相信.
舊約裡面所說的.
在耶穌身上應驗的話.
你就知道信靠祂是可能的.
信靠祂是有盼望的.
所以去到第32節的時候.
他們很清楚就是.
原來當我們一路一路走.
一路走 走到二萬五師的路上的時候.
他說 在路上.
跟他們說話的時候 一路講解的時候.
原來他們的心裡面.
已經出現了一團火.
真是很愁的時候.
裡面是濕的.
是點不到火的.
但原來不知不覺.
藉著明白神的說話.
原來那些柴.
已經慢慢乾了.
可以重新再點燃一個火.
讓我們火熱.
所以我們要明白.
我們真的要認真去讀神經.
特別我們知道.
我自己在教會裡面.
去講道的時候.
我都期望就是.
懂得經講道.
解釋多一點.
聖經的講道.

$^{1121}$因為有時候我都面對一個問題.
就是很多弟兄姊妹說.
穆斯理講道.
講了很多的道理.
很多的知識.
但好像應用不到.
但我想告訴你.
你不明白那些知識.
去決定你能不能夠應用得到.
有時候我們的盼望搞錯了.
因為我們對於神的計劃.
我們誤解了.
唯有我們能夠清楚.
上帝在聖經舊約.
和新約的大圖畫.
我們能夠有合適的理解的話.
我們才能夠點燃裡面的火.
當你回教會回得太久.
還是沒有一團火的話.
你就要問.
究竟你是不是明白聖經呢.
在整個神的救贖裡面.
有一個計劃.
而你是可以參與上帝的計劃裡面.
你是有一個角色.
你不會是一個閒人.
很多時候我們想著回教會.
就是我在這裡坐著就可以了.
我坐兩個小時.
就可以了.
這個星期就夠力量了.
這個可能是你一個經歷.
但這個不是一個足夠.
支持你生命.
有一團火的經歷.
所以高明謙博士.
他有一本書叫《識經識出火》.
其中說的是.
當你認識聖經.
能夠清楚解釋聖經的話.

$^{1161}$你就明白聖經.
如何能夠幫助你生命裡面有火.
求主幫助我們.
好好花點時間.
了解聖經說什麼.
神的計劃是什麼.
我們如何能夠參與.
上帝的計劃.
所以我自己在加勒團裡面說.
我們加勒團.
不只是吃喝玩樂.
我們有一個使命.
他們知道我們十一月.
有什麼使命.
盼望他們準備.
我們會開始逐步準備.
我們十一月.
真的有一個使命.
我們有一個感恩節.
的報道聚會.
去年因為一些原因我們停了.
我們盼望今年.
會有另一個.
你不要說不關你的事.
因為神的救世計劃是預了你的.
盼望你都能夠.
預備好自己團伙.
弄清楚你為什麼.
去參與這個報道事情.
第二個.
跟大家分享的重點就是.
從眼睛模糊.
到眼睛開了.
因為當你看到.
第十四節的時候.
他們一直在談論.
談論談論談論.
耶穌走過來的時候.
跟他們一起走.
但是十六節說.

$^{1201}$可視.
他們的眼睛迷糊了.
認不出耶穌.
直到去到.
第三十一節.
他們的眼睛才能夠開了.
才認得出耶穌.
當之後.
耶穌就在他們面前消失.
但我們要想的就是.
為什麼之前.
他們的眼睛模糊了呢.
為什麼之後眼睛又能夠開了呢.
發生什麼事呢.
關於耶穌走近的時候.
他們沒有認出.
其中有些可能.
有些人就說.
可能門徒.
忙於彼此談論.
交談議論.
他們不斷討論.
說完又說 說完又說 說完又說.
都是說耶穌不見了.
耶穌死了 不斷地說.
說完之後又無法解釋.
繼續說.
就好像有時我們.
說得是非太多 說得太多.
就令我們看不到神.
因為那些事其實不是最重要.
第二個可能是什麼呢.
在八月方十六章十二節.
可能有另一個.
比較近的解釋就是.
因為耶穌以另一個形象.
去顯現.
因為當耶穌復活之後.
其實是一個新的形象.
這個新的形象.

$^{1241}$不是一般人能夠理解得到.
所以你記得.
莫大拉的瑪利亞.
在約翰福音的時候.
見到耶穌的底下.
一開始都不知道耶穌的.
事實上耶穌這個新的形象.
門徒有時可能不太看得到.
不太明白這就是耶穌.
而第三個可能.
是更加多學者認同的.
可能是什麼呢.
神使他們眼睛迷惑.
神使他們當下這一刻.
眼睛迷惑.
因為他們的心理狀況.
其實是不準備好.
耶穌讓他們認得出.
對他們沒有幫助.
所以我們無論怎樣都好.
我們知道就是.
他們就是看不到.
認不到耶穌.
原來不知道耶穌和他們同行.
所以他們以為.
耶穌是什麼.
記不記得這段經文說.
他們以為耶穌是什麼人.
知道嗎.
是旅客.
經文很清楚說.
當他是一個旅客.
旅遊人士離開耶路撒冷.
離開的人士.
一個旅客.
看不到他是我們生命的主.
接著我們就要問.
為什麼在結束的時候.
他的眼睛又開了.
究竟什麼讓他們的眼睛開了.

$^{1281}$發生了什麼事呢.
按照經文我們可以看到.
在第28節開始.
其實說了有兩件事發生了.
兩件事發生了.
之後就令到.
這兩個門徒開了眼睛.
什麼事呢.
看一下 第一件.
第一件就是當他們走近.
所要去的村子的時候.
耶穌好像還要往前走.
他們的強流他說.
事後晚了天快黑了.
請你和我們住下吧.
耶穌就進去.
要和他們住下.
因為當這兩個人.
離開耶路撒冷.
去以馬五師的時候.
現在特別去到這個村子.
但耶穌似乎還要.
繼續走一些路途.
去另外一個地方.
而當時這兩個門徒.
做了一個很特別的邀請.
就是什麼呢.
不是邀請那麼簡單.
是強流.
強流耶穌.
好像強硬來一樣.
無論如何都要抓住耶穌.
他說事後晚了天快黑了.
請你和我們住下吧.
原來當他們一路.
其實你記得剛才說.
他一路一路走這條路的時候.
他已經心裡火熱著.
他火熱著.
去到和耶穌分手的時候.

$^{1321}$他這團火熱.
令到他切切想留下耶穌.
切切想留下耶穌.
所以說強流耶穌.
和他同住.
這個同住的字.
讓我們想起什麼畫面呢.
想起這幅圖畫.
這幅是什麼圖畫.
葡萄樹與枝子.
是來自.
約翰福音第十五章的段落.
因為.
當原來這幫門徒.
這兩個門徒要.
強流耶穌的時候.
是叫耶穌住下的時候.
同住這個字.
其實就是Abide.
這個字就是.
約翰福音裡面常在的字.
同住就是.
約翰福音裡面常在的字.
約翰福音裡面怎麼說呢.
耶穌說耶穌和門徒的關係.
是什麼呢.
常在的關係才能夠持久.
他說你們.
要常在我裡面.
我也常在你們裡面.
枝子若不常在葡萄樹.
自己是.
不能夠結果子.
如果你們門徒.
若不常在耶穌裡面.
也是這樣.
我就是葡萄樹.
你們是枝子.
常在耶穌裡面的.
常在祂裡面.

$^{1361}$這個人就多結果子.
因為利了我.
你們就不能做什麼.
當你心裡火熱的時候.
還不夠.
當你明白聖經.
明白上帝的深夜的時候.
還不夠.
你仍然要搞清楚.
你和耶穌的那份關係.
而這裡就是說.
你和耶穌的關係是什麼關係呢.
一種常住常在的關係.
要留下耶穌和你一起住.
原來我們生命裡面.
我們就需要和耶穌.
有一份這樣熱切的渴望.
就是耶穌你一定要.
在我這裡.
你一定要與我同在.
我一定要與你同在.
我要堅持無論怎樣.
我都要留你在這裡.
原來這裡就是說.
這兩個門徒經歷了火熱的時候.
忍不住要留住耶穌.
要和耶穌一起.
這就是那份渴望.
就是那份渴望.
要常在耶穌裡面.
是一種怎樣的關係呢.
你有沒有這份.
這樣期盼著.
常常都和耶穌在一起.
這裡很清楚說.
就是.
如果你願意常在.
耶穌裡面.
耶穌一定常在你裡面.
一定常在你裡面.

$^{1401}$所以當這些門徒願意.
讓耶穌同住的時候.
耶穌就留在這裡.
這是第一件.
發生的事.
要求同在同住.
留下耶穌之後.
接著要發生什麼事呢.
有時候我們都會邀請耶穌.
進入我們生命裡面.
但我們還欠一些東西.
就是第二件事.
第二件事是什麼呢.
第三十節.
他坐下來和他們用餐的時候.
耶穌拿起餅袋.
祝福你要抹開第給他們.
是不是很熟悉這些對白.
熟悉得很厲害.
耶穌進去了.
要和他們住下.
就是留下了.
就是叫同在.
那其實耶穌住在哪裡呢.
這個經文沒有說.
究竟為什麼這兩個門徒.
來到二晚五師的時候.
為什麼要去這個地方.
為什麼去這個地方的時候.
會有地方.
邀請耶穌來呢.
其中一個可能.
除非就是.
好像我們去旅行那樣.
租了一間B and B.
我們明知道去哪裡了.
租了一間B and B.
然後耶穌來了.
我們去B and B.
你進來住吧.

$^{1441}$現在就上網預約.
但那時候不是這樣想的.
所以學者說.
其中一個可能.
就是其實這兩個門徒.
回到二晚五師的時候.
這條村莊就有其中一個門徒的屋.
其中一個門徒就是.
慶至蘭山.
離開了耶路撒冷的門徒.
而回到其中一個門徒.
二晚五師這個落後一點的村莊裡面.
的屋裡面.
所以他們能夠邀請耶穌.
來到其中這個門徒的屋裡面.
如果是這樣.
那就很不同了.
那間屋是誰的.
那間屋就是其中一個門徒的.
誰是屋主.
就是其中一個門徒.
是不是.
懂了嗎.
所以如果是這樣.
整件事就變得很不同了.
屋主.
是有屋主的責任的.
但是.
我們看得到.
坐下.
跟他們一起用餐的時候.
是誰拿起餅.
祝福抹開分給人呢.
是耶穌.
原來在猶太人的習慣裡面.
就是誰負責做這件事的呢.
就是屋裡面的主人.
屋裡面的主人.
就是帶領著屋裡面的其他人.
包括賓客.

$^{1481}$包括家人.
一起去抹開那個餅.
然後分給他們吃.
所以你見得到.
在最後晚餐的時候.
就由耶穌去負責.
將餅抹開第二級門徒.
因為耶穌就是.
這十二個門徒的主人.
主的老師.
是最大的.
現在去到這個.
其中一個門徒的屋裡面.
現在耶穌做了.
這個屋主人的工作.
代表什麼呢.
代表這兩個門徒.
一定要認定了耶穌是主.
耶穌是主.
就由主.
邀請他們一同團啟.
有table fellowship.
所以耶穌就和他們.
抹開這些餅分給他們.
這個主人邀請他們一同去坐席.
一同去吃這個晚餐.
耶穌是那個主人.
我們需要的是.
參與他的fellowship.
參與他的團啟.
所以你見得到.
這幫門徒.
不但要心裡面.
有團火.
而且什麼呢.
有一份熱切的渴求.
要耶穌和他們一起.
同在同住.
以致到.
甚至乎要讓耶穌作為主.

$^{1521}$與他們一同團啟.
進入一個.
和耶穌更深的關係裡面.
神就會開他的眼睛.
你和神的關係都不好.
神怎樣開眼睛你都看不到.
唯有你心裡面.
有一份渴望.
甚至乎一種祈求.
以致到你一開眼睛的時候.
你就會經歷到那種.
莫大的興奮.
莫大的雀躍.
弟兄姊妹.
我們究竟有多願意.
我們真的花時間.
讓耶穌常常和我們一起.
甚至乎讓耶穌是我們生命的主.
求他開我們的眼睛.
看得見我們眼前的.
事物.
眼前令你受苦的事.
其實當你睜大眼睛.
看清楚的時候.
你會知道並不一樣.
最後講一個人物.
大家都可能認識一下.
這個人物叫做.
王國倫.
大家聽過嗎.
王國倫.
就是其中一個.
作曲家.
簡單來說.
王國倫.
在他年輕的時候.
都因為他媽媽的原因.
能夠帶他去信耶穌.
帶他去信耶穌之後.
其實他都有不同的經歷體會.

$^{1561}$他都願意去認識神.
其中他有一個心願.
就是什麼呢.
他很想在音樂上.
有好的發展.
他在音樂裡面.
希望有一個好的發展的時候.
因為他之前試過.
參加不同的歌唱比賽.
不同的活動等等.
後來.
有一天.
他經歷到.
有兩個的牧者.
跟他說.
你要等待神.
跟你說話.
以致你才好.
開始再一次.
作一些作曲.
作一些歌.
讓別人感動你.
你才開始作這些歌曲.
於是他就.
因為這個事件.
他自己就.
學習去等候神.
以及一直去.
尋求上帝的心意.
在這個過程裡面.
他其實.
當中他要走音樂的事業.
他做了四年的時間.
希望求神在這四年的時間.
希望能夠闖出名堂.
結果怎樣呢.
第一年他還鬥志很旺盛.
他覺得.
因為他在歌唱比賽裡面.
能夠贏過很多人.

$^{1601}$希望有好的.
音樂事業.
結果.
遇上很多閉門羹.
第二年情況更加惡劣.
沒有收入.
結果就好像麻醉一樣.
經歷上帝的養活.
到了第三年.
他發現很多的.
當兵的同學朋友.
其實都已經開始.
開始儲到一些錢.
買樓 開始賺到很多錢.
但是他仍然.
坐著一輛爛的電單車.
第四年.
就是他跟神說.
給他四年時間的最後一年.
他就不斷問神.
為什麼我的歌.
比外不受歡迎呢.
我現在三餐不計.
我怎樣去度過呢.
雖然在這個.
艱難的時刻裡面.
他仍然去抓緊上帝.
他在這四年裡面.
他仍然追求神.
參加不同的聚會.
一些講座 一些訓練.
去預備自己.
就是這樣 他一直去.
等待 等待 等待.
直至到有一天.
原來神真的要他.
去做事了.
這個時候是怎樣呢.
原來這個時候.
其實那時候他已經經歷了半年.

$^{1641}$沒有任何寫作.
任何的歌曲.
有一間唱片公司.
邀請了他去為一個新人.
一個新的歌手叫黃靜雯.
去寫一首歌.
於是他回到家裡.
就拿起結他 心裡面.
當時聽到兩把聲音.
寫什麼呢 你不會紅的了.
第二把聲音.
很微小的聲音.
他說 跪下.
將自己獻上.
跪下 將自己獻上.
他有一份心裡面的渴求.
他要做一個.
人生最重要的禱告.
他說 神啊.
我這四年來 一無所得.
但如果我還能夠.
在流行音樂創作裡面.
尋求你賜給我的靈感.
讓我用這首歌.
來榮耀你的名的話.
我就要將這首歌奉獻給你.
讓每一個人聽到.
你的美好.
就是這樣.
當他願意親近神.
以神作主.
在神的裡面的時候.
神就開他的眼睛.
也開他的耳朵.
他心裡面.
即時出現了一個旋律.
願意為你.
我願意為你.
聽過這首歌了嗎?.
這首歌是三十多年前了.

$^{1681}$沒聽過?.
碧姐聽過吧?.
碧姐沒聽過.
原來這首.
這首歌.
是三十多年前.
由王國倫這個人寫的.
就是這首歌.
令王妃爆紅了.
王妃.
我願意為你.
神就是這樣.
讓他在腦裡面.
就閃出這個旋律.
他就為神寫了這首歌.
這首歌其實是去獻給神.
希望去為神而寫的.
結果這首歌.
一直都確實.
紅了三十多年.
到今天都仍然有人聽過這首歌.
但是他自己.
很經歷到就是.
當他想起這首詩歌的時候.
他跟神說.
耶穌啊我願意.
我什麼都願意.
結果他自己.
看見這首歌紅得.
很厲害的時候.
他含著眼淚.
但是他聽到主耶穌.
怎樣跟他說呢?.
他說孩子.
我願意這首歌.
其實是我對你唱的.
因為當你還作罪人的時候.
我就為你死.
且死在十字架上.
我的愛就在此顯明了.

$^{1721}$所以黃國磊.
很明白我願意這首歌.
原來我們以為我們願意.
去為神去擺上的時候.
他聽到就是.
神說他願意為人的罪擺上.
為我們犧牲一切.
神有無比的愛.
以致他能夠體會到.
上帝的愛是何等豐盛.
以致到他心裡面.
心裡面火熱.
他知道神是真實的.
他那團火不單止.
在流行音樂裡面去見證神.
他也在過往的日子裡面.
去到不同的地方裡面.
去做見證.
去講神的作為.
所以這幅照片其實就是說.
他應該是.
不知道是今年還是去年.
就是說他來到美國.
做了三十場的見證報導會裡面的照片.
弟兄姊妹.
你現在是眼睛模糊.
還是眼睛開了呢.
求神幫助你.
你自己向神祈禱.
求神幫助你開你的眼睛.
能夠看得到.
最後的總結就是.
這兩個門徒在灰心失望的時候.
就由耶路撒冷.
灰心地離開了的團隊.
一班一起事奉的團隊.
耶路撒冷是一個宣教的中心.
一個宣教的發源地.
但是他們的灰心喪志.
使他們離開了二萬五仙.

$^{1761}$但是耶穌親自.
為什麼在這個下午.
或者這個早上和晚上這個時段.
要在這兩個人身上出現呢.
因為他要將這兩個人.
帶回去到耶路撒冷.
我們有時候都是這樣.
我們灰心喪志.
我們想躲在一個.
只有自己認識自己的地方裡面.
沒有人認識的村落裡面.
我們獨自去過我們的餘生.
但是耶穌不是這樣.
耶穌要將這兩個人.
帶回去耶路撒冷.
所以第三十三節.
我們剛才沒有讀過的.
竟然說什麼呢.
於是他們立刻起床.
沒有再遲疑.
回到家裡應該休息了.
因為他們走了十一公里的路了.
他們叫耶穌留下來.
因為已經慢慢開始夜了.
立刻起床.
回到耶路撒冷.
看見十一個使徒.
和他們正在一起聚集.
在馬可樓聚集的時候.
他說主果然復活了.
你們一直困擾的事情.
現在告訴你.
主真的復活了.
已經顯現給西門漢了.
而且他們將他們路上.
親自怎樣遇見主的事情.
去再說一次.
你留意.
因為.
一個人的見證.

$^{1801}$是不確認的.
在猶太人裡面.
兩個人以上的見證.
是一個有效的見證.
而這兩個的弟兄.
願意一起並肩同行.
回去做一個有效的見證.
去告訴耶路撒冷的人.
我們不用灰心了.
我們有力了.
我們要見到主.
所以在之後的經文.
我們見到耶穌在晚上.
再一次在他們面前.
去印證祂是活的.
祂是真的.
求主幫助我們.
當我們生命裡面.
面對一些心灰意冷無助.
令我們有不同的emoji.
一種低沉的情緒狀況的時候.
我們回到主身邊.
好好地慢慢看看.
上帝的作為.
看看神的說話.
我們再在靈裡面.
神聽 神啊.
你與我同在.
你成為我生命的主.
帶領我吧.
我不知道怎麼走.
不過你帶領我吧.
我們一起祈禱.
主我求你帶領我們.
因為有時我們真的不知道怎麼走.
左右無路 右又無路.
前面又塞死.
總是走不到一條山的路.
自己走不到.
家人走不到.

$^{1841}$無路可行.
但求主你幫助我們.
我們真的回到你面前.
我們學習.
剛才說.
幾位弟兄裡面的生命.
一個又一個.
回到神的面前.
安靜在你面前.
以你作主.
求主你親自.
為我們抹開這個餅.
分給我們.
讓我們得到這份福氣.
我們有力立刻.
起程.
為主你增戰.
我們受用不再.
我們火熱仍在.
祈禱是奉耶穌的名字而求.
阿們.
\newpage



\section{馬可福音 8:31-38}
\label{sec:031954}
\textbf{我要與基督同釘十字架 (粵/華)}~~~~許永勝傳道
\newline
\newline
連結: \href{http://www.church.com.hk/acms/content.asp?site=cdc\&op=show\&type=product\&code=031954\&layout=sermon}{\texttt{031954}} ~~~~ 語音日期: 2026-04-05
\newline
\newline
\hyperref[sec:031953]{\small{< < < PREV SERMON < < <}}
~
\hyperref[sec:index_chronic]{\small{[返順時目]}}
~
\hyperref[sec:031955]{\small{> > > NEXT SERMON > > >}}
\newline
\newline
馬可福音 8:31-38
\newline
\begin{longtable}{cl}
\hline
\hline
章節 & 經文 (和合本修訂版)\\
\hline
8:31 & \begin{tabularx}{0.7\textwidth}{X} 從此，他教導他們說：「人子必須受許多的苦，被長老、祭司長和文士棄絕，並且被殺，三天後復活。」 \end{tabularx} \\ \\ \relax
8:32 & \begin{tabularx}{0.7\textwidth}{X} 耶穌明白地說了這話，彼得就拉著他，責備他。 \end{tabularx} \\ \\ \relax
8:33 & \begin{tabularx}{0.7\textwidth}{X} 耶穌轉過來看著門徒，斥責彼得說：「撒但，退到我後邊去！因為你不體會神的心意，而是體會人的意思。」 \end{tabularx} \\ \\ \relax
8:34 & \begin{tabularx}{0.7\textwidth}{X} 於是他叫眾人和門徒來，對他們說：「若有人要跟從我，就當捨己，背起自己的十字架來跟從我。 \end{tabularx} \\ \\ \relax
8:35 & \begin{tabularx}{0.7\textwidth}{X} 因為凡要救自己生命的，必喪失生命；凡為我和福音喪失生命的，必救自己的生命。 \end{tabularx} \\ \\ \relax
8:36 & \begin{tabularx}{0.7\textwidth}{X} 人就是賺得全世界，賠上自己的生命，有甚麼益處呢？ \end{tabularx} \\ \\ \relax
8:37 & \begin{tabularx}{0.7\textwidth}{X} 人還能拿甚麼換生命呢？ \end{tabularx} \\ \\ \relax
8:38 & \begin{tabularx}{0.7\textwidth}{X} 凡在這淫亂罪惡的世代，把我和我的道當作可恥的，人子在他父的榮耀裡與聖天使一同來臨的時候，也要把那人當作可恥的。」 \end{tabularx} \\ \\
[1ex]
\hline
\hline
\end{longtable}
$^{1}$復活節就是必須要慶祝和紀念的.
但是不管怎樣不守節旗也好.
但是復活節還是必須紀念跟慶祝的.
因為這個是普世信徒.
都是重視的一個基督教節日.
因為這個是普世信徒都重視的節日.
那究竟怎樣去慶祝又怎樣去紀念呢.
那到底要怎樣去慶祝跟紀念呢.
才可以符合慶祝復活節的真正的意義呢.
才能符合慶祝復活節的意義呢.
下面我想給大家看四組的圖片.
下面我想要給大家看四組圖片.
是關於那些人是怎樣去慶祝和紀念復活節的.
是關於人們怎樣去慶祝跟紀念復活節.
第一組圖就是世俗化的慶祝.
第一組圖是世俗化的慶祝.
世俗化的慶祝很大程度.
已經跟耶穌復活變得一點關係都沒有了.
世俗化的慶祝已經大部分都跟.
耶穌復活的紀念都已經沒有任何關係了.
變成享受連續的假期.
家庭樂,購物消費或者是旅遊.
只是變成單純的連續假期.
家庭樂跟購物消費.
第二組圖就是我們最熟悉的慶祝方法.
第二組圖是我們最熟悉的慶祝方法.
就是我們會舉行受苦節和復活節崇拜.
就是我們會舉行受苦節跟復活節的崇拜.
我們教會每年都是這樣做的.
我們將會每年都這樣做的.
其實大家知不知道.
基本上每一個禮拜的主日.
都是紀念耶穌復活的.
其實大家知道.
其實每一個禮拜天的主日.
都是在紀念主耶穌的復活.
因為耶穌在七天的第一天是復活的.
因為耶穌是在七天的第一天去復活的.
也就是禮拜天.
所以主日又被稱為小復活節.

$^{41}$所以主日其實也被稱為小復活節.
為什麼我們每年仍然要特別去慶祝復活節呢.
為什麼我們每年還要特別去慶祝復活節呢.
我又給大家看第三組的圖片.
關於某些國家或是宗派慶祝復活節的方法.
那我們也看一下第三組圖片.
是關於某一些國家跟宗派去慶祝復活節的方法.
為了凸顯復活節的重要.
有些教會就會慶祝.
就是跟紀念一系列的活動.
為了凸顯復活節的重要性.
他們會特別去舉行一些活動.
比如40天的遇苦期或者叫大災期.
復活節前一個禮拜的中之主日.
受難節當然是必然的紀念活動.
還有一些國家或者教會他們為了加深人.
對耶穌受難或者復活的感受.
就會用逼真的話劇去重演耶穌受難.
死亡和復活.
還有一些國家的教會為了加深.
大家對耶穌受難的復活的感受.
他們會用比較逼真的話劇.
重演耶穌受難,死亡和復活.
比如用真的皮鞭去抽打.
比如說用真的皮鞭去抽打.
用真的荊棘帶在扮演耶穌的人頭上.
用真的荊棘放在扮演耶穌的人頭上面.
甚至有人真的用釘去釘穿手掌.
甚至有人用真的釘去釘穿自己的手腳.
給第四祖徒一些個人對復活節的紀念和慶祝.
第四祖徒是個人對復活節的紀念和慶祝.
有人會在這段時間用各種靈修方法去認罪.
有人會在這段時間用不同的靈修方法去認罪.
思考主的受難和復活.
有人會去傳福音或者參加一些短宣.
有人會去傳福音或者參加短宣.
有人會選擇去關心社會上有需要的人.
利用假期去做祝福人的事.
利用假期去做祝福人的事情.
也有人會透過特別的奉獻.

$^{81}$其實每一個信徒都是有自由.
去用自己願意的方法去紀念的.
其實每一個信徒都是有自由.
去用自己願意的方法去紀念的.
紀念為他們受難和復活的主.
不過正如我上次講道.
講信仰也都是好像愛情一樣.
正如我上一次講到.
說到信仰就像愛情一樣.
如果我們的守節不能夠令到神.
感受到我們對祂的真心.
如果我們守節並不能讓神.
感受到我們對祂的真心.
那這樣守節其實沒什麼多大的意義.
那這樣守節其實沒有太大意義的.
弟兄姊妹.
今天我希望跟大家分享的.
不是耶穌有多愛我們.
弟兄姊妹.
我今天想要跟大家分享的.
不是耶穌有多愛我們.
又或者耶穌真的復活了.
又或者是耶穌真的復活了.
而是想分享關於紀念.
耶穌受苦和復活的真正的意義.
而是分享關於.
今天耶穌受難和復活的真正意義.
還有我們怎樣回應主.
才能夠讓主感受到我們那份真心.
那我們要怎樣才能真正回應主那份真心.
所以我為大家預備了一個講題.
是我要與基督同釘十字架.
所以我為大家預備了講題是.
我要與基督同釘十字架.
因為與基督同釘十字架.
是紀念主受難和復活最重要的意義.
因為與基督同釘十字架.
是紀念主受難和復活最重要的意義.
也是我們對主的犧牲.
所能夠作出的最好的回應.

$^{121}$也是我們對主做出的犧牲最好的回應.
經過經文我希望分享.
與基督同釘十字架是最好回應的.
有三個原因.
經過經文.
對與基督同釘十字架是最好的回應.
有三個原因.
第一個原因.
與基督同釘十架是代表真正的跟隨.
和體貼主的心意.
第一個原因.
與基督同釘十字架.
代表真正的跟隨體貼主的心意.
當我們看聖經的時候.
是經常要注意上下文.
當我們看聖經的時候.
是經常需要注意前文後理.
另外經文如果是出現在四福音的.
《馬太》《馬可》和《路加福音》的話.
另外經文如果是出現在四福音的.
《馬太》《馬可》和《路加福音》的話.
同一件事可能很多時候都會同時出現在三卷的福音書裡面.
同一件事情很多時候會出現在這三卷經文裡面.
所以大家可以用對觀的方法去閱讀.
這樣對了解經文是有很大的幫助.
八章三十一至三十八節的內容.
就是耶穌向門徒預告自己的受苦,死亡和復活.
八章三十一到三十八節的內容是.
耶穌向門徒預告自己會受苦,死亡和復活.
因為這是祂作為基督的使命.
不過上文的經文在二十七節到三十節.
耶穌是故意這樣問的.
但是前面的經文裡面提到就是耶穌特別這樣問的.
別人是怎樣看待祂是誰呢?.
又問跟隨自己的門徒看祂是誰呢?.
有人說耶穌是死了的施洗約翰.
有人說耶穌是死掉的施洗約翰.
有人說祂是以利亞.
有人說祂是先知中的其中一位.
唯獨門徒彼得就認出耶穌真正的身份.

$^{161}$唯有門徒彼得認出耶穌的真正身份.
彼得說你是基督.
如果耶穌問祂是誰是一次突擊的考試的話.
如果耶穌問祂是誰是一個突擊考驗.
彼得就應該拿第一名.
彼得應該是第一名.
因為他很清楚地說出.
主就是神所命令的基督.
因為他非常清楚地說出.
主就是神所命令的基督.
在馬可福音沒有顯示耶穌對祂的回答的態度.
在馬可福音裡面沒有描述耶穌對祂回答的態度.
但是馬太福音十六章十七節.
耶穌稱讚彼得的回答是有福的.
但是在馬太福音第十六章十七節.
耶穌是有稱讚彼得的回答是有福的.
今天的經文就是彼得認出耶穌基督這個背景之下發生的.
今天的經文是在彼得認出耶穌是基督的背景下發生的.
耶穌不是單單只希望門徒知道祂的身份.
耶穌並不是單單希望耶穌知道祂的身份.
也希望祂知道作為基督祂的使命和責任是什麼.
也是希望祂知道祂作為基督的使命和責任是什麼.
三十一節耶穌說.
人子必須受許多的苦.
被長老祭司長和民事棄絕並且被殺.
過三天復活.
三十一節說到人子必須受許多的苦.
被長老祭司長和民事棄絕並且被殺.
過三天復活.
知道主是基督的彼得.
聽到主說要受苦被殺.
祂第一個反應是拉著祂勸祂不要.
知道主是基督的彼得.
聽到主要受苦被殺.
祂第一個反應是拉住祂還有勸祂.
耶穌聽見就很生氣地跟彼得說.
耶穌聽到就非常生氣地跟彼得說.
撒旦退我後面去.
因為你不體貼神的意思.
只是體貼人的意思.

$^{201}$因為你不體貼神的意思.
只是體貼人的意思.
按道理彼得才剛剛認出.
耶穌是基督.
以及剛剛被稱讚.
道理上來說.
彼得才剛剛認出耶穌是基督.
還有剛剛被稱讚.
但現在就被主嚴厲地責備.
大家想一想.
那大家想一下.
彼得明明知道耶穌是基督.
那彼得明明知道耶穌是基督.
為甚麼想攔住祂呢.
為甚麼我想要攔住祂呢.
不讓祂受苦受死呢.
那是因為彼得的心裡面.
所想所期待基督要做的事.
是和耶穌說要做的事是完全不同的.
那是因為彼得心裡所想的.
其他的耶穌基督做的事.
跟耶穌說祂將要做的事情.
是完全不一樣的.
當我們查看馬可福音八章之前的經文.
當我們查看馬可福音八章以前的經文.
跟隨耶穌的門徒.
其實是已經見證了.
主行了很多神蹟其事的.
跟隨耶穌的門徒.
其實是已經見證了.
主行了很多神蹟其事的.
包括給四千人吃飽.
醫好耳聾舌結的人.
包括讓四千人吃飽.
治好耳聾舌結的人.
又包括他替婦人的女兒講鬼.
以及給五千人吃飽.
治好血漏婦人等等.
還有替婦人的女兒趕鬼.
讓五千人吃飽.

$^{241}$治好血漏婦人等等.
那麼有能力的基督.
跟隨祂的人生裡面.
一定是這樣想的.
這次發財了.
這樣有能力的人.
跟隨在祂身邊的人.
一定會心想這一次發財了.
國語的大英字幕可能不明白.
這句廣東話這次發財的意思.
國語的弟兄姊妹可能聽不太明白.
這一次發財的真正意思.
這句說話經常用來表達自己.
你做對了決定.
這句話其實常常是用來表達自己.
做對了決定.
這個決定將會為我帶來好處的.
這個決定會為我帶來好處的.
對猶太人來說.
舊約所應許的基督.
那對猶太人來說.
舊約所應許的基督.
是要令到以色列民族得復興.
是要讓以色列民族得復興.
是要成為大衛般的君王.
是要成為像大衛一樣的君王.
那跟隨主的門徒心裡面.
就大概會期望耶穌作王之後.
那跟隨主的門徒心裡面.
大概會期待耶穌成為王之後.
那自己就可以成為開國公神.
自己能成為開國君臣.
最起碼耶穌也可以用神蹟騎士.
為他們解決生活上各種各樣的.
需要和問題.
那最少可以讓耶穌用神蹟騎士.
去解決他們生活上各種需要和問題.
那我問大家.
那我們期待基督要做的事.
是否好像彼得和門徒那樣想的一樣呢.

$^{281}$那我們期待耶穌做的事.
是不是像彼得期待耶穌.
所做的事情一樣呢.
我們每日向主祈求.
就是向主祈求的事.
跟祂說要受苦的事.
受死的事,復活的事有沒有關係呢.
我們每天向主祈求.
是跟祂要受死,受苦,復活的事有關係嗎.
那我們不能夠不承認自己.
其實我們心裡面所期待基督要做的事.
和彼得,門徒期待基督要做的事.
是沒有甚麼分別的.
那我們不得不承認自己心裡面所期待的.
跟彼得,門徒期待基督要做的事是沒有差別的.
那麼我請世祟為我們讀.
約翰福音十四章十二節到十四節.
這邊約翰福音十四章十二到十四節.
我實實在在的告訴你們.
我所做的事.
信我的人也要做.
並且也要做比這更大的事.
因為我往父那裡去.
你們奉我的名無論求甚麼.
我必成就.
叫父嬰兒子的榮耀.
你們若奉我的名求甚麼.
我必成就.
那麼在楊福音三節的經文裡面.
很多人只是關注十四節.
那在約翰福音這三節經文裡面.
很多人都只會關注第十四節.
甚至將這個十四節當成金句去背.
甚至會把十四節當成金句來背誦.
但是十四節不能當成獨立的經文去理解和應用.
但是十四節的經文是不能獨立去理解跟應用的.
因為它是接著十二到十三節說的.
因為它是接著十二到十三節去說的.
耶穌說他要做的事信他的人都要做.
耶穌說他所做的事信他的人也要做.

$^{321}$那麼這句說話的意思.
和今天馬後福音八章三十四節所說的是一樣的.
那麼這句說話的意思.
和今天馬後福音八章三十四節所說的是一樣的.
若有人跟從我.
就當社己背起他的十字架來跟從我.
若有人要跟從我就當舍己.
背起他的十字架來跟從我.
在跟隨主的前提之下.
我們奉主的名去求天父幫助.
在跟從主的情況下.
我們奉主的名求天父來幫助.
天父就會應允我們.
所以耶穌的受難受死和復活.
不是為了滿足我們個人的需要和期待.
所以耶穌的受難受死復活.
並不是滿足我們個人需求而期待的.
而是為了體貼天父的心意.
完成天父所交付他的使命.
而是為了體貼天父的心意.
完成天父所交付的使命.
同樣我們被拯救.
也是要追隨主.
同樣我們被拯救也是需要跟隨主.
體貼神的心意.
完成神給我們的使命.
弟兄姊妹.
耶穌向彼得和門徒透露.
祂將要受苦受難.
耶穌向門徒彼得透露祂將要受苦受難.
是因為不久的將來.
祂們都要跟隨祂們的主.
學 祂們都要學祂一樣.
是因為在不久的將來.
也需要跟隨主 像祂一樣.
主對信去的人的期待.
和我們對主的期待其實相差很遠.
主對我們的期待.
跟我們對主的期待其實相差很遠.
同樣今天慶祝受難節和復活節.

$^{361}$我們究竟在期待什麼呢.
同樣今天慶祝受苦節和復活節.
我們期待的到底是什麼呢.
是不是好像彼得和門徒一樣呢.
是不是跟彼得跟門徒一樣呢.
為了自己的人生可以得到主更加多的祝福.
為了自己的人生可以得到主更多的祝福.
如果是這樣的話.
我們就仍然是體貼著自己的心意.
而不是體貼著神的心意.
如果是這樣的話.
我們依然是體貼著人的心意.
而不是主的心意.
唯獨我們真的願意與主同釘十字架.
這樣在紀念和慶祝受難節復活節的時候.
這樣在紀念跟慶祝受難節和復活節的時候.
我們的心才能真正體貼著天父的心意.
第二個原因.
與基督同釘十字架才是對生命真正有益的.
與基督同釘十字架才是對生命真正有益處的.
我們的主是看透人的心.
我們的主是能看透人的心.
祂知道我們的心裡面在想什麼.
當彼得去攔阻耶穌去送死而被主責備的時候.
當彼得攔阻耶穌去受死而被主責備的時候.
我們都嘗試代入祂當時的思緒和立場.
我們也嘗試去代入祂的思緒和立場.
你會怎樣去回答主的責備呢?.
你會怎樣回答主的責備呢?.
我也想像過這個問題.
如果我是彼得的話我的辯解大概是這樣.
我也想像過這個問題.
如果我是彼得的話我會大概是這樣回答.
主啊!我知道你是基督.
但是你是我所愛的主啊.
主啊!我知道你是基督.
但是你是我所愛的主啊.
我不捨得你受苦啊!不捨得你離開啊!.
我捨不得你受苦啊!捨不得你離開啊!.
難道你要我袖手旁觀你被害被殺嗎?.

$^{401}$我做不到!.
不好意思各位我的演技太差了.
對不起各位我的演技太差了.
因為感情戲是最難的.
我的感情表達可能比彼得還差一點.
我的感情表達可能比彼得還要差.
我們一起看馬太福音十六章二十二節.
彼得就拉著祂勸祂說.
主啊!萬不可如此.
這事必不臨到你身上.
從馬太福音記載的版本.
彼得的話是充滿了不想耶穌受害受苦的感情.
彼得的話是充滿了不想耶穌被害的感情.
是嗎?.
對嗎?.
是嗎?.
如果耶穌不是能夠洞察人心的神呢?.
我相信彼得的情感輸出.
肯定贏得主的稱讚不是責備.
我相信彼得的感情輸出.
肯定能贏得主的稱讚而不是責備.
其實聖經裡面這種情感的表演.
是有不少的例子的.
譬如在撒母耳記上的十五節.
譬如說在撒母耳記上的十五節.
蘇羅為了掩飾自己因為貪念背叛神的行為.
故意表現自己完全是為了獻祭好東西給神.
另外在撒母耳記下十二章五節.
先知約拿丹講了一個跟大衛奪妻有關的故事.
先知約拿丹講了一個跟大衛奪妻有關的筆錄故事.
大衛假裝聽不明白.
就用情緒來表達自己是一個義人.
還用情緒來表達自己是一個義人.
那我們是不是經常有些藉口來辯解我們不遵守神的使命的行為呢.
那我們是不是經常有些藉口來辯解自己不遵守主的行為呢.
譬如做十一奉獻.
誰不知道這是基督徒的責任呢.
但是我們總找到合理的理由去不去履行.
譬如大使命.
有誰不知道這是我們的責任呢.

$^{441}$誰不知道這是我們的責任呢.
但是我們有千萬的理由不做.
再譬如彼此相愛.
有誰不知道主特別吩咐我們做的呢.
有誰不知道這是主特別吩咐我們做的呢.
但是我們總有一些原因讓我們不能夠愛弟兄姊妹.
但是我們總有一些原因讓我們不能愛弟兄姊妹.
回到馬可福音的版本.
沒有像馬太福音那樣記載.
沒有像馬太福音那樣記載彼得豐富的情感.
馬上迎來耶穌擇備彼得.
從經文的35到37節.
耶穌看出彼得極力阻止的原因不是因為愛祂.
耶穌看出彼得極力阻止祂的原因不是因為愛祂.
而是為了彼得自己.
35節的救.
36節的賺.
37節的換.
都是人生中為了想得到的東西而願意付出的努力和手段.
都是人生中為了得到想要的東西而願意付出的手段.
彼得生命中最想得到的是生命中的益處.
彼得生命中最想要得到的是生命中的益處.
這也可能是當時他和其他門徒跟隨主的目的.
我們何嘗不是這樣呢.
我們一生都為了自己想要得到的東西.
不斷去救.
不斷去賺.
不停地用自己的努力去換成果.
大家有沒有試過努力去追求某些東西.
但是想達成某些的成就.
大家有沒有試過努力去追求一些東西.
或者是想要達成一些成就.
但是你所得到的最終是導致跟你心願相反的結果.
但是你最終得到的是跟你所想的是相反的結果.
我舉一個例子.
大家可能知道我父母退休前都是中醫.
大家都可能知道我父母退休前是一個中醫.
他們做了幾十年中醫.
他們做了幾十年的中醫.
我爸爸經常跟病人說.

$^{481}$不要為了賺錢而失去健康.
因為到頭來錢在銀行就人去了天堂.
他真的經常這樣說.
他經常這樣說.
我從小到大有時間都會去診所幫父母做事.
我從小到大有時間都會到診所去幫忙.
其中有一個病人我記憶很深.
其中有一位病人我非常深刻的印象.
他是一位三十幾年前的病人.
他是一個很成功的皮革商人.
他是一位非常成功的皮革商人.
他的生意越做越大.
整天為了談生意要陪客人去應酬喝酒.
有時候他累了或者不舒服.
他就會來到我爸爸的診所做一些治療.
有時候他累了或者不舒服.
就會到我爸爸的診所去做一些治療.
我記得很清楚爸爸不止一次勸他要多注意身體.
我清楚記得我爸爸不是不止一次勸他注意身體.
不要為了賺錢去應酬喝酒傷身體.
不要為了賺錢去應酬傷身體.
他的回答都是一樣.
說自己應酬喝酒是身不由己不可以不做.
他的回答總是一樣.
說自己應酬喝酒是身不由己的不能不做.
但是當他隔了很久沒有出現之後.
有一次來到診所看病呢.
但是當他隔了滿久一段時間沒出現.
之後再來診所看病的時候.
他整個人都瘦得好像排骨一樣.
他整個人瘦得跟皮包骨一樣.
原來他就患了末期肝癌.
原來他患了末期肝癌.
西醫已經放棄了治療.
西醫已經放棄治療了.
他想找我爸爸希望有奇蹟發生.
他找我爸爸希望有一些奇蹟發生.
但是可惜結果他沒多久就病逝了.
但是可惜在沒多久他就病逝了.
我相信好像這樣的個案.

$^{521}$在世界不同的地方都不是少數.
我相信這種個案在世界不同的地方都不是少數.
其實從主對彼得雜備的話來看呢.
其實從主責備彼得的話來看.
主只是講出一個人之常情的道理.
主只是說出一個人之常情的道理.
就是每一個人都為了自己生命中想得到的東西.
或者想重視的東西去努力.
就是每個人為了自己生命中.
想要得到或者重視的事情而去努力.
但是我們所重視的東西.
是不是真的對生命是有益的呢.
但是我們做自己所重視的事情.
是不是真的對生命有益處呢.
大英字媒主說人就是賺得全世界.
賠上自己生命有什麼益處呢.
大英字媒主說人就是賺得全世界.
賠上自己生命有什麼益處呢.
這句話是不是人不可以努力爭取自己想要的東西呢.
這是不是說人不可以努力去爭取自己想要的.
譬如你努力工作希望得到老闆賞識.
升職加薪改善生活.
譬如說你努力工作希望老闆欣賞你.
升職加薪改善生活.
是不是不行呢.
是不是不可以了.
問題的焦點其實不是應不應該賺得全世界.
問題的焦點不是應不應該賺得全世界.
而是賠上自己的生命.
我們活在這個世上當然要和不信的人一樣.
為生活努力奮鬥.
我們活在這個世界上當然要跟不信的人一樣.
為了生活而努力奮鬥.
但問題是我們知道世上的生命是短暫的.
但是問題是我們知道世上的生命是短暫的.
真正對生命有價值和重要的是關乎永恆的生命.
真正對生命的價值和重要的是關乎永恆的生命.
所以如果我們只是願意努力去追求短暫生命上的東西.
所以如果我們只是願意去努力去追求短暫生命上的東西.
但不願意為永恆生命去努力.

$^{561}$但不願意為永恆生命而努力.
那這個就是耶穌所說你賺得全世界最後卻賠上自己生命的意思.
這就是耶穌所說你賺得全世界最後卻賠上自己生命的意思.
所以耶穌說為我和福音喪掉生命的必救了生命.
與耶穌同釘十字架.
這種捨己的態度和行為.
就是代表我們看重神的使命比自己生命所想要的一切更重要.
主說為了祂的福音是因為主和福音是同義詞.
主說為了祂和福音是因為主和福音是同義詞.
主就等同福音福音就是主.
主就是福音福音就是主.
今天藉著復活節的日子.
我們跟隨主的人要思考自己為了福音為了主究竟願意付出多少的努力呢.
我們跟隨主的人要思考自己為了福音為了主願意付出多少的努力呢.
比起我們在世上想追求生活而付出的努力.
比起我們在世上追求生活所付出的努力.
那為了福音我們又做過什麼努力呢.
最後第三個原因.
與基督同釘十字架代表已主的復活為榮.
最後一節我們一起讀三十八節.
凡在這淫亂罪惡的世代.
把我和我的道當作可恥的.
人子在他父的榮耀裡.
同生天賜更臨的時候.
也要把那人當作可恥的.
我曾經在網上聽到一位外國的牧師講道.
那他問了惠中一個問題.
我覺得很有趣.
我覺得蠻有趣的.
他問你們知不知道世界上最多臥底聚集的地方在哪裡.
他說你們知不知道世界上最多臥底聚集的地方是在哪裡呢.
大家猜猜.
大家猜一猜.
對 就是教會.
為甚麼呢.
大家覺得牧師是想稱讚還是想諷刺信徒呢.
大家覺得牧師這句話是想要稱讚還是諷刺.
稱讚的舉手.
沒有人吧.
這句話的意思就是基督徒在世上好像臥底一樣將自己信徒的身份隱藏起來.

$^{601}$這句話的意思其實是基督徒在世界上像臥底一樣把自己的信徒的身份隱藏起來.
努力地融入這個世界.
令到身邊的人覺得你是跟他們一樣.
讓身邊的人覺得你是跟他們一樣的.
經文中淫亂罪惡的世代對比同性仙事降臨的時候.
很多詩篇書都會解釋焦點放在世界的墮落.
很多詩篇書都會把焦點放在世界的墮落.
強調世上沒有一個異人都需要主的救贖.
這樣的解釋表面聽起來是不錯的.
但是跟經文的信息是沒有關係的.
但是跟經文的信息是沒有任何關係的.
其實最後一節經文耶穌是很準確地講出彼得當時對於耶穌受難受死復活背後的心態.
其實最後一節經文是說出耶穌.
耶穌非常精準地說出彼得對於耶穌受難受死和復活背後的心態.
因為在馬可福音的十四章就記載耶穌預言彼得三次.
將有三次在人面前不認祂.
因為馬可福音第十四章就有記載.
耶穌預言彼得將會有三次在人的面前不認祂.
經文說主的道就是耶穌的受難受死和復活.
是代表了神對人的拯救屬.
是代表了神對人的救贖.
是關乎幕後主在臨.
是關乎迎接榮耀且永恆的生命.
是關乎迎接榮耀並且永恆的生命.
赴的榮耀裡同聖天使降臨的時候.
是主的道價值所在.
因為耶穌的受難受死復活將會為我們帶來無限的祝福和榮耀.
因為耶穌的受難受死和復活是會為我們帶來無限的祝福和榮耀.
問題是彼得當時根本不明白耶穌的道有多麼寶貴.
他心裡已經選擇在淫亂罪惡的世代去追求自己所認定的福.
他心裡其實已經選擇在淫亂罪惡的世代追求自己所認定的福.
第一姊妹你覺得自己是不是臥底的基督徒呢.
第二姊妹你覺得自己是不是臥底的基督徒呢.
今天是復活節的主日.
我們是不是真的明白耶穌為我們受苦受死和復活的真正目的和意義呢.
你知不知道主的復活能為我們生命帶來多麼寶貴的祝福嗎.
你知不知道主的復活為我們生命帶來多麼寶貴的祝福嗎.
如果我們真的明白.
我們每年慶祝和紀念復活節的心態就不會完全一樣了.
如果我們真的明白的話我們每年慶祝復活節的心態就不會完全一樣了.

$^{641}$最後我想用保羅在《肥肋比書》三章七到十一節一番信仰告白作為總結.
那麼我請西施幫我讀出.
只是我先前以為與我有益的.
我現在因基督都當作有損的.
不但如此我也將萬事當作有損的.
因我以認識我主基督耶穌為自保.
我為祂已經丟棄萬事看作糞土.
為要得著基督並且得以在祂裡面不是有自己因律法而得的益.
乃是有信基督的益.
就是因信神而來的益.
使我認識基督.
曉得祂復活的大能.
並且曉得和祂一同受苦.
效法祂的死.
或者我也得以從死裡復活.
如果我們真的願意我要與基督同釘十字架.
作為對耶穌為我們受苦受死受難的回應.
那麼我邀請大家和我一起.
那我邀請大家跟我一起.
以保羅的告白為榜樣.
以主的救贖為自保.
不以福音為恥.
在世上不是做臥底的信徒.
在世上不做臥底的信徒.
而是勇敢的和別人去傳福音.
而是勇敢的跟別人傳福音.
我相信如果我們這樣的心態去紀念復活節的話.
我相信我們用這種心態去紀念復活節的話.
主一定能感受到我們這份真心.
我們一起祈禱.
我們一起禱告.
親愛的主.
我們為到今天的復活節崇拜.
我們先上感恩.
主啊我們不是紀念.
把紀念復活節成為一個習慣.
成為一個儀式.
主啊我們真的是知道你為我們受苦受死.
可復活的意義有多重要.
你的拯救是多麼的寶貴.

$^{681}$因為你的拯救使我們有永生的盼望.
主啊我們在這個世界所追求的都是短暫的.
唯獨我們所盼望的那份永恆的生命才是最寶貴.
主啊就讓我們每一位弟兄姊妹都得到這個寶貴的祝福.
也讓我們身邊我們還沒有認識你.
還沒有信靠你的人.
我們有勇氣.
有人力.
有決心去跟他們分享福音.
讓今天我們的紀念我們的真心主你都感受到你都知道.
我們這樣子禱告奉主耶穌基督的名求.
Amen.
\newpage



\section{詩篇 25:1-22}
\label{sec:031955}
\textbf{建立健康的使命教會 (11) 辨識並遵從神的引導}~~~~江山偉牧師
\newline
\newline
連結: \href{http://www.church.com.hk/acms/content.asp?site=cdc\&op=show\&type=product\&code=031955\&layout=sermon}{\texttt{031955}} ~~~~ 語音日期: 2026-04-12
\newline
\newline
\hyperref[sec:031954]{\small{< < < PREV SERMON < < <}}
~
\hyperref[sec:index_chronic]{\small{[返順時目]}}
~
\hyperref[sec:031956]{\small{> > > NEXT SERMON > > >}}
\newline
\newline
詩篇 25:1-22
\newline
\begin{longtable}{cl}
\hline
\hline
章節 & 經文 (和合本修訂版)\\
\hline
25:1 & \begin{tabularx}{0.7\textwidth}{X} 耶和華啊，我的心仰望你。 \end{tabularx} \\ \\ \relax
25:2 & \begin{tabularx}{0.7\textwidth}{X} 我的神啊，我素來倚靠你；求你不要叫我羞愧，不要叫我的仇敵向我誇勝。 \end{tabularx} \\ \\ \relax
25:3 & \begin{tabularx}{0.7\textwidth}{X} 凡等候你的必不羞愧，惟有那無故行奸詐的必要羞愧。 \end{tabularx} \\ \\ \relax
25:4 & \begin{tabularx}{0.7\textwidth}{X} 耶和華啊，求你將你的道指示我，將你的路指教我！ \end{tabularx} \\ \\ \relax
25:5 & \begin{tabularx}{0.7\textwidth}{X} 求你指教我，引導我進入你的真理，因為你是救我的神。我整日等候你。 \end{tabularx} \\ \\ \relax
25:6 & \begin{tabularx}{0.7\textwidth}{X} 耶和華啊，求你記念你的憐憫和慈愛，因為這是亙古以來所常有的。 \end{tabularx} \\ \\ \relax
25:7 & \begin{tabularx}{0.7\textwidth}{X} 求你不要記得我幼年的罪愆和我的過犯；耶和華啊，求你因你的良善，按你的慈愛記念我。 \end{tabularx} \\ \\ \relax
25:8 & \begin{tabularx}{0.7\textwidth}{X} 耶和華是良善正直的，因此，他必教導罪人走正路。 \end{tabularx} \\ \\ \relax
25:9 & \begin{tabularx}{0.7\textwidth}{X} 他要按公平引領謙卑人，將他的道指教他們。 \end{tabularx} \\ \\ \relax
25:10 & \begin{tabularx}{0.7\textwidth}{X} 凡遵守他的約和他法度的人，耶和華都以慈愛信實待他。 \end{tabularx} \\ \\ \relax
25:11 & \begin{tabularx}{0.7\textwidth}{X} 耶和華啊，求你因你名的緣故赦免我的罪，因我的罪重大。 \end{tabularx} \\ \\ \relax
25:12 & \begin{tabularx}{0.7\textwidth}{X} 誰敬畏耶和華，耶和華必教導他當選擇的道路。 \end{tabularx} \\ \\ \relax
25:13 & \begin{tabularx}{0.7\textwidth}{X} 他要安然居住，他的後裔必承受土地。 \end{tabularx} \\ \\ \relax
25:14 & \begin{tabularx}{0.7\textwidth}{X} 耶和華與敬畏他的人親密，他要將自己的約指示他們。 \end{tabularx} \\ \\ \relax
25:15 & \begin{tabularx}{0.7\textwidth}{X} 我的眼目時常仰望耶和華，因他必將我的腳從網裡拉出來。 \end{tabularx} \\ \\ \relax
25:16 & \begin{tabularx}{0.7\textwidth}{X} 求你轉向我，憐憫我，因我孤獨困苦。 \end{tabularx} \\ \\ \relax
25:17 & \begin{tabularx}{0.7\textwidth}{X} 我心裡愁苦甚多，求你救我脫離我的禍患。 \end{tabularx} \\ \\ \relax
25:18 & \begin{tabularx}{0.7\textwidth}{X} 求你看顧我的困苦、我的艱難，赦免我一切的罪。 \end{tabularx} \\ \\ \relax
25:19 & \begin{tabularx}{0.7\textwidth}{X} 求你察看我的仇敵，因為他們人數眾多，並且痛恨我。 \end{tabularx} \\ \\ \relax
25:20 & \begin{tabularx}{0.7\textwidth}{X} 求你保護我的性命，搭救我，使我不致羞愧，因為我投靠你。 \end{tabularx} \\ \\ \relax
25:21 & \begin{tabularx}{0.7\textwidth}{X} 願純全、正直保護我，因為我等候你。 \end{tabularx} \\ \\ \relax
25:22 & \begin{tabularx}{0.7\textwidth}{X} 神啊，求你救贖以色列脫離他一切的愁苦。 \end{tabularx} \\ \\
[1ex]
\hline
\hline
\end{longtable}
$^{1}$各位影子妹,平安!.
十幾年前,有一套電影,不知道大家看過了嗎?.
這套電影叫做《功夫熊貓》.
有沒有人看過?舉手吧.
證明我們都是那個年代的,對吧?.
這套電影,當我第一次看的時候,其實對我有很多啟發.
沒看過的,我很快地說吧.
故事裡其實是講及這隻熊貓,叫阿寶.
阿寶很想去學功夫.
夢想是很想成為一個武林的高手.
當他拜師學藝,找到一個很好的師父.
他心裡很快地,很希望很快地能夠學到.
原來大家在座有沒有人學功夫呢?.
問一下,有沒有人學功夫的?.
懂得幾招的.
如果你大概知道的話,你會知道學功夫需要一段很長很長的時間.
阿寶很希望能夠學到一個秘密的武功.
很想學到一招半式去持平.
當時他的環境裡,都有很多作惡的事情,很想去抱不平.
但你也知道他的身形,就像這樣.
所以師父教他一些很普通,甚至是一些很悶的基本功.
練習平衡,練習他的站姿,耐力.
所以阿寶一開始的時候,其實很困惑,很失望.
心裡覺得自己是來學武功的.
你教來教去都是一些很枯燥的練習.
但師父知道一件事,如果沒有基本功的話.
其實再厲害的招式,他也不能夠發揮出來.
其實當我一直看這部電影的時候,我都在想.
因為當時我剛剛出道,做傳道人.
我就在想,屬靈的生命其實都是這樣.
當我們信了耶穌之後,很快我們都會跳到一個問題.
我怎樣能夠聽到神的聲音呢?.
怎樣知道神的心意呢?.
其實我們都會很容易,很想學到一些招式,技巧.
但其實我們很多時候都忽略了一件很重要的事情.
當我們說我們要去學習去辨識神的帶領的時候.
其實我們必須要先建立一些屬靈生命的基本功.
當我們說我們很想去尋求上帝的旨意,準確性的時候.
其實是跟一個人的屬靈的質素是有關的.
而不是在於他用什麼方法.

$^{41}$所以今天很想跟大家去想這兩樣東西.
第一,就是我們怎樣去辨識神引導的屬靈的心法.
就好像阿寶一樣,不是教你什麼武功招式.
先從心法開始.
第二方面,我們怎樣去遵從神在我們生活中的引導呢?.
所以第一樣東西我們很想你去看.
辨識神引導的屬靈心法.
其實在剛才的詩篇二十五篇.
其實在裡面都帶給我們一些很寶貴的訊息.
如果你有閱讀過詩篇二十五篇,其實就是大衛的一篇祈禱詩.
其實是一篇個人的哀求詩.
從內容來看,大衛當時正可能面對著很多逼迫,敵人的攻擊.
生命都受到很大的威脅.
在外在的環境裡面,甚或乎可能遇到一些危險,困境.
在他的內心其實都帶著很多的羞愧和罪的重擔.
所以大衛就來到神的面前.
一方面求神保護他,拯救他.
另一方面他都求神指教他當行的路.
所以心法的第一.
我們需要從一個起點開始,我們怎樣去轉向神.
剛才Eliza都為我們讀出這五節.
大衛一開始的時候其實不是要去分析目前的形勢.
也不是要去衡量得與失.
也不是要依賴過去自己的經驗.
他做的第一件事,他就是提到我的心仰望神.
換句話說,真正的屬靈的分辨.
是要由我們的心先轉向神來開始.
很多時候我們都會祈禱.
我們都會有很多代討的清單.
好像品神一樣,陰一陣,一直在說.
其實我們是不是在尋求神的帶領呢?.
還是其實是在尋求神的承諾.
給你一個印,神啊,給你一個印.
拜託你跟著這個我想的方法去應許.
我們不是求問神,究竟你的路是什麼.
而是我們不斷地在想.
神啊,求你祝福我的路.
祝福自己心裡面所想.
大衛在這五節裡面,不單止他的心轉向神.
同樣一路看他的禱告的時候.

$^{81}$他講了這幾樣東西.
他說求神指教我,指示我,引導我.
我整天等候你.
大衛不是先有自己的答案.
然後去找神去同意自己.
而是他願意先放下自己的想法.
讓神一步一步帶領他走前面的路.
有很多時候我們開車.
我們都很依賴GPS的,是不是?.
即是我自己還是,OK.
GPS有時候它帶你走一條.
可能你覺得,嘩,為什麼繞這麼大的彎.
這樣,是不是?.
即是比較迂迴,是不是?.
當有時候看的時候,是不是?.
你心裡面可能都會想.
唉,走這條不是更快嗎?這樣,是不是?.
於是乎我們很多時候都會,唉,不太相信它.
就走我心目中想走的路.
這條所謂的捷徑.
不過可能你開車的時候就會發現.
咦,原來前面又塞車,又收路.
根本就像封了路,是不是?.
那時候可能你才明白.
原來這個GPS導航.
其實已經一早考慮了.
一些我們可能預計不到的情況.
所以我們的生命其實都好像是這樣的,是不是?.
我們其實都很想神.
唉,你給我一條捷徑就好了,是不是?.
你帶我走一條最快,最容易,最順利的路.
不過很多時候,神的帶領.
往往都是一條比較.
走得比較遠一點點.
在人的角度看,為什麼要這樣走,是不是?.
當神的道路.
跟我們心目中的捷徑不一樣的時候.
不知道你會怎樣選擇呢?.
我們是否相信自己的判斷.
還是我們的心.

$^{121}$好像大衛一樣.
轉向神.
定義的跟隨神的帶領.
神的引導.
不只是讓我們知道方向.
神的引導更加需要一件事情.
就是我們願意.
交出這個方向盤.
今天.
第一件事我們要去問的.
就是我們的心.
是不是真正地仰望神.
第二方面.
就是屬靈分辨.
一個的基礎.
就是親近神.
如果我們一直看.
再看詩篇第八節和第九節的時候.
我們一起讀出來,好嗎?.
預備,一二三.
耶和華是良善正直的.
因此他必教導罪人走正路.
他要按公平引領謙卑人.
將他的道指教他們.
這兩節提醒我們一件很重要的事情.
神不是隨便將祂的道路.
顯明給所有人.
神得必引導的是誰呢?.
不是那些最聰明的.
不是那些最有經驗.
也不是最有能力.
而是甚麼?.
而是謙卑的人.
為甚麼?.
因為謙卑的人.
我們都知道.
如果你是謙卑的人.
你都知道的.
你有些特質.
你的特質是甚麼呢?.

$^{161}$你是有一個願意受教的心.
只有謙卑的人.
才會承認自己未必知道所有東西.
只有謙卑的人.
他不會將自己的想法.
當作是真理.
今天我們最大的困難.
往往不是我們知不知道.
而是我們很多時候都不肯放下.
很多時候我們.
其實都隱約知道神的心意.
或許我們都知道.
某一段的關係需要修補.
某一個的態度需要去修改.
某一個的決定.
可能其實是出於自己的私慾.
或者某一步應該要停一停.
去想一想.
再去禱告.
但很多時候我們都不聽.
因為神的路.
有時候未必是我們所想的.
所以謙卑的意思.
就是要承認.
承認甚麼?.
承認神比我們更加清楚.
神比我們看得更加遠.
神的道路.
神的意念.
都是高過我們的道路.
和我們的意念.
詩篇在二十五篇繼續看的時候.
第十二和第十四節這樣說.
它說「遂經威耶和華」.
「耶和華必指示他當選擇的道路」.
「耶和華予敬畏他的人親密」.
「他必將自己的弱指示他們」.
我們這兩節經文又再一次看到.
神的引導.
其實是有一個關係的.

$^{201}$不是一個指令式的.
很多時候我們以為神的大令.
就是要給我一個答案.
去還是不去.
應該不應該.
做還是不做.
其實神更加關心的.
很多時候不是要你才知道的答案.
而是要你不是和祂建立一個緊密的關係.
十四節這裡說.
「耶和華予敬畏他的人親密」.
原文有一種很深的意思.
好像朋友之間.
一個你想想.
一個你很熟悉的人.
那種親密的對談.
那種沒有防衛的.
很坦誠的去分享.
換句話說.
神不單止要給我們一個方向.
神更加很想我們和祂一起去同行.
我們固自然當然很想知道神的旨意.
但是有時候原來我們其實都不太想花時間去親近神.
我們很想神帶領我們前面的道路.
不過我們又不太想和神一起去同行.
我不知道你會如何形容此刻你和神的關係.
你覺得它是否很親密.
你和上帝的關係是否親密.
因為你對神其實都很冷淡.
有事就忠無厭.
究竟你會如何形容你和上帝的關係呢.
就是很親密.
如膠似漆.
每天沒有祂就不行.
我一直在想.
其實今天你和什麼最親密.
你知不知道.
有沒有想過.
你和你的電話最親密.
我告訴你.

$^{241}$基本上你每天都機不離手.
當你離開了電話.
你就周身不聚財.
好像那天吸毒那樣.
心都在亂了.
我不知道這是不是.
我也覺得這是事實.
當然我們不只是電話.
當然現在有很多用途.
現在智能電話.
的確有很多很重要的東西在.
但是其實我們又有多少時間.
與神親近呢.
透過你的智能電話.
其實我們提供了很多的零售材料.
在Church App裡面.
你有多可按下去看看.
按下去聽聽.
透過這個方式來親近神呢.
或許你和某一個頻道的YouTuber比較熟悉.
所以我希望.
透過聖經裡面再一次提醒我們.
回到一個基本功.
今天.
我們和上帝的關係.
有多接近.
真正明白神道路的人.
聖經裡面說得很清楚.
就是那些敬畏神.
親近神的人.
就好像.
和你的配偶一樣.
很多時候其實你還沒說完.
他已經明白你在說什麼.
很多時候師母.
當我和她對話的時候.
有時候我說一兩個字.
她已經知道.
不是說她懂得什麼讀心術.
什麼特別的技巧.

$^{281}$而是因為.
當我們相處越久.
溝通越多.
關係越親密的時候.
就越來越熟悉對方的心意.
屬靈的分辨都是一樣.
當我們越親近神的時候.
其實就會越容易去認出.
上帝的聲音.
就好像這幅操場一樣.
很多小朋友.
媽媽在魂魂裡面叫.
「林小明」.
在座有沒有「林小明」?.
沒有.
純粹去求而已.
「林小明」.
那個小孩.
當那個是他的兒子的時候.
很快他就認出這是我媽媽的聲音.
又要叫我走.
很多很不情願.
為什麼他能夠在這麼吵的環境裡面.
都能夠認出他媽媽的聲音呢?.
因為很熟悉.
很多時候.
為什麼我們聽不清楚神的聲音?.
其實不是神沒有和我們說話.
只不過我們不熟悉祂而已.
我們聽了很多.
剛才所說的.
很多不同的頻道.
不同的聲音.
那些你也很清楚祂在說什麼.
這個星期這個YouTuber又說這句話.
真棒.
但是很多時候.
我們卻是聽不到.
原來神已經在和我們說話了.
這種的熟悉.

$^{321}$不是一朝一夕的建立.
而是很需要每天.
我們如何與神同行.
慢慢的建立.
送靈的分辨.
也都需要一件很重要的事.
就是我們的生命.
需要被結織.
在這篇詩篇裡面.
大衛其實多次這樣.
求神赦免自己的罪.
大衛不要假裝自己沒有問題.
聖人不是.
當他一路去尋求神的時候.
他自知.
他需要來到神的面前.
去認罪悔改.
在這裡也給我們提醒一件事情.
有時候我們分辨不了神的帶領.
不單止是方向的問題.
其實更加關鍵就是.
我們生命裡面的問題.
可能是裡面一些的驕傲.
苦毒.
自我忠心.
隱藏了的罪.
不願意去饒恕.
不肯去馴服.
其實這些都會影響我們.
那個屬靈的敏銳度.
罪我們知道.
上帝是聖潔的.
當我們有一些污穢的時候.
其實是會難了我們去親近神.
罪會使我們屬靈的感官變得遲鈍.
就好像這一幅鏡子一樣.
如果鏡子的面積滿了灰塵.
固知人當然看不清楚.
是不是這塊鏡子失去了功能呢.
不是.

$^{361}$而是它的表面.
被那些塵遮住了.
同樣如果我們的心被罪.
被施辱或者苦毒遮蓋的時候.
往往都會令到我們很難去看清楚.
神的道路.
最後.
辨識神引導的屬靈心法.
我們怎樣去實踐呢.
我們需要去依靠神.
很多時候我們都很希望.
神向我們顯明祂的道路.
在人生的旅程裡面.
不一定每一步都很清楚.
所以屬靈的分別不單止是尋找答案.
更加重要就是要在一個不確定的裡面.
持續地學習去依靠神.
真正的關鍵不是我們知道多少.
而是我們願不願意.
將生命交給神.
讓祂來帶領.
詩人在二十五篇第十節裡面說.
凡遵守祂的約和.
祂法度的人.
也和華都以慈愛信實待他.
這次的經文提醒我們.
神帶領與人的信實.
是息息相關.
當我們願意照著神的話語去生活的時候.
神就以經文說.
以慈愛信實與我們同行.
在第十五節裡面又說.
我的眼目是上仰望耶和華.
因祂必將我的腳從網裡拉出來.
人生的道路.
其實都充滿著很多的網絡和陷阱.
大衛在這裡提醒我們.
屬靈的分別是需要一種.
持續的態度.
事事仰望神.

$^{401}$尤其是當我們面對一些困惑.
一些選擇.
十字路口裡面.
我們更加需要不斷地.
將目光轉向神.
經文裡面提到.
神就會在我們看不見的地方裡面.
去保守我們的腳步.
另外在二十和二十一節裡面.
祂說求你保護我的性命.
你答救我.
使我不自羞愧.
因為我投靠你.
願純全正直保護我.
因為我等候你.
在這裡也描繪了一種屬靈心靈的態度.
投靠神.
內心地來等候神.
有時候.
很多時候.
那個的帶領.
都不明顯的.
在座坐在這裡的每一位.
我們都經歷過.
也都未必是符合我們的時間表.
我們固然當然想快.
好像吃這個即食麵那樣.
三分鐘你就給我可以吃.
但是屬靈的成熟.
就是我們去學習.
怎樣願意在一個不確定的裡面.
我們仍然去信靠神.
就好像經文裡面說.
讓純全和正直成為我們生命的保護.
我們呢.
我都想起.
當一路想的時候.
其實都好像我們很想.
當然想.
神好像一個探射燈.

$^{441}$一射就射遠一點.
最好呢.
將未來的五年.
你想不想五年之後.
你自己變成怎樣.
你想不想知道.
對不對.
五年之後.
純全家變成怎樣.
就是.
你很想知道.
越清晰越好.
不過很多時候.
神就不是好像一個探射燈那樣.
其實是一個電筒來的.
往往的帶領就好像電筒.
只是照亮眼前的一步.
你就走這一步.
不要走太多不要走太少.
當你走了這一步的時候.
前面的路又會再照亮一點.
當再走下一步的時候.
神又會帶你看多一點.
想起過去這四年移民.
來到多倫多這個經歷.
其實都是一種走步見步的過程.
來到加拿大之前.
其實我們一家都知道了.
其實我們當時是未符合.
這個救生艇計劃的申請資格.
直到來了之後一年.
政府的政策改變.
我們終於可以申請了.
作為一個新移民.
其實都面對著很多未知的.
對不對.
住又怎樣.
家裡怎樣安頓.
究竟我們國人能不能夠適應.
經濟搞不搞定.

$^{481}$這裡會不會成為我們.
未來的惡劣禍等等等等.
很多的問題都是沒有答案.
神沒有一次將所有的答案.
告訴我.
告訴我們一家人.
而是一步一步的帶領.
當我們願意走出眼前一步的時候.
神就為我們預備下一步.
所以回顧這幾年的裡面.
其實我們一家都很深體會.
神的手是一直拖帶著我們.
去走前面的路.
這就是屬靈生命上.
上經歷的一個節奏.
行步見步.
我們固然當然想.
見步 行步 好吧.
但是原來這個屬靈的法則.
往往是相反的.
行步見步.
辨識神的引導其實不是技巧.
而是我們怎樣一起與神同行的一個過程.
先轉向神.
每天都去親近神.
並且在每天的生活學習去依靠神.
第二件事我們很想的.
每天都一起看.
我們怎樣能夠遵從神在生活的裡面.
給我們的引導呢.
當我們面對人生.
大家研究顯示.
每天我們其實有三萬多個抉擇.
大大小小.
即是當一些你自己會懂得分辨.
又懂得過濾的.
那些很重要的東西.
你就會來到神的面前.
究竟我們怎樣去.
有一些錦囊.

$^{521}$去分辨呢.
第一件事.
我覺得我們很需要去看.
那個決定究竟是否合乎神的話語.
神給我們的教導.
神永遠不會帶領我們.
做一些違背祂的事情.
違背祂的法則.
違背祂一直我們所學習.
聖經裡面給我們的教導.
因為神的靈和神的話.
永遠都一致的.
很多時候我們就說.
聖靈的感動.
必須要對應這個感動.
和你在聖經裡面看.
是否一致的呢.
聖靈不會感動一件.
和聖經真理相矛盾的事情.
如果那個選擇明顯是違背了.
神的原則.
其實我們就不需要再問.
究竟這個是不是神的心意.
例如你要做一件事情是不誠實的.
不公義的.
不聖潔的.
傷害到其他人的.
或者只是為了滿足自己.
這個很肯定.
都不是神的帶領.
聖經雖然沒有在.
每日我們生活的每一個處境.
提供一個很具體細節的指示.
所以呢.
當我們實際應用神的話語的時候.
可能有時候都會有一些不確定性的.
不過只要我們沒有忽視到.
拒絕神的教導.
而是我們正在學真誠地願意去信服.
其實在這個尋求的過程裡面.

$^{561}$就已經是.
你是在行神的心意.
與神同行.
聖經更加提醒我們.
就好像帕薩利奧.
就是幾個星期前他提到.
就夜思想神的話語.
讓神的話語.
不斷地來去塑造我們的心思和價值觀.
當我們持續地去.
就夜思想神的話語的時候.
在愛和智慧裡面作出抉擇.
即使那個處境不斷地來改變.
其實我們仍然在.
每日的生活和決定裡面.
都是在學習如何去活出神的心意.
一個人如果他沒有.
與神有一個深切的認識的話.
即是他固然就很想.
幾日之內去知道神的心意.
其實往往都沒有結果.
相反.
當你每一日.
你去到很深切去認識.
很認真地去讀神的話語.
其實我們很容易會找到那個答案.
第二樣東西.
就是這個的決定.
是否幫助我們更加去靠近神呢.
有些選擇可能是挺好的.
不過實際上可能會令到我們.
慢慢來到去遠離神的時候.
我們就需要留意.
即可能這個的抉擇.
帶來給我們.
會吞噬了我們的時間.
心力.
甚至乎影響我們的敬拜.
影響我們的家庭.
影響我們個人屬靈的生活.

$^{601}$我們就必須要小心.
我們不是只是問.
究竟這件事做不做得了.
做不做得了 會不會成功呢.
更加重要的問題就是.
究竟這個選擇.
會不會使我們更親近神.
抑或是會令到我們更加分心.
過去我也有機會跟一些弟兄.
在香港的時候.
你也知道.
香港的一些獵頭公司.
尤其是在高峰的時候.
每個星期都會打電話給你.
有些弟兄來跟我說.
有一份很吸引的人工.
薪水好.
即全道理覺得.
我應不應該去呢.
當時跟他一邊聊天.
一邊分析.
這份工作.
原來星期天你又回不了教會.
接著又會跟你的家人的相處時間.
又沒什麼.
是人工好.
不過你的犧牲可能會更大.
正如《聖經》裡面這樣說.
你贏得全世界卻賠上了的生命.
值不值得呢.
很感恩這個弟兄.
最後他也在實驗裡面.
我們一起去祈禱.
他選擇也是不要剔.
究竟今天.
當我們要去作出一個屬靈的分辨的時候.
一個很重要的指標.
就是這個選擇的背後.
究竟是令我們更加去愛神愛人.
更不更加像基督.

$^{641}$還是相反呢.
第三方面.
就是我們能不能夠在一個屬靈的群體裡面得到引證.
神的帶領很多時候.
不單只是透過我們個人的感覺來顯明.
神也會很經常透過屬靈的群體.
幫助我們去分辨.
過去我們也不斷地鼓勵大家.
在你的教會裡面.
有沒有一位屬靈的摯友呢.
能不能夠你可以跟他傾心道義.
一起去禱告,去守望呢.
其實這個也是上帝透過一個這樣的方式.
在聖經裡面.
你看耶穌和門徒那種關係的裡面.
其實也給我們一個提醒.
一種彼此去守望的途徑.
因為很多時候當我們自己來到作決定的時候.
很容易會受到情緒,壓力.
或者個人的期待影響.
可能我們會太急,太主觀.
甚至乎只聽到自己的聲音.
但是在一個屬靈的群體.
或者在一個你信任的弟兄或者姊妹裡面.
有時候可能他一句話.
你這樣做會不會有一個這樣的後果呢.
可能這一番話.
或者他為你的代討去守望的時候.
往往讓你都能夠可以察覺到一些的盲點.
可能你之前沒有想過.
在朕信會的傳統裡面.
我們特別強調.
「信徒皆為祭司」.
意思是我們每一位弟兄姊妹.
都能夠可以直接在神的面前.
能夠在神的話語裡面一起去尋求神的心意.
同時之間我們也都相信教會群體.
在神的面前共一起去分辨的重要性.
所以我們教會在治理上面.
會養大會就是最高決策的群體.

$^{681}$因為我們相信神會透過整個群體的禱告.
以及在討論的裡面去帶領我們教會.
這個提醒.
神的帶領不再是個人的事情.
而是群體一起去尋求的過程.
當一個決定越來越重要的時候.
我們就越需要帶到當中.
一起去禱告 一起去分辨.
在彼此的提醒和印證裡面.
能夠看得更加清楚神的帶領.
最後我覺得也是最重要的.
就是我們是否有一個願意馴服神的心意呢.
神的旨意即是什麼.
其實即是神的主權.
不是一個意見 而是命令.
任何面對神命令的人.
都只有一個合適的回應.
就是馴服.
馴服不是因為神的旨意.
可以帶來一些利益.
或者是可以帶來給你個人的享受.
一個更加舒服的生活.
我們不要用一個這樣的動機.
馴服只有一個原因.
就是因為這個是神的旨意.
一個真正屬靈的分辨.
最後一定會來到馴服這兩個字.
我們很多時候不自覺地.
為神的旨意加了一層糖衣.
好像剛才所說.
好像神的帶領一定會帶來更多的利益.
更好的條件 更少的困難.
當事情變得困難的時候.
我們就開始懷疑.
神風怒.
這個不是神的心意.
會不會.
不過其實神的旨意.
剛才一直都說.
祂更加關心的不是我們舒不舒服.

$^{721}$而是我們的生命.
有沒有被塑造呢.
好像上兩個星期.
吳牧師跟我們分享羅馬書八章.
萬事都互相效力.
要愛神的人能得益處.
在我們此刻所面對的不同的事.
人.
的裡面.
其實上帝要給我們甚麼功課.
當我們看回耶穌基督.
祂馴服天父的心意.
最後走向十字架的路.
保羅祂馴服神的呼召.
經歷的是逼迫.
被監禁.
患難.
當我們看歷世歷代.
每一個跟隨耶穌的人.
他們的道路其實都是一條不容易的路.
神正使用我們人生不同的經歷.
磨練我們的生命.
使我們能夠更加堅壯.
所以.
真正的馴服是.
主啊.
縱使這一條是不容易走的路.
縱使這一刻我未能夠完全明白.
只要是出於你.
我們願不願意.
緊緊地去跟從.
所以當我們去看.
第三個G.
勤成長的靈命.
去到今日.
當我們說我們要去辨識.
和遵從神的引導.
即是可能我們都會想.
勤成長的靈命.
當我們一路這樣聽到.

$^{761}$繼續吸收很多聖經的知識.
繼續參加不同的聚會.
這些都重要.
不過真正的屬靈成長.
不是去尋找秘訣.
而是說我們如何能夠與神同行.
聖經在舊約裡面說.
以樂與神同行.
不過就沒有再說多一點.
究竟如何與神同行.
是否因為他每天得到神特別的指示.
來與神同行呢.
聖經沒有再說.
不過我們相信以樂.
他每日的日常生活裡面.
他持續地去親近神.
聽神的話.
願意去馴服神的時候.
當我們選擇這樣的生活的時候.
就是在每日的經歷裡面.
慢慢體會到神的帶領.
所以結束的時候.
讓我們都好像私人大衛一樣.
第四和第五節.
這個都成為你和我的討聞.
我們一起讀出來好嗎.
預備 一二三.
耶和華.
求你將你的道指示我.
將你的路指教我.
求你指教我.
引導我進入你的真理.
因為你是救我的神.
讓我們一起低頭禱告.
在這一刻.
可能你正在人生的十字路口裡面.
可能你正面對一個很重要的決定.
心裡面很想去知道神的帶領.
今天神跟我們說.
我們不是要先去求答我.

$^{801}$而是我們要先來到神的面前.
又說神的心意已經向我們顯明.
但可能我們仍然遲遲都未馴服.
可能你在裡面已經知道.
神正在某方面提醒你.
光照你 呼召你的時候.
今天神跟你說.
不要再拖延.
憑著信心踏出那一步.
可能這一刻.
你覺得自己也很混亂.
很多的不確定.
甚至乎你感覺不到神的同在.
今天這一首詩篇再一次提醒我們.
在混亂的裡面.
我們仍然可以仰望神.
在困難的裡面.
我們仍然選擇等候.
在不明白的裡面.
我們仍然可以去投靠神.
求聖靈向我們說話.
我們一起禱告.
主啊 我們承認很多時候.
我們很想知道你的帶領.
但卻未必真正願意順服你的帶領.
求你赦免我們裡面的驕傲.
自我.
那種的焦急.
甚至乎對你的不順.
求你再次給我們一顆謙卑.
受教 敬畏你的心.
主啊 求你把你的道指示給我們.
把你的路教訓我們.
求你的真理繼續成為一道光.
引導我們走在你的心意道路當中.
主啊 讓我們不是單單去尋找一些.
我們可能找了很久的答案.
願意我們去學習.
成為一個跟你很親密.
很敬畏你.

$^{841}$與你一起去同行.
聽你的聲音.
順服你引導的疑慮.
讓我們的生命能夠在你裡面得著根深.
得著那份的穩妥.
甚至乎你要去帶領我們往更高處行.
願你的旨意成就.
聽我們的祈禱.
奉耶穌基督得勝的名求.
阿們.
\newpage



\section{啟示錄 2:1-7}
\label{sec:031956}
\textbf{以弗所:信仰純正卻愛心冷淡的教會}~~~~林穎豪傳道
\newline
\newline
連結: \href{http://www.church.com.hk/acms/content.asp?site=cdc\&op=show\&type=product\&code=031956\&layout=sermon}{\texttt{031956}} ~~~~ 語音日期: 2026-04-12
\newline
\newline
\hyperref[sec:031955]{\small{< < < PREV SERMON < < <}}
~
\hyperref[sec:index_chronic]{\small{[返順時目]}}
~
\hyperref[sec:031957]{\small{> > > NEXT SERMON > > >}}
\newline
\newline
啟示錄 2:1-7
\newline
\begin{longtable}{cl}
\hline
\hline
章節 & 經文 (和合本修訂版)\\
\hline
2:1 & \begin{tabularx}{0.7\textwidth}{X} 「你要寫信給以弗所教會的使者，說：『那右手拿著七顆星，在七個金燈臺中間行走的這樣說： \end{tabularx} \\ \\ \relax
2:2 & \begin{tabularx}{0.7\textwidth}{X} 我知道你的行為、勞碌、忍耐，也知道你不容忍惡人。你也曾察驗那自稱為使徒卻不是使徒的，看出他們是假的。 \end{tabularx} \\ \\ \relax
2:3 & \begin{tabularx}{0.7\textwidth}{X} 你能忍耐，曾為我的名勞苦而不困倦。 \end{tabularx} \\ \\ \relax
2:4 & \begin{tabularx}{0.7\textwidth}{X} 然而，有一件事我要責備你，就是你把起初的愛心拋棄了。 \end{tabularx} \\ \\ \relax
2:5 & \begin{tabularx}{0.7\textwidth}{X} 所以你要回想你是從哪裡墜落的，並且要悔改，做起初所做的工作。你若不悔改，我要到你那裡去，把你的燈臺從原處挪去。 \end{tabularx} \\ \\ \relax
2:6 & \begin{tabularx}{0.7\textwidth}{X} 然而你還有一件可取的事，就是你恨惡尼哥拉派的行為，這種行為也是我所恨惡的。 \end{tabularx} \\ \\ \relax
2:7 & \begin{tabularx}{0.7\textwidth}{X} 凡有耳朵的都應當聽聖靈向眾教會所說的話。得勝的，我必將神樂園中生命樹的果子賜給他吃。』」 \end{tabularx} \\ \\
[1ex]
\hline
\hline
\end{longtable}
$^{1}$早陣子我太太Karen看了一套韓劇.
大家… 報劇很正常.
不過很快看著看著.
她就沒有看了.
我心裡面很疑惑.
為什麼很疑惑呢?.
因為不是很正常.
因為她平時看什麼.
看什麼劇都會一直繼續看下去.
又或者不要看什麼劇.
我很喜歡看動漫我這個人.
我開什麼動漫跟她一起看.
她都會陪我看完它.
但是為什麼這次韓劇.
她竟然不看完它呢?.
我疑惑之際.
她就很簡單一句回覆我.
男女主角沒有火花.
可能我們都看過一些電視劇或者電影.
當男女主角之間沒有火花.
或者沒有化學作用.
沒有那個化學.
縱使那個電視劇的劇本寫得再好.
電影的畫面再亮麗.
我們總是覺得.
就差一點點.
慢慢我們就對她失去了興趣.
今天《啟示六二章》一至七節.
同樣給我們看到一間.
在人的眼裡面近乎是完美.
不過差一點點的教會.
這間教會有著一個很.
即是台面上很亮麗.
很令人讚嘆的一個CV是一個履歷表.
這間教會是位於當時文化經濟中心的地方.
是伊弗所的地方.
這間教會的歷史也很悠久.
也有很多名目曾經在那裡目過會.
大家知不知道.
有誰在伊弗所目過會?.

$^{41}$有保羅曾經在那裡目過三年的會.
有提摩泰.
有阿居拉 有伯基拉 有阿波羅.
這些我們也在聖經裡面聽過他們的名字.
甚至今天我們讀的《啟示六》.
寫的那個使徒約翰晚年的時候.
也是在那裡服侍.
這麼多個出名的牧者在那裡服侍.
是一間像面這麼亮麗的教會.
但是就是差一點點.
所以當那位右手拿著七星.
在七個金燈台中間行走的耶穌基督.
寫信給這間教會的時候.
祂的評價就像一盆凍水一樣.
提醒了這間教會裡面.
所有自以為很敬虔的人.
今天就讓我們透過這封.
耶穌寫給伊弗所教會的信.
同樣我們也誠實地反思一下.
我們的心會不會也好像.
伊弗所教會的人一樣.
雖然我們的外殼裡面.
堅守著純正的信仰.
但是內裡卻是一早已經失去了起初的愛心.
失去了教會應有的味道.
失去了信徒應有的溫度.
首先主耶穌在信的開頭.
肯定了伊弗所教會三個的優點.
第一他們有什麼做得好呢.
他們是持久奴婦的服侍.
主耶穌說我知道你的行為.
你的勞碌 你的忍耐.
原來伊弗所教會的信徒.
不是一群只是坐在教會口裡面說.
這裡真好的信徒.
他們是積極服侍.
為著福音為著教會的事工.
他們擺上時間.
努力的一群實幹家.
他們盡心竭力的服侍.

$^{81}$令到耶穌要讚他們讚到寫在聖經上.
第二 耶穌看到他們.
肯定他們做得好的.
就是他們是一群堅守真理.
並且不容隱惡人和異端的教會.
耶穌說我知道你們不容隱惡人.
你們也曾經試驗自稱為使徒.
卻不是使徒.
你看得出他們是假的.
原來伊弗所教會在神學上.
是極度嚴謹的一間教會.
在真理上他們完全不妥協.
他們不但不容忍.
教會裡面出現任何的惡行.
如果有講員來到他們中間.
傳講一些似是而非的道理.
他們也會立刻拿出聖經去檢驗.
耶穌更加在第六節說到.
他們恨護尼哥拉一黨人的行為.
尼哥拉一黨就是指當時伊弗所教會.
他們面對著一群很有影響力的異端.
不過 伊弗所教會的信徒.
有著極高的屬靈的分辨力.
得憑力可以將這些假教師.
趕出他們的教會.
第三個優點.
耶穌說他們是一群老虎.
有忍耐的見證.
耶穌說 你也能忍耐.
曾為我的名老虎並不拔券.
伊弗所教會的信徒.
甚至為著福音的緣故.
付上了極大的代價.
即使面對當時羅馬帝國的逼迫.
當時是很多偶像崇拜的.
一個龐大壓力之下.
但是他們沒有退縮.
他們咬緊牙關.
他們不怕吃苦.
為主的命忍耐到底.

$^{121}$聖經還說 並不拔券.
這個可能也是今天.
我們很多信徒忠心服侍的縮刑.
我們每個禮拜就在教會裡面.
很多事情忙很多事情做.
我們忙著大敬拜 練歌.
教主也學 參加茶經.
擔任小組的組長.
如果耶穌今天來到這裡.
可能祂也會點點頭跟我們說.
嗯 真的見到你的老虎.
我知道你很努力.
去支持著這間教會.
然而 信仰的危機.
往往不是發生在.
我們什麼都不做的時候.
而是發生在.
我們做盡了一切正確的事.
卻是忘記了我們為了什麼而做.
可能我們會想.
一群這麼盡忠的教會.
一群這麼努力的信徒.
這樣也有信仰危機.
回到一開始的第一句.
兩個主角之間沒有火花.
第四節.
耶穌的話鋒一轉.
心裡面原本盛載著對教會的期待.
但是此刻剩下的是失望.
耶穌說什麼?.
耶穌對教會說.
你變了.
你以前不是這樣的.
你以前很愛我的.
耶穌忽然之間.
從那個天上尊貴君王的樣式.
他變成了戀愛片裡面的情侶一樣.
就好像申命記四章二十四節說到.
因為耶穌說你的神乃是烈火.
是忌邪的神.

$^{161}$忌邪.
又或者更加貼切的翻譯.
是嫉妒的神.
他是一個 Jealous 的神.
神對你的愛劇烈如火.
神亦都嫉妒妒忌任何阻隔.
在你和神之間所有的人事物.
為什麼你要和祂一起?.
為什麼你花這麼多時間在祂身上.
而不是用那些時間來陪我?.
你知不知道我多愛你?.
我創造你.
我供應你生活一切所需要.
我甚至為你承受人世間肉體的自痛.
心靈的自負.
為你死在十字架上.
你還牽著祂?.
所以啟示錄這裡所說的愛心所說的.
不是教會裡面信徒之間那種.
關懷互相相愛的愛心.
而是更加純粹作為一個.
相信一個跟隨主的人.
對主耶穌基督應有的愛.
應有的渴慕.
耶穌說.
你變了.
你以前不是這樣的.
你以前很愛我的.
耶穌在這裡用的字是.
起初的愛心.
就是當我們剛剛認識主的時候.
那種單純.
那種熱烈.
那種毫無保留的愛.
親愛的弟兄姊妹.
你當初信主的時候.
是怎樣的你記不記得?.
那時候我們為了自己的罪得赦免.
而痛哭淚流.
我們渴望讀丁.

$^{201}$不是因為我們要預備查經.
預備領思.
而是我們單純.
我們想聽到神的聲音.
我們想認識主更加多.
我們參與事奉.
不是因為那個鄰庚表輪到我們當值.
更加不是因為我們背著什麼身份.
所以我們一定要服侍.
而是因為我們愛主.
所以我們很想為主做些什麼.
昔日我們受洗的時候.
我們在台上興奮地分享.
我們怎樣經歷神.
怎樣感受到神的愛.
是因為我們心裡面很開心.
很雀躍.
很想分享我們怎樣經歷神.
很想讚美神.
很想多謝神.
那時候我們很激動.
捉住身邊的朋友跟他們說.
很想跟他們介紹這位主.
很想他們快點信主.
因為我們心底裡真的相信.
信主是生命裡面最好最好的事情.
昔日我們如此如此.
但是今天.
今天我們是不是真的仍然.
從心底裡覺得.
信主是最好的事.
我們是不是從心底裡覺得.
耶穌是真的最愛我們的主.
今天我們發生了什麼事.
耶穌說耶穌所教會.
離棄了這份敬愛.
離棄.
不是不小心跌了.
而是在忙碌裡面.
任由這份敬愛去到枯萎.

$^{241}$任由這一點的火花去到熄滅.
當他們的手.
還是為著主.
做著各樣的事工.
腦海裡面裝著很多很正統的神託.
但是他們的心麻木了.
弟兄姊妹.
以忽所教會的光景.
是不是都反映著我們生命的景況.
我們每個主日回來教會崇拜.
忙著煉歌.
忙著煮飯.
忙著洗碗.
忙著預備茶經.
忙著開會.
忙著寫會議紀錄.
但是當耶穌.
當主耶穌問我們.
你愛我媽的時候.
我們好像.
不是很敢正面回答主耶穌.
我們只會說.
主啊 你知道的了.
我為你做這麼多事情.
難道還不夠嗎.
特別我們華人.
我們經常都以為.
為神做牛做馬.
就代表我們很愛主.
我們把責任當成了愛.
我們以事工.
取代了關係.
但是耶穌這一刻.
又變回那個.
戀愛當中.
戀愛片中的情侶.
耶穌說.
耶穌說我不是要這些東西.
我不是要你買多少東西給我.
我不是要你.

$^{281}$做完全屋的家務.
我不是要你做這些事.
我要的是你的生物.
信仰不是服侍.
其實信仰不是發展事工.
信仰不是死做爛做.
信仰甚至都不是.
維持教會體制的運作.
真正的信仰.
是你和主耶穌之間的關係.
是你和主耶穌之間的火花.
耶穌接著給了一個很嚴厲的警告.
你若不悔改.
我就臨到你那裡.
將你的燈台從遠處下去.
燈台代表的是教會發光的見證.
如果一間教會失去了對神的愛.
即使祂的神學再純正.
聚會人數再多禮堂再莊嚴再亮麗.
在神的眼中.
祂都失去了為主發光的資格.
神寧可收回這個見證.
都不怨街外的人世人.
見到一個只有冷酷真理.
卻是毫無愛主的心的假象.
而歷史也都很確切跟我們說.
伊夫索這座的城市最後真是沒落.
伊夫索的教會的燈台最終都被拿去.
伊夫索教會剩下的.
只是一些斷壁.
一些的石柱.
一些的殘門.
今天伊夫索教會已經沒落了.
啟示六二章.
這封寫給伊夫索教會的信.
對於今天的普世教會.
或者對於我們來說.
還有什麼意思?.
是不是就只是對教會一個強烈譴責.
嚴厲警告呢?.

$^{321}$有時候我們讀聖經.
我們會很著重於經文的用字或者內容.
但是我們會容易忽略了.
原來聖經.
經文裡面有很多的編排.
有很多不同的結構.
就好像我們人與人之間的溝通.
說話的內容當然很重要.
但是語氣和態度.
可能反映了一些超越文字超越內容的意思.
啟示六二章這封信有什麼編排呢?.
我們看到耶穌對教會的說話是怎樣的.
耶穌對教會的說話.
祂不是一口砂糖.
一口…….
耶穌的說話.
不是先說了兩句好聽的說話.
然後再一味地訓示教會.
耶穌的說話是怎樣的?.
耶穌的說話是一口砂糖.
一口…….
最後再多加一口砂糖.
其實你也有做得好的地方.
你也有做得可取之處.
這個編排讓我們看到什麼呢?.
讓我們看到.
原來耶穌不是想打沉這間教會.
耶穌也不是放棄這間教會.
放棄這段感情.
祂不是想教會灰心喪志.
反而是希望教會.
未來在往後的日子.
能夠做得更加的好.
雖然我們離棄了對耶穌那份的愛.
失卻了和耶穌之間的火花.
但是耶穌仍然很想和你重新開始.
祂接下來指明了三個.
帶著行動力的明確方向.
第一個 耶穌說要回想.
所以應當回想你是從哪裡墜落的.

$^{361}$回想不是進入對過去的懷緬,懷舊.
而是找回我們軟弱跌倒的坐標.
我們要停下來.
停下來我們忙碌的腳步.
安靜下來問自己.
其實我是什麼時候?.
我是什麼時候開始.
讀經讀到好像交差一樣?.
我是什麼時候開始.
不想再認識這個神更加多?.
我是什麼時候開始.
不再相信禱告是帶著權柄的?.
然而 耶穌要我們回想.
不是陷在那個罪疚之中.
耶穌接著要我們悔改.
悔改不是輕描淡寫.
說一句 對不起.
我最近忙.
悔改不只是心思意念的改變.
是因為愛不是一種虛無飄渺的感覺.
而是一種行動.
所以要有實質的行動.
才算是悔改.
要有什麼行動呢?.
耶穌繼續說.
行起初所行的事.
行起初所行的事意味著.
重新建立與主耶穌親密的關係.
主耶穌很想與你重新開始.
所以我們重新建立不斷的關係.
就必須要花時間.
花時間在信仰的基本功裡面.
我們去禱告 讀經 靈修.
經歷神的同在.
是我們放下過去.
我們出於習慣.
出於責任的行為模式.
我們回到我們的內心.
單單與神去親近.
特別是我們靈修 親近神.

$^{401}$不一定要只在房裡.
當我們走在街上.
遇到不同的人 不同的事.
我們都可以繼續祈禱.
特別是我們身處澳洲.
這麼美麗的地方.
昨天我們一群年輕的弟兄姊妹.
出去了行山.
是不是行山 還是行海邊.
看到一個世界很美.
當我們看到這個世界這麼美的時候.
我們豈不是可以隨時隨地.
為著上帝的創造去讚美祂.
早幾個星期.
我認識了另一間教會的一個姿態.
他說他親近神的時候是什麼時候.
當他駕車 上班 下班.
那四十分鐘 那一個小時.
就是他禱告親近神的時間.
有時候.
當我們信主的年日久了.
我們開始覺得自己是頒典的信徒.
我們覺得自己開始上岸了.
我們不用再去那些賢經班.
那些訓練班.
叫那些初信的年輕人去就對了.
讓他們認識神就可以了.
但我 我回教會很久了.
但是 綺弗所教會.
正正是一個很好的例子.
告訴我們.
人對神的愛.
和人回教會的年日.
不是一個正比的關係.
不過 要做些什麼.
才可以重建和主耶穌基督的關係.
和耶穌重新立國.
其實說到底.
不是我可以給你一個很確切的清單.
不是我可以給你一個清單.

$^{441}$你做了這些事情就可以了.
隔天你就跟他說我愛你.
定時定後跟他約會 出去吃飯.
大時大節送花.
不時預備一些小禮物.
我講不了給你聽.
因為這個是你和神的關係.
唯有你才知道.
起初你是怎樣經歷神.
所以經文說.
耶穌說 你要回想.
是你要回想第一次的時候.
你遇見主的時候.
那種恩典的時候.
那種被傳言接納的感動.
是你要回想第一次.
你親眼見證神大能的那份震撼.
是你要回想第一次.
聽到神親自跟你說話的那份安慰.
那份提醒對你心靈的行動.
然後 讓這一份的感動.
這一份的感受.
這一份的愛.
這一份的渴慕.
再次帶回來進入同主耶穌基督的關係.
再次帶回這個愛.
這個感動去到復仕.
再一次渴望為主拜上.
耶穌基督深愛我們.
祂也渴望我們同樣的愛祂.
祂最後在第七節.
用一個盼望作為一封信的結尾.
祂說得勝的.
我就將神樂園裡面生命樹的果子賜給祂吃.
當我們願意悔改.
當我們願意找回起初的愛.
找回跟神之間的火花.
祂就應許賜給我們永遠豐盛的生命.
並且享受跟祂最深切的連結和團契.
親愛的弟兄姊妹.

$^{481}$耶穌今天不是用一個很嚴苛老闆的身份.
去檢視我們教會的業績.
反而耶穌是用一個.
深愛我們救主的身份.
去尋找我們生命的起初的火花.
就好像戀愛中戀愛片的情侶一樣.
期待著對方毫無保留的愛慕之情.
不過.
不過好像用情侶去類比.
我們跟耶穌之間的關係好像不太合適.
因為在舊約聖經裡面的二塞亞書.
在新約聖經裡面的啟示錄.
往後我們再會看到的篇章講到.
所引述的人和神的關係是什麼.
是比情侶更加情深.
更加火熱.
更加激進的夫妻關係.
甚至是新婚夫妻的關係.
新婚夫妻的關係是什麼.
是彼此眼中已是得對方一個.
尊一而終.
並且是以盟誓結連的關係.
親愛的弟兄姊妹.
主耶穌基督跟我們的關係.
不是單單好像情侶之間.
很盼望跟祂朝夕相對.
經常跟祂黏在一起.
經常看到祂可以分享彼此的情感.
思想那種親近.
不僅如此.
我們跟主耶穌之間.
更加應該是夫妻之間那種.
從今日去到死.
我的眼中只有你一個.
永遠都至得對方.
患難與共.
共度悲喜.
歷盡所有.
這一種夫婦之間關係那種的偽身.
盼望啟示律而將這一封.

$^{521}$寫給爾弗所教會的信提醒我們.
不要讓日積月累的習慣.
以及從身份而來的責任.
掩蓋了我們對神那份.
單純的愛慕和偽身.
即使可能我們已經丟失了.
或者離棄了.
我們對神那個單純的愛.
那點的火花.
可能我們跌了很久了.
但是從今天起.
從現在起.
要我們再一次回想.
找回我們與神之間.
那個起初的愛心.
與主耶穌基督重新來過.
讓我們再一次.
深深愛上這位創造我們.
為我們寫記的主耶穌基督.
我們一起同心吐槽.
親愛的主耶穌基督.
我們多謝你.
你是那個愛我們.
為我們寫在十字架的主.
主啊求你赦免我們.
很多時候我們在教會裡面.
忙著做各樣的事工.
我們以為這樣就是代表我們愛你.
但是原來我們不知不覺間.
用責任取代了關係.
原來我們心裡面.
在忙碌的時候.
心裡面就慢慢失去了.
起初對你那種的熱誠.
那種的火花.
那種的熱切.
求聖靈今天光照我們.
幫助我們回想.
初信的時候.
那份單純的感動.

$^{561}$那份的渴望.
主耶穌我們都願意悔改.
求你幫助我們.
放下出於習慣的宗教行為.
卻是重新建立和你親密的關係.
讓我們每一天的生活.
每一個的服侍.
都是因為被你的愛激勵而發出的.
我們獻上我們的禱告.
立志應聚.
奉主耶穌的聖名祈求.
\newpage



\section{出埃及記 13:8}
\label{sec:031957}
\textbf{聖經人聖經事:要告訴你的兒子}~~~~黃克勤牧師博士
\newline
\newline
連結: \href{http://www.church.com.hk/acms/content.asp?site=cdc\&op=show\&type=product\&code=031957\&layout=sermon}{\texttt{031957}} ~~~~ 語音日期: 2026-04-08
\newline
\newline
\hyperref[sec:031956]{\small{< < < PREV SERMON < < <}}
~
\hyperref[sec:index_chronic]{\small{[返順時目]}}
~
\hyperref[sec:031958]{\small{> > > NEXT SERMON > > >}}
\newline
\newline
出埃及記 13:8
\newline
\begin{longtable}{cl}
\hline
\hline
章節 & 經文 (和合本修訂版)\\
\hline
13:8 & \begin{tabularx}{0.7\textwidth}{X} 當那日，你要告訴你的兒子說：『這樣做是因為耶和華在我出埃及的時候為我所做的事。』 \end{tabularx} \\ \\
[1ex]
\hline
\hline
\end{longtable}
$^{1}$各位弟兄姊妹,我們一起讀聖經.
今次的主題是.
「要告訴你的兒子,告訴你的兒子」.
生命說書網絡「聖經人,聖經事」.
幫我們一起去看聖經.
去學習一下,認識一下.
上帝是怎樣的.
今次的主題「要告訴你的兒子」.
這個就是上帝說的.
要叫以色列人告訴你的兒子.
繼續看儒字的故事.
然後他們就離開了埃及出來.
出入埃及的第十三章.
摩西對百姓說.
「你們要記念從埃及圍爐之家.
出來的這一天」.
記住這一天.
摩西對百姓說要記住這一天.
其實摩西當時是代表著上主跟我們說話的.
主耶和華跟我們說.
「我要記念你們從埃及圍爐之家出來這一天」.
那麼這一天發生了甚麼事呢.
是耶和華用大能的手.
將你們從這地方領出來的.
這一天記住這一天.
「記念這一天」.
這一天是為甚麼要記念呢.
是因為耶和華用大能的手.
將你們從這地方領出來.
接著這一句很有趣的.
好像有關又好像沒有關.
「有教的餅都不可吃」.
這一天耶和華帶他們離開埃及.
就說「有教的餅都不可吃」.
其實這個有一個背景的.
就是他們當天將麵粉用水拌勻.
可能有些人加了一些塑膠.
即是發粉.
有些人加了一些沒有加的.
然後就放在這裡.

$^{41}$大家都知道了.
做過麵包做過餅的都知道.
那些塑膠要好幾個小時才能夠將麵團發起.
然後才去烤餅.
未有機會他們就生的.
未烤的也有.
未落教的也有.
於是就一起帶出來離開.
放在缸裡面拿著走了.
因此他們出來之後餓了.
第二天早上就吃「無教餅」了.
而「無教餅」他們平時都有吃的.
是因為他們在余節的晚餐.
是指定吃「無教餅」的.
因此「無教餅」是象徵著「沒有罪」.
那個沒有殘疾的羊羔.
是象徵我們的救贖主耶穌是一個沒有罪的主.
而吃「無教餅」和苦菜.
是象徵著我們這個人要潔淨.
以及在人生裡面經歷了很多苦難.
亦都可以說我們的主是無罪.
再次說是無罪又經歷了很多苦楚.
這個「烤羊」是表示經歷過苦楚.
比較是象徵十價痛苦多一些.
「無教餅」是象徵著無罪.
苦菜象徵著一個罪惡困苦中的人生.
「無教餅」是要紀念逾越節.
帶他們出來這一天.
特別是拿回「無教餅」.
有教的都不可以吃.
即是說要吃「無教」了.
再下去「啊不悅」.
「艱的謝日是你們出來的日子」.
記住這個月份是叫做一月 正月.
是「啊不悅謝日是你們出來的日子」.
又是說這個「日子」.
大家看到「日子」紀念這一天.
這一天是神用達林的手帶你出來.
有教的都不可以吃.
用吃「無教餅」來紀念神帶他們出去埃及.

$^{81}$另外要紀念的就是這一天是什麼時候.
什麼時候是「啊不悅」.
不是上次說離散月嗎 是正月.
本來離散月不是正月.
不過由雨節開始.
神就說這個月份是正月.
開始歷法計過了.
又是「啊不悅」.
「啊不悅」是「啊不」.
「啊不」意思是初熟大墨水.
即是說大墨長起.
然後很多水紙這樣向下.
未熟的是剛剛出.
譬如我們看到有蘋果,橙在樹上.
但是它未熟.
有芒果,青色未熟.
未熟的時候大家已經看到整個果.
看到大墨水了.
是初熟還是青的.
這個「啊不」就是初熟大墨水的意思.
即是他們要來慶祝初熟節的時候.
感恩 已經有墨水見到.
不過未熟.
即是說再過多四,五十天就熟了.
那時候就變成黃了.
那時候就是五巡節.
或者叫七七節的時候.
這個月份他們除了叫離散月.
又叫做初熟大墨水月.
即是叫「啊不」的意思.
是另一個名字而已.
正名是叫做離散月.
「啊不月」是紀念農田的需要.
農作物的時日.
就是這個大墨水初熟的時候.
你們就初一級了.
將來耶和華領你入迦南.
黑人,亞摩利人,希美人,耶布斯人之地.
就是神向你們祖宗起誓.
應許給你那塊流瀨與物之地.

$^{121}$你們進去都還要在這個月份.
即是離散月或者叫做「啊不月」.
來搜借禮.
借禮是什麼?.
就是照那晚吃魚子羔羊那樣.
要做這個魚子的禮.
記得 特別要記得.
吃魚子的禮除了羊.
最特色的地方就是無餃餅.
不可以看到有餃的.
你會知道當時的迦南地.
就是有迦南人.
有黑人,亞摩利人,希美人,耶布斯人等等.
他們住在迦南地.
迦南這個地方是一個低地.
意思是低的平地.
流瀨與物之地就是神帶你們去那個地方.
我將它神稱為「流瀨與物之地」.
就這樣看「流瀨與物之地」.
通常我們會認為是很多牛奶和很多蜜糖.
其實不是的.
流瀨是象徵著畜牧業比較適合.
特別是羊多的.
而蜜是指蜜棗.
以前都講過了.
蜜棗Date.
就不是指蜜糖.
就是有很多特別的食物.
很有營養的.
這個Date這個蜜棗.
以及畜牧.
那麼你就要在進去那裡.
這塊地迦南地.
你都還要守這個月節.
你們要吃無教餅七日.
到第七日要向耶和華守節.
這七日之久要吃無教餅.
在你四景之內不可見有教的餅.
也不可見發酵的物.
這裡應該讀「發酵」.

$^{161}$發酵的物.
這個第一要大家留意.
來到這裡才講的.
其實就是逾越節.
逾越節就是一個晚上來的.
吃羊肉.
然後就吃無教餅.
吃苦菜.
然後就有些血在門尾門腔上.
然後就走了.
一聽到一聲令下.
大家就收拾東西走了.
走出來第二天.
就繼續開始吃無教餅了.
出來之後.
摩西就跟他們講這番話了.
記住.
大家去到應許之地.
都是要守這個節日的.
都在阿筆月的時候守.
阿筆月就是離散月.
阿筆就是農業的月份.
離散月就是正式數曆法.
在這裡.
出來之後.
在逾越節之後.
接著那七天.
要吃無教餅七天.
所以.
在猶太人來講.
以色列來講.
逾越節那一晚一過.
就是接著守七天的.
叫做無教節的了.
那七天無教節.
就吃無教餅.
以及所有會發酵的東西.
都丟掉了.
那今天的猶太人都是這樣守的.
那七天.

$^{201}$就會將所有廚房裡面.
麵粉,發粉全部丟掉.
麵包,餅乾全部不要.
因為早點吃光它.
要是完全無教的.
所以他們每一年.
都是將他的房子,廚房清理一次.
餐廳停業七天.
酒店廚房停業七天.
就吃無教餅.
這樣的話.
預備先弄好.
那七天就清潔了.
現在有些酒店餐廳大的.
他要找一位.
那些叫拉比尼.
就是給他檢查.
穿好手套這樣檢查.
然後寫一張證明給他檢查.
他是無教了.
他才重新開業.
將證書要擺出來的.
那猶太人很守規矩.
那個才去吃的.
他就死守著這個無教節.
就七天.
那如日一圓就是無教節了.
他就叫做除教節了.
這樣吧.
好.
除教節的來歷就是這裡的.
無教節或者除教節.
好了 第八節.
當那日你要告訴你的兒子.
謝氏因爺說.
在我出埃及的時候.
為我所行的事.
在這個吃無教餅.
在除教節的七天.
就要跟子女們說.

$^{241}$以色列人逾越節.
在埃及圍爐逾越節摩西.
神透過摩西.
帶他們出來的故事.
在那裡就要說.
第二點.
你要告訴你的兒子.
這個就是今天的說話.
今天的主題.
你要告訴你的兒子.
告訴他什麼呢.
是耶和華在我們出埃及的時候.
為我們所行的事.
整個逾越節是怎麼過的.
又進入到今天.
無教了.
就是這樣.
好了 記得我要告訴你.
而且要時時將那些經文.
或者一些裝飾的東西.
放在你的手上做記號.
或者你額上做紀念.
使耶和華的律法尚在你口中.
因為耶和華曾經用大量的手.
將你從埃及領出來.
為什麼這個七七.
不是七一節是月節.
這個無教節.
或者叫除教節.
那七日要做什麼呢.
跟子女不停讀聖經.
講述出埃及 逾越節.
以及除教節的故事.
要講述.
而且要將一些經文.
小小一條的放進盒子裡.
皮的盒子.
綁在額頭上.
今天他們都會看到的.
有些就繞一條帶.

$^{281}$綁在手上.
在盒子裡放一些經文.
用來唸聖經.
唸什麼經文呢.
唸剛才讀的那些.
神叫他們怎麼離開埃及.
怎麼讀聖經.
有時候我們去以色列.
有時候看到一些小朋友.
很有趣的.
就一條很長的黑色帶.
布皮帶.
上面有個盒子.
裡面放一些經文.
然後在街上請人.
會不會想唸聖經.
好啊.
他就幫他繞著繞著.
盒子裡有些經文.
繞著繞著.
那個人就在唸.
唸盒子裡的經文.
唸完就取回給他.
這樣就給錢他.
原來那個小朋友.
是可以走出來叫人唸經.
不是唸經.
唸那些經文而賺錢的.
哇 多好.
那些人就很好的.
又會坐在那裡唸.
唸那些經文.
然後就取回給他.
就是戴在頭上.
唸那些經文.
然後給他.
我覺得他們以色列人.
有些傳統習慣挺好的.
很喜歡唸聖經.
因為神是透過摩西教.

$^{321}$你要記得將那些有關的.
在手上做記號.
在額上做紀念.
律法常在你口中.
記得耶和華曾用大能的手.
將你從埃及地領出來.
所以你每年要按著日期.
來守這條例.
變成條例.
一定要守的.
將來耶和華照祂.
向你和你祖宗所起的誓.
是將你領到迦南人之地.
將這個地賜給你.
那個時候第三.
你就要做多一件事.
除了繼續守彌散月.
或者亞不月這個日子.
守逾越節和隨交節.
講返這段故事.
戴在頭上做經文紀念之外.
還有一件事告訴子女聽.
講故事.
另外到那時候第三件事要做的.
要將一切投生的.
並生出中投生的歸給耶和華.
去到迦南地英去之地.
留禮與物之地的時候.
不單止要繼續守彌散月.
逾越節,隨交節.
整齊經文.
時時講給子女聽.
還要做多一件事.
去到那裡才做.
到那時候你要做什麼呢?.
要將一切投生的歸給耶和華.
連生出都要.
什麼叫做投生的歸給耶和華呢?.
就是所有長子都說.
是屬於耶和華的.

$^{361}$歸耶和華為姓.
就是神啊 我的大兒子奉獻給你.
你帶領他的人生.
我這隻牛.
投生了一隻牛仔.
公的.
這隻公的小牛.
公的小羊.
屬於神你的.
你要我就會將牠獻給你.
什麼時候我就會將牠獻祭獻給你.
這一隻是投生的屬於耶和華.
人就不會獻了.
人是神啊 奉獻給你了.
屬於你了.
帶你去殿那裡.
去禱告告訴神.
這是我大兒子.
他獻給你了.
這樣的關係挺好的.
這樣就挺好的.
我們今天很多時候不是的.
以前特別多一些.
以前有一個牧師.
他說很不開心.
為什麼呢?.
有個婆婆.
那個婆婆都不是太婆婆了.
其實她的子女都大了.
都差不多了.
都十幾歲了 二十歲了.
以前的小朋友很多.
五個兒子 有些女兒.
牧師啊.
我五個兒子.
這一個是奉獻給神的.
為什麼奉獻給神呢?.
那個做醫生.
這個做律師.
這個教書.

$^{401}$這個做什麼.
這個沒什麼特別出色.
奉獻給神.
你帶他做傳道人.
不是吧?.
真的很不開心.
最沒什麼表現的那個.
什麼叫沒表現呢?.
是你認為他沒表現.
做醫生他有表現.
做律師他有表現.
做教書他有表現.
這個沒什麼特別.
讀書多或者讀書好.
沒什麼特別出眾的東西.
奉獻給主做傳道.
這樣就對了.
不是這樣的.
為什麼不將最好的.
奉獻給主呢?.
是要將他認為.
賺錢最少他的.
奉獻給主做傳道.
不是吧?.
為什麼會這樣呢?.
以前聖經裡面規定.
將一切投生的獻給耶和華.
大洗的那個 奉獻給耶和華.
這樣就對了.
適合神的心意了.
以前特別講的神祕就是這樣的.
好了.
不過多數以色列人.
是牛羊多的.
如果有驢 怎麼辦?.
投生的驢.
需不需要獻給神呢?.
不太能獻給.
為什麼呢?.
因為驢要用.

$^{441}$牛羊可以用來獻祭.
一歲的公牛.
一歲的公羊是投生的.
就會獻祭了.
驢又不能獻.
因為驢不潔淨.
不可以獻祭.
駱駝又怎能獻祭呢?.
不能獻.
兔子又不能獻.
用驢做例子.
如果是投生的驢.
照道理也是屬耶和華的.
將它獻給神.
又不能獻.
怎麼辦呢?.
你可以用羊羔代替.
譬如我有一隻驢.
投生的要獻給神.
就應該是做獻祭.
不過我又不獻.
怎麼辦呢?.
驢一來不能獻.
因為不潔淨的動物.
又越來越再誇有性.
你可以用一隻羊羔來代替它.
用羊羔來代替.
代贖它.
如果不代贖.
就要打折它的頸項.
什麼意思呢?.
如果你不肯用一隻羊羔來代贖它.
獻祭給神.
你就要打斷它的頸.
即是什麼意思呢?.
殺死它.
這麼固執.
其實意思是想特別表達.
投生的屬耶和華.
如果那樣東西不能獻做祭物.

$^{481}$你就要用羊羔代替.
如果不是那一隻獻給神的.
雖然不能獻.
但都要用它的生命來交給神.
而始終投生的都要屬出來.
如果是兒子.
就打斷了兒子的頸.
不會獻祭的.
人不能獻祭的.
於是你又用一隻羊羔來代替它.
獻給神.
第一次這兩個字是表示什麼呢?.
生出的投生.
如果是牛羊.
就可以獻祭給神.
屬於它的.
兒子都屬於神.
那怎麼辦呢?.
要麼就擺好似撒母耳那樣.
就在聖殿那裡生活.
那又不行.
又沒有位置又沒有祭司肯收.
那怎麼辦呢?.
那你就用一隻羊羔代替你的兒子.
將那隻羊羔獻給神.
所以投生的都要屬出來.
用一隻羊羔來屬它.
這是大家要留意的.
都要叫它這樣做的.
好了.
如果日後你的兒子問你.
什麼意思啊?.
為什麼要獻這一隻羊羔來代替我的?.
你就要告訴他.
因為你是大兒子.
大兒子的長子是屬於耶和華的.
但是耶和華不歡喜人.
是用人獻祭的.
所以我用一隻羊羔來代替你.
不過你要記得.

$^{521}$這隻羊羔代替你獻了給神.
其實你是屬於神的.
你要告訴他.
什麼意思啊?.
你就告訴他.
是紀念耶和華.
用大力的手.
將我們從埃及圍爐之家領出來.
因此.
這裡幾點要大家記得就是.
神想我們世代傳下去.
將神大能的作為告訴子女.
子女每一年都聽.
每一年都聽這個故事.
其實對它很有用的.
而且還有一些實質的做的行動來代替.
吃無教餅的.
就記起我原來是逾越節.
我就要更加聖潔無罪.
另外.
長子是屬於耶和華.
是要更加跟隨祂的.
其實這個很聰明.
上帝當然更聰明了.
就是要意識的人是這樣.
因為長子是獻了給神.
他自己又聽了這個故事.
他更加有意識.
將來他爸爸不在.
由他做一家之主.
他就會帶領著全家人繼續跟隨主.
然後他的長子.
又會帶領著全家人一起跟隨主.
所以這個是非常之重要的.
大家要記得這件事了.
要記得一代一代傳下去.
將這個神大能的手的故事.
講給你的子女聽.
講給別人聽.
我們一起祈禱.

$^{561}$求主幫助我們.
有很多時候我們很善忘.
那些恩典經歷了.
哇 神給我這麼大的恩典.
我一定要記得.
不過轉過頭來忘記了做.
去到有事情的時候.
哎呀 又慌忙.
又重新問過.
其實神大能的手從來沒有改變.
以色列人他們是有主的指示.
每一個世代.
每一個兒子問這件事.
是代表什麼的時間.
就要將故事跟他講.
每一年的餘節.
每一年的除教節.
就要將這些故事.
餘的故事重新再講.
是指明你大能的手.
帶領了他們出埃及.
幫助我們都記得.
你自己大能的手.
願你得到榮耀讚美.
時時傳揚你的救恩.
禱告奉耶穌名求.
阿姆.
\newpage



\section{撒母耳記上 1:1-11}
\label{sec:031958}
\textbf{一心•倚靠}~~~~李玥坤傳道
\newline
\newline
連結: \href{http://www.church.com.hk/acms/content.asp?site=cdc\&op=show\&type=product\&code=031958\&layout=sermon}{\texttt{031958}} ~~~~ 語音日期: 2026-04-12
\newline
\newline
\hyperref[sec:031957]{\small{< < < PREV SERMON < < <}}
~
\hyperref[sec:index_chronic]{\small{[返順時目]}}
~
\hyperref[sec:031959]{\small{> > > NEXT SERMON > > >}}
\newline
\newline
撒母耳記上 1:1-11
\newline
\begin{longtable}{cl}
\hline
\hline
章節 & 經文 (和合本修訂版)\\
\hline
1:1 & \begin{tabularx}{0.7\textwidth}{X} 以法蓮山區有一個拉瑪的瑣非人，名叫以利加拿，他是蘇弗的玄孫，託戶的曾孫，以利戶的孫子，耶羅罕的兒子，是以法蓮人。 \end{tabularx} \\ \\ \relax
1:2 & \begin{tabularx}{0.7\textwidth}{X} 他有兩個妻子：一個名叫哈拿，另一個名叫毗尼拿。毗尼拿有孩子，哈拿卻沒有孩子。 \end{tabularx} \\ \\ \relax
1:3 & \begin{tabularx}{0.7\textwidth}{X} 這人每年從本城上到示羅，敬拜萬軍之耶和華，向他獻祭。在那裡有以利的兩個兒子何弗尼和非尼哈當耶和華的祭司。 \end{tabularx} \\ \\ \relax
1:4 & \begin{tabularx}{0.7\textwidth}{X} 每逢獻祭的日子，以利加拿把祭肉分給他的妻子毗尼拿和毗尼拿所生的兒女。 \end{tabularx} \\ \\ \relax
1:5 & \begin{tabularx}{0.7\textwidth}{X} 他給哈拿的卻是雙分，因為他愛哈拿。耶和華卻不使哈拿生育。 \end{tabularx} \\ \\ \relax
1:6 & \begin{tabularx}{0.7\textwidth}{X} 她的對頭毗尼拿因耶和華不使哈拿生育，就常常惹她發怒，要使她生氣。 \end{tabularx} \\ \\ \relax
1:7 & \begin{tabularx}{0.7\textwidth}{X} 年年都是如此。每當她上到耶和華殿的時候，毗尼拿就這樣惹她發怒，以致她哭泣不吃飯。 \end{tabularx} \\ \\ \relax
1:8 & \begin{tabularx}{0.7\textwidth}{X} 她丈夫以利加拿對她說：「哈拿，你為何哭泣？為何不吃飯？為何傷心難過呢？有我不比有十個兒子更好嗎？」 \end{tabularx} \\ \\ \relax
1:9 & \begin{tabularx}{0.7\textwidth}{X} 他們在示羅吃喝完了，哈拿就站起來。祭司以利坐在耶和華殿門框旁邊的位子上。 \end{tabularx} \\ \\ \relax
1:10 & \begin{tabularx}{0.7\textwidth}{X} 哈拿心裡愁苦，就痛痛哭泣，向耶和華祈禱。 \end{tabularx} \\ \\ \relax
1:11 & \begin{tabularx}{0.7\textwidth}{X} 她許願說：「萬軍之耶和華啊，你若垂顧你使女的苦情，眷念不忘你的使女，賜你的使女一個子嗣，我必使他終生歸給耶和華，不用剃刀剃他的頭。」 \end{tabularx} \\ \\
[1ex]
\hline
\hline
\end{longtable}
$^{1}$曾經看過一本關於兒童的書籍.
名為《童言童語話圖告》.
當中提到的問題和一些有趣的插圖.
我覺得很有意思.
所以也很想在開始的時候跟大家分享.
說它都是我們在信仰過程面對的情況或者疑惑.
第一個問題他說.
為什麼要向神禱告?.
他的插圖是一個插頭.
就是吸塵機.
基督徒的祈禱就好像吸塵機的插頭插到電源一樣.
說到祈禱就是開啟了天上的能源.
因為神是那位無所不知的神.
祂不單止知道我們最深層的需要.
連我們不知道的事情神也能理解.
所以禱告使我們更認識父神.
也更認識我們自己.
第二個問題他說.
既然神已經知道我們要禱告什麼.
那為什麼我們還要禱告呢?.
這幅圖說禱告就是耶穌發明的一個交通工具.
他說禱告其實就是承認我們需要主帶領我們的生命.
邀請耶穌進入我們的心裡面.
我們不需要靠自己的力量.
我們可以完全來交托給神.
禱告之所以重要.
是因為禱告可以改變禱告人的心.
第三個問題.
他為什麼神不把我們所禱告的某些東西刺激我們?.
這幅圖說的是小傑的一個想像.
父母很驚訝.
但是小朋友很開心看到.
收到他所期待的一隻小象.
而這隻小象是已經經過訓練受訓好的.
變成一個大象.
其實都是說神對我們的生命有計劃.
他希望給我們的是最好的.
他給我們更加有智慧.
有時神沒有答應.
可能那些東西是會傷害我們.

$^{41}$又或者是會讓我們走錯路.
有時神沒有立刻立即去答應.
是祂要我們去耐心等待祂的時間.
今天經文記載到哈拿向耶和華的禱告.
在我們的人生當中.
有沒有曾經發出好像哈拿一樣的禱告呢?.
今天早上我們就從哈拿身上學習一心依靠.
剛才主席幫我們讀出了經文.
亦都交代了在二百年山地.
有一個二百年人命叫以利迦南.
是蘇忽的原孫.
托胡的真孫.
以利戶的孫子.
耶諾漢的兒子.
我們從家族的背景來看.
連祂的真真祖父的名字.
都很清楚地去記載.
顯然祂是一個有名望的家族.
身份非常顯赫.
祂有兩名妻子.
一名哈拿.
一名皮利南.
從牌序上來說.
哈拿是第一位的妻子.
但因為祂不能夠生育.
所以為了傳宗接代.
立了第二位的妻子.
接著經文就說.
按摩西律法規定.
祂每年要去到士羅這個地方.
去敬拜和獻祭.
按人頭計算.
每人能夠得育一份.
但因為哈拿沒有孩子.
但丈夫又愛她.
所以她能夠得雙份.
皮利南見丈夫對哈拿雙份的愛.
她就窒怒.
她就激怒哈拿.
在當時的社會.

$^{81}$沒有生孩子.
在那個這麼大的家族來說.
一定是沒有地位.
也有很多閒言閒語.
加上面對皮利南的說話.
哈拿內心裡必定有.
說不出的委屈和痛苦.
我們在經文裡.
第七節看到描述哈拿.
當時她為了這件事情.
她不開心到一個地步.
是她哭泣不吃飯.
但到第十八節.
描述到她走去吃飯.
面上不大受容了.
到底在這個過程.
她經歷了些什麼.
讓哈拿有這麼大的轉變呢.
以下從兩方面.
讓我們一同去學習.
第一方面就是從寄望得子.
到完全的信靠.
第二就是從她不被滿足.
到深中釋懷.
我們看到在經文裡.
第五節來形容哈拿不能生育.
是因為耶和華不使哈拿生育.
我們就會好奇.
為什麼耶和華要使哈拿不能生育呢.
在舊約傳統裡.
兒女是耶和華所賜的產業.
得兒女是上帝賜福的標記.
所以不能夠生育的婦人.
彷彿好像被神遺棄一樣.
第二就是連義房都有兒子.
自己又沒有兒子的時候.
一定不是丈夫的問題.
而是自己出問題.
所以多少哈拿的內心裡.
她都會疑惑.

$^{121}$是不是我不夠好.
是不是我不夠屬靈.
或者我不夠勁淺.
加上她每年上耶和華殿的時候.
比利拿都會激怒她.
在《新譜及譯本》是說年復一年.
《和修本》更加說年年都是這樣.
我們更加能夠去理解.
哈拿長期面對比利拿的刺激.
歧視的眼光和說話.
她一定流不少的眼淚.
難道上帝不顧念哈拿.
難道上帝看不見哈拿的仇虎.
我們看看經文這裡連續兩節.
都提到耶和華不使哈拿生育.
不是說哈拿無法生育.
而是耶和華不使哈拿生育.
強調哈拿的不孕是出於神.
有譯本更加翻譯是.
耶和華關閉了祂的子宮.
這個是神的計劃.
那到底上帝在做什麼呢.
我們看到每年獻祭.
丈夫都會給哈拿雙份的肉.
但也不能使哈拿心得滿足.
以利家甚至跟哈拿說.
哈拿你為何哭泣.
你為何不吃飯.
你為何傷心難過.
有我不是比有十個兒子更好嗎.
哈拿長期內心的傷痛.
是否真的可以藉著.
以利家那這幾句說話.
和雙份的肉就能夠除去和減少呢.
加上面對以防長期.
一些不好聽的說話.
這些都是觸及了哈拿心裡最在意的事.
往往人最在意的事.
都可能是上帝在考驗我們生命的時候.
就像亞伯拉罕一樣.

$^{161}$他老來得子.
上帝卻要他獻上唯一的兒子一樣.
我們再看看經文第十到第十一節.
哈拿怎樣祈禱.
他說萬軍的耶和華.
你若誰顧婢女的苦情.
顧念不忘婢女.
賜我一個兒子.
他心裡最在意.
最寄望就是求上帝賜他一個兒子.
我們看到追求得著兒子.
一個很重要的原因.
就是使哈拿心裡得著滿足.
甚至是挽回他在人前.
或者他在人生裡的價值.
但我知道哈拿不能夠生育.
是上帝的計劃.
因為上帝想哈拿將他生命的焦點.
價值和他寄望.
不是只是放在他所求的事情上.
而是放在他所信的神那裡.
上帝要哈拿的生命.
他的重心再一次從自己身上轉向神.
很多時候上帝容許我們面對一些困難.
說許我們禱告很久都沒有答案.
我想我們知道其實都不是偶然.
上帝是要引導我們的心去轉向祂.
就好像開始引燃第二幅圖畫和那個問題一樣.
既然神已經知道我們要禱告些什麼.
那為什麼我們還要禱告呢.
禱告之所以重要.
是因為當我們向神去禱告.
其實禱告就是提醒我們.
在我們所求的事情裡面.
耶穌我們的神.
才是我們生命裡面真正的大領和主人.
我看回哈拿的丈夫以利加拿的祖先.
其實在經文裡面剛才讀到.
有記載到他五代的祖先.
名字都很清楚地列出.

$^{201}$就證明到這個家族.
他們是一直夢耶和華的紀念.
同樣很多時候我們向上帝的祈求.
其實都是沒有錯的.
都是我們的需要.
但是若果我們祈求的事情不經不覺.
在我們心裡面佔一個很重要的位置.
會不會有時可能因此而高過上帝呢.
我們想一想.
我們的情緒有沒有會被一些事情.
或者覺得是一些我們很在意的事情.
常常牽動.
我們能不能夠升職.
或者我們會不會得到某一段的關係.
或者有一個大好的前途.
或者一個受認定的地位.
或者其他的事.
當我們生命的重心.
和我們生命的滿足.
是放在上帝以外的事情的時候.
這些可能就成為我們生命的主.
哈拿的心因為沒有孩子而不能夠滿足.
上帝藉著這個環境.
我們看到是上帝的計劃.
要重新調教哈拿生命的焦點.
哈拿的仇苦除了是來自皮利納之外.
更重要的是來自她內心最在意.
她要得著的事.
相信我們每一個人在不同的階段.
心裡面我們都會有一些.
我們很想追求的東西.
今天你最迫切祈求的事情是些什麼呢.
神是那位全知的神.
如果我們已經將事情藉著禱告去交托的時候.
我們所需要的神也都知道的時候.
或許我們就不需要去問.
為什麼上帝你沒有聽我禱告.
為什麼上帝你要容許我經歷這些困難.
為什麼你不直接將我所求的賜給我.
又或者可能我們會問.

$^{241}$為什麼你將好像皮利納這樣的人.
放在我們的生命當中.
而是要問上帝現在你要透過這些環境.
這些事情你要向我顯明些什麼呢.
要我學習些什麼呢.
我們看看經文第九到第十一節.
他們在侍羅學完之後.
哈拿就站起來.
我們看到這個動作.
其實就代表他心裡面已經做了一個決定.
他不讓自己停留在痛苦裡面.
他會重新轉向耶和華.
我們再留意第十一節哈拿怎樣通告.
他是萬軍的耶和華.
你是誰顧皮利的苦情.
顧念不忘皮利賜我一個兒子.
我必使他終身歸於耶和華.
不用剃頭都剃他的頭.
哈拿的禱告好像跟上帝進行一場交易.
他用獻上的兒子換得一個兒子.
今天跟我們的禱告會不會有時候有點相似呢.
主啊 如果你給我升道職.
我會翻團替.
主啊 如果你給我一份新的工作.
我會去服侍你.
主啊 如果你給我這些祝福.
我會奉獻十分之二給你.
如果你給我一段婚姻.
我就會怎樣怎樣.
不得不否認.
我們有很多時候.
我們的禱告好像跟上帝進行一場交易.
但是看回哈拿的禱告.
不是交易.
從撒母耳的上往後經文來看.
撒母耳三歲的時候他斷奶之後.
他已經離開了哈拿.
到了聖殿當中.
在撒母耳記上二章記載.
伊利加拿往拉瑪回家了.

$^{281}$那時候孩子就是撒母耳.
就在祭司伊利面前來事奉耶和華.
接著繼續形容.
他的母親每年為他做一件小抱.
他跟丈夫上來.
天練祭的時候才來帶給他.
所以看到哈拿.
他每年只有一次機會見到撒母耳.
所以從二章經文推算到第一章經文.
其實哈拿祈求.
反而令他失去兒子.
當撒母耳斷奶之後.
他必須離開自己而去.
那為什麼哈拿仍要做這個祈禱呢.
我們看到.
他從必須有孩子.
到他現在說我必須獻上給神.
這一刻哈拿將他渴望得兒子這個寄望.
轉向投靠在上帝的身上.
他不再只是想有沒有孩子這件事上.
而這一刻他更加渴望他向上帝獻上他的兒子.
要上帝得著榮耀.
今天我們最寄望的.
最在意的是什麼.
是不是有些東西我們還沒得著.
我們人生不感到滿足呢.
有沒有曾經因為一些我們得不到的事.
我們心裡很愁煩難過呢.
哈拿在這裡做了一個決定.
哈拿就站起來.
這個站起來是一個動詞.
「起來」行動的意思.
標誌著他的轉變.
他不再沉溺在他的情緒裡面.
而是他主動地站起來.
他要轉向神.
他要將他的苦情帶到去神的面前.
向耶和華來到去禱告.
說他今天我們正是好像哈拿一樣.
面對一些難處.

$^{321}$我們可以繼續的來到去受苦埋怨.
但是同時我們都可以好像哈拿一樣.
做一個決定.
就是「起來」禱告.
從寄望得著兒子.
到完全來到去信靠神.
第二方面想和大家分享的.
就是從不被滿足到心中釋懷.
我們再看第十二到第十八節的經文.
哈拿在耶和華面前不住地祈禱.
以利定睛看他的嘴.
原來哈拿心中密討.
只動嘴唇不出聲音.
因為以利以為他喝醉了.
以利就對他說.
「你要醉到幾時呢?.
你不應該喝酒」.
哈拿回答說.
「主啊,不是這樣的.
我是心裡受苦的婦人.
清酒濃酒都沒有喝.
但在耶和華面前傾心吐意.
不要將婢女看作不正經的女子.
我因被人激動受苦太多.
所以祈求到如今」.
以利就說.
「你可以平平安安地回去了.
願意識你的上帝允准.
你向祂所求的」.
哈拿說.
「願婢女在你眼前蒙恩.
於是婦人走去吃飯面上.
再不大受容了」.
我看到哈拿真的受苦很多.
心裡受煩.
但他沒有向人發洩.
他選擇的.
將他的痛苦誠實的.
帶到神的面前.
以利見哈拿的祈禱.

$^{361}$嘴巴在動.
沒有出任何聲音.
還以為他醉了.
以利聽到哈拿的解釋之後.
以利就說.
「你可以平平安安地回去了.
以利說完這番說話.
其實哈拿是不是確定.
上帝一定會賜給他兒子呢」.
我們知道答案當然不是.
環境沒有改變.
皮里娜仍然他都要去面對.
奇妙的轉變就是.
哈拿他怎樣來向耶和華去禱告.
讓我們再仔細看看這幾節的經文.
第九節形容到哈拿痛痛的哭泣.
他向耶和華去禱告.
第十二節形容到.
他在耶和華面前不住的祈禱.
第十三節他說.
他心中來到蜜桃.
第十五節.
他在耶和華面前傾心 討厭.
我們看不看到.
其實當哈拿他心裡面痛苦到一個位.
他那一刻很難受.
他只是單單想去到神的面前去祈禱.
從心裡面的深處去呼求主.
十六節說祈求到如今.
在NIV的譯本他說他一直在這裡祈禱.
有一個意思就是說.
他擺脫巨大的痛苦和悲傷.
這個就是哈拿禱告的轉捩點.
哈拿一直的祈禱.
他密密的祈求.
直至擺脫那些困擾他的負能量和情緒.
在十一節他開聲的祈禱.
他的呼求是說來宣認承認萬軍的耶和華.
是掌管他的生命的主.
他的意思就是要再一次承認.

$^{401}$主才是他的帶領.
我們看到在經文繼續描述.
以你定睛看他.
哈拿的祈禱還以為他醉了.
的確有的時候我們基督徒.
可能我們面對難處的時候.
或者在我們同事裡面.
在朋友裡面.
有時候我們選擇合上眼睛.
我們安靜.
我們可能不說話.
我們密密的來到他不說話.
我們只是想到神面前不斷的祈禱.
可能旁人不明白的時候.
還以為我們很傻.
我們看到哈拿在這一刻.
其實他應該沒有理會環境.
外在是怎樣.
他一心只是想向神祈禱.
他毫無保留的表達.
他心裡面一切的苦情.
只是想向耶和華來傾心吐血.
接著兩岸的奇妙的事就發生了.
接著耶和華藉著祭司以禮.
向他來說話.
你可以平平安安的回去了.
有些譯本是說你可以安心回去了.
願以色列的上帝.
允許他向你所祈求的.
在這一瞬間.
哈拿聽到以禮的話.
哈拿就立即回應.
願你的婢女在你眼前.
你到去蒙恩.
於是婦人走去吃飯.
面上不再帶愁容.
其實這段經文最吸引我的地方.
就是這一句.
於是婦人走去吃飯.
面上不再帶愁容.

$^{441}$究竟什麼令到哈拿的情緒.
能夠有這麼大的轉變呢.
他明明之前還是痛痛的哭泣.
但竟然可以得到祭司的祝福之後.
他就面上不再帶愁容.
在舊約的時代.
我們知道祭司的角色.
就是人與人上帝之間的橋樑.
一方面代表人向神的憲制.
另一方面也代表神向人作出一個越立的宣告.
我看到以你作為祭司.
當他了解哈拿原來不是醉酒的時候.
他只是從心裡面悲痛的發出祈禱.
他就給哈拿一個肯定.
你可以安心平安的回去.
願以色列的上帝允准.
你向祂所求的.
他祝願哈拿的禱告得蒙神的應允.
但我們細心地去想.
以你的說話只是一個祝福.
哈拿的禱告其實仍未得到應允.
他仍然不知道自己會不會有懷孕.
或者得著孩子.
但為什麼他的心情能夠180度轉變.
由此可見.
我們知道哈拿的轉變.
不是因為祈禱已經蒙英韻.
而是因為他將一切苦情向神來傾訴.
在他禱告親近神.
承認耶和華是他生命主的時候.
他的心開始轉化.
他重新恢復信心.
加上藉著耶和華藉以來的話.
平安已經進入他的心了.
他內心得著滿足.
從上帝而來的平安.
他得到釋懷了.
所以他在面上不再大受用.
他如今亦不再只是看自己的處境.
這一刻又沒有疑似.

$^{481}$已經不再牽動他的心了.
因為他最在意的事.
他可以放下.
這一刻他降服.
他甘心完全的交託給神.
他就可以走去吃飯.
面上不再大受用.
最近我為了一些事情去祈禱.
都給我很多煩惱.
因為為了搬屋的事.
都叫我很煩.
我在想究竟不撞兒子考試好.
還是不撞先生來出差.
因為我先生工作經常都要仿孫來出差.
但如果要不撞這兩樣事.
等到暑假的時候可能我又比較忙碌.
我就在祈禱究竟什麼時候才最合適呢.
然後我就在想.
還是不撞兒子考試應該是最重要的.
是最好的.
但要兼顧生碌.
又要兼顧收拾的時候.
面對工作這麼多忙碌的事.
我覺得很沉重.
所以在這個時候我就選擇.
先不要煩惱.
不如我放下這些事情.
我選擇安靜去到神面前去祈禱.
在一放下這些事情.
不再用我自己的方法去想的時候.
我說神啊你才是我的帶領.
神啊你才是我幫助.
我求你赦免我.
我以為自己可以處理很多事情.
當我這樣開始祈禱的時候.
那一刻我就感受到神很大的安慰.
神的擁抱.
神的愛能夠臨到.
那一刻我很大提醒.
就是神告訴我.

$^{521}$祂知道我面對的難處.
神也提醒我.
其實我有很多事情可以感恩.
原來在那段時間.
其實我忘記了妹妹也給我很多協助.
裝修安排的事.
那一刻我又開始不再去看自己的事.
我開始想起家人的難處.
家人患病的需要.
我開始為家人祈禱.
很奇妙的一個月之前.
突然先生有一天跟我說.
我那段時間的訪孫轉了日期.
我那一刻的心情開心到不得了.
也就是又不會撞我兒子考試.
又不會撞他出差的時候.
當時我很大的體會就是.
原來其實神可以直接來扭轉環境.
祂可以改變一些局面.
因為神從來沒有留下一樣好處.
不給每一個求祂的人.
但神的能力和恩典.
是更加能夠藉著我們的祈求.
叫我們再一次將主權交給神.
並相信神是有最好的帶領.
我再看經文.
次日清早.
他們在第十九到第十二十節.
次日清早.
他們起來.
在耶和華面前去敬拜.
就回到拉瑪.
到了加侶.
以利加拿和妻哈拿同房.
耶和華顧念哈拿.
哈拿就懷孕了.
日期滿足.
生了一個兒子.
就給他起名叫撒母耳.
說這是我從耶和華那裡求來的.

$^{561}$這一節開始就說次日清早.
這裡更加看到神祂的奇妙.
只要是上帝的時間.
其實神可以立即來工作.
如果上帝沒有等哈拿的心.
和他的生命的重心轉向神之前.
就賜給他兒子.
我們能夠估計.
哈拿可能非常之寵愛撒母耳.
甚至寵壞了他.
又或者因為撒母耳被哈拿的期待.
壓到頭不過氣來.
上帝很愛哈拿.
等到哈拿的心轉向神.
先給他來懷孕.
我們再看之後的經文.
哈拿之後再生了.
三個兒子.
兩個兒女.
耶和華人愈發喜愛撒母耳.
撒母耳漸漸長大.
耶和華與他同在.
使他一句都不落空.
我們看到當哈拿的心.
完完全全轉向上帝的時候.
上帝給他的祝福.
是超出他所想所求.
撒母耳日後也成為一位很偉大的先知.
我們回到.
在隱形最初那幾個問題.
為什麼我們要向神祈禱呢.
既然神已經知道我們要禱告的事情.
那為什麼我們還要禱告呢.
為什麼神不將我們所祈求的事情.
直接賜給我們呢.
今天盼望.
從哈拿身上我們學到一個重要的功課.
禱告是一個操練.
往往生命的復興更新.
是由我們禱告開始.

$^{601}$禱告首要改變的不是環境.
而是我們的心.
在參與週五的祖禱會裡面.
對我最大的幫助.
就是提醒我將焦點.
從我自己的需要轉回父神的身上.
讓我看到的不是我自己的需要.
有沒有被滿足.
而是我願不願意再一次向神.
來到獻上我自己.
在安靜裡面.
我能夠去認罪悔改.
我也期盼能夠藉著禱告.
再一次聆聽神的聲音.
以致使我的心能夠緊緊跟隨神而去行.
在這時候我想跟大家一起有個祈禱.
親愛的天父.
您不僅只是愛我們.
您也是認識我們.
我們坐下.
我們起來.
您也知道.
您從遠處已經知道我們的意念.
我們走路.
我們肚子餓.
主你都誓詞.
您更深知.
我們一切所行.
所想.
所需要的.
求親愛的聖靈幫助我們.
在面對難處.
和等候您回應我們禱告的時候.
環境或許未有改變.
但我們求您賜給我們有哈拿的決定和信心.
我們願意起來.
我們向您你去禱告.
因為您就是那位愛我們.
愛到底的神.
您的眼目從來沒有離開我們.

$^{641}$您有更美的計劃.
並且您會用著每一個環境.
使我們再一次靠近您.
求聖靈引導我們的心.
不要被環境影響.
讓我們能夠堅定緊緊去跟隨您.
一心依靠您.
我們能夠坦然來到您的面前.
向您去呼求,向您去禱告.
求恩主紀念我們的禱告.
讓我們經歷到哈拿的祝福.
成為一個蒙大恩的人.
我們同心祈禱.
奉主耶穌基督寶貝名字而求.
阿們.
\newpage



\section{歌羅西書 4:1-6}
\label{sec:031959}
\textbf{從歌羅西書看成熟生命表現}~~~~梁永善牧師
\newline
\newline
連結: \href{http://www.church.com.hk/acms/content.asp?site=cdc\&op=show\&type=product\&code=031959\&layout=sermon}{\texttt{031959}} ~~~~ 語音日期: 2026-04-12
\newline
\newline
\hyperref[sec:031958]{\small{< < < PREV SERMON < < <}}
~
\hyperref[sec:index_chronic]{\small{[返順時目]}}
~
\hyperref[sec:031960]{\small{> > > NEXT SERMON > > >}}
\newline
\newline
歌羅西書 4:1-6
\newline
\begin{longtable}{cl}
\hline
\hline
章節 & 經文 (和合本修訂版)\\
\hline
4:1 & \begin{tabularx}{0.7\textwidth}{X} 你們作主人的，待僕人要公正，因為知道，你們也有一位主在天上。 \end{tabularx} \\ \\ \relax
4:2 & \begin{tabularx}{0.7\textwidth}{X} 你們要恆切禱告，在禱告中警醒感恩。 \end{tabularx} \\ \\ \relax
4:3 & \begin{tabularx}{0.7\textwidth}{X} 同時，也要為我們禱告，求神給我們開傳道的門，能宣講基督的奧祕， \end{tabularx} \\ \\ \relax
4:4 & \begin{tabularx}{0.7\textwidth}{X} 使我能按著所該說的話將這奧祕顯明出來，我為此而被捆鎖。 \end{tabularx} \\ \\ \relax
4:5 & \begin{tabularx}{0.7\textwidth}{X} 你們要把握時機，用智慧與外人來往。 \end{tabularx} \\ \\ \relax
4:6 & \begin{tabularx}{0.7\textwidth}{X} 你們的言談要時常帶著溫和，好像用鹽調味，讓你們知道該怎樣應對每一個人。 \end{tabularx} \\ \\
[1ex]
\hline
\hline
\end{longtable}
$^{1}$我自己在過去的墓會的時候.
教會有洗禮.
當年我們教會洗完禮之後.
就會送一些東西給他們.
有三個選擇給他們選擇.
一個就是聖經.
一個就是舊約的附讀.
一個就是新約的附讀.
任由他們選擇.
無論這三種的禮物.
無論是聖經也好.
新舊約附讀也好.
我都會寫一些東西.
來祝福他們.
我寫的時候一定是用一段經文.
這段經文就在歌羅西書的第二章第六到第七節.
我們一起來讀書好嗎.
歌羅西書第二章第二章到第七節.
我們一起來讀書.
不好意思.
我做的拋光能力都不是很好.
這裡是考大家的眼力.
你看得不清楚.
你就可以再配一個眼鏡.
請大家描述我.
當我們信主的時候.
不是一個完結.
很開心我今天收到一封相片.
在Warrington裡面有一個老人家.
我上兩個月的時候.
他趕到利物浦.
他們家人就通知我.
原來他的爸爸很喜歡聽我的影視演講.
原來我去探望他.
他已經信了主.
就是在近期.
我剛剛上個星期洗了禮.
說親我.
我就鼓勵他.
記住洗禮不是一個完結.

$^{41}$是另一個開始.
六月的時候我都會負責一對年青人的婚禮.
都是會講這件事的.
結婚不是一個完結.
是另一個開始.
所以為什麼我洗禮的時候.
送禮給他們會寫這段經文.
就是你們要更加盡心去一個新的地方.
你們要怎樣呢.
就是要來到他們裡面三根建築.
你信心堅固.
你的生命才會成長.
但是生命的成長是會有表現的.
不是說你講就是了.
很簡單舉個例子.
小朋友他們的生命教育是有亂的.
就需要我們去照顧他.
但是他們很可能.
他們累或者他們有需要.
就會跟媽媽說.
他們不會理爸爸媽媽那時候的困擾.
但是如果你是生命成熟的一個成年人.
可能你雖然有需要.
但是你都會去體諒.
我發出這個呼求.
但是對方有沒有那個時間去處理呢.
所以原來生命成熟.
不是一個講字.
是有表現的.
所以我會跟大家從《歌羅西書》第四章裡面.
第一到第六節裡面.
就來看到究竟生命成熟有什麼樣的表現.
聖經是很好的.
它是對每一種不同的人.
都有一個教導的.
第一個就是一個生命成熟的主人.
他應該要做些什麼.
我發覺在我們當中.
很多人都是這樣.
基本上.

$^{81}$要麼你就是家庭主婦.
或者家庭主公.
因為在英國有些人.
是男性負責家務的.
女性出來工作的.
負責家務.
沒有出來工作.
要麼就退休了.
要不然的話大部分人.
要麼你就工作.
打工的.
要麼你就做老闆.
聖經裡面也提醒我們.
原來做一個老闆.
或者做一個工作者.
你應該持守什麼態度.
一個成熟的.
有生命力.
依從神豐富的一個老闆.
他要做些什麼呢?.
我們一起來看看.
《哥洛西書》第四章的第一節.
我們一起讀出來.
他說你們做主人.
要公公平平地對待那些打工的人.
就是說有些人很可能.
他很懂得討那個老闆歡喜.
懂得擦鞋.
很可能會得到那個老闆的重視.
喜好.
但是可能一個很勤奮的人.
甚至很可能.
覺得老闆有些事情做得不好.
提醒他.
那就慘了.
那個老闆就不喜歡他了.
但是這樣的時候就不是公平了.
他說你要帶去公平地對待每一個人.
因為我們的神都是公平和公義的.
是不是?.

$^{121}$這個公平還是指什麼呢?.
就是你要好好的.
不要欺負他.
我們先帶他去看《理美記》第十九章的第十三節.
我們一起讀出來.
就是說要發薪水給他.
拜託你就按時發薪水.
特別是那時候有些機械低下階層的人.
在舊時代也是.
他只是日薪.
就是他所說的散工.
你每天就要怎樣呢?.
他做完事情就要給他工資.
不可以扣起他的工資.
如果是這樣的時候.
這是觸犯了神所定的法例.
我們再帶他去讀《申命記》第十四章十四,十五兩節.
這裡不單單是指你自己的同胞.
你的外勞.
你也要好好的去對待他.
我們一起帶他去再讀.
《申命記》第十四章十四,十五兩節.
(念經).
很清楚.
如果他是做完工作.
你就要按時發薪水給他.
不應該扣起他的薪水.
如果是日薪就每天發薪水給他.
他說如果你不發給他的時候.
他就有困難.
他會向神呼求.
神就聽他的禱告.
會支持你的.
當我們過來英國的時候.
我女兒剛剛一來到英國.
就開始找工作.
做一些散工.
在我們住的區裡有一間外賣店.
只是做禮拜五,禮拜六的事情.
是禮拜五,禮拜六最忙的.

$^{161}$只是日薪而已.
已經是最低工資了.
但我女兒就跟我說.
做完工作了.
應該要發薪水了.
老闆娘就站在收銀機前面.
等她發薪水.
哎呀不要忙.
你遲點叫她來拿吧.
不是馬上給的.
那他也要坐車過來.
那時候他又不懂得駕車.
不是一天是這樣.
經常都是這樣.
我想他其實是沒有事情做的.
他就這樣去拿錢.
也不是很多.
做幾個小時.
後來我跟他說.
不要再去泰國了.
我暫時都養得起你.
是呀.
不要給他這麼多錢.
雖然過幾天他去拿.
但我覺得是很過分的.
你不應該這樣對他.
所以那時候.
現在有些人就這樣.
差一點就不要跟他這樣.
跟大家傻.
我發覺教會也會是這樣.
因為有些教會.
我來英國是很得意的.
是很不好意思.
要我去追港元費的.
我講完道之後.
你們很好的.
你們馬上給我的.
有些是沒有給你.
我等.

$^{201}$三個星期.
四個星期都沒有.
於是就問.
你們會不會是忘記了.
或者是你給了我.
我不見了呢.
是呀.
忘記了.
我們會叫師父轉給你的.
不是一間.
是有幾間教會都是這樣.
但我講道也不是為了港元費.
但其實如果這個原則.
你照樣套下去的時候.
其實教會是有問題的.
所以教會都應該要很清晰.
就是這些的事.
而最貴就是有些人.
你看到現在在一個社會裡面.
不是打壓英國.
英國反而沒有那麼多.
你問香港裡面.
或者國內.
有很多人就剝削那些工人.
我們怎樣做呢.
就是不給正規的工資.
甚至我見過.
有些基督徒很不幸.
還是教會去執事.
去請人用.
但是你知道政府會定了那個工資.
它沒有依足那個工資給人.
多給一點.
這幾年是要多給一點.
就是要搶.
但十幾年前就不是了.
它是少給一點.
我要給多一點工資.
它拿回印尼.
已經很好了.

$^{241}$這個就是剝削.
其實是犯法的.
但有時基督徒都犯了這個錯.
但一個真正成熟的人.
是一個遵照神吩咐的人.
絕對不會做這些事.
我們再讀最後耶利米蘇第二十二章第十三節.
我們一起讀出來.
哪吾以為奉行不公諸侯.
罷使人不公.
不公家的猶我.
這裡很清楚說.
請人來幫你家建屋.
但是你竟然走數.
這是神要審判的.
但這個耶利米的年代.
在香港照樣有這種情形.
你見到很多利盤.
大公司判出來大判.
大判給二判.
二判給三判.
勞動體力.
甚至有時找人剝削.
最近你有沒有留意到幾個的天秤.
跌下來壞了.
被人自發了.
原來是拿舊的料.
然後安新的招牌上去.
所以力氣不夠.
結果就死了人.
這時候這些人有血有汗的.
但竟然是被剝削.
做得這些體力勞動的階級.
一定是很辛苦的.
但你都剝削他的時候.
有汗出沒糧出.
用廣東話是很節度的.
還有最近大家知道.
去年大班.
西餅.

$^{281}$還未結束.
你開油鹽公司其實都無所謂.
但老闆自己是有豪宅的.
就說不要緊.
有勞工保障.
你們就向政府申請.
海蜈蚣粥麵店.
我都很光顧過.
那時豪宅怪幾好食.
但問題是怎樣.
又是同樣的情形.
所以這是一個很大的問題.
我們仍然相信神是公義.
我們應該要呼求神審判這些.
致肥而不剝削的人.
如果我們作為一個老闆.
我們可能開店.
我們會不會犯了這個錯呢.
要很小心.
基督徒這三個字我經常說.
沒有任何保證的.
一樣可以犯錯的.
聖經裡面很清楚講.
工人得工價是應該的.
我們是應該要給他.
但是聖經裡面是會兩方面的.
每節是講做老闆的.
做工人的會怎樣呢.
一個成熟有生命力的工人.
他應該有甚麼態度呢.
我們一齊來讀歌羅西書.
第三章22到25節.
我們一齊來讀.
您要守好人的.
以你的眼前的.
和你的目中所指的.
以你的視窗的眼前的視圖.
將使好人圓滿的.
從你的視窗的眼睛對著.
無論做甚麼都要從心去做.

$^{321}$將使好處做的.
不是不能做的.
因為你的視窗.
擋住了好人的眼前的.
以你的視窗的眼前的.
大視主的地圖.
難為以敵.
受受得起的報應.
以平凡的運作.
記不記得這一段.
如果主人行不義的時候.
實際會審判主人.
但是你不要只要求主人要.
叫我們準時出來.
我們都要檢討.
如果我們做事的.
我們用甚麼態度去工作.
一個很重要的原則.
就是要從心裡做.
不要只是去做表面的工作.
有些人很適合.
上司在這裡.
老闆在這裡的時候.
這樣就很勤奮.
但是一不在的時候會怎樣呢.
在英國.
很容易可以犯這個錯誤.
因為大部份是.
Work from home.
Work from home的好處在哪裡呢.
就是你表面在做事.
我見過啊.
我去探訪的時候.
是Work from home的.
但是原來開著另一邊.
在看綁架.
在看電視.
因為沒有人監管.
但是聖經告訴我們.
原來我們不只是要向人交代.

$^{361}$我們是要向神來帶來交代.
因為是監察我們的.
是我們的主.
剛才你唱最後那首的詩歌.
是取材自詩篇139篇.
我自己都剛才彈了一整場.
我很少唱歌.
其實很喜歡唱歌.
因為我唱歌.
大家知道是有些創意的.
我就不想侮辱那首歌.
但是剛才我都唱得很大聲.
因為我都很觸動那首歌.
我怎樣可以留下神的體會.
你想想現在在這個社會裡面.
特別是香港.
有多少的閉路電視人員辨識.
有人跟我說現在回大陸.
你掉了錢包.
都不用怕.
就是不會有人撿.
是不是啊.
真的這麼真.
是啊.
因為現在很多閉路電視.
隨時拍到你.
撿了別人的錢包.
我以為在日本才會有這種事.
我不知道這個真的假的.
你想想人都可以這樣來有監控.
我們的神是無所不知的.
你表面裝作做勤力.
但不是的時候.
可能你騙到你的上司.
但是你騙不了神.
我們要向神了交代.
我自己作為一個牧者.
當年我也是這樣牧起小孩.
都跟父母說.
你不只是向那裡執事交代.

$^{401}$你是向神了交代.
你不要在大陸裝作做事.
其實某些事情我們都知道.
所以今天.
如果我們每一個人都向著.
有一個心態.
我做每一件事.
不只是為對方.
不是為上司.
而是我要向神交代.
你整個人的視角.
是完全不同的態度.
我記得.
我有一個回憶.
我很欣賞.
在政府.
包括部門裡面做事.
原來在政府裡面.
他們那時候.
是有彈性工作.
就是上班時間.
有些是8點15分.
有些是8點30分.
最遲的是8點45分.
原來那個年代是要簽到的.
但是他們是很好的.
行下一氣的.
有一個術語就是.
大家用廣東話叫「蛇鼠一窩」.
是什麼呢?.
譬如回到8點30分.
遲回來.
但是要簽到的.
8點半那個回來的時候.
大家看到他.
他就是上面那一行留空了他.
他就簽在那裡.
等他回來的時候.
就簽到上面.
你明白我的意思嗎?.

$^{441}$簽到假時間.
大家都可以.
但是我這個基督徒.
基督徒就被人覺得他成功不快.
他一回去的時候就簽名了.
他又不會留空上面那一格.
那些人不應該早過去的.
那就慘了.
你簽了.
你8點35分回來.
但是他應該回到6點79分的.
但是他8點35分已經回到了.
但是8點35分都還沒回來.
那你簽到後面的時候.
他不可以說謊的.
我經常跟他說.
你是不是搞錯了?.
你怎麼不留空位給我?.
他說沒有.
我真的回到那裡了.
好了.
9點鐘設置.
才對市民.
他們就會馬上開了電腦.
但是繼續.
他們佈置吃早餐.
但是已經是上班時間了.
好了.
如果上司來.
就馬上不同了.
你明白嗎?.
但是這些就是表面的工作.
所以就被人討厭.
這些人罵他.
搞錯了.
一定有一星的.
做了這麼長時間的時候.
他真的當時都不偷吃.
他們真的離譜了.
這句說話真的有問題.

$^{481}$有誰喜歡單笑偷食?.
我不會相信有人會做這件事.
結果就被人杯葛了.
所以後來他申請調另外一個員工.
後來他跟我分享.
當他調另外一個員工的時候.
工作量是重很多.
辛苦很多.
但是他做得開心.
因為每個人都是這樣做事的.
就不會這樣被帶去蛇鼠一窩.
所以今天很重要的就是.
我們記住.
無論是牧者還是工作者.
我們做事的.
不只是做表面功夫.
表面功夫可以騙一個人一輩子.
但是長時間你都騙不了他.
別人是會知道的.
正如我當年.
請工人的時候.
師傅都要上班.
我們要知道哪些工人是好哪些工人是不好.
有一個工人.
我們出街.
都不用帶鎖匙.
因為有鎖匙他可能很悶.
他馬上在街上摔.
就是摔著.
我經常跟他說你先不要摔著.
你摔著了我不知怎樣處理.
怎樣幫你.
為什麼他要摔門呢.
因為一定要按門鈴.
你按門鈴他就知道你回來了.
好了.
我回家的時候我記得.
他打開一個燙板.
有個燙斗在那裡.
都燙了三件衣服在那裡.

$^{521}$還有一件衣服在那裡.
那時候我按門鈴.
我都多謝他.
幸好他關了電掣.
我回家的時候他脫光浸水.
我摸一摸那個燙板.
是冷的.
然後我再摸一個燙斗.
都是冷的.
但我再摸一摸家裡.
那時候電視機是大牛龜.
後來拉手.
還有一張梳化椅.
梳化椅的皮.
跟現在的皮差不多.
是沒有彈性的.
我看到梳化椅的位置.
是鎖上了的.
即是什麼呢.
他開電視看.
他趕不及調整電視台.
如果我開電視機.
是英文台來的.
即是說他會燙好衣服.
好了你出去.
他就休息了.
他回來的時候.
他就會假裝工作.
我不是不知道.
我知道.
但我都不敢揭穿.
因為有人知道在他手上.
我的電視台很小.
但我不會用這些工人.
兩年後再換.
阿May有一個非常好的基督徒.
不讓我有壓力.
他會留意我有沒有靈修.
一整天不見你讀聖經.
我很害怕.

$^{561}$是工人問我有沒有靈修.
真是很大件事.
但是很OK的.
即是你做得很好.
也是一個好的基督徒.
阿May我太太不用再上班了.
所以我們沒有用她.
我們都用完了那個合約.
我轉介給我的執事.
用了很多年.
阿May很幸福.
嫁了去多倫多.
我剛剛去年十月多倫多的時候.
我都繼續聯絡你.
很幸福.
就貓在你肩膀那裡.
video call.
一個好的工人.
到現在我們都保持聯絡的關係.
所以不要想著你的上司.
不會被你騙到.
可能某個程度.
他暫時不會動你.
所以聖經裡面很清楚說.
原來我們作為基督徒.
我們所面對的.
不只是我們的上司.
你成熟的時候.
就是你要怎樣呢?.
我要向神有交代.
監管的是我們的神.
有一句話我很喜歡.
It's not the hours you put in that account.
but what you put in that hours.
不是你做了多少個小時的事.
而是那幾個小時裡面.
你做了些甚麼事.
明白我說甚麼嗎?.
很簡單舉個例.
當年有個童工.

$^{601}$傳統童工來的.
他很喜歡搞音響.
經常都去搞音響.
幫教會又談這又談那.
我說你不是負責做音響的.
你是負責關心人的.
你不要放這麼多時間在搞音響.
做牧者最重要是要做甚麼?.
就是要去關心弟兄姊妹.
你用了多少時間關心弟兄姊妹?.
可能你是不斷在做事.
但可能你做你自己喜歡做的事.
但不等於你是要去做的事.
所以也要求神來幫助你.
懂得怎樣去面對.
還有.
當我們做事的時候.
我們會不會.
剝削我們的老闆?.
剛才就說做老闆不可以剝削工人.
但我們做工人會不會剝削老闆呢?.
有個調查.
就是很多時候人會在公司時間裡面.
做私人事.
我甚至很多年之前.
即是那時候的會議.
我們教會的執事做文書.
他出來的會議記錄.
是有皇冠在頂上那張紙.
不是這個意思.
是用政府的資源.
那些紙.
你帶去做筆記.
可能是公司的時間.
我跟他說.
第一就是你不應該拿政府的紙來用.
甚至有時候會拿到文具回家用.
第二我說你不應該在你辦公時間做教會的工作.
雖然你說是為神做事.
但都錯.

$^{641}$不可偷渡.
某個程度你是在偷竊公司.
他們批的員工給你.
這就是違反了十戒令.
不可偷竊這件事.
所以今天你說他剝削我.
我剝削下他.
你記不記得我剛才說.
那個人不是即時出糧給我的女兒.
我叫他不要做.
如果你覺得他剝削你.
你不應該給他剝削.
你就辭職.
不要做.
但為什麼你不辭職.
因為我還沒有第二份工作.
那就剝削他.
這就錯了.
當年我在香港的時候.
都有很多機會去一些學校裡面.
講一些老師培訓人.
我很深印象一間基督教的學校.
在將軍那裡.
是講老師培訓人.
比率上幾元.
老師是教徒來的.
校長一開始會坐在那裡.
為什麼呢.
唱完歌之後.
校長就要弄一隻錦旗.
給我拍一張照片.
完了之後校長就走了.
校長一走的時候.
這班老師.
好.
很多時候怎樣.
看手機.
即是做其他事.
我都很生氣.
我跟他說.

$^{681}$如果你上課教學生.
學生是這樣的態度.
你會怎樣.
你作為一個老師.
你現在這樣是很錯的.
你怎會要求你學生這樣.
我說是.
我會明白的.
你們這麼多的工作壓力.
況且我還要上那些老師培訓班.
其實講些什麼你都要知道.
是很沒趣的.
但是你做得這樣.
你就要好一點.
我說如果你不喜歡的.
你就辭職.
就不要這間學校.
為什麼你不辭職.
因為你還沒有找到另一個更好的工作去做.
但我說如果你不想被帶入去.
即是你不願意.
但你都要承擔責任.
《聖經》裡面的弱勢就是一個最好的例子.
弱勢是不自願地被賣成為奴隸.
但是他不會棄罪你去做.
他被迫做奴隸.
但在柯提法的家裡.
是非常認真的工作.
所以被柯提法重視他.
柯提法是皇帝總管的護身的頭.
所以他很敏銳觸覺.
所以他要把所有的工作交給弱勢去處理.
弱勢沒有犯罪而要坐牢.
他是不自願去坐牢的.
但是他亦都幫獄長打理得非常好.
所以獄長都怎樣呢.
就是把所有的工作交給他.
你記住他們兩個的處境.
都不是他願意去的.
但他仍然盡他的責任.

$^{721}$他的焦點就是他向著神.
記不記得柯提法的太太.
喜歡小鮮肉.
每天向著這個小鮮肉.
敲明眼.
但是弱勢怎樣說.
後來還要拔她.
要和她發生不正常的關係.
我怎可以犯這樣大的罪.
得罪神.
他不是說我得罪我的老闆.
是他這樣犯當然會得罪老闆.
但他背後就知道我是這樣得罪神.
所以一個成熟基督徒.
應有的生命表現.
如果我是工作者.
我就會盡我的責任.
因為我不是只是向上司交代.
我是要向神來帶去交代.
如果你覺得他剝削你.
你不讓他剝削.
你便走.
但還不能走的時候.
仍然要堅持做得好.
這個就是聖經裡面給我們的提醒.
一個成熟生命基督徒的表現是甚麼呢.
我們再讀這段經文.
很短的.
《哥羅世書》第四章第二節.
我們一起讀出.
「行徹禱告,謹聖感恩」.
兩者是有關的.
「行徹祈禱」.
不是有事才去求神.
今天很可悲.
很多很幼嫩的基督徒的生命.
將神變成另類的菩薩.
是甚麼意思呢?.
人們怎樣拜菩薩呢?.
是有事才去求的.

$^{761}$臨急抱甚麼?.
佛腳.
我們是臨急才抱主腳.
有事才去求神.
將神當作菩薩.
放在一旁備用.
有事便去找祂.
但絕對不是.
其實很渴望和祂有好的交往.
我剛過去的星期.
是第一次全家一起去旅行.
和我的孫女.
我們有很大的感受.
你可以上去聽全家人的歌.
因為和我兩個孫女.
溝通愈來愈少.
因為他們五年前已經來了.
大家全部都在說英文.
她很少和我聊天.
她不懂廣東話.
但當我看著她們時.
我都很開心.
看到她們時很開心.
她想和我說幾句話.
我已經想哭了.
但她最喜歡說.
「爺爺,我喜歡這個公仔」.
我們的神是很樂意.
很享受和祂有交往.
但我們沒有祈禱.
祈禱不僅是求.
祈禱是溝通.
所以我們要行切禱告.
不是有事才行.
在此警醒感恩.
為何祈禱令我們警醒呢.
我們一起讀多少前書.
第四章第七節.
我們一起讀出來.
.

$^{801}$彼得說了三方面.
在世界中國將來的時候.
很多的誘惑.
有三件事要說.
第一是自己謹慎自守.
警醒禱告.
第二是要切實傷害.
第三要好地事奉.
這三節經文是對證.
耶穌港馬大綱二十五章末世的三個比喻.
因為你禱告.
你會與神越來越親近.
你就會越來越敏銳於.
神在你身上的計劃,思維和思想.
你就會順著神的心意去做.
當你有警醒的時候.
你便會感恩.
我分享一件事.
剛才不是說我和我家人一起去旅行.
我們去的時候.
全部都是坐泰航.
要轉機的.
但我回程.
我就用我的James Irons.
換了兩張飛機票.
坐英航回來.
我想我們全部都在成年機場.
到最後那一天.
我應該是1點05分的飛機回到英國.
已經收拾好了.
早一天已經上網已經check in了.
給了聲明提醒.
那天是早上8點45分.
我想究竟我們的行李是23KG.
還是25KG.
我再click進去看.
一看就是.
不是.
我原來不是在成年機場離開的.
我在雨田.

$^{841}$但那時候我還在成年機場附近的酒店.
不是吧.
糟糕了.
那時候已經9點多了.
怎麼搞的.
如果我再坐半小時去成年.
然後再去成年去雨田.
已經趕不上了.
於是.
馬上要call uber.
不管了.
總之由哪輛車最靠近.
是一輛小車.
我們兩個人.
四件大行李.
因為英航每一個人有兩件23KG行李.
再加兩件scan carry.
大家可不可以想到一個普通的的士.
可不可以等四件大行李.
再加兩件小行李.
終於都迫到了.
嘩.
千幾元港幣.
終於去到雨田機場.
剛剛趕得及check in.
所以後來回到英國的時候.
我們是第一個拿到行李的.
因為我們是最後拿到的.
雖然用了一筆錢.
但我要感恩.
我一直都沒有醒覺.
原來我是坐雨田機場的機班人.
我覺得這個是神提我.
早上8點多.
如果不是重大事件.
去到.
我初時還以為在城內還可以去吃壽司.
然後才去.
來到Terminal 2.
很划算可以吃壽司.

$^{881}$原來一些這樣的事.
我要去感恩.
原來是神的力量.
所以每一次我坐飛機的時候.
長途機.
多點時間去祈禱.
為什麼呢?.
第一就是在飛機上沒有事要做.
第二就是我坐飛機去祈禱的時候.
應該神收得快很多.
因為是復活的.
所以我們今天有沒有能夠去.
行切禱告.
以致我們警醒感恩.
但是很可悲地.
今天很多基督徒.
都知道祈禱是什麼一回事.
知道祈禱有好處.
但是我們有沒有好地樂於.
去禱告呢?.
有時候我們輕忽了.
我們禱告的對象.
是一個尊貴榮耀的君王.
我們輕忽了我們禱告的內容.
有時候隨便出來說.
但最大問題是我們輕忽了.
我們自己的需要.
我們可能只是求一些表面的東西.
只是求平安.
求健康.
求五加一順利.
縱使不用給快訊.
都可以三日批准.
可能我們就會求這些.
但我們說這些不行.
也應該要求的.
正如耶穌教我們祈禱.
我們每天用飲食給我們思考過門.
這是實際的需要.
但我們有沒有求我們的靈命要成長.

$^{921}$我們有沒有求.
我們要對神有更深的認識.
所以作為一個基督徒.
我們知道.
祈禱的原則.
原來不單單是我們個人的作為.
然後我們是跟神一起合作.
所以請你記住.
原來祈禱不單單是求.
是跟神建立關係.
一個成熟生命的人.
是會行切禱告.
然後警醒感恩.
你越行切禱告.
你的生命自然就會更加成熟.
因為直接禱告.
你可以內省你自己的生命.
在別人面前騙不了神.
很多人讚賞你.
但你是否真的值得被讚賞.
神對你的評價是怎樣.
你透過禱告你會知道.
不要再被別人.
甚至自己騙了自己.
所以求神有幫助你.
還有就是當我們去禱告的時候.
我們將你的焦點不再只是放在我們自己的事上.
你有沒有留意到你們今天的程序表.
有甚麼代禱事項.
我欣賞你們.
這三天至少.
我看到這三個代禱事項.
我是感恩的.
我們黃福縣聽證會.
看到很多的部門.
互相推廣責任.
其實這星期雷扇頭講的這些.
對一些黃福縣的災民.
是一個很大的傷害.
第二為了烏克蘭.

$^{961}$四年了.
要求平安.
第三就是美伊之戰.
我們究竟只是打在程序表.
還是實力有沒有去做.
我們的禱告是不是只是為了自己.
尤其是我們有這個擴闊的視野.
這是一個成熟一副的表現.
當我們禱告的時候.
我們會內省.
去做給世界的事.
你帶我去甚麼.
你帶我去困擾的場地.
不是只是祈禱.
是為自己.
一個成熟生命的人.
是會為別人來禱告的.
正如我一開始講.
一個嬰兒不會去關注別人.
不會去關注別人的需要.
但是一個成年人是會的.
一個成熟的基督徒.
不會只是為自己祈禱.
是會為別人來禱告的.
我們再來讀《哥羅洗書》第四章第三第四兩節.
《哥羅洗書》第四章第三第四兩節.
《哥羅洗書》第四章第三章第三章.
《哥羅洗書》第四章第三章.
我很喜歡這一段經文.
在啟示錄裡面在三章第八節.
我們一齊讀出.
我思考著你的行為.
你仍然繼續承擔.
假想相信我.
我仍然繼續承擔.
可能我在你面前.
曾考驗過我當中的思想.
我還未明白.
在啟示錄裡面.
有七間教會的書上.

$^{1001}$有兩間教會是完全.
責備而沒有讚賞的.
其中一間大家都很熟悉.
就是「無底家教會」.
但有兩間教會.
是完全有讚賞.
但一句責備都沒有.
其中一間是「善妹娜教會」.
但最短的得四個字.
第二間就是.
「飛拉鐵飛」的教會.
是在邊界.
是非常弱小.
用人體來看.
是沒有什麼能力.
但我來到時就說到.
雖然你們只是略有一點力量.
其實很微小.
如果給「無底家教會」.
他們是微不足道.
但真是不有心.
我就為你們開一個傳動的門.
開你們的門.
關你們的門.
這個就很無謂.
只要你有心.
神就會給你機會.
我最深看不到的就是.
只要你有這個心.
神就會給你機會.
很簡單的.
應該四月最後一星期.
我帶隊的意思.
但現在打仗不能去.
就取消了.
但原來神都很好.
當我決定取消的時候.
幾個小時後.
我收到北愛爾蘭一位姊妹.
會找我一年多.

$^{1041}$原來我們北愛爾蘭.
開了一個「香港人團契」.
我們看看有沒有時間.
過來港島.
原本就一定沒有.
因為已經滿座.
我們飛過北愛爾蘭.
當然有時間.
但剛好我取消了.
就五月第一個禮拜.
五月第一個禮拜.
他們禮拜六聚會.
我就禮拜五去港島.
禮拜五還有一個.
做外賣的.
星期六那個就回禮拜五那個.
星期六我們在設新人教會.
那時候.
我問新人教會的教徒.
他們說OK的.
我用英文講道.
他們用國語講道.
聽到英文講道應該OK的.
但我說.
我不會了.
為甚麼.
我們北愛爾蘭.
那天我去南愛.
南愛教會我都熟.
是一個主日夜.
不同的國度.
原來.
當你願意為神.
你就擺上.
神就為你開路.
雖然要自己自費過去.
但是祝福你.
所以今天保羅就說.
你要為我祈禱.
讓我有更多機會去傳福音.

$^{1081}$保羅已經很熱切.
去傳福音.
已經左撲右撲.
但他並不滿足於他所做的事.
他希望有更大的托子.
這也是我自己.
今年年初的祈禱.
我希望我有一些.
補充能力去我的教會.
可以讓我獨善無欲的人.
擴闊我自己屬靈的領域.
神就是這樣特定的.
所以.
祈禱我們有沒有祈禱.
神是給我們教會.
可以有更多機會去傳福音.
不單單這樣.
還有就是求神.
我們聆聽者.
有宣講的能力.
保羅就說.
你要為我祈禱.
以致我能夠去講基督的奧秘.
這個很重要.
不單單是去講道.
這個是很重要.
今天原來我們需要養皇神.
我自己很清楚一句說話.
因為神的國不在乎言語.
乃在乎權力.
是靠著聖靈.
我們一起去猜度.
第一章第十到第十一次.
我們一起去猜度.
保羅很清楚.
他的宣講不是為了討人喜悅.
今天有很多人.
有時見到.
有些人喜歡聽這些.
那就好了.

$^{1121}$我五月第二的星期.
應該是母親節.
我已經答應了.
一間教會去講道.
但保羅早不及跟我說.
保羅那天你不要講母親節.
好嗎?我說好啊.
你想講什麼?我說講奉獻.
講奉獻?OK.
好啊沒問題.
很多教會的人請我講奉獻.
我問為什麼是我說的?.
他說.
你在這裡就是敏感.
敏感要我講?.
他說是啊.
你是港元嘛.
你講完就是港元.
沒問題.
但如果我們在這裡.
那時候有些人很喜歡.
我自己說我當年.
我每兩年四月.
第一個星期或是最後一個星期.
一定是講奉獻.
我一定會好一點祈禱.
為什麼我不可以教好一點奉獻?.
但有些人就怕.
講這裡可能敏感.
甚至在英國講道.
我都被人審查那些講章.
那些不行這部分可以講.
太敏感了.
但保羅說.
我講道不是為了討人歡喜.
是要你怎樣呢?.
按照聖經的真理.
分界真理的道.
信息是有時代感.
所以保羅說請你.

$^{1161}$為我去禱告.
天下都為我祈禱.
今天星期三開始.
我查時思記.
因為今早六點多.
已經起床.
我在我妹妹家住.
我今早繼續.
做時思記那些筆記.
第四章叫三思會.
即是講三個時思.
我到你家去做.
不容易.
原來是時思記.
因為大家知道現在.
用語音來講.
但他沒有這些聖經的terms.
古,山,地,薩,田.
甚麼來的?.
底波拉.
大家都會知道.
底波拉即是看波.
無謂的.
很多這樣的東西.
但我都要認真去做.
以及不容易講.
所以.
都要求神.
叫我有宣講的能力.
所以求神幫助我們.
承受基督徒行切禱告.
謹慎感恩.
為別人祈禱.
不單單為別人祈禱.
為牧者去禱告.
可以打開傳福音的門.
祝福他宣講的能力.
最後.
承受基督徒是懂得管理自己的生命.
管理自己的世界.

$^{1201}$我們一起讀.
歌羅西書的第四章.
第五節.
我們一起讀.
要愛惜剛恩.
用智慧與為交往.
時間是非常寶貴.
是不可以浪費的.
我們有沒有善用.
我們的時間?.
這個就是好做的.
每個人都有24小時.
但有些人是懂得.
管理自己的時間.
很簡單.
我昨晚上來化名漢.
住在梅友家.
化名漢的火車站.
然後轉去另一個地方.
不知道在談什麼.
15分鐘車程.
剛剛錯過了班車.
看著山路.
等下班車.
15分鐘後有班車都無所謂.
一上了車12號月台.
New Street.
走了不夠幾秒鐘就停在那裡.
後來就說.
前面有小火所以我們要停在這裡.
弄了整個小時.
都開不到.
竟然還要破壞另一個計劃平台.
再轉另一個班車.
本來15分鐘車程就弄了一個多小時.
在英國坐車是這樣的.
時間就沒有了.
但我不會.
我有書.
我在這裡看書.

$^{1241}$做筆記.
不只是在這裡等.
所以我們要怎樣.
求神給我們智慧.
去善用自己的時間.
還有.
要有智慧與外人交往.
什麼叫與外人交往.
不相信那些人.
可能他們經常罵我們.
講假了沒什麼可以做.
大家開台吧.
幫忙一下智力上的運動.
我經常說.
不應該打他.
因為耶穌愛麻雀.
還要打他.
沒有問題.
有時時間就這樣浪費了.
今天.
我們不是不去享受.
不是不去休息.
問題是你怎樣去用你的時間.
你用你的時間.
在電視上網.
那個比例和你親近神.
或者做一些事.
那個比例是怎樣.
我們一齊了讀傳道書.
第十一章第八章第九節.
我們一齊了讀.
有沒有特別提到聖經說.
少年人.
你喜歡做什麼就做什麼.
因為少年人經常浪費自己的時間.
浪費得起.
都沒有問題.
但不要忘記你做的事.
有一天要向神交還.
我們去到這一百年.

$^{1281}$那個感受是不同的.
過一天就少一天.
你不知道還有多少日子.
所以我會珍惜.
每一天.
所以我會珍惜每一天.
希望不會浪費.
這個曾經在大家當中.
分享過.
天真風趣那個.
大家都記得.
到我們六七十歲的時候.
就已經風濕不痛風雨飄搖.
我們有沒有好一點.
來帶去善良的時間.
不要想著一定很多的機會.
你看看一個.
醫生.
適值壯年.
突然間跌倒.
這樣就死了.
一個預備入大學的年輕人.
去赤柱做義工.
竟然會摺一棵樹.
跌下去.
勾頭.
摔死了.
所以今天我們要有智慧.
去選擇.
我們要選擇去做什麼.
可能是兩樣東西都很正確.
但你要去選擇.
哪樣東西是更加有效.
正如當耶穌.
要被釘十字架之前.
瑪利亞找了高高耶穌.
尤特曼罵祂.
說拿錢這麼浪費.
拿來救濟窮人是多好.
其實耶穌沒有否定這些窮人的需要.

$^{1321}$祂說常有窮人和你不同在.
只是你不常有我.
祂說為我安眾之人.
你帶去做.
所以我們要在正確的事上.
你帶去.
怎樣去選擇.
剛才.
我在會友.
即是在會友家.
跟他說.
現在正在找當地的教會.
有很多香港人.
他說現在有一個西人牧師.
也很關心我們香港人.
他說遲下.
看看會不會叫我們.
搞一些茶經班.
可以認罪他們.
牧師知道你們會兩個月.
都是來辦命案的.
我說沒有問題.
如果我在命案教會.
不來幫你.
牧師真的很好.
這麼爽.
我想清楚就答應.
因為是示範.
當然不會拒絕.
因為示範的機會越來越少.
所以一個誠實身份的人.
就是你去選擇.
正確的事.
可能兩者都很正確.
但哪樣更有價值.
所以其實是幫助我們.
小心用你的時間.
你要怎樣去處理.
然後就是誠實身份的人.
懂得怎樣跟別人有交往.

$^{1361}$我再讀歌羅西書.
大家可以看第四章第六節.
我們一起讀出.
一個誠實身份表現的人.
他的言語.
他的對答.
他的回應是很有智慧的.
不會隨便的來勒著別人.
聖經裡面也是這樣說.
有時候一句說話.
說得好就如金瓶子.
金瓶果落在文王子裡.
一言可以輕放.
一句說話就觸碰對方.
特別是在香港.
很厲害的.
一句說話就觸碰對方.
可能很多人都很鬱結.
很暴躁.
問題是懂得怎樣去跟別人有交往.
很簡單.
我曾經說過.
我女兒跟媽媽說.
你爸爸很甜絲.
我被他騙了.
想嫁給他.
如果我生氣.
就死蠢.
第二件事.
你女兒.
不要說謊.
說謊是沒有好處的.
是你自己受苦.
我受了苦.
你明白嗎.
你懂得怎樣回應.
我就贏了她.
所以問題是你有沒有智慧去回應.
這是很重要的.
所以記住.

$^{1401}$說話是要有.
從神而來的智慧.
我們一起讀《真言》第十二章第十八節.
我們一起讀.
他說.
如果你說話說得不好的時候.
只會傷害別人.
但如果你能夠說話說得好的時候.
就好像一些風物.
能夠滋潤人的心.
所以求神要幫助我們.
我們還有一個成熟心靈的人.
你不斷地跟神親近.
神自然會給你智慧.
懂得怎樣去回應.
最後一次經文我作為結束.
我們一起讀第一情書第三章第十五節.
很多未信的人可能會問你.
關於信仰的問題.
你有沒有隨時有這個準備.
你跟神有一個.
好的關係的時候.
神自然會教你怎樣去回應.
是用一個溫柔的心去回應.
各人.
所以我們有一個成熟心靈的人.
作為一個老闆.
他要公平對待員工.
不要嚇求員工的工資.
不要剝削他們.
作為一個成熟的工人.
我們是要盡責.
我們不只是做表面的功夫.
因為我們要向神祈交代.
我們有一個成熟心靈的基督徒.
是會銜切禱告.
會經常跟神建立一個好的關係.
敏銳於神的教導.
警醒感恩.
不單單為自己祈禱.

$^{1441}$更加會為牧者被派入去禱告.
為必人來代叩.
為教會打開全套的門.
來禱告.
為牧者宣告有能力來禱告.
一個成熟基督徒的人.
會懂得怎樣管理自己的生命.
管理自己的社交.
不會浪費寶貴的光陰.
懂得怎樣與愛人來交往.
更加懂得怎樣與別人來應對.
所以.
仍然是那一句說話.
在歌羅西書.
就是我們一開始的時候所讀的.
你們既然接受了主基督耶穌.
就當進他而行.
在他裡面.
生根建築.
於是你信心堅固.
正如你們所領的教訓.
求神幫助我們.
做一個成熟的基督徒.
相信當你有成熟的表現.
會影響到周遭的人.
一起禱告.
多謝你給孩子們機會了解和宣告.
願意在接受訪問的訊息.
都成為眾大影展會的幫助.
都成為我們的提醒.
請求主力捐助.
祝福我們洋教會.
能夠榮耀主力.
祈禱拉動主義特性的民.
\newpage



\section{路加福音 11:1-13}
\label{sec:031960}
\textbf{禱告蒙應允的秘訣?}~~~~鄺偉志牧師
\newline
\newline
連結: \href{http://www.church.com.hk/acms/content.asp?site=cdc\&op=show\&type=product\&code=031960\&layout=sermon}{\texttt{031960}} ~~~~ 語音日期: 2026-04-12
\newline
\newline
\hyperref[sec:031959]{\small{< < < PREV SERMON < < <}}
~
\hyperref[sec:index_chronic]{\small{[返順時目]}}
~
\hyperref[sec:031961]{\small{> > > NEXT SERMON > > >}}
\newline
\newline
路加福音 11:1-13
\newline
\begin{longtable}{cl}
\hline
\hline
章節 & 經文 (和合本修訂版)\\
\hline
11:1 & \begin{tabularx}{0.7\textwidth}{X} 耶穌在一個地方禱告。禱告完了，有個門徒對他說：「主啊，求你教導我們禱告，像約翰教導他的門徒一樣。」 \end{tabularx} \\ \\ \relax
11:2 & \begin{tabularx}{0.7\textwidth}{X} 耶穌對他們說：「你們禱告的時候，要說：『父啊，願人都尊你的名為聖；願你的國降臨； \end{tabularx} \\ \\ \relax
11:3 & \begin{tabularx}{0.7\textwidth}{X} 我們日用的飲食，天天賜給我們。 \end{tabularx} \\ \\ \relax
11:4 & \begin{tabularx}{0.7\textwidth}{X} 赦免我們的罪，因為我們也赦免凡虧欠我們的人。不叫我們陷入試探。』」 \end{tabularx} \\ \\ \relax
11:5 & \begin{tabularx}{0.7\textwidth}{X} 耶穌又對他們說：「你們中間誰有一個朋友半夜到他那裡去，對他說：『朋友！請借給我三個餅； \end{tabularx} \\ \\ \relax
11:6 & \begin{tabularx}{0.7\textwidth}{X} 因為我有一個朋友旅途中來到我這裡，我沒有東西招待他。』 \end{tabularx} \\ \\ \relax
11:7 & \begin{tabularx}{0.7\textwidth}{X} 那人在裡面回答：『不要打擾我，門已經關了，孩子們也同我在床上了，我不能起來給你。』 \end{tabularx} \\ \\ \relax
11:8 & \begin{tabularx}{0.7\textwidth}{X} 我告訴你們，雖不因他是朋友起來給他，也會因他不顧面子地直求，起來照他所需要的給他。 \end{tabularx} \\ \\ \relax
11:9 & \begin{tabularx}{0.7\textwidth}{X} 我又告訴你們，祈求，就給你們；尋找，就找到；叩門，就給你們開門。 \end{tabularx} \\ \\ \relax
11:10 & \begin{tabularx}{0.7\textwidth}{X} 因為凡祈求的，就得著；尋找的，就找到；叩門的，就給他開門。 \end{tabularx} \\ \\ \relax
11:11 & \begin{tabularx}{0.7\textwidth}{X} 你們中間作父親的，誰有兒子求魚，反拿蛇當魚給他呢？ \end{tabularx} \\ \\ \relax
11:12 & \begin{tabularx}{0.7\textwidth}{X} 求雞蛋，反給他蠍子呢？ \end{tabularx} \\ \\ \relax
11:13 & \begin{tabularx}{0.7\textwidth}{X} 你們雖然不好，尚且知道拿好東西給兒女，何況天父，他豈不更要把聖靈賜給求他的人嗎？」 \end{tabularx} \\ \\
[1ex]
\hline
\hline
\end{longtable}
$^{1}$各位兄弟姐妹早晨.
我們再一次祈禱.
天父多謝你讓我們在一起.
在主的面前.
我們一群的弟兄姊妹.
集體敬拜.
我們可以透過詩歌.
透過聽道.
透過禱告.
來去敬拜你.
主我們將我們的心.
交給你.
我們將我們今天.
我們一起去敬拜.
願你得著當得的榮耀.
今日是聽到.
聽到禱告的功課.
主求你給我們再一次.
重視禱告.
亦都去更新我們禱告的生活.
生命.
以致我們能夠與你更加親近.
更加明白.
禱告在我們的人生裡面.
是不可或缺.
求聖靈幫助我們去.
明白你的話語.
幫助孩子.
講的很有限.
但你聖靈去幫助孩子.
講得明白.
與及弟兄姊妹.
亦都聽得明白.
擺在心裡面.
我們一同緊守進行主你的道.
我們禱告.
奉耶穌基督的名字.
阿們.
我以前在教會很多年.
他們都經常去背主禱文.

$^{41}$我就是這樣.
懂得背主禱文.
我自己開始.
去信主.
教會的背景.
是一些福音堂.
是不會背主禱文的.
但背主禱文.
是非常之好的.
考考大家會不會背.
講任何東西.
大家不要緊張.
背不到的話.
動一動就可以了.
好嗎.
我們一起試試背.
一二三.
我們在天上的父.
願人都尊你的名為聖.
願你的國降臨.
願你的旨意行在地上.
如同行在天上.
我們日用的飲食.
今日賜給我們.
免我們的債.
如同我們免了人的債.
不叫我們遇見試探.
救我們脫離凶惡.
因為國度權柄榮耀.
全是你的.
直到永遠.
阿們.
都不錯.
懂得背.
但有沒有去運用呢.
今日大家一起探討一下.
這一張照片.
是其中我在網上拿出來的.
因為我去聖地.
沒有去到這個地方.

$^{81}$其實聖地有兩個地方.
我想和大家分享.
和主禱文有關.
這個很明顯.
是一個中文版本的主禱文.
其實就在這個地方.
這個地方.
你見到這個拉丁文的字.
就是阿們的父.
其實就是.
這一間教堂.
曾經開了又拆了的一個地方.
可以說是主禱文堂.
在鑒欖山.
鑒欖山在哪裡呢.
如果你有印象.
現在那個地方.
有一個叫神殿.
不是清真寺.
一個金頂的神殿.
那個位置.
就是以前聖殿的位置.
對面就是鑒欖山.
鑒欖山其實有幾間教堂.
其中一間教堂.
就是這一間.
叫做主禱文堂.
裡面是怎樣的呢.
裡面有很多主禱文.
有140個不同國家翻譯的主禱文.
是景點.
除了說是一個景點.
為何在那裡有一間教堂呢.
值得和大家提一下.
其實它在地牢那裡.
你見到有新的裝修.
也有舊一點的情況.
相傳這個地方.
就是耶穌經常在鑒欖山對面.
對著聖殿.

$^{121}$在那個地方來教導.
因為陰涼一點.
在山洞裡.
相傳今天的經文.
門徒問.
究竟我們怎樣來禱告.
可不可以教我們禱告.
就在這個山洞來教他.
不知道是不是.
但是一直傳下來.
也值得大家知道.
這是第一個地方.
第二個地方.
如果你去過聖地猶撒冷.
就最有名的.
就算你沒有去過.
你也看過這個叫什麼.
哭牆.
哭牆我們去看什麼呢.
其實進去看看.
怎樣去祈禱.
在哭牆不詳細說很多歷史背景.
一大堆很多東西要解釋.
但是這邊是男的.
這邊是女的.
男的進去還有一個很大的地方.
很多傳統的猶太人.
就會在那裡禱告.
不用開聲.
如果有聲音.
這裡有一個我拍的片.
原來沒有動.
不要緊.
其實你要見到那些人在牆那裡.
當然那些紙很有趣.
寫了那些土文.
塞進去.
我們也是寫了塞進去.
你說有沒有人清理.
是沒有人清理的.

$^{161}$怎樣清理.
天然清理.
下雨.
一直這樣.
其實很骯髒的.
不知道為什麼會這樣做.
但是這個哭牆前面.
你看到猶太人.
就會這樣祈禱.
站著在那裡祈禱.
特別沒有片.
不過你看到.
這個手刮氣墊.
它其實就是一條帶.
這裡不看了.
你看看下面.
這個就叫Tiffany.
就是一個叫經文盒.
經文盒是怎樣呢.
就是猶太人.
十三歲成年禮的時候.
他就會開始戴這個盒.
就會祈禱.
因為他們是這一類傳統的猶太人.
他們是每天都會祈禱的.
一會兒我也會講一點.
但是他們就會有這個經文盒.
兩個的.
一個就是戴在頭上的.
你祈禱的時候.
就好像額頭鑿了一些東西.
凸了出來.
經文盒.
顧名思義就是盒裡面.
寫了一些經文在裡面.
戴在頭上.
你想想.
為何戴在頭上.
是在這裡.
第二個.

$^{201}$就是你如果是右手.
慣用手.
你就左手綁上去.
這個另外一個經文盒.
你見到為何要一張相.
因為它是對著哪個位置.
對著心.
所以為何有兩個.
為何有這個方位.
因為他在祈禱的時候.
一方面就是.
這兩個經文盒.
這些經文.
一會兒告訴你大概是什麼經文.
就是對著你的思想.
你的頭.
你祈禱.
你的人.
是要在思想上.
來敬畏神的.
一切都是向著神的.
所以在頭上提醒自己.
另外就是.
這個心.
你對著你的感情.
你的意志.
所以很有意思.
還要綁在手上提醒.
你的祈禱.
你是有行動的.
為何經文盒呢?.
所以要說裡面.
究竟猶太人會怎樣祈禱.
傳統猶太人.
有三種祈禱.
有不同的圖文.
最出名的一個就是Shema.
Shema的意思就是聽.
如果你記得《申命記》.
我以前也好像在講道講過.

$^{241}$《申命記》第六章.
第四至九節.
Shema.
其中一件事就是.
以色列電台.
每一天都是Shema.
問候的時候都是Shema.
以色列你要聽.
聽不只是耳朵聽.
聽就是聽從的意思.
你有行動的.
以色列你要聽.
耶和華主是獨一的神.
接著是怎樣呢?.
這個很熟悉.
盡心盡性盡意愛主.
你的神.
這個就是Shema.
每天祈禱都會.
背著Shema.
還有幾段經文.
在《申命記》和《出阿俠記》的地方.
這個就是他們祈禱的內容.
其中經文盒裡面.
一定有Shema.
另外有兩個.
很重要的.
圖文.
在猶太人傳統裡面.
現在他們也在用.
跟我們主圖文有關.
現在不是講猶太人.
主圖文其實是.
那個圖文是從猶太人的祈禱而來.
所以我一定要這樣讓大家知道.
你自己Google.
你自己AI慢慢搜尋.
慢慢看.
讓大家知道怎樣做.
說到AI.

$^{281}$如果你要去研究聖經.
AI你要檢查.
你要懂得問問題.
麻煩就上余大豐傳道.
那個《齊用AI學聖經》.
順便幫他宣傳一下.
的確不要全部都相信.
但是你要檢查.
第二種的圖文是什麼呢?.
它叫做Kaddish.
Kaddish是什麼意思呢?.
聖潔.
也就是聖潔師.
其實在主圖文.
稍後我們再看.
你讀的時候.
我們很順.
因為我們只是讀主圖文.
我們在天上的父.
願人都尊你的名為聖.
願你的國國臨.
願你的旨意行在地上.
如同行在天上.
三願.
這個三願其實是從哪裡來的呢?.
Kaddish其實很類似.
你直接一看到.
原來真是猶太人.
他根本就是這麼祈禱.
他大概是怎麼樣呢?.
他第一句.
願你的名被高舉.
被尊為聖.
是不是很相似?.
這個就是Kaddish.
這個聖潔師.
後來他也有很多行.
另外就說.
願你的國快快臨到.
所以呢.

$^{321}$從這裡你看到.
耶穌說的主圖文.
其實是一個.
再精簡的.
或者是再優化.
改良.
教導門徒去禱告的.
一個猶太人祈禱的版本.
另外一個圖文.
就叫做阿彌陀.
阿彌陀的意思.
這些音你不用理他.
其實他的意思就是站立.
站立.
所以你看到那些猶太人.
為什麼站在那裡祈禱呢?.
就是這樣.
剛才忘記說.
那個Shamma說到後面.
那個申命記錄章.
4-9節你看到.
他說.
你要將那些聖經.
帶在額上.
還有帶在手上.
所以為什麼.
Diphne就是做一個經文盒.
就是要這樣做.
他簡直就是.
他怎麼寫.
就怎麼做.
留下傳統.
去提醒自己.
怎樣去遵行神的話語.
怎樣去祈禱.
所以就有這樣的一個做法.
所以你要看到.
站起來禱告.
阿彌陀.
是在祈禱什麼呢?.

$^{361}$我們在看主圖文.
不過你會記得.
譬如說.
日用的飲食.
赦免人的罪.
罪債.
或者是.
不叫我們遇事.
試探.
這些字眼.
這些說法.
其實都是站起來的禱告.
站起來的禱告.
他有18項.
後來就加到19項.
你對一對照.
所以其實.
就是說.
簡單來說.
這個你慢慢再去查考.
大家去了解多一點.
在猶太人.
他一直都是祈禱的.
所以你會看回.
這個經文.
這個很特別的.
為什麼.
一來我要在講路加福音.
但是我們剛才背的是什麼版本呢?.
其實是馬太福音版本.
有一點點不同.
兩個重點有不同.
但是這裡.
路加福音講到一件事.
為什麼有主圖文出來.
還有主圖文.
很多人.
一直在聽.
主圖文.
主耶穌的圖文就不是.

$^{401}$你看清楚一點.
主耶穌教導我們怎麼祈禱.
所以有些人就說.
不是Lord's Prayer.
是Disciple's Prayer.
是門徒的圖文.
就是說我們應該怎麼祈禱.
所以我們一定要學習.
耶穌怎麼教我們祈禱.
因為.
什麼時候教人是最好的.
什麼時候教小朋友是最好的.
什麼時候教學生是最好的.
就是他問問題.
門徒就是有一個問題.
他看到施洗約翰.
尤其是.
其實這裡有背景的.
施洗約翰帶著一群門徒.
他也有教人祈禱.
所以.
耶穌就在那個情況下.
就被門徒問.
他看到人家都.
施洗約翰就是耶穌的表哥.
施洗約翰是耶穌的先鋒.
他這樣教我們祈禱.
就是教門徒祈禱.
不如耶穌你教我們祈禱.
最好問問題的時候.
耶穌回答.
所以很重要.
大家順便說一下.
大家準備要教導的話.
不要不斷地說.
OK.
你要準備一下.
當有人問的時候.
你不要不回答他.
小朋友如是.

$^{441}$大人如是.
因為.
這件事是最.
他這個人是最吸收.
其實耶穌已經準備了.
應該怎麼祈禱.
所以他就.
告訴門徒.
你們應該是這樣祈禱的.
其實.
我們會唸.
其實主禱文是用來唸口往的.
是用來背的.
你背完.
你知不知道自己背了什麼.
有時候我們不知道背了什麼.
不懂背就糟了.
什麼都沒有.
裡面.
說個秘密給你知道.
我背主禱文.
就.
什麼時候背.
我每天都背幾次.
怎麼背幾次呢.
我家裡.
如果你來.
歡迎你來.
不過不要一起來.
在露台那裡.
有個架子.
那個叫單槓.
在地上很穩固.
用來做什麼呢.
我不是練習.
而是.
拉著.
透過身體的力量.
這是醫生教的.
拉鬆腰.

$^{481}$腰不是很好.
拉鬆.
每一次就吊.
吊在那裡.
很悶的.
背經.
我有幾樣經文一定背的.
第一個就是主禱文.
第二個詩篇23篇.
最近加了詩篇第8篇.
背完這3篇.
分幾宗.
已經很辛苦了.
一路去背.
這個其實是.
學習在.
有一本書叫做原子習慣.
詳情不說了.
他教其中一種建立習慣.
就是要.
將一些事情.
拉在一起.
以致.
你能夠可以.
每次一做到那件事.
你就拉在一起.
譬如你一刷牙洗臉.
你就記得.
去畫眼線.
你就可以將兩件事拉在一起.
就會成為一個習慣.
希望.
這個Point我覺得很重要.
大家去建立一支習慣.
還有我們已經.
有時候不要說讀經.
不要說背經.
讀經也不是很讀.
背經其實是.
很重要.

$^{521}$但是.
這個經文.
這些經文是不是用來背的呢.
念口往背了.
你學了禱告了.
當然不是.
是要用你的心思.
你要明白.
在做什麼.
你背的時候.
是給你有些東西在裡面.
但是你是怎樣運用呢.
今天就是想說.
一會說到最後.
有些實用的方法.
幫我們一起去用主禱文.
因為很重要.
我們.
看到這個.
在禱告的時候.
應該怎樣來祈禱呢.
就有這個主禱文出現.
我們輕輕讀一讀它.
你會看到.
剛才三願.
願.
符.
或者說一個比較.
其實括號.
就是那些經文.
有一個版本.
這個是和合本修正版.
其實括號.
經文都在這裡.
只不過.
古卷.
就會有這個括號.
你發覺.
馬太福音就有這些內容.
有這些內容.

$^{561}$還加上最後一句.
國度權並傳是你的.
這一句.
這裡就沒有了.
所以這個版本.
沒有括號.
可能是耶穌.
馬上在門徒問的時候.
第一個版本.
就沒有了括號.
不過有括號.
都沒有問題.
你說我們背開登山寶訓.
那個就可能是一個.
後來再完善一點的版本.
耶穌的教導.
不只是教導一次.
可能教導很多次.
所以我們背的那個是OK的.
就是多一點.
所以我想給大家知道.
還有登山寶訓的情況.
將主禱文放在哪裡.
就是登山寶訓.
耶穌就是教人.
來怎樣.
我們所謂天國門徒的約章.
怎樣做一個基督徒.
就是登山寶訓.
其中一件事.
就是要禱告.
那個禱告就拿著禁食.
拿著施捨的.
這個是猶太人.
最重要的那幾件事.
那幾件善行.
所以在登山寶訓.
這三件事裡面.
他用主禱文來講.
一來就是特別去講.

$^{601}$你們應該怎樣禱告.
這個也是.
一會兒也會解釋.
但是他的重點.
前文後理.
就是那個主禱文說.
你們不要像外邦人那樣.
在十字路口那裡祈禱.
那時候有些人是這樣的.
以及扯到假冒為善.
以及不要用重複的話來禱告.
所以就不是背.
不斷背不斷重複.
你又叫他祈禱.
一定不是.
而且這個是那些.
當時那些外邦人.
拜其他神明的人.
有些咒語.
其實我們中國人也會這樣.
很多宗教.
我不是詆毀你們的宗教.
但是有些是.
唸而已.
他不是很知道要唸什麼.
他覺得唸多一點.
好像一個行為.
就可以取悅到神.
或者得到一些東西.
或者個人就會提升.
其實耶穌已經是反對這件事.
不是重複這樣唸.
你就有用的.
這些禱告你不明白究竟發生什麼事.
你這樣唸是沒有意思的.
以及馬太方特別說.
你不要假冒為善.
祈禱不是做show.
有些人不知道為什麼.
我認識一個神學院的同學.

$^{641}$很有趣的.
說話平時很悶悶的.
祈禱的時候說話瘦乎乎的.
很厲害的.
不知道為什麼這麼有文采.
我覺得神給他.
祈禱裡面.
有些特別能力.
我不知道為什麼.
但是他當然不是這樣.
這位同學.
但是有些情況就是.
當時那些法利賽人.
祈禱好像很敬虔.
讓人看到他祈禱多麼敬虔.
這個好了.
這些耶穌非常反對.
這叫做假冒為善.
所以那個主禱文.
第一.
如果補充馬太方.
就不是用來不斷背.
這樣就可以的.
要明白地跟著禱文去禱告.
第二就是不要假冒為善.
祈禱的時候都用來裝模作樣.
使人覺得你敬虔一點.
這個絕對是不行的.
因為耶穌已經說了.
這裡的重點就有不同.
在解釋每一節.
你會發覺前文後理.
第一.
這裡說了原因.
為什麼耶穌要教.
接著你會發覺.
後面那兩段的故事.
是完全跟主禱文有關的.
原來耶穌說了主禱文.
接著後面我們一定要解釋.

$^{681}$這兩個的比喻.
以致去明白耶穌為什麼要說主禱文.
還有它的方向是怎樣.
是另外一個方向.
講完這兩個比喻的意思.
你就會知道方向是什麼.
第五節.
耶穌先說第一個比喻.
你留意這個比喻裡面.
有沒有一些字眼.
是跟主禱文有關的.
考考大家.
前文後理的問題.
故事是怎樣呢.
有一個人.
他接待了一個朋友.
當時你要想像.
是一個比較鄉村的地方.
不是我們這些大城市.
當然你也可以代入.
你們住house住apartment.
也有隔離鄰舍.
關係是怎樣的.
你可以嘗試代入.
當時那些村.
鄉村.
一條村什麼事都發生.
都知道的.
我現在也是住house.
很靜的.
有什麼事發生.
有隻狗走過.
也會吵的.
你發覺當時也是這樣.
有個人接待了一個朋友.
他可能是預備的問題.
不夠東西吃.
他就要問隔離屋.
去借三個餅.
但是隔離屋.

$^{721}$其實他就說.
走去隔離屋.
他知道有餅.
因為當時是鄉村.
隔離屋有什麼.
投還申請.
他也知道有餅.
敲門.
但是那個人就說.
關了門.
因為酬態傳統的背景.
平時就.
日光日白打開了屋子.
以前住新界.
隔離所有都是打開了門.
就算我住屋村以前.
有時候也不鎖閘.
一個比較這樣的情況.
好了.
但是關了門.
關了門就不要.
禮節不應該去騷擾人.
但是他就.
要去騷擾人.
那個人也說關門.
我兒子睡了.
我不行.
但是那個人怎樣.
耶穌說什麼.
第八節.
雖然.
不是因為他的朋友.
就會給他.
但是也會因為那個人.
不顧面子.
修訂版.
和合本身的版本.
情詞迫切.
不顧面子.
他一定會給他.

$^{761}$這裡有兩件事.
這個祈禱.
應該是怎樣的.
第一.
他和主禱文有什麼關係.
三個餅.
我們日用的飲食.
那個其實就是餅.
其實他已經解釋了.
他正在解釋.
這裡有兩件事.
第一件事.
其實第九節.
開始.
在第二個比喻中間有幾句.
我們很熟悉的.
就是你祈求就給你門.
尋找就尋見.
叩門就給你開門.
其實這裡叩門就給你開門.
是說前面那個例子.
為什麼敲門.
開窗是不行的.
因為前面那個例子是跟著.
但是在說祈禱.
就是說那個人他都知道.
但是你堅持.
三個動詞都是有種.
持續的意思.
你堅持的時候一定他會開門給你.
這是第一點.
第二點.
因為是他朋友.
其實他朋友.
他怎樣都會給他.
這是說一國的關係.
你看到還有一個.
情況.
就是下面那個.
例子.

$^{801}$還有一個字眼.
我們一會再說.
第二個例子.
十一節.
你們中間作父親.
有誰兒子求魚.
給他蛇就當魚.
求雞蛋就給他茄子.
其實耶穌是用一個很誇張的手法.
來說這件事.
你們雖然不好.
這裡比較難看一點.
跟主禱文的關係.
你們雖然不好其實可以譯為邪惡.
你們是邪惡的.
尚且都知道.
有好東西給兒女.
何況天父更加不將聖靈.
給求他的人嗎.
他在說這件事.
這個例子很誇張.
兒子和父親.
不舉這麼多例子.
比如說.
車.
是不是這樣.
正常父親不會.
所以在說什麼.
那個關係.
所以無論朋友.
那個朋友不願意.
但因為他是你的朋友.
你試過向人借錢嗎.
你會向什麼人借錢.
如果你真的要借錢.
去到一個地步.
其實正常來說.
我們會分辨.
大家不要亂借錢.
但我們會分辨.

$^{841}$通常一般不會向人借錢的人.
都向你借錢.
即是說他有需要.
可以這樣看.
你一個陌生人向你借錢.
你真的要小心一點.
為什麼那些人會向人借錢.
一來是有需要.
二來是問朋友.
這是一個關係.
第二個又是一個關係.
父親和兒子的關係.
連結到.
所以去解釋前面.
那個主禱文.
我們在天上的父.
我們和主.
因為有那個關係.
我們會祈禱.
或者倒過來.
就算我們.
有些可能是初信主耶穌.
或者我們信了耶穌很多年.
和神的關係.
若即若離.
這樣就不太行.
但如果你和神的關係.
有什麼可以試到呢?就是你的祈禱.
你如果和神的關係好.
你就會和祂祈禱.
亦都因為神和你關係好.
祂是天父.
祂看到兒子有需要.
祂怎會不去照顧你呢?.
祂怎會不去.
不聽你的祈禱呢?.
加上剛才那兩個例子.
你情詞迫切.
你不要面.
你那種.

$^{881}$不斷的祈求敲門.
那種堅持.
一定會給你.
解釋主禱文的內容.
待會很快講一個應用的版本.
你們可以去做.
但是.
耶穌在前文後理的解釋.
就是你這樣去祈禱.
你那種態度.
馬太福音就是說.
你不要虛偽.
你祈禱是真誠.
但是.
這裡是說.
這個祈禱你又要注重.
那個關係.
你和神的關係是這樣.
神是怎樣去.
看著你.
只要我們堅持這樣去祈禱.
我們是一定.
知道神是聽我們祈禱的.
現在實在太多的.
問題,我們人生有很多問題.
世界有很多問題.
如果基督徒.
自己想自己的祈禱生活.
如果基督徒.
是少祈禱的.
不懂祈禱的,或者很多時候.
很可惜.
我們基督徒是生活.
做事奉.
甚至做很多事.
但是祈禱是永遠.
不行的時候才來祈禱.
人的盡頭就是神開始.
但經常都是盡頭才開始.
不覺得盡頭就不開始.

$^{921}$祈禱每一天.
都沒有一個.
關係的建立.
有時候我們說夫妻的比喻.
或者做一個好朋友的比喻.
你不會到問他借錢的時候.
才是朋友,你明不明.
請你吃飯,其實問你借幾萬.
不是這樣的.
是一直有一個關係.
你才會很自然的.
我不想.
我不願意,但我都要.
問他借,因為我有需要.
而那個人也肯借給你.
天父怎麼會.
不聽我們的祈禱.
問題是我們跟神沒有關係.
我們就不會堅持.
有時候很奇怪,祈一次禱.
不靈的.
就更加去做.
其他事,不祈禱.
你說是不是有問題.
其實是不斷的,持續的.
好像耶穌教,你要.
來到這樣的態度.
去祈禱,所以.
建立一個.
一天去祈禱,用主禱問.
或者是你的心態.
就是有什麼.
你都可以祈禱,最難的,最易的.
其實是一個關係,祈禱很簡單.
其實祈禱是什麼意思,就是跟神.
聊天,但我們不聊天.
經常都不聊天,突然問他.
很大的事幫我解決.
你是不是覺得,你自己也覺得很怪.
所以那些人說,為什麼我祈禱.

$^{961}$不行,就是因為平時.
你根本都沒有建立一個跟神的.
關係去祈禱,這個真是很大.
問題,如果我們祈禱.
生活出了問題,我們.
整個基督教生活其實出了很大.
很大問題,很容易很虛假.
很容易.
沒有真是.
倚靠神,我們經常說,怎麼倚靠神.
就是祈禱.
所以這個是很重要.
最後,講清楚究竟.
這一段怎樣應用.
裡面每一句都是.
很精要.
三願.
我們用.
傳統,全面的版本.
我們在天上的父.
關係,你很清楚.
我們天上的父,這裡不解釋.
有些東西可以很詳細解釋,但先不解釋.
但願人都尊你的名為聖.
是什麼,如果用一個詞語.
兩個詞語就是.
尊你的名為聖,就是說.
我們要尊敬神,這一句.
是尊敬神.
每次你背主統文.
然後你想一想,每一句再.
願人都尊你的名為聖.
神啊,我要尊敬你.
我今天有沒有尊敬你,我今天求神.
給我,在我生活行事為人.
要尊敬你,另外一件事.
願人都尊你的名為聖.
就是有一個傳福音.
見證的導向.
作見證,所以這一句.

$^{1001}$當你背完之後.
你再說,願人都尊.
你的名為聖,神啊.
願更多的人來相信你.
又會反省自己.
我們會不會用我們的見證.
或者有機會可以來傳福音.
有時候有些人說,我沒有恩賜.
傳福音,這個不是恩不恩賜的問題.
有沒有恩賜都要作見證,都要傳福音.
願人都尊你的名為聖.
你這樣祈禱,接著到第二句.
願你的國降臨.
這是盼望,現在末世.
這句更加重要,神啊,你的國降臨.
我要有這個盼望.
好了,接著,願你的旨意.
行在地上如同行在天上.
神的旨意是用來怎樣的.
我們明白了之後,要信服.
神啊.
我信不信服你呢.
你有什麼旨意在我身上呢.
可能你一直都知道.
你今天再要提醒自己.
神啊,我要遵行你的旨意.
好了,我們日用的飲食.
今日賜給我們,或者天天賜給我們.
這個就是一個支足.
一個倚靠.
大家.
手不手緊.
大家發不發覺現在油價升得很厲害.
接著呢.
你說我不開車的,但是呢.
因為物流啊,各樣東西.
現在是不是很大的經濟的難題.
是不是啊.
真的個個都是.
你會發覺這句很重要.

$^{1041}$每一天就提醒.
我不是說你這樣就解決了問題.
但是你跟天父說的.
天父會給最好的你,剛才說完了.
我們日用的飲食,天天賜給我們.
主啊,就是一個感恩.
神啊,你仍然給我能夠有天天的飲食.
並且我的需要.
肉身的需要.
家人的需要,經濟的需要.
我的工作很危險.
那些的需要.
我都有倚靠你的心.
好了,接著呢.
祈禱到這裡了,赦免我們的罪.
因為我們也赦免.
犯規欠我們的人.
有兩個角度,第一就是認罪.
這個是提醒我們向神認罪.
我們真的是罪人,每天都要認罪.
另外一件事呢.
因為我們也犯.
因為我們也赦免.
犯規欠我們的人.
一句到尾就是說.
你是在說.
這個提醒我們,我們要愛人.
而愛人呢,最難的就是.
別人得罪了你.
你要赦免他,或者包容他.
所以這個是.
人際關係,這一句.
提醒我們,你一輩子的時候.
赦免我們的罪.
向神認罪.
接著你再想一想這一句.
因為我們也赦免犯規欠我們的人.
或者免我們的債.
如同我們免我們的人的債.
人和你的關係.

$^{1081}$這個提醒我們.
這一天我們要這樣面對.
不叫我們陷入試探.
華盤修訂版.
或者不叫我們.
遇見試探.
或者有另外一個說法.
就是在試探裡面.
你不會犯罪.
這幾個都是一個提醒.
有很多試探.
這個是求主.
保守我,不要遇到試探.
或者在試探裡面.
其實試探這個字.
和試煉是一樣的.
在試煉裡面,在試探裡面.
求你讓我能夠.
去面對到.
我現在面對的難題.
以及在.
我已經在面對這件事.
但是不要陷入這個試探.
落到一個田地去犯罪.
所以呢.
這一句就提醒我們.
怎樣去生活.
我們自己的心去求神保守.
最後救我們脫離那惡者.
或者救我們脫離兇惡.
我們.
不知道會面對什麼.
我們經常說.
我們都不知道.
今天不知道明天的事.
隨時都會有意外.
又有很多惡者的攻擊.
又有很多很凶險的事情.
這個世界.
大求神保守我們.

$^{1121}$這個是一個precaution.
就是說我們要.
有個預防.
最後那句.
其實都可以放進去.
因為國度榮耀全是你的.
直到永遠也滿.
將所有的東西.
我們的生命.
我們這個人.
都是屬於神的.
我們用這個結束都可以的.
如果我們這樣去祈禱.
你會發生什麼事.
你天天這樣去記住.
這樣去背.
然後這樣去祈禱.
會發生什麼事.
你真的在祈禱.
你會發生什麼事我不知道.
但是神一定會聽你祈禱.
並且在你身上工作.
但是更加這個祈禱.
如果你體會登山寶訓.
門徒是要有這個祈禱.
其實是騷擾我們門徒的生命.
你試想一下一個人.
如果天天這樣去祈禱.
這樣被提醒和求神.
支取力量.
他的品格會改變.
所以這句話就是祈禱不但能夠.
改變那些祈禱的事.
那個題目是什麼.
祈禱是不是有秘訣.
是不是一定誰聽.
但是有時候可能表面.
祈禱可以改變一些事.
但是就算祈禱.
你好像看不到改變的事.

$^{1161}$但是那句是這樣的.
祈禱不但能夠改變你要祈禱的事.
更加是改變.
祈禱那個人.
我們基督徒的生命.
我們基督徒的品格如何塑造.
祈禱.
我們一起祈禱.
讓我們一起思考主禱文.
求主幫助我們.
在主禱文裡面.
我們就是要學會祈禱.
求你幫助我們.
透過主禱文的提醒學習.
我們又這樣去進行.
就像猶太人一樣.
其實他們有一個很好讓我們學習的地方.
就是每一天不同的時段.
每一天都是這樣穩定的.
向你祈禱.
讓我們願意這樣去依靠你.
建立這個關係.
也都信靠你.
以致我們在祈禱裡面.
是經歷主與我們同行.
改變我們的人生.
改變我們的品格.
更加改變一些我們改變不了的事.
我們面對的難處.
甚至透過祈禱去幫助其他人.
領其他人歸主.
我們就是這樣的.
求你讓我們進行你的禱.
好好地去祈禱.
我們這樣的禱告交託奉耶穌基督的名字.
阿們.
\newpage



\section{創世記 22:15-19}
\label{sec:031961}
\textbf{海沙}~~~~陳崇基牧師
\newline
\newline
連結: \href{http://www.church.com.hk/acms/content.asp?site=cdc\&op=show\&type=product\&code=031961\&layout=sermon}{\texttt{031961}} ~~~~ 語音日期: 2026-04-12
\newline
\newline
\hyperref[sec:031960]{\small{< < < PREV SERMON < < <}}
~
\hyperref[sec:index_chronic]{\small{[返順時目]}}
~
\hyperref[sec:031962]{\small{> > > NEXT SERMON > > >}}
\newline
\newline
創世記 22:15-19
\newline
\begin{longtable}{cl}
\hline
\hline
章節 & 經文 (和合本修訂版)\\
\hline
22:15 & \begin{tabularx}{0.7\textwidth}{X} 耶和華的使者第二次從天上呼喚亞伯拉罕， \end{tabularx} \\ \\ \relax
22:16 & \begin{tabularx}{0.7\textwidth}{X} 說：「耶和華說：『你既行了這事，沒有留下你的兒子，就是你的獨子，我指著自己起誓： \end{tabularx} \\ \\ \relax
22:17 & \begin{tabularx}{0.7\textwidth}{X} 我必多多賜福給你，我必使你的後裔大大增多，如同天上的星、海邊的沙。你的後裔必得仇敵的城門， \end{tabularx} \\ \\ \relax
22:18 & \begin{tabularx}{0.7\textwidth}{X} 並且地上的萬國都必因你的後裔得福，因為你聽從了我的話。』」 \end{tabularx} \\ \\ \relax
22:19 & \begin{tabularx}{0.7\textwidth}{X} 於是亞伯拉罕回到他僕人那裡。他們一同起身，往別是巴去，亞伯拉罕就住在別是巴。 \end{tabularx} \\ \\
[1ex]
\hline
\hline
\end{longtable}
$^{1}$我們先祈禱.
主愛多謝你.
因為我們都是不堪一擊的人.
我們都很脆弱.
很容易倒下.
主愛你的生命.
是因為你的恩典去到成災.
我們經歷到你的恩典.
你的恩典和你的力量.
你的信實.
是恆久不變.
主幫助我們仰望和信靠你.
特別當我們踏入新的一年.
我們在生命裡面去到更高之處.
更高之處不是因為我們更加厲害.
而是更加在你面前.
我們匍匐承認.
我們是軟弱是需要你的.
你是那位最偉大的上帝.
我們去尊崇你.
去讚美你.
我們這樣祈禱逢尊名金.
甚麼叫元氣大上.
我真的有少少元氣大上.
氣不是很夠.
不過都要說.
剛剛村妹還說跳舞唱歌.
今年新的一年.
我們有一個新的.
連題.
就叫NEXT.
下一步是甚麼呢.
我們想挑戰大家.
我們整個教會每一個人的個人生命.
我們怎樣向前一點.
向前一點不是說我們再強一點.
而是我們更加見到上帝.
是一位值得信靠的上帝.
所以在這次.
我們整年裡面.

$^{41}$我們都會說NEXT這個topic.
我自己會.
其他staff不會.
但我自己會.
我選了兩段經文.
創世記十二章和二十二章.
今次開始來說這個系列.
甚麼叫NEXT.
NEXT的意思.
就是我們需要去了解.
到底這位上帝.
是一位不同的上帝.
我們所讚美的上帝.
我們所認識的這位上帝.
到底是怎樣的上帝.
我覺得最好的方法開始.
就是回到第一本書.
就是創世記.
這個簡圖就給你.
創世記最重要的故事.
就是阿巴拉罕的故事.
所謂的結構.
所以我選了十二章一段經文.
和二十二章一段經文.
說到阿巴拉罕的時候.
第一次接觸聖經也好.
阿巴拉罕是我們每個人.
現代人都需要認識的.
因為你知道.
在世界裡面.
三大個獨神的宗教.
三大個獨神的信仰.
只是承認一位神.
ONE GOD ONLY.
三個大的信仰.
猶太教,基督教和伊斯蘭.
都是以.
阿巴拉罕.
作為他們一個.
最尊崇的.

$^{81}$其中一個人.
三個.
大的宗教.
都是以阿巴拉罕作為基礎.
包括基督教.
所以算起來.
所有這些.
去到認識阿巴拉罕.
以阿巴拉罕為他們的.
所謂的父的話.
這個人數可以有四十億.
包括三個.
一神宗教所有的信徒.
所以如果你今天.
在這個世界裡面生活.
如果你不認識阿巴拉罕.
沒有聽過阿巴拉罕的話.
有些追不上.
所以我們要知道阿巴拉罕是什麼.
這裡我們第一點.
我想看到阿巴拉罕的.
故事的時候.
他的故事告訴我們什麼呢.
就是認識這一位.
宣教宗的上帝.
GOD ON A MISSION.
原來這個上帝藉著阿巴拉罕的故事.
去呈現出.
這位上帝是一個不斷.
去到做宣教工作的.
我們很多時候想到.
宣教工作的時候.
就想起我們有沒有傳道部.
我們有沒有去傳福音.
或者某一個項目等等.
聖經從頭到尾都介紹.
這位上帝是一個怎樣的上帝呢.
就是一個GOD OF MISSION.
是一個.
向普世去到宣教.

$^{121}$在做工作的一位神.
你如果看創世記.
由開始的話.
我們很簡單的說.
由開始創造這個世界的時候.
上帝的心意.
就是這個世界是一個.
美好的世界.
還有.
是一個祂願意賜福給.
全人類的一個狀態.
這是創世記第一章.
到第三章發生什麼事呢.
最闖入世界.
然後有邪惡,有暴力,有很多驕傲的事.
三至十一章就告訴我們.
如果有最闖入這個世界.
搞亂了這個世界的時候.
這個世界會是一個怎樣的地方呢.
到巴別塔第十一章.
就是列國分散.
然後人類試圖離開.
或者控制上帝.
自己建立自己的命.
所以一至十一章.
告訴我們.
最搞亂這個世界怎麼辦.
上帝是不是.
好像我們平時那樣電腦不運作.
我們就按reset.
或者有些東西不運作.
我們就刪除它.
重新再來.
不是在介紹上帝.
是見到上帝.
見到他的世界.
被罪人所搞亂之後.
他是沒有放棄過的.
他是進入這個世界.
並且去到.

$^{161}$展開他自己.
要萬國得福的救恩的計劃.
這個就是.
神的回應一步.
創世記第十二章.
上帝在這個.
搞亂了的世界.
你今天看看我們的世界是不是搞亂了.
搞亂到不可以再搞亂了.
但是上帝沒有放棄過.
沒有give up過.
他呼召阿巴漢出來.
他去展開一個新的計劃.
所以這個就是.
介紹這位的上帝.
描述這位的上帝.
就是一個不斷的.
在世界裡面去展開他的工作.
並且去拯救萬國.
十二章這裡發生什麼事呢.
由一至十一章我們知道.
只有死亡.
樣樣東西都game over.
十二章上帝.
就突然間講說話.
好像他開始的時候.
從無到有.
他講說話.
要有什麼就有什麼.
上帝又再一次講說話.
第十二章這裡.
他講說你要離開你本地本族本家.
往我所指示你的地方.
我使你能夠成為大國.
也都賜福給你.
使你的名為大.
要使你成為別人的祝福.
為你所祝福的我就賜福他.
救助你的我必救助他.
就是一個保護你的.

$^{201}$上帝保守他.
地上的萬族都要因你.
阿巴漢和你的子孫.
你的國度.
你所有的後代來得福.
這就是上帝在黑暗裡面.
去到講說話.
的第二次.
第一次是創造的時候.
第二次就是十二章這裡.
上帝.
God speaks into chaos.
在我們生命裡面.
最亂的時候.
世界裡面最亂的時候.
上帝原來說我不禁reset.
我禁new plan.
新的工作.
要展開.
你都可以展開.
上個禮拜一個最大的新聞.
當然就是我們美國.
射了一架飛船上去.
然後繞過月球.
一周就回來了.
安全回來了.
繞過一周.
發現沒有嫦娥.
見了月亮.
從來沒有見過.
人們沒有見過.
的一些相片.
就好像這一張.
太空人說了什麼.
他說見到一個地球.
一個地球這麼獨特.
是一個真空的狀態.
nothing的狀態裡面.
居然有一個藍色的一點.
藍色的一個地球.

$^{241}$他就說.
we always choose earth.
上個禮拜三我們祈禱會的時候.
麗仙也有分享這句說話.
但我想告訴你.
這個perspective.
看回這個地球.
we will always choose earth.
其實這句說話.
是.
上帝對亞伯拉罕說.
I will always.
choose earth.
我要拯救世人.
於是我要去.
呼召你.
要使你能夠成為祝福.
並且祝福不是一個死水.
不只是拿了上帝的祝福.
你是一個channel.
能夠藉著你去使.
萬國能夠得福.
很多時候我們看亞伯拉罕的故事.
第十二章到二十二章.
我們的焦點.
都放在那個人身上.
亞伯拉罕有多好.
我們說亞伯拉罕是什麼.
後期的文獻說.
他是信心之父.
信心之父.
或者以色列的國父.
我們看這些.
founding father的story.
founding story的時候.
都有一種很仰慕的感覺.
誰開創了什麼.
他是多麼的了不起等等.
誰開始了某些工作.
我們現在承載他.

$^{281}$承傳他的好處.
或者他的基礎.
我們覺得那個人.
特別好一點.
教會也是很多時候.
一個很不好的文化.
覺得教會的founding father.
一定是他了不起一點.
其實沒有這回事.
其實我們看亞伯拉罕.
都有這樣的錯覺.
在我們某程度上.
某一些的迷思上面.
我們覺得他特別好.
但是你看十二章到二十二章.
整個亞伯拉罕.
整個narrative.
其實描述一個亞伯拉罕是一個很真實.
很脆弱的人.
有饑荒他就逃跑.
然後想撿回他的命.
他就說謊.
或者他覺得等不及了.
他就give in.
去到人性的軟弱.
自己找自己的方法等等.
有戰爭有很多延遲.
有很多考驗等等.
其實很多時候他都是fail.
你不要看錯.
他是一個不堪一擊的人.
我想如果亞伯拉罕再生.
可能他很老.
如果你和他做過訪問.
我們的幻想就是.
這個亞伯拉罕是一個強人.
他講的話.
或者他做的事.
很厲害.
其實很可能.

$^{321}$如果你去訪問他.
他都是一個滿有病態.
不堪一擊很脆弱的一個人.
但是他口中告訴你的.
其實不是我可以.
是上帝可以.
所以整個亞伯拉罕的故事.
他敘述的時候.
雖然人很多軟弱很多脆弱.
很多懼怕跌倒等等.
但是背後的上帝.
他不斷地推展他的計劃.
會不會有挫敗.
有威脅的時候.
上帝的應許沒有變.
上帝可以work around it.
上帝可以克服所有的困難.
所有的難阻.
無論那些難阻是世界給的難阻.
或者人的破碎.
造成的難阻.
到最後的時候.
上帝可以get it done.
是上帝可以fulfill.
他的計劃.
所以我們看亞伯拉罕的故事.
我們不要說這個人是多麼英雄.
founding father是多麼值得我們仰慕.
NO! 值得我們仰慕的是上帝.
是亞伯拉罕背後的上帝.
如果你.
比較一下.
你知道舊約和新約之間.
兩約之間.
那些.
京外的文獻.
都有很多說.
亞伯拉罕以撒亞國.
甚至說他們的老婆.
他們有幾個老婆等等.

$^{361}$讚到他們.
讚口不絕.
連說撒拉.
亞伯拉罕的原配.
是多麼漂亮.
從頭讚到尾.
有這樣的女人.
真是沒有瑕疵.
這就是後期對那些founding people.
那些的幻想.
但在聖經裡面.
你看清楚一點.
原來不是這樣.
是上帝可以克服所有人的軟弱和脆弱.
他可以克服所有的困難.
以致他能夠成就他的計劃.
推進他的計劃.
使萬國能夠得福.
這是第一點.
我們要知道這位上帝是一個怎樣的上帝.
他可以克服我們所有的脆弱.
克服我們教會所有的不足.
克服我們生命裡面可能會fall and stumble.
fall and stumble again.
但我們所信靠的上帝.
就是一個可以繼續推進他的計劃.
他講過的不會出爾反爾.
他應許過的.
要萬國得福.
他就會做得到.
這就是亞伯拉罕的故事.
第一 介紹這位上帝.
第二 既然這位上帝是這樣.
我們的身份是什麼呢.
就是我們被邀請.
有這樣的恩典.
可以被邀請做他的子民.
我們能夠.
和上帝同供.
帶著使命.

$^{401}$來到他同供.
這一點在十二章已經很清楚.
你記得我們看那篇經文.
上帝和亞伯拉罕說什麼.
我要使你成為blessing.
並且藉著你.
使萬國能夠有.
blessing 藉著你.
讓人認識不是你多厲害.
是上帝是可以的.
上帝才是真神.
這就是一個新的.
一個的身份.
不是.
選亞伯拉罕不是因為他好.
選亞伯拉罕不是因為他行.
而是.
給他一個新的identity.
要藉著他.
雖然他不行也好.
使他能夠有一個使命.
全新的身份 使萬國.
能夠得福.
我這樣說不代表.
這個身份是容易的.
因為在亞伯拉罕的故事裡面.
我們不詳細交代.
有很多他自己.
跌倒的地方 失敗的地方.
一次又一次.
但是這個的應許.
因為他的失敗 因為他不行.
我們看下去.
他的故事的時候.
就好像我們看長篇劇.
我覺得主角為什麼這麼不行.
又不帥.
說話又不行.
然後他做的事.
又傻乎乎.

$^{441}$是一個很不濟的人.
什麼水彈散.
就好像你看長篇劇.
編劇搞什麼.
但最重要的就是告訴你.
背後的編劇才是濟那個人.
這個就是上帝.
他可以承擔所有人的軟弱.
到最後亞伯拉罕的故事.
由十二章開始到多少章.
二十二章.
中間大概.
你數一數大概有十個考驗.
十個.
對亞伯拉罕的考驗.
有可能有饑荒 有可能有軟弱.
可能有等候延遲等等.
到第二十二章發生什麼事.
就是以殺的故事.
獻以殺的故事.
獻以殺的故事怎麼看都是爛的.
但最重要的一件事就是.
上帝說給亞伯拉罕聽.
你要獻上你最寶貴的那個.
你最心肝寶貝的那個.
但你要獻上.
我們怎麼看都覺得.
有沒有搞錯.
上帝這麼沒人情.
但你看整個故事.
其實上帝不是真的要.
亞伯拉罕失去以殺.
他最主要想表達什麼.
就是想告訴亞伯拉罕.
你要信靠我的時候.
要有一個怎樣的信心.
就是不要再跟我說這麼多條件.
你對我的信靠.
應該是全心全意的.
這個才是最後的考驗.

$^{481}$你願不願意放下你的控制.
你讓我去控制.
駕車的人知道.
有誰很喜歡駕車.
一上車就是你駕車.
就不喜歡別人駕車.
有誰是這樣.
有誰是有這樣的病態.
就是你非要駕tie 盤的那個.
別人駕你總是坐不安樂.
別人煞車你又走去煞車.
但你沒東西按.
你明不明白我們有這種病態.
我們想控制.
我也是這樣.
我駕不是別人駕.
上帝在天意撒的故事.
就是在講這件事.
你走開.
我開車.
並且我開的車.
走的那條路.
是比你任何的路更加美好.
因為可能是有些adventure.
可能有很多驚喜.
你覺得失去了掌控.
但上帝.
是比你.
更好的司機.
上帝是更好的司機.
所以亞伯拉罕.
故事的結尾.
天意撒的故事結尾.
上帝又再一次講說話.
上帝講什麼.
二十二章 這裡完了.
天意撒 上帝也不是要你的.
上帝會告訴你 你對我的信心.
就是要你走開一邊.
讓我拿著tie 盤.

$^{521}$讓我去控制.
二十二章十五節.
你既行了者是沒有留下.
你的兒子 你的獨子.
你的獨生子.
我指著自己起誓 神說什麼.
我必多多賜福給你.
使你的後裔大大增多.
如同天上的星 海邊的沙.
你的後裔必得著.
受敵的城門.
地上的萬國.
必因你和你的後裔而得福.
因為你聽從我的話.
如果亞伯拉罕.
經過所有這些跌跌碰碰.
最後的時候 他明白.
這位上帝是正在宣教的上帝.
這位上帝是.
值得我們敬佩.
值得我們信靠的.
值得讓他拿著我們的tie 盤.
這個就是結果.
上帝說 我真的會.
I will fulfill my plan.
I will make it work.
You see I have made it work.
And I will make it work.
And I am making it work.
你的後裔必定會.
如同天上的星 海邊的沙.
一樣度.
你可以想像那個情況.
獻完以殺的故事.
是一個怎樣的意境.
亞伯拉罕可能在海邊.
沙灘.
上帝對他講說話.
你數一下.
海灘.

$^{561}$有多少粒沙在那裡.
如果你數得到的話.
我就是會賜福你.
這樣賜福你.
賜福到另一個地步.
你是數不盡的.
這個就是亞伯拉罕.
所介紹的上帝.
是一位.
願意去祝福我們.
只要我們願意.
願意去信靠他.
上帝對萬國的祝福.
上帝對萬國的旨意.
是堅定不移的.
這幅圖畫.
天上的星這麼多.
海邊的沙這麼多.
你試一下數.
數不到.
你見到亞伯拉罕的故事.
的交代.
由一二三四五六.
你就發現原來上帝.
真的很信實.
由一個很小很小的家庭.
還是沒有子女的家庭.
一個死亡的家庭.
變成一個家庭.
接著變成一個更大的家族.
以撒雅各.
接著到十二支派.
以色列國.
成為一個民族.
到最後.
是耶穌基督降生.
耶穌說藉著我的教會.
要使萬國.
能夠成為他的門徒.
萬國能夠得福.

$^{601}$你見到從亞伯拉罕直到.
耶穌基督.
上帝就是在實行.
祂自己的計劃.
在耶穌基督裡面就真正得到.
我們有的祝福.
你看看我們.
你看看我們教會.
你看看普世教會.
每一代的教會.
我們都會跌跌碰碰.
我們會散亂.
我們會有散有盡.
但是我們要見到的是.
我們背後原來這位上帝.
是一個值得我們.
事奉的上帝.
因為有亞伯拉罕.
有耶穌基督.
因為有耶穌基督.
有我和你.
誰會想像到.
在亞伯拉罕的時候.
可以想像到一個這樣的海沙.
這麼多的祝福.
不是因為人不行.
不是因為以色列國德.
不是因為祖父列祖.
好東西.
不是因為我們好東西.
而是上帝好東西.
我們在這一年.
新的一年.
我們希望和你.
去到一起經歷.
去經歷這個.
信實的上帝.
正在宣教工作的上帝.
我們怎樣能夠和祂同共.
所以在這一年裡面.

$^{641}$開球的時候.
我們希望送一個.
瓶子給大家.
你拿著這個瓶子.
我們今天這個瓶子.
只有一樣東西.
我們將來再說這個.
系列的時候.
我們會加一些東西.
進去裡面.
我想請你今天拿這個瓶子.
拿一個OK.
你不是一個人拿十個.
一個夠 為什麼.
你拿回來的時候.
你數一下有多少.
沙在裡面.
你數一下有幾粒.
你下個禮拜告訴我.
我保證你數到下個禮拜.
你還沒數完.
這個就是上帝的應許.
就是上帝自己的本性.
縱使我們.
是軟弱.
我們不大.
我們不堪一擊.
但是上帝.
但上帝.
先是無堅不摧.
他值得我們信.
值得我們.
去事奉他.
所以請你拿一個瓶子回去.
這裡有我們的UV的蓋印.
你可以回去寫你的名字.
寫日子.
在這一年裡面我們都會用這個瓶子.
就好像你去遠行的時候.
我們去hiking的時候.

$^{681}$我們拿著這個瓶子.
或者你的裝備.
裡面每一樣東西可以提醒我們.
這位上帝.
值得我們全心全意去愛.
全心全意去事奉.
對嗎?.
海沙可以提醒我們.
上帝最早的.
應許.
提醒我們教會.
我們希望能夠參與上帝.
宣教中的工作.
我們成為海沙的一部分.
最終我們願意更多海沙.
能夠參與到.
享受到上帝自己的同在和祝福.
好不好?.
我們同心協力去祈禱.
主祈求你教導我們.
信靠的不是我們自己.
因為我們實在太渺小.
因為我們實在太.
不堪一擊.
我們實在是脆弱的人.
主愛我們的生命.
我們只不過是爛命.
一條.
但你卻以恩典.
你的愛.
來拯救我們.
不是因為我們行.
而是因為你.
你是行的.
主要讓我們認識你.
是一位宣教中的上帝.
要使萬國的心.
歸向你.
主要我們已經歸向你.
在基督裡面領受祂的祝福.

$^{721}$讓我們這個.
軟弱的生命.
能夠成為一個帶著使命的生命.
去認識你,去事奉你.
以致我們個人生命.
我們的家庭.
我們每一個的每一方面.
能夠去認識你的恩典.
並且叫人去認識你自己的名字.
奉耶穌基督的名字.
願上帝祝福各位.
\newpage



\section{出埃及記 13:18}
\label{sec:031962}
\textbf{聖經人聖經事:神領百姓繞道}~~~~黃克勤牧師博士
\newline
\newline
連結: \href{http://www.church.com.hk/acms/content.asp?site=cdc\&op=show\&type=product\&code=031962\&layout=sermon}{\texttt{031962}} ~~~~ 語音日期: 2026-04-08
\newline
\newline
\hyperref[sec:031961]{\small{< < < PREV SERMON < < <}}
~
\hyperref[sec:index_chronic]{\small{[返順時目]}}
~
\hyperref[sec:031963]{\small{> > > NEXT SERMON > > >}}
\newline
\newline
出埃及記 13:18
\newline
\begin{longtable}{cl}
\hline
\hline
章節 & 經文 (和合本修訂版)\\
\hline
13:18 & \begin{tabularx}{0.7\textwidth}{X} 神領百姓繞道而行，走曠野的路到紅海。以色列人出埃及地，都帶著兵器上去。 \end{tabularx} \\ \\
[1ex]
\hline
\hline
\end{longtable}
$^{1}$各位靈姐妹.
以色列人在摩西帶領出埃及入迦南.
現在來到一個新的篇章.
就是他們出了來埃及.
都走了一段路.
快要來到一件很大轉折.
很大件事了.
新穎樹宣告網絡跟大家一起看聖經的人與故事.
看背後上帝大能的手.
今次的題目是神領百姓妖道.
為了繞路.
妖道.
為什麼神很多時候.
我們都不帶祂直行直直.
一條直路就到了.
如果我們認識這件事.
你會發覺以色列人在曠野漂流了40年.
這條路這麼遠嗎.
真的整個地球都走了40年.
但是祂在一個很小的地方.
神為什麼要祂妖道不直行呢.
這裡有解釋了神領百姓妖道的故事.
以色列人出埃及入迦南的路線圖.
這個尼羅河.
去到近海的地方變成一個三角洲.
自然的.
水流得比較寬.
這塊地很肥沃的尼羅河三角洲.
在當日叫做戈山.
就是戈山地.
是塊土地.
當時以色列人住在戈山.
這麼好.
因為當時約瑟在做宰相.
法老皇上.
我的家人來了.
他們牧羊找塊地居住.
可不可以讓他們住.
沒問題.
你選最好的地住.

$^{41}$在哪裡住.
我說可以.
其實你說已經夠了.
於是他們住在戈山.
後來約瑟死了.
不認識約瑟的新王起來.
新的民族入侵了.
可能是埃及裔的.
叫做Hicksos.
做了埃及的王.
朝代改變了.
他們不喜歡這些以色列人.
但當時以色列人已經有近兩百萬人.
男男女女.
挺好.
難怪以前這麼旺.
有二百萬免費的人在做事.
做國民也挺好.
開始將他們剝削.
令他們變成奴隸.
逐步逐步變成奴隸.
免費嘛.
這些外族人.
於是奴役他們.
要蓄貨蓄城.
收藏糧食,器具.
即貨倉.
其中一個叫蘭塞,比東.
兩個積貨城.
蓄得很辛苦.
還要殺死他們.
種族滅絕.
要將男生殺死等等.
這些故事大家知道了.
後來摩西終於帶他們離開埃及.
他們在蘭塞起行.
到比東一起走到疏國.
蘭塞,比東,疏國.
疏國一併集合.
開始向南走.

$^{81}$來到今天以唐的地方.
接著在麥奪,比哈希祿,巴力洗勳對岸.
在這裡過紅海.
過紅海在這裡.
所以看到這裡很窄.
是紅海最窄的部分.
為什麼不在.
其實他們要去的目的地是.
耶利哥.
再上就是希伯倫.
再上就是耶路撒冷.
耶路撒冷大概這個位置.
為什麼不在這裡.
疏國,比東一行就這樣走.
這樣走海邊的路.
來到加薩即是非利士.
這裡叫非利士路.
走過非利士這個民族.
加薩即是今天以色列的加薩地帶.
全部都是非利士人.
在這裡很近.
但是他們這樣繞道.
然後這樣才進.
為什麼這樣走了四十年呢.
神有祂的心意.
神沒帶他們這樣.
反而帶他們在這裡過紅海.
在西乃半島走到河以山.
即是西乃山或叫河列山領受十佳.
過程是這樣的.
路線是這樣.
十三章第十七節.
在法老容伯勝去的時候.
非利士的土地雖然近.
但是神卻不領他們從那裡走.
出來的時候其實非利士在這裡.
加薩.
在這裡的路雖然近.
確實是雖然近.
但是神沒領他們走那裡.

$^{121}$非利士地道路雖近.
神卻不領他們從那裡走.
因為神說恐怕百姓遇見打仗.
後悔就會回去埃及.
因為如果他們走到那裡.
很有可能他們會跟非利士民族打.
當地的民族海邊這個民族打.
他們走到那裡會遇到戰爭.
打仗了,不想打了.
總之有人傷亡都不想怎麼辦.
走回去埃及住吧.
這裡講得很清楚.
恐怕百姓遇見打仗後悔就回埃及去.
但是其實神就領百姓繞道而行.
第一點要講的.
神是特意領百姓繞道而行.
走紅海曠野的路.
其實以色列人出埃及地都帶著兵器上去.
都預了打.
雖然預了會打.
但是其實他們叫耶和華的軍隊.
但是其實他們不善戰爭.
個個都有一些刀劍.
有一些武器.
但是其實少打是可以的.
但是你說軍隊打是他們不熟的.
沒有練過兵.
沒有練過陣.
其實是不行的.
很難對他們打的.
除非神在那裡出手.
神是帶百姓繞道而行.
是怕他們這麼快一出就遇到戰事.
其實接下來都會有很多戰事.
不過你一出就遇到戰事.
你就會後悔馬上走.
神知道他們.
所以帶他們走紅海曠野的路.
特別是預帶他們過紅海.
特別預帶他們過紅海.

$^{161}$想了計劃.
想了路線.
所以第一點你要記得.
神為什麼不讓我走直路呢.
我老老實實說.
我這個人其實不是一個讀書人.
我小時候喜歡踢球.
喜歡玩.
喜歡到處走.
坐不定.
但是我真的完全沒想過.
我想認識我的人.
包括我家裡.
我所有團友.
當日中學的團友.
這個瘦瘦猛猛的瘦柴.
又不喜歡讀書.
英文又不厲害.
沒有一件事可以.
你覺得他很差.
連我家人.
我弟弟.
甚至當時團契的導師.
都看不起我.
覺得這個不行的.
懶得說.
唯一是什麼都是高.
比別人高一點.
不過我小時候很瘦.
但是路一直走.
帶我走了很多路.
讀神學.
怎樣再讀大學.
怎樣再走石屎.
這樣上來.
到我拿到兩個博士的時候.
我已經五十幾.
近六十歲了.
繞路的時候.
為什麼帶我繞路.

$^{201}$一早就帶我讀書.
馬上上去讀博士.
讀兩個博士.
博士後出來工作.
服侍你不就好了.
為什麼繞這麼遠.
才去到應許之地.
我不知道.
是我一路走一路走.
神就這樣帶我.
我就跟著祂走.
有些是自己猛猛地.
有時是這樣.
有些是想自己走各樣的路.
神是會領.
這個詞是領他的百姓.
有時是繞到你不知道.
因為要避了某些事.
或者要遇到某些事.
要避了這些事.
要遇到紅海.
可能是這樣的.
要避了這些事.
跟菲律賓打仗.
但是要遇法老的軍兵追來.
然後要過紅海.
有時候神這樣帶我們繞路.
你不知道.
神啊你開的士嗎.
繞路嗎.
對不起.
有很多的士不繞路.
那些是繞路的.
我也坐過很多的士繞路.
繞道的.
其實這個繞道.
如果是神帶.
不差的.
有時繞一繞道.
學多了很多東西.

$^{241}$是我們生命必須的.
摩西將約瑟的孩軍一同帶去.
因為約瑟曾經叫這些人嚴嚴地起誓.
對他們說.
神會祝福眷顧你們.
不過你們回去故鄉的時候.
要記得將我的孩軍從這裡一起帶去.
我不屬於埃及的.
雖然我在這裡做宰相.
我不屬於這裡.
我屬於我的故鄉.
神給我的地方.
記得帶我去.
這個心在約瑟來說.
雖然是帶著約瑟的孩軍的棺木.
但是對以色列人也是一個鼓舞.
約瑟雖然死了很久.
但是仍然是一個鼓舞.
見到約瑟的遺骸的棺木就會記得.
他也要回去.
他在這裡做宰相.
但這裡也不是他的故鄉.
他要回去.
這樣就會令他們很有激勵.
他們從疏國起行.
在曠野邊的耳塘安逸.
耳塘這個字.
耳塘就是與他們同在的意思.
在那個地方神說去耳塘安逸.
帶他們去那裡紮營停留.
為什麼去這個地方呢?.
我覺得神是有旨意的.
我自己的解釋就是因為這個地名.
意思就是與他們同在.
神是會帶我們繞路的.
有時候走一條路.
你沒想過是這樣.
不直的.
那就對了.
就去走一條不直的才學到東西.

$^{281}$不過無論怎樣直與不直.
最重要是神與他們同在.
這個地方的名字.
意思就是與他們同在.
可能你以前有些事情在那裡一起.
一起的意思.
挺好的.
神跟他說我們一起.
無論去到哪裡.
繞路不繞路.
直路不直路.
彎路不彎路都好.
總之一起走.
我與你同在.
我與你一起住.
在那個地方安逸.
有一個神同在的感覺.
這裡叫什麼?.
我們在這裡紮營.
這個地方叫什麼?.
我們還沒來過.
當然還沒來過.
曠野來的.
這個地方叫耳塘.
是什麼意思?.
就是與他們同在.
神與我們一起.
好,就在這裡紮營.
日間耶和華在雲柱中領他們的路.
夜間在火柱中光照他們.
使他們日日都可以行走.
由這裡開始.
日間雲柱.
夜間火柱.
帶著他們白天看這條雲柱.
跟著走就行了.
有火柱晚上照著.
永遠在他們周圍旁邊有光.
雲柱火柱是怎樣的?.
我告訴你.

$^{321}$我都沒見過.
我見過的頂兩屋都是別人在想.
一個圖畫畫出來.
雲柱多數都是捲捲捲的雲.
火柱紅紅的像火一樣在燒.
照到周圍都光了.
就在四面安營.
這些一堆堆的營地.
那些十二支牌在那裡住.
他們在中間有火.
回望在旁邊.
多數都是看這樣的圖畫.
當時來到這裡.
回望又沒有.
那些百姓都是自己喜歡怎樣住怎樣住.
未曾分地區的.
不過有雲柱有火柱逐步來.
雲柱走他們跟著走.
火柱起程走.
令到他們日夜都可以行走.
這個雲柱火柱有意義的.
在生命的故事裡面就是.
你不論日夜.
只要見到神的同在.
你都可以行走的.
放膽行走.
第22節.
日間雲柱夜間火柱.
總不離開百姓的面前.
這個第三點要留意的就是.
神總不離開你.
神會叫你在一個地方安營.
那個地方叫以堂.
意思是什麼?.
就是與一齊同在.
與你同在這樣的意思.
而雲柱火柱都是日夜的同在與帶領.
我們的神就是這樣的.
但我們行人生的路永遠不會離開.
永遠不會丟下裡面的孤兒.

$^{361}$哎呀!爸爸媽媽在哪裡?.
不會.
白天用雲柱.
晚上用火柱.
叫我們知道神總不離開你.
離開他的百姓的面前.
即是看到.
面前即是你帶著他走.
哎呀!雲柱火柱在哪裡?.
神在我後面.
那我怎樣走呢?.
不是,是在你前面讓你看到的.
所以你放心.
神會這樣帶領你.
一定要順靠主他的帶領.
留意著雲柱火柱.
留意著他怎樣帶領.
相信他的同在.
有時路好像與我心目中不同.
照我知去流禮物之地應該是那邊的.
為什麼要走這邊呢?.
神有心意的.
如果知道雲柱火柱這樣走.
你儘管跟著去就行了.
心意就是這樣.
一起祈禱吧.
求主幫助我們更加進行你的心意.
知道你每一個時間.
不論日夜你都帶領我們.
你是與我們同在的.
雖然路有時未必跟我們想的直線.
但是有你一起.
我們就知道是最好了.
求主祝福給我們對你更多的信任.
讓我們跟著你走人生的路.
有時路有時兜兜地.
為什麼會這樣打轉.
不過知道是你的帶領.
是你有心意的.
願意順靠你.

$^{401}$順聽祈禱.
奉耶穌基督名求.
阿們.
\newpage



\section{馬可福音 6:31-34}
\label{sec:031963}
\textbf{無處不在的基督}~~~~周弘毅牧師
\newline
\newline
連結: \href{http://www.church.com.hk/acms/content.asp?site=cdc\&op=show\&type=product\&code=031963\&layout=sermon}{\texttt{031963}} ~~~~ 語音日期: 2026-04-11
\newline
\newline
\hyperref[sec:031962]{\small{< < < PREV SERMON < < <}}
~
\hyperref[sec:index_chronic]{\small{[返順時目]}}
~
\hyperref[sec:031964]{\small{> > > NEXT SERMON > > >}}
\newline
\newline
馬可福音 6:31-34
\newline
\begin{longtable}{cl}
\hline
\hline
章節 & 經文 (和合本修訂版)\\
\hline
6:31 & \begin{tabularx}{0.7\textwidth}{X} 他就說：「你們來，同我私下到荒野的地方去歇一歇。」這是因為來往的人多，他們連吃飯的時間也沒有。 \end{tabularx} \\ \\ \relax
6:32 & \begin{tabularx}{0.7\textwidth}{X} 他們就坐船，私下往荒野的地方去。 \end{tabularx} \\ \\ \relax
6:33 & \begin{tabularx}{0.7\textwidth}{X} 眾人看見他們走了，有許多認識他們的，就從各城步行，一同跑到那裡，比他們先趕到了。 \end{tabularx} \\ \\ \relax
6:34 & \begin{tabularx}{0.7\textwidth}{X} 耶穌出來，見有一大群的人，就憐憫他們，因為他們如同羊沒有牧人一般，於是開始教導他們許多事。 \end{tabularx} \\ \\
[1ex]
\hline
\hline
\end{longtable}
$^{1}$領治梅伍安.
今天雖然我要跟大家分享那段經文.
在《馬可福音》六章三十一至三十四節.
好像剛才主席所說.
不過我在開始分享之前.
都很想邀請大家留意一下.
《馬可福音》其中一個很值得我們注意的地方.
就是在成卷《馬可福音》當中.
有一個希臘文的單字.
叫做Euthyos.
這個字其實在成本新約聖經加起來.
都出現了五十一次.
不過只是在《馬可福音》就出現了四十一次.
即是你見到這個字.
八成的機會都在《馬可福音》出現.
這個字本身的意思寫給你了.
是「立即」的意思.
所以如果按照《馬可福音》的技術.
主耶穌服事每一件事完了第二件事.
好像每件事都要馬上做完.
如果你們有聖經在手的話.
你可以立一立《馬可福音》第六章整章聖經.
你就會見到這些事情的發展其實是很緊湊的.
第六章開始的時候.
就說到主耶穌在祂自己的家鄉.
人們不相信祂.
拒絕相信.
於是去到第七節.
祂就馬上就將門徒拆遣出去.
去其他地方傳.
這裡的人都不聽了.
在完成了十二至第十三節所說.
那裡叫他們做些什麼呢?.
去外面醫病,趕鬼,傳福音之後.
接著就發生「五餅二魚」.
這個大家很熟悉的神蹟.
接著說到祂在海面上行走.
之後就在格尼沙勒這個地方醫病.
可以見到主耶穌祂自己關於天國的服事上很趕急.
馬不停蹄.

$^{41}$所以有聖學者就幫馬可福音起了一個別名.
叫做「The Quick Gospel of Mark」.
我將它翻譯成「快D 」的福音.
什麼都要快D .
每樣事情都要馬上做.
若然在主耶穌祂生活的年代.
生活都這麼忙的話.
其實二千多年前.
今天我們就更加不用說了.
到底忙碌是否二千年都解決不了的問題呢?.
一直都是這樣的.
如果我們要找一條出路.
是否我們要放下我們手頭上所做的事情.
然後天天去受苦節日營.
這樣就可以了.
夠安靜了.
才體會到主耶穌的臨在呢?.
今天我跟大家分享的.
馬可福音的第六章三十一至三十四節.
其實正正是在描述一個這樣的情況.
四節聖經而已.
我想先請大家跟我們一起讀了它.
好,預備起.
「他就說:你們來.
跟我暗暗地到曠野地方去歇一歇.
這是因為來往的人多.
他們連吃飯也沒有功夫.
他們就坐船暗暗地往曠野地方去.
眾人看見他們去.
有許多認識他們的.
就從各城步行.
一同跑到那裡.
被他們先趕到了.
耶穌出來.
見有許多的人.
就憐憫他們.
因為他們如同羊.
沒有牧人一般.
於是開口教訓他們許多道理」.
經文是開始.

$^{81}$第三十一節上半節那裡這樣說的.
「他就說:你們來.
跟我暗暗地到曠野地方去歇一歇.
為什麼?.
因為來往的人多.
他們連吃飯也沒有功夫」.
其實先前所發生的是什麼事呢?.
為什麼看成這樣?.
我們看一看三十節.
「原來當其時試圖聚集到耶穌那裡.
將一切所作的事.
所傳的道.
傳告掃他」.
因為他們已經完成了.
早在第六章第七節.
主耶穌差遣他們去工作.
做些不同的工作.
於是回來之後.
就和耶穌來匯報.
你可以想像一下當時的場景是怎樣.
他們圍著耶穌.
然後七嘴八舌.
滔滔不絕地告訴耶穌.
「剛才我們施行過神跡醫治.
驅趕邪靈.
很多人悔改.
願意歸向神」.
你講一句 我講一句.
他們爭先擁後.
其實就是想耶穌.
來聽到他們那些非常令人震撼的經歷.
既然是這樣.
應該就打鐵趁熱.
我們也是這樣的.
剛剛缺字.
當然是邀請他參加福音班.
難道等到他癱瘓了.
兩個月之後才邀請他嗎?.
不會的.
所以匯報完之後.

$^{121}$其實最好是再接再厲.
再出去傳福音.
不過主耶穌聽完他們的匯報.
這張很亮麗的成績單交出來之後.
耶穌就留意到他們很忙.
為什麼呢?.
三十一節下半節.
因為來往的人多.
吃飯也沒空.
看到這樣的情況.
主耶穌給了門徒一個反高潮.
Anti-climax的呼喚.
叫他們靜下來.
跟我暗暗地到曠野地方去歇一歇.
如果我問你們覺得自己的生活很忙嗎?.
我相信這個問題有點多此一舉.
還用說嗎?.
當然很忙.
但如果我再追問.
為什麼這麼忙?.
忙什麼?.
我不知道你會怎樣回答.
因為工作很麻煩.
時間很長.
我又要進修增值.
家裡又有家人需要照顧.
撇開這些沒得選擇的範疇.
必須做的事情.
靜下來的時間似乎真的不多.
但不知道為什麼.
當我去問那些已經退休的弟兄姊妹.
他們也一樣忙.
非常忙.
他們說時間多了.
但仍然不夠用.
很簡單.
其實我們實際上真的覺得很難有多出來的空間.
原因是如果我們有多出來的空間.
我們就會放東西進去.
令到它塞滿為止.

$^{161}$就算是一些休閒的活動.
想告訴大家.
我們也會令自己很忙.
休閒的活動包括娛樂.
剛剛我們才過了清明和復活節重疊的假期.
因為兩者重疊.
所以大家都知道.
就算你不用請假.
你也有五天的假期.
如果你要自我製造一個長一點的假期.
你就請三天假.
加起來就有十天.
你也會計算的.
如果在十天這麼長的假期.
你會做些什麼呢.
很多人會選擇去旅行.
新聞就告訴我.
我看完之後就覺得挺厲害的.
原來在這個假期裡面.
是有212萬的人離開香港外遊.
212萬是什麼數字呢.
就是說三個裡面就有一個人去了外遊.
不管是長還是短.
總之就是離開香港.
第一時間離開香港.
然後去到上班之前那一晚.
才晚一點回來.
是不是很忙.
你還不要計算.
現在旅遊是輕自由行的.
就是之前你要做功課的.
然後你就去哪裡玩.
訂那些地方住.
那就搞一輪了.
然後回來之後.
也需要花點時間撿回你的美照.
那也是時間來的.
所以其實在我們的生活當中.
是忙得不得了.
我再舉多一個例子.

$^{201}$我不知道像不像你.
現代人基本上都相當晚睡.
十二點多就叫做正常.
有些人就去到半夜兩三點.
放了工 放了學.
甚至去完街.
回到家要休息一下.
玩完也要休息一下.
很悠閒 漫無目的來滾一下手機.
在現在這個資訊爆炸的年代.
社交媒體越來越多.
我們不可以這麼落後.
這麼無知的.
一定要和我們身邊的世界保持聯繫才可以.
有一家專注為手機.
加電科技產品提供維修服務的公司.
我在網上看到的.
它在2022年 也幾年前了.
我相信現在不止這個數也不止.
做過一個調查.
它說二千個美國的成年人.
想了解一下他們用手機的習慣.
結果發現差不多有四分之三的人.
都會帶著手機上床.
然後滾一下手機才睡覺.
可能你未必帶上床.
不過滾一下手機才睡覺.
這個也是很普遍的事情.
甚至有百分之.
拿著手機的那些.
百分之九十三的人.
當你滾完手機之後.
你會將你的手機放在一個.
你的手伸到可以拿到的地方.
你滾完之後放在床邊.
你的手可以拿到的.
百分之九十三.
甚至有百分之十的人.
會將他的手機放在枕頭底.
這個真的不太妥當.

$^{241}$大家都聽過很多的輻射研究等等.
所以.
希望不是在說你.
拿著手機滾下滾下.
一兩個鐘頭.
我想告訴你.
很快過的.
不過不少人覺得.
這個其實是在他們一日的生活當中的黃金時間.
因為這個是me time.
是能夠可以自己的時間.
甚至有人這樣說.
如果人生當中.
只有上班和睡覺的話.
其實不是很值得活下去.
所以無論是工作.
無論是進修家庭.
甚至是消閒活動也好.
通通都令到我們非常忙.
可能大家都有聽過.
有人將忙碌英文busy.
這個字就引申成為.
being under satan's yoke.
即是忙到令到我們被壓制在惡者魔鬼的壓下.
太忙不行.
所以主耶穌就說.
一定堅持要帶領門徒.
抽時間去曠野.
剛才那段經文是這樣說的.
即是四野無人的地方.
來休息一下.
事實上休息在某一福音當中.
這個也是其中一個主耶穌不斷強調的主題.
因為至低限度有九次.
主耶穌是明確地說到.
要抽出時間來休息.
安息和祈禱.
所以這個屬靈功課.
不是單單針對最初祂所呼召的門徒.
這個屬靈功課.

$^{281}$也是我們今天基督徒生活.
一個非常重要的部分.
據說昔日有一個宣教士.
他去到偏遠的地方.
非洲的森林偏遠的部落想去傳福音.
因為他要帶上他的行李.
所以就僱用了當地的一些人來幫他搬.
第一天他們行程是比較快.
進度很好.
所以已經是走了比較遠.
宣教士滿懷希望.
可以快點去到他想去的地方.
上帝帶他去的目的地來傳福音.
夜晚休息過後.
第二天又起行了.
不過到第二天起床的時候.
當地人就沒有反應.
不知為何.
無論如何都坐在地上.
無論如何都不肯動.
擾攘了一段時間.
翻譯就解給宣教士聽.
他說我們昨天走得太快了.
所以要等我們的靈魂追上來.
才趕得及.
這樣啊.
如果我們的工作和生活.
都好像剛才所說的那種忙碌狀態.
看來我們真的需要休息一下.
等我們的靈魂追上來才可以.
偉大的報道家施暴真.
他曾經這樣說.
"Rest time is not waste time".
"It is economy to gather fresh strength".
"It is wisdom to take occasional follow".
"In the long run, we shall do more by something".
"We shall do more by sometimes doing less".
休息的時間不是浪費時間.
而是為了積蓄力量.
偶爾放假是明智之舉.

$^{321}$從長遠來說.
我們可以做少一點.
反而可能做到更加多.
我不知道你同不同意這個說法.
一味不停地做.
其實是不太行的.
在《馬可福音》裡面.
主耶穌多次說到.
他需要祈禱,安靜,休息.
其實他的門徒也是一樣.
如果門徒是沒有歇息的機會.
沒有時間和主耶穌單獨安靜相處.
其實他們是沒有力量來迎接未來的挑戰.
經過了興高采烈的彙報之後.
門徒願意聽,叫他們停下來.
經文三十二節這樣說.
他們就坐船暗暗地往曠野地方去.
但怎知眾人看見他們去.
有許多認識他們的.
就從各城步行.
一同跑到那裡.
被他們先趕到了.
人生的經歷很少能完全按照我們的計劃發生.
不是這樣的.
計劃歸計劃.
但我們往往會遇到一些出奇不意的事情.
一些爛jor ,一些挫折.
我們又試想像一下.
當門徒的船差不多靠近岸邊的時候.
其實他們已經看見有一大堆人.
在岸上等待他們.
你猜猜門徒當時的心情和表情是怎樣的.
可能他們心裡想.
有沒有搞錯.
為什麼那些人走路都比我們坐船快.
他們可能很渴望見到耶穌.
於是他們跑來好像馬拉松一樣.
跑得快過耶穌和這班門徒.
預早去等他們.
我們都想像得到.

$^{361}$這個情景是在他們成群人的意料之外.
他們面上的表情可能覺得很無奈.
原本想離開休息一下的.
可以有片刻寧靜.
結果事與願違.
回到第六章.
其實前面有一個段落.
我剛才沒有提到.
就是講到希律王.
他提到在自己生日的時候.
他擺設筵席.
將施洗約翰斬首.
當中有一段是有這樣的技術.
聖經沒有直接講.
不知道這班門徒.
是否因為收到這個消息.
令到他們的傳道工作.
都是一些令人心痛的事情.
所以被打斷了.
不知道.
但在我們的生命裡面.
其實往往都會有一些類似的事情發生.
我們可能突然間遇到健康會出現問題.
人際關係出現了挫折.
生活上可能更加是家常便飯.
明明好好的在做一份工作.
公司說要裁員.
又或者不跟你續約了.
正在住的那間屋.
業主說要收回.
他說要賣了.
大幅加租了.
還不算好像早幾個月前.
大埔那單的大火.
這些災難性事情.
所以有人為了苦難.
取了這樣的定義.
他說苦難是沒有選擇.
沒有目的.
沒有解釋.

$^{401}$完全是偶發性的.
這些看來好像很無奈的看法.
相信一定會令到人感到相當疲累.
就算你一日復一日.
你撐得住也好.
相信心裡也會覺得很沉重.
因為我們的心理.
跟我們的身體健康是彼此關連的.
所以或多或少都會影響的.
我不知道你們今天會不會面對一些很艱難的處境.
你覺不覺得累.
我們會怎樣面對這些突然發生的事呢.
我們的反應.
可能會受到我們的性格特質.
和我們的生命狀態影響.
有些人無可否認.
他承受壓力的能力是不同的.
有些人是能夠受壓多些.
有些人是不能夠受壓的.
高些能夠受壓的.
處變不驚.
沒那麼好的.
很緊張.
不行.
除了將壓力內化.
影響我們的情緒和健康之外.
也有人會將壓力.
用一些外在的方式來釋放.
這些忙得倒捨羅蟹.
但一定要交的功課.
一定要做.
你最好不要來惹我.
否則我就發火了.
究竟是不是這樣呢.
或者我會選擇逃避.
走投了.
不做了.
算了.
若然你之前曾經面對過這樣的處境.
又或者你現在是一個這樣的困局當中.

$^{441}$其實都值得想一想自己的反應.
在這樣的情況下.
在這個受壓的情況下.
究竟我們是否能夠體會到.
耶穌的同在.
是不是真的不行呢.
我們再回到《馬福音》六章段落的場景.
喂.
主耶穌.
又是你叫我們去曠野去歇一歇.
剛才也是你說的.
累死了.
現在這班人竟然不斷追著我們.
沒有辦法.
怎麼搞.
要不你就忍住.
不要發灰氣.
要不就不如這樣吧.
船掉頭了.
去別的地方吧.
當年可能也有很多曠野.
這個曠野不行就去別的吧.
就不要理它了.
走吧.
就算你真的上岸了.
也可以打發那些人走.
我不一定要在當中.
來回應它的需要.
我再說一次.
這個去曠野歇一歇的做法.
是耶穌提出的.
但是耶穌提出的方法.
竟然是行不通的.
是你叫我做的.
不過不可行.
耶穌祂自己怎麼反應.
這次是不是輪到.
耶穌不耐煩發火了.
不是.
我們看看34節.

$^{481}$這裡怎麼說.
耶穌出來.
見有許多的人就憐憫他們.
因為他們如同羊沒有牧人一般.
於是開口教訓他們許多道理.
一會兒就有五餅二魚大家野餐.
不過野餐我不打算請那麼多人來.
而在過程當中.
這些是不測之客.
但是耶穌的取代是什麼.
憐憫他們.
這樣的做法.
是不是跟祂當初提議.
叫他們去曠野歇一歇.
其實是互相矛盾的.
憐憫這個字.
其實在《馬福音》一章41節出現過.
那裡是另外一個翻譯.
叫做「動了慈心」.
這個字其實是一個相當強烈的字.
在新約只是用在耶穌身上.
這個字根是來自另外一個希臘文的字.
叫做「慧」.
即是「慧部」的「慧」.
「長慧」的「慧」.
新約聖經就在這個字上.
讓我們想像得到.
進入了肚腹裡面.
加了一份情感.
進到耶穌的內心當中.
因為詩篇116篇15節.
其實都說過給我們知道.
耶護說:有因為有公義.
我們的神以憐憫為懷.
所以這個憐憫.
正正是上帝的屬性.
為什麼耶穌要憐憫他們呢?.
34節下半節就說了.
因為他們如同羊沒有牧人一般.
「如同羊沒有牧人」.

$^{521}$其實是7節文述的27章第17節.
是說這群以色列民.
是需要一個顧念他們.
以他們的好處為著想.
去帶領他們的領袖.
沒有牧人的結果就是會失去方向.
羊不是盲忠忠.
不知走到哪裡掉死了.
都會成為其他猛獸的獵物.
結果都是死路一條.
主耶穌在《約翰福音》10章11節.
說自己是一個好牧人.
他願意為羊捨命.
而這群一直追隨著他的群眾.
就真的很需要耶穌.
回到我們今天的處境.
事實上我們是否真的忙到.
無法體會到耶穌的同在.
想休息又不行.
有功課要教 有事要做.
那怎麼辦.
平衡了兩件事之後.
結果努力一點.
努力去做就打贏了.
去曠野揭一揭.
是不是這樣.
整天都沒辦法.
唯有馬死落地走.
所以我們真的需要間中揭一揭.
去參加壽苦節日營.
去退休一下.
去打一支針.
補一補.
福音書裡面.
我們不時都會見到.
這裡引了兩段經文給大家.
馬福音一章35節講到.
次日早晨天未亮的時候.
耶穌起來到曠野地方去.
在那裡禱告.

$^{561}$路加方五章講到.
耶穌卻退到曠野去禱告.
很多這些描述.
這些經文是提醒我們.
耶穌很喜歡祈禱.
祈禱是在祂的服事生命當中.
是很優先的一個需要處理的項目.
是十分十分之重要.
不過請你留意一下.
這些描述的上下文.
這些場景.
全部都是講到.
耶穌忙到脫光衣服的時候去做的.
耶穌不是因為祂很休閒.
所以祂就祈禱.
而是很忙碌的時候這樣做的.
耶穌很希望人能夠爭取休息.
安靜親近祂的機會.
在過程當中重要.
不過同樣祂也會回應.
那些尋求祂幫助的人.
在這方面祂會讓步的.
祂憐憫每一個需要祂的人.
竭盡所能去幫助他們.
祂珍惜每一個休息和祈禱的機會.
但是祂真正的優先次序.
仍然是實踐祂的福音使命.
所以講得白一點.
帶著耶穌的福音使命忙碌的生活.
才是真正效法基督.
你聽清楚.
是忙的.
不過是帶著耶穌福音的使命去忙.
如果你是生活得休閒的話.
對不起.
這並不是耶穌的真門徒.
耶穌的真門徒不是空閒到僵硬.
沒有這樣的事.
我們需要彼此提醒.
我們不要將我們所有的精力.

$^{601}$都放在工作.
放在我們的事業上.
我們要夠休息.
好好照顧身體.
我們還有家.
還有社會等等的責任.
我們需要肩負的.
除此之外.
如果我們信了耶穌.
我們要努力去操練我們的屬靈生命.
去參與教會的聚會和服事.
是不是這樣呢?.
但是請我們不要將忙碌的工作.
和靈命放在一個對立面.
兩個不是互相對立的.
忙碌和親近神.
並不是水火不相容.
忙碌不是靈命的敵人.
而是實踐靈命操練的場景.
你忙得來.
仍然可以親近上帝.
甚至是一個必要的前提.
沒有問題的.
你在公司.
如果你一忙就亂.
就發脾氣.
那你能夠在工作上的表現.
我相信是極有限的.
就是這樣.
但是如果你很忙.
你仍然能夠從容不迫.
做好你要做的工作.
就行了.
屬靈生命是一樣的.
忙碌的時候.
上帝在不在?.
值得思考的問題就是.
為什麼我們在日常生活.
忙碌的時候.
我們就沒有辦法親近上帝.

$^{641}$然後我們的靈命就低落.
哎呀 回教會都沒有時間.
讀聖經很累.
晚上要睡覺了.
如果我們有這個前設的話.
你推論出來.
很簡單.
我們在教會.
我們有家庭服務中心.
兩個教育中心.
在裡面服侍的.
牧者和童工.
我們每天都在服侍上帝.
每天都在親近上帝.
我們的靈命就會低落.
是不是這樣?.
讀神學的更好.
潘姑娘坐在後面.
不用上班.
拉長了屬靈操練的生活.
每天都要入營.
是不是這樣?.
跟上帝的關係是否一定緊密一點呢?.
真實的信仰是關乎你真實的生活.
生活的部分.
理應跟我們的信仰是放在一起的.
我們每天.
無論你是上班多少個小時.
你要花多少心力去照顧家庭.
要陪爸爸媽媽去覆診.
連長有很多覆診.
全部都是我們親近上帝的靈性活動的一部分.
梁家倫牧師說話有些尖銳.
他說如果我們每天只是朝九晚七.
可能有些不止.
朝九晚八朝九晚九一跨宗.
忙碌生活是跟神絕緣的話.
而只是能夠在忙裡偷閒.
我們才間歇性地去親近神的話.
其實我們是把上帝變成了一個那軀之神.

$^{681}$God of the Gap.
我們所信的神其實跟我們的生活無關.
只是我們間中的交往對象.
他還說了一個很該死的類比.
偶爾做靈命操練的人.
跟其他異教人士飛到日本參禪.
又或者短期出家.
其實有什麼分別呢?.
沒有啊.
這是一個跟我們生命沒有關係.
虛假的信仰.
梁家倫牧師說.
作為一個總結.
今天我定了這個講道題目是無所不在的基督.
聖經對於生活忙碌的教導好像很兩極.
一方面.
忙得太多.
我們要回去找回安息.
以免我們掉進一個徒勞虛空的生命當中.
另一方面.
聖經又提醒我們.
因勤不可懶惰.
西方中世紀的諺語.
有一句這樣說的.
清閒是魔鬼的枕頭.
不要太得閒.
又太忙又不行.
沒錯.
從創造的秩序來看.
神在六天創造之後.
就在第七天.
揭了他一切的功.
安息了.
做完之後.
安息了.
那就是要安息了.
不過請你留意.
是做完才安息.
工作是要先完成後安息.
第七天的創造是一個高峰.

$^{721}$因為神看他自己所做的一切.
非常好.
所以停下來慶祝.
聖經對於神的描述.
是一個工作的神.
從來沒有說是一個休息的神.
所以儘管耶穌和門徒有計劃想要休息.
但結果還是被群眾打亂了安排.
意料之外.
雖然私人空間又再一次被人侵佔了.
但是耶穌仍然以憐憫的心腸.
來回應他們的需要.
調整自己的計劃.
去教導幫助他們.
最後我想用一個故事.
來結束今天的訊息.
故事是這樣的.
有一個外國生活的年輕人.
其實他和我們在香港的生活很相似.
工作很忙.
時間很長.
一早起床上班的時候.
也是剛剛天亮.
不到多久.
到了下班.
他開車回家.
基本上看不到太陽.
已經黑了.
所以他每天下班的時候.
其實開車回家也很累.
有一天.
當他開車回家.
突然間看到一個黑影.
在他左手邊閃出來.
他嚇了一跳.
他以為有人在行人路衝了出來.
其實本能反應就是踩煞車.
煞車吧.
當他看清楚一點.
沒有撞到人.

$^{761}$輸了一口氣.
但是他明明聽到剛才.
啪一聲.
很大聲的.
什麼來的呢.
於是他就把車子退後.
下車看一看.
看到有塊磚頭在地上.
他的車身很明顯被磚頭擲凹了.
他的心情很快就由輸一口氣.
火都來了.
有沒有搞錯.
為什麼會有人把磚頭擲出來呢.
他看到有一個小朋友.
在路邊不知道在做什麼.
就認定這個人一定是擲磚頭的小朋友.
於是他就嚇著這個小朋友.
他就走過去罵他.
你在搞什麼.
你這樣把磚頭擲出來.
其實是很危險的.
你現在有空嗎.
還把我的車子擲凹了.
這個小朋友很害怕.
就哭著和這個年輕的司機道歉.
他說先生先生請你原諒我.
對不起.
其實我不知道應該怎樣幫我哥哥.
因為路太斜.
我就負責推輪椅幫他下斜路.
太斜所以輪椅失控打翻了.
哥哥就跌到頭破血流.
他就沒有足夠的力度把哥哥扶起來.
旁邊經過的車子沒有人理他.
於是他心生一計就把磚頭擲出去.
這個年輕人就停下來.
這個年輕人聽完這個故事之後.
他不知道怎樣去安慰這個小朋友.
但是他馬上的怒火就沒有了.
他知道他可以做的就是打電話報警.

$^{801}$到他見到救護人員送走了這兩兄弟之後.
他就繼續拖著他疲乏的身軀駕車回家.
在路途當中.
他仍然想著剛才發生的事情.
他決定了一件事.
就是到往後他不會再弄他的車子.
這個凹位是很明顯的.
被磚頭弄凹了.
不過他不再弄.
為什麼呢.
因為他要記住這個凹位.
被磚頭擲出去的凹位.
是提醒他的.
無論多忙.
上帝都在.
我們除了可以安靜的時候.
是退休是好的.
對的.
向神傾心吐意.
和神傾下計.
其實我們都需要將我們每一日忙碌的生活.
成為我們靈性活動的一部份.
我們不要只顧著向前衝.
直至有塊磚頭擲在你面前.
不用這樣.
我們都可以醒一醒.
我們知道.
我們的上帝是一個無所不在.
無處不在的基督.
在每一個場景.
每一個時間.
祂都在.
這個調教.
不單止幫助我們在日常生活當中.
更加多經歷神.
其實我們身邊很多未信耶穌的人.
要等你帶他回教會.
很久才一次.
是不是.
或者他會推你.

$^{841}$他不肯來.
你請他十次一次.
給他面子.
不過他每一日都在你身邊.
所以你每一日都在傳福音.
每一日都在見證上帝.
告訴他.
這個上帝在你生命當中.
有多真實.
所以很期盼.
這個不單止成為我們生命當中的幫助.
這個也能成為福音的出口.
\newpage



\section{以西結書 36:26-27}
\label{sec:031964}
\textbf{全人更新}~~~~林羅麗斯牧師
\newline
\newline
連結: \href{http://www.church.com.hk/acms/content.asp?site=cdc\&op=show\&type=product\&code=031964\&layout=sermon}{\texttt{031964}} ~~~~ 語音日期: 2026-04-12
\newline
\newline
\hyperref[sec:031963]{\small{< < < PREV SERMON < < <}}
~
\hyperref[sec:index_chronic]{\small{[返順時目]}}
~
\hyperref[sec:031965]{\small{> > > NEXT SERMON > > >}}
\newline
\newline
以西結書 36:26-27
\newline
\begin{longtable}{cl}
\hline
\hline
章節 & 經文 (和合本修訂版)\\
\hline
36:26 & \begin{tabularx}{0.7\textwidth}{X} 我也要賜給你們一顆新心，將新靈放在你們裡面，又從你們的肉體中除掉石心，賜給你們肉心。 \end{tabularx} \\ \\ \relax
36:27 & \begin{tabularx}{0.7\textwidth}{X} 我必將我的靈放在你們裡面，使你們順從我的律例，謹守遵行我的典章。 \end{tabularx} \\ \\
[1ex]
\hline
\hline
\end{longtable}
$^{1}$在開始講道之前我們低頭有個禱告好嗎?.
天父我們再一次在您的面前去頌讚您.
再多謝您將耶穌基督成為人.
活在地上走了一段日子.
再多謝您明白我們.
亦了解我們.
您有無限的溫柔.
在我們身邊幫助我們.
讓我們緊緊跟隨您的時候.
我們的人生不再一樣.
願您真摯的將聖靈.
在我們心裡幫助我們.
讓我們有一個柔軟的心.
此刻的,願意的.
您都可以跟從主.
聽我們禱告.
奉教主耶穌基督的名.
Amen.
今天我們講題叫做全人更生.
因為已經進入了26年第二個季度.
其實更生的題目常常在教會都會出現.
都會配搭不同的經文.
特別是羅馬書是很多時候都會用的.
在去年一直去預備.
今年是市民人生的時候.
見到整條Line.
上帝很清晰地給了這兩節的經文.
我們在第二季那裡成為主題的經文.
這段經文亦有機會聽到.
有機會聽過.
因為將錫心變成玉心.
都會是常常講的講壇的訊息.
不過整卷以西結書.
就普遍比較少講的.
因為都不是常常主要學熱門的課題.
但是因為凌空在以西結書裡面抽了出來.
做我們的分題.
我都想花少少的時間.
讓大家去理解.
為什麼會有這段經文的出現.

$^{41}$以西結書是先知書的其中一卷.
我都講過是少講的.
特別兒童教會就算講故事.
都是沒有什麼特別的.
是透過以西結書講.
原因就是以西結書本身不是一個好聽的訊息.
是一個警告的訊息.
當時的背景是以色列人拜偶像.
尼契耶和華.
而先知的訊息是一個宣佈.
他們會接受審判和亡國的消息.
所以不是一個好消息.
甚至他們一直很引以為傲的聖殿.
所羅門所建的很輝煌的聖殿.
竟然會被拆毀.
這個在以色列人來說.
是想都沒有想過.
但是聽到先知的警告.
都是一個很大的恐慌和灰暗.
但是神的話語是不會用這個做結尾.
往往上帝和人說話.
背後都是建立的.
我們有很多人在做父母的.
小朋友太頑皮的時候.
是一定要管教的.
可能都會有懲罰的.
但是所有這些做法的背後.
都是為了愛.
所以最尾的以色列書.
都是一個充滿著盼望的書.
因為他做了一次宣告.
神是一個守施慈愛的神.
有朝一日.
沒錯他們是承受了犯罪的結果.
他們會亡國.
但是上帝仍然會招聚他們一起.
復興是會臨到.
拆毀的聖殿是會拆毀了.
但是聖殿是可以重建的.
最重要的.

$^{81}$他是將一個很大的安慰的信息.
就是無論在什麼狀態裡面.
上帝的榮耀.
都會一直與他的子民同在.
因此其實以色列書.
我自己會感覺.
其實都很適合我們現在.
做基督徒去看的.
因為疫情之後.
面對的世界.
截然不同.
已經不是我們.
希望可以回到以前的狀態.
是不能回到以前的.
而且某程度來說.
世界是越來越瘋狂.
無論在很多經濟的壓力.
或者是我們對於前面.
是沒有把握的.
讀的專業.
很快很快.
越專業那些.
反而很快就沒事做.
所以某程度來說.
以色列書的背景是差不多的.
大地是一片灰暗的.
是絕望的.
是好像找不到一條出路的.
但是你能不能夠.
緊緊的去相信.
你的上帝仍然坐著為王.
坐在高高寶座之上.
這是很重要的.
兒童教會有一首歌.
叫做《雲上太陽》.
我鼓勵你不懂唱應該學一下.
是的.
我們現在面對的生活.
是有很大的挑戰.
有很大的難處.

$^{121}$以往你用很多時間.
去讀的書.
培養的專業.
竟然原來用不了半世.
很坦白的說.
然後整個世界就已經是.
你都好像沒有讀過書一樣.
好像找不到工作.
很多事情.
我們以為有把握的事情.
都沒有再把握.
在這樣的狀態裡.
究竟你是在埋怨上帝.
還是在質疑世界上有沒有神.
還是你真的相信.
行雷暴雨的裡面.
背後的太陽仍然是高掛的.
這一世的黑暗絕望.
是一定會成為過去的.
在正經書裡.
更加絕望的就是.
在以色列裡.
沒有人可以用一樣.
但是在先知書裡提醒我們.
我們的神是一位創造主.
不要將人類看得太大.
以為好像我們讀專一樣.
但是上帝是仍然創造的.
四活物究竟是什麼呢.
其實現在都沒有人知道的.
但是當人不能被上帝使用的時候.
神會創造另外一些東西.
來驅使自己使用.
同樣在聖經裡說的.
如果你不開口讚美主.
神可以使石頭高歌讚美.
因此到三十六章的時候.
神要將石心變成肉心.
接著就是神要將枯骨復活.
枯骨是什麼呢.

$^{161}$是死到剩下的棚骨.
而那棚骨仍然是枯的.
是扶不起的.
拔起就會散的.
但是在這麼死到絕的狀態裡.
神卻帶出一個很大的安慰.
就是讓以色列人知道.
他仍然是神.
生命在上帝的裡面.
創造力在上帝的裡面.
就算枯骨.
他仍然能夠復活.
復活不是做一個普通人那麼簡單.
他還能夠成為大君.
為上帝所用.
所以結局其實不是灰慘的.
是上帝應許他會再來到.
憐憫他所揀選的人類.
也醫治和平安.
會與他的子民永遠同在.
而以色列人的榮耀將會重現.
所以在這樣的氛圍裡.
這兩句話就是正正.
上帝一手要審判以色列人.
但是卻又同時告訴他們知道.
就算你們被審判.
被絕望的當中.
但你仍然是有希望.
所以是一個很弔詭的兩節經文.
我們去面對著上帝的帶領.
一直去禱告每年的主題.
然後我們在分題裡面.
慢慢慢慢去學習.
在第一季度裡面.
我們都知道.
當你要想你自己使命的人生.
你想更加愛主.
你想你活得更加有價值.
你需要的是明白上帝的旨意.
因為我們往往為主瘋狂.

$^{201}$或是為神瘋狂.
我們一線之間都會有混亂.
當我們明白上帝有這樣的心意.
我們就會面對下一個問題.
就是我們到底是否願意跟進.
有時候知歸知.
但要自己去參與的時候.
其實有無限的拉扯.
在我們一群很愛主的弟兄姊妹心裡.
因此我很明白.
上帝在第二季度裡面.
想我們去討論一個更新的問題.
如何能夠更新自己.
以致我們能夠開開心心.
去跟從主的心意.
展開我們真正合乎上帝心意的.
所謂使命的人生呢.
在這兩節節目裡.
我特意調了一些你和我.
好像神親自和我們去講說話一樣.
所以我想大家用這樣的字眼.
讀一次,1,2,3.
神也要賜給我們一個新心.
將心靈放在我們裡面.
又從我們的肉體中除掉錫心.
賜給我們肉心.
神必將神的靈放在我們裡面.
使我們信服神的律例.
謹守遵行神的典像.
這是上帝今天要跟我們講的說話.
也要將這個英渠.
不單賜給我們個人.
也會賜給我們的小組和教會.
上帝真是很愛我們.
和我們說上帝很愛你.
其實我們已經在教會走了三十年.
本來應該是一個成年的時間.
和現在的青年一代一樣.
三十歲他們就說自己很老.
所以我們教會三十歲都應該很老.

$^{241}$其實在肉身上我們大部分很愛主.
奔波了三十年的頂智慧.
我們都已經是退休年期.
但在這個時間上帝反而叫我們再想.
你們的人前面還有路.
你們應該怎樣走呢.
如果長命百二歲.
哇,你還有六十年.
你應該怎樣去想呢.
我們是否真的接受自己.
就是在現在這個狀況.
直到永遠厭主免呢.
這是一個新的挑戰.
為什麼要更新呢.
有時候我們都會像阿A這樣想.
其實那些東西多好.
是不需要更新的.
不好的東西就換掉.
每樣東西都覺得OK的時候.
為什麼還要想這個情況.
我繼續持續這樣下去就可以了.
特別是在精密教會.
一直由第一次來到現在都還在.
我們覺得我們終於是很老了.
我們的二代都覺得我們已經很努力了.
沒有青年人是這樣的.
又回到小組,又回到崇拜.
然後他們又學習了奉獻.
每樣東西都做齊了.
他們覺得自己已經是挺好的了.
又說更新.
沒什麼更新的了.
但我不知道.
我們如果是這個年紀.
七十年代在香港拼搏的那班人.
我們當時的氛圍不是這樣的.
整個香港是好的.
但是可以更好的.
大家還有這個情懷嗎.
跟隔壁的人說.

$^{281}$可以是更好的.
香港已經沒有這種感覺很久了.
我記得小時候我們住武區.
武區很開心.
就是可以上樓.
住到一些高樓大廈.
但是那時候在屋村長大的小朋友.
每個人都有一個夢想.
就是我要衝出屋村.
我不做屋村妹.
每個人都很努力.
不受政府資助.
我們要努力自己賺錢.
我要住私家樓.
喜歡恥女.
很多時候不能讀書的女孩子.
很早出來工作.
當有夜校一出現的時候.
你不要管她.
回家十二個小時也好.
十個小時也好.
然後就拿著飯壺.
就去讀夜校.
心裡面是滿有希望.
滿有盼望的.
總之有人我讀到書.
我畢業了.
我就不用再做工廠妹.
我不用做藍領.
我要做白領.
但是現在這個世代.
從來都沒有看到這些.
只會說我想不做事.
已經不是這個世代了.
因此我在想.
我們教會有多好.
其實真的挺好的.
頂尖的外主.
我們很多東西已經很好了.
為什麼神要我們更新.

$^{321}$我們更新些什麼呢.
今天在觀念上.
第一課.
想大家有些調整.
第一.
是不是每樣東西.
都可以更新.
我們很怕神跟我們說更新.
但是你想清楚.
是不是每樣東西.
都可以更新呢.
原來不是.
剛剛過了一年一段時間.
大部分的花.
一個月都捱不了.
已經扔掉了.
我每年都有水仙.
水仙不能留太久.
就是說.
有些東西.
更新不了.
它已經死了.
沒有一個可以再種的機會.
但是.
有一樣東西.
我們叫蘭花.
人家送給蘭花.
通常就這樣放著.
通常我們都不知道.
原來蘭花是可以.
繼續生下去的.
我很少插這些花.
很多年前.
在監獄大一點的監獄.
人家送了一盤給我們.
因為家裡大.
那些花那麼重.
當然不想扔掉.
放在騎樓.
騎樓看到它在那裡.

$^{361}$那盤已經死了.
沒有事.
又在下衣.
隨便吧.
潑水給它.
不知道為什麼.
我突然發現.
原來它會出花lum .
原來它是還沒死的.
原來.
你一式上傳.
它竟然是可以.
再一次出花lum .
原來一盤蘭花.
那盤蘭花.
直到我走的時候.
它太重.
我沒有帶它.
它仍然可以令你有驚喜.
仍然可以有花lum .
當這樣去想的時候.
弟兄姊妹.
你有沒有很珍惜.
今天上帝跟你說.
叫你更生.
跟你說.
原來你還沒死.
我們調一調思維去想.
我們經常說.
上帝你要我更生.
要我改變.
原來今天上帝叫你.
因為你還沒死.
你還有一絲希望.
如果你已經死了.
我相信你聽不到上帝跟你說話.
上帝都不會理你.
都不會叫你傳福音.
上帝不會再叫你彼此相愛.
因為你已經死了.

$^{401}$所以我們第一件事.
去預備這一季學習的時候.
我們要在神面前.
真是有感恩的心.
原來今天上帝還挑戰我們.
還叫我們努力.
不是我們個個不夠努力.
但原來我們還可以改變.
我們還可以更生.
頂智慧我們都常常說.
在香港.
尖墨都一天到晚叫人去傳福音.
為什麼會這樣呢.
但當你真的方言去看的時候.
你真的要好好珍惜.
我們還可以傳福音.
我們剛剛去完KL.
兩個星期來回了兩次.
很坦白.
在KL裡面.
其實教會都可以有自己的自由.
但絕對不可以有街頭報道.
你絕對更加不可以向回教徒.
去做傳福音的工作.
我們現在看到很多戰士.
很多難民.
但在崇靈的角度.
我們看不到上帝的心意.
我們想都沒有想過.
我去到KL的時候.
發現他們原來有很多難民.
當然他們油價和我們一樣.
都很高.
但他們多了很多難民.
很多難民去到的時候.
自然就要找自願團體去托管照顧.
你知道基督教也是很喜歡小朋友的地方.
都會給很多小朋友在當中來幫忙.
他們不可以直接傳福音.
但來到吃飯.

$^{441}$會借飯.
他們會和他們唱聖詩.
又會講自己的故事.
因為講自己的故事.
那個師妹在後台幫忙做耳機.
我們聊天她才說.
她說她沒有想過.
因為他們是難民.
他們不是馬來西亞本地的居民.
所以和他們分享耶穌.
不算犯法.
因為他們是護照的人.
她說在她自己服務的機構.
七成的小朋友.
在他們當中親眼看到耶穌.
我們要將名歸給主.
我們聽見都呆呆.
她說真的.
她走過來和你說.
姐姐 耶穌剛才在你身邊.
祂拍你肩膀.
一開始以為是眼花.
但很多小朋友都見到.
七成的小朋友都見到.
所以她說我們都彼此真是代禱.
亦不想太誇張.
因為她怕太多有問題.
真的有驚動.
但我想告訴弟兄姊妹.
原來上帝今天仍然在做事.
就算我們每個人不做事.
上帝都親自去拯救靈魂.
因為這本身是上帝唯一的心意.
因此當上帝呼召我們.
有很多事要想辦法.
我們做過很多事.
有效果.
但又沒有效果.
是否方法要更新呢.
還是我們覺得很累了.

$^{481}$又學一件事.
然後又學另一件事.
很累了.
這可能是我們的想法.
不過我鼓勵你今天重新去想.
如果上帝還給你新的東西去學.
上帝還給你創意去服侍的時候.
珍惜這個機會.
因為我們仍然有這個機會去事奉主.
第二個很重要的觀念.
在經文裡寫得很清楚.
無論你怎樣去看.
橫看怎麼看.
看原文去看都好.
這兩節的經文.
單單去看.
其實人是沒有事做的.
我在想.
在應用程式裡和我老公聊了整段經文.
我們是沒有事做的.
因為你想想.
上帝在更新的工作裡.
全部都是他做的.
人已經是石心.
沒有心跳的.
死了.
然後那棚是枯骨.
想扶起來都碎的.
怎樣能夠興起變成大軍呢.
經文所有東西都是上帝自己做的.
上帝是為他們做一個新的心.
然後將新的靈放在他們裡面.
然後要將這個石的心都換了.
要換一個肉的心.
後面那兩句話才令我們震驚.
我們真的不要驕傲.
原來我們沒有事可以做的.
他說我要將我的靈放在你們裡面.
使你們信服我的律例.
遵行我的典章.

$^{521}$我們常常教弟兄姊妹.
教小朋友.
你要聽上帝的話.
要跟著上帝的話去行.
好像我們自己很厲害.
我們可以.
我們努力.
我們拼除萬難.
我們體貼肉體.
我們要事奉我們主.
我們要跟聖經說.
剛才中堂說我們要行十戒.
我們以為我們做了很多事.
原來經文的揭示說.
我們能夠信服上帝的律例典章.
能夠謹守遵循上帝的話.
都是聖靈的工作.
今天是朕禮的日子.
我們大家都很清晰.
如果不是聖靈的感動.
沒有人可以口稱耶穌基督為主.
意思就是說.
當我們由信主去願意.
說我信耶穌那一刻.
其實已經是上帝的工作.
神的靈充滿我們.
幫助我們.
以至我們能夠爭奪魔鬼一切的爪牙.
我們能夠有自由說.
主是我們願意成為你的救主.
然後我們還要求我們去學習.
我們以為我們靠我們自己.
讀書厲害.
看得明白.
能夠聽得清楚.
還是專心.
原來不是.
我們能夠聽得清楚.
我們能夠登一聲.
有頓悟.

$^{561}$然後我們去實踐那個決心.
這一切.
都是神的靈在背後.
推動著我們 幫助著我們.
當我們這樣去想的時候.
我們知道這一切全是恩典.
今天你和我還能夠回到教會.
不要再去瘋瘋癲癲.
不要再想怎麼辦了.
有些人說不要關我事.
膚致不要緊.
隔壁那個吧.
如果你要更新就更新隔壁那個吧.
為什麼會這麼怕呢.
小心一件事.
什麼事情才會沒有興趣.
什麼事情才會沒有反應.
什麼事情才會覺得無所謂呢.
我們留意一下什麼叫錫心.
錫心是說一個對上帝剛硬冷漠.
抵擋真理.
最恐怖的這一句.
甚至對上帝的愛與憐憫.
都沒有感覺.
那你就知道什麼叫死寂寂了.
當一個人.
當一個基督徒.
你記住.
這個說話不是和非基督徒說的.
這段經文是和以色列民說的.
是和信主的人說的.
原來曾經一個很愛主.
很熱心的人.
覺得上帝很棒的.
但是有一天.
如果他不為意的時候.
他的心是會冷的.
冷到一個地步.
他變成了錫心.
他對上帝是剛硬的.

$^{601}$對上帝說的話是冷漠的.
就算明知是真理.
他們都會抵擋的.
最恐怖.
你想擁抱他.
愛他.
都沒有用的.
因為是沒有感覺.
我相信我們生活了這麼多年.
我們都有好朋友.
我們都有親戚.
我們有愛的小朋友.
你都知道.
當他們發脾氣的時候.
你想去愛他.
擁抱他.
都不行的.
他會推開你的.
就是這個狀態.
表面上錫心有一個好處.
為什麼人會用錫心去對待所有的事情.
因為我們想不受傷害.
小朋友很生氣.
很不開心.
他都不要你的.
他不想放軟錫心.
接受你的愛.
因為接受你的愛.
就代表我不對.
媽媽你剛才罵我.
罵得對.
我不想.
所以怎樣都要裝硬.
然後什麼都頂得住.
不關我的事.
我沒有問題.
我沒有所謂.
你怎樣說都沒有所謂.
我不會受傷.
好像在那個狀態裡面.

$^{641}$我們還有一些理直氣壯的能量.
但是有沒有想過.
表面上一個錫心是不容易受傷.
因為表面上是硬的.
不是那麼容易受傷.
但是一個錫心才會不堪一擊.
當人在那個狀態裡面.
一敲.
整個就會碎.
所以你看到.
很多時候.
姐妹之間.
如果在一個狀態.
他很硬.
他不會聽你的勸.
他不會聽你的安慰.
就算你道歉.
他也不會接受你的時候.
你就知道.
他就是一個錫心的狀態.
沒有什麼可以改變他.
但是其實他不是堅硬.
他裡面是非常非常脆弱.
因此我相信.
每一個信徒.
以色列民不是一朝一夕.
馬上就變心.
不是柔柔軟軟.
馬上就變成錫心.
我們要小心留意自己的心.
因為聖經說.
我們一生的果效.
是由心發出的.
大家有爭執.
有不開心.
是很生氣的.
究竟你是讓自己拿到實.
和他死過.
談都不談.
我永遠都不會接受你解釋.

$^{681}$我永遠都不會接受你道歉.
還是我仍然冒險地.
幫助自己放軟心呢.
以致我們還可以談.
我們還可以和解.
我們還可以尋求方法呢.
其實人的關係.
如果有不同的態度.
就會有不同的效果出現.
但我明白的.
經常保持自己的肉心.
是柔軟的.
是有感情的.
但我會覺得很傷.
我們經常說沒有玻璃心.
玻璃心又不是肉心.
玻璃表面是硬的.
一敲就碎了.
現在說的不是玻璃心.
是因為一個真正有生命力的心.
是一個肉體.
真的像我們皮膚一樣.
一碰就會流血.
是會化的.
但只要我們相信.
我們有這個生命力.
這個身體仍然活著.
上帝的醫治聯盟仍然存在.
我們是會有痊癒的機會.
是一個選擇.
究竟你要讓你的心剛硬.
還是讓自己能夠放軟呢.
你的生命力是完全不同的.
有很多的爭執.
大家都能夠放慢.
大家都能夠放軟.
慢慢再傾過再游過.
可能事情帶出來的結果.
就不再一樣.
小心保守自己的心.

$^{721}$不要以為在那個堅強.
什麼都不能改變我.
我不會再理會.
我不會讓自己受傷害.
這不一定是一件好事.
在這個思想裡.
我們就會想.
上帝如果你真的呼召我們.
去行一個使命的人生.
是合你的心意.
又讓我們真的領受.
你為我們預備在人生裡.
一切的福氣的時候.
我們除了要珍惜自己.
能夠有改變的機會之外.
我們除了去相信.
你是那位能夠幫助我們.
改變我們的神之外.
我們要做些什麼呢.
華人特別是很壞的.
我們總是覺得.
沒有事情做.
好像不知怎樣去達成.
我們都要去做些什麼.
我在這裡很安靜.
因為我在想這篇講章都很久了.
因為都很難講.
很難講也很難落地.
我都問神.
其實一樣東西都不是我們做.
我們不用做事.
是不是很不用做.
第二就是說.
如果說要做.
那做多少呢.
要讀經讀多少呢.
要靈修靈修多久呢.
如果要上課.
是不是上完基督生平就不用上課了.
突然間又有保羅.

$^{761}$保羅之後又有摩西.
接下來你們說要學希伯來書.
有好事騙人的.
是不是每次都要上課呢.
所以都沒有很清晰去處理.
但很明顯.
我知道不是這些.
不是你做了一二三四.
你就會更新.
不是做了一二三四.
你就會柔軟的心.
然後一二三四你又不會受傷害.
然後你就永遠都幸福快樂下去.
不是在我們這種所謂的思維裡面.
因為我覺得.
上帝是我們的上帝.
祂是我們的創造主.
上帝的創意和貼心.
我想是沒有人能及的.
祂做每一個動物.
同樣都是人類的.
有眼耳口鼻.
但你看到旁邊那個.
就算是兄弟姐妹.
雙胞胎也好.
都不完全一樣.
同一棵樹生出來的樹葉.
你都不可以搭在一起.
用剪刀剪得齊齊整整.
一樣大小的.
所以上帝的特性是.
很貼心的.
簡單地說.
你就是你.
上帝為你預備了一條路.
是要你和上帝去經歷的.
今天可能我們面對了.
有不同的際遇.
有不同的幸福的地方.
或者不幸的地方.

$^{801}$但總是不同的人.
有不同的組合.
我們寫出來的國章.
就不再一樣.
所以在這麼吊詭的裡面.
我只可以覺得.
我們需要祈禱.
但我不用祈禱的字.
因為我覺得不只是祈禱.
我選了在經歷神會上用的字.
就是「舉手向主呼求」.
我們說一次好不好.
「舉手向主呼求」.
在經歷神會上就知道了.
如果你真的很想很想.
去經歷主的時候.
我們要真的打開我們的心.
打開我們的手.
死都向著上帝.
屌到有為止的那種心態.
我要經歷到你.
經歷不到你.
我就不走了.
那種心態.
今天同樣的.
在說使命的人生.
你是和上帝一起走的.
不是隨便.
看一看.
隨便吧.
因為上帝是萬變的神.
祂是無限的創意.
不是上帝沒有rule.
而是上帝愛你.
所有的計劃.
是為你tailor-made的.
即是每一分每一秒.
你的選擇.
都是影響了上帝的計劃.
你有沒有看到這個微妙.

$^{841}$這樣才是神.
無論你是什麼款式.
無論你是什麼性格.
神都可以為你tailor-made的.
這就是一個有能力.
而且愛你的神.
才可以做得到的.
爸爸媽媽都很愛我們的子女.
我們都很想和他們tailor-made的.
但有時候真的頂不住.
我們沒有這個能力.
沒有這個胸襟.
真的完全按你的子女的心意.
陪著他們走.
又能夠帶他們回正軌.
真的沒有這個能力與智慧.
但是.
我們的上帝.
就是那位萬有的神.
祂創造我們.
世界上所有的東西都是祂的.
所以無論你.
每一個moment的決定.
無論你走錯路.
走歪路.
上帝都與你同行.
只要你always與祂同在.
一些很差的日子.
很哭到媽媽都不認得的日子.
但是當你跟著祂走的時候.
回頭你會多謝那段時間.
因為那段時間是令你成長.
那段時間是令你真正的堅強.
所以弟兄姊妹.
我們要有一個心智.
向主呼求.
如果你不為自己的生命呼求的時候.
上帝是一位.
等你情願的神.
所以你要明白.

$^{881}$為什麼大家都回到神明會.
大家都回到小組.
有些生命就很快地向上飆.
以前有什麼問題.
祂很快就講見證.
祂又怎樣過渡.
祂又知道自己怎樣衰化.
祂又認罪.
又離過.
總之大家都走十年的癡迷路.
但是你發覺.
十年的時候.
大家的分野是可以很遠的.
我們怎樣來證實更新的問題呢.
如果更新不是上帝有空沒事做.
而是上帝想我們變得更好.
如果更新不是上帝找你幫忙.
要你很辛苦.
而是上帝愛你.
祂知道你還可以救.
祂還想幫你.
當這樣去想的時候.
我們真的能否去領受.
上帝對我們這一份的好意.
以致我們真的去渴求.
祂將我們知道.
祂為我們預備前面的路.
可能我們想都沒有想過.
但是卻一定是蒙福的.
因為祂帶領我們的路.
是一定我們自己都會喜歡的.
所以弟兄姊妹.
在禱告的時候.
我想你今天學習一件事.
正式的你也跟神說.
既然上帝要設新的心.
新的靈給我們.
我們就按這個方法去禱告.
讓神給聖靈去掌管我們的心.
心這麼重要.

$^{921}$是影響我們一生的果效.
我們有沒有試圖將自己的心思念.
交給神去掌管呢.
以致我們不要變成石心.
石心不是一朝一夕就變成石心.
這個石心影響了我們對上帝的反應.
如果你在石心的狀態是很難改變的.
除非神即其事.
好像二世傑書一樣.
但是在我們正常的生活裡面.
我們要小心.
每一個在教會的弟兄姊妹.
你一定不是石心的.
因為你還想回教會.
你還想親近主.
在小祖雖雖地仍然坐在那裡.
已經是很難得了.
所以給在聖所的組員.
你這麼悶.
你坐在那裡.
我都覺得不得了.
但是為什麼會這樣.
你要留心.
原來我們一下子不小心.
不為意的事.
我們的心是會漸漸冷淡的.
如果我們手指尾有條刺刺在那裡.
我們不理它.
算了.
有鬼只有它.
接著它就會發炎.
它就會含糊.
一個人由一個肉心.
變成了石心.
是一定有無數小的事情.
來影響.
來變成的.
以致今天回教會的弟兄姊妹.
為什麼對上帝的呼召是這麼冷淡呢.
因為已經聽過了.

$^{961}$短孫也去過了.
基督的神明會.
已經說過很多次了.
又是說那些東西.
接著放星祈禱會.
還有什麼呢.
什麼都見過了.
躺著也見過.
滾地也見過.
什麼都見過.
還有什麼呢.
好像我們真的.
還沒死光.
但是回來.
坐著看看.
看看隔壁那些.
我們還好笑.
上一次有一個.
有一個姐妹的禱告.
很搞笑的.
她很想經歷.
但又不想自己跌.
她有一個這樣的祈禱.
她說主人我又想見一下.
但我又不想再經歷.
我已經跌過了.
最好你給我早掌跌.
真是陰毒.
我想說.
接下來這個月.
不要搞這些事情.
但你不要來.
上帝真的聽了我的祈禱.
嚇到他瘋了.
才說出來.
真的他的掌.
砰一聲就躺在他前面.
嚇死他.
不是因為他跌.
沒事的.

$^{1001}$是他突然覺得.
隨便說一下.
上帝這麼靈.
令他整個人都敬畏了主.
弟兄姊妹.
是的.
是有一些事情.
我不是覺得自己很差.
但是有些事情.
當上帝再次想叫我.
想和他挑戰我的時候.
我不知道為什麼.
我有些事情不是很想.
不是很願意.
覺得這個人很卡.
弟兄姊妹.
留心這個狀態.
因為這個狀態.
是代表我們有些事情.
我們以為沒事的.
但是偏偏.
他暗地裡.
就在影響我們.
究竟錫心可以是什麼.
什麼令我們變成錫心的呢.
我都問上帝.
有沒有見證可以分享.
超級頂字都沒有.
但是在回程的飛機裡.
突然上帝令我想起了另一件事.
這件事我放下很久.
差點忘記了.
是剛剛上去美景花園.
我安納了牧師母來.
可能是2005-2007之間的階段.
關於卡片的問題.
關於我的稱號的問題.
其實本來就武士武幹.
我一天到晚在兒童教會.
叫林師母.

$^{1041}$大家都武士武幹.
無端端就上來大堂.
做主任牧師.
接著又要見人.
所以始終要做卡片.
很正常的.
我們這些傳統美德的中國婦女.
貫膚性是正常的.
難道你還沒嫁嗎.
所以證人會寫林羅麗斯牧師.
但是就吵架了.
因為大部分自己有就叫師母就算了.
外面的人又不知是叫你林師母還是羅師母.
總之大家都明白.
崇拜事事都讀錯了.
只能笑一笑.
我都沒有理會.
但我覺得堅持是正常的.
這個稱呼是正常的.
直到我回大陸.
其實我第一次遇到挑戰.
很多年前.
有些人有興趣問我.
不好意思.
我也想問.
這麼有趣的.
香港有福姓林羅的嗎.
我真的不知道什麼叫福姓.
想不起這件事.
歐陽那些.
不是不是.
我不是姓歐陽的那些林羅.
我是姓羅的.
我老公是姓林的.
所以我就叫林羅麗斯牧師.
哦.
我們不是這樣的.
我們這裡不會灌呼姓的.
中國不是灌呼姓的嗎.
我們完全沒有.

$^{1081}$這裡沒有人會灌呼姓的.
我們所有教會那些.
你是什麼名字就什麼名字.
這個就是范牧師.
她老公也是范牧師.
他們兩個姓范.
我們不會叫范范牧師.
這倒是.
突然間我的腦袋覺得.
很不同文化.
其實幾年了.
我聽了就算了.
但我們香港嘛.
也正常.
台灣和我們一樣.
朱詩未必非牧師.
我們經常這樣叫.
我就沒有特別的留意.
但是崇拜就忍耐.
主席讀到亂七八糟就算了.
自己有無所謂.
牧師有時候都說我們.
我老公說你們直播的.
到處聽著.
改個名字都覺得不正.
我不理他.
但是到了我去弘恩.
開始真的要和外面的人.
去接觸.
派咭片.
他們又真的讀錯了.
個個走過來.
大學教授來的.
不好意思.
都說林牧師.
還是羅師母.
亂七八糟.
我說不用了.
跟他們叫就叫師母就行了.
沒所謂的.

$^{1121}$直到剛剛沒多久.
他們二月份.
參加他們週年校慶.
因為代表教會.
領土的.
都穿著博士袍.
出來做主席的.
都讀錯了我的名字.
你明不明白.
請林牧師.
接著下午有人說.
請羅師母.
又是這樣.
亂七八糟.
接著就拿張紙出來.
讀一次就回去.
接著我心裡就想.
算了.
是不是要時代已經不同了.
現在都不流行複姓.
那就不複了.
因為我們本身.
不是把身份證轉了名.
因為以前.
匯豐工作的人都知道.
結婚之後.
身份證上要轉了自己的名.
如果英文來說.
我不可以保留Jerima Low.
我一定會變成Jerima Lam.
所以換了名之後.
我真是一個林牧師.
因為我是林姓.
不是羅姓.
但香港又不是這樣.
很吵.
但神為何提起這件事.
我就想.
為何要為了卡片.
原來我都擾攘了很多年.

$^{1161}$為何不下定決心.
做什麼就做什麼.
叫師母就全部叫師母.
叫牧師就全部叫牧師.
當我這樣想的時候.
在飛機裡.
突然閃了一下.
我想起我在200幾的階段.
一件很小的事.
其實是很小的事.
只是半小時就發生了.
有個指揮.
因為那時有很多新人來教會.
因為在美景很大.
數量又很大.
很多新人來了.
有個姐妹來了我們教會一段時間.
加入了小組.
但都挺霸氣的.
都覺得有點風範.
有一次她突然約我.
她叫我羅牧師.
她說我想約你.
我想見你.
我覺得可以.
我初初以為她是婚姻問題.
因為那時都做很多婚姻輔導.
坐下來見.
誰知坐下來.
她第一句說話問我.
她說我很想知道.
為何你要我們叫你林師母.
我心想.
我想不到別人問我這些問題.
突然不懂回答.
我說沒有.
因為他們一向都叫.
習慣了叫師母就叫師母.
她說但不是.
你已經按納了.

$^{1201}$你是羅麗絲牧師.
你是牧師.
你不是師母.
我說不是.
我是牧師.
我也是師母.
我老公也是牧師.
她說不是這樣說.
因為你是上帝按納的遺訓.
所以你是應該牧師.
你不是應該叫師母.
來來去去地說這些.
接著她用了很偉大的篇章.
說了半小時.
簡單來說.
我明白她的心意.
她非常不滿我.
因為你好像貶低了我們女性的身份.
明明你好像羞辱了上帝給你這麼好的遺訓.
明明你已經按納了牧師.
你和林以樂是平輩.
你做什麼還是這麼沒出息.
你還要叫自己做林師母.
她處處逼人到一個位置.
真的不知道怎麼說.
難道跟她吵架嗎.
突然她說了一句話.
真是激火了我.
我火氣上也挺令人害怕的.
很坦白.
她跟我說什麼呢.
總之我覺得這樣是不OK的.
我和我老公一起做生意.
我從來不會讓我的員工叫我老闆娘.
我是老闆.
我老公也是老闆.
為什麼要叫我老闆娘.
她這麼強硬就火起了.
我不知道.
我是衝口而出的.

$^{1241}$不是我自己有退白想的.
也都我以前沒想過.
我衝口而出去反駁她.
你知道我說什麼嗎.
我明白.
但是我喜歡冠夫姓.
我喜歡叫林羅麗絲.
是因為我想讓人知道我有老公支持我.
其實我突然發她懶的那一刻.
我都呆了一會.
但心底裡我真的想知道我有老公.
我不是還沒嫁的.
所以感覺說出來也很擔心.
她都張開嘴巴掌摑著我.
但是她說不出話.
她沒有說話.
接著她說.
我說了就說了.
你自己好好反省一下.
接著她就敲門走了.
OK.
上帝讓我想起這個片段.
真的有小火.
但這個片段不是結局.
而是我想起她敲門.
我做什麼呢.
原來我沒有離開房間.
我坐在房間裡.
原來那天我坐在房間裡哭了一會.
哭的感覺.
讓我想起那時很不開心.
不是不開心被人罵.
而是覺得很委屈.
覺得我有沒有做錯什麼.
不是中國傳統教人這樣做的嗎.
是對的.
聖經也是叫人的丈夫是頭的.
我究竟做錯了什麼.
無端端被人罵了半小時.
而我自己又不是很憤氣的那種感覺.

$^{1281}$在飛機那段時間.
我突然感受到那份很委屈.
很委屈.
有時候突然會是.
我們在事奉主的路裡.
在做教徒的路裡.
我們都會很盡心盡力去事奉.
或者我們都很努力的.
吃虧 愛人多一點.
跟從耶穌多一點.
不要那麼玻璃心.
或者不要那麼計較.
但是有時候我們會靠自己剛強.
算了.
由得它吧.
我不會跟它計較的.
我不會跟它說的.
跟你說聲對不起.
不用跟我說對不起.
不要再跟我說這件事了.
我不會再跟你談這件事了.
我們覺得我們OK.
我們不再糾纏就行了.
我們沒事的.
但是上帝真的突然間.
令我的心突然間很軟.
我想起自己那種委屈.
那種不開心.
我就明白.
原來往往我們心裡.
我們手上拿著的那些東西.
可能不是對還是錯.
那些委屈.
那些傷害.
或者那些遺憾.
在小組教會裡面.
很多頂尖姐妹做過組長.
她們離開.
她們不知為何會帶著遺憾.
因為我們覺得要成功.

$^{1321}$不可以不成功.
但是我想說.
其實從來沒有人要求我們成功.
教會也沒有要求我們成功.
我們FAST而已.
當有一天不可行的.
就另一個頂上.
我做了那麼久輪到你了.
其實我們能夠去明白.
那個核心的價值.
不用自己逞強.
但就是因為我們不留意.
有一點點的刺.
黏在手上.
我們以為可以的.
不用看的.
想不到.
慢慢慢慢.
它就在這裡.
含濃.
所以再叫我做事.
不叫不做.
你叫我不做.
我不會再做了.
我坐在小組裡面.
看著死的.
我都不說話的.
我已經覺得我最好了.
小心啊頂尖姐妹.
原來這些是錫心.
是我們一個剛硬的心.
有人愛你.
你都不會接受的.
我不用照顧的.
我OK的.
你照顧那些吧.
不用約我吃飯.
我OK的.
我沒問題的.
我和上帝的關係很好.

$^{1361}$靈修很好.
你約那些吧.
我們很多時候是這樣.
我們沒有想過.
原來這個不願意.
給上帝再一次使用.
再一次改變.
所謂剛硬的心.
正是今天上帝要我們.
重新更新和面對的.
一個很重要的位置.
頂尖姐妹.
當我們面對很多事情.
一生真的有很多遺憾.
我們曾經為主瘋狂.
小時候立志又做過媽媽.
又做過媽媽.
但好像每件事都不成的時候.
我們有時候都會覺得灰心失意.
今天已經這麼老了.
你還叫我去傳福音.
不用了.
我都預備退休了.
我們有很多給自己的想法.
上帝要去除掉這個錫心.
是實要我們願意.
如果不是.
上帝硬來就斬水.
上帝愛我們.
讓我們是情願的神.
所以上帝用無數的方法.
去鼓勵我們.
去安慰我們.
又給我們不同的機會.
去試不同的東西.
為的士都想我們能夠重新.
找回那份事奉主.
愛主的那種經驗.
今天我們怎樣去面對呢.
整節經文我們沒什麼可以做.

$^{1401}$我們只能謙卑的來到上帝的面前.
不斷的禱告.
如果我們裡面真的有傷口.
求上帝親自的醫治.
面對的時候是有痛的.
想起來是不開心的.
但是當你願意打開的時候.
上帝是有能力讓你將那條刺拿走.
你就可以重新的出發.
當我們說.
這件事我又做不成.
那件事又不做成.
那些讀寫障礙的都做不成.
其實不是的.
神做每一個人都有自己的恩賜.
都有自己的能力.
我們只不過可能是因為.
按照世人的需要去滿足.
很多時候弟兄姊妹的事奉.
就是有什麼需要我們來事奉.
今天我很想你開一個新的眼界.
不是因為有什麼需要.
應該是你問.
主人我可以怎樣去事奉你.
上帝創造你.
祂知道你適合什麼的態度.
適合什麼的光去事奉.
事奉不只是在教會裡面.
天大地大.
全地都是上帝的.
只要你願意被祂使用的時候.
上帝是可以去使用你.
耶穌基督一生裡面.
都沒有被宗教人士接納.
祂都沒有辦法在聖經裡面.
在會堂裡面做過什麼.
但是耶穌在曠野.
在山間.
祂仍然可以完成祂的使命.
弟兄姊妹.

$^{1441}$不要被那個框框來限制你.
怎樣能夠有這個突破.
一定是禱告.
個人的禱告.
我不知道你的房間還在不在.
很多時候我們教了很多方法.
弟兄姊妹做的.
可能做三個月就不做了.
我們有守望小組.
是不是只顧著吃喝玩樂.
生日快樂呢.
我們都是在屬靈裡面.
真的彼此去支持一下.
三四個人這個群組很小.
如果試了一輪.
我們都投契的時候.
躺開心扉.
我們在靈裡面.
是需要一起彼此去守望.
彼此去鞭策.
在我們的小組裡面.
每個禮拜既然都見了.
不要只問有什麼感恩事項.
真是人生我們有很多事想做.
或者我們人生不知道想做什麼都好.
弟兄姊妹為我祈禱.
我現在很空閒.
我現在退休了.
除了跳賊尾操.
我們還有什麼可以做呢.
我們不如去尋求一下上帝的心意.
正經說你要尋找的.
你就會尋見.
叩門的.
就會一定為你開門.
問題究竟我們有多少持續的能力.
你堅持就可以了.
神是切身和身心的上帝.
在這段時間.
頂智慧都知道.

$^{1481}$我很坦白的人.
我在廣台很容易.
分享我自己的經驗.
當愛麗絲木區在2020年end了.
2023年開始.
十個不情願的.
去接受男士木區.
戰戰兢兢.
所有板斧用在男士木區.
都不OK的.
請他們來我家吃飯.
令我非常憔悴.
所以完全不知道怎麼做的時候.
但我常常都提醒.
神我不管的.
我很努力的.
就是按照聖經原則.
服侍三年.
如果三年裡面不見果子的.
要不你就做下我.
要不你就做下他們.
所以現在橫斤.
我都是用這樣的心態.
我不管的.
靈人.
又說有朋友帶來.
我見不到朋友出現.
還是靈人的時候.
我和上帝說.
我不管的.
如果你說小朋友是需要的.
我們現在有義工坐在那裡.
吃冷板凳的.
沒有人在那裡.
我們的義工士工都在那裡.
沒有小朋友來.
我們都繼續預備.
你明白我的意思嗎.
我們盡我們本份.
我們守三年.

$^{1521}$但如果三年都搞不定.
上帝你就做下我.
我就做下他們.
去年2025年.
其實已經是第三年.
在男士木區.
我由很想做事.
到不是很敢做事.
但我就跟別人說.
這樣下去不行的.
一定是太大.
太大很難做.
我希望能夠有三個區域.
將我們分成三區.
好好的去處理.
小一點的組別才能重疊.
但我心目中.
一邊討估一邊看.
一邊陪著區長去走.
其實我看到有兩個.
但兩個不成事.
我和上帝說.
兩個不成事.
一定要三個.
但數算日子.
上帝2025年.
那怎麼辦.
明年就由得他了.
你解僱我就算了.
但真的很有趣.
我經歷神經很奇妙.
在12月的時候.
我突然有個區長.
走進來.
崇拜.
走進來我那邊.
師母.
我有事找你和你談.
我以為.
我們男士木區發生事.

$^{1561}$馬上進去.
一坐下.
他跟我說.
他說不行了.
我已經被上帝催逼.
我想和你說.
我準備做區木.
真是打冷顫.
你明不明.
因為我不知如何選擇.
很怕要選擇.
你明不明.
竟然有人自己拐個頭出來.
說我準備了.
如果我再不事奉.
我會死的.
我覺得.
頂曉梅.
不是十年前的見證.
是活在我們當中的見證.
看著.
枯骨一片.
不知如何是好.
但我們見著.
男士木區.
喜歡有人讀聖經的.
有課程的.
我們喜歡有人受浸的.
我說今年.
2026年.
我們的核心訓練.
32個人裡面.
有11個是來自.
男士木區的.
我們將榮耀歸給主.
頂曉梅.
人可以做甚麼.
沒事做.
所有事做不到.
板斧出不到.

$^{1601}$但是.
仍然工作.
所以無論你怎樣去看.
你自己的工作.
現在的環境.
糟了.
學了個專業.
沒事做.
然後.
我們所依賴的.
股票.
又沒收入.
現在開車.
貴到媽媽都認不出來.
然後.
接下來買魚很貴.
因為沒人去拿魚.
我們的環境灰暗.
但是.
記住.
我們的上帝.
是仍然掌權.
只要我們去親近他.
你的人生.
還可以有更精彩的日子.
Amen.
讓我們一起站立.
在神的面前.
真的好好切切的去禱告.
求上帝給我們.
聖靈去掌管我們的心.
拿走我們一切的錫心.
讓我們謙卑的.
去親近他.
沒人可以幫你的.
你的生命.
你的人生.
是要你和上帝一起去走.
稍為低頭.
有一段安靜的時間.

$^{1641}$當敬拜隊唱的時候.
我鼓勵你.
都可以有一段禱告.
真的打開你的心.
跟上帝說.
上帝你掌管我的心.
你幫助我.
原來我有些東西.
是很硬的.
求你拿走.
讓我敏感於你.
體貼你的聲音.
\newpage



\section{路加福音 14:1-24}
\label{sec:031965}
\textbf{筵席上的教導}~~~~俞大鳴牧師
\newline
\newline
連結: \href{http://www.church.com.hk/acms/content.asp?site=cdc\&op=show\&type=product\&code=031965\&layout=sermon}{\texttt{031965}} ~~~~ 語音日期: 2026-04-12
\newline
\newline
\hyperref[sec:031964]{\small{< < < PREV SERMON < < <}}
~
\hyperref[sec:index_chronic]{\small{[返順時目]}}
~
\hyperref[sec:031966]{\small{> > > NEXT SERMON > > >}}
\newline
\newline
路加福音 14:1-24
\newline
\begin{longtable}{cl}
\hline
\hline
章節 & 經文 (和合本修訂版)\\
\hline
14:1 & \begin{tabularx}{0.7\textwidth}{X} 安息日，耶穌到一個法利賽人的領袖家裡去吃飯，他們就窺探他。 \end{tabularx} \\ \\ \relax
14:2 & \begin{tabularx}{0.7\textwidth}{X} 這時在他面前有一個患水腫病的人。 \end{tabularx} \\ \\ \relax
14:3 & \begin{tabularx}{0.7\textwidth}{X} 耶穌回答律法師和法利賽人，說：「安息日治病合不合法？」 \end{tabularx} \\ \\ \relax
14:4 & \begin{tabularx}{0.7\textwidth}{X} 他們卻不說話。耶穌扶著那人，治好了他，叫他走了。 \end{tabularx} \\ \\ \relax
14:5 & \begin{tabularx}{0.7\textwidth}{X} 耶穌對他們說：「你們中間誰有兒子或有牛在安息日掉在井裡，不立刻拉他上來呢？」 \end{tabularx} \\ \\ \relax
14:6 & \begin{tabularx}{0.7\textwidth}{X} 他們對這些事不能反駁。 \end{tabularx} \\ \\ \relax
14:7 & \begin{tabularx}{0.7\textwidth}{X} 耶穌見所請的客人選擇首位，就用比喻對他們說： \end{tabularx} \\ \\ \relax
14:8 & \begin{tabularx}{0.7\textwidth}{X} 「你被人請去赴婚宴，不要坐在首位上，恐怕主人請了比你尊貴的客人， \end{tabularx} \\ \\ \relax
14:9 & \begin{tabularx}{0.7\textwidth}{X} 請了你和他的那人前來，對你說：『請讓座給這一位吧。』你就羞羞慚慚地退到末位去了。 \end{tabularx} \\ \\ \relax
14:10 & \begin{tabularx}{0.7\textwidth}{X} 你被請的時候，去坐在末位上，好讓主人來對你說：『朋友，請上座。』那時，你在同席的人面前就有光彩了。 \end{tabularx} \\ \\ \relax
14:11 & \begin{tabularx}{0.7\textwidth}{X} 因為凡自高的，必降為卑；自甘卑微的，必升為高。」 \end{tabularx} \\ \\ \relax
14:12 & \begin{tabularx}{0.7\textwidth}{X} 耶穌又對請他的人說：「你準備午飯或晚餐，不要請你的朋友、弟兄、親屬和富足的鄰舍，免得他們回請你，你就得了報答。 \end{tabularx} \\ \\ \relax
14:13 & \begin{tabularx}{0.7\textwidth}{X} 你擺設宴席，倒要請那貧窮的、殘疾的、瘸腿的、失明的， \end{tabularx} \\ \\ \relax
14:14 & \begin{tabularx}{0.7\textwidth}{X} 你就有福了！因為他們沒有甚麼可報答你。到義人復活的時候，你要得到報答。」 \end{tabularx} \\ \\ \relax
14:15 & \begin{tabularx}{0.7\textwidth}{X} 同席的有一人聽見這些話，就對耶穌說：「在神國裡吃飯的有福了！」 \end{tabularx} \\ \\ \relax
14:16 & \begin{tabularx}{0.7\textwidth}{X} 耶穌對他說：「有人擺設大宴席，請了許多客人。 \end{tabularx} \\ \\ \relax
14:17 & \begin{tabularx}{0.7\textwidth}{X} 到了坐席的時候，他打發僕人去對所請的人說：『請來吧！樣樣都已齊備了。』 \end{tabularx} \\ \\ \relax
14:18 & \begin{tabularx}{0.7\textwidth}{X} 眾人異口同聲地推辭。頭一個對他說：『我買了一塊地，必須去看看。請你准我辭了。』 \end{tabularx} \\ \\ \relax
14:19 & \begin{tabularx}{0.7\textwidth}{X} 另一個說：『我買了五對牛，要去試一試。請你准我辭了。』 \end{tabularx} \\ \\ \relax
14:20 & \begin{tabularx}{0.7\textwidth}{X} 又有一個說：『我才娶了妻子，所以不能去。』 \end{tabularx} \\ \\ \relax
14:21 & \begin{tabularx}{0.7\textwidth}{X} 那僕人回來，把這些事都告訴了主人。這家的主人就發怒，對僕人說：『快出去，到城裡大街小巷，領那貧窮的、殘疾的、失明的、瘸腿的來。』 \end{tabularx} \\ \\ \relax
14:22 & \begin{tabularx}{0.7\textwidth}{X} 僕人說：『主啊，你所吩咐的已經辦了，還有空位。』 \end{tabularx} \\ \\ \relax
14:23 & \begin{tabularx}{0.7\textwidth}{X} 主人對僕人說：『你出去，到大街小巷強拉人進來，坐滿我的屋子。 \end{tabularx} \\ \\ \relax
14:24 & \begin{tabularx}{0.7\textwidth}{X} 我告訴你們，先前所請的人沒有一個可以嘗到我的宴席。』」 \end{tabularx} \\ \\
[1ex]
\hline
\hline
\end{longtable}
$^{1}$猶太人都好鍾意飲飲食食.
佢地如果擺婚宴.
食足七日.
就唔係好似華人咁食一餐.
並且好多時候.
佢地亦都係會樂意款待遠方ge 來客.
因為佢地先祖亞伯拉罕就係好好客.
並且曾經因為招待一D 陌生人.
原來gwo 幾位係天使黎.
招待客人就成為jor 猶太人中其中一種善行.
佢地認為係一種善行.
跟住耶穌一路慢慢行去耶路撒冷.
係佢地行路ge 路徑ge 時陣.
佢地一般都會去會堂gwo 度祈禱講道.
所以亦都會有當地ge 熱心宗教人士.
法利座人請佢去食飯.
不過猶太人請食飯好多時候係公開.
左鄰右里都可以不請自來.
佢地未必食飯.
但係佢地會企埋一邊.
或者係企係窗外面.
黎到聽主人同客人ge 對話.
今日呢一段ge 經文.
就發生係耶穌又比人請去食飯ge 場合.
邀請耶穌ge 人係一位法利座人ge 領袖.
佢唔單止係法利座人.
就係領袖級ge 人物.
咁一般情況之下.
做客人ge 都會客氣.
你好我好大家好.
不過耶穌係上帝ge 兒子.
係真理ge 化身.
法利座人雖然請佢食飯.
耶穌有時都會不留情面咁責備佢地.
係十一章三十七至五十三節gwo 段經文.
我地已經睇到耶穌.
亦都係比法利座人請食飯.
gwo 次耶穌用jor 六字和哉.
或者遭殃jor 責備法利座人同法律教師.
今次又係法利座人.

$^{41}$並且係領袖請食飯.
耶穌亦都責備或者教導在場ge 人.
不過責備ge 方式今次用jor 問題ge 方式.
教導方式或者係比喻.
咁樣我地從耶穌對法利座人ge 責備裡面.
去學習一D ge 真理.
今次經文可以分為三段.
從時間上黎講.
係未食飯之前.
係預備坐席或者埋位ge 時候.
同埋坐低jor 之後.
第一段係十四章一至到六節.
呢段經文係有關安息日ge 責備.
對猶太人黎講.
安息日係一個好重要ge 日子.
因為猶太人同上帝建立關係.
係要通過納藥.
使佢地能夠成為上帝ge 子民.
而上帝就成為jor 佢地ge 上帝.
咁我地一般聽猶太人同上帝納gwo 藥.
就叫舊約.
舊約有幾個唔同ge 藥.
而安息日就成為jor .
上帝同猶太人納藥ge 一D 記號.
我地納jor 藥有個記號.
我地結婚多數戴個戒指.
今日我都有戴返個戒指.
猶太人同上帝納藥用乜ye 做記號呢.
其實就係用jor 安息日做記號.
上帝規定猶太人係一定要遵守安息日.
唔守安息日ge 猶太人點呀.
好大件事ge .
要被處死ge .
死刑黎ge .
我地睇兩節經文.
係《出埃及記》31章14至15節.
你們要守安息日以他為聖日.
凡干犯者日的必被處死.
凡在這日做工的那人必從百姓中剪除.
六日要做工.

$^{81}$但第七日是向耶和華守完全安息的安息聖日.
凡在安息日做工的必被處死.
好嚴重呀 講兩次.
你係呢日千萬唔好做工.
你做工要處死.
安息日係聖日.
就係話將gwo 一日變成奉獻俾上帝ge 日子.
但人總係好難唔犯錯.
守安息聖日慢慢慢慢就變成形式化.
gwo 日雖然唔做工.
但守節ge 人又唔係自己紀念上帝.
因此係以賽亞先知年代.
先知責備當時的猶太人.
係以賽亞仙58章13節.
佢話你又禁止自己D 腳踐踏安息日.
好嚴重喎.
唔係當gwo 日聖日 你D 腳踐踏安息日.
不在我地聖日做自己高興的事.
稱安息日為可喜樂的.
稱耶和華的聖日為可尊重的.
尊敬這日.
不走自己的道路.
不求自己的喜悅.
也不隨意說話.
那麼你地就會以耶和華為樂.
我地留意到上帝責備gwo D 猶太人.
係安息日係非常自我中心.
佢地可能係停止jor 工作 佢地真係唔做ye .
但係佢地內心係做咩呢.
佢地內心係做自己高興做ge 事.
佢地尋求自己ge 喜悅.
亦都隨意講說話.
先知責備佢地.
唔能夠只係係外表停止工作.
守安息日.
而內心卻係拜自己.
做自己喜歡做ge 事.
換成我地今日ge 環境黎講.
我地係主日可能唔洗返工.
唔去返工咁就黎o下教會.

$^{121}$但係心思意念上係諗住自己ge 事.
你崇拜我地要趕住做D 咩工程.
或者有咩灑落.
或者有時今晚有場重要ge 球賽.
掛住去預備睇球賽.
有大球賽ge 日子我發覺晚堂D 人會少D .
有D 人索性先做自己ge 事.
然後一日頭唔黎崇拜.
教會有晚堂就行啦.
我搞掂自己ge ye 晚上先去崇拜啦.
如果我地用呢種態度黎到守主日.
我地真係應該自我檢討.
猶太國後來被巴比倫滅國之後.
又被巴比倫王強逼移民去其他地方居住.
過jor 幾十年.
佢地先至比波斯王准許佢地回歸故鄉.
有好多猶太人就明白.
佢地國家ge 慘痛遭遇.
係上帝對佢地ge 懲罰.
因此當佢地有機會返到故鄉之後.
有部分人就好嚴謹遵守舊約界明.
其中一派ge 人就係法利賽人.
對佢地黎講.
安息日係唔能夠醫病.
唔好做工呀嘛.
醫病都係一種工作啦.
耶穌知道佢地ge 規定.
但仍然挑戰佢地ge 諗法.
當佢地遵守上帝吩咐ge 時候.
佢地變成jor 教條主義ge 人.
佢地變得好不近人情.
失去jor 人性.
仲以為自己好神聖.
前幾次耶穌係安息日.
醫好人之後再同佢地理論.
呢一次呢.
耶穌去到法利賽人ge 屋企.
見到有一位有病ge 人.
就先問下gwo D 法利賽人.
同埋法律教師一個問題.

$^{161}$我地ge 法律准唔准佢係安息日治病ge 呢.
事實上呢耶穌ge 問題呢.
係屬於明知故問ge 問題.
法利賽人.
亦都唔知道應該點樣回答耶穌.
因此佢地唔出聲唔回答.
嘗試以靜制動.
耶穌ge 反應.
耶穌ge 反應就首先醫好gwo 個有病ge 人.
然後話呢.
再問佢地.
你地當中誰有兒子或牛.
係安息日調進井裡.
不立刻拉他上來呢.
耶穌呢個問題呢.
就係指責gwo D 法利賽人ge 虛偽ge 行為.
佢地對自己ge 財產.
對自己ge 生足親人.
唔理人咩日子有危險ge 時候一定會拯救ge .
但係對其他人呢就缺乏憐憫之情.
上帝雖然要人遵守安息日.
但gwo 個呢係為人ge 身心靈健康而定ge 例.
而唔係要人變成冷酷無情.
教條主義ge 人.
實際上有一位叫梅加ge 先知就曾經話.
西人啊,耶華爾子是你,何謂善?.
他向你所要的是什麼呢?.
只要你行公義,好憐憫.
全謙卑的心與你的上帝同行.
上帝要求跟隨ge 人要孝乏懼.
要有愛心同埋憐憫ge 心.
對待我地周圍ge 人.
主要我地要伸張正義.
心存謙卑與神同行.
呢段經文大概就發生係埋位食飯之前.
跟住落黎ge 經文7-11節就可能預備埋位.
大家坐埋去.
主耶穌就留意到有D 客人.
佢地真係好唔客氣.
自己就搵jor 賢直上ge 貴賓位.

$^{201}$我就唔清楚佢地點樣分.
邊D 係主人位,邊D 係貴賓位.
我地華人好多時候就隨便坐.
而家有D 唔同ge 新時代.
就全部偏好.
邊D 貴賓位就梗係坐近主人gwo 邊.
但係猶太人有D 位原來係貴賓位.
並且貴賓通常會遲到.
顯示佢地尊貴ge 身份.
幫助我地明白點解主耶穌話.
恐怕有比你更受尊重ge 客人.
也在被邀請之列.
即係有時D 貴賓會遲到,仲未到.
呢D 早到ge 人自己走去坐上貴賓位.
而個位貴賓後黎到ge 話.
主人就會請原先擺jor 位.
「唔該你讓開,有位貴賓要坐呢個位」.
如果咁ge 情形發生.
原先坐jor 上去擺貴賓位ge 人就會好尷尬.
實上主耶穌呢個教導.
唔係一個咩新鮮ge 教導.
唔係主耶穌自己諗出黎ge ye .
而係傳統ge 智慧.
因為係《尊賢書》二十五章七節就有講.
「寧可讓人家說請你上到這裡來.
強於在你近見的貴人面前令你退下」.
即係尊賢ge 智慧.
我地先一排思想個尊賢.
尊賢裡面其實有好多智慧.
可唔可以留意到呢.
係《新約聖經》裡面呢.
gwo D 實踐ge 部分.
好多時都係來自尊賢ge 教導.
只不過呢今次呢度ge 客人呢.
有幾位可能係發你錯人.
佢地係有知識ge .
但係呢佢地自己係實踐上就有D 問題.
佢地習慣jor 係會堂上就坐上台.
咁我唔知大家有冇留意.
如果係係舊《新約》時代.

$^{241}$會堂係冇椅ge .
冇我地今日咁ge .
邊D 人係坐係度ge .
教師坐係度ge .
D 宗教領袖坐上台ge .
得幾張椅ge .
會眾做乜ye .
會眾就企係度聽ye 囉.
係咁樣ge .
同我而家調轉.
而家講完要企係度.
會眾就係度.
以前唔係架.
以前係調轉你架.
佢地就有知識冇乜行為.
我地信jor 主人地要小心.
唔好單有知識而冇行為.
我地要記得主耶穌.
係呢度講下十一節.
凡至高的必降為卑.
至今卑微的必升為高.
十二節至到十四節ge 教訓.
我地相信主耶穌.
係順帶將上帝所要求ge 善行.
教導gwo 位請佢食飯ge 法利賽人領袖.
真正ge 公義唔係你對我好.
我又回敬你.
你請我食飯.
我又請返你食飯.
gwo 個只係朋友.
朋友或者親戚之間往來.
甚至乎有時人地請食飯.
係有其他目的ge .
包括係為jor 自己ge 方便.
將來做事ge 方便.
我記得我未移民之前.
係香港曾經係一間大車行ge 修理部.
做文員.
gwo 時職位好低.
但久唔久銷售部有D 經紀或者經理.

$^{281}$會請我地食飯.
點解佢地咁樣做呢.
有目的ge .
請jor 我地食飯之後.
將來如果佢有D 客仔.
要收車ge 時候.
我地gwo 度好多時係要排隊.
要預約ge .
預約可能兩三個星期之後先有得收.
佢地就想請我地食飯.
有咩事ge 時候打電話.
呢度我熟客仔.
唔該你幫佢排快D 啦.
係咁啦.
呢個社會上ge 人.
請食飯多數都有D 目的.
但係耶穌就係呢度教導佢地.
我地請人食飯.
如果係屬於真正ge 義行.
係應該照顧gwo D 貧窮ge .
殘疾ge .
跛腳ge .
失明ge 人.
即係話社會上ge 弱勢群體.
佢地冇能力報答你.
咁你係義人復活ge 時候.
上帝要親自報答你.
講jor 咁多ye .
仲未開始食飯.
第一段就arm arm 去到法理上係屋企.
第二段就係入席埋位.
最後一段就係可能開始食ge 時候.
有一位客人聽見jor 耶穌ge 話.
就回應話.
能夠在上帝ge 國裡面享受筵席ge 人.
多麼有福呀.
耶穌從呢句話ge .
講呢句話ge 人ge 表情或者語氣.
感覺到呢個人可能呢.
係佢自己其實好有把握.

$^{321}$覺得自己係能夠係上帝國裡面享受筵席.
所以佢自覺都好有福氣.
於是耶穌又講一個大牌筵席ge 比喻.
黎去教授佢.
從呢個比喻裡面呢.
我地知道猶太人請客呢.
係先指定日期呀.
講好jor 邊日.
但係幾分鐘又唔講架喎.
冇講呀到時等通知喎.
到日期之後係呢個比喻當中呢.
我相信呢個筵席呢.
開始可以黎ge 時候主人就派僕人啦.
佢請gwo D 預先請jor ge 客人.
佢地都答應jor 要黎參加筵席ge 人.
而家可以過黎啦請你黎啦咁樣.
但係呢我地發覺原先答應jor 黎ge 人呢.
個個都唔黎喎.
就搵藉口啦拒絕啦.
咁事實上呢咁樣拒絕呢.
真係好唔比面子比主人.
因為呢主人要知道jor 有幾多人黎.
佢要預備足夠ge 食物架嘛.
可能要tong 三隻羊呀幾隻牛呀咁樣.
係預備好啦.
你到筵席ge 時候先至話唔能夠黎呀.
係好唔禮貌.
變成呢係浪費食物.
所以呢個主人呢就非常之嬲呀.
吩咐僕人你去請gwo D 貧窮殘疾.
失明跛腳ge 都黎.
最後呢仲有空位.
主人就吩咐僕人啦.
你盡量出去街上面搵人黎坐到滿為止.
呢個比喻ge 含義係乜ye 呢.
就有D 先前被請答應要黎.
但又唔能夠黎ge 客人.
對佢地黎講呢就係.
其實就係責備gwo D 法利賽人.
開始ge 時候呢gwo 一位話.

$^{361}$能夠係上帝國裡享受賢職ge 人.
多麼有福呀.
係讚揚自己同埋其他在場ge 法利賽人.
佢地自己認為佢地真係有福ge 人.
都認為能夠進入上帝國ge 人.
就係好似法利賽人咁樣.
謹慎遵守上帝ge 法律ge 人.
而耶穌ge 比喻呢係話.
能夠參加天國賢職ge 人.
係受到主人邀請ge 人.
同埋係當賢職預備好之後呢.
收到通知肯黎參加賢職ge 人.
佢地先係有福ge 人.
猶太人係先收到上帝邀請ge 人.
上帝係先同佢地納藥.
好可惜佢地時時都係想要靠自己ge 行為.
爭取進入天國.
佢地以為自己有好ge 行為.
就忽略jor 天國主人ge 心意.
唔願意接受上帝慷慨ge 恩典.
堅持要靠自己ge 行為黎到進天國.
於是天國ge 僕人就係耶穌黎.
佢大聲嗌:「天國近了,你們要悔改,順福音」.
就係咁簡單je .
賢職已經預備好架啦.
黎參加啦.
法理賽派猶太人呢.
就唔肯接受耶穌.
所以上帝就邀請gwo D 貧窮ge 殘疾ge 失明ge 跛腳ge .
黎到參加天國ge 盛宴.
gwo D 人呢就係我地呢D 肯悔改,相信耶穌ge 人.
我唔知當中有冇人係被勉強,被人勉強信耶穌.
但即使你被人勉強呀.
好似二十三節gwo 度所講呢.
被人強拉進來ge .
不過呢個係天國ge 盛宴.
能夠參加到ge 人都係真正有福ge 人.
實際上今日世界有好多宗教都強調呢.
要靠行為先至可以得救.
包括佛教同回教.

$^{401}$唯獨基督教係強調.
得救係同人ge 行為無關.
完全係上帝ge 恩典.
就好似我地出世咋嘛.
每一個人出世之前.
有冇比父母ge 工作過架.
其實冇架喎.
我地每一個人父母生我出黎.
係父母比我地ge 恩典黎架喎.
唔係因為我地為父母做jor 好多ye .
係你應該生我出黎養我下.
照顧我,唔係咁架.
係父母比我地ge 恩典.
最近睇到一個人生ge 比喻.
人生好似一張有去無回ge 單程車票.
你係冇機會再重頭黎過.
當你去到尾站ge 時候.
發覺gwo 度係地獄ge 時候.
就已經太遲.
因為你冇回頭車票.
冇回頭車搭.
又有D 人就話人生好似演戲咁樣.
不過你出黎係冇機會彩排.
唔似得D 演員.
係演出之前係可以練習.
人生ge 戲每一場都係現場直播.
所以每個人都要把握好你現在有ge 機會.
接受耶穌ge 邀請.
準備參加天國ge 盛宴.
對我地已經信jor 主ge 弟兄姊妹.
呢個比喻提醒我地.
我地能夠參加天國ge 盛宴.
唔係因為我地有幾好.
唔係依靠我地ge 行為.
係出於上帝ge 恩典.
唔係我配受邀請.
我地本來就係罪人.
係貧窮殘疾失明跛腳.
係上帝ge 恩典讓我地.
能夠參加到天國ge 盛宴.

$^{441}$因此我地係一無可跨.
只能夠全心謙卑ge .
黎到跟隨主耶穌.
但願我地都樂意去邀請.
其他人黎到參加天國ge 盛宴.
因為係天國裡面仲有好多空位.
讓我做一個總結.
耶穌係受邀參加法利賽人ge 筵席.
佢arm arm 係安息日.
耶穌就醫好一位病人.
違反jor 猶太人ge 規則.
但佢冇違反到律法ge 精神.
安息日係上帝定落黎要人學習.
去紀念上帝.
但安息日唔只係應該停止日常ge 工作.
而係要有內心ge 反應.
去內心ge 反省.
黎到去建立自己同上帝ge 關係.
以賽亞先知ge 提醒係好重要.
我地今日首主日.
有冇一主日.
只係去做一D 自己喜歡做ge 事.
隨便講說話.
尋求自己開心.
第二點就係要謙卑.
唔好自高.
有機會請客ge 話.
就請一D 冇能力回報ge 人.
先至係真正ge 善行.
最後就係我地已經接受天國盛宴ge 人.
我地要黎.
亦都要邀請我地親戚朋友.
一齊黎參加呢個充滿上帝恩典ge 盛宴.
同我一齊祈禱.
你要請我地黎到天國.
同你建立和諧ge 關係.
求你幫助我地將主日睇為聖日.
係與你有約ge 日子.
讓我地係主日多思想你ge 話.
同埋天國ge 事務.

$^{481}$奉耶穌基督ge 名祈禱.
阿們.
\newpage



\section{士師記 1:1-2:5}
\label{sec:031966}
\textbf{士師記析讀 (1) 鳥瞰士師記}~~~~梁永善牧師
\newline
\newline
連結: \href{http://www.church.com.hk/acms/content.asp?site=cdc\&op=show\&type=product\&code=031966\&layout=sermon}{\texttt{031966}} ~~~~ 語音日期: 2026-04-15
\newline
\newline
\hyperref[sec:031965]{\small{< < < PREV SERMON < < <}}
~
\hyperref[sec:index_chronic]{\small{[返順時目]}}
~
\hyperref[sec:031967]{\small{> > > NEXT SERMON > > >}}
\newline
\newline
士師記 1:1-2:5
\newline
\begin{longtable}{cl}
\hline
\hline
章節 & 經文 (和合本修訂版)\\
\hline
1:1 & \begin{tabularx}{0.7\textwidth}{X} 約書亞死後，以色列人求問耶和華說：「我們中間誰當首先上去攻打迦南人，與他們爭戰呢？」 \end{tabularx} \\ \\ \relax
1:2 & \begin{tabularx}{0.7\textwidth}{X} 耶和華說：「猶大要先上去。看哪，我已將那地交在他手中。」 \end{tabularx} \\ \\ \relax
1:3 & \begin{tabularx}{0.7\textwidth}{X} 猶大對他哥哥西緬說：「請你同我上到我抽籤所得之地，與迦南人爭戰；我也同你去你抽籤所得之地。」於是西緬與他同去。 \end{tabularx} \\ \\ \relax
1:4 & \begin{tabularx}{0.7\textwidth}{X} 猶大就上去，耶和華把迦南人和比利洗人交在他們手中。他們在比色擊殺了一萬人。 \end{tabularx} \\ \\ \relax
1:5 & \begin{tabularx}{0.7\textwidth}{X} 他們在比色遇見亞多尼‧比色，與他爭戰，擊敗了迦南人和比利洗人。 \end{tabularx} \\ \\ \relax
1:6 & \begin{tabularx}{0.7\textwidth}{X} 亞多尼‧比色逃跑，他們追趕他，捉住他，砍斷他大拇指和大腳趾。 \end{tabularx} \\ \\ \relax
1:7 & \begin{tabularx}{0.7\textwidth}{X} 亞多尼‧比色說：「從前有七十個王，大拇指和大腳趾都被我砍斷，在我桌子底下拾取零碎食物。現在神照著我所做的報應我了。」他們把亞多尼‧比色帶到耶路撒冷，他就死在那裡。 \end{tabularx} \\ \\ \relax
1:8 & \begin{tabularx}{0.7\textwidth}{X} 猶大人攻打耶路撒冷，奪取了它，用刀殺城內的人，並且放火燒城。 \end{tabularx} \\ \\ \relax
1:9 & \begin{tabularx}{0.7\textwidth}{X} 後來猶大人下去，與住山區、尼革夫和低地的迦南人爭戰。 \end{tabularx} \\ \\ \relax
1:10 & \begin{tabularx}{0.7\textwidth}{X} 猶大去攻打住希伯崙的迦南人，殺了示篩、亞希幔、撻買。希伯崙從前名叫基列‧亞巴。 \end{tabularx} \\ \\ \relax
1:11 & \begin{tabularx}{0.7\textwidth}{X} 猶大從那裡去攻擊底壁的居民。底壁從前名叫基列‧西弗。 \end{tabularx} \\ \\ \relax
1:12 & \begin{tabularx}{0.7\textwidth}{X} 迦勒說：「誰能攻打基列‧西弗，奪取那城，我就把我女兒押撒嫁給他。」 \end{tabularx} \\ \\ \relax
1:13 & \begin{tabularx}{0.7\textwidth}{X} 迦勒的弟弟基納斯的兒子俄陀聶奪取了那城，迦勒就把女兒押撒嫁給他。 \end{tabularx} \\ \\ \relax
1:14 & \begin{tabularx}{0.7\textwidth}{X} 押撒來的時候，催促丈夫向她父親要一塊田。押撒一下驢，迦勒就對她說：「你要甚麼？」 \end{tabularx} \\ \\ \relax
1:15 & \begin{tabularx}{0.7\textwidth}{X} 她對迦勒說：「求你賜我福分；你既然把尼革夫給了我，求你也給我水泉。」迦勒就把上泉和下泉都賜給她。 \end{tabularx} \\ \\ \relax
1:16 & \begin{tabularx}{0.7\textwidth}{X} 摩西的岳父是基尼人，他的子孫與猶大人一起上到棕樹城，往亞拉得以南的猶大曠野去，住在百姓當中。 \end{tabularx} \\ \\ \relax
1:17 & \begin{tabularx}{0.7\textwidth}{X} 猶大和他哥哥西緬同去，擊殺了住洗法的迦南人，將城徹底毀滅。因此，那城名叫何珥瑪。 \end{tabularx} \\ \\ \relax
1:18 & \begin{tabularx}{0.7\textwidth}{X} 猶大攻取了迦薩和所屬的領土，亞實基倫和所屬的領土，以革倫和所屬的領土。 \end{tabularx} \\ \\ \relax
1:19 & \begin{tabularx}{0.7\textwidth}{X} 耶和華與猶大同在，猶大取得了山區，卻不能趕出平原的居民，因為他們有鐵的戰車。 \end{tabularx} \\ \\ \relax
1:20 & \begin{tabularx}{0.7\textwidth}{X} 以色列人照摩西所說的，把希伯崙給了迦勒。迦勒從那裡趕出亞衲的三支後裔。 \end{tabularx} \\ \\ \relax
1:21 & \begin{tabularx}{0.7\textwidth}{X} 至於住耶路撒冷的耶布斯人，便雅憫人沒有把他們趕出。於是，耶布斯人與便雅憫人同住在耶路撒冷，直到今日。 \end{tabularx} \\ \\ \relax
1:22 & \begin{tabularx}{0.7\textwidth}{X} 約瑟家也上到伯特利去，耶和華與他們同在。 \end{tabularx} \\ \\ \relax
1:23 & \begin{tabularx}{0.7\textwidth}{X} 約瑟家去窺探伯特利，那城起先名叫路斯。 \end{tabularx} \\ \\ \relax
1:24 & \begin{tabularx}{0.7\textwidth}{X} 探子看見一個人從城裡出來，就對他說：「請你把進城的路指示我們，我們會厚待你。」 \end{tabularx} \\ \\ \relax
1:25 & \begin{tabularx}{0.7\textwidth}{X} 那人把進城的路指示他們。他們就用刀擊殺了城中的居民，卻放走那人和他的全家。 \end{tabularx} \\ \\ \relax
1:26 & \begin{tabularx}{0.7\textwidth}{X} 那人往赫人之地去，建造了一座城，起名叫路斯。那城到如今還叫這名。 \end{tabularx} \\ \\ \relax
1:27 & \begin{tabularx}{0.7\textwidth}{X} 瑪拿西沒有趕出伯˙善和所屬鄉鎮的居民，他納和所屬鄉鎮的居民，多珥和所屬鄉鎮的居民，以伯蓮和所屬鄉鎮的居民，米吉多和所屬鄉鎮的居民；迦南人仍堅持住在這地。 \end{tabularx} \\ \\ \relax
1:28 & \begin{tabularx}{0.7\textwidth}{X} 以色列強盛的時候，就叫迦南人做苦工，沒有把他們全然趕走。 \end{tabularx} \\ \\ \relax
1:29 & \begin{tabularx}{0.7\textwidth}{X} 以法蓮沒有趕出住基色的迦南人。於是迦南人仍住在基色，在以法蓮中間。 \end{tabularx} \\ \\ \relax
1:30 & \begin{tabularx}{0.7\textwidth}{X} 西布倫沒有趕出基倫的居民和拿哈拉的居民。於是迦南人仍住在西布倫中間，成了服勞役的人。 \end{tabularx} \\ \\ \relax
1:31 & \begin{tabularx}{0.7\textwidth}{X} 亞設沒有趕出亞柯的居民和西頓的居民，以及亞黑拉、亞革悉、黑巴、亞弗革和利合的居民。 \end{tabularx} \\ \\ \relax
1:32 & \begin{tabularx}{0.7\textwidth}{X} 亞設人因為沒有趕出那地的居民迦南人，就住在他們中間。 \end{tabularx} \\ \\ \relax
1:33 & \begin{tabularx}{0.7\textwidth}{X} 拿弗他利沒有趕出伯˙示麥和伯˙亞納的居民。於是拿弗他利就住在那地的居民迦南人中，而伯˙示麥和伯˙亞納的居民卻成了為他們服勞役的人。 \end{tabularx} \\ \\ \relax
1:34 & \begin{tabularx}{0.7\textwidth}{X} 亞摩利人強逼但人住在山區，不准他們下到平原。 \end{tabularx} \\ \\ \relax
1:35 & \begin{tabularx}{0.7\textwidth}{X} 亞摩利人仍堅持住在希烈山、亞雅崙和沙賓；然而約瑟家權勢強盛的時候，他們成為服勞役的人。 \end{tabularx} \\ \\ \relax
1:36 & \begin{tabularx}{0.7\textwidth}{X} 亞摩利人的地界是從亞克拉濱斜坡，從西拉延伸而上。 \end{tabularx} \\ \\ \relax
2:1 & \begin{tabularx}{0.7\textwidth}{X} 耶和華的使者從吉甲上到波金，說：「我領你們從埃及上來，帶你們到我向你們列祖起誓應許之地。我曾說：『我永不廢棄我與你們的約。 \end{tabularx} \\ \\ \relax
2:2 & \begin{tabularx}{0.7\textwidth}{X} 你們不可與這地的居民立約，要拆毀他們的祭壇。』你們竟沒有聽從我的話。你們為何這樣做呢！ \end{tabularx} \\ \\ \relax
2:3 & \begin{tabularx}{0.7\textwidth}{X} 因此我說：『我必不將他們從你們面前趕出。他們必作你們肋下的荊棘，他們的神明必成為你們的圈套。』」 \end{tabularx} \\ \\ \relax
2:4 & \begin{tabularx}{0.7\textwidth}{X} 耶和華的使者向以色列眾人說這些話的時候，百姓放聲大哭。 \end{tabularx} \\ \\ \relax
2:5 & \begin{tabularx}{0.7\textwidth}{X} 於是他們給那地方起名叫波金，並在那裡向耶和華獻祭。 \end{tabularx} \\ \\
[1ex]
\hline
\hline
\end{longtable}
$^{1}$我們開始《時詩記》.
在《時詩記》讓我們看到一個很重要的題目.
《時詩記》是一本很特別的書.
涵蓋的時期是以色列被撈之前.
最黑暗的一段歷史.
今晚其實會較為悶.
我記得我少年時代已經開始去研讀《聖經》.
也會看《釋經》書.
當我買《釋經》書的時候.
前面當然有個《言曹德順》.
《引論》.
就會講《聖經》的作者,背景,寫作的特色.
當中有些什麼神學思想.
其實這些東西是很悶的.
所以我多數不看的.
一聽就跳下去.
立刻開始看第幾章,第幾節開始.
來看《釋經》.
但當我再成熟一點的時候.
就發覺到我們應該對每一本書的背景要有所瞭解.
你瞭解了背景的神學思想.
於是你看那本書.
你纔會更加深入.
而不會有偏差.
所以今晚第一節講《時詩記》.
我們就先了解一下《時詩記》的內容是怎樣.
剛纔我已經講了.
《時詩記》涵蓋的時期.
是以色列被捕之前的一段最黑暗的歷史.
為什麼呢?.
後來又有了君王之後.
蘇羅王有一段風光的時間.
當然後來他失敗了.
到了大魏.
大魏作王的時候就開始穩定下來.
然後傳位給所羅門.
至少在王國時期.
他們就會有一個王管理著他們.
但是在沒有王之前.
和約書帶領他們進入迦南地這段歷史裡.

$^{41}$可以說是他們的黑暗時期.
很多人都忽略了這一本書.
想想你們教會的講道有多少時一些講道.
會引用《時詩記》.
我發現到年來.
我們最熟悉的就是《參孫》.
不知道大家有沒有看過一部電影.
我少年時代看的.
叫《霸王妖姬》.
那時候我好像還沒有信耶穌.
那個男主角叫做衛托米曹.
女主角就不知道了.
《參孫》這個故事那時候我就是這樣認識.
但原來我在美恩堂那30年.
我都沒有講過《參孫》.
今天很多人會忽略了這一本書.
又或者當《時詩記》.
就當那些英雄人物來看.
基甸.
我有講的基甸.
或者耶佛他.
或者參孫.
當是一個英雄故事.
這些《時詩記》很好.
但當我和大家一起再研讀《時詩記》的時候.
我就發覺原來作者絕對不是將這群人視為英雄.
如果和其他的書卷來比較.
其實這本書提醒我們.
我們的屬靈生命的歷程.
會不會就好像《時詩記》的記載一樣呢?.
《時詩記》最記載的特色是什麼呢?.
就是原本他們是給神祝福的.
但當在平靜安穩的時候.
他就離棄神了.
於是神就容讓這群子民.
被外敵來入侵.
被外敵來壓迫他們.
人就是這樣犯賤的.
在困難當中.
在痛苦當中.

$^{81}$他纔有醒覺.
於是他就向神有呼求.
我經常覺得我用這個術語可能是錯的.
其實神很容易哄的.
神很愛我們.
正如我剛剛說我去旅行.
第一次和我的家人全家一起去.
我已經說我的孫子和我的溝通越來越少了.
因為他們是Twins.
兩個都是講英文的.
你知道他們講英語的時候.
好像我們講廣東話講得很順的時候.
其他人都聽得懂的.
我聽不懂.
但只要我看著他們兩個很開心地喫東西.
看著他們兩個很開心地撿禮物.
我已經自動現身了.
我就想起原來我們的神也是這樣的.
當我們來到天上去親近神.
當我們來到天上向神來到天上去呼求.
或者我們來去讚賞神的時候.
神就很開心了.
好了.
當這群以色列人.
他們被壓迫的時候.
他纔想起神.
於是向神來到呼求.
於是神的憐憫就發動了.
結果神就會興起時司.
來解救他們.
幫他們脫離這個困壓.
或者透過這些時司.
來戰勝那些外敵.
然後就記載他們在未來的日子裡面.
就有太平幾多年了.
但是可能一段長時間之後.
他們又再重蹈覆轍.
又再犯錯.
神又再讓另外一些外邦人來對付他們.
於是他們又向神呼求.

$^{121}$神又再興起時司.
時司就是這樣不斷不斷循環下去的.
這豈不是好像我們的生命一樣嗎?.
我們檢討一下我們與神的關係.
會不會有時候我們在平安穩妥.
較為順利的情況當中.
我們就沒有好好的親近神.
甚至把神像藥丸一樣備用.
放在一邊.
有事就去求祂.
可能當我們困難的時候.
我們向神呼求.
神答應了我們的呼求.
幫我們解除了困壓.
我們也會有一段時間好的.
但是又久了的時候.
會不會我們又再不斷地重複複設呢?.
但是不要忘記.
原來當這樣的時候.
我們所面對的痛苦.
其實會越來越沉重的.
而時司記的生活應用.
似乎又不是那麼容易.
但我一定會在當中找到.
一些熟悉的教導和應用給大家.
這就是我一直講道.
我的查經和我的理念都是這樣.
神的話語不僅是對當代的人而言.
一定是對當代的人而言.
但也可以應用在我們今天的生活上.
但我絕對不會穿鑿附會.
就好像今年年初.
劉業強去車公廟求籤.
他自己竟然拿著籤.
來帶去解釋經文.
就講到香港要怎樣來帶去做得好.
怎樣配合國內.
我想車公廟如果真是顯靈的話.
會夾著頭.
各配凡都是這樣.

$^{161}$真是他有一個概念.
就是將這些東西拿上去.
這個就叫穿鑿附會.
我絕對不會.
我會先解釋聖經.
如何應用在我們生活中.
所以我們不是單單看一些歷史故事.
是鑒古詩今.
你也確信神藉著他的話語.
向我們講說話.
教導我們怎樣面對現在這個世代.
《時詩記》的作者.
向我們揭露了以色列民的悲劇.
我們見到子民.
及後對神律法那種輕視.
沒有很嚴謹.
沒有很認真地遵守神的律法.
因此使他們人際關係受到損害.
民族也受外敵的入侵.
所以我們這次《茶經》的進度.
當然是以士師作為一個主角.
但並不要當他們是英雄.
他們不是英雄.
如果我們深入研究的時候.
每一個士師.
除了頭兩三個叫做好些之外.
其實是每放月下.
是越來越差的.
最後詳細記載的就是參孫.
參孫是一個不知所謂的人.
所以每一個士師都有他的問題和缺陷.
而《時詩記》也提醒我們.
其實是我過去沒有留意到的.
就是要著重家庭.
家庭是很重要的.
家庭教育是很重要的.
因為如果我們很草率地去視一系列的英雄故事.
就失去了《時詩記》的作者所記載的真正意思.
也失去了《時詩記》的統一性和連貫性.
尤其是.

$^{201}$我們留意到《時詩記》最後那幾章.
就是說到當日他們的社會出現了什麼問題.
包括連神的僕人都出現了很多的問題.
所以這些的事.
譬如參孫和伯利恆事件是混為一體的.
《時詩記》的作者很精心地刻畫了幾個人物.
一會兒我們會看PowerPoint就會看到.
用一種很獨特的方式描繪他們.
然後引出他所講的論點.
作者很有技巧.
藉著這些人物的刻畫出《時詩》.
又帶出一個觀點.
好像《時詩記》第18章.
就講到一個神的僕人利未人.
為了賺錢.
不是去事奉神.
這個地方的經濟比我好一些.
於是我就跳槽了.
甚至不是真的帶出神的教導.
甚至他會帶著偶像去幫那些人.
都很可悲.
這就和第17章參孫的故事.
其實有很多相似的地方.
所以如果我們將這兩段經文來比較的時候.
就發覺到作者想帶出.
當時以色列人的領袖是怎樣的失敗.
子民是怎樣的有犯錯.
我們先看看作者.
《時詩記》的作者和《時詩記》的寫作者是怎樣.
希伯來文聖經的正典《時詩記》.
是屬於前先知書的部分.
《時詩記》的卷名來自武加大譯本.
就是《時詩》的生平.
或者是《七十事》譯本的《時詩》.
大家知道《七十事》譯本.
就是在亞歷山大的時候.
去請一幫精通希伯來文的學者.
去翻譯當時的聖經.
所以叫做《七十事》譯本.
聖經的書卷的名字.

$^{241}$其實是後人加上去的.
去表達當時人對聖經對這卷書的看法.
好了.
《時詩記》的作者是誰呢?.
不清楚.
不是好像《先知書》那樣.
是很清楚的說明.
先知是誰去寫這卷書.
我們是不清楚.
是誰來記下了《時詩記》.
就連寫書的日期.
都很難以理解去確定.
由於《時詩記》用的文字和主題兩方面.
和《若書也記》是有很多重疊和交叉之處.
所以按正常的推測.
或者合理的推測.
《時詩記》是在《若書也記》的後面.
當中的教訓.
可以應用在以色列的歷史年代裡面.
也指向王國朝期的歷史.
所以《時詩記》的事件大部分.
發生在若書也死了之後的時期.
我們看看《時詩記》第一章的第一節.
若書也死後.
以色列人求問耶和華說.
我們中間誰當首先想去攻擊迦南人.
與他們爭戰?.
這一節的經文.
我們下一個星期的《茶經》就會講到.
但這裡很清楚地講到.
當若書也死了之後.
他們就開始要面對這些事了.
所以《時詩記》應該是在若書也死了之後.
然後才發展的.
但作者也一路回憶追溯到若書也年代.
甚至引伸到大魏的年代.
我們看看《時詩記》第一章的下一節.
汴雅敏人沒有講出駐耶路撒冷的耶布斯人.
耶布斯人仍在耶路撒冷與汴雅敏人同住.
直到今天.

$^{281}$如果我們有過去《茶經》參加.
我講識讀大魏的時候就會知道.
大魏要建自己的首都.
要找什麼地方呢?.
就是耶布斯人.
就是耶路撒冷.
所以你留意到這裡一路追溯到大魏的時期.
所以我們很相信《時詩記》是介乎在.
若書也記和撒慕爾記之間.
如果將《時詩記》的年份加在一起.
一定不會早過大魏王朝.
所以很簡單做一個結論.
就是若書也死了之後.
到大魏王朝建立之前.
接著我們看一看《時詩記》的結構.
如果你留意到《時詩記》是有不斷的循環.
剛纔我不是講了嗎?.
而每一個循環都出現.
之前作者一些小的事實去記錄.
很簡單舉個例子.
我們看到那些事實.
剛纔我已經講了.
不是英雄.
他們其實是越來越差的.
每一個主題都有一個方向性.
譬如很強調.
他們的信心衰敗了.
就是巴拉.
我們稍後就會看巴拉.
巴拉你可能不太記得.
但有一個叫底波拉.
你一定會記得.
Deborah 是不是?.
巴拉神父召他去拯救以色列人.
但他很沒有信心.
一定要底波拉和他在一起.
你就會看到他信心那種衰敗.
第二我們看到.
接著就是基甸.
基甸很厲害.

$^{321}$以小勝多.
只是300人就可以打到很多人.
但如果你看到.
其實基甸的信心很缺乏.
一次又一次.
神向他顯現.
向他行神蹟.
他仍然沒有這一份信心.
所以就是信心的缺乏.
還有就是信心上的錯誤.
耶佛他是一個很特別的人.
不知道大家知不知道耶佛他的背景是什麼?.
他是妓女的兒子.
草莽之夫.
但後來他和神說.
如果你幫我打聖戰的時候.
我就會怎麼樣呢?.
誰一出來我家的時候.
第一個出來的.
我就會將他獻祭.
將他殺了.
想不到第一個出來迎見他的.
是他的女兒.
究竟他有沒有將女兒獻祭呢?.
這個我們當然會在日後來探討.
你看到他信心上的錯誤.
最後.
我們看不到參宣有任何的信心.
他全部都是憑自己血氣來行事.
所以你看到的時候就會知道.
原來作者列出的事實.
就是一械不如一械.
越來越差.
藉著這樣帶出當時的社會.
是否一個很混亂的世代呢?.
除了信心的主題之外.
記不記得我剛纔說.
是我以前都沒有留意到的.
原來都很著重講到關於家庭.
重點就是你要好好的重視你的家庭.

$^{361}$蒂波拉是一個母親.
但她的職責是不應該參與戰爭的.
但是後來無奈地參與了.
機電.
古時我們的故事裡面.
就是機電是怎樣來的呢?.
他求神去彰顯他自己的心意.
把羊毛弄得要溼或者要幹.
接著帶著幾百人就贏了.
就是米電人.
我們就停在這裡了.
我曾經講過都是停在這裡.
但是原來那個後著.
我們一直很多人都很少去看的.
機電的家庭裡面很失敗.
他有很多的兒女.
即是說他有很多的妻妾.
接著他的兒子就來同室抽窩.
殺死很多很多自己的兄弟.
他的錯誤就帶來了很多的不幸.
還有就是耶穌他.
耶穌他剛纔我已經講了.
就是他的錯誤的誓言.
就令到他的女兒.
甚至可能因為這樣的緣故而死亡.
參宣.
他的父母沒有好好地管教參宣.
沒有好好地教育參宣.
結果參宣一生沒有一個正常的婚姻關係.
最終也是以悲劇收場.
所以《詩詩記》不單單是記載某一個的詩詩.
他如何打贏仗.
這是我們很多弟兄姊妹的想法.
這個沒錯的.
但絕對不只是這樣.
更加原來在《詩詩記》裡面.
是很著緊我們要有家庭上的好的帶領.
在這方面是描繪得很淋漓盡致的.
我自己這次旅行也有個期望.
就是會和我的孫子.

$^{401}$基本上每一晚都可以一起分享,祈禱.
但很可惜做不到.
但也有幾晚我可以和他們一起.
有十多分鐘和他們講聖經故事.
和他們一起禱告.
很著緊宗教教育.
就是這個意思.
《詩詩記》的主題.
就是由《詩詩記》第三章第七節.
到第十六章第三十一節.
這就是描繪很多的詩詩.
由六個大詩詩來組成.
大家要看這個部分.
由南到北.
有《額陀列》.
有《以弗》.
有《底玻拉》和《基甸》.
由東到南就有《耶佛他》和《鏟宣》.
剛纔我說他們是一個循環.
這個循環出現了七次.
在《泡芳》裡面有.
就是他們犯罪.
離棄神.
於是受苦.
透過外族的入侵.
他們纔是反省.
於是他們向神祈求.
神祈求答應了他們.
於是他們得到拯救.
但過了一段時間又再犯罪.
你看到《泡芳》裡面有這個循環.
我們看看有哪些詩詞.
大家看到這個表.
我列出來了.
我們再留意.
原來聖經裡面.
詩詩紀的焦點放在幾個人身上.
詩詩有四個詩詞.
就是《底玻拉》.
《基甸》.

$^{441}$《耶佛他》和《鏟宣》.
但你有沒有留意到.
《底玻拉》之前.
有三個詩詞.
這些是很簡單的記載.
然後就帶出一個大的詩詞.
然後沒多久.
完了之後.
又有另外一些小的詩詞出現.
我們看看.
有《額陀列》.
《以弗》.
和《山格》.
你知不知道.
《山格》只用了一節經文.
還是兩節經文.
是描繪它的.
很少的.
沒有說什麼.
在《泡芳》裡面.
我們都會看到《額陀列》所面對的就是.
「米索波大米王古山尼撒」.
然後.
《以弗》所面對的就是.
「摩厄王伊基倫」.
《山格》來帶去面對非利士人.
然後.
我們就看一個大的詩詞.
就是《底玻拉》.
《底玻拉》帶領著以色列人.
戰勝了迦南王耶賓.
我對耶賓也有一個很大的體會.
為什麼呢?.
因為我小時候.
很喜歡喫一個麵包.
不知道你們有沒有喫過.
叫做耶賓.
就是椰絲加在一起.
很好喫.
好了.

$^{481}$《底玻拉》開始.
接著就是《基甸》.
《基甸》要面對米殿人.
接著就會有三個.
就是有兩個的詩詞.
其中一個呢.
有沒有留意到顏色很特別.
就是阿比米勒.
阿比米勒就是基甸那些孩子.
但是就帶來了很多家庭裡面的慘劇.
他要作王.
就將自己的兄弟.
就是同父異母的兄弟.
全部殺死.
這就是基甸的家庭生活做得不好而引致.
接著就有兩個.
就是《陀拉》和《賣爾》.
這些《陀拉》和《賣爾》.
根本是怎樣呢?.
他們就會帶去很小的描繪.
就不是一個很大的描繪了.
接著就出了《耶佛他》.
《耶佛他》之後就出《以彼贊》.
也是很短的描繪.
接著就出了《以倫》和《厄頓》.
也是三個人.
最後就出了《非利士人》.
叫他壓迫以色列人蔘宣.
所以你留意到.
原來《事實記》的作者.
就是這樣來描繪.
因為我跟大家查經.
雖然這些小的事實.
可能很多人都不會知道.
連我自己都不是很熟悉.
《陀拉》,《賣爾》,.
《以倫》,《以彼贊》,.
《厄頓》這些我們都會查的.
所以為什麼我們這次查經.
就不會像《港道》那樣.

$^{521}$有很多經文會用在崇拜《港道》裡面.
所以我們要從頭到尾去看.
好了.
讀《事實記》的故事的時候.
我們必定要著重.
他所描繪的是什麼呢?.
就是他們背若.
受欺壓,悔改.
然後拯救循環.
你留意到以色列人.
在《聖經》裡面就說到.
恆耶華眼中害民惡的事.
耶華就將他們交在外邦人的手裡.
讓他們來被欺壓.
以色列人就服侍外邦的王.
達到一個數目的年日.
可能是很多次.
多少年多少年.
以色列人向神來呼求.
耶華的靈降臨在這些拯救者的事實裡面.
然後外邦的王.
或者敵對的勢力被剷除.
以色列人所受的欺壓就結束了.
於是國中太平.
多少年就是這樣.
你看到這個圖就是這樣.
這樣的循環是有七次的.
而在這六個大的事實記載裡面.
只有額陀列.
是完全按照這個循環來記載.
越是晚期的事實.
就越講得壓要和簡略.
只有一兩句就講了.
被簡化的內容包括了.
講明誰是拯救者.
譬如好像是巴拉或者是提波拉.
在耶佛他的時代裡面.
以色列人的墮落就描繪到.
更加的歷代去詳細和具體.
而當我們向神呼求的時候.

$^{561}$甚至拒絕拯救他們.
在耶佛他和參宣的故事裡面.
是不再以平安作為結束.
到最後的一個事實就是參宣.
整個循環特徵就完全消失了.
參宣甚至以外邦欺壓者作為陪葬.
所以剛纔這幾樣東西很重要.
事實記一開始的時候就描繪了.
最後的結束是第17章到第21章.
指出以色列人他們背棄神.
信奉了其他的宗教.
又拜偶像.
17,18章就講到他們怎麼拜偶像.
而導致社會非常混亂.
事實記19到21章.
就講到當時的社會如何混亂.
道德如何淪亡.
就看到以色列人在內戰當中.
這就是事實記的結構.
好了.
我們看回當時的世代是怎樣.
原來當時的世代宗教非常混亂.
以色列人不要忘記.
他們離開埃及.
然後在曠野飄擾的時候.
他們是什麼呢?.
他們是遊牧民族.
但當他們定居在加拿大之後.
他們就落地生根了.
那種生活模式就大大改變了.
以前在曠野的時候.
朝陽晚策或者借住幾天.
可能民主保持他們走.
他們要走了.
他們是居無定所的.
但當他們去到加拿大的時候.
他們就穩定了.
於是慢慢慢慢就變成以農業為主.
他們在曠野飄擾.
不會種植的.

$^{601}$為什麼呢?.
種植要長時間的.
在加拿大他們就要開始種植了.
馬來亞也停止了來下降.
他們很需要雨水的資源.
受加拿大宗教的影響.
就是拜巴力.
巴力是生產力的神.
是農作業的神.
也會拜亞瑟拉.
如果你留意到巴力和亞瑟拉.
原來一早就有了.
直到耶綏別.
亞克王的妻子.
正式引入全國.
陷入敬拜巴力和亞瑟拉當中.
所以這些農民他們很想怎麼樣呢?.
得到雨水的資源.
於是就拜巴力和亞瑟拉.
就好像中國人.
三行就要拜老班.
老班師傅就是三行的輩祖.
有不同的神他們要拜.
打雨的那些出海的那些.
就拜什麼呢?.
媽祖.
即是娘媽.
戲劇呢?不知道拜誰呢?.
好像是華光.
每一個神都有獨特的地方.
他們認為.
這些是民間宗教.
所以他們就拜巴力.
他們拜巴力為什麼這麼嚴重呢?.
因為他們拜巴力的時候.
就會隨同巴力的境位.
他們會很淫亂.
一會兒我們會再講.
為什麼神要滅絕這些民族.
就是這個原因.

$^{641}$我們再看他們宗教上的混亂.
剛纔我們看了.
接著我們看軍事上的軟弱.
留意到他們進入加拿大之後.
他們要面對的是什麼呢?.
就是很多高大衛猛的民族.
當年他們見到民族這麼高大衛猛.
於是埋怨神我們進去就死了.
結果這群人就死在抗議那裡.
因為他們沒有這份信心.
但事實上他們欠缺了支援.
由於他們佔的地圖.
大部分是山區.
如果我們去過以色列都知道.
平原的地方不多.
山區的地方就很多.
這個山勢將支派分隔開.
身處比較遠的支派.
有戰事的時候.
其他的支派就很難支援他.
很簡單.
但這個支派是在最北面.
在敘利亞的旁邊.
如果他們有事的時候.
南面猶大支派是最強的.
怎樣上去幫他們呢?.
山長水遠.
所以他們就欠缺了支援.
還有一件事是什麼呢?.
就是他們不得同心.
沒有一個事實.
曾經帶領著他們.
統一了整個以色列民族去打仗.
甚至他們如果我們看時事記.
是會有內訌的.
打完勝仗回來.
有個支派叫以法聯.
到時我們會解釋.
為什麼以法聯會經常做這件事呢?.
去罵基地.

$^{681}$你做什麼?.
你打仗不叫我?.
就好像現在中東那樣.
特朗普去對付伊朗.
但其他那些.
法國甚至英國.
都說我們不會去幫忙的.
到贏的時候.
就出來做事了.
英國首相又說.
我會去做調停.
要拿個尾彩.
要分享那種榮耀.
你留意到原來沒有一個事實.
曾經帶領著12個支派的人.
一起去打仗.
沒有這個統一的戰爭.
直到王朝之後.
纔有短暫的戰爭.
可以統領到.
譬如大衛後來統一了南北.
跟著所羅門.
纔可以統一整個民族.
還有他們指手不攻.
除了事實底波拉的勝利之外.
其餘的戰士.
都沒有看伊斯蘭贏得新的領土.
他們只不過是阻擋了外敵.
只是去入侵而已.
我們看看事實的簡介.
這個表就比較小.
但是沒辦法.
所以大家可以向我拿這個powerpoint.
我們看看事實的分佈.
autolib.
是屬於什麼支派呢.
是屬於猶大支派.
記載在事實記第三章第七到第十一節.
記不記得我說.
事實記真真正正的本位.

$^{721}$就在第三章第七節開始.
他就是在米索多敗米那裡.
這些人來核制了他們四十年.
令他們很慘.
跟著就是爾弗.
爾弗就是汴雅敏支派.
第三章十二到三十節.
被摩壓人來控制.
但是後來我剛纔說.
四十年得到太平.
後來有八十年的太平.
跟著就是山加.
山加你沒有留意到.
沒有說什麼支派.
原來山加不是希伯來人.
不是神的子民.
但是神都用外邦人來幫助他.
解救以色列人.
就好像當日大衛來幫助基爾拉.
當時在基爾拉城裡的人不是希伯來人.
但是大衛就幫助他.
就是山加幫助他.
就是對付非利士人.
只是殺了幾百個非利士人.
跟著就是底波拉.
我們一直都是講底波拉.
沒有講巴拉.
底波拉是以法聯支派.
就是《史詩紀》第四到第五章.
他被迦南王來帶去對付.
迦南王耶賓.
四十年.
後來國中太平.
解決了迦南王耶賓.
這個很特別的.
特別姐妹一定要聽.
有兩個女士.
來帶去解救了這幫以色列人.
不單單是底波拉.
還有迦南王耶賓的元帥.

$^{761}$都是寫在那個女人手上.
跟著我們熟悉的就是基甸.
基甸就是嗎哪西支派的.
第六到第八章.
因為基甸那個年代被米電人來欺負.
後來基甸解決了他的時候.
就有國中太平四十年.
但是你想一想基甸有多少個兒子?.
七十個兒子.
你想一想他一定不是一個太太生的.
他自己是沒有遵照神的吩咐.
因為神的吩咐其實就是一夫一妻.
但是他其實有很多的妻妾.
所以七十個兒子.
帶來了家庭裡面的悲劇.
跟著就是陀拉.
陀拉就是以撒加支派的.
就是第十章第一到第二節.
記不記得?.
有些是間長的.
這個就是間長的.
只有兩節經文.
跟基甸有三章聖經描繪.
差很遠.
後來國中太平二十三年.
跟著就是埃爾.
埃爾就是基列支派的.
埃爾就是第十章第三到第五節.
後來就是國中太平二十二年.
但是埃爾也不差.
竟然有三十個兒子.
也挺厲害的.
跟著就是耶佛他.
耶佛他就是面對阿滿人的挑戰.
這個時候他們就是在第十一章到第十二章.
後來他戰勝了阿滿人.
就是六年的太平時日子.
以比贊就是猶大支派的.
就是十二章第八到第十節.
後來國中太平就是七年.

$^{801}$但是以比贊也是有三十個兒子.
跟著就是以倫.
以倫就是西布倫支派的.
第十二章第十一到第十二節.
都是隻有兩節的描繪.
國中太平就是十年.
跟著就是厄頓.
厄頓就是以法聯支派的.
第十二章第十三到第十五節.
國中太平八年.
但是都有四十個兒子.
你沒有留意到很特別的.
在《史詩篇》裡面會記載他們有多少個兒子.
最後就是參宣了.
參宣就是彈的支派.
有三章聖經描繪參宣的.
他們給非理士人來帶去對付.
去對抗非理士人.
說明國中太平是二十年.
如果我們真的這樣來帶去計算的話.
《史詩篇》的時代大概是四百一十年.
這是一個結論.
我們會看他們是如何欺壓了他們.
最後今天我們會提一個問題.
當我們研讀《史詩篇》的時候.
或者申命記載的一個問題困擾著我們.
就是神要以色列人滅絕迦南人.
某個程度是種族屠殺的.
為什麼要這麼殘忍呢?.
原來迦南的文化是很深遠.
是很強的影響力.
改變了人的信仰能力.
他們的宗教生活影響著周圍的民族和他們的文化.
由格里蒂和希臘移民去到迦南地的非理士人.
完全留意到.
其實某個程度是由歐洲而來的.
但已經受迦南影響.
用了迦南的名字雜俗.
所以可見當時迦南的文化影響得很厲害.
我們有沒有想過為什麼神母完全毀滅埃及?.

$^{841}$埃及是拜很多神的.
太陽神,甲蟲都是他們的神.
蛇也是他們的神.
你留意到實哉是針對他們的假神.
甚至埃及人虐待以色列人.
要逼他們做奴隸,剝削他們的基本人權.
甚至殺他們出生的男嬰.
但為什麼神仍然留埃及在這裡.
而沒有去毀滅他們所有的民族呢?.
反而地域相近的民族就是迦南的民族.
神就要嚴厲的對付.
為什麼呢?.
原來神的作為就是防止以色列人.
受迦南的文化影響.
而產生了信仰的背叛.
因為如果他被迦南文化影響的時候.
他就會離棄神了.
如果我們看埃及人都拜不少假神.
包括甲蟲,牛.
但是記不記得我剛纔說.
他們現在在迦南拜的是什麼?.
是巴力和阿斯塔羅.
這個不單單是農作業的神.
剛纔說是生產力的神.
是讓他們生而育女的.
在這些神的過程當中.
他就會放縱情慾.
男女之間是習交.
隨便和不同的異性發生性行為.
請記住.
不單是妓女.
可能是一些婦女.
他們懷孕.
那個小孩子屬於誰的?.
整個民族就會很混亂.
所以不單單影響他個人.
是影響他的後代.
影響整個民族.
所以神要毀滅這些迦南地的民族.
就好像毀滅索多瑪,俄摩拉一樣.

$^{881}$再者他們敬拜的儀式當中.
是包含了什麼呢?.
將兒童獻祭.
這是一個很嚴重違反道德的行為.
神內心等待他們迴轉.
但他們不肯悔改.
所以最終受到重重的刑罰.
我們看看《創世記》第15章13到16節.
這裡就記錄得很清楚.
原來神不是一下子就要滅絕他們.
而是要給他們四百多年.
好 謝謝.
耶和華對亞伯蘭說.
「你要的確知道.
你的後裔必居別人的地.
及又服侍那地的人.
那地的人要苦待他們四百年.
並且他們所要服侍的那國我要懲罰.
後來他們必帶著許多財物從那裡出來.
但你要在大丑受訴.
平平安安地歸到你列祖那裡被人埋葬.
到了第四代他們必回到此地.
因為阿摩利人的罪孽還沒有滿盈」.
「因為阿摩利人的罪孽還沒有滿盈」.
這裡很清楚說到.
原來那四百多年.
神是給他們機會去悔改的.
有沒有福金傳給他們呢?.
我們絕對相信有.
如果你看《阿摩司書》的時候.
是很清楚說到原來有不同的民族.
神都會有信息傳給他們.
問題是他們願不願意去接受.
但因為聖經主力是放在猶太民族裡面.
所以就沒有詳細說.
但我們很相信神是給他們機會.
但他們沒有.
所以後來神就要滅絕他們.
所以神滅絕他們是要保護以色列的文化.
和彰顯他們的公義.

$^{921}$是要賞善罰惡.
這次聖戰裡面獨一無二的行動.
在以後的歷史裡面.
我們再也看不到有這個要求.
要以色列人有類似這樣的種族屠殺.
是沒有的.
所以滅絕的行為.
是神公義的彰顯.
也是唯一一次.
今天就和大家看了.
史詩紀一些的背景.
我們一起留下來跟我們禱告.
天王多謝你給孩子有機會.
服侍弟兄姊妹.
聖誕求你來帶去眷顧.
透過你自己的話語.
給我們有更多的反思.
聖誕叫我們不要重鋪.
史詩紀裡面所記載的循環.
有時當我們在平靜安穩的時候.
我們沒有好好的親近你.
當我們遇到困難的時候.
我們纔會來帶去向你呼求.
聖誕叫我們不是利用你.
叫我們去愛你.
特別最重要叫我們能夠遵守你的法道.
叫我們能夠信服神裡自己的指引.
我們的生命才能得到完美.
求神裡自己來帶去使用我們的茶經.
讓你的話語成為我們的幫助.
也不單祝福在現場參加的弟兄姊妹們.
在YouTube裡面.
可能有些弟兄姊妹或朋友.
他們也會來帶去收看或收聽的時候.
你要同樣地祝福他們.
都願意我們的信息成為他們的鼓勵.
阿父求你來帶去眷顧我們.
聽耳者面前耶路教堂.
奉主德聖的名.
阿們.

\newpage



\section{路加福音 24:13-21}
\label{sec:031967}
\textbf{以馬忤斯路上的驚喜 (華)}~~~~曾興才牧師
\newline
\newline
連結: \href{http://www.church.com.hk/acms/content.asp?site=cdc\&op=show\&type=product\&code=031967\&layout=sermon}{\texttt{031967}} ~~~~ 語音日期: 2026-04-12
\newline
\newline
\hyperref[sec:031966]{\small{< < < PREV SERMON < < <}}
~
\hyperref[sec:index_chronic]{\small{[返順時目]}}
~
\hyperref[sec:031968]{\small{> > > NEXT SERMON > > >}}
\newline
\newline
路加福音 24:13-21
\newline
\begin{longtable}{cl}
\hline
\hline
章節 & 經文 (和合本修訂版)\\
\hline
24:13 & \begin{tabularx}{0.7\textwidth}{X} 同一天，門徒中有兩個人往一個村子去；這村子名叫以馬忤斯，離耶路撒冷約有二十五里。 \end{tabularx} \\ \\ \relax
24:14 & \begin{tabularx}{0.7\textwidth}{X} 他們彼此談論所發生的這一切事。 \end{tabularx} \\ \\ \relax
24:15 & \begin{tabularx}{0.7\textwidth}{X} 正交談議論的時候，耶穌親自走近他們，和他們同行， \end{tabularx} \\ \\ \relax
24:16 & \begin{tabularx}{0.7\textwidth}{X} 可是他們的眼睛模糊了，沒認出他。 \end{tabularx} \\ \\ \relax
24:17 & \begin{tabularx}{0.7\textwidth}{X} 耶穌對他們說：「你們一邊走一邊談，彼此談論的是甚麼事呢？」他們就站住，臉上帶著愁容。 \end{tabularx} \\ \\ \relax
24:18 & \begin{tabularx}{0.7\textwidth}{X} 兩人中有一個名叫革流巴的回答：「你是在耶路撒冷的旅客中，惟一還不知道這幾天在那裡發生了甚麼事的人嗎？」 \end{tabularx} \\ \\ \relax
24:19 & \begin{tabularx}{0.7\textwidth}{X} 耶穌對他們說：「甚麼事呢？」他們對他說：「就是拿撒勒人耶穌的事。他是個先知，在神和眾百姓面前，說話行事都大有能力。 \end{tabularx} \\ \\ \relax
24:20 & \begin{tabularx}{0.7\textwidth}{X} 祭司長們和我們的官長竟把他解去，定了死罪，釘在十字架上。 \end{tabularx} \\ \\ \relax
24:21 & \begin{tabularx}{0.7\textwidth}{X} 但我們素來所盼望要救贖以色列民的就是他。不但如此，這些事發生到現在已經三天了。 \end{tabularx} \\ \\
[1ex]
\hline
\hline
\end{longtable}
$^{1}$每一個人對驚喜都有不同的反應..
好的驚喜,正面的驚喜,會給人帶來加倍的快樂;.
負面的驚喜,會給人帶來加倍的挫敗..
你知道嗎,我們的神是常常給我們驚喜的神..
當耶穌基督復活的早晨,耶穌基督就給那幾個婦女一個驚喜..
他們來到墳墓那裡,有兩位天使告訴他們說,.
你們為什麼在死人當中找活人呢?.
祂不在這裡,已經復活了..
然後在馬太福音28章第10節告訴我們說,.
這兩個婦女,這些婦女又害怕又歡喜,.
就跑去告訴耶穌的門徒..
老張不是一個很羅曼蒂克的人,.
也不是一個很會給太太驚喜的人..
經過朋友們的調教之下,他開始送玫瑰花給老婆,.
而且每次都送十一朵..
有一天,太太在CHPT上找到花語背後十一,.
乃是代表一生一世..
心裡想,哇,老公真的下了功夫了..
但是她非常願意聽到她老公親口告訴她,.
這十一朵代表什麼意思..
就問她老公,老張就吱吱嗚嗚了半天,.
終於說,這是花店的促銷活動了,買十送一..
告訴旁邊的人說,今天神要給你一個驚喜..
是的,我們的神是會給我們驚喜的神..
今天我們來思想有兩個門徒,.
他們在非常憂愁,難過,灰心,失望的途中,.
耶穌給他們一個特別大的驚喜..
這個驚喜我稱它為以馬五師路上的驚喜..
我們一起來看《路加福音》第24章,.
我們從第13節,我們今天要講到36節,.
我們先看13節到21節..
《路加福音》第24章第13節說,.
同一天,門徒中有兩個人往一個叫以馬五師的村子去,.
那村子距離耶路撒冷約有11公里..
在《和爾本》裡面說25公里,其實不正確,正確應該是11公里..
他們彼此談論所遇見的這一切事,.
正在交談議論的時候,耶穌親自走近他們,與他們同行..
然後16節說,但他們的眼睛卻給矇蔽了,認不出他來..
耶穌說,你們走路的時候,彼此討論的是什麼事呢?.
他們就站住,面帶愁容..

$^{41}$18節說,一個名叫格留巴的回答他說,.
在耶路撒冷做客的難道只有你,不知道這幾天在那裡所發生的事嗎?.
他說,什麼事呢?.
他們說,就是拿撒勒人耶穌的事,他是先知,.
在神和眾人面前說話,行事都有大能..
20節說,我們的祭司長和官長竟把他交出去,判了死罪,釘在十字架上..
但我們樹萊所盼望要贖以色列的就是他..
不但如此,這些事發生到今天,已經是第三天了..
這是在伊瑪烏斯上的神給他們的驚喜..
我們先來做幾個非常重要的觀察..
這裡第一句話說,同一天..
同一天指的是當時他們走上伊瑪烏斯路上的時候..
那同一天代表什麼呢?是在之前早上,也就是復活天的清晨,.
婦女到墳墓裡面看見耶穌的墳墓是空的,.
然後天使對她們說,祂不在這裡,已經復活了..
最後她們跑回去告訴門徒,就有彼得和約翰跑到墳墓裡面,.
一樣看見墳墓是空的,而且耶穌不在那裡了..
這件事情就是同一天所發生的..
所以這兩個門徒不是十二個使徒當中的門徒,.
因為這裡面他不在這使徒當中..
他們就從耶路撒冷,可能是下午了,.
就開始朝向他們的鄉村,就是伊瑪烏斯來走去..
這裡告訴我們,這個伊瑪烏斯是一個非常小的鄉村,.
聖靈學者都無法確定它真正的地點在哪裡,.
只知道他們是距離耶路撒冷有十一公里,相等於大概七英里..
如果一般上走路的速度的話,大概需要三個小時左右,.
才能夠走到他們的村莊..
他們在那裡正在討論議論,非常激烈地討論當天,.
這幾天所發生的事情的時候,.
耶穌就忽然間出現在他們的身邊,與他們一同行走..
這是一般在猶太人的習俗,.
當你在同一條路上走向同一個方向的時候,他們都會接伴而行..
但是呢,他就問他們一句非常有趣的話,.
你們在這裡議論些什麼事情呢?.
聖經告訴我們說,其中一個名叫格留巴的這個人,.
格留巴這個人的名字只在這裡出現,.
所以聖經學者也不知道他真正到底是何方人物,.
只知道格留巴的名字乃是希臘人的名字..
我們可以想象他是外族人,.
可能是一個外地來的人來到耶路撒冷一同來慶祝玉月節的慶典,.

$^{81}$他也是跟隨耶穌的一個門徒,但是不在這個使徒當中..
耶穌就問他說,格留巴就回答耶穌說,.
"難道你是唯一在耶路撒冷做客的,不知道這幾天所發生的事情嗎?.
這幾天所發生的事情太轟動了,.
這位大家眾人以為是神的兒子彌賽爾竟然被釘死在十字架上,.
然後埋葬了,然後有婦女跟門徒們看見房子是空的,.
難道只有你不知道這些事情發生在這幾天之中嗎?".
耶穌非常有趣地,非常幽默地問他們說,.
"到底是什麼事呢?".
他們就講出大概四五個非常正確的事實,.
他說,這個納撒勒的耶穌他是個先知,.
這是大家公認了他是個先知,.
然後他說話,行事都大有能力,.
像就位裡面的摩西一樣,不但是如此,.
然後他說,我們的祭司長還有官長,.
包括這位文士法利薩人把他交給外邦人羅馬,.
被釘在十字架上,這都是非常正確的..
然後呢,他其實已經埋葬了多少天?.
第三天了,換句話說沒有指望了..
其實他是我們眾人所期待要來救贖以色列百姓脫離羅馬政府的捆綁,.
所以我們得自由的那位救主,沒想到他已經死了,.
而且死了三天,根本是沒有希望了..
所以這非常有趣,而且非常詳細,.
透過路加醫生記載在路加封印第24章,.
其他福音書,馬太,約翰沒有記載,.
馬可稍微提了一下,我們看看馬可封印第16章,.
稍微提這個事情,但沒有詳細的介紹,.
他說,這次以後,門徒中間有兩個人往鄉下去,.
就是以馬五師的鄉下,正走路的時候,.
耶穌用另外的形象,就是復活的形象向他們顯現,.
他們就去告訴其他人是誰,就是耶穌的門徒們,.
他那些人也不相信他..
所以我們可以一起來思想,在以馬五師的路上,.
耶穌給這兩個門徒帶來的驚喜,.
到底給我們上了什麼功課呢?三方面..
首先第一,耶穌靠近憂愁的人,讓他不再孤單,.
耶穌靠近憂愁的人,使他不再孤單..
你可以想象,這些門徒,好像其他的門徒,.
包括色門徒一樣,他們都期待要來的彌賽亞,.
要拯救以色列,脫離這麼多年的羅馬政府的挾持,捆綁和逼迫,.

$^{121}$要使他們得自由,重新復國以色列..
但是他們卻萬萬想不到,.
耶穌基督竟然在耶路撒冷被門徒出賣,.
結果被釘死在十字架上..
所以他們真的是非常非常的失望,傷心難過..
就在這個時候,耶穌基督就走到兩個門徒當中,.
與他們一同同行..
我們想象耶穌基督看到他們心裡的憂愁,.
到底他們心裡為什麼有憂愁呢?.
我猜他們的憂愁是什麼?錯誤期待的結果..
一起來講一遍,憂愁是錯誤期待的結果..
憂愁是錯誤期待的結果..
因為他們所期待的,但是所發生的,.
剛好與他們所期待的不能符合,.
因此他們就落入到憂愁,失望,難過當中..
很多時候我們生活當中也不是如此嗎?.
夫妻的關係,結婚以後做丈夫的,.
我們有錯誤的期待,.
希望我們的太太能夠滿足我們一切的需要,.
結果失望極了..
做太太也希望丈夫能夠了解你一切的需要,.
結果也非常的失望..
你可不知道只有耶穌基督真正瞭解你心中心靈深處的需要,.
只有祂能夠滿足你一切的需要..
因此錯誤的期待只有帶來心裡面的憂愁,.
甚至失望,甚至難過..
其實在這個之前,這些門徒們都早已經知道耶穌基督復活的消息,.
雖然沒有看見耶穌基督復活的樣子,.
都已經聽婦女說,也聽門徒說了,怎麼知道呢?.
我們繼續看《馬太入教福音22章》的22節..
而且我們中間有幾個婦女清晨到墳墓裡面去,.
使我們驚奇,她們清早到墳墓那裡去卻找不著她的身體,.
她們回來說看見天使顯現,天使說什麼?.
她活了!.
又有我們中間的幾個人,包括門徒彼得,約翰特別是這兩個人,.
到墳墓那裡去所看見的,正如婦女們所說的,只是沒有看見祂..
她們因為沒有看見覆活的耶穌,跟她們心中所期待的完全不能符合,.
因此她們落入到憂愁,失望,難過的觀景裡面..
沒有想到,耶穌竟然給她們一個這麼大的驚喜,.
出現在她們的旁邊,跟她們一起行走,走向野馬霧獅兩三個小時的路程當中..

$^{161}$親愛的網友姐妹,你知道嗎?.
耶穌會在你意想不到的時候給你一個驚喜..
在旁邊說,耶穌在你意想不到的時候給你一個驚喜..
這兩個門徒真的是非常意外的,.
從來沒有想到這位復活的耶穌竟然跟她們接伴而行,.
一同走向野馬霧獅的路上,這是何等大的驚喜..
親愛的網友姐妹,你可不知道我們的神會給你驚喜嗎?.
我自己個人經歷神給我的驚喜在2024年5月份的時候,.
我母親在睡覺的時候就被主接回天家,.
所以我趕回去馬來西亞,參加追思禮拜,按贊禮拜之後,.
我就從吉隆坡坐飛機回到舊金山..
在飛機上的時候,我忽然間想到母親的離去,.
想到95年我父親也走了,.
頓時感到心中非常難過,感覺到這裡好孤單,.
好像沒有父母親的孤兒一樣,.
真的非常難過坐在那裡..
忽然間神就給我一個驚喜,讓我想起詩篇27篇第10節,.
我父母離去我,耶和華必收留我..
我那時候馬上感覺到好像耶穌就坐在我旁邊,.
祂說你父母有一天會離開你,真的是離開我了,.
但是耶穌說我會收留你..
家裡的弟兄姊妹,你成佛也一樣經歷了神給你的驚喜呢?.
如果還沒有的話,告訴你,你可以期待,.
神也會給你一個驚喜,阿們..
告訴旁邊的人,神也會給你一個驚喜..
第二,我們看見耶穌基督不但是靠近憂愁的人,.
也體恤軟弱的人,而且完全接納..
一起講一遍來,耶穌體恤軟弱的人,完全接納..
你現在想想看,這兩個門徒可能以前也聽過耶穌的講道,.
也看過耶穌的神蹟奇事,現在聽了父女從墳墓回來,.
告訴他們天使顯現,告訴他們說他不在這裡,已經復活了..
也聽到這兩個門徒,包括耶穌的門徒使徒們,.
告訴他們說耶穌基督已經不在墳墓裡面,復活了..
他們心裡又怎麼樣?還是不能相信..
你看看耶穌基督怎麼看待他們信心的軟弱呢?.
我們繼續看第二十二章,二十四章的二十五節..
耶穌說:無知的人啊!.
先知所說的一切話,你們心裡信得太遲鈍了..
基督這樣受害,然後進入他的榮耀,不是應當的嗎?.
二十七節:於是他從摩西和眾先知起,.

$^{201}$把所有關於自己的經文都給他們做什麼?.
解釋明白了..
每次讀到這裡,我好羨慕這兩個人啊!.
你可以想象最偉大的老師,從最偉大的聖經,.
跟他們講解最偉大的神救贖的計劃,哇!.
這一席的談話,比你上過十年二十年的神學課還要寶貴..
耶穌說:你們無知的人啊!.
你們的心為什麼信得這麼遲鈍呢?.
你們心裡信得太遲鈍了..
信得太遲鈍,原本的意思是你們心裡信得太慢了..
他們是信,信得太遲鈍,太慢了..
已經告訴你那麼多事情,已經看見那麼多神蹟奇事,.
為什麼你們還遲遲不能相信呢?.
心裡信得太遲鈍了..
其實早在耶穌基督來到世上之前,受苦受害之前的七百多年前,.
所以您透過先知,不但先知,眾先知,已經告訴他們要發生的事情..
這是後來彼得這裡所做的見證..
看看彼得前書第一章第十節第十一節,.
彼得說論到這救恩,耶穌基督的救恩,.
那預先說你們要得恩典的什麼?.
眾先知,不止一個先知,眾先知啊!.
早已詳細地尋求考察..
就是考察在他們心裡基督的靈,就是聖靈,.
預先證明基督受苦難,後來得榮耀..
基督受苦難,後來得榮耀..
這早就在眾先知,包括以賽亞,包括和稀啊,.
眾先知已經告訴你們,折耳折到世上他要受苦受害,被釘死,.
第三天從死裡復活,得榮耀..
換句話說,在榮耀的冠冕之前,必須要有荊棘的冠冕..
但是他們因為心裡對耶穌基督的期盼,.
跟耶穌基督要來所做的工作不一樣,.
所以他們心中裡面就感覺到失望,灰心..
所以耶穌說,你們心裡為什麼性德那麼遲鈍呢?.
耶穌在伊瑪烏斯路上給他們的驚醒,.
是告訴他們一個好消息,什麼好消息呢?.
受苦的基督復活,得榮耀啦!.
阿們,一起講一遍來!.
受苦的基督復活,得榮耀啦!.
耶穌基督必須為我們眾人的罪,受苦受害,為我們眾人的罪,.
丁死在十字架上揹負我們的重擔,擔當我們一切的憂患,.

$^{241}$為我們死,買熟我們一切的罪孽過方,第三天從死裡復活..
其實在這個之前,耶穌基督不但一次,兩次,至少三次聖靈記載說,.
他已經預先,還沒有上十字架之前,已經預先告訴他十二個門徒要發生的事情..
他們心裡真的是遲鈍了,信得太慢了..
我們看看第三次耶穌跟他們所發的預言,.
在馬太福音第二十章第十七節到十九節,.
耶穌上耶路撒冷去的時候,就是他受死的地方,.
在路上把十二個門徒帶到一邊,他是不讓其他人知道,.
把十二個門徒帶到旁邊,對他們說什麼呢?.
看哪,我們上耶路撒冷去,做什麼呢?.
人子,就是耶穌基督,要被交在祭司長和文士,他們要怎麼樣?.
釘他死罪.清不清楚啊?非常清楚,要釘他死罪..
然後十九節又交給外辦的羅馬人將他戲弄鞭打,釘在十字架上..
一起講來,第三日他要復活.清楚嗎?非常清楚啊..
他要被釘死在十字架上,第三日怎麼樣?他要復活..
不只講了一次哦,講了兩次哦,這裡說至少三次,.
這裡是第三次耶穌清楚地告訴他們將要發生在他身上的事情..
十二個門徒就好像這兩個迷茫無事路上的門徒一樣,.
他們心裡性格太遲鈍了,太緩慢了..
親愛的基督徒姐妹,我們是否也是像他們一樣呢?.
性格太慢,太遲鈍了.這裡很清楚地告訴我們,.
腓立比書第四章十三節,我靠著他嫁給無力量的凡事都能做,你信嗎?.
如果你信了,為什麼你現在還常常靠著你自己的力量去做身上要你做的事情呢?.
這裡很清楚地告訴我們說,馬太福音二十八章十七節十八節,.
天上地下所有的全面都給你賜給我的,所以你們要去使萬民做我的門徒,.
我就藏你們同在,製造世界的末了..
那我們為什麼還有這麼多人遲遲不敢走出我們的舒適圈,.
能夠為神做見證去宣教呢?我們是否信心也性格太遲鈍了?.
求神幫助我們,加點給我們信心,讓我們快快地信,阿們!.
其實耶穌基督受苦得榮耀,彼得非常清楚明白這件事情..
我們這個四月份所背誦的經文,大家記得嗎?.
菲利比書第二章二十一節,這裡說什麼呢?.
你們蒙昭原是為此,因基督也為你們受了苦,.
給你們留下了榜樣,叫你們跟隨他的腳蹤行..
耶穌基督給我們留下了榜樣啊,我們蒙昭原是為此啊,.
因為基督已經為你們,也為你們受了苦,.
好讓我們可以跟隨他的腳蹤行..
因此當你真的願意為神的緣故,基督的緣故,.
為傳福音做見的緣故,為主受苦的話,那是何等榮幸的事情?.
主必定與你同在,必定加給你力量,哈利路亞!.

$^{281}$因為受苦復活的耶穌基督已經復活得榮耀了,阿們!.
然後最後第三,耶穌不但是靠近憂愁的人,.
不但體恤軟弱的人,耶穌也喜悅可悟的人,.
重新得力,一起講一遍來,耶穌喜悅可悟的人,重新得力..
哇!你可以想象這兩個門徒走在路上的時候,.
耶穌走在旁邊,跟他們從舊約聖經,摩西啊,.
還有眾先知志裡面,一一給他們講解明白啊..
我相信耶穌肯定是從篡世紀就告訴他們神的心意,.
亞當夏娃犯罪後預備的一個救主,帶罪的羔羊,.
一直到以賽亞書53章,祂為我們眾人的罪壓在祂的肩膀上,.
因為我們都向羊走迷的路,.
耶和華神將眾人的罪孽都壓在祂的身上..
哇!為他們講解明白!.
你可以想象他們心中馬上信心油然升起,哇!真的太棒了!.
原來聖經裡面這麼清楚,早已經告訴他們,.
彌撒亞來到世上,那榮美的計劃,.
他們心中就更加渴慕多知道關於耶穌基督的事情..
接下來我們看看最後一段,24章..
24章我們從第28節到32節..
他們走近快要去的村子,一馬五十的時候,.
耶穌基督好像還要往前行..
耶穌就是故意的,好像要繼續往前走..
他們就強留他,說什麼天晚了,太陽下山了,.
請我們同住吧!.
他們眼睛開了,才認出是耶穌來,祂卻從他們前面不見了..
32節說什麼,一起來講..
他們彼此說,在路上,祂對我們說話,給我們解釋聖經的時候,.
我們的心不是火熱的嗎?記得嗎?.
當在路上,祂為我們從舊約,聖經,一直到眾星之,為我們來解釋,使我們明白所發生的事情,.
我們心裡是否好像一團火在那裡燃燒呢?.
聖經裡面講,我們的心是火熱的,.
在原文的意思是,我的心是在我裡面繼續燃燒..
不是因為他們有了神的話語,不是因為他們聽了神的話語,.
而是耶穌基督為他們清楚地解釋,使他們明白神的話語,哈利路亞..
所以我們可以下一個這樣的結論說,明白神的話語會使你信心恢復..
一起講一遍來,明白神的話語使你信心恢復..
換句話說,你有神的話語還不夠,背誦神的話語還不夠,.
你要怎麼樣?你要明白神的話語,明白神的心意..
如此你的信心纔會被恢復,你的心纔會像一團火在裡面燃燒啊!.
這是為什麼我們牧者一直鼓勵你參加這禮拜六的專輩日,還有諸如專輩課程..

$^{321}$好,讓有人替你講解,解釋聖經,.
但你明白神的話語的時候,你忍耐的信心再次能夠被恢復起來..
你那忍耐的心,冷卻的心再次像火一樣在裡面燃燒啊!.
這是為什麼鼓勵你每個禮拜要參加諸如專輩來替牧師講解聖經,.
讓你明白神的話語,太重要了,明白神的話語..
當他為我們解釋聖經的話的時候,我們的心裡不是火熱的嗎?.
家裡的弟弘姐妹,你的心現在是否有點熱呢?.
你的心是否現在很火熱呢?.
如果還不是讓求生幫助你,.
那你願意下功夫,花時間,上專輩課程,好好地考察,明白聖經,.
當你願意明白聖經,願意明白神的心意的時候,.
你的心就再次能夠火熱起來,.
不再是一個不冷不熱的基督徒,因為你明白主的心意..
當耶穌基督一一地為他們解釋,使他們明白神話的時候,.
這兩個門徒從非常灰心,難過,憂愁到非常興奮,.
到什麼地步呢?我們看看第三十三節到三十六節..
他們就歷史起身,回耶路撒冷去,.
正遇見十一個使徒和他們的同仁聚集在一處,.
說主果然復活,已經顯給西門看了..
三十五節,兩個人就把路上所遇見的和波餅的時候,.
怎樣被他們認出來的事都述說了一遍..
正說話的時候,耶穌親自站在他們當中,.
說什麼一起說來,「願你們平安」..
回到三十三節,三十三節說什麼呢?.
他們歷史起身,馬上起身,你知道那時候是什麼時候了嗎?.
天晚日落了,他們強留耶穌基督住下來..
耶穌基督是考驗他們的信心,故意要往前走,.
但是被他們強留下來,好吧,我就跟你們一起留下來..
那時候已經天晚日落了,通常猶太人晚上不出門的,.
當時沒有像我們這裡這麼多路燈,特別走寬野路,.
你看從伊瑪烏斯大概十一公里的路要走回到耶路撒冷至少三個小時,.
晚上走更困難,可能三四個小時的時間,多麼危險的事情,.
一般猶太人是不會做的..
但是這裡他們就歷史離開,說我們要回去耶路撒冷,.
告訴死一個使徒,還有他的門徒們說他已經復活了,哈利路亞..
這裡講得很清楚,說他已經復活了,其實在原文的意思是說,.
這是真的,it is true,這是真的,耶穌基督果然復活了..
看到他們的信心不斷被恢復,我們也看到他們因為復活的耶穌使他們怎麼樣?.
重新得利,阿們..
剛纔我們說復活的耶穌會使你重新得利,.

$^{361}$所以復活的耶穌會使你重新得利..
他們回到耶路撒冷,一五一十地就把他們所經歷的路上耶穌基督如何為他們解釋明白聖經..
到家裡面的時候,耶穌那時候拜餅,其實那時候拜餅不是在他最後晚餐,.
設立聖餐的時候,因為那時候這兩個門徒不在其中,只有十二個門徒在其中,記得嗎?.
那是耶穌可能在無病惡愈的神蹟的時候,.
耶穌習慣性地拿起餅來,然後注視他們,.
看到這個舉動的時候,眼睛就被打開了..
當時耶穌是故意不讓他們知道他是誰,故意讓他們眼睛矇蔽,.
讓他們眼睛被打開了,認出是耶穌..
耶穌怎麼樣? 耶穌就立刻在他們面前,聖經說:消失了,不見了..
乃是證明說這位復活的耶穌的身體,乃是一個榮耀的身體,.
不受時間,不受空間,不受物質所控制的..
他可以從以馬無師消失了..
當這兩個門徒開始告訴一個使徒的門徒說:耶穌就真的復活了以後,.
講了一五一十非常詳細地講述他們路上所遇到的驚喜之後,.
聖經告訴我們說,耶穌親自站在他們中間,對他們說:願你們平安..
我相信這是十二個,是一個門徒跟其他門徒在一起給他們一生當中最大的一個驚喜..
他們到墳墓裡面看見墳墓是空的,.
但是這兩個門徒卻告訴他們說,其實他已經顯現給西門看..
請問這兩個門徒怎麼知道西門彼得去看過這個空墳墓,然後耶穌顯現給他看呢?.
因為是耶穌親口告訴他的嘛..
記得嗎?耶穌花三個多小時的時間在路上為他們講解所發生的事情..
所以他們因為復活的主,他們就怎麼樣?在神面前重新得力了..
我們看到彼得的生命,因為耶穌基督的復活,完全被翻轉,.
從一個三次不認主的彼得,後來變成教會重要的領袖..
我們看看他寫的彼得前書..
彼得前書第一章第三節第四節,彼得自己做見證..
他說,願頌讚歸於我們的主耶穌基督的父神,.
他承照自己的大憐憫,.
藉助什麼?基於耶穌基督從死裡復活,彼得親自見證了..
從死裡復活,重生了我們,叫我們有什麼?活潑的盼望..
Living hope,不是一個死的盼望,可以得著,.
那不能血壞,不能玷污,不能衰殘,為你們存留在天上的基業..
哇!彼得說願頌讚,詐美歸於我們在天上的父神,.
因為他藉助耶穌基督,死裡復活,才重生了我們,.
給我們活潑的盼望,這是唯一的一個宗教..
救主耶穌基督是從死裡復活的,.
世界上其他四大宗教沒有一個教主是從死裡復活的,.
你可以還在他們的墳墓當中找到他們的屍體,.
只有耶穌基督從死裡復活,使我們重生,有復活,而且有活潑的盼望..

$^{401}$因此彼得繼續鼓勵,.
當時第四集,那些跟隨耶穌基督的門徒,.
在水深火熱,隨時遭到逼迫殺害的這些信徒當中,.
怎麼鼓勵他們呢?.
第六節,第七節,彼得前述第一章,.
他說:因此你們大有喜樂,.
但如今在百邦的失戀中,暫時憂愁..
為什麼?因為有活潑的盼望,有在天主嘛..
叫你們的信心既被試驗,就比那被火試驗的仍然能壞的金紙更顯寶貴,.
可以在耶穌基督顯現的時候,.
祂再來的時候,得到稱讚,榮耀和尊貴..
他鼓勵信徒:你看,耶穌就從死裡復活了,.
你們是大有喜樂的,雖然現在在百邦的失戀中,暫時憂愁..
換句話說,不要憂愁太久,暫時憂愁..
為什麼?因為當耶穌基督顯現的時候,.
你們的信心經過火的考驗,比那個金紙更加寶貴..
當耶穌基督第二次再來的時候,你們得到稱讚,榮耀和尊貴..
我們生活在這麼自由,暫時還不受到逼迫的美國,.
我們很難體會彼得這裡所講的,.
在百邦的失戀中,暫時憂愁,要以為是大喜樂..
但是在那些不是那麼自由的國家,.
傳福音的牧者傳到基督徒們,.
他們的確明白彼得所講的這些話..
上個禮拜我剛從緬甸回來,.
我們在緬甸的分堂的牧者,通黑牧師非常有心,.
召集了五十位牧者跟領袖從緬甸各地來到陽光首都,.
有兩天的牧者訓練會..
對很多牧者來說是第一次他們參加這麼大的一個聚會,.
對他們來說是一個很好的一個驚喜..
他們盼望我們明年再去,要三天兩夜,四天三夜,.
他們覺得這個太棒了..
你可以想象這些牧者在緬甸服侍那麼困難..
我們看看這些牧者,我們訓練的情況..
然後下一張圖畫,.
我們有五十位牧者跟領袖,我們一起合照..
你看在緬甸的服侍,不像在美國服侍那麼容易,那麼輕鬆..
因為緬甸百分之九十是佛教的一個國家,.
而且佛教在那裡是有特殊的地位,特殊的待遇的..
再加上2021年政變的時候,.
現在緬甸在軍事政府的管轄之下,.

$^{441}$因為內戰的緣故,許多的教堂因此被毀壞,被火焚燒..
所以他們在緬甸從服役,不是那麼容易,不是那麼簡單的..
所以很多的挑戰隨時會面臨這些危險,甚至逼迫..
當我們在一起聚集的時候,我告訴他們說,.
下個禮拜,就是上個禮拜,就是復活的主日啊..
因為耶穌基督復活了,所以彼得可以說,你們大有喜了..
你們在百般的失戀中暫時憂愁..
因為你們信心經過失戀以後,要比寶貴的金子更加寶貴..
要在耶穌基督再來的時候,在現現的時候,.
要得到稱讚,榮耀與尊貴,Amen..
親愛的弟兄姊妹,耶穌基督真的復活了..
就好像這兩個門徒告訴,是一個使徒告訴門徒們說,.
這是真的,It is true,他真的復活了..
所以使他們能夠重新得力,.
也是緬甸的信徒們,牧者們能夠重新得力,.
知道耶穌基督復活了,有一個活潑的盼望,.
不是一個死的盼望,是一個有確擊的盼望,.
有一天耶穌基督要再來..
我們的信心經過失戀,要比金子更加的寶貴..
親愛的弟兄姊妹,你真的相信耶穌基督從此裡復活了嗎?.
只有一些人,你真的相信耶穌基督從此裡復活了嗎?.
我們剛剛慶祝了復活,.
你如果真的相信耶穌基督從此裡復活的話,.
我們就不再過一個軟弱,.
沒有能力,不忍不忍的基督徒生活了..
我們要重新得力,因為復活的大能要硬逼在我們身上,.
我們要每天向仇敵來誇勝,.
而且我們要得勝,保證說得勝有餘..
無論在什麼環境,在什麼情況之下,.
我們都可以靠著耶穌基督復活的大能大聲的宣告,.
是的,耶穌基督已經復活了,.
祂復活的能力如今靠著聖靈住在我們的心中,.
因此我們知道我們隨時都可以重新得力,.
繼續為主而活,為主做美好的見證,.
並且期待我們的耶穌隨時在你意想不到的時候也會給你一個驚喜..
Amen!我們一起禱告..
何等的寶貴!主啊,禰允許禰的愛子耶穌基督為我們受苦受害,.
因此在十字架上第三天重新復活,戰勝過黑暗,死亡,魔鬼,撒旦一切的權柄,.
讓我們能夠重新得力..
是的,真的,It is true,耶穌基督已經復活了,.

$^{481}$祂就活在我們每一個人的心中,因為我知道我的救贖者祂活著,.
因此我的心靈不再沉睡,我的心也不再憂愁,.
因為我知道我的救贖主活著,我的生命每一天就不再一樣..
謝謝我們的主,聽我們的禱告,奉告主耶穌基督,寶貴的民囚,Amen!.
\newpage



\section{約翰福音 21:1-11}
\label{sec:031968}
\textbf{開往水深之處}~~~~吳天凱牧師
\newline
\newline
連結: \href{http://www.church.com.hk/acms/content.asp?site=cdc\&op=show\&type=product\&code=031968\&layout=sermon}{\texttt{031968}} ~~~~ 語音日期: 2026-04-12
\newline
\newline
\hyperref[sec:031967]{\small{< < < PREV SERMON < < <}}
~
\hyperref[sec:index_chronic]{\small{[返順時目]}}
~
\hyperref[sec:031969]{\small{> > > NEXT SERMON > > >}}
\newline
\newline
約翰福音 21:1-11
\newline
\begin{longtable}{cl}
\hline
\hline
章節 & 經文 (和合本修訂版)\\
\hline
21:1 & \begin{tabularx}{0.7\textwidth}{X} 這些事以後，耶穌在提比哩亞海邊又向門徒顯現。他怎樣顯現記在下面。 \end{tabularx} \\ \\ \relax
21:2 & \begin{tabularx}{0.7\textwidth}{X} 西門‧彼得、叫低土馬的多馬、加利利的迦拿人拿但業、西庇太的兩個兒子，和另外兩個門徒，都在一起。 \end{tabularx} \\ \\ \relax
21:3 & \begin{tabularx}{0.7\textwidth}{X} 西門‧彼得對他們說：「我打魚去。」他們對他說：「我們也和你一起去。」他們就出去，上了船；那一夜並沒有打著甚麼。 \end{tabularx} \\ \\ \relax
21:4 & \begin{tabularx}{0.7\textwidth}{X} 天剛亮的時候，耶穌站在岸上，門徒卻不知道他是耶穌。 \end{tabularx} \\ \\ \relax
21:5 & \begin{tabularx}{0.7\textwidth}{X} 耶穌就對他們說：「孩子們！你們有吃的沒有？」他們回答他：「沒有。」 \end{tabularx} \\ \\ \relax
21:6 & \begin{tabularx}{0.7\textwidth}{X} 耶穌對他們說：「你們把網撒在船的右邊，就會得到。」於是他們撒下網去，竟拉不上來了，因為魚很多。 \end{tabularx} \\ \\ \relax
21:7 & \begin{tabularx}{0.7\textwidth}{X} 耶穌所愛的那門徒對彼得說：「是主！」那時西門‧彼得赤著身子，一聽見是主，就束上外衣，跳進海裡。 \end{tabularx} \\ \\ \relax
21:8 & \begin{tabularx}{0.7\textwidth}{X} 其餘的門徒因離岸不遠，約有二百肘，就坐著小船把那網魚拉過來。 \end{tabularx} \\ \\ \relax
21:9 & \begin{tabularx}{0.7\textwidth}{X} 他們上了岸，看見那裡有炭火，上面有魚和餅。 \end{tabularx} \\ \\ \relax
21:10 & \begin{tabularx}{0.7\textwidth}{X} 耶穌對他們說：「把剛才打的魚拿幾條來。」 \end{tabularx} \\ \\ \relax
21:11 & \begin{tabularx}{0.7\textwidth}{X} 西門‧彼得就上船，把網拉到岸上，網裡滿了大魚，共一百五十三條；雖然魚這樣多，網卻沒有破。 \end{tabularx} \\ \\
[1ex]
\hline
\hline
\end{longtable}
$^{1}$在Pleasanton跟大家分享神的話語.
在過去幾個禮拜我身體有點軟弱.
每次以為好了.
但是手尾長還在拖.
所以請大家包容.
今天開始的時候我想問大家一個問題.
你不要開聲回答.
為什麼不要開聲回答呢?.
因為你的答案可能你說了出來會覺得尷尬.
所以你不要回答出來.
心裡面回答就可以了.
我問你.
你每個禮拜都來教會聽道.
一年你又聽了52篇道.
兩個禮拜放假你一年聽了50篇道.
你記得有多少呢?.
再簡單一點.
上個禮拜是復活主日.
你記得上個禮拜的編導說了些什麼嗎?.
你記不記得誰在那裡說道?.
很多時候我在想.
如果我們一個禮拜可以記得一句.
或者一個概念.
一年我就記得52句52個概念.
我們再進一步來說的時候.
如果我們有52個的思想在我們心裡面.
我們去實踐我們去走出來的時候.
我們生命會有多大的改變呢?.
我們每天應該做靈修的.
如果我們每天靈修的時候.
可以記得一樣東西的時候.
我們屬靈生命會有多大的長進呢?.
說到這裡大家可能覺得會灰心.
聽了聽了很多.
但是靈命一點長進都沒有.
我們年紀比較大一點就說.
我是個老人了.
我沒什麼記性.
所以聽了一些東西.
這裡聽了進去接著出去了.

$^{41}$我裡頭沒什麼的.
我年紀老了.
老人之外記不了多少.
但是自己知道自己的事.
很多時候這個只不過是一個藉口.
但是有些時候.
我們面對一些事情的時候.
突然間靈光一閃.
茅塞頓開.
好像突然間通了.
明白過來了.
原來我聽了這麼久.
牧師說了這麼久.
就是這樣的東西.
我明白了裡面的意思.
原來就是這麼簡單.
為什麼這麼久我都不明白呢?.
為什麼我這麼蠢呢?.
一燈一光.
什麼事都清楚了.
你有沒有這樣的經驗?.
我有.
有時候看聖經的時候.
看來看去就是這個意思.
但是過了一段時間.
一年五年十年.
再回去看的時候.
突然間明白了.
原來裡面是這樣的意思.
可能說到這裡.
牧師你這麼說.
我也覺得舒服一點.
原來我不是自己是這樣的.
我跟你說.
不單止我是這樣.
不單止你是這樣.
耶穌的門徒和我們都差不多.
你不要以為.
約翰彼得,聖約翰,聖彼得.
他們是聖人.

$^{81}$特別又可以什麼.
其實他們和我們差不多.
你看看聖經就可以看出來.
所以門徒和我站在一起的時候.
我也不怕認為.
我自己也不記得很多東西.
我們看看.
在馬太福音16章20節說.
耶穌帶門徒去以色列的最北邊.
一個地方叫蓋薩尼亞菲勒比.
當他去到的時候.
耶穌問門徒.
人家說我是誰.
門徒說.
我聽到有些人叫你施浸約翰.
有些人說你是先知以利亞.
有些人說你是先知耶利米.
這些都是偉大的人.
耶穌再問他們.
你們呢.
人家說我是誰.
你們說我是誰.
彼得爭著說.
大聲地回答.
你是基督.
你是永生神的兒子.
耶穌看著彼得說.
厲害.
100分.
你拿了.
你真是有福氣.
接著我們看到.
當下耶穌祝福門徒.
不可對人說他是基督.
從此耶穌才子是門徒.
他必須上耶路撒冷去.
受長老,祭司長,文士許多的苦.
並且被殺.
第三日復活.
說到這裡.

$^{121}$耶穌突然嚴肅地吩咐門徒.
你們聽著.
從現在開始.
你不要隨便跟人說我是基督.
千萬不要亂說.
明不明白.
我和你說清楚.
我們現在要下去耶路撒冷.
下去的時候.
我會受很多的苦.
那些長老,文士,祭司長.
會一齊來整我.
讓我受盡折磨.
最後他們還要殺死我.
但是不用怕.
第三日之後.
我就會復活.
耶穌在這裡.
說得這麼正經的說話.
在祝福那些門徒.
門徒.
有沒有聽到?.
他們有沒有記得?.
耶穌和他們說了那些偏道.
他們有沒有記得?.
一點都沒有.
好像耳邊風一樣.
和我們差不多.
他們走到耶路撒冷的時候.
一邊走一邊聊天.
這12個門徒.
他們在聊什麼?.
他們說我們一定要向耶路撒冷去.
這是首都.
耶穌說他可以走神蹟,騎士.
說不定他可以做皇帝.
他做了皇帝之後.
誰做了宰相? 誰坐在他右邊? 誰坐在他左邊?.
哇!.
這些門徒.

$^{161}$就是想著這些.
耶穌是不是說他要去耶路撒冷做皇帝?.
沒有啊! 耶穌說得很清楚.
他下去要受苦被殺.
第三日復活.
門徒聽了之後.
有沒有聽進耳裡? 有沒有聽進腦裡?.
一句都沒有.
你覺得我現在好不好一點?.
不記得說了.
你會不會覺得好一點呢?.
甚至.
在耶穌被賣的晚上.
耶穌再一次的來教他們.
和他們說.
愛邦人有君王為主來治理他們.
那些掌權的人.
要叫他們做恩主.
但是你們不要學他們的樣子.
你們裡面誰要做大的.
就要做年幼的.
你做首領的.
你就要去那裡服侍人.
然後耶穌就拿一盆水.
給那些門徒洗腳.
做一個好的榜樣給那些門徒來學習.
你不知道.
到最後晚餐的時候.
我們看見這班門徒.
他們都不明白耶穌所說的話.
甚至到耶穌復活之後.
他們看見耶穌受死.
受難,受苦,受死.
到復活之後.
他們都不明白耶穌說.
我第三日會復活的.
他們都記不住.
他們的腦袋好像塞住了一樣.
所以你不記不住.
上個禮拜那個牧師說了些什麼.

$^{201}$你記住.
彼得,約翰,雅各.
他們都不好了太多.
這就來到另外一個問題.
就是聖經的說話.
很重要是神的話.
耶穌在跟門徒說話.
不是說閒話.
耶穌教我們不是在說閒話.
我們怎麼辦呢.
神的話語給我們.
我們聽完之後.
我們不記得.
我們怎樣好呢.
記住.
神知道我們每個人都是軟弱的.
所以祂不停地提醒我們.
會用經驗,經歷.
一次又一次地提醒我們.
來教導我們.
聽到就好像吃飯一樣.
不是一聽就有的.
是一個過程來的.
完了之後你請我喝茶.
去喝茶的時候就叫一個蝦餃.
蝦餃來到的時候你怎樣吃呢.
你放進口裡.
然後慢慢召它.
召完之後你吞進肚子裡.
吞進肚子裡之後怎樣呢.
消化然後才吸收.
肚腩才不會大.
是不是.
記不記得.
如果你不記得請我去吃.
不單止蝦餃.
請我吃燒賣.
請我吃叉燒包.
過程怎樣呢.
都一樣要放進口裡.

$^{241}$召啊召啊召啊.
召完之後吞進肚子裡.
吞進肚子裡之後怎樣呢.
慢慢地消化吸收.
肚腩就再大一點.
OK.
聽到都是一樣的.
你聽完之後.
進去了.
你就要反覆思想.
在那裡消化.
想啊想啊想啊.
吸收了.
那你.
可以做出來.
你實踐.
你去應用.
那你生命呢.
就會成長了.
我們看看耶穌就是用這樣的方法.
一次又一次又一次地去教導門徒.
讓他們循循善誘.
慢慢地來教導他們.
讓他們學習可以命運.
上個禮拜我們紀念耶穌的復活.
2000年前耶穌復活之後.
門徒很開心.
因為他們看到耶穌.
知道他在他們中間.
但是過了兩三個禮拜之後.
好像無聲無息.
耶穌顯現了幾次.
之後就沒什麼事情發生.
門徒那時候在耶路撒冷島.
這個不是他們的地頭.
人生路不熟.
所以整天坐在那裡沒事做.
還有11條大漢.
有吃粟.
吃飯.

$^{281}$住宿吃飯這些洗費.
大城市的生活費用很高.
他們就想.
不如這樣.
我們離開耶路撒冷.
回去老家.
回去北部.
去到加里尼海的老家.
這個就是今天經文21章.
約翰福音的21章.
第一節說.
這些事以後.
耶穌在提比利亞海邊.
又向門徒顯現.
他怎樣顯現.
記載下面.
這裡說提比利亞海.
其實所說的是加里尼海.
因為那裡有幾個大城市.
所以有時候用不同的名稱.
加里尼海.
格尼薩勒海.
或者提比利亞海.
有時候叫海.
有時候叫湖.
因為是湖.
但是大到你看不到另外一邊.
就好像海一樣大.
所以有時候就叫海.
這個故事就在提比利亞海邊發生.
當時在場的.
有彼得.
多瑪.
拿丹葉.
約翰雅各.
可能安德烈和肥力都在當中.
這七個門徒走在一起.
是很重要的.
為什麼呢.
耶穌被賣的晚上.

$^{321}$他說.
今晚你們為我的緣故都要跌倒.
因為經上記著說.
我要擊打牧人.
羊就分散了.
當耶穌受難的時候.
釘死十字架上高的時候.
那些羊.
耶穌的十二隻羊.
就全部走散了.
現在擊打過了.
羊又走在一起了.
門徒走在一起.
就可以彼此照顧.
一起去面對事情.
這裡跟我們說.
好像教會的小組都是一樣.
我們大家走在一起的時候.
你面對的問題.
可能和其他人面對的都很相似.
所以我們大家可以分享彼此的扶持.
一起並肩行走天路.
很重要的.
所以你沒有參加小組.
你都要加入一個小組.
大家可以彼此的幫助.
門徒雖然有看見復活主.
雖然他們走在一起.
但是在這個時候.
他們對自己未來的方向.
一點都不清楚.
他們不知道將來會有什麼事情發生.
有一個晚上.
幾個門徒坐在海邊.
點起營火.
可能一起來宵夜,一起聊天.
一起懷舊,一起聊天.
說到耶穌.
他們跟從耶穌的經歷.
我記得耶穌在登山部分說的比喻.

$^{361}$我們的生命就像路邊的花草一樣.
我們不要為太多事情擔心.
想起耶穌說的道.
大家一起分享.
說到耶穌走的神蹟.
耶穌治好盲宮.
耶穌治好麻瘋病人.
耶穌趕鬼.
最後說到最後晚餐的時候.
他們怎樣爭拗.
真是不好意思.
怎樣去到看到耶穌被賣.
看到耶穌釘死.
看到耶穌復活.
耶穌復活了.
祂出現了幾次.
影子都不見了.
我們接下來又要怎樣呢.
現在毫無方向.
整天坐在這裡游手好閒.
好像一個無業流氓一樣.
這樣不是路.
他們在談越談越失望.
談著談著心情就越低落.
他想一想.
我什麼都沒有.
跟了耶穌三年.
我懂得什麼呢.
除了懂得打魚之外.
什麼都不懂.
那我不如回去以前的職業再做漁夫.
讓他看看天.
天都差不多到打魚的時間了.
他又彈起身子.
說我出去打魚了.
其他的人聽到.
覺得我們都沒事做了.
你去我都跟你去.
他們就一起出去.
一起出去打魚.

$^{401}$大家一起上船.
把艇撐到外面去撒網.
那晚天黑黑.
風又不是很大.
應該是打魚最好的時候.
誰知道他們.
從天黑撒網撒到天亮.
網撒出去拉回來一次又一次.
半條魚都沒有.
網都是空空如也.
門徒他們覺得.
哎呀我真是.
黑得沒得再黑了.
跟耶穌又跟不上.
我回去打魚呢.
哎呀都不懂打魚.
我的打魚技能完全都沒有了.
哎呀算了算了.
回岸那裡以後再看看.
再打算過了.
在這個時候.
天剛剛的光.
把艇撐到差不多近岸的時候.
岸邊有個人影站在那裡.
大聲問他.
喂小伙子.
你們有沒有魚?.
我要魚在那裡煮魚片粥.
這個聖經說的.
不過這個就是廣東話的譯本.
廣東話譯本.
這個小字的意思就是小朋友.
或者小伙子.
很親切的稱呼.
問他有沒有抓到魚來給我煮早餐?.
給得很沒意思.
他說沒有不好意思.
他說沒有.
說什麼都沒有.
那個人就說.

$^{441}$你們試一下再一次出海.
這次你出海的時候.
就將網撒在.
不要撒在左邊.
撒在右邊.
你就一定有魚了.
門徒雖然不知道這個是誰.
但覺得試一試都沒有吃虧.
就照著去做.
這次網一撒下去.
網就馬上沉下去.
魚多到拉不上來.
網快要爆開.
在那個時候.
耶穌愛的門徒就是約翰.
他又醒目.
心水又清.
他小聲的跟彼得說.
那個你不知道嗎.
是主啊.
約翰怎麼知道是主耶穌啊.
我們要回到這個.
路加福音五章.
那裡再看.
有聖經的時候.
這隻手指夾著這個約翰福音.
揭去路加福音.
我要再回去.
在路加福音第五章那裡.
就說耶穌剛剛出來傳道的時候.
他四圍廣道.
十分的受歡迎.
去到哪裡呢.
一群人圍著他.
耶穌沒辦法廣道.
那天剛剛去到海邊.
彼得剛剛打完魚.
在那裡洗網.
耶穌看到彼得就跟他說.
那我可不可以借你的小艇來用.

$^{481}$當然可以了.
就叫他把小艇撐離岸一個距離.
然後耶穌坐下.
就在船上對岸邊的人廣道.
當時.
彼得撐小艇出去.
他沒得走.
所以他一定要坐下.
聽耶穌廣道.
西比太的兩個兒子.
約翰雅各和安德烈也在場.
可能有其他漁夫也來聽.
耶穌廣道之後就跟彼得說.
你把小艇撐到水深的地方.
你在那裡洗網.
但是彼得說.
Sir.
你知不知道.
我們捉魚捉了一整晚.
一條也沒有.
不過你這樣說.
我就聽你的話.
就出去那裡洗網.
去捉魚.
其實這句話很重要.
為什麼呢.
彼得可以跟耶穌說.
喂.
你是拉比.
我是漁夫.
你知不知道.
我就在這個艇上.
媽媽生我出來.
就在這個艇上.
我就在這個沙灘上.
在這個海灘上.
我在這裡學走路.
我一認識就上船出去捉魚.
我這麼大個人就在這裡.
哪裡有魚哪裡沒有魚.

$^{521}$我都知道.
我是漁夫.
你是教書.
你是拉比.
你又懂得多少.
你對打魚又懂得多少.
彼得可以這樣說.
的確又是真的.
彼得他打魚的經驗.
比耶穌打魚的經驗多很多.
但是他願意聽耶穌的話.
他受信服.
他就跟從.
他就去到水深的地方.
就把網殺出去.
就捉到很多魚.
網差一點都要裂開.
在那個時候.
這群漁夫.
看見魚這麼多.
他們心都驚.
看見他們心都驚.
捉了這麼多魚.
又是這麼大的魚.
我們試了一整晚一點都沒有.
突然間這麼多.
彼得看到這個情景.
馬上跪在耶穌的面前.
他說主啊你離開我.
我是個罪人.
我們看看.
彼得他未落網之前.
他叫耶穌做夫子.
叫他做阿信.
但是當他經歷了主之後.
他看見耶穌行的神蹟之後.
他知道耶穌不是普通人.
他不是一個普通的拉比.
他知道耶穌.
可以行神蹟.

$^{561}$他是神.
他的口說你是主.
他也知道自己是一個罪人.
看見自己的罪惡.
看見自己的誤判.
你知不知道只有認識.
耶穌是主的人.
才會看見自己是有罪的人.
這些人才會被神的呼召.
來使用他們.
被呼召的人.
不是一個完全人.
乃是神可以使用的人.
耶穌跟著和彼得說.
你不要怕.
從今以後.
你就要得人了.
他們就把船撞上岸.
就將一切放下.
來跟從耶穌.
這個就是他們的經歷.
這個就是他們信主的見證.
有人問他們.
你們為什麼要做耶穌的門徒.
為什麼要跟從耶穌.
他的答案就是.
有一天我出去打魚.
一條魚都捉不到.
但是耶穌一來的時候.
就叫我殺網出去.
嘩 裡頭魚多到.
網要爛.
但是又沒有爛.
捉到所有的魚.
約翰彼得說這個見證.
也說了很多次.
所以在那一晚.
在約翰福音那裡.
整晚都沒有魚.
現在網裡頭有這麼多的魚.

$^{601}$多到要爛.
但是又沒有爛.
約翰心水清.
還想起.
在岸上.
講話的那個.
一定是耶穌.
他就跟彼得說.
彼得 你記不記得.
幾年前我們出去打魚.
打不到魚.
但是耶穌一來.
我們捉到很多魚.
但是網沒有爛.
你知不知道.
岸上那個人.
耶穌.
是主.
嘩 彼得一聽見.
彼得的心.
是一個火熱的心.
因為他當時打魚.
為了方便就把外衣脫了.
在那裡捉魚.
但是這一天說到是主.
嘩 要見主 要整齊一點.
就穿起外衣.
然後跳進水裡.
游水到岸邊.
去見主.
當他上岸的時候.
他就.
看見那裡有炭火.
上面有餅.
又有魚.
主耶穌已經為他們準備好.
早餐等他.
耶穌跟彼得說.
你知不知道.
這裡五餅二魚不夠吃.

$^{641}$你就多拿些魚出來.
拿幾條.
剛捉到的魚過來.
我們一起煮來吃.
彼得回到船上.
他們數一數.
有153條.
魚 還是大魚.
不是小魚.
153條大魚.
這麼多的魚.
網竟然沒有爛.
一條都沒有跑掉.
今天我們聽了兩個.
相似的故事.
我們聽了.
兩個很相似的故事.
一個是耶穌在地上.
工作開始的時候.
一個是耶穌.
工作結束的時候.
這兩個.
很相似的故事裡面.
耶穌是用來.
教這班.
善忙的門徒.
第一次.
他做完之後.
需要再一次.
來教他們.
我們可以在這兩個故事裡面.
學習到.
三個功課.
第一個功課.
就是我們生命失落的時候.
就是神要.
工作的時候.
我們看見彼得.
和門徒坐在海灘.
在營火會面前.

$^{681}$覺得失落無望.
準備放棄.
跟從主.
準備重操舊業.
再回去打魚.
打不到魚.
不是.
他們很倒楣.
乃是主耶穌.
為他們有其他的.
預備.
在我們生命裡面.
很多時候都好像門徒一樣.
我們會覺得我們.
失落了方向.
有時我們想放棄.
有時我們會.
面對失敗覺得很失意的時候.
我們會說.
神啊你是否離棄我.
還有.
教會.
有時我們都會覺得.
失落.
我來的時候.
聽見弟兄姊妹的分享.
就覺得我們接下來要怎樣.
接下來的前途.
是怎樣.
下一步會怎樣.
會覺得有點失落.
但我們要記得清清楚楚.
一件事情.
覺得和事實.
是兩件很不同的事.
當我們覺得失落的時候.
我們就更加要去到神的面前.
祈禱.
等候.
因為這個時候.

$^{721}$如果我們放棄的話.
我們就不會看見神的工作.
但當我們去到神的面前.
神啊我知道這個時候.
就是你要工作的時候.
我願意在這裡等你.
我願意被你使用.
神啊求你.
在這裡復興我們.
再一次的.
使到教會復興.
使到我們生命.
清楚前面的前途是怎樣.
在我們生命失落的時候.
在神要工作的時候.
我們要.
對神有忍耐.
我要坐在這裡.
等神的工作.
我要去到神的面前.
祈禱和等候.
神知不知道我們處境.
神是無所不知.
怎會不知道你的情形是怎樣.
神怎會不知道會點的情形是怎樣.
神怎會不知道.
門徒當時的感覺.
是怎樣呢.
神在準備.
神在工作.
準備將更大的事情.
可以發生在我們身上.
神是要我們.
一同的.
祈禱等候.
好像打風一樣.
是不是.
當那個.
旋風.
這個颱風.

$^{761}$在海上越來越久的時候.
做什麼.
在那裡.
吸收水份.
當他來到的時候.
風力就更加的大.
同一樣的.
神在這裡工作.
我們祈禱等候的時候.
神不停的工作.
當來到的時候.
我們得著的.
就會更加的大.
接著我們要知道.
當我認識耶穌.
基督的時候.
就是人.
要悔改的時候.
彼得尊重.
耶穌是拉比.
所以他聽他的話.
當他眼見耶穌是主的時候.
他知道自己.
是一個罪人.
在我們當中很多人來參加崇拜.
聽聖經的話語.
可能你認識耶穌.
耶穌是你的拉比.
真的.
聖經裡面有很多教訓.
有很多道德的思想在這裡.
真的好.
也是真的.
但是我們.
從這個角度.
去認識神的話.
我們.
生命的改變.
是一個有限的改變.
只不過在思想上高.

$^{801}$在哲理上高.
這些是皮毛的改變.
我們得的祝福.
是很少的.
除非我們認識耶穌.
是我們個人的救主.
否則的話我們仍然背著.
這個罪擔在這裡.
我們生命一樣沒有改變.
我們是沒有辦法得到永生的福氣.
所以我們去到耶穌面前.
認耶穌是主.
認自己的罪.
耶穌將我們的罪擔.
除下.
讓我們在祂面前得到最大的福氣.
還有.
當我們重複經歷的時候.
就是神要.
祝福我們的時候.
神不是要來.
囉嗦我們.
是因為我們.
頭腦不好.
善忘失憶.
老人痴呆.
所以神來到我們這裡.
一次又一次來.
教我們.
要我們有重複的經歷.
使我們可以明白.
使我們更加清楚.
考驗一次又一次來的時候.
就是羅馬書講的.
老練.
老練使我們的經驗更加豐富.
使我們生命.
爐火純青.
這就是神要.
祝福我們的時候.

$^{841}$準備我們去承受.
神無限的恩典.
我們剛才看到.
門徒捉了.
153條大魚.
那為什麼153條.
不是152條.
又不是154條.
有沒有想過這個問題.
解經家用很多方式.
去解答.
但我的解釋.
是很簡單的.
簡單牧師解釋.
為什麼呢.
因為如果網裡.
有154條的話.
網就會裂開.
魚全部跑掉.
如果152條的話.
神給我們的就不夠清.
差一條.
神給我們的.
豐盛.
是充充足足.
不會多一條.
使我們的網承受不到.
也不會少了一條.
不夠清.
這就是神給我們的.
所以我們去到神面前.
我們要聽神的話.
我們的生命.
可以充滿神給我們的驚喜.
不單止被呼召的時候.
是最豐盛的祝福.
在我們的生命.
低落的時候.
神一樣.
要來祝福我們.

$^{881}$我們得的.
祝福是多到.
極點.
但也不會超太多.
我們承受不起.
所以我們今天來到神的面前.
我們知道神給我們的.
一定是最好.
當神給我們有重新經歷的時候.
我們去神面前.
獻上感恩.
神給我們的是最好.
你給我們匯點的也是最好.
我們來到你面前.
我們一起低頭祈禱.
神啊你真是一位信實的神.
很多時候我們覺得你離棄了我們.
但其實你就是在我們身邊.
準備來祝福我們.
在我們不明白你的作為的時候.
我們就好像那些叛徒一樣.
我們需要你帶我們.
舊地重遊.
再一次來提醒我們.
使我們可以對你有一個.
更深入的體會.
讓我們更加認識你.
明白你的心.
願我們的心像約翰一樣的清.
願我們的心像彼得一樣的熱.
願我們更加的來愛你.
奉耶穌聖名祈求.
阿們.
\newpage



\section{出埃及記 14:4}
\label{sec:031969}
\textbf{聖經人聖經事:我在法老和他全軍得榮耀}~~~~黃克勤牧師博士
\newline
\newline
連結: \href{http://www.church.com.hk/acms/content.asp?site=cdc\&op=show\&type=product\&code=031969\&layout=sermon}{\texttt{031969}} ~~~~ 語音日期: 2026-04-10
\newline
\newline
\hyperref[sec:031968]{\small{< < < PREV SERMON < < <}}
~
\hyperref[sec:index_chronic]{\small{[返順時目]}}
~
\hyperref[sec:031970]{\small{> > > NEXT SERMON > > >}}
\newline
\newline
出埃及記 14:4
\newline
\begin{longtable}{cl}
\hline
\hline
章節 & 經文 (和合本修訂版)\\
\hline
14:4 & \begin{tabularx}{0.7\textwidth}{X} 我要任憑法老的心剛硬，他要追趕他們。我必在法老和他全軍身上得榮耀，埃及人就知道我是耶和華。」於是以色列人照樣做了。 \end{tabularx} \\ \\
[1ex]
\hline
\hline
\end{longtable}
$^{1}$各位電影姐妹大家好.
又和大家在網上見面了.
這次也是用Zoom的方式.
和大家分享《聖經人》和《聖經事》.
是《生命說書教》的網絡.
每天都和大家讀《聖經》的一個活動.
這次的題目就是.
《我在法老和他全軍得榮耀》.
這個內容你一聽就會知道.
由於一直在說摩西,法老.
你就知道現在他們已經去到紅海旁邊.
這個地方.
我們先看看今天的經文.
看一張地圖.
從這張地圖你會看到.
哥山上次提過了.
哥山地在左手邊這個地方.
是一個很多條河.
其實是一條河.
就是尼羅河由南向北流下來.
然後進入地中海.
由於去到平原地方.
河流就會分岔.
和形成一個三角洲.
這個就是尼羅河三角洲.
很肥沃的.
是全埃及最好的地方.
由於以前約瑟和法老的關係.
就讓他們在以色列住.
後來新王興起.
不認識約瑟.
200萬的以色列人成為奴隸.
在這裡要他們建兩個城.
一個就是蘭塞.
另一個就是比東.
建兩個城都很大.
要來擺放貨物.
很辛苦.
後來摩西就帶他們去昌埃及.
由這一點開始走蘭塞.

$^{41}$然後向南走.
就一直下來.
去到梭葛.
比東城的人從以色列出來.
然後大家一直向南走.
向南走的時候.
我用一些大字讓大家看到.
你看到這裡有一個點.
是一個城.
一個地方.
就叫做比哈希祿.
一會兒我們讀聖經就看到了.
然後再向南走一點.
就叫做麥奪.
然後就看到這個.
這個就是紅海.
紅海其實這裡只不過去到一個很小的部分.
就是一個最窄的地方.
就預備在那裡過海.
他們由蘭塞起行.
就向南走了.
就匯合了比東的人.
比東的人就在梭葛那裡集合齊了.
就向南走了.
照聖經他們一路走到二塘.
去到二塘這個地方.
其實就是這個位置.
然後神就帶他們不要再向北走.
就去了麥奪和比哈希祿這個地方.
就在梭葛安營.
安營的地方就是說他們在那裡紮營.
在那裡停留.
而對著這個海很窄的就是紅海.
紅海對面這個地方就叫巴力洗分.
所以其實他們由蘭塞起行.
去到比東匯同了.
比東的以色列人.
這兩個城的人就在梭葛集合.
然後向南走.
就來到二塘.

$^{81}$二塘然後向北走回頭.
就去到這個海邊.
然後紅海邊集合.
雖然我們現在看這裡很窄.
但是我相信在當地來看.
就已經是很寬闊的海.
就是一段海路.
結果就在這裡.
紅海發生就在這裡.
就是他們在這裡.
於是由於他們是由蘭塞向南走到二塘.
然後繞回頭.
這個就令到埃及人認為他們.
當然是在曠野走迷了路了.
如果不是怎麼會走到.
走過了隆幣繞回頭.
其實他們真的以為是的.
原因就是.
大家看看這個地方叫巴加薩.
這個別氏巴.
已經是現在的以色列地方.
迦南地就是這一帶.
看到耶利哥.
耶利哥附近大概這個利治附近.
就已經是耶路撒冷了.
加薩就是現在的非利士地.
就是經常在加薩地帶和以色列打仗.
別氏巴這些都是以色列地方.
其實如果由蘭塞起行.
到梭葛匯和比東裡面的以色列人.
如果他們走海邊的路.
非利士路來到迦南是很近的.
不過神是帶他們向南走去二塘.
再走回北.
然後就在那裡預備.
預備有些東西在那裡安營.
大家都莫名其妙.
然後聖經說.
埃及法老追兵下來.
追他.

$^{121}$他們的神很厲害.
殺長子的.
後來都回心轉意.
不是啦.
二百萬免費勞工.
追他.
於是就追到那裡.
由於他們發覺.
他的神力一定是失去了.
因為帶他們去了南.
然後再回去北.
於是以為他們迷路了.
迷路即是說他們很害怕.
當然很害怕.
在抗疫又不適應.
不如我們去攻擊他.
這是引了埃及士兵追他們的.
聖經說這次毒的就是比哈希路.
和墨奪的對面.
巴力西芬對著海.
他們就在這個地方安營.
大家看到這條線.
你們知道過了紅海之後.
就下去順的抗疫.
即是西太半島.
然後再下去走回上來.
又在這裡繞一圈.
繞了一圈之後.
才去到河裂山.
即是西乃山.
然後才拿到誡命.
然後才進去.
經過摩押進入耶利哥.
這裡他們在抗疫住了四十年.
好了.
繼續看看今天.
《昌合記》第十四章第一節.
「耶和華曉遇摩西船」.
其實應該說是神對摩西船.
由於我們以往的中國人.

$^{161}$對於上帝和我們說話.
或者皇上說話.
那就是叫做「曉遇」.
「曉遇」意思就是說.
上從下講.
就叫做「曉遇」.
其實就是耶和華對摩西船.
你吩咐以色列人.
轉回安營在比哈希祿.
墨奪和海的中間.
轉回.
甚麼叫轉回呢?.
轉回的意思就是說.
他們本身已經去了南面的二唐.
他們回頭去哪裡?.
就是我剛才說的.
就是比哈希祿.
去這個地方.
去比哈希祿.
的前面.
比哈希祿的前面.
其實叫他們去墨奪.
去墨奪那裡.
剛剛對著海.
大家知道了.
我想解釋一下比哈希祿是甚麼意思.
「Piha」.
「Ch」讀K.
「Pihakiroth」.
「Pihakiroth」呢?.
「Pihakiroth」.
是沙草生長之處.
沙草是甚麼呢?.
對不起,我也不知道.
我想是一種草.
這種草可能在曠野比較多.
在那裡有一大片沙草.
名字就叫「Pihakiroth」.
它就在那裡.
前面的墨奪.

$^{201}$「Midot」.
「Midot」是高塔的意思.
可能以前有人在那裡蓋過一個高的塔.
就用來做記號.
你見到巴比塔也是.
為甚麼要蓋一個高塔呢?.
我要大家都可以見到.
招聚那些人一起.
那一片曠野太寬太遠.
看不到.
於是就用一個高的塔.
讓人可以看到.
作為一個地方.
就是墨奪.
就是「Midot」.
就是高塔.
高塔就在海邊.
這個墨奪在海邊.
對著對面.
對面就是巴力洗分.
巴力洗分.
這個「Ba」就是巴力.
「Chephon」.
「Ba Chephon」.
就是北方之王的意思.
「Chephon」就是北方.
「Ba」就是巴力.
是這樣的意思.
將它譯成「巴力洗分」.
譯成「北方之王」.
是靠近海邊安逸.
第一點我和大家去思想.
這個藍色的螢光.
就帶著他們靠近海邊安逸.
在甚麼海呀?.
紅海旁邊.
你見到由出埃及.
整件事.
你會發覺到上帝.
是有祂的手在那裡工作的.

$^{241}$祂在那裡工作.
是預備好一個階段一個階段.
怎樣做怎樣做.
嗯.
他們的哀聲我聽到了.
祂在那裡已經四百年了.
我要帶祂們出去.
然後就找一個人.
就找摩西.
這個摩西早就預備了.
他成為一個埃及王子.
他學了埃及所有的東西.
知道他們想些甚麼.
然後摩西安排了在那裡有四十年.
這個都很重要的.
因為這四十年有甚麼作用?.
是非常之有作用的.
這四十年的作用就是令到.
你說帶了二百萬人入曠野.
有誰有曠野求生的技能?.
有誰在曠野懂得怎樣做?.
只有一個,就是摩西.
因為他在那裡住了四十年.
熟悉了.
就帶了他們出來.
帶了他們去兩個城.
蘭塞和比東.
大家看到吧.
兩個城.
他們就要集合在一個地方.
就是在疏國集合.
這樣想.
然後向南走到以唐.
然後向北走回頭.
就在巴黎.
對著巴黎洗坑的對面.
就是在這個墨奪那裡集了營.
你看到吧.
蘭塞起行.
比東集合點.

$^{281}$疏國向南走.
南走之後回頭.
這個整個都是一個計劃.
是一個計劃.
就是上帝的計劃.
是去幫以色列人的.
就是帶他們出來.
於是就靠近海邊安營.
靠近海邊安營就背上了箭.
第三節 法老必說.
就是這一段說話.
這段說話是誰說的呢?.
仍然是上帝曉如摩西所說的.
你帶他們這樣走.
走到那裡就繞回頭.
然後在這裡集營.
這個地方被希律和墨奪對著對面.
就是巴黎洗坑.
在那裡集營.
這樣的時候.
我是想那些埃及人.
他們是認為哈哈哈哈.
以色列人在地中妖魅了.
迷路了.
抗議他們不懂得走.
已經將他們困住了.
他們又沒有地圖.
又沒有指南針.
又不是熟悉的地方.
他們你知道.
剛才所說的這些走的路線.
全部都是埃及地.
所以埃及政府.
是完全懂得這些地的.
哈哈哈哈,他們迷路了.
這時正好.
我,耶和華.
要使法老的心肝硬.
解釋過了.
哎呀,又是你搞的.

$^{321}$又不讓人放.
然後你又洗他的心肝硬.
你洗他的心肝硬.
不要罰他了.
這個「使」.
是指「接受」.
或者是指「讓」.
「容讓」.
容讓他的心肝硬.
有時候你會知道.
有時小朋友都是這樣的.
一定要這樣,一定要這樣.
好吧,隨便你.
你自己小心.
你不拿雨傘打到,你就打到了.
你穿這雙鞋子.
一會兒你去旅行就會滑倒.
我不是喜歡穿這雙.
然後去到,鞋又爛,人又滑.
就讓你試一次.
讓你的心肝硬.
好,給你.
讓你見識一下.
法老是「讓」,神是「讓」,法老的心肝硬.
於是他就說.
哈哈哈哈,迷路,我要追趕他們.
當他們要做這個追趕的動作的時候.
我,耶和華,就會第二點.
就是在法老和他全軍身上得榮耀.
我不單止帶你出來.
我還要在他的全軍上得榮耀.
法老來說,最厲害的地方.
除了是他的國土.
他的財力.
他的王位.
國土夠大.
都是重要的國家.
他最仗賴的是什麼呢?.
就是他有很強的軍隊.
譬如今天世界上.

$^{361}$強國都是有很強的軍隊.
當然,譬如像埃及這麼大的國家.
又這麼富裕.
又這麼文明,那些叫文明古國之一.
跟今天所說的一樣.
就是他的貨幣很重要.
他的技術很重要.
他的軍力很重要.
這三樣東西.
多數都是他的金融.
加上他的技術.
加上他的軍力.
這三合一.
就會造成他的國家強盛.
可以在別人身上掠奪等等.
這次上帝就說.
我特意要帶他來.
讓他追趕你.
我由他的心肝硬以為自己很厲害.
我要在他,在法老.
和法老的全軍身上.
奪得榮耀.
讓他看看,原來他還是不行的.
埃及人就知道我是耶和華.
他們都知道了.
滅長子之災已經知道了.
厲害了.
你們去拜你的耶和華.
快點走,快點走.
要金要銀,要衣服就給你.
他們更加認識耶和華了.
這個也挺有趣的.
我會加多兩個字.
更加知道我是耶和華.
於是以色列人就這樣走.
就是說下了以堂.
向北走回去勿奪.
有人就告訴埃及王法老.
百姓逃跑.
這個很明顯就是假消息.

$^{401}$有時候傳得挺傳.
為了討好皇上.
各樣臣服.
於是就作出一些信息出來.
一些假消息.
讓有親打仗都會有假消息.
百姓就跟著摩西不跟.
在那裡擴散了.
明明摩西向南走.
誰知道他們向北走.
摩西要跟他回頭.
現在亂了.
百姓逃跑.
法老和他的臣服.
就向百姓變心.
這個百姓是指以色列.
向以色列變心.
這個變心是談戀愛.
變心不是.
是改變自己.
向他們改變了自己的想法.
怎樣改變自己呢?.
就說我們用以色列人去.
不再服侍我們.
這做的是甚麼呢?.
為甚麼我們這麼傻.
讓他們走也有.
我們是否傻了?.
他們是否瘋了?.
為甚麼會這樣?.
傻傻的就肯放他們走.
不行.
他們不再服侍我們.
怎會放他們走.
他們見到現在.
以色列人弱.
當然他們已經逃跑.
當然不聽摩西的話.
摩西已經失勢了.
這個時候是最好.

$^{441}$去抓他們回來.
法老就預備他的車輛.
帶領軍兵和他.
並帶了六百輛特選的車.
這個要留意.
六百輛特選的車.
即是說他的特種部隊的車輛.
其實都是戰車.
不過是他最精銳的.
六百輛特選部隊.
這多厲害.
因為每一輛車都有車兵長.
車兵長.
不是拉著一條韁繩.
那個是司機.
那個是師傅.
車兵長就在上面.
喊著走哪裡.
左一點 右一點.
他就拿著武器.
另外就有士兵一兩個在那裡.
那些車很大.
還有一個士兵去幫車兵長.
即是你可以想像.
用六百輛今天的坦克.
坦克裡面就有一個司機.
有一個車長.
那個車長很多時候.
會把頭伸出來看.
把頭伸回去.
然後有一個士兵在裡面.
就入炮彈.
另外有一個就做幫手.
一輛坦克車裡面有三,四個人.
其中一個就是車長.
是長官.
指揮話事.
叫車行轉左轉右.
瞄準哪裡.
人就馬上換炮彈.

$^{481}$入炮彈.
其實在當時來說.
這六百輛坦克.
都是特選的.
裝甲部隊.
另外埃及所有的車.
你不要以為.
所有埃及的奴仔.
拉的木頭車.
甚麼都來.
不是 走到半路爛了.
所有的車.
其實全部軍隊的車.
其中有六百輛.
是特種部隊坦克車.
另外有普通坦克車.
都來.
每一輛車都有車兵長.
所以不要誤會.
所有的車.
連運貨的牛車.
甚麼都來.
不是.
是戰車.
所有出動到的戰車.
和六百部.
特別的戰車.
因為每一輛都有車兵長.
所有的車不是普通車.
是戰車.
即是說.
全國精英進出.
這裡可能有一千人.
二千人都不一定.
一千輛車 兩千輛車.
都不一定.
第三點要說的是.
埃及皇法老的心肝牙.
又來了.
屎他的心肝牙.

$^{521}$是讓他的心肝牙.
其實.
全國的精英一齊.
和士兵出去打仗.
如果你的心不是紅.
如果不是自己肝硬.
以為自己很厲害.
你是不會這樣做的.
因為你當然要留一些士兵在國內.
鎮守.
防守其他敵人.
出來見你打的時候.
他會過來動你.
當然會留一些.
一條心.
要帶二百萬人回來.
沒有二千輛戰車.
只剩下特種部隊六百架坦克.
怎麼辦.
於是他就打起來.
追趕以色列人.
由於他覺得不是很遠.
所以都可以.
趕快追.
是因為以色列人士.
昂然無懼地出埃及.
為什麼法老的心肝牙.
一定要追他回來.
因為回想起來真的生氣.
為什麼呢.
他們竟然昂然無懼地.
離開埃及.
這樣傳出去.
我還要出來走.
竟然一群.
以色列人做奴隸.
說走就走.
當時.
第十個災那麼厲害.
第九個災有光沒有光.

$^{561}$我們沒有光.
以色列地有光.
十個災真是沒得頂了.
長子所有死了.
無論人與畜的長子.
都死了.
那麼厲害.
當然要放他走.
但是其他國家不知道.
怎麼辦.
我們埃及都不怕.
我們還要服眾.
於是不是.
走了都要追他回來.
這就是埃及王法老的心肝.
是神讓他的心肝.
好了.
你見到就是這樣.
埃及人追趕他們.
馬匹車輛馬兵.
和軍兵就在海邊上.
那你見到.
馬匹車輛.
是一樣東西.
就是戰車.
另外還有馬兵.
這些是騎馬.
沒有車的.
軍兵就步兵.
就在後面跑啊跑啊.
跑上來.
傳那些士兵出來.
就去到海邊.
比哈希勒.
和巴力西芬.
巴力西芬是對面海.
在他們安然的地方.
就追上他們.
嘩.
怕不怕呢.

$^{601}$以色列人真的臉都青了.
終於追到了.
去到這些地方追上了.
今天和大家分享的.
就去到這裡.
因為追上了.
被他們追到.
回到線.
停了.
我很想和大家.
留意一件事.
就是上帝的工作.
是很奇妙的.
你不知道.
神會做些什麼工作.
不過祂是有計劃.
將那些工作做好.
你要知道.
信靠祂.
祂會一步一步有祂的心意.
我們一起祈禱.
基本上我們多謝你.
讚美你.
因為你是有奇妙的作為.
你自己做一件事.
是預先都鋪排好.
都處理好.
一步一步.
有個成就.
每一個人甚至仇敵.
都在你的掌握底下.
你是掌管他們的心.
你要他們怎麼做就怎麼做.
你容讓他們怎麼做就怎麼做.
你任憑他們怎麼做.
都是托著以色列人.
都是托著我們.
讚美你,一切盡在你手中.
幫助我們將自己的重擔.
卸給耶和華.

$^{641}$因為祂顧念我們.
讓你得到榮耀稱讚.
祈禱奉耶穌名求,阿們.
今天暫時和大家講到這裡.
最後的powerpoint.
有些處理不到,不好意思.
我們明天再見.
再見,拜拜.
\newpage



\section{申命記 4:44-5:33}
\label{sec:031970}
\textbf{上帝也講KPI?(1) 重述十誡}~~~~沈炳華傳道
\newline
\newline
連結: \href{http://www.church.com.hk/acms/content.asp?site=cdc\&op=show\&type=product\&code=031970\&layout=sermon}{\texttt{031970}} ~~~~ 語音日期: 2026-04-11
\newline
\newline
\hyperref[sec:031969]{\small{< < < PREV SERMON < < <}}
~
\hyperref[sec:index_chronic]{\small{[返順時目]}}
~
\hyperref[sec:031971]{\small{> > > NEXT SERMON > > >}}
\newline
\newline
申命記 4:44-5:33
\newline
\begin{longtable}{cl}
\hline
\hline
章節 & 經文 (和合本修訂版)\\
\hline
4:44 & \begin{tabularx}{0.7\textwidth}{X} 這是摩西在以色列人面前頒佈的律法。 \end{tabularx} \\ \\ \relax
4:45 & \begin{tabularx}{0.7\textwidth}{X} 這些法度、律例、典章是摩西在以色列人出埃及後對他們說的， \end{tabularx} \\ \\ \relax
4:46 & \begin{tabularx}{0.7\textwidth}{X} 在約旦河東伯‧毗珥對面的谷中，在住希實本的亞摩利王西宏之地；這西宏是摩西和以色列人出埃及後所擊殺的。 \end{tabularx} \\ \\ \relax
4:47 & \begin{tabularx}{0.7\textwidth}{X} 他們得了他的地，又得了巴珊王噩的地，就是兩個亞摩利王，在約旦河東，向日出方向的地： \end{tabularx} \\ \\ \relax
4:48 & \begin{tabularx}{0.7\textwidth}{X} 從亞嫩谷旁的亞羅珥，直到西雲山，就是黑門山， \end{tabularx} \\ \\ \relax
4:49 & \begin{tabularx}{0.7\textwidth}{X} 還有約旦河東的整個亞拉巴，向日出方向，直到亞拉巴海，靠近毗斯迦山斜坡的山腳。 \end{tabularx} \\ \\ \relax
5:1 & \begin{tabularx}{0.7\textwidth}{X} 摩西召集以色列眾人，對他們說：「以色列啊，要聽我今日在你們耳中所吩咐的律例典章，要學習，謹守遵行。 \end{tabularx} \\ \\ \relax
5:2 & \begin{tabularx}{0.7\textwidth}{X} 耶和華－我們的神在何烈山與我們立約。 \end{tabularx} \\ \\ \relax
5:3 & \begin{tabularx}{0.7\textwidth}{X} 這約耶和華不是與我們列祖立的，而是與我們，就是今日在這裡還活著的人立的。 \end{tabularx} \\ \\ \relax
5:4 & \begin{tabularx}{0.7\textwidth}{X} 耶和華在山上，從火中，面對面與你們說話。 \end{tabularx} \\ \\ \relax
5:5 & \begin{tabularx}{0.7\textwidth}{X} 那時我站在耶和華和你們之間，要將耶和華的話傳給你們，因為你們懼怕那火，沒有上山。他說： \end{tabularx} \\ \\ \relax
5:6 & \begin{tabularx}{0.7\textwidth}{X} 「『我是耶和華－你的神，曾將你從埃及地為奴之家領出來。 \end{tabularx} \\ \\ \relax
5:7 & \begin{tabularx}{0.7\textwidth}{X} 「『除了我以外，你不可有別的神。 \end{tabularx} \\ \\ \relax
5:8 & \begin{tabularx}{0.7\textwidth}{X} 「『不可為自己雕刻偶像，也不可做甚麼形像，彷彿上天、下地和地底下水中的百物。 \end{tabularx} \\ \\ \relax
5:9 & \begin{tabularx}{0.7\textwidth}{X} 不可跪拜那些像，也不可事奉它們，因為我耶和華－你的神是忌邪的神。恨我的，我必懲罰他們的罪，自父及子，直到三、四代； \end{tabularx} \\ \\ \relax
5:10 & \begin{tabularx}{0.7\textwidth}{X} 愛我、守我誡命的，我必向他們施慈愛，直到千代。 \end{tabularx} \\ \\ \relax
5:11 & \begin{tabularx}{0.7\textwidth}{X} 「『不可妄稱耶和華－你神的名，因為妄稱耶和華名的，耶和華必不以他為無罪。 \end{tabularx} \\ \\ \relax
5:12 & \begin{tabularx}{0.7\textwidth}{X} 「『當守安息日為聖日，正如耶和華－你神所吩咐的。 \end{tabularx} \\ \\ \relax
5:13 & \begin{tabularx}{0.7\textwidth}{X} 六日要勞碌做你一切的工， \end{tabularx} \\ \\ \relax
5:14 & \begin{tabularx}{0.7\textwidth}{X} 但第七日是向耶和華－你的神當守的安息日。這一日，你和你的兒女、僕婢、牛、驢、牲畜，以及你城裡寄居的客旅，都不可做任何的工，使你的僕婢可以和你一樣休息。 \end{tabularx} \\ \\ \relax
5:15 & \begin{tabularx}{0.7\textwidth}{X} 你要記念你在埃及地作過奴僕，耶和華－你的神用大能的手和伸出來的膀臂領你從那裡出來。因此，耶和華－你的神吩咐你守安息日。 \end{tabularx} \\ \\ \relax
5:16 & \begin{tabularx}{0.7\textwidth}{X} 「『當孝敬父母，正如耶和華－你神所吩咐的，使你得福，並使你的日子在耶和華－你神所賜給你的地上得以長久。 \end{tabularx} \\ \\ \relax
5:17 & \begin{tabularx}{0.7\textwidth}{X} 「『不可殺人。 \end{tabularx} \\ \\ \relax
5:18 & \begin{tabularx}{0.7\textwidth}{X} 「『不可姦淫。 \end{tabularx} \\ \\ \relax
5:19 & \begin{tabularx}{0.7\textwidth}{X} 「『不可偷盜。 \end{tabularx} \\ \\ \relax
5:20 & \begin{tabularx}{0.7\textwidth}{X} 「『不可做假見證陷害你的鄰舍。 \end{tabularx} \\ \\ \relax
5:21 & \begin{tabularx}{0.7\textwidth}{X} 「『不可貪戀你鄰舍的妻子；也不可貪圖你鄰舍的房屋、田地、僕婢、牛驢，以及他一切所有的。』 \end{tabularx} \\ \\ \relax
5:22 & \begin{tabularx}{0.7\textwidth}{X} 「這些話是耶和華在山上，從火焰、密雲、幽暗中，大聲吩咐你們全會眾的，再沒有加添別的話了。他把這些話寫在兩塊石版上，交給我。 \end{tabularx} \\ \\ \relax
5:23 & \begin{tabularx}{0.7\textwidth}{X} 山被火焰燒著，你們聽見從黑暗中發出的聲音，那時，你們各支派的領袖和長老都挨近我。 \end{tabularx} \\ \\ \relax
5:24 & \begin{tabularx}{0.7\textwidth}{X} 你們說：『看哪，耶和華－我們的神將他的榮耀和他的偉大顯給我們看，我們也聽見他從火中發出的聲音。今日我們看到神與人說話，人還活著。 \end{tabularx} \\ \\ \relax
5:25 & \begin{tabularx}{0.7\textwidth}{X} 現在這大火將要吞滅我們，我們何必死呢？若再聽見耶和華我們神的聲音，我們就必死。 \end{tabularx} \\ \\ \relax
5:26 & \begin{tabularx}{0.7\textwidth}{X} 凡血肉之軀，有誰像我們一樣，聽見了永生神從火中講話的聲音還能活著呢？ \end{tabularx} \\ \\ \relax
5:27 & \begin{tabularx}{0.7\textwidth}{X} 求你近前去，聽耶和華－我們神所要說的一切話，將耶和華－我們神對你說的話都傳給我們，我們就聽從遵行。』 \end{tabularx} \\ \\ \relax
5:28 & \begin{tabularx}{0.7\textwidth}{X} 「你們對我說的話，耶和華都聽見了。耶和華對我說：『這百姓對你說的話，我聽見了；他們所說的都對。 \end{tabularx} \\ \\ \relax
5:29 & \begin{tabularx}{0.7\textwidth}{X} 惟願他們存這樣的心敬畏我，常遵守我一切的誡命，使他們和他們的子孫永遠得福。 \end{tabularx} \\ \\ \relax
5:30 & \begin{tabularx}{0.7\textwidth}{X} 你去對他們說：你們回帳棚去吧！ \end{tabularx} \\ \\ \relax
5:31 & \begin{tabularx}{0.7\textwidth}{X} 至於你，可以站在我這裡，我要將一切誡命、律例、典章傳給你。你要教導他們，使他們在我賜他們為業的地上遵行。』 \end{tabularx} \\ \\ \relax
5:32 & \begin{tabularx}{0.7\textwidth}{X} 所以，你們要照耶和華－你們神所吩咐的謹守遵行，不可偏離左右。 \end{tabularx} \\ \\ \relax
5:33 & \begin{tabularx}{0.7\textwidth}{X} 你們要走耶和華－你們的神所吩咐的一切道路，使你們可以存活得福，並使你們的日子在所要承受的地上得以長久。」 \end{tabularx} \\ \\
[1ex]
\hline
\hline
\end{longtable}
$^{1}$早上好晚安.
講道.
講完這篇道.
一個多月後我就死了.
最後我就選擇了十誡.
今晚.
看你們的樣子都呆了.
今晚我們講的是申命記.
你知道申命記的背景的確是這樣.
申命記是說摩西大利以色列人40年後.
準備進入迦南地.
在摩亞平原對美國約旦河之前的三篇講論.
申命記是有三篇的講論.
最後一個結尾.
就是在摩西臨死之前.
大約有些不同的學者說.
摩西大利以色列40年後.
最後大約在十一,十二月.
講完那幾篇的講道.
兩個月後就死了.
約書亞就帶領以色列進入迦南.
的確真是這樣.
所以今晚我們是新一個系列.
是講上帝也 講KPI.
如果你有印象.
這是第一篇講十誡.
重述十誡.
在約旦河東摩亞平原.
整個申命記有三篇的講論.
今天開始講第五章.
十誡第五章.
其實是由第四章開始到二十八章尾.
申命記是一個很長的講論.
所以今天所講的.
稍後會點題講一講.
四章四到四十九是一個宣告.
關於律法.
上帝也講KPI.
講律法 我們今天是怎樣理解聖經講律法的呢.
接著由五章一到三十三節重新再講一次十誡.

$^{41}$如果你在教會內.
第一次有印象聽十誡.
通常都是講Tri-Cup記.
Tri-Cup的二十章.
今天是講申命記第五章.
有些人說十誡是以色列律法的憲法.
即是一個outline或大綱.
接著由第六章一直講.
其實有很多法律的內容.
條文 社會生活 民生條例.
很細節的.
因為當時的以色列.
他們是講他們的法律和生活是整合在一起的.
我們就不是.
你明不明白意思.
即是我們今天.
你在香港.
我們叫宗教自由.
你宗教自由你喜歡拜甚麼就拜甚麼.
不是 以色列.
你敬拜哪一個在你的生活是其中一部份.
很多的律法 律例是黏在一起的.
所以到27,28歲.
小傑就講了祝福和奏哉.
也是申命記其中一個重要的主題.
好 既然如此.
就和大家很簡單.
讀不讀完.
嘩 讀完這麼多經文.
我想這樣讀.
頭尾那些我讀.
笑甚麼.
你聽了 嘖 講了這麼久.
但中間那十誡我想你們讀.
可以嗎.
頭尾不是十誡那些我讀.
十誡那些你讀.
可以嗎 好.
新民記四章44節.
摩西在以色列面前所陳明的律法.

$^{81}$就是摩西在以色列人出埃及後.
所傳給他們的法道律例典章.
在約旦河東伯比爾對面的谷中.
住在希什本阿摩利王西宏的地方.
就是這個西宏是摩西和以色列人.
出埃及後所擊殺的.
你記得上一個系列是重述歷史回顧.
他們得到他們的地.
又得到巴山王岳的地.
就是兩個阿摩利王在約旦河東日出之地.
從阿亂谷的阿羅爾直到西雲山.
就是黑門山.
環遊約旦河東的全阿拉巴.
直到阿拉巴河靠近比斯加山根.
好 正式進入第五章.
摩西將以色列眾人召了來.
對他們說.
以色列人啊.
我今日曉遇你們的律例典章.
你們要聽.
生命的一個很重要的字.
使你們可以學習謹守遵行.
耶和華說我們的上帝.
在河列山和我們立約.
這個約不是我們列祖所立的.
是與我們今日在這裡存活的人所立的.
耶和華在山上從火中面對面跟你們說話.
那時我站在耶和華跟你們中間.
要將耶和華的話傳給你們.
因為你們懼怕那火沒有上山.
就說了.
請.
我自己雕馬頭像.
也不可做什麼形象.
可得上天下地.
我地底下水中的寡民.
不可為大那些象.
也不可羨慕他.
因為我耶和華的上帝是基石的上帝.
恨我的.

$^{121}$我必追討他的罪.
義父及子.
直到三世代.
愛我受我皆命的.
我必向他努力善待.
直到善待.
我求求王請耶和華給上帝命令.
因為王請耶和華命令.
耶和華必不變化為無罪.
先休息一會兒.
還有兩張泡沫.
我求求耶和華.
給上帝所賓會的受訪者為聖人.
每日要勞碌.
做你一切的勞.
但第一次向耶和華給上帝講授的安息.
第一日.
你要我祈念對徒.
被牢牢尊重.
永遠要去尋對.
要去客許.
無論何苦都不可做.
所以要我祈念.
可以沒有你一樣的安息.
第二日要紀念天上的祂.
被祖國祝福.
要我給上帝.
用大量的收穫新創的光輝.
讓你永不去領袖.
因此耶和華給上帝.
吩咐你受訪聖人.
我求耶和華給上帝所吩咐的.
孝敬父母.
使你得福.
並使你的日子.
在耶和華給上帝所賜給你的上身.
以長久.
不可殺人.
不可姦淫.
不可勞碌.

$^{161}$不可作奸淫證侵害人.
不可貪污人的妻子.
也不可貪污人的河奴.
便你復起.
有道.
並拜一世所有的天.
好 到我了.
這些話是耶和華在山上.
從火中雲中幽暗中.
大聲曉喻你們傳會中的.
此外並沒有天別的話.
他就把這話寫在兩塊石板上.
交給我了.
那時火焰瀟山.
你們聽見從黑暗中出來的聲音.
你們支派中所有的首領和長老.
都來就近我.
說 看啊.
耶和華我們神.
將祂的榮光和祂的大能.
顯給我們看.
我們又聽見祂的聲音.
從火中出來.
今日我們聽見.
得見上帝與人說話.
人還存活.
現在這大火將要消滅我們.
我們何必冒死呢?.
若再聽見耶和華我們.
上帝的聲音就必滅亡.
凡屬血氣的.
曾有何人聽見.
永生上帝的聲音.
從火中出來.
將我們聽見.
還能存活呢?.
求你近前去聽.
耶和華我們上帝所要說的.
一切話.
將祂對你說的話.

$^{201}$都傳給我們.
我們就聽從遵行.
28節.
你們對我說的話.
耶和華都聽見了.
耶和華對我說.
這百姓的話我聽見了.
他們所說的都是.
唯願他們存在的.
這樣的心敬畏我.
常遵守我一切誡命.
使他們和他們的子孫.
永遠得福.
遵守和得福.
都是《新命紀》.
另外兩個很重要的詞語.
你去對他們說.
你們回將彭去吧.
至於你可以站在我這裡.
我要將一切誡命律例.
典章傳給你.
你要教訓他們.
使他們在我賜他們遺孽的地上遵行.
所以你們要照耶和華.
你們上帝所吩咐的謹守遵行.
不可偏離左右.
耶和華你們上帝所吩咐你們的.
你們都要去行.
使你們可以存活得福.
並使你們的日子.
在所要承受的地上得以長久.
先透口氣.
但我覺得這樣讀是很有意思的.
你有沒有試過在家裡讀經呢?.
讀出來會有些印象.
有時未必完全明白.
或者會有些混亂.
但你不斷地讀.
讀得多你對於上帝的說話.
起碼會有印象.

$^{241}$這是很重要的.
剛才的經文其實不是很複雜.
44-49節是一個宣告.
整個大段落.
20章經文的大段落裡面.
是宣告關於律法.
如何與歷史扣上.
然後再引落去.
重新說一次十誡.
這個十誡如果回到教會.
簡單來說是我們與上帝的關係.
與我們與其他人的關係.
今天很想在十誡裡.
歸納為三個重點.
想與大家分享.
關於上帝.
關於安息.
關於我們與不同人的關係.
到最後剛才說了.
懼怕和勉勵.
上帝說.
「惟願他們全者的心敬畏我.
常遵守我的一切誡命.
使他們和他們的子孫永遠得福」.
我們一起祈禱.
永活的上帝.
我們來到你的面前敬拜你.
祈求真理的聖靈.
這一刻幫助我們每一個弟兄姊妹.
不是只聽到台上.
一個叫Ivan的人所分享的訊息.
而是我們的心向你打開.
去明白你的說話.
明白你透過你的說話.
想今日我們所領受的是什麼.
你按照我們每一個弟兄姊妹.
今日生命的地步和光景.
去餵養我們.
叫我們的靈能夠得到飽足.
多謝主聽我們祈禱.

$^{281}$奉耶穌的名.
阿們.
在基督教信仰或是整個猶太文化背景裡.
上帝很特別的是.
祂和以色列人是立約的.
祂不是純粹說.
就這樣交下這本寶.
全部都請你自己處理.
祂一開始就說.
我是耶和華 你的上帝.
是有一個關係.
這個創天造地的主.
是來和以色列人.
或者和我們今日.
建立一個個人的關係.
所以當舊約裡面.
很多時候說群體.
新約會比較多一點.
其實兩者都是互通.
但這個同樣的道理.
很重要的就是.
耶穌基督來到這個世界.
舊約和新約一樣.
都是來到我們的生命裡面.
和我們有一個連繫.
和我們有一個關係.
好像剛才所說的五章二席.
耶和華 我們的上帝.
摩西稱呼.
摩西和一群以色列人.
在摩亞地平原裡面.
一路分享.
耶和華 我們的上帝.
在荷列山.
即是齊合記二十章.
和我們立約.
有一個連繫.
有一個關係.
這個約不是和我們列祖立的.
是整個群體和我們今日.

$^{321}$here and now.
我們今日.
是有一個關係.
無論我們認識上帝多少.
無論我們透過哪一個認識上帝.
這個上帝最終是很想和我們.
建立一個個人的關係.
是和我們.
無論我們有多準備.
所以新約說得很有意思.
耶穌基督為我們死後.
升上天.
聖靈來到.
很想住在我們心裡.
以致這個很偉大的上帝.
在我們的人生.
在我們個人的裡面.
能夠有一個連繫.
以致當我們走這條路的時候.
祂想和我們一路同行.
一路走下去.
而這個上帝是一個什麼上帝呢?.
祂說曾將你從埃及地圍爐之家領出來.
這個上帝是一個救贖的上帝.
我們新約就真的好了.
當時以色列人未必明白.
但不是的.
出埃及其實一樣.
其實律法的頒佈.
其實一開始是先救贖.
後講律法.
這個觀念在聖經裡面其實很重要.
但如果你今天很多時候不是.
等價交換 你上班.
你上班是怎樣?.
先和你講R and R.
你的JD做什麼的.
你做到東西才有薪金出.
是不是這樣?.
你一開始是怎樣?.

$^{361}$你上班的時候.
拜託Victor.
我想先墊兩個月薪金.
你今天剛剛簽約.
約未簽的.
你可以平平安安離開.
救贖是另一回事.
怎樣去頒佈律法.
祂說先救了你.
先將你從圍爐之地拯救出來.
然後在曠野裡頒佈律法.
救贖主.
接著這個救贖主是怎樣呢?.
這個救贖主是定界線的上帝.
你讀十誡很多.
不可不可不可.
上帝這樣又不准綱又不准.
遲些打幾多都不行.
那些都寫下去了.
不過好彩是以前寫下去的.
定界線其實不是不好的.
如果什麼都可以.
其實未必是好事.
這個上帝是怎樣呢?.
是先祝福.
後管教的上帝.
我講幾個比較日常生活的例子.
這個上帝是怎樣呢?.
就好像.
你爸爸.
我當你爸爸媽媽.
當你爸爸媽媽在家的時候.
知道你考得不太好.
他第一件事是怎樣呢?.
我不知道你那些.
一來到就即照.
我給你那麼多錢讀書.
你考個99分不是吧.
99分.
這個上帝如果是一個爸媽會怎樣?.

$^{401}$他不是先罵.
而是先會講.
無論你考幾多分.
你都是我最愛的子女.
這件事完全不同.
納若的上帝.
我是和你有個關係的.
你知道我是單親大的.
我從小到大沒有爸爸.
我是這樣教我的女兒的.
我在聖經裡面學.
我就想上帝是這樣的.
所以我和女兒.
我以前是這樣說的.
很生氣.
你知道無論我多生氣.
無論我多生氣.
我第一句.
她小時候.
我第一句和她說.
無論你做得多不好.
多少事.
我都那麼愛你.
這個就是我們的上帝.
他是怎樣呢?.
先去肯定那個關係.
然後再和你講管教.
還有管教是怎樣呢?.
另一個例子.
如果你在家.
小朋友不斷上網打機.
遊戲不斷玩.
到了晚上.
原來路由器不開啟.
設定了時間表.
12點就不能打機.
然後爸媽就說.
其實我不是不讓你打機.
我只是不想你晚睡.
所以休息一下.

$^{441}$明天早上我和你一起打.
嘩 正呀.
你有沒有這樣的爸爸媽媽.
你會怕有什麼秘訣.
一起看YouTube.
你明白我的意思嗎.
上帝就是這樣.
他定的規條.
不是純粹什麼都不讓你做.
是為了保護你.
你休息完之後才去做.
我在聖經裡面是這樣學的.
我怎樣教女兒.
什麼叫節制呢.
我學一下.
什麼叫節制.
節制不是什麼都不能做.
節制是讓你可以做.
對你真正有益的事.
你能夠做得到.
叫節制.
你明白我的分別嗎.
節制.
嘩 又不能吃炸雞.
又不能喝汽水.
又不能喝啤酒.
嘩 這樣節制.
人生有什麼樂趣.
不是 可以.
不過你的量要怎樣.
你明白我的意思嗎.
上帝一樣.
上帝說不讓你做那些事.
其實是為了保護你.
他是想你能夠真正做到你值得做的事.
另外 給多兩個例子.
就是公司老闆.
公司老闆.
剛剛才說完.
一來到.

$^{481}$找多個特映.
Max剛回來.
Max來到.
一來到.
嘩 大件事.
Ivan 召廢.
你們知道什麼叫召廢嗎.
不說了.
Ivan一來.
不是 一來.
Max.
你之前對公司的貢獻.
你在公司面對一些困難的問題.
你先有很多的建議.
或者你都不怕做多一點.
你就專程去Cross Department.
提很多的意見.
跟很多同事去討論完這個問題.
可以怎樣處理.
然後才說.
不過這些地方.
另外那些地方.
可能是應該可以這樣考慮一下.
改善一下.
他先肯定你在公司裡面的貢獻.
先去告訴你.
原來他可以察覺到.
原來你做一些.
beyond你工作以外的東西.
他看得到的.
這些老闆.
這些老闆在哪裡找呢.
你回教會跟阿南.
你都看得到的.
我想應該是吧.
這些是這樣的.
接著.
剛剛說到阿南.
如果教會怎樣呢.
如果教會.

$^{521}$你一見到阿南會怎樣.
嘩 阿南又很棒.
不是.
啊 什麼.
阿南 我有些事想跟你說.
一聽.
嘩 這麼糟糕.
他聽完你講的故事.
不知道什麼很複雜.
他第一件事是什麼.
第一件事.
搭著你肩膀.
不如我們祈禱先.
有時我們做基督徒.
說真的.
誰沒有古靈精怪的問題.
可能是上一代.
可能自己行差踏錯什麼路.
有沒有人.
你夠膽在私底下.
適當的環境.
受保護的情況下.
你可以揭露你的軟弱.
你可以甚至揭露你的罪惡.
然後你講完之後.
他第一件事做什麼呢.
他先接納你.
不是說你的事是對的.
他先接納你.
然後跟你一起祈禱.
怎樣拿到力量.
一起去面對.
這個上帝就是這樣.
他先有恩典再去管教.
我剛才講的幾個例子.
無論他是你的父母.
或是職場裡面.
他是你的老闆.
甚至你在教會上所遇到的同工.
這個上帝就是這樣.

$^{561}$這個上帝就是先祝福.
後管教的上帝.
十誡所講的上帝.
他會定的.
不過他定是為了祝福.
但是他先救贖了你.
這個就是聖經所講的上帝.
好 接著第二個 安息.
如果安息.
其實在十誡裡面是第四誡.
安息日.
但其實有很多節的經文.
預備的時候.
為什麼我會這樣打這個出來呢.
其實我就想.
我就拿了這一條.
和稍後都是.
《申命記》第五章.
和《初一及記》二十章.
兩段的經文是有少少不同的.
如果你看得到.
綠色和藍色那兩段.
你有沒有注意到那兩個分別在哪裡.
有沒有注意到.
其實都看得到D字.
分別是.
在《初一及記》講安息.
是從上帝的創造來講的.
《申命記》講安息.
是從救贖的角度來講.
那有什麼特別呢.
對呀 有什麼特別呢.
《申命記》所講的安息.
我想串起來一起講.
《申命記》所講的安息.
因為你知道.
以色列人在埃及維勞之地生活了多少年.
四百年.
習慣了當奴隸.
習慣了當丫鬟.

$^{601}$一見到松木屍.
馬上就倒杯茶.
習慣了自動自覺.
不用別人教的.
習慣了.
你帶我出去逛街.
即是怎樣.
什麼安息.
你有沒有試過不做事.
習慣了做事.
習慣了沒有停下來.
但是他帶領這些人出去了.
埃及在抗疫裡面就說.
你要安息.
你多久沒有試過好好停一停.
還是你覺得停低三十分鐘.
都有一份罪疚感.
你知道.
其實摩西帶領以色列人出埃及去那裡.
英雄之地.
流浪漁民之地.
但他去到的其實是面對著要進入迦南這個地方.
所以其實他是一個文化改變.
是他在幫他們.
在準備他們進入一個新的文化.
所以他其中一個很重要的就是他重新說安息.
今天香港我想你大概都是.
如果你有功返.
有功返就好了.
但你到處走的時候你都不敢停下來.
但是這裡摩西就說.
在一個舊屬的背景裡面.
在遺老之地走出來之後.
你要停下來.
有人說舊約.
講十誡其實是像小蚊子那樣.
他要講很多很清晰很具體的東西.
所以你今天講安息是講不了.
你今天去到巴勒斯坦.
耶路撒冷.

$^{641}$猶太人是怎樣呢.
他們的升降機有沒有去過?我去過幾次.
他們的升降機進去之後是層層停的.
他不想你做事.
但你真的要停下來.
所以你有多久沒有停下來?.
你問自己你敢不敢停下來?.
上帝說在忙碌.
在受壓制的生活環境裡面.
停下來.
這個安息能夠叫我們從遺老之地被釋放.
這個是自由.
當我們今天很多的自由選擇.
吃不同的好東西.
有自由選擇去不同的旅行.
有自由選擇買不同的名牌.
但是很多時候我們沒有一個自由.
夠膽停下來 因為我們不知道什麼是安息.
另外 初一級的同樣在講安息的位置.
也很有意思的是什麼呢.
創造.
創造的安息是.
你回去看一看.
初一級的經文是.
上帝用六天創造了這個世界.
第七天安息.
接著有一些解經家就在解釋.
但是這個上帝是不是真的以後不做事?.
讀一段經文給你聽.
《約翰福音》五章十七節.
耶穌就對他們說.
我不講這些故事 只是讀.
耶穌就對他們說.
我父做事 直到如今 我也做事.
我們所信的上帝.
其實不是創造天地完結之後就翹起雙手不做事.
還在做事.
如果你讀新約的聖經.
耶穌說今天.
耶穌仍然在父上帝的右邊為我們祈禱.

$^{681}$創造的安息有什麼意思?.
創造的安息是.
上帝看著是好的.
創造完結之後 祂就停下來.
今天有時我們的工作.
除了瘋狂.
或者未必是一定的工作.
可能是家人的張力.
有不同的關係.
可能有很多事都停不下來.
有些長輩長期病患.
我沒法停下來.
有些解經家說.
創造的安息.
其實在新約的經文裡說.
這是上帝仍然的工作.
創造的安息給我們一個很大的意識.
其實有些事就算仍然要做事.
我們今天都可以安息.
我的工作做到一半.
我還未完成.
怎樣安息?.
比較理想的是.
求上帝給我們智慧.
安排得好一點.
以致我們有空間停下來.
再說得刁鑽一點.
就算事情不是百分百安排得好.
我都能夠在未完成工作的張力下.
可以安息.
因為安息在上帝的創造裡.
是秩序的一部分.
你明白我的意思嗎?.
這個創造神學是很值得大家去思考的.
今天可能有些事你還未完成.
或者有些事情有些壓力.
你可能還未處理得到.
但是安息日.
你能夠安息.
你能夠交託.

$^{721}$因為先有一個立約的上帝.
先有一個在乎我們的上帝.
我可以在祂面前安息.
不是在我老闆面前安息.
你明白我的意思嗎?.
我老闆只是娶了我.
安什麼息啊?.
你在老闆面前怎樣安息呢?.
但是在上帝的面前就可以安息.
在這個創造的上帝.
在這個立約的上帝.
在這個救贖的上帝.
在我人生多麼顛簸也好.
上上落落.
峰迴路轉.
我在祂面前.
我可以安息.
這是祂的應許.
這是祂的創造.
第三.
從最後的六誡.
很多時候人們都說.
和人的關係.
在這個關係中有什麼特別呢?.
我想帶出兩個重點.
這個關係.
兩段經文.
《初一及記二十章》.
和《申命記第五章》.
這兩段經文.
有兩個比較明顯的分別.
有些解迦就說.
這兩個分別指的是什麼呢?.
剛才我說了.
《初一及記》很多時候是強調上帝的創造.
而《申命記》就強調上帝的救贖.
而第一段是說孝敬父母.
也是說在地上的日子長久.
不過《申命記》多了.
除了長久之外.

$^{761}$還會得福的.
第一個.
剛才你聽我說過.
《申命記》說.
讓神的子民透過進行上帝的說話.
得福是一個重要的主題.
這裡也有說.
第一個.
得到祝福.
第二.
就是說.
不可以貪戀人的房屋.
是《初一及記》記載.
不可以貪戀人的妻子.
但是去到《申命記》.
將不可以貪戀人的妻子.
抽了上來.
不可以貪圖人的房屋.
田地.
多了東西.
還有.
《申命記》的那兩個貪戀.
原來的文字是兩個字.
兩種不同的貪戀.
有些解經家說.
其實這裡指出.
去到《申命記》的時候.
其實凸顯了的是.
一個是心態.
一個是實質的慾望.
你今天上班.
誰不想賺多點?.
今天上班.
誰不想升職加薪?.
不過最好不要做那麼多事.
我們也有些實質的慾望.
但是《申命記》說的.
那個也是一部分.
但是.
還有一樣就是心態.

$^{801}$是內心貪戀別人的妻子.
實質的貪圖人的房屋.
你明白我的意思嗎?.
是說在社會倫理當中.
將我們心的東西.
特別是在現今的時代.
我們會明白得多很多.
其實人的心真的有很多東西.
你看我的樣子.
好像正正常常.
也OK.
但如果捐到我的心裡.
嘩 艾芬 你是這樣的.
但是我想說的是.
你看人的關係有很多的不可不可.
嘩 有很多事情不准做.
但是當你了解當時的背景.
當時的背景.
如果你有少許熟悉的舊約.
迦南地他們父母的關係.
那些父母是可以怎樣對待子女都可以的.
任意去殺.
你知道在迦南地其中一個.
我們今天很難說匪夷所思的是什麼呢?.
他們那些叫什麼呢?.
將嬰兒獻祭.
我發覺打仗的時候.
搞不定了.
皇帝怎樣呢?.
捉了自己的女兒.
捉了自己的兒子.
出來獻祭.
獻巴黎.
然後獻完之後.
這個神就會保護我.
不可殺人.
什麼不可以姦淫.
當時是很多的廟祭.
很多的亂倫.
很多的聖淫亂.

$^{841}$你看回舊約的故事都是這樣.
是吧 猶大.
去廟祭那裡留了一支像.
然後就出事了.
它是很普遍 很開放的.
甚至那些偷渡 巧取豪奪.
作假見證就更加.
說兩句話就冤枉你.
所以說的是什麼呢?.
所以說的是.
當十誡一直說很多.
不可以 不可以 不可以的時候.
其實它說很多.
不可以 不可以的時候.
這些東西不是.
純粹給我們很多的負擔.
不是純粹告訴我們.
抽空地說.
你什麼都不能做.
不是.
因為上帝透過摩西.
帶領以色列人出埃及.
向著應許之地出發.
是為了進入迦南地.
上帝之前已經有說過.
叫你進入迦南地.
不是因為你特別人多.
你原本在百姓裡面.
人很少的就是以色列人.
只不過他們那班人太壞.
所以就想你進去.
所以上帝給的律法 律例.
其實不是給我們.
抽空無理的一堆的規條.
是為了預備 當時為了預備.
這班以色列人.
面對迦南的挑戰 你明不明白.
當他們在摩押平原聽到說.
孝敬父母 不可殺人 不可姦淫.
不可偷渡 不可作假見證.

$^{881}$不是說 你要做一班聖潔的子民.
什麼都不好.
不是 因為那班在迦南的人.
就是不是這樣.
所以他要幫這班以色列人.
Getting ready.
不要進去的時候.
嘩 在什麼地方 你明白我的意思嗎.
所以上帝擺這些律例出來的時候.
其實是要準備他們.
面對將來這些的挑戰.
教會講教導 講律例 一樣.
姐姐們 你出去上班.
我不知道你做哪一行.
我不知道你做什麼職位.
有沒有同事收回佣.
有沒有私底下互相篤背脊.
自己添.
有沒有石膊 洗完膊之後還說你不對.
作假見證.
有些人看老闆的西裝有多好.
背後就說是鮑爾拉.
不過你不知道的.
我們都面對著這個世界的衝擊.
和當時的以色列人其實是一樣的.
聖經不是說.
給很多的規條 不讓我們做很多事情.
聖經是要將人性的黑暗.
或者將這個世界的挑戰列出來.
以致我們有這個Awareness.
知道至少在什麼的Perspective裡面.
我們要Getting ready.
當我們回到家 上班後已經很累了.
還要被媽媽罵 怎樣祝福她.
但這是我們的上一代.
我不是說 爸媽說什麼都對.
她說什麼我都要聽.
但我們怎樣學去Get along with 他們.
怎樣學去Appreciate他們所付出的.
譬如我媽過世了 早幾年.

$^{921}$我媽是很暗沉的.
我從小到大 成長其中一個最大的困難.
就是很怕聽一些critical的東西.
因為我從小到大被我媽罵.
但我相信了耶穌之後.
我都覺得這是對我人生很大的困難.
直到有一天 我後來才突然有些位置.
知道我媽說 為什麼她這麼暗沉.
說什麼不行 說什麼不行 我去捐血都說不行.
捐血有什麼不行 我心想.
我們這些出來聽別人說 捐血不是.
捐血好像對你的血液循環好些.
或者你成長好些等等.
捐血不行 後來我長大之後.
我媽突然有一天跟我聊天.
她說因為我小時候是換過血的.
我不知道 可能我小時候身體不好之類.
所以我媽在她有限的知識裡.
她覺得捐血對我來說是有傷害的.
所以你明白嗎.
有時我們父母很煩人.
但你怎樣去明白在這個鴻溝裡面.
在這個代溝裡面.
你有沒有嘗試了解過你爸媽.
你有沒有知道他當年的辛苦.
你有沒有知道他嘴巴在罵你的時候.
其實他怎樣用他的方法來愛錫你.
你想一下摩西寫了多少年前.
大衛王朝是在說公元前一千年.
耶穌已經死了 升了天二千年.
最少三千多年前.
是三千多年前在寫這些律法回來.
我們今天都覺得很難做.
你明不明白我的意思.
所以這件事不是在說一個時代.
而是在說怎樣準備我們面對這個世界的挑戰.
我們其實已經活在一個充滿挑戰的時代.
好 我仍然很記得.
每周少功課.
我想鼓勵大家.

$^{961}$這個星期差不多完結了.
星期一我想具體一點.
下個星期一 早上十三號.
你早點起床 十分鐘而已.
你醒了之後.
你讓自己靜一靜在上帝的面前.
從上帝的裡面拿安息.
你放慢自己的呼吸五至十次.
你嘗試感受一下上帝的同在.
安息 學而已.
那你而已.
你這樣很難下台.
你可以做十個鐘.
我這裡短一點.
所以我是說五至十次呼吸.
你不會五至十次呼吸十分鐘.
不是我的氣袋.
那些是例外的.
接著你用一首短詩敬拜上帝.
唱一唱.
如果你沒有甚麼詩歌在腦內.
你回去找.
In His Time.
很短的詩歌.
In His Time.
He Makes All Things Beautiful.
Including You.
Including Your Adversity.
唱出來.
但你唱的時候是敬拜上帝.
接著在祈禱裡面.
將來的這個星期.
下個星期.
你會遇到一些關係很難搞.
很困難.
你很想為他祈禱.
可能是你的父母.
可能是你的上司.
可能是你的工作夥伴.
可能是你身邊的朋友.

$^{1001}$總之你在祈禱裡面.
提名教託而已.
不用說太長.
五至十分鐘.
你讓今日這個訊息.
成為你和上帝生命的encounter.
我想請敬拜他.
弟子們.
領袖會過去.
時代會轉變.
上帝的話語卻歷久常新.
當新一代的以色列人.
聚集在摩訶的平原.
面對著約旦河.
上帝的僕人摩西臨終之前.
很詳細地頒佈律法律例典章.
講解之前.
他重新講一次十誡.
十誡不是十條僵化冰冷的規條.
而是能夠讓上帝.
讓神的子民在不同的文化世代當中.
可以不會迷失.
不單止不迷失.
而且還可以捉緊上主的.
祝福的十條生命的原則.
弟子們.
今日你和我不再需要去巴勒斯坦.
去尋求這個應許之地.
再入去迦南地.
我們可以仍然去領受.
上帝帶領我們.
在我們的人生軌跡.
在我們生命的profile裡面.
進入一個滿足和豐富的人生.
十誡.
今日我分享的.
用了三個角度.
歸納了三個大原則.
第一.
敬拜尊崇唯一值得我們.

$^{1041}$全心交託給他的上帝.
就是那位創天造地.
創造萬物的主.
也是那位為我們寫十字架上.
而且復活了的耶穌基督.
無論我們去到那裡.
心情高低.
環境順逆.
都陪著我們的聖靈.
第二.
進入真正的安息.
領受真正的平安.
在這個世界.
充滿著競爭和生存的壓力裡面.
被釋放.
我們仍然在上帝面前努力工作.
但是學習.
不去被這個世界的成就和眼光綁死.
或者我們可能外在的環境很困難.
甚至可能我們的內心有很多忐忑.
但是我們的心靈.
仍然能夠安穩在上主的應許.
和所賜下的安息裡面.
第三.
健康的關係.
無論是父母,長輩,叔叔.
有時代鴻溝的上一代.
或者是鄰舍,同事,工作夥伴.
在這些相處當中有很多張力.
或者是我們很想有親密的關係.
但是在這些掙扎當中.
仍然選擇光明.
選擇懂得按上帝的教導.
如何做一個有智慧的選擇.
不同的時代有不同的勢力.
會有不同的生活方式去主導.
但無論如何.
願十誡上帝所頒佈的這十條誡命.
幫助我們.
讓我們每一個頂尖在上帝的裡面.

$^{1081}$繼續被祂帶領前行.
\newpage



\section{約翰福音 17:21-23}
\label{sec:031971}
\textbf{與同輩同行的起點:合而為一}~~~~唐郭麗婷牧師
\newline
\newline
連結: \href{http://www.church.com.hk/acms/content.asp?site=cdc\&op=show\&type=product\&code=031971\&layout=sermon}{\texttt{031971}} ~~~~ 語音日期: 2026-04-12
\newline
\newline
\hyperref[sec:031970]{\small{< < < PREV SERMON < < <}}
~
\hyperref[sec:index_chronic]{\small{[返順時目]}}
~
\hyperref[sec:031972]{\small{> > > NEXT SERMON > > >}}
\newline
\newline
約翰福音 17:21-23
\newline
\begin{longtable}{cl}
\hline
\hline
章節 & 經文 (和合本修訂版)\\
\hline
17:21 & \begin{tabularx}{0.7\textwidth}{X} 使他們都合而為一。正如父你在我裡面，我在你裡面，使他們也在我們裡面，好讓世人信是你差我來的。 \end{tabularx} \\ \\ \relax
17:22 & \begin{tabularx}{0.7\textwidth}{X} 你所賜給我的榮耀，我已賜給他們，使他們合而為一，像我們合而為一。 \end{tabularx} \\ \\ \relax
17:23 & \begin{tabularx}{0.7\textwidth}{X} 我在他們裡面，你在我裡面，使他們完完全全合而為一，讓世人知道是你差我來的，也知道你愛他們，如同愛我一樣。 \end{tabularx} \\ \\
[1ex]
\hline
\hline
\end{longtable}
$^{1}$我們剛才說,耶穌在離開這個地球上.
為我們,你跟旁邊的人說.
我們是耶穌禱告的對象.
我們每個人都是耶穌當日.
上十字架之前禱告的對象.
那禱告是說什麼呢?.
它說到要我們合而為一.
然後說到這個合而為一的根源.
我們能夠合一的根源是什麼?.
就是我們雖然每個人都不同.
真的每個人都不同.
但是我們能夠合一是因為.
我們全部同在基督當中.
我們一起來讀17章21-23節.
1,2,3.
使他們都合而為一.
正如你附在我裡面.
我在你裡面.
使他們也在我們裡面.
叫世人可以信你差了我來.
你所賜給我的榮耀.
我已賜給他們.
使他們合而為一.
將我們合而為一.
我在他們裡面.
你在我裡面.
使他們完完全全的合而為一.
叫世人知道你差了我來.
也知道你愛他們.
如同愛我一樣.
我現在讀好像很自然.
其實當我預備這個信息的時候.
我在家裡大聲朗讀.
耶穌這個禱告.
但是我讀完之後不知道為什麼.
我心裡有很大的感觸.
然後我就哭了.
我跟主說.
你為我們祈禱合一.
我心裡有很大的感觸.

$^{41}$當我們回到歷世歷代.
職教會所發生的事情.
當我看到我們自己裡面的困難.
當我看到我自己.
多麼不能夠守護這個合一的時候.
我就哭了.
跟神說.
主啊求你寬恕.
但是求你幫助.
因為這是你的禱告的時候.
你一定會幫助我們.
也是你的心意.
所以我就說.
我要由自己開始.
因為要講這篇講宗的時候.
我要由我自己開始.
來學習合而為一.
這裡說到.
三次要叫信徒合而為一.
所以你看到多麼懇切.
一叫他們合而為一.
叫他們完完全全的合而為一.
你感不感受到.
當你今天讀這段主耶穌的禱文.
這真是主禱文.
之前那些都是教門徒怎麼禱告.
這真是主自己的禱文.
你感不感受到.
主耶穌的那個心懷.
還有你感不感受到.
祂對你的呼召.
你要合而為一.
你要去合而為一.
其實在這裡.
主耶穌就用了祂和父神.
祂沒有用其他東西做例證.
祂就用祂和父神的關係做一個例證.
就告訴我們.
其實祂期望我們信徒.
我們信祂的人之間.

$^{81}$那個關係的合一是怎樣.
就正如祂說的.
好像祂和父神一樣.
完完全全的合一.
那種合一不是說.
我和你做朋友算了.
其實都是很表面的.
而是說生命的那種連結.
我們怎樣能夠連結.
其實我們每個人都不同.
但是我們能夠連結起來.
因為我們好像天父和耶穌一樣.
是連結.
生命是連結在一起的.
所以如果我們只見到彼此.
其實真的很不一樣.
但是我們能夠見到我們同一樣的東西.
就是我們都是同歸主.
我們全部都是一起來信神的.
所以耶穌說的合一.
是一個生命的連結.
與神的聯合.
不是說我們性格像不像.
我們做事的方法像不像.
我們興趣一不一致.
其實完全不是在說這件事.
是在說我們有同一個神.
那你跟別人說.
彼此鼓勵.
我們有同一個主.
這件事是非常重要的.
因為我們曾經.
可能大家都有參加過小組.
我也有參加過.
我去過一個小組.
其實開頭的時候真的很難.
因為大家真的不一樣.
有些人很多東西要說.
整場都要他說.
有些人你叫他.

$^{121}$你分享一下.
沒事的.
每次都說沒事的.
有些人可能信了主很久.
每件事都說主是這樣說的.
我們應該是這樣的.
但是有些人.
新信的人不太知道.
說些東西出來.
有些人就說.
你這麼不熟悉.
不知道怎麼熟悉.
都還沒學習.
有些人就說.
很熟悉.
分享到上天.
有些什麼很厲害的經歷.
有些人就說.
你上了天.
我不知道你在說什麼.
因為我沒有經歷過.
其實很艱難.
整組人混在一起.
看到組長爭執.
怎麼辦.
搞了這麼久.
人家的關係都很疏離.
有些人回不回.
我不喜歡這個.
底屍底下說.
我不喜歡說話.
不想讓我知道.
搞來搞去.
怎麼搞都不行.
組長後來祈禱.
想想辦法.
怎麼可以做.
不如這樣.
我們回去主的教導.
為什麼要有小組.

$^{161}$為什麼我們這群人.
要在一起.
第一.
在一起是因為.
我們同屬一個基督.
因為我們都是信主的.
我們沒有什麼次.
唯一次就是信主.
我們信同一個神.
所以我們能夠走在一起.
是因為我同屬一個主.
當大家有這個概念.
我們同屬一個主.
其實我們接受.
其實你和我都不一樣.
你又不一樣.
我們個個都不一樣.
但是我們怎麼可以.
一起去追求.
我們怎麼一起去.
認識基督.
我們一起去盡心.
學習彼此怎樣一起走.
這個就是合一的開始.
就是你aware.
你知道.
原來我們是同屬一個主.
特別是同輩當中.
然後我們.
我發覺很開心.
雖然沒有一下子.
整組人的關係.
馬上變得很好.
都很難的.
個個都不一樣.
但當我們有同一個目標.
追求的時候.
你會發覺大家的心開始近.
由這個點開始越來越近.
然後大家就會開始.

$^{201}$我們怎樣一起去追求神.
然後你會見到的東西.
就是很特別.
也是非常重要.
我們的心.
我們的生命開始改變.
因為本身我看到別人的不一樣.
你這樣不喜歡.
你這樣不合適.
但我們開始見到.
我們怎樣要去接納人.
學習.
每一個人的生命結果都是有進步.
我們對人的接納.
然後你會發覺.
其實有差異是很好的.
不同真的可以和和的.
如果你發覺不同的人.
好像一個orchestra.
一個樂隊.
每個樂器都不同.
但我們如果一起聽著主.
來指揮我們.
其實我們可以很厲害.
就是奏出來的音樂.
是可以很好.
所以.
這個就是主叫我們看.
你不要看著大家的不一樣.
你要看著大家的一樣.
是甚麼.
而在當中學習.
然後.
主耶穌說.
其實我們合一是一個使命.
就是讓人先來到.
去見到.
祂是副差來的.
回到剛才那一章.
讓世人知道.

$^{241}$從我們的合一見到.
耶穌真的來過這個地上.
是真的把這個世界上不一樣的愛.
賜給了我們.
在我們當中.
所以合一其實是一個見證.
如果我們這班人是一班信徒.
我們說自己是信徒.
但經常不停吵架.
吵架.
你不信我 我不信你.
你想想其實信主的人.
見到我們這班人.
其實他會問.
這班人有甚麼特別.
又說自己是信主.
又說耶穌是愛.
我也試過有人跟我說.
見到這種情況.
他跟我說.
其實我們都比你們好.
我們都比你們好.
多傷心 他這樣說.
但我又心想.
我們的確被人這樣見到.
是不是.
我的確被人這樣見到.
那怎麼辦.
所以我們要學習.
我們還是在學習的過程當中.
所以 弟兄姊妹.
我們當說合一的時候.
一定不是看看我跟你有多合.
旁邊的人跟我像不像.
我喜不喜歡他.
不是說這些.
是說從基督開始.
我們都是同屬基督.
第二件事.
對不起 下一張.

$^{281}$然後真正的合一.
不是說我喜不喜歡.
我喜不喜歡.
旁邊的人.
而是我們信服.
我們願意一起.
信服在主的心意當中.
合二為一.
你說:是的 我也知道要信服.
這段聖經我讀了很多遍.
為什麼總是做不到呢.
其實真的很難.
真的很難.
我發覺我自己去研究.
或者看回聖經.
其實為什麼我們這麼難做到.
最主要都是.
我們生命未被神完全改變.
同不同意呢.
我們每一個人都有自己的驕傲.
包括我自己.
我們當然有些事覺得是對的.
我們當然覺得自己是對的.
有些事應該是這樣的.
我們每個人都是這樣的.
誰舉手覺得自己沒有驕傲.
嘩 真的沒有人舉手.
你可能覺得.
不過我問誰覺得自己謙卑.
一開始我都以為自己很謙卑.
但想一想.
我真的不是很謙卑.
因為去到某個位置.
我都會起哄的.
是不是.
所以呢.
我們真的要祈禱.
神治死我們的驕傲.
我們才能夠合二為一.
求主幫助我們.

$^{321}$因為驕傲.
令我們不單止聽不入耳別人的話.
更加神的話都聽不入耳.
自我膨脹.
是不是.
神你說什麼就你說.
我是不行就不行.
神完全明白我們是不行的.
所以我們要祈求神的幫助.
好了.
其實聖經裡面.
神都有透過保羅來告訴我們.
鼓勵我們.
怎樣守住這個合一.
其實真的要守住這個合一.
其實是怎樣去守呢.
其實我們先來讀.
以弗所書 保羅所講的東西.
他講了幾樣態度.
我們來讀.
1 2 3.
我為主被囚的勸你們.
既然蒙召行事為人.
就當與蒙召的恩相稱.
凡事謙虛 溫柔 忍耐.
用愛心互相寬容.
用和平彼此聯絡.
竭力保守聖靈所賜.
合而為一的心.
這個就是剛才我們有詩歌的歌詞.
其實我很喜歡這些詩歌.
因為將神的說話用來敬拜祂.
並且將祂引入我們的心中.
作為我們聖靈裡面時常的引導.
這裡就講到.
保羅告訴我們.
其實神提醒我們.
我們今天信了祂的人.
你記住.
不要將信仰與你的生活分開.

$^{361}$我們查經 查經 查經.
知道很多聖經的知識.
但是就將它停留在這裡.
然後在生活裡面.
其實是沒有應用到的.
或者我的生命根本沒有被這些說話.
神的說話.
來影響到 改變到.
以致我們講一套.
然後就做一套.
做的時候.
其實都是用世界的方法.
來去想 去做.
這裡有講到.
我們要竭力保守.
聖靈所賜 合而為一的心.
有時候我們祈禱.
主啊 讓我們更加合一.
你賜下合一的心給我們吧.
其實更加正確的.
應該去求.
求你幫助我們.
竭力.
能夠竭力去保守.
這個合而為一的心.
因為其實這裡講.
聖靈已經賜下了.
這個合而為一的心給我們.
神已經做了祂的工作.
但是問題就是.
我們這群信祂的人.
有沒有負我們自己的責任.
來去竭力保守.
這份應該在我們裡面.
那份的心.
你跟隔壁的那個人鼓勵他.
努力都不夠.
要竭力.
我們要竭力.
你以為我很努力.

$^{401}$但其實我們已經很努力了.
但為什麼努力都不行呢.
其實是要竭力.
神說要竭力.
因為其實.
合一真是不容易的事.
你越多人就越難合一.
越多意見.
總之凡是有人的地方.
都會有不同的意見.
是不是.
所以一定要竭力.
那好了.
保羅這裡有幾樣東西.
凡事要謙虛 溫柔 忍耐.
用愛心 互相寬容.
那我在預備這篇讀的時候.
我就自己去檢視自己.
因為我覺得.
如果我在說這篇訊息的時候.
第一個要改變的是我自己.
不然說來幹什麼呢.
那我就跟神說.
你幫助我.
一直去檢視一下我自己.
是不是謙卑.
是不是溫柔.
是不是忍耐.
那好了.
我自己檢視了.
我不告訴你結果是什麼.
我只是求神幫助我.
更加進步.
但在這一刻.
我想跟大家.
我們每講一樣的時候.
你可以求聖靈幫助你.
想想你自己是不是一個.
有這些態度的人.
好不好.

$^{441}$第一 謙卑.
謙卑就不是說要自卑.
告訴別人什麼都不懂.
其實神給了我們很多東西.
我們都懂很多東西.
但謙卑的意思是.
我不看自己過高.
不會說我什麼都是對的.
所以你就跟別人爭辯.
其實有時候.
我們同輩的相處特別下.
其實你覺得我這麼高這麼大.
我就是要證明我自己是對的.
所以吵架都是說.
證明我自己是對的.
我是對的.
如果到最後.
其實你想證明自己是對的.
但別人始終覺得.
我又對.
認同不了你的時候.
你就會想.
你不尊重我.
我不被尊重.
我走了.
但你會想.
是不是一定你全對呢.
其實神眾我們每個人都不同.
其實想一件事.
大家有千萬個方法想出來.
其實謙卑的意思是.
我們要願意聽別人的說話.
承認其實我自己也有限.
當我們要聽別人的意見的時候.
不立刻反駁.
別人做一件事.
你立刻說不是的.
我是這樣的.
什麼什麼.
其實不需要.

$^{481}$我們聽聽別人怎麼說.
當別人給一些做法跟你不同的時候.
不用立刻反駁.
不用了.
聽聽別人怎麼說.
然後從中去學習.
原來真的有人這樣想.
原來可能真的有些方法.
是比我的方法好.
縱然可能最終是你的方法是最好.
但你都可以去學習去聽.
這個謙卑就是一樣.
其實謙卑說的就是.
放下自己.
不用爭取.
不用告訴別人我是對的.
接受別人也可能是對的.
你要發覺我們經常說什麼屬靈人.
其實真正屬靈的人.
是真的沒有什麼自己.
有些人會說.
我是這樣的.
你不要叫我改.
改不了的.
有沒有?.
有時我也會這樣說.
改不了的.
是 改不了.
如果耶穌叫你改.
你改不改?.
是不是?.
其實真的放下自己.
耶穌是生命的主.
好了 溫柔.
溫柔很多時候.
會被人誤解是很軟弱.
很靜 沒什麼主見.
不是說很多話.
每次問你意見你都可能.
哎喲 說著說著.

$^{521}$沒有堅持.
其實溫柔不代表沒有主見.
溫柔不代表沒有自己的想法.
自己的意見.
不過是為了上帝的緣故.
他不會不斷地狂拋自己.
告訴別人.
唯我獨尊 我是對的.
溫柔的人.
可以懂得控制自己的情緒.
反應 對人的反應.
所以很多時候溫柔的人.
譬如有些事情大家不對.
那個人會說.
那你想怎樣?.
你有沒有試過這樣跟別人說話?.
其實也有.
那你想怎樣?.
如果你不是溫柔的話.
你會怎樣回答?.
那你想怎樣?.
通常都是這樣的.
但溫柔的人可能會說.
其實大家先談一下 好嗎?.
其實會用柔和.
他不會那一刻情緒大爆發.
不會因為別人的反應.
我就去頂撞別人 報復式.
而溫柔的人是.
打下來 唔.
然後慢慢跟你說.
用柔和來化解一些剛硬.
其實就是這樣.
同輩之間有時候.
其實我們沒有什麼.
大家好好的.
但有時候就是因為說話.
我大聲不代表我沒有禮貌.
但我想告訴你.
這個是你自己覺得而已.

$^{561}$大聲不代表沒有禮貌.
因為一般人都覺得.
大聲是沒有禮貌.
所以當你大聲跟別人說話.
別人就會覺得你沒有禮貌.
你不覺得有什麼.
你自圓其說.
但別人已經覺得受傷.
你這樣跟我說話.
所以在教會當中.
弟兄姊妹.
我們說話的態度.
聲線 語氣.
其實都非常重要.
大家看看自己有沒有忍耐.
覺得自己有忍耐的舉手.
有人舉手 有人鼓掌.
耶.
我們要有膽量.
我真的有就有.
忍耐是什麼呢.
哥林多前書十三章說.
愛是恆久忍耐.
忍耐即是愛.
有一個英文譯本.
把這句譯成.
Love is long suffer.
Long suffer.
聽到之後都想走.
誰想long suffer呢.
但我想告訴你.
這是事實.
有哪種愛不是long suffer.
養兒一百歲.
長憂九十九.
長憂.
Long suffer.
是不是.
其實你發覺.
原來你正在long suffer.

$^{601}$有兒女.
老師對學生都是.
各樣.
其實所有的愛.
都是在說一個long suffer.
所以主說.
今天我這個愛.
是要你們接受到長期的受苦.
你才可以忍到.
其實很辛苦.
做基督徒.
但不是.
當你這個long suffer去忍耐.
你是有神的愛在背後.
去back up的時候.
你會覺得.
其實都不是那麼辛苦.
因為你的生命被改變了.
你不會那麼容易被人激怒.
那麼刺激.
還有你會發覺.
主啊因為你那麼愛我.
我可以去愛人.
我這麼多的問題.
你都愛我.
其實我都可以愛人.
是不是.
當你有主的愛去back up的時候.
你會發覺.
這個愛.
這個suffer.
其實是沒有想像中那麼suffer的.
甚至你會開心.
因為我們是可以跟隨到主的心意.
是不是.
我們信服神的心意.
在我的生命裡面.
這樣.
我們說.
其實忍耐的意思就是說.

$^{641}$我們不是一有什麼事.
就馬上.
別人覺得別人做得不好.
馬上黑臉.
馬上說別人.
其實這個是說.
願意給時間別人去改變.
更加就是給神去做事.
因為其實神要改變一個人的時候.
都要花時間.
一下子.
一刻就可以改變.
好了.
大家又看看你有沒有寬容.
因為說.
用愛心互相寬容.
用和平彼此聯絡.
寬容是什麼呢.
有沒有人說你寬宏大量.
哎呀 你真是寬宏大量.
有沒有.
說沒有的.
有沒有人說你小器.
哈哈哈.
是的.
寬容是什麼呢.
就是.
其實如果我們認識這個福音.
你是經歷過福音的時候.
我覺得你做人會比較寬容一點.
我信了主之後.
都比我未信主之前.
寬容.
因為.
其實當我們知道這個福音.
就是主無限的愛.
無限的寬容我們.
是不是.
有沒有一些事情是神叫你改的.
你到今天都沒有改.

$^{681}$都可能有的.
其實祂就在寬容我們.
是不是.
今天神說我不積毀滅你們.
因為我要人人都信主.
不如一人沉淪.
其實神都是寬容.
但不是說縱容罪.
寬容完全都不是縱容罪.
只不過是知道.
其實人都是不完美的.
我們經常都會.
其實你不要說別人.
你經常捉住人.
你又這樣.
你又錯.
這裡錯.
那裡錯.
但你有沒有發覺自己都經常做錯事.
我們經常都不完美的.
是不是.
我們都做錯事.
我們都做漏.
做得不好.
要明白其實.
寬容的意思就是.
大家都清楚.
其實我們所有人.
生命都在成長的當中.
我們不是完美的.
只不過每個人要進步的地方是不一樣的.
所以這個就是成長.
所以我覺得在教會裡.
我很被寬容接納的.
特別是我事奉敬拜隊.
有時候沒有很好的音樂根基.
有時候唱音又唱歪.
唱得不好.
入音也不好.
有時候拿捏不到詩歌的節奏.

$^{721}$我覺得弟兄姊妹對我很大包容.
我這個服侍一直都不是一個壓力.
如果你想一下.
她經常在挑剔.
整隊人你又這樣你又那樣.
搞錯了.
你怎麼做Worship Leader.
你有什麼資格.
慘了.
我想我都真的.
不知道怎樣能夠去服侍.
會變成很大很大的壓力.
事奉.
但是他們的包容提醒.
讓我越來越去學習去進步.
知道自己有什麼做得不好.
就一路去做.
結果那個感覺就是.
我們互相的保護.
我們不是不提意見.
但是我們怎樣去提.
怎樣能夠包容彼此.
讓我們越來越合一.
用和平彼此聯絡.
我記得有一次.
我們有一班人要去解決一件事情.
解決的事情出現了不同的意見.
有人說.
可以了.
不要這麼快.
出戰速決.
最好快快搞定它.
但是有人就說.
不要吧.
好像這樣做比較不近人情.
其實我們也要顧一下別人的感受.
其實爭論多久也爭論不完.
但我們基督徒有一個很好解決問題的方法.
是什麼呢.
就是一起祈禱.

$^{761}$之後大家都變了.
又不想一拍兩散.
就一起祈禱.
祈禱完之後.
有一個人就說.
我還是不堅持我的意見.
我放棄.
你們想做都可以的.
那些人就會想.
其實我那一刻也想.
哈哈.
當你祈禱的時候.
主告訴你.
你錯了.
是你想錯了.
你的方法是錯的.
主光照你.
我們會這樣想.
但是後來又問他.
其實你這麼堅持.
之後會怎樣.
祈禱之下.
然後他就說.
其實不是.
是因為當他在禱告的時候.
主真的光照他.
光照他就跟他說.
你覺得你贏了比較對.
還是合一比較好.
你要什麼.
他說他選擇合一.
所以他之後.
他放下自己的意見.
但神很奇妙的.
其實合一永遠是有一方面願意謙避下來.
去保守.
來守住這個關係.
主動修補.
明知已經開始有裂痕.
主動封了裂痕.

$^{801}$這就是用和平彼此聯絡.
竭力保守.
聖靈所賜.
合二為一心.
不被魔鬼殺.
但在裂痕那裡.
有地步.
讓我們因此更加分裂.
這就是竭力保守.
這裡每一個人見到他的舉動.
他這麼謙避下來.
其實大家本來都很支持的.
都軟化下來.
其實最奇妙的是神.
當合一臨在的時候.
不是說表面合一.
像這些真正的合一臨在的時候.
神結果沒有用任何一方之前提的方法.
是用他再次下一個方法.
很完美地解決了這件事.
你會知道我們不是要想方法.
而是想怎樣討好神.
然後神會將他的旨意降臨在我們當中.
教我們怎樣做.
所以爭是沒有用的.
合一才有用.
然後神說.
你們不要看著自己個個都不一樣.
合一不是要排除異己.
只是剛剛談過的人在一起.
而合一是要接納差異.
還有你要記住.
我們有一些東西是一模一樣的.
我們再看看.
《以弗所書》四章四至六節.
那裡怎樣說.
身體.
我們這班信主的人.
我們是一個身體.
聖靈是一個.

$^{841}$一個指望.
一主.
一信.
一洗.
一神.
這裡說的七個一.
靈姐妹你看看旁邊那個.
這就是我們七個一樣的東西.
是遠超過我們人與人之間的不同.
其實神說.
我們每一個都是奇妙可畏去創造的.
你要求統一.
沒有可能的.
不會統一的.
他也不會統一我們.
他不是要機械人.
我們每一個有生命.
我們有氣息.
他給我們不同的性格.
但他說你們是可以合一的.
因為你們有這七個一.
你們要記住.
因為這七個一的緣故.
你們就要竭力去保守.
竭力來接納彼此之間的差異.
其實真的很難.
其實我相信真的很不容易.
但是我們相信.
如果當我們這班群體.
我說跟未信的各樣.
可能有少少難.
但如果我們這班人.
都知道有這七個一的時候.
如果我們全部都回到聖經裡面.
我們回到神的教導.
我們每一個都說.
因為主的緣故.
我願意去做的時候.
弟兄姊妹我們真是合一在望.
不是一件完全不能做到的事.

$^{881}$雖然我們今天看回這個情況.
其實距離真正合一仍然很遠.
所以記住弟兄姊妹.
合一不是叫口號.
我們要合一.
然後說你合一 你不合一.
你不合一.
其實你問問自己.
我合不合一.
不是說主求你幫助我們.
而是你真是願意擺上我們的責任.
我們的代價.
來竭力保守.
聖靈已經做了它的工作.
賜下這個合二為一的心.
責任保守這個合二為一的心.
是在我們所有這班人一起去做.
這是我們的份.
所以求主幫助我們.
一起來合一.
抵擋魔鬼在當中分裂的工作.
最後.
我們為什麼這麼肉緊地說.
我們要追求合一呢.
有什麼這麼好.
對我有什麼益處.
這麼辛苦.
是不是.
有什麼益處.
其實合一從來不是在說我們自己的益處.
合一從來不是在說我舒不舒服.
而是我們要養神得榮耀.
合一是同心榮耀神.
我們是神的創造.
我們被造的目的.
是要養神得榮耀.
所以我們一起讀.
這段羅馬書十五章五至六節.
但願賜引來安慰的神.
要你們彼此同心.

$^{921}$效法基督耶穌.
一心一口榮耀神.
我們主耶穌基督的父.
這裡說我們要一心一口榮耀神.
雖然我們有不同的想法.
不同的立場想事情.
但是我們要學習.
最終我們是怎樣能夠.
一心一口.
來榮耀我們的主.
其實你要記住.
合一真的不是我們自己舒不舒服.
合一是說要放下自己.
真的不舒服.
但是合一我們要知道.
最終是讓創造我們的神.
我們的主耶穌基督.
祂的名被高舉.
耶穌是主.
很立體很具體地.
在我們這群人的同行當中.
來讓人們看到.
然後將福音.
是好人親眼看到.
福音就在這群人當中.
所以我們就要努力.
另外高林多前書一章十節我們就讀了.
弟兄們.
我直我們主耶穌基督的名.
勸你們都說一樣的話.
你們中間也不可分黨.
只要一心一意彼此相合.
同輩之間不可以分黨.
是很重要的.
因為有時候我們年紀相若.
有些事情經常一起相處.
變了很多時候.
其實都出現了一些事情.
分黨或者拉黨.
結派.

$^{961}$其實有時候.
不是刻意的.
但不知為何.
搞著搞著就覺得.
我們這群人.
我們自己這群人.
其實他說.
我們教會有這麼多牧區.
那是不是分黨?.
這麼多黨.
這裡有這麼多名.
其實不是.
是因為牧養.
行政的緣故.
是不是?.
意思是我們要將人們分開.
方便牧養.
行政上.
但是我想告訴你.
我們不是屬哪個區.
只是屬哪一個隊.
而是我們全部都是.
同屬基督的.
其實我們.
不同的牧區之間.
我們都彼此相愛的.
我們是弟兄姊妹.
不同的隊伍.
我們都是彼此配搭的.
我們不是屬於這個區.
屬於那個黨.
屬於外面已經有很多黨.
但我們這些不是黨.
我們是彼此相愛的.
雖然有不同的名字.
不同的區.
但我們是同屬基督.
好了.
如果我們真的.
不是這樣想.

$^{1001}$裡面都是結黨結派.
我們這班人自己為內.
我們那班人自己為內.
那班人自己為內.
外面的人去見到的時候.
他們說.
其實你們口講合一.
一進來其實都四分五裂的.
各自各的.
其實都是做不到見證.
還有.
其實很多人.
見到我們之間的不一樣.
但他們很開心見到我們.
原來你們個個都不一樣.
但你們可以彼此配搭.
你們可以彼此補位.
很厲害的人又很謙卑.
但有些人做不到一些事.
你又不會說.
他只會走過去補位.
其實人就是這樣見到福音.
他不是我們最後那個結果出來.
多麼冠冕堂皇.
而是真正在我們裡面.
見到福音的人.
是見到我們中間的人是怎樣相處的.
所以.
靈姐妹.
福音本身就是宣教.
福音.
這個合一本身就是宣教.
合一其實就是福音的一部份.
今天我們合一.
其實就是全靠福音.
大家明白嗎.
所以主說.
我們在講的合一.
不是那種表面.
不說了不說了.

$^{1041}$但其實裡面.
每個人的心很拉扯.
互相防範.
不喜歡他.
其實我不喜歡他.
而是我們真正.
我們的生命.
被神去改變.
被福音去改變.
當我們的生命.
每個人都被福音去改變的時候.
合一就會出來.
是不是.
好了.
講到這裡.
大家會不會覺得.
其實我們的信仰.
和這個世界很不一樣.
當這個世界.
人們越來越分化.
你看到很多戰爭.
每個人都要自己贏.
看到很多人.
甚至乎互相踐踏.
人越來越自我.
我們有時候很難過.
但其實.
原來我們都參與其中.
所以.
求主幫助我們.
我們能夠.
和這個世界不一樣.
當這個世界越來越分化.
我們的主告訴我們.
我們要謙虛.
溫柔.
仁愛.
用愛心互相寬容.
用和平彼此聯絡.
我相信當這個世界.

$^{1081}$就算怎樣分裂都好.
我相信一個真正.
能夠竭力保守聖靈合一.
而唯一的群體.
我們在這個世上唯一.
能夠真正有愛.
還能夠見到愛的群體.
Amen.
我很想.
我們真的能夠在合一的當中.
來到.
活出.
在這個分裂的世界裡面.
活出一個合一的見證.
你鼓勵一下隔壁那個.
我們要一起.
在這個分裂的世界裡面.
活出這個合一.
我們彼此鼓勵.
最後.
有幾樣東西給大家帶走.
記住我們都在基督裡面.
我們所有一切做的事.
要土神氣月.
不要爭東西.
不要覺得隔壁那個跟你不一樣.
你不願意和他在一起.
基本上我們個個都不一樣.
求主幫助我們.
可以和不一樣的人去同行.
我們要操練.
我們的生命未完全被改變.
我們要繼續操練我們的謙卑.
溫柔 仁愛.
操練來溫柔控制我們的情緒.
不要讓情緒破壞我們和人之間的關係.
之後記住.
我們的同行.
是要一切所有人榮耀神.
不是為自己建立一些什麼勢力.

$^{1121}$為自己建立一些什麼王國.
然後大家記住.
最後一點我想說.
因為可能在信仰的過程裡面.
我們不多不少.
可能都經歷過一些挫敗.
是不是.
關係上的挫敗.
其實我自己都經歷過.
有時候你會想.
好難.
但是我求主醫治我們.
我們經過創傷的.
求主醫治我們.
並且讓我們有個盼望.
知道今天我們這班人.
因著回到主的真理的緣故.
我們是能夠合一的.
Amen.
\newpage



\section{出埃及記 19:1-6}
\label{sec:031972}
\textbf{信徒皆具雙職的身份}~~~~丘放河牧師博士
\newline
\newline
連結: \href{http://www.church.com.hk/acms/content.asp?site=cdc\&op=show\&type=product\&code=031972\&layout=sermon}{\texttt{031972}} ~~~~ 語音日期: 2026-04-12
\newline
\newline
\hyperref[sec:031971]{\small{< < < PREV SERMON < < <}}
~
\hyperref[sec:index_chronic]{\small{[返順時目]}}
~
\hyperref[sec:031973]{\small{> > > NEXT SERMON > > >}}
\newline
\newline
出埃及記 19:1-6
\newline
\begin{longtable}{cl}
\hline
\hline
章節 & 經文 (和合本修訂版)\\
\hline
19:1 & \begin{tabularx}{0.7\textwidth}{X} 以色列人出埃及地以後，第三個月的初一，就在那一天他們來到了西奈的曠野。 \end{tabularx} \\ \\ \relax
19:2 & \begin{tabularx}{0.7\textwidth}{X} 他們從利非訂起程，來到西奈的曠野，在那裡的山下安營。 \end{tabularx} \\ \\ \relax
19:3 & \begin{tabularx}{0.7\textwidth}{X} 摩西到神那裡，耶和華從山上呼喚他說：「你要這樣告訴雅各家，對以色列人說： \end{tabularx} \\ \\ \relax
19:4 & \begin{tabularx}{0.7\textwidth}{X} 『我向埃及人所行的事，你們都看見了， 我如鷹將你們背在翅膀上，帶你們來歸我。 \end{tabularx} \\ \\ \relax
19:5 & \begin{tabularx}{0.7\textwidth}{X} 如今你們若真的聽從我的話，遵守我的約，就要在萬民中作屬我的子民，因為全地都是我的。 \end{tabularx} \\ \\ \relax
19:6 & \begin{tabularx}{0.7\textwidth}{X} 你們要歸我作祭司的國度，為神聖的國民。』這些話你要告訴以色列人。」 \end{tabularx} \\ \\
[1ex]
\hline
\hline
\end{longtable}
$^{1}$各位親愛的弟兄姊妹們 早安.
沒想到在三十多年前.
能夠有機會來到史雲頓華人基督教會.
那時應該還是團契吧 Andrew.
那時還是團契的時刻.
你們應該在Freshport Evangelical Church的時候.
就有這個團契.
我又有機會帶著我兒子.
和我父母一起來這裡.
我這次來到這裡算是懷舊了.
為什麼呢?.
因為完全不知道自己來過這裡.
我都忘記了.
如果不是梁銓道有機會.
來跟陳銓道說這件事.
我也不知道自己曾經來過這個地方.
這個城市.
所以在Doncaster第一次坐火車.
再坐濕圍 再坐火車下來這裡.
很感恩.
真是難以想像.
原來人生裡真的有很多很多的改變.
在我們的人生裡.
那時梁銓道和我都很年輕.
眨眨眼睛 白頭髮來了.
所以看到.
我們今天如何在屬靈生命裡成長.
如何在生活裡繼續見證神的恩典.
正是你們教會今年的主題.
如何屬靈恩典來發展.
我坐火車下來 坐地鐵下來.
但很多時候都是在一個地方.
這個地方是什麼地方呢?.
對了 機場.
因為我經常一年跑四次回香港.
當然要回香港坐飛機.
然後又去台灣 去其他地方的時候.
很多時候我都會在這些機場.
在機場的時候我都會有一個感覺.
究竟這些人在做什麼呢?.

$^{41}$會不會像我這樣出差呢?.
還是去旅遊呢?.
究竟他們的身份是什麼呢?.
我們基督徒是不是只有一種身份呢?.
當我們一路去想的時候.
我們就說究竟我在這個世界裡.
是不是只在教會裡出現的一種身份.
究竟我在其他地方有沒有另外一些身份呢?.
我想一下 這是我們基督徒有的一些不同的身份.
所以我今天來跟大家分享的時候.
說相織的時候.
就是說我們有兩個不同的織法.
不同的身份.
如果我們在生活裡有醫生護士會計師.
做不同的職場裡的不同的職業.
這些都是叫做不同的身份.
我自己和太太每個星期六在紐約.
就是叫做NYPD 做輔警的.
我們一路做輔警.
你看到要出勤.
要來幫助社區的民眾.
這個漂亮嗎?.
漂亮啊.
這個是師姐.
就是師母.
我們一路來到不同的地方.
來幫助民眾.
來做的過程裡.
有時候我在街上碰到一些基督徒.
他們看著我.
這個人很像邱牧師.
真的邱牧師啊.
然後他們說.
邱牧師你轉了行嗎?.
他們覺得這是一行.
做警察就是一行.
我做牧師又是另外一行.
究竟是不是我們不同的行.
不同的職業.
就是我們所謂的身份呢?.

$^{81}$記得2024年11月.
當我看到今天的經文.
出埃及記第19章的時候.
我就有一個很不同的感覺.
突然之間我發覺.
這個經文裡其實都很熟.
看了很多次了.
但是這次突然之間.
好像眼睛明亮的.
有一個很特別的現象.
我就在想.
他們出埃及.
出埃及的第三個月.
在這三個月裡經歷了什麼呢?.
他們出埃及的時候.
當然是離開了法老.
離開埃及比喻是世界.
然後藉著主耶穌的寶血.
我們現在知道主耶穌的寶血.
當時就是羊羔的血.
就天使如月.
來進行了第十個災.
然後他們就如月之出了西乃.
去到西乃的時候.
都不是一下子很開心.
因為沒得吃.
沒得吃就給嗎哪他們吃.
沒水喝.
神就叫摩西打石頭出水.
然後又沒肉吃.
神又從東風吹了鵪鶉來給他們.
有肉吃.
然後就遇到阿瑪利人.
他們就和他爭戰.
然後神叫摩西舉起手.
他就贏.
垂下手就輸.
在這個出埃及的第十七章.
但是到了第十九章這裡.
有一個很特別的現象.

$^{121}$神就和摩西說.
他們如果遵守我的約.
守我的誡命的話.
我給他們一個不同的身份.
是什麼身份呢.
祭司的國度.
什麼叫國度呢.
你來到這裡叫什麼國度呢.
United Kingdoms.
Kingdom就是國度.
而這裡說的是什麼呢.
祭司的國度.
祭司的國度裡面的所有人是什麼呢.
祭司.
明白嗎.
那你想想原來他們只要跟從神遵守他的誡命.
就擁有一個很特殊的身份.
叫做祭司的國度.
聖潔的國民.
是神的子民.
而在這個過程裡我們就看到.
祭司是一個目的的身份.
為什麼這樣說呢.
因為你一看到祭司.
特別是在舊約.
舊約有很多祭司.
在舊約的祭司你就會發現.
這些祭司是一種目的的身份.
如果這樣說.
所有跟從神的人.
聽他的誡命的人.
他就擁有一個很特殊的身份.
叫做祭司.
我稱之為目的.
這個目的的身份和今天的牧師職是不同的.
牧師職是有一個全職的工作.
我們有義務的.
好像梁銓道是義務的.
但是有壽身的.
這些是牧師職.

$^{161}$但是每一個信徒.
記住每一個跟從神.
聽他的誡命的人.
都擁有這個祭司的國度的其中一員.
是祭司.
是目的的身份.
所以當我們想到這個時候.
我就開始看完你自己的經文.
我就覺得不是的.
那就是怎樣的.
那祭司是怎樣來的.
我再看28章.
就說到神對亞倫說.
你和你的後代.
全部做祭司.
這是我們很熟悉的一個情況.
我們看到原來舊約裡面.
有一種似乎叫代言人.
他們來擁有這個目的的身份.
而那些以色列人.
那個目的的身份似乎怎樣.
是被隱藏了.
我稱之為被擱置了.
好像由這班代言人.
來做了那個以色列人.
應該擁有的身份.
這個很特別的一個感覺.
是我當時看聖經的時候.
我就看到這個改變的時候.
我說不是的.
本來應該每個人都有份的.
只要你聽從他的話.
守他的誡命.
但突然之間就由一班代言人來代言.
那我就說為甚麼會這樣呢.
甚麼原因呢.
那我又想一下.
舊約裡面的大祭司很有趣.
大祭司一年只上一次宮.
是屬庶人那天.

$^{201}$要進至聖所.
帶著香.
還要有條腰帶.
不知道那條腰帶是用來做甚麼的.
他死了要從裡面拔出來.
我發覺大祭司的工作.
多簡單.
拿著紅繩子繞著角.
然後趕他出去抗議.
就如表他們的罪.
就離開了他們.
屬罪.
一年只上一次宮.
在那個幕裡面.
進入至聖所.
這個很特別的一個現象.
我就一直思考.
祭司和以色列人.
擁有木職的身份.
究竟是怎樣的呢.
一直想的時候.
我就想到去了新約.
原來在新約裡面.
這個祭司的木職身份.
竟然人人都有份.
而似乎就是擱置了之後.
現在去到新約就激活了.
為甚麼這樣說呢.
因為當我們看到.
這個舊約裡面.
祭司的工作.
回望.
有聖所至聖所的過程裡面.
突然之間.
這個讓我看到.
一個很特殊的現象.
就是當耶穌棄一團的時候.
有些景象.
發生.
聖殿裡面的回望.

$^{241}$由上到下裂開了.
如果你在當時.
你會不會進去看看.
平時是不能看的.
只有大祭司可以進去.
但是現在裡面裂開了.
不要浪費了.
進去看看吧.
當然是看到約櫃.
約櫃有兩個基督派天使.
拱著.
翼翼在那裡.
有一個位置.
你知不知道那個位置是甚麼呢.
那個位置就是.
叫做私恩族.
私恩族在哪裡呢.
就是基督派.
拱著的翼.
下面的空間.
就是私恩族.
這個地方.
很特別的地方.
希伯來書的作者.
怎樣來描述.
我們可以坦然無懼地進入.
私恩族前.
來得連宿.
蒙恩衛.
作敘述的幫助.
弟兄姊妹們.
有一個目的身份.
激活了每一個信徒.
都擁有.
這個目的身份.
就是祭司的身份.
這個身份以前是有.
大祭司做代言人.
現在希伯來書的作者告訴我們.
不是的.

$^{281}$每個人都可以這樣進去.
厲不厲害.
這個身份突然之間.
很奇妙的.
就是因為主耶穌的寶血.
主耶穌的生命.
熟了我們.
我們就激活了我們應有的.
這個身份.
就是目的的身份.
所以我們想.
我擁有目的的身份.
我怎樣去履行.
我這個目的的身份.
我又不是傳道.
又不是牧師.
又不是全職.
帶領教會的人.
我怎樣履行.
這個身份.
每一個信徒.
只要你跟從神.
守他誡命的人.
都擁有這個目的.
這個身份.
如果我們從.
加里里胡編的一幕.
我們可以想到一個很特殊的現象.
你知道.
加里里胡編.
就是彼得當時怎樣.
《約翰福音》.
第20章.
描述.
當時耶穌復活了.
向他的門徒展現了兩次.
然後描述.
彼得帶著.
另外的門徒.
一起去打雨.

$^{321}$你想想.
這個場景.
這樣的描述.
耶穌復活了兩次.
但是彼得.
竟然還是帶著他們.
帶著他們去打雨.
是什麼意思.
就是說當時彼得的心境.
是怎樣的.
耶穌是他的夫子.
他跟了他三年.
然後不認他.
然後悔改.
悔改之後耶穌又復活.
展示第二次給他知道.
這一刻對彼得來說.
會怎樣想.
他就在想.
耶穌說拜拜了.
我走了.
去哪了.
耶穌說我是道路真理生命.
如果不直走.
沒有人能到庫去.
那是怎樣的.
說來說去都不知道他去哪了.
你明不明白.
彼得和那些門徒很茫然.
究竟我今天跟了他三年.
他又顯神蹟.
死的都可以復活.
但是現在是怎樣.
他說我是預備地方.
怎樣了.
說一半就說一半.
是很茫然的.
一茫然的時候.
一個人不知道要做什麼的時候.
做什麼最好呢.

$^{361}$做回本行.
做回老本行.
去打魚.
去打魚的過程裡面.
你知道.
有一個神蹟出現.
耶穌說打去右邊.
就有153條大魚.
耶穌怎樣.
在那一刻裡面.
展現了一個神蹟.
再展現一個神蹟給他們看.
彼得一下子跳進水裡.
走到岸邊.
有東西吃了.
耶穌在旁邊就預備了食物給他.
耶穌問了他三次.
同一個問題.
差不多.
你愛不愛我.
他的答案是.
愛.
耶穌給的答案是.
你餵養我羊.
牧養我羊.
這一幕.
很特別.
怎樣特別.
當一個信徒.
遵守他的約.
守他的誡命.
加上愛神的時候.
愛主的時候.
你就可以牧養.
我的羊.
各位.
今天我們來問問自己.
為什麼我展現不到.
這個目的呢.
這個目的已荊棘活了.

$^{401}$我們可以坦然無懼.
進到施恩座前.
有這樣的能力.
但是我們仍然未能去.
履行我們的目的.
以致我們要問自己.
What's going on.
發生了什麼事.
彼得.
和耶穌的.
邂逅在那裡.
警方告訴我們.
當一個人真心愛神的時候.
你才可以餵養.
牧養我的羊.
在這裡給我們一個很清楚的意思.
如果今天我們不愛神的時候.
我們沒有辦法.
去展現.
那個生命.
所以當你說.
靈命更新的時候.
當你的靈命不更新的時候.
你是沒有辦法.
去做牧養的工作.
所以在這裡.
我們就問另外一個問題.
如果在這樣的過程裡.
要愛主才能夠牧養.
餵養主的羊.
那麼我們的目的是什麼.
今天的目的是什麼呢.
保羅很清楚的告訴我們.
特別是彼得.
在彼得前書第四章已經馬上說出來了.
說了一個很特殊的.
就是.
屬靈恩賜.
究竟神給你什麼.
屬靈恩賜呢.

$^{441}$你就去做吧.
然後彼得無論在.
哥林多前書第十二章.
他講了很多不同的恩賜.
然後在這個.
裡面講到.
每個恩賜.
是一個肢體的發揮.
我們每個肢體有不同的功能.
而這個身體裡.
有很多的肢體.
這個肢體.
就是屬靈恩賜.
所以在這個哥林多前書第十二章.
講到很多肢體的.
運作的問題.
然後在羅馬書第十二章.
也都講了有很多不同的恩賜.
這些不同的恩賜.
是聖靈隨著.
他的心意來給我們.
如果我們擁有這些恩賜.
我們需要發揮出來.
以弗所書第四章第十一節.
也都講了有五個屬靈恩賜.
有使徒.
先知.
傳福音.
還有牧師和教師.
但是這個牧師特別要.
highlight來講一講.
這個牧師.
不是牧師織的牧師.
而是牧師.
是pastoring.
有牧羊的那種屬靈恩賜.
當這樣講的時候.
就出現了一個現象.
這個現象是什麼現象呢.
如果他是屬靈恩賜.

$^{481}$他不是一個職份.
不是一個職業的職份.
這個恩賜代表什麼呢.
更加能夠顯示到.
整本新約聖經裡面.
除了這裡講了中文的牧師之外.
沒有別的地方講牧師.
這表示什麼.
原來牧師這個職份.
是要屬有職業的.
特別是恩賜.
來去牧養.
但是我剛才講過.
我們每一個信徒都有牧職.
這個牧職.
這個祭司的身份.
是藉著屬靈恩賜.
來彰顯.
來表現在.
這個身體.
教會裡面.
所以這樣講的時候.
我們每一個人不一定.
有很多很多的恩賜.
但是.
我們錯誤的理解.
就是說.
將牧師職的人物.
是覺得他.
方方面面全面都可以.
你有沒有想過這樣的情況.
一個做牧師.
在教會牧養的時候.
他要上台.
下廚房去餵養.
又要給屬靈恩賜.
又給靈量.
又給聖體.
屬靈恩賜又給靈量.
還可以開車做司機.

$^{521}$還可以彈琴.
神學院.
舊式的神學院.
差不多每一個都要.
彈琴才可以畢業.
彈琴不畢業.
我都沒有屬靈恩賜.
音樂的細胞.
你叫我彈.
我連唱都會走音.
你怎樣可以來牧養呢.
如果要全方位的牧養.
這個導致了.
有一個數據.
就是說.
我看過一本書.
你也可以看我那本書裡面說了.
那本書裡面其中一個內容.
就是說牧師傳道人有.
兩到四成是有憂鬱症.
他不是有神經恩典嗎.
他不是應該很有滿有靈高的嗎.
為什麼他會有憂鬱症呢.
就是因為.
大家搞成我這樣.
你要我全方位.
樣樣都可以.
但是聖靈又說.
給兩個你吧.
我只有兩個恩賜.
那兩個恩賜我最厲害.
但是其他的都是很平庸.
但是你要全方位要求我.
我又不可以說自己不行.
那我又怎樣.
我就很委屈.
那就憂鬱症.
我舉一個例子.
你有五個屬靈的恩賜.
其中一個叫先知.

$^{561}$先知的恩賜是什麼.
先知是什麼恩賜.
有什麼表現.
是沒錯了.
先知好像預先知道一些東西.
但是也有其他功能.
譬如哥林多前書第十三章.
先知講道之能.
聽過沒有.
先知可以有講道之能.
先知能夠講出神的話語.
先知講道之能.
你看舊約就是.
那些先知除了預言之外.
還可以講神的話語.
譬如你知道有一個肅骨先知.
叫若拿.
他就縮到魚肚裡.
叫若拿他就縮到魚肚裡.
他就不肯去尼尼未成.
結果怎樣.
神就要在餘福裡.
令他反省之後.
他願意去了.
套他去岸上.
他就走去尼尼未成.
你知道他講了什麼嗎.
如果你看中文聖經是十四個字.
十四個字講完就收檔.
這個就是先知.
講道之能.
假設你有先知.
講道之能.
除了有能力去預知之外.
你還要講道.
但如果我沒有先知講道之能.
我測試過沒有.
我去教會申請.
我來申請做牧師職.
你說好啊.

$^{601}$歡迎啊.
不過有一樣東西我是不講道的.
你請不請他.
請不請.
當然馬上.
走啦.
你就明白.
我為什麼說憂鬱症.
因為原來我恩賜.
我有敬拜的恩賜.
有教道的恩賜.
有音樂的恩賜.
知而處理的恩賜.
就是沒有先知的恩賜.
但是.
牧師不講道.
你是很難接受.
牧師要全方位.
但是如果一個牧師.
沒有這樣的恩賜.
行不行呢.
我自己在香港的教會的牧師.
正正是遇上這個問題.
他是對老人家很有恩賜.
他是對老人家關懷他們.
帶他們去看醫生.
陪著他們.
然後帶他們信耶穌.
一之百.
很好啊.
教會裡很多老人家就是因為這樣來教會.
但是因為他的屬靈恩賜.
對老人家.
在執事會的時候.
問牧師我們以後怎麼辦.
就是說.
面前一年三年五年要做什麼.
他不是行政恩賜.
他就說.
不如我們一起禱告.

$^{641}$你又不可以駁他嘴.
突然禱告.
但是.
不是吧.
明年怎麼樣.
後年怎麼樣.
給個指示帶領.
他就是對著老人家.
搞到會友很難受.
因為年輕人.
中年人都沒有人關心.
只是關心老人家.
那些中年人多數是執事.
不關心我.
覺得牧師不行.
但是又說不出口.
他就說禱告.
結果就請一個行政的人來幫忙.
結果就請一個行政的副牧師來幫忙.
結果怎麼樣.
行政那個.
越來越.
處理得很好.
那個就被解僱了.
憂不憂鬱.
肯定他在憂鬱.
所以你會發覺.
原來屬靈恩賜.
是要彼此配搭.
而不是.
獨當一面.
一個人最大.
一個人最厲害.
每個人都有各個功能.
全身都靠他聯絡得合適.
百折各按各職.
照著各體的功用.
彼此相助.
在愛裡面建立自己.
這個是以弗所書第四章.

$^{681}$第十六節告訴我們.
所以當我們這樣看的時候.
我們就看到.
這個目的是這樣.
但是我又.
信徒皆具備雙職.
第二個職是什麼.
就是專職.
專職不是專業的職業.
是說一個專有的職份.
當然當保羅面對著.
帖撒羅尼加的教會的時候.
他們就不做工.
只顧著享受.
他們那時候叫做.
販物公用.
販物公用這件事.
我去享受.
販物公用.
有些大食懶在那裡.
這些大食懶坐食三崩.
結果只顧著在那裡.
吃東西.
還要講閒話.
然後給保羅怎麼說.
不做工不得飯吃.
我說我都出去打工.
你知道保羅都是雙職.
他又是大職事奉.
他又積將彭.
又去傳福音.
如果按照我出來.
做牧羊的工作.
吃你們的飯是正常的.
對不對.
工人得工價.
牛在場搶曲是可以的.
但是.
為了你們的榜樣.
我都出去打工.

$^{721}$表示一個人.
他有一個專有的職份.
但這個專有的職份.
不單單是職業上的職份.
他說.
做什麼都要來榮耀神.
你想想.
當我們面對著.
今年的主題.
就是共證主因的時候.
你怎麼去證明.
神的恩典在你的生命裡.
你要表現出.
你自己是真的.
屬於神.
是被人看到你是怎樣的樣子.
所以當耶穌在最後晚餐的時候.
他是這樣.
跟那些門徒說.
我賜給你一個新的命令.
要你們彼此相愛.
當你們彼此相愛的時候.
就證明.
神的恩典在你的生命當中.
人家就看出.
你是我的門徒.
這一段的經文.
是有些背景在裡面的.
就是去吃.
如意節晚餐.
吃如意節晚餐的時候.
誰坐在耶穌旁邊最尊貴的位置.
就是約翰.
但約翰是.
小伙子.
小孩子.
彼得就欽他.
就設一個位置.
讓他坐在耶穌最尊貴的位置.
彼得是大師兄.

$^{761}$他就想著.
我就坐在最後的位置.
最後的位置.
就是當晚.
需要來幫人洗腳.
他心想我是大師兄.
坐在那位置.
耶穌就說大師兄.
坐在最尊貴的位置.
約翰就去下面幫人洗腳.
是這樣的想法.
雖然聖經沒有說.
我猜測.
但也不會離得開很遠.
總之彼得就坐在那位置.
還有十個門徒.
你喜歡坐那位置就坐吧.
你去洗腳吧.
沒有人出聲.
耶穌不出聲.
就坐在那位置.
結果彼得就站在那裡.
大哥.
我坐在這裡難道我洗腳嗎.
他就不肯洗腳.
耶穌怎麼做.
耶穌拿起毛巾.
去幫門徒洗腳.
然後洗完腳的時候.
我怎樣愛你們.
你們都要怎樣相愛.
然後賜給他們.
一條新的命令.
叫你們彼此相愛.
讓人看出.
你是我的門徒.
這個新的命令.
和舊的命令有什麼分別.
知道舊的命令是什麼嗎.
有新的命令.

$^{801}$當然有舊的命令.
舊的命令知道在哪裡嗎.
當然在舊約裡面.
就是要盡心盡性盡意.
愛你的上帝.
其次也相仿.
就是要.
愛人如己.
愛自己.
絕對沒問題.
你愛不愛自己.
當然愛.
愛人愛自己不是一個問題.
彼得也是這樣做了.
所以新的命令是什麼.
是彼此相愛.
如果在彼此相愛的時候.
一個教會裡面.
弟兄姊妹彼此相愛的時候.
外人走進來就看到.
哇.
真是很好.
這個教會很有愛心.
我想留下.
當你作為一個基督徒.
滿有主的愛的時候.
你走出去.
你們自己的工作裡面.
來去行事的時候.
你就會展現.
你被人看到.
你是基督徒.
是主的門徒.
我記得有一次有個姊妹告訴我.
她說我有一次在公司裡面.
就看到.
我的同事.
桌子上有本聖經.
我就問她你是基督徒嗎.
她說是呀我是基督徒.

$^{841}$那姊妹就說我也是.
那人怎麼說.
哎呀看不出來.
弟兄姊妹.
你現在的頭沒有寫著.
基督徒.
你走出去看看別人看不看得出來.
如果別人看不出來.
這句話.
正是告訴我們.
我們的尊有積分.
等於主的門徒.
你失去了.
你沒辦法展現.
你是主的門徒.
所以我們有目的在教會裡面.
展現我們的屬靈恩賜.
在教會以外.
我們要來去展現.
讓別人看到.
我是主的門徒.
你如果能夠讓人看到.
你有這個相職的身份.
目的和專職的話.
這就是我們信徒.
應該擁有的相職的身份.
而這個相職的身份.
教會裡面和教會外的.
展現是兩回事.
教會裡面是因為我們是肢體.
我們的肢體的功能.
就是屬靈恩賜.
所以我們在教會裡面展現出來.
但在教會外面.
無論是社區職場或者是學校.
或者是家庭特別是家庭.
原來法國有很多人.
不信耶穌就是因為我家裡有個基督徒.
明不明白.
那個基督徒就是你.

$^{881}$慘不慘.
就是因為你信了耶穌.
你的家人不信耶穌.
因為我們的表現.
沒辦法.
讓別人看出我是基督徒.
我的生命沒有真正的改變.
我沒辦法.
讓別人看到這件事.
所以當我們寫了這本書的時候.
書名的題目是.
我都擁有目的.
不是吧.
明不明白我的意思.
就是針對著我們很多時候都不知道.
我已經一出埃及第三個月.
就已經擁有這個.
祭司的國度.
的身份.
由那個時候一直整個舊約裡面.
有祭司來展現.
祭司的展現就是要講神的話語.
要教導神的話語.
要為人代求.
做很多的事.
我們都沒辦法.
來表現.
我們那種的目的.
我們在教會裡沒辦法表現.
在教會以外.
更加看不出.
所以我們這個.
相職是沒辦法展現.
領事們們.
信徒都有相職.
特別是在大使命裡.
我們每個信徒都說.
領命就是大使命出去.
但是其實大使命和我們這個相職.
又有什麼關係呢.

$^{921}$如果每個信徒都說我有大使命的時候.
你究竟明白不明白.
大使命和我們的相職有什麼關係呢.
如果這樣我們就要問.
究竟大使命的真正意義是什麼.
你知道.
當一個人去讀神學的時候.
很多時候都要讀.
這個玄文的.
陳全圖有沒有讀玄文.
有吧.
梁全圖呢 都有.
讀了希臘文和希伯來文.
每個人都要讀.
原來讀神學.
要讀玄文.
為什麼玄文好呢.
因為玄文幫助你更加了解某些內容.
是更加醒覺了.
現在就看看.
大使命.
這段經文的玄文.
現在很容易.
因為現在買一本.
希臘文的中文英文.
希臘文的聖經.
那你就懂得.
去看這段經文.
有很多不同的內容.
不過.
希臘文的文法裡有一個很重要的現象.
就是.
如果那句話.
那個動詞.
是那句話的核心.
那我們就問大使命的真義是什麼.
就是問.
這段經文的動詞.
在哪裡 如果找到動詞.
就知道大使命的核心在哪裡.

$^{961}$那我們就一起來.
看這段經文.
我們看到.
要去.
然後又看到.
要施洗.
然後又看到.
教訓.
然後又看到.
做門徒.
使萬民作主的門徒.
做門徒.
這裡有四個用詞.
請問這四個用詞.
好像.
做得不太好.
going, practicing, teaching.
和make disciples.
請問哪個是動詞.
going的舉手.
ok.
practicing的舉手.
還有 teaching的舉手.
為什麼你們每個都舉.
make的舉手.
好像都是這幾個舉.
你們四個都舉.
我剛才不是說了.
只有一個是動詞.
為什麼你們四個都舉.
不可能的.
是一個動詞.
所以我們看到大使命的真正意義.
這個.
ing going practicing teaching.
有ing.
英文叫做分詞.
participle.
有沒有英文老師在這裡.
如果有.

$^{1001}$你就知道是這樣教的.
但是我們在看的是希臘文.
希臘文的.
participle.
就是助動詞.
聽得明白了嗎.
going, practicing, teaching.
是什麼.
助動詞.
動詞是使萬民.
作我的門徒.
所以今天大使命原來是.
使萬民作我的門徒.
不是去傳福音.
不是去只是施洗.
也不是去教導.
而是令他做門徒.
然後有三樣東西.
going, practicing, teaching.
來去助他成事.
這個就是大使命的.
真正核心.
所以我們看到原來這三個的動詞.
分詞就是助動詞.
來幫助這個.
做門徒的核心的意思.
所以大使命是.
叫我們來使萬民.
作我的門徒.
然後我們就開始想.
為什麼是going呢.
因為保羅說.
你不傳福音就有禍了.
嚇得我們快快快快就短孫了.
又說短孫常孫.
不孫就不行了.
所以我們就經常說大使命.
就是傳福音.
為什麼就被保羅嚇成這樣.
其實耶穌怎麼說.

$^{1041}$耶穌來跟他們說.
去傳福音.
不是這樣的.
他還要告訴我們.
我們要去施洗.
Baptism.
為什麼Baptism呢.
這個.
施洗.
有一個很特殊的.
中國意味在這裡.
是奉父子.
聖靈的名.
這個.
還要大聲一點.
奉父子聖靈的名.
為什麼要大聲一點.
因為我奉大上老君.
奉子奉聖子.
還不要大聲一點.
好像大力一點就有效一點.
所以我們.
按他下水.
不是這樣的.
你看回聖經有些細節.
細節怎麼說.
合作給他們施洗.
歸入父子聖靈的命.
這個才是真正的意義.
是進入了.
父子聖靈的命.
下面.
所以你會發覺耶穌怎麼說.
耶穌就說那些法利賽人民士.
是怎樣.
勾引一個人入教.
跑遍洋海陸地.
做短孫中孫.
然後勾引一個人入教.
既入了教.

$^{1081}$就作地獄之子.
給你們眾臣.
給你們法利賽人民士眾臣.
這個.
給我一個提醒是什麼.
當我們來做傳福音.
我們做宣教的過程裡.
如果我們只是想不排成.
來帶人入到.
不是入到.
父子聖靈.
就是入到某個角度.
我記得有一次我在.
華福會參加聚會.
大型的聚會.
那大型的聚會有好幾千人.
然後我跟.
旁邊的人聊聊天.
你好嗎.
你從哪來的 貴姓.
他說我是某某某牧師的會友.
某某某牧師的會友.
失敬失敬.
然後我再問他.
你貴姓.
他說我是某某某牧師的會友.
說完這兩次之後.
我就關了喇叭走了.
為什麼 連爸爸姓什麼都不知道.
為什麼 我是某某某牧師的會友.
我是來到.
史雲頓華人基督教會的會友.
我是跟著.
陳嘉麗傳道的會友.
來 勁勁.
不是這樣的.
當.
法利賽人民士.
勾引人入教.
教會有很多人.

$^{1121}$不過這些人個個都是虛偽的.
假目為善.
比你們更加配.
原來長老可以這樣做 牧師可以這樣做.
因為長老回到公司的時候.
看不出他是基督徒.
在教會裡長老執事為至.
但是出去.
看不出.
這就是讓我們看到一個很大的危機.
還有教導是什麼.
教導就是說.
不論你辛苦還是教訓他去遵守.
弟兄姊妹們.
耶穌教導了你們什麼.
你是不是都可以.
你想想.
當我們看到.
雙職的身份.
我們要使萬民.
作我們的門徒才是真正的.
大使命的意思.
我們就使萬民作我們的門徒.
接著我下鏡一照.
這是什麼水.
你有幾斤幾兩.
耶穌教導你的.
你懂多少.
你如果這樣去教.
使萬民作我們的門徒.
我們中國人就說.
師父教徒弟留下一招.
我們老師.
他又是學生.
我又收一招.
他又收一招.
收了二千多年剩下多少.
屬零.
零雞蛋的零.
你想想.

$^{1161}$師父不能高過老師.
所以我們就發覺.
我自己來到去.
什麼都不知道.
照鏡子看看.
自己不行.
沒有去裝備.
你怎麼使萬民作我的門徒.
那些是劣質的門徒.
不是優質的門徒.
只有我好好裝備自己.
提升自己.
知道凡我所吩咐你的.
都教訓遵守.
教訓我什麼.
我才能使門徒.
是優質的門徒.
今天我們太多人被目擁.
在教會裡面.
八成人都是觀眾.
沒有將.
我們的屬靈恩賜彰顯出來.
沒有將.
主耶穌的愛.
彰顯出來.
所以我們這個目的.
專職的身份.
都沒有辦法去顯現出來.
當我們說.
靈明更新.
共正主恩的時候.
是口號.
每年都有口號給你的.
向上紮根.
向下紮根 向上結果.
每年都說這些.
有什麼用.
口號而已.
我在教會裡面做觀眾.
我是知體.

$^{1201}$不過我不選擇.
看東西.
我的知體是眼睛.
我的耳朵是知體.
我不選擇聽.
不選擇看.
我的血管選擇.
不流通.
不流通的話.
會變成什麼.
心疼塞.
你還去哪裡.
當然是醫院還是殯儀館.
所以弟兄姊妹們.
你擁有相職.
你已荊棘活了.
這個目的的身份.
你要好好表現出來.
而且.
你有一個專職.
專有的職份是主的門徒.
有沒有人被人看得到.
還是.
看不到.
求主幫助我們.
今天要好好.
將我們.
每個信徒都有的相職.
來彰顯出來.
一起低頭禱告.
當你給我們你的話語.
提醒我們的時候.
我們真的心裡要.
跟從你.
來走.
原來你離開埃及.
已經給每一個信徒.
擁有這個目的.
你也讓我們激活了之後.
要真是在愛裡面.

$^{1241}$來彰顯.
讓我們在彼此相愛當中.
讓人看出.
我們是主你的門徒.
求天父你幫助我們.
提醒我們.
真是要好好發揮自己的恩賜.
好好的來將.
神你的愛.
彰顯在人群當中.
無論在哪個境況裡面.
都能夠彰顯.
奉主耶穌基督聖名而求.
阿們.
\newpage



\section{約書亞記 24:20-24}
\label{sec:031973}
\textbf{在應許地活出「雙職」的使命}~~~~丘放河牧師博士
\newline
\newline
連結: \href{http://www.church.com.hk/acms/content.asp?site=cdc\&op=show\&type=product\&code=031973\&layout=sermon}{\texttt{031973}} ~~~~ 語音日期: 2026-04-12
\newline
\newline
\hyperref[sec:031972]{\small{< < < PREV SERMON < < <}}
~
\hyperref[sec:index_chronic]{\small{[返順時目]}}
~
\hyperref[sec:031974]{\small{> > > NEXT SERMON > > >}}
\newline
\newline
約書亞記 24:20-24
\newline
\begin{longtable}{cl}
\hline
\hline
章節 & 經文 (和合本修訂版)\\
\hline
24:20 & \begin{tabularx}{0.7\textwidth}{X} 你們若離棄耶和華去事奉外邦的神明，耶和華在降福之後，必轉而降禍給你們，把你們滅絕。」 \end{tabularx} \\ \\ \relax
24:21 & \begin{tabularx}{0.7\textwidth}{X} 百姓對約書亞說：「不，我們要事奉耶和華。」 \end{tabularx} \\ \\ \relax
24:22 & \begin{tabularx}{0.7\textwidth}{X} 約書亞對百姓說：「你們選擇耶和華，要事奉他，你們自己作證吧！」他們說：「我們願意作證。」 \end{tabularx} \\ \\ \relax
24:23 & \begin{tabularx}{0.7\textwidth}{X} 「現在，你們要除掉你們中間外邦的神明，專心歸向耶和華---以色列的神。」 \end{tabularx} \\ \\ \relax
24:24 & \begin{tabularx}{0.7\textwidth}{X} 百姓對約書亞說：「我們必事奉耶和華---我們的神，聽從他的話。」 \end{tabularx} \\ \\
[1ex]
\hline
\hline
\end{longtable}
$^{1}$第一次就是我的兒子一歲零十個月.
趁著他還沒有兩歲可以有10\%坐飛機.
由我太太當時要上班.
我就和爸爸媽媽和兒子來了英國.
這是第一次我來到英國.
來到英國有機會來參加Swindon的一個.
好像是春令會 不是應該是秋令會.
一起坐巴士一起去玩.
當時陳全道說我曾經來過這裡的一個教會.
其實我忘記了.
但是他就將梁全道的那個.
我們當時有幅相傳給我看.
有相為證 有我 有我爸爸.
有我媽媽還有我兒子.
還不是 沒得說了.
原來真的來過這裡的.
這次第二次來這裡.
也是很特別的.
這次有機會來到.
第一次坐英國的火車和坐英國的地鐵.
由Dunkester來到這裡和大家有機會見面.
這一次我一直將過去差不多一年半前.
一個信息 神給我的一個信息.
就在《出海及記》第19章.
第一節到第六節的那個啟示.
來給我看到.
今天早上我就說了.
原來每個信徒都有相職的身份.
我今天來到第二堂.
我就想和大家來講另外一個議題.
但是離不開相職.
也都是這麼說.
從2024年11月.
我領悟了這個相職的身份之後.
這個信息之後.
我差不多每一堂都有這兩個字相職.
這兩個字是很重要的一個信息.
對我來說.
也都對每一個信徒是很重要的信息.
今天我們看到《約書亞記》第24章的時候.

$^{41}$當我們看到這段經文的時候.
我們看到約書亞和那些以色列人.
好像有一種誓神拚命的經歷在那裡.
為什麼這麼說呢?.
當我們一路看的時候就說.
你信不信神?跟不跟從神?.
一定跟 一定跟.
做不做見證?.
見證 真的見證.
是真的 不是真的降禍你的.
是一定會跟從神.
這段經文似乎是.
約書亞記第24章就是成卷約書亞記的一個終結篇.
在這個終結篇裡面.
我們看到.
原來整個約書亞記是在說.
他們去到這個應許地.
所謂的迦南美地.
去到這個地方.
來到去有一個很特別的一個經歷在這裡.
而一開始整個約書亞記第一章.
就說到他們已經進入了這個應許地了.
然後進入了應許地之後.
原來第一章就說到一個很重要的.
剛強壯膽.
知不知道什麼意義?.
原來當應許這個地方給你的時候.
是要你打回來的.
不是送給你的.
這一刻.
你要將整個迦南地的所有民族征服了.
才是屬於你的地方.
怕不怕?.
真的怕.
因為打仗.
誰不怕?.
現在世界各地.
中東打到這樣的時候.
有時候我們會說.
為一些無辜的百姓去祈禱.

$^{81}$你不為那些兵士祈禱.
他也很無辜的.
他只是打工而已.
誰知道走出去要死的.
然後他的媽媽沒有了兒子.
他的老婆婆沒有了老公.
他的兒女沒有了爸爸.
在這一刻裡面.
我們發覺原來打架的過程裡面.
真的約書給神來說.
你要剛強壯膽.
他們接著就找了一些探子進去看這個地方.
探子差點被抓了.
幸好有拉合妓女.
讓他們有機會逃脫了.
然後他們進入英女地.
過這個約旦河.
這個過約旦河很特別.
因為過約旦河就像過紅海一樣.
過紅海.
摩西的腳一踏進水裡.
水就分開.
他們走乾淨地過紅海.
但是原來這件事是重複的在這裡出現.
就是約書帶他們過約旦河.
都是這樣的.
這個給我們看到.
原來有一個很特殊的現象.
就是摩西帶領以色列人出埃及.
離開了世界.
離開了法老這個在預表魔鬼.
預表世界的地方.
進入迦南地.
預表天堂的地方.
但是摩西.
他是代表.
律法不能夠帶人進入天堂.
迦南.
誰帶他們去.
約書.

$^{121}$約書的字根就是救恩的意思.
所以救恩帶人進入迦南地.
所以他們就過了迦南地.
遇到第一個大城市叫做耶利哥城.
在耶利哥城這麼堅固.
怎麼打.
第一場就打到硬仗.
神就叫他們圍著這個城七次吹號吹角.
這個城自動倒下了.
如果你去以色列地.
你仍然可以看到耶利哥城.
那個城牆.
毀壞在那裡.
還沒翻轉.
還沒重建.
在這個過了之後.
他們跟著一個一個一個的民族.
來征服.
征服了之後就分地了.
十二個支派分了地之後.
有些就抽籤.
來分了地.
分了之後.
在這個事情裡面.
約書怎麼跟他們說.
約書說.
我和我家.
必定事奉神.
你想想.
當在第24章一開始的時候.
就告訴我們知道.
約書帶領那幫以色列人.
去到一個事件的地方.
來到那裡去祭神祭怨做一些事情.
你想想事件是什麼地方.
事件就是昔日神童亞伯拉罕他們的先祖.
來到去說你的子孫後代.
好像天上的星.
地上的沙.
這麼多的一個地方.

$^{161}$然後在那裡就允許他們.
來到去可以得到一個很特別的神的眷顧.
而約書帶領他們去到這個地方.
終於怎麼樣.
征服了所有民族.
然後在這一刻裡面.
就在事件這個地方.
和他們再重新.
好像立約一樣.
你們現在要怎麼做.
他們一次又一次的說.
行 做見證.
跟從神.
但是約書怎麼跟他們說.
你們要除掉外邦神.
在你們中間的外邦神要除掉他.
當這一刻這樣說的時候.
給我們一個提醒就是說.
原來當我們去到應許地.
已經征服了所有民族的過程裡面.
但是約書要跟他們說一樣東西.
你們將你們的外邦神來挪開他.
然後要專心來歸向神.
你想一下我們看到一樣東西就是說.
說就容易.
大道理誰不會說.
所以道德經說.
大道自揀.
知易行難.
今天我們基督徒都是這樣.
說些道理容易什麼.
我們能不能做到.
是很大的一個問題.
經常我都說我們每一年的那個叫做什麼.
今年的主題是什麼.
口號而已.
說一下而已.
因為為什麼.
立志為善由得我.
只是行出來由不得我.

$^{201}$弟兄姊妹們.
說就容易.
但是走出來並不是容易.
當我們在婚姻裡面說I do的時候.
是不是.
我願意.
結了婚之後.
You do.
你做了是不是.
垃圾不倒.
款又不用.
買菜又不買菜.
煮飯又不煮飯.
你做了這樣.
哇結過婚的人都有這樣的經歷.
對不對.
真是生不如死.
你想想結婚的時候多浪漫.
但是結完婚之後.
真是很艱辛的來維持這個婚姻.
因為立志容易.
所以你會發覺原來立志只是一個什麼.
開始而已.
當那些約書和這幫人來做不做.
做.
跟不跟從.
跟從.
做不做見證.
做見證.
立了石頭做見證.
立志容易.
但是走出來才是最最大的挑戰.
所以在這一刻裡面我們看到.
完全跟在抗議的過程裡面完全不同的.
你想想昔日在抗議的時候.
他們走的時候.
哇多好啊.
早上一出去就撿東西吃.
你現在早上趕著上班的.
但是他有東西吃.

$^{241}$有水喝.
又有鵪鶉吃.
接著打仗.
簡單的舉起手就贏.
坐下來就輸了.
舉著他都撐到贏了.
你會看到抗議的經歷.
跟真真正正要面對著在英他們進入之後.
是完全不同的.
在這裡他們需要面對著.
他的生命.
他的職份.
他擁有什麼身份.
所以當然雙職是需要在這裡展現出來.
所謂雙職就是說有目的和專職.
這個目的讓我們看到我們有一個祭司的身份.
當我們看祭司的身份就是在他們出埃及的第三個月.
在利菲定去到抗議的地方.
神就怎樣跟他說.
神說如果你遵守我約手我誡命的時候.
我就給你這個祭司的國度.
聖傑的國民.
你是屬於我的子民百姓.
祭司的國度表示什麼.
每一個跟從神的人都有祭司的身份.
祭司的身份就是什麼.
目的.
所以每一個願意跟從神的人.
都擁有這個目的的身份.
如你擁有這個目的的身份.
你就要彰顯這個目的的身份出來.
還有第二個身份叫專職.
這個專職不是我們打工的專有職業.
而是專有的職份.
就是當主耶穌在聖餐的時候.
來跟那幫門徒說.
我要賜給你一條新的命令.
叫你們彼此相愛.
如果你彼此相愛.
眾人就知道.

$^{281}$你是我的門徒.
你會發覺.
今天我們基督徒的夫妻.
能不能夠相愛到讓人看到.
我們是主的門徒呢.
有時候那些基督徒.
都一樣到處吵架.
我見到一個師母在街邊.
戳著牧師的頭.
像死人一樣.
這樣戳.
嘩 你想想.
原來夫妻是基督徒.
就算做到全道做到牧師.
都不代表不吵架.
嘩 我們在這裡開始想.
究竟為什麼.
特別是在.
約翰福音第17章.
那裡寫什麼.
分離的禱告.
分離的禱告就是表示.
耶穌基督快要離開世界了.
他在之前有一個禱告.
在這個第17章裡.
五次的.
耶穌說了.
合而為一.
五次.
整本新約聖經裡出現了七次.
有兩次是保羅所說.
但是五次在這個.
分離的禱告裡說了.
合而為一.
分離就是耶穌快要死.
中國人有一句話.
人之將死.
其言也善 其言也真.
就是你臨死的時候說的話.
都特別真.

$^{321}$你到時不會騙人的了.
為什麼.
肺腑之言.
究竟我的子女.
我的後代.
我們的人要怎麼做.
快點說遺言.
耶穌的遺言是合而為一.
意思是什麼.
意思是耶穌已經預示到.
他將來的門徒.
是搞搞震.
亂七八糟.
不合而為一.
今時今日我們看得到.
不同的宗派.
不同的教會的鬥爭.
不同的表現.
是沒有合而為一.
是爭鬧不斷.
而在這個過程裡.
合而為一.
要怎麼樣.
眾人就認出.
你是我的門徒.
你是有一個專有的職份.
是主的門徒.
所以當我們看到.
特別我寫了一本書.
我都擁有木職.
不是吧.
在那裡有機會可以下一本.
你會發覺這本書講了一個很重要的.
我們擁有這個雙職.
是專職和木職.
木職我是一個.
屬於神的一個祭司.
我要怎麼樣在教會裡.
展現出來.
我要怎麼樣將我的恩賜.

$^{361}$彰顯出來.
一會兒我們會講為什麼.
下午我們會再講.
這個屬靈恩賜的測試.
怎樣去應用這個木職的身份.
但是在教會以外.
無論你是工作場所.
在家庭裡.
在社區裡.
或者在任何地方.
你能不能給人看出你是主的門徒.
比如你去超市買東西.
給了他20磅.
他當100磅就給你.
你拿著80多磅.
你怎麼做.
又不是我的錯.
那就走了.
或者是你說.
你搞錯了我給你20磅.
你現在當作100磅.
這樣找錢給我.
那個收銀員怎麼樣.
看著你.
真的謝謝你.
不然他要虧80磅.
你想想.
你能不能給人看出你是主的門徒.
這是很重要的.
一種專職.
專有的職份.
所以當約書客對著這些.
以色列人在那裡誓神疲軟.
在事件這個地方.
是很有歷史價值的地方.
來跟他們說你們要將你們的外邦神.
除掉的過程裡.
為什麼.
在江山之後.
在分好地之後.

$^{401}$在這個應許地才是正式的.
信仰的戰場的開始.
原來.
我們很多時候都不想到.
我們是面對著一個.
征戰.
而且那裡聖經告訴我們.
我們是用血氣來征戰.
魔鬼撒旦並沒有離開.
他繼續的搞擾.
所以當我們看到應許地.
其實就是一個信仰.
要實踐我們的使命.
我們擁有相織的使命的地方.
在這裡.
當這樣離開的過程裡.
我們就更加看到.
今天的信徒要面對著什麼的.
情況.
當日要將.
那個外邦神.
除掉.
外邦神除掉是一個很重要的.
現象在這裡.
因為在抗議裡面.
剛才都說過.
有黑暗又有暗春色.
所面對的所有方方面面.
都是耶和華神.
和摩西.
是這兩個.
摩西又代言光神說話.
又是神.
沒有別的東西.
沒有挑戰.
只有神.
沒有其他的偶像.
所有的生活模式.
都是很純良的.
但是當我們看到.

$^{441}$各個地.
準備在應許地.
生活下去的時候.
就完全不同的局面.
為什麼呢?.
有地方就種下地.
我們中國人來到英國.
有後院做什麼?.
種下花之外.
當然是種下菜.
一有菜.
你知道種了菜你能吃完嗎?.
吃不完就拿回教會送給弟兄姊妹.
但是有些人.
就去超市買.
你想想在那時候.
農耕的時代.
有羊有牛.
有種菜.
那又怎樣? 做生意.
當在這個.
進到應許地.
拿好地之後.
就開始做生意的過程裡面.
你就會變成一個很特別的現象.
就好像我以前.
去美國讀書.
我讀工商管理.
(笑聲).
就是工商管理.
奸商管理都一樣.
因為商就是奸.
那你想想.
在應許地裡面.
安逸? 不是安逸.
去發展.
有商業行為.
奸的東西就出現了.
基督徒都一樣奸.
你想想.

$^{481}$有個姐妹跟我說.
她說我很不喜歡.
那個教會的長老.
為什麼? 我就跟她打工.
她刻薄你? 不是.
不過在教會她是一套.
在公司又是另外一套.
那裡是天使.
那邊是魔鬼.
看不出她是長老.
或者是執事.
或者是基督徒.
當我們來到看到.
原來在應許地裡面.
正正是我們自己的生命.
能不能彰顯到.
神的榮美.
我擁有的兩個職位.
是專職和目職.
當我們這個說法.
正正是我們今天新約的信徒.
都有面對的.
我們信主的時候.
特別是.
洗禮的時候都一樣.
面對著 牧師說.
你信了主之後 你願不願意一生事奉神?.
事奉.
誰不事奉都不會.
上到祠了.
你說一生事奉神.
不就和結婚的時候.
矮肚子一樣.
結婚之後.
做了基督徒 洗禮之後.
都不刀你.
這個是我們今天很大的問題.
為什麼?.
因為我們很多的外邦神.
還沒來取調.

$^{521}$這個約書叫他們將外邦神取調.
就是說.
不是真的我們家裡擺了個佛像.
偶像擺在那裡.
我們心裡的人生觀.
價值觀.
和道德觀.
有沒有改變過.
原來在聖經裡.
有很多的原理.
很多的原則.
很多的真理.
叫我們去學習.
但是我們有沒有來學習.
這個才是重要.
當我們在生命裡.
有很多的偏離的思想.
帶著我們走.
從小到大.
媽媽爸爸教導我們.
你越賺得多錢.
就越成功.
所以我們在物質主義裡.
要好好地去做.
有個弟兄.
準備要去讀神學.
要讀神學的時候.
他兒子跟他說.
以後我們每年去旅行.
沒有了.
原來讀了神學.
要做牧師.
很慘的.
收入很低的.
爸爸本來做領導階層.
每年去兩次旅遊.
都有的.
一做了牧師.
一讀了神學.
就沒有了.

$^{561}$這種物質主義.
是充斥著小孩子都知道的.
甚至我們舊約裡.
都有成功神學.
成功神學是什麼意思.
我們來到.
越高的位置越被尊崇.
所以我有一個註冊部的同工.
就告訴我.
有個弟兄來到想讀神學.
讀什麼?讀道學碩士.
為什麼你要讀道學碩士?.
做牧師.
為什麼要做牧師?.
做牧師很威風的.
好像一言堂牧師說了.
就聽了.
一個來自中國的弟兄.
為什麼?.
中國大陸牧師是很少的.
牧師在課堂裡.
他講什麼.
就像聖主一樣.
所以我經常在想.
當提摩太前書第三章第一節.
保羅這樣說.
他說.
如果你要做監督的時候.
你要善務善工.
善務善工.
就是做好事.
我們很多人善務.
只是聖職.
做好事.
做牧師,傳道人,告慰.
以為這樣才是成功.
你想想.
當我們來看.
有一個會議.
叫尼西亞會議.

$^{601}$尼西亞信經.
你聽過.
尼西亞信經就是尼西亞大會.
在主後325年.
開的一個大會.
尼西亞大會.
尼西亞信經出來了.
還有66卷書成為正典.
都是那一次會議.
來決定的.
去參加的人是什麼人?.
全部都是主教.
但我們翻查聖經裡.
沒有主教這個位置.
有長老.
有執事.
有監督.
但沒有主教.
一開始.
不僅不合而為一.
還是善務.
成績.
每個人都想在高位.
更加去到一個地步.
有教皇,洪雨主教.
舒基主教.
各樣不同的職位.
成功神學.
充斥著我們.
基督徒的場合.
還有誰是你的老闆?.
當我們來看.
我們的生命裡.
去到這個應許地裡.
當我們去到工作裡.
誰是我們的老闆?.
很多時候當我們.
回到家裡的時候.
老婆走過來跟我們聊天.
我還沒搞定.

$^{641}$兒女跟他玩.
都沒空.
你的工作.
你的老闆.
真的是你的老闆.
你的工作就是你的主人.
在過程裡充斥著.
這些外邦神.
在過程裡.
價值觀.
衡量是怎樣的衡量呢?.
還有我們有.
物質主義.
也有崇拜.
品牌.
我們帶著.
名牌包包.
身份高了很多.
這是我們很多時候.
都有的現象.
這個品牌的崇拜.
成為我們的外邦神.
還有一樣更厲害的.
就是自我為中心.
自我為中心.
表示什麼?.
整個世界都是在我周圍轉的.
我是中心.
你們每個聽我說.
來了解我.
我就是世界最核心的一個.
你想一下.
今天你去想一下.
我怎樣我想什麼.
我們每一樣事情.
都以我的想法為最終.
這個外邦神.
充斥著我們的世界.
所以當我們看到.
這個約書夏說.

$^{681}$要將外邦神除去的時候.
我們千萬不要有一個迷思.
不要緊.
我會好好管理外邦神.
我懂得怎樣管理.
我保留著他不要緊.
但是約書夏說.
將這個外邦神.
除掉.
不可以留在當中.
當這個說法.
讓我們知道.
原來我們在出埃及記第19章.
和約書夏記第24章.
裡面.
有兩節經文作為對比.
這兩節經文.
當我們說.
我們如果來遵守神的約.
守祂的誡命.
我們就擁有.
祭司的國度.
和祭司的身份的時候.
那裡有一個很重要的說法.
若實在聽從我的話.
弟兄姊妹們.
不要將我們每年的.
主題作為一個口號式.
來看.
我們實在的.
來聽從神的話.
不是只有「說」字.
到了約書夏記第24章.
特別是第14節.
也說.
你們想在大河那邊.
去信那些敬拜的神.
還是要誠心實意的.
事奉耶和華神呢.
今天我們來到神的教會.

$^{721}$我們在這裡的時候.
究竟我們信神的真實度.
去到什麼地步呢.
你就會發覺.
這兩節經文的對比.
是很清晰的.
無論摩西.
來跟神說.
對他們來告訴他們.
你們要怎樣真真實實.
事奉神.
或者是約書夏.
在這裡取得所有的地方.
進入了應許地之後.
對他們說.
你們要怎樣去面對.
是實在的.
誠心實意的.
事奉神.
所以當我們這樣看的時候.
這兩節的經文.
對比摩西.
叫我們有目的.
然後約書夏記.
進入應許地.
實踐.
我們擁有的雙職身份的時候.
這就在令到.
我們今天要思想.
一個很重要的議題.
就是叫做.
受洗就等同.
進入應許地.
你會不會發覺.
洗禮之後.
剛才說的.
在這裡我願意一生跟從神.
其實不就正正像約書夏.
和這些以色列人.
就在那裡事件的地方.

$^{761}$你們願不願意.
去誠心實意的跟從.
不是在這裡.
要做什麼.
我們廣東人.
有一句話.
什麼當吃生菜.
沒錯了.
看到他們嗎.
一定願意.
做不做見證.
做見證.
我們看約書夏記.
第24章.
這幾節的經文.
說了23節.
約書夏說.
將你的外邦神除掉.
一心歸向神.
這才是我們真實要面對的.
所以保羅.
有一個.
保羅三論.
是值得我們學習的.
保羅三論是說什麼呢.
其實很簡單.
也你們聽過.
就是在.
哥林多前書第12章那裡.
教會論.
肢體論.
和屬靈恩賜論.
兩節的經文.
是在12章27.
28節.
所謂教會論.
這裡說你們就是基督的生子.
大家認不認同.
你們就是基督的生子.
教會就是基督的生子.

$^{801}$所以.
頭就是基督.
教會就是他的身.
所謂教會是他的身的時候.
我們可以想起.
合二來說.
就是我們這裡的教會.
就是.
史雲頓華人基督教會.
講二來說.
就是古往今來的所有的基督徒.
是教會.
但是我們不要說這麼遠.
就說回這裡.
教會就是基督的生子.
然後第二點.
就是肢體論.
你們各自作肢體.
每個人都作肢體.
那你就會問.
肢體有什麼.
肢體有功能.
每個人都有功能.
那你會發覺.
你如果有這個功能的時候.
你會不會發揮呢.
願不願意發揮出來呢.
先問一下你是什麼功能.
心肝脾肺腎是哪一種.
還是盲腸.
盲腸也很特別.
盲腸.
盲腸有什麼作用.
要來割的.
盲腸.
是不是要來割的.
人們的研究就說.
盲腸好像沒什麼作用.
有空.
塞了就割了算了.

$^{841}$原來當疫情發生之後.
人們就開始發覺.
盲腸可以來生發益生菌.
可以防疫的.
保住盲腸.
不要隨便割.
原來我們每一個身體裡面的器官.
肢體都有用的.
就是說如果它.
是有用而你不去用它.
有沒有問題.
你是眼睛不看.
不看.
問一下自己.
你存在的價值是什麼.
沒有.
對不對.
如果你是眼睛不看.
你是沒有價值的.
所以我們發覺原來.
在這個過程裡.
我們看到保羅.
將這三個論.
教會論,肢體論.
還告訴我們屬靈恩賜論.
他在第十一節那裡.
告訴我們神聖靈隨著他的記憶.
將屬靈恩賜分給我們每一個信徒.
每一個信徒.
都有這個屬靈恩賜.
屬於肢體裡面的發揮功能.
所以我們的屬靈恩賜.
最少有一個.
可以多過一個.
二個,三個,四個.
有些人說.
很有恩賜的,很有才華的.
他可能聖靈給多他幾個.
但是你要記住一件事.
當給了的時候.

$^{881}$耶穌說一個比喻是很特別.
他就說有一個人.
拿了一千兩銀子之後.
他將它埋在地裡.
埋在地裡.
很特別.
埋在地裡.
如果用肢體論.
和屬靈恩賜論.
肢體不發功能.
屬靈恩賜不發揮.
埋在地裡.
很像什麼呢?.
牧師傳徒人跟弟兄姊妹說.
你出來幫忙.
事奉一下,發揮一下你的恩賜.
他說不行的.
多謙卑.
謙卑不代表.
屬靈恩賜不發揮.
因為那個.
拿著一千兩銀子.
埋在地裡.
他的結局是什麼呢?.
四個字.
是嗎?.
你現在說出來.
大聲一點.
被人收起來.
四個字.
不是,左邊四個字.
一無所有.
又不是.
聖經說.
哀哭切齒.
在黑暗裡.
哀哭切齒.
什麼叫哀哭切齒呢?.
就是.
落地獄.

$^{921}$嚴重到這個地步.
我再說一個例子.
路加福音第十六章.
有個財主和拉薩路.
拉薩路財主一起死了.
財主在下面.
哇.
拉薩路在阿伯拉赫的懷裡.
很舒服.
我這條肝一滴水都沒喝.
很辛苦.
然後呢.
他的心裡有什麼想法.
財主.
他沒想過.
好像包青天那樣.
拉一些.
判案的時候.
拉了一個人.
要去狗頭閘,龍頭閘.
虎頭閘.
通常拉出去的時候.
那些人叫什麼.
冤枉.
冤枉啊.
我沒做啊.
是冤枉的.
財主有沒有說冤枉啊.
沒有.
他關心的是什麼.
我五個弟兄還沒信主.
快點拉薩路.
阿伯拉赫叫拉薩路復活.
叫我五個弟兄.
不要走到我這個地步.
我這個地步.
是兩個字而已,又是什麼字.
辛苦?不是.
還要更濃稠.
對了,活該.

$^{961}$活該.
不是冤枉的.
弟兄姊妹們.
我告訴你,今天你來這個教會.
你不發揮你的屬靈恩賜.
活該,將來去敲凹設址.
嚴不嚴重.
嚴重到什麼地步.
那個財主.
是沒有怨言的.
因為他知道過去這麼長的歷史.
裡面神的恩典.
在我們生命裡一點又一點.
來到去彰顯.
主恩.
同行.
但是我們.
漠視.
我的屬靈恩賜.
漠視我應有的表現.
我是有相職.
在這個應許地裡.
我有相職的使命.
我沒有去彰顯出來.
當沒有去彰顯的時候.
你想想怎麼樣.
在教會裡.
我們應該表現.
屬靈恩賜.
剛才這個保羅三論.
是在教會的範圍.
為什麼?身體就是教會.
所以你在教會裡.
你要將你的屬靈恩賜發揮出來.
但是在教會裡外.
你有很多不同的場合.
家庭,社會.
或者是社區.
或者是學校.
有不同的境界裡.

$^{1001}$你又沒有三個字基督徒在額頭.
人家不知道你.
會不會將主的愛.
表現出來.
如果你將主的愛.
表現出來的時候.
眾人就認出你是我的門徒.
所以在教會外.
教會裡有一個很不同的身份.
一個叫目職.
一個叫專職.
在這個專職的過程裡.
正正是我們很多時候.
看到宣教的議題.
為什麼今天.
很多人不信耶穌呢?.
就是因為他看到有很多的基督徒.
這樣的樣子.
他就不信.
看到嗎?.
我也是基督徒.
我也是屬於那一部分.
他看到我,所以他不信.
這是很嚴重的指控.
但是又很嚴厲的.
一個真實的現象.
今天我們擁有這個商職.
但是我們沒有很好的.
來彰顯.
所以當我們看到我們有這兩樣東西的時候.
約書客是很清楚的.
除掉外邦神.
將所有不屬於聖經的.
所有倫理.
道德.
價值觀.
拿開.
你要專心的來歸向耶和華神.
要活出我們擁有的商職的身份.
這個正正是今天.

$^{1041}$我來跟大家說.
在應許地.
活出商職的使命.
如果你都認同的時候.
你真的要明白.
究竟神給了你什麼.
屬靈恩賜.
你不知道.
下午來聽.
我們先來看頭禱告.
多謝天父.
你給我們有屬靈的氣質.
你告訴我們我們是擁有目的.
因為我們是你的君尊的祭司.
當我們今天來到你的面前.
在這個教會裡面.
我們屬於哪一個肢體.
有什麼功能.
求你展現給我們知道.
讓我們知道.
怎樣來彰顯.
你又告訴我們.
我們要將你的愛.
要來彼此相愛.
特別是基督徒要彼此相愛.
才能夠讓人認出.
我們是你的門徒.
讓我們真的能夠有這種表現.
將我們的專職.
在這個應許地裡面.
活出來 求主幫助.
奉主耶穌基督的聖名而來.
\newpage



\section{詩篇 51:1-19}
\label{sec:031974}
\textbf{內裡的更新}~~~~張健庭牧師博士
\newline
\newline
連結: \href{http://www.church.com.hk/acms/content.asp?site=cdc\&op=show\&type=product\&code=031974\&layout=sermon}{\texttt{031974}} ~~~~ 語音日期: 2026-04-12
\newline
\newline
\hyperref[sec:031973]{\small{< < < PREV SERMON < < <}}
~
\hyperref[sec:index_chronic]{\small{[返順時目]}}
~
\hyperref[sec:031975]{\small{> > > NEXT SERMON > > >}}
\newline
\newline
詩篇 51:1-19
\newline
\begin{longtable}{cl}
\hline
\hline
章節 & 經文 (和合本修訂版)\\
\hline
51:1 & \begin{tabularx}{0.7\textwidth}{X} 神啊，求你按你的慈愛恩待我！按你豐盛的憐憫塗去我的過犯！ \end{tabularx} \\ \\ \relax
51:2 & \begin{tabularx}{0.7\textwidth}{X} 求你將我的罪孽洗滌淨盡，潔除我的罪！ \end{tabularx} \\ \\ \relax
51:3 & \begin{tabularx}{0.7\textwidth}{X} 因為我知道我的過犯；我的罪常在我面前。 \end{tabularx} \\ \\ \relax
51:4 & \begin{tabularx}{0.7\textwidth}{X} 我向你犯罪，惟獨得罪了你，在你眼前行了這惡，以致你責備的時候顯為公義，判斷的時候顯為清白。 \end{tabularx} \\ \\ \relax
51:5 & \begin{tabularx}{0.7\textwidth}{X} 看哪，我是在罪孽裡生的，在我母親懷胎的時候就有了罪。 \end{tabularx} \\ \\ \relax
51:6 & \begin{tabularx}{0.7\textwidth}{X} 你所喜愛的是內心的誠實；求你在我隱密處使我得智慧。 \end{tabularx} \\ \\ \relax
51:7 & \begin{tabularx}{0.7\textwidth}{X} 求你用牛膝草潔淨我，我就乾淨；求你洗滌我，我就比雪更白。 \end{tabularx} \\ \\ \relax
51:8 & \begin{tabularx}{0.7\textwidth}{X} 求你使我得聽歡喜快樂的聲音，使你所壓傷的骨頭可以踴躍。 \end{tabularx} \\ \\ \relax
51:9 & \begin{tabularx}{0.7\textwidth}{X} 求你轉臉不看我的罪，塗去我一切的罪孽。 \end{tabularx} \\ \\ \relax
51:10 & \begin{tabularx}{0.7\textwidth}{X} 神啊，求你為我造清潔的心，使我裡面重新有正直的靈。 \end{tabularx} \\ \\ \relax
51:11 & \begin{tabularx}{0.7\textwidth}{X} 不要丟棄我，使我離開你的面；不要從我收回你的聖靈。 \end{tabularx} \\ \\ \relax
51:12 & \begin{tabularx}{0.7\textwidth}{X} 求你使我重得救恩之樂，以樂意的靈來扶持我， \end{tabularx} \\ \\ \relax
51:13 & \begin{tabularx}{0.7\textwidth}{X} 我就把你的道指教有過犯的人，罪人必歸順你。 \end{tabularx} \\ \\ \relax
51:14 & \begin{tabularx}{0.7\textwidth}{X} 神啊，你是拯救我的神；求你救我脫離流人血的罪！我的舌頭就高唱你的公義。 \end{tabularx} \\ \\ \relax
51:15 & \begin{tabularx}{0.7\textwidth}{X} 主啊，求你使我嘴唇張開，我的口就傳揚讚美你的話！ \end{tabularx} \\ \\ \relax
51:16 & \begin{tabularx}{0.7\textwidth}{X} 你本不喜愛祭物，若喜愛，我就獻上；燔祭你也不喜悅。 \end{tabularx} \\ \\ \relax
51:17 & \begin{tabularx}{0.7\textwidth}{X} 神所要的祭就是憂傷的靈；神啊，憂傷痛悔的心，你必不輕看。 \end{tabularx} \\ \\ \relax
51:18 & \begin{tabularx}{0.7\textwidth}{X} 求你隨你的美意善待錫安，建造耶路撒冷的城牆。 \end{tabularx} \\ \\ \relax
51:19 & \begin{tabularx}{0.7\textwidth}{X} 那時，你必喜愛公義的祭和燔祭，全牲的燔祭；那時，人必將公牛獻在你壇上。 \end{tabularx} \\ \\
[1ex]
\hline
\hline
\end{longtable}
$^{1}$主持人:好,弟兄姊妹,早晨!我們再一次低頭祈禱..
天父願你的話語在我們中間來做一次跟我們說話..
願我們的心靈打開的時候,你聽到的不單止是人的聲音,.
你聽到,讓我們也聽到你自己的聲音..
你的呼喚不單止是讓我們知道一些內容,.
讓我們知道一些知識,也是呼喚我們來跟從..
願你讓我們看到我們自己真實內心的情況的時候,.
我們向你面前來求恩典,求憐憫,.
來禱告奉耶穌基督的名其求.阿們..
好了,弟兄姊妹,今天我們就一起來思想第四篇,第五十一篇..
之前在壽苦節那個時間,你也知道我們一起來思想,.
就是在教會歷史當中稱為七重要的七宗罪的內容,.
說到人心裡面的那些黑暗的情況,.
這個也是一個很真實的情況..
最近我就看了一本書的內容,.
不知道你有沒有發覺在現在這個世界,特別是在美國,.
最近這一兩年,其實也不止一兩年,也有一段時間,.
很多的政治的爭拗,政治的分歧,.
去到大家都是你死我活的一個地步,.
就是大家都是各不相讓,甚至乎看對方是一個仇人的情況..
當然也不是理想,但是我都問了一個問題,.
究竟其實背後人的政治取向究竟是怎麼來的呢?.
或者哪一類的人會選擇哪一方面的政治傾向呢?.
那本書其實是在說這件事,.
說到最後的時候,其實整本書有一個框架,.
我從來都沒有這樣想過,.
就是原來在這個世界,.
人怎麼看這個世界,其實就是你選擇什麼政治的取向,.
這個很合理,你看這個世界是這樣的,.
那你就選擇參加什麼政治的政黨,.
或者你的政治的取向是怎麼樣的..
但是那本書的框架有一個很特別的地方,.
它說其實不同的,簡單來說就是左派右派,.
自由派或者保守派,.
它說其實最重要的一個框架就是兩種人看世界的不同,.
這兩種人是什麼人呢?.
有一種人不是看世界,而是看人性的不同,.
有一種人他看人性是美好的,.
這個世界慢慢慢慢是藉著建立一些美善的制度,.
政府的幫助,外在的立法,一些的規條,.

$^{41}$可以在地上建立一個越來越好的世界..
當然去不去到烏托邦就很遠的路了,.
但是最低限度就是這樣看世界,.
人性是美好的,所以要藉著,.
有些人現在可能是落到一些很慘的位置,.
譬如可能他做錯了一些事要坐牢,.
或者他做了一些錯的決定落到一個,.
譬如可能要睡街或者無家可歸,.
但是基本上他都是美好的,.
所以藉著一些外在的幫助,.
就可以幫助他慢慢慢慢去建立起來,.
這個世界就會越來越美好了..
這是其中一種看世界的方式,.
或者看人性的方式,.
你就知道他選了哪一邊,.
就是進步主義的方向..
但是另外一種的世界觀是不同的,.
另外一種的人就說看世界看人性是邪惡的,.
人是黑暗的,.
在這個世界都不用指望是會建立一個.
人類美好的烏托邦,.
你打算要建立一些用一些政府的權力,.
將我給一些權力給一些機構,外在的建立一些制度,.
給那些人打算他會幫你建立一個美好的世界嗎?.
你不用指望,他拿了那些權力他就腐敗了,.
他就會有貪污各樣東西,.
所以政府的位置你就不要給他那麼多權力,.
你甚至乎要限制他,.
因為你也是污衊,他也是污衊,.
他不會做了政府官員的時候突然間變成聖人,.
所以其實就是這兩種不同的世界觀,.
或者不同的人性觀來給我們看到,.
其實就是一個最基本的分別..
而我問大家,.
為什麼在美國當中還有那麼多的人.
仍然都是覺得政府需要限制呢?.
或者仍然都是那麼支持保守主義的呢?.
這個我之前很少想,.
原來因為在美國當中還有很多基督徒,.
不知道你有沒有這樣想過,.

$^{81}$譬如你看一些現在在歐洲那邊的,.
無論是政治,無論是什麼,.
你就發覺,因為當信仰的影響力越來越少的時候,.
人對於人性有很多很不切實際的美好的一個憧憬,.
他覺得人都是好的,.
即使他吸毒,也不是他選擇的,.
我們給一些輕量的毒品他,.
慢慢慢慢他都會做一個,.
你不要限制他不要吸,.
因為他吸了也會有一些快感的時候,.
可能解脫了他心靈一些的困境,.
他慢慢慢慢其實他都是一個好的,.
諸如此類,.
因為人對於信仰的,那種以前在西方社會,.
對於基督教信仰那種很明白人性的那種黑暗,.
現在這個世界很多人就不是這樣看,.
特別如果你和那些年輕人去聊天的時候,.
慢慢慢慢講到尾,.
因為你知道年輕人很多時候他們去讀大學,.
他們受了教育,.
他們不斷有一個觀念就是人性都是好的,.
我想這個就是一個最基礎的底線,.
但在我們的信仰當中,.
其實我們不是這樣看這個世界的,.
不知道你有沒有想過,.
今天早上和剛剛主一學的時候,.
我就開始了一個新的課程,.
就是講舊約概論,.
在整個舊約我就很快地用了幾十分鐘的時間,.
來講整個舊約的一個發展,.
整個舊約的發展如果今天早上有來過聽主一學的同學,.
我想你都會掌握到一點,.
就是整個舊約的發展就是說,.
神的創造原本是美好的,.
這個世界原本是一個好的世界,.
但是當罪進入了這個世界的時候,.
這個世界就再不是那麼美好,.
整個舊約的發展就是說,.
神用多次多方的不同的方法,.
來希望回復這個世界的美好,.

$^{121}$這個就是整個舊約的故事,.
但是歸根結底,.
有一樣東西要開始做的呢,.
神要革命的呢,.
不是從一個世界外面開始革命開始,.
是要由人的心開始做一個革命,.
因為唯有如果人的心不變的時候,.
沒得變的,很多東西沒得變,.
這個世界都沒得變,.
這個就是我們信仰當中一個很基本對人性的一種判斷,.
因為人的心要先開始變,.
如果人的心不變,.
外在一切給你的無論制度,無論幫助,.
無論怎樣給你不斷有鼓勵,支持,.
有很好的支持系統,沒用的,.
因為最終那個底層要打穩了,.
上面才可以起樓的嘛,.
而這個就是我們信仰怎樣去看人性,.
看自己,看世界,.
所以今天跟大家來到思想詩篇第51篇的時候,.
就是一篇很出名的一篇舊約的詩篇,.
就是大衛他自己的經歷,.
從此反省到我們究竟人活著的時候,.
究竟遇到一個什麼樣的問題呢?.
在還沒開始講這個詩篇第51篇的時候,.
首先講一點背景,.
大衛可能大家都知道了,.
他就是舊約當中一個很出名的一個君王,.
大家都聽過大衛打哥利亞的故事,.
很記得大衛是一個很厲害的戰士,.
或者一個很合神心意的一個君王,.
事實上聖經對於大衛的評語都是這樣說的,.
在大衛離開世界的時候,.
神對於大衛的評語就是說,.
「現在你的王位必不長久,.
耶和華已經尋著一個合他心意的人,.
立他做百姓的君王,因為沒有人遵行耶和華的心意.」.
我想這裡就是一開始的時候,.
在大衛還沒做王的時候,.
神就跟撒母耳說,.

$^{161}$就叫大衛告訴訴羅,.
神已經找到一個合他心意的人了,.
就是大衛..
到了大衛終結了一生的時候,.
神對於大衛的評價都是跟這裡一樣的,.
就說大衛是合他心意的人..
這樣的評價在聖經當中是很高的,.
很少有人經過了一生的時候,.
神可以對他的評價說他是一個合我心意的人,.
因為神的標準很高..
大衛就是一個這麼少數的,.
得到這樣的評價的君王..
但是你不要忘記,.
大衛其實在人生當中有一個很重要的污點,.
那個污點是什麼?.
就是我們今天所讀的那篇詩篇當中的背景..
我們今天一開始讀詩篇的時候,.
第一節就說,.
神求你,按你的慈愛廉恕我,.
按你豐盛的慈悲,.
塗抹我的過犯..
在我們的聖經一開始第一節就說這一句,.
但是如果你有聖經的時候,.
你揭開你的聖經來看,.
在這一句之前,第一節之前,.
就有一行的細字,.
那一行的細字寫什麼呢?.
那一行的細字就說,.
這首詩篇就是大衛犯罪和別氏巴犯罪之後,.
所寫的一篇詩篇..
這個是很有趣的,.
因為很多的聖經的詩篇,你都不知道它的背景是什麼,.
也沒有一個事件告訴你,.
這篇詩篇發生了這件事,.
但是少數有的,.
其中一篇就是詩篇51,這裡就有說,.
很清楚的說,.
大衛寫這篇詩的時候,.
他就做了一件很羞家的事,.
這件事可能大家都知道,.

$^{201}$根本你說出來,你都是覺得他很羞家,.
但是你作為一個公眾人物就更加羞家,.
那是什麼事呢?.
大衛見到別人的老婆美色當前,.
於是就勾引了別人的老婆,.
然後就跟她發生了一些關係,.
然後就搞大了別人的肚子,.
搞大了別人的肚子,然後就出事了,.
那個人是別人的老婆,那怎麼辦?.
後來最終就殺了她的老公,.
殺了她的老公之後,就名正言順,.
就變成寡婦,就搶了她過來,就做了她的老婆,.
然後就生了個兒子,.
就搞大了別人的肚子,.
但是生了個兒子,很快就夭折死了,.
這個就是大衛所做的事情,.
但是通常這些事情都是藏起來的,.
特別你做到皇帝的時候,.
如果你誰敢將這些以前的時代,.
將皇帝的醜事公佈出來,.
那當然是不會得到好結果的,.
皇帝當然會罰的,.
但是在聖經裡面就不是這樣,.
為什麼呢?.
因為在聖經裡面是誰將這件事情說出來的?.
是拿丹先知,.
就是說是神者說出來的,.
大衛又不敢殺神的先知,.
那就通了天了,.
於是後來就全世界都知道了,.
到了今時今日,過了這麼多年我們都知道,.
你想一下多羞恥,.
你想一下如果有時候你想你自己,.
可能有時候你做了一些事情,.
你都不體面的,.
你都很想快快地粉筆字抹掉,.
沒有人記得就算了,.
但是現在通了天了,怎麼辦?.
找洞鑽,.
如果在現在來說,如果是政治人物就要下台了,.

$^{241}$那些公眾的明星,.
又有這些就怎麼辦?.
被人杯葛,下架,.
不用自己有工作,.
一定是這樣的,.
但是大衛就是面對著一個這樣的國家級的醜聞,.
國家級的醜聞,.
而這一首的詩篇,.
就是在一個國家級的醜聞的正值黃金時期的時候,.
來發佈的,.
寫出來的,.
都給我們一個很深的反思,.
有時候其實這些事情本身沒有人想發生的,.
當發生了的時候,也都想盡量去掩飾的,.
這個是人性的一部份,.
因為這個就是很矛盾的,.
一方面我們就知道罪那些事情我們覺得不想見到,.
但是一發生的時候我們就想將它掩飾,.
一方面見到的時候我們又不想承認,.
當它公佈了出來的時候就將它掩飾,.
但是大衛在這裡就做了一個非常之反常態的事情,.
他不單止被人知道了,然後逃不掉之後,.
他還寫了詩篇,將他的日記,.
怎樣懺悔的日記,.
寫了成一首歌,.
然後給人來唱,.
成為了今天我們所讀的詩篇第51篇,.
這裡一開始他怎麼說呢?.
他說:「神啊,求你按你的慈愛廉恤我,.
按你豐盛的慈悲,淘沒我的過犯.」.
在這首詩篇一開始的時候,.
他就說出了一件事,.
他知道他現在已經跌到谷底了,.
怎樣有面見人呢?你想一下,.
沒有面見人,.
而怎麼辦呢?.
有什麼出路呢?找誰求救呢?.
譬如公司有醜聞的時候,.
就找一些公關公司來求救,.
看看怎樣去解決這個醜聞,.

$^{281}$但是現在沒有了,.
公關公司也幫不了他,.
這樣的時候,他就知道唯一一個出路,.
就是求神..
但是你想一下,這個其實也很矛盾的,.
為什麼?.
是誰將大衛的醜聞公開的?.
是神自己,.
而是誰見到大衛這樣的罪,.
而覺得這個是非常之不能容忍的?.
是神,公義的是神,.
將他的醜聞暴露的也是神,.
但是現在出路的也只有一個,.
就是回到神那裡..
我想這就是大衛這首詩篇一開始.
給我們看到整個框架就是這樣,.
最不能夠容許罪惡的發生,.
最不能夠容許人心裡面的污穢醜惡的,.
是神自己..
可能你自己有多污穢,你自己也不知道,.
但是神看得一乾二淨,.
神比你看得更加清楚,.
所以神對於你的污穢是比你自己更加珍惚的,.
有時候你說我都很恨自己,.
這樣的事情都想,這樣的事情都會做,.
但是神比你想讓你做的看得更加清楚,.
知道你的動機,知道你各樣的想法的時候,.
可能你還有更壞的,你還沒做出來,.
但是神都知道..
那你說神是不是最….
你自己都說恨自己的神,.
神不是更恨你,是不是?.
但是這就是聖經給我們一個很特別的出路,.
有時候面對罪惡的時候,.
我們在這個世界就說不聞不問,.
將它擺在地氈底下,.
好像出了事的時候找公關公司,.
好像粉飾太平那樣,.
然後好像沒有發生,.
然後不斷地跟自己說人性都是美好的,.

$^{321}$現在只是一時的掉在地上,.
現在慢慢地站起來就可以了,.
我們在這個世界當中不斷地用這種方法來處理,.
來對待自己裡面也好,那些黑暗的污穢,.
但是能不能處理呢?根本處理不了,.
就好像大衛一樣,現在已經去到一個國家級的醜聞,.
甚至殺人的老公搶別人的老婆,.
兒子自己也死了,.
怎麼處理呢?.
根本處理不了人性的罪惡,.
甚至你的軟弱,.
你自己的失敗所帶出來的結果也處理不了的時候,.
但是聖經給我們一個框架,.
唯一的出路就是回到那個可以定理罪的神的面前,.
一方面祂是最公義的神,.
但是在祂那裡你可以跟慈愛和憐恤相會,.
神的公義和慈愛就好像當祂面對罪人的時候就在那裡相會,.
神自己就是那個唯一可以塗抹祂過犯風聖慈悲的源頭,.
一開始的時候就是這樣去將整個詩篇定調,.
然後在這個詩篇裡面用了不同的名詞去講罪惡,.
第一個名詞就講到過犯,.
中文就叫做過犯transgression,.
神求你按你的慈愛憐恤我,.
按你風聖的慈悲塗抹我的過犯,.
第三節就說,我知道我的過犯罪常在我面前,.
過犯是什麼意思呢?.
過犯就是說越過了一些界線,.
譬如你的後丫或者你的前丸,.
那個是你的property line來的,.
那個是你的物業範圍來的,.
如果隔離屋放了一些東西過來,那就過了界線,.
這個就是過犯的意思,.
我們越過了神所為我們設定的那條界線,.
神做人那時候原是美好的,.
聖經講到神做人原本是美好的,.
但是我們在裡面就設計了很多的絞計出來,.
我們原本的美好的那條界線,.
我們不斷地踩了過去,去了過犯那裡,.
這個就是聖經第一個形容人心的黑暗情況,.
有時候你踩界是不經意的踩界的,.

$^{361}$但是有時候踩界你是知道都會踩的,.
小孩子頑皮的時候,有時候明明那裡是有一條路走的,.
我在那些屋村那裡住的,.
明明那裡有一條路走,要去一處去另一處,.
沿著那條路走一些樓梯,.
但是我很多時候都不走那條路,.
我很多時候就爬過欄杆,.
就喜歡爬那條斜坡,.
爬那條斜坡第一就好玩,.
第二就覺得近一些,.
但是其實別人爬那條斜坡就是不想你人走斜坡有時候會有危險,.
但是我的小孩就不理,.
我其中一個就是transgression,就是踩界,.
我不知道這樣踩界,我是知道的,.
知不知道很危險呢?我是知道的,.
知不知道我不應該走呢?.
知不知道那些草不應該踩呢?.
知道,.
但是我不喜歡聽,.
這就是transgression,.
就是踩界,我明知道有時候都是踩界,.
第二種可以形容的詞就是罪孽,.
Guilt,.
求你將我的罪孽洗除淨盡,並潔除我的罪,.
我是在罪孽裡生的,.
在我母親懷胎的時候就有了罪,.
這裡所說的罪孽,guilt,.
是有多重的意思的,.
其中有一種的意思就是,.
是一個心理上的重擔來的,.
罪孽或者罪咎感,.
都可以這樣叫guilt的,.
就是當你踩過界,.
或者你做了一些不應該做的事情的時候,.
你有一種的重擔,有一種的壓力,.
好像將你壓住了,.
或者好像你的原本白色的衣服,.
弄了一塊污漬進去,.
就是污穢了,.
所以他說罪孽和洗乾淨,.

$^{401}$就是這樣的意思,.
這個是說到罪的第二個面相,.
第二個層面,.
就是當我們去犯罪的時候,.
不單只是我們自己一時的行動方面,.
踩過界這麼簡單,.
其實是有後果的,.
後果就是你的心靈裡面有很多的重擔,.
你的心靈裡面有很多的重擔,.
你有你的guilt,.
你原本像白色的衣服,.
現在是弄得一塊塊的污漬垃圾,.
你自己都知道你的心都不是那麼乾淨,.
就是現在很容易,.
特別是在西方的社會的時候,.
很容易,.
你想人家閉嘴的時候,.
你說人家是racist,.
他就閉嘴,.
你說人家是種族主義就閉嘴,.
我們自己不是白人就沒有這些guilt,.
我們都不會閉嘴,.
我們會繼續說,.
但是那些白人那些人,.
他們一被人說是種族主義,.
他們就馬上閉嘴,.
好像按了一個按鈕被人按了,.
為什麼呢?.
因為他就有很多這些guilt,.
很多的罪的重擔,.
甚至乎他一個年輕人,.
只有十幾二十歲,.
他哪裡有買過奴隸,.
他奴隸都沒有見過,.
我們何來有奴隸給他買呢?.
但是因為不斷地覺得我們這些白人,.
以前賣過奴隸,.
很多這樣的罪,.
所以他就很怕人家說他種族主義,.
一說他是種族仇就閉嘴,.

$^{441}$因為就是很多guilt在人的心裡面,.
很多的罪孽,.
好像罪不單止是一個愛罪的行為,.
它是將我們的心好像重重的枷鎖,.
你是白的衣服也是將它污穢了,.
你是很輕鬆的,.
肩膀沒事的,.
現在有很多的重擔在你裡面,.
這是第二個他的形容是guilt,.
第三就是惡行,.
我向你犯了罪,唯獨得罪了你,.
在你眼前行了這惡,.
以致你責備我的時候顯為公義,.
判斷我的時候顯為清淨,.
第三個惡行的意思,.
evil doing,.
就是我所做出來的事情是evil的,是惡的,.
這個也是很有背後很有大衛很明白的重擔在其中,.
我們沒有人會看自己是evil的,坦白說,.
為什麼?.
如果你小時候玩過那些遊戲,.
角色扮演那些遊戲,.
沒有人想做那些壞蛋的,.
我記得我小時候跟附近的鄰居玩的時候,.
看那些卡通,.
大家就扮誰做好人,.
沒有人都想做壞蛋的,.
甚至乎就算猜輸了要做壞蛋的時候,.
玩完之後都很憤怒別人,.
玩完你不要再叫我那個名字了,.
你明白嗎?.
因為我們知道evil,.
做一個evil的人是很難接受的,.
即使是一個真正的evil的人,.
都不承認自己是一個evil的人,.
為什麼?.
因為你承認自己是一個evil的人,.
一個惡的人是很難接受的,.
根本你是不容許的,.
所以即使是那些殺人大魔頭,.

$^{481}$看過那些劇也好,.
看過那些小說也好,.
那些變態的人,.
他都不覺得自己是evil的,.
他心裡面都有一些機制,.
將他自己那些evil邪惡掩飾或者美化,.
但這裡大衛就說,.
他是做evil,.
evildoing,.
他做出了一些是evil的東西,.
惡恨就是這個意思,.
做出了一些是evil的東西,.
某個程度上我承認他是一個evil的人,.
一個evil的人就做evil的東西,.
所以大衛說到罪的時候,.
就有三個不同的面向,.
第一就是踩界,踩出了界,.
第二就是他將你的人污染,.
他將你的心靈加鎖,加重擔,.
但第三方面他明白到他所做的事情,.
好像是發自一個evil,.
一個邪惡的狀況之下去表現出來,.
這三個判斷都不是那麼容易接受,.
一般人不會這樣看自己,.
你想想是不是?.
有時候即使我們有這樣的體會,.
有一些自己的錯誤失敗,.
軟弱,.
我們都不想這樣去看自己,.
因為這樣是很難受的,.
你去看別人的黑暗,.
就好像看電視劇一樣,.
你覺得很順暢,很無所謂的,.
但如果你這樣去看自己的時候,.
你很難受的,.
所以我們在成長的過程當中,.
我們心靈裡面已經設立了很多自我防衛的機制,.
那些自我防衛的機制不是幫你去對付敵人的,.
當然有些敵人你要對付,.
譬如人家怎樣壓你,.

$^{521}$你從小到大自我防衛,.
惡人曲出來的時候,.
你懂得怎樣去踩人,.
但那些是少數的,.
很多自我防衛的機制,.
是你防備著自己不去看自己的黑暗,.
所以你明明見到這個人,.
為甚麼他會這樣,.
他講人的時候就講到人家錯的那些,.
講到那麼義正言善,.
但到他自己又做那件事的時候,.
他一點都看不到,.
他就說為甚麼做人做到這樣,.
兩把尺都可以的,.
但那個人完全都可以不知道的,.
為甚麼呢,因為我們每個人都是建立了這樣的機制,.
看別人就是這樣,.
但到自己的時候就好像落了沙下去,.
或者去洗白了,自我洗白了,.
在這樣的時候,.
但是有沒有出路呢?.
是沒有出路的,.
沒有出路的,.
為甚麼呢?.
因為根本就是沒有變的,.
你只是自己覺得好過了,.
你洗白了自己之後令自己覺得好過了,.
但是和神的關係,.
和人的關係,.
對於自己的罪疚的重擔的關係,.
全部都沒有處理過的,.
那真正的出路是怎樣呢?.
大衛很明白,.
他說首先第一步就是誠實地去面對自己,.
所以在詩篇第五十一篇第六節他說,.
「你所喜愛的是內裡的誠實,.
你在我隱密處必使我得智慧.」.
第一個大衛很明白的時候,.
他是需要去誠實地去面對自己,.
人的改變本身是需要對話的,.

$^{561}$是需要對話的,.
有時候你和一些親密的朋友,.
很深入地去對話的時候,.
一邊說一邊將你的心事打開的時候,.
說完之後你會發覺有些事情是需要改變的,.
說完之後你會發覺你的看法不同了,.
是不是?.
可能你也有這樣的經歷,.
但是現在人也很少有這樣的朋友,.
那怎麼辦呢?.
那就付錢去取消,.
取消者很多其實只是收錢和你聊聊天而已,.
就是因為在一個對話的過程當中,.
你打開自己,.
那是你改變的第一步,.
而事實上很多的對話,.
去到最深入的時候,.
甚至你最親密的朋友,.
可能你也未必說得到,.
或者你也不敢揭露出來去說,.
那時候最終那個真實,.
那個最深入的對話對象是誰呢?.
是誰呢?.
是神自己,.
在禱告的時候,.
是你和神對話的時候,.
即使你也不敢將自己打開的時候,.
但是你不怕你在神面前,.
因為神知道,.
你的每一個想法,.
甚至你不斷掩飾的那種恐懼,.
神完全都知道,.
你那些的痛苦,甚至你對自己的憎恨,.
神都知道,你也是在神面前跟祂說,.
那時候,.
當你真實的在神面前,.
願意真誠打開的時候,.
這個是改變的第一步,.
而事實上,所以這種改變,.
很多時候是從人來不了的,.

$^{601}$因為人與人之間,坦白說有多少真誠,.
而真誠去到什麼尺度呢?.
而去到神面前的時候,.
你完全打開的時候,.
那個就是改變真實的開始,.
接著就說,.
其實還有這裡說,.
「你在我隱密處必使我得智慧.」.
當我們真誠的打開自己在神面前的時候,.
神真的讓我們得到智慧,.
即是你跟一個很有智慧的人聊天的時候,.
有時候慢慢慢慢,.
你就好像一直打開自己去聊的時候,.
你看法各樣東西不同的時候,.
你好像真的得到智慧那樣,.
這個就是第一步,.
所以我們是在神面前,.
讓神來跟我們對話,.
祂一光光照我們,.
第二步祂說,.
而神的對話給我們智慧,.
當然就帶來了真理,.
而這樣的時候,.
幫助我們脫離罪的重擔和果子,.
祂說「求你使我得聽你歡喜快樂的聲音,.
使你所壓傷的骨頭可以踴躍.」.
當我們在神面前打開自己的時候,.
有時候我們被罪壓傷的傷害,.
是再一次的正面的去面對,.
有時候我們都不太願意正面的去面對自己被罪壓傷的傷害,.
好像有時候你明明跌得很痛,.
然後站起來的時候就裝沒事,.
「哎,沒事沒事.」.
為什麼呢?.
尤其是小孩子的時候,.
因為很瘀傷,跑來跑去跌在地上,.
然後站起來的時候,.
明明已經跌得很痛,.
站起來的時候裝沒事,.
希望還能跑前一點,.

$^{641}$讓人覺得你很厲害,很厲害,.
其實只是死撐而已,.
其實我們就是這樣,.
但是內心裡面被罪壓傷很多的重擔,.
我們是不斷的自己背著,.
當打開的時候,.
神就幫我們正面去處理這種傷口,.
「求你使我得聽你歡喜法律的聲音,.
使你所壓傷的骨頭可以踴躍.」.
我們在神面前去承認我們被壓傷,.
承認我們是需要醫治,.
承認我們是沒有喜樂的,.
承認我們是很久沒有聽過神的聲音,.
而讓神再一次將我們醫治,.
而當這樣的時候,.
我們再一次可以親近神,.
「不要丟棄我,使我離開你的面,.
不要從我收回你的聖靈.」.
大衛很明白當他將他的關係跟神打開的時候,.
他知道他原本是失去了一個關係,.
跟神的關係,.
當你很久都沒有能夠聽到神的聲音的時候,.
其中可能是有些罪你沒有處理,.
其中你沒有處理的罪,.
使你跟神的面離得很遠,.
你很久都沒有見過神的面,.
崇拜的時候本來是親近神的,.
朝見神的,.
好像都見不到神的面,.
禱告的時候都聽不到神的聲音,.
因為罪將這個關係斷絕了,.
特別當他打開的時候,.
大衛就再一次承認現在處於這樣的狀況,.
而且大衛這樣求他不要從我收回你的聖靈,.
在整個《舊約聖經》講到聖靈的時候,.
在不同的地方我隱隱約約講到神的靈,.
但是用「聖靈」這個字只出現了兩次,.
就是這裡,.
以及在《耳塞夏書》,.
講到神在將來復興祂的子民的時候,.

$^{681}$會將聖靈賜給祂的子民,.
所以這裡所講的聖靈真的是住在人心裡面,.
當人和神的關係,.
神的靈住在人心裡面的那個聖靈,.
但是因為罪的緣故,.
現在大衛都覺得勝靈好像被隱藏起來一樣,.
當他打開的時候出路是怎樣的呢?.
他說「求你用牛膝草洗潔淨我,我就乾淨,.
求你洗淨我,我就比雪更白,.
求你掩面不看我的罪,.
塗抹我一切的過犯.」.
這裡呢,.
他就求一件事,.
這裡你看到連續用那個「求」字,.
就出現了三次,.
「求你用牛膝草洗淨我,求你洗淨我,.
求你掩面不看我的罪.」.
這裡就求了三次,.
這裡三次其實就.
說出了一個景象,或者.
這個好像是描述了一幅圖畫一樣,.
「用牛膝草洗淨我」,為什麼我說用牛膝草洗淨我呢?.
在舊約的時候,用牛膝草很多時候是做什麼的呢?.
就是做一些潔淨的禮儀,.
如果你看《利未記》的時候,.
就是那些潔淨的禮儀,祭司就用一些血,.
去潔淨那些祭壇,.
或者潔淨一些不潔的地方的時候,.
就用牛膝草,.
但是這個牛膝草本身用潔淨的禮儀,.
用在人那裡,就只有一個地方是用的,.
在哪裡用呢?.
聖經那裡就說,.
在舊約的人,.
患了大麻風的時候,.
患了麻風病的時候,.
你知道在舊約的時候,其實.
直到近代,麻風才可以治療的,.
以前一直都不能治療麻風,.
是一種很恐怖的病,.

$^{721}$就是它在皮膚不斷擴散,.
接著又會傳染人,.
接著呢,恐怖到你的肉會爛,.
接著耳朵又掉下來,鼻子又掉下來,.
是這樣的,很恐怖的,.
所以人們要隔離的,.
因為不能治療,就把他困在一些.
沒有人去的地方,.
讓他自己慢慢去死,是不是?.
那麻風就是這樣的病,.
但是聖經那裡就說,.
如果有些人是麻風,.
得到潔淨的時候,.
祭司就用牛膝草,.
其中就是幫他潔淨,.
就代表這個人痊癒了,.
大衛在那裡求你用牛膝草,.
潔淨我,使我乾淨而已,.
我想其實他也可能想到的罪,.
就好像一個患了麻風病的人那樣,.
患了麻風病的人,.
是等死的,.
是感覺到很孤獨的,.
沒有人可以就近的,.
也都是自己是絕症,沒有出路的,.
也是會傳染的,.
但是他就求神說,.
用牛膝草使我潔淨,.
這個是他可以知道唯一的出路,.
誰可以將他的麻風病來潔淨呢?.
人的心好像發了風一樣,.
好像麻風病入侵,入心那樣的時候,.
就唯有神,.
而這個所謂的多次說,.
潔淨,洗滌,掩面不看的罪,.
這個是一個很圖像化的形象,.
就是罪好像將一些東西污染了,.
而那些東西就是我們的心,.
罪將整個人污染了,.
但是這些污染是擦不掉的,.

$^{761}$唯有怎樣呢?要用水洗,.
不知道你有沒有試過呢?.
舉例來說,有時候你行山到郊外,.
沒有水,但是踩了一些東西,.
或者弄了一些骯髒的東西,.
你不斷想擦,越擦,怎麼擦都不掉呢?.
你真的容易掉,有水就真是好了,.
現在就是這樣的情況,.
人的心,我們沒有水,.
但是擦不掉,越擦越黑,.
有時候有些骯髒的東西,.
你可能弄在手上,想著擦一下就擦掉,.
怎麼越擦,整隻手都黑了,.
那怎麼辦呢?我們現在就是活在這樣的狀況,.
我們的心靈的狀況,.
這不是說一些外在的骯髒,.
是說一個心靈的不潔淨的狀況,.
而這種心靈的不潔淨的狀況,.
我們根本沒有水,不斷地擦,擦不掉,.
我聽過一個真實的故事,.
不過是很噁心的,.
不過也不怕告訴你,.
以前你知道那些人去大陸旅遊,.
去大陸旅遊還沒有發展得那麼好,.
那些廁所下面有個木桶,.
因為要回收那些東西去做施肥,.
農業施肥,有個木桶在那裡,.
有一次有一個女士去完廁所,.
一個不經意踩錯腳,.
整個人就跌下去了,.
跌下去的時候她就尖叫一聲,.
外面的人就知道她有事,.
但她尖叫一聲之後就暈了,.
外面的人聽到她尖叫,.
所以她都可以尖叫,.
在暈之前尖叫,不是暈了就沒聲音,.
聽到她尖叫就破門進去,.
就扯她上來,.
扯她上來之後就沒什麼事了,.
只是骯髒一點而已,.

$^{801}$也不是一點,只是骯髒而已,.
沒事的,沒有受傷的,.
就洗澡等等,.
你想像一下如果是這樣的情況,.
給你些廁紙,任你抹,行不行?.
肯不肯?.
我們就是一個這樣的情況,.
水都沒有,只有廁紙,怎麼辦?.
那個女士就說她洗澡洗了很多次,.
洗了一整天都仍然覺得自己身上有味道,.
就是這一種,.
所以她說唯一可以幫她洗淨的就是神自己,.
大衛就說神啊,你洗淨我,你潔淨我,.
我們就好像染了這些掉進桶裡的情況,.
沒辦法來洗淨自己,.
而唯有神,唯有耶穌基督的羔羊的寶血,.
才可以將我們洗淨,.
也唯有神可以掩面不看我們一切的罪孽,.
祂說罪常在我的面前,.
我自己都看到,.
我自己有時候你犯了這個罪,.
你都不想去想,或者你傷害了一些人,.
你都不想去發生,.
但是發生了,你都不想去想,.
就好像不斷的那些畫面重新的出現在你的面前,.
我都洗不掉,.
我想不看,我想不想,.
好像不斷的畫一幅畫的畫面重新的出現,.
她說,但是神啊,求你掩面不看我的罪,.
然後她去求什麼呢?.
她說,神啊,求你為我做清潔的心使我裡面重新有正直的靈,.
神所要的滯就是憂傷的靈,.
神啊,憂傷痛苦的心,你必不輕按,.
所以大衛就求一樣東西,.
去求一個新的心,.
他不是說修修補補的,.
將你以前有什麼不足,.
有時候有什麼一時錯失錯手,改正下去,.
或者有時候想錯了,矯正一下就可以了,.
他不是,大衛知道我們落了一個情況,.

$^{841}$跟神的關係,自己心靈裡面那種污穢是不能改的,.
不能改的時候,去到最後的時候,.
他說,神啊,不如你給我做一個新的心,.
我要一個新的心,我不要一個舊的心去修修補補,.
我要一個新的心,.
而這個新的心,在新約的時候,.
高林多後書就說,.
若有人在基督裡,他就是新造的人,.
舊事已過,都變成新的了,.
他就是新造的人那個字,.
原本的意思不是說新造的人,.
而是新創造,是一個新的創造,.
好像在新天新地的時候,那個就是一個新的創造,.
舊的創造已經被罪污染了,.
要等到新天新地的時候,就是新創造,.
保羅就是在說這件事,.
將來整個救贖的計劃就是,.
當罪進入了這個世界的時候,.
這個世界最終會進入一個新的創造,.
整個救贖的過程,.
就是一個回復這個創造的歷史,.
新的創造的一個過程,.
但是這個世界未得到新創造之前,.
先要新創造的,.
想要新創造的,.
就是你和我,我們要有一個新的心,.
是一個新的心,.
所以大會說,神求你為我做一個清潔的心,.
使我裡面重新有正直的靈,.
而這個清潔的心,.
是怎樣可以換回來呢?.
就是說由一個憂傷痛悔的心那裡換的,.
有時候,我們以前買汽水喝呢,.
就要按樽的,是不是?.
就是你其實給了一些錢,.
然後你按樽的時候,你給一個樽,.
然後就拿回一些錢,是不是?.
就是說,我給一個罐子你行不行?.
我給一個罐子,可不可以給一些按樽的錢?.
不行的,為什麼?.

$^{881}$是要那個樽才能換回那些錢的,是不是?.
同樣的,那個新的靈怎樣換回來呢?.
如果當是換,怎樣換回來呢?.
其他的東西都換不了,.
是由一個憂傷痛悔的心,.
你將那個憂傷痛悔的心給神那裡,.
神才能夠將新的靈換來裡面,.
你沒有一個憂傷痛悔的心,其他的東西是換不了的,.
所以大衛是很知道的,.
而當我們有一個憂傷痛悔的心的時候,.
神的聖靈做一次的,來保證和我們同在,.
「求你使我認得救恩之樂,使我樂意的靈扶持我.」.
當我們有一個憂傷痛悔的心的時候,.
我們就再一次好像經歷了那種喜樂,.
前面說到當人的心被罪重壓的時候,.
人背著自己很多的重擔,.
人背著自己很多骨頭都好像被壓傷那樣,.
但是他說現在可以歡呼喜樂了,.
我的心可以再一次輕散了,.
我不需要自己背著那些罪的重擔,.
因為神為我們有一個新的心,.
神的靈再一次給我們,讓我們與我們同在,.
當我們這樣的時候,我們自己經歷了一條出路,.
我們也可以告訴別人這條出路是這樣走得通的,.
所以我要將你的道指教,指教有過分的人,.
罪人必歸順理,因為這條路是走得通的,.
我們是經歷過的時候,我們是這樣走過,.
唯有一條這樣的路,我們心靈的重擔才可以拿開,.
唯有一條這樣的路,這個世界是需要恢復的,.
這個世界現在是罪惡,污穢,紛爭,破碎,.
什麼都齊全,是需要恢復的,.
這個絕對不是一個神創造的完美世界,.
但是怎樣恢復呢?是要由人的心開始,.
而當這樣的時候,我們心裡面就可以湧出讚美,.
而當這樣的時候,我們的敬拜也可以恢復,.
神閱立我們的敬拜,我們與神的關係重新恢復,.
由神幫助我們,讓我們看到真實心靈的情況的時候,.
我們願意將一個憂傷痛苦的心帶到神的面前的時候,.
換回一個喜樂,與神恢復乾淨的心..
\newpage



\section{腓立比書 1:3-11}
\label{sec:031975}
\textbf{同心合意,興旺福音 (華)}~~~~陳明正牧師
\newline
\newline
連結: \href{http://www.church.com.hk/acms/content.asp?site=cdc\&op=show\&type=product\&code=031975\&layout=sermon}{\texttt{031975}} ~~~~ 語音日期: 2026-04-12
\newline
\newline
\hyperref[sec:031974]{\small{< < < PREV SERMON < < <}}
~
\hyperref[sec:index_chronic]{\small{[返順時目]}}
~
\hyperref[sec:031976]{\small{> > > NEXT SERMON > > >}}
\newline
\newline
腓立比書 1:3-11
\newline
\begin{longtable}{cl}
\hline
\hline
章節 & 經文 (和合本修訂版)\\
\hline
1:3 & \begin{tabularx}{0.7\textwidth}{X} 我每逢想念你們，就感謝我的神， \end{tabularx} \\ \\ \relax
1:4 & \begin{tabularx}{0.7\textwidth}{X} 每逢為你們眾人祈求的時候，總是歡歡喜喜地祈求， \end{tabularx} \\ \\ \relax
1:5 & \begin{tabularx}{0.7\textwidth}{X} 因為從第一天直到如今，你們都同心合意興旺福音。 \end{tabularx} \\ \\ \relax
1:6 & \begin{tabularx}{0.7\textwidth}{X} 我深信，那在你們心裡動了美好工作的，到了耶穌基督的日子必完成這工作。 \end{tabularx} \\ \\ \relax
1:7 & \begin{tabularx}{0.7\textwidth}{X} 我為你們眾人有這樣的想法原是應當的，因為你們常在我心裡；無論我是在捆鎖中，在辯明並證實福音的時候，你們都與我一同蒙恩。 \end{tabularx} \\ \\ \relax
1:8 & \begin{tabularx}{0.7\textwidth}{X} 我以基督耶穌的心腸切切想念你們眾人，這是神可以為我作證的。 \end{tabularx} \\ \\ \relax
1:9 & \begin{tabularx}{0.7\textwidth}{X} 我所禱告的就是：要你們的愛心，在知識和各樣見識上，不斷增長， \end{tabularx} \\ \\ \relax
1:10 & \begin{tabularx}{0.7\textwidth}{X} 使你們能分辨是非，在基督的日子作真誠無可指責的人， \end{tabularx} \\ \\ \relax
1:11 & \begin{tabularx}{0.7\textwidth}{X} 更靠著耶穌基督結滿仁義的果子，歸榮耀稱讚給神。 \end{tabularx} \\ \\
[1ex]
\hline
\hline
\end{longtable}
$^{1}$葛牧師,還有眾位牧者,同工,弟兄姊妹,大家平安!.
很榮幸,四五十年前認識葛牧師的時候,.
沒有想到有一天會來到洛杉磯..
我是厚的臉皮,請葛牧師讓我可以來這個地方朝聖,順便講個道..
我在這裡也遇到了靈異建牧師,也遇到了楊長老..
好像發現我們,還有別人,發現我們華人基督徒的世界其實很小,.
我們有很多共同的朋友,但是我們最好的朋友就是我們的主耶穌基督,.
阿們嗎?我們最好的朋友也是教會,教會是我們屬靈的母親,.
是我們在地上最美好最榮耀的家,阿們嗎?.
大家聲音就小了一點,覺得說我們在這個地上碰到很多很多的難處,.
有的時候有些難處也是我們在教會生活裡面必然會有的,.
但是我們今天要特別的講究,凡是在這個世界上,.
我們所碰到許多跟這個世界的人,都有一樣的困難,問題和難處,.
但是我們要求上帝給我們一個恩典,.
在教會裡面經歷憂愁之後的喜樂,患難之中的平安,.
還有所有這個世界所經過的一切一切,.
最終我們都認為是出於上帝的恩典和帶領,.
你的生命,你的教會生活,你的這個世界的生活可以在地如同在天,阿們嗎?.
所以我們今天要來講同心合一,興旺福音或者團隊事奉..
在這一段經文裡面,我們常常都會引用,.
我們會說上帝必定保守,我深信上帝的保守直到那日..
我們都會說這是講到應許的結果,各位,上帝給我們很多應許,.
教會是上帝創建,那麼他一定保守,我們會講到這個部分..
可是那是應許的結果,可是聖經95\%以上的應許都是有條件的..
華人教會最大的問題就是只講應許的結果,而沒有講當遵行的條件..
所以當我們碰到問題的時候,我們會禱告,.
可是我們不懂得我們要遵行這個條件,上帝就與我們同在..
舉個例子來說,聖經說,你們彼此相愛,別人就認出你們是我的門徒;.
你們彼此相愛,別人就認出我在你們中間,阿們嗎?.
對,這是條件,如果我們滿足了,就可以有這個結果..
這一段經文,我們要把它濃縮來看,因為否則很複雜,對不對?.
我把應許的條件怎麼說呢?就是同心合一,心望福音..
結果就是上帝就必成全這宮,可是這還太複雜..
我再簡念一下,就是同心合一,禱告必蒙垂聽,事奉必蒙成全..
好不好,我希望我們可以跟旁邊的人說,同心合一,禱告必蒙垂聽,事奉必蒙成全..
各位,我的教會也不大,我的教會在鄉下一點,大概有三百多人逐日崇拜..
在我這個宗派裡面,我是進信會,我是第一個華神畢業到進信會,募會的人..
特別是到有一點鄉下,又不是真的鄉下,又不是真的都市..
最後我們可能會看到我們教會的照片,所以我們教會既不在山區,也不在市區..
所以我常常跟他們說,我所在的教會叫做不山不市的教會,對吧?.

$^{41}$真正在市區很好發展,真正在鄉村可以得到別人的關注和支持,.
可是我在的地方是不山不市.可是我來到這個教會的時候,教會的光景是在荒涼低谷當中..
這也是我和師母今天一直覺得我們陪這個教會走過荒涼低谷,是我們的榮幸,是上帝對我們的恩典..
同心合意是多麼的大,禱告並蒙垂聽,你要知道你的事奉可以清神很多..
如果你藉著禱告能夠成全世公,各位,你的福是不知道好像裝上翅膀一樣,因為事奉必會成全..
我在教會裡面我們就經歷了這樣的恩典,我要鼓勵大家,最後我們會再說一些希望整個教會可以一起蒙福的話語,.
這叫做同心合意,心望福音.好的,我們讀過聖經..
第一個我們要有一個同心合意,要用恩賜為事奉的原則,什麼叫恩賜呢?.
舉個例子來說,我們教會如果有愛宴,一定要讓會做菜的人,像我們教會每一年有兩個梯隊的福音茶會,.
而五月底是我們的週年慶,所以我們就有週年慶的收割莊稼佈道會..
各位,不是那個時候才佈道,是那一天收割莊稼,前面一個半月每一個小組都辦福音茶會,.
茶會當中有很多形式,有餐會,有茶會,有音樂見證,有編織,有剁肉植物,有變魔術..
我們教會有一位年輕人是台大魔術社社長,他的爸爸現在做台灣道生神學院的院長,.
之前是企業家,把企業放掉,跑到愛丁堡去拿到神學博士回來,很不容易,.
他的兒子就變魔術,在小組裡面邀請新朋友,一個小組12個人,我們的福音茶會可以來到30個人..
因為我們在鄉下,如果我們在市中心離我們5公里,麥當勞的對面,我們租下活動中心,.
我們辦闖關活動可以來300多人,可是那邊去到教會,我們活動的人很難來到我們教會,因為我們教會在比較鄉下..
所以我們辦這些闖關,或辦福音茶會,或者這些,你要讓這個有恩賜的人來一起辦理..
每一個人都有恩賜,你不要說我老了沒恩賜,你老了仍然有恩賜,而且你的恩賜和經驗很夠..
你年輕一樣有恩賜,所以我們做剁肉植物,我們做很多很多很多的東西,.
但是教會愛宴,比較跟我們有關,我們一定要找會做愛宴的人..
那個香噴噴的菜味,會讓大家的胃被征服,心也被留下來,所以你知道做傳福音,很要對別人的胃口..
我兩年半以前,我在我們的城市辦了一個全台灣第一次的客家福音活動,客家很難傳福音的地方,對不對?.
我們教會在那個地方29年,大約有20幾間教會,沒有一間教會增長,因為我們在客家地區..
我們那天去辦客家文化館,館長,一個純粹的客家人,跟我說,"陳牧師,我們文化館成立20年來,從來沒有任何基督教機構來借用過,.
全台灣都沒有.".
那一天我們辦客家的群體的福音活動,很多的一半客家的老鄉親進來,結果480個人,我們坐了621個人..
然後外面的階梯上面坐的都是人..
我最感恩的是,那一天在開始之前的半個小時,外面排隊排了差不多50公尺..
各位,你說哪一個講員這麼有吸引力?是唐崇榮嗎?還是葛國光?還是哪一位?.
當然不可能嘛,對不對?.
沒有人會知道,沒有人會因為你的步道會就衝過來,對不對?.
就是這麼大一個袋子的禮物,裡面放了八落的衛生紙,八落!.
人家以為我們跟政府申請款項,不是,我們教會出的,多少萬..
所以大家都來了,來了之後,第一次我把眾教會聯絡起來,聯合師班,然後活動,讚美抄..
你不能在那邊一直傳福音,最後我上去做見證,呼召..
那個是一個突破,可是最美好的事情是,完畢之後我們發現,.
我們那邊的很多村裡辦公室,附近的好多村裡辦公室都發文,.
說他們在課文館辦了這個活動,所以這些禮物都是他們送的,都是我們送的..
可是他們就進來說,哇,這個活動這麼成功,從來沒有過的,.

$^{81}$但是是我這個里長,我這個村長辦的..
以恩賜為事奉的原則,在小組的裡面,在教會的裡面,在對外的福音事工裡面,我們都可以做..
當然很小的教會,不能夠都用恩賜事奉為原則,我只有20個人..
你說,弟兄你完畢之後幫忙打掃,他說我沒恩賜,對吧?.
當你晚餐在家裡面吃完飯,你的太太做非常豐盛的晚餐,.
我們家只有我和師母,永遠至少有三菜一湯,甚至四菜一湯..
有的時候我們常常也打開家庭,請弟兄姊妹進來,從我們進入教會開始,.
進入教會,就我們常常打開家庭,你不能吃完飯之後,然後師母說,.
牧師,你去洗個碗,你說,噢,我的護照不在這個地方,那就完蛋了,對不對?.
你要小心,你說的下一句話,會決定你一生的命運,明白嗎?.
跟你旁邊說話很重要,所以我們要用恩賜的事奉,為我們服侍的原則..
可是,有一個原則更重要,就是同心,就是愛心,對不對?.
所以你們以後,如果葛牧師安排你們服侍,你不能說,葛牧師,我沒這恩賜,對吧?.
你應該說,我要盡量的配合,但是在你的恩賜上能夠發揮,是事半功倍,對不對?.
事半功倍,所以你要知道在門口的招待很重要,你說,我很難笑,.
可是你做招待,你就要笑,對不對?.
而且你不是笑到那種笑是一個線條,然後都憋屈在裡面,憋成內傷,對吧?.
所以你要笑出聲音來,我最怕來到教會,大家都必須笑,可是都沒有笑出來..
我跟你說,這種笑沒有帶來造就,帶來內傷,所以你們以後要學會很多這樣的文化,要笑好嗎?.
我們很多牧者,我是講我自己,我們有的時候服侍久了,又不能把自己的軟弱跟別人說,對不對?.
可是碰到弟兄姊妹,我們又必須很好,但是我們的眉宇之間有一種濃得化不開的哀愁,對吧?.
這不可以的,我們的生命要改變,我們有受苦的心智來擺上,可是我們的心裡面卻充滿喜樂..
所以有的時候我看見弟兄姊妹服侍很認真,很努力,真的很棒很棒,可是我就問他們說累不累?.
弟兄姊妹,不好意思說累吧,對不對?.
弟兄姊妹總是覺得說,牧師你這麼忙,你這麼累,你看起來這麼瘦,你這麼累我們幫忙好嗎?.
不會累啦,我說累還是不累?跟牧師說沒關係,我們再累我們也要擺上嘛,不能讓牧師孤單..
我說累還是不累?繼續講下去就很累,很累又不能退,這很累,對不對?.
好,我都告訴大家,後來就有同工跟我說,牧師,我越累我越歡喜快樂,我越擺上,我越覺得可以重新得力..
最怕的不是累或不累,最怕的很多人說,不管怎麼累,我們靈陽堂這麼大,這麼多的需要,我還是得擺上..
可是在說的當中,他沒有說我服侍到失敗,我服侍到受傷都沒有,他很願意擺上,可是他散發出一種叫做淡淡的哀怨..
懂我意思嗎?這個比受傷還嚴重,因為淡淡的哀怨很像感冒病毒,看不見摸不著,每個人都中標,對吧?.
你知道教會的氛圍是很重要的,我們在服侍當中,你在服侍會不會累?會..
有個人的因素,身體,時間,體力,家庭的因素,還有整個很多很多的因素,可是沒問題的,你需要依靠主,你需要信靠教會,你需要牧者,小族長的幫助,讓你的心能夠真正的恢復,真正的有力量..
不要散發一種淡淡的哀怨,一定要,所以如果你敏感一點的話,你的太太坐在那邊,你的先生坐在那邊,你的主人坐在那邊,你看他一眼,你就知道他心靈裡面的故事,對不對?.
因為光的速度,光纖,傳輸,很大很大量啊,對不對?他不用講,他只要睫毛一低下來,哇,世界大事,對吧?.
所以我們要求上帝幫助我們,每一個人我們都能夠不但以恩賜來事奉,而且我們以受苦的心智來擺上,當你能夠受苦,當作你事奉中間,呼召當中更好的另外一半..
更棒的一面,如果你沒有這種觀念,價值觀,你在事奉當中,一遇到挫折,馬上就跌倒..
一個剛剛信主,剛剛來到教會的人,突然大發熱心,我們都要很小心,對吧?.
因為他一大發熱心之後,他就開始有很多的意見,他會跟牧師,跟童宮講,應該這樣,應該那樣,應該這樣,應該那樣,如果你沒有照著做,他就跌倒,他就不滿,不高興,對不起,不要對號入座..
我是說,我們應該讓每一個人學習先適應教會的文化,融入教會的文化,不要讓教會去適應你,不讓教會去融入你好嗎?.

$^{121}$我們要學會尊重整個教會的群體,和教會的領袖團隊的帶領,然後用受苦的心智來擺上..
因為只有這麼長的歲月的積澱,才能夠帶來一種新的文化,而且這個榜樣,在這個當中,是非常美好的幫助..
我講我的教會,可以嗎?我去到這個教會的時候,比較驚訝的事情是,我們的鋼琴彈下去,很多鍵都上不來..
然後我們兒子去參加兒童主日學,我們進去會有兒童主日家日光會,前後兩個半小時,上去彈琴的兒童不會彈,上去領唱的小朋友不會唱..
所以好幾個新來的家庭,因為我去的時候,這個教會前面三十年沒有一位新朋友,我去的時候,他們說三十年沒有一個新朋友加入這個教會,那我去之後開始有新朋友,新朋友的孩子就說不參加兒童主日學,我說為什麼?.
他說因為彈琴的不會彈,唱歌的不會唱,所以參加了兩個月,沒有唱過,不會唱任何一首兒童詩歌,沒有用N次啊,那我就說為什麼呢?.
我們的童工就說,我們要栽培超練我們的下一代,我說你也不能這樣栽培超練嘛,很麻煩對不對?鋼琴彈下去上不來,教會財政赤字,三十年沒有一個新來的人,請問你這樣子,新來的人怎麼會留下?不會..
到了1997,我到這個教會的那一年的12月24號,下午我們買了一台全新的YAMAHA三號鋼琴,當天調音完畢,晚上就開始用,而且當天來了一個非常有名的音樂家夫婦,他們說讓這個聖誕節聚會完全由我們來帶新的鋼琴,很棒..
這個人在台灣也在美國都稍有名氣,就是這一對音樂家,在我們很破落的教會平房裡面,他來N次的事奉,當天我們大家都很感動,可是受苦的心智,因為這一個音樂家他沒有嫌我們的場地,沒有嫌我們的鋼琴不好..
各位,那個時候教會很荒涼,我每天早上8點上班,8點1分會接到弟兄姊妹關心的電話,每天都接到,後來他發現我8點都到了,以後就不再關心了,其實我7點就到,各位,我7點就到,我10點半就出去探訪,因為教會都是老人家,我就一個一個探訪..
到12點回來,下午不能探訪,因為老人家在睡覺,我們教會都是老人家,下午3點半之後我又去探訪,探訪到5點回來,師母下班,我們就一起聊天,吃飯,洗碗,洗完碗之後,7點半我又出去探訪..
我們的慈善會主席兼董事長,也是我前任女傳道的先生,就跟我說,那時候茶經禱告會很特別,6點半到7點半就結束,7點半一到,必須立刻結束,不管聖靈有沒有在場,我們都必須結束..
我講到的時候,當我看到我的同工在第一排看手錶的時候,我就知道我必須在一分鐘之內結束,因為他們是學科技的,我有一天茶經禱告會完畢,7點半我要出門,他看我要出門,他說,陳傳道,我們這個地區晚上7點半是沒有人出門的,你還要出門嗎?.
我說,我去探訪,因為我們這個教會30年沒有新來的人,我就出去探訪,去到石門水庫邊邊很黑的社區,一家一家探訪,因為這些男生都沒有信主,我每天晚上,我原來我的身體不好,我30年前,29年前來到這個教會之後,我才50公斤..
你們如果看到葛牧師現在頭髮逐漸變白,你們除了覺得他很有智慧之外,你們也要多憐憫他,對不對,年紀慢慢稍長嘛..
我那個時候,大家一看,就是化療做到一半,還沒有恢復的樣子,吸毒還在戒毒的過程當中,50公斤啊,對不對,然後一副言語粗俗,氣貌不揚,又沒有保羅的能力的那種人..
但是我就很殷勤各位,我就這樣子每天探訪,我去探訪的時候,附近那個時候有一個神學院,那我這個人是這個樣子嘛,我就騎著腳踏車,白天腳踏車探訪,到最後我就去那個神學院坐在那邊,跟大家聊天啦,哈啦哈啦,反正師母在上班嘛,然後完畢再隨便吃個飯或回家吃個飯..
我去的時候呢,那個神學院的師母是一個比較有想法的人,我一去就跟我說,陳傳道我跟你說,華神的畢業生,沒有人會牧養的,你們都沒有學牧養,都只是讀書,不瞭解民間疾苦,不瞭解民間疾苦..
我說是是是是,因為我們剛到那個地方嘛,你只能說yes yes yes,你說no no no就完了,結果我每天都來耶,我每天都來,不到一個月我的名聲就傳遍整個地區,說這個教會來了一個很喜歡探訪的人,所有的地方都一直看到我的身影穿梭在其中,又瘦又醜..
一個月後,這個神學院的人就跟我說,我們敢說,將來你一定是這個地區最大的教會,才一個月.果真當然我們的教會就成長了比較多倍,我覺得受苦的心智是用喜樂的心在運行了..
後來我去這些每一家,我不能喝茶,可是我必須跟他們喝茶,我不能喝咖啡,可是我必須跟他們喝咖啡,我一家一家,結果每一個家庭的男主人都幸福了,包含唯一的醫生也都幸福了..
那個時候有沒有受苦的心智?有.我這樣子講,希望你們不要這樣來衡量,當時教會問你要不要住在教會,一個16坪的平房,是中式建築,所以裡面有很多動物,不是他侵犯我們,是我們侵犯了他..
我們進去住的時候不能睡覺,因為裡面都是螞蟻,都是螞蟻,還有外面都是有蟒蛇,把門一開,這一條蟒蛇兩公尺長,慢慢從你面前經過,你要知道,你就說,平平安安的去吧,你千萬不要怕到牠,好幾次..
那麼我們一下雨,每個地方都會漏水,會堂上面被燒過,鋼琴彈下去起不來,而我一去之後,我們現在華神的學務長,才子高,是我一起長大,比我小一些,帶著他一起長大的大專生,還有很多這樣的人,華神都是很多好朋友,現在新進去的老師,也是我們帶的下一代李修齊老師..
那麼他就來找我,還有很多從紐約來的傳道人都來找我,叫我不要留在這地方,沒前途,他們說這只是退休的人待的地方,我說沒問題,我到這個地方上帝帶領,我原則上一輩子不走,哪裡找我,我都不走,薪水是很低的,但是我們的想法就是留在這邊..
所以你知道,有一天,我們的一個長輩,也是葛慕斯認識的,可能就叫藏西梅,是我唯一從屏東巡禮會到這個地方的,所以他就算我的長輩,他就來看我,坐在我的側房,側房的椅子上,對面就是牆壁,這個牆壁就跟我們的螢幕差不多,那邊螢幕差不多大,他坐在那邊..
因為我們不能睡嘛,因為螞蟻太多,房間太小,他們裝了一個天花板,上面每天都有貓來活動,貓來尿尿,之後才走,貓尿很毒,你開冷氣,你就在毒的空氣當中,你不開冷氣,那個平房溫度35度到40度不能睡覺,不能睡覺,那別的地方又都是螞蟻,我的兒子就因此一直住院,因為貓尿很毒..
不過這個環境真的很差,所以我的那個,就是以前巡禮會的姐姐,有一天晚上來看我,差不多一年了,她來看我,坐在那邊,她一看,清潔衛生公司來過,地上的螞蟻,這一面牆地上的螞蟻,死掉的,可以掃成三分之一個本機呀,螞蟻耶,螞蟻不是蟑螂,蟑螂掃成三分之一個本機,.
就很恐怖了對不對,可是是螞蟻,可是在牆壁上的螞蟻呢,就像燒餅上的芝麻一樣多,所以我很長的時間不太敢吃燒餅,開玩笑,照吃不誤..
你看我的太太是不是很像公主,對不對,可是就很願意進到教會裡面生活下去,那當然太多的事情我就不講,可是今天都是我們的勳章,今天都是覺得上帝給我們特別的恩典,.
我的太太說,如果你要錢多事少離家近,你就繼續做公務員,委員是國家的公務員,而且做到一定的位階,她說如果你希望輕鬆一點,你繼續做公務員,你不要做傳道人,.
要做傳道人就去你最需要你的地方,最複雜的地方,她還加了一句話,錢少一點的地方,鄉下的地方,因為這才是需要你的地方,她說華神畢業生都去都會區,當然都會區更大的需要,.
她說我們兩個去鄉下一點好不好,我們一起商量,不是她要我去,按著我的頭去,不是,我是歡喜快樂的一同去到這個地方,我在華神的畢業生裡面有兩次的分享,.
我就說我和師母可以到一個遙遠,不為人所知的角落同心服事,這是上帝的祝福,也是我們最大的幸福,不過我說這種境界非你們所能了解,.
同學就一直笑,學校的老師一直希望我可以常回華神跟他們分享,分享就是做一個真正的牧養者,特別是到鄉下教會牧養,.
沒有這麼悲慘,可以讓以後的畢業生願意到鄉下去牧養傳福音,建立教會,栽培門徒,也可以成為榮耀的教會,但是前提是你要有一個受苦的心智來擺上,.
我去到這個教會的時候是赤字,教會很紛擾,教會還有兩家自己請了一個傳道人進來,因為他們不滿意前任傳道,我一去嚇了一跳,.
我不是明媒正娶進來的嗎,怎麼旁邊的小門又進來一個呢,而且他不上班領全職薪水,這個很複雜,所以我一輩子唯一一次住院,就是來到這個教會吃不下睡不著,.
又在鄉下就突然來了二十位非常棒的弟兄姊妹,沒辦法想像,他們從來都不認識任何一個人,十一月十二月來了二十個人,.
我們把裝備更新這些,教會從赤字變成政治,還買了這麼多東西,第二年各位,我們結餘了超過十萬美元,十萬美元,.

$^{161}$第三年結餘了八萬美元,另外兩萬美元為九二一大地震奉獻,第四年我結餘了二十萬美元,一個很小的教會,都是年長的居多,但是整個教會很大的改變,.
每一年只增長二十個人,不多,因為搬走的速度很快,搬走的速度很快,你還要增長二十個人,表示你要來將近五十個人,.
因為每個人聚會的比例可能只有百分之六十,所以神的恩典,我覺得上帝從來沒有從天上跟我說,明正我的愛子,我喜悅你,聽到我會嚇死,.
直接上天,可是上帝會用拐彎我腳的方式,用你沒有辦法用肉眼或實質看見的方式,從側面一直告訴你說,孩子,你所做的我知道,.
你捨棄了高薪,你願意來到鄉下,來到鄉下一點都沒有什麼特別,你願意埋在這個地方,我也願意祝福你,雖然別人看不出來,但是受苦的心智很重要,.
師比獸更為有福,我們服侍上帝不是一直得,我們服侍上帝應該背十字架,其實這個意義就是說,鐵器放下這個世界上從父母兒女田產,.
所有一切,但上帝應許說,沒有不在今生並來世得到百倍的,你們還敢阿悶,有經過深思熟慮嗎,受苦的心智奉獻來擺上,很重要,.
第三個我們要講,教會一定要建立團隊的事奉,這個很重要,第一個要講的是同心合意,分工合作,就是我們一定要有這個概念,.
有人專注一個工作,但他要跟別人合起來,很多人一做一個工作,誰都不可以碰,我覺得這個態度是不對的,我們要有界限,也要有開放,.
我們要尊重服侍的人,也不要去彼此搶,就是分工但是要合作,讓教會各有所施,各施其職,但是又能夠合在一起,我覺得我在我的教會沒這麼忙,.
在這些文化建立之後,凡是別人可以做的,都是別人做,我很輕鬆,我就是走來走去,看看大家,讓大家看看我,大笑幾聲,江湖一笑明恩仇,.
很快,我們教會是一個很健康溫暖的教會,很多人來都說你的教會為什麼會這樣,因為超過兩百人還這麼健康溫暖,比較少,比較少,.
當然在我們台灣的比較小的區,超過五十人的進信會就算是水準以上,超過一百人的我沒有聽過,可能是我孤陋寡聞,超過一百二十人的進信會可能只有我們一間,.
對不起,這個不是自誇,神的恩典,超過到三百多人,這個比較少見,而且我又是華神的畢業生,第一個到進信會的,我覺得我的歷史包袱很重,.
我可能是,我覺得我也不容易,但是求上帝幫助我,第二個有效率感,有效率感很重要,就是你的福祉一定要有果效,.
你不能說上帝的施工沒有果效,不是,上帝交託給我要忠心全力擺上,一定要有一種信念,教會是可以增長的,阿們嗎?.
阿們,聲音小很多,阿們嗎?對,否則我的腳會踢出來,阿們嗎?要有效率感,各位你知道,這很重要,你有效率感,這個包含了我們的生命會改變,.
各位,生命會改變,我記得,我當時我們家每天去,我那時候還在一個進信會,我們禮拜六,我就是去讀華神,但是是代職,然後我叫實習,可是我就問這個教會說,我可以做什麼?.
做義工嗎?不是說帶大專生,我說那大專生有多少?他說,零,陳哥,零,你知道我怎麼做呢?我每個禮拜六,我政府機關下班之後,我就趕到另外一個城市,然後我就趕到中央大學,我們台灣一個中央大學去站崗,.
我去找一個我認識的人,然後請他把我帶到每一個宿舍去,把認識有接觸過一點點的人全部帶出來,晚上一起去吃晚餐,晚餐,然後晚上帶他們聚會,聚完會,全部到我們家住另外一個城市,帶到我們家一起吃宵夜,然後洗澡,他們都很高興,.
然後第二天再全部帶到教會,我開成人主義學,會前禱告,我是崇拜主之士,我們那個時候的禱告都要跪在前面,以前就不用跪,因為只有主任牧師私會跟私請三個禱告,就我現在一帶二十多個人,全部跪在前面禱告,.
中午我們再帶他們吃飯,一個大專生就在短短的時間從零變成二十多個人,但是我覺得這是我們的福分,因為我們最大的改變是我很節儉的一個人,在成功嶺的大內窟我穿了十年,師母最後命令我丟掉,偷偷丟掉,我才沒穿,.
那麼我洗澡的時候呢,洗澡的時候瓦斯熱水器一開,離我的浴室有六公尺嘛,我就進去沖,沖水一分半就沖了,然後我就開始抹肥皂,然後我就跟他說趕快把瓦斯爐關掉,然後這六公尺的水沖在我身上,就剛好用完,我用的又是南橋水晶肥皂,從頭頂到腳底,.
當他們出來的時候非常優雅,把浴室的門打開,那個水氣跟肝病一樣沖出來,我的心臟都碎了,各位,可是你為了要服侍,你要不要改變,你要啊,我就慢慢的改變了,就慢慢從吝嗇或者節儉變為慷慨大方,然後我的教會現在也建立慷慨文化,.
所有來的人,講員,貴賓都高興到不行,所以我們教會裡面還有一個四星級旅館,旅館房,大家都很驚訝,我們的地下室是全台灣我可以說是最好的下一個世代的敬拜中心,很多人都來參觀,不管是燈光,我們的螢幕是202吋的LED,.
我們有兩套全台灣,我們的音控是跟心電行道會一樣的等級,可是我們兩套,一個鄉下教會,所以我覺得我們是可以的,求上帝來幫助我們,受苦的心智,求上帝幫助我們,同時呢,我們在團隊的裡面分工合作,.
你要學會一些功課,分工合作,就是不是一個人做所有的工作,是每一個人做一部分的工作,而且是配搭,像我們教會開福音車的人有30位,30位,因為我們那個地方沒有公共交通系統,禮拜三要開福音車,配搭計程車,禮拜六要開福音車,還要開四個方向,每個方向大約10公里的路線去接送青少年,.
以前我也開,到禮拜天也是一樣,福音車加上計程車,四個方向,包含所有的童工,我們有30位來開福音車,但是我們要分工合作,要學會我們宣教之父William Carey給大家的勸告,.
那個勸告就是,對不起,我有時候弄太慢了,這個已經是結論,第一個就是對人的過失不要耿耿於懷,不要耿耿於懷,要寬大,第二個很重要,對人的優點常常肯定有加,務必做到用愛心互相寬容,.
以禱告將和他人的誤會可以解開,分工合作,主要的領袖都要用這些原則來做一些事情,最後我要說,其中一個原則就是常常相處,常常交通,自然而然達到共識,教會的意象常常是這樣產生出來的,不是憑空真空裡面冒出來的,.
第三個,常常一個人宣稱自己得到意象,而且別人毫無感動,這是沒有意義的,真正的意象會感動別人的,需要時間的等待來印證上帝的作為,.
第四個,這沒有第四個,對不起,我的眼睛看不太清楚,所以這些我就不講多,我要講最後的原則,就是一起禱告,常常一起禱告,我看了你們的禱告會有三十多人,.
我鼓勵,我們中間是小組長的請舉手,小組長都特別知道我要來,講到都沒來,你是執事的請舉手,你們可能在別堂崇拜,.
我鼓勵,如果你是領袖,你是小組長,你是副小組長,我鼓勵你要排除萬難來參加實體禱告會,多久一次各教會來決定,因為疫情之後,.
但是你知道我在那個鄉下,青少年住在這個城市,可是他們晚上不能出來,社青都住在外地,大專生都住在我們台灣的中部南部,老年人都不能出來,.
可是我的實體禱告會,我們大約這樣,一百人,在台灣我有問過別人說,比較鄉下的一個教會,禱告會可以多少人,附近的教會大概人都很少,.
整個台灣,像這樣的鄉下,當然我們有時候不一定一百人,你要知道都從遠方,甚至下班回來,但是有相當的人數,有時候一百多人,有的時候是這樣,.
當然我不可能每個禮拜,因為我們教會太偏遠了,可是可以到這樣,在晚上,所以我覺得是很棒的事情,希望大家,.
我們禱告會做了設計,讓每一個人來可以有感覺到上帝的同在,我們不用購物清單式的禱告,而是敬拜詩歌或讚美與禱告呼求交織的方式,.
那我要說就是,禱告用力,事奉神力,禱告神力,事奉用力,禱告可以蒙垂聽,事奉可以蒙成全,是和這個有很大的關係,.

$^{201}$第五個,我們要持守同宮中間,起初的愛,這一點我就是還有四五分鐘可以做結論,因為我覺得這個很重要,.
同宮中間有的時候因為工作服飾種種的配搭,有的時候會有不同的意見,我要鼓勵大家,在不同的意見,意見可以不同,.
愛心卻要相同,我們在教會要做社會的典範,不要做社會的壞的榜樣,要做這種榜樣,所以我們同宮之間應該有一種情愛的關係,.
超越服飾,可以彼此探訪,可以有年度的同宮退休會,不是聚焦在工作上,而是有歡樂的時光,各位這一點再怎麼強調也不為過,.
我們每個人都會受傷,每個人都在過程裡面被誤會,但是你仍然散發出這種光芒是很重要的,.
我最後因為實在有很多,我覺得需要講的是重燃異象,重回共識,可是我可能最想講的一點,.
就是你要知道你在教會做一個同宮弟兄姊妹,比天使更重要,比魔鬼更重要,.
你說我從來沒聽過這個道,我把這個告訴你,魔鬼攻擊教會正常不正常,正常不正常,正常,告訴我正常不正常,正常,.
攻擊是他的本分,是他的生命的特質,他不攻擊活不下去,魔鬼攻擊教會正常,.
唯有今天你在教會裡面做重要同宮,做會友,你跟葛牧師,跟我們這些人相處,.
深知我們的軟弱,深知教會的不夠,深知教會的虧欠,可是你在核心當中,你看到了這一切,.
你出去說,你不要看這個牧師,那個長老,那個執事道貌岸然,我告訴你其實他怎樣怎樣怎樣,.
這個破壞力比魔鬼說的高一百倍一千倍,因為你是跟他在一起五年十年二十年,.
假設說今天葛牧師說如果活到兩百歲,搞不好你是他一百八十年的同宮,.
可是你說我可以出來做見證,他多差多差多差,我告訴你這個嚴重性是不得了的,.
可是如果你說我跟這個同宮,同宮雖然有軟弱可是滿有上帝的恩典,.
教會有軟弱正是我們可以擺上的機會,他是上帝好像在摩西保羅身上,.
他們看見以色列百姓逼迫他們,毀謗他們,對抗他們,.
可是他們對神的百姓那種驚人執著的愛沒有改變,同樣的我們看見神的僕人為我們流淚歎息,.
為我們深夜睡不著,白天吃不下,你都沒有看到,.
可是當你看見神的僕人或教會的領導者有不夠的時候,你卻說他很棒,.
他是上帝給我們的神的僕人,是帶領我們的,我們說他很棒,.
其實這個力量是遠比天使說的更重要,對天使跟你我不一樣,.
同樣的如果你毀謗他,這個是比魔鬼更嚴重,.
你以前都沒聽過,所以跟你旁邊說你很重要,.
葛慕斯的命運在你手中,真的很重要,.
如果我們靈陽堂的弟兄姊妹出去都說我們教會的牧者多棒多棒多棒,.
我們的教會多棒多棒多棒,我告訴你復興會快快來到,.
你不信主的家人會很快信主,你不信主的朋友會很快信主,.
因為這樣的福音傳揚比天使更重要,比魔鬼的破壞我們更超越,.
所以讓我們的教會可以成為這樣一個蒙福的教會,.
他會帶來一種特質是人間所有的規條,所有的鼓勵都達不到的,.
因為這個是超越成敗得失缺點失敗而帶來的上帝超越性的祝福,.
永恆性的喜悅可以澆灌在你我的身上,阿們嗎?.
跟你旁邊,你的教會如果不會說阿們,不會復興,.
所以跟你旁邊大聲說阿們阿們阿們,.
\newpage



\section{撒母耳記上 24:1-19}
\label{sec:031976}
\textbf{以善勝惡}~~~~李思聰牧師
\newline
\newline
連結: \href{http://www.church.com.hk/acms/content.asp?site=cdc\&op=show\&type=product\&code=031976\&layout=sermon}{\texttt{031976}} ~~~~ 語音日期: 2026-04-12
\newline
\newline
\hyperref[sec:031975]{\small{< < < PREV SERMON < < <}}
~
\hyperref[sec:index_chronic]{\small{[返順時目]}}
~
\hyperref[sec:031977]{\small{> > > NEXT SERMON > > >}}
\newline
\newline
撒母耳記上 24:1-19
\newline
\begin{longtable}{cl}
\hline
\hline
章節 & 經文 (和合本修訂版)\\
\hline
24:1 & \begin{tabularx}{0.7\textwidth}{X} 掃羅追趕非利士人回來，有人告訴他說：「看哪，大衛在隱‧基底的曠野。」 \end{tabularx} \\ \\ \relax
24:2 & \begin{tabularx}{0.7\textwidth}{X} 掃羅就從全以色列中挑選三千精兵，往野山羊磐石的東邊去，尋索大衛和他的人。 \end{tabularx} \\ \\ \relax
24:3 & \begin{tabularx}{0.7\textwidth}{X} 到了路旁的羊圈，在那裡有個洞，掃羅進去大解。大衛和他的人正藏在洞裡的深處。 \end{tabularx} \\ \\ \relax
24:4 & \begin{tabularx}{0.7\textwidth}{X} 大衛的人對大衛說：「看哪，這日子到了！耶和華曾對你說：『看哪，我要將你的仇敵交在你手裡，你可以照你看為好的對待他。』」大衛就起來，悄悄地割下掃羅外袍的衣角。 \end{tabularx} \\ \\ \relax
24:5 & \begin{tabularx}{0.7\textwidth}{X} 隨後大衛心中自責，因為他割下了掃羅的衣角。 \end{tabularx} \\ \\ \relax
24:6 & \begin{tabularx}{0.7\textwidth}{X} 他對他的人說：「耶和華絕不允許我對我的主，耶和華的受膏者做這事，伸手害他，因為他是耶和華的受膏者。」 \end{tabularx} \\ \\ \relax
24:7 & \begin{tabularx}{0.7\textwidth}{X} 大衛用這話勸阻他的人，不許他們起來害掃羅。掃羅起來，從洞裡出去，預備上路。 \end{tabularx} \\ \\ \relax
24:8 & \begin{tabularx}{0.7\textwidth}{X} 然後大衛也起來，從洞裡出去，呼喚掃羅說：「我主，我王！」掃羅回頭觀看，大衛就屈身，臉伏於地下拜。 \end{tabularx} \\ \\ \relax
24:9 & \begin{tabularx}{0.7\textwidth}{X} 大衛對掃羅說：「你為何聽信人的讒言，說『看哪，大衛想要害你』呢？ \end{tabularx} \\ \\ \relax
24:10 & \begin{tabularx}{0.7\textwidth}{X} 看哪，今日你親眼看見，在洞中耶和華將你交在我手裡。有人要我殺你，我卻愛惜你，說：『我不敢伸手害我的主，因為他是耶和華的受膏者。』 \end{tabularx} \\ \\ \relax
24:11 & \begin{tabularx}{0.7\textwidth}{X} 我父啊，請看，看你外袍的衣角在我手中。我割下你外袍的衣角，卻沒有殺你。你知道，並且看見我沒有惡意要悖逆你。你雖然要獵取我的命，我卻沒有得罪你。 \end{tabularx} \\ \\ \relax
24:12 & \begin{tabularx}{0.7\textwidth}{X} 願耶和華在你我中間判斷，願耶和華在你身上為我伸冤，我卻不親手加害於你。 \end{tabularx} \\ \\ \relax
24:13 & \begin{tabularx}{0.7\textwidth}{X} 古人有句俗語說：『惡事出於惡人。』我卻不親手加害於你。 \end{tabularx} \\ \\ \relax
24:14 & \begin{tabularx}{0.7\textwidth}{X} 以色列王出來要尋找誰呢？你要追趕誰呢？不過是一條死狗，一隻跳蚤而已。 \end{tabularx} \\ \\ \relax
24:15 & \begin{tabularx}{0.7\textwidth}{X} 願耶和華作仲裁者，在你我中間判斷。願他鑒察，為我伸冤，救我脫離你的手。」 \end{tabularx} \\ \\ \relax
24:16 & \begin{tabularx}{0.7\textwidth}{X} 大衛向掃羅說完了這些話，掃羅說：「我兒大衛，這是你的聲音嗎？」於是掃羅放聲大哭， \end{tabularx} \\ \\ \relax
24:17 & \begin{tabularx}{0.7\textwidth}{X} 對大衛說：「你比我公義，因為你以善待我，我卻以惡待你。 \end{tabularx} \\ \\ \relax
24:18 & \begin{tabularx}{0.7\textwidth}{X} 今日你已顯明是以善待我，因為耶和華將我交在你手裡，你卻沒有殺我。 \end{tabularx} \\ \\ \relax
24:19 & \begin{tabularx}{0.7\textwidth}{X} 人若遇見仇敵，豈肯放他平安上路呢？願耶和華因你今日向我所做的，以善回報你。 \end{tabularx} \\ \\
[1ex]
\hline
\hline
\end{longtable}
$^{1}$上一次我們講到.
撒慕爾記的二十二,二十三章.
當時大衛正在逃避蘇俄追殺的過程.
在這個曠野裡面他體會到神是他真正的倚靠.
亦明白到他是我們人生的避難所.
今天來到二十四章.
我們看到這個圖畫.
大家猜一下是講什麼的.
大衛割蘇俄的衣襟在哪裡.
在山洞裡面 很明顯是在山洞裡面.
我問問你們當時蘇俄在做什麼.
這個很記得 蘇俄正在大解.
衣襟都不要在他後面.
所以這個是撒慕爾記第二十四章.
我們看看經文 第一節我們一起讀.
蘇俄追趕菲律賓人回來.
有人告訴他說大衛在隱基底的曠野.
隱基底 這是一個死海沿岸的地方.
我們知道猶大曠野沒有什麼水源.
但是去到死海旁邊有一個水源.
這是一個豐富的水源.
而且因為有水的地方.
有一個綠洲 有植物的生長.
一個非常古老的地方.
今天這個地方還存不存在.
存不存在 存在吧.
因為我們去過.
認不認得這些人.
很模糊 看不到.
這個清楚一點 比較大.
看到那個水泉.
其實隱基底這個隱字.
N就是希伯來文水泉的意思.
基底其實只是野山羊.
所以這個地方就是野山羊.
剛才我們也看到有些野山羊在那裡.
為什麼呢 因為水源嘛.
所以這個叫做野山羊的水泉.
大家看到這個水泉 但是更加看見就是這裡.
就是什麼 就是很多山洞.

$^{41}$所以當時這個蘇羅人有三急.
進去解急.
誰知道剛好誰躲在山洞裡面呢.
就是大衛和跟隨他的人.
所以我們看到這個就剛好.
也是好機會來了.
所以我們讀完第二節 好不好 請大家一起讀.
蘇羅就從以色列中挑選三千精兵.
率領他們往野羊的盤石.
岑藏著大衛和跟隨他的人.
到了路旁的羊圈在那裡有洞.
蘇羅進去大解.
大衛和跟隨他的人正藏在洞裡的深處.
跟隨的人對大衛說.
要話曾應許你說.
我要將你的仇敵交在你手裡.
你可以任意待他.
如今時候到了.
大衛就起來悄悄地割下蘇羅外袍的衣襟.
所以這個就是剛才我們所描寫的情況.
但是大衛有沒有割呢.
有 有.
割完之後 其實他不單可以割他的外袍.
還可以割他的什麼.
割他的頭 其實是可以殺他的.
不過他沒有.
他猶豫不過割了外袍.
不過只是割了外袍.
大衛的反應都是這樣.
以下我們一起讀.
用啟認的方式讀 好不好.
請弟兄姊妹你讀一節.
然後我讀一節.
我們交替讀一下.
第五節請你們讀出.
請聽我讀.
對庚從他的人說.
我的主乃是耶和華的受告者.
我在耶和華面前萬不敢伸手害他.
因他是耶和華的受告者.

$^{81}$隨後大衛也起來從洞裡出去.
呼叫 蘇羅啊 我主我王.
大衛就屈身 俯伏於地下拜.
今日你親眼看見在洞中.
耶和華將你交在我手裡.
有人叫我殺你.
我卻外釋你 說我不敢伸手害我的主.
因為他是耶和華的受告者.
十一節請讀.
我父啊.
我愛你愛我的人在我心中.
我看你的言行滿有的善意.
你有的知我已知道.
我滿有的樂意滿了你.
你隨意脫罪我的命.
我卻不再犯罪於你.
冤耶和華在你我中間判斷是非.
在你身上為我伸冤.
我卻不親手加害於你.
每人有一個女兒喜歡你.
我才喜歡我的人.
我卻不再親手加害於你.
十四節.
以色列王出來要尋找誰呢?.
追趕誰呢?.
不過追趕一條死狗.
一個.
這個字怎麼讀?.
這個讀「吉祖」.
一個吉祖就是了.
好像失禮一樣.
十五節請讀.
冤耶和華在你我中間判斷是非.
大衛向蘇羅說完這話.
蘇羅說:我兒大衛.
這是你的聲音嗎?.
就放聲大哭.
大衛說:你給我貢獻,因為你以善待我,我卻以惡待你.
你今日顯明是以善待我.
因為耶和華將我交在你手裡.

$^{121}$你卻沒有殺我.
你今日顯明是以善待我,我卻以惡待你.
相反蘇羅很清楚地說:以惡報善.
你想做誰?.
你想做大衛還是蘇羅?.
其實報復是很自然的心態.
你質我?質回你.
你罵我?罵回你.
你害我?不敢說.
心裡很想整你.
這就是我們人的血肉,我們的天性.
有一位前蘇聯的領袖對一位新聞記者說.
其實我們共產黨和耶穌基督的教訓有很多相似的地方.
凡物公用,大家都是平等.
但唯一不同的是,如果有人打我右邊的臉.
我就狠狠地打他左邊的臉.
所以我們知道這個事實上是人的天然心態.
事實上是,心裡有種怒氣仇恨的時候.
很容易以牙還牙.
你有半斤,我也有八兩.
但羅馬書說,十二章二十一節.
你不可為惡所聖,反要以善成惡.
意思是說,如果以惡報惡.
就比惡先聖過自己.
因為惡先聖過自己.
但如果我們能夠以善去面對的時候.
我們能夠以善驅走惡.
就好像光驅走黑暗一樣.
內心本來充滿黑暗,像火一樣.
但一道光將它照明,一道水將它滅色.
然後可以在心靈裡得到重新的釋放和光明.
當然,我們知道主耶穌說.
有人打你右邊的臉.
你連左邊都給他打.
這其實是主耶穌強調的,不是一種規條式.
你打我右邊,給你打左邊.
心中會很憤怒,不過不是.
意思是以愛化解衝突.
以恩慈去對待那些得罪我們的人.
最後,耶穌也被人打.

$^{161}$但祂也沒有說你打我左邊.
反過來,祂在十字架上的禱告就是.
主啊,赦免那些得罪你的人.
求你赦免他們,因為他們所作的,他們不了得.
祂為敵人來祈禱,求天父赦免.
祂的心裡無恨.
這就代表以愛勝過恨.
以善去驅走惡.
當然,我們不應該混淆了社會公義.
或者是整體的倫理.
例如有人犯法,所以就應該受制裁.
有人做一些殺人枉法的事,就應該要執行法紀.
有人搗亂整個群體,就應該要攔阻和懲治.
不可以說我們還要多些因待他.
這就要看情況,看事件而定.
也要看個人的倫理.
和社會的公義有時會有衝突的時候.
我們需要有神給我們的智慧.
如何用哪一方面的真理去面對.
也不代表不告訴他對和錯,縱容他犯罪.
大衛也跟訴羅說.
你覺得你這樣待我對不對?.
所以大衛也很勇敢指出訴羅的不是.
也表明神的公義是怎樣應該進行的.
所以我們也不要極端化他,或者誤解了當中的意思.
所以我們明白到被人惡待的感受是不好受的.
我們相信一定在人生當中試過.
有些人無理取鬧,根本就是.
惡言惡語,無理取鬧,根本都不是這樣的,強行來的.
有些人我們會發覺會被人陷害,被人出賣.
為了他的利益,我們做了犧牲品.
我們也經歷過人生當中很多不公平的事.
這個世界沒有公平.
因為在罪惡的世界裡,事實上也有不公平的事發生.
如果我們是受害者,我們心裡一定會非常憤怒.
無端端來受虧損.
又或者有人佔我們便宜.
我們明知道他想騙我們,或者他想佔我們便宜.
所以我們也需要分辨.
不過我們怎樣應對是另一回事.

$^{201}$所以我們明白到我們要拒絕這些惡在我們心裡.
就唯有在神的善裡尋求幫助.
以致我們能夠真正以善來成惡.
怎樣才能以善成惡呢?.
在這段經文裡,我們看到大衛有幾方面值得我們學習.
怎樣以善勝過素羅的惡.
第一,我們可以看到不要聽蠶言.
第四節,第五節,第七節,我們剛才讀過了.
我們知道是跟隨他的人跟大衛說.
如今時候到了,這是個好機會.
他說神為你預備的機會.
所以大衛也心動,真的聽他們說就去前面.
不過他在下手殺保素羅之前.
只割了他的衣襟.
所以後來馬上大衛心中自責.
然後第七節,大衛用這些話攔著跟隨的人.
因為他不殺,後面的人想殺.
大衛反而攔著他們,不用他們起來害素羅.
所以我們看到,其實最初是因為有人建議.
這個很多時候是這樣的.
真的有些人專門去做什麼?.
專門去潑火.
所以你要懂得在面對這些惡的時候.
要留心究竟聽人的聲音還是聽神的聲音.
究竟我要得到人的歡心.
還是我要來順從神的心意.
所以在這個情況下.
周圍的人都說機會來了.
大衛真的動心.
但是卻最後真的有好的靈性.
所以他攔了自己,也攔了別人.
他沒有聽從這些人的推崇.
而做一些不合神心意的決定.
比對,我們看到.
素羅在第九節.
大衛對素羅說,你為何聽信人的慚言.
說大衛想要害你.
所以這個反過來,素羅就聽別人的慚言.
明明大衛很好的.
明明大衛是應該去縱容他的.

$^{241}$明明大衛沒有什麼要害你的.
不過你為什麼要害他呢.
就是聽別人的慚言.
所以我們發覺素羅聽人的慚言做出糊塗的事.
如果大家有時間回去讀第二十二章.
我們發覺他聽了一個人.
這個人叫做多益.
他名字中文翻譯跟他的人完全不同.
這個人應該叫做什麼.
應該叫做武益.
應該叫做多益.
他聽了這個多益,真是血氣之輩.
他竟然殺了大祭司希米勒.
不單止,殺大祭司也有膽量.
然後殺了他全家85人.
全部殺了,還有其他跟他同城的男女.
嘩,真是大錯特錯.
素羅為什麼會這樣.
殺了自己,以色列人自己的大祭司.
還有那些祭司仔,還有他家人,還有整個男女.
就是聽了這個多益,豈有此理.
我也試過有一次,真的有人勸我.
是時候了.
因為這個時候,趁他病又怎樣.
他真的這樣勸我.
那個本來會受害的人還蒙昧無知.
還稱讚他好,那才厲害.
全靠我沒有聽他講.
他今天還保留在這裡.
所以有很多人在旁邊的時候.
會有一些慚言.
所以你要懂得分辨,什麼是正確,什麼是不正確.
記得我們要分辨你聽人的聲音.
要討人的贊成.
還是要聽神的聲音,討神的喜悅.
這是大衛第一樣,很重要.
第二,我們看到敬畏上帝.
這就是他的原因.
第五節,第六節,第七節.
隨後大衛心中自責.

$^{281}$因為割了素羅的衣襟.
然後他說什麼.
我的主乃是耶和華的受告者.
我在耶和華面前不敢伸手害他.
這就是他的原因.
大衛對神敬畏.
所以他懂得不聽人的插言.
去信服上帝的心意.
第十節他再說一次.
他跟素羅都這樣說.
我不敢伸手害我的主.
即是你.
因為他是耶和華的受告者.
這是神的心意.
與剛才那個不同.
剛才是聽插言.
這個是因為敬畏神的心.
去聽上帝的話語.
他不殺素羅的原因.
不是因為素羅不值得死.
真的值得死.
乃是因為他敬畏上帝.
萬不敢伸手害耶和華的受告者.
這個不單止出現一次.
我們不會再說.
在二十六章又出現一次.
所以二十四章 二十六章.
兩次都有機會大衛可以殺素羅.
兩次都沒有.
原因都是這個.
兩次都是這樣說.
中間有一章 遲些我們說.
所以在神學的意義上.
有兩次可以殺不殺.
中間有一次大衛都要殺.
不過不是素羅.
是殺另一個人 那個叫拿伯.
但是又被攔阻.
他又沒有殺他.
所以這三章聖經有神學意義.

$^{321}$今天暫時不交代這件事.
但我們說的是二十四章和二十六章.
我們看到有兩次.
第二次二十六章.
大衛的隊伍在西忽的曠野.
當時素羅帶著三千精兵.
去到這個曠野要尋索大衛.
結果他們在曠野搭營睡了.
剛剛大衛和跟隨他們的人到那裡.
他們都睡著了.
所以他們可以潛入素羅所睡的帳幕裡面.
而這個將領叫阿比西.
阿比西說等我來.
大衛你這麼手軟 等我來.
可以一下子殺了他.
但是大衛再一次攔阻阿比西.
不行 我們不要伸手害他.
因為他是耶和華的受告者.
所以兩次機會.
兩次大衛都是這樣.
又看到素羅的人是這樣.
素羅一次給他機會.
但是他沒有把握這個機會.
又來過.
所以是很抵死的.
但是卻因為神的憐憫.
和大衛對敬畏神的心.
所以大衛為何敬畏神的心.
因為他知道審判不在他的手中.
他交給上帝.
審判的主權交給神.
讓主伸冤最好.
所以第十二節這樣說.
我們留意剛剛我們讀過的.
再讀一次好不好.
十二 十五節.
願耶和華在你我中間判斷是非.
在你身上為我伸冤.
我卻不親手加害於你.
願耶和華在你我中間施行審判.

$^{361}$斷定是非.
並且監察.
為我伸冤.
救我脫離你的手.
所以大衛他敬畏神的心.
就是因為他知道.
上帝作主最好.
他判斷是非.
他會為我伸冤.
他會監察.
救我脫離你的手.
我們知道在舊約的《創世記》.
有一個人物遭人刻薄對待.
陷害的是誰?.
約瑟.
約瑟被他的兄弟陷害.
被賣到埃及做奴隸.
後來又被主母陷害.
叫他無故受冤屈.
落了監.
最後在監獄裡.
他有機會善待了一個賤長.
結果提醒他出監時要平反.
出了去又怎樣?.
沒了這件事.
豈有此理.
我對你這麼好.
你也忘恩負義.
帶約瑟一次又一次.
來不記仇對家人也好.
上司也好.
幫助他的人也好.
但是他沒有一出監.
沒有升為宰相之後.
嘩 事後多了.
豈有此理逐個逐個出來.
全部殺光.
約瑟沒這樣做.
相反他在整個過程裡.
有一個很強調的.

$^{401}$如果留意聖經裡的強調.
就是因為很多次他.
面對試探也好.
面對人也好.
他說因為約瑟敬畏上帝.
所以我們看見.
在這裡讓我們看到.
約瑟他跟他的哥哥這樣說.
他說現在不要因為.
把我賣到這裡自憂自恨.
這是神差我在你們耳先.
來到這裡為要保存成名.
約瑟不但不念舊惡.
還安慰他們.
不用怕.
不要自怨自恨.
這種真是以德報怨.
以善成惡.
所以當後來那些兄弟有些不信.
以為遲早.
阿爸在這裡還好.
阿爸死了.
雅各死了.
那怎麼辦.
這次我們一定死了.
一定約瑟會報仇.
結果當雅各也死了.
約瑟再一次承諾.
叫他們放心.
然後說出原因.
我不會再追究你們.
你們放心.
因為他說了一句話很重要.
就是唱世紀五十章第十九節.
他說:我豈能代替神.
這跟大衛真的很相似.
所以我們看到.
他知道神掌管一切.
就算被惡待也好.
神看見的.

$^{441}$我不用自己血氣的方法.
不用惡報惡.
我將它交給上帝.
他是公義的.
他一定會處置.
所以我們看到.
訴羅反覆無定.
大衛卻是很堅持.
一直都是這樣.
羅馬書十二章也是這樣教導我們.
我們一起讀.
十九節請讀.
親愛的弟兄.
不要自己申冤.
寧可讓報.
聽憑主路.
因為經上記主說.
申冤在我.
我必報應.
生氣是很生氣的.
不過由主那邊.
聽憑主路.
所以如果是不公義的.
是不對的.
是壞的.
主一定發怒.
那我就告訴神.
我不自己申冤.
由你來.
所以敬畏神的表現.
幫助我們以善成惡.
因為知道善在神手中.
惡神別罪徒.
所以我就安然.
交託給這位上帝.
以敬畏神的心.
進行祂的心儀.
是嗎?.
第三.
大衛能夠以善成惡的原因.

$^{481}$是因為他心中有愛.
第十一節.
今日你親眼看見在洞中耶和華將你交在我手裡.
有人叫我殺你.
我卻.
我卻愛惜你.
說我不敢伸手害我的主.
因為他是耶和華的受告者.
然後他稱祂為.
我父.
看看你外婆的衣襟.
所以我們看到.
這裡大衛.
向蘇羅.
表達.
沒有殺祂的原因.
他自己說.
因為他說他愛惜祂.
他當祂是可尊敬的父.
是祂的親人.
事實上心中有愛.
是一個最大最大能夠以善成惡的力量.
越有耶穌基督的愛.
就越不容易以惡還惡.
大衛雖然看見蘇羅的卑劣.
心胸的狹窄.
事實上反覆無定.
但他卻紀念祂是.
不單止是耶和華的受告者.
也是以色列的君王.
所以他愛國.
也愛王.
祂是他軍中的主.
所以他開始來到祂面前的時候.
他稱祂為我的主.
我的王.
這是最初祂的稱呼.
所以他稱祂為我的主.
因為祂是軍中的主.
稱祂為我的王.

$^{521}$因為祂是以色列的王.
但其實他們有一個關係.
大家記得嗎.
蘇羅是大衛的外父.
所以他稱祂為我父.
這種就表示那種親情.
愛勝過了恨.
所以不單止稱祂為我主.
我王.
更加稱祂為我的父.
怎樣能夠以善勝惡呢.
其實這個很重要.
就是去想祂的好方面.
不要只想祂的壞方面.
想祂的好方面的意思是.
去想想怎樣用神的愛.
看見祂好的地方.
可以尊敬和尊重的地方.
因為祂是耶和華的受告者.
以色列的君.
是我軍中的主.
但更多一點就是.
我和祂有親.
所以很多時候我們罵人都是這樣.
我和你什麼.
我和你有親.
弟兄姊妹就好一點.
有沒有親.
家人呢 有沒有親.
又容易一點.
如果想想其實我們也有親.
有親密的時候.
所以有些人說.
有關係就沒有關係.
沒有關係.
就有關係.
所以找一些關係.
關係就是建立在主的裡面.
所以你還沒有做基督徒.
真的要快點做基督徒.

$^{561}$因為你夫婦經常吵架.
子女經常不和.
朋友經常衝突.
見到這個不對那個又很生氣.
信了耶穌就不同了.
信了耶穌.
你的配偶就看到好一點的地方.
如果大家都信耶穌更好.
看見其他人就想到主耶穌都愛他.
所以信主真的很重要.
耶穌都這樣教導我們.
《馬太福音》五章四十四節.
這裡是這樣說的.
我們一起讀一讀.
「只是我告訴你們要愛你們的仇敵.
為那逼迫你們的禱告」.
請問重點在哪裡.
愛.
重點在愛.
沒有愛就為他祈禱.
所以我們見到主耶穌基督.
面對這些逼迫他的人.
他能夠寬恕他們.
是因為有愛.
中國大陸曾經有件事賣上報紙.
就是一個叫做沙河口.
有沒有人去過.
沙河口市.
有一家叫做姓張家.
有一個感人故事.
還沒有賣上報紙.
是因為七年之前.
有一位姓李的女士.
因為丈夫死了.
第二次結婚嫁給姓張.
這個姓李的女士帶著她的十五歲女兒.
一起嫁給姓張.
誰知嫁了後才後悔.
因為姓張是酒徒.
後來一直喝酒.

$^{601}$整天都喝醉.
可見家庭非常坎坷.
生活很困難.
整天都喝醉酒之後.
又怎樣.
她沒有說打女人.
只是說發脾氣.
你們這麼厲害.
姓李就怕她搬出去躲一躲.
帶上女兒.
但她又不帶.
自己走出去.
只剩下女兒在.
慘了.
所以姓張在發脾氣的時候.
比打女人還要厲害.
就用煙頭.
燙她.
壓她.
整她.
然後用滾水潑她.
很殘忍.
在一年.
十五歲到十六歲這兩年.
這個女子的身上有疤痕.
從煙頭燙到的有十幾個.
水燙傷的有四五個.
她跟同學說.
這個很痛苦.
我永遠不會忘記.
我繼父所做的壞事.
不過很奇妙.
後來神有公義.
姓張在過馬路的時候.
忽然被車撞到.
撞到之後兩隻腳.
都有不同程度的殘廢.
所以她沒有辦法照顧自己.
然後就令太太姓李出去打工.
沒時間照顧她.

$^{641}$只剩下這個女兒.
當時她已經二十一歲.
這個女兒見到繼父這麼可憐.
心裡面就起了憐憫之心.
這就是愛.
所以她就服侍她.
幫她洗臉洗腳.
照料她起居.
這個養女沒有恨.
反而想到繼父有恩.
所以她就願意以善勝惡.
結果當那些政府.
社會工作人員來探望她的時候.
探望家人的時候.
這個姓張的爸爸.
就在這些社會工作人員面前.
來讚揚他這個女兒的孝心和善心.
甚至跪在他女兒面前求他饒恕她.
所以就賣相報紙被刊登出來.
當然今天如果我們發現有這樣的事.
就應該要立即報警.
如果有虐待兒童的事.
但是這個例子給我們看見.
就是一種因為愛連續能夠勝過仇恨和惡的.
很具體的事例.
愛的力量很大.
如果普通人的愛都是這麼大.
何況在主的愛裡面不是這麼大嗎?.
我們所唱的詩歌.
我們體會到耶穌基督被懸掛在木頭上.
祂為我們擔當祂的罪.
為我們犧牲祂的生命.
本來我們應該受罪帶來的刑罰.
但是主為我們擔當.
所以這種大愛.
幫助我們想到主的愛的時候.
我們所受的都不會算得上什麼.
體會主的連續.
心裡面就有連續.
體會主的愛.

$^{681}$心裡面就充滿到神的愛.
就可以用愛成為力量.
化解苦毒,憤怒,仇恨.
所以最後我們看見.
為什麼這麼重要有主的愛.
就是我們需要謙卑,仗著.
大衛有一個謙卑的心.
這個我相信是幫助他以善勝惡.
一個很重要的其中一個原因.
十四節這樣說.
「以色列王出來要尋找誰呢?」.
他追趕誰呢?.
「不過追趕一條死狗」.
即是誰呢?.
他自己在說.
「為什麼你要追得這麼辛苦呢?」.
「你為什麼要為一條好像我這些」.
「沒什麼價值的死狗窮追猛打?」.
「死狗都不夠」.
還要說什麼呢?.
「一個」.
怎麼讀呢?.
又忘記了.
「一個吉祖」.
一個吉祖呢?.
「死狗都大大隻」.
「吉祖」.
「為什麼你要失腦呢?」.
「這麼辛苦來追殺你」.
「我有什麼價值呢?」.
「為什麼你要追殺我呢?」.
這顯出了大衛在訴羅面前的謙卑.
所以他謙卑.
十五節說.
「願耶和華在你我中間施行審判」.
「斷定是非」.
「並且監察」.
「為我伸冤」.
「救我脫離你的手」.
所以人何時不能消除憤恨報復呢?.

$^{721}$因為他覺得自己很重要.
這很明顯.
他看自己太重要的時候.
自尊心就令他要報復.
有人少少得罪我也好.
看得來很大.
因為他覺得自己很重要.
「叫我失威豈有此理」.
「叫我難堪」.
「我跟你死過」.
這就是他的心態.
但相反當人自覺謙卑.
在神面前.
其實不太重要.
他就容易寬恕.
沒有用惡去還擊.
這不是說我們輕薦自己的看法.
而是大衛將他的價值建立在耶和華內.
他認為耶和華會為他伸冤.
耶和華愛我.
耶和華會為我秉行公義.
因為我將這事交給上帝.
我屬於祂.
這是一個很健康對自己的看法.
就是他知道我其實無憂.
不過在神眼中很重要.
所以他在神眼中很重要的時候.
他就學校主耶穌自己.
今天我們學校主耶穌.
就好像我們不會用血氣.
用鬥爭來使祂死去活來.
因為要看自己這麼重要.
反而我們將重點放在神身上.
聖經今天讓我們知道.
我們作為基督徒真的很蒙恩.
為什麼呢?.
在加拉太書第二章第二十節.
我們一起讀一讀這節.
「我已經與基督同釘十字架.
現在活著的不再是我了.

$^{761}$乃是基督在我裡面活著.
並且我如今在肉身活著.
是因順上帝的兒子而活.
他是愛我.
為我寫記」.
當他明白耶穌的愛.
十字架的愛的時候.
祂為我們寫記.
我們就與主同釘十字架.
將這個舊我釘死.
不再是我.
但是這個我就是在基督耶穌裡面活著.
上帝愛我.
這就是我們基督徒對自己的看法.
所以我們不需要聽人的讒言.
我們敬畏上帝.
我們也不去以惡還惡.
因為我們有神的愛.
更加不看我們自己這麼大.
我們知道我們已經與主同釘十字架.
現在活著的不再是我了.
所以我們看到這就是大衛的喜方面.
從聖經裡面看見他能夠面對騷擾.
但是卻以善勝惡的原因.
中世紀有一次.
日耳曼的勳爵希利他班.
與另一個勳爵有很大的衝突.
非常憤怒.
當時很容易決鬥.
決鬥不是你死就是我死.
所以預備決鬥之前.
他就想起不知道這次結果如何.
可能殺了他.
又可能是我自己死.
所以他就說不如先去教堂.
所以他就去教堂.
誰知道進入教堂看見三幅畫.
第一幅畫是什麼.
第一幅就是耶穌在希律王面前.
被人戲弄的畫面.

$^{801}$在畫下寫著他被罵不還口.
第二幅就是兵丁鞭打耶穌時的畫面.
在畫下寫著受害不說威嚇的話.
這是彼得前書第二章的說話.
他的心一直很感動.
看見耶穌被罵不還口.
受害不說威嚇的話.
看見第三幅更感動.
第三幅就是耶穌釘在十字架的形象.
下面有一句說話寫著.
父啊赦免他們.
因為他們所作的他們不孝道.
當時他感動到淚流滿面.
跪在十字架的面前.
求主赦免他.
神的愛充滿他的心.
他馬上去到敵人那裡.
伸出雙手與他和解.
不但沒有決鬥.
還告訴他這個見證.
就是神怎樣感動他.
所以後來竟然這兩個人從那天開始.
成為了最好的知己.
真是難得.
我相信大衛的心性不斷培養.
可能在曠野的時候更加親近神.
看見乾旱的沙漠.
卻有忍饑底的綠洲.
成為他讚美神.
主光重露的時候.
看見神創造的美麗.
看見在乾旱之地.
疲乏之地有神的傳源.
以致他能夠再一次敬畏上帝.
再一次生出愛心.
再一次體會到怎樣.
以德報怨.
以善成惡.
奧古斯丁這樣說.
奧古斯丁是初期的教父.

$^{841}$一位屬靈的偉人.
他說以惡報善.
就像蘇羅一樣.
就是屬魔鬼.
人家好好的對你.
那麼多次給你機會.
不但不珍惜.
還不斷以惡報善.
是不是屬魔鬼呢.
但是他說以善報善.
就是屬於人.
人也做到.
以惡報惡.
以善報善.
你對我好.
那我當然要對你好.
還有一句.
你猜一下.
一定猜得到.
猜不猜得到.
是.
以善報惡.
就是屬於上帝.
頂子妹.
不知道你心裡有沒有憤恨.
或者只是在一些很不滿的地方.
或者有些人對你好不好.
你選擇.
屬於哪樣.
我們低頭祈禱.
認同神道.
為我們受難.
受死.
帶卻復活得勝.
賜給我們新的生命.
我們看到主耶穌基督的榜樣.
祂以愛感動我們.
讓我們來學習祂.
以愛代替恨.
以德報答怨.

$^{881}$以善勝過惡.
讓仇恨和黑暗驅走.
在曠野乾涸的地方.
見到引機底的綠洲.
我們知道.
這就是你要給我們得到的.
永恆的賞賜.
內心的平安.
和作為基督徒的福分.
多謝你.
聽我們的祈禱.
奉主耶穌基督的名.
\newpage



\section{出埃及記 14:13}
\label{sec:031977}
\textbf{聖經人聖經事:看耶和華施行的救恩}~~~~黃克勤牧師博士
\newline
\newline
連結: \href{http://www.church.com.hk/acms/content.asp?site=cdc\&op=show\&type=product\&code=031977\&layout=sermon}{\texttt{031977}} ~~~~ 語音日期: 2026-04-11
\newline
\newline
\hyperref[sec:031976]{\small{< < < PREV SERMON < < <}}
~
\hyperref[sec:index_chronic]{\small{[返順時目]}}
~
\hyperref[sec:031978]{\small{> > > NEXT SERMON > > >}}
\newline
\newline
出埃及記 14:13
\newline
\begin{longtable}{cl}
\hline
\hline
章節 & 經文 (和合本修訂版)\\
\hline
14:13 & \begin{tabularx}{0.7\textwidth}{X} 摩西對百姓說：「不要怕，要站穩，看耶和華今天向你們所要施行的拯救，因為你們今天所看見的埃及人必永遠不再看見了。 \end{tabularx} \\ \\
[1ex]
\hline
\hline
\end{longtable}
$^{1}$各位我們又一起讀聖經.
《生命書》書架網絡.
和大家一起看聖經人.
聖經事背後的上帝.
他在做什麼.
今次仍然是講摩西的故事.
他的故事實在很長.
現在來到他人生八十歲以後.
大以色列人出埃及.
去到曠野.
過紅海.
今天是看耶和華施行的救恩.
看耶和華施行的救恩.
是今天和大家一起看的故事.
《出埃及記》第十四章第十節.
法老臨近的時候.
以色列人舉目.
看見埃及人趕來.
就甚懼怕.
就向耶和華哀求.
讀多一次.
為什麼要讀多一次呢.
慢慢讀多一次.
會令到我們對於經文的感受強一些.
如果你發覺讀聖經有時不太入味.
你要自己慢讀.
而且是連續重複讀它幾次.
你就會進入到那個場景裡.
多很多的.
法老臨近的時候.
以色列人舉目.
看見埃及人趕來.
就甚懼怕.
向耶和華哀求.
他們那時候就去到海邊對著對面海的地方.
叫巴力西芬.
他們就在這裡對著對面海.
這個時間.
法老臨近.
他看見埃及的追兵.

$^{41}$嘩.
有很多塵.
什麼事呢.
快點派人去看一看.
一看.
嘩.
很遠就看到知道是法老,埃及的兵.
心裡面就知道.
死了.
來追我們了.
以色列人舉目.
看見埃及人趕來.
自然是怎樣呢.
就甚懼怕.
死了.
為什麼呢.
他究竟來追殺我們.
還是要抓我們回去呢.
無論追殺或者抓我們回去.
都很害怕.
這個時間就甚懼怕.
所以看見一些事.
就甚懼怕.
這是第一點要講的.
是一點都不出奇的.
因為你看到.
看到你就害怕.
其實看到之後.
害怕也有的.
其實聽到更害怕.
譬如.
我沒有看到.
不過我派探子去.
然後回來告訴我.
嘩.
不得了.
很多埃及人數不到那麼多.
戰車很多.
嘩.
很闊.

$^{81}$隆隆聲.
追著來.
多大煙.
煙塵就知道很多兵馬.
我們這次沒有了.
唉.
慘了這次.
我們一定來不及了.
你聽到.
有時候驚恐還大過見.
看見就甚懼怕是正常的.
因此.
咦.
你基督徒也害怕嗎?.
為什麼基督徒會不害怕?.
我是一個人.
見到危險的事.
當然害怕.
害怕得震得不懂得走.
很害怕的人.
是害怕到震驚.
被震一震.
震撼到.
我真的很震撼.
這樣說而已.
真的震撼是怎樣呢?.
真的震撼是人呆了.
雙腳不動.
嘩.
突然之間爆炸之後.
離了體.
或者有些事很厲害.
家人忽然遭遇意外不在.
傷痛欲絕.
那一刻接受不了.
接著人呆了.
震撼到不能走路.
跌倒.
那些真是害怕了.
這裡就甚懼怕.

$^{121}$死了死了怎麼辦?.
由於不知怎麼做.
向耶和華哀求.
哀求.
向耶和華哀求.
好對,非常好.
向耶和華哀求.
這個以色列人遭遇這樣的問題.
前無去路.
後有追兵.
很害怕.
背水一戰.
但其實他是專業軍隊.
我們只不過是平民百姓.
我也有刀.
不過不夠他打.
戰車馬兵.
盔甲整齊.
向耶和華哀求.
對了.
但是.
他們有一件事不好.
同時這個懼怕.
令他們除了神救我們.
哀求之外.
他們對摩西法怨言.
難道在埃及沒有墳地嗎?.
說話很尖酸刻薄.
難道在埃及沒有墳場嗎?.
你要將我們帶來死在曠野嗎?.
為什麼你這樣對待我們?.
為什麼你要將我們從埃及領出來?.
自己一點責任都沒有.
沒有人鎖住你走.
難道埃及沒有墳地嗎?.
為什麼你這樣對待我們?.
為什麼你要將我們從埃及領出來?.
為什麼要將我們從埃及領出來?.
請不吝點贊訂閱轉發打賞支持明鏡與點點欄目.
\newpage



\section{民數記 9:1-23}
\label{sec:031978}
\textbf{生命更新•新旅程}~~~~鄭玉兒傳道
\newline
\newline
連結: \href{http://www.church.com.hk/acms/content.asp?site=cdc\&op=show\&type=product\&code=031978\&layout=sermon}{\texttt{031978}} ~~~~ 語音日期: 2026-04-11
\newline
\newline
\hyperref[sec:031977]{\small{< < < PREV SERMON < < <}}
~
\hyperref[sec:index_chronic]{\small{[返順時目]}}
~
\hyperref[sec:031979]{\small{> > > NEXT SERMON > > >}}
\newline
\newline
民數記 9:1-23
\newline
\begin{longtable}{cl}
\hline
\hline
章節 & 經文 (和合本修訂版)\\
\hline
9:1 & \begin{tabularx}{0.7\textwidth}{X} 以色列人出埃及地以後，第二年正月，耶和華在西奈的曠野吩咐摩西說： \end{tabularx} \\ \\ \relax
9:2 & \begin{tabularx}{0.7\textwidth}{X} 「以色列人應當在所定的日期守逾越節。 \end{tabularx} \\ \\ \relax
9:3 & \begin{tabularx}{0.7\textwidth}{X} 你們要在本月十四日黃昏的時候，在所定的日期守這節，按照一切的律例典章守節。」 \end{tabularx} \\ \\ \relax
9:4 & \begin{tabularx}{0.7\textwidth}{X} 於是摩西吩咐以色列人守逾越節。 \end{tabularx} \\ \\ \relax
9:5 & \begin{tabularx}{0.7\textwidth}{X} 正月十四日黃昏的時候，他們就在西奈的曠野守逾越節。凡耶和華所吩咐摩西的，以色列人都照樣做了。 \end{tabularx} \\ \\ \relax
9:6 & \begin{tabularx}{0.7\textwidth}{X} 有幾個人因屍體成了不潔淨，不能在那日守逾越節。當天他們到摩西、亞倫面前。 \end{tabularx} \\ \\ \relax
9:7 & \begin{tabularx}{0.7\textwidth}{X} 那些人對他說：「我們因屍體而不潔淨，為何禁止我們，不能和以色列人在所定的日期獻供物給耶和華呢？」 \end{tabularx} \\ \\ \relax
9:8 & \begin{tabularx}{0.7\textwidth}{X} 摩西對他們說：「你們稍等，讓我去聽耶和華對你們有甚麼吩咐。」 \end{tabularx} \\ \\ \relax
9:9 & \begin{tabularx}{0.7\textwidth}{X} 耶和華吩咐摩西說： \end{tabularx} \\ \\ \relax
9:10 & \begin{tabularx}{0.7\textwidth}{X} 「你要吩咐以色列人說：你們和你們後代中，若有人因屍體成了不潔淨，或出外遠行，仍然要向耶和華守逾越節， \end{tabularx} \\ \\ \relax
9:11 & \begin{tabularx}{0.7\textwidth}{X} 他們就要在二月十四日黃昏的時候守節，要吃羔羊，以及無酵餅和苦菜。 \end{tabularx} \\ \\ \relax
9:12 & \begin{tabularx}{0.7\textwidth}{X} 他們不可留一點食物到早晨；羔羊的骨頭一根也不可折斷。他們要照逾越節的一切律例守這節。 \end{tabularx} \\ \\ \relax
9:13 & \begin{tabularx}{0.7\textwidth}{X} 但潔淨又不出外遠行的人若不守逾越節，那人要從百姓中剪除，因為他沒有在所定的日期獻供物給耶和華，必須擔當自己的罪。 \end{tabularx} \\ \\ \relax
9:14 & \begin{tabularx}{0.7\textwidth}{X} 若有寄居在你們那裡的外人要向耶和華守逾越節，他要照逾越節的律例和典章做。無論是寄居的或是本地人，都用同樣的律例。」 \end{tabularx} \\ \\ \relax
9:15 & \begin{tabularx}{0.7\textwidth}{X} 立起帳幕的那日，有雲彩遮蓋帳幕，就是法櫃的帳幕；從晚上到早晨，雲彩在帳幕上，形狀如火。 \end{tabularx} \\ \\ \relax
9:16 & \begin{tabularx}{0.7\textwidth}{X} 經常都是這樣：雲彩遮蓋帳幕，夜間雲彩形狀如火。 \end{tabularx} \\ \\ \relax
9:17 & \begin{tabularx}{0.7\textwidth}{X} 雲彩幾時從帳幕上升，以色列人就幾時起行；雲彩在哪裡停住，以色列人就在哪裡安營。 \end{tabularx} \\ \\ \relax
9:18 & \begin{tabularx}{0.7\textwidth}{X} 以色列人遵照耶和華的指示起行，也遵照耶和華的指示安營。雲彩在帳幕上停留多久，他們就留在營裡多久。 \end{tabularx} \\ \\ \relax
9:19 & \begin{tabularx}{0.7\textwidth}{X} 雲彩在帳幕上停留許多日子，以色列人就遵照耶和華的吩咐不起行。 \end{tabularx} \\ \\ \relax
9:20 & \begin{tabularx}{0.7\textwidth}{X} 有時雲彩在帳幕上只停了幾天，他們就遵照耶和華的指示留在營裡，也遵照耶和華的指示起行。 \end{tabularx} \\ \\ \relax
9:21 & \begin{tabularx}{0.7\textwidth}{X} 有時雲彩從晚上留到早晨；早晨雲彩上升，他們就起行。無論是白天是黑夜，當雲彩上升的時候，他們就要起行。 \end{tabularx} \\ \\ \relax
9:22 & \begin{tabularx}{0.7\textwidth}{X} 雲彩停留在帳幕上，無論是兩天，一個月，或更長的日子，以色列人就留在營裡不起行；但雲彩一上升，他們就起行。 \end{tabularx} \\ \\ \relax
9:23 & \begin{tabularx}{0.7\textwidth}{X} 他們遵照耶和華的指示安營，也遵照耶和華的指示起行。他們遵守耶和華的吩咐，是耶和華藉摩西所指示的話。 \end{tabularx} \\ \\
[1ex]
\hline
\hline
\end{longtable}
$^{1}$各位鼎姐妹晚安.
今日我會同大家繼續去講這個門訓.
就是這個文掃記的九章.
整章有二十三節這麼多.
因為這段經文其實都是很重要.
我們很想大家一齊去重讀這二十三節.
以致我們在去理解這章經文的時候.
有一個概念.
我們一起重讀 預備起.
以色列人守逾越節.
他們就在西內的曠野.
正有黃昏的時候守逾越節.
凡耶和華所吩咐摩西的.
以色列人都照樣行了.
有幾個人因死屍而不潔淨.
不能在那日守逾越節.
當日他們在摩西亞倫面前說.
我們雖因死屍而不潔淨.
為何被阻止不得同以色列人.
在所定的日期獻耶和華的供物呢?.
摩西對他們說.
你們暫且等候.
我可以去聽耶和華指著你們是怎樣吩咐的.
耶和華對摩西說.
理曉於以色列說.
你們和你們後代中.
若有人因死屍而不潔淨.
或在遠方行路.
還要向耶和華守逾越節.
他們要在二月十四日黃昏的時候.
守逾越節.
要用無餃餅和苦菜和逾越節的羊羔同吃.
一點不可留到早晨.
羊的一斤也不可截斷.
他們要照逾越節的一切律例而守.
那潔淨而不行路的人.
若推遲不守逾越節.
那人要從民中剪除.
因為他在所定的日期.
不獻耶和華的供物.

$^{41}$應該擔當他的罪.
若有外人寄居在你們中間.
願意向耶和華守逾越節.
他要照逾越節的律例典章行.
不管是寄居的.
是本地的.
同歸一例.
立起帳篷的日.
有雲彩遮蓋帳幕.
就是發跪的帳幕.
從晚上到早晨.
雲彩在旗上.
形狀如火.
像是這樣.
雲彩遮蓋帳幕.
夜間形狀如火.
雲彩何時從帳幕收上去.
以色列人就何時起行.
雲彩在那裡停住.
以色列人就在那裡安營.
以色列人遵耶和華的吩咐而行.
也遵耶和華的吩咐安營.
若在帳幕上停住何時.
他們就住營何時.
雲彩在帳幕上停留許多日子.
以色列人就守耶和華所吩咐的不起行.
有時雲彩在帳幕上幾天.
他們就照耶和華的吩咐住營.
也照耶和華的吩咐起行.
有時從晚上到早晨.
有者雲彩在帳幕上.
早晨雲彩收上去.
他們就起行.
有時晝夜雲彩停在帳幕上.
收上去的時候.
他們就起行.
雲彩停留在帳幕上.
無論是兩天是一月是一年.
以色列人就住營不起行.
但雲彩收上去.

$^{81}$他們就起行.
他們遵耶和華的吩咐安營.
也遵耶和華的吩咐起行.
他們守耶和華所吩咐的.
都是憑耶和華吩咐摩西的.
九章完了.
不過想大家讀完後面那幾字.
好,一起讀,預備起.
第二年二月二十日.
雲彩從法桂的帳幕收上去.
以色列人就按站往前行.
離開西奈的曠野.
雲彩停住在巴蘭的曠野.
是他們照耶和華.
摩西所吩咐的.
初次和前行.
這裡是十章的十一到十三節.
好了,九章是我們今天的重點.
為什麼說完第十章的這幾節呢.
因為很想給大家一個.
很重要的Picture.
他們首先怎樣去.
為什麼會在西奈的曠野那裡.
去停留.
這個表是我理解完之後.
去和大家做了一個這樣的表.
你見到.
以色列人其實.
停在西奈的曠野這個地方.
其實就是在他出埃及.
滿了三個月的那一天.
就在西奈的曠野那裡安營.
這個是出埃及記的.
十九章一節那裡.
可以和我們說的.
什麼時候他們會離開曠野.
就是我們剛才讀到的民數記.
第十章那裡說.
他就在第二年的二月二十日.
算一算時間.

$^{121}$中間其實有十一個月的時間.
以色列人停留在西奈曠野的地方.
我們看一看.
在西奈曠野有什麼重要的標誌.
我們是要知道的呢.
第一就是.
在第二年.
通常在第二年發生的.
你見到.
在民數記有日期去標記的.
就是這些日子.
第二年的正月初一.
在民數記第七章.
或者是出埃及記的四十章.
都去記載.
他們是立起那個帳幕的.
而立起帳幕就代表著.
帳幕裡面就有神的法鬼在當中.
後面有讀到.
有神的雲彩去遮住帳幕.
第二年的正月十四.
就是剛才我一開始讀的九章的時候.
那個日期.
耶和華就和摩西說.
要以色列人去守逾越節.
在正月十四.
第二年的正月十四.
去守逾越節.
中間我們剛才讀到.
有一些不潔的人.
他們去問.
為什麼我們不守逾越節呢.
摩西就說.
我問一下神.
就得到答案在後面.
就是第二年的二月十四號.
他們都要去守逾越節.
中間當然要做一些潔淨的禮儀.
他們要潔淨了自己.
才可以繼續守逾越節.

$^{161}$我們又看到.
神給了一個月的時間.
那些不潔的人.
去做預備潔淨的工作.
中間做些什麼呢.
原來在民數記的一章的時候.
一開始的時候.
就告訴我們.
第二年的二月初一的時候.
是做些什麼.
數點意識裂人的將士和那些人數.
我們如果看回這個章節的時候.
我們看到.
原來我們當逾越節.
應該潔淨了的人.
在正月做完.
正月十四號做完.
隔了一個月.
那些不潔的人.
都做完潔淨禮的時候.
過了如節的時候.
第二年的二月二十號.
他們就離開西來的曠野.
我們又會知道.
原來這個是一個重要的標記.
原來逾越節過後.
意識裂人要展開一個新的旅程.
暫新的旅程.
如果我們看回第十章.
剛才讀到的時候.
那個雲彩停留在巴蘭的曠野.
他就停留在巴蘭.
第二站就去到巴蘭的曠野.
我們又看看.
原來在整個裡面.
之前那幾章講什麼.
其實是一個beefing.
他們起行之前的一個神給他們的beefing.
如果你還有印象的時候.
民數記二章講什麼.

$^{201}$就是講十二支派怎樣去安營.
安營在圍著帳幕周圍去安營.
第三第四章和第八章.
就講到利未人的特別的角色.
和他們的職責.
第五第六章就講那些祭司的責任.
和他們的權利.
和他們怎樣獻祭.
而第七章剛才就講到.
那個帳幕就被立起.
摩西在帳幕裡面分別圍聖.
然後就去獻祭.
而第八章剛才講的就是利未人.
然後就是第九章.
就是神要以色人去過愉悅節.
在過愉悅節的時候.
就是一個預備起行一個很重要的標記.
所以之前所做的一切功夫.
都是一個beefing.
我們有了這些背景的時候.
我們要再看.
看到原來在第二年的正月的時候.
神叫以色列人要去過愉悅節.
在第四節的時候.
摩西就吩咐那班以色列人.
你們要守節.
愉悅節的背景.
這次是他們出埃及第二次去守愉悅節.
第一次是甚麼時候守.
也是在講起源在哪裡.
當然是在講出埃及的時候.
是第十災.
第十災的時候.
神就去告訴摩西和阿倫.
在初一及十二章就開始記載.
愉悅節的起源.
我們重溫一下.
第一次的愉悅節.
為甚麼神要這樣去設立.
原來祂是帶領以色列人去離開埃及.

$^{241}$那個也是在展開一個新的旅程.
一班以色列人在埃及維勞茲地.
在430年的時間在維勞茲地去過.
他們要展開新的旅程的時候.
神就設立了愉悅節.
他們說你們要在本月正月.
成為你們一年之首.
在初十的時候.
你們要預備一隻羊羔.
每家一隻.
然後在正月十四號黃昏的時候.
你們就要將羊羔烤了.
殺了.
然後將血塗在門框和門尾上.
然後第八節就說到.
過節的時候要預備的.
有愉悅節的羔羊.
要用火烤.
其實在當時來說是烤全羊.
甚麼叫烤全羊.
大家去旅行都會知道.
很美味.
但當時烤全羊不是這樣的.
是一個憲制.
那些羔羊的血其實是代表著一個袋贖的意思.
神藉著羔羊的血.
越過那些在屋裡面的人.
亦都在屋裡面是一個避難所.
是一個拯救.
在屋裡面的人避過被殺長子的災難.
無教餅就是象徵著出埃及那些以色列人要去除罪.
而苦菜其實是神要他們去紀念.
他們在埃及裡面圍爐的那些受苦的日子.
最後的highlight我們就要記住.
他說這個是耶和華的愉悅節.
愉悅節就是這樣來的.
愉悅節就是這樣去設立的.
後面的內容就告訴我們.
原來神當晚就會巡行埃及遍地.
然後會擊殺所有投山的.

$^{281}$無論是人或者是牲畜.
後面又有一個很重要的字眼.
又要敗壞埃及一切的神.
我是耶和華.
這個是耶和華神自己的宣告.
他不單是藉著殺長子這個十災的懲罰.
或者這個災去拯救那些以色列人.
更加是要彰顯神的大能.
神說他們要去敗壞埃及裡面一切的神.
我們有沒有記住.
原來在中埃及的三章四節那裡.
摩西問你是誰的時候.
神怎樣去介紹他自己.
那裡就說 I am who I am.
而和合本就說我是自有永有.
而最新的新漢語譯本就告訴我們.
直譯那句說話就是.
我是就是我是.
就是說我就是我這樣的意思.
耶和華告訴他們.
我就是我.
我就是那位大能的神.
全地都是我的.
祂是那位世界上唯一的神.
祂要將埃及所拜的假神全部毀滅.
所以這個第十三節告訴我們.
原來這個十災裡面.
是告訴我們.
神要拯救一班以色列人.
他們要按照神的吩咐.
將神羔羊的血畫在門眉上.
以致這個災亡不會臨到他們.
甚至在這個十災之後.
將以色列人帶領離開埃及.
這個民數記八章那裡.
剛才我提過.
原來利未人的設立.
是一個很特別的角色.
利未人就是神所特別揀選.
特地去歸祂.

$^{321}$並且在第十八節告訴我們.
我揀選利未人代替以色列人中一切投生的.
祂告訴我們原來利未人的角色.
就是祂要被神去使用.
讓一班長子可以避屬出來.
所以利未人在以色列人.
他們沒有地分.
因為他們要在世世代代.
都在聖殿裡面去幫助.
去協助祭司去做一些獻祭的工作.
特別在儒節那裡.
都需要他們有特別的角色在當中.
我們知道原來神怎樣去設立這個儒節.
儒節怎樣對以色列人有這麼重要的標記.
以至他們出埃及的時候.
他要重新起行離開西奈的曠野的時候.
神再一次去告訴摩西和阿倫.
這班以色列人.
我要帶領你們去展開一個新的旅程的時候.
你們要守儒節.
再一次紀念我怎樣在埃及維羅斯尼.
去帶領你們出來.
所以在民數記裡面.
一個很特別的.
我覺得是一個重點.
因為不是所有的以色列人.
當時的以色列人.
都能夠可以在正月十四.
第二年的正月十四號的黃昏的時期.
可以守儒節.
就是有一班因為屍體不潔症.
即是不潔的人.
他們不能夠去守這個儒節.
他們覺得很無辜.
因為我們不守儒節.
我們的身份是怎樣呢.
我們可不可以受到神的保護.
我們可不可以跟隨大隊.
去離開西奈的曠野呢.
有很多問題.

$^{361}$其實對於摩西和阿倫來說.
其實他們都爭拗了.
因為他們在出埃及之前.
才是第一次過儒節.
出埃及的第三個月.
去到西奈的曠野的時候.
他們基本上是第二次去過儒節.
很多的禮儀.
他們都不知道.
不潔的人.
他們是怎樣去守儒節的呢.
然後摩西就說.
我先幫你問問神.
第八節講得很清楚.
但是我們知道一件事.
民數記的五章.
其實是提過那些.
摸過屍體不潔的人.
是怎樣處理的.
摸過屍體不潔的人.
是要怎樣呢.
是要離開營外的.
因為祂說.
這個營裡面所住的是我.
我是住在耶和華.
耶和華和摩西說.
在營裡面是我所住的.
是聖潔的.
那些不潔的人要在營外.
這個講了怎樣處理他們.
講了原來在營裡面.
是有神分別為聖的意義在當中.
而民數記.
去到民數記第十九章的時候.
我們才看到摸過屍體不潔的人的處理.
這個應該是再記述.
不潔的人要怎樣去潔淨自己.
如果我們看民數記.
看長一點的時候.
原來摸過屍體不潔的人.

$^{401}$怎樣去處理.
怎樣潔淨自己.
是要用七天的時間去處理.
要找一些灰.
要一些隨序的灰.
沒有很詳細的解釋.
總之有些灰加一些活水.
然後第三天和第七天.
要找一個潔淨的人.
去幫那些不潔的人去灑.
然後到第七天的時候.
他們就會成為聖潔.
不潔的人就會被潔淨.
而那些人就要去洗衣服.
就要洗澡.
令到自己去潔淨.
用七天的時間去潔淨自己.
耶和華怎樣去回應摩西呢.
因為剛才摩西說要問神.
你們這些不潔淨的人.
怎樣過逾越節.
神就去回應摩西.
祂就說你們和你們的後代.
這個以下的處理.
其實就講明這個例子.
你們以後的人都會這樣做.
所以這裡講到.
不單止是因為死屍不潔淨.
去講那些遠方走路的人.
什麼叫遠方.
簡單來說就是出遠門的人.
出遠門的人.
他們都要守逾越節.
你潔淨了.
潔淨自己.
或者你按時回來的時候.
再潔淨自己.
你就可以繼續守逾越節.
所以剛才我給你們看一個chart.
其實這裡都講到.

$^{441}$他們要在下一個月.
第二月十四號黃昏的時候.
他們那些不潔的人.
或者在遠方回來的時候.
他們要再守.
我要看看那些走遠方出遠門的人.
為什麼.
因為剛才講到.
營裡面是聖潔的.
營外面是不潔的.
所以那些出了營外面的人.
是屬於不潔的.
所以那些不潔的人.
回來營裡面的時候.
要守節的時候.
他們要再潔淨自己.
神就告訴你.
你用一個月的時間.
無論你出遠門出多久.
你自己按時間.
回來營裡面去守逾越節.
其實意義就是這樣.
一樣.
要吃苦菜.
要吃無教餅.
和吃羔羊.
一起吃.
同一樣.
和剛才十二節.
猜及記十二章那一套.
一點都不可以留到第二天.
羊羔羊的骨頭.
一斤都不可以截斷.
這個就是如節的定例.
一會兒再解釋這個Point.
要講到第十三節裡面.
他告訴我們.
那班屬於潔淨了.
屬於可以去守如節那班.
又不用出遠門那班的人.

$^{481}$又推遲.
推遲是什麼意思.
他們不願意.
他們不守如節的話.
就要從民中剪除.
不歸屬耶和華的意思.
他們要擔當自己的罪.
因為如節的意義就是.
原本羔羊的血獻祭.
給耶和華的羔羊的血.
塗在門尾上.
如果神見到這個記號的時候.
這個血的記號的時候.
就越過.
就不會去擊殺屋裡面的人.
那些長子.
但因為他們不願意守如節的時候.
就沒有了這個記號.
沒有了這個效應.
他們就要擔當自己的罪.
所以羔羊的血.
是幫助他們救贖.
是贖罪的血.
十四節告訴我們.
如果有愛人.
愛人是什麼意思.
原來一些藝館就告訴我們.
是寄居者.
寄居在以色列人那些.
有些人願意成為.
自己成為以色列人的時候.
跟他們一起的時候.
他們又走過來.
又願意去守耶和華的節日.
都按著以色列人守如節的定例.
去守的時候.
都是可以備受保護的.
這個是當時我們已知的.
原來寄居者.
都能夠去受到耶和華的保護.

$^{521}$再看這個記號.
我們看到了.
為什麼我說.
要處理不潔的人那麼特別.
是因為.
其實你有沒有想過.
原來他們要.
要等候所有的人.
無論你是潔淨了.
又或者你原本是不潔的.
你要守這個如節的時候.
你去先見.
然後你守完這個如節的時候.
我們先去起行.
神是一個都不會撇下的.
神是會給時間.
他們去潔淨自己.
去守這個如節.
以致他們一起去起行.
去展開新的旅程.
這個就是最特別的記號.
原來神是在等候.
藉著逾越節.
守節.
守如節的標記.
讓到祂去帶領一班的以色列人.
不單是開始的時候.
離開維勞之家.
埃及.
更加在西奈抗日的時候.
他又要展開另外一個新旅程.
就是神去指示他的.
應祂之地繼續進化.
所以對於以色列人來說.
逾越節是一個相當重要的節期.
也都在很多舊約的聖經裡面.
我們看到原來.
很多以色列人守如節的時候.
都是一個重要的標記的新旅程.
剛才都提過.

$^{561}$他第一次守節的時候.
他們離開埃及.
離開430年的維勞之家.
展開新的旅程.
當然這個旅程.
他不知道在抗日裡面.
還要怎樣.
還要有40年的時間.
當然都有很多故事.
但是他離開之後.
三個月.
神要他們再展開另外一個.
停留在西奈抗日一段時間.
預備好自己的時候.
再起步一個新的旅程的時候.
他們要向新的旅程進化.
然後在西奈的時候.
又是剛才說了.
他們要離開西奈抗日的時候.
也都是新的旅程的開始.
他們要走一個未知前面.
是怎樣的路程.
是困難.
是怎樣.
他們不知道.
發生什麼事.
他們不知道.
所以為什麼.
神在他們繼續起步的時候.
神要逼他們.
要他們預備好自己.
所以我們一路再看.
其實還有一個很重要的.
就是他們離開抗日.
走出抗日的時候.
在一甲的耶利哥平原的時候.
以色列人又守如二節.
那時候不知道第幾次.
但是那個重要的標記又是什麼.
約斯亞帶領一班的以色列人.

$^{601}$進入應許之地.
因為他們要攻打應許之地.
還有一些拜偶像的民族.
摩西帶著一班出埃及的以色列人.
他們死在抗日.
在四十年裡面.
那班我們叫舊的以色列人.
他們死在那裡.
因為神說他們是叛逆的百姓.
所以他們不可以進入神的應許之地.
新一批由約斯亞帶領進入.
應許之地的迦南這個地方的時候.
他們要在耶利哥平原裡面.
那批新的以色列人.
他們還沒有受國禮.
他們要重新受國禮.
然後重新去守如二節.
然後再由約斯亞帶領這班的以色列人.
進入耶利哥去攻打.
進入迦南的地方.
神又給了一個新的旅程的標誌.
原來對於以色列人來說.
其實是一個分別為性.
和歸屬於耶和華神的一個很重要的標記.
那班人是我的子民.
是我要保護.
是與我同在的一班子民.
我要帶領他們去繼續前行.
在舊約裡面除了一個性別的子民.
或者分別為性的一個記號的時候.
也標記著以色列人民的復興和更新.
原來在這個王國分裂的時期.
希西加王和約西亞王.
推行這個宗教改革的時候.
要復興聖殿的時候.
他們又守如二節.
因為如二節代表他們要再一次記起.
誰是我們的神.
誰帶領我們去出埃及.
去成為這個王國.

$^{641}$簡單來說是一大段的歷史.
去到這個王國的時期.
他們要再守如二節.
再重新一次他們自己的身份.
透過守如二節.
重新自己的身份.
跟著以色列.
標記著什麼.
就是被擄回歸的時期.
以色列帶領一班被擄回歸的.
以色列人民去守如二節.
告訴我們這班是屬神的子民.
他們回來再一次回歸.
所以在舊約的時期.
如二節代表著一班以色列人民.
進入一個新時代的開展.
生命需要神去幫他們.
去聖潔,復興和更新.
所以守如二節對於舊約的以色列人民來說.
其實是一個重要的標記.
在新約的時期又如何.
這個更加重要.
新約的時期.
耶穌基督就是如二節的羔羊.
有一個很象徵性的意義.
在臨前告訴我們.
因為我們的如二節的羔羊.
基督已經被殺獻祭了.
很明顯.
原來如二節的羔羊.
其實就是耶穌基督自己.
耶穌基督設立聖餐的時候.
在上十字架的時候.
跟門徒在最後晚餐.
設立聖餐的時候.
都自己承認.
祂自己是如二節的羔羊.
祂說這個餅就是我的身體.
這個血就是我立約的血.
所以這個如二節的聖餐.

$^{681}$今天我們叫聖餐.
其實是一個很重要的如二節的標記.
今天我們可能.
現代人我們都不過如二節.
可能猶太人都會過.
但是我們今天用什麼去告訴我們.
再一次紀念耶穌基督為我們所做的.
因為耶穌基督已經是如二節的羔羊.
所以祂被懸掛在十字架上面的時候.
告訴我們一斤的骨頭都不能夠截斷.
約翰福音十九章告訴我們.
他說猶太人原本就想.
吩咐那些兵丁將耶穌的腳打斷.
但是耶穌因為死了.
那些兵丁就說不用打斷那些腳筋.
所以就應驗了經上所說的.
祂的骨頭一斤都不可以截斷.
這裡告訴我們.
在舊約這句話裡面.
應驗這個經上說.
就是出自《出埃及皮》和《民數記》這裡.
如二節的羔羊是整隻的.
是敲全羊的.
也都預表著耶穌基督犧牲的整全性.
我們都明白到.
都紀念到就是說.
原來約翰福音三章受節.
祂怎麼說.
耶穌基督的救贖.
我們一起讀好不好.
預備起.
這裡其實就告訴我們.
耶穌基督救贖的普世性.
一切順他的.
我們看回《民數記》九章十四節的時候.
我們會看到.
若有外人.
如果願意去守這個御節的時候.
他們的童歸一例.
所以一開始在《民數記》.

$^{721}$在舊約的時候.
已經預表了這個救贖的普世性.
所有願意相信.
接受耶穌基督的人.
都是耶穌基督救贖的子民.
這裡又告訴我們.
接著我們三七十六節就很清楚.
十七十八節告訴我們什麼.
祂說.
神的兒子.
他的目的是什麼.
是要拯救世上所有的人.
接著十八節說.
送他的人不被定罪.
不信的人罪已經定了.
就在他們自己身上.
因為他們不信.
或者他不願意接受耶穌基督.
成為他的救主.
《民數記》九章十三節說什麼.
祂就說那些能夠潔淨的.
他們又不是出軟門的.
他們原本應該可以守禦節.
但是他們不願意守這個節的時候.
就要擔當自己的罪.
所以.
其實原來在舊約裡面.
《民數記》給出一個很重要的訊息.
讓我們都很清楚.
原來耶穌基督是世界上唯一的救贖.
史徒行傳四章十二節告訴我們什麼.
我們一起讀一次.
除他以外便無拯救.
因為在天下人間.
沒有次下別的名.
我們可以靠著得救.
在新約的時期.
逾越節代表著耶穌基督是獻祭的羔羊.
祂是世界上唯一的救主.
耶穌基督在新約的時期.

$^{761}$這個救贖標誌著救贖的一個新紀元.
不會好像已經不再是舊約那樣.
需要敲全羊.
一個很實在的羔羊去獻祭.
因為耶穌基督已經自己成為獻祭的羔羊.
所以要每年一次去做一個贖罪.
因為耶穌基督已經擔當了我們的罪.
耶穌自己就是獻祭的羔羊.
一次性的獻上救贖了我們一切.
所有願意相信祂的人.
這是一個很重要的標記.
讓我們繼續看下去.
雖然時間不是很多.
或者已經過了一點.
但也讓我們繼續看後面.
因為這個也很重要.
這裡告訴我們.
原來在當時來講.
帳幕是一回事.
立起帳幕那天.
有雲彩遮蓋帳幕.
然後發跪在帳幕當中.
代表著什麼.
原來雲彩遮蓋帳幕.
代表著神同在的記號和保護.
在初一及記四十章的時候.
完全有提到發跪.
講到發跪被立起的時候.
而民數記七章一至講到分別為聖.
初一及記四十章裡.
也提到雲彩遮蓋帳幕.
而耶和華的榮光在雲彩裡.
告訴我們原來耶和華的同在.
與一群以色列人同在.
也提到雲彩何時收上去.
它就起行.
何時遮蓋的時候.
它就停住.
代表著神一直引領著.
以色列人前行的方向.

$^{801}$何時停何時上.
都是有耶和華的引領在當中.
原來在民數記九章.
十七到二十三節裡.
其實我們看到紅色的標記.
一直告訴我們.
原來出現了九次這麼多.
在這短短的幾節的經文.
或者七節的經文裡.
其實就告訴我們.
出現了九次就是說.
我們要按照耶和華的吩咐.
起行和作息.
這是一個很重要的.
在他們臨行離開西來的曠野的時候.
耶和華與他們同在.
耶和華要指引他們的路.
他們要學習的是什麼.
其實他們要學習的就是這個.
藍色字那裡.
就是有時是這樣.
有時就是神會這樣.
有時就會停.
二十二節告訴我們.
雲彩停留在帳幕上.
無論是兩天.
是一月.
或者是一年.
以色列人就住刑不起行.
直至到雲彩收上去就起行.
這裡告訴我們是什麼.
這個就是神要以色列人學習的.
就是等候.
無論神要他們等一天.
等兩天.
等一個月.
一年的意思就不是七年一年.
原文的意思是一段時間.
所以無論神要他們等多久.
安營多久.

$^{841}$或者要走多久.
都是看神說了算.
這是一個很重要的標誌.
當他們要起行.
臨離開西來的曠野的時候.
他們被神提醒和警告.
你們要按照耶和華的吩咐行事.
今天我們是不是都按照耶和華的吩咐行事呢.
在這段第九章的成章裡面.
告訴我們什麼呢.
就是在舊約的時期.
預節代表著一個新旅程的展開.
以色列人是進入新的時代.
生命要聖潔.
要復興和要更新.
新約的時期預節代表著.
耶穌基督救贖的新紀元.
耶穌基督就是那個憲制的羔羊.
祂是世界上唯一的拯救.
是一個救贖的新旅程.
新的標記.
運彩代表著神的同在引領.
和神的子民同在.
我們要按照神的吩咐行事.
今天我們要怎樣過我們的預節呢.
今天我們要怎樣去守我們的聖餐呢.
很有儀式感.
每個月一次.
但是很多的教會.
很多其他的教會.
他們不單止是一個月一次.
他們是每個禮拜都守聖餐.
守聖餐的意義.
代表著我們真的要重新紀念.
我們的生命裡面有基督.
十字架是耶穌基督為我們犧牲的標記.
是我們每天都要去.
重新得力的一個象徵的標記.
讓我們真的在我們的生命裡面.
我們每天都向神認罪.

$^{881}$悔改.
更新我們的生命.
讓我們的生命都得到復興和聖潔.
我雖然不可以繼續和大家講我的例子.
但是我深信大家在九章裡面.
很明白到.
原來今天我們在生命裡面.
我們要展開每天新的旅程的時候.
我們都要依靠耶穌基督與我們同在.
我們才能夠有一份的能力.
繼續向前行.
讓我們同心去禱告.
慈悲人愛天父.
從我們願意去認你是我們生命的主的一刻開始.
你就與我們同在.
和我們的心一起去引證.
你是我們的主.
我們的神.
願主你繼續帶領著我們每一個.
我們願意跟從你的.
都願意緊守進行神的吩咐.
求神你給我們有明確的指引.
讓我們生命能夠透過你繼續更新.
繼續行我們每一天新的旅程.
我們同心去禱告.
奉主明求.
阿們.
\newpage



\section{民數記 9:1-23}
\label{sec:031979}
\textbf{走一條沒走過的新路!(華)}~~~~董家驊牧師
\newline
\newline
連結: \href{http://www.church.com.hk/acms/content.asp?site=cdc\&op=show\&type=product\&code=031979\&layout=sermon}{\texttt{031979}} ~~~~ 語音日期: 2026-04-12
\newline
\newline
\hyperref[sec:031978]{\small{< < < PREV SERMON < < <}}
~
\hyperref[sec:index_chronic]{\small{[返順時目]}}
~
\hyperref[sec:031980]{\small{> > > NEXT SERMON > > >}}
\newline
\newline
民數記 9:1-23
\newline
\begin{longtable}{cl}
\hline
\hline
章節 & 經文 (和合本修訂版)\\
\hline
9:1 & \begin{tabularx}{0.7\textwidth}{X} 以色列人出埃及地以後，第二年正月，耶和華在西奈的曠野吩咐摩西說： \end{tabularx} \\ \\ \relax
9:2 & \begin{tabularx}{0.7\textwidth}{X} 「以色列人應當在所定的日期守逾越節。 \end{tabularx} \\ \\ \relax
9:3 & \begin{tabularx}{0.7\textwidth}{X} 你們要在本月十四日黃昏的時候，在所定的日期守這節，按照一切的律例典章守節。」 \end{tabularx} \\ \\ \relax
9:4 & \begin{tabularx}{0.7\textwidth}{X} 於是摩西吩咐以色列人守逾越節。 \end{tabularx} \\ \\ \relax
9:5 & \begin{tabularx}{0.7\textwidth}{X} 正月十四日黃昏的時候，他們就在西奈的曠野守逾越節。凡耶和華所吩咐摩西的，以色列人都照樣做了。 \end{tabularx} \\ \\ \relax
9:6 & \begin{tabularx}{0.7\textwidth}{X} 有幾個人因屍體成了不潔淨，不能在那日守逾越節。當天他們到摩西、亞倫面前。 \end{tabularx} \\ \\ \relax
9:7 & \begin{tabularx}{0.7\textwidth}{X} 那些人對他說：「我們因屍體而不潔淨，為何禁止我們，不能和以色列人在所定的日期獻供物給耶和華呢？」 \end{tabularx} \\ \\ \relax
9:8 & \begin{tabularx}{0.7\textwidth}{X} 摩西對他們說：「你們稍等，讓我去聽耶和華對你們有甚麼吩咐。」 \end{tabularx} \\ \\ \relax
9:9 & \begin{tabularx}{0.7\textwidth}{X} 耶和華吩咐摩西說： \end{tabularx} \\ \\ \relax
9:10 & \begin{tabularx}{0.7\textwidth}{X} 「你要吩咐以色列人說：你們和你們後代中，若有人因屍體成了不潔淨，或出外遠行，仍然要向耶和華守逾越節， \end{tabularx} \\ \\ \relax
9:11 & \begin{tabularx}{0.7\textwidth}{X} 他們就要在二月十四日黃昏的時候守節，要吃羔羊，以及無酵餅和苦菜。 \end{tabularx} \\ \\ \relax
9:12 & \begin{tabularx}{0.7\textwidth}{X} 他們不可留一點食物到早晨；羔羊的骨頭一根也不可折斷。他們要照逾越節的一切律例守這節。 \end{tabularx} \\ \\ \relax
9:13 & \begin{tabularx}{0.7\textwidth}{X} 但潔淨又不出外遠行的人若不守逾越節，那人要從百姓中剪除，因為他沒有在所定的日期獻供物給耶和華，必須擔當自己的罪。 \end{tabularx} \\ \\ \relax
9:14 & \begin{tabularx}{0.7\textwidth}{X} 若有寄居在你們那裡的外人要向耶和華守逾越節，他要照逾越節的律例和典章做。無論是寄居的或是本地人，都用同樣的律例。」 \end{tabularx} \\ \\ \relax
9:15 & \begin{tabularx}{0.7\textwidth}{X} 立起帳幕的那日，有雲彩遮蓋帳幕，就是法櫃的帳幕；從晚上到早晨，雲彩在帳幕上，形狀如火。 \end{tabularx} \\ \\ \relax
9:16 & \begin{tabularx}{0.7\textwidth}{X} 經常都是這樣：雲彩遮蓋帳幕，夜間雲彩形狀如火。 \end{tabularx} \\ \\ \relax
9:17 & \begin{tabularx}{0.7\textwidth}{X} 雲彩幾時從帳幕上升，以色列人就幾時起行；雲彩在哪裡停住，以色列人就在哪裡安營。 \end{tabularx} \\ \\ \relax
9:18 & \begin{tabularx}{0.7\textwidth}{X} 以色列人遵照耶和華的指示起行，也遵照耶和華的指示安營。雲彩在帳幕上停留多久，他們就留在營裡多久。 \end{tabularx} \\ \\ \relax
9:19 & \begin{tabularx}{0.7\textwidth}{X} 雲彩在帳幕上停留許多日子，以色列人就遵照耶和華的吩咐不起行。 \end{tabularx} \\ \\ \relax
9:20 & \begin{tabularx}{0.7\textwidth}{X} 有時雲彩在帳幕上只停了幾天，他們就遵照耶和華的指示留在營裡，也遵照耶和華的指示起行。 \end{tabularx} \\ \\ \relax
9:21 & \begin{tabularx}{0.7\textwidth}{X} 有時雲彩從晚上留到早晨；早晨雲彩上升，他們就起行。無論是白天是黑夜，當雲彩上升的時候，他們就要起行。 \end{tabularx} \\ \\ \relax
9:22 & \begin{tabularx}{0.7\textwidth}{X} 雲彩停留在帳幕上，無論是兩天，一個月，或更長的日子，以色列人就留在營裡不起行；但雲彩一上升，他們就起行。 \end{tabularx} \\ \\ \relax
9:23 & \begin{tabularx}{0.7\textwidth}{X} 他們遵照耶和華的指示安營，也遵照耶和華的指示起行。他們遵守耶和華的吩咐，是耶和華藉摩西所指示的話。 \end{tabularx} \\ \\
[1ex]
\hline
\hline
\end{longtable}
$^{1}$港府堂的弟兄姊妹 逐日早安.
非常的抱歉 我不會講廣東話.
所以麻煩各位了.
那今天呢 先代表世界華服中心.
謝謝港府堂一起來對華服運動的支持.
那今天要講的這個題目叫做.
走一條新路.
我想問大家一下 有多少人你是.
如果可以選你喜歡走新路的.
可以舉一下你的手嗎.
有嗎.
好 我是超愛的 但結婚之後.
當太太在旁邊的時候.
我就比較不敢走新路.
今天 昨天收到林承鑫牧師的信息.
他很貼心 就問我說 董牧師.
今天來港府堂 講到.
你知不知道怎麼走.
需不需要我派弟兄去接你.
其實呢 我每一次來都是有人帶我來的.
但這一次我就想 我的題目叫做走一條新路.
那.
我就憑著信心跟林牧師說.
林牧師沒事的 我來過幾次.
應該可以 我自己來.
所以今天早上 我就拎著我的行李箱.
從油麻地 就是走一條新路.
來到這裡.
很多人喜歡走舊路.
因為舊路有安全感.
我們熟悉 我們知道.
有些人喜歡走新路.
因為在新路當中 似乎是.
有新鮮的事情發生.
但其實 本質上.
我們都喜歡走比較安全的路.
如果我們覺得這條新路不安全.
我們大概不會走.
如果我們覺得這條舊路已經也不安全了.
我們可能也想找新路.

$^{41}$所以在新路跟舊路的背後.
根本的問題是.
我們的安全感 到底建立在什麼之上.
在我們今天讀的經文 明書記第九章.
講到的是以色列人.
他們在上帝的拯救下.
離開了埃及.
他們要去到上帝應許他們的迦南美地.
他們要走一條他們沒有走過的路.
如果你打開地圖.
你會看到 其實從埃及到迦南地.
就是今天的以色列.
距離其實不太遠 還蠻近的.
我估計 如果我來走.
一個月一定走得到.
就好像今天早上從尤瑪第走來這裡.
我估計如果有人可以走的通道的話.
應該一天一定走得到.
再怎麼慢 一天大概也走.
划船可能都可以 是來得到的.
可是上帝卻帶以色列百姓走了四十年.
上帝為什麼帶上帝的百姓.
繞道在曠野走了四十年呢.
在一個生活在講求效率的香港.
我想我們都很困惑對不對.
如果老闆交辦的事情.
一個月可以做完 你做了四十年.
我想你大概等不到走完的那一天了.
你就已經被let go.
我們從小也被教導說要贏在起跑線.
我們的社會 我們的文化.
整個世界全球資本主義的大環境.
都不斷的在告訴我們 效率 效率 效率.
我有一個朋友在美國的科技業矽谷工作.
他跟我說他們的公司這幾年叫做效率之年.
意思是什麼呢.
意思就是沒有效率的就準備裁員.
效率似乎把我們推著走.
但是面對以色列的百姓要進到上帝.
應許他們的迦南帝.

$^{81}$上帝更看重的不是效率.
上帝更看重什麼呢.
上帝更看重他的百姓是成為怎麼樣的一群人.
以色列人在埃及其實要進到迦南帝之前.
已經住了幾百年的時間.
坦白講在埃及的當中.
他們已經很習慣埃及的偶像崇拜的生活了.
那是他們的日常.
事實上就在今天讀的經文不久之前.
以色列人才因為找不到摩西.
就為自己雕刻了一個金牛圖.
然後跟人說.
看啊 這就是帶我們出埃及的耶和華上帝.
今天的我們何嘗不是如此.
很多時候我們以為我們信主之後.
我們就在敬拜上帝.
但是會不會很多時候我們跟當年的以色列百姓一樣.
我們只不過是把過去敬拜的東西換了一個名字而已.
我有一個好朋友在台灣的偏鄉牧羊教會.
那裡是偶像崇拜非常盛行的.
而且在台灣的偏鄉其實信主的比例可能不到0.1\%.
傳福音是非常困難的.
當我有一次跟這位牧者在聊天的時候.
我問他說在偏鄉傳福音最關鍵的是什麼的時候.
他跟我說是要認識上帝的話.
我很驚訝.
但是這位牧者他跟我說.
因為在偏鄉偶像崇拜民間宗教太普遍了.
以至於如果他們只是來教會受洗的話.
很多人只是把過去拜媽祖關公的那一套現在來拜耶穌而已.
除非認識上帝的話.
認識上帝是怎麼樣的一位上帝.
生命才能夠真正的更新.
因此相較於多快可以走到應許之地.
上帝更在意的是怎麼樣在這個旅程當中.
改變轉化上帝的百姓.
在明初紀第九章一開始一節到十四節的時候.
其實上帝在先告訴以色列百姓要怎麼守愉悅節.
怎麼守一個敬拜的生活.
然後在講完了怎麼過敬拜的生活之後.

$^{121}$第十五節立起帳幕的那日.
有雲彩遮蓋帳幕.
就是法櫃或者是一般講叫約櫃.
我用的是和本修訂本.
就是法櫃的帳幕.
從晚上到早晨雲彩在帳幕上形狀如火.
經常都是這樣.
雲彩遮蓋帳幕.
夜間雲彩形狀如火.
在要啟行的時候.
上帝最在意的首先要建立是以色列百姓的敬拜生活.
首先教導他們怎麼過愉悅節.
然後立起帳幕.
帳幕是什麼.
帳幕是上帝同在的地方.
帳幕是上帝同在的可見的印記.
為什麼立帳幕這件事情在啟行的時候這麼重要呢.
因為不要忘記在不久之前.
以色列百姓才為自己造了一個金牛賭.
犯罪得罪了上帝.
因此上帝本來是說.
摩西你帶他們去就好了.
我不跟你們去.
我派我的使者保證你們安安穩穩的進到應許之地.
但是我不要與你們同去.
而這個百姓悖逆犯罪.
我看不下去.
可是摩西不斷的為百姓代求.
我們也看到上帝是有恩典.
有憐憫的上帝.
因此最後當上帝說我要與你們同在.
當帳幕立起雲彩降下遮蓋的時候.
意思是什麼.
意思是上帝選擇饒恕他的百姓.
意思是什麼.
意思是上帝的同在不離不棄.
意思是什麼.
意思是上帝的恩典比我們的過半更大.
我很喜歡的一位牧師是美國的提摩泰凱勒牧師.
他已經過世了.

$^{161}$他在生前他有一次他說.
如果用簡單一句話來總結福音是什麼的話.
他是這樣總結的.
他說福音就是.
我們比我們敢想像的更糟糕.
而上帝比我們敢想像的更愛我們.
什麼意思.
就是不論你覺得你是多糟糕的一個人.
對不起你對自己的認識還不夠深刻.
不論我覺得我是一個多壞的人.
對不起我對我的壞的理解還太膚淺.
我們比我們想像的更糟糕.
然而上帝卻比我們敢想像的更愛我們.
因此當我們犯罪的時候.
因此當我們感到自己離上帝很遠的時候.
不論是什麼事情.
我們要記得一件事情.
上帝的愛提醒我們.
我們在最終不是要逃離上帝.
而是要逃向上帝.
不是逃離上帝.
而是逃向上帝.
當立起帳目的那日雲彩遮蓋了帳目.
就提醒了上帝的百姓一件事情是.
上帝的同在在他的百姓的當中.
在一群有罪的百姓當中.
可是上帝的恩典大過我們的罪.
而當這個雲彩降下之後.
接著第十七節.
雲彩幾時從帳目上升.
以色列人就幾時啟行.
要啟行了要啟程了終於要啟程了.
雲彩在哪裡停住.
以色列人就在哪裡安穎.
蠻好的有啟程有安穎.
以色列人遵照耶和華的指示.
啟行也遵照耶和華的指示安穎.
聽起來都很好.
但後面就不太好了.
雲彩在帳目上停留多久.

$^{201}$他們就留在營裡多久.
雲彩在帳目上停留許多日子.
以色列人就遵照耶和華的吩咐.
不啟行.
以色列人啟程了.
而上帝也藉著雲彩來引導以色列的百姓.
他們最終的目的也很清楚.
就是要進迦南帝.
然而這條路的路上.
卻是一條充滿不確定性的路.
所以走一條這條新路.
首先它是一個信心的路.
目的跟終點很清楚.
但是過程時間時機怎麼走.
非常的不清楚.
其實這也如同我們的人生.
其實每一個耶穌基督的門徒.
每一個把生命交給耶穌基督的人.
我們都有一個確據.
是我們的結局是怎麼樣確定的.
我們都有個確據是.
當耶穌基督再來的時候.
我們要與他一同坐王掌權.
當耶穌基督再來的時候.
他要恢復萬物時一切都更新了.
當耶穌基督再來的時候不再有眼淚.
當耶穌基督再來的時候.
一切的不公不義將被伸張.
當耶穌基督再來的時候.
一切的黑暗都將過去.
我們有一個確實的盼望.
知道那是我們的終點.
然而我們怎麼去到那裡.
坦白講我們並不知道.
但這段經文首先提醒我們.
或許重點不是我們知道怎麼去到那裡.
重點是我們如何在過程當中.
抓住上帝的同在.
其實就算你知道未來的路怎麼走.
若是沒有上帝的同在.

$^{241}$我們仍舊不會有平安.
正如同一開始我問的.
你喜歡走新路嗎.
關鍵其實不是新或舊.
而是那是不是我們有安全感.
就算我們知道舊路又怎麼樣.
如果舊路不安全.
如果沒人保護我們.
那又怎麼樣.
但若是有上帝的同在.
不論上帝帶我們現在走的路.
是我們熟悉的還是我們不熟悉的.
我們可以有一個確距.
是這條路最終將帶我們.
進到上帝那確定的未來跟結局.
所以首先這條路.
是一條對上帝信心的路.
而接著第二十節這邊說.
有時候雲彩在帳幕上只停了幾天.
他們就遵照耶和華的指示留在營裡.
也遵照耶和華的指示起行.
有時雲彩從晚上留到早晨.
雲彩上升他們就起行.
無論是白天是黑夜.
當雲彩上升的時候.
他們就要起行.
這條新的路不單單是一條信心的路.
而且也是一條需要有耐心的路.
因為你不知道什麼時候要起行.
你也不知道扎營要扎多久.
想像一下.
我聽說在香港很多時候.
大家很喜歡旅遊.
台灣也一樣.
如果你去參加一個旅行團.
然後在機場集合的時候.
旅行團的領隊說.
我們這次的行程不太知道有多久.
我們去的第一站不確定會住幾天.
但是我會告訴大家.

$^{281}$什麼時候.
那天早上你就知道要不要起行了.
當我說要起行的時候.
大家有一個小時的時間可以打包.
當我說要繼續住下來的時候.
不論你是不是打包好.
你都要繼續住下來.
我不知道有多少人會參加這樣的旅行團.
我們可能都會嚷嚷著要退費了.
可是這就是上帝對他百姓的帶領.
說我會帶領你們.
你們該扎營的時候就扎營.
該起行的時候就起行.
但是要待多久不知道.
你們不會事前知道.
所以你想像一下.
如果你是以色列的百姓.
當上帝的雲彩降臨.
你要扎營的時候.
請問你帶的行李是要.
只要打開一點還是打開全部.
不知道.
如果我只是去一個地方住一晚.
我大概是最簡單的.
只是把牙刷拿出來.
如果住的地方.
已經有所有的衛浴的東西了.
我盡量就不要去動我的行李.
因為隔天一早就要走了.
但如果我要在那邊住一個月.
那我一住進去.
我就是全部行李箱打開來.
可是那是因為我知道.
我是住一天還是住個月.
但如果你是每天都不知道呢.
以色列人不是只是過了一個月這樣的生活.
他們過了四十年這樣的時間.
我很難想像.
去年上帝讓我有一個機會.
讓我稍微的體會到一點這種感受.

$^{321}$去年我有一趟服飾的行程.
是需要繞地球一圈.
先到西班牙再到英國再到加拿大.
然後再回到亞洲.
所以我去的第一站西班牙.
我帶了很多的行李.
然後我也都算好了.
我大概需要的換洗衣物.
然後規劃好要怎麼洗衣服.
當我一到西班牙的時候.
到了機場我非常的開心.
看到老朋友來接我.
我們要一起事奉.
然後就在我們這個.
把行李放上車的時候.
正要上車的時候.
突然一個西班牙的本地的人.
就很慌張的指著我們的車子底下.
我們想怎麼了.
我第一個反應是.
是有炸彈嗎.
立刻就彎下腰去看.
沒有東西啊.
然後我們就起身.
你看我我看你.
然後那個人也不見了.
我們就想說.
那可能他看錯了.
我們就上車.
然後司機一發動車子的時候.
突然說糟糕.
後車廂被打開了.
然後我還沒有反應過來.
我想那就去關就好了.
然後這個司機就立刻繞到後面.
跟我說 董事長你的包沒了.
我想哪一個包.
他說就是你的背包.
我立刻腦筋一片空白.
我的電腦.

$^{361}$我講到的iPad.
我的護照.
我的錢包.
我的現金.
我的非現金 信用卡.
全部在那個包裡面.
我立刻慌了.
而且那個時間是.
復活節之前的時間.
也就是西班牙休假的時候.
然後我一上車.
然後這個司機就跟我說.
董事長 你可能要有心理準備.
你在西班牙可能要住上一兩個月.
我說為什麼.
去報案啊.
他說你不了解.
這裡是西班牙不是亞洲.
我們是很放鬆的.
他說而且你現在遇到的是.
復活節前的這一週叫聖週.
大家是不上班的.
我說那報警把他找回來啊.
他說董事長.
我們在這邊住了那麼多年.
沒有聽說過找回來的.
我在那一刻.
我真的可以體會到一點點.
以色列百姓要過40年的生活.
我不知道我在西班牙待多久.
而且我後面是有很清楚.
我要去英國.
服侍那裡的一個.
香港人的青年的聚會.
然後我要去加拿大.
然後後面這些行程都安排好了.
很多都是大家飛到那個地方的.
我突然之間我不知道該怎麼辦.
也不知道該怎麼規劃.
也不知道該怎麼跟我的太太講.

$^{401}$我就只能打電話跟我太太說.
請為我禱告.
然後突然想起來.
我還有電話.
然後我在想說.
都是Apple的話.
iPhone跟iPad.
是不是可以手機定位.
然後我就看一下我的定位.
發現我的這個iPad.
被定位在一個很偏遠的地方.
我就跟那位司機說.
可不可以帶我去那邊找.
然後那個司機跟我說.
童牧師 你太天真了.
你太天真了.
你以為這個方法之前沒有人用過嗎.
沒有人找到過.
可是我還是不死心.
我說你可不可以拜託帶我去.
他說沒關係.
我們先回去到服侍的地方.
然後我請另外一個人.
他說因為你要去報案.
你要去弄這些東西.
至少要一天的時間.
我們先回去.
然後我們請另外一個人帶你去.
後來有一個姐妹就開車.
帶我去報案.
然後我就再跟姐妹說.
可不可以求求你帶我去這個地方.
她說這個地方是中途無路.
我說對啊.
通常作案要丟掉的東西.
不是就是拋棄在中途無路嗎.
所以很合理.
她就經不住我的要求.
也幫我打好預防針.
說童牧師.

$^{441}$我們那麼多次沒有人找到的.
甚至有一次也定位好在一個房子裡面.
可是我們就是進不去那個房子.
一看就是幫派的房子.
他們也不讓你進去.
你也不能夠破門而入.
長話短說我們就到了那個地方.
然後去之前我就禱告.
一開始禱告當然是上帝啊.
讓我能夠找到我的包包.
因為我是要事奉你的啊.
我後面的去英國不是去旅遊的啊.
我是要服侍的啊.
我去加拿大不是要去觀光的啊.
我是要事奉的啊.
主啊為了你的榮耀為了你的事奉.
為你的百姓你好歹幫我把包包找到吧.
你知道很多時候我們禱告.
跟上帝很多的這種情緒勒索的.
可是不知道為什麼禱告到後來.
突然之間我覺得上帝引導.
我就開始禱告說上帝.
但是吉貨不然.
我相信你有美意.
我相信你把我留在西班牙.
兩個禮拜一個月.
你有你的美意.
雖然我不想.
你有你的美意.
我們就一直禱告一直禱告一直禱告.
弟兄姊妹.
這會不會很多時候.
也是我們的人生.
你現在在的位置處境.
跟你原本想的不一樣.
讀書的時候.
你發現你申請上的學校.
跟你想的不一樣.
工作的時候.
你發現上帝給你的帶領.

$^{481}$跟你的職涯規劃不一樣.
退休的時候.
你以為你的退休會是什麼樣子.
可是實際上退休之後又是不一樣.
然而這段經文提醒我們.
重點不是我們知道會怎麼樣.
重點不是快.
而是上帝仍舊在當中行動.
上帝在每一個過程當中.
其實都在行動.
為了要讓我們一起去經歷它.
也許你很好奇.
我的包到底找到沒有.
我賣個關子.
但是我要說那一趟旅程.
我經歷上帝太多太豐富了.
原本是一個好像計畫好的旅程.
我覺得我就是按照我的計畫.
一站一站的去服侍.
但在那一趟旅程當中.
我經歷到包包失竊.
但也經歷到上帝是又真又活的上帝.
而這不是比有包更寶貴的嗎.
就是在我們的生命當中.
我們經歷到上帝是真的.
我還記得有一次.
那時候我還在讀神學院.
然後上了一門課.
那門課叫做Power Encounter.
全能對抗.
我是從很保守的教會出身的.
所以我從來沒有上過這種課.
我就跟我媽媽說.
媽 我這學期選了一門課.
叫做Power Encounter 全能對抗.
我媽媽聽了很興奮跟我說.
兒子啊 你們神學院那麼先進啊.
還教你怎麼搞教會鬥爭啊.
我說那個Power Encounter.
不是只教人與人之間的全能對抗.

$^{521}$而是指跟靈界的.
那在上那門課的當中.
真的開了我的眼界.
但是最讓我印象深刻的是.
在上課的過程的那個學期當中.
有一天我去探訪一個學生的時候.
這個學生竟然就有靈彰顯了.
然後我是完全沒這種經驗的.
不知道該怎麼辦的.
可是我覺得好像聖靈就提醒我.
我在課堂上所學到的東西.
就一步一步引導我.
然後最後花了四個小時.
把那個邪靈給趕走.
我記得那一天我回到家裡的時候.
我的老婆看到我.
那時候我二十幾歲的一個神學生.
我的老婆看到我 看到我臉色慘白.
就說 你還好嗎.
然後我跟她說.
老婆你知道嗎.
魔鬼是真的.
接著我跟她說.
但你知道嗎.
耶穌也是真的.
上帝也是真的.
我想這不就是我們這一生最寶貴的.
就是在生命中的高低起伏當中.
去經歷到原來上帝是真的.
原來上帝對我們的愛是真的.
原來復活的大能是真的.
因此我們不用害怕靈界.
因此我們不用害怕我們的人生失控.
因為上帝是真的.
第二十二節.
雲彩停留在帳幕上.
無論是兩天一個月或更長的日子.
以色列人就留在營裡不起行.
但雲彩一上升他們就起行.
他們遵照耶和華的指示安寧.

$^{561}$也遵照耶和華的指示起行.
他們遵守耶和華的吩咐.
是耶和華借摩西所指示的話.
我們剛剛講到這條新的路.
首先是一條信心的路.
其次是一條耐心的路.
但第三也是一條根心的路.
走一條他們沒有走過的路.
以至於他們必須要緊緊的依靠上帝.
很多時候我們走一條.
我們熟悉的路的時候.
我們都很老練了.
甚至覺得有沒有上帝都無所謂了.
其實這是非常危險的.
你知道學開車最容易出車禍的時候.
不是剛剛學會的時候.
因為剛學會的時候都很謹慎很小心.
頂多就是不小心碰到一下後視鏡而已.
最容易出事的時候.
是你覺得你剛剛覺得自己很會了.
那是最容易出事的時候.
然而上帝帶領以色列百姓走一條新路.
重點不是繞遠路.
重點是塑造他們.
要根心這一群百姓.
同樣的今天不論面對我們個人的生命.
還是面對教會的群體.
我們也需要不斷的心意根心而變化.
我們需要對上帝忠心.
但是在華人教會當中.
我們必須要常常的反思.
我們忠心的對象到底是什麼.
我們忠心的對象是我們自己的經歷呢.
是我們熟悉的文化傳統制度.
還是我們忠心的對象.
是上帝和上帝在今日的工作.
我們忠心的心態.
是一種捍衛擔心憂慮防衛的心態呢.
還是我們的忠心的心態是一種.
不斷的對上帝敞開.

$^{601}$期待被上帝驚喜.
我們的忠心到底是怎麼表達.
今年2026年是華扶運動的第50週年.
華扶是1976年.
王永信牧師,藤靖輝牧師.
他們那一代的領袖.
在50年前.
他們看到華人教會的存在.
不應該只是自保求存.
而是應該要把福音帶到萬國萬民當中.
他們看到華人教會.
我們不應該是各自有各自的山頭.
而是能夠彼此連結.
為著福音的使命.
他們在他們那一代.
回應了上帝給他們的託付.
50年後的今天.
其實今年華扶運動50週年.
我們有一群的華扶的同工.
都一直在禱告說.
到底我們要怎麼樣來過50週年.
我想在華人的文化當中.
50週年就是慶典.
50週年就是要數算過去的50年.
上帝成就了什麼.
這些都很好.
這些都很對.
可是兩年前.
當我們一群華扶的同工領袖.
在香港我們有三天的退休會.
一起在禱告尋求說.
上帝要我們去紀念這50年的時候.
第一天坦白講.
我們每個人都有不同的想法.
我們覺得好像不可能有共識.
但到了第二天的時候.
馮浩留醫師Patrick Fong.
OMF的國際前總主任.
他突然就分享到.
他就說我們需要悔改.

$^{641}$從悔改開始.
他講完之後.
我們在場所有的人.
都感到上帝在對我們說話.
我們需要從悔改開始.
紀念50年.
重點不是這50年我們做了什麼.
也不是50年前的前輩做了什麼.
而是50年前的前輩.
他們當年如何聆聽上帝的工作.
上帝的聲音.
然後回應了.
那是我們在今天50年後要效法的.
這是最好紀念的方式.
就是要聆聽上帝在此時此刻的工作.
如果我們覺得我們今天所活在的時代.
是一個困難的時代.
50年前那是石油危機的時代.
那是冷戰的時代.
如果我們活在50年前.
如果我還沒出生.
但我們在座有些人當年已經在了.
會知道其實當時也是一個風雨飄搖.
我們不知道未來要怎麼走.
一個混亂的時刻.
可是上帝在那個時代仍舊有工作.
同樣的今天50年後的今天.
上帝也仍舊在工作.
我們的問題不是我們的出路在哪.
我們的問題是上帝你在做什麼.
我們的問題不是我要如何的做決定.
我才能夠有一個美好的未來.
我們的問題是.
上帝我們如何能夠與你同工.
而我們會發現.
當我們在走這一條仰望跟隨的心路的時候.
其實上帝的每一個帶領每一個突破.
都是讓我們對他有一個新的認識.
教會今年的主題.
愛主更深,門旅同行.

$^{681}$門徒的旅程也是如此.
我們過去對上帝經歷的經驗.
不是要用來複製,貼上.
而是要因著過去的經歷.
而建立對上帝的信心.
你知道這個很重要.
我們很多時候.
我們很容易想要複製過去.
然後把它貼上.
在自己或別人的身上.
我記得當我二十幾歲.
剛剛結婚的時候.
二十四歲的時候.
我跟我太太那時候都相對年輕.
二十四歲結婚.
我們覺得生孩子不急.
所以也沒有太急.
等到二十七歲.
神學院畢業.
開始穩定事奉的時候.
就覺得可以開始嘗試了.
沒想到要懷孕原來沒那麼簡單.
不是說有就有 名利就利.
試了一個月 試了兩個月.
試了半年 試了一年.
都懷不上.
後來終於鼓起勇氣去看西醫.
西醫檢查完.
然後說你們都沒有問題.
應該只是壓力.
我們又去看中醫.
中醫講得更玄了.
說你們那個線還沒有接上.
我說那個線要怎麼接上呢.
那位中醫師說.
你多來我這邊扎針灸.
針個半年應該就接上了.
我們就針了半年.
但是還是沒有接上.
後來教會有一個姊妹.

$^{721}$一位阿姨很熱心.
就跑來跟我分享說.
董事長 我跟你講.
我年輕的時候也無法懷孕.
可是後來.
當我開始食醫奉獻之後.
我就懷上孩子了.
我當下不知道該怎麼回應.
這位姊妹.
你是懷疑我沒有食醫奉獻嗎.
還是說.
What are you trying to say.
好 弟兄姊妹.
我要講.
這位姊妹她因著食醫奉獻.
經歷上帝.
這是不是真實的.
是.
她因著食醫奉獻.
經歷上帝.
建造她對上帝的信心.
這是不是好的.
是 她的經歷是好的.
她的經歷是美的.
她的經歷是真實的.
可是當我們把個人的經歷.
當作是要複製.
貼上在別人身上的時候.
其實這就變成律法主義.
不是嗎.
所以很多時候在教會.
我們不是故意.
要變成法利賽人的.
我們都是好意.
我們都是出於善意.
你的孩子是這個樣子.
我的孩子以前怎麼樣.
但是怎麼樣怎麼樣之後.
他就怎麼樣怎麼樣怎麼樣.
意思就是.

$^{761}$那你是不是沒有怎麼樣怎麼樣.
我們都是善意.
真的都是善意.
而且我也做過類似的事情.
都是善意.
可是可能在聽的人的耳中.
就是律法主義.
可是今天的這段經文提醒我們.
過去的經驗.
過去走過的路.
是有價值的.
但它的價值不是在於.
未來可以複製貼上.
它的價值是.
讓我們去認識經歷上帝.
是怎麼樣一位上帝.
以至於在面對.
前面的未知的時候.
我們可以有一個確距.
往前走.
因為上帝是一位.
不斷做心事的上帝.
我們也可以有一個盼望.
就是上帝最終.
會做成他的事情.
其實我們既然知道.
我們的終點在哪裡.
因此在每一步路上.
雖然我們當下看不懂.
上帝到底在做什麼.
但是我們都可以有一個確距是.
我就算不懂.
上帝仍舊有美意.
八年前.
那時候我在美國服侍.
那在美國我是服侍華語堂.
就是講普通話的.
那當時我就已經決定說.
我遲早是要回到亞洲來事奉的.
所以我就開始去規劃.

$^{801}$怎麼樣慢慢地.
把我在美國.
所牧養的華語的群體.
找到合適的人.
變成學生可以接棒.
我要準備回到亞洲.
可是當時在預備的過程當中.
教會發生一件事情.
就是教會的英語堂突然出事了.
教會兩位英文牧師都先後離職.
那教會的董事會主席跟主任牧師.
就左看右看.
不知道該怎麼辦.
就突然看到.
那時候我三十幾歲.
董牧師是相對年輕的.
英文也應該可以.
就問嘉華你要不要來當英文堂牧師.
坦白講我是覺得說.
這關我什麼事.
而且如果我是要留在美國長遠牧會.
我覺得這是很好的一個歷練.
對不對.
你華語堂牧養了十年了.
然後去英文堂歷練幾年.
那將來在美國的服侍.
好像都非常的順.
可是我是一個已經決定.
要回到亞洲服侍的人.
而且我不是要來亞洲的.
International Church服侍.
可是因為在教會夠久了.
我已經很圓滑了.
我知道要怎麼回答.
我就跟我的主任牧師說.
請你給我四十天讓我禱告.
四十天後我再回覆.
人很詭詐的.
其實當我講這句話的時候.
我都已經擬定好策略了.

$^{841}$就是四十天的當中.
我會為這件事情禱告一次.
然後四十天之後.
當他們再來找我的時候.
我就可以說.
經過四十天的禱告.
反正中文沒有分單數或複數.
所以經過四十天的禱告.
一次也是 兩次也是 多次也是.
我就說經過四十天的禱告.
沒有負擔 沒有感動.
沒有 上帝沒有說話.
我都想好了.
正當我在這四十天當中.
決定就是在休息的時候.
有一次我在慢跑.
我就想說.
我至少要禱告一次.
才能夠憑著良心說誠實話.
於是那次我在慢跑的時候.
我就跟上帝禱告說.
好 上帝我現在要認真尋求了.
你有沒有要我去英文堂.
如果你數到三.
我數到三 你不說話的話.
意思就是沒有要我去囉.
你看多敷衍.
一邊跑步一邊跟上帝這樣講.
一 上帝沒有說話.
二 上帝沒有說話.
就在我心裡面要數到三的時候.
突然之間有一個很微小的意念.
就問我說.
佳華如果今天.
你的兩個兒子在英文堂.
你會不會去.
當時我的兩個兒子.
都還是在兒童主日學.
還不在英文堂.
我覺得如果我的兒子在.

$^{881}$那應該會去吧.
突然之間那個意念又來.
第二個意念.
難道今天只是因為.
他們不是你親生的兒子.
你就不去嗎.
我的眼淚立刻就流下來.
那不是感動的眼淚.
那是一種羞愧的眼淚.
我突然發現.
原來我在想這件事情的時候.
完全不是去問上帝的心意.
也不是一個牧者的心腸.
我只是從我的career path的角度來想.
就是我是一個要回亞洲.
牧養華語教會的人.
英文堂出事關我什麼事.
而且又不是我搞出來的事.
可是上帝卻讓我看到他的心.
於是我就明白.
上帝要我去英文堂.
坦白講.
在那兩年的英文堂的服事當中.
我真的覺得我是在繞路.
上帝帶我繞路.
好幾次我在台上講道的時候.
我覺得台下是完全disengage.
不知道我在講什麼.
然後當我在用華語講道的時候.
即便在香港.
你們聽我的華語很努力了.
可是至少我還知道.
怎麼跟你們engage.
可是當我用英文講道.
看到大家disengage的時候.
我想要努力地講一個笑話.
都發現我的英文笑話一點都不好笑.
而他們真正會笑的時候.
都是我講英文講得口音有問題.
發音發錯 文法錯的時候.

$^{921}$他們台下就會有些青少年.
還沒有社會化.
就會在那邊笑得很開心.
我裡面都覺得說.
我好像是一個小丑一樣.
好像我每個禮拜來主日講道.
是在講脫口秀 talk show.
我覺得很痛苦.
可是我當時就覺得說.
好吧 上帝你既然這樣帶我.
我雖然不明白.
可是我就當作這是一個順服的操練.
兩年後找到了英文的牧師.
我邀請這位英文牧師來接.
我也跟教會正式地提出.
說我要回到台灣.
於是我就回到台灣牧會.
我很開心覺得我終於解脫了.
覺得上帝啊.
我為了你繞這個路了.
我也不明白上帝的意思是什麼.
但回到台灣之後.
三年之後.
上帝把我放在華府總幹事的事奉.
我接了華府總幹事.
我才突然明白.
原來那兩年的操練.
是一點都沒有浪費.
因為很多時候我在各地服侍的時候.
可能有百分之十到二十的場合.
我需要用英文講道.
在那一刻我才突然明白.
原來當年我覺得繞路.
根本不是繞路.
是上帝在預備.
弟兄姊妹.
你覺得你的生命在繞路嗎.
你覺得你的生命卡住了嗎.
或者你覺得你好像.
知道上帝要帶你走一條.

$^{961}$你不明白的路.
你很害怕.
因為那不是你熟悉的.
也許是在我們的事業.
也許是在我們的事奉.
也許是在我們跟親子的關係上.
很多時候我在.
教養我的兩個兒子的時候.
我也覺得我好害怕.
上帝藉著他們兩個.
帶我走一條我沒走過的路.
他們個性跟我不一樣.
他們的喜歡的東西.
跟不喜歡的東西.
跟我也很不一樣.
很多時候我看著他們很焦慮.
想說你們這樣.
以後要怎麼活下去.
可是講到的時候.
又說要對上帝有信心.
我覺得自己好虛偽.
好兩面人.
但是不正是就在一次一次.
我們在很誠實的.
面對我們裡面的軟弱當中.
一次一次去經歷上帝的真實.
因此今天的早上.
我想邀請大家.
若你覺得你正在.
感到你的人生在繞路.
或者你正感到.
上帝要帶你走一條新路.
是你沒有把握.
是你沒有經歷.
是你不明白的.
我們願不願意.
有一個姿態.
伸出我們的手.
好像一個小孩對爸爸說.
爸爸我不明白.

$^{1001}$但求你牽我的手.
帶我走前面的路.
並且有一個盼望.
我們的父親不會鬆手.
我們天上的父.
祂既然帶領我們.
就必帶領我們走到終點.
我們一起來禱告.
天父謝謝你.
謝謝你今天早上.
藉著以色列的旅程提醒我們.
不論我們現在人生.
在什麼樣的處境.
不論我們自己覺得.
我們在繞路.
原地踏步.
還是被邀請進入一個.
我們懼怕的新的處境.
但今天的早上.
我們伸出我們的手.
對主說主啊.
求你牽我們的手.
引導我們走前面的路.
我們放下.
我們想要掌控未來的慾望.
我們放下.
我們想要消除一切.
不確定性的渴望.
而是學習單單的信靠你.
主我相信.
你要在我們身上所成就的.
超過我們的所求所想.
而且不單單.
使我們自己.
經歷你的根心祝福.
你還要透過.
我們每一個的順服.
使我們成為他人的祝福.
我們這樣同心的禱告.
奉考耶穌基督的聖名.

$^{1041}$Amen.
\newpage



\section{以賽亞書 66:1-6}
\label{sec:031980}
\textbf{蒙主恩待的敬拜者}~~~~黃茂林傳道
\newline
\newline
連結: \href{http://www.church.com.hk/acms/content.asp?site=cdc\&op=show\&type=product\&code=031980\&layout=sermon}{\texttt{031980}} ~~~~ 語音日期: 2026-04-12
\newline
\newline
\hyperref[sec:031979]{\small{< < < PREV SERMON < < <}}
~
\hyperref[sec:index_chronic]{\small{[返順時目]}}
~
\hyperref[sec:031981]{\small{> > > NEXT SERMON > > >}}
\newline
\newline
以賽亞書 66:1-6
\newline
\begin{longtable}{cl}
\hline
\hline
章節 & 經文 (和合本修訂版)\\
\hline
66:1 & \begin{tabularx}{0.7\textwidth}{X} 耶和華如此說：「天是我的座位；地是我的腳凳。你們能為我造怎樣的殿宇呢？哪裡是我安歇的地方呢？ \end{tabularx} \\ \\ \relax
66:2 & \begin{tabularx}{0.7\textwidth}{X} 這一切是我手所造的，這一切就都存在了。我所看顧的是困苦、靈裡痛悔、因我言語而戰兢的人。這是耶和華說的。 \end{tabularx} \\ \\ \relax
66:3 & \begin{tabularx}{0.7\textwidth}{X} 「至於那些宰牛，殺人，獻羔羊，打斷狗頸項，獻豬血為供物，燒乳香，稱頌偶像的，他們選擇自己的道路，心裡喜愛可憎惡的事； \end{tabularx} \\ \\ \relax
66:4 & \begin{tabularx}{0.7\textwidth}{X} 我也必選擇苦待他們，使他們所懼怕的臨到他們；因為我呼喚，無人回應；我說話，他們不聽從；反倒做我眼中看為惡的事，選擇我所不喜悅的事。」 \end{tabularx} \\ \\ \relax
66:5 & \begin{tabularx}{0.7\textwidth}{X} 你們因耶和華言語而戰兢的人哪，當聽他的話：「你們的弟兄，就是恨惡你們，因我名趕出你們的，曾說：『願耶和華彰顯榮耀，好讓我們看見你們的喜樂。』但蒙羞的終究是他們！ \end{tabularx} \\ \\ \relax
66:6 & \begin{tabularx}{0.7\textwidth}{X} 「有喧嘩的聲音出自城中！有聲音來自殿裡！是耶和華向仇敵施行報應的聲音！ \end{tabularx} \\ \\
[1ex]
\hline
\hline
\end{longtable}
$^{1}$弟兄姊妹 早安 平安.
多謝主席簡單介紹我.
你坐在這裡聽別人介紹自己很有趣的感覺.
好像有一個更加懂你的人在這裡幫你.
介紹一下你.
剛才大家看的那段聖經.
和我可能展示的slide那個不同.
我會用的是和學本.
可能你是用其他的新的譯本也不一定.
自己為了這個訊息也有禱告了一段時間.
而事實就是當我讀以賽亞書的時候.
以賽亞書給我的感覺是很害怕的.
因為它講到神的審判.
講到那些百姓那種虛假.
那種得罪神又不知死.
然後依然顧我.
但中間又好像有些地方令我們覺得有些興奮.
看到有安慰的說話.
特別是在以賽亞書的四十章開始.
一直往後又講到受苦的僕人的預言.
然後我就在那裡越來越覺得比較容易看.
又好像有很多連繫到新約.
不知道是不是你的感受.
然後我很深刻到讀到最尾的時候.
六十六章其實正常來說的心態.
一個人看以賽亞書去到那時候.
就好像你看預場片的結局.
結局是怎樣?齊歡喜同慶賀.
走出來拉手仔唱歌仔.
但是他突然間轉了主題.
無論之前講到那些審判也好.
講到那些對百姓的宣告他們怎樣不對.
又甚至我們看到那個救恩這位救主.
如何在預言裡將會來去成就.
可能有些人會在想自己和這位主的關係.
所以其實以賽亞書的六十六章.
我相信以賽亞是在說.
他自己可能同樣在聖殿那裡.
他有一個很大的感受.
就是我們和神的關係究竟是怎樣.

$^{41}$所以其實六十六章裡面是在說敬拜.
以賽亞是在釐定自己.
究竟在神面前應該是怎樣.
如何觸摸這一位這麼超越.
但是同他又時時給他說話.
而他又是會親近他的.
其實這個就是精髓.
所以讀到這一章.
如果你只是很機械式的看.
你就會讀過了.
其實看聖經就是這樣.
一定要心停留.
有一個秘訣.
有一個秘訣就是.
聖經是怎樣寫的呢.
聖經原來就是拿著那些.
他們肯定不是原子筆.
他拿著一支蘆葦等等.
不是叫藥水.
總之是有些像墨一樣的東西.
在一個紙張那裡寫.
不會快的.
是要停留的.
如果我讀過希伯來文.
每一個字都是逐個字母寫的.
逐個那些字.
是停留在那裡去思考的.
讓我們都有這樣的心情去讀他的說話.
我們一起祈禱.
前外的天父我們主耶穌基督.
願你的令帶領我們去思想.
默想你的話.
我們去思想敬拜.
幫助我們.
帶領我們這段時間.
幫助說的.
聽的.
交在你的恩手.
奉主耶穌基督的聖名.
阿們.

$^{81}$看一些記載.
遠到加州.
大概在20年前.
有個人分享到他去教會的經驗.
他說.
他一去到禮拜堂的時候.
你們都有一個來賓的接待.
他就去到那裡.
那個是西人教會.
他去到那裡就叫一個Hub.
就是一個接待和聯絡的中心.
然後就說 先生 你喜歡傳統敬拜.
傳統聖詩.
現代敬拜.
現代音樂.
還是喜歡爵士樂.
還是其他不同的音樂的敬拜.
就有不同的通道去那裡參加.
彷彿不同的語言.
不同的方法表達.
去滿足當時代不同的音樂.
個人喜好需要.
就是其中敬拜的表現.
當然今天更加有網上的模式.
有不同的語言.
有實體.
甚至我們也有網上播放.
大家好 你在網上.
我雖然見不到你.
但我盼望我們一同跟著這個崇拜.
但我自己的思想.
可能最重要的就是.
想想我們應該怎樣去敬拜這位神.
怎樣去親近祂呢.
是我們自己一廂情願.
我們覺得有這個方式更加好.
還是更加按神自己給我們看到的模式呢.
在66篇裡面剛才我們讀.
有講到這位主.
是一位什麼的主呢.

$^{121}$是一個很超越的.
還未按.
我先跳到那裡.
現在實況就是.
這位主是一個怎樣的主呢.
有人說是有威嚴充滿天地的創造主.
沒有了 大家要背了.
第二件事講到他.
另一件事就是坐在寶座上的王.
今天當我們來到敬拜的時候.
我們記錯了.
我作為一個先知.
我都提醒自己一件事.
神的啟示自己就是.
祂是一個很超越的.
祂是創造萬有的.
祂是在寶座上的王.
我們要以祂這個身份.
來參與敬拜.
弟兄姊妹.
問問我們自己.
在你來的時候.
你會覺得你在去哪裡呢.
你有沒有想過.
我去的這個地方.
是很大的事.
因為是創造主和坐在寶座的主.
要我們去敬拜祂.
我很記得.
初初年輕的時候.
西裝都沒有見.
我媽媽給我一件什麼呢.
一塊東西掛在頭上.
跟著就有一件.
很像襯衫的東西.
有條領帶不用打的.
跟著就給我一件不知道什麼.
衣服蓋起來.
就可以拍照了.
我就參加一些.

$^{161}$什麼什麼表演.
或者是大的.
宴會.
都不是真的好像.
整個襯衫打領帶和西裝.
我心裡面就覺得.
何必搞這麼多呢.
原來那次是參加一些.
什麼表演.
要拍照的.
或者某次去參加.
一個新年的團拜.
弟兄姊妹.
我們去一個比較重要的會很講究.
因為.
有一個不同的對著.
到我大一點要去聽演唱會.
我就留意.
要穿一些.
比較隆重的.
如果你.
去某些餐廳.
你不是.
很認真地穿.
他不讓你進去.
我試過去紐約.
我進去.
外面等什麼.
我沒有領帶.
那個人要打領帶.
他打了領帶.
他才可以進去.
有些餐廳還誇張.
你沒有外套.
西裝.
他就去那裡挑.
有些都很高級.
他都顧著你.
我就覺得何必這麼麻煩.
不是的.

$^{201}$有些情況.
嚴重性.
敬拜神是不能夠.
不知道你敬拜的這位神是誰.
所以我們要記住.
是有威嚴充滿天地的創造主.
他是坐在寶座上面.
反而我們今天不是.
最主要講這裡.
因為我最驚訝.
當我不斷去看這段聖經的時候.
我是被第二節所吸引.
因為第一節就告訴你.
你不要以為有一個地方.
巴比爾節神就會覺得.
朱尊降貴來咒我們.
其實不會的.
不會因為建了一個聖殿.
聖殿就會大過這位神.
這位神仍然永遠超越.
這個敬拜的地方.
或者參與的.
那些活動.
所以最重要的那個中心是神自己.
所以我們.
無論我們去到人生的什麼階段都好.
我們知道敬拜.
就不是只聽聽一個故事.
聽聽一些講道訊息.
所以有些人以為敬拜就是.
開一條youtube片.
講道的.
就算王傳道也有些片.
可以去網上聽.
聽完那一堂.
那就是我都參加完敬拜.
那我就跟著去打golf呢.
和總統一起.
在附近打球那樣行不行呢.
如果他侵你的話.

$^{241}$我們知道敬拜不是這樣的.
敬拜就是有一個這樣的主.
我們去親近他.
所以我最被吸引的就是.
這位主既然這麼超越.
根本可以說我們沒辦法.
這麼容易滿足到他.
所以很想分享的就是.
第二節裡面講到.
究竟這位主所.
look for什麼.
主自己這樣講.
祂是我所看顧的.
我是用和學盤的.
祂look for什麼呢.
什麼叫看顧啊.
看顧是有一個定義的.
看顧就不是看到.
就好像人家叫你.
拜託你看著我的小孩.
是什麼意思啊.
行 我就看著他.
哎呀 掉到地上了.
看著他 看著他.
看著他怎麼死.
是不是這樣看啊 當然不是.
並不是就這樣看.
那個意思就是說.
祂是帶有一種要.
要恩待他 要關照他.
要對他好的.
這個就是說.
神會看顧.
這些這樣的人.
所以也可以這麼說.
神真的會看得上.
神是會恩待.
這樣來敬拜他的人.
你說重要嗎 第二節.
我覺得很重要.

$^{281}$重要到其實你應該背了它.
我所看顧的就是.
虛心痛悔.
因我華而.
賤敬的人.
很重要.
那我就查過.
究竟是說.
虛心痛悔.
再去定義.
要因祂的華賤敬啊.
還是說三樣東西啊.
說一樣東西啊.
究竟原來的文字.
是怎麼表達的呢.
那我分析過了.
其實是說神是.
正在找.
三樣東西.
是三樣東西.
神是在我們敬拜者的身上.
如果要恩待他.
就要看這三樣.
最主要的成分.
所以弟兄姊妹.
我不要將注意力.
走到那些外圍的東西.
有一次有人問.
王傳道 究竟.
這樣的敬拜的歌.
和這樣的敬拜的歌 哪個比較好呢.
我就馬上抓他回來.
那你說什麼是敬拜呢.
他就說.
他就怕了.
其實我們真的要想想什麼是敬拜.
神所看什麼的敬拜者.
是關鍵的.
第一.
經文很清楚說是虛心.

$^{321}$虛心在神的面前.
有些翻譯就是.
很貧窮啊.
如果用.
馬太福音五章三字.
虛心的人.
在靈裡面貧乏.
靈裡面很貧乏.
很窮啊.
我們都很怕很窮啊.
但是有一樣東西.
貧乏又知道.
在神面前是這樣.
是好的.
就是在神面前虛心.
在神面前看到自己的貧乏.
在敬拜的.
背景裡面.
有些人認為.
他能夠在神面前謙卑自己.
我們會覺得自己有很多東西.
事實我們真的有很多東西.
在.
幾貧乏的.
發達世界.
有不少人所擁有的東西.
比起一些發展的國家.
可能有的物質是多很多.
但是.
他不知道的是什麼呢.
他心靈的貧乏.
有些朋友分享到.
他頭幾天.
一退休.
他就發覺.
他的生活好像窮了很多.
沒有一個固定的地方.
去上班.
去開會.
去談什麼.

$^{361}$沒有那些場景.
或者他做管理的.
沒有那些手下跟著他.
要給他指令.
去研究.
再不會有機會去發表他的研究.
等等.
就發覺自己好像沒有了很多東西.
弟兄姊妹.
來到神的面前.
我們有沒有去想一下.
其實我們在神面前.
沒有什麼.
是真的我們.
是很卑微.
很少的.
虛心不只是他覺得自己很窮.
其實就是內心那種的貧乏.
是神能夠自己滿足到自己.
是神真的能夠滿足到我們.
正面一點來說.
有人認為這一節聖經.
當我們去敬拜神的人.
我們不能夠離開.
一樣東西就是.
我們會覺得自己很需要神.
很需要神.
有一個作者.
叫陶樹.
他在他寫的書.
有一個祈禱.
他感受到.
他知道神很好.
神啊我已經嘗到你的美善.
你的恩惠使我滿足.
也令我更渴望你.
他的缺乏.
他知道他缺乏這位主.
他渴望這位神.
他有極其的迫切要求.

$^{401}$更深的恩典.
我因為自己缺少.
這個願望.
很抱愧.
很慚愧.
神啊我要更加渴望你.
更仰望你.
我渴望你比我更渴慕的心.
求你將你的榮耀.
顯給我看.
使我認識你.
弟兄姊妹.
我相信以蔡霞所提到.
當我們去敬拜神的人.
我們是覺得自己很需要神.
我們有時也想想.
自己生命裡是否需要神.
有些人真的.
可能在生活裡.
慢慢覺得自己好像不太需要神.
很多事情都是.
我們用自己的想法想好.
我有一次.
挑戰一個群體.
他們全部都是退休人士.
我說你們有沒有這個想法.
你們所謂.
準備的退休是怎樣退休.
你們只是.
準備自己在財務上.
自由.
或者獨立.
是不是.
他們個個都靜了.
個個說怎樣退休.
就是認為.
你能不能夠退休後繼續生活.
有沒有人想過.
自己退休後的事奉生活會不同.
有沒有人想到.

$^{441}$自己退休後和神的關係.
其實是要繼續加強.
有沒有人想到.
時間多了不等於和神更親近.
可能玩遊戲機多了.
可能看影片多了.
煲爆電視.
我們究竟.
是不是需要神.
當我們要敬拜的時候.
我們這個心.
是否很想參與.
這件事.
有些人去參加.
講座是告訴你.
怎樣去求生.
有些叫什麼獎章.
不知道有沒有人.
比較年輕的參加過.
那些年輕人.
我知道很緊張.
為什麼.
我講的以下的東西.
當我扔了你在荒島的時候.
也不是真的荒島.
扔了下去的時候.
你真的要記得那些秘訣.
譬如怎樣選擇.
哪些可以吃不可以吃.
怎樣生火等等.
你真的要學.
你真的緊張.
為什麼.
有時我們覺得和神.
好像很遙遠.
真的覺得心底.
不是真的需要神.
大衛就有這個感受.
這位私人.
這位牧者陶樹.

$^{481}$他的感受就是.
他很多次提到.
渴慕這位神.
這位主有很多事我們很羨慕.
很想知道.
要有這種切勿他的心.
人覺得自己.
很足夠的話就不會有期望.
你覺得自己夠.
你就不會再.
期望什麼.
當你覺得你自己看聖經也看得很足夠的時候.
甚至在聖經老弟家的.
書信裡面說.
主要責備這個教會.
的使者.
就是他們不知道自己.
困苦可憐.
貧窮.
赤眼赤身.
簡直是傷殘人士一樣.
有沒有一個真的傷殘的.
不知道自己是傷殘的.
不會的.
這件事.
是享用不到那個功能.
是不是.
他不能夠看得見.
他是知道的.
我們近視深了.
或者我腦花厲害了.
我們都感覺到.
我們要檢查他.
弟兄姊妹.
但在淑齡來說.
如果她連貧窮她不知道.
是很可憐.
所以在二賽亞裡面.
神要看顧.
神要看到哪些人.

$^{521}$這些人.
他們親近神是不同的.
因為他們虛心.
他們很需要神.
他們有一種切無的心.
不單止這樣.
有些人可能覺得.
我去到知道多點東西.
認識多點東西也好.
貴教會將我的履歷登了出來.
我不知道有沒有人會覺得.
讀熟了那兩板子.
就認識我.
很認識我.
照片也看了.
知道我是誰.
你們是否這樣認識一個人.
你不會吧.
你會多點認識他.
你會和那個人接觸.
所以我相信.
真的老友不會拿著你的照片.
回家就這樣.
我和你很老友.
有你的照片我以後不用見你.
有沒有人會這樣去看.
親近神.
很舊了.
我未退休前都幾十年不去敬拜神.
我退休後就周圍走.
不再回崇拜.
你覺得好像很瘋狂.
我見到有些人真的像.
因為不斷旅行.
旅行到一個地步.
五十二個星期.
三分二.
都是在路上.
在路上.
而在路上.

$^{561}$他時時都沒有空去崇拜.
我覺得很嚴重.
所以我覺得有時候.
怎樣去看崇拜.
在他未退休前.
都要解決神學觀.
但我不是說那個問題.
我想說第二個大點.
就是痛悔.
除了虛心這個.
忍神所看得上的就是.
他有一種痛悔的心.
痛悔在.
舊約聖經裡面.
主要有兩個字眼.
原來文字.
一個就是.
我們看大衛的.
懺悔詩.
當他說神要的祭.
就是憂傷的靈.
憂傷痛苦的心神不興旱.
講到他那種被壓傷.
被crush 弄碎了.
就是一個主要的字.
包括剛才讀的第二段.
五十七章十五節.
他和什麼人住在一起.
時時和他在一起就是痛悔的人.
那裡的痛悔就是壓碎.
在罪的狀態很慘的人.
我就在研究.
但是.
理賽亞書六十六章二節.
在舊約只用過一次的另外一個字.
我們杜希伯來文.
就叫做.
這個字.
就不是打到他碎.
就是打到他傷.

$^{601}$打到他傷底下.
他在神面前痛悔.
是什麼意思呢.
我想了很久.
可能就是他不肯謙卑.
就打到他服了.
是不是.
就好像那些比賽.
是都不.
給白毛巾.
是都不拍那幾下.
投降.
要KO他為止.
可以是這樣.
但是更加重要的就是.
那裡是說到.
某一個敬拜者.
那裡是單數的.
就是每一個敬拜者.
原來我們去到神面前.
我們覺得自己無憂.
我們謙卑自己.
但是在我們裡面要有些東西.
神是知道要打他.
打低那樣東西.
來去親近神.
用.
再去解釋多一點.
其實有時候我們去到神的面前.
有些東西阻礙我們去敬拜神.
神是要打低那樣東西.
每個人都有些不同的.
打到那樣東西傷.
讓他真的可以去敬拜神.
我自己是一個.
想靠自己的人.
神給我一信耶穌.
一年都不夠.
我就有個大病.
我就要進醫院.

$^{641}$在過程裡面.
神就是去.
其中一個磨練就是.
時間很多.
怎麼用好呢.
開了信主.
都很短時間就投意到自己.
原來其實讀聖經讀很少.
不是很看重.
讀神的話.
在那段時間.
時時在醫院又要檢查.
然後又要做這樣那樣.
結果就多了很多時間.
神就給自己.
這件事真的很阻礙.
自己去親近神.
我從那時候開始.
有一個更加強的心智.
要讀聖經.
我記得我在.
做大手術.
有兩次進手術室.
第二次進去之前.
橫跨了幾個月.
我就答應神要完成我第一次的聖經閱讀.
在我進第二次手術室.
之前的早上.
我剛剛就讀到啟示錄的二十二章.
為什麼是二十二章.
為什麼不是二十三章.
為什麼.
因為最後那章就去到.
二十二章.
剛剛讀完.
如果你在想其他原因.
你錯了.
我想說的.
每個人都有點不同.
有些東西卡住你去親近神.

$^{681}$有個姐妹很難參加.
什麼聚會.
因為她要看著某隻寵物.
我不是說寵物不要.
千萬不要這樣仇視我.
其實我很喜歡動物.
我真的很喜歡動物.
不過我家裡.
不是所有成員都喜歡.
我就沒有了.
但我覺得很可惜.
很好的夏令會.
或者很好的一個陪靈會.
沒有人看了我的寵物.
我還以為是一個人.
所以.
寵物沒有人看.
這件事就攔住他.
但我告訴大家.
如果這件事真的在你裡面.
那件事要被打破.
那件事要被對付.
我不用說你真的不讓他.
但你要對付這件.
卡住你的東西.
好不好.
有時候真的有些.
做.
以前我在銀行做的時候.
有些交易員.
睡不著覺.
你不要說.
他做其他事他都睡不著覺.
看著三四個市.
日本美國.
還有英國.
那你不用睡了.
那他怎麼崇拜.
你這裡是白天.
再晚一點.

$^{721}$就有事了.
所以呢.
他知道有些事是攔了他自己.
他要對付他.
所以你心靈裡面.
會為罪難過.
有時候我們想不到.
想一下會不會有些事.
就像第二節所說.
就是神要去打傷他.
對付他的東西.
在我們生命裡面.
以致我們看到這件事.
其實是分隔了我們和神.
令到我們在罪之中.
都不留意.
我們可能做父母就會留意到.
一千年來經常打機.
但你有沒有留意你的子女.
或者我們的小孩.
都會留意到做父母的.
你又經常這樣經常這樣.
你都有.
你前類的罪.
影響你的敬拜.
但我們只會留意到.
其他的人.
你還在玩這個遊戲.
還在計算你的等級.
你沒有心機去讀書.
或者你沒有心機去敬拜.
我們每個人都要反省自己.
原來一個能夠去敬拜神的人.
真是神看顧的.
就是他能夠在神面前真的謙卑.
在他裡面看到.
自己虧欠神.
攔了他的罪.
那件事被打傷他.
激到他.

$^{761}$激中他.
有些人可能不是說他很多.
是那樣.
我們讀聖經.
那一個少年的官.
和主耶穌說.
主耶穌列出的那些誡命我從小都遵守了.
你不要跟我說謊.
但當主耶穌說出那一樣.
就是他對金錢的.
看法.
你將你所有的變賣.
分給窮人.
這個是他的致命.
他就和他死過.
他不是和主耶穌死過.
他只是.
悠悠愁愁地走了.
太緊要了.
如果要敬拜神.
要這樣放下.
他寧願繼續抓住.
所以他人就算在耶穌旁邊.
他並不能真的敬拜.
我們當中.
可能有人.
只是一個人在這裡.
心還是纏繞在其他的事.
不去.
真誠在神面前敬拜.
所以.
這樣的人有一個好處.
就是他這樣破碎.
這樣被打倒.
他就知道這是罪.
他真的會歸向神.
因為他知道罪有多可怕.
今天.
有很多詩歌.
很奇怪.

$^{801}$很開心.
我真的想問問那個人.
你這麼開心是為了什麼呢.
他又不說.
看那首詩歌.
知道他為什麼這麼開心.
知道他為什麼這麼快樂.
主耶穌在他的生命裡.
做了什麼.
我再說.
在聖經裡很多時候.
能夠有真正的心靈釋放.
就是有那種神的對付.
在裡面.
在《士徒行傳》說.
有三千人悔改.
那三千人悔改受針.
之前發生了什麼事呢.
就是有責心.
責心就像新約.
表達同樣的意思.
就是有些東西被神傷害.
如果神的說話.
是不會給你任何反應.
應該有反應是好的.
如果沒有任何反應.
不是一件好事.
所以.
為什麼有很多人來到教會.
聽很多道.
未信耶穌.
有一次我是一個長者團契.
跟他們聊福音.
有時候有些長者.
會在一個狀態.
你跟他說.
剛才才說完世人都犯了罪.
長者就會說.
是的很多人都犯了很多罪.
很和應你的說道.

$^{841}$是的.
好像那個XX.
他就犯了很多罪.
甚至得罪他.
人都是這樣的.
他又說了五分鐘.
接著我繼續說我的信息.
過了一段時間.
有一天.
他突然間.
跟他聊天.
我那一次.
我講到什麼呢.
我講到亞國.
舊約的亞國很狡猾.
很詐.
這麼狡猾又說謊.
是不是罪.
我想問你.
究竟你有沒有這樣說謊.
或者你覺得虛荒的罪.
你覺得你會不會都像亞國.
當然不是一下子.
第一眼見到他就問他.
是跟他都有些關係.
說多了.
很奇妙.
同樣一個長者.
這次他就說.
是的我也有罪.
什麼罪.
他就說了他的例子.
是不是很不同.
當你遇到一個人.
原來他覺得自己.
不需要拯救因為他沒有罪.
沒有病的人為什麼要醫生呢.
不同以前.
有些人心理問題.
不斷看四五六個醫生.

$^{881}$配了五六種藥.
有些吃有些不吃.
心理毛病.
但正常來說.
你有那個毛病你要去醫生.
一個人經常覺得自己.
沒有要對付的罪.
沒有得罪神他為什麼要拯救他呢.
所以.
哪些人看重神呢.
就是當他覺得自己沒有.
很貧窮的心靈.
在罪的狀態.
在神的面前.
很痛悔.
受傷.
最後.
這樣的人.
他親近這位神.
他會聽到神的每一句說話.
聖經說.
因主的說話.
是驚到震.
這樣的意思.
大概.
原來我們對主的說話有反應是好的.
我做牧者.
做傳道者.
我都希望我講完.
讓那些人有反應.
當然希望你沒有反應.
就是扔東西過來那些.
但是你的反應.
就是同意.
或者有些疑惑.
或者不明白.
或者你有你的一些.
高見 我都覺得很歡迎.
但是你知道有不少人.
的反應是什麼.

$^{921}$我有一次派單張.
那個人就是.
很.
四平八穩的站在那裡.
拿了我的單張.
又不笑.
又沒有生我的氣.
又沒有做任何事.
接著就.
是一張單張.
他就走了.
這是什麼反應.
不理會 莫不關心.
有些人就會反感.
不喜歡.
覺得你妨礙他.
侵犯他自由.
有些就生你的氣.
我認識一個弟兄.
他去一個地方.
做差傳 即是短宣.
在一個停車場派派下單張.
突然間.
有一群年輕人.
手拿著石頭.
衝過來.
你以為什麼 你玩十放豆袋.
扔他們.
他就走了.
那時他還坐得很爽.
他現在70多.
不知道有沒有跑得這麼快.
有不同的人有不同的反應.
有些人用理性跟你搏.
有些就是.
有不同的反應.
但有一種反應是神特別看上.
他對神說的每句說話.
有一個反應.
他會戰驚.

$^{961}$什麼叫戰驚.
戰驚可以反映出幾方面.
他覺得神的說話.
是有權威的.
是可以觸碰到他的生命.
是有權柄在我們身上.
你知道有很多人.
公然侮辱耶穌基督的名字.
用耶穌的名字亂說話.
他肯定不是戰驚.
是面對神的說話.
但是弟兄姊妹.
我們不需要害怕.
我有一次讀約翰福音.
主耶穌說我講的說話.
在審判那一天.
就是去審判他們.
就是說主耶穌.
要去審判這些人.
要審判的人.
包括我們和你.
早早都已經講了.
不用再作篇文.
來審判他們.
就是說已經定實.
但是.
有這樣的體見.
覺得主的說話.
我會很害怕.
我聽了真的會戰驚.
另外一個表現就是.
他會戰驚因為心存敬畏.
離敬拜主.
我覺得大大小小.
就好像舊約的.
認識神的秘訣.
要敬畏神.
我鼓勵我們當中.
做父母或祖父母.
在我們的生命.

$^{1001}$在我們的教導裡面.
注入去敬畏神的元素.
要給自己.
和身邊.
學習認識神的人.
知道去敬畏神.
從來在聖經裡面.
你見到那些敬畏神的人.
神很體重.
在敬畏神裡面.
這些人.
就是神調轉.
他竟然很有反應.
對神的說話.
講了那句說話.
很緊張.
聖經有什麼人會緊張主耶穌的說話.
不是每個人.
有些人聽到主耶穌的說話.
是沒有反應的.
但有些人是沒有很大反應.
之後是有反應.
我舉一個例子.
當主耶穌說.
你將會三次不認他.
他還裝聰明.
覺得自己信祂都不會.
到他真的三次不認主.
主耶穌.
坐在受審的位置.
看到他.
看他一眼.
他知道主知道.
他知道主講過的說話.
他又真的這麼驕傲.
現在他又真的跌倒.
他很緊張.
他已經害怕.
而且他心裡很難過.
他心裡有一種.

$^{1041}$如果沒有那種敬畏神的心.
他不會有那種反應.
弟兄姊妹.
原來就是這麼重要的.
最基本的三樣元素.
虛心.
痛悔.
有罪在裡面.
你知道在護教的時期.
在馬丁路德.
1517年改教的時候.
那時候他們重寫.
一些詩歌.
他們就很強調提到.
要畏罪痛悔.
要知道那個拯救.
不是那些聖禮.
而是耶穌基督的救恩.
但是那些人不會欣賞這個拯救.
除非他裡面真的覺得.
需要這位救主.
在新約.
路加呼音18章13節.
我很喜歡.
也很受感動.
當神祈禱的瑞麗.
祂的禱告.
在路加呼音18章13節.
很簡單.
祂說.
神啊 開恩可憐我這個罪人.
當我們來到去敬拜神的時候.
有沒有這個心呢.
當我們落在一些光景.
我們會不會.
只是介意被人看到我是怎樣.
而不去看到.
很需要在神面前.
真誠 求神憐憫.
開恩可憐.

$^{1081}$反而他旁邊.
這個.
以為自己是很有義的.
這個法利賽人.
你會看到神.
好像完全不理會這個.
以為自己在祈禱.
但其實是跟自己說話的一個人.
神是很.
看重一個這樣.
向他謙卑下來.
歸向他的人.
有一天我繼續去看這段的時候.
我說這類人.
在敬拜 親近神.
很看重.
他們心靈裡面.
看到自己空虛貧窮.
痛悔 看到自己的罪.
他們因為主所說的一言一語.
他們裡面會漸漸緊張.
我盼望我們都緊張.
帶給他們什麼呢.
如果在新約來說.
這些豈不是真誠悔改.
得到拯救的人.
弟兄姊妹.
有一天我們輪到.
自己傳福音的時候.
我們都掌握到這個秘訣.
有些人是真的沒去到.
信耶穌的 為什麼呢.
他還覺得自己跟神.
差不多高 差不多大.
他還覺得自己沒什麼錯.
他從來沒虛心.
他從來沒看到自己.
在神面前是徹底貧窮.
是需要神的拯救.
甚至會覺得.

$^{1121}$神的說話有什麼地方.
要這樣跟隨.
是不是這麼有權威.
不知道大家有沒有被人問過.
你怎麼知道聖經是真的.
有沒有想過.
覺得 糟了 不知道怎麼回答.
交第一招給大家.
為什麼你認為.
它不是真的.
扔給他看看.
再有別的方法.
可能就是.
你認不認證法.
你懂不懂.
他跟你講理性.
你就要給他講道理.
很多根據.
有機會再詳細講.
有很多人.
他都沒去到看重神的說話.
他可能.
很近都好.
他沒去到相信.
但是當他心靈準備好.
我們千萬不要.
不懂得脫國.
有些人已經有了這幾個元素.
你不幫他去信耶穌.
我們就失去了.
那個寶貴的時刻.
所以我們在報道裡面.
有一樣東西叫那個Precious Moment.
千萬不要錯過.
那個寶貴的時刻是什麼.
那個人準備好要去信耶穌.
當他真的虛心痛悔.
因神的懷疑賤經.
所以在《士多人傳》第二章.
那些被扎心的人.

$^{1161}$我們怎麼辦.
他們準備好了.
我們怎麼辦.
他們要得著耶穌基督的拯救.
你講清楚.
這個救恩的幾個重點.
他就可以.
決志了.
這樣的人有什麼好處呢.
我一定要講這個.
他有神的同在.
這個特別在57章15字那裡.
神就是要.
專門和這些人在一起.
我們來敬拜.
我們想不想神和我們在一起.
當然想啦.
是不是.
我告訴大家.
我照顧我的孫子的時候.
其中一個感覺就是玩具是次要.
請問做婆婆的.
你能不能回答.
什麼比較重要.
你不要跟我說退休金.
什麼重要.
你和他在一起.
你知不知道.
有一次我挺開心.
我有一張BB椅子.
是我兒子的時代坐的.
我想扔掉.
後來我兒子說不要扔.
歷史文物.
我孫子留意到.
他從此以後喜歡坐.
有一天騰空了旁邊的位置.
要求外公和他一起坐.
你會不會覺得很厭惡.
不要啦.

$^{1201}$小孩子坐在那裡幹嘛.
你是不是這樣.
我很開心.
因為他沒有找外婆坐.
先找外公坐.
你明白那種.
你這樣在神面前.
虛心痛悔.
去賤驚.
神體寵你.
神最後有一個心願.
我要和你一起.
我不知道你羨慕不羨慕.
其實我們人生.
有很多人人生很苦.
你走過的路.
就算有人想陪你.
也不能陪你.
我有一次要做檢查.
整個人進了一個洞.
做一個MRI或CT.
其實很恐怖.
好像進了一個山洞.
如果那個人按錯了按鈕.
沒有頭沒有腳.
你病的時候就什麼都想.
但我那次.
沒人陪你.
我祈禱主耶穌一定要陪我.
那個洞真的很黑.
還注射了一種不知什麼的藥.
好像搞肚子一樣.
主啊.
你真的要陪著我.
我真的盼望.
我們真的看到.
原來我們敬拜神.
有這位神在我們身邊.
是很寶貴.
最後我用我剛才說的.

$^{1241}$我在醫院那裡的見證.
那一朝.
看史錄22章的時候.
我很感恩.
如果這個手術不成功.
我就有一個快的方法.
去到神那裡.
我真的覺得很羨慕.
看到21章22章.
天家的美好.
我就將自己的祈禱交給神.
我記得那一朝很早會上來.
那一朝五點左右我已經醒了.
預備好.
靈修員.
還有最後那一套聖經.
我覺得很開心.
如果我現在有這個快的方法.
去見主就好了.
當然主的計劃不是我想的.
結果是.
出人意外地成功的.
所以接著我就再活多幾十年.
甚至今天還可以在這裡.
講到見證神的恩典.
但你說那段時間.
我真的覺得.
很寶貴.
我每天去親近神.
神真的是.
很甜蜜.
很寶貴.
所以在那個時候開始.
我就做了一樣很對不起我媽的事.
跟我媽說.
我們雖然很窮.
但暑假我真的不想去做主義工.
賺那些.
有一點錢.
但沒了那麼多時間.

$^{1281}$我看我人生有暑假的日子很少.
我就用那些日子.
就是我去操練禱告.
和讀聖經.
那些日子.
對於富裕的家庭.
是很容易的.
對於貧窮的家庭.
我做了這樣的決定.
我媽或我姐妹.
可能是有意見的.
但她沒有告訴我.
我自己也有一點內疚.
但我立了這個心智.
那幾年的暑假我真的這樣做.
我沒有後悔.
我很多的熟練的功課.
和神親密的體會.
是在那些時候.
你不是這樣定一些時間.
你和神一起.
你也沒有看到自己那種.
很敗壞.
或者很不妥的心態.
屬靈.
那種驕傲才恐怖.
告訴別人.
我看了多少本書.
我看了什麼書.
他就教我謙卑.
我真的覺得.
你明明人怎樣敗壞.
但敗壞.
不是人的理性.
是需要神的恩典.
你要親近神多一點.
你的力量才會強大.
我的時間差不多了.
我的總結.
就是.

$^{1321}$很簡單.
你願不願意.
做一個真正蒙主恩代的敬拜者.
可能你的答案.
是願意.
但如果你是願意的話.
你就不要skip那三個元素.
我們來溫習一下.
第一.
虛心.
在神面前貧窮.
很切無.
很需要祂.
第二 痛悔.
究竟什麼在你裡面.
你覺得要為祂痛悔.
真的需要神.
砌下祂.
最後 因著主的說話.
來到《戰經》.
當你這樣看《聖經》.
你就會.
什麼春夏秋冬.
什麼狀態都出現.
有時候你會很開心.
好像你看到.
啟示錄的那種歡欣.
有時候你會很懼怕.
你知道神是多麼的公義.
聖潔.
好像意菜下他自己感受到.
禍災我滅亡.
那個不是口號.
那個是體會到神有的反應.
我盼望我們一起去祈禱.
親愛的主耶穌基督.
願你來吸引我們.
當我們來敬拜你.
我們願意.
好像做了.

$^{1361}$有崇拜的活動.
我們祈求主讓我們.
真是.
虛心 痛悔.
在你的面前.
因著你的說話來到《戰經》.
因為是你跟我們說的.
你去指引我們.
我們當中還未信耶穌的主.
我求你在我們當中的.
朋友中間工作.
讓我們看見神是真是.
那位超越的神.
但祂願意這樣看顧我們.
就近我們.
願意這麼卑微歸向祂的人.
施恩.
向我們每一個都豐豐富富的.
來去恩代.
祈禱奉主耶穌基督的聖名.
\newpage



\section{列王記上 1:1-53}
\label{sec:031981}
\textbf{王位誰屬:人看智慧,神看應許}~~~~袁惠鈞牧師博士
\newline
\newline
連結: \href{http://www.church.com.hk/acms/content.asp?site=cdc\&op=show\&type=product\&code=031981\&layout=sermon}{\texttt{031981}} ~~~~ 語音日期: 2026-04-12
\newline
\newline
\hyperref[sec:031980]{\small{< < < PREV SERMON < < <}}
~
\hyperref[sec:index_chronic]{\small{[返順時目]}}
~
\hyperref[sec:031982]{\small{> > > NEXT SERMON > > >}}
\newline
\newline
列王記上 1:1-53
\newline
\begin{longtable}{cl}
\hline
\hline
章節 & 經文 (和合本修訂版)\\
\hline
1:1 & \begin{tabularx}{0.7\textwidth}{X} 大衛王年紀老邁，雖然蓋著外袍，仍不夠暖和。 \end{tabularx} \\ \\ \relax
1:2 & \begin{tabularx}{0.7\textwidth}{X} 臣僕對他說：「不如為我主我王找一個年輕的少女，侍立在王面前，照顧王，睡在王的懷中，好使我主我王得暖。」 \end{tabularx} \\ \\ \relax
1:3 & \begin{tabularx}{0.7\textwidth}{X} 於是他們在以色列全境尋找美貌的少女，找到了一個書念女子亞比煞，帶到王那裡。 \end{tabularx} \\ \\ \relax
1:4 & \begin{tabularx}{0.7\textwidth}{X} 這少女極其美貌，她照顧王，伺候王，王卻沒有與她親近。 \end{tabularx} \\ \\ \relax
1:5 & \begin{tabularx}{0.7\textwidth}{X} 那時，哈及的兒子亞多尼雅妄自尊大，說：「我要作王」，就為自己預備座車、騎兵，又派五十人在他前頭奔跑。 \end{tabularx} \\ \\ \relax
1:6 & \begin{tabularx}{0.7\textwidth}{X} 他父親從來沒有責怪他，說：「你為何這麼做？」他非常俊美，生在押沙龍之後。 \end{tabularx} \\ \\ \relax
1:7 & \begin{tabularx}{0.7\textwidth}{X} 亞多尼雅與洗魯雅的兒子約押和亞比亞他祭司商議；他們就順從亞多尼雅，幫助他。 \end{tabularx} \\ \\ \relax
1:8 & \begin{tabularx}{0.7\textwidth}{X} 但撒督祭司、耶何耶大的兒子比拿雅、拿單先知、示每、利以，以及大衛自己的勇士都不順從亞多尼雅。 \end{tabularx} \\ \\ \relax
1:9 & \begin{tabularx}{0.7\textwidth}{X} 亞多尼雅在隱‧羅結旁瑣希列磐石那裡獻牛羊、肥犢為祭，請了他的眾兄弟，就是王的眾兒子，以及所有作王臣僕的猶大人。 \end{tabularx} \\ \\ \relax
1:10 & \begin{tabularx}{0.7\textwidth}{X} 但他沒有邀請拿單先知、比拿雅和勇士們，以及他的弟弟所羅門。 \end{tabularx} \\ \\ \relax
1:11 & \begin{tabularx}{0.7\textwidth}{X} 拿單對所羅門的母親拔示巴說：「哈及的兒子亞多尼雅作王了，你沒有聽見嗎？我們的主大衛卻不知道。 \end{tabularx} \\ \\ \relax
1:12 & \begin{tabularx}{0.7\textwidth}{X} 現在，來，我給你出個主意，好保全你和你兒子所羅門的性命。 \end{tabularx} \\ \\ \relax
1:13 & \begin{tabularx}{0.7\textwidth}{X} 你去，進到大衛王那裡，對他說：『我主我王啊，你不是曾向使女起誓說：你兒子所羅門必接續我作王，他必坐在我的王位上嗎？亞多尼雅怎麼作了王呢？』 \end{tabularx} \\ \\ \relax
1:14 & \begin{tabularx}{0.7\textwidth}{X} 看哪，你還在那裡與王說話的時候，我會隨後進去，證實你的話。」 \end{tabularx} \\ \\ \relax
1:15 & \begin{tabularx}{0.7\textwidth}{X} 拔示巴進入內室，到王那裡。那時，王很老了，書念女子亞比煞正伺候著王。 \end{tabularx} \\ \\ \relax
1:16 & \begin{tabularx}{0.7\textwidth}{X} 拔示巴向王屈身下拜，王說：「你要甚麼？」 \end{tabularx} \\ \\ \relax
1:17 & \begin{tabularx}{0.7\textwidth}{X} 她對王說：「我主啊，你曾向使女指著耶和華－你的神起誓：『你兒子所羅門必接續我作王，他必坐在我的王位上。』 \end{tabularx} \\ \\ \relax
1:18 & \begin{tabularx}{0.7\textwidth}{X} 現在，看哪，亞多尼雅作王了，你，我主我王卻不知道。 \end{tabularx} \\ \\ \relax
1:19 & \begin{tabularx}{0.7\textwidth}{X} 他獻許多牛羊、肥犢為祭，請了王的眾兒子和亞比亞他祭司，以及約押元帥，他卻沒有請王的僕人所羅門。 \end{tabularx} \\ \\ \relax
1:20 & \begin{tabularx}{0.7\textwidth}{X} 但你，我主我王啊，以色列眾人的眼目都仰望你，等你告訴他們，在我主我王之後誰坐你的王位。 \end{tabularx} \\ \\ \relax
1:21 & \begin{tabularx}{0.7\textwidth}{X} 若不然，我主我王與祖先同睡的時候，我和我兒子所羅門必列為罪犯了。」 \end{tabularx} \\ \\ \relax
1:22 & \begin{tabularx}{0.7\textwidth}{X} 看哪，拔示巴還與王說話的時候，拿單先知也進來了。 \end{tabularx} \\ \\ \relax
1:23 & \begin{tabularx}{0.7\textwidth}{X} 有人奏告王說：「看哪，拿單先知來了。」拿單進到王面前，臉伏於地，向王叩拜。 \end{tabularx} \\ \\ \relax
1:24 & \begin{tabularx}{0.7\textwidth}{X} 拿單說：「我主我王，你果真說過『亞多尼雅必接續我作王，他要坐在我的王位上』嗎？ \end{tabularx} \\ \\ \relax
1:25 & \begin{tabularx}{0.7\textwidth}{X} 他今日下去，獻了許多牛羊、肥犢為祭，請了王的眾兒子和軍官們，以及亞比亞他祭司；看哪，他們正在亞多尼雅面前吃喝，說：『亞多尼雅王萬歲！』 \end{tabularx} \\ \\ \relax
1:26 & \begin{tabularx}{0.7\textwidth}{X} 至於我，就是你的僕人，和撒督祭司、耶何耶大的兒子比拿雅、王的僕人所羅門，他都沒有請。 \end{tabularx} \\ \\ \relax
1:27 & \begin{tabularx}{0.7\textwidth}{X} 這事果真出於我主我王嗎？王卻沒有告訴僕人們，在我主我王之後誰坐你的王位。」 \end{tabularx} \\ \\ \relax
1:28 & \begin{tabularx}{0.7\textwidth}{X} 大衛王回答說：「召拔示巴到我這裡來。」拔示巴就來，站在王面前。 \end{tabularx} \\ \\ \relax
1:29 & \begin{tabularx}{0.7\textwidth}{X} 王起誓說：「我指著救我性命脫離一切苦難的永生的耶和華起誓。 \end{tabularx} \\ \\ \relax
1:30 & \begin{tabularx}{0.7\textwidth}{X} 我既然指著耶和華－以色列的神向你起誓說：你兒子所羅門必接續我作王，他必繼承我坐在我的王位上，我今日必這樣做。」 \end{tabularx} \\ \\ \relax
1:31 & \begin{tabularx}{0.7\textwidth}{X} 於是，拔示巴屈身，臉伏於地，向王叩拜，說：「我主大衛王萬歲！」 \end{tabularx} \\ \\ \relax
1:32 & \begin{tabularx}{0.7\textwidth}{X} 大衛王又說：「召撒督祭司、拿單先知、耶何耶大的兒子比拿雅到我這裡來！」他們就都來到王面前。 \end{tabularx} \\ \\ \relax
1:33 & \begin{tabularx}{0.7\textwidth}{X} 王對他們說：「要帶領你們主的僕人，讓我兒子所羅門騎我自己的騾子，送他下到基訓。 \end{tabularx} \\ \\ \relax
1:34 & \begin{tabularx}{0.7\textwidth}{X} 在那裡，撒督祭司和拿單先知要膏他作以色列的王；你們也要吹角，說：『所羅門王萬歲！』 \end{tabularx} \\ \\ \relax
1:35 & \begin{tabularx}{0.7\textwidth}{X} 你們要跟隨他上來，使他坐在我的王位上，他要接續我作王。我已立他作以色列和猶大的君王。」 \end{tabularx} \\ \\ \relax
1:36 & \begin{tabularx}{0.7\textwidth}{X} 耶何耶大的兒子比拿雅回應王說：「阿們！願耶和華－我主我王的神這樣說。 \end{tabularx} \\ \\ \relax
1:37 & \begin{tabularx}{0.7\textwidth}{X} 耶和華怎樣與我主我王同在，願他照樣與所羅門同在，使他的王位比我主大衛王的王位更大。」 \end{tabularx} \\ \\ \relax
1:38 & \begin{tabularx}{0.7\textwidth}{X} 於是，撒督祭司、拿單先知、耶何耶大的兒子比拿雅，以及基利提人和比利提人都下去，讓所羅門騎上大衛王的騾子，送他到基訓。 \end{tabularx} \\ \\ \relax
1:39 & \begin{tabularx}{0.7\textwidth}{X} 撒督祭司從帳幕中取了盛膏油的角來，膏所羅門。他們就吹角，眾百姓都說：「所羅門王萬歲！」 \end{tabularx} \\ \\ \relax
1:40 & \begin{tabularx}{0.7\textwidth}{X} 眾百姓跟隨他上來，吹著笛，大大歡呼，地被他們的聲音震裂。 \end{tabularx} \\ \\ \relax
1:41 & \begin{tabularx}{0.7\textwidth}{X} 亞多尼雅和所有的賓客剛吃完，聽見這聲音；約押聽見角聲就說：「城中為何有這響聲呢？」 \end{tabularx} \\ \\ \relax
1:42 & \begin{tabularx}{0.7\textwidth}{X} 他正說話的時候，看哪，亞比亞他祭司的兒子約拿單來了。亞多尼雅說：「進來吧！你是個賢明的人，必是來報好消息的。」 \end{tabularx} \\ \\ \relax
1:43 & \begin{tabularx}{0.7\textwidth}{X} 約拿單回答亞多尼雅說：「我們的主大衛王已經立所羅門為王了！ \end{tabularx} \\ \\ \relax
1:44 & \begin{tabularx}{0.7\textwidth}{X} 王派撒督祭司、拿單先知、耶何耶大的兒子比拿雅，以及基利提人和比利提人和所羅門一起去，叫他騎上王的騾子。 \end{tabularx} \\ \\ \relax
1:45 & \begin{tabularx}{0.7\textwidth}{X} 撒督祭司和拿單先知已經在基訓膏他作王了。他們從那裡歡呼著上來，城都震動，這就是你們所聽見的聲音。 \end{tabularx} \\ \\ \relax
1:46 & \begin{tabularx}{0.7\textwidth}{X} 所羅門也已經登上國度的王位了。 \end{tabularx} \\ \\ \relax
1:47 & \begin{tabularx}{0.7\textwidth}{X} 王的臣僕也來為我們的主大衛王祝福，說：『願神使所羅門的名比你的名更尊榮，使他的王位比你的王位更大。』王在床上屈身敬拜， \end{tabularx} \\ \\ \relax
1:48 & \begin{tabularx}{0.7\textwidth}{X} 王也這樣說：『耶和華－以色列的神是應當稱頌的，因他今日賞賜一個人坐在我的王位上，我也親眼看見了。』」 \end{tabularx} \\ \\ \relax
1:49 & \begin{tabularx}{0.7\textwidth}{X} 亞多尼雅所有的賓客都戰兢，起來，各走各路去了。 \end{tabularx} \\ \\ \relax
1:50 & \begin{tabularx}{0.7\textwidth}{X} 亞多尼雅懼怕所羅門，就起來，去抓住祭壇的翹角。 \end{tabularx} \\ \\ \relax
1:51 & \begin{tabularx}{0.7\textwidth}{X} 有人告訴所羅門說：「看哪，亞多尼雅懼怕所羅門王。看哪，他抓住祭壇的翹角，說：『願所羅門王先向我起誓，必不用刀殺死僕人。』」 \end{tabularx} \\ \\ \relax
1:52 & \begin{tabularx}{0.7\textwidth}{X} 所羅門說：「他若作賢明的人，連一根頭髮也不致落在地上；他若作惡，必要死亡。」 \end{tabularx} \\ \\ \relax
1:53 & \begin{tabularx}{0.7\textwidth}{X} 於是所羅門王派人叫亞多尼雅從壇上下來，他就來向所羅門王下拜。所羅門對他說：「你回家去吧！」 \end{tabularx} \\ \\
[1ex]
\hline
\hline
\end{longtable}
$^{1}$「列王記」的第二講.
要講「列王記」上的第一章.
又講回舊約.
又要講故事,這個講古佬.
由故事當中我們找到神學的思想.
這個就是舊約.
新約就很直接了.
我今天給他的題目叫做.
「王謂屬誰,人看智慧,神看應許」.
其實「列王記」上的文筆是非常精彩.
精彩到不得了.
在「列王記」裡的第一章.
有兩個字是重複出現的.
一個字就是「Yada」.
就是「認識」,「認知」.
甚至可以解釋做親密的關係.
記不記得我跟大家都講過.
在創世記,亞當認識夏娃就出現什麼事.
搞出人命,生了蓋引.
這個就是「認識」,「認知」.
就是希伯來文的「Yada」這個字.
但是另外有一個字在第一章.
重複出現的是什麼.
就是「Shaba」.
就是「世」,「起世」.
這個字都是重複出現.
我今天的題目就是.
「人看智慧,神看應許」.
人來講是靠我們的智慧,我們的知識.
甚至與人之間有親密的關係,搞關係.
但是神是看什麼的?.
神是看祂的應許.
如果祂曾經「起世」的話.
祂曾經是神的誓言,就一定會發生.
這個就是我們今天的題目.
說開又說.
有沒有看到我今天的封面,漂亮嗎?.
我也懂得用AI來幫我設計我的封面.
我絕對不會用AI來幫助我有關講道的內容.
AI來講,你叫它弄一幅圖給你.

$^{41}$弄一個片段給你.
沒得談.
甚至你要查一查.
在聖經裡面這件事有沒有發生過.
這些是很準確.
但是神學來說,實在是不敢恭維.
你要什麼神學觀,AI都懂得給你看.
所以有很多時候我們要很小心.
現在在網上,你聽很多的講道.
都可能是AI弄出來的.
你說牧師英,我就不會了.
我會選那些講員是我喜歡的.
都會很危險的.
我告訴你,我很喜歡聽John MacArthur講道.
有一次我在網上看到一篇道.
我還沒聽過.
我就去聽.
那個視頻是John MacArthur.
聲音也是John MacArthur.
為什麼我說他講出來的東西.
完全跟他平時講的不同呢?.
相違反的呢?.
後來才發覺原來有人用他的樣子.
用他的聲音弄了視頻出來,弄了一篇道.
不是他,是AI.
所以我們現在在網上聽道都要很小心.
你不要以為這個就是我喜歡聽的牧師.
分分鐘是AI,是假的道理.
現在聽道你要很小心.
你要知道認識那個網頁.
你要認識那個頻道是信得過.
你知道不會有假的訊息,你去聽.
這樣就對了.
但是需要有些智慧.
列王記一開始的時候.
大衛已經年紀老邁.
他就快離開這個世界.
而他的王朝正面臨一個很大的危機.
就是有兩股勢力要爭取他的王位.
在列王記裡面所記載的.

$^{81}$不單只是宮廷裡面的明爭暗鬥.
也讓我們看到因為以色列人不聽神的話.
神在申命記第六章給人民有他的教訓.
在申命記第七章給以色列的王有他的教訓.
但是以色列人都不聽.
最後就應驗了在聖經裡面的預言.
以色列人會遭到滅國.
亡國會被擄.
這一切都會被實現.
當大衛的王朝正面臨一個很大的危機.
通常當我們面臨危機的時候會怎樣?.
會反映出一個人的本性.
就好像太陽曬泥土.
可以曬到泥土很堅硬.
同樣的太陽也可以曬到冰上,令到冰融解.
究竟我們遇見危難的時候.
我們的心會被軟化,好像冰變成水.
還是我們的心會變得更加堅硬.
好像泥土越來越曬得乾,變得堅硬.
究竟會怎樣呢?.
以色列人就是變得堅硬,很可惜.
其實今天的列王記上的第一章.
就有一個非常精彩的首尾結構.
我不會用時間一點一點和大家來看這個首尾結構.
如果你們喜歡的話,我會留在powerpoint那裡.
你們可以自己download來看.
但是在這個首尾結構,我想指出有兩件很重要的事.
就是經文最中心的地方.
就是列王記上的一章28至31節.
大衛向拔士巴重申他的誓言.
記不記得我說,在列王記上有兩個字很重要.
Yada,認識,認知,還有一個Shava,起誓.
他再重申,告訴拔士巴他曾經向他起的誓.
這個首尾結構的第二層,是說出兩個重要的人物.
這兩個重要的人物的名字都有連帶的關係.
一個就是拿丹先知,希伯來文的拿丹的意思是禮物.
第二個重要的人物,在這個首尾結構的下半部.
是若拿丹,這個名字就是耶和華所賜的禮物.
其實無論是拿丹先知,無論是若拿丹這個使者.
都是神賜給大衛,神賜給以色列人的很重的禮物.

$^{121}$今天我的大綱也是很簡單,一共是五點.
你可以說我的大綱是基於剛才的首尾結構.
所以我的大綱也是一個首尾結構.
一頭一尾都是說亞當尼亞,他開始的時候是一個投機主義者.
他要爭奪王位,但是去到最後第五點,又是說亞當尼亞.
東窗事發,他的野心被發覺,於是怎樣?.
由一個投機主義者變成一個乘機的求生者,他要求生.
這個大綱的第二點和第四點就是我剛才所說的.
拿丹先知和若拿丹這兩個人.
所以我剛才已經解釋了.
拿丹先知是要提醒拔仕巴大衛曾經許下的誓言.
這個就是Shabah這個字.
而若拿丹成為大衛的使者,將所羅門成為王的事實.
宣告給以色列人,讓他們知道.
知道這個字就是Yada.
所以你會見到這兩個字重複出現在今天的經文裡面.
而最中心的經文是甚麼?.
大衛本來不知道亞當尼亞是要搶奪王位.
現在他知道了,他就向拔仕巴重申他的誓言.
說他將會將所羅門放在他的王位上,繼承他的王位.
所羅門就是拔仕巴和他所生.
我們今天首先看第一點.
就是這個投機主義者亞當尼亞.
有人說當一個屬靈的領袖.
面對危機的時候,他會想甚麼?.
我怎樣可以為我的人民做些甚麼?.
但是一個機會主義者.
當他面對危機的時候,他會怎樣?.
他會問一個問題,就是我會為自己可以做甚麼?.
這個就是投機主義者和一個屬靈領袖不同的地方.
一個投機主義的領袖,多數是無垂自薦.
他在自己的面上貼金,他將所有的利益歸於自己.
這個就是亞當尼亞.
亞當尼亞是誰?.
亞當尼亞是大衛第四的兒子.
記不記得亞當尼亞在今天的經文事件發生的時候,他應該是35歲.
他是大衛第四的兒子.
大衛的大兒子是誰?.
是暗聯,因為他強姦了自己的妹妹他媽.
被他的弟弟阿沙隆將他殺死.

$^{161}$他第二個兒子是基利亞,又叫做達爾尼.
學者一般都認為這個基利亞是英年早逝.
因為在聖經裡面根本沒有提到他的事.
但是我就有另外一個看法.
我經常都注意到他的媽媽就是那位既美麗又聰明的阿比蓋.
這個女人很聰明,她可能一早已經提醒兒子.
你不要和你的哥哥爭權奪位.
你要明哲保身,你要潔身自愛.
所以在聖經裡面沒有提到基利亞的事.
可能他就是聽他媽媽說沒有和他的哥哥爭奪王位.
以致他沒有提到他有糾紛的事.
大衛第三個兒子就是阿沙隆.
我們都講過阿沙隆背叛大衛.
是被大衛的元帥約納所殺死.
現在大衛背叛最高的兒子是誰?.
就是亞當尼亞,因為他的哥哥已經死了.
他覺得自己應該坐在大衛的王位上.
他也覺得大衛應該指定他就是繼承人.
他覺得如果大衛一旦駕崩.
如果群龍無首,以色列沒有指定的王.
就會有很大的紛爭.
亞當尼亞和他的哥哥一樣.
趁大衛病的時候就起來爭奪王位.
亞當尼亞低估了他爸爸大衛的能力.
低估了大衛的智慧.
而為他的驕傲付出了很沉重的代價.
經文一開始的時候就說大衛年老怕冷.
於是他的大臣在全國物色一個很美貌的年青童女.
來侍候大衛王.
和合本就說是照顧.
其實是要他睡在王的懷中.
在聖經裡面睡在懷中是很親密的關係.
曾經撒萊亞伯拉罕的妻子.
是叫他的婢女夏甲.
將夏甲放在亞伯拉罕的懷中.
出現了什麼事?.
就生了以撒瑪利出來.
在《創世記》第十六章第五節.
所以我們見到睡在懷裡是非常親熱的一個字眼.
在聖經裡面有另外一個字是一樣的意思.

$^{201}$就是認識,認知這個字.
但我剛才跟大家說「Yada」這個字.
意思就是在聖經裡面同房.
也都解釋做同房.
就是我剛才說「阿當認識夏娃生了蓋人」.
所以認識這個字是很親密的認識.
可見這個阿彼撒不是.
我講了找到阿彼撒嗎?.
是找到阿彼撒.
所以阿彼撒不是一個普通的僕人.
他其實是大衛的妾侍.
他睡在大衛的懷中其實是名正言順.
經文就很有趣了.
所以我說「列王記」是文筆非常精彩.
在這裡似乎作者是對比兩個睡在大衛懷中的女人.
還有一個是誰呢?.
是拔仕巴.
當拿丹仙芝來找大衛講出他的罪的時候.
拿丹仙芝用了一個比喻.
她說有一個很貧窮的牧羊人.
他只有一隻母羊.
這隻母羊就睡在他的懷裡.
結果就被這個牧羊人的鄰居.
一個很有錢的人搶了這隻母羊.
這個有錢人是誰呢?.
就是大衛.
那隻母羊是誰呢?.
就是拔仕巴.
在《薩姆異常》的十二章.
現在這個作者又微妙地講出.
阿彼撒是睡在大衛的懷中.
有沒有看到他比較這兩個女子.
在其中暗示出來.
就是因為大衛和拔仕巴的罪.
現在引起皇室有紛爭,有悲劇.
將會有一場風波會發生.
在大衛的後人的身上.
大衛死了之後.
這個阿彼撒將會成為一個很重要的人物.
在列王記上的第二章.

$^{241}$亞當尼亞和所羅門又會有一場爭奪.
今次不是爭奪王位.
第一章是爭奪王位.
第二章爭奪甚麼呢?.
爭奪阿彼撒.
以致亞當尼亞會連生命都失去.
經文說大衛沒有和阿彼撒親近.
有沒有看到又是這個字「yada」.
認識,親近.
大衛沒有和阿彼撒親近.
在古代近東.
一個王的統治能力和他的性機能是成正比.
作者告訴我們.
大衛沒有和阿彼撒有關係.
即是說他已經沒有性機能.
亦告訴我們他沒有統治能力.
這亦是原因.
所以亞當尼亞認為他的父親已經沒有影響力.
他便起來爭奪王位.
但他錯了.
如果他可以多點關心他的父親大衛.
如果他可以和他的兄弟搞好關係.
又贏得領袖的支持.
可能他真的能夠成為下一任的以色列王.
但他沒有這樣做.
亞當尼亞效法他的哥哥亞沙隆.
他妄自尊大.
在列王記上的第一章第五節.
這個字在聖經裡面出現.
都是負面的意思.
亞當尼亞的方法和亞沙隆一模一樣.
叫戰車和五十個士兵在他的面前奔跑.
這個就是王的架子.
在撒母耳夏的十五章第一節.
亞沙隆都曾經有這樣做.
作者亦暗示.
大衛其實沒有說過亞當尼亞就是他的繼承人.
這一切都是亞當尼亞自己做出來的.
我們從這裡就見到亞當尼亞的野心.
因為他在聖經裡面第一句說的話就是.

$^{281}$我必作王.
亞當尼亞是一個非常英俊的人.
就好像亞沙隆一樣.
所以他得到他父親大衛的寵愛.
在述列王記上的一章六節.
經文說大衛素來沒有使他憂悶.
你怎樣能夠令到你的子女憂悶?.
你去管教你的子女他們就會憂悶.
經文即是告訴我們甚麼?.
無論亞沙隆,無論是亞當尼亞.
大衛都沒有管教他們.
這個就是大衛的問題.
他不是一個好的爸爸.
他寵愛亞當尼亞.
就好像他寵愛亞沙隆一模一樣.
亞當尼亞就招募了將軍約納.
和大祭司阿比亞他來支持自己.
也就是說他得到部分的軍事支持.
甚至是宗教的支持.
這兩個人約納和阿比亞他.
都曾經事奉大衛多年.
尤其是當大衛落難而逃.
是最艱難的時候.
他們和大衛是同甘共苦.
但是現在因為要納下一個王位的繼承人.
而居然和大衛倒戈相向.
亞當尼亞去到引羅傑的旁邊.
一塊叫做所希烈的盆石那裡.
做什麼?宰殺牛和羊.
將肥牛毒獻給神.
來到獻祭.
是有一個很盛大的宴會.
經文最初的時候沒有告訴我們.
這個宴會的原因是什麼.
目的是什麼.
其實這個和亞沙隆.
他騙人說自己要去希伯倫獻祭.
是一模一樣.
去到希伯倫記不記得他做什麼.
自稱為王.

$^{321}$其實和亞當尼亞一模一樣.
都是有政治的動機.
經文是很有意思.
用了一個雙重的名稱.
來稱呼亞當尼亞的兄弟.
叫他們做什麼.
首先是朱弟兄.
然後是王的眾子.
這個雙重的稱呼.
使這個聚會更加懸疑.
令我們明白這個宴會其實有兩個目的.
不單止這個是一場家庭的聚會.
這些是朱弟兄.
這個也是有政治的動機在後面.
因為這些都是王的眾子.
有沒有看到.
這個就是一個雙重的稱呼.
在列王記上的一章十節開始.
就讓我們看到這個宴會的目的.
是透過那些沒有被亞當尼亞邀請人的名單.
有什麼人沒有被亞當尼亞邀請.
有大衛的先知拿丹.
也有祭司薩督.
還有護衛長比娜雅.
還有被稱為王的朋友的侍媒.
大衛的衛兵勇士都沒有人被邀請去參加這個聚會.
但是最重要沒有出現的是誰.
是所羅門.
因為所羅門所有的兄弟都有出席.
只有所羅門自己沒有出席.
在列王記是第一次提到所羅門這個名字.
這些人沒有參加這個宴會代表什麼.
這些人是忠心於大衛.
他們都不是擁納亞當尼亞的人.
這一場將會成為一個加冕的盛宴.
所有參加這一場聚會的賓客.
他們都願意擁納亞當尼亞成為以色列人的王.
或者他們都認為.
雖然大衛沒有出現.
但是可能因為大衛病了.

$^{361}$其實可能大衛已經按守在亞當尼亞的身上.
任命他成為下一任以色列的王.
而所有亞當尼亞的兄弟.
除了所羅門都參加.
為什麼.
因為他們要讓亞當尼亞知道.
我們根本沒有心跟你爭奪王位.
這個王位是你要就行了.
我們不會跟你爭奪.
在這裡我們看到.
當薩都這個祭司和拿丹沒有出現的時候.
就代表什麼.
代表大衛根本沒有支持這件事.
因為薩都和拿丹都是大衛的心腹.
亞當尼亞也應該知道.
所羅門就是神所揀選成為下一任的以色列王.
為什麼我這樣說.
因為他沒有請所羅門來這個宴會.
所以他應該知道所羅門會不高興.
因為神已經揀選所羅門成為下一任以色列的王.
其實亞當尼亞,亞比亞他和約納都背叛了神.
因為是神的旨意要所羅門承繼大位.
他們都忘記了詩篇的三十三篇.
耶和華的籌算永遠立定.
在聖經裡面當我們讀到有人謀朝篡位.
或者做這些用手段來達到目的的時候.
我們就好像覺得以賽亞書的五十九章十四節所說.
經文說什麼.
公平轉移退後.
公平都退後了.
公義站在遠處.
離開我們很遠.
還有誠實在街上撲倒.
我們經常說忠忠直直就忠雖黑食.
原來這句話有聖經出典.
在哪裡?.
在以賽亞書的五十九章十四節.
但是當這些事發生的時候.
我們要記得詩篇的九章十四節.
惡人必因自己手所做的事陷入網羅.

$^{401}$神仍然在他的寶座上來在掌權.
就是因為這件事提醒了大位.
是時候要任命所羅門成為以色列下一代的王.
其實那丹先知一早就知道.
所羅門是神所揀選下一任的以色列王.
他怎麼知道?.
因為神曾經差派那丹.
將所羅門的另外一個名字帶給大位.
這個名字是什麼?.
耶底底亞.
意思就是主所愛的.
在撒慕爾夏的十二章二十四至二十五節.
所以所羅門就是神所愛.
所以那丹知道.
所羅門就是神所揀選下一任的王.
我們講完了第一點.
就是投機主義者亞當尼亞.
第二點就是這位公義的先知那丹.
那丹立刻知道了這件事.
就去找拔士巴.
他說:你知不知道亞當尼亞已經成為王了?.
這個對你和所羅門將會造成很大很大的威脅.
如果亞當尼亞真的坐正的話.
你們要小心了.
而亞當尼亞日後都會去找拔士巴.
可見得拔士巴是一個非常有影響的女人.
亞當尼亞會去找他做什麼?.
第二章我剛才說.
他要拔士巴答應將比薩許配給自己.
他又去找拔士巴.
你說拔士巴為什麼會這麼有影響力?.
他的丈夫是誰?.
是大衛.
而拔士巴是大衛最愛的女人.
所以每個人都知道你要什麼.
你去找拔士巴,走後門.
拔士巴於是就去見大衛.
見到大衛年紀老邁.
同女亞比薩正在他旁邊侍候他.
那丹就教拔士巴怎樣和大衛講話.

$^{441}$在列王記上的一章十三節.
你進去見大衛王對他說.
我主我王,你不曾向婢女起誓說.
你兒子所羅門必接續我作王.
坐在我的位上摩.
現在亞當尼亞怎麼作了王呢?.
很好的問題.
你說大衛何時向拔士巴發誓.
說所羅門會坐在他的王位上?.
有可能是因為神曾經和大衛立約.
在撒母耳夏的第七章.
說他的後人當中有一個會坐在他的王位上.
會為神起聖殿.
這是誰?就是所羅門.
而大衛在歷代至上的二十二章和二十八章.
曾經宣佈要為神起聖殿.
他指名門聖說所羅門會為神起聖殿.
所以很多學者說.
就是在這裡拔士巴認為這就是對他的一個誓言.
另外有一個可能性.
就是我剛才說過神曾經給所羅門有一個很特別的名字.
這就是耶底底亞,就是主所喜悅.
所以從這一點我們也可以看到.
所羅門就是大衛的繼承人.
拔士巴去找大衛.
告訴他.
王啊,你好無知啊.
無知這個字又是誰?.
又是丫達這個字.
在聖經裡面已經被我們看到.
這個字用在大委身上兩次.
第一次說出什麼?.
他是無能.
他沒有統治的能力.
現在他是無知.
亞當尼亞成為王他都不知道.
所以拔士巴就指出.
亞當尼亞已經舉行了宴會.
你的元帥約鴨.
你的祭司亞比亞他都參加了.

$^{481}$還有你所有的王子除了所羅門都參加.
你居然不知道這件事.
這裡被我們看到.
亞當尼亞真的有很大的野心.
拔士巴就告訴大衛.
亞當尼亞真的有野心成為王.
現在全國的人民在做什麼?.
在等待你的宣佈.
在等待你的決定.
你要告訴人.
誰將會繼承你的王位.
拔士巴當然暗示給大衛聽.
你不搞清楚這件事.
你將來去世.
我和你的兒子所羅門.
我們就不得了.
命都沒了.
在這個時候.
拿丹才知道自己也去見大衛.
說出所羅門設立宴會的事.
拿丹和拔士巴兩個都再說出.
所有被邀請人的名單.
和沒有被邀請人的名單.
而最後拿丹問了一個很重要的問題.
在列王記上的一章27節.
他問大衛.
你有沒有讓你的僕人知道.
誰會繼承你成為以色列的國王.
有沒有看到「也達」這個字又出現了.
不單止經文指出大衛是無能的.
又是「也達」.
不單止經文指出大衛是無知.
根本對亞當尼亞這件事.
現在拿丹是要大衛作出一個決定.
讓人知道「也達」這個字又再出現.
你要讓僕人知道誰會繼承你.
否則將來會有很大的問題.
拔士巴也好像在責問大衛.
我和你已經沒有一個兒子了.
記不記得大衛和拔士巴行奸淫.

$^{521}$已經一個嬰兒沒有了.
拔士巴好像在暗示給大衛.
你是否想沒有多個兒子.
這就是所羅門.
而的確是因為大衛和拔士巴的罪.
皇室現在又出現問題.
記不記得拿丹曾經告訴大衛.
因為你和拔士巴的罪.
刀劍將會不離你的家.
現在就是這個問題又再重複出現.
我們說了兩點.
亞當尼亞就是一個投機主義者.
拿丹就是那位公義的先知.
現在是我們經文最中心的部份.
就是本來不知情的大衛.
你現在知道了有甚麼回應.
這就是我們第三點.
我說過經文圍繞著兩個字.
「Yada」我們已經見到了.
重複出現在大衛的身上尤其是.
但是現在大衛知情之後.
他就要起誓.
就是「Shava」這個字.
另外一個在經文重複出現的字.
他要起誓將所羅門放在他的皇位上.
來繼承他成為下一任的意識烈王.
「Shava」這個字在列王記上.
一共出現有十二次這麼多.
有八次都是在列王記上的第一章.
和第二章.
而這個字明字的名詞是甚麼.
起誓是動詞.
出現十二次.
名詞只出現一次.
在列王記上的二章四十三節就是誓言.
而這個字和拔士巴的名字是息息相關.
拔士巴的名字甚麼意思.
「Bat Shava」甚麼意思.
是誓言的女兒.
我以前時時都不明白.

$^{561}$為何這個人的名字叫誓言的女兒.
因為在聖經裡面.
一個人的名字通常和他息息相關.
現在我知道為何他叫拔士巴.
誓言的女兒.
就是大衛將會向他起誓.
所羅門一定會坐在他的皇位上.
繼承他成為王.
大衛指著救贖他脫離一切的耶和華起誓.
他一定會指定所羅門成為他的繼承人.
而他這個誓言.
不單止是他的誓言.
還是誰的誓言.
是耶和華的誓言.
在十五二下的第七章.
和在列王記上的第一章三十節.
大衛知道他應該當機立斷.
因為他已經見到亞當尼亞的野心.
他會搶奪皇位.
如果他有一天成為了以色列人的王.
這樣拔士巴和所羅門就一定是實死無生.
而兩個兒子當中.
有沒有見到亞當尼亞就是反叛的那個.
所羅門是比較聽大衛的話.
所羅門不單止是王子.
不單止是大衛的繼承人.
他和大衛都是有親密的關係.
為什麼他和大衛有親密的關係.
因為他的媽媽就是拔士巴.
是大衛最愛的女人.
所以他和大衛當然有親密的關係.
大衛於是照他的忠心僕人.
先知那丹,祭司撒督和他的護衛長比那亞.
這些都是他信得過的人.
指示他們要去羈奮.
公開宣佈所羅門將會成為以色列人的王.
原來耶路撒冷是有兩個水源.
一個就在隱羅傑.
隱羅傑我剛才說是亞當尼亞在那裡自立為王的地方.
但是還有一個水源在哪裡.

$^{601}$在羈奮.
大衛都很心碎清.
為什麼選擇羈奮.
因為耶路撒冷有兩個水源.
而羈奮才是最主要的水源.
是重要過隱羅傑.
去那裡做什麼.
要所羅門坐在大衛的驢驅子的上高.
證明他就是大衛所揀選.
就好像耶穌基督也是坐著驢驅子進入耶路撒冷城.
隱羅傑其實在石安山東面的山坡上高.
距離羈奮是很近很近.
不到一里的路程.
所以很快所羅門成為王的消息.
就會傳到亞當尼亞那裡.
薩督和拿丹都用會幕裡面的高遊來高沒所羅門.
然後士兵就吹響號角.
宣佈有一件很重要的事已經發生了.
所羅門現在已經成為以色列和猶大的國王和統治者.
為什麼牧師把這兩樣東西分得這麼清.
你是國王你不是已經成為統治者了嗎.
其實這個大衛搞得很清楚.
在聶王紀上的一章35節他宣佈所羅門是以色列的統治者.
也是王子這個字.
但是在聶王紀上的一章34節.
他說所羅門已經被任命成為以色列和猶大的國王.
Malak這個字.
為什麼要分得這麼清.
因為無論是素羅或者是大衛.
當他們是被高納為王的時候.
他們只不過是被高納成為統治者.
未是王子.
大衛是後來他的人民才擁納他為王.
在薩姆以下的第五章第六章和第七章.
大衛和素羅不同的地方是什麼.
大衛一向都知道以色列真正的王是誰.
是耶和華自己.
他只不過是被神授高差派.
來成為一個統治者.
成為一個王子.

$^{641}$所以他知道王一向都是耶和華.
為什麼在這裡他要這麼指名門姓.
要指定所羅門要成為以色列的國王.
因為亞當尼亞已經被他見到有很大的野心.
他知道他應該要當機立斷.
事不言辭.
他一定要立克納所羅門成為全以色列人的王.
經文接著就有四個萬歲在三個人的身上.
這三個人是誰.
第一個就是亞當尼亞.
在列王記上的一章二十五節.
當他自稱為王之後.
人民就叫他做萬歲.
這是第一次.
第二次在列王記上的一章三十一節.
是拔事巴掌去見大衛.
稱呼大衛做萬歲.
還有第三次.
就是在列王記上的第一章三十四節.
大衛說我要所羅門成為以色列人的萬歲.
而在第三十九節.
人民真的稱所羅門為萬歲.
有沒有見到王位正確的傳承應該是怎樣.
是透過大衛這個在位的萬歲爺.
按立下一任的萬歲爺.
這個才是真正傳承的途徑.
而不是像亞當尼亞那樣自立為王.
經文在這裡讓我們看得很清楚.
我剛才說到大衛心腹之一是誰.
比拿亞.
非常之重要.
因為這個人重握兵權.
很重要.
我在說大衛勇士名單的時候.
我也說過這個比拿亞.
他其實是祭司的兒子.
通常祭司都是繼承他爸爸成為祭司.
但是這個比拿亞.
他希望是他揀選了軍事的生涯.
他希望成為一個軍人.

$^{681}$他終於成為了大衛的勇士.
還成為了大衛的護衛長.
他本來是一個沿海居住的人.
就好像斐理士人.
來自地中海.
但是他後來就投靠了大衛.
加入了大衛的軍隊成為大衛的勇士.
當比拿亞聽見大衛的吩咐之後.
他立刻派兵保護大衛和保護亞當尼亞.
保護這個所羅門.
因為他知道亞當尼亞可能對他們父子不利.
後來當約納他支持亞當尼亞.
反叛了大衛之後.
被所羅門處決了約納這個元帥.
就把元帥這個職位給了比拿亞.
薩都,拿丹和比拿亞在大衛的指示之下.
宣佈所羅門成為以色列人的王.
人民都很高興.
歡天喜地.
彈奏樂曲.
還怎樣?.
高聲呼喊.
神保守所羅門王.
聲音之大.
是怎樣?.
好像地震一樣.
地都裂了.
這些聲音在山谷裡回音.
傳了去引羅傑那裡.
引羅傑那裡又發生什麼事?.
那裡的人民又在大聲呼喊.
亞當尼亞王萬歲.
兩邊都說王萬歲.
但似乎大衛那邊的聲音比較響.
他們都不知道耶路撒冷發生什麼事.
是不是大衛死了?.
為什麼這麼大聲?.
是不是有人來攻打以色列?.
攻打耶路撒冷?.
這樣就帶領我們去到第四個重要的人物.

$^{721}$就是目擊的使者約拿丹.
Jonathan.
約拿丹真的可以說不可以說.
我剛才說.
祭司亞比亞他已經投靠了亞當尼亞.
就是他的爸爸.
他的爸爸就是亞比亞他.
居然投靠了亞當尼亞.
但約拿丹他棄暗投明.
他知道所羅門才是神所揀選.
所以他支持大衛.
他支持所羅門.
當亞當尼亞正在慶祝.
人民正在呼喊亞當尼亞王萬歲.
突然間聽到很大的聲音.
約納問城裡面發生什麼事?.
就在這個時候.
約拿丹就來到帶信息給那裡的人.
作者又巧妙地.
將亞當尼亞和亞沙隆反叛的事跡放在一起.
為什麼?.
因為這個約拿丹.
在兩件事當中都是一個重要的使者.
而兩件事發生都在引羅傑這個地方.
當亞當尼亞反叛大衛趕他出耶路撒冷.
當亞沙隆反叛大衛趕大衛出耶路撒冷的時候.
曾經約拿丹成為大衛的使者.
將很重要的信息帶給大衛.
就在這裡引羅傑.
叫他要逃命.
因為亞沙隆已經在追殺他.
現在是同一個地方引羅傑.
同一位的使者帶出來的信息很不同.
這個信息是由大衛傳給大衛的敵人.
在亞沙隆的時候是傳給大衛.
但是現在就是傳給大衛的敵人.
所以有沒有看到.
其實列王記尚文采是非常非常之的美麗.
亞當尼亞以為約拿丹幫他帶來好消息.
但是結果帶給他是最壞的消息.

$^{761}$透過約拿丹作者第三次講出.
大衛怎樣安納所羅門成為以色列人的王.
所以約拿丹真是一個目擊的證人.
他親眼見到所羅門騎著大衛的驢驅子.
他親眼見到所羅門是被薩都和拿丹用油來高抹他.
講得非常之仔細,非常之準確.
所以約拿丹真是一個目擊的證人.
但是有一件事有很多學者都提出異議.
是在列王記尚的一章47節至48節.
其實是描述皇宮所發生的事.
很多學者都說約拿丹有沒有去過皇宮.
他怎知道皇宮發生什麼事.
首先我們要知道這兩節的經文是講什麼.
一節在列王記尚的一章47節.
就是說那些大臣來祝賀大衛.
他們說「願王的臣使所羅門的名比王的名更尊榮.
使他的國位比王的國位更大」.
大衛就在床上屈身下拜.
接著在列王記尚的一章48節經文說.
「王又說耶和華以色列的臣是應當稱重的.
因他賜我一人今天坐在我的位上.
我也親眼看見了」.
這些事都發生在宮殿裡.
約拿丹怎會知道.
如果我們注意這兩節經文一開始.
都是一個希伯來字「監」這個字.
「監」這個字是什麼意思.
就是英文的Moreover.
又或者是中文的載者或者還有.
但是在很多其他的經文的意思是什麼.
是報信者,通信者.
所以極有可能就是有報信者將這些消息.
講給約拿丹知道.
所以他知道在皇宮裡面發生什麼事.
我也說過列王記真的是一部文筆非常精彩的一部巨著.
列王記是將歷史,文學,神學融會貫通.
放在一起.
我們今天的經文絕對是歷史的事實.
但是作者也是透過很精彩的文筆.
讓我們看到一開始就有一個很美麗的首尾結構.

$^{801}$接著是用兩個字來指出整個文章的神學觀在哪裡.
第一個字是什麼?.
Yada.
「認知認識」.
就好像這個公義的先知拿丹.
他真的走去質問大衛.
你有沒有讓你的僕人知道.
誰會繼承你成為以色列的國王.
這個就是列王記上的一章二十七節.
也是知道這個字.
你有沒有作好決定誰會繼承你成為王.
然後下一個字就是「起誓」.
是透過使者若拿丹.
我剛才也說過拿丹和若拿丹也是有很微妙的意思.
兩個人的名字都是出於同一個字根.
透過若拿丹說出大衛是指定所羅門成為他的繼承人.
於是大衛在列王記上的一章四十八節說.
神是應當稱重的.
因他賜我一人.
今日坐在我的位上.
我也親眼看見了.
接著的經文就是我們最後的一點.
又是說亞當尼亞.
他已經不再是那一個滿有野心的投機主義者.
但他已經成為落難而逃的乘機求生者.
當若拿丹告訴大家.
所羅門已經成為以色列人的王.
亞當尼亞和他的客人都知道.
大能臨頭了,今次大禍了.
他的賓客和那些天真的王子都落難而逃.
跑回耶路撒冷城來逃命.
亞當尼亞更厲害.
他走到那個帳棚.
這個帳棚我相信不是會幕.
因為會幕是機編.
這裡是大衛為藥櫃所建的帳棚.
用來儲藏藥櫃.
亞當尼亞相信跑去那裡.
那裡有個祭壇.
他就抓住祭壇的兩隻角.

$^{841}$這是猶太人的一個傳統.
我們知道神叫以色列人.
在民數記的三十五章.
建了有六個屠城.
讓那些人伸張正義.
任何人被人追殺都可以走去屠城.
雖然屠城已經被建立.
但是亞當尼亞不夠膽跑那麼遠.
耶路撒冷這個帳棚比較近.
所以他立刻去那裡抓住兩隻牛角.
什麼意思?.
是以色列人的傳統.
任何人抓住兩隻牛角.
當他被人追殺的時候.
都可以延遲判決.
可以讓這個人有一個公平的聆聽.
這個就是這個意思.
角是代表什麼?.
是代表權柄.
這個就讓我們看到.
耶和華最終是掌權那個.
你們任何誰都要讓人有一個公義的聆聽.
這個就是亞當尼亞.
抓住兩隻牛角不肯放.
直至到要所羅門起誓不殺死他為止.
「起誓」這個字又出現了.
有沒有看到?.
又是Shabaa這個字.
所以我說.
列王記的文筆真的很精彩.
所羅門就憐憫他的哥哥.
他誓言不會殺死他.
讓他存活.
叫他回去他自己耶路撒冷的家裡.
其實這個就等於軟禁.
所羅門就派了衛兵監視著亞當尼亞.
所羅門也說.
如果你有叛亂越界.
你有野心.
你將會被處決.

$^{881}$所羅門對亞當尼亞說.
你若作忠義的人.
連一根頭髮也不至落在地上.
你若行惡必要死亡.
亞當尼亞就向所羅門下拜.
但是他是不是真心服從這個造王的弟弟?.
當然不是.
他甚至不服從耶和華.
我們怎麼知道?.
下一章的書他又會出來.
跟所羅門爭奪阿比薩.
結束的時候很簡單.
這是我的信息.
就是兩個字.
Yada.
就是認識,認知,親密的關係.
甚至可以解釋為智慧.
Shaba.
是喜誓,誓言.
甚至我們可以說是應許.
我們先從兩個角度來看.
首先救贖.
是不是靠我們自己的智慧.
甚至靠我們的行為?.
不是.
是就是認識.
不是認識人.
認識那位神.
還要有親密的認識.
我們要知道認識這位神.
我們要走在祂的應許裡面.
祂的應許是什麼?.
所有接受耶穌基督的人都會得救.
這就是神的應許.
所以是靠神的應許.
祂的喜誓.
是靠我們認識這位萬佑的真神.
這是救贖.
你說處事呢?.
我們處事又如何?.

$^{921}$都是一樣.
不是靠你自己的智慧.
不是靠你自己的知識.
或者你認識哪一個人.
是我們認不認識這位神.
所以認識這個字仍然有它的意思.
這就是我們處事.
我們要認識神的旨意.
我們才能夠成功.
我們才能夠走在神的應許當中.
我們不認識神的旨意.
我們沒辦法走在神的應許.
因為當我說我們要在神的應許.
很多人就問.
我們怎樣知道神的應許?.
你要認識神.
認識神的話語.
認識神的旨意.
願意我們今天.
這兩個字給我們是有鼓勵.
我們要認識我們的神.
跟我們的神有一個親密的關係.
然後我們要知道祂的起誓.
祂的誓言.
知道祂的應許.
這就是從祂的話語我們能夠知道.
願神祝福祂自己的話語.
讓我們低頭禱告.
在列王記上精彩的文筆.
給我們看到有兩個很重要.
很有意思的字.
一個就是認識.
Yada.
就是親密的關係.
願意我們今天能夠繼續學習認識你.
跟你是有一個親密的關係.
願意我們今天能夠認識你的應許.
就好像另外那個字.
Shabbat.
在聖經裡面有你的誓言.

$^{961}$有你的應許.
願意我們今天能夠透過你的話語.
能夠認識這一切.
以致我們能夠走在你的旨意當中.
願意我們今天能夠像施暴真一樣.
能夠渴望成為像耶穌基督一樣完美.
一樣聖潔.
願意我們今天能夠在有生之年.
繼續追求你的話語.
追求你的聖潔.
以致當你回來的時候.
能夠被你稱為一個忠心良善的僕人.
將那榮耀還賺得歸給你.
我們今天禱告.
奉主耶穌基督的命而求.
阿們.
\newpage



\section{列王記下 6:8-23}
\label{sec:031982}
\textbf{請開我雙眼}~~~~曹子光牧師
\newline
\newline
連結: \href{http://www.church.com.hk/acms/content.asp?site=cdc\&op=show\&type=product\&code=031982\&layout=sermon}{\texttt{031982}} ~~~~ 語音日期: 2026-04-12
\newline
\newline
\hyperref[sec:031981]{\small{< < < PREV SERMON < < <}}
~
\hyperref[sec:index_chronic]{\small{[返順時目]}}
~
\hyperref[sec:031983]{\small{> > > NEXT SERMON > > >}}
\newline
\newline
列王記下 6:8-23
\newline
\begin{longtable}{cl}
\hline
\hline
章節 & 經文 (和合本修訂版)\\
\hline
6:8 & \begin{tabularx}{0.7\textwidth}{X} 亞蘭王與以色列作戰，他和臣僕商議說：「我要在某處某處安營。」 \end{tabularx} \\ \\ \relax
6:9 & \begin{tabularx}{0.7\textwidth}{X} 神人派人到以色列王那裡，說：「你要小心，不要從某處經過，因為亞蘭人下到那裡去了。」 \end{tabularx} \\ \\ \relax
6:10 & \begin{tabularx}{0.7\textwidth}{X} 以色列王派人到神人告訴他的地方去。神人警告他，他就在那裡有所防備，不止一兩次。 \end{tabularx} \\ \\ \relax
6:11 & \begin{tabularx}{0.7\textwidth}{X} 亞蘭王因這事心裡氣憤，召了臣僕來，對他們說：「我們當中有誰幫助以色列王，你們不告訴我嗎？」 \end{tabularx} \\ \\ \relax
6:12 & \begin{tabularx}{0.7\textwidth}{X} 有一個臣僕說：「不，我主，我王！只有以色列中的先知以利沙，把王在臥房所說的話告訴以色列王。」 \end{tabularx} \\ \\ \relax
6:13 & \begin{tabularx}{0.7\textwidth}{X} 王說：「你們去查看他在哪裡，我好派人去捉拿他。」有人告訴王說：「看哪，他在多坍。」 \end{tabularx} \\ \\ \relax
6:14 & \begin{tabularx}{0.7\textwidth}{X} 王就派遣車馬和大軍往那裡去，夜間他們到了，圍困那城。 \end{tabularx} \\ \\ \relax
6:15 & \begin{tabularx}{0.7\textwidth}{X} 神人的僕人清早起來出去，看哪，車馬軍兵圍困了城。僕人對神人說：「不好了！我主啊，我們該怎麼辦呢？」 \end{tabularx} \\ \\ \relax
6:16 & \begin{tabularx}{0.7\textwidth}{X} 神人說：「不要懼怕！因與我們同在的比與他們同在的更多。」 \end{tabularx} \\ \\ \relax
6:17 & \begin{tabularx}{0.7\textwidth}{X} 以利沙禱告說：「耶和華啊，求你開他的眼目，使他能看見。」耶和華開了這年輕人的眼目，他就看見了，看哪，滿山有火馬和火焰車圍繞以利沙。 \end{tabularx} \\ \\ \relax
6:18 & \begin{tabularx}{0.7\textwidth}{X} 亞蘭人下到以利沙那裡，以利沙向耶和華禱告說：「求你擊打這國，使他們眼目失明。」耶和華就照以利沙的話，擊打他們，使他們眼目失明。 \end{tabularx} \\ \\ \relax
6:19 & \begin{tabularx}{0.7\textwidth}{X} 以利沙對他們說：「這不是那條路，也不是那座城。你們跟我走，我必領你們到你們要尋找的人那裡。」於是他領他們到了撒瑪利亞。 \end{tabularx} \\ \\ \relax
6:20 & \begin{tabularx}{0.7\textwidth}{X} 他們進了撒瑪利亞，以利沙說：「耶和華啊，求你開這些人的眼目，使他們能看見。」耶和華開了他們的眼目，他們就看見了，看哪，是在撒瑪利亞城中。 \end{tabularx} \\ \\ \relax
6:21 & \begin{tabularx}{0.7\textwidth}{X} 以色列王看見他們，就對以利沙說：「我父啊，我真的可以擊殺他們嗎？」 \end{tabularx} \\ \\ \relax
6:22 & \begin{tabularx}{0.7\textwidth}{X} 他說：「不可擊殺！這些人豈是你用刀用弓擄來給你擊殺的呢？當在他們面前擺設飲食給他們吃喝，讓他們回到他們主人那裡。」 \end{tabularx} \\ \\ \relax
6:23 & \begin{tabularx}{0.7\textwidth}{X} 王為他們預備了盛大的宴席。他們吃喝完了，王就送他們回到他們主人那裡。此後，亞蘭的軍隊不再侵犯以色列地了。 \end{tabularx} \\ \\
[1ex]
\hline
\hline
\end{longtable}
$^{1}$在我們日常的生活中.
往往都充滿了很多很多的威脅.
可能是來自身體的檢查.
突然間看一個醫生.
醫生告訴你報告不是很靚仔.
你的心情馬上跌了下來.
可能是公司裁員.
又或者是你學校裡面被人排擠.
很多不開心的事.
甚至家庭裡面關係.
是有一些的衝突.
當我們陷入了這些困難的處境的時候.
我們的眼睛.
我們的眼目看到的是些什麼東西呢.
是現實裡面的殘酷.
還是看到上帝的同在呢.
今天想和大家看一段比較長的經文.
不過不會全部照讀的.
就是來自獵王幾下第六章第八節到第二十三節.
這裡提到一個故事.
這個故事在開始的時候.
就提到亞蘭王和以色列人在爭戰.
故事的結束就是亞蘭王不再犯境了.
就是說這個故事叫做Happy Ending.
這個故事裡面大概可以分為三個部分.
第一個部分就和亞蘭王有關的.
就說到他見到一件事件.
但是他完全沒有敬畏上帝.
第二個事件是.
以利沙的僕人一個年輕的僕人.
和以利沙的一些對話.
第三件事件就是以色列王的反應.
也是見到一些事.
但是他完全沒有任何敬畏的心.
兩個王對神和神的工作沒有半點的敬畏.
而上帝的僕人以利沙和以利沙的僕人.
和以利沙的對話就成為了今天的故事的焦點.
神人以利沙開了少年人的眼.
讓他看見屬靈的事情.
今天我們想透過這段經文.

$^{41}$去學習信心的幾樣很重要的事.
第一樣就是信心是必須有一個觸摸到.
見到屬靈視野的擴張.
你要見到屬靈視野的擴張.
第一樣就是看見上帝才是真正的掌控者.
阿難王在皇宮那裡密室有一個很重要的會議.
這個會議他以為沒有人可以知道.
因為只有內閣的人才可以聽到.
所以經文用「臥室」這個字.
就是說一個很秘密的地方.
最安全的地方.
但是想不到神人以利沙.
竟然可以預早知道裡面到底說了些什麼.
他們想著這班人開個會.
有一個秘密的消息.
就是說以色列的王在哪裡的時候.
他們就會突襲當地附近的環境.
這個的線報基本上是密不透風的.
但是竟然之間.
以利沙可以知道.
打發一兩次告訴以色列王.
提醒他不要經過那個地方.
所以敵人沒辦法可以去陷害到以色列的王.
阿難王這個密室的秘密.
基本上是很周傳的.
沒有人會知道.
但是竟然之間給以利沙去知道.
舉個例子.
舉個例子就是.
在一個臥室裡面.
你可以想一下.
如果我們今時今日在這裡開一個會.
只有這裡的人才知道.
這個是一個絕對的秘密來的.
沒有人可以知道我們在說什麼.
除非我們中間有人洩漏.
所以你要明白.
為什麼這個以利沙.
為什麼她這麼厲害在一個密不透風的裡面.
都能夠知道這些秘密呢.

$^{81}$而且我想你留意的就是.
這一段的經文裡面.
提到的侵略不是一些什麼大戰.
只不過是搶一下東西.
干擾一下以色列周邊的小地方.
搶一下貨 搶一下人.
分分鐘不需要動刀動兵.
阿難王的目的基本上不是要買起以色列的王.
而是想掃蕩一下.
以色列一些細小的地方.
做一下擄掠.
令到他們的經濟蒙受一些大的損失.
這樣也是另類去打擊這個王的一些做法.
以色列的王當然聽從了以利沙的說話.
受到警惕.
而且做了防備.
所以經文就說提防未受這個的其害.
可能就是王派了一些兵部署.
加強了保安.
又或者叫了當地的百姓.
趕快撤出都未定.
好 我們看看第六章第十一節.
阿難王又怎麼樣呢.
我們一齊讀一讀第六章第十一節.
好嗎 請讀.
我們這類有誰幫助以色列王.
你們不只給我嗎.
這個經文提到.
這個皇帝有一個反應就是心裡敬異.
敬異這個字是非常特別的.
這個字的意思就是打颶風.
打大風.
所以他心裡面刮起颶風.
這一股的颶風似乎又驚但又很憤怒.
因為他估計裡面有二五仔.
我們都不難明的 是不是.
你想想 讀廣東話就明白了.
沒有鬼不死得人.
你明白嗎.
假設我們在這裡開一個大會.

$^{121}$剛才我們也說了.
我們一些秘密說出來.
我們的秘密就是.
一會兒我們決定去城寄飲茶.
不去這個home.
這個秘密只有五堂的人才知道.
但是想不到連祖堂的弟兄姊妹.
都會知道.
竟然我們不去home就去城寄.
你說這裡是不是有人洩密.
一定有人洩密的.
可能已經傳了一些信息出來了.
Benny說的.
今天我們去城寄.
我們不去這個home.
所以沒有鬼是不死得人的.
所以一定有二五仔.
所以他就跟那幫內閣說.
你們交人出來.
肯定是你們自己一個做的.
如果你是他們內閣的時候.
你們會怎樣.
肯定是個個震騰騰的.
六章第十二節就有一個神僕說.
沒有人幫他.
只有以色列中有一個先知.
以利沙.
將王的臥房所說的話.
告訴了以色列王.
這裡很明顯.
因為矛頭指向有內鬼.
肯定是這幫內閣的人.
肯定是害怕的.
個個都有可能.
於是肯定是推一個外人出來.
受靶的.
以利沙這麼出名.
推一個外人出去受靶.
就是最對的了.
所以很有可能是這樣的原因.

$^{161}$當然我們不可以排除.
也是神利用了他們的口.
去表達出以色列中有神和先知.
上帝自己.
當然重敬個人造衛星.
他知道了人心裡面的秘密.
甚至你們怎樣開會.
都準確知道在哪裡.
還有說了些什麼.
今時今日我們都容易明白.
因為我們經常聽到.
某些國家又給人斬首元首.
準確知道你住哪個房間.
你住哪間酒店哪個房間.
他都知道就是炸那個房間.
這個世界真的有一些這麼高的科技.
不知道你有沒有.
我以前的年代.
以前我到現在.
最近我才知道原來可以這樣.
以前我都沒聽過這樣的事件發生.
但是神重敬.
上帝重敬的地方就是.
連他們說什麼都知道.
他們心裡面想什麼都可能知道.
所以你要見到上帝重敬.
上帝真正是掌管一切的.
真正掌控的是上帝.
不是亞蘭王.
有些人以為他可以掌控到一個海峽.
不是可以將整個世界的世界.
可以控制了.
有些人以為有錢.
有人就可以控制了新聞的媒體.
有些人以為有權力.
就可以控制不同的人.
令到他跪低,將勁牽.
有些人以為掌握了科技,軍力.
可以橫行無忌.
到底哪一個才是掌權者呢?.

$^{201}$屬靈的信心.
是要建立在上帝對現實世界.
絕對的主權之上.
以利沙的介入.
不是因為以色列的王求他.
絕對不是.
而是上帝主動的憐憫.
上帝不單止記得他的子民的苦況.
也讓整個世界的人知道.
整個世界和宇宙的主宰.
誰是真正掌權者.
還有第二樣東西.
我們要看到的就是.
上帝才是真正的主.
我想你留意經文裡面.
有一個動詞經常出現.
就是「差遣」.
今天很正確.
請「差遣」我.
「差遣」這個字重複出現了三次.
以利沙「差遣」.
以色列王「差遣」.
亞蘭王「差遣」.
如果一個人「差遣」另一個人.
這是在說權力的問題.
有權去做這件事.
但是最終的掌權者是誰呢?.
不是人.
而是上帝.
不是地上任何一個君王.
而是天上的主.
還有第三樣東西.
看到什麼呢?.
我們還要看到上帝的同在.
還有保護.
我們一起看看第六章的第十四節.
我們一起讀出來.
「夜間到了,惠昆」.
這裡很明顯.
亞蘭王他很憤怒.

$^{241}$下令要捉拿以利沙.
就派出他的精銳部隊.
去到多灘這個地方去捉他.
咦?.
其實亞蘭王會相信是以利沙做的.
那就是了.
不然捉來做什麼呢?.
相信以利沙竟然可以做到一件很神奇的事情.
就是他們在密室裡說的事情.
他竟然知道.
這個神真是神通廣大.
否則就不會下令去捉人.
不過.
當你聽到這些經文的時候.
你就會感受到.
他雖然知道有一股神奇的力量.
但是他完全對於上帝沒有半點敬畏的心.
有些人經常覺得.
看到神蹟我就相信了.
看蹟是不是這樣?.
這段經文告訴我們.
並非如此.
有些人看到明明有一股神奇的力量.
甚至他可能了解到真的有一個神存在.
但是他們仍然都可以相信.
仍然都可以做一些很無恥的事情.
這個皇帝就是這樣.
亞蘭王他沒有帶有半點的心.
沒有產生半點敬畏的態度.
反而用軍事的力量去捉拿以利沙.
有些人似乎背後明白到有一些很難明的力量.
甚至有些真的可能經歷到神都未定.
但是就好像這個亞蘭王一樣.
沒有產生對信仰或者對神的敬畏.
亞蘭王彷彿就好像法老王一樣.
那種心硬.
下令現在就去捉以利沙.
想去生擒他.
嘗試去游說他.
還是想給多點錢去辯解.

$^{281}$游說不成功才鋸了他.
經過神目的調查之後.
明查暗訪.
結果就知道鎖定了以利沙住的地方.
就叫做多丹這個地方.
將這個城市重重圍困.
現在可以說以利沙插翅難飛.
不過大家有沒有想過.
有時候有些國家的防空系統真的很弱雞.
橫衝直撞.
完全好像沒有染雞籠一樣.
以色列就奇怪了.
不知道你有沒有留意.
長驅直進.
他的精銳部隊掃入以色列的多丹這個地方.
亞蘭軍落以色列這個地方都有條路行.
有段時間都要先去到.
還要長驅直進進入多丹這個地方.
完全沒有任何系統攔截.
沒有大軍去阻止.
為什麼.
證明了以色列的軍隊有多廢.
所以當時以色列的軍隊是完全起不到任何作用.
甚至可能是任由它進來都未定.
沒有任何抵抗之下.
就被亞蘭大軍.
外來的大軍包圍著多丹和以色沙的住處.
六章第五十五節第二天的早上.
服侍神的僕人.
服侍以利沙的僕人就起來起身.
拉開窗簾一開.
全部都是敵人的車馬.
嘩!這次大件事了.
所以經文就告訴我們.
他說看一看.
是一種很緊張的表達.
他看見的就是亞蘭的大軍壓境.
所以這位年輕人.
他就跟著跟神人說.
大鑊了.

$^{321}$死了這次.
我想問大家.
你什麼時候說過大鑊.
這三個廣東話字.
有沒有說過.
經常都說的.
大鑊了.
什麼時候會說.
你會不會說.
升職的時候.
大鑊了.
升職了這次.
都是的.
哦!OK!收到.
如果是學生.
我想問.
如果你考到A.
你會不會大鑊了.
考到A.
不會吧.
嗚~.
我要A+.
什麼家長來的 要A+.
什麼時候會說大鑊.
什麼時候會說死了.
當然是一些絕望的表達.
原來這六章第五字出現過一次.
就說到有一個先知的僕人.
借了一個斧頭回來.
掉了到海裡.
於是就說大鑊了.
那是貴的東西.
賠不起的.
所以現在耶利沙這位的僕人.
同樣發出一個這麼絕望的反應.
神人的僕人.
其實他是信主的.
不要忘記他是信主的.
但是他一開槍點看到的東西.
和耶利沙看到的東西不同.

$^{361}$為什麼呢.
大家都有一雙眼.
為什麼他看到的東西.
和別人看到的東西.
會不一樣呢.
他的僕人看到的是大軍壓境.
人的眼光往往是會被威脅來佔據.
當恐懼,當亞蘭人的大軍佔滿了他的視網膜的時候.
他的屬靈的眼光就關閉了.
僕人看到的是環境.
耶利沙看到的是掌管環境的上帝.
兄弟姐妹.
你的眼睛看到的是什麼呢.
所以六章第十六字.
神人就這樣回應.
不如讀一讀這節這麼好的經文.
請讀.
神人說不要懼怕與我們同在的比與他們同在的更多.
這句話真是精彩.
耶利沙沒有痛罵這個僕人一番.
說他真的沒有鬼使用.
又或者說他一點信心都沒有.
比那些不信更壞.
沒有 完全沒有.
反而他說不用懼怕.
這個字和看見的原文是押韻的.
要解除僕人的懼怕.
就是要令他看見.
而耶利沙強調與我們同在的比與他們同在的更多.
我們要記得.
信心不是要否認現實.
而是要加上一層屬靈的視線去看現實.
信心的眼睛會和我們說.
除了敵人之外.
我還看到我們的神.
不知道我們的人生有沒有像多灘城市一樣.
這個清早這麼糟糕呢.
那種壓力你正在經歷的東西.
會不會比鐵騎這個圍城更厲害呢.
而他忘記了.

$^{401}$我們還有屬靈的軍隊.
已經一路守護著我們.
僕人彷彿就代表了大多數信徒的狀態.
我不是不相信神的存在.
但是現實困難嘛.
債務纏身.
卡數不斷來.
一會又有病.
又有很多不同的不開心的事件.
考試又失敗.
大多的程度.
搞到我現在暫時.
靈性失明.
直到以利沙祈禱.
神開了他雙眼.
才體會到屬靈的真實性.
原來有火車火馬一直都在這裡.
只不過是他被肉眼.
那種感受遮住了.
我想你留意.
單單只是和這個僕人說不要懼怕.
肯定是不足夠的.
所以以利沙為了去堅固祂.
為了祂就去祈禱.
所以第十六節就說去祈禱.
然後第十七節就說垂聽了祈禱.
於是僕人抬頭一看.
就看到滿山都是火車.
滿山都是火馬圍繞著以利沙.
我想你留意這個僕人.
在第十五節他不是說看看.
車馬圍繞著這個城.
也是圍繞著他家.
但是第十七節的下半節.
其實原文一開始都是說看看.
所以僕人和以利沙說你看看.
滿山原來有火車火馬圍繞著.
所以艇姐妹.
敵人從來都沒有消失過.
環境從來都沒有改變過.

$^{441}$但是僕人現在看到的.
上帝派來的火車火馬.
塞滿了整個城市圍著他們住的地方.
這就是信心的奇蹟.
信心的神蹟不是要改變環境.
而是改變我們對環境的眼看體見.
而是改變我們對環境的體見.
六宗第十七節提到.
用以利沙祈禱的字眼.
開,眼睛和看見.
證明了神聽了和應允了以利沙的祈禱.
而且更加肯定了一件事.
就是神一定會解救以色列人.
這位僕人現在出之.
看得見上帝的同在和保護.
在十九世紀的時候.
有一位蘇格蘭的宣教士名叫Patton.
他被神呼召到南太平洋一個極為危險的地方做宣教.
當時裡面都是一些未開發的民族.
他們的民族有一個甚至是食人族的稱呼.
對於這一對夫婦.
因為他和他的老婆一起去的.
對他們來說是一種很大的威脅.
敵人,不是敵人.
這個部族其實都是對他們充滿了敵意的.
但是當然這對宣教士對他們沒有敵意.
但是他們就對這對宣教士充滿了敵意.
想他們,現在他們死了.
於是就燒他們的屋.
就在晚上燒他們的屋.
想他們死.
面對著四面楚歌.
毫無逃生的情況之下.
也沒有任何武器.
他們兩夫婦只能夠跪在地上.
整晚祈禱.
祈禱睡著了.
他將他們完全交給神.
到了早上一起來.
沒事.

$^{481}$不止沒有死.
他打開門發現圍著他的人全部走了.
這件事令他們很難理解.
直到一年之後.
其中一個酋長信了耶穌之後.
他就告訴他們.
原來那一晚這班酋長或者這班土著.
就見到有一班很巨大的人.
又發光的人.
手持著武器包圍著宣教士的家.
所以他們極其驚慌.
於是這班酋長和土著馬上就去走路.
弟兄姊妹我們見不見到.
我們見不見到當我們面對困難的時候.
上帝其實真的很真實去守護著我們.
當然我們不要忘記.
神的保護未必好像這樣.
保護其實有可能有另一個層次.
你記不記得《詩篇》23篇有提到.
我雖然行過死刃的幽谷.
原來我們有死刃的幽谷.
上帝沒有保證我們不會有死刃和幽谷.
上帝保證就是你經過死刃幽谷的時候.
不需要害怕.
因為我與你同在.
這個才是聖經的教導.
我們要留意神母答應不會有幽谷.
神答應讓我們在幽谷裡面與我們同在.
神的保護首先不是將你帶離開危險的困境.
而是在困境裡面去托著你.
信心怎樣可以看見上帝的保護呢.
大概有三個很簡單的東西.
讓我們學習.
第一就是神限制了困難的程度.
神母拿去了困難.
但是祂至少令到你還能站在那裡.
至少也沒有令你完全崩潰.
甚至至少你仍然存留在這個世界上.
還有第二.
神托著你不讓困難吞噬你.

$^{521}$困難仍然在.
但是你沒有被它毀滅.
這也是一個很重要的事情.
當然第三.
神會利用困難來塑造我們.
這是最深層的保護.
有時候神保護的不只是你人生的狀況.
更加保護了你的生命沒有被浪費.
沒有被罪扭曲.
有時候我們以為神的保護是將困難拿走.
但是聖經再一次提醒我們.
神更深的保護是將困難的裡面托著我們.
限制困難傷害我們.
令到我們有更深的傷害.
而且我們不會被困難打倒.
信心不是因為我見到神蹟才相信.
是因為就算我看不到神蹟.
但我仍然相信神比我眼前的困局更加珍貴.
這是最重要的.
還有第二.
信心就是要看見上帝的引導.
上帝如何帶領我們的人生.
這是非常重要的.
方向錯誤.
我們的人生也會搞錯.
當亞蘭大軍來到以利沙的時候.
要活捉他們.
以利沙做了第二個祈禱.
記載在第十八節.
我們一起讀一讀.
第十八節請讀.
求你使這些人眼目昏迷.
耶和華就照以利沙的話.
使他們的眼目昏迷.
有幾樣東西值得我們留意.
首先第一.
錄了第十八節原文其實是沒有敵人這個字.
只不過是說他們下到以利沙那裡.
還有第二.
這次也是跟眼睛有關的.

$^{561}$但是是叫這些人的眼目昏迷.
還有第三.
中文聖經用這些人其實是在說國民或者是國家.
而且是單數的.
就是說是表達一個整體的國家.
或者整體的國民.
還有第四.
中文聖經翻譯漏了一個字.
就是擊打這個字.
所以以利沙求的是什麼.
求耶和華擊打這些國民.
又或者擊打這些國家.
還有第五.
昏迷不是在說腦殘.
是在說眼殘.
是你的眼失明.
看不到.
什麼時候你的眼睛看不到.
你打了手機多就看不到.
晚上用了手機很厲害的.
突然之間你發覺你的眼睛很模糊.
你就發覺用了手機太多.
是時候要關掉電話了.
有時候失明就是看不到方向.
看不到前面.
六章第十九至以利沙就和敵軍說.
他說不是這條道路.
也不是那個城.
你跟著我.
你走錯地方了.
以利沙說我現在帶你去一個對的地方.
一帶就帶了二十公里的以色列北角的首都.
撒瑪利亞.
就是皇帝住的那個城市那個地方.
這一班的軍隊亞蘭人.
他們恃著自己軍力很強.
戰鬥力雄厚.
但是想不到神幽默的樣子.
利用以利沙玩弄他們.
其實真是的.

$^{601}$正確的信心是需要看見上帝的引導.
我們才有方向.
否則的話我們會迷路.
很多的學生以為自己很勤力讀書.
那就可以讀很好的書.
但是其實未必一定是這樣.
勤力當然是重要.
最重要的其實就是.
你跟著有沒有對的人生的方向.
都是很重要.
所以方向是非常重要.
就好像GPS那樣.
你輸錯了地址.
你會不會去到那個地方.
真的去不到.
就算你輸對地址也未必去到.
更何況你輸錯地址.
但是你要知道.
當你跟著GPS走的時候.
你錯了也好.
如果GPS是對的話.
它會導航你去一個對的方向的位置.
有時人生都是這樣.
神一定會讓都給答案我們.
但是他會慢慢一步一步地.
帶領著我們的路.
有一個大學生他畢業沒多久之後.
他找工作都找了幾個月.
都很不順利.
最後終於進了一間公司.
頭幾個月覺得自己不斷地衝.
又OT又應酬又博表現.
還跟一班朋友天天去PUB.
去一些地方去周圍玩.
為了應酬他們.
搞到自己玩得很盡.
因為這個緣故.
開始發覺自己很累.
身體開始出現一些問題.
情緒越來越差.

$^{641}$經常覺得自己開始.
這條路到底是否應該這樣走.
不知道自己應不應該繼續這樣走下去.
因為自己越來越少回教會.
後來跟一些前輩聊天的時候.
就告訴他.
你會不會有些地方.
需要去糾正一下.
你的方向是不是搞錯了.
那一刻開始.
他才重新去檢視自己.
到底是不是有些東西搞錯了.
他開始慢慢調整自己生活的節奏.
調整自己跟哪些人去接觸.
於是乎他就能夠糾正自己的方向.
方向很重要.
我們的方向搞錯了的話.
我們很多時候會有更大的問題出現.
當亞蘭大軍入了以色列的首都.
君王住的薩瑪利亞城之後.
六章二十節.
以利沙又祈禱.
我們一起讀一讀二十節.
請讀.
不料是在薩瑪利亞的城中.
這次又是以利沙的祈禱.
這番話幾乎跟以利沙的僕人所說的話一模一樣.
當他們睜開眼睛之後.
看到自己竟然在敵軍的城市.
被以色列的軍隊重重包圍著.
這是第三次以利沙的祈禱.
第一次是為年輕的僕人祈求他們.
祈求開他們的眼.
第二次祈禱就是.
叫這些國民突然之間失明.
第三次的祈禱就是.
開這班人的眼.
讓他們看見.
當他們看見的時候.
就嚇得他們目瞪口呆.

$^{681}$因為原來已經落了去敵方的核心的城市.
薩瑪利亞.
我想你留意.
一開始故事的時候.
亞蘭王設局.
他打算設局.
所以神人就叫他小心提防.
現在神人設一個局.
亞蘭人卻不能夠提防.
結果他們整批大軍.
就去到以色列人所.
以色列的首都.
薩瑪利亞這個城市.
這個是多麼諷刺.
而更諷刺的地方就是.
敵人原來是想找以利沙的.
但原來找到的一條路.
是被引到去一個敵軍的首都.
當我們方向錯.
就會帶我們人生走錯路.
所以我們必須要神的引導.
神的手能夠將敵人的路.
變成救恩的路.
將圍困的局面.
變成祝福的局面.
信心不是要求神立即拿走困難.
而是讓我們相信.
在困難的裡面.
神不停不停在引導我們.
開一條路給我們走.
信心就是我們必須要認定.
祂是引領我們的路.
引證我們正確的方向.
昨天我有機會去到萬城一個地方.
聽一個敬拜的聚會.
那裡面認識了一位來自倫敦的香港人.
他就說他自己來了英國五年.
他是一個音樂樂手.
他一邊聊天的時候就說.
他和他的太太來了英國這個地方.

$^{721}$但是來到英國這個地方就很辛苦.
因為有很多實在太難適應.
甚至這個不適應.
存在在他們的婚姻關係裡面.
甚至太太最後離開了他.
他們兩個都是一個很愛主的基督徒.
甚至太太都離開了他.
他覺得人生不知怎算好.
他原本是想著.
有一個可以用他的恩賜來服侍神.
但是來到英國.
很多地方都沒有辦法可以實踐到這件事.
不過他後期經歷了很多不開心之後.
他仍然繼續尋求上帝的引導.
結果他就去做了一間傢俬公司.
他專門運送傢俬.
他就拿著他的車全英國送傢俬.
他不單止要送傢俬.
他更加要幫人「撼傢俬」.
你知道英國很多家裡都很空閒.
經常都在家裡面.
於是乎他經常去「撼傢俬」的時候.
就和那些人在那裡聊天.
原來他每一次透過「撼傢俬」.
就和當時那些無論是窮.
或者是獨居.
或者是殘障人士.
就聊了很多天.
裡面久而久之就建立了一些橋樑.
給他傳遞福音.
他相信神叫他來英國.
有他的意思.
當然在裡面有很多困難的地方.
但是仍然都相信神會引導他.
今日相信我們每一個人.
都是有很多不同的路經歷.
有時仍然都覺得不容易行.
但是不要忘記.
上帝是會引導我們.
弟兄姊妹.

$^{761}$最後第三樣東西.
信心是經歷生命的轉化.
在這個故事裡面提醒我們.
真正對神的信心.
必須至少有兩重生命的轉化.
首先是「赦免」.
我想你留意讀二十一節.
請讀二十一節.
這段經文.
首先我們要學習的是「饒恕」.
當以色列王.
見到被俘虜的亞蘭軍的時候.
兩眼發光.
他就說我富.
我可以擊殺擊殺.
兩次擊殺.
可能是一個強調的作用.
又或者.
這裡是反映他心裡的掙扎或者興奮.
以色列這個王的反應.
就代表著我們掌握到主權.
主動權的時候.
多麼容易去利用自己的能力.
去做一些不是上帝要我們做的事情.
當我們有這樣的空間.
有這樣的主動權的時候.
這個王竟然之間用他自己的心意.
想去報復.
去證明自己.
去彰顯他的權力.
但是六章二十二節.
以尼沙馬上就回應.
他說不可以擊殺他們.
我想你留意第二十二節.
二十二節是比較難處理的.
我們讀一讀二十二節.
請讀.
當在他們面前擺設飲食.
使他們吃喝回到他們的主人那裡.
這裡以尼沙馬上回應.

$^{801}$在原文其實這句的意思是說.
這些人起羨你用你的刀.
用你的弓擄來.
讓你去擊殺的嗎.
簡單一點說.
意思就是說.
這班人又不是你打仗打回來的.
又不是你用刀殺的.
又不是你用你的能力做出來的.
你有什麼權利.
有什麼能力去擊殺他們.
根本就不是你的努力.
所以你這個字是一個很重複.
很強調的位置.
以尼沙馬用了兩次擊殺這個字.
去回應王兩次擊殺這個字.
以尼沙馬給出的原因.
很簡單.
根本不是你的努力.
所以你沒有資格殺他們.
以尼沙似乎借這個機會.
讓這個王要認真去相信一件事.
如果你對上帝有信心的話.
就應該要學習一個生命的品格.
就是要饒恕.
而這個饒恕真正的能力.
就是在饒恕那裡.
所以甚至在六章二十二字.
以尼沙更加要王行前多一步.
他說要擺設飲食.
使他們吃喝回到他們的主人那裡.
對以色列的王來說.
一舉消滅亞難人.
你說這個豈不是更好.
當然以尼沙給出的原因.
剛才我們說.
不能夠殺.
因為你根本沒有努力過.
所以你根本沒有資格去殺.
既然是這樣.

$^{841}$不如感激他們.
讓他們感激我們.
饒過他們一世.
這樣讓他們將來不會打我們.
令以色列人都有喘息的機會.
其實我們要學的.
就是這個王要學的東西.
饒恕.
我們經常說我們對神有信心.
饒恕一個功課.
就是一個最有信心的基督徒.
應該要學習做的東西.
饒恕就是一個要學習的東西.
在一九九四年.
非洲國家盧旺達.
發生了震驚世界的盧旺達大屠殺.
短短一百日之內.
死了八十萬人.
有一個基督徒的母親.
親眼見到自己的丈夫.
和幾個小朋友被殺.
幾年之後.
兇手被捕.
在監獄裡面.
這些兇手來到去悔改信主.
這個兇手悔改信主.
後期有人安排了這個母親.
和這個兇手會面.
現場的人都會覺得.
這個母親可能會控訴.
會咒罵.
甚至會情緒崩塌.
但是她做了一件令人難忘的事.
就是她竟然之間和這個人說.
她說我沒有辦法忘記.
但我仍然會選擇去饒恕你.
因為基督已經饒恕了我.
然後就伸手去抱著這位的人.
令到全場都安靜了.
之後她還簡單說了一句話.

$^{881}$她說如果我不饒恕你.
我就仍然被你去捆綁我.
還有第二件事.
除了饒恕學習.
是一個生命的轉化的學習之外.
就是要學習馴服的功課.
值得留意的地方就是.
以利沙只是說食物和水.
食物是指什麼.
是指麵包或者餅乾.
水當然是水.
等他們吃完之後就讓他們回家.
這是恩典比令人記仇好.
以利沙是這樣想的.
但是這個王.
以色列的王竟然為亞蘭的大軍.
預備更多的食物.
因為經文用盛大的筵席.
什麼叫盛大.
什麼叫筵席.
筵席就是說他吃得很厲害.
喝到有酒在裡面.
換句話說.
以色列的王.
在飲食這麼少的事上.
都有自己的主見.
在這麼少的事上.
他都有自己的主見.
作了一個加強豪華版出來.
以色列王就好像覺得.
神的吩咐不夠力.
或者太簡單.
要自己加些料進去.
例如我們都可能會這樣.
饒恕不是不行的.
要這個人做某些事才行.
我和他做朋友不是不行的.
要他做某些事才行.
我們都會加了些料進去.
王是不是想透過厚待這班人.

$^{921}$或者用收買的角度.
社交的手段.
去迎來更大的操縱呢.
想換來更長久的太平.
不過如果真的這樣的話.
恐怕這種和平只是很脆弱.
只是很短暫.
六章第24節果然是.
沒多久就再打以色列.
這位的以色列王.
似乎是一位權力的導向機會主義者.
我們見不到這位的王.
見到這樣的局面之後.
他有半點感恩的心.
你找一找.
經文從來都沒有這樣說.
也沒有見到這個王.
有半點敬畏上帝的心.
以色列的王沒有.
他掛在嘴邊的.
就是擊殺殺.
他做的是自己想做自己的事.
而不是神人想吩咐他的完整版.
弟兄姊妹 今天我們要學習的.
生命轉化.
就是必須要信服我們的上帝.
有一個很出名的運動員.
叫做Tim.
他是美國大學一個足球界裡面.
幾乎是一個很傳奇的人物.
他拿過很多榮耀的獎項.
前途可以說是一片光明.
對一個年輕人來說.
這個簡直是夢想.
全部集在這個年輕人身上.
但當他進入職業聯賽之後.
他選擇公開自己的信仰.
甚至在比賽裡面.
他去禱告 去祈禱.
結果引來大量的媒體笑他 批評他.

$^{961}$甚至有人說他太誇張 太極端.
但他就說.
我只是為神而活.
這個就是信服的一個表現.
弟兄姊妹.
信服這個功課是一個經歷生命轉化.
不可以缺少的東西.
弟兄姊妹 信服有很多東西.
可以是轟轟烈烈的.
可以是很多大場面的事.
但我們不需要想得這麼誇張.
我們在生活的細小的每一個部分.
小的部分 我們去學習.
上帝你想我做些什麼.
而想我做的這件事.
我不嘗試去加任何的料.
我不嘗試去用我自己的看法放在裡面.
我只是簡簡單單 很單純的.
上帝你這樣說 我照去做.
希望這段經文都提醒我們.
我們求主開我們雙眼.
藉著亞蘭王 以利沙的僕人.
和以色列王的這些舉動.
讓我們開我們雙眼.
看到我們真正要學習的信心是怎樣.
我們一起禱告.
透過這段聖經.
從這三個人物裡面.
我們看到他們對於神.
對於神的工作的反應.
態度 行動.
我們去學習.
主幫助我們去學習信心這個功課.
主幫助我們透過.
以利沙在這段經文裡面的角色.
他怎樣深入這三個人物裡面.
去提醒我們.
我們要怎樣去學習信心.
願你將這段經文放在我們的心裡.
奉耶穌基督恩主的名 阿們.

\newpage



\section{使徒行傳 11:27-12:25}
\label{sec:031983}
\textbf{越獄風雲 (2)}~~~~林鍵牧師
\newline
\newline
連結: \href{http://www.church.com.hk/acms/content.asp?site=cdc\&op=show\&type=product\&code=031983\&layout=sermon}{\texttt{031983}} ~~~~ 語音日期: 2026-04-11
\newline
\newline
\hyperref[sec:031982]{\small{< < < PREV SERMON < < <}}
~
\hyperref[sec:index_chronic]{\small{[返順時目]}}
~
\hyperref[sec:031984]{\small{> > > NEXT SERMON > > >}}
\newline
\newline
使徒行傳 11:27-12:25
\newline
\begin{longtable}{cl}
\hline
\hline
章節 & 經文 (和合本修訂版)\\
\hline
11:27 & \begin{tabularx}{0.7\textwidth}{X} 當那些日子，有幾位先知從耶路撒冷下到安提阿。 \end{tabularx} \\ \\ \relax
11:28 & \begin{tabularx}{0.7\textwidth}{X} 內中有一位，名叫亞迦布，站起來，藉著聖靈指示普天下將有大饑荒；這事在克勞第年間果然實現了。 \end{tabularx} \\ \\ \relax
11:29 & \begin{tabularx}{0.7\textwidth}{X} 於是門徒決定，照各人的力量捐錢，送去供給住在猶太的弟兄。 \end{tabularx} \\ \\ \relax
11:30 & \begin{tabularx}{0.7\textwidth}{X} 他們就這樣做了，託巴拿巴和掃羅的手送到眾長老那裡。 \end{tabularx} \\ \\ \relax
12:1 & \begin{tabularx}{0.7\textwidth}{X} 約在那時候，希律王下手苦待教會中的一些人， \end{tabularx} \\ \\ \relax
12:2 & \begin{tabularx}{0.7\textwidth}{X} 用刀殺了約翰的哥哥雅各。 \end{tabularx} \\ \\ \relax
12:3 & \begin{tabularx}{0.7\textwidth}{X} 他見猶太人喜歡這事，也去拿住彼得。那時候正是除酵節期間。 \end{tabularx} \\ \\ \relax
12:4 & \begin{tabularx}{0.7\textwidth}{X} 希律捉了彼得，押在監裡，交給四班士兵看守，每班四個人，企圖要在逾越節後把他提出來，當著百姓辦他。 \end{tabularx} \\ \\ \relax
12:5 & \begin{tabularx}{0.7\textwidth}{X} 於是彼得被囚在監裡，教會卻為他切切禱告神。 \end{tabularx} \\ \\ \relax
12:6 & \begin{tabularx}{0.7\textwidth}{X} 希律將要提他出來的前一夜，彼得被兩條鐵鏈鎖著，睡在兩個士兵當中；門前還有警衛看守。 \end{tabularx} \\ \\ \relax
12:7 & \begin{tabularx}{0.7\textwidth}{X} 忽然，有主的一個使者顯現，牢房裡有光照耀；天使拍彼得的肋旁，叫醒了他，說：「快起來！」鐵鏈就從他手上脫落下來。 \end{tabularx} \\ \\ \relax
12:8 & \begin{tabularx}{0.7\textwidth}{X} 天使對他說：「束上腰帶，穿上鞋子。」他就照著做了。天使又對他說：「披上外衣，跟我來。」 \end{tabularx} \\ \\ \relax
12:9 & \begin{tabularx}{0.7\textwidth}{X} 彼得就出來跟著他走，不知道天使所做是真的，以為見了異象。 \end{tabularx} \\ \\ \relax
12:10 & \begin{tabularx}{0.7\textwidth}{X} 他們經過了第一層和第二層監牢，就來到往城內的鐵門，那門就自動給他們開了。他們出來，走過一條街，忽然天使離開他去了。 \end{tabularx} \\ \\ \relax
12:11 & \begin{tabularx}{0.7\textwidth}{X} 彼得清醒過來，說：「現在我真知道主差遣他的使者，救我脫離希律的手，和猶太人所期待的一切。」 \end{tabularx} \\ \\ \relax
12:12 & \begin{tabularx}{0.7\textwidth}{X} 他明白了，就到那稱為馬可的約翰的母親馬利亞家去，在那裡已有好些人聚集禱告。 \end{tabularx} \\ \\ \relax
12:13 & \begin{tabularx}{0.7\textwidth}{X} 彼得敲外門時，有一個使女，名叫羅大，出來應門， \end{tabularx} \\ \\ \relax
12:14 & \begin{tabularx}{0.7\textwidth}{X} 認出是彼得的聲音，歡喜得顧不了開門，就跑進去報信，說彼得站在門外。 \end{tabularx} \\ \\ \relax
12:15 & \begin{tabularx}{0.7\textwidth}{X} 他們對她說：「你瘋了！」使女堅持真有其事。他們說：「那是他的天使。」 \end{tabularx} \\ \\ \relax
12:16 & \begin{tabularx}{0.7\textwidth}{X} 彼得不停地敲門；他們開了門，一見是他，就很驚奇。 \end{tabularx} \\ \\ \relax
12:17 & \begin{tabularx}{0.7\textwidth}{X} 彼得做個手勢，要他們不作聲，就告訴他們主怎樣領他出監；又說：「你們要把這些事告訴雅各和眾弟兄。」然後，他離開往別處去了。 \end{tabularx} \\ \\ \relax
12:18 & \begin{tabularx}{0.7\textwidth}{X} 到了天亮，士兵中起了不少騷動，不知道彼得到哪裡去了。 \end{tabularx} \\ \\ \relax
12:19 & \begin{tabularx}{0.7\textwidth}{X} 希律找他，找不著，就審問警衛，下令帶走他們處死。後來希律離開猶太，下凱撒利亞去，住在那裡。 \end{tabularx} \\ \\ \relax
12:20 & \begin{tabularx}{0.7\textwidth}{X} 希律向推羅和西頓的人發怒。他們那一帶地方是從王的土地供應糧食的，因此就託了王的內侍大臣伯拉斯都的情，一心來求和。 \end{tabularx} \\ \\ \relax
12:21 & \begin{tabularx}{0.7\textwidth}{X} 希律在所定的日子，穿上朝服，坐在位上，對他們演講。 \end{tabularx} \\ \\ \relax
12:22 & \begin{tabularx}{0.7\textwidth}{X} 民眾一直喊著：「這是神明的聲音，不是人的聲音。」 \end{tabularx} \\ \\ \relax
12:23 & \begin{tabularx}{0.7\textwidth}{X} 希律不歸榮耀給神，所以主的使者立刻擊打他，他被蟲咬，就斷了氣。 \end{tabularx} \\ \\ \relax
12:24 & \begin{tabularx}{0.7\textwidth}{X} 神的道日見興旺，越發廣傳。 \end{tabularx} \\ \\ \relax
12:25 & \begin{tabularx}{0.7\textwidth}{X} 巴拿巴和掃羅完成了供給的事，就回到耶路撒冷，帶著稱為馬可的約翰同去。 \end{tabularx} \\ \\
[1ex]
\hline
\hline
\end{longtable}
$^{1}$歡迎大家來到我們的周六崇拜.
上星期是復活節.
我們當中的小朋友,大朋友.
都為我們上演了一場非常精彩的月玉風雲.
故事是出於《少爺行傳》第十二章.
是《少爺行傳》裡面一個很核心的故事.
精英學者告訴我們.
《路加福音》和《仕途行傳》其實是一本書.
是一個分為兩卷合一的著作.
而兩本書加起來可以分為三大部分.
第一部分是關於基督耶穌.
所謂基督就是上帝所揀選,所認許的.
舊約的彌賽亞以色列的君王,救主.
第二部分是他差遣了聖靈來到人間.
第三部分因為聖靈的緣故.
於是有《仕途行傳》的出現.
你會發現在每個部分底下都有兩個空格.
是邀請大家去填寫的.
當你填寫了那些空格的時候.
你就會更加清楚看到.
《路加福音》和《仕途行傳》的大綱和主題.
第一部分A Part A.
開始的時候是基督耶穌.
基督耶穌為我們做了什麼呢?.
祂受死,祂復活,祂升天,祂作王.
祂為我們受死,為我們復活.
祂替我們支付了我們的罪債.
祂又為了我們得勝了死亡.
然後祂升天,祂作王.
然後祂差遣了祂自己的聖靈.
既是聖靈又是祂自己的靈.
聖靈降臨是為了什麼呢?.
聖靈降臨是為了我們有一個以馬五思的經歷.
這裡說到《路加福音》第24章.
如果你在我們當中一段時間.
你已經很熟悉的.
因為我們經常說.
是《路加福音》的一個核心的重點.
在這個以馬五思路上.
耶穌受死復活.

$^{41}$然後有兩個眼睛不明亮.
心裡快要死的門徒.
耶穌和他們一起有兩,三個小時一起去查經.
有藥查到他身弱.
然後他們起死回生.
他們有一個生命的改變.
他們眼睛明亮,心裡火熱.
然後還做了一個大的轉彎.
於是我們說.
耶穌受死復活,升天作王.
然後祂差遣祂的聖靈來.
最主要的目的是.
讓罪人有一個以馬五思的經歷.
以致他們眼明.
以致他們心火熱.
於是他們成為了一個復活的聯盟.
唯有一個復活的聯盟.
才可以帶給我們一個仕途的行轉.
而仕途行轉的方向,目的.
就是要地極見證,宣揚天國.
《紹行傳》第一章第八節.
整個《紹行傳》的主句.
「但聖靈將要降臨在你們身上」.
那群門徒身上.
「你們就必得著能力.
要在耶路撒冷,猶大,泉地,撒瑪利亞.
直到地極為耶穌作見證」.
《紹行傳》可以分為兩半.
頭一半是一章至十二章.
第二半是十三至二十八章.
頭一半的主角或主要的配角是彼得·彼得·彼得·彼得·彼得·彼得·彼得·彼得·彼得·彼得·彼得·彼得·彼得·彼得.
而福音的事公是圍繞著.
耶路撒冷的教會猶太人的信徒.
而後一半的主角或主要的配角.
是索羅或保羅.
第十三章的名字改了由索羅變成保羅.
福音的事公是圍繞著耶路撒冷,猶大,撒瑪利亞以外的地方.
外邦的信徒.
在這裡第十二章.
尤其是非常重要.

$^{81}$是整本書的轉接點.
亦在第十二章裡彼得受死復活的故事.
上星期我們看,今星期我們繼續看.
同時去到最後一節你會發現.
索羅和彼得是同框的.
在同一章聖經裡索羅都出現了.
於是我們開始明白這是兩個主角主要配角交棒的時候.
去到第十三章開始索羅成為了主角.
在第十二章裡.
表面上我們看的故事.
是彼得的個人越獄的救贖.
他基本上死定了.
誰知上帝竟然救贖了他.
但實質上你會發現.
在彼得的故事背後有一個更偉大的人物是耶穌基督.
實質上是耶穌基督受死復活.
然後走出墳墓的影子.
上星期我們去看這兩個故事的對比的時候.
我們發現好像是在說彼得.
其實是在說耶穌.
所有的一切都是聖靈的工作.
為要使主道興旺.
我們看到最後的終結是主道興旺.
越發廣傳然後地極見證.
去宣揚天國.
上星期我用了一個圖表來幫大家列舉.
使徒彼得越獄救贖的故事.
和耶穌基督受死埋葬復活的故事.
有多少項很多很多不同的平衡.
我想再和大家很快的溫習一下.
我想你看一下耶穌基督受死復活是如何引發了.
使徒彼得的越獄救贖.
一開始的時候兩個故事都是.
除巧逾越節耶穌受死埋葬.
然後除巧逾越節彼得被押在監牢.
就算他被押在監牢那個圖畫都很特別.
他說彼得睡在兩個士兵當中.
在古典文獻裡面睡覺通常是代表了死亡.
所以他睡覺你不要說他有沒有信心.
很安心的明天都被人砍頭.

$^{121}$他今天還會好心不是那不是重點.
重點是他睡覺的圖畫代表了他死亡.
和他死亡的時候死在兩個士兵中間.
是不是好像耶穌.
耶穌被釘死的時候都是在兩個人中間.
所以整個故事整幅圖畫是.
映射著耶穌死的那個圖畫.
然後天使滾開墳墓石頭.
天使又打開了監獄的鐵門.
夫人見證耶穌復活.
羅大聽見彼得的聲音.
仕途不信耶穌的復活.
眾人不信彼得逃生.
門徒以為是見了雲.
眾人以為是他的天使.
沒有身體大家都不相信有身體.
門徒以為要把這事告訴使徒.
耶穌說的.
然後彼得說要把這事告訴雅各.
雅各是耶穌的弟弟.
不是之前死去的那個雅各.
士兵收錢復活受禁.
士兵被殺逃生受禁.
臣使耶穌復活作王.
臣使彼得逃脫希律.
最後十家跨姓權勢解除.
希律斷氣主導興旺.
上星期我開始引領大家去看.
一個很重要最後的元素.
就是希律王.
這個希律其實是帶著魔鬼撒旦的影子.
你記得他最後死之前.
他好像穿了皇服.
然後別人在讚他.
是上帝的聲音.
於是他很高傲地.
沒錯是上帝的聲音.
撒旦就是想取締上帝.
想做上帝.
於是這裡的希律帶著魔鬼撒旦的影子.

$^{161}$這個撒旦去送耶穌上十字架.
但是卻因為耶穌的受死復活救贖成全.
撒旦最終被擊敗.
被天使所擊打.
受了上帝的審判.
從某個角度來說.
我們可以這樣說.
由始至終.
魔鬼撒旦的終極命運是沒有改變的.
無論耶穌是否想上十字架.
其實魔鬼撒旦最後還是落地獄的.
撒旦最終的命運.
就是被陰盡地獄的火紅熊裡面.
但是撒旦其實是有一個計劃.
而他的計劃是被耶穌的十字架.
所徹底破壞了.
撒旦的計劃就是要你.
撒旦的計劃就是要我.
撒旦的計劃就是要地上每一個罪人.
你很有份.
撒旦對你很有興趣.
撒旦對我很有興趣.
他的計劃就是要地上每一個罪人.
都和他一起叛逆上帝.
和他一樣高傲自義.
和他一樣貪愛世界.
最後和他一起被墮進地獄.
他就死定了.
我不知道他信不信他真的死定了.
但是他就想你陪他一起死.
你就不用死.
因為你有耶穌.
在耶穌未上十字架之前.
我們就死定了.
但是因為耶穌受死復活.
因為耶穌升天作王.
因為耶穌差遣了聖靈.
因為今天我們有機會聽到福音的見證.
所以現在耶穌通過他的受死復活.
使彼魯的德釋放.

$^{201}$我們可以有機會脫離撒旦的捆綁.
使半億的罪人.
我們半億的罪人.
從前半億的罪人.
有可能轉向基督十字架.
重生我們成為上帝的兒女.
耶穌未去十字架之前.
每一個都和他一鑊熟.
全部和撒旦一鑊熟.
但是耶穌上了十字架之後.
今天我們有一個不一樣的可能.
有一個不一樣的結局.
當然你不選擇耶穌.
你的結局和撒旦一樣.
和你以前這麼多年來是沒有分別的.
但是這裡若有人相信耶穌.
若有人願意悔改.
信靠耶穌跟隨他.
再不跟隨撒旦這個世界去走的時候.
我們今天可以有另外一條路.
而這個是撒旦不想見到的東西.
所以我想告訴你.
在路加的思想裡面.
撒旦是一個很重要的人物.
我想帶你去看一下路加福音兆行傳.
這兩本書其實是一本.
當中有一些關於撒旦的經文.
因為耶穌的受死復活.
其中一個很重要的結局.
很重要的結果.
就是可以打敗撒旦.
我想帶你去先看路加福音第十章.
我想大家都很熟悉這段經文.
耶穌派十二個使徒.
他的門徒出去.
去趕鬼去傳福音.
於是那七十二個.
歡歡喜喜地回來.
第十章第十七節說.
主啊因你的名.

$^{241}$就是鬼也服了我們.
接著第十八節耶穌對他們說.
我看見撒旦.
就不是鬼仔了.
這個是那隻大鬼.
撒旦從天上墜下.
像閃電一樣很快跌下來.
我已經給你們權柄.
可以踐踏蛇禍蠍子.
有勝過仇敵一切的能力.
絕沒有什麼能害你.
這裡耶穌雖然說.
我已經看見.
你看英文其實是一個過去式.
我已經看到了.
我看見撒旦從天上墜落.
但其實這一切.
只是將來勝利的預演.
它還沒發生.
為什麼呢.
因為耶穌終極的勝利.
是在十字架上.
如果耶穌沒有上十字架.
是沒有勝利的.
唯有十字架才可以救贖罪人.
唯有十字架才可以使我們.
爭奪撒旦的捆綁.
唯有十字架才可以粉碎撒旦的計劃.
所以我想你將十字架.
放回撒旦墜落的那幅圖畫裡.
所以不單止是撒旦墜落.
而是因為十字架的緣故.
所以撒旦墜落.
耶穌說.
因著十字架的緣故.
我看見撒旦墜落.
然後耶穌再說.
因著十字架的緣故.
最後那句.
你們可以勝過仇敵一切的能力.

$^{281}$絕沒有什麼能害你們.
說得很大.
這個當然不代表.
你信了耶穌之後.
百毒不侵.
身體健康.
出入平安.
什麼叫做因著我的德性.
因著十字架的緣故.
絕沒有什麼東西能害你們.
其實他指的是下一節聖經.
下一節聖經是第十章第二十節.
耶穌這樣說.
然而不要因靈服了你們就歡喜.
因為他覺得自己很厲害.
去趕鬼也可以趕鬼.
醫治就醫治.
耶穌說.
不要因靈服了你們而歡喜.
而是要因你們的名記錄在天上而歡喜.
耶穌說.
為什麼你沒有東西可以傷害到你.
沒有任何的力量.
絕對沒有什麼能害你.
因為你的名字.
因著你跟隨我的緣故.
被記錄在天上.
不是因為告訴你.
你百毒不侵身體健康.
其實耶穌很清楚告訴我們.
大家都記得.
第十二章第四節這樣說.
那最多只能殺人身體的.
不能再做甚麼的不要怕他們.
其實你身體可不可以被人殺.
當然是可以的.
但是你不用怕.
為什麼.
因為你的名字寫在天上.
如果你名字寫在天上的話.

$^{321}$我提醒你.
該怕的是誰.
該怕的是那殺了以後又有權.
並把你陰在地獄裡的.
但是你的名字.
如果你跟隨了耶穌.
你的名字已經在天上的話.
那你還有什麼好怕呢.
於是這是耶穌告訴我們.
他十字架的特性.
於是我們看到.
耶穌和撒旦不斷的爭鬥.
耶穌和撒旦的爭鬥.
最後帶我們去到.
耶穌受難前的前夕.
很出名的.
路加福音第二十二章第一節.
第二節第三節.
除巧節又叫做逾越節.
近了.
祭司長握民事.
在想法子怎樣去殺害耶穌.
這是撒旦.
撒旦.
入了猶大的心.
接著第三十一節.
主又說西門.
西門撒旦.
撒旦要得著你們.
但我已經為你祈禱.
使你不至於失去信心.
你回頭以後要堅固你的弟兄.
這裡說到撒旦入了猶大的心.
正如我之前所說撒旦送耶穌上十字架.
但他其實很傻.
他不知道耶穌的受死復活.
成全救贖.
最後會令到他被擊敗.
當然耶穌擊敗了撒旦.
不代表所有人都會得救.

$^{361}$因為每個人都要作一個個人的決定.
於是第二十二章第三十一節.
主說西門西門.
西門就是彼得.
撒旦要得著你們.
你們是重數.
不是說只有彼得.
是說所有十二門徒.
他要得著你.
怎樣得著你呢.
像西密子這樣.
西密子是劇烈的震盪.
You shake your life.
但是耶穌和彼得說.
我已經為你.
這個你是單數.
但他和彼得說.
也和其他人說.
只要你有一天.
你會回頭.
我知道你們.
我釘死的時候.
你們全部都跑了.
但是只要.
你有回頭的那一天.
只要你回頭仰望我.
你就有盼望.
沒有十字架.
彼得回頭都沒得救.
你明白嗎.
但是有十字架.
如果彼得有一天回頭的時候.
他可以仰望耶穌的十字架.
他要跟隨這個.
為他釘十字架的基督的話.
他就有得救.
然後耶穌再加一句.
你回頭以後.
我知道你會回頭的彼得.
你回頭之後.

$^{401}$你要兼顧你的弟兄.
我們知道.
彼得後來是成為了.
耶路撒冷和別處教會的.
教會建築師.
如果你聽Kevin說.
彼得前書.
他經常這樣說.
他不是一個漁夫這麼簡單.
他懂得建築教會.
教會是由耶路撒冷開始.
他第一個建築項目.
然後彼得繼續去建立教會.
但是我想多說一件事.
就是彼得除了做教會建築師以外.
他同時間也是兼任教會的保護者.
教會的驅魔人.
很有趣的.
以前你會發現.
彼得也是被撒旦所利用.
所以有一次.
耶穌和他說.
撒旦退去吧.
要幫彼得趕鬼.
因為耶穌說.
他要上十字架.
彼得就被撒旦引誘.
就說不可以.
然後耶穌說要撒旦退去吧.
但是當彼得最後回轉的時候.
彼得成為了教會的驅魔人.
你說林穆斯你在說什麼呢.
我想帶你現在.
由路加布音轉移到.
仕途行傳.
當彼得最後信主回轉的時候.
他會成為了教會的驅魔人.
我想帶你去看.
耶穌基督復活的仕途.
和撒旦的一場激戰.

$^{441}$激戰在哪裡.
激戰在仕途行傳第五章.
撒旦再次出現.
這次已經被耶穌十字架打倒了.
但是他又再次出現.
很熟悉的故事.
有一個人名叫阿南尼亞.
和他的妻子撒菲拉.
買了田產.
把錢撕至樓下一部分.
他的妻子也知道.
其餘的部分拿來.
放在仕途的腳前.
彼得說阿南尼亞.
你為什麼撒旦充滿了你的心.
為什麼撒旦充滿了你的心.
使你欺騙聖靈.
把田地的錢撕至樓下呢.
撕至樓下也只是一部分.
這段是一個我們很熟悉.
但是又很不明白.
很看不明白的故事.
這裡說的不是一對夫婦.
一時不小心失腳.
於是犯了小小的罪.
這裡說的很明顯.
是一個大大的罪.
這裡說的是撒旦.
撒旦被受死撲擊的耶穌.
重拳擊倒之後.
如果你可以想像.
我拍這部電影.
拍這個故事的背後.
腦修成奴對基督教會.
展開全力的攻擊.
你明白嗎.
撒旦擊倒耶穌一拳.
他已經暈了.
撞到他八,九.
然後他起來.

$^{481}$現在他腦修成奴.
他要向基督的教會.
展開一個全力的反擊.
我們問安娜尼亞和撒菲拉.
犯了什麼罪.
如果你是一個律師.
你用什麼罪名去起訴他們呢.
通常答案第一個人想起的.
就是第九誡不可說謊.
因為他欺騙聖靈.
但是你要記住.
第九誡其實是這樣的.
不可做假見證謀害你的論社.
其實他們沒有謀害論社.
他只是將榮耀吸取自己自身.
他又不是害人.
所以更加貼切的罪名是第八誡.
第八誡是更加大的罪名.
第八誡是不可偷渡.
我想問問你.
他們偷了誰的東西.
他們偷了上帝的東西.
他們偷了上帝的什麼.
他們偷了上帝的榮耀.
將上帝的榮耀變成他們自己的榮耀.
這條是死罪.
做了什麼.
他們賣了田產.
將錢私自留下一部分.
因為聖經沒有告訴我們多少部分.
是10\% 20\%還是50\%.
跟著跟彼得說.
我們所有的錢都在這裡.
簡單來說就是報大數.
報大數的原因是為了什麼.
要得到人的讚美.
他想全世界傳消息出去.
這對夫婦賣了田地.
所有的田地都給了教會.
他想得到人的讚美.

$^{521}$但同時又捨不得.
因為他們貪愛這個世界.
如果我總結他們的罪行.
你開始害怕了.
尤其我們當中有夫婦的時候.
你開始害怕了.
一方面他們恣意求稱讚.
另一方面他們回頭望世界.
他們已經做了標圖.
但他們回頭望世界.
放下不了這件事.
在路加福音裡面.
路加有一句很毒的一句.
叫做Remember Lost Way.
這是路加獨有的耶穌說的話.
你們想想羅德的妻子.
她就是這樣.
這兩個人她回頭望世界.
她已經跟隨了耶穌.
她已經進了教會.
但原來她不是.
她回頭望世界.
第一她恣意求讚美.
求稱讚.
第二她回頭望世界.
兩條都是死罪.
我們知道死罪是因為.
給了聖靈充滿有透視眼.
給了一句點出他的問題.
撒旦充滿了你的心.
撒旦使你們成為了他的工具.
他的奴隸.
使你們成為了撒旦的工具.
與聖靈的教會爭戰.
怎樣爭戰.
要用恣意稱讚和貪愛世界.
來攻擊上帝的教會.
基督的身體.
你明白嗎.
他們想做些什麼.

$^{561}$他們要放一些病毒進去教會.
他放什麼呢 恣意求稱讚.
他放什麼呢.
回望貪世界.
貪愛世界.
恣意稱讚.
他要將這兩樣東西注入耶穌基督的身體.
這裡說的是一個終極的爭戰.
是生死之戰.
不是小小的測驗.
而是大大的考試.
如果給得在這裡手軟的話.
教會就不保了.
這個嬰兒生出來.
你注入什麼.
嬰兒死的.
於是你發現給得上帝毫不手軟.
如果你繼續讀下去.
我讀第五節和第十節給你聽.
安娜尼亞一聽見者話.
就撲倒斷了氣.
接著他的太太一會兒又來了.
又立刻撲倒.
給多一次機會他.
你知不知道這件事.
她立刻撲倒.
再給得腳前斷了氣.
兩個撲倒斷了氣像誰.
像希律.
希律就是第十二章的希律.
那個希律給上帝處罰.
有天使一戳他.
他就撲倒斷了氣.
所以這個故事不是你想像中那麼簡單.
也不是上帝特別嚴謹.
而是貪愛世界.
自以稱讚.
求稱讚.
是教會的死敵.
還有更恐怖的事情.

$^{601}$請教授告訴我們.
這個故事是帶著一個舊約故事的影子.
是另外一個令人恐懼震驚的舊約故事.
也是說一對舊約的夫婦密謀偷取上帝的榮耀.
聶王記上帝二十一章.
你大概主要學到想過這個故事.
但你永遠看不到這個故事究竟有幾奸.
這個故事有幾黑.
這個故事說的是.
阿蝦王和他的皇后.
他們洗別皇后.
他們做假見證.
他們謀害那八.
你記不記得.
為了要拿那八那塊地.
要拿那東西是第八屆.
偷人的東西.
然後犯了第九屆做假見證.
謀害你的鄰舍.
最後殺了他第六屆.
他最後的目的是要拿那八的葡萄園.
我讀兩次聖經給你聽聽.
聶王記上帝二十一章第二節第三節.
阿蝦對那八說.
把你的葡萄園給我作菜園.
因為它靠近我的皇宮.
你若要銀子.
我就按著價錢給你.
開始的時候好像很nice.
想跟他買.
但他忘記一樣東西.
地是不可以賣的.
那八就知道.
因為那八是一個愛耶和華的人.
那八對阿蝦說.
不准我把祖先留下的產業給你.
不單祖先留下.
是上帝給他的.
但我想你想一下.
究竟阿蝦要做什麼.

$^{641}$阿蝦要將他的葡萄園.
變為菜園.
你大概看不明白這個故事.
看不明白這個故事究竟有何等邪惡.
原來是這樣的.
在舊約聖經裡面.
葡萄園是代表了以色列.
上帝和耶穌.
上帝和他的子民.
上帝的園主.
他們是葡萄園.
菜園呢.
原來菜園是代表埃及.
你猜都猜不到.
原來在申命記第十一章第十節裡面.
說到埃及好像一個菜園.
我讀給你聽聽.
你要進去得遺孽的那地.
即是說迦南.
不像你出來的埃及地.
在埃及你殺種後.
要用腳挑冠.
像挑冠菜園一樣.
埃及是菜園.
以色列人是上帝的葡萄園.
所以其實他買了這塊地.
將葡萄園變成菜園.
就是要將整個以色列國.
變成埃及的模樣.
所以當你將舊約的故事.
讀回來新約的故事的時候.
你會驚訝發現.
撒旦的計謀就是.
要將基督的福音教會.
變成埃及的偶像之家.
你明不明.
舊約故事和新約故事其實是有連貫的.
撒旦怎樣可以將教會.
變成埃及的偶像之家.
他的陰謀就是.

$^{681}$將致意求稱讚回頭望世界這兩樣東西.
引入耶穌用十字架贖回來的教會.
你不用攻打那間教會.
你不用攻那間教會.
你只要將兩樣東西放在那間教會裡.
它就變成埃及了.
你只要將致意求稱讚.
你只要將回頭望世界這兩樣東西.
放在教會裡面.
少少少少放進那間教會裡面.
它就變成埃及.
撒旦就掌管.
然後你和他一起去.
於是彼得成為了教會的保護者.
彼得成為了教會的驅魔人.
他以前要耶穌幫他驅魔.
現在他回轉了之後.
他變成了教會的驅魔人.
一定要趕.
這個撒旦一定要趕.
這裡可能會有人問我.
我們教會從未有這樣的事情發生.
是否代表我們教會沒有.
致意求稱讚回頭望世界的夫婦呢.
你是想的.
我想這是一個很特別的故事.
上帝將一個終極的審判.
提前在當時處理了這個問題.
這兩個人馬上撲倒斷氣.
我們今天再沒有這些事情發生.
你不要心裡想.
因為我們教會沒有這樣的夫婦.
那你就錯了.
我覺得上帝只不過在容忍我們.
包容我們.
將終極的審判留到末日復活的時候才執行.
彼得後書第二章第三章第九節這樣說.
主沒有言辭他得應許就如有人以為他言辭了.
其實他是寬容你們不願意一人沉淪而人人都來悔改.
我們需要悔改.

$^{721}$如果我們當中有一對這樣的夫婦.
都不用一定是夫婦.
個人都可以.
如果你自己知道.
你有這個致意求稱讚回頭望世界的話.
你要把握這個時機去回轉悔改.
怎樣回轉悔改.
你可以好好地讀我們的終極問答那本書.
你有沒有發現整本書都是在說求上帝的國求上帝的義.
而整本書都是告訴我們.
我們不可以致意求稱讚.
我們不可以回頭望世界.
我們要回轉仰望十字架上的君王救主.
我們要單求他的國和他的義.
去到最後整本書的總結就是.
求他的國和他的義.
我們回來兆候行傳.
當我們從兆候行傳的頭一半部.
一至十二章進入下半部.
第十三章至二十八章的時候.
你會發現撒旦還有一次出場的時候.
這次因為彼得已經交棒了保羅.
所以從前彼得是一個驅魔人.
現在去到後半部.
你猜保羅要做些什麼.
保羅都是做驅魔人.
於是我們讀兆候行傳的時候.
我們可以將兆候行傳看為耶穌兩個使徒的驅魔傳.
前面就是彼得在驅魔.
後面就交棒給保羅繼續驅魔.
當然彼得也是繼續驅魔.
兆候行傳第二十六章第十六節.
撒旦再次會出現.
他說起來站著.
我要派你作僕役.
我也要救你脫離百姓和外邦人的手.
這個很像彼得的經驗.
是相對和保羅說的.
但是經驗和彼得是一樣的.
我會救你的直至我不救你.

$^{761}$因為我們知道最後彼得和保羅都是被人斬首.
所以最後有一次就不救你.
應該說有一次終極救你.
前面就不救你.
我知道我們的想像是很不一樣的.
我們救不救就是有沒有條命出來.
不是.
最後有一次是會救你.
今次就暫時不救你.
所以說我會救你脫離百姓和外邦人的手.
以至你可以繼續去為我做建制.
直至有一天我會召你回來.
第26章第18節.
我差你到他們那裡去.
要開他們的眼睛.
使他們從黑暗中轉向光明.
從撒旦的權下歸向上帝.
當那些人未信主的時候.
每一個人都是在撒旦的權下.
在黑暗的當中.
我差你去做什麼呢.
就是將他們從撒旦的權柄下歸向上帝.
使他們因信我而得蒙社稅.
和一切誠信的人同得基業.
上帝呼召保羅要以基督的福音.
十字架的道理和撒旦的謊言爭戰.
撒旦的謊言是什麼呢.
這個世界很好.
繼續去愛這個世界.
撒旦的謊言很好.
你很厲害.
你很棒.
你是一個好人.
這些全部都是撒旦的謊言.
聖經告訴我們.
完全是兩樣東西.
撒完毒化了罪.
虧缺了上帝的榮耀.
耶穌又不斷說.
地上的財寶都是糞土.

$^{801}$唯有天上的財寶才是真實的.
地上的東西全部都會失去.
上帝會審判這個世界.
唯有存到永遠的東西才是真寶.
用上帝的話語.
用十字架的道理.
用與撒旦的謊言爭戰.
使罪人從黑暗中轉向光明.
使撒旦從撒旦的權下歸向上帝.
這個是保羅的使命的總綱.
於是上帝要預備保羅.
怎樣預備保羅呢.
要預備保羅承擔這個既沉重又偉大的福音.
上帝就特別安排他去耶路撒冷.
在《哲學行傳》第十二章的時候.
你會發現保羅是來了耶路撒冷.
還有保羅是親自目睹彼得的奇妙救贖.
彼得的受死復活.
如果你上個星期在這裡.
你已經看過這個ABCDCBA的故事架構.
但其實我們看少了一層.
這一層裡面有七個.
其實原來是九個的.
一套是九個的.
外面還有一層.
外面這一層是什麼呢.
外面那一層是叫S.
外面那一層原來是跟彼得無關的.
外面那一層是跟掃羅有關係的.
你會發現在第十一章尾的時候.
是講到掃羅因為一些原因來到耶路撒冷.
我讀一點給你聽.
你會記得剛才Wilson也有讀過.
當那西日子有幾名先知從耶路撒冷下到安提河.
保羅那時候在安提河.
內中有一位叫做阿加布.
站起來直聖靈的啟示將有大饑荒.
在耶路撒冷.
猶大那裡會有大饑荒.
於是門徒決定照各人的力量去捐錢.

$^{841}$然後送去給猶大的弟兄.
他們就派了巴拿巴和掃羅的手.
送到耶路撒冷的長老那裡.
所以整個故事開始的時候.
是保羅來到耶路撒冷.
然後講完故事之後.
就是第十二章第二十五節.
巴拿巴和掃羅完成了供給的事.
就從耶路撒冷回去安提河.
我想你留意一件事.
從時序上來看.
其實希律王的死都不是死在耶路撒冷的.
希律王的死是後來的事.
應該保羅是沒有親眼看到的.
但是從寫作的編排寫作的手法來講.
做的彷彿要我們知道.
保羅在耶路撒冷所吸收到的.
所經歷到的包括了希律王的死.
你明白嗎?.
當時的保羅相對來說仍然是年輕的.
他應該還不夠四十歲.
只有三十多歲.
還有他還沒出去做宣教.
還沒出去.
於是這個相對年輕的掃羅.
順著沒多久的掃羅.
他在耶路撒冷親眼目睹.
用心學習.
甚至我會想像.
保羅其實有參加耶路撒冷當中的切切禱告.
你明白嗎?.
他剛好在耶路撒冷的時候.
被人抓了.
於是他又和他們一起學他們的禱告.
這一切都是要預備他很快.
第十三章被差遣到地極作見證.
去萬國做宣教.
於是我們繼續讀的時候.
就是巴拿巴和掃羅完成了工作.
就從耶路撒冷回去安提俄.

$^{881}$在安提俄的教會.
在司空時,敬拜時,禁食的時候.
聖靈對他們說.
要為我分派巴拿巴和掃羅.
去做我照他們所做的工作.
於是他們禁食禱告後.
按手就派遣他們去了.
有關掃羅的呼召和使命.
其實不單是任眾而道遠.
而且是極其艱鉅和危險重重.
我想帶你們回去看看.
掃羅信主那一刻.
起初上帝已經將他將來的使命宣告給他.
這是《紹行傳》第九章.
那時在達曼的市郊.
有一個門徒名叫阿南利亞.
上帝叫阿南利亞去找盲了眼的掃羅.
阿南利亞第一個反應是.
那個人是不行的.
他經常攻擊教會.
然後上帝說已經搞定了.
於是上帝的訊息就是.
起來去找掃羅.
他是我所揀選的器皿.
揀選去做什麼呢.
要在外邦人,軍王和以色列人面前宣揚我的名.
在軍王面前宣揚好像很厲害.
通常軍王都死的.
後面那句才是主句.
我也要指示掃羅.
為我的名必須受許多的苦難.
掃羅的一生.
被上帝的選召.
就是要受許多的苦難.
於是上帝要預備好這位弟兄.
去承受苦難.
令我想起我們最近楊福音的講道系列裡.
真光耶穌派遣十二門徒宣告猜興他們去Into Darkness.
你記不記得.
因為耶穌是真光.

$^{921}$光來到世界裡面.
世界是黑暗的.
現在耶穌要將十二門徒猜派去.
去哪裡呢.
這個世界沒有光的地方.
只有黑暗的地方.
於是Into Darkness送他們進去.
為了好好地預備十二門徒.
在耶穌做完五餅二魚.
剩下十二男子碎餅.
每個門徒要帶著碎餅.
將福音傳到地上的時候.
耶穌再給他們一個為門徒度身訂做的神蹟記號.
叫做水上行走.
要告訴他們.
耶穌會與他們同在.
耶穌猜派他們去一個黑暗的地方.
是一個抵擋福音的地方.
有水有黑夜.
但耶穌將與他們同在.
這是耶穌對他們的應許.
我必天天與你們同在.
直到世界的終結.
這個應許.
如果你是追求世界.
對不起這個應許跟你無關.
唯有當你領受上帝的使命.
為祂而活.
單單尋求祂的國 祂的義的時候.
我必與你同在.
必須要有我與你同在.
因為你不會有好日子過.
於是上帝給一個神蹟記號.
給他們做印證.
我們回到仕途行轉的保羅也是一樣.
因為保羅需要一個印證.
去幫助他知道上帝會與他同在.
你要記得 後來索羅才蒙召.
後來才加入做仕途.
所以索羅沒有親自跟隨耶穌.

$^{961}$和十二門徒不一樣.
十二門徒跟隨耶穌幾年.
索羅沒有機會親自跟隨耶穌.
索羅也沒有親眼看見耶穌的受洗復活.
沒有 他沒有找.
他也沒有經歷過初期教會的誕生.
他所有重要的事件都錯過了.
於是到仕途行轉第十二章.
上帝通過彼得的受洗復活.
這個救贖故事一次過給回他.
你明不明.
他不夠訓練.
於是在他被差遣出去的時候.
上帝為索羅度身訂造了一個福音補習班.
使命加油站.
因為他沒有親眼看見耶穌的受洗復活.
我現在就通過彼得給你看看.
你沒有機會跟初期教會一起.
我給你一個禱告會.
經歷一下初期教會的禱告會是怎樣的禱告會.
你沒有看見雅各被斬.
你現在又看雅各被斬.
你沒有看見希律被上帝審判.
你現在又看那些東西.
於是成為了他的福音補習班.
成為了他的使命加油站.
你沒有看見撒旦的罪行.
不要緊.
我現在補回一個教訓給你.
我們可以用三大點.
來總結年青保羅在通過紹行傳第十二章裡面.
所學到一生受用的寶貴功課.
那段日子我不知道有多長.
他所學到的功課.
我將他濃縮為三大點.
他學到甚麼呢.
第一切切禱告.
放膽傳講.
主導興旺.
越發廣傳.

$^{1001}$一會兒回來我們再說.
第二雅各殉道.
彼得逃生.
主必拯救.
僕人脫離.
這是很重要的.
第三基督上升.
撒旦墜落.
他錯過了.
因為他沒有機會跟隨耶穌.
不要緊現在補回給你.
新約教會.
要逃出埃及.
新約教會的使命就是要帶領眾人逃出埃及.
我們先看第一大點.
第一大點是切切禱告.
放膽傳講.
主導興旺.
越發廣傳.
我想帶你回到上個禮拜.
上個禮拜.
我和大家一起做了反思.
我說究竟他們.
耶路撒冷教會.
切切禱告禱告些甚麼.
全世界的牧師.
所有人都說他們禱告.
當然是禱告彼得放出來.
我說我不同意.
我就和你看.
召喚行傳前面後面.
我相信當時教會所.
做的切切祈禱.
不是祈求彼得的釋放.
而是召喚行傳第四章.
第29節的祈禱.
那個祈禱是甚麼.
就是他們被人捉了被人恐嚇.
放出來的時候他們祈禱.
主啊現在求你監察.

$^{1041}$他們的威嚇.
我們不用怕他們.
使你僕人放膽講你的道.
其實他們要.
祈禱的是使你的僕人.
得時不得時.
放膽傳講你的道.
他們禱告完要聚會的地方.
震動他們被聖靈充滿.
就放膽傳講上帝的道.
OK.
這個是我上個禮拜的重點.
我想告訴你.
我相信年輕的素羅.
他在那個禱告會裡.
學了東西.
明顯學了東西.
我怎知道他明顯學了東西.
因為在他後來的書信裡告訴我們.
他真的學了東西.
這個是已不鎖書.
已不鎖書是其中一個監獄書信.
在監獄裡.
他寫了一個prayer request.
大家為他祈禱.
為他祈禱甚麼.
為他放監.
彼得可以活那麼久我也想活.
已不鎖書第六章第十八節.
要靠著聖靈.
隨時多方禱告祈求.
要為此警醒不倦.
為眾聖徒祈求.
也要為我祈求.
來了來了.
讓我有口才能放膽開口.
講明福音的奧秘.
你要記得他現在不是在香港.
自由世界.
他在坐牢.

$^{1081}$他祈甚麼.
放膽講明福音的奧秘.
我為這福音的奧秘.
作了帶鐵鏈的使者.
See.
他不是叫他鬆鏈.
讓我能放膽宣講.
他是不是學了東西.
是不是希耶路撒冷學的.
不單止是已不鎖書.
還有一本歌羅西書.
又是另一個監獄書信.
你們要.
懇切的禱告.
再禱告遭警醒.
感恩.
同時也要.
為我們禱告求上帝.
給我們開全道的門.
能宣講基督的.
奧秘使我能按著.
所該說的話將這奧秘.
顯明出來因我為此.
而被.
困鎖.
又是坐牢的.
祈甚麼 祈上帝開.
全道的門.
以至主道.
興旺.
廣傳大地.
你有沒有看過保羅無論是已不鎖書.
還是歌羅西書.
他也是這樣祈禱他的禱告是怎樣來的.
耶路撒冷學的.
真是厲害.
去一次耶路撒冷以後.
就醒目了知道怎樣祈禱.
這是第一點.
接著有第二點.

$^{1121}$保羅在耶路撒冷還學了甚麼.
他見到雅各遜道.
他見到彼得圖森.
他不是說很不公平.
不是的.
主必拯救僕人脫離.
他發現兩個上帝.
都是拯救了.
不過不同方法一個馬上回家.
另一個繼續去.
宣講.
但我們回到我上星期所說的.
我引用.
但以理書.
但以你三個朋友他們所說的話.
尼布加尼撒的回應.
去.
映射到《使徒行傳》裡面.
其實他們無論是.
雅各直接回天家.
又或者彼得.
救出監牢.
繼續見證.
直到他回家的時候.
兩個其實上帝都拯救了.
兩個都是上帝的拯救.
這個和保羅有甚麼關係呢.
因為保羅後來的書信.
又是反映了他真的學了東西.
我們經常讀.
菲臘比書.
菲臘比書裡面.
我們很喜歡引這一段.
從來沒有想過.
他會不會在耶路撒冷學的呢.
因為他在耶路撒冷的經歷.
見到雅各和.
彼得的經歷.
跟著他歷歷在目.
到他在菲臘比書的時候.

$^{1161}$他說.
第一章第十九節.
我知道藉著你們的禱告和耶穌的靈.
終必使我得到釋放.
他沒有說哪一種釋放.
因為下面他說我處於兩難之間.
但兩難兩種.
都是釋放.
無論他死了.
或者他放出來繼續傳揚福音.
都是釋放.
於是第一章第二十三節.
我處在兩難之間.
我情願離世與基督同在.
這是好得無比的.
我很喜歡雅各和.
然而.
我為著你們.
好像彼得這樣活著更加要緊.
都可以.
既然我這樣心順.
我知道仍要留在世間.
與你們眾人一起.
使你們在所信的道上.
又長進.
又喜樂.
無論如何.
我終必得到釋放.
一是直接.
跪天下.
繼續見證直到回家的時候.
Where did he learn this from?.
他哪裡學的?.
他不是自己自創的.
他上過.
少林的.
他在少林學的.
這是第二點.
最後第三點.
基督上聖.

$^{1201}$撒旦墜落.
新約教會.
逃出埃及.
因為當耶穌說.
去看見撒旦墜落的那一刻.
保羅不在那裡.
他要通過紹行傳第十二章.
才看到希律墜落.
撒旦墜落那一件事.
於是讓他.
有福音的.
加油站.
我想先和你們去.
看一段路加福音的經文.
然後再回去試圖.
行傳第十二章.
思想蘇羅在紹行傳第十二章.
所領受有關.
逃出埃及的訊息.
我想帶你們去.
看一個很熟悉.
也是很重要.
路加才有的一節聖經.
路加福音第九章.
二十九節這樣說.
正土地的時候.
他的面貌改變了.
衣服潔白.
講到我們很熟悉的登山面像.
耶穌將近.
去到耶路撒冷.
然後他帶了三個門徒.
彼得,約翰和雅各.
上山登山面像.
然後第三十節.
忽然有摩西和以利亞兩個人.
和耶穌說話.
他們在榮光裡顯現.
談論耶穌去世的事.
就是他在耶路撒冷.

$^{1241}$將要完成的事.
去世的事.
所有聖經學家.
都告訴你.
那個字叫Exodus.
Exodus.
這個字如果字面就這樣解出來.
耶穌和.
摩西和以利亞.
聊天.
聊什麼什麼的.
那個物體叫出埃及記.
當然你可以將它演繹.
去世.
就是要走.
出埃及記就是要走.
耶穌要離開這個世界.
要回家.
但是他們也可以翻為.
耶穌和他們聊的是有關.
一個新出埃及記的事情.
一個終極出埃及記的事情.
什麼是關於出埃及記的事情.
我們可以數到幾樣.
第一.
耶穌要成為逾越節的羔羊.
耶穌要為上帝的子民.
贖罪上帝的羔羊.
Behold the Lamb of God.
who will take away the sin of the world.
就是施帝約翰那句說話.
上帝的羔羊.
耶穌要成為逾越節的羔羊.
第二.
法老要被擊打.
撒旦要被擊打.
第三.
上帝的子民要被領出埃及.
\newpage



\section{出埃及記 14:18}
\label{sec:031984}
\textbf{聖經人聖經事:知道我是耶和華了}~~~~黃克勤牧師博士
\newline
\newline
連結: \href{http://www.church.com.hk/acms/content.asp?site=cdc\&op=show\&type=product\&code=031984\&layout=sermon}{\texttt{031984}} ~~~~ 語音日期: 2026-04-12
\newline
\newline
\hyperref[sec:031983]{\small{< < < PREV SERMON < < <}}
~
\hyperref[sec:index_chronic]{\small{[返順時目]}}
~
\hyperref[sec:031985]{\small{> > > NEXT SERMON > > >}}
\newline
\newline
出埃及記 14:18
\newline
\begin{longtable}{cl}
\hline
\hline
章節 & 經文 (和合本修訂版)\\
\hline
14:18 & \begin{tabularx}{0.7\textwidth}{X} 我在法老和他的戰車、戰車長身上得榮耀的時候，埃及人就知道我是耶和華。」 \end{tabularx} \\ \\
[1ex]
\hline
\hline
\end{longtable}
$^{1}$第一節目又讀聖經了.
生命樹書的網絡.
和大家一起研究聖經人物.
研究聖經故事.
最重要看到背後上主的工作.
和祂的手.
今天就知道我是耶和華了.
這一句是神講的.
我將它不想太長.
就將它簡化一點.
知道我是耶和華了.
這樣他們埃及人就知道.
我是耶和華.
好 留意一下今天的經文.
《出一及記》第十四章.
第十五節.
耶和華對摩西說.
你為什麼向我哀求呢?.
你吩咐以色列人往前走.
你舉手向海生象把水分開.
嘩 這個考驗也大啊.
很有趣的.
神是向摩西說.
你為什麼向我哀求呢?.
其實這個「你」.
摩西沒有哀求.
是百姓而已.
似乎「你」和「摩西」和「百姓」.
是沒有分開的.
摩西就是百姓.
百姓就是摩西.
你為什麼向我哀求呢?.
你吩咐以色列人繼續向前走.
繼續向前走.
向前走又怎樣?.
向前走就是海了.
好 各位群眾大家聽著.
現在神叫我們向前走.
嘩.
摩西sir 你不是啊.

$^{41}$向前走是海來的.
雖然斜斜地下去.
像個灘這樣下去而已.
海來的怎麼走啊.
向前走.
接著他最重要的就是第一點.
你做一件事.
你舉手向海生象把水分開.
原來就是說.
你手裡拿著的那一隻象.
你的手和那一隻象.
是有能力把海水分開的.
嘩 這樣也行.
好 現在我想告訴大家.
你的手和你手裡拿著的那一隻象.
是有分開海水的能力.
你現在試試.
我想我手都震了 不敢試.
我擔心一指.
哎呀 沒分開.
再指一下也沒有.
那怎麼辦呢.
那我想那些石頭會扔在我那裡.
那些人就會殺死我.
那些鞋子就飛過來.
那些拳頭就打過來.
那你舉手向海生象.
哎.
我相信我會有這個恐懼.
不敢舉.
但是摩西敢.
為什麼呢.
他經歷過很多.
你知道他由.
我帶著子女 太太.
要回埃及帶他們走的時候.
他們是騎路人的.
但是摩西手裡拿著的那一隻象.
一直這隻象就是神同在的象徵.
那你要一直拿著.

$^{81}$那你如果今天我們就沒有那一隻象拿著.
拿著一本聖經行不行呢.
那些聖經在手機裡面.
為了拿著手機很奇怪.
那拿著一本聖經.
那就沒有都行.
拿著十字架行不行.
都行.
那其實就是你知道.
它是一樣的證據或者信物.
就是上帝要透過它.
做一些事情.
而曾經也做過很多大事.
什麼事呢.
他就舉手向海生象.
這個憑信心舉起手一伸.
分開.
這樣.
把水分開.
要不要講就不知道.
分開還是說心想分開.
這樣不知道.
總之一字分開.
那些水就分開了.
哇厲害.
這是第一點講的.
就是你要憑信心舉起神的象.
叫海水分開.
這個象與你的手是有這樣的能力的.
這個能力不出於我自己的.
練不到的.
整天在海邊練.
練你的水分開練不到的.
是上帝有這個能力在我們身上.
當我們知道上帝給了我們這種能力的時候.
就可以了.
哎呀.
帶一個人信主很難.
但聖靈降臨在你們身上.
你們就要得著能力.

$^{121}$由猶沙冷撒瑪利.
猶太傳地撒瑪利.
一直到地極.
要作我的見證.
叫我們作見證而已.
沒有一定要帶到祂信主.
是給我們有能力作見證.
於是我們作見證.
我向海生象分開.
水會分開.
我是有上帝的能力的.
好了.
我要使埃及人的心肝壓.
又是這個字.
使.
這個字應該叫做「容陽」.
我任由祂.
這麼厲害.
我現在追你.
我派了這麼多我的士兵一起.
全國軍力來追你.
還追不到你.
咦.
落海了.
照追.
你的心肝壓.
我會讓埃及人的心肝壓.
他們就會跟著下去.
我就要在法老和.
他的全軍車輛馬兵尚得榮耀.
法老的全軍.
全軍特別是一個.
當時的軍隊最威武的是什麼呢.
最威武的是戰車.
騎兵.
那些陸軍很龐大.
那些步兵很龐大.
但是他不出眾.
而且步兵多到不得了.
戰車馬兵就少一點.

$^{161}$就威武一點.
坐得高.
那些衣服又漂亮一點.
我要在他們身上得到榮耀.
全軍車輛馬兵得榮耀.
就是說他這麼威.
其實我要讓他知道.
我是巴閉過他的.
我會讓他們的心肝硬夠膽.
不怕死就死.
一場都來了.
跟馬開.
好了.
我在法老和他的車輛馬兵尚得榮耀的時候.
埃及人就知道.
我是耶和華了.
看到這兩字.
我特別多紅色字.
很多紅色字是因為重要.
我在法老和他的車輛馬兵.
身上得榮耀.
這個就是法老最威的地方.
法老本身.
再加權力.
再加上他的軍力.
我就要在他這樣最榮耀的上面.
我勝過他.
他們埃及人就會知道.
我是耶和華了.
我才是真神.
他們知不知道我是耶和華的呢.
一開始不知道不相信.
後來知道摩西和希伯來人都是信耶和華的.
是啦 知道啦.
好.
不過在這裡.
我想你了解到一件事.
就是經過十災到最尾.
撒掌子之災的時候.
他們就說你去敬拜你的耶和華.

$^{201}$走啦 耶和華.
你去敬拜他 真是的.
行了 走了 走了.
他們都開始知道.
哇 耶和華真厲害.
希伯來人那裡就會光.
我們會黑三天.
黑到逼過來 很黑.
然後撒掌子.
啊 你這個耶和華真厲害.
要金嗎 拿去吧.
要衣服嗎 拿去吧.
要銀子嗎 拿去吧.
拿幾隻羊嗎.
拿一瓶 儘管拿.
快點走 快點走.
哎呀 你和你的耶和華快點走.
不然我們真的沒有了.
我們死定了.
於是呢.
我要在他身上得榮耀.
其實那次已經得榮耀了.
不過法老要不甘心 要追.
那神就說 好.
我要再得榮耀.
在以色列營前行走.
在以色列營的前面.
在以色列營地.
成群人前行走.
那個神的使者.
就轉到他們後面去.
運主也從他們前面.
轉到他們後面立住.
留意一下 再讀一次.
在以色列營前行走的.
是有一個神的使者的.
這個神的使者.
就是神的天使在這裡帶領著.
一個就夠了.
帶二百萬人一個.

$^{241}$因為這個神的天使.
神的使者有什麼呢.
有些解釋因為沒有說清楚.
很多時候說神的使者.
就真的是一個使者的天使.
有很多時候說神的使者救約.
是神自己本身.
可能是這樣的.
那這次不要緊.
既然用神的使者.
我們就用神的使者.
就是說有一個神的使者.
是在以色列營前面行走的.
而跟這個神的使者一起.
就有運主.
很有可能有另一個使者.
又一起行走 就是火柱.
晚上我不知道這個神的使者.
會不會需要換班.
還是說他收起運主.
就拿一條火柱出來.
於是晚上有火柱.
就在這裡帶路照亮等等.
是一個使者還是兩個使者會換班.
不需要理會 不需要計較.
總之在以色列營那裡.
是有神的使者.
OK嗎 OK啦.
那他整個以色列很多.
最前的就是神的使者.
那他一起行的就是運主.
所以就帶著以色列一起行.
現在去到海邊停了.
那這個時間這條運主.
和他的使者就去了以色列人的後面.
那這裡再後面是什麼呢.
再後就是埃及軍隊了.
明白嗎.
所以已經迫到紅海了.
那在前面帶他們去到海邊.

$^{281}$一他就來了後面停定.
那以色列人在他前面.
埃及人在後面.
那些士兵就在後面立住停定了.
在埃及營和以色列營中間有運主.
因為前面帶路嘛.
那是運主的神的使者.
現在去了後面停定.
在運主一邊是黑暗 一邊是發光.
在以色列那邊就是發光.
埃及那邊黑的.
那中夜兩下就不得相近.
那黑的不會去到光那裡.
光不會去到黑那裡.
那大家就說這個使者站在中間.
就分開了.
光的時候他們是看運主的.
所以以色列人看就看到運主在後面.
那埃及人看見什麼就不知道了.
總之就是前面全黑了 看不到.
那全黑了他就不敢前進.
本來有一點點星光那樣照一下.
哎呀 全黑了 看不到東西.
那他當然不敢動.
全暗了.
看得不夠清楚.
兩下不得相近.
就是說光和暗分開了.
摩西就向海生象在最前.
向海生象 紅海那邊向海生象.
耶和華變用大東風.
那時候他們是向著東的.
而他生象的方向就是東 向著東.
紅海生象.
那東風就這樣吹了.
就紙這邊的風在這個方向吹出來.
多向東.
那就使大東風使海水一夜退去.
水便分開.
海就成了乾地了.

$^{321}$留意一下.
這個神蹟真是超級超級偉大的.
就是一紙就開始有大東風.
哇.
有沒有吹到陸上那些以色列人頭髮亂了呢.
是有的.
吹到衣服發出聲音.
已經是這樣 紙 風吹了.
風吹到能夠使海水一夜吹去.
吹到兩邊的水都不能埋回來吹.
而且風大到水分開了.
還是吹成乾地.
海底是濕的.
起碼吹了表面那一層的水氣.
哇 都很厲害.
於是以色列人.
這個神蹟真是很大.
下海中走乾地.
乾的.
不會像泥漿一樣.
嘩 怎麼走啊.
海底濕的.
水在他們的左右作了長援.
那些水呢.
如果你多數看圖畫.
那些水很高的.
哇 鳳凰聲這樣.
很高很高.
那些人就像人仔在下面這樣走過.
不是啊.
如果是中文字.
照中文字解.
長援呢.
援呢.
頹援敗瓦這樣.
援是指矮長.
矮的長.
就是說繼續吹了那些水.
兩邊都還有海水.
不過吹了.

$^{361}$這裡中間沒有.
那些水呢.
就站起來成為一面牆.
不過這面牆未必說.
嘩 很高.
因為用長援這個字呢.
就表示那些水未必很高.
他們中間就走過了.
喂.
二百萬人都走成晚的.
真的是晚上來的.
埃及人追趕他們發露一切的.
馬匹車輛和馬兵都跟著下到海中.
就以為可以走了.
我們追下去了.
如果不追讓他們過了海.
沒有的了.
而且去到別人國境的了.
應該都是他的國境.
快點追過去.
如果現在不追.
他們過去了.
如果海合環環.
我們就要繞一條很遠的路.
他走了的了.
那麼追下去.
到了晨更的時候.
晨更就是說.
就來到天亮那一間了.
耶和華從雲火柱中.
就在雲柱火柱中.
就是說呢.
其實神的使者.
這個剛才那個神的使者.
雲柱去到後面站.
去到後面停了.
其實這個神的使者.
如果你看這裡.
其實這個神的使者.
應該是指神本身的聯想和邏輯是高一點.

$^{401}$那麼神當然是成為使者.
不過祂就是雲火柱這樣出現.
日間雲柱夜間火柱.
那個時間要分開光暗.
就雲柱.
以色列人回頭看.
光的看到雲柱.
埃及官丁看前面.
黑的看到火柱.
於是神就做事了.
令到埃及的軍兵混亂.
令他混亂.
混亂了.
大龍.
怎麼混亂呢.
就搞到他們.
哎呀.
哪個方向.
哎呀.
馬不聽話.
叫他走這邊.
他不斷扯.
又或者呢.
車輪脫落.
輪胎掉了.
真的不知道怎麼辦.
輪胎壞了.
警察又搞不定.
那麼車兵長搞不定.
那些士官搞不定.
不知道怎麼辦.
難以行走.
以前埃及人就說.
不行了.
我們從以色列人面前走.
因為耶和華為他們攻擊我們.
其實士兵更聰明.
不是了.
這次走了.
因為耶和華攻擊我們.

$^{441}$無緣無故車輪弄得很好.
怎麼會脫落呢.
哇.
很危險啊.
馬亂了.
車又脫落了.
不行了.
這次我們快點走了.
因為耶和華攻擊我們.
士兵都看到了.
這個時候耶和華對摩西說.
你向海深象.
叫水營合載埃及人.
並他們的車輛馬兵身上.
那麼這個時間.
耶和華向摩西說.
向海深象.
令到水合番.
就是說摩西已經是最後一個人.
都過了海.
反過來看.
士兵還在海底.
哇.
又脫落了.
拿著車輪不知怎麼辦.
爛了.
馬又不聽話了.
很害怕的時候.
指.
一指.
海水合番.
摩西就向海深象.
第三點要大家看的.
就是到天一亮.
到了天一亮.
海水營救復原.
埃及人避水逃跑的時候.
哇.
不行了.
水衝回來了.

$^{481}$走啊.
想跑回上岸.
耶和華就把他們推回在海中.
就是說用海水將他們沖到全跌了.
全滾在水底.
哇.
在水裡面浸死了.
第三點特別藍色字要大家看.
是到了天一亮.
那麼這個到天一亮.
為什麼要到天一亮呢?半夜不行嗎?.
這個天一亮可能有些屬靈的意義.
因為在詩篇第四十六篇.
都有這一句中文.
不過沒有一個念字.
到天一亮.
那裡說.
有一道河.
死神的城.
去到神的城.
就在那個城的分岔.
令到那個城很受安慰.
到天一亮神必幫助者城.
詩篇第四十六篇第五節.
到天一亮神必幫助者城.
咦.
這個城受慰困.
不過有道神的河圍著它.
保護著.
令它很舒服.
那麼另外呢.
就是到天一亮神必幫助者城.
那為什麼要到天一亮啊?.
天黑不幫得啊?.
正午不幫得啊?.
不是.
到天一亮很多時候是會.
起碼這兩道.
我覺得是指.
是可以隱喻指.

$^{521}$神指定的時間.
神安排的時間.
到了.
海水營救復原.
於是他們圖畫的時候.
就被推翻在海中.
因此.
哎呀神啊.
我很害怕啊.
不知道怎麼辦啊.
你站著看神為你做些什麼救恩.
然後呢.
你閉上嘴巴看.
然後神就說.
你將你信心的象指向紅海吧.
紅海就分開走乾地.
然後埃及人追過來的時候.
時間到了.
你要知道.
神有祂的時間.
到天一亮.
就是神的時間到了.
你就好像看到明亮成神星那樣.
時間到了.
神的時間到了.
海水就會合起來.
將敵人掩蓋.
那麼這個呢.
你要信靠神有祂的時間的.
為什麼還沒做啊?.
為什麼你不用地震.
列在地上.
軍兵跌落就可以嗎?.
我們都不用那麼怕過紅海了.
你做意識的人過紅海怕不怕?.
當然怕啊.
兩邊水啊.
那什麼時候合起來.
就沒有了.
一樣怕啊.

$^{561}$那不如你遠遠地.
就在地震列在地上.
所有埃及軍兵跌落.
哎呀.
那就很開心了.
不行啊.
但是神是要這樣做.
祂有祂的時間.
祂有祂的方法.
到天一亮.
神就會做了.
水就回流.
淹沒了車輛和兵馬.
那些跟意識的人.
落海法老的全軍.
連一個也沒有剩下.
沒有.
如果我看過這套電影.
講摩西呢.
就只有法老一個自己走而已.
他都沒有下去.
在旁邊的山上面看.
哎呀.
悲傷到不得了.
那就唯有走了.
那麼一個也沒有剩下.
意識的人卻在海中走乾地.
水在灘門的左右.
鑿了長門.
鑿了矮牆.
啊.
厲害.
那你知道呢.
就是這句今天要跟大家講的.
特別第三點.
記得.
到天一亮.
時間一到.
就會發生的了.
有信心的.

$^{601}$我們一起禱告吧.
多謝我們的主.
剛才讀了很長的聖經.
就給我們看到一個.
以往發生的事.
是你怎樣.
你的手一步一步一步.
將埃及軍兵帶出來.
讓他們以為.
啊.
以色列人走錯路迷路了.
我們可以追他.
然後他們真的去回心轉意.
覺得我們容讓他們走錯.
追過來的時候.
就將他們引到海中.
你又帶領了.
以色列人是順靠你.
順靠摩西的領導.
主求你自己祝福.
我們真的需要你自己.
給我們的信心.
我們等候主.
到天一亮.
你的時間一到.
你就為我們進行解救.
你的救恩就會來到.
讓我們站定閉嘴.
去看.
去瞻仰你的榮美.
看到你自己的救恩.
心領神會.
知道你是耶和華.
你要在埃及的身上.
在敵人身上.
顯出你是耶和華.
叫人敬畏你.
要叫仇敵都敬畏你.
我們更加要敬畏你.
主啊你的名為聖.

$^{641}$讚美你.
祈禱奉耶穌基督的名字.
阿不得!.
\newpage



\section{尼希米記 4:1-11}
\label{sec:031985}
\textbf{爭戰到底•祈禱包底}~~~~顏沛恩牧師
\newline
\newline
連結: \href{http://www.church.com.hk/acms/content.asp?site=cdc\&op=show\&type=product\&code=031985\&layout=sermon}{\texttt{031985}} ~~~~ 語音日期: 2026-04-12
\newline
\newline
\hyperref[sec:031984]{\small{< < < PREV SERMON < < <}}
~
\hyperref[sec:index_chronic]{\small{[返順時目]}}
~
\hyperref[sec:031986]{\small{> > > NEXT SERMON > > >}}
\newline
\newline
尼希米記 4:1-11
\newline
\begin{longtable}{cl}
\hline
\hline
章節 & 經文 (和合本修訂版)\\
\hline
4:1 & \begin{tabularx}{0.7\textwidth}{X} 參巴拉聽見我們建造城牆就發怒，非常惱恨，並嗤笑猶太人。 \end{tabularx} \\ \\ \relax
4:2 & \begin{tabularx}{0.7\textwidth}{X} 他對他的弟兄和撒瑪利亞的軍兵說：「這些軟弱的猶太人做甚麼呢？要為自己重建嗎？要獻祭嗎？要一日完工嗎？要使土堆裡火燒過的石頭再有用嗎？」 \end{tabularx} \\ \\ \relax
4:3 & \begin{tabularx}{0.7\textwidth}{X} 亞捫人多比雅在一旁說：「他們所修造的石牆，就是狐狸上去也必崩裂。」 \end{tabularx} \\ \\ \relax
4:4 & \begin{tabularx}{0.7\textwidth}{X} 我們的神啊，求你垂聽，因為我們被藐視。求你使他們的毀謗歸於他們自己頭上，使他們在被擄之地成為掠物。 \end{tabularx} \\ \\ \relax
4:5 & \begin{tabularx}{0.7\textwidth}{X} 不要遮掩他們的罪孽，不要使他們的罪惡從你面前塗去，因為他們在修造的人前面惹你發怒。 \end{tabularx} \\ \\ \relax
4:6 & \begin{tabularx}{0.7\textwidth}{X} 這樣，我們修造城牆，整個城牆就連接起來，到一半高，因為百姓一心做工。 \end{tabularx} \\ \\ \relax
4:7 & \begin{tabularx}{0.7\textwidth}{X} 參巴拉、多比雅、阿拉伯人、亞捫人和亞實突人聽見耶路撒冷城牆正在修造，破裂的地方開始進行修補，就非常憤怒。 \end{tabularx} \\ \\ \relax
4:8 & \begin{tabularx}{0.7\textwidth}{X} 大家同謀要來攻打耶路撒冷，使城混亂。 \end{tabularx} \\ \\ \relax
4:9 & \begin{tabularx}{0.7\textwidth}{X} 然而，我們向我們的神禱告，又因他們的緣故，就派人站崗，晝夜防備他們。 \end{tabularx} \\ \\ \relax
4:10 & \begin{tabularx}{0.7\textwidth}{X} 但猶大有話說：「扛抬的人力氣衰弱，瓦礫太多，我們自己不可能建造城牆。」 \end{tabularx} \\ \\ \relax
4:11 & \begin{tabularx}{0.7\textwidth}{X} 我們的敵人說：「趁他們不知道，看不見的時候，我們進入他們中間，殺了他們，使工作停止。」 \end{tabularx} \\ \\
[1ex]
\hline
\hline
\end{longtable}
$^{1}$我弟兄姊妹早晨.
不知大家還記不記得上個月講到.
我用尼稀米見到以色列民.
他們被擄終於歸回,可以回家了.
回到耶路撒冷定居.
應該是很興奮的時候.
因為他們終於可以回家.
可以得著自由,不用再被擄.
誰知回到家的時候.
就發現耶路撒冷城門發生什麼事.
被火燒.
然後城牆就沒有了,倒塌了.
城門,城牆被毀.
成為了家不成家.
完全失去了一份安全保護保障.
好像完全沒有了.
就好像你回到家的時候.
發現這扇門沒有了的時候.
所有人都走進來你家.
這個家就會被人霸佔了.
他們也是一樣.
耶路撒冷城門,城牆被毀的時候.
回到家竟然發現.
那些敵人可以隨時做什麼.
走進來攻打那些以色列民.
所以當時尼稀米很難過.
很傷心的時候.
那些百姓繼續下去的時候.
他們仍然沒有好日子過.
仍然過著一個被擄的生活.
但是尼稀米很愛他們.
他亦很知道上帝給他的呼召和呼召.
尼稀米就說.
「我要決心做什麼?我要重建這個城牆」.
當這份熱誠帶著決心.
帶著祈禱.
帶著行動.
經過當時的「黃波斯王」.
亞達斯西王的批准.
把照書給你.

$^{41}$預備好所有的木材.
去吧!.
尼稀米和以色列民.
在那個時候終於可以什麼?.
鹹魚翻山了.
終於有救了,有希望了.
哇!終於不是這次鹹魚翻山.
哇!接著發生了什麼事呢?.
大英姊妹.
給你一條常識題想一想.
你知不知道為什麼鹹魚不能夠翻山呢?.
(笑).
為什麼鹹魚...現在翻山了.
為什麼鹹魚不能夠翻山呢?.
當你將鹹魚煎的時候.
越煎越香,煎得很香.
當你一翻的時候.
翻到哪裡都翻不到.
為什麼呢?.
沒有理由不知道的在座的婦女們.
(笑).
這是一個常識題而已.
為什麼呢?.
因為....
因為你....
(笑).
怎麼了?.
為什麼會黏鍋呢?.
因為你沒有下油.
你看看這個答案.
所以為什麼鹹魚不能夠翻山呢?.
就是因為你沒有下油.
導致鹹魚黏鍋.
所以怎麼翻山呢?.
離稀米都是一樣的.
我想說.
在這個熱鍋上面.
是一個難以翻山的景地裡面.
我們一起看看當時的背景.
第二章.

$^{81}$二章十節是這樣說的.
它說「和倫人參巴拉」.
以及作神僕的阿滿人.
叫做「多比亞」.
他們聽見有人來.
為以色列人爭取利益.
於是他們做了什麼?.
很生氣.
很老老.
接著到了十節.
「和倫人參巴拉」.
還有這個阿滿人.
「多比亞」.
再加上一個.
叫做阿拉伯人.
叫做「基善聽見」.
於是就恥笑我們了.
藐視.
我們就說.
你們所做的....
你們所做的這些事.
是什麼呢?.
你們在搞什麼呢?.
你們是不是在背叛王啊?.
接著四章七節就這樣說.
有「參巴拉」.
有「多比亞」.
「阿拉伯人」.
再加上阿滿人.
「阿薩達人」.
聽見耶路撒冷城牆.
正在修造的時候.
破裂的地方開始進行修補.
於是做什麼?.
很生氣.
就非常憤怒了.
當時尼希米.
受著一個很大的攻擊.
以及反對.
當那幫仇敵.

$^{121}$聽見耶路撒冷要重建的時候.
那幫人非常生氣.
當那幫人聽見的時候.
去恥笑.
尼希米在這裡搞什麼呢?.
藐視他們所做的一切.
聽見.
以及憤怒的敵人.
是越來越多的.
由這裡第十節.
只有兩個人而已.
「參巴拉」.
以及「多比亞」.
一到十九節的時候.
再加上一個人.
就是「基善」.
到了剛才讀的.
四章七節的時候.
變成五幫人了.
反對的人越來越多.
外來的攻擊勢力.
是四方八面而來的.
你想想.
尼希米當時的場景.
根本是四面楚歌.
當然我們知道.
背後有一個最重要的目的.
就是要迫使尼希米重建.
你停工了.
故意要破壞神的計劃.
破壞神的善工.
所以頂只有一句金句.
你都會背的.
就是「盜賊來」.
然後做什麼呢?.
無非要做什麼?.
偷竊,殺害,毀壞.
所以究竟黑暗的勢力.
是怎樣攻擊人的呢?.
怎樣很有策略地.

$^{161}$去破壞神的計劃呢?.
黑暗勢力.
怎樣去攻擊一個人?.
怎樣很有策略地.
去破壞神的計劃?.
一個屬靈的計劃.
我們今天在尼希米記裡面.
稍後有幾章的經文.
我們就在經文裡面.
一起來辨別.
四章一至四節.
我們一起看.
剛才我們讀的經文.
他說「參巴拉聽見我們建造城牆發怒」.
「非常怒恨」.
「恥笑猶他人」.
「他對他的弟兄」.
「和撒瑪利亞的軍兵」.
說什麼呢?.
你們這些軟弱的猶他人.
你們在幹什麼呢?.
你們為自己重建嗎?.
要限制嗎?.
要一天員工嗎?.
要使土堆裡.
火燒過的石頭再熔嗎?.
只是這兩節.
那惡者.
利用人的嘴巴.
來打擊人的靈性.
一開始的時候.
他和他的弟兄.
還有撒瑪利亞的軍兵.
是一個公開羞辱人的方式.
你們是什麼?.
你們只是軟弱的人.
英文說「feeble」.
即是說什麼?.
只是有氣無力的人.
你們在幹什麼?.

$^{201}$只是這句說話.
已經令人跌入在一個.
失去了尊嚴和價值的裡面.
當你再看第四節.
他們做什麼?.
藐視他們.
毀謗他們.
恥笑他們.
在眾人面前做什麼?.
公開貶低他們的價值.
嘲笑這群有熱誠的人.
他所做的一切.
都是枉言的.
你白費心機.
接著參巴拉一句反問.
說你們這群軟弱的猶太人.
在幹什麼?.
你們為自己重建嗎?.
新標點和合本很到肉地說.
你們做什麼?.
你以為自己大發熱心.
你以為自己重建為主嗎?.
你再看清楚一點.
你只是在保護自己.
你知道這句說話.
正在擾擾他們所做的.
你們在獻祭.
要一天員工.
參巴拉其實是在挑肌.
在問他們.
你們在幹什麼?.
說是正義.
說是為主做.
其實你在保護自己.
還說什麼?.
一天內要做得怎樣?.
然後還沒說完.
接著他說.
要屎土堆裡.
為什麼要這樣說?.

$^{241}$那些是石頭.
在土堆裡.
什麼會在土堆裡?.
不要的東西就扔掉.
丟在土堆裡.
所以那些全部都是被棄.
廢掉過的東西.
而那些是什麼?.
火燒過的石頭.
為什麼要這樣形容?.
火燒過的石頭會怎樣?.
會脆.
會裂.
然後結構會受損.
所以是非常不安全的.
你們這些人在幹什麼?.
竟然用那些.
拿出來的.
火燒過的.
沒用的東西.
變脆的東西.
可以扔掉的東西.
來用.
來重建.
你們不是吧?.
沒有果效的.
浪費力氣.
當你聽到這些說話的時候.
第一.
被貶低的時候.
第二.
打擊你.
很自然就懷疑.
咦?.
我有沒有聽錯神的心?.
我有沒有.
聽錯了上帝給我的召喚呢?.
究竟我是否要重建呢?.
是不是做錯決定了?.
好了.

$^{281}$接著就到第三節了.
第三節到了阿們人多比亞.
他說了一句話.
他說什麼?.
他們所修造的石牆.
那是什麼小動物?.
你知道什麼有意思嗎?.
這裡說.
就算狐狸上去也會怎樣?.
啪!.
做什麼?.
會崩跌.
你知道嗎?.
這一句.
兩句話.
一盆冷水倒下來.
還要踩到你一腳.
他說什麼?.
狐狸小小隻而已.
很輕巧.
扭一下尾巴就會怎樣?.
彈了上去城牆那裡.
你剛才說.
你用的是什麼石頭?.
那些會裂開的.
火燒過的.
那隻狐狸一跳上去.
城牆那裡.
那城牆會怎樣?.
立刻就會倒下.
多大的一個縫隙.
一下子就已經.
一盆冷水倒下來了.
你在做什麼?.
你做的一切白費心機.
是枉然的.
李希斌聽到之後.
真的灰.
真的灰爆.
你知道嗎?.

$^{321}$當你試試進入經文的時候.
你就會明白他當時的心路歷程.
很失望.
沒有了盼望和堅持.
那我還有什麼信心.
繼續去建造這個城牆呢?.
什麼心機.
什麼熱誠.
都沒有了.
沒有信心和鬥志.
弟兄姊妹.
你知不知道.
這些不負責任的說話.
是那惡者常常用的伎倆.
他說完之後.
不用負責任的.
但是你聽到之後.
誰負責任呢?.
就變成你自己要負責任了.
惡者很有計劃.
去打擊有熱情.
有熱情.
你很想服侍主的那些.
他就來打擊你.
他貶低你的身份.
我們看看這個表你就會知道.
他貶低你的身份.
他羞辱你.
這個軟弱的人.
你沒什麼用的了.
就好像.
你有沒有聽過這句話.
可能我們小時候也經歷過.
生狗什麼?.
說得這麼快.
莫非你也經歷過.
生狗什麼?.
叉燒.
好過生你.
說什麼?.

$^{361}$你也不是一個人.
人也不如.
然後開始質疑.
究竟我是誰?.
我做錯了什麼?.
我是否聽錯了神的心意?.
我開始懷疑神給我的照明.
究竟我做這個服侍.
是否用這個方法?.
很多東西來懷疑自己.
最後會怎樣?.
不要再做了.
那怎樣?.
爛尾.
那惡者最想看到我們的侍宮是怎樣的?.
爛尾.
沒用.
沒效果.
那惡者攻擊人的靈性.
由貶低人的價值開始.
一旦我們信心動搖.
我們信心的根基是哪裡來的?.
是我們跟神之間的關係.
當這個關係被動搖的時候.
信心被打擊的時候.
很自然就會懷疑神.
也會懷疑自己.
當失去信心.
人就沒有動力.
不知道在做什麼.
那惡者就會偷走了.
人對神那份渴慕的心.
那惡者就會偷走了.
人對神那份熱情.
敷衍敷衍就算了.
最後.
你就會變成一條.
什麼?.
鹹魚.
你還記得周星馳說了一句名言嗎?.

$^{401}$記不記得?.
如果人什麼?.
沒看過?.
如果人沒有什麼?.
如果人沒有夢想的話.
你.
然後怎樣?.
你跟一條鹹魚.
都沒什麼分別.
今天我說.
如果人沒有信心和熱情.
所以其實.
我們都會跟一條鹹魚.
都沒什麼分別.
所以你明白.
為什麼上次在講道的時候.
我們定的主題是.
當我們這個生命裡面.
沒有了或者少了.
失去了主那份的同在.
生命沒有牆的時候.
沒有了到成牆.
那你會怎樣?.
你會很快收檔.
先跟旁邊的人說一遍.
生命沒有牆.
然後怎樣?.
很快收檔.
所以今天要重新建立成牆.
這就是那惡者.
他怎樣打擊人的心理戰術.
這是內在的.
接著我就要說外在.
在環境裡面.
剛才我們讀的第四章.
也讓我們看到.
敵人裡面用了最毒的方法.
亦都是魔鬼最大的目的.
就是要怎樣?.
使城來混亂.

$^{441}$當你看到這個地圖的時候.
你就發現.
原來參巴拉,多比亞,阿拉伯人.
阿姆人,阿薩達人.
全部都是怎樣?.
東南西北五組的勢力.
東南西北全部的敵人.
都圍著了.
接著聽見耶路撒冷.
正在來收做的時候.
破裂的地方開始進行修補.
於是做甚麼?.
很生氣.
非常憤怒.
在你重新決定.
我要站起來的時候.
重新站立起來的時候.
開始補破口.
要重建的時候.
說真的還是危危乎乎的.
很容易倒下的.
在軟弱的時間.
攻擊是最猛烈的.
在危危乎乎走來倒下的那一刻.
第八節.
祂說.
大家同謀.
要來攻打耶路撒冷.
使城來混亂.
祂不是一個炮彈炸下來.
祂是聯手,合力,同謀.
再加上剛才說的心理戰術.
製造內部一些恐慌,混亂.
你自然會失去方向.
你自然會恐懼.
你開始懷疑神.
究竟上帝是否在這裡?.
就好像之前我們在COVID-19時經歷過.
搶米,搶廁紙.
甚麼都搶.

$^{481}$連去銀行禁機.
都沒有錢.
錢都被搶走了.
在這個恐懼焦慮之中.
第十一節.
祂說.
我們的敵人說.
趁他們不知道.
趁他們看不見.
做甚麼?.
我們就進入中間.
涉過了.
接著做甚麼?.
殺了他們.
使工作停止.
這裡說.
趁他們看不到.
不知道的時候.
在一個長期恐懼不安之中.
在一個內憂外患之下.
慢慢煎熬到你的身心靈.
去到一個內耗的狀況.
就是經歷著一個很大的恐懼.
一個壓迫,威嚇.
甚麼動力.
甚麼熱誠.
呼召都沒有.
最後就變成了一條.
放棄掙扎的鹹魚.
當內心.
在這個交織著心理戰.
再加上外在的壓迫.
當強功無效的時候.
那惡者就會用一些更加卑劣的方法.
這個手段是甚麼呢?.
我們一起看第六章.
第六章是這樣說的.
他說.
參巴拉多比亞.
阿拉伯人羈善.

$^{521}$還有我們其餘的手的.
聽見我已經建造了城牆.
沒有破裂之處在其中.
那時候我還未在城門安門扇.
第二節.
參巴拉和羈善就派人來見我.
然後他說甚麼.
「喂!請你來約一個地方」.
「這個地方在哪裡?」.
「就是叫做阿羅平羽村莊見面」.
「為甚麼?」.
「因為他知道會有一班人要殺他」.
於是我派使者到他們那裡說.
「我正在進行大工程」.
「我在建造」.
「你在弄我幹甚麼?」.
「不能下去」.
「我怎能離開」.
「下去見你們」.
「讓這個工作停頓」.
於是第四節.
他們這樣派人.
嘩!一次,兩次,三次.
到了第幾次?.
到了第四次.
都有同一番說話.
「過來吧!約你在這裡見見面」.
尼希米也一樣跟他說.
「我正在重建」.
「我不會下去」.
然後第五節.
嘩!鍥而不捨.
參巴拉做甚麼?.
「第五次」.
「同樣派僕人來見我」.
手裡面還拿著甚麼?.
一個未封的信.
打開信先看看.
第六節寫著.
「列國中有風聲」.

$^{561}$基善也是這樣說.
尼希米.
原來你是跟猶太人謀反.
所以你來到建造城牆.
據說原來你做甚麼?.
你要做他們的王.
並且你派先知.
在耶路撒冷.
指著你們對孫子說.
「嘩!這次真棒」.
「再猶大有王了」.
如今這些話.
必傳到波斯王那裡.
現在請你來.
跟我們一起商議.
然後尼希米派人去他那裡說.
「你所說的這些事」.
「最抵過去」.
「一概沒有」.
「是你心裡甚麼?」.
是假的.
「是你心裡捏造的?」.
完全是說.
「捏造壞話」.
「來陷害尼希米」.
「老屈他」.
「你來作王而已」.
「所以你才這麼熱心」.
「要來重建而已」.
「你想一班以色列民」.
「來叩拜你」.
「嘩!這還不是乘機陷害他嗎?」.
你想想.
如果這個謠言.
真是傳到波斯王那裡.
他會有甚麼反應?.
如果你看回第二章.
是誰讓尼希米去重建的?.
是這個王.
他看到尼希米很傷心,很難過的時候.

$^{601}$「好,你去吧」.
「你去重建,回去耶路撒冷」.
「原來你背後的目的」.
「你是為了想做王」.
「嘩!你完全是謀反」.
「信他一成就怎樣?」.
「傷目失明」.
「簡直是一個大的欺騙」.
「嘩!分分鐘會怎樣?」.
「這個王馬上派人做甚麼?」.
「殺了尼希米了」.
這是離間的手段.
弟兄姊妹.
這裡短短幾節.
其實我們要學習一個功課.
這個功課是甚麼呢?.
是「小心聽」.
罵惡者要破壞人的關係.
他只要做一招就可以了.
這招是甚麼呢?.
就是「以當眼」.
你明白我說甚麼嗎?.
耳朵是用來做甚麼?是聽的.
眼睛是用來做甚麼?是看的.
但當你的耳當眼的時候.
「喂!你只是聽而已」.
「不是真的」.
「那些東西可能是假的」.
「你沒有事實查證」.
「你沒有看清楚」.
「你聽完就當是真的」.
那在說甚麼呢?.
其實背後罵惡者.
要破壞彼此之間的關係.
彼此之間的離間.
去破壞當中的關係.
「以當眼」.
第十節.
第十節在這裡說.
「我到了,尋找了一個人」.

$^{641}$「離稀米,去尋找了一個人」.
這個人叫做「事馬雅」.
去了他家.
接著「敲門啦,唔該開門啦」.
誰知道他不開門.
「閉門不出」.
他就說「事馬雅說的」.
「我約你去一個地方」.
「我們可以去神的殿那裡」.
「在神的殿那裡」.
「鎖住那個殿門」.
「我知道外面那些人做甚麼」.
「要殺你啊」.
「還要晚上來殺你啊」.
接著離稀米就說.
「像我這樣的人會逃跑嗎」.
「像我這樣的人」.
「豈能進入殿裡保命嗎」.
「我當然不去啦」.
接著十二節.
離稀米很清醒.
屬靈的辨別.
「看清楚了」.
「看啊,神並沒有派他」.
「是他自己說甚麼」.
「預言來做甚麼」.
「攻擊我」.
接著再說.
「是多比亞和參巴拉做甚麼」.
「收買我回來的」.
「收買他,那個人」.
「收買他的目的」.
「叫我懼怕」.
「令我犯罪」.
「留下一個壞的明星」.
「好讓他們可以繼續做甚麼」.
「毀謗我啦」.
為了要殺離稀米.
今次罵惡者用了一招.
叫做借刀殺人.

$^{681}$他收買了一個.
看起來很像一個屬靈人.
如果你還記得.
施瑪雅其實是九月的先知.
同名.
誰知這是收買回來的.
假先知施瑪雅.
他說的話.
是一種偽裝的神的聲音.
以為是神的旨意.
其實是被收買的謊言.
是謊言.
偽裝成整個屬靈的場景.
「來吧」.
「我們去聖殿一起見面」.
「不用怕」.
「神同在的地方很安全」.
「把門也鎖上」.
聽姐妹.
再留意.
施瑪雅說甚麼.
剛才她約離稀米去哪裡.
哪裡見面.
在神的殿見面.
接著我就問.
有誰才可以去到.
進入神的殿.
是誰.
是利未人.
是祭司.
離稀米你是誰.
你只是一個普通人.
你只不過是一個王的酒正.
替王來試酒.
來服侍祂.
你不是利未人.
如果.
離稀米真的聽了.
施瑪雅說.
走進去聖殿的時候.

$^{721}$會發生甚麼事.
他便犯了.
這個律法上所說的事.
律法說甚麼.
民數記.
犯埃近的愛人.
觸碰回眸的.
必備做甚麼.
處死.
所以.
如果.
他離稀米.
聽到施瑪雅說.
來吧.
在聖殿見面的時候.
我真的走進了聖殿.
參巴拉.
多比亞.
理直氣壯地.
一下子煽動群眾.
你離稀米.
你褻瀆神的罪名.
我可以合法地做甚麼.
處決你.
馬上可以立即.
暗殺他都可以.
於是.
如果是這樣.
就回應到這段經文.
十三節說.
離稀米就會留下.
壞的名聲.
離稀米就犯罪.
好讓他們.
跟著做甚麼.
往後的日子.
不斷不斷地說.
不斷不斷地.
來到去.
毀謗離稀米.

$^{761}$剛才說.
是借刀殺人.
借施瑪雅的刀.
去殺離稀米.
但如果.
當離稀米真的.
走進了聖殿.
是自己走進了聖殿.
那把刀變了誰.
就是變了.
自己殺自己.
最後這把刀是.
離稀米的手.
自己殺自己.
所以整個.
策略性地.
一個破壞神的計劃.
一個屬靈的破壞.
最狠毒的地方.
一起看下一頁.
罵惡者.
他不需要直接.
去傷害我們.
或者殺害我們.
他用迷惑.
誘惑我們.
離間我們.
給錢的關係.
叫我們做甚麼.
走進去死路.
是自墮墳墓.
是自己殺自己.
自墮墳墓.
當離稀米.
他很清楚知道.
自己站對位.
做對神的功.
知道上帝很喜悅這件事的時候.
罵惡者.
總是會找方法.

$^{801}$令你失敗.
打擊你.
折騰你.
最終的目的.
你收手.
不用做了.
爛尾了.
罵惡者就開心了.
拍爛手掌吃花生.
接著我做了一個表.
這個表就是整個表.
讓你看到整個策略性.
屬靈破壞的狀況.
就是我剛才說到的經文.
那些重點都在這裡.
這套策略裡面.
其實是一個很恐怖的地方.
你有沒有發現.
是一層一層遞上去的.
是先讓你覺得.
你不行的了.
你靈性裡面退步了.
你不是上帝愛的孩子.
越來越自卑自憐.
於是我給一個盲字.
去形容這個.
靈性裡面的墮落軟弱.
一個人裡面.
好像披上面紗一樣.
我盲了.
我看不到原來我是上帝所愛的孩子.
我看不到我是一個軍專的弟子.
我看不到在這個聖潔的國度.
是屬神的子民.
我什麼都看不到.
我只看到自己.
多麼差多麼軟弱.
多麼被人奚落.
盲了.
接著再讓你覺得.

$^{841}$環境很混亂.
很多外在的東西.
一波未停一波又起.
處理完這一樣.
又要處理那樣的東西.
令你很忙.
身心都軀疲.
自己都不知道自己是誰.
接著又要搞很多的東西.
要處理很多的東西.
最後讓你發現.
讓你覺得他選擇那條.
你以為是對的道路.
但其實那條是錯的道路.
那我只用了詭詐.
令你忙言模糊不清.
令你沒有方向.
跌入迷惘的裡面.
弟兄姊妹.
這三個格.
我們曾經跌入.
哪一個陷阱的裡面呢.
又或者.
我們現在又在.
哪一個陷阱的裡面呢.
這個圖就幫助我們.
不管在我們屬靈生命的裡面.
甚或是.
在我們所發生的經歷.
一個屬靈的辨別.
這個表.
幫助我們一個屬靈的辨別.
叫我們小心.
不要跌入陷阱的裡面.
好了.
講一個故事給你聽.
有一個人做夢.
他做夢見到一間很大的教會.
建築很漂亮.
金碧輝煌.

$^{881}$裡面的擺設很精緻.
再看一看.
他發現有一隻魔鬼在裡面.
在做什麼呢.
原來是在睡覺.
然後又做夢.
他見到另外一間小教會.
走進去時嚇到.
很多隻魔鬼.
睜大眼睛.
四圍盯著他.
這個人百思不得其解.
當他再看看那些迴縱.
大教會真的很漂亮.
但裡面的信徒.
彼此之間是很冷漠的.
一打開聖經.
或者牧師講道的時候.
廚師在做什麼呢.
在釣魚.
在欺負.
在熬夜.
無心來到他手上.
所以一隻魔鬼就完成了.
一隻魔鬼還要睡覺.
都沒有影響力.
多棒.
但小教會就不一樣了.
火熱敬拜.
信徒立志.
每天讀十章聖經.
很投入的生活.
彼此之間有很多愛的關係.
又去三福.
去報道.
去傳福音.
所以我們的努力.
魔鬼同樣都很努力.
很勤力.
剛才說盲.

$^{921}$亡.
亡言.
不斷地引誘人.
所以騙徒的手法是怎樣的呢.
騙徒的手法是.
跟你旁邊的人說一次.
既然.
對於這三樣.
又盲.
又令你亡.
然後又亡言.
迷亡.
另外一個亡.
誰能夠可以在.
這個盲.
中間這個亡.
亡言的裡面.
迷惑的裡面.
清醒.
暫穩呢.
當尼希米經歷.
參巴拉.
多比亞的嘲笑.
尼尼希米.
馬上做了一件事.
每一段的聖經都是.
他立刻.
就轉向上帝.
四章四節他這樣說的.
我們的神.
求你垂聽.
因為我們被藐視.
求你使他們的委託.
歸於他們的頭上.
令他們的被擄之物.
成為這個掠物.
接著.
去到.
當尼希米.
經歷聯手同謀.

$^{961}$要來攻打耶路撒冷的時候.
四章九節.
又來一個禱告.
我們向我們的神.
祈禱.
當他聽到.
捏造謠言.
去陷害尼希米的時候.
六章九節.
他說.
主啊.
求你堅固我的手.
當他面對.
經歷這個.
假先知的時候.
六章十四節.
我的神.
求你記得.
那些人所做的事情.
接著.
下面一個表.
同樣所有的經文.
都是出於這句說話.
我的神.
求你紀念我.
求你因這事.
來紀念我.
求你記得.
他們的罪.
而整個尼希米記.
的最後一句.
傳卷書的最後一句.
就是紅色這句的金句.
就是我們剛才背的金句.
我們一起讀.
準備起.
我的神.
求你紀念我.
施恩於我.
再來多一次.

$^{1001}$肯定一點.
準備起.
我的神.
求你紀念我.
施恩於我.
你知道嗎.
尼希米每一個經歷.
面對這個屬靈的破壞.
他立刻做的事情.
是馬上立刻.
立刻轉向一個.
最強大的力量和保護.
他不是一個.
很冗長的祈禱.
主啊救我啊.
這個人什麼什麼.
我現在怎樣怎樣.
不是.
他是什麼.
很簡短.
全部我印出來的經文.
是不是很短.
只有一兩句.
三或者四句.
簡短的祈禱.
第一個的祈禱重點.
就是.
我的神.
求你紀念我.
施恩於我.
就成為他生命裡面.
最強大的力量和反擊.
單單只是這一句.
求你紀念我.
背後在說什麼.
紀念.
主知道.
主看到.
主明白.
當主看見.

$^{1041}$明白的時候.
主就會怎樣.
施恩.
當我們進入到.
尼希米的世界.
內外交鋒的裡面.
他還要照顧這班百姓.
嘩.
又要很有理智.
又要很有吉士.
又要好好安排人.
怎樣擺在哪裡.
來看守城牆.
要怎樣建.
他專心集中.
不受外敵影響.
要站穩的時候.
因為尼希米他相信.
他知道自己做對.
真戰到底.
所以尼希米他相信.
這刻真戰到底.
但最重要是什麼.
祈禱包底.
和你旁邊的人說一次.
這八個字.
真戰到底.
祈禱包底.
在說到最後的時候.
我想和你分享一個故事.
這個故事叫.
懸崖上面最後的一根火柴.
究竟在說什麼呢.
他說有一位老伯伯.
他一輩子都守著.
這個懸崖上面那座燈塔.
所以他是一個.
很忠心的守望者.
所以他只有一個任務.
任務就是很重要的.

$^{1081}$他要確保燈塔上面那盞燈.
一定要亮.
因為他知道這盞燈很重要.
所有船隻都是靠這盞燈.
平平安安回家.
誰知有一個晚上.
有一場很恐怖的暴風雨.
突然間殺到.
燈塔上面的玻璃.
全部扔爛了.
老伯伯整個人濕透了.
一直顫抖.
一直拼命地去修補.
偏偏在這個時候.
家裡傳來一個消息.
你的孫子病了.
重病.
你快點回家.
這一刻.
老伯伯真的內外交鋒.
很疲倦.
外在的裡面.
看見這樣狂風暴雨的時候.
都不知道可以怎樣.
這個時候.
那惡者的聲音就來了.
在他耳邊響起.
伯伯.
算了.
這麼大雨.
你在這裡搞什麼.
漆黑一片.
你見不見到有船.
你都看不見了.
你快點回家.
你的孫子在等你.
浪費氣力.
你現在在這裡.
於是這個老伯伯.
看著自己已經皺了皮的手.

$^{1121}$心裡面很自憐.
他就在想.
我都這麼老了.
我還可以熬到多少個晚上.
不如算了.
不如放棄這個使命.
當他手上握住最後一根火柴.
他看著這個深淵裡面的黑暗.
正想放下這根火柴的時候.
他突然聽到一個很微弱.
是很尖銳的笛聲.
是什麼聲音.
是求救信號.
於是這個老伯伯.
突然好像清醒了.
他知道如果他放棄.
外面.
大海上面.
那些人命.
不知道會發生什麼事.
他馬上之後.
不再看環境.
不再看自己的軟弱.
他很大聲地來說.
主啊!求你紀念我.
施恩於我.
老伯伯這個祈禱.
他不是求主醫好我的孫子的病.
他不是求主止息這個風浪.
他只是求主啊!.
求你什麼.
紀念我.
神就會施恩於我.
當他望著自己.
一個皺了皮.
顫抖的手.
最後他擦亮了.
他最後一根火柴.
點燃起這個燈芯.
光就出現.

$^{1161}$外面那些漁船.
就是因為這裡光.
避開了那些礁石.
終於獲救.
主啊!求你紀念我.
施恩於我.
這個施恩.
點燃起我心裡面的火.
擦上我最後一根火柴.
主啊!點燃起我心裡面的火.
他重新看清楚自己的身份.
我是一個忠心的守望者.
丁姊妹.
無論惡者.
怎樣打擊人.
怎樣部署打擊你.
無論環境.
怎樣叫我們.
做一條鹹魚.
但你也要記得.
今天離希米記最後的討問.
是什麼?.
我的神.
接著是什麼?.
求你紀念我.
施恩於我.
你知不知道.
紀念 Remember這個字.
就是說神記得你.
當神記得你的時候.
就是說祂看見你.
祂明白你現在的狀況.
於是祂就會施恩.
就好像離希米的經歷一樣.
當你被欺壓.
被打倒.
感覺人生失去方向.
喪志.
失去信心.
盼望到谷底.

$^{1201}$甚至你說我都抬不起頭.
可能你在職場裡.
可能你在家庭裡.
你覺得很孤單.
或者在壓力的底下.
你發現人與人的關係.
很多的分化.
很多的結黨.
有很多是非的裡面.
或許你每一天.
忙過不停的處理.
很多的事.
令到你精神都衰弱.
令到你生命裡那種內耗.
請你不要忘記.
我的神啊.
求你紀念我.
施恩怨.
所以請你在活日的時候.
擦亮你這一根的火柴.
你向主說.
主啊.
求你點燃起我心裡的火.
\newpage



\section{約翰福音 6:1-15}
\label{sec:031986}
\textbf{你生命可以參與神蹟}~~~~鄭偉誠牧師
\newline
\newline
連結: \href{http://www.church.com.hk/acms/content.asp?site=cdc\&op=show\&type=product\&code=031986\&layout=sermon}{\texttt{031986}} ~~~~ 語音日期: 2026-04-19
\newline
\newline
\hyperref[sec:031985]{\small{< < < PREV SERMON < < <}}
~
\hyperref[sec:index_chronic]{\small{[返順時目]}}
~
\hyperref[sec:031987]{\small{> > > NEXT SERMON > > >}}
\newline
\newline
約翰福音 6:1-15
\newline
\begin{longtable}{cl}
\hline
\hline
章節 & 經文 (和合本修訂版)\\
\hline
6:1 & \begin{tabularx}{0.7\textwidth}{X} 這些事以後，耶穌渡過加利利海，就是提比哩亞海。 \end{tabularx} \\ \\ \relax
6:2 & \begin{tabularx}{0.7\textwidth}{X} 有一大群人因為看見他在病人身上所行的神蹟，就跟隨他。 \end{tabularx} \\ \\ \relax
6:3 & \begin{tabularx}{0.7\textwidth}{X} 耶穌上了山，和門徒一同坐在那裡。 \end{tabularx} \\ \\ \relax
6:4 & \begin{tabularx}{0.7\textwidth}{X} 那時猶太人的逾越節近了。 \end{tabularx} \\ \\ \relax
6:5 & \begin{tabularx}{0.7\textwidth}{X} 耶穌舉目看見一大群人來，就對腓力說：「我們到哪裡去買餅給這些人吃呢？」 \end{tabularx} \\ \\ \relax
6:6 & \begin{tabularx}{0.7\textwidth}{X} 他說這話是要考驗腓力，他自己原知道要怎樣做。 \end{tabularx} \\ \\ \relax
6:7 & \begin{tabularx}{0.7\textwidth}{X} 腓力回答他：「就是兩百個銀幣的餅也不夠給他們每人吃一點點。」 \end{tabularx} \\ \\ \relax
6:8 & \begin{tabularx}{0.7\textwidth}{X} 有一個門徒，就是西門‧彼得的弟弟安得烈，對耶穌說： \end{tabularx} \\ \\ \relax
6:9 & \begin{tabularx}{0.7\textwidth}{X} 「這裡有一個孩子，帶著五個大麥餅和兩條魚，但是分給這麼多人還算甚麼呢？」 \end{tabularx} \\ \\ \relax
6:10 & \begin{tabularx}{0.7\textwidth}{X} 耶穌說：「你們叫大家坐下。」那地方的草多，人們就坐下，男人的數目約有五千。 \end{tabularx} \\ \\ \relax
6:11 & \begin{tabularx}{0.7\textwidth}{X} 耶穌拿起餅來，祝謝了，就分給坐著的人，也同樣分了魚，都照他們所要的來分。 \end{tabularx} \\ \\ \relax
6:12 & \begin{tabularx}{0.7\textwidth}{X} 他們吃飽後，耶穌對門徒說：「把剩下的碎屑收拾起來，免得糟蹋了。」 \end{tabularx} \\ \\ \relax
6:13 & \begin{tabularx}{0.7\textwidth}{X} 他們就把那五個大麥餅的碎屑，就是大家吃剩的，收拾起來，裝滿了十二個籃子。 \end{tabularx} \\ \\ \relax
6:14 & \begin{tabularx}{0.7\textwidth}{X} 人們看見耶穌所行的神蹟，就說：「這真是那要到世上來的先知！」 \end{tabularx} \\ \\ \relax
6:15 & \begin{tabularx}{0.7\textwidth}{X} 耶穌知道他們要來強迫他作王，就獨自又退到山上去了。 \end{tabularx} \\ \\
[1ex]
\hline
\hline
\end{longtable}
$^{1}$平安.
好,感恩來到大家當中.
我今早的時候見到抱架.
我就記起一件事.
我就說.
其實.
我真是.
這個也是一個神蹟.
因為從小到大.
我家是一個武系社會.
即是女人說了算.
因為我爸是不在家的男人.
我四個姊姊.
四個姊姊都很強.
所以從小到大.
基本上我被訓練.
我不用說的.
他們全說就可以了.
所以我說.
我是從小到大.
我都可以拿一個最佳的案證.
最初其實是.
這是我們米飯班主.
我們來吃飯的時候.
我們就招呼這些學生.
當時他們很節儉的年代.
這些學生.
那些女生通常就叫一個正面.
我還記得一元.
那些男生大食一點.
就叫兩個正面.
就是兩元.
我們就趕著回去幫忙.
所以我記得有一次.
我有很多功課.
不知道要做些什麼.
我跟媽媽說不行了.
我要做功課.
做不來 我沒有幫忙.
然後我媽媽那天整天都不理我.

$^{41}$她覺得我不乖.
因為你不幫家人.
所以我的童年就是.
經常在店裡幫忙.
很不喜歡學.
那些朋友.
那些朋友有時候很羨慕我們的.
你也好.
喜歡吃什麼就吃什麼.
相反我說你也好.
你喜歡玩什麼就玩什麼.
就是你對著那些東西.
其實你是對得多.
你都不喜歡吃的.
童年就是這樣過.
然後到我大一點.
到我中學畢業的時候.
環球盒都清拆了.
我媽就一直覺得做生意是最好的.
所以當時我們就搬進新市鎮.
去大埔裡面.
我們就繼續做這個.
那時候就不是叫冰室了.
潮流不同了叫茶餐廳.
我們就開了一間茶餐廳在大埔.
不過在大埔的時候.
這間茶餐廳的位置不太好.
沒什麼生意.
我們很短時間就轉型了.
做一個小菜館.
香港那段過程裡面.
因為進入大埔.
我們在環潭的時候.
都是星期天就放假.
都回到教會.
進入大埔.
真的做了一天生意都不行.
我記得在過程裡面.
真的朝七晚十一.
幾個月都沒有假期.

$^{81}$我當時就覺得.
上帝是不是做了一個行業.
是一個受戶助的行業呢.
我回不了教會.
大英展會他們山長水遠.
進來大埔探望我們.
我們又沒有空招呼他們.
久而久之人家來過一兩次都.
怎樣都疏遠.
你放假的時候.
真的有些時間你去找人.
人家又沒有空.
覺得好像被遺棄.
更加在這個過程裡面.
我還記得.
為了要怎樣去省錢.
為了省三百元那時候.
自己每一晚要去倒垃圾.
推垃圾去垃圾站.
但是接著就被人罵了.
被一個倒垃圾的姐姐就說.
三百元下欄錢你都不給我賺.
那時候我是二十歲出頭.
有很多事情我覺得.
真的很灰暗.
接著又被人怎樣呢.
被人家收保護費收拖地.
我們在環潭這麼大的時候.
當時那些小飛.
是我們從小一起帶的.
所以他們不會騷擾我們.
大家都認識的.
但是去到大埔新的地方裡面.
面對被人家收拖地.
接著剛才說轉型做小菜館的時候.
很多事情都不懂得做.
沖奶茶咖啡.
餐蛋治沙嗲牛肉麵那些.
我懂得做的.
因為那些小時候就被訓練.

$^{121}$但是要埋鍋炒東西.
拿著大鍋炒東西.
真的不懂得.
當時因為轉型要做小菜館的時候.
就請一些不同的廚師.
或者廚房老.
我很討厭那些廚房老為什麼呢.
因為當時如果生意好的時候.
他們就會用那些碼兜.
不知道大家知不知道什麼是碼兜.
就是煮東西他們會撿一些.
他們會在廚房冰冰地說髒話.
他們都不懂得做什麼什麼.
就是很多很難聽的說話.
沒有生意的時候.
又在那裡說.
不如你們不要做了.
讓我們做吧.
一批廚師也是這樣.
兩批也是這樣.
可不可以有一班來到.
比較好一點.
他們就說.
你們的貨源不太好.
不如這樣吧.
我帶你們去街市賣貨.
帶他們去街市.
當時跟著廚房老去街市賣貨.
他就說.
阿成你在這邊拿些什麼.
我去那邊買些什麼.
久而久之.
原來知道為什麼他這麼熱心.
現在香港也是這樣的.
你去街市賣東西.
你問人家拿單.
老闆請給單.
你知不知道他會怎樣問的.
怎樣回答你.
你想寫多少.

$^{161}$跟著或者給張白頭單.
你自己寫就行了.
原來在街市賣東西.
有很多這些什麼.
有很多這樣的好處.
當我們發現他的時候.
他就說.
我不做了.
香港我不知道是怎樣.
香港到今時今日飲食都是這樣的.
我們叫推工即是辭工.
七天.
即是七天通知就可以了.
無論你做了多久都是.
但是你這樣的情況.
他說不做.
你明天還敢不敢讓他上班.
你都不敢讓他上班.
但是沒有人做事.
你開不開到檔呢.
很多的壓力.
不過在這同時又發生一件事情.
就是有一天有一個外國人.
來到我們的店鋪吃東西.
進入大門的時候.
當時我坐在收銀位.
這兩個樓面就看到有個外國人進來.
他們又想著我的英文好一點.
其實我都不太行.
他就叫我招呼他.
我走過去的時候.
他說 勞駕.
即是他會說.
廣東話.
OK 搞定.
寫完單 入單.
吃完.
他過來給錢.
他就說.
我都知道你們信耶穌的.

$^{201}$我想和你們聊聊天.
我不知道你們怎麼知道我信耶穌.
我們沒有掛十字架.
沒什麼.
我說.
好好好.
下午我們兩點半下場的時候你才來.
現在很忙.
他說 好.
他真的兩點多三點就來了.
我們店鋪前面.
是一個停車場.
當他遠遠.
我看到他走過來的時候.
我其實那時候在賣帳.
在賣中午的帳.
你猜我當時是什麼心情.
是不要學武和傳道人聊天.
那就錯了.
我看到他來的時候.
我馬上從後門撩走他.
為什麼呢.
我不是討厭他.
不要說.
我都信了耶穌.
不用和我說耶穌.
甚至當時我都覺得.
你都不知道.
你這個鬼佬我都不知道你信什麼.
總之就是你不要騷擾我.
因為當時的心情是.
我面對太多問題了.
你不要來煩我.
不過.
我那次撩走了之後.
因為我們打開門做生意.
他就繼續來.
久而久之我們熟絡了.
剛才說當我遇到那些問題的時候.
我們都會和他分享.

$^{241}$我記得.
被人收陀地那次是怎樣呢.
有一天.
我們跟別人說.
我們面對這些困難.
他就說.
當有朋友遇到問題找你的時候.
我相信你都會說怎麼樣.
我們祈禱加備主.
是很實在的.
祈禱很重要的.
祈禱加備主.
他說你祈禱.
他說不要緊.
明天我一早來陪你去等這個人.
他那天很早來.
一直等.
等到下午的時候.
這個收陀地的人真的出現了.
我們當天就坐在這張桌子上.
當然.
這是我們近來借回那個地方的場景.
我們已經沒有做生意了.
就坐在這張桌子上.
當時.
那個人來到的時候.
看到我旁邊有一個外國人.
你看三十年前他比現在肥一點點.
款式也是差不多.
你覺得他會想到什麼.
你覺得他像不像警司.
我估計他那時候想到.
為什麼你找了警司來跟我談判呢.
當然.
他不知道我是一個牧師的傳道人.
他來到的時候.
他也做完他做的事.
因為他大哥叫他來收帳.
他就說.
真的要收保護費.

$^{281}$我當時一聲也沒有說.
這個西教士就跟他說.
他們有耶穌保護就很足夠了.
不需要你們保護.
這樣也可以.
接著又怎樣.
拍桌子.
接著你就是不能給.
不給.
看看我們接下來要做什麼.
這個人真的是徒弟.
後來我真的知道他大哥是那種人.
不過很奇妙.
那次之後.
我在那裡三年沒有一個耶穌有料到.
他真的是那個先知.
他又怎樣.
要推他為王.
聖經說.
耶穌是不是王.
耶穌是王.
不過耶穌不是要做這件事.
當時說到逾越節接近.
逾越節接近的時候.
其實這幫以色列人.
我記得逾越節.
紀念逾越節就是.
當時他們的祖先在埃及圍爐的時候.
神藉著摩西拯救他們.
這件事情就是逾越節.
他們紀念神的拯救.
所以這個不是一個好的時機嗎.
我們現在一起上耶路撒冷.
我們就推耶穌做王.
來一次革命.
推回羅馬政府.
這一定是好事.
當時是他們的心願.
要脫離羅馬這個統治.
這正是一個大好的機會.

$^{321}$耶穌你去吧.
那些歡呼聲多大.
耶穌有沒有被歡呼聲推去.
應該沒有.
耶穌是退去後面.
去祈禱.
去不容讓.
這些聲音.
這些歡呼聲,鼓勵聲.
去逃避神的心意.
去完成人的心意.
我讀書的時候.
有一本書對我來說是很大的提醒.
那本書叫做.
《神的將領兩冊》.
它說到在美國很多次復興的當中.
神興起了一些將領.
這些人可能是一些無名的小子神.
他們靠著神去帶來一些教會.
一些復興.
在過程中.
他沒有滅退一些陰暗面.
就是這些神這麼重用的人.
在路上走著走著的時候.
有些可能會走.
走出自己的意向去走.
他沒有跟著神的路去走.
雖然他帶來復興.
但可能他在那件事中.
他自己要面對很多其他的跌倒.
當然當中的一些在跌倒的情況下.
他再經歷的神是很矮的.
這就是一個很大的提醒.
今天我們在一個新的侍功.
或者我們面對一些復仕的情況下.
我們開始的時候.
我們一定.
我相信每一個基督徒一定.
我們會祈禱得很清楚.
坐在家裡祈禱.

$^{361}$領受.
好了 我的神開路了.
我開始了侍功了.
一邊做的時候.
哇 你做得真是好啊.
你繼續去吧.
哇 你就開始心亡.
是的 你怎樣做.
很好的聲音.
或者說你做得真是差.
你不要繼續這樣做.
很多周圍.
可能有不同的聲音在你身上.
這些聲音可能就影響了你自己.
或者你有一些經驗.
你覺得是.
我繼續用我的經驗去做就可以了.
我們忘記了.
要一步一步緊緊跟隨神.
因為我們第一步.
我們本來對著終點.
走向上帝的心意裡面.
但原來我們不斷不斷走.
有時候我們真的聽到聲音.
走歪了的時候.
我們不在禱告裡面修正.
我們就會越走越過.
所以耶穌教導就是.
我們要在過程裡面不斷.
我們要尋求神的心意.
去走向他的心意裡面.
我們要除去很多的雜聲.
那些雜聲是在影響我們.
可能是在影響人的心意.
在聖經裡面還有一個例子.
就是說當日約書.
他要帶領一班的以色列人.
進去迦南.
其實約書是很荒謬的.
因為之前一直跟著誰.

$^{401}$跟著摩西.
聖經說摩西是很厲害的.
摩西是埃及王子.
他學了埃及裡面很多的知識.
他有行軍打仗什麼的知識.
約書是誰.
約書是做奴隸的.
是那段時間跟著摩西那四十年.
學到一些東西.
但是摩西死了.
他要帶領以色列人.
去迦南那部入迦南.
對著耶利哥.
一個這麼高大的城怎麼打呢.
聖經記載.
他上了山上看著耶利哥.
他當時應該禱告.
他不知道怎麼做.
神經當時經常跟他說什麼.
你要剛強 大大 壯膽.
所以神不斷地加油.
那次裡面.
突然間他看到一個人出來.
他說你幫我們還是幫敵人.
我要看耶和華的懸殊.
他要帶他們怎麼去打這場仗.
怎麼打這場仗呢.
那場仗就經典了.
打耶利哥.
我們今天有時候看聖經的時候.
覺得真的很精彩.
但是你細心想一想的時候.
怎麼會這樣打仗的.
他們怎麼打仗的.
你記不記得.
祭司先走.
所以做祭司的先走先死.
先走 拿著角.
圍著那裡.
第一天繞一圈.

$^{441}$不要出聲.
敵人在城牆上.
看著這幫人在做什麼呢.
第二天又這樣繞一圈.
裡面的人不會想一想.
他們打仗的嗎.
扔那些石頭都扔死你了.
是不是.
每天都這樣繞.
到第七天.
繞了七個圈.
接著那場仗又怎麼樣.
就贏了.
因為約書我相信他沒有方法.
所以他跟著耶和華懸殊的方法.
因為他知道這才是最可靠的方法.
那場仗打贏了.
不過人打贏仗很危險的.
這場仗不是最精彩.
最精彩是第二場仗.
第二場仗.
大家記不記得那個城叫什麼城.
那個叫做艾城.
他們贏了那場仗.
很心虛.
這個城很小的.
我們派二三千人去搞定它.
不過在那場戰爭裡面.
他們一敗塗地.
當然我們看聖經的時候.
他們說因為阿光犯罪.
但是最重要的時候.
就算阿光犯罪.
約書也去帶領一幫以色列人.
他們好好尋求上帝的時候.
我相信他們一早就已經知道.
不用反過來才知道.
贏了.
我們經歷過成功之後.
如果我們忘記繼續依靠神.

$^{481}$是依靠自己的方法.
依靠自己的智慧.
依靠自己的經驗的時候.
我們很容易就進入.
這個網絡裡面.
所以主耶穌在這裡.
很大的提醒我們就是.
我們跟隨神是要.
一步一步去尋求神的心意.
鼎姐妹.
今天同樣.
上帝要在你和我當中.
我一同去經歷.
這個挑戰就是.
讓神的神蹟.
在我們生命上面.
一同的彰顯.
在過程裡面.
耶穌.
你記不記得.
我們剛剛講過復活節.
耶穌已經復活.
當天耶穌復活的時候.
祂向門徒顯然.
聖經說.
在大使命之前.
大使命我們都記得.
耶穌眾前來對他們說.
天上地下所有的權兵.
都給了我們.
所以我們要去.
使萬民作祂的門徒.
奉父子聖名的名.
給他們施洗.
凡我吩咐他們的.
都教訓他們遵守.
我就常與你們同在.
直到世界的末了.
我們很熟悉大使命.
不過在大使命之前.

$^{521}$那一節的聖經怎麼說呢.
祂說耶穌來的時候.
門徒就拜去.
那些門徒就拜去.
然而他們.
有人依然.
疑惑.
我不知道在你.
跟隨主義的過程裡面.
你服侍主義的過程裡面.
你有沒有曾經疑惑過呢.
我告訴你我有.
我有疑惑的.
神好也知道.
我們也有時會疑惑的.
不過你看這段經文的時候.
後面沒有說.
耶穌如何處理他們的疑惑.
耶穌就是這樣.
祂就告訴他們.
天上地下所有的權兵.
都給了我.
你們要去.
使萬民作祂的門徒.
祂說你們去吧.
奉父子聖名的名.
給他們施洗.
凡我吩咐你們的.
都教訓他們遵守.
我就常與你們同在.
原來.
除去疑惑的就是怎樣.
就是去的.
當你走出去的時候.
他們就會看到.
神的同在.
就看到神蹟.
在他們的生命上面發生.
今天.
怎樣破除疑惑呢.

$^{561}$是要冒險.
走出去.
相信.
我相信這裡.
我們每一個人.
你是不是已經相信耶穌了.
是的.
你舉手給我看看.
你已經相信了耶穌了.
多謝大家.
相信.
原來不夠的.
相信.
不夠.
什麼叫相信不夠呢.
是的.
我告訴你.
你信不信得舉一張椅子.
可以承載一個.
二百多磅的人呢.
那你說.
我也相信.
這個牧師起碼也八九磅.
坐下去.
椅子也承受得了.
這個叫做相信.
好的.
我現在跟大家玩個遊戲.
如果你可以的話.
你跟我一起站起來.
你可以的話.
你就站起來.
我們來玩一個.
挺高難度的挑戰給大家.
大家都站穩了.
你可以就跟著我這樣做.
好 我現在想請你.
坐在椅子上.
不過屁股不要碰著椅子.
坐在椅子上.

$^{601}$不過屁股不要碰著椅子.
但你要坐到.
好像剛才那樣的形態.
剛才的形態.
但你的屁股不要碰著椅子.
你做到你就做.
做不到你就棄權.
我不做了.
可以嗎.
好像三秒也不夠.
好 大家請坐.
我剛才說相信很容易.
我聽這麼多見證.
我相信.
不過當你自己.
要進入這個相信的時候.
剛剛那個叫做什麼.
我不單只相信椅子.
我要信靠椅子.
我就要靠著它.
靠著它其實是怎樣.
很舒服的.
不靠著它是怎樣.
你剛才無影椅的時候.
你試過.
不靠著它.
其實很辛苦的.
如果我今天說.
我相信耶穌.
但我不信靠耶穌.
我是不是很辛苦.
很辛苦的.
不過這個信靠.
是要冒險的.
因為你可能會怕跌下去.
不過你再考慮一次的時候.
你就怎樣.
你就會更加信靠.
丁姐妹.
今天你生命.

$^{641}$上帝說.
很想你和祂.
參與在祂的足跡裡面.
靠著祂.
今天你知道.
你來到這個城市.
這個地方裡面.
你不是要逃避什麼事情.
是上帝要你來這裡當中.
將一個很重要的信息.
就是你已經得到的.
耶穌基督的恩典.
去和那些未得到的人分享.
就好像當日.
耶路撒冷的信徒.
好像受逼迫要四散.
但四散的目的是什麼.
四散的目的不是逃難.
四散的目的是.
四散去到不同的地方裡面.
將福音去和那些人分享.
我們看《仕逗人傳》的時候.
你看著.
當他們去到哪裡的時候.
聖靈的能力就帶著他們去到那裡.
將福音和他們分享.
所以鼓勵你.
無論你什麼年紀.
我們一起怎麼樣.
我們一起去上路.
去經歷耶穌.
和我們一起同行的時候.
見證神的神蹟.
\newpage



\section{約翰福音 21:15-25}
\label{sec:031987}
\textbf{當失敗遇上復活主}~~~~禤凱榮牧師
\newline
\newline
連結: \href{http://www.church.com.hk/acms/content.asp?site=cdc\&op=show\&type=product\&code=031987\&layout=sermon}{\texttt{031987}} ~~~~ 語音日期: 2026-04-19
\newline
\newline
\hyperref[sec:031986]{\small{< < < PREV SERMON < < <}}
~
\hyperref[sec:index_chronic]{\small{[返順時目]}}
~
\hyperref[sec:031988]{\small{> > > NEXT SERMON > > >}}
\newline
\newline
約翰福音 21:15-25
\newline
\begin{longtable}{cl}
\hline
\hline
章節 & 經文 (和合本修訂版)\\
\hline
21:15 & \begin{tabularx}{0.7\textwidth}{X} 他們吃完了早飯，耶穌對西門‧彼得說：「約翰的兒子西門，你愛我比這些更深嗎？」彼得對他說：「主啊，是的，你知道我愛你。」耶穌說：「你餵養我的小羊。」 \end{tabularx} \\ \\ \relax
21:16 & \begin{tabularx}{0.7\textwidth}{X} 耶穌第二次又對他說：「約翰的兒子西門，你愛我嗎？」彼得對他說：「主啊，是的，你知道我愛你。」耶穌說：「你牧養我的羊。」 \end{tabularx} \\ \\ \relax
21:17 & \begin{tabularx}{0.7\textwidth}{X} 耶穌第三次對他說：「約翰的兒子西門，你愛我嗎？」彼得因為耶穌第三次對他說「你愛我嗎」，就憂愁，對耶穌說：「主啊，你無所不知，你知道我愛你。」耶穌說：「你餵養我的羊。 \end{tabularx} \\ \\ \relax
21:18 & \begin{tabularx}{0.7\textwidth}{X} 我實實在在地告訴你，你年輕的時候，自己束上帶子，隨意往來；但年老的時候，你要伸出手來，別人要把你束上，帶你到不願意去的地方。」 \end{tabularx} \\ \\ \relax
21:19 & \begin{tabularx}{0.7\textwidth}{X} 耶穌說這話，是指彼得會怎樣死來榮耀神。說了這話，耶穌對他說：「你跟從我吧！」 \end{tabularx} \\ \\ \relax
21:20 & \begin{tabularx}{0.7\textwidth}{X} 彼得轉過身來，看見耶穌所愛的那門徒跟著，就是在晚餐時靠著耶穌胸膛說「主啊，出賣你的是誰」的那門徒。 \end{tabularx} \\ \\ \relax
21:21 & \begin{tabularx}{0.7\textwidth}{X} 彼得看見他，就問耶穌：「主啊，這個人怎樣呢？」 \end{tabularx} \\ \\ \relax
21:22 & \begin{tabularx}{0.7\textwidth}{X} 耶穌對他說：「假如我要他等到我來的時候還在，跟你有甚麼關係呢？你跟從我吧！」 \end{tabularx} \\ \\ \relax
21:23 & \begin{tabularx}{0.7\textwidth}{X} 於是這話在弟兄中間流傳，說那門徒不死。其實，耶穌不是說他不死，而是對彼得說：「假如我要他等到我來的時候還在，跟你有甚麼關係呢？」 \end{tabularx} \\ \\ \relax
21:24 & \begin{tabularx}{0.7\textwidth}{X} 這門徒就是為這些事作見證、並且記載這些事的，我們知道他的見證是真的。 \end{tabularx} \\ \\ \relax
21:25 & \begin{tabularx}{0.7\textwidth}{X} 耶穌所行的事還有許多，若是一一都寫出來，我想，就是全世界也容不下所要寫的書。 \end{tabularx} \\ \\
[1ex]
\hline
\hline
\end{longtable}
$^{1}$上個星期我們看過彼得的故事.
他三次不認主.
可以說是跌入了人生的低谷.
那段時間我相信彼得真是自怨自艾.
他嘗試回到以前那種固有生活.
嘗試忘記以前的失敗或者痛苦.
不過似乎不是很有用.
當耶穌在加利利的海旁.
第三次向門徒顯現的時候.
奇怪在於耶穌沒有責備這幫門徒.
反而在岸邊為這幫人預備早餐.
彷彿在彼得跌倒的時候.
耶穌伸出他的手幫彼得站起來.
這個是整個故事的上集.
我們今天來看這個故事的下集.
這個故事說他們吃完了早餐.
耶穌問彼得要向他兒子西門.
你愛我比這些更深嗎?.
聖經這裡所說的這些.
不是很清楚說到耶穌在說什麼.
什麼這些.
有些解經家就說.
會不會是耶穌在說那些魚啊.
餅啊等等呢.
又或者會不會是耶穌在說.
他們謀生的船啊網啊等等呢.
又或者耶穌會不會在說.
你對耶穌的愛.
大不大得過深不深得過.
其他門徒對耶穌的愛啊.
聖經沒有說彼得有追問下去.
所以我們大概可以想像.
當時面對面的交流裡面.
彼得已經知道會議到.
耶穌基督說的是什麼.
不過無論這個書所說的這些.
是指到什麼東西.
耶穌基督的這個問題呢.
可以說說.
一下子就刺中.

$^{41}$戳中彼得最痛的地方.
為什麼.
因為我上次也說過.
很有可能彼得其實呢.
想過無數次.
他自己都問過自己.
究竟我愛不愛這個.
我的主耶穌基督呢.
我愛不愛他呢.
我相信彼得自己知道的.
他是愛這個耶穌基督的.
不過行動做出來.
卻是三次不認主.
大家對這個情況.
有沒有一些似曾相識的感覺.
記不記得.
主耶穌自己都這樣說.
他說.
你們一定要警醒禱告.
免得入了迷惑.
因為為什麼.
你們心靈.
固然願意肉體卻軟弱了.
就是我們不是沒有心.
是有心不過沒有力.
也是使徒保羅的經驗.
使徒保羅在羅馬書第七章這樣說.
他說.
立志為善.
我想做好人就由得我.
不過行出來.
卻由不得我.
所以我所願意的善.
我又不會做.
我不想做的惡.
我反而去做.
保羅在那裡.
將他人生最無奈的事情.
就像一種吶喊一樣說出來.
不是我們不想去愛.

$^{81}$不過很多時候我們發覺.
我們得完美的生命.
是沒有辦法能夠活出完美的愛.
使徒約翰在約翰一書那裡.
跟我們說.
人能夠去愛.
為什麼.
因為神先愛我們.
神對人的愛彰顯了.
究竟愛是什麼一回事.
神的愛也給我們有能力.
去活出一個耶穌基督的愛.
如果我們留心看看.
耶穌和彼得這個對話當中.
其實耶穌基督是告訴我們.
愛的能力是什麼.
如果我們看這個十五到十七節.
這三節的經文.
耶穌是問了彼得三次.
你愛我嗎.
三次.
第一次彼得怎麼回答.
彼得說主啊是的.
你知道我愛你.
大家我想大家留意.
在耶穌發明的時候用的那個愛.
和彼得回答的時候那個愛.
在希臘原文是兩個不同的字.
耶穌用的那個是阿加皮.
我們翻譯中文是愛.
不過彼得回應的.
用那個希臘文那個愛.
翻譯成愛中文.
不過用的字是完全另外一個字.
兩個字有什麼不同.
如果說阿加皮.
耶穌基督用的那個愛.
是指到一個很深沉.
但是有犧牲.
是一種沒有條件的愛.

$^{121}$是一個很高層次的愛.
不過彼得用的那個愛.
是泛指一些兄弟之間深厚的愛.
大家看到分別了.
耶穌問彼得你願不願意對我.
是一種沒有保留.
沒有條件的愛.
而不是單是朋友兄弟之間的愛.
不過彼得的回應.
就讓我們看到.
當時彼得內心的心境是怎麼樣.
在彼得三次不認耶穌的時候.
我們可以看到.
似乎彼得從一個什麼都說出來.
什麼都第一時間就有反應的人.
慢慢變成一個很內斂的人.
他的答案.
你可以說是一個很誠實的答案.
彷彿他好像跟耶穌說.
主啊 我的失敗已經證明了.
我沒有能力對你付出一種.
沒有條件沒有保留的愛.
我現在只能夠給你.
主啊 就是一種破碎了的友誼.
很明顯彼得.
那個時候的彼得.
很明顯他沒有寬恕自己.
他仍然好像一個監牢那樣.
被他的內疚被他的羞辱所困住.
出不來.
我相信彼得是知道.
他好像在坐牢那樣.
不過無能為力.
他做不到任何事.
為什麼.
坐牢這條鑰匙都開了門.
他沒有這條鑰匙.
那你說.
這條是什麼鑰匙.
這條鑰匙其實.

$^{161}$如果我們認真看耶穌的問與答.
其實就在耶穌的問題與答案裡面.
好了 你說.
耶穌的問題問什麼.
耶穌說.
你愛我 比這聲更深.
沒錯.
愛其實就是這條鑰匙.
可以將彼得.
從這個監獄.
這個羞辱與內疚的監獄裡面.
釋放出來.
不過他沒有用到這條鑰匙.
沒有用到.
那怎樣才可以用到這條鑰匙.
怎樣才可以.
將這條鑰匙.
將上帝的愛.
從他的內疚裡面釋放出來.
有一個女孩子.
很愛她媽.
知道她媽那天生日.
打算做一頓豐富的晚餐給她.
趁她媽出門的時候.
用盡她的百般的武藝.
來到廚房裡面.
打算做一頓好的晚餐給她.
不過我從來沒有煮過東西.
這個小女孩.
也不懂得煮東西.
當然我們可以想像.
廚房弄得一塌糊塗.
當我母親回來的時候.
發覺整個廚房一股燒焦的味道.
那些煲煲餐餐全部燒焦了.
那個女孩子站在那裡.
不知道怎麼辦.
心裡就想.
這次慘了.
搞得廚房這樣.

$^{201}$一定母親回來見到.
一定會很生氣很生氣的.
後悔為什麼要試呢.
正當她在徬徨無助.
不知道怎麼辦的時候.
她媽就回來.
進屋見到.
接著她就跟女兒說.
接著就說.
你想做一個晚餐給媽媽吃.
那個女兒就說是.
接著母親沒有責備她.
走過去就抱著女兒.
就跟女兒說.
這個是我最好的晚餐.
母親沒有讚揚到她的失敗.
弄得廚房一塌糊塗.
不過.
母親確是很確認.
這一個努力背後的女兒.
對母親的愛.
接著她就跟女兒說.
不要緊.
我們一起清理廚房.
然後我教你怎麼煮東西.
當母親見到女兒的努力.
背後的愛.
而當女兒在她失敗的時候.
通過母親的擁抱.
而感受到母親的愛的時候.
女兒的失敗.
不能夠像彼得那樣.
囚禁在她的內疚與羞辱之中.
因為愛.
因為她的失敗.
不單止不能夠囚禁這個女兒.
反而成為一個原因.
讓她以後繼續的嘗試.
所以丁姐妹.
彼得沒有用這條鑰匙.

$^{241}$去將他自己從.
內疚與羞愧當中釋放出來.
為什麼.
因為她在那一刻.
仍然不能夠肯定.
究竟她的主有沒有接受她.
她的主有沒有寬恕她.
那三次的不認主.
不過.
一個個早餐改變了.
在岸上耶穌基督.
為他們所做的早餐.
顯示了證明了.
主已經寬恕了彼得.
已經接受了彼得對她的愛.
當愛被確認.
當愛被感受到的時候.
愛就能夠去恢復.
失敗所做過的破壞.
愛永遠都是.
醫治破碎最好的良藥.
大家同意嗎.
我們再看下去.
好了.
彼得確認了.
他說.
你知道我愛你.
接著主怎麼回答.
主就說.
你餵養我的小羊.
這個問題和答案.
差不多類似這樣.
重複發生了三次.
如果我們看第十五十六.
和十七這三節的時候.
三節都記載了問題.
和答案.
當耶穌問你愛我.
彼得說愛.
接著耶穌就說.

$^{281}$你餵養我的羊.
當然了.
大家要明白.
耶穌不是叫彼得.
你幫我搞個牧羊生意.
包你發財.
耶穌不是這樣的意思.
耶穌是作為一個比喻.
你要照顧我的跟隨者.
為什麼三次.
有些解經家就說.
三次就是對照耶穌基督.
對照彼得三次不認耶穌.
三次問彼得.
我們要留意.
就算是對照彼得的.
三次不認主.
耶穌三次的發問.
並不是要挖苦.
或者羞辱彼得.
你那次令我這麼難受.
我現在就令你這麼難受.
不是這樣的意思.
三次是代表著.
耶穌基督願意.
完全處理這個問題.
接納和寬恕彼得.
很明顯.
耶穌基督很清楚.
彼得內心的掙扎.
究竟是什麼.
他三次的問題.
跟著三個的命令.
是告訴彼得.
他對耶穌基督的愛.
已經將他從.
羞辱與內疚的困境裡.
釋放出來.
讓彼得可以準備好.
去接受耶穌基督的命令.

$^{321}$準備好來事奉他的主.
這三個命令是什麼.
我們說了三個問題.
三個答案.
這三個命令是什麼.
其實很相似.
如果我們看的時候.
一個是餵養我的小羊.
第二個是牧養我的羊.
第三個是餵養我的羊.
三個都很相似.
不過是有一點點不同.
如果我們看第一個命令的時候.
就是說你餵養我的小羊.
小羊是指什麼.
小羊這裡的英文是lamb.
大家可以更加相信.
這裡所說的小羊.
是說一些初信的信徒.
或者寧命還是很幼嫩的基督徒.
他們通常是撒旦最容易.
也是最想攻擊的對象.
為什麼.
因為這些是.
對真理命未有好心的認識.
他們還不能夠靠自己.
對聖經真理的認識來判斷.
真假和對錯.
記不記得在希伯來書的作者.
曾經有一番這樣的說話.
希伯來的作者在第五章是這麼說的.
他說看看你們學習的功夫.
你們應該學了這麼久.
應該做師傅了.
誰知道是什麼.
誰知道你還要人.
將聖經的真理.
最簡單的真理.
要另外教導你們.
就成為.

$^{361}$必須吃奶.
你們現在應該是在吃乾糧.
但你們仍然是要吃奶.
因為你們還不認識.
仁義的道理.
就好像小孩子一樣.
不能像成人一樣吃乾糧.
因為你們的心竅雜念得.
還未夠通達.
所以不能夠分辨很大.
其實屬靈的生命.
和我們地上的生命很類似.
大家都知道財生的嬰兒是怎樣.
不能夠吃乾糧.
他們只能夠吃奶.
所以初信主的基督徒.
其實都是需要我們成熟的基督徒.
一匙一匙的靈糧餵給他們.
因為他們還未對真理有足夠的認識.
不能夠分辨真正的對與錯.
所以最容易受到其他異端邪說影響的信徒.
就是初信的信徒.
而耶穌就叫彼得.
你要幫我牧養這些幼嫩的信徒.
讓他們成為成熟的信徒.
簡單來說就是.
你要把神的話語.
清晰地正確地教導這群初信的信徒.
第二和第三個命令.
這不是小羊.
是 sheep.
翻譯是羊.
羊代表著一些比較成熟的基督徒.
甚至是整個的信仰群體.
而這兩個命令也有點不同.
第一個命令是牧羊.
第二個是餵養.
簡單來說就是把神的話語真理講解明白.
讓這些信徒得著熟齡的餵養.
無論是初信的信徒.

$^{401}$還是信了很久.
很成熟的信徒.
都需要吃東西.
人一樣.
無論你是初生.
還是現在我們七老八十.
都需要吃東西.
只不過是吃不同的東西而已.
那你有什麼分別.
財生的嬰兒就喝奶.
成人就吃乾糧.
不過無論什麼年紀.
我們都需要熟齡的真理來餵養我們.
如果餵養是指神的話語得到講解明白.
牧羊就代表什麼.
牧羊就是指我們需要一個好像牧羊人一樣.
帶領基督徒走正確的人生方向.
當有危難來到的時候.
有牧羊的人幫我們保護.
我們走錯路.
我們做錯事的時候.
有人提點我們.
矯正我們.
當我們受傷的時候.
有人為我們包紮傷口.
這個就是牧羊的意思.
這個就是牧羊人要盡的責任.
簡單來講.
耶穌基督就是叫彼得.
你要照顧我的信徒.
好像牧羊人照顧他的羊群一樣.
大家留意.
也值得留意.
這裡耶穌的形容是什麼.
牧羊我的羊.
大家有沒有留意.
不是彼得你的牧羊.
彼得我給你的羊.
不是這個意思.
是牧羊我的羊.

$^{441}$意思就是說.
羊群真正屬於誰.
屬於耶穌基督.
所以教會的領袖.
只不過承擔著屬靈的責任.
為耶穌基督在地上牧養他的信徒.
所以彼得很清楚這個教導.
為什麼我知道他很清楚.
因為如果我們看彼得前書.
彼得自己這樣說.
他那時候教導其他的教會領袖.
就是當我們牧養我們中間神的群羊的時候.
我們是需要愛著神的心意.
不要出於勉強.
也不要貪財.
也不要濫巴斃地要核製你的羊群.
反而要怎樣.
要作群羊的榜樣.
當耶穌三次確認了彼得.
很清楚地說.
你知道我愛你的.
你還問.
好了.
確認了.
因著耶穌就叫他來牧養他的羊.
牧養他的羊.
不過要留意一件事.
千萬不要搞錯.
耶穌基督叫彼得給他一個工作.
給他一個使命.
叫他牧養他的羊.
不是要懲罰彼得.
要弄清楚.
為什麼.
因為今天我們大家都很清楚.
當我們做錯事到法庭的時候.
你一會兒就坐牢.
不坐牢又怎樣.
有些什麼.
社會服務令.

$^{481}$你要幫社會做些事.
這是懲罰.
但是耶穌基督叫彼得做的事.
不是教會服務令.
耶穌基督要彼得去投入做牧養的工作.
是委任彼得有一個使命.
在神的國度裡面建立的時候有他的份.
這是彰顯了神對彼得的愛.
也讓彼得有一個機會去經歷.
或者甚至領受更深的祝福.
做牧者是優差.
做牧者是不是優差.
大家覺得做牧師容易嗎.
或者不要說牧師.
做教會的領袖是不是一個很容易的工作.
如果說容易的通常都沒做過.
我告訴大家其實不容易.
一點都不容易.
特別是當教會有不同的意見.
特別是當教會大部分的信徒.
都並不是成熟的信徒.
教會有時候特別是做領導的工作是不容易.
想像一下.
在教會裡面當你做一個決定的時候.
你有多可見是100\%完全的人都是贊成的.
很少的.
通常有些人說不是.
我應該覺得走東邊是對的.
另一班就說不是.
應該走西邊是對的.
那怎麼辦.
都要做決定的.
任何一個決定.
我告訴大家.
都有一班人開心.
有一班人是不開心的.
有時候牧師去聽.
你這麼這樣的提議.
我接受這個提議.
這個牧師這麼耳朵軟.

$^{521}$別人說他們聽什麼.
好了.
到牧師說.
我覺得你的意見都不是很適合.
這個牧師獨斷獨行.
不聽意見.
大家看不看到.
其實做領袖的工作.
不是一個簡單的工作.
剛才我所舉的例子.
只不過是眾多的困難.
眾多的挑戰裡面的其中一兩樣.
不過奇怪在於.
你以為給得不知道嗎.
就算你當時不知道.
給得一路做下去的時候都知道.
但是奇怪還有一件事.
給得沒有轉行.
沒有做了幾年牧師之後.
搞不定.
轉行太辛苦.
這班人搞不定就搞不定.
我寧願回去打魚.
沒有.
給得在整個人生遇下.
到他完結生命那一天.
都仍然是忠於神的使命.
忠於神的教堂.
為什麼他可以這樣做.
並不是因為他有個人的能量.
並不是因為他有個人的耐力.
留意.
留意一件事.
當耶穌叫給得做這件事的時候.
通常如果我們叫人做事會怎樣.
大家有沒有試過.
幫公司招聘工人.
我就做過聘牧小叔.
在另一個教會做過很多年.
第一件事要看的是什麼.

$^{561}$你拿一個履歷表來看看.
我看看你的履歷表.
有沒有這樣的工作經驗.
我看看你的履歷表.
有沒有這樣的學歷.
都夠齊了.
然後我想面試你的.
看看你能不能夠勝任.
那我們要問的就是.
為什麼耶穌基督.
沒有問過給得的能力與經驗.
是不是能夠勝任這個.
做牧者的工作.
有沒有問到.
有.
不過不是問他的經驗與學歷.
問的是什麼.
你愛我比這些更深.
耶穌基督唯一的條件.
你不需要愛我.
是的時候.
你就要承擔這個責任.
所以其實很簡單.
整個故事讓我看到一個.
很機要的真理.
因為什麼真理.
對耶穌基督的愛.
是我們侍功的基礎.
沒有了這個基礎的時候.
一切侍功都是枉然.
也都由於我們對耶穌基督的愛.
讓我們能夠走在一個.
從失敗可以得著治療的路上.
是不是.
我們再看下去.
好了.
說完這個使命之後.
耶穌在第十八節就跟彼得說.
老友.
將來有這樣的事情會發生.

$^{601}$什麼事情.
就是說.
你年老的時候.
你要伸出手來.
別人要把你束上.
帶你到不願意去的地方.
這個是他的預言.
預言將來你年紀大了.
老了的時候會有什麼廢事發生.
去到一些你不想去的地方.
還要被人托住.
就是說什麼.
將來很有可能是.
代表到將來他會被逮捕.
被關在監牢裡面.
或者甚至是流放.
而寫這本書的使徒約翰.
更加補充多一句.
這話是指著彼得要怎樣死.
就是預言彼得的殉道.
而事實上根據這個教會的歷史.
彼得是真的刀掛在十字架上離開世界的.
你說一個這樣的消息.
是不是令人雀躍.
聽到情緒立刻低沉.
奇怪的是.
耶穌說完之後.
沒有跟他說彼得不用怕.
沒有這些事情.
耶穌接著說的是什麼.
你跟從我吧.
說完之後.
直接說你跟從我吧.
我問你.
如果你是彼得.
聽完之後.
將來會發生你身上的事情.
接著耶穌說.
你跟從我吧.
我想問你.

$^{641}$你會不會跟從.
我相信彼得那時候心大心小.
想想落落.
如果我接受這個使命.
很大一件事.
正在他囚障的時候.
正在囚障的時候.
他看一看周圍.
看到耶穌所愛的門徒約翰.
接著就說.
主啊.
這人將來如何.
不知道大家有時候.
會不會真的很喜愛.
彼得那種性格.
想到什麼就說.
自己不要看.
只要不要太.
不要跟自己的關係太親就行了.
保持一些距離.
但是彼得的性格真的有點可愛.
他想到什麼就什麼.
他沒有隱藏他的感受.
你覺得.
彼得為什麼要問這個問題呢.
主啊.
此人.
你所愛的門徒.
他將來如何.
為什麼彼得要問這個問題呢.
我相信很簡單.
如果代入去想的時候.
很簡單.
神說你給我一條這麼難走的路.
將來要遜渡.
那麼這個呢.
你不是給同樣困難的路去走的嗎.
我相信如果.
當時耶穌答的答案是說.
不是.

$^{681}$他走的路舒服很多.
的時候.
你以為給的反應是什麼呢.
說不是吧.
主啊.
你給我一條這樣的路.
那個這麼容易走.
不公平.
我想問大家.
你覺得這個世界有沒有絕對的公平.
這個世界有沒有絕對的公平.
如果你覺得.
這個世界有絕對的公平的時候.
也許你活在很大的痛苦當中.
如果你經常問的問題就是說.
為什麼我做那麼多.
他做那麼少.
為什麼派給我的那一份是小那麼多.
派給他的那一份是大那麼多.
為什麼我受的對待.
是比起另外那個人所接受的待遇.
是差那麼多的.
如果你經常問這個問題的時候.
我告訴大家.
你不能夠活在喜樂當中.
因為這個世界.
很現實的一個事實.
這個世界不會有絕對的公平.
包括教會.
中國經常都說.
人比人比死人.
當我們經常比較自己.
你覺得不幸的事情.
來去比較另外一個人.
為什麼他那麼舒服.
那麼安舒的時候.
我告訴大家.
你一定不會在你的事工裡面.
找到任何的喜樂.
比較.

$^{721}$絕對會奪走.
你事工的喜樂.
那麼你說.
耶穌基督對於公平.
有些什麼看法.
說得出來.
二十二節.
我若要他等到我來的時候.
與你何干.
這人將來如何.
與你何干.
這幾個字我告訴大家.
是我們基督徒要持守一個.
如果你想我們的精神健康的話.
你要記住這幾個字.
如果你不記得.
或者沒有實踐這幾個字.
不明白這幾個字的時候.
很多時候我們會精神不健康.
要不然就完全不做.
視乎.
簡單.
與人的比較.
就是我們的喜樂.
被奪走的時候.
在神的角度裡面.
我告訴大家.
神的角度裡面.
不是說公平.
神的角度裡面.
是說愛.
耶穌基督對彼得的呼召.
與對約翰的呼召.
是完全不同的呼召.
耶穌基督對某一個人的祝福.
與對另外一個人的祝福.
也是完全不同的祝福.
今天耶穌基督.
要我們走的十字道路.
並不是一個標準化的十字道路.

$^{761}$人人走得上這條十字道路.
所面對的都是一模一樣的.
沒有這回事.
十字道路.
從來都不是一個.
標準化了的事業.
我們所面對.
我們所領受.
是一個完全不同的十字路.
是耶穌基督.
我們的主親自.
為我們預備的十字道路.
所以我們很多時候.
都花了很多的精力.
來經常地問.
神啊.
為什麼某某基督徒.
他有這麼好的婚姻.
他有這麼成功的事業.
他的呼召這麼清晰.
他的兒女這麼孝順.
他的父母這麼健康.
我就不同.
我一直在那裡掙扎.
搞錯了.
我們要問的問題.
從來都不是為什麼.
我們要問的問題很簡單.
這裡告訴我們.
我們愛我們的主嗎.
愛的話.
跟從他.
到最後的時候.
約翰用一個很特別.
我覺得很有趣的方式.
來完結了這個書卷.
在這個二十三節他說.
者話.
就是說這個信徒不死.
者話存在弟兄中間.

$^{801}$說那門徒不死.
其實耶穌不是說他不死.
乃是說我若要等.
我若要他等到我來的時候.
與你何干.
留意.
者話存在弟兄中間.
就是說什麼.
就是你一句我一句.
傳來傳去.
最後傳到.
就是我們玩的遊戲.
以訛傳訛.
最後傳到一個.
講法就是說.
這個信徒.
不用死.
耶穌說不用死.
約翰自己知道.
不是不用死.
只不過說.
不關他的事而已.
在早期教會.
現在教會.
每一代的教會.
都不能夠逃避.
傳來傳去.
以訛傳訛這個問題.
不過縱然當時的話.
將耶穌基督沒有說過的話.
好像弄錯了他的意思.
耶穌從來沒有說過.
約翰不用死.
不過說傳來傳去.
但是約翰在這裡講得很清楚.
整個福音.
或者他寫過這本約翰福音的最終極.
不是教會裡面傳的謊言.
也都不是講到彼得的信徒.
也都不是講到約翰的長壽.

$^{841}$而是講到什麼.
有一個很簡單.
很直接的命令.
你愛我嗎.
愛的話.
跟從我吧.
別人如何.
與你無干.
與我無干.
我們走在耶穌基督.
為我們每一個人所預備的十字架道路的時候.
我們就將我們的生命.
萎生在這個十架道路上面.
一直跟隨主.
直到見主面.
這個就是耶穌基督.
簡單直接的命令.
願意我們都能夠明白.
願意我們有勇氣承擔.
請各位同低頭祈禱.
耶穌基督我們多謝你.
因為通過這個故事.
你告訴我們.
在這個世界裡面.
你已經為我們預備了十架道路.
這一條路.
與其他人的十架道路.
是不盡一樣.
但實際上我們知道.
你是為我們預備最好的.
所以給我們有單純的信心.
踏上這條十字架道路的時候.
我們再一次確認.
主啊.
我們愛你.
我們願意來到這十架道路上面.
跟隨你.
直到生命的完結.
直到見你的面.
多謝你給我們一切.

$^{881}$祈禱鳳爺.
蘇冥至以求.
阿們.
\newpage



\section{啟示錄 2:8-12}
\label{sec:031988}
\textbf{士每拿:貧窮卻又富足的教會}~~~~趙崇健牧師
\newline
\newline
連結: \href{http://www.church.com.hk/acms/content.asp?site=cdc\&op=show\&type=product\&code=031988\&layout=sermon}{\texttt{031988}} ~~~~ 語音日期: 2026-04-19
\newline
\newline
\hyperref[sec:031987]{\small{< < < PREV SERMON < < <}}
~
\hyperref[sec:index_chronic]{\small{[返順時目]}}
~
\hyperref[sec:031989]{\small{> > > NEXT SERMON > > >}}
\newline
\newline
啟示錄 2:8-12
\newline
\begin{longtable}{cl}
\hline
\hline
章節 & 經文 (和合本修訂版)\\
\hline
2:8 & \begin{tabularx}{0.7\textwidth}{X} 「你要寫信給士每拿教會的使者，說：『那首先的、末後的，死過又活了的這樣說： \end{tabularx} \\ \\ \relax
2:9 & \begin{tabularx}{0.7\textwidth}{X} 我知道你的患難和貧窮—其實你卻是富足的，也知道那自稱是猶太人的所說毀謗的話，其實他們不是猶太人，而是撒但會堂的人。 \end{tabularx} \\ \\ \relax
2:10 & \begin{tabularx}{0.7\textwidth}{X} 你將要受的苦，你不用怕。看哪！魔鬼要把你們中間幾個人下在監裡，使你們受考驗，你們要遭受苦難十日。你務要至死忠心，我就賜給你那生命的冠冕。 \end{tabularx} \\ \\ \relax
2:11 & \begin{tabularx}{0.7\textwidth}{X} 凡有耳朵的都應當聽聖靈向眾教會所說的話。得勝的必不受第二次死的害。』」 \end{tabularx} \\ \\ \relax
2:12 & \begin{tabularx}{0.7\textwidth}{X} 「你要寫信給別迦摩教會的使者，說：『那有兩刃利劍的這樣說： \end{tabularx} \\ \\
[1ex]
\hline
\hline
\end{longtable}
$^{1}$在開始之前 我都首先想和各位弟兄姊妹打個招呼.
但我都想大家都和我打個招呼 好嗎.
天氣冷冷的 我們都或者精神少少好些.
我和大家說聲主內平安 大家都大聲點回應 好嗎.
各位弟兄姊妹主內平安.
啊 好好啦.
今日都涼涼的 但都有溫暖的.
希望神給我們有好好的一個主日.
亦都盼望 神給我們有一個學習神話與經文的機會 好嗎.
家中的弟兄姊妹都想說聲主內平安.
你都可以大聲回應我 好嗎.
各位弟兄姊妹 我們今日繼續去看啟示錄.
這一點我們需要留意.
其中一點很重要 就是啟示錄本身其實是一封書信.
而在這封書信裡面.
是會由當時第一世紀不同的教會去傳閱.
他們只有一份手本.
會由不同的人帶到不同的教會讀出來.
給弟兄姊妹去聽.
而啟示錄這個書信裡面.
裡面有很奇妙地包括了七封不同教會的小書信.
所以我們不難理解.
啟示錄最少會在這七個教會裡面傳閱.
並且當啟示錄傳閱的時候.
這七個教會每一間都會聽到那一封寫給自己的書信.
所以這七封信並不是一些抽象的道理.
而是真正確實直接回應著七間教會當時的處境.
今日我們看的是第二封書信.
是寫給士美娜教會的.
這個教會當時面對的情況.
我們不會陌生的.
只要我說你就明白.
因為他們面對的就是貧窮和富足這個很普遍的問題.
如果你今天打開經濟的報紙.
去看經濟的新聞.
有時候你會見到這個詞彙.
叫做堅尼系數.
什麼來的呢.
就是當這個數字越大的時候.
就代表社會的財富就越不平均.

$^{41}$80年代的時候.
澳洲這個堅尼系數的數值大概是0.27.
以一分為滿分.
去到2019年 即是大概過了40年.
這個數值現在就升到0.325.
即是都升了大概有兩三成的地步.
一個社會的財富分佈差距很大的時候.
其實好像現在這樣.
又有一個新的詞彙.
叫做K型社會.
之所以叫K型不是因為社會K型.
是因為社會財富的分佈是真的K字的形狀.
意思是富有的人越來越富.
貧窮的人越來越窮.
從前最大部分的是中產.
中產就慢慢消失.
有一次我和我的兒子看卓京新聞.
基本上是我看的時候他坐在我旁邊.
新聞就在說現在的樓價真的很貴.
看完之後他就看著我.
跟我說Daddy I am so quote.
如果你不明白什麼是I am so quote.
就問問我們當中的年青.
並不是解他煮熟的意思.
他說我以後都不會有自己的地方住.
跟著我就跟他說笑.
我說兒子不用擔心.
家裡有地方住.
你準時交租就可以.
但是當我想了一下.
我就發覺到在家裡.
都會有一個K型的社會.
我就想問一個問題.
如果你是窮那邊.
你是不是就是失敗.
如果你是富那邊.
你又是不是真的是成功.
我相信我們再助大家.
其實我們都是想向上的.
即使不是剛剛上月球那支火箭.

$^{81}$直飛向上.
最少都慢慢向上.
甚至乎最起碼不是向下.
但是我又想問.
我們為了向上.
我們願意付上什麼代價.
去到什麼地步.
我們願意妥協到什麼地步.
啟示錄第二章.
是耶穌親自寫給.
是梅拿教會的書信.
祂就說.
你是貧窮的卻是富足.
這句話有問題.
哪有可能同時又貧窮.
同時又富足.
這句話其實顛覆了.
我們現在的價值.
史梅拿教會.
就是需要這種顛覆的價值.
去面對他們的境況.
現在的社會是畸形社會.
其實我們都需要.
這套顛覆的價值.
去面對我們這個世界.
讓我們低頭.
求神幫助我們.
親愛的恩主.
讓我們看到.
我們在這個社會當中.
我們是何等的有福份.
不是單單因為我們擁有多少.
而是因為我們有主理.
我們擁有主耶穌基督.
無論我們的境況如何.
我們深深地相信.
要是我們信得不足.
求主教我們相信.
主你必供應.
你必定看顧.

$^{121}$有時候我們會走過高低.
但主你必定陪同我們.
教導我們今天藉著主你.
親自寫信給史梅拿的書信.
去教導我們.
我們如何全心的跟隨你.
致死忠心.
因為神你說我們是你的僕人.
你不是要求我們作僕人的能力.
要去到多高.
因為我們的能力是神的所持.
神你但願,盼望.
我們這班作僕人的.
就是作一個有良善.
有忠心的僕人.
求主你幫助我們.
祈禱奉耶穌基督你的名求.
阿們.
閱讀啟示錄.
我們要明白.
其實啟示錄.
是在第一世紀.
真真實實教會的處境.
而史梅拿教會.
就是活在史梅拿這個.
很強大的城市裡面.
這個城市是怎樣.
簡單地說.
有錢,繁榮.
還有文化.
對著羅馬的政權.
是極度的忠誠.
所以他們有一套很清楚的價值觀.
就是在這個城市裡面住的人.
你要融入這個城市.
我們這個城市講忠誠.
你要忠誠.
要和其他人一樣.
否則其實不是很歡迎.
甚至乎有後果.

$^{161}$詳情怎樣呢.
當我將歷史和經文放給大家看.
你就會感覺到當時的人.
其實有多大壓力.
我們試試這樣看一看.
我們在經文裡面這樣說.
我說耶穌說的原因是.
這封信是耶穌說的.
耶穌說是先脫手的.
是抹後的.
這句話對我們基督徒來說.
就是在說神.
阿拉法,奧梅加.
但當時的人.
其實他們會想到第二件事.
就是當時羅馬帝國進入了和平的階段.
他們一直在說.
羅馬已經為這個世界帶來新的時代.
過去二百年的內戰已經結束了.
現在是和平的開始.
這一刻.
現在和平這一刻.
才是真正這個世界開始的時代.
甚至乎當時我們要理解.
羅馬已經在說.
羅馬是永恆的.
我們現在的和平.
只要我們所有人一起維持的話.
我們現在的生活.
是一個永恆的國度.
而甚至乎當時的羅馬皇帝.
對於現在這個時代.
定了一天是開始的日子.
就是羅馬皇帝自己出生的那一天.
他的生日就是這個時代.
永恆時代開始的一天.
所以這裡當耶穌說.
首先的 末後的.
對當時的人來說是很衝擊的.
因為不是羅馬.

$^{201}$耶穌基督說 是我.
歷史的書寫.
不是由當時的權貴去決定.
面對當時社會的境況.
耶穌又說.
死過又活的.
一聽下去就是.
受苦節 復活節.
但這句話對士梅拿的信徒.
有另外一種感受.
當然一方面說耶穌基督死裡復活.
另一方面說士梅拿這個城市.
一個死過又活的城市.
這個城市曾經被摧毀.
在公元前600年.
荒廢了一段很長的時間.
200年 之後又再被重建.
用了五六十年的時間.
慢慢恢復過來.
所以這個城市裡面的人.
他們是很明白.
士梅拿就是一個復活了的城市.
但是耶穌在說什麼.
你們看到這個城市死裡復活.
但是真正死裡復活.
不是這個城市.
真正死裡復活的是主耶穌基督自己.
接著耶穌又說.
我知道你的患難 你的貧窮.
問題就來了 畸形社會.
在一個如此富裕的城市裡面.
教會竟然是貧窮.
不是說靈女貧窮.
是說真實貧窮.
可能他們本來是不窮的.
但是不知道什麼原因.
他們慢慢變窮.
又可能他們一直都很窮.
窮了很久.
沒辦法翻身 沒辦法向上.

$^{241}$但重點不是他們有多窮.
是他們為什麼這麼窮.
經文就說.
有人毀謗他們.
我們不知道具體的內容.
到底是毀謗什麼.
但是在這個部分.
我們可以肯定一件事情.
他們不單止窮.
還要得罪人.
還要得罪有錢人.
耶穌就說.
你們中間有幾個人.
會被下到監獄裡面.
坐牢現在對我們來說是懲罰.
但當時坐牢其實不是懲罰.
在羅馬的制度裡面.
不會判你坐牢的.
只會判你體罰.
或者其他的死刑.
坐牢在一個地方.
嫌疑的人等候審判.
不讓你走 即是拘留.
換句話說.
當時他們這個教會.
是有人在告他們的.
他們在等待進入審判的當中.
當時羅馬做了一件事.
很有趣的事.
就是在很大的羅馬帝國裡面.
要選一個城市.
一個而已.
去興建一個廟宇.
廟宇就只是供奉羅馬的皇帝.
很多城市都想爭這個榮譽.
因為有個地標在這裡很重要.
到最後選了士梅拿.
原因很簡單.
他們願意表達.
從今至永遠.

$^{281}$這個城市都會對羅馬.
絕對的忠誠.
不會出賣羅馬.
從此以後.
當這個殿.
供奉皇帝的殿.
立了在士梅拿之後.
整個城市的運作.
就是敬拜皇帝.
表示忠誠 表忠.
城裡面充滿了有錢的人.
有地位的人.
趕著去這個城市裡面表忠.
你可以想像.
在這個環境裡面.
教會被毀謗.
都未必需要他們做錯了什麼.
或者只需要他們沒有做什麼.
已經是一個問題.
就是他們沒有跟著社會價值走.
他們沒有推崇.
大家都推崇的東西.
我們要忠誠.
你是不是搞事.
教會當時.
沒有活出.
社會所謂的.
正常.
因為他們沒有這樣做.
他們貧窮.
他們被毀謗.
被針對.
被當成是異類.
士梅拿這個地方.
要求的是.
妥協.
同化.
跟著走.
但我想問今天的世界.
不是都要我們.

$^{321}$妥協.
跟著走嗎.
是不是.
有一本書你一定看過.
沒有看過都聽過.
《安徒生童話》.
認識吧.
它有一個故事叫《丑小鴨》.
故事很簡單.
我幾句講到.
就是在一個農場裡面.
有一個鴨媽媽.
孵了一堆我相信都很可愛的小鴨子.
其中最後出生的小鴨子.
不知道為什麼長得特別大.
又灰色又醜怪.
跟其他的小鴨子不同.
所以這隻小鴨子自出生開始.
就被排擠.
連他自己的親生兄弟姐妹.
甚至最後他媽媽都嫌棄他.
這隻丑小鴨實在受不住.
受不住鴨的嘲笑.
他離開了農場.
他就去流浪.
日子很艱苦.
他就被貓和雞笑他.
你一無是處你做隻鴨.
到了冬天.
又餓又孤獨.
故事講句差點死了.
直至到春天來臨.
過了不知多久.
有一天丑小鴨很傷心.
他就去到湖邊.
看著自己的倒影.
他發現自己不再是那隻.
又醜又灰的怪鴨怪鳥.
原來他是一隻白天鵝.
很熟悉這個故事.

$^{361}$小時候都會問這個故事.
你有什麼反省.
要做讀書報告.
但是我們未必知道.
你知不知道這個丑小鴨的故事.
不是故事來的.
大家知道嗎.
安徒生晚年有個朋友.
幫安徒生寫了一本自傳.
這個朋友就問安徒生.
為什麼你不寫一本自傳.
安徒生說我已經寫過了.
我那本自傳就是丑小鴨.
所以後世一直都流傳.
其實丑小鴨的故事.
本身就是安徒生一生經歷的寫照.
這個丑小鴨的故事.
對於那些被排擠.
被針對的人.
又或者像我自己這樣.
對特殊需要有特別理解的人.
聽下去是會特別有感受.
因為這個世界.
一直都要求我們妥協.
同化 跟大隊.
不跟就會被針對 被排擠.
仕梅拿這個世界.
安徒生的世界.
今日這個世界.
其實都不是很大分別.
當時仕梅拿的人.
我想深一層.
難道全部都真心.
恭奉敬大皇帝嗎.
我相信未必.
當中有不少的人.
其實是表忠.
賣堆.
拆鞋.
拿著數.

$^{401}$所以很多人就選擇了.
妥協 同化 跟大隊.
是為了利益.
其實今日都是一樣.
如果你跟別人說.
今天逢星期日早上.
我會早點起床.
別人問你做什麼.
我回教會 敬拜神.
有些人心裡想.
可能沒說出口.
真是浪費時間.
不如睡晚一點更好.
星期一至五 我還不夠忙嗎.
如果你有用十分鐘靈修.
每一天的話.
你跟別人說.
可能別人會覺得.
我不會用十分鐘學實用的東西.
如果我學外語.
我想我已經成功了.
最明顯的是奉獻.
如果你有十一奉獻.
你告訴別人.
或者跟別人分享.
我有十一奉獻.
有些人 社會的人可能覺得.
是牧師騙你的錢嗎.
第二就是.
我不如慢慢儲起來.
放在超級市場.
不就更實際.
基督徒信仰有句話.
在世.
而不屬世.
人活在世界上.
我們沒有選擇.
哪個世界哪有選擇.
我們活在一個怎樣的世界.
我們都沒有選擇.

$^{441}$但我們可以選擇的是.
怎樣去活.
我們可以選擇.
我們是否屬於這個世界.
這兩個層面我們是有選擇的.
是每個教會面對的世界.
就是要他們妥協.
和忍耐之間做選擇.
到底舍梅拿應該妥協好.
還是忍耐好.
我們今天其實都是面對一樣的選擇.
剛才我說了.
舍梅拿是一個怎樣的世界.
但是耶穌寫信給教會.
是要他們知道一件事情.
這個世界.
按著你.
擁有多少.
能力怎樣.
你的地位去定義你.
但是神.
是按著我們屬於.
哪一個來定義.
我們.
經文是這樣說.
耶穌說我知道你的.
患難和貧窮.
教會自己.
都覺得自己貧窮.
世界都是這樣看的.
可能還會嘲笑.
但是耶穌說你卻是.
富足的.
原來教會.
看不到自己是富足的.
教會亦都不知道.
自己是富足的.
耶穌要.
提醒他們.
看回.

$^{481}$他們之所以富足不是因為他們擁有.
多而是因為他們沒有.
出賣到在這一刻.
他們都還沒有出賣.
他們作基督徒的身份.
耶穌說過一句說話.
值得我們記得.
人若賺得全世界.
卻賠了.
自己的生命.
有什麼益處呢.
人還能拿什麼.
換生命呢.
接著經文就說.
你將要受的苦你不用怕.
這句說話.
很難聽.
很難接受.
既然將要受苦.
是不是.
為什麼神你不拿走那個苦.
神正在叫他們.
你不用怕.
你將要受苦你不用怕.
如果你找我.
跟我說你很苦.
我跟你說不用怕.
將要受苦不用怕.
耶穌不是說他們不用受苦.
而是說.
你會受苦.
但是你真的不用怕.
沒有人喜歡受苦.
神當然有能力.
去挪開痛苦.
但是有一點我們常常忘記.
神有能力挪開痛苦.
但神更大的能力是.
將痛苦變成為.
光強和喜樂.

$^{521}$不需要改變環境.
環境一樣.
但改變了所有經歷環境的人.
和特質.
經文裡面就說.
我叫你們被試探.
同一個環境.
受苦.
可以是試探可以是試煉.
有能力的人會覺得是一樣.
但本質是不同.
讓我解釋一下.
試探.
是出於撒旦.
目的.
是想你跌倒.
試探是出於撒旦.
目的是想你跌倒.
試煉是出於神.
目的是想你.
站立得穩.
外表看起來很相似.
但是就好像你去健身房.
你想練健身.
差的教練就會給你.
一個很重的你.
讓你舉不起.
甚至令你受傷失敗.
差的教練.
好的教練就會給你.
適合的重量去舉.
並且扶著你去舉.
所以事每拿教會來說.
試探和試煉很相似.
他們覺得自己被試探.
但神說我是在試煉你.
魔鬼想他們跌倒.
但是神想將試探.
化為試煉.
叫他們經歷神.

$^{561}$經文又再說.
你要受難十日.
每次看啟示錄.
對大部份的數目字.
都要存在一種小心.
因為大部份的數目字.
其實都是象徵性.
十是代表短暫.
代表苦難是很短暫.
就好像十戒.
很簡短.
不過以色列人都做不到.
十就代表很完全.
又簡短又完全.
又完整.
苦難十日.
是代表會受苦.
但是有限有界線.
不會永遠.
所以神就跟教會說.
要藉著這十日去學習.
致死忠心的功課.
我就賜你生命的冠冕.
苦難是短暫的.
不至於死.
但是要藉著這次的機會.
去演習.
去rehearse.
致死忠心.
德勝的功課.
在羅馬時候.
的運動場上.
贏了的德勝者.
頭上會戴著冠冕.
所以這裡所說的冠冕.
不是皇帝頭那個.
是現在見得到.
用黃金鑄成的一種.
而滋梅拿這個地方.
有一個很著名的形象.

$^{601}$滋梅拿這個城市.
是建在山坡上.
並非一直起上去.
是沿著山腰.
打橫去起.
就好像遠看的時候.
好像冠冕一樣.
也因為這個城市在海邊.
所以坐船做貿易的人.
很遠就會見到滋梅拿這個地方.
就好像彎彎地.
一個冠冕一樣.
所以當耶穌在說.
生命的冠冕.
聽的人會想起我們這個城市.
多厲害 多好.
耶穌說.
我賜到一個是生命的冠冕.
冠冕不是滋梅拿.
或者這班信徒.
很歡喜滋梅拿.
但是耶穌說真正的冠冕.
不是在這個世界裡面.
你有沒有看到神.
耶穌在做什麼.
耶穌是要教會知道.
兩個世界裡面拉扯著.
一個世界告訴你.
你很貧窮.
耶穌的世界告訴你.
你很富足.
問題是教會要選擇.
屬於哪一個的世界.
今天在澳洲.
我們大部分都是移民.
有部分的.
我們曾經試過宣誓.
入籍.
如果你宣誓屬於一個國家.
你就會成為.

$^{641}$那個國家的.
公民 citizen.
今天很多人移民.
所以就出現了一個問題.
就是一個人可不可以.
同時屬於兩個的國家.
你覺得可不可以.
有些國家是容許的.
例如英國.
有些國家.
沒有寫明的默許.
例如加拿大.
有些國家是完全禁止.
有些雙重國籍.
例如新加坡.
但其實重點我想說.
不是制度.
不制度的問題.
重點回到那個問題.
就是一個人.
可以向多少個國家.
真正忠誠呢.
公民權.
和永居權.
是不同的.
明白嗎 公民權是 citizenship.
永居權是 permanent residency.
永居權.
就是容許你住在這裡.
但你不一定要完全忠誠.
你住在這裡OK.
但公民權就關乎你的身份.
關乎到忠誠.
而且公民權可以.
一代一代地傳下去.
我們當中有不少的人.
或者你會有兩本.
甚至乎兩本以上的護照.
當你拿著.
三幾本不同的護照.

$^{681}$在你面前的時候.
你會想起什麼呢.
passport這個字.
是一本文件.
去旅行的時候.
很重要.
是用來通過港口.
現在就是機場.
passport這個字.
直譯.
原來的意思就是.
passport.
通過一個關口.
港口需要的文件.
就會慢慢變成.
我們今天所說的passport.
代表著一件事.
你承認.
你屬於.
這本護照的國家.
並且.
你進入了一個國家.
都承認.
你的身份.
兩者都屬於大家.
永居權PR.
沒有passport.
因為它只是給你.
在這個地方住.
並不代表你真正屬於這個地方.
我們基督徒.
在這個世界.
我們基督徒在這個世界.
其實我們只有PR.
OK嗎.
我們基督徒在這個世界.
只有居留權.
我們不是這個世界.
的公民.
有一天我們都要.

$^{721}$拿著天國那本.
passport.
去到那個我們真正屬於.
我們屬於神的世界.
這個才是我們真正的身份.
公民權就關乎忠誠.
永居權的忠誠.
可以很有彈性.
就引申到我們今天.
在這個世界居住.
即是居留.
我們怎樣生活.
其實我們有彈性.
在這個世界教會.
今天的基督徒.
就是面對一個問題.
在這種彈性中間.
我們有多少屬於神.
有多少屬於世界.
就好像黑色.
和白色混在一起.
可以是灰色.
50/50.
可以是偏黑.
可以是偏白.
問題是我們想自己更加.
像屬神.
還是我們更加.
像屬於世界.
知死忠心.
就是叫我們永遠不要忘記.
我們最根本.
最首先的身份.
我們是基督徒.
無論是永居.
還是公民.
你知道在澳洲都要.
負上責任.
都要付出.
就是交稅.

$^{761}$但從神的角度.
人就從來沒有先付出.
人只會得領受.
所有的領受.
都是在神.
首先付出之後.
發生的回應.
我們去研讀今天的經文.
我很想.
我們感受到一種安心.
帶著一種.
平安去過著.
屬神.
的生活和生命.
因為今天的經文.
我們看到.
神是很愛護.
侍梅拿教會.
其實神也很愛護.
我們每一個人.
你要留意幾點.
當教會.
面對困難.
神親自.
跟他們說話.
人的說話.
我們聽得多了.
但恐怕神的說話.
聽一次.
就已經充滿能力.
今天神仍然說話.
神會跟你說話.
神會跟我說話.
神之所以要說話.
是因為.
神愛我們.
神說話是因為神愛.
並且神兩次說到.
我知道.
兩次我知道.

$^{801}$這一個跟我們說話的神.
對我們的處境.
他不是不清楚.
而是更清楚.
如果神知道.
你就可以放心.
他不會誤會你.
他不會忽略你.
他不會不明白我們.
神知道.
神容許人.
被試煉.
很多人都不喜歡壓力.
好像.
我知道.
我知道運動對身體好.
但我仍然是一個.
可以坐的說.
不會動的人.
但人就是需要被試煉.
被訓練.
才能夠成長.
神已經為我們.
預備好一切.
如果你在健身房裡舉重.
神在地上鋪滿軟墊.
神選擇最適合的重量.
去問你.
這個可以嗎.
試一下.
我跟你一起試.
有神在你旁邊扶著你.
你覺得好重 好辛苦.
神幫你撐著.
這樣是試煉.
是跟神一起的試煉.
神要人忠心.
耶穌自己.
在十字架上沒有妥協.
所以我們今天都要背負十字架.

$^{841}$我們可不可以.
不妥協呢.
當耶穌背負十字架的時候.
神知道.
耶穌辛苦.
派了一個鼓勵內人西門去幫他.
聖經叫我們.
去背負十字架.
神不知道我們辛苦嗎.
神當然會親自跟我們一起背負十字架.
派人幫我們.
所以背負十字架不要忘記.
前提是.
神幫我們的.
神賜人 冠冕.
頭上.
耶穌的荊棘.
耶穌頭上的冠冕是荊棘的冠冕.
我們今天領受.
是生命的冠冕.
完全不對稱.
完全不合理.
但正正.
彰顯著神對人.
何等的愛.
所以弟兄姊妹.
神做了五件事情.
在四節經文裡面.
神說話.
我們要聽.
神知道的話.
我們就要信靠.
思念實有.
神思念.
我們要靠主忍耐.
神要求忠心.
史姆蘭教會.
就不要妥協.
神賜生命的冠冕.
我們不是純粹付出.

$^{881}$我們是領受.
我們要持守盼望.
經文的重點.
不是要我們.
缺字與否.
不是說我們有沒有信仰.
而是我們活得像不像.
我們真正屬於.
那個世界.
基督徒應當在世.
卻不屬世.
因為神.
我們是真的可以.
活得不一樣.
經文裡面最後說.
聖靈向宗教會所說的話.
還有意的就應當聽.
讓我們低頭祈禱.
親愛的恩主.
我們明白你的心意.
因為我們得救.
是本乎因.
但父神你也將我們.
分辨為聖.
適合主用.
叫我們不要忘記.
分辨為聖這四個字.
因為這也是神.
你將我們安放於.
現在的世界.
你還未將我們.
召回見主面的原因.
願意我們現在這個世界.
無論是K型.
還是什麼型.
但你父神願意我們.
能夠明白.
我們這個世界不過是.
短暫的居住.
有一天.

$^{921}$我們要去見主你的面.
所以我們應當活出.
我們屬於那個.
世界.
求神你加力給我們.
我們有軟弱求主你扶持.
我們有疑惑.
求主你賜我們信心.
我們面對挑戰.
面對困難.
求主你叫我們持守盼望.
父神我們.
今天禱告.
表明我們的心職.
是我們願意在世上作鹽作光.
但我們更清楚.
沒有了主你的幫助.
我們做不到.
所以父神我們求你幫助我們.
給我們心.
給我們有信.
給我們有實踐.
在這個世界上.
我們活出神你賜給侍母.
侍母留教會的說話.
我們祈禱奉耶穌基督的名.
阿們.
\newpage



\section{出埃及記 15:3}
\label{sec:031989}
\textbf{聖經人聖經事:我的神,他名是耶和華}~~~~黃克勤牧師博士
\newline
\newline
連結: \href{http://www.church.com.hk/acms/content.asp?site=cdc\&op=show\&type=product\&code=031989\&layout=sermon}{\texttt{031989}} ~~~~ 語音日期: 2026-04-13
\newline
\newline
\hyperref[sec:031988]{\small{< < < PREV SERMON < < <}}
~
\hyperref[sec:index_chronic]{\small{[返順時目]}}
~
\hyperref[sec:031990]{\small{> > > NEXT SERMON > > >}}
\newline
\newline
出埃及記 15:3
\newline
\begin{longtable}{cl}
\hline
\hline
章節 & 經文 (和合本修訂版)\\
\hline
15:3 & \begin{tabularx}{0.7\textwidth}{X} 耶和華是戰士；耶和華是他的名。 \end{tabularx} \\ \\
[1ex]
\hline
\hline
\end{longtable}
$^{1}$好了,又讀聖經了,弟兄姊妹.
生命說書的網絡和你一起看聖經人,聖經事.
今次的題目我的神,他名是耶和華.
這一句是一陣子的經文裡面出現的.
不過呢,就比這個長一點.
我將它做題目呢,就短了一點.
我的神,他叫什麼名啊?.
我的神,他就是耶和華了,他名是耶和華.
耶和華呢,是"I am" 自有永有者,這樣叫的.
好,《出一及記》第十四章三十節.
當日,耶和華這樣拯救以色列人.
脫離埃及人的手.
以色列人是看見埃及人的死屍都在海邊.
他們就在這邊,就是這個叫做密奪.
過海,去到那邊巴力西芬.
自己走了海底的乾爐窩.
然後埃及士兵追下來,他們回望.
哇,死了,又追過來了,怎麼辦呢?.
這個時候摩西一伸脹,海水合番.
哇,這樣我想他們嘔吐了很多次.
就是哇,埃及官兵來了.
哇,摩西你這麼壞,帶我們來.
然後摩西說向前走.
哇,你不是說向前走,是海來的.
然後一指,海分開了.
然後他們說哇,海水都分開了,地都乾了.
於是他們下去了.
一路走的時候,哇,海水變矮了.
哇,好驚啊.
走上去,走上去,就回望了.
站好之後回望.
哇,埃及官兵又下去追我們了.
哇,沒有了,這次還不趕快走.
哇,我想他們是嘔吐了無數次.
但是當他們看到海水合起來之後.
哇,那些埃及人的屍體衝過來岸邊.
死屍都在海邊了.
他們三十一節,以色列人看見耶和華.
向埃及人所行的大事敬畏耶和華.
又信服他神和他的僕人摩西.

$^{41}$第一點呢,藍色螢光.
留意了,有幾個動詞.
看見,敬畏,信服.
有時候都是這樣的.
只是深信不行,要看.
看,看一下,做點神蹟看一下.
那我才能夠信的.
這樣就看見.
以色列人呢,很有恩典.
因為他看見耶和華所行的大事.
然後才全服了.
哎呀,對不起啊,上帝,我們不信靠你.
哎呀,真是對不起,摩西大哥.
哎呀,我們還要去埋怨你呢.
對不起,我們會跟你的了.
你做大哥了,繼續說你怎麼說,我們怎麼做.
看見耶和華向埃及人所做的大事.
就敬畏耶和華,又信服耶和華.
與他的僕人摩西.
這個就是上帝很想他們有的.
就是要敬畏與信服.
敬畏與信服就是我們今天.
上帝對每一個信徒的要求.
是敬畏與信服.
信,我信,不過不服而已.
他叫我這樣做,我不做.
敬畏與信服.
他叫僕人摩西.
好了,最後,第十五章.
摩西和以色列人向耶和華唱歌.
哇,唱歌.
我要向耶和華歌唱.
因他大大戰勝.
將馬和騎馬的頭在海中.
他們信靠,看見,敬畏,信靠,信服之後.
他們開始讚美.
就好了.
不是說好彩.
是神真厲害,摩西真是我們的領袖.
我們信靠,敬畏.

$^{81}$看見耶和華大大戰勝.
大大這兩個字都是中文.
如果希伯來文,我沒有寫出來.
大大是升起,升起的意思.
因為他升起,升起戰勝.
本來我都不知道他在那裡.
突然出來,升起,再升起.
哇,很厲害,看見.
因為他升起,升起戰勝.
將馬和騎馬的頭在海中.
馬和騎馬是他們很敬畏的.
他們被大海消滅.
他們覺得,哇,真是厲害.
第二點.
第一點就是看見,敬畏,信服.
第二點呢.
他們唱的第二點是什麼呢.
耶和華是我的力量.
是我的思歌,是我的拯救.
這一節,十五章第二節.
我們很多時候都要記得.
一起讀吧.
耶和華是我的力量.
我的思歌.
也成了我的拯救.
耶和華是我的力量.
我的思歌.
也成了我的拯救.
三樣東西.
耶和華是什麼.
我的力量.
我的思歌.
我的拯救.
力量是我另外有力量能夠面對.
是我的讚美.
祂是我拯救的人.
也就是說,我是安全.
耶和華是我的力量.
是我的思歌.
是我讚美的對象.

$^{121}$是我拯救的來源.
是祂拯救我的.
第三點呢.
最重要的就是.
耶和華是祂的.
祂很厲害.
最重要的就是.
我的關係.
這是我的神.
到了第三節.
你看見,全部字都轉了.
我的神.
我的力量.
我的思歌.
我的拯救.
我的神.
我要讚美祂.
是我父親的神.
就是我們歷代的.
我要尊崇祂.
很多「我」.
有幾個「我」呢.
1,2,3,4,5,6,7個.
我的力量.
我的思歌.
我的拯救.
我的神.
我要讚美我父親的神.
我要尊崇祂.
全部都是.
神啊,我要向你做出這些敬拜.
再次認定那七種的關係.
力量,思歌,拯救,神.
讚美的對象.
父親的神.
歷代都相信的了.
是我尊崇的對象.
這個很重要.
第三點.
是我的神.

$^{161}$他的名是耶和華.
耶和華是戰士.
很好打.
是我的神.
他的名字叫耶和華.
第三點跟大家說.
當你遭遇很多困難.
很軟弱,很沒力的時候.
你要記得第三點.
這是我的神.
你的神都不能幫助你嗎.
不是的.
到天一亮.
祂就會幫助你.
這是我的神.
他的名是耶和華.
記得不要忘記這個名字.
很多人的名字.
我都不記得.
你都知道.
最怕是人家忽然之間問我.
認不認得我啊.
唉.
我多數的答案都是.
熟面口.
不過名字叫不出來.
他就暗示說他的名字.
我叫陳大文.
他就會說.
我要認.
舊街坊.
不可以的.
你要是.
我就會說.
你樣子很熟面口.
不過我一時間叫不出來.
不好意思.
你再說一次我聽.
我叫陳大文什麼什麼名字.
那時候你來我這裡.

$^{201}$說到我決志了.
真的有人這樣跟我說.
我怎麼認得.
講報道會決志那個.
有一個在一個會裡面.
看到華福會.
四年一次以前.
最近都少開了.
就開會了.
見到他.
又見到你.
四年前見到你了.
現在在做什麼.
我讀神學.
這麼好讀神學.
然後他說.
你呼召我的嘛.
忘記了名字.
很多時候.
久沒見過名字.
就生疏了.
記不起了.
不過有一個名字.
你必須要記得.
他的名是耶和華.
耶和華這個名字.
不要忘記了.
所以求主依靠.
憐憫我.
即是將來我.
去到Dementia.
老人失智的時候.
失憶的時候.
仍然記得.
耶和華是他的名就好了.
最重要記得.
耶和華是他的名.
當你記得.
耶和華是他的名.
你就會找出一部.

$^{241}$《十四十五章》.
就會看到耶和華是戰士.
是我的力量.
是我的思覺.
是我的拯救.
是我的神.
我要讚美他.
是我父親的神.
我要尊崇他.
耶和華是戰士.
他的名是耶和華.
I am.
be.
至有永有.
is.
他的名是is.
簡單.
現在式.
任何時候都是上帝.
任何時候都是.
現在式.
沒有改變.
式在今在.
以後永在.
永不改變.
全能的上帝.
就是我們的神.
我雖然可能會有一天.
死了.
或者我有一天.
忘記了所有東西.
就老人失智了.
不過.
耶和華.
仍然是式在今在.
以後永在.
他的名是耶和華.
主幫助我們.
記得這個名字.
我們同心禱告.

$^{281}$讚美主.
你的名.
叫做耶和華.
是你的名字.
摩西說那些人如果問.
叫我帶你們走.
那個神叫什麼名.
你的答案就是.
我是耶和華.
主幫助我們.
記住耶和華這個名.
直到我們有一天.
完了我們世上的路程和呼吸.
我們仍然要記得耶和華.
因為那個時間.
就是我們去見你的時候.
讚美你.
你是我們的詩歌.
是我們的拯救.
是我們的神.
是我們的戰士.
是我們尊崇的.
是我的神.
也是父親的神.
歷代的信靠.
願你得著榮耀.
讚美耶和華.
禱告奉耶穌的名.
阿們.
\newpage



\section{歷代志上 29:10-19}
\label{sec:031990}
\textbf{一切屬你樂意獻上}~~~~莊明偉牧師
\newline
\newline
連結: \href{http://www.church.com.hk/acms/content.asp?site=cdc\&op=show\&type=product\&code=031990\&layout=sermon}{\texttt{031990}} ~~~~ 語音日期: 2025-09-21
\newline
\newline
\hyperref[sec:031989]{\small{< < < PREV SERMON < < <}}
~
\hyperref[sec:index_chronic]{\small{[返順時目]}}
~
\hyperref[sec:031991]{\small{> > > NEXT SERMON > > >}}
\newline
\newline
歷代志上 29:10-19
\newline
\begin{longtable}{cl}
\hline
\hline
章節 & 經文 (和合本修訂版)\\
\hline
29:10 & \begin{tabularx}{0.7\textwidth}{X} 大衛在全會眾眼前稱頌耶和華；大衛說：「耶和華－以色列的神，我們的父，你是應當稱頌的，直到永永遠遠！ \end{tabularx} \\ \\ \relax
29:11 & \begin{tabularx}{0.7\textwidth}{X} 耶和華啊，尊大、能力、榮耀、勝利、威嚴都是你的；天上地下的一切都是你的；耶和華啊，國度是你的，並且你為至高，為萬有之首。 \end{tabularx} \\ \\ \relax
29:12 & \begin{tabularx}{0.7\textwidth}{X} 豐富尊榮都從你而來，你也治理萬物。在你手裡有大能大力，你的手使人尊大強盛。 \end{tabularx} \\ \\ \relax
29:13 & \begin{tabularx}{0.7\textwidth}{X} 我們的神啊，現在我們稱謝你，讚美你榮耀之名！ \end{tabularx} \\ \\ \relax
29:14 & \begin{tabularx}{0.7\textwidth}{X} 「我算甚麼，我的百姓算甚麼，竟然能夠如此樂意奉獻？因為萬物都從你而來，我們把從你的手得來的獻給你。 \end{tabularx} \\ \\ \relax
29:15 & \begin{tabularx}{0.7\textwidth}{X} 我們在你面前是客旅，是寄居的，與我們的列祖一樣。我們在世的日子如影子，沒有盼望。 \end{tabularx} \\ \\ \relax
29:16 & \begin{tabularx}{0.7\textwidth}{X} 耶和華－我們的神啊，我們預備這許多材料，要為你的聖名建造殿宇，都是從你的手而來，都是屬你的。 \end{tabularx} \\ \\ \relax
29:17 & \begin{tabularx}{0.7\textwidth}{X} 我的神啊，我知道你察驗人心，喜悅正直；我以正直的心樂意獻上這一切。現在我歡喜見你的百姓在此樂意奉獻給你。 \end{tabularx} \\ \\ \relax
29:18 & \begin{tabularx}{0.7\textwidth}{X} 耶和華－我們列祖亞伯拉罕、以撒、以色列的神啊，求你使你的百姓心中常存這樣的心思意念，堅定他們的心歸向你， \end{tabularx} \\ \\ \relax
29:19 & \begin{tabularx}{0.7\textwidth}{X} 又求你賜我兒子所羅門全心遵守你的命令、法度、律例，成就這一切的事，用我所預備的建造殿宇。」 \end{tabularx} \\ \\
[1ex]
\hline
\hline
\end{longtable}
$^{1}$此錄音或錄影的版權是屬於基督教宣導會北角堂所有.
我們的屬靈傳統之一是以真理教導為本.
歡迎你將此錄音或錄影的連結傳給親友.
懇請注意未經授權請勿轉錄或剪輯本錄音或錄影.
心願,如《仕途行傳》二十章三十二節下所說,這道能建立你們,叫你們和一切聖聖的人同得基業.
留意神的作為,去配合祂的工作.
留意不是我們怎樣去做,選擇怎樣去做.
應該更加留心神帶領他的工作.
不知道這個差泉月對你又有什麼提醒呢?.
我們今天又會透過另一段經文去看看.
當我們向神獻相,獻情,奉獻的時候.
我們有什麼事情是需要去留意的呢?.
我們一起低頭有個祈禱.
主人,你是至高,尊榮,偉大的神.
求主,你去閱讀我們向你的敬拜.
你也去闖入我們的心,讓我們可以聽到主你自己的信息.
你對我們要說的話.
求主,你與我們同在,聖靈在你當中自由的運行.
我們這樣祈禱,交託奉主耶穌基督的名求,阿們.
今天這本書是來自於歷代志.
其實歷代志是一本很明顯在述書中神的作為的書卷.
如果大家記得歷代志一開始的時候有很多的家譜.
其實我們只要看看家譜.
在創世的亞當,神去呼召亞伯拉罕.
說到波亞斯,說到大衛,耶穌的家譜.
其實透過這個家譜,神在世上的作為是很明顯的.
耶穌的祖先.
一切都是神在歷史當中一步一步地預備.
去實現了歷史當中的救恩.
領子妹,其實我們今時今日活在2025年.
有時候我們沒有留意到神的作為.
但是只要你花開聖經,花開歷史.
你就真的看到神的救恩的預備,神的作為.
其實是很明顯地實踐了.
而這段經文之前.
其實神已經設立了當時聖殿的制度.
吩咐了以色列利未知派.
他們之後怎樣在聖殿當中去服侍.
很快聖殿就要被建造了.
我們知道是所羅門王去建造聖殿.

$^{41}$但是其實在經文我們看到.
原來大衛王已經預備好這一切了.
在歷代至上二十九章二節.
他說我為我神的殿已經努力了.
已經盡力了.
預備了金紙,銀紙等等.
去做這些器皿.
而第三節也說到.
我自己積蓄的金銀獻給我神的殿.
神是大衛的神.
為什麼大衛可以這麼盡力.
為神獻相呢?.
我相信是因為他真是真切的.
認識這位上帝是誰.
正如我們今天的經文.
第一節開始.
第十節其實是一個頌讚的禱告.
由第十節到第十三節.
都是一個頌讚.
第十節說到耶和華以色列的神.
我們的父你是應當稱眾的.
直到永永遠遠.
大衛一開始就在表達.
自己跟這位神的關係.
以色列的神,我們的父.
大衛看到神跟祂的關係.
跟祂民族之間的關係.
我們翻開聖經我們真的看到.
神其實在以色列民當中.
做了很多很多奇妙的事.
然後大衛也繼續的讚頌.
他說耶和華專大能力榮耀.
勝利威嚴都是你的.
不單止神祂本身的尊榮.
甚至十二節下說到.
在你手裡有大能大力.
你的手使人尊大強盛.
是被神所用的人都一樣很有能力.
領子們我們所認識的神.
其實跟大衛接觸他所認識的神.

$^{81}$是同一位神.
但是我們心底有沒有這份認同呢.
神的確是這位的專大能力.
威榮的上帝.
大衛他能夠這樣頌讚的時候.
他在生命當中.
他自己有深切體會過.
經歷過神各樣的大能.
大家記不記得在大衛落難的時候.
他被追殺.
神多次將他從騷擾,敵人等等.
當中拯救出來.
他深切經歷神的保守.
以至於他能夠頌讚的時候.
是有很清楚的一個根基.
記得有一次我去一個月.
在海外服侍.
當時我還沒有做教務同工.
當時要代表某個機構.
要跟當地的官員簽約.
其實真的膽子粗粗的去做.
當時都很年輕.
機構主管信任.
就要代表他跟別人簽約.
其實現在回想起來.
也不是很適合我當時去做.
整個行程.
約官員去談,溝通.
最後能夠達成合約.
竟然很出奇,很順利.
以前從來沒有做過這件事.
是神在當中很奇妙.
一路一路帶領.
而且在整個過程中.
接觸到不同的人.
還有機會分享福音,見證.
姐妹們,其實我們所認識的神.
是不是你親身經歷過的?.
以致你能夠由衷地去頌讚祂.
也以致你能夠.

$^{121}$在你生命有好的事情的時候.
你很想給神為祂所用.
我相信大衛在頌讚的時候.
是想起祂在神裡面.
很多很多一切的經歷.
我們又有沒有常去回想.
神在我們裡面的作為.
或者神是去帶領我們.
經過一些困壓的.
神也成就了一些.
我們以為好像很難去成就的事.
我們頌讚的時候.
會不會也想到神在我們的群體裡面的作為.
或者甚至是神在不宣家裡面的作為.
又或者你是認識宣教士的.
神是透過宣教士.
在世界各地很多不同的地方.
做了很多奇妙的工作等等.
以致我們去頌讚神的時候.
其實是有這個根基.
是願意獻情給祂.
而且我們的確可以在生活當中.
不同的範疇都可以去經歷神.
因為經文告訴我們.
天上地下的一切都是你的.
也就是說國度是你的.
神啊 你是萬有之手.
你也是治理萬物.
經文重重複複告訴我們.
神是一切都是屬於神的.
沒有任何的遺漏.
所以我們必定可以在不同的範疇.
都可以去經歷到這位信神的神.
我們就是這樣常常要去頌讚祂.
因為祂配得一切的頌讚.
這位神才是這位至高偉大的神.
我想祂是先去邀請我們.
是想去獻賞我們可以對祂最深切的頌讚.
我們很想的去頌讚祂.
然後我們才在想.

$^{161}$我們如何去獻情我們的生命.
獻情我們任何的資源.
去給祂 為祂所用.
好了 當我們看到一位這麼偉大的神.
至高的神 全能的時候.
我們就很自然又會望望自己.
我們當望望自己的時候.
會看到自己那一份的渺小.
你看看經文大衛的禱告那裡.
說我算什麼 我的百姓算什麼.
真的 當你相比起尊大榮耀的神的時候.
我們真的不算什麼.
當大衛說到我算什麼的時候.
其實他明白到他一開始.
其實真的是一個完全不起眼的小子.
大家記不記得.
他只不過是一個很年輕的牧羊人.
甚至如果你記得的話.
先知薩姆爾他來到以色列民當中的時候.
他想尋這個新君王的時候.
他父親耶西是怎樣的.
他甚至沒有把大衛算在眼內.
他好像不記得自己有這個兒子那樣.
你記得那段經文.
大衛是這樣一個渺小的人.
他真的是Nobody.
其實我們也可能是這樣.
只不過是神 耶和華神對他的揀選.
現在大衛他已經成為一個之君.
但是他知道這一切一切.
只不過是神所賜給他.
不是自己的能力所能夠奪取回來的.
大衛知道自己的本相.
不能夠在神面前變得謙卑.
丁姊妹同樣我們知不知道.
經文說萬物是從神而來的.
就是我們所得到的.
其實都是從神而來的.
其實很吊詭的.
好像我們擁有.

$^{201}$但是其實又不是真的我們所擁有.
是神所賜 神託管給我們的.
當然神要我們做好中心的管家.
好好運用金錢.
運用在有意義的事情上.
但是既然說的是代管的話.
其實就是可以隨時後命.
應該可以用在神要用的地方.
這裡當然有說到是財富.
但是不單止是在說財富.
其實可以說是包括你的財降.
你的知識 你的能力.
我們有時候以為那些是我的興趣而已.
那些是我自己所學得回來的技能.
一些的經驗.
我的語言能力.
專業知識等等.
但是其實.
我們以為跟神的國度沒有關係.
但是不是的.
在神的國度裡面.
祂不是這樣去分的.
不是這樣去分到這個屬靈的.
這個屬世.
因為根本一切一切都是從神而來的.
我們所得到的一切是從神而來.
所以你已經是有一個開放的態度.
可以任憑主去用.
好像一位宣教士.
其實他原本是讀會計的.
他做會計都多年了.
但是原來他很不喜歡會計的.
他做了很多年.
但是仍然都不喜歡會計.
但是到了宣教工場的時候.
因為有一個需要.
他就願意去擺上.
用這些會計的技能.
就能夠成為當地宣教團隊的祝福.
一個很大的祝福.

$^{241}$而另一個相關的事實是.
相比於永恆的神.
永恆的國度.
我們其實在地上的日子是很短暫的.
其實是很短暫的.
經文說我們在你面前是客旅寄居的.
與我們的列祖一樣.
在世的日子如影子.
好像沒有盼望.
說的是人在地上是客旅.
我們現在身處在香港.
你知不知道香港已經是.
全球人均壽命最高的地區.
不是日本.
知不知道原來2025年.
在香港的女性.
平均壽命是多少歲.
大家猜一下.
原來2025年.
女性的平均壽命是88歲.
很厲害的.
其實我在北宣街也見到不少有些九十幾歲.
一百歲的長者都很精神.
他們都很厲害的.
我們這十幾二十年.
是越來越講得多.
有養生,營養,運動等等.
都是好的.
盡量讓地上的日子延長.
又過得好一點.
但是其實相比起永恆的神.
永恆的天家.
其實這些年月日的長度.
相比之下.
我會覺得好像接近於零.
即是相比於永恆.
其實這些的年日.
多少歲都好.
都是近乎零.
少到好像微不足道那樣.

$^{281}$要留意.
我不是說要否定地上身體好的重要.
因為沒有好的身體.
你想做有意義的事.
要去服侍人都做不到.
但是相比之下.
真的.
人生相對於永恆.
其實只是很短很短.
我們要知道人的限制.
不要以為我們有時很有能力.
有時好像很無限那樣.
在地上盡量拿得越多越好.
其實真的沒有去認清楚.
聖經對於人生的真相.
說明的本相是怎樣.
這個星期正好.
我剛剛病了.
又發燒又頭昏腦脹.
渾身又骨痛又肌肉痛.
真的有時說章寫得很累.
都不知道自己在寫什麼.
這個正正是人的限制.
其實正正是人的限制.
其實已經是小病了.
如果有什麼大病的時候.
人其實就更加無能.
好像什麼都做不到那樣.
這個是人的本相.
人生就是短暫.
但是神供應足夠的.
祂說我們可以用祂所給我們的.
託管我們的.
可以放在永恆的國度上.
在這個極其短暫的人生.
我們是願意將這些資源財降.
用在地上最終會燒毀的事情.
可能甚至我們發明了一些東西.
我們設計了一個很好的計劃等等.
還是我們是願意將這些資源財降.

$^{321}$放在永恆國度上面.
是有永恆價值的事情上面.
第十六節祂在說.
一切是從你手而來的.
都是屬於你的.
但是為祂認定將這些資源.
用在神國的事情上面.
是更加值得的.
其實差元奉獻有不少部分.
是用於差派宣教士身上的.
雖然現在網絡發達.
人要在網絡認識福音是容易了.
但是其實對於一些.
很少基督徒的國度裡面.
非基督徒他們不只是要知道理論.
怎樣可以去得救.
而是他們真的要見到.
有耶穌生命的人.
以致可以被耶穌.
以致可以被福音去吸引.
這是在永恆國度裡面重要的事.
有價值的事.
之前我去了一個南亞國家.
當地的宣教士每個星期日.
都去到不同山區的村落.
去支持不同的教會.
支持那裡的崇拜,兒童主日學.
這些教會其實都是宣教士所建立的.
你想想其實他每個星期.
怎樣駕車在那些很崎嶇,不平.
很多泥濘的道路去不同山區的地方.
其實是很辛苦的.
但是宣教士願意在永恆國度上面.
擺上自己的生命去傳福音.
去扶貧等等這些工作.
教會就在這件事上面去支持.
這份有永恆價值的工作.
一切都是從神而來.
當我們認清自己的本相.
就更加能夠去謙卑自己.

$^{361}$更加知道將生命,將這些資源.
獻在有永恆價值的事情.
永恆價值的國度上面.
好了,去到第十七節.
經文說:我的神,我知道你察驗人心.
喜悅正直,我以正直的心樂意獻上這一切.
現在我歡喜見你的百姓在此樂意奉獻給你.
這裡他說到察驗人心的上帝.
這次是在奉獻的事情上面去說.
弟兄姊妹我們是存著一個怎樣的心去奉獻的呢?.
會不會奉獻是為了好像想得到神額外的恩典.
神額外的祝福.
其他宗教都有說到奉獻.
譬如伊斯蘭教,佛教都有.
但是這些奉獻有時候好像是要換取一些東西.
譬如伊斯蘭教說的奉獻好像是去賺取功德.
就可以去到天堂,他們說的樂園.
佛教奉獻也是一種修行.
也是可以去得到一些東西.
雖然聖經有經文說到.
神是會賜福這些奉獻者的.
但是聖經說得更加多的是.
一切都是來自於神的.
就好像十四節這裡所說的.
萬物都是從你而來的.
所以奉獻其實理所當然.
不是存在一種去換取的心態.
不是想得到有什麼益處.
交換,交易這樣.
而是的確認定一切是從神而來的.
想去獻情給神.
是有一份樂意獻賞的心態.
不單止是自己一個人.
甚至這裡說的是整體的百姓.
其實很重要的是整體的群體.
一切的樂意獻賞.
接著經文也有說到.
神去喜悅正直.
這裡特別說的是奉獻者的正直.
經文沒有說神去喜悅那個憲制.

$^{401}$它是說神去喜悅這個人.
這個人的正直.
世上萬物本於神屬於神.
所以其實神不會因為.
這些本應屬於他的物品.
再去還給他.
然後他會很悅樂去喜悅嗎?.
其實不是的.
但是人就不同了.
人是神所看重的.
人可能會選擇跟不跟從神.
而神就是很喜悅.
這些願意跟從神的人.
是心中正直的人.
奉獻是一種外在的表現.
一種外在的行為表現.
但是聖經說神是喜悅正直.
即是整個人的內心.
都是要向著神.
其實有不同舊約的經文.
都有說神是看重內心.
而重於外在的表現.
神喜悅這個人的時候.
他就圓了他的獻賞.
譬如在《創世記》四章有提到.
神就是看重了亞伯和他的共物.
神是先看上亞伯這個人.
是他圓立的人看他為好.
其實在經文裡.
喜悅這個詞也有在這裡出現.
二十八章的時候說神是喜悅大衛.
就立他作全以色列的王.
神是去喜悅一個人.
我們要反思的是.
我們自己本身能不能得到神的喜悅.
重點是我們本身是怎樣.
並不是我們獻賞的物品.
我們的心向著神的時候.
樂意獻賞給上帝的時候.
神就越南.

$^{441}$於是大衛就繼續向神祈求.
你看他祈求的.
其實都是在說內在的方向.
是眾人的心.
是要由內而外的去向著神.
以至於得到神的喜悅.
同樣的.
當大衛為到去兒子所落門.
去祈求的時候.
你看看最主要也不是說.
祈求聖殿快快的去喜好.
他在這裡分了兩大部分.
前面那部分是比較詳盡.
是在說比較重要.
就是祈求所落門能夠全心的.
去遵守你的命令,法度,律例.
成就這一切的事.
最重要的仍然是全心去遵守神的命令,法度,律例.
大衛王知道重要的是所落門王.
他內心是否能夠正直.
是否能夠遵從神.
這才是最重要.
他是否整個人走在神的道路上.
之後才是在說見點.
優先次序其實很明顯.
領子妹.
我們當然是應當將神所交託給我們.
供應我們的資源.
甚至我們很懵福的.
神給我們的是足夠有餘的.
我們是將這些獻承給神,給上帝.
但我們切忌不要只是外在的去獻上.
我們更加要反照自己的內心.
自己的狀況.
我們是否在神眼中看為正直.
我們的生命有沒有全心去遵守神的命令.
我們是要內外,裡外一致的去向神獻上.
就好像我們所熟悉的羅馬書.
十二章一致的經文所說.
我們去獻上是全人的去獻上.

$^{481}$是活的祭來的.
是生命的獻上.
我們平時生活都應該是在呈現給我們的神.
這是神所喜悅的活祭.
神喜悅的真是主要不是祭物.
神喜悅的是這個人.
是我們在座的人.
願意我們丁子梅是全人的去獻上.
去真正得神的喜悅.
丁子梅今天透過大衛王祈求.
我們看到其實我們願意向神有多少的獻情.
是財幹,金錢,生命等等.
其實首先是和我們所認定這一位的上帝.
是一位怎樣的上帝有關.
聖經很清楚告訴我們.
也是我們經歷的神.
祂真是至高偉大的神.
祂是尊榮的.
祂是掌管萬有的位置.
祂配得我們對祂一切的獻上.
而見到神這麼偉大永恆的時候.
其實我們自然會回望自己.
其實是很渺小的.
但是當我們願意去獻上給神使用的時候.
的確可以成就了一些神要去成就的事情.
可能是我們都沒有想像過的.
而是在短暫的生命連日.
就可以將我們的生命,我們的資源.
真正投放在有永恆價值的事情上.
而最後.
神看重的始終不是獻上的物質是怎樣的.
神看重的是你.
是你這個人.
神去喜悅的,神去看重的.
是我們在座的每一個人.
我們是不是女女愛愛是向著神的.
然後神去悅納你的獻上.
神得著喜悅.
我們先低頭.
有個簡單的祈禱.

$^{521}$你可以先有一個墨土向著神的.
然後我帶領大家有個禱告.
我們一起祈禱.
創天造地,偉大的真神.
你的確是至尊,偉大,尊榮的那一位.
我們真的要去獻上我們的頌讚.
我們的生命,我們的資源.
去為你所願.
我們的生命,我們的資源去為你所用.
主啊,我們在你面前其實都真的很渺小.
但是你竟然去悅納我們.
你竟然願意去使用我們.
讓我們能夠參與在其中.
主啊,我們知道你願意我們傳人的.
是真的裡裡外外都是向著神你自己.
我們不是只有外在的去獻情.
卻是能夠我們整個的生命能夠向著你.
你是看重我們.
你是去喜悅我們這裡的人.
我們這樣去交上給你.
求你去幫助我們.
去建立我們.
賞與我們同在.
這樣祈禱交託奉主耶穌基督的名求.
阿們.
\newpage



\section{馬太福音 7:1-12}
\label{sec:031991}
\textbf{最好的東西}~~~~王正傑牧師
\newline
\newline
連結: \href{http://www.church.com.hk/acms/content.asp?site=cdc\&op=show\&type=product\&code=031991\&layout=sermon}{\texttt{031991}} ~~~~ 語音日期: 2026-04-19
\newline
\newline
\hyperref[sec:031990]{\small{< < < PREV SERMON < < <}}
~
\hyperref[sec:index_chronic]{\small{[返順時目]}}
~
\hyperref[sec:031992]{\small{> > > NEXT SERMON > > >}}
\newline
\newline
馬太福音 7:1-12
\newline
\begin{longtable}{cl}
\hline
\hline
章節 & 經文 (和合本修訂版)\\
\hline
7:1 & \begin{tabularx}{0.7\textwidth}{X} 「你們不要評斷別人，免得你們被審判。 \end{tabularx} \\ \\ \relax
7:2 & \begin{tabularx}{0.7\textwidth}{X} 因為你們怎樣評斷別人，也必怎樣被審判；你們用甚麼量器量給人，也必用甚麼量器量給你們。 \end{tabularx} \\ \\ \relax
7:3 & \begin{tabularx}{0.7\textwidth}{X} 為甚麼看見你弟兄眼中有刺，卻不想自己眼中有梁木呢？ \end{tabularx} \\ \\ \relax
7:4 & \begin{tabularx}{0.7\textwidth}{X} 你自己眼中有梁木，怎能對你弟兄說『讓我去掉你眼中的刺』呢？ \end{tabularx} \\ \\ \relax
7:5 & \begin{tabularx}{0.7\textwidth}{X} 你這假冒為善的人！先去掉自己眼中的梁木，然後才能看得清楚，好去掉你弟兄眼中的刺。 \end{tabularx} \\ \\ \relax
7:6 & \begin{tabularx}{0.7\textwidth}{X} 不要把聖物給狗，也不要把你們的珍珠丟在豬面前，恐怕牠們踐踏了珍珠，轉過來咬你們。」 \end{tabularx} \\ \\ \relax
7:7 & \begin{tabularx}{0.7\textwidth}{X} 「你們祈求，就給你們；尋找，就找到；叩門，就給你們開門。 \end{tabularx} \\ \\ \relax
7:8 & \begin{tabularx}{0.7\textwidth}{X} 因為凡祈求的，就得著；尋找的，就找到；叩門的，就給他開門。 \end{tabularx} \\ \\ \relax
7:9 & \begin{tabularx}{0.7\textwidth}{X} 你們中間誰有兒子求餅，反給他石頭呢？ \end{tabularx} \\ \\ \relax
7:10 & \begin{tabularx}{0.7\textwidth}{X} 求魚，反給他蛇呢？ \end{tabularx} \\ \\ \relax
7:11 & \begin{tabularx}{0.7\textwidth}{X} 你們雖然不好，尚且知道拿好東西給兒女，何況你們在天上的父，他豈不更要把好東西賜給求他的人嗎？ \end{tabularx} \\ \\ \relax
7:12 & \begin{tabularx}{0.7\textwidth}{X} 所以，無論何事，你們想要人怎樣待你們，你們也要怎樣待人，因為這就是律法和先知的道理。」 \end{tabularx} \\ \\
[1ex]
\hline
\hline
\end{longtable}
$^{1}$大英師妹平安.
今天我們要看的一段經文.
是一段可能我們回到教會.
若干年的時間 我們都聽了很多次.
不要論斷人.
今天我們都一起再嘗試.
透過整個登山的部分.
一起去了解 去認識.
究竟這裡說的不要論斷人.
對我們是有什麼提醒.
以及教導.
今天我們要講.
最好的東西.
我不知道對大家來說.
信耶穌最好的東西.
是什麼 可能是.
生活安舒等等.
一切都順利.
有些人可能對我們來說.
最好的東西是這些.
今天我們更加要看.
究竟天父給了我們什麼最好的東西.
大英師妹我們要知道一件事.
在我們的生命.
我們每一個基督徒.
其實每一個都有一套.
我們自己的神學.
我們信耶穌的人.
我們每一個都對自己的信仰.
是有自己一套的想法.
是有自己一套的做法.
而這套想法.
這套做法.
就不一定是建基於聖經.
而是很多時候.
我們個人對聖經的一些詮釋.
我們怎樣去理解一本聖經.
我們怎樣理解聖經的教導.
甚至可能我們在教會的經驗.
我們經歷過一些事情.

$^{41}$幫助我們建立了一套.
我們自己的想法和做法.
甚至可能是一些.
我們個人的喜好.
可能我們至少有些.
我們的偏好.
我們自己喜歡的東西.
甚至也有可能是一些.
來自教會的傳統.
以前教會的長輩.
是這樣教的.
所以我們都會建立了一套.
自己在信仰裡面的想法和做法.
所以我們每個人都有一套神學.
其實我們每個人有一套神學.
其實不是問題.
但是很多時候.
我們對這套神學.
就不單止是用來要求自己.
去評價自己.
是不是一個好的基督徒.
或者稱職的基督徒.
虔誠的基督徒之外.
我們同時間都會用.
我們自己這套的神學.
來去評價我們身邊的人.
身邊的弟兄姊妹.
甚至可能我們要求他們這樣做.
當我們見到他們不是這樣做的時候.
我們會怎樣呢.
我們覺得他是很不稱職的基督徒.
或者我們覺得他是很不敬虔.
所以我們每個人.
都有這一套的神學.
所以在網上.
有人這樣去分享.
亦都在某個團契.
都有分享這一段.
在網上的法文.
這段法文在說什麼呢.

$^{81}$就是說.
我聽過教會最離譜的東西.
OK.
都是源於一些人.
自己的神學.
對信仰的理解.
然後對別人的要求.
有一個他這樣分享.
他說發生在他二年級的時候.
他就說被教會學校的老師.
去誣陷他偷東西.
老師說.
為什麼他認為這個同學.
去偷東西呢.
就是跟他這樣說.
因為神有聲音告訴我.
是你偷的.
另外亦都有一個這樣分享.
他說他阿姨回教會三十多年.
很熱心幫教會去做義工.
亦都有洗禮.
亦都有奉獻.
過身前的三年.
因為中了風.
行動不便.
所以就回不了教會.
當這個阿姨去過身之後.
他表哥就通知教會.
請教會的牧師幫他們.
做一個葬禮的安排.
牧師當時就跟這個阿姨.
就是跟這個表哥說.
他的阿姨怎樣呢.
就說他因為太久沒有回教會了.
就幫不了他.
連會友證明都出不了給他.
安排不了申請.
基督教的墳場的位置.
最後家庭就只能夠.
選擇海葬.

$^{121}$另外亦都有.
基督徒去醫院做教會.
做一個醫院的探訪.
這個教會的義工.
就跟那個.
在隔壁床的婆婆說.
耶穌愛你啊婆婆.
然後這個婆婆就說.
不用啦.
我信佛的.
然後義工就跟婆婆說.
你想不想出院.
我不知道弟兄姊妹.
在你生命裡頭有沒有聽過.
在教會的一些教導.
一套變人的神學.
甚至可能是.
我們有一套我們自己的神學.
我們都會要求人.
甚至用來評論人.
但這些很多時卻是.
帶來很多的破壞.
甚至可能很多的傷害.
今天我們去看的.
這一段經文講到論斷人.
我們如果真的要去明白.
究竟講論斷人是什麼意思的時候.
我們必須要.
去看完整段登山寶訓.
我們才能夠完全理解.
明白究竟論斷人.
是在講什麼意思.
當我們看看登山寶訓.
其實由一開始.
主耶穌是透過登山寶訓.
去教導那些跟隨祂的人.
他們在悖改了之後.
其實生活應該是怎樣的.
他們怎樣次世代人.
耶穌在這個時候.

$^{161}$在登山寶訓就要去教導.
那些跟隨祂的人.
主耶穌先講的是.
祂說祂來.
祂不是要去廢掉律法.
而是要成全.
祂不廢掉律法是一件很重要的事.
因為律法是要讓我們知道什麼.
知道罪.
所以主耶穌說祂來.
祂不是要去廢掉律法.
而是要成全.
然後祂更加吩咐信徒.
在地上作鹽作光.
我們要發揮我們信徒的影響力.
更加要求的是.
我們跟隨祂的人.
要勝過民事和法律賽人的義.
我們要留心的是.
很多時候我們一聽到.
法律賽人和民事.
我們都覺得祂很壞.
我們見主耶穌經常都罵他們.
但是我們.
要知道的就是.
其實民事和法律賽人.
他們對信仰是很有要求的.
他們其實是要求自己.
是要怎樣去.
遵行律法.
遵行律法書所講的東西.
當然.
他們最大的問題是因為他們怎樣.
做不到.
做不到的時候.
就會假冒違善.
有時候要假裝.
但是這裡主耶穌是要求什麼.
要求我們其實跟隨祂的人.
我們的義.

$^{201}$是要勝過民事法律賽人的義.
我們要做得比祂更加好.
民事和法律賽人.
他們是假冒違善的.
我們真是要怎樣.
從心而作.
我們真是為主作的.
所以主耶穌更加說.
民事法律賽人的義.
他們是假冒違善.
是要將善事行在人的面前.
故意叫人看見.
我們專門做給人看的這些民事和法律賽人.
所以去到今天我們看的這段經文.
馬太福音7章1至12節.
主耶穌就更加進一步的.
去命令跟隨祂的人.
不要去論斷人.
免得都被論斷.
在這裡我們看經文的時候.
經文一開始.
7章1至12節.
主耶穌就命令.
其實他是在用一個命令.
他命令跟隨祂的人不要去論斷人.
免得被論斷.
論斷的意思是什麼呢.
當然我們經常回教會.
我們都可能會去.
有些評價有些評論.
今天你回來可能你坐在這裡.
你都會去評價一下.
可能你去評價一下.
可能見到鏡輝今天穿得這麼帥.
你都會想一下今天.
二領全道他究竟有沒有走音.
跟著我們又可能會評價一下.
我覺得今天的冷氣.
不是很夠.
我們都會評價一下這些.

$^{241}$我們有很多評價.
但這些評價只是一個評價.
就沒有去到一個論斷.
為什麼我說沒有去到論斷呢.
因為論斷是什麼呢.
論斷其實是一個審判.
我們用廣東話去形容.
看死了一個人.
他就是這麼壞.
你回來崇拜.
你睡覺.
他很不正經.
我們會加了這些罪名.
我們將一個人定罪.
這個在教會裡面.
經常都會很容易出現的情況.
主席叔叔這裡就命令我們.
跟隨祂的人不要去論斷人.
免得被論斷.
所以說.
我們論斷人.
我們同樣都會被別人論斷.
被審判被定罪.
剛剛很簡單.
你說得人人又可以怎樣.
說回來了.
這裡我們就字面理解就可以這樣解釋.
不過我特別提一提.
我特別要提一提.
我們不是說我們在教會.
不去分辨對錯.
不去提醒彼此提醒.
我不是說這件事.
聖經告訴我們不要去.
將人去審判.
將人去定罪.
看死了一個人.
回到教會.
我們仍然可以去分辨.
對與錯.

$^{281}$是否合乎聖經.
還要我們給予彼此提醒.
用愛心挽回.
在這方面稍後有機會的話.
我會再跟大家講多一些.
關於我們怎樣可以在主裡面.
我們能夠分辨對錯.
可以給予彼此提醒.
我同時再講.
為什麼不要論斷.
第二節他這樣說.
因為你們怎樣論斷人.
也必會怎樣被論斷.
你用什麼良氣去良給人.
也必用什麼良氣良給你們.
所以我們很明白.
我們純粹理解.
你怎樣去論斷人.
就會被人論斷.
你怎樣說人就會被人說.
還要你用什麼尺度.
去量度人的時候.
你也是這樣被人量度.
我們就這樣字面去理解.
但對於猶太人來說.
在猶太的文獻.
其實這是猶太文獻的.
一種格式的表達.
是一種格言式的表達.
因為當時對於.
那些聽眾和讀者.
其實他們有一個很重要的觀念.
就是對他們來說.
在這個世界上.
只有誰有權並去審判.
就只有神.
所以在這裡其實原文的表達.
它說你怎樣去論斷人.
也必怎樣被論斷.
你用什麼良氣良給人.

$^{321}$也必用什麼良氣良給你們.
這裡其實是在說什麼.
就是在說當你怎樣去論斷人.
你都會被遷附.
好像你論斷人一樣去論斷你.
當你用什麼良氣.
去良給人的時候.
你的遷附就會怎樣.
用什麼良氣.
用什麼良氣去良人的時候.
這樣去量度你.
所以當我們去理解這裡.
其實原來當我們論斷人的時候.
我們都同樣怎樣.
神會用我們怎樣去論斷人的方式.
來去論斷我們.
定我們的罪.
當我們用什麼良氣去良給人的時候.
天父都要怎樣.
用我們用什麼良氣去量度人的時候.
這樣去量度我們.
所以耶穌說我們不要去論斷人.
而當時.
民視和法理賽人.
他們在人的面前.
只是做給人看他們的假冒違善.
主耶穌在這裡就引用了.
他在木工房.
我們都知道主耶穌他本身做什麼.
木匠來的.
所以他用他在木工房.
那些事物來做一個比喻.
他就指出.
民視和法理賽人.
他們都是假冒違善的.
他就用別人眼中的刺.
就是一些很小的木片.
一些叫木舌.
其實在希臘話來講.
是比喻一些很微小.

$^{361}$很微小的東西.
微小到很小的東西.
然後他再形容.
自己眼中的良木是怎樣.
是支撐防止華人良.
如果我們用我們華人的理解.
大概是這樣.
我們看到一些華人良.
這裡其實他要講的是.
原來當我們.
好像民視和法理賽人.
他們這樣假冒違善.
就是他們看不見.
自己眼中的良木.
只看到別人眼中有事.
他們看不到自己的大錯.
只看到別人的小錯.
所以主耶穌在這裡要指出.
民視和法理賽人.
他們都是假冒違善.
就是他們只看到別人.
微小的過錯.
就用綠法.
將別人審判.
去定罪.
這不是他們沒有過錯.
而是他們不去看自己.
好像良木那麼大的過錯.
還要很自然地指出別人的錯.
將別人去定罪.
所以我們廣東話有一句.
很多時候有些人只會說別人.
就不懂得說自己.
對不對.
我們只看到別人有問題.
看不到自己有問題.
什麼問題都是別人的問題.
所以.
主耶穌說不要將別人.
去審判 去定罪.

$^{401}$因為當你怎樣去量度別人的時候.
你的天賦亦都會怎樣去量度你.
更加重要的就是.
我們在神的面前.
最重要的一件事情 弟兄姊妹.
其實我們都是蒙恩的罪人.
你看看你旁邊那個.
他不給你好.
對不對.
我也好不過你.
我們都是罪人.
簡單說.
其實我們都是好壞的.
好不好你幫我和你旁邊那個說.
其實你好壞的.
在昨天早上的時間.
我就和兩位同工.
一起去神學院.
上了一個課堂.
課堂是在說.
七種的內在批判家.
其實我們每個人都在內在批判.
有些人可能.
對自己有很高要求.
有些人是完美主義的.
所以要求自己很高.
在課堂的中間.
我就去問一問我的老師.
我當然承認.
我都有內在批判.
我都是一個.
其中一個我叫貶低者.
我都是一個貶低者.
我經常覺得自己很差 我覺得別人很厲害.
我都是一個這樣的人.
我見到.
我自己有內在批判的時候.
我見到更多的是.
很多時候身邊的人都好喜歡去批判人.
我問他其實很多時候.

$^{441}$是不是因為我們內在有.
有這個自我的批判.
所以我們都好容易用我們這種.
內在的自我批判來去批判.
我們身邊的人.
甚至批判社會批判這個世界等等.
我老師就告訴我.
是呀 沒錯.
正正是因為我們有這種內在的自我批判.
我們很多時候同樣用這種批判.
來去批判.
我們身邊的人.
甚至回到教會.
我們可能都對自己有一種內在批判.
我們有自己一套承諾.
但同時間我們都在批判.
我們身邊的人.
在這裡有七種的.
內在批判家.
不知道第一種是哪一種.
第一種是完美主義者.
完美主義者就是怎樣.
什麼都很高要求.
還可以做得更好的.
所以你回到教會的時候.
如果完美主義者就可能在這裡留意.
哪裡有錯 哪裡有不好.
看看power point 你打錯字了.
然後看看.
剛才按了一首歌.
你的power point 你按得這麼慢.
慢了 我唱不到歌詞.
等等.
我們可能完美主義者就回來.
留意著這些.
另外一種是.
內在的控制者.
這種人又怎樣呢.
對自己有很多壓抑.
對自己有很多壓抑的時候.

$^{481}$你都不可以.
為什麼我自己回來崇拜.
我坐在這裡這麼安靜.
為什麼你又這麼多話說.
為什麼你又這麼有反應.
為什麼你又這麼興奮.
為什麼你在笑.
很多人其實就會可能.
不知道是不是這種內在批判者.
控制者就是怎樣.
我們自己要求自己很壓抑.
很希望要求身邊的人都這麼壓抑.
第三種就是任務大師.
任務大師是什麼來的呢.
就是一些人怎樣.
不停在這裡工作.
做完一樣接一樣.
做完一樣接一樣就是他為什麼呢.
因為他唯有這樣.
他才覺得自己叫競田.
但同時間這些批判者.
同時間會批判身邊的人.
就會怎樣呢.
為什麼我做得這麼辛苦.
怎麼回來這麼舒服呢.
還可以睡覺呢.
這種就是任務大師.
另外一種就是貶低者.
就是覺得自己怎樣做都是不好的.
就算自己做到的時候.
都會跟別人說什麼呢.
就是人人都做到了.
阿謀阿壽都做到了.
我覺得這些很小事而已.
另外還有一種是摧毀者.
摧毀者是怎樣呢.
是一些很小的事.
弄得很大.
嘩.
今日國編講到.

$^{521}$講得這麼差.
教會明天就會倒閉.
這種是摧毀者.
將很小的事化得很大.
然後跟著.
一個最惡感的製造者.
就將你做錯一些事.
怎樣呢.
加個罪名.
嘩.
你今天穿成這樣.
幫我看看你旁邊那個穿得好不好.
你今天穿成這樣.
你很不聖潔.
是不是.
你很不尊重神.
你沒有敬畏神的心.
是不是.
這種是最惡感的製造者.
還有一種是信徒者.
信徒者的意思是什麼呢.
就是有一種.
有些人很喜歡有些框框框著.
最好你就是.
但是好啦.
當我們見到旁邊那個人.
他沒有做點點點點點的時候.
我們會覺得怎樣.
嘩這個不行.
我們就說這些人真的不愛主.
是不是.
他不是基督徒來的.
他都不遵守規則.
我們很容易會這樣.
等於我不知道你是.
哪一種的批判家.
但Jesu.
他命令我們.
不要論斷人.
我們要小心.

$^{561}$我們都會被論斷.
然後跟著.
等於姊妹.
去到這裡.
你幫我看看你現在的狀態.
因為接下來這裡.
這一部分是有點煩的.
如果你現在的狀態是.
已經開始發現.
你氣數已盡的話.
你隨便.
竭力進入安息就可以了.
OK 我們預備了.
我們開始進入第六節.
當我們看到第六節.
我們是很不明白的.
前面在說論斷人.
是不是 弟兄姊妹.
我們不要論斷人.
因為如果論斷人.
同仁被纖腐會這樣論斷.
但到了第六節就無端端.
提到你不要將聖物.
給狗.
也不要拿珍珠去雕在珠前.
到了這裡.
其實我們不太明白.
為什麼第六節無端端.
說了一件事.
一句我們完全不太理解.
也不要緊.
我們看到七節到十一節.
說了什麼.
是不是在教導我們.
禱告蒙英文.
我們就會想.
是不是開了一個新的主題.
說另一件事.
是不是主題鬆無端端.
說完論斷人.

$^{601}$所以當我們要去解明.
七節到十一節.
我們必須用整段的登山寶訓.
馬太福音五節到七章.
去理解.
首先我們要知道.
當我們看馬太福音的五節到七章.
其實主要是教導.
教導那些跟隨他的人.
他們的生活要怎樣.
其實主要是針對當時的民事.
和法律採人的信仰.
所以到了馬太福音的七章.
一節到五節.
我們不要去論斷人.
到了馬太福音的七章.
十二節.
你會看到主耶穌說什麼.
你要人怎樣待你.
你就要先怎樣待人.
我們反而看到的一件事是什麼.
顯然主耶穌是在講同一個主題.
第十二節.
其實這個好像是.
整段經文的一個.
全個主題的一個結論.
所以主耶穌他在.
七節到十一節的時候.
他不是開了一個新的課題.
他並不是.
講講下論斷人.
就轉過去教我們祈禱.
他不是在做這件事.
而是他在講同一樣東西.
我們就要問究竟他在講什麼.
當我們去看七節八節.
這裡主耶穌在講什麼.
他在講.
祈求的就會得著.
尋找的就尋見.

$^{641}$叩門就給他開門.
其實這兩節是包含了.
的具型是相同的.
所以指出是怎樣.
指出當我們去尋求天父.
就必定會尋找.
神給我們的應許.
包括祈求就會得著.
尋找就會尋見.
叩門就給你開門.
但我們更加要問的是.
祈求尋找叩門.
我們得到什麼呢.
這裡主耶穌就用了.
一個很簡單的比喻.
他說而只去求餅.
求魚就不會給你.
接著說世間的父親不好.
其實他在講父親不好.
其實如果在原文理解.
其實意思是什麼呢.
就是父親是邪惡.
為什麼這樣要表達呢.
因為當時的父親.
教導他們的子女.
可以嚴厲扎打.
嚴厲扎打到一個地步是什麼呢.
打死他們都可以.
他們的父親是可以操生死大權.
但是儘管.
父親是不好是邪惡的.
但是為父的.
仍然會將好東西給兒女.
更何況我們天上的父親.
怎會不將好東西給我們呢.
那我們就再進一步.
再來我們就會問.
在主耶穌這裡他要指出的就是.
當猶太人去向天父尋求的時候.
他們就會得著.

$^{681}$必然得著的時候.
其實是在講什麼呢.
就是猶太人一直祈求的.
那位彌賽亞.
並且天父要將這個最好的東西.
賜給他們就是什麼呢.
耶穌基督.
這個賜給他們耶穌基督.
是遠超過他們的所求的.
以至.
他們這一班的人.
可以怎樣.
不再去被.
定罪.
不再去靠律法清義.
不再去靠行為清義.
並且得救可以去被寬恕.
所以今天弟兄姊妹.
我們要明白.
天父給我們最好的東西.
是在給什麼給我們.
就是在給耶穌基督給我們.
耶穌基督來到這個世界上.
就是要讓我們不再去被定罪.
並且得救.
同樣的在我們在教會的時候.
我們更加要做一件事是什麼.
我們都不要去定.
人的罪.
不要按照人的行為.
去定人的罪審判人.
並且我們要.
能夠彼此.
去寬恕.
在.
《仰向福音》三章十六至十八節.
那裡這樣說.
神愛世人.
甚至將他的獨生子賜給他們.
叫一切信他的不至滅亡犯得永生.

$^{721}$因為神差他的兒子降世.
不是要定世人的罪.
乃是要叫世人因他得救.
信他的人不被定罪.
不信的人罪已經定了.
因為他不信神獨生子的名.
主耶穌來到這個世界.
不是要定我們的罪.
乃是叫我們要怎樣.
要因為他而去得救.
所以在這裡我們回頭.
我們就明白到第七章.
六節那裡在說什麼.
這一節其實是延續著前面.
不要論斷人.
就是主出豬和狗.
是被視為不潔的.
為什麼他們被視為不潔.
是因為他們不接受真理.
所以你將祭物去餵狗.
那些珍珠那麼貴重的東西.
給那隻豬是沒有意義的.
因為他們都不懂得欣賞.
不明白那個價值.
反過來還要怎樣.
咬那個給他們的人.
所以主耶穌在這裡.
要指出的是什麼.
那些民事和法律採行.
他們去論斷人.
按照法律去定人的罪.
是因為他們都不願意接受.
也都不明白.
那個寶貴的價值.
仍然以這個守律法來去清義.
所以不接受天父給他們.
最好的東西.
就是耶穌基督.
在這個有控福音的八章一至十二節.
我們就見到.

$^{761}$主耶穌為什麼親身去示範.
不去定人的罪.
派出人和民事.
去將一個只是行淫的婦人.
帶到耶穌的面前.
按照摩西的律法.
要用石頭將她扔死.
但是當時的羅馬政府.
是不容許猶太人.
在他們的法庭有處死的權力.
所以在這裡.
民事和法律採行就要去試探耶穌.
要得罪一些告她的把柄.
主耶穌就跟他們說.
你們中間誰沒有罪.
誰就可以用石頭去打她.
當他們聽到這番說話的時候.
全部人都走了.
主耶穌就對那個婦人說.
沒有人去定你的罪嗎.
我都不定你的罪.
她說從此不要再犯罪了.
弟兄姊妹.
今天我們來到神的面前.
我們中間.
究竟誰是沒有罪的呢.
我們又有誰真的可以去.
定人的罪呢.
主耶穌.
對那個犯罪的淫婦說了.
沒有人去定你的罪.
我都不定你的罪.
去吧.
最重要的一句是.
從此不要再犯罪了.
所以弟兄姊妹.
今天對我們來說.
誰是不接受耶穌基督的人呢.
不接受天父給我們.
這個最好的東西的人呢.

$^{801}$就是當我們都好像民事.
或者法律採人那樣.
用一些所謂的.
律法.
甚至可能是我們.
個人的神學.
你去定人的罪.
去審判人.
去看死人的時候.
我們就成為了主耶穌.
所說的那些怎樣.
將聖物給狗.
真之雕在豬前.
我們都不懂得欣賞.
不懂得去接受.
跟大家分享.
近年在教會其實.
有一些的媒體.
去反映了教會都出現.
一些的情況.
包括有很多人慨嘆.
就是教會的年輕人.
不起來.
不參與服侍.
也都有些人就說.
請了個全清的牧者回來.
教會就.
一定興旺但是就.
對教會有很多的埋怨.
支持不夠等等.
支持不夠.
也都有些人就去.
討論的就是.
教會的年長的人就.
嫌歌太吵.
年輕人就嫌歌太老.
今天我們剛才.
唱的那些詩歌都沒有太吵.
是不是.
很多這些的.

$^{841}$所謂的爭拗在.
教會裡面有很多的指責.
有很多的審判.
有很多的定罪.
但是等於我們都要.
認真去面對的就是.
等於我們不要去論斷人.
也都不要去.
隨便的去定人罪.
其實等於我們要留心的就是.
其實這些問題.
其實都不是一些.
很嚴重的問題.
但是我們必須要面對的就是.
等於當我們慨嘆教會的年輕一代.
未必肯出來.
或者不容易出來.
承擔事奉的時候.
我們都要問的就是.
當年輕一代肯出來事奉的時候.
其實我們又給了多少支持.
給了多少支援呢.
還是我們都很多的要求.
可能有時.
我們甚至.
有些挑剔.
當你有這麼多要求.
有這麼多挑剔的時候.
又哪有人隨便願意.
出來承擔事奉呢.
是不是.
其實我們在職場.
我們很明白的.
例如上班.
其實對你來說.
除了賺錢之外.
最影響我們的是什麼.
人事問題.
昨天我聽講座的時候.
就認識了.

$^{881}$我前面的一個弟兄.
那個弟兄很主動.
為什麼呢 是因為他做健康顧問.
我知道他遲些.
都會聯絡我來關心我的健康.
跟他聊聊天.
說起.
工作辛苦不辛苦的時候.
我也跟他說了一番話.
我說其實從來.
從來事是做不死人的.
對不對.
很多事做是不會做死人的.
但是回到公司.
你最辛苦的那件事是什麼.
是人事問題.
對不對.
同樣回教會都是這樣.
事會不會做得死你.
教會的事會不會做得死.
做不死我們的.
但是很多時最磨我們的是什麼.
是人事的問題.
並且其實.
我覺得頂多我們要去.
學習有一個開放的心.
我們要去.
祝福年青一代成功.
當他們出來.
事奉的時候我們支持他們.
鼓勵他們.
就算有錯的時候.
就給他們一點空間成長.
有哪個人出來.
第一次做事奉就做得很好.
對不對.
當然是怎樣.
經過時間他不斷去鍛煉.
不斷累積經驗.
他慢慢去成長.

$^{921}$越做越好.
同樣的我都不可以要求.
我的同工.
新入職的.
出來事奉沒多久.
我不可能要求他.
所做的事跟我一樣.
如果.
他一入職已經做到.
跟我一樣做到的事.
其實是什麼來的.
對不對.
所以我們要給.
空間年青一代.
教會.
我們經常討論歌曲的問題.
在教會.
老實說有舊歌.
有新歌是很平常的.
對不對.
什麼青年性歌.
頌主新歌.
頌主新歌好像丟光了.
我們換世紀頌讚的時候.
就把頌主新歌送給弟兄姊妹.
自己拿回家留念.
我還記得.
我當年還在土層.
有新歌有舊歌.
老實說我們很多人都會說.
教會唱太多新歌.
我們不會唱.
我很不開心.
但是弟兄姊妹.
我們就說.
我們唱那麼多新歌.
不懂唱.
不如我們多唱舊歌.
誰不知原來對某些人來說.
這些舊歌都是他的新歌.

$^{961}$對不對.
你沒唱過的.
就是你的新歌.
我有些舊歌真的沒唱過.
那就變成我的新歌.
但是弟兄姊妹.
無論舊歌也好.
新歌也好.
都是唱著唱著就唱得到.
那我們就一起去學.
唱新歌也好.
唱舊歌也好.
所以弟兄姊妹.
回到今天的經文是.
我們不要去論斷人.
我們要建立一個群體是什麼.
是一個天父給了我們最好的東西.
就是耶穌基督.
祂是要拯救我們.
我們要去得罪堯叔.
祂都不去定我們的罪.
我們都不要彼此的去定罪.
然後.
去到最後.
主耶穌祂說的.
正因為主耶穌祂來.
不是要定人的罪.
而是要人因為祂得救.
所以主耶穌祂說什麼.
你希望別人怎樣待你.
你也要怎樣待人.
所以在這裡說什麼.
你希望天父怎樣待你.
你也要怎樣的去待人.
所以如果你希望天父是寬恕你.
你也要怎樣的去待別人.
所以這裡說.
我們願意天父怎樣去寬恕我們.
我們就要怎樣去寬恕.
你去待我們身邊的人.

$^{1001}$而我們怎樣能夠做得到.
我們能夠可以寬恕待人.
經文告訴我們是什麼.
因為這個就是律法和先知的總綱.
律法和先知道理的總綱是什麼.
就是登山寶佛整個的核心信息.
耶穌跟他們說什麼.
說你要怎樣.
你要盡心盡性盡意愛主你的神.
就是誡命中的第一.
就是最大的.
其次也相仿就是愛人如己.
將這條誡命是律法和先知一切的道理的總綱.
所以簡單說.
我們怎樣能夠可以.
以寬恕去待人就是因為什麼.
我們愛神我們愛人.
當我們愛神.
愛人的時候我們就能夠.
可以以寬恕來去待人.
所以丁姐妹.
今天我們去看.
這些信息很簡單.
我們不要去論斷人.
我們要以寬恕去待人.
我們最重要的一點.
就是愛神愛人.
當我們真的愛神.
真的愛人的時候.
我們就可以怎樣.
可以去不去論斷人.
不去定人的罪.
可以以寬恕來去.
待我們身邊的人.
因為愛神愛人如己.
以致我們都能夠在無論任何事.
我們願意別人怎樣待我們.
我們就先怎樣去待人.
不會用法來去論斷人.
審判人定人的罪.

$^{1041}$而是以寬恕來去待人.
和你分享一個見證.
是我.
自己親身的經歷.
這個見證是這樣的.
就是在很多年前.
當時我有一位.
童工.
因為有些工作上的事.
我要和他分享.
我要去提醒他.
讓他知道他要怎樣去面對.
怎樣去處理.
和他分享完之後.
他當晚六點鐘打電話給我.
他和我說.
他把我和他說的話.
和老闆說了.
然後他說.
你要有心理準備.
我聽完之後就很驚訝.
我在跟你說.
我和你分享的東西.
為甚麼我要有心理準備呢.
老闆叫他.
把他的東西給我做.
我想不到他想甚麼.
不要緊.
隨便吧.
等明天回去.
看看見老闆會怎樣.
第二天回到教會的時候.
我和這位童工.
都一起見老闆.
既然老闆真的可以.
真的和他.
做回這個澄清.
因為這些工作.
你真的需要怎樣做.
原來我這位童工.

$^{1081}$他誤會了.
我在整鬼他.
在害他.
所以他誤會了.
所以他和老闆.
都是同樣地表達.
但我老闆都做了些事.
他在背後都找其他人問.
究竟我是否可以.
整鬼這個童工呢.
所以他了解了之後.
其實一般來說.
我們應該都很生氣.
這些傻子.
我們應該不要理他.
不要和他做朋友.
應該和他翻臉.
以後看你死就算了.
我們應該正常這樣想.
但我做了一件事.
我和這位童工.
欠缺的是甚麼呢.
就是信任.
所以我要和這位童工.
建立一個信任.
用最好的方法.
來建立信任.
沒錯.
我還在吃飯.
我還要請他吃.
燒牛肉.
我還記得那天燒牛肉.
一個人三百多元.
結果.
由那頓飯開始.
我真的和這位童工.
可以建立信任.
到今天.
我很想告訴你.
在主裡面.

$^{1121}$是很好的童工.
很好的朋友.
很好的弟兄.
我們現在有甚麼困難.
或是有甚麼事情.
都可以分享和提醒.
兄弟姊妹.
教會我們不是要建立一個地方.
要去.
論斷證人的罪.
而是我們要建立一個群體.
我們要怎樣呢.
不是去吃飯.
這麼簡單.
吃飯是要的.
我們是要建立一個地方.
是我們能夠給予包容.
給予饒恕的一個地方.
以致我們真的建立一個愛的群體.
弟兄姊妹.
今天你來到神明面前.
你願不願意.
都放下.
自己那個很牢固的一套神學.
我們不去隨便.
去論斷人.
而是建立一個.
寬恕.
給予包容的群體.
由我們一同低頭.
我們有一段禱告的時間.
差主爺叔呢.
因為你的愛.
我們德蒙你的救贖.
因為你的犧牲.
讓我們罪得赦免.
從前我們被罪捆綁.
今天我們得到真正的釋放.
真正的自由.
當我們回來教會.

$^{1161}$主我們懇求你幫助我們.
教導我們.
讓你的愛激勵我們.
以致我們真的能夠學習.
不要去論斷人.
不去定人的罪.
讓我們真的在教會當中.
縱使可能我們有一套.
我們自己的神學.
我們有一套我們自己.
怎麼去跟隨耶穌的方式.
讓我們.
不要隨便的.
去論斷人定人的罪.
以致我們真的能夠在教會.
我們去學習建立一個.
彼此寬恕.
彼此包容.
彼此接納的一個群體.
讓你的愛.
在我們當中.
你能夠得著當德的榮耀.
讓我們因為這份愛.
我們能夠去祝福.
我們身邊的人.
主我們就懇求你.
親自幫助我們.
讓我們來到你面前.
得到你的娛樂.
和使用.
讓我們真的能夠成為.
你明亮的燈台.
可以為你作炎作光.
可以榮耀你.
見證你.
主我們多謝讚美你.
不過使用我們.
我們飲祈禱交托.
是奉求我主耶穌基督得勝名字祈求.
阿們.

\newpage



\section{馬可福音 8:29-9:1}
\label{sec:031992}
\textbf{只有一位受苦的上帝能救我們}~~~~鄧瑞強牧師博士
\newline
\newline
連結: \href{http://www.church.com.hk/acms/content.asp?site=cdc\&op=show\&type=product\&code=031992\&layout=sermon}{\texttt{031992}} ~~~~ 語音日期: 2026-04-19
\newline
\newline
\hyperref[sec:031991]{\small{< < < PREV SERMON < < <}}
~
\hyperref[sec:index_chronic]{\small{[返順時目]}}
~
\hyperref[sec:031993]{\small{> > > NEXT SERMON > > >}}
\newline
\newline
馬可福音 8:29-9:1
\newline
\begin{longtable}{cl}
\hline
\hline
章節 & 經文 (和合本修訂版)\\
\hline
8:29 & \begin{tabularx}{0.7\textwidth}{X} 他又問他們：「你們說我是誰？」彼得回答他：「你是基督。」 \end{tabularx} \\ \\ \relax
8:30 & \begin{tabularx}{0.7\textwidth}{X} 於是耶穌切切地囑咐他們不可對任何人說起他。 \end{tabularx} \\ \\ \relax
8:31 & \begin{tabularx}{0.7\textwidth}{X} 從此，他教導他們說：「人子必須受許多的苦，被長老、祭司長和文士棄絕，並且被殺，三天後復活。」 \end{tabularx} \\ \\ \relax
8:32 & \begin{tabularx}{0.7\textwidth}{X} 耶穌明白地說了這話，彼得就拉著他，責備他。 \end{tabularx} \\ \\ \relax
8:33 & \begin{tabularx}{0.7\textwidth}{X} 耶穌轉過來看著門徒，斥責彼得說：「撒但，退到我後邊去！因為你不體會神的心意，而是體會人的意思。」 \end{tabularx} \\ \\ \relax
8:34 & \begin{tabularx}{0.7\textwidth}{X} 於是他叫眾人和門徒來，對他們說：「若有人要跟從我，就當捨己，背起自己的十字架來跟從我。 \end{tabularx} \\ \\ \relax
8:35 & \begin{tabularx}{0.7\textwidth}{X} 因為凡要救自己生命的，必喪失生命；凡為我和福音喪失生命的，必救自己的生命。 \end{tabularx} \\ \\ \relax
8:36 & \begin{tabularx}{0.7\textwidth}{X} 人就是賺得全世界，賠上自己的生命，有甚麼益處呢？ \end{tabularx} \\ \\ \relax
8:37 & \begin{tabularx}{0.7\textwidth}{X} 人還能拿甚麼換生命呢？ \end{tabularx} \\ \\ \relax
8:38 & \begin{tabularx}{0.7\textwidth}{X} 凡在這淫亂罪惡的世代，把我和我的道當作可恥的，人子在他父的榮耀裡與聖天使一同來臨的時候，也要把那人當作可恥的。」 \end{tabularx} \\ \\ \relax
9:1 & \begin{tabularx}{0.7\textwidth}{X} 耶穌又對他們說：「我實在告訴你們，站在這裡的，有人在沒經歷死亡以前，必定看見神的國帶著能力臨到。」 \end{tabularx} \\ \\
[1ex]
\hline
\hline
\end{longtable}
$^{1}$各位頂小廟,早晨.
神的恩典尚要你們同在.
在這個蒼茫大地.
群雄俗綠中原.
拜的都是戰神.
當大家不再講規矩.
回到叢林法則.
最後就看看誰能夠在這個施殺裡面取勝了.
一生追求溝通理性的.
那位哲人.
哈伯瑪斯.
近日離世.
他屍骨都未寒.
各國的行動的重點.
就是儲備軍火.
大家都在想.
只有戰神.
才能夠救自己.
第一世紀.
羅馬帝國的鐵蹄.
踐踏在巴勒斯坦的土地上.
猶太人苦不堪言.
他們祈求.
救世主.
好像一個.
無敵的戰神那樣降臨.
驅散敵人.
救助他們.
他們望穿秋水.
終於來了一個耶穌.
他們看到祂行神蹟.
對祂滿有期望.
他們在想.
這個是不是我們期望的戰神呢.
我們聽聽耶穌基督的話.
就是今天我們講到的經文.
記載在《馬可福音》第8章.
我們由第29節讀到第9章第1節.
我們的經文印製在我們的程序表.
教後的部分.

$^{41}$大家可以打開.
我們由8章第29節開始讀起.
耶穌問他們.
祂的門徒.
你們說.
我是誰.
彼得回答說.
你是基督.
耶穌就禁戒他們.
不要告訴人.
從此.
祂教訓他們說.
人只必須受許多的苦.
被長老祭司長和民事氣絕.
並且被殺.
過三天復活.
耶穌明不明白的說這些話.
彼得就拉著祂.
勸祂.
勸這個字.
和下一節.
第33節那個責備那個字.
其實完全是同一個字.
耶穌斥責污鬼的時候.
也是用這個字.
所以.
彼得不是勸祂.
彼得在這裡.
簡直就是斥責耶穌.
罵祂.
第33節.
耶穌轉過來.
看著門徒.
就責備彼得說.
和剛才那個勸字同一個字.
責備彼得說.
撒旦.
退我後邊去吧.
因為你不體貼上帝的意思.
只體貼人的意思.

$^{81}$這個是非常罕有的場面.
使徒互相斥責.
耶穌責備彼得.
你站在撒旦的那一邊.
然後耶穌面對一個更加大的群眾.
祂這樣說.
第34節.
於是叫眾人和門徒來.
對他們說.
若有人要跟從我.
就當社己背起他的十字架來.
跟從我.
因為.
凡要救自己生命的.
必喪掉生命.
凡為我和福音喪掉生命的.
必救了生命.
人就是賺得全世界.
賠上了自己的生命.
有什麼益處呢.
人還能拿什麼換生命呢.
在這個淫亂罪惡的世代.
把我和我的道當作可恥的.
人只在他負的榮耀裡.
和聖天使降臨的時候.
也要把那人當作可恥的.
耶穌又對他們說.
我實實在在告訴你們.
站在這裡的.
有人在末常死未以前.
必要看見上帝的國.
大有能力臨到.
最後這一節.
當然是講到耶穌.
祂自己的復活.
祂的復活更生了天地.
讓這個被死亡氣息.
籠罩的世界.
能夠窺見生命的光明.
在這裡耶穌說的.

$^{121}$和當日門徒的期望.
和今日信徒的期望.
落差都非常大.
當時給得很激動.
給得很激動.
他們對耶穌說.
喂!我們跟著你.
是因為你厲害.
你強.
你好打.
好像那些黑幫大佬那樣.
在這個講拳頭的世界.
能夠殺出一條血路出來.
誰知道你說什麼.
你說受苦.
被氣絕.
被死.
被殺.
喂!如果你真的有能力復活.
喂!拜託你.
不要死了之後才復活.
你在之前來顯出你這個大能.
這樣效果就更加厲害.
在十字架上.
你大搖大擺這樣走下來.
然後殺盡那些釘你十字架的羅馬兵.
哇!那個宣傳效果多震撼.
跟著你就是不想受苦.
今天.
信徒對耶穌的想像.
仍然希望祂是一個戰神.
無堅不摧.
無往不利.
病的時候我們希望耶穌醫治.
而不希望祂叫我們承受病的痛苦.
體會那些受苦的意義.
然後用這個體會去鼓勵其他的病人.
我們投資的時候.
我們是祈求耶穌祝福.
一帆風順一本萬利.

$^{161}$我們不希望耶穌提醒我們要知足.
要親手勞碌做一些有意義的事.
不少學生現在做功課.
是用那些人工智能.
叫人工智能把那些文章或者書本的徵要抽出來.
然後按照這些觀點.
輸出成為一篇論文.
學生的祈禱就是這篇論文.
哇!一定要通過大學的票舌的檢測系統.
千萬不要讓人知道我是抄回來的.
能夠拿高分.
我們從來都不希望耶穌吩咐我們.
就算時間多麼短促.
也都要一篇文章一篇文章來讀.
也不要計較分數的高低.
要憑自己的努力寫一篇論文出來.
我們不期望耶穌這樣提醒我們.
我們不希望耶穌說.
凡走捷徑救了自己分數的.
最後必定會喪掉這個學術的生命.
凡是放下拿分數的旁門左道.
最後才會救到自己的學術生命.
有些神學生畢業的時候.
找教會入職做一個傳道.
他們的祈求就是.
哇!人工高,福利好,責任少.
他們想的是他們心中渴望的教會.
而不是去思考.
哪一道才是最需要他們寫幾福事的教會.
不少人的人生規劃都很渴望.
蘇軾所說的就是.
無災無難到功興.
耶穌說的卻是要有天國的軟象.
為了讓上帝的恩典改變這個世界.
要背起十字架.
你要甘心多背負一些仁義的責任.
你要給其他人多作一點犧牲.
哇!耶穌所說的這些人生.
根本就是一盤蝕本的生意.
說這些話在今天很多信徒都會罵他.

$^{201}$不過我知道大家是不會罵耶穌的.
不過肯定你們會罵我.
我們看到耶穌背著十字架前行.
祂的苦難.
來自世界的人對祂的反對.
因為祂不滿足世界的人和一般信徒的期望.
反而祂想改變這個.
令大部分人都活得舒舒服服.
開開心心的世界.
俄國有一個大問號.
妥爺夫斯基.
名字都有點難讀.
他有一本巨著.
《卡拉馬佐夫兄弟》.
這本巨著裡面有一個常常被人引用的.
一個很出名的情節.
話說在十六世紀.
耶穌再回到這個世界.
祂說我要給人自由.
或者用現在的說法就是.
祂渴望帶領那些迷失了在這個世界.
種種虛假價值裡面的人.
幫他們重新找回生命的永恆意義.
十六世紀有宗教裁判所.
當時的宗教的大法官.
就是類似聖經的大祭司那種人物.
這個宗教的大法官對這位耶穌說.
你走吧.
這裡的信徒其實要的是麵包.
並不是你說的那些自由.
你說的那些生命的意義.
反而令到習慣了虛假價值的這些信徒混亂.
你這樣是會搞亂這個世界的.
你走吧.
如果你不走.
我還是會釘死你.
只要彼得和信徒要的是世界的價值.
就一定會和耶穌爭論衝突.
只要這個世界敬拜的是戰神.
而耶穌在努力背緊十字架.

$^{241}$這樣世界和耶穌就一定會衝突不休.
二十世紀在二次大戰期間.
在納粹的集中營.
住了一個囚犯.
一個反納粹的德國人.
一個德國的神學家.
他的名字叫做潘霍華.
在這個監倉裡面.
他天天在想幾個他覺得很關鍵的問題.
在監倉裡面他天天在想關乎將來的問題.
他在想戰後的德國怎麼重建的呢.
他在想基督信仰在戰後歷史前景的問題.
他在想主耶穌對這個世界.
還意味著什麼呢.
在他的監獄書信裡面他這樣寫.
在世上.
上帝是很軟弱的.
無能的.
正正就是這樣.
也唯獨是這樣.
祂才在我們的身邊幫助我們.
基督幫助我們.
不是靠祂的全能.
而是靠祂的軟弱和受苦.
在這個世上.
人類一般的宗教渴望.
就是指引人去追求上帝的大能.
一個懂得在困境裡面變法術的上帝.
但是聖經卻是指引人.
面向上帝的無能和受苦.
只有這位受苦的上帝能夠救我們.
潘佛華在想歷史前景的問題.
一個遲早都會戰敗的德國.
其實可以憑什麼來重建呢.
支撐著整個歐洲的文明的基督教.
還有沒有前途呢.
前途從何說起.
對著這個殘垣敗瓦的世界.
基督還意味著什麼呢.
人還會不會找基督呢.

$^{281}$戰火仍然在燃燒.
潘佛華在想.
就是人類的歷史要有前途.
就不能夠再敬拜戰神.
就是要跟從那一位.
背著十價前行的耶穌.
人人都敬拜戰神的話.
每個人都期望要比其他人厲害.
這樣人人都希望得到最大的利益.
如果是這樣的話.
就無可避免地.
大家都在叢林裡面混戰.
以力取勝.
以暴易暴.
在武力的輪迴裡面無限輪迴下去.
他殺了我的爸爸.
我就要上山學武.
有朝一日.
我下山要為父報仇.
終有一日你會殺了他.
然後他的兒子又會上山學武.
為的是有朝一日.
他都能夠報殺父之仇.
歷史就是無盡的後人.
去練武.
而練的武功越來越厲害.
殺人的能力也越來越強.
未來的歲月.
我們的歷史的將來.
核彈會越來越多.
射程會越來越遠.
殺傷力也會越來越大.
來到這樣的光景.
是因為無論國家 世人 信徒.
信拜 戰神.
對耶穌的排斥力度越來越大.
潘福華在想.
人類要有前途的話.
只能夠靠那位背著十字架前行的上帝.
這位上帝沒有用核彈.

$^{321}$去炸毀他看不順眼的世界.
他放下身段.
以自己的尊嚴.
與無能的人同在.
與受苦的人同行.
這位上帝以脆弱與無能.
去反抗這個暴力的世界.
去揭露神聖的天國.
就因為上帝甘願這樣受苦.
他就真的明白.
那些軟弱而受苦的人.
是完全明白他們.
那個好像離棄了世人的上帝.
他原來不但沒有離棄我們.
還在苦難的深淵.
緊緊地與我們結連.
正因為他在苦難裡面的同在.
我們才能夠在苦難的深處.
找到來自上帝的希望.
耶穌對每一個人說.
經過十架的黑夜之後.
他會見到天國的黎明.
如果他是一個戰神.
他當然會威風凜凜.
大有能力.
但是一個這樣的上帝.
他只能夠高高在上.
他沒有辦法親自體驗我們的苦難.
沒有辦法在苦難的深淵與我們同在.
一顆子彈能夠解決惡人.
但是這顆子彈.
不會為我們的痛苦流淚.
不能夠給苦難的我們一點同情,安慰.
也不能讓我們對天地間的人愛.
有多一點的信任.
唯有那位為我們擋子彈的耶穌.
那位承受苦難而與苦難的我們不離不棄的耶穌.
才能夠在我們的心靈深處.
新發出對人愛的信任.
把捉到來自上天的救贖.

$^{361}$燃點起一丁點的照亮前路的火光.
哪怕是一丁點的火光.
哪怕是猶豫不堪的人愛的幼苗.
也能夠給人勇氣.
讓我們放下武器而愛前行.
耶穌背著十字架.
緩緩的前行.
跟著祂.
反而能夠看到歷史的出路.
這些說法.
必然引來世人的嬌笑和謾罵.
在這個叢林的世界.
說的是現實政治.
人們不拿起武器.
下一刻死的就是你了.
現實真是這樣的.
畢竟耶穌真是死了.
但是這樣是不是證明耶穌不對呢.
或者大家倒轉過來想一想.
如果沒有了耶穌.
這個這樣的叢林世界.
還有沒有希望呢.
或遲或早.
你總是要抉擇的.
在這個叢林世界裡面.
以力去決戰.
然後就看看誰能夠留到最後.
你這樣選擇.
還是你甘願吃一些虧.
放下一些尊嚴.
多承受一些責任.
能夠終止這個暴力的循環.
聖樂家潘霍說.
只有這位受苦的上帝能夠救我們.
耶穌毫不含糊.
也絕無妥協.
要求祂的門徒捨己.
背起十字架來跟從我.
或者是時候我們每一個都檢視一下.
整個的人生.

$^{401}$十字架跟世界的人追求的價值.
和目標是相反的.
世人的目標是權力和名利.
採取的手段是暴力.
十字架卻是放下身段.
與悲下的人同在.
與受苦的人同行.
堅拒以這個世界的價值.
作為人生的目標.
當時的彼得.
拒絕耶穌的這種想法.
以耶穌在祂面前強大的堅持.
他說這條才是正路.
最後讓我們靜下來.
再次聽一聽耶穌的呼喚.
馬後福音第8章第34第35節.
耶穌說若有人要跟從我.
就當捨己.
背起十字架來跟從我.
因為凡要救自己生命的.
必喪掉生命.
凡為我和福音喪掉生命的.
必救了生命.
願神賜福給你們.
也願榮耀歸於聖父聖子.
主席.
\newpage



\section{彼得後書 1:5-8}
\label{sec:031993}
\textbf{建立健康的使命教會 (12) 勤成長的靈命}~~~~吳磐德牧師
\newline
\newline
連結: \href{http://www.church.com.hk/acms/content.asp?site=cdc\&op=show\&type=product\&code=031993\&layout=sermon}{\texttt{031993}} ~~~~ 語音日期: 2026-04-19
\newline
\newline
\hyperref[sec:031992]{\small{< < < PREV SERMON < < <}}
~
\hyperref[sec:index_chronic]{\small{[返順時目]}}
~
\hyperref[sec:031994]{\small{> > > NEXT SERMON > > >}}
\newline
\newline
彼得後書 1:5-8
\newline
\begin{longtable}{cl}
\hline
\hline
章節 & 經文 (和合本修訂版)\\
\hline
1:5 & \begin{tabularx}{0.7\textwidth}{X} 正因這緣故，你們要分外地努力。有了信心，又要加上德行；有了德行，又要加上知識； \end{tabularx} \\ \\ \relax
1:6 & \begin{tabularx}{0.7\textwidth}{X} 有了知識，又要加上節制；有了節制，又要加上忍耐；有了忍耐，又要加上虔敬； \end{tabularx} \\ \\ \relax
1:7 & \begin{tabularx}{0.7\textwidth}{X} 有了虔敬，又要加上愛弟兄的心；有了愛弟兄的心，又要加上愛眾人的心。 \end{tabularx} \\ \\ \relax
1:8 & \begin{tabularx}{0.7\textwidth}{X} 你們有了這幾樣，再繼續增長，就必使你們在認識我們的主耶穌基督上，不至於懶散和不結果子了。 \end{tabularx} \\ \\
[1ex]
\hline
\hline
\end{longtable}
$^{1}$各位弟兄姊妹早晨.
在我們教會的廣告系列裡面.
不知道大家是否記得.
我們之前有三個主題.
第一個主題就是.
活出敬拜的人生.
敬拜人生其實最重要的關鍵就是.
究竟我們有沒有去回應神.
信服神.
其實我們在神的面前.
每一個的毀身.
每一個的犧牲.
每一個的獻上.
那個都是一個的敬拜.
所以不只是我們在這裡去敬拜神.
大家記得我們在這裡.
神是看著我們.
神是唯一的觀眾.
我們每一個都是參與者.
是一個演員.
究竟我們的心靈是怎樣去敬拜神呢.
神就看著.
所以讓我們一起都是.
這樣去過我們的生命.
另外.
其實我們如果要過一個.
毀身的人生不容易的.
因為我們信服神.
很多時候我們自己有很多掙扎.
所以要有質素的相交.
而這個相交的生活.
不是一個只是社交的.
好朋友的關係.
而是有守望的.
是有愛心說誠實話的.
是能夠互相的去鼓勵成長的.
激發開心免禮行善.
所以這些有質素的相交.
可以幫助到我們.
內在生命的成長.

$^{41}$能夠活出到這個敬拜的人生.
我們其實在過去這四個星期.
今天是最後一講關於.
這個勤成長的靈命.
其實勤成長的靈命.
和之前那兩個的主題都很相似.
因為你的靈命成長的目的.
是要像基督.
你的靈命成長的目的.
就是要毀身.
和降服在神的面前.
去回應祂.
但是我們要回應祂的同時.
就是能夠互相的守望.
幫助對方去這樣回應神.
我們今天從後書的第一章裡面.
更加清楚知道.
究竟我們怎樣為一個勤成長的靈命.
那個圖畫是怎樣的呢.
你又可以怎樣去投入參與呢.
或者我們一起進入聖經話語之前.
我們同心低頭祈禱.
聖靈我們求你帶領我們.
求聖靈打開我們的心門.
取去我們的錫心.
賜給我們肉心.
以致我們的心靈是柔軟的.
我們是謙卑在神的面前.
降服的.
以致你的話語就好像兩人的利劍一樣.
刺透我們的心.
更新我們的生命.
讓我們過討你喜悅的生活.
不斷的在主理的裡面.
在你的帶領底下.
我們去成長.
我們將以下言錄你話語的時間.
交給天父求聖靈親自的.
向我們每一個人說話.
祈禱感恩.

$^{81}$奉著宋明裘.
我不知道如果我問你.
你現在的靈命狀況是怎樣呢.
你會怎樣回答.
有時候我們每年都去身體檢查.
我們看看究竟有沒有三高.
或者現在的生命狀況是怎樣.
但我不知道你有沒有常常.
或者恆常去做你自己靈命的檢查.
如果我問你這個問題.
不知道你怎樣回答.
我沒有答案.
只能你自己填個洞.
我的靈命都不錯.
我也有靈修.
我也有回教會.
你會怎樣回答.
我給你十五秒時間去想一想.
這三個不同的顏色.
可能代表你不同的人.
其實我都挺好的.
我又稍一奉獻.
我又有事奉.
我是不是正在成長呢.
我自己覺得有三個歸類.
當然包括不了所有.
但總結了我自己過去幾十年.
跟隨神的一個循環.
不是因為我做了牧師.
就會神聖了.
一樣牧師都會靈命停滯.
所以你也會.
所以我形容有些階段.
會是靈命遊蕩.
你知道遊蕩是什麼嗎.
我們現代人特別.
不是很給自己時間放空.
遊蕩.
可能子女會遊蕩.
我們自己未必去.

$^{121}$遊蕩是什麼呢.
就是萬無目的.
去一個地方.
總之有事做.
不要坐在家裡那麼悶.
總之就渡過一天.
靈命有時都會是這樣.
有時靈修有時不靈修.
有時祈禱有時不祈禱.
或者最近看靈修材料很悶.
不如看看別的樣子.
跳來跳去.
看來看去都覺得很悶.
或者很乾.
那是靈命的遊蕩.
好像沒什麼突破.
或者經歷神.
又或者是靈命停頓.
我想我自己過去一些.
這麼多年都曾經試過.
停頓的時候是怎樣.
我不知道你的停頓是怎樣.
停滯的時候是滯留.
新冠肺炎大家都經過滯留.
哪裡不用去.
滯留在家裡.
滯留在一個地方.
尤其是如果你回到香港.
忽然間又說不准走.
你就滯留在香港.
其實我們靈命都會滯留.
為什麼.
因為有時我們會覺得.
我是聽的聽得到.
其實我今天.
牧師講的題目或者經文.
我讀出來你已經知道我會講什麼.
已經知就知過.
聽就聽過.
視乎我什麼崗位.

$^{161}$牧師問.
我什麼崗位從倒垃圾到港島都做過了.
我整個生命這麼多年.
跟隨主很多年.
我就已經什麼都試過.
什麼都做過.
有什麼沒見過.
不好意思.
其實這是一個靈命停頓的階段.
因為沒有什麼新鮮感.
沒有什麼引起你那種激情.
是不是.
覺得一件事都試過了.
都聽過了.
不過就習慣了這樣.
當然要來崇拜.
要去主教學.
要去小組關了.
其實就停滯了某個階段.
雖然有很多活動.
但那些活動並不是我們靈命的成長.
所以我們說我們要追求靈命成長.
但是這個人生的靈命的周期.
往往你要問自己.
誠實地答這個問題.
究竟這個周期你在哪個地方.
今天.
你不要給準答案.
如果你要出聲要舉手.
當然每個人都舉靈命成長.
但是你自己問自己.
現在你真是在遊蕩的階段.
靈修完又好像不覺得自己有什麼特別.
最怕在小組裡面就問了.
上個禮拜有什麼經歷神.
前個禮拜有.
今個禮拜沒有.
行不行.
或者有些新的事奉的機會.
但又覺得真的有很多難搞的.

$^{201}$已經有很多老油條的言論.
我說試過了.
我就知道怎樣不行.
怎樣一定可以.
你說些什麼.
試過了.
究竟你在哪個階段.
如果你是靈命遊蕩的階段.
你現在是怎樣遊蕩呢.
又或者你是靈命停滯的階段.
你滯留在那裡多久呢.
我想知道.
有沒有一些事刺激到你興奮一些.
在屬靈生命上面.
還是現在你是一個靈命成長的階段.
每一天都是新鮮的.
因為神告訴我們.
每一天都是新的.
我們的經歷.
我們的信心.
我們的愛心.
所有東西都是新的體會和經歷.
新的機會我們又會去掌握.
因為我們靈命是不斷地被塑造.
不斷成長.
聖經裡面.
希伯來書第六章一節.
大家很熟悉的.
怎樣說.
所以我們應當離開基督道理的基礎.
從這個《易本》講開端.
竭力進到成熟的地步.
在這次經文裡面.
在原文裡面講基礎是講起初.
開端或者起點.
在這個起點裡面是包括什麼呢.
就可能我們認識悔改.
認識信心.
認識朕禮.
認識復活.

$^{241}$認識將來的審判.
這些基本的基礎的信念和信仰.
我們都懂得.
首先我們有事奉的.
這些是基礎入門的道理.
但是留意經文不是叫我們停留在那裡.
而是要離開.
而是要竭力.
竭力就不是.
靈命成長不是自己會發生.
自然會發生.
什麼都不用做就發生.
而是需要一些計劃.
有些是刻意的.
所以我想問大家的.
就是你靈命成長的下一站要去到哪裡呢.
可能你會覺得.
我也不知道下一站.
不如更長基督吧.
那是什麼意思呢.
更長基督的哪一部分呢.
哪一方面呢.
我們自己有沒有一個心目中的理想目標.
或者我們下一步怎樣.
其實我們每個人都不同.
可能你正在掙扎的是一個關係.
可能你正在掙扎的是.
入籍的問題.
等一等都沒有收到入籍通知.
或者居留通知.
以至於在鍛鍊,磨練,折騰你的信心.
每一個人的狀況我們都不同.
可能是生病.
可能是我們的信心軟弱.
各種.
但是對你來說.
靈命的下一站又要去到哪裡呢.
如果我們不清晰的時候.
如果我們沒有一個渴求的時候.
就只是做了崇拜主義學小組.

$^{281}$日日的靈修.
可能我們都在遊蕩著.
可能我們在停滯著.
我們都不知道自己要去到哪裡.
所以今天彼得後書.
第一章幫助我們.
就是怎樣從信心開始.
以至到我們要去到一步一步的成長.
以至到能夠結出生命的果子.
整個勤成長的靈命.
一幅圖畫.
今天在彼得後書裡面.
就展現在我和你的眼前.
我們第一句話.
在彼得後書第一章五節已經說了.
正因者緣故.
所以我們就要努力去成長.
就好像希伯來書六章一字一樣.
竭力成長.
這裡就說要份外努力.
在希伯來書就說竭力成長.
都是有意識的.
都是刻意的.
都是用力的.
但是因為什麼緣故呢.
在彼得後書的前兩節.
第一章第三第四節就這樣說.
我們一起讀.
我們一起讀這段經文.
神的神能.
已把一切關乎生命和乾淨的事.
賜給我們.
因我們認識那用自己榮耀.
和美得照我們的神.
因此祂已把寶貴又極大的應許賜給我們.
使我們既脫離世上從情慾來的敗壞.
就得分享神的本性.
這次是讚美主.
你見不見這段經文說.
有一件事你不要忽略.

$^{321}$賜給你神的本性.
是不是.
開心嗎.
賜給你了.
好像中了樂彩一樣.
賜了還得到了.
所以這裡特別黃色的字說什麼.
神的神能已把一切.
不是一點.
是一切.
關乎生命和乾淨.
即是關於神的事.
賜給我們.
兩次說賜給我們.
賜給我們應許.
賜給我們極大寶貴的應許.
是什麼呢.
就是脫離了我們的漏腦.
以致我們能夠分享神的本性.
於是你說.
好像經文本身是相矛盾.
賜了給我們就得到了.
等於你中了樂彩.
拿獎拿回家.
為什麼又好像拿得了又拿不得.
又要我們努力拿.
的確要努力拿.
因為就好像.
在我們科恩夫裡面說.
耶穌叫彼得把網子放進去釣魚.
但那些魚已經在水裡.
你也要拉上來.
所以很關鍵就是我們明白.
什麼是分外努力.
這裡是在說靠我們自己的努力.
還是神給我們的.
在原文裡面.
這個字真的很重要.
因為幫助我們明白.
其實那個分外地努力.

$^{361}$是在說讓自己持續的.
被帶進去成熟.
換句話說.
你已經得到神的本性.
所以你已經有這個潛質.
但是你要竭力去配合神.
在你生命裡面的工作.
以至於在配合的時候.
可以得著.
令到你的潛質.
可以被彰顯.
顯露.
發揮出來.
這個關鍵就在這裡.
我們要配合神.
這兩段經文.
希伯來書第六章一節.
和彼得後書一章五節.
關鍵就是告訴我們有一種張力.
就是人的回應.
就是離開基礎.
基督道理的基礎.
這個是我們人要做的.
你需要有這個決心.
離開不是說.
去放棄那些基礎.
離開是叫你不要停留在那裡.
而是要開始走新的一步.
去有新的生命的經歷.
你不要停留.
這個是我們首先要在腦袋裡決定.
然後我們有這個決心去做.
然後.
在彼得後書第一章五節講到.
奮愛地努力就是神的工作.
我們是被動的.
是讓神帶我們去到成熟.
所以不只是靠我們自己的努力.
但是我們要努力是讓神帶我們.
那你說什麼意思呢.

$^{401}$其實.
或者說回建堂.
用建堂這個例子最容易就是.
去年可能我們看到.
五百萬要籌一年才籌五百萬.
沒可能的.
你和我都覺得很難.
你問我真實的答案都覺得很難.
但是.
神就是在那一刻.
如果我們不配合神的工作.
就放棄吧.
大過我們.
其實我記得十年前.
Spencer跟我說.
你記不記得我們回香港.
第一個籌款目標是五百萬.
就應該覆蓋我們整個建堂目標.
我多開心啊.
香港回來.
已經籌完建堂的錢.
當然不關Spencer的事.
只是告訴大家.
那是第一期的錢.
沒說什麼是全部的錢.
但是.
無論如何.
但是在那一刻.
你可以決定說我不做了.
就不配合神了.
但是如果你期圓土.
我們共同一起在我們的月會裡面.
去決定這件事.
認定是神的帶領的時候.
我們就走.
那我們就會一步一步.
見到神的帶領.
我們由沒有信心變成有信心.
我們由埋怨到不埋怨.
我們由.

$^{441}$其實這件事是不行的.
說這麼容易嗎.
說這麼玩嗎.
到.
咦 真的.
我不敢說了.
又近一點.
什麼時候得到.
又好像我還沒做事.
但是英神在做事.
我正在做事.
神又幫我雙倍做事.
總之就是.
我不做事又會發生.
我做事又更多的發生.
就是神究竟怎麼做.
透過什麼人做.
我全部不知道.
但是又增加了我們的信心.
是不是.
所以我們是給神帶領.
就是你回應神的時候.
信服祂.
就是你還沒看到的事.
這些不只是建堂的事.
我想用這個例子.
不過容易明白.
在你自己的生命裡面都是.
無論夫婦關係.
你自己儲錢買房子.
無論入籍等候.
等極了.
加拿大政府好像將那個門檻.
或者將那個龍門又移遠一點.
但是.
只要你知道.
神帶你來這裡的.
你又信得過祂.
跟第一次一齊同行的時候.
你那個心態.

$^{481}$你的回應就很不同.
神就可以透過這些事情.
大家記不記得.
在我們過去這三個禮拜.
說關於勤誠長的靈命第一.
要愛慕神的話語.
首先要有神的話語.
是不是.
然後第二堂.
我跟大家說的就是.
在我們生活裡面.
所有的事情.
萬事都互相效力.
就是我們所遇到的.
所有的事情.
在生活裡面.
神的話語.
透過我們生活的事情.
就能夠塑造我們的生命.
所以我們要跟神配合.
然後上一堂.
江牧師就說.
你要懂得辨識.
為什麼.
因為神的話語在這裡.
你也跟生活跟不上.
又或者你生活上.
需要神的話語.
怎樣去實踐.
能夠幫助你的問題.
你又看不到那個彼此的關係.
又不願意信服的時候.
你辨識不到.
這個是神已經在跟你說話.
那個就是問題.
所以勤誠長的靈命.
就是當你有神的話語.
你看到你的遭遇.
而神的話語幫到你.
解決你面對的問題.

$^{521}$和看到神在你生命裡面.
在做什麼工作.
你就在靈命的成長裡面.
所以你需要配合神的.
如果在每一個神給你的機會.
和生命的經歷裡面.
你都絕不拒絕神.
不信服祂.
只是依靠自己.
或者自己的努力.
難怪我們沒有辦法成長.
所以是需要努力地依靠神.
你說依靠神的努力.
不用做的.
當然不是.
要做事的.
依靠神去做事.
以致你見到神.
在你透過你的生命去成就奇妙的事.
所以要竭力.
去信服去回應.
所以這兩段經文幫助我們.
就是說.
我們靈命的成長.
是要不停步的.
因為離開道理的開端.
就是不要停下來.
而是開始起步.
你起步了沒有.
而那個份外的努力.
就是我們要順著神的帶領.
讓祂在你的生命當中工作.
我們是要和祂合作的.
以致我們能夠殷勤的.
就是勤力的.
去回應神的恩典.
神的恩典就是.
祂已經和你分享了祂的神性.
祂已經給了你恩典.
救贖的恩典.

$^{561}$神本性的恩典.
你的責任就是去發掘.
去開放.
去配合神.
在你生命當中的工作.
以致能夠彰顯神的榮耀.
在你的生命當中.
所以我問你.
份外努力.
你有沒有完全的降服給神.
讓祂帶領你.
憑信心前行.
是有計劃的.
不是等運到的.
是有計劃是刻意的.
而是去到成熟的地步.
我自己在過去十五年.
我都覺得.
十五年前.
我已經事奉了神十年.
就是全職事奉.
但是從小開始事奉.
就是三十多年了.
因為少年團已經開始事奉了.
但是到了那個階段.
就是不太懂得勝靈.
在心靈裡面說的話.
是知道聖靈有人的印證.
但是經常聽到.
究竟怎樣聆聽聖靈呢.
怎樣去安靜呢.
沒有聲音也不只是安靜.
安靜不只是沒有聲音.
而是你自己心靈裡的聲音都要安靜.
聽起來真的好像很炎妙.
都不知道這麼虛無.
但這就是我要黎明成長的地方.
我記得十年前.
我在天竺遊讀神學博士的時候.
我其中一個同學說是牧師.

$^{601}$那個先生叫他.
我們每個禮拜上五堂.
每堂頭半小時安靜默讀.
他多生氣.
他說我這麼多學費來天竺.
你讓我坐在這裡聽.
你沒有東西聽.
沒有聲音.
但是就算牧師都會成長的.
他沒有試過.
他只懂得外在的靈修.
總之就是多看一些資訊.
多看一些聖經.
一年就要看三次整本聖經.
塞到自己都聽不清.
那些資訊不是不容易.
我絕對不是開玩笑說的那些.
其實看聖經.
主要是學我們明白聖經.
識經是重要.
但不是唯一.
對我來說.
那些識經.
明白聖經.
明白了.
夠不夠呢.
不夠.
就要懂得如何去領受.
懂得如何去聆聽聖靈.
懂得如何在生命裡面.
更深層的領受.
所以才會去到前年.
真是四十天的安靜靈修裡面.
很多朋友都說.
四十天行不行.
原來四十天很快過的.
不是因為風涼話.
而是.
如果你再給我四十天.
我都願意的.

$^{641}$不是因為四十天沒什麼事做.
而是.
四十天沒什麼事做其實很難的.
你試一下.
兩天不工作.
在家裡.
但如果你已經有了這個操練的時候.
你很嚮往.
但我的成長就是要突破.
我過去要忙.
要做.
而是能夠完全放下.
經歷神.
這個就是我要成長的地方.
所以每個人都有.
你的成長是下一步.
你的成長下一步是哪一種.
而你有沒有停滯不前.
不要停步.
弟兄姊妹.
每一個人的答案.
這條問題都不一樣.
接著經文裡面又很精彩的說.
有了信心又要加上德行.
有了德行又要加上知識.
有了知識又要加上節制.
有了節制又要加上忍耐.
有了忍耐又要加上獻勤.
有了勤勁的心.
又要加上愛弟兄心.
又要愛弟兄心.
又要加上愛眾人的心.
真是很多東西加.
所以就很怕這些.
信耶穌有這麼多事情要做.
不是.
又是.
的確有很多事情要做.
不過不是經常你做.
我剛剛才說.

$^{681}$你被神帶著做就很容易.
好像浮在水面.
抽你上來.
但是你要合作才行.
不然神在抽你上來.
你會不斷游不斷亂搞.
所以這段經文.
雖然很多東西要加.
有很多東西.
有八樣屬靈的特質.
但是其實.
他是在說.
你看到.
由信心開始.
以愛去結束.
靈命成長.
真正的屬靈的成熟.
必須要由信心開始.
我們相信神.
信服神.
然後經歷過中間.
然後有愛.
就是你會去愛人愛神.
是會有愛在的行為.
所以這是一個進程.
我們可以看得到.
第二就是.
第二方面我們就會見到.
就是.
有了什麼又要加上什麼.
這段經文裡面.
最多的「有了」和「又要加上」.
就是整段經文.
如果少了那兩大字.
你會看到只剩下兩行.
短了很多.
但是為什麼要這麼長氣呢.
為什麼不說.
有了信心.
得行知識.

$^{721}$節制忍耐.
先說完這些.
然後才說又要加上什麼.
很快.
減少很多字.
沒有效率.
正正因為.
正經的作者.
想我們明白一點.
就是.
要加上.
其實不是說一步一步來的.
而是原文那裡要加上.
是說再豐富一點.
就是有了這些.
你又會令到另外一些東西.
會令你更豐富.
有了這些你又要更豐富.
換句話來說.
是一個不斷的.
不是說有做了信心.
還有第二步第三步.
不是說一步一步來的.
而是一個整全的生命圖畫給你看見.
就是不斷添加的.
持續的成長的.
是每一部分.
整全生命的一個進程.
所以不斷的神是會供應你.
而你是要累積一直成長的.
是同時發展的.
是彼此影響的.
這些元素並不是說有一才二.
而是同時間進行.
是彼此配搭.
互相的.
是幫助我們有一個整全的建立.
你的生命有沒有越來越完整.
這就是這個經文裡面想提醒你的.
你已經滿足於那些已經一九的東西.

$^{761}$閉家火.
那就停滯了.
滯留在那裡了.
有了你都要再想怎麼添加.
添加什麼.
然後我們就不會遊蕩.
因為我們就要知道自己有什麼.
然後知道我們屬靈的下一步是什麼.
於是我們就會知道.
我們要去到什麼地方.
讓神開我們的眼睛.
這段經文也很豐富的告訴我們.
就是屬靈生命裡面的這八個元素.
那個進程包括信心是根基.
然後藍色的德行知識節制忍耐.
就是塑造我們的品格.
我們的character.
我們的屬靈品格.
讓我們有這些品格去見證神.
人家看到我們的生命就看到基督.
然後乾淨.
第三方面就是生命對準神.
乾淨的事就是我們要做乾淨的人.
是對準神的生命.
最後就是愛弟兄和愛眾人.
然後就是愛展.
愛是流向其他人的生命當中.
所以你會看到一幅什麼圖畫.
你會看到一幅整存的生命成長.
勤成長的靈命的圖畫.
就是建立一個由內而外的屬靈生命.
我很大力推薦大家.
其實過去這幾年.
我們很開心有宗旨超傳道.
幫我們創建了一個靈修材料.
就是靈穴的靈修材料.
三月底我們會暫時過一段落.
因為已經存了幾年.
如果大家沒有用過.
請大家打進去.

$^{801}$可以看到以前這麼多年的.
但我很想幫助大家進入.
在我們的手機程式裡面.
我有一個新的推介.
幫助大家真的一步一步.
能夠照顧自己內在生命的屬靈生命.
為什麼.
其實有時候很多弟兄姊妹說.
我聽宗旨超傳道靈穴.
聽著聽著睡著了.
其實睡覺是很正常的.
我聽也會睡覺.
不是因為他說會睡覺.
而是因為如果你不懂得去默想.
很快你就會睡覺.
或者你太累.
屬靈很疲倦.
屬靈裡面也很疲倦.
你又會睡覺.
所以其實要睡醒才聽.
不要叫你不睡才聽.
所以你要學習靜.
這個新的資源是十分之好的.
是香港一個很受尊敬的.
已故的譚佩泉牧師.
他就做了這個網頁.
他這個裡面.
如果你進去我們的app.
你打進去今日靈修.
你就會見到一個新的圖畫.
那個就叫做靜觀心土.
按下去你就會見到那個app.
你就下載免費的.
你下載之後.
你就會見到三個材料.
第一個材料最划算的叫做.
心靈心語.
太小的字.
心靈心語.
那個是祈禱的.

$^{841}$其實你不要想著裡面有很多東西.
它就是教你祈禱.
每天八分鐘至十分鐘.
學會怎樣安靜祈禱.
慢慢祈禱.
如果你嫌它慢.
就是說你太快了.
如果第二天你又聽到同一個祈禱.
就是說你太沒耐性了.
因為它就是要用半年的時間.
來幫你重複聽完.
你慢慢就由很急到慢.
由不懂得慢到真的懂得慢.
所以如果我們用我們的速度來聽.
就覺得又第二天講.
如果三天都是講同一樣東西.
不如我按就按第四天.
那你就全錯了.
難怪你硬是弄來弄去都是聽.
寧願都是睡覺.
因為你從來都沒有靜過.
你從來都沒有懂得怎樣安靜自己的心靈.
所以我就幫大家錄了一個介紹.
推介的短片.
二十二分鐘的.
但是你嫌我長氣.
你就跳過了頭十一分鐘.
頭十一分鐘之後.
就教你怎樣用那個APP.
所以為什麼要講二十二分鐘呢.
除了我平時都長氣之外.
就幫助你可以有這個忍耐.
以及可以進入這個慢的階段.
所以你一天.
所以你問自己.
香港靈命城的下一步是什麼.
如果你覺得自己什麼都懂.
請你入定修材料.
那你會發覺你什麼都要重新學.
因為其實你真的沒有經歷過.

$^{881}$半年之後我擔保你.
因為它第一個就是學祈禱.
第二個就是靜心靜肚.
就是CP 139篇.
一篇的CP用半年時間看.
真的逐個字逐個句去默想.
有很長一段時間.
慢慢幫你安靜.
然後呢.
第三個重土深谷就是三百六十五天體意忽所書.
所以我介紹大家.
如果你進去看這個三百六十五天的時候.
頭十四天都還沒進正題.
那你就會發覺你太急了.
你太急.
你不懂怎樣去重複咀嚼之後.
原來神要跟你講一些你過去沒聽過的東西.
你明白嗎.
好 那我希望這句真是鼓勵大家.
這個是重點極力推介你靈命成長.
內心生命裡面的操練進路.
我不叫你聽完所有.
因為這個八分鐘有三個不同的.
如果你已經比較熟悉默想.
你直接去三百六十五天那個意忽所書.
但如果你根本沒試過.
請你去祈禱那個.
那你慢慢學安靜學祈禱.
然後才進入 所以這是三個不同的.
那你是每天八分鐘十分鐘.
但你可以繼續用你現在用開的靈修材料.
沒問題的.
但你可以另外一個時間八分鐘做這個.
那我又回到正題.
因為這個推介都是正題.
不過不是這個經文裡面特別去推介這個apps.
但是這個apps可以幫助我們去實踐.
這段經文裡面所講的東西.
就是回到基本層.
原來我們靈命成長不是亂走的.

$^{921}$不是遊蕩的.
是要知道你去哪裡的.
是想去哪裡的 追求的.
所以第一個根基層.
就是同神是有一個關係的起點.
信心就是一個起點.
但是信耶穌.
不只是由你接受耶穌那一刻開始.
做了那個所謂的確句的祈禱.
我們都有個conversion prayer.
確句的祈禱.
那個只是一個起點.
但是你其實信心需要建立信靠和衛生.
不是頭腦裡面的認知.
而是生命的依靠.
就是說我們知道神是要我們建堂.
我們知道神是會供應.
但是就是我們不信.
但是我們就是不依靠.
還是靠著自己的努力.
那個依靠就是要學習.
由依靠自己到依靠神.
這個都是在第一層基本的.
所以我們過去一年恭喜大家.
因為大家一起在建堂的時候.
在學習那個根基層.
因為是一個信心的功課.
是一個依靠的功課.
是一個信服的功課.
不只是在建堂.
就在你生命裡面的任何事情一樣.
然後第二層就去到品格層.
再一次提醒不是一層一層去到第二層.
這四層都是一個互動的過程進程.
無論信了主多少年你的靈情在哪裡.
這四方面其實我想大家留意.
這個品格層就是靈命成長的核心操練.
所以我們要靈命成長.
就要操練這四方面是重要的.
因為這個內在生命的塑造.

$^{961}$德行是什麼.
是卓越.
就是在神的裡面卓越.
不是我們在人的面前卓越.
就是做好神的心意的東西.
就是那個卓越.
是道德的勇氣是正直的.
就是說神是正直的.
神是有道德的.
神想什麼.
合神心意的生活和要求.
和那些原則.
就是你有這個德行.
就是勇氣去選擇這個去做.
第二方面就是知識.
明白了對.
明白什麼是對的.
這個知識就是屬靈的洞察力.
和能夠明辨.
聖經的話語在這些事情上.
在你的生活當中是怎樣意思.
那個是知識.
第三方面就是節制.
節制需要約束自己.
有一個自律的能力.
因為你有了德行知道對的事.
你又辨識到.
你有這個勇氣去做.
但是你也有知識知道.
那樣東西是對的.
是神的心意.
但是你沒有約束自己.
去做也不行.
很感恩的.
其實很多弟兄姊妹約束自己.
去少一次旅行.
捐了建堂.
我鼓勵大家.
其實你是已經在做.
我知道很多弟兄姊妹.

$^{1001}$用你各種不同的方式.
按著神感動你的方式.
你就約束過一些事情.
以前可能你隨便一點花錢.
現在小心一點花錢.
這個也是一個約束.
然後忍耐就是持續的堅持.
信得過神是會帶領.
你堅持而不放棄.
那種持續.
其實在生命裡面.
任何東西.
你每天的靈修.
如果你在這四個導向裡面.
去問自己四個問題.
你的靈命必定會成長.
因為德行就是你每天問自己.
是的 現在我靈修.
聽到神的聲音了.
神的話語了.
你有沒有這份勇氣活出來呢.
可能是和社會的文化價值觀不同.
和你家人的意見都不同.
但是就是神這樣和你說.
你有沒有這個勇氣活出來.
你又或者會不會明白到.
神的心意是怎樣的呢.
在你的生命裡面.
怎樣回應這件事情呢.
在夫婦關係上.
在入籍的事情上.
在工作的事情上.
事業的事情上.
都是究竟聖經有什麼指引呢.
然後既然明白了神的指引.
又知道你需要這份勇氣的時候.
你需要有一個決心去約束自己.
以至有一個自律去完成.
其實然後最後才是不斷的.
在壓力裡面仍然持守.

$^{1041}$這四方面你覺得很難的時候.
我有一個鼓勵給大家.
就是你不要忘記這段經文.
第五節就說.
其實神已經將他的神的本性.
賜給你了.
所以不是說基督徒有很多事情要做.
不是的.
神已經賜給你了.
不用做都賜了.
你要做的唯一事情.
就是配合神的工作.
會應神.
配合他.
以至於你能夠被帶進新的境地.
其實我自己最體會這個.
就是我十五六年前.
大家都不會想.
不是十五六年前.
是二十年前.
二十年前我大約接近二百磅.
現在我大約一百六十磅.
我當時在加拿大讀完神學.
回香港事奉這樣的身形.
在那時候四十歲左右.
就真的要去做運動.
加拿大有一個好處.
就是因為它曾經舉辦過冬季奧運會.
所以它有一些很便宜價的奧運運動員.
二十五元就可以做一個小時.
這個個人教練.
真的很划算.
所以教你鍛鍊身體.
你真的要配合他.
因為你不配合他也沒辦法.
很多人都說.
教練有什麼用.
但我是那些很乖的.
他說什麼我就做什麼.
然後真是死五仔.

$^{1081}$他半個小時跟著我.
我已經回家躺著動也不行.
第二天痛得要命.
然後又要去再去.
想著再去.
是不是真的要去.
但既然自己願意就去.
讓他折磨一下.
其實你要配合他.
如果你不配合他.
他多厲害的教練也沒用.
其實回到加拿大.
於是我二千年就回到香港事奉.
十年之後回到加拿大.
二零一零年重大劑.
因為那時候瘦了.
回到香港服事.
服了十年又胖了.
於是又回來又去減肥.
因為脂肪乾了.
真的很有動機.
因為有脂肪乾再不然就不太好了.
所以就真的有減肥.
家庭醫生在加拿大真的很好.
介紹你去營養師.
不要說哪個名字.
因為賣廣告的.
跟那個營養師.
又不用錢.
因為我有脂肪乾.
家庭醫生寫紙.
你就有拿證牌.
每兩個星期三個星期就去見營養師.
他真的很好做.
他問政府拿錢很容易.
坐在前面幫我磅一磅.
多重.
然後就說在這裡吃單.
你不要吃什麼.
不要吃什麼.

$^{1121}$說完之後.
三個星期回來.
又是跟我說話.
其實跟靈月差不多.
不是跟靈月.
跟淨土差不多.
其實都是說多一點.
所以我很有經歷.
所以神是預備了我去做靈修.
因為你聽完就是.
你廢話.
說完.
等於沒說我就知道.
有很多人說你教他都失敗.
真的.
有很多人說去醫生那個營養師不行.
當然不行.
你瘦了五磅之後.
馬上吃十磅東西.
所以你要乖.
就需要配合.
我是很配合的.
所以當時師母樂翠常常跟我說.
她吃東西很難吃.
其實我覺得沒問題.
你吃什麼都無所謂.
總之就是配合那個營養師.
我真是如是者.
減到我自己很有理想的磅數.
所以寧願成長.
同樣就是.
你有勇氣去做決定.
你要做的事之後.
知道是神的心意.
好像健康就是你的目標.
你就可以對那些誘惑說不.
向耶穌或者教師告訴你.
說是.
然後堅持的時候.
你就會達到那個目標.

$^{1161}$這個就是我的經歷.
第三層就是神向層.
就是乾敬.
是敬畏神.
以神為中心.
這個其實就是敬拜人生.
第四方面才是關係層.
所以我們將整個生命對準神.
就是敬拜人生.
而我們最後一層.
才去到那個關係層.
就是生命流露出來.
是愛弟兄.
就是先在教會.
他說愛弟兄.
在教會內的人.
我們有團契的愛.
就是那個淡成相交.
互相守望的愛.
然後在教會外的.
就是愛眾人.
所有的人.
包括不可愛的人.
無條件犧牲的愛.
對外面的人.
所以其實這段經文要求這麼多.
你簡單去理解的時候.
它有四層.
而每一層一個圈.
其實是互動的.
是一起是一個整合.
最關鍵就是.
我們要明白.
它是由內而外的.
是先建立你的生命.
然後再有關係.
然後才去影響別人.
如果你吊亂這個基礎的時候.
你還在影響別人.
哪怕你有多少的事奉.

$^{1201}$但是你的生命.
其實是不能夠勝在的話.
都是很危險的.
所以要刻意的整存.
建立靈命.
是要配合神.
在你生命裡面的工作.
就是去聆聽神.
透過聖經的話語.
藉著聖靈.
回應祂給你遭遇的事情.
以致你能夠作合適的決定.
去回應祂.
而你就會不斷的成長.
是有方向的.
有進路的.
你自己是不是一個.
這樣的持續的成長.
以下我很想找一位弟兄.
馮其進弟兄.
他是去年在十二月.
聖誕節的時候.
才受浸.
所以每一個人.
都有他不同的歷程.
但是馮弟兄他就是.
信了一段初信主.
當年十二月受浸.
我也想聽一下他.
神在他生命裡面.
怎樣經歷神呢.
神怎樣在他生命裡面更新呢.
請他分享他的見證.
大家好.
多謝牧師今天邀請我上來.
去分享.
真的由十二月到現在.
四個月大概的這個.
靈命上的一個經歷.
以及成長在裡面.

$^{1241}$其實在上年十二月.
受浸的前後.
其實我那時候.
跟太太問太太請教一下她.
究竟一些屬靈的問題就是.
究竟怎樣知道自己.
在經歷神呢.
又或者是究竟神.
當神在說話的時候.
究竟是怎樣的呢.
怎樣聽得到.
怎樣聽明白祂說什麼呢.
那時候太太就說.
神其實跟每個人.
祂說話的方式.
都可以很獨特.
可以很不同.
讓經歷祂的方式.
都可以很不同.
如果你想知道的.
就去禱告祈求問一下神.
究竟經歷祂的神是怎樣的呢.
如是者我就向天父去禱告.
跟祂禱告就說.
我想知道究竟.
經歷你的那個情況.
那個感覺是怎樣的呢.
很感恩的就是.
短短四個月.
由十二月到四月.
現在這一刻.
真是每一個月.
我都跟太太是.
真真實實的經歷神.
祂是怎樣去工作的.
祂首先讓我們知道.
真是要放下自己的心思意念.
祂讓我們知道.
我們應該先祈求他的國和他的義.
祂其他所需的.

$^{1281}$祂都會交給我們.
又讓我們透過歌詞去領受得到.
即使是激風雨風浪的時候.
只要我們仍然是.
甘心願意屈膝在祂面前.
繼續去清揚祂的美善的時候.
祂的恩典必然會帶我們跨過.
亦都真是很令我們.
很真實的去體會得到.
或者經歷得到.
就是當我們是甘心樂意.
去獻上自己微小的一些.
微小之物的時候.
祂是會樂意去接受.
亦都是喜悅的.
亦都真是.
甚至是去供應更加多給我們.
正正是這件事的時候.
今個星期其實在思考.
分享哪一段的經歷呢.
而神真是很奇妙地.
真是知我心思意念.
俯瞰著什麼的時候.
祂身心熱辣.
讓我再一次.
今個星期去經歷祂.
一個小小的經歷跟大家分享.
首先要補充一點的背景就是.
我這一刻是在尋求工作的.
在工作上尋求工作.
在幾個星期之前.
其實在朋友介紹之下.
就介紹了一份工作.
就是一份比較我心儀的工作.
亦都真是交了申請.
去申請.
那又很有趣的就是.
今個星期的星期一星期二的時候.
又有姊妹.
另一個姊妹就說有一份工作.

$^{1321}$我說你要試一下.
那我就去申請一下.
那原本想著.
通常加拿大的速度.
都是很慢的.
要等待申請的申請日.
之後.
她才會開始去申請.
去process(過程).
那誰不知就是.
我申請的第二天.
突然間我在煮飯.
拿著一把菜刀的時候.
電話在響.
接著我聽.
我以為是.
看得到是那間公司打過來.
那我純粹以為.
是confirm(確認)一些細小的資料.
那誰知她突然間就.
跟我做了一個電話面試.
那我手忙腳亂之下.
又做了.
不知道是不是我.
這樣做完那個電話面試.
那她跟我說.
哦.
明天要做一個virtual interview(實際面試).
她也去做吧.
那.
那我就好了.
沒什麼問題.
那在mean time的時候.
原本幾個星期前申請的那份工作.
突然間又收到email(電郵)說.
哦.
禮拜五要做一個面試.
可以嗎.
那就變了.
好像突然間.

$^{1361}$本個星期想著.
想一下上台分享.
什麼經歷的時候.
變了我突然間.
全部都在那裡.
prepare(準備)面試.
那就在這個準備的期間.
那我就跟太太去分享.
分享一些.
關於屬靈的掙扎在裡面.
就說.
哦.
如果天父這次真的.
為我開路.
無論哪一份工作都好.
為我開路的時候.
假若自己.
哦.
做了自己的本分.
遞進了不足.
或者覺得自己做得最好.
仍然都不夠的時候.
或者是很想做完美主義者.
才能夠.
去令到天父.
去成就這件事的時候.
是怎樣的呢.
當時太太突然一句就說.
你為什麼跟神去想這麼多事.
祂的道路高過我們的道路.
祂的意念也高過我們的意念.
你放手讓祂.
做祂自己的事就可以了.
那時候我又一起去祈禱.
心裡面是放鬆了一點點.
但是當天父.
真的要跟祂說話.
祂要告訴你的時候.
你基本上是沒有辦法.
去忽略或者懷疑.

$^{1401}$祂真的實際在你身邊.
在說話給你聽.
那就在這一段準備面試的期間.
有跟我的朋友.
在電話上面去做禱告.
他用的經文是一幅.
在我工作台前面的一個海報.
是每天都會見到.
但是每天有時候.
都會忽視得到的一句經文.
但是接著.
之後做日常的靈修的時候.
我又再一次看到這句經文.
甚至再過一陣子.
跟我的朋友.
即是介紹那份工作的那個朋友.
去測試訊息.
那個朋友都仍然是用.
同一句經文去鼓勵我.
天父真的四次不約而同的.
都用同一句經文去跟我說話.
還要很風趣的.
怕我看不明白.
他用兩中兩應的跟我說.
是的.
但是那一刻看完之後.
我的心真的放下了.
放開看得到.
感受得到天父真的去支撐著我.
也都真的讓我去提醒我.
那麼至於星期五的那個面試.
那個結果是怎樣.
不知道的.
但是這一刻我很明白的就是.
天父這一句話語已經刻在我心裡面.
你要專心養賴耶和華.
不可依靠自己的聰明.
要在你一切所行的路上都要認定它.
它必使你的道路平靜.
曾言三章五節六節.

$^{1441}$多謝.
你的新的經歷是什麼.
你屬靈生命的成長的經歷又是怎樣.
最後給了後書一章八節提醒我們.
在結束這段經文的時候.
他說你們有了這幾樣.
再繼續增長.
所以無論是Kenning.
無論是我和你每一個人.
有了這些都要繼續的.
是一個繼續豐富的增長.
就必使你們在認識我們的主耶穌基督上.
不至於懶散和不結果子.
意思就是認識.
在這裡是在說經歷.
在生命裡面的遭遇.
就像剛才Kenning所分享的.
和你一樣.
在生活裡面你經歷了神.
怎樣和你說話.
這個就是這樣去認識神的.
不只是頭腦上的正式.
不只是看一些知識的認識.
而是經歷神.
主耶穌.
你越經歷得神多.
在你的生活.
透過神的話語和聖靈.
你越經歷得多.
你會認識神更多.
認識耶穌更多.
你在主耶穌學童.
在小組和在崇拜當中.
知識上認識更多.
你又會在經歷上更加增多.
其實是一個互動的關係.
所以我鼓勵大家.
不要放空.
你有一個很清晰的目標.
就是要結果子.

$^{1481}$因為如果你成長的時候.
你不會不結果子的.
這裡說繼續增長.
以致能夠結果子的.
不要讓他不結果子.
空轉就是打卡.
打卡回來也是在原地跑步.
其實是有得著有結果的.
是要活出愛和生命的果效.
所以我想請問你.
你每天的靈修.
你每個星期的崇拜.
小組主耶穌學.
幫助到你活出基督的生命多少.
你又結出了哪方面.
榮耀神的果子.
讓你數算.
讓你計劃.
讓你刻意地去追求.
去追求成長.
我很期盼你用未來幾分鐘.
你做一個立志.
現在有兩分鐘.
我想你做一個立志.
就是在勤成長.
不斷去靈命成長的事情上.
你立志你不停步.
不亂走.
是不空轉.
而是要因勤回應神的恩典.
刻意地建立靈命.
活出愛和生命的果效.
其實我們整個使命人生.
那五方面.
就是在幫助你和我.
去活出一個這樣結果子.
有果效的生活.
這五方面其實就好像一個.
檢查身體的表格.
這五方面我們活成怎樣.

$^{1521}$有哪些靈命成長的下一步.
在這五方面.
如果你是新來殷傳的弟兄姊妹.
你未拿這張杯墊.
你可以在出門口的時候.
招待師會給你的.
如果你不見了.
弟兄姊妹不見了.
你在門口你也可以拿回來.
讓這五方面是你在神面前問他.
我怎樣去不停步.
因勤的回應.
怎樣整存的建立.
以至我每天的靈修問四個問題.
起碼四個問題.
就是在這個靈修完結之後.
在這個崇拜之後.
在主日學之後.
在小組之後.
我知道什麼是神的心意.
對我來說.
我有沒有勇氣去活出.
祂對你發出的邀請.
你能不能夠對那些誘惑.
和阻礙神的事情說不.
然後就算有壓力.
你都能夠仍然持守.
可能是一個新的愛心.
可能是一個關係的重整.
可能是一個新的事奉崗位.
可能是延續現在的事奉崗位.
在你的事業.
在你的加拿大入籍的進程裡面.
你有哪一方面神現在要跟你說話.
我想你閉上眼.
我想你做一個立志.
你自己向神禱告.
主啊 我不願意停步 不願意亂走.
不願意空存.
我要追求成長.

$^{1561}$我要刻意讓主你.
建立我整存的生命 靈命.
主啊 我願意配合.
如果這個是你的回應.
你說主啊 我願意配合你.
在這類日子.
我也想你舉手.
這個是你和神立志.
不是給別人看的.
如果你今天說你願意立志.
在神的面前.
你說我願意竭力的去成長.
願意在主的裡面.
是配合他在生命的工作.
你有沒有舉手示意一下.
然後我們有個結束的祈禱.
神感動你 你就舉手.
在神的面前立志.
其實這一刻不是做給別人看.
而是讓神知道.
也是你自己生命裡面.
一個的回應.
你就記得你這一刻所作的.
回應和邀請.
我們同心低頭祈禱.
親愛的主我們求你憐憫我們.
求你感動我們.
求你兼顧我們每一個人.
在你面前的立志.
有聲無聲的.
或許我們已經舉手.
或者我們仍然.
未能夠作出這個決定.
我也求你給我們.
有一個令命成長的新一步.
以致我們的生命.
是活得有意義價值.
禱你的喜悅.
祈禱奉靠主耶穌基督的名靠.
阿們.

\newpage



\section{使徒行傳 10:9-16}
\label{sec:031994}
\textbf{求主藉異象激動我}~~~~何俊文牧師
\newline
\newline
連結: \href{http://www.church.com.hk/acms/content.asp?site=cdc\&op=show\&type=product\&code=031994\&layout=sermon}{\texttt{031994}} ~~~~ 語音日期: 2026-04-19
\newline
\newline
\hyperref[sec:031993]{\small{< < < PREV SERMON < < <}}
~
\hyperref[sec:index_chronic]{\small{[返順時目]}}
~
\hyperref[sec:031995]{\small{> > > NEXT SERMON > > >}}
\newline
\newline
使徒行傳 10:9-16
\newline
\begin{longtable}{cl}
\hline
\hline
章節 & 經文 (和合本修訂版)\\
\hline
10:9 & \begin{tabularx}{0.7\textwidth}{X} 第二天，他們走路將近那城，約在正午，彼得上房頂去禱告。 \end{tabularx} \\ \\ \relax
10:10 & \begin{tabularx}{0.7\textwidth}{X} 他覺得餓了，想要吃。那家的人正預備飯的時候，彼得魂遊象外， \end{tabularx} \\ \\ \relax
10:11 & \begin{tabularx}{0.7\textwidth}{X} 看見天開了，有一塊好像大布的東西降下，四角吊著縋在地上， \end{tabularx} \\ \\ \relax
10:12 & \begin{tabularx}{0.7\textwidth}{X} 裡面有地上各樣四腳的走獸、爬蟲和天上的飛鳥。 \end{tabularx} \\ \\ \relax
10:13 & \begin{tabularx}{0.7\textwidth}{X} 又有聲音對他說：「彼得，起來！宰了吃。」 \end{tabularx} \\ \\ \relax
10:14 & \begin{tabularx}{0.7\textwidth}{X} 彼得卻說：「主啊，絕對不可！凡污俗和不潔淨的東西，我從來沒有吃過。」 \end{tabularx} \\ \\ \relax
10:15 & \begin{tabularx}{0.7\textwidth}{X} 第二次有聲音再對他說：「神所潔淨的，你不可當作污俗的。」 \end{tabularx} \\ \\ \relax
10:16 & \begin{tabularx}{0.7\textwidth}{X} 這樣一連三次，那東西隨即收回天上去了。 \end{tabularx} \\ \\
[1ex]
\hline
\hline
\end{longtable}
$^{1}$弟兄姊妹早晨,好開心今日係當中同大家分享神ge 話語.
我都唔好習慣見到呢張相,可能下一張.
我黎jor Melbourne18年架啦.
咁我介紹我自己ge 屋企先啦.
其實我講到習慣大家可以比少少反應我架.
咁我今朝都幾早起身架.
咁如果你呢可以比少少反應呢.
我就會做D 醒架,得唔得?.
你同佳麗都講啦,唔該幫幫手咁樣.
你同佳麗都講,好唔好?.
好,預備,我黎jor 18年,我ge 家庭.
咁我太太而家繼續係Cosway教會事奉.
我而家要去計一計我個眼ge .
因為我有時我行得前呢,我好似望唔到你地.
我而家要望一望你地先.
我太太繼續係Cosway教會事奉架.
我而家都係Cosway教會事奉架.
咁就我大學畢業呢,我讀社工ge .
我大學畢業呢就去jor 一間福音界度事奉.
jor 五年,咁你睇我ge 裝束.
你又知道我留住一D 比較基層呀.
服侍社區呀咁樣.
所以我係教會一般呢就.
弟兄姊妹都叫我Alexge .
我14年按jor 木ge ,但我教會ge 文化都係叫我Alex.
咁因為係大家心目中呢.
抗牧師ge 形象呢,你起起身呀,幫幫手,唔該.
你幫幫手呀.
好有個性,我知道佢ge .
咁佢就似D 牧師ge 樣架嘛,係咪呀.
係啦,咁我係咁ge 樣架嘛.
我係咁ge 樣架嘛.
唔好意思呀,因為我第一次係度講到.
我唔係好慣,我通常會快少少架.
okok,呢個教會ge .
cosplay教會先,ok.
係啦,咁我.
我係呼音誡讀裡面.
似乎gwo 五年啦,有冇人.
聽過呼音誡讀架.

$^{41}$通常聽過神稀會啦.
我太太呢第一個工作就係.
係神稀會教書ge .
gwo 陣時都攀拖,所以神好大ge .
恩典比我地有機會.
好細個開始呢其實.
已經去接觸社區.
有需要ge 人啦,福音係神ge .
大能啦,要救一切相信ge 人啦.
如果我講ge 時候.
你認識你ge .
朋友係曾經被.
毒海捆綁.
經歷耶穌ge 改變,根深改變.
突破呢,我頭先.
講gwo 句,福音係神ge 大能.
要救一切相信ge 人呢.
就唔係口腔腔啦,唔係好似.
我地講完gwo D 經文.
其實未入心就講jor 出黎.
福音係神ge 大能.
要救一切相信ge 人.
如果你係座當中你仍然.
有一D ge 朋友.
甚至乎係你ge 家人.
係經歷.
係生命裡面ge 捆綁.
難關,痛苦.
福音係神ge 大能.
要救一切相信ge 人.
願意你同我都係一個信心ge 宣告.
雖然我地未睇見.
但我地相信我地ge 神係信實.
係慈愛,係憐憫,係有能力.
福音係神ge 大能.
要救一切相信ge 人.
咁我係gwo 五年裡面.
有好多學習啦.
有好多眼睛開啟.
個視野開啟啦.

$^{81}$頭先.
即係而家睇到相片裡面啦.
我自己每一個睇到ge 時候.
都好多ge 畫面.
裡面出現啦,咁樣.
咁我為jor 保護佢地ge 原故.
所以睇唔到佢ge 樣啦,咁就睇到.
睇到我ge ,大家睇唔睇到呀.
係咪其實都差唔多樣呀.
係呀,我係Cosway.
似乎我十八年.
之前讀jor 四年神學,再之前做jor 五年.
社工,咁其實我都唔細架啦.
咁就,但係我.
習慣jor 架啦,我早.
arm arm ,我同另外一個同工.
一齊黎ge ,咁佢今日係.
兩個教會講道,咁佢本身.
都係華傳ge 同工啦.
咁昨日就同阿鄭歐牧師.
大家唔識識阿鄭歐牧師,原來今日係.
差傳,孫教主日黎架喎.
今日係咪呀.
多謝阿鄭牧師遷就我.
就改jor 叫,從孫教.
坐係度差傳.
係咪,咁.
其實差傳孫教.
對我黎講呢,係mission.
我地通常呢.
好耐好耐以前呢,就會譯.
孫教嘛.
但係對我黎講呢,係使命.
係咪,如果你同我.
冇一個使命,mission life呢.
唔會.
好似我地Cosway教會.
張牧師,我地教會有七十幾個孫教士.
咁.
我地其中一個要求呢.

$^{121}$就講話本身,係本地ge 裡面.
都係一個missionge 人.
先至能夠去考慮成為孫教士.
因為,主持人牧師就講呢.
即係gwo 個人呢.
唔會搭jor 飛機落jor 去呢.
就突然間本身唔孫教.
mission,搭jor 飛機落jor 去.
就變jor 孫教,你聽得明架,o下.
所以今日大型節目,無論.
今日個主題講緊孫教.
無論講緊差傳,其實.
最返到核心ge 就係.
我同同神ge 關係,有冇一個聆聽.
而順服,有冇一個回應.
去過緊一個.
孫.
使命ge 人生.
o下.
係呢張圖裡面呢.
就有一個人呢.
就冇蓋住ge ,係咪.
即係唔係兩兩個,一個係我啦.
咁,我都為到.
呢個弟兄ge 感恩ge ,咁其實.
現在佢.
佢好多年前.
吸毒.
去jor 中心裡面戒毒,跟住.
留返低,係成為一個.
過來ge 童工,咁你呢D 名稱.
我諗如果你去過方戒毒,你知道大概.
過來ge 童工即係以前係度吸毒ge .
跟住重新得救改變.
成長,留返去服侍啦.
跟住讀神學啦.
結一婚啦,現在呢.
就去同.
唔同ge 組織.
機構呀,敬拜團隊呀.

$^{161}$就去分享佢ge 見證啦.
咁我都為到佢ge 生命.
好感恩,係咪,福音係神ge .
大能,要救一切.
相信ge 人,呢D 就好似雲彩ge .
見證人,去鼓勵我地.
神.
有陣時可能好抽象.
有D 知識性.
但原來當我地見到一D .
如雲彩ge 見證人.
佢地係唔同ge .
生命裡面,去活現佢ge .
信仰,活現佢.
經歷神ge 能力ge 時候呢.
嘩,好有能力架嘛.
集合埋一齊呢,其實就係一個.
Mission ofge 群體囉.
就係一個去回應.
宣教ge 群體囉.
咁今日我同大家分享ge 題目呢.
係求主直意象.
激動我,咁我相信.
大家都好熟悉呢段經文.
咁呀.
意象同埋使命.
呀.
睇你心目中.
到底你諗到意象係D 乜ye 呢.
ok.
你地返緊工ge .
機構有遠景ge .
有Visionge ,舉一舉手.
或者你以前啦.
以前都有返工架嘛.
有遠景ge .
講比你聽,係咪.
對我黎講,我地記憶圖嘛.
意象係從神而來ge .
係持久ge .

$^{201}$係讓神得榮耀ge .
意象係永恆ge .
人要配入神ge .
遠景係乜ye ,遠景係從人而來.
短暫ge ,係.
為到人ge 益處ge ,係屬地ge .
人想操控神.
ok,諗一諗.
今日當你聽.
今日ge 訊息,求主直意象激動我.
想唔想gwo 個意象,係從神而領受ge 時候.
咁我問一問自己囉.
或者你都有D 諗法,無論係諗法.
意象,遠景,使命,咩都好啦.
問一問自己.
左邊定右邊,係咪.
講個例子.
我唔知道.
感覺弟兄姊妹就應該.
都黎jor .
都一段時間ge ,係咪.
黎jor ,我就十八年,不如用十八年.
去衡量啦,超過十八年ge .
即係你senior過我ge .
係澳洲.
嘩,好多senior,ok,ok.
可能呢個故事就.
已經係遠D 啦.
你唔係好記得啦,當你賣屋ge 時候.
見到一間心儀屋.
ge 時候,係咪通常都.
O主啊,如果你.
按你ge 志成就呢間屋呢.
我會去開放啦,被你.
使用啦,開屋啦.
讓神你的工作.
能夠係屋企裡面.
能夠發生,係咪?.
通常都係咁過,係咪?.
但係有冇D 記憶圖呢?.

$^{241}$買jor 間屋,你唔記得神架喎.
買一間好ge 屋,係咪?.
唔記得jor .
你,大英之妹,你,現在.
諗緊ge ye ,係你自己.
個人ge 遠景啊,還是.
係神ge .
異象,還是.
係,當你見到.
神ge 異象ge 時候呢,你甚至乎.
放低你一D 過去ge .
傳統,你一D ge 習慣.
你一D 生活形態.
甚至乎,你覺得.
gwo 樣係你.
好緊要ge 一種.
一種信仰ge .
執掌,但係講D .
唔係真理黎ge ,真理就唔.
可以放低架嘛,但係D 信仰ge 傳統.
甚至乎好重要ge ye .
但係如果為jor 神ge 緣故.
你願意能夠調整.
改變同配合,我地今日睇.
睇呢,彼得ge 故事.
頭先大家讀jor 啦.
都讀jor D .
經文,我比大家有個大綱.
先啦,呢段經文呢其實.
係一個非常之重要.
好特別ge 轉折啦.
如果你講,使徒行傳.
又係靈性行轉.
神ge 聖靈係當中好多ge 工作.
呢段經文ge 上文裡面.
就講到哥利劉聽見.
神ge 異象,哥利劉係.
邊個呀?係一個.
虔誠ge 人.
係一個官,佢.

$^{281}$周濟窮人,佢.
禱告神,咁就哥利劉.
呢,就三次.
聽到神ge 聲音,跟住就.
派jor 三個僕人去搵.
彼得,三三三你都知道.
對於彼得黎講就係.
好紮實ge ,你知唔知?.
即係彼得呢,係呢段經文裡面.
就當聽到三.
如果咁講就係.
呢段經文呢就十七.
至到三十三次就聖靈.
三次同哥利劉.
講,就話.
同彼得講,哥利劉.
出現,跟住彼得.
就見到.
今日頭先讀gwo 個經文,見到.
一個不潔正ge 意象.
三次ge .
對於彼得黎講呢.
點解三咁重要?.
佢三次不認主.
跟住返返去提比利亞海.
去重操.
故業,去打魚ge 時候.
耶穌.
去到彼得ge 面前.
用三個.
你愛我比.
這些更心魔,所以對於.
彼得黎講呢.
三次.
去經歷一D .
聖靈ge 提醒,係好.
特別啦.
呢段經文ge 裡面講到.
當彼得.
見到從天.

$^{321}$上黎,一個.
不潔正ge 意象.
就結jor ,結果就改變jor .
佢ge 人生.
丁子梅,今日同大家.
分享,求主直意象激動我.
諗一諗,你而家諗緊ge .
係神ge .
意象,還是.
係你個人ge 想法.
夢想遠景.
當.
唔係,我而家好清楚.
神有對我一個ge .
帶領,意象.
我要去放.
低自己ge 時候.
你又點樣去為到神ge 角度.
去擺上呢.
今日,呢個不潔正.
意象,係咪,今日同大家.
分享,好簡單,三個ge .
角度啦.
意象呢,係會人令人.
意象呢,係會.
改變人,意象呢亦都呢.
唔單止帶領.
人ge ,唔單止改變ge .
個改變呢,係.
持久ge .
係.
呢段經文ge 裡面,有講到.
有好多聖靈ge 工作.
譬如係第十章.
第三節,我地一齊.
讀一讀D 經文先,好嗎.
第三節,預備,一二三.
第十九節都大家讀啦.
彼得還說者話ge .
時候,聖靈降在一切.

$^{361}$聽到ge 人身上.
第四十五節,第四十七節.
裡面都包括jor .
係係士堂主裡面.
有好多聖經ge 記載.
講到.
神去,透過聖靈去啟動.
去同歌利流講啦.
歌利流就跟住.
就派三個僕人去搵.
彼得啦,跟住呢.
彼得去聽.
呢個.
歌利流派三個.
僕人去到.
一個外邦人家裡面.
講到,呢個過程ge 裡面.
係彼得.
三次見到頭先.
gwo 個圖裡面講gwo 個.
神ge 異象啦,咁樣.
所以,但係你聽到.
呢D 故事裡面講到.
有咩特別je ,咁.
歌利流去到.
彼得.
去到.
見到一D 人.
係見到佢地.
相信上帝.
跟住奉聖靈ge 名.
去受洗,有D 咩.
有D 咩特別je ,係第二章.
三十八節到三十九節都咁講啦.
彼得說,你門各人要悔改.
奉耶穌基督ge 名.
受洗,叫你門ge 罪.
得赦必領受所賜ge .
聖靈,第三十九節一齊讀.
預備,一二三.

$^{401}$整個故事.
Alex,到底你講乜ye .
咩歌利流.
三個僕人.
有咩聖靈引領.
彼得就去jor 一個外邦人.
ge 家裡面,跟住就.
同佢傳福音,見到佢地.
領受聖靈.
佢地ge 形態就同.
彼得不嬲服侍開.
gwo D 猶太人,一樣ge ,差唔多ge .
所以就叫佢地受洗,歸入.
神ge 家庭.
大家庭啦.
有D 乜ye 特別je ,其實.
好似講來講去都係我地.
一路平時成長講開話,信耶穌.
跟住得永生,跟住就受洗.
見證耶穌基督ge 信仰.
第二章三十八節,三十九節.
都係咁講啦.
但我想話比大家聽,係.
《使頭行轉》第二章裡面講ge 時候.
彼得講ge 時候呢,其實.
係佢ge 世界觀ge 裡面呢,其實係.
冇外邦人ge ,佢講ge 時候.
個對象呢,係講緊.
佢所同基,好叔呀.
好親朋戚友,好似大家咁講.
你望一望隔離,gwo D 自己人黎架.
因為係當時.
彼得,佢gwo 個信仰觀.
黎講呢.
外邦人呢,係同佢.
勢不兩立ge .
咁講多少少比大家聽,你就了解.
多少少ge .
當時呢,其實呢外邦人.
係被,即係D 猶太人.

$^{441}$講外邦人ge 時候呢,係形容.
佢地係一個動物黎ge ,大家.
想寵物架,又可以舉舉手.
ok.
你隻狗嘛.
有D 同佢一齊訓架嘛.
有架,ok.
當時外邦人.
試猶太人呢,係狗黎架.
ok.
唔單止係咁喎.
你想像呢.
地獄係乜ye .
顏色架.
估下,唔複雜ge .
大聲少少,唔該.
黑色,仲有呢.
仲有咩.
咩咩咩.
紅色,點解紅色.
我黎呢邊下.
有火架.
猶太人.
試外邦人.
係地獄裡面ge 燃料.
所以你可以想像呢.
當時彼得其實講呢個說話ge 時候.
佢根本冇外邦人ge 角度架.
因為當時猶太人.
同外邦人呢,係.
制不兩立ge .
咁所以你睇返呢段經文呢.
彼得.
聽到.
被聖靈引領.
去進入jor 一個.
佢從前.
根本唔會進入.
唔會回應ge .
環境ge 裡面.

$^{481}$就係神.
帶領佢.
由佢一個比較安舒.
熟悉.
自己人.
格吉勒ge 環境.
進入到去接受.
哥利留ge .
邀請去到一個外邦人.
ge 群體.
佢講到.
挺姊妹,原來今日.
整個裘主席.
異象激動我,第一個訊息想同大家.
講,就係原來.
彼得係當.
佢被聖靈引領.
ge 時候,就係.
慢慢整個過程.
無論同佢思維.
傳統價值觀.
佢ge 習慣.
不停咁轉變ge 過程.
裡面,被主ge 引領.
但係個過程裡面.
係聖靈.
不停咁引導.
我地有冇一種對聖靈.
敏銳ge .
操練呢.
講一D 比較實際ge 例子.
今次我同.
有個同工,生命導師.
同工,係唔同ge 教會裡面.
有機會去到唔同ge 樹梨教會.
去分享.
同工前日.
去行.
早上行,做運動.
佢就.

$^{521}$有兩條路選擇,咁佢就.
聽到一個好微小ge 聲音.
行gwo 邊.
佢有少少掙扎ge ,行埋去gwo 邊.
佢就.
去敏銳.
行去左邊,係聖靈.
提佢ge ,咁佢就行出去gwo 邊.
咁我都為佢感恩.
佢信服jor 聖靈ge .
引領呢,結果佢就.
去到公園另一邊,就見到.
一班中國人,咁其中一個.
朋友呢,跟住同佢.
傾下計,本身佢諗住講唔到.
福音ge .
差唔多,差唔多要走ge 時候呢,另外.
gwo D 人都走jor ,跟住之後呢.
結果就同gwo 個.
朋友呢,就去.
分享jor 耶穌,跟住.
就帶jor 佢去見耶穌.
佢返黎同我分享,就話.
阿力手,感恩,感恩.
gwo 個敏銳就係行左行右.
佢其實.
話比我聽呢,佢gwo 一刻呢.
其實佢感覺到神要.
提佢要去行去gwo 邊.
gwo 個就係gwo 個被聖靈.
引導,gwo 個.
過程就係要放低自己.
一D ge 習慣,一D ge .
傳統,一D ge 想法,去.
去學習,去.
被上帝去引領.
頂尖,我諗我地生活ge 裡面.
好多ge .
情況呢,裡面呢,我地都.
需要呢,被上帝.

$^{561}$ge 靈去引導.
呢個係第一,第二樣我想同大家.
分享呢,就係對於.
彼得.
gwo 個性格.
到底你對於彼得.
呢個人呢,你有.
對佢ge 印象係點呢.
頭先我講jor 啦,係咪.
我一路郁,你地都郁架喎.
ok,係喇,ok.
彼得三次.
不認主,係咪,呢個最出名,係咪.
仲有呢,你諗起彼得.
諗起乜ye ge 故事呀.
或者D 乜ye 你深刻架.
彼得,我提jor 你架喇.
一般我就應該叫佢禁住先.
唔好出架.
咁呀.
呢段經文裡面講到.
彼得就拉住他勸他說.
主呀,萬不可如此.
這事必不臨到你身上.
同邊個講架,彼得.
同耶穌講架嘛.
耶穌講乜ye ,耶穌話佢要走上.
十字架ge 道路.
跟住彼得.
斬釘截鐵咁樣講.
咩意思,彼得就同耶穌講.
呢D 事呢,唔可以.
臨到你ge 身上,好.
激動架,跟住.
耶穌係最後ge 晚餐.
到門徒洗腳ge 時候.
跟住聖經裡面講.
你永不可.
洗我ge 腳.
你留意gwo D 詞語.

$^{601}$有冇當中係夫婦一齊.
黎崇拜架,可唔可以舉舉手.
留意下你地講ye ,同夫婦.
意思,你永遠都係咁架.
gwo D 呢.
好似對,講到好似.
gwo 個就係真理,彼得.
就係咁ge 人囉.
斬釘截鐵,耶穌係神ge 兒子.
耶穌係夫子.
耶穌話比佢gwo D 門徒聽.
彼得亞國約翰.
話我要走上十字架ge 道路.
gwo D 門徒.
聽就聽得到.
但又聽唔盡,所以就好多.
自己ge 諗法,係咪.
知唔知我地,其實講jor .
好似Supression咁,左定右.
講jor 囉.
其實係Supression話.
其實佢知道呢,一行jor 呢就要開工.
佢話我想休息,行下.
行下散步je ,係咪.
同上帝係度拗數ge .
你明我意思o下.
你呢,你同彼得.
你會唔會relate到彼得呢.
斬釘截鐵.
係,仲有裡面講.
係好掘牆啦,係咪.
我就是必須和你同死.
也總不能不認,你呢度講咩.
三次不認主之前,係咪.
耶穌就講話當中有門徒.
係會唔認我.
跟住又係講住.
其他人.
其他人會,我都.
唔會.

$^{641}$係咪.
最經典,你有冇記得一個圖像.
圖畫,係咪,彼得.
就係當耶穌.
被兵兵.
同埋D 人捉拿ge 時候呢.
有一幕呢.
你記唔記得其實都幾血腥架,有冇印象呀.
當人呢.
要捉拿耶穌ge 時候呢.
gwo 段經文講架嘛.
呢度冇寫架嘛,做比大家睇啦.
就係.
有人,有個兵兵.
捉耶穌ge 時候呢.
有一個人呢,就lor 把刀仔呢.
就,咁.
你想像下聖經,gwo D 耳仔.
彈jor D 血出黎架.
你知唔知gwo 個邊個黎架.
gwo 個彼得黎架.
點解我要同大家.
去進入.
彼得ge 性格gwo 個世界.
佢講ye 好掘牆.
佢係跟從耶穌ge .
彼得,雅各,約翰.
跟耶穌去到.
見到耶穌.
登山變象ge 時候,變形象ge 時候呢.
就gwo 班人就話,嘩.
呢個神ge 形象,我地想同留.
係神ge 同在裡面呀,好似頭先.
唱敬拜ge 時候,係咪呀,大家好.
享受弟兄姊妹.
無論透過安靜,透過試歌.
透過幫我地進入神ge 同在.
有時,唔好講到住,我想佢敬拜.
即係gwo 種個狀態呀,gwo 個性格.
我永遠不會不認你呀.

$^{681}$其他人唔會唔認你呀,即係好驕傲.
好覺得自己不可一世ge .
好覺得自己掌握曬.
神ge 心意ge 一個咁ge 人呀.
如果咁ge 心路歷程呢.
再返返去今日gwo 個故事呢.
彼得見到.
從天上黎面ge .
異象,三次呀.
跟住就原來.
係呢個過程裡面呢,就見到.
哥利劉派三個ge 僕人.
去到彼得.
ge 眼前.
即係老實講,如果彼得.
冇gwo 個從天上黎面ge 異象.
神比佢ge 時候呢.
我諗佢一見到個僕人呢.
僕人未開聲呢.
就老實講呢,肯定硬開佢.
走呀!因為佢係狗.
佢地獄裡面火,係咪?.
la ,個故事係咁樣嘛.
ok,當.
彼得,即係人啦.
被聖靈引領底下.
見到神ge 呼聲,跟住呢.
就開始呢,被神引導.
改變,個性格開始要.
改變呀嘛,ok,進入去呢.
你睇返個經文就好有趣啦.
結果呢,就.
即係彼得,我唔知佢當時聖經冇講呀.
唔知到底佢係咪死死地起.
定係點樣啦,但係佢當佢進去到.
哥利劉.
見到哥利劉ge 時候呢.
個經文裡面就26節gwo 度話.
彼得拉住我,說你.
起黎,我也是人.

$^{721}$其實講緊乜ye 呢?即係你都想像呢.
當彼得見到哥利劉ge 時候呢.
哥利劉係點樣架.
成個,原來第十節呢.
25節呢,係講緊哥利劉.
係跪低ge .
係,即係.
好尊重彼得,多謝你黎.
咁彼得係點呀.
彼得當時呢,我相信.
投入佢當中呢.
你緊係跪啦,其實我真係唔想黎架.
你班外邦人呢.
我同你真係誓不兩立ge .
與不相干.
但係唔係喎,彼得.
拉住哥利劉,叫佢起身.
我也是人.
等一陣你睇唔睇到.
gwo 個成個.
被聖靈異象.
改變整個.
個過程裡面,gwo 個變化.
當一個人.
願意擁抱上帝ge 異象.
佢ge 性格,佢ge 行為.
佢ge 生活模式,願意.
完全改變,佢ge 性格都開始.
逆向佢ge 性格架,佢原本.
覺得我真係唔想黎架.
你跪,當然要跪啦.
起身,你起身啦.
哥利劉,我也是人.
結果呢.
當彼得進到哥利劉家.
見到哥利劉同gwo D 外邦人呢.
同信耶穌ge .
猶太人都差唔多ge 時候呢.
佢就口快快呀.
佢地受洗歸主啦.

$^{761}$點解口快快呀.
因為係當時gwo 個處境呢,佢地要受割禮架.
彼得.
口快快其實出於jor 佢對.
gwo 班人,gwo 種ge 憐憫.
gwo 種一視同仁呀.
信仰能夠超越jor 種族.
超越jor 傳統,超越jor .
價值觀,因為係耶穌基里面.
就係一家人呀.
但係姐妹今日求主.
直意象激動我,你呢.
你睇唔到神ge 意象呢.
我講講點解我今日.
有機會黎到呢度囉.
我其實.
繼續係Cosplay事奉.
五年前呢.
神讓我有一個機會就係.
其實教會裡面.
有一個coach ministryge 侍工.
其實做jor .
gwo 陣時做jor 十五年,即係而家二十年啦.
Cosplay係一個比較.
有規模ge 教會,即係英文.
英文部,有九個分堂.
當時其實已經有十五年ge .
生命導師做英文ge .
咁就已經.
同澳洲.
英國,大概一百個.
英文教會去合作.
訓練一D 生命導師.
願意一年,一個星期,一個鐘.
關心社區裡面有需要ge 人.
咁當時我就.
去long story short.
就係gwo 個童工.
知我參加華福會.
完jor 之後佢問我.

$^{801}$話Alex啊,你有冇.
同gwo D 中文牧者呢.
好似啊,個牧師一見.
sorry啊,冇意思,對唔住.
爛晒.
ge 牧者呢.
去分享,我其實.
答佢,你都知我個性格啦,因為我同佢太熟啦.
十幾年同事嘛,冇喎.
唔會喎,因為我冇時間.
因為唔係我ge 角色.
但結果呢.
gwo 個星期呢,神就.
call我,神就感動.
咁我就去回應啦,結果我.
好感恩.
其實我都好感恩就係.
透過呢幾年,有機會神.
讓我有機會同唔同ge 中文教會.
咁我都,現在.
大概有十幾個中文教會啦.
有D 係雪梨,有D 係香港,有D 係.
主要係Melbournege .
而家大概有四十幾個.
中文ge 個案啦.
不如我比大家睇一睇個故事先啦.
頭先佢好stress架,而家做得好好.
大聲少少得唔得.
係,係,耶,耶,耶,唔該晒.
係,對唔住.
我知道,而家我差唔多要落台架啦.
唔緊要,我知架,因為通常.
你講四十呢,係official.
時間,我有五分鐘.
不凡,我知道.
我係澳洲都差唔多三十年啦.
當時.
我覺得我.
無得揀ge .
有一支手.

$^{841}$比我,我就拉住ge .
我覺得.
無方向啦.
我覺得.
好孤單係呢度.
我都唔知我點解.
要黎澳大利亞ge .
我第一次見到佢呢.
我覺得佢係有個.
好勇人.
我好想知道.
佢.
可以幫到D 我D 咩呢.
我第一次見到Jenny.
好孤單啦.
好無方向啦.
好需要人陪伴佢,聽佢講.
咁所以.
我同佢見到年之後,好願意.
同佢一齊同行囉.
佢真係好似.
我ge 媽媽.
電視ge 打電話啦.
問我.
好唔好呀.
有咩ye 發生啦.
你家人好唔好呀.
你有冇打電話.
返屋企啦.
覺得平安囉.
聲音.
咁我自己係基督徒啦.
我地大清智會都有愛ge ,我地教會有愛ge .
咁但係呢就唔夠囉.
咁我地身邊,我地社區.
我地身邊好多人都需要.
幫助啦.
同埋好需要有人去.
陪佢聆聽佢啦.
支持佢啦.

$^{881}$咁所以我就去參加生命導師.
我感覺我好想.
同佢講多D 囉.
遇到過咩人啦.
講D 咩ye 啦.
有咩感覺啦.
我覺得我.
搵到我ge 方向囉.
我感覺.
係耶穌底下.
我感覺有一個.
屋企啦.
神話比我知.
我要.
重活一次.
之前.
做錯過好多ye .
我睇錯過.
好多事,好多人.
而家我重新.
睇番自己.
需要改進囉.
我想講呢,即係我地.
生命導師呢個program呢.
係真係好好ge ,即係唔係淨係.
我大英姊妹啦,我關心大英姊妹.
就係呢個生命導師ge program.
比我知到呢,更加即係.
我地出面ge 社會,我地周周呢.
係好多需要幫助.
亦都好需要耶穌愛係裡面ge .
黎講係個使命人生.
重生時ge 人間.
都係有真情.
我都感覺.
從我屋企到.
社會呢,我都係感覺都係.
有條件ge 愛ge .
就算父母.
朋友都係有.

$^{921}$有條件ge ,我感覺.
係信得愛.
係無條件ge .
當我睇返呢個故事ge 時候.
我心裡面好感恩,其實係透過.
生命導師ge 平台.
結連jor 唔同ge 教會.
MCBC係Mentalge 教會啦.
跟住Participant.
就去jor City住ge .
所以Cityge .
福音堂ge 牧師就接觸過佢.
之後佢就係係Cityge .
Mentalge 教會,佢就係.
Mentalge 教會,佢就係Cityge .
以祈禱靠耶穌基督的姓名祈求.
Amen.
\newpage



\section{馬太福音 7:13-27}
\label{sec:031995}
\textbf{二揀一}~~~~石偉漢牧師
\newline
\newline
連結: \href{http://www.church.com.hk/acms/content.asp?site=cdc\&op=show\&type=product\&code=031995\&layout=sermon}{\texttt{031995}} ~~~~ 語音日期: 2026-04-19
\newline
\newline
\hyperref[sec:031994]{\small{< < < PREV SERMON < < <}}
~
\hyperref[sec:index_chronic]{\small{[返順時目]}}
~
\hyperref[sec:031996]{\small{> > > NEXT SERMON > > >}}
\newline
\newline
馬太福音 7:13-27
\newline
\begin{longtable}{cl}
\hline
\hline
章節 & 經文 (和合本修訂版)\\
\hline
7:13 & \begin{tabularx}{0.7\textwidth}{X} 「你們要進窄門。因為通往滅亡的門是寬的，路是大的，進去的人也多； \end{tabularx} \\ \\ \relax
7:14 & \begin{tabularx}{0.7\textwidth}{X} 通往生命的門是窄的，路是小的，找到的人也少。」 \end{tabularx} \\ \\ \relax
7:15 & \begin{tabularx}{0.7\textwidth}{X} 「你們要防備假先知。他們到你們這裡來，外面披著羊皮，裡面卻是殘暴的狼。 \end{tabularx} \\ \\ \relax
7:16 & \begin{tabularx}{0.7\textwidth}{X} 豈能在荊棘上摘葡萄呢？豈能在蒺藜裡摘無花果呢？憑著他們的果子，就可以認出他們來。 \end{tabularx} \\ \\ \relax
7:17 & \begin{tabularx}{0.7\textwidth}{X} 這樣，凡好樹都結好果子，而壞樹結壞果子。 \end{tabularx} \\ \\ \relax
7:18 & \begin{tabularx}{0.7\textwidth}{X} 好樹不能結壞果子，壞樹也不能結好果子。 \end{tabularx} \\ \\ \relax
7:19 & \begin{tabularx}{0.7\textwidth}{X} 凡不結好果子的樹就砍下來，丟在火裡。 \end{tabularx} \\ \\ \relax
7:20 & \begin{tabularx}{0.7\textwidth}{X} 所以，憑著他們的果子就可以認出他們來。」 \end{tabularx} \\ \\ \relax
7:21 & \begin{tabularx}{0.7\textwidth}{X} 「不是每一個稱呼我『主啊，主啊』的人都能進天國；惟有遵行我天父旨意的人才能進去。 \end{tabularx} \\ \\ \relax
7:22 & \begin{tabularx}{0.7\textwidth}{X} 在那日必有許多人對我說：『主啊，主啊，我們不是奉你的名傳道，奉你的名趕鬼，奉你的名行許多異能嗎？』 \end{tabularx} \\ \\ \relax
7:23 & \begin{tabularx}{0.7\textwidth}{X} 我要向他們宣告：『我從來不認識你們，你們這些作惡的人，給我走開！』」 \end{tabularx} \\ \\ \relax
7:24 & \begin{tabularx}{0.7\textwidth}{X} 「所以，凡聽了我這些話又去做的，好比一個聰明人把房子蓋在磐石上。 \end{tabularx} \\ \\ \relax
7:25 & \begin{tabularx}{0.7\textwidth}{X} 風吹，雨打，水沖，撞擊那房子，房子總不倒塌，因為根基立在磐石上。 \end{tabularx} \\ \\ \relax
7:26 & \begin{tabularx}{0.7\textwidth}{X} 凡聽了我這些話而不去做的，好比一個無知的人把房子蓋在沙土上。 \end{tabularx} \\ \\ \relax
7:27 & \begin{tabularx}{0.7\textwidth}{X} 風吹，雨打，水沖，撞擊那房子，房子就倒塌了，並且倒塌得很厲害。」 \end{tabularx} \\ \\
[1ex]
\hline
\hline
\end{longtable}
$^{1}$丁字幕到平安.
我們一起去學習神的說話.
我們一起祈禱.
天父,我們求你去幫助我們.
還有你讓我們一路去思考登山寶訓.
讓我們明白耶穌出來傳道的時候.
所給我們的每一個提醒.
讓我們今天在結束的部分.
讓我們都能夠整理.
或者我們認真去思考.
我們要如何回應主理在山上的部分.
求你讓我們今天聽到的時候.
我們都能夠明白.
或許我們有很多疑問.
但求你讓我們有聰明的智慧能夠曉得.
幫助我們.
聽我們祈禱.
奉耶穌名字而求.
阿們.
今天的題目叫「二揀一」.
不知道你有沒有聽過一句說話.
叫做「有得揀,你才是老闆」.
聽過沒有?沒聽過?.
這裡其實不止只有兩個老闆.
因為你們都有得揀.
所以有時候你們都是老闆.
有時候我們以為我們有得揀.
但有時候我們以為有很多選擇給我揀.
但不知道你有沒有發現有時候.
有很多選擇的背後.
其實最後只有「二揀一」.
今天我們都藉著這個題目.
去思考著我們在瑪特芬登山寶訓的結束裡面.
我們總結一下.
首先要看看大家.
你看過這個節目沒有?.
看過沒有?看過的舉手先確認一下.
大家知道他在做什麼.
看過嗎?知道他在做什麼.
這個就是其中一個被稱為永恆的電視台.

$^{41}$因為這個電視台一直都倒閉不了.
這個電視台裡面其中一個很重要.
成為他們翻本的節目.
就叫做「八萬富翁」.
「八萬富翁」裡面是透過十五條問題.
一直提升難度.
由最少的題目開始.
每贏一條就會晉級.
可以贏得更多獎金.
直到最後.
最後的題目會關乎你能否成為「八萬富翁」.
能否得到一百萬港幣的一條問題.
所以這個節目很受歡迎.
亦都令到平時不看這個台的都特意去看這個台.
所以今天給你機會做「八萬富翁」.
看你能否成功做到「八萬富翁」.
其中「八萬富翁」的主持就是陳啟泰.
他其中在他的節目裡面.
第一個得到一百萬的觀眾.
他就面對著這一條問題.
看你懂不懂得回答.
中國名勝蝴蝶泉位於哪一個地方.
誰懂得回答.
一百萬.
要看看是否.
不要馬上給ABCD.
陳啟泰通常說一輪話才會給ABCD.
既然大家都想知道.
我們給你四個答案選擇.
第一個在A廣東.
B在四川.
C在雲南.
D在新疆.
請問你認為哪一個會是蝴蝶泉所在的地方呢.
一百萬.
雲南C嗎.
你們為何如此肯定.
你去過嗎.
沒有.
沒有嗎?你這麼肯定.

$^{81}$知道那裡很漂亮.
一個很漂亮的地方.
好.
有些人選C.
不過在那天的節目裡面.
他最後還有一個錦囊未用的.
就是五十五十.
所以我們都給他用五十五十的錦囊.
看看有誰選擇.
五十五十錦囊.
最後剩下B四川C雲南.
一百萬在你手裡面.
究竟在你身邊飛走.
還是留在你的袋裡面呢.
認為B的請你舉手.
認為B的.
B的.
只有兩個.
B很出名,狗仔溝也很厲害.
C,雲南.
嘩,有這麼多個雲南.
你們只是碰巧吧.
還是真是認識的.
只是碰巧,真是碰巧.
有沒有人真是認識的.
不是碰巧的.
真是認識的.
因為他不是雲南人.
是嗎.
貓屁股想不想選C.
是.
哦,這樣.
雲南四季而出.
這些就是真的.
如果真的贏.
這些是完全有智慧.
有知識的去贏.
好,確定答案.
好,看看答案是甚麼.
Final answer.

$^{121}$雲南.
真是雲南.
原來這條一百萬的題目.
真是難你不到的.
所以兩個選一個.
你選到雲南了.
是不是很爽呢.
做一百萬富翁.
如果真是一百萬的話.
不過今天跟你玩另一樣東西.
你也是做一百萬富翁.
不過我們會玩另一個等級.
不過先看一下.
我們這個隱形先.
其實我們已經看了登山寶訓.
由第五章看了第七章.
我們今天嘗試直接跳看多些經文.
就完成了整個登山寶訓的教導.
登山寶訓的教導.
其實就來自一開始說甚麼.
有福了.
是吧,八福.
不斷說有福了.
有福了.
一直就說怎樣去成為.
跟隨耶穌的門徒其實是有福的.
於是耶穌就說到.
怎樣認真去理解聖經.
和新舊約去比較.
用一個耶穌的角度.
去理解舊約聖經的意思.
也都解釋了很多錯誤的觀念.
也都調整我們和神的禱告的關係等等.
都看了很多.
所以整個登山寶訓的重點.
不是要讓我們成為百萬富翁.
不過可以引導我們成為一個有福的人.
所以今天你已經聽了這麼久.
登山寶訓.
所以今天我們就要去到一個.

$^{161}$有福人生大問答.
看看你能不能夠順利拿到這個大獎.
在這個遊戲裡面.
今天這個遊戲裡面.
我們就會有錦囊的.
這個有福人生大問答裡面.
有三個錦囊.
第一個叫五十.
五十.
剛才大家明白了.
一半機會.
五十五十.
第二個就是打電話問親友.
第三個就是問現場觀眾.
有三個錦囊.
待會你嘗試去.
能不能夠你成功拿到這個有福人生.
準備好了嗎.
我真的問了.
認真的.
好.
我們先看看第一條有福人生的問題.
看你能不能夠懂得回答.
第一條有福人生的問題就是.
以下哪一條路是通往生命或者永生的門呢.
有四個答案給你選擇.
第一個.
門寬.
路大.
人多的.
那條路就是.
B.
門窄.
路小.
人少的.
C.
條條大路都得永生.
D.
以上皆不是.
好.

$^{201}$這個是什麼.
D啊.
D for delta.
以上皆不是啊.
OK 好.
好.
有什麼.
好.
我們這個有福人生大問答裡面.
它的錦囊是隨機發放的.
每一條題目都會隨機使用一個錦囊.
好.
我已經收到風了.
開始啟動第一個錦囊是什麼呢.
五十 五十.
這個錦囊就是五十 五十.
好 看看最後現在剩下哪兩個答案給我們選擇.
剩下兩個.
就是.
A 門寬.
路大.
人多.
B 門窄.
路小.
人少.
不知道你覺得是哪一個是什麼.
不過我們試試看.
分析一下每一個答案.
看看大家能否猜到答案是什麼.
門寬 路大 人多.
是怎樣呢.
原來它說.
這個門.
原來它說什麼呢.
就是通往滅亡的門是寬的.
路是大的.
進去的人都多的.
OK.
原來是這樣的.
不過什麼叫門寬呢.

$^{241}$我不知道你明不明白什麼叫門寬.
我們想想.
正常一道門是怎樣的.
怎樣叫正常的門呢.
就夠一個人過.
就叫正常的門.
不過這段聖經裡說的門寬不是那麼簡單.
是怎樣的呢.
是很寬的.
這裡所說的門寬是說一些很大的門.
大到怎樣呢.
有足夠的空間符合我們一般人的期望.
甚至這道門是足以讓一會兒很多人進入的時候.
其實很多時候說城門 電雨的出入口那些大門.
就好像現在這幅圖.
這幅圖就是大馬士革裡面其中一個的門.
是一個城門裡面的門是很寬的.
足以幾個人同一時間進去.
這個就是在耶路撒冷裡面的一個門.
所以門寬是這樣.
所以當你想想道門要這麼寬.
你去道門那條路都要寬一點的.
要寬很多.
而且因為很方便.
所以很多人都會走這條路的.
原來A的答案就是在說這件事.
門寬路大人多.
就是代表什麼呢.
很多人都會用.
很多人都會選擇.
今天原來我們有很多時候都是這樣的.
做事就是大家都這麼做.
有什麼問題.
個個都是這麼做.
那我就跟著做.
就算犯法.
個個都是這樣.
那我就照做.
所以你知道在加州裡面.
就是那個之前我們常常提到的法例.

$^{281}$就是什麼.
個個都去搶劫.
因為搶劫都不用拘捕.
搶了都沒有問題.
所以人人都這麼做.
那我就這麼做.
越做的人就越多人參與.
大家都是這樣 有什麼問題.
有些人又說.
賺錢最重要的.
信仰可以慢慢來.
個個都想著賺錢.
那為什麼我不想著賺錢.
個個都要發大財.
為什麼我不可以努力去賺很多很多的錢.
信仰可以放在一旁.
這些都是很多人的選擇.
有些人說.
不用太認真.
做人輕輕鬆鬆.
鄧子豐.
聽過嗎.
沒聽過.
聽過輕輕鬆鬆鄧子豐.
就是什麼.
一天飲飲食食就過一天.
就是說.
有時候都不用做事那麼認真.
隨便隨便懶懶閒閒.
只吃一頓就吃一頓.
這些都是很多人做的方法.
很多人做事的方法都是這樣.
原來有時候我們都會有這些做事方法.
可以讓我們跟隨世俗的做法.
跟隨大隊的做法.
跟隨最多人的做法.
就好像我們走最多人去的路.
選一條最寬的路.
選到門是最寬的.
不過剛才我們看到.

$^{321}$可能給你提示.
選答案的時候.
可能會通往滅亡.
通往滅亡.
第二條路B的答案是怎樣.
B的答案就是門窄.
路少人少.
因為通往生命的門是窄的.
路是少的找到的人也少.
正常的門大家剛才想像得到.
而窄的門是代表什麼.
就是窄到可能僅僅夠一個人出入的門.
甚至乎窄門其實就在說.
是一些很不顯眼.
很容易被忽略的地方.
一些的門.
正如有時候我們開車或者走路.
走過某個地方我們也不知道.
因為那個門口太小了.
也不知道那個是門口.
所以原來這個窄門.
也是這樣的意思.
而到門少是什麼意思.
到門少代表可能受到擠壓.
或者很狹窄.
狹窄到一個地步就是.
你要過就好像要被壓住一樣.
會很辛苦很痛苦.
那是什麼意思呢.
原來有時候我們做人的一些決定.
有些的路向有時候都很窄.
很艱難的.
不知道你覺不覺得你的人生是很艱難.
有時候我們也會面對這種艱難.
甚至乎要為神去受苦的.
這個似乎都是笑不滅.
而這些門往往是怎樣呢.
讓你看看在香港裡面.
不知道你有沒有去過這個地方.
香港裡面有個地方叫杜峰山.

$^{361}$杜峰山就有一道門.
這道門叫做生命門.
生命門如果你去過你就知道.
其中它有一門口.
就是剛剛夠一個人經過的.
門口裡面寫著兩句話.
歡樂行人多並無增樂.
窄門盡者少內有永生.
所以原來有時候有些路都是有點窄.
難走一點的艱苦一點的.
吃苦要多一點的.
付出的代價要多一點的.
但這條路會不會就是你找到永生.
找到生命的答案呢.
事實上在這個杜峰山上面.
經過這道門之後你會見到什麼呢.
你會見到前面豁然開朗.
正所謂柳暗花明.
又一串了.
經過了這些地方之後.
見到一個巨大的十字架.
成鳥就在眼前.
而這道門的另一面寫著什麼呢.
寫著你們往普天下去.
全福音給萬民聽.
當你願意走窄路的話.
你就要承擔著往普天下去.
全福音給萬民聽的使命.
這個使命或許不容易.
或許很辛苦.
或許想要付上代價.
或許要你在不同的崗位裡面.
要做生命的見證.
被人挑剔.
被人挑戰.
被人挑機.
但這條路會不會就是你認為.
可以幫助你去找到生命的門呢.
要追求神的心意.
這裡說我們可能會變成獨行俠.

$^{401}$值不值得呢.
因為那裡人太少了.
和你的戰友可能都不會很多.
或者你的戰友會一個一個倒下.
但你要走一條容易的路.
還是走一條正確的路呢.
兩者你選誰呢.
在這個地方裡面.
你見到有人為了教會的原因.
有人為了基督身體的原因.
他願意堅持開教會.
維持著教會的運作.
但有時候那個地方未必容許你.
甚至乎要去坐監.
在監獄裡面要受一個長期的痛苦.
很多時候有些地方就是這樣.
就是你要坐監.
不知道什麼時候可以出監.
不知道什麼時候完結.
甚至乎你可能要被扣上手銬.
要被剃頭髮.
甚至乎你日漸消瘦.
沒有人能夠認得出來.
但窄的門.
就是怎樣呢.
能夠幫我們透過患難生忍耐.
忍耐生老練.
老練能夠生盼望.
我們能夠有一個更堅穩的盼望.
這位牧師他自己要進監獄裡面.
在監獄外面的教會弟兄姊妹怎樣呢.
他們嘗試去幫助牧師打官司.
去爭取他們應有的權利.
結果教會裡面好幾位律師.
都被吊銷牌照.
停止執業.
甚至乎要關門大吉.
一個一個的隊友似乎被打下來.
但你能不能夠繼續堅持你的呼召呢.
究竟堅持你的呼召的話.

$^{441}$可能這條不是人多的路.
不過這條路.
它另有一番新景象.
就是通往生命的門.
好了.
說了這麼久.
是時候決定了.
有福人生.
以下哪條路是通往生命的門呢.
兩個答案.
一是A.
一是B.
選什麼.
選B.
不單止選B.
你是不是選B.
你的生命取向你是不是選B.
是不是.
好 看看對不對.
好 答案是B.
門窄 路少 人少.
這條路是通往生命的路.
你剛剛選了一條通往生命的路.
恭喜你.
你成功過到第一關.
慢慢你應該都開始越來越接近大獎的時刻了.
好.
第二條題目來了.
我們看看第二條題目.
看你能不能夠繼續過關.
第二條題目.
怎樣能夠分辨出.
哪一棵樹才是好樹呢.
怎樣能夠分辨出以下哪一棵樹.
才會是好樹呢.
分辨的方法A.
結壞果子的樹就是好樹了.
B.
結好果子的樹就是好樹了.
C.

$^{481}$不結果子的樹就是好樹了.
D.
荊棘和質禮就是好樹了.
能不能想到.
能不能決定你的答案.
B.
真的是B.
好.
既然你有這樣的想法.
我們又用錦囊.
我們看看今天開什麼錦囊.
開錦囊問現場觀眾.
哪個現場觀眾想回答.
改聲一下你為什麼會選擇這個答案.
選擇B啊.
真的嗎.
真的覺得是B.
等等.
我收到WATCHER.
有警告.
警告.
他告訴我.
我知道什麼事了.
看看發生什麼事.
他說.
你們要防備假先知.
他們到你們這裡來.
外面披著羊皮.
裡面卻是殘暴的狼.
不是每一個稱呼我.
主啊主啊的人都能進天國.
唯有遵行我天父旨意的人才能進去.
在那一天.
必有許多人對我說.
主啊主啊.
我們不是奉你的名傳道.
奉你的名趕鬼.
奉你的名行許多義能嗎.
我要向他們宣告.
我從來不認識你們.

$^{521}$你們這些作惡的人.
吸我走開.
這個警告告訴我們.
原來現場觀眾裡面有些人有黑心的.
原來問現場觀眾目的是什麼.
是拿貼士.
但是原來這裡有現場觀眾喜歡拿貼士.
知道什麼是拿貼士嗎.
就是毒死你的那些.
原來參加這些問答比賽.
現場的觀眾有些人是想你輸的.
因為你快點輸了出局.
他就有機會被抽在現場裡面.
成為下一個答問題的人.
所以有些人說答案給你聽.
不一定是對的.
甚至乎他知道真正的答案.
他會說假的答案給你聽.
我們叫「fake」他.
給你山埃貼士.
你死了他就開心了.
原來有些人是這樣的.
原來這裡也提醒我們.
當我們玩這個人生有福的.
有福人生的問答比賽時也要留意.
原來有些人給你山埃貼士.
那些人是什麼人呢.
這裡說那些人是什麼.
假先知.
就是說跟你坐在一起.
大家跟你在一起教會生活.
但那些人掛著先知的名字.
但其實是假的先知.
就是說他可能掛著一個基督徒的名字.
但可能是假的基督徒.
他帶著一個愛主的面容.
但其實他是假愛主.
很多人都是假假的.
那些人是想陷你一鍋.
所以這裡提醒我們.

$^{561}$原來也有這樣的情況.
甚至這些人是怎樣的特徵呢.
就是這樣.
當你以為你是一隻羊.
你披著狼皮的羊進入這個社會裡.
但這個社會裡有很多披著羊皮的狼.
就是說你以為你懂得賺錢.
懂得在這個世界玩.
裝作自己好像狼一樣.
但其實外面有更多懂得玩的.
他是一隻狼.
他裝作羊混在我們當中.
那究竟他是羊還是狼呢.
有時候也很難去分辨.
所以這裡原來也是這樣說的.
因為原來很多人都知道.
有時候先知都會披著像羊皮一樣的東西去保暖.
去比較簡約的生活.
原來有些人裝作先知都會披著羊皮.
所以你要知道這個世界有很多人會先知你一件事.
最近我也試過以為很安心都被人先知了一個.
為什麼呢.
因為我去eBay購物.
我去eBay購物以為eBay都是大商戶.
應該做很多審慎的篩選.
結果那些eBay的騙徒.
他就是在這個eBay賺錢騙你的錢.
話說我去購物.
就是上網購物.
見到那個東西的價錢.
你當一個正常的價錢是十元.
見到他說.
因為可能他未有些評論.
可能剛剛進來做生意.
為了吸引更多人購物.
好像薄利多斯.
他說給你六折.
給六折大約這樣的概念.
見到這麼划算.
可能真的想推銷一下.

$^{601}$做好一個商譽.
我就以為是真的.
這麼划算就買了.
買完之後.
他馬上很快幫我發貨.
他馬上發貨就由日本寄到美國.
我就很心急地等他的貨來了.
來到美國.
結果看著追蹤電郵.
因為有追蹤地址.
如果你買了一些貨物.
追蹤追蹤.
追蹤到不知道美國哪個州.
我都不知道那個叫什麼州.
接著說貨送到了.
那個州不是科華達.
他說貨送到了.
原來網上有些人這樣的方法.
叫假追蹤.
就是說他故意送一些無謂的東西.
送到一個不是你的地方的地方.
而是去假裝他的貨已經送了.
因為如果在eBay裡面.
你的貨已經成功送了.
那他就可以拿錢.
但其實你又追不到那件貨.
因為那個地址不是你的.
他就騙了幾百個人.
我見到那個評論裡面.
個個都說 scammer scammer.
接著我都寫 scammer.
因為我都被騙.
你打算很安心去買東西.
原來很多騙徒.
在教會裡面你都打算很安心.
但是我們要生命成熟.
才可以避免.
原來教會裡面都可能有些假才知道.
教會那些叫什麼.
瞞徒.

$^{641}$是不是.
教會有很多瞞徒的.
那些瞞著你.
所以要小心.
小心.
不要被人騙.
所以待會你選答案要小心.
那怎樣去選得好呢.
這裡有方法的.
待會幫你教你找答案.
方法是什麼呢.
耶穌怎麼看人呢.
耶穌說不要只看外表.
不要只看著他披著羊皮.
就當他是羊.
他說要看裡面.
就是沒有遵行天父的旨意.
原來有方法的.
分別是不是假才知道.
不是只看外表.
有敬虔的外表.
就好像民事法理賽人.
都很敬虔的外表.
不過要看裡面.
有沒有遵行天父的旨意.
有沒有愛神愛人的表現呢.
你記得在聖經裡有些民事法理賽人說.
我都做了奉獻了.
我就不用照顧父母了.
是不是.
那就有問題了.
所以耶穌提醒我們.
原來不要只看外表.
要看裡面.
所以幫到你了.
耶穌的提示.
看著也是這樣的.
不要只看外表.
還要看什麼.
看裡面.

$^{681}$看樹要看什麼.
有沒有結果指.
如果耶穌說看人.
就不要看外表.
要看裡面.
如果耶穌說看樹.
不可以看外表.
要看裡面.
如果耶穌看你.
耶穌是不是只看你外面.
耶穌看不看你裡面.
所以耶穌怎麼看你.
你怎麼看自己呢.
你給人什麼印象呢.
在你的生活裡有沒有行天賦的旨意呢.
你站出來像不像有好果子呢.
或許表面沒那麼容易分得出.
時間一來.
看到你有沒有果子.
就知道你的本質是什麼.
所以很重要的.
這三樣東西是不是好樹.
第一個荊棘.
第二個是疾泥.
荊棘和疾泥其實都很相似.
都是那些東西.
其實三個是什麼.
髮毛的橙.
這三個是不是好樹.
不是.
所以既然是這樣.
你應該都會開始找到答案了.
以下哪個答案能夠幫你分辨出哪棵樹是好樹呢.
答案是什麼.
肯定.
那怎樣分辨出你是不是一個好基督徒.
要什麼.
你被人看到有沒有好的果子.
是不是.
是不是遵行神的旨意.

$^{721}$是不是好的果子.
一眼看得出其實都騙不了多久.
騙了一會騙不了多久.
所以看看我們的答案對不對.
恭喜你你又猜中了.
又答中了.
那你現在又晉級了.
接下來我們要去挑戰最後最難的題目了.
準備好了嗎.
最後最難的題目.
第三條問題.
看看是什麼問題.
他說以下哪一種根基的房屋是最穩陣呢.
哪一種的房屋是最穩陣呢.
A 蓋在盆石上面.
B 蓋在沙土上面.
C 蓋在大海上面.
D 蓋在月球上面.
哪一種的房屋是最穩陣呢.
A.
很心急很肯定.
不過不用心急.
我們有錦囊未用.
還有誰的錦囊未用.
打電話給親友.
好我們用打電話給親友的錦囊.
打給我們的親友.
好我們就試試打電話.
我試試打電話看看聽不聽到.
喂聽不聽到.
聽到了是吧.
我現在問你一條問題.
我想你幫我解答一下.
我就是想問你以下哪一種的根基的房屋是最穩陣呢.
凡聽了我這些話又去做的.
好比一個聰明人.
把房子蓋在盆石上.
風吹雨打水沖.
撞擊那房子.
房子重不倒頭.

$^{761}$因為根基建立在盆石上.
凡聽了我這些話而不去做的.
比一個無知的人.
把房子蓋在沙土上.
風吹雨打水沖.
撞擊那房子.
房子就倒塌了.
並且倒塌得很厲害.
所以耶穌剛才給我的提示.
我應該找到答案嗎.
應該找到答案吧.
所以我試試想想耶穌給我們有什麼提示.
他說原來有不同的根基.
他說剛才耶穌提醒我們.
原來凡聽耶穌剛剛跟我們說的提示.
應該找到答案了.
因為他說一個聰明的人.
會把房子蓋在盆石上.
所以風吹雨打水沖.
撞擊那房子.
都不會倒塌.
因為根基建立在盆石上.
所以早前在香港有大火.
其中有7棟樓被火燒.
但燒完之後.
那棟樓有沒有倒塌下來.
沒有.
雖然受到一些損毀.
但是很安全.
原來因為香港的樓是怎樣建立的.
要打樁的.
要有一部打樁機打了很多的地樁.
所以不能短樁.
如果短樁的話.
那棟樓就危險.
所以打得穩.
打得樁.
打得多的話.
那整棟樓在上面建上去的.
就安全了.

$^{801}$香港二十多三十層樓都沒有問題.
又不怕風.
又不怕雨.
又不怕水.
所以很安全.
所以這裡可能都告訴我們.
待會的答案是什麼房屋最安全.
但是原來有些房屋是不同的.
有些不是香港的樓.
是什麼樓呢.
是美國的樓.
現在的房屋可能都是這樣的.
這裡說原來洪水襲擊的時候.
這裡曾經有一宗新聞說到.
美國的新墨西哥州.
是整棟房屋沖走了.
有沒有發生的.
都有發生的.
特別如果去到五級.
五級的時候.
你的房子可能都會沒有了屋頂.
可能沒有了.
所以呢.
哪一類房屋最安全呢.
是不是心裡有數.
哪一類房屋最不安全.
是不是心裡有數.
但是心裡有數.
跟你聽不聽一回事.
你心裡有數.
就知道哪一類房屋安全.
哪一類房屋不安全.
不過可能有些人仍然覺得.
我做的選擇是最好的.
所以你即使知道.
很多時候你沒有地基.
沒有打樁.
沒有一個穩固根基的房子.
是很容易催走的時候.
但是有些人仍然都是用一些.

$^{841}$自己認為好的方法去做.
剛剛最近上網都看到.
有一個家庭特意在自己家門前.
有一塊土地.
就這樣用一架鏟泥車鏟平它.
接著怎樣呢.
就在中國淘寶買了一間屋回來.
這間屋是剛剛一個貨櫃的長度的屋.
不知道你見過這些屋嗎.
這些叫摺疊屋.
剛剛一個貨櫃的長度.
接著當你運到美國的時候.
只需要在貨櫃裡面整個抽出來.
接著放到地面.
接著就打開盒子.
打開底部.
打開側門.
接著中間.
不是側門.
先打開上面.
接著在下面有道牆.
你就推上去.
推上去就頂住上面那塊.
這樣就可以了.
明白我的意思吧.
就是這樣.
所以你見到全部都好像右邊那樣.
就推來頂住頂住.
那就是一間屋了.
這間屋弄出來就是這樣的.
多少錢呢.
就是整個包在這裡.
就是來到美國整個費用多少錢呢.
兩萬五.
這麼貴是一萬多塊的.
加上運費又不聽.
還要留心聽聽.
加上運費.
加上要搬書.
加上要找一輛車幫你搭建這間屋.

$^{881}$兩萬五好像很划算.
不過他買完這間屋之後.
他發現很多問題.
因為他整間屋本身的地基又不平.
所以當他放了這間屋的時候.
又不穩固.
原來那些位置因為地不平.
因為要翻出來.
翻出來再穩固的時候本身別人的設計是坑坑好的.
就因為他說地不平.
有些就穩固了又不能弄.
他要強行弄到穩固了門邊的塑膠.
結果弄到這間屋就真的不是那麼理想.
如果大風大雨.
你覺得這間屋能不能頂住.
所以你要知道.
怎樣的根基是好一些.
似乎你都有方向了.
對於剛才耶穌給你的提醒.
不過重要的問題就是問.
你聽不聽得進耶穌的說話呢.
耶穌說答案就給了你.
這個錦囊你用了.
你打了電話給耶穌.
但耶穌說原來有個情況就是.
有些人聽耶穌說.
凡聽了我這些話去做的.
有些人聽耶穌說.
有些人凡聽了耶穌的話就不去做.
有時候耶穌說了答案給我們聽.
我們聽就照聽.
做呢?.
又做不做?.
有時候又不做.
答案是顯而易見.
問題在於你信不信得過耶穌.
祂和你說的話.
你就聽祂的話.
然後去做呢?.
我們很多時候有個很有趣的現象就是.

$^{921}$我們聽耶穌的話.
但我們不聽耶穌的話.
我們聽耶穌說話.
但我們聽完之後就不聽耶穌的說話.
聽得明白嗎?.
是不是就是這樣?.
耶穌告訴我們.
有福人生的道理.
講就講得很明白.
透過五六七三章聖經.
講就講了怎樣成為一個有福的人生.
有福了有福了.
八福擺在我們面前.
但我們聽完耶穌的話.
有沒有真的聽祂的話呢?.
有福的人生就要相信.
你打電話給耶穌.
你打電話給祂.
你就要相信祂.
你不相信祂就不要打給祂.
你相信祂知道答案.
你就打給祂.
所以我們一定要以耶穌的話作為根基.
你又會不會要打電話給耶穌.
但又信不過耶穌的答案呢?.
在你人生很多的選擇裡面.
你認不認定真的透過耶穌跟你說的話.
你又去遵行的時候.
就是一個有福人生的表現呢?.
到了第三條問題.
最後決定你能不能夠成功得到大獎的問題.
以下哪一種根基的房屋最穩陣呢?.
決定A.
你們是否都願意相信耶穌給你的提示.
你是否確定耶穌說給你的提示一定是對的.
是.
那你願不願意之後跟著耶穌的提示去行.
願意.
看看耶穌給你的提示能不能拿到大獎.
開答案.

$^{961}$A.
建在盤石上面.
恭喜你們了.
今天經歷了三條有福人生問題.
你不是百萬富翁.
不過你成為有福之人.
因為你明白主的心意.
懂得去分辨如何讓自己生命能夠結出果子.
甚至乎就是讓你的根基在盤石上面.
在神的說話上面.
所以恭喜你.
你拿獎了.
相信耶穌有福了.
你今天就是有福之人.
你就是有福之人.
說到這裡.
我們這個問答比賽都差不多完結了.
不過在這個問答比賽完結的時候.
我想告訴你.
我這個主持就沒有什麼特別.
但剛才那個百萬富翁主持就很不同.
剛才那個百萬富翁主持是誰.
陳啟泰.
百萬富翁主持陳啟泰.
原來在這幾年裡患了抑鬱症.
他說這是他人生最最最黑暗的階段.
為什麼呢.
因為之前你知道.
他透過百萬富翁成為一個很當時得令的主持.
但是自從在疫情期間.
你知道這幾年有疫情.
在這幾年疫情期間.
他就患上抑鬱症.
因為他本來有很多工作要去做.
但是後來就因為這個疫情.
也有很多工作停止了.
他開始在家裡發呆.
然後嚴重失眠.
甚至乎要吃安眠藥.
每晚都要吃.

$^{1001}$沒有辦法消除安眠藥.
甚至乎他也因為當時是一個疫情.
所以他很怕外出.
他怕外出就會染病.
一染病就會被隔離.
一隔離他就覺得他會發瘋.
所以他長時間都委直寄意.
一方面他知道有問題.
他又睡不著.
很多問題他要看醫生.
或者要穩定的治療.
但是他又很怕外出.
甚至乎要躲在家裡.
覺得一外出就一定會有事.
一外出就會有問題.
當你天天又不外出的話.
那你的情緒會怎樣.
就會再有問題.
於是他的情緒就越來越低落.
甚至乎他太太都為了要陪伴她的丈夫.
陳啟泰他就停止了他的工作.
半年的時間.
一路有半年的時間.
一路都停止工作.
而太太就每天出門去買東西吃.
24小時.
一個星期七天去陪伴陳啟泰.
然後陳啟泰才覺得有一點安全感.
一個好像什麼都懂.
好像做一個百萬富翁的主持.
好像什麼都懂的人.
不過原來面對自己人生的低谷的時候.
在一個患病的處境裡面.
他都找不到最好的答案似乎.
所以是令到他很大的困擾.
但是他後來說到他怎樣能夠康復呢.
其中一個重要的狀況是什麼呢.
有三樣東西能夠幫到他後來康復的.
哪三樣東西呢.
第一就是靠藥物.

$^{1041}$因為他後來真的被鼓勵去看醫生.
真的吃一些抗抑鬱的藥.
慢慢吃一些藥幫他調整.
第二就是他的太太.
這裡說他有八個月的時間.
就是24小時每天七天的陪伴.
不斷的陪伴.
第三順養.
為什麼他能夠經歷順養呢.
原來報道裡面說.
原來他太太在他曾經沒信耶穌之前.
原來已經信了耶穌了.
他太太信了耶穌.
他甚至在教會裡面都有做思琴的工作.
一些的事奉.
他自己在信耶穌的日子裡面.
因為有一個穩定的信仰.
一個穩定的根基.
讓他太太有一個真正的盼望.
導致他能夠面對著丈夫患重病的時候.
都能夠陪伴.
以及度過其實他自己都艱難的歲月.
所以就是這樣.
他太太也說.
如果他沒有信仰.
我想我也會崩潰.
正正因為他有信仰.
又信耶穌.
所以他沒有崩潰.
能夠陪伴他的丈夫.
而這個陪伴的過程裡面.
也讓陳啟泰認識到.
主是會幫助他.
於是他就是一路一路慢慢去認識主.
信耶穌.
以致到他自己最經歷的是什麼呢.
他說.
我覺得主是給了我力量.
信仰絕對給了我力量.
幫到我有智慧去康復.

$^{1081}$他解釋原來他人生裡面.
百萬富翁.
他這個不是百萬富翁.
不過他做了一個好像比百萬富翁更重要的決定.
當他在人生的抉擇.
二選一.
究竟他要選死亡.
還是選擇生存呢.
因為當時他自己已經有輕生的念頭.
只不過他已經望到二十二樓.
即是他住二十幾樓.
他望到樓下之後.
他就退下房間.
但是其實他有輕生的念頭.
這個是他其中一個的選擇.
但是他也因為他太太.
因為他信耶穌.
慢慢信耶穌.
而他有另外一個選擇.
他的選擇就選了一個最有智慧.
最能夠幫到他的人生抉擇.
有福的人生其實就在於我們怎樣去選擇.
我們懂不懂得去選擇.
人生有很多問題.
可能這次這條題目選得對.
可能另外一條題目選得不對.
有時候每次都不錯.
選得不錯.
但有時候每次都一次又一次錯.
但是最重要的還是.
我們如果要成為一個有福的人生.
我們要選擇最終的決定.
是不是剛才我們所認定的.
是.
聽耶穌話.
跟著耶穌.
相信耶穌不會騙我們.
相信耶穌是給我們最好的.
我們願意就是這樣去學校耶穌.
求主幫助我們今天.

$^{1121}$藉著我們這個的有福人生大問答.
幫我們想一想.
我還要怎樣選擇.
第一說其實我選對了.
我不需要再猶疑.
我博取我現我的人生在耶穌上面.
求主幫助我們.
我們一起祈禱.
天父我們求你去幫助我們.
因為我們人生裡面.
我們真的實在很多無奈的時刻.
或許我們人生的路.
有時候太窄太少同行的人.
甚至有時候找不到那道門.
好像沒有出路.
大耶穌我們多謝你.
因為你告訴我們.
你就是道路真理和生命.
藉著你我們能夠去到那道門裡面.
就是那道永生的門裡面.
求主你幫助我們.
堅定我們所選擇的答案.
我們願意聽主你給我們生命的提示.
走得對走得正.
走一個你喜悅的結果.
讓我們認識你更多.
讓我們堅持一生的信仰.
求主你就是這樣幫助我們.
聽我們祈禱.
奉耶穌名字而求.
阿們.
\newpage



\section{出埃及記 14:2}
\label{sec:031996}
\textbf{聖經人聖經事:吩咐以色列人回轉}~~~~黃克勤牧師博士
\newline
\newline
連結: \href{http://www.church.com.hk/acms/content.asp?site=cdc\&op=show\&type=product\&code=031996\&layout=sermon}{\texttt{031996}} ~~~~ 語音日期: 2026-04-14
\newline
\newline
\hyperref[sec:031995]{\small{< < < PREV SERMON < < <}}
~
\hyperref[sec:index_chronic]{\small{[返順時目]}}
~
\hyperref[sec:031997]{\small{> > > NEXT SERMON > > >}}
\newline
\newline
出埃及記 14:2
\newline
\begin{longtable}{cl}
\hline
\hline
章節 & 經文 (和合本修訂版)\\
\hline
14:2 & \begin{tabularx}{0.7\textwidth}{X} 「你吩咐以色列人轉回，要在比‧哈希錄前面，密奪和海的中間，巴力‧洗分的前面安營。你們要在對面，靠近海邊安營。 \end{tabularx} \\ \\
[1ex]
\hline
\hline
\end{longtable}
$^{1}$各位電影姐妹 大家好.
又跟大家讀聖經.
今天讀這段聖經就有趣了.
是早幾天讀過的.
生命樹宣教網絡.
和你一起看聖經人 聖經事.
看看上帝背後.
在做很多工作的.
神是吩咐以色列人轉回.
神吩咐以色列人轉回.
不過這次轉回.
聖經時時都是說神叫人回轉.
回轉 回轉.
就不要這麼倔強.
不要離開神 要回轉.
但是這個轉回.
是說他們出埃及之後.
是走一條路.
繞回頭的意思.
這個轉回.
讀過的了.
跟大家再讀一下.
今天是看一個挺有趣的.
可能你沒有想過.
也挺重要 或者真的這樣發生.
而極可能是這樣發生的.
說了這麼大一輪 說什麼呢?.
留意一下經文.
《出埃及記》十四章第一節.
《耶和華曉遇摩西說》.
第一點要跟你們說的.
你吩咐以色列人轉回.
安寧在比哈希祿前.
就是蜜奪和海的中間.
對著巴力洗芬.
靠近海邊安寧.
第二點.
第一就是你吩咐以色列人.
要轉回 繞圈.
不是心回轉.

$^{41}$是他們出埃及條路的路線.
走著走著就要繞圈.
《耶和華曉遇摩西說》.
你吩咐以色列人轉回.
去一個地方紮營.
停在那裡.
停在哪裡呢?.
安寧在比哈希祿前.
蜜奪和海的中間.
有三個地方.
有一個叫比哈希祿.
一個叫蜜奪.
另外在海的中間.
對著海的對面.
就叫巴力洗芬.
是靠近海邊停下.
三個地方中間.
哪裡呢?.
比哈希祿.
蜜奪和海中間.
那裡安寧.
紮營.
這裡是靠近海邊的.
而海的對面.
就是巴力洗芬這個地方.
你會知道.
有幾個神叫他們這樣做.
在海邊安寧.
在海邊安寧有什麼好處呢?.
不知道.
神有他的心意.
有時候你會發現.
神為什麼帶我們走路.
走著走著.
叫我們走回頭的呢?.
很多基督徒是這樣想的.
就是.
這件事有沒有神的心意呢?.
祈禱了之後.
發現是神叫我這樣做.

$^{81}$我開始走路了.
一直以神叫我這樣走.
我走著走著.
那條路好像有一點不通.
就要轉一轉彎.
沒理由的.
神叫我向這個方向走.
一定是直走的.
直走直走.
走不了也要走.
頂住也要走.
他不轉彎的.
但是有時候.
神會帶我們走一條路.
遇到事轉一個彎.
像駕車一樣.
前面有人塞住.
有車塞住.
那怎樣呢?.
唯有塞得太老.
就繞另一條路.
可能現在原來有一條新路.
快一點.
不要緊.
就可以轉.
不過有些人認為.
上帝的心意是不轉的.
因為上帝的心意.
他叫我這樣做.
他一定是開通達的道路.
有時候神可能是帶我們轉彎.
有時候叫我們回國.
有時候叫我們爬山.
有時候叫我們下水.
不出奇的.
到了海邊安營危險.
因為如果你被人攻擊.
你沒得走.
是不是?.
背水一戰.

$^{121}$人家在那裡圍著你.
只是餓都餓死了.
又推你下海.
所以呢.
其實為什麼神叫我們會轉回.
走著走著會繞圈呢?.
第二.
為什麼走著走著.
會叫我們去這麼危險的地方.
在海邊安營呢?.
海邊好.
風景好.
又可以打雨.
哎呀不行的了.
在海邊安營危險.
就是要這樣做.
海這個字.
就是叫Yan.
希伯來文叫Yan.
特別為大家介紹一下.
好.
今天要講的.
就是在海邊安營.
然後呢.
好了.
法律就說.
哈哈哈哈.
我收到情報.
以色列人在地中繞迷了.
抗議將他們困住了.
他們在我們這裡埃及出去.
走著走著.
誰知道.
走著走著.
繞回頭.
不會的.
他們應該要走的.
不會繞回頭的.
繞回頭.
哈哈哈.

$^{161}$迷路了.
這樣的想法.
原來呢.
這個是上帝出的有敵之計.
好了.
我要使法律的心剛硬.
上帝說我會讓法律的心剛硬.
他就會追趕那些百姓.
我便在法律和他全軍身上得榮耀.
這一句呢.
就是我們做過一次的聖經人.
聖經事實是早幾天的了.
埃及人就知道我是耀華.
於是以色列人就這樣走了.
神做這件事是要在法律身上.
顯出我就是耀華.
我想我又說回這段講過了.
沒什麼說又說.
不是的.
因為今天要跟大家說一些很重要的事情.
有人就告訴埃及王.
這個是假料.
你知道一打仗就有很多假的信息.
他是逃跑.
他是投降.
他是怎樣.
他是很差.
虐待俘虜.
他是打敗仗.
他是打勝仗.
一個打死四十個.
很多這些的了.
有人就告訴法律.
這個以色列人是走啊.
跑啊.
逃跑啊.
法律還說逃跑啊.
哇 這麼巧.
法律和他的臣服.
就是向百姓變心.

$^{201}$百姓是指以色列.
本來我們放他走.
現在不要放他走.
為什麼呢.
我們容讓以色列人去.
不用他再服侍我們.
放二百萬的奴隸.
哎呀.
我們是不是瘋了.
我們為什麼會這樣做.
這麼傻.
不行.
我們起來追他.
法律就預備他的車輛.
帶領軍兵同去.
帶著六百輛特選的車.
我上次說就帶著六百輛.
特別的戰車.
特別的坦克.
六百輛特選的車.
當時是這樣圍著車.
上面就有人拉著姜繩.
在騎馬.
不是騎馬.
是鞭策馬.
那個車兵長.
車兵長就負責這個車.
有一個人在保護他等等.
通常以前的戰車.
馬拉著的戰車.
很多時候都有三個人.
或者四個人的.
有人在射箭.
有人做什麼.
輔助.
很多時候都是這樣.
這個是法律自己.
可能特選的車.
可能親衛隊那些.
特別是特種部隊強一點的.

$^{241}$埃及所有的車.
這裡要補充.
埃及這麼大.
怎麼會一下子叫到.
全國埃及的車一起出來.
其實應該是說.
他們省裡面調動到的所有戰車.
都帶了出來.
特別是自己六百輛特選的車一起.
而且每車都有車兵長.
就是說那裡都有很多人.
一起去追了.
看看這個圖.
我今天跟大家說的就是這裡.
好.
留意一下.
這個就是一個衛星圖.
大家看得出來.
我要跟大家講學.
先講埃及在哪裡.
埃及在這裡.
或者現在這個西奶半島.
西奶半島現在都是埃及的.
是到這裡這個三角位.
這個三角位以上一直上去就是以色列.
這個圖是今天的.
算是近期的.
這個加薩在這裡.
加薩地帶在這裡.
就是說菲律師的地方.
就是說現在的加薩地帶.
就是以色列.
哈馬斯組織在那裡.
這個是地中海.
埃及就是這裡.
這個埃及很大的.
尼羅河也很寬.
你看衛星.
因為拍得這麼寬.
水流下去.

$^{281}$就去到尼羅河三角洲.
於是河就會分了很多條出去.
有些地方就會積了一坨水.
好了 尼羅河三角洲.
尼羅河.
這個是紅海.
以色列人過紅海.
我們經常都在想.
以色列人過紅海在哪裡過呢?.
今天我就想跟大家研究一下.
以色列人過紅海在哪裡過.
於是我們想.
哇 當然是在這裡過了.
它是從這裡出來的.
大概蘭塞在這裡.
蘭塞在這個位置.
他們是在這裡走的.
現在一路出來.
就過了西賴山.
去到西賴半島.
那些山.
看不看得到這些深色一點.
是山嶺來的.
西賴山大概在這裡.
在這裡繞一繞.
然後最後出去.
就進入耶利哥.
大概是這樣.
紅海.
這裡出去就是印度洋.
下去印度洋.
紅海有兩條像角一樣.
兩個角.
那這裡.
這個彎.
就叫蘇伊士彎.
這個彎叫阿卡巴彎.
有兩個彎.
然後下去紅海.
這個位置當然守著這裡.

$^{321}$這個長長的.
現在蘇伊士運河就是這裡.
蘇伊士運河.
這裡有一個大了一點點的.
叫虎湖.
是一個湖來的.
但是那些水是苦的.
所以叫虎湖.
現在蘇伊士運河就是這樣.
在這裡開通.
經過虎湖.
然後進來蘇伊士灣.
然後這樣.
過紅海出去印度洋.
就不用繞了一個非洲.
那麼大路那麼遠.
省了很多路.
雖然蘇伊士運河這裡大概163公里.
163里.
163公里在這個位置.
那些船就在這裡走.
有點狹窄.
去到虎湖當然是寬一點.
好.
那我們在這裡想.
以色列人過紅海.
法律追兵.
剛才讀過經文.
我帶他們繞回頭.
然後在這個地方安逸.
在海邊.
在哪裡呢.
如果過紅海.
照理說應該最窄的紅海都在這裡.
那當然了.
由上面就是蘭塞.
出來梳割耳塘.
就來到這裡了.
在這裡過.
其實今天的發現.

$^{361}$很有可能不是的.
今天的發現.
很有可能以色列人.
是並沒有過紅海的.
你是不是異端來的.
以色列人.
我們一直都是聽.
摩西帶領以色列人過紅海.
你說沒有過紅海都可以.
他可能是的.
我帶大家看了這個圖.
尼羅河三角洲.
尼羅河.
埃及.
蘭塞這個位置.
地中海來的.
蘇巖市灣在這裡.
這裡是最短的.
最容易過的紅海.
但是如果你要知道.
如果河是這樣的.
已經很闊了.
已經一望無際了.
如果這樣也很闊了.
走不了的.
這樣的神蹟.
也是的.
但是似乎不是那裡過了.
今天快點看看.
看這個圖.
這個圖曾經在前幾天出過.
戈山地.
戈山地是很肥沃的.
就是尼羅河三角洲.
他們建了兩個城.
以色列奴隸.
被埃及逼著建.
一個叫蘭塞.
一個叫比東.
他們出發在疏國.

$^{401}$所以蘭塞的人又來疏國.
比東的人來疏國.
周圍做奴隸的都來疏國.
他們就在那裡集合.
然後一起出發.
但是路可能是這樣走的.
來到這裡伊唐.
他們第一站到伊唐.
伊唐已經南面了.
這裡是紅海.
這裡有個字寫著.
小小一個叫虎湖.
這裡很窄的.
於是他們就去到.
比哈希陸和密奪中間.
河海邊那裡安營.
這三個位置.
海 比哈希陸.
和密奪三個位置的安營.
對著就是比利巴利西芬.
就是這樣.
於是他們的安營就在這裡.
其實他們已經下來了南面.
然後轉回回到北.
所以令到埃強覺得.
哈哈哈哈.
走迷了路了.
我們去殺他.
還有一些假消息.
以色列人在逃跑.
亂了.
於是他們就在戈山地.
跑出來追他們.
就是這樣的.
大概是這樣.
然後回到上面.
然後安營.
然後過紅海.
西拉半島.
然後就進入耶利哥.

$^{441}$去到應許之地.
這個就是一般人.
想的地方來的.
這個虎湖.
現在你看是通的.
不過今天.
經過了多久呢.
摩西大以色列人出安.
大概是主耶穌之前.
公元前的一千三百年.
哇.
一千三百年才到公元.
再加上現在2022年.
即是再加上2000年.
即是說.
三千三百年.
或者三千四百年了.
已經.
那些地貌有沒有改變了呢.
改變了.
那如果今天呢.
看回這個圖就這樣了.
已經.
那麼這裡的虎湖.
已經不再是黏著了.
如果看以前呢.
那個虎湖是黏著這個.
蘇爾士灣的.
那麼現在這裡已經斷了.
乾了.
變了濃地了.
那個虎湖呢.
是零零丁丁很小了.
就變成這樣了.
虎湖.
在這裡一個了.
是嗎.
這裡就很幼了.
以前呢.
很有可能是這樣的.

$^{481}$又或者其實是這樣的.
不過呢.
是以前的地圖畫得不對.
好了.
又看回這裡.
啊.
看這一張圖呢.
就有趣了.
留意這一張圖.
很有可能是對的.
將它放大了地中海.
它們呢.
這個位置就是.
歌山.
它們出來的地方.
那它們出來呢.
就在埃及地出來.
那出來呢.
就這樣走.
然後呢.
不正常的路線.
它們走著走著.
應該通常是.
如果它們要去迦南地.
怎麼走呢.
它們在埃及出來呢.
就會這樣走.
所以呢.
疏割那些是在這附近的.
那些城是在這裡的.
走出來呢.
就會這樣走.
就最近就去到迦南地了.
那這一條路呢.
就是.
叫做.
非利士路.
這個非利士地方嘛.
非利士路.
非利士路呢.

$^{521}$就是它們不走.
因為《聖經》說.
上次都讀過的.
就是那裡是有埃及的士兵.
在那裡守著.
那摩西就沒有帶它們走.
神就沒有帶它們走那條路了.
沒有帶它們走這條.
最短的路去迦南.
而是帶它們呢.
是預備繞到下面的.
但是繞到下面之前呢.
竟然剛剛就繞了向北.
所以說它們走錯了路了.
應該是這樣走的.
原因就是.
這條路上面.
有很多埃及的軍營.
在這裡駐守.
因此如果走這裡呢.
雖然百姓是有刀劍.
但是呢.
就遇到埃及士兵在這裡守著.
它們都會怕可能會打仗.
於是就不要了.
沒有走.
走一條不正常的路.
繞回去上面.
去了這裡.
去這裡呢.
原來.
在這個現在的人的研究呢.
是發覺什麼呢.
原來剛才說的那些密奪.
比哈希陸.
原來是在這個海邊.
這個海邊.
那這個是什麼呢.
這個地中海來的.
不是紅海.

$^{561}$紅海是在這裡.
很低的紅海.
原來是地中海.
地中海這裡就是埃及了.
那地中海就進來.
這些這樣的位置.
所以這些位置呢.
就是鹹淡水交界.
這裡有一條河.
鹹淡水交界的地方.
是有很多盧偉在這裡.
所以這個海呢.
這一塊海灣叫盧偉海.
而比哈希陸密奪呢.
就在這些地方.
那因此他們.
他們過埃及追兵.
他們過紅海呢.
是在這裡過的原來.
過這裡.
這裡又叫紅海.
有一個問題.
這裡叫盧偉海的.
當時原因就是.
希伯來文呢.
「em」就是海.
那紅海囉.
如果你看.
我們初一及記.
讀過這幾次的經文呢.
只是講了一個字.
就是海.
有沒有講紅海.
是沒有的.
只是講他們過海.
鹹乾地.
沒有講紅海.
其實那些考古文獻呢.
就發覺那裡叫盧偉海.
em-souf.

$^{601}$souf呢就是盧偉.
em就是海.
em-souf呢.
應該英文譯做Read Sea.
Read就是盧偉.
很有可能就是.
這個字是讀錯了.
Red Sea.
而紅海.
Red Sea這個字呢.
是去到.
很後期.
主前.
有叫紅海尾.
後尾才叫.
所以他們是不會.
經過紅海的.
他們可能經過叫em-souf.
em-souf呢.
意思就是Red Sea.
Read Sea.
可能.
在這個翻譯.
或者記錄抄寫的時候.
少了一個e字.
變成Red Sea.
其實是過盧偉海.
盧偉海就是這個地方.
那用不著過.
走這邊就可以了.
但如果有士兵跑來追你.
那隊兵呢.
你這樣走.
就會久了.
就會被他們追到了.
在這裡我直接過.
比較快一點.
而還有一件事.
就是上帝要引他們下海.
引他們下海.

$^{641}$而這個位是虎海.
虎湖記得嗎.
其實這樣一直連下去.
有一條河.
這條河呢.
就有一些水等等.
鹹淡水交界的地方.
特別多盧偉.
所以他們是在這裡過的.
可能是這樣.
就追過來引他們下去.
那這裡呢.
要是一個都容易過島.
那原來呢.
蘇維士運河每一年裡面.
都要好幾次.
是有好大的東風.
不知為甚麼會有這麼大的東風.
是會吹到那些.
蘇維士運河要封閉.
那些水是吹開的.
於是呢.
就好像過紅海那樣.
吹開的.
那些水是會淺了的.
於是呢.
風壓好大令到那些水.
是會分開是會淺一點.
那些船是不能行的.
就看那些書.
看那些文件就這樣講.
因此他們可能.
過的是這裡.
盧偉海而不是下面的.
拉斯洪海.
大家知道呢.
他們行這條正常的路.
是最好的.
因為他們是埃及士兵的.
駐防地.

$^{681}$所以他們不想行.
以免出事打仗.
所以他們就帶他們行這裡.
兜 是為了甚麼呢.
是為了過紅海.
令到他們信服.
上帝的大力與摩西.
另外是引埃及的軍兵.
落了去.
在他們身上得榮耀.
然後呢過完這裡.
他們再落下面西冷半島.
又看看.
這個圖的字.
可能不太清楚了.
埃及地.
在蘭塞出來.
蘭塞出來呢.
蘭塞出來 比東出來.
其他地方的人都出來.
就在疏角這個地方集合.
疏角就行去伊唐.
伊唐應該在下面.
不過呢如果看這個圖.
伊唐在這裡.
他們一齊出發了.
怎知他們就純帶他們.
兜上去密奪.
密奪比哈希盧.
對岸呢.
就是巴力西芬.
原來在這裡過的.
於是埃及的士兵就追落這裡.
在這裡染死了.
被海水染死.
他們在這裡過.
如果這個海呢.
沿出去就是地中海.
這個位就是這裡過.
於是沿出去地中海.

$^{721}$他們就說.
聚集在這裡.
這個就是另一個想法.
因此呢.
我們很多時候就會想.
以色列人出埃及.
是摩西帶他們過紅海.
很有可能.
是這個Ritzy.
是在這個七十四譯本.
開始譯錯.
然後一直流傳下來.
都是錯了變成紅海.
但是在以色列出埃及的時候.
是沒有紅海這個名稱的.
是完全沒有Red Sea.
紅海這個名稱的.
紅海後尾才叫的.
那你知道呢.
變成了其實.
他們很有可能.
是沒有過紅海.
是過這裡的.
那今天很多人又說在紅海的海底.
看到有些潛水的.
在這裡考究.
就看到有些車輪.
以前戰車的車輪.
金屬.
沒有可能.
原因就是他們那個時候.
車輪很有可能是木製造的.
未去到.
做金屬的器具.
那麼厲害的時候.
可能做到一些刀劍.
鋤頭那些這樣的.
但是未厲害到可以做戰車.
鐵的金屬的車輪.
而且三千幾年浸在海水.

$^{761}$什麼都沒有了.
如果是木的車輪.
很有可能如果是陶瓷器.
還有機會見到一點.
因此.
他們不是在紅海那裡過.
不是在下面那個紅海.
蘇伊士灣的尖.
那裡過的,不是的.
應該是在盧偉海.
這個機會是比較大的.
好了.
這個就是想和大家.
今天裡面講這一段.
就是他們是在這裡.
過盧偉海.
即是說是地中海.
的一部份.
而不是下面南邊的紅海.
而今天這裡.
是變了蘇伊士運河.
他們是在.
尖端這裡過的.
我想和大家講一下這件事.
為什麼這樣講呢.
因為我後來想了一下.
算了,不要講紅海了.
有些爭議性.
不過我覺得給大家.
多一個資料都好.
盧偉海.
比下面那個.
現在叫做紅海.
的呢.
很有可能是更加真實的.
因此我個人是.
傾向信盧偉海.
的.
就是一個這樣的地方.
地中海的南.

$^{801}$的一小部份.
很窄的.
也有人認為是什麼呢.
為什麼水會突然間.
退去呢.
原因就是.
很有可能是因為.
以前那麼講,希臘.
有火山大爆發.
而造成地中海.
有海嘯.
海嘯呢,海水就會.
退去,然後才湧回來.
他們剛剛那個時間.
海水退去的時候.
他們就走過了.
然後埃及軍兵追下去的時候.
就是海水回來匯流的時候.
把他們沖走.
由於這裡呢.
盧偉海不是很深.
怎麼可以浸死.
埃及軍兵.
如果他是海嘯匯流的話.
水就不用很深.
不過可以浸死人.
那麼你要二百萬人.
通過急急腳.
那麼很多人.
所以是要那些水.
全部退去才行.
怎麼會有一個大自然的環境.
那些水全部退去.
於是有人想就是海嘯.
那裡有海嘯呢.
有人研究就說.
很有可能是相同的時間.
希臘有海底.
有火山爆發.
然後造成海嘯.

$^{841}$海水退去.
埃及軍兵下去的時候.
海水匯流.
而後面這個火山.
也是沉降在海底裡面.
而剩下的部分.
變成一個島.
今天希臘有很多藍色白色的屋.
那個就是這個.
火山的一部分.
那麼講.
我今天都是講.
一些我能夠看到聽到的東西.
就不是我親自去過.
考察過.
甚至我沒有潛過海底.
見到圓形的東西.
好像車輪那樣.
我沒有見過.
不過我想今天跟大家分享.
有一個方法.
做它的示意.
而紅海這東西.
紅海.
Yen.
整段經文都是講海.
是沒有講過紅海.
而.
只是講海.
有些古卷就認為有.
是盧偉海.
譬如呂振宗.
譯本聖經.
他就會譯盧偉海.
就是這樣.
那其他的古卷.
都是講Yen Su.
盧偉海.
就不是下面的紅海.
很有可能是READ SEA.

$^{881}$後來就是.
拼錯了少了個字.
出現的.
好了今天跟大家分享到這裡.
不是很多屬靈的分享.
不過是給大家有一個新的概念.
其實.
很有可能不是經過.
素然是彎的紅海.
今日的紅海.
而是經過盧偉海.
READ SEA.
我們一起祈禱吧.
天父上帝雖然我們.
知道一些逐漸的.
更多考古的發掘.
發現到好像傾覆了我們的思想.
因為我們一路小時候.
聽起以色列人出埃及.
就是過紅海的.
於是我們周圍去紅海那裡找.
但是都不知道哪裡過.
原來今天我們越來越多的.
資料告訴我們.
很有可能就是這個盧偉海.
是地中海南邊的.
一個沼澤的地區.
主要是求你自己.
真是憐憫幫助我們.
帶著一個謙卑的心去等候.
可能有越來越多的.
認識主要無論哪個海.
而你是做了.
一個很大的神蹟.
救了以色列人.
又將埃及.
法律的軍兵.
染死在海裡面.
因為海水回流.
他們走乾路.

$^{921}$主要是我們讚美你.
求你自己親自加強我們的信心.
今天這些新的思考.
雖然是.
好像跟我們一路相信.
認識的是不同.
不過我們都帶著一個開放的心.
知道其實.
哪一個海都好.
上主仍然可以在那裡工作.
阿彌陀佛 耶穌名求 阿們.
\newpage



\section{歌羅西書 2:1-8}
\label{sec:031997}
\textbf{牧人的掙扎}~~~~羅錫堅牧師
\newline
\newline
連結: \href{http://www.church.com.hk/acms/content.asp?site=cdc\&op=show\&type=product\&code=031997\&layout=sermon}{\texttt{031997}} ~~~~ 語音日期: 2026-04-19
\newline
\newline
\hyperref[sec:031996]{\small{< < < PREV SERMON < < <}}
~
\hyperref[sec:index_chronic]{\small{[返順時目]}}
~
\hyperref[sec:031998]{\small{> > > NEXT SERMON > > >}}
\newline
\newline
歌羅西書 2:1-8
\newline
\begin{longtable}{cl}
\hline
\hline
章節 & 經文 (和合本修訂版)\\
\hline
2:1 & \begin{tabularx}{0.7\textwidth}{X} 我要你們知道，我為你們和老底嘉人，和所有沒有與我見過面的人，是何等地勤奮； \end{tabularx} \\ \\ \relax
2:2 & \begin{tabularx}{0.7\textwidth}{X} 為要使他們的心得安慰，因愛心互相聯絡，以致有從確實了解所產生的豐盛，好深知神的奧祕，就是基督； \end{tabularx} \\ \\ \relax
2:3 & \begin{tabularx}{0.7\textwidth}{X} 在他裡面蘊藏著一切智慧和知識。 \end{tabularx} \\ \\ \relax
2:4 & \begin{tabularx}{0.7\textwidth}{X} 我說這話，免得有人用花言巧語迷惑你們。 \end{tabularx} \\ \\ \relax
2:5 & \begin{tabularx}{0.7\textwidth}{X} 雖然我身體不在你們那裡，心卻與你們同在，很高興見你們循規蹈矩，對基督的信心也堅固。 \end{tabularx} \\ \\ \relax
2:6 & \begin{tabularx}{0.7\textwidth}{X} 既然你們接受了主基督耶穌，就要靠著他而生活， \end{tabularx} \\ \\ \relax
2:7 & \begin{tabularx}{0.7\textwidth}{X} 照著你們所領受的教導，在他裡面生根建造，信心堅固，充滿著感謝的心。 \end{tabularx} \\ \\ \relax
2:8 & \begin{tabularx}{0.7\textwidth}{X} 你們要謹慎，免得有人用他的哲學和虛空的廢話，不照著基督，而是照人間的傳統和世上粗淺的學說，把你們擄去。 \end{tabularx} \\ \\
[1ex]
\hline
\hline
\end{longtable}
$^{1}$好多信徒都好奇 究竟牧師每日是做什麼的呢.
我記得是2008年.
我和太太被邀請回來香港.
在一個堂會變牧會.
有一位信徒就跟我說.
牧師你今年五十五歲進入退休的年紀.
我看到你回來香港為了一份好工作 做牧師.
慢慢進入休息的境界 真的很聰明.
這位信徒是誰我就忘記了.
不過他說的話準不準確是另外一回事.
但是他對我的祝福和開心是我很感激的.
間中有信徒問我牧師星期日講道.
一至六他做什麼呢 是不是很空閒呢.
今天的信息就可以給你一些明白.
究竟一個牧人是做什麼的.
我們首先看看保羅寫這封信.
這段話在他的字旅行間你會感受到他的情況是怎麼樣.
他說我願意你們知道我為你們和在老底家的人.
以及所有沒有跟我見過面的人士這樣竭力奮鬥.
英文的翻譯就是struggle.
保羅為他的信徒 甚至一些沒有見過面.
但是因為他的信徒而信主的人.
struggle 竭力奮鬥.
保羅寫這封信的時候 他是坐牢的 在羅馬坐牢.
他第一次坐牢 第一次坐牢的環境是好一點.
他可以自己租一間房子.
不過每一刻他都有一個士兵和他的鍊子鎖在一起.
他是被囚禁的.
那個士兵是每六個小時換間的.
然後第二個士兵也是和他鎖在一起.
24小時他都是這樣.
究竟他這個奮鬥這個struggle是怎樣struggle呢.
他不是和士兵struggle 因為沒得struggle的.
他的struggle就是在禱告裡面 在牧養裡面禱告.
他為哥羅西老底家那些信徒有祈禱.
而且也靠牧養 他寫信牧養他們.
還有有些人可以探望他.
他有自己的房子 瑞士有人探望他.
有些信徒來探望他 他也牧養他們.
但是他就用這個竭力奮鬥這個struggle.

$^{41}$他的生活就是竭力奮鬥.
我可以說如果一個忠心的傳道人.
一個牧師 好像我們教會 江惠明牧師那樣.
是struggle 每一天都是struggle.
那保羅struggle是什麼呢.
他說是要為叫那些信徒心裡得著勉勵.
在愛中彼此聯繫.
可以得著憤怨性和確信而有的豐盛.
認識上帝的奧悲.
那如果我們畫一幅圖畫看.
勉勵的愚民的意思是.
就是將一個人帶到你的身邊.
來培育他 來鼓勵他.
英文就翻譯做encourage.
愛中結連 英文的翻譯就是need in love.
原文就是銷行 兩個銷行在一起.
你想像我和你的愛銷行在一起 黏在一起.
勉勵和愛心的聯繫.
就以致我們有悟性.
understanding 和有信心.
他們的心得著勉勵.
在愛中彼此聯繫.
可以得著憤怨性和確信而有的豐盛.
當你有悟性有信心的時候.
你就豐盛.
豐盛這個詞很奇妙.
你覺得信徒的生命有沒有豐盛.
沒有了 我有茶有米.
又怕沒有工作.
又怕自己身體不好.
又怕兒子女兒不孝順.
然後爸爸不給我錢我花.
這些東西.
但是原來保羅希望他的信徒.
透過悟性有信心.
以致得著豐盛.
最終的結果就是認識神的奧秘.
就是基督.
原來整個天地間的奧秘.
最大的奧秘就是基督.

$^{81}$因為一切的智慧知識的寶庫.
都隱藏在基督裡面.
我們當中有年輕人讀書.
有些人就進修.
年紀大的都仍然進修.
是為了知識.
你說我買股票.
希望有智慧買得好一點.
原來所有這些東西.
都在基督裡面.
有了基督有齊這些東西.
而保羅就想他的信徒.
能夠有這些東西.
從保羅的信裡面.
我們看到一個牧者的工作.
並不是星期天講的那麼簡單.
背後有很多東西.
做很多東西.
是神徒不知道的.
包括保羅所提的.
牧養和祈禱.
是為他們勞心勞力.
今天阿光牧師不在.
我就說說他就對了.
一個傳恩堂的堂主人.
他怎樣牧養的呢.
一方面他要帶領教會的發展.
譬如我們搬堂.
我們教會的大方向.
宣教等等.
不單止這樣.
他還做他日常的工作.
關心弟兄姊妹.
有些牧人照顧羊仔.
有些弟兄姊妹身體軟弱.
或者心靈憂傷.
他會親自去探望他們.
或者透過電話短訊問候他們.
有些弟兄姊妹.
就怕著牧師來找他們.

$^{121}$就避開.
所以這些情況都有出現.
而且阿光牧師也會.
和其他的教會和機構聯繫.
建立非常好的夥伴關係.
這個很難得.
我們傳恩堂不是很大.
但阿光牧師的人面很廣.
從柴灣到北角.
到整個播道中派裡面.
其實很多人都很喜歡和他一起.
很喜歡和他一起事奉.
他做樹人大學廣東話.
有幾十人的導師.
十幾間教會來和他一起合作.
他有聽過空襟.
其他的牧師說.
為什麼你們傳恩堂這麼厲害.
你不介意我們將樹人大學.
那些學生帶來我們教會嗎.
不來傳恩堂嗎.
他說無所謂.
去哪裡都不要緊.
每個教會都是上帝的.
你帶他去信主.
他來到你那裡就好了.
所以很多人都很欣賞他.
我記得使徒永富博士說過.
他和洪福堂的總裁.
他說阿光牧師叫我做什麼.
我都會幫他.
他真是得到人心.
這個是對我們整個傳恩堂.
都有好處的.
不單止在過去六年的時間.
他在播道環球差會.
做董事.
最後那四年做主席.
他現在就沒有做了.
去年年底海選.

$^{161}$他又退下來休息了.
做了六年的.
他不單止推動到中派的差傳工作.
而且將我們傳恩堂的差傳.
帶到一個新的境界.
所以我們其實很感恩.
他的人脈.
他的眼界.
他神給他不同的崗位.
令到我們整個傳恩堂.
發揮的影響力很大.
很多基督徒.
甚至有些傳道人都覺得.
一個牧師最好不要在外面搞那麼多.
在教會裡面四道牆.
怕牧師出去了.
就會做少了教會的事.
這個是非常之錯誤的.
其實江牧師.
或者一個中心的傳道人.
他要將他的堂會的頂子.
提升到宣教的境界.
為什麼我去台灣宣教呢.
就是因為台灣的教會跟我說.
羅牧師我們台灣的教會.
就太耐罪了.
看自己.
所以我們很多教會都很小間.
我們沒有一個宣教的觀念.
沒有一個廣闊的空間.
這個我們傳恩堂就沒有這個問題.
我想到我有五年的時間.
2017到2021那五年.
就在報道總會做總幹事.
總幹事就是CEO.
記得有些人不知道總幹事是什麼.
前任總幹事跟我說.
莫斯基牧師.
我女婿那時候跟我說.
爸爸為什麼你以前做黨主.

$^{201}$現在做了general secretary.
一個秘書.
而且還是general的.
普通的秘書.
其實又不是什麼意思.
general secretary總幹事.
其實是要負責整個中派的.
我的工作就是帶領中派的發展.
以及督導總會那30個同工.
但是江牧師做的事比我多.
他不單止帶領傳恩堂的發展.
他還照顧了我們幾十人的宿營生命.
就是一個牧人牧羊群.
好像耶穌基督對彼得華.
你餵養我的小羊.
而且牧羊就好像媽媽照顧嬰兒.
是很貼身而且很上心.
很牧人餵羊舍命.
我聽過一個報道是真事.
印度孟加拉有一個野生動物的保護區.
裡面有些野生動物老虎.
印度有老虎的.
但是有些非法獵人就設在陷阱裡面.
想抓老虎.
其中一隻老虎就在陷阱裡面.
讓牠受傷重傷腹部和後腿都流血.
本來要死的.
但這隻老虎不肯死.
牠用前腿一路爬一路捲.
把自己整個身體爬了幾公里.
去到那個洞裡.
因為那個洞裡有三隻老虎.
在等待媽媽餵奶.
怎麼知道有這件事呢.
就是在野生動物的保護區.
那個巡邏員他們開車周圍巡邏.
就見到一個非法陷阱.
陷阱旁邊有血跡一直拖著向前.
他們就跟著血跡.
去到幾公里以外那個洞裡.

$^{241}$就見到媽媽已經死了.
而那三隻老虎還捉著媽媽的身體.
喝完最後一次奶.
這個巡邏員就把他們.
這三隻老虎帶回去辦公室.
人工飼養.
養大了之後就放回去外面.
這隻老虎媽媽的最後掙扎.
救回自己三個孩子的性命.
這個就是一個牧人的掙扎.
你看保羅那封信的字旅行間.
你看加拿大書.
你看俄羅斯書.
伊弗索書.
你看他的字旅行間.
你看到他的那個苦難.
那個土局.
那個掛念.
他是很埋身的.
很想他的信徒得著勉勵.
能夠在愛中相結.
以致他能夠認識基督.
這個就是一個牧人的心.
所以耶穌說.
我是很牧人.
很牧人為陽寫明.
我們的光牧師.
大家知道了.
認識他這麼多年.
其實他很上心.
很埋身.
但因為這樣.
他就辛苦.
所以做一個傳道人.
一個忠心的傳道人.
就不是一份信功.
有些傳道人不是忠心的.
但光牧師是忠心的.
保羅在這裡.
他接著就提醒信徒.

$^{281}$我雖然不在你那裡.
但我心和你們一起.
光牧師今天不在我當中.
他的心都是和你一起.
我看見你們循規蹈矩.
並且對基督有堅定的信心.
就歡喜了.
循規蹈矩這個言文.
是一個軍事用語.
是discipline有紀律.
堅定的信心.
也是一個軍事用語.
勇敢.
所以保羅提醒的信徒.
你們要循規蹈矩.
堅定.
所以有一個軍隊的紀律.
是堅定不屈.
頂主媒我們要提醒.
我們信徒的時候.
不是很有軍隊的紀律.
我們星期日睡晚了.
就不回崇拜.
我們星期六早點睡覺.
我們又不回.
或者我們北上有些便宜東西買.
我們又從SUN走了.
有些頂尖應該休息.
不是今天不在的.
其實我們應該有軍隊的紀律.
現在我都提過.
我不知道為什麼.
回到香港十八年之後.
竟然有兩份工.
兩份full time的義工.
不知道什麼時候進入.
那位信徒跟我說的退休狀態.
但是這個怎麼做到.
因為我學到有軍隊的紀律.
就是discipline.

$^{321}$我七點多就回辦公室.
七點多.
這個就是discipline.
當我們這樣做的時候.
我們就能夠在基督裡面.
扎根建造.
信心堅定.
有這樣的心.
我們就能夠抵抗人的傳統.
世俗言論和哲學的兇談.
這個世界有很多的theory.
有很多的philosophy.
有很多的religion.
保羅就說.
免得有人不著著基督.
而著著人的傳統.
世俗的言論.
藉著哲學和偏人的兇談.
把你們擄去.
如果我們把保羅的話.
變回現在的處境.
我們就看到.
這些世俗言論.
哲學.
偏人的兇談.
仍然存在.
人民主義.
將人類的理性置於上帝之上.
看重人類的理性.
多於看重上帝.
律法主義.
遵循規則.
獲得職名地位.
而不是依靠恩典.
融合主義.
將基督教.
和其他宗教新時代哲學.
混合在一起.
這個對我們香港.
很多所謂福音派教會的信徒.

$^{361}$就是懶惰哲學.
玩.
玩樂.
躺平.
不思考.
玩樂就不用說了.
躺平就是什麼都不做.
總之就是.
休息多少就多少.
不思考這個就很嚴重.
人家說什麼你就信什麼.
沒有一個嚴肅的去想.
究竟這個所說的.
究竟這個社會發生的那件事.
是對還是不對.
我這個基督徒有什麼看法.
所以我對於政治.
對於文學.
哲學.
我都是很醒的.
很警醒.
因為我們做基督徒要醒的基督徒.
不是做個笨笨的基督徒.
如果基督徒不思考.
亂亂這個簡直是.
不明耀神的.
我是發光發燃的.
所以.
很多時候有些信徒.
我現在教神學.
色經學有些信徒.
都問我.
黃牧師你怎麼看這件事.
我有看法.
蔡玲一回說我有看法.
我有一個批判性的想法.
這個不是只有我才可以做.
每一個信徒都應該做.
當我們這樣的時候.
我們就能夠.

$^{401}$真是.
透過信心.
透過我們的.
understanding.
和faith.
我們的悟性.
我們認識基督.
也能夠防備世俗的哄談.
所以.
我們怎麼將這些.
今天的信息.
應用在我們現在的處境裡面.
有五點想跟大家分享.
第一.
我們對牧者有一個欣賞的態度.
欣賞我們牧者的事奉.
尊重神給牧者的屬命權利.
遵從牧者的教導.
如果你說.
教會的牧者是懶的.
是不盡忠.
是胡混混的.
那你不要幫他教.
全因堂不是這樣.
全因堂是江牧師領.
不是我領的.
有些人.
我昨天就碰見一個.
不見了十九年前的.
以前的一位弟兄.
以前那時候三十幾歲.
他說.
黃牧師你現在是不是.
在柴灣領教會.
我說不是.
我一直都是在做總會.
差會.
宣教士.
全因堂是江牧師領.
我的其中一個會友.

$^{441}$欣賞牧者的事奉.
尊重神給牧者的屬命權利.
我尊重江牧師.
他是我的牧者.
我是遵從他的教導.
以愛和其他肢體結合.
Need together.
蕭行這樣.
這個保羅教導我們.
努力紮根建道.
經歷基督.
保羅說.
我看萬事如糞土.
以經歷基督為至寶.
因為所有的知識.
所有的智慧.
所有的豐盛.
都是在基督裡面.
這個很厲害.
我下個禮拜.
都會講多一點這方面.
所以如果你不讀聖經.
你不經歷基督的話.
你就損失了.
你浪費了很多機會.
尚存感恩的心.
全因堂這個好的教育.
我們很欣賞執持.
他們按照他們的能力.
盡心來服侍主.
這個很好.
這是一個好的教育.
上帝將你放在這裡.
祝福你.
我希望你在這裡發揚光大.
叫教會.
叫整個宗派.
整個基督教都得益.
我們全因堂這麼小.
但我們有這麼多宣教基地.

$^{481}$支持這麼多宣教士.
這個很厲害.
你每奉獻一元.
在我們的宣教基金.
就會影響到全世界.
這個真的沒有浪費.
提防世俗的哲學.
和觀念.
世俗來了很多東西.
搞亂我們的思想.
所以如果你讀聖經.
在教會裡彼此相愛.
當你這樣的時候.
你就醒了.
你就在基督裡得到悟性.
信心.
你就處變不驚.
你就知道做一個怎樣的基督徒.
在這個混亂的世界.
這個變幻莫測的時代裡.
你能夠做一個真的有福氣的.
有堅定的喜樂的基督徒.
我們一起低頭祈禱.
親愛的主席.
我們多謝你.
我們為全因堂的弟兄姊妹感恩.
也為我們的目的感恩.
為每一位執事的奴婦感恩.
每一個信徒的師傅感恩.
我們知道每一件事都是上帝賜給我們.
我們在這個黑暗渤茂的世代.
這個末世的時候.
我們能夠彼此相愛.
我們有牧者來招待我們.
我們又學習去欣賞我們的牧者.
我們又彼此之間激勵.
求你祝福全因堂.
叫我們能夠在這個黑暗的世代裡做光做焰.
叫你的名透過我們得出榮耀.
我們祈禱奉耶穌基督的名.

\newpage



\section{列王記上 19:9-16}
\label{sec:031998}
\textbf{只剩下我一個人}~~~~文麗芳傳道
\newline
\newline
連結: \href{http://www.church.com.hk/acms/content.asp?site=cdc\&op=show\&type=product\&code=031998\&layout=sermon}{\texttt{031998}} ~~~~ 語音日期: 2026-04-19
\newline
\newline
\hyperref[sec:031997]{\small{< < < PREV SERMON < < <}}
~
\hyperref[sec:index_chronic]{\small{[返順時目]}}
~
\hyperref[sec:031999]{\small{> > > NEXT SERMON > > >}}
\newline
\newline
列王記上 19:9-16
\newline
\begin{longtable}{cl}
\hline
\hline
章節 & 經文 (和合本修訂版)\\
\hline
19:9 & \begin{tabularx}{0.7\textwidth}{X} 他在那裡進了一個洞，在洞中過夜。看哪，耶和華的話臨到他，說：「以利亞，你在這裡做甚麼？」 \end{tabularx} \\ \\ \relax
19:10 & \begin{tabularx}{0.7\textwidth}{X} 他說：「我為耶和華－萬軍之神大發熱心，因為以色列人背棄了你的約，毀壞了你的壇，用刀殺了你的先知，只剩下我一人，他們還要追殺我。」 \end{tabularx} \\ \\ \relax
19:11 & \begin{tabularx}{0.7\textwidth}{X} 耶和華說：「你出來站在山上，在耶和華面前。」看哪，耶和華從那裡經過。在耶和華面前有烈風大作，山崩石裂，耶和華卻不在風中；風後有地震，耶和華也不在其中； \end{tabularx} \\ \\ \relax
19:12 & \begin{tabularx}{0.7\textwidth}{X} 地震後有火，耶和華也不在火中；火以後，有輕微細小的聲音。 \end{tabularx} \\ \\ \relax
19:13 & \begin{tabularx}{0.7\textwidth}{X} 以利亞聽見，就用外衣蒙臉，出來站在洞口。聽啊，有聲音向他說：「以利亞，你在這裡做甚麼？」 \end{tabularx} \\ \\ \relax
19:14 & \begin{tabularx}{0.7\textwidth}{X} 他說：「我為耶和華－萬軍之神大發熱心，因為以色列人背棄了你的約，毀壞了你的壇，用刀殺了你的先知，只剩下我一人，他們還要追殺我。」 \end{tabularx} \\ \\ \relax
19:15 & \begin{tabularx}{0.7\textwidth}{X} 耶和華對他說：「去吧，從原路回去，往大馬士革的曠野去。到了那裡，你要膏哈薛作亞蘭王， \end{tabularx} \\ \\ \relax
19:16 & \begin{tabularx}{0.7\textwidth}{X} 又膏寧示的孫子耶戶作以色列王，並膏亞伯‧米何拉人沙法的兒子以利沙作先知接續你。 \end{tabularx} \\ \\
[1ex]
\hline
\hline
\end{longtable}
$^{1}$各位早晨.
我們開始的時候一起低頭祈禱.
我們一起低頭祈禱.
慈悲引領的天賦我們求你.
求你今天我們因你的名聚集在這裡的時候.
求主的靈充滿我們每一個人.
打開我們每一個的心.
讓使你為主你預備的廣章.
可以聽進每一個人的心裡面.
讓被主你的道教導.
也讓我們的生命得著轉化.
主你求你與我們同在.
我們這樣禱告.
奉主耶穌基督得勝的名祈求.
阿們.
今天的題目是只剩下我一個人.
是不是只剩下我一個人呢.
不知道你們有沒有一些寂寞的經驗呢.
記得有一次我有一個朋友.
我認識了很多年了.
一個姐妹她結了婚有兩個小朋友.
她就問我.
你沒有結婚.
其實沒有結婚的感覺是怎樣的呢.
情況是怎樣的呢.
我就跟她說.
有時候想找人陪也沒有.
也挺寂寞的.
她想了一會她就回覆我.
其實結了婚也是一樣.
不知道大家有沒有經歷過寂寞的感覺呢.
我相信在現在這個世界裡面.
寂寞也是大部分人共同的經驗.
寂寞似乎不是單純一個人.
就是我們就算身邊很多人.
我們都可能有這個寂寞的感覺.
我就去問AI.
人怎樣才會感到寂寞呢.
AI就告訴我.
寂寞是一種渴望的連結.

$^{41}$和實際擁有的連結之間出現的落差.
所產生的痛苦感受.
亦都是這樣說.
即使你身邊有人.
但如果你的說話只不過是說.
我們今天吃什麼飯.
天氣有多好.
你那種渴望被理解.
被關心的需要.
未得到滿足的時候.
你都會產生寂寞.
換句話說.
人與人的相處品質比數量更加重要.
香港有機構在二零二五到二六年.
做了一個調查.
叫做諾寧身心靈健康調查.
數據顯示有兒孫同住的長者.
他的生活的沉悶感是3.7分.
原來比獨居的長者高.
獨居的長者只有3.2分.
反映在實際上同住.
未必能夠轉化成心理上的陪伴.
另外在41.6\%受訪者中表示.
有心事沒有人跟他分享.
當中有五成都是由家人跟他同住.
有人同住.
有人跟他們一起都會感到很寂寞很孤單.
當然寂寞為什麼會出現.
原因很錯綜複雜.
可能是因為有些人覺得被團體排斥.
或者忽視.
或者覺得自己的感覺沒有被接納的時候.
寂寞的感覺就會很強烈地出現.
或者當面對人生很重大的轉變.
切斷原來很熟悉的連結的時候.
寂寞也會湧現.
或者感受到身邊的人不明白你的感受.
可能讓你覺得在人群中自己好像孤獨.
不知道你有沒有這樣的感覺.
曾經有這樣的感覺.

$^{81}$香港大學的研究發現.
如果人不斷在想自己有多孤獨.
他很容易會幻想抑鬱症.
他的意思就是說其實孤獨並不可怕.
可怕的就是你擔心你自己孤獨.
說了這麼多.
我今天想說什麼呢.
今天說的主題就是只剩下我一個人.
第五姊妹.
在你信仰的生活當中.
你有沒有感覺.
有時候感覺到剩下你一個人呢.
今天我說的經文就是.
《列王記上》十九章九到十六節.
是先知以利亞的故事.
以利亞是一個最差的年代的先知.
經人的背景是以色列分了南北兩國的時候.
當時猶大國是阿薩作王.
而以色列是阿哈作王.
在政治上.
阿哈是一個很有能力的王.
然而在宗教上.
阿哈做盡了耶和華眼中看為惡的事.
特別是他娶了一個西頓的公主耶西.
耶洗別為他的太太.
把外邦的偶像敬拜都一併帶入了以色列.
他自己也跪拜巴力.
他所作所為都激怒了上帝.
這是一個怎樣的年代呢.
我想這是一個很黑暗的年代.
表面上以色列國是一個強國.
國富家強.
但仔細去看的時候.
國內全部都是偶像敬拜.
滿載著不公義的事.
上帝並沒有無動於衷.
面對一心信巴力的阿哈王.
上帝妻派了先知以利亞.
向阿哈王宣告信息.
他這樣說.

$^{121}$我指著所事奉永生耶和華.
以色列的神起誓.
這幾年我若不禱告.
必不降露不下雨.
之後上帝就要以利亞.
到曠野裡躲起來.
他吩咐烏鴉供應食物給他.
直至到溪水都乾.
因為沒有下雨.
所以溪水都乾.
之後上帝就吩咐以利亞到西敦.
撒勒法那裡去.
住在一個寡婦的家.
經歷了神奇妙的供應.
並且令寡婦的兒子從死裡復活.
三年半之後.
當時以色列的旱災已經很嚴重.
上帝再吩咐以利亞去見阿哈王.
要在加密山決戰.
在眾以色列人的面前.
去測試誰才是真神.
巴力是真神.
還是耶和華才是真神.
這個故事相信大家都耳熟能詳.
如果你未聽過的話.
你回去看《列王記上》第十八和十九章.
裡面有精彩的記載.
我這裡就不會詳細去講.
結局是怎樣呢.
在全以色列人面前.
以利亞證明了耶和華是真神.
以利亞下令以色列人殺死巴力的先知.
神更加配合以利亞的禱告.
降雨在災旱的以色列地.
如果你是以利亞.
你這個時候的心情會是怎樣呢.
我相信一定是很高漲.
三年半以來.
神藉著他的口跟亞哈王說不下雨.
就是三年半都不下雨.

$^{161}$而他自己是得到神奇妙的供應.
他更加在眾以色列人面前.
擊敗了八百五十個.
巴力和阿舍拉的先知.
神就藉著他的禱告.
叫三年多都沒有下雨的.
以色列地下雨.
我想在以色列人的面前.
此時此刻.
以利亞是一個很觸目的人.
很偉大的人.
經文記載他很快就趕去耶穌列.
為什麼呢.
因為他很想看.
當亞哈把事情告訴耶穌別.
有關山上發生的事情的時候.
他很想知道有什麼反應.
我們可以想像.
他期望耶穌別一致很懼怕耶和華.
他可能會期望和亞哈一起去請教.
怎樣才可以歸向耶和華.
放棄巴力的敬拜.
他期望宗教復興在以色列國.
他期望以色列可以改變有革新.
但結果怎樣呢.
結果不像他的預期.
耶穌別還說要殺害他.
耶穌別說.
明日若在這個時候.
我若不使你的性命.
像那些人的性命一樣.
願神命重重地降伐於我.
發生的事和以利亞的預期.
很大的差別.
他知道之後很震驚.
他恐懼於是他就逃亡了.
他逃亡的行動除了出於恐懼和害怕.
我相信都因為他很失望.
他對加密山的事件很失望.
他對自己很失望.

$^{201}$他對神也很失望.
他以為事件可以讓亞哈回轉.
讓耶穌別回轉.
讓整個以色列人回轉.
但一切希望都幻滅.
得到的卻是一個死亡的恐嚇.
他要逃命.
他一直從以色列的耶穌列.
一直走走走.
走到別士巴.
你看看地圖走了多遠.
雖然以利亞要逃離耶穌別的追殺.
其實離開了以色列國就已經可以了.
但他一直走一直走.
走到150公里以外的別士巴.
而他的情形.
大概用了五天的路程才停下來.
他將僕人留在城裡.
自己就繼續在曠野裡多走一天的路程.
我們可以想像他有多害怕.
他有多失望.
當他的生命受到威脅的時候.
他感到他似乎毫無意義.
毫無希望.
他陷入情緒的低落.
墮入了一個很大的深淵裡.
他的信心被絕望掩蓋.
他覺得他做的事很徒勞無功.
他的事奉好像再沒有力量.
情感被耗盡.
他腦裡好像在說.
我沒有辦法繼續前行.
這超出我的能力.
我不能再這樣了.
死更好.
我一個人走到曠野裡.
在一棵羅藤樹下之下.
他就跟上帝說.
耶和華我罷了.
求你取我的性命.

$^{241}$因為我不勝於我的列祖.
我們可以想像.
如果以利亞這樣在敵人面前大獲全勝.
為神做了這麼大的事.
曾經在無數困難裡面.
靠賴耶和華可以面對.
可以克服的時候.
一個人應該是一個屬靈的偉人.
應該不會有靈性枯竿的時候.
不應該有意志消沉的一天.
但原來不是.
他也是一個很普通的人.
被人的條件所限制.
都因為恐懼而怯懦.
灰心和沮喪.
以利亞求死的祈禱.
更加好像向神發揮氣一樣.
他說我罷了就好像說算了.
事奉耶和華.
要忍受的我都受夠了.
我不要再受.
不要再要我事奉你了.
叫我去死好過了.
其實他渴望的可能不是死.
而是他不再想事奉.
他給神的理由是.
因為我不勝於我的列祖.
他沒有清楚跟誰比較.
但是我們知道.
他感覺到自己所做的徒勞無功.
對以色列國也沒有所做的.
也沒有引起什麼作用.
看來在以利亞的心裡面.
他認為.
有他心中預期的效果就是好.
沒有出現自己預期的效果就是失敗.
他判斷自己的事奉是沒有意義.
神的僕人他也不勝任.
當一個人沮喪的時候.
就很容易陷入自憐.

$^{281}$自憐就很容易叫人扮演受害者的角色.
究竟以利亞的心.
其實是為了滿足他自己.
還是滿足神的心呢.
我記得有一次講道.
聽到一個牧師這樣講.
他說他有一個牧者.
他因為事奉的效果不如理想而很沮喪.
他覺得自己很沒用.
牧師就跟他說.
如果你因為一個事奉的工作不成功.
就歸咎於自己的話.
當有一天.
你的事奉有美好的成果.
你就會將榮耀歸給你自己.
這個說話對我很大的提醒.
神交託了給我們做的每一件事.
只要我們盡心去做就好了.
結果怎樣就交給神.
或者上帝就是使用我們不完美的結果.
來成就祂的工作.
我們都不知道.
我們不需要用結果來判斷.
我們自己的事奉是否有成果.
更加不需要用結果來判斷.
別人的事奉有沒有成果.
有沒有意義.
以後以利亞身心疲勞疲避.
他就在羅騰樹下睡著了.
神知道以利亞的需要.
在以利亞最低層的時候.
神沒有跟他講道理.
也沒有斥責他沒有信心.
神就跟他說.
神也沒有跟他說你快點回去事奉.
神也沒有吩咐以利亞做什麼.
只是命令他要睡覺和吃飯.
睡覺和吃飯.
和睡覺和吃飯.
你的經文都是你看《電王記》就看到的.

$^{321}$神容許以利亞休息休息再休息.
當以利亞疲勞消除之後.
神就有一個指示給他.
就是要跟以利亞見面.
他引領以利亞走了四十天.
到了神的山就是荷列山.
別士巴和荷列山距離有四百公里.
平常人用十天就應該會到.
但以利亞用了四十天.
我想他不是用一個趕路的速度.
去走這個路.
而是帶著一個很憂心忡忡.
心事重重每幅惆悵地走.
這個時候他在想什麼呢.
之後他進了一個洞.
他就住在洞裡.
當時的以利亞仍然是心灰意冷.
失去了事奉的能力.
神知道以利亞的心情.
他也明白他.
神沒有責怪他信心軟弱.
他是鼓勵他要他再面對自己.
面對神.
在十九章九節裡面.
上帝跟他說.
以利亞你在這裡做什麼.
上帝究竟知不知道.
以利亞在恐懼裡戰鬥呢.
知不知道以利亞對未來的方向.
一點都沒有信心.
上帝問這個是為了什麼呢.
以利亞就把心裡的不快.
告訴神他向神投訴.
神在聽.
他說我為耶和華萬軍之神.
帶發熱心.
因為以色列人背棄了你的藥.
毀壞了你的壇.
用刀殺了你的先知.
只剩下我一個人.

$^{361}$他們還要尋索我的命.
由別士巴到荷列山那四十天.
以利亞得到一個結論.
就是只剩下我一個人.
他告訴上帝.
他孤身作戰.
很寂寞.
只剩下他一個人.
面對這個大惡人夜洗別.
神怎樣回應他呢.
神看到他沮喪.
看到他很自戀.
但是沒有放棄他.
神實在知道他的需要.
知道他身心疲勞.
要休息的時候.
要讓他休息和吃東西.
身體恢復了.
神就讓他能夠有安靜的時間.
有四十天這麼多.
經文沒有交代.
為什麼以利亞去荷列山.
是不是神引領他呢.
但是神卻是讓以利亞.
在這四十天裡面.
可以有思考.
過去發生的一切事的機會.
但是以利亞的回答.
似乎不是回答神的問題.
只是向神投訴自己的經歷.
自己的苦況.
但是神仍然沒有走.
他還在聽.
神沒有責怪他.
或者和他解釋什麼.
而是叫他.
你出來站在山上.
來到我面前.
以利亞從山腳的山洞裡走出來.
去到山上.

$^{401}$要和神見面.
那時候耶和華從哪裡經過.
在祂面前有烈風大作.
崩山碎石.
耶和華卻不在風中.
風後地震.
耶和華卻不在其中.
地震後有火.
耶和華也不在火中.
火後有微小的聲音.
神用一些自然界的情景.
象徵性地.
向以利亞顯明祂的心意.
好像很多其他人一樣.
以利亞可能只相信.
神會出現在一些.
很轟烈很戲劇性的場景裡.
當然上帝有時會這樣形式出現.
但是神更加經常出現的.
不是很轟烈.
不是很戲劇化的場景裡.
神就在安靜裡.
用微小的聲音.
和以利亞說話.
就好像告訴他.
人不一定在轟烈的作為裡.
才能見到神.
這才叫做事奉神.
細微沉靜的事情.
同樣都代表著神的同在.
從翻天覆地的經歷.
去到平淡無奇.
安靜的景況是很不容易的.
上帝是用微小的聲音出現.
這個聲音是很有力量的.
之前以利亞在山洞裡.
神就問以利亞.
你在這裡做什麼.
當以利亞聽到神微小的聲音.
和他說話的時候.

$^{441}$以利亞就用愛依夢想念.
出來站在洞口.
表示他對神的敬畏.
上帝用相同的問題.
再問以利亞一次.
以利亞你在這裡做什麼.
或者以利亞沒有足夠的時間.
但他領悟神想跟他講什麼訊息.
或者被他自己的情緒捆綁得太緊要.
以致他跟之前的回答一樣.
面對以利亞的抱怨.
上帝繼續的寬容他.
他要以利亞繼續的事奉他.
神將新的工作交託給他.
以利亞的情況就像我剛才所說.
他面對他的期望.
和實際的差異落差.
當以利亞反覆他自己去想的時候.
他覺得他自己很孤獨.
就好像跌入了抑鬱的境況裡面.
神卻要以利亞繼續的向前行.
知不知道寂寞的相反是什麼.
是連結.
神給以利亞新的工作.
是為了找到他自己同工的人.
耶和華跟他說.
你回去從曠野往大馬士革去.
到了那裡就要高哈舌作亞蘭王.
由高凌時的孫子耶和作以色列王.
並高阿伯米可拉人沙法的兒子.
以利沙作先知接族理.
神要以利亞回去.
不是回去他來的地方.
而是回去他事奉的路.
以利亞曾經拒絕.
想放啟示奉神.
現在神反而再呼召他.
叫他再踏上事奉的路.
他新的任務是要高立神揀選的人.
神的心意實在跟以利亞的心意不同.

$^{481}$以利亞以為神呼召他.
是叫阿哈王和耶洗別.
因為他所做的事認罪悔改歸向神.
但不是.
神已經預備了耶和作以色列王來接族.
來代替阿哈王.
以利亞以為自己孤單一個人去事奉神.
神他是為他預備一個徒弟以利沙.
將來繼承他的工作.
這些都不是以利亞想到的.
後來上帝更向他揭示一個重要的真理.
在十八節裡耶和說.
「但我在以色列人中.
為自己留下七千人.
是未曾向巴力屈實的.
未曾與巴力親嘴的」.
神是在告訴以利亞.
你不是孤單的.
神知道那些對他忠心的人.
除了以利亞之外.
在以色列境內還有七千個.
未曾向巴力屈實的人.
以利亞知道各宗還隱藏著.
這些對神忠心的人.
應該對於他的事奉有一股動力.
就像告訴他.
你不是孤單的.
以利亞你不是孤單的.
之後以利亞就重新回到自己事奉的崗位.
繼續完成神交託給他的工作.
弟兄姊妹.
以利亞今天講這個故事.
告訴我們什麼呢?.
當我們被過去的經驗.
或者慘痛的經歷捆綁著我們的時候.
孤單自然的感覺.
通常給我們發出一些負面的訊息.
但我們要知道.
我們不會只剩下我一個人.
有主和我們同在.

$^{521}$無論什麼的境況.
是高是低.
祂都和我們同在.
我們得力的泉源.
不是身邊的歡呼聲.
卻是神自己.
神在我們身上有祂的美意.
以往上帝可能用這種方法.
那種方法來引領我們.
去按照我們的禱告去施展大能.
或者轟轟烈烈地去使用我們.
這些美好的經驗我們不需要忘記.
但是我們不要讓這些經驗.
去摯礙著我們聽到上帝的帶領.
神的帶領有很多可能性.
就好像上帝早知道.
以利亞加密山上面的工作.
絲毫都不會影響.
阿哈和耶洗別的心.
上帝自有祂自己的心意和計劃.
神的計劃不一定按照我們的預期去實行.
我們算是什麼呢?.
我們要做的.
我們只是要自我降服.
聽從主的聲音.
信服基督的帶領.
特別是在這個充斥著很多噪音.
很混亂的世界裡面.
面對世界裡面很多的不公義.
令人灰心的事每一天都發生.
情況可能比以利亞的時代更加差.
聆聽主的聲音幫助我們不會迷失.
不會被這個世界的事情佔據.
這個世界太多的聲音.
我們裡面可能也有很多東西跟我們說.
寂寞孤單的感覺.
經常纏繞著我們.
但是我們需要記住.
不是上帝有沒有說話.
而是我們有沒有聆聽.

$^{561}$不只是外在的聆聽.
不只是外在的安靜.
更加是內在的安靜.
放下自己操縱的心.
專心仰望神.
聆聽神的帶領.
靠著主的力量跟隨主的腳蹤.
有一個這樣的故事.
有一個木匠.
他在工作室.
他不小心把他的錶掉進了一堆木舌裡面.
他大聲抱怨.
他說我倒霉了.
一直在找.
在撥開地上的木舌.
想找回自己的錶.
也有些人幫他帶提燈.
打開了眼燈照.
想找回.
但是找了半天都找不到.
他的兒子突然間悄悄地走進去工作室裡面.
一陣子他就找到錶了.
木匠很高興就問他.
為什麼你做得到的.
兒子就跟他說.
我只是靜靜地坐在地板上.
之後我就聽到滴答滴答的聲音.
我就知道錶在哪裡了.
當我們問.
怎樣才能聽到上帝的聲音.
通常不是上帝有沒有說話.
而是我們有沒有足夠的安靜.
去聽到上帝的聲音.
我們需要操練安靜.
我們要知道上帝是上帝.
我們不要只顧著埋怨.
顧著自憐.
我們要聽到上帝微小的聲音.
仰望著豐富的恩典.
才能確定自己事奉的心志.

$^{601}$不要做時代的受害者.
我們要繼續為神增展.
剛才說的寂寞什麼時候會產生呢.
人會感到寂寞.
往往不是因為沒有人在身邊.
而是因為沒有人真正的在場.
神一直都在場.
但神也期望我們在場.
當上帝問以利亞.
你在做什麼的時候.
是什麼意思呢.
難道神不知道.
以利亞在做什麼嗎.
我也有類似的經驗.
我經常都跟別人分享.
在我讀神學的時候開始.
上帝時常都擺了一條問題在我心裡.
就是問祂你在做什麼.
這條問題伴隨了.
我經歷了很多困難和迷失的時候.
可能是風光明媚的日子.
或者是死忍一居的時候.
神都提醒我要反省.
知道我為什麼要做一個牧者.
為什麼要做我現在正在做的事.
這條問題幫助我找回我自己的初心.
叫我不要計較果效是不是我所預期的.
或者不要去想我的努力是不是白費.
可能在世人眼中.
好像我所做的事都是微不足道.
或者徒勞無功.
但是當我時常思想.
我自己為什麼要做一個牧者的時候.
我就可以在神裡面重新得到力量.
繼續向前行.
事奉的路可能真的會讓我們感到孤單.
感到沒有人明白.
沒有人接納和重視.
但是當我們建立.
和神和人那一份的連結的時候.

$^{641}$我們就不會感到寂寞.
我也很想問大家.
你在這裡在做什麼.
你為什麼在這裡.
當你知道你在這裡在做什麼的時候.
你的行動就有意義.
弟兄姊妹.
神很多時候都為我們預備了痛恨者.
雖然可能獨來獨往是你的個性.
好像以利亞一樣.
我們可以想像.
如果你以利亞聽到耶穌要殺他的消息的時候.
他不是拔足逃亡.
而是和神那七千個中心的信徒.
連結在一起的時候.
故事會不會被改寫呢.
你記不記得敬畏神那個俄巴底.
你可以回去看看《聖經》.
他是阿哈王的家宰.
但是當耶穌要去殺害耶和華先知的時候.
他偷偷地收買了一百個.
並且在旱災的時候供養他們.
其實以利亞早就認識了俄巴底.
俄巴底跟他講過這件事.
如果他知道俄巴底和他同一陣線的時候.
以利亞知道不是剩下他一個人.
神已經為他預備了很多痛恨者.
只是以利亞看不到.
我們可以呢.
可能我們也是一樣.
神給了我們的家人.
我們的朋友.
我們教的弟兄姊妹.
相信在他們當中也為我們預備了痛恨者.
叫我們總不孤單.
神呼召人永遠不是呼召一個人.
總是呼召一群人.
或者你有時感到寂寞.
但你要相信神有為你預備痛恨者.
當你去找這些痛恨者的時候.

$^{681}$他們就會出現跟你一起去事奉.
一起去打這場屬靈的真戰.
當中神自己更加是我們隨時的幫助.
事上與我們同在.
我們永遠都不會剩下自己一個.
神的心意是要我們參與祂的工作.
將天國降在人間.
但我們要做的不一定是一些轟轟烈烈的事.
而是要我們帶著謙卑的心事奉神.
與神同工.
將我們的心獻上.
無論世界的反應是怎樣.
當我們願意甘心事奉的時候.
神都會喜悅我們的事奉.
並且感到心滿意足.
並且因著我們的事奉.
祂會感到心滿意足.
頂著我們一起祈禱.
就是天父我們多謝你.
你呼召我們作為你的兒女.
作為你事奉的同工.
主啊求你給我們感覺的不是孤單.
而是我們知道.
上帝你永遠都與我們同在.
你又派很多同行者與我們一起.
叫我們在寂寞低落的時候.
我們可以起來.
我們可以繼續走我們信仰的路.
主啊求你叫我們.
懂得如何去時常知道我們不孤單的時候.
更加願意輔助和幫助我們身邊.
可能感到孤單的人.
主啊求你與我們同在.
用我們這個群體.
成為一個有力的群體去見證主.
在這個這麼黑暗的年代.
我們繼續為主發光.
我們今天禱告.
奉主耶穌基督得勝的名祈求.
阿們.

\newpage



\section{箴言 9:10}
\label{sec:031999}
\textbf{認識至聖者便是聰明}~~~~蕭壽華牧師
\newline
\newline
連結: \href{http://www.church.com.hk/acms/content.asp?site=cdc\&op=show\&type=product\&code=031999\&layout=sermon}{\texttt{031999}} ~~~~ 語音日期: 2025-09-28
\newline
\newline
\hyperref[sec:031998]{\small{< < < PREV SERMON < < <}}
~
\hyperref[sec:index_chronic]{\small{[返順時目]}}
~
\hyperref[sec:032015]{\small{> > > NEXT SERMON > > >}}
\newline
\newline
箴言 9:10
\newline
\begin{longtable}{cl}
\hline
\hline
章節 & 經文 (和合本修訂版)\\
\hline
9:10 & \begin{tabularx}{0.7\textwidth}{X} 敬畏耶和華是智慧的開端，認識至聖者便是聰明。 \end{tabularx} \\ \\
[1ex]
\hline
\hline
\end{longtable}
$^{1}$此錄音或錄影的版權是屬於基督教宣導會北角堂所有.
我們的屬靈傳統之一是以真理教導為本.
歡迎你將此錄音或錄影的連結傳給親友.
懇請注意未經授權請勿轉錄或剪輯本錄音或錄影.
甚願如《仕途行傳》二十章三十二節下所說,這道能建立你們,叫你們和一切成聖的人同得基業..
就是我們要作出合適的選擇.
我們大家一起祈禱,我們仰望主義帶領我們.
人算什麼,你竟然眷顧我們.
我們算什麼,你竟然稱我們為你的兒女.
主,我們今天在世上體會自己很多很多的軟弱,甚至錯失.
也體會整個世界裡面充滿著各種的邪惡.
你要幫助我們行在光明的路上.
如今我們等候你,等候你自己,直到你的聖靈賜下你的恩言.
讓我們心靈領受智慧,知道一生跟隨你.
求你私恩,賜封這段時間,與你的兒女們同在.
我們仰望祈禱,奉耶穌聖名求下美.
認識智聖者,便是聰明.
我記得早兩個星期教會有猜拳聯會.
在那段時間教會一路推介一本書.
一本薄薄的書,叫《異鄉的跑道》.
是講述一班很知名的運動員.
他們後來也變成了宣教士的不同的故事.
一開始就講一個人物,這個人物大家都會熟悉的.
就是Eric Liddell,李愛蕊先生.
大家看過《烈火戰車》電影裡面的主角就是這位李愛蕊先生.
講述到他1924年,他拿到奧運400公尺的金牌.
1925年,他便來到中國做宣教士.
1941年,被日冠軍捉進了集中營.
1945年,死在集中營裡面.
我去留心他的傳記的時候.
沒有留意一件事情,就是這本書所講的.
在英國的集中營裡面,李愛蕊先生發現有很多很多的兒童少年.
他們都是從自己的父母身上分開出來.
他們在集中營裡面很辛苦,很擔憂,很焦慮.
因為父母不在身邊.
但是李愛蕊先生走出來照顧他們,關心他們.
也做他們的老師,教他們很多方面的功課.
其中有一位叫Steve McCarff.
這個麥凱夫弟兄,他當時是少年.
他講述在集中營期間,李愛蕊先生送了兩份禮物給他.

$^{41}$第一份禮物是將一對破舊的球鞋送給他.
你說破舊也送給別人?.
在集中營裡面,不要說鞋子.
連拖鞋也未必可以穿.
但是你要將他自己僅有的那對球鞋送給他.
他很感動.
第二件事情,李先生給他的一個很重要的禮物.
就是教他,要為那些日本軍人每天祈禱.
這個教導很不容易.
你可以想像一個小孩子,一個少年人.
在日本的集中營裡面,他每天面對著日本人對他的虐待.
對他的傷害.
但是李愛蕊就教他,你每天為他們祈禱.
你聽從耶穌給我們的吩咐.
要愛你的仇敵,為那逼迫你們的祈禱.
所以史蒂夫就一直每天學習,為日本軍人祈禱.
想不到李愛蕊後來自己死了在集中營.
因為史蒂夫後來就離開了,出來後被釋放.
釋放之後,當他慢慢長大之後,他決心去日本.
去日本幹什麼?.
要找那些日本人算帳.
他去日本做宣教士,一早做了四十年.
當我看到這段分享的時候,我的心很感動.
是一件絕對不簡單的事情.
曾經受到日本人各種的傷害.
但是竟然可以為他們祈禱.
有些有神給他的力量.
可以跟日本人生活,在他們之間傳揚著福音四十年.
各位電影節目,你會發覺.
在人生的路上,有很多路給你選擇.
有些時候到了十字路口.
你不知道要走一條寬路還是走一條窄路.
但是你會發覺,在聖經中有一個很清楚的方向.
如果你是認識自性者,你會選擇他的路.
而你會發覺,當你選擇他的路的時候.
你的心裡面有一種力量.
可能過去曾經受過傷害.
或者將會面對更多的傷害.
但是知道是主的路.
你就放心選擇主的路.

$^{81}$聖經說,這個就是「智慧人」.
他認識神.
因此他按神的心意選擇神的路.
他生命領受智慧.
我們讀《真言書》的時候一直說到.
《真言》提醒我們要尋求智慧.
來到第九章的時候你會發覺.
聖經帶給我們一個很重要的呼喚.
呼喚我們要選擇回應智慧.
回應智慧的呼籲.
大家一起看回聖經.
或者你歡喜今天可以拿出聖經來.
你會看到整本聖經裡面第九章的編排情況.
在《真言》第九章.
當你看的時候.
你會看到在第九章開始.
很工整有十八節的經文.
這裡十八節的經文分成三個小段.
每一段六節.
第一段第一到第六節.
《真言》第九章第一到第六節.
你看到是用一個比喻.
是用一個女子去比喻智慧.
然後去到第三個小段的時候.
第十三到第十八節.
就用另一個女子來描述.
比喻愚昧.
而中間那個小段另外的六節經文.
就講述給我們知道.
如何去分辨智慧和愚昧.
以及為自己的選擇而負責任.
很工整的排列.
但也給我們很清楚的知道.
有一個很重要的選擇.
擺在我們的面前.
當你讀到第九章.
其中第三節的時候.
你會看到智慧.
她去到城中至高處去呼喚.
但同樣去到第三個小段裡面.

$^{121}$其中也提及到.
愚昧這個聖經描述的另一個女子.
她坐在自己的家門口.
在城中高處的座位上呼喚.
同樣講到兩位都是在一些高地.
來向人發出呼喚.
也都意味著在一些大家都看到的地方.
大家都會留意到的地方.
他們高調地來發出呼喚.
跟隨我 跟隨我 進來享受我的宴席.
另一個角度看.
你會發覺二者其實在競爭中.
競爭要得到你和我的回應.
進入他們所預備的不同的宴席.
而第九章講到一位智慧的女子的時候.
師兄就提到智慧建造了房屋.
造成七斤的柱子在薩森竹吊好尾酒.
又擺設賢席.
這個智慧她建造了美屋的房屋.
請提她有七斤的柱子.
當時的房屋一般來說.
最多一斤至兩斤的柱子.
七斤的柱子表示這個智慧.
她在一些很廣闊的地方.
建築了宏大寬敞的房屋.
用七斤柱子才能夠支撐得住.
這個宏大的房屋裡面.
很歡迎你和我進入.
而智慧為我們預備了很豐盛的宴席.
宰殺了生畜 調好了尾酒.
盼望我們能夠進入.
智慧預備了很豐富的宴席.
盼望我們會進入.
智慧很主動地呼籲.
你看到聖經第三節 第四節.
一直在說 你們來吧.
你們轉過來吧.
誰是愚夢的人.
讓他轉到這裡來.
很熱切地呼籲.

$^{161}$特別是那些愚夢的人.
無知的人.
當然一個人.
他要承認自己的愚夢.
他會接受自己的無知.
他才會回應這個呼召.
但是智慧不斷向他們提出.
你們轉過來.
到我這裡吧.
我預備了很多的美食.
你過來 你過來 你過來.
一直在呼籲.
但是你知道.
有很多人不願意回應這個呼喚.
大家都會記得.
當主耶穌在地上的時候.
祂說了很多的比喻.
《盧高芳》十四章說到一個宴席.
天國宴席的比喻.
你記得嗎.
就是主人邀請了很多的客人.
來享用宴席.
到時候預備好了.
吩咐那些僕人出去.
你們過來吧.
享用我的宴席.
因此僕人就出去.
就跟那些被邀請的客人說.
預備好了.
但是那些客人怎麼反應.
有人就說 對不起.
剛剛買了一塊地.
我是趁低吸納的.
現在非常有價值.
我在想怎麼去發展.
我家而已.
又遇到一些人就說.
我買了五對牛.
你知道牛當時來說是財富.
這對牛怎麼能夠發展他們的能力.

$^{201}$一直要去看看.
有人說不行.
我剛剛娶了太太.
很多事情都不管了.
僕人會去匯報給主人知道.
主人聽見很憤怒.
然後就說了一句話.
先前所請的人.
沒有一個得償我的宴席.
他們沒錯曾經被邀請的.
但是他們不願意回應.
因此沒有一個能夠得償.
天國的主人的宴席.
各位頂智媒.
其實你都會體會.
在今天這個世界裡面.
有很多不同的呼聲呼喚你.
呼喚你 來吧.
跟著我 你有前途了.
來吧 你這樣決定.
包你前面沒有任何缺乏.
很多不同的事情.
在呼喚著你的注意.
但是你要選擇去回應.
智慧的呼喚.
免得你失去真實美善的宴席.
這正是聖經對我們的吩咐.
當聖經說完這件事情之後.
直接再說第六節的主題.
你們來吃我的餅.
喝我調的酒.
如果你要去進入這個宴席.
選擇智慧的時候.
你就要離棄愚昧.
離棄的意思.
含有一種悔改的意思.
為了過去錯誤的選擇悔改.
離棄那些錯誤的選擇.
然後你就會得到存活.
得到生命.

$^{241}$並要走.
繼續走在智慧的路上.
繼續讓你堅持在神的路上.
你就能夠領受智慧.
領受智慧不只是一刻的決定.
更是一個長久持定的決定.
一生願意去行智慧.
說完這六節經文之後.
剛才說到最後的六節經文.
聖經說到另一件事情.
聖經說到你要去留心拒絕愚昧.
在九章十三十四節.
說到另一個女子.
這個女子在裡面有很多呼喚.
聖經說她圈養.
圈養這個字有很多含義.
除了說她很大聲.
四處呼喚去爭取你的注意之外.
這個字含有一種就是.
不守若束 荒蕩無恥的意思.
不管你會怎樣.
一直要她拉著你.
扯著你.
要你去跟從她 付從她.
當聖經提到.
坐在她自己的家門口的時候.
很多當時的人會想起.
在當時的社會.
有些妓女會坐在她的家門口.
呼喚別人進入她的家.
跟她有勾當.
因此你發覺聖經用一個邪淫的女子.
去形容愚夢 愚昧.
而她不斷地發出呼喚.
她呼喚誰呢.
呼喚那些剛剛路過的.
聖經接著說到.
呼喚過路的 向那些在路上直走的人.
他們未必特意要去找她.
但是剛剛路過的時候.

$^{281}$剛剛看到.
其實無意去回應她的.
但是被她叫住.
就受到吸引.
我們平時在街頭上.
有時候都遇到一些.
街頭做一些推銷的人士.
推銷某些物品.
我也試過一兩次不期然地走過去.
一直用著一些抗音器大聲說.
這些東西不用買的.
今天是送的.
全部是贈品.
你用一個好用的.
你回頭再幫我買.
接著我走過去之後.
站在那裡.
想著她送一些東西.
誰知道她繼續解釋.
拿著那份物品一直說.
說到最後目的只是上來買.
有太多太多的事情你會明白.
用很多不同的方法.
要呼喚你來接受.
而且她告訴你.
偷來的水是甜的.
暗藏的餅是美的.
就這樣吃那些餅不好吃的.
你要偷回來.
靜悄悄沒有人知道的時候吃.
就最美味了.
另一個的翻譯.
環聖利翻譯就是.
到來的水才是甜的.
偷來的餅才是香的.
我記得一件事情.
我很深刻的.
很多年前有一位.
在外籍的弟兄.
他去到柬埔寨.

$^{321}$經歷了一些事情.
告訴我.
他在柬埔寨一個傍晚的時候.
他經過一個地區.
那個不熟悉的地方.
他去的時候.
後來才發覺是一個紅燈區.
一路走的時候.
突然有一個男士出現.
見到他的時候.
就微微笑就跟他說.
用英文說.
little boys.
little girls.
他一開始都不明白.
後來他再聽那個人說的時候.
那個人一直告訴他.
用他的方法表達就是.
很刺激的.
比其他更滿足的.
大家知道是什麼.
在那些紅燈區裡面.
不單單要人有.
一些妓女的活動.
然後推銷一些兒童.
讓一些人去種他們的情慾.
當我聽的時候.
我感覺很噁心.
但是同時體會到.
當時的人用什麼方法.
讓你認為.
這些更加滿足的.
更加香甜的.
聖經描述到撒旦.
在啟示錄第十二章講述到.
魔鬼有一些名稱.
其中的名稱.
魔鬼就叫撒旦.
撒旦的意思是敵對者.
是敵對神.

$^{361}$撒旦是敵對光明的事情.
如果你走在光明當中.
祂要將你拉出來.
走在黑暗裡面.
撒旦又是迷惑普天下的聖經講.
去粉飾那些邪覺.
告訴你很美妙的.
以致你能夠進入那種迷惑裡面.
各位電影節目的聖經提醒我們.
如果你不是警醒禱告.
你就會進入迷惑.
當你面對著很多不同的選擇的時候.
你要祈禱.
你要小心地去求問神.
哪一條是智慧的路.
哪一條是智慧的選擇.
如果你不小心選擇.
你會被迷惑.
你會中了詭計.
進入了愚無人的家.
這只是聖經不斷地提醒我們要留心.
祂告訴你.
是香的餅.
是甜的水.
但接著是第十八節.
聖經說.
人卻不知道有陰魂在他那裡.
他召喚的人士在陰間的深處.
祂帶你去的地方.
不是什麼美善的宴席.
是陰間的宴席.
結果是永遠的死亡.
是永遠的離開光明的上帝.
他一定會面對著試探的時候.
某些處境你感覺到有一些問題.
當你要做決定的時候.
你要想一想.
如果我選擇了這件事情.
有什麼後果.
對我的生命有什麼影響.

$^{401}$當你小心想.
原來祂帶你去陰間.
你會得到死亡的時候.
你就會有更大的力量去拒絕這些試探.
這些危險.
沒錯.
智慧在呼喚邪惡的愚昧的.
又在呼喚你.
但是到了聖經中間的小段.
其中第七到第十二節.
聖經提醒我們要留心.
當你選擇的時候.
你會承擔你的選擇的後果.
在第七節開始說到一件事情.
就是你如何去分辨智慧和愚昧.
聖經說到.
糾正傲慢人的必招羞辱.
責備惡人的必備侮辱.
不如責備傲慢人免得他恨你.
聖經說到一件事情就是.
你如何去分辨智慧和愚昧.
其中一個方法就是.
透過他們如何回應別人對他們的勸告.
他們不同的回應方式.
你會知道他是愚昧還是智慧.
這裡提到一件事情就是.
智慧人他如何去回應別人對他的勸告.
聖經說他是傲慢的.
他是驕傲的 自以為是.
愚昧的人他認為自己是最正確的.
看不見自己的錯失.
也不承認自己的錯失.
甚至別人勸告他的時候.
他會用惡言去反擊他.
甚至去侮辱對方 仇恨對方.
我想大家可能都記得我曾經說過.
很多年前自己一個經歷.
我又想起.
就是很多年前我做社工.
在一個社區裡面.

$^{441}$一個社區中心裡面.
常常有很多小朋友來到這個中心裡面.
那件事情.
去到中心裡面.
這些朋友很小的.
四五歲 三四歲 差不多.
竟然他們在中心裡面大罵.
是大家互相罵人.
用粗口罵對方.
我聽見 嘩 這麼小的罵粗口的.
粗口很純熟.
當時我就跟他說 不行的.
小朋友不可以說粗話.
為什麼不可以說.
差不多想罵我.
接著我就說 好.
你們下次什麼時候來.
我幫你們辦一個辯論比賽.
然後題目就是.
小孩子應否說粗口.
到時他們就真的來了.
分成兩隊.
一隊正 一隊反 就要辯論.
我很記得在辯論過程當中.
有一位社區的叔叔就進來.
進來的時候.
這麼有趣 小朋友全都在.
聽他們說話.
嘩 說這些話的.
然後就用粗口罵他們.
一直罵他們.
為什麼這麼小就說粗口的.
我旁邊聽到.
我就很明白.
在這樣的環境裡面成長.
自然就學會了.
然後我也忍不住.
走過去跟這位叔叔說.
叔叔 你這樣說粗口.
他們怎麼會不說呢.

$^{481}$然後我永遠不會忘記.
這位叔叔看著我.
我沒有說過粗口.
是我第一次.
這麼清楚看到一個人.
自以為是.
不願意承認自己所做的事情.
其實你和我.
有些時候.
好像都不是這樣.
做一些自己知道不應該做的事情.
有人問及的時候.
不知道是誰撞跌的呢.
你發覺聖經裡面提及一件事情.
一個自慰人是願意去接受勸告.
願意面對自己的錯失.
在過程裡面聖經就說.
這些自慰人.
聖經跟著就說.
他受到責備的時候.
他去愛那些責備的人.
而當他接受人的責備.
在當中自己謙卑的反思.
因此從那個過程當中.
有新的學習.
聖經就說他愈有自慰.
也說給我們聽.
這些異人.
聖經將自慰人等同異人.
只是異人他就增長學問.
他會更加有對自慰的認識.
進深的學習自慰.
我想請問一下聖經的針形.
其他地方提到.
一個人當接受到別人勸誡的時候.
是何等的甘美.
聖經針形的二十七章.
其中的第六節說到.
有些的敵人很欣賞你.
和你親密.

$^{521}$大家知道當時表示的一種.
很親密的一種態度.
但是這些是多餘的.
但是朋友家的傷痕.
好像是傷痕.
但是出於忠誠.
卻是帶給你祝福.
高油和香料使人心喜.
朋友誠心的勸勉也是如此的甘美.
是帶給你有滿足.
帶給你真正的祝福和學習.
最近大家可能都聽到.
就是我們所敬愛的.
曾立華牧師過世的消息.
當有機會來參與追思禮拜.
聽到很多不同的同學.
弟兄姊妹來說說他們的經歷的時候.
我聽到一位曾經是.
曾牧師神學院的學生.
現在已經是傳道人.
他說回的情況就是當年.
自己在一次的事奉的經歷裡面.
覺得很沉重.
很大壓力.
心裡面覺得很不想繼續.
甚至想放棄.
在一個很沉重的心情的時候.
遇到曾牧師.
因此就和曾牧師.
傾訴自己所面對的困難.
一邊說的時候.
曾牧師一邊小心聽.
然後和他說了一句話.
他說弟兄.
你當跑的路還沒跑完.
只是簡單一句話.
這位弟兄這位牧師就說.
他當時很深刻地體會一件事情.
無錯有艱難.
無錯有挑戰.

$^{561}$但是神給我的路.
我要忠心去走完它.
因此這句話就鼓勵他.
繼續面對那些艱難.
繼續的事奉神.
你會發覺.
能夠有人給你勸告.
是你的祝福.
智慧的人.
是會在其中小心的分辨.
小心去聆聽.
還在其中領受智慧.
當聖經說到.
當你有這種智慧的表現的時候.
你就更加明白九章十節.
也是我們上上讀的經文.
重要的意思.
不如大家又再一起讀出這則經文.
敬畏耶和華是智慧的開端.
認識自性者便是聰明.
請提到敬畏耶和華是智慧的開端.
提出智慧的根源在哪裡.
智慧的本質是什麼.
智慧不是說你的智商有多高.
你的智慧是否爆棚.
而是說你是否尊崇上帝.
是否你的心真的知道.
一位是至高者.
是可重可畏的全能的上帝.
你不會輕忽他對你的旨意.
不會輕視他對你的吩咐.
你會聽從你去尊崇他.
這是敬畏也是智慧.
當你去認識他的時候.
你知道這位的智性者.
他喜愛公義和公平.
你就知道自己可以.
按公平公義的原則.
去選擇你的道路.
你就變得有智慧.

$^{601}$你記得我們上上都會提到.
耶利米書九章這一節的經文.
你如何去誇厚呢.
聖經提醒我們.
先知提醒我們.
智慧人不要因他的智慧誇厚.
我比別人聰明一點.
不是.
勇士不要因他的力氣誇厚.
那些是他自己最強的項目.
他不要因為自己力氣強就誇厚.
財主不要因為他的財富誇厚.
要按什麼誇厚呢.
誇厚的卻要誇自己有聰明.
什麼聰明呢.
是認識我是耶和華.
知道我喜悅在世上施行.
慈愛公平和公義.
你認識到這位上帝.
他很重視慈愛公平公義.
他所做的一切都是公義的.
當你認識的時候.
你就跟隨著這位上帝.
他的屬性.
你去作人生的選擇.
你就變得智慧.
當你認識智慧至高者.
你就有真正的聰明.
或者用一句話說.
什麼叫聰明.
聰明的定義是什麼.
聰明就是認識至聖者.
當然你知道這個認識.
不單單是說頭腦上的認識.
是你在生命當中.
你去跟隨的時候.
一步步去經歷這位神的真實.
而在當中更加去信靠他.
按他的心意來生活.
我想說一件很簡單的經歷.

$^{641}$有一年我去南美.
參加一些宣導會的國際會議.
我搭飛機去到紐約.
然後再轉機去南美.
去到紐約的時候.
飛機延遲了時間.
去到很遲了.
我立刻要接駁那班飛機.
但當時已經是傍晚.
晚上的時間.
我拖著行李一直跑.
因為我要先把行李清關.
再拿去新的地方再入機.
還有差不多一個鐘頭.
我打算也夠時間的.
當我跑到去的時候.
全部關門了.
因為他說這裡很多人帶大行李.
所以一早要清關.
才可以進來.
關門了.
人家都進去了.
我看著飛機怎麼辦.
我心裡很焦急.
而且很焦慮.
那怎麼辦呢.
我怎麼辦呢.
在當時才能夠默默地祈禱.
那個官員看到我.
似乎很焦慮.
我們可以安排酒店給你住一晚.
這樣就好了.
住一晚.
第二天早上休息完之後.
就去到機場再進去搭飛機.
進去的時候.
見到昨天的接待員.
就是你昨天晚上上不了飛機.
好吧.
給你一個好機位.

$^{681}$就把我帶到第一階段.
我就發覺.
竟然會這樣優待我.
但是當我一邊坐著.
享受那個機程的時候.
我突然間就告訴自己.
為什麼當時要這麼焦慮.
其實我是不需要這麼早去到.
我是去到當地.
還有一點時間.
才開始那個會議.
而且經過一段長程去到紐約的時候.
其實我身心靈很疲倦.
天父知道我其實需要先休息一晚.
你才再搭飛機.
一個很微小的經歷.
我就體會到.
在人生很多的事情上.
不需要急.
你失去了某一些事情.
不一定不好.
上帝掌管你的人生.
你掌管你的時間.
你放心.
每一步順靠祂.
你會經歷神是那位信實的上帝.
在人生很多的路上.
留心去聽從祂的旨意.
去依靠祂.
你會嘗到主人的滋味.
就知道他事未善緹.
投靠他的人有福了.
而你更加知道.
在人生的路上.
放膽去選擇神的路.
走進那條路.
看起來好像很艱難.
那條是尾線的路.
我自己年輕的時候.
很喜歡唱《青年聖歌》.

$^{721}$其中一首歌是我很喜歡唱的.
這首歌叫做《神的路最尾線》.
首詩歌我一邊唱的時候.
講到神的路最尾線.
雖我不明瞭.
然後講到.
為何憂愁試煉常把我環繞.
冒著跟隨主的路上.
不知為何有些試煉會出現.
為何這些事情令我這麼憂心.
這麼憂慮.
是不明白的.
但一邊唱的時候.
富哥一邊提醒.
神的路最尾線.
神的路最美好.
我願常依靠祂.
為祂能引導.
當你依靠的時候.
你就經驗到.
祂其實是在引導你.
讓那些看起來像是擇路的路.
成為叫你蒙福的路.
神的路是最尾線.
等於是選擇行神的路.
這正是《針燕第九章》.
不斷地呼喚我們.
不論那條路你看起來是怎樣.
你記得.
那條如果是神的路.
你分辨到是神的路.
你就放膽走上去.
因為那條是尾線的路.
聖經再說.
當讀下去的時候.
我們回到剛才經文裡.
第十章十字講到.
敬畏業禍罷.
以及認識至聖者.
接著聖經就說.

$^{761}$藉著我 藉著智慧.
你的日子必增多.
你的生命的年數也必加添.
提到在我們人生的路上.
神賜福給我們.
這裡表達的是一種蒙福的人生.
而且有永恆生命的恩惠.
神將祝福帶給那些智慧的人.
那些敬畏神的人.
然後最後來到.
第十二節這小段的最後的時候.
這個小段也提出一個.
很總結性的一個提醒.
大家一起讀一讀這節經文好嗎.
你若有智慧.
是自己有智慧.
你若傲慢.
就自己承擔.
是什麼意思.
先講了一件事情就是.
你要為自己的選擇負責.
你要去自己承擔.
你做的決定的後果.
有人說這節經文.
是聖經裡面.
一節最個人主義的經文.
個人主義不是一個負面的講述.
是在告訴你.
你做的決定是你自己要承擔的.
當你憑著智慧去選擇神的路.
你就發覺你是那個最大的受益者.
你可能所做的事情對別人有祝福.
但是最大的祝福臨到你的身上.
你是那個最大的受益者.
但如果你選擇了傲慢.
選擇了愚昧.
你就要自己承擔愚昧所帶來的一切的後果.
很多人做錯事情就會說.
不關我的事.
是社會造成的.

$^{801}$不關我的事.
我父母自小給我很壞的榜樣.
影響我做這些決定.
沒錯.
可能別人對你有不好的影響.
但仍然是你自己決定.
是不是跟隨這些影響.
當你要決定怎麼做這件事的時候.
你自己要承擔責任.
神仍然是沉默去追問你.
是不是你自己決定走的路.
請你們很清楚把兩條路擺在我們面前.
各位電影節目.
當我們讀完這段經文的時候.
你會知道智慧不斷大增地呼喊著你.
你要留心回應祂的呼喊.
每一天小心地作選擇.
就算要付出代價.
就算你經歷一些損失.
你放心去選擇神的路.
你最終會體會到是美善的路.
是給你永恆地蒙福的路.
我們一起祈禱.
親愛的主 我們同心仰望你.
我們知道這個世界.
有太多迷惑人的事情.
有太多似是而非的事情.
不斷衝擊我們的心.
我們在你面前謙卑地求告主.
祝賀賜我們心靈的那種敏銳.
能夠認識你 敬畏你.
以致我們知道分辨邪惡.
知道行在光明的路上.
主啊 求你施恩.
多謝你給我們過去這一季.
一路的讀真言.
我們今天在你面前.
痛心向你祈禱.
求主將一顆敬畏你的心.
放在我們心的深處.

$^{841}$讓我們在你裡面.
領受從天而來的智慧.
以致人生在很多的抉擇上面.
我們都懂得去選擇你的路.
主啊 你幫我們.
幫我們每一個人.
行在一條蒙福的.
真正快樂的路上.
奉主耶穌聖名祈禱.
行門.
\newpage



\section{詩篇 139:1-24}
\label{sec:032015}
\textbf{認識神無所不在的生命轉化}~~~~張健庭牧師博士
\newline
\newline
連結: \href{http://www.church.com.hk/acms/content.asp?site=cdc\&op=show\&type=product\&code=032015\&layout=sermon}{\texttt{032015}} ~~~~ 語音日期: 2026-04-19
\newline
\newline
\hyperref[sec:031999]{\small{< < < PREV SERMON < < <}}
~
\hyperref[sec:index_chronic]{\small{[返順時目]}}
~
\hyperref[sec:032016]{\small{> > > NEXT SERMON > > >}}
\newline
\newline
詩篇 139:1-24
\newline
\begin{longtable}{cl}
\hline
\hline
章節 & 經文 (和合本修訂版)\\
\hline
139:1 & \begin{tabularx}{0.7\textwidth}{X} 耶和華啊，你已經鑒察我，認識我。 \end{tabularx} \\ \\ \relax
139:2 & \begin{tabularx}{0.7\textwidth}{X} 我坐下，我起來，你都曉得；你從遠處知道我的意念。 \end{tabularx} \\ \\ \relax
139:3 & \begin{tabularx}{0.7\textwidth}{X} 我行路，我躺臥，你都細察；你也深知我一切所行的。 \end{tabularx} \\ \\ \relax
139:4 & \begin{tabularx}{0.7\textwidth}{X} 耶和華啊，我舌頭上的話，你沒有一句不知道的。 \end{tabularx} \\ \\ \relax
139:5 & \begin{tabularx}{0.7\textwidth}{X} 你前後環繞我，按手在我身上。 \end{tabularx} \\ \\ \relax
139:6 & \begin{tabularx}{0.7\textwidth}{X} 這樣的知識奇妙，是我不能測的；至高，是我不能及的。 \end{tabularx} \\ \\ \relax
139:7 & \begin{tabularx}{0.7\textwidth}{X} 我往哪裡去，躲避你的靈？我往哪裡逃，躲避你的面？ \end{tabularx} \\ \\ \relax
139:8 & \begin{tabularx}{0.7\textwidth}{X} 我若升到天上，你在那裡；我若躺在陰間，你也在那裡。 \end{tabularx} \\ \\ \relax
139:9 & \begin{tabularx}{0.7\textwidth}{X} 我若展開清晨的翅膀，飛到海極居住， \end{tabularx} \\ \\ \relax
139:10 & \begin{tabularx}{0.7\textwidth}{X} 就是在那裡，你的手必引導我，你的右手也必扶持我。 \end{tabularx} \\ \\ \relax
139:11 & \begin{tabularx}{0.7\textwidth}{X} 我若說「黑暗必定壓碎我，我周圍的亮光必成為黑夜」， \end{tabularx} \\ \\ \relax
139:12 & \begin{tabularx}{0.7\textwidth}{X} 黑暗對你不再是黑暗，黑夜卻如白晝發亮。黑暗和光明，在你看來都是一樣。 \end{tabularx} \\ \\ \relax
139:13 & \begin{tabularx}{0.7\textwidth}{X} 我的肺腑是你所造的，我在母腹中，你已編織我。 \end{tabularx} \\ \\ \relax
139:14 & \begin{tabularx}{0.7\textwidth}{X} 我要稱謝你，因我受造奇妙可畏，你的作為奇妙，這是我心深知道的。 \end{tabularx} \\ \\ \relax
139:15 & \begin{tabularx}{0.7\textwidth}{X} 我在暗中受造，在地的深處被塑造；那時，我的形體並不向你隱藏。 \end{tabularx} \\ \\ \relax
139:16 & \begin{tabularx}{0.7\textwidth}{X} 我未成形的體質，你的眼早已看見了；你所定的日子，我尚未度一日，都在你的冊子寫上了。 \end{tabularx} \\ \\ \relax
139:17 & \begin{tabularx}{0.7\textwidth}{X} 神啊，你的意念向我何等寶貴！其數何等眾多！ \end{tabularx} \\ \\ \relax
139:18 & \begin{tabularx}{0.7\textwidth}{X} 我若數點，比海沙更多；我睡醒的時候，仍和你同在。 \end{tabularx} \\ \\ \relax
139:19 & \begin{tabularx}{0.7\textwidth}{X} 神啊，惟願你殺戮惡人；你們好流人血的，離開我去吧！ \end{tabularx} \\ \\ \relax
139:20 & \begin{tabularx}{0.7\textwidth}{X} 他們說惡言頂撞你，你的仇敵妄稱你的名。 \end{tabularx} \\ \\ \relax
139:21 & \begin{tabularx}{0.7\textwidth}{X} 耶和華啊，恨惡你的，我豈不恨惡他們嗎？攻擊你的，我豈不憎惡他們嗎？ \end{tabularx} \\ \\ \relax
139:22 & \begin{tabularx}{0.7\textwidth}{X} 我恨惡他們到極點，以他們為我的仇敵。 \end{tabularx} \\ \\ \relax
139:23 & \begin{tabularx}{0.7\textwidth}{X} 神啊，求你鑒察我，知道我的心思，試煉我，知道我的意念； \end{tabularx} \\ \\ \relax
139:24 & \begin{tabularx}{0.7\textwidth}{X} 看在我裡面有甚麼惡行沒有，引導我走永生的道路。 \end{tabularx} \\ \\
[1ex]
\hline
\hline
\end{longtable}
$^{1}$丁姐妹平安.
我們再一次低頭祈禱.
主持人:天父,願你自己的說話跟我們說.
有時候我們對於自己都不是那麼明白,那麼清楚的時候.
但是你卻是知道我們的心腸肺腑.
你比我們更加明白我們自己.
你也是不斷時刻的去監察,眷顧我們的神.
願你藉著今天的詩篇來跟我們說話.
奉耶穌基督的名祈求,阿們.
好了,你知道現在來說.
很多的觀念,有些觀念是現代人有的,以前的人未必有的.
其中一個觀念,不知道你有沒有想過.
就是私隱的觀念.
以前的人其實是未必有私隱的觀念.
現在的人就很注重.
譬如父母,你不要亂看你子女的東西.
他很不高興的,就算很小的時候也有這樣的觀念.
更加不會說,如果他寫了一本日記.
你看了他,你犯了彌天大罪,他會生你的氣很久.
這個私隱的觀念,不知道是不是近代的.
不知道是否現代人有的.
以前的人,我想,其實真的未必有這樣的觀念.
你想想,以前,舉例來說.
一個家庭,八九個人住在一個很小的房間裡.
哪裡有私隱可言?.
甚麼都被人知道了.
無論是工作的,譬如以前的農業社會.
一家人可能是一起去耕種同一個田的.
你今天做過甚麼,你都不用跟別人說.
大家都知道你做過甚麼.
回到家裡一起吃飯,一起睡覺.
家裡所有的東西都不多.
你的所有的生活,所有人都知道的.
沒有任何私隱觀念.
但現代人對於私隱的觀念是很有介懷的.
為甚麼呢?.
我相信其中一個原因就是.
因為很多時候被人監察.
這是一個很不自在的觀念.
你想想,有句說話叫.

$^{41}$looking over the shoulder.
即是有個人在後面監視著你工作.
你都很不開心.
可能你老闆經常在你後面監視著你工作.
你會發覺很不自在,我不是不懂得做.
有時有些員工都很不客氣.
你走開一點,你不要看著我.
你看著我,我不喜歡的.
因為我們被人監察著工作.
有一個監控是不會覺得舒服的.
我不知道你知不知道.
其實以前的監獄和現在的監獄是很不同的.
當然現在的監獄有很多科技.
但其實在18,19世紀的時候.
對於企監獄有了很大的變革.
或者有了很大的改革.
是怎樣呢?.
就是因為以前的監獄需要很多玉卒.
來監視囚犯.
否則他們造反時失控很難處理.
但後來在18世紀的時候.
有一個建築家想了一個觀念.
這個觀念就把整個監獄的建築完全改革了.
自此之後就減少了很多玉卒.
只用兩三個玉卒就能監視幾十個囚犯.
這是很厲害的.
有什麼方法呢?.
如果幾十個囚犯造反.
為什麼用兩三個玉卒就能處理呢?.
你知不知道有什麼方法?.
那個建築師就想出了.
其實如果囚企的監獄每間房間.
是互相監視對方的.
可能對面房間看到你.
或者你看到對面房間的時候.
被監禁的人覺得有人時時刻刻.
24小時有人監視.
那時候還沒有監控鏡頭.
我相信現在有監控鏡頭的時候.
監房就不會這樣建造.

$^{81}$因為監獄根本都能監視.
但以前沒有監控鏡頭.
所以建築師的想法就是.
在一間房間對著一間對著一間.
囚犯就能互相監視對方.
發現原來真的可以.
囚犯發現有人監視.
不論是玉卒或是室友.
都知道有人監視.
不敢做一些很離譜的事.
這樣就真的可行.
自此之後建監獄就有了改革.
不同的建法.
現在的監獄就沒有這樣的觀念.
沒有需要這樣做.
為什麼?.
因為監控鏡頭已經搞定了.
何必要建在對面.
但是.
就是一個這樣的觀念.
人覺得被監控的時候.
是會很不自在的.
而事實上我們現在.
就是活在一個無時無刻被監控的時代.
所以私隱這個觀念.
我相信就是隨著監控的.
越來越普及的時候.
人就越來越覺得需要重視私隱.
因為人什麼都可以找出來的時候.
就發覺很不舒服.
所以政府有時候就要立一些私隱法.
就是因為.
人裡面有一種.
對於自己的生命.
躺開出來.
或者將自己的生命.
揭露出來.
顯露在別人面前.
有一個不是那麼安全的感覺.
這個就很有趣.

$^{121}$就是為什麼人會有這樣的感覺呢?.
而事實上.
在聖經今天我們所看的這篇蕭篇.
正是去探討這種.
人覺得被人監控的時候.
那個感覺不安全或者不舒服的感覺.
在耶和華.
在舊約聖經的時代.
當然就沒有監控鏡頭.
也沒有說24小時有人看著你.
但是在猶太人的觀念當中.
有一樣東西他們是很知道的.
就是耶和華神是無處不在的神.
耶和華神也是監察.
所有人的心思意念和眼目.
所以詩人就用這樣的進入點.
其實他就反省到.
如果被耶和華神看著.
那又如何呢?.
所以一開始的時候.
他就說耶和華.
你以監察我.
認識我.
我坐下我起來.
你都曉得.
你從遠處知道我的意念.
我行路我躺臥.
你都世察.
你也心知我一切所行的.
耶和華.
你舌頭上的話.
我舌頭上的話.
你沒有一句不知道的.
你在我前後環繞我.
你按手在我身上.
詩人就一系列地列出了.
如果活在神面前的時候.
有什麼是神知道的.
他說神知道什麼呢?.
監察認識我的時候.

$^{161}$我做的每一個動作.
我坐下.
我起來.
有時候你做老師就知道.
有些頑皮的學生.
很喜歡坐後面.
因為不想被老師看著.
動一動手的時候.
老師的眼睛已經看過來.
就不太想.
但是現在就說.
我動一動手.
站一站.
坐下.
我起來.
神就知道.
你從遠處知道我的意念.
而且說.
我不用說廣東話.
動一動尾巴就知道你做什麼.
不是的.
他說連意念都知道.
有時候有些事.
你想的你都沒做.
有時候你想都覺得.
不是太好了.
那就不要做了.
但是你說.
想不好的時候.
神都已經知道.
你從遠處知道我的意念.
然後他說.
我行路我躺臥你都細察.
甚至神是很詳細的觀察.
你怎麼走路.
今天不一樣了.
或者你很累嗎?.
這麼快睡?.
有時候父母很留意小孩子.
就是這樣的觀察.

$^{201}$我行路我躺臥都細察.
你心知我一切所行的.
我所做過的一切.
某一件事神不知道.
24小時都所做的一件事.
有時候你自己都忘記了.
有時候隨著時間過.
有時候你做慣了的事.
好像等於沒做.
你不知道有沒有這樣的感覺.
有時候譬如舉例來說.
有些事隨手你擺慣了在那裡.
你用完之後.
很自然的動作就擺在那裡.
你是沒有刻意去做這個動作.
可能你旁邊有些人留意的.
就問你.
你為什麼要擺在那裡的時候.
可能你都忘記了你擺在那裡.
你說是嗎?我擺在那裡嗎?.
有時候我們自己做的事.
我們自己都不留意.
我們自己都不知道.
但是神完全知道.
你想想.
然後再說.
我舌頭上的話沒有一句你不知道.
這樣就糟糕了.
有時候說的話.
無心之失也好.
衝口而出也好.
有時候你說完後都後悔.
你都希望說的時候沒有人聽到.
是不是?.
然後說.
我所說的每一句.
你都知道.
有時候說好話.
有時候你還小聲的跟別人說中文.
說悄悄話.

$^{241}$或者說別人背後的說話的時候.
或者說一些很私隱的事情的時候.
神都知道.
然後說.
你在我前後環繞我.
按手在我的身上.
然後.
他就再說到.
那種包圍感.
就好像三十六度全方位包圍一樣.
好像被神抱著一樣.
他的意思其實就是說.
你在我前後環繞我.
按手在我的身上.
就好像被神抱著一樣.
就是.
我養貓.
那些剛開始養貓的時候.
那些小小的貓.
那些野貓.
抓回來的.
其實是被人 kidnap 的.
那些是街貓.
被人抓的.
不是我們抓的.
是有一個朋友.
他在後院生了一群.
抓了那些貓.
但是那些是野貓.
不習慣被人養的.
我們最初去領養那些貓的時候.
他教我們一招.
他說怎麼樣呢?.
你用一條毛巾包著牠.
就好像包密西哥卷餅一樣.
包著那些貓.
他說你就把牠關在籠子裡.
一出來.
一抓.
一開始養的時候.

$^{281}$牠是野性的.
然後就把牠關在籠子裡.
然後.
到你給牠吃的時候.
你抓牠出來.
你就用毛巾.
就像包圍一樣包著牠.
包到牠動都不能動.
只有牠的口.
這樣.
就讓牠吃東西.
那.
你想像一下.
如果你是那隻貓.
都很恐怖的.
被人.
就是24小時在籠子裡.
所有東西被人知道了不僅如此.
然後牠出來的時候.
就是吃東西的時候放出來的時候.
是被人好像抱著.
就是其實那塊毛巾包著.
就好像被人強硬地抱著.
動彈不得.
就是.
不同的動物有不同的特性.
狗可能很喜歡被主人抱著.
但是其實貓是不喜歡被人抱著.
就是貓被人抱著是很覺得受威脅的.
但是沒辦法.
這麼小的貓.
人又大力過牠.
然後毛巾包著牠.
手腳都伸不出來.
然後就被人包圍著.
這裡其實有些事情像是這樣.
他說神前後圍繞我.
按手在我身上.
好像強制地抱著我.
那.

$^{321}$在私人的角度一開始的時候.
當然你可以從一個正面的角度來演繹.
這樣就很好.
就是很親密.
很有那個.
就是很有那個親密感.
覺得神很近.
但是.
但是從另一個角度來講.
就是.
神是對人這種親密的認識.
人.
是不是就是.
會很舒服呢?.
那我不說.
所有的人.
就是說你自己.
你想想你自己.
當你想一下.
如果你每一個的意念.
你神知道.
每一句話.
神又知道.
你走路躺我.
有時候你動一動.
神又盯著.
神又看著.
然後.
現在好像四面八方包圍著我.
不說其他人了.
就是說你自己.
你會有什麼感受呢?.
你會不會想.
神啊.
不如你不要看我那麼多.
會不會呢?.
還是你說.
神啊,看我多一點.
盯著盯著.
就是.

$^{361}$你的眼不要離開我.
專注著我最好.
我不知道你是哪一種.
但是我相信很多人來講.
其實這個.
都是一種.
有威脅感的一種.
心理狀況.
就像剛才說的.
你做事.
人家在後面.
不斷好像有眼睛在看著.
很有威脅.
心理上很有威脅的狀況.
接著.
所謂詩人就說.
這人的知識奇妙是我不能測的.
至高是我不能及的.
他說.
他就.
由然而生出了一種.
敬畏的一種感覺.
究竟是什麼意思呢?.
就是說.
有一個神是完全超出了我自己.
我沒辦法.
想像得到究竟是怎樣.
神是.
不是.
好像說.
有時候我們很容易會想像神是.
好像跟我們一樣.
很親近.
耶穌是我最好的朋友.
那種.
好像.
很近的.
好像一個.
磁場的長者也好.
或者很親密的朋友也好.

$^{401}$那種.
雖然他知道我.
但是也是一種沒什麼威脅感.
好像朋友一樣.
我要告訴他.
我就隨便.
我就隨我的心意去說.
但是我不說也可以的.
有時候我們會有一種.
有時候我們會有一種這樣的感覺.
或者有一種這樣的想法.
想像神.
好像都是選擇性的.
我想跟神說什麼.
心事就說.
但是我不說也可以的.
或者我想神看到我什麼.
我就.
向神面前躺開.
但是如果我的心關起來.
神也沒辦法管的.
有時候我們會這樣想.
但是這裡就說.
私人就發覺其實根本就不是這樣.
他發覺有一種敬畏.
神對我們的認識太奇妙了.
用奇妙這個詞.
可能翻譯不是那麼準確.
是一種.
A-W-E.
就是一種.
很有威嚴之下.
那種奇妙.
而.
我真的至高不能下.
我完全都沒辦法.
好像掌握得到.
或者完全沒辦法理解得到.
神竟然可以這樣.
給我一種關注.

$^{441}$而且這種關注度.
我完全逃避不了.
這個就是私人來到這裡的時候.
有一種體會.
有沒有想過.
神是讓你認識你自己更加的.
更加的認識.
神是讓你每一個的心思.
你自己有時候都自欺欺人也好.
自己都搞不清自己也好.
自己有時候都動機混淆也好.
神完全知道一切.
神對你.
神對你的那種透徹程度.
有時候你自己都驚訝.
為什麼我講了一句這樣的東西.
你自己都搞不清楚.
神完全知道一切.
甚至還沒講之前已經知道.
所以私人就發覺.
我真的要屬於神那種偉大的面前.
除了那種不安之外.
我簡直是五體投地去.
去說服於神面前.
神啊,你真是太奇妙了.
我不能夠用我們這個人的觀念去參測神.
我們有時候某個程度.
覺得自己畫了一條界線.
圈了一個圍欄.
我就覺得很有安全感.
這就是我自己.
然後神就沒辦法侵入.
或者神根本沒辦法去干涉任何東西.
私人體會到.
根本我們的神不是那麼有限.
祂所有都知道了.
祂所有都監測了.
祂所有都是透徹地去明白了.
甚至有時候我們自己都不明白自己的時候.
神完全都知道.

$^{481}$於是私人有一個什麼反應呢?.
祂就說.
我往那裡去逃避你的靈.
我往那裡去躲避你的面.
我若升到天上你在那裡.
我若在陰間下達你在那裡.
我若展開清晨的翅膀飛到海極居住.
就是在那裡你的手必引導我.
你的右手也必扶持我.
我若說黑暗必定遮蔽我.
我周圍的亮光必變為黑夜.
黑暗也不能遮蔽我使你不見.
黑暗卻如白晝發亮.
黑暗和光明在你看到是一樣.
私人的感覺就.
然後祂的反應就是.
我很想躲起來.
因為不太舒服.
不太有安全感.
如果神完全都知道.
祂知道我每一個意念的時候.
不只是有好的意念.
會有很多不好的意念.
甚至可能不好的意念比好的意念多.
只不過我說出來的時候.
我篩選過.
別人聽到覺得好.
覺得我真是一個很好的人.
如果神看得到.
我表面上是一個很好的人.
但原來我每一個裡面有這麼多.
冒出來的原來都是不好的.
所以我想躲起來.
我不想去面對.
我怎麼能夠面對一個這樣的神呢.
所以他說.
我可以逃避你呢.
我可以躲到哪裡去.
他用誇張的文藝走法來說.
我好像小鳥一樣.

$^{521}$飛到天上.
飛到很遠的.
飛到天涯海角.
神就在那裡.
那怎麼辦.
飛到地獄.
或者飛到陰間死亡的地方.
那裡人都不會去的.
但怎麼辦.
神又在那裡.
我好像小鳥一樣.
撐開翅膀.
從一邊飛到海極的另一邊.
他說.
但那裡的小鳥飛.
其實都是在神的管理之下.
神的手都是引導著小鳥飛的.
那怎麼辦.
那我不如躲在黑暗那裡.
好像小朋友玩遊戲.
遇到很害怕的時候.
用布蓋著自己.
躲在床下.
那裡全部都黑的時候.
就沒有人看到.
他說.
但是.
全部都黑的時候.
光明都好像沒有了.
到處都黑黑的.
但是在神看.
黑暗和光明都是一樣的.
神難道要用光照明才可以看到嗎.
黑暗和光明在神看都是一樣的.
根本神又看得到.
哇 斯人就.
有這一種的困局.
他覺得.
那怎麼辦呢.
我怎麼可以逃避神呢.

$^{561}$我每一個的意念.
每一個神都知道.
每一個動機有時候不是那麼善良的.
有時候表面覺得好像很善良.
但是.
你自己也知道.
有時候有些動機其實都是很邪惡的.
表面做出來好像很好.
但是其實背後有很多私心的.
神又知道.
有時候說的話背後有骨頭的.
人家那個人不懂聽的時候還以為你讚他.
原來你在踩他.
有時候在人與人之間說的話.
有時候還傷了別人.
或者有時候你自己的驕傲.
在踩低別人.
每一個說話.
你不想別人聽到.
神又聽到了.
所以他就面對一個困局.
就是這樣的時候.
但是又躲不了.
怎麼辦呢.
躲不了怎麼辦呢.
就是面對一個這樣的神.
他不是一個山神.
不是一個海神.
不是一個地獄神.
是全宇宙的神.
因為你山神.
你就躲到海裡去.
海神就跑到山上去.
地獄神就跑到人間去.
他有一個範圍.
你超出範圍就行了.
但是我們的神就沒有範圍.
沒有人可以超出他的範圍.
那怎麼辦呢.
你面對一個這樣的神.

$^{601}$他的眼目是公義的.
他的眼目是明亮的.
監察一切.
這個比極權國家的監控鏡頭還要恐怖.
24小時.
這個是真實的.
在文化大革命的時候.
你知道那些資源都很少.
那些報紙都會物盡其用.
不會像現在隨時丟掉.
有時候報紙要來偷火.
或者占煲底.
會來用的.
有一次.
那些人都很小心.
因為毛主席是很自告無償.
有一次有個人拿報紙來占煲底.
然後有些人舉報他.
說他不尊重毛主席.
我哪有不尊重.
我24小時都很小心.
因為他占煲底的報紙.
其中背後印了一句毛主席的.
引用了一句毛主席的說話.
就是因為這樣.
然後被人抓去坐牢.
如果是這樣的話.
是避無可避的.
在這樣的情況下.
24小時監控之下都不夠.
有時候你自己都不會知道自己會做錯.
你不會知道自己做錯.
那你說你媽媽是很恐怖.
很恐怖.
但是接下來就很有趣.
詩人就在第13節開始說.
我的肺腑是你所做的.
我在無福中你已浮避我.
我要稱謝你.
因我受造奇妙可畏.

$^{641}$你的作為奇妙.
這是我心知道的.
我在暗中受造.
在地的深處被聯絡.
那時我的形體並未向你隱藏.
我未成形的體質.
你的眼早已看見了.
你所定的日子.
我尚未到一日.
你都寫在你的冊上了.
這裡改變了.
前面那裡.
他想躲起來.
想飛到海極那裡.
在哪裡可以逃避神呢.
因為神什麼都知道.
我又這麼污穢.
我又這麼不堪.
很難受.
所以就躲起來.
但這裡你發覺他改變了.
他說他稱謝神.
為什麼會有這樣的改變呢.
有沒有想過.
他發現了一件事.
我的肺腑是你所做的.
我在無福中你已浮避我.
後面說.
你的作為奇妙.
這是我心知道的.
我在暗中受造.
在地的深處被聯絡.
那時我的形體並不向你隱藏.
我未成形的體質.
你的眼早已看見了.
你所定的日子.
我尚未到一日.
你都寫在你的冊上了.
這裡斯人有一個什麼的發現.
或者有一個什麼的醒悟.

$^{681}$以至到他有一個改變.
他覺得.
你看到他改變了.
他覺得神就不是那麼可怕.
反而神的無所不在.
是令他要多謝神.
要稱謝神.
你覺得這裡有什麼改變呢.
他這裡說什麼呢.
他這裡其實是說一件事.
他體會到一件事.
就是原來.
神雖然是什麼都知道.
但是神知道.
不單止是知道.
而且是神主動的參與地來知道.
就好像一個嬰兒.
他說到一個嬰兒在母親的肚裡.
神已經慢慢慢慢來到去.
不單止是在遠方來監察.
是在我後備.
這個後備的意思.
就是保護令他很安全.
就像一個嬰兒.
在母親懷胎的時候.
特別是在古代的時候.
你知道有時候都會夭折的.
不是每一個嬰兒都會順產的.
在母覆胎裡成形的時候.
在深處被聯絡.
我的形體並不向你隱瞞.
意思就是說.
嬰兒在母覆胎裡長育的時候.
可能根子和骨頭.
在他想像中要聯合在一起.
慢慢成為一個人的階段的時候.
神不單止是暗中監察.
其實神是積極的保護.
去賜福他.
令他可以出生.

$^{721}$可以健康.
他體會到這一點.
當他體會到這一點的時候.
他說我還沒有成形的體質.
但神已經看見了.
那時候還沒有出生.
我還像一顆蝌蚪一樣.
慢慢成為一個人.
後來出生的時候.
整個過程神不單止是看著.
神後備著.
神保護著.
神祝福著.
神也是親自引導著.
以致後來成為一個生命.
誕生在這個世界裡.
然後他慢慢長大成為了現在的私人.
他體會到這一點.
以致他讚嘆神.
這個改變在哪裡呢?.
這個改變就是.
私人體會到神的無所不在.
神的無所不知.
神自己監察的眼目.
就由負面變成正面.
為什麼可以這樣.
由負面變成正面呢?.
第一.
就是神的監察.
神的知道.
不是想傷害的.
是想保護的.
這一點就很不同.
你想想.
如果你去到哪裡.
有一個人不斷跟著你.
在後面.
不斷地打醒十二個精神.
到處看,到處監察著.
你會不舒服的.

$^{761}$但是如果那個人跟著你.
是你的保鑣.
他到處看.
你還沒看到那裡.
他已經幫你看了.
你去到哪裡.
他保護著.
那你就很舒服了.
這就完全將.
神的無所不在的監察.
顛覆了.
因為.
神的無所不在的監察.
神的周圍.
那樣東西知道.
不是想傷害你的.
是想保護你的.
神對我們所有的意念.
都是善的.
不是審判的.
當然.
我們的罪惡神會審判.
但是最終審判完之後.
神還是想.
我們所說的意念都是善的.
都是愛的意念.
都是美好的意念的時候.
那你就很放心地讓神知道.
因為神的知道.
不是那種invasive.
好像那些黑客.
黑客那些入侵你的電腦.
說知道你很多東西.
那你就大件事了.
當然要避免.
但是神不是黑客.
但是問題是.
很多時候我們把神想成黑客.
說入侵.
知道會傷害我們.

$^{801}$那你就會覺得.
很不舒服.
但是原來神的意念是保護的.
神的意念是保護的時候.
你就不是這樣看.
第二.
他體會到那個BB仔.
在母胎裡成形的過程.
然後後來.
來到他出生的時候.
神已經知道監察.
他體會到另外一樣東西.
就是.
神的完全的接納.
有時候.
你.
不是那麼舒服.
你的意念.
舉例來說.
如果你有寫日記.
你的日記.
也不是那麼容易想被人看.
為什麼.
因為如果你被人看了.
你不知道那個人.
我不說他公開的那些東西.
只是他自己看.
你也不知道那個人對你有什麼評價.
是不是.
他看了你的日記.
可能你在那篇日記說.
今天很懷疑自己.
考試又不合格.
我也不知道有沒有價值.
有時候你把自己的心情寫下來.
如果那個人.
你也不知道他看完之後會怎樣.
可能他看完之後.
這個人原來這麼自卑.
看他表面好像.

$^{841}$開朗.
好好健談.
開放.
原來他背後這麼自卑.
你也很怕.
那些.
這樣的評價.
別人加在你身上.
是不是.
這就是另外一個.
覺得被人知道的一個危機.
或者.
不安全的成分.
就在其中.
這個世界人知道你越多.
你不會知道他對你的想法是怎樣的.
是正面是負面.
是傷害是接納.
你完全不知道的時候.
你有這個危機的時候.
你當然不想讓人知道這麼多.
但是.
聖經這裡.
詩人就明白到.
神是完全的接納.
神完全知道.
他在我沒有胎的時候.
他已經接納我.
他已經引導我了.
他怎會不接納我呢.
這個.
也是一個.
很大的釋放.
你可以放心.
神知道你一切.
但是他知道.
第一他剛才所說的.
他是善的.
更加你可以放心.
他知道你一切.

$^{881}$他不是用你所知道的.
用他所知道的.
來負面的評價你.
人在世界上.
是需要有一種的接納.
人的自我形象本身.
是很大部分.
是和.
怎樣被人接納.
去建造起來的.
這個大家都明白.
如果一個人.
從小到大都不被人接納的.
可能她是一個美女.
但是人家.
一讚她美的時候.
她都覺得.
人家是在揶揄她的.
都會有的.
反過來說.
一個普普通通的.
甚至都不是美的.
她一走出街.
她就自信心爆棚.
她自己都是一個.
尋魚落雁.
西施下凡.
都可以的.
為什麼?.
因為從小到大.
她父母.
特別是她爸爸.
不停地讚她.
哇,你是美女.
人見人愛,車進車載.
她信以為真.
因為她的自信心.
就是這樣建立起來的.
我們每個人都需要.
有人的接納.

$^{921}$但是最終我們也都.
有很多的保留.
因為我們都不敢完全的開放.
因為.
你就知道.
未必是可以得到完全的接納.
但是在神的面前就不是這樣.
私人都體會到.
那位完全知道我們的神.
也是完全接納我們的神.
那個完全.
來到由始至終.
去看著我們的神.
也是完全的保護我們的神.
所以當她明白到這一點的時候.
她就去讚賞.
去稱謝神.
誰可以比神更好呢?.
在這個世界上.
有時候你都想找到一個知心的朋友.
中國人說.
得一知己,死而無憾.
但是知己就是.
那個朋友.
為什麼叫知己?.
就是完全的知道我.
但是你有時候找不找得到?.
你想找都未必找得到.
但是神就是這位這樣的知己.
神就是這位完全的知道.
所以她說.
我尚未到一日.
你都寫在你的冊上.
完全不在乎我的行為.
可以這麼說.
有時候我們活在世上的時候.
我們的行為.
我們的表現.
我們能夠達到.
有什麼人生的成就.

$^{961}$好像都成為了我們自己對自己的評價也好.
別人對我們的評價也好.
但是私人在神面前就說.
完全都不是這樣的一回事.
我都未做任何的事情.
如果我的每一天.
就好像神寫在冊子上.
那個冊子都還沒開始寫.
我尚未到一日.
但是神就已經知道了.
神已經都是接納了.
所以你會看到.
接下來她就說.
神啊,你的意念向我何等寶貴呢?.
她之前想躲起來,想走的.
說怎樣逃避神的.
當她看到神.
她真正明白到神的時候.
她就發覺不是了.
她說你的意念向我何等寶貴?.
其數何等眾多?.
我日數點比海沙更多.
我睡醒的時候迎和你同在.
多麼好呢?.
就好像一個愛人或戀人的意念.
就好像初初拍拖的時候.
想著對方.
你知不知道我今天想了你多少次?.
日朝又想著你.
我做工又想著你.
開車又想著你.
我都不知道有多少次對你的想法.
就是這一種觀念.
如果你真心愛的人.
你都不想他想你那麼多.
你不要想我.
你out of sight, out of mind 最好.
但你對於你愛的人.
你就說你越想我越好.
而且你想我的每一個意念都很寶貴.

$^{1001}$我都很珍惜.
我想起你這樣想.
我已經覺得很甜蜜.
而且不單是想像出來的.
我睡覺你又與我同在.
我起來的時候你也與我同在.
所以這個就完全改變.
顛覆了之前所說的那種對神的逃避.
有時候我們不明白神的時候.
我們就覺得神甚麼都知道是很威嚇性的.
很受到威脅的.
很不自在的.
但你明白到神知道你多麼的醜惡.
知道你多麼的不濟.
知道你多麼的軟弱.
但你不怕躺開.
因為神既然知道.
但他也完全的接納.
而且他每一個意念就是想著你.
多麼的寶貴.
神對於你的每一個意念.
都是善,都是愛,都是美的意念.
多麼的寶貴.
所以他說你對於我每一個意念是何等的寶貴.
好像海沙一樣多.
多麼的美.
然後這裡就說.
好像突然間話氣一轉的時候.
很多人就不明白究竟這裡說甚麼呢.
為甚麼無端端說說.
好像愛人拍拖,戀人拍拖的時候.
就想起神那麼愛我.
很多意念的神.
然後現在就說殺人.
神啊,你必要殺戮我人.
所以你們好流人血的離開我去吧.
因為他們說我然頂撞你.
你的仇敵也罔稱你的名.
而我恨惡你的.
我豈不恨惡他們嗎.

$^{1041}$攻擊你的.
我豈不憎厭他們嗎.
我要切切地恨惡他們.
以他們為仇敵.
這裡究竟是做甚麼呢.
他無端端說戀人拍拖.
想起神好寶貴.
現在就走來罵人殺人.
為甚麼呢.
其實這裡某個程度都是一個很合理的轉變.
這個不是轉變.
其實是一個很合理的延伸.
當施人體會到神對他的意念是多麼美.
他的心向神開.
他下一步就發現.
他需要和神同步.
神所喜惡的他也要喜惡.
神所喜歡的他也要學習喜歡.
神所憎恨的他也要學習憎恨.
我們的神是公義的神.
我們的神是恨惡邪惡的神.
所以神說你憎恨的那些人我也憎恨.
你不喜歡的邪惡我也不喜歡.
你不喜歡的那些人我也和他一刀兩斷.
你不喜歡的那些人我也不想和他們走在一起.
這個就是很順著心理狀況的一種延伸.
當你經常感受到神的同在.
當你經常欣賞和經歷到神與你同在的時候.
你是很自然地想去愛神所愛的.
很想去珍神所珍的.
一方面你說神對我多麼好.
對我多麼的美善.
但另一方面你不斷去想那些很骯髒邪惡的東西的時候.
行不行呢?.
根本是格格不入的.
是不是?.
根本是格格不入.
而事實上我們就是這樣.
我們一方面越是經歷神的同在.
越是經歷到神對我們的愛的寶貴.

$^{1081}$越是經歷到神對我們無窮無條件的接納的時候.
我們也會開始去愛神所愛.
去恨神所恨.
那個不是一個交換式的條件.
說我這樣這樣這樣神才會愛我.
不是這樣.
而是說神根本是已經愛我在先.
但我的心現在就好像願意與他同步.
就好像你與很親密的朋友.
你與他越走越近的時候.
那些愛好那些話題.
好像都是現在同步起來的一樣.
所以最後的那個禱告就顯得完全合理.
所以神求你監察我.
知道我的心思.
試煉我知道我的意念.
看在我裡面有什麼惡行沒有.
引導我走永生的道路.
一開始的時候他只是說出了一個敘述.
就是說神監察一切的東西.
但現在結束的時候.
他是祈禱求神監察他的內心.
你看到這個是很不同的.
一開始的時候他只是說出了神就什麼都知道.
我又沒辦法逃避.
想走都走不了.
但去到這裡的時候.
當我明白到神的愛神的接納.
神那種無微不至的關注.
眼目的不離開的時候.
他更加想神看我多一點.
監察我的內心.
究竟我的心思意念裡面有沒有一些不合理心意.
我想走永生的路.
跟你同步.
求神幫助我們.
求神監察我們的內心.
你做基督徒.
你永遠都是要從心裡面開始的.
你在外面可以做很多的.

$^{1121}$人做的事.
人參與的活動.
但是你跟神的關係永遠都是從你的內心開始.
而你的內心神完全知道.
唯一你可以求的.
不是求神看不到.
是求神看你多一點.
因為你這樣的時候.
你跟神同步的時候.
我們就是走在永生的路上.
求神幫助我們.
用這個禱告成為我們的禱告.
我們一起讀吧.
神啊 求你監察我.
知道我的心思.
試煉我.
知道我的意念.
看在我裡面.
有什麼惡行沒有.
引導我走永生的道路.
\newpage



\section{使徒行傳 27:1-28:1}
\label{sec:032016}
\textbf{征服世界的復活救主}~~~~林鍵牧師
\newline
\newline
連結: \href{http://www.church.com.hk/acms/content.asp?site=cdc\&op=show\&type=product\&code=032016\&layout=sermon}{\texttt{032016}} ~~~~ 語音日期: 2026-04-18
\newline
\newline
\hyperref[sec:032015]{\small{< < < PREV SERMON < < <}}
~
\hyperref[sec:index_chronic]{\small{[返順時目]}}
~
\hyperref[sec:032017]{\small{> > > NEXT SERMON > > >}}
\newline
\newline
使徒行傳 27:1-28:1
\newline
\begin{longtable}{cl}
\hline
\hline
章節 & 經文 (和合本修訂版)\\
\hline
27:1 & \begin{tabularx}{0.7\textwidth}{X} 既然非斯都決定要我們坐船往意大利去，就將保羅和別的囚犯交給御營裡的一個名叫猶流的百夫長。 \end{tabularx} \\ \\ \relax
27:2 & \begin{tabularx}{0.7\textwidth}{X} 有一隻亞大米田的船要開往亞細亞沿海一帶地方去，我們上了那船，就起航了；有馬其頓的帖撒羅尼迦人亞里達古和我們同去。 \end{tabularx} \\ \\ \relax
27:3 & \begin{tabularx}{0.7\textwidth}{X} 第二天，我們到了西頓。猶流寬待保羅，准他往朋友那裡去，受他們的照應。 \end{tabularx} \\ \\ \relax
27:4 & \begin{tabularx}{0.7\textwidth}{X} 我們又從那裡開船，因為遇到逆風，就貼著塞浦路斯的背風岸航行， \end{tabularx} \\ \\ \relax
27:5 & \begin{tabularx}{0.7\textwidth}{X} 渡過了基利家、旁非利亞一帶的海面，就到了呂家的每拉。 \end{tabularx} \\ \\ \relax
27:6 & \begin{tabularx}{0.7\textwidth}{X} 在那裡，百夫長找到一隻亞歷山大的船要往意大利去，就叫我們上了那船。 \end{tabularx} \\ \\ \relax
27:7 & \begin{tabularx}{0.7\textwidth}{X} 一連多日，船行得很慢，我們好不容易才來到革尼土的對面；又因被風攔阻，我們就貼著克里特島背風岸，從撒摩尼對面航行。 \end{tabularx} \\ \\ \relax
27:8 & \begin{tabularx}{0.7\textwidth}{X} 我們沿岸前進，十分艱難，來到一個名叫佳澳的地方，離那裡不遠有拉西亞城。 \end{tabularx} \\ \\ \relax
27:9 & \begin{tabularx}{0.7\textwidth}{X} 航行的日子久了，已經過了禁食的節期，行船又危險，保羅就建議， \end{tabularx} \\ \\ \relax
27:10 & \begin{tabularx}{0.7\textwidth}{X} 對眾人說：「諸位，我看這次航行，不但貨物和船要受損傷，大遭破壞，連我們的性命也難保。」 \end{tabularx} \\ \\ \relax
27:11 & \begin{tabularx}{0.7\textwidth}{X} 但百夫長信從船長和船主，不信保羅所說的。 \end{tabularx} \\ \\ \relax
27:12 & \begin{tabularx}{0.7\textwidth}{X} 且因在這港口不適宜過冬，船上大多數的人都主張開船離開這地方，或者能到非尼基去過冬。非尼基是克里特的一個港口，一面朝西南，一面朝西北。 \end{tabularx} \\ \\ \relax
27:13 & \begin{tabularx}{0.7\textwidth}{X} 當南風微微吹起時，他們以為對目的地已有了把握，就起錨，貼近克里特開去。 \end{tabularx} \\ \\ \relax
27:14 & \begin{tabularx}{0.7\textwidth}{X} 過了不久，有一股叫「友拉革羅」的東北巨風從島上撲來， \end{tabularx} \\ \\ \relax
27:15 & \begin{tabularx}{0.7\textwidth}{X} 船被風抓住，無法頂風航行，我們只好任它漂流。 \end{tabularx} \\ \\ \relax
27:16 & \begin{tabularx}{0.7\textwidth}{X} 我們貼著一個叫高大的小島的背風岸急航，好不容易才保住了救生艇。 \end{tabularx} \\ \\ \relax
27:17 & \begin{tabularx}{0.7\textwidth}{X} 既然把救生艇拉上來，他們就用纜索捆綁船底，又恐怕在賽耳底淺灘上擱淺，就落了篷，任船漂流。 \end{tabularx} \\ \\ \relax
27:18 & \begin{tabularx}{0.7\textwidth}{X} 我們被風浪逼得很急，第二天眾人就把貨物拋在海裡。 \end{tabularx} \\ \\ \relax
27:19 & \begin{tabularx}{0.7\textwidth}{X} 第三天，他們又親手把船上的器具拋棄了。 \end{tabularx} \\ \\ \relax
27:20 & \begin{tabularx}{0.7\textwidth}{X} 許多天都沒有看到太陽和星辰，又有狂風大浪催逼，我們獲救的指望都放棄了。 \end{tabularx} \\ \\ \relax
27:21 & \begin{tabularx}{0.7\textwidth}{X} 眾人已有好幾天沒有吃東西，保羅就出來站在他們中間，說：「諸位，你們本該聽我的話不離開克里特島，就不致遭到這樣的損失和破壞。 \end{tabularx} \\ \\ \relax
27:22 & \begin{tabularx}{0.7\textwidth}{X} 現在我勸你們放心，除了損失這條船，你們中間沒有一人會喪失性命。 \end{tabularx} \\ \\ \relax
27:23 & \begin{tabularx}{0.7\textwidth}{X} 因為昨夜，我所屬所事奉的神的使者站在我旁邊， \end{tabularx} \\ \\ \relax
27:24 & \begin{tabularx}{0.7\textwidth}{X} 說：『保羅，不要害怕，你必定站在凱撒面前；並且神已把安全賜給與你同船的人了。』 \end{tabularx} \\ \\ \relax
27:25 & \begin{tabularx}{0.7\textwidth}{X} 所以，諸位可以放心，我信神怎樣對我說，事情也要怎樣成就； \end{tabularx} \\ \\ \relax
27:26 & \begin{tabularx}{0.7\textwidth}{X} 只是我們必須在一個島上擱淺。」 \end{tabularx} \\ \\ \relax
27:27 & \begin{tabularx}{0.7\textwidth}{X} 到了第十四天夜間，船在亞得里亞海漂來漂去。約在半夜，水手以為漸近旱地， \end{tabularx} \\ \\ \relax
27:28 & \begin{tabularx}{0.7\textwidth}{X} 就去探測深淺，探得有十二丈；稍往前行，又探深淺，探得有九丈。 \end{tabularx} \\ \\ \relax
27:29 & \begin{tabularx}{0.7\textwidth}{X} 恐怕我們撞到礁石，他們就從船尾拋下四個錨，盼望天亮。 \end{tabularx} \\ \\ \relax
27:30 & \begin{tabularx}{0.7\textwidth}{X} 水手想棄船逃走，把救生艇縋下海裡，假裝要從船頭拋錨的樣子。 \end{tabularx} \\ \\ \relax
27:31 & \begin{tabularx}{0.7\textwidth}{X} 保羅對百夫長和士兵說：「這些人若不留在船上，你們就不能獲救。」 \end{tabularx} \\ \\ \relax
27:32 & \begin{tabularx}{0.7\textwidth}{X} 於是士兵砍斷救生艇的繩子，由它漂去。 \end{tabularx} \\ \\ \relax
27:33 & \begin{tabularx}{0.7\textwidth}{X} 天快亮的時候，保羅勸眾人都用餐，說：「你們一直捱餓等候，不吃甚麼，已經十四天了。 \end{tabularx} \\ \\ \relax
27:34 & \begin{tabularx}{0.7\textwidth}{X} 所以我勸你們吃點東西，這是關乎你們獲救的，因為你們各人連一根頭髮也不至於掉落。」 \end{tabularx} \\ \\ \relax
27:35 & \begin{tabularx}{0.7\textwidth}{X} 保羅說了這話，就拿起餅來，在眾人面前祝謝了神，然後擘開來吃。 \end{tabularx} \\ \\ \relax
27:36 & \begin{tabularx}{0.7\textwidth}{X} 於是他們都放心，就吃了。 \end{tabularx} \\ \\ \relax
27:37 & \begin{tabularx}{0.7\textwidth}{X} 我們在船上的共有二百七十六個人。 \end{tabularx} \\ \\ \relax
27:38 & \begin{tabularx}{0.7\textwidth}{X} 他們吃飽了，為要使船輕一點，就把船上的麥子拋到海裡。 \end{tabularx} \\ \\ \relax
27:39 & \begin{tabularx}{0.7\textwidth}{X} 天亮的時候，他們不認得那地方，只見一個有岸可登的海灣，就想法子看能不能把船靠岸。 \end{tabularx} \\ \\ \relax
27:40 & \begin{tabularx}{0.7\textwidth}{X} 於是他們砍斷纜索，把錨丟到海裡，同時也鬆開舵繩，拉起頭篷，順風向著岸行去。 \end{tabularx} \\ \\ \relax
27:41 & \begin{tabularx}{0.7\textwidth}{X} 但碰到兩水夾流的地方，就擱了淺，船頭膠住不動，船尾被浪的猛力衝壞了。 \end{tabularx} \\ \\ \relax
27:42 & \begin{tabularx}{0.7\textwidth}{X} 士兵的意思要把囚犯都殺了，免得有游水脫逃的。 \end{tabularx} \\ \\ \relax
27:43 & \begin{tabularx}{0.7\textwidth}{X} 但百夫長要救保羅，不准他們任意而行，就吩咐會游水的，跳下水去，先上岸； \end{tabularx} \\ \\ \relax
27:44 & \begin{tabularx}{0.7\textwidth}{X} 其餘的人則用板子或船的碎片上岸。這樣，眾人都獲救，上了岸。 \end{tabularx} \\ \\ \relax
28:1 & \begin{tabularx}{0.7\textwidth}{X} 我們既已獲救，才知道那島名叫馬耳他。 \end{tabularx} \\ \\
[1ex]
\hline
\hline
\end{longtable}
$^{1}$歡迎大家來到我們的周六崇拜.
今天跟平時有點不一樣.
因為今天的讀經和講道會融會在一起.
因為今天的整段經文.
其實都是有關保羅前往羅馬海上的旅程.
可能比較讓我們明白.
一邊是經文,一邊是地圖.
Wilfred會幫我讀幾節聖經.
然後我會有一點講解.
Wilfred繼續讀,我繼續講解.
開始的時候,時間大約是公元59至60年左右.
保羅當時因為猶太人的逼迫.
在耶路撒冷聖殿被捕.
後來轉押到凱撒利亞這個地方被囚.
有兩年時間.
所以開始的時候我們是在這裡.
Caesarea,我不知,希望在這裡比較快.
我們在凱撒利亞這個地方開始.
於是Wilfred幫我讀頭兩節.
既然菲斯都決定要我們坐船往意大利去.
就將保羅和別的囚犯交給預迎他的一個名叫猶留的白夫長.
有一隻亞大米歇的船要開往亞細亞沿海一帶地方去.
我們上了那船就起航了.
有馬其頓的鐵薩羅里嘉人阿里達古和我們同去.
所以我們開始的時候我們在凱撒利亞這裡.
他要去意大利,意大利的左邊還沒出來的地方.
如果你大概知道歐洲地理的時候.
他們就上了一艘船.
這艘船叫做在亞大米田.
亞大米田在這裡上面.
一艘船在那裡.
當然現在這艘船在凱撒利亞.
他們要回去亞細亞不是我們的亞洲的Asia.
總之這艘船應該是回去亞細亞那邊的地方.
如果你留意你會發現一開始的時候.
路加這樣寫的.
他說我們坐船往意大利去.
原來路加也有份的.
這個是聖經裡面很出名的.
我們叫做We Passages.

$^{41}$我們的經文.
當時路加醫生是陪伴著保羅.
而且除了他們兩個之外.
還有一個在馬其頓的鐵薩羅里嘉人叫阿里達古和他們一起去.
好我們繼續讀.
第二天我們到了西頓.
由留歡待保羅.
準他王朋友那裡去.
受他們的照應.
我們又從那裡開船.
因為遇到逆風.
就貼著塞浦路斯的背風岸旁行.
渡過了基里加.
龐菲尼亞一帶的檯面.
就到了旅加的貝拉.
這裡有些複雜.
我們第一天一去就去了西頓.
就是上去.
這是一天很小的路程.
然後他們就去了塞浦路斯.
裡面上基里加.
然後是龐菲尼亞.
這裡是龐菲尼亞.
基里加就在這裡.
塞浦路斯就在這裡.
就到了這裡.
大概知道在哪裡就好了.
繼續去的時候.
就到了梅加的這個地方.
梅拉這個地方.
好繼續.
在那裡.
白夫長找到一隻亞歷山大的船要往意大利去.
就叫我們上了那船.
一年多日船行得很慢.
我們沒什麼意義才來到格里普的對面.
由於被風擋住.
我們就貼著克里特島背風岸.
從薩摩尼對面航行.
OK.

$^{81}$原來當時你沒辦法坐一艘船去哪裡就去哪裡.
沒得 Uber 它.
於是它到了這裡就要.
找另一艘船.
就去意大利.
這艘船就由亞歷山大這裡來的.
就是亞歷山大.
其實是非洲那裡.
下來的.
但是它是前往意大利.
OK.
如果我們繼續看.
有幾個地方.
它總之.
就這裡來到這裡左右.
就是現在是在海面上.
來到克里特島的薩摩尼.
好我們繼續.
(來到一個名叫佳奧的地方).
(離那裡不遠有拉西亞城).
(航行的日子夠了).
(已經過了金城的節期).
(行船很危險).
(普羅就建議對峙人士).
(諸位我看這次航行).
(不但貨物和船要受損傷).
(大遭破壞).
(連我們的性命也難保).
OK.
所以它繼續在這裡.
來到一個地方叫佳奧.
但是到了這裡.
天氣就已經開始變得很差了.
這裡他說.
他們經過了一個禁食的節日.
相傳禁食的節日.
通常在舊約聖經裡面.
只有一天需要禁食.
我不知道你知不知道.
其實以色列人只有一天必須禁食.

$^{121}$就是贖罪日.
也有學者說.
以色列人到了現代的傳統.
在贖罪日的下午.
會讀一本舊約的書信叫做《若拿記》.
這個很有趣的.
因為如果當日普羅也有讀《若拿記》.
就很有趣了.
因為之後的故事.
普羅在海面上的故事.
很像《若拿記》.
到了這裡.
普羅當時就跟眾人說.
他說諸位我看這次旅行.
不但貨物和船都要受損傷.
大遭破壞.
連我們的性命也難保.
這個不是性靈感動.
而是普羅自己的經驗老道.
如果你有機會去看.
過不多後書第十一章第十五節.
普羅說他曾經三次遭海難.
三次沉船.
所以他一看當時的情況.
他知道今天不行了.
要麼你不要走.
你一走我們死定了.
我們的貨物都沒有了.
之外.
我們性命也難保.
繼續.
但伯夫長順從船長王船主.
不信普羅所說的.
且因在這港口不適而過冬.
船上大多數的人.
都主張開船離開這地方.
或者能到菲利基去過冬.
菲利基是克里特的一個港口.
一面朝西南.
一面朝西北.

$^{161}$OK.
他們已經到了一個氣候改變的時候.
其實.
這個是二千多年前的事情.
你知道當時要航海是很困難的.
基本上每年有幾個月是沒有辦法航海的.
他們已經很問尾了.
海上的情況非常不穩定.
他們要找一個地方去過冬.
OK.
他們已經到了這裡.
但他們想再移動一點.
去到另外一個地方.
他覺得在菲利亞那裡過冬會好一點.
就是因為這個決定.
結果使他們最後翻船.
我們繼續讀.
他們開始的時候以為微微風應該不是很遠.
可能半天一天可以到達這裡.
誰知道一開船之後.
就八號風球十號風球.
於是他們的船現在被人一直吹著走.
飄離了他們的岸邊.
被他抓住.
都不知道飄流到哪裡了.
繼續.
我們貼著一個叫高大的小島的背風岸急航.
很不容易才抓住了救生艇.
既然把救生艇拉上來.
他們就用纜索綑綁船底.
又恐怕再坐在二底淺灘上擱淺.
就下了風任船飄流.
我們被風浪迫得很急.
第二天眾人就把貨物拋在海裡.
OK.
他們本來想去北一點的這裡.
誰知道現在去了高大的小島.
就是說他們已經來了.
他們很怕會來到這裡擱淺.
那個地方就叫做塞爾底擱淺.

$^{201}$當時他們連帆布都下來了.
為什麼呢?.
因為帆布不斷吹.
就越吹得快.
於是下航了.
希望不會飄流得那麼遠.
但是基本上是聽天由命.
到了第二天.
第十八節.
告訴你.
第二天已經開始把貨物扔進海裡.
你要知道當時所有的航運都是因為送貨的緣故.
今天我們會坐飛機.
然後我們的東西是搭班的.
但不是的.
那時候所有東西其實都是為了.
那些財物是最重要,最值錢的東西.
但是他現在都不管了.
為什麼呢?.
因為性命攸關.
隨時死人了.
第三天.
他們又親手把船上的器具拋棄了.
許多天都沒有看到太陽和星辰.
又有狂風大浪吹逼.
我們獲救的指望都放棄了.
眾人已有好幾天沒有吃東西.
保羅就出來說.
諸位.
你們本該聽我的話不離開克里特島.
就不至遭到這樣的損失和破壞.
第三天連其他器具都扔了.
其實是最後的希望.
但是他已經很多天沒有看到太陽和星辰.
這次十號颶風.
前所未有.
你想想當時的東西.
他沒有任何航運的儀器.
他只可以看天,看星.
但是現在天,星,太陽什麼都看不到.

$^{241}$那就等死了.
現在我勸你們放心.
除了損失者條船.
你們中間沒有一人會喪失性命.
因為昨夜我所事奉的上帝,的師者.
站在我旁邊說.
保羅不要害怕.
你必定站在凱撒面前.
並且上帝要以把安全之吉與你同船的人了.
這次不是靠自己的經驗.
這次是上帝啟示.
告訴你儘管是如此無望的情況.
這艘船上不會有任何一個人喪失生命.
稍後你會看到.
他連人頭都數了.
所以諸位可以放心.
我信上帝怎樣對我說.
事情也要怎樣成就.
只是我們必須在一個島上擱淺.
到了第十四天夜間.
船在亞德里亞海飄來飄去.
約在半夜.
水手以為漸近港地.
就去探測深淺.
探得有十異象.
稍往前行.
又探深淺.
探得有九象.
好,繼續.
恐怕我們撞到礁石.
他們就從船尾拋下四個盲.
盼望天亮.
水手想棄船逃走.
把救生艇追下海裡.
假裝要從船頭拋下一樣子.
保羅對白夫長和士兵說.
這些人若不留在船上.
你們就不能活救.
於是士兵砍斷救生艇的繩子.
由他飄去.

$^{281}$有人想自己走.
保羅說不行.
要麼我們一起.
或者其實他的意思是.
你不跟我一起你就死定.
待會我們有機會去看看.
背後的一些神學思想.
保羅基本上說.
你離開我你就死定了.
這樣.
不可以離開.
天快亮的時候.
保羅勸眾人都用餐吃.
你們一直捱我等候.
不吃什麼.
已經十四天了.
所以我勸你們吃點東西.
這是關乎你們活救的.
因為你們各人連一根頭髮.
也不至於掉落.
保羅說了這話.
就拿起餅來.
在眾人面前祝謝了上帝.
然後張開來吃.
很有耶穌的影子.
但是很誇張.
他說你們連一根頭髮也不會掉落.
在一個.
基本上大家覺得無忘.
一定會死人的時候.
他說一個也不會少.
這樣也可以.
於是他們都放心就吃了.
我們在船上的共有276個人.
他們吃飽了.
為要使船輕一點.
就把船上的墨紙拋到海裡.
天亮的時候.
他們不認得那地方.
只見一個有按摩燈的海灣.

$^{321}$就想法子能不能把船拋下去.
276人.
這個船也很大.
接著他們吃最後一餐.
吃完之後連餘下的食物也丟掉.
現在就是.
聽天由命.
等上帝去拯救他們.
知道船一定玩完.
但是保羅說沒有一個人會滅亡.
這樣也可以.
於是他們撼斷纜索.
下游丟到海裡.
同時也鬆開繩.
拖繩.
拉起頭縫.
順風向著岸走去.
但到兩水隔流的地方就割了繩.
船頭交轉不動.
船尾被浪的猛力沖壞了.
是必的意思要把囚犯都殺了.
免得有游水逃脫的.
最後.
但白夫長要救保羅.
不准他們任意而行.
就吩咐會游水的跳下水去.
先上岸.
其餘的人則用板子或船的碎片上岸.
這樣眾人都獲救上了岸.
我們既已獲救.
才知道那島名叫馬爾他.
所以所有人都得救.
276位.
接著他們現在去到.
馬爾他.
羅馬就在上面.
OK.
好.
多謝Wilfred.
我們現在開始.

$^{361}$重播我們這個講道.
雖然你會發覺到這個位置的時候.
他們還沒有去到羅馬.
但好像你看電影別人會說.
病人已經渡過了危險期.
保羅已經上了岸抵達羅馬.
此日可待.
我想你回想我們現在讀的這本書.
我們讀的是《仕途行傳》.
更加準確來說.
我們正在讀一本叫做.
《路加福音》加《仕途行傳》.
是一個一分為兩卷的合一作品.
其實是一本書.
這本書有上冊有下冊.
這本著作一書兩冊的著作.
如果你記得.
早兩星期我已經開始說.
可以被分為三大部分.
第一部分說的是耶穌基督.
第二部分是祂差遣聖靈.
第三部分是成為了《仕途行傳》.
每個部分的底下都有兩個空格.
當你填好這些空格的時候.
你會更加清楚.
明白這兩冊.
《路加福音》和《仕途行傳》的大綱和主題.
我們開始前衝.
耶穌基督為我們做了些什麼.
祂為我們受死復活.
祂跟著升天作王.
祂的受死是為我們贖罪.
祂的復活是勝過死亡.
跟著祂升天祂作王.
於是祂差遣祂的聖靈來做些什麼.
不是幫你考試.
不是幫你做公.
祂的聖靈要創造以馬五師的一個經驗.
使罪人可以得到重生.
使罪人可以死裡復活.

$^{401}$成為一個復活的聯盟.
如果你在我們當中.
你知道對於我們教會來說.
我們一向都很重視以馬五師路上.
這兩個門徒.
核眼心死的門徒.
最後他們成為了眼明心火熱的門徒.
以致我們今天會有《仕途行傳》.
如果你沒有這一群眼明心火熱的門徒.
是不會有《仕途行傳》的.
你只是一群Catafalque在那裡.
創造不知道什麼群體.
但是要真正有力量的教會.
是需要眼明心火熱.
於是《仕途行傳》開始.
他們最後的目的是要地極作見證.
宣揚上帝的天國.
《仕途行傳》第一章第八節.
「但聖靈將要降臨在你們身上.
你們必要得著能力.
並要在耶路撒冷,猶大,全地,撒瑪利亞.
直到地極為我作見證」.
正如我剛才告訴你們.
《仕途行傳》可以分為兩半.
第一半是一至十二章.
第二半是十三章至二十八章.
頭一半主角,主要配角是彼得.
沒有人是主角.
上帝是主角.
頭一半主要配角是彼得.
福音事公圍繞著耶路撒冷的教會.
猶大的信徒.
然後到後一半.
第十三章到二十八章.
主要配角是掃羅,保羅.
福音事公圍繞著耶路撒冷.
猶大,撒瑪利亞以外外邦人的信徒.
在這裡第十二章尤其重要.
所以我們花了兩個主日.
去明白《仕途行傳》第十二章的重要性.

$^{441}$在《仕途行傳》第十二章裡.
我們看到彼得的受死復活.
接著到尾聲的時候.
我們看到保羅又入框.
於是掃羅和彼得同框.
可以說是兩個主角.
兩個主要配角交棒的時候.
在《仕途行傳》第十一章的尾聲裡.
我們看到掃羅來到耶路撒冷.
因為要送捐款的原因.
年輕當時還三十多歲.
未夠四十歲的掃羅來到耶路撒冷.
接著當他最後離開的時候.
第十二章第十四,二十五節的時候.
他親眼目睹以上所發生的一切.
他學了很多東西.
這個旅程我不知道要去多久.
但他學到他看到彼得的受死復活.
他看到初期教會如何的禱告.
他開始明白究竟他將來的福音事供是怎樣的.
我們用了三點來總結年輕掃羅.
在第十二章裡所學到一生受用的寶貴功課.
第一個他學的是切切禱告放膽傳講.
主道興旺 越發廣傳.
當彼得被人抓走的時候.
他們的祈禱不是快點放他出來.
他們的祈禱是使主道興旺 越發廣傳.
第二他親眼目睹.
雅各的殉道 彼得的逃生.
但他知道主必拯救所有他的僕人.
脫離敵人的手.
這是我們過去兩個星期都很努力學習可以記得的事.
接著上個星期我們開始看這個很重要的事供.
就是基督上升 撒旦墜落.
新約的聖經有一個使命是逃出埃及.
我不知道你還記不記得上個星期我們去看整個十二章.
是充滿著出埃及記的影子.
我想大家特別留意第三點.
無論是彼得或是保羅的事供.
他們福音事供的核心就是出埃及記.

$^{481}$無論是彼得或是保羅的事供.
就是要帶一群迴轉悔改.
信靠跟隨耶穌的上帝子民.
拋棄地上的偶像.
去事奉得到自由.
去事奉一個獨一的真神.
這個福音事供的核心出埃及記.
令我想起保羅在以弗所城裡面兩三年的事供.
以及基督福音對於以弗所這群基督徒.
以及整個以弗所城所創造的巨大影響.
我想帶大家去看看以弗所的出埃及記.
《紹行傳》第十九章第十節.
這樣有兩年之久.
再加上前面幾個月後面幾個月.
差不多有三年.
保羅很少在一個地方留那麼久.
但是保羅在以弗所.
我們相信他留過差不多有三年的時候.
這樣有兩年之久.
許多已經信的人.
來承認並公開自己所行的事.
又有許多評述行邪術的人.
把他們的書都拿來焚燒.
他們計算書的價值.
共值五萬塊銀錢.
這樣主的道大大的興旺.
而且普遍傳開了.
這裡說許多已經信主的人.
他們公開承認自己所行的事.
是甚麼事呢.
我們不知道那個指示.
但我們相信都是和拜偶像.
拜假神的事有關.
他們曾幾何時自以為而.
曾幾何時他們貪愛世界.
但現在福音來到他們的人生.
在顛覆他們的人生.
於是第十九節.
他說有許多人.
評述行邪術的人.

$^{521}$把他們的書都拿來了.
堆積在眾人面前焚燒.
然後他們計算這些書.
在當時的市價值多少錢.
答案是值五萬塊銀錢.
五萬塊銀錢值多少錢呢.
其實就等於五萬日薪水.
五萬日薪水你開始計算一下.
有美國北美的神學家估計.
等於大概六百萬美金.
六百萬美金等於港幣四千七百萬左右.
倒水.
燒了.
你是不是覺得很可惜.
為何燒了它.
可否把它賣了.
然後捐給教會都可以.
但賣不掉.
因為那些是邪術的書.
邪術的書都是歷史文化遺產.
可否當歷史文物去賣.
現在應該很值錢.
不可以.
所以你說這班人很捨得.
這樣六百萬美金.
四千七百港幣.
這樣倒水.
不可思議的捨棄.
亦都是因為這樣.
所以你才有第二十字.
如果他們不是.
這樣去捨棄.
這麼願意去捨棄.
你怎會有著的到大大的興旺.
而普遍傳開了呢.
四千七百萬港幣.
不可思議的捨棄.
令我想起三個現實的場景.
第一個場景.
2003年的路德傳.

$^{561}$不知道大家有沒有看過.
中間那個馬丁路德.
後面有幾個是配角.
我特別喜歡這個配角.
這個右手邊的配角.
這個配角在歷史上.
我們叫他做Frederick the third.
他叫做菲德烈的三世.
他就是當年保護路德的人.
他以前是羅馬天主教.
現在他因為路德的教導.
開始扭tie .
他成為宗教改革.
一個很重要的政治人物.
如果不是他保住路德的命.
馬丁路德就死了.
根本沒有後面的事.
我記得在電影裡的其中一幕.
他有一個僕人來問他.
你現在真的拜路德的教導嗎.
如果是這樣的話.
我們聖人的東西就假了.
我們只會唯獨耶穌.
然後他問他.
你存了很多聖人的遺物.
你知道嗎.
當時很流行的.
這是彼得睡覺的時候.
那對襪子.
你明白嗎.
有很多或是誰剪指甲剩下來的.
有很多聖人的遺物.
而這個人.
存了很多聖人的遺物.
你知道那些東西多值錢.
那些遺物怎麼辦.
就扔掉.
扔掉.
600萬美金.
拋湯.

$^{601}$不賣了.
因為那些不是.
合乎聖經的東西.
因為那些東西.
其實是異端邪說.
那些東西是妖言惑眾.
一定要走.
第二個場景讓我想起.
如果你有機會去.
泰國.
你見到如果有一群泰國的.
弟兄姊妹.
從前拜佛的.
買賣落落.
那裡有很多大佛小佛.
金佛銀佛.
銅佛什麼都有.
現在你信了主你怎麼辦.
很值錢的.
還要某某高僧去開光.
那你怎麼辦.
難道賣了它不行.
600萬美金.
一下一鋪.
完全拋湯.
第二個場景.
第三個場景讓我想起香港的場景.
林牧師.
我沒有拜偶像.
我從來沒有用錢買過偶像.
所以沒事沒問題.
然而我們必須要想.
考量這件事.
我們真的要跟隨耶穌.
我們要離開埃及.
出埃及記.
我們要脫離巴別的隱憂.
其實是否很有可能.
我們都要付出這樣的代價.
我們最近開始了.

$^{641}$MASS 2.0.
MASS 2.0 有個口號.
Be Mad, Go Waste Your Life.
做個瘋子.
浪費一生.
你可以回家仔細去計算.
如果我們當中.
有一位醫生.
有一位律師或專業人士.
有幸.
有幸夢照成為傳道人.
成為宣教士.
不是提早退休.
做幾年那一種.
是將他最黃金的歲月.
獻給上帝那一種.
如果是有這樣的事.
你回去計算一下.
我告訴你.
600萬美金.
有多貴呢.
閒閒地.
閒閒地就是這個價錢.
不過.
你一定要想.
你地上少了600萬.
有沒有想過天上.
會多了多少個零.
有沒有想過.
上面的咪錶和下面的咪錶.
其實是撐天拱地.
上面是笑笑聲.
你下面頂多幾百萬.
上面笑笑聲.
的士一坐進去.
就幾笑.
你想一下你地上所失去的.
你天上.
會得到多少.
不可能計算的數字.

$^{681}$更重要的.
地上的600萬美金.
可以買你今生的自由.
可以讓你成為神奇女俠.
加入復活聯盟.
從此再沒有任何的束縛.
沒有任何的束縛.
從此再沒有任何的束縛.
沒有上限沒有底線自由自在.
單單為上帝而活.
是不是很抵呀.
600萬美金可以買自由.
你覺得是不是抵呀.
人家問自由何價.
答案是自由無價.
但我告訴你.
今日自由特價.
只要600萬美金.
只要4700萬港幣.
你就可以得到自由.
而且你不用馬上付錢.
只要你放棄追求.
4700萬你就有自由.
是不是很抵呀.
What a deal.
今日特價.
你要不要.
換取.
這個自由.
這個是保羅帶領上帝的子民.
逃出埃及.
逃出巴別.
逃出偶像之城的救贖士工.
Incredible.
保羅的士工很深.
很深.
但是保羅的士工同時之間很廣.
Deep and wide.
很多人都知道.
在少行傳裡面.

$^{721}$記載了保羅的三個宣教旅程.
我先和你.
看看第一個旅程.
第一個宣教旅程是少行傳第十三章.
由安提柯教會出發.
大概是公元46至47年左右.
保羅的初試提升.
你看到那個地方.
你都知道不是去得很遠.
雖然對於他當時來說.
已經去得很遠了.
接下來是第二個宣教旅程.
你有沒有發覺.
保羅現在貪心了很多.
第二次再去的時候.
去得遠很多.
少行傳第十五章.
又是從安提柯教會出發.
應該是公元48至51年左右.
這次越走越遠.
他的野心.
上次北米.
這次可以去遠一些.
但我想你特別留意.
一個就是上方的斷線.
有沒有留意到.
斷了線.
保羅當時.
如果你知道少行傳.
是想向東方行走.
我們經常說.
如果當時保羅來了東方.
得到福音的我們中國人.
但是因為.
聖靈不讓他來東方.
叫他回到西方.
結果他有一個馬其頓的異象.
於是就去了.
肥臘比.
我們就失迷.

$^{761}$他就得迷了.
我們回來再看.
他在肥臘比裡面.
上帝為他預備了什麼.
是上帝特意叫他去.
馬其頓去肥臘比.
之後保羅有第三個.
宣教旅程.
由少行傳第十八章.
開始.
又是安提拉教會出發.
時間是公元52至57年.
這次去得長一些.
因為他包括了.
他三年在以弗所的工作.
從地理形勢.
上面都是差不多.
因為他要去那些教會.
去堅強他們.
我想你特別留意.
這次回不到家.
他回到耶路撒冷的時候.
他被抓了.
他沒有辦法回到安提拉教會.
他在安提拉教會出發.
但他最後在耶路撒冷.
那裡被捕.
後來轉押至凱撒利亞.
然後被困兩年.
然後終於兩年之後.
可以前往羅馬.
上訴凱撒利亞.
就是今天我們讀的經文.
前往羅馬的旅程.
其實是保羅的.
第四個宣教旅程.
你有沒有看到這次.
傳得多遠.
那時候第三個宣教旅程.
都是這裡.

$^{801}$嘩.
要去地極.
我們一直說他有三個旅程.
然後被人抓了.
搞了很久.
差點死了.
然後才去到羅馬.
不是的.
有沒有想過.
一定要被傳揚地極.
你看這次.
很遠.
但除了距離以外.
有沒有留意.
就是頭三個宣教旅程.
這裡全部是頭三個.
這裡是第四個.
頭三個宣教旅程.
大多是陸地上進行.
但第四個旅程.
第四個宣教旅程.
大多都是海上行的.
而且是荊濤海浪中行的.
就是我們剛才讀的.
為什麼讀了這麼多.
有沒有想過.
為什麼你去到四大行轉.
最後的.
第27集.
為什麼整集都是海上.
究竟路加要我們知道些什麼.
我想你知道一件事.
路加不單歷史.
他還在寫神學.
時間關係.
我要馬上幫你開估.
我想用五點.
來總結.
四大行轉下半部.
史奴保羅在四大行轉裡面.

$^{841}$所領受的福音使命.
我們怎樣明白.
保羅的福音使命.
五點.
第一.
史奴保羅會成為陸地的征服者.
一至三.
他第一第二第三個旅程.
他成為海洋的征服者.
在海洋的征服的同時.
他成為一個更美的若拿.
在海洋的經歷裡面.
他是一個救贖的持有者.
解釋多些.
最後.
當他死裡復活之後.
若拿也是經過死裡復活.
他成為世界的征服者.
這五點.
是史奴保羅的使命.
但這五點同時也是.
復活救主耶穌.
升天.
聖靈降臨.
在他的使徒身上.
所引發出來的五點.
所以最後保羅所做的所有事.
其實都是耶穌所做.
所以我們可以說.
復活的救主耶穌基督.
是陸地的征服者.
是海洋的征服者.
是更美的若拿.
是救贖的持有者.
是世界的征服者.
剩下的.
講到我想和你快快去看.
這五點.
我們先看.
陸地的征服者.

$^{881}$但我們回到剛才.
馬其頓的故事.
大家都很熟悉.
或者你記得有這樣的故事.
《最後行傳》第十六章.
第九節.
夜間有異象.
向保羅顯現.
有個馬其頓人站在.
求他說請你過來.
馬其頓幫助我們.
看見這異象.
我們就立即.
設法往馬其頓去.
認為上帝呼召我們.
傳福音.
給那裡的人.
在我們還沒有討論.
馬其頓人之前.
我想你留意一樣.
馬其頓人所需要的幫助.
馬其頓人說.
到馬其頓來.
幫助我們.
幫助我們什麼呢?經文是沒有直接說明.
但從保羅的回應.
保羅的回應是什麼呢?.
他認為上帝呼召我們傳福音.
給馬其頓.
馬其頓需要的是.
福音.
第一件事你要知道.
福音不是一個基督徒.
創造出來的詞句.
原來在保羅傳教的時代.
福音.
已經是一個.
普遍運用的詞句.
Evangelion這個希臘的詞組.
當時已經存在.

$^{921}$不是基督徒創造出來我們有福音要傳.
不是的,是當時已經有.
一件事叫做福音.
什麼叫做福音?聽著.
福音字面上叫做好消息.
但是福音.
不是一個普通的好消息.
在古希臘.
羅馬的文化世界裡面.
福音代表了一個.
君王的好消息.
是關於一個君王誕生.
一個君王登基.
君王戰勝敵人.
海船而歸的公告.
所以當你說.
馬其頓需要福音.
即是以下的意思.
馬其頓需要一個君王.
馬其頓需要一個.
救主.
如果你知道.
當時的歷史背景.
你會想馬其頓都需要.
一個君王.
這是一個何等大的諷刺.
為什麼?.
因為馬其頓已經有.
他.
馬其頓已經有一個君王.
而且不是.
隨便普通的一個君王.
是當時歷史裡面.
其中一個最偉大的君王.
他的名字叫Alexander the Great.
他是什麼人?.
他是馬其頓人.
所以馬其頓需要君王.
馬其頓已經有.
亞歷山大大帝.

$^{961}$他是古希臘.
馬其頓國的君王.
你看看時間.
你已經知道他死了三百多年.
然而.
他是當時公認的.
世界的征服者.
在他三十歲的時候.
已經建立了一個史無前例.
人類歷史裡面.
其中一個最龐大的帝國.
有一個傳說說.
他有一次打完仗的時候.
他開始哭.
他的將軍問他.
為什麼你哭?.
他說因為我再沒有土地可以去征服.
他是一個名副.
其實我們叫西方不敗.
因為他從東邊打到西邊.
從來沒有輸過.
三十歲的時候.
已經站在世界的高峰.
然後三十三歲就死了.
想都想不到.
原來馬其頓已經有一個君王.
而且.
亞歷山大不單止是世界的征服者.
他最有名的.
他是福音的傳遞者.
他每去到一個地方.
他的侍從就會說.
Good News! Good News!.
有好消息 有福音.
馬其頓君王.
亞歷山大大帝.
帶給眾人的福音.
是古希臘的福音.
是古希臘的智慧.
你可能有聽過.

$^{1001}$我們今天一想起.
古希臘的智慧.
我們會想到這三個代表的人物.
蘇格拉底.
然後他的學生叫.
柏拉圖.
然後柏拉圖的學生叫.
亞里士多德.
你知道亞里士多德的學生是誰嗎.
亞歷山大大帝.
他是第四個.
他是最崇拜的.
就是他前人.
這三個人.
他去到哪裡.
把希臘的福音.
希臘的智慧.
傳到那裡.
開始你想起.
在哥林多前書第一章裡面.
保羅去到哥林多的時候.
保羅說.
猶太人要的是神蹟.
希臘人要的是智慧.
就是這個.
但我們傳一個不一樣的福音.
猶太人有猶太人的福音.
希臘人有希臘人的福音.
我們傳一個不一樣的福音.
希臘人有希臘人的福音.
我們所傳的是十字架的道理.
亞歷山大大帝三百多年前已經死了.
希臘也不在.
現在是羅馬統治.
但你有沒有留意.
這個世界是沒有羅馬文.
這個世界只有希臘文.
去到羅馬帝國.
新約聖經寫的也是用希臘文.
因為儘管希臘不在.

$^{1041}$羅馬統治.
而希臘的語言.
智慧文化.
福音繼續在這個世界宣傳.
所以保羅要面對的是.
希臘人要的是智慧.
是他們的福音.
所以這一切.
都是馬其頓.
異象的歷史背景.
我想你想想路加.
到底要跟我們說什麼.
路加在說.
他在畫一個屬靈的圖畫.
這個圖畫非常諷刺.
這個圖畫也非常真實.
圖畫諷刺真實的地方.
就是馬其頓人需要福音.
意味著什麼.
君王的家鄉.
其實是沒有福音的.
如果這個世界有一個地方.
是不需要福音.
不需要君王.
應該就是馬其頓.
因為他已經有這個世界裡.
其中一個最偉大的君王.
然而馬其頓是假的.
原來馬其頓完全在一個黑暗的當中.
有一個救主君王去拯救他們.
這個是他路加給我們的圖畫.
西方如是.
東方也是一樣.
如果當日保羅沒有去西面.
去了東邊.
相信在異象裡面.
呼召保羅的.
很可能是孔子家鄉裡一個儒家的高人.
你明白嗎.
孔子的家鄉需要福音.

$^{1081}$為什麼.
孔子不可以.
西方有西方的黑暗.
東方有東方的黑暗.
路加想告訴我們.
當在天國福音的光照下.
最優良的傳統.
最昌盛的國度.
全都是自以為聰明的智慧.
自以為光明的黑暗.
現在呢.
自以為光明的黑暗.
自以為聰明的智慧.
世人都犯了罪.
虐了上帝的榮耀.
罪的公敬乃是死.
唯有上帝的恩賜在耶穌基督裡面.
乃是永生.
唯有耶穌基督才有智慧.
才有真正的救贖.
於是.
當福音傳到任何一個地方.
福音會挑戰當地的福音.
福音.
十字架的福音會挑戰.
也要得勝當地的福音.
但我們繼續去看多一個.
在肥腩的地方.
發生的故事.
這是你必須要明白的故事.
約在半夜.
你記得保羅坐在監獄裡.
和西拉.
保羅和西拉.
正在禱告.
暢思讚美上帝.
忽然地帶震動.
好像彼得出獄那一刻.
獄警一醒.
以為囚犯已經逃走.

$^{1121}$就拔刀要自殺.
保羅大聲呼叫.
不要傷害我.
我們都在這裡.
獄警令他們出來說.
兩位先生我必須要.
做什麼才可以得救.
大家都知道.
大家記得這個.
肥腩的獄警.
有時候我們叫他肉卒.
我們聽這段聖經很容易誤解.
因為我們一想.
這個獄警.
我們說他是卒仔.
他是被人.
托一托嚇一嚇就信了耶穌.
擺個保險.
這樣的卒仔又要自殺.
怕這個怕那個.
於是很快就信了耶穌.
這是完全錯誤的理解.
再讓我們知道當時的羅馬.
監獄的.
獄警或者監獄官.
其實是.
之前新京百戰.
終生至死的退伍勇士.
不是隨便一個人.
會做了監獄長.
這是新京百戰.
終生至死.
退休了的.
勇士.
不是普通的士兵.
他是一個位高權大的將領.
如果他不位高權大.
你怎麼可以解釋到後來.
隨便把保羅西拉帶出來.
回家吃飯.

$^{1161}$還要幫他們療傷.
整晚去哪裡了.
你老闆不找你.
為什麼他敢做這些事.
不是因為他本身.
權高位重.
他想拔刀自殺.
不是因為他.
貪小.
懼怕.
而是因為他忠心耿耿.
他覺得他一生人.
就是要服侍羅馬的凱撒.
於是當他見到.
逃犯走了的時候.
他污辱了他的主人.
於是他沒辦法接受.
他自己的失敗.
他願意自殺.
因為他一生裡.
唯一的主就是凱撒.
然而保羅.
看著他.
當保羅看著他的時候.
他開始驚訝.
他驚訝保羅所傳的福音.
因為在他們唱歌裡.
充滿著他的福音.
他驚訝保羅在監獄裡的喜樂.
他驚訝保羅沒有逃走.
於是他.
他問保羅.
我必須.
做什麼才可以得救.
這個世界沒有多少人問.
這樣的問題.
上次聖靈降臨的時候.
耶路撒冷的人聽完彼得那一場.
刀切入心中.
然後問我該.

$^{1201}$做什麼可以.
得救.
而保羅的回答.
我沒有打出來但大家都知道.
是整個故事的高潮.
當信主耶穌.
你和你的一家必得救.
我想很多人都有.
想去哪裡.
買東西.
放在床頭.
當信主耶穌你和一家必得救.
保羅正在說.
對不起大家.
我想你們會錯意.
不是一個廉價福音.
什麼叫當信主耶穌.
有幾個字的.
當信耶穌.
主.
主應該很巨大的.
主耶穌.
你當知道唯有耶穌是主.
這個人未信.
耶穌之前.
他腦裡面一生的座右銘.
就是凱撒是主.
現在保羅說你怎可以得救.
保羅說很簡單.
凱撒再不是你人生的主.
誰取締了凱撒的地位.
凱撒是耶穌.
在眾多人裡面.
這個肉警是最能明白.
那句說話.
你給第二個.
什麼叫耶穌是主.
主什麼.
主辦.
沒所謂你叫我信耶穌.

$^{1241}$簽個紙馬個保險.
沒問題.
我們今天你說耶穌是主.
沒人知道你在說什麼.
但是這個肉警完全知道.
明白這件事.
他的人生意義.
十字架的基督.
將要成為他生命中的至寶.
十字架的基督將要成為.
他生命裡面的唯一寶貝.
Christ is King.
耶穌是.
君王是救主.
Not Caesar.
再不是凱撒.
我現在重新來過.
我的人生.
要跟隨耶穌.
所以我們會說耶穌征服了.
這個.
菲勒比.
的監獄長.
耶穌征服了他們.
征服了他.
和他的一家.
但我們繼續去想.
耶穌不單止是征服.
陸地上的人.
耶穌不單止是一個陸地的征服者.
耶穌也是一個海洋的征服者.
這個很重要的.
因為對於舊約去到新約的時候.
唯有上帝可以.
征服海洋的.
而如果耶穌的福音.
要傳到地極.
必須要征服海洋.
因為這個世界大部份的地方.
都是海洋.

$^{1281}$當時你又沒有坐飛機.
你要將福音傳到地極.
你必須要征服.
海洋.
所以去到這裡可能我們當中.
有人說我想起林錫祥.
我知你說甚麼.
我們剛剛在一樣福音裡面學過.
耶穌有一晚.
在大雨之後.
派遣十二個門徒.
到加利利海那裡.
然後狂風大浪.
黑夜狂風.
耶穌在海上行走.
然後營救門徒.
是不是.
好像是.
但是路加說不算數.
路加說不算數.
我怎知道路加說不算數.
我想帶你們去看路加福音的版本.
路加福音第五章第一節.
耶穌站在.
格尼撒勒湖邊.
格尼撒勒其實就是.
加利利海.
但你有沒有留意他怎樣寫.
耶穌站在格尼撒勒湖邊.
他就是不想寫海.
眾人擁擠他要聽上帝的道.
他見有兩隻船靠在海邊.
如果你讀馬太.
你馬河是海邊的.
但他又要強調不是海邊.
是湖邊.
當然我明白理由.
因為格尼撒勒湖不是海.
難道維多利亞海.
不是很對的.

$^{1321}$但當時他還是要講海的.
但他偏偏要告訴你是湖.
打雨的人就離開.
然後第八章第二十二節.
耶穌和門徒上了船.
這個就是平靜風浪的故事.
說我們到湖的對面岸去.
他又是跟你爭辯.
不是海是湖.
船行的時候.
湖上忽然起了狂風.
到底是誰.
連風和.
你看其他的馬太馬河.
是會說風與海的.
但他故意不想講海字.
他說風與水.
路加說不算數.
不算數.
不是海邊是湖邊.
不是海上狂風.
是湖上狂風.
不是風與水.
不是風與海.
聽從他.
是風與水聽從他.
什麼都不算.
不算.
不夠勁.
如果.
福音要傳到地極.
如果萬國要聽到這個福音.
要征服世界.
我們要的是.
海上狂風.
不是湖上狂風.
一定要有個勁一點的.
可以這樣寫.
因為他知道最後出場的.
是海上狂風.

$^{1361}$他早就預了.
第27章紹行傳的時候.
海上.
前面跟你說不算數.
你看看後面.
第27章第22節.
現在我勸你們放心.
除了損失.
這條船.
這條船就沒有了.
你們中間沒有一個人會喪失性命.
因為昨夜.
我所.
所事奉的上帝.
的使者.
站在我身旁說保羅不要害怕.
你必定在.
凱撒面前.
並且上帝已把.
安全賜給你和你同船的.
所有人有多少人.
船上共有276個人.
我不明白.
如果你讀這個故事.
你說.
上帝是否能夠應付.
如果能夠應付.
為何會損失船上的所有人.
如果他夠強.
船上和人上都會沒事.
但現在不是.
命是保住了.
我當你算了.
但是.
保不住船上.
是否有點勉強.
上帝是否有點失控.
有否危危乎乎的感覺.
復活的救主.
是否真的能夠.

$^{1401}$掌管大海的狂風.
復活的救主.
是否有足夠的掌控.
以至他的門徒.
可以地的見證萬國宣揚.
征服世界.
這就是問題的所在.
我想.
整段征服海洋的經文.
我稱之為征服海洋的經文.
關鍵在哪裡.
就是第37節.
37節.
276人.
這不是一個.
大概的數目.
如果大概數目是300人.
是一個很精準的人頭.
一個都沒有少.
一個都沒有少276人.
與此同時.
那艘船.
就徹底破壞.
不是一點點破壞.
是徹底破壞.
船沉大海.
飛機墜落.
276位乘客.
無一倖免.
通常我們是這樣讀.
但現在飛機墜落.
竟然276名乘客.
沒有一個傷亡.
這是不是神蹟.
這是不是更厲害的神蹟.
你可以想像一下.
上帝需要有多高超的掌控.
才可以創造這件事.
保住整艘船.
人沒事一件事.

$^{1441}$你要.
船沉.
海底.
276人一個都沒事.
行不行.
需要的掌控.
有多厲害.
你想起很多年前.
我曾經和幾個.
中學的同學討論.
我們討論的題目是.
會拿9A 10A容易.
還是拿9B 10B容易.
有沒有想過這個問題.
其實當然.
拿9B 10B.
更加困難.
拿9A 10A其實很容易.
渾身解數全力以赴就對了.
懂得做就做.
但你要拿9B.
我有個朋友.
拿2A 7B.
公規一據.
但你要拿9B.
拿10B.
要很精準的計算.
要完全掌控.
我哪裡放水.
要放多少水.
以至我僅僅拿不到A.
有沒有想過.
船沉大海.
毫無損傷.
是一個不可思議的神蹟.
船沉大海.
下一句不是無一生還.
不是無一倖免.
而是毫無損傷.
這代表什麼.

$^{1481}$上帝完全掌控住.
連一根頭髮.
也不至於掉落.
完全掌控.
完全征服.
於是耶穌復活的救主.
是陸地的征服者.
海洋的征服者.
第三點是什麼.
是更美的.
若拿.
精英學者告訴我們.
航海前往羅馬的故事.
帶著很多很多.
若拿的影子.
而若拿的影子裡面.
有很相似的地方.
也有與若拿不同的地方.
我只是列出.
六個大點.
給你看.
保羅與若拿.
兩個都是帶著使命.
一個不願前往.
另一個熱切渴望.
兩個都去外國人的首都.
一個宣告審判.
他心裡面是想審判.
另一個要宣揚救恩.
兩個都經歷狂風大浪.
一個是連累他人.
另一個是拯救眾生.
兩個故事結局.
都是無人喪命.
第一是因為若拿被拋下海.
但是保羅的原因.
是因為上帝應許他.
兩個都經歷一個.
受死而復活的人生.
若拿之所以.

$^{1521}$受死後可以復活.
如果你回頭聽我以前講的.
約拿書的時候.
是因為他沉落水底的時候.
在耶路撒冷的聖殿.
做了一個祈禱.
是因為他回轉得生.
而保羅之所以會受死復活.
是因為保羅是一個受苦的僕人.
正如耶穌是一個受苦的僕人.
兩個結果都在愛邦那裡.
去見證上帝.
一個在城外.
講完之後在城外發怒.
另一個在城裡不斷繼續的去教導.
你看到嗎.
保羅是一個更美的若拿.
The true and the better Jonah.
為什麼.
因為耶穌自己就是一個.
True and better Jonah.
耶穌就是一個更美的若拿.
路加福音第十一章.
耶穌已經講了他自己的身份.
當眾人越來越擁擠的時候.
耶穌說.
這世代是一個邪惡的世代.
他們求神蹟.
除了若拿的神蹟以外.
再沒有別的神蹟.
及其神蹟.
再沒有別的神蹟.
這個耶穌無盡的恩典.
如果你問我除了索多瑪.
和俄羅斯的神蹟.
沒有神蹟給你看.
你要就去地獄一起死.
你一群人死.
耶穌給他的不是索多瑪和俄羅斯.
耶穌給他的是若拿的神蹟.

$^{1561}$若拿怎樣為尼尼美人成了神蹟.
人子也.
為這世代的人成了神蹟.
My God.
他會受死復活.
向這個世界的尼尼美人.
傳放天國的福音.
看押比若拿更偉大的在這裡.
於是若拿受死復活.
耶穌受死復活.
現在我們看到保羅受死復活.
下海出海.
全部都是為著要使救贖和生命.
帶給那些需要救贖需要生命的眾人.
在整個航海的旅程裡面.
無論周遭環境有多危險.
保羅是一個救贖的持有者.
保羅是一個生命的持有者.
保羅死不去.
而和保羅一起的人.
每個人都會生.
不和保羅一起每個人都會死.
在聖經第27章第20節.
他說很多天都沒有看到太陽升起.
狂風大浪獲救的指望都放棄了.
沒有了獲救救贖.
沒有了沒有希望.
然而第31節.
保羅對伯夫長和士兵說.
這些人若不留在船上.
你們就不能獲救.
如果他們留在船上.
他們必然得救.
保羅的意思是.
保羅在船上.
如果他們和保羅一起.
他們一定得救.
第34節.
所以我勸你們吃點東西.
這是關乎你們獲救的事.

$^{1601}$聽我的話.
聽我的話你們就會獲救.
因為你們各人連一根頭髮也不至於掉落.
保羅是救贖的持有者.
保羅是生命的持有者.
你與保羅同在.
你就可以得救.
你與保羅同在你就可以得救.
這是耶穌的影子.
有耶穌就有永生.
沒有耶穌什麼都沒有.
只有地獄.
有耶穌什麼都有.
第27章第43節.
伯夫長要救保羅.
當然上帝會救保羅.
不准他們任意而行.
會游水的跳下水去.
先上岸.
用板子和船的碎片上岸.
這樣眾人.
276個人.
全部都獲救.
第28章第1節.
我們既已獲救.
才知道那島叫做馬爾他.
有保羅.
聽保羅.
有得救.
不聽保羅.
沒有保羅.
276個.
只要聽保羅的話.
與保羅一起.
都會有得救.
令我們想起約瑪書那句話.
salvation is of the lord.
救恩出於耶和華.
救恩出於耶穌基督.
救恩出於使徒保羅.

$^{1641}$我不是說保羅可以救我們.
但因為耶穌與保羅同在.
所以與保羅一起就等於與耶穌一起.
聽保羅的話就等於.
聽耶穌基督的話.
你不聽保羅的話.
與保羅一起.
就沒有得救.
這也是一個.
現代我們今天教會所領受的.
一個很偉大的呼召.
如果我們傳揚的福音是真實的.
是準誠的.
是忠心的.
你聽這個福音你就有得救.
你不聽這個福音你就沒有得救.
這是我們的使命.
不是我們要認叻還是什麼.
而是實際上.
是使徒行傳第29章的訊息.
告訴我們.
今天的福音在我們手中.
耶穌與我們同在.
我們要忠心傳揚這個訊息.
聽這個訊息的人.
就有得心.
不聽的話就要受到懲罰.
其實這個故事還有一幕.
我們沒有時間去看.
但你一定會記得.
他們去到馬爾他小島上.
之後.
保羅拾起一個.
一棍柴.
放在火中.
有一條毒蛇出來.
因為熱的緣故轉上來.
纏繞著他的手.
然而保羅一甩.
將蛇放在.

$^{1681}$火中.
而他沒有受傷.
這是什麼來的.
這又是告訴我們.
保羅是一個生命的持有者.
他是一個救贖的持有者.
有一點點.
你傷不了我.
而且你會承受上帝的審判.
有一點點.
在《肖行傳》第十二章裡面.
那個希律王最後.
彼得復活.
然後希律.
窒息而死.
因為蛇代表了.
撒旦.
代表了撒旦的權勢.
保羅是一個復活的.
救主的代表.
保羅的力量.
生命的力量遠遠超過了.
撒旦的權勢.
最後.
保羅這個.
救贖持有者.
這個生命持有者終於來到羅馬.
在羅馬他成為了.
世界征服者.
羅馬書第一章.
第十三節.
我屢次計劃.
往你們那裡去.
到如今.
仍有阻難.
願意盡我的力量.
把福音也傳給你們在羅馬的人.
保羅寫這本書的時候.
羅馬書.
他的西遊記.

$^{1721}$大概是公元57年.
他要等三年.
公元60年.
他終於到了羅馬.
保羅在自己所租的房子裡面.
他當時受軟禁.
又有冰釘.
跟他鎖鏈鎖在一起.
但是他有自己的地方.
租的房子裡面足足住了兩年.
放膽傳講上帝的話.
並教導主耶穌基督的事.
沒有任何的禁止.
公元57年.
他寫羅馬書.
公元60年.
他去到羅馬寓院.
以上.
地的見證萬國宣揚.
征服世界.
雖然手上帶著鎖鏈.
軟禁在家裡.
他放膽傳講上帝的國.
並教導主耶穌的事情.
沒有人可以禁止到他.
保羅的.
復活救主.
耶穌基督.
今天仍然活著.
耶穌要帶領我們.
今天.
明天.
教會的福音使命.
我們要寫的是.
第29章.
時間已經夠了.
結束的時候.
我想用一首歌來.
總結應用.
這篇講道.

$^{1761}$和我們三篇復活節的.
迷你講道系列.
我選的歌是Jesus.
Shall Reign.
主字萬方.
這裡寫出來的是.
頭三節的歌詞.
我只是想跟你讀第二節.
我很喜歡第二節.
Blessing abound where he reigns.
The prisoner leaps to lose his chains.
The weary find eternal rest.
And all the sons of want are blessed.
在主旨下.
多蒙福樂.
罪人跳出.
捆綁束縛.
倦者能享永恆安息.
風刑負足.
恃與貧者.
上帝來.
要施行出埃及記.
我們要從這個世界的.
捆綁裡面逃脫出來.
我們要單單的.
服侍上帝.
Jesus shall reign.
主字萬方.
但這個主字萬方必須要從我們開始.
我們第一個要問自己的問題.
主是不是在治理我們的人生.
如果上帝不是在治理你.
你怎麼可以參加一個.
主字萬方的使命呢.
所以我們求主幫我們.
與他同死.
同復活.
讓我們.
棄絕這個世界的偶像.
讓我們放下那六百萬美金.

$^{1801}$讓我們.
好像監獄長說.
耶穌是主.
凱撒不是.
世界不是.
唯獨基督.
讓我們自由的去奉獻給上帝.
讓我們歡喜獻上.
我們的今生我們的最好.
讓我們仿效保羅.
渴慕參與這個使命.
以致我們可以.
與眾聖徒一起.
地極見證.
萬國宣揚.
征服世界.
讓我們.
在上帝的面前.
在為我們受死.
復活的耶穌的面前.
讓我們用這首歌.
從心裡.
去讚美他.
這是一個英文版.
一個比較modern的版本.
我希望你一邊聽的時候.
一邊心裡去思想.
也請留意.
我幫你做了一個很特別的video.
我也可以坐下了.
(歌的反應越來越厲害).
多謝您 我們敢祈禱 奉主耶穌基督的名求 阿們.
\newpage



\section{出埃及記 16:8}
\label{sec:032017}
\textbf{聖經人聖經事:你們的怨言是向耶和華發的}~~~~黃克勤牧師博士
\newline
\newline
連結: \href{http://www.church.com.hk/acms/content.asp?site=cdc\&op=show\&type=product\&code=032017\&layout=sermon}{\texttt{032017}} ~~~~ 語音日期: 2026-04-17
\newline
\newline
\hyperref[sec:032016]{\small{< < < PREV SERMON < < <}}
~
\hyperref[sec:index_chronic]{\small{[返順時目]}}
~
\hyperref[sec:032018]{\small{> > > NEXT SERMON > > >}}
\newline
\newline
出埃及記 16:8
\newline
\begin{longtable}{cl}
\hline
\hline
章節 & 經文 (和合本修訂版)\\
\hline
16:8 & \begin{tabularx}{0.7\textwidth}{X} 摩西又說：「耶和華晚上必給你們肉吃，早晨必給你們食物得飽，因為耶和華已經聽見你們向他所發的怨言。我們算甚麼呢？你們的怨言不是向我們發的，而是向耶和華發的。」 \end{tabularx} \\ \\
[1ex]
\hline
\hline
\end{longtable}
$^{1}$各位電影節目大家好.
生命樹又和大家一起看聖經了.
今天聖經人物和聖經故事.
仍然集中在摩西和以色列人.
我們看這些故事.
其實是想看看背後上帝的手和想法.
這是最重要的.
認識摩西多一點也好.
認識以色列人多一點也好.
但最重要是認識上帝.
今次的題目剛剛這麼長.
就是你們的怨言是向耶和華發的.
這句話是誰說的呢?.
是摩西說的.
你們的怨言向耶和華發的.
似乎發怨言我們每個人都會做的.
好,看看段經文.
以色列人出埃及.
我將地圖大概看得多.
你也有些概念了.
戈山這個地方.
尼羅河向北流.
流入地中海.
去到平地就變成一個尼羅河三角洲.
非常肥沃.
而且每一個雨季都會泛濫.
大家很喜歡泛濫.
泛濫大家都知道.
泛濫就搬走了.
他們就會浸了土地.
然後水一退就很肥沃.
埃及就靠這個戈山地.
以色列人是住在裡面的.
因為很早年約瑟在這裡做宰相的時候.
接了以色列一家來.
當時的法老.
你喜歡選一個地方來住.
你喜歡住哪裡就住哪裡.
沒所謂了.
他就選了一個地方.

$^{41}$話說回來.
埃及人是喜歡種田的.
是不喜歡畜牧的.
看低畜牧的人.
所以他要選一個地方.
是不會被埃及人不喜歡的.
於是在戈山選了一個地方.
後來到不認識約瑟的新王起來.
就開始說.
這班人不是我們埃及子民.
不過二百萬.
個個又身強力壯.
如果有外敵.
跟他們通在一起.
外憂內患.
一齊我們埃及就沒有了.
於是就決定.
將他們逐漸強硬地變成奴隸.
沒有身份.
還有要殺那些男孩子.
誰一出生就殺死他.
這些故事大家都知道的.
建了蘭塞和比東這兩個織火城.
在這個城裡面.
到了某一天摩西來帶他們走.
十災之後.
殺長子之災如月之後.
他們聚集到梭國.
當然有其他人住在那裡.
一齊來梭國.
是他們第一個集合的地方.
梭國出去後.
下了二唐.
二唐之後就調回北上去到麥奪.
於是埃及人就說.
哈哈哈哈.
他們在曠野迷路了.
我們有機會追他們回來.
於是就派兵在這些高山地追他們.
誰知他們背水一戰.

$^{81}$就過了海.
去了巴黎洗婚.
終於埃及人死在紅海裡.
上一集就說到紅海.
其實照現在看回頭.
很有可能不是我們所說的紅海.
而是苦海之上有些地方.
他們走上去.
是地中海進來的一些.
比較闊的海面.
他們就過了那裡.
其實現在的紅海.
這個就是蘇伊士灣.
這裡一直有河流.
有湖.
照現在這個苦湖已經封了.
只剩下水.
蘇伊士運河就在這裡.
其實水就是地中海的水.
所以很有可能他們過的.
叫做一個蘆葦海.
叫蘆葦海.
Yem Suv 蘆葦海.
希伯文是.
好了不要緊.
他們就用回這張圖.
然後去到第一站巴拉.
水很苦.
上一次我剛剛提到.
神就指示他有一棵樹.
斬一棵樹放進水裡.
水就會變甜.
我的象徵靈異的解釋就是.
一棵樹被斬了.
死了放進水裡.
而將苦水變成甜.
這棵屬靈異上.
靈異上是可以象徵十字架主耶穌.
犧牲生命.
而令到我們的生命變得甜.

$^{121}$然後去到以林.
以林就有十二個泉水.
全眼.
然後有七十棵棕樹.
可能是象徵靈異.
不是字面的解釋.
不是字意的解釋.
靈異令我聯想有十二個仕途.
他們好像將他們的生命流出來.
像泉水一樣甜.
有七十棕樹.
在十二個門徒之外.
有七十個門徒.
跟主耶穌很接近.
他們都可以提供水.
提供遮陰的地方.
以林.
然後他們再向前走.
一到以林.
就來到這個地方.
信的曠野.
信的曠野再南就是西乃山.
這裡是西乃半島.
西乃山是一個高的山脈.
來到這裡還沒去到西乃山.
曠野就非常乾旱.
於是他們遭遇法院賢.
補充.
這裡就是耶利哥.
迦南.
其實人們如果去耶利哥.
去迦南地.
他們是這樣走的.
很快就到.
已經是別時巴了.
這裡是非利士路.
加薩地帶.
非利士路.
為什麼不走呢?.
因為這條路是埃及的駐軍在這裡.

$^{161}$所以他們不敢走.
怕那裡會打.
所以他們向南走.
就說他繞錯路了.
信的曠野.
法院賢的地方.
出埃及的十六章第一節.
以色列傳會眾從以林起行.
在出埃及第二個月.
十五日就到了以林.
和西乃的中間.
叫做信的曠野.
那個地方名叫信.
以色列第一點.
藍色的字.
傳會眾在曠野.
向摩西亞林發願.
大家留意.
傳會眾.
全部人.
有沒有這麼誇張.
傳會眾200萬人.
傳會眾當然是很容易的.
因為你收到的訊息.
一傳十十.
一傳百.
大家心裡的不開心.
一起出來.
為什麼呢?.
我都有多少痛情.
不過這樣是不對的.
就是我去過很多次以色列.
最好就是你八月去.
因為八月去是全年最熱.
我帶著團友有一次很好笑.
坐在旅遊巴裡面.
這裡就是主耶穌.
鬼入豬群.
趕鬼的地方.
大家下去看看.

$^{201}$不下去了.
看到了.
其實看不到.
還要走很遠.
不下去了.
很熱.
有些老人家去.
他就說.
以色列人在曠野四十年.
當然發怨言.
我們下車四分鐘都發怨言.
是明白的.
很辛苦.
不夠水.
曠野嘛.
他們就在曠野發怨言.
傳媒還向摩西亞倫發怨言.
怨領袖.
怨什麼呢?.
說話以色列人非常刻薄.
說什麼?.
巴不得我們早死在埃及地.
在埃及地耶和華的手下.
在埃及的時候.
耶和華的手下照著我們.
那時候我們坐在浴窩旁邊.
吃得飽足.
誰知你們就將我們領出來.
到遮曠野.
是叫全會眾都餓死.
很刻薄.
很尖銳.
很有趣.
有個分開的.
這個就是怨言.
怨言.
怨言有些特別的地方.
怨言人說的都是.
怨言很多時候是誇張的.
針對性的.

$^{241}$或者有時會說得很普遍性的.
全部人都這麼壞.
其實怎麼會全部人呢?.
全個教會的人都是這樣.
所有牧師都是這樣.
所有老師都是這樣.
所有同事都是這樣.
所有.
全部.
死在埃及耶和華手下.
那時候我們在浴窩旁邊.
這句話真不真?.
當然不真.
誇張.
埃及做奴隸.
怎麼會每天坐在浴窩旁邊任食呢?.
當然不會.
應該不會.
那時候不出來.
就有耶和華照顧.
那時候吳香瑩就發哀聲.
說多好啊.
現在神帶我們當日在埃及.
食肉都吃到鼻子流出來.
多厲害啊.
現在呢.
死在抗爭日期.
餓死我們全會中.
你把我們帶出來想害我們.
這些怨言.
是怨言.
你要懂得處理怨言.
第一.
人一定有怨言.
只要多過一個人就有怨言.
有時自己都要埋怨自己.
那時真是傻了.
如果兩個人一家人就會怨言.
二百萬人怨言多到不得了.
而且去到大家都同樣問題.

$^{281}$肚餓.
不夠肉食.
上次食肉都不知幾時.
而且大家是靠自己帶出來的牛羊.
食下一隻.
我帶少一點.
當時窮.
你帶十隻.
給你一隻.
他又不肯.
於是我要跟他買.
總之生活很困難.
於是說我們在埃及肉窩旁邊.
這是誇張法.
如果怨言對他看得太過認真.
就慘了.
你就很受苦.
到第二點.
神對摩西說.
我明白.
我要將糧食從天降給你.
請你有信心.
神會將糧食從天降給你.
你留意一下.
被人怨到這樣.
罵到這樣.
但摩西亞倫未曾跟上帝說.
上帝是自己說.
耶和華對摩西說.
我要將糧食從天降給你們.
神是有看顧的.
聽到的.
你有信心.
祂會看顧.
主耶穌說.
上帝看顧養活麻雀.
野地的花.
野地的草.
祂沒有做過什麼.
神會給它每天顏色的裝飾.

$^{321}$麻雀如果不是神許可.
是不吊在地上.
頭髮也是.
所以神是體會的.
是體諒照顧的.
我要將糧食降給你們.
百姓可以出去.
每日收每天的墳.
我試驗一下他們.
是否遵守我的法道.
這裡有一點解釋.
其實神當時未給法道.
未到西乃山.
信的抗耶在下才是.
但其實我會告訴他們.
每天有東西吃.
要拿那一份.
不要拿多.
我想看看他們.
這樣這麼簡單的事.
聽不聽我說.
這就是他的法道.
每天收每天的墳.
我試試他們.
是否聽我說.
到第六天.
他們可以將收來的.
是預備好一點.
多一點.
比每天所收的多一倍.
到第七天.
安息日.
不可以再出去拿糧食.
可以在早一天.
拿多一倍.
到現在還未說是麻辣.
我會給糧食他們.
請你記得.
第二點要說.
是神將糧食供應我們.

$^{361}$如果神都沒有供應.
我終於餓死了.
這也是神安排我們的生命.
這樣結束.
我只有接受了.
願然不要發這麼多.
摩西亞倫和會眾說.
到了晚上你就知道.
是耶和華將你們從埃及領出來.
到了早晨.
晚上你會知道早晨說什麼.
其實各位不要稍安毋躁.
今晚各自回家.
記得一件事.
是耶和華神將你們從埃及帶出來.
今晚回去想想.
然後明天早上.
你就會看到耶和華的榮耀.
因為耶和華聽見你們.
向他所發的怨言.
留意一下.
今晚稍安毋躁.
回家再紀念神帶你出埃及.
明天早上你就會看到.
上帝的恩典榮耀.
不要發怨言.
其實神聽見了.
繼續思想.
我們算什麼.
你向我發怨言是沒有用的.
因為我也是神的僕人.
帶你們出來.
摩西亞倫說.
晚上會給你肉吃.
早上會給你吃飽.
你們向耶和華發的怨言.
神也聽見了.
以前在埃及.
發什麼呢.
發向耶和華.

$^{401}$發愛聲.
神啊.
我們很慘.
出來後就發怨言.
很有趣.
人要麼發愛聲.
要麼發怨言.
以色列就像我們這個人版.
聽見了.
我們算什麼呢.
你們的怨言不是向我們發的.
是向耶和華發的.
第三點跟大家說.
你的怨言向人發.
如果向領袖發.
特別是摩西亞倫.
你的怨言是向耶和華發的.
即是神說.
你向他發怨言.
罵他其實是罵我.
是我叫他帶你出來.
其實如果你是上帝.
你開心嗎.
如果我是上帝.
真的很慶幸.
但神的心很闊.
很容忍.
發怨聲救他.
現在怨言.
沒東西吃.
還說以前在玉窩旁邊那麼厲害.
我又幫幫你.
這樣神真的做.
給一點關子.
放一點關子給大家.
神接著做什麼呢.
又用東風吹來很多鵪鶉.
多到鋪在地上.
一尺.
那麼多.

$^{441}$隨便捉.
他們吃鵪鶉肉.
吃到鼻子流出來.
晚上吃到飽.
然後早上起床.
露水過去.
看到一點點白色.
白色是什麼呢.
就是麻辣.
下一次說.
神是晚上給他們吃鵪鶉肉.
早上開始吃麻辣.
麻辣很神奇.
是每天的神蹟.
和每周的神蹟.
因為到了第七天.
第六天要多執一倍.
明天四天不能執.
平時執定一個星期.
帶兒子出去執.
沒用的.
因為明天麻辣就壞了.
每天只能吃當天.
給我們日常飲食.
到第六天執明天.
這樣又不難.
真奇怪.
所以你看到.
上帝沒有大腦.
下次我們更加看到.
在這裡要小心.
你要知道.
你發的怨言.
耶和華是聽到的.
我記在心裡.
他很警惕.
不要發怨言.
不要向耶和華發怨言.
因為他是好的.
虐失神體中.

$^{481}$他真的很好.
為什麼呢.
他雖然經歷這樣的事.
又被哥哥賣了等等.
又被人陷害.
又坐牢.
他沒有發怨言.
然後他遇到哥哥.
他說其實這件事是好的.
你們想害我.
不過在神面前是好的.
差了我來.
預備到老.
挺好.
我們要懂得.
欣賞上帝有祂奇妙的工作.
不要每次發怨言.
怨言是破壞關係的.
我們跟人相處.
跟神那麼大.
更加不要發怨言.
跟人不要發怨言.
少一點.
多一點體諒.
多一點明白.
發怨言.
他可能有背景傷害.
他發怨言多一點.
有些人在這件事上.
一碰到就爆炸.
因為他以往有不同的經歷.
所以要知道.
小心.
不過最重要是.
向神發怨言.
是很危險的.
結束了.
在這個時候.
記得.
當耶穌時代.

$^{521}$一班百姓要將耶穌釘他十字架的時候.
彼拉多說.
我查不出他有什麼罪.
他那些不是搞遊毒.
不是反叛.
他只是說.
他的國不屬於這個世界.
他是說天國的事.
現在我要將他責打.
警告他就放了.
不行 釘他十字架.
彼拉多就說.
留這個人的血.
歸在誰身上.
不關我的事.
我要放他.
那班猶太領袖就說.
歸到我們和我們子孫身上.
哎呀.
我讀到這裡.
手裡就真的.
刺一刺.
好痛.
為什麼呢.
其實跟著這一班以後.
以色列的後裔.
就受了很多血的傷害.
歸到我們和我們子孫身上.
所以猶太人.
在二次大戰.
其他地方.
都受了很多傷害.
死的人非常多.
是害得很慘的.
不要以為二次大戰.
希特勒才害猶太人.
其實以往歐洲很多地方.
都害猶太人.
英國殺死不少猶太人.
甚至宗教改革的領袖.

$^{561}$馬丁路德.
他開頭是支持這班人的.
猶太人.
但後來發覺到.
他們很得意.
不行.
終於馬丁路德.
是很反對猶太人的.
當然他沒有親手殺他們.
不過他是支持.
這班猶太人受苦的.
所以其實呢.
千萬不要向上帝.
貿然講說話.
很危險的.
我們一起祈禱吧.
天父我們.
向你首先認罪.
很多時候我們.
遇到一件事情不順利.
我們就會問.
為什麼會這樣.
為什麼神你不祝福.
我們都會埋怨.
當然我們會先埋怨人.
然後心裡不太敢.
不過都說了.
神為什麼你這樣.
為什麼你做事這麼安排.
主要我們在.
《初一及記》裡學到的.
就是不要發怨言.
摩西說.
你向我發怨言算什麼.
你向神發怨言.
很危險.
讚美主.
因為你是.
心胸很闊的主.
這些以色列人.

$^{601}$雖然二百萬人.
很多成年人.
不過都好像小朋友一樣.
時時亂講說話.
時時發出怨言.
或者發出哀聲.
很多都在自憐與自恨裡.
主要幫助我們脫離.
這些自憐與自恨.
脫離對人的怨恨.
我們更光明地.
走我們人生抗疫的路.
求主幫助我們更喜樂.
祈禱奉耶穌名求.
阿們.
\newpage



\section{哥林多前書 1:26-29}
\label{sec:032018}
\textbf{神國的精英主義}~~~~林可豐傳道
\newline
\newline
連結: \href{http://www.church.com.hk/acms/content.asp?site=cdc\&op=show\&type=product\&code=032018\&layout=sermon}{\texttt{032018}} ~~~~ 語音日期: 2026-04-19
\newline
\newline
\hyperref[sec:032017]{\small{< < < PREV SERMON < < <}}
~
\hyperref[sec:index_chronic]{\small{[返順時目]}}
~
\hyperref[sec:032019]{\small{> > > NEXT SERMON > > >}}
\newline
\newline
哥林多前書 1:26-29
\newline
\begin{longtable}{cl}
\hline
\hline
章節 & 經文 (和合本修訂版)\\
\hline
1:26 & \begin{tabularx}{0.7\textwidth}{X} 弟兄們哪，想一想你們的蒙召，按著人的觀點，有智慧的不多，有能力的不多，有尊貴地位的也不多。 \end{tabularx} \\ \\ \relax
1:27 & \begin{tabularx}{0.7\textwidth}{X} 但是，神揀選了世上愚拙的，為了使有智慧的羞愧；又揀選了世上軟弱的，為了使強壯的羞愧。 \end{tabularx} \\ \\ \relax
1:28 & \begin{tabularx}{0.7\textwidth}{X} 神也揀選了世上卑賤的，被人厭惡的，以及那一無所有的，為要廢掉那樣樣都有的， \end{tabularx} \\ \\ \relax
1:29 & \begin{tabularx}{0.7\textwidth}{X} 使凡血肉之軀的，在神面前，一個也不能自誇。 \end{tabularx} \\ \\
[1ex]
\hline
\hline
\end{longtable}
$^{1}$好 姐妹平安.
話說昨天就.
聖事連崇拜嘛.
兩個 即是近來有些人.
post了一些什麼十年前的.
即是Facebook片啊 IG相片出來.
那些小朋友看到.
然後有個中醫生這樣.
見到我 嘩 海峰.
原來你以前是瘦的.
嘩 多謝你.
多謝你讚我以前.
證明了我結婚之後很幸福.
那麼思想回到十年前的時候.
我是這樣看的.
有些人經常說減肥是節食啊.
大家知道如果認識我的話.
我現在很小吃.
不過十年前我超大吃的.
平時如果出去大家樂.
我以前可以吃兩個焗豬扒飯.
可以吃光的.
還是吃旁邊的人吃剩的那些.
但是也這麼瘦.
就是因為我天天都做運動.
經常打球.
所以我的減肥哲學就是.
你做得夠運動.
你的輸出夠多.
其實你是不用節食的.
所以大家都要做運動.
那麼.
那時候我經常打球.
大概十年前就是中學畢業.
讀大專左右.
天天都去打球的.
那麼我打球的時候.
可能大家以前有做過運動.
或者現在也在做運動.
看到有時候星期六.

$^{41}$有些叔叔阿姨和三味傳道Anna.
他們一起打羽毛球的.
很有火的.
有時候打球出去.
打比賽鬥牛.
打籃球有個term叫.
好像和牛對抗一樣的.
很激烈的.
那麼一定打比賽要想著贏的.
以前有一段時間我在觀塘.
不是.
不是觀塘.
調景嶺的HKDI讀書的.
那時候.
調景嶺在籃球學界來說.
都挺厲害的.
其中一間寶國中學.
很厲害的.
這十年都是經常去到學界.
最後四強.
又在調景嶺的.
調景嶺的街場是飛天遁地的.
中學生都進了瓶子.
很厲害的.
那麼.
HKDI那裡.
都是有很多打籃球很厲害的人.
有一個同學.
很厲害的.
他可以進瓶子的.
就是進瓶子.
就是.
用手把球丟進籃球框裡.
然後射進去那些.
就是跳得很高的.
很大隻的.
就是.
不是那麼容易的.
現在經常有很多中學生做到的.
不知道為什麼.

$^{81}$十年前的時候沒有那麼多人.
體能那麼厲害的.
那麼那時候我們一起打球.
剛才說了.
年輕人打球當然有火.
想贏的.
所以那時候打球.
經常都很想.
全部人都是.
說到學校.
我們都是籃球隊.
特別想跟他一隊的.
因為有.
他在我們那隊的時候.
就代表什麼呢.
代表我們有全校最好的球員在.
我們呢.
最容易可以去贏球.
這個隊友.
這個可以進瓶子的同學.
就是在學校裡面.
最有價值的隊友.
最有價值的球員.
那麼其實呢.
一個這樣的思維呢.
在世界很多不同的層面.
都有的.
什麼思維呢.
例如在學校裡面.
其實讀書成績最好的學生.
會最被高舉的.
或者有時候在公司裡面.
去請人.
這些最有能力的人.
就最搶手了.
或者在社會裡面.
哪一些人說話最有說服力呢.
可能就是.
在社會上面最有成功地位的人.
其實這種是很自然的價值觀.

$^{121}$但是這種價值觀.
是什麼價值觀呢.
就是一種精英主義的價值觀.
那些優秀的人呢.
才是重要.
才是有價值.
才是值得大家投資在他們身上.
但是.
今時已經不同往日了.
十年時間過去了.
當我剛剛在.
二月的時候.
新年啦.
農曆新年啦.
有些時間.
那很好可以.
跟一些年輕人一起去打球.
我再一次踏上球場.
我發覺.
原來呢.
人是會改變.
是會進化的.
當我再次踏足球場的時候.
我完全對勝利毫無慾望了.
我只是發覺.
原來我幾個月沒有做過運動了.
現在腰很痛.
三十歲就.
慘了.
那對了.
剛剛說到.
上個月新年假的時候.
有一次.
唐慶聯合崇拜.
那星期六.
沒有聚會.
有一個大專生就.
興起就說.
不如跟一些中一生.
一起去打一下球.

$^{161}$這樣.
那我和那個大專生就分開了.
跟他們一起去.
比球.
那我想起.
以前我打球的時候.
就是.
很熱血的時候.
我覺得打球是很好玩的.
這件事.
那.
就是當我剛剛那次星期二.
不是同.
不是星期二.
就是二月的時候.
和一些年輕人打球的時候.
我就在想.
啊.
就是.
我經常都.
就是經常突然間想很多事情.
打球的時候就在想.
其實我現在是不是覺得很好玩呢.
跟他們一起玩這樣.
在想我在做什麼呢.
是不是純粹在打工呢.
要陪陪他們這樣呢.
這樣.
在想的時候呢.
我就發覺.
其實我都真的覺得不是很好玩的.
就是和這些小朋友.
就是.
說真的.
超煩.
這些小朋友.
就是.
不是這兩個.
這兩個很乖的.
就是可能認得出來.

$^{201}$很乖這兩個.
還有一些在畫面外面那些超煩.
就是.
有.
打籃球世界裡面有一個叫Trash Talk的說話.
就是Trash Talk是.
就是.
不是說廢話.
就是Trash Talk的意思是.
就是其實都是廢話來的.
但是類似就是.
就是挑釁對方.
就說你是不是廢的.
你懂不懂的.
這樣.
你.
就是類似這樣.
不斷地挑釁.
那.
但是.
這些小朋友.
不斷地說這些.
就是.
但是.
就是.
他怎麼會救我們兩個打呢.
太高度也差三個頭.
這樣.
不斷地射球.
一砰.
射不到球.
但是.
極多Trash Talk.
極煩的這些小朋友.
雖然呢.
其實是不好玩的.
就是在籃球世界裡面.
就是和他們鬥.
什麼都.
球都拍不到兩下.

$^{241}$射球就全部進不了.
都進得了.
進得了一下球.
但是就是.
不是.
我以前.
和同學一起打球那麼好玩.
這樣.
但是再摸摸心.
其實不是很好玩的.
但是開心嗎.
那天.
我想起.
那天其實很開心.
那很開心的那個位置.
就是.
就是發現.
成長了.
不用贏才行的.
其實做一下運動.
可以動一下.
這樣.
又看到小朋友一起.
挺開心的.
其實大家都笑得很開心.
其實那天.
特別呢.
當我就是.
誒.
特意拍了這張照片.
這樣.
幸好拍了背面.
這樣.
那天.
不然就沒想用.
那麼.
很開心的其中一個看見.
就是.
當我看到那個大專生.
他也很用心地和小朋友在那裡打球.

$^{281}$這樣.
又.
就是.
是啊.
他也被小朋友打了很多次.
不過很厲害.
就是.
就是被人打.
就馬上射進去.
這樣打到他們沒聲音了.
打了一下之後.
這樣.
但是呢.
就是.
很好這個年輕人.
就是.
知道.
他們嘴巴很臭.
但是呢.
看到他們.
其實有些位置.
你可以再進步一點的.
這樣.
將來在校隊這樣打就厲害了.
那麼.
很用心地和他們打.
教他們打球.
有幾個小朋友.
那天不懂得打球.
他也.
很完全.
就是.
沒有.
因為.
就是.
平時打球.
你當然想射一下球.
鬥一下球的.
但是有些小朋友.
不懂得打球.

$^{321}$他沒有嫌他們阻礙到.
很用心地去.
教他們打球.
我就看到這個.
弟兄姊妹這個圖畫.
不是因為那些小朋友.
就是.
都可愛的.
其實.
但不是他們很乖.
這樣.
或者和他們一起的時候.
我們的心情會很好.
有時候會很頂癮的.
不是因為誰厲害.
是因為.
我們大家願意.
這樣去走在一起.
願意這樣去生命.
躺開.
大家一起去交流.
成為弟兄姊妹.
為什麼願意這樣做呢.
我看到是因為.
我們都知道.
特別是那個大轉身.
我看到.
他知道.
我們這樣可以.
一起成為弟兄姊妹.
是因為耶穌基督的救贖.
我們今天可以一起.
成為教會.
一起成為一家人.
我們有關係.
不是各樣的工作.
逼大家.
你是大人.
你是教會組長.
你就一定要和小朋友一起.

$^{361}$不是.
而是因為.
我們知道.
大家都是耶穌很重視的人.
大家都是一家人.
是弟兄姊妹.
所以這個年輕人.
這個大轉身.
願意這樣.
將他的時間.
衛生在年輕人的身上.
剛才我們一起去讀的經文.
我們今天所看的經文.
保羅就是在說這件事.
在神國的.
價值觀裡面.
上帝看人的價值裡面.
和世界看人的價值.
是不同的.
袁勝寧今天.
也去教訓我們的眼光.
讓我們可以有.
上帝看人一般的眼光.
去看我們身邊每一個弟兄姊妹.
我們一起去祈禱.
可以經歷.
你愛之耶穌基督的救恩.
也讓聖靈在今天.
這個世界.
我們可以明白上帝你的心意.
你讓我們可以在地上.
這個破碎的世界當中.
有教會群體.
有我們身邊的弟兄姊妹.
也有上帝你華語在當中.
充滿我們教導.
讓我們明白認識你.
有聖靈引導我們.
去看我們身邊的弟兄姊妹.
我們可以明白上帝的心意.

$^{401}$引導我們的心.
讓我們願意去親近上帝你.
我們去獻上感恩.
教會是上帝你天國的養子.
讓我們可以去經歷.
讓世界去看見.
讓我們可以預想到.
我們將來永恆.
和上帝你一起的美好的地方.
是一個有形的群體.
讓我們去認識到上帝你自己.
去經歷到上帝你自己.
但我們在教會群體當中.
我們仍然有很多人的軟弱.
在裡面.
我們有時候會看不到上帝你自己.
有時候我們會忘記了上帝你的教導.
有時候我們沒有用上帝你喜悅的方式.
去對待我們身邊的人.
去看大家成為弟兄姊妹.
但願上帝你華語今天都去教導.
讓你的靈去充滿我們.
讓我們今天可以在你的教導.
你的提醒當中.
我們重新拿回一個上帝你的眼光.
去看我們自己.
我們在你眼中是一個怎樣的兒子.
也用你的眼光去看我們身邊的弟兄姊妹.
祈禱奉主耶穌基督得勝 聆聽者者 求.
阿們.
哥林多前書我不是第一次去講.
我之前也有用過哥林多前書.
也有簡短講過哥林多前書的一些背景.
哥林多前書的一些背景我再講一次.
就是保羅在他的宣教工作當中.
有聽到他曾經去建立服侍過的教會.
哥林多的教會有些弟兄姊妹向他投訴.
就說在教會裡有些人分黨分派.
有些人覺得自己是跟保羅出身的.
就比較輕鬆.

$^{441}$有些人跟彼得出身就比較好.
有些人跟其他人.
有些人跟耶穌 我最大.
分黨分派.
之前也有講過.
其實是不是很頑皮.
這不是的.
哥林多教會裡的人其實都很乖.
都很竭力去服侍.
但有時候當我們很竭力去服侍.
但走在一起的時候.
我們人的不同的靈覺就全都飛出來了.
然後我們就經常互相比較.
互相增敬.
沒有了弟兄姊妹的樣式.
所以今天簡宣的經文.
就哥林多前書一章.
其實就是保羅聽到你們教會的一些事情.
講了一大堆事情.
就是講了一大堆事情之前的一個引言.
保羅嘗試去捉緊哥林多教會的人的焦點.
提醒他們你們有些事情有些問題.
你們有些價值觀 有些觀念要調整.
就是這個引言做了的工作.
今天我們去看的經文一開始去講.
保羅就說.
按著肉體的有智慧的不多.
有能力的不多.
有尊貴的都不多.
很明顯就是在說.
剛才他在說什麼分黨分派.
分黨分派不是只是跟誰的.
而是背後他們在說.
這個老師這個我們所跟的牧者.
他有一些特點的.
這些才是對的.
這些才是正的.
可能行動派是好一點的.
學術派是好一點的.
各樣各樣.

$^{481}$那麼保羅就是很特別的.
你們可以想一下.
保羅其實在哥林多教會.
或者在普遍人的眼中.
他其實有沒有智慧的呢.
有沒有能力的呢.
有沒有地位的呢.
你們大概一想.
你們回內教會都知道.
保羅不是那麼簡單的.
這個人厲害的.
我讀書都知道保羅厲害的.
讀到我多了很多書要讀.
保羅在教會界裡面.
在歷史界裡面.
在猶太人裡面.
他是一個什麼背景的人呢.
他是一個法利賽人.
法利賽人有時候看耶穌會話.
他們說得多.
但是他們其實是什麼很差的人.
不是的.
法利賽人是很厲害的人.
他們是學者來的.
所以保羅.
其實在說.
他在鎮壓.
在下馬威.
現在我是保羅.
你們曾經的牧者.
我跟你們說.
我雖然是你們的牧者.
但是我怎麼看我自己呢.
我不是很有智慧的.
我都不是很有能力的.
我都未必很有尊貴地位.
但實際他是不是沒有這些東西呢.
可能其實未必是.
其實是有.
他都這樣謙卑.

$^{521}$自己自嘲.
你們自己在那裡.
大家自己覺得自己比我厲害.
大家謙卑一點.
可能是在表達這件事.
但也就是說.
不是一個普通的謙虛.
而是在說.
在上帝的教會裡面.
在神國的眼中.
不是用這些東西.
來定義價值.
不是一個精英主義.
我們以前的人.
看敵人的價值和我們今天看人的價值都很相似.
當時的人或今天的人也好.
定義人的價值.
那個邏輯是怎樣呢.
能力越高就越重要.
表現越好就越被人看見.
成績越好就越被尊重.
這種世俗的價值觀.
其實不知不覺地.
滲入了教會裡面.
很多時候其實大家不是故意的.
但是我們的心.
每天都在外面生活.
世界就是這樣運作.
我們就用了這樣的方法.
去定義去評價去看我們身邊的弟兄姊妹.
例如可能我們有時候在教會裡面.
大家就去一點點檢視.
我們有沒有這樣的眼光.
我們會覺得有某一些.
教會服事的崗位.
是特別重要的.
特別輕香的.
或者在教會裡面某一些服事的人.
會覺得這個人是一個服事很好的人.
他是一個值得我愛的人.

$^{561}$這些服事.
就沒有那麼輕香.
沒有那麼有價值.
或者某些人就覺得.
你服事的質素有點差.
不是做得那麼好.
你就不是一個那麼有價值的人.
我們會覺得某一些崗位.
或者某一些人.
他們表現不好的時候.
我們就會覺得他們沒有那麼有價值.
或者在另一些情境.
有沒有一些人.
我們的組員就會覺得.
跟他很容易溝通.
他講話很中聽.
他修養很好.
我就願意跟他成為弟兄姊妹.
願意跟他一起去相處.
多一點聊天.
但是有些人可能.
他是安靜一點的.
慢熱一點的.
沒那麼會說話的.
話題終結者來的.
經常只顧著自己說話.
不聽別人說話.
盡量不要跟他說那麼多話.
他說完話就不回應他.
會不會有些.
身邊有這些人.
我們是這樣看他們的呢?.
這樣想他們?這樣對待他們呢?.
但是保羅在.
經文繼續說.
在經文之後的.
段落裡面.
保羅三次說了同一句說話.
「神是揀選了愚拙的」.
「神是揀選了軟弱的」.

$^{601}$「神是揀選了卑賤的」.
三次用了同一個字.
「揀選」.
上帝是怎樣看待人的呢?.
上帝.
看待人.
在這裡如果純粹去看的話.
他是揀選的.
如果你很用邏輯.
你會說揀選也有些條件的.
沒有啊,你不知道.
你不知道上帝是揀選誰.
上帝怎樣定義.
哪些人有價值呢?.
是他揀選的.
你看看你身邊很多人.
我們有沒有些東西共通呢?.
有的,有頭髮,有眼耳口鼻.
背景很多不同的,讀不同的書.
不同性別,不同性格,高矮肥瘦都不同.
但上帝就是這樣揀選我們.
上帝看人的價值,他揀選人的價值.
他叫人得到的價值.
不是看誰有厲害一點.
誰有什麼才能,誰有什麼條件.
單單是他看中了.
那個人就有價值了.
就有揀選了.
這樣說的話.
如果你回教會久一點,讀過一點神學的話.
你開始蠢蠢欲動了.
揀選這麼簡單嗎?.
這個字,預定論還是普及論.
這個世界上所有人都可以經歷耶穌基督.
還是上帝一早在創世二線.
已經揀選了.
那些人想很多東西.
我們今天不要這麼複雜.
不要搞這些爭論.
哥林多我們繼續去讀.

$^{641}$之前有提到.
哥林多前書十二章一開始這樣說.
耶穌是主.
我們可以這樣去有這個宣認.
我們可以說得出.
耶穌是主,為什麼呢?.
是因為有聖靈在我們生命的工作.
聖靈在我們的心.
就是神揀選了我們的一個狀況.
所以今天.
在座各位.
誰覺得自己是一個基督徒呢?.
誰都會說.
耶穌基督是你們的主呢?.
如果你們心裡這樣想.
耶穌基督是我們主的時候.
你就是上帝揀選的人.
不要管那些什麼預不預定.
普及論這些.
你今天說得出耶穌基督是你的主.
你就是被上帝揀選的那些人.
所以在神的群體裡.
教會裡.
基督徒群體裡.
所有的人.
你覺得他們不顯眼的.
或者尖銳的再去說.
在教會裡.
不太喜歡的弟兄姊妹.
他們不是多餘的.
他們不是沒有價值的.
在教會裡.
你看不到服侍崗位的.
做後勤的那些人不是次等的.
那些走得比較慢的.
不應該是被忽略的.
因為.
在教會裡.
每一個弟兄姊妹.
都是神揀選了的人.

$^{681}$都是神很寶貴的兒女.
你不知道.
你身邊的人的生命.
特別是你看不起的人的生命.
他們會怎樣被上帝運用.
在我們教會裡.
每個禮拜有很多不同的聚會.
很多不同的事工.
有很多弟兄姊妹在服侍大家.
也有很多弟兄姊妹在出去.
這個世界服侍這個世界.
讓福音被人看見.
被人經歷.
不是因為.
服侍很輕鄉.
就是.
按powerpoint在那裡彈琴.
有什麼輕鄉呢?沒有人供收的.
出去玩音樂不好嗎?.
坐在那裡享受一下冷氣不好嗎?.
九點才回來不好嗎?.
為什麼要八點七點回來呢?.
相信就是因為.
每一個在那裡服侍的.
或者出去參加.
不同的事工在那裡服侍的.
珍惜.
大家的生命.
看大家是弟兄姊妹.
在座的各位.
有沒有當大家是弟兄姊妹.
去服侍呢?.
我們其實是在用一個世界的眼光.
在用一個公私的角度.
論資排輩的.
用能力的.
去看我們身邊的弟兄姊妹.
想清楚一點.
就是.
剛才他講的.

$^{721}$人其實是會有轉變的.
以前可能我打球的時候.
都會覺得.
要打贏才會開心.
好玩的.
但今時今日就發覺.
其實不是的.
很多時候人生能夠經歷到.
那種喜樂那種滿足.
或者去經歷到上帝的同在.
經歷到上帝在當中.
弟兄姊妹那種團契.
就是一件很美好的事.
這件美好的事情.
是怎樣成就呢?.
不是在說打球好不好.
不是在說一起相處的時候有沒有質素.
而是在說我們.
有沒有主內的關係.
我們和上帝有沒有關係.
我們和身邊的弟兄姊妹有沒有關係呢?.
剛才在引言那裡我講了.
有位大專生和我們一起去和中一打球.
說真的不是那些中一很聽話.
那些中一真的.
是很激氣的.
但為什麼仍然願意.
這樣去花時間.
在這些中一身上呢?.
我相信是因為.
那個大專生.
他願意用上帝的眼光.
去看這些中一生.
他很珍惜上帝都很愛的生命.
他看大家是弟兄姊妹.
他重視他們.
很想進入他們的世界.
很想和他們建立關係.
讓他們在牧養的經歷裡.
去經歷了.

$^{761}$除了一起去教導.
一起去學習.
更加是經歷到那種同行.
那種同在.
那種很重視的生命.
相信很多在教會當中.
我們一起去服事的人.
我們都是這樣經歷了上帝的愛.
經歷了弟兄姊妹的服事.
我們先嘗試去效法基督耶穌的教導.
我們都嘗試去服事人.
保羅繼續在29節說.
為什麼上帝會這樣做呢.
當上帝這樣做的時候.
就令到一切有血氣的.
在神面前.
一個都不可以自誇.
神這樣去揀選人.
神用這樣的眼光.
去對待人.
透過保羅這樣去教導人.
喂,在上帝的國裡.
看著我們.
喂,在上帝的國裡.
看人的價值和世界看人的價值.
應該是不同的.
就是要用到一個.
天國的價值觀.
去拆毀一個屬世.
甚至有魔鬼在當中使用的價值觀.
這個價值觀.
是什麼價值觀呢.
一個自我驕傲的價值觀.
一個人覺得可以自己.
去令自己有價值.
自己掌管自己生命.
一個這樣的價值觀.
驕傲的價值觀.
這個驕傲的價值觀.
大家不要小看.

$^{801}$當時在創世記當中.
亞當和夏娃吃分別善樂樹果子的時候.
最最大的問題.
不是吃了那顆蘋果.
或者吃了那顆分別善樂樹的果子.
而是在說.
他們想成為和上帝同等的人.
什麼叫和上帝同等.
一開始創天造地的時候.
誰定義什麼是好還是不好.
什麼叫好還是不好.
是上帝定義的嗎.
不是,是上帝定義什麼是好還是不好.
這個世界的律是由上帝去定義的.
分別善惡,什麼叫分別善惡.
如果說回原文.
就是定義什麼是好還是不好.
吃了分別善樂樹的時候.
人就有和上帝一樣的權柄能力.
去定義什麼是好.
和不好.
所以亞當和夏娃犯的罪是什麼罪.
他們想和上帝一樣.
不聽上帝的話.
想自己成為神.
他們想自己定義自己的價值.
他們想自己定義世界的價值.
這個就是驕傲.
也是今天當我們.
充斥著很多成功主義.
精英主義的時候.
背後隱藏著的一個心態.
我們要自己定義自己的價值.
所以當我們今天都是.
用一個這樣的價值觀.
去成為別人的弟兄姊妹.
在教會當中去服侍的時候.
我們就會有這樣的思維.
就是我可以這樣去服侍.
我可以在教會有某一些的崗位.

$^{841}$是因為我好.
是因為我厲害.
教會可以成為今天這樣的地方.
是因為有我在.
不過這句話可以正面看.
或者反面看.
今天教會搞成這樣就是因為我.
anyway.
大神的國系怎麼說呢.
我們今天可以成為弟兄姊妹.
我們今天可以在教會服侍.
是不是因為我厲害呢.
不是.
我們今天可以成為弟兄姊妹.
是因為上帝的簡選.
完全是因為上帝的恩典.
上帝的簡選有沒有條件呢.
沒有條件的.
上帝看上了每一個.
祂所創造的.
看上今天在座的大家.
每一個坐在這裡.
你說耶穌基督是你的主.
你都是上帝所簡選.
完全是上帝的恩典.
所以上帝透過保羅的教導.
跟我們說.
上帝的國裡面真正的精英.
不是那些最厲害的人.
而是那些.
毫無驕傲.
謙卑的知道自己也是蒙恩的罪人的人.
這些人.
就是最知道.
你的價值的人.
我再說一次.
在上帝的國裡面大家都有價值的.
但是是不是真的大家知道自己有價值呢.
你是不是真的知道身邊的人.
都是上帝很著緊的人呢.

$^{881}$哪些人才知道你自己有價值的人呢.
哪些人才知道.
上帝怎麼看你.
怎麼看身邊的人.
那些忠心在上帝的路裡面的人呢.
這些人應該是那些.
毫無驕傲.
謙卑的知道你今天可以坐在這裡.
都完全是因為上帝的恩典.
我們每一個.
都是蒙恩的罪人.
唯有這樣看自己.
這樣看身邊的人.
你才真的知道.
你在上帝裡面是有價值的.
當我們這樣去看自己.
又這樣去看身邊的人的時候.
我們又應該.
成為一個怎樣的教會呢.
我們應該怎樣一起去服事呢.
我們怎樣會一起去議事.
一起去開不同的會議.
一起去做不同的事工呢.
我們應該就是.
不是去互相比較.
不是去看那些事工做得好不好.
說真的.
教會的事工做得好不好.
我們有沒有能力呢.
真的計劃得好不好呢.
如果你這樣想的時候.
你就不是相信上帝是教會的主.
如果我們真的這樣去看.
自己是蒙恩的罪人的時候.
大家都是上帝很珍愛的兒女.
每個人都是我的弟兄姊妹的時候.
我們一起成為教會的時候.
比起那些事工做得好不好的話.
我們應該更加著緊的.
就是.

$^{921}$我們願不願意和上帝.
願不願意和人一起去建立關係.
我們今天去看教會的時候.
更多又會是著緊在教會的事工.
或者聚會做得好不好.
還是著緊我們這個群體裡面的人.
大家有沒有真的成為弟兄姊妹呢.
拿自己的心出來.
去和身邊的人.
生命去有交流.
去建立關係呢.
總結了.
所以當我們今天再一次去思想.
世界怎樣定義人的價值呢.
是一種精英主義.
那些最強的人.
最優秀的人才有價值.
但神國的精英主義.
講的是什麼呢.
神國的精英主義.
單純就是講神愛世人.
剛才敬拜主令都用了.
約翰一書的經文.
講我們人類的價值.
講我們.
今天是經歷了上帝的愛.
而我們應該.
都這樣去愛我們身邊的弟兄姊妹.
這個是耶穌基督的教導.
是大解明當中的內容.
所以我們今天都成為一個.
真的一起去跟從耶穌基督.
效法耶穌基督.
去聽耶穌基督教導的一個群體.
我們不是要成為一個很完美.
很漂亮的階級.
而是我們.
要知道.
我們每一個都是蒙人的人.
我們每一個.

$^{961}$身邊的每一個.
都是上帝很愛的人.
因此我們不會去.
互相憎恨比較.
甚至我們要去.
好像耶穌基督去服侍人.
去為人洗腳.
去為人抹餅.
去為人寫生一樣.
去走向你身邊的弟兄姊妹.
特別主動.
跟那些你身邊不熟悉的.
甚至是跟你.
同一個小組.
可能已經五年十年二十年了.
但你其實從來都沒有真的.
用心去認識他.
一起去跟他相交的弟兄姊妹.
這些人.
我們應該要走向他們.
否則我們是否還是用上帝的眼光.
去跟我們身邊的弟兄姊妹一起團契呢.
我們還是不是一個教會呢.
去到結尾.
其實之前在長青.
已經講過這篇訊息.
有個姨姨跟我說.
當你講完這篇訊息的時候.
令她想起有一位在我們這裡.
的大前輩.
就是鄭一妹婆婆.
即是禮鄭一妹會座.
尊重地稱呼.
禮鄭一妹.
鄭一妹婆婆.
是誰呢.
就是大家在座很多人的.
太婆太媽.
即是很長輩.
我們今年.

$^{1001}$五十六年了.
而當我們未成為.
筲箕灣浸信會之前.
我們這個群體當初是怎樣被建立的呢.
就是當時的傳道宣教士.
以前筲箕灣.
西灣河這裡是一個什麼地方.
漁村來的.
傳道人宣教士.
上船這樣去傳道.
上船這樣去崇拜.
而鄭一妹婆婆.
就是在當中的時候.
在船上這樣去信主.
而她就帶領她的家族.
去信主.
令到今天在座的你們當中.
我們偶爾會聽到的黎氏家族.
很大的家族.
都是這樣去信主.
鄭一妹婆婆帶領了她的家人.
一起去信主.
甚至就建立了一個很大的群體.
甚至到今天我們仍然看到.
在我們教會裡.
很多的弟兄姊妹.
都是因著鄭一妹婆婆.
當日她這樣忠心.
被上帝去揀選 被上帝去使用.
祝福了很多的後代.
也都因著很多後代的事工.
今天很多的弟兄姊妹.
是這樣被邀請回到我們這裡.
而那位姊妹.
她跟我說.
她在想的時候.
就是整天人們想起鄭一妹婆婆的時候.
會怎樣去評價她.
她是一個很賢淑的婦人.
很賢淑的婆婆.

$^{1041}$做人十分正直.
很竭力愛主.
參與教會的聚會.
很投入.
再去想的時候.
她是一個福音大使.
因著她而回教會.
信耶穌的人有多少呢.
數一數小則的.
有三位數.
再去想一下教會這五十多年.
有沒有人帶了三位數這麼多人的信主呢.
不知道.
我想起來真的沒有.
再去想深一層.
這位姊妹婆婆.
是不是一個讀過很多書的人呢.
是不是傳道人呢.
會在執事的時候.
她是一個信徒.
她不是傳道人.
她沒有讀過神學.
但今天回想.
我又想自己 我也是傳道人.
我應該要傳道傳福音.
為什麼這麼多人信主呢.
大個人回教會也有的.
數起來十個也很多.
三位數的人信主.
我真的做不到這件事.
但上帝就是這樣使用.
鄭一妹婆婆 不是她讀了多少書.
又再想近一點.
如果再想.
今時今日.
在我們的群體裡.
有沒有人很大大地被上帝使用.
走到哪裡都傳福音.
整個西灣河都可以認識她.
立即又想起一個人.

$^{1081}$就是想起Anna.
Samuel的太太.
有一次我們一起開玩笑.
很怕和Anna一起搭電梯.
特別是.
上她家玩的時候.
因為我這個人比較內斂.
通常在街上.
不會突然間跟人聊天.
但Anna是會做這些事.
搭電梯時會和身邊的.
小朋友聊天.
有些中年夫婦.
以為她認識的.
之前和她聊過天.
不是呀 第一次見.
第一次見好像見了很多次.
突然間.
喂 吃飯了嗎 今天吃飯吃什麼.
超勁.
但Anna不只是那種外向.
教會我們經常做很多外展活動.
無論那些外展活動.
有時候我們覺得.
我們身邊有沒有朋友.
或者年青人.
怎會喜歡這些活動.
但Anna是不會這樣去想教會的活動.
她只會想每一次教會有任何聚會.
都是機會帶一些人回來.
認識主耶穌.
我記得有一次.
聖誕節晚會.
她搭電梯見到外婆.
外婆 聖誕節有什麼做.
就邀請她來.
外婆真的來了.
我們就想起.
在我們的群體裡面.
Anna就是這樣被上帝使用.

$^{1121}$亦是這樣展示可以這樣和人.
有生命的交流 可以這樣認識人.
可以這樣傳福音.
但Anna經常會怎樣說自己.
她就會說自己.
她不喜歡讀書.
不喜歡成為領袖.
不想站在某些崗位裡.
但是當我們在青少年群體裡.
經常想.
我們要怎樣傳福音.
我們群體裡應該怎樣去做.
某一些見證.
去傳福音 去做外展的時候.
我們一定會想起親愛的Anna姊妹.
上帝使用人.
讓人去聽到福音.
剛才想起.
請一位外婆.
想起Anna.
不知大家怎樣去定義她.
我很欣賞.
我覺得Anna是一個很厲害的人.
甚至經常說Anna你應該去讀神學.
你做傳達人應該做得好過我很多.
為何不讀呢.
上帝有她的心意.
有她的帶領.
但在當中我們看到上帝揀選人.
上帝使用人.
不是用我們的價值觀.
教會請人可能.
就是要看這些東西.
我們要選哪些領袖可能要看這些東西.
讀過什麼課程.
上過什麼裝備.
什麼神學院出身 什麼背景.
上帝不是這樣揀選人.
我們可以從耶穌基督.
昔日的服事中.

$^{1161}$我們可以知道上帝是怎樣揀選人.
他的群體他使用的人.
是怎樣的人.
耶穌基督.
昔日來到這個世界的服事.
他對待的 他走向的.
他服事的.
他看上的.
就是一顏一忽的人.
不是只是那些一顏一忽.
應該是每個人都是一顏一忽的.
但他都會和一些正正老老的人.
會和法理錯人聊天.
有辯論.
但是他特意去服事的人.
很多都是.
是妓女 是罪人.
耶穌會和他們.
一起去相處.
建立關係.
去吃飯.
所以我們今天.
我們又沒有用.
耶穌基督的教導.
他所展示的眼光.
去看這個世界的人.
去看我們身邊的弟兄姊妹.
耶穌這樣去服事我們.
我們有沒有一樣.
用這樣的眼光.
願意這樣去服事我們身邊的弟兄姊妹呢.
\newpage



\section{民數記 10:1-36}
\label{sec:032019}
\textbf{起行與安營:聆聽神的信號,跟從祂的節奏}~~~~陳傳華牧師博士
\newline
\newline
連結: \href{http://www.church.com.hk/acms/content.asp?site=cdc\&op=show\&type=product\&code=032019\&layout=sermon}{\texttt{032019}} ~~~~ 語音日期: 2026-04-19
\newline
\newline
\hyperref[sec:032018]{\small{< < < PREV SERMON < < <}}
~
\hyperref[sec:index_chronic]{\small{[返順時目]}}
~
\hyperref[sec:032020]{\small{> > > NEXT SERMON > > >}}
\newline
\newline
民數記 10:1-36
\newline
\begin{longtable}{cl}
\hline
\hline
章節 & 經文 (和合本修訂版)\\
\hline
10:1 & \begin{tabularx}{0.7\textwidth}{X} 耶和華吩咐摩西說： \end{tabularx} \\ \\ \relax
10:2 & \begin{tabularx}{0.7\textwidth}{X} 「你要用銀子做兩枝號筒，把它們錘出來，給你用來召集會眾，拔營起行。 \end{tabularx} \\ \\ \relax
10:3 & \begin{tabularx}{0.7\textwidth}{X} 吹號的時候，全會眾要到你那裡，聚集在會幕的門口。 \end{tabularx} \\ \\ \relax
10:4 & \begin{tabularx}{0.7\textwidth}{X} 若只吹一枝，眾領袖，就是以色列部隊的官長，要到你那裡聚集。 \end{tabularx} \\ \\ \relax
10:5 & \begin{tabularx}{0.7\textwidth}{X} 你們大聲吹號的時候，東邊安營的要起行。 \end{tabularx} \\ \\ \relax
10:6 & \begin{tabularx}{0.7\textwidth}{X} 第二次大聲吹號的時候，南邊安營的要起行。起行的時候，要大聲吹號； \end{tabularx} \\ \\ \relax
10:7 & \begin{tabularx}{0.7\textwidth}{X} 但召集會眾的時候，你們要吹號，卻不要吹出大聲。 \end{tabularx} \\ \\ \relax
10:8 & \begin{tabularx}{0.7\textwidth}{X} 亞倫子孫作祭司的要吹這號筒，作為你們世世代代永遠的定例。 \end{tabularx} \\ \\ \relax
10:9 & \begin{tabularx}{0.7\textwidth}{X} 當你們在自己的土地上，與欺壓你們的敵人打仗時，要用號筒吹出大聲。你們就在耶和華－你們的神面前得蒙記念，也必蒙拯救脫離仇敵。 \end{tabularx} \\ \\ \relax
10:10 & \begin{tabularx}{0.7\textwidth}{X} 在快樂的日子，節期和初一，獻燔祭與平安祭的時候，你們要吹號筒，在你們的神面前作為紀念。我是耶和華－你們的神。」 \end{tabularx} \\ \\ \relax
10:11 & \begin{tabularx}{0.7\textwidth}{X} 第二年二月二十日，雲彩從法櫃的帳幕上升。 \end{tabularx} \\ \\ \relax
10:12 & \begin{tabularx}{0.7\textwidth}{X} 以色列人離開西奈的曠野，一段一段地往前行，雲彩停在巴蘭的曠野。 \end{tabularx} \\ \\ \relax
10:13 & \begin{tabularx}{0.7\textwidth}{X} 他們遵照耶和華藉摩西所指示的，初次往前行。 \end{tabularx} \\ \\ \relax
10:14 & \begin{tabularx}{0.7\textwidth}{X} 按照隊伍首先起行的是猶大營旗幟下的人，帶隊的是亞米拿達的兒子拿順。 \end{tabularx} \\ \\ \relax
10:15 & \begin{tabularx}{0.7\textwidth}{X} 以薩迦支派帶隊的是蘇押的兒子拿坦業。 \end{tabularx} \\ \\ \relax
10:16 & \begin{tabularx}{0.7\textwidth}{X} 西布倫支派帶隊的是希倫的兒子以利押。 \end{tabularx} \\ \\ \relax
10:17 & \begin{tabularx}{0.7\textwidth}{X} 帳幕拆卸了，革順的子孫和米拉利的子孫就抬著帳幕往前行。 \end{tabularx} \\ \\ \relax
10:18 & \begin{tabularx}{0.7\textwidth}{X} 按照隊伍往前行的是呂便營旗幟下的人，帶隊的是示丟珥的兒子以利蓿。 \end{tabularx} \\ \\ \relax
10:19 & \begin{tabularx}{0.7\textwidth}{X} 西緬支派帶隊的是蘇利沙代的兒子示路蔑。 \end{tabularx} \\ \\ \relax
10:20 & \begin{tabularx}{0.7\textwidth}{X} 迦得支派帶隊的是丟珥的兒子以利雅薩。 \end{tabularx} \\ \\ \relax
10:21 & \begin{tabularx}{0.7\textwidth}{X} 哥轄人抬著聖物往前行。他們未到以前，帳幕已經立好了。 \end{tabularx} \\ \\ \relax
10:22 & \begin{tabularx}{0.7\textwidth}{X} 按照隊伍往前行的是以法蓮營旗幟下的人，帶隊的是亞米忽的兒子以利沙瑪。 \end{tabularx} \\ \\ \relax
10:23 & \begin{tabularx}{0.7\textwidth}{X} 瑪拿西支派帶隊的是比大蓿的兒子迦瑪列。 \end{tabularx} \\ \\ \relax
10:24 & \begin{tabularx}{0.7\textwidth}{X} 便雅憫支派帶隊的是基多尼的兒子亞比但。 \end{tabularx} \\ \\ \relax
10:25 & \begin{tabularx}{0.7\textwidth}{X} 作全營後衛，按隊伍往前行的是但營旗幟下的人，帶隊的是亞米沙代的兒子亞希以謝。 \end{tabularx} \\ \\ \relax
10:26 & \begin{tabularx}{0.7\textwidth}{X} 亞設支派帶隊的是俄蘭的兒子帕結。 \end{tabularx} \\ \\ \relax
10:27 & \begin{tabularx}{0.7\textwidth}{X} 拿弗他利支派帶隊的是以南的兒子亞希拉。 \end{tabularx} \\ \\ \relax
10:28 & \begin{tabularx}{0.7\textwidth}{X} 以色列人就這樣按著隊伍往前行。 \end{tabularx} \\ \\ \relax
10:29 & \begin{tabularx}{0.7\textwidth}{X} 摩西對他岳父米甸人流珥的兒子何巴說：「我們要往前行，到耶和華所說的地方；他曾說：『我要將這地賜給你們。』現在請你和我們同去，我們必善待你，因為耶和華已經應許賜福氣給以色列人。」 \end{tabularx} \\ \\ \relax
10:30 & \begin{tabularx}{0.7\textwidth}{X} 何巴對他說：「我不去，我要回本地本族去。」 \end{tabularx} \\ \\ \relax
10:31 & \begin{tabularx}{0.7\textwidth}{X} 摩西說：「請你不要離開我們，因為你知道我們要在曠野安營，你可以當我們的眼目。 \end{tabularx} \\ \\ \relax
10:32 & \begin{tabularx}{0.7\textwidth}{X} 你若和我們同去，將來耶和華以甚麼福氣恩待我們，我們也必這樣善待你。」 \end{tabularx} \\ \\ \relax
10:33 & \begin{tabularx}{0.7\textwidth}{X} 以色列人離開耶和華的山，往前行了三天的路程。耶和華的約櫃在前面行了三天的路程，為他們尋找安歇的地方。 \end{tabularx} \\ \\ \relax
10:34 & \begin{tabularx}{0.7\textwidth}{X} 他們拔營往前行，日間有耶和華的雲彩在他們上面。 \end{tabularx} \\ \\ \relax
10:35 & \begin{tabularx}{0.7\textwidth}{X} 約櫃往前行的時候，摩西說：「耶和華啊，求你興起！願你的仇敵潰散！願恨你的人從你面前逃跑！」 \end{tabularx} \\ \\ \relax
10:36 & \begin{tabularx}{0.7\textwidth}{X} 約櫃停住的時候，他說：「以色列千萬人的耶和華啊，求你回來！」 \end{tabularx} \\ \\
[1ex]
\hline
\hline
\end{longtable}
$^{1}$今日講題是「喜行與安營」.
是文掃記第十章.
我不知道大家有沒有喜歡去日本旅行.
我們和家人都喜歡的.
上個月尾我適逢兒子放假去了幾天.
日本有些很有趣的文化.
我不知道你們有沒有體驗過.
我懂日文的不得幾句.
和大家一樣「阿諾」,「歐嗨喲」,「撒優那啦」,「阿利嘎多」.
就是這幾樣.
但每次去到日本.
常常我覺得日本人都覺得我們懂日文.
會講很長的東西.
我記得今次也是.
一去到其中一個酒店.
很好的接待,很有禮貌.
「皆さん,よこそ」.
我當然不知道他在說什麼.
我只是亂說.
他繼續說「今夜,よこそ」.
我就很有禮貌地扮明.
他當然知道我不明白.
接著他講到重點.
重點是「Welcome」.
突然聽到我明白的東西.
接著他再說一句「Dinner」.
我明白了,剛剛肚子餓.
接著他又繼續說日文.
很有趣「よこそ」,「のう」,「しんかん」.
我當然不知道他說什麼.
「Dinner time」.
「6.30」.
解決了我的問題.
很開心,我繼續看.
但他知道我給的反應不夠好.
他開始轉移看我兒子.
或者我兒子比較年輕.
他繼續在說的時候.
繼續說早餐.
「Breakfast」「8.30」.

$^{41}$結束之後.
我問我兒子.
其實你懂日文嗎?.
為什麼他開始看你而不看我?.
因為日本人的文化.
通常都是看男士和比較年長的.
我兒子說「不是的」.
「我只是扮懂而已」.
不過重點是我聽到.
「Keywords」.
「Breakfast」.
「Dinner」.
在我們人生路上.
可能很多東西我們都不明白.
我們解答不了.
不過我想告訴大家.
上帝都給很多重點的「Keywords」.
而上帝從來都不是沉默的.
上帝在我們人生中.
總會給我們很多關鍵字.
特別可能在我們迷失,沉溯的時候.
祂的提醒,祂的引導,祂的號聲.
今天去到文數第十章.
正正是講到上帝的記號.
在我們人生中.
當我們有聲音,環境有聲音.
究竟我們怎樣才能聽到神的聲音.
是我今天的重點.
而在今天的重點.
我們知道第十章講到一個很重要的原則.
就是當他們出了埃及.
在西乃山的地方經歷了年多.
他們領受了十屆.
他們有回望.
但是他們未進入應許之地.
就是說當我們離開遺牢之地.
不等於我們自動經歷豐盛的生命.
我再講一次.
在以色列文他離開遺牢的埃及地.
他不等於他自動進入了應許之地.

$^{81}$就好像今天我們信主一樣.
我們經歷了認識主的真實.
不等於我們自動進入豐盛的生命.
因為從蒙拯救到進入應許.
中間是有一段信心的旅程.
這個旅程我們要聽得清楚神的聲音.
並且要跟從.
所以今天我會用三點,三個方向.
有三個英文字母.
我們一起學習.
但願我們學習在這個經文裡面.
帶領我們怎樣走我們人生的路.
怎樣走我們信仰的路.
怎樣跟緊跟隨主.
我們先讀民數記頭一到十節.
我們分段讀.
我們先讀四節.
預備起.
耶和華曉如摩西說.
你要用銀紙做兩支號.
都要除出來的.
( 音樂 ).
就是以色列軍中統領.
要聚集到你們那裡.
一開始我們已經看到.
銀紙做兩支號不同.
我們已經看到.
單吹一次.
我們慢慢會看到.
不知單吹大聲,細聲.
吹完一次又一次.
連續吹.
我們先看完.
繼續看下去.
第五到第八節.
預備起.
吹出大聲的時候.
東邊.
( 音樂 ).
你會看到.

$^{121}$對不起.
繼續.
( 音樂 ).
這要祝你們世世代代永遠的平安.
太刺激.
對不起.
又會有大聲,細聲兩次.
不要吹大聲.
召聚會眾.
有不同的作用.
我們再看第九個十節.
預備起.
你們在自己的地.
與欺壓你們的敵人打仗.
( 音樂 ).
並月數.
( 音樂 ).
也要吹號.
( 音樂 ).
我是耶和華.
你們的神.
第一件我們要學的.
不是懷疑神會不會出聲.
我要學的是.
怎樣分辨神的聲音.
神一定會出聲.
一定有.
第一個英文就是.
Design.
中文是分辨.
剛才我們說到有不同的聲音.
大聲,細聲.
一次長的.
或者斷軸的.
( 音樂 ).
又或者不同的時候.
打仗的時候.
聚集的時候.
又或者是獻煩祭感恩的時候.
而在本身原文.

$^{161}$它都有兩個不同的字.
例如.
吹出大聲的用Tagua.
用吹號.
它用Hazora.
我們不用理會原文背後.
不過我想告訴大家.
是有不同的字.
也代表我們人生裡面.
在行的時刻.
我們需要分辨那些聲音.
非常重要.
而我們現在往往在日常生活.
都非常繁忙.
無論工作的聲音.
坐車的聲音.
別人的聲音.
世界文化的聲音.
如果我們一旦分辨不到.
我們內裡的聲音.
和上帝的聲音.
我們很容易就混亂.
我記得.
我不知道大家有沒有提款機提款.
我記得早一段時間.
我又做這個動作.
現在都很少.
因為我們用電子貨幣.
我也是.
但都需要有少少現金傍身.
但我家中的太太其實是提款機.
我沒有錢就用錢包拿.
我有錢會放進去.
不過我沒有什麼時候有錢.
又去提款.
提款的時候.
那次有兩個提款機.
我正在排一條.
在另一條.
又有人在排隊.

$^{201}$我見到有一個很忙的人.
拿著電話.
又在發WhatsApp排隊.
突然間.
我排這邊的時候.
聽到那邊的提款機.
有聲音發出.
是嘟嘟嘟嘟嘟.
大家知不知道什麼事.
(拿錢).
是呀,沒有拿錢.
當時的情況是這樣的.
我聽到嘟嘟嘟嘟.
我立刻望過去.
沒有人在提款機.
排隊後面的人就立刻向後退.
很害怕.
在我這條隊立刻向左動.
但大家知不知道為什麼.
你不能碰.
你碰到拿了錢就要坐牢.
有CCTV的.
有新聞說過.
曾經有人提款.
沒有拿錢.
後面那個人拿了錢.
結果被人抓了.
不過我想說重點就是.
應該聽到聲音的那個.
沒有聽到就走了.
不需要聽那些人.
就在聽.
並且反應是.
害怕.
有時我們都會這樣.
本來應該聽到上帝的聲音的時刻.
聽不到.
走了.
旁邊的人就知道的時候.
很緊張.

$^{241}$並且都找不到.
不單止這樣.
我們剛才聽到一些細節.
就是東邊.
南邊聽到的起行.
是代表什麼呢.
我們回一點重溫.
在文宿記第一章.
我們知道.
應該大家都已經看過了.
重溫一點.
就是當時.
那個描述的.
會幕和四圍人的安寧.
我們知道當時是有.
多少萬的議事列民.
二百多萬.
男士六十多萬.
接著有女士和小孩.
估計二百多萬.
二百多萬人如果吹號.
起來.
嘩.
很噓憾.
大混亂.
不過上帝不是這樣.
上帝是叫人井井有條.
是這樣的.
我們看看.
當時第二章.
他這樣描述.
會幕在這個位置.
約跪這個位置.
接著.
我們不用讀.
我們看看黃色的字.
東邊.
日出的地方.
接著七萬人.
五萬人.

$^{281}$三個支派.
總共十八萬人.
就在東邊.
可以嗎.
OK.
所以剛才說.
一吹東邊就起營了.
可以嗎.
不會一齊.
很混亂.
慢慢.
很清楚了.
接著.
南邊.
第二章.
看到了.
右邊的營.
四萬人.
西面五萬人.
加德四萬人.
可以嗎.
接著加起來.
特別號碼.
十五萬人.
接著南邊.
可以嗎.
看到了.
OK.
當然有東南還有什麼.
西.
沒錯.
西見到.
西邊.
二百年.
四萬.
三萬.
東南西之後.
是什麼.
北.
沒錯.

$^{321}$我們見到.
北邊.
蛋形.
六萬.
阿徹.
四萬.
等等.
接著加起來.
十五萬.
好了.
我們見到.
原來.
是有佈局的.
二百多萬人.
井井有條.
加起來.
六十多萬男士.
在這個情況.
我們知道.
並且知道.
在整個以色列文的中心.
不是摩西.
是上帝的回眸.
我們今天.
可能我們人生裡面.
都在混亂的狀態.
我們值得問我們的中心.
是不是上帝的回眸同在.
是不是上帝的恩典和救恩.
我們聽不聽到神的聲音.
所以怪不得我們明白.
原來回眸正正在中間.
當我們吹號的時刻.
大家都要聽到並且回應.
上帝的吹號.
不是單是要起行.
也會安然.
我記得當我在職場工作的時候.
有時都會有獵頭公司.
請我去不同的公司工作.

$^{361}$有一次有一個獵頭公司.
請我去做.
Buy side.
當時我都覺得Buy side.
因為我做Sell side.
我做金融.
Buy side我都很喜歡.
基金.
但當時同事朋友都說好呀去呀.
但那一刻.
在我安靜靈修的時刻.
神話安然.
不要聽別人說.
在那個時候.
還有福音工作.
每個人處境不同.
但我們要聽得清.
不是我們裡面想聽的東西.
或者我們裡面的野心.
或者我們裡面想要的人的榮耀.
乃至要安靜靈聽神.
我又說一下.
我剛剛去了日本玩.
剛才說了.
和我兒子去日本玩.
是很辛苦的.
我兒子25歲.
我55歲.
和一個25歲的人去玩.
他安排的東西.
其實應該是25歲的人去的.
這五天我就和兒子說.
所有的安排.
吃的東西你安排.
去看的櫻花你安排.
然後結果.
他就這樣安排.
我有一個例子.
你明白.
有一晚我們吃完晚飯.

$^{401}$我打算回家休息.
都一百年多.
他說爸爸還沒住.
日本有一間是全亞洲Top50的Cocktail Bar.
很便宜的.
Cocktail Bar.
55歲.
他說爸爸去了.
當然去了.
你不去.
下次他不和你去玩.
好了.
十點才開始.
喝到十二點.
談一下工作.
談一下人生.
笑得很開心.
然後走了.
回到酒店.
洗完澡都一點多.
然後打算睡一覺.
明天起床.
起床就起.
起不了就不起.
然後我兒子就說.
爸爸.
明天我已經預約了.
七點半.
日本早餐.
我們六點起床.
我看一看.
只有四個小時睡.
我再看一看他.
取消不了的.
給了訂金.
當然去了.
我打算去完早餐.
回到酒店睡一會.
他說不是.
然後我們就看櫻花.

$^{441}$開車.
一路玩到晚上.
再吃一個Tempura.
很便宜的Tempura.
好像吧檯一樣.
一格一格坐.
兩百元港幣一個人.
吃了很好的東西.
不過那一刻.
有些特別的事情發生.
當我疲於奔波.
當我疲於奔命.
在想何時睡覺的時候.
我們在排隊.
就看到那個吧檯.
有一個露宿者.
在吃東西.
一代二代.
我兒子說.
他的腳在流血.
估計他是糖尿病.
我兒子是見習醫生.
他的位置旁邊.
是沒有人坐的.
日本人都不坐.
他都不叫日本人坐.
我們三個人排隊.
我和我太太是日本人.
跟著空位旁邊兩個人站起來.
他就說.
你三個人去吧.
我覺得都幾有趣.
日本人他不叫去.
香港人叫我去.
不過很感恩.
開謝我那一頓飯.
我專做露宿者.
如果大家知道.
雅庚堂當日出來.
就是聽到整個聲音.

$^{481}$去深水Po .
當時我望一望我兒子.
望一望我太太.
誰坐啊.
我兒子和我太太說.
當然是你啦.
當時我坐在那裡.
上帝和我說.
我都有少少覺得辛苦.
因為很擠迫.
他的位置.
他的身體不停撞到我.
並且公共一起用浸水的地方.
他用了食物的碗再浸水.
可以嗎.
但神很清楚和我說.
安寧.
我就想起馬太福音第25章.
耶穌講得很清楚.
因為他名的緣故.
服侍不能穿的.
不能吃的.
就等於服侍誰.
耶穌.
我馬上又想起.
你見到嗎.
讀墓者很煩的.
經常想起這些.
想起腦子.
我又想起德蘭修女.
我自己真的很尊重.
很仰慕她.
她講一件事.
就是當在印度的時候.
那些沒有人要的人.
那些人都丟了.
印度人都不要的.
每一次去到她那裡.
她望見他.
她就和主講.

$^{521}$多謝你給我服侍耶穌的機會.
在那一刻.
我本來想幫他付錢.
給了錢才坐下.
那些餐廳.
但是我可以坐在他旁邊.
是我的榮幸.
我竟然可以在日本.
坐在耶穌旁邊.
現在都有點觸動.
他們不要.
我們要.
弟兄姊妹.
我不知道你最近聽到的聲音是什麼.
是你焦慮的時候.
是你盼望的時候.
還是神有事提醒.
我們不如刻意的.
將我們的手提電話.
或者我們一切的響鬧裝置.
不是習慣去震聲.
是關掉它.
15分鐘.
30分鐘.
當我們的內心的焦慮.
內心的野心.
別人對我們的評價.
內心的慾望.
前來的焦慮擺下.
我們就單單打開聖經.
或者回想今天.
我們崇拜聽到.
Amazing Grace.
聽到這段聖經第十章.
我們靜下來.
看看神有什麼引導.
有什麼聲音.
又或者當我們在裡面.
感恩的時刻.
在節期的時候.

$^{561}$我們帶著喜樂.
又或者我們有需要的時候.
我們就禱告.
讓上帝拯救.
第一點.
分辨神的聲音.
諧相.
第二點.
接著慢慢去到更加深入.
去到第十一到十三節.
我們一起讀好不好.
1,2,3.
第二年2月20日.
.
第二個重點是什麼呢.
第二年2月20日的意思.
就是他離開了圍爐之地.
大概一年多.
他在西乃山下面.
其實二百多萬人已經.
一年多應該都很舒服了.
那裡領受了十誡.
那裡知道回膜的運作.
在那一刻.
當上帝呼喚他們起行.
我們就要起行.
第二點.
Follow.
有時候我們已經在一個安書區.
那個地方.
大家不要搞錯.
有人以為在曠野是一個.
受新奴的地方.
No.
離開埃及去曠野是敬拜主的地方.
當時上帝呼召摩西.
跟他們人民說.
我們去到曠野是敬拜上帝.
那裡不是受難的地方.
那裡是離開圍爐得自由的地方.

$^{601}$在那裡.
可以這樣說.
是一個短暫經歷平靜的時刻.
我可以稱為.
那個等於西乃山的時刻.
但是大家要謹記.
那個時刻不是我們的終點.
那個時刻是帶著我們去終點.
但是重點問題就是.
很多時候.
我們去到那個中轉站.
就以為是終點.
以色列民正正發生這件事.
我們知道的.
本來幾十天應該去到應許之地.
結果在曠野這個中轉站就停留了多久.
四十年.
搞錯了.
雲彩移動的時候.
你就要移動.
意味著信心之旅的開始.
西乃山不只是一個舒適圈.
Comfort Zone.
當然我們知道那裡經歷了.
律法經歷了和神立約的時刻.
但是那個不是終點.
換言之我們要離開那個.
多好的地方.
才可以去到更好的地方.
很可惜.
太多的基督徒.
去不到那個真正豐盛之地.
剛才我講到露宿者.
有些原因.
因為當日我本來在你們的子堂.
我的舞會和福堂經營了二十年.
像不像呢.
差不多.
很舒服.
我在那裡結婚.

$^{641}$做執事.
做組長.
大敬拜.
讀神學.
但是在那裡服事的時刻.
神在我裡面有一個不滿足的呼召.
有一次我在禱告裡面.
神就在禱告裡跟我說.
我要呼召你去到全香港最窮的地區.
當時我只是想在哪裡.
難道是天水圍找.
不是.
原來是深水Po .
我們開始2011年.
第一次和有家無家一家親.
和露宿者去服事.
當時很害怕.
很不舒服.
第一次擁抱露宿者.
第一次看到他插著針.
在暈的時候.
我心想應該怎麼做.
我就馬上問當時的關懷無家者協會.
他很簡單.
你不要搞他.
他打完針.
他就會跟你聊天.
你不要嚇到他突然打歪了.
這樣就行了.
當然其實很多都是很困難的時候.
他才要去到露宿.
我記得當時去到服事的時刻.
和幾個人.
10個20個去到.
最後我們很多人在那裡經歷.
很記得我們開始分享的時候.
第一次去到聚會.
第一次去到外展.
第一次去到嘉年華.
就總共有300多人信主.

$^{681}$街坊露宿者.
第一次我很興奮.
和你們一樣.
當時都不習慣.
都是中產背景.
但是上帝的心意.
比我自己裡面的呼聲更大.
很奇妙.
我們就這樣去服事.
很想用這個領袖去說一句話.
我覺得很好.
Good times blind us.
Bad times reveal truth to us.
意思是什麼呢.
很美好的時候.
有時候我們見不到更美好的東西.
我初時去到都很害怕.
但是我發現原來上帝可以打破我自己服事的範疇.
不單止是某一個方向.
是可以更廣泛.
不單止是這樣.
但已經是一樣.
艱難的時候揭示真相.
原來上主可以有更廣闊.
我們服事在我們面前.
我知道港福堂.
KFTC在香港仔.
也都有服事社區的中心.
我自己很感動.
當林牧師去分享的時候.
當孫牧師去講到不同的團隊.
我覺得很感恩.
教會就是去服事不同的群體.
一個都不能少.
那裡有很多耶穌.
所以我都鼓勵.
當你聽到裡面有一種感動.
你都可以去回應.
可能不是完全去.
一個月去一次.

$^{721}$兩個月一次.
去探訪.
一起崇拜.
一起服事.
我記得當我女兒在COVID去讀大學的時候.
我女兒已經畢業了.
現在在大學做研究.
當時當然是在香港讀.
在美國讀.
在英國讀.
當時我們一直禱告.
但我們知道上帝的啟示揭示在美國.
但在COVID的時刻.
我最擔心作為牧者.
就是他能否找到教會.
當時已經關閉了.
香港關閉了.
美國關閉了.
當時美國就快點開.
我不知道怎樣可以有渠道.
可以在他附近找到教會.
但很奇妙.
我和林牧師在總會都有服事.
我其中一個地方.
我們都要開會.
要有些事情要處理.
有些外展等等.
支票我都是簽.
我記得林信我女兒去外國讀書.
前三個星期.
我都完全無頭緒.
去到怎樣呢.
可以怎樣去做.
但很奇怪.
那天去到總會.
就是要簽支票.
我不知道你知不知道我.
不過你未必認識我.
我都很喜歡漏人經濟.
你好嗎.

$^{761}$我經常都這樣.
我就漏那個會計經理.
我問你為何來到總會.
他說是呀.
都很奇怪.
我本來在美國工作.
美國馬上叮叮叮叮叮叮叮.
美國.
跟住我就問.
你在美國哪裡.
美國西岸.
西岸.
叮叮叮叮.
我女兒去西岸.
神的吹號響起.
跟住我再問西岸哪裡.
他說灣區.
灣區.
叮叮叮叮叮叮.
跟住我再問.
灣區哪裡.
帕克萊.
帕克萊就是我女兒那裡.
叮叮叮叮叮.
我就問他.
你那裡有沒有好的教會.
他說當然有.
他說他的小侄女.
就是回那裡教會.
我說你小侄女在哪裡.
在哪裡.
就是帕克萊.
他剛剛在那間大學.
剛剛那年畢業.
於是我飛去.
介紹了.
結果他的小侄女.
就成為我女兒的mentor.
安心了.
最安心其中一樣東西.

$^{801}$就是當我女兒拍拖的時候.
作為Daddy.
你知道Daddy.
Daddy見到女兒和兒子不同.
我不知道你是否這樣.
我見到兒子很豪氣.
我兒子以前在家裡.
和女兒讀書.
我回家.
我兒子通常是這樣.
抬高頭.
我回家.
Daddy.
繼續垂頭.
在那裡工作.
我女兒就不同了.
我回家.
我女兒就說.
Daddy.
走去門口抱著.
你明不明.
我還說.
嘩.
糟了錄了音.
給他們聽.
的確擔心一點.
但很有趣.
我女兒拍拖.
我們分享禱告.
那個mentor跟她說.
你知不知道什麼是拍拖.
你知不知道怎樣畫界.
作為一個教徒.
她是美國人.
我都嘩.
真是跟隨主的門徒.
你知不知道.
拍拖怎樣跟隨主.
嘩.
我自己聽的時候.

$^{841}$非常感動.
各位電視姊妹.
無論是我們自己的生命.
無論我們的服事.
我們都要跟隨主的聲音.
而過程必須有這幾樣東西發生.
第一.
從我們的comfort zone.
我們的安書區.
開始.
我們會跳出去.
第一個層面.
我們跳出去.
就會去到fear zone.
什麼是comfort zone.
我們覺得我們搞掂.
I got it.
沒問題.
不用怕.
萬事是我控制底下.
但是那裡就只是那裡.
但是去到下一步.
你發現.
我不可以自己做得到.
我可能會失控.
我為你感恩.
你進入另一個zone.
叫做fear zone.
在那裡.
你開始知道自己的能力有限.
你會尋求上主.
亦尋求別人的禱告.
帶領和幫助.
在那一刻.
你開始就進入第三個zone.
叫做reliance zone.
就好像我和我女兒一樣.
主啊.
你何時吹號.
主啊 你何時準備.

$^{881}$在那一刻.
你知道原來你所有的資源.
就算exhausted.
都幫不了你.
我為你感恩.
因為fear zone.
不是就這樣拆毀我們.
它要reshape我們.
重塑我們.
讓我們的範圍更大更廣.
讓我們知道.
我們倚靠的不是自己.
我們倚靠的是上主.
接著我們要進入learning zone.
再次學習和上帝的相遇.
他怎樣帶領.
接著我們進入growth zone.
我們真的成長.
西乃山是一個暫時逗留的地方.
我們要聽得清楚.
不要急著走.
亦都不要留戀那裡.
今天在第二點.
我們要follow.
要反思你生命中.
有沒有一個西乃山的時刻.
那個可能是祝福你是真的.
但是那個可能是要你離開的地方.
我不知道你有些什麼.
你離開不了.
又或者你和人際的關係裡面.
有些什麼張力.
你要打破.
又或者你服事裡面.
上帝呼召你.
有些不同的範疇.
你去體驗.
又或者你工作裡面.
上帝有些不滿足.
告訴你.

$^{921}$起來了.
若你的魂采移動.
求你賜我信心勇氣.
行出第一步.
並且知道第一步是什麼.
學習二.
follow.
第三.
第29到32節.
我們一起讀.
123.
摩西對他岳父米顛人留爾的兒子阿爸說.
學習三.
你發現一個很有趣的事.
就是摩西這個屬靈的偉人.
他需要人幫.
很多時候我們有個誤解.
覺得很屬靈的人可以單打獨鬥.
不是不是.
他需要人.
那你讀下去.
你繼續看完31,32節.
你更加明白.
為什麼他需要什麼人.
我們讀好嗎.
你不要離開我們.
因為他需要人作他眼目.
意思是什麼.
摩西看到雲住火住.
不過在他很多細則.
在這個曠野.
哪裡有水.
哪裡危險.
他需要米顛人.
一直在曠野住.
帶著他.
所以當何爸說不去.
他說不要啊.
我要你去.
向他表達你的需要.

$^{961}$也都在那個過程.
更加重要.
就是你明白.
你要學習謙卑.
尋求同路人.
第三個英文字是seek.
seek help.
尋求同路人.
剛才我說.
我們不要覺得自己可以搞定.
我們不要覺得.
我自己可以掌握一切.
不是不是.
也是上帝的心意.
和上帝的設計.
是我們需要團隊.
需要隊工.
我們需要mentor.
我們需要信仰群體.
我們需要牧者.
需要組長.
需要大家.
當我小背景解釋.
這個是不是摩西的岳父葉蒂羅.
解英家有不同的看法.
可能是葉蒂羅的兒子.
不過無論如何.
都是說一件事.
他需要一個很熟悉地形的人幫他.
而在那個時刻.
他清楚就要表達.
我不知道你人生有沒有這一刻.
或者曾經需要人幫助.
但你沒有去尋索.
又或者可能你是別人的幫助.
你是別人的河壩.
在我們一生中.
我們深信出現過很多個河壩.
但我們要緊要的.
我們要把握.

$^{1001}$我們不要找一些合聽的人.
這個說話合聽.
這個最喜歡我.
我說的話他都說是.
這個不是你的河壩.
你要找一些他有熟齡的眼光.
有不同的智慧.
可以為你分辨.
如果不是我們只是找人去.
endorse我們自己的看法.
這段聖經整個說就是.
我們是跟著上帝的號聲.
走神的看法.
我每次去不同的地方講道.
我都會組幾個字去.
一方面我會尋思那段經文.
同時我們會為那些會眾禱告.
上帝有什麼揭示.
讓我有些不同的重點.
也都為大家禱告.
昨天也是.
第一次.
昨天晚堂來到廣福的時候.
我們組了四五個字.
我就在月台.
在那裡尋索.
走圈禱告.
大家如果有時見到我.
在你附近在走路.
我不是迷路.
你來打個招呼都可以的.
當時我剛剛下車.
在月台.
一兩分鐘.
很快就發生了一些很有趣的事.
我相信.
別人一定不會知道我是牧師.
不過我覺得.
別人應該覺得我是一個好人.
好像你不是很覺得.

$^{1041}$你笑.
為什麼呢.
因為三分鐘左右.
有六個人問路.
第一個問路是一個西人.
他一來到說.
是中央嗎.
我說不是.
這是廣東話.
他說不是.
有人告訴我是中央.
有人告訴我.
不是不是.
是廣東話.
中央是下一個站.
去那裡坐車.
然後他很懷疑.
不是不是.
是中央.
我說不是.
這是廣東話.
OK.
好 我去吧.
然後很快.
我想有十秒左右.
就有一個說廣東話的.
先生.
我想問你.
哪裡有廁所.
我知不知道有廁所.
但我真的知道.
東節線的尾友.
東節線那條線走到篤.
我說你去到東節線那裡.
去到篤那裡有廁所.
他說真是好了.
原來有廁所.
然後我想著一祈禱.
接著有一個說國語來.
先生想問你.

$^{1081}$.
怎麼去.
我要去香港大學.
瘋的.
因為那時四點多.
黃昏.
小紅書館那裡拍照很美.
我就不太懂.
我說你坐地鐵去到中環再轉.
接著我說不是.
這裡說金鐘轉.
是喔.
.
去堅尼地城線.
是喔.
接著我想著靜.
又有一個西人來.
Where is Central.
這個是女士.
說完之後.
接著又有兩個香港人夫婦.
廁所在哪裡.
我綜合了.
我真的綜合了.
香港人喜歡去廁所.
西人就找中環.
國內來就去香港大學拍照.
五分鐘reserve.
不知道準不準.
但我想說.
我們真的有需要問路.
我們人生很多時候.
我們真的要找人問路.
我自己也是.
在我的屬靈上.
我有很多.
在神學院有老師.
在美國或新加坡有不同的老師.
他在我人生要考慮的時刻.
他未必很準繩.

$^{1121}$但他給了很多的看法.
我記得.
當我神學院快畢業.
我快出來工作的時候.
我二哥癌症末期.
發現一年就走.
當發現癌症的時候.
他小兒子還在肚子裡.
第三個掣.
當他離開.
他小兒子還未一歲.
我當時媽媽跟我說.
你不要做傳道了.
出來工作.
當時我覺得很為難.
她說你哥有三個.
你負責照顧第一個.
我姐一個.
跟著嫂一個.
我當然想.
跟著我找神學院楊醫.
我跟他說楊醫.
很困難.
楊醫跟我說.
我知道你很困難.
但你怎樣走.
你內心知道的.
跟著我想應該有第二句.
是沒有的.
說完.
其實我真的知道.
我很想找個人告訴我.
不如出來工作.
做幾年.
但我內心知道.
神的號聲並不是如此.
結果我也有照顧第三個兒子.
我也有繼續做傳道.
我想問問第三.
在你人生的曠野路上.

$^{1161}$神有沒有為你準備荷包.
一個可以成為你眼目.
幫你的同路人.
這個星期.
如果你真的有需要.
把握機會去找荷包.
表達你內心的迷惘和曠野.
或者這個星期.
可能有人找你幫忙.
可能他也迷失.
你也豁達出去.
成為和你一起走一個曠野路.
一個荷包.
去到最後.
也是一個高峰.
總結.
這三節經文很有趣.
我們看看33到36節.
我們一起讀好不好.
1,2,3.
以色列人離開耶和華的山.
為他們尋找安潔的地方.
走了三天.
跟著就安潔.
有人覺得跟隨主.
做死嗎?.
錯.
你對上帝的看法是錯誤的.
跟隨主.
永遠坐在那裡安安.
天天有平安.
錯.
跟隨主.
又會前行.
又會有安潔.
跟隨主.
還有什麼在我們前面.
藥櫃.
我們看不到.
看不到怎麼走.

$^{1201}$不要緊,我看到藥櫃就行了.
在那裡.
我們就知道.
藥櫃在前行的時候.
摩西就說.
耶和華求你興起.
願你的仇敵四散.
願恨你的人.
從你面前逃跑.
藥櫃停住的時候.
他們就說.
耶和華求你回到以色列.
千萬人中.
說什麼呢.
就算我們正在走主的路.
同樣會有仇敵.
哎呀,跟隨主.
無痛分娩的.
錯.
跟隨主什麼事都沒有錯.
你信錯了.
跟隨主從來都不是天色上藍.
不過就在那一刻.
我們看到.
耶穌就在那裡與我們相遇.
帶領我們.
所以我們要分辨好.
當我們在跟隨主的過程.
讓我們知道跟隨主都會有艱辛.
不過藥櫃就在我們前面.
不過我們要小心.
不是所有的困難.
都是因為跟隨主.
很多就是因為你沒有跟從主.
那裡你說祈禱主你幫我.
祈禱主你幫我.
神不知道怎麼幫你.
因為本身那條路.
是與神的路分離.
在以色列民的處境你知道.

$^{1241}$這四十年就是.
很多時候他不是走神的路.
神都不知道怎麼幫.
神很想你跟隨他.
他幫助你.
今天.
在最後可以提一提.
我們最重要的是要堅定相信.
我們的態度就是.
以神為先.
以藥櫃在我們前面.
當我們審閱一下.
我們人生重要的決定和計劃.
你目前.
是否確認走在主的路.
還是你走在神的前面.
叫神祝福你.
這不可行.
當我們走在主前.
緊緊跟隨他.
將我們的主權交給神.
我們必然會見到.
第一.
分辨神的聲音.
第二.
跟隨他的聲音.
第三.
seek help.
DFS.
Duty Free Shop.
不是.
都可以這樣想.
Duty Free Shop 裡面有很豐富的東西.
但我們在神的Duty Free Shop.
更豐富.
帶領我們進入.
豐盛的生命.
分辨.
跟從.
謙卑.

$^{1281}$讓我們一起祈禱.
天母我們謝謝你.
你真的這樣帶領我們.
那個的聲音.
從來都不間斷.
讓我們在混亂中.
聽到你的指引.
那個的魂.
那個魂處.
火處在我們人生裡.
讓我們清晰.
見到.
讓我們不會亂衝亂撞.
忽視了.
你的.
號聲.
讓我們不要留戀在.
安書的西乃山.
也不是留戀在那個時刻.
讓我們願意.
緊緊跟隨.
你的雲彩前行.
讓我們願意.
將生命的主權交給你.
讓我們願意.
見到你的約會.
成為我們家庭.
事業.
事奉的中心.
我們謝謝你.
聽我們祈禱.
\newpage



\section{約翰福音 21:20-23}
\label{sec:032020}
\textbf{你跟從我吧}~~~~黃健羒牧師
\newline
\newline
連結: \href{http://www.church.com.hk/acms/content.asp?site=cdc\&op=show\&type=product\&code=032020\&layout=sermon}{\texttt{032020}} ~~~~ 語音日期: 2026-04-19
\newline
\newline
\hyperref[sec:032019]{\small{< < < PREV SERMON < < <}}
~
\hyperref[sec:index_chronic]{\small{[返順時目]}}
~
\hyperref[sec:032021]{\small{> > > NEXT SERMON > > >}}
\newline
\newline
約翰福音 21:20-23
\newline
\begin{longtable}{cl}
\hline
\hline
章節 & 經文 (和合本修訂版)\\
\hline
21:20 & \begin{tabularx}{0.7\textwidth}{X} 彼得轉過身來，看見耶穌所愛的那門徒跟著，就是在晚餐時靠著耶穌胸膛說「主啊，出賣你的是誰」的那門徒。 \end{tabularx} \\ \\ \relax
21:21 & \begin{tabularx}{0.7\textwidth}{X} 彼得看見他，就問耶穌：「主啊，這個人怎樣呢？」 \end{tabularx} \\ \\ \relax
21:22 & \begin{tabularx}{0.7\textwidth}{X} 耶穌對他說：「假如我要他等到我來的時候還在，跟你有甚麼關係呢？你跟從我吧！」 \end{tabularx} \\ \\ \relax
21:23 & \begin{tabularx}{0.7\textwidth}{X} 於是這話在弟兄中間流傳，說那門徒不死。其實，耶穌不是說他不死，而是對彼得說：「假如我要他等到我來的時候還在，跟你有甚麼關係呢？」 \end{tabularx} \\ \\
[1ex]
\hline
\hline
\end{longtable}
$^{1}$首先會問大家兩個問題.
請你舉手 不用太高.
舉一舉就好.
如果你的手不方便舉起.
伸手指都可以.
在座當中誰是基督徒?.
誰是耶穌的門徒?.
OK 謝謝.
如果不是舉兩次手.
一定是對門徒這個字有些疑問.
是什麼意思呢?.
因為通常我們說.
耶穌的門徒可能指那時候.
才是跟耶穌的人.
那我們今天.
我們是不是耶穌的門徒呢?.
其實我們就比較少用這個名詞.
如果你說看士徒行傳第十一章二十六節.
這裡是很清楚的說到.
門徒稱為基督徒.
是從安提阿開始.
安提阿是今天在敘利亞的一個城市.
那時候除了耶路撒冷.
安提阿的基督徒是最多的.
這裡說是最先開始.
那些跟耶穌或信耶穌的人.
是叫做門徒的.
在那個地方開始.
才改為叫做基督徒.
其實門徒就是一群.
跟著老師學東西的學生.
或者徒弟跟著師父學東西.
那些叫做門徒.
耶穌的門徒就是基督徒.
基督徒就是耶穌的門徒.
雖然我們不是用門徒這個字.
沒有使用這個字來自稱.
但是還是用門徒訓練這個字.
教會很多時候都說.
去跟從耶穌 參加門徒訓練班.

$^{41}$其實就是學習怎樣去做耶穌的門徒.
學習怎樣去跟從耶穌.
所以根據福音書的記載.
耶穌復活之後.
有四十天的時間.
向門徒顯現.
今天離開復活節幾多天?.
我想只有兩個星期.
好像時間過得很快.
不夠四十天.
所以我們可以想像.
耶穌在這段時間裡.
在當時用不同的機會.
向他的門徒顯現.
有時是幾個人.
有時是一群人.
哥林多前書第十五章.
保羅特別提到有一次.
是向五百多人的聚會當中顯現.
其中有一大半在保羅寫信的時候.
還是在生的.
所以耶穌從死裡復活是一個事實.
如果是讀哥林多前書的時候.
有些人不相信的話.
你就找那些還在生的來問問就知道了.
所以那段時間沒有人問這個問題.
沒有人懷疑這個問題.
因為耶穌復活是一個事實.
所以今天我們就其中看一下.
耶穌的門徒就是跟從耶穌的人.
相信他是基督的人.
所以今天如果相信.
耶穌是基督的人.
我們基督徒也有耶穌的門徒.
在約翰福音剛才我們讀的那段經文.
特別提到.
耶穌是第三次向他的門徒顯現.
希望通過耶穌和西門彼得的對談當中.
我們多了解.
跟從耶穌究竟是什麼一回事.

$^{81}$首先耶穌呼召彼得說.
你跟從我吧.
其實是願意付代價的跟從.
耶穌對彼得說.
你年輕的時候自己要速上大子.
隨意往來.
但年老的時候你要伸出手來.
別人要把你速上.
帶你去不願意去的地方.
耶穌說這句話是指彼得.
會怎樣死來榮耀神.
然後耶穌對彼得說.
你跟從我吧.
換作你和我可能也要想一想.
年老的時候.
退休.
所謂享清福.
享兒女福.
彼得有沒有兒孫.
聖經沒有記載.
但他肯定是已經結婚.
大家記得在福音書有記載.
耶穌有一次去他家.
醫治他的岳母.
他結了婚.
但是有沒有兒女呢.
不知道.
不過通常來說我們猜測有.
根據教會歷史的記載.
彼得殉道的時候.
大概是63至67歲左右.
極有可能有兒孫.
年老了.
退休了.
享受一下兒孫福.
多好.
但耶穌的預言是.
彼得要為他的名義的緣故被處死.
彼得沒有猶豫.
他就跟從耶穌.

$^{121}$教會的歷史有傳說.
彼得被羅馬皇帝尼祿.
判處死.
為了信仰的緣故.
他被釘十字架的時候.
他自己要求.
我不配像我的主耶穌一樣.
這樣釘.
如果你要釘我十字架.
就將我倒轉.
頭下腳上這樣釘十字架.
這是教會傳統的一個事蹟.
所以如果大家看以前的畫.
一看到倒轉釘的那個.
就是彼得.
知道自己年老的時候都要殉道.
還要跟從主?.
真正的廣東話就是說.
是呀 我要跟 死跟.
跟定了.
沒有其他原因.
為什麼他會這樣?.
因為他深深感覺到.
耶穌是為了他而死.
這個愛感動他.
我也願意為主而死.
每一個感到.
耶穌是為你而死的人.
都願意付出代價來跟從主.
跟從主 得到主.
那個才是生命.
過去的星期一.
我在西南區團契分享了.
我提到這一節的經文.
是《老教福音》第九章二十三節.
《馬太》《馬可》福音都有記載.
第二十三節.
耶穌對眾人說.
若有人要跟從我.
就當捨己.

$^{161}$天天背起自己的十字架來跟從我.
因為凡要救自己生命的.
必喪失生命.
凡為我喪失生命的.
他必救自己的生命.
雖然二十四節沒有很清楚.
充分明白究竟生命是什麼.
但起碼我們知道.
耶穌一定是說兩種不同的生命.
我用不同的顏色來表明.
一個是黃色的.
一個是綠色的.
其實是比較清楚一點.
救自己的生命.
就要失去生命.
凡是為我失去生命的.
就必定救我的生命.
當然是說不同的生命.
有人問會否是肉體的生命.
我們現在可以行動.
可以說這個生命是否死了.
就得到另外的生命.
有些人會這樣想.
我現在死了.
離開這個世界.
就得到永生.
永恆的生命.
兩段好像不同的生命.
是否這樣說呢.
但如果我們看二十三節.
天天背起十字架.
天天跟從.
不可能天天死的.
當然不是這個意思.
其實背起十字架.
是作為一個犯人.
走向刑場的路.
是走向死亡.
當時被判刑的犯人.
根據考究.

$^{201}$只是背起橫的大木.
已經夠重了.
如果要天天背起十字架.
天天寫記.
天天走向死亡.
就不只是身體幾十年的生命.
而是講到心靈的問題.
講到生命方式的問題.
如果你愛一個人.
但他不愛你.
開始的時候很傷心.
失戀.
國語就說.
歸咪.
廣東話好像不是這樣說.
歸蜜.
歸房的蜜友.
廣東話可能叫好姊妹.
男生那些.
死黨.
來安慰你.
他不愛你不要緊.
當然有一個更加合適的人愛你.
死了心.
不要再想他了.
經過一段時間果然死了心.
Facebook就刪了他.
電話電郵都取消了.
情信卡片都扔了.
現在很方便.
電子相片都可以刪了.
兩個人一起拍.
你不要他.
可以刪了.
死了心.
當他不存在.
在你的生命裡.
他真的死了.
對你一點情感都沒有影響.
他曾經是你生命的一個部分.

$^{241}$那一段的生命就死了.
當然認識新的朋友.
甚至將來戀愛成熟.
那是新的生命.
你仍然是你.
但這兩段的生命.
就不同了.
所以耶穌說.
你要喪失的.
要失去的.
要天天處死的生命.
跟後面所說的生命.
剛才不同顏色的生命.
一定是相反的.
這個道理很簡單.
其實是一個得失的問題.
今年是馬年.
大家都記得.
過了新年.
會不會真的像我們所說的.
又要馬而好.
又要馬而不吃草.
行不行啊.
你又要馬健康.
跑得快.
又不給牠吃草.
又不給牠營養.
有沒有這樣的事啊.
很難啊.
又要馬而好.
又要馬而不吃草.
其實耶穌所說的.
你要喪失的.
要失去的.
要天天處死的.
跟後面的真的不一樣.
你不能兩樣都行.
必須要付代價.
必須要有所取捨.
有得和有失.

$^{281}$每個人的時間就是這樣.
你的精力就是這樣.
畫一個圓圈來代表.
你要藍色的那裡多一點.
紅色的那裡當然少一點.
你要紅色多一點.
藍色少一點.
這是得失的問題.
我們在有限的時間.
有限的精力裡面.
一定要有所取捨.
如果你說我要多點認識神.
那你就少點認識那些明星,歌星.
不看他們的新聞.
如果你想多唱一些聖詩.
那些流行曲就當然要少聽一點.
你需要這麼多時間.
你要背那些歌詞.
我要和某某歌星一起唱.
你花了時間背那些歌詞.
那你就背不了聖詩.
你不要說兩樣都行的人.
是少的.
不是說沒有.
是少一點.
你的時間就是這麼多.
你要多讀聖經的話.
就當然要少看一些新聞.
少看一些小說.
少看一些煲劇.
不是叫你不看.
就是少一點.
這方面就多一點.
有些人就說.
我現在趁還有力氣的時候.
多賺點錢.
打一份工都不夠.
打兩份工.
那你身體一定是累的.
累的時候就算你來得團契.

$^{321}$來得崇拜.
你都要合眼訓.
精神不夠.
就是這麼多.
你這麼累的時候.
怎麼有時間來到.
有精力來到參加教會的活動呢?.
所以跟從主是需要有付代價的.
你要認識祂.
你要明白祂的話語.
你要記得並且實行祂的教訓.
是需要付代價的.
每個門徒都需要自己反省.
你願意放棄什麼.
你願意喪失什麼.
你願意不受哪一方面的掛慮.
去拖著你.
以至你可以得到最有價值.
最寶貴的生命.
要付代價.
你就要自己計一計.
第二方面.
我們就說.
不和別人比較的那種是跟隨.
雖然都是跟隨耶穌.
但是你不要和別人比較.
剛才我們所選讀的經文.
就給得很好奇.
耶穌和他說.
你跟從我吧.
我們一轉頭就看到.
原來有一個約翰在這裡.
這是講到福音書的作者約翰.
他說:那他又怎樣?.
我將來年老可能要被殉道.
這個比我年輕的約翰.
他將來會怎樣?.
我就知道不會善終.
不會死得好.
那他又怎樣?.

$^{361}$耶穌還說.
他怎樣跟你有什麼關係?.
假如我要他等到我來的時候.
還在這裡.
跟你有什麼關係?.
如果耶穌已經離開世界.
升天之後再來.
到現在還未再來.
如果耶穌再來的時候.
約翰還在的時候.
即是說他不用死.
如果是我.
第一個反應是什麼?.
就是公平.
為什麼我不能好死.
他不用死?.
是不是他給我好東西?.
或者是不是耶穌愛他多一點?.
一定會這樣問.
跟別人比較的時候.
就會覺得是這樣.
我們慣性.
跟別人比較.
當然是跟那些比較好的人比較.
春天來了.
雖然今天早上有點雪.
不過大概是春天.
英文有句說話.
The grass is always greener on the other side.
即鄰居的草永遠是綠一點.
種得多好又綠.
我們自己的草地.
這裡有黃.
我們都是看比我們好的人.
有時有些妒忌.
或者有些不舒服.
為什麼呢?.
因為別人的種比較好.
上一代的教育.
其實父母和先生的時候.

$^{401}$都將我們跟其他同學比.
或者跟鄰居的小朋友比.
甚至跟你自己的兄弟姐妹來比.
他們希望刺激我們.
使我們發奮圖強.
賭力.
目的是這樣.
通常都是反效果.
更加憤怒.
有時傷害了我們自己.
或者是自己放屁.
你說我不行,我就不行.
免得生氣.
早前大家有看福音話劇.
剛剛那幾天我們出埠.
所以沒有出席.
但後來Winnie就發了一個copy.
連結給我們看.
真的很精彩.
話劇裡面的媽媽.
怎樣罵女兒.
將她跟別人比.
女兒跟兒子比.
跟弟弟比.
跟別人的女兒比.
就罵她.
傷害了她的心.
我記得我小時候.
我媽媽在那些.
有時是那些阿姨.
那些阿姨就說.
我忘記了多少歲.
她說我這個兒子.
讀書又不夠姐姐厲害.
又是我姐姐.
經常都考前面那幾名.
我沒有她那麼勤力.
讀書又不夠姐姐厲害.
頭腦又不夠哥哥聰明.
哥哥真是聰明一點.

$^{441}$我反應慢,他反應比較快.
如果這樣比的時候.
真的很傷.
又比不上姐姐.
幸好那時候媽媽就說多一句.
不過這個兒子都很喜歡幫人.
這樣也好,起碼一樣東西都不錯.
跟兄弟姐妹.
哥哥姐姐這樣比的時候.
其實是很傷細蚊仔心.
後來我做了爸爸之後.
就提醒自己.
跟兒子互動或聊天的時候.
千萬不要跟他比.
就算他長大了.
都不要說你讀完大學.
他又讀到博士.
他賺錢多少.
人家這樣年紀都結婚了.
有子女了.
不要這樣比,不能比.
我們說人比人就氣死人.
或者比死人.
我們當中不少人信主的年日.
不算太長.
或者信了幾年.
當你信主越久的時候.
你就越能體會到耶穌說.
進去天國的門是窄的.
能夠真正進入的人是少的.
我牧會的時候.
有一個弟兄非常熱心.
他曾經帶領一個朋友.
信了耶穌.
這個朋友後來成為教會的長老.
但反而帶他信主的弟兄.
後來有婚外情.
連教會都不去了.
現在怎樣就不知道了.
隔了很久沒有聯絡.

$^{481}$如果你要比的話.
你要跟誰比?.
你跟信主大人的傳福音.
有人經過他信耶穌.
跟他比,但他後來離開了教會.
還是跟這個後來信主的.
他又去教會裡事奉.
後來成為了長老.
你跟誰比?.
兩個都不好比.
一般人眼中所謂成功的教會.
或者興旺的教會是什麼?.
人多,錢多,活動多.
有些人說.
如果T4C好像那個教會就不錯了.
幾千人聚會.
什麼節目都有.
每天教會都開放.
因為每天教會都有活動.
不用等到星期六,日.
這樣不錯.
幾千人的教會.
有七,八個牧師.
有人離開都無所謂.
有很多牧師繼續去牧養.
事實上一些所謂超級教會.
Mega Church.
教會幾萬人.
不是幾千人.
是幾萬人.
每個週末就一,兩萬人聚會.
崇拜都有五,六堂.
可以坐幾千人的大堂.
有些還有電視轉播.
全世界都有.
曾經風光一時.
但後來因為各種原因.
傳道人下台.
教堂破產.
要買.

$^{521}$有些人是貪錢.
或者挪用公款.
或者是瞞稅.
有婚外情.
濫用權柄等等.
有時我們每次聽到這些消息.
不要幸災樂禍.
我們的傳道人到今日.
仍然在恩典當中事奉祂.
更重要的是不要跟他們比.
只要我們忠心跟從.
不要看別人說.
又可以說得很淑女.
但有時會跌倒.
人而已.
傳道人都是人.
都會有跌倒的時候.
軟弱的時候.
不跟別人比.
或者不理會別人.
不是各走各路.
但不是因為別人的制遇.
你自己跌倒.
不要因為別人的制遇.
你自己退後.
不要因為別人不信.
你就不信.
有些人說.
這樣信耶穌沒有用.
你不信就吃虧了.
你只看別人跌倒的時候.
不跟從耶穌是你的損失.
所以耶穌會說.
他們以後和你有什麼關係.
你跟從我吧.
這是你和耶穌之間的關係.
你跟從我吧.
不用跟別人比.
第三是持續的跟從.
耶穌復活後對彼得說.

$^{561}$你跟從我吧.
這是不是第一次跟他這樣說?.
不是.
因為根據福音書的記載.
耶穌第一次呼召彼得的時候.
做門徒的時候.
彼得和兄弟安德烈還在打魚.
耶穌呼召.
來,跟從我.
我要叫你們得人如得魚一樣.
那是第一次呼召.
所以從文字上來看.
跟從不一定是呼召缺芝.
舉手信耶穌.
未必是指那樣.
而是主說你們要天天跟從我.
天天跟從我.
明顯是有另外一種意思.
意思是要跟我來.
跟我來看.
或者跟去這個地方.
不要到處走.
要貼近我.
所以在加利利海邊向門徒顯現.
主對彼得說.
你跟從我吧.
其實重點是.
不要在乎別人的看法.
自己要做自己的門徒.
要效法我.
效法耶穌基督.
當耶穌在世界上的時候.
跟從他意味著什麼呢?.
比較簡單.
耶穌在那裡.
可以走可以走.
你就跟著他.
去到哪裡就跟到哪裡.
去到哪裡門徒就跟著他.
他見什麼人你就見什麼人.

$^{601}$他行神蹟趕鬼.
門徒看見的.
耶穌是怎樣的.
是依字.
他講到或者他怎樣解釋.
門徒留意聽.
祈禱.
門徒在旁邊聽著.
學習怎樣祈禱.
就跟從.
很簡單.
因為耶穌就在那裡.
但我們今天.
耶穌不在身邊.
你看不到.
你不能說耶穌去哪裡.
你只是想像.
或者你以為.
其實看不到.
那怎樣去跟從呢?.
今天我們要跟從耶穌.
是意味著什麼呢?.
下面我會講四方面.
第一方面就是要放棄一切.
特別是放棄一切的擁有權.
耶穌說過.
如果你不撇下一切.
不能夠作為門徒.
當然我們今天的屋.
你的屋契仍然是你的名字.
或者是聯名.
你的車的標誌也是你的名字.
但那個真正的擁有權.
是在神那裡.
所以如果你那間屋.
是你的名字.
但你願意為主的名.
開放這間屋.
給人來.
在這間屋裡有小組聚會.

$^{641}$或者甚至你有個房間.
如果我有外地來的港元.
或者有些宣教士.
剛剛回來.
沒有地方住.
你給他住.
因為這間屋是屬於神的.
祂才是最後擁有的.
所以如果你這樣用這間屋.
或者用車.
載人去聚會.
或者車去買菜.
這輛車的標誌是你的名字.
但擁有權是在神那裡.
所有的財產都是屬於神的.
那你就是一個放棄一切.
你可以跟從.
甚至我們的兒女.
我們的後代.
都是屬於神的.
如果神要接走.
我們要很痛苦地接受.
是很難過的.
我知道有些弟兄姊妹.
可能都是兒女那一代.
或者是孫子那一代.
都有些是被神接走.
有學校.
約伯.
約伯那時候.
我們記得舊約裡面說.
他有十個兒女.
七個兒女 三個兒女.
同一天之間.
十個都不在.
多傷心.
不過他講了一句話.
賞賜的是瑤華.
收取的是瑤華.
瑤華的名是應當稱頌的.

$^{681}$因為他知道.
雖然這些兒女是他的.
但最後是屬於神的.
如果神要拿走的時候.
他才能夠接受.
他仍然是稱頌神.
所以我們要想想.
無論是財產也好.
兒女也好.
都是屬於神的.
當然最基本就是寫記.
如果我們連命都沒有.
捨棄生命的話.
還有什麼是捨不得的呢?.
如果有些東西是捨不得給神的話.
真的說不通.
因為你要捨棄自己.
連生命都要給神.
生命是屬於神的話.
這樣你就可以跟從耶穌.
第二方面是經歷耶穌的同在.
當然這最先是要進行他的話.
耶穌自己也說過.
我是葡萄樹.
你們是枝子.
如果你看《約翰福音》十五章.
有時也很亂.
我也試過畫一個圖.
越畫越亂.
他在我裡面.
我在他裡面.
他在父裡面.
父在他的裡面.
我們在誰裡面.
亂七八糟.
其實你拿著這句話就可以了.
是進行他的話.
我的話在你的裡面.
你慢慢就會在想法上.
在說話上.

$^{721}$你會效法耶穌基督.
你跟著行動就對了.
聖靈是叫我們想起耶穌的說話.
你經常讀他的說話.
在我們裡面的時候.
神就在我們的裡面.
我們可以行出來.
特別是進行耶穌的話.
要彼此相愛.
除了所謂大命令.
之後就是大使命.
大使命是你們要去.
洗萬民作我的門徒.
奉父子的命.
聖靈的命.
看他們施洗.
讓他們加入教會.
我教訓你們.
你們也教訓他們.
一個傳承.
大使命.
叫我們來到傳福音.
繼續栽培信徒.
另外一個就是彼此饒恕.
同心祈禱.
有時我們只看這聖經.
若有兩個人在地上.
同心合意求什麼.
我在天上的符都替他們來成全.
因為在哪裡有兩三個人.
奉我的名罪會.
我就在他們當中.
有時我們只說.
我們奉主的名罪會.
兩三個人在一起.
耶穌在我們當中.
其實在這段聖經之前.
是說到彼此饒恕.
即是說你彼此饒恕之後.
能夠同心祈禱的時候.

$^{761}$我就一定在你們當中.
所以如果你要耶穌在我們當中.
我們要學習.
要彼此饒恕.
有些事我們真的很困難.
特別有些情況下.
是不容易的.
第三我們要做耶穌所做的事.
當然我們要履行神的命令.
要愛神,愛人.
愛神包括敬拜神,奉獻,稱頌.
亦包括照顧我們自己的健康.
這也是愛神的一種表現.
因為這個身體是神賜給我們的.
是聖靈居住的地方.
如果沒有好的健康.
沒有好的身體.
怎樣去事奉呢?.
要履行神的命令.
愛神其實也是愛我們自己.
不要分得太開.
愛人也是.
怎樣去服務社會.
怎樣去貢獻人類.
這些都是神的命令.
另外就是《約翰福音》第十四章十二節.
耶穌說:我所做的.
信愛人都是要做.
而且做比這些事更大.
當年耶穌算了醫治了多少人.
可以數,很有限.
但今天我們這個時代有醫院.
從事醫療服務.
可以幫到很多人.
可以醫治到很多人.
還有陪伴他們.
教導他們.
很多東西可以參與.
是做耶穌當日做的事.
甚至一些比較闊的.

$^{801}$在法律制度之下.
如果做律師.
可以幫一些人去爭取公義.
照顧窮人等等.
很多事我們今天可以做.
是做耶穌所做的事.
第四方面就是.
組成一個耶穌屬靈的身體.
即是教會.
祂是頭.
我們彼此成為祂的身體.
所以聖經裡面有很多的.
彼此相愛.
彼此關懷.
彼此照顧.
或者互相尊重.
互相接納等等.
很多發揮教會整體的功能.
這亦成為跟從耶穌的表現.
所以總括來說.
我們要放棄一切.
放棄一切.
放棄一切的擁有權.
我們要經歷耶穌的同在.
我們要做耶穌所做的事.
組成耶穌屬靈的身體.
雖然我們今天肉眼看不到耶穌.
但如果我們能夠在這幾方面提醒自己.
我們就是跟從著.
亦知道祂在那裡.
我們就去到那裡.
耶穌復活後就在海邊向門徒顯現.
數一數.
海邊裡面有七個門徒.
特別是和彼得說.
你晚年的時候要受苦.
之後就叫彼得跟從祂.
如果以時間計算.
之後彼得有三十年左右.
是繼續耶穌的工作.

$^{841}$繼續跟從祂.
直到為主殉道.
耶穌那時候已經不在世界上.
但他依然繼續做耶穌所做的事.
之後真的殉道.
又過了二,三十年.
約翰才因為年老而離開世界.
約翰是十二個門徒當中.
唯一一個不用經過特別苦難離開世界.
其他的使徒都是因為信仰的緣故殉道.
使徒都是因為耶穌的緣故.
無論怎樣的結局.
他們特別是彼得和約翰和其他門徒.
他們都是一生跟隨著.
今天耶穌也和我們一樣說.
「跟從我吧,你跟從我吧」.
有一首舊歌.
有些比較年長的人懂得唱.
今天我不準備唱.
「向前走,努力向前走,一生不回頭」.
「向前走,努力向前走,前進莫退後」.
「守乎著泥,向後看的人,不配盡天國」.
「向前走,努力向前走,主就在前頭」.
我們同心祈禱.
天父我們知道在這個世界上.
我們今天跟隨你不是一個簡單的事.
是一個生命的轉變.
是一個每一天都要付上代價.
每一天背著我們的十字架來跟從你.
天父幫助我們不斷有操練.
知道我們所有的其實都是你的.
你今天給我們有生命.
給我們能夠用上你.
給我們恩賜是你的恩典.
天父幫助我們在餘下的光陰當中.
無論在世界上有多久.
我們都是每一天跟從你.
每一天效法你.
求你的聖靈在我們當中引導.
指導我們向哪個方向來.

$^{881}$幫助我們了解今天我們能夠為你做什麼.
以致我們在生命當中真的得到這個生命.
一個永遠的生命.
\newpage



\section{哥林多前書 4:1-13}
\label{sec:032021}
\textbf{使徒的真面目}~~~~黎炳鍵弟兄
\newline
\newline
連結: \href{http://www.church.com.hk/acms/content.asp?site=cdc\&op=show\&type=product\&code=032021\&layout=sermon}{\texttt{032021}} ~~~~ 語音日期: 2026-04-19
\newline
\newline
\hyperref[sec:032020]{\small{< < < PREV SERMON < < <}}
~
\hyperref[sec:index_chronic]{\small{[返順時目]}}
~
\hyperref[sec:032022]{\small{> > > NEXT SERMON > > >}}
\newline
\newline
哥林多前書 4:1-13
\newline
\begin{longtable}{cl}
\hline
\hline
章節 & 經文 (和合本修訂版)\\
\hline
4:1 & \begin{tabularx}{0.7\textwidth}{X} 人應該把我們看為基督的執事，為神的奧祕的管家。 \end{tabularx} \\ \\ \relax
4:2 & \begin{tabularx}{0.7\textwidth}{X} 所求於管家的，是要他忠心。 \end{tabularx} \\ \\ \relax
4:3 & \begin{tabularx}{0.7\textwidth}{X} 我被你們評斷，或被別人評斷，我都以為是極小的事；連我自己也不評斷自己。 \end{tabularx} \\ \\ \relax
4:4 & \begin{tabularx}{0.7\textwidth}{X} 雖然我不覺得自己有錯，卻也不能因此判為無罪；審斷我的是主。 \end{tabularx} \\ \\ \relax
4:5 & \begin{tabularx}{0.7\textwidth}{X} 所以，時候未到，在主來以前甚麼都不要評斷，他要照出暗中的隱情，揭發人的動機。那時，各人要從神那裡得著稱讚。 \end{tabularx} \\ \\ \relax
4:6 & \begin{tabularx}{0.7\textwidth}{X} 弟兄們，為你們的緣故，我拿這些事應用到我自己和亞波羅身上，讓你們從我們學到「不可過於聖經所記」這話的意思，免得你們自高自大，看重這個，看輕那個。 \end{tabularx} \\ \\ \relax
4:7 & \begin{tabularx}{0.7\textwidth}{X} 使你與人不同的是誰呢？你所有的有哪一個不是領受的呢？若是領受的，為何自誇，彷彿不是領受的呢？ \end{tabularx} \\ \\ \relax
4:8 & \begin{tabularx}{0.7\textwidth}{X} 你們已經飽足了，已經富足了，用不著我們，自己就作王了。我願意你們果真作王，讓我們也可以與你們一同作王！ \end{tabularx} \\ \\ \relax
4:9 & \begin{tabularx}{0.7\textwidth}{X} 我想，神把我們作使徒的明顯地列在末後，好像定死罪的囚犯，因為我們成了一臺戲，給世界、天使和眾人觀看。 \end{tabularx} \\ \\ \relax
4:10 & \begin{tabularx}{0.7\textwidth}{X} 我們為基督的緣故成為愚拙的；你們在基督裡倒是聰明的。我們軟弱，你們倒強壯；你們有榮耀，我們倒被藐視。 \end{tabularx} \\ \\ \relax
4:11 & \begin{tabularx}{0.7\textwidth}{X} 直到如今，我們還是又飢又渴，又赤身露體，又挨打，又到處漂泊， \end{tabularx} \\ \\ \relax
4:12 & \begin{tabularx}{0.7\textwidth}{X} 並且勞碌，親手做工；被人咒罵，我們就祝福；被人迫害，我們就忍受； \end{tabularx} \\ \\ \relax
4:13 & \begin{tabularx}{0.7\textwidth}{X} 被人毀謗，我們就勸導。直到如今，人還把我們看作世上的污穢，萬物中的渣滓。 \end{tabularx} \\ \\
[1ex]
\hline
\hline
\end{longtable}
$^{1}$今日想同大家講到呢個題目叫做「使徒的真面目」.
當睇到呢個咁樣ge 題目ge 時候.
咁可能有D 頂展會有朋友會心諗.
係咪有D 咩秘密要講呢?.
所以要有咩真面目.
好似以前D 電視台有D 節目叫做咩「睇真D 」咁樣.
咁黎到去揭開D ye 啦.
咁我諗即係會更加會諗下會唔會好似.
達文新密碼咁樣帶比一D 有趣ge .
題目黎到係去思考呢咁樣.
咁首先講呢就係.
即係我唔識得講D 揭秘呀gwo D 咁ge ye .
呢個唔係我ge 專長黎架.
咁呢個唔係我.
咁另外就係呢.
我同好多大部分.
即係大家ge 人一樣.
我唔係話牧師或者傳道人.
我都係打份工ge 平常日子.
咁我唔係D 業內人士有D 咩ge 秘密想同大家講.
咁如果要問就問牧師啦.
咁係咁但係呢我想講ge 就係話.
昔日ge 教會裡面.
如果我地頭先多謝Sandy幫我地讀出呢段經文.
幫我地讀出關於一D 使徒ge 一D ge 事情.
昔日ge 使徒就好似彷彿今日ge 一D ge 教會ge 領袖.
包括牧者或者一D ge 長輩.
令到我地對呢方面有D 咁樣ge 聯想.
亦都以為呢.
我地再睇返就係話.
今日我地再睇呢段經文.
頭先Sandy同你讀過啦.
關於使徒亦都關於.
即係我地頭先聽ge 時候.
亦都包括保羅自己佢自身.
一D ge .
即係去述說返.
即係有D 咩發生過呀.
或者係點樣ge 事.
咁我地從呢段經文裡面.

$^{41}$高林多前書ge 第四章一至第十三節裡面.
我地希望黎到係去睇一睇.
當其時呢一班ge 使徒.
保羅講到自己講到D ge 使徒.
究竟佢地係做緊D 咩ye .
佢地係D 咩ge 身份.
佢地從何而黎.
佢地係做緊D 咩ye .
黎到係對應返我地今日.
對於教會領袖.
以至於一般ge 領袖.
我地ge 睇法.
有D 咩ge 唔同呢.
或者係話我地對領袖呢樣ye ge 睇法.
同保羅係度所寫ge .
係有幾大ge 落差呢.
咁.
首先我地就從第一節gwo 度黎到係去睇一睇.
四章ge 第一節.
咁如果你有聖經ge 呢都可以打開我地一齊跟住黎睇.
咁首先第一節gwo 度保羅呢度咁樣寫架和合本咁樣寫架.
人應當以我地為基督的執事.
為上帝奧庇時的管家.
第一節呢度就已經係一個好明顯ge 一個ge .
講jor 一樣ye 就係話原來.
係保羅眼中呢一個ge 使徒係有兩重ge 身份.
第一個就係執事.
咁關於執事有一D ge 如果你係睇第二D 版本呢有D 係譯做servants.
即係譯做僕人.
呢個執事有D 咩ge 特點呢就係如果我地再睇返之前第三章我地之前第三章有講過啦第二即係第二二至第23節.
有講到其實係一切ge 事都係屬於上帝屬於基督.
基督係屬於上帝咁因此一切受造ge 事呀.
睇到睇唔到過去ge 事現在ge 事都因著係基督ge 底下.
基督屬於上帝咁而一切ge 事呢D 受造之物都係屬於上帝ge 咁而我地作為人如果我地屬於基督ge 話呢我地相信主相信基督ge 話呢呢D ge 事都係屬於我地ge .
咁呢度我地睇返就係話僕人呢個角色或者執事呢個角色就係話教會唔係比佢地係度話事ge 而係佢地都係屬於僕人ge 一部分屬於黎到係去創造ge 一部分係屬於呢個群體ge .
咁但係呢咁如果話呢D ge 執事係為神ge 奧秘黎到係去做事咁即係話佢地唔自主啦咁即係係咪好似而家有D 有體制ge 教會就係教會有講D 咩ye 或者做D 咩ye 呢個ge 執事就要跟住去做呢咁樣咁亦都但係我地睇返係第一章呢度.
亦第一節呢第四章第一節呢度亦都有講呀呢一D ge 使徒亦都係神奧秘ge 事ge 管家咁有關於奧秘ge 事gwo D ge 詳情呢咁其實大家應該係去睇返哥林多前書ge 第二章就有去好詳細ge 解釋不過係呢度亦都簡單黎到講呢就係話呢個ge 奧秘ge 事就係指到耶穌基督同埋佢ge 十字架.
當好多人以為要去贏要去得勝就係要去有擁有好多ge 權力好多ge 武力好多ge 財力黎到係去先至黎到係去得勝到ge 時候呢第二章哥林多前書ge 第二章gwo 個講基督卻係以十字架ge 事黎到係去戰勝死亡.
係好多人黎到講係覺得呢件係一件不可思議ge 事點樣係用一個軟弱用死亡十字架係一件凝聚黎到去jor 咁呢個就大家可以睇返第二章之前同大家睇過下或者自己去睇ge 時候你都可以睇得返但係呢件事係人所唔明白ge 事黎架.
呢一個ge 死ge 人黎到係去當然佢會復活啦但係黎到係去戰勝呢一個ge 死亡戰勝呢個罪呢個係大家唔明白ge 事呢個就係gwo 個奧秘咁我地係過去呢個幾月以至過去呢幾個星期我地都見證過呀我地係世界ge 政治或者gwo 個ge 軍事我地睇過呢個ye ge 時候呢個ye 係一個ge 奧秘ge 事.

$^{81}$gwo 個ge 軍事我睇到一個擁有絕對武力絕對軍力ge 國家佢地可以將人地第二個國家ge 總統阿山擒到自己ge 國家裡面等下受審啦亦都睇見佢可以出動絕對ge 軍力黎到係去救返佢一個ge 士兵咁樣我睇見係呢個現代ge 世界裡面我地都睇見一種絕對ge 能力.
你都係去壓制另外一個一D ge 能力黎到係達成佢地所做ge 事咁而家咪話大家咪要一同黎到係做一個ge 升級啦你ge 能力大D 咩我就再加自己ge 能力總之係一個ge 升級希望呢你都係透過呢個咁樣ge 自我能力gwo 個增加黎到係去你都係去贏對方你都係cum 個對方你都係去.
使到自己要達成ge 目標黎到係去達成但係我地睇見上帝出自逐基督同埋佢ge 十字架一個體諒好似係失敗者你都係去勝過死亡並且將所有萬事黎到係收歸佢ge 氣下呢個係方法係有為人ge 常理所明白保羅係度再次重申.
只存基督同埋佢ge 十字架而佢係呢件事上係作為一個官家係呢個ge 奧秘呢D ge 上帝奧秘ge 事係引領人黎到係去真係同神修和.
先能夠係要上帝ge 靈先係明白咁呢個就係gwo 個管家gwo 個ge 角色咁我地睇見呢一個ge 保羅講述到呢個ge 使徒ge 呢個雙重ge 角色呀既係執事亦都係上帝奧秘ge 管家呢件事ge 時候咁我地就要睇啦咁究竟佢地係要做D 咩ye .
咁我地就睇下第二節.
第二節gwo 度咁講所求於管家的事要他要忠心呢個講就係講呢D ge 管家神ge 奧秘士ge 管家所要ge 係忠心.
原來呢個管家唔係擁有教會呢D 管家佢地係受上帝ge 委派委託你都係去看守呢個教會確保教會係全港呢個基督同埋十字架.
咁呢個呢就其實可能我地都唔係一D 即係得罪D 講句我地唔係D 咩大富人家.
我地都係從一D 閱讀呀或者係從一D 甚至咁講啦片集呢.
即係即係D 古裝片或者D 管家咁樣呢.
佢係打工ge .
咁老細要佢去看住呢個屋企咁打理得井井有條咁樣.
但係好似一個僕人咁樣好似冇自主性.
但係管家另外一個ge 角色就係話佢要維持呢個家ge 秩序.
甚至如果你係一個皇室你要維持整個皇室ge 管家皇室ge 運作.
咁會講呢戲裡面係咪有D 少主ge 即係未做即係儲君呀或者係未做皇帝即係大支.
管家ge 工作都係係佢地未長大成人之前呢都係去教導佢地.
即係無論你係睇青裝片又好或者係睇英國宮廷片都係咁樣架啦.
個少主呀或者係個皇帝gwo D 公主呀都會見到個管家都好驚呀驚佢地都係去做錯ye 黎到係罰佢呀咁樣.
咁我地睇到一樣ye 就係話原來管家雖然佢gwo 個身份上係好似係一個被.
係一個好似打工ge 或者係一個唔係一個自主ge 人係一個甚至講話係一個僕人.
但係佢亦都係因為佢管家呢個身份呢係被授予一D ge 權力.
黎到去管束一D 未來會管返佢ge 人黎到係去確保佢地係有一個正常健康.
係呢個ge 正路黎講黎到係去成長ge 道路咁我地睇見原來呢個管家係有呢個權利.
咁而呢度第二節我arm arm 讀到就係對呢D ge 管家係需要有要求係忠心.
咁忠心D 咩呢又係返返去頭先一路我不露全幅就係究竟係咪忠心於去傳基督同埋佢十字下ge 道理呢.
咁呢個就係你睇返之前ge 經文你即係唔係憑空想出黎架呢個係一路講緊ge .
即係即係我地今日有個教會有好多唔同ge 班有好多唔同ge ye 呀烹飪班跳舞班空手道台權度乜乜好多ye 呀補習班gwo D 都即係以現有社區黎講唔係唔係大問題ge .
只不過就係話係做曬呢D ye ge 時候究竟基督ge 道係咪係咪得到傳開呢.
還是係我地係呢D ge 處理呢D 康樂文娛呀community service呢D ye ge 時候呢我地慢慢都唔記得jor 最主要ge 角色我地就係讓人係知道基督同佢十字下ge 道理呢.
咁呢個我地就即係作為係教會無論我地係咩位置我地都係有呢個警醒不過睇到呢兩節經文ge 時候我地有一個要搞清楚我地有兩方面可能ge 誤會.
第一個就係話我地以為教會ge 領袖不論係牧師或者係傳道人我地都以為佢地係教會ge 老闆乜ye 都係話世事ge 永遠呢佢講ge ye 都係要一錘定音佢講ge ye 你就係要去跟要去做.
咁但係我地係當中有冇審審測到就係話究竟呢個傳道人或者牧師佢係教ge 事好似講一錘定音ge 事究竟係帶領弟兄會更加對基督更加ge 認識呀對於十字架對於教會更加明瞭了解.
還是話keep the ball rolling有D 時工咁大家即係教會運作呢咁呢個我地係要第一方面我地第一個誤會我地要搞清楚對傳道人對教會ge 身份第二個呢誤解呢就係以為牧師傳道都係打份工ge je 咁教會有錢請ge 咁咪洗到就洗囉見jor 個堂咯咁咪請個位啦.
咁咪請個牧師黎即係衰D 講句連個連個janitor都慳返ge 咁有D 位置下咁但係我地睇見唔係呀.
神ge 僕人係係授權係教會裡面黎到去教導你都係去執行佢地ge 職責黎到去教導人認識基督認識十字架係呢方面正如好似gwo 個頭先講過gwo D ge 咁ge 古裝片呀少主都要聽gwo 個ge ge ge 管家話even though eventually.
最後就算最後呢一個ge 少主長大後係話係話事係掌管呢個家庭呢個皇室ge 人都好係佢未真正成為主人之先佢係需要服從管家ge 教導佢有好多唔成熟ge 地方係需要呢個管家黎到去看護保護咁呢個我地要明白呢個ge 管家gwo 個ge 道理.
基於保羅對呢一個ge 身份ge 明白.
第三節.

$^{121}$我自己也不論自己.
我雖不覺自己有錯.
卻也不能因此得以稱義.
但判斷我的乃是主.
原來明白到一樣ye .
使徒.
最終要終身ge 最終ge 對象終身ge 對象係上帝佢自己.
而判斷呢件事ge 事係上帝.
所以係呢度黎到黎到講呀講到呢度佢受到gwo 唔多人ge 論論斷或者係點樣點樣講佢.
佢都係睇jor 一D 極少ge 事.
今日我地面對別人ge 批評係點樣我地會好緊張.
我地尤其是一個social mediage 年代.
我地有陣時擺個comment係度.
有D 人舉個手指公呀比個讚你呀.
好開心呀咁樣.
但有D 人比D 唔好話講埋D ye 傻ge .
你冇腦架.
咁你又會好情緒好低落呀咁樣.
我地係情緒我地係比呢D ye 牽制住.
比別人ge 說話牽制住.
人地點睇我係最重要.
但係呢度講使徒gwo 個ge 情形就係話.
你呢班人講D 咩ye 佢都係好興旱.
係睇一D 極少ge 事唔係話冇ye .
有多ye 要處理.
咁所以即係保羅先寫jor 十幾章ge 經典呀.
咁大有都要處理但問題係最主要就係.
上帝點睇你呀.
你係咪忠心gwo 個呀.
但係上帝.
又冇say呀.
係我地自己ge 心中又冇say呀.
甚至保羅亦都講佢自己都唔論斷自己.
保羅講呢個就係話唔係話我.
我即係叫盤皮呀.
即係盤皮明唔明呀.
即係殺賴.
唔.
我地要睇下貓子呀.
佢叫抵賴.

$^{161}$唔關我事呀呢D 上帝你ge 事.
我唔理.
唔係講呢D ye 黎到係抵賴.
而係話.
而保羅亦都好自己好清心直率咁講.
我根本即係我自己黎講.
我都唔覺得自己有犯錯.
我唔係唔係唔係.
因為呢D 犯錯.
我我我我真係有犯錯.
於是我就話我唔理你D D 人地.
你話我地而家D 人最輕架嘛.
我唔care你點睇.
因為我覺得自己做ge ye 你做咩要管束我.
關你咩事.
保羅唔係呢個態度.
而係話.
佢自己都冇犯錯呀.
佢自己都講自己冇犯錯.
但係佢自己檢醒自己ge 時候.
佢都覺得佢自己都冇犯錯.
黎到係去作為一個提醒.
不過佢想講ge 就係話.
呢D ye 唔係靠呢樣ye 黎到係去.
係面係神面前黎到係去.
清移.
好似呢度講啦第四節gwo 度講架嘛.
卻不能因此得以清移.
即係唔係靠人地點睇我.
或者係我自己點睇自己.
黎到係去決定我係唔係係神面前黎到係清移呀.
la 係話最終審判gwo 個係神.
最終判斷gwo 個係神.
點樣去面對呢個上帝先至係最重要.
我係唔係忠心呢樣ye 唔係.
A呀B甚至我自己黎到係判斷.
呢種時候到ge 時候神黎到決定囉.
因此佢將呢D 事係睇為.
好別人ge 批評論斷.
黎到係去當為一件好.

$^{201}$好輕ge 事.
到到時候到了呀呢度.
呀頭先再睇下我地睇第五節.
所以時候未到.
什麼都不要論斷.
只等主來.
他要照出暗中ge 忍情.
顯明人心ge 意念.
那時各人要從上帝那裡得著稱讚.
我地睇到原來話.
保羅唔係話識時靈人呀唔好論斷呀唔好成呀唔好指.
即係唔論斷還唔論斷.
就唔係話講不分青紅皂白呀gwo D 咁ge ye .
但係而係話我地唔去.
將一件事係當係一錘定音講話.
呢D 係某某人係好某某人係壞.
呢個係上帝面前自有判斷.
但係我地自己作為.
我地自己黎講我地就係係面前黎到係去.
等待上帝對佢委託ge 人.
係時候滿足ge 時候黎到係去作出一D ge 判斷.
而保羅係度亦都提及到呢D ge 人.
都會係因著佢地受委託ge 事.
而係會得著上帝ge 稱頌.
原來使徒最終要負責ge 係上帝自己.
而上帝亦都係會按著每一個人.
係佢面前有幾多ge 忠心.
黎到係去稱讚佢地.
咁其實呢一個概念.
係有一個將來末世審判ge 概念.
亦都承接返之前我地睇第三章.
提及到無論係耕種ge 比喻.
或者係建築gwo 個ge 比喻.
如果大家唔記得可以睇返第三章.
即係講到種好jor D ye .
或者係個樓起好jor 比火試驗過.
佢地已經survive到啦 存在到啦.
咁樣上帝係會因著呢D ge 事.
黎到係去有獎賞.
我地睇見原來我地人唔係話.

$^{241}$我地做自己ge ye 黎到係去.
為自己負責.
一人做事一人當.
咁樣黎到係去做.
上帝最終會因著我地做ye 黎到係去肯定.
透過佢ge 稱讚黎到係肯定.
如果你係忠心係會得著上帝ge 肯定.
咁樣引申到就係話.
如果你要諗.
即係講賺得全世界ge 稱讚都好.
但係最終我地得唔到上帝ge 稱讚.
其實係一件最大件ge 事.
最不濟ge 事.
基於四章呢度第一至到第五節第四節.
上帝對佢ge 僕人對佢ge 管家ge 呢個中心ge 要求呢.
上帝點樣ge 判斷佢.
保羅提醒佢地就係話唔好.
提醒佢地就係話上帝係最終判斷gwo 個.
係呢段期間我地唔好互相ge 論斷.
互相ge 黎到係去踐踏.
因為未來ge 事自然有上帝黎到係去顯露出黎.
黎到係上判.
回過頭黎講我地就係話.
我地唔係話某某人帶領我信主傳福音.
我欠jor 佢.
呢個恩典係因為佢傳比我先知.
於是我地就話佢講D 咩arm ge .
某某牧師因為佢講ge 道.
報道會我聽佢信主.
佢講ge 係arm ge .
佢arm 哂ge .
而係保羅係呢度就係話.
按照保羅ge 教導就係話.
某D ge 傳道人向我個人傳福音也好.
或者係.
港獨會領導我信主.
佢地都係源於出於對主ge 中心.
佢地gwo 個ge .
結果佢地gwo 個中國係點樣上帝自然有.
判斷.

$^{281}$正如好似.
我諗我地作為一個環基督徒.
我地都要成日提醒自己.
即係.
百多二百年前.
當福音呀.
由西方傳入教會ge 時候.
好多人因此黎到係開始有接觸福音呀咁樣.
我地又唔好睇話.
我地好似欠jor D 西人D 咩ye 咁樣.
欠jor 佢地.
於是呢我地於是睇愛國的月光特別圓ge .
西人ge 教會一定係好過華人教會呀咁樣.
咁我地發展到今日我地睇到唔係啦.
西方教會都有好多ge ye .
問題係.
判斷一間教會好與好.
唔好.
唔係話.
佢有冇傳福音俾你.
佢有冇對你好呀.
而係話佢咪都係講緊.
基督同埋佢ge 十字架呢.
今日西方ge 教會.
都有好多做好多.
社會ge 服務.
有好多community service.
有好多包容呀有咩.
呢D 係咪都係十字架.
之下所做ge 事呢.
呢個係.
還是話.
我地唔係話.
因為西方教會比我地先信.
佢地有好多神學發展得好.
所以我地唔好駁佢呀.
唔係呀.
我地真係.
用返保羅呢度個個ge 判斷就係話.
究竟佢地係咪都係講緊基督ge 十字架呢.

$^{321}$咁呢個就係.
即係整體.
保羅係一個ge .
同大家講.
呢一個ge 使徒究竟gwo 個ge .
角色.
佢gwo 個身份.
同埋佢對佢要求係D 咩ye .
咁呢個就.
呢度第一至五節呢度gwo 個段落.
ge .
保羅要同大家講ge 說話.
同哥林多講ge 說話.
好咁跟住呢我地就睇第二部分.
今日講到ge 第二部分就係講到.
我比個題目叫一台戲.
一台戲係D 咩呢就係六至十三節.
咁呢台戲究竟係.
係講緊D 咩ye 呢.
我地睇見呢其實.
呢一個亦都係一個ge 延續呀其實係從第三章第五節開始.
頭先亦都有提及過啦.
即係係第三章保羅係有.
災眾呀僥倖呀.
黎到係去比喻神呢D ge 工人.
執事佢地gwo 個ge 工作呀亦都講到.
一個建造房屋ge 比喻黎到係去.
講呢D 佢地ge 角色呀咁樣點樣黎到係去.
忠心呀咁樣亦都講到就係.
最終佢地係會按著佢地所做ge 係.
黎到係.
黎到係去.
即係黎到係去賞賜佢地呀咁樣.
呢一切都係想講就係話.
呢一D ge .
工人所做ge 都係由上帝所委派.
真正令.
源於係上帝佢自己.
正如呢個ge 救恩.
十字架訊跡亦都唔係話因為人做jor D 咩ye .

$^{361}$去賺取返黎.
去earn返黎.
係神施予神比.
而呢D ge 神ge 工人黎到係去係當中ge 角色係.
gwo 種ge 參與.
我地唔好將D 人ge 角色.
睇得過高.
將佢同神同等咁樣.
咁我地黎到呢度ge 時候就係話咁究竟我地.
點樣ge 黎到係去.
去.
咁我地讀一讀第六章俾大家聽.
第六節.
弟兄們.
我為你們的緣故.
拿這些事轉俾自己和阿波羅.
即係之前講gwo D 事而家就apply落去.
應用落去佢自己同阿波羅身上.
叫你們效法我們不可過於聖經所記.
免得你們自高自大.
歸眾者過.
輕看哪個.
第七節.
是你們與人不同的是誰呢.
你們有什麼不是領受的呢.
若是領受的為何自誇.
彷彿不是領受呢.
呢度第六同第一節.
我地黎到係俾個一個標題.
使徒真正的身份和角色.
使信徒不要再彼此去論斷.
過去我地.
睇哥林多前書時.
我地睇見.
神的僕人.
上帝的執事.
管家.
佢地角色去中心傳道.
因此我地個焦點.
唔係話將功過.

$^{401}$或者將判斷落係呢D 人身上.
我地將呢D 人應該去看成為.
佢地都係領受神託付.
黎到去工作.
我地唔可以因為呢D 人.
黎到去因為自己屬邊派.
黎到去互相踐踏.
有D 人覺得阿保羅講ye arm 聽D .
覺得佢勁D .
就要去聽佢講.
有D 人覺得阿保羅講ye 比較好D .
阿保羅gwo D 差D .
黎到去互相踐踏.
因為佢地都係.
正確黎講.
佢地都係上帝所委託.
佢地傳ge 訊息.
就唔係話.
黎到去比我地.
聽完之後.
黎到去比較選擇.
黎到去互相踐踏.
呢度亦都係好似第六節.
呢度講.
叫佢地唔好貴重.
這一個 輕瀚那一個.
互相ge 黎到去.
唔妥對方.
係我自己成長ge 年份.
大家仲記唔記得.
邊個叫譚詠麟.
仲記唔記得邊個叫張國榮.
arm arm 四月死jor .
廿三周年.
通常張國榮gwo D .
Fans.
就會覺得.
譚詠麟.
唔夠靚仔.
唱歌.

$^{441}$雖然係好聽.
你知邊個譚詠麟.
我理你.
會覺得.
佢係唱歌好聽.
但係成日唔記得jor 歌詞.
唱唱完唱.
唱唱完就叫唱歌.
呢個就一般.
張國榮ge Fans.
對譚詠麟ge 睇法.
反過來講.
譚詠麟ge Fans.
可能鍾意到.
將自己ge 英文名.
改做Allen.
咁又會覺得.
張國榮係.
由頭粉面.
即係.
花花公子.
咁當然啦.
佢本身家境唔錯.
花花公子咁樣.
同埋唱歌就.
麻麻地je .
呢個就係講.
貴重者過.
漢重那過.
互相黎到係去.
唔係話.
我鍾意呢個.
我都欣賞你呢個.
我鍾意呢個.
你就係要差.
黎到係互相判斷.
好似呢度.
咁我地睇返歌咁多.
一樣.
鍾意保羅ge .

$^{481}$其他ge 都好易差D .
覺得阿彼得係好ge .
耶穌大弟子.
咁應該好ye 黎ge .
無出其右.
阿阿波羅.
亦都係一個有知識.
ge 背景.
佢地互相比較.
互相ge 黎到係去.
將自己體重ge 黎到係去比較.
黎到係壓低對方.
增加jor 自己gwo 種ge 優越感.
gwo 種ge .
比起人地叻.
其實呢一種ge 心態.
就好似我地而家gwo D 咁ge .
Fans黎到係去.
抬高別人踩低.
即係抬高自己黎到係去.
踩低別人.
gwo 種ge 心態.
所以.
波羅係度講得講得.
有D 咩ye 黎到係去.
使到你覺得自己係.
與別不同ge 呢.
第六節gwo 個講.
弟兄們我為你地ge 緣故.
那者.
第六節.
第七節啦sorry.
使你與.
人不同的是誰呢.
使你與人不同的是誰呢.
喜歡.
吹捧呢D 人.
抬舉人.
將人.
推到一個地位就係lor 黎去.

$^{521}$去追捧ge .
其實最主要就係話.
我係與別不同.
即係頭先講開張句.
但係推後少少啦.
鍾意呀.
到四大天王年代.
鍾意阿歌神gwo 個.
就會覺得.
另外三個唔識唱歌.
但係呢.
跳舞好叻gwo 個.
你鍾意gwo 個.
跳舞好叻gwo 個.
又覺得另外三個.
跳舞好似傷殘人士.
即係會有咁樣ge 比較.
咁就係話.
我鍾意呢個偶像.
佢gwo 個ge 特點.
佢gwo 個ge .
個人之處.
我都鍾意呢個偶像.
佢個人之處.
我都因此黎到去.
話.
我揀得呢個arm .
我揀得呢個合理.
我揀呢個識唱歌.
我揀呢個識做戲.
我揀阿波羅.
佢地都有佢地個人之處.
呢一種慢慢就係形成jor 一種.
就係話.
反映出就係話.
表面上好似係追捧唔同ge 人.
但係其實佢地係睇見自己.
透過呢個選擇.
話俾人知我係幾咁與眾不同.
但係保羅係呢度.

$^{561}$反問就係話.
使你與人不同ge 係誰呢.
我地所信ge 福音.
都係從.
七神而黎.
我地唔係話.
我地自己揀.
因為呢個對呢個.
信jor 基督教.
咁可能有D 人會係.
自己一係去講.
同自己講.
有一種產生jor 一個驕傲.
佢可能會同自己講.
呢個信仰係我自己揀ge .
我要信呢個神.
我本來唔係信主ge .
我係第二D 思想.
第二D ge 宗教ge .
第二D 無神論ge .
我係持另外一D ge 概念ge .
但係我係轉信基督教.
呢個基督教.
係我經過自己.
好深思熟慮.
比較過好多ye .
讀過好多ye 黎都係去揀.
我地睇見呢種.
係信主ge 群體裡面.
都有種咩驕傲.
甚至我地好多時都會講.
我已經.
立志決志去.
跟隨基督.
好熟齡呀.
去跟隨基督表面上.
但其實心裡面係一種驕傲.
就係話.
呢樣ye 係我揀ge .
呢個信係我揀ge .

$^{601}$我選擇跟隨基督.
我決志.
我立志.
黎跟隨基督.
個主體仍然係我.
當睇到呢度.
保羅gwo 個經文.
呢度亦都.
保羅可能會去問.
點解你講D ye 好似.
唔係你領受.
從神而領受.
好似你自己賺返.
呢個ge .
自我的一種心態.
其實係直根於.
係我地每一個人ge 裡面.
亦都係.
哥林多教會裡面.
有出現.
於是我地見到係之前經文有提及到.
要分黨分派.
先至黎到.
去搞.
有咁樣分裂ge .
情況.
如果我地略為對於哥林多教會.
個背景略為有少少.
認識ge 話.
佢地呢一個城市因為之前ge .
戰亂.
被二位平地.
咁大概係.
保羅寫哥林多前書.
一百年之前.
咁呢個.
不無之地.
因為頭先講.
之前都有講過.
哥林多地理位置.

$^{641}$既係南北交通要塞.
亦都係東西海路.
ge 一個重要.
ge 節點.
而家有條運河都係百幾年之前.
但係之前就係話.
唔使繞過南邊.
海域去運輸.
經過哥林多.
南邊就.
落返貨咁繼續.
雖然呢個.
當其時因為戰爭.
而二位平地呢個地方.
羅馬政府睇.
仲好有價值.
於是就再次.
殖民係呢個.
城市gwo 度.
用我地.
加拿大例子去諗.
五十年前.
多唔多其實係.
少人口.
滿地可.
到1976年先至.
爬jor 上頭.
之前不嬲係滿地可.
加拿大最大城市.
滿地可.
但係因為政治緣故.
因為法語.
慢慢好多.
大商業機構.
好多活動搬離多倫多.
於是多倫多就慢慢大.
即係.
如果你睇返一D 舊ge .
文獻.
或者係.

$^{681}$紀錄片.
多倫多係好多停車場.
好多.
工廠係度烏煙瘴氣.
如果講好住.
佢地會覺得滿地可會好住.
過多倫多.
係一個.
大恩著時移世易.
有多ge 政治經濟因素.
慢慢多倫多就爬.
50年前爬過jor 頭.
咁.
多倫多好似.
都咁樣理解.
佢雖然好似話.
唔係好似羅馬.
或者瓦典咁樣擁有.
咁深厚ge 歷史.
或者係有咁深厚.
政治影響力.
但係因為佢地理位置.
令到羅馬政府不得不重視.
不得不重見.
慢慢呢一班.
從羅馬黎ge 移民.
慢慢佢地.
開始又發財.
亦都係得到一個.
重見.
亦都成為一個好重要ge .
軍事同埋貿易.
ge 樞紐ge 時候慢慢成為佢地區.
一個好重要ge 地方.
係呢D 城市裡面好大ge 特色.
就係話.
有D 咩ye 唔係我自己賺返黎.
今日ge 地位.
今日我能夠搵到ge 錢.
今日我能夠比人地睇得起.

$^{721}$係我自己賺返黎.
用著呢個心態.
黎到去睇福音ge 時候.
就係話呢個信仰都係.
我自己揀返黎.
當我金錢滿足到.
我事業滿足到.
我家庭滿足到.
我咪信耶穌.
黎到去.
令我地有一個心靈上高高高ge 滿足.
為我ge 生命黎到.
係錦上添花.
對應返就係話.
呢一個ge 心態.
其實係好多.
有發達ge 地方.
都會.
大城市ge 地方ge 信徒.
都係一個有咁樣ge 情況.
但係呢一個ge .
宗教信仰.
永遠都只不過係一種ge .
錦上添花ge 形式.
係我自己揀ge .
我自己選擇ge .
從來都唔係我地.
或者呢一班ge 信徒.
gwo 種ge 核心ge 價值.
引導我地黎到去.
ge 一舉一動.
但係呢個係真正信仰ge .
原因啦.
呢個係保羅黎到去.
提出黎ge 就係話.
你地.
有邊個叫你地唔同呢.
無論你地大城市.
或者你地小地方.
你其實都係神揀選你地.

$^{761}$去相信佢.
點解呢D ye 你地好似覺得.
係自己去賺返黎ge 呢.
而呢種.
好似賺返黎咁樣.
自己選擇ge 方法就係.
透露jor 係佢地.
哥唔多教會裡面gwo 種ge 分裂ge 狀況.
咁往後我地再有機會.
睇呢段睇哥唔多前書ge 時候.
我地會更加睇到佢地呢一種ge .
心態黎到係去.
點樣ge 黎到係去.
影響到.
佢地所犯ge 錯.
同埋.
所對人對事ge 呢一D ge .
情況呢個我希望.
往後我地可以再睇ge 時候.
更加明白.
咁去到講完呢度.
一D ge 情形之後呢.
我地黎到係.
去睇一睇第九節.
第九節咁樣講.
我想上帝把我地.
使徒冥冥列在末後.
好仗.
定死罪ge .
囚犯因為我地成了.
一台戲.
給世人和天使.
就係話講到.
保羅自己.
即係係面對住呢一個.
咁ge 情況.
即係人自大自滿.
自己覺得係自己.
好滿足ge 時候.
佢自己同埋其他ge .

$^{801}$使徒成了一台.
戲黎到係去.
成了一台戲黎到係.
首先我想同大家講.
呢度就係好有趣ge 講到.
就係話世人同埋天使.
呢度講.
即係好多人講天使有好多.
唔同ge 解釋呀咁樣但係我諗呢個.
我唔詳細但係我地睇返一樣ye .
就係話原來成個ge .
過程運作ge 裡面.
唔係淨係比人睇.
亦都係天使睇.
係整個創造ge .
比我地睇ge 就係.
有人ge .
有D 唔係人ge .
我地.
同我地而家即係所謂一個.
啟蒙後ge 人.
我地成日睇就係話.
好多ye 都係事在人為ge je .
好多ye 都係人做ge .
但係呢度比到另一個角度.
原來呢一個ge .
使徒佢地呢一台戲.
係比天使同埋人睇.
係有好多唔同ge 人.
有好多唔同ge .
受造ge 黎到去睇.
即係呢個比我地一個.
新ge .
資訊或者係新ge .
一個提醒啦.
咁另外一樣ye 呢就係.
呢度講呀第九節一開頭.
gwo 度講ge 就係話.
就像上帝把我們使徒.
中蘇呀命命ge .

$^{841}$列在末後.
好像死罪ge 囚犯.
la 呢度呢好多時D 人呢就會話.
末在列後.
咁即係保羅認.
自己細啦.
認自己少啦咁好多人就會係去.
聯想到就係話.
哦梗係啦因為.
計十二個.
使徒門徒.
黎講咁佢地.
gwo D 係.
最先接觸主ge 咁樣.
咁然後到保羅係.
大馬識路先至再.
接觸主咁佢咪計呢個.
咁ge 時間呀計先後.
次序呀咁樣咪.
排名最後囉咁樣.
咁黎到係去解.
咁呢個即係聽jor 好多年咁樣講.
咁但係呢係呢度想.
同大家分享一個.
即係一個角度.
就係講呢度講.
角度黎到係去了解呢個經文.
呢度講就係話.
成jor 一台戲.
我地好似睇.
一齣歌劇.
或者戲咁.
一台戲但係我地如果.
睇返其他ge 翻譯本.
有英文呀.
或者叫Theater.
係一個舞台.
舞台.
即係唔係話一個劇.
由九點.

$^{881}$八點公映.
九點半完.
咁樣一台戲完.
係講緊係一個ge 舞台.
咁我地.
有少少了解就係呢個舞台.
係一個.
公眾ge 地方.
同我地而家ge 舞台可能.
有D ge 唔.
咁尤其我地北美洲.
凍ge 地方D 舞台呢都係.
即係戲院呀.
咁樣即係可能.
係室內ge .
除非你去過D .
我zar 住架車睇住.
露天咁.
但係古代ge 舞台呢.
係一個.
公開ge 場面.
係.
公眾ge 地方.
係人人都.
可以睇到ge 公眾.
所以咩呀.
同埋如果我地.
根據返《使徒行傳》.
第十九章gwo 度講.
如果大家有印象ge 會記得.
即係講到gwo 度.
去到已忽所.
保羅ge 同伴.
咪比人地拉去.
因為同個即係打爛jor .
簡單黎講打爛jor 銀像.
個飯碗於是.
班班咁個銀像.
就發起D 人黎到.
去捉jor 保羅gwo D 同伴.

$^{921}$黎到係去一個地方.
就係舞台.
黎到係去公審佢地.
想去.
叫人去.
自死佢地.
叫地方官.
我地睇返就係話.
舞台唔係話我地純粹話.
lor 去尋找娛樂咁簡單.
睇出戲.
比廿蚊睇出戲.
哈哈哈哈咁樣ye 就返黎.
係古時ge 舞台就係話.
一個大家聚集ge 地方.
有少少似.
中國人ge 祠堂.
即係一個公眾ge 地方.
拉埋係去.
當中有好多判斷.
有好多ge .
要同.
大家gwo 個城或者係.
gwo 個鎮要去公開.
要去處理ge 事.
亦都係好多ge .
劇目.
都會係.
可能係一個.
鬥獸.
即係你無話.
單同話睇演員演出戲咁簡單.
係呢D ge 劇.
劇院裡面亦都會可能牽涉到.
一D 最後係要.
即係將呢D ge 人要去.
拉去.
死.
係互相ge 廝殺.
或者係去比野獸.

$^{961}$黎到係去.
阿熊又好獅子又好黎到係去.
殺死gwo 個情況.
如果我地.
用呢個咁樣ge .
背景黎到去睇呢個舞台ge 時候.
我地就知道.
呢度講.
保羅講到.
試圖係呢D ge 情況.
底下係末後ge 意思.
你唔會一開套劇或者.
一開始就dump D 人去.
殺死佢地.
梗係搞一輪.
演一輪戲.
佢地諷刺gwo D 死囚.
先至黎到係去.
最後先.
割佢地又好D 咩.
呢度就係話擺到末後.
似乎係同.
呢一班ge .
試圖需要面對死亡.
呢一個ge 情景.
係有一種咁樣ge 構連.
唔係一件.
純粹娛樂ge 事咁簡單.
我地睇到呢個時候.
係用一句講.
唔係講玩ge .
係一件.
嚴肅ge 事.
亦都因此係咁樣.
我地都係去睇.
保羅就係講到.
作為試圖ge 佢地.
就係有一個最終.
唔係得著一個ge .
優越ge 地位.

$^{1001}$相反佢地係因為.
呢個信.
因為佢地忠於呢一個.
道.
忠於呢一個十字架同基督ge 道.
佢地黎到係去.
最終可能都要面對死亡.
唔光彩ge 事情.
所以佢地咪.
所以保羅咪係呢一度.
係第十節開始咁樣講.
係第十節咁樣講.
我地為基督ge 緣故.
算是如絕的.
你們在基督裡.
倒是聰明的.
我們軟弱.
你們倒強壯.
你們有榮耀.
我們被藐視.
三組唔同ge 對比.
即係講到.
試圖自己點樣點樣.
即係又.
又.
遇著.
軟弱.
同埋呢一個被藐視.
相對於哥林多教會.
ge 人黎講佢地聰明.
強壯又有榮耀.
咁我地其實呢個可以歸納返.
之前gwo 九節講ge .
哥林多教會ge .
佢地呢一句信仰係令到.
佢地更加ge 圓滿.
覺得自己選擇得正確.
選擇得醒我乜都做得到.
又重建返個城.
佢又有份又成為.

$^{1041}$軍事貿易ge .
大中心.
冇人睇唔起佢.
咁而家.
我信宗教.
唔係我要救恩.
而係基督教.
都因為我信而韜光.
但係相反.
呢一班ge 試圖.
佢係因為呢一個.
被世人捍衛如注ge 十字架.
自己都變得如注.
自己都變得軟弱.
自己都係被藐視.
一個咁樣ge .
強烈ge 對比.
世界對成功.
對出類拔萃.
ge 標準.
同有基督生命ge .
人所睇ge .
優越.
所睇ge 會得著.
上帝ge 讚賞.
係兩件唔同ge 事.
我地成日聽見有D 基督徒講.
我要.
為主做見證.
我要係我ge 專業裡面.
要話俾人地知.
基督徒都可以做到好好.
咁我想講一樣ye .
我想問一個問題.
人地係因為你做得個專業好.
你都係信主.
還是人地真係聽見基督ge 救恩去信主.
我係一個.
成功ge .
專業人士.

$^{1081}$係因為人地睇見你係成功ge .
專業人士去信耶穌.
可能一個觸發點.
係去俾佢去.
去探討.
但係個問題.
還是睇見真正.
呢個十字架.
同埋基督.
俾佢所吸引.
去盡心去認識住呢.
呢個我地要去.
即係.
我地要俾.
呢個寶來提醒.
寶來係呢度甚至講.
呢度第十一節.
呢度最結尾.
第十一節至第十三節.
直到如今.
我地還是又饑又餓.
又赤身老體.
又捱打.
又沒有一定住處.
並且勞苦.
親手做工.
被人咒罵.
我地就祝福.
被人罵我地就忍受.
被人毀破我地就善勸.
直到如今.
人患.
把我地.
看作世上的污穢.
萬物中的渣滓.
大家留意十一節.
直到如今.
呢度講緊ge 就係.
保羅唔係話.
唔係講緊D 舊屎ge 事.

$^{1121}$係直到佢寫呢段書信.
比哥林多.
前書ge 時候.
都係繼續係發生緊.
唔係話.
老人家過jor .
某個年紀 退休.
回想過去前塵往事.
發表ge 感觸.
就算到係呢一刻.
保羅都仍然.
係係呢一個ge 環境裡面.
係受著呢D ge .
被人捍衛.
被人譴責.
軟弱.
被藐視ge 事.
去度日.
但係仍然係.
繼續去工作.
仍然係去.
為主的道.
做事.
但係.
佢地有所不同ge 就係.
好多人.
去講.
保羅應該係一個.
係度一路講.
好負面ge ye .
但係我地要睇返第十二節.
係呢D 咁百般ge .
困難ge 裡面.
保羅仍然係有工作.
仍然係係呢個奴婦裡面.
仍然係去.
祝福人.
仍然係呢個ge .
忍耐.
係度係去煽軒人.

$^{1161}$即係第十二同十三節.
保羅唔會因為呢D ge .
苦.
繼續將呢D ge 苦.
去無限化.
擴大.
有D 人有時講見證.
強調自己好慘.
賣慘賣可憐.
相反佢係呢D ge .
環境裡面話比人知.
仍然係主ge 裡面.
佢知道自己.
係要向主忠心.
要使人更加.
要使人明白到十字架.
基督並拒ge 十字架ge 道理.
於是佢願意裡面.
祝福願意裡面.
繼續ge 勸人.
願意ge .
去係呢D 事.
忍受,唔係出於.
人gwo 種ge 意志力.
唔係出於.
自己擁有.
一個所謂道德高地.
黎壓倒人.
睇返前文後理.
我明白到佢因為明白到.
基督ge 十字架.
並佢ge 救恩.
於是有呢個咁樣ge .
動力.
去明白到.
恩典係真正ge 動力.
咁係呢度.
我地睇見呢一個ge .
舞台唔係一個純粹.
一個ge .

$^{1201}$娛樂觀眾ge 一個地方.
相反ge 就係話呢度.
睇到基督.
ge 奧秘ge 執事.
同埋佢ge 十字架.
受死.
試圖願意.
行係呢D 道上.
呢個.
當然我地.
唔能夠話lor 一卷聖經.
解釋一卷聖經,但我地.
都係去想返.
肥納比書第二章.
第六至九節.
講到基督ge .
心.
我地一一去望一望.
肥納比書後少少.
肥納比書第二章.
呢個.
我呢個版本佢好似用.
詩咁樣黎到去寫.
好多人話.
傳聞D 都係當時流行.
ge 一種ge 歌.
咁我地又睇.
肥納比書第二章.
第二章.
第六至九節.
他本有.
呢個講到基督呀.
理門呀,應該睇埋第二節.
理門當以基督的心為心.
他本有上帝的形象.
不以自己.
與上帝同等.
為強奪的.
反倒虛己.
取了奴僕的形象.

$^{1241}$未有人的樣子.
既有人的樣子.
就自己謙卑.
全心信服.
以至於死.
且死在十字架上.
所以上帝將他.
升為至高.
又賜給他那超乎萬名之上的名.
叫一切在.
天上的.
地下的.
和地底下的.
因耶穌的名.
而無不口稱.
耶穌基督為主.
使榮耀.
歸與乎上帝.
當然不能用.
一卷聖經.
解到另外一卷聖經.
雖然都是保羅寫.
寫作的對象,時分,時期.
都有不同.
目標可能都會有少少不同.
但我們看見都有一個pattern.
簡單來說就是先降.
比後升高.
用這個pattern.
來看那不少.
前書我們剛剛看過的經文.
好像使徒.
他們保羅說到.
他們使徒是.
很卑微.
被人看為愚蠢.
被世界唾棄.
但這正正不是.
基督的路嗎.
使徒.

$^{1281}$作為基督的.
追隨者.
只不過走回.
他的主要走的路.
保羅亦有提及.
這些路不是.
白白的走.
主會.
因著各人的忠心.
來賞賜.
所以我們有看見.
保羅在他.
寫過那麼多.
提摩太后書之後.
有觀念.
傳承.
我們很多時候.
不求無功但求無過.
信了主之後.
上天堂有個位留給我.
天堂有個位留給我.
這句很熟悉.
幾年前聽人說過.
不是.
等那時.
而是我們信了主之後.
認識了十字架的道理之後.
我們繼續在這件事上.
為主忠心.
希望得到主的獎賞.
縱然世界唾棄我們.
縱然世界不要我們.
嫌我們這樣.
嫌我們那樣.
最終我們要得到上帝.
對我們的讚賞.
這才是.
重視上帝的生命.
今天和大家分享的.
就在這裡.

$^{1321}$這段經文讓我們看到.
上帝對僕人的要求.
不是要他很厲害.
很有智力.
要有不知多少個博士學位.
碩士學位.
院士什麼的.
又不是要他強壯高大帥.
身高180厘米.
然後要怎樣.
又不是要他有些名望.
真正要求基督.
真正要求的.
是我們能夠.
試圖能夠對基督.
和十字架.
去忠心.
這個.
映照我們今天坊間.
對教會領袖的.
要求.
究竟有什麼不同呢.
我們的領袖我們要求他有能力.
又有口才.
又要有行動力.
最好走出來又帥.
又有魅力.
這些都是好的.
有很多幫助.
但是都不能夠取代.
向基督和十字架.
信息忠心的.
時尚高度取代.
今天我們的教會是怎樣呢.
我們要在這裡講.
看完這段經文.
可能都會心想.
這些只是關於試圖.
形容試圖的事.
我是一個評論圖.

$^{1361}$關我什麼事呢.
這些都是.
給牧師聽.
給教會主席聽.
關我什麼事呢.
下次留待.
第四章十四至二十一節的時候.
再和大家看.
究竟這些事.
關我們什麼事呢.
讓我一起祈禱.
親愛的天父.
給我們能夠在.
一同有敬拜的時間.
也讓我們明白.
你是那位.
唯一那位我們值得.
敬拜的主人.
求主你幫助我們.
今天我們在這個世代有很多價值.
有很多觀念.
告訴我們.
什麼是對什麼是錯.
什麼是好什麼是壞.
去分曉我們.
但是其實我們.
有多少.
是用你的眼光.
來看事物呢.
求主你幫助我們.
世界對領袖.
對帶領的人.
有不同的要求.
和我們.
在教會裡.
有什麼分別呢.
我們所要求.
我們看到.
保羅說到.
你對人的要求.

$^{1401}$是要求是忠心.
我們有多少是忠於主你的道理.
我們有多少是.
信了主之後.
仍然去尋找.
主你所需要的事.
求主你引導幫助.
讓我們能夠繼續.
在前面的道路上.
高,有你的指引.
祈禱是奉主耶穌基督的名.
求,阿們.
\newpage



\section{出埃及記 16:35}
\label{sec:032022}
\textbf{聖經人聖經事:以色人吃嗎哪四十年}~~~~黃克勤牧師博士
\newline
\newline
連結: \href{http://www.church.com.hk/acms/content.asp?site=cdc\&op=show\&type=product\&code=032022\&layout=sermon}{\texttt{032022}} ~~~~ 語音日期: 2026-04-18
\newline
\newline
\hyperref[sec:032021]{\small{< < < PREV SERMON < < <}}
~
\hyperref[sec:index_chronic]{\small{[返順時目]}}
~
\hyperref[sec:032023]{\small{> > > NEXT SERMON > > >}}
\newline
\newline
出埃及記 16:35
\newline
\begin{longtable}{cl}
\hline
\hline
章節 & 經文 (和合本修訂版)\\
\hline
16:35 & \begin{tabularx}{0.7\textwidth}{X} 以色列人吃嗎哪共四十年，直到進入有人居住的地方；他們吃嗎哪，直到迦南地的邊境。 \end{tabularx} \\ \\
[1ex]
\hline
\hline
\end{longtable}
$^{1}$大家好,又讀聖經了.
《生命書》和大家一起看聖經人和聖經的事.
這次是以色列人吃嗎哪四十年.
這件事也不簡單.
其實你詳細想下去,想深一層.
真的不簡單.
因為四十年.
一年三百六十五天,一日三餐.
四十年上帝供應嗎哪給他們.
直到他們進入了應許之地.
嗎哪才停止.
這也是每一日的神蹟.
以色列人吃嗎哪四十年.
開始時是有一段故事.
今天跟大家講一下這個故事開始.
仍然是摩西亞倫.
帶著以色列人出埃及.
當然真正帶著的是神.
耶和華神在做的.
《出埃及記》第十六章.
這次我先讀完聖經.
一邊讀一邊講,然後再講重要的要點.
第十節,亞倫正對以色列傳會眾說話的時候.
你會看到,摩西沒有說是亞倫代表摩西說.
有兩個領袖,摩西和亞倫.
摩西說:我不認識,我口才不好.
神說:我叫你哥哥一起吧,亞倫口才很好.
如果我個人認為,如果摩西當日沒有推.
摩西就做了一個獨一的領袖.
不過現在,他覺得口才不好.
神說:我叫亞倫一起.
於是就有摩西和亞倫一起.
分了工作給他的哥哥.
也分了領導的力量給他的哥哥.
亞倫正對以色列傳會眾說話的時候.
他們向曠野觀下看.
他們是指以色列傳會眾.
說的時候不料.
耶和華的榮光在雲中顯現.
嘩!一看就覺得是一種感覺榮耀.

$^{41}$有時候我們也會發生.
有些人拍攝天空的雲的時候.
就發覺有些光線在雲的縫隙上射下來.
嘩!很漂亮,好像上天跟我們說話一樣.
光在雲中射下來.
有線條的,直線的.
這些攝影人士就說.
這叫耶穌光.
很多東西都叫耶穌光.
現在微信也叫耶穌光.
已經傳說是這樣的.
榮光在雲中顯現.
上天跟我們說話一樣.
很厲害,讚美主.
如果你看到耶和華的榮耀顯現.
你的腳有點軟,想跪下.
心裡很敬畏.
耶和華曉如摩西說.
我已經聽見以色列人的怨言.
你告訴他們.
到黃昏的時候你們要吃肉.
早晨必有食物得飽.
你們就知道.
我是耶和華,你們的神.
這裡開始看到了.
以色列人向神又發怨言.
他們發怨言不直接向神說.
向領袖說.
摩西阿倫,你是不是想餓死我們.
昨天才說.
我們在埃及那時候.
是坐在肉鵝旁邊生活.
不知有多少肉吃.
出來這裡什麼都沒有.
想餓死我們.
想我們餓死在曠野.
埃及沒有墳墓.
說得很難聽.
很尖酸黑薄.
阿倫就跟我們聊天.

$^{81}$說話的時候.
神的榮光顯現.
透過摩西說.
透過摩西不是透過阿倫.
因為摩西才是上帝選擇的人.
是最直接選擇的人.
好了.
到黃昏.
你們的怨言我聽到.
到黃昏你們會吃肉.
早晨也有另有食物吃得飽.
你們就會知道.
我是耶和華,你們的神.
到了晚上.
真的有鵪鶉飛來.
他們說沒有肉吃.
哎呀沒有肉吃.
這樣就餓了.
鵪鶉飛來多到怎樣呢.
遮滿鳥營.
早晨遮滿鳥營.
晚上大家出去捉鵪鶉.
然後拿剪皮.
剝皮.
燒鵪鶉.
火燒烤鵪鶉吃.
早晨在營的四圍地上有露水.
正常啦.
露水上升之後.
不料野地面上有如白霜的小原物.
咦.
一塊塊圓形的東西.
以色列人看見不知道是什麼.
就彼此對問說.
這是什麼呢.
摩西對他們說.
這就是耶和華給你們吃的食物了.
耶和華所吩咐的是.
你們要按照個人的飯量.
即每人吃多少.

$^{121}$為帳篷裡的人.
按照人數收起來.
收起來即是說.
撿那些小原物在地上.
各拿一鵝梅爾.
不要拿多.
以色列人就這樣行.
有多收的有少收的.
有些人拿多些.
有些人撿多些.
及至用鵝梅爾量一量.
咦 這麼奇怪.
神說一人一鵝梅爾.
跟著你家裡有多少人.
就拿多少鵝梅爾.
有些人就拿多些.
有些人就拿少些.
結果回去一量.
都是一鵝梅爾.
五個人就五鵝梅爾.
沒有多.
多收沒有餘.
少收沒有缺.
個人按照自己的飯量收取.
個人按照自己的飯量收取的意思.
就當每個人都吃一鵝梅爾.
摩西對他們說.
所收的不許什麼人留到早晨.
今天的飲食.
今天吃了它.
然而他們不聽摩西的話.
來中有留到早晨的.
哎呀 誰知生蟲變臭了.
摩西 你看那些東西變生了蟲爛了.
摩西便向他們發怒.
叫了你們不要留到明天.
那你又留到明天.
生蟲發臭了.
他們每天早晨.
按照個人的飯量.

$^{161}$一鵝梅爾收取.
白色的小圓餅.
就會消失.
消化了 用光了.
到第六天他們就要收雙倍的食物.
那些小圓餅要收多一倍.
每人兩鵝梅爾.
因為到第六天.
每人兩鵝梅爾.
會中的官長便來告訴摩西.
摩西對他們說.
明天是星期日 安息日.
那摩西便說.
耶和華是這樣說的.
為什麼要收兩鵝梅爾.
因為那些代表來了.
明天是聖安息日.
是向耶和華守的聖安息日.
你們要敲就敲.
煮就煮 隨你了.
所剩下的都留到早晨.
即是說今天星期六.
叫我們收兩份.
那些官長便來問.
是不是收兩份 為什麼收兩份.
那你爛的.
摩西便說.
明天是聖安息日.
那你今天收了回來.
你喜歡怎樣吃就怎樣吃.
留一些給明天.
剩下的要留到早晨.
即是明天安息日.
其實因為明天是安息日.
是不能收的.
他們便照摩西的吩咐.
留到早晨.
這麼有趣.
留到安息日早晨.
也不臭 裡頭也不生蟲.

$^{201}$摩西便說.
你們今天要吃這些.
因為今天是向耶和華守的安息日.
你們吃這些便可以了.
其實即是說.
到第六天.
便要收多一倍.
於是今天吃了第六天的.
明天安息日.
便不用收.
於是便不爛.
這個神蹟也很偉大.
就是第一天至到第六天.
收的只是每人一個人頭一鵝眉耳.
夠一天吃.
我收兩天.
明天放假不用收.
會生蟲 會臭.
不行.
你也要倒掉.
一曬便溶掉.
但是到第六天便可以收第七天.
便是安息日.
大家會知道了.
不爛.
你們安息日不要出去填耕.
有些人除了撿這些.
一鵝眉耳的白色圓餅.
他們便吃.
有些人也會出去捕捉.
有沒有野鳥.
有沒有動物可以捕捉回來吃.
野果也會.
其實也沒有什麼.
填東西你們找不到.
一定是吃這些.
第六天可以收取.
第七天是安息日.
那一天便沒有小圓餅.
而你出野外也捕捉不到東西.

$^{241}$第七天百姓中竟然有些人.
又是不聽話.
不守安息便出去收.
看看有沒有小圓餅.
沒有.
什麼也找不著.
很有趣.
真的沒有.
又說對摩西船.
你們不肯守我的戒命和律法.
要到什麼時候呢.
其實戒命和律法現在還未給.
即是說.
第一至六天.
出去捕捉.
早上不要晚起.
捕捉了.
然後第六天再捕捉一份安息日用.
誰知道那些人又不聽.
有些人捕捉多些.
平時弄得臭了.
有些人雖然.
第六天捕捉多了一份.
但安息日還要出去看看有沒有.
叫了你不要.
這麼簡單的事.
那些以色列人都不跟著做.
只是給了這麼小的戒命和律法.
他們都不守.
你們這麼不聽話.
要到什麼時候呢.
李滿漢又說.
既將安息日賜給你們.
所以第六天賜給你們兩天的食物.
第七天個人要住在自己的地方.
不許什麼人出去.
住在自己的棚裡.
帳篷裡.
不要出來.
在裡面守安息日.

$^{281}$休息.
於是百姓第七天就安息了.
其實先教他們第一件事就是安息日.
可能他們以前在埃及圍爐有沒有安息日.
沒有的.
一個星期可以七天供的.
現在出埃及之後.
開始第一件事教.
就是安息日.
不要走路.
不要出來找那些小圓餅.
要怎樣呢.
在帳篷裡休息.
就是這樣.
也不要出去.
不要出去.
第七天安息日是沒有的.
又有人出去.
借食物.
以色列家叫做嗎哪.
樣子像芫荽子.
樣子像芫荽.
我也不太明白是怎樣.
顏色是白色小圓餅.
滋味就好像材物的薄餅.
味道甜甜的.
好像蜜糖做出來的餅.
薄薄的.
所以叫薄餅.
很多塊小的圓形白色薄薄的.
有蜜糖味的餅.
樣子像芫荽子.
我就不太明白.
葉子又不太像.
芫荽又不太像.
因為說是小圓餅.
於是第八聖第七天安息了.
借食物以色列家叫做嗎哪.
樣子像芫荽子.
顏色是白色的.

$^{321}$滋味就像材物的薄餅.
可能那個餅叫小圓餅.
但其實可能像芫荽的葉子.
不一定是正圓的.
好像一片葉子.
顏色是白色的.
味道就像混合了蜜糖的薄餅.
摩西就說.
所吩咐的是這樣的.
要將一碗俄梅爾嗎哪.
留到世世代代.
使後人可以看見.
我當日將你們領出埃及地.
再曠野所給你們吃的食物.
耶和華說的.
特別給我收起一俄梅爾的嗎哪.
那些嗎哪是不會爛的.
不會變的.
是我的神跡騎士.
令到它變成這樣.
後人就會知道.
我耶和華帶以色列民族出埃及的時候.
就吃這些了.
你看到這是一個很大的神跡.
摩西亞倫說.
你拿一罐子滿一俄梅爾的嗎哪.
傳到耶和華面前.
要留到世世代代.
耶和華吩咐摩西亞倫就這樣行了.
就將嗎哪放在法季前存留.
法季有些疑問.
當時不知道有沒有法季.
應該有的.
如果不是不會這樣寫.
他們有一個季.
特別神聖的地方.
要來放一些重要的東西.
就將嗎哪放在季前存留.
這個季很有可能是一個四方形的季.
有木桿可以抬.

$^{361}$很有可能.
上帝已經叫了摩西這樣做.
如果沒有這個法季兩個字出現.
未必像聖殿或是會幕裡的法季.
樣式可能未發展好.
但法季的頂.
就叫做施恩座.
上帝坐在那裡.
施下恩典.
以色列人在嗎哪共四十年.
偉大的神蹟.
還有一個人居住之地.
就是迦南的境界.
摩西解釋一下.
阿梅爾是伊法的十分之一.
伊法就是二十二公升.
今天就讀完經文.
很快跟大家解釋一下.
很餓啊.
沒有肉吃.
我們以前在埃及.
在肉窩旁邊.
有很多肉吃.
如果我是摩西就說.
你不要吹太大聲.
以前你們是奴隸.
哪有在肉窩旁邊任食放題.
自助餐嗎?.
其實有肉吃都很開心.
你記得嗎?.
不要說太大聲.
好了.
不要再說了.
現在神就說.
我聽到.
當你受到一些壓力委屈的時候.
你要知道.
其實神知道我聽到他們的怨言.
那些人這樣說話.
我聽到的.

$^{401}$所以我會做事.
做什麼呢?.
到黃昏我會安排肉給他們吃.
清晨我會有特別的食物給他們.
讓你知道我是耶和華.
是你的神.
晚上就是鵪鶉飛來了.
鵪鶉是野鳥.
很多.
片滿鳥營.
吃肉吃到飽.
哇,肉吃到飽了.
早晨就有露水.
露水之後就有小原物.
這個原.
你不要想著.
哇,真的很原.
可能因為它像芫茜籽.
可能像芫茜的葉.
一種香料的葉.
一塊塊的.
圓圓的.
那些人就問.
這麼有趣的是什麼?.
摩西就說.
全部拿起來吃吧.
每人一鵝眉耳拿回家.
吃?可以吃的嗎?.
這些是什麼?.
拿著什麼?.
這是什麼?.
什麼?.
翻譯成英文.
加了一個A.
M-A-N-A.
或者叫Mala.
其實就是Mala.
Mal不是讀Man.
Mal就是Mala.
其實就是.

$^{441}$什麼來的?.
這樣的意思.
究竟這是什麼?.
大家就問.
什麼來的?.
Mala就跟中文有點像.
Mala.
所以Mala.
吃什麼都不知道的東西.
吃了四十年.
這是上帝的工作.
每天吃了四十年.
三百六十五天.
另一個神蹟的地方在哪裡呢?.
當中有一個很複雜的神蹟.
是什麼呢?.
各人照自己的飯量拿的.
這個故事真有趣.
就是每人計人頭一鵝眉耳.
鵝眉耳.
Omer.
Omer是什麼意思呢?.
就是兩公升.
兩公升多不多呢?.
也不是很多.
如果我們現在在超級市場.
那些一公升的牛奶.
小房子紙盒那些.
兩升一天.
這樣很少的.
如果大隻一點就不夠吃.
原來可能都飽肚的.
兩公升一天.
照著自己的飯量拿.
每人一鵝眉耳.
以色列就是這樣.
多收的有少收的有多收.
有些人收多些人收少些.
誰知花的一糧是一樣的.
多收的沒有餘.

$^{481}$少收的沒有缺.
第二點要說的是什麼呢?.
就是.
看回第一點是什麼呢?.
不是說完.
大家忘記了第一點是什麼.
第一點就是.
當你有困難的時候.
耶和華的榮光會顯現.
幫我們解決問題.
神在百姓的餓的時候.
神向他們顯現.
給他們解決了問題.
給他們食物.
然後到第二點.
就是上帝在這裡做了一個神蹟.
是什麼呢?.
就是多收的沒有餘.
少收的沒有缺.
即是說神的供應是夠的.
我貪心拿多些.
他又不會特別多.
拿少了.
又不會拿少了.
起床遲了一點點.
因為白天一曬.
他溶了就沒有了.
所以一定要早起來拿.
少收有沒有缺.
是按照你的飯量去撿.
撿多些撿少些.
差不多吧.
撿到回去.
一糧都是兩.
一個是俄梅爾.
一個是奧梅爾.
好了.
這個神蹟很大.
大家記住還有一個神蹟是什麼呢?.
就是這個關於瑪雅的神蹟.

$^{521}$奇妙之處就是.
你.
很早起床.
我不想這麼早起床.
今天撿多些.
明天不用撿了.
不行.
這樣就會生蟲變臭.
是每一天都要拿的.
是我們日用的飲食.
而且每天早上拿的.
每天早上拿的.
這個叫我們每天早上起床.
就要撿糧食.
拿好你的飯量.
一開始發熱就消化了.
沒有了.
溶了.
蒸發了.
很有趣.
一定要記得.
早上就要拿了.
很多傳道人會說.
好像我們的靈修.
要早上就要做了.
不然.
遲些就不好了.
到第六天就要拿兩倍.
因為第二天是安息日.
安息日.
就拿兩天.
這個魔式.
自己加的.
其實你拿了的.
今天第六天.
你烤也可以.
煮也可以.
這個是什麼呢.
很有趣.
又是一個神蹟.

$^{561}$媽.
馬蘭媽.
這個是什麼呢.
真的要問一下是什麼.
是怎樣的.
火燒.
燒烤.
串燒.
又可以.
薄薄脆.
如果我要煮.
煮是什麼呢.
用水.
又可以.
煎也可以.
又可以.
焗也可以.
可以.
我不知道.
中國人說.
燜媽.
可以嗎.
我猜是可以的.
這裡說的烤和煮.
任你怎樣做都可以.
剩下的可以留到明天.
因為今天.
明天是安息日.
不能收拾.
多收拾一點.
或者今天有剩留給明天.
都可以.
都不壞.
也不臭.
裡頭都沒有蟲.
從今天起.
就要吃這個.
你填東西.
明天.
安息日你找不到.

$^{601}$收拾不到.
今天.
就要多收拾一倍.
第六天可以收多一點.
明天安息日沒有了.
百姓中.
真的不理你.
就出去找找.
結果什麼都找不到.
唉.
真是.
叫你多收拾一點.
就多拿一點.
第六天.
就吃完了.
然後第七天.
在帳幕裡休息.
不要走來走去.
這樣你都不聽.
又出去想收拾.
唉.
神就說.
唉呀.
這班人不肯守我的律法誡命令.
要到什麼時候呢.
第三點要說的就是.
這班人真的不聽話.
我不想.
神會跟你說.
唉呀.
為什麼你這麼不聽話.
我也不想神跟我說.
為什麼你這麼不聽話.
唉 頭鈴鈴.
你這麼不聽話.
怎麼會這樣.
這就是這個意思.
為什麼你這麼不聽話.
為什麼這樣.
你什麼時候才學到.

$^{641}$要到什麼時候呢.
我想這個就不要傷到上帝的心.
他叫我們這樣做.
就跟著這樣做.
他已經叫了.
你要留在帳幕裡.
不要出去.
好了.
到提案式.
食物以色列家叫做嗎哪.
是什麼來的.
再說一次.
元衰子.
白色.
味道很像薄餅.
是混合了蜜糖的.
這樣都挺有趣.
第七天.
他們都是這樣吃的.
摩西就說.
神叫摩西.
收一罐做永遠的紀念.
這一罐永遠的紀念在哪裡呢.
後來神叫他們做藥櫃之後.
藥櫃裡面就放了十屆重造的兩塊石板.
那支杖.
還有一罐嗎哪.
這一罐嗎哪.
由現在開始.
一直放到藥櫃裡.
都沒有破.
如果藥櫃繼續在.
就繼續看到那罐嗎哪.
沒有破.
嘩,多厲害.
這又是另一個關於孖吾的神蹟.
嗎哪的神蹟.
藥櫃那罐是不破的.
法櫃那罐.
這罐一盒煤油放在裡面.

$^{681}$於是亞倫就做了.
好了.
最後一點.
以色列吃嗎哪四十年.
嘩,這個神蹟大.
四十乘三百六十五日.
都很多日子吃嗎哪.
直到一日加難就沒有了.
這也是一個很重要的神蹟.
阿梅爾就是告訴我們.
兩星.
大概是一個熱花.
熱花的十分之一.
二二星.
兩星的十分之一.
跟大家分享到這裡.
以色列人是吃嗎哪四十年的.
而且是每日的.
六日加上第七日.
而且怎樣煮都可以.
而且又有營養.
又可以頂肚.
又吃四十年的.
有一罐放在這裡又不破.
星期六那罐又不破.
平時那罐拿多一點就破了.
遲些出去就沒有了.
嘩,這個慢什麼.
是一個神蹟.
所以我結束的時候要說的是.
上帝有時的工作真的很大.
很厲害.
到一個地步我們只會說一個字.
慢.
嘩,神啊,這是什麼.
這是什麼神蹟.
是什麼慢.
什麼什麼慢.
什麼.
漢語什麼.

$^{721}$希伯來文慢.
就是嗎哪.
就是這個意思.
所以當你有困難.
你記住嗎哪.
你會記得上帝是這麼大能的供應者.
在同期徒.
主剛才亂世經很長.
不過讓我們記得.
慢.
這些是什麼呢.
有時我們都不知道原來這個就是.
供應我們四十年.
每一日吃的食物.
這個神蹟這個供應.
實在是奇妙無比.
是主自己大能的彰顯.
是你向以色列人當日顯現你自己.
你從榮耀中彰顯你的大能.
主你忍耐人那種不聽話.
但是每日仍然供給他們.
多謝你的愛,多謝你的照顧.
多謝讚美你的大能.
禱告奉耶穌名求.
阿們.
\newpage



\section{出埃及記 2:23-25}
\label{sec:032023}
\textbf{改變現狀的禱告——嘆息哀求}~~~~馮永和博士
\newline
\newline
連結: \href{http://www.church.com.hk/acms/content.asp?site=cdc\&op=show\&type=product\&code=032023\&layout=sermon}{\texttt{032023}} ~~~~ 語音日期: 2026-04-18
\newline
\newline
\hyperref[sec:032022]{\small{< < < PREV SERMON < < <}}
~
\hyperref[sec:index_chronic]{\small{[返順時目]}}
~
\hyperref[sec:032024]{\small{> > > NEXT SERMON > > >}}
\newline
\newline
出埃及記 2:23-25
\newline
\begin{longtable}{cl}
\hline
\hline
章節 & 經文 (和合本修訂版)\\
\hline
2:23 & \begin{tabularx}{0.7\textwidth}{X} 過了許多年，埃及王死了。以色列人因做苦工，就嘆息哀求；他們因苦工所發出的哀聲達於神。 \end{tabularx} \\ \\ \relax
2:24 & \begin{tabularx}{0.7\textwidth}{X} 神聽見他們的哀聲，就記念他與亞伯拉罕、以撒、雅各所立的約。 \end{tabularx} \\ \\ \relax
2:25 & \begin{tabularx}{0.7\textwidth}{X} 神看顧以色列人，神是知道的。 \end{tabularx} \\ \\
[1ex]
\hline
\hline
\end{longtable}
$^{1}$親愛的主,求你在這個時間裡面,.
無論我們年紀有多大的差距,.
求你讓我們一起進入你更深的同在裡面,.
主,求你的說話,甦醒我們,.
因為你是大有能力,聽人禱告的神..
所以當你禱告的時候,求你的靈魂走在這裡,.
主,感動我們的心,.
求你歸到我們的面前,.
進入你的說話當中,.
求你甦醒,用你的說話甦醒我們每一個人,.
聽我們的禱告,奉耶穌基督的名,求..
阿們.謝謝..
我想如果我說一些名字出來的時候,.
多數你都不會認識的,不如我說吧..
其實,我想大家都認識中神,對嗎?.
中神即是中國聖經,十號院,是嗎?.
我全名都不太認識..
就是近時,.
中神其實是第一任的院長,.
彭國輝牧師,其實很推崇,.
對禱告有很深認識的方知,.
這個人其實是很多年前的,.
不過在聖經,如果你是神學界的,.
其實一定會認識他,.
叫E.M.Barnes,方知,.
他曾經說過一個很出名的話,.
這個人對禱告是一個很有經驗的人,.
他說有效的禱告一定是一個改變生命的祈禱,.
但其實如果我們看回的時候,.
其實如果我們仔細看聖經的時候,.
就會發覺聖經裡面,.
是記載了很多有效禱告的狀態,.
很多時候可能我們不太認識,.
其實詩篇充滿著這樣的狀態,.
如果你不去到這樣的地步,.
經歷不到這些的時候,.
你是不會真正明白詩篇的,.
詩篇第23篇,第46篇,90篇,.
我想得到很多很出名的詩篇的時候,.
其實都是一些禱告的狀態,.

$^{41}$其實它形容了一些狀態,.
怎樣去改變那時基督徒的生命,.
同時我們也看到舊的歷史裡面,.
無數的寄述,其實他們都記載了,.
一些改變生命有效禱告的不同的狀態和層次,.
但是當我十幾年前回到香港的時候,.
當我開始在浸神教神學,.
做他的副教授的時候,.
助理教授的時候,.
我就開始發現,.
其實那時候我發覺那些學生是不太懂的,.
然後呢,.
總之我覺得震撼的東西就是,.
我在教會,也有幾多大的教會的時份過,.
我發覺原來很多教會其實都不太知道,.
他們不太知道原來,.
真正的禱告是可以去到什麼地步,.
和禱告有什麼能力的,.
所以其實呢,.
我發覺更明顯的時候,.
我都不太為意,.
知道的時候可以慢慢學,.
但我發覺更加重要的就是,.
當2019年,.
我們有社會事件,.
然後經濟有不景,.
一段很長的時間,.
香港有一個很低層的時候,.
很多人在很困難的時候,.
在這六七年裡面,.
我發覺不是很多人,.
是知道怎樣進入這些狀態,.
和去運用神給我們的能力,.
和權柄去禱告,.
這使得我很震撼,.
因為這個其實我覺得是基督徒生命裡面最重要的一個環節,.
所以,近年,.
我想特別是去年開始,.
我開始講多些禱告有關的事,.
其實禱告是一個很深的subject,.

$^{81}$其實可以是一個十堂,.
每堂三個小時的course,.
所以講的,.
所以其實,.
以45分鐘的時候,.
根本講不到什麼,.
我想和你分享,.
其中禱告一個的狀態,.
這個狀態我覺得重要的,.
因為你們人生裡面,.
如果你們追求神,.
是一定會經歷過的,.
這個就是聖經裡面較為出名的狀態,.
它的名字叫彈色哀求,.
我想用一段經文,.
很短的,.
去和你講這個狀態,.
如果你看出埃及記的時候,.
我只讀三節的聖經,.
很短,.
第23至第25節,.
當中裡面,.
你聽我讀,.
就是過了很多年,.
埃及王死了,.
這個講摩西的故事,.
和以色列的故事,.
以色列人因為作苦工,.
就彈色哀求,.
他們的哀聲達於神,.
神就聽到他們的哀聲,.
就紀念他和亞伯拉罕以撒雅各所立的約,.
神看顧以色列人,.
也知道他們的苦情,.
第23節裡面,.
講到彈色哀求,.
其實這是一個很特別的希伯來文,.
如果在這個希伯來文,.
如果你認識原文聖經的時候,.
你就會發覺,.

$^{121}$這個原文,.
在舊約聖經裡面,.
一共出了大約100次,.
其實全部都是形容個人,.
或者群體,.
祈求的狀態,.
但是如果仔細去研究,.
彈色哀求的運用的時候,.
特別是當彈色哀求的對象,.
是神自己的時候,.
你就會發覺,.
彈色哀求不是一般的祈禱,.
不是像你早上起床禱告,.
吃飯禱告,.
然後臨睡又跟家人禱告,.
然後為你的子女禱告,.
不是,.
而是什麼狀態呢?.
是當一個人,.
去認識到,.
有很大的困難淋到他,.
而他用出他的方法,.
都沒有辦法去處理這個困難,.
去到走投無路,.
去到絕境的時候,.
在強大的壓力下的時候,.
是由心底裡發出一個很自然,.
也很迫切的一個吶喊,.
如果你根據我們歷史書裡面的書卷,.
好像《時詩記》,.
《列王記》上下,.
《歷代節》上下,.
你就會發覺,.
當每一次以色列人被仇敵壓榨的時候,.
或者陷在絕境裡面,.
跟神彈色哀求的時候,.
是每一次,.
神都會幫助他們,.
是應運他們的禱告,.
拯救他們脫離仇敵的手,.

$^{161}$以及帶給他們有贖靈的復興,.
就好像有一條贖靈的方法,.
就是說,.
當你彈色哀求的時候,.
就有拯救,.
當你彈色哀求的時候,.
就會有復興,.
然後我們就會問,.
那什麼才叫彈色哀求呢?.
是不是當我們祈禱大聲一點,.
當我們帶著強烈一點的情緒去跟神說,.
神你一定要聽我的禱告,.
那就是彈色哀求了,.
當你將這100節的經文,.
分類研究,.
以及把它分類之後,.
其實彈色哀求基本上是有三個特徵,.
第一,.
彈色哀求是承認,.
神是我們唯一的希望,.
這是一個講義造難的狀態,.
也很老實說,.
從一個過來人的角度,.
是十分困難的,.
為什麼呢?.
因為當你有困難的時候,.
你想想你自己有困難的時候,.
你怎麼做?.
有困難的時候,.
你當然要解決困難,.
你父母和學校的老師,.
現在的學校通常都給你很多功課,.
給你很多壓力,.
你的子女就知道了,.
你的年代沒那麼重要,.
你的子女就慘了,.
他的子女不用想了,.
所以你發覺一件事,.
學校給你那麼多東西的時候,.
其實不是單單要小孩解決困難,.

$^{201}$要家長幫小孩一起解決困難,.
不然小孩會死硬,.
所以其實一件事,.
在我們很小的時候,.
其實我們已經被人訓練,.
一定要自己解決困難,.
我們中國人有個思想,.
就是天助自助之人,.
所以其實解決困難是天公地義,.
當我們在工作,在家庭,.
在人際關係遇到困難的時候,.
我們會想盡一切方法去解決困難,.
而這個社會雖然沒有明顯,.
但其實潛規則告訴你,.
你越有能力解決困難,.
你的工資就會越高,.
這是一個事實,.
你的職位也會越高,.
你坐在總經理的位置的時候,.
你就是一個整間公司的困難都可以解決的人,.
他就獎勵你坐在這個位置,.
是基本上是憑實力去拿得到的,.
這個實力是要很長時間的訓練,.
很艱苦,經過很多困難,.
累積過來,.
所以對很多人來說,.
其實當我們遇到困難的時候,.
我們是自己去解決,.
神不是我們唯一的希望,.
神是什麼呢?.
神是當我們完全用自己的方法,.
都解決不了困難的時候,.
不過是我們最後的希望,.
這個很不同,.
其實你不要覺得,.
嘩,真慘,.
其實這是一個很難達到的目標,.
但你要發覺一件事,.
就是如果你看回聖經的例子,.
你會發覺,.

$^{241}$聖經所說的,.
基本上一開始的時候,.
都不會覺得神是他唯一的希望,.
你看看以色列人,.
就是一個很明顯的例子,.
你看看這個例子,.
你知不知道以色列人,.
在埃及一共被奴隸幾年,.
他才開始嘆息哀求呢?.
如果你看看,.
出埃及記第一第二章,.
然後再拍攝《使徒行傳》第七章的時候,.
你就會發覺,.
摩西在未出世之前,.
他已經被奴役了,.
而摩西在八十歲的時候,.
他才開始嘆息哀求求神幫助他們,.
你知不知道以色列人,.
過了多少年才嘆息哀求呢?.
是八十年,.
超過了八十年,.
同一個狀態,.
他們在掙扎,.
過了八十年才開始嘆息哀求,.
你知不知道為什麼,.
他到那時候才嘆息哀求呢?.
聖經裡面最明顯說,.
因為那時候,.
近時奴役他們的埃及王死了,.
有新的埃及王起來,.
他那時候可能,.
先有想,.
我們盡力而為,.
我們已經成為奴隸,.
怕人來欺負,.
也沒有什麼辦法,.
只能盡力而為,.
在當中博一下好感,.
做多點事的時候,.
別人認同我有價值的時候,.

$^{281}$對我好一點,.
但發覺不是,.
做什麼都沒有用,.
然後終於等到了,.
這個埃及王終於死了,.
我做到他死了,.
終於有新的埃及王興起,.
會不會有些改變呢?.
原來沒有的,.
比近時更差,.
然後,.
其實在當中,.
是到了那個時間,.
以色列人才會發覺一件事,.
原來我做什麼,.
做多少,.
做得多辛苦,.
多默默的去做,.
有多大的貢獻,.
原來是沒有用的,.
原來是完全沒有用的,.
基本上是救不了自己,.
是直到那個時候,.
以色列人才會發覺,.
耶和華原來是他們唯一的希望,.
是不是耶和華之前沒有那麼大能力,.
現在能力大了呢?.
不是,.
而是因為人的認識,.
終於發覺他自己不行,.
其實,.
在我第一次接觸這些經驗的時候,.
是當我差不多和你一樣大的時候,.
我開始追求神,.
我開始去對付我自己的罪的時候,.
那,弟兄,.
你知道會有什麼罪的了,.
其實我也會有這樣的試探,.
有這樣的經驗,.
很多時候,.

$^{321}$當我很多時候,.
有一些的,.
那時候我和你差不多大年紀,.
那時候當我有些色情的意念進來的時候,.
我發覺,.
很難抗拒,.
很大的困難,.
我很想追求神,.
但我發覺常常都跌倒在這些罪裡面,.
然後,.
在當中我開始慢慢意識到,.
真的不行,.
然後,.
有一兩次的時候,.
我有些很特殊的屬靈經驗,.
就是有一次,.
我發覺,.
它試探強大到一個地步,.
我發覺我一定會跌倒,.
然後那時候,.
我在跌倒之前,.
我做了一個禱告,.
我說,.
神,.
救我,.
然後在這個禱告之後,.
我做了這個禱告之後,.
我發覺,.
突然之間,.
那個很強大的誘惑和勢力,.
是突然間消失的,.
是沒有了的,.
是直到那時候,.
我才開始明白,.
原來所有的addiction後面的受,.
原來有撒旦的勢力在裡面,.
所有我想到的addiction,.
這個是到了我再深入認識神,.
以及聖經的時候,.
開始知道,.

$^{361}$就是說,.
原來那些勢力,.
在當中你覺得無可抗拒的時候,.
只有當你承認是不行的,.
專心依靠神的時候,.
神是很神跡地將那件事拿去,.
然後你就會,.
那時候我第一次經歷的時候,.
我發覺,.
我出了一身冷汗,.
為什麼呢?.
因為我第一次看到有邪靈的力量在後面,.
這個我年紀差不多跟你一樣,.
你不可能的,.
如果不是的時候,.
上一秒還是,.
那個試探打了一秒,.
一定會跌倒的,.
下一秒完全消失的,.
你真的不可能,.
最後這件事不是一次,.
而是每一次當我這樣禱告的時候,.
它都這樣發生,.
然後我就明白一件事,.
當我知道,.
專心單單放下自己,.
承認自己不行,.
去依靠神的時候,.
我才開始經歷一件事,.
就是超乎尋常的改變,.
就會開始發生,.
所以其實真正有效果的刀是真的,.
是一個可以改變生命的禱告,.
接著我們再看看,.
這個彈色哀求,.
如果你看回聖經裡面的時候,.
第二個category的意思是什麼呢?.
彈色哀求是指向一件事,.
就是你是不自知的,.
要求神介入你的生命裡面,.

$^{401}$當你認識到,.
神是你唯一的希望的時候,.
當你看到你是無能為力的時候,.
當你在絕對的困境裡面,.
不能夠被拯救的時候,.
接著你的禱告就會開始改變,.
你的禱告開始改變,.
不會再是像你吃飯之前,.
那十五秒你說神啊,.
求你潔淨這些食物,.
使我有能力,有精神,.
今天走你讓我走的道路,.
藉著這些食物,.
不會的,.
晚上禱告就是神啊,.
多謝你今天幫助我,.
渡過這麼多的難關,.
讓我睡得好,.
一覺睡天光,.
不會再是這樣的禱告,.
你的禱告在當中會被改變,.
是一件很自然的事,.
當你開始意識到,.
當你是一個很絕對的困難,.
你發覺你自己不能夠得勝,.
但是你很迫切的時候,.
你就開始有突破,.
你開始會認識到一件事,.
當你認識到神是你唯一的幫助的時候,.
你會進入日以繼夜的禱告裡面去,.
你可能晚上睡不著,.
因為那件事壓在你的身上,.
無論在家裡,.
在你工作,.
在你當中對罪的對抗也好,.
壓到你呼吸不了的時候,.
你晚上醒了,.
你不是在當中立刻睡著,.
而是你暗暗的時候,.
在你家裡不為意的時候,.

$^{441}$在你當中下櫃,.
在你床邊,.
你進入禱告裡面去,.
你開始上班的時候,.
你搭地鐵,.
搭巴士,.
搭minibus的時候,.
你不是再看你的手機,.
而是開始的時候,.
你進入禱告裡面去,.
神是你幫助我,.
神是你拯救我,.
當你下午的時候,.
當公司有午飯時間的時候,.
你會發覺你不和那時候,.
你盡量不和你的同事一起吃飯,.
而是你盡快,.
很快吃完之後,.
你玩一塊沒人知道的地方,.
你進入禱告裡面去,.
神啊,你救我,.
這些禱告會慢慢組織起來,.
由開始很疏變成很密,.
因為你開始意識到,.
只有神才能夠幫助你,.
在當中裡面,.
這些禱告慢慢會形成一條很長的河,.
在當中裡面成為你的血,.
你的肉,.
你不可以分割的一部分,.
成為你的生命,.
你發覺你一直會這樣禱告下去,.
直到神介入你的生命為止,.
在《出海及記》第二章裡面的時候,.
如果你仔細看回聖經的時候,.
你會發覺,.
以色列人在困難裡面,.
是一直不止息的禱告,.
直到神賜下摩西,.
當你在《事詩》裡面,.

$^{481}$你就會看到,.
當以色列人在困難裡面,.
在仇敵的壓榨裡面的時候,.
他們是一直的禱告,.
直到神賜下以蜜和姬澱,.
這些事詩去幫助他們,.
接著你在《歷代刺》下面,.
你看到,.
希西加和先知爾賽亞,.
當亞述的西拿基納王,.
貴昆要殺難的獸,.
滅國好像是他唯一的選擇的時候,.
他們在當中是很切,.
不絕的禱告,.
直到師行神職騎士,.
一晚裡面,.
是可以消滅了亞述的軍隊18萬5千人,.
而拯救了耶路撒冷,.
免於滅亡為止..
在基督教的歷史裡面,.
有一個人叫David Blainer,.
他的詩跡是充分解釋了,.
什麼是改變生命,.
和不停求神介入的禱告..
在18世紀1740年代的時候,.
David Blainer是騎馬,.
超過了3000公里的路,.
走了3000公里的路,.
獨自在荒野裡面去住宿,.
為了傳福音給當時的印第安人,.
他當時已經有一種致命的疾病,.
其實就是肺癆,.
他當時是日以繼夜,.
和這種疾病去搏鬥,.
其實你知道肺癆的時候,.
你會呼吸不了,.
所以你會知道他多辛苦,.
大部分的晚上,.
他只能夠在寒冷的郊野裡面過夜,.
而且是沒有毛纖寶聯,.

$^{521}$而且絕大部分的時間,.
當他向印第安人傳福音的時候,.
他只有一個未信主,.
而且在上層出錯的翻譯去幫助他,.
他當時其實已經意識到,.
他根本就不能夠做些什麼,.
他根本是什麼都做不了,.
但因為他的信念是神,.
要他向印第安人傳福音,.
所以基於他對印第安人那種愛的時候,.
他開始進入日以繼夜的祈禱裡面,.
他單單祈求一件事,.
就是施行神蹟騎士,.
去介入這個絕對不可能的情況裡面,.
使他們得救,.
神就聽他的禱告,.
介入他的事奉裡面,.
很多人當時認為是完全不可能的神蹟,.
曾經有一個情況,.
當他為他翻譯的時候,.
是醉到完全不知道他自己在說什麼,.
前言也不對後語,.
但是聖靈的能力,.
是好像火焰一樣,.
和David Blainard是同才,.
在那次的聚會裡面,.
無數的印第安人認罪悔改,.
接受耶穌那一刀,.
做個人的救主,.
David Blainard在1747年,.
當他29歲的時候,.
死了,因為肺癆死了,.
他事奉當地的印第安人,.
只有四年的日子,.
很短的時間,.
你想想你去YAS多久,.
他真的去了四年,.
又靈召又是四年,.
而且大部分的時間他在禱告裡面,.
但是四年裡面,.

$^{561}$當他死的時候,.
教會,家庭的祈禱會已經set up了,.
而主日被遵守,.
最重要的是,.
這班印第安人本來是一班很頑蠻,.
很野蠻,.
完全是較為殘忍的一班人,.
但是其實在四年裡面,.
被改變成一班很敬虔,.
和有愛心的一群,.
所以當Jonathan Edward,.
美國那時候最偉大的復興家,.
他有一些人看了他的事跡的時候,.
他曾經有一個這樣的comment,.
就是說,.
有很多比David Blainer更加有才能,.
恩賜和智慧的人,.
不能夠留下這樣的痕跡,.
他們的工作,.
他們雖然很叻,.
很有恩賜,.
但是他們的工作像什麼呢?.
像沙灘上的竹蜻,.
很快就不見了,.
David Blainer藉著他熱切如火焰一樣的禱告,.
將他的工作刻在神的金剛石上面,.
將會永遠在世人和在神裡面,.
作為一個紀念..
當歎息哀求這個希伯來文,.
如果在先知書裡面出現的時候,.
就出現了第三個characteristic,.
就是歎息哀求,.
包括切實的凝聚悔改..
當我們研究先知書的時候,.
我們就發覺,.
當人在絕望,.
在困難裡面,.
和神歎息哀求的時候,.
先知通常在當中有一個要求,.
就是首先要求以色列人要凝聚悔改,.

$^{601}$和離開偶像的崇拜,.
這是神介入人生命裡面一個先決的條件..
如果明顯的例子,.
就好像以賽亞書第57章第13節,.
耶利米書第11章第11節,.
和荷西亞書第7章第14節..
在我們中國人的歷史裡面,.
在中國人的教會歷史裡面的時候,.
在20世紀頭40年的時候,.
中國經歷過三次復靈的大復興..
第一次稱為滿洲大復興,.
是由加拿大一個宣教士叫Johnathan Goldforth,.
叫古約翰牧師,中文譯做..
大靈之下在1900年代發生的..
第二次是山東大復興,.
是一個荷蘭的宣教士在1920年,.
大靈之下發生的..
而最後一次是宋尚哲博士,.
在1930年代大靈發生的復興運動..
這三次的復興運動,.
使得中國的新教徒,.
即Protestant Christians,.
的數目是由1900年代的18萬人,.
大幅增加到1930年代的80萬人..
這三次的復興運動,.
是很有效地預備了中國的信徒,.
在將來要面對的苦難,.
先是中日的戰爭,.
後來是國共內戰,.
以及一連串的逼迫等等..
而當我們研究這三次的復興運動的時候,.
我們就會發覺有一個共通點,.
就是當人向神嘆息哀求的時候,.
聖靈是會以超賤的能力,.
去迫使人認罪悔改的..
我現在想讀一段,.
由Dr. Walter Phillips,.
對古約翰牧師,.
當他帶領的復興聚會,.
其中一個他的觀察,.

$^{641}$讓你明白,.
究竟聖靈是怎樣工作的..
在彭建輝牧師學時,.
即是中神第一任的院長,.
他曾經翻譯了一本書,.
這本書的英文是《By My Spirit》,.
這本書很有名,.
因為當時我讀神學,.
這本書是差不多150年之前寫的,.
但我讀神學的時候,.
這本書仍然在書架裡面找到,.
那你就知道,.
他沒有outdate(更新)到..
我找到英文的書,.
《By My Spirit》,.
彭建輝牧師當時,.
我想比我早50年,.
他將它翻譯成中文,.
就放在他那本書裡面,.
那本書叫做《飲於能力之源》,.
我相信可能你們未聽過,.
因為這本書已經絕版了,.
但我相信在宣道會的書架裡面可以找到..
當中他這樣說,.
Dr. Water Phillips這樣說,.
他說我在錦州第一次接觸到滿洲復興的情形,.
我到錦州的時候,.
奮興會已經開始了一個星期,.
我在心情沒有準備的情況之下,.
走入聚會之中,.
坦白來說,.
我的性格使我對於任何形式的奮興會,.
那些人好像很歇斯底里的反應,.
是有強烈的偏見和反感,.
所以我的觀察至少是客觀..
然後他就繼續寫下去,.
當我一進到禮拜堂的時候,.
我就立刻感覺到有不尋常的事在進行,.
堂裡面異常的擁擠,.
聽眾坐到門口,.

$^{681}$每個人的臉上都帶著緊張,.
但肅靜的表情,.
唱詩的聲音裡面充滿了新的喜樂和力量,.
會眾首先是跪在地上祈禱,.
先是脈吐,.
後來這裡一個,那裡一個,.
開始有人出聲禱告,.
會眾禱告的聲音越來越大,.
合成了禱告的巨浪,.
然後慢慢縮小下去,.
變成絕音之聲,.
現在我終於明白,.
為什麼地上是濕的,.
因為滴滿了眼淚,.
氣氛之中好像有電力,.
我嚴肅地說,.
奇異的微震在我身體上下不停地移動,.
後來,有一個人,.
抽搐而壓抑的聲音,.
超出了會眾哭泣的聲音之上,.
開始凝聚祈禱,.
我的說話實在不能夠形容,.
這個凝聚的禱告裡面,.
所包含,.
也所引起的震驚和同情,.
使人震驚的並不是凝聚之中所顯露的大罪,.
而是凝聚者心靈裡面的傷痛,.
嘆息,.
泣求,.
振戰的聲音,.
他們經過掙扎,.
本來不想說的,.
但是被一種力量,.
催逼他,.
不得不起身,.
將他的心事去傾討,.
而這種的能力,.
這種的情況,.
使人不禁流淚,.
我從來沒有經驗過,.

$^{721}$比這個更加動人心神的情況,.
每個人的心,.
是赤露塘開,.
在神的面前,.
這種的情況,.
是一個小時,.
一個小時,.
不停地重覆下去,.
親愛的弟兄姊妹,.
有一個我會說,.
你和我說的,.
都是百多年前發生的事,.
香港的歷史裡面,.
從來沒有屬靈的福氣,.
我們試過很多次,.
但是都不行,.
究竟我們有沒有這樣的可能,.
神有沒有聽我們禱告呢,.
如果我們追求的時候,.
會不會有這樣的可能性呢,.
親愛的弟兄姊妹,.
我想和你分享一個,.
更加近期的例子,.
在70年代的時候,.
上個世紀70年代,.
1970年,.
距離現在我想50年,.
那時候的美國學生,.
對那時候美國的政府,.
十分之不滿,.
他們反對美國政府,.
介入越南戰爭裡面,.
這個我想你們不太熟,.
當時美國裡面,.
四處都有暴亂動盪和示威,.
但是在這個困難的時間裡面,.
教會要同時面對兩個很大的問題,.
一個是性解放運動的開始,.
一個是同性戀運動的開始,.
美國的屬靈的狀況,.

$^{761}$好像在很短的時間裡面,.
急劇的變壞,.
屬靈的復興神介入,.
這些想法,.
好像變得很遙遠,.
但是就在這個很困難的情況下,.
在1970年2月的一個早上,.
在美國中部Kentucky的一個,.
Esprit College裡面,.
當學生進行早會的時候,.
有一位同學問准了老師上課,.
要他講幾句說話,.
於是他就在台上,.
分享他的見證,.
和神怎樣改變他,.
接著他做了一件事,.
就是他每一個得罪過的人,.
他向他認罪悔改,.
他的分享當時,.
在那個70年代,.
仍然是很保守的年代,.
美國的中部更加的保守,.
是產生了很大的迴響和震撼,.
當這個同學做完見證的時候,.
下台之後,.
馬上有一個女同學,.
接著他上台,.
去認罪和做見證,.
在聖靈的感動之下,.
學生是一個接著一個,.
上台認罪悔改和做見證,.
那時候神的同在這麼真實,.
以致那時候,.
其中一個Professor,.
轉頭和他坐在旁邊的School of Dean,.
教導主任,訓導主任說話,.
God is here,.
復興的火焰一直燃點下去,.
那個早會沒有在那個早上結束,.
到下午吃飯的時候,.

$^{801}$其他學生去飯堂,.
他就發覺,.
所有去早會的人,.
並沒有去吃午飯,.
所以吃午飯的飯堂裡面,.
是很少人,.
於是他就開始不禁問,.
究竟他們去了哪裡,.
結果就是,.
這些人都去了那個早會裡面,.
當天,.
本來晚上有個籃球比賽,.
一個很重要的節目取消,.
所有課堂暫停,.
那個早會並沒有在那天結束,.
而是一直維持了八天八夜,.
超過一百七十個小時,.
當這個奇特的復興現象,.
傳出社區的時候,.
就吸引了電台和電視做採訪,.
當消息經過廣泛的報導之下,.
就吸引了很多美國和加拿大,.
學生院校的注意,.
他就有很多去請Asprey College的學生去做見證,.
結果很奇妙,.
當Asprey College的學生在台上做見證的時候,.
同樣性的能力降臨在會場上,.
而這個事件,.
是引起了北美一個很大的復興,.
這個就是很著名,.
你們可能還不知道,.
是新一波學生運動的開始,.
Navigators是那個例子,.
有一個人本來是想去佛羅里達度假的,.
但是他有心機知道Asprey College的消息,.
結果他最後的時候轉了機票,.
想去Fayetteville, Kentucky,.
他想了解一下Asprey College發生了什麼事,.
但是當他的車來到Asprey College的門口停下的時候,.
他開始聽到好像火車一樣的聲音,.

$^{841}$然後他發覺原來這些聲音是教堂傳出來讚美和禱告的巨鱷,.
他就被吸引走向教堂,.
一路走到教堂的時候,.
聖靈就感動他,.
他就一直開門,.
走到教堂裡面的時候,.
去到祭壇的面前,.
就在那裡跪下,.
認罪悔改,.
接受耶穌基督做他們的個人救主,.
你想不想知道這個復興運動的秘密呢?.
Asprey College的校長Dr. Keno,.
曾經錄製一個錄影帶,.
解釋過這個現象,.
這個錄影帶我看過,.
是在Gordon Conwell的圖書館,.
你要用VHS看,.
那時候他還有機器,.
現在我不知道了,.
他說一切緣起都是一個女同學,.
是一個人,.
他懇切的禱告,.
他對復興有特別的負擔,.
起初是一個人,.
為美國的情況,.
和他學校的情況祈禱,.
接著在69年10月的時候,.
他開始邀請其他的學生一起禱告,.
這個禱告的小組開始越變越大,.
有越來越多的人去同心的呼求神,.
在4月之後,.
神聽了他們的禱告,.
行了神跡祈禱,.
挽回了那個的世代,.
所以今天你問我,.
究竟學講得不得,.
我從來都未試過,.
是否不得呢?.
那我問你一個問題,.
你相不相信希伯來書第十三章第八節所講的話,.

$^{881}$當中他講到耶穌基督,.
昨日今日一直到永遠,.
是不改變的,.
如果你相信這句聖經,.
你相信神是不改變的,.
祂的能力沒有改變,.
祂的心意沒有改變的時候,.
那我可以答你,.
是可以的,.
因為神是一個可以將不可能變成可能的神,.
當你進入到有效的禱告裡面的時候,.
這必然是一個改變生命的禱告,.
當你認識到,.
神是你唯一的希望的時候,.
當你不再靠自己的力量,.
容許這個認識去改變你禱告的狀態的時候,.
當你進到一個地步,.
是不止息的求神加入你的生命的時候,.
當你真真正正切實認罪悔改的時候,.
容許你一起禱告..
親愛的主,.
我眼目定睛在你的身上,.
主啊,.
我知道你是一個可以將不可能變成可能,.
可以由沒有變有,.
可以叫死人復活的神,.
神啊,求你今天用這種能力去改變,.
我們每一個坐在這裡等於只會尋求你的心,.
讓他們曉得又明白,.
你是大能的神,.
沒有東西不在你的控制下,.
你能夠做任何的事,.
神啊,求你帶領我們更深入你的同在,.
讓我們認識到在當中,.
不只是我們去看聖經,.
而是讓我們進入到聖經的裡面,.
進入到你的心意裡面,.
讓我們經歷聖經的大能,.
讓我們看到你自己的榮耀,.
在一個不可能的世代,.

$^{921}$在一個彎曲薄茂的日子,.
在我們有時不知何去何從,.
我們怎樣走下去的時間裡面的時候,.
主啊,我們可能都會遇到困難,.
我們遇到一些困難,.
我們力不能勝,.
我們不知道怎樣跟別人說,.
但是神啊,求你伸出你的手,.
帶領我們積著困難發出你的邀請,.
讓我們進入到你的同在裡面去,.
你在當中就好像說,來啊,來啊,.
你在困難裡面嗎?.
你在你勝不過的罪裡面嗎?.
你在你自己敵擋不過的環境裡面,.
你在當中已經沒有能力,.
不能夠再去對抗的時候,.
來啊,來到我面前,.
當你真真正正認識,.
我是獨一的神的時候,.
當你認識到在當中,.
你既然靠著聖靈入門,.
今天不是靠著肉身去成聖的時候,.
今天你仍然接受,.
神是你唯一的希望的時候,.
神就能夠盡力為生命去改變你..
\newpage



\section{撒母耳記上 15:17-23}
\label{sec:032024}
\textbf{始終一僕人}~~~~曾敬賢傳道
\newline
\newline
連結: \href{http://www.church.com.hk/acms/content.asp?site=cdc\&op=show\&type=product\&code=032024\&layout=sermon}{\texttt{032024}} ~~~~ 語音日期: 2026-04-19
\newline
\newline
\hyperref[sec:032023]{\small{< < < PREV SERMON < < <}}
~
\hyperref[sec:index_chronic]{\small{[返順時目]}}
~
\hyperref[sec:032025]{\small{> > > NEXT SERMON > > >}}
\newline
\newline
撒母耳記上 15:17-23
\newline
\begin{longtable}{cl}
\hline
\hline
章節 & 經文 (和合本修訂版)\\
\hline
15:17 & \begin{tabularx}{0.7\textwidth}{X} 撒母耳說：「你雖然看自己為小，你豈不是作了以色列諸支派的元首嗎？耶和華膏你作了以色列的王。 \end{tabularx} \\ \\ \relax
15:18 & \begin{tabularx}{0.7\textwidth}{X} 耶和華差遣你，吩咐你說：『你去除滅那些犯罪的亞瑪力人，攻打他們，直到把他們完全滅盡。』 \end{tabularx} \\ \\ \relax
15:19 & \begin{tabularx}{0.7\textwidth}{X} 你為何沒有聽從耶和華的話呢？你為何急著撲向掠物，行耶和華眼中看為惡的事呢？」 \end{tabularx} \\ \\ \relax
15:20 & \begin{tabularx}{0.7\textwidth}{X} 掃羅對撒母耳說：「我聽從了耶和華的話，行了耶和華派我行的路，擒了亞瑪力王亞甲來，滅盡了亞瑪力人。 \end{tabularx} \\ \\ \relax
15:21 & \begin{tabularx}{0.7\textwidth}{X} 百姓卻從掠物中取了牛羊，是當滅之物中最好的，要在吉甲獻給耶和華－你的神。」 \end{tabularx} \\ \\ \relax
15:22 & \begin{tabularx}{0.7\textwidth}{X} 撒母耳說：「耶和華喜愛燔祭和祭物，豈如喜愛人聽從他的話呢？看哪，聽命勝於獻祭，順從勝於公羊的脂肪。 \end{tabularx} \\ \\ \relax
15:23 & \begin{tabularx}{0.7\textwidth}{X} 悖逆與占卜的罪相等，頑梗與拜偶像的罪孽相同。因為你厭棄耶和華的命令，耶和華也厭棄你作王。」 \end{tabularx} \\ \\
[1ex]
\hline
\hline
\end{longtable}
$^{1}$在一月的時候,我去了一個按摩禮.
下午也有按摩禮,是嗎?.
一月的那個就是我同學的一個按摩禮.
當中有一位很有份量的牧者.
尤其是在我們播道家裡面.
郭文慈牧師就去做了一個訓練.
對兩位將會按摩的弟兄去做訓練.
到最後的時候有一句話對於我來說很深刻.
這句話就像今天的主題說的.
「始終一步人」.
不知道你們有沒有人聽過這句話.
有…沒有.
我每天都聽到這句話「始終一步人」.
其實這句話我當時還沒有聽過.
但當郭牧師跟我們分享的時候.
原來他說當中有兩個意思.
第一個就是我們要始終如一.
始終的意思是我們要始終如一.
特別是在說什麼呢?.
可能是…因為他對牧者說.
由我們去呼召做一個牧者.
去到現在的狀態.
甚至將來我們的服侍.
我們都要始終如一.
我們對上帝的心要始終如一.
我們要…可能甚至是一個僕人.
一個謙卑的心.
我們都要始終如一.
所以一個僕人的意思就是謙卑的身份.
我們是上帝的僕人.
無論在任何時候.
我們要做一件事都要聽命於上帝.
上帝要我們做的事我們就做.
上帝不要我們做的事我們就不要做.
所以那一刻對我來說是一個很好的提醒.
因為那一刻我都想起.
我作為…做牧會十年.
我都想起剛出來牧會的時候到現在.
我的事奉是不是都是始終如一.
那句話好像是在對牧者說.

$^{41}$但另一方面我都想其實作為一個信徒.
我們都要去想一下.
由我們初信那一刻到現在.
我們是不是對上帝都是始終如一.
甚至將來我們回去要見耶穌基督的時候.
我們是不是都是始終如一.
今天我們去看一個人物.
就是素羅王.
素羅王開始的時候.
其實他是一個很謙卑的人.
因為他謙卑.
他被上帝揀選了他為一個王.
也是以色列裡面第一個的王.
他的生命.
但是就是因為成了王之後.
他開始去改變.
由一個謙卑變成用自己的方法.
白天說的就是驕傲.
用自己的方法去處理這些事.
就是叫驕傲.
驕傲會令到人家怎樣呢.
會忘記上帝.
忘記了我們.
原來我們的生活我們是要依靠著上帝.
我們今天看撒慕爾記上的第十五章.
我們一起讀一讀.
第十七節.
因為我們用的是新譯本.
我們一起讀.
預備請.
撒慕爾說.
從前你雖然自以為微小.
不是作了以色列眾支派的首領嗎.
耶和華高納了你作以色列的王.
撒慕爾眼中.
或者其實撒慕爾這一刻.
他的評論去比素羅是什麼呢.
他說你以前是自以為微小.
其實說什麼呢.
他以前的確是一個謙卑的人.

$^{81}$待會我們再看看一個謙卑的人.
也是因為這個謙卑的原因.
或者這個特質.
上帝揀選了他做什麼.
以色列的王.
素羅怎樣謙卑呢.
我們可以看回撒慕爾記上的第九章.
我們讀一讀.
一到二節.
預備請.
有一個變眼敏人名叫基時.
他有一個兒子名叫素羅.
既年輕又英俊.
在以色列人中.
沒有一個能比他更英俊的.
他比人高出一個頭.
素羅是怎樣的.
每一場我問到這個問題.
弟兄姊妹回答出來都不同.
昨天星期六.
我一說的時候就很一致.
他怎樣.
英俊.
我這個一定是姐妹回答的.
但我更加期望大家看聖經.
由上看到下.
聖經說他第一個是什麼.
年輕.
我們先看年輕.
不要一看就看英俊.
我的姐妹很喜歡.
還要英俊說兩次.
重要.
說兩次.
第二個才英俊.
第三個高出一個頭.
我都說我為同福堂去感恩.
因為我們都有一位主任牧師.
又年輕又英俊.
又高人一個頭.

$^{121}$Ken哥今天不在.
去了九龍東.
所以我可以說一下他.
如果各位弟兄.
你是一個又年輕又英俊.
又高大的人.
換上現在這個年代.
你應該是一個很有自信的人.
走出來的時候應該.
嘩.
即是怎樣.
應該都不會自卑.
但蘇羅其實他那時候是怎樣.
我們再看一看.
第21節我們一起讀.
預備123請.
蘇羅回答.
我不是以色列眾支派中.
最微小的汴瓦文支派嗎.
我的家族不是汴瓦文支派的.
家族中最微小的嗎.
你為什麼對我說這樣的話呢.
蘇羅說了一件事.
他是最微小.
微小是有兩方面.
第一他在他的支派裡面是微小.
他在哪個支派.
汴瓦文.
不是因為他是最小的兒子.
而是汴瓦文的確他的人口.
即是那個年代.
人口上和他的地理上.
都是真的小.
所以對比於其他的支派.
所以他說.
這個就是說他支派裡面.
他真的是一個微小.
第二.
他說到他的家族.
在這個微小的支派裡面他怎樣.

$^{161}$再微小一點.
即是微小中的再微小.
所以當其實.
薩姆爾去跟蘇羅說.
呼召他的時候.
其實他看自己.
就是說我一個這麼微小的人.
怎樣去做.
其實你會看到.
原來薩姆爾.
蘇羅最初的時候.
他都是一個謙卑的人.
各位姐妹.
如果我有時候去想.
到底我們.
什麼時候是跟上帝的關係.
是最緊密.
或者是.
白說.
最聽話的呢.
可能有些人說.
我經歷了一件事情.
之後上帝幫助了我.
跟上帝有一個很緊密的關係.
但有時候我想最簡單的.
就是我們初信.
因為為什麼.
我們初信什麼都不懂.
別人說什麼的時候.
我們就去聽.
應該是這樣的.
因為如果初信都不聽.
應該都不會相信的.
所以如果你初信的時候.
應該是最聽話.
我想起我自己.
我初信的時候.
我是中學的時候.
我又信耶穌.
我的教會是團契的.

$^{201}$要回到團契.
回到團契的時候.
傳道人就跟我說.
喂你要零修.
我就說零修啊.
一個禮拜多少次.
好像看醫生一樣.
一個禮拜吃藥多少次.
一個禮拜多少次.
他說每天都要的.
我每天都要.
又不敢拒絕的.
一個禮拜後.
回到團契.
他第一件事就把我帶進房間.
他說.
Castor.
你這個禮拜.
有沒有零修啊.
你猜我會怎麼回答.
有些人就說.
有有.
死都說有.
初信嘛.
初信就信良如格子的.
當然要跟他說.
沒有啦.
那時候都不知道零修是什麼.
信了沒多久.
然後他說.
沒有啊.
你進來.
在房間裡.
關上門.
然後被人禁錮了一個多小時.
在裡面做什麼.
沒有事情做的.
不斷地跟你說.
為什麼你不零修啊.
零修很重要的.

$^{241}$你為什麼不零修啊.
為什麼這麼不聽話.
我在想.
如果我現在做牧者.
我夾過大英姐妹.
問她有沒有零修.
然後就不斷地跟她說.
你猜她會不會理我.
我不知道.
我不說了.
可能有些人會理.
有些人不會理.
但是那一刻對我來說.
真的這麼信服.
下個星期.
你零不零修啊.
零.
一定零修.
不是下次就兩個小時.
為什麼你這麼不聽話.
所以一定零修的.
但是怎麼做呢.
不知道.
但是照做.
大英姐妹.
但是有時候都真的.
原來當我們什麼都不懂.
或者其實我們覺得自己不足的時候.
原來我們反而是最聽話的時候.
真的是最聽話的時候.
就是那個時候了.
到了現在的時候.
都有零修的.
不是說我.
因為我們屯門堂.
我們都有門裡.
門裡有一個要求.
每天不單止零修.
每天是要將零修的筆記.
就是分享.

$^{281}$扔到群組那裡.
然後那些人就牙痛.
我自己零修就好了.
還要扔過來.
那一開始是牙痛的.
現在被我鞭策到他們.
都是要扔出來.
但是的確是想回.
可能是最單純的時間.
是我們初信.
上帝要我們做什麼.
我們就去跟.
蘇羅是怎樣呢.
我們看看.
看回第十五章.
蘇羅的時候.
我們一起讀第十八十九節.
請.
差遣你去說.
你去徹底毀滅那些犯罪的亞馬利人.
攻打他們.
直到把他們全部毀滅.
你為什麼沒有聽從耶和華的話.
急於奪取戰利品.
作了耶和華看為惡的事呢.
原來這次其實.
上帝要蘇羅做什麼.
他要蘇羅去攻打亞馬利人.
上帝要求蘇羅是.
第十八節最後那四個字.
什麼.
全數毀滅.
OK.
什麼叫全數毀滅呢大家.
不用解釋都知道.
全部東西.
滅盡.
很清晰.
OK.
上帝給我們的一個要求.

$^{321}$或者一些要我們去做的.
其實有時很清晰的大英姐妹.
OK.
不過蘇羅做什麼呢.
第十九節.
撒末爾說他沒有聽從.
原來當他沒有聽從的時候.
就是在做一些.
耶和華看為惡的事.
這個徹底滅盡.
說的就是不可以去保留.
等一下蘇羅有解釋的.
為什麼他要去保留呢.
他會有解釋.
為什麼蘇羅會這樣呢.
有很多時候就是說.
原來當我們的信仰久了.
我們有一些自己的經歷.
慢慢有時我們反而會忘記了上帝.
會放下了上帝.
甚至乎在我們生活中.
我們會忙碌.
忙碌也是我們忘記上帝其中一個有機會發生的事情.
很感恩我們在屯門堂.
在受苦之前一天.
受苦之前一天.
我們很難得我們跟屯門區五間播道會的堂會.
一起有一個聯合的祝福行動.
祝福什麼呢.
就是我們屯門堂旁邊有一條村叫做補田村.
很想在那裡有一個祝福行動.
我們戶外的.
我們可能派一些禮物給他們.
跟他們去敬拜.
去講信息.
其實以前有的.
疫情前我們每個月.
屯門堂我們自己都有.
但疫情之後沒有了.
後來有其中一間堂會的人打給我.

$^{361}$我們不如一起去搞一個活動.
我說好啊沒問題.
很好啊.
講完之後.
他說什麼時候.
他說兩個星期後.
兩個星期.
會不會急了一點.
他說不急了.
因為我們當然是有意義的.
因為復活節.
受冷戰之前.
我們祝福街坊很有道理.
我說好啊.
我們就決定.
然後我另一個同工就說.
那怎麼做啊.
這麼急.
我說不要緊.
我們之前做過.
搬那堆東西出去就行了.
我講完這句話的時候.
突然聖靈提醒.
祂說你做漏了東西.
知道做漏了什麼嗎.
祂說我做漏了祈禱.
我一想.
反正事情跟以前這麼像.
這麼一樣.
就要搬回來就行了.
尤其是時間這麼擠.
這麼忙碌.
這麼緊迫.
很多原因.
然後一連串去合理化.
搬回以前那樣東西就行了.
但是聖靈提醒我說你還沒祈禱.
你沒問上帝這次想怎麼做.
其實看聖經也是.
這次是這麼做.

$^{401}$下次未必是這樣做.
我們要祈禱.
問上帝.
所以認罪.
立刻我們祈禱.
我跟同工一會兒祈禱.
祈禱就是上帝你要怎麼做.
祈禱之後.
有一間教會就打來.
就是其中一個合作的教會.
不如我們大敬拜.
我說好啊.
就大敬拜.
但更重要的是什麼呢.
是我最頭痛的事情.
我們有很多禮物.
我們那天排隊.
然後我想禮物.
其實我預備禮物.
是我最弱的一個環節.
他幫我處理了.
但更感恩的是什麼呢.
人.
我們去預備了這個聚會.
上帝為我們預備了天氣.
如果你記得在壽苦節前一天.
應該是星期四.
我相信大家應該在九龍和港島區.
多數.
那天其實下大雨.
上午的時候.
我們有一個同位.
因為他上班.
他在九龍區.
就馬上給我電話.
喂 卡薩.
你下午能不能搞定.
我問他做什麼.
他說下雨.
然後我下雨脫毛帽.

$^{441}$然後他說.
不過全黑了.
我請你們努力為我祈禱.
下午不要下雨.
然後再過一會兒.
到下午.
喂 真的不停.
九龍區真的不停.
我真的沒有雨.
我們結果全程都沒有雨.
更加很感恩.
跪給耶穌.
更加感恩的是什麼.
那次在我們無心插柳的裡面.
有街坊信耶穌.
弟兄姊妹.
上帝要我們跟隨祂的心意去走.
是跟隨祂的方式.
雖然在我們信仰.
我們信得久了.
我們會覺得.
咦 應該是這樣.
但是上帝從來不是要我們應該.
是要跟隨祂去做.
所以其實上帝的訊息很清晰.
上帝說要蘇羅是要怎樣呢.
徹底滅盡.
全部毀滅.
但是蘇羅沒有做到.
蘇羅用了什麼的原因呢.
我們如果看回前少少撒母耳記上的第十五章第九字.
我們也讀一讀好嗎.
預備請.
但蘇羅和眾人連飾雅各.
也愛惜上好的牛羊.
肥牛毒.
羊羔和一切美好的東西.
他們不願徹底毀滅這些.
凡是卑賤的.
沒有價值的.

$^{481}$他們就徹底毀滅.
再重複看一次.
上帝要求的是什麼.
徹底毀滅.
四個字而已.
徹底毀滅.
全部我們都要毀滅.
但是蘇羅做了什麼.
連飾雅各.
雅各是什麼呢.
雅各是一個王.
白說就是生擒了他.
拿了他回來.
這個一會兒再說.
但是更重要的是.
他愛惜上好的牛羊.
肥牛毒.
羊羔.
一切美好的.
所以弟兄姊妹.
當我們看到美好的.
我們要小心.
美好的有時.
我們都要看會不會讓我們跌倒.
亞當 夏娃.
都是看見美好的.
但也是讓他們跌倒的一件事.
所以其實他們不願徹底毀滅.
他們毀滅了什麼.
一切是卑賤沒有價值.
這個完全違反了上帝.
在蘇羅的心目中.
或者其實要蘇羅做的事.
所以蘇羅其實在這一部分.
最大的問題是什麼呢.
第一件事就是.
蘇羅他本身是一個謙卑的人.
但是因為他的經驗.
或者其實他做王.
做了一段時間.

$^{521}$他遺忘了上帝.
所以剛才第一個.
其實說始終一僕人.
就是由始到終.
我們都是要成為一個聽從上帝的一個人.
他因為王位.
他越來越自大.
更加重要的是什麼呢.
他驕傲開始.
不再需要依靠神.
他依靠的是自己.
他喜歡做什麼就做什麼.
所以有時候我們看回罪的根源.
有時候說行為.
蘇羅做了什麼.
他不聽話.
他拿了一些牛羊回去.
但其實有時候不是.
其實是由不聽話開始.
好像一個罪的根.
已經在蘇羅裡面去發生.
我們再看第三部二記.
第十五章的二十節.
我們一起讀.
預備請.
蘇羅回答撒母耳.
我實在聽從了耶和華的話.
行了耶和華差遣我行的路.
我把亞馬利人的王雅各帶了回來.
又把亞馬利人都徹底毀滅了.
蘇羅有沒有聽.
有.
有些.
有沒有些.
我們做信徒也是的.
我有時也是的.
上帝說的話我們有些.
有沒有些.
聽些不聽些.
有些毀滅了.

$^{561}$有些他帶了回來.
其中一個就是我剛才說的雅各.
雅各在這裡說什麼呢.
原來是說亞馬利人的王.
為什麼要帶雅各回來呢.
其中有一個原因.
有人說蘇羅.
其實你想想.
我去打一個國家.
打贏了.
擄掠了一個王回來.
其實挺威風的.
行不行.
其實也挺威風的.
所以其實有時.
這個其中一個原因就是.
他將個人的榮譽.
凌駕了上帝的說話之上.
他是一個驕傲的表現.
第二件事很明顯就是.
違反了上帝對他的要求.
我們要聽從耶和華的話.
但是我們不要將一些錯誤合理化.
其實蘇羅你會看到.
就是他將一些東西.
他就將它合理化了.
上帝說要徹底毀滅.
蘇羅他沒有徹底毀滅.
而且更加重要就是說.
他已經聽從了.
他覺得自己已經聽從了.
還要是已經跟了耶和華的話去行.
我們再看第21節.
我們再讀.
預備請.
「可視眾民卻從那些戰利品中.
取了應當毀滅之物中.
上好的牛羊.
要再吉甲獻給耶和華.
你的神為祭」.

$^{601}$將它合理化到一個什麼程度呢.
上帝再重申.
我們要記住.
上帝要他怎樣.
徹底毀滅.
不過他這次帶了牛羊回來.
然後跟撒慕爾說.
我帶回來有原因的.
什麼原因.
憲制.
所以撒慕爾會用憲制.
去跟蘇羅說.
其實他在做的是什麼.
但很明顯.
這不是上帝要做的事.
但蘇羅他用了一個解釋.
他去合理化他自己的錯誤.
首先我們要看.
到底什麼叫做當滅之物.
其實有三方面的.
當滅之物.
第一樣.
原來這個物件.
是要歸給耶和華.
這東西不是屬於任何人.
人不可以拿.
這個是要歸給耶和華的.
用什麼方法呢.
只有一個方法.
就是徹底毀滅.
所以為什麼上帝要蘇羅.
將所有的東西毀滅.
如果你想和記得.
以色列.
當他進到迦南.
打第一場仗是誰.
耶利哥.
耶利哥的時候.
上帝都是有同一個要求.
將耶利哥所有的東西消滅.

$^{641}$不過有一個人不聽話.
阿干拿了.
都是因為貪戀.
漂亮的東西先拿了.
結果犯罪.
當滅之物有一樣很重要.
就是要徹底毀滅.
這個是上帝的吩咐.
這個是一個方式.
我們說這個是一個方式.
內容是什麼呢.
內容是要完全獻給神.
所以你會看到當滅之物.
不應該落在人的手上.
當滅之物我們要用上帝的方法.
就是要除滅.
我們要獻給上帝.
所以不是獻祭.
我們有時候.
用一些藉口或原因.
我們要說服我們自己做的事.
在穆會裡.
過去都有人問我一個問題.
到底是否有一個.
基督徒可以去旅行.
一定可以去的.
我也有去的.
說清楚一點.
基督徒應不應該星期天去旅行.
我經常都被人問這個問題.
尤其是剛剛過了.
復活節.
然後看到有一堆.
突然之間我講屯門堂.
也有些人去了旅行.
有人問我.
Castor到底應不應該去旅行.
不去崇拜.
很多人想我做牧者.
都可以去.

$^{681}$很多人想我做牧者.
想牧者怎樣畫一條線.
說行不行.
我經常都說.
有些事不是行不行.
是你怎樣去做.
我認識有一位.
是信二代還是三代的弟兄.
他很堅持的.
很堅持什麼呢.
就是他星期天一定要回崇拜.
所以他去旅行怎樣呢.
星期一飛.
星期六回來.
星期五飛.
什麼時候回來.
星期六回來.
星期六飛.
都是星期六回來.
不是飛的了.
星期六那些.
飛了過去都馬上坐回來.
又飛幾何.
他很堅持的.
跟我說.
嘩 我欣賞.
真的很欣賞.
但是我也有一些弟兄姊妹.
跟我說.
Castor 不好意思.
這個禮拜不能回崇拜.
什麼事.
然後我去旅行.
因為我媽媽年紀大了.
很難得我整個家族.
都可以一起去旅行.
所以每個人其實環境背景.
都有一些不同.
我不可以畫一條線.
告訴你對不對.

$^{721}$但是我們的內心.
我們的心態.
我們要知道.
我們到底在做什麼.
我有一個要求.
我就跟他說.
雖然我說你這個禮拜不能回崇拜.
不過請你都要回崇拜.
這個是我唯一的要求.
因為我們是基督徒.
這個是我們對信仰一個的.
信仰對我們一個的要求.
所以當我說你去旅行不要緊.
請你星期日.
你在當地找一間教會去.
你找一間華人教會也好.
你找一間.
如果去日本找一個日文教會.
你聽得懂的也可以.
你去韓國的.
請你找一間正統的.
你確定是正統.
也可以.
那你說不行.
找不到教會怎麼辦.
不要緊.
我們同福堂也好.
網上崇拜.
都請你.
你們要有一個地方.
你們要敬拜上帝.
這個是我的要求.
他說不是的.
我去到一些地方.
連WIFI都沒有.
我心想.
現在是不是去野餐.
我不知道.
我說那你先預備好.
家庭崇拜你也要.

$^{761}$但僧兄姊妹有時我們都是.
不要給一些藉口.
遺忘了我們做基督徒.
我們要做的本分.
不要給一些藉口.
合理化了我們一些事情.
當是上帝要我們做的事.
我們就要去聽.
因為這個是我們的態度.
我們對信仰的態度.
剛才素羅就是說.
他拿這些東西就是為了憲制.
所以去到最後的時候.
撒慕爾都是用憲制去跟素羅說.
他就說聽命勝於憲制.
我們讀一讀第22節.
預備123請.
撒慕爾說.
耶和華喜悅凡濟和平安濟.
喜將喜悅人聽從耶和華的話.
聽從勝過憲制.
聽命勝過公羊的自由.
有一些人將宗教和生活分開.
簡單來說.
我們經常說有一些叫做.
星期日的基督徒.
星期六的基督徒.
我在讀神學之前.
我是做銀行.
我知道同福堂很多牧者都是銀行界出身.
我是其中一個.
當然沒有他們那麼高級.
我是做服務的.
那時候就去了.
剛剛去了一條新的團隊.
跟一個上司.
那個上司看到我.
咦 你吃飯祈禱啊.
你是基督徒來的.
我是啊 我是基督徒來的.

$^{801}$我也是啊.
接著我就覺得.
哇 你給了我禮物.
你會怎樣 你會很開心.
因為上司也是基督徒.
是不是 應該很棒才是.
接著之後.
我就很自然地問他.
你去哪一間教會啊.
我沒有去教會.
我的心裡面就嚇一跳.
沒有去教會的基督徒.
但是問題是.
我又不敢.
因為那個是上司.
怎樣也要給一點面子.
我就說 哦 是啊.
沒有去教會 哦 好.
直至有一次我走過他的位置.
他的位置有一個飾物.
放在桌面上.
我一看這個飾物我就知道.
應該是去哪裡旅行之後.
不知道求了什麼.
帶回來的東西.
我就真的有點忍不住.
我就跟他說.
咦 我說作為基督徒.
這個是不是不應該放在這裡.
甚至不應該去求呢.
接著他就說了一句話.
哦 不要緊.
上帝明白的.
我也不知道怎樣回答他.
接著他再說.
還有 真的不要緊.
我每年聖誕節都有回教會.
那我說你回去幹什麼.
馬上就很好奇.
你回去幹什麼.

$^{841}$回去認罪.
我回去認罪.
我會將那一年.
我上帝不喜歡的事.
我會藉著那一次回教會.
求耶穌去寬恕.
耶穌一定會原諒我的.
弟兄姊妹.
對於我們來說.
上帝的教導是不是這樣.
是不是可以用這個方法去合理化.
我也很明確地說.
哦 不過我知道.
每個星期都要回教會.
靈修也要被迫.
回教會也是.
我也說你每個星期都要回教會.
當然到了最後.
後來我沒有做.
接著我不知道他怎樣.
弟兄姊妹.
我們有很多時候.
真的將我們的生活.
將我們的信仰.
我們分開.
覺得好像平時做的事.
不要緊的.
世俗的事.
我們都是.
每個人都是這樣的.
上帝會知.
上帝會寬恕.
但是上帝其實要我們.
做一個二四七的基督徒.
我們無論任何事.
任何問題.
其實我們都是要跟上帝的心意去做.
時間的關係.
我可能一會兒.
兩個繁祭和平安祭.

$^{881}$我可能簡潔一點.
我跟大家去說一下.
好了 繁祭其實是什麼呢.
如果我們看回利未記的第一章.
第二章和第三章.
你會看到.
其實是有一個規定.
但是這個規定其實是挺有趣的.
為什麼呢.
在第二章那裡.
原來聖經說.
他們如果想獻繁祭.
我們知道繁祭是燒掉的.
這是最基本的.
但是問題是.
如果我們有個人想獻繁祭的話.
其實他就可以獻.
意思是什麼呢.
其實是甘心.
如果有個人願意甘心付出.
把祭獻給上帝獻繁祭的話.
他就可以獻.
第二個就是說.
這個祭.
繁祭只有一個目的.
就是討神喜悅.
做完這件事之後.
神會如立的.
OK.
所以就是這個繁祭.
去到最後的時候.
繁祭還有一個特點.
他不會因為一個人的經濟能力有影響.
而影響了他不能夠討神喜悅的.
意思是什麼呢.
記不記得繁祭其實是用什麼來獻的.
其實繁祭是很厲害的.
很多東西都可以的.
隨便一件事.
不是隨便一件事.

$^{921}$是上帝要求的.
上帝要求很多種不同的東西.
你都可以獻繁祭.
第一.
你有錢可以獻牛.
你沒那麼多錢就獻羊.
再沒錢又怎樣呢.
頒扣.
頒扣都可以.
所以其實.
這個祭.
上帝是說.
你如果獻的時候.
我就會喜悅.
原來繁祭是不會因為經濟能力影響了我們.
去討好上帝.
或者我們讓神去喜悅.
最重要是什麼呢.
我們是甘心樂意.
我們真的想去討神喜悅.
其實我們的生命都是的.
弟兄姊妹.
當我們去將我們生命當作活祭獻上的時候.
其實我們都是.
我們的目的只有一個.
應該就是討神喜悅.
這個.
好了.
平安祭又有一點點的特別.
平安祭有些人.
我們都會知道這件事.
甘心祭.
詳情不在這裡說.
有機會再跟大家分享.
我們再看.
你未記的那裡再看.
也是按能力的.
牛或者羊都可以.
看情況.
但是這個平安祭最特別的地方.

$^{961}$就是說.
原來祭物.
就不像繁祭那樣燒光.
是可以分的.
分給誰.
第一是分給祭司.
第二個.
連獻祭者.
OK.
即是我獻祭者.
我都可以去分.
但是分完有一個要求的.
甘心.
還願.
有一點點不同.
有些第二天之前要吃光.
OK.
有些是要隔一天.
你可以隔多一天的.
OK.
但是就一定不是隔兩天的.
就是隔兩天你要燒光.
你都不能吃的.
所以其實.
如果你真的有隻牛這麼大.
你又分了一點的話.
其實你都吃不完的.
其實鼓勵是什麼呢.
鼓勵是跟你身邊的人一起去分享.
所以其實.
平安祭有一個意思的.
就是.
神願意在人的中間.
跟我們一起去分享.
這是一個平安祭.
所以如果總合來說.
如果說到底什麼叫做獻祭.
獻祭如果簡單來說.
就是人跟神之間的一個禮儀.
就是我們對神.

$^{1001}$去表達一樣東西.
我們甘心樂意.
以及我們是要.
就是我們甘心樂意.
以及我們是.
用心去呈獻.
更加重要的是.
我們是表示我們有關係.
這個才重要的.
弟兄姊妹.
原來我們跟上帝有關係.
所以當我們生活的時候.
如果我們可以將上帝放在一邊的時候.
表示一樣東西.
原來我們跟上帝的關係.
可能都有一些的問題.
所以我們要跟上帝有關係.
最後很重要的是.
按神的方式去獻.
雖然獻祭是土神喜悅.
可以是.
所以素羅說.
我獻一個祭給上帝.
上帝喜不喜悅呢.
上帝不喜悅.
因為他不是在用上帝的心意.
他不是在聽上帝的吩咐.
雖然獻祭本來是土神喜悅.
但是我們要用得對.
跟上帝的心意.
最後說一個小小的故事.
我老婆.
我老婆9月將會去讀書.
都是入實學院.
三年.
她是一個很怕青的人.
所以其實她在三年要讀的課程.
她已經在那裡.
她還沒有報名.
但問題是她竟然拿了別人的課程出來.

$^{1041}$寫下要讀什麼.
這個要讀什麼.
那個要讀什麼.
寫了寫了寫了.
然後就想著跟隨她的時間表.
她就平平安安地去讀.
誰知道最近出了.
時間表有些不同.
其實她有點不安.
因為原來我們也很想.
突然之間問心一句.
我們也很想.
我們有一個時間表.
其實上帝是跟隨我們的時間表去做的.
但有一次她在祈禱的裡面.
上帝給了她三句話.
第一個就是希望她.
在這三年的裡面好好地去操練.
其實操練是什麼呢.
操練其實就是我們不確定的事情.
我們要繼續去相信.
操練是我們一些我們沒聽過.
或者我們不明白的.
我們是順從跟隨上帝的吩咐去做.
第二就是好好地去度過這三年的旅程.
上帝給了一個字叫旅程.
弟兄姊妹我們的生活.
都是有高有低.
這個一定.
但我們人生是旅程.
無論高低.
我們知道上帝都跟我們一起去走.
詩篇23篇大家都很熟.
最特別或者其實最重要的.
我想不是在青草地溪水旁.
而是在死淫幽谷.
上帝仍然為你和我擺設筵席.
弟兄姊妹這是很難得的一件事情.
我們沒有人想去死淫幽谷吧.
但問題的確有時候看回我們的生活.

$^{1081}$我們都有死淫幽谷的時候.
但是弟兄姊妹不要忘記我們的神是怎樣.
不是跟我們一起這麼簡單.
是在敵人面前為我們去擺設筵席.
最後更加重要的就是.
上帝要塑造他成為一個合神心意的人.
不是自己.
不是他去計劃三年的事情.
上帝你這樣給我我就會很平安.
原來不是.
上帝是要讓他去成為合神心意.
即是神喜悅的人.
神要他想做的事情.
所以到23節的時候很簡單不讀了.
就是撒母耳最後說的訴羅.
他用了一個搏逆和行邪術.
還梗和拜偶像的一個對比.
因為訴羅他有這些原因或者特質之後.
他不再需要上帝.
他仍然不信服上帝的時候.
上帝最後就說氣著你作王.
上帝離開了訴羅.
弟兄姊妹這是很可惜的一件事.
有一個謙卑上帝揀選的人.
去到最後因為他的不聽不信服.
因為自己的驕傲.
而被耶華去氣著.
弟兄姊妹對於我們的生命到底是怎樣.
我們都會有高有低.
但我們仍然會不會都聽上帝的吩咐.
最後我們都用這句話.
我們作我們今天的一個提醒.
我們一起讀.
預備123請.
聽命聖於憲制.
順從聖於公羊的自由.
始終一瀑人.
弟兄姊妹.
由始到終.
甚至將來.

$^{1121}$到耶穌基督回來的時候.
我們的心是否仍然都是.
這麼單一這麼單純.
我們聆聽信服在上帝的說話裡面.
親天父我們要多謝我們讚美你.
因為你是一位時刻與我們同在的神.
你是一位無時無刻都看顧我們的神.
我求上帝你.
將你的恩典傳言覆蓋在我們身上.
給我們生活有力.
給我們生活裡面無論遇到任何困難挑戰.
我們都要堅守我們的信仰.
給我們無論面對任何困難挑戰.
不要成為我們的藉口.
而將上帝你放在一邊.
因為我們知道.
當我們將你放在一邊的時候.
我們的驕傲就會出來.
就會做一些你不願立的事情.
上帝我們要成為一個祭.
一個輕香的祭.
為的只是討你喜悅.
所以上帝無論我們面對什麼情況.
讓我們生活裡面.
真的做一個信徒的本分.
可能很艱難.
我們仍然每個星期在教會裡面崇拜.
可能我們有艱難.
上帝當你說要我們做十一奉獻的時候.
我們就願意去擺上.
或者會有艱難.
我們生活裡面仍然信靠你.
聽從你.
因為我們知道你是獨一的神.
你是我們的主.
等著我們一起張開雙手.
領受上主的賜福.
願賜平安的神親自使你們傳言成聖.
由於你們靈魂和身子得夢保守.
在我主耶穌基督再降臨的時候.

$^{1161}$完全無可自責.
由於我主耶穌基督的恩惠.
天父上帝的慈愛.
使你們感動.
尚於你們眾人同在.
從今時直到永遠.
阿們.
\newpage



\section{路加福音 18:15-30}
\label{sec:032025}
\textbf{神國的真諦:單純的信靠和全心的跟隨}~~~~梁成裕牧師
\newline
\newline
連結: \href{http://www.church.com.hk/acms/content.asp?site=cdc\&op=show\&type=product\&code=032025\&layout=sermon}{\texttt{032025}} ~~~~ 語音日期: 2026-04-19
\newline
\newline
\hyperref[sec:032024]{\small{< < < PREV SERMON < < <}}
~
\hyperref[sec:index_chronic]{\small{[返順時目]}}
~
\hyperref[sec:032026]{\small{> > > NEXT SERMON > > >}}
\newline
\newline
路加福音 18:15-30
\newline
\begin{longtable}{cl}
\hline
\hline
章節 & 經文 (和合本修訂版)\\
\hline
18:15 & \begin{tabularx}{0.7\textwidth}{X} 有人甚至連嬰孩也帶來見耶穌，要他摸他們，門徒看見就責備那些人。 \end{tabularx} \\ \\ \relax
18:16 & \begin{tabularx}{0.7\textwidth}{X} 耶穌卻叫他們來，說：「讓小孩子到我這裡來，不要阻止他們，因為在神國的正是這樣的人。 \end{tabularx} \\ \\ \relax
18:17 & \begin{tabularx}{0.7\textwidth}{X} 我實在告訴你們，凡要接受神國的，若不像小孩子，絕不能進去。」 \end{tabularx} \\ \\ \relax
18:18 & \begin{tabularx}{0.7\textwidth}{X} 有一個官問耶穌說：「善良的老師，我該做甚麼事才能承受永生？」 \end{tabularx} \\ \\ \relax
18:19 & \begin{tabularx}{0.7\textwidth}{X} 耶穌對他說：「你為甚麼稱我是善良的？除了神一位之外，再沒有善良的。 \end{tabularx} \\ \\ \relax
18:20 & \begin{tabularx}{0.7\textwidth}{X} 誡命你是知道的：『不可姦淫；不可殺人；不可偷盜；不可作假見證；當孝敬父母。』」 \end{tabularx} \\ \\ \relax
18:21 & \begin{tabularx}{0.7\textwidth}{X} 那人說：「這一切我從小都遵守了。」 \end{tabularx} \\ \\ \relax
18:22 & \begin{tabularx}{0.7\textwidth}{X} 耶穌聽見了，就對他說：「你還缺少一件：要變賣你一切所有的，分給窮人，就必有財寶在天上；你還要來跟從我。」 \end{tabularx} \\ \\ \relax
18:23 & \begin{tabularx}{0.7\textwidth}{X} 他聽見這些話，就很憂愁，因為他很富有。 \end{tabularx} \\ \\ \relax
18:24 & \begin{tabularx}{0.7\textwidth}{X} 耶穌見他變得很憂愁，就說：「有錢財的人進神的國是何等的難哪！ \end{tabularx} \\ \\ \relax
18:25 & \begin{tabularx}{0.7\textwidth}{X} 駱駝穿過針眼比財主進神的國還容易呢！」 \end{tabularx} \\ \\ \relax
18:26 & \begin{tabularx}{0.7\textwidth}{X} 聽見的人說：「這樣，誰能得救呢？」 \end{tabularx} \\ \\ \relax
18:27 & \begin{tabularx}{0.7\textwidth}{X} 耶穌說：「在人所不能的事，在神都能。」 \end{tabularx} \\ \\ \relax
18:28 & \begin{tabularx}{0.7\textwidth}{X} 彼得說：「看哪，我們已經撇下自己所有的跟從你了。」 \end{tabularx} \\ \\ \relax
18:29 & \begin{tabularx}{0.7\textwidth}{X} 耶穌對他們說：「我實在告訴你們，凡是為神的國撇下房屋，或是妻子、兄弟、父母、兒女的， \end{tabularx} \\ \\ \relax
18:30 & \begin{tabularx}{0.7\textwidth}{X} 沒有不在今世得更多倍，而在來世得永生的。」 \end{tabularx} \\ \\
[1ex]
\hline
\hline
\end{longtable}
$^{1}$今天分享的題目是.
神國的真諦.
單純的信靠.
以及全心的跟隨.
我不知道大家有沒有玩過一個團體建立的活動.
叫做TrustFall.
有些人就將它直譯.
中文叫做信心跌落.
這是一個團隊精神建立的活動.
不過我先在這裡講講.
因為有錄音一定要先講講.
這個活動是有危險性的.
所以大家不要就這麼自己玩.
一定要有專業受過訓練的教練帶著你們玩.
才好玩.
那麼這個活動是怎麼樣的呢?.
就會有一群人.
那麼有一群人就會按照導師的指導.
去組成一個人肉的安全網.
然後有一個人就憑著信心.
跳落這個人肉的安全網.
如果大家都是按照著教練的指導.
其實這個人肉的安全網.
是可以安全地接到跳下去這個人的.
那麼這個活動其實我也玩過不少次的.
那麼我曾經就是做過安全網的成員.
是的,就好像現在這樣.
安全網的成員.
那麼我記得我做安全網成員的時候.
我就在想.
是不是真的按照這個導師的指示.
這樣做一個安全網.
當那個人跳下來的時候.
我們就真的很安全地去接住他呢?.
那麼你可能會問.
那麼你有沒有試過做跳下來的那個人呢?.
我試過啊.
好啊.
有一次呢.
我參加一個是有幾十位牧者一起參加的活動.

$^{41}$那麼這幾十位牧者我們一起參加一個團隊建立的運動.
那麼其中一個的活動就是這個 trust fall (信任摩擦).
那麼我們三十多位牧者在這裡在聽指示.
於是那個教練就說.
現在你們記住.
待會你們會抽籤.
抽一個人出來做跳下來的那個.
其他的就請按照指示一起去組成這個安全網.
於是我們三十多位牧者同心協同.
包括我自己.
千萬不要抽到是姚珍啊.
千萬不要抽到我.
千萬不要是姚珍啊.
那麼原來有時候禱告是美夢英雲的.
那這三十多位牧者包括我自己一起同心祈禱主啊.
跳下來那個千萬不要是我的時候.
一抽.
就抽到姚珍啊.
然後呢.
就說不是啊.
真的是姚珍啊.
那麼大家很害怕是不是?.
其實我也很害怕啊我想說.
於是我就站在一旁.
我站在一旁就看著這三十多位牧者.
就按照這個教練組成了一個安全網.
我真的那時候很明白什麼是全心的信靠.
因為我不斷不斷祈禱主啊.
我待會怎麼跳下去這個安全網呢?.
我都深信下面那三十多位牧者同心了祈禱主啊.
我們待會怎麼能夠接著姚珍這樣跳下來呢?.
但是當三十多位牧者真的按照導師的指示組成了一個安全網.
那麼我就站在上面了.
站在上面之後教練就說.
姚珍跳吧.
真的跳啊?.
不是吧?.
跳吧.
接著我明明聽到我問真的跳而已.
誰知道原來下面安全網的人都說.

$^{81}$不是吧教練?.
真的跳啊?.
大家都在問是不是真的跳啊.
於是教練就說姚珍.
憑信心跳.
接著我就說.
啊!我真的跳了.
於是跳了下來.
原來真的厲害了.
這個安全網真的很輕易就接住我了.
然後我跟那個牧者聊天.
他很坦白的說.
剛才我以為我的手是會斷的.
但是居然.
原來我們做成一個安全網.
你掉下來的時候.
又不是真的這麼重而已.
感恩啊.
那麼跳下來會是怎麼樣呢?.
就是這樣.
跳.
那麼他們就接住我了.
但是原來有時候我們的生命.
當我們說跟隨神的時候.
其實我們有時候也很像這個例子的.
我不知道你有沒有試過.
你也問神.
神啊.
是不是真的這麼走的?.
這麼走是不是真的可以的?.
當聖經說不是啊.
你憑信心按著聖經的真理走.
你可能也會問.
神啊.
是不是真的這樣的?.
是不是真的去的?.
是不是真的這麼做的?.
生命裡面我們有時候也面對不同的掙扎.
可能我們也會問神.
我們是不是真的單單憑著信心.

$^{121}$單單按著你的心意.
單單按著聖經的教導和價值觀走就可以了.
也就是說我們其實可以跟隨這個世界.
所以今天我們要學習的一個功課就是.
今天我們的生命是不是單單的信靠.
而今天我們的生命.
又是不是全心的跟隨呢?.
今天的題目.
神國的真諦.
單純的信靠.
全心的跟隨.
讓我們一起在《路加福音》十八章那段聖經裡面.
我們一起去學習三個生命很重要的功課.
今天我們學習的第一個功課.
是我們要有好像小孩子那樣的信靠.
我們一起由第十八章的十五節.
讀到第十七節好不好呢?.
一起讀.
一二三.
有人抱著自己的嬰孩來見耶穌.
要他摸他們.
門徒看見就擲避那些人.
耶穌卻叫他們來說.
讓小孩子到我這裡來.
不要禁止他們.
因為在神國的正是這樣的人.
我實在告訴你們.
凡要承受神國的.
若不讓小孩子.
斷不能進去.
當我們看這段聖經的時候.
我們看到有一些人.
將那些嬰孩帶來見耶穌.
這嬰孩除了是說那些初生的嬰孩.
也包括一些年幼的小孩子.
人們抱著他們的小孩子來見耶穌.
很想耶穌摸他們.
很想耶穌賜福給他們的小孩子.
在《路加福音》其他的地方也講到.
耶穌伸手去摸人的時候.

$^{161}$代表著他能力的彰顯.
但是當人們將小孩子帶到耶穌的時候.
這裡告訴我們.
門徒看見就擲避那些人.
為什麼門徒會擲避那些人呢?.
因為當時小朋友在社會上的地位.
是很低微的.
現在小朋友有些地位很高的.
我曾經試過喝茶.
喝茶很吵的.
但是旁邊的桌子.
有一家人男女老友.
即是奶奶,爺爺一起.
突然間.
那個媽媽跟整個桌子的人說.
噓!.
不要走.
我說怎麼回事?.
寶寶睡著了.
你沒有人可以吵醒我的寶寶.
你看看現在小朋友的地位多高.
讓他大聲地噓.
讓他噓過來我的桌子.
我心想你是不是也在噓我?.
現在小朋友的地位很高.
但是當時的小朋友.
地位是很低微的.
所以在這裡有一個價值觀.
就是門徒覺得這些地位這麼低微的人.
他們是不配來到耶穌面前的.
我們之前看過法利賽人.
和瑞利禱告的比喻.
法利賽人覺得自己配得.
就站得很前.
然後就藐視別人.
那些世界上覺得最沒有地位.
最沒有功德的人.
你帶他去耶穌面前做什麼呢?.
不要浪費耶穌的時間.
在這裡反映一個價值觀.

$^{201}$就是人覺得不重要的事情.
就不要佔用耶穌的時間.
但是在這裡耶穌要顛覆他們的價值觀.
如果我們看第十六節.
耶穌顛覆了這些人的價值觀.
耶穌說什麼呢?.
「讓小孩子來到我這裡來」.
「不要禁止他們」.
但是接著耶穌說了一個更加顛覆的話.
他說為什麼呢?.
他說因為在神國裡面的.
就是好像這些小朋友一樣.
門徒覺得這些沒有地位.
沒有功德的小朋友.
耶穌反而說.
神國裡面的人就好像這些小朋友一樣.
那究竟什麼是「讓小孩子」呢?.
小孩子是會全然依靠父母的.
而「讓小孩子」就是代表我們的生命要全然依靠上帝.
我們不可以靠自己的行為.
或者我們的美德.
去得著神的悅拿.
原來今天我們要在神面前很謙卑.
我們來到神的面前.
我們承認我們自己的無能.
但是我們就單單領受.
單單信靠.
單單跟隨耶穌.
原來單單信靠.
全然領受.
單單跟隨耶穌.
是進入神國的模式.
所以告訴我們.
神國不是靠資歷,靠表現,靠地位.
而是單單靠著被主接納的恩典.
所以到了第十七節.
耶穌再一次強調.
他說「我實在告訴你們.
凡要承受神國的.
若不像小孩子.

$^{241}$斷不能進去」.
這裡不是你看完之後就突然吹噓.
「我是小朋友,現在我們一起進入神的國」.
不是這樣的意思.
這裡他講了很重要的東西.
凡承受神的國.
神國是我們承受.
是我們領受.
因為我們知道我們不配.
不是我們爭取回來.
我們沒有做什麼去爭取神國回來.
這裡不是叫我們變回一個很幼稚的小朋友.
這裡要我們變回一個很單純的心.
我們單純去信靠.
單純地去跟隨耶穌.
所以原來進入神國的關鍵.
不是成人式的自信.
控制,計算.
而是孩子式的謙卑.
單純信靠.
今天我們要問我們自己一個問題.
就是我們的生命是不是很單純地去跟隨耶穌呢?.
當我們面對生命的選擇.
面對生命的掙扎.
我們是單純地去選擇跟隨耶穌.
還是我們都是沒有信心.
所以我們有時跟隨神.
但有時我們覺得神不知道是否跟得過.
還是跟隨這個世界.
我曾經分享過.
也分享過不少次.
我信耶穌的時候是我在讀大學二年級.
當時暑假我有機會去北京讀一個法律的課程.
我知道我有很多時間看書.
所以那時候我就想.
那個月裡面我看一本什麼書好呢?.
當時有兩個人介紹了兩本書給我.
我記得當時.
就是我律師行做暑期工的老闆介紹了一本書給我.
叫做《職場上位秘笈》.

$^{281}$這本書很精彩的.
它可以教你如何在職場上.
如果你背了鍋.
如何將鍋子射向你身邊的那個.
另外它會教你如何拆老闆的鞋.
拆得十分自然.
讓你老闆不覺得他在拆你的鞋.
其實我講過很多次之後.
第一次他會說那本書在哪裡.
我找不回來.
但是我另外一本書.
就是當時我英國傳道人給我的就是《基本基督教》.
我們是不是單單信靠我們的主耶穌.
但是我想講多一點.
這個我還沒講過.
其實你可能覺得.
那你一個月的時間.
看一本書有什麼重要?.
其實當時我去選擇那本書的時候.
代表我心裡面一個價值觀的選擇.
我剛剛信了耶穌.
我剛剛在人生的交叉點.
我快要大學畢業.
我快要出來工作了.
我的價值觀究竟是什麼呢?.
我的價值觀是基本基督教.
還是職場上位秘笈呢?.
曾經有一個人聽完我這樣講之後.
他說很容易而已.
你星期一至星期六就是職場上位秘笈.
星期日就是基本基督教.
到了你回到星期一至星期六的時候.
又是職場上位秘笈.
是不是有鑊就卸.
有鞋就拆.
然後你回到星期日的時候.
你就是基本基督教.
但是今天我們的生命不是這樣的.
我們的生命是要單單跟隨耶穌.
但是什麼叫做單單跟隨耶穌呢?.

$^{321}$就是在你生命的掙扎.
在你生命的抉擇裡面.
究竟你的價值觀是什麼呢?.
你是不是單單相信耶穌?.
我曾經跟一位肢體聊天.
他說很多年前.
當他幫他的小朋友入學的時候.
昔日如果你住一個地址.
你就可以進到那一區的學校了.
於是有人跟他說.
買一個地址就可以了.
你買一個通訊地址.
你用那個通訊地址.
當然有人賣地址出來給你的.
你買一個地址.
用了那個通訊地址.
你就進到那一區的學校了.
但是那時候他們有一個掙扎.
他說但是這個是不合神心意.
是不誠實的.
於是他朋友跟他說.
搞什麼神的心意.
個個都是這樣做的.
你不做你就吃虧了.
但是那一刻你要問自己.
我是不是單純地去跟隨神.
就是我跟隨這個世界呢.
但是當我們想深一層.
其實這個不只是一個行為的抉擇.
其實當你去想你的小朋友.
讀什麼學校的時候.
其實都提醒我們.
我們人生的預備和計劃裡面.
究竟我們是不是按著神國的價值觀.
去做我們人生的抉擇和預備.
當你想一個小朋友入學的時候.
其實不只是說一個地址這麼簡單.
究竟你為他選學校的時候.
你是用屬世的價值觀.
還是用屬靈的價值觀呢.

$^{361}$當然你用的方法是世界的方法.
還是神的方法.
但是最後其實你都想.
當你的小朋友成長的時候.
你想他的生命.
他的前路.
是世界的價值觀.
還是神國的價值觀.
所以原來今天很重要.
我們要問我們自己.
我們的生命.
有沒有好像小朋友那樣去單單信靠神.
而當我們說主啊我願意.
我的生命好像小朋友那樣.
單單信靠你的時候.
我們就去到今天分享的第二點.
就是原來我們要清除.
難阻跟隨耶穌的障礙.
我們一起由第十八讀到第二十五節.
好嗎?一起讀.
一,二,三.
有一個官問耶穌說.
良善的夫子,我該做什麼事.
才可以承受永生?.
耶穌對他說.
你為什麼稱我是良善的?.
除了神一鬼之外.
再沒有良善的.
誡命你是曉得的.
不可姦淫.
不可殺人.
不可偷盜.
不可作假見證.
當孝敬父母.
乃人說.
這一切我從小都遵守了.
耶穌聽見了.
就說.
你還缺少一件.
要變賣你一切所有的.

$^{401}$分給窮人.
就必有財寶在天上.
你還要來跟從我.
他聽見這話.
就甚憂愁.
因為他很富足.
耶穌看見他.
就說.
有錢財的人.
盡神的國.
是何等的難?.
駱駝穿過針的眼.
給財主盡神的國.
還容易呢?.
這裡開始有一個人問耶穌.
這裡聖經稱呼他為一個官.
除了官之外.
我們沒有其他的背景.
不過當時的官.
是一個猶太人的領袖.
這條問題表面是挺熟悉的.
良善的夫子.
我該做什麼事.
才可以承受永生.
表面很熟悉.
但其實前設錯了.
因為其實他是在問.
我可以做些什麼.
來承受永生.
他的前設就是永生是可以靠我做的事情來換取回來的.
是一個我可以用我做的事情所完成的任務.
所以在這裡耶穌要去拆除他的觀念.
所以耶穌在第十九節跟他說.
他說.
你為什麼稱呼我是良善的.
除了神一位之外.
再沒有良善.
這裡耶穌不是否定自己.
因為我們在路加福音裡面.
我們已經看到.

$^{441}$路加福音強調.
耶穌就是神的兒子.
而耶穌和上帝是同長王權.
耶穌也是末世救贖計劃的承傳者.
所以在這裡其實我們搞清楚.
耶穌不是否定自己.
相反耶穌說如果得神是良善.
而我和神有這個關係.
你唯一的出路就是來跟隨我.
在以賽亞書裡面這麼說.
唯有我是耶和華.
除我以外沒有救主.
而在《仕途行傳》路加這麼說.
除他這個耶穌.
除他以外別無拯救.
因為在天下人間.
沒有次下別的名.
我們可以靠著得救.
其實耶穌這裡說的.
不是你做什麼能夠有永生.
你只有一種方法可以得著永生.
就是跟隨我.
並且你要去撇下一切跟隨我.
才能夠進入神的國.
因為對於一個猶太人來講.
良善的夫子.
就是遵守律法的人.
這個人也告訴耶穌我遵守律法的.
但是耶穌這一段話.
就是要將他的心思意念.
從去追求律法的遵守.
轉去你的生命.
要跟隨耶穌.
如果你說神是良善.
耶穌是良善.
你認定了耶穌和神的獨一無二的關係.
你就要撇下一切跟隨耶穌.
這是唯一進入神國的途徑.
所以在第二十節.
耶穌就說.

$^{481}$律法你是知道的.
不可姦淫.
不可殺人.
不可偷渡.
不可作假見證.
當孝敬父母.
這五條的誡命都是來自十誡的.
分別是第七誡,第六誡.
第八誡,第九誡和第五誡.
耶穌將孝敬父母放在最後.
可能強調就是孝敬父母是重要的.
但是今天孝敬父母不代表我們敬拜獨一的真神.
這個官他做了一切人際關係可以守的律法.
但是他忘記了十誡最重要的一誡是什麼?.
你要跟隨敬拜真神.
所以其實耶穌說第二十節是還沒說完的.
但是那個人已經搶著說了.
你說的這些事情.
我已經所有都做了.
代表這個官他還沒明白怎麼才能進入神國.
人不可以靠自己的行為進入神國.
人唯有靠跟隨獨一的真神.
撇下一切.
跟隨獨一的真神.
才可以進入神國.
所以耶穌說了誡命.
他急不及待已經說這些誡命我守了.
然後耶穌在第二十二節就說了.
他說你還缺少一件.
你要變賣你一切所有的分級窮人.
就必有財寶在天上.
你還要來跟從我.
這裡耶穌不是在說理財之道.
這裡是告訴他你在欠缺什麼.
這裡的重點不是要求我們每一個人都變賣我們所有分比窮人.
但是因為耶穌知道在這個人的心裡面.
他的偶像就是他的財富.
為什麼這個人的偶像是他的財富呢?.
因為當時猶太人有一個錯誤的思想.
他們認為如果有一個人有很多錢的話.

$^{521}$這個就是代表神對他的眷顧.
一個人有很多錢.
代表他有神的祝福.
所以我不可以沒有錢.
因為沒有錢就等同沒有了祝福.
這個觀念是錯的.
所以耶穌要去搞清楚他的觀念.
今天你已經當這些是你的偶像.
而這些東西阻礙你去跟隨耶穌.
因為在舊約第十誡.
就是說我們不可以貪戀別人的妻子.
不可以貪戀別人的財產.
為什麼?因為你貪這些的時候.
這些東西成為了你的偶像.
你就犯了裡面最重要的誡命.
就是除了我以外不可有別神.
所以其實耶穌說.
當財富成為了你的偶像的時候.
你就要除掉這個跟隨我的障礙.
但是在路加福音第十九章我們之後會看的.
撒該將他一半的財產給了窮人.
他也得著救恩.
所以將所有財產給窮人.
不是一個絕對性的要求.
在路加福音第十章也有其他人問耶穌同一個問題.
我為什麼得著永生.
耶穌是用另外一個誡命回應他的.
為什麼?因為耶穌是按著每一個人的情況去回應.
但是對於這個人來說.
如果財富成為了他的偶像.
他就要去清除這個障礙.
讓他的生命能夠真正去跟隨耶穌.
所以今天.
原來我們要移除障礙.
全心全意去跟隨耶穌.
因為跟隨耶穌是進入神國的唯一的途徑.
我曾經聽過一個比喻.
有一個人掉了一條金條在一個瓶子裡.
於是他伸手進入那個瓶子裡想拿回那條金條.
但是不知道是不是他的手跟我一樣比較胖.

$^{561}$他塞進了瓶子裡.
他抓住了一條金條.
但是他的手吹不回來.
於是他搖來搖去.
怎麼辦?.
然後他旁邊的朋友跟他說.
他說你現在要先放下那條金條.
是不是?.
你不放下那條金條.
其實你的手是吹不回來的.
但是原來今天我們的生命也是這樣.
我們有很多東西抓得很緊.
而緊到那樣東西成為了我們生命的偶像.
那樣東西就難阻我們去跟隨耶穌.
所以今天如果你抓住的.
我不知道抓住的是什麼.
你抓住的可能是金條.
你抓住的可能是其他東西.
你要放手.
你要清除這個障礙.
你才能夠真正跟隨耶穌.
但是今天當我們看二十三節的時候.
當這個官聽了這番話的時候.
聖經說什麼?.
就甚憂愁.
因為他很富足.
這裡是一個很大的諷刺.
因為這個人的憂愁程度.
跟他的富足成正比.
我們以為人經常覺得財富應該為我們帶來滿足.
但是這個人的財富不但沒有為他帶來滿足.
反而成為了進入神國的障礙.
成為他的憂愁.
這個人其實不是不明白道理的.
但是他割捨不了.
所以原來今天最大的悲劇.
可能不是說我們的生命公開反對耶穌.
可能最大的悲劇就是我們明白真理.
但是我們不捨得放手.
所以耶穌在二十四節就說.

$^{601}$他說有錢財的人進神的國是何等難.
駱駝穿過針的眼.
比財主進神的國還容易呢.
這裡是一個很強烈很誇張的形容.
你想想駱駝.
再大隻一點.
針的眼是針的洞.
現在我有時候老花也看不到針的洞在哪裡.
但是他說駱駝穿過針的洞.
比有錢人進神的國更容易.
當然耶穌用這麼誇張的手法.
他是想震撼我們的心.
他再一次說今天你不要靠自己.
如果今天你以為靠你的財富.
靠你的財幹.
靠你的智異.
你以為可以跟耶穌.
不行,你要放手.
唯有你放手.
你才能夠讓耶穌去帶領你的生命.
所以今天我們要問.
究竟我們的生命裡面有些什麼.
是難阻我們跟隨耶穌的呢.
我很喜歡那個吃自助餐的比喻.
如果今天晚上你來金會吃自助餐.
你幻想一下今天晚上你訂了城中最好吃的自助餐.
你下午會怎樣的.
你會不吃東西的.
你下午會不會去嘟華吃個燒賣.
然後再嘟記又再吃一個燒賣.
然後再去吃一個腸粉.
再去吃一碗燒鵝奶你會不會.
你不會的.
因為當你吃完這些東西的時候.
晚上你就吃不到自助餐的了.
但有時候我們人生就是這樣.
我又抓住這一樣.
我又抓住那一樣.
這一樣東西我又放不下.
那一樣東西我又放不下.

$^{641}$到最後反而我們生命不能夠跟隨耶穌.
所以我們要問自己.
究竟今天什麼是我們生命的障礙呢.
會不會是我們的工作.
我們的物質.
我們的財富.
我們的投資.
我們的消費.
會不會是我們的家庭關係.
我們將配偶.
將子女放在生命的首位.
又或者是我們生命有一些的恐懼.
我們很怕失去.
或者我們對舒適的追求.
又或者會不會是社交媒體.
是不是.
在我成長的年代.
我們不會看著一個手機看這麼久的.
但是現在我們人可以看著一個手機看很久的.
還有我也試過你看了很久.
你不覺得你看了這麼久的.
你看了這麼久之後.
你又發現原來他過了一個小時.
我們被這些東西吸引了.
可能不單止是這樣.
網上的娛樂.
網上的遊戲.
我發現我成長的年代.
是年輕的一輩才被網上的遊戲去沉迷.
現在原來是什麼年紀都能夠被網上的遊戲去沉迷的.
我想說.
真的.
又或者是煲劇.
我成長的年代.
煲劇是一晚的事情.
因為一晚.
因為以前的劇最多都是二十幾三十集而已.
不是你一晚你怎麼都煲完了.
但是現在是一百幾二百集的.
你可以將你的面貼在螢幕上.

$^{681}$然後出不來的.
所以我們要問自己究竟什麼.
是我們跟隨耶穌的障礙.
有些什麼我們需要去清除.
因為如果不清除的時候.
我們不能夠跟隨耶穌.
當我們決定我們的生命要單單跟隨耶穌.
當我們說我們要清除跟隨耶穌的障礙.
就到了今天我們的第三個功課.
就是我們要校正生命的優先次序.
我們讀最後二十六至三十節好不好.
我們一起讀.
一二三.
聽見的人說.
這樣誰能得救呢.
耶穌說.
在人所不能的事.
在神卻能.
彼得說.
看!我們已經撇下自己所有的根從理了.
耶穌說.
我實在告訴你們.
人為神的國撇下房屋.
或是妻子弟兄父母兒女.
沒有在今世不得百倍.
在來世不得永生的.
在一開始聽到的人其實很詫異.
當他們聽到財主不能夠進入神的國的時候.
這些人很詫異,為什麼呢?.
因為這裡的群眾對於救恩有兩個誤解.
因為我剛才說.
猶太人覺得錢就是神的眷顧.
這個人這麼有錢.
神這麼眷顧他.
如果神這麼眷顧的人都進不了神的國.
誰能進入神的國呢?.
另外他說他所有的誡命都遵守了.
一個人覺得神這麼眷顧他.
他又遵守誡命的人都進不了神的國.
誰能夠進入神的國呢?.

$^{721}$另外一方面.
人們就想.
如果人不能夠靠自己的身份.
不能夠靠遵行律法去進入神的國.
那麼誰能夠得救呢?.
所以這一刻的群眾是很詫異的.
但是耶穌面對他們這個問題就回答了.
在人所不能的是在神卻能.
這裡也有兩個很重要的意義.
第一個是神恩典的原則.
人不是靠自己的能力能夠親近神.
但是神可以拯救那些不配的人.
哇!這是一個好消息.
人不靠自己的能力可以進入神的國.
但是神可以去救那些不配的人.
你說誰是不配.
我們每個人都不配.
所以我們看看自己.
你發達了.
因為神會拯救那些不配的人.
但是今天我們怎樣能夠進入神的國呢?.
是我們要靠著神自己主動的工作.
是耶穌來到這個世界.
為我們釘十字架.
所以我們要承認我們的不配.
承認我們的無能.
但是因為知道有耶穌.
我們的生命就單單去跟隨耶穌了.
因為這一段第二個重要的意義.
就是耶穌告訴這些猶太人.
神和你的祖宗亞伯拉罕立約的時候的應許.
就藉著耶穌成就了.
是靠著耶穌成就對亞伯拉罕所納的約.
所以昔日亞伯拉罕回應神的呼召.
今天我們都要撇下一切.
生命跟隨耶穌.
在這個時候彼得就問.
他說主阿罕啊.
我們已經撇下自己所有去跟隨你了.
彼得這裡說得對的.

$^{761}$這個官不願意捨棄自己的財富.
但是門徒已經撇下一切去跟從耶穌.
但是彼得這樣說.
這個官不願意跟隨你.
但是我們已經撇下一切跟隨你的時候.
其實彼得這裡有兩個很重要的問題.
第一就是.
主啊我們撇下一切之後.
我們能不能進入神的國.
第二.
我進入神的國的時候.
我們能夠得到什麼.
所以第29節耶穌就這樣回答.
他說我實在告訴你們.
人為神的國撇下房屋.
或是妻子弟兄父母兒女.
今日人要從定他的優先次序.
為神的國撇下房屋.
妻子弟兄父母兒女.
房屋不只是說一層樓或者一個建築物.
房屋其實也是映作一個家.
家就是人際網絡關係的中心.
這裡耶穌對門徒有一個絕對性的要求.
但是這裡你先不要搞錯.
你不要突然間就說.
這樣就好了.
你現在走出去崇拜之後.
就跟你的父母說以後不用養你了.
因為我為了神的國撇下了父母.
然後有些人走出去就跟他的兒子說.
嘩!真是感動了.
以後不用給你零用錢了.
因為我為了神的國撇下了兒女了.
不是這樣的意思.
這裡是說你要搞清楚你的優先次序.
將神擺在你生命的首位.
面對神國或者最親密的關係.
神要在生命裡面居首位.
就是說當神國和家庭的價值觀有衝突的時候.
神國的價值觀要居首位.

$^{801}$但是我想說你絕對不需要擔心.
因為其實過去做了這麼多屆門女.
有門徒的配偶和家人來問我的.
那是不是以後他不用理我的了?.
他是不是用光所有的時間去做門徒了?.
是不是他以後不用理家人了?.
他用光時間做門徒.
我說不是.
每一次我都說我保證一件事.
當你的家人你的配偶做了門徒的時候.
他會更加愛你.
為什麼?.
因為他會懂得用耶穌的愛去愛你.
那你說肥蛇什麼是耶穌的愛?.
耶穌的愛就是社稷的愛.
今天如果你的配偶用社稷的愛去愛你.
你發不發達了?.
發達了是不是?.
今天你的家人用社稷的愛去愛你.
你發不發達了?.
發達了.
所以如果坐在你旁邊的配偶或者家人.
你做門女.
你快點鼓勵他去做門徒.
他會學習用社稷的愛去愛你.
你發達了.
所以在這裡不是在說我們不理我們的家人.
而是我們要較正我們生命的優先次序.
但是可能你就會問.
那我怎麼去較正我生命的優先次序呢?.
我覺得有一個很重要的關鍵.
這個重要的關鍵是什麼呢?.
就是我們要捉住我們今天最重要的身份是什麼.
今天我們要捉緊我們作為門徒的身份.
我們就會搞清楚我們生命的優先次序了.
你想想彼得.
彼得跟耶穌之前他做什麼的?.
他是漁夫來的他打魚的.
你可以想像一下這個虛擬的場景.
彼得在打魚.

$^{841}$魚網很重.
然後耶穌走過來說彼得.
你要撇下一切跟隨我.
然後彼得就說.
主啊.
我真的很想跟隨你的.
不過我是一個漁夫.
作為一個漁夫我是不可以放下漁網的.
因為如果放下漁網我就不是一個漁夫了.
所以我雖然很想跟隨你.
但是我不能夠放下漁夫的身份.
所以我是不會放下漁網的.
裡面還有很多聖彼得魚.
其實聖彼得魚是紀念聖彼得才有的.
所以我經常不是叫聖彼得魚.
我們很多時候就是這樣.
我們沒有搞清楚什麼才是最重要的.
所以雖然今天耶穌叫我去做門徒.
但是今天我做漁夫很重要的.
所以我是不會放下漁網的.
所以我放下不了漁網.
因為對於我來說漁夫這個身份就是最重要.
但是如果今天我們搞清楚我們的身份.
我們就會搞清楚我們生命的優先次序.
如果我搞清楚我作為門徒的身份.
我會帶著門徒的身份回到我的家庭.
我就知道我會用聖經的價值觀去愛我的家人.
愛我的父母.
愛我的兒女.
今天我帶著門徒的身份進入職場.
我就知道我是要按聖經的價值觀去做我的工作.
我就會在職場裡面去活出一個為主做的職場的見證.
今天我抓住這個門徒的身份進入社區.
我就知道今天神叫我們要去照顧社區裡面的孤兒寡婦.
所以我就會去幫助那些有需要的人.
今天我抓住了這個身份.
我就能夠重訂我生命的優先次序.
所以原來今天當我們抓住門徒這個身份.
我們就會校正好我們生命的優先次序.
在結束的時候我們請敬拜隊.

$^{881}$今天我們學了三個很重要的功課.
這三個功課是有次序的.
我們首先跟自己說.
我們今天我們願不願意去單單單純地去跟隨耶穌.
你要問自己我是不是已經決定了.
最近我又唱了一首很舊的歌.
跟那些少年人唱.
我願意跟隨著耶穌.
我已經決定跟隨著耶穌.
有人唱過嗎.
有的話你應該回到舊的一段時間.
然後永不回頭.
永不回頭.
搞得我好像石器時代那樣.
應該有人懂的.
當你決定好像小孩子那樣去信靠的時候.
接著你就要開始清除你裡面的障礙.
有些什麼現在是攔阻著我去跟隨耶穌的呢.
我講過這個例子.
曾經有一次有一個肢體.
他就申請入門旅.
但是他就跟我說對不起Eugene.
我不能夠參加你這個門旅.
我說為什麼.
他說因為過去幾十年.
我都是星期二打網球的.
我不能夠放棄網球的.
然後他跟我說一句英文.
Tennis is life.
有時候我們就是捉住了那件事.
Tennis is life.
我不能夠放下那件事.
所以我不能夠.
但我說不是.
其實我們有班弟兄很好.
他們現在逢星期六晚上打網球.
Alan.
星期六跟他打.
星期二去門旅就行了.
但是有時候我們放下不了.

$^{921}$但是當我們放下了的時候.
我們就要捉緊那個身份.
當你捉緊那個身份的時候.
你就會懂得怎麼去校正生命次序.
彼得如果還覺得自己是一個漁夫.
他不會放下漁網的.
但是他搞清楚.
他是跟隨耶穌的門徒.
他就撇下一切跟隨耶穌.
弟兄姊妹今天讓這個成為我們每一個人的提醒.
也讓我們再一次去立志.
邀請大家再一次起立.
我們就用這一首詩歌.
去回應今天神給我們的呼召.
神給我們的教導.
我們的生命就單單.
將我們的心思意念定精在耶穌.
就讓我們的生命全然為你.
因為你是我們的主耶穌.
\newpage



\section{創世記 25:19-34}
\label{sec:032026}
\textbf{長子或紅湯}~~~~梁智程傳道
\newline
\newline
連結: \href{http://www.church.com.hk/acms/content.asp?site=cdc\&op=show\&type=product\&code=032026\&layout=sermon}{\texttt{032026}} ~~~~ 語音日期: 2026-04-18
\newline
\newline
\hyperref[sec:032025]{\small{< < < PREV SERMON < < <}}
~
\hyperref[sec:index_chronic]{\small{[返順時目]}}
~
\hyperref[sec:032027]{\small{> > > NEXT SERMON > > >}}
\newline
\newline
創世記 25:19-34
\newline
\begin{longtable}{cl}
\hline
\hline
章節 & 經文 (和合本修訂版)\\
\hline
25:19 & \begin{tabularx}{0.7\textwidth}{X} 這是亞伯拉罕的兒子以撒的後代。亞伯拉罕生以撒。 \end{tabularx} \\ \\ \relax
25:20 & \begin{tabularx}{0.7\textwidth}{X} 以撒四十歲時娶利百加為妻。利百加是巴旦‧亞蘭地的亞蘭人彼土利的女兒，是亞蘭人拉班的妹妹。 \end{tabularx} \\ \\ \relax
25:21 & \begin{tabularx}{0.7\textwidth}{X} 以撒因他妻子不生育，就為她祈求耶和華。耶和華應允他的祈求，他的妻子利百加就懷了孕。 \end{tabularx} \\ \\ \relax
25:22 & \begin{tabularx}{0.7\textwidth}{X} 胎兒們在她腹中彼此相爭，她就說：「若是如此，我為甚麼會這樣呢？」她就去求問耶和華。 \end{tabularx} \\ \\ \relax
25:23 & \begin{tabularx}{0.7\textwidth}{X} 耶和華對她說：兩國在你腹中；兩族要從你身上分立。這族必強於那族；將來大的要服侍小的。 \end{tabularx} \\ \\ \relax
25:24 & \begin{tabularx}{0.7\textwidth}{X} 到了生產的日期，看哪，腹中是對雙胞胎。 \end{tabularx} \\ \\ \relax
25:25 & \begin{tabularx}{0.7\textwidth}{X} 先出生的身體帶紅，渾身有毛，好像皮衣；他們就給他起名叫以掃。 \end{tabularx} \\ \\ \relax
25:26 & \begin{tabularx}{0.7\textwidth}{X} 隨後，以掃的弟弟也出生，他的手抓住以掃的腳跟，因此給他起名叫雅各。兩個兒子出生時，以撒六十歲。 \end{tabularx} \\ \\ \relax
25:27 & \begin{tabularx}{0.7\textwidth}{X} 兩個孩子漸漸長大，以掃善於打獵，常在田野；雅各為人安靜，常住在帳棚裡。 \end{tabularx} \\ \\ \relax
25:28 & \begin{tabularx}{0.7\textwidth}{X} 以撒愛以掃，因為常吃他的野味；利百加卻愛雅各。 \end{tabularx} \\ \\ \relax
25:29 & \begin{tabularx}{0.7\textwidth}{X} 有一天，雅各熬了湯，以掃從田野回來，疲憊不堪。 \end{tabularx} \\ \\ \relax
25:30 & \begin{tabularx}{0.7\textwidth}{X} 以掃對雅各說：「我累死了，請你讓我吃這紅的，這紅的湯吧！」因此以掃又叫以東。 \end{tabularx} \\ \\ \relax
25:31 & \begin{tabularx}{0.7\textwidth}{X} 雅各說：「你今日把長子的名分賣給我吧。」 \end{tabularx} \\ \\ \relax
25:32 & \begin{tabularx}{0.7\textwidth}{X} 以掃說：「看哪，我快要死了，這長子的名分對我有甚麼用呢？」 \end{tabularx} \\ \\ \relax
25:33 & \begin{tabularx}{0.7\textwidth}{X} 雅各說：「你今日對我起誓吧。」以掃就向他起誓，把長子的名分賣給了雅各。 \end{tabularx} \\ \\ \relax
25:34 & \begin{tabularx}{0.7\textwidth}{X} 於是雅各把餅和豆湯給了以掃，以掃吃喝以後，起來走了。這樣，以掃輕看他長子的名分。 \end{tabularx} \\ \\
[1ex]
\hline
\hline
\end{longtable}
$^{1}$弟兄姊妹平安.
再一次一齊敬拜.
盼望上主喜悅我們對祂的尊崇.
又期待上主向我們說話.
仍然在我們中間有祂自己的工作.
由我們開始講道和聽道的服事之先.
我們一齊同心向上主禱告.
自我啟示的三一上帝.
求你使用孩子的口.
宣講你自己的說話.
又求你祝福在座每一位在聽的弟兄姊妹的耳朵.
都能夠聽得清楚你自己的聲音.
幫助我們整個窩福群體.
每一個人的心都向你開放.
樂意去回應你的話語.
禱告奉主明求.
不知道大家最近有沒有玩一個手機遊戲.
叫做皮克文Pikmin.
這個遊戲近日在年青人.
或者你覺得比較年青的群體中間是頗為流行的.
打開IG,Facebook這些社交媒體的時候.
周圍都是關於這個遊戲的影片攻略.
皮克文這個遊戲其實很簡單.
就是你每天透過行路,即是步行.
在你行路的過程中.
遊戲中的角色皮克文.
就會幫你去種花,探索,收集資源.
同時他會慢慢累積了一定的友好度.
當這些皮克文的友好度成長到某一個階段.
他們就會穿上不同類型的飾品.
即是會換上衣服.
例如有些會有飛機的形狀.
有些會有漢堡包薯條的形狀.
有些會有拉麵的形狀.
弟兄姊妹,如果你到現在都不知道這個遊戲有多熱門.
我可以告訴你,在我們教會的初職團體中.
有超過九成的弟兄姊妹都在玩.
甚至連一些在海外的團友都一起玩.
這個遊戲就成為我們彼此之間一種另類的聯繫.
久而久之,我都受到我的太太和這班弟兄姊妹的薰陶.

$^{41}$結果我都下載了這個遊戲,一起種花.
初初其實我完全不知道自己在做什麼.
有花就去種,有這些資源就去撿.
但玩得久,你就會慢慢發現.
其實這個遊戲需要玩家不停去做不同的選擇.
例如當你看到路邊有些很小的資源.
你就會想,我要不要停下來慢慢去撿他們呢?.
還是我繼續走去完成另一個任務呢?.
又或者我現在出門口或者回家的時候都會想.
好不好不要坐車呢?改去走路或者踩單車.
方便自己去這些偏遠地區去種花.
或者去做不同的任務.
有一次我和我太太在逛街.
她在路上突然大叫.
她說出事了.
都要關心一下發生什麼事.
原來在她的遊戲地圖出現了一個限時的稀有飾品花苗.
很簡單,你去撿就可以了.
但她同時在培養另外一些花苗.
導致她的袋子沒有足夠的空間去收集.
結果她就眼白白要錯過這些平時比較難見的稀有花苗.
之後她就說了一句話,早知就不種這個盤.
其實遊戲裡面的經歷都相當之真實.
因為我們每一天的生活都同樣面對著大大小小類人的選擇.
有時候我們會為了一樣東西而錯過了另一些的機會.
甚至乎是一些更加大,更加長遠,更加珍貴的事情.
而今天的經文正正都是在說一個關於選擇的故事.
讓我們回到這段經文裡.
看看上主透過這個訊息對我們有什麼提醒.
創世記二十五章十九到三十四節.
可以說是一個序曲.
是整個關於雅各和爾素的故事的一個開端.
經文一開始,作者用「這是亞伯拉罕的兒子,以撒的後代」作為標記.
這是創世記裡反覆出現的一個重要結構標記.
去標示一個新的敘事單元的開始.
這裡不單止是一個家譜性的過渡.
更加是上帝應許歷史的一個延續.
顯示故事由上一代的亞伯拉罕.
已經轉移到下一代的人物,以撒身上.
但是有趣的是.

$^{81}$就算名義上是在說以撒的後代.
但實際看來敘事的焦點根本就不在以撒身上.
而是很快就轉到他兩個兒子身上.
就是爾素和雅各.
接著,作者在二十字的經文交代了.
以撒娶了利伯加為妻子.
並且很特別提到他來自亞蘭地,是拉班的妹.
這些細節不單止提供一些背景資料.
而是為了雅各將來的故事鋪陳.
為了雅各日後逃離到外家,在拉班家中生活多年的片段.
埋下一些伏線.
值得我們注意的是.
這段經文延續了創世記一個反覆出現的主題.
就是一個族長,一個一家之主.
他的老婆面臨著不育,即是生不到孩子的問題.
正如撒拉曾經不育,利伯加同樣面對著這個困境.
似乎這不只是個人身體條件的問題.
而是更加關乎到上帝的應許究竟是否真的能夠實現的問題.
因為上帝曾經應許過阿巴拉罕的後裔要成為一個大國.
但若果沒有下一代的出身,這個應許就會中斷.
因此,不育的背景凸顯了一個神學主題.
就是應許的成就並不是自然就會發生的事.
而是完完全全需要依賴上主的介入才能成事.
然後,去到第21至26節.
經文描述了以素和雅各兩兄弟還未出世時的情況.
以撒因他妻子不生育就為他祈求耶和華.
耶和華應運他的祈求,他的妻子利伯加就懷了孕.
故事由以撒的祈禱開始.
作者讓我們知道,因為利伯加不能夠生育.
以撒就為他向上主祈求.
而上主垂聽了他的禱告,讓利伯加成功懷孕.
這一點其實很重要,因為他再次強調.
生命的誕生並不是偶然,而是上主恩手的介入.
以撒和亞伯拉罕一樣扮演了一個代求者的角色.
而上主也再次顯明他是那位聽禱告和施恩的主.
不過,利伯加的懷孕並沒有為他們帶來一個純粹的喜樂.
反而伴隨著一個極大的痛苦.
22節,託兒們在他福中彼此相爭.
他就說:「若是如此,我為什麼會這樣?」.
他就去求問耶和華.

$^{121}$作者形容利伯加託中的託兒彼此相爭.
相爭到令利伯加發出一個強烈的疑問.
「若是這樣,我為何活著?」.
這句話反映了他內心的困惑和掙扎.
因為這一刻他正在經歷的是一個身體上極大的痛苦不息.
同時也是一種生命的困擾.
以至他質疑繼續活下去是否有意義.
這種痛苦是否真的有價值.
於是利伯加去問上主,而耶和華上帝在此給予他一個很關鍵的啟示.
「兩國在你福中,兩族要從你身上分立.
這族必強於那族,將來大的要服侍小的」.
這一番來自上帝的說話.
首先預言指出利伯加託中的不只是兩個孩子.
而是兩個民族的起源.
其次他也預示了當時的一個社會常規.
就是長子理應擁有優先的地位和繼承權.
將來要在這個家庭裡打破.
耶和華上帝宣告大的要服侍小的.
顯示上主的揀選並不受人類的制度文化傳統所限制.
是完完全全出於他自己的主權和旨意.
當生產的日子來到.
二十四到二十六節的經文講述了兩兄弟出生的經過.
到了生產的日期,看,福中是對雙胞胎.
先出生的身體大雄,渾身有毛,好脹皮衣.
他們就給他起名叫耳掃.
隨後耳掃的弟弟也出生.
他的手抓住耳掃的腳筋.
因此給他起名叫雅各.
兩個兒子出生時已達六十歲.
耳掃首先出了泥,身體發紅,全身多毛.
因此得到這個名.
然後弟弟雅各也出生.
他的手相當戲劇性地抓住耳掃的腳筋.
結果原本在媽媽肚裡的爭鬥延續到肚外.
雅各抓住哥哥的腳筋的動作.
象徵著他拼命想追趕甚至取代哥哥的位置.
而雅各這個名字本身帶有抓住腳筋的意思.
某程度上作者在這裡也想預示了.
雅各這個人物將來的性格和行動.
依照作者的記載.

$^{161}$我們知道當利伯加生兒子的時候.
以撒已經六十歲.
這個時間性的標示顯示出.
由以撒最初的代求到利伯加真的成功懷孕.
中間其實是經歷了很長的時間.
是足足二十年.
而耳掃和雅各的出生.
亦成為了以撒禱告的答案.
或許這一刻這個提醒也告訴我們.
他們出生的結局所帶來對他們家庭的喜悅.
已經完完全全蓋過了.
太太在懷孕期間所經歷的痛苦和疑慮.
當我們看到這裡 作者也在告訴我們一個訊息.
就是上主的狀為確實不到人的表現去限制.
而是按照著上主自己的旨意去推進.
人類的歷史說到底都是在上帝的主權之下去開展.
接著創世記的作者描述到.
耳掃和雅各兩兄弟的出現.
為了以撒這個家庭帶來的轉變.
二十七到二十八節.
兩個孩子漸漸長大.
耳掃善於打獵 常在田野.
雅各為人安靜 常住在帳棚.
耳撒愛耳掃 因為常吃他的野味.
李伯嘉卻愛雅各.
當這兩兄弟慢慢成長後.
他們之間的差異也慢慢浮現.
原本隱藏在這個家庭裡的張力也開始揭露出來.
經文很簡潔但鮮明地描繪了兩人的分別.
耳掃善於打獵 是一個常在田裡出沒的人.
而雅各是一個安靜的人 住在帳棚裡.
很明顯作者不是純粹想交代這兩個人在職業上的不同.
而是很想反映出這兩兄弟.
在性格上 在生活方式上存在著很大的分別.
耳掃這個人性格外向 行動主導.
過著一個貼近大自然的生活方式.
而雅各性格相對內向.
並且與他家庭的聯繫相信比較深.
有些聖經學者根據常住在帳棚裡這句話.
推斷雅各很可能與他的父親和祖父一樣 是一個牧羊人.

$^{201}$但真正令到這個家庭裡張力加劇的 不單是這兩兄弟性格和工作上的不同.
而是他們父母的偏愛.
經文說 以撒愛耳掃 因為他常吃耳掃所打的野味.
利伯加則愛雅各.
作者也特別展示了耳掃性格的另一面.
他不單是一個比較被動或會禱告的人.
同時他是一個喜歡吃東西的食家.
因為以撒對食物的偏愛.
竟然可以導致這兩兄弟之間的緊張關係加劇.
尤其是當他老婆利伯加偏愛雅各的時候.
經文沒有詳細說明為什麼利伯加偏愛雅各.
可以讓我們聯想揣測一下.
會不會是因為雅各比較多留在家裡?.
會不會是因為雅各比哥哥更易受控制?.
無論利伯加實質的原因是什麼.
這種家庭內部的分裂.
令原本已經存在在兩兄弟之間的差異變得更尖銳.
在這個家裡 這對父母沒有成為一個調和的力量.
反而無意間站在不同的陣營.
讓這兩兄弟的對立去生發.
靚姐妹 其實這段對家庭描述的經文.
也沒有什麼神學大道理.
但作者對以撒家庭這麼寫實的描繪.
反而在提醒我們.
很多屬靈問題的出現不是突然之間.
而是在日常的關係 情感上的偏差慢慢累積出來.
甚至我們可以說.
當人的喜愛喜好凌駕在上主自己的心意之上的時候.
這些張力都會變得無可避免.
經過作者對以撒家庭一番描繪之後.
鏡頭一轉 來到整個敘事的高潮.
就是第29到30節.
有一日 雅各熬了湯 以素從田野回來.
疲憊不堪 以素對雅各說.
我累死了 請你讓我吃這紅的湯吧.
因此以素又叫以東.
一個看似平常 家庭生活的場景.
卻成為了決定命運的關鍵時刻.
當時雅各正在煲湯.
而剛好他哥哥又從田野回來.

$^{241}$他哥哥看起來很累 經文說他疲憊不堪.
在古代近東地區 這種精疲力竭的狀況.
通常都是源於乾渴或辛勞.
所以以素當時最需要的是喝水和休息.
因此他和弟弟要求他拿一碗湯去喝.
但是以素說的話又暗示.
其實他沒有他所說的那麼虛弱.
因為他說「請你讓我吃這紅的湯吧」.
以素竟然可以在他的說話中.
用「請你」,「求你」這些客套的說話.
但他又同時有些粗魯 重複地喊著「紅的紅的」.
來迫切地表達他當刻的渴求.
作者在這裡的敘事手法.
特別去突出「紅」這個字.
有可能是在玩一些相關語.
因為在希伯來文裡「紅」這個字.
一方面是直接連繫到.
或者是等同於以素另一個名移動.
而另一方面.
以素這一刻其實是在「紅紅」亂叫.
其實是在反映他的性格.
「紅的紅的」不是在描繪湯的品質.
濃郁的顏色那麼簡單.
而是表達以素那種強烈的感官渴望.
去體貼肉體的需要.
他那一刻就在喊著「紅色的紅色的」.
他很想喝這一碗湯.
相對於哥哥以素對著這個湯流口水自自浸.
弟弟雅各的回答就來得比較生硬.
三十一節雅各說.
「你今日把長子的名分賣給我吧」.
在雅各的說話裡.
他就完全省略了請求這一類客套說話.
反而用了「今日」「即刻」.
並且強調了「給我」.
這一類很直接近乎命令的字眼.
在雅各口中的長子名分.
所指的是在以色列裡長子能得到的特別尊崇.
作為爸爸第一個的兒子.
他們被視為分別為聖歸給上帝的那一位.

$^{281}$長子在爸爸還在生的時候就享有特權.
而當爸爸離世之後.
在分配遺產的時候.
也能夠享有特別的待遇.
根據《申命記》二十一章十七節的經文.
告訴我們長子能夠獲得雙倍的家業.
由經文去看.
雅各向他哥哥二嫂提出要求的方式.
反映他不是突然想起這個計謀.
而是他早有盤算.
他早就想好需要利用哥哥這個軟弱的時刻.
去奪取長子的名分.
面對著弟弟這麼重大的要求.
二嫂在此的回應顯得令人相當驚訝.
因為他竟然好像毫不在意地說.
「看吧,我快要死了,這長子名分對我有什麼用?」.
其實他這裡的說話都頗為輕佻.
反映他內心的狀態.
第一,似乎他真的還沒到快要死的時候.
否則他怎會有力氣這樣說話.
第二,似乎他對當時社會中很重視的長子身份和應許.
表現出一種漠不關心,毫不在意的態度.
雅各隨即對二嫂說「你今日對我起誓」.
雅各很簡潔,急促地用三個字回應.
「對我起誓今日」.
雅各的簡潔和急促再一次證實了他的冷靜和精明.
他那份要從哥哥的鬆懈中獲得利益的決心.
完完全全地表現出來.
結果在雅各的堅持下.
二嫂就向他起誓.
把長子的名份賣給了雅各.
於是雅各把餅和豆湯給了二嫂.
二嫂吃喝以後起來走了.
在第十三節下半部分到第三十四節上半部分.
經文呈現了一個交叉對稱的結構.
就是二嫂賣了一些東西給雅各.
而雅各給了一些東西二嫂.
告訴她這場交易是雙向的.
但同時也讓她知道這場交易是不平等的.
因為二嫂賣了一樣東西.

$^{321}$但雅各卻只給了一些東西.
值得留意雅各到底給了什麼東西二嫂.
似乎雅各給二嫂的並不是三十節二嫂所要求的.
紅的紅的那個濃湯.
而是他給了餅和一碗豆湯.
相信是一碗小的扁豆湯.
一種在當時很廉價平民的食物.
敘事者,創世記的作者似乎在這裡用一個精妙的寫作手法.
透過用字的改變去強烈呈現一個戲劇的落差.
為的是讓我們了解這件事究竟有多荒謬.
很難想像竟然會有人為了一碗普通的豆湯.
賣了他這麼珍貴的長子繼承權.
這裡是一個很大的諷刺.
就是二嫂原本以為他用他的長子名分換的是一份紅色的濃湯.
結果揭曉的時候.
其實那一碗只是很平庸,廉價的豆湯.
用我們今天的理解可能是一個人用了他的流契.
去換了一張二十元的快餐券.
任何人聽到都覺得匪夷所思.
於是二嫂吃了喝了便起來走了.
經文接著用了一連串很簡單四個連貫的動詞.
讓我們讀者去反思二嫂的行為.
用前面長篇大論很多東西說的二嫂相比.
二嫂這一刻顯得很沉默.
難道她真的不在意她的長子名分?.
還是當她拿到這個餅和那碗湯的時候.
她忽然之間理解到她失去了什麼.
導致她內心那種苦毒.
沒辦法跟她弟弟去交談.
我們不知道為什麼作者要呈現二嫂的沉默.
但作者很刻意去展示了事情發生前和後.
二嫂這個人物很大的落差.
最後經文.
創世記的作者給了一個很罕見但很直接的評價.
就是這樣二嫂輕看他長子的名分.
顯而易見,作者這句說話很重要.
因為我們讀聖經的時候.
很少出現這一類這麼明確的評論.
這句評論導出了整段經文的核心問題.
那個問題不是交易的本身.

$^{361}$而是二嫂對極具價值,神聖事物的態度.
雖然雅各在故事裡被描繪成一個冷酷的剝削者.
但作者沒有對雅各多加評論.
固然我們知道在後面故事的發展裡.
雅各要為他引起的問題敵對的狀態付上代價.
但這裡更重要的問題是.
二嫂對神聖的揀選表現出一種輕佻的態度.
二嫂一方面對自己這個寶貴身份藐視.
另一方面實際上她也在對耶和華上帝藐視.
因此二嫂最終必須和雅各交換位置.
她必須失去這個長子的身份.
弟兄姊妹.
當我們再一次細閱二嫂和雅各的故事.
我們會發現故事的重點不在於道基督教化.
也不是在爭論誰更詭詐,狡猾.
作者真正的用意是透過二嫂的遭遇.
向我們發出一個挑戰.
就是究竟我們是否懂得作出正確的抉擇.
今日的我們會否也同樣輕看了上帝那份寶貴的揀選.
事實上這場發生在以撒迦的古老衝突.
早已成為早期教會弟兄姊妹的嚴厲警告.
在新約希伯來書的作者曾經引用這個片段.
去斥責那些在迫害中動搖.
想去回歸猶太教的信徒.
那些想透過回歸猶太教而逃避迫害的猶太基督徒.
在經文中提醒.
「免得有人淫亂或不敬錢如二嫂,她因一點點食物,把自己的長子名分賣了.」.
在新約的作者眼中.
二嫂就是一個悲劇的標記.
她是一個體貼肉體,輕看永恆價值的典範.
或者今天我們在座的每一個人.
我們沒有人真的會為了一碗湯去賣掉長子的名分.
亦相信沒有人會荒唐到用一個樓契去換一張二十元的快餐券.
但是這個經文並不是和我們無關.
因為經文帶給我們一個更加震撼的提醒.
就是人面對屬靈上的危機.
往往不一定在於很重大的事情或危難.
而是可以發生在一個極其平常的日子.
一場看似微不足道的選擇中.
或許我們都難免試過將一些事情.

$^{401}$例如學業的成績,工作的成就,人際的關係.
或者娛樂等等看得太過重要.
就好像二嫂被當下的饑渴主導.
繼而被眼前所看到的紅的紅的湯蒙蔽了雙眼.
輕視了上帝揀選她的身份.
我們會否都這樣忽略了上主在我們生命中的一個宏大計劃.
結果在不經不覺之間我們都犯了二嫂的錯誤.
竟然用永恆的價值去換了平庸,廉價.
而且轉信即使的短暫滿足.
弟兄姊妹,今天的經文二嫂的故事.
就好像一面影照我們的鏡.
要我們去反思我們有沒有讓一些事情成為了那碗紅的紅的湯.
以致我們輕看了上帝給我們寶貴的名分.
盼望我們都懂得去謹慎,去警醒.
能夠作出正確的選擇.
做蒙上帝喜悅的決定.
免得我們好像二嫂那樣失落了那份給我們寶貴的名分.
願上主幫助我們每一個.
\newpage



\section{使徒行傳 1:8}
\label{sec:032027}
\textbf{與你同行,直到地極}~~~~李文康傳道
\newline
\newline
連結: \href{http://www.church.com.hk/acms/content.asp?site=cdc\&op=show\&type=product\&code=032027\&layout=sermon}{\texttt{032027}} ~~~~ 語音日期: 2026-04-19
\newline
\newline
\hyperref[sec:032026]{\small{< < < PREV SERMON < < <}}
~
\hyperref[sec:index_chronic]{\small{[返順時目]}}
~
\hyperref[sec:032028]{\small{> > > NEXT SERMON > > >}}
\newline
\newline
使徒行傳 1:8
\newline
\begin{longtable}{cl}
\hline
\hline
章節 & 經文 (和合本修訂版)\\
\hline
1:8 & \begin{tabularx}{0.7\textwidth}{X} 但聖靈降臨在你們身上，你們就必得著能力，並要在耶路撒冷、猶太全地和撒瑪利亞，直到地極，作我的見證。」 \end{tabularx} \\ \\
[1ex]
\hline
\hline
\end{longtable}
$^{1}$各位弟兄姊妹平安.
神透過黎加禮牧師.
與我們分享保羅,提摩太和菲律賓教會之間的結連.
讓我們看見在大使命的服事當中.
關係建立的重要.
並且鼓勵我們以基督的心為心.
效法基督的榜樣來服事.
今天我和大家延續猜拳分享的訊息.
就是我們今年關愛年會的主題.
「與你同行,直到地極」.
願意我們從使徒行傳去思想.
主不單止去猜派我們.
主也願意與我們同行.
還有就是我們怎樣去回應主的呼召.
「直到地極」作主的見證呢?.
第三就是教會如何與被猜的使者同行.
在開始之時讓我們一起同心去禱告.
借助你的愛子主耶穌基督曾經為我們受死.
為我們復活.
也為我們升天.
主耶穌成就了救恩.
但我們明白主耶穌仍然在工作.
也願意藉著聖靈與我們同行.
去幫助我們.
去成就主的旨意.
在這一刻,我們將無論是說的,聽的,傳言的.
恭敬,交託.
求聖靈親自引導和幫助.
讓我們不單止領受.
也從主得著力量去進行.
我們誠心禱告.
奉靠主耶穌基督的名字而求.
阿們.
在《使徒行傳》一章八節裡說到.
「但聖靈過臨在你們身上.
你們就必得著能力.
並要在耶路撒冷.
猶太傳遞和撒瑪利亞直到地極.
作我的見證」.
這段經文很清楚提示我們.

$^{41}$主耶穌基督是透過聖靈.
繼續願意與我們同行.
耶穌在升天前曾經吩咐門徒.
不要離開耶路撒冷.
要等候父所應許的聖靈.
這個等候非常重要.
主並不是差派門徒獨自出去.
靠著自己的能力計劃去成就大使命.
乃是應許聖靈會先降臨.
然後門徒就透過聖靈.
可以得著能力去為主作供.
隨後在《使徒行傳》第二章.
我們也看到五重節聖靈降臨.
標誌著教會的時代開始.
主耶穌雖然已經升天.
但主確實藉著聖靈.
繼續與教會同行和同工.
在猜拳事奉的過程當中.
確實很需要人用行動去回應.
但我們在此看到.
原來猜拳事奉的起點.
並不是靠人的意志或努力計劃.
普世猜拳工作的開始.
原來是透過聖靈的同在.
聖靈的帶領和聖靈能力的推動.
而我們要做的就是配合聖靈的工作.
如何做到呢?.
我們看看一章細節裡.
耶穌祝福門徒不要離開耶路撒冷.
要等候父所應許.
第一步我們需要有等候聖靈的心態.
讓聖靈主導我們的工作.
並不是由我們去主導神的工作.
正如昔日的門徒.
當耶穌復活之後.
門徒問耶穌.
「你復興以色列國是否就是這個時候呢?」.
耶穌的回應是.
「神的國度,祂的計劃,時間.
並不是人所能知道的」.

$^{81}$神並不是單單要復興以色列國.
神國度的擴展是普世性的.
神要賜福給萬族,萬民.
祂的心意是要願意繼續透過聖靈與我們同工.
第二步.
在等候聖靈指引的期間.
我們就好像一章十四節昔日的門徒一樣.
需要同心合意,行切地去禱告.
在禱告當中我們能夠調節自己.
調節自己去領會神的心意.
從而去配合.
最後當我們有了等候,有了禱告.
我們明白神的心意之後.
就要有信服聖靈的回應和行動.
用行動去具體回應.
我們知道猜拳並不是我們為神做偉大的事.
而是神邀請我們一起參與祂的服事.
在泰國有一些關愛服事的機會.
在當中也曾經學習如何等候聖靈的引導和幫助.
我們最先在教會附近一間幼稚園.
我們有一些接觸,建立了關係.
就可以去到這間幼稚園教授中文.
與老師,校長溝通.
接著向一群小朋友.
透過教授中文的過程中慢慢分享福音.
接著我們禱告.
「很想有本地的童工彼此配合」.
接著禱告了一段時間.
隔了大約兩年.
來到我們教會能夠彼此配搭.
雖然他們只能夠在星期六,星期日.
才能夠來到.
但是我們多了機會.
星期六可以一起到公園裡.
接觸人,派單張,傳福音.
很感恩.
接著也有多一位關愛事者加入了我們的團隊.
他本人對一些貧窮的村子的服務工作很負擔.
於是他就常常提起.
「我們會不會去發掘一些村子,傳福音.

$^{121}$去接觸一些村民呢?」.
接著我們發覺.
「我們當然慢慢開始人手多了.
但也有不同的事奉兼顧.
我們究竟行不行呢?.
有沒有這樣的空間呢?」.
於是我們就不斷地禱告.
求主去引證和帶領.
這是一次我們接待從香港到泰國的一隊關愛的短期體驗隊.
其中有一個環節.
就是出到一些社區.
鼓勵弟兄姊妹一邊去走這個社區的時候.
所看到的人與事.
都不斷地在內心裡面去為當中禱告.
在過程中我們帶隊伍經過一條河道.
這條河道有一個特色.
越走越黑.
越走越有味道.
雖然有很多味道.
但河道兩邊仍然住了很多村民.
接著我們去到一個分叉位.
通常我們的做法會轉左.
因為轉左之後.
我們沒有多遠.
就會去到一些預定的其他點.
準備帶領隊員去了解,去認識.
但在那一次.
心裡面突然有一個感動.
「不如向右轉吧」.
但向右轉的話.
我們自己沒有走過.
又帶著一群人.
會不會比較冒險呢?.
但最後還是按著這份感動.
向右轉.
走了沒多久.
過了一座橋.
經過了一塊荒地.
沒多久.
接著看到很多小朋友跑來跑去.

$^{161}$接著發覺原來我們慢慢進入了一條.
很長的木屋的村子.
當中的村民是很貧窄的.
都是很貧窮的.
我們跟弟兄姊妹一邊走一邊走.
發覺將來我們會不會.
就是可以在這條村子裡面.
去傳福音做一些訪問工作呢?.
於是就打聽一下居民.
有沒有一些教會已經在這裡.
做一些訪問工作呢?.
接著知道原來很久以前有.
現在就沒有了.
去遇到村長.
接著我們表明來意.
我們是教會.
可以免費提供一些服事.
幫助一些孩子去學習英文,中文.
問可不可以有這個機會呢?.
感恩.
村長竟然願意讓我們在村口.
一間比較大的房子.
來舉辦活動.
這間房子就是當大人.
每天早上出去工作.
這些小朋友就會帶到這間房子當中.
好像一個托兒所.
他們也請了一些老師.
做一些教導,托兒.
接著逢星期六下午.
這個地方就會空置了.
很多小朋友可能爸爸回來了.
媽媽回來了.
於是就可以回到自己家裡.
村長願意讓我們使用這間比較大的房子.
最先我們每個月一個星期.
接著兩個星期.
接著三個星期.
接著每個星期.
星期六下午.

$^{201}$都可以進到這個房子裡面.
去教導這些小朋友.
我們一步一步.
最初玩遊戲.
唱一些詩歌.
接著慢慢一步一步.
跟他們講聖經的故事.
真理的教導.
我們都很感恩.
但很寶貴的是.
這個服事的開始.
並不是單單只是人的熱心.
或者一些計劃.
而是我們體會到.
我們透過禱告,等候.
也體會到聖靈的指引,感動.
以致一步一步去開路.
更加寶貴的是.
我們這個經歷.
是跟從香港而來的一些隊員.
一起去發現這個地方.
一起去慢慢服事.
之後有隊員第二年會繼續來.
一起去幫助.
雖然我們跟太太.
回到香港一段長時間.
也知道這群孩子.
其實很多他們慢慢成長.
現在都已經變成青少年了.
甚至他們願意繼續在教會裡面.
受栽培,信主,成長.
這也是我們很感恩的.
弟兄姊妹.
讓我們一起去學習.
昔日的門徒.
他們當進行大使命的時候.
他們所要配合的.
不只是靠自己的能力,計劃.
而是受著聖靈而來的引導和幫助.
所以就算來到今天.

$^{241}$我們同樣在這個大使命.
同一個大使命底下.
在同一個聖靈的帶領底下.
其實我們同樣的原則.
同樣不是靠我們有什麼財綱,計劃,力量.
其實都是靠著聖靈的引導和幫助.
所以求主幫助我們.
當我們在做一些重大的決定.
或者一些服事之先.
讓我們先靜下來.
我們去禱告,去等候神的引導.
也在學習當中.
透過禱告,聖經的真理.
來辨識前面所做的.
是否神的心意.
以致我們能夠配合聖靈的工作.
既然主並不是單單叫我們獨自出去.
主耶穌願意藉著聖靈與我們同在.
賜給我們能力.
那麼這個能力究竟是想達到什麼目的呢?.
在經文裡很直接地說.
「直到地極作我的見證」.
在一章第八節後半段.
主耶穌勾畫出一幅福音的地圖.
並要在耶路撒冷,猶太傳遞和撒瑪利亞.
「直到地極作我的見證」.
這並不只是地理上遠近的問題.
更加是要跨越心理文化上的一些限制.
我們知道耶路撒冷.
是當時主耶穌基督受死復活的地方.
也是福音的起點.
是適合門徒安殊的地方.
接著是猶太地.
猶太傳遞就是以耶路撒冷為中心.
周邊的猶太地區.
是屬於同文化的群族.
因為他們能夠是同語言.
同種族的延伸.
接著撒瑪利亞就是猶太北方.
濃近的一些地區.

$^{281}$地極就是指近遠的地方.
是萬邦,是跨越不同的地域.
不同的語言,不同的文化.
所以算是一個跨文化的服事.
在整卷《仕徒人傳》的序詞當中.
我們會見到福音是從耶路撒冷.
一直擴展開去.
從門徒的身邊.
但不只是停留在他們身邊的同族.
神賦予給教會有普世的視野.
神的心意是要門徒.
將福音由本族傳到萬邦.
是同步的擴展.
當中值得我們留意的地方.
就是撒瑪利亞這個地方.
它並不只是單純較遠的地方.
我們知道撒瑪利亞在當時對猶太人來說.
他們會看為是一個外邦人的地方.
是鄙視他們的.
因為這個地方從古至今.
是充滿著歷史上的衝突.
文化的隔閡.
和宗教上的敵意的地方.
但耶穌特別提到撒瑪利亞.
其實我們心信不單只是提醒昔日的門徒.
甚至是提醒今日的我們.
福音並不只是傳給同族人.
並不是傳給我們喜歡的人.
福音是會傳給一些.
我們不容易接近的人.
甚至我們本來也不願意靠近的人.
所以主耶穌所給我們的大使命.
不單只是要我們突破我們的舒適區.
要突破民族的界限.
宗教的隔閡和限制.
將福音不斷地擴展開去.
上星期我們關愛年會.
我們成人有訊息.
由我們教會所支持的其中一間猜拳機構.
「新生資源中心」.

$^{321}$他們有同工來到我們教會.
在我們兒童部聚會中分享聚會.
他們設計了一套叫做「Fatima」的.
兒童動畫材料.
配合關愛年會的訊息.
讓兒童了解一些伊斯蘭宗教文化生活.
和基督信仰有什麼差異.
提醒兒童要謹記自己有基督徒的身份.
鼓勵他們能夠學習和穆斯林做朋友.
和傳福音給他們.
當中的目標是要讓孩子從小對自身的信仰.
能夠站立得穩.
甚至盼望我們教會的下一代.
他們能夠傳遞關心穆斯林的福音需要.
有很多機構福音的資源.
我們能夠彼此配搭,合作.
其實穆斯林的世界在不同的地區裡.
都是增長最快的民族.
所以當我們不快些守護著我們下一代的話.
其實我們下一代.
他們身邊的同學,朋友,同事.
其實都已經可能是其他非華語的群族.
他們來到香港.
不單止帶著他們不同種族的文化.
也帶著他們不同宗教的信仰.
如果我們不能夠守護我們自己.
甚至進一步讓孩子知道自己基督徒的身份的話.
很容易就會被同化.
所以求主幫助我們.
我們不要輕忽.
就算是小朋友也好.
關愛普世萬俗這個使命.
我們可以從小就放在他們的心中.
最近聽到一個很感恩的見證.
就是一位姊妹.
有一位兒新女讀五年級.
在學校裡發生了一件事.
兒新女的同學是其他種族的群族.
她的皮膚很黑.
然後發生了一件事.

$^{361}$非華語群體的同學很不開心.
她說自己的皮膚很黑.
擔心被人取笑.
事實是同班裡沒有太多人刻意地取笑她.
但這位姊妹的兒新女告訴她.
「原來神創造了我們世界上不同種族的人.
神也比我們獨特.
其實你也很漂亮」.
我們想不到一個讀五年級的小朋友.
能夠透過聖經所認識的.
神創造萬族這份心意.
並且他能夠肯定.
每一個人在神面前都有價值.
這是很寶貴的.
所以我們千萬不要輕忽.
神也會使用小孩子.
在他們的年紀裡.
同樣可以為主作見證.
並且說出一些當講的說話.
去安慰不同種族的人.
在這節經文當中.
主耶穌提到「作我的見證」.
原來「見證」這個字.
在原本的希臘文的文法當中.
並不是一個動詞.
而是一個名詞.
我們知道動詞所關注的.
就是做一些動作,做一些事.
名詞就是說你是誰.
是關注在身份上的.
所以我們深信當這一處所表達的時候.
就是要當見證人.
並且直譯為.
「將是我的見證人」有未來式的味道在當中.
所以在新漢語的翻譯.
是「直到地極作我的見證人」.
在這一處我們知道.
原來耶穌的本意.
並不只是單單鼓勵門徒.
我們懂得去講見證.

$^{401}$或者懂得去分享福音.
其實這也是重要的.
但是主耶穌很想我們門徒.
當我們去到人群當中的時候.
我們這個生命就是一個活生生的見證.
以致人能夠在我們身上.
知道耶穌基督就是那一位能夠拯救人.
甚至改變人生命的救主.
我們看回《少唐人傳》四章十三節裡.
都有記載彼得和約翰站在公會面前.
聖經記載說.
眾人見彼得和約翰的膽量.
又看出他們是學問的平民.
就稀奇認出他們是跟過耶穌的.
為什麼公會的人會認出.
彼得和約翰就是跟過耶穌的呢?.
原來並不是他們宣講了的說話.
原來就是透過彼得和約翰.
一群門徒他們的生命的改變.
還有他們有一份無法解釋的膽量.
還有神透過他們身上.
藉著他們顯出一些神蹟,奇事.
以致周邊的人會知道.
這群人是跟過耶穌的.
弟兄姊妹.
當我們有時舉辦一些關愛體驗隊.
就算短期也好.
我們鼓勵隊員.
我們可以有些不同的服侍.
首先我們要預備好自己的得救見證.
或者信仰的生活見證.
以致當我們去到工場的時候.
可能有機會關愛使者隨時會邀請我們.
可以去分享見證.
就算去到不同種族語言的地方也好.
可能有關愛使者幫我們翻譯.
這是很好的.
我們懂得去分享見證.
懂得去傳福音.
但在此經文提醒我們.

$^{441}$我們更加要很看重.
我這個人能夠用生命.
去見證耶穌基督的真實.
求主幫助我們.
讓我們能夠藉著我們的愛心.
藉著我們的忍耐.
藉著我們的一些選擇.
藉著我們一些為人的品格.
讓人在我們身上看到耶穌是真實的.
讓我們成為耶穌的證據.
正如我們今年的教會主題.
是「走出去傳遞愛 活出來見證神」.
我們深信當我們在神面前領受了很多恩典.
領受了很多真理的時候.
我們能夠將這份祝福真理活現出來.
讓我們身邊的人都能夠體會我們神的真實.
讓我們對不同文化種族的人.
不要排斥學習去接觸他們.
聆聽他們.
有時弟兄姊妹會問我.
「為什麼你和太太昔日會有感動.
去到泰國的地方服侍呢?」.
我都很直接地跟他們說.
「其實我一開始對泰國人沒有負擔的.
不是說看到泰國人就馬上流口水.
是沒有的.
但為什麼會去到泰國這個地方服侍呢?」.
以前我很喜歡去旅行.
去不同的地方.
但原來都沒有去選擇泰國.
因為會覺得泰國有很多偶像.
有很多邪靈的工作.
屬靈裡面好像一塌糊塗.
很骯髒 不喜歡.
但自從試過有機會透過短期關愛體驗.
去探訪教會.
探訪當地的基督徒.
探訪關愛使者.
真的有機會看到一個真實的場景.
正正因為這個地方有這麼多偶像.

$^{481}$這麼多邪靈工作.
其實更加需要福音臨到他們當中.
以致當中的泰國人能夠得到釋放.
得到主耶穌基督福音的自由.
所以求主去幫助我們.
願意我們都踏出我們的宣誓圈.
我們都嘗試多去了解不同種族群體的需要.
其實不單止遠去到泰國有泰國人.
我們在上個星期日下午.
我們教會裡面都很感恩.
有一群弟兄姊妹報名.
大概有四十多人.
一起去到九龍城這個地方.
去參與一個認識泰國人新年文化的活動.
當中去了解他們.
發現原來在我們教會.
坐巴士而已.
坐幾程車去到九龍城區.
已經有一群不同種族的群族在這裡生活.
他們很努力地去生活.
去投入融入在香港這個地方裡面.
我們出發之前.
有些簡介會都有教一點點泰文.
但其實都給大家有心理準備.
其實很多泰國餐廳或者店舖的店主.
已經操得一口流利的廣東話.
在當中我們嘗試去體會.
很多弟兄姊妹都發覺很特別.
原來都看到神開我們的視野.
原來在這麼近的地方.
我們都有一群群族.
他們很努力地在香港去生活,工作.
甚至他們的下一代.
都融入在我們的生活當中.
求主幫助我們.
我們深深知道.
並不是單單泰國人.
在香港這個地方裡面.
有很多非華裔的群族.
他們帶著自己的文化.

$^{521}$帶著他們的信仰來到香港當中.
盼望他們在香港生活.
不單止能夠有一個穩定.
能夠他們生活可以有工作.
可以他們的下一代有正常的教育.
穩定地生活下去.
更加盼望他們能夠在香港這個地方裡面.
可以有聽到福音,信主,成長的機會.
願意將來我們教會繼續舉辦.
不同的本地跨文化活動的時候.
都鼓勵弟兄姊妹一起來參與.
在《士徒人傳》一章八節.
剛才的經文的焦點.
就是聖靈設下能力.
使門徒能夠成為主的見證人.
這個「你們」很直接.
在當時來說是一群門徒.
他們從主耶穌基督領受這個大使命.
但是主耶穌基督的大使命.
從一開始都不只是單單一個人去承傳.
或者一小撮人去承就,去承擔.
整卷的《士徒人傳》.
我們看到後來的發展.
很快已經擴展成為整個教會.
共同承擔的使命.
第三點想和大家分享的是.
教會如何與被猜的使者同行呢?.
第一段經文在《士徒人傳》.
八章十四到十五節裡.
我們看到教會的關心和前往.
當中經文記載.
在耶路撒冷的使徒.
聽見撒瑪利亞人領受了神的道.
就打發彼得和約翰到他們那裡去.
兩個人下去就為他們禱告.
要讓他們領受聖靈.
我們知道最先是由肥利門徒.
使徒下去撒瑪利亞城.
去傳講耶穌基督.
耶路撒冷教會知道這件事.

$^{561}$知道撒瑪利亞城有人領受了神的道.
於是教會的領袖當中就被猜派.
由彼得和約翰前往.
這個行動代表著初期教會.
對新的工場,新的群體沒有袖手旁觀.
他們願意很實際地去參與.
去確認,去關心,去牧羊.
並不是只留於遠方的祝福.
他們願意走近.
願意一起去承擔.
我們看到當中教會的關心和前行.
第二段經文在九章二十六到二十七節.
我們看到教會的接納和承傳.
經文記載.
「蘇羅到了耶路撒冷.
想與門徒結交.
大家卻都怕他.
不信他是門徒.
只有巴拿巴接待他.
領他去見使徒.
把他在路上怎麼看見主.
主怎麼向他說話.
他在大馬士革.
怎麼奉耶穌的名放膽傳道.
都熟說出來」.
我們知道保羅.
之前的名字蘇羅.
初期是很迫不及待教會的.
所以早期的教會.
信徒對蘇羅迴轉信主.
都有很多顧慮和疑惑.
在過程當中.
原來只有巴拿巴第一個出來.
願意去接待保羅.
為他作見證.
幫助保羅融入在教會的群體當中.
所以我們看見同行.
是意味著能夠幫助別人.
去被接納,被建立,被成長.
發掘別人的恩賜.

$^{601}$引導他能夠進入在服事的行列當中.
回想我自己過去.
裝備自己關愛服事的時候.
原來真的不是單憑自己的意志.
或是自己個人的負擔.
就能夠成全.
在過程當中.
在不同的階段.
都回想看到神.
有差派不同的道士.
童工,牧者.
與我同行.
有些願意與我進入退休營.
進入一些營會.
一起尋索神的呼召.
有時未必是別人對自己不認同.
有時是自己對自己不認同.
思考會疑惑.
我可以做關愛事者?.
可不可以呢?.
是不是呢?.
是否要有前線呢?.
其實後方都可以吧.
在過程當中.
是需要被人同行.
給予意見,肯定和幫助.
弟兄姊妹.
你願不願意呢?.
都預備自己.
可能在整個大使命.
整全的服事當中.
神都會願意使用你.
成為其他弟兄姊妹的巴拿巴.
在服事當中.
願意我們一同尋求神的心意.
尋求神賦予給我們獨特的角色.
我們就在這個角色中心服事.
第三段經文.
我們看到教會的差派和代徒.
在《使徒行傳》十三章一至三節.

$^{641}$當中記載.
「在安提安的教會中.
有幾位先知和教師.
就視巴拿巴和稱為利潔的西面.
古利內人路求.
與希魯分封王一起長大的馬廉和蘇羅.
他們在事奉主和禁食的時候.
聖靈說.
要為我分派巴拿巴和蘇羅.
去做我召他們做的工作.
於是他們禁食禱告後.
給巴拿巴和蘇羅按手.
然後派遣他們走了」.
安提安的教會禁食,禱告,按手.
差派巴拿巴和蘇羅出去.
我們看到教會的差派.
並不是只是因為某一小撮人.
他們按照個人的志向.
或者個人的熱心,熱血衝出去.
而是整個教會.
一起在聖靈裡面禱告,印證,回應.
教會將工人差出去的時候.
同時擺在禱告當中.
擺在關懷裡面.
往後繼續有祝福的行動.
我們現今有支持關愛使者.
有十七個單位.
他們分別在不同的較為開放一點的地方.
或者在一些創啟的地方.
當我們公開傳福音的時候是不可以的.
有機會是被逼迫和被禁止的地方.
這些關愛使者的服侍.
都不斷激勵著我們.
但不知道大家對他們的認識有多少.
我盼望當每一次有關愛使者.
他們回來香港述職的時候.
我們都很願意,很踴躍地.
去聽他們的分享會.
我們聽完之後接觸他們.
不但送上一句祝福.

$^{681}$「願神大大使用你」.
我們更加跟他們說.
「我願意成為你的同行者」.
我們學習以禱告,奉獻,探訪,關懷.
實際的行動來與關愛使者同行.
我記得當我和太太在泰國服侍的時候.
其實我們最掛心,最擔心的是什麼呢?.
除了在工場上的日常事奉.
弟兄姊妹的需要.
其實很多時候都會擔心在香港的家人.
甚至特別當他們有病的時候.
會特別擔心.
我記得當爸爸肺癌的時候.
當然有機會回來香港.
我自己探望.
但總會有機會需要回到工場.
但發覺有一些支持教會的弟兄姊妹.
其實我也跟他們不太熟悉.
但他們來過工場做過短宣隊.
我們認識過.
接著他們看到我們《代讀信》的需要.
就問「可不可以探望你的家人?」.
我們說「當然可以」.
之後看到他們買了一些健康的食物.
去探訪我患病的爸爸.
透過他們的相片.
透過他們的服侍.
我深深地很受安慰.
因為我很想在親人身邊服侍.
但做不到.
但是一群弟兄姊妹.
他們願意代替我服侍當中的家人.
當中的感動是很深刻的.
所以弟兄姊妹.
我們對關愛使者的服侍.
不單可以走到前線.
去陪伴他們.
去聽他們的分享.
就算在後方.
我們同樣可以有很多不同的服侍.

$^{721}$按照關愛使者的需要來擺上.
求主幫助和感動我們.
所以讓我們知道.
在猜拳的事奉當中.
有人出去.
有人在後方作支援.
有人在前方宣講.
有人在後方不斷地代讀來守望.
大家在整個大使命中都是團隊.
沒有分別.
都是成為神所使用的.
所以弟兄姊妹.
與你同行直到地極.
並不是單單一句口號.
從聖經中我們看到.
是神與人的一個文約.
神邀請我們.
藉著聖靈加添給我們力量,指引.
以致我們能夠按照神的呼召.
從本土直到地極.
為主作見證.
藉著生命來活出基督的榜樣.
神也呼召教會猜拳,同行,守望.
願意我們教會繼續被神大大的使用.
求主使用我們每一位.
雖然我們未必每一位.
被神呼召成為走上前線.
作為一個關愛使者.
但是作為一個關愛使者的同行者.
我們每一個人都可以做到.
求主幫助我們.
讓我們一起同心去禱告.
親愛的主耶穌.
你願意藉著聖靈賜給我們能力.
賜給我們引導和幫助.
以至我們去進行大使命的時候.
我們知道我們與耶穌基督同行.
讓我們的生命得著改變.
讓我們領受從主而來的祝福.
但主你更願意我們透過生命的流露.

$^{761}$讓我們身邊的人認識.
耶穌基督的拯救是真實的.
你的祝福臨到今天.
直到永遠.
所以主求你去眺望.
我們眾弟子們的心智.
讓我們在教會當中.
我們一起成為主你所使用的.
不單只是一間差遣的教會.
我們甚至被每一位被差遣的使者同行.
我們願意在主你面前.
我們願意去回應.
我們願意與主同行.
甚至與主你所差遣的使者同行.
直到地極.
求主你親自的幫助和帶領.
誠心禱告奉靠主耶穌基督的名字而求.
阿們.
\newpage



\section{歌羅西書 1:1-4:18}
\label{sec:032028}
\textbf{歌羅西書全書概論:超越的基督}~~~~劉少康牧師
\newline
\newline
連結: \href{http://www.church.com.hk/acms/content.asp?site=cdc\&op=show\&type=product\&code=032028\&layout=sermon}{\texttt{032028}} ~~~~ 語音日期: 2026-04-15
\newline
\newline
\hyperref[sec:032027]{\small{< < < PREV SERMON < < <}}
~
\hyperref[sec:index_chronic]{\small{[返順時目]}}
~
\hyperref[sec:032029]{\small{> > > NEXT SERMON > > >}}
\newline
\newline
歌羅西書 1:1-4:18
\newline
\begin{longtable}{cl}
\hline
\hline
章節 & 經文 (和合本修訂版)\\
\hline
1:1 & \begin{tabularx}{0.7\textwidth}{X} 奉神旨意，作基督耶穌使徒的保羅，和我們的弟兄提摩太， \end{tabularx} \\ \\ \relax
1:2 & \begin{tabularx}{0.7\textwidth}{X} 寫信給歌羅西的聖徒，在基督裡忠心的弟兄。願恩惠、平安從我們的父神歸給你們！ \end{tabularx} \\ \\ \relax
1:3 & \begin{tabularx}{0.7\textwidth}{X} 我們為你們禱告的時候，常常感謝我們主耶穌基督的父神， \end{tabularx} \\ \\ \relax
1:4 & \begin{tabularx}{0.7\textwidth}{X} 因為聽見你們對基督耶穌的信心，並對眾聖徒有的愛心。 \end{tabularx} \\ \\ \relax
1:5 & \begin{tabularx}{0.7\textwidth}{X} 這都是因著那給你們存在天上的盼望，它就是你們從前所聽見真理的道，就是福音； \end{tabularx} \\ \\ \relax
1:6 & \begin{tabularx}{0.7\textwidth}{X} 這福音傳到你們那裡，也傳到普天下，並且繼續增長，不斷結果，正如自從你們聽見福音，真正知道神恩惠的日子起，在你們中間也是這樣。 \end{tabularx} \\ \\ \relax
1:7 & \begin{tabularx}{0.7\textwidth}{X} 這福音是你們從我們所親愛、一同作僕人的以巴弗學到的。他為我們作了基督的忠心僕役， \end{tabularx} \\ \\ \relax
1:8 & \begin{tabularx}{0.7\textwidth}{X} 也把聖靈賜給你們的愛告訴我們。 \end{tabularx} \\ \\ \relax
1:9 & \begin{tabularx}{0.7\textwidth}{X} 因此，我們自從聽見的日子就不住地為你們禱告和祈求，願你們滿有一切屬靈的智慧和悟性，真正知道神的旨意， \end{tabularx} \\ \\ \relax
1:10 & \begin{tabularx}{0.7\textwidth}{X} 好使你們行事為人對得起主，凡事蒙他喜悅，在一切善事上結果子，對神的認識更有長進。 \end{tabularx} \\ \\ \relax
1:11 & \begin{tabularx}{0.7\textwidth}{X} 願你們從他榮耀的權能中，得以在一切事上力上加力，好使你們凡事歡歡喜喜地忍耐寬容， \end{tabularx} \\ \\ \relax
1:12 & \begin{tabularx}{0.7\textwidth}{X} 又感謝父，使你們配與眾聖徒在光明中分享基業。 \end{tabularx} \\ \\ \relax
1:13 & \begin{tabularx}{0.7\textwidth}{X} 他救了我們脫離黑暗的權勢，遷移到他愛子的國度裡。 \end{tabularx} \\ \\ \relax
1:14 & \begin{tabularx}{0.7\textwidth}{X} 藉著他的愛子，我們得蒙救贖，罪得赦免。 \end{tabularx} \\ \\ \relax
1:15 & \begin{tabularx}{0.7\textwidth}{X} 愛子是那看不見的神之像，是首生的，在一切被造的以先。 \end{tabularx} \\ \\ \relax
1:16 & \begin{tabularx}{0.7\textwidth}{X} 因為萬有都是在他裡面造的，無論是天上的、地上的，能看見的、不能看見的，或是有權位的、統治的，或是執政的、掌權的，一概都是藉著他為著他造的。 \end{tabularx} \\ \\ \relax
1:17 & \begin{tabularx}{0.7\textwidth}{X} 他在萬有之先；萬有也靠他而存在。 \end{tabularx} \\ \\ \relax
1:18 & \begin{tabularx}{0.7\textwidth}{X} 他是身體（教會）的頭；他是元始，是從死人中復活的首生者，好讓他在萬有中居首位。 \end{tabularx} \\ \\ \relax
1:19 & \begin{tabularx}{0.7\textwidth}{X} 因為神喜歡使一切的豐盛在他裡面居住， \end{tabularx} \\ \\ \relax
1:20 & \begin{tabularx}{0.7\textwidth}{X} 藉著他，神使萬有與自己和好，無論是地上的、天上的，都藉著他在十字架上所流的血促成了和平。 \end{tabularx} \\ \\ \relax
1:21 & \begin{tabularx}{0.7\textwidth}{X} 從前你們與神隔絕，心思上與他為敵，行為邪惡； \end{tabularx} \\ \\ \relax
1:22 & \begin{tabularx}{0.7\textwidth}{X} 但如今，他藉著他兒子肉身的死，已經使你們與他自己和好了，把你們獻在他的面前，成為聖潔，沒有瑕疵，無可指責。 \end{tabularx} \\ \\ \relax
1:23 & \begin{tabularx}{0.7\textwidth}{X} 只要你們持守信仰，根基穩固，堅定不移，不致動搖，離開了你們從前所聽見的福音的盼望；這福音也是傳給天下一切被造之物的，我—保羅作了這福音的僕役。 \end{tabularx} \\ \\ \relax
1:24 & \begin{tabularx}{0.7\textwidth}{X} 現在我為你們受苦，倒很快樂；並且為基督的身體，就是為教會，我要在自己的肉身上補滿基督未盡的苦難。 \end{tabularx} \\ \\ \relax
1:25 & \begin{tabularx}{0.7\textwidth}{X} 我照神為你們所賜我的職分作了教會的僕役，要把神的道傳得完滿； \end{tabularx} \\ \\ \relax
1:26 & \begin{tabularx}{0.7\textwidth}{X} 這道就是歷世歷代所隱藏的奧祕，但如今向他的聖徒顯明了。 \end{tabularx} \\ \\ \relax
1:27 & \begin{tabularx}{0.7\textwidth}{X} 神要讓他們知道，這奧祕在外邦人中有何等豐盛的榮耀；就是基督在你們心裡成了得榮耀的盼望。 \end{tabularx} \\ \\ \relax
1:28 & \begin{tabularx}{0.7\textwidth}{X} 我們傳揚他，是用諸般的智慧，勸戒各人，教導各人，要把各人在基督裡完完全全地獻上。 \end{tabularx} \\ \\ \relax
1:29 & \begin{tabularx}{0.7\textwidth}{X} 我也為此勞苦，照著他在我裡面運用的大能盡心竭力。 \end{tabularx} \\ \\ \relax
2:1 & \begin{tabularx}{0.7\textwidth}{X} 我要你們知道，我為你們和老底嘉人，和所有沒有與我見過面的人，是何等地勤奮； \end{tabularx} \\ \\ \relax
2:2 & \begin{tabularx}{0.7\textwidth}{X} 為要使他們的心得安慰，因愛心互相聯絡，以致有從確實了解所產生的豐盛，好深知神的奧祕，就是基督； \end{tabularx} \\ \\ \relax
2:3 & \begin{tabularx}{0.7\textwidth}{X} 在他裡面蘊藏著一切智慧和知識。 \end{tabularx} \\ \\ \relax
2:4 & \begin{tabularx}{0.7\textwidth}{X} 我說這話，免得有人用花言巧語迷惑你們。 \end{tabularx} \\ \\ \relax
2:5 & \begin{tabularx}{0.7\textwidth}{X} 雖然我身體不在你們那裡，心卻與你們同在，很高興見你們循規蹈矩，對基督的信心也堅固。 \end{tabularx} \\ \\ \relax
2:6 & \begin{tabularx}{0.7\textwidth}{X} 既然你們接受了主基督耶穌，就要靠著他而生活， \end{tabularx} \\ \\ \relax
2:7 & \begin{tabularx}{0.7\textwidth}{X} 照著你們所領受的教導，在他裡面生根建造，信心堅固，充滿著感謝的心。 \end{tabularx} \\ \\ \relax
2:8 & \begin{tabularx}{0.7\textwidth}{X} 你們要謹慎，免得有人用他的哲學和虛空的廢話，不照著基督，而是照人間的傳統和世上粗淺的學說，把你們擄去。 \end{tabularx} \\ \\ \relax
2:9 & \begin{tabularx}{0.7\textwidth}{X} 因為神本性一切的豐盛都有形有體地居住在基督裡面； \end{tabularx} \\ \\ \relax
2:10 & \begin{tabularx}{0.7\textwidth}{X} 你們在他裡面也已經成為豐盛。他是所有執政掌權者的元首。 \end{tabularx} \\ \\ \relax
2:11 & \begin{tabularx}{0.7\textwidth}{X} 你們也在他裡面受了不是人手所行的割禮，而是使你們脫去肉體情慾的基督的割禮。 \end{tabularx} \\ \\ \relax
2:12 & \begin{tabularx}{0.7\textwidth}{X} 你們既受洗與他一同埋葬，也就在此禮上，因信那使他從死人中復活的神的作為跟他一同復活。 \end{tabularx} \\ \\ \relax
2:13 & \begin{tabularx}{0.7\textwidth}{X} 你們從前在過犯和未受割禮的肉體中死了，神卻赦免了你們一切的過犯，使你們與基督一同活過來， \end{tabularx} \\ \\ \relax
2:14 & \begin{tabularx}{0.7\textwidth}{X} 塗去了在律例上所寫、敵對我們、束縛我們的字據，把它撤去，釘在十字架上。 \end{tabularx} \\ \\ \relax
2:15 & \begin{tabularx}{0.7\textwidth}{X} 基督既將一切執政者、掌權者的權勢解除了，就在凱旋的行列中，將他們公開示眾，仗著十字架誇勝。 \end{tabularx} \\ \\ \relax
2:16 & \begin{tabularx}{0.7\textwidth}{X} 所以，不要讓任何人在飲食上，或節期、初一、安息日等事上評斷你們。 \end{tabularx} \\ \\ \relax
2:17 & \begin{tabularx}{0.7\textwidth}{X} 這些原是未來的事的影子，真體卻是屬基督的。 \end{tabularx} \\ \\ \relax
2:18 & \begin{tabularx}{0.7\textwidth}{X} 不要讓人藉著故作謙虛和敬拜天使奪去你們的獎賞。這等人拘泥在所見過的幻象，隨著自己的慾望無故地自高自大， \end{tabularx} \\ \\ \relax
2:19 & \begin{tabularx}{0.7\textwidth}{X} 不緊隨元首；其實，由於他全身藉著關節筋絡才得到滋養，互相聯絡，靠神所賜的成長而成長。 \end{tabularx} \\ \\ \relax
2:20 & \begin{tabularx}{0.7\textwidth}{X} 既然你們與基督同死而脫離了世上粗淺的學說，為甚麼仍像生活在世俗中一樣，去服從那「不可拿、不可嘗、不可摸」等類的規條呢？ \end{tabularx} \\ \\ \relax
2:21 & \begin{tabularx}{0.7\textwidth}{X} 見上節 \end{tabularx} \\ \\ \relax
2:22 & \begin{tabularx}{0.7\textwidth}{X} 這些都是根據人的命令和教導，論到這一切都是一經使用就都敗壞了。 \end{tabularx} \\ \\ \relax
2:23 & \begin{tabularx}{0.7\textwidth}{X} 這些規條使人徒有智慧之名，用私意崇拜，自表謙卑，苦待己身，其實在克制肉體的情慾上毫無功效。 \end{tabularx} \\ \\ \relax
3:1 & \begin{tabularx}{0.7\textwidth}{X} 所以，既然你們已經與基督一同復活，就當求上面的事；那裡有基督，坐在神的右邊。 \end{tabularx} \\ \\ \relax
3:2 & \begin{tabularx}{0.7\textwidth}{X} 你們要思考上面的事，不要思考地上的事。 \end{tabularx} \\ \\ \relax
3:3 & \begin{tabularx}{0.7\textwidth}{X} 因為你們已經死了，你們的生命與基督一同藏在神裡面。 \end{tabularx} \\ \\ \relax
3:4 & \begin{tabularx}{0.7\textwidth}{X} 基督是你們的生命，他顯現的時候，你們也要與他一同在榮耀裡顯現。 \end{tabularx} \\ \\ \relax
3:5 & \begin{tabularx}{0.7\textwidth}{X} 所以，要治死你們在地上的肢體；就如淫亂、污穢、邪情、惡慾和貪婪—貪婪就是拜偶像。 \end{tabularx} \\ \\ \relax
3:6 & \begin{tabularx}{0.7\textwidth}{X} 因這些事，神的憤怒必臨到那些悖逆的人。 \end{tabularx} \\ \\ \relax
3:7 & \begin{tabularx}{0.7\textwidth}{X} 當你們在這些事中活著的時候，你們的行為也曾是這樣的。 \end{tabularx} \\ \\ \relax
3:8 & \begin{tabularx}{0.7\textwidth}{X} 但現在你們要棄絕這一切的事，就是惱恨、憤怒、惡毒、毀謗和口中污穢的言語。 \end{tabularx} \\ \\ \relax
3:9 & \begin{tabularx}{0.7\textwidth}{X} 不要彼此說謊，因為你們已經脫去舊人和舊人的行為， \end{tabularx} \\ \\ \relax
3:10 & \begin{tabularx}{0.7\textwidth}{X} 穿上了新人，這新人照著造他的主的形像在知識上不斷地更新。 \end{tabularx} \\ \\ \relax
3:11 & \begin{tabularx}{0.7\textwidth}{X} 在這事上並不分希臘人和猶太人，受割禮的和未受割禮的，未開化的人、西古提人、為奴的、自主的；惟獨基督是一切，又在一切之內。 \end{tabularx} \\ \\ \relax
3:12 & \begin{tabularx}{0.7\textwidth}{X} 所以，你們既是神的選民，聖潔、蒙愛的人，要穿上憐憫、恩慈、謙虛、溫柔和忍耐。 \end{tabularx} \\ \\ \relax
3:13 & \begin{tabularx}{0.7\textwidth}{X} 倘若這人與那人有嫌隙，總要彼此容忍，彼此饒恕；主怎樣饒恕了你們，你們也要怎樣饒恕人。 \end{tabularx} \\ \\ \relax
3:14 & \begin{tabularx}{0.7\textwidth}{X} 除此以外，還要穿上愛心，因為愛是貫通全德的。 \end{tabularx} \\ \\ \relax
3:15 & \begin{tabularx}{0.7\textwidth}{X} 你們要讓基督所賜的和平在你們心裡作主，也為此蒙召，歸為一體。你們還要存感謝的心。 \end{tabularx} \\ \\ \relax
3:16 & \begin{tabularx}{0.7\textwidth}{X} 當用各樣的智慧，把基督的道豐豐富富的存在心裡，用詩篇、讚美詩、靈歌，彼此教導，互相勸戒，以感恩的心歌頌神。 \end{tabularx} \\ \\ \relax
3:17 & \begin{tabularx}{0.7\textwidth}{X} 你們無論做甚麼，或說話或行事，都要奉主耶穌的名，藉著他感謝父神。 \end{tabularx} \\ \\ \relax
3:18 & \begin{tabularx}{0.7\textwidth}{X} 你們作妻子的，要順服自己的丈夫，這在主裡面是合宜的。 \end{tabularx} \\ \\ \relax
3:19 & \begin{tabularx}{0.7\textwidth}{X} 你們作丈夫的，要愛你們的妻子，不可虐待她們。 \end{tabularx} \\ \\ \relax
3:20 & \begin{tabularx}{0.7\textwidth}{X} 你們作兒女的，要凡事聽從父母，因為這是主所喜悅的。 \end{tabularx} \\ \\ \relax
3:21 & \begin{tabularx}{0.7\textwidth}{X} 你們作父親的，不要惹兒女生氣，恐怕他們會灰心。 \end{tabularx} \\ \\ \relax
3:22 & \begin{tabularx}{0.7\textwidth}{X} 你們作僕人的，要凡事聽從你們肉身的主人，不要只在眼前服事，像是討人喜歡的，總要心存誠實，因為你們敬畏主。 \end{tabularx} \\ \\ \relax
3:23 & \begin{tabularx}{0.7\textwidth}{X} 你們無論做甚麼，都要從心裡做，像是為主做的，不是為人做的； \end{tabularx} \\ \\ \relax
3:24 & \begin{tabularx}{0.7\textwidth}{X} 因為你們知道，從主那裡必得著基業作為賞賜。你們要服侍的是主基督。 \end{tabularx} \\ \\ \relax
3:25 & \begin{tabularx}{0.7\textwidth}{X} 行不義的人必受不義的報應；主並不偏待人。 \end{tabularx} \\ \\ \relax
4:1 & \begin{tabularx}{0.7\textwidth}{X} 你們作主人的，待僕人要公正，因為知道，你們也有一位主在天上。 \end{tabularx} \\ \\ \relax
4:2 & \begin{tabularx}{0.7\textwidth}{X} 你們要恆切禱告，在禱告中警醒感恩。 \end{tabularx} \\ \\ \relax
4:3 & \begin{tabularx}{0.7\textwidth}{X} 同時，也要為我們禱告，求神給我們開傳道的門，能宣講基督的奧祕， \end{tabularx} \\ \\ \relax
4:4 & \begin{tabularx}{0.7\textwidth}{X} 使我能按著所該說的話將這奧祕顯明出來，我為此而被捆鎖。 \end{tabularx} \\ \\ \relax
4:5 & \begin{tabularx}{0.7\textwidth}{X} 你們要把握時機，用智慧與外人來往。 \end{tabularx} \\ \\ \relax
4:6 & \begin{tabularx}{0.7\textwidth}{X} 你們的言談要時常帶著溫和，好像用鹽調味，讓你們知道該怎樣應對每一個人。 \end{tabularx} \\ \\ \relax
4:7 & \begin{tabularx}{0.7\textwidth}{X} 推基古是我親愛的弟兄，忠心的僕役，和我一同作主的僕人；他要把我一切的事都告訴你們。 \end{tabularx} \\ \\ \relax
4:8 & \begin{tabularx}{0.7\textwidth}{X} 我特意打發他到你們那裡去，好讓你們知道我們的情況，又讓他安慰你們的心。 \end{tabularx} \\ \\ \relax
4:9 & \begin{tabularx}{0.7\textwidth}{X} 我又打發一位親愛忠心的弟兄阿尼西謀同去；他也是你們那裡的人。他們會把這裡一切的事都告訴你們。 \end{tabularx} \\ \\ \relax
4:10 & \begin{tabularx}{0.7\textwidth}{X} 與我一同坐牢的亞里達古問候你們。巴拿巴的表弟馬可也問候你們。關於他，你們已經得到指示；他若到你們那裡，你們要接待他。 \end{tabularx} \\ \\ \relax
4:11 & \begin{tabularx}{0.7\textwidth}{X} 稱為猶士都的耶數也問候你們。奉割禮的人中，只有這三個人是為神的國與我作同工的，也是使我心裡得安慰的。 \end{tabularx} \\ \\ \relax
4:12 & \begin{tabularx}{0.7\textwidth}{X} 有一位你們那裡的人，作基督耶穌僕人的以巴弗問候你們。他禱告的時候常為你們竭力祈求，願你們能站穩而成熟，充分確信神一切的旨意。 \end{tabularx} \\ \\ \relax
4:13 & \begin{tabularx}{0.7\textwidth}{X} 他為你們、老底嘉和希拉坡里的弟兄多多勞苦，這是我可以為他作見證的。 \end{tabularx} \\ \\ \relax
4:14 & \begin{tabularx}{0.7\textwidth}{X} 親愛的醫生路加和底馬問候你們。 \end{tabularx} \\ \\ \relax
4:15 & \begin{tabularx}{0.7\textwidth}{X} 請問候老底嘉的弟兄以及寧法，和她家裡的教會。 \end{tabularx} \\ \\ \relax
4:16 & \begin{tabularx}{0.7\textwidth}{X} 你們宣讀了這書信，也要交給老底嘉的教會宣讀；你們也要宣讀從老底嘉轉來的書信。 \end{tabularx} \\ \\ \relax
4:17 & \begin{tabularx}{0.7\textwidth}{X} 你們要對亞基布說：「務要完成你從主所領受的職分。」 \end{tabularx} \\ \\ \relax
4:18 & \begin{tabularx}{0.7\textwidth}{X} 我—保羅親筆問候你們。要記念我在捆鎖中。願恩惠與你們同在！ \end{tabularx} \\ \\
[1ex]
\hline
\hline
\end{longtable}
$^{1}$各位電影小妹,平安!.
我要用一次的時間介紹保羅書信裡面.
一卷很重要的書,就是《歌羅西書》.
《歌羅西書》最主要就是保羅教.
哥羅西教會去辨識真假的信仰.
因為當時有些假先知起來.
講一些似是而非的道理.
還有《歌羅西書》很重要講到.
基督是教會的元首.
這本書是講教會論.
很重要的書.
講到教會是什麼?.
在神的計劃裡面.
怎樣安排教會.
還有教會我們應該怎樣看它.
我們信徒在教會裡面.
應該有的位分是什麼?.
保羅很特別.
他原來沒有去過哥羅西這個城市.
哥羅西這個城市.
是在爾弗所這個城市的.
東邊大概一百里的地方.
只是一個很小很小的城市.
這個城市有什麼特別呢?.
他們出產一種很特別的黑色的羊毛.
羊毛很多時候我們都認為是白色的.
但是這裡出產一種黑的羊毛.
所以非常之出名.
這間教會是由爾巴弗.
保羅很親密的同工所建立的.
保羅沒有去過探訪的.
因此大概在後來的時候.
保羅托推基古和歸順主的一個奴隸.
很出名的阿尼西姆.
就是《肥利門書》的主角.
阿尼西姆去探訪這間教會.
阿尼西姆在和合本修訂版.
就譯為「阿尼西謀」.
因為阿尼西姆好像是誰的媽媽.
所以和修版就改了一個名字叫做「阿尼西謀」.

$^{41}$所以其實和修版有一個好處.
就將好像「尼哥底姆」這些「尼哥底謀」.
翻譯成像一個人名多一些.
不會誤會了是誰誰的媽媽那一類.
所以這個都是有幫助的.
結果保羅就在這本書裡面一開頭.
就說你們要分析,分辨.
甚麼是真信仰,甚麼是假信仰.
《歌羅西書》是有四章的.
一個主題就是.
基督是教會的元首.
基督要居於首位.
我們要高舉基督.
在第一章保羅就提到.
愛子的身份和工作.
即是基督是神的愛子.
神差他的愛子來到這個世界.
他要愛子為他來工作.
起一切的屬靈悟聖上.
保羅有一個禱告.
願你們起一切屬靈悟聖上.
滿心知道神的旨意.
你們要明白神的旨意是甚麼.
滿心這個字即是說.
你要盡你的理解去明白神的旨意.
所以各位弟兄姊妹.
我們實在很需要時事尋求神的旨意.
而且我們要按神的旨意去行.
不是靠我們自己想出來的.
不是自己去猜的.
而是等候,聆聽,仰望神.
在我的人生當中有甚麼差遣?.
神對我的生命有甚麼安排?.
行事為人.
我們必須要對得起主.
今日你的行事.
你的為人.
是否對得起主?.
彰顯主自己.
然後第二章就講到保羅的身份.

$^{81}$和他的工作.
保羅就特別提到.
要防備那些不合真理的學問.
有些人是傳一些假的學問.
我們應該以基督為主.
而且我們要不要在世俗裡活著.
要以基督為我們生命的中心.
到了第三章.
他就提到我們時時都要讓基督的平安.
在我們的心裡作主.
這個就好像一場拳擊比賽.
大家都知道兩個拳手都差不多.
所以那個裁判就很重要.
你們兩個打在一起的時候.
將他們分開.
然後誰打敗對方.
還繼續打下去就制止他.
不准他再打.
所以讓基督的平安.
在我們心裡做裁判,作主.
就是以一個判決.
基督的平安是我們的判決.
你時時都要有基督的平安.
在你心裡面.
在一個很高的位置.
平安在我心.
基督的平安.
在我心裡面作主,作王.
所以我們基督徒的生活.
應該是以基督為重心.
因此我們時時都要忠於基督.
我們要行切祈禱.
所以到了第四章.
就講到要行切祈禱.
在此警醒感恩.
保羅知道主耶穌回來的日子很近.
我們必須要警醒,禱告,感恩.
然後我們要用智慧跟外人交往.
其實我們跟弟兄姊妹的交往.
要憑愛心.

$^{121}$要用和平彼此聯絡.
竭力保守聖靈所賜合而為一的心.
但我們跟外人交往.
必須要有智慧.
懂得怎樣去回答各人.
然後第四章又講到.
要拓展神的國度.
努力的傳呼音.
讓神的國不斷得到開展.
所以推基古和亞尼西牟.
就是這兩個人帶這封信去到哥羅西.
也有人說可能推基古和亞尼西牟.
都是帶信去到菲臘比教會的一個帶信人.
所以這兩本書排起保羅書信.
以弗所.
加拿大以弗所.
然後就到菲臘比哥羅西.
好像連著一樣.
所以這兩封信有些近似的地方.
所以歌羅西書實在給我們很多很寶貴的提醒.
這個基督論就講到基督的至高無上.
特別是第一第二章.
講到基督的威憤.
基督的至高無上.
然後第三第四章.
就是我們要信服基督.
我們要在基督裡面.
我們尊崇基督.
信服基督.
跟隨基督的帶領.
各位弟兄姊妹.
耶穌基督就是我們生命的目標.
我們要向著目標前進.
我們要高舉基督.
讓基督在我們生命裡面.
作主,作王.
我們時時要讓主居首位.
我自己印象很深刻.
當我小時候在教會成長的時候.
時時唱一首聖詩.

$^{161}$「讓主居首位」.
我們要讓主居首位.
願主居首位.
我生命中.
讓主居首位.
不單止願主居首位.
要讓主居首位.
是一步一步的.
最先你可能學怎樣尊主為大.
然後慢慢你願基督居首位.
最後你讓.
讓出你的位.
讓基督居首位.
求主都同樣幫助我們.
藉著歌羅西書.
一卷很特別的書.
提醒我們.
我們真的謙卑下來.
尊主為大.
教會在基督的計劃裡面.
是很重要的.
同樣我們每一個人.
都要讓主居首位.
尊主為大.
主是我們生命的中心.
求主幫助我們.
叫我們每一個人.
都在我們的生命裡面.
高舉耶穌.
尊主為大.
為主而活.
願神賜福給各位.
\newpage



\section{羅馬書 15:1-6}
\label{sec:032029}
\textbf{不求自己的喜悅}~~~~陳雅莉傳道
\newline
\newline
連結: \href{http://www.church.com.hk/acms/content.asp?site=cdc\&op=show\&type=product\&code=032029\&layout=sermon}{\texttt{032029}} ~~~~ 語音日期: 2026-04-26
\newline
\newline
\hyperref[sec:032028]{\small{< < < PREV SERMON < < <}}
~
\hyperref[sec:index_chronic]{\small{[返順時目]}}
~
\hyperref[sec:032030]{\small{> > > NEXT SERMON > > >}}
\newline
\newline
羅馬書 15:1-6
\newline
\begin{longtable}{cl}
\hline
\hline
章節 & 經文 (和合本修訂版)\\
\hline
15:1 & \begin{tabularx}{0.7\textwidth}{X} 我們堅強的人應該分擔不堅強的人的軟弱，不求自己的喜悅。 \end{tabularx} \\ \\ \relax
15:2 & \begin{tabularx}{0.7\textwidth}{X} 我們各人務必要讓鄰人喜悅，使他得益處，得造就。 \end{tabularx} \\ \\ \relax
15:3 & \begin{tabularx}{0.7\textwidth}{X} 因為基督也不求自己的喜悅，如經上所記：「辱罵你的人的辱罵都落在我身上。」 \end{tabularx} \\ \\ \relax
15:4 & \begin{tabularx}{0.7\textwidth}{X} 從前所寫的聖經都是為教導我們寫的，要使我們藉著忍耐和因聖經所生的安慰，得著盼望。 \end{tabularx} \\ \\ \relax
15:5 & \begin{tabularx}{0.7\textwidth}{X} 但願賜忍耐和安慰的神使你們彼此同心，效法基督耶穌， \end{tabularx} \\ \\ \relax
15:6 & \begin{tabularx}{0.7\textwidth}{X} 為使你們同心同聲榮耀我們主耶穌基督的父神！ \end{tabularx} \\ \\
[1ex]
\hline
\hline
\end{longtable}
$^{1}$過去一個星期我不知道你又經歷了什麼呢?.
如果我想你用一個天氣來形容.
你過去一個星期的經歷.
我不知道你會用一個什麼天氣呢?.
會是晴天 好像這個星期的天氣一樣.
陰天 還是橫風橫雨呢?.
不知道你的心情 你的情緒又是怎樣呢?.
如果要你用一個圖像來形容你的心情.
不知道你會用什麼樣子呢?.
哈哈笑 微微笑 還是愁眉苦臉呢?.
無論過去你經歷了什麼.
我盼望你仍然看到.
在你每一個情緒 每一個經歷裡面.
上帝都與你同在.
上帝感受你的感受 上帝知道你的經歷.
上帝都知道你的心情.
不過有時是我們看不見祂.
或者是我們刻意忘記了祂.
為了今天我們仍然能夠在神的殿裡面.
還能夠回到教會.
如果你手上有聖經.
無論是電子也好 紙本也好.
我都希望你今天可以翻開聖經.
翻開你的聖經來到《羅馬書》第十五章.
今天我們看的經文在《羅馬書》第十五章.
雖然我們今時今日可能有時都很方便.
聽到的時候又有PowerPoint.
但我都很想偶爾叫大家翻一下聖經.
翻到《羅馬書》.
你們一直找的時候 我一直說.
《羅馬書》是寫給當時兩批不同的信徒看的.
《羅馬書》是寫給保羅說是致羅馬人的書信.
其實當時在羅馬有兩批的信徒.
有猶太的信徒 也有羅馬的外邦信徒.
有兩邊的人 你能夠想像.
兩邊的人 他們當然就會有兩個不同的意見.
兩種不同的生活模式 兩種的看法.
可能無論在價值觀 生活 很多的儀式上都不一樣.
今天我們來到看《羅馬書》第十五章的時候.
我的題目是「不求自己的喜悅」.

$^{41}$可能你會想 為什麼無緣無故要不求自己的喜悅呢?.
人人都是求自己的喜悅的事.
為什麼會不求自己的喜悅呢?很奇怪.
如果我們看一看《羅馬書》第十五章的前一章.
是《羅馬書》第十四章.
當時它就說到 原來在羅馬裡面有兩批的信徒.
一批就是羅馬的外邦信徒 一批就是一路以來的猶太信徒.
原來在第十四章說到 羅馬的外邦信徒.
他們是神信主的 他們就百無禁忌.
什麼都能吃 信了耶穌 憑信心怎樣都行.
所以他們就吃肉又吃菜 他們是堅固的人.
他們覺得自己是堅固的人.
而猶太的信徒在第十四章說 他們只吃蔬菜 不吃肉.
當時聖經就形容 他們覺得這些人是不堅固的人.
是軟弱的人 因為你們這麼多規矩 這麼多禁忌.
信了耶穌的嘛 這麼沒有信心.
就是因此吃肉不吃肉 你吃菜我不吃肉.
就開始在第十四章形容他們彼此論斷.
甚至是叛島人.
所以在第十四章的時候 保羅就提醒他們.
不要因為這樣論斷別人 不要因為這樣叛島人.
有時我們都會這樣的.
好像我們守規矩守得嚴謹一點.
我們跟聖經走得足一點 就覺得厲害一點.
我是一個虔誠一點的人.
但有些人是反過來看的.
覺得信了耶穌 很多東西都是憑信心的.
百無禁忌 不要那麼多規矩.
搞那麼多東西 多餘的.
可能今天我們都會落在這兩種心態裡面.
覺得我的看法是對的.
有人就覺得你的看法很多餘.
就是因為這樣 大家都有自己的一套標準.
有自己的一個想法 有自己的一些背景.
大家都有自己的意思.
所以大家就按著自己的意思去做不同的事情.
又見到別人不跟自己的意思的時候.
不跟自己所想的規矩的時候.
就開始彼此論斷.
甚至聖經形容他們彼此輕看大家 看不起大家.

$^{81}$所以保羅提醒他們.
不論吃還是不吃.
各人都是為主的緣故.
在14章的時候 保羅已經提醒他們是為主的緣故.
因為我們已經信了主了.
我們是為主而活的.
所以他說其實我們不可以論斷彼此.
只有主可以論斷我們.
來到15章 保羅不僅要他們不要不論斷別人.
不要叛倒別人 還要他們額外做多一點.
保羅進一步要求他們.
你們不僅不准論斷別人.
你們還要擔當別人.
你們還要擔當別人的軟弱.
你聽到這裡 你覺得這是一個什麼樣的要求呢?.
他和我有不同意見 他和我的標準不一樣.
我不罵他也給他面子.
我不和他爭辯 不和他批評.
已經很給面子了.
還要我擔當他?.
還說要接納他?.
不知道你們會不會這樣想呢?.
覺得保羅這個要求是什麼樣的要求.
真的很奇怪.
讓我們一起 今天再用一顆謙卑的心.
再一次來到聖經的面前.
再一次來到神的面前.
去篩讀羅馬書給我們今天的教導 好嗎?.
保羅在15章一節說.
我們堅強的人.
要分擔不堅強的人的軟弱.
不要求自己的喜悅.
如果你單從字面來看 很合理.
堅強的當然擔當不堅強.
但你要記得 剛才我們的背景是說什麼?.
什麼是堅強的人?.
不是真的有錢一點 有力一點的那些叫堅強的人.
他們是說有信心一點.
我不吃 我吃肉又吃菜 我有信心一點.
這些是羅馬的外邦信徒.

$^{121}$他覺得自己厲害一點 威一點.
他覺得猶太的信徒沒有信心.
只是吃菜 他們是軟弱的.
保羅說 如果你既然這樣覺得.
那堅固的人就要去擔當不堅固的人的軟弱.
再想深一層這個什麼道理呢?.
我堅固是我的事.
你軟弱是你的事.
為什麼我要擔當你?.
你軟弱關我什麼事?.
我堅固就是我的厲害 是不是?.
為什麼我要來擔當你?.
保羅還說不求自己的喜悅.
我為什麼要為你不求我自己的喜悅?.
前一章保羅很消極的去說.
你們不要叛島人 不要興旱人 不要論斷人.
來到這裡 保羅是積極的說.
你們堅固的要去擔當不堅固的.
再細想一下 為什麼我們要這樣擔當?.
如果你擔當 是因為你堅固.
你能夠擔當 是因為神使你的信心堅固.
當時也是一樣.
是因著神讓愛邦的信徒去明白.
他們已經在主裡面.
不需要再搜救那些葛禮 猶太的律法.
所以他們是堅固的在信心裡面.
以致他們可以跳出那些律法的規限.
但是否就是因為你自己跳出了.
你自己明白了.
見到那些還未明白的人 就任由他呢?.
你就甚至看不起他呢?.
擔當 不是因為對方值不值得我們擔當.
不是因為他聰明不聰明 好不好.
所以去擔當他.
保羅說是單單因為你是堅固.
所以你就要去擔當.
原來聖經教導我們看到軟弱的人的時候.
不是去踩多一腳 或者假裝看不到.
聖經教導我們是你要去擔當.
在《迦拉太書》是這麼說的.

$^{161}$它說如果你們各人的重擔要互相擔當.
這樣就成全了基督的律法.
擔當 不是因為對方很值得我擔當.
擔當是因為當見到別人的軟弱.
別人有不明白 別人有未理解的時候.
我們願意去擔當他 為成全基督的律法.
擔當 是不求自己的喜悅.
保羅勸勉信徒 你不要求自己的喜悅.
你不要堅持自己的己見 不要以自己為中心.
你要以別人的益處行先.
堅固的人 好像很不值得.
要去不求自己的喜悅.
難道我堅固 我不可以就單是顧自己嗎?.
我為什麼要不求自己的喜悅?.
不過我想在這裡說 我們不要走極端.
不要覺得保羅在說.
你從今以後不可以有自己的喜好.
不可以有自己的中意.
不可以有自己的堅持.
然後你就要詛咒別人 不是的.
保羅不是在說這樣.
意思不是說你不可以有自己的意思.
有自己的想法.
聖經也不是要我們這樣.
如果我們看聖經 不要脫離上文下理.
不要斷章取義.
還記得上一章是說一些很具爭議性的議題.
就是吃不吃肉 是不是?.
在這裡是說一些在教會裡面.
甚至在社會裡面一些很具爭議性的議題上面.
我們能不能夠尋求上帝的旨意.
而不是只求自己的意思.
擔當 要不求自己的喜悅.
是要用愛去接受和忍耐.
因為聖經說.
我們各人務要使鄰者喜悅.
使對方得益處 得造就.
這個是恩著愛 是恩著愛.
可能你會想 其實吃不吃肉很小事.
今天如果我們隔壁的弟兄姊妹跟你說.

$^{201}$我不吃肉很小事 是不是?.
你不要跟他爭拗.
你不明白為什麼當時的人這樣也要爭拗?.
但如果你有一點點歷史的背景.
你就知道在猶太人的背景裡面.
這個不是小事.
對他們來說是一些很核心的價值觀.
對他們來說 在猶太人的背景裡面.
吃不吃肉是甚至牽涉到結與不結的問題.
不是一個容易妥協.
不是一個容易接納的議題.
是非常具爭議性的 在當時.
因為在猶太人的背景裡面.
他們實在覺得.
這個是我們跟外邦人最大的分別.
就是他們有很多不結.
昔日是吃不吃肉?.
今天呢?.
今天其實我們也有很多這樣的標準.
今天我們也有很多這些.
林林總總不同的看法.
跟對方有很多不同的意見.
我相信有些人可能也聽過一些極端的例子.
就是說當有信徒他病了不舒服的時候.
應該不用看醫生的.
你有沒有聽過這些說法呢?.
只要關上房門 用信心禱告就可以了.
如果你不停去求醫.
其實你是沒有信心的.
這些是很極端的看法.
甚至乎我也聽過一些信徒.
為了幫不幫襯一家店鋪.
都吵了一架.
因為有人覺得.
要支持跟我們基督教信仰.
立場一致的商鋪.
如果那些跟我們基督教立場不一致的商鋪.
我們就應該抵制它.
不要幫襯它.
有人就會覺得.

$^{241}$抵制這個行為本身就不符合真理的.
耶穌基督不喜歡我們這樣抵制人的.
但又有人會覺得.
不是 你用你的金錢去支持這些商鋪.
就等於你支持邪惡.
支持它們繼續營運.
所以應該要抵制.
很多這些林林總總的看法.
很多這些很強烈的意見.
今天每一天都仍然發生.
其實我們是不是不可以有自己的意見呢?.
我們是不是不可以有自己的看法呢?.
其實有自己的看法,有自己的意見.
本身是沒有錯的.
但關鍵是你怎樣用正確的態度.
用正確的態度.
有智慧的去聆聽彼此對方的意見.
過去幾年.
如果有弟兄姊妹在這裡.
已經超過十年,二十年的弟兄姊妹.
你一定會看到過去幾年教會的轉變.
跟三十年前,二十年前.
甚至乎十年前的教會.
今天已經毫不一樣了.
短短幾年裡面.
一大批的新移民來到教會裡面.
從前可能二十年來.
紫雲頓教會都是三,四十人.
我們很和諧 是不是啊?.
我看到有些鴨頭鴨頭.
近這幾年.
突然間帶了很大批的新移民來到我們中間.
我們都彷彿好像昔日.
羅馬書裡面所說的猶太人和外邦人的信徒.
新來的弟兄姊妹.
我們都帶著自己昔日的教會來到 是不是啊?.
如果你不是在這裡信主的.
你過去在香港有舞會的.
你都帶著你的舞會來到.
其實我們每個人後面都有一間教會的.

$^{281}$帶著自己很喜歡的教會.
因為我很感恩 為你很感恩.
你很愛你昔日的教會.
你很懷念你昔日的教會.
所以你帶進來的不單止是昔日的制度.
昔日的系統,昔日的模式,昔日你的經驗.
你還是帶了一份你昔日對你的舞會很深的感情.
這些都是我們很喜歡的.
想用過去我的經驗,我的方法 帶進來這裡.
所以其實現在我看下來.
這裡其實是坐了二,三十間教會的弟兄姊妹的.
要放棄一些我們自認為對的,好的.
我很懷念的,我很喜歡的方法.
這個不只是少許的遷就和改變那麼簡單.
我明白這個是在說你裡面有一些底線.
如果要你放棄這些底線.
裡面是包含了一些情感的捨棄.
甚至你會懷疑你身份.
你事奉上身份的認同的問題.
所以不是簡單說一句「遷就一下」.
那就行了.
而是你在裡面有一些很核心.
你認為正確的,你認為對的.
你認為美好的價值.
你沒有辦法移動.
就好像昔日的猶太信徒都一樣.
他們有一些守了多少年.
他們祖宗一直守下來的律法.
他沒有辦法移動.
又或者我都知道有一些弟兄姊妹是另外一種.
她在過去昔日的教會就是很不開心,很不快樂.
可能在過去的舞會裡面受過很多傷害.
甚至她就帶著很多傷痛和失望.
希望來到英國重新開始.
有一個新的景象,新的教會.
希望過去所有傷害我的人,傷害我的事.
在這間教會裡面從此不再發生.
我憧憬一間完美的教會.
希望所有過去的東西都沒有了.
不過可惜進來沒多久又發現好像差不多而已.

$^{321}$問題又出現了.
教會確實常常都是這些狀況.
我們很容易就輕看,論斷那些和我們不同意見的人.
當然有很多問題在聖經裡面是很清晰的.
應該做還是不應該做.
牽涉到真理,牽涉到罪.
我們很清晰的,大家都很一致的.
但是在一些小小模糊,稍微具爭議性的問題時.
我們就在那裡爭議.
開始造成彼此的不同心,彼此的不合一.
國語裡面有一句說話我覺得說得挺好的.
不知道你們有沒有聽過.
叫做「有關係就沒關係」.
有沒有人聽過呀?.
什麼意思呢?.
有關係就什麼都沒關係.
只要大家有好的關係.
就什麼都無所謂了.
沒關係.
小小事沒關係沒關係.
有關係就很多事都無所謂,無關係.
容易包容得過去,容易接納得過去.
不用那麼執著.
大家關係好,彼此信任,彼此包容.
最近有一個姐妹跟我分享.
你們不用猜是誰,不在你們中間的.
有一個姐妹跟我分享.
她說近幾年覺得來英國之後事奉很困難.
她是新移民.
她說從前她在教會裡面事奉的時候.
都是很忙的,很多事奉.
但因為大家彼此的信任,彼此的認識.
當每一次她做錯一件事.
或者她很忙,做漏了一些事.
忘記了提醒,忘記了通知的時候.
弟兄姊妹總是會私下問候她,關心她.
總是會問她:你是不是忙得太多?.
有沒有名字我可以幫你?.
你是不是忘記了?我提醒你.
但她說近幾年不是這樣.

$^{361}$但近幾年裡面她發現.
她每一次在教會裡面做錯一件事.
忘記了通知一句,忘記了寫多一隻字.
大家就會覺得她是失職.
大家會在群組裡面公開地問她.
為甚麼你不通知?為甚麼你不安排好一點?.
為甚麼你不通知大家?為甚麼你會忘記了?.
她說她感覺到對方一開始已經認定.
她是不尊重人,她不懂得怎樣跟人溝通.
她甚至會自作主張.
甚至她跟我說:有時其實真的不是甚麼錯.
甚至有時她是出於好意.
她做的事都是為教會好的.
她說:但只不過因為一時忘記了,一時漏了.
甚至一時跟對方的做法有點不同.
大家只是關心她有沒有做好一件事.
有沒有按規矩辦事.
但沒有人過問過她為甚麼她會忘記了.
為甚麼她會做漏了.
甚至沒有人關心她為甚麼她會這樣做?.
她告訴我.
在近年裡面她不停在事奉上要說的是.
對不起,對不起和對不起.
到底從何時開始.
教會只是變得規矩,制度.
從何時開始教會沒有愛,沒有包容,沒有接納,沒有關心.
當然我不是說真的有大錯的時候.
甚至有罪在我們中間我們都要容忍.
不是.
教會不是一個不可以指證錯誤的地方.
不是.
教會更加不是縱容罪和惡的地方.
絕對不是.
不過即使有罪.
也不要說不是罪,即使有罪.
即使有罪在我們中間.
在這一切之前.
我們是先指責,批評和責備.
還是先去關心,了解對方在發生甚麼事?.
真的.

$^{401}$當沒有關係的時候.
甚麼都有關係.
各人都從自己不同的切入點去思想這些事的時候.
各人都對神的話語有不同的理解,不同的認識的時候.
保羅提醒我們.
你們要在神的話語上面紮根.
先回去好好讀聖經.
你才能夠分辨究竟是甚麼才是真的重要.
究竟甚麼其實根本無關重要.
不需要這麼執著.
不需要這麼認真.
究竟有甚麼在真理上面我們必須要抓緊.
不能放棄的.
我們必須要認真討論.
今天同樣.
在教會裡面一定會持不同的意見.
上帝不是不容許的.
否則上帝很容易的.
將你們全都打了一針就變得很乖了.
每個人都很一致.
你明白我,我明白你.
大家一起朝著上帝走.
其實很簡單的.
上帝要做這件事用甚麼呢?.
今晚就打暈你.
明天醒來就變成一樣的想法.
容易得不得了.
用不著寫這麼多聖經來叫你們合一.
上帝給我們自由意志.
給我們選擇.
給我們在真理上面去追求.
給我們去經歷它.
給我們去體會它.
我們能不能夠.
當我們聽到對方跟我們的意見不一樣.
當我們看到別人做法跟我們不一樣的時候.
我們不是覺得對方是跟我們對立.
我們能不能夠覺得.
原來對方有一個不一樣的看法.
我們能不能夠覺得其實是對方開闊了我狹隘的眼光.

$^{441}$我們能不能夠覺得.
原來對方這樣看東西是我看不到的.
其實對方是在幫我看這個世界闊一點.
最重要的是.
當我們彼此討論.
彼此講出大家的感受的時候.
我們不是憑著血氣去講.
我們能夠是靠著基督聖靈的指引.
去講出.
求神給我們智慧.
保羅很清楚.
其實基督的信仰不是教大家飲不飲食的問題.
所以你看到保羅.
他沒有在羅馬書要討論究竟.
保羅是沒有結論的.
他最後沒有說你們應該吃肉還是不應該吃肉.
沒有說猶太的信徒你們就應該憑信心吃肉.
外邦的信徒你們就要嘗試吃菜.
沒有的 完全沒有.
保羅是指向一個更高的.
他說當你們有這樣的看法.
他們有這樣的看法的時候.
你們能不能夠各人務必讓鄰人喜悅.
使人得益處得早就.
你們能不能夠以這個位置高.
保羅不單止叫人去擔當.
你看到保羅他是說我們.
他不是叫人去擔當.
是你們去擔當 你們要使人喜悅.
他是說我們 他包括了自己在裡面.
保羅寫這封書的時候.
他覺得他都需要去擔當.
他都要使人得喜悅得益處.
今天當我們 如果你自認是堅固的.
當你看到不堅固頂尖的時候.
是不是覺得對方很煩 很奇怪 這樣看事情.
這樣做事.
還是你願意去擔當他的軟弱呢?.
他做漏了什麼不要緊.
他有什麼幫手.

$^{481}$你是不是忘記了.
我趕快幫你補位.
可以這樣嗎?.
保羅還進一步提醒我們.
為什麼要不求自己的喜悅.
他講出了兩個原因.
第一個他說是因為基督都不求自己的喜悅.
如經上所記 肉麻你的肉麻都落在我身上.
這是耶穌基督.
保羅提醒我們能夠不求自己的喜悅.
擔當別人是因為我們的主都不求自己的喜悅.
堅造教會是用愛心 不是用規條的.
是靠生命 不是靠手段.
當然 我不是說教會可以無規無矩.
不是這個意思.
而是在規矩之前 先有愛.
這令我想起世世代代的信徒.
都欣賞不同的立場.
無論是政治立場 社會立場 經濟立場.
教育等等的議題上不停爭論.
互相攻擊 彼此不尊重.
或者我們都經歷過這些不信任 不認同.
造成了很多傷害.
今天如果在座還有未信主.
未認識耶穌基督的人坐在我們中間.
真的會覺得你們這些人還叫我來信耶穌.
我們真的要好好反省.
基督的信仰從來都是顛覆世界的價值觀.
要效法主 不求自己的喜悅.
我再一次說 這是選擇.
耶穌自己做了榜樣 但祂竟然讓你選擇.
是你的選擇.
其實從《羅馬書》十二章一開始.
一路寫給這兩班信徒.
成卷《羅馬書》都是寫給他們.
很多時候我們會覺得《羅馬書》只是講恩信清義.
不是 不是這麼簡單.
其實《羅馬書》為什麼要寫恩信清義.
就是因為當時有猶太的信徒.
他們覺得不是恩信清義 是要靠行為清義.

$^{521}$所以保羅斯要這麼辛苦 寫這麼多篇章.
《羅馬書》由第一章至第十一章.
都是寫給當時的猶太信徒看.
讓他們知道今天不需要再行國禮了.
耶穌他們已經來了 我們是恩信清義.
但其實由十二章開始 一路去到尾.
保羅是寫給外邦的信徒看.
你不要以為自己恩信清義就很厲害.
你不要看不起猶太的信徒.
就是因為他們當時覺得自己已經有耶穌基督了.
我已經看通了 我不用再守律法了.
就看不起猶太人.
保羅就是寫給他們 你們不要這樣做.
你恩信清義是恩著神的恩典.
你沒有什麼這麼厲害.
你是要效法基督的榜樣.
因為基督都不求自己的喜悅.
還有第二個原因.
第二個原因是 保羅說因為聖經.
是為了教導我們而寫的.
因為聖經已經寫下了要教導你們.
要藉著你們的忍耐.
和聖經所生的安慰.
讓你們得著盼望.
如果你只是看到我要效法主.
你就會覺得很辛苦.
真是苦 又要不求自己的益處.
又要不求別人的喜悅.
又要擔當別人.
你會覺得這個信仰很苦.
但保羅不是停在那裡.
保羅說你是會得著盼望.
你會有安慰 你會得著盼望.
你能不能夠看到這一切的最終.
你是要得安慰 你要得盼望.
聖經寫給不是我們一個人看的.
聖經是寫給我們看的.
是我們這個共同體.
不只是我 是我們.
而我們是能夠一起共享這個盼望的頂子.

$^{561}$只有在基督裡面.
我們才能夠一起看到這個盼望.
未信的人看不見.
這個是我們共同擁有的.
所以保羅提醒我們.
我們要一起留心.
我們要一起留心聖經的教導.
我們一起要常常查經 為什麼呢?.
並且不只是查.
你要實踐在你的生活裡面.
你怎樣要一起和你的弟兄姊妹.
一起去實踐信仰.
並最終能夠因著忍耐得安慰.
最後得盼望.
在受苦的時候.
保羅提醒我們.
受苦的結果 最終你會得著盼望.
我們持守這個盼望.
我們就能夠繼續健全地合而為一.
為著這個盼望 我們一起努力.
當信徒越捨棄你自己.
當時你是會覺得越苦的.
但其實是樂來的.
因為你望著前面那個榮耀的盼望.
反過來一個人.
如果越損人利己.
當時好像很快樂 是不是? 利己嘛.
但其實.
有聖靈提醒你的時候 你是很苦的.
你想起自己的過犯.
你就難過了.
所以連主自己都放下他一切的權利.
這個就是我們的目標.
效法主耶穌基督.
在他裡面找到合一.
不過保羅再說.
合而為一都未是最終的目標.
合而為一最終最終目標是為了什麼?.
羅馬書15章第5第6節這樣說.
「但願賜忍耐和安慰的神.

$^{601}$使你們彼此同心.
效法基督耶穌.
為使你們同心同心.
榮耀我們主耶穌基督的父神」.
這個是最終目標.
合而為一最終是為了什麼?.
不是為了開心一點.
為了你好我好沒有爭拗.
是為了一同來榮耀我們的父神.
我很喜歡不得心牧師Eugene Peterson.
他在信息版這個聖經裡面.
他這樣翻譯.
這個翻譯很生動.
他說那時候我們將成為一個合唱團.
不僅僅是我們的聲音.
而是我們整個生命在和諧裡.
為我們主耶穌基督.
我們的父神獻上壯麗的讚美詩.
他形容我們將會好像合唱團.
合唱團讓我想起這幅圖畫.
我從中學開始就被選入朗誦隊.
不過就不是個人朗誦.
而是集體朗誦.
你們有沒有人參加過?.
中一當時每個學生一進去那間學校.
都要參加各方的活動.
我運動又不行.
這樣又不行 那樣又不行.
最後被老師選了你來朗誦吧.
就這樣入了宗師集誦隊.
當時四五十個同學一起去練習.
我的老師很嚴格.
所以我學校的集誦隊是很出名的.
我們連續很多年.
在全葵清區的校際朗誦節.
宗師集誦都是拿冠軍的.
很厲害吧.
練習得很辛苦 很嚴格.
表情管理 笑多一點也不行.
嘴巴合起來一點也不行.

$^{641}$朗誦要誇張.
口要很放鬆 面部表情.
咬字要很準確.
大聲一點也不行 小聲一點也不行.
到我大學的時候.
又被選入了大學的合唱團.
到後來信主神就呼召我加入師班.
一唱就十多年了.
每個星期都練習.
當時我的師班是每個星期四晚都練習.
我們是除了八號風球之外.
一年三百六十五天都要練習的.
不論是朗誦 唱詩.
幾十人走在一起.
是一件非常不容易的事.
我還記得以前我師班的指揮非常嚴格.
很出名的.
很嚴 要求很高.
我們整個師班都對她又敬又畏.
但我們整個師班都很尊重她.
因為在她身上學了很多東西.
她是很嚴格的.
她的耳朵很靈敏.
你稍微沒有練習 她一定聽得到.
她會走過來指著我們說.
「你走音」.
四十多人面前.
然後就問「你這個星期回去有沒有練歌?」.
她常常用心教導我們.
師班不只是用口唱的.
你要用心.
用心是包括了你擺相的時間.
用心是包括了你在家中用心地練習.
她經常叫我們「你不單要用心」.
「你還要用耳朵聽一下」.
「你要看 你要觀察」.
「你要聽你前後左右的人」.
「他們有沒有需要幫忙的地方」.
「不然就是你自己唱得好」.
「你有沒有聽到你旁邊那個人的走音」.

$^{681}$「你願不願意完了再坐下來幫他」.
「你有沒有聽到你前面後面那個」.
「他的節拍有少許不對」.
「願不願意再花時間和他一起練多一會兒」.
「你有沒有用心看指揮的指示」.
「有沒有留心聽樂器和你的配搭」.
「有沒有聽到你自己比別人大聲還是小聲」.
「如果只是你一個人唱得很好」.
「音調正確,高音準,低音甜」.
「只是你一個人正確是沒有用的」.
第一個師班裡是講求和諧合一.
每個人都要自己回家下功夫.
我以前也是這樣 要下功夫.
我由一粒豆豉都不會看入師班的第一天.
到一路一路學 一路花心機,花時間.
整天要在家練運氣 不是Lucky那些運氣.
就是呼吸 丹田那些運氣.
練咬字,練發聲,下不少功夫才回來.
不是隨便唱 小聲唱了一會兒.
「不要緊啦 唱都沒人聽到」.
又不是大大聲唱到只聽到你的聲音.
主是這樣形容.
他很想見到我們好像一個師班和諧合一.
如果人人都自我中心 只顧自己.
不謙卑,不看其他人.
不會有美妙的聲音 更加不能榮耀神.
我在敬拜隊裡也講過這個比喻.
我經常形容我們好像一群茶葉蛋.
一群茶葉蛋.
茶葉蛋要怎樣才入味呢?.
「剝碎」它 是嗎?.
是的 茶葉蛋要「剝碎」它才入味.
你要打爛殼 你的外殼要「破碎」.
茶葉的味道才能浸入蛋裡面.
如果我們的生命要有一些味道.
要在主裡面多一點成熟.
我們才能「破碎」我們的自己.
我們的生命才不會原地踏步.
你看茶葉蛋上面的裂痕.
一點都不漂亮 是不是?.

$^{721}$很噁心 一點都不好看.
不完美.
但是打開來的時候.
入味了 很香.
如果我們看不到這個神榮耀的景象.
我們很難彼此接納.
如果我們仍然不「破碎」就這樣浸入茶葉裡.
是的 你的殼很漂亮 是不是?.
但是打開來 你裡面還是白雪雪的.
唯有以心 以合一的靈.
來取代我們信徒之間彼此的排斥.
我們才能真正榮耀神.
其實你留心看.
其實保羅沒有祈求當時兩班的信徒.
要那一班就那一班 沒有.
保羅沒有叫猶太信徒.
你要聽完外邦信徒的話.
更加沒有叫外邦信徒.
你從此就要守猶太人的律法.
他反過來給他們一個新的向導.
保羅祈求神.
不是要他們毫無爭論 毫無意見.
要大家完全一致.
保羅反而是禱告.
他們可不可以有一個新的目標.
來克服彼此之間的差異.
這個新的目標是.
你能不能夠看見耶穌基督的榮耀.
神其實不是希望我們冷漠地一致.
有些人可以說.
我什麼都不看 什麼都不聽 什麼都不參與.
那就沒有交涯了 是不是?.
我都沒有意見 隨便吧.
這是冷漠.
當人人都有自己的意思.
當很想貢獻給教會的時候.
我們難免會有不一致的時候.
但是你願不願意在這個不一致之中.
仍然打開你的耳朵彼此聆聽.
願不願意求主給你足夠的智慧和膽量.

$^{761}$有你的看法 有你的意見的時候.
能夠讓對方聽得見明白.
關鍵不是在血氣下進行.
而是在神的靈的庇蔭下進行.
如果每一個人都是朝著讓教會得最大的益處.
不僅僅只是看到自己看到的東西.
而是看到更闊更大的圖畫.
我相信不同的觀點.
不僅不會削弱教會的合一.
反而會讓教會更堅強,更寬闊.
歸根究底.
其實聖經讓我們看到.
神在召我們每一個不一樣的人在一起.
堅持守信仰最基要的元素.
心理在你裡面是不可動搖的.
但是在一些沒有定論的事上.
我們能不能夠開放彼此聆聽.
保持忍耐,包容的態度.
仍然能夠相信.
大家都可以在自己的空間上榮耀神.
我們能不能夠有這樣的體見呢?.
我要總結了.
還記得今天我一開始的時候問大家兩個問題嗎?.
不知道你過去一個星期的心情是怎樣呢?.
不知道你經歷了些什麼呢?.
如果我們仍然覺得我所有的經歷.
其實沒有上帝在我中間的話.
你仍然看不見.
你仍然看不見上帝要怎樣去塗造你的生命.
祂怎樣幫助你每一天走你人生的路.
無論在生活上,家庭上,事奉上.
如果你能夠看見.
其實榮耀的主每天都與我們同在.
每天都與你一起經歷.
並且祂感受你一切的感受.
我們就不會常常把問題放得很大.
而看不見其實上帝的能力是更大.
你要看到一切的問題上.
即使問題再大.
沒有一個問題比上帝的能力更大.

$^{801}$如果保羅今天仍然在生.
他會寫一封什麼信給我們教會呢?.
不知道他會不會像寫羅馬書一樣.
寫給我們教會呢?.
不知道他會不會跟寫羅馬書的教導一樣呢?.
從一開始我們問.
我怎樣擔得起別人?.
我為什麼要擔當他?.
他的軟弱為什麼不是他自己來擔?.
為什麼要我擔?.
為什麼我要不求自己的喜悅?.
願這些問題不是讓我們去想.
我不願意.
我不要.
而是想我要效法基督.
我要同心同性榮耀神.
不是為自己.
而是單單為榮耀神.
保羅不是說你們.
保羅說是我們.
我也一起.
願主繼續賜忍耐給我們.
賜安慰給我們.
讓我們彼此同心.
彼此接納的裡面.
我們最終能夠一起歸榮耀彼實.
我們一同禱告.
你們本來就在愛中有團契.
你們就是愛.
聖父,聖子,聖靈三位一體的神.
本來就有內在的團契.
彼此相愛.
你成為我們最好效法的榜樣.
聖經不是說神有愛.
聖經說神是愛.
你本來就是愛.
當我們因信進入你的生命裡面的時候.
我們就得著這一份從你而來的愛.
使我們能夠效法你的樣式.
跟隨基督的能力.

$^{841}$跟隨基督從死而復活的能力.
已經在我們每一個人的身上.
使我們能夠活出基督.
求主幫助我們.
我們今天仍然有很多的限制.
有很多的不能夠.
甚至當我們每一天回想的時候.
我覺得在當中有很多的傷害.
有很多的傷痛.
但求主讓我們看見.
這些不是對方要為我造成的傷痛.
我不是要繼續在這裡忍耐對方.
那是讓我們看見.
原來主將我們擺在一起.
有祂的心意.
我們究竟怎樣.
能夠與我們的弟兄姊妹同心同心.
這個今天是我們每一個人都要學的功課.
願主你繼續賜盼望給我們.
讓我們看到那榮耀的盼望.
我們就能夠俯伏在你面前.
重新謙卑再學習.
願主你的德性在我們中間.
叫那惡者不能夠在我們的教會中搞擾.
願主你的德性.
你的榮耀始終降臨在我們中間.
奉耶穌基督名字祈求.
阿們.
\newpage



\section{約翰福音 4:23}
\label{sec:032030}
\textbf{父神在尋找:真正的敬拜者 (華)}~~~~王亞辰牧師
\newline
\newline
連結: \href{http://www.church.com.hk/acms/content.asp?site=cdc\&op=show\&type=product\&code=032030\&layout=sermon}{\texttt{032030}} ~~~~ 語音日期: 2026-04-26
\newline
\newline
\hyperref[sec:032029]{\small{< < < PREV SERMON < < <}}
~
\hyperref[sec:index_chronic]{\small{[返順時目]}}
~
\hyperref[sec:032031]{\small{> > > NEXT SERMON > > >}}
\newline
\newline
約翰福音 4:23
\newline
\begin{longtable}{cl}
\hline
\hline
章節 & 經文 (和合本修訂版)\\
\hline
4:23 & \begin{tabularx}{0.7\textwidth}{X} 時候將到，現在就是了，那真正敬拜父的，要用心靈和誠實敬拜他，因為父要這樣的人敬拜他。 \end{tabularx} \\ \\
[1ex]
\hline
\hline
\end{longtable}
$^{1}$這是一個豐富又榮耀的旅途.
記得週五的時候我們剛抵達貝南.
當地的主任牧師Pastor Corneli才告訴我們.
隔天的禮拜天啊.
就是貝南總統大選.
哦我想說這是非常特別的日子.
我們可以代表生命和.
一起為貝南這個國家來禱告.
所以呢.
貝南的主任牧師就說是的.
逐日當天讓我們一起來求神.
為貝南預備一位神所約納所檢選.
能夠用神的智慧來治理國家的總統.
之後呢.
Corneli的牧師就靠近我說.
希望你禮拜天講到可以注意一下時間.
這樣大家才有時間去投票.
這個準時結束.
而且當天有十位新朋友.
接受耶穌成為個人的救主.
哈利路亞.
這次行程非常的豐富.
我們拜訪了六個村鎮.
那也參加了.
貝南非洲質詢學校第一屆畢業典禮.
我們總共服侍了超過1000人.
並且帶領了超過80個人.
接受耶穌成為個人救主.
我們將榮耀歸給我們的神.
哈利路亞.
謝謝大家的禱告和這個支持.
相信我們不論在前面打仗和最後面.
看守兵器的都要同得獎賞.
Amen.
讓我們一起來讀一段.
呃經文是在約翰福音4章23節.
這是新一本的翻譯.
我們一起來.
然而時候將到現在就是了.
那用心靈按真理敬拜父的.

$^{41}$才是真正敬拜的人.
因為父在尋找這樣敬拜的人.
這邊說那用心靈.
按真理敬拜父的.
才是真正敬拜的人.
弟兄姊妹們我們每周.
都會來教會敬拜讚美神.
我們在小組四個w.
其中有一個也是 worship.
相信我們大家在個人生活當中.
也有與神親近的.
敬拜讚美的時間.
可是在這邊這個經文卻說.
有一種敬拜.
才是真正的敬拜.
而有一種敬拜的人.
才是真正敬拜的人.
而且父神在尋找這樣的人.
聽姐妹今天我們許多人坐在這裡.
請問我們怎麼樣才能夠確定.
我們是真正敬拜神的人呢.
我們怎樣才知道.
我們是父神在尋找的人.
今天好不好我們跟旁邊聽姐妹說.
祝你能夠成為真正的敬拜者.
我想這是一個非常.
serious的一個話題.
我們都在這邊敬拜讚美神.
我們都在事奉神.
可是有一天突然我們發現.
哇.
父神找的不是我們.
父神看到是別人的敬拜.
那不是虧大了.
我相信神的心意.
是讓我們都能夠成為.
摸著他心意的真正的敬拜者.
所以我們先來看一下敬拜的意義是什麼.
敬拜神事實上是我們.
被造的目的.

$^{81}$出埃及記八章一節.
這是耶和華吩咐摩西對.
法老王的一個宣告.
我們大家都很熟悉.
這邊說.
容我的百姓去.
好.
事奉我.
而我們看到英文NIV的翻譯.
這邊說.
Let my people go.
so that they may worship me.
我們看到這邊事奉和敬拜.
worship and service.
是可以interchangeable.
所以在父神的心中.
敬拜不只是音樂.
敬拜不只是聚會.
敬拜是一個生命的流露.
敬拜是一個態度.
那麼這邊講到說.
容我的百姓去.
好事奉我.
所以事實上神.
拯救以色列人出埃及的目的.
事實上所以.
耶穌基督拯救他們的百姓.
從罪惡裡出來的目的.
弟兄姊妹.
不是讓.
以色列百姓生活會更好.
不是讓以色列人能夠.
有好的對象.
有好的工作.
他的孩子能夠上更好的學校.
更不是.
要發掘自己的夢想.
活出你的獨特追求.
你的夢想.
你的這個人生的標竿.

$^{121}$弟兄姊妹這些都很重要.
但是這邊講的最重要的.
容我的百姓去.
好怎麼樣.
事奉我.
弟兄姊妹原來神創造我們的目的.
原來神賜下救恩的目的.
是要來事奉他.
今天我們在.
西方的這樣一個人本主義.
自我主義的這樣一個薰陶之下.
好像所有事情都是為了我們自己.
甚至連信仰.
連救恩.
連神的拯救.
都是為了我好.
是的神愛我們.
神在我們生命中有計劃有命令.
Yes.
但是這個信仰.
的中心.
是為了神.
這個信仰的目的.
是為了來事奉.
來敬拜他.
那什麼叫做事奉神呢.
就代表我們要來關心.
神的事.
如果今天我們.
得了救恩的好處.
我們仍然只關心.
自己的事.
那很抱歉.
我們的敬拜.
不是神在尋找的敬拜.
如果今天我們接受了救恩.
但是我們只想要來高舉我們自己.
而沒有想到高舉耶穌的名.
那麼很抱歉.
我們的事奉.

$^{161}$不是神所尋找的事奉.
所以今天真正的敬拜.
真正的事奉.
遠遠超過.
台上的敬拜.
遠遠超過我們所謂的事奉的.
那些時間.
事奉的這些項目.
原來真正.
父神在尋找的敬拜者.
原來父神在尋找真正的那那事奉者.
是要來先求神的國.
和神的義的人.
弟兄姊妹.
真的我們必須要來.
面對這樣的一個真理.
一個事實.
今天我們的信仰.
不是為了自己.
今天我們的信仰.
是為了神.
神的國.
高舉耶穌基督的名.
但是如果我們現在這樣的一個.
自我中心的陷阱裡面.
我們就.
我們的信仰我們的事奉.
是會大大虧損的.
我們看到因此父神在尋找.
真實.
敬拜的人 弟兄姊妹.
約翰福音四章二十三節這邊說.
然而時候將到.
現在就是那用心靈.
按真理來.
敬拜父的.
這邊又講到剛才說.
容我的百姓去好事奉我 好敬拜我.
這邊講到神在尋找的.
不只是來事奉來敬拜.

$^{201}$而是用真理.
按心靈按真理來敬拜.
那我們就來看.
在.
舊約和新約時代的人.
他們心目中.
他們閉上眼睛所想像的敬拜和事奉.
是什麼樣子.
可能跟我們現在所想像的.
真的很不一樣.
首先我們來看舊約的.
這個敬拜這個字.
是shakal這個字.
shakal這個字.
我們來看這個圖.
shakal就是謙卑.
屈膝的意思.
原來在舊約的敬拜.
是一個.
謙卑屈膝服服貼於地上的一個動作.
就好像.
這個我們看古裝劇當皇上駕到的時候要怎麼樣.
大家都要服服跪拜.
如果你就是那個硬跪不下去.
或者是剛好腿發麻跪不下去就完蛋了.
你的頭就落地了.
所以事實上在舊約的shakal.
敬拜的一個觀念裡面.
是一個極度的敬畏的一個態度.
十篇九十五章六節這邊就是說.
來啊我們要屈身敬拜.
在造我們的夜華面前.
跪下.
這圖可以先收起來.
所以屈膝代表一個.
就是shakal.
敬拜代表一個謙卑服服屈膝的動作.
但是這個動作後面代表著三個.
這個意思第一個.
就是代表著.

$^{241}$你所有的.
都是屬於他的.
當你跪拜在一個君王面前的時候.
你的財產.
你的家業.
你的家族.
甚至你的性命.
都是屬於你所跪拜的君王的.
哇.
所以我這個.
不知道你們有沒有聽過一個笑話.
就是在中國古代.
是沒有講閩南語或台語的太監的.
為什麼呢.
因為曾經有一個是講閩南語的太監.
當他要吃飯的時候他就說.
皇上.
家崩了.
台語吃飯就是家崩了.
哇怎麼可能說皇上家崩了哇.
所以之後就再也沒有講台語的太監了.
的確我們在跪拜在一個君王面前的時候.
我們連生命.
都是屬於他的.
所以在舊約時代我們看到這個shakal.
是一個.
極其敬畏的一個態度.
一切都是屬於神的.
而第二呢.
就是這位.
君王所說的.
聖旨.
你是不能違背的.
還記得.
又是古裝劇.
聖旨到.
你不能說我不買單.
我沒興趣不行啊.
哇那個如果你.
如果是皇上對你說話你要說什麼.

$^{281}$渣呀是不是.
要渣下去.
這讓你也順便講有趣的事情我發現.
蒙古文的好.
就是渣.
就是渣渣.
像是西班牙話就是ccc.
French就是we we we.
那蒙古話就是渣渣.
我那時候一開始我說哇塞蒙古的.
牧師也太權威了吧.
他一講話旁邊就渣渣渣.
我說哇.
後來發現原來是好好好的意思.
事實上我們在.
shakal在全然降服的.
情況就是.
哇神所說的一切話.
都是要來遵守的.
而第三.
雖然剛才前面講還是很權威啊.
好像沒什麼選擇性啊.
但是第三當你跪拜在神的面前.
這個君王面前的時候.
代表什麼呢.
代表是.
他會.
照顧你一切的需要.
哇這就是令人感到非常的安心啊.
因為當你降服於他.
是一切都是屬於他的.
但是一個好的君王.
他必定會照顧百姓的需要.
確保百姓衣食無缺.
可以平安度日.
所以這個時候弟兄姊妹我們.
降服於誰.
我們將自己的主權交給誰是何等的重要.
你如果降服於一個暴君.
一個昏君.

$^{321}$一個只.
利用你來.
滿足自己的.
對象的時候.
哇那你這個shakal.
你還是要來順服.
但是這個人不能保護你.
我們的神.
當我們全然順服於他降服於他.
透過我們的.
將一切都歸給他.
透過我們完全的順服他所說的話.
我們可以有個確距.
他會照顧我們一切的需要.
他會.
擔我們一切.
所有的責任.
他會來帶領我們.
會保守我們.
帶領我們向前.
所以弟兄姊妹shakal.
就約的這個.
敬拜.
哇原來並不是拿吉他在那邊彈啊.
也不是.
在小組.
在這樣子一個.
放youtube.
但這跟這些其實是很有關係的.
當一個人有一個shakal的一個態度的時候.
無論他在彈吉他.
他在放youtube.
或是他在公司工作.
他在為神而活.
哇神就.
與他同在了.
那我們看見.
就約的時代這個shakal呢.
一個很特別的事情.
就是我們.

$^{361}$地上的王.
是.
要求人要跪拜.
但是我們的神卻給人自由意志.
就是要人甘心樂意來回應.
他的邀請.
成為他的百姓.
我們看到先知以賽亞.
他在.
這個經歷寶座意象的時候.
就是一個非常完整的shakal的一個概念.
以賽亞書六章.
烏西亞王崩的時候.
那時候.
整個政治.
整個世界的情勢非常混亂.
但是以賽亞卻看見一個天空被打開來.
意象有神的寶座在其上.
又有天使.
在高唱.
聖哉聖哉聖哉.
萬軍之耶和華.
他的榮耀.
充滿全地.
這個時候以賽亞看見.
有在這個動盪的時代.
有一個王.
是可以保護他的.
有一個王.
是仍然掌權不能震動的.
而這個王.
又是極其的榮耀.
極其的聖潔.
因此他就自然而然的.
甘心樂意的.
俯伏在耶和華的面前.
全然.
擺上.
而這個時候.
以賽亞就聽見.

$^{401}$神說.
有誰.
我可以差遣誰呢.
有誰肯為我們去呢.
哇 當聖旨到的時候.
以賽亞他其實.
我想他可以說不啦.
因為神是給人自由意志的.
但是在這種狀況之下.
如此的榮耀.
如此的聖潔.
又如此的.
神的.
這個充滿.
這時候以賽亞就回應了.
我在這裡.
請.
差遣我.
這是一個shakara的一個.
敬拜的一個時刻.
我們看見.
以賽亞就成為一個.
神所找到的.
真實的.
敬拜的人.
所以一個真正的敬拜者.
第一個根基.
就是完全的降伏.
將自己交給神.
並且一個為神而活的.
一個態度.
接下來我們看到在新約的.
敬拜這個字呢.
是.
proskuneo.
也是一個俯伏跪拜.
但是又多了一個親吻.
或是一個.
一個親近的關係.
那我們看到.

$^{441}$這個.
這個原文的字在.
三博士.
也不是三博士.
東方的博士.
要來這個.
尋找耶穌的時候.
他就問西律王說.
那生下來做猶太人之王的.
在哪裡.
我們在東方看見祂的心.
特來.
拜祂.
這時候他們用的詞就是.
proskuneo.
我們要來拜祂.
不但要來敬拜降伏於祂.
我們也要來親近祂.
我們.
願意跟祂建立一個.
親近的.
個人的關係.
而我們看到這三.
這個東方的博士.
他們的確.
找到在.
在這個馬槽找到了.
Baby Jesus.
向祂跪拜.
我相信那個時候.
不但是很敬畏的跪拜.
其實也是個非常親密的關係.
最後我不知道有沒有.
像我們平常看小孩子.
好可愛.
耶穌基督未來的君王.
我不是很確定.
但是他們的確就.
歡喜快樂的.
離開了.

$^{481}$弟兄姊妹在新約的敬拜.
仍然是一個.
福福.
一個降伏.
但是多了一個.
親近的關係.
因此我們看到撒旦.
在試探耶穌的時候.
他也同樣用這個字.
他說你若福福.
拜我.
你若.
Pruskunio.
來拜.
叫耶穌來拜撒旦.
我就將一切榮華富貴都給你.
原來撒旦所要的.
不只是耶穌基督.
跪下去.
他更要奪取耶穌基督的.
心啊.
其實是個非常.
令人作嘔的這樣的一個試探.
那時候耶穌怎麼回答.
撒旦推我後面去吧.
因為經上記得說.
我們當單單的.
敬拜我們的神.
Amen.
所以一個真正充滿的.
一個真正的敬拜者.
他是.
有一個敬畏的心.
但同時呢.
他有一個與主親近的.
一個關係.
我們怎麼樣能夠.
看到.
與主的關係是能夠親近的呢.
特別在新約.

$^{521}$父神真的在尋找.
真正的敬拜者.
他不只是在天上尋找.
他竟然拆開了耶穌基督.
他的獨生愛子.
來到了世上.
來親近人.
33年半的時間.
與人相處.
道成了肉身.
住在我們中間.
充充滿滿有恩典.
有真理.
不但如此.
他看見了人的軟弱.
他看見了人的污穢.
他看見了人的失望.
自私與黑暗.
他不但沒有來責備來定罪.
他卻來接納.
他卻來產果.
他卻賜下醫治.
與救恩.
世上約翰福音4章23節這個經文.
這個敬拜的真正的意義.
耶穌是對誰說的.
是對緊編的婦人說的.
緊編的婦人.
他有.
五次失敗的婚姻.
第六次的婚姻.
他也是沒有什麼信心.
他對他對.
社會完全的失望.
社會也排擠了他.
他成為一個非常極度哀的一個人.
但是當耶穌遇到了他.
事實上耶穌刻意的.
detour.
繞路去找他.

$^{561}$耶穌.
不但.
認識他.
去尋找他.
知道他的過去.
但是耶穌卻接納他.
並且向他顯明了.
天國的奧秘.
而這位緊編的婦人.
因為.
這位.
尼賽亞最榮耀的君王.
竟然.
不經意的時刻來到他的身邊.
不但.
認識他知道他過去所做的一切的.
傷心懊悔的事.
不但不定罪他.
還接納他.
陪伴他.
並且向天國的奧秘告訴他.
這是.
父神的愛.
神樂意與我們建立何等親密的關係.
就好像耶穌基督去找撒該一樣.
撒該.
是一個矮矮.
應該有胖吧.
聖經我不知道有沒有講胖.
反正是矮就是了.
一個矮矮的瑞麗.
他是賣國賊.
大家都很痛恨他.
我相信他不是因為.
從小.
寫作文就是我志向是做賣國賊.
肯定不是的.
他一定是被bully.
被霸凌被欺負被排擠.
到後來有許多的怨恨.

$^{601}$最後成為一個刺蝟般的人.
但是他內心深處他何等渴望被人接納.
沒有想到.
被大家所.
稱為賣國賊被排擠的人.
耶穌基督.
來到世上.
走在街道上.
往他所.
隱藏的.
這個桑樹走過去說撒該.
快下來.
今天我要住到.
你家裡.
弟兄姊妹這位.
偉大的神.
他是.
佩德我們所.
敬拜.
所降服.
但是他卻倒成了肉身.
如此的愛我們接納我們.
我們也看見彼得.
曾經異性風發覺得自己是最厲害的.
領袖.
可是沒有想到他三次.
不認主.
最講.
義氣的人.
竟然背叛了他的主人三次.
背叛他的大哥三次.
這樣講更貼切.
但是當耶穌基督復活之後.
耶穌不但沒有.
摒棄他.
還幫他做了世界上第一個.
這個.
趙少瑜.
你是魚和麵包嗎.
而且呢.

$^{641}$耶穌基督.
完全的接納挽回了他.
弟兄姊妹.
當這樣子的一個.
君王.
的關係.
主人與僕人關係父親與兒女關係.
這樣所帶出來的一個敬拜.
不只是一個全然的降服完全的.
海上.
不只是一個知道你會.
供應我需要的這樣的一個.
transactional的關係.
更是一個.
蠻有溫度.
蠻有恩典的一個關係.
今天我們.
的敬拜讚美.
是否是.
建立在一個敬畏神的根基上.
今天我們的敬拜讚美是否是建立在.
耶穌基督的接納與饒恕.
這樣的一個愛的根基之上呢.
還是我們仍然覺得是.
我要在乎我的事情.
我比較重要.
我比別人優秀.
我想要做的事情.
我相信我們當中有許多的.
基督徒.
我們的.
主內的弟兄姊妹.
我們曾經很單純.
火熱愛主.
我們真的拋下所有就是跟隨神.
可是曾幾何時我們開始比較.
曾幾何時我們開始斤斤計較.
曾幾何時我們失去了.
敬拜讚美當中或是服事.
甚至跟隨神那樣一個單純的心.

$^{681}$我們當中有許多的弟兄姊妹.
你曾經覺得.
有主.
一切.
都不怕了.
有主.
我就.
不怕風浪了.
可是現在.
你.
你門心自問你心中仍然充滿了恐懼.
可能比以前更加的.
黑暗.
說穿了.
你跟你旁邊不信主的同事沒什麼兩樣.
為什麼.
是否因為我們.
沒有.
已經.
不再是曾經全然的將自己擺上.
單單的為主而活.
先求神的國和神的義.
在乎主的心意呢.
今天我們.
需要能夠.
成為那父神所尋找.
真正的敬拜者.
我們當中也許有些弟兄姊妹.
你現在正在.
一個不討神喜悅的生活的一個.
狀態.
你也知道.
但是你不知道.
該怎麼樣來走出來.
或是你很懊悔.
覺得主.
對你失望.
今天.
主在這邊.
很清楚的.

$^{721}$透過聖經.
的教導讓我們知道主.
他樂意接納你.
他不定你的罪他不責備你.
他要來.
擁抱你.
他要來扶持你.
他的恩典是夠你用的.
因為耶穌來不是.
為健康的人.
而是為有病的人.
是為有需要的人.
弟兄姊妹求主讓我們心裡再次柔軟下來.
好讓我們可以來接受主的.
包容.
他的接納.
我在大學的時候有一個經歷.
讓我清楚的知道接納的力量.
那時候我快要畢業了.
在期末考的前幾天.
我開的車.
那個要換機油的燈已經亮了.
可是我一直在念書.
我會念書對我一直在念書.
可能不念書就沒有這些問題了.
總之我是在念書.
那我想說考完期末考我再去換機油吧.
沒想到在考完期末考前兩天.
我開在highway.
這個680上面.
突然間蹦了一聲.
這個車的引擎就炸開了.
因為完全沒有機油了.
我說天吶怎麼辦.
那時候我還.
禱告說我該怎麼辦.
我還真的然後去.
還真有感動要從這個地方下.
然後下面右轉還真的有個修車廠.
OK這不是重點.

$^{761}$重點是我車子又.
車子壞了.
而且換引擎是.
是個大錢.
我心裡非常的.
恐慌因為.
我覺得我我我闖禍了.
但是我必須要面對.
我要跟我在.
當時在臺灣的父母親講這件事情.
我打電話去的時候.
我已經期待.
電話的對方是非常.
大聲的這個責備.
可是沒有想到當我.
認罪悔改的時候.
對方.
我父親的回應竟然是說.
唉.
本來我們想存些錢.
讓你畢業的時候幫你買一輛好車.
但是現在因為你車子壞了.
我們只好用這錢先幫你修舊車了.
各位姐妹你可以想像.
如果當時我的父親是.
真的就是他也有權利可以來謾罵.
但是他這樣的一個接納.
哇我跟我父親關係就更近了.
這又讓我想到.
雖然我在生命和服飾已經.
已經許多年.
我回想.
我在年輕.
作為年輕傳道的時候.
年輕氣盛.
不可一世覺得自己還算.
蠻有恩賜當時又有頭髮可以撥.
覺得說.
自己很了不起.
看別人都覺得他們.

$^{801}$不夠好.
唉現在回頭看.
真的是非常的不成熟.
很膚淺.
看事情都不全面.
可是在這樣的一個覺得很羞.
羞愧啊.
覺得自己很.
很糟糕的時刻又讓我想到.
我又看見.
我們教會的幕長始終如一的.
是用笑臉在看著我.
哇原來他們.
都一直在.
接納一直在包容著我.
因為這樣子的接納和包容.
讓我今天.
更加的.
覺得.
這個.
是我要好好來繼續來服飾.
來回應神.
的一個.
一個回報的這樣的一個心態.
當然現在每天.
每次又看到幕長用微笑的臉看著我.
我真的已經不知道他們在想什麼.
弟兄姊妹接納.
是讓我們心可以柔軟.
接納使我們可以更親近.
我們的神.
我們的主完全的接納我們.
所以弟兄姊妹.
敬拜.
就是一個.
全然的擺上.
說主啊我所有都是屬於你的.
敬拜.
就是對主有一個信任.
主啊你會.

$^{841}$照顧我的一切.
敬拜.
就是.
在跪拜當中.
在祥福當中順服當中.
讓接納.
讓包容.
讓主的愛.
成為你親近祂的動力.
和回報祂.
回應祂的.
一個力量.
當一個人成為真正的敬拜者.
我們是可以.
真的得享平安的.
因為我們知道.
這位.
不能震動的國的君王.
祂要替你負責.
當我們成為真實的敬拜者的時候.
我們的心.
是不用懼怕的.
因為我們知道.
我們這位敬拜的神.
祂全然愛我們.
完全的接納.
而愛裡.
沒有懼怕.
今天我們需要.
成為真實的.
敬拜者.
我們需要.
成為真實的.
敬拜者.
腳踏兩條船.
你跪拜兩個君王.
是沒有用的.
只會讓你更加的迷糊.
更加的疲憊.
但是今天父神.

$^{881}$祂在尋找.
真實.
敬拜祂的人.
我就發現.
原來真實的敬拜.
不但是一個敬畏.
一個順服.
更是愛的回應.
所以呢.
我今天早上.
又再一次在思想這個信息的時候.
突然感覺上帝點歌.
我才發現原來有一首情歌.
也是一首詩歌.
是一個敬拜很深的敬拜的詩歌.
在這邊唱給你們聽.
等一下回應時候我們也可以一起來唱.
這是我們黃國倫弟兄寫的.
我願意.
這個副歌副歌.
聽聽看.
願意為你.
我願意為你.
忘記我姓名.
就算多一秒.
停留在你懷裡.
失去世界也不.
可惜.
願意為你.
我願意為你.
被放逐天際.
只要你真心.
拿愛與我回應.
我什麼都願意.
什麼都願意.
為你.
我想我們可以請敬拜團到台上來.
我們可以預備我們的回應.
所以今天弟兄姊妹.
就好像羅馬書十二章一節保羅說的.

$^{921}$為神的慈悲勸你們.
將身體獻上.
當作活祭.
是聖潔的.
是神所喜悅的.
你們如此事奉.
乃是理所當然的.
原來我們的神不是一位很嚴格的教練.
好像強迫我們一定要考一百分.
更不是一個暴君.
要我們成為他的奴隸.
No.
因為我們被造的目的.
就是敬拜他.
就是事奉他.
而當我們在一個對的一個態度的時候.
我們才能夠得享.
他一切的保護.
才能夠得享.
他在你生命中.
一切的供應和計劃.
而當我們先求神的國.
當我們先求神的義.
這一切.
都要嫁給你.
讓我成為真正的.
真實的敬拜者.
才能如樣.
而在這位愛我們的主人面前.
我們可以享受他的完全的接納.
可以放下自我的防衛.
武裝.
面具.
就算我們現在仍然在罪惡之中.
他仍然接納我們.
愛裡沒有懼怕.
因此我們可以全然的回應.
主啊我在這裡.
仆人敬聽.
我在這裡.

$^{961}$請拆遷我.
我們的一生.
就是個敬拜的一生.
我們一起來禱告.
天父謝謝你.
你在我們當中.
你說你在尋找真正敬拜的人.
其實也許你的意思是.
你希望我們都能夠成為真正敬拜的人.
所以你才能夠.
將你的愛.
你的衣襟.
你的磅壁.
賜給我.
這才是真正的敬拜.
請你的磅壁賜給我.
主啊今天幫助我們.
能夠在安全感的裡面.
能夠將我們自己再次獻上給你.
好讓我們的事奉不是只是每週一次.
我們敬拜不是只是.
在日曆上面的那幾個時間點.
而我們能夠成為一個敬拜的器皿.
事奉者.
討你的喜悅.
讓我們這一生就成功成為摸著你心的人.
照你執意而行的人.
而主啊謝謝你.
我們可以安心.
因為這一切你都會take care.
奉耶穌的名禱告.
Amen.
\newpage



\section{馬太福音 18:1-20}
\label{sec:032031}
\textbf{一個都不能少}~~~~駱瑞麟傳道
\newline
\newline
連結: \href{http://www.church.com.hk/acms/content.asp?site=cdc\&op=show\&type=product\&code=032031\&layout=sermon}{\texttt{032031}} ~~~~ 語音日期: 2026-04-26
\newline
\newline
\hyperref[sec:032030]{\small{< < < PREV SERMON < < <}}
~
\hyperref[sec:index_chronic]{\small{[返順時目]}}
~
\hyperref[sec:032032]{\small{> > > NEXT SERMON > > >}}
\newline
\newline
馬太福音 18:1-20
\newline
\begin{longtable}{cl}
\hline
\hline
章節 & 經文 (和合本修訂版)\\
\hline
18:1 & \begin{tabularx}{0.7\textwidth}{X} 當時，門徒前來問耶穌：「天國裡誰是最大的？」 \end{tabularx} \\ \\ \relax
18:2 & \begin{tabularx}{0.7\textwidth}{X} 耶穌叫一個小孩子來，讓他站在他們當中， \end{tabularx} \\ \\ \relax
18:3 & \begin{tabularx}{0.7\textwidth}{X} 說：「我實在告訴你們，你們若不回轉，變成像小孩子一樣，絕不能進天國。 \end{tabularx} \\ \\ \relax
18:4 & \begin{tabularx}{0.7\textwidth}{X} 所以，凡自己謙卑像這小孩子的，他在天國裡就是最大的。 \end{tabularx} \\ \\ \relax
18:5 & \begin{tabularx}{0.7\textwidth}{X} 凡為我的名接納一個像這小孩子的，就是接納我。」 \end{tabularx} \\ \\ \relax
18:6 & \begin{tabularx}{0.7\textwidth}{X} 「凡使這些信我的小子中的一個跌倒的，倒不如把大磨石拴在這人的頸項上，沉在深海裡。 \end{tabularx} \\ \\ \relax
18:7 & \begin{tabularx}{0.7\textwidth}{X} 這世界有禍了，因為它使人跌倒；絆倒人的事是免不了的，但那絆倒人的有禍了！ \end{tabularx} \\ \\ \relax
18:8 & \begin{tabularx}{0.7\textwidth}{X} 如果你一隻手或是一隻腳使你跌倒，就把它砍下來扔掉。你缺一隻手或是一隻腳進入永生，比有兩手兩腳被扔進永火裡還好。 \end{tabularx} \\ \\ \relax
18:9 & \begin{tabularx}{0.7\textwidth}{X} 如果你一隻眼使你跌倒，就把它挖出來扔掉。你只有一隻眼進入永生，比有兩隻眼被扔進地獄的火裡還好。」 \end{tabularx} \\ \\ \relax
18:10 & \begin{tabularx}{0.7\textwidth}{X} 「你們要小心，不可輕看這些小子中的一個；我告訴你們，他們的天使在天上，常見我天父的面。 \end{tabularx} \\ \\ \relax
18:11 & \begin{tabularx}{0.7\textwidth}{X} 人子來，為要拯救失喪的人。 \end{tabularx} \\ \\ \relax
18:12 & \begin{tabularx}{0.7\textwidth}{X} 「一個人若有一百隻羊，其中一隻走迷了路，你們的意見如何？他豈不留下這九十九隻在山上，去找那隻迷路的羊嗎？ \end{tabularx} \\ \\ \relax
18:13 & \begin{tabularx}{0.7\textwidth}{X} 若是找到了，我實在告訴你們，他為這一隻羊歡喜，比為那沒有迷路的九十九隻歡喜還大呢！ \end{tabularx} \\ \\ \relax
18:14 & \begin{tabularx}{0.7\textwidth}{X} 你們在天上的父也是這樣，不願意失去這些小子中的一個。」 \end{tabularx} \\ \\ \relax
18:15 & \begin{tabularx}{0.7\textwidth}{X} 「若是你的弟兄得罪你，你要去，趁著只有他和你在一起的時候，指出他的錯來。他若聽你，你就贏得了你的弟兄； \end{tabularx} \\ \\ \relax
18:16 & \begin{tabularx}{0.7\textwidth}{X} 他若不聽，你就另外帶一個或兩個人同去，因為『任何指控都要憑兩個或三個證人的口述才能成立』。 \end{tabularx} \\ \\ \relax
18:17 & \begin{tabularx}{0.7\textwidth}{X} 他若是不聽他們，就去告訴教會；若是不聽教會，就把他看作外邦人和稅吏。 \end{tabularx} \\ \\ \relax
18:18 & \begin{tabularx}{0.7\textwidth}{X} 「我實在告訴你們，凡你們在地上所捆綁的，在天上也要捆綁；凡你們在地上所釋放的，在天上也要釋放。 \end{tabularx} \\ \\ \relax
18:19 & \begin{tabularx}{0.7\textwidth}{X} 我又實在告訴你們，若是你們中間有兩個人在地上同心合意地求甚麼事，我在天上的父必為他們成全。 \end{tabularx} \\ \\ \relax
18:20 & \begin{tabularx}{0.7\textwidth}{X} 因為，哪裡有兩三個人奉我的名聚會，哪裡就有我在他們中間。」 \end{tabularx} \\ \\
[1ex]
\hline
\hline
\end{longtable}
$^{1}$主席.
陳鑑林議員.
主席.
梁國雄議員.
主席.
馬逢國議員.
主席.
陳鑑林議員.
主席.
馬逢國議員.
主席.
馬逢國議員.
主席.
馬逢國議員.
主席.
陳鑑林議員.
主席.
馬逢國議員.
主席.
馬逢國議員.
主席.
陳志全議員.
主席.
陳志全議員.
主席.
馬逢國議員.
主席.
馬逢國議員.
主席.
陳志全議員.
主席.
陳志全議員.
主席.
馬逢國議員.
主席.
陳志全議員.
主席.
陳志全議員.
馬逢國議員.
陳志全議員.

$^{41}$馬逢國議員.
陳志全議員.
主席.
陳志全議員.
馬逢國議員.
陳志全議員.
馬逢國議員.
陳志全議員.
馬逢國議員.
陳志全議員.
主席.
\newpage



\section{出埃及記 17:6}
\label{sec:032032}
\textbf{聖經人聖經事:我必站在你面前}~~~~黃克勤牧師博士
\newline
\newline
連結: \href{http://www.church.com.hk/acms/content.asp?site=cdc\&op=show\&type=product\&code=032032\&layout=sermon}{\texttt{032032}} ~~~~ 語音日期: 2026-04-19
\newline
\newline
\hyperref[sec:032031]{\small{< < < PREV SERMON < < <}}
~
\hyperref[sec:index_chronic]{\small{[返順時目]}}
~
\hyperref[sec:032033]{\small{> > > NEXT SERMON > > >}}
\newline
\newline
出埃及記 17:6
\newline
\begin{longtable}{cl}
\hline
\hline
章節 & 經文 (和合本修訂版)\\
\hline
17:6 & \begin{tabularx}{0.7\textwidth}{X} 看哪，我要在何烈的磐石那裡，站在你面前。你要擊打磐石，水就會從磐石流出來，給百姓喝。」摩西就在以色列的長老眼前這樣做了。 \end{tabularx} \\ \\
[1ex]
\hline
\hline
\end{longtable}
$^{1}$粉絲姐妹又來讀聖經啦.
生命樹和大家一起看聖經人聖經時.
今次是「我必站在你面前」.
這個「我必站在你面前」是誰說的呢?.
是主耶和華說的.
不用怕.
我會站在你面前.
即是說我會在場在此.
哎呀 糟糕了.
那個人還沒回來.
那怎麼辦?.
哎呀 講完還沒來 聚會怎麼開呢?.
那個人還沒來 說不出話來.
很惆悵徬徨的時候.
那個說得是有力量的人說.
「我必站在你面前」.
「我到時會在場」.
嘩 定了.
看看經文.
初級十七章.
以色列全會眾都遵守耶和華的吩咐.
就按著站口從信的曠野往前行.
在利菲定安營百姓沒有水渴.
一章章說.
以色列全會眾就按著神的吩咐.
怎樣神的吩咐呢?.
好 見到雲柱動.
大家就動了.
雲柱一條雲直落到地下.
雲柱的地下收起.
在前面一塊雲.
即是一塊雲.
向前走 大家就跟著它.
一見到雲柱收起.
大家收拾東西.
收拾帳篷.
我相信收拾帳篷.
綁好 放在身上.
揹起來.
都一段時間.

$^{41}$雲柱可能會等一等大家.
動了 大家就起行.
有時夜晚就火柱收起.
顯示要動了.
大家又收拾東西.
嘩 這麼晚睡.
夜晚不要緊.
因為夜晚行涼爽兼冇咁熱.
好 不理它 入口夜晚.
以色列全會眾跟住耶和華的吩咐.
就按著「站口」.
由他們去到「信」的曠野.
信是一個曠野的名.
那裡就向前行.
這個「站口」這個字.
我特別要解釋一下.
什麼叫「站口」呢?.
這個是比較舊的中文.
「站口」Messah.
Messah.
不是Mekah嗎?.
希伯來文如果你見到它譯音出來.
英文C是讀S的.
所以Messah.
Messah重音在後面.
Messah.
就不是Messah.
Messah.
Messah是營地 駐紮地.
由他們紮了營.
大家都在旁邊.
在生活中住一段時間才再起行.
你天天走真的太疲累了.
你那個行軍什麼行軍都頂不住.
都要休息一下.
他們去一個地方停一下.
又走一下 停一下 又走一下.
雲柱火柱就帶著他們.
耶和華又會跟摩西說.
叫他們準備了.

$^{81}$起行了.
神就吩咐他們.
就在他們所駐紮的營地那裡.
就預備起行了.
那就起行了.
就由「信」的曠野住了一段時間.
就向前走了.
在那裡就說.
沒東西吃了.
好 先給你一些鵪鶉吃.
走了一段時間.
就去到另一個新的地方.
其實是比較向南一點走的.
就叫利菲定.
這個利菲定就叫Refidim.
Refidim.
Refidim就是休息的意思.
行完曠野去到一個地方.
休息.
利菲定.
Refidim.
Refidim的意思.
休息.
其實很有可能就是.
有一點水.
一個Oasis.
一個綠洲.
有些水.
或者有些樹木.
在那裡住.
休息一下才行.
起碼有樹蔭才行.
當然沒有自己的帳篷.
可以休息.
不過辛苦.
去到那裡.
就很多人會在利菲定那裡停留.
因為是一個休息的地方.
所以叫Refidim.
已經叫休息了.

$^{121}$去到那裡.
又吵.
哇!吵什麼呀?.
百姓們有水害.
沒有水呀.
水庫都沒有了.
所以就向摩西爭鬧.
這些爭鬧.
Refidim.
是爭論.
跟他爭論.
為什麼你帶我來這裡.
爭論什麼呢?.
給我一些水來喝.
又說什麼呢?.
為什麼你跟我爭鬧.
為什麼你跟斯坦爺說呢?.
其實就是說.
沒有水.
沒有水就大家一起禱告.
求神給我們.
神行這麼大的神蹟.
十災帶我們出來.
然後過紅海.
然後埃及軍兵死在水裡面.
這樣還不夠?.
接著你說沒有肉吃.
接著大家有鵪鶉吃.
有鰻.
有嗎哪吃.
很大的神蹟神一路帶著我們.
沒有水喝.
哇 又來.
根本不想開罷.
一直爭論.
爭論沒有水喝很慘.
你們跟我爭論都沒有用.
埋怨神都沒有用.
因為這樣.
你是斯坦爺和華.

$^{161}$C讀S.
NASA.
為什麼NASA爺和華.
斯坦爺和華.
斯坦爺和華即是測驗試驗.
行不行.
買這東西回來.
先測測它.
試試它行不行.
即是說.
是沒有信心.
所以要測.
看你行不行.
為什麼不要斯坦爺和華.
百姓在那裡口口開口喝水.
向摩西發怨言.
其實是表達這件事.
為什麼.
又很刻薄.
你為什麼將我們從埃及領出來.
使我們我們的兒女並生畜.
都渴死在這裡呢.
你帶我們出來是否想渴死我們.
現在是否想渴死我們.
嘩這麼難聽.
不如我逼你出來.
摩西呼求耶和華.
神啊.
百姓這樣.
我應該怎樣做呢.
他們剛才生氣得差點用石頭咪我.
要扔我扔死我.
說我帶他們出來要渴死他們.
我應該怎樣做呢.
神你教教我吧.
耶和華對摩西說.
你手裡拿著你先前擊打河水的杖.
河水呢.
Ye'or.
這個Ye'or.

$^{201}$是河大家要留意.
河或溪.
即是說.
以色列人過海.
但其實這裡是說過河.
很闊的河.
那些很闊的水.
溪應該不怕.
是闊的河是一條河.
而未必是海.
早幾集我們見到以色列人過紅海.
我們一直有這個概念.
但其實很有可能他們不是過今天的紅海.
而是過這個海.
這個海可能叫做蘆葦海.
是一個海口的地方.
比較闊.
不過有一個彎彎的好像一個海.
是很多蘆葦的.
是鹹淡水交界的地方.
以至可能有很多蘆葦在那裡.
這樣也不一定.
全鹹水是沒有蘆葦的.
淡水才有蘆葦的.
而那裡很有可能有養分.
所以蘆葦長得很旺盛.
叫做蘆葦海.
不會太深.
不過又不是很淺.
擊打河水那支仗.
你記得帶著以色列幾個長老.
有幾個不要緊.
以色列長老即是十二支派裡面有些代表人物.
從百姓面前走過去.
即是說.
哎呀 渴死我們了.
你是不是想渴死我.
現在我們圍你.
摩西就說等等.
我去祈禱先.

$^{241}$神說怎麼辦.
他們幾乎扔死我.
石磨過來.
幸好我閃了.
很危險.
現在怎麼辦.
沒有水喝.
我都一籌莫展.
我都拿不出水給他們喝.
神就說不要緊.
你拿著你擊打河水那支仗.
然後埃及人全部浸死那支仗.
很重要的.
神的仗.
你就帶著它.
神的憑證在那裡曾經掉落地變成蛇.
就帶著這支仗.
還有找幾個百姓支派的代表.
在他們面前走.
即是說展示給他們看.
展示時間來了.
表演的時間給你看看.
我就會在那裡呢.
河裂的盤石那裡站在你面前.
我會在你去哪裡呢.
其實他們離非定之後.
就很近河裂山了.
河裂山是比較在西乃半島沙漠地帶裡面.
是比較有少少山向西南一點.
東南一點是有少少山.
河裂山在那裡.
去到那裡河裂的盤石.
去到那裡.
然後你要擊打盤石.
盤石就有水流出來.
百姓就可以看.
摩西就在以色列的長老面前.
就這樣走了.
擊打.
然後大家看到.

$^{281}$嘩.
盤石真的出水.
這個故事很聞名.
所以我特別要講一下這個故事.
讀這段經文這麼長都讀.
不過我想你留意一件事.
河裂呢.
河裂這個字呢.
Korib.
Korib是沙漠.
河裂是沙漠.
河裂山就是沙漠的山.
盤石這個字很難發音.
Chur.
盤石這個TS讀.
Chur.
就是盤石的意思.
盤石呢.
你不要想著先找盤石.
頭一向下看就看不到盤石了.
這麼大舊種不是盤石.
不是.
看盤石的頭要這樣向上.
因為盤石的意思.
是峭壁.
或者盤石.
是一個山很高很直.
如果大家有機會去以色列.
就去到巴尼亞斯.
這個Pen.
即是耶穌時代的該塞尼亞菲勒比.
那裡有很多人拜不同的神.
如果今天去.
就會去到一個很大的盤石.
一個山.
去到那裡大家找一找.
指給我看哪一塊是盤石.
去到那裡.
這塊是不是盤石.
很大.

$^{321}$三四個人都抱不緊.
那叫做一塊石仔.
盤石是要.
好像向山舉木一樣.
整個峭壁在後面.
那個才是一塊盤石.
盤石是這個意思.
去到那裡.
你用那支杖擊打一個盤石.
好像你去擊打一個峭壁一樣.
當然你不是在峭壁頂.
是去到底.
上面是一個峭壁.
你站在峭壁的底.
然後你擊打這塊石.
哄哄聲有水出.
二百萬人在那裡拿水喝.
所以那個山很大.
一擊打下去.
走虎為出水.
很厲害.
所以長老在眼前這樣走.
不要說你不知道.
這個是純粹耶和華.
我們神的能力.
本來大家乾到渴死.
那幾個長老都很索氣.
不過一擊打整個盤石.
走虎為出水.
幾百萬人,二百萬人都喝得到.
於是長老.
你們的長老都看到的.
武士很厲害的.
第二點就是.
我必站在你面前.
這個是今天的主題.
當你有什麼困擾困惑的時候.
你記得.
神是在場.
我必站在你面前.

$^{361}$所以你不會空打.
一打就彈回那支杖上來.
差點打到自己.
不會.
你一擊打.
我就會叫他出水.
嘩 怎麼這麼厲害.
泥土就說擊打容易藏水.
盤石石怎麼會有水.
所以充分是一個神蹟.
打 很多水出來.
這個呢.
我必站在你面前.
即是說.
我會在場撐你.
我會支持你.
我會做事.
放心.
他們把那個地方.
這塊盤石的地方叫馬殺.
又是S.
馬殺.
馬殺就是誘惑.
誘惑.
或者叫米利巴.
米利巴.
不是 叫米爾巴.
米爾巴.
米利巴就是糾紛.
你就將那個地方叫誘惑和糾紛.
那裡有時你缺乏.
它是一個誘惑.
令你很想要做些事.
或者在那裡引起糾紛.
因為以色列人爭鬥.
好了.
誘因他們私貪.
誘惑說.
誘惑是不是在我們中間.
問這些都有.

$^{401}$這一句千萬不要問.
這一句是私貪誘惑的說話.
誘惑是不是真的在我們中間.
教會問.
是不是真的有神在.
是不是真的有耶穌在.
這樣問的時候.
是一個很大的侮辱與試探.
千萬不要.
有報應的.
神很不高興.
請你記得.
不要試探我們的主.
而且你要知道.
我必站在中間.
我會站在那裡.
我必站在面前.
我會參與在當中.
神什麼時候都站在那裡看著.
放心.
一起禱告.
主求你幫助我們.
真的很需要你站在我們當中.
和我們同在.
和我們一起去經歷.
和我們去做出你的支持.
你在場的時候.
我們的心是很平安.
因為知道你在場.
幫助我們對你更多的信靠.
而不會去試探你.
你是不是真的在這裡.
你是不是在教會當中.
這個世界這樣你去了哪裡.
主我們不敢問這個問題.
我們信靠你.
你是在場的.
你是支持.
你是看到的.
你知道的.

$^{441}$你的榮耀要在這裡彰顯.
暫利主.
祈禱奉耶穌名禱.
阿們.
\newpage



\section{以西結書 47:1-12}
\label{sec:032033}
\textbf{活水江河}~~~~郭張玉清姊妹
\newline
\newline
連結: \href{http://www.church.com.hk/acms/content.asp?site=cdc\&op=show\&type=product\&code=032033\&layout=sermon}{\texttt{032033}} ~~~~ 語音日期: 2026-04-26
\newline
\newline
\hyperref[sec:032032]{\small{< < < PREV SERMON < < <}}
~
\hyperref[sec:index_chronic]{\small{[返順時目]}}
~
\hyperref[sec:032034]{\small{> > > NEXT SERMON > > >}}
\newline
\newline
以西結書 47:1-12
\newline
\begin{longtable}{cl}
\hline
\hline
章節 & 經文 (和合本修訂版)\\
\hline
47:1 & \begin{tabularx}{0.7\textwidth}{X} 他帶我回到殿門，看哪，有水從殿的門檻下面往東流出，因為這殿是朝東的。水從殿的側面，就是右邊，從祭壇的南邊往下流。 \end{tabularx} \\ \\ \relax
47:2 & \begin{tabularx}{0.7\textwidth}{X} 他帶我出北門，又領我從外邊轉到朝東的外門，看哪，水從右邊流出。 \end{tabularx} \\ \\ \relax
47:3 & \begin{tabularx}{0.7\textwidth}{X} 他手拿繩子往東出去，量了一千肘，使我涉水而過，水到腳踝。 \end{tabularx} \\ \\ \relax
47:4 & \begin{tabularx}{0.7\textwidth}{X} 他又量了一千，使我涉水而過，水就到膝；再量了一千，使我過去，水就到腰； \end{tabularx} \\ \\ \relax
47:5 & \begin{tabularx}{0.7\textwidth}{X} 又量了一千，水已成河，無法過去；因為水勢高漲成河，只能游泳，無法走過。 \end{tabularx} \\ \\ \relax
47:6 & \begin{tabularx}{0.7\textwidth}{X} 他對我說：「人子啊，你看見了嗎？」他帶我回到河邊。 \end{tabularx} \\ \\ \relax
47:7 & \begin{tabularx}{0.7\textwidth}{X} 我回到河邊時，看哪，河這邊與那邊的岸上有極多的樹木。 \end{tabularx} \\ \\ \relax
47:8 & \begin{tabularx}{0.7\textwidth}{X} 他對我說：「這水往東方流，下到亞拉巴，直到海。所流出來的水，一入海就使水變淡。 \end{tabularx} \\ \\ \relax
47:9 & \begin{tabularx}{0.7\textwidth}{X} 這兩條河所到之處，凡滋生的動物都必存活；這水流到那裡，使那裡的水變淡，因此裡面有極多的魚。這河水所到之處，百物都必存活。 \end{tabularx} \\ \\ \relax
47:10 & \begin{tabularx}{0.7\textwidth}{X} 必有漁夫站在河邊，從隱‧基底直到隱‧以革蓮，全都成了曬網的場所。那裡的魚各從其類，好像大海的魚甚多。 \end{tabularx} \\ \\ \relax
47:11 & \begin{tabularx}{0.7\textwidth}{X} 但是沼澤與池塘的水無法變淡，只能作產鹽之用。 \end{tabularx} \\ \\ \relax
47:12 & \begin{tabularx}{0.7\textwidth}{X} 河這邊與那邊的岸上必生長各類樹木，可作食物；葉子不枯乾，果子不斷絕。每月必結新果子，因為這水是從聖所流出來的。樹上的果子必作食物，葉子可以治病。」 \end{tabularx} \\ \\
[1ex]
\hline
\hline
\end{longtable}
$^{1}$各位弟兄姊妹早晨.
外面好陽光.
不如我們將陽光帶進來我們當中.
我們用一個很輕快.
很有陽光.
很健康的聲音.
和你旁邊的那位朋友.
說一聲早晨.
祝福他好嗎.
早晨.
其實早上跟別人說一聲早晨.
真是很醒神的 是不是.
是一種祝福.
記得幾十年前的香港.
即是當我還是很小的時候.
一家人去旅行.
不是全家人搭飛機.
去日本,去泰國,去韓國,去新加坡.
遠一點就去歐洲,去美加.
我們去哪裡呢.
我們去九龍公園.
遠一點就去動植物公園.
有人說為甚麼不是摩士公園.
我們那時候都不知道摩士公園是甚麼.
只記得有一個九龍公園.
和有一個香港的動植物公園.
即是俗稱甚麼.
我們的長輩知道了.
就是冰頭花園.
冰頭花園可能我們年輕的小朋友.
都未必知道這個名字的由來.
旁邊就是總督府.
總督是統領三軍.
所以它是冰頭.
所以叫冰頭花園.
我記得媽媽那時候帶著麵包,罐頭.
還拿著一張很大很大的草脊.
一家人就去去咸咸地.
搭車又搭船.
好像經過千山萬水才可以去得到.

$^{41}$冰頭花園那時候有一個地標.
可能你都會記得.
是一個很大的噴水池.
很多家庭例如放在那裡.
拍一張家庭大合照.
如果你翻箱倒籠.
找媽媽或者找婆婆.
找奶奶那些珍藏舊照.
可能你都會找到一張.
有這個大噴水池的相片.
我們和很多小朋友一樣.
都會在那裡噴水池旁邊玩水,潑水.
媽媽還會給幾個毫子.
不是一元銀,一元已經很厲害了.
封利是收到一元已經很開心.
我們給幾個毫子.
就扔進水池.
她說扔下去吧.
會快高長大.
還有讀書會好一點.
幾十年後的今天.
好像一來一點都沒有變.
你依然會看到一般人去旅行.
去到外地也好,去到國內也好.
去到一些景點.
只會有水流出來的地方.
不要管它是羅馬噴泉.
還是甚麼噴泉.
或者有隻石龜.
有一條鯪魚的水池.
你都會在那裡拍一張漂亮的相片.
那些水你都會用來洗手.
或者甚至乎喝一口.
或者你都會拋個銀下去.
許個願.
因為一般人都覺得.
這些流出來的水.
是帶有生命力的.
可以為自己.
或者你身邊所愛的人.

$^{81}$帶來祝福.
所以你會見到.
你裝下池底.
會有很多銀在那裡.
今日跟大家分享的經文.
《以西結書》47章1至12節.
請大家翻開唐容的《聖經》舊約.
1307頁.
《以西結書》47章1至12節.
是描述以西結先知.
看到有聖殿湧流出來的泉水.
可以令到四周萬物生長的一個意象.
這個意象都和流出來的水有關係.
《聖經》裡面所講的意象是甚麼呢.
幻覺?.
發夢?.
不是.
《聖經》所講的意象.
英文是vision.
是指神主動向人顯現的.
一些超自然的影像或者是啟示.
是發生在人清醒的狀態.
而不是睡覺發夢的時候.
而意象的目的.
就是神向人傳達他的心意.
他的計劃.
預言.
或者是給人一些指引.
在進入以西結這個意象之前.
不如我們首先去看一看.
了解一下《舊約聖經》.
以西結書他成書的一些背景.
以西結書的背景.
是由大亡國.
亡國之前.
至到被擄的初期.
大約是公元前.
593年.
至到571年之間.
當時巴比倫王攻陷耶路撒冷.

$^{121}$他擄走了約瓦根王.
還有他們的貴族.
公將.
甚至乎是祭司.
去到巴比倫.
當然當中是包括了以西結.
以西結出生在一個祭司的家庭.
對聖殿和一些祭祀的禮儀非常熟悉.
他在被擄到巴比倫之後的第五年.
他是夢召成為先知.
每一個先知都有他的責任.
而以西結的責任.
是承擔了守望以色列家的責任.
他負責提醒和警戒百姓.
提醒他們不要行差踏錯.
從這些背景資料可以看到.
以西結他寫書的時候.
以色列人正處於一個很黑暗困苦的時期.
當時耶路撒冷的聖殿被毀.
他們的家園也是荒廢.
但神在這個異象裡面.
是帶領以西結才看到一座新的聖殿.
他告訴我們.
神的恩典不會留在一個建築物裡面.
是滿有生命力.
好像泉水一樣從聖殿湧流出來.
目的是甚麼呢.
讓他的子民在絕望裡面看到盼望.
今天講到經文的第一和第二節描述說.
他帶我回到殿門.
看啊.
有水從殿門的門檻下面往東流出.
因為這殿是朝東的.
水從殿的側門就是右邊.
從祭壇的南邊往下流.
他帶我出北門.
又領我從外邊轉到朝東的外門.
看啊.
水從右邊流出.
這裡提到.

$^{161}$水是從殿的門檻下面經過祭壇的右邊流出來.
在神學的意義裡面來看.
聖殿是代表了神的同在.
祭壇是代表了基督的救贖.
在聖壇流出來的水.
是神恩典的源頭.
耶穌基督為世人犧牲在十字架上面.
他右邊立旁所受傷傷.
即是被士兵刺傷的位置.
是流出了寶血.
洗淨了我們.
救贖了我們.
恩典不是我們賺回來的.
是神伯伯賜給我們的.
當我們領受了恩典之後.
我們的生命又有什麼改變呢?.
經文第三至第五節.
是描述了信徒成長所經歷四個階段.
以西結在意象裡面.
那個人是帶他在水裡面行走.
那個人是用繩來量了四次.
每一次是一千爪.
以當時量度的單位來說.
一爪大概是五十厘米.
即是我們舊制.
即是英制的大概是二十吋左右.
一千爪即是大概是五百公尺.
五百公尺有多長呢?.
如果以一個國際標準泳池來說.
大家都知道是五十公尺.
即是來回要由十次那麼長.
第一次量一千爪的時候.
以西結在水裡面行走.
水是去到他的腳裸.
即是去到他腳眼的位置.
他仍然可以抓緊地面.
走起來很治癒.
這個階段很讓我們初信主.
我們剛剛接觸到信仰的時候.
我們嘗到主人的滋味.

$^{201}$但是我們很容易按著自己的思想去行事.
自己掌握了自己的主權一樣.
而第二次量一千爪的時候.
以西結在水裡面行走.
水是到膝蓋的位置.
相信大家都有水浸去澗水的經驗.
當水去到膝蓋的時候.
你會感受到水的阻力.
但是我們仍然都會行得到.
而這個階段很讓我們在信仰上面去追尋.
信仰生命是盡心的一個階段.
我們會學習去研讀聖經.
會學習多些落膝蓋的功夫.
即是我們會跪下去祈禱.
跪下去祈禱.
在我們初信的時候.
我們真是很常有.
現在好像忘記了這回事.
我們面對聖經的教導.
我們會感受到它的提醒或者阻力.
但在這個時間我們都會按自己的意思來行事.
到第三次量一千爪的時候.
水是去到以西結的腰部.
你想想水去到腰部的時候.
那個水力是會強很多的.
我們是用腳來很用力去抓緊地下.
才不會被水沖跌.
如果你要行的時候.
你會加倍要發力.
又要抓緊地下.
又要向前去移動的時候.
是比較艱難的.
而這個時間.
這個階段.
就很讓我們在信仰操練的階段.
是信徒在生活裡面.
去經歷神的管教和引導的時間.
大家都看過電影.
當兩個高手過招的時候.
他們會先做一件什麼事呢?.

$^{241}$搓一搓自己的腰帶.
「嘻」一聲.
當他搓這個腰帶的時候.
他的勇氣.
他的力度.
好像突然間會得到提升.
在《耳窟所書》六章十四節這樣說.
所以要站穩了.
用真理當作大指速要.
用公義當作護心鏡遮控.
沒錯.
我們在信仰操練裡面.
真的要幫自己.
去用真理來去搓腰.
讓自己不要被試探和引誘來去激導.
而搓腰.
也都有另一種的含義.
是代表了服事.
在那個時期的工人.
他們為了方便工作.
都會用一些布去纏住自己的腰.
而在服事當中.
我們可以學習謙卑的功課.
面對別人的提醒.
批評.
甚至乎一些要求.
我們有機會學習忍耐.
包容.
甚至乎饒恕別人.
所以信徒要在服事的積分上面.
去盡心和操練自己.
當梁第四次一千肘的時候.
經文第五節是這樣記載.
「水已成河沒法過去.
因為水勢高漲成河.
只能游泳沒法走過」.
在這個時期.
腳已經不可以碰到地下.
因為水已經是高過人.
這個是代表一個成熟的基督徒.

$^{281}$他不再靠自己的力量.
或者是按自己的意思來生活.
而是全然的信服在神的帶領.
讓聖靈的江河去盛載著我們向前進.
這條江河是源自聖殿.
經過祭壇集結而成.
是滿有恩典的活水江河.
在第六節裡面.
那個人又和以西結說.
「人子啊,你看見了嗎?」.
如果用廣東話來說.
這句話就是說.
「老友,你看見了嗎?」.
是提醒我們要留意這一切.
我們信了基督.
領受了恩典之後.
應該是要離開基督道理的開端.
我們要下定決心來操練自己.
要竭力去進入一個完全的地步.
我們又再看看第七至十二節.
那個人是將以西結帶回到河邊.
河兩邊很多樹.
河是向東流.
經過阿拉巴進入海.
阿拉巴的意思是荒野.
是寸草不生的地方.
是指由加利利海.
去到阿哈巴灣之間的約旦河谷.
而海是指死海.
和本聖經所說的「嚴海」.
大家可能都看過一些旅遊節目.
去介紹死海.
甚至乎可能你都會去過死海旅遊.
死海是位於以色列和約旦之間.
那個鹹度是非常之高.
是一般海的八至十倍.
這麼高的鹹度.
讀物理的都知道.
那個密度這麼高.
浮力是會很大.

$^{321}$海裡面沒有生物.
而周圍的地方都是寸草不生的.
但是在意象裡面.
以西結是看到.
當河水流入死海的時候.
河水經過所有地方.
滋生的動物都可以得到存活.
水流到哪裡.
就令到哪裡的水變淡.
水裡面有很多魚.
多到好像海裡面的魚那麼多.
而漁夫是站在岸邊.
由隱磯底至隱爾甲連.
周圍的地方.
全部都成為了漁夫.
捕完魚之後曬網的曬網場.
隱磯底是死海裡面的一個綠洲.
而隱爾甲連有解經家說.
可能是位於死海北部的困難地區.
死海古卷就在這個地方發現.
什麼叫做死海古卷呢.
死海古卷是1947年.
有一個牧羊人在困難洞穴裡面發現.
距離現在是超過2000年.
是最早的聖經的文獻.
是用希伯來文寫在羊皮紙上面.
現在這個聖經文獻又放在哪裡呢.
就在以色列的死海古卷的博物館裡面.
希望打完仗之後.
我們有機會去到以色列旅行.
都去看看這個死海古卷.
我們看看英文聖經現代籍譯本的譯法.
在以西結先知的意象裡面.
這條經過聖殿流出來的河水.
不單止令到死海變為活海.
可以有生物.
在十二節裡面.
我們更加看到四周圍的地方.
吸收了聖殿流出來.
經過祭壇源源不絕的水.

$^{361}$植物都生長.
而且很茂盛 定期每一個月.
都可以長出果實給人食用.
成為別人的供應和祝福.
而植物的葉都可以用來醫人.
即是用自己的生命力去撫平別人的傷口.
去活化了別人的生命.
今天短短十二節的經文展示出來.
以西結先知看到的意象裡面.
我們可以總結出四個重點.
第一個重點是.
聖殿流出來的泉水是強而有力的.
是源源不絕的.
是代表了神的恩典是極之大.
不會停 是白白地給我們.
而第二個重點.
是提醒我們信徒成長.
要經歷的四個階段.
就像那個人用繩子去量四次.
第一次水是去到腳眼.
接著去到膝蓋.
接著去到腰.
接著是高過一個人.
這裡提醒我們.
我們領受了救恩之後.
我們要努力的.
逐步去進入一個完全的地步.
在信仰上面我們要紮根.
要勤讀聖經.
要多些落膝頭歌的功夫學習祈禱.
要用真理去束腰成為力量.
要參與在服事裡面.
學習謙卑 忍耐 關心.
和體諒別人.
我們更加要放下自己.
去信服聖靈的帶領.
我們要明白神的心意.
我們要活在神的旨意裡面.
第三個重點.
是讓我們明白到.

$^{401}$神是有醫治的大能.
祂可以令到死海變活.
百物可以重生.
而我們的生命.
有多枯萎.
有多枯乾.
有多腐敗.
只要我們願意重新去靠近神.
去吸取從祂而來恩典之傳.
我們枯死的生命.
必定可以重新去活化.
而第四個重點.
是一個領受了神恩典的生命.
可以成為一道活水江河.
成為別人的供應和祝福.
我小時候聽過一個故事.
我到今時今日都記得很清楚.
那個故事是講到一個媽媽.
她生了一個很懶很懶很懶的兒子.
有一次媽媽必須要出門遠行.
在她臨離開的時候.
她將一個很大很大的乾酥餅.
掛在兒子的頸上.
掛好了才出門.
誰知過了幾天回來.
兒子居然是餓死了.
難道兒子很大吃.
這麼大的乾酥餅都不夠他吃.
看清楚一點.
原來掛在他頸上的乾酥餅.
只是咬了最近他口的那一口.
就是說兒子是懶得放在他眼前.
為他預備好可以養活他的食物.
他都不去拿.
這樣就餓死了.
弟兄姊妹.
我們是不是都可以自己耐性一下.
究竟我們現時我們的生命.
我們的靈命狀態又如何呢?.
我們有沒有感到很枯乾.

$^{441}$我們的生活是不是很困難.
我們是不是內心裡軟弱無力.
面對生活裡不同的挑戰.
我們會感到困惑.
甚至乎會跌倒.
擺在我們眼前的恩典.
就好像這個大乾酥餅的故事.
我們有沒有願意去靠近著.
讓神所賜的活水存源.
來滋養我們的生命.
讓我們的生命可以重新活起來.
成為一道活水江河呢?.
最後和大家分享一幅圖畫.
這幅圖畫如果去過.
在協和街上面的基督教聯合醫院的朋友.
或者你都會覺得眼熟.
這幅畫曾經在基督教聯合醫院S座大堂.
即是邵逸夫大樓的大堂.
擺放了近30年的時間.
這幅畫是1995年.
醫院經基督教協進會.
邀請了當時的南京宗教藝術家何麒博士.
特別來到香港.
畫了一個星期.
而牆底我們找了工業學院.
幫我們做好整個牆身.
做好牆底.
讓這幅畫可以顏色和裂紋盡量保存減少.
真的足足擺放了30年.
雖然最近這幾年.
因為擴建的緣故.
這幅畫已經沒有了.
但是這幅畫裡面的意思和祝福.
它的意境是不會被遺忘.
這幅很有中國色彩的宗教圖畫裡面.
你可以看到救贖主.
祂生化萬物.
救贖和醫治的恩典.
基督耶穌舉起雙手.
在應召和祝福世人.

$^{481}$在祂身上右邊肋旁.
那個槍傷的地方.
流出來的寶血.
匯聚成為一個溪河.
這條溪河是潔淨和哺育了不同的生命.
那裡還有方舟.
伊甸園.
還有河邊飲水的鹿.
還有被餵飼的小鳥.
整幅圖畫都揚入了一個很祥和的景象.
還有一個慈愛淪捨.
無私伸手助人的好撒瑪利亞人.
象徵著對群眾的承諾.
我們願意用基督的愛.
去祝福守望我們的社區.
相信大家都會覺得是一幅很有意思.
很有意境的圖畫.
而當中表達的景象.
也跟我們今天講到的經文主題.
都是相當之有呼應.
今天是教會所定的健康主日.
我們每一個人都不可以避免.
我們都會因為年紀.
或者是疾病的緣故.
我們的肉身是會毀壞.
但是當我們經常親近主.
吸取源於基督的生命之泉.
我們的肉身雖然是會毀壞.
但是我們的內心是會一天生此一天.
會享有一個健康活潑的屬靈生命.
盼望在座每一位我們接受了救恩的朋友.
都可以活出有生命力的生命.
成為活水江河.
讓我們身邊的人.
因為你而認識主耶穌.
也都因為你而得了福氣.
我們同心祈禱.
主啊.
你在約翰福音七章提醒我們說.
人若學了.

$^{521}$可以到我這裡來學.
信我的人.
就如經上所說.
從他的福中要流出活水的江河來.
這是你給我們的應許.
求你讓我們知道.
如果我們感到生命枯乾.
生活無指望.
軟弱無力.
往往是因為我們離開了生命之泉的源頭.
即是離開了主耶穌你自己.
願我們都能夠靠近你.
願我們稽渴枯乾的生命.
跟著你恩典的滋潤.
可以重新活化.
我們有美好的見證.
去祝福別人.
自己亦都蒙受祝福.
阿們.
\newpage



\section{約翰福音 10:1-10}
\label{sec:032034}
\textbf{得更豐盛的「生命」}~~~~林津牧師
\newline
\newline
連結: \href{http://www.church.com.hk/acms/content.asp?site=cdc\&op=show\&type=product\&code=032034\&layout=sermon}{\texttt{032034}} ~~~~ 語音日期: 2026-04-26
\newline
\newline
\hyperref[sec:032033]{\small{< < < PREV SERMON < < <}}
~
\hyperref[sec:index_chronic]{\small{[返順時目]}}
~
\hyperref[sec:032035]{\small{> > > NEXT SERMON > > >}}
\newline
\newline
約翰福音 10:1-10
\newline
\begin{longtable}{cl}
\hline
\hline
章節 & 經文 (和合本修訂版)\\
\hline
10:1 & \begin{tabularx}{0.7\textwidth}{X} 「我實實在在地告訴你們，那不從門進羊圈，倒從別處爬進去的，就是賊，就是強盜。 \end{tabularx} \\ \\ \relax
10:2 & \begin{tabularx}{0.7\textwidth}{X} 那從門進去的才是羊的牧人。 \end{tabularx} \\ \\ \relax
10:3 & \begin{tabularx}{0.7\textwidth}{X} 看門的給他開門，羊也聽他的聲音。他按著名叫自己的羊，把羊領出來。 \end{tabularx} \\ \\ \relax
10:4 & \begin{tabularx}{0.7\textwidth}{X} 當他把自己的羊都放出來，就走在前面，羊也跟著他，因為牠們認得他的聲音。 \end{tabularx} \\ \\ \relax
10:5 & \begin{tabularx}{0.7\textwidth}{X} 羊絕不跟陌生人，反而會逃走，因為不認得陌生人的聲音。」 \end{tabularx} \\ \\ \relax
10:6 & \begin{tabularx}{0.7\textwidth}{X} 耶穌把這比方告訴他們，但他們不明白他所說的是甚麼。 \end{tabularx} \\ \\ \relax
10:7 & \begin{tabularx}{0.7\textwidth}{X} 所以，耶穌又對他們說：「我實實在在地告訴你們，我就是羊的門。 \end{tabularx} \\ \\ \relax
10:8 & \begin{tabularx}{0.7\textwidth}{X} 凡在我以前來的都是賊，是強盜；羊沒有聽從他們。 \end{tabularx} \\ \\ \relax
10:9 & \begin{tabularx}{0.7\textwidth}{X} 我就是門，凡從我進來的，必得安全，並且可進出，找到草吃。 \end{tabularx} \\ \\ \relax
10:10 & \begin{tabularx}{0.7\textwidth}{X} 盜賊來，無非要偷竊、殺害、毀壞；我來了，是要羊得生命，並且得的更豐盛。 \end{tabularx} \\ \\
[1ex]
\hline
\hline
\end{longtable}
$^{1}$大家好 主內平安.
很感恩今天可以在這裡跟大家一起敬拜.
也分享上主的道.
今日約翰福音的經課.
是我們大家都很熟悉的羊圈的比喻.
當中都包含了一句我們很喜歡的聖經章節.
一個應許.
就是主耶穌基督跟我們說.
我來了是要叫羊得生命.
並且得的更豐盛.
我相信沒有人不想擁有一個豐盛的生命.
特別我想在大學的環境當中.
很多人都假設.
認為我們應該有一個比別人更豐盛的生命.
為什麼呢.
首先我們在大學裡面.
我們都擁有比一般人更多的知識和智慧.
今天在這個知識型的經濟當中.
其實有知識有智慧.
其實就代表著有財富有經濟的能力.
當然我都知道現在在大學的環境裡面.
都競爭很大.
教席都是不容易的.
不過無論如何我都相信.
我們都是比一般人來說.
都是有一個相對穩定的經濟情況.
都算是OK.
起碼不會餓壞.
特別我想除了我們有知識.
有經濟的能力穩定之外.
更加重要我們是能夠得到別人的欣賞.
敬重.
其實這個都很重要.
在我們人生當中找到成就.
找到別人的喜愛.
其實這些都是豐盛生命.
很多人都以為是這樣.
但事實上.
大家覺得我們的生命是不是真的這麼豐盛呢.
如果是的話我就可以不用再講下去了.

$^{41}$因為大家都已經得到很豐盛的生命了.
不過有時候我們都會覺得有點疑惑.
好像知識,財富,地位,成就.
都似乎擁有.
不過好像還不是一個真正豐盛的生命.
為什麼呢.
我相信因為這套經文裡面.
耶穌講得很清楚.
或者說祂講得很特別.
祂說:我來了是要叫楊得生命.
並且得得更豐盛.
所以這裡是有一個先後次序的.
我們要先得到生命.
然後才繼續豐盛下去.
今天我們很多時候就是為了追求人生的豐盛.
有時候是反過來本末倒置了.
反而是失去了我們真正的生命.
因為這些屬世的豐盛.
其實和真正的生命可能是無關的.
甚至可能在我們不斷去追求成功.
不斷去追求豐盛的時候.
可能是失去了我們真正的生命.
究竟什麼才是真正的生命呢.
就讓我們再仔細去一起看.
《約翰福音》十章一至十節這一段經文.
我們首先要明白這個經文的背景.
其實就是第九章.
耶穌醫治了一個生出來是黑眼的孩子.
但是當日是在安息日當中.
所以就被很多的宗教領袖所排斥.
耶穌接著在第十章就講這個比喻.
其實就是想讓當時的門徒或者信徒知道.
其實信仰是否真的能夠帶給我們真正的生命呢.
都是有一些的巧妙之處.
因為在這個世界裡面有真的先知.
也有假的先知.
信仰其實都有真的信仰和假的信仰.
如果是虛假的信仰不能夠帶給我們真正的生命.
只有真正的信仰才能夠帶給我們真正的生命.
所以耶穌就在這段《約翰福音》的比喻當中.

$^{81}$告訴我們什麼才是真正的生命.
整段經文是很豐富的.
我只是抽了其中四個元素.
跟大家一起去思考.
首先耶穌說.
真正的牧羊人或者真正的信仰是必須從門進去的.
這道門就是耶穌基督祂自己.
我們當然也知道耶穌是好牧人.
為什麼突然間從好牧人又變成了羊的門呢?.
原來這個也是和當時的處境有關.
因為當這些牧羊人帶著羊到野外的地方.
晚上休息的時候他們沒有一些很正式的羊棚或羊圈.
只是在郊野地方有一些很簡陋的圍欄圍著羊.
甚至連門也沒有.
牧羊人就會躺在門口充當了這個羊圈的門.
所以耶穌說我是羊的門.
其實大家都明白當時的背景.
這裡的意思是什麼呢?.
原來耶穌就是說一個真正的信仰.
一個真正的生命必須要有一個重要的核心.
就是以耶穌基督以上帝為我們的核心.
正如在約翰一書五章十二節也這樣說.
約翰說有了上帝的兒子就有生命.
沒有上帝的兒子就沒有生命.
也相反的在馬太福音十六章二十五至二十六節.
耶穌基督說凡要救自己生命的就會喪失生命.
凡為我喪失生命的就會得著生命.
人又賺得全世界賠償自己生命有什麼益處呢?.
人還能夠拿什麼去換生命呢?.
所以在這裡耶穌再一次告訴我們.
當我們追求豐盛的時候.
其實我們要常常留心.
我們這個追求.
我們所做的一切其實是為了什麼而做的呢?.
很多時候我們沒有察覺.
甚至很多時候當我們埋頭苦幹的時候.
是會很容易迷失自己的.
我自己在教會裡面服侍.
經常都要為教會做一些很重要的決定.
一些很重要的事情.

$^{121}$但當一些這麼重要的事情發生的時候.
我們就會很緊張.
我們就會很埋頭埋腦地去處理這些事情.
但在處理的過程裡面.
其實真的很容易就會迷失了方向.
不知道為什麼要做.
就令到自己很痛苦,很辛勞,很擔心.
今天其實我們每一個人你可以想一下.
當你進入大學.
讀到你心儀的科目.
或者當你能夠得到一個你很理想的工作.
一個教席的時候.
其實當刻我相信每一個人都是很興奮.
很雀躍,很期待.
只不過當我們進入了這個崗位之後.
就可能被很多的工作,很多的責任.
令到我們迷失了自己.
令到很多的事情變得非常痛苦和艱辛.
法國的哲學家迪卡爾說過一句很重要的話.
他說「我思故我在」.
其實是真的.
如果我們沒有理解,明白我們在做什麼的時候.
其實一切就會變得沒意義,變得很痛苦.
我們作為上主的受造物.
是應該要警醒究竟我們在做什麼.
如果離開了我們的上主.
其實我們就什麼都不是.
也做不到什麼.
在教會當中我們都經常要去見一些申請人.
當他申請想成為一位教務.
當他要申請進入神學院的時候.
當然在這個過程當中我們很嚴謹的.
我們會看很多的事情.
看他很全面究竟他的人是不是合適.
不過我相信有些人可能知道.
其中一個我們很看重的元素.
就是我們會嘗試了解.
他是不是真的有從主而來的一個呼召.
就是說他們很確信上帝是要他走這條路.
要他踏上這一個的目的.

$^{161}$當然其實他講完之後我們都沒有辦法去核實的.
這個是他很個人和上帝之間一個的經驗.
但是我們很看重.
因為我們知道目的這條路是不容易的.
如果他沒有這一份的心智.
沒有這一份的領袖.
沒有這一份的初心.
其實他是很難可以走到最後.
今天我們未必每一個人都是在目的當中.
但我都相信其實上帝在我們每個人的生命裡面.
都有祂的奇妙的安排.
都有祂的感召.
讓我們常常都能夠回到這個初心當中.
回到上帝這個感召裡面.
以致我們生命能夠常常都得著那份方向和力量.
找到生命當中的意義.
和得到更豐盛的生命.
這段經文接著就會講到.
這些羊是能夠聽到牧羊人的聲音.
這裡有三個條件.
第一個就是牧羊人原來是會按著那些名字去叫自己的羊.
原來在當時真的很多的牧羊人.
都會為他自己所愛的羊去改名.
改名最基本最簡單是怎樣改呢?.
當然就是按著那些羊一些外在的特質.
好像我們小時候都有很多花名.
叫大狗,亂毛,黑仔,告妹.
我不知道有些什麼花名.
這些花名都是很直接的.
但是每個人都不同.
就是因為每個人都有他的獨特性.
今天當耶穌說每一個牧羊人都會按名去叫自己的羊.
其實就是告訴我們我們要明白我們生命當中的獨特性.
為此而感恩.
因為我們都相信我們這些的獨特都是上帝所賜.
必定有祂奇妙的心意.
今天你知不知道你的獨特之處在哪裡呢?.
你懂不懂得欣賞呢?.
我除了是一個牧師之外.
其實我也是一個潛水的發燒友.

$^{201}$很多人都知道.
但是在學潛水的過程裡面.
其實給我自己有一個很深的體會.
就是很多人問潛水需不需要懂得游泳呢?.
是需要的.
我其實是懂得游泳.
不過我至少也不擅長游泳.
游得不快.
更加游得不遠.
每一次游泳我都覺得很辛苦的.
很花力氣的.
雖然懂得游泳 但是游得不快.
我一直都不明白為甚麼.
直至我學潛水的時候 我才知道.
因為學潛水的人.
我們是需要背著一些懸塊在身上.
令到自己才可以下沉得到.
因為正常人應該都是浮的.
要沉下去一定要背一些懸.
當我背懸的時候 我就發現.
原來我比平常正常的人背的懸.
可以少很多.
也可以下沉得到.
換言之 我本身的身體密度是很高的.
我也試過後來在泳池當中.
放鬆了整個人.
吸滿一口氣 不動.
都會沉到池底.
我就明白 原來我就是因為這個原因.
我一直都不擅長游泳.
所以 如果我要和別人比較.
鬥游泳 我是一定輸的.
但是如果和別人鬥潛水 我就會贏.
所以這個就是我們每一個人.
都體會到是很獨特的.
今日我們的社會傾向.
只是用一個標準去衡量.
人的價值和成敗.
導致我們經常都看不到自己的獨特性.
更加不懂得為我們自己的獨特去感恩.

$^{241}$這樣就很容易埋沒了我們自己的潛能.
也不能夠按著上帝的心意.
去做回我們真正的自己.
唯有當我們能夠從上帝的眼光.
去看到自己的獨特性.
知道自己在上帝的眼中.
是多麼寶貴 是獨一無二的.
我們才能夠活出真正的生命.
得到更豐盛的生命.
第三 就是羊會跟著牧羊人.
牧羊人會領牠們出來 在前頭走.
然後羊就會跟著牠們.
今日我們有沒有緊緊的去跟隨著基督.
跟隨著上帝的腳步呢?.
今日基督宗教仍然是全世界信徒最多的宗教.
很感恩 不過都有點危機.
很快可能就會被另外一些宗教去蓋過.
不過縱使是我們很多基督徒.
但究竟信徒之中有多少真是認真對信仰的呢?.
你對信仰有多認真呢?.
大家自己心中有數.
我見到很多信徒 雖然他知道上帝是真的.
上帝是好的.
但是在生活上卻沒有足夠對上帝的認識.
和有那份願意實踐信仰的心.
我們一方面很希望得到上帝很多的祝福.
但另一方面又不願意去學習上帝的道理.
有時我們覺得這些真理好像已經不合時宜.
要真是這樣做的話.
會不會被人覺得我自己很傻,很蠢呢?.
今日我們如果是這樣的話.
就會很容易好像聖經裡面所講.
耶穌所說「手扶著泥,向後看的人,不配盡上帝的國」.
我們永遠沒有辦法經歷到這份豐盛的生命.
剛才我有提到游泳.
我不知道在座有多少人是不懂游泳的.
我每一次問都有不少.
其實游泳和踩單車都一樣.
有些人說你懂了就懂了.
你不懂就永遠都不懂.

$^{281}$你懂了就懂了 不懂你就學.
但大家要記住一件事.
沒有一個人是能夠先學懂游泳才跳下水.
也沒有人會學懂踩單車才騎上單車.
我們永遠都是要下水才能學懂游泳.
上了單車才能學懂踩單車.
其實信仰上帝的道都是一樣.
我們很難先明白完然後我們才去行.
反過來當我們願意去嘗試.
願意去實踐的時候.
我們就會明白.
所以主耶穌基督在馬大福音第七章.
登山寶訓的總結裡.
說道凡聽了我這些話.
有去做的就好比一個聰明人.
將房子蓋在盤石上.
風吹雨打水沖撞擊個房子.
房子都不會倒塌.
因為根基立在盤石上.
反過來聽了我這些話不去做的.
就好像一個無知的人.
將房子蓋在沙土上.
風吹雨打水沖撞擊個房子.
房子就會倒塌.
並且倒塌得很厲害.
但願我們都能夠藉著上帝的真道.
好好的學習.
好好的實踐.
能夠得著主所賜更豐盛的生命.
最後呢.
這裡說到羊是認得牧羊人的聲音.
認得牧羊人的聲音.
這個認得.
其實是指經過一段長時間的相處.
羊已經熟悉了牧羊人的聲音.
或許這些羊都不知道牧羊人在說什麼.
所以不太清楚牧羊人的意思.
不過牧羊人會認得出.
這把聲音是屬於牧羊人的.
因為羊和牧羊人已經建立了一份信任.

$^{321}$和一份深厚的關係.
羊聽到牧羊人的聲音.
就會感到平安,快慰.
在我人生當中.
當然沒有養過羊.
不過我養過很多隻狗.
可能當中有些朋友都有養狗的經驗.
每隻狗的性格都很不同.
有些比較活潑.
有些比較安靜.
有些比較黏人.
有些比較害羞.
但我養過的幾隻狗.
雖然有不同的性格.
但都很親主人.
我自己工作比較忙.
不是很多時間可以帶牠們出去玩.
多數都在家裡一起.
當我做了一件事 看到牠 叫牠.
或者和牠說一些話的時候.
你都會發現牠的尾巴馬上會很興奮地動.
我相信這就是一隻狗聽到,認得我的聲音.
雖然牠不知道我說什麼.
但是只要一聽到.
牠心裡就會感覺到安舒,感覺到快慰.
今天究竟我們有沒有欣著我們和上帝的那份關係.
有沒有因為我們和上帝的那份親密.
而感到快樂,感到雀躍呢?.
很多時候如果我們將我們的歡樂,我們的滿足.
仍然只是放在很多熟悉的東西上面.
其實是很不穩定的.
大家都知道很多東西我們來的時候會快樂.
但是當我們擁有了之後.
有一天它失去就會帶給我們更大的痛苦.
這一切都不是永恆的.
都不是長久的.
唯有上帝祂自己的愛.
祂和我們的同在.
是永遠不離不棄,是永遠長存.
能夠帶給我們那份永恆的快樂.

$^{361}$所以很盼望今天我們在這段經過當中.
能夠學習這些羊得到更豐盛的生命.
而當我們得到這份豐盛的生命的時候.
其實是會帶給我們兩個很重要的結果的.
在第九節那裡經文說.
你就會得救並且是出入得草吃.
這兩樣東西我們都可以檢視一下自己的.
什麼叫得救呢?.
我們常常都提醒弟兄姊妹.
得救不只是將來的事情.
是今天其實我們就已經在上帝的當中.
享受到那份安穩,那份平安.
任何事情發生.
我們都知道掌管在上帝的手中.
上帝已經寬恕我們,赦免我們.
與我們同在,愛我們.
為我們預備永恆的生命.
所以一個真正得到生命.
得到豐盛生命的人.
他是不會害怕的.
無論遇到什麼事情.
他都能夠是欣然的面對.
能夠有力量的面對.
知道靠著上帝是能夠有出路的.
今天生命就能夠帶給我們這份的安穩.
是我們更加值得追求的.
得草吃就更加有趣.
很多朋友就會覺得得草吃就是發達了.
有很多物質的享受.
但其實是不是呢?.
作為一隻羊.
有草吃.
其實有多物質呢?.
有多富裕呢?.
其實沒什麼的.
反過來你想想一隻羊.
其實牠可以做些什麼呢?.
牠又不可以教書.
又不可以玩電腦,玩手機.
又不可以去吃很多好的東西.

$^{401}$其實牠一生可以做的就是吃草.
牠跟著牧羊人出出入入.
為了什麼呢?.
都是為了吃草.
牧羊人帶著羊.
為了什麼呢?.
就是找草給牠吃.
所以在這裡得草吃.
其實就是在說.
今天我們每一個人.
其實在上帝的心意當中.
其實是很簡單的.
我們不需要想得太複雜.
按著上帝的心意.
按著上帝的指引.
我們就走我們在上帝當中的道路.
其實就會很容易,很輕散.
正如主耶穌基督給我們的應許.
就是說我的擔子是容易的.
我的額是容易的.
我的擔子是輕散的.
只要我們全心全意去事奉上帝.
我們就不需要擔心任何的困難.
任何的挑戰.
任何的痛苦的事情的出現.
這是一份很美麗的圖畫.
希望我們都能夠一起努力.
希望我們都在上帝的當中.
一起得著這份豐盛的生命.
\newpage



\section{詩篇 136:1-26}
\label{sec:032035}
\textbf{因祂的慈愛永遠長存}~~~~張健庭牧師博士
\newline
\newline
連結: \href{http://www.church.com.hk/acms/content.asp?site=cdc\&op=show\&type=product\&code=032035\&layout=sermon}{\texttt{032035}} ~~~~ 語音日期: 2026-04-26
\newline
\newline
\hyperref[sec:032034]{\small{< < < PREV SERMON < < <}}
~
\hyperref[sec:index_chronic]{\small{[返順時目]}}
~
\hyperref[sec:032036]{\small{> > > NEXT SERMON > > >}}
\newline
\newline
詩篇 136:1-26
\newline
\begin{longtable}{cl}
\hline
\hline
章節 & 經文 (和合本修訂版)\\
\hline
136:1 & \begin{tabularx}{0.7\textwidth}{X} 你們要稱謝耶和華，因他本為善；他的慈愛永遠長存。 \end{tabularx} \\ \\ \relax
136:2 & \begin{tabularx}{0.7\textwidth}{X} 你們要稱謝萬神之神，因他的慈愛永遠長存。 \end{tabularx} \\ \\ \relax
136:3 & \begin{tabularx}{0.7\textwidth}{X} 你們要稱謝萬主之主，因他的慈愛永遠長存。 \end{tabularx} \\ \\ \relax
136:4 & \begin{tabularx}{0.7\textwidth}{X} 稱謝那惟一能行大奇事的，因他的慈愛永遠長存。 \end{tabularx} \\ \\ \relax
136:5 & \begin{tabularx}{0.7\textwidth}{X} 稱謝那用智慧造天的，因他的慈愛永遠長存。 \end{tabularx} \\ \\ \relax
136:6 & \begin{tabularx}{0.7\textwidth}{X} 稱謝那鋪地在水以上的，因他的慈愛永遠長存。 \end{tabularx} \\ \\ \relax
136:7 & \begin{tabularx}{0.7\textwidth}{X} 稱謝那造成大光的，因他的慈愛永遠長存。 \end{tabularx} \\ \\ \relax
136:8 & \begin{tabularx}{0.7\textwidth}{X} 他造太陽管白晝，因他的慈愛永遠長存。 \end{tabularx} \\ \\ \relax
136:9 & \begin{tabularx}{0.7\textwidth}{X} 他造月亮星宿管黑夜，因他的慈愛永遠長存。 \end{tabularx} \\ \\ \relax
136:10 & \begin{tabularx}{0.7\textwidth}{X} 稱謝那擊殺埃及凡是頭生的，因他的慈愛永遠長存。 \end{tabularx} \\ \\ \relax
136:11 & \begin{tabularx}{0.7\textwidth}{X} 他以大能的手和伸出來的膀臂，因他的慈愛永遠長存。領以色列人從埃及人中出來，因他的慈愛永遠長存。 \end{tabularx} \\ \\ \relax
136:12 & \begin{tabularx}{0.7\textwidth}{X} 見上節 \end{tabularx} \\ \\ \relax
136:13 & \begin{tabularx}{0.7\textwidth}{X} 稱謝那分裂紅海的，因他的慈愛永遠長存。 \end{tabularx} \\ \\ \relax
136:14 & \begin{tabularx}{0.7\textwidth}{X} 他領以色列從其中經過，因他的慈愛永遠長存； \end{tabularx} \\ \\ \relax
136:15 & \begin{tabularx}{0.7\textwidth}{X} 卻把法老和他的軍隊推落紅海裡，因他的慈愛永遠長存。 \end{tabularx} \\ \\ \relax
136:16 & \begin{tabularx}{0.7\textwidth}{X} 稱謝那引導自己子民行走曠野的，因他的慈愛永遠長存。 \end{tabularx} \\ \\ \relax
136:17 & \begin{tabularx}{0.7\textwidth}{X} 稱謝那擊殺大君王的，因他的慈愛永遠長存。 \end{tabularx} \\ \\ \relax
136:18 & \begin{tabularx}{0.7\textwidth}{X} 他殺戮威武的君王，因他的慈愛永遠長存； \end{tabularx} \\ \\ \relax
136:19 & \begin{tabularx}{0.7\textwidth}{X} 殺戮亞摩利王西宏，因他的慈愛永遠長存； \end{tabularx} \\ \\ \relax
136:20 & \begin{tabularx}{0.7\textwidth}{X} 殺戮巴珊王噩，因他的慈愛永遠長存。 \end{tabularx} \\ \\ \relax
136:21 & \begin{tabularx}{0.7\textwidth}{X} 他賞賜他們的地為業，因他的慈愛永遠長存； \end{tabularx} \\ \\ \relax
136:22 & \begin{tabularx}{0.7\textwidth}{X} 作為他僕人以色列的產業，因他的慈愛永遠長存。 \end{tabularx} \\ \\ \relax
136:23 & \begin{tabularx}{0.7\textwidth}{X} 我們身處卑微，他顧念我們，因他的慈愛永遠長存。 \end{tabularx} \\ \\ \relax
136:24 & \begin{tabularx}{0.7\textwidth}{X} 他搭救我們脫離敵人，因他的慈愛永遠長存。 \end{tabularx} \\ \\ \relax
136:25 & \begin{tabularx}{0.7\textwidth}{X} 凡有血有肉的，他賜糧食，因他的慈愛永遠長存。 \end{tabularx} \\ \\ \relax
136:26 & \begin{tabularx}{0.7\textwidth}{X} 你們要稱謝天上的神，因他的慈愛永遠長存。 \end{tabularx} \\ \\
[1ex]
\hline
\hline
\end{longtable}
$^{1}$主席.
今天我們一起來思想詩篇第136篇.
剛才已經讀過了.
其實如果你讀這篇詩篇的時候.
你必然會留意到.
不斷重複「因他的慈愛永遠長存」這句話.
其實講詩篇有一個困難.
就是在哪裡呢.
就是根本詩篇是一些歌詞.
有時候你聽歌.
你沒得聽那首歌.
就有個人讀那首歌詞給你聽.
你就發覺其實失去了一半的味道.
可能不止一半.
可能馬上掉落了.
只有兩三成.
因為是要配合那首歌一起唱的時候.
你就發覺很有韻味.
詩篇第136篇.
很明顯一定是的.
就是他不斷重複那句話.
如果你這樣讀的時候.
你會說用不用這樣.
如果有人寫文章寫成這樣.
就真的不合格了.
不斷重複老師說.
你重複再寫過吧.
怎麼會這樣寫.
但是這首歌就不是這樣說的.
詩篇第136篇就是一首這樣的歌.
但是沒辦法.
我們又不懂《希伯來文》.
又已經失傳了那些歌詞.
失傳了歌譜.
唯有我們按照我們發揮的想像力.
我們真的要發揮我們自己的想像力.
一邊去聽或者一邊去讀這些詩篇的時候.
讓它好像帶我們去成為一個土文.
因為這個其實是寫詩篇的一個作用.
就是你雖然不能夠去唱出來了.

$^{41}$因為你可能都沒有那些歌譜.
但是你一邊去讀詩篇的時候.
有時候你自己可以讓那些詩篇的內容.
成為你自己一個的禱告.
當你自己成為你自己一個禱告的時候.
你心裡面也可以配合你的韻律.
這個就是我覺得其中一個讀詩篇的.
一個失落了的藝術.
就是有時候.
我不知道你們有沒有.
有時候小時候.
有些詞或者一些譜.
你自己也哼一哼 作一些音下去.
自己嘗試去唱 有沒有人會這樣做的.
我想我都不是很特別的.
我想可能你們都會有的.
就是你唱得好不好聽是另外一回事.
但是你沒有所謂.
因為你自己唱 貪得意玩一下而已.
也是唱給自己聽的.
沒有人叫你去表演的.
但是其實有時候唱.
讀詩篇其實我們要是這樣.
因為你自己將一些韻律加在其中.
成為你自己一個的禱告.
這樣來讀詩篇.
我相信會完全不同的.
因為其中一個困難就是.
很多人讀詩篇就當是讀教科書來讀.
你就會發覺又重複了.
說來說去只有一兩點就很大篇.
長篇大論.
因為你完全搞錯了方向.
不是這樣來讀的.
這是一個導言.
說到詩篇的時候.
說回這篇詩篇有136篇.
它不斷地重複.
「因他的慈愛永遠長存」這句話.
剛才我們也有讀.

$^{81}$當你去讀「因他的慈愛永遠長存」的時候.
其實「慈愛」這個字是很難翻譯的.
在中文的聖經多數是.
「因他的慈愛永遠長存」.
但是其實在原本的希伯來文.
其實「慈愛」這個字.
是有一個忠誠之愛的意思.
在整個的舊約聖經這個字.
其實是不斷地出現的.
出現了超過100次.
如果我沒有記錯應該是130多次.
Hazard這個字就算是英文聖經也很難翻譯.
它是有faithful.
就是有一個忠誠之愛的意思.
所以在整個的舊約聖經.
就算是英文聖經也很難翻譯的.
有一個mercy就是有一個憐憫.
kindness就是有一個恩慈.
或者是strength就是很有力量的一個扶持.
是有多方面的意思的.
所以根本這個字Hazard這個字.
其實是沒有一個固定的單一的名詞可以表達得到.
但是在原文其實就是有多方面的意思.
你也可以形容神自己對人的那一種的愛.
所以就有一個忠誠之愛的一個表達.
所以就算是不同的英文聖經那裡也有不同的.
就是在NASB.
就是New American Standard的Bible那裡.
它就翻譯做Nothing kindness就是愛當中表現到的恩慈.
LSV和ESV就翻譯做Steadfast love.
就是一個很堅定的愛.
King James就是mercy就是憐憫.
NIV就翻譯做Faithful love.
就是一個很忠誠的愛.
這個就是在舊約聖經所要去表達.
神對人那種關係的一個很特別的一面.
就是我們很多時候說愛就是愛.
就是一個就是愛.
就是我們大概有一個這樣的觀念.
但是在聖經那裡所說到.

$^{121}$特別是在舊約聖經用Hazard這個字去說的時候.
就不是單是說一種普通的愛.
神對人之間不是一種普通的愛.
而是一種很堅實.
很忠誠.
很不會反悔.
也是充滿著憐憫.
充滿著恩慈的一種不變的愛.
所以這個Hazard這個字某個程度.
就好像將整個的舊約聖經.
神對人的情懷一種的心腸.
很完美的去表達了出來.
整個舊約聖經.
特別是我們現在在主日學那裡說舊約概論的時候.
當你一直看的時候你就發覺.
那些故事不是一個單獨的故事.
即是有亞伯拉罕,以撒那,雅各,約瑟那.
這些故事本身完全不是一個單獨的故事.
其實他是在說一個神自己的Hazard.
神自己一個忠誠的愛.
怎樣來在人間或者與人的相遇當中.
不斷地去表達.
以致那個舊屬的故事可以不斷地去推進.
而在整個舊約聖經.
根本就是在說神在人間怎樣來表達祂愛的一個故事.
以色列人有他的失敗.
神有他的嚴厲.
人有他的忠誠.
有的時候就是背叛.
就是這一切都成為了交織於.
神怎樣去表現他的Hazard.
他的忠誠之愛的故事發展的元素.
這一切完全不是一個偶然的出現.
舊約所揀選的故事.
無論人物,無論歷史的取材.
其實全部都是在一個神的忠誠之愛的眼光之下.
來將它表達出來.
所以如果你看聖經.
你不去用一個去愛神的眼光去看這些事情.
你當是看故事,看歷史.

$^{161}$多看兩看你就覺得很沒趣味了.
因為說真的.
如果以寫小說的角度來說.
聖經是很失敗的.
你看金庸小說那些.
哇!起承轉合多麼厲害.
都是好看的,你舊約那些.
你看兩看你已經猜得到.
這個都是壞的.
這個王都是注定失敗的.
就是以色列人都沒什麼好收場的.
就是你已經猜到結局了.
寫小說最失敗就是被人猜到結局.
但是你說.
如果你以這樣的角度看故事.
看歷史.
看人物發展.
這些.
你看兩看你就覺得沒什麼興趣了.
但是原來奇異博物館不是以這樣的角度去看聖經.
聖經本身挺好的.
所展現出來就是在不同的.
無論人與神交集.
那個相處之中的各樣的元素.
而表現出神對人的忠誠之愛.
那個hazard.
那個無論人多麼的失敗.
無論人多麼的軟弱.
無論事情的發展多麼的不如人意.
但是神的救贖計劃仍然要向前推.
無論人多麼的背逆.
甚至乎多麼的傷神的心.
好像那個.
俄西亞書的那個不忠的婦人.
但是神仍然可以將那個不忠的婦人再一次收納.
就是這一切就是.
神要讓我們看到祂對人的情懷.
所以奧古斯丁講得很清楚.
就是如果讀聖經.
你不去愛神和愛人.

$^{201}$這樣的角度去看的話.
你就會惡毒聖經.
這句話惡毒聖經就不是奧古斯丁的原文.
但是意思也是差不多的.
我們讀中國古文的時候.
可能你看電視劇《包青天》的時候也說.
《包青天》看到有一個可能是壞的.
不知道是什麼秀才或者什麼貪官的時候.
你讀讀聖賢書所為何事.
你枉讀聖賢書.
你做這樣的事情.
但是這種觀念.
其實在讀聖經那裡.
其實奧古斯丁也是有的.
就是讀聖經你不懂得用愛神愛人的角度.
來看讀聖經的時候.
你就枉讀了聖經.
所以今天選這篇詩篇的時候.
其實他就是從這樣的角度去演繹.
他所看到的事情.
他就是用神對人的忠誠之愛來去看.
整個世界究竟是什麼一回事.
一開始的時候.
他就說到.
神的主權的彰顯.
也是表現了神的忠誠之愛.
他說.
「你們要稱謝耶和華,因祂本位善,祂的慈愛永遠長存.」.
「你們要稱謝萬神之神,因祂的慈愛永遠長存.」.
「你們要稱謝萬主之主.」.
「因祂的慈愛永遠長存.」.
「稱謝那獨行其事的因祂的慈愛永遠長存.」.
這裡就用了一個多方面的去表達.
就是神是至高無上的權能.
神的主權是永遠長存的.
神的權能是永遠長存的.
甚至他是獨行其事的.
但是神的各樣的主權.
神的偉大.
神的權能.

$^{241}$神的屬性.
是要表達.
祂的慈愛永遠長存.
祂的「Hazard」.
祂對人不變的愛.
是永遠長存的.
這一點是非常重要的.
當然我們信耶穌的人.
我們會明白到神是至高的.
神是絕對的有權能的.
神是獨行其事.
神的偉大.
神是很有能力的.
但是在這一切的表達的背後.
詩人再一次的提醒我們一點就是.
是在神的愛的框架之下來去演繹的.
神的無所不在.
我們上一次上個禮拜看到.
詩篇第139篇.
神的無所不在.
往往乃類去逃避你的面.
往往乃類去逃避你的靈.
往往乃類去逃避你的面.
講到最終是回到神對人那種無微不至的認識.
背後都是一個愛的框架.
上一次我們講.
同樣都是一樣.
所以當我們不能單獨去看神的權能和神的主權.
因為這一切都是需要在愛的框架來做一次的重新來去反省.
如果你單單去想.
好像哲學家那樣去想.
神的權能又是怎樣怎樣.
神的主權又是怎樣怎樣.
你很容易得出一些很偏離聖經的結論.
如果神是無所不能的時候.
神要誰生就要誰生.
要誰死就要誰死.
就好像那些大獨裁者那樣.
這樣的演繹和詮釋是完全偏離了聖經那個神的慈愛永遠長存的框架.
無論神的揀選也好.

$^{281}$神的主權也好.
神自己的獨行其事的能力也好.
所有的一切不能夠單獨的去思想.
一切都是要在神的愛的框架當中去詮釋的.
弟兄姊妹.
今天都是一樣.
有時候可能在你的生命當中遇到的困難.
遇到的挑戰.
遇到的苦難或挫折.
可能你單獨的去看那些事情發生的時候.
你可能會想神的主權又怎樣.
神的權能為何不幫我.
神的能力是否因為這件事而做不到.
你不是這樣去詮釋不需要是這樣.
因為一切都是在神的愛的框架之下來演繹的.
我們是首先認識神的愛.
神的忠誠不變的愛.
然後才去認識神其他的屬性.
這個次序是要有的.
不是所有的屬性是平排的.
你知道在聖經當中.
神去表達他自己的時候.
當你一路讀聖經的時候.
你看到神怎樣去表達他自己的時候.
有些屬性是放在其他屬性之下的.
譬如他講神的不變.
神的永恆不變.
如果以哲學家來說.
神的不變應該是放在很高的層次的.
因為神的不變影響了神的B型存在.
究竟神是怎樣存在的.
是很高的.
但是在聖經來說不是的.
聖經說神有時候是會變的.
是不是?.
神改變了他的心意.
你看約拿書已經是了.
神看到那些尼尼眉人悔改.
改變了他的心意.
約拿就講了一個假的.

$^{321}$也不是假的.
是失落的預言.
約拿明明是在打鑼打鼓.
說三天之後尼尼眉人就會傾覆.
三天之後沒有發生.
所以他很生氣.
因為神改變了.
在舊日的聖經多次都講神改變了他的心意.
因為神的不變.
神自己怎樣永恆存在.
B型存在本身的框架.
是超過人可以想象得透.
而神要給我們唯一一個合法的理解.
就是從神的慈愛.
神的Hazard.
忠誠的愛當中去理解.
神的愛.
那種愛是不變的.
是長久是堅實的.
就是Hazard這個字所要表達.
無論以色列人有多麼的背弱.
神都是忠誠的.
無論我們有多麼的軟弱.
但是神都不會離棄.
因為神就是忠誠的Hazard.
他是堅實的.
仍然選擇站在我們的邊.
我們剛剛今天講朱一學的時候.
在創世記有一個很奇特的故事.
就是講到神和亞伯拉罕顯現的時候.
然後神就叫亞伯拉罕也把動物劈成兩邊.
接著亞伯拉罕又打瞌睡.
然後有一個大黑暗臨到.
接著亞伯拉罕又打瞌睡.
接著睡到半醒的時候.
看到在黑暗當中有一個火把頭.
有一個火爐在動物當中走過.
《安聯聖經》沒有什麼很詳細的解釋.
但是就很詳細的講發生了一件這樣的事情.
很多人就不知道究竟講什麼.

$^{361}$他們無端端講完叫神和亞伯拉罕說.
你看看天上的星,海邊的沙.
你看到天上的星的時候.
就看到你的子孫有這麼多.
萬俗都要因你得福了.
最後一次講了三重英海的時候.
就有一個這樣的故事發生了.
是不是?.
這個故事本身其實就是想要反映神對人的忠誠愛.
因為在那個時代來說.
如果兩個國家打仗.
接著兩個皇帝立和約的時候.
就要劈開動物.
接著兩個皇帝就在動物當中走過.
一起走過的時候就代表立約.
那誰背約的就不得好結果.
就好像動物一樣被劈開兩邊.
那是一個當時立約的儀式.
但是這個立約的儀式的故事就很特別.
就講到本來應該神和亞伯拉罕立約.
神就拉著亞伯拉罕的手.
兩個人在動物當中走過就行了.
因為神叫亞伯拉罕斬開動物.
很明顯就是要履行這個儀式的.
是不是?.
但是接著那個故事就講到.
亞伯拉罕就睡了一覺.
接著大黑暗來到了.
其實這個故事本身就講到人那種黑暗和軟弱.
而亞伯拉罕半夢半醒.
醒來的時候就看到有一個火把頭.
有一個長官在動物當中走過.
那個是誰來的呢?.
那個就是神自己.
意思就是本來那個立約的應該兩方面的契約.
是神自己獨立的承擔.
即使人在黑暗的當中沒有履行那個約的條件.
神自己都是這樣走過.
那個故事本身其實就是預演了整個舊約後期的發展.
就是這個約本來應該兩方面來立的.

$^{401}$但是人是多次的失敗.
但是即使在黑暗的當中.
耶和華神自己單獨的去承擔.
甚至乎承擔了那個背若的結果.
就是你去到《耳塞俠書》那個異牧.
第四十章後面那個異牧的出現.
五十三章那裡那個.
因他受的刑罰我們得平安.
因他受的鞭傷我們得醫治.
那個異牧那個故事.
那麼.
簡單來說就是所有的一切.
都是表現了神對人的那種忠誠.
即使我們的軟弱.
即使我們的失敗.
但是神仍然都是很堅實的一個去決定.
他跟我們的關係在慈愛的框架之下是不會變的.
而接下來.
所謂詩篇第五十.
136篇第五至第九章就講到神的創造.
要清者那用智慧做天的.
因他的慈愛永遠長存.
清者那鋪地在水以上的.
因他的慈愛永遠長存.
清者那造成大光的.
因他的慈愛永遠長存.
他做日頭管白晝.
他的慈愛永遠長存.
他做月亮管黑夜.
因他的慈愛永遠長存.
這裡就講到神的創造.
整個我們人是活在神的創造之下的.
可能你平時很少這樣去想.
但是如果你.
譬如你想像一下自己.
譬如你去到大自然那裡.
譬如你去到一個很偉大的一個.
國家公園自然環境的時候.
你對比一下整個宇宙萬物.
星宿窮瘡那種偉大.

$^{441}$你就發覺我們是多麼渺小.
我們都是在神的宇宙萬物的創造當中.
好像一粒微塵一樣.
我們無足輕重的坦白講.
短暫的人生在這個世界上.
你可以活到多少年.
對比宇宙浩瀚的歷史.
你覺得就算是一百年.
一百五十年.
根本在整個宇宙的時鐘裡面.
零點零零零幾秒都不夠.
如果宇宙是二十四小時.
根本是無足輕重.
但是在這個宇宙萬物的創造之下.
詩人就看.
我們是活在神的愛的框架之下而創造.
宇宙萬物是因為神的愛而創造.
而我們也是被安置在.
這個愛的創造之下.
可以這麼說.
我們有那個安身立命的居所.
這個世界稱為我們的家.
是因為神的愛設計的.
在舊約的故事當中.
講到神就是為了愛而創造了宇宙萬物.
最終愛的高潮.
就按著祂自己的形象做人.
將人放置在伊甸園當中.
這一切是愛的設立.
你想想如果你去到一個地方.
譬如你去旅行去探望朋友.
朋友知道你來住的時候.
又特別為你預備了一個房間.
整理得很漂亮很舒適.
你已經覺得很貼心.
因為那個居住的地方.
居所雖然是短暫.
但是有愛在其中安置.
整理為你預備的.
神是這樣為人預備的.

$^{481}$人活在世上雖然是短暫的存在.
但是整個宇宙萬物.
當然現在的世界也有很多的罪惡.
醜陋,不足.
當然是有的.
因為這個世界已經不是伊甸園了.
已經墮落了.
但是某個程度.
今天我們所享用的一切.
因為都是在愛的框架之中.
神為我們預備.
所以在這一段的時候.
詩人就說「因他的慈愛永遠長存」.
你看看創造的天地.
你看看神所創造的太陽,月亮.
各樣東西好像.
某個程度都挺自我中心.
但是也是真的這樣.
神好像將人放在中心之下.
來去如珠如寶的設計.
在舊約有句話說.
「太陽還存,月亮還在」.
神的愛就不會止息.
所以這裡就提到神造月亮和太陽.
神的慈愛永遠長存.
在洪水毀滅大地之後.
神也是這麼說的.
「我用彩虹在天上和你們立約」.
當你們每一次看到彩虹的時候.
你就知道.
我不會再用洪水來毀滅大地.
所以地上仍然都可以有四季.
春夏秋冬不斷地運作.
這個宇宙之所以不完全崩潰.
因為神的愛仍然在其中維持.
而祂用愛來維持這個宇宙的核心.
就是因為祂愛我們.
仍然為我們預備.
這個就是我們的家鄉.
接著他說得夠熟.

$^{521}$「清謝那擊殺埃及人之長子,因他的慈愛永遠長存」.
「他令以色列人從他們中間出來,因他的慈愛永遠長存」.
這一段的歷史當然以色列人就很熟悉.
說到出埃及的殺長子十災.
那段的歷史,接著出埃及.
但是在舊約的框架來說.
出埃及的事件就成為了整個救贖的一個典範.
當每一次他們說到出埃及的時候.
其實他就想到神的行奇妙的事.
來拯救他們脫離一個極大的苦難.
今天我們未必是做奴役的.
我們未必是做奴隸的.
不是好像以色列人在埃及那樣.
但是今天我們也可以想像到.
多方面我們需要神的救贖.
你的軟弱需要神的救贖.
有些罪惡的捆綁需要神的救贖.
你自己在黑暗當中.
你都好像沒有出路的時候.
你需要神的救贖.
這一切的救贖都是源於神的慈愛永遠長存.
出埃及的事件好像就是那個救贖的一個典範.
告訴我們神的救贖.
神的慈愛的救贖仍然的長存.
今天在你的困境當中.
神的傍臂沒有縮短.
神的救贖仍然的存在.
可能你有你自己的各樣的掙扎.
無論是身體的軟弱.
心靈的困苦.
或者是你自己面對一個實際的困難的時候.
你可以仰賴神的救贖.
因為神的慈愛.
神的堅實的愛仍然是長存.
接著他說.
他施展大能的手和伸出來的傍臂.
因他的慈愛永遠長存.
稱謝那分開紅外的.
因他的慈愛永遠長存.
他令認識的人從其中經過.

$^{561}$因他的慈愛永遠長存.
他卻把法老和他的軍兵推翻在紅海裡.
因他的慈愛永遠長存.
這裡當然是說回以色列人出埃及的時候.
後來被法老追.
追到山窮水盡.
前面是紅海.
後面是法老追上來.
但是神就行了一件很奇妙的事.
就是分開紅海叫他們過去.
接著就把紅海合起來.
殺了那些埃及的追兵.
這個故事大家都很熟悉的.
但是斯人就從另一個角度來說.
是因為神的慈愛.
我們在困境當中神仍然會保護的.
好像以色列人.
已經被人追到山窮水盡的時候.
神的慈愛保護著他們.
今天你的困境.
可能你都不知道.
你都絞盡腦汁.
你都想像不到有什麼出路的時候.
神的保護仍然都在其中.
神仍然都在保護你.
因為他堅實的愛.
仍然都沒有改變.
所以無論什麼困境.
神都在保護著你.
接著他就再說.
「稱謝那引渡自己的民.
行走曠野的.
因他的慈愛永遠長存」.
接著你知道.
接下來那段故事過了紅海之後.
接著他們叛逆.
接著有40年在曠野漂流.
那個故事大家都很熟悉.
但是在私人的角度來說.
他就看到因為神的慈愛永遠長存.

$^{601}$所以他引渡他的自己的民走曠野.
這裡是很有情感上的牽絆.
他說他引渡自己的民走曠野.
他這麼說的時候.
是沒有一個審判的語氣的.
是不是?.
本來以色列人在曠野40年行走.
其實是一個審判的.
可以這麼說是罪有應得的懲罰.
他不肯進去.
但是在這裡他不是這麼說.
詩人不是這麼說.
他說他引渡自己的民走曠野.
意思在愛的框架之中.
神用另外一個角度.
或者神想我們來看.
我們的人生有另外一個角度來看.
有時候沒錯我們好像是走在曠野的路上.
有時候走一些冤枉路.
甚至乎走一些我自己都後悔不應該這樣走的路.
但是沒辦法已經走到那裡了.
但是到了這樣的時候.
那又怎樣呢?.
神就告訴我們.
因為我堅實的愛.
因為我忠誠的愛.
我的悲情的愛.
我陪著你們走.
我仍然都看你是屬於我的.
我引渡你.
這個就是神.
怎樣來詮釋我們的人生.
有時候我們人生當中.
我們自己都發覺.
哎呀不應該這樣.
或者有得再選我也不敢選.
好像現在沒辦法.
自己承受了那些曠野走冤枉路的後果.
但是神說不是這樣看.
你仍然都是我的子民.

$^{641}$我的慈愛.
我堅實的愛.
仍然都在其中的時候.
我就會繼續引渡你.
我繼續引渡你.
於是他就再說.
撒鹿阿摩利王西雲因他的慈愛永遠長存.
又撒巴山和岳因他的慈愛永遠長存.
撒撒撒撒撒撒撒撒撒撒撒撒撒撒撒撒撒撒撒撒撒撒撒撒撒撒撒撒撒撒撒撒撒撒撒撒撒撒撒撒撒撒撒.
現在打仗都覺得很殘忍.
但我們不要用這樣的角度去看.
因為以色列人那個時代.
打仗來說.
不是你打人人打你.
所以撒鹿本身不是重點.
撒鹿本身就是說.
神幫助他們去克服要去殺他們的敵國.
其實這裡是說出了一種勝利.
神帶領以色列人走過抗夜路.
最終不單止是走完這段路.
我們人生的路好像抗夜路.
我們不單止是要走完.
我們是要得勝地走完.
是走到得勝的地步.
就是這裡所說.
戰勝了那些本來要殺他們的皇帝.
當時來說那些人跟他們打仗.
是因為他們要殺以色列人.
你看回那段故事.
他們走過那些路的時候.
那個地頭蟲可以這麼說.
那個地域的霸主.
不讓他們走過還要殺他們.
以色列人就跟他們打仗.
就是說我們人生在愛的框架之下.
全息的時候.
不單止是苟延殘喘.
撿回一條命就算了.
我們是要勝利的.
在神的帶領之下.

$^{681}$我們是可以勝利的走完人生的路.
這個就是這裡所說.
接著不單止是這樣.
他將他們的地賜他的百姓為業.
因他的慈愛永遠長存.
就是賜他的僕人以色列為業.
因他的慈愛永遠長存.
因為神的堅實的愛.
我們不是走完就算了.
人生的路不是走完就算了.
我們是有所得的.
得基業.
我們得到天上的賞賜.
得到神所為我們預備的福氣.
我們是有所得的.
那些路我們不是撿回一條命.
這次不死就算了.
神是一路引導一路祝福給我們.
得勝的時候是要我們有所得的.
我們將來說回天家見主的時候.
我們是有所得著的.
我們不是像保羅在《歌吾多前書》那裡說.
有些人就好像著火.
家中失火走出來.
好像空條條.
就是你們見過有些人走火警.
不是走火警.
是真實火災走出來.
穿著一件睡衣.
什麼都沒有了.
本來就形容有些人一生.
渡過一生的時候.
他們沒有去預備好見主的時候.
就這樣燒了.
用那些草木和街建造人生.
但是他說.
我們不是.
我們可以得基業.
我們用金銀寶石來建造的時候.
我們是可以得到主的賞賜的.

$^{721}$接著他就再說.
他顧念我們在卑微的地步.
因他的慈愛永遠長存.
他救拔我們脫離敵人.
因他的慈愛永遠長存.
有時候我們都是很軟弱的.
我們是很卑微的.
但是因為神的愛.
神顧念我們.
祂顧念我們的軟弱.
祂看出我們不過是塵土.
很多時候.
有時候我們都懷疑.
究竟我這麼軟弱.
究竟前面是怎樣呢?.
我們人生的路.
見到神的時候.
神又會怎樣看我們呢?.
但是有一樣東西你知道.
就是祂堅實的愛.
是長存的時候.
祂會顧念我們.
顧念我們的本體.
思念我們不過是塵土.
在軟弱的時候.
祂會提拔我們.
救我們脫離敵人.
接著他就再說.
他賜糧食給凡有血氣的.
因他的慈愛永遠長存.
你們要稱謝天上的主.
因他的慈愛永遠長存.
去到最後的時候.
你當然如果用回那段.
詩人那段.
他用回那段聖經來說.
可能你都會回應.
可能他說的是天賜嗎哪.
走抗夜路神給他們糧食.
但是我覺得.

$^{761}$也未必需要這樣看.
似乎去到最後的時候.
詩人讓我們想起一件事.
今天你有衣有食.
其實都是神的慈愛.
有時候你想到人生大圖畫的時候.
有時候低谷.
有時候軟弱.
有時候困難.
有時候危險.
是很多的挑戰.
但是想到這麼大的時候.
神顧念你的軟弱.
神的慈愛永遠長存.
這個當然是很實在的.
是不是.
但有時候我們也未必每時每刻.
都是想這些人生大圖畫.
有時候就在很細微的.
每天的生活當中.
你的工作.
你的家庭.
你所面對的很小的事情.
但是你想想.
就是吃一碗飯.
就是喝一碗湯.
就是吸一口呼吸的空氣.
都好像是神將糧食賜給我們一樣.
因為在神愛的框架之下.
祂的慈愛永遠長存.
因為神堅實的愛.
永遠都沒有離開屬祂的人.
\newpage



\section{以弗所書 2:1-10}
\label{sec:032036}
\textbf{天父之恩,聖徒之行}~~~~歐陽萬廷弟兄
\newline
\newline
連結: \href{http://www.church.com.hk/acms/content.asp?site=cdc\&op=show\&type=product\&code=032036\&layout=sermon}{\texttt{032036}} ~~~~ 語音日期: 2026-04-26
\newline
\newline
\hyperref[sec:032035]{\small{< < < PREV SERMON < < <}}
~
\hyperref[sec:index_chronic]{\small{[返順時目]}}
~
\hyperref[sec:032037]{\small{> > > NEXT SERMON > > >}}
\newline
\newline
以弗所書 2:1-10
\newline
\begin{longtable}{cl}
\hline
\hline
章節 & 經文 (和合本修訂版)\\
\hline
2:1 & \begin{tabularx}{0.7\textwidth}{X} 從前，你們因著自己的過犯罪惡而死了。 \end{tabularx} \\ \\ \relax
2:2 & \begin{tabularx}{0.7\textwidth}{X} 那時，你們在過犯罪惡中生活，隨從今世的風俗，順服空中掌權者的領袖，就是現今在悖逆的人心中運行的邪靈。 \end{tabularx} \\ \\ \relax
2:3 & \begin{tabularx}{0.7\textwidth}{X} 我們從前也都生活在他們當中，放縱肉體的私慾，隨著肉體和心中的意念去做，和別人一樣，生來就是該受懲罰的人。 \end{tabularx} \\ \\ \relax
2:4 & \begin{tabularx}{0.7\textwidth}{X} 然而，神有豐富的憐憫，因著他愛我們的大愛， \end{tabularx} \\ \\ \relax
2:5 & \begin{tabularx}{0.7\textwidth}{X} 竟在我們因過犯而死了的時候，使我們與基督一同活過來—可見你們得救是本乎恩— \end{tabularx} \\ \\ \relax
2:6 & \begin{tabularx}{0.7\textwidth}{X} 他又使我們在基督耶穌裡與他一同復活，一同坐在天上， \end{tabularx} \\ \\ \relax
2:7 & \begin{tabularx}{0.7\textwidth}{X} 為要把他極豐富的恩典，就是他在基督耶穌裡向我們所施的恩慈，顯明給後來的世代。 \end{tabularx} \\ \\ \relax
2:8 & \begin{tabularx}{0.7\textwidth}{X} 你們得救是本乎恩，也因著信；這並不是出於自己，而是神所賜的； \end{tabularx} \\ \\ \relax
2:9 & \begin{tabularx}{0.7\textwidth}{X} 也不是出於行為，免得有人自誇。 \end{tabularx} \\ \\ \relax
2:10 & \begin{tabularx}{0.7\textwidth}{X} 我們是他所造之物，在基督耶穌裡創造的，為要使我們行善，就是神早已預備好要我們做的。 \end{tabularx} \\ \\
[1ex]
\hline
\hline
\end{longtable}
$^{1}$我今日的題目是天父之恩 聖徒之恆.
其實我有一個大膽的嘗試.
怎樣大膽的嘗試呢?.
就是希望今天的早晨.
我能夠看完三篇聖經.
所以需要有些….
希望神幫助我們.
在開始的時候我想讀一段翻譯的畫.
是從Charles Winder寫的一本書.
Paul 保羅這本書裡面的一小段.
為甚麼要讀這段呢?.
就是他去梵蒂岡的時候.
去到一個世監獄.
這是他自己寫的一段說話.
我將其中一部分折錄了出來.
請大家聽我讀.
在距離梵蒂岡掃經里的一條偏僻的小巷裡.
樹立著一座細細的建築物.
據說保羅正是在這個建築的牢房裡度過他生命最後的時光.
當然究竟是不是當年的牢房.
至今仍然是一個爭議的話題.
我們順著樓梯下去下面一個很狹窄幽暗的地下的穴洞.
在這個幽暗的牢房裡面足足逗留了大概半個小時.
那裡很陰寒很潮濕.
瀰漫著一種沉附的霉味.
在天花板上開了一個細細的窗格.
更加能夠透進一支些微的天光.
當我獨自坐在那個冰冷石質的地面上的時候.
我不禁在我腦袋裡勾勒出.
保羅當時在生命裡面最後那段的日子那種遭遇.
我們在那個牢房中站了低聲交談.
全程來到灌注來到聽.
彷彿聽到那個沉默的石塊能夠開口.
會向我們美美道出他當日所見到歷史的故事.
我自己想想.
這裡就是那位撰寫很多筆漏的書卷.
並且被譽為第一世紀最大的宣教士.
與教會建立在保羅是長眠之地.
讓我們低頭祈禱.
這個早晨我聽到很多弟兄姊妹都有些少咳嗽.

$^{41}$懇求主理自己的特別的恩手來按摩.
撫摸他們讓他們能夠依次他們.
我亦都將今日講道的時間恭敬教給你.
因為話語是從你而來解釋話語的是聖靈.
最後願我們藉著話語來向你敬拜.
藉著話語來向你歌頌.
打開我們的心智打開我們的意念.
挖到我們心裡面最深的深處.
讓我們摸到你自己的心意.
奉主耶穌賜給我們.
有一個問題是你們當中誰是聖徒.
有沒有人再舉多一次手.
好像少了很多人.
為何要提到聖徒這個詞呢.
聖徒在希臘文我不敢讀.
因為我不認識希臘文.
聖徒這個詞在整本新約聖經裡出現了61次.
其中八次在保羅寫的羅馬書出現.
但在以弗所的六章裡面.
我們看到他出現了九次.
甚麼是一個信徒呢.
而是說信徒所擁有的一種新的身份.
是藉著耶穌基督被神稱為分別為聖的人.
我們可以用一個很普通話來形容.
如果神站在這裡.
Anissa站在旁邊.
我不是故意的.
神就對大家說.
這個人是屬於我的.
說得比較俗一點.
我包住你.
你是屬於我的.
這是地位上的一種分別.
是一種身份的不同.
就好像父母說這個是我的兒子.
這個是我的女兒.
意思是甚麼呢.
他是屬於我的 是我的一部份.
無論有時候多頑皮.
他仍然是屬於神的子女.

$^{81}$聖徒就是一個這樣的概念.
雖然聖徒是說到一個位份上的分別.
但聖徒這兩個字也帶著一個意念.
是甚麼呢.
是我們在生活裡面.
有一個要求 有一個標準.
是我們需要去達到.
當我想到聖徒這兩個字的時候.
有兩個影像在我腦袋裡面.
一個是一道門.
我希望大家不要笑.
如果你去到一道門.
譬如去星巴克.
去到一道門.
外面寫著是男廁的話.
你知道只有男生才能夠進去.
意思是甚麼呢.
外面 一道門的外面和一道門的裡面.
兩樣應該相對.
另外一個影像讓我看到的.
一面鏡.
鏡是反映到實物真實的樣子.
而不是扭曲的樣子.
所以我希望聖徒這兩個字.
表達了不單止你內裡面的那種真實.
也表達了你外面的人看到的那種形象.
我將耳筆鎖書分為兩部分.
很多人都是閱讀耳筆鎖書的時候都提到.
耳筆鎖書可以分為兩部分.
在很多的英文聖經裡面.
都是用 Therefore.
所以你可以翻譯為因此或者是所以.
第一部分是一至三章.
很容易記住是一至三章.
然後下一部分是講到第四章到六章.
我給兩段的分段有兩個不同的題目.
我第一個的分段一到三章.
是講到聖徒在基督裡面救贖的豐盛.
而在四到六章是聖徒的行事為人.
當與所信的福音相承.

$^{121}$我們今天只有時間去分享第一部分.
就是第一章到第三章.
以後有機會再和大家分享第四到第六章.
其實主席剛才帶領我們讀那段經文.
就是從第三章轉接到第四章最後那一段.
但我們今天所希望讀到的.
就是全部的一到三章.
所以如果你有帶聖經來.
或者你自己有手機的話.
你先打開.
我們讀到的經文.
如果你想看多一點的話.
你可以參考一下你們所看的經文.
我怎樣開始來講.
以弗所書一到三章呢.
我希望用一個我們十部《釋經法》.
我們都提到的一個搜查訊息的方法.
就是看一些重複的字.
我們嘗試在以弗所書的一些重複的關鍵的詞裡面.
嘗試去明白作者在這幾章的聖經裡面.
要表達一些什麼的中心思想.
第一個出現最多的次數在以弗所書.
就是「在基督裡」.
或者是「在基督耶穌裡」.
或者是「在他裡」.
「他」當然是指著基督.
在整本的以弗所書裡面.
我們看到出現的次數有二十三次.
所以這個提醒我們.
保羅要跟我們講一個很重要的訊息.
就是聖徒的生活,行為.
無論他每一個部份.
都應該以耶穌基督為他們的中心.
我希望大家能夠用這個詞.
當你重新再讀以弗所書的時候.
能夠作為一個工具來打開書裡面的訊息.
所以我們看到.
我們一會兒會探討到.
現在稍微講一講.
就是說.

$^{161}$信徒的生命.
跟信徒所表達.
在生活裡面所表達的.
應該是一致的.
應該是沒有差別的.
意思是什麼呢?.
當你認識神很淺的時候.
你所表達的生命就很膚淺.
當你如果是盡心到更加認識耶穌基督.
在你的生命裡面.
有何等浩大的恩典的時候.
你就懂得應該去回應.
神所付出這麼大的恩典.
我們有很多的部分.
在整本的聖經裡面.
講到生活.
講到生命.
同時在後半段.
我們看到.
講到教會.
也都是講到個人.
需要的一種回應.
我們首先要看.
另外其他的字.
我們可以用作為.
我們查考這本書的工具.
另外就是講到靈.
或者是聖靈.
講到行善.
或者行事.
Walk.
或者是講到豐富.
或者是豐盛.
這些都是一些.
我們要考慮的東西.
我們要學習的東西.
或者是講到豐富或者是豐盛.
這些都是一些重複的詞.
我們可以用來作為一種.
查考這本書的一種工具.

$^{201}$我們進入.
已忽所書的第一章.
第一章裡面.
除了第一章一到三節是一個文後.
他接著是從第四節到十四節.
是講到基督徒那種生命.
是一種的頌讚.
這裡說他頌讚.
這個父神在基督耶穌裡面.
所賜給我們的屬靈的福氣.
如果我們聽見那個福音的時候.
我們很容易聯想起.
究竟什麼叫福音.
我可以在五分鐘或者一分鐘之內講完.
什麼叫福音.
我們都很熟悉.
福音就是講到人是有罪的.
人不能夠改變自己的罪.
不能夠改變自己的最終的終局.
只有神耶穌基督來到這個世界的時候.
才能夠改變.
因為耶穌為我們的罪.
是讓人將他釘在十字架裡面.
廢去所有罪帶來的後果.
耶穌在十字架上面所流出的寶血.
讓他們代替.
那種生命的代替.
讓我們得到一個永恆的生命.
靠著相信耶穌基督.
我們能夠進入永恆的裡面.
能夠與神來和好.
這是很簡單地描述一下.
信耶穌究竟是什麼一回事.
可能我們有五分鐘的版本.
或者有七分鐘,十分鐘.
甚至半個小時的版本.
但這個基礎就是我們有罪.
我們需要耶穌基督來拯救.
然後我們會去到因為耶穌的救贖.
我們能夠與神和好.

$^{241}$永遠得到一個永恆的生命.
但我們很少想到我們.
信仰這個福音是一種福氣.
有沒有?.
誰曾經這樣想過?.
我希望大家.
我見到Velma.
很好.
有時候我們拿著一杯咖啡坐在那裡.
我們安靜的時候.
我們想一想.
原來我們今天.
能夠得到這個平安.
得到這個喜樂.
能夠得到那種安靜.
原來是一種福氣.
因為我們得到的訊息是什麼?.
我們不再被定罪.
我們看看第三節.
讓我讀給大家聽.
第三節提到.
是父神在耶穌基督裡.
將天上各種屬靈的福氣賜給我們.
在基督耶穌裡.
將各樣屬靈的福氣賜給我們.
容許我.
很長的經文.
容許我讀一讀.
因為很重要.
一到三章.
因為祂在創世以前.
在基督裡揀選了我們.
令我們成為聖潔.
無有瑕疵.
滿有愛心.
祂按著自己旨意喜悅的.
預定我們.
藉著耶穌基督得而至的名分.
讓使祂的榮耀的恩典得著稱讚.
這恩典使祂在愛子裡白白賜給我們.

$^{281}$我們藉著這愛子的血得蒙救贖.
過凡得以赦免.
這是照著祂豐富的恩典.
充充足足地賞給我們的.
祂以諸般的智慧聰明.
照著自己在基督裡所立定的旨意.
使我們得知道祂旨意的奧秘.
讓我們細心看看這幾節的經文.
原來我們今天這麼多位坐在這裡.
已經是信了耶穌的人.
是不是你當時那段時間.
你信主那段時間.
我們也有信主不同的時間.
但這裡很清楚地告訴我們一件事.
因為祂在創世以前.
已經在基督裡面揀選了我們.
這是極大的福氣.
你還沒出生之前.
你還沒說我要信耶穌之前.
你還沒聽福音之前.
神已經在基督耶穌裡面揀選了你.
讓你有機會去相信祂.
讓你有機會去接受這麼寶貴的生命.
不僅如此.
祂預見我們是稱為神自己的兒女.
第五節.
我們要藉著耶穌基督得到兒子的名分.
這個發生在很久之前.
甚至我們還沒出生之前.
祂所有的安排.
是我們能夠白白得到這個生命.
我們能夠經過祂愛子的血得無救數.
而且得到祂豐滿的恩典.
這是一個奧秘.
我們不明白.
有一部份我們很大力.
很多人都不明白.
怎麼可能神在我們還沒出生之前.
已經揀選了我們.
叫我們成為祂自己的兒女.

$^{321}$讓我讀一讀這幾節的經文.
「要照著所安排的.
在時機成熟的時候.
使天上地下一切所有的都在基督裡.
同季如日.
我們也在祂裡面得了基業.
這是賢士那位隨己而行萬事的.
神照著自己的旨意所預定的.
為要使我們.
這些首先把希望寄託在基督裡的人.
頌讚祂的榮耀」.
第十三節.
「在基督裡你們聽見真理的道.
就是那些使你們得救的福音.
你們也信了.
就接受了所應許的聖靈為印記.
這聖靈是我們得基業的憑據.
直到神的子民得救熟.
使他榮耀得道稱讚」.
我們信耶穌不能夠跳過那個步驟.
第十三章.
我們要聽到福音.
我們要信那個福音.
我們要接受那個福音.
信不是一個口頭上的回應.
信是你的口裡面承認.
是心裡面相信.
是一致性的相信那個福音會帶來.
你自己的救恩裡面所得到的救贖.
要接受.
就是神當我們信那一刻的時候.
將聖靈賜給我們.
讓我們能夠有那個確信說.
我是屬於神的.
神也是屬於我的.
聖靈在裡面與我們.
聖靈與我們的心同征.
我們是神的兒女.
這個是聖靈的印記.
這個是我們得到基業的那個憑據.

$^{361}$所以我們今天採取了每一個人.
甚至是在當日耶穌在地上生活那段時間.
每一件事情.
無論耶穌在加勒略海邊傳道的時候.
無論祂在井旁和那個婦人傳福音的時候.
無論祂去到曠野和四五千人的老人傳福音的時候.
講到天國的問題的時候.
所有的事情的發生.
是不是偶然的.
是神自己所安排的.
神給他們有機會聽到天國的福音.
給他們有機會能夠對這個福音有所回應.
這個就是預定的意思.
對於所有今天在座的人.
信的已經是相信主耶穌的人.
這個是一個福氣.
但是對不起.
從反面來看的時候.
如果你今天坐在這個位置.
仍然未相信耶穌基督.
我們的願望就是你快點相信.
因為對於你們來說.
曾經聽到福音又拒絕福音的時候.
這個聽到福音可能成為你自己的咒詛.
有一天你上到天堂.
通常是說彼得站在天堂的門口.
今天我用約翰.
如果約翰站在門口說.
為什麼我會讓你進入這個天堂.
你會怎樣回應?.
神可能告訴你.
你已經聽到福音了.
很多次了.
很多人向你傳福音.
但是你沒有一種回應.
你有機會聽到福音成為你自己的咒詛.
所以我問大家一個問題.
你在神的眼中重不重要?.
特別是對一些已經是信主的人.
你在神的眼中重不重要?.

$^{401}$我希望我們的答案是重要.
我們跳到第一章的下半節.
對不起,我剛才漏了一些東西.
我自己怎樣回答這個問題.
我記得我昨天回去.
我沒有帶我的舊聖經.
我在1970年信耶穌.
受贈的時候是9月13日.
教會送了一本聖經.
將我受洗的日子寫在上面.
我們更重要的是.
明白這個恩典需要我們有回應.
在基督裡是一個什麼的意義?.
沒錯,可能今天你聽到的訊息是.
你在神的眼中很重要.
但我可能需要跟進一個問題.
你的信心,你對神的認識.
是否停留在當天你信耶穌的時候?.
今天的你.
有沒有進到一個更深的地步?.
更加深入的認識.
愛我們的主耶穌基督?.
如果你沒有的話,你仍然是一個嬰孩.
在聖經裡鼓勵我們要吃乾糧.
要明白真理,要明白.
神在我們生命裡.
甚至在這個世界裡,在我們周圍裡.
做了一個何等奇妙的工作?.
我們每一個信耶穌的人都有一個永遠的盼望.
我們問一句話.
在第一章15到第23節提到一件事.
這是保羅的禱告.
保羅如何為已乏所的人禱告.
第17節我們注意幾個字.
其中一件事是你要真正的認識神.
你要明白神呼召你來的那種指望.
The hope of your calling.
要明白神給你們的基業是何等的豐盛.
要明白神在你們身上和所有人.
相信和不相信的人身上.

$^{441}$祂的能力是何等的豪大.
17節求我們主耶穌基督的神.
榮耀的父將那種賜人智慧和啟示的靈.
使你們認識祂,真正認識祂.
我們如何才能真正認識神.
每次我們讀聖經的時候.
我們求神打開我們的心.
求聖靈幫助我們.
做我們個人的導師.
讓我們看到一些我們從來未明白的東西.
讓我們認識到神經過所有的作者.
讓我們傳遞的訊息究竟是什麼.
我們需要那種智慧.
不是知識,我們需要那種智慧.
我們需要那種啟示.
啟示是Revelation,啟示.
神的靈在我們裡面讓啟示給我們看.
神的能力是何等的豪大.
讓我們在這個思考的過程裡面.
我們真真正正能夠認識我們的神.
我們的主耶穌基督.
我們需要問自己不同的問題.
我們是否知道神呼召你的目的.
我們有沒有更深去認識神自己的性情.
你所付出的代價是何等的豪大.
記得在新約裡面,馬可福音講到.
一個被鬼附的人住在格拉森的墳墓裡面.
人們用鎖鏈鎖都鎖不住.
他會爭奪那些鎖鏈.
大聲叫喊.
耶穌將那個鬼趕了出來的時候.
他變成一個正常人.
開始穿衣服.
以前是赤身露體的.
你不要跟著我.
你不需要跟著我.
你回去告訴你所認識的人.
神在你身上做了何等大的事.
所以我們是被神召來.
我們能夠成為神的兒女.

$^{481}$同樣在同時間,神給我們有一個任務.
是需要我們自己去回應.
是需要我們留意.
是需要我們接納那個任務是屬於你的.
神擴大的能力表現在哪裡呢.
我們很熟悉那個故事.
拉薩路,馬大,瑪利亞的弟弟.
拉薩路死了三天.
耶穌去到已經是第四天.
在約翰福音上一章.
復活是神向人表達最大的能力.
因為祂可以叫一個死到發臭的人來復活.
神同樣的大能在我們今天的人身上.
無論那個人在我們人的眼中.
在道德上,行為上是多麼的敗壞.
神仍然可以靠著祂的大能.
將他從死裡面帶出來.
這個就是神自己的大能.
神第二樣東西是甚麼呢.
就是耶穌基督從死裡復活.
是神彰顯祂自己的大能.
沒有人能夠做到這件事.
我已經講完第一章.
很快出來了.
我們還有很多時間.
我們看看第二章的時候.
我給一個標題是甚麼呢.
我給一個標題是.
這個是起死回生的大愛.
第一點似乎很模糊.
但是當我解釋的時候.
大家就會明白.
寫著你們我們.
當我們翻開第二章的時候.
我們很容易就混亂了.
本來應該說你們.
應該說我們.
應該說你們.
究竟在搞甚麼呢.
究竟你們是你們呢.

$^{521}$還是我們是你們呢.
還是我們是我們呢.
我們要搞清楚這件事.
你們呢.
當保羅寫了一封信的時候.
你們是指一些外邦人.
但是我們那時候.
開始的時候.
我們那時候是說一些猶太人.
第二章開始第一節.
從前你們因著自己的過犯.
罪惡死了.
那時你們在過犯罪惡中生活.
隨著今世的風俗.
信服空中掌權者的領袖.
就是今天.
現今在博奕的人.
心中運行的邪力.
這個是你們.
說到一些外邦人犯罪.
做很多不好的事.
壞事.
但第三節他就改變了那個名詞.
他說甚麼呢.
我們從前也在生活中.
在他們當中放縱肉體的私慾.
隨著肉體和心中的意念去.
和別人一樣生來就該受罰.
所以如果你們是外邦人.
我們是猶太人.
他們在這裡有沒有分別.
沒有分別.
第四節說到.
然而神有豐富的憐憫.
因著祂愛我們的大愛.
竟然在我們因著過犯而死的時候.
令到我們與基督一同活割.
這裡又轉做我們.
沒有分別.
這個我們也說到是猶太人.

$^{561}$神給猶太人一個機會來相信.
大家相信關鍵的人物是甚麼.
是與基督一同活割.
猶太人不能夠因為摩西的律法.
遵循律法有好行為.
能夠成為一個活的人.
能夠親近神.
只有他們相信耶穌基督.
才能夠一同活割.
第八節說到.
你們得救由本父恩.
大家很熟悉這段經文.
也因著信者並不是出於自己.
而是神所賜的.
即是你們我們.
但我覺得在第十節的時候.
可能這個我們有少許分別.
我們是他們所造之物.
在耶穌基督裡創造的.
為要使我們行善.
就是神已經預備我們應該走的.
每一個在耶穌基督裡被造的人.
都是一個新的人.
都是一個熟神的人.
並沒有分別.
所以這是一個轉捩點.
是一個關鍵詞.
以前無論如何.
當耶穌基督來的時候.
我們只能夠靠著相信耶穌基督.
我們只能夠在耶穌基督裡面.
我們才能夠得到新生命.
我們在第十一節到第十四節.
他們講到這兩種人的分別在哪裡.
猶太人是有優先權.
有機會成為神的選民.
唯一的是他們是亞伯拉罕的子孫.
因為亞伯拉罕相信神的應許.
臨到以色列民.
外邦人沒有這個機會.

$^{601}$他們不知道.
不認識神.
不知道神對每一個人都有一個要求.
但是現在是怎樣呢.
第十三節講到.
從前你們是遠離神的人.
你們是外邦人.
如今卻在耶穌基督裡.
靠著他的血得以親近.
保羅為何要講這麼多.
他要表達一件事.
無論你是猶太人還是外邦人.
只有一個條件.
能夠令到猶太人和外邦人.
都經過同一件事.
第十六節.
記載十字架上消滅了淵神.
就藉著十字架使雙方歸為一體.
只有一個方向.
無論你是甚麼人 甚麼種 甚麼國籍.
只有一個方法能夠達到進到神的國.
就是通過主耶穌基督.
十九節.
讓你們不再是外邦人.
你們不是外邦人.
不是外面的人.
而是與聖徒同國.
同是神家裡面的人.
接著他講到.
用一個房子的樣式.
讓我們明白這件事.
基督耶穌是房閣石.
使徒才是建築物的根基.
我們每一個信徒.
要連接緊湊.
漸漸成為神所住的聖殿.
神的殿完成了嗎?.
在經文裡面.
他用的字「漸漸」.
仍然是建造中.

$^{641}$但重要的東西都存在.
基督是我們的房閣石.
以前的使徒.
以前的先知.
他們留下的說話.
我們今天都有機會看到.
問題是我們今天如何建造神自己的殿.
這個殿是神藉著聖靈所居住的地方.
你,我如何建造這個聖殿.
我們不能只坐在這裡.
我們信了耶穌.
神的計劃不是要我們坐在這裡.
甚麼都不做.
我們要成為這個殿的一部份.
我們要成為神藉著聖靈居住的地方.
你不能說無所謂了.
我沒甚麼能力.
讓某某人做這件事.
我們每個人都可以做一些事.
可以成為這個聖殿的一部份.
亦都給我有機會去成長.
看到弟兄姊妹那種愛心.
多謝Tony送我們去機場.
我們每次都很驚奇.
會有麻煩.
會很細心地包一包東西給我們上飛機.
讓我們可以吃.
剛才我咳嗽.
Amy會給一些喉糖.
希望我不要咳得太多.
我們在整個殿建造的過程.
我們需要那種磨合和連接.
那種相交.
才能夠神的連接緊湊.
不然的話.
一大塊石頭.
地震一會兒.
整座都會倒在地上.
我們需要學習那種磨合.
需要那種相交.

$^{681}$需要那種彼此相愛的活動.
才能夠建立神自己的殿.
另外一件事我希望提一提.
使徒和先知的根基.
我們的建造不是建造在空白的地方.
我們是建造在使徒和先知的基礎上.
是我們所讀的聖經.
福音書.
四本福音書.
使徒所記錄的.
很多的先知書.
他們講述了神的奧秘.
神的訊息 神的福音.
讓我們今天能夠體驗到.
明白神是多麼愛我們的一個神.
我們已經可以跳到第三章.
我想今天我可以按時完成今天的講道.
第三章我給一個標題.
奧秘 第一部分.
我為什麼要服侍你啊 外邦人.
對不起 可能用一個比較通俗的標題來寫.
在第三章的第一部分.
從第三章的第一節到第十三節.
常常出現一個名詞.
就是奧秘 奧秘 奧秘 奧秘.
究竟奧秘是什麼呢.
奧秘這個字不是說很神奇.
普通人不會明白的東西.
只有某某人才會明白.
我們不是說奧秘這兩個字.
在這段經文裡面不是這個意思.
奧秘只是說到有一些東西.
很珍貴的東西.
以前是隱藏 現在是被發明出來.
這段經文寫得很清楚.
第六節 外邦人在基督裡.
藉著福音得以同為後製.
同為一體 同為同為盟.
可能我打錯了.
同為盟 應許的人.

$^{721}$這個就是奧秘.
以前外邦人沒有機會接觸到神.
沒有機會明白神 相信神.
但今天 在神自己安定的時間裡.
我們今天能夠知道這個事情.
而且保羅說他自己是奧秘的管家.
是奧秘的福音的獨譯.
這個也是神賜給他的恩典 恩賜.
我們很難明白一個.
在地球的一個人.
在一個黑暗的地方.
他所寫的 所表達的東西.
他所表達神的恩典.
是有一種 另外一種層次在這裡.
第八節 雖然我給聖眾聖徒.
都是最少的還少.
他仍然賜給我這個恩惠.
就是傳福音給外邦人的一個恩惠.
讓我張爺基督那種不能測透的風聲.
傳給外邦人.
一個被囚的人.
當他想起福音的時候.
他不是說 為什麼我因為福音的時候.
會被坐牢.
他所體驗到這個福音是豐富的.
是豐富的.
是值得傳給外邦人聽的.
因為這個而且是一個神.
在第十一節.
這個是照著神在我們主耶穌基督裡.
所完成 永恆的計劃.
所以無論是猶太人或者外邦人.
這個是神自己計劃.
是一個永恆的計劃.
是需要他自己在基督裡完成的一個計劃.
十二節.
我們因信耶穌.
就在他裡面放膽無懼.
滿有自信的來到神的面前.
所以我求你 第十三節.

$^{761}$不要因為我所受的患難喪膽.
這懸是你們的光榮.
這個在第十三節是接續.
待會第三章的下半節.
不要因為我為福音的緣故坐牢.
你們就感到喪膽.
我們會的.
當我們看得起某某講員 某某牧師的時候.
我們都有一個期待在.
當那個期待被破滅的時候.
我覺得神為什麼會這樣.
這麼好的一個牧師.
這麼好的傳道人.
為什麼會因為某愛在的原因被監禁.
我們都會喪膽.
保祿說你們不要喪膽.
你看到這些事情不要喪膽.
我們都知道.
如果我們看回中國基督教歷史的時候.
王明道被監禁了很多年.
我自己的經歷是什麼呢.
我自己的經歷就是.
我在中國工作那幾年.
有機會在鄭州去一些家庭教會.
那天不知道為什麼.
可能是那個星期我出差.
回來為什麼沒有教會聚會呢.
應該是家庭教會.
被公安局封了.
他們很知道聚會的是誰.
誰是長老 誰是執事.
誰是牧師 誰是牧師娘.
全部抓了.
牧師娘被監禁了一陣子.
很快就放出來.
其他人被監禁的時間就比較長.
牧師被監禁了好像是多過一年.
我們都不明白.
在這樣的環境.
愛主 為主 傳道.

$^{801}$冒著很多不同的危險.
神為什麼這樣來待他.
後來過了一年多.
他出來的時候.
聽見弟兄姊妹.
我沒有直接去見他.
但是聽見弟兄姊妹說.
他出來瘦了.
是一件好事.
瘦了.
他說在監獄裡面.
他也帶了一些人來信耶穌.
這就是福音的大能.
我們從人的角度會喪膽.
遇到事情的時候.
會很灰心.
為什麼神會這樣對待.
為什麼神會容許事情發生.
但是在神的計劃 神的角度裡面.
有祂自己的美意.
無論你在哪裡 有修葛的時候.
這是永恆的計劃.
所以在第三章的下半節.
從十四節到尾 二十一節.
是另外一個禱告.
還沒有進入第四章之前的一個禱告.
這個禱告是什麼呢.
要希望當時已忽所的信徒.
不會因為他所受的患難而喪膽.
所以第一件事.
他求神的靈在你的內心裡面.
令你能夠剛強起來.
有時候我們走不掉的.
有時候遇到一些事情的時候.
我們就會很自然地灰心 退縮.
甚至我們到了一個很低的低潮.
但是保羅說.
我希望神藉著祂的靈.
將你內心那種力量剛強起來.
這是我們可以學習如何禱告的一件事.

$^{841}$我遇到一件事的時候.
不要尋找太多的原因.
我們第一件事就是說.
在這個情況底下的人.
我們如何為他們祈求.
讓神的靈來支持他們.
不是我們.
不僅是人.
人的力量到了某個程度的時候就沒有用了.
第二件事是說到.
基督在我們裡面生根.
第十七節.
又要使基督因著你們的信住在你們的心裡面.
使你們既在愛中生根.
就是說神在這個情況底下.
有沒有在你裡面.
因為你相信祂的信心.
有一個生根.
我很喜歡生根.
因為我很喜歡種東西.
就是有一個活性的東西.
我們不是扔一塊石頭在那裡.
而是一個活的東西.
基督可以因著我們的信心.
有一個動態的東西.
表達了生命的東西.
是能夠生長的東西.
根不生上面的東西就長得不好.
我們要讓基督因為我們的信心.
不要因為那個時勢.
無論是經濟還是政治影響到我們.
我們要相信神在所有的時代.
所有的角度裡面.
祂有自己的美意.
我們要在這段動盪的時間.
學習如何在神的耶穌基督裡面.
因為那個信心更加扎實.
「與眾聖徒一同明白神的愛是何等長闊高深.
並且知道這愛是超過人知識所能測讀的.
為要使你們充滿神一切的豐盛」.

$^{881}$我記得在主學裡也常常提到這件事.
我突然發覺那時候.
讀這段經文太熟悉了.
我們經常祈禱都會說.
我們要明白神的長闊高深.
但我們發覺這裡保羅特別有兩個字.
有四個字不是兩個字.
有四個字我們要注意.
他不是說你們要明白神的愛是何等長闊高深.
他這裡說你們能夠和眾聖徒一同.
這個分別就很大了.
他們所要求的不是我們個人如何去盡心.
好像一個屬靈生命的盡心.
而是說你們能夠明白神的長闊高深.
你們要一體來明白.
或者我嘗試來解釋這件事.
或者我們當中有些弟兄姊妹病了很久.
得到神的醫治.
他的見證就成為我們明白神的愛是如何表達出來.
是不是?.
可能我們當中也有其他的人.
為某人祈禱了很久很久.
那個人一直堅持不相信.
終於有一天開竅了.
願意接受耶穌基督.
這就是我們看到神的愛的長度.
神沒有放棄我們為他祈禱的人.
是不是?.
或者我們見到某個人真的到了井底的地步.
吸毒,嫖賭,飲蕩吹.
他什麼都做了.
他的父母為他祈禱.
終於有一天見到這個人從井底浮起來.
成為一個新族的人.
我們就看到神愛的深度在那裡.
是不是?.
所以我們需要頂照.
我們不能夠只是以為我們自己盡心.
多讀幾科的神學課程.
我們要成為一個更加有深度的基督徒.

$^{921}$我們需要在我們頂子妹那種相交裡面.
在神自己的話語裡面.
在神聖靈的痛在底下.
我們才能夠明白神自己的長闊高深.
我們剛才也提到神的大能.
特別是講到神自己復活的大能.
神自己復活的大能.
能夠改變一切.
我們眼前所看到的.
可能是我們的親戚.
可能是我們的子女.
他們離開神.
他們遠離神.
但我們必須相信一件事.
就是神自己的大能.
能夠將他們挽留回來.
能夠將他們帶回到我們堂中.
我們要相信那種復活的大能.
最後可能是一種保羅對神的歌頌.
但是他說.
願他在教會中並在基督耶穌裡.
得著榮耀直到世世代代.
永永遠遠.
這可能是一個頌讚.
同時表達了我們如何回應愛我們的神.
可能我們今天的講道時間長了一點.
我最後一個頑童篇.
是希望大家一同有思想的問題.
你需不需要重新檢討自己的信仰.
我們不要停留在初信的地方.
我們是否問自己一個問題.
我們是否真真信耶穌.
如果我們真的信耶穌.
我們是否坐在那裡甚麼都不做.
教會不是牧師的.
教會不是執事的.
教會不是長老的.
是我們每一個人.
我們問我們生活的問題.
我們是否活在基督裡.

$^{961}$活在基督裡是一個甚麼概念.
我們是否明白神的復活大能.
我們是否真真體驗神自己的大愛.
活在低頭的基督.
最後多謝你給我們這份所書.
最後給我們靠聖靈引導我們重新思考我們的信仰.
讓你自己在我們生命裡得到當得的榮耀.
奉主耶穌的名禱.
阿們.
\newpage



\section{帖撒羅尼迦前書 1:1-5:28}
\label{sec:032037}
\textbf{帖撒羅尼迦書 (1) 主快再來}~~~~劉少康牧師
\newline
\newline
連結: \href{http://www.church.com.hk/acms/content.asp?site=cdc\&op=show\&type=product\&code=032037\&layout=sermon}{\texttt{032037}} ~~~~ 語音日期: 2026-04-15
\newline
\newline
\hyperref[sec:032036]{\small{< < < PREV SERMON < < <}}
~
\hyperref[sec:index_chronic]{\small{[返順時目]}}
~
\hyperref[sec:032038]{\small{> > > NEXT SERMON > > >}}
\newline
\newline
帖撒羅尼迦前書 1:1-5:28
\newline
\begin{longtable}{cl}
\hline
\hline
章節 & 經文 (和合本修訂版)\\
\hline
1:1 & \begin{tabularx}{0.7\textwidth}{X} 保羅、西拉、提摩太寫信給帖撒羅尼迦在父神和主耶穌基督裡的教會。願恩惠、平安歸給你們！ \end{tabularx} \\ \\ \relax
1:2 & \begin{tabularx}{0.7\textwidth}{X} 我們為你們眾人常常感謝神，禱告的時候提到你們， \end{tabularx} \\ \\ \relax
1:3 & \begin{tabularx}{0.7\textwidth}{X} 在我們的父神面前，不住地記念你們因信心所做的工作，因愛心所受的勞苦，因盼望我們主耶穌基督所存的堅忍。 \end{tabularx} \\ \\ \relax
1:4 & \begin{tabularx}{0.7\textwidth}{X} 神所愛的弟兄啊，我知道你們是蒙揀選的； \end{tabularx} \\ \\ \relax
1:5 & \begin{tabularx}{0.7\textwidth}{X} 因為我們的福音傳到你們那裡，不僅在言語，也在能力，也在聖靈和充足的確信。你們知道，我們在你們那裡，為你們的緣故是怎樣為人。 \end{tabularx} \\ \\ \relax
1:6 & \begin{tabularx}{0.7\textwidth}{X} 你們成為效法我們，更效法主的人，因聖靈所激發的喜樂，在大患難中領受了真道， \end{tabularx} \\ \\ \relax
1:7 & \begin{tabularx}{0.7\textwidth}{X} 從此你們作了馬其頓和亞該亞所有信主的人的榜樣。 \end{tabularx} \\ \\ \relax
1:8 & \begin{tabularx}{0.7\textwidth}{X} 因為主的道已經從你們那裡傳播出去，你們向神的信心不只在馬其頓和亞該亞，就是在各處也都傳開了，所以不用我們說甚麼話。 \end{tabularx} \\ \\ \relax
1:9 & \begin{tabularx}{0.7\textwidth}{X} 因為他們自己已經傳講我們是怎樣進到你們那裡，你們是怎樣離棄偶像，歸向神來服侍那又真又活的神， \end{tabularx} \\ \\ \relax
1:10 & \begin{tabularx}{0.7\textwidth}{X} 等候他兒子從天降臨，就是神使他從死人中復活的那位救我們脫離將來憤怒的耶穌。 \end{tabularx} \\ \\ \relax
2:1 & \begin{tabularx}{0.7\textwidth}{X} 弟兄們，你們自己知道我們來到你們那裡並不是徒然的。 \end{tabularx} \\ \\ \relax
2:2 & \begin{tabularx}{0.7\textwidth}{X} 我們從前在腓立比蒙難受辱，這是你們知道的，可是我們還是靠著神給我們的勇氣，在強烈反對中把神的福音傳給你們。 \end{tabularx} \\ \\ \relax
2:3 & \begin{tabularx}{0.7\textwidth}{X} 我們的勸勉不是出於錯誤，也不是出於污穢，也不是用詭詐。 \end{tabularx} \\ \\ \relax
2:4 & \begin{tabularx}{0.7\textwidth}{X} 但神既然認定我們經得起考驗，把福音託付我們，我們就照著傳講，不是要討人喜歡，而是要討那考驗我們的心的神喜歡。 \end{tabularx} \\ \\ \relax
2:5 & \begin{tabularx}{0.7\textwidth}{X} 因為我們從來沒有用過諂媚的話，這是你們知道的，也沒有藏著貪心，這是神可以作證的。 \end{tabularx} \\ \\ \relax
2:6 & \begin{tabularx}{0.7\textwidth}{X} 我們作為基督的使徒，雖然可以受人尊重，卻沒有向你們或向別人求榮耀， \end{tabularx} \\ \\ \relax
2:7 & \begin{tabularx}{0.7\textwidth}{X} 反而在你們當中心存溫柔，如同母親哺乳自己的孩子。 \end{tabularx} \\ \\ \relax
2:8 & \begin{tabularx}{0.7\textwidth}{X} 既然我們這樣愛你們，不但樂意將神的福音給你們，連自己的性命也樂意給你們，因為你們是我們所疼愛的。 \end{tabularx} \\ \\ \relax
2:9 & \begin{tabularx}{0.7\textwidth}{X} 弟兄們，你們記念我們的辛苦勞碌，晝夜做工，傳神的福音給你們，免得你們任何人受累。 \end{tabularx} \\ \\ \relax
2:10 & \begin{tabularx}{0.7\textwidth}{X} 我們對你們信主的人是何等聖潔、正直、無可指責，這有你們作證，也有神作證。 \end{tabularx} \\ \\ \relax
2:11 & \begin{tabularx}{0.7\textwidth}{X} 正如你們知道，我們待你們好像父親待自己的兒女一樣。 \end{tabularx} \\ \\ \relax
2:12 & \begin{tabularx}{0.7\textwidth}{X} 我們勸勉你們，安慰你們，囑咐你們，使你們行事對得起那召你們進他自己的國、得他榮耀的神。 \end{tabularx} \\ \\ \relax
2:13 & \begin{tabularx}{0.7\textwidth}{X} 為此，我們也不斷地感謝神，因為你們聽見我們所傳神的道的時候，你們領受了，不以為這是人的道，而以為這確實是神的道，而且在你們信主的人當中運行著。 \end{tabularx} \\ \\ \relax
2:14 & \begin{tabularx}{0.7\textwidth}{X} 弟兄們，你們與猶太地區神的各教會，就是在基督耶穌裡的各教會，有同樣的遭遇，因為你們也受了同胞的迫害，像他們受了猶太人的迫害一樣。 \end{tabularx} \\ \\ \relax
2:15 & \begin{tabularx}{0.7\textwidth}{X} 這些猶太人不但殺了主耶穌和先知們，又把我們趕出去。他們令神不悅，且與眾人為敵， \end{tabularx} \\ \\ \relax
2:16 & \begin{tabularx}{0.7\textwidth}{X} 阻撓我們傳道給外邦人，使他們得救，以致常常惡貫滿盈，但神的憤怒終於臨到他們身上。 \end{tabularx} \\ \\ \relax
2:17 & \begin{tabularx}{0.7\textwidth}{X} 弟兄們，我們被迫暫時與你們分離，身體離開，心卻沒有；我們極力想法子，渴望見你們的面。 \end{tabularx} \\ \\ \relax
2:18 & \begin{tabularx}{0.7\textwidth}{X} 所以我們很想到你們那裡去。我－保羅有一兩次要去，只是撒但阻擋了我們。 \end{tabularx} \\ \\ \relax
2:19 & \begin{tabularx}{0.7\textwidth}{X} 當我們的主耶穌再來，我們站在他面前的時候，我們的盼望、喜樂和所誇的冠冕是甚麼呢？不正是你們嗎？ \end{tabularx} \\ \\ \relax
2:20 & \begin{tabularx}{0.7\textwidth}{X} 你們就是我們的榮耀和喜樂！ \end{tabularx} \\ \\ \relax
3:1 & \begin{tabularx}{0.7\textwidth}{X} 既然我們不能再忍，就決定獨自留在雅典， \end{tabularx} \\ \\ \relax
3:2 & \begin{tabularx}{0.7\textwidth}{X} 於是差派我們在基督福音上作神同工的弟兄提摩太前去，在你們所信的道上堅固你們，勸勉你們， \end{tabularx} \\ \\ \relax
3:3 & \begin{tabularx}{0.7\textwidth}{X} 免得有人被這些患難動搖。因為你們自己知道，我們受患難原是命定的。 \end{tabularx} \\ \\ \relax
3:4 & \begin{tabularx}{0.7\textwidth}{X} 我們在你們那裡的時候，曾預先告訴你們，我們必受患難；你們知道，這果然發生了。 \end{tabularx} \\ \\ \relax
3:5 & \begin{tabularx}{0.7\textwidth}{X} 為此，既然我不能再忍，就差派人去，要知道你們的信心如何，恐怕那誘惑人的果真誘惑了你們，以致我們的勞苦歸於徒然。 \end{tabularx} \\ \\ \relax
3:6 & \begin{tabularx}{0.7\textwidth}{X} 但是，提摩太剛從你們那裡回來，將你們信心和愛心的好消息報給我們，又說你們常常記念我們，切切想見我們，如同我們想見你們一樣。 \end{tabularx} \\ \\ \relax
3:7 & \begin{tabularx}{0.7\textwidth}{X} 所以，弟兄們，我們在一切困苦患難中，因著你們的信心得到鼓勵。 \end{tabularx} \\ \\ \relax
3:8 & \begin{tabularx}{0.7\textwidth}{X} 如今你們若靠主站立得穩，我們就得生了。 \end{tabularx} \\ \\ \relax
3:9 & \begin{tabularx}{0.7\textwidth}{X} 我們在神面前，因著你們滿有喜樂。為這一切喜樂，我們能用怎樣的感謝為你們報答神呢？ \end{tabularx} \\ \\ \relax
3:10 & \begin{tabularx}{0.7\textwidth}{X} 我們晝夜切切祈求要見你們的面，來補足你們信心的不足。 \end{tabularx} \\ \\ \relax
3:11 & \begin{tabularx}{0.7\textwidth}{X} 願我們的父神自己和我們的主耶穌，為我們開路到你們那裡去。 \end{tabularx} \\ \\ \relax
3:12 & \begin{tabularx}{0.7\textwidth}{X} 又願主使你們彼此相愛的心，和愛眾人的心，都能增長，充足，如同我們愛你們一樣， \end{tabularx} \\ \\ \relax
3:13 & \begin{tabularx}{0.7\textwidth}{X} 好堅固你們的心，使你們在我們的主耶穌同他眾聖徒來臨的時候，在我們父神面前，成為聖潔，無可指責。阿們！ \end{tabularx} \\ \\ \relax
4:1 & \begin{tabularx}{0.7\textwidth}{X} 末了，弟兄們，我們靠著主耶穌求你們，勸你們，既然你們領受了我們的教導，知道該怎樣行事為人，討神的喜悅，其實你們也正這樣行，我勸你們要更加努力。 \end{tabularx} \\ \\ \relax
4:2 & \begin{tabularx}{0.7\textwidth}{X} 你們原知道，我們憑主耶穌傳給你們甚麼命令。 \end{tabularx} \\ \\ \relax
4:3 & \begin{tabularx}{0.7\textwidth}{X} 神的旨意就是要你們成為聖潔，遠避淫行； \end{tabularx} \\ \\ \relax
4:4 & \begin{tabularx}{0.7\textwidth}{X} 要你們各人知道怎樣用聖潔、尊貴控制自己的身體 ， \end{tabularx} \\ \\ \relax
4:5 & \begin{tabularx}{0.7\textwidth}{X} 不放縱私慾的邪情，像不認識神的外邦人。 \end{tabularx} \\ \\ \relax
4:6 & \begin{tabularx}{0.7\textwidth}{X} 不准有人在這事上越軌，佔他弟兄的便宜；因為這一類的事，主必報應，正如我預先對你們說過，又切切警告過你們的。 \end{tabularx} \\ \\ \relax
4:7 & \begin{tabularx}{0.7\textwidth}{X} 神召我們本不是要我們沾染污穢，而是要我們聖潔。 \end{tabularx} \\ \\ \relax
4:8 & \begin{tabularx}{0.7\textwidth}{X} 所以，那棄絕這教導的不是棄絕人，而是棄絕那把自己的聖靈賜給你們的神。 \end{tabularx} \\ \\ \relax
4:9 & \begin{tabularx}{0.7\textwidth}{X} 有關弟兄間的手足之情，不用人寫信給你們，因為你們自己蒙了神的教導要彼此相愛。 \end{tabularx} \\ \\ \relax
4:10 & \begin{tabularx}{0.7\textwidth}{X} 你們向全馬其頓的眾弟兄固然是這樣行，但我勸弟兄們要更加努力。 \end{tabularx} \\ \\ \relax
4:11 & \begin{tabularx}{0.7\textwidth}{X} 要立志過安靜的生活，管自己的事，親手做工，正如我們從前吩咐你們的， \end{tabularx} \\ \\ \relax
4:12 & \begin{tabularx}{0.7\textwidth}{X} 好使你們的行為能得外人的尊敬，同時也不依賴任何人。 \end{tabularx} \\ \\ \relax
4:13 & \begin{tabularx}{0.7\textwidth}{X} 弟兄們，至於已睡了的人，我們不願意你們不知道，恐怕你們憂傷，像那些沒有指望的人一樣。 \end{tabularx} \\ \\ \relax
4:14 & \begin{tabularx}{0.7\textwidth}{X} 既然我們信耶穌死了，復活了，那些已經在耶穌裡睡了的人，神也必將他們與耶穌一同帶來。 \end{tabularx} \\ \\ \relax
4:15 & \begin{tabularx}{0.7\textwidth}{X} 我們照主的話告訴你們一件事：我們這活著還存留到主來臨的人，絕不會在那已經睡了的人之先。 \end{tabularx} \\ \\ \relax
4:16 & \begin{tabularx}{0.7\textwidth}{X} 因為，召集令一發，天使長的呼聲一叫，神的號角一吹，主必親自從天降臨；那在基督裡死了的人必先復活， \end{tabularx} \\ \\ \relax
4:17 & \begin{tabularx}{0.7\textwidth}{X} 然後我們這些活著還存留的人必和他們一同被提到雲裡，在空中與主相會。這樣，我們就要和主永遠同在。 \end{tabularx} \\ \\ \relax
4:18 & \begin{tabularx}{0.7\textwidth}{X} 所以，你們當用這些話彼此勸勉。 \end{tabularx} \\ \\ \relax
5:1 & \begin{tabularx}{0.7\textwidth}{X} 弟兄們，關於那時候和日期，不用人寫信給你們， \end{tabularx} \\ \\ \relax
5:2 & \begin{tabularx}{0.7\textwidth}{X} 因為你們自己明明知道，主的日子來到會像賊在夜間突然來到一樣。 \end{tabularx} \\ \\ \relax
5:3 & \begin{tabularx}{0.7\textwidth}{X} 人正說平安穩定的時候，災禍忽然臨到他們，如同陣痛臨到懷胎的婦人一樣，他們絕逃脫不了。 \end{tabularx} \\ \\ \relax
5:4 & \begin{tabularx}{0.7\textwidth}{X} 弟兄們，你們並不在黑暗裡，那日子不會像賊一樣臨到你們。 \end{tabularx} \\ \\ \relax
5:5 & \begin{tabularx}{0.7\textwidth}{X} 你們都是光明之子，都是白晝之子；我們不屬黑夜，也不屬幽暗。 \end{tabularx} \\ \\ \relax
5:6 & \begin{tabularx}{0.7\textwidth}{X} 所以，我們不要沉睡，像別人一樣，總要警醒謹慎。 \end{tabularx} \\ \\ \relax
5:7 & \begin{tabularx}{0.7\textwidth}{X} 因為睡了的人是在夜間睡，醉了的人是在夜間醉。 \end{tabularx} \\ \\ \relax
5:8 & \begin{tabularx}{0.7\textwidth}{X} 但既然我們屬於白晝，就應當謹慎，把信和愛當作護心鏡遮胸，把得救的盼望當作頭盔戴上。 \end{tabularx} \\ \\ \relax
5:9 & \begin{tabularx}{0.7\textwidth}{X} 因為神不是預定我們受懲罰，而是預定我們藉著我們的主耶穌基督得救。 \end{tabularx} \\ \\ \relax
5:10 & \begin{tabularx}{0.7\textwidth}{X} 他替我們死，讓我們無論醒著、睡著，都與他同活。 \end{tabularx} \\ \\ \relax
5:11 & \begin{tabularx}{0.7\textwidth}{X} 所以，你們該彼此勸勉，互相造就，正如你們素常做的。 \end{tabularx} \\ \\ \relax
5:12 & \begin{tabularx}{0.7\textwidth}{X} 弟兄們，我們勸你們要敬重那些在你們中間勞苦的，就是在主裡面督導你們、勸戒你們的人。 \end{tabularx} \\ \\ \relax
5:13 & \begin{tabularx}{0.7\textwidth}{X} 又因他們所做的工作，要以愛心格外尊重他們。你們也要彼此和睦。 \end{tabularx} \\ \\ \relax
5:14 & \begin{tabularx}{0.7\textwidth}{X} 弟兄們，我們勸你們，要警戒不守規矩的人，勉勵灰心的人，扶助軟弱的人，對眾人要有耐心。 \end{tabularx} \\ \\ \relax
5:15 & \begin{tabularx}{0.7\textwidth}{X} 你們要謹慎，無論是誰都不要以惡報惡，彼此間和對眾人都要追求做好事。 \end{tabularx} \\ \\ \relax
5:16 & \begin{tabularx}{0.7\textwidth}{X} 要常常喜樂， \end{tabularx} \\ \\ \relax
5:17 & \begin{tabularx}{0.7\textwidth}{X} 不住地禱告， \end{tabularx} \\ \\ \relax
5:18 & \begin{tabularx}{0.7\textwidth}{X} 凡事謝恩，因為這是神在基督耶穌裡向你們所定的旨意。 \end{tabularx} \\ \\ \relax
5:19 & \begin{tabularx}{0.7\textwidth}{X} 不要熄滅聖靈； \end{tabularx} \\ \\ \relax
5:20 & \begin{tabularx}{0.7\textwidth}{X} 不要藐視先知的講論。 \end{tabularx} \\ \\ \relax
5:21 & \begin{tabularx}{0.7\textwidth}{X} 但凡事要察驗：美善的事要持守， \end{tabularx} \\ \\ \relax
5:22 & \begin{tabularx}{0.7\textwidth}{X} 各樣惡事要禁戒。 \end{tabularx} \\ \\ \relax
5:23 & \begin{tabularx}{0.7\textwidth}{X} 願賜平安的神親自使你們完全成聖！願你們的靈、魂、體得蒙保守，在我們的主耶穌基督來臨的時候，完全無可指責。 \end{tabularx} \\ \\ \relax
5:24 & \begin{tabularx}{0.7\textwidth}{X} 那召你們的本是信實的，他必成就這事。 \end{tabularx} \\ \\ \relax
5:25 & \begin{tabularx}{0.7\textwidth}{X} 弟兄們，請也為我們禱告。 \end{tabularx} \\ \\ \relax
5:26 & \begin{tabularx}{0.7\textwidth}{X} 用聖潔的吻向眾弟兄問安。 \end{tabularx} \\ \\ \relax
5:27 & \begin{tabularx}{0.7\textwidth}{X} 我指著主囑咐你們，要把這信宣讀給眾弟兄聽。 \end{tabularx} \\ \\ \relax
5:28 & \begin{tabularx}{0.7\textwidth}{X} 願我們的主耶穌基督的恩惠與你們同在！ \end{tabularx} \\ \\
[1ex]
\hline
\hline
\end{longtable}
$^{1}$各位電影姐妹平安.
很開心今天能夠跟大家分享.
大家都知道現在這個世界.
很多打仗的事正在發生.
美國和伊朗互相較量.
表面上好像美國有很多東西限制著伊朗.
封鎖這個海峽.
禁止他們擁有核武器.
而且將他們很多軍事設施.
全部都炸毀了.
但是伊朗卻仍然充滿信心.
要和美國這個強國來較量.
今天打仗的風聲.
在我們四周圍包圍著我們.
我們在想不到的時候.
主耶穌就要回來了.
所以在帖撒羅里加前書補下去就說.
主回來的日子.
是好像賊來一樣.
在你想不到的時候.
突然間就來了.
各位.
你有沒有想過主耶穌很快回來.
你預備好自己了沒有?.
所以帖撒羅里加前書提醒我們.
我們要警醒緊守.
醒啊.
不要再睡覺了.
我們要守望.
我們要站穩自己的崗位.
今天我會用一堂的時間.
介紹帖撒羅里加前書.
讓我首先介紹一下.
這間帖撒羅里加教會.
保羅大約在他第二次宣教行程的時候.
他就和西拉和提摩泰.
去了帖撒羅里加這個小城市.
他就在那裡工作.
根據《使徒行傳》的記載.
一年三個安息日.

$^{41}$他都去到當中.
向猶太人講道.
聖經說保羅很辛勤地工作.
結果就帶了一幫人信主.
保羅就很掛念這幫信徒.
就兼顧他們將一些很重要的道理教導他們.
然後保羅由於當時他被人反對.
因為當地的人就很不喜歡這個新道理.
因此起來反對的人就用很多的聲音.
慫恿那些人.
結果保羅就不得已去了雅典.
他在雅典等候的時候.
心裡面仍然惦記著這間小教會.
剛剛成立沒多久的教會.
於是他就派了提摩泰.
你不如回去探望一下帖撒羅里加.
看看他們情況怎麼樣.
結果提摩泰去了一段短時間.
就回來向保羅作報告.
他說這間教會他們信心堅固.
信徒們很好.
繼續有堅定的信心.
但是很可惜.
有一些人就傳了一些不合真理的假道理.
而且對保羅有些批評.
有些說他閒話.
結果保羅就很掛念他們.
於是就馬上寫這封信給他們.
而且表明他充滿喜樂.
原來你們惦記著我.
他心裡面很流露出那份喜樂.
因此在帖撒羅里加這本前書裡面.
有很多不同很豐富的表情流露出來.
帖撒羅里加前書一共有五章.
主題就說我們要成聖.
我們一定要離開道理的開端.
去到一個很穩固的地步.
而且我們要成為聖潔.
合乎神的心意.
我們要懂得遵行神的旨意而行.

$^{81}$帖撒羅里加前書一共有五章聖經.
第一章聖經就說到保羅為教會感恩.
常常發出讚美的使徒.
今天我們都應該學這種態度.
神啊,你真好,賜我生命氣息.
你賜我健康.
你賜我有呼吸.
你賜我有空氣.
特別在紐西蘭這裡.
空氣非常好.
外面充滿花草樹木.
很美的環境.
而且天亮的時候就看到太陽出來.
黃昏的時候就看到日落那種美景.
每一天我都讚美不盡.
神的恩,神的愛.
寫在他所創造的宇宙當中.
大家都知道美國最近發射了一個新的太空船.
環繞月球飛行.
是人類有史以來最長的飛行途徑.
原來在這四個太空人裡面.
有很多位都是基督徒.
我看過一個YouTube.
在說這些太空人.
怎樣在太空上面.
他們是直接感覺到神與他們同在.
其中一個太空人帶著他的聖經.
他說他在很寧靜的時候.
他好像見到神一樣.
神在寧靜中與他同在.
他感覺到神那種真實性.
他看見神怎樣將地球懸掛在太空當中.
然後這個地球怎樣環繞太陽來轉.
他看見神創造的奇妙.
他心裡面就讚嘆不已.
各位,原來我們很需要學會感恩.
因此保羅在《貼薩萊傑加傳說》第一章.
就教我們要感恩.
而第二章就講到.
他為一些事奉來申辯.

$^{121}$去辯護.
可能是那些貼薩萊加的信徒誤會了保羅.
提摩太鄉報告了.
所以他就澄清這些誤解.
今天在教會裡面.
特別是華人教會裡面.
很多是是非非.
就是我們發生了很多誤會.
因此我們必須澄清一些誤解.
要表明自己很清晰的立場.
然後到了第三章.
保羅就表明.
他很想再去探訪貼薩萊加這間教會.
他很想去探望他們.
而且他鼓勵他們要彼此相愛.
彼此相愛是今天教會很需要的一個訊息.
你去到教會做崇拜.
有沒有跟其他人打招呼呢?.
只是跟你認識的那些你才打招呼呢?.
其他坐在那裡的那些好像不關我的事.
好像三寶人一樣.
大家都是弟兄姊妹.
為什麼你不主動去介紹一下自己.
我是什麼名字?你好嗎?.
怎樣稱呼你?.
大家見過面?.
因為有時我們在教會裡面見到一些人.
熟口熟面的.
但是就不認識.
沒有聊過天.
因此就好像陌生人一樣.
既然我們坐在同一間禮堂.
一同來到敬拜.
讓我們拉近彼此之間的關係.
而且他說我們不單要彼此相愛.
我們還要有愛眾人的心.
我們要增長.
要充足.
那個愛眾人的心是要增加的.
是要很充足,很豐富的.

$^{161}$求主今天都幫助我們.
我們的愛心是不斷地成長.
不斷地增長.
到了第四章.
就是保羅對信徒行為的一些教導,指引.
他說我們需要按神的旨意來生活.
我們很需要成為聖潔.
避免淫亂的行為.
這是保羅給我們一個很好的提醒.
然後他就說到.
在末後的時候.
當號筒吹響的時候.
所有死人都要復活.
這個很震撼的訊息.
我記得很多時候在安葬禮的時候.
我都會讀這段經文.
當號筒末次吹響的時候.
主耶穌就回來.
所有死了的人都要從死裡復活.
然後他們就要接受主再來迎接著.
到了第五章.
就是有關教會生活的一些勉勵.
他說主回來的日子.
好像夜間的賊一樣.
你想不到的時候主耶穌就回來了.
所以各位我們必須要警醒.
警醒謹守自己.
謹守你的行為.
保守你的心智.
堅定站穩你的信仰.
而且我們不要消滅聖靈的感動.
聖靈真的很好.
祂時時在我們心裡面提醒我們.
叫我們知道我們自己怎麼走.
而且我們凡事要測驗神的旨意.
善美的你要去持守.
所以保羅實在是一個很有經驗的使徒.
他對教會裡面的弟兄姊妹.
那種呵護,那種關懷,那種愛心.
那種想他們千萬不要走錯.

$^{201}$他們一定要站立得穩,信心要堅固,要牢固.
然後他們的愛心增長.
在各方面成為土主喜悅.
遵行神旨意的人.
我作為一位的牧者.
我也很想對各位弟兄姊妹.
有這樣一個期望.
你是不是一個願意遵行神旨意的人.
你是不是實際上土主喜悅?.
善美的事你要持守.
堅定地去行.
讓我們在今天這個末後的日子.
我們更樂意持守主的道.
渴望我們能夠土主的喜悅.
願天父賜福給各位.
\newpage



\section{出埃及記 17:5}
\label{sec:032038}
\textbf{聖經人聖經事:耶和華尼西}~~~~黃克勤牧師博士
\newline
\newline
連結: \href{http://www.church.com.hk/acms/content.asp?site=cdc\&op=show\&type=product\&code=032038\&layout=sermon}{\texttt{032038}} ~~~~ 語音日期: 2026-04-20
\newline
\newline
\hyperref[sec:032037]{\small{< < < PREV SERMON < < <}}
~
\hyperref[sec:index_chronic]{\small{[返順時目]}}
~
\hyperref[sec:032039]{\small{> > > NEXT SERMON > > >}}
\newline
\newline
出埃及記 17:5
\newline
\begin{longtable}{cl}
\hline
\hline
章節 & 經文 (和合本修訂版)\\
\hline
17:5 & \begin{tabularx}{0.7\textwidth}{X} 耶和華對摩西說：「你帶著以色列的幾個長老，走在百姓前面，手裡拿著你先前擊打尼羅河的杖，去吧！ \end{tabularx} \\ \\
[1ex]
\hline
\hline
\end{longtable}
$^{1}$弟兄姊妹平安!大家好!.
今次是講耶和華利西.
你一定聽過了.
生命樹向大家提供機會.
一起讀聖經,多謝你的參與.
講聖經的人,聖經的故事.
看看上帝做些什麼.
耶和華利西是他們叫的名字.
看看是什麼事.
十七章出一及記,那時阿瑪利人在利菲定和以色列人爭戰.
摩西屯爵書講了.
你為我們選出人來,出去和阿瑪利人爭戰.
明天我手裡要拿著神的杖站在山頂上.
這裡是第一場仗.
以色列人喜歡第一場仗.
第一場是埃及人追它,但它們全都浸死在海裡.
或者河裡.
盧偉海.
這時正式有人來攻擊它,就是阿瑪利人.
阿瑪利人經常搞事.
阿瑪利人為何會在利菲定攻擊以色列人.
因為你知道,忽然有二百萬人這樣行.
全世界都知道,有一個很大的隊伍在行.
向著我們的方向來,你會很驚.
因為二百萬人一行過,你的地方就被它踩平.
如果它要攻擊你,你只有十多萬人,你不夠它來.
如果你的民族只有一千幾百人,三兩萬人,十萬八萬人.
嘩!兩百萬.
你只要看它走過,都可能看幾天.
如果它不是走得像過海般,你真的要看很久.
嘩!很厲害.
它們發覺,這些人的衣服又不錯.
很多穿著埃及式的衣服,內套也穿得出來.
而且有些金器,有些金飾,這些人有錢.
阿瑪利亞人大概想在這裡攻擊它拿到多少.
看它們不是很會打仗,沒有什麼軍隊.
衝上來的衛兵都不是很厲害,看它們的樣子.
沒有什麼章法.
阿瑪利亞人.
阿瑪利亞是山谷的居住者,即是一群住在山谷裡的民族.

$^{41}$來到利菲定休息的地方,去那裡搶它,跟它爭戰.
你知道,應該不是阿瑪利亞人第一次.
他們住在山谷裡面,何列山.
是一個很大的山,很多山谷.
他們住在山谷裡面的人,過來利菲定很近.
在那裡休息的人當中,搶東西,打它.
搶到人,搶到錢也好.
這個時候摩西就說.
你為我跟約書亞說,你找些人出來,出去打.
我拿著整個像站在山頂.
第一件事就是,如果你是約書亞你怎樣.
那不如調位,你出去打,我上山頂拿著一支像.
老人家又不是很能打,約書亞是個戰士.
他就說好,我帶人出去打.
其實約書亞是一個猶太民族或者以色列人.
他們很尊敬的一個戰士.
有幾個,約書亞,基澱,大衛和兩約之間的瑪加比.
都是他們很尊敬的打得的人,以往重要的戰士.
約書亞第一次打仗.
你出去跟阿瑪利亞人爭戰,我明天拿著神的像.
一見到摩西拿著神的像,大家就放心.
有神同在,我會上去山頂看你.
好像聖經所說,我在場,你放心出去.
在山頂上做什麼呢?再看下去.
於是約書亞就帶著摩西照著摩西對他說的話.
帶著人馬出去跟阿瑪利亞人爭戰.
這個時間有三個人,摩西,阿倫和胡爾上山頂.
胡爾是第一次作為一個領袖出現.
跟了一段路,一個新的領袖幫手出來.
胡爾應該是一個文人.
摩西是什麼意思呢?.
摩西就是拉上來,拉到他.
在水裡拉上來.
阿倫就是光明的人,有陽光的.
發光的意思.
胡爾就是山洞,一個洞.
這個名字真好笑.
很有可能他生的時間在那個洞裡.
也說不定,於是就改他的名字.
以前他們改名很有趣.

$^{81}$譬如有些人一個晴朗的日子出世.
可能就改做曉晴.
曉就是早上,晴朗的早上出世就是曉晴.
改這些名字,可能在那個洞出世就叫胡爾.
改他兒子的名字叫洞,山洞.
如果將來上課老師問.
小朋友你叫什麼名字?.
我叫洞,我叫洞,山洞,洞.
都挺有趣.
去到山頂,做什麼呢?.
去到山頂不是觀戰.
都是觀戰,不過去到山頂是祈禱.
摩西第二點要跟大家說的.
第一點,這已經是第二了.
其實我想說的第一點,忘記了highlight.
就是在你最安息的時候.
休息一下,又吃完鵪鶉,又在吃麻辣.
又沒水喝,現在又喝飽水.
這個時候在休息.
以為自己平安穩妥的時候就遭遇攻擊.
人生不是這樣一帆風順,是過山車的.
是上帝讓他一個過山車的人生.
令到我們時時都有一些危機感.
可以懂得依靠他,以及更多應變的能力.
是上帝特意的.
讓阿瑪利人攻擊他.
所以在你的利非定那裡,你就會出事.
人家就會來攻擊你.
不過這個時候有一個應對的辦法.
第一,自己落力出去打仗.
第二,舉手祈禱.
剛才上一節就說,約書一下出去打仗.
就在安書的時候遭遇問題.
好我要出力,找一個對的人自己承擔起.
第二,就是祈禱,何時舉手以色列人就得勝.
何時垂手阿瑪利人就得勝.
其實這個十七章十一字很重要.
摩西何時舉手禱告以色列人得勝.
手好累,一跌倒,何時阿瑪利人就得勝.
其實這是一個很特別的聖經例子.

$^{121}$是強調祈禱以靠神的能力不得勝.
這個例子,摩西何時舉手何時得勝.
這個舉手是什麼呢?.
到,我推測他的舉手就是.
神啊,你的大能要彰顯.
哎呀,他們現在正在打仗,很辛苦很危險.
神啊,可能你的手這樣.
其實在近東的人,他們多數的祈禱都是不閉眼的.
舉目望天,兩隻手是這樣的.
好像神啊,我真的需要一些東西.
因此在這裡,摩西何時舉手.
他的手是這樣舉,還是這樣舉.
都算舉,神啊,我真的很需要你的能力,我需要你的幫助.
可能雙目望天,兩隻手都是這樣的.
現在他們很多時候祈禱都是這樣的.
要嘛,當他們讚美的時候,手才是這樣的.
否則,他的手是這樣的.
神啊,很高.
這樣叫你舉一天,就算是這樣,這樣更加累.
這樣行不行啊,神啊,我們真的很需要你的能力.
很需要你的恩典.
他們打仗的人,很需要你的智慧.
主啊,求你給我們.
這樣都累的.
於是手一垂低,全避開.
見到約書亞他們在輸,馬利人得勝.
神啊,我們輸了,很危險,有人遭殃了,死了,受傷了.
求你醫治他,求你幫助我們,軍隊有力量.
又贏了,拉鋸了.
所以摩西何時舉手,以色列人得勝.
何時垂手,馬利人得勝.
這是一個很特殊的例子.
是告訴我們,祈禱會得勝.
第十二節,摩西的手累了.
不是吧?這樣的時間累?.
那麼摩西多少歲?.
他見異象的時候已經八十了.
再大衛出來,現在九十多歲了.
可能八十多九十歲了.
手真的不是很耐用了.

$^{161}$於是他們搬石頭來,讓摩西坐在上面.
搬石頭來放在他以下,然後他就坐在上面.
搬石頭來讓他坐在上面就可以了,這個寫法挺有趣的.
然後阿倫與胡爾第二樣就是扶著他的手.
要知道一個很困難的場面,很大的爭戰的時候.
一個人未必夠,因為你的手會軟.
一個人應付不了.
於是有阿倫與胡爾扶著他的手.
一個人扶一邊,一個人扶一邊,捉住他的手,托住他.
於是他的手就穩住,直到日落的時候了.
於是他的手一直企圖的時間,這樣就得勝了.
約書亞就得勝了,戰鬥一天之後.
就用刀殺了阿瑪力王和他的百姓.
約書亞對摩西說,你看看我今天做了這件事.
我要將阿瑪力的名號從天下全面塗抹.
阿瑪力的名號我將會將阿瑪力的人全部都塗抹.
即是說沒有人再認識他,即是說將他變得很弱或者絕了.
你要將這句話寫在書上做紀念.
唸給約書亞聽,即是說你要將這一次的祈禱得勝.
努力戰爭加上祈禱而得勝穩住手.
這樣的得勝,寫在書上給約書亞聽.
即是說已經開始有一個提示,有一個hint.
就是約書亞會承繼你,你用書傳給他,寫在書上給他聽.
這件事就是這樣,所以在這裡跟大家說.
何時舉手,何時得勝.
耶和華是支持我們的,用禱告來托住,看回今天的主題.
他們叫耶和華利西,是怎樣來的呢?.
摩西捉鳥一座壇,勝利了,起名叫耶和華利西.
又說耶和華已經起了勢,必世世代代和亞瑪力人爭先.
耶和華利西這個字在我們很多時候都會說的.
特別是讚美主,因為你是耶和華利西,很多時候有人這樣祈禱.
其實耶和華利西,重音在後面,耶和華這個不解釋了.
這個C在S,利西,利西,利西,利西是旗號的意思.
耶和華我的旗號,打仗,無千無萬,個個在背後看上去.
穿上盔甲都差不多樣,誰是司令都看不到,軟軟的一匹馬.
有人舉著旗,很多時候這套旗令我知道我的隊人是這支旗.
他的隊人是那支旗,於是我可以跟著這支旗去打仗.
耶和華就是我的旗號,他向前一揮,我就會跟著他.
他停,我就會停,他在這裡,我保護這支旗.
到最後我都要扶著這支旗,輸我都要扶著這支旗才行.

$^{201}$很多時候都是這樣的,這個旗號的意思就是我要忠貞的對象.
而且他是帶領我的,耶和華利西,很多時候有人說耶和華使我得勝.
沒有這個意思,耶和華利西是我的旗號,旗號就是我跟著.
是我的精神所在,是我獻身的地方,而耶和華都幫助我們得勝.
因為是他的旗,如果我跟你的旗,你跟我的旗.
不太行,很弱,手軟,這支旗就爛了,斷了.
耶和華的旗,嘩,厲害,這樣就一定很穩固,而且得勝.
所以我們跟著耶和華的旗,耶和華利西就是他們第一次正式面對面戰鬥.
而勝利,透過賣力戰鬥與禱告,而獲勝的一個紀念.
所以就築了一座壇,這座壇就叫做耶和華使我的旗.
是一座壇的名字,這座壇就是感恩讚美上主的能力.
而是我們一齊,好像祂是我們的旗號,我們向祂跟隨,屈身,帶領,敬拜.
這樣的意思,讚美,而且更加妙,就是耶和華已經起誓.
他說:阿瑪利人,我一定會跟祂,代表你們跟祂爭戰到底.
我今天的題目講到這裡,耶和華利西是一個壇的名字.
在這個壇,我們讚美上主得勝,而我們又用祂做我們的旗號.
為此付出努力,為此付出禱告,而是祂的旗,所以必勝的.
我們讚美主,我們一起祈禱,主我們讚美你,我們也要築一座壇.
我們再看回在我們人生裡面,在我們教會裡面,在我們的服事裡面.
很多時候都是因為靠你,靠我們又下力,而你又透過你的手幫助了我們.
我們又透過禱告,一齊來做這個工作,我們見到有果效,見到主的祝福.
主我們再次記得這座壇的名字叫做耶和華利西,是你,是我們的旗號.
我們跟隨,我們信服,求主帶領我們,顯出你的能力.
順天祈禱,奉耶穌基督的名字,阿們.
\newpage



\section{馬太福音 17:1-9}
\label{sec:032039}
\textbf{起來!下山去吧!}~~~~謝鵬輝傳道
\newline
\newline
連結: \href{http://www.church.com.hk/acms/content.asp?site=cdc\&op=show\&type=product\&code=032039\&layout=sermon}{\texttt{032039}} ~~~~ 語音日期: 2026-04-26
\newline
\newline
\hyperref[sec:032038]{\small{< < < PREV SERMON < < <}}
~
\hyperref[sec:index_chronic]{\small{[返順時目]}}
~
\hyperref[sec:032040]{\small{> > > NEXT SERMON > > >}}
\newline
\newline
馬太福音 17:1-9
\newline
\begin{longtable}{cl}
\hline
\hline
章節 & 經文 (和合本修訂版)\\
\hline
17:1 & \begin{tabularx}{0.7\textwidth}{X} 過了六天，耶穌帶著彼得、雅各和雅各的弟弟約翰，領他們悄悄地上了高山。 \end{tabularx} \\ \\ \relax
17:2 & \begin{tabularx}{0.7\textwidth}{X} 他在他們面前變了形像，他的臉明亮如太陽，衣裳潔白如光。 \end{tabularx} \\ \\ \relax
17:3 & \begin{tabularx}{0.7\textwidth}{X} 忽然，有摩西和以利亞向他們顯現，與耶穌說話。 \end{tabularx} \\ \\ \relax
17:4 & \begin{tabularx}{0.7\textwidth}{X} 彼得回應，對耶穌說：「主啊，我們在這裡真好！你若願意，我就在這裡搭三座棚，一座為你，一座為摩西 ，一座為以利亞。」 \end{tabularx} \\ \\ \relax
17:5 & \begin{tabularx}{0.7\textwidth}{X} 說話之間，忽然有一朵明亮的雲彩遮蓋他們，又有聲音從雲彩裡出來，說：「這是我的愛子，我所喜愛的。你們要聽從他！」 \end{tabularx} \\ \\ \relax
17:6 & \begin{tabularx}{0.7\textwidth}{X} 門徒聽見，就俯伏在地，極其害怕。 \end{tabularx} \\ \\ \relax
17:7 & \begin{tabularx}{0.7\textwidth}{X} 耶穌進前來，拍拍他們，說：「起來，不要害怕！」 \end{tabularx} \\ \\ \relax
17:8 & \begin{tabularx}{0.7\textwidth}{X} 他們舉目，不見一人，只見耶穌獨自一人。 \end{tabularx} \\ \\ \relax
17:9 & \begin{tabularx}{0.7\textwidth}{X} 下山的時候，耶穌囑咐他們說：「人子還沒有從死人中復活，你們不要把所看到的告訴人。」 \end{tabularx} \\ \\
[1ex]
\hline
\hline
\end{longtable}
$^{1}$兩弟兄姊妹平安.
在講到之前先問大家一個問題.
大家有沒有試過在一個地方.
那個地方你很想留在那裡.
覺得很開心.
很多回憶.
很愉快.
跟一些很開心的人在一起.
有沒有這樣的經歷.
有啦 見到別人點頭.
我相信這些可能是跟好朋友去旅行.
跟家人去一個很好的樂園.
去玩.
對我自己.
大概這兩年之間.
我都去了一些很美的地方.
在東南亞的國家.
更重要的是跟一些.
我很想一起去的.
年輕人一起去旅行.
連續兩年.
都去東南亞的地方.
那些地方很美.
環境很好.
更重要.
剛才說了.
就是跟一些很想.
而這班是我自己.
認識了十年以上的年輕人.
這班年輕人願意接待我去一起去玩.
所以在這些環境.
我覺得這裡真是好.
如果能夠留長一點.
多留幾天.
就更加好了.
不過.
去旅行.
總是要回家的.
去了幾天.
五天四夜.

$^{41}$四天三夜都好.
你都買了機票.
你住宿都有時間.
最近還要.
還要回去上班.
這個不能走.
所以還是要回家去.
所以今天.
都和大家去看這段經文.
這段經文就是.
一班門徒.
上到山上有經歷.
但是上得山.
自然要下山.
這段經文我大概分為三個部分去說.
第一個部分就是.
起來了不要再留在山上.
第二就是下到山上.
我們看清這是一個怎樣的世代.
是一個不信.
也是薄貌的世代.
第三部分我會讓大家去思考就是.
我怎樣活出一個.
知而行的信仰.
我們先看一看這段經文.
的故事發生什麼事.
這段經文.
很清楚有三個門徒上山.
比德,雅各和約翰.
上到山上.
有一個很大的情況發生就是.
耶穌變了形象.
不單止是耶穌.
甚至出現了兩個他們很認識的.
歷代的一些.
以色列的一些很偉大的領袖.
一個是摩西.
另一個是以利亞顯現.
這班門徒的回應很直接.
是很害怕的.

$^{81}$但是更重要他們說出一個感受就是.
我們在這裡.
真好.
我們很期望為這三個人.
耶穌,摩西,以利亞.
去搭一座的棚去.
去敬拜他們,讚美他們.
但若然我們只看到這段經文的時候.
從物來看.
我們應該會停在這裡.
會覺得.
真的不錯.
在這個山上的敬拜.
但若然我們要理解這段經文.
我相信.
更加要從.
這段經文的前文後理去理解它.
我們去看回前文後理的時候.
那些很小的字其實就是.
我告訴你是甚麼來的.
大概是馬太福音,馬賀福音和盧格福音.
如果大家看這三段的福音其實.
這個故事.
即是現在說.
耶穌認耶穌為基督.
耶穌預言自己的受難復活.
耶穌改變形象.
治好了鬼的附身的孩子.
甚至耶穌第二次預言.
其實這三段的福音書.
都是同一個次序出現的.
為甚麼要這樣做呢.
正正是想說出一個的.
內容給我們知道.
我們看一看說些甚麼.
第一.
讓他.
耶穌問我是誰的時候.
讓他說.
你是基督.

$^{121}$耶穌說.
你欣賞祂的.
讚賞祂.
但是很快出現一個情景.
就是當耶穌去給門徒說.
我要.
受死.
三天後復活.
彼得第一個跟耶穌說.
你不要去拉住他.
耶穌就說你難了我的.
你魔鬼退去吧.
彼得好像明白耶穌在做甚麼.
但原來一點都不明白.
接著.
耶穌就帶了三個他喜愛的門徒上山上.
這幫門徒的回應就說.
我在這裡真的好了.
這個感覺.
這個榮耀.
很開心.
很愉快.
但是上山要下山.
耶穌帶他們去看到一個很真實的處境.
這個是一個.
不信有薄貌的世代.
但這幫門徒明不明白呢.
當耶穌再跟這幫門徒說.
第二次預言祂自己要死.
三天後復活的時候.
門徒的回應是.
憂憂愁愁.
是不明白的.
從這個片段這個脈絡去看.
這幫門徒其實對耶穌所說的事.
其實是知叛戒.
不清不楚的.
因為我們對耶穌的期望.
是甚麼呢.
他們期望耶穌是一個軍事家.

$^{161}$來復興以色列.
使以色列離開羅馬帝國.
重新回到最鼎盛的王國時期.
但耶穌說了甚麼呢.
耶穌來地上是做些甚麼呢.
耶穌來地上就是受死.
三天後復活.
但門徒真的有一個停止的情景就是.
山上的這種經歷是很震撼.
甚至令他們很懼怕.
但門徒只能夠看到.
享受著山上那種榮耀.
而忘記了耶穌來地上最終的目的.
是要拯救這個世代.
最後的榮耀是在地上.
在釘身在十字架上.
而不是停留在山上.
我講一下我自己的一些經歷.
我自己的事 讓大家認識一下我.
我自己曾經有些團體分享.
我講我有自己十年.
是一個很荒誕的人生.
有機會可以請我分享.
我可以講一下我怎樣荒誕.
但更重要的是.
很感恩我有一個好朋友.
帶我回教會.
希望我能夠去迴轉.
所以我回了教會.
更重要的是他帶我去聽唐崇明牧師的解經.
如果大家有印象.
大概在零幾年開始.
唐牧師在尖沙咀浸信會逢星期二晚上.
在這裡講解羅馬書 約翰福音.
我當時正正在聽著約翰福音.
因為我這個朋友是報道團.
我經常去報道團看唐牧師的書.
對信仰的基礎建立了很多.
很快不如受浸.
因為我覺得自己曾經很早信了.

$^{201}$不過離棄了.
再次信了.
接著繼續聽唐牧師.
結果怎樣呢.
那段時間就很開心.
很滿足覺得.
這樣的生活是最理想最美好.
但其實是不是呢.
原來不知不覺發現了自己.
和這群山上的門徒一模一樣.
其實只不過是在享受一種.
我自己期望的那種信仰生活.
其實我不知道耶穌來地上.
忘記了耶穌來地上做什麼.
我就是為了自己.
在享受自己那種頭腦上的認知.
因為唐牧師講了很多信仰的解釋.
很豐富 很吸引.
頭腦上知道很多很多很多.
繼續很想聽.
但是奈何這些知識停留在什麼層面.
知識的層面.
我就是知道.
但其實不太明白耶穌來做什麼.
而這種感覺其實只不過是自己的滿足.
自己的一種感受.
覺得這種信仰體驗才是最好.
但這樣久而久之.
自己就變得停了下來.
停在哪裡.
停在這個山上.
我慢慢覺得我很滿足當時的情況.
我希望這種就夠了.
這種學習很足夠.
我不需要再找更多.
我覺得從星期二聽到.
星期日回來崇拜.
滿足了.
但發現了久而久之.
其實只有另外一些問題.

$^{241}$除了因為滿足現在的現況.
另外一些情況可能是生活的處境.
因為有很多人的工作很忙碌.
你說要花更多時間去學很多東西.
去神學院上課.
去學一些信徒課程.
都不容易.
因為今時今日大家工作時間很長.
又或者要照顧家庭.
或者是因為自己有很多忙碌.
一靜下來.
其實再做很多事已經很有限.
當久而久之.
如果長期留在山上.
有什麼問題呢.
慢慢這種就變成了一個習慣.
是最舒服的習慣.
就是逢星期一.
做好要做的事.
星期二我去聽到.
就OK的了.
很滿足的了.
我就會停在這裡.
很怕改變.
不要變了.
不要搞那麼多事了.
一變就會花很多時間.
有很多未知之數在裡面.
慢慢自己對這種靈性的渴求.
只是有那兩次就足夠.
反而變得冷淡.
我剛才說我回轉.
因為一個朋友.
這個朋友.
後期就正正.
因為他要結婚.
就是結婚生小朋友.
原先的他呢.
剛才說在報道團幫忙.
很熱心.

$^{281}$又去教年輕人去足球.
去踢球帶他們去比賽.
做很多的事.
不過結婚之後.
當然有家庭的都知道.
又要照顧家人.
又照顧小朋友.
重心放在家裡.
不是錯.
正常.
久而久之.
變成怎樣呢.
我問他 你近期怎樣.
都是這樣 有回教會的.
星期天都有.
我再問他.
你還有沒有其他的事奉.
有沒有幫手.
他說沒有了 教兩個小朋友.
幫他們去找學校都忙死了.
我記得這個朋友曾經.
很有心.
很想去做傳道.
我都是被他感動的一個.
結果呢.
我就做了傳道.
就停了在這裡.
停在他的山上.
我覺得是可惜的.
因為到今天.
可能他最重要的就是他的家庭.
不是錯.
但慢慢久而久之再問他.
有沒有想過再做多些什麼.
為教會做些什麼.
他說都沒有了.
變得有點冷淡.
變得好像現在的習慣.
生活就足夠了.
所以.

$^{321}$我們真的要去想一想.
我們是否要再起身.
不要再留在這個山上.
我們下山.
我們不要再讓自己.
停留在一個安書的環境.
我們下山上.
去看看這個世界.
發生什麼事.
耶穌正正下到山上.
就發生一件事.
他說.
有一個人來找耶穌.
跟耶穌說.
他的兒子.
他的癲癇病很嚴重.
嚴重到.
他掉到火裡面.
掉到水裡面.
經常這樣發生.
他跟耶穌說.
我帶他去你的門徒那裡.
都幫不到他.
耶穌很感嘆地說.
這個不信又薄貌的世代.
他要忍耐到什麼時候呢.
當然耶穌最後.
用他的能力.
去斥責這個鬼.
鬼就出來.
這個孩子就康復了.
這個片段讓我想起.
耶穌在.
之前的一個教導.
這個經文大家都熟悉.
在馬太福音.
九章35至36節.
這樣說.
耶穌走遍各城各鄉.
在會堂裡教訓人.

$^{361}$宣講天國的福音.
有醫治各樣疾病.
他看見許多人.
就憐憫他們.
因為他們困苦流離.
如同羊沒有牧人一樣.
這段經文.
很提醒我.
耶穌來地上.
做些什麼呢.
耶穌來地上就是拯救.
顛沛流離的.
迷失的失喪的人.
不知道大家有沒有留意.
你身邊有沒有這些人呢.
有沒有一些需要.
主的福音.
去幫助你的人呢.
這個世代是一個什麼世代.
這個世代是一個.
不順又薄貌的世代.
我相信在新年期間.
有很多畫面.
一定會出現.
先說左邊.
左邊先問下大家.
有沒有親人.
有沒有見到.
自己的家人.
仍然是裝香拜神.
有些人說.
我相信大家都很想.
他們傳福音.
傳給他們.
但是對我來說.
我自己小時候是裝香的.
因為我爸媽都是.
拜傳統信仰.
當時.
做什麼呢.

$^{401}$每逢早晚我爸媽裝香.
下午跟我說.
裝香.
數數數數數.
每個香爐插三支.
七個香爐插二十支.
每次都如是.
爸媽叫.
或者不叫.
差不多時間就裝香.
但你問我知不知道做什麼.
不知道.
爸媽叫嘛.
不做就怎樣.
還不做.
一定這樣.
但為了做.
做完不明白.
確實就持續去做.
不過感恩就是.
今天我家沒有這些.
因為我媽我妹都相信.
我爸就未信.
繼續為他禱告.
另外那幾張.
大家何時會見到.
新年前後.
都很熟悉的人物.
但更重要是.
新年的習俗.
去莊頭住鄉.
去長洲長包山.
拿個平安包.
右邊那棵樹.
今天變成膠樹.
被人扔爛了.
去林村.
我自己都去過.
很多年前第一次拍拖.
就跟女朋友去.

$^{441}$習俗嘛.
新年就去一下.
其實這些習俗.
是什麼呢.
很多時不知道.
不太清楚.
家人要你做.
傳統就跟著做.
大概六七年前.
有件事挺震撼我的.
我住的地方有個巴士站.
我經過商場的時候.
一定要經過巴士站.
巴士站旁邊.
有個站頭.
巴士司機坐的站頭.
有個叫右邊的.
有些香爐.
神主牌.
我自己一次去商場的時候.
就見到一個.
大概我自己估.
三四歲左右的妹妹.
在這張桌子.
我記得這張桌子.
她做什麼呢.
走過去神主牌.
我見到的時候.
我心裡很不開心.
因為在我心目中.
我見到這個妹妹.
如果你問她.
知不知道她在做什麼.
不會知道.
對她來說 為什麼會這樣做.
因為父母做過.
做開 那就跟著做.
那我在想.
她長大一點的時候.
怎麼辦呢.

$^{481}$如果感恩 有人跟她傳福音.
她有機會去信主.
就真是好了.
我當然不會捉住她跟她說.
小妹妹小妹妹.
嚇死她.
唯有為她祈禱.
希望她能夠明白.
她在做什麼.
能夠去信主.
這個世代.
有很多事情.
是很混亂.
很迷失.
甚至很荒謬.
AI近期很熱.
左邊那張圖.
是AI.
jam出來的.
小妹妹.
會跳舞.
用AI做的.
現在的世代就是這樣.
你要用AI想做什麼.
就做什麼.
甚至去到右邊那張照片.
大家都覺得很震撼.
為什麼會這樣.
不是AI是邪惡.
是人的心.
是因為這個世代.
就是這麼邪惡.
這個世代就是有很多.
你不公不義.
很混亂很混沌.
很彎曲的事情.
會發生.
剛才我講到我自己.
繼續講一下我自己.
發生什麼事.

$^{521}$其實中間要加插一些東西.
我自己讀神學之前.
一些事件的經歷.
其實很早已經有夢召的經歷.
原來比受浸還早.
受浸之前已經有這個經歷.
因為當時有一次.
唐木屍的一個暴徒會.
問有沒有人.
有夢召的心.
想去擺上.
可能當時的情感.
感動舉了手.
但沒有走出去.
收拾好放在那裡.
我自己.
剛才講了.
對這個信仰一知半解.
放在那裡.
有跟別人分享過.
甚至去到.
尖浸的時候.
當時尖浸有夢召團契.
去到夢召團契.
我經常問牧師.
有夢召但好像不太肯定.
是不是真的上帝呼召我.
我都不太明白.
當時有個牧者跟我說.
是的,去吧.
不要想那麼多.
我就在想.
是不是那個傳到牧者也是這樣.
有人應了就去吧.
但是.
他說上帝是很幽默的.
上帝真的給你這個呼召.
上帝一定會有方法提醒你.
我自己正正就是.
在讀神學之前.

$^{561}$他給了我一個經歷.
重新去思考.
究竟發生什麼事.
有一個年輕人.
這個年輕人怎樣認識的呢.
就是我一個同事.
他的表弟的同學.
認識了他.
這個年輕人很特別.
這個年輕人一歲換肝.
當時是找了.
因為換肝很難.
所以出了很多報道.
找人招一個肝回來.
最後換到肝.
一歲換了肝.
這個年輕人因為換了肝.
父母就很疼他.
很寵愛他.
疼到怎樣呢.
什麼都給他.
什麼都給最好的他.
結果這個小朋友有個問題.
就是很喜歡說謊.
我這樣形容他.
我認識他這麼久.
他起碼.
說了超過一百個謊言.
十句裡面.
起碼有兩三句是假的.
但我仍然陪著他.
因為他很喜歡和我一起.
又喜歡問事.
又喜歡打乒乓球.
我陪他練乒乓球.
這個小朋友很厲害.
你看到細節.
他在08,10年.
因為換過肝.
所以打到殘疾運動會.

$^{601}$去到哪裡打呢.
國內打.
國內很多精英.
拿到獎牌.
其中一個最厲害的就是金牌.
這個年青人的事.
我為他很驕傲.
這個年青人真的很厲害.
這個年青人繼續努力.
真的可以做到很多事.
但我先說.
這個年青人喜歡什麼.
說謊.
說謊慢慢變成怎樣呢.
越來越喜歡騙人.
騙到怎樣呢.
開始網上越來越普及.
網上的騙錢.
騙門票.
不停地.
不停地騙人.
你看到下面的報道.
騙八千元.
其實算起來不止.
騙了很多個幾千.
最大個八千.
騙門票.
我看他的Facebook.
很多都是和錢有關的.
當時我跟他認識的時候.
我跟他說.
他問我很多事.
都問他會不會回來教會.
他問我耶穌的事.
我都回答他.
跟他聊天都會明白一點.
我就覺得.
你有回教會就好了.
但我沒有很用心.
希望他轉變.

$^{641}$結果大家看到報道.
就是坐牢.
坐了一年大概三個月在壁屋.
我大概探了他兩三次.
我探他的時候.
隱約都看到.
一些異象.
一個肚子.
很瘦的.
好像一條柴一樣.
後期我看到有個八個月十個月的肚子.
出來更明顯.
大家應該猜到.
發生什麼事.
他就是肝臟出現問題.
因為他不珍惜他自己.
在外面的時候曾經花天酒地.
亂喝.
不照顧自己.
喝酒.
做了不應該做的事.
大概不夠半年的時間.
就離開了.
對我來說.
這個年輕人.
我覺得他是欠缺了一些東西.
欠缺了什麼.
欠缺了跟他講福音.
我沒有很積極.
去看到.
他的需要.
他的需要就是給主的肝.
去照亮他.
當時的我來說我仍然是迷迷糊糊.
我就覺得陪著你.
打乒乓球教你一些東西.
你會學好.
但原來是不足夠.
因為這個世代.
就是這麼彎曲.

$^{681}$這麼薄茂.
是讓人.
離開一個正道.
走一條歪路.
唯有上帝的肝才能改變他.
但是我沒有做到.
所以這件事.
令我很大的激動.
就是我要做好要做的事.
我要做什麼.
我要將上帝的道.
分享給人.
我要去見證他.
讓人看到主的肝.
這是上帝.
給我們的使命.
也是上帝耶穌.
來地上的使命.
去慰藥拯救失喪的靈魂.
我身邊有很多.
不知道大家有沒有看到.
不知道大家有沒有留意.
我們有沒有為這些人.
去緊張呢.
為這些人.
去擺上禱告呢.
所以對我來說.
我不再.
我不可以再一知半解.
我會更加清楚.
耶穌想我們做些什麼.
我們只知道.
要有行動.
如果只知道沒有行動.
我們所信的.
都只是在頭顱上面.
我們能不能夠透過行道.
或者透過一些.
生命的根深.
去活現上帝的教導.

$^{721}$讓更多的人看見主的肝.
上帝教我們做些什麼.
耶穌說他來地上.
就要釘上十字架.
去成就救恩.
我們的回應就是去選擇.
去跟從他.
去捨己.
背起自己的十字架.
去跟從他.
捨己是什麼意思呢.
放下自己.
雖然中國人有句說話.
人不為己 天誅被滅.
但上帝說去到他面前.
你要放下自己的心.
你要去仰望他.
讓他去改變你.
讓他去帶領你.
而這個背起不是.
就是星期日或者某一個時刻背起.
每一刻 每一天.
在路加福音講得更全面.
就是天天背起的十字架.
這個也是.
我自己路加團的團訓.
當時我不太明白.
但我更加越來越明白.
是要天天去背起.
這個十字架是什麼.
有些人覺得好像很難.
十字架.
就是你願意去跟從主.
不單止是順境的時候.
跟從他.
任何警方.
主想帶你去哪裡.
你就讓主帶領你去哪裡.
讓自己的生命.
放在主的手裡.

$^{761}$效法主.
主想我們做些什麼.
讓人看見.
主的實在.
主的愛 主的美善.
耶穌曾經講到.
在山下的時候.
這個是一個不信的世代.
耶穌不單止.
說是這個世代.
也跟門徒說.
因為他們對耶穌.
都是一知半解.
究竟我們是不是.
信是否不足夠.
是不是這個不信.
也是跟我們說.
希望我們可以反思.
我們在這個世代裡.
這個世代很混亂.
我們更需要的是.
上帝的智慧.
帶領我們生活.
否則我們很容易迷失.
很容易跌倒.
我們要有主的教導.
天天讓他的話.
更新我們.
我們才能夠抵擋.
魔鬼的引誘.
魔鬼不會告訴你.
作著我是魔鬼.
去到你面前想你怎樣.
魔鬼很聰明.
魔鬼會變成一個糖衣.
變成一顆毒藥.
你也未必知道.
所以我們要起來.
我們要有行動.
我們要在靈性上不停地操練.

$^{801}$讓自己的生命.
去被更新.
如果你說最簡單的方法.
我會說.
從禱告開始.
禱告是一個很好的操練.
因為禱告.
可以做很多的事.
我們不單止求.
我們也去到上帝面前.
我們也將我們的心意.
求主去更新.
有時我們得罪了上帝 得罪了人.
我們也未必知道.
去求神的光.
去光照我們.
先從禱告開始.
仰望上帝的公認.
去經歷上帝那份.
大令.
所以教會.
每個主日.
都會有報告.
我們有周三祈禱會.
周五早禱會.
接下來這個禮拜.
禮拜六.
還有一個叫敬拜祈禱會.
因為這樣禮拜三的祈禱會就沒有了.
我希望.
透過禱告.
讓我們的生命去改變.
禱告是什麼呢.
禱告就是你願意.
去到主的面前.
求主去改變.
我自己.
第一次的禱告.
禱告是怎樣呢.
我小時候.

$^{841}$小學和中學都是在基督教學校.
對我來說聽很多.
因為大時大節.
就會有很多的報道會.
可以講很多.
你聽得多.
就知道.
知道就祈禱.
第一次的禱告是什麼情況呢.
就是.
我家是一個.
經常吵架的家.
我父母會因一些.
很小的事吵架.
因為我父母很喜歡打麻將.
打麻將.
我父母就會.
發脾氣.
發脾氣就會扔麻將.
扔麻將我母親就覺得.
很害怕.
打的時候.
接著回家的時候.
一邊罵.
一邊打.
接著回家.
吵架的時候.
我母親就說離婚吧.
接著她就.
走回家.
走回外公外婆那裡.
對我來說.
我是這樣成長的.
我父母就不停地.
有空就吵架.
對我來說我沒什麼平安.
我父母的祈禱.
我一次罵得很厲害.
罵到鄰居也有空來投訴.
是這麼誇張的.

$^{881}$再吵就報警了.
是很誇張的情況.
對我來說我很害怕.
我唯一那次很害怕的時候.
我走了出去冷巷.
我在冷巷的時候.
我唯一想到的就是.
祈禱.
我跟天父說.
天父我受不了了.
我母親經常這樣.
我不知道家裡什麼時候受不了.
求你去幫助我.
接著我就硬門.
我在冷巷.
再逗留多五六分鐘.
平靜了.
哭了一場就回家.
回到門口.
很靜.
我想我母親應該走了.
回到我父母那裡.
但不是.
我開了門.
坐著很安靜.
我馬上感覺到.
上帝真的聽禱告.
一個出於.
真心.
渴慕的祈禱.
上帝怎會不聽.
很多時候我們禱告.
都是為了自己.
希望上帝幫我什麼.
但很多時候我們忘記了.
上帝希望帶來平安.
帶來安穩.
一個求安穩的禱告.
帶領的禱告.
上帝一定會聆聽.

$^{921}$今年教會.
講什麼主題.
結年全城復興.
講到復興.
什麼是復興.
很多時候人會覺得.
復興要轟轟烈烈.
一講復興.
大家很澎湃.
其實復興.
不是要翻天覆地.
不是要很震撼.
風雲變色.
第一步是先由內在的生命.
去改變.
怎樣內在的生命改變.
多些禱告.
讓自己去仰望上帝.
更重要是.
讓上帝改變我們.
不要停留在這個山上.
如果我們看刀.
我們是明白上帝的教導.
讓這份不寒淡的.
靈性再一次.
被更新.
讓上帝帶領我們.
看見這個世代的需要.
為自己.
為教會祈禱.
上帝說教會是地上的鹽.
我們有沒有多些.
為教會的祈禱.
希望教會成為這個地區.
成為香港的燈台.
我們需要每一個人.
去獻上自己.
這個一定能夠復興.
所以我們.
下來下山.

$^{961}$去回應上帝.
我們去用愛.
用禱告.
為這個不信.
由薄默的世代禱告.
願上帝幫助我們.
我們成為一個.
明白祂的心意的人.
能夠聽道.
行道.
透過禱告.
上上去顧念教會.
顧念教會的牧者.
顧念教會怎樣去.
實踐信仰.
讓我們能夠去成為.
這個地上的乾.
我們一起禱告.
七天上帝.
因為你犧牲.
自己為要的.
是成就這個救恩.
我們得救是因為.
你的呼召.
讓我們不要停留在.
只是對你的認知上.
而願意將生命.
去塘開.
我們願意走到人前去見證你.
我們願意.
生命去更新去改變.
我們從禱告.
從領修.
從讀經開始.
讓我們的生命能夠去轉化.
我們看到自己的不足.
看到我們願意.
去到你面前.
效法你.
我們願意去見證你.

$^{1001}$讓未信的人從我們身上.
看到主的乾.
我們禱告.
阿們.
\newpage



\section{路加福音 18:31-43}
\label{sec:032040}
\textbf{看見和跟隨耶穌的道路}~~~~梁成裕牧師
\newline
\newline
連結: \href{http://www.church.com.hk/acms/content.asp?site=cdc\&op=show\&type=product\&code=032040\&layout=sermon}{\texttt{032040}} ~~~~ 語音日期: 2026-04-26
\newline
\newline
\hyperref[sec:032039]{\small{< < < PREV SERMON < < <}}
~
\hyperref[sec:index_chronic]{\small{[返順時目]}}
~
\hyperref[sec:032041]{\small{> > > NEXT SERMON > > >}}
\newline
\newline
路加福音 18:31-43
\newline
\begin{longtable}{cl}
\hline
\hline
章節 & 經文 (和合本修訂版)\\
\hline
18:31 & \begin{tabularx}{0.7\textwidth}{X} 耶穌把十二使徒帶到一邊，對他們說：「看哪，我們上耶路撒冷去，先知所寫的一切事都要成就在人子身上。 \end{tabularx} \\ \\ \relax
18:32 & \begin{tabularx}{0.7\textwidth}{X} 他將被交給外邦人；他們要戲弄他，凌辱他，向他吐唾沫， \end{tabularx} \\ \\ \relax
18:33 & \begin{tabularx}{0.7\textwidth}{X} 並要鞭打他，殺害他；第三天他要復活。」 \end{tabularx} \\ \\ \relax
18:34 & \begin{tabularx}{0.7\textwidth}{X} 這些事門徒一點也不明白，這話的意思對他們是隱藏的；他們不知道所說的是甚麼。 \end{tabularx} \\ \\ \relax
18:35 & \begin{tabularx}{0.7\textwidth}{X} 耶穌將近耶利哥的時候，有一個盲人坐在路旁討飯。 \end{tabularx} \\ \\ \relax
18:36 & \begin{tabularx}{0.7\textwidth}{X} 他聽見許多人經過，就問是甚麼事。 \end{tabularx} \\ \\ \relax
18:37 & \begin{tabularx}{0.7\textwidth}{X} 他們告訴他，是拿撒勒人耶穌經過。 \end{tabularx} \\ \\ \relax
18:38 & \begin{tabularx}{0.7\textwidth}{X} 他就呼叫說：「大衛之子耶穌啊，可憐我吧！」 \end{tabularx} \\ \\ \relax
18:39 & \begin{tabularx}{0.7\textwidth}{X} 在前頭走的人就責備他，不許他作聲，他卻越發喊叫：「大衛之子啊，可憐我吧！」 \end{tabularx} \\ \\ \relax
18:40 & \begin{tabularx}{0.7\textwidth}{X} 耶穌就站住，吩咐把他領過來，他到了跟前，就問他： \end{tabularx} \\ \\ \relax
18:41 & \begin{tabularx}{0.7\textwidth}{X} 「你要我為你做甚麼？」他說：「主啊，我要能看見。」 \end{tabularx} \\ \\ \relax
18:42 & \begin{tabularx}{0.7\textwidth}{X} 耶穌對他說：「你看見吧！你的信救了你。」 \end{tabularx} \\ \\ \relax
18:43 & \begin{tabularx}{0.7\textwidth}{X} 那盲人立刻看得見了，就跟隨耶穌，一路歸榮耀給神。眾人看見這事，也都讚美神。 \end{tabularx} \\ \\
[1ex]
\hline
\hline
\end{longtable}
$^{1}$今天我們分享的題目是.
看見和跟隨耶穌的道路.
我們會繼續《路加福音》十八章裡面的經文.
我不知道如果大家去一個你沒去過的地方旅行.
你會怎麼去預備這個旅程呢?.
面對一個你沒去過的地方.
可能你會想一下.
咦?我坐飛機去啊.
怎麼坐?.
或者是坐火車.
坐哪一班去呢?.
到了那個地方你可能會想.
那我坐什麼交通工具.
去到酒店.
或者去到我的目的地呢?.
你可能會想.
那裡的天氣是怎麼樣的呢?.
我要預備什麼衣服呢?.
或者要預備什麼物資.
或者那裡的氣候.
或者那裡的水土是怎麼樣的呢?.
我們需不需要預備一些藥物呢?.
在四月初的時候.
我就很感恩.
因為我和太太.
和一群弟兄姊妹有機會去雲南旅行.
因為其實是弟兄姊妹為我預備和計劃這個旅程的.
所以我什麼都不用想的.
無論我們是坐什麼飛機去.
到了的時候坐什麼交通工具.
去到酒店.
住什麼酒店.
每一天的行程是怎麼樣.
去哪裡吃.
以及吃什麼.
我也多謝弟兄姊妹為我預備周詳的.
但是我發現當我快要出發的時候.
我發現有一些東西.
是我自己需要去預備.
但是我還沒預備.

$^{41}$我記得和我的中學同學吃飯.
我就說很開心啊.
接下來我去雲南了.
然後我其中一個同學就說.
你在雲南也預備好了嗎?.
我就說預備什麼呢?.
他說你知不知道會有高原反應的.
我說這是什麼呢?.
原來你去比較高的地方可能會有高原反應.
他說你沒有預備的嗎?.
我說沒有啊.
然後我問另一位同工.
他就跟我說.
原來你要小心.
因為他說八十年代有一位歌星.
去了一個很高的地方.
然後面對高原反應身體出現了狀況.
我說糟糕了.
我沒有想過.
然後立刻問一些團友.
團友說是啊.
你要預備好的.
你要吃一種藥叫做紅莖天.
他說你要先吃紅莖天才去.
不僅如此.
他說你要預備下飛機之後.
又要有藥物和氧氣.
我說蛤?.
為什麼要預備這麼多東西呢?.
完全不知道.
他直接去買了紅莖天.
跟我說你現在要開始吃了.
你要開始預備好了.
不僅如此.
我另一位同學就說.
你知不知道那裡的紫外光很厲害的.
我說是嗎?.
很厲害的嗎?.
他說你要預備好.
你有沒有防曬.

$^{81}$遮住你全身的帽子.
我說沒有啊.
他說你什麼都沒有預備就去啊.
我以為是就這樣去而已.
於是他就說你快點去找那些帽子.
還有你小心啊.
你有沒有那些保護眼睛的鏡子.
因為那些紫外光可能傷害到眼睛的.
我說是嗎?.
於是我有一次就去買了這些所有保護自己的衣物.
之後我的樣子就變成這樣了.
昨天有知體就說.
這張照片很像一個姓蔡的過氣歌手.
我說不過我跟他不一樣.
我的生命就絕對不空虛.
還拿出來比較.
但是我想說.
原來有時候我們跟隨耶穌的生命也是這樣的.
我們口說很開心啊.
我現在踏上跟隨耶穌的路了.
很開心啊.
我的生命已經是耶穌的門徒了.
但是我們沒有想過我們的生命是怎麼跟隨耶穌的.
我們沒有想過我們會遇見什麼.
譬如說我在雲南.
我可能有高原反應.
我可能會有紫外光.
但是今天我們跟隨耶穌.
你有沒有想過跟隨耶穌你會遇見什麼呢?.
你有沒有想過你怎麼去預備好自己去跟隨耶穌呢?.
我們只是一個口號.
哇!我跟隨了耶穌了.
你跟你的中學同學說.
然後你的中學同學問你.
那你預備好了沒有?.
要預備什麼呢?.
跟隨耶穌會面對什麼呢?.
會面對什麼呢?.
但是今天我們要搞清楚.
當我們的生命說我們跟隨耶穌.

$^{121}$我們會面對什麼呢?.
今天當我們的生命說跟隨耶穌.
我們又預備好了沒有呢?.
上個星期我們說路加福音.
十八十五至三十節的經文.
那裡說到我們的心.
要好像一個小孩子一樣.
單單跟隨耶穌.
我們的生命要去清除跟隨耶穌的障礙.
我們的生命要撇下一切去跟隨耶穌.
但是今天我們有沒有想清楚.
究竟跟隨耶穌的生命是怎樣的呢?.
而我們有沒有預備好.
我們的生命去跟隨耶穌.
所以今天當我們去看.
三十一到四十三節的經文的時候.
讓我們學習三個很重要的功課.
今天我們看的第一個功課.
是我們要走背起十字架的道路.
原來今天跟隨耶穌.
是要背起十字架去跟隨耶穌的.
我們一起讀三十一至三十四節.
一起讀一二三.
耶穌帶著十二個門徒對他們說.
看啊!我們上耶路撒冷去.
先知所寫的一切事都要承就在人子身上.
他將要被交給外邦人.
他們要戲弄他.
凌辱他.
吐唾沫在他臉上.
並要鞭打他.
殺害他.
第三日他要復活.
這些是門徒一樣也不懂得.
意思乃是隱藏的.
他們不曉得說的是什麼.
在這裡耶穌帶著十二個門徒.
跟他講解耶穌的使命.
和耶穌在救贖歷史上所扮演的角色.
耶穌特別將自己將要受苦的事.

$^{161}$聯繫於之前他跟門徒的教導.
耶穌說先知所寫的一切事.
都要承就在人子身上.
在這裡其實他是想告訴門徒.
我要走十字架的道路.
而當我走十字架的道路.
當你跟隨我的時候.
你也是要走十字架的道路.
看啊!我們上耶路撒冷去.
代表耶穌不是偶然經過耶路撒冷.
耶穌是有目的.
有方向地走向十字架.
先知所寫的一切事.
都要承就在人子身上.
是在講耶穌的受苦不是失控.
也不是失敗.
而是要應驗舊約的預言.
承就舊恩.
耶穌很清楚地告訴門徒.
我的路就是一條十字架的路.
所以在第32節.
耶穌就說他將要被交給外邦人.
他們要戲弄他.
凌辱他.
吐唾沫在他臉上.
在這裡耶穌說猶太的領袖.
將他會交給外邦人.
而將他交給外邦人的時候.
那些人會戲弄他.
凌辱他.
吐口水.
在他的生命那裡.
在這裡耶穌不單是受肉身的痛苦.
更加要放下他這個尊崇的地位.
耶穌將要為我們進入一個最深的羞辱和苦難.
而在33節耶穌繼續說.
他們要鞭打他.
殺害他.
鞭打是囚犯釘十字架之前所受的酷刑.
所以鞭打和殺害是連在一起的.

$^{201}$他會首先被鞭打.
接著被殺害.
如果聖經讀到這裡就完了.
那我和你都沒有希望了.
但是耶穌接著說了很重要的一句.
第三天他要復活.
耶穌在這裡說得很清楚.
他的死將要被上帝紀念.
也證明耶穌之前和門徒說的話是真實的.
如果大家記得上個星期耶穌說.
我實在告訴你們人為神的國撇下房屋.
或是妻子弟兄父母兒女.
沒有在今世不得百倍.
在來世不得永生.
耶穌說得很清楚.
為什麼你們會有永生.
因為我首先會復活.
所以耶穌說得很清楚.
我會走十字架的道路.
但是我將要復活.
耶穌說我會為你們受苦走到死亡.
但是我將要復活.
所以耶穌告訴他們.
是的,我會走十字架的道路.
但是十字架的道路不是以死亡作為收場.
是以復活作為終結.
十字架是要和復活連在一起看的.
今天我們都要知道.
我們背十字架.
也要和耶穌的復活連在一起看.
當我們願意走十字架的道路.
受苦不是結局.
因為基督的榮耀才是終結.
當耶穌告訴門徒.
我要走十字架的道路.
門徒明白嗎?.
門徒不明白.
34節告訴我們.
這些事門徒一樣也不懂得.
以斯乃是隱藏的,他們不曉得所說的是什麼.

$^{241}$門徒不明白耶穌.
以斯乃是隱藏代表門徒心裡面.
他們仍然期待一位在地上得勝.
地上興盛.
在地上趕走羅馬軍隊的一個彌賽亞.
所以對他們來說.
一個受苦的彌賽亞超出了他們的想像.
當耶穌告訴他們的時候.
當時他們不明白.
但原來之後他們會明白.
如果我們看看路加福音24章這段經文.
之後他們會明白.
因為當耶穌復活了之後.
耶穌開他們的心竅.
就叫他們明白聖經.
今天耶穌復活了嗎?.
復活了.
原來大家不是很確切.
耶穌復活了嗎?.
復活了.
我們明白聖經嗎?.
明白.
門徒也明白.
你明白嗎?.
所以今天我們要搞清楚.
原來跟隨耶穌的路.
是背十字架的路.
但今天我們要問自己.
當我們說我們跟耶穌的時候.
我們的生命願不願意去背十字架.
去跟隨耶穌.
耶穌不是意外受苦.
祂更加沒有逃避苦難.
祂是甘心走十字架的道路.
去完成上帝救贖的計劃.
所以十字架的道路不是失敗的.
十字架的道路不是失敗.
是神榮耀的得逞.
而今天原來我們說我們作為門徒.
我們的生命都要以十字架為我們生命的中心.

$^{281}$我們不是去追這個世界眼中的成功.
我們也不是去逃避苦難.
而是我們都願意去背起十字架.
願意為主受苦和付代價.
原來真正的門徒.
是願意背十字架.
和生命願意去跟隨耶穌.
我認識一位少年人.
他在學校經歷了被同學欺凌的經歷.
但是他沒有選擇逃避.
也沒有求神.
「神啊,你幫我轉去另外一間學校吧」.
他是選擇面對.
他學習如何去面對那些欺凌他的同學.
他繼續學校的生活.
並且積極去認識新朋友.
不僅如此.
他更加用他過去被欺凌的經歷.
去關心那些都只是被欺凌的同學.
帶他們走出這段困難的道路.
他在學校繼續認識新的朋友.
他認識了新的朋友.
很積極地邀請他們回來教會跟他們傳福音.
原來面對苦難.
這個少年人沒有選擇逃避.
我們的生命都會有困難和挑戰.
但是我們都不是去逃避.
而是在當中去活出一個跟隨神的生命.
所以今天當我們面對困難和挑戰.
我們會問.
我們願不願意去活出神給我們的生命.
背起十字架.
在我們的困難和挑戰裡面.
繼續去跟隨耶穌.
當我們跟隨耶穌的時候.
有時候都會被人凌辱.
我想問一下.
有哪一個傳福音被人拒絕過?.
有的.
應該全部都有的.

$^{321}$大家記得.
我在公園跟婆婆聊得很開心.
聊到差不多她突然跟我說.
走啦肥仔.
是會被人拒絕的.
有沒有人去洗手間.
和你一起洗手間的時候.
被人關上門叫你離開的?.
有的.
我經常說其實去洗手間.
是WIN WIN的.
為什麼WIN WIN呢?.
因為當你敲門的時候.
如果他拒絕你.
閉上你的門.
你就為基督背十字架了.
是不是?.
但是如果他開門呢?.
那你又為神開了福音的門.
所以原來他開門或者關門.
你都WIN WIN的.
是不是?.
他關上門.
你就為主背十字架了.
他開門讓你進去.
你就開了福音的橋樑.
哇!正不正啊?.
有什麼是WIN WIN的?.
下次洗手間你去不去啊?.
去啊!是不是?.
你跟大約說WIN WIN啊!.
但是今天很重要.
我們的生命都可能被凌辱的.
我們的傳福音會被人拒絕.
甚至會被人取笑.
我又不知道當你跟隨耶穌的時候.
當你生命改變.
有沒有人說.
你怎麼會這麼傻啊?傻子.
有沒有人這樣說過你?.

$^{361}$你真是傻啊!傻子.
有沒有?因為他說我傻女.
是不是?.
有人會說你傻啊.
有人會說你.
你不要這麼迷啊.
有沒有人叫你不要這麼迷啊?.
有!是不是?.
但是你要問自己.
今天你願不願意.
去背起十字架跟隨耶穌.
但是很有趣的聖經說.
我們每個人要背起自己的十字架跟隨耶穌.
就是說我們的生命.
我們都會面對不同的狀況.
我們的生命會面對不同的處境.
我們會有不同的困難.
我們會有不同的挑戰.
但是今天.
無論你遇見的是什麼.
你願不願意去付代價去跟隨耶穌.
你願不願意為耶穌去走十字架的道路.
而今天跟著我們分享的第二點.
是我們要認清耶穌是誰.
並且呼求耶穌去改變我們的生命.
我們一起從35讀到39節.
我們一起讀.
1,2,3.
耶穌將近耶利哥的時候.
有一個孩子坐在路旁吐飯.
聽見許多人經過就問是什麼事.
他們告訴他是拿撒勒人耶穌經過.
他就呼叫說.
大位的子孫耶穌啊.
可憐我吧.
在前頭走的人就責備他.
不許他作聲.
他卻越發喊叫說.
大位的子孫可憐我吧.
在這裡我們看到一個很重要的一段聖經.

$^{401}$告訴我們耶穌將近耶利哥.
有一個盲人在路邊乞食.
當時的古城耶利哥.
我們看舊約的耶利哥已經變成廢墟.
這個耶利哥是近著古城的廢墟附近的一個新城.
而它是通往耶路撒冷的通道上.
就是說耶穌正在去耶路撒冷.
正在去耶路撒冷就經過耶利哥這個新城.
在當時的希羅文化裡面.
盲眼是一個很嚴重的疾病.
而這個人不單止他有肉身的疾病.
他更加坐在那裡乞食.
代表他是一個被社會遺棄的人.
但是這個人他聽到有很多人經過.
就問什麼事啊.
這個人他的眼睛看不到.
但是他有耳朵.
還有他有好奇心.
更加重要他有一個對自己生命轉化的敏感.
他不是只坐在那裡.
他是留意著那些人群.
留意著那些動靜.
留意著他有一個生命的轉機.
當他問什麼事的時候.
其實他是在敏感自己的需要.
並且為自己的生命去找答案.
今天我們有沒有敏感我們生命轉化的一個機會.
我們有沒有為自己的生命去找答案.
當他問的時候.
第37節那些人就跟他說.
哦!是拿撒勒人耶穌經過.
這個經文告訴我們.
那些群眾他們還沒有真正了解耶穌真正的身份.
因為這是一個很普通的稱呼.
拿撒勒人耶穌.
拿撒勒就是耶穌的家鄉.
你問他那是誰.
哦!那是拿撒勒的耶穌.
但是對於這個盲的人來說.
他抓住了一個希望.

$^{441}$我們從他怎樣去稱呼耶穌.
我們就發現.
原來他看到一樣那些有眼的人看不到的東西.
他怎樣稱呼耶穌呢.
我們看38節.
我們一起讀好不好.
1,2,3.
他就呼叫說大衛的子孫耶穌.
可憐我吧.
這個盲的人看到那些不盲的人看到的東西.
那些不盲的人稱呼他拿撒勒人耶穌.
但是這個盲的人.
他說大衛的子孫.
大衛的子孫代表什麼.
在舊約裡面將來要來的救主.
是來自大衛家的.
所以當他稱呼他大衛的子孫.
其實他知道這個就是在舊約裡面將要來的救主.
這個人一定聽過耶穌做過的事.
他一定聽過耶穌行的神蹟騎士.
而他認定這個就是上帝施行拯救那個人.
所以他稱呼這個人做大衛的子孫.
他未必知道他是一個大能的醫生.
但是他知道這個人是拯救我生命.
大衛的子孫是一個很重要的迎信宣言.
當這個盲眼的人去稱呼耶穌為大衛的子孫.
然後他說可憐我吧.
這個盲的人知道我不能夠靠自己得著救恩.
於是他向這個大衛的子孫去呼求.
希望他的生命得著拯救.
原來今天我們的生命有三樣很重要的東西.
我們要正確地認識耶穌.
但是我們也要正確地去認識我們自己.
我們要知道耶穌是救主.
我們要知道我們不能夠靠自己得救.
而我們要對耶穌作出一個正確的生命回應.
就是今天我們究竟是否懂得呼求耶穌去改變我們的生命.
去到第39節.
當這個人大聲叫的時候.
那些在前面走的人就責備他.

$^{481}$說你不要出聲.
但是當那些人不讓他出聲的時候.
他會怎樣?.
他更大聲.
是不是?.
就好像你.
你不要說話.
大聲叫大衛的子孫.
可憐我吧.
今天我們有沒有認清耶穌是誰.
並且呼求耶穌去改變我們的生命呢?.
前幾年我跟神學院一班朋友.
一班好朋友一起去台灣旅行.
這個台灣旅行是自駕遊.
我們去到一個比較偏僻的地方.
就爆了胎.
我們七個人正在坐一輛九人車.
不知道為什麼.
我坐那邊的車胎就爆胎了.
我的朋友就認定.
是我的體重讓這輛車爆胎.
爆了胎之後.
我們就想著.
因為我們有幾個人都懂得換胎.
我們就說自己換吧.
但是有一個同學就很聰明.
說不是不是.
你看那個租車合約.
我們就在車上拿那個租車合約來看.
一看裡面寫得很明白.
如果你的車有任何故障.
你不要自己去修.
你一定要打去租車公司.
等他來處理.
因為我想他可能怕你弄壞了他的車.
於是在一個很偏僻的地方.
我們就打去租車公司.
將我們的定位告訴他.
跟他說我們在這裡爆胎.
他說你不用擔心.

$^{521}$我們會派人來修的.
但是接著問等多久.
六個小時.
因為一個地方很偏僻.
就要等六個小時.
六個小時做些什麼好呢.
於是我們就開始上Google Map看.
然後我們就發現附近有一個.
足部反射治療的地方.
於是我們就.
啊有足部反射.
有六個小時.
就去那個足部反射那裡看看.
有沒有足部的治療.
可以去等這六個小時.
當我去到的時候.
我有一個朋友很雀躍.
他跟我說.
哇.
原來這個是吳藥石神父.
足部反射健康法的地方.
有沒有人聽過吳藥石神父.
足部反射健康法.
也有些人聽過.
原來也挺出名的.
但我不知道.
但是就很開心.
然後我就上網查.
原來他說吳藥石神父.
是鑽研腳底按摩二十多年.
他最初就是想治療自己的病痛.
但是當他鑽研越深的時候.
他就發現一些更加新更加有用.
以及真的能夠保障人體健康.
這個足部反射健康法.
於是這位吳神父.
他就一邊傳教.
一邊教台灣人.
如何用這套健康法去自我保健.
以及預防疾病.

$^{561}$你知道我一無所知.
我就說我們七個人.
能不能做治療.
他說七個人要等很久.
我說不要緊.
我有六個小時.
他說那你一個一個等吧.
我的朋友就先去做.
我就在這裡等.
沒什麼事做.
那裡這麼偏僻.
有棵樹.
我就坐在那棵樹旁邊.
旁邊有一位老人家.
他就這樣穿著毛衣.
沒有在看著前面.
我坐在他旁邊.
沒什麼事做.
就跟他聊天.
他就問我.
當然他用國語問我的.
我用我的國語跟他嘗試溝通.
他就問你為什麼會在這裡.
我說沒有.
因為我們爆了胎.
爆了胎等六個小時.
他說你怎麼找到這裡.
我說沒有.
上Google Map說有足部反射健康法.
就來了.
但是來到因為不是很多位.
於是我的朋友就去了.
我們有六個小時.
然後就在這裡聊天.
聊了一段時間.
然後去到一個位置.
我就問這位老人家.
他一直都是這樣坐著.
沒有在看著前面.
然後我就問他.

$^{601}$其實怎麼稱呼你.
他說我是吳藥石神父.
我說什麼.
你是吳藥石神父.
你就是這個健康法的始創人.
然後對於我來說.
其實你告訴我.
你叫做吳藥石神父.
對於我來說.
你只是一位老人家.
但是我旁邊的朋友.
是開心到不得了.
吳神父啊.
我想學了這套健康法很久了.
你可不可以教我.
怎麼學這套健康法.
吳神父說好.
其實我們有不同的訓練課程.
你可以來台灣上這個課程.
我最近再跟這個朋友吃飯.
他已經去了台灣三次去學這個課程.
他已經是初級,中級,高級都學了.
然後我就問.
其實學了有什麼用.
原來很有用的.
原來他用這個足部健康法.
去服侍教會的老人家.
不僅如此.
他說這是一個很有效的傳福音工具.
因為他去探望那些未信的老人家.
那些老人家怎麼開門給他呢.
他就說.
我幫你做足底按摩.
那些老人家說.
需不需要錢的.
不用的,免費的.
他就進去幫人做這個傳福音.
真的厲害.
不僅能夠祝福老人家.
還更加傳福音.

$^{641}$所以長者崇拜的義工.
可以考慮一下去台灣學.
你聽到一件事嗎.
當你真的想生命改變的時候.
你見到那個人.
你會跟他說我想生命改變.
但如果你不想生命改變.
那個人跟你是沒有關係的.
對於我來說.
我坐在吳約石神父旁邊這麼久.
就算我知道他是吳約石神父.
我的反應是什麼.
Nice to meet you.
對於我來說.
吳約石神父只是Nice to meet you.
但對於我的朋友.
這是他學習的好機會.
是他生命轉化.
服侍老人家的機會.
所以今天我們要問自己.
我們知不知道耶穌是誰.
當你知道耶穌是誰的時候.
你是否會說.
耶穌Nice to meet you.
還是你說.
不是啊耶穌.
我想生命改變.
請你改變我的生命.
所以今天你要問.
耶穌是你的誰.
祂是不是只是一個Nice to meet you的人.
還是你說.
耶穌啊.
我向你呼求.
我想生命改變.
所以今天我們要問自己.
一個很重要的問題.
我們有沒有認清楚.
耶穌是誰.
耶穌對於我們來說是誰.

$^{681}$而我們今天.
如何對耶穌.
作出一個生命的回應.
今天我們需不需要耶穌的憐憫.
我們每一個人都需要耶穌的憐憫.
因為我們沒有一個人.
可以靠自己來得救.
今天我們的生命.
每一個人有很多需要改變的地方.
當我們遇見耶穌的時候.
我們有沒有求耶穌.
去改變我們的生命.
或者繼續改變我們的生命.
今天我們有沒有和耶穌呼求.
耶穌啊求你幫我.
甚至有些人可能跟你說.
你不要求耶穌了.
你用得著求耶穌嗎.
耶穌能幫到你嗎.
甚至有些人說.
你不要出聲.
但是今天究竟.
當人們都叫我們不要找耶穌的時候.
我們的生命有沒有繼續去呼求耶穌.
然後求耶穌.
去繼續改變我們的生命.
而今天我們去到第三點.
就是當我們去呼求耶穌.
我們的生命經歷了耶穌的醫治.
我們就歸榮耀給神了.
我們讀最後四十字.
四十三字好不好.
一二三.
耶穌站住.
吩咐把他領過來.
到了庚前就問他說.
你要我為你做什麼.
他說主啊我要能看見.
耶穌說你可以看見.
你的信救了你.

$^{721}$何止立刻看見了就跟隨耶穌.
一路歸榮耀與神.
眾人看見這事也讚美神.
當我們看當一個人呼求耶穌的時候.
耶穌的回應是什麼呢.
耶穌在那裡告訴我們.
他站住了.
他停了站在那裡.
我們不要小看這一句的句子.
因為耶穌本來在那裡.
在耶路撒冷.
他有沒有使命去那裡.
有的.
今天你有沒有試過很趕去一個地方.
要完成那些事情.
有人問你事情的時候.
你會怎樣.
我很忙.
我現在去一個地方.
很多事情要做.
很重要的.
你不要搞我.
你會這樣說的.
但是當一個人呼求耶穌的時候.
耶穌為他停下來.
不單止為他停下來.
更加叫人將他帶到自己的面前.
他不是很遠地說.
為什麼.
為什麼幫到你.
為什麼.
他不是這樣.
他是將他帶到他的面前.
耶穌原來和我們是面對面近距離的.
我們不要忘記在當時的文化.
這個乞丐是被社會忽略的.
這個社會沒有人理他的.
但是耶穌叫人將這個乞丐.
帶到自己的面前.
面對面近距離的接觸.

$^{761}$接著耶穌在那裡問他一個問題.
你要我為你做什麼.
這個問題.
不是耶穌不知道他需要什麼.
但是很想這個人.
能夠很明白地將自己的需要.
帶到耶穌的面前.
今天其實很多人問我一個問題.
為什麼我們要祈禱.
耶穌不知道我的需要嗎.
耶穌讓我更加知道我的需要.
但是為什麼今天我們要祈禱.
因為我們看不看到我們自己的需要.
我們看不看到自己的需要.
並且將我們的需要帶到耶穌的面前.
所以我們看看這個盲眼的人.
他怎樣回應.
他說:主啊!我要能看見.
當他這樣去呼求耶穌的時候.
耶穌的回應就是.
你可以看見.
其實很奇妙的.
今天如果他說想看見.
耶穌回答他你可以看見就可以了.
但是耶穌沒有停在那裡.
耶穌說你的信救了你.
救這個字的原文.
不只是代表一個肉體上的醫治.
它包括了我們全人.
我們的靈魂.
我們的全人得救.
但是這個盒子.
怎樣能夠得到這個全人得救的醫治呢.
耶穌說是因為你的信.
但是如果我們看回這個盲眼的人的生命.
他的信代表什麼呢.
是他首先聽到耶穌是誰.
然後他認清楚耶穌是誰.
然後向耶穌呼求.
今天我們有沒有這種信呢.

$^{801}$今天我們有沒有聽到耶穌是誰呢.
應該有的.
因為不然的話你不會回教會.
但是我們有沒有認清楚耶穌是誰呢.
有的.
因為大家都回來聽到.
希望牧師有沒有解鎖經給你聽.
那你就知道耶穌是誰了.
但是你有沒有呼求耶穌去改變你的生命.
是唯有這個盲眼的人.
他聽到耶穌是誰.
他認清楚耶穌是誰.
並且他呼求耶穌.
這個就是他的信.
今天我們有沒有這種信呢.
今天我們是不是停留.
耶穌,很高興認識你.
只是我們跟耶穌說.
主啊,請你救我.
當我們有這種信的時候.
原來這裡耶穌就說.
你的信救了你.
這個救不只是肉體的醫治.
是全人的得救.
然後你說這個盲眼的人.
怎麼全人的得救呢.
我們看回43節就知道了.
43節這裡說.
這個盲眼的人立刻看見.
如果只是肉身得著醫治.
聖經停在這裡就可以了.
但是接著他不只是停在這裡.
這個人不只是他的眼睛立刻看見.
他的生命是跟隨耶穌.
跟著一路歸榮耀於神.
所以我們看到這個盲眼的人.
他不只是他的眼睛得到醫治.
他是全人整個生命翻轉了.
他由一個路邊的乞丐.
變成一個耶穌的跟隨者.

$^{841}$不只是這樣.
最後群眾見證的是.
眾人看見這事也讚美神.
原來其他人見到這個人.
他這樣經歷了耶穌之後.
群眾也歸榮耀給耶穌.
原來今天我們也要承認.
有時我們的心是合了眼的.
有沒有人現在眼睛是有問題的.
我幫你試一試.
肥成肥不肥.
如果你說不肥.
你就去檢驗眼睛.
你說肥.
你的眼睛就沒問題了.
今天大部份的人.
你的眼睛應該也沒問題.
但是我們要問自己.
我們的心會不會盲了眼.
今天我們可能不會處理我們的金錢觀.
今天我們可能不會處理我們的感情關係.
今天我們可能不會為主作見證.
今天我們可能攪亂了我們的自我價值.
但是今天我們有沒有回到耶穌的面前.
說主啊.
我要看見.
一個真正的醫治.
不是單單問題的解決.
我們很多時候回到教會.
我們也想耶穌解決我們的問題.
當耶穌解決我們的問題的時候.
但原來一個真正的醫治.
不是我們頭痛醫頭.
腳痛醫腳.
而是我們的生命有一個真正的轉向.
一個盲眼的人.
他不單止眼睛看得見.
他經歷一個生命的轉向.
他的生命跟隨耶穌.
就歸榮耀給神.

$^{881}$原來今天.
當我們經歷神的改變.
經歷神的醫治.
我們的生命都歸榮耀給神.
在《門女》裡面.
我們每一次差不多去到完成的階段.
我們都會問門徒.
你去問身邊的人.
當你經歷了《門女》之後.
你的生命有什麼改變.
而每一次當我們聽到不同的人.
分享的見證.
我們都很感恩.
你都真的歸榮耀給神.
大家知道我們現在不是只有雅庚堂有《門女》.
陳牧師都帶領我們牧者.
在外面有牧者《門女》.
就是說我們同一班的牧者.
一班的領袖一起都做《門女》.
而在最近有一個.
我有份參與的牧者《門女》裡面.
我們都經歷了這個階段.
因為我們差不多完成了.
於是我們就問牧者和領袖.
你回去.
你最親的人問你有什麼改變.
我們聽回來很感動.
有一位牧者.
他說他的配偶說.
我覺得我的配偶更加懂得愛我.
他更加甜蜜.
他更加懂得關心我.
了解我的需要.
所以我上個星期說.
如果你的配偶說做《門女》.
你不要擔心他不理你.
你不要擔心他沒有時間.
你鼓勵他去做.
因為他做了他會更加愛你.
他更加懂得關心你.

$^{921}$更加懂得聆聽你的需要.
另外一個就說他的兒女.
他說以前.
作為母親的人.
不懂得關心我們.
但是現在他很懂得和我們聊天.
很懂得關心我們.
連脾氣都變了.
所以你不要不鼓勵你的家人做《門女》.
因為他會更加懂得關心你.
對你好.
脾氣都變了.
當我們發現這一班牧者是領袖.
我們都繼續經歷改變.
歸榮耀給神.
我們今天的生命.
有沒有繼續經歷神的帶領和醫治呢?.
正正昨天.
我鼓勵一位少年人.
去參加我們的《門女》簡介會.
我們成人也有《門女》誓師大會.
少年也有《門女》簡介會.
我鼓勵他.
他說他最初真的沒有考慮過.
但是他說他後來很認真地考慮.
我說為什麼?.
他說因為我的父母做了《門女》.
我看到他們有很大的改變.
是的.
原來我們的生命改變的時候.
我們的生命都歸榮耀給神.
所以今天我們要問自己一個問題.
就是我們的生命.
如何繼續被神帶領改變和醫治.
讓我們的生命都繼續歸榮耀給神.
我們作為牧者.
我們作為領袖.
我們的生命都需要繼續被神改變.
我們仍然有很多需要改變的地方.
求神去醫治去改變我們的.

$^{961}$而當我們繼續的生命被醫治被改變.
你身邊的人見到就歸榮耀給神.
所以今天我們有沒有首先呼求耶穌.
繼續去改變我們的生命.
但是當我們的生命繼續被耶穌改變.
我們的生命就能夠歸榮耀給神.
在總結的時候我們請勁大隊.
這一段聖經給我們一個很重要的功課.
當我們首先看耶穌的生命.
耶穌祂不是走一條逃避受苦的道路.
今天祂是堅定去走上十字架.
所以今天我們都要問自己.
當我們說我們的生命跟隨耶穌.
我們預備好了沒有.
跟隨耶穌是要背十字架去跟隨耶穌.
跟隨耶穌不是追求舒適.
也不是逃避苦難.
在這世上我們會有苦難.
耶穌說的.
不過他說你不用擔心.
為什麼.
耶穌說我已經勝過世界.
但是今天當我們跟隨耶穌的時候.
我們有沒有預備好背十字架去跟隨耶穌.
我們每一個人背的也許都不一樣.
但是今天我們願不願意去背十字架.
去跟隨耶穌.
當我們看這個盲眼的人.
這個盲眼的人他認清楚耶穌是誰.
並且呼求耶穌去改變他的生命.
當他這樣呼求的時候.
他不只是問題解決.
頭痛醫頭腳痛醫痛.
他的眼睛看得到.
更加重要的是他整個的生命反轉.
今天原來我們的生命不是知道耶穌.
是要跟隨耶穌.
並且在背十字架的路上去學習.
信靠耶穌讓我們的生命.
都是一個榮耀神的見證.

$^{1001}$這裡很奇怪.
一個盲眼的人.
看到一些不盲眼的人看到的東西.
但是這個盲眼的人看到.
一些不盲眼的人看不到的東西的時候.
叫我們的眼睛看到.
但是心可能看不到的人.
我們的生命都改變.
所以今天我們要問問自己.
可能在某一些層面上.
我們都是伶侶盲眼的.
可能我們不懂得怎樣去處理我們的家庭.
可能我們都不懂得處理我們的婚姻.
我們不懂得怎樣去運用我們的金錢.
或者我們不懂得怎樣在職場上去做見證.
又或者我們不懂得怎樣去搞我們的人際關係.
我們可能都被我們自己的驕傲.
被我們的妥協.
或者被我們的恐懼所捆綁.
但是今天我們都看清楚我們自己個人的需要.
但是都讓我們一起跟耶穌講主啊.
我要能看見.
當我們這樣的時候.
我們發現我們的敬拜不再只是停留在崇拜裡面.
而是我們每一天.
我們的生活.
我們都在經歷神的掌權和意志.
我們想不想我們每一天都在經歷神的掌權和意志.
唯有我們每一天去認清耶穌.
並且向祂呼求.
並且願意去跟隨十字架的路.
我們的生命就能夠一路歸榮耀給神.
\newpage



\section{出埃及記 17:8-16}
\label{sec:032041}
\textbf{托手zoung}~~~~梁永善牧師
\newline
\newline
連結: \href{http://www.church.com.hk/acms/content.asp?site=cdc\&op=show\&type=product\&code=032041\&layout=sermon}{\texttt{032041}} ~~~~ 語音日期: 2026-04-26
\newline
\newline
\hyperref[sec:032040]{\small{< < < PREV SERMON < < <}}
~
\hyperref[sec:index_chronic]{\small{[返順時目]}}
~
\hyperref[sec:032042]{\small{> > > NEXT SERMON > > >}}
\newline
\newline
出埃及記 17:8-16
\newline
\begin{longtable}{cl}
\hline
\hline
章節 & 經文 (和合本修訂版)\\
\hline
17:8 & \begin{tabularx}{0.7\textwidth}{X} 那時，亞瑪力來到利非訂，和以色列爭戰。 \end{tabularx} \\ \\ \relax
17:9 & \begin{tabularx}{0.7\textwidth}{X} 摩西對約書亞說：「你為我們選出人來，出去和亞瑪力爭戰。明天我要站在山頂上，手裡拿著神的杖。」 \end{tabularx} \\ \\ \relax
17:10 & \begin{tabularx}{0.7\textwidth}{X} 於是，約書亞照著摩西對他所說的話去做，和亞瑪力爭戰。摩西、亞倫和戶珥都上了山頂。 \end{tabularx} \\ \\ \relax
17:11 & \begin{tabularx}{0.7\textwidth}{X} 摩西何時舉手，以色列就得勝；何時垂手，亞瑪力就得勝。 \end{tabularx} \\ \\ \relax
17:12 & \begin{tabularx}{0.7\textwidth}{X} 但摩西的雙手沉重，他們就搬一塊石頭來放在他下面，他就坐在上面。亞倫與戶珥扶著他的手，一個在這邊，一個在那邊，他的手就穩住，直到日落。 \end{tabularx} \\ \\ \relax
17:13 & \begin{tabularx}{0.7\textwidth}{X} 約書亞用刀打敗了亞瑪力和他的百姓。 \end{tabularx} \\ \\ \relax
17:14 & \begin{tabularx}{0.7\textwidth}{X} 耶和華對摩西說：「你要把這事記錄在書上作紀念，又念給約書亞聽：我要把亞瑪力的名字從天下全然塗去。」 \end{tabularx} \\ \\ \relax
17:15 & \begin{tabularx}{0.7\textwidth}{X} 摩西築了一座壇，起名叫「耶和華尼西」。 \end{tabularx} \\ \\ \relax
17:16 & \begin{tabularx}{0.7\textwidth}{X} 他說：「我指著耶和華的寶座發誓，耶和華必世世代代和亞瑪力爭戰。」 \end{tabularx} \\ \\
[1ex]
\hline
\hline
\end{longtable}
$^{1}$我用了一百多年的歷史和歷史經驗.
來寫文章.
我用了一百多年歷史和歷史經驗.
來寫文章.
我用了一百多年歷史經驗.
來寫文章.
我用了一百多年歷史經驗.
來寫文章.
我用了一百多年歷史經驗.
來寫文章.
我用了一百多年歷史經驗.
來寫文章.
我用了一百多年歷史經驗.
來寫文章.
我用了一百多年歷史經驗.
來寫文章.
我用了一百多年歷史經驗.
來寫文章.
我用了一百多年歷史經驗.
來寫文章.
我用了一百多年歷史經驗.
來寫文章.
我用了一百多年歷史經驗.
來寫文章.
我用了一百多年歷史經驗.
來寫文章.
我用了一百多年歷史經驗.
來寫文章.
我用了一百多年歷史經驗.
來寫文章.
我用了一百多年歷史經驗.
來寫文章.
我用了一百多年歷史經驗.
來寫文章.
我用了一百多年歷史經驗.
來寫文章.
我用了一百多年歷史經驗.
來寫文章.
我用了一百多年歷史經驗.
來寫文章.

$^{41}$我用了一百多年歷史經驗.
來寫文章.
我用了一百多年歷史經驗.
來寫文章.
我用了一百多年歷史經驗.
來寫文章.
我用了一百多年歷史經驗.
來寫文章.
我用了一百多年歷史經驗.
來寫文章.
我用了一百多年歷史經驗.
來寫文章.
我用了一百多年歷史經驗.
來寫文章.
我用了一百多年歷史經驗.
來寫文章.
我用了一百多年歷史經驗.
來寫文章.
我用了一百多年歷史經驗.
來寫文章.
我用了一百多年歷史經驗.
來寫文章.
我用了一百多年歷史經驗.
來寫文章.
我用了一百多年歷史經驗.
來寫文章.
我用了一百多年歷史經驗.
來寫文章.
我用了一百多年歷史經驗.
來寫文章.
我用了一百多年歷史經驗.
來寫文章.
我用了一百多年歷史經驗.
來寫文章.
我用了一百多年歷史經驗.
來寫文章.
我用了一百多年歷史經驗.
來寫文章.
我用了一百多年歷史經驗.
來寫文章.

$^{81}$我用了一百多年歷史經驗.
來寫文章.
我用了一百多年歷史經驗.
來寫文章.
我用了一百多年歷史經驗.
來寫文章.
我用了一百多年歷史經驗.
來寫文章.
我用了一百多年歷史經驗.
來寫文章.
我用了一百多年歷史經驗.
來寫文章.
我用了一百多年歷史經驗.
來寫文章.
我用了一百多年歷史經驗.
來寫文章.
我用了一百多年歷史經驗.
來寫文章.
我用了一百多年歷史經驗.
來寫文章.
我用了一百多年歷史經驗.
來寫文章.
我用了一百多年歷史經驗.
來寫文章.
我用了一百多年歷史經驗.
來寫文章.
我用了一百多年歷史經驗.
來寫文章.
我用了一百多年歷史經驗.
來寫文章.
我用了一百多年歷史經驗.
來寫文章.
我用了一百多年歷史經驗.
來寫文章.
我用了一百多年歷史經驗.
來寫文章.
我用了一百多年歷史經驗.
來寫文章.
我用了一百多年歷史經驗.
來寫文章.

$^{121}$我用了一百多年歷史經驗.
來寫文章.
我用了一百多年歷史經驗.
來寫文章.
我用了一百多年歷史經驗.
來寫文章.
我用了一百多年歷史經驗.
來寫文章.
我用了一百多年歷史經驗.
來寫文章.
我用了一百多年歷史經驗.
來寫文章.
我用了一百多年歷史經驗.
來寫文章.
我用了一百多年歷史經驗.
來寫文章.
我用了一百多年歷史經驗.
來寫文章.
我用了一百多年歷史經驗.
來寫文章.
我用了一百多年歷史經驗.
來寫文章.
我用了一百多年歷史經驗.
來寫文章.
我用了一百多年歷史經驗.
來寫文章.
我用了一百多年歷史經驗.
來寫文章.
我用了一百多年歷史經驗.
來寫文章.
我用了一百多年歷史經驗.
來寫文章.
我用了一百多年歷史經驗.
來寫文章.
我用了一百多年歷史經驗.
來寫文章.
我用了一百多年歷史經驗.
來寫文章.
我用了一百多年歷史經驗.
來寫文章.

$^{161}$我用了一百多年歷史經驗.
來寫文章.
我用了一百多年歷史經驗.
來寫文章.
我用了一百多年歷史經驗.
來寫文章.
我用了一百多年歷史經驗.
來寫文章.
我用了一百多年歷史經驗.
來寫文章.
我用了一百多年歷史經驗.
來寫文章.
我用了一百多年歷史經驗.
來寫文章.
我用了一百多年歷史經驗.
來寫文章.
我用了一百多年歷史經驗.
來寫文章.
我用了一百多年歷史經驗.
來寫文章.
我用了一百多年歷史經驗.
來寫文章.
我用了一百多年歷史經驗.
來寫文章.
我用了一百多年歷史經驗.
來寫文章.
我用了一百多年歷史經驗.
來寫文章.
我用了一百多年歷史經驗.
來寫文章.
我用了一百多年歷史經驗.
來寫文章.
我用了一百多年歷史經驗.
來寫文章.
我用了一百多年歷史經驗.
來寫文章.
我用了一百多年歷史經驗.
來寫文章.
我用了一百多年歷史經驗.
來寫文章.

$^{201}$我用了一百多年歷史經驗.
來寫文章.
我用了一百多年歷史經驗.
來寫文章.
我用了一百多年歷史經驗.
來寫文章.
我用了一百多年歷史經驗.
來寫文章.
我用了一百多年歷史經驗.
來寫文章.
我用了一百多年歷史經驗.
來寫文章.
我用了一百多年歷史經驗.
來寫文章.
我用了一百多年歷史經驗.
來寫文章.
我用了一百多年歷史經驗.
來寫文章.
我用了一百多年歷史經驗.
來寫文章.
我用了一百多年歷史經驗.
來寫文章.
我用了一百多年歷史經驗.
來寫文章.
我用了一百多年歷史經驗.
來寫文章.
我用了一百多年歷史經驗.
來寫文章.
我用了一百多年歷史經驗.
來寫文章.
我用了一百多年歷史經驗.
來寫文章.
我用了一百多年歷史經驗.
來寫文章.
我用了一百多年歷史經驗.
來寫文章.
我用了一百多年歷史經驗.
來寫文章.
我用了一百多年歷史經驗.
來寫文章.

$^{241}$我用了一百多年歷史經驗.
來寫文章.
我用了一百多年歷史經驗.
來寫文章.
我用了一百多年歷史經驗.
來寫文章.
我用了一百多年歷史經驗.
來寫文章.
我用了一百多年歷史經驗.
來寫文章.
我用了一百多年歷史經驗.
來寫文章.
我用了一百多年歷史經驗.
來寫文章.
我用了一百多年歷史經驗.
來寫文章.
我用了一百多年歷史經驗.
來寫文章.
我用了一百多年歷史經驗.
來寫文章.
我用了一百多年歷史經驗.
來寫文章.
我用了一百多年歷史經驗.
來寫文章.
我用了一百多年歷史經驗.
來寫文章.
我用了一百多年歷史經驗.
來寫文章.
我用了一百多年歷史經驗.
來寫文章.
我用了一百多年歷史經驗.
來寫文章.
我用了一百多年歷史經驗.
來寫文章.
我用了一百多年歷史經驗.
來寫文章.
我用了一百多年歷史經驗.
來寫文章.
我用了一百多年歷史經驗.
來寫文章.

$^{281}$我用了一百多年歷史經驗.
來寫文章.
我用了一百多年歷史經驗.
來寫文章.
我用了一百多年歷史經驗.
來寫文章.
我用了一百多年歷史經驗.
來寫文章.
我用了一百多年歷史經驗.
來寫文章.
我用了一百多年歷史經驗.
來寫文章.
我用了一百多年歷史經驗.
來寫文章.
我用了一百多年歷史經驗.
來寫文章.
我用了一百多年歷史經驗.
來寫文章.
我用了一百多年歷史經驗.
來寫文章.
我用了一百多年歷史經驗.
來寫文章.
我用了一百多年歷史經驗.
來寫文章.
我用了一百多年歷史經驗.
來寫文章.
我用了一百多年歷史經驗.
來寫文章.
我用了一百多年歷史經驗.
來寫文章.
我用了一百多年歷史經驗.
來寫文章.
我用了一百多年歷史經驗.
來寫文章.
我用了一百多年歷史經驗.
來寫文章.
我用了一百多年歷史經驗.
來寫文章.
我用了一百多年歷史經驗.
來寫文章.

$^{321}$我用了一百多年歷史經驗.
來寫文章.
我用了一百多年歷史經驗.
來寫文章.
我用了一百多年歷史經驗.
來寫文章.
我用了一百多年歷史經驗.
來寫文章.
我用了一百多年歷史經驗.
來寫文章.
我用了一百多年歷史經驗.
來寫文章.
我用了一百多年歷史經驗.
來寫文章.
我用了一百多年歷史經驗.
來寫文章.
我用了一百多年歷史經驗.
來寫文章.
我用了一百多年歷史經驗.
來寫文章.
我用了一百多年歷史經驗.
來寫文章.
我用了一百多年歷史經驗.
來寫文章.
我用了一百多年歷史經驗.
來寫文章.
我用了一百多年歷史經驗.
來寫文章.
我用了一百多年歷史經驗.
來寫文章.
我用了一百多年歷史經驗.
來寫文章.
我用了一百多年歷史經驗.
來寫文章.
我用了一百多年歷史經驗.
來寫文章.
我用了一百多年歷史經驗.
來寫文章.
我用了一百多年歷史經驗.
來寫文章.

$^{361}$我用了一百多年歷史經驗.
來寫文章.
我用了一百多年歷史經驗.
來寫文章.
我用了一百多年歷史經驗.
來寫文章.
我用了一百多年歷史經驗.
來寫文章.
我用了一百多年歷史經驗.
來寫文章.
我用了一百多年歷史經驗.
來寫文章.
我用了一百多年歷史經驗.
來寫文章.
我用了一百多年歷史經驗.
來寫文章.
我用了一百多年歷史經驗.
來寫文章.
我用了一百多年歷史經驗.
來寫文章.
我用了一百多年歷史經驗.
來寫文章.
我用了一百多年歷史經驗.
來寫文章.
我用了一百多年歷史經驗.
來寫文章.
我用了一百多年歷史經驗.
來寫文章.
我用了一百多年歷史經驗.
來寫文章.
我用了一百多年歷史經驗.
來寫文章.
我用了一百多年歷史經驗.
來寫文章.
我用了一百多年歷史經驗.
來寫文章.
我用了一百多年歷史經驗.
來寫文章.
我用了一百多年歷史經驗.
來寫文章.

$^{401}$我用了一百多年歷史經驗.
來寫文章.
我用了一百多年歷史經驗.
來寫文章.
我用了一百多年歷史經驗.
來寫文章.
我用了一百多年歷史經驗.
來寫文章.
我用了一百多年歷史經驗.
來寫文章.
我用了一百多年歷史經驗.
來寫文章.
我用了一百多年歷史經驗.
來寫文章.
我用了一百多年歷史經驗.
來寫文章.
我用了一百多年歷史經驗.
來寫文章.
我用了一百多年歷史經驗.
來寫文章.
我用了一百多年歷史經驗.
來寫文章.
我用了一百多年歷史經驗.
來寫文章.
我用了一百多年歷史經驗.
來寫文章.
我用了一百多年歷史經驗.
來寫文章.
我用了一百多年歷史經驗.
來寫文章.
我用了一百多年歷史經驗.
來寫文章.
我用了一百多年歷史經驗.
來寫文章.
我用了一百多年歷史經驗.
來寫文章.
我用了一百多年歷史經驗.
來寫文章.
我用了一百多年歷史經驗.
來寫文章.

$^{441}$我用了一百多年歷史經驗.
來寫文章.
我用了一百多年歷史經驗.
來寫文章.
我用了一百多年歷史經驗.
來寫文章.
我用了一百多年歷史經驗.
來寫文章.
我用了一百多年歷史經驗.
來寫文章.
我用了一百多年歷史經驗.
來寫文章.
我用了一百多年歷史經驗.
來寫文章.
我用了一百多年歷史經驗.
來寫文章.
我用了一百多年歷史經驗.
來寫文章.
我用了一百多年歷史經驗.
來寫文章.
我用了一百多年歷史經驗.
來寫文章.
我用了一百多年歷史經驗.
來寫文章.
我用了一百多年歷史經驗.
來寫文章.
我用了一百多年歷史經驗.
來寫文章.
我用了一百多年歷史經驗.
來寫文章.
我用了一百多年歷史經驗.
來寫文章.
我用了一百多年歷史經驗.
來寫文章.
我用了一百多年歷史經驗.
來寫文章.
我用了一百多年歷史經驗.
來寫文章.
我用了一百多年歷史經驗.
來寫文章.

$^{481}$我用了一百多年歷史經驗.
來寫文章.
我用了一百多年歷史經驗.
來寫文章.
我用了一百多年歷史經驗.
來寫文章.
我用了一百多年歷史經驗.
來寫文章.
我用了一百多年歷史經驗.
來寫文章.
我用了一百多年歷史經驗.
來寫文章.
我用了一百多年歷史經驗.
來寫文章.
我用了一百多年歷史經驗.
來寫文章.
我用了一百多年歷史經驗.
來寫文章.
我用了一百多年歷史經驗.
來寫文章.
我用了一百多年歷史經驗.
來寫文章.
我用了一百多年歷史經驗.
來寫文章.
我用了一百多年歷史經驗.
來寫文章.
我用了一百多年歷史經驗.
來寫文章.
我用了一百多年歷史經驗.
來寫文章.
我用了一百多年歷史經驗.
來寫文章.
我用了一百多年歷史經驗.
來寫文章.
我用了一百多年歷史經驗.
來寫文章.
我用了一百多年歷史經驗.
來寫文章.
我用了一百多年歷史經驗.
來寫文章.

$^{521}$我用了一百多年歷史經驗.
來寫文章.
我用了一百多年歷史經驗.
來寫文章.
我用了一百多年歷史經驗.
來寫文章.
我用了一百多年歷史經驗.
來寫文章.
我用了一百多年歷史經驗.
來寫文章.
我用了一百多年歷史經驗.
來寫文章.
我用了一百多年歷史經驗.
來寫文章.
我用了一百多年歷史經驗.
來寫文章.
我用了一百多年歷史經驗.
來寫文章.
我用了一百多年歷史經驗.
來寫文章.
我用了一百多年歷史經驗.
來寫文章.
我用了一百多年歷史經驗.
來寫文章.
我用了一百多年歷史經驗.
來寫文章.
我用了一百多年歷史經驗.
來寫文章.
我用了一百多年歷史經驗.
來寫文章.
我用了一百多年歷史經驗.
來寫文章.
我用了一百多年歷史經驗.
來寫文章.
我用了一百多年歷史經驗.
來寫文章.
我用了一百多年歷史經驗.
來寫文章.
我用了一百多年歷史經驗.
來寫文章.

$^{561}$我用了一百多年歷史經驗.
來寫文章.
我用了一百多年歷史經驗.
來寫文章.
我用了一百多年歷史經驗.
來寫文章.
我用了一百多年歷史經驗.
來寫文章.
我用了一百多年歷史經驗.
來寫文章.
我用了一百多年歷史經驗.
來寫文章.
我用了一百多年歷史經驗.
來寫文章.
我用了一百多年歷史經驗.
來寫文章.
我用了一百多年歷史經驗.
來寫文章.
我用了一百多年歷史經驗.
來寫文章.
我用了一百多年歷史經驗.
來寫文章.
我用了一百多年歷史經驗.
來寫文章.
我用了一百多年歷史經驗.
來寫文章.
我用了一百多年歷史經驗.
來寫文章.
我用了一百多年歷史經驗.
來寫文章.
我用了一百多年歷史經驗.
來寫文章.
我用了一百多年歷史經驗.
來寫文章.
我用了一百多年歷史經驗.
來寫文章.
我用了一百多年歷史經驗.
來寫文章.
我用了一百多年歷史經驗.
來寫文章.

$^{601}$我用了一百多年歷史經驗.
來寫文章.
我用了一百多年歷史經驗.
來寫文章.
我用了一百多年歷史經驗.
來寫文章.
我用了一百多年歷史經驗.
來寫文章.
我用了一百多年歷史經驗.
來寫文章.
我用了一百多年歷史經驗.
來寫文章.
我用了一百多年歷史經驗.
來寫文章.
我用了一百多年歷史經驗.
來寫文章.
我用了一百多年歷史經驗.
來寫文章.
我用了一百多年歷史經驗.
來寫文章.
我用了一百多年歷史經驗.
來寫文章.
我用了一百多年歷史經驗.
來寫文章.
我用了一百多年歷史經驗.
來寫文章.
我用了一百多年歷史經驗.
來寫文章.
我用了一百多年歷史經驗.
來寫文章.
我用了一百多年歷史經驗.
來寫文章.
我用了一百多年歷史經驗.
來寫文章.
我用了一百多年歷史經驗.
來寫文章.
我用了一百多年歷史經驗.
來寫文章.
我用了一百多年歷史經驗.
來寫文章.

$^{641}$我用了一百多年歷史經驗.
來寫文章.
我用了一百多年歷史經驗.
來寫文章.
我用了一百多年歷史經驗.
來寫文章.
我用了一百多年歷史經驗.
來寫文章.
我用了一百多年歷史經驗.
來寫文章.
我用了一百多年歷史經驗.
來寫文章.
我用了一百多年歷史經驗.
來寫文章.
我用了一百多年歷史經驗.
來寫文章.
我用了一百多年歷史經驗.
來寫文章.
我用了一百多年歷史經驗.
來寫文章.
我用了一百多年歷史經驗.
來寫文章.
我用了一百多年歷史經驗.
來寫文章.
我用了一百多年歷史經驗.
來寫文章.
我用了一百多年歷史經驗.
來寫文章.
我用了一百多年歷史經驗.
來寫文章.
我用了一百多年歷史經驗.
來寫文章.
我用了一百多年歷史經驗.
來寫文章.
我用了一百多年歷史經驗.
來寫文章.
我用了一百多年歷史經驗.
來寫文章.
我用了一百多年歷史經驗.
來寫文章.

$^{681}$我用了一百多年歷史經驗.
來寫文章.
我用了一百多年歷史經驗.
來寫文章.
我用了一百多年歷史經驗.
來寫文章.
我用了一百多年歷史經驗.
來寫文章.
我用了一百多年歷史經驗.
來寫文章.
我用了一百多年歷史經驗.
來寫文章.
我用了一百多年歷史經驗.
來寫文章.
我用了一百多年歷史經驗.
來寫文章.
我用了一百多年歷史經驗.
來寫文章.
我用了一百多年歷史經驗.
來寫文章.
我用了一百多年歷史經驗.
來寫文章.
我用了一百多年歷史經驗.
來寫文章.
我用了一百多年歷史經驗.
來寫文章.
我用了一百多年歷史經驗.
來寫文章.
我用了一百多年歷史經驗.
來寫文章.
我用了一百多年歷史經驗.
來寫文章.
我用了一百多年歷史經驗.
來寫文章.
我用了一百多年歷史經驗.
來寫文章.
我用了一百多年歷史經驗.
來寫文章.
我用了一百多年歷史經驗.
來寫文章.

$^{721}$我用了一百多年歷史經驗.
來寫文章.
我用了一百多年歷史經驗.
來寫文章.
我用了一百多年歷史經驗.
來寫文章.
我用了一百多年歷史經驗.
來寫文章.
我用了一百多年歷史經驗.
來寫文章.
我用了一百多年歷史經驗.
來寫文章.
我用了一百多年歷史經驗.
來寫文章.
我用了一百多年歷史經驗.
來寫文章.
我用了一百多年歷史經驗.
來寫文章.
我用了一百多年歷史經驗.
來寫文章.
我用了一百多年歷史經驗.
來寫文章.
我用了一百多年歷史經驗.
來寫文章.
我用了一百多年歷史經驗.
來寫文章.
我用了一百多年歷史經驗.
來寫文章.
我用了一百多年歷史經驗.
來寫文章.
我用了一百多年歷史經驗.
來寫文章.
我用了一百多年歷史經驗.
來寫文章.
我用了一百多年歷史經驗.
來寫文章.
我用了一百多年歷史經驗.
來寫文章.
我用了一百多年歷史經驗.
來寫文章.

$^{761}$我用了一百多年歷史經驗.
來寫文章.
我用了一百多年歷史經驗.
來寫文章.
我用了一百多年歷史經驗.
來寫文章.
我用了一百多年歷史經驗.
來寫文章.
我用了一百多年歷史經驗.
來寫文章.
我用了一百多年歷史經驗.
來寫文章.
我用了一百多年歷史經驗.
來寫文章.
我用了一百多年歷史經驗.
來寫文章.
我用了一百多年歷史經驗.
來寫文章.
我用了一百多年歷史經驗.
來寫文章.
我用了一百多年歷史經驗.
來寫文章.
我用了一百多年歷史經驗.
來寫文章.
我用了一百多年歷史經驗.
來寫文章.
我用了一百多年歷史經驗.
來寫文章.
我用了一百多年歷史經驗.
來寫文章.
我用了一百多年歷史經驗.
來寫文章.
我用了一百多年歷史經驗.
來寫文章.
我用了一百多年歷史經驗.
來寫文章.
我用了一百多年歷史經驗.
來寫文章.
我用了一百多年歷史經驗.
來寫文章.

$^{801}$我用了一百多年歷史經驗.
來寫文章.
我用了一百多年歷史經驗.
來寫文章.
我用了一百多年歷史經驗.
來寫文章.
我用了一百多年歷史經驗.
來寫文章.
我用了一百多年歷史經驗.
來寫文章.
我用了一百多年歷史經驗.
來寫文章.
我用了一百多年歷史經驗.
來寫文章.
我用了一百多年歷史經驗.
來寫文章.
我用了一百多年歷史經驗.
來寫文章.
我用了一百多年歷史經驗.
來寫文章.
我用了一百多年歷史經驗.
來寫文章.
我用了一百多年歷史經驗.
來寫文章.
我用了一百多年歷史經驗.
來寫文章.
我用了一百多年歷史經驗.
來寫文章.
我用了一百多年歷史經驗.
來寫文章.
我用了一百多年歷史經驗.
來寫文章.
我用了一百多年歷史經驗.
來寫文章.
我用了一百多年歷史經驗.
來寫文章.
我用了一百多年歷史經驗.
來寫文章.
我用了一百多年歷史經驗.
來寫文章.

$^{841}$我用了一百多年歷史經驗.
來寫文章.
我用了一百多年歷史經驗.
來寫文章.
我用了一百多年歷史經驗.
來寫文章.
我用了一百多年歷史經驗.
來寫文章.
我用了一百多年歷史經驗.
來寫文章.
我用了一百多年歷史經驗.
來寫文章.
我用了一百多年歷史經驗.
來寫文章.
我用了一百多年歷史經驗.
來寫文章.
我用了一百多年歷史經驗.
來寫文章.
我用了一百多年歷史經驗.
來寫文章.
我用了一百多年歷史經驗.
來寫文章.
我用了一百多年歷史經驗.
來寫文章.
我用了一百多年歷史經驗.
來寫文章.
我用了一百多年歷史經驗.
來寫文章.
我用了一百多年歷史經驗.
來寫文章.
我用了一百多年歷史經驗.
來寫文章.
我用了一百多年歷史經驗.
來寫文章.
我用了一百多年歷史經驗.
來寫文章.
我用了一百多年歷史經驗.
來寫文章.
我用了一百多年歷史經驗.
來寫文章.

$^{881}$我用了一百多年歷史經驗.
來寫文章.
我用了一百多年歷史經驗.
來寫文章.
我用了一百多年歷史經驗.
來寫文章.
我用了一百多年歷史經驗.
來寫文章.
我用了一百多年歷史經驗.
來寫文章.
我用了一百多年歷史經驗.
來寫文章.
我用了一百多年歷史經驗.
來寫文章.
我用了一百多年歷史經驗.
來寫文章.
我用了一百多年歷史經驗.
來寫文章.
我用了一百多年歷史經驗.
來寫文章.
我用了一百多年歷史經驗.
來寫文章.
我用了一百多年歷史經驗.
來寫文章.
我用了一百多年歷史經驗.
來寫文章.
我用了一百多年歷史經驗.
來寫文章.
我用了一百多年歷史經驗.
來寫文章.
我用了一百多年歷史經驗.
來寫文章.
我用了一百多年歷史經驗.
來寫文章.
我用了一百多年歷史經驗.
來寫文章.
我用了一百多年歷史經驗.
來寫文章.
我用了一百多年歷史經驗.
來寫文章.

$^{921}$我用了一百多年歷史經驗.
來寫文章.
我用了一百多年歷史經驗.
來寫文章.
我用了一百多年歷史經驗.
來寫文章.
我用了一百多年歷史經驗.
來寫文章.
我用了一百多年歷史經驗.
來寫文章.
我用了一百多年歷史經驗.
來寫文章.
我用了一百多年歷史經驗.
來寫文章.
我用了一百多年歷史經驗.
來寫文章.
我用了一百多年歷史經驗.
來寫文章.
我用了一百多年歷史經驗.
來寫文章.
我用了一百多年歷史經驗.
來寫文章.
我用了一百多年歷史經驗.
來寫文章.
我用了一百多年歷史經驗.
來寫文章.
我用了一百多年歷史經驗.
來寫文章.
我用了一百多年歷史經驗.
來寫文章.
我用了一百多年歷史經驗.
來寫文章.
我用了一百多年歷史經驗.
來寫文章.
我用了一百多年歷史經驗.
來寫文章.
我用了一百多年歷史經驗.
來寫文章.
我用了一百多年歷史經驗.
來寫文章.

$^{961}$我用了一百多年歷史經驗.
來寫文章.
我用了一百多年歷史經驗.
來寫文章.
我用了一百多年歷史經驗.
來寫文章.
我用了一百多年歷史經驗.
來寫文章.
我用了一百多年歷史經驗.
來寫文章.
我用了一百多年歷史經驗.
來寫文章.
我用了一百多年歷史經驗.
來寫文章.
我用了一百多年歷史經驗.
來寫文章.
我用了一百多年歷史經驗.
來寫文章.
我用了一百多年歷史經驗.
來寫文章.
我用了一百多年歷史經驗.
來寫文章.
我用了一百多年歷史經驗.
來寫文章.
我用了一百多年歷史經驗.
來寫文章.
我用了一百多年歷史經驗.
來寫文章.
我用了一百多年歷史經驗.
來寫文章.
我用了一百多年歷史經驗.
來寫文章.
我用了一百多年歷史經驗.
來寫文章.
我用了一百多年歷史經驗.
來寫文章.
我用了一百多年歷史經驗.
來寫文章.
我用了一百多年歷史經驗.
來寫文章.

$^{1001}$我用了一百多年歷史經驗.
來寫文章.
我用了一百多年歷史經驗.
來寫文章.
我用了一百多年歷史經驗.
來寫文章.
我用了一百多年歷史經驗.
來寫文章.
我用了一百多年歷史經驗.
來寫文章.
我用了一百多年歷史經驗.
來寫文章.
我用了一百多年歷史經驗.
來寫文章.
我用了一百多年歷史經驗.
來寫文章.
我用了一百多年歷史經驗.
來寫文章.
我用了一百多年歷史經驗.
來寫文章.
我用了一百多年歷史經驗.
來寫文章.
我用了一百多年歷史經驗.
來寫文章.
我用了一百多年歷史經驗.
來寫文章.
我用了一百多年歷史經驗.
來寫文章.
我用了一百多年歷史經驗.
來寫文章.
我用了一百多年歷史經驗.
來寫文章.
我用了一百多年歷史經驗.
來寫文章.
我用了一百多年歷史經驗.
來寫文章.
我用了一百多年歷史經驗.
來寫文章.
我用了一百多年歷史經驗.
來寫文章.

$^{1041}$我用了一百多年歷史經驗.
來寫文章.
我用了一百多年歷史經驗.
來寫文章.
我用了一百多年歷史經驗.
來寫文章.
我用了一百多年歷史經驗.
來寫文章.
我用了一百多年歷史經驗.
來寫文章.
我用了一百多年歷史經驗.
來寫文章.
我用了一百多年歷史經驗.
來寫文章.
我用了一百多年歷史經驗.
來寫文章.
我用了一百多年歷史經驗.
來寫文章.
我用了一百多年歷史經驗.
來寫文章.
我用了一百多年歷史經驗.
來寫文章.
我用了一百多年歷史經驗.
來寫文章.
我用了一百多年歷史經驗.
來寫文章.
我用了一百多年歷史經驗.
來寫文章.
我用了一百多年歷史經驗.
來寫文章.
我用了一百多年歷史經驗.
來寫文章.
我用了一百多年歷史經驗.
來寫文章.
我用了一百多年歷史經驗.
來寫文章.
我用了一百多年歷史經驗.
來寫文章.
我用了一百多年歷史經驗.
來寫文章.

$^{1081}$我用了一百多年歷史經驗.
來寫文章.
我用了一百多年歷史經驗.
來寫文章.
我用了一百多年歷史經驗.
來寫文章.
我用了一百多年歷史經驗.
來寫文章.
我用了一百多年歷史經驗.
來寫文章.
我用了一百多年歷史經驗.
來寫文章.
我用了一百多年歷史經驗.
來寫文章.
我用了一百多年歷史經驗.
來寫文章.
我用了一百多年歷史經驗.
來寫文章.
我用了一百多年歷史經驗.
來寫文章.
我用了一百多年歷史經驗.
來寫文章.
我用了一百多年歷史經驗.
來寫文章.
我用了一百多年歷史經驗.
來寫文章.
我用了一百多年歷史經驗.
來寫文章.
我用了一百多年歷史經驗.
來寫文章.
我用了一百多年歷史經驗.
來寫文章.
我用了一百多年歷史經驗.
來寫文章.
我用了一百多年歷史經驗.
來寫文章.
我用了一百多年歷史經驗.
來寫文章.
我用了一百多年歷史經驗.
來寫文章.

$^{1121}$我用了一百多年歷史經驗.
來寫文章.
我用了一百多年歷史經驗.
來寫文章.
我用了一百多年歷史經驗.
來寫文章.
我用了一百多年歷史經驗.
來寫文章.
我用了一百多年歷史經驗.
來寫文章.
我用了一百多年歷史經驗.
來寫文章.
我用了一百多年歷史經驗.
來寫文章.
我用了一百多年歷史經驗.
來寫文章.
我用了一百多年歷史經驗.
來寫文章.
我用了一百多年歷史經驗.
來寫文章.
我用了一百多年歷史經驗.
來寫文章.
我用了一百多年歷史經驗.
來寫文章.
我用了一百多年歷史經驗.
來寫文章.
我用了一百多年歷史經驗.
來寫文章.
我用了一百多年歷史經驗.
來寫文章.
我用了一百多年歷史經驗.
來寫文章.
我用了一百多年歷史經驗.
來寫文章.
我用了一百多年歷史經驗.
來寫文章.
我用了一百多年歷史經驗.
來寫文章.
我用了一百多年歷史經驗.
來寫文章.

$^{1161}$我用了一百多年歷史經驗.
來寫文章.
我用了一百多年歷史經驗.
來寫文章.
我用了一百多年歷史經驗.
來寫文章.
我用了一百多年歷史經驗.
來寫文章.
我用了一百多年歷史經驗.
來寫文章.
我用了一百多年歷史經驗.
來寫文章.
我用了一百多年歷史經驗.
來寫文章.
我用了一百多年歷史經驗.
來寫文章.
我用了一百多年歷史經驗.
來寫文章.
我用了一百多年歷史經驗.
來寫文章.
我用了一百多年歷史經驗.
來寫文章.
我用了一百多年歷史經驗.
來寫文章.
我用了一百多年歷史經驗.
來寫文章.
我用了一百多年歷史經驗.
來寫文章.
我用了一百多年歷史經驗.
來寫文章.
我用了一百多年歷史經驗.
來寫文章.
我用了一百多年歷史經驗.
來寫文章.
我用了一百多年歷史經驗.
來寫文章.
我用了一百多年歷史經驗.
來寫文章.
我用了一百多年歷史經驗.
來寫文章.

$^{1201}$我用了一百多年歷史經驗.
來寫文章.
我用了一百多年歷史經驗.
來寫文章.
我用了一百多年歷史經驗.
來寫文章.
我用了一百多年歷史經驗.
來寫文章.
我用了一百多年歷史經驗.
來寫文章.
我用了一百多年歷史經驗.
來寫文章.
我用了一百多年歷史經驗.
來寫文章.
我用了一百多年歷史經驗.
來寫文章.
我用了一百多年歷史經驗.
來寫文章.
我用了一百多年歷史經驗.
來寫文章.
我用了一百多年歷史經驗.
來寫文章.
我用了一百多年歷史經驗.
來寫文章.
我用了一百多年歷史經驗.
來寫文章.
我用了一百多年歷史經驗.
來寫文章.
我用了一百多年歷史經驗.
來寫文章.
我用了一百多年歷史經驗.
來寫文章.
我用了一百多年歷史經驗.
來寫文章.
我用了一百多年歷史經驗.
來寫文章.
我用了一百多年歷史經驗.
來寫文章.
我用了一百多年歷史經驗.
來寫文章.

$^{1241}$我用了一百多年歷史經驗.
來寫文章.
我用了一百多年歷史經驗.
來寫文章.
我用了一百多年歷史經驗.
來寫文章.
我用了一百多年歷史經驗.
來寫文章.
我用了一百多年歷史經驗.
來寫文章.
我用了一百多年歷史經驗.
來寫文章.
我用了一百多年歷史經驗.
來寫文章.
我用了一百多年歷史經驗.
來寫文章.
我用了一百多年歷史經驗.
來寫文章.
我用了一百多年歷史經驗.
來寫文章.
我用了一百多年歷史經驗.
來寫文章.
我用了一百多年歷史經驗.
來寫文章.
我用了一百多年歷史經驗.
來寫文章.
我用了一百多年歷史經驗.
來寫文章.
我用了一百多年歷史經驗.
來寫文章.
我用了一百多年歷史經驗.
來寫文章.
我用了一百多年歷史經驗.
來寫文章.
我用了一百多年歷史經驗.
來寫文章.
我用了一百多年歷史經驗.
來寫文章.
我用了一百多年歷史經驗.
來寫文章.

$^{1281}$我用了一百多年歷史經驗.
來寫文章.
我用了一百多年歷史經驗.
來寫文章.
我用了一百多年歷史經驗.
來寫文章.
我用了一百多年歷史經驗.
來寫文章.
我用了一百多年歷史經驗.
來寫文章.
我用了一百多年歷史經驗.
來寫文章.
我用了一百多年歷史經驗.
來寫文章.
我用了一百多年歷史經驗.
來寫文章.
我用了一百多年歷史經驗.
來寫文章.
我用了一百多年歷史經驗.
來寫文章.
我用了一百多年歷史經驗.
來寫文章.
我用了一百多年歷史經驗.
來寫文章.
我用了一百多年歷史經驗.
來寫文章.
我用了一百多年歷史經驗.
來寫文章.
我用了一百多年歷史經驗.
來寫文章.
我用了一百多年歷史經驗.
來寫文章.
我用了一百多年歷史經驗.
來寫文章.
我用了一百多年歷史經驗.
來寫文章.
我用了一百多年歷史經驗.
來寫文章.
我用了一百多年歷史經驗.
來寫文章.

$^{1321}$我用了一百多年歷史經驗.
來寫文章.
我用了一百多年歷史經驗.
來寫文章.
我用了一百多年歷史經驗.
來寫文章.
我用了一百多年歷史經驗.
來寫文章.
我用了一百多年歷史經驗.
來寫文章.
我用了一百多年歷史經驗.
來寫文章.
我用了一百多年歷史經驗.
來寫文章.
我用了一百多年歷史經驗.
來寫文章.
我用了一百多年歷史經驗.
來寫文章.
我用了一百多年歷史經驗.
來寫文章.
我用了一百多年歷史經驗.
來寫文章.
我用了一百多年歷史經驗.
來寫文章.
我用了一百多年歷史經驗.
來寫文章.
我用了一百多年歷史經驗.
來寫文章.
我用了一百多年歷史經驗.
來寫文章.
我用了一百多年歷史經驗.
來寫文章.
我用了一百多年歷史經驗.
來寫文章.
我用了一百多年歷史經驗.
來寫文章.
我用了一百多年歷史經驗.
來寫文章.
我用了一百多年歷史經驗.
來寫文章.

$^{1361}$我用了一百多年歷史經驗.
來寫文章.
我用了一百多年歷史經驗.
來寫文章.
我用了一百多年歷史經驗.
來寫文章.
我用了一百多年歷史經驗.
來寫文章.
我用了一百多年歷史經驗.
來寫文章.
我用了一百多年歷史經驗.
來寫文章.
我用了一百多年歷史經驗.
來寫文章.
我用了一百多年歷史經驗.
來寫文章.
我用了一百多年歷史經驗.
來寫文章.
我用了一百多年歷史經驗.
來寫文章.
我用了一百多年歷史經驗.
來寫文章.
我用了一百多年歷史經驗.
來寫文章.
我用了一百多年歷史經驗.
來寫文章.
我用了一百多年歷史經驗.
來寫文章.
我用了一百多年歷史經驗.
來寫文章.
我用了一百多年歷史經驗.
來寫文章.
我用了一百多年歷史經驗.
來寫文章.
我用了一百多年歷史經驗.
來寫文章.
我用了一百多年歷史經驗.
來寫文章.
我用了一百多年歷史經驗.
來寫文章.

$^{1401}$我用了一百多年歷史經驗.
來寫文章.
我用了一百多年歷史經驗.
來寫文章.
我用了一百多年歷史經驗.
來寫文章.
我用了一百多年歷史經驗.
來寫文章.
我用了一百多年歷史經驗.
來寫文章.
我用了一百多年歷史經驗.
來寫文章.
我用了一百多年歷史經驗.
來寫文章.
我用了一百多年歷史經驗.
來寫文章.
我用了一百多年歷史經驗.
來寫文章.
我用了一百多年歷史經驗.
來寫文章.
好 謝謝大家 再見.
\newpage



\section{羅馬書 8:18-30}
\label{sec:032042}
\textbf{從廢墟到盼望 (華)}~~~~張路加牧師
\newline
\newline
連結: \href{http://www.church.com.hk/acms/content.asp?site=cdc\&op=show\&type=product\&code=032042\&layout=sermon}{\texttt{032042}} ~~~~ 語音日期: 2026-04-26
\newline
\newline
\hyperref[sec:032041]{\small{< < < PREV SERMON < < <}}
~
\hyperref[sec:index_chronic]{\small{[返順時目]}}
~
\hyperref[sec:032043]{\small{> > > NEXT SERMON > > >}}
\newline
\newline
羅馬書 8:18-30
\newline
\begin{longtable}{cl}
\hline
\hline
章節 & 經文 (和合本修訂版)\\
\hline
8:18 & \begin{tabularx}{0.7\textwidth}{X} 我認為，現在的苦楚，若比起將來要顯示給我們的榮耀，是不足介意的。 \end{tabularx} \\ \\ \relax
8:19 & \begin{tabularx}{0.7\textwidth}{X} 受造之物切望等候神的眾子顯出來。 \end{tabularx} \\ \\ \relax
8:20 & \begin{tabularx}{0.7\textwidth}{X} 因為受造之物屈服在虛空之下，不是自己願意，而是因那使它屈服的叫他如此。 \end{tabularx} \\ \\ \relax
8:21 & \begin{tabularx}{0.7\textwidth}{X} 但受造之物仍然指望從敗壞的轄制下得釋放，得享神兒女榮耀的自由。 \end{tabularx} \\ \\ \relax
8:22 & \begin{tabularx}{0.7\textwidth}{X} 我們知道，一切受造之物一同呻吟，一同忍受陣痛，直到如今。 \end{tabularx} \\ \\ \relax
8:23 & \begin{tabularx}{0.7\textwidth}{X} 不但如此，就是我們這有聖靈作初熟果子的，也是自己內心呻吟，等候得著兒子的名分，就是我們的身體得救贖。 \end{tabularx} \\ \\ \relax
8:24 & \begin{tabularx}{0.7\textwidth}{X} 我們得救是在於盼望；可是看得見的盼望就不是盼望。誰還去盼望他所看得見的呢？ \end{tabularx} \\ \\ \relax
8:25 & \begin{tabularx}{0.7\textwidth}{X} 但我們若盼望那看不見的，我們就耐心等候。 \end{tabularx} \\ \\ \relax
8:26 & \begin{tabularx}{0.7\textwidth}{X} 同樣，我們的軟弱有聖靈幫助。我們本不知道當怎樣禱告，但是聖靈親自用無可言喻的嘆息替我們祈求。 \end{tabularx} \\ \\ \relax
8:27 & \begin{tabularx}{0.7\textwidth}{X} 那鑒察人心的知道聖靈所體貼的，因為聖靈照著神的旨意替聖徒祈求。 \end{tabularx} \\ \\ \relax
8:28 & \begin{tabularx}{0.7\textwidth}{X} 我們知道，萬事都互相效力，叫愛神的人得益處，就是按他旨意被召的人。 \end{tabularx} \\ \\ \relax
8:29 & \begin{tabularx}{0.7\textwidth}{X} 因為他所預知的人，他也預定他們效法他兒子的榜樣，使他兒子在許多弟兄中作長子。 \end{tabularx} \\ \\ \relax
8:30 & \begin{tabularx}{0.7\textwidth}{X} 他所預定的人，他又召他們來；所召來的人，他又稱他們為義；所稱為義的人，他又叫他們得榮耀。 \end{tabularx} \\ \\
[1ex]
\hline
\hline
\end{longtable}
$^{1}$各位親愛的弟兄姊妹.
大家族內平安.
平安.
非常的感恩 也非常的開心.
在自己屬靈的娘家.
能夠跟親愛的弟兄姊妹們.
一起來交流.
交通上帝的話語.
今天我想特別用這樣一個題目.
來跟各位思想.
從廢墟到盼望.
就是在剛好40年前的今天.
1986年的4月26號.
人類歷史上發生了一場巨大的災難.
就是在當時的蘇聯.
也就是今天的烏克蘭的.
切爾諾貝利.
發生了核爆炸.
那一場的災難.
禍及整個的歐洲.
它所釋放的輻射.
相當於數百克原子彈.
所以造成了巨大的傷害.
當場死亡的有幾十人.
之後不斷的被發現.
許多人遭受輻射的傷害.
直到如今.
那一片土地.
是一片廢墟.
我想不但發生在那裡.
其實就看現在.
這個核威脅.
事實上好像也從來沒有.
遠離過我們人類.
不但有核的威脅.
各處發生的戰爭.
各處發生的災難.
後來發現從昨天下午開始.
我們Torrance變得全美有名.
不曉得大家有沒有注意到.

$^{41}$我收到好幾個關切電話.
電話說.
路加布斯聽說你的鄰居.
跑到華盛頓去殺人了.
我們知道原來就是.
那位叫做Cole Allen的.
他的工作場所離這裡.
大概只有開車不到十分鐘.
在Hawthorne跟Torrance的.
交界的地方.
那個地方叫做C2 Education.
這樣的一位教師.
我女兒的好多同班同學.
都接受他的課後輔導.
一位畢業於加州理工學院.
Caltech的高材生.
一位剛剛在去年十二月獲得了.
閱讀最佳教師的.
這樣一位年輕人.
攜帶多把武器.
在華盛頓犯下了驚天大案.
事實上我相信.
很多時候對我們來講.
今天世界有哪一個地方是安全的呢.
今天世界有哪一個地方.
不像是一場歷經浩劫的廢墟呢.
我想這個就是.
剛剛我們所念的經文.
一開始使徒保羅就提醒我們.
他說我想現今的苦楚.
The present suffering.
這對每一個世代.
活在世上的每一個人.
他都遭遇到他當時代的那一種的.
suffering 苦楚.
保羅提醒我們不單單是.
有信仰的人會遭遇各樣的逼迫.
他特別強調要提到.
一切受造之物.
不管你的背景.

$^{81}$不管你的文化.
不管你的出身.
不管你的意識形態.
不管你屬於哪個黨派.
不管你政治觀點如何.
所有的人.
男女老幼蓋莫例外.
一切受造之物.
我們都面對的叫做現今的苦楚.
保羅進一步展開講.
這個現今的苦楚到底包含什麼呢.
首先他就提到.
一切受造之物都浮在虛空之下.
我想這是人心靈裡面一個揮之不去的.
一個陰影.
這個虛空.
其實他原來的意識非常的豐富.
不但是提到人的那種.
感覺沒有意義.
感覺沒有希望.
而且那種的無能為力.
不但如此.
他也包含了保羅後面特別展開講.
就是都在敗壞的瑕疵當中.
這個敗壞事實上他用的是一個化學名詞.
這個化學在敗壞這個詞.
這個英文 decay 在化學上是指衰退.
任何原子我們知道隨著時間的推移.
都在進入那個衰退 衰退 衰退當中.
他說我們人類好像活在這個敗壞的瑕疵當中一樣.
難道不是嗎.
不但我們的外貌.
我們的外在好像漸漸隨著歲月而改變.
我們的行動力.
我們的思考力.
我們的記憶力.
不斷不斷的在衰弱當中.
但是更可怕的.
他提到整個人類被這種衰退所瑕疵.
無論我們認為我們科技多麼發達.

$^{121}$我們的文明多麼的昌盛.
但事實上鐵一般的事實.
就是我們都在這個衰退當中.
好像被他緊緊的包裹起來.
所以二十二節保羅提到了第一種的嘆息.
就是一切受造之物.
一同嘆息勞苦直到如今.
我想這個詞用的多麼精準啊.
當我們看看不但是周邊的環境.
看看世界上一片的亂象.
嘆息勞苦.
我們一面嘆息.
我們一面還不敢停啊.
所謂要討生活啊.
所謂要往前走啊.
而且保羅說這種情形自古至到如今.
事實上他用那個嘆息這個詞.
是一個非常深刻的一個詞.
描述的是一個產婦.
一個媽媽在要生小孩的時候.
最艱難的慘難當中的身影.
這是保羅要告訴我們的這種的嘆息.
離開了上帝.
我們有盼望嗎.
我們有前途嗎.
我們有未來嗎.
所以事實上更準確的說.
不但世界好像一片廢墟.
很像四十年前那場災難造成的廢墟.
今天你去看無論是伊朗一片廢墟.
無論是加沙一片廢墟.
其實世界的各處地方.
包括我們的人心.
都像坍塌的廢墟.
為什麼會這樣呢.
為什麼我們用嘆息.
成為我們人類的主旋律呢.
其實詩篇九詩篇.
老早告訴我們一個原因.
摩西說上帝啊.

$^{161}$我們因你的怒氣而消滅.
因你的憤怒而驚慌.
你將我們的罪孽擺在你的面前.
將我們的隱惡.
擺在你的面光之中.
因為上帝看人不像人看人.
不但我們有做出來的行為.
包括我們的心思意念.
這些所謂隱藏的惡.
都赤裸敞開在上帝的面前.
所以摩西告訴我們.
我們經過的日子.
都在上帝的震怒之下.
我們度盡了年歲.
好像一生嘆息.
這是何等真實.
人類概莫能外.
無論你活在亞洲.
活在美洲.
活在歐洲.
活在哪裡.
我們都在嘆息當中.
我們本來都在上帝的震怒當中.
所以摩西告訴我們說.
我們一生的年日是七十.
若是強壯可到八十.
我想現在過了三千五百年.
大概我們最多也就是.
在七十八十上面.
再加十年二十年.
人類沒有辦法有太多的改變.
而且內容也沒有改變.
摩西說就算七八十年的人生.
但是這個過程裡面能夠誇口的.
就是十二個字.
勞苦愁煩.
勞苦指到外在的.
愁煩指到內心的.
然後轉眼成空.
就是一眨眼就過去了.

$^{201}$第三個叫做如飛而去.
有時候你仔細思想.
時間怎麼這麼快速.
我想到我剛到琳瑯堂的時候.
一九九二年.
怎麼忽然間三十四年已經過去了呢.
多少的前輩們 多少的長輩們.
我常常想到以前在講台上.
常常看到這位坐在這邊.
這位坐在那邊.
現在好像我們剎那間.
我們的歲月如飛一般的過去.
有人特別畫了兩條曲線.
就是我們出生到二十.
這個叫做教育期.
我們上各種學校.
我們拼命讀書.
我們華人尤其看重讀書.
學而優則試.
所以我們一個學位一個學位念.
學試完了念碩士.
碩士完了博士.
博士以後變烈士.
我們就.
這叫活到老學到老.
活到老他一直學.
那麼當然要花費很大的費用.
誰都知道在美國教育也相當的昂貴.
其實不但在各處都是一樣.
中間大概你能夠所謂用命去換錢的.
也不過就四十來年.
你大學畢業以後.
可能開始工作了.
慢慢到你退休以前.
這一段時間所謂的奮鬥打拼.
那麼最後這一段呢.
進入晚年進入老年.
又需要花錢.
很多人就在奮鬥期的時候拼命賺錢.
為了使將來好像好一點的養老院.

$^{241}$好一點的將來這個醫療的照護.
因為現在大家都知道以前叫養兒防老.
現在叫養老防兒.
所以基本上沒辦法了.
都沒有安全感.
但是這兩條曲線大部分人是拉不平.
也許我們華人稍微好一點.
我們很努力很拼搏.
為的就是希望這兩條支出的曲線.
跟收入的曲線可以打平.
但是這樣的人生.
你覺得怎麼樣.
常常一個案例一個案例.
我看到去年九月份.
忽然中國有個很年輕的演員.
這個叫于蒙龍.
我想大概他爸媽看了太多瓊瑤的小說.
蒙蒙龍龍.
所以給孩子取名也叫蒙龍.
不過他在蒙龍當中就跳樓.
才三十多歲.
另外一個是這個所謂教育網紅.
頂流專家.
擁有非常龐大的這些的受眾群體在網路上.
他叫張雪峰.
他告訴人家就是人生怎麼能夠成功.
教你怎麼能夠去高考.
可以去進入大學.
不到四十一歲.
忽然間倒下走了.
其實常常這些的案例都提醒我們一件事實.
意義在哪裡.
難道不就像聖經告訴我們.
一切受造之物都浮在虛空之下.
很虛無很虛空.
突然大概是對虛空最好的一個描述.
我們人類如果意識不到.
不要神 轉背向神 離開神 褻瀆神 抵擋神.
這是我們必然的後果.
這是我們必然的結局.

$^{281}$我們活在神的震怒之下.
我們都記得所羅門王.
這是一個從世界上人的角度看起.
權力 財富 聰明智慧於一生的人.
各種的享受.
但是他講的最多的就是兩個詞.
虛空的虛空 一切都是捕風.
以前我常常把這兩個.
覺得大概就是一個排比句.
虛空跟捕風.
後來仔細查考發現有區別.
虛空是講一個人的內心空的.
空蕩蕩的 抓不住意義.
捕風呢.
因為風這個詞在希伯來文裡面.
它同時有兩個意思.
一個叫風 一個叫靈.
我們記得耶穌跟一個弟母談到的時候.
特別用風來比喻.
其實是要來提醒一個弟母.
我們人裡面那個靈.
所以當所羅門說.
一切都是捕風的時候代表什麼呢.
代表對一個有意義的活物.
他本來要抓住那個靈.
這是我們人一生的渴望 追尋.
所謂生之追尋.
就是我們這樣一個有創造力.
有思想 有思考 有意義.
追尋意義的人.
到底為什麼我們人要往前跑呢.
我們一定要問價值和意義.
這個是我們裡面的靈所催發的.
所催動的.
所以我們人的一生變成在抓靈.
但是我們知道從亞當夏娃.
他們違背上帝自己當神那一刻開始.
他們裡面的靈已經與上帝分離了.
所以我們人今天活在一個沒有靈的狀態.
我們只有一個身體的狀態.

$^{321}$包括我們的思想的狀態.
但是我們的靈.
那是直接上帝本來造人的時候.
不但造我們的身體.
而且把他的靈吹在我們裡面的.
那一部分不再在我們裡面了.
我們變成與上帝隔絕的人.
所以我們一生苦啊.
我們在追啊.
一方面感覺空.
一方面又想要抓到那個靈.
但是若不是藉著耶穌基督.
重新把這個靈放到我們裡面.
藉著祂的生命.
信靠祂的重新在祂裡面活過來.
我們仍然是一生在追靈.
在抓靈.
我們的狀態叫做活在上帝的正路當中.
今天我們的重點是廢墟當中的盼望.
人類的盼望在哪裡呢?.
哈巴古書告訴我們.
雖然你方言看一切都好像沒有希望.
無花果樹不發旺.
葡萄樹不接果.
橄欖樹不效力.
田地不出糧食.
圳中也絕了羊.
棚內也沒有牛.
這叫做廢墟.
通常戰爭以後.
地震以後.
災難以後.
這就是一個活生生的場景.
但是祂說我要因耶和華.
歡欣.
以救我的神.
喜樂.
歡欣和喜樂.
這是這個世界何等缺少的.
人們歡欣不起來.

$^{361}$喜樂不出來.
還能喜樂嗎?.
早上我們出門的時候看到.
路邊有一個.
這個.
應該是homeless晃來晃去.
然後再翻撿很多的.
這個垃圾桶.
面目猙獰.
很多時候你會很自然的感覺.
哇.
這是Torrance嗎?.
這是一個過去華人覺得還不錯.
還比較安全的地方嗎?.
這裡竟然還會有一個殺手跑出來嗎?.
我女兒的很多同學都非常震驚啊.
說這個人.
這個我們課後都去找他輔導的.
這個怎麼忽然間會變成這樣.
我想.
人類的歎息.
除非我們找到了神.
得到了靈.
要不然我們永遠在歎息當中.
第二種的歎息.
這裡告訴我們.
就是我們這些有聖靈出接國旨的.
這是指到已經重生得救的基督徒.
也會心中歎息.
哇.
我實在覺得勝經真的是太真實.
上帝對我們太了解.
神從來不講假大空的話.
基督徒不是超人.
基督徒不是所謂的刀槍不入的鋼鐵之軀.
以前中國大陸有本書說.
鋼鐵是怎樣煉成的.
我試學材料的.
我以為這個大概是講金屬材料.
後來發現不是.

$^{401}$講某一個人.
這個一個一個共產主義者.
怎麼樣.
這個最後好像就是完全像個超人一樣.
所以叫做鋼鐵.
怎麼把它煉成.
都希望把每個人都煉成這樣一個.
不太像人的人.
聖經講得很清楚.
我們就算有了聖靈出接國旨在我們裡面.
我們也會心中歎息.
但是請你注意這個歎息後面馬上.
保羅要告訴我們.
這個歎息內容是什麼呢.
就是 渴望 切忘.
等候我們身體完全的得贖.
這是何等真實.
當我們被病痛纏身的時候.
當我們遭遇患難的時候.
很多時候真的我們說.
主啊 願你快來 願你快來.
我們何等盼望.
這個要朽壞的身體.
能夠改變成一個榮耀的身體.
上帝應許我們一定會實現.
因為血肉之軀不能承受上帝的國度.
這個軀體好像很多器官地上缺貨.
但是天上沒有用.
不能承受將來上帝榮耀的國度.
所以這裡特別給我們一個應許.
我們在歎息中是在.
等候的身體完全的得贖.
就是直到我們的主耶穌基督第二次再來.
號筒吹響的時候.
一切都要改變.
我們的身體會剎那間改變形狀.
那麼這裡又提到另外第三種的歎息.
特別的寶貴.
就是我們的軟弱當中.
聖靈在幫助我們.

$^{441}$我們本不曉得當怎樣禱告.
只是聖靈親自用說不出來的歎息.
替我們禱告.
哇!.
你有沒有想過.
受造之物歎息.
神的兒女歎息.
聖靈也歎息.
但是這個歎息.
我去查了一下.
這個有一個神學上的解釋非常妙.
他要把這個詞特別要描畫出來.
他說好像一個媽媽.
抱著一個新生的嬰孩.
他自己的孩子.
抱在懷中愛他.
眼神裡面充滿了各樣的愛.
但是輕輕的歎息.
為什麼呢.
因為對這個孩子有期待.
期待這個生命.
不但是生出來.
媽媽期待用最好的來養育他.
媽媽更期待的是.
他能夠健康成長.
他能夠長大成人.
他能夠成為一個有用的一個.
棟樑.
一個人才.
帶著很大的一個期許.
就是要描述在這裡.
聖靈親自用說不出來的歎息.
這是已經地上你沒有辦法去形容.
聖靈住在我們裡面.
成為我們的珠寶.
聖靈安慰我們.
聖靈鼓勵我們.
聖靈帶領我們更深的進入上帝的裡面.
所以如果我們留意到剛剛的經文.
我們一共看到的是三種歎息.

$^{481}$但是很不一樣.
受造諸的歎息.
好像在分娩的當中那種痛苦.
那種慘難裡面的身影.
蒙恩得救之人的歎息.
是因為內心充滿了上帝啊.
今天你已經一個果子接在我的裡面.
我裡面靈已經被你拯救得救了.
我何等盼望這個果子能夠完全的成熟.
就是直到有一天.
我的身體完全的得贖.
那麼第三種就是聖靈的歎息.
他站在珠寶的立場.
非常體貼,熱為的替我們代求.
當我們覺得禱告不出來的時候.
當我們覺得軟弱無力的時候.
你千萬不要忘記.
上帝沒有忘記你,沒有忘記我.
他在我們裡面無聲的,無言的.
這樣的來為我們代禱,代求.
但是你留意這三種的歎息.
都是指望著一個未來.
所以十八節保羅一上來就說.
現今的苦楚如果跟將來的榮耀來一比.
就不足介意.
弟兄姊妹們為什麼我們能夠歡欣,能夠喜樂.
這個叫做盼望.
就是將來的榮耀.
所以在這段聖經裡面.
你看到接下來保羅要提醒我們的.
就是這個盼望.
這是我們的緣由,這是我們的底氣.
這是我們的把握.
那裡說神眾子的生命要完全的顯現出來.
那裡說我們必將脫離敗壞的瑕疵.
進入上帝兒女自由的榮耀.
那裡特別應許.
我們完全得到神兒子的名分.
就是我們的身體完全的得所.
清楚嗎?何等清楚?.

$^{521}$親愛的弟兄姊妹.
上帝的話從不會改變.
一點一話都不會改變.
所以回到當下,回到現在.
讓我們改換一下我們的perspective.
改換一下我們的觀念.
世界上的觀念是什麼呢?.
現在決定未來.
人生要贏在起跑線.
明天不知道.
我們活在當下.
不知不覺當中.
我們基督徒也會受這樣的影響.
我真是看到有些媽媽為了.
這個孩子還在肚子裡面.
為了讓孩子贏在起跑線.
然後就在肚子上放一個錄音機.
就教他這個.
有的是希望孩子將來成為音樂家.
就一遍遍放音樂.
就希望.
有的希望孩子在中國.
希望他們將來出國留學.
就一遍遍的放那個音語的版本.
我想這小孩真的生出來第一句話.
Good morning.
嚇他媽一跳.
怎麼會這樣.
就是贏在起跑線.
到一個程度.
現在這個叫做捲.
捲的不行.
就是每一個pin.
就是我只有今天.
明天我也不知道.
反正今天我就要拼.
我就要贏.
但是請你留意.
聖經的觀點剛好相反.
現今有苦楚.

$^{561}$無論你是誰.
川普總統終於講了一句實話.
大家看到他昨天晚上的記者會.
他說我早知這個總統這麼高的風險.
我乾脆不要選算了.
其實上到苗塘之上.
下到草根階層.
我們每一位都是如此.
有什麼是我們抓得住的呢.
有什麼是我們可以誇口的呢.
但是將來.
這是我們切實的盼望.
因為有將來.
所以現在不足介意.
因為有將來.
所以我們現在蠻有盼望.
因為有將來.
所以我們現在能夠毫不懼怕.
我們毫不畏懼.
前不久我去北加州.
參加了一位基督五家的長老.
叫做阮雄長老.
他的安息禮拜.
他從被查出來.
得了胰腺癌.
到被主接回天家.
大概只有不到一個月的時間.
但是當我聽到這個消息.
去看他的時候.
非常的淡定.
跟我談的都是歐洲的情況怎麼樣.
現在那邊宣教的狀況怎麼樣.
Albert 叫我是劉德華.
其實我們教會好幾位劉德華.
光能跟勇力都在德國.
我們留在德國的.
留學德國的.
或者是留在德國服侍的.
都叫自己劉德華.
就是留在德國的華人.

$^{601}$簡稱劉德華.
那麼.
他到了最後.
被主接去的前一天.
還在把他的主日學的課程錄音下來.
我想為什麼沒有懼怕.
很坦然.
因為知道.
盼望馬上要從眼睛看不見.
變成實際可以看見.
人間有盼望嗎.
我想到去年.
中國大陸發生了三塌.
就是一個叫許家印.
這個是曾經亞洲首富.
那麼以為錢可以他們許家隨便印.
叫許家印.
他做房地產叫恆大房產.
結果塌了.
到現在一大堆最近被判刑了.
那麼還有一個也是當過幾屆中國的首富.
叫鍾慶厚.
其實你看他名字就蠻已經.
蠻有伏筆了.
那麼他一貫的人設叫做.
一妻就一個太太.
一女就一個孩子一個女兒.
一雙布鞋.
你看我多麼清廉.
結果現在發現這個墳也塌了.
因為他人雖然走了.
結果查出了一大堆私生子女.
所以很多人說他的品牌叫哇哈哈.
那個礦泉水應該叫哇多多.
一連串的哇.
那麼還有一個有道高僧.
這個少林寺的主持.
中國佛教協會的副會長.
叫做釋永信.
後來說這個名字你千萬不能信.

$^{641}$怎麼可能信.
也被判刑了.
那個廟塌了.
所以山塌了.
房塌了.
墳也塌了.
廟也塌了.
我想其實世界上.
其實你常常能夠看到的.
這個叫做在這個 decay.
這個衰退的.
朽壞的這個 bondage 裡面.
無人能夠逃脫.
若不是耶穌基督來拯救我們.
我們誰能夠比別人更誇口呢.
我們誰能夠逃脫上帝的憎怒呢.
藉著耶穌基督的拯救.
所以使徒保羅接下來要特別.
跟每一位已經蒙恩得救的.
主的門徒們.
主的神的兒女們來說.
我們需要有正確的心態.
我們需要有一個真正的把握.
我們也需要開始行動.
心態是什麼.
剛剛我們念了聖經.
如果你注意.
中文有三個詞來描述.
這個盼望.
第一個叫做切忘.
切忘等候.
就是你是帶著一個.
熱切的一個 passion.
一個 expectation.
一個期待在盼望.
不是隨便.
這個盼望有沒有無所謂.
不是的.
我們真知道有這個盼望的時候.
我們裡面好像會被火點起來一樣.

$^{681}$充滿了 passion.
充滿了熱忱.
第二個叫做指望.
指望的中文也很妙.
跟手指有關係.
就是你手指指向的那個盼望.
換句話說這個盼望.
不是一個無邊無際的.
發散的沒有目的的.
是一個指向特定的目的的盼望.
所以中文叫做指望.
那麼第三個.
中文把它翻作盼望.
這個詞前面還有一個定貫詞.
是 the only.
獨一的.
就是我們的耶穌基督.
基督成了我們心中那獨特的唯一的盼望.
這樣的心態.
是我們神的兒女.
雖然上帝跟我們每一位都是同樣的人生.
基督徒不會比非基督徒.
好像我們擔子就輕了一點.
我們人生遭遇就少了一點.
我們的苦難就好像輕了一點.
都一樣.
上帝很公平.
但是心態可以不一樣.
同樣兩個擔子好像一個人挑扁擔一樣.
一個人一路埋怨腳步沉重.
另外一個人心情還不錯.
但是步履蹣跚.
只有第三個人.
一路歌聲漸步如飛.
為什麼會這樣.
我想這兩個擔子.
一個叫做快樂.
一個叫痛苦.
人生概莫能外.
無論是誰.

$^{721}$你總有喜樂的時候.
高興的時候.
你也一定會遇到難處.
挫折困難的時候.
但是你怎麼看待.
是因為我們的心態來決定.
第一個人把痛苦這個擔子放在前面.
快樂的放在後面.
所以一面走.
常常一低頭看到都是苦難難處.
充滿了哀怨.
第二個人把兩個擔子換一個個兒.
把快樂放到前面.
所以想到了上帝的恩典.
高興.
痛苦可以稍微忘記一點.
但是第三個人最聰明.
他用上了.
剛剛我們講那個真切的盼望的這把剪刀.
他在痛苦.
苦難的這個擔子上面.
扎一個洞.
剪一個洞.
所以苦難好像一路走來都漏掉了.
我常常想到我們當中的汪姑姑.
二十幾年的牢獄生涯.
但是你看到她很多苦難嗎.
你聽到她很多嘆息嗎.
好像沒有.
為什麼.
這樣的聖徒太多了.
因為他們有一把真實的盼望的剪刀.
所以.
接下來的經文是我們很熟悉的八章二十八節.
我們曉得神使萬事互相效力.
要叫愛神的人的益處.
不過我常常要問各位.
有時候也要問我自己.
我們真信得過這句話嗎.
當我們遭遇苦難.

$^{761}$遭遇挫折.
遭遇很大的那種壓力的時候.
我們感受.
我們曉得神使萬事互相效力.
叫愛神的人的益處.
有時候我們會說上帝啊.
益處在哪裡.
我看到的都是痛苦.
我看到的都是難處.
我看到的都是走不下去.
但是請你留意.
二十八節.
恰好就是夾在二十七節和二十九節之間.
所以剛剛我們念了聖經.
其實一共出現了四次.
知道.
兩次被翻譯做知道.
兩次被翻譯做曉得.
都是同一個意思.
同一個原文.
第一個知道就是.
我們曉得一切受到之物.
都浮在虛空之下.
這是世界和歷史的真相.
基督徒比世人更為清楚的看見.
這個世界非我家.
地和其上都要過去.
我們已經知道了這個世界的真相.
我們就不會全情的.
全身心的在裡面捲啊.
在裡面轉啊.
在裡面拼搏啊.
我們已經知道了這是世界歷史的真相.
但是接下來.
第二十七節剛剛講.
聖靈曉得神的意思.
所以我們需要留意的就是.
當我們說我們二十八節我們敢講.
神使萬事互相效力的時候.
好像一個三明治.

$^{801}$是因為有二十七節和二十九節.
二十七節告訴我們.
聖靈住在我們裡面.
這位聖靈完全知道上帝對我們整全的安排.
知道我們的軟弱.
知道我們什麼時候好像禱告都禱告不出來.
我們本來也不曉得怎麼禱告.
但是聖靈曉得.
那麼二十九節就更清楚的告訴我們.
原來我們這一生這一切.
神在萬古以先已經認識我們.
美國有一條高速公路上.
有一次有一個司機被攔下來.
他有點莫名其妙.
幹嘛把我攔下來.
那麼停到邊上.
警察看了他的駕照.
然後就說你知道我為什麼把你攔下.
他說我不知道.
我不覺得我犯了哪一條交通規則.
警察說我注意你很久了.
你這個車後面兩個後面的煞車燈都不亮.
這個很危險.
所以今天我只給你一個警告.
你趕快找一個路口下去.
找一個Garage.
找一個修理的.
把那個燈修好才能上路.
下次就不會給你警告單.
是給你直接罰單.
結果這個警察還沒講完.
那個人一聽就急得不得了.
馬上把門打開要跳下來.
警察說我已經跟你講不是給你罰單.
你不用那麼緊張.
趕緊你找一個路口開下去.
這個司機說警察先生.
你不緊張我很緊張.
因為你說那個後面煞車燈不亮.
代表後面那個拖車不見了.

$^{841}$我把兩個燈裝到那個拖車後面.
那個拖車上還有我的老婆小孩都在上面.
那現在完了他下車一看.
果然後面拖車不知道什麼時候.
那個中間那個他大概沒有接得太好.
那個拖車不見了.
還在往前開.
我想我們的人生.
唯有交給上帝最可靠.
你千萬不要以為我們很有經驗.
我們老司機了.
我們可以掌握我們的方向盤.
我們對前面的道路一清二楚.
上帝的話說.
婦人會忘記她生來的嬰孩.
雖然一般不會.
但是就會忘記.
神不會忘記.
神說看啊.
我將你銘刻在我掌上.
神把我們每一個人的名字刻在他的手掌上.
你的強顏.
這是指到我們的周邊包括你的住處.
包括你的一切.
藏在我眼前.
人百密總有一疏.
掛一總會漏萬.
但是上帝不會忘記.
所以最後這部分.
這個聖經特別提醒我們.
就是需要行動.
在現在和將來之間.
我們一方面心態上因為熱切的盼望.
催逼我們許多的行動.
我們教會到如今幾十年在上帝的恩典當中.
多少神的僕人使女從這裡走向四面八方.
多少前輩們他們敬前愛主.
雖然如今離開我們回到天家.
但是他們同樣有一天我們要與他們再相會在神家.
我們需要結實的忍耐的結實.

$^{881}$就是忍耐再忍耐.
輕看如今一切的苦楚.
帶著信心注目我們的盼望.
我們需要與神同工.
聖靈在我們裡面的禱告提醒.
安慰帶領.
我們與聖靈同工.
我們效法基督靠主恩典.
特別當托倫斯現在出現這麼多事情的時候.
我們同心禱告.
不但是在國家祈禱日的時候我們同心禱告.
我們為我們的社區禱告.
為我們的靈社禱告.
為我們的周邊禱告.
為美國禱告.
為世界禱告.
唯有禱告.
是我們最可以做的.
也一定看見上帝親自來成就這事.
我最後要講這個見證.
常常難忘.
這位老媽媽住在以前我上海家旁邊不遠的地方.
那裡有一條胡同叫做錢家堂.
蠻諷刺的就是.
因為錢家堂就是很多姓錢的人早期到了上海聚集在那裡.
所以那裡大部分人姓錢.
不過呢一窮二白一分錢都沒有.
那個年代.
那麼這個老奶奶呢.
我們叫她錢奶奶.
我父母後來因為信仰去勞動改造.
就囑咐我一件事.
路嘉你要替我們常常去看看那個錢奶奶.
這個錢奶奶無兒無女一生也沒有結婚.
她靠什麼過日子呢.
就靠在那個胡同口擺一個蔥薑攤維生.
那個年代啊.
七十年代八十年代那個蔥薑都是一分錢兩分錢.
我有時候經過那個胡同口看到錢奶奶一頭白髮被風吹起來.
冬天啊非常刺骨的人.

$^{921}$那麼因為父母的嘱咐.
所以我就這個過一段時間會到她家去看看她.
一到她家的更小.
她家就在那個胡同上面搭一個上海人叫做過街樓.
一個小小的大概只有幾個平方一個閣樓.
爬那個樓梯上去都顫顫巍巍的.
我想這個錢奶奶不容易啊.
要帶著這個蔥薑攤每天爬上爬下.
在那個弄堂口靠這個維生.
但是最奇妙的每次我去看她.
我說錢奶奶啊這個有什麼我可以幫忙的.
她都是笑容滿面.
她說路嘉你趕緊回去你們很忙的讀書.
奶奶沒事哎呀我如釋重服.
因為我覺得好像一個任務一樣.
因為爸媽交代.
那麼我說既然沒事我要走了.
結果呢剛要走她一把拉住我.
來禱告而且她禱告一定要跪下來.
一跪呢我發現她奶奶知道太多了.
為世界教會禱告中國教會禱告上海教會禱告.
我是跪得腿腳發麻.
所以呢我就聲音比較大.
我說阿們意思提醒該結束.
結果呢我越說阿們.
她越高興啊繼續禱告.
後來發現奶奶說阿們才可以.
你說阿們沒有用.
一個禱告的人.
好幾次我要走的時候.
他把一個小紙包塞在我手裡.
路嘉回去交給你外婆.
我拿回家給外婆.
外婆打開一看一塊錢.
有一次甚至兩塊錢.
外婆說路嘉怎麼可以.
怎麼可以錢奶奶這麼辛苦.
她的錢我們怎麼可以拿.
趕緊給她送回去.
我心裡想送回去倒容易啊.

$^{961}$又得跟她禱告一回.
你想一個老人家.
他要賣出多少蔥薑才能夠賺到一塊錢.
在那個年代.
直到有一天有個姐妹經過她家.
聖靈感動.
你要上去看一看錢奶奶.
這個姐妹那天有事情.
她就走過去.
又聖靈感動又回來.
來回幾次.
她不敢消滅聖靈的感動.
就上到樓上.
姐妹自己的見證.
她說當我推開那個門的時候.
發現裡面那種光我從來沒有看到過.
奶奶躺在床上.
看到姐妹進來就跟她招手.
姐妹過去扶在她的床邊.
奶奶聲音很低.
我剛剛一直禱告秋神.
派一個天使來.
神把你派來.
因為你看看我這張桌子要送給誰.
我這張椅子你幫我送給誰.
他交代的就是最後家裡面.
只有剩的這幾樣東西.
他想到就是誰有需要.
然後對姐妹說我們一起禱告吧.
姐妹說禱告禱告.
當我睜開眼發現.
奶奶已經回到主的家中.
被主接回天家.
她說那一張臉我從來沒有看到.
那麼樣的漂亮.
像天使一樣.
為什麼這些先輩們.
能夠這樣的喜樂.
因為這個盼望對他們是何等的真實.
今天你我廢墟上我們有希望有盼望.

$^{1001}$我們要再次回應保羅的禱告說.
我想現在的苦楚.
若比起將來要賢於我們的榮耀.
就不足介意.
但願最後神的僕人奧古斯丁.
在幾千年前.
在將近兩千年前告訴我們的話.
繼續成為我們的激勵.
他說將我們的過去.
放在神的憐憫之下.
將我們的現在.
放在神的慈愛當中.
將我們的未來.
放在神的信使之上.
\newpage



\section{約伯記 19:23-29}
\label{sec:032043}
\textbf{我的救贖主活著}~~~~黃馮潔霞師母
\newline
\newline
連結: \href{http://www.church.com.hk/acms/content.asp?site=cdc\&op=show\&type=product\&code=032043\&layout=sermon}{\texttt{032043}} ~~~~ 語音日期: 2026-04-26
\newline
\newline
\hyperref[sec:032042]{\small{< < < PREV SERMON < < <}}
~
\hyperref[sec:index_chronic]{\small{[返順時目]}}
~
\hyperref[sec:032044]{\small{> > > NEXT SERMON > > >}}
\newline
\newline
約伯記 19:23-29
\newline
\begin{longtable}{cl}
\hline
\hline
章節 & 經文 (和合本修訂版)\\
\hline
19:23 & \begin{tabularx}{0.7\textwidth}{X} 「惟願我的言語現在就寫上，都記錄在書上； \end{tabularx} \\ \\ \relax
19:24 & \begin{tabularx}{0.7\textwidth}{X} 用鐵筆和鉛，刻在磐石上，存到永遠。 \end{tabularx} \\ \\ \relax
19:25 & \begin{tabularx}{0.7\textwidth}{X} 我知道我的救贖主活著，末後他必站在塵土上。 \end{tabularx} \\ \\ \relax
19:26 & \begin{tabularx}{0.7\textwidth}{X} 我這皮肉滅絕之後，我必在肉體之外得見神。 \end{tabularx} \\ \\ \relax
19:27 & \begin{tabularx}{0.7\textwidth}{X} 我自己要見他，親眼要看他，並不像陌生人。我的心腸在我裡面耗盡了！ \end{tabularx} \\ \\ \relax
19:28 & \begin{tabularx}{0.7\textwidth}{X} 你們若說：『我們怎麼逼迫他呢？事情的根源是在於他』， \end{tabularx} \\ \\ \relax
19:29 & \begin{tabularx}{0.7\textwidth}{X} 你們就當懼怕刀劍，因為憤怒帶來刀劍的刑罰。這樣，你們就知道有審判。」 \end{tabularx} \\ \\
[1ex]
\hline
\hline
\end{longtable}
$^{1}$各位弟兄姊妹平安.
看完這麼精彩的slide show.
現在回來看神的話語.
領了幾首關於主活著的.
耶穌基督是復活的主.
所以今天的題目是.
我的救贖主活著.
是救贖主.
這位主活著.
祂不是只是一位審判的主.
或者說懲罰的主.
祂是一位救贖的主.
所以當我預備這篇講道的時候.
是因為每一年來到復活節.
大家都會思想一下主復活.
我們應該怎樣去更加愛耶穌.
為我們這位復活的主來活.
所以這節聖經浮現在我的腦裡.
大家都很熟悉的.
這節聖經是在《約伯記》第十九章二十五節.
一起來讀一下好嗎?.
「我知道我的救贖主活著.
活了必站在地上」.
這句話是誰說的呢?.
是約伯 因為是《約伯記》.
而且他和一班幾個朋友在討論.
人生怎會有苦難呢?.
苦到無法再苦了.
磨難磨到無法再苦了.
有沒有可能一天之內所有的子女都沒有了?.
一天之內所有的財產都沒有了.
僕人也沒有了.
什麼都沒有了 只剩下一個太太.
還有就是全身塞滿了瘡.
從頭到腳.
痛苦到她的地步 用啞片來刮自己.
你想想人生去到這樣的時候.
怎麼辦呢?.
很多朋友來安慰她.
那些朋友是出於好心的.

$^{41}$想開解她.
誰知越開解就越令到約伯很煩.
好像他們一直說.
你是不是犯了罪?.
你是不是做錯了事?.
壞人罰你.
罰你生病 罰你沒有子女.
還有沒有家產.
所以他一直覺得很難受.
身體難受 心靈更加難受.
所以這段聖經其實是約伯記的高峰.
約伯在這裡扭轉了.
因為其實在聖經裡面講到約伯記的時候.
他提到一件事就是.
認為善有善報 惡有惡報.
但是約伯是善人來的.
為什麼會受這麼多苦呢?.
所以約伯在這段聖經裡面.
23至29節.
約伯認為真相終必顯明.
朋友們你們要小心有報應.
他還是相信有因有果 有報應.
一會兒我和你們讀經文.
但是我先給你們一個背景.
約伯本來希望用人的方法.
來將他的冤屈永恆地保留下來.
讓人知道我很冤枉.
他要將他這樣的經歷寫在.
不是一本書 而是寫在石頭上.
一會兒你們讀聖經就知道了.
後來他轉念一想.
穿過腦袋一想.
他說不好了 不如我抓住神.
祂才是我救贖者.
所以我希望神終究為他伸冤.
所以這一段聖經我一會兒會讀.
就是約伯做了一個信心的跳躍.
他跳過了他的掙扎懷疑.
舊約時代早期好像約伯那樣.
約伯記裡面其他地方都顯示.

$^{81}$對於永生和復活的事.
不是知道很多.
但約伯就像一個預言那樣告訴我們.
認定他死後會見到神.
而且神也會為他伸冤.
當然在苦難裡面的約伯.
難免有些反反覆覆.
如果你遇過疾病 你遇過痛苦.
遇過艱難 苦難.
你有時會信心強一些.
有時信心軟弱.
但為何要選擇我呢? 是否我真的有問題呢?.
自我矛盾經常都會出現.
但這一段經文和約伯記之前的章.
人死如燈滅.
這個死亡觀是剛剛相反的.
天淵之別.
人是否一死了就了之呢?.
好像燈熄了似的.
不是 原來我們還要見神.
親自見神 不是派代表.
是自己去見神.
所以這幾個聖經給我們看到.
約伯受苦.
反而叫他促成了一個信心.
一個觀念的條約.
約伯因為確信上帝.
一定會為他伸冤.
雖然他以為自己死後.
才會得到公義的判斷.
所以他回頭警告他的朋友.
不要隨便逼迫我.
因為公義的神.
都會審判你們那些不公義的論斷.
或者是批判.
關於約伯的錯.
所以其實約伯在這裡.
好像一個先知似的.
他宣告了.
原來我們的救贖主是活著的.

$^{121}$而且最後 他會站立在地上.
來審判我們.
就好像「善有善報 惡有惡報」.
那個時候就開始了.
我們一起來讀這幾節聖經.
這張slide請弟兄們幫我讀.
19章23到25節.
請 惟願我的言語現在寫上.
都記錄在書上.
用鐵筆盡刻.
用鉛灌在盤石上.
直到傳到永遠.
我知道我的救贖主活著.
沒了必站立在地上.
這幾節就告訴我們.
約伯痛苦到一個地步.
他真的很想將他所說過的話.
寫下 記錄在書上.
但他希望用鐵筆在石頭上挑剔.
然後用鉛灌在盤石上.
以前有人倒紅油上去.
讓它不會沒了顏色 很難看.
有時你們 如果我們去哪裡呢.
我記得我有一次去萬里長城.
那裡一進去就有些東西寫在那裡.
或者有時你去行山.
看到山石那裡.
是有些東西刻在上面.
最多見當然是一兩個心.
串在一起.
經常都說海枯石爛.
永遠都不改變那些.
都是在石頭上試下可否刻.
或者在木頭上 樹那裡.
很多的 樹那裡.
石那裡用些紅油寫上去.
這裡約伯痛苦到一個地步.
他很想有些紀錄.
我不知道你在你人生最痛苦的時候做些甚麼.
是否用口來罵人.

$^{161}$用口來埋怨神.
還是你將你心裡的悶氣.
鬱氣 甚麼氣都好 怨氣.
寫下它.
這件事其實很多輔導員都教我們.
你將你心裡要說的話寫下.
就是寫你的日記.
我寫了很多本日記.
早上 晚上都寫.
特別在我的人生要轉另一個階段.
一個人的時候.
就有兩個小朋友.
那時做甚麼好呢.
寫 寫下你跟神講說話.
還有你裡面跟自己講說話.
另外 去個海邊走.
走 你會看到天地 海這麼澎湃.
這麼偉大 自己這麼渺小.
你怎會敢再去駁嘴.
去跟神鬥 沒得鬥的.
你只能讚美祂 稱頌祂 尊崇祂.
下一章就請姊妹們幫我們讀.
又不動 Lawrence 幫我一手 轉去第二章.
十九章二十六節.
好 姊妹們請你讀.
我這皮肉滅絕之後.
我必在肉身之外.
我自己要見祂.
親眼要看祂.
並不障礙人.
我的心牆在我裡面消滅了.
你看看這裡有一個隱藏者的訊息.
我是會再活的.
因為他說皮肉是會滅絕的.
但我在肉體以外是會見到神的.
而且我要自己去見祂.
是親眼要見到祂.
因為祂不是外人.
祂認識神 祂是敬畏神的人.
所以無論裡面外面都會消滅.

$^{201}$都會消化 都會沒了.
一個人走了之後.
撞在地上也好 火撞也好 土撞也好.
都會變成灰 變成塵.
塵歸於塵 土歸於土.
但若把他相信.
就算身體沒了 我仍然要去朝見神.
下一章.
你們讀吧 請.
「你們若說我們逼迫他要何等地寵呢?.
野事的根乃在乎他.
你們就當懼怕刀劍.
因為憤怒惹動刀劍的刑罰.
使你們知道有報應」.
原文是審判.
我們的救贖主是公義的.
他是審判萬人的主.
所以我們要懼怕他.
因為他是公義的神.
我們要敬畏他.
因為他是審判的主.
他要審判每一個人.
所以來到這裡的時候.
我們看到約伯.
他真的在苦難當中找到出路.
在苦難當中他的信心跳了一步.
他相信世界的人事物.
無論順的環境逆的環境.
都會過去的.
但是約伯的救贖主.
是永遠都活著的.
我想有一個故事大家聽過很多次.
但當我預備的時候.
仍然在我腦裡面.
這是幾十年跟著我的.
就是講到馬丁路德.
幾百年前 五百多年前.
一位宗教改革的先鋒.
他在威登堡教堂門口.
發表了著名的95條抗議文.

$^{241}$正式宣佈與羅馬教廷的決裂.
於是羅馬教廷在各處搜索他.
要將他置於死地.
那些日子裡面.
他身心都受到很大的煎熬.
不知道何時自己會沒命.
時時為這件事愁眉苦臉.
有一天.
馬丁路德回到家裡.
見到自己的太太.
穿了一身.
不是你們扮什麼以色列皇帝那些.
穿了一身那些什麼.
麻衣 披麻蒙灰.
真是黑衣服.
好像穿了喪禮服.
心裡面就很奇怪.
就問他太太.
你為什麼會穿這些喪服呢.
是不是有人離開世界.
你要紀念他.
他太太就回答她的丈夫.
是你的上帝死了.
這麼大膽.
他再說是你的上帝死了.
所以我要穿這些喪服.
馬丁路德就很害怕地問他.
你為什麼這樣說.
太太就說.
因為你的面容.
和你的生活態度.
都告訴我.
你所信靠的上帝已經死了.
所以我就穿這件喪服來悼念他.
馬丁路德那時才發覺.
宗教改革的工作.
不是靠自己的力量.
而是神在動工.
如果這是神的旨意.
神一定會成為他的後盾.

$^{281}$托住他.
如果是這樣.
還要怕什麼呢.
有什麼好怕呢.
於是他就用更加積極的態度來面對.
勝過各樣的困境.
對基督教的產生.
你知道它起了很大的作用.
弟兄姊妹.
你的上帝是活的嗎.
還是死的呢.
那就要由你的生活.
你的面容.
還有由你的對人的態度.
做事的方法.
來證明給人看.
我們的神是活的.
昨天下午我開了一個會.
在吃飯的時候.
就是我們教會在2008年開始的.
祈禱醫治事工.
已經來到第三十個cycle.
三十個cycle即是一個月兩次.
然後連續做十三個禮拜.
十三次不是十三個禮拜.
是十三個session.
然後又再開始一個cycle.
所以我想全南灣的教會.
是沒有一間教會這樣做的.
很多是怎樣呢.
他們直落做十三個禮拜.
然後休息.
看看有沒有人手.
因為做完十三個禮拜之後.
大家很累的.
因為你每個禮拜都會見到一班人來.
都是愁眉苦臉的.
都是有很多怨氣的.
因為他離婚.
帶著子女在掙扎生活.

$^{321}$因為他們沒有了父母.
沒有了配偶.
沒有了子女.
沒有了兄弟姊妹.
總之一定有些東西失去他才會來.
所以很多大的教會.
過千人教會都不能夠這樣不斷地開.
只能夠一段時間.
所以有人打電話.
譬如說我很想找這些.
他們就要找一些是有的.
即是這段時間是有的.
我曾經幫人問過一間教會.
因為他們很想去in person.
我們現在就在網上.
我們不是親身有人在教會等你.
我現在在網上.
以前做了很多年.
都是在實體的.
十幾年都在實體.
在教會.
但因為疫情.
我們轉了做網上.
一直到現在都未曾回來.
但問過那幾間大教會.
我們沒有.
因為我們沒有人手做.
但昨天.
當我們開會的時候.
我望下這一班.
幫了我們.
真是十幾二十年的.
的facilitators.
他們叫主持的小組.
我們分幾個小組.
他們在幫忙.
我望下他們.
這個沒有了兒子.
還是adult.
新來的.

$^{361}$即是四十歲了.
突然間才找到他已經走了.
這個就離婚很多年了.
嘩.
個個都是破碎的人.
但是可以堅持到十幾年.
還在我們教會的事工那裡.
我們教會只有我和Cindy Kwok.
兩個人做.
以前就是Helen.
Helen Chow是一開始的.
但是還有幾個.
其他堂會人來幫一下的.
都有.
但做得不久.
但這幾個來.
和我們一起開會的人.
他們是做了十幾年.
是沒有停過的.
所以我真是為了他們.
很感恩.
我們大家一起祈禱的時候.
都很感恩 很感動.
就因為他們將破碎人生.
轉過來.
不是一路碎下去.
他們說我要找回意義.
找回方向.
走回一條.
有道路真理生命的方向.
找到耶穌.
全部都是好的基督徒.
他們擺上自己.
所以他們越事奉.
去幫其他那些.
以前就是這樣的人.
他們自己就是這樣.
很感恩 很喜樂.
我又可以扶回一個人上來.
他可以來教會.

$^{401}$他可以繼續生活.
去做事.
那種滿足 那種喜樂.
真的不是錢可以買到的.
所以這個事奉.
是群策群力的.
是教會有異象開始.
然後有其他教會的弟兄姊妹來幫忙.
因為他得益.
他說我要給予別人一些益處.
弟兄姊妹.
你的人生有沒有破碎過呢?.
或者有.
又或者你從來沒有破碎過.
但都不要緊.
我們需要停下我們的腳步.
做一個引路的人.
我們不是將人帶到人的面前.
也不只是將人帶到教會.
因為我們教會不是完美 不是完全的.
我們真的要將人帶到神面前.
十字架.
就是叫他朝向十字架.
不是看著人.
因為看著人他很快就走了.
求神幫我們.
真的要明白.
原來我們的救贖主活著是這麼重要的.
因為我們會得到祂每天的供應.
扶持.
幫助.
最重要的是.
約伯說他是聽到 看到.
他將來要審問我們.
他會賞賜給我們.
也會懲罰我們.
這是約伯所認識的救贖主.
我們看看.
在私人筆下的救贖主是怎樣的呢?.
這是從《詩篇》裡面.

$^{441}$找到兩節聖經.
請姊妹來讀.
耶和華 我的盤石 我的救贖主.
願我口中的言語.
心裡的意念在你面前蒙月立.
他們也追念神是他們的盤石.
至高的神是他們的救贖主.
如果你找找神的屬性.
盤石是堅固可靠的.
他是真是大能的 他是至高的.
我給姊妹讀是因為我們很喜歡說話.
很喜歡說話.
有些人統計過.
一天姊妹會說多少句話.
突然間說多少句話.
一比較起來.
嘩 原來姊妹多這麼多倍.
突然間可以一天都不說話.
但這裡說.
當我們的神是我們的救贖主.
是我們的盤石的時候.
我們要依靠祂 仰望祂的時候.
我們口中的言語.
心裡的意念都要在神面前蒙月立.
口中心裡思念.
其實這些全部都是我們這個人.
每天在運作中口說的話.
腦想什麼 心又想什麼.
而且我們的神是至高的神.
我們真的不去敬畏神.
不去尊重別人也要敬畏神.
因為神是我們至高的神.
另外透過《箴言》的作者.
他也是在詩歌書裡面.
看看神是怎樣的.
弟兄請你讀.
因他們的救贖主大有能力.
他必向你為他們辯鬱.
弟兄很喜歡咬頸.
經常不是辯論就是爭論 討論也好.

$^{481}$但為了某些範圍的議題.
有時會咬到面紅耳熱.
甚至拍台走了.
但是你不要再這樣了.
因為其實我們的救贖主是大有能力的.
你和我都是他救贖的.
整個世界的人都是他救贖的.
有人說他是基督徒.
是甚麼大的公會中派出來的.
但做些事出來你很不忿.
你不要不忿.
我們有位救贖主是大能的救贖主.
他看到一切.
他會明白你心裡的那種矛盾那種掙扎.
而且他會為我們辯鬱.
誰對誰錯沒有人可以做一個判斷.
就好像有人說.
清官難審家庭事.
清官也難審教會事.
是真的.
有些事兩邊跟你說.
你也不知道哪一邊是真哪一邊是假.
而且你不知道怎樣回答這兩方面.
最好不要做審判官.
交給神來做.
這裡提到先知傳揚的救贖主.
剛才是約伯認識的救贖主.
又提到私人筆下的救贖主.
現在看看先知所傳揚的救贖主是怎樣的呢?.
以塞瓦書48章17節.
弟兄們請你們讀.
耶和華 你的救贖主.
以色列的聖者如此說.
我致耶和華 你的神教訓你.
使你得益處.
引導你所當行的路.
這裡又提到神的屬性是聖者.
祂是聖潔的神.
所以聖經你會看到神說.
我是聖潔的.

$^{521}$你們也要聖潔.
做神的子民百姓兒女.
也要聖潔.
也要追求聖潔.
這個教導是從舊教導到新教導.
所以新教導你看到保羅的書信.
告訴我們神的旨意是甚麼.
我們追求聖潔.
遠避那些偶像.
神是我們的救贖主.
是我們聖者.
聖潔的救贖主.
祂就會教訓我們.
使我們得益處.
引導我們行當行的路.
所以我們在營會的時候.
小組討論問過一個問題.
當你有危急的事發生.
危難來了.
你第一個會求誰幫助.
第一個會求誰幫助你.
神.
你們真的好.
你們的答案全對.
有時候我們口是這樣說的.
但你說我要求神.
突然拿了電話來.
打給誰呢.
上網了.
電郵給誰呢.
一定是用自己的方法.
等待神的時候.
我要找誰來談談.
忍不住了.
但這裡提到.
我們需要相信.
我們這位聖者.
我們的救贖主.
祂會藉著聖經教訓我們.
藉著神的話.

$^{561}$自己的話來提醒我們.
引我們行一條.
應該行的路.
畢業生主日快到了.
如果你們家裡有小畢業生.
中畢業生.
大畢業生.
大大畢業生.
請你把他們的大名提給我們.
畢業生主日提醒我們什麼.
不只是讓人知道.
我有四個孫五個孫畢業.
或者我家裡有多少人畢業.
不只是告訴別人.
你想把榮耀歸給神.
神讓我有兒子畢業.
有女兒畢業.
有孫兒畢業.
有媳婦畢業.
你說榮耀神是對的.
但你記得不只是停在那裡.
我們的畢業生主日.
主要真的提醒我們.
要繼續為我們的子女.
宣佈祈禱.
而且.
心言三章五道六字教我們什麼.
你要專心仰賴耶和華.
不可以靠自己的聰明.
在你所行的一切事上.
都要認定祂.
祂必指引你的迫路.
從父母開始.
我們不能夠用自己方法.
我們要專心仰賴耶和華.
不是依靠自己的聰明.
你多聰明都好.
我們沒辦法將一條正的路.
指示給我們下一代.
再下一代.

$^{601}$你有這麼長的生命.
你都沒有這樣的智慧.
我們這次去遊船河.
真的有人認不出路回來.
那張餐桌.
有一個侍應要扶著他.
拖著他.
撬著他.
帶他回來.
這個姐妹就認了他做乾兒子.
因為她真的細心到不得了.
將他帶回來.
有人說電話不見了.
有人說名牌不見了.
很多東西不見了.
連自己都不見了.
所以我們都要彼此照應.
我們身邊有很多年輕一點的.
他們願意停低腳步.
放緩他們的腳步.
等那些老人家.
還有用他的雙臂救力.
幫他們推行李.
真是活出彼此相愛的團契生活.
愛和關懷.
所以我們都不可以依靠自己.
要認定他.
就是要將神放在第一位.
英文聖經有個譯本說.
Put him first.
將神放在第一位.
誰是我們優先的次序呢?.
我記得我拍拖的時候.
我問黃牧師.
那時還不是黃牧師 是傳道仔.
我問他.
我和你拍拖.
究竟我在你第幾位呢?.
如果我沒有記錯.
我是排到第四,第五位.

$^{641}$我忘記了可能是他爸爸,他媽媽.
或者當然是神先.
爸爸,媽媽還有其他.
我心想我也是這麼後.
人生有很多優先次序.
那時可能還沒有求婚.
當然是放你後一點.
人生都是這樣的.
有時將你的子女放在第一位.
先是父母.
子女,孫子,媳婦.
那些都不可靠.
只能將神放在第一位.
所以我真的很感恩.
就是倚創下.
才知道他寫到這一句說話.
你要真的依靠神.
因為神才可以教訓你.
讓你得益處.
引導你所行的,當行的路.
才知道耶利米書.
耶利米在耶利米書中告訴我們.
好,姐妹們你讀一下.
「他們的救贖主大有能力.
萬軍之耶和華是他的名.
他必申清他們的冤.
好使全地得平安.
並搞擾巴比倫的居民」.
這裡提到有兩個字.
我真的很搖眼.
申清.
你知道申冤有時申不清.
有時去找法官去幫你申冤.
特別是最慘就是.
罵離婚的時候.
我真的試過勸一對.
主內弟兄姊妹.
不是這些教會.
是以前的教會.
他真的連一張床都計算.

$^{681}$他說我要和他離婚.
這張床是我的.
我心想.
用不著計算得這麼清吧.
是真的.
我們有時很想清那些恩恩怨怨.
但清不清.
因為你有你講.
他有他講.
但這裡只有神.
這位大有能力的神.
可以幫我們申清我們所有的冤.
然後使我們得到平安.
好,去快點.
申約裡的救贖主.
上主眷顧私行救贖.
這是路加福音一章六十八節.
主以色列的神是應當稱頌的.
因為祂眷顧祂的百姓.
為他們私行救贖.
這句話大家知不知道是誰說的?.
在路加福音一章.
一章六十八節未到耶穌出世.
一定不是耶穌說的.
是耶穌先鋒的背景.
誰說的?.
祂未出世.
可能祂出世了.
祂的父親稱頌神.
祂之前未出世就不相信會有兒子.
已經一把年紀了.
後來神給祂啞,啞到祂的兒子出世了.
然後祂開口了.
你知道嗎?.
祂就說了這一段話.
我讀給你聽,因為很重要.
「祂父親撒迦利亞被聖靈充滿.
就預言說.
主以色列的神是應當稱頌的.
因為祂眷顧祂的百姓.

$^{721}$為他們私行救贖.
在祂僕人大衛家中.
為我們興起了拯救的國.
正如主席著從創世以來.
聖先知的口所說的話.
拯救我們脫離仇敵.
和一切恨我們之人的仇.
向我們列祖施憐憫.
紀念祂的聖約.
我們的神是守約,思持愛約.
特別祂幫我們和我們立的聖約.
是不能改的.
我們人會改,會離棄,毀約.
但是神是永遠都持守著祂的聖約.
什麼聖約?.
就是祂對我們祖宗亞伯拉罕所起的誓.
叫我們既從仇敵手中被救出來.
就可以終身在祂面前.
坦然無懼地用聖潔公義事奉祂.
祂再說:.
「孩子,(祂為祂的兒子祝福).
你要稱為至高者的先知.
因為你要走在主的前面.
預備祂的道路.
所以叫做開路先鋒.
叫祂的百姓因罪得赦.
就知道救恩.
因我們神憐憫的心腸.
叫清神的日光從高天臨到我們.
要照亮坐在黑暗中死忍裡的人.
把我們的腳引到平安的路上」.
司祭約翰說:「那孩子漸漸長大.
心靈強健.
住在曠野.
直到祂顯明在以色列人面前的日子」.
司祭約翰一生人都是做一把聲.
因為他在曠野告訴別人.
我們的主就來了.
是救贖主就來了.
是神的兒子就來了.

$^{761}$但是人生一過了他就沒有了.
所以他有一句名言.
「主必興旺我必衰微」.
我們不能一路擴大自己.
我們只能一路收細自己.
甚至消失了自己.
但是被人看到的是我們的救贖主.
我們的主耶穌基督.
我們的父神.
我們的聖靈.
今天我編導是cut了三分一的.
因為我預備好我才進去的那個營會.
所以後來接到知道會有slide.
我就把它cut了三分一.
不過不要緊.
剛才你唱的詩歌.
等下你再唱那首詩歌.
都是有一些我要講的東西.
其實我的救贖主活著.
和我有什麼關係呢?.
最後我要用一首詩歌的歌詞.
和大家來勉勵.
既然他是活著.
我們就「生命充滿了希望」.
這首歌大家猜到是什麼歌?.
「生命充滿了希望」.
哪一首?.
「印他活著」.
剛才不是唱過了.
好了.
因為我救贖主活著.
和他有什麼關係?.
「有他同在同行」.
這首是什麼歌?.
剛才不是唱過了.
他和我說話.
他陪我走路.
就是主活著.
我沒時間和大家講主活著的寫作背景.
真的很感動.

$^{801}$你上網查一查.
你會看到主活著.
這個作詩寫詞的人.
他受過很多難處.
他講的那一刻.
你有沒有試過有人站起來.
說你講這個不對.
耶穌是死了的.
他埋在地上變成了塵土.
但真理是會繼續進行的.
在教會他都信真理.
但他就是不信耶穌復活的.
第二.
他也受了很多逼迫.
被人都一樣說不信耶穌.
他為了要澄清這件事.
他很有感動.
所以就寫了這首很偉大的歌.
就是主活著.
你回去看看他的背景.
他和他太太都是很賣力的.
而且他一直唱的時候.
每個人都是學習將音樂和敬拜融入裡面.
最後.
我的救贖主活著與我什麼關係?.
天天為主而活.
我們的主每天都陪著你走.
陪著你聊天.
他也使你有希望.
人生多苦都有希望.
那我要怎麼做?.
待會你會唱這首歌.
我願獻身心為基督而活.
在一切事上討我主喜悅.
寬心信服主.
憑主指導引.
願背起十架一路隨主行.
我 我主耶穌.
我救主.
我心奉獻給你.

$^{841}$因你在十字架上.
為贖我的罪釘死.
我要你作我救主.
來住我心裡.
我一生一世.
到永遠.
我 主.
我完全歸你.
完全歸你.
為主而活.
正在低頭我們祈禱.
天父我們真的很感恩.
你派耶穌基督來.
第一次是要拯救世人.
第二次是要審判世人.
但我們就是在這一段.
等候耶穌基督回來.
要受審判之前.
我們信了耶穌.
我們得到救贖.
可以成為神的兒女.
今天我們的救贖主是活著的.
是與我們同行.
伴我們走.
甚至與我們聊天.
安慰我們 扶持我們.
所以我們願意將我們的身心靈.
都獻給我們的救贖主.
願你帶領我們.
幫我們討你的喜悅.
幫我們信靠信服你.
奉耶穌基督得勝的名頭.
阿們.
\newpage



\section{出埃及記 20:2}
\label{sec:032044}
\textbf{聖經人聖經事:我曾領你出來}~~~~黃克勤牧師博士
\newline
\newline
連結: \href{http://www.church.com.hk/acms/content.asp?site=cdc\&op=show\&type=product\&code=032044\&layout=sermon}{\texttt{032044}} ~~~~ 語音日期: 2026-04-21
\newline
\newline
\hyperref[sec:032043]{\small{< < < PREV SERMON < < <}}
~
\hyperref[sec:index_chronic]{\small{[返順時目]}}
~
\hyperref[sec:032045]{\small{> > > NEXT SERMON > > >}}
\newline
\newline
出埃及記 20:2
\newline
\begin{longtable}{cl}
\hline
\hline
章節 & 經文 (和合本修訂版)\\
\hline
20:2 & \begin{tabularx}{0.7\textwidth}{X} 「我是耶和華－你的神，曾將你從埃及地為奴之家領出來。 \end{tabularx} \\ \\
[1ex]
\hline
\hline
\end{longtable}
$^{1}$各位電影姐妹.
今天讀的聖經很緊張.
因為是十誡.
聖經人聖經事由生命樹宣教網絡為你預備.
今次的題目是.
「我曾領你出來」.
這是十誡的開始.
大家要留意一下.
特別跟大家講到十誡.
嘩.
很厲害,很嚴厲.
其實原來不是.
《出埃及記》第二十章第一節.
「神吩咐這一切的話說」.
神有個吩咐.
祂的說話是這樣說的.
「我是耶和華你的神」.
「曾將你從埃及地遺奴之家領出來」.
「除了我以外」.
「你不可有別的神」.
這句大家都已經熟悉了.
由於有些綠色字介紹.
大家要看的是黑色字.
紅色字.
和藍色字.
就是和合本本身有的中文字.
綠色字是我的註解.
藍色是什麼呢?.
藍色就是原文.
應該是沒有這個字的.
為什麼要寫下去呢?.
就是幫我們讀者容易明白一點.
紅色字就是想吸引眼球.
所以有紅有藍有綠.
讓你看得很花.
不過你慢慢看.
又可以翻看.
這樣也挺好的.
如果是螢光了的.
就表示是我的大題目.

$^{41}$通常我都會歡喜找三個大題目.
讓我們在一個訊息裡面記得.
「除了我以外,你不可有別的神」.
再讀完它.
「不可為自己雕刻偶像,也不可做甚麼形像」.
「彷彿上天,下地和地底下水中的白物」.
「不可跪拜那些像,也不可事奉他」.
「因為我耶和華裡的神,是忌邪的神」.
「恨我的,我必追討他的罪」.
「至父及子,直到三,四代」.
「愛我,守我誡命的」.
「我必向他們發慈愛,直到千代」.
這個就是今天要跟大家讀的經文了.
我們先回到最開頭開始.
「我曾領你出來」.
這裡神吩咐這一切的話.
神吩咐很多東西,不只是這幾節.
這一切就是講以下很重要的話.
「我是耶和華裡的神」.
是我們一般叫「十誡的開始」.
為什麼叫「十誡」呢?.
是不是聖經所說的「十誡」呢?.
不是,不過我們知道這些叫「誡命」.
「誡命」其實就是上帝給的說話要守著.
我們通常叫「誡命」.
人給的,我們通常叫「命」.
不過是「命令」.
神會不會給「命令」呢?.
都有,聖經有很多「命令」.
「命令」,「律例」,「典章」,「誡命」等等都有.
其實「誡命」是神給的.
由於有十條,因此一般來講教會叫「十誡」.
我們一聽到「十誡」就很害怕.
不可以做,不可以犯.
犯「十誡」很嚴重.
其實用這個「誡」字.
要先講,用這個「誡」字.
我個人覺得不算是很好的字.
不算是很好的字,為什麼呢?.
是因為它很惡.

$^{81}$很有壓力.
很嚴.
很有趣的,我去過一間教會.
一進去.
那間教會也很大,一座建築物.
一進去就不會馬上看到紅地毯的走道.
兩邊椅子,是看不到的.
因為有一個很大的屏風.
好像你,大家看電視.
你望進來,就是這個教會.
有一道牆遮住.
來到要分兩邊進去.
才能進到課堂,課堂在後面.
在這個位置寫著「十誡」.
12345678.
不是,是打直的.
「十誡」12345678910.
一進到課堂,第一件事就是「十誡」.
嘩,很嚴.
是很嚴的.
其實信仰是嚴厲的.
不過,我覺得上帝跟我們說的時候.
說「十誡」不是這麼嚴.
不是用這麼惡的態度.
反而是一個很柔和.
很有愛意的態度.
所以,我結論就是不要叫它「十誡」.
有些人喜歡叫它「十言」.
「十句說話」.
或者叫「十言」,講十樣東西.
第一誡何時開始?.
「除了我以外,你不可有別的實」.
第一誡,其實不是.
第一誡的開始在第二節.
留意,一開始就說「講明身份」.
「我是耶和華,你的實」.
「跟你講的東西」.
所以是「耶和華的話」.
「十話」或者「十言」.
第一個要留意的第一點.

$^{121}$不是第一誡,第一點是甚麼呢?.
「曾將你從埃及地為奴之家領出來」.
「從埃及地為奴之家領出來」.
「曾將你從埃及地為奴之家領出來」.
「我就是那個在埃及」.
「在為奴之家救你出來的人」.
「我就是那個在埃及」.
「在為奴之家救你出來的人」.
「自然在我以外」.
「以外」是指「之上」.
「所以在我面前以上」.
「你不可有別的神」.
這個「你」是沒有的.
「你不可有別的神」.
這個「你」是沒有的.
「你不可有別的神」.
「別的神」.
我,甚麼名?是你先生,是你丈夫.
你的神,這是一個關係.
我爺爺說,是你的神,是忌邪的.
quena,這個希伯來文的字.
英文可以翻譯為妒忌,嫉妒,jealous.
忌邪由於用在跟其他偶像假神的對比.
所以用了忌邪的中文.
其實我忘了寫中文,就是嫉妒.
我是一個嫉妒的神.
怎樣嫉妒?堂堂上帝這麼偉大.
嫉妒,這麼小心眼,這麼狹窄.
條條大路通羅馬,這樣是不行的.
總之虔誠的就是神,不行的.
為甚麼不行?條條大路通羅馬,不要這麼狹窄.
當然不可以.
如果你生了個兒子,生了個女兒.
有一天你見到他叫另一個做爸爸.
叫另一個做媽媽,那你怎樣?.
那不要緊,沒問題,他總之叫,當然不行.
因為你的爸爸或者媽媽只有一個.
就是我,就是你的父母.
大家明白嗎?所以不可以.
老公有多少個?不要緊,男的叫做老公.

$^{161}$有鬍子的叫做老爸,這樣不行.
只有一個,叫另一個,我會很歡喜.
當然不會,因為你叫另一個,我會生氣.
我會嫉妒,我會妒忌.
為甚麼你會叫另一個?這樣我會很生氣.
是這樣的意思,而我認為.
如果你會叫另一個做你的神.
你就是恨惡我,你就是不喜歡我.
我會追討他的罪.
他不喜歡你,你追討別人的罪?.
是的,因為我是將你贖身出來的那個.
所以你是要跟我,我是愛你跟我.
不是我贖你身出來,你都聽我話,不是我鞭你.
不是這樣的奴隸,而是愛的跟從.
我贖他出來,還是愛他.
其實神是想跟以色列有這個關係.
就好像一個情人,我將他救出來.
所以你會見到,其實這個關係.
聖經有沒有表達的呢?有,舊約裡面都很明顯.
舉個例子,賀西亞書,神是叫賀西亞仙子.
你去娶一個淫婦為妻.
這個仙子真的暈倒了,不過這樣叫,我就做吧.
主要就是,他是用以色列人.
好像一個淫婦那樣,又會叫第二個做老公.
除了賀西亞仙子之外,搞到他真的很慘.
個個都看著他這樣.
他象徵神很愛屬於他的人跟他結婚.
誰知那個人就走了,跟別人說.
忌邪,Jealous,其實神是想有這個關係.
所以你不要製造,不要跪拜.
製造跪拜是否不好.
只有愛神,如果你覺得其他是神.
這樣你就令到我真的很生氣,忌邪.
我是會心胸疾的,是嗎?是的.
其實這個我就算為恨我,即是不歡喜我.
既然我救了你出來,所以我會追討你的罪.
你的王是罪惡不公,不對.
這樣不對,救了你出來,你又不跟他走路.
我會追討他的罪,致父及子直到三死.
沒有個代字的,致父及子直到三死.

$^{201}$三死是甚麼呢?代.
不是三死日嗎?應該不是.
直到三死,沒有說是代.
不過推測應該是代,所以中文就譯為代.
致父及子直到三死.
有父子,於是想有孫,就是三代.
有父子孫失,就是四代.
由於是父子,所以就認為是代.
而不是日,年,月,星期,大家可以理解.
好了,再看下去.
第三點到第六節就出現第三點.
愛我的阿孝.
愛我的呢?.
甚麼人是愛我的?.
守我誡命的,誡命呢?.
那就是誡命,神給的就是誡命.
人給的是命令.
守我的命令,守我的誡命.
我就會向他發慈愛.
發慈愛又是這個字,阿孝是製造.
我會,甚麼叫製造呢?.
沒有的,我做出來.
來到這裡又沒有桌子,又沒有椅子,又沒有床.
那怎麼住?可以,等等.
我製造出來,由沒有變成有.
不過就在這個空間裡面.
已經拿些東西給他,做做做給他.
就不是那個創造.
就是起初神創造,是製造.
我就會製造慈愛,念文.
CH讀K,C讀S.
這個讀K,這個讀S.
所以Keset,我就會向他製造慈愛.
其實Keset我就會向他製造念文.
直到千,Eleph.
又沒有袋子.
直到千,沒理由千日吧?.
沒理由千小時吧?.
於是上面,這裡是父子.
到三四,於是父子.

$^{241}$算,三是袋子.
父子宣室,在這裡父子,根據下來.
第四代,父子宣室四代.
於是想這個Eleph.
千,都是袋子.
直到千袋.
其實聖經只有直到千,直到三四.
不過應該是袋子.
神的慈愛,神的Keset大一點.
還是神的Jealous大一點呢?.
大一點呢?.
神的Avon,Avon這麼大.
神的恨,Sane,Sane.
Low,大一點還是神的愛大一點呢?.
很明顯是神的愛大一點.
因為祂是會直到千袋的祝福.
如果不守祂的誡命,那也很麻煩.
因為至父及子直到三四袋.
有時我們看到神很愛我們.
好了,重複結語.
十誡,不如我們立刻將我們的腦袋串一串.
不要叫它做誡.
是愛耶和華神,我們的神.
愛的十句.
愛的十言.
愛的十說.
愛的十句,這樣的意思是更加好一點.
第一句從哪裡開始呢?.
我耶和華你的神,不是祂的神.
不是天上的神,是我的神.
是你的神,耶和華你的神.
你要記得就是我將你們在埃及那裡取出來的.
或者贖出來的.
幫你贖身的那個.
所以你不再需要有其他人的了.
有我就可以了.
不要再搞那些又說是又說他又像我那些.
你搞一些我會很生氣的.
因為我只有一個.
你叫第二個做老公,你叫第二個做爸爸.

$^{281}$是不可以的.
我會心胸窄的.
我會生氣的.
不過當你這樣跟著我.
我是會愛直到千代的.
這個就是我今天和大家做這個十誡的分享.
第一句,第一誡.
不要叫誡,叫做第一言.
第一句,第一個愛的言語.
祝福,我們一起祈禱.
多謝我們在天上的神.
雖然我們肉眼沒有能夠親眼看到.
不過我們由這麼多世界上所做之物的奇妙.
更加由以色列人的身上.
你見到你將他們這樣帶出埃及.
你好像一個丈夫娶了太太一樣.
帶他們走完整個人生路程.
去到回到家鄉.
主你就是這樣愛他.
而你就跟我們表明.
有你就夠了.
不需要再有第二個了.
就幫助我們專心愛你.
求你的慈愛更多臨到我們.
讓我們堅定專心跟隨你,愛你.
祈禱奉耶穌基督的名字.
阿們.
\newpage



\section{申命記 6:1-25}
\label{sec:032045}
\textbf{上帝也講KPI?(2) 點解要做?}~~~~符氣松牧師
\newline
\newline
連結: \href{http://www.church.com.hk/acms/content.asp?site=cdc\&op=show\&type=product\&code=032045\&layout=sermon}{\texttt{032045}} ~~~~ 語音日期: 2026-04-25
\newline
\newline
\hyperref[sec:032044]{\small{< < < PREV SERMON < < <}}
~
\hyperref[sec:index_chronic]{\small{[返順時目]}}
~
\hyperref[sec:032046]{\small{> > > NEXT SERMON > > >}}
\newline
\newline
申命記 6:1-25
\newline
\begin{longtable}{cl}
\hline
\hline
章節 & 經文 (和合本修訂版)\\
\hline
6:1 & \begin{tabularx}{0.7\textwidth}{X} 「這是耶和華－你們的神所吩咐要教導你們的誡命、律例、典章，叫你們在所要過去得為業的地上遵行， \end{tabularx} \\ \\ \relax
6:2 & \begin{tabularx}{0.7\textwidth}{X} 好叫你和你的子孫在一生的日子都敬畏耶和華－你的神，謹守他的一切律例、誡命，就是我所吩咐你的，使你的日子得以長久。 \end{tabularx} \\ \\ \relax
6:3 & \begin{tabularx}{0.7\textwidth}{X} 以色列啊，你要聽，要謹守遵行，使你可以在那流奶與蜜之地得福，人數極其增多，正如耶和華－你列祖的神所應許你的。 \end{tabularx} \\ \\ \relax
6:4 & \begin{tabularx}{0.7\textwidth}{X} 「以色列啊，你要聽！耶和華－我們的神是獨一的主。 \end{tabularx} \\ \\ \relax
6:5 & \begin{tabularx}{0.7\textwidth}{X} 你要盡心、盡性、盡力愛耶和華－你的神。 \end{tabularx} \\ \\ \relax
6:6 & \begin{tabularx}{0.7\textwidth}{X} 我今日吩咐你的這些話都要記在心上， \end{tabularx} \\ \\ \relax
6:7 & \begin{tabularx}{0.7\textwidth}{X} 也要殷勤教導你的兒女。無論你坐在家裡，走在路上，躺下，起來，都要吟誦。 \end{tabularx} \\ \\ \relax
6:8 & \begin{tabularx}{0.7\textwidth}{X} 要繫在手上作記號，戴在額上作經匣； \end{tabularx} \\ \\ \relax
6:9 & \begin{tabularx}{0.7\textwidth}{X} 又要寫在你房屋的門框上和你的城門上。 \end{tabularx} \\ \\ \relax
6:10 & \begin{tabularx}{0.7\textwidth}{X} 「耶和華－你的神必領你進他向你列祖亞伯拉罕、以撒、雅各起誓要給你的地。那裡有又大又美的城鎮，不是你建造的； \end{tabularx} \\ \\ \relax
6:11 & \begin{tabularx}{0.7\textwidth}{X} 有裝滿各樣美物的房屋，不是你裝滿的；有挖成的水井，不是你挖的；有葡萄園、橄欖園，不是你栽植的；你吃了而且飽足。 \end{tabularx} \\ \\ \relax
6:12 & \begin{tabularx}{0.7\textwidth}{X} 你要謹慎，免得你忘記領你從埃及地為奴之家出來的耶和華。 \end{tabularx} \\ \\ \relax
6:13 & \begin{tabularx}{0.7\textwidth}{X} 你要敬畏耶和華－你的神，事奉他，奉他的名起誓。 \end{tabularx} \\ \\ \relax
6:14 & \begin{tabularx}{0.7\textwidth}{X} 不可隨從別神，就是你們四圍民族的眾神明， \end{tabularx} \\ \\ \relax
6:15 & \begin{tabularx}{0.7\textwidth}{X} 因為在你中間的耶和華－你的神是忌邪的神，恐怕耶和華－你神的怒氣向你發作，把你從地上除滅。 \end{tabularx} \\ \\ \relax
6:16 & \begin{tabularx}{0.7\textwidth}{X} 「你們不可試探耶和華－你們的神，像你們在瑪撒那樣試探他。 \end{tabularx} \\ \\ \relax
6:17 & \begin{tabularx}{0.7\textwidth}{X} 要謹慎遵守耶和華－你們神的誡命，和他所吩咐的法度、律例。 \end{tabularx} \\ \\ \relax
6:18 & \begin{tabularx}{0.7\textwidth}{X} 耶和華眼中看為正直和美善的事，你都要遵行，使你得福，可以進去得耶和華向你列祖起誓應許的美地， \end{tabularx} \\ \\ \relax
6:19 & \begin{tabularx}{0.7\textwidth}{X} 可以從你面前趕出你所有的仇敵，正如耶和華所說的。 \end{tabularx} \\ \\ \relax
6:20 & \begin{tabularx}{0.7\textwidth}{X} 「日後，你的兒子問你說：『耶和華－我們神吩咐你們的法度、律例、典章是甚麼意思呢？』 \end{tabularx} \\ \\ \relax
6:21 & \begin{tabularx}{0.7\textwidth}{X} 你要告訴你的兒子說：『我們在埃及作過法老的奴僕，耶和華用大能的手將我們從埃及領出來。 \end{tabularx} \\ \\ \relax
6:22 & \begin{tabularx}{0.7\textwidth}{X} 在我們眼前，他施行重大可怕的神蹟奇事對付埃及、法老和他的全家。 \end{tabularx} \\ \\ \relax
6:23 & \begin{tabularx}{0.7\textwidth}{X} 他將我們從那裡領出來，為要領我們進入他向我們列祖起誓應許之地，把這地賜給我們。 \end{tabularx} \\ \\ \relax
6:24 & \begin{tabularx}{0.7\textwidth}{X} 耶和華又吩咐我們遵行這一切的律例，敬畏耶和華－我們的神，使我們一生得福，得以存活，像今日一樣。 \end{tabularx} \\ \\ \relax
6:25 & \begin{tabularx}{0.7\textwidth}{X} 我們若照耶和華－我們神所吩咐的，在他面前謹守遵行這一切誡命，這就是我們的義了。』」 \end{tabularx} \\ \\
[1ex]
\hline
\hline
\end{longtable}
$^{1}$好 平安 平安 平安 各位.
感恩今天可以在這裡再跟大家一起分享上帝話語的時間.
今天的題目 為什麼要做.
為什麼要做.
我想這個問題在生活不同層面都會出現.
特別當我們面對一些很艱鉅的工作.
很不想做的功課 或者不喜歡的東西等等.
我們都會問這個問題.
不僅是在我們身上的事.
或者我們見到其他人做一些我們覺得不合理.
或者我們不喜歡的事.
我們都會問為什麼要這麼做.
在這些情況之下.
通常這個問題是一些對於這件事情不合理的諷刺.
不過在另一個角度.
我覺得去想這個問題的話.
其實這個問題也是在探究事情背後的邏輯.
思考 或者原因.
因此當我們要解答或者回應這些問題的時候.
我們需要適合的理由去回應.
當別人不理解的時候.
如果他沒有一個適合的原因.
他就會不知道如何去做.
所以我們是需要知道如何去回應和回答的.
在今天的經文當中.
其實我們都會見到上帝說要透過摩西所寫的東西.
去解釋給異式聯盟聽究竟他們要做些什麼.
也解釋給他們聽為什麼要這麼做.
在我們思想經文之前.
我想邀請大家一同再做一個祈禱.
今天我們有一同在這裡敬拜的時間.
聚會的時間.
也可以有一起去思想上帝的話語.
願主能夠幫助我們.
願主能夠帶領我們每一個.
在這裡能夠從中有所得著.
從而我們可以得著改變.
讓我們的生命被主所使用.
讓我們成為祝福.
讓我們能夠去祝福更多的人.

$^{41}$願主能夠帶領我們.
奉主的聖命祈求.
Amen.
今天經文就在《申命記錄》章1到25節.
不過我會分很多部分去和大家讀.
也邀請大家一起去讀這些部分.
讓我們可以有一同去學習.
我們一起先來讀第一個部分.
就是《申命記》的6章1到3節.
我們一起讀.
1,2,3起.
這是耶和華你們神所吩咐教訓你們的誡命律例典章.
使你們在所要過去得為業的地上進行.
好叫你和你子子孫孫一生敬畏耶和華你的神.
謹守他的一切律例誡命.
就是我所吩咐你的.
使你的日子得以長久.
以色列你要聽.
要謹守遵行.
使你可以在流乃餘物之地得以享福.
人數極其增多.
正如耶和華你祖的神所應許你的.
當我們看這段經文當中.
其實我們就會見到.
很明顯的就是摩西儘量去解釋給以色列民聽.
究竟他們要做些什麼.
也告訴他們為什麼要這樣做.
其實這件事我覺得在讀《申命記》的時候都挺經常出現.
因為當我們讀到《申命記》的時候.
其實我想上個禮拜Ivan有提到.
其實《申命記》是一個摩西給當時以色列民要入迦南地之前.
給他們一些最後的祝福或者一些要講的提醒.
當我讀到《申命記》的時候.
感覺就好像一個很沉戲或者是很長戲的一個長輩.
當他提到的時候他會告訴你你要做這件事.
然後告訴你為什麼要做這件事.
再加上想當年是怎樣怎樣的.
所以要這樣做.
如果我們看回第五章.
就是講十誡.

$^{81}$就是上個禮拜我們去思想的東西.
後面就是講當時頒佈戒命的時候是有這樣的發生.
所以你記得要這樣進行 要遵守.
很多這樣的沉迷的說話其實都會出現在不同的地方.
包括我們今天所讀的東西也是一樣.
當摩西去講的時候就說你要這樣做.
就是因為你會得到祝福 你會享福.
也不僅是來你的後代.
也會得到很多的福分.
也會住在應許之地的時間長久.
這個就是第一段所講的一件事.
我覺得在這裡可以說是比較重要的一個篇章.
之後後面就會繼續解釋.
你要聽什麼呢 你要做什麼呢.
在那裡摩西就直接畫出一個重點.
當我們繼續看四到九節 我們一起來讀.
《申命記》六章的四到九節.
一二三起.
以色列 你要聽.
耶和華 我們臣是獨一族主.
你要盡心 儘性 盡力.
這是愛人的 你的臣.
我所吩咐你的話都要記在心上.
在你正在教訓你的子女 不論你所在哪裡.
行在路上 躺下 起來都要談論.
也要細制受傷的記號.
帶在額上為經文.
要寫在你房屋的門框上.
並利的城門上.
這一段如果我們熟悉福音書.
當時我們會記得可能其中一個事情.
就是當有一個律法師去找耶穌.
問最大的誡命是什麼.
當時他也問回那位律法師是什麼.
那位律法師就提出了兩件很重要的事情.
他提出了在《摩西五經》當中.
其中兩處的經文.
第一處 當講第一條誡命.
就是我們剛纔所讀的第五節.
就是講盡心儘性盡力愛耶和華 你的臣.

$^{121}$這件事其實很重要.
甚至他們不斷地要去背.
甚至說要有一個記號在手上或頭上.
這個記號是不是直接寫在那裡.
甚至文在那裡呢.
當然不是.
因為當時的文化當中.
當我們講有民生的時候.
其實成為了別人的奴隸.
所以在當時的文化.
甚至我們讀《理美記》.
那邊也講過不要去民生那類的東西.
所以在當時的經文來講.
一個記號就弄了一個盒子出來.
這個盒子就是.
咦 沒有圖片 給個圖片來看看.
圖片 圖片.
這個圖片特別在.
對了 這裡 沒錯.
這個盒子裡面.
其實就是講.
特別右下角.
我的樣子可以飛過去.
就是那個盒子.
那個盒子會放在頭上.
也會綁在手上.
那個盒子裡面會裝著.
就是其中一個寫著我們剛纔讀的第五節的經文.
到今天 很多保守派的.
保守的猶太教徒.
也都會仍然帶著這些東西.
甚至他們作為軍人的時候.
仍然都會帶著這些東西.
到今天我們仍然會見到.
就是為了讓他們可以不斷記得.
我們要敬畏上帝.
其中一些專家就說.
當以色列今天能夠復國.
其中一個原因就是.
因為他們不斷弄了那個記號.

$^{161}$或者那個盒子.
放在頭上不斷記得要敬畏上帝.
所以是一個這樣的事情.
直到今天我們仍然會見到那個習慣.
剛纔說的房屋的門框等等.
就另外有一個類似裝飾盒的東西.
就會掛在門框裡.
就是這一段的一些重點.
然後我們繼續下去.
就是說在第十節到第十二節.
就會繼續告訴以色列民聽.
你要這麼做就會得到很多的祝福.
也會告訴他們是怎樣.
我們看第十到十二節一起來讀.
一二三.
耶和華你的神領你.
他的耶穌.
亞伯拉罕以撒雅各.
起世應許給你的地.
哪裡有城邑有大有微.
非你所建築.
有房屋裝滿各樣美物.
非你所裝滿的.
有作成的水井.
非你所作成的.
還有葡萄園.
橄欖園.
非你所栽種的.
你喫了而且飽足.
哪時你要謹慎.
免得你忘記從你同埃及地.
圍爐之家領出來的耶和華.
去到這裡我們看到就是說.
再一次摩西解釋給他們聽.
提醒他們.
接著下來你們會進入那個艱難地.
那個艱難地會有很多東西.
你會得到很多的祝福.
那些全部給你.
但那些祝福當然是令他們開心.

$^{201}$但開心之餘千萬不要忘本.
千萬不要忘記.
那是上帝所給的.
來自上帝的.
而值得注意的是.
其實以色列.
如果我們看回它的歷史當中.
它曾經在這件事上有一個很大的教訓.
在之前.
由一月到三月.
在我自己服侍的教會當中.
其實我們在思想舊約的歷史書的經文.
什麼叫歷史書呢.
在舊約當中我們分為五大部分.
有《律法書》或者《摩西五經》.
之後有《歷史書》.
然後《詩歌智慧書》.
還有《大先知書》和《小先知書》.
歷史書就是由《若司雅記》講到《以斯帖記》.
其中當我去思想有關歷代志.
或者《列王記》裡面的經文.
我們就會見到當時有的皇帝.
在以色列王朝之後.
大魏去到所羅門.
特別是所羅門就開始有些驕傲的跡象.
而這些跡象慢慢也都影響之後.
慢慢開始他們忘記上帝.
也都可以說是他們之後滅國的原因.
因為得到祝福.
但慢慢就忘本.
而這件事.
當這件事回頭去看《摩西》.
當時他所給的提醒其實非常之重要.
尤其重要.
就是說他們不可以忘本.
一定要記得是上帝所給的.
因此我們接著下來.
再提醒就是說.
我要敬畏上帝.
我們一起來讀13到15節.

$^{241}$123.
你要敬畏耶和華你的神.
事奉他.
只著他的名起誓.
不可隨從別神.
就是你們四位國民的神.
因為在你們中間的耶和華.
你神是記者的神.
唯恐耶和華你神的怒氣向你發作.
就把你從地上燒.
我們見到摩西再一次提醒.
要敬畏耶和華.
免得被除滅.
但是我們就好像剛纔所提及.
之後的以色列就是因為不敬畏耶和華.
導致之後的多諾滅亡.
摩西這個一而再再而三的提醒.
其實再一次是非常非常重要.
甚至他不斷地提醒整個申命記.
都是講同樣的東西.
接著下去.
就去到第十六節.
到第十九節了.
又講些什麼呢.
我們一起來看.
123.
你們不可試探耶和華你們的神.
將你們在馬薩那樣試探他.
要留意遵守耶和華你們神所吩咐的誡命.
法道律例.
耶和華眼中看為正看為善的.
你都要進行.
使你可以幸福.
並可以進去得耶和華向你列祖.
起世應許的那美地.
照耶和華所說的.
從你面前揮出你的一切受的.
在這裡再一次提醒一些以色列民.
要謹慎自己.
不要像在馬薩一樣.

$^{281}$馬薩發生了什麼事呢.
如果我們回到《出埃及記》就有提到.
在講以色列人沒有水喝.
所以就和摩西吵.
和摩西在罵.
然後我們回到《出埃及記》.
特別是第十七章七節的經文.
最後就提到.
當那件事結束了就說.
他給那地方起名叫馬薩.
又叫米利巴人.
因以色列人爭罵.
又因他們試探耶和華說.
耶和華是在我們中間不是.
所以當我們回到剛纔的經文.
如果說不可試探.
好像在馬薩一樣.
其實就是提醒以色列人.
千萬不要像在馬薩一樣.
就是說當上帝不在.
或者懷疑上帝究竟在不在.
這個就是說試探上帝.
當說這件事的時候.
就是提醒以色列人要敬畏上帝.
更加不要懷疑上帝是否同在.
這件事之後就是說要建立一個敬畏上帝的生活.
基本上去到這裡.
已經有一個很完整的個人的生命.
不過他的提醒不停在這裡.
我們繼續讀.
然後我們之後再思想.
我們繼續讀.
之後的經文就是我想大家讀到尾.
由二十節到二十五節.
有兩版的經文.
我們一起來看.
一二三.
日後你的日子.
耶和華我們臣吩咐你們的這些法道律例典章是什麼意思呢.
你就告訴你的兒子說.

$^{321}$我們在埃及作個法老的奴僕.
耶和華用大能的手將我們從埃及領出來.
在我們眼前將重大可怕的神蹟騎士.
施行在埃及的魔法佬並他全家的身上.
將我們從那裡領出來.
要領我們進入他向我們列祖.
起世應去之地.
把這地賜給我們.
耶和華又吩咐我們進行這一切律例.
要敬畏耶和華.
我們的神使我們嘗得好處.
蒙他保存我的生命.
像今天一樣.
我們約照耶和華我們臣所吩咐的一切誡命.
謹守遵行.
這就是我們的意料.
在講完個人敬畏上帝之後.
不可事奉上帝之後.
就提到有關若果兒子聞的時候.
我想我們不是每個人都有兒子.
我兒子也未去到那個可以清楚講話的時間.
不過我想我們都曾經經歷過那一段我們什麼都想問的時間.
或者我們都想問.
為什麼這樣東西有意義.
或者為什麼要這樣去做.
好像剛纔我們開頭所問的.
或者跟我們的題目一樣.
為什麼要這樣做.
我們以前可能問過.
為什麼要上學.
為什麼要做功課.
為什麼要做這樣.
為什麼要做那樣.
當然我覺得不只是孩子或兒子兒女會問.
其實我們身邊的人都會因著我們所做的事情.
覺得好奇而問同樣的問題.
為什麼要這樣做.
我自己都被人問過.
為什麼要信上帝.
為什麼要上教會.

$^{361}$為什麼要做傳道人.
等等這些問題都是在問原因.
在這段經文用兒子.
是因為兒子最容易受自己影響.
也是最會問自己的人.
摩西在這段經文其實是提醒以色列民.
當有人問為什麼要這樣做的時候.
就要懂得解釋給他們聽.
主是帶領以色列民出埃及的主.
也是次應他之地的主.
所以要敬畏他.
好了.
整件事.
整段經文.
讀完.
去到這裡我想問大家.
一起總結一下.
邀請大家去回想一下.
如果我們做一個總結.
今天的經文會有什麼主題.
有什麼致詞.
我們會怎麼去看今天的經文呢.
如果我們要講今天的主題或者經文的話.
我自己會講的是.
我們要敬畏上帝.
可能我們會講到要傳承.
因為講到有子女.
我們要怎麼去分享給其他人聽.
可能是在講這樣的東西.
甚至如果我們看回.
剛纔一到三節都是在講那個主題.
甚至是我覺得很重要的一點.
就是我們要敬畏上帝.
要謹守遵行他所講的東西.
然後我們要告訴子女.
你們也要謹守遵行.
但是是不是隻會停在這裡呢.
其實當我們看回這段經文的時候.
要回想一下.
這段經文.

$^{401}$當摩西寫出來.
要寫給這群以色列民的時候.
他們是要進入英雄之地.
他們是進入一個新的環境.
他們是進入一個新的地方.
就像上次提到.
其實當你進入一個新的環境的時候.
當然會有很多文化習慣.
是當地本身的人帶出來的東西.
為了不要和他們一樣.
所以有十節去解釋.
我究竟不要做些什麼.
而當摩西提醒以色列民的時候.
就是在講希望他們.
不只是去敬畏等等.
是希望去到新的地方的時候.
能夠建立一個敬畏上帝的環境和文化.
而這件事是其實.
我想到今時今日.
我們作為一個基督.
其實我們都應該去想一下.
要去做一下.
要去思想一下.
究竟怎麼去做.
我們今天傳福音.
或者我們今天希望別人來到教會去信主的時候.
其實都是在講同樣的問題.
就是我們怎麼幫助人能夠一起來敬畏上帝.
跟著我們從而建立一個敬畏上帝的環境和文化.
不過很多時候.
當我們面對這個任務.
或者挑戰的時候.
我們都會想的一個問題就是.
我不是一個有能力的人.
我只不過是一個small potato.
我怎麼會做到這些東西.
我都沒有那個經歷.
那我怎麼去講呢.
因為我們都會說.
在當時的經文來說.

$^{441}$告訴那幫以色列民聽.
他們不同.
因為我們回想起來.
我們講回申命記的那段經文.
我問大家.
其實是給哪些人呢.
我們都會想到是以色列民.
是一幫正在進入英雄之地的以色列民.
但是他們是不是一幫有經歷的以色列民.
因為當我們讀申命記六章的這段經文當中.
又提到馬沙 又提到出埃及.
suppose他們因為有這些經歷.
所以摩西解釋給他們聽的時候.
他們纔會能夠有一個共鳴.
所以他們纔會做到.
我不知道大家會不會和我一樣.
在當時第一次去思想的時候.
會有這樣的想法.
不過我想告訴大家.
事實不是這樣的.
當摩西提醒這幫人的時候.
原來這幫人是完全不是同一幫人.
當講出埃及和馬沙的時候.
其實這兩件事.
在經文當中提到的兩件事.
其實全部都是在出埃及記記載.
特別是開頭.
但是去到民訴記就有些變化了.
因為民訴記如果大家都知道的話.
為什麼叫民訴記呢.
就是因為它有兩次的人口普查.
第一次的人口普查是在民訴記的第一章.
第二次的人口普查是在民訴記的二十六章.
那裡就提到這件事.
這些就是被摩西和祭司爾利亞沙所掃的.
他們在網瓦平原與耶利哥相對的.
約旦河邊掃點以色列人.
但被掃的人中沒有一個是摩西和祭司亞倫.
從前在西內的曠野所掃的以色列人.
因為耶和華論到他們說.

$^{481}$他們必要死在曠野.
所以除了耶福福利的兒子.
迦勒和亂的兒子約宣雅以外.
連一個人也沒有存留.
簡單來說就是以前那群曾經經歷過出埃及.
曾經經歷過馬殺武水的事件的人.
一個都沒有.
除了摩西 除了迦勒和除了約宣雅.
當在這段經文當摩西作出提醒的時候.
其實這群人是完全沒有經歷過這些事情的人.
因此當講到這點.
我就開始更加明白.
原來摩西當作出這樣的提醒.
不只是叫他們敬畏上帝或者傳承.
其實是在幫助這群人再一次理解.
為什麼我要這樣做.
再次幫助這群人去理解.
為什麼要敬畏上帝 為什麼要建立這些生活.
而當明白這個原因.
我們纔有能力去做上帝所給我們的任務.
因為當不明白的時候.
其實我們很難去做.
我想對於新一代的以色列民非常重要.
因為他們只不過按照以前的人所留下的習慣生活.
和大部分信二代一樣.
我不知道我們當中有多少是信二代.
我自己是因為我父母信主.
然後我就信主.
甚至我父母也是一個全職的傳道人.
因為我父母是在家裡第一個信主.
他也非常經歷上帝的救恩.
所以他很熱心地去服侍.
他很熱心地去不斷做不同的事情.
去回報上帝的恩典.
但對我以前來說.
我沒有這個經歷.
我也不明白.
變成我被逼每個星期要跟著他去聚會.
開頭小時候還好.
但到了中間青少年時期.

$^{521}$我開始問為什麼要這樣做.
為什麼要跟著他去聚會.
為什麼要一起去做這些事情.
特別是我那時候是星期天聚會.
如果是星期天聚會.
大家都知道星期天有很多很多的活動.
我以前喜歡打籃球.
除了學界學校的比賽之外.
大部分學校外的比賽或者訓練.
都是在星期天之後就舉行.
而我自己星期天就是因為要去聚會.
所以不被允許去參加這些事情.
在那些時候我更加問為什麼要這樣做.
不明白 不理解.
但當今天我能夠明白理解當中的意義之後.
做這些事情就很不一樣了.
去聚會就很不一樣了.
去服侍就很不一樣了.
去傳福音就很不一樣了.
去講道就很不一樣了.
為什麼 因為明白當中的意義.
明白為什麼我今天要這樣做.
我不知道大家有沒有理解到.
其實今天我們所有的這些習慣.
為什麼要這樣做.
當我們能夠明白裡面的意義的時候.
其實是幫助到我們有一個能力.
去建立敬畏上帝的環境文化.
也是剛纔在經文檔中所提到的.
所給意識聯盟也是給我們的任務.
但今時今日很可惜的是.
其實很多人可能習慣了回來.
但是不明白當中的意義.
或者沒有想過當中的意義.
因此我很想邀請大家.
我記得今次是小功課.
其中一個要做的功課就是.
我們試一下思想一下.
今天我們所擁有的習慣.
我們為什麼要這樣去做.

$^{561}$不只是包括我們在普通生活上的習慣.
但包括我們為什麼要回教會.
為什麼要祈禱.
為什麼要讀經.
等等的東西.
我們試問回頭.
為什麼要這樣去做.
我不是說問完為什麼之後不做.
不是.
就是問完為什麼.
讓我們明白當中的意義.
再讓我們有能力去做.
很多時候我們不明白.
所以我們選擇了不做.
或者我們選擇了沒有能力去做.
而明白當中的意義尤其重要.
因為摩西也是解釋給這些人聽.
他們本身不明白.
沒有經歷過這一切.
但解釋給他們聽之後.
就是希望他們能夠自己去找回.
去理解.
去體驗.
去經歷上帝給他們的那份救恩.
同樣今天我們也要付出.
我們經常說.
如果我們像以前那些使徒就好了.
我們就可以經歷所有東西.
我們就有能力.
是不是代表今天我們就經歷不了呢.
不是.
我們仍然要付出努力.
去明白理解當中的意義.
我們要明白當中.
究竟上帝要我們做的.
為什麼要這樣做.
就是因為上帝是給我們救恩的上帝.
上帝是看管我們的上帝.
上帝也是保守我們的上帝.
上帝也是帶領我們能夠出困苦的上帝.

$^{601}$當我們回頭.
或者我們回到上帝裡面.
去思想這些東西的時候.
我們就會得著那種意義.
去做上帝要我們今天做的事情.
為什麼要做.
這個問題.
在屬靈上我覺得是不斷要去尋求.
不斷要去想.
我們要去祈禱.
當我們思想在生命當中.
每一件事.
每一個祝福.
每一個恩典.
主人我們都知道是來自禰的.
禰供應我們.
讓我們的需要得著.
有所得著.
讓我們能夠被禰所看顧.
禰也不斷保守我們讓我們有所平安.
禰帶領我們的生命.
讓我們到今天仍然能夠見到上帝的恩典.
主人當我們回想這些一切.
我們就明白.
我們今天應該為禰而活.
讓我們真的明白當中的意義.
去活出敬畏上帝的人生.
讓我們能夠建立一個敬畏上帝的生命.
從而讓很多人問我們為什麼要這樣活著.
我們也能夠有一個機會.
讓我們可以去見證,去分享.
因為主耶穌在我的生命當中.
做了很多的作為.
因為主耶穌救了我.
主啊!.
願禰能夠帶領我們.
願禰能夠幫助我們.
不斷讓我們在主裡面成長.
讓我們今天可以同步思想你的話語.
願主能夠祝福我們.

$^{641}$願主能夠讓我們不斷在主裡面.
能夠明白.
也能夠不斷地活出一個敬畏禰的生命.
Amen.
\newpage



\section{馬太福音 18:21-22}
\label{sec:032046}
\textbf{得饒人處且饒人}~~~~唐志濠牧師
\newline
\newline
連結: \href{http://www.church.com.hk/acms/content.asp?site=cdc\&op=show\&type=product\&code=032046\&layout=sermon}{\texttt{032046}} ~~~~ 語音日期: 2026-04-26
\newline
\newline
\hyperref[sec:032045]{\small{< < < PREV SERMON < < <}}
~
\hyperref[sec:index_chronic]{\small{[返順時目]}}
~
\hyperref[sec:032047]{\small{> > > NEXT SERMON > > >}}
\newline
\newline
馬太福音 18:21-22
\newline
\begin{longtable}{cl}
\hline
\hline
章節 & 經文 (和合本修訂版)\\
\hline
18:21 & \begin{tabularx}{0.7\textwidth}{X} 那時，彼得進前來，對耶穌說：「主啊，我弟兄得罪我，我當饒恕他幾次呢？到七次夠嗎？」 \end{tabularx} \\ \\ \relax
18:22 & \begin{tabularx}{0.7\textwidth}{X} 耶穌說：「我告訴你，不是到七次，而是到七十個七次。 \end{tabularx} \\ \\
[1ex]
\hline
\hline
\end{longtable}
$^{1}$首先我們要學習.
我想再一次和各位說.
大家有不同.
所以我們不是去學習避免衝突.
因為是避不了的.
反而有時激烈的討論.
甚至乎有些叫做激烈的相交.
可能如果處理得好.
會成為一個我們的關係之間的幫助.
所以如果中間有嚴規.
或者有誤會.
或者有一些事情的時候.
我們千萬不要表面和諧.
但往往會將一些矛盾和不滿.
收藏在心裡.
好像我們這樣說.
掃到地氈底裡.
表面就江湖升平.
內裡就破爛不堪.
很多人不會.
我們現在的人都很有品.
不會想當面去指責別人的過錯.
有什麼就不想說.
收藏就算了.
收藏在心裡.
然後再收藏.
自己受不了的時候.
就慢慢選擇離開.
轉小組.
轉牧區.
轉教會.
結果都是重複又重複.
解決不了問題之後.
最後就說和自己說.
我相信神.
我不用回教會的.
這個也是我在教會牧養了很多年.
見了很多不同的狀況.
所以我在鼓勵大家.
其實我們要學習的是.

$^{41}$我們要去面對衝突.
當有矛盾有衝突發生的時候.
我們真的要學習.
要活出基督的愛.
去處理.
而更加重要的時候.
我們不是神沒有教過我們.
怎樣去處理對方.
有的.
但不是最重要.
最重要的是處理我們自己的心.
聖經其中一個最重要的教導.
就是我們要學習.
饒恕和被饒恕.
還有我們要學習.
怎樣在主裡面勇敢的去認錯.
成為我們一個屬靈生命的美好的品格.
當我們教會能夠活出這個美好的品格的時候.
我們肢體之間很多的.
嫌隙很多的仇恨都可以消除.
我們能夠可以.
因著主的愛.
我們可以繼續在基督裡面.
我們可以所謂的同心往前.
去前行.
繼續的去活出.
神所吩咐我們在地上應該要活出的使命.
我們首先讀第一段的聖經.
《馬太福音》.
十八章二十一,二十二節和三十五節.
一二三.
那時彼得進前來對耶穌說.
主啊 我弟兄得罪我.
我當饒恕他幾次呢.
到七次夠嗎.
耶穌說.
我告訴你.
不是到七次.
而是到七十個七次.
你們各人若不從心裡饒恕你的弟兄.

$^{81}$我天父也要自一樣待你們.
彼得就來問耶穌.
他就說想為這個饒恕設一個上限.
他就問耶穌.
七次夠不夠呢.
想著很聰明.
三次就少一點.
不如七次吧.
七也是一個好的數字.
想著耶穌會讚他的.
因為其實饒恕在人性上是一個很難的課題.
憑自我感覺.
我深信如果有人得罪過你.
你一次也很難饒恕他.
單憑感覺其實真的很難饒恕.
因為感覺總會告訴你.
你是那個受害者.
它總會讓你記得之前所被人傷害的那份傷痛.
得罪了我又不認錯.
又不知錯.
怎樣原諒他.
原諒了他不是很划算嗎.
你們是不是這樣想.
但是弟兄姊妹我想跟大家說.
其實神說饒恕真是跟我們說的.
苦毒是會累積的弟兄姊妹.
你越不在最初期的時候.
你做好饒恕的工作.
你就越難饒恕.
你越難去做得到.
只是繼續傷害的是自己.
不肯饒恕就好像你拿著一杯毒酒.
在搖來搖去好像很輕鬆.
一邊搖一邊喝一邊想.
我看你什麼時候死.
為什麼呢.
人家傷害你.
晚上睡不著的人會是你.
每晚都睡不著.
每晚都在輾轉反側.

$^{121}$在想為什麼自己要被傷害.
為什麼別人怎樣得罪你.
怎樣傷害我.
怎樣辛苦.
值不值得呢.
好.
我們信主的.
祈禱.
我們祈禱求神幫我解決.
所有得罪過我的人.
教會可能不見了一半人.
放心.
看看旁邊那個.
你還在.
大家都還在.
你們應該是另外那一半.
不肯饒恕帶來了很多憤怒和偏激的思想.
會令到你的人越來越苦毒.
好像剛才所說.
越積越多.
好像雪球一樣.
越滾越大.
溫偉耀博士曾經解釋過什麼是天見.
跟大家分享一下.
當你去生氣一個人.
當你去討厭一個人的時候.
你再看不到他任何的優點.
你每一次見到他.
每一次想起他的時候.
就總是慣性挑剔他的缺點.
挑剔他的過失.
他做錯過的事.
你再沒有一些容量去看他做任何一件好事.
如果別人稱讚他.
說他好的時候.
你聽不入耳.
但如果有人說他不好的時候.
你不僅信到十足.
還不停記得十足.
這就是偏見.

$^{161}$有沒有點頭?.
大家點頭了吧.
我自己也有很多偏見.
大家也有很多偏見.
於是耶穌回答彼得.
他說不是到七次.
是到七十個七次.
好像很難明白.
其實意思就是不設限的一個饒恕.
饒恕是不會設限的.
你要不停去實踐饒恕.
其實耶穌很愛我們.
耶穌是為我們好.
其實饒恕不是益對方.
而是釋放自己的一個很大的功課.
這位女士叫Erica Curt.
她37歲.
是美國人.
丈夫是一個美國政治家.
叫Charlie Curt.
在右邊的位置.
31歲.
很年輕.
Charlie在去年2025年9月的一場公開辯論裡.
中槍被槍擊身亡.
殺他就是中間的位置.
在葬禮裡.
右邊的Erica在她的分享當中.
她公開原諒了兇手叫Tyler Robinson.
是一個一點都不容易的決定.
是一個公開的原諒.
但Erica她認為是對自己生命的一個解脫.
她這麼說.
世人或許認為寬恕是一個軟弱的表現.
但事實其實剛剛相反.
敵人拿走了我的丈夫.
也拿走了他的性命.
我不寬恕我將會失去更多.
我不願意敵人長期佔據了我的心.
我不想整天想來想去都是想起他.

$^{201}$我謹記每一天神都在寬恕我.
所以選擇原諒其實是一種自由.
是一個很實際很實在的一個分享.
面對被傷害其實人可以選擇去恨別人.
去恨所有的人去恨神.
然後將自己病了.
但也可以選擇繼續去信靠神.
而我們要看到Erica她選擇了後者.
她願意將自己交託給神.
去醫治她失去了丈夫這個傷口.
她願意靠神去繼續生活.
不讓仇恨使得她身上的傷害更加加深.
所以能夠寬恕能夠饒恕.
能夠忘記我們所說的.
Forgive and forget的人是有福的.
放過人其實就是放過自己.
讓自己得到自由.
是一個不容易的功課.
但我們憑著信心我們依靠神.
我們深信我們所信的神是在人不能.
在神凡事都能.
問題是我們信不信神是神的能力.
當然當我們願意去坦誠面對.
我們承認自己軟弱的時候.
不是說叮一聲就會立刻你由很憤怒變成不憤怒.
完全可以你不是一個這樣的神跡騎士.
而是一個過程裡面是一個意志釋放的過程.
你願意去跟隨神的時候.
你就願意教託.
讓神的能力成為了你的幫助.
這是第一點.
我們接著讀以弗所書四章三十二節.
要仁慈相待全憐憫的心彼此饒恕.
正如神在基督裡饒恕了你們一樣.
我們看到其實我們能夠懂得可以饒恕.
是因為神先饒恕了我們.
我們本身都是一個罪人.
其實撒旦最厲害的招數去騙我們是什麼呢.
就是用罪疚感.
我們很多時候.

$^{241}$我們知道我們經歷了耶穌基督的寬恕.
但是我們仍然有罪疚感去纏繞我們.
很多時候我們不能夠去做到寬恕人.
是因為我們連自己都寬恕不了.
隱藏的內疚變成了對人的仇恨.
將我們困在監獄裡.
弟子們我們要深信.
耶穌釘在十字架上的能力.
祂為我們釘在十字架上.
我們的罪赦免.
我們就是真的被赦免.
我們能夠活出寬恕.
因為基督已經寬恕了我們.
寬恕是神一個很偉大的屬性.
是對我們世人一個無條件愛的彰顯.
正如我們剛才最初舉的例子.
該隱殺了阿伯.
但是神最終也寬恕了該隱.
祂讓該隱存活.
祂沒有處死該隱.
在《創世記》四章十五節的時候.
耶和華會給該隱立了一個記號.
免得人遇見到祂就去殺祂.
而在新約的時代.
耶穌為猶大洗腳.
和他一起領聖餐.
盼望猶大回轉悔改.
但結果猶大最終出賣了他.
祭司長文氏和法利賽人.
羞辱去釘死耶穌.
耶穌在十字架上仍然去求父.
父說赦免他們.
因為他們所做的.
他們不知道.
我們的神就是這麼偉大.
祂的屬性就是要讓我們知道.
讓我們先去經歷.
祂這份無條件的寬恕.
我們人總是有雙重標準.
很多時候我們都有雙重標準.

$^{281}$被人得罪的時候很容易生氣.
對方為什麼不立刻跪下斟茶認錯道歉.
到你得罪別人又怎樣呢.
沒有的,對方不會介意的.
給點時間,過了就沒事了.
然後給點藉口自己.
甚至乎說.
他不會介意我們的.
是不是這樣呢.
但我們看一段聖經.
約翰一書一章八至九節.
我們若說自己沒有罪.
就是欺騙自己.
真理就不在我們裡面了.
我們若認自己的罪.
神是信實的是公義的.
必要赦免我們的罪.
洗淨我們一切的不義.
所以如果做錯就要認錯.
我們不要輸.
是因為我們不懂得被饒恕.
所以我們要先學習被饒恕.
我們要經歷被饒恕.
我們首先要學習認錯.
懂得說對不起.
而不是說對不起.
請你原諒我.
這句話無論對神對人.
都是一個勇敢靠著耶穌德性的行動.
對不起.
其實認錯是基於你對對方的重視.
通常你知道人與人之間.
有這些情況都是很親的人.
和你經常聯絡的人.
經常溝通的人.
包括你的配偶,你的家人.
你的同事,你的弟兄姊妹,你的老友.
所以認錯是基於你對對方的重視.
基於你對所在的群體的關係的重視.
其中這個群體裡面包括有神.

$^{321}$包括有人在我們當中.
我們不肯認錯很多時候是因為我們要面子.
因為我們驕傲.
但是我們究竟在地上.
誰覺得自己真的會完美.
你一點都沒有錯.
你不會做錯事.
所以我們要在主裡面經常學習的是.
我們要承認我們自己的不足.
我們要謙卑在主裡面.
我們因為神的愛的緣故.
我們先經歷了神的愛.
所以我們可以藉著這份愛.
我們可以選擇認錯.
我們可以選擇道歉.
其實對自己是一個很大的釋懷.
其實懂得認錯.
有錯就要認,打就要站定.
是一個很大的釋懷.
你不再被這些罪咎感纏繞著你.
當我們能夠經歷了被赦免被寬恕的時候.
這個就會成為我們的能力.
我們自己本身經歷了被赦免被寬恕.
我們就能夠知道對方被赦免被寬恕是多麼美麗.
我們就有能力去寬恕曾經傷害過我們的人.
加拉太書說,基督釋放了我們.
為使我們得自由.
所以弟兄姊妹得自由的是我們.
所以我們學習被寬恕.
不是叫你到處作惡.
做錯事,做壞事.
然後去得罪人,到處去求別人寬恕你.
而是我們要真誠地.
來到耶穌基督的面前.
我們承認我們自己是一個罪人.
願意我們認罪悔改.
我們接受耶穌基督的寬恕.
接受祂赦免的這份恩典.
邀請聖靈進到我們當中去提醒,督責我們.
願意每一天去掌管我們.

$^{361}$去讓我們去審查我們自己.
我們每一天自己一路的思想.
我們的言語,我們的行為,我們的態度.
從而我們能夠可以這樣去拒絕罪.
和犯罪得罪人所帶來的罪疚感.
就正如哥林多後書五章十七節所說.
若有人在基督裡,他就是新造的人.
舊事已過,都變成新的了.
看看你旁邊左右的那位.
我們在主裡面都是一個新造的人.
Amen?.
好,我們繼續讀第三段經文.
在彼得前書四章八節.
最要緊的是彼此切實相愛.
因為愛能遮掩許多的罪.
要恕的原因完全是因為神的愛.
要恕不是因為對方好.
不是對方肯改,不是對方認錯.
當然他們肯改肯認錯是最好.
但不是神的愛仍然足夠帶領我們要恕.
神的愛在哪裡彰顯?.
其實神是聽我們禱告的神.
神的愛就是聽我們的禱告.
我們有什麼將我們內裡傾心吐意.
有什麼不開心,有什麼苦讀,有什麼困難.
我們都可以寫給祂.
在《馬太福音》裡教導關於禱告的權柄之後.
耶穌緊接著說,就是在說搖樹的比喻.
所以其實禱告和搖樹的關係非常密切.
我們要靠著我們的禱告能夠達到神的根前.
神讓神加添我們的能力去學習搖樹.
主禱文大家都很熟悉,是吧?.
是的,但是大家究竟只是念.
還是我們的生命真的願意活出主禱文的教導呢?.
其中兩句說:免我們的債,如同我們免了人的債.
究竟我們是念口往,還是我們真的這樣向神祈禱呢?.
人家欠你多少呢?.
其實你想想你自己欠神多少呢?.
原來這句的禱告是說.
我們去求的時候,叫神去免我們的債.

$^{401}$如同我們免了人的債.
哪樣先行?是我們先去搖樹.
我們不肯搖樹,我們搖樹得少.
神就免得我們少.
正如耶穌在後面之後,他接著說.
就是這段經文《馬太福音》六章十四至十五節.
他說:你們搖樹人的過犯.
你們的天父也必搖樹你們的過犯.
你們不要樹人的過犯.
你們的天父也必不搖樹你們的過犯.
所以耶和華其實很有憐憫,很有恩惠的.
搖樹其實是他的屬性.
他基本上是會搖樹我們.
甚至我們將他的兒子釘在十字架上.
他都將世人搖樹.
但他這裡唯一很清楚,他在這段經文裡表明的話.
就是他不會搖樹.
如果我們不搖樹.
天父也不會搖樹我們.
所以耶穌在後面也跟著他說.
如果我們向弟兄發怒.
或者向弟兄懷恨.
我們都免不了會受審判.
所以我們拒絕搖樹.
就是我們自己決定和神的關係.
要斷絕和神的關係.
我們會失去神的恩典.
神肯先搖樹我們.
所以我們肯搖樹並不是因為一些法律的規條.
不是強行要我做.
我多不願意我都要做.
而是我們願意在神裡面.
我們去回應神這個無條件的愛.
在教會裡面我們既然知道.
相處,剛才也說了.
有些事我們真的避免不了.
很容易發生言語誤會,發生矛盾.
當然一方面我們要謹慎自己的言談.
好好管理自己的舌頭.
不要以為熟就什麼都可以說.

$^{441}$我們都是成熟的人.
我們應該要學習什麼叫做有分寸.
我們不要以為什麼說什麼都可以.
弟兄姊妹不會介意.
另外當然我們也求主幫助我們.
有一樣東西我們很需要求主幫助我們.
是我們能夠建立一個健康的自我的形象.
不要做容易受傷的男人.
或者容易受傷的女人.
因為我們在主裡面有一個尊貴身份.
當我們去肯定我們去確信.
我們是主裡面的一個寶貴的子女.
我是神寶貴的兒女的時候.
就算身邊發生什麼事.
在外界發生什麼事.
我們都不會容易受到影響到我們自己的形象.
不會那麼容易影響到我們的情緒.
很多年前我初信主的時候.
認識了一個我覺得他非常愛神又愛人的弟兄.
他是我初信主時候的屬靈導師.
其實我欣賞他的地方.
是他無論在待人接物方面.
我都見到他是一個可以活出聖經教導裡面.
謙虛溫柔忍耐.
用愛心寬容.
用和平聯絡.
不輕易發怒.
也不計算人的惡的一個人.
有時候我們在同一個團契.
當時很多時候相交很多時候見面.
當我們覺得他有時候被人用言語攻擊.
被人bully被人tease的時候.
又被人欺負的時候.
看他好像若無其事.
好像很naive.
我問他是不是傻的.
你有沒有感覺的.
好像沒有什麼感覺.
他都是用一些穩定的情緒.
用穩定的言語來回應.

$^{481}$有時候甚至我們身邊的弟兄姊妹愛他.
親他.
想保護他.
不值得想幫他撐.
然後問候他的時候.
為什麼你不回應.
你是不是忍得這麼辛苦.
但是他的回答都很特別.
他說我根本不覺得對方在攻擊我.
是嗎.
他在攻擊我嗎.
我不是很覺得.
他真是有一個在神裡面很獨特的性情.
就是他不容易受影響.
不容易受傷害.
這件事我記得很清楚.
也跟了我.
也記了到現在今時今日.
我想我憑我的性情很難做到.
我很想去求.
如果神給了我這個性情.
一個寬宏亮度的恩賜.
我覺得很感恩.
這個很難去學到.
但是我相信.
是要在神裡面.
跟神有一個很好的關係.
可以建立到這位弟兄.
《真言》十章十二節裡說.
恨能夠挑啟爭端.
愛能夠遮掩一切的過錯.
跟《新約彼得前書》.
我們剛才讀過四章八節一樣.
他特別在《真言》也提到.
愛能夠遮掩罪.
愛能夠遮掩過錯.
在這裡的意思不是包庇.
不是隱藏.
不是說神的愛就是.
你面對一些犯罪的事情.

$^{521}$你都不會計他的罪.
不會計較.
而是在裡面我們能夠.
像剛才那樣有一個寬宏的亮度.
我們不會張揚.
不會到處跟所有人說他錯.
但是就沒有跟那個人說.
也不是急於去定罪.
而是很希望.
對方能夠在主裡面.
真誠的悔改.
能夠在主裡面被主挽回.
在我預備這篇信息的時候.
三月底的時候.
就有同工告訴我.
你知不知道台中發生了一件街頭命案.
一個26歲的女子.
因為拒絕復合.
就被她的前男友拿著長刀.
當街捅她的喉嚨.
捅死了她.
刺她的頸.
那個女子的父母.
原來報道說她在高雄的一間.
叫做天傳611靈良堂的牧師和師母.
在4月5號的追思禮拜裡面.
牧師和師母就面對著失去了一個愛女.
也面對一個這麼殘暴.
這麼兇殘的行為的時候.
她仍然好像剛才我講的Erica那樣講.
她願意放下憤怒.
她說用基督的愛和饒恕.
去勸這個兇手.
希望他能夠真正的悔改能夠勇於面對.
所以不是饒恕並不是認同他的罪.
他仍然要勸這個兇手.
去面對他將要面對地上的法律.
或者地上的制裁.
但是這對的父母選擇去饒恕.
是為了要活出基督的愛.

$^{561}$第一他為自己好.
另外也很希望對方能夠被主去挽回.
可以在主裡面認罪悔改.
所以我們看到總括來說.
我們看到的這個饒恕.
是完全是因為從神的愛而來.
神的愛和恩典可以足夠我們去饒恕一些.
甚至我們認為是不值得饒恕的人.
我自己家庭裡的一個分享.
大家認識我太太.
教會裡面的國牧師.
我叫她Lillian.
我最初認識她的時候.
她和我分享她年輕時代的照片.
她不笑 每張照片都不笑.
我們問她為什麼.
她說我有隻爆牙.
所以笑起來不漂亮.
後來她說她也處理了爆牙.
但是還是不笑.
相信了耶穌之後.
開心了一點.
但是好像一直還有一些東西放不下.
當然認識了她.
知道她多一點的時候.
當然知道其實是因為她媽媽和她自己.
饒恕 原諒不了.
拋棄了他們在菲律賓有另一頭家的爸爸.
她媽媽四十歲的時候患了抑鬱病.
病了二十多年.
一直抑鬱.
因為自己沒有工作.
又照顧著幾個子女.
一直抑鬱 患了二十多年.
病了二十多年.
很感恩.
我們去和她傳福音.
她相信了耶穌.
她經歷了神的醫治.
但是仍然都因為丈夫的傷害.

$^{601}$她仍然都不是那麼開心.
整天都不開心.
子女對她多孝順.
對她多乖.
好像都沒有什麼用.
甚至有一次和女兒說.
將Lillian摔在地上.
因為媽媽和她說.
女兒 你很乖 你對我都很好.
不過你對我多乖多好多孝順都沒有用.
一單歸一單.
你對我好都不能填補.
我都會繼續不開心.
就是因為我不能忘記.
你這個壞爸爸怎樣對我.
這也讓Lillian很不甘心.
我們很不甘心.
為什麼自己和媽媽不開心.
一直都不開心.
是不是因為一件過去的事情.
就一直帶到一直下去都不開心.
特別到媽媽越來越年老的時候.
明明說信了耶穌.
明明信了耶穌.
將來回到天家的時候.
上天堂是永遠的喜樂.
那怎麼可以帶著一件地上的不開心的事.
離開這個世界呢.
想做些什麼.
剛才就說對她好又沒有用.
對她怎樣好都沒有用.
人真的幫不了人.
無奈我們只可以同心一切在神面前祈禱.
神就吩咐我們.
神跟我們說.
正如那時候說.
我們要學習在聖經裡面的教導.
去學習寬恕.
尤其那時候我們剛正值零幾年的時候.
我們領受呼召.

$^{641}$我們準備進入教會全職事奉的時候.
我們心想沒有理由之後做了傳道人.
都經常苦口苦面.
那怎麼面對弟兄姊妹呢.
但是心裡面又很明白.
真的不容易.
怎樣去寬恕呢.
怎樣去跟爸爸說.
是不是就這樣說我要寬恕就寬恕得到呢.
所以一定要帶著行動.
但是行動又怕再度受傷.
真的很難做到.
但是作為她的先生.
我當然會全力跟她一起祈禱.
全力去支持她.
在2007年的時候.
我們就特別安排了一個假期.
應該是復活節.
安排了一個假期.
我們就去了一次宿霧旅行.
大家知道宿霧是菲律賓一個度假勝地.
各位旅遊達人都知道.
如果香港人去宿霧旅行.
通常都是直航機去直航機返的.
哪會這麼轉折的呢.
但是我們就很特別.
我們是專門訂機票的時候.
飛去宿霧玩幾天.
然後回程我們選擇了飛上去馬尼拉.
然後再在馬尼拉轉機回香港.
就希望藉著有半天的時間.
可以約爸爸出來跟他說一聲.
我們不再生他的氣.
其實當時說不是因為他做了什麼好事.
他不是認錯 他不是悔改.
他不是說回來要負他的責任.
我相信完全都是因為.
耶穌的愛 耶穌的吩咐 耶穌的教導.
也是對自己的一個解脫.
是對自己將來的一個幫助.

$^{681}$如果解決不了.
於是我們就決心去解決這個問題.
但當然了.
如是者我們去宿霧玩完.
準備第二天就坐飛機上去馬尼拉.
在宿霧打個電話給爸爸.
我們會來的.
我們再一次經歷的是再一次傷害.
我們被放飛機.
爸爸說.
我不得閒 我又很忙 有客人要見.
我們不見面了 你們要來嗎.
好像當沒事一樣.
哇 直接將我們再一次落入一個傷心失望的境界.
當然沒有什麼 已經安排好了.
我們就照樣去馬尼拉.
自己吃一頓飯 睡一晚.
第二天早上就回到香港.
是不是這樣就了結呢.
但神仍然鼓勵Lillian沒有放棄.
我們再繼續祈禱的時候.
祂最終決定了.
回到香港之後 祂真誠的打了個長途電話給爸爸.
不是打算罵爆他.
而是真的跟他談了一個很長的長途電話.
說了很多的往事.
大家分享了很多的情況.
或者內心的分享.
到最後祂終於可以在結束的電話裡跟爸爸說.
我愛你.
我們做了我們的事 我們最終沒有見過外父.
而外父也在我們07年過了一段時間之後.
他在菲律賓離開了.
但我深信神的恩典 神的神蹟騎士.
就是臨在Lillian和我們的家庭裡面.
我看到她慢慢地去釋懷.
在過去被拋棄 被傷害的仇恨.
不再纏繞祂.
祂開心了很多.
你現在看到祂是不是很開心呢.

$^{721}$我感恩 而且不單只是這樣.
神的神蹟騎士更加臨在媽媽的身上.
媽媽信了主 有返向日葵的崇拜.
也是被我們叫她信了.
一個信徒 每個禮拜返崇拜.
但是自從發生了之後.
她竟然會慢慢跟我說.
不知為何我返崇拜越來越開心.
我以前返崇拜時不懂得笑.
不想跟弟兄姊妹相交.
常常想躲在一旁.
但近來不知為何.
我現在敬拜主的時候.
我初初敬拜的時候是這樣.
但到後來敬拜的時候是這樣.
到我今早返向日葵崇拜的時候.
我的敬拜可以是這樣.
她真的很開心.
好像小朋友般返來跟我分享.
神的恩典是不是很大.
大神的能力是不是很大.
不單止是這樣.
未完.
她到後期78歲的時候.
她開始有認知障礙.
開始有些東西不記得.
原來認知障礙神都有祂的恩典.
有祂的神蹟騎士.
人生有很大的片段.
由出世到她認知障礙.
但神就很奇妙地.
拿走她被丈夫拋棄不開心的記憶.
她可以中間抽了這段東西出來.
沒有了.
她跟我們還有回應的時候.
她跟我們說.
都是她記得她年輕的時候.
怎樣在農村,鄉村的時候成長.
很多開心的日子.
還有她年老的時候.

$^{761}$跟子女,孫子,孫女.
一起的開心事.
神是不是很奇妙呢.
張榮耀還贊歸給她.
弟兄姊妹.
耶穌基督與我們同在.
是我們隨時的幫助.
記住饒恕是一種態度.
我們的態度就決定了我們的高度.
求神給我們一個寬廣的心.
讓神的話語,讓聖靈.
和祂賜給我們的平安.
給我們的喜樂.
我們所思考的是神的美善.
我們學習去活出寬恕.
給神帶領我們.
放心交給神帶領我們.
去走出過去的傷害.
我們去經歷醫治.
進入祂的豐盛.
進入祂的自由當中.
記不記得以前有一個人物叫約瑟.
約瑟被他的哥哥賣了去埃及.
他哥哥不喜歡他.
賣了他去埃及.
但結果他做了埃及的丞相.
結果要接濟他的兄弟.
接濟兄弟的時候.
有一句說話他跟他的兄弟說.
我覺得很值得大家去學習.
就是從前你們的意思是要傷害我.
是跟那些害他的人說的.
但神的意思原是好的.
有一首歌叫《迎能讚美》.
大家有幾句歌詞.
我覺得很有意思跟大家分享.
他說全心相信你已編好結局.
萬事出必有因.
靠主永不抖震.
當我們相信一切都有神的旨意.

$^{801}$神的美好的安排的時候.
這是我們的信心.
我們可以相信經歷神將一些傷害化為祝福.
讓我們的生命能夠變得更加寬廣.
我們能夠經歷到神的喜樂和神的盼望.
所以我在這裡鼓勵大家.
我們能夠靠著耶穌.
去活出得遙人處且遙人.
這個美好的屬靈質素.
讓我們能夠更加成為耶穌基督美好的見證.
說就是這麼說.
我還是覺得很懶惰.
所以我們要祈禱.
好不好?我們一起唱喜樂好不好?.
祈禱,因為神是聽我們禱告.
這就是神的愛.
很難做得到的功課.
求神去幫助我們.
是一個階段,是一個過程.
是一個進程.
讓聖靈親自的充滿我們.
去彰顯祂的意志釋放.
讓我們專注在祂的愛有多麼長闊高深.
讓我們浸淫在祂救贖的恩典裡面.
讓我們思想我們是何等的罪人.
祂仍然願意寬恕我們.
願意捨棄祂兒子的生命去救贖我們.
幫助我們擴闊我們的胸襟.
我們能夠靠著主去實踐,寬恕.
甚至乎是一些不值得我們寬恕的人.
我們心裡面有沒有想到一些呢?.
如果有的話,我們可以把手放在我們的心裡面.
我們一起和神祈禱.
求主幫助我們.
讓我們專注在神的愛裡面.
耶穌基督的救贖恩情.
聖靈的帶領,聖靈的獨澤.
幫助我們.
我們坦誠地告訴神.
我們軟弱.

$^{841}$我們不行.
我們不能饒恕.
但是神的能力在我們不能.
在神凡事都能.
不是我所說的任何一件事.
我相信的是聖靈的能力.
我相信是聖靈可以在我岳母身上所做的工作.
她沒有做過什麼.
但是神可以釋放她.
我也相信神今天可以釋放我們.
在座任何一個弟兄姊妹.
幫助我們.
主我們求你以我們同在.
\newpage



\section{約翰福音 4:43-54}
\label{sec:032047}
\textbf{你是耶穌的粉絲嗎?}~~~~譚偉明宣教士
\newline
\newline
連結: \href{http://www.church.com.hk/acms/content.asp?site=cdc\&op=show\&type=product\&code=032047\&layout=sermon}{\texttt{032047}} ~~~~ 語音日期: 2026-04-26
\newline
\newline
\hyperref[sec:032046]{\small{< < < PREV SERMON < < <}}
~
\hyperref[sec:index_chronic]{\small{[返順時目]}}
~
\hyperref[sec:032048]{\small{> > > NEXT SERMON > > >}}
\newline
\newline
約翰福音 4:43-54
\newline
\begin{longtable}{cl}
\hline
\hline
章節 & 經文 (和合本修訂版)\\
\hline
4:43 & \begin{tabularx}{0.7\textwidth}{X} 過了那兩天，耶穌離開那地方，往加利利去。 \end{tabularx} \\ \\ \relax
4:44 & \begin{tabularx}{0.7\textwidth}{X} 因為耶穌自己作過見證說：「先知在自己的家鄉是沒有人尊敬的。」 \end{tabularx} \\ \\ \relax
4:45 & \begin{tabularx}{0.7\textwidth}{X} 到了加利利，加利利人都歡迎他，因為他們也上耶路撒冷去過節，曾經看過他在節期間所做的一切事。 \end{tabularx} \\ \\ \relax
4:46 & \begin{tabularx}{0.7\textwidth}{X} 耶穌又到了加利利的迦拿，就是他從前變水為酒的地方。有一個大臣，他的兒子在迦百農病了。 \end{tabularx} \\ \\ \relax
4:47 & \begin{tabularx}{0.7\textwidth}{X} 他聽見耶穌從猶太到了加利利，就來見他，求他下去醫治他的兒子，因為他兒子快要死了。 \end{tabularx} \\ \\ \relax
4:48 & \begin{tabularx}{0.7\textwidth}{X} 耶穌對他說：「若不看見神蹟奇事，你們總是不信。」 \end{tabularx} \\ \\ \relax
4:49 & \begin{tabularx}{0.7\textwidth}{X} 那大臣對他說：「先生，求你趁著我的孩子還沒有死就下去吧。」 \end{tabularx} \\ \\ \relax
4:50 & \begin{tabularx}{0.7\textwidth}{X} 耶穌對他說：「回去吧，你的兒子會活！」那人信耶穌所說的話，就回去了。 \end{tabularx} \\ \\ \relax
4:51 & \begin{tabularx}{0.7\textwidth}{X} 正下去的時候，他的僕人迎面而來，說他的兒子活了。 \end{tabularx} \\ \\ \relax
4:52 & \begin{tabularx}{0.7\textwidth}{X} 他就問甚麼時候見好的。他們對他說：「昨天下午一點鐘熱就退了。」 \end{tabularx} \\ \\ \relax
4:53 & \begin{tabularx}{0.7\textwidth}{X} 他就知道這正是耶穌對他說「你的兒子會活」的時候；他自己和全家就都信了。 \end{tabularx} \\ \\ \relax
4:54 & \begin{tabularx}{0.7\textwidth}{X} 這是耶穌從猶太回到加利利後所行的第二個神蹟。 \end{tabularx} \\ \\
[1ex]
\hline
\hline
\end{longtable}
$^{1}$弟兄姊妹平安.
今日真係好高興可以返到自己屋企.
呢個唔係隨便ge 教會.
呢個係我自己屋企.
因為我認識好多弟兄姊妹.
好多弟兄姊妹係度長大ge .
同埋好多導師都係度教導我ge .
咁樣啊.
真係好開心今日有機會係度同大家分享我地宣教ge 經歷.
旅程.
但係我都同你講先.
你大部份都認識我係邊個啦.
我都好心跳.
因為我今日要講廣東話.
同埋你地知道我廣東話ge 文化水平唔高.
但係先要打下底.
因為有D 弟兄姊妹唔認識我.
咁我呢D 中文ge 水平.
基本上係交TBB學費ge je .
所以有D 詞呢.
比較地道D 呢.
即係英文比較ghettoD 呢.
請大家見諒啊.
咁樣.
已經去jor 三年半啦.
咁我都估唔到咁快已經一任響到台灣.
咁我響台灣真係有D 習慣ge .
最唔習慣返到黎第一樣ye 就係.
點解咁會動ge 芝加哥.
咁樣啊.
咁我響台灣宣教呢.
基本上呢我係響.
冇好日響高雄.
我係響直堂ge .
響屏東市.
一個好細ge 城市.
響高雄ge 東邊.
大概四十五分鐘ge 車程.
咁教會你見到.
我右邊gwo 張呢就係.

$^{41}$大概教會前面ge 門口啦.
閂jor 閘而家.
但係左邊gwo 張就係條街係咁樣.
咁但係我地教會大概係三十至四十人啦.
咁呢度一家人響到影幅相比你知道.
大部分呢弟兄姊妹呢都係大過我ge .
咁我響到其中ge 職位呢.
就係做一個年青牧師.
咁我其實我發覺.
我去到三十幾四十歲.
而家先去做年青牧師.
其中一樣ye 呢.
我以前都做jor 十三年架啦.
即係gwo D 團體導師.
我發現我唯一一樣ye 做唔到.
而家就係我唔可以同佢玩全程.
我真係好gui 而家比較.
係但係.
另外做年青牧師呢個職位呢.
就係我會搞好多.
我同太太會搞好多D 夏令營呀冬令營囉.
咁你見到左邊gwo 張呢.
就係我第一次搞夏令營ge 時間.
係gwo 時教一個細路都冇ge .
咁我地點樣做呢.
我同太太同埋一個留學生.
係去一個小學gwo 度呢排擔仗呢.
我大大聲聲咁樣同佢地.
大家用英文對話.
比佢地見到我地D 標準ge 英文.
咁啦.
第一ge 暑假又大概四個學生.
加埋我兩個仔女有六個.
咁第二呢第二年呢就D 家長.
見到呢度幾好喎咁樣.
咁唔需要我再去派傳單.
咁就家長幫我推銷就容易好多.
除jor 做年青牧師.
我地建堂呢.
唔係話我淨係做小學ge .

$^{81}$淨係細路仔ge 工作ge .
其實基本上呢乜都需要做ge .
我無論係招待呀.
淨係去serve午餐呀.
呢D 全部都需要做ge .
其實呢呢D 咁ge 策略呢.
同埋點樣做法呢.
我完全都係響UV呢學番黎ge ye .
出面要去人入黎呀.
淨係招待呀.
定點樣有人serve呀.
有cafe呀.
呢D 全部我呢D 教弟兄姊妹.
都係響到學番黎ge .
有都要帶成人ge 呢個bible study啦.
同埋小組啦.
最後呢我有機會主持一個婚禮添.
有個弟兄叫我主持婚禮.
咁我都未做過ge 其實.
第一次ge je .
我見過好多次.
跟住我問AI.
通常應該點樣ge .
一個baptistge 婚禮.
但係跟住同我D 經驗.
我話大概係咁樣啦.
應該冇問題ge .
咁就係咁樣ge .
真係發覺到.
去到gwo 度.
你冇人幫.
冇人.
你乜ye 都需要做ge 其實.
但係其中一樣ye .
弟兄姊妹成日都.
我去到gwo 度呢.
其實弟兄姊妹好尊敬我地.
因為我地係全教師ge 身份.
同埋係牧師ge 身份.
但係我都成日.

$^{121}$聽佢講一句.
我就係咁樣話佢講啦.
我現係我ge 職位.
係一個牧師.
但係我跟你將來ge 關係.
係乜ye .
我跟你係弟兄姊妹.
所以你唔應該稱呼.
我係Steve牧師.
只有稱呼我Steve.
還是.
佢唔知道我叫阿明.
咁你地呢其實.
我都叫阿明仔得架啦其實.
大家返黎係咪先.
我D 導師係度.
好似Jimbo咁樣睇住我大ge .
係咪先.
唔好意思呀Jimbo.
爆jor 你年齡先.
所以.
咁今日ge 主題.
想學你講兩種ge 信心.
第一種呢就係.
好淺ge 信心.
我地睇一睇.
黎到話.
係呢度ge 背景就係基本上.
耶穌呢就係從耶路撒冷.
御治ge 緣故.
返到黎.
佢地自己ge 家鄉.
但係佢特登係傾過.
呢個薩瑪利.
因為佢要讓薩瑪利亞人.
薩瑪利亞人.
知道佢自己ge 身份.
係表面.
咁佢最後呢就返到黎呢度呢.
但係好奇怪第44章.

$^{161}$跟住你見到佢話.
記得耶穌講.
呢個先知.
係無尊敬ge .
但係45呢跟住佢返到黎.
嘩D 加里利亞人就好開心.
歡迎佢.
但係我地見到佢歡迎佢呢.
因為耶穌做jor 好多ge .
好多ye .
係耶路撒冷當時.
但係我見到佢做jor D 乜ye 呢.
如果你返返去2章.
你唔需要返去我講.
你返去2章23節呢其實呢.
耶穌做jor 好多神蹟.
咁你相信佢.
咁你發現44同45.
點解約翰爾係咁樣寫呢.
係咪有D 衝突呢.
好似玩ye 喎.
但係佢講呢一句呢.
如果你睇返其他D 福音書呢.
通常講呢一句ge 背後呢.
就係.
通常D 人呢係拒絕耶穌.
唔相信佢呢佢就講呢句.
佢回應呢句.
但係係度冇喎.
佢係冇任何ye .
都未返到黎屋企約翰爾已經講呢句.
點解呢.
其實佢想同你講.
佢想做一個介紹俾你知.
呢個成個故事呢.
其實呢D 加利利人呢.
佢地要歡迎耶穌呢.
其實佢ge 信心真係好膚淺ge .
佢好表面去歡迎佢.
咁我地去到46節呢.

$^{201}$咁而家耶穌真係返到黎.
自己屋企.
特登返到加拿大呢個地方.
咁你見到佢.
約翰就.
講就話.
佢呢個地方係將水變走ge 地方.
第一間神職.
咁即係話佢想俾我地一個印象就話.
其實好多人知道呢.
耶穌係度做jor 呢件事.
但係同一時間好多人都聽聞.
佢係耶路撒冷做ge 事.
即係話.
耶穌呢而家係.
比較出名啦而家開始.
跟住.
黎緊47節你見到.
46節.
下半節gwo 度呢你見到.
開始有個衝突就.
開始發生就係呢個官.
有個黎求耶穌去.
醫治佢仔因為佢仔就係死啦.
但係你發現.
呢個官其實唔可以命令耶穌叫佢.
醫治咩.
去逼迫佢咩.
其實佢可以喎我覺得.
但係我覺得去到一個地步呢.
我覺得佢已經.
用jor 好多唔同ge 方法呢佢先會.
肯去求耶穌一個猶太人.
去醫治佢仔.
佢發現.
即係你發現一個人有財有勢.
去到最後gwo 個.
終極點呢都係死亡.
都係打唔過死亡呢個最後ge 敵人.
唯有我地所知道.

$^{241}$耶穌復活先係.
可以能夠打贏呢個死亡.
響歷史記載只有.
耶穌基督係有.
唯一ge 人係一個空ge 墳墓.
響48字呢你見到.
耶穌去回應.
佢回應就話.
你地班人你係相信我.
因為你見到呢D 神蹟奇妙.
咁佢話你地呢.
英文個U呢響.
希伯來文.
響希臘文就話唔係單數係雙數.
佢唔係單單同呢個官講ye .
同其他所有ge 觀眾講.
話你地係相信.
你地係相信.
我係因為你見到呢D .
神蹟奇妙.
所以耶穌就話.
你地成日都係好鍾意睇呢D 好漂亮ge ye .
好犀利ge ye .
好似D 魔術咁ge ye .
佢地見到耶穌係一個咩.
係一個好偉大ge miracle worker.
同埋而家就係出jor 名.
返返黎家鄉.
咁耶穌帶jor 好大個光榮返黎家鄉.
因為點解呢.
如果你知道.
猶太地南部.
耶路撒冷京城咁.
佢點睇加利利北邊D 人呢.
基本上睇佢地就係鄉下佬.
無文化.
咁耶穌而家就係變jor 一個.
好出名ge gwo D .
名人.
咁呢一班加利利人.

$^{281}$就好似D 粉絲咁樣.
係追住呢個名人好似潮流咁樣.
好開心.
咁你諗諗而家.
最潮流係乜ye 呀.
其實我都唔係好清楚.
我見過D 後生ge 一輩.
好似我D 仔女gwo D .
最鍾意係Taylor Swift.
我諗好似有D 大人都可以cancerge .
所以我覺得最潮流係gwo D 但我又唔好識.
咁你追粉絲.
但我唔知道.
但我只係知道我以前追係乜ye .
就係我地工牛隊23號.
Michael Jordan.
Michael Jordan佢係邊度出世.
佢大學讀邊度呢.
佢ge 阿爸阿媽係咩名呀.
佢D 仔女係邊個ge 呢.
個名我都講得出.
但係我同佢係完全冇關係.
我認識佢但佢唔認識我.
就係呢樣ye .
咁第二樣ye 就係.
我想再返返去44字gwo 度.
就係講話一個先知係冇尊敬ge .
響指ge 家鄉.
咁其實.
我都覺得好難ge .
你一個人.
你一班人睇住耶穌.
甚至乎去認識耶穌ge .
係好難去.
服從係個權柄之下ge .
其實我覺得.
就好似我以前呢.
我所有個電腦gwo D .
我D 智慧呢我都係教我阿叔ge .
其實佢一D 電腦都唔識ge 之前.

$^{321}$但返到黎佢而家話有新ge software.
要去教我點用.
我成日覺得有D 怪怪地.
無可能呢D ye 咁都發生.
所以好難我覺得.
但係你發現.
之前D 薩瑪利亞人.
冇呢個問題.
因為佢冇呢個包袱.
佢就基本上耶穌係冇咩.
做咩神蹟.
佢聽見耶穌ge 話.
佢已經信.
佢信係乜ye ?響42字佢話.
佢信耶穌係救世主.
呢D 神蹟奇妙呢.
其實係幫助我地.
去指證.
俾我地知道耶穌係邊個係彌賽亞.
係救主係神的兒子.
讓我地將.
我地gwo 個信心可以擺係我地gwo 度.
呢D 係幫助我地ge .
信心ge 其實.
唔係去滿足我地ge 人.
你見到我地ge 關係.
呢度冇真正ge 關係.
同埋佢地唔需要做任何commitment.
呢D 人.
佢ge 信心真係好淺.
但係我可以同你講.
大部份ge 人我響台灣.
我遇到都係咁樣.
人黎到教會.
基本上.
唔係話想認識耶穌.
我自己係一樣.
佢地唔見到耶穌係主.
佢地知.
大部份響台灣就係.

$^{361}$見到耶穌係一個好偉大ge 醫生.
好犀利ge 醫生.
佢地有一個神燈.
佢地唔想同耶穌有關係.
去到呢個地步.
我有一次去到.
三幾個月前.
有個.
女士同我講.
我地一齊有個茶經組.
佢話我好多年都未.
訓覺過一個好.
睡眠睡得唔好.
佢話如果你幫我可以.
祈禱比我.
訓得好呢我就相信佢.
我心諗我滴哂汗.
大佬你想我點.
我唔知我點.
但我都幫你祈禱.
我照幫你祈禱.
第二日真係.
佢話我訓得好好.
咁係咪真係信耶穌.
咁我話.
巧合黎ge .
咁我又玩ye .
但我睇一睇其他D .
example就係.
有D 人.
好直接.
黎教會為jor 餐飯.
有D 為jor 關係.
有D .
黎學英文.
如果你要講耶穌就唔好.
咁我要同佢講.
我地係真係學英文.
但係我生命.
係因為耶穌.

$^{401}$我先有今日.
我唔可以唔講耶穌.
好難搞有時有D 人.
另外D 家長.
好有goal oriented.
即係好有目標.
黎教會.
禮拜二日.
淨係玩.
學品格.
點做好人.
呢sai 時間.
佢梗係冇咁同我講.
但佢ge 意思係咁.
我一至六話我要去補英文.
禮拜六仲要夜晚去補課.
雖然禮拜六唔需要返學.
但我地又點算.
我地有半個鐘頭英文.
我地學對話.
跟住仲要學樂器.
好彩我老婆識少少.
另外就係.
有機會去.
唔係成日都有得做.
去拆神台.
有位弟兄.
信jor 耶穌.
叫我去拆神台.
我都未做過呢D ye .
但係發現台灣人.
如果信jor 耶穌.
拆jor 佢家裡面ge 神台.
唔代表拆jor 佢心ge 神台.
有D 台灣人.
弟兄姊妹.
將佢以前gwo 個偶像搬過黎基督教.
耶穌變jor 佢偶像.
你見到所有ge ye .
ge 目標黎教會.

$^{441}$唔係因為耶穌.
係咪有其他ge 好ge 事情.
佢可以得到D 乜ye 先黎教會.
佢ge 生命係冇改變.
佢只係想上帝祝福佢.
但佢唔想同上帝有關係.
弟兄姊妹想問.
咁我地呢.
咁我地係咪咁樣做.
我地係咪黎教會.
因為ge 兒童士工好.
講道好.
當然講道好.
我響Markgwo 度學到好多ye .
但係我唔想你地做粉絲.
弟兄姊妹.
我想你做門徒.
所以我地睇一睇下一點.
點做真正ge 信心係點樣.
49至54.
咁但係唔似.
呢個官就唔似gwo D 粉絲咁.
佢再去.
要求耶穌.
你快D 過黎啦,我仔就黎死架啦.
但係耶穌話點呀.
你唔.
你ge 仔會.
即係唔會死.
咁你諗一諗.
你會點做呢.
你相唔相信.
我地邊個餐館.
我地個個人都睇過D google review.
睇下好唔好食.
睇下其實D 人點講.
佢就係咁樣講,gwo 時有無google review.
你點知呀.
但係呢.
你見到其實係好多地方.

$^{481}$下一節就係咁樣.
gwo 個人去呼求耶穌好多次.
耶穌做過ge 神蹟去醫治好gwo 個人.
今次你見到無.
其中一個原因就係.
佢唔想去滿足呢班粉絲.
第二呢.
佢就係需要呢個人去揀.
我冇俾奇蹟.
你見唔到奇蹟.
你會唔會相信我個話語.
所以我真係奇蹟呢.
唔係我地基督徒.
我地可以去build upon我地ge 根基.
係上帝ge 話語.
咁但係你最後你見到.
呢個官員係無得揀.
佢最後呢.
係需要接受耶穌所講.
而家好表面講ge 話.
跟住離開jor .
至少你見到佢肯去接受離開.
你見到佢呢個信心係慢慢.
即係有增加ge .
跟住去51節.
去52節就係佢返屋企.
我覺得佢返到屋企.
佢差唔多返到屋企gwo 時佢個僕人就見到佢.
佢話.
你ge 仔而家冇事喎.
咁正常黎講.
我地個個都好開心.
其實你搞掂gwo 單ye 呢.
你唔會搵返gwo 個人去幫助.
你唔會搵返gwo 個醫生.
但係佢呢就係冇.
佢就諗起.
52就諗起.
佢幾時好返ge .
我通常唔會諗我仔幾時發燒好.

$^{521}$我唔會去諗返幾時gwo D ye .
但係佢有諗過.
因為我覺得佢記得耶穌講ge ye .
跟住佢去對一對.
原來真係咁arm 耶穌講gwo 時候.
佢就真係醫治好.
你發覺我地基督教gwo 個信心.
唔係淨係得個feel.
我覺得耶穌今日好好.
其實都要有根據.
要考ge 要用邏輯去諗.
我地.
我地gwo 個.
Faith.
但係佢見到佢跟住點樣做.
佢真係將呢個信.
擺係耶穌gwo 度佢唔係淨係信耶穌.
佢係一個好偉大ge .
Miracle Worker.
佢係信耶穌係呢個彌賽雅呢個救世主.
佢唔單止信.
佢將呢個福音.
傳俾佢ge 家人.
佢全家人相信.
但你諗一諗弟兄姊妹.
佢話.
佢途中係一點鐘.
前一日佢先知.
佢個仔好返喎.
從迦納去到加伯隆.
大概係二十個mile.
要用一日ge 時間.
係gwo 度.
佢過夜有機會係過夜一晚.
你諗一諗.
你過一晚.
你ge 心情係點樣.
其實我覺得好多弟兄姊妹.
我地大家都經歷過個心情.
上帝ge 應許得唔得架.

$^{561}$即係有冇咩料架呢個.
即係呢個上帝.
我覺得gwo 個掙扎呢.
真係好難過過一晚.
但係我地個個都經歷過呢樣ye .
但係呢樣ye .
上帝係要我地ge struggle.
係呢方面去依靠.
佢ge 話.
之後佢引證jor .
真係推動jor 佢ge 信心.
我想講.
我地真正ge 信心.
唔係一次搞掂就搞掂.
呢個救恩就係.
信心每一次.
我地同上帝ge 關係.
佢每一次每一刻.
佢每個季節每一年.
佢都將呢D trials.
呢D 試煉.
擺係我地ge 身上.
讓我地去認識佢多D .
去信教多D .
因為你諗一諗.
我地幾時最唔去搵上帝ge 時候.
係我地最安穩ge 時候.
有時都唔記得佢係邊個.
大部份人.
好少數ge 人.
會將.
真係想慢慢將呢個生命.
咁樣係完全.
交比上帝.
完全相信佢.
呢個係我地每日.
弟兄姊妹每日記錄我ge 功課.
都係我ge 功課.
但係我係台灣雖然大部份人係咁.
但係我都好開心.

$^{601}$我見到一班.
好少數ge 人.
會肯將呢個生命.
即係佢地ge 信心慢慢去生長.
其中一個就係.
我ge 大查經.
我gwo 時都唔知點.
大查經.
都係大概五六個弟兄姊妹肯黎.
呢D 全部都帶我.
我黎到話查咩好呀.
你查咩都得啦Steve.
若何福音.
我查jor 五六次架啦我聽過.
OK咁你查jor 五六次.
咁我再試一次睇一睇.
跟住呢我就.
就係度所學架啦我係神學院.
同埋係教會學ge .
咁教導下唔係淨係咁睇.
同我講聖經ge 背景.
咁點樣.
gwo D 神蹟.
即係點systematically.
咁去睇.
咁佢話我未學過呢D ye 架喎.
咁就係.
原化於法言.
佢明白地理ge 背景呢.
佢幫助佢ge 信心.
佢知道呢度.
若何講咩ye 幫助佢信心多好多.
跟住佢成日都上黎.
去到gwo 地步呢我其實同佢講.
我2月22號係我最後一日.
係呢個教會.
咁我諗住係1月.
最後gwo 個禮拜六呢就算啦.
停啦因為2月中係新年.
咁佢地話.

$^{641}$Steve唔得喎可以繼續開.
我話點解呀你地新年唔係好多ye 搞架咩.
佢預備搞乜鬼呀.
而家仲有出去食飯架你仲煮架咩.
搞咁多ye .
跟住佢話佢想繼續去.
我做宣教士其中最開心.
ge 一份.
即係我開心ge 時候就係呢樣ye .
D 濟英之輩想明白上黎ge 話語.
第二呢.
呢兩個呢.
係我第一批ge .
gwo D 下令營呢.
gwo D 學生gwo 次有4個學生加埋我仔.
呢6個呢佢係其中兩個.
叫阿Sam同阿Tina.
呢個仔呢特別係呢個仔呢.
即刻我桌底呢.
有桌布ge .
佢即刻捐jor 入去啦.
佢唔肯見人ge 好驚你點樣ge .
咁我成日都哄佢哄佢.
成個東令營呢都係咁樣架啦.
佢阿媽呢之後都繼續去.
擺佢響我地ge 賴皮日ge .
團契gwo 度.
過jor 一年咁樣.
就開始同我friendjor 啦.
咁樣就慢慢有機會.
即係同佢講聖經故仔啦.
但係一難ge 我覺得.
大部份台灣屏東市ge 家庭.
我所認識全部係單身家庭.
佢地已經唔認識阿爸.
佢以為阿爸已經同阿媽.
分開jor 好耐.
所以我好開心另外一樣ye 就係.
可以響gwo 度做一個男性ge 榜樣.
我覺得上帝ge 呼召係我生命.

$^{681}$其中一樣ye 就係呢樣ye .
咁樣我地你見到.
佢gwo 幅相其實呢佢都對我.
都okge 而家.
所以呢個有個progression.
係需要時間架.
係需要比佢知道.
我真係關心佢.
呢兩個就係.
黎到我地教會大概兩年前.
佢係響台北.
做工ge 但係burn outjor .
因為台北好辛苦.
返返黎屏東市.
我發現好多鄉下都係屏東市.
返黎搵一日教會.
搵到我地ge 相遇教會.
咁佢黎到.
咁你知D .
通常大過我.
同我咁大ge .
我通常mark住佢地.
咁我都要做好我ge 工作.
咁其中一樣ye 我地一齊分享.
我成日都同佢地出去食飯.
同佢地傾計.
佢話響以前ge 教會.
成長gwo 時好多傷害.
咁樣我大家傾下D dialogy.
其中一樣ye 其實.
台灣其實唔係冇教會.
其實有時我地覺得dialogy.
未必係最arm ge .
佢覺得好多台灣人覺得.
我浸jor 泥先係相信耶穌.
我未浸泥就唔係相信耶穌.
咁其中一樣ye 就係.
佢呢個女仔就係好.
特別女仔就係好.
我覺得係一個perfectionist.

$^{721}$佢覺得乜ye 都係要.
我將所有100\%.
全部俾晒上帝.
我唔論係我工作上我教會上.
我都要俾佢最好ge .
我都要做到最好.
我都要同佢傾一傾.
你呢樣ye 真係好.
但係你係現實有時.
你唔要自己抖下.
所以你成日之前burn out.
我唔係話你應該係工作做得唔好.
係應該做得好.
但係你點樣去balance your life.
你係俾到你ge 生命.
過個禮拜四ge gwo 一日.
俾到上帝就係gwo 一日.
我唔係想discourage.
弟兄姊妹唔去事奉.
你自己去搵個balance.
去搵個平衡.
但係見佢慢慢去肯去做呢樣ye .
同埋去.
真係好認真咁樣.
去寫低佢D 筆記.
我地上.
上主學個時.
同埋.
最後去幫助我地翻譯.
係我其中一個翻譯.
因為之前我可以用英文講到.
咁最後呢一個家庭係一個好.
貧窮ge 家庭.
咁佢又因為係.
下令營ge 緣故第二年.
佢就聽聞我地有下令營.
佢就黎參加我地下令營.
咁樣呢個好貧窮.
因為佢地全家.
包括個丈夫都好多.

$^{761}$D 先遷性ge 病.
所以lor 晒D 錢去.
去睇病.
咁佢地其實都唔係話未去過教會.
佢地都去過.
唔同ge 教會.
佢有一次.
同我講.
我地成個月已經食緊粥.
佢女兒話可唔可以唔食粥.
我話發生咩事呀咁大鑊.
台灣餓唔死.
你買個便當都平.
但係發現.
D 錢已經去哂gwo 度.
佢跟住有一日同我講.
你地教會雪櫃有咩呀.
我可唔可以lor 埋黎去食呀.
gwo D 雞塊你lor 去啦.
如果你需要.
即係料我姐係夾過ge .
我地可唔可以lor .
我話唔係咁大鑊呀嘛.
跟住認識jor 佢.
佢慢慢同我分享.
我以前有教會ge 傳道人.
同我地講.
我地全家有呢D 病.
因為你唔相信耶穌.
嘩大佬咁.
呢個咩傳道人黎架我心諗.
跟住都同我講.
佢D 朋友都話.
你過黎做呢個教會ge 會員.
你睇做會員有幾多ye lor .
又有米lor 又有gwo 樣ye lor .
佢話如果我只係為jor .
去教會做會員可以lor 到ye .
咁我呢個信係咩黎架.
咁我真係過jor 一晚.

$^{801}$我真係好傷心呀我打邊爐呢.
我都打到火滾呀真係.
我真係知道gwo 個傳道人係邊個.
我真係想搦架撐去打鬼死佢.
但係我唔知呀嘛.
真係好氣我個人.
我跟佢出黎講.
得我地一定要.
去而家去ge 超市食完飯.
我地.
成卡ye 我地full up佢.
將所有D 油鹽.
米醋呀乜ye .
都擺晒整滿佢.
我即刻lor 去佢屋企我同佢講.
呢個係耶穌俾ge 教會俾ge .
我地冇條件ge .
所以佢慢慢你會明白.
真係幾難去拆解呢D .
傳統呢D 信仰.
跟住呢.
佢慢慢.
佢會問我福音係咩黎.
有時發覺.
弟兄姊妹係需要時間同埋關心.
所以弟兄姊妹.
我唔知你點睇耶穌.
佢係一個好名人.
定係我地ge 主.
咁我都希望呢.
耶穌係我地ge 主啦.
咁我希望.
既然如果耶穌係我地ge 主.
我希望我地ge 生命同福音.
係相清ge .
弟兄姊妹呀.
我地ge 基督徒ge 生命呢.
其實同宣教係走唔甩ge .
你係基督徒呢.
你係傳教士黎架啦.

$^{841}$我同你ge 崗位唔同.
就係我係台灣你地係呢度.
你無論係你ge 工作你地ge 家人gwo 度呢.
你邊度你都係一個傳教士.
你係代表緊耶穌ge .
因為你記唔記得大使命呀.
gwo 個人都記得大使命.
你要去張曼文.
做我門徒.
但係你唔可以講話.
我其實好渣ge ye 我乜都唔識ge 其實.
我同你講你唔可以咁樣講ye 架.
但係因為你唔可以.
渣得過我.
我去到大學gwo 時我先arm 學識廣東話.
講都講唔正gwo 陣我地而家咪係咁樣.
跟住你話連個夢都無.
做乜ye 呀咁大ge 仔.
佢去college係三年.
去做個老師去做個好榜樣.
男性ge 榜樣所以我教小學.
之後我覺得個生活好好.
呢一世真係好好.
跟住耶穌話你去做傳教.
我話唔係呀嘛D ye 唔arm 我.
點做牧師呢D ye .
跟住話咁傳教係要讀神學院.
我就諗緊.
如果我班友知道我背景佢點收我.
佢真係收jor 我喎我發現.
跟住呢我又畢到業喎.
我都唔知點畢到業.
跟住個傳教機構又收jor 我去台灣.
而家係度慢慢考緊國語.
我真係.
我而家係堅咁講.
喂爛坦黎ge 大佬.
大佬呀我都唔知我自己.
有咩恩賜真係.
即係音樂又唔得.

$^{881}$電腦又唔係最靚.
我都唔知有咩恩賜其實而家.
我都未知有咩恩賜但係我唯一一樣ye 就係.
我就話上帝.
我係度你需要ge 話.
你可以用但你諗一諗.
如果你冇恩賜你須咩驚.
如果上帝係真係呼召你去.
呢個崗位佢會俾呢個資源.
你去做好佢啦係咪呀.
你諗諗你係你個公司你個老細叫你做ye .
佢唔俾部電腦黎架嘛大佬你自己lor 佢返黎架.
或者有D 係都唔頂我唔知.
得你明白我講咩呀.
但係我想你記得就係.
我話你你記得.
大使命但係唔好唔記得.
大誡命係乜ye 好重要.
記唔記得大誡命係咩呀.
你要愛.
主利益神.
另外gwo 個係咩呀.
你要愛人如己係咩呀.
好多時呢我地大誡命呢.
我唔記得jor .
gwo 個人呢係變jor 一個.
福音ge 推銷對象.
你愛全員福音呢.
佢唔信呢.
拜拜.
我冇關係架啦.
但係人呢係可以見到你係咪真心對佢架.
你想.
同人傳福音呢.
你真係要真心對gwo 個人.
其中一樣ye 我響孫九宮長學識.
佢雖然呢我地.
真係好難先可以去.
呢個關係好耐先可以.
建立好但係.

$^{921}$佢真係見到你真係真心對佢.
唔係因為你要推銷呢個福音比佢.
呢個好大ge 分別呀.
咁而家將來呢.
我地將來會做D 咩ye 呢.
我地而家做係響.
高雄南部ge .
響高雄南部住緊但係我會去.
新ge 教會去嘉義.
嘉義縣一個.
小ge 地方叫普治.
一個好細ge 城市.
好鄉下仲鄉下而家.
佢大概芝加哥十倍咁多啦.
芝加哥大概十倍啦.
鄉下黎架真係好ge 鄉下.
呢度呢對面呢就係而家.
我住館住gwo 度呢佢樓上呢.
就係一個好似有個.
細路仔傳團係咁.
我地而家已經去jor 一個月架啦.
咁同佢地玩啦結了關係.
呢一班細路仔呢.
都好慘ge 其實好似.
你諗一諗留守兒童係點樣.
佢地基本上係咁樣.
佢有D 父母唔係度.
去jor 第二度做工但係你諗一諗台灣ge 島好細ge .
佢地唔係好遠.
但係唔得閒去理佢地.
還是佢有D 父母呢.
就係係個廟gwo 度做ge .
咁你諗一諗佢父母係個廟gwo 度.
將來係點樣咪就係係個廟gwo 度做.
佢就話個父母就話.
呢度係咪教會黎架拍教會gwo 度唔讓我D 細路仔黎.
咁唔係咪我地唔係係度講耶穌ge .
我地呢間呢度囉.
我真係唔係講耶穌我真係同佢好似團契.
咁樣去愛佢地建立關係.

$^{961}$有時講一個.
教會即係一個故事.
一個聖經故事.
好開心我都可以繼續同佢地一齊.
有機會打波同佢教佢點打籃球.
另外第二個地方就係.
好近ge 係呢個我地叫.
恩典中心.
因為我地呢個教會係要.
恩典中心因為頭先我都講囉.
如果係教會我地唔黎架啦.
跟住有D D 人入到黎跟住.
即係崇拜呀.
一齊去傾計呀食飯.
跟住D 人就開始發覺o下D 唔係教會.
教會唔係咁做架嘛係呀其實教會就係咁ge 咋囉.
佢未入過黎佢地驚呀嘛唔知道呀嘛.
咁我地教會.
其實真係三十至四十人.
一個細教會.
我呢其實.
第一任我返去呢就去呢個地方.
其實呢係學點樣.
做一個好ge 牧者.
同一時間.
去練好我國語囉我黎緊呢間教會呢.
係完全講中文ge 即係國語.
我臨走之前.
我嘗試.
講我自己ge 見證係度.
大概用jor 十三分鐘咁樣全國語ge .
我就話我之前去台灣一個國語.
都唔識咪係咁.
夾硬黎架je 我.
但我覺得上帝係係我呢個生命中.
係有一個好大ge .
mission.
我覺得幾難呢我都會做好佢ge .
我盡量因為呢.
我做基督徒呢.

$^{1001}$我同我上帝ge 關係呢.
唔係話靠feel架.
突然之會.
係睇下我服唔服從呀.
Am I obedient or not.
呢個就係我其中一個信心.
比較增長ge 其中一個.
那個manifestation.
好啦咁將來呢.
因為我.
我真係將來我主管同我講.
Steve Shirley 你不如去.
席堂啦.
我冇諗過喎而家專教宣告都以為.
做唔到而家係席堂ge 一天.
我就話好我有一天我將來.
我目標係台灣我就要去席堂.
但我唔知點樣但我覺得因為.
即係.
又有一個異象係度.
我想再比無論係神學上.
定係我地ge 品格上我要比佢地知道.
天家係真正.
個弟兄姊妹.
應該係點樣相處ge .
多謝弟兄姊妹.
\newpage



\section{約書亞記 24:14-15}
\label{sec:032048}
\textbf{家庭文化2.0}~~~~招世超牧師
\newline
\newline
連結: \href{http://www.church.com.hk/acms/content.asp?site=cdc\&op=show\&type=product\&code=032048\&layout=sermon}{\texttt{032048}} ~~~~ 語音日期: 2026-04-26
\newline
\newline
\hyperref[sec:032047]{\small{< < < PREV SERMON < < <}}
~
\hyperref[sec:index_chronic]{\small{[返順時目]}}
~
\hyperref[sec:032049]{\small{> > > NEXT SERMON > > >}}
\newline
\newline
約書亞記 24:14-15
\newline
\begin{longtable}{cl}
\hline
\hline
章節 & 經文 (和合本修訂版)\\
\hline
24:14 & \begin{tabularx}{0.7\textwidth}{X} 「現在你們要敬畏耶和華，誠心誠意事奉他，除掉你們列祖在大河那邊和在埃及事奉的神明，事奉耶和華。 \end{tabularx} \\ \\ \relax
24:15 & \begin{tabularx}{0.7\textwidth}{X} 若你們認為事奉耶和華不好，今日就可以選擇所要事奉的：是你們列祖在大河那邊所事奉的神明，或是你們所住這地亞摩利人的神明呢？至於我和我家，我們必定事奉耶和華。」 \end{tabularx} \\ \\
[1ex]
\hline
\hline
\end{longtable}
$^{1}$可能有些人都知道,上星期對我們招家來說是個大日子.
因為上星期是師母生日,很大件事,師母生日是二十七十四.
兒子都回來了,兒子帶兩個孫回來,吃什麼,雖然是師母生日.
都要看兒子和孫子喜歡什麼,結果就吃韓國燒烤.
吃得很開心,我慣常地會做什麼呢?.
孫子和媳婦回來的時候,通常都是我招呼他們.
喝什麼,吃什麼,飯前果,飯後果,一大堆的東西.
他們喜歡玩什麼,跟他們鬥鬥開心,因為不是經常有機會見到他們.
吃完飯之後就切好生日蛋糕,我大兒子就坐在一旁.
開始拿他的手機出來查一下電郵,我兩個孫子就去客廳.
坐在地上玩Woblox,小孩子很流行的遊戲.
不過半小時之後,鈴鈴鈴鈴鈴鈴,計時鬧鐘響了.
他們很會做的,馬上關機,自己找別的東西玩.
我和師母通常吃完飯之後,會跟Parsnip和他的媳婦.
他們很喜歡跟我們同在一起,我們有這樣的習慣.
講故事,講什麼故事呢?.
講一下過去一星期的情況,我們弄得怎麼樣,你們弄得怎麼樣.
分享一下我們在事奉上的一些經驗,一些體會.
這個其實是一個很標準的,每次我們一家人一起吃飯的時候.
都是這樣的情況,不知道大家怎麼看呢?.
我們一家人吃完飯之後,特別是對我們來說是一個很典型的習慣.
我不知道你們怎麼去看或者去評估,這樣好不好呢?.
為什麼我會問這個問題呢?.
因為我知道有些家人,吃完飯之後,大家都是一家人一起吃飯.
都是坐在客廳裡,不過那些年輕的青年,少年.
就會去掃TikTok,小的那些,不像他們兩個孫子那樣.
玩Roblox,或者是其他一些流行的電玩,好像他們兩個孫子一樣.
不過旁邊就沒有計時鬧鐘的,就進去玩.
父母,公公,婆婆,爺爺,奶奶在做什麼呢?.
各自在那裡,有些在看新聞,有些在看YouTube,看什麼不知道.
不知道大家有沒有留意到,我們的朋友,家人當中.
有這樣的情況出現呢?或者會不會出現在我們家當中呢?.
如果有這樣的情況的時候,會不會我們在不知不覺當中.
將我們的子女,不經意地,變成將來的一些宅男,宅女呢?.
雖然很多時候我們仍然是一家人在一起,特別是會一起吃飯的.
間中都會,就算不是住在一起.
但是會不會仍然會覺得有些孤獨的感覺?.
家裡滿了人,到處看,這個在這裡,那個在這裡,大家坐在一起.
但是各人做自己的事,有時候都會覺得,好像很冷漠.
我希望沒有,希望家裡沒有這樣的情況出現.

$^{41}$但我知道有些人的情況是這樣的.
當然在美國,有另外一些極端的.
我們感染了另外一些美國的家庭文化.
來一來的時候,我們會習染到的.
其中一個比較普遍的,是什麼呢?.
就是對子女地獄式的教育.
什麼叫地獄式的教育呢?.
為什麼呢?背後有一個理念.
我們不知不覺地很希望我們的子女和很多.
是不是很多,我不知道,其他一樣.
要將他們打造成為金牌子女.
有沒有這樣的想法呢?.
可能我們沒有這樣想的.
不過我們的行動很多時候是有這樣的情況出現.
現在我發覺我們的中學生真的很忙.
他們家裡的生活有很嚴格的規矩,規律.
和密密麻麻的時間表.
譬如你要上advanced placement,AP的課程.
有些巡迴的體育比賽.
還有你要充實你的履歷,要申請大學.
還有很多校外課程,課外的活動.
所以很多時候你晚上吃飯的時間.
不經意地變成一個board meeting.
父母帶著你開會,開什麼會?.
講一下學業成績怎麼樣,校外課程要讀什麼.
可以爭取進入哪間名校.
還有這個星期,1月3,4,5.
每個星期哪天做什麼,在哪個時候做什麼事情.
所以很多時候吃飯就會談這些事情.
不知道我們會不會不自覺地,不是有心的.
造就了新一代成就很高,high achievement.
但是很容易出現焦慮的年輕人.
其實影響可能會更加深遠也說不定.
因為當我們這樣做的時候.
他們不會思考的,不過他們吸收了訊息是什麼呢?.
他們的價值是取決於他們的成就.
而不是取決於他們的個性和品格.
於是,說嚴重一點.
很容易就有一些脆弱的自我認同.
什麼意思呢?.

$^{81}$人的自尊心很弱的,很多時候是玻璃心.
一遇到困難,挫折,失敗,有什麼不行.
很容易就depress,崩潰.
所以上個禮拜還是前個禮拜我忘記了.
跟大家說過,在美國每十個年輕人當中.
有一個是認真想過會自殺,就是這麼嚴重.
不知道大家聽完之後.
有沒有想過自己家裡面有什麼家庭文化呢?.
今天我跟大家說的是家庭文化2.0.
你自己想一想,現在我們家的文化是怎麼樣的?.
我們的家庭文化是什麼呢?.
大家有沒有想過,我們的子女在我們現時的環境氣候.
在這個家庭文化中長大的時候.
對他們以後的人生會有什麼影響?.
說到家庭文化,可能不是常常說的family culture.
其實簡單來說是什麼呢?.
family culture就是在日常生活當中,大家不用說的.
不言而喻的一些默契.
家裡大家知道怎麼做事的,怎麼運作的.
是一些我們不用說了,大家都會做的.
一起都會去守的守則,生活上的守則.
如果我們將家庭文化比喻為一家公司,當然不是.
家庭文化就是這家公司裡的企業靈魂.
我們都不太明白,如果我們將家庭文化比喻為一台電腦.
比喻一台電腦的時候,家庭文化就是一台電腦的作業系統.
如何運作你的電腦.
所以如果我們有這樣的觀念,其實家庭文化最低限度有三種不同的功能.
第一種功能是什麼?.
防禦的功能,就好像一道防火牆.
大家玩電腦就知道有防火牆.
一道防火牆可以保護家裡的成員在一個混亂的價值觀.
很容易被動搖的.
但是一道防火牆在家裡長大.
可以保護你,就不會那麼容易被不同的價值觀衝擊影響他們.
有個保護的功能,防禦的功能.
家庭文化也有一種導航的功能.
當小朋友或家裡任何一個人都好.
當我們面對人生的誘惑,挑戰.
或者要面對一些重大的抉擇的時候.
家庭文化就會成為我們的北極星.

$^{121}$應該如何定位呢?.
應該如何做決定呢?.
應該如何走呢?.
所以家庭文化很多時候都能夠幫助我們.
平時沒事就沒事,在有事的時候.
就會主導著我們的抉擇.
家庭文化也要凝聚在一起.
這是一個功能.
我希望我們家裡不只是一個回家睡覺的地方.
以前我們很忙碌,要工作十幾個小時才能回家睡覺.
工作完就回家睡覺.
家庭文化讓我們知道.
家裡不是一個我們回家休息睡覺的地方.
是一個有獨特的氣息,能夠讓人感到安全.
感到被接納的一個地方.
所以人們說這是一個安全的地方.
躲在家裡就覺得安全.
不是因為家裡特別堅固,用石屎水來做.
而是因為裡面的人.
他們如何互相保護大家.
所以我剛才說.
家庭文化包括有防禦,有導航.
也有凝聚的功能.
其實對家裡每一個人.
不單是小孩子.
對家裡每一個人如何在世上生活.
影響是相當深遠的.
而對基督徒來說.
我相信大部分都是基督徒.
無道的朋友,我都希望你留心聽一聽.
對基督徒來說.
更是我們經常說的信仰.
如何傳承給我們下一代.
不可以忽略的一個課題.
所以今天跟大家說.
家庭文化2.0是說到.
如何建立一個基督徒的家庭文化.
當然任何的文化.
包括基督徒的家庭文化.
就不可以一朝一食.

$^{161}$忽然之間就搞定.
特別我知道我們在座的人.
不同於以前.
很多都不是小時候在基督徒家庭長大的.
很多人來美國讀書才相信耶穌.
所以有時我們沒有在那個環境裡.
經歷過,長大過.
所以現在你叫我們去建立一個家庭文化.
有時真的有些挑戰.
我們真的不知道從何開始.
甚至我們根本不知道要這樣做.
我們就去教會崇拜.
帶小朋友回來.
我以為已經做了.
很多時候都是這樣.
但是如果我們相信上帝.
我們相信神.
我們是基督徒.
我們經歷過神的恩典.
有沒有經歷過神的恩典?.
有沒有經歷過神在我們生命當中給我們那種祝福?.
我們有沒有希望.
我們真的希望,我們真的想.
我們的子女,我們的下一代.
成為基督徒去信耶穌.
以致他們不單單有永生的盼望.
當他們在世上生活的時候.
他們能夠活得有意義.
能夠活得精彩.
能夠覺得不枉此生.
我們想不想我們的下一代.
有這樣的體會和經歷呢?.
如果想的話.
弟兄姊妹.
我們就必須要在家裡.
建立起基督徒的文化.
我們可以怎樣開始呢?.
首先我們要知道什麼才可以.
首先我們要問.
我們家裡面.

$^{201}$什麼是我們的核心價值呢?.
我們的核心價值.
如果我們大家都不知道.
在家裡什麼是重要的.
什麼是寶貴的.
在家裡我們最看重的是什麼來的.
在我們家裡面.
是不是每個人都知道.
什麼是我們的指路明燈.
是我們的北極星.
在我們家裡面.
我們有沒有防火牆.
我們有防火牆是什麼呢?.
我們的平安是來自什麼呢?.
有時候我們沒有問這些問題.
現在我們就要問.
要知道.
到底家裡的文化.
到底是怎樣的呢?.
《聖經》裡面其實記載.
今天跟大家說一下.
約書亞記.
約書亞帶領以色列人.
進入迦南.
其實進迦南說得很輕鬆.
其實是帶著以色列人.
征服迦南地.
迦南裡面滿佈了很多.
不同的種族.
不同的神明.
當他們大獨已定.
已經佔據了迦南.
很有把握了.
約書亞年紀大了.
於是就招聚了所有的以色列人.
來到他們面前.
做什麼呢?.
幫他們重溫耶和華上帝.
拯救以色列人.
離開埃及的歷史.

$^{241}$約書亞提醒他們.
從前你們埃及做奴隸多久了?.
400年了.
上帝怎樣用十個災禍.
讓你們能夠擺脫了.
埃及人的統治和虐待.
雖然去到曠野.
但是神在曠野當中.
用火柱,雲柱帶領你們.
你們經過曠野.
遇到很多的仇敵.
但是神一一將他們打敗.
使用你們.
而且夜夜鎖衫.
一無所缺.
然後走乾地.
通過約旦河.
攻佔了迦南這個肥沃的地方.
將他們不是自己建造的城一切.
將不是他們自己栽種的葡萄園,橄欖園.
送給他們.
可以吃到裡面的果子.
所以約書亞提醒他們.
我快要走了.
回天閣.
我快要離開這個世界了.
我要將最重要的事.
吩咐你們.
什麼事呢?.
十四節說.
你們現在要敬畏耶和華.
誠心實意地去事奉祂.
將你們列祖在大河那邊.
和在埃及所事奉的神除掉.
去事奉耶和華神.
若是你們以為事奉耶和華不好.
今天就可以選擇.
你所要事奉的.
是你們列祖在大河那邊.
所事奉的神.

$^{281}$是你們所駐這地的阿摩利人的神.
至於我和我家.
我們必定事奉耶和華.
剛才我稍微提過.
當日以色列人.
和我們今天一樣.
他們生活的環境.
是滿天神佛的.
到處都是偶像和假神.
無論在埃及.
圖然是.
出來之後.
一直抗夜那四十年走過很多地方.
約旦河的東西嶺岸.
包括迦南地裡面.
都是滿天神佛.
完全一樣.
很多假神和偶像.
所以.
若雙俠離開世界之前.
有件很重要的事要告訴你們.
你們要做一個明智的抉擇.
做一個明智的抉擇.
若雙俠鼓勵百姓.
要敬畏耶和華.
誠心實意的去事奉神.
千萬不可以有疑心.
這是第一件事.
若雙俠在《聖經》當中.
他帶頭表態.
不過大家有沒有留意到他怎麼表態呢?.
他的表態是什麼呢?.
他說不單單是他自己會敬拜神.
事奉神.
從他的子女行間.
他更告訴我們.
他立志要自己家裡面每個人.
不單單是他的子女.
包括他的伯伯,奴隸.
都必定要事奉耶和華神.

$^{321}$這就是他在他們面前的表態.
指的是我的家.
整個家裡面.
大大小小的所有人.
上上下下.
他們都要事奉耶和華神.
所以對若雙俠來說.
他的核心價值是什麼呢?.
他家裡的核心價值是什麼呢?.
就是要敬畏神.
事奉神.
若雙俠家裡的紫露明燈.
他的北極星是什麼呢?.
耶和華上帝.
他所相信的神.
若雙俠在地上的時候.
在上天教之前.
他朝聚以色列人.
大家有沒有想過.
若雙俠在地上的時候.
他有他的工作.
他的工作是什麼呢?.
承接摩西一族之長.
他的工作是族長.
帶領以色列人.
還有.
他是開疆拓土的民族英雄.
進來之後.
打敗這個打敗那個.
加拿大就是我們的了.
他是帶領以色列人.
征服加拿大的偉大領袖.
這是若雙俠在地上的工作.
他之所以能夠完成這些艱巨的任務.
做一族之長.
不容易的.
還有.
要開疆拓土.
開創新天地.
又要征服加拿大所有的人.

$^{361}$以致他們可以在裡面安居樂業.
任何一個的工作.
都是很艱難的事.
而他只能夠完成這些工作.
為什麼呢?.
因為若雙俠敬畏神.
專心一意事奉祂.
在座的弟兄姊妹.
如果你們記得.
上個星期我跟大家說過.
如果你錯過了.
可以聽網上的說法.
如果我們自己不愛神.
不尊重神.
我告訴你們.
我們很難將我們的信仰.
傳給我們的下一代.
所以第一件事.
如果我們真的想我們的子女信耶穌.
我們自己才要做一個好基督徒.
什麼叫做好基督徒?.
敬畏神.
因為今天有太多不敬畏神的人.
太多不敬畏神的基督徒.
不單要敬畏神.
專心一意.
誠心實意.
這是若孫毅用的說話.
去事奉神.
所以今天在這個時候.
我想向大家呼籲.
如果大家真的經歷過神的恩典.
經歷過神怎樣帶領你.
幫助你,祝福你.
以致到今天.
你雖然不是大富大貴.
但不單可以一年多次旅行.
放假.
生活無缺.
開車.

$^{401}$可以吧.
不單可以.
其實是鬼車.
可以住大屋.
雖然不是大富大貴.
但我們可以享受美食.
喜歡吃什麼.
從來不收多少錢.
讓我鼓勵大家.
這是神的恩典.
大家知不知道?.
如果是的話.
是否應該敬畏神?.
是否應該專心一意去事奉祂?.
這是我們基督徒.
特別是經歷過.
嘗過主恩之味的基督徒.
當時的本分.
如果我們有機會.
安靜下來.
再去想一想.
你的家庭文化是怎樣?.
有沒有想過?.
如果你的家庭文化.
想一想.
如果你發覺.
在我家裡.
誰是北極星?.
哪些東西是我的指路明燈?.
我的良心.
你發覺你的北極星.
你的指路明燈.
原來是你自己的良心.
是你自小從家人.
朋友.
老師身上.
累積過濾後.
整理了.
這是我的價值觀.
於是成為我的良心.

$^{441}$我就根據這些東西.
去打理我家.
比如說.
做工.
要追求上進.
和Excellence.
可能你沒有寫出來.
不過在心裡.
這是你的價值觀.
做人.
要誠實正直.
家人不會這樣教你.
要奉公守法.
家人之間.
要互相照顧.
這些可能是我們自小.
如果你的家庭背景好的話.
通常你都會感染到.
類似這些價值觀.
這些都是好東西.
不要誤會.
我不是說這些不好.
最少.
這些不像類似的.
有人說人不為己天誅地滅.
這些都是價值觀.
人為財死了為食亡.
有沒有聽過.
這些都是人的價值觀.
這是自私自利的心態.
我們比這些價值觀好很多.
但是正如我所說.
如果你檢討一下.
如果現時我的北極星.
我的指路明燈.
就是我所謂自己認同的良心.
我們要小心.
如果我們真的.
正如我剛才所說.
你嘗過主人的致命.

$^{481}$我們是否應該像若遜一樣.
一生都要學習去敬畏上帝.
誠心實意去事奉神.
我做了牧師幾十年.
我見到很多教會衰落.
我聽到很多基督徒說.
牧師我裡面覺得很枯乾.
裡面很乾旱.
我不是沒有東西.
但是不知為何.
時時都有一種不足.
不夠的感覺.
為何會這樣.
如果你問我的話.
我相信其中一個主要原因是.
我們是選擇隨從自己的心意.
尋求滿足自己覺得.
比信仰,比順從神更重要的事.
簡單說.
我們沒有將神放在我們生命當中的首位.
所以求聖靈幫助我們.
在這個時間我們來敬拜上帝.
希望你的心向上帝打開.
求聖靈幫助我們看到.
到底我們家裡所把持的.
所寶貝的核心價值到底是甚麼.
上帝在我們家裡到底扮演了甚麼角色.
我們有多敬畏神.
有多信服神.
沒錯,每個人都有自己的工作.
所以剛才我們說.
若雙雅有工作.
你記得嗎,若雙雅有工作是甚麼.
一族之長,民族英雄.
做很多事的.
我們有自己的工作.
我們有自己的生活處境.
但作為一個信耶穌的人.
我們有沒有將我們的工作.
我們生活的處境.

$^{521}$包括我們家裡的生活.
看作為我們事奉神.
讓神得到榮耀的地方.
我們能否讓我們的生活.
每一個環節.
無論是家裡以內或家裡以外.
都可以讓人看到我們信神的人.
我們有見證.
我們的信仰讓我們在生活上.
讓人看到的時候.
上帝得到榮耀.
八零年代的時候.
有一套很出名的基督教電影.
叫做《Chariot of Fire》.
有沒有看過.
《烈火戰車》有沒有聽過.
裡面說到英國兩個主人翁.
一個是英國的運動員.
跑得很快的.
叫做Eric Liddell.
另外一個叫做Abraham,是猶太人.
一個為了要排除猶太人的恥辱.
被人歧視.
拼命去跑.
但這個Eric Liddell他出名的地方是甚麼.
他參加比賽那天是星期天.
他說我不去參加比賽.
因為我是基督徒.
他有一句很出名的說話.
是甚麼.
我看看鏡頭在說甚麼.
他怎樣說.
他說.
I believe God make me for a purpose.
他說我相信神做我一個目的.
But he also make me fast.
那個神做得我跑得很快.
所以當我去跑的時候.
When I run.
I feel God's pressure.

$^{561}$他是一個運動員.
但他做運動員是為了榮耀上帝.
無論神做我們甚麼也好.
大家有不同的專業.
但是在你的專業當中.
在你的工作裡面.
特別是在你回到家的時間.
有沒有辦法讓我們的同事,家人,朋友看到.
We do this for the glory of God.
我們是為了神的榮耀.
因為神做我們就是這樣.
神wire我,神裝備我.
就是要做這些工作.
而神將我放在家庭裡面.
我在裡面.
我們的Daddy,Mummy,公公,爺媽,婆婆.
我們在裡面.
要建立一個能夠讓神榮耀的家庭文化.
以致到我們的下一代.
我們的子女.
在裡面成長的時候.
不單單有機會認識神.
我們真的能夠成為神的來源.
我們可以怎樣建造神的家庭.
怎樣建立一個以基督為中心的家庭文化呢?.
首先我們要知道甚麼呢?.
經常說家庭文化是一個很特別的term.
裡面包含了甚麼?.
到底講了甚麼呢?.
聖經有些地方給我們參考的.
我找到一段.
不單單是這一段.
我覺得這一段也很有代表性.
歌羅西書第三章.
十二,十七節.
大家留心聽我讀.
所以你們既是神的選民.
聖潔蒙愛的人.
就要存憐憫.
即是甚麼呢?.

$^{601}$要同情心.
要恩慈.
即是要友善.
謙虛溫柔.
要有禮貌.
和忍耐的心.
這些是講到我們待人的態度.
十三節.
倘若這人與那人有憐憫.
包括他裡面也是.
有衝突有矛盾.
大家不和睦.
要怎樣呢?.
總要彼此包容.
即是甚麼意思呢?.
要自我克制.
包容要自我克制.
彼此饒恕.
我上個禮拜解釋過這一點.
不要時時想著要報仇.
不要時時想著要怎樣擺平.
不要時時想著要fair.
主怎樣饒恕了你們.
你們要怎樣去饒恕人.
十四節.
在這一切之外要存著愛心.
包括愛神和愛人的心.
因為愛心就是我們的動力.
愛心就是聯絡全得的.
又要叫基督的平安在你們心裡作主.
就好像中國人怎樣說呢?.
中國人說家和萬事興.
類似的意思.
不過.
我們追求和睦.
是因為神.
竟然願意藉著耶穌基督.
主動和我們和好.
你們也要為此.
你們也是.

$^{641}$你們也是為此蒙召.
讓我們可以和神和人和好.
歸為一體.
所以呢.
我們當以各樣的智慧.
把基督的道理.
豐豐富富地存在心裡.
用詩章頌詞.
靈歌彼此教導.
互相勸誡.
不知道我們家裡.
有沒有經常這樣做.
如果有.
我們就會心被恩感.
歌頌神.
最後,普羅提醒我們.
無論作什麼.
或說話或行事.
都要奉主耶穌基督的命.
第一次我們代表了耶穌.
並且被耶穌賦予權並能力.
去做事和說話.
我們在地上就是祂的代表.
我們應該為有這樣的機會去事奉神.
我不知道會不會讀得太快.
我估計大家回家之後.
今晚上網重溫我剛才.
怎樣去讀這段聖經給大家聽.
然後呢.
自己打開聖經.
讀幾次.
聽完我再讀多次之後.
上網聽,不是現在聽.
回家自己讀多幾次.
一面讀.
一面祈禱.
求聖靈幫助你.
將神的話存在心裡.
讓神的話變成你的信念.
變成你的價值觀.

$^{681}$讓聖經裡面的教導.
能夠在我們每天的日常生活當中.
自然地流露出來.
讓你的家人,朋友,同事.
很容易看到.
怎樣能夠建立一個基督徒的家庭文化呢.
第一樣就是界定.
到底我們所說的是什麼.
剛才提到.
歌羅西書三章.
給我們一些啟發.
這是第一步.
如果我們不知道培養什麼.
你就不知道怎樣做.
所謂界定家裡的基督徒文化.
這是怎樣建立家庭文化的.
這是第一步.
第二步.
還沒去到這麼快.
第二步就是家庭會議.
不知道你們有沒有開家庭會議.
大家一起吃飯會聊天.
家庭會議比較荒蕪.
比較少開.
不過有時候也要荒蕪一下.
一會兒我會說為什麼.
家庭會議做了什麼.
我建議你們做.
一起討論一下.
到底我們家人裡面.
有哪三樣東西是最重要的.
大家覺得我們家人裡面.
有哪三樣東西是最重要的.
然後知道這些東西之後.
有沒有可能我們可以自己去創造一下.
有些創意.
找一兩句正經的家庭格言.
我知道很多人家裡有雪櫃.
貼了很多東西.
很多的金句,格言都是別人的.

$^{721}$可不可以自己覺得.
這個是我們家人裡面.
我們覺得重要的.
譬如說.
我們的家不追求第一.
但求凡事都能榮耀神.
舉例而已.
譬如說.
我們的生活就是要事奉神.
很簡單一句.
大家認同的.
這個是我們的格言.
所以這些就是我們的指路明燈.
就是我們的北極星.
有什麼決定的時候.
這些會幫助我們.
第一件事是家庭會議不是經常開的.
通常是有重大的事.
我們才會開家庭會議.
不過如果我們能夠把握機會.
有時候不需要很重大的.
一家人坐在一起.
討論一下家裡一些比較重要.
一起決定的事.
有什麼帶來的後果呢.
你會在家裡每個人.
無論是老與戀.
覺得被尊重.
因為一起討論一起談.
他們都被尊重.
他們有一個聲音在家裡.
當他們能夠參與家庭會的時候.
他們更加會有歸屬感.
更加會尊重父母.
這是第二步.
第三步.
建立基督徒信仰的家庭傳統.
家庭儀式或者一些習慣.
譬如說.
不知道大家每天早上.

$^{761}$如果有小朋友的話會很忙.
又要早餐.
預備午餐.
又要送他上學.
如果兩個人一個送這裡一個送那裡.
你一頭煙.
可不可以培養一個習慣.
總之出門口之前.
去到門口大家就圍著在門口.
一起做一個禱告.
求神帶領.
保守我們無論是上學也好.
做工也好.
去玩也好.
C 一起祈禱.
求神保守帶領.
或者我們可以利用.
家人一起吃飯的時間.
好像我們剛才也提及過.
一開始分享當日.
我們最開心的.
或者是我們最具挑戰的一些事情.
聖經有些事情.
怎樣會提醒和教導我們去面對這些問題呢?.
或者我們在睡覺之前.
我舉例.
一起祈禱.
可能大家都有.
睡覺前和小朋友一起祈禱.
將一切憂慮煩惱交給上帝.
答應神.
不會陷路到日落.
可以嗎?.
試試吧.
第四步.
改善父母長輩在家中的榜樣.
榜樣是很重要的.
以後我們可能會再講也不一定.
但你知道嗎?.
我們的子女.

$^{801}$是不會心悅誠服地.
按照我們吩咐他們做的事情去做的.
但是他們會不自覺地.
按照我們經常做的事情.
跟著去做.
當小朋友的課外活動.
和教會的聚會有衝突的時候.
我講一些很敏感的問題.
父母的決定.
就是用行動告訴我們的子女.
信仰在家中的優先次序.
是排第幾.
我也經歷過這些事情.
有時候小朋友要去比賽.
或者有些特別的事情.
學校有活動要舉行.
那怎麼辦?.
跟教會崇拜有衝突.
我就帶他們去夜晚崇拜.
帶他們去西洋教會.
他們是聽英文的.
當時有Pastor Mac做College Pastor.
是星期一下午四點崇拜.
這樣最好.
有辦法的.
不過我們要想.
傳遞給我們下一代的訊息是什麼.
第五步.
家裡的設計和你怎樣用時間.
是一個訊息.
我們怎樣能夠讓我們的家.
無論是佈置也好.
大家生活在一起也好.
感覺受保護.
在裡面很安全.
這是安樂窩.
大家去想一想.
最後一步是什麼?.
經常把握機會.
分享一些個人在信仰上的經歷.

$^{841}$見證.
對四經的看法.
我覺得我們要習慣跟小朋友.
討論人生信仰的問題.
告訴他們.
上帝怎樣祝福我們.
幫助我們.
跟他們分享我們的見證.
讓他們知道.
信上帝是一回事.
所以剛才說了六件事.
我很快跟大家Recap.
現在可以出來了.
這六件事是什麼呢?.
第一件事是什麼呢?.
什麼是基督徒的家庭文化.
參考《哥羅世書》.
第二件事.
界定在家裡什麼是最重要的.
第三件事.
建立基督徒的家庭傳統.
第四件事.
改善家長在家裡的榜樣.
榜樣很重要.
第五件事.
讓自己的家成為一個安樂窩.
最後.
分享蒙福的見證和故事.
小朋友喜歡聽故事.
大人也喜歡.
多點講故事給他們聽.
我們一起祈禱.
因為你與我們同在.
你讓我們有機會去敬拜和事奉.
我們祈求你的靈與我們同在.
因為有時我們說要建立一個文化.
我們以為這是牧師的工作.
但在家裡是我們每個人都要做的事.
天父我們知道我們沒有用.
我們軟弱.

$^{881}$我們沒有能力.
但你既然這麼愛我們.
幫助我們不單單自己學習愛你.
我們也很希望.
我們的下一代.
我們的家人也愛你.
祈求你加給我們力量.
智慧.
怎樣開始建立基督徒的家庭文化.
我們可以從一小部分開始.
不需要一次就到位.
但起碼給我們一個開始.
以致我們能夠讓我們的家.
成為一個真正屬於你的家.
我們這樣的祈禱 奉耶穌基督的聖名.
\newpage



\section{約翰福音 21:15-17}
\label{sec:032049}
\textbf{陀飛輪:從睇化,到睇透,再到激活}~~~~秦潤祥牧師
\newline
\newline
連結: \href{http://www.church.com.hk/acms/content.asp?site=cdc\&op=show\&type=product\&code=032049\&layout=sermon}{\texttt{032049}} ~~~~ 語音日期: 2026-04-26
\newline
\newline
\hyperref[sec:032048]{\small{< < < PREV SERMON < < <}}
~
\hyperref[sec:index_chronic]{\small{[返順時目]}}
~
\hyperref[sec:032050]{\small{> > > NEXT SERMON > > >}}
\newline
\newline
約翰福音 21:15-17
\newline
\begin{longtable}{cl}
\hline
\hline
章節 & 經文 (和合本修訂版)\\
\hline
21:15 & \begin{tabularx}{0.7\textwidth}{X} 他們吃完了早飯，耶穌對西門‧彼得說：「約翰的兒子西門，你愛我比這些更深嗎？」彼得對他說：「主啊，是的，你知道我愛你。」耶穌說：「你餵養我的小羊。」 \end{tabularx} \\ \\ \relax
21:16 & \begin{tabularx}{0.7\textwidth}{X} 耶穌第二次又對他說：「約翰的兒子西門，你愛我嗎？」彼得對他說：「主啊，是的，你知道我愛你。」耶穌說：「你牧養我的羊。」 \end{tabularx} \\ \\ \relax
21:17 & \begin{tabularx}{0.7\textwidth}{X} 耶穌第三次對他說：「約翰的兒子西門，你愛我嗎？」彼得因為耶穌第三次對他說「你愛我嗎」，就憂愁，對耶穌說：「主啊，你無所不知，你知道我愛你。」耶穌說：「你餵養我的羊。 \end{tabularx} \\ \\
[1ex]
\hline
\hline
\end{longtable}
$^{1}$我們一起讀約翰福音21章14至17節.
1,2,3.
耶穌從死裡復活以後.
向門徒顯現.
這次是第三次.
他們喫完了早餐.
耶穌對西滿彼得說.
約翰的兒子西滿.
你愛我比這些更深嗎?.
彼得說.
主也是的.
你知道我愛你.
耶穌對他說.
你慰養我的小羊.
耶穌第二次又對他說.
約翰的兒子西滿.
你愛我嗎?.
彼得說.
主也是的.
你知道我愛你.
耶穌說.
你慰養我的羊.
第三次對他說.
約翰的兒子西滿.
你愛我嗎?.
彼得因為耶穌第三次對他說.
你愛我嗎?.
就憂愁對耶穌說.
主啊.
你是無所不知的.
你知道我愛你.
耶穌說.
你慰養我的羊.
大兄姐妹 我相信大家都不是第一次.
看這一節經文.
當我們看見.
彼得因為耶穌第三次對他說.
你愛我嗎?就憂愁的時候.
很多解經家或者我們看《聖經》的.
都直接去參考.

$^{41}$看回彼得三次不認主那段經文.
我相信彼得.
記起這件事的時候.
是充滿羞恥感.
和罪疚感.
以我所認識的彼得.
在福音書裡面.
他是一個大情大性的人.
當聖上善壽谷起的時候.
他就很自我.
眾人認你.
我都是不會不認你的.
甚至乎為你死都可以.
是嗎?.
另一方面.
他也是很自卑的.
記不記得當耶穌第一次呼召他的時候.
在船上.
耶穌打算給回.
一些佣金.
借了一艘船去港島.
你洗船出去的時候.
彼得他說什麼.
他說我們昨晚打了一整晚雨都沒有.
我是專家.
不過見是你我就出去試一下.
誰知道很多雨.
是嗎?.
船差點沉.
接著他見到耶穌的神聖的時候.
他就跪下.
他說主人你離開我.
因為什麼?.
我是罪人.
他覺得自己很不配.
弟兄姊妹.
我不知道你今天是否覺得自己很不配.
好像彼得這個大情大聖的人.
他經歷了很多的事.
我相信他那種自卑感和那種自傲感.

$^{81}$是源自他很多的原生家庭.
世界的價值觀.
有很多疑似的傷害也好.
很多不同的東西組合起來的.
當我看這段經文.
看回彼得的生平的時候.
我想起我自己之前.
神給我的一些責怒.
我發覺我自己也是相當自卑.
也是相當自大.
當信警的時候.
我還是很自大.
當我逆境的時候.
我會很自卑.
我記得我在畢業之後.
接著我就回中國做生意.
結果是失敗了.
當我回來美國之後.
那種羞恥感.
和那種罪疚感很強.
當時我是在聽一首歌.
有沒有人聽過《陀飛輪》這首歌.
我現在說《陀飛輪》好像很搞笑.
但當時我看《陀飛輪》這首歌的時候.
哇 真的每次都迷上我.
有沒有人會唱這首歌.
我不會唱 我會讀出來.
這首歌其實充滿人生的哲學.
充滿人生的智慧.
對於我來說.
這首歌在未信主之前.
是一種有秋天的感覺.
秋天的落葉.
你看到時日無多.
但是你卻懂得珍惜餘下的時間.
雖然之前發生的事情已經發生了.
但這也算是一種悽美的智慧.
我記得這首歌有兩個年紀.
第一個就是過去十八歲.
沒大標不過有時間.

$^{121}$記不記得.
十八歲那種燦爛的天真浪漫.
我十八歲的時候.
我敢說是人見人愛.
車見車長.
因為很單純.
挺可愛的.
做什麼都很順利.
做兼職每個人都找你.
好像很順利.
你慢慢去到二十七歲的時候.
踏入這個社會的人開始真靈.
開始你想賺錢.
開始喜歡賺錢.
但其實你是做什麼呢.
真的不知道.
我還狂到回到中國去創業.
覺得自己一定行的.
誰知道後來是不行的時候.
當中我做了很多很蠢的事.
入了這個世界的時候.
有時候我都不知道是在應酬.
上班.
還是在玩.
當我什麼都沒有的時候.
我回來的那種罪疚感.
那種羞恥感.
是很濃烈的.
看回陶飛倫這首歌.
後來我說都看化了.
我記得這首歌最後的那一段.
昂貴在這刻.
我事而這刻我覺悟了.
人生的覺悟都算是一種看化.
但弟兄姊妹.
基督徒不只是看化.
我們要從看化到.
看透再被激活.
這個就是基督徒的生命.
看回那首歌.

$^{161}$我當時看到他說.
曾付出多少心跳來換取一堆堆的發票.
我當時不知道自己為什麼這麼喜歡賺錢.
我家裡也不是很窮.
究竟我為了什麼賺錢.
後來我最後回來之後.
我想起來.
或者是我的原生家庭.
令我想證明我自己.
究竟我為什麼要這麼拼命去做一件事.
我不知道為什麼.
生命其實是為了什麼.
人值得命中減少幾秒多買一隻鞭.
誰知道那隻鞭是什麼牌子.
Brolex.
應該是努力士.
就是你賺了很多錢.
你買努力士那又怎麼樣呢.
弟兄姊妹.
你今天有些什麼呢.
或者你有一層樓.
一千萬.
十萬.
八萬.
so what.
是不是.
秒速抓得緊了.
你看回頭的時候.
皮膚就.
已經偷偷鬆了.
我真的回來的時候.
慢慢發現已經不同了.
已經不是十八歲的時候.
做什麼都可以了.
現在出去不是人見人愛那種了.
變了變了老了.
是不是.
那些小朋友見到你.
那時候叫什麼.
以前叫哥哥.

$^{201}$現在叫什麼.
叔叔了.
是不是.
真的完全不同了.
還剩下多少心跳呢.
當時我就聽見歌詞的時候.
我就很有感觸.
如下的一生我怎麼做呢.
我究竟應該怎麼看人生呢.
我應該怎麼去好好利用我的餘生.
去覺得更加有意義呢.
當我見到人面跟水晶標面對照的時候.
我更加有感觸.
原來一隻rolex.
標面裡面的玻璃.
它因為考慮天宮的製造.
你看到裡面的齒輪在動的時候.
你好像很掌握到時間.
但是對於時間的流逝.
你完全掌握不到.
那隻標很貴.
但是那又怎麼樣呢.
弟兄姊妹.
我們不要停留在看畫.
聽一首歌.
或者看看人生的格言.
是可以幫助到我們去看畫.
覺得是覺悟了.
但是這種都是因為時間的沖洗.
令我們無奈地妥協.
弟兄姊妹.
我們要從看畫到看頭.
我們需要恩典的依次.
我們靠的不是自己.
我們靠的是什麼呢.
復活.
耶穌基督的復活.
當我們看剛纔那段經文的時候.
經文寫得很清楚.
彼得因為耶穌第三次對他說.

$^{241}$你愛我媽就憂愁.
作者寫得很清楚.
是彼得覺得的.
耶穌不是想他這麼覺得.
耶穌想怎樣呢.
復活之後向他第三次顯現.
他不是想他看畫.
彼得那時候他用人生的格言.
人生那一套.
世界那一套來看他的背叛.
令他很有羞恥感.
罪疚感.
令他慢慢想離開.
他回去打魚.
但是耶穌基督親自回到他的身邊.
在他身邊顯然告訴他.
你靠的不是人生裡面的智慧.
你靠的是我的復活.
弟兄姊妹.
耶穌用一個早餐.
早飯.
來介紹兩個場景.
來去把彼得醫治.
希望今天這兩個場景.
都是醫治我們的場景.
兩個場景.
一條生魚.
一條熟魚.
行不行.
早餐喫什麼.
喫魚.
要喫魚就當然要打魚.
我們一起看看.
剛纔我們看那段經文之前的上文.
發生了什麼事.
他說.
西滿彼得對他們說.
我打魚去.
他們說.
我們也和你同去.

$^{281}$他們跳出去上了船.
那一夜並沒有打著了什麼.
耶穌就對他們說.
小子.
就是耶穌突然間出現了.
你們有喫的沒有.
他們回答說.
沒有.
耶穌說.
你們把網殺在船的右邊.
就逼得著.
他們便殺下網去.
竟拉不上來.
因為魚甚多.
他們上了岸看見那裡有炭火.
上面有魚.
又有餅.
西滿彼得就去把網拉到岸上.
那網滿了大魚.
共153條.
魚雖這樣多.
網卻沒有破.
弟兄姊妹.
當你看到那麼多條魚的時候.
你在彼得會想起什麼事?.
應該就是路加福音第五章.
當耶穌呼召他的時候.
耶穌用他的船來講道.
之後他說.
你把船開出去.
我給你一點賠償.
接著彼得說.
我打了一整晚魚.
都沒有魚.
我是一個漁夫.
我是一個專家.
但是因為你這樣說.
我就出去了.
結果怎麼樣?.
打了很多魚.

$^{321}$接著他跪在那裡說.
主啊 離開我.
我是個罪人.
耶穌怎麼對他說?.
耶穌說.
不要怕.
從今以後.
你要得人魚得魚.
這是耶穌對彼得的呼召.
一個新的身份.
結果當耶穌上了十字架.
死亡之後復活.
他因為羞恥感.
因為罪疚感.
因為很多很多的原因.
他離開了.
他走回去打魚.
你見不見到.
當彼得跳下水走去耶穌的時候.
他望見的倒影.
是多麼可笑.
如果你看回這首歌的MTV.
我開頭不是很明白.
我不知道你明白不明白.
他有個穿著兔子的衣服.
那個costume.
有個穿著太空人的衣服.
如果你沒聽那首歌.
你回去聽一下 挺好聽的.
去到最尾那裡.
你會見到有些影子.
那個人就脫了頭盔.
在那裡透氣.
他穿著那套costume.
是很辛苦.
出汗了.
弟兄姊妹.
這個不就是我們的倒影.
我們是不能被釋放.
那個羞恥感.

$^{361}$困了我們在costume裡面.
彼得竟然.
本來有個新的身份.
他跳了出來.
因為羞恥感和罪疚感.
他跳回去這個costume.
以至他見到自己的倒影.
是多麼可笑.
弟兄姊妹.
我不知道今天你是什麼情況.
但是如果你有羞恥感.
或者有罪疚感.
你有很多過去發生的傷害也好.
原生家庭的傷害也好.
弟兄姊妹.
神很想你將這個costume拿走.
因為那個不是真正的你.
因為你成為基督徒之後.
神已經給你一個新的身份.
弟兄姊妹.
記不記得這兩個字.
炭火.
知道為什麼要有炭火嗎.
炭火這個字在約翰福音.
甚至整個新約出現了兩次.
所以這兩次非常之緊要.
耶穌基督告訴彼得.
你的倒影多麼可笑.
你已經走回去打雨.
我要你回去.
你認為你羞恥的場景裡面.
去醫治你.
炭火.
你想起什麼.
我就想起約翰福音十八章.
我們今天是讀二十一章.
十八章裡.
那開門的侍女對彼得說.
你不也是這人的門徒嗎.
你都不是耶穌的門徒嗎.

$^{401}$彼得說.
我不是.
僕人和差役.
因為天難就身了什麼.
炭火站在那裡.
敲火.
彼得也同他們站著.
敲火三次.
不認主.
當彼得見到炭火的時候.
他想起的是他的罪疚感.
羞恥感.
但對於耶穌來說.
祂不是要他徘徊在這個罪疚感.
羞恥感當中.
而是在炭火上.
有餅也有魚.
耶穌要藉著這個場景.
來醫治彼得.
耶穌不是要用這個場景.
來令他覺得不安.
弟兄姊妹.
你見不見到.
今天真的看不到.
見不見到炭火.
有很多時候.
我們將我們定格在過去的炭火裡面.
過去我們的.
你做得不好的事也好.
你受傷害也好.
你的童年陰影也好.
這些全部是定格在過去.
人就是這樣的.
經常把自己放在過去成為一個終點.
你經常回望過去.
你向前行.
都是過去所有羞恥感的一個投影.
這就是活死人.
有很多人不明白這件事.
信了主.

$^{441}$都好像沒有盼望.
經常都好像很不喜樂.
就是因為你經常記住過去的羞恥感.
罪疚感.
耶穌不是想我們停留在過去.
耶穌想我們.
是望著天國永恆的盼望.
當你這個觀念一轉變的時候.
從回望過去成為一個終點.
以至你望見未來.
永生盼望的時候.
你每一個當下.
都是有天國的剪影.
我希望今天我們每一個當下.
都是有天國的剪影.
我們面對我們自己的倒影的時候.
我們會見到真正的自己.
這也是耶穌基督想見到我們的模樣.
Amen?.
如果我們徘徊在羞恥感和罪疚感的時候.
我們得不到醫治.
我們不知道聖經其實是想鼓勵我們的時候.
羞恥感就好像一個拉扯力.
將我們拉到兩邊.
第一個可能性就是你離開神.
好像彼得那樣.
你離開我 我是個罪人.
我不配 我回不了頭.
我不知道你有沒有試過這樣的場景.
如果有的時候 你聽見這個信息的時候.
神想你回來.
另一方面有第二個可能性就是.
他將你拉到一個場景.
就是被宗教利用.
你成為了被人控制的一個人.
用你的羞恥感 用你的罪疚感來控制你.
有沒有聽過?.
邪教很喜歡的.
甚至有很多所謂的政教.
都是在做這件事的.

$^{481}$弟兄姊妹.
聖經我告訴你.
是用來釋放你的.
聖經是使你得自由的.
所以弟兄姊妹.
你來到這間教會.
我們是先走B型的.
我們要被醫治之後.
我們纔可以做.
我們不要用我們的羞恥感.
罪疚感.
我奉獻了 我要補償.
我做了 因為我做得不好.
耶穌不是要這樣的門徒.
耶穌要是被醫治之後.
將他愛抹去的門徒.
弟兄姊妹.
我記得我有一次去夏威夷.
我去了一個摩門教參觀.
進去的時候.
有些很年輕的年輕人來賣.
我進去聽聽.
裡面挺漂亮的.
他跟我說.
他說英文.
他說.
你知不知道.
如果我們做了十一奉獻.
如果我們真的禁肉 禁食.
做了所有要求的東西之後.
我們就可以穿著白色的衣服.
在聖殿那裡一起去敬拜.
我見到他的時候.
他的樣子都不太確定.
因為他都覺得自己不配.
他覺得自己不行.
那時候我覺得他真的很慘.
那你不就是穿著兔子的服裝之後.
你再穿白色衣服.
那不是更辛苦嗎.

$^{521}$那你順利做什麼呢.
弟兄姊妹.
我希望我們都可以被醫治.
如果你過去有很多的罪疚感.
羞恥感.
我希望藉著這個信息.
很親切地再三地跟你說.
你是已經被醫治了.
耶穌基督想你專注在哪裡呢.
祂第三次顯現.
祂的復活當中.
你看回經文.
耶穌基督給很多提示.
又給炭火.
又煮東西給他.
和他怎麼問他.
他說約翰的兒子西門.
說了多少次.
三次.
有沒有想過.
為什麼他說約翰的兒子西門.
為什麼不說彼得呢.
知不知道約翰的兒子西門.
這個名字是什麼時候有的.
是耶穌呼召他之前.
耶穌之後說你叫彼得.
是不是.
耶穌基督藉著這個名字.
告訴彼得.
你已經回到以前的服裝.
我看到的不是彼得.
是一隻兔子.
OK.
你要醒醒了.
OK.
我不是要你憂愁.
我是要你重新學習.
你愛我比姐姐更深.
你回去打魚.
你喜歡魚多一點.

$^{561}$還是你喜歡我多一點.
弟兄姊妹.
今天耶穌基督.
和我們講同一個邏輯.
祂的要求很簡單.
祂不是要我們先奉獻.
何必奉獻.
又禁肉禁食呢.
你才賺得到這個救恩.
祂不是.
祂只是說.
你愛我就可以了.
弟兄姊妹.
你今天愛不愛神.
如果你是經歷過神的請求的時候.
我們今天早上讀聖經也是.
神在得時不得時都是要我們同在的.
你經歷過祂請求的時候.
你是很應該去學習.
怎樣去愛神.
以至我們可以愛人.
和憂愁絕對是無關的.
弟兄姊妹.
我們是需要耶穌基督十誡的救恩.
祂說我們.
釘身受死.
讓我們的罪得而赦免.
祂的復活.
讓我們可以重新活多一次.
弟兄姊妹.
你信不信你可以重新活多一次.
你信的話.
你可不可以給我一個Amen.
Amen.
我們屢次都犯錯.
好像剛纔那首歌.
我要將身體獻作活祭.
很難的.
但不要因為你犯錯.
有那個罪咎.

$^{601}$你不回到神的身邊.
神是想你很積極地看著天國的盼望.
以至你今天有天國的剪影.
有很多時候這些負面的東西.
成為了我們的一個諾印.
就好像一個紋身一樣.
揮之不去.
讓這個紋身成為你基督徒人生的起點.
我以前聽過一個故事.
真人真事.
以前還是在公元后一千年左右.
有一個人偷羊.
偷羊被人抓了.
他是壞人.
你知不知道那些人怎麼罰他呢?.
他就在他的額頭那裡.
彈了兩個字.
叫做Shit Thief.
即是偷羊的人.
他要他這一世都有一個恥辱感.
他以後不敢偷羊了.
這個就是人的邏輯.
但這個偷羊的人.
後來他信了主.
慢慢他可以活出主的愛之後.
整個town見到他的時候.
都見到他就是神的門徒了.
他就是耶穌基督的門徒的時候.
當那些小朋友問.
為什麼他這麼好人.
額頭有個St兩個字呢?.
那些人說.
因為他是聖人.
聖而已.
你們明白嗎?大英姐妹.
讓這個疤痕.
這個傷口成為人生的起點.
就可以成聖了.
就可以在基督裡面有份了.
大英姐妹今天的你.

$^{641}$還有沒有羞恥感和罪疚感呢?.
如果你有的話.
我奉主耶穌基督的名.
我告訴你.
Everybody在這裡.
全部都被赦免.
全部都被醫治.
以至我們可以繼續走上去.
只要我們被醫治.
靠著神的耶穌基督的復活.
我們就可以真的用我們的愛心來回應.
我們以後就不是靠自己.
不是靠人間的邏輯.
不是靠聽諾託菲倫.
我們是靠聖靈的帶領.
很難免看《約翰福音》21章.
你不看過阿加皮和弗爾尼奧.
就是說愛是有兩個希臘文.
但是其實這兩個希臘文.
我就不會用太多篇幅去解釋.
因為這兩個希臘文在《約翰福音》裡面是一樣的.
它的意思都是愛.
就這麼簡單.
但是為什麼作者要用兩個不同的愛.
來表達給我們看呢?.
第一個就是阿加皮.
第二個就是弗爾尼奧.
就是兩個不同的字.
很簡單地告訴你.
你當第一個字是.
耶穌說:你喜歡我.
比這些更甚嗎?.
接著Peter怎麼說?.
主啊 你知道我是歡喜你的.
那為什麼他不用同一個字呢?.
接著之後怎麼樣?.
耶穌再問他.
你是不是喜歡我啊?.
接著比特怎麼說?.
你知道我是歡喜你的.

$^{681}$行不行?.
那看第三次了.
耶穌基督用回祂的字.
祂講下來.
那你歡不歡喜我?.
你看他怎麼說.
你知道我是歡喜你的.
當時他不明白.
他走去憂愁那邊.
但我相信比特之後回想.
他明白耶穌基督要他再一次復活.
耶穌基督在天上降下來成為人.
去愛我們.
用回我們明白理解到的愛.
去愛我們的時候.
讓我們慢慢學習.
如何去愛祂.
如果不是耶穌基督上十誡.
如果不是耶穌基督復活.
聖靈是不會來到我們當中.
讓我們明白領略到.
神對我們的厚愛.
我希望聖靈與我們同行的時候.
我們學習如何去愛神.
神歡喜我們歡喜祂.
神也喜歡我們喜歡祂.
我們學習神如何去愛我們.
我們如何去愛神.
如果我們懂得這樣回應的時候.
你就會明白餵養和牧養是說什麼.
對於我的冥想.
當比特聽見耶穌基督說.
你餵養我們的時候.
這件事是非常立體的.
知不知道為什麼.
因為剛剛耶穌才喂完他.
喂他的時候他醫治了他.
當他喂他的時候.
他會想起他之前的羞恥感.
對耶穌不住.

$^{721}$但耶穌說既往不咎我餵你.
你現在要喂那些小羊.
那就是說當比特出去傳教.
去本地教會的時候.
當弟兄姊妹犯錯的時候.
他應該怎麼樣.
他應該做同樣的事.
給他機會.
醫治他.
餵養他.
兩個餵養夾著什麼.
牧養.
應該怎麼去牧養.
怎麼去牧養呢.
透過聖經的教導.
你想想.
如果比特是大弟子的話.
沒有人這麼三不四時說他三次庇佑主.
一定是他自己經常說的.
他要告訴別人他還是不行.
我們要跟著耶穌基督去走.
其實我有一個小小的陰影.
我不知道是我不走運.
還是香港的教育我接受不了.
我小時候不知道我讀了什麼.
去補習社.
他將一些很簡單的東西搞得很複雜.
令我沒興趣去讀.
慢慢我就不讀了.
又被人打手板.
我就更加不讀了.
我是靠小聰明大學畢業的.
總算畢業了.
你不讀書也畢業.
很厲害的.
但其實我有陰影.
如果當年我肯讀書的話.
現在真的是不得了.
一直埋在我心裡.
真的是一個陰影.

$^{761}$直至讀神學醫治了我.
他說你就用你自己的方法重新讀一次.
我讓你復活.
不要將這個成為你的疤痕.
結果我真的用我的方法去讀.
原來有些東西唱首歌就會了.
英文用不著學.
你去聽一下布努馬斯.
聽一下那些歌.
你都會了.
原來是可以很輕鬆.
可以很有趣味.
就可以了.
我不是讀得很厲害.
但我總算有個交代了.
弟兄姊妹.
神一呼召我去牧養的時候.
我想將這些方法給你們.
不要當看聖經好像讀書一樣.
看聖經是去領略神的愛.
是可以很輕鬆.
很自然.
就可以明白神想和我們說些什麼.
我們早上有讀聖經的時候.
如果你有空的話8點搭9.
ASB那個燈好嗎.
OK 我們現在在讀.
詩篇和什麼.
《厥書雅記》.
不知道為什麼我忘記了.
總之就是在讀.
那你過來吧.
希望我們可以一起被牧養.
最後.
有件事我們要非常小心.
小心什麼呢.
經文講得很清楚.
餵養誰的羊.
我的羊.
餵養誰的羊.

$^{801}$我的羊.
很多時候我們被醫治之後.
我們想大發熱心為神去事奉的時候.
我們就會take it personal.
我以前犯過這個錯.
我怎麼纔可以奉獻我所有的東西來回應神呢.
當這個是我的不就行了.
這個是大錯特錯的.
弟兄姊妹.
這個是會讓你用理智的方法去管神的羊.
今天我們所有的氣息.
今天我們所有的奉獻.
今天我們所有的ministry.
全部都不是我們的.
弟兄姊妹.
我們這間教會要是健康的.
這間教會不是我的.
不是pastor Jackie的.
不是你的.
是誰的.
是神的.
所以我們要有deacons.
我們遲些會有membership meeting.
講deacons的.
來吧.
好不好.
我們一起去.
藉著耶穌基督的名去管理教會.
好不好.
所有的東西都是誰的.
是神的.
不然你很容易驕傲.
如果你覺得是自己的方法.
去成就這個ministry.
去成就全福音的時候.
當你成功的時候.
你會驕傲.
當你不成功的時候.
你會發脾氣.
有沒有見過人試服的時候發脾氣.

$^{841}$我見過.
知道是誰嗎.
就是我.
很害怕吧.
我們要慢慢學習.
如何去成為一個有良善.
有忠心的僕人.
不要下來做一個陀飛輪.
不要穿這個costume.
我們不要像撒旦一樣.
耶穌說.
「塵看見撒旦從天上墜落.
像閃電一樣」.
弟兄姊妹.
希望今天這個信息.
可以讓你成為一個有良善.
有忠心的僕人.
從體化.
無奈的妥協.
你有很多羞恥感.
揮之不去的時候.
希望你可以看透.
被醫治.
藉著復活的能力.
我們一起.
用愛心來回應.
牧養耶穌基督的羊.
成為見證.
靠著聖靈.
我們as one body.
一起走出去好不好.
Amen.
我們一起祈禱.
鼓勵我們.
讓我們可以被醫治的時候.
主人.
讓我們過去的羞恥感.
罪疚感.
藉著十字架一筆勾銷.
我們前途走的.

$^{881}$是充滿天國的剪影的路途.
讓我們一群人一起很健康.
一起藉著你的靈.
和真理去走的時候.
我們每一個moment.
都是敬拜著你自己.
讚美你.
土貴奉主.
耶穌救命祈求.
Amen.
\newpage



\section{出埃及記 32:13}
\label{sec:032050}
\textbf{聖經人聖經事:你曾指著自己起誓}~~~~黃克勤牧師博士
\newline
\newline
連結: \href{http://www.church.com.hk/acms/content.asp?site=cdc\&op=show\&type=product\&code=032050\&layout=sermon}{\texttt{032050}} ~~~~ 語音日期: 2026-04-22
\newline
\newline
\hyperref[sec:032049]{\small{< < < PREV SERMON < < <}}
~
\hyperref[sec:index_chronic]{\small{[返順時目]}}
~
\hyperref[sec:032051]{\small{> > > NEXT SERMON > > >}}
\newline
\newline
出埃及記 32:13
\newline
\begin{longtable}{cl}
\hline
\hline
章節 & 經文 (和合本修訂版)\\
\hline
32:13 & \begin{tabularx}{0.7\textwidth}{X} 求你記念你的僕人亞伯拉罕、以撒、以色列。你曾向他們指著自己起誓說：『我必使你們的後裔像天上的星那樣多，並且我要將所應許的這全地賜給你們的後裔，讓他們永遠承受為業。』」 \end{tabularx} \\ \\
[1ex]
\hline
\hline
\end{longtable}
$^{1}$又來讀聖經了.
眨眼就一天.
第一天節目來平安.
生命說書教網絡.
又和大家聖經人聖經時.
看看上帝是怎樣的.
今天用的「你曾指著自己起誓」.
這句說話是摩西說的.
摩西和誰說?.
和主耶和華說.
「你曾指著自己起誓」.
嘩 即是甚麼意思?.
「臣啊 你不要忘記」.
嘩 這麼大膽說這些話.
跟神說「你發過誓不要忘記」.
這個很多時候.
我們在禱告裡面.
這一類型.
靠著上帝的應許來做禱告.
很多時候神是特別聽的.
因為祂說「是呀 是我應許來的」.
「好 我會做」.
「你曾指著自己起誓」.
「一夜跳到初一及二三十二章」.
就是神叫摩西和約書亞.
上了西乃山.
然後把祂手上寫的法板.
交給摩西解釋了.
教了他很多東西.
四十天.
嘩 這麼久.
那些不回答了.
那摩西吃甚麼?.
有沒有水喝?.
怎樣睡覺?.
有沒有帳篷?.
不知道.
如果神要令他那四十天沒有睡都可以.
神都會令到他可以的.
因為神是神.

$^{41}$那些不提了.
上了山四十天.
就把那塊法板.
兩面的.
兩面都是字.
於是推測是一塊板.
寫了十誡在裡面.
可能是不是一面五誡呢?.
一面五誡.
一面五言.
不知道.
那塊板現在也看不到了.
上面又教了他很多東西.
告訴他很多東西.
山下就出事了.
山下四十天不見了摩西.
摩西就是帶領他們的屬靈領袖.
不見了約書亞.
約書亞就是.
幫他們打仗.
就是軍事的將領.
又不在了.
剩下誰在呢?.
亞倫.
亞倫也是祭司.
就是帶領他們宗教上面的.
那幫以色列人也知道故事了.
他們就找亞倫.
亞倫.
那個你看聖經.
那個摩西在山上這麼久不下來.
都不知道他發生什麼事了.
不知道會不會被獅子叼了.
不知道會不會掉下山坑死了.
這個我講的.
他上了山這麼久.
不知道他遭遇什麼事了.
現在我們怎麼辦?.
怎麼辦?繼續等嗎?.
他又不在了.

$^{81}$以前神透過他講話的.
也有透過亞倫講話.
不過神透過他講話多.
現在怎麼辦?.
其實他們這樣問是問.
耶和華在哪裡?.
不見了耶和華代言人.
耶和華在哪裡?.
這裡就沒辦法了.
現在回看聖經.
我們會說是亞倫的軟弱.
他就說不如這樣吧.
不如這樣吧.
可能不是亞倫說的.
可能是亞倫最後也不敢.
既然大家都想這樣就這樣吧.
這樣也不一定.
妥協了.
很有可能是那些以色列眾人.
甚至可能有些長老.
不如我們做一個耶和華出來吧.
怎樣做一個耶和華出來?.
他們就想起他們腦裡.
你知道人是不能夠想到一些.
創作到一些.
我們的創作力是超級這麼低.
這麼有限.
就是我們創作不到一些我們沒見過的東西.
一定要有經驗.
或者其他東西我們見過.
經驗過才能創作到.
才能做到創作.
他們想起如果是神.
馬上想起四百多年都在埃及.
神是怎樣的?.
埃及最大的神就是金牛.
真是一隻全用金做的牛.
他們想起不如我們做一隻金牛吧.
埃及的神我們又拜.
不,我們不是拜埃及的神.

$^{121}$其實祂是象徵耶和華而已.
你看到十言裡面是說不要雕刻.
不要製造.
雕刻的意思是製造.
不要製造,不要跪拜,不要事奉.
他們就想不如做一隻金牛來代替象徵耶和華.
不是代替,象徵耶和華.
有認真的成份.
要認真嗎?.
於是亞倫說既然這樣,那好.
不如大家回去問你太太.
還有很多東西.
那些金原來在太太那裡.
聖經說回去問你太太.
各人問你太太拿她的金耳環.
拿她的金環,金戒指.
金鏈,金手鐲拿來.
我們用了它做一隻金牛.
做好了明天我們要向耶和華守節.
向耶和華守節.
他們做一隻金牛是象徵耶和華.
而且大家認真的.
你試試回家問你太太.
拿些金戒指,金耳環回來看她行不行.
她真的每個人都拿出來.
所以他們都認真的.
只不過用了金牛象徵耶和華.
這裡錯得很嚴重.
神都很不歡喜.
已經和他們以前說過不要了.
而十言是重新再說.
有些人認為摩西一下來就生氣了.
整塊法板扯到那些人那裡.
扔下去了.
那就真的不好了.
法例在山上一定下來就執行.
其實不是的.
之前已經說過不要拜那些神.
而且十災那些很多都是象徵埃及地的神.
我是和他們有所分別的.

$^{161}$不痛的不要拜那些.
好了.
耶和華就吩咐摩西說.
行了四十天了下去吧.
下去拜.
因為你的百姓就是你從埃及地領出來的.
第一點要說的就是已經拜壞了.
他們出來跟你走看過十災.
殺長子之災.
然後出來跟你走這麼遠路.
過盧偉海.
紅海.
過盧偉海.
一直走到現在.
嘩.
經歷到埃及人在水裡死了.
然後屍體衝上岸邊.
嘩.
很厲害.
然後雲柱火柱一起.
沒東西吃了餓死了.
又有鵪鶉又有嗎哪.
然後又有好口乾.
於是摩西擊打盤石.
又有水飲.
二百萬人夠飲.
嘩.
厲害厲害.
於是你會知道.
其實這些人全部都認識耶和華.
他們甚至乎.
死啦.
沒有一樣具體的東西跟我們真的害怕.
我要靠一個祭司.
靠一個中間人.
現在摩西不知道發生什麼事了.
那怎麼辦.
我們做一個耶和華.
我們做一個象徵他.
誠心做就會是.

$^{201}$這個就是很多時候.
我們敗壞的地方.
千萬不要想著.
嘩.
我多厲害.
多多人生經歷.
我做靈修部那時你都未信主.
我在教會事奉的時候說.
人少少我什麼都做.
那時靈性都不知多好.
你現在這樣新入教.
不行的.
不聽我講.
其實自己原來已經變得靈性很低落.
都不知.
因為覺得自己很厲害.
其實記得有時候.
我們不知不覺之間.
已經敗壞了.
由埃及出來.
經過這麼多東西.
經過這麼多上帝的神蹟.
不過可惜靈性已經敗壞了.
要反省.
很多時候我們都會這樣.
自己病了.
靈性病了都不知.
他們快快偏離我所吩咐的道.
吩咐了他們.
自己鑄了一隻牛獨.
向他下拜獻祭說.
以色列.
謝是領你出埃及地的神.
他們沒有拜其他的.
他們說這隻是耶和華.
這隻是耶和華.
他是做了一隻真金的牛.
是付出代價而且認真地獻祭敬拜.
不過走錯路敗壞了.
於是耶和華對摩西說.

$^{241}$我漢者百姓真是硬著頸汗的百姓.
你奢猶在我.
我要向他們發烈怒.
將他們滅絕.
使你的後裔成為大國.
於是神跟摩西說.
你不要攔著我.
這班敗壞了.
看過我過紅海,十災,殺長子之災,逾越節.
然後過紅海.
然後吃過鵪鶉嗎哪.
見過雲柱,火柱.
喝過盆石出來的水.
這班人竟然這樣不聽話.
所以他們叫做硬著頸汗的百姓.
硬著頸汗的百姓在神來看是抵死的.
是我抵罰的.
所以你不要再攔著我.
我要向他們發烈怒,滅絕.
沒有我的百姓怎麼辦?.
我會叫你摩西的後裔成為大國.
你的後裔也是亞伯拉罕的後裔.
所以沒問題,可以的.
摩西便懇求耶和華.
他的神說.
耶和華.
你為什麼向你的百姓.
向你在埃及地帶出來的人發露呢?.
你剛剛才第一句跟我說.
除了我之外.
我就是在埃及帶你出來的那個.
除了我之外你不需要有別的神.
對呀,就是了.
現在他弄了一隻牛.
製造了這個牛杖.
又敬拜,又事奉,又跪拜.
所以我起火.
你說.
他們這班傻子,你為什麼要發露呢?.
其實他們是你用大力和大能的手.

$^{281}$從埃及地領出來的.
主人你記得,你做了很多事情.
摩西的談判術真的高.
他跟神說.
你為什麼要發露呢?.
不要生氣.
其實你是施展你的大能.
辛苦帶他們出來的.
下了很多功夫的.
而且你把他們全部殺了.
你就會在埃及人面前失威被他們說了.
是的,說什麼?.
耶和華領他們出去是要降禍於他們.
是要把他們殺死在山裡面.
然後把他們在地上除滅.
還有沒有人說你耶和華是慈愛的神?.
沒有了,埃及人都會笑的.
你就用水淹死我們.
然後你去到山那裡,你殺死他們.
所以求你轉移不法你的烈怒.
你後悔不降禍於你的百姓.
摩西的祈禱真的厲害.
如果這個時間聽到這樣說.
殺了他們,然後我那族人做了你的選民.
好啊,真是活該的.
那我的族人多好啊,我教得不錯啊.
自己拿了功夫.
沒有啊,不要啊.
你很辛苦帶他們出來.
也不要讓埃及人笑.
還有就是求你紀念你僕人亞伯拉罕爾薩·以色列.
這就是以色列人,不要叫雅各.
叫以色列,你曾指著自己起誓的.
我要讓他們的後裔好像升了那麼多.
在我應許之地永遠承受違約.
你說過的,不能不算的.
你起過誓的.
於是耶和華後悔不把所說的災禍降禍於他的百姓.
經常出現這句詞.
後悔的,神為何這麼差的.

$^{321}$早就知道了,你現在才知道.
後悔,後悔呢.
Lakam,Lakam是遺憾.
遺憾和後悔是不同的.
後悔就是,要做過,遺憾就是.
他們這樣做,令我很遺憾.
即是很不開心.
而耶和華的神又很激動.
雖然很不開心,但我回心轉意不再降災,降禍.
這就是我們的神,祂真是愛祂的百姓.
一下叫這隻金牛起火,你快點下去壞祂.
挺有趣的.
看這聖經,我們看到原來我們很不可靠.
已經敗壞了.
第二,我們經常要硬著頸汗.
我們是一個硬著頸汗的百姓.
不要啊,我們很感恩.
神是記得他的應許.
好在這一點非常有效.
一講,神就說好.
好,我紀念我和亞伯拉罕以撒以色列說過的話.
於是祂就收回祂的災禍.
繼續帶以色列走了四十年,非常感恩.
讚美我們的主,和我們同心祈禱.
多謝主,你是一位好僕人摩西.
祂很能夠摸著你的心意,在你面前有很好的禱告.
祂亦都深愛這群帶出來.
又硬著頸汗,又已經敗壞了的百姓.
我們不會取笑那些以色列人,因為我們都是這樣.
我們要警醒自己,因為我們都是那麼軟弱,那麼脆弱.
小小事,不見了一個領袖,我們就亂了.
我們想起我們以前擺了什麼,不如我們擺回吧.
很多時候我們都是這樣的.
最後求你更新我們這個人.
讓我們更認識你,更信靠你,更多的等候你.
很多時候我們未有答案,等久了,未見有人給我們提供答案.
我們就會偏行幾路,我們的頸汗就會硬,我們就會敗壞.
求主祝福我們,求你將一個摩西賞賜給我們.
又求你自己更多紀念你的應許.
以示你的福氣,你的祝福,你的引導帶領.

$^{361}$常常發生在我們身上.
我們仰望主,等候你,祈禱奉靠耶穌,基督的明救.
好嗎?.
\newpage



\section{約書亞記 21:1-45}
\label{sec:032051}
\textbf{應許}~~~~方兆雄牧師
\newline
\newline
連結: \href{http://www.church.com.hk/acms/content.asp?site=cdc\&op=show\&type=product\&code=032051\&layout=sermon}{\texttt{032051}} ~~~~ 語音日期: 2026-04-26
\newline
\newline
\hyperref[sec:032050]{\small{< < < PREV SERMON < < <}}
~
\hyperref[sec:index_chronic]{\small{[返順時目]}}
~
\hyperref[sec:032052]{\small{> > > NEXT SERMON > > >}}
\newline
\newline
約書亞記 21:1-45
\newline
\begin{longtable}{cl}
\hline
\hline
章節 & 經文 (和合本修訂版)\\
\hline
21:1 & \begin{tabularx}{0.7\textwidth}{X} 利未人的眾族長近前來到以利亞撒祭司和嫩的兒子約書亞，以及以色列人各支派父系的領袖那裡， \end{tabularx} \\ \\ \relax
21:2 & \begin{tabularx}{0.7\textwidth}{X} 在迦南地的示羅對他們說：「從前耶和華曾藉著摩西吩咐給我們城鎮居住，以及城鎮的郊外供我們牧養牲畜。」 \end{tabularx} \\ \\ \relax
21:3 & \begin{tabularx}{0.7\textwidth}{X} 於是以色列人照耶和華的指示，從自己的地業中，把這些城鎮和城鎮的郊外給了利未人。 \end{tabularx} \\ \\ \relax
21:4 & \begin{tabularx}{0.7\textwidth}{X} 哥轄族抽了籤。利未人中亞倫祭司的子孫，從猶大支派、西緬支派、便雅憫支派的地業中，抽籤得了十三座城。 \end{tabularx} \\ \\ \relax
21:5 & \begin{tabularx}{0.7\textwidth}{X} 哥轄其餘的子孫，從以法蓮支派、但支派、瑪拿西半支派宗族的地業中，抽籤得了十座城。 \end{tabularx} \\ \\ \relax
21:6 & \begin{tabularx}{0.7\textwidth}{X} 革順的子孫，從以薩迦支派、亞設支派、拿弗他利支派、住巴珊的瑪拿西半支派宗族的地業中，抽籤得了十三座城。 \end{tabularx} \\ \\ \relax
21:7 & \begin{tabularx}{0.7\textwidth}{X} 米拉利的子孫，按著宗族，從呂便支派、迦得支派、西布倫支派的地業中，得了十二座城。 \end{tabularx} \\ \\ \relax
21:8 & \begin{tabularx}{0.7\textwidth}{X} 以色列人照耶和華藉摩西所吩咐的，把這些城鎮和城鎮的郊外，抽籤給利未人。 \end{tabularx} \\ \\ \relax
21:9 & \begin{tabularx}{0.7\textwidth}{X} 他們從猶大支派和西緬支派的地業中，給了以下所記名字的各城， \end{tabularx} \\ \\ \relax
21:10 & \begin{tabularx}{0.7\textwidth}{X} 就是給利未人哥轄宗族的亞倫子孫，因為他們抽到第一籤： \end{tabularx} \\ \\ \relax
21:11 & \begin{tabularx}{0.7\textwidth}{X} 把猶大山區的基列‧亞巴，就是希伯崙，和四圍的郊野給了他們。亞巴是亞衲族的祖先。 \end{tabularx} \\ \\ \relax
21:12 & \begin{tabularx}{0.7\textwidth}{X} 但是，這城的田地和所屬的村莊卻給了耶孚尼的兒子迦勒為業。 \end{tabularx} \\ \\ \relax
21:13 & \begin{tabularx}{0.7\textwidth}{X} 他們把希伯崙，就是誤殺人的逃城和城的郊外，給了亞倫祭司的子孫；又給了立拿和城的郊外、 \end{tabularx} \\ \\ \relax
21:14 & \begin{tabularx}{0.7\textwidth}{X} 雅提珥和城的郊外、以實提莫和城的郊外、 \end{tabularx} \\ \\ \relax
21:15 & \begin{tabularx}{0.7\textwidth}{X} 何崙和城的郊外、底璧和城的郊外、 \end{tabularx} \\ \\ \relax
21:16 & \begin{tabularx}{0.7\textwidth}{X} 亞因和城的郊外、淤他和城的郊外，以及伯‧示麥和城的郊外，共九座城，都是從這二支派中分出來的。 \end{tabularx} \\ \\ \relax
21:17 & \begin{tabularx}{0.7\textwidth}{X} 又從便雅憫支派的地業中給了基遍和城的郊外、迦巴和城的郊外、 \end{tabularx} \\ \\ \relax
21:18 & \begin{tabularx}{0.7\textwidth}{X} 亞拿突和城的郊外，以及亞勒們和城的郊外，共四座城。 \end{tabularx} \\ \\ \relax
21:19 & \begin{tabularx}{0.7\textwidth}{X} 亞倫子孫作祭司的共有十三座城，以及城的郊外。 \end{tabularx} \\ \\ \relax
21:20 & \begin{tabularx}{0.7\textwidth}{X} 利未人哥轄的宗族，就是哥轄其餘的子孫，抽籤所得的城是從以法蓮支派來的。 \end{tabularx} \\ \\ \relax
21:21 & \begin{tabularx}{0.7\textwidth}{X} 他們把以法蓮山區的示劍，就是誤殺人的逃城和城的郊外給了哥轄其餘的子孫；又給了基色和城的郊外、 \end{tabularx} \\ \\ \relax
21:22 & \begin{tabularx}{0.7\textwidth}{X} 基伯先和城的郊外，以及伯‧和崙和城的郊外，共四座城。 \end{tabularx} \\ \\ \relax
21:23 & \begin{tabularx}{0.7\textwidth}{X} 又從但支派的地業中給了伊利提基和城的郊外、基比頓和城的郊外、 \end{tabularx} \\ \\ \relax
21:24 & \begin{tabularx}{0.7\textwidth}{X} 亞雅崙和城的郊外，以及迦特‧臨門和城的郊外，共四座城。 \end{tabularx} \\ \\ \relax
21:25 & \begin{tabularx}{0.7\textwidth}{X} 又從瑪拿西半支派的地業中給了他納和城的郊外，以及迦特‧臨門和城的郊外，共兩座城。 \end{tabularx} \\ \\ \relax
21:26 & \begin{tabularx}{0.7\textwidth}{X} 哥轄其餘的子孫共有十座城，以及城的郊外。 \end{tabularx} \\ \\ \relax
21:27 & \begin{tabularx}{0.7\textwidth}{X} 利未人宗族中革順的子孫，從瑪拿西半支派的地業中所得的是巴珊的哥蘭，就是誤殺人的逃城和城的郊外，以及比‧施提拉和城的郊外，共兩座城。 \end{tabularx} \\ \\ \relax
21:28 & \begin{tabularx}{0.7\textwidth}{X} 從以薩迦支派的地業中所得的是基善和城的郊外、大比拉和城的郊外、 \end{tabularx} \\ \\ \relax
21:29 & \begin{tabularx}{0.7\textwidth}{X} 耶末和城的郊外，以及隱‧干寧和城的郊外，共四座城。 \end{tabularx} \\ \\ \relax
21:30 & \begin{tabularx}{0.7\textwidth}{X} 從亞設支派的地業中所得的是米沙勒和城的郊外、押頓和城的郊外、 \end{tabularx} \\ \\ \relax
21:31 & \begin{tabularx}{0.7\textwidth}{X} 黑甲和城的郊外，以及利合和城的郊外，共四座城。 \end{tabularx} \\ \\ \relax
21:32 & \begin{tabularx}{0.7\textwidth}{X} 從拿弗他利支派的地業中所得的是加利利的基低斯，就是誤殺人的逃城和城的郊外、哈末‧多珥和城的郊外，以及加珥坦和城的郊外，共三座城。 \end{tabularx} \\ \\ \relax
21:33 & \begin{tabularx}{0.7\textwidth}{X} 革順人按著宗族共有十三個城，以及城的郊外。 \end{tabularx} \\ \\ \relax
21:34 & \begin{tabularx}{0.7\textwidth}{X} 其餘的利未人，就是米拉利的子孫，按著宗族從西布倫支派的地業中所得的是約念和城的郊外、加珥他和城的郊外、 \end{tabularx} \\ \\ \relax
21:35 & \begin{tabularx}{0.7\textwidth}{X} 丁拿和城的郊外，以及拿哈拉和城的郊外，共四座城。 \end{tabularx} \\ \\ \relax
21:36 & \begin{tabularx}{0.7\textwidth}{X} 從呂便支派的地業中所得的是比悉和城的郊外、雅雜和城的郊外、 \end{tabularx} \\ \\ \relax
21:37 & \begin{tabularx}{0.7\textwidth}{X} 基底莫和城的郊外，以及米法押和城的郊外，共四座城。 \end{tabularx} \\ \\ \relax
21:38 & \begin{tabularx}{0.7\textwidth}{X} 從迦得支派的地業中所得的是基列的拉末，就是誤殺人的逃城和城的郊外、瑪哈念和城的郊外、 \end{tabularx} \\ \\ \relax
21:39 & \begin{tabularx}{0.7\textwidth}{X} 希實本和城的郊外，以及雅謝和城的郊外，共四座城。 \end{tabularx} \\ \\ \relax
21:40 & \begin{tabularx}{0.7\textwidth}{X} 利未宗族其餘的人，就是米拉利的子孫，按著宗族抽籤所得的，共十二座城。 \end{tabularx} \\ \\ \relax
21:41 & \begin{tabularx}{0.7\textwidth}{X} 利未人在以色列人的地業中所得的城，共四十八個，還有城的郊外。 \end{tabularx} \\ \\ \relax
21:42 & \begin{tabularx}{0.7\textwidth}{X} 這些城的四圍都有郊野，每個城都是如此。 \end{tabularx} \\ \\ \relax
21:43 & \begin{tabularx}{0.7\textwidth}{X} 這樣，耶和華將從前向他們列祖起誓要給他們的全地賜給以色列人，他們就得了為業，住在其中。 \end{tabularx} \\ \\ \relax
21:44 & \begin{tabularx}{0.7\textwidth}{X} 耶和華照著向他們列祖起誓所應許的一切，賜給他們全境安寧。他們所有的仇敵，沒有一個能在他們面前站立得住。耶和華把所有仇敵都交在他們手中。 \end{tabularx} \\ \\ \relax
21:45 & \begin{tabularx}{0.7\textwidth}{X} 耶和華應許賜福給以色列家的話，一句都沒有落空，全都應驗了。 \end{tabularx} \\ \\
[1ex]
\hline
\hline
\end{longtable}
$^{1}$各位弟兄姊妹平安.
大家有印象我們講約書亞記整個系列.
前面我們看過以色列人怎樣過約旦河.
看過打仗的時候耶利哥城怎樣神跡地倒塌.
又見他們爭戰當中神讓日頭停住.
南征北討.
但到了後半卷書第十三章之後.
就開始講到地方名 支派界線.
開始未必那麼容易讀.
尤其是我們去到第二十一章的時候.
更加好像一個利未人的列表.
主編成音這樣.
這個會不會只不過是聖經裡面的歷史記錄呢.
或者有沒有圖畫其實要我們看到的呢.
我相信可以看到裡面有些深一層的意義.
有三方面的.
包括神怎樣去紀念服侍祂的人.
神怎樣為祂的百姓建立一個敬虔生活的制度.
又怎樣神的恩典.
應許一步一步地落實.
而全章最重要的高峰就在最後第四十五節.
我這裡都印得出來.
讓我們一齊去重讀.
耶和華應許賜福給以色列家的話.
一句也沒有落空.
都應而不了.
我相信這節的經文也是作為解釋.
也章裡面.
甚至整卷約書亞記的一條約士.
在當中究竟神對祂所分派出來服侍的人有什麼安排呢.
怎樣維持保證在應許之地的以色列民.
能夠事實上得到神的律法的提醒,印證呢.
最後神的應許怎樣一句也沒有落空呢.
就是我們一齊來探索.
我們一齊來祈禱.
因為你賜下聖經.
你的話語.
天父願以你的聖靈開啟我們的心竅.
明白你的聖經.
也讓我們有這份心智樂意進行.

$^{41}$我們祈禱.
奉主耶穌基督的聖名祈求.
阿們.
這次的經文也會有些分散的.
大家可以.
假如有聖經.
不論電話也好.
教會所提供的也好.
舊約的283年就是約書亞記第21章的位置.
我們會看.
分幾段.
首先開始的時候我們看一下.
利未知派有什麼獨特的事情發生了.
我們一到二節會印在你的螢光幕裡面的.
我們一起去讀.
「那時利未人的眾族長來到祭司爾利亞薩和聯的兒子約書亞.
並以斯列國支派的族長面前.
在迦南地的侍羅對他們說.
從前耶和華席著摩西.
吩咐給我們誠一居住.
並誠一的郊野可以目讓我們的生畜」.
這裡利未人向當時的領袖.
族長,祭司長提出要求.
「耶和華曾經藉著摩西吩咐給我們誠一居住.
並且誠一的郊野」.
他們要求的理據是神早在民數記35章裡面所定下的安排.
在這裡先補充一些背景.
以色列民有12個支派.
而利未知派是一個很獨特的支派.
因為神給他們一個獨特的呼召.
在民數記第8章裡記載.
利未知派是神親自分別出來.
要來事奉耶和華的.
不論是參與會幕獻祭各種敬拜支援的事項.
當其他11個支派在分配土地建立地業的時候.
利未人卻沒有得到大片完整的土地作為他們的產業.
而民數記第18章也這樣記載.
耶和華這樣說「我就是你的墳,是你的產業」.
這個「產業」在希伯來文所用的字眼不單止是指財產.
更加是承傳下來.

$^{81}$可以讓他們生活裡面安心立命的地業.
也提到「墳」.
對一般的支派來說.
他們的產業是甚麼呢?.
就是土地,可以耕作,可以畜牧.
但對於利未人來說.
他們的產業卻不是地.
而是耶和華自己.
雖然身份很特別.
但同時也帶出一些很現實的問題.
當他們沒有擁有土地的時候.
他們怎樣生活呢?.
他們的家庭,兒女,生產日常怎樣安排呢?.
但經文裡面給我們看見.
不單止揀選,呼召他們的子民.
也會供應他們所需要的.
利未人雖然沒有地圖上的版圖.
但他得著48個城邑和附近的郊野.
翻查一些參考資料.
並不是指48個城邑只給利未人居住.
而是仍然居住在支派當中.
不過他們有權在城邑居住.
而各支派就照著耶和華藉著摩西所吩咐的.
將城邑給了他們.
這裡值得我們思想留意的屬靈原則.
就是神既然將他們分別出來服侍神.
亦會負責他們的生活.
在世俗的眼光來講.
利未人好像少了些甚麼.
沒有土地就沒有了當時帶來穩定,資源.
或者安全感.
但神的設計裡面.
他沒有看輕他們僕人的需要.
主耶穌後來差派門徒出去.
主耶穌吩咐他們要先求神的國和神的義.
而其他的天父是知道他們所需要用的.
今日我們的教會場景也是這樣.
有人被神呼召全職服侍.
在教會裡面工作.
在機構裡面工作.

$^{121}$亦有人被神呼召在職場和崗位裡面服侍.
或者在平時不用工作的時間參與教會的服侍,教導.
探訪,祈禱,關懷的事情.
每一個跟隨主耶穌基督的信徒.
都有祭奠這個身份.
而在姐妹多方面的付出.
時間,心力,作出的犧牲.
的確未必在世間裡面有甚麼回報.
你不會說成長班增長了10\%,20\%.
有獎勵給你,沒有這回事.
但神所分派的慰訓.
並不是用世間的尺度去衡量.
有時這個訓可以是產業,的確是.
我們在神家裡面事奉.
認識很多弟兄姊妹.
譬如啟德那裡有人住豪宅,有人住公屋.
有時是恩賜.
為甚麼這個人做這件事情的時候.
那麼上手,那麼容易,那麼厲害.
為甚麼我弄多厲害都不行呢?.
有時亦都可以是兒女.
和神所安排的時間.
好像在新約的時代.
當神要為主耶穌基督.
預備一位先鋒的時候.
就讓到年老的祭司撒加尼亞.
他的太太以利沙白.
年老的時候才生出施贈約翰出來.
聖經給我們看到最重要的問題.
並不是我有或者沒有.
或者我擁有的多和少跟別人比較.
而是我們是否活在神賜給我們的墳裡面.
當然這個墳亦不是一成不變的.
按照神的時間.
不同的需要.
情況變化.
神會賜多.
神亦都會收回.
利未人的見證亦都是這樣.
他們雖然沒有像其他支派那樣.

$^{161}$有地.
沒有世人眼中的保障.
但是他是神親自為他們準備這個居所.
希望在這段經文裡面給我們看到.
神是紀念服侍他的人的.
亦都會供應這些人的需要.
第二點.
可能想講長一點的.
就是我們留意到這個名單.
第九節到第四十二節.
很長的名.
很多的地方.
很多的成音.
這段的名單裡面.
很簡單來說.
分成四個大段.
第一段就是利美的大兒子.
哥核宗族裡面的.
屬於亞倫的家族的支派.
因為他們專門祭司的職份.
他們就是第一部份.
第二部份就是大兒子其餘的宗族.
哥核的其餘的子孫.
第三部份就是二兒子格遜的宗族.
第四部份就是米拉尼.
第三兒子的就是各人的宗族.
他們各自分配了不同的成音.
總共四十八座城.
值得留意的.
就是這些成音是很分散的.
差不多鋪滿了約旦河東西的不同地方.
好像很零碎.
好像不集中.
好像在地圖上.
沒有一個屬於他們的顏色.
但是.
神這樣安排.
他的心意會不會就是用屬神的經脈.
法律的教導.
神與子民立約的關係.

$^{201}$時常在神的子民當中提醒他們呢?.
利未人的職責.
不知道大家會不會有些誤解.
覺得他們只是在憲制裡面幫忙.
祭司.
那些什麼湯牛羊那些.
其他利未人就幫忙.
可能是抹東西 洗東西.
或者起會幕之類.
當然 祭司的制度.
會幕的事奉.
他們其中一個極重要的使命.
但並不是說.
他們不用輪班的時候.
就在他們的城邑裡面躺著.
不是的.
因為在《不道德摩西五經》裡面.
記載歷史書裡面所給我們看到.
利未人他們的角色.
就是要在神的子民中.
教導律法.
幫助他們分辨.
潔淨和不潔淨.
要引導百姓認識神的心意.
正如《申命記》三十三章提到.
他們要將神的典章教訓雅各.
雅各的意思就是以色列民.
將神的律法.
來教導.
讓人知道神的心意.
在舊約應許之地.
利未人.
被安置分散在各支派當中.
就是神很有意 刻意的.
使他們的話語.
那種生活的秩序.
屬靈的要求.
讓到人需要有幫助.
或者有疑問的時候.
可以找到他們.

$^{241}$這裡.
神不讓利未人去分散.
亦都提醒百姓.
你不只是每年過時過節的時候.
上去會望.
不論那時候的時羅.
或者後來的耶路撒冷.
守節期.
才去敬拜神的.
以日常生活裡面.
各樣的.
責任.
排難解分.
田地使用權.
審判案件.
上星期曾牧師提到屠城.
你說你不殺人.
你說你真是無辜.
但是審過才知道.
哪些人有權力呢.
哪些人懂得當時的律法呢.
利未人.
商業的買賣.
公不公平.
各類人際關係當中.
都需要在當中有一個指引.
而利未人分散在以色列當中.
就是讓到人可以.
相對方便地.
尋求到神的旨意.
神的律法.
的確.
到今時今日.
我們人生很多問題上.
很需要神的話語.
給我們指引.
很需要神的同在.
給我們安慰.
其實都不單止是平常的.
就算資深的領袖.

$^{281}$同樣都需要.
在這裡分享一個曹牧師.
曹院長.
在他二二年.
浸神院信裡面所分享的經歷.
也都輯服了在他那本.
《深化神學智慧》.
剛剛出的書.
250頁那裡.
在09年曹牧師.
曹院長他分享.
他在浸神擔任院長第一年.
在6月的時間去美國.
探訪當地的教會.
探訪當地的校友.
期間他也都去拜訪.
年近一百歲的林道亮牧師.
這位林牧師曾經在台灣.
中華科音神學院.
擔任院長十年的.
曹院長很希望林牧師為他祈禱.
並且過了一些年.
面授一下作為院長的心得.
有什麼要注意留意.
當他們開始交談的時候.
都很期望林牧師.
關於可不可以提醒一下.
他有什麼提點.
作為院長有什麼行政.
有什麼要留意之類.
但是.
林牧師只是跟他說聖經.
曹院長憶述.
林牧師用了三節經文.
勉勵他要留心聆聽神的心聲.
起初曹院長有點失望.
心想這些經文我都懂.
我都懂得解釋.
教過.
但是當謙卑下來.

$^{321}$仔細聆聽的時候.
他感到十分感動.
感受到神的令和他們同在.
接著當林牧師為他按手祈禱.
也都祈求他令命強健.
很少提及院務如何處理.
在祈禱按手期間的時候.
曹院長記得他整個人都在顫抖.
林牧師本來對於他來說.
是一個需要拜候的長者.
前輩 陌生人.
但是卻成為神的使者.
作為他當時的屬靈的導引.
甚至林牧師送別他的時候.
讓曹院長看見好像天父的微笑一樣.
這個豈不是好像利未人的模式.
當人需要神的指引.
祈禱的時候.
神的話語的時候.
去到利未人的前面.
感受到這種見證.
用新約教會的概念來說.
這是神如何差遣我們每一個人.
進入世界.
今日的信徒我們不像舊約利未人那樣.
有血統.
怎樣才能被服侍.
沒有這回事.
但是彼得前書這樣說.
我們在基督裡面.
全教會都是有君尊的祭司.
而我們的責任就是在社會裡面.
不同的範疇裡面.
如何活出見證.
不論是學校教學的.
做支援的.
如何按神的形象和樣式去培育.
年輕的生命呢.
在醫院裡面做醫護的.
或者行政的.

$^{361}$如何透過醫療.
透過制度上幫助別人去康復呢.
在藝術上如何能夠表達出.
神所創造的美.
如何去批判罪的惡.
和對人性的扭曲呢.
在商業上如何能夠將人的需要.
和世界的資源.
人的才能當中結合呢.
這個商而通之呢.
在日常生活裡面.
可能我們都覺得自己只不過是普通人.
但是神將我們放在不同的地方的時候.
就是要讓人可以見到神和他們同在.
在當中發揮或大或小的影響力.
不論是安慰啦.
祈禱啦.
神的話語啦.
指引啦等等.
在昨天.
都發生一件小事.
不過都是很特別的事.
是在師班室附近.
有位小朋友哭得很厲害.
爆哭.
我們關心他問他什麼事.
那時候正在練師.
我剛剛在那邊走過才見到他.
原來他不見了電話.
很傷心.
他說他電話裡面有很重要的資料.
當然我心想.
其實不應該這樣想的.
有什麼重要資料.
但是他真的哭得很傷心.
於是我和他一起祈禱.
交託給神.
求神看看能不能找到.
當然師班完練師的時候.
事務完.

$^{401}$拿回閉路電視看看之類.
接著沒多久那個小朋友就走了.
後來問.
原來他找到他的電話了.
在哪裡找到呢?.
在他自己的袋子裡面.
這位小朋友很有禮貌.
他透過導師傳訊息給我.
和他一起祈禱.
我所想到的不僅是小朋友.
他不見了手機.
或者遊戲裡面很珍貴的武器.
或者和朋友的照片之類.
但我想到的是弟兄姊妹.
或者失去親人.
或者發現自己有頑疾.
或者可能在投資上.
同事稱呼他為「古神」.
但是他損失了大半的積蓄.
他能不能夠在教會城鎮裡面爆哭.
有弟兄姊妹.
導師和他一起祈禱.
給他安慰.
這豈不就是利未人的福事呢?.
最後這段經文裡面.
43到45節.
有一個短的段落.
我們之前讀到45節.
耶和華應許賜福給以色列家的話.
一句都沒有落空.
應驗了.
而怎樣實現他的應許呢?.
在第43節我們一起讀.
耶和華將從前向他們列祖.
起誓所應許的全地.
賜給以色列人.
到這裡為止.
這裡清楚講到是耶和華賜的.
應許之地本質上就是恩典.
不是他們的努力.

$^{441}$不是他們的報酬.
當然以色列人都是要打仗的.
都是要死人的.
在當中也都經歷了神的各種保守.
假如這個地圖是人憑自己力量所爭取的.
人很自然地也都應該垮口.
但是假如承認是神所賜的時候.
才有這份謙卑敬畏的心來承受.
經文也都特別提到是向他們列祖.
起誓所應許的.
一下子將這個時間軸拉回很長了.
所謂列祖是哪些人呢?.
就是神對亞伯拉罕.
以沙,雅各所許解應許.
在近之後四百多,五百年的時候.
到約書合時代.
才具體地圖居住上得著應驗.
很長的時間.
經文提到神向他們列祖起誓.
起誓這個詞當然有立約.
有保證的意思.
是以神自己的信實作為保證的.
跨越世代四,五百年.
終於成為可見的現實.
不單止這樣.
第四十四節經文也都提到.
神使他們四境平安.
神帶他們的百姓去到應許之地.
不單止留下他們有地方住.
更加讓他們享安息,四境平安.
不代表從此他們沒有仗打.
因為聖經裡面都不是這樣記載的.
之後他們犯罪,得罪神.
於是神要管教.
亦都用其他外邦去成為他們的審判.
但起碼在那段時間.
從年年的征戰,漂流當中.
現在可以有安居.
可以耕田,牧羊,建家立業.
可以有敬拜的生活.

$^{481}$這就是神給他們恩賜,賜給他們平安.
更加神讓他們在其他仇敵當中.
能夠站得住,因為將一切的仇敵都交在他們手中.
讓他們知道是耶和華為他們征戰.
以至最後第四十五節.
好像我們開始讀到一樣.
神的應許所說的話一句都沒有落空.
整個時間線讓我們看回.
從亞伯拉罕,蒙召.
你要離開自己的本族,本家.
往我指示你的地方去.
到了約瑟,整個家族.
到了埃及,從饑荒中得到拯救.
到了後來又不認識約瑟的法老.
紮莊了.
於是就欺負以色列民.
讓他們做苦工,在埃及受苦.
到了出埃及.
由摩西帶領他們過紅海.
到了,快要進入加林地了.
但是因為他們不信.
於是再在曠野飄來四十年.
然後新的一代才進入到迦南美地.
四百多年,五百年.
在人的角度去看.
這個是不是神的拖延呢?.
擔延呢?.
或者神的應許是否真的會成就呢?.
四十五節給了我們很確切的答案.
就是會,神的應許一句都沒有落空.
很多時候,我們很多年都未必看見神的安排,傳授.
或者甚至對於很多代人來說.
可能他們的世代都過了.
他們自己的生命已經完結了.
都未看見神的應許怎樣實現.
但是往往當人可能是回頭看.
才看到,原來神怎樣運用他的生命.
在這個時空當中.
為人帶來祝福.
其中想介紹一位歷史人物.

$^{521}$十九世紀英國的George Buller.
喬治·莫勒.
他怎樣來經歷神的信實.
在他的生命當中,其中最為人熟悉的.
就是他在Bristol英國的地方.
開辦照顧孤兒的工作.
他受到先人的感召.
亦都是神給他的感動.
在1836年的時候.
莫勒和他妻子開始了時工照顧孤兒.
他們先在布里斯托城市其中一條街.
Wilson Street 預備了一個房子.
當時收容三十個小朋友其實都不簡單.
不過之後因為他們都辦得不錯.
於是根據社會的需要.
第二間,第三間,第四間的房屋.
都陸續投入使用.
後來更加在鄉郊建立了.
一系列專門為照顧孤兒而設的院舍群.
莫勒他不是一個有錢的人.
亦都不是一個大慈善家.
但他有一個很鮮明的原則.
他從來不向人募捐.
亦都不借貸.
所有的需要都是用祈禱帶到神的面前.
他自己亦都被路斯篇.
八十一篇第十節的經文感動.
「要大大張口我就給你充滿」.
莫勒深信.
這個侍工是神呼召他做的話.
神必然會供應.
他其中一次經歷神的應許.
就是當孤兒院的資金用盡了.
相信都是初期的時候.
孤兒院裡面甚至食物都沒有了.
莫勒依然安排小朋友早上的時候.
坐在餐桌前面.
做藉飯祈禱.
但甚麼都沒有.
那次祈禱完之後.

$^{561}$門外有人敲門.
原來是一位麵包師傅.
他說自己夜晚裡面一直睡不著.
心裡面有聲音催促他們.
要為孤兒院焗幾爐麵包.
他凌晨兩點鐘起身準備.
到剛剛早餐時間之前送到.
接著馬上有一位送奶工人來到.
因為原來他們的運奶車.
那些不知是牛車還是馬車之類.
在孤兒院前面壞了.
他都修補不了.
那些牛奶不像現在有pasteurized.
於是假如不送出去的話.
就會變壞了.
與其浪費.
不如全部送給孤兒院.
很奇妙的經歷.
而George Muller他的經歷.
亦不單止是神跡式的供應.
他在Bristol幾十年的服事裡面.
在他相關的院舍裡面.
受照顧的孤兒總數達到一萬人.
很誇張的.
因為在當地都有紀念他時功的博物館.
而他自己在七十歲開始.
對當時來說已經是很老年紀了.
就去周遊列國去暴動.
足跡去過四十二個國家.
傳主耶穌基督的福音.
將神奇妙的供應.
各樣奇妙的英語英文.
來跟人分享.
而在他自己的記錄裡面.
在他一生裡英文過超過五萬次的祈禱.
有時孤兒院的經費.
在門口不知哪裡.
洩過信封進來.
在街上有人給他一些硬幣.
這樣去籌回來.

$^{601}$而他的工作也成為了一代人的祝福.
要知道十九世紀.
剛剛.
尤其是英國工業革命.
很多時候都找一些小朋友去做工.
尤其是那些無依無靠的孤兒.
最好剝削他們.
又低薪.
一天有記載.
做十二個小時.
甚至做十四個小時.
十六個小時.
亦都在這些環境之下.
孕育了馬克思.
馬克思.
見到這個社會不公平的形象.
於是引發他去研究.
去提倡共產主義.
但是.
Mullah.
他當中成為了一代人.
一個城市的祝福.
但是祝福也不是沒有批評的.
那些人批評.
你教育小朋友這麼好.
搞得我們沒有廉價勞工.
工廠怎麼運作之類.
有些商家會這樣批評他.
但是他祝福了超過一萬個生命.
讓我們有書讀.
學聖經.
懂得怎樣敬拜.
成為一個時代的祝福.
到了今時今日.
我們怎樣了斃提醒呢?.
若書一記第二十一章.
我們所看到的.
可能很多的文字裡面.
很多地名,人名.
或者這個歷史記錄.

$^{641}$卻又給我們看見.
第一方面.
神揀選出來.
要事奉祂的人.
未必有世俗裡面.
給到祂安全感.
但是神必然紀念.
有時我們因為忠於信仰.
在職場上面.
在社會上有些事不肯做的.
我們有自己的立場.
有自己的方向.
可能會吃虧.
大家記得剛才我們崇拜.
讀啟應經文.
行善受苦.
有時我們花很多時間.
在事奉上.
頂尖本身上.
都被人覺得.
做得不好.
批評這樣.
但是很希望.
我們知道.
神看見.
神分派我們.
神會準備給我們的供應.
知道我們的需要.
這是第一點.
第二方面.
我們看到神也很看重.
怎樣去持續地.
讓祂的子民.
能夠得到律法.
在屬靈上各樣的幫助.
約書日記前半卷書.
講他們怎樣去到.
國有大河,攻城.
打的戰役.
畫面很震撼.

$^{681}$成世家同樣過去.
經過很多歲月.
不同的風雨.
社會的挑戰.
在早期時向.
也都經歷過打仗.
到今天.
我們仍然各位頂尖姐妹.
都可以繼續了敬拜神.
並不是因為偶然一次.
什麼大型成功的報道會.
說幾千人信主.
或者什麼驚天動地的計劃.
而我相信是長年累月的.
不論透過我們一起的.
去崇拜.
聽聖經,講道.
在以前的主日學助道會.
現在的成長班裡面.
來查經,祈禱.
接受牧羊.
去探訪,關懷.
師班來講.
我們每個星期的練師.
是細水長流的功夫.
不一定好像打仗.
也都不一定很亮麗.
但十分重要.
而約書日記裡面.
讓我們看到.
神安排你們人.
在不同的地方的時候.
你都維持著他們敬拜的生活.
到今時今日.
教會會不會也都成為了.
社會的那些利味成泣呢?.
我們的兒童成長班.
我們的長者祈壯中心的活動.
我們青少年的栽培.
或者我們背後行政.

$^{721}$或者是採用財務的安排.
這些未必有人看得見的.
但是我們同樣要建立在神的子民.
成為這個社會裡面登台.
當人有需要.
和有迷惘.
生活裡面失去方向的時候.
我們可以指向.
在神裡面.
一個有神的印象.
有樣式的生活是怎樣的.
第三方面.
也都讓我們看到.
神的應許.
未必一定讓我們立刻看得見的.
但是神卻是從不失信.
祂所講的必然會成就.
我們在45節那裡.
一個最有力的宣告.
神應許賜福給議事家的說話.
一句都沒有落空.
今時今日.
我們活在這個很快速的時代.
搭電梯也好.
先按了關門鍵.
才按哪層樓的.
總之每件事都要快.
一定要有果效.
立竿見影.
但是我們的屬靈的眼光.
被這種影響節奏所影響的時候.
我們不太習慣.
祈禱的時候.
為什麼祈禱這麼久.
神都沒有應允的.
或者是自己的健康.
或者家裡.
或者未信的朋友親戚來到祈禱.
我們的前路.
很多都不知道怎樣走的時候.

$^{761}$或者會開始懷疑.
神會不會忘記了呢?.
但是約書亞記給我們看到的.
就是神所應許的說話.
一句都沒有落空.
以色列文.
由阿巴拉罕到約書亞那裡.
前後五百年.
寄居過.
承受過苦難.
做過苦工.
有等待的時間.
有漂流的時間.
征戰的時間.
但是到了神的計劃.
又得到彰顯的時候.
一句說話都沒有落空.
而我們在約書亞那裡面的時候.
不單止看見神的信神.
更加相信祂是在預表著.
將來主耶穌基督所帶來的救恩.
起碼我們讀聖經都知道.
在約書亞記.
他們雖然得到這個加寒地.
但是沒多久.
士師記馬上對他們犯罪了.
然後就是神要去審判了.
甚至後來亡國分裂.
亡國被擄等等.
這個安息.
這個應許之地都是短暫的.
但是.
祂指向一個更大的安息.
就是在主耶穌基督裡面所得著的生命.
主耶穌在十字架上面.
勝過罪惡的權勢.
徹底打敗罪對人的影響.
以致我們能夠可以藉著主耶穌基督.
與神有和好的生命.
而主耶穌所帶給我們的.

$^{801}$也都不是一段.
或者只是一生的平安.
而是這份永恆的平安.
都在我們身上.
約書亞記.
有一章四十五節所講到.
一句都沒有落空.
到了新約.
就是史圖保羅在哥林多後書一章二十節所講的.
神的應許.
不論有多少.
在基督裡面.
都是事的.
Yes.
希望我們讀這段經文的時候.
縱使我們未體現結局.
或者我們在這個過程當中的時候.
都深信神會用我們.
神讓我們經歷祂個人的恩典.
到最後了.
希望在這章聖經裡面.
雖然我們見到很多人名和地名.
我們都會記得神紀念服侍他們的人.
神亦都讓敬拜的生活.
滲透在.
不論是我們自己生活上.
社會裡面不同的角落.
讓到神的說話.
神的心意.
神的同在安慰.
都能夠讓人.
願意尋求神的時候.
可以借助我們找到.
而第三方面神所應許的.
亦都沒有一句落空的.
希望今時今日.
或者可能你正面對若干的困難.
若干的迷惘.
都深信.
神的旨意.

$^{841}$在你身上的旨意.
會成就.
最後讓我們再讀一次第45節.
成為我們的提醒.
預備起.
耶和華應許賜福給以色列家的.
一句也沒有落空.
都應而不了.
讓我們一起祈禱.
慈愛信息的天父.
因為你是守施慈愛的神.
神亦都讓我們看見剛才的經文.
信息裡面.
有甚麼提醒我們的.
有甚麼要我們學習的.
願意你的聖靈感動引領.
主耶穌基督.
我們當中有不少弟兄姊妹.
處於掙扎.
心靈過度疲乏.
很大的重擔.
在等候當中.
或者迷惘的裡面.
求神你親自的來承托.
支持.
賜下安慰.
平安.
更加藉著弟兄姊妹的關心.
祈禱.
讓他們經歷到.
神的同在和保守.
在不容易的世代當中.
讓他們衷心的倚靠你.
不單止讓我們自己得著幫助.
更加讓神的名得到榮耀.
天父願你隨天預立我們祈禱.
因為我們祈禱奉主耶穌基督的聖名祈求.
阿們.
\newpage



\section{馬太福音 5:3-10}
\label{sec:032052}
\textbf{真正有福的人}~~~~朱卓慈牧師
\newline
\newline
連結: \href{http://www.church.com.hk/acms/content.asp?site=cdc\&op=show\&type=product\&code=032052\&layout=sermon}{\texttt{032052}} ~~~~ 語音日期: 2026-04-25
\newline
\newline
\hyperref[sec:032051]{\small{< < < PREV SERMON < < <}}
~
\hyperref[sec:index_chronic]{\small{[返順時目]}}
~
\hyperref[sec:032053]{\small{> > > NEXT SERMON > > >}}
\newline
\newline
馬太福音 5:3-10
\newline
\begin{longtable}{cl}
\hline
\hline
章節 & 經文 (和合本修訂版)\\
\hline
5:3 & \begin{tabularx}{0.7\textwidth}{X} 「心靈貧窮的人有福了！因為天國是他們的。 \end{tabularx} \\ \\ \relax
5:4 & \begin{tabularx}{0.7\textwidth}{X} 哀慟的人有福了！因為他們必得安慰。 \end{tabularx} \\ \\ \relax
5:5 & \begin{tabularx}{0.7\textwidth}{X} 謙和的人有福了！因為他們必承受土地。 \end{tabularx} \\ \\ \relax
5:6 & \begin{tabularx}{0.7\textwidth}{X} 飢渴慕義的人有福了！因為他們必得飽足。 \end{tabularx} \\ \\ \relax
5:7 & \begin{tabularx}{0.7\textwidth}{X} 憐憫人的人有福了！因為他們必蒙憐憫。 \end{tabularx} \\ \\ \relax
5:8 & \begin{tabularx}{0.7\textwidth}{X} 清心的人有福了！因為他們必得見神。 \end{tabularx} \\ \\ \relax
5:9 & \begin{tabularx}{0.7\textwidth}{X} 締造和平的人有福了！因為他們必稱為神的兒子。 \end{tabularx} \\ \\ \relax
5:10 & \begin{tabularx}{0.7\textwidth}{X} 為義受迫害的人有福了！因為天國是他們的。 \end{tabularx} \\ \\
[1ex]
\hline
\hline
\end{longtable}
$^{1}$自從林誠信牧師請我降道時,有些同工問我如何介紹窩福堂的前堂主任..
我想如果說窩福堂,即是窩福,即是窩到露露山,即是知道一個地方,.
不過原來也有個問題,即是太長,所以很開心今天能夠來到大家一起分享神的話語..
首先我們一起有個禱告:.
「慈悲天父,我們多謝你,我們在主裡面有愛,我們一同分享真理,願你今日將你的話語去教導我們,.
以做我們聽的得著幫助,祈禱奉主耶穌,祈禱祂的名,阿們.」.
今天跟大家講這個題目是「真正有福的人」,.
言下之意,很多人都覺得有福,不過什麼是真正呢?這是值得我們去想的..
或者我問大家什麼叫做「有福的人」呢?什麼叫做「福」呢?.
我相信大家都會想到,即是我都會想到,.
有好本事,找到工作,有穩定的收入,.
還有一個是很有成就,即是如果能夠做到別人知道你的名字,那就更加厲害了..
還有在生活上,我們有好的家庭,好的配偶,好的婚姻,.
有時間我們都很喜歡出去旅行,然後到了晚年,我剛剛就落任了,我退休了,.
能夠生活無憂,這個不是福呀?.
有多少人說這個福,你舉舉手給我看?.
是呀,我都是這樣說的,這個沒錯,但是我們要知道,.
剛才我講到那些種種的福,都是一些物質上的一些愛在的東西來的,.
這些能夠滿足到我們的心靈呢?或者我們人生是不是就這樣得到一個滿足的人生呢?.
我相信大家都聽到最近很轟動的一件事情是Jeffrey Esten,.
我們翻譯為愛潑斯坦,他是一個很出名的一個,因為出名犯罪,.
在幾年前其實是很轟動,他是一個很聰明的人,很會賺錢,他是一個投資家,.
但是還是一個罪的罪犯,性的罪犯,特別牽涉到一些人口販賣,.
或者專門將一些未成年少女逼他們做一些不道德的一些賣淫的工作,.
後來審了他很久,結果在2019年他被捕,在7月被捕,然後在8月就死亡,.
為什麼死亡呢?就是官方說是自殺,但是後來那些人很多的這件事情牽連甚廣,.
可能都有些境廣在裡面,是謀殺的居多,很多的嫌疑,在很多輿論的底下,.
在2025年,就是過了幾年之後,美國的司法部就揭發一個文件叫做Esten的檔案,.
這個檔案為什麼這麼轟動呢?就是原來發過這個檔案將所有和Esten有關聯的人的名單披露出來,.
其中包括了一些億萬的巨富,有些商界,文化界,教育界的知名人士竟然在裡面,.
甚至乎一些國家的政要,或者一些王室的人士都在當中,.
當然這些名字在名單上不是代表你一定犯罪,但是有不少的人在這個檔案當中的人物是被夾擊,.
以及被處分,有些現在還在審判中,這件事情告訴我們什麼事情呢?.
原來這個世界很發達的背後,那些領袖高層原來有很多的黑幕,.
人生以為我們很羨慕的名譽,地位,權力,金錢,原來背後有很多人的軟弱,敗華在當中隱藏著,.
什麼是找到人生的滿足呢?什麼是福分呢?是不是百萬富翁呢?.
我們看到這段經文很震撼,主耶穌基督在一個當時都很窮困,很多人很艱難,疾苦的當中說了八個福分,.
這個八個福分是很顛覆的,一般人我們剛才看到的福分是這樣的,是一些外在,.
而這個福分,這八種的福分是人內裡生命的,是超越環境的一個福分來的,.
所以我們好好地看,可能你會看到很奇怪,這些都是福,或者我們先看看開頭那四個福,.

$^{41}$我們經文這裡,或者我再讀出來給大家一起去學習,第五章三節,這裡是第三節至第六節,.
他說心靈貧窮的人有福了,因為天國是他們的,哀動的人有福了,因為他們必得安慰,.
謙和的人有福了,因為他們必承受地土,飢渴無意的人有福了,因為他們必得飽足,.
這裡講到四個福,我們先看著這四個福,我們發覺這八個福是緊緊相連的,.
很有一個權柄在當中,或者我們首先看看心靈貧窮的人有福了,心靈貧窮,即是知道自己不足,很貧乏的,.
以至於人願意謙卑在生命裡尋出一些答案,原來很吊詭,當你覺得沒有的時候,你就很有福,.
但人往往不是這樣看,或者我們看看這個福是怎樣的福,他說天國是他們,很豐富,.
天國即是有著神的國度進入,永恆當中有神的眷顧,神的塗造,就是因為你承認你是一個一窮二白的人,.
有一位教授在伊利諾州做哲學的教授叫Hughes Moorhead,他在研究一個問題,究竟人生的意義是甚麼,.
於是他用了很多年的時間,分別發給250個傑出之士,包括了哲學家,科學家,音樂家,作家,問他們一個問題,.
簡單來說,甚麼是人生的意義,結果回覆說,我儘管給一個我最好的猜想,有些說我儘管作一個人生的目的意義給你,.
有些坦白說,我毫無頭緒,其中一個我想你們一定認識的Carl Jung,一個心理學家,非常出名,他的答案是這樣說,.
他說,事實上我不知道甚麼是人生的目的和意義,就是這樣,然後有一個很出名的美國作家叫Joseph Heller,他怎樣說呢,.
他說,人生的意義我沒有答案,我不想再去找,就是這樣,原來沒有人,也有一批是這樣寫的,很謙卑地說,.
魔吃,如果你找到人生的意義,請你告訴我,這本書大家可以去看,1988年魔吃將這些回覆編輯成書,叫做人生的意義,很有意思,.
原來人生的意義在乎你看到你自己,你沒辦法只用你自己的智慧很有限,事實上我們的生命存留在世上多少我們都不知道,.
若然我們人願意去謙卑,去尋求神的時候,我們發覺天外有天,天國是他們的,這是第一個福,有了這個福會帶來很多福,.
然後第二個福,哀動的人有福,哀動的人是甚麼呢,當一個人願意去尋求神,找答案的時候,他會檢視一下自己,.
發覺自己原來有很多問題,很多過錯,以致他能夠懺悔痛哭,也為了周邊的罪惡,周邊的人犯罪而難過流淚,.
我不知道大家有沒有哭過,人最慘是在現代社會我們麻木了,我們不懂得哭,見到自己一些問題的時候,我們也不再去想,.
我們會去痛哭,原來這個是福利,我們會看到人有兩個危機,第一就是我們不認自己錯,錯的只有別人,因為人都是自我中心,這是第一個危機,.
第二個危機就是即使我們知道自己有罪,我們不去面對,我們也不懂得處理,有良心的人更慘,.
活在罪疚當中不原諒自己,很痛苦,但這裡說哀動的人有福利,因為他們必得安慰,因為在神那裡有出路,.
神會向神去認罪悔改,要神赦免,要讓我們得到釋放,並且在神的恩典裡面,我們要彌補我們的過錯,.
我有一個很熟悉的親人,很聰明,我覺得他真的很聰明,他想的東西真的很好,現在說機械人,他小時候說他想的都是機械人,我覺得他很聰明,.
但很可惜,聰明反被聰明誤,當然讀不成書,後來群組了一些壞人,還沒畢業,還沒畢業就出去做工,做營業代表,.
因為有些聰明真的發跡,在行內有名,挺好,但後來虧空供款,名聲差了,爸爸媽媽幫他,後來他自己做生意,.
但又理財不善等等,變成又失敗,父母又幫他,後來沒了,撈來撈去,開的士,開的士的時候,會落監的時候,一班一班的同事就說不如玩玩,.
越賭越大,或者很多大耳窿知道搞不定,父母又去幫忙,到爸爸死的時候,他告訴我,我是無動於衷,有一次我收到他電話,.
我說你在哪裡,我聽到有點不對勁,他原來在山頂預備把車開到山崖,我說你不要,因為走投無路,說人的絕路就是神的出路,.
結果他沒有做這件事,長話短說,後來有一次他坐一個客人去到一個地方,停了車,剛好時間對面有間教會,其實是窩福黨,走來聚會,.
牧師講道理,講到人生很多道理,原來人生是這樣,他覺得很有意思,繼續定期回來教會,然後有一次報道會,牧師邀請人決志順著他的手走出來,眼裡面充滿著淚水,.
他為了這麼多年的罪過哭泣,從此生命改變,爸爸雖然離開,媽媽還在,買些好東西給媽媽,媽媽喜歡吃雲吞粉,經常買些雲吞粉,生命改變,哭泣的人有福了,.
哭泣能夠令一個人的生命在主裡面流轉,第三個福,謙和的人有福了,如果一個人願意在神面前悔改認罪,他必定會馴服神在他生命的帶領,行事為人,待人切勿,願意跟隨神,.
以致他能夠待人有一種謙和的態度,處處看到人的需要,而自己願意去謹謹慎慎行事,以致別人得著好處,溫柔不代表軟軟弱弱,別人說什麼就說什麼,.
也有著自己的神的帶領,然後他就謙謙和和去實行神所起的事,在今天的很競爭性很大的社會,如果你謙和,我想看你如果成為謙和的時候,你會不會很吃虧,.
我們很多時候在競爭的時候,說你搶先人的說話,你就會比較有好處,因為先講你的理,更加用一些壓力,強權的方式,往往是得到成就的,.
這些成就大概是短暫,別人口服心不服,若然你溫溫和和,以理服人,這種正面的影響是很深遠,這裡更加講到神給他的福,因為他們必承受土地,我們不需要搶奪,.
神是話事的,神會將我們應得的那一份給我們,我們不需要那麼狼狽,所以人能夠謙謙和和,面對很多,就算是多少的一些壓力,我們仍然可以按著神的心意,可能是得罪人,.
但我們照樣走,我們就得到神給我們的那一份,很有把握的,第四個福份,第六次主耶穌說,饑渴無意的人有福了,.

$^{81}$即是公義的生活,能夠追求到一個很渴很飢,之前講到謙和的人是按著神的心意去走,很自然,神的公義他很想進行,.
是一路貫穿下去,人是把持著神認為合宜的東西,做出公義的行為,這個就是饑渴無意的人的福,我想每個人都想做好人,.
但是否滿足呢?這個就另計了,每個人都有自己的道德標準,不過這個是相對的,你覺得這樣好,各個國家都有不同的道德標準法例,.
有些地方,你吃大麻,我不是說大家吃,你吃大麻是不犯法的,其實是真的,我女兒現在在美國,有次有朋友跟她一起,你吃糖果,那些糖果是大麻,.
我女兒說不好吃,我寧願吃小朋友那些維他命的甘蜜啤,你隨時可以有自己的標準,這裡這樣那裡這樣,但神是有美善的絕對標準,.
不是任何一方,任何一個看法,絕對美善,你說這是一個福,如果我們知道道德標準,但也不是,走出來才是一個福,.
但聖經裡面很清楚告訴你,人能夠靠自己可以行出公義,聖經裡面說你立志行善你可以說一下,但走出來不由得你,.
聖經裡面說我們人有罪性,我不知道大家會不會很掙扎,就算你做基督徒,幾何無疑我是很想,但有時我們好像軟弱無力,.
但我們知道這個公義,神是得到給我們完滿的滿足,但要依靠他,有一個讀哲學的人,不知道大家有沒有聽過范孝德,.
他是一個來自中國背景的學者,後來他成為牧師,他說了一段話很有意思,.
他說在我少年的時候我犯多少罪我忘記了,不過我看我現在這個小朋友這個兒子,我就能夠看到我以前的過去,.
他說三年來我帶大這個小朋友,我看到他令我很驚嚇,他又會說謊話,他又懂得執導,他又懂得貪心,他無私自通,.
那我又好得給他多少呢?人總是有罪性,在基督徒裡面我們能夠去幾何無疑,我們能夠找到真正的道德標準,.
並且我們有能力,神給我們能力去實踐出來,這個是滿足的,能夠令我們滿足的,剛才我們說了四個方面的福,.
主要是人因為跟神有關係,生命得到轉化的福,跟著我們再看第四個福,非常重要,特別是說到人際之間有了生命帶來的震撼性很大,.
我們讀一讀第七節到第十節,這個講話,憐憫的人有福了,因為他們必蒙憐憫,清心的人有福了,因為他們必得見神,.
締造和平的人有福了,因為他們必稱為神的兒子,為易受迫害的人有福了,因為天國是他們的..
我們首先看第五個福,憐憫人的人有福了,憐憫人的人是看到別人的需要,當然在艱難的當中願意伸出援手,.
又或者如果一個人有錯的時候,不會執著執怪,有個連率在寬恕原諒,以致到不會只是執著定人死罪的一種態度,.
人最容易會這樣,這個錯那個錯,我們會執著,但是人往往是喜歡做審判,我們會執怪..
我舊裡認識一位伯伯,我有一次去醫院探望他,為什麼入到醫院呢?原來他心情非常之鬱結,.
他說有什麼事呢?原來他兒子忽然間有一天告訴他我要結婚,.
他說什麼?這麼大的事現在才告訴我?連談戀愛我都不知道,你還尊重我這個爸爸?.
於是激起來,心臟頂不住要入醫院,很鬱結,有個院務去慰問他,開解他,.
他說伯伯你有沒有試試寬恕,寬恕你的兒子?這個伯伯說在我人生的字典裡面沒有這句寬恕這件事,.
院務走了之後他自己想想,是啊,為什麼不聽院務的建議,寬恕自己的兒子,後來平靜下來,.
接受我兒子,多些體恤,人家犯錯的時候我們都是軟弱,你和我一樣的軟弱,.
基督徒特別是基督徒我們蒙受神的赦免的時候,我們特別體會到連恤的重要,.
逼人連恤的重要,我們更加願意連恤人,而這段經文說連恤的人有福了,因為他們必蒙連恤,.
我們是誰呢?我們經常定人罪做什麼?我們見到這個不對那個不對,自己最對,.
看我們不放下連恤的態度,因為這樣我們必蒙連恤,第六個的福,清心的人有福了,.
清心的人雖然很簡單,我們中文也叫清心,非常不簡單,覺得一個人很單純,忠誠於神,.
以至於他願意他的行事為人乾乾淨淨,清清潔潔,真是不簡單,聖經裡面怎麼說我們的人生,.
我們的心裡面,他說人的心比萬物都詭詐,誰能夠想得到呢?你跟我日日的思想,你如果真是世察的時候,.
有時你自己也嚇一驚,有時我們想到清清潔潔的生命我們都可以做的,不過可能是表面,.
我們恭恭敬敬,別人覺得大好人,這個好人沒什麼犯罪,但這裡說到他的心,他的心是要我們自己審查,.
我們外表一套,裡面是怎樣,只有自己知道,中國人經常說那句偽君子,.
聖經裡面經常說那句,假冒為善,我們經常說那些偽君子,或許我們經常都成為了偽君子,.
心裡面清潔,唯有神才做得到的事,唯有一個人坦蕩蕩,我裡面骯髒垃圾,主幫我,.
我能夠清心的人有福了,因為他們必得見神,這個很厲害,見神,神都見得到,當然不是說看到,.

$^{121}$而是很體現到神的存在,因為當我們這樣去生活的時候,神很喜悅,神的作為在我們的生命很活然很真實,.
我們和神能夠很親近,這是莫大的福份,我們的信仰是真確的,然後第七個福,.
他說締造和平的人有福了,締造和平說到人際關係,人往往有些嫌隙,令我們拒人於千里,.
這個我以後不要碰他,這個壞人弄得我這麼慘,締造和平是主動和人和好,.
可能是自己的錯,可能是別人的錯,可能是某些處境令我們彼此仇恨,.
締造和平不是代表抹殺對方的錯或自己的錯,而是將這些事情交給神讓神去處理,.
而自己學習放下怨恨去和人和好,這個生命的本質是非常寶貴的,.
我們人也去抗拒別人,我們不需要再和別人打甚麼交道,.
有時這些是和我們很為仇的,我們想放下,但突然又想起這個人,.
又再生氣一次,我不罵他也好,還和他和好,不要搞了,.
唯有我們真正理解締造和平的人有福,原來神有很多祝福在背後,.
我認識一位弟兄,有時我們的仇人或敵人,可能我們想抗拒的,.
很親的人也不一定,一位弟兄,我暫時叫他林弟兄,.
林弟兄告訴我一些他以往的往事,爸爸是一個非常嚴格的父親,.
很厲害的,但他讀書很不厲害,小時候考得不好,爸爸很生氣,.
要打他,如何打他呢?睡著了,揭開被子打,有時他在洗澡,.
走進浴室打他,他睡覺時握著拳頭,因為很害怕,很緊張,.
有一次聽到爸爸乘飛機,聽聞被人劫機,他內心說,.
好了,爸爸不再回來,當然爸爸回來,後來林弟兄心靈很大創傷,.
並且身體也有病患,曾經經歷過死亡的邊緣,一條命,.
很多的摧殘在世界裡面,後來跟著太太回教會,認識了主,.
生命改變,願意好,有勇氣地跟爸爸和好,跟爸爸傳福音,.
爸爸信了主,後來爸爸跟他說,我以前做的事錯了,對不起,和好,.
我們可以逃避這個締造和平,不過我們失去了人間很多的美善,.
締造和平的人有福,因為他們必稱為神的兒子,剛才我們說見神,.
神的兒子即是說似神,你直接有神那種形象在裡面,.
我們中國人也說,有其父必有其子,就是這個意思,.
神的兒子就是似神,一個人能夠似神,很大的福氣,第八個福,.
為易受迫害的人有福了,你一聽到這個福,可能很奇怪,.
跟之前的福很不同,很特別,這個福,第一你想想,.
之前的福全部是講人的態度,品格,這裡有沒有講他什麼品格,.
沒有,似乎不是,他講到他有什麼遭遇,這個遭遇是什麼,.
就是迫害,第二個問題,迫害也有福,怎麼會,奇怪,.
但原來最精彩的就是這個福,我們再想想,.
如果我們之前看什麼叫有福的人,有福的人是謙卑的人,.
他能夠在神的角度得到神的塗造,神的眷顧,.
這個福就是一個擁有這些福分的人,他會面對的結果,.
你說為什麼會,又好人又和好,你再想想,.
如果一個人願意尋求神的心意,神的道德標準成為他的標準,.
他的價值觀成為他的價值的時候,.

$^{161}$相信唯一這個世界獨一的真神,又願意去堅持這個理念,.
我們從世人就很脫節,你試想如果我持守我的價值觀,.
我會怎樣,神說不要講謊話,不講謊話,.
那你在辦公室裡面,有時你不多不少都要講謊話,.
人家就做這些事,你不做違法,不做一些不好的事,.
有一次我聽到一個弟兄說,他一班人出外公幹,.
半路就收工,但未到下班時間,他說我要回公司工作,.
其餘的人都走了回家,當是提早下班,他都不敢告訴別人,.
他只有靜悄悄回到寫字樓,乖乖地回去,.
到下班的時候才走,基督徒,如果你要堅持你的信仰,.
如果去到極權的國家,甚至你真的會受到迫害,.
甚至會危害你自己的生命,但你留意寶貴在於,.
一個人經歷了神的眷顧,生命的改變,.
正是付代價我仍要堅持我這個信仰,.
那就證明這個人找到人生真正的意義和目的價值,.
我不妥協,我不隨波逐流,有誰願意一些不值得的事付代價,.
沒有,唯有知道這個天國的福分是這麼真的時候,.
我不會放棄,甚至你要拿走我的生命,.
這種生命是很光彩,這裡說到福是什麼,.
他說因為天國是他們的,你要發覺第一個福也是天國是他們的,.
簡單來說是什麼意思,我估計天國不是一些遙不可及的東西,.
也不是要你做到很多很多你才有,當你單純謙卑的時候,.
你就進入神的角度,神就塗造你,並且勝在你一生,.
無論你被迫害,無論你逆境順境,.
神帶你同走,看到祂的活現,而當我們為祂犧牲,為祂付代價,.
將來有永恆的獎賞,基督徒的福分是能夠看到永恆的,.
這個福是莫大的福,因為天國是他們的,.
由開始謙卑,去到一個的生命,能夠和人的關係,能夠去順逆境,.
我能夠去滿足,這個是真正的幸福,這個福是超越環境,是內在,.
如果我們這樣看上來,這八個的福,是人生能夠給我們真正的福分,是滿足到我們,.
事實上是神一早給了我們,如果我們看這些的品格,.
神起初創造人的時候是按照他的形象做的,人是有美善的,.
但是當人離開他的時候,我們偏行幾路,我們自我中心,.
咦,我們沒有這些福,我們不覺得這些是福,我們反而覺得那些揹得緊的那些是福,.
剛才我一開始說的那些,以為那些是福,我們看不到原來,.
當我們尋求福樂的時候,我們脫離了神,.
聖經說我們的問題是什麼?世人都犯了罪,虧欠了神的榮耀,.
我們沒有神的形象,我們活在自己的罪惡當中,.
與神遠離,最終面對死亡,在這樣的情況之下,.
我們活在罪裡面,我們還去尋找我們的福樂,.

$^{201}$有時候基督徒都要小心,我們都走了這個方向,.
我們漠視了我們人的罪性,漠視了神的存在的時候,.
我們尋找的福樂只不過是暫短和外在,.
在我們的世界裡面,我們看不看到罪惡在周邊都發生,.
在校園裡面,青少年的校園有欺凌的事件,.
在辦公室裡面我們看到很多道德行為操守的一些亂奴,.
辦公室政治,科技越來越發達,.
我們現在有沒有很多時候我們找AI幫忙找答案,.
我發覺自從有了AI,我什麼都找AI,.
科技越來越發達的時候,人的道德的罪惡越來越大,.
最近有沒有聽到很多詐騙案,我聽到這些我很生氣,.
血汗錢沒有了,一班犯罪集團就打電話來,.
弄到人家很大的損失,我昨天和太太吃飯的時候,.
突然間收到一個電話,因為我做牧師我不能夠不聽電話,.
我不知道什麼事,有個人說Andy,叫到我的名字,.
Andy,然後跟我說一番,又不說自己的名字,.
我又不好意思,好像認不到人,後來我也迫不得已,.
你是哪位,你這樣也不認識我,那我想一想,.
我不認識,我說我不認識,這樣也不認識,.
接著又說明天十點鐘我約你,你來二樓見我,.
我說二樓,我說什麼二樓,二樓,我說什麼二樓,.
那我就知道,立刻收錢,那我太太在前面,.
我很緊張,你收錢,我做牧師不可以這樣,.
我怎知道誰是教友,最後我知道這個一定是,.
我才收錢,我試著想想,如果他用了AI,.
拿了我的聲音去騙我的親友,我當時擔心起上來,.
並且很激起豈有此理,這個世界我們不知道什麼是真假,.
你越厲害越發達,你犯罪的能力越高,.
這個世界人與人之間的紛爭,各與各之間的爭鬥,.
我們今天看到的戰爭的處境是什麼光景,.
我想不用大家說也知道,聖經說中了,.
聖經正是這樣的情況,耶穌基督來到這個罪惡的世界,.
祂憐憫我們,神差派祂來到這個世界,.
為了我們的罪,為了我們的死亡的問題,.
祂負上代價,溫溫柔柔的上十字架,.
聽從神的命令,一生公義的,祂沒有罪,.
以至到彰顯神的公義,死在十字架上,.
第三日後又復活,以至到祂贖去人的罪,.
罪的功架,聖經說是死,祂袋贖了我們,.
以我們不用死,以我們能夠有一個生命的轉機,.

$^{241}$能夠得到八個福,剛才所說的八個福,.
藉著相信耶穌基督,人生就是這樣,.
當我們願意去看到這個福的時候,.
認定耶穌基督是唯一的救主的時候,.
我們就會看到人可以有出路,.
我們能夠找到真正的福,這個福是唯有什麼?.
是唯有相信耶穌基督,將我們生命的主權,.
以往是這樣走,我現在交給神,.
由祂成為我們人生的主宰,帶領我們走前面的路,.
以至我們享受真正的福分,成為一個真正有福的人,.
願意在我們在座當中,如果你還未認識主的時候,.
我很邀請你們去想一想,.
耶穌基督是唯一能夠醫治我們破碎心靈,.
唯一醫治整個人類的疾苦,唯一的答案,.
去開出給我們,給我們有一個真正的福分,.
願意你能夠成為一個真正有福的人,.
我們一起去禱告,.
至於天父,將我們交代給你,.
我們願意我們的生命是被你所掌管,.
能夠成為一個真正有福的人,.
在座當中,如果你還未認識主的時候,.
我都去邀請你,你如果願意接受耶穌基督,.
相信耶穌基督的時候,.
讓祂帶領你的生命的時候,請你把手舉起來,.
有沒有這樣的朋友呢?.
有沒有這樣的朋友把手舉起來呢?.
我們在這個世界裡面,我們很多聽到驕傲的聲音,.
但是少聽到溫柔謙卑的言辭,.
我們看到人生很多人際關係的破裂,.
但是我們看不到人有勇氣去和好,去寬恕,.
唯有神是醫治我們,.
這招給我們有能力去實踐真正人生的美善,.
是償到永生的福分,.
在座當中,如果你是有感動,.
你覺得我以前都聽過主耶穌基督,.
我又願意在這一刻將我的生命交給主,.
請你都舉手,.
又或者有些在我們在座當中,.
以前你缺過志,.
你願意更加委身在基督立志,.

$^{281}$實踐神的真理,願意歸向神,服從神的,.
請你都舉手,我們都為你祈禱,.
還有沒有呢?.
好,我們一起禱告,.
慈悲的天父,我們多謝你,.
在這個世界裡面,都是紛亂,.
我們很需要你,.
唯有知道你是唯一拯救的主,.
去醫治我們破碎心靈的人,.
我們求主你去看顧,.
掌管我們的生命,.
讓我們真正成為一個有福的人,.
我們求主去帶領,.
剛才有位姐妹信了主,.
求主你去帶領她,.
保守她,帶領她一生,.
我們這樣的祈禱,.
是奉主耶穌基督的名,.
阿們..
\newpage



\section{約翰福音 13:1}
\label{sec:032053}
\textbf{愛到底的意思}~~~~余德淳博士
\newline
\newline
連結: \href{http://www.church.com.hk/acms/content.asp?site=cdc\&op=show\&type=product\&code=032053\&layout=sermon}{\texttt{032053}} ~~~~ 語音日期: 2026-04-26
\newline
\newline
\hyperref[sec:032052]{\small{< < < PREV SERMON < < <}}
~
\hyperref[sec:index_chronic]{\small{[返順時目]}}
~
\hyperref[sec:032054]{\small{> > > NEXT SERMON > > >}}
\newline
\newline
約翰福音 13:1
\newline
\begin{longtable}{cl}
\hline
\hline
章節 & 經文 (和合本修訂版)\\
\hline
13:1 & \begin{tabularx}{0.7\textwidth}{X} 逾越節以前，耶穌知道自己離世歸父的時候到了。他一向愛世間屬自己的人，就愛他們到底。 \end{tabularx} \\ \\
[1ex]
\hline
\hline
\end{longtable}
$^{1}$各位早晨,很高興見到你們,特別有些人問我記不記得叫甚麼名字,每次都令我有點擔心,因為你們實在太多人了,但今日的題目很好,我想你走的時候會有些感覺不同,.
剛才我喝杯奶茶時問剛才參加的人,我覺得這個題目很值得你聽完再聽,就是愛到底的意思,.
星期三晚我和大埔火災義工開會講哀傷輔導,我問他們你見到那些人時用甚麼辦法他們肯跟你傾談,原來雖然有心理輔導室,但約見的人是很少的,.
用甚麼方法令他們可以給人關心呢,方法就是一個服務叫按摩,很多人都有很多傷,甚至是身體上的痛苦,有些基督徒按摩師義務去做這些工作,.
而我很感動有些只是大學講師來的基督徒,他們在椅子入按摩室,那班人站在椅子上,有沒有興趣給我機會傾談,有些人是肯的,.
於是就這樣傾談,差不多機會率是十次有四次得手,六次是沒有人理會,但那四次都是五分鐘左右,於是我看到如果你愛一個人,關心一個人,你要接受到未必他領你情,.
不過你一定要有忍耐,可能有機會,但這個機會出現後又要想下次怎樣呢,我跟他說我以前做哀傷輔導講明你沒空不要做這些工作,.
因為這些工作如果一個很重好像大埔那些差不多是年半至五年那種叫創傷後遺,也有些是五年以上,我做精神科那些是十年,.
很多人經常想起舊時的事都是哭得很辛苦,我說如果你做義工關心人你要愛到底,差不多你預了排第一就是做這一範,其他的盡量不要應約,.
因為這一範要找你真的要靠你,所以可能你看到今日香港這個處境,事實上很多基督徒見證很好,他們不怕別人不答應他,總給機會別人,但要愛到底,要很久很久時間,.
我馬上想起這段聖經在楊潔篪13章第一節,聖經說他既然愛世間屬自己的人,即是屬耶穌的人,他愛他們到底,今日你坐在港福堂就看到這個意思,.
他們到底,譬如你看看隔壁所有人,有些樣子幾可愛,我沒有說有些不可愛,有些幾可愛,你想想公平,.
小朋友其實我做很多年兒童工作,小朋友最不值得就是有些人很受人歡迎,樣子很可愛,簡直是明星樣子,.
個個都爭著跟他拍照,但有些小朋友沒有,其實很多小朋友都很不開心,覺得為何我生出來他比我可愛,很不開心,特別家裡,.
如果你有兒子有女兒,通常女兒都比兒子靚,兒子經常覺得為何那麼醜,是否學到爸爸,但問題是等於個個都喜歡抱著妹妹,很不公平,.
但這句聖經很重要,愛他們,即是個個人都是一樣,是平等的,你在教會跟教會門口外面很大不同,.
在門口外面的時候,人會用他的主觀去看我喜歡那個多些,或者我喜歡找那個多些,但在教會裡面是個個人一樣,.
如果你第一次來教會,你放心,這個經驗是耶穌已經做了要求,就是人人都要平等,如果你有這個經驗,你放心,.
你曾經被人看不起,或者覺得很自卑,就是因為這句說話,你屬於耶穌的人,你被耶穌愛到底,愛到最尾,.
所以我覺得今日我們經常很擔心,人世間有沒有真愛到底的東西,原來真的找不到,但今日在教會,在聖經你找到,.
在講的時候我就想起一個遭遇,我很少講這個遭遇,就是我小學的時候在中文小學讀書,.
後來去了一間傳英文的聖Joseph那裡讀,傳之後都要講英文,連同學和同學聊天都不鼓勵講廣東話,.
我突然發現到,我上中一開始我懷疑自己有些心臟病,所以經常心跳不規則,原因是什麼呢?.
就是因為上生物課,我英文差,所以經常覺得老師叫我,就要講reproduction,我叫余德順,我真的被他嚇死,.
經常都reduction,我回家我和我爸爸媽媽講,我真的很害怕,經常心跳,數學都有時講reduction,.
我就想死了,我怎樣熬到上中五呢?我英文又差,讀書又差,經常都被人叫,有一次還自己站起來,同學在笑,什麼時候叫你?.
我想起我和我阿公講,我阿公是做南華早報的,英文很好,我和他講,我都很害怕,怎樣解決呢?.
我阿公想到個辦法,他真的很聰明,他買了一個禮物給我,這個禮物都幾貴的,是一個手提,手動的打字機,.
他說你學打字,你打字可以拆開英文字,那些發音你會警覺到,譬如reproduction,你試試拆開打字,一來你英文串字會準很多,.
二來你聽到這些音你會掌握得好些,他就買了打字機給我,從此我就是這樣練,其實這個經驗很難忘,.
我覺得我很幸福,有個年紀這麼大的公明心理,和他找到個辦法怎樣幫我媽媽幫我,他找到個辦法,.
直至到他臨死的時候,我記得在聖卜醫院去探望他,他還跟我說,我明天很快出院,我繼續和你研究英文,第二天他就死了,.
肺積水,死了,我都暈一件事,這個世界有沒有一個愛,一個愛是持續的,是不會完的,其實完是很痛的經驗,.
一個很愛你的人,愛到無微不至,幫你想怎樣解決你的憂愁,但突然間他不在了,那時十幾歲我就開始尋找,.
我怎樣找到一個我覺得可以長期和我的關係不會完的,所以後來我就看到聖經,基督教,如果不是我想我的遺憾感都很高,.
我們第一節聖經讓人回想十六章二十二節,就是這種愛的條件很厲害,這個條件是能夠解憂愁,解決憂愁,憂愁是很難解決的,憂愁是會增加的,.
在我們的右額頁,只要你重覆想憂愁的事,憂愁會增加到一個地步是離譜的憂愁,譬如假設你多了一些白頭髮,.
你一路重覆想,你很快會想到很老,再想下去就快死了,沒理由幾條頭髮就快死,你問非法佬,很久都未死,但他每天都見白頭髮,.
如果你自己單獨憂愁,你想都離譜,所以你說人的腦最不可靠就是在憂愁面前,你會自己嚇自己,所以病人,醫生最怕你憂愁,.
覺得很難去輔導幫你積極好,因為你憂愁,但你看聖經說你們現在也是憂愁,耶穌看穿了,門徒聽到他說快要離開世界了,憂愁,.

$^{41}$但接著那句說話就是但我要再見你們,再見你們,這句說話不是任何人沒資格說的,特別是老人家在醫院裡面,.
你有沒有寫包單會再見你的家人,但耶穌說我要再見你們,結果是什麼呢,你們的心就喜樂了,你就會開心,.
我們個個都試過接機的經驗,接機特別是你很想再見他,他一出來堆行李,你跟他笑到不得了,再見就開心了,.
但耶穌那種開心不止,這句說話沒有人能奪去的,沒有人能搶走的,沒有人能偷走你的,你說信耶穌是一個絕對式的東西,.
不是一些可能都幾穩陣的八成,希望九成,不是這些東西,是絕對式的,這個喜樂是沒有人能奪走的,果然你看看那些門徒,.
他見到耶穌復活之後,你看看整個基督教,就是因為那些見到耶穌復活的基督徒,後來你看看全世界傳到那裡,.
甚至那麼多人去宣教,就是因為那種喜樂他趕著要告訴很多人聽,好像今天講福堂,福音注入,你講給一些未聽過的人聽,.
你會聽一下這些東西,特別當你很失望的時候,你覺得這個世界還有沒有可能有希望呢?是有的,這個肯定的,你看到這段很好,.
不過當然很多人沒有信耶穌,是因為很理性,或者覺得我的根據不夠,或者你講見證一個因,你叫多一千個來,我就會信,.
或者你一會兒這樣講,你給我看多十年我就信,所以今天我準備福音注入,我想拿個證據,這個證據是大過剛才講的,.
這個證據是多少人呢?二百萬人,是多少年呢?如果你做research longitudinal,這個是四十年,四十年,.
二百萬人的證據,如果你聽完之後你覺得是真的,我覺得你不信耶穌,是你拒絕了你的智慧,是可以想到的,.
這個聖經大家研究一下,我們幾分鐘看出埃及記,二十三章,二十八至三十節,很長,不過我拆開來讀就不長,.
第一句,這個我就是上帝和摩西,以色列人講的,我要打發什麼呢?黃蜂,黃蜂飛在你面前,.
然後將三種人,希美人,迦南人,黑人拼出去,即是趕走三種人,因為這三種人當時在迦南地的時候是佔領了那裡的,.
用黃蜂趕走他們,我不在一年之內將他們從你面前拼出去,我不在一年之內,恐怕什麼呢?.
為什麼我不這麼快做這件事呢?我有原因的,恐怕這個地,這個地方成為了荒涼,一荒涼沒有人,.
人都趕走了,就有東西來了,就是野地的獸,那些野獸就會來了,來了之後怎樣呢?多到可以害你,多到可以害你,.
這個解釋得很好,為什麼我不馬上聽你祈禱,不馬上聽你說呢?因為這樣做法,如果馬上這樣做,就會馬上有個後遺症走出來,.
但是上帝說我有辦法,我的辦法就是我要漸漸地將他們從你面前拼出去,.
最後重要了,等到,等到,有個時間劑的,等到你的人數加多,就可以承受這個地方的威脅了,這句話很厲害,聽的人是多少呢?.
我算一算,聖經說有六十萬成年的男人,六十萬,就是未成年的,婦女的,弟弟的,加起來是超過二百萬,.
二百萬,這句話講的日子是什麼時候呢?是耶穌降生前,主前,一四四六年,一千四百四十六年,是出埃及之後的第三個月,.
研究那五樣東西,黃蜂,你生命裡面都可能有,就是你祈禱,你希望家裡那個人的病,或者那個經濟困難,.
或者那個心理有問題的人,上帝會幫你,不過未必用你想的明天,或者叫個醫生,或者叫個輔導員來幫你,.
他派另一個幫你的方式,可能你不是很習慣,這個黃蜂,我想大家都試過,你去醫院,那些醫生看到你望著他很懷疑,.
醫生第一句說話就說什麼?我不是新人,因為他很怕你來實習,可能還告訴你,我昨天剛剛被人治死了,嚇到你走也來不及,.
可能有時幫你那個不是你很心水,有時我們自己一個主觀覺得誰才能幫我,那個不行,但上帝可能真的找個你認為不行的,.
黃蜂,你覺得是骯髒,或者很難看,很惡看,但不在一年內,你要等,因為原因地會變成荒涼,很多野獸,.
即是你的環境我幫你改了,但如果你想現在改,可能更惡劣,真的,你看教育弟弟,如果你馬上幫他代他做他困難的事,.
他永遠都是找你,你更加證明他沒能力做事,同樣地如果你稱讚他很厲害,他以後怕不厲害,所以以後不要做那麼多事,.
為了保持繼續厲害,即是有很多方法不是方法,然後漸漸趕,這些漸漸你不喜歡的,譬如你說你老公很差,.
那個人說我幫你漸漸改,你知不知道講漸漸的人有個代價,代價是他不會離開你,因為他漸漸,.
所以證明他日日都在,他日日都在,他日日都動一點點,證明他日日都找過你,有時做基督徒常常感覺到上帝出手又慢又不見影,.
但如果你留意他做在你屋企的事,日日都有些事發生,是你遇到的人找你,跟你講一句,雖然不是多,但證明了上帝日日都在找資源在你屋企,.
是漸漸,一點都不是壞事,漸漸證明他沒有離開你的最好方法,等到你人數多,這個是真相,就是他要你都要負責,.
等到你,上帝沒有說不幫你,而是他想你準備好我的你,是等你,等到你人數增加,等到你配合到的人出現,這件事就會成功,.
我看公送給我打字機那種感覺,如果我不去練英文,我一世都會對reproduction搞不掂,直到我,等到我完全練得很熟,我就不怕聽有些發音誤會了我,.
這個是我要負責的,我不去練,這件事幫不到我的,你看到上帝的愛和寵愛的不同,寵愛就是你這麼煩我搞掂他幫你,.
上帝愛是無微不至的,他細到想出這一步會出現什麼會發生,不適合你的,不要出這一步,.

$^{81}$好了,不止,下一段更厲害,在申命記第七章二十二節,今日這兩段是我們福音主人一個重頭戲,你研究下,.
上帝怎樣愛我們呢?耶和華禮神必將這些國的民從你面前漸漸趕去,又講這句,.
你不可把他們速速滅盡,千祈不要,你千祈不要,講給你聽,你不要貪快,.
然後又一樣那個字,恐怕野地的獸多起來害你,這段申命記一查是幾時呢?.
是在上一段,四十年後講的,四十年後講的,是對即將入迦南,摩西對新一代的以色列人講的,.
這件事值得想一想,已經講了四十年前,為什麼還講呢?我記得我爸爸信了耶穌,.
跟我回教會崇拜,一聽到主席說今日經文選自救藥,我爸爸馬上跟我說,.
這麼舊不要拿出來講,他覺得這麼舊的事幾千年,講來沒什麼意思,.
但如果你問我們做輔導,我覺得是真的,有時我做過兩代輔導,三十多年前被我教過,.
就好像你最近的兒子,爸爸這麼壞,兒子也一樣這麼壞,其實找上一代去教下一代未必成功,.
因為他自己都不是很榜樣,就算你講給他爸爸聽,他媽媽聽,不代表他的兒子是聽到的,.
因為他父母傳錯消息下去,也可能是其身不正,令兒子覺得聽了也不是真的,他也做不到,.
所以上帝要容我世,再講給下一代,再講,你留意再講,申命記那種再講,有什麼好處呢?.
就是上帝的愛,沒有後顧之憂,他看到你今日的處境,還看到你遲些的處境,.
遲些你搞不定的時候,他仍然是保養期,還是對你家沒有放棄,還要救你下一代,.
因為我們做不到,我們做不到,你就看到人性的軟弱,很多人都很討厭人家說,講了多少次都還是錯,.
的確是的,人講了很多次都還是錯,但是你留意那幾句話,漸漸講,是學什麼呢?忍耐,.
最勁的生命力度就是你的耐性,耐性,忍耐,第二是什麼呢?警告,不可速速滅盡,.
不要私下加快你的時間,其實我覺得很多人都很聰明,閉了一件事就是太急,太急,.
這個急,誤事,譬如講話,快了講一些不對的東西,遲一些,然後後悔了,講錯話了,.
太急,行動太急,太辣,對方感覺到很難受,如果想一想他受不起,.
可不可以換個方式好一點呢?可能以前好很多,所以上帝經常叫我們不要太急,.
不要比我快,始終野獸多,會害你,愛是什麼呢?愛就是留意到你不要受傷害,.
所以讀詩篇二十三篇,我雖然行過死刃的幽谷,也不怕傷害,因為你和我同在,.
他都要陪你行,他知道你一個人行不行,你不行,所以信耶穌不是一個弱者的行為,.
信耶穌是一個有自知之明的行為,就是你一個人不行,你不可以靠你自己解決,.
最重點最大安慰就是頭幾個字,耶和華,你讀主義學就知道,耶和華為什麼呢?.
就是自有永有的神,在世界未有一些先,神已經存在,他自己存在,他永遠存在,.
所以整個空間那條線,耶和華就專插在歷史上,由頭到無窮盡這麼厚,.
這個字對我們是什麼意思呢?就是你放心,你看不到的未來,但你和神有好的關連,.
這個未來你放心,他怎樣你都是怎樣的,他怎樣你都是怎樣的,這種安慰真的太大了,.
你連幾十年來的創傷痛苦,在無限年裡面,幾十年不是很長的日子,眨眼就過了,.
但你以後最好的日子是會出現的,就是永生,這個觀念其實我想我們今日香港人,.
經常想的問題就是我們看事物很近期,遲些怎樣呢?所以很多年長者都很好,.
有人問他你健康怎樣,他通常回答現在沒事,遲些就很難說了,這個很對的,.
現在沒事,遲些就很難說了,誰能說遲些呢?沒有人有資格,反而這個長者很明白他自己,.
但耶穌超越了你憂慮的事,在約翰福音十六章三十三節,你看到這段聖經很厲害的,.
我將這件事講定給你聽,是要叫你們在我那裡有平安,講定給你聽,叫你們在哪裡有平安呢?.
什麼地址呢?在我裡面,這三個字很厲害,我裡面,很多人都找,找個名醫在哪裡呢?.
或者那本書在哪裡買呢?講到尾就是在耶穌裡面有平安,這個說話如果反轉講,.

$^{121}$耶穌外面有沒有平安呢?是沒有的,為什麼沒有呢?再詳細講下去就是在世上你有苦難,.
但你可以放心,我已經贏了這個世界了,能夠講這個說話的沒有人講到,可以贏了這個世界,.
這個苦難是什麼呢?這個苦難到盡頭是死亡,是死亡,但是這個盡頭我已經贏了,.
這個很厲害,你最怕第二天不知如何是好,第二天不知如何是好,原來是贏了,因為這個永遠,你仍然是永生裡面有份的,.
因為你在神裡面,裡面,看聖經,很多時候看約翰福音,譬如看葡萄樹,支子十五章,.
看羅馬書,很多時候都用一個字叫住,你住在神裡面,其實這個字很靚,這個解釋到就算基督徒,.
有些都很癱瘓的,有些都很沒信心的,就因為他不是住在神裡面,他有個名份,他是基督徒,.
但是他買了一間屋,但是他不去那裡睡,真的很浪費,他不去那裡睡,但是住在那裡,你不止有一間屋,.
你還住在那間屋裡面,那個才是真正快樂的基督徒,如果不是,你擁有了,但是你享受,你不去享受那件事,是浪費的,.
所以今天如果你是非基督徒,你留意一下,有些基督徒他經常都很快樂的,因為他曉得在哪裡有平安,他經常住在那裡,.
所以他經常都開開心心知道那裡是帶給他很穩陣,很幸福,我做這行,我最矜貴的字就是鎮定,.
鎮定是一個指標,看到一個人他的生命質素有多少,看你鎮定,.
為什麼我們不鎮定呢,因為不可知的太多了,譬如明天怎樣,下個禮拜怎樣,不知道,.
但是對不可知的事仍然鎮定,一個可能就是他傻的,所以他什麼都鎮定,一個可能就是他有些把炮,.
他有些肯定在那裡,譬如看新聞就知道了,每個一兩歲那些都鎮定的,只要媽媽抱著就可以了,.
他有把炮,他媽媽抱著就可以了,但是我們大人就麻煩了,我們的鎮定為什麼沒有了,後來變成了抑鬱,.
狂躁,精神病,因為鎮定你沒有把炮,所以到最後結果就是你失控,你控制不了,.
所以今天你想一下,我將這些事講給你聽,是叫你們在我這裡有平安,.
我試過一個很沒有平安的經驗,就是我在理工大學教,92,93,94,95,一年比一年差就是聲音,.
終於看了一個英國籍的西醫,一個大學的醫生,他就告訴我,我那些是哮喘類,但是就不能吃哮喘藥,.
你的情況很特別的,所以他說你的情況最沒有事的了,我說為什麼這樣,他說因為你可以做一些不講話的事就可以了,.
我都覺得很厲害,這樣就醫好了我,我不講話,你想一下如果我在大學教書,我申請不講話,.
我都不知道校長怎樣出糧給我,我覺得很辛苦,我辛苦到吃我的學生日,我久不久拿著個熱水壺去教書,.
一打個噴嚏,聲音很低,我就用最熱的滾水滾我的喉嚨,就出八成聲,但是很辛苦很辛苦,.
那些學生跟我說,醫生我不捨得你走,不過我們都看見你很辛苦,我說我一定走,.
上帝怎樣幫我呢,有人找我,在一間神學院找我,他跟我說我們需要一個講廣東話的心理學的兼職先生,.
你來教,我心想我在那邊辭工,你又馬上請我,兼職,兼職又不怕,我說我的聲音隨時沒有,.
他說不要緊,你一沒聲音就休息一下,等他們講一下,我說這樣也可以,當然最後一句很重要,.
我沒有糧出給你,我說okokok,我很內疚,如果你出糧給我,我想乘機看看我何時可以好,.
我真的要講話才知道我有沒有好法可能,如果我不做又不行,所以我說這個方法幾好,我就去教,.
那些短書中心又請我教心理學,拼拼埋埋都好像做了兩份兼職,兩份兼職,你說我是否吃虧,我吃虧,.
我的好朋友跟我說,你92年拿到香港第一屆優秀社工獎,跟住再回來,你以前讀過文憑理工教大學,.
真的很開心,但我問為何上帝拿走了我那件事,教了三年就沒有了,是上帝計劃行錯路,.
你給我這麼好的位讓我入大學讀研究院,跟住就玩完了,我開始都懷疑,但我不敢,.
我覺得總有個原因,但我不知道是什麼原因,我就學鎮定,我不知道什麼原因,為何我要行這樣的路,沒有工資,.
都不知道做到何時,沒有以前的事,以前努力的事全部沒有了,為何呢,但我學一件事就是鎮定,你相信他沒有騙你,.
我看聖經就看到這句,他對屬於他的人愛到底,後來他怎樣對我,我竟然第一日看到答案,.
在神學院教書時,他請我食午餐,每人派一張紙張,紙張很細,紙張用來做什麼,你不知道,.
因為他們沒有工友,所以每個人食那些渣那些骨要自己包,所以院長說自己包自己的骨,自己包就自己包,自己扔,.
那頓飯沒有sung ,其實很快食完,差不多兩分鐘食完,有人食一碗,但他笑,這班同學笑,.

$^{161}$這班同學不是開玩笑,有些是讀完建築,有些是讀完會計,有學位,賺的錢不少,但出來讀,可能出來人工低幾倍,.
起碼讀書時沒有收入,但怎會笑,我在理智上我明白,基督徒應該知足,我明白,我全部明白,但我見到真正的人笑,.
當然我第一個懷疑他傻,但問題他不是我以前病人的樣子,他不是傻,他不是傻,怎會食這些飯會笑,.
人工那麼差,將來他這樣會笑,這個喜樂是沒有人可以奪去,搶不到他,他要快樂時,你甚麼理由都搶不到他,.
我做了25年,這25年好像讀了6個學位,讓我看完一年又一年,沒有例外,沒有人後悔退學,沒有,.
為何基督徒是這樣,我終於覺得不用再解釋,上帝愛他愛到底,當這25年我發現一個離奇那麼開心的事,.
竟然我痊癒了,我沒有食任何藥,後來我打一次,我的聲音無低了,到今日我都沒事,我可以無咪都很大聲,就痊癒了,但我賺到甚麼呢,.
我賺到信心,我相信這個神是真的,我相信他的方法比我的方法好,後來我看到這25年我寫了十幾本書,.
EQ,由細路仔到大人,這十幾年就是中了醫生說的話,你不用一定要說話,你做些不用說話的就行,於是我就寫東西,.
他真的沒有騙我,我只是用筆,後來我成為了訓練機構,我有很多教材在自己十幾本書,如果不是走這條路,.
我沒理由有那麼多教材,我沒理由去著作,我沒理由看到那麼肯定,這個快樂不能奪去,這些全部都是真人,讓我見的,.
所以我相信一件事,上帝的智慧如果量度,是高過我們人一萬倍以上,都一定不止,如果由一個一萬倍以上的智慧來教你怎樣處事,.
一定好過你一萬倍以上,聽完這樣,是不是放心呢,你真的放心,你怎樣想辦法,都不夠他幫你找一個開始點,.
走下去就可以了,所以因為這個緣故,沒事了,我覺得沒事了,原來在他裡面有平安,一有平安就有鎮定,這個世界沒有困難可以嚇到你,.
信耶穌不是等於讀大學,只是一些知識,而是你經驗到那麼厲害,你完全放心,但我再補充一個證據,.
這個證據就是我阿公,他對我很訓練,我是他最大的孫,我媽媽是大女兒,有八兄弟姊妹,我有十四個表弟表妹,不過他未出世之前,.
只有我一個恩,他阿公生日走,端節中秋節都是我一個恩去,我阿公就預備了一個朋友陪我一起玩,這個朋友也是我阿公的鄰居,.
這個鄰居是開車,是一個司機,這個司機生了一個兒子,跟我一起童年,為什麼他預備這個小朋友跟我玩呢,原因就是這個小朋友是啞的,.
啞就經常笑,對我來說是沒有壓力的,只要我看著他就沒事了,我開始留意一下,所以你看到我每次說話都笑,就是我被他影響了五年,.
我跟這個小朋友經常都啞笑,我開始留意他笑什麼呢,我就學觀察,我問他媽媽他是不是這樣,他媽媽說不是是這樣,.
我發現原來我阿公已經預備了我,教弟弟做兒童輔導,又教我觀察做評估,就是我小時候學回來的,.
奇怪了,原來上帝可以說有些人說你做錯了啞的,不是的,他做來幫人的,但奇怪的就是這個啞的,為什麼他沒有自卑感呢,.
我也在懷疑,如果我是啞的,我真的不會去別人那邊喝,被人看著說不了話,但原來這個爸爸永遠跟這個兒子坐在旁邊都是拍著肩膀的,每逢他說話,爸爸就會笑,.
我就明白到世界最有威力的東西就是愛得令人有歸屬感,歸屬感,他是他的兒子,他有這樣的愛他的爸爸,所以他有勇氣望任何人他都開心的,.
因為他爸爸很疼他,他屬於他爸爸,後來我就明白到為什麼很多人有問題,就是他對家庭沒有歸屬感,他對婚姻沒有歸屬感,他對事業沒有歸屬感,他對將來沒有歸屬感,.
所以全部出事,就是他患得患失,見到很多東西都覺得令人眼淵,他沒有勇氣去面對新的人生,然後歸屬感是最厲害的,.
所以回看剛才第一句說話,耶穌已經說了這句說話在第一節裡面,十三章第一節裡面,他說他既然愛世間,屬自己的人就愛他到底,.
所以今天我們全部人坐在這裡,我們是屬於耶穌的人,他是愛我們到底,這個歸屬感是令我們鎮定,沒有東西是失去信心,.
世界怎樣亂也好,我們仍然放心活下去,很厲害,贏過你可以賺到多少錢,多少地位,甚至多少人支持你,.
但是一個歸屬感,你已經擁有了世界最大的東西,就是自有永有的神,祂是你的神,是你的神,這個歸屬感就是令你擁有一切,.
所以你今天離開港福堂的時候,你覺得自己原來富有到這樣,你擁有了一切,所以今天如果你是第一次來教會,或者來過很多次,.
但是你未覺得自己是基督徒,你只要做一件事,就是你歸屬於神,祂就愛你到底,.
如果你說這樣的條件,我不歸屬,祂不愛我,正如你要吃藥,人家給你一粒藥,你就是不吞,是治不了你的,.
但是藥是白白給你的,但是你不肯要,所以你還沒有好,所以有時神不是要求你很高,只是要你歸屬,原因是什麼呢?.
耶穌自己說,最後我們今天的聖經就是約翰福音三章十八節,信耶穌的人不被定罪,不信耶穌的人罪已經定了,.
因為那個人不信神獨生子的名,不信耶穌是神獨生子,就是這個,不是你有沒有犯法,作奸犯科的罪,.
原來那個罪是你不信耶穌是神獨生子,如果你是信,你就會歸屬於祂,祂就愛你到底,講完,.
所以希望今天這半小時,你看完覺得很真,我總結一次,耶穌請我們屬於祂,祂愛我們到底,祂能夠解決憂愁,.
祂的愛是無微不至,祂的愛是沒有後顧之憂,你是珍貴的,只要你歸屬祂,這個天國居留權就是你的了,.
希望今天大家聽完之後多點去想想,二百萬人四十年的證據,神怎樣愛,正如一樣是今天愛我們一樣,.

$^{201}$對付我們最慘的弱點就是我們心急,我們看不到神的計劃是這麼漂亮,還用自己的計劃來生活,.
所以很多人都很不開心,不如我們歸屬於祂,由祂帶領我們,所以你信了耶穌之後,你繼續回廣福堂,.
你聽到很多聖經的教訓,每一句話是舊的,但是有力量的,這個力量是幫我們繼續有平安有喜樂,一起祈禱,.
這個祈禱是代表我們去多謝上帝,所以我想你這一分鐘去感恩,很感恩,上帝怎樣漸漸地幫助你的家,.
祂從來沒有離開過你,同樣地你不要快速速要走快過上帝,祂這個漸漸地從來沒有離開過你,.
只是你要準備好,準備好祂來到你的生命那裡,你就得到這個快樂..
\newpage



\section{路加福音 15:1-32}
\label{sec:032054}
\textbf{理應歡喜}~~~~俞大鳴牧師
\newline
\newline
連結: \href{http://www.church.com.hk/acms/content.asp?site=cdc\&op=show\&type=product\&code=032054\&layout=sermon}{\texttt{032054}} ~~~~ 語音日期: 2026-04-26
\newline
\newline
\hyperref[sec:032053]{\small{< < < PREV SERMON < < <}}
~
\hyperref[sec:index_chronic]{\small{[返順時目]}}
~
\hyperref[sec:032055]{\small{> > > NEXT SERMON > > >}}
\newline
\newline
路加福音 15:1-32
\newline
\begin{longtable}{cl}
\hline
\hline
章節 & 經文 (和合本修訂版)\\
\hline
15:1 & \begin{tabularx}{0.7\textwidth}{X} 許多稅吏和罪人都挨近耶穌，要聽他講道。 \end{tabularx} \\ \\ \relax
15:2 & \begin{tabularx}{0.7\textwidth}{X} 法利賽人和文士私下議論說：「這個人接納罪人，又同他們吃飯。」 \end{tabularx} \\ \\ \relax
15:3 & \begin{tabularx}{0.7\textwidth}{X} 耶穌就用比喻對他們說： \end{tabularx} \\ \\ \relax
15:4 & \begin{tabularx}{0.7\textwidth}{X} 「你們中間誰有一百隻羊，失去其中的一隻，不把這九十九隻留在曠野，去找那失去的羊，直到找著呢？ \end{tabularx} \\ \\ \relax
15:5 & \begin{tabularx}{0.7\textwidth}{X} 找到了，他就歡歡喜喜地把羊扛在肩上。 \end{tabularx} \\ \\ \relax
15:6 & \begin{tabularx}{0.7\textwidth}{X} 他回到家裡，請朋友和鄰舍來，對他們說：『你們和我一同歡喜吧，我失去的羊已經找到了！』 \end{tabularx} \\ \\ \relax
15:7 & \begin{tabularx}{0.7\textwidth}{X} 我告訴你們，一個罪人悔改，在天上也要這樣為他歡喜，比為九十九個不用悔改的義人歡喜還大呢！」 \end{tabularx} \\ \\ \relax
15:8 & \begin{tabularx}{0.7\textwidth}{X} 「同樣，哪一個婦人有十塊錢，若失落一塊，不點上燈，打掃屋子，細細地找，直到找著呢？ \end{tabularx} \\ \\ \relax
15:9 & \begin{tabularx}{0.7\textwidth}{X} 找到了，她就請朋友和鄰舍來，對她們說：『你們和我一同歡喜吧，我失落的那塊錢已經找到了！』 \end{tabularx} \\ \\ \relax
15:10 & \begin{tabularx}{0.7\textwidth}{X} 我告訴你們，一個罪人悔改，神的使者也是這樣為他歡喜。」 \end{tabularx} \\ \\ \relax
15:11 & \begin{tabularx}{0.7\textwidth}{X} 耶穌又說：「一個人有兩個兒子。 \end{tabularx} \\ \\ \relax
15:12 & \begin{tabularx}{0.7\textwidth}{X} 小兒子對父親說：『父親，請你把我應得的家業分給我。』他父親就把財產分給他們。 \end{tabularx} \\ \\ \relax
15:13 & \begin{tabularx}{0.7\textwidth}{X} 過了不多幾天，小兒子把他一切所有的都收拾起來，往遠方去了。在那裡，他任意放蕩，浪費錢財。 \end{tabularx} \\ \\ \relax
15:14 & \begin{tabularx}{0.7\textwidth}{X} 他耗盡了一切所有的，又恰逢那地方有大饑荒，就窮困起來。 \end{tabularx} \\ \\ \relax
15:15 & \begin{tabularx}{0.7\textwidth}{X} 於是他去投靠當地的一個居民，那人打發他到田裡去放豬。 \end{tabularx} \\ \\ \relax
15:16 & \begin{tabularx}{0.7\textwidth}{X} 他恨不得拿豬所吃的豆莢充飢，也沒有人給他甚麼吃的。 \end{tabularx} \\ \\ \relax
15:17 & \begin{tabularx}{0.7\textwidth}{X} 他醒悟過來，就說：『我父親有多少雇工，糧食有餘，我倒在這裡餓死嗎？ \end{tabularx} \\ \\ \relax
15:18 & \begin{tabularx}{0.7\textwidth}{X} 我要起來，到我父親那裡去，對他說：父親！我得罪了天，又得罪了你， \end{tabularx} \\ \\ \relax
15:19 & \begin{tabularx}{0.7\textwidth}{X} 從今以後，我不配稱為你的兒子，把我當作一個雇工吧。』 \end{tabularx} \\ \\ \relax
15:20 & \begin{tabularx}{0.7\textwidth}{X} 於是他起來，往他父親那裡去。相離還遠，他父親看見，就動了慈心，跑去擁抱著他，連連親他。 \end{tabularx} \\ \\ \relax
15:21 & \begin{tabularx}{0.7\textwidth}{X} 兒子對他說：『父親！我得罪了天，又得罪了你，從今以後，我不配稱為你的兒子。』 \end{tabularx} \\ \\ \relax
15:22 & \begin{tabularx}{0.7\textwidth}{X} 父親卻吩咐僕人：『快把那上好的袍子拿出來給他穿，把戒指戴在他指頭上，把鞋穿在他腳上， \end{tabularx} \\ \\ \relax
15:23 & \begin{tabularx}{0.7\textwidth}{X} 把那肥牛犢牽來宰了，我們來吃喝慶祝； \end{tabularx} \\ \\ \relax
15:24 & \begin{tabularx}{0.7\textwidth}{X} 因為我這個兒子是死而復活，失而復得的。』他們就開始慶祝。 \end{tabularx} \\ \\ \relax
15:25 & \begin{tabularx}{0.7\textwidth}{X} 「那時，大兒子正在田裡。他回來，離家不遠時，聽見奏樂跳舞的聲音， \end{tabularx} \\ \\ \relax
15:26 & \begin{tabularx}{0.7\textwidth}{X} 就叫一個僮僕來，問是甚麼事。 \end{tabularx} \\ \\ \relax
15:27 & \begin{tabularx}{0.7\textwidth}{X} 僮僕對他說：『你弟弟回來了，你父親因為他無災無病地回來，把肥牛犢宰了。』 \end{tabularx} \\ \\ \relax
15:28 & \begin{tabularx}{0.7\textwidth}{X} 大兒子就生氣，不肯進去，他父親出來勸他。 \end{tabularx} \\ \\ \relax
15:29 & \begin{tabularx}{0.7\textwidth}{X} 他對父親說：『你看，我服侍你這麼多年，從來沒有違背過你的命令，而你從來沒有給我一隻小山羊，叫我和朋友們一同快樂。 \end{tabularx} \\ \\ \relax
15:30 & \begin{tabularx}{0.7\textwidth}{X} 但你這個兒子和娼妓吃光了你的財產，他一回來，你倒為他宰了肥牛犢。』 \end{tabularx} \\ \\ \relax
15:31 & \begin{tabularx}{0.7\textwidth}{X} 父親對他說：『兒啊！你常和我同在，我所有的一切都是你的； \end{tabularx} \\ \\ \relax
15:32 & \begin{tabularx}{0.7\textwidth}{X} 可是你這個弟弟是死而復活，失而復得的，所以我們理當歡喜慶祝。』」 \end{tabularx} \\ \\
[1ex]
\hline
\hline
\end{longtable}
$^{1}$多謝各位平安.
我們先一切祈禱.
讓我們可以一同來到這裡聚集見拜你.
由你的聖靈在我們當中.
帶領我們.
幫助我們時常在你的面前歡喜快樂.
奉耶穌基督的名祈禱.
阿們.
今天我們繼續看《路加福音》.
《與我的王同行一路》的系列.
今天我想先問問大家.
你今天來參加崇拜心情是怎樣的呢?.
有沒有人覺得歡喜快樂的請舉手.
歡喜快樂.
也有的.
歡喜快樂.
有沒有人覺得憂愁煩悶的.
今天來到很煩.
心裡受煩.
沒有嗎?.
是.
有沒有人覺得平常心.
又不是很歡喜.
又不是很受煩呢?.
有沒有?.
有沒有這樣?.
又剛剛去完旅行回來.
今天的講題叫做「理應歡喜」.
不是「我理他個理」.
不過跟「理」字是同音的.
我也鼓勵大家要歡喜.
我們有生命氣息.
有機會來到這裡跟大家一起參加崇拜.
真的應當歡歡喜喜.
我記得去年聖誕節前夕.
因為搶賭室進了醫院.
結果聖誕節也不能夠跟大家一起崇拜.
今年二月二十七號星期五.
又搶賭室進了醫院.
三月六號又不能夠跟大家一起崇拜.

$^{41}$只能夠參加網上崇拜.
這樣就不是很歡喜了.
在醫院裡面一個人.
這次被我單人房.
因為我體弱.
今天的講題.
「理應歡喜」.
就是來自今天的經文最後一節.
最後一句話.
所以我們理當歡喜慶祝.
我們是上帝的兒女.
上帝藉主耶穌救贖了我們.
使我們脫離罪惡.
使我們恩順正義.
又給了我們上帝兒女的身份.
又賜下聖靈與我們同在.
引導我們的人生.
保證我們永遠與上帝同在.
與主耶穌同享永生.
我們真是理當歡喜慶祝.
有這麼多好東西.
是不是理當歡喜慶祝呢?.
尤其是主日.
是主耶穌復活的日子.
如果有人不知道.
基督徒為什麼星期日崇拜.
我簡單講一講.
主要的原因不是由於主日不用上班.
大家有空就星期日崇拜.
不是這個原因.
由於主日是主耶穌復活的日子.
因此我們不再跟猶太人的舊約.
去守安息日.
而改為守主日.
在主日裡面崇拜.
有些基督徒不了解.
他們與上帝建立關係.
是透過主耶穌的血所納的新約.
以為既然是新舊約.
都是聖經上帝的啟示.

$^{81}$我們應該守一下舊約的規定.
應該守安息日為聖日.
有些人還搞了一個安息日會的宗派.
認為他們是好的基督徒.
遵守聖經的教導.
遵守聖經的教導固然好.
但是清楚認識信仰的基礎.
然後遵守聖經的教導.
就是更好.
又要一提思想主耶穌.
在前往耶路撒冷的路上.
所說的這三個比喻.
我們要思想比喻的內容.
以及我們從中學到的教訓.
《路加方》十五章記載了三個比喻.
耶穌一口氣說了三個比喻.
這三個比喻其實都有些共同點.
就是「失而復得」.
並且三個比喻的結尾.
都是主角和其他人.
一起歡歡喜喜地結束.
為何主耶穌這麼嘮叨呢?.
要說三次的嗎?.
十五章一到二節.
就交代了原因.
在第十四章裡.
我們看到有一大群的群眾.
跟著耶穌一起去耶路撒冷.
其中也有法利賽人.
更加有法利賽人的領袖.
大排筵席.
請耶穌吃飯.
今天的經文也有提及法利賽人.
不過第一節才提到.
「稅利和聚人」.
「稅利和聚人」不是兩種人.
而是從兩個角度去形容同一類的人.
「稅利」是形容這些人的工作.
「聚人」就是他們社會上的地位.
從前中國古代的社會.

$^{121}$也有將人分級.
分為士農工商.
商人的地位是最低的.
在耶穌時代的社會裡.
是羅馬政府掌權.
是羅馬的時代.
所有政府.
無論古今中外.
都一定要人民交稅.
你要交稅.
不需要覺得奇怪.
到處都是一樣的.
從來沒有一個政府或者政治家.
出錢來管理人民.
洋務一定是出自洋人身上.
不過.
羅馬政府的收稅辦法.
是非常聰明的.
他們不像我們現代人的政府.
請一批公務員來負責收稅.
而是將收稅的權.
用招標的形式拍賣.
給私人去經營的.
大概是先統計某個地方的人口.
然後公開招標.
有錢投標的人就估計一下.
你收到多少稅就出多少錢去招標.
羅馬政府就節省一批公務員.
何必養一批公務員那麼麻煩呢.
每個月要出糧.
可能每天都要出糧給他們.
很麻煩.
不僅如此.
一早就收了稅.
誰能夠投標就馬上給錢政府.
那些稅本就馬上植入國庫.
但是中了標負責收稅的人.
他是商人.
他們要賺錢的.
他們就欺負老百姓.

$^{161}$收高了一些稅.
於是那些百姓就叫那些稅來.
這些是罪人.
好像我們叫商人.
奸商來的.
無商不奸.
在猶太人的眼中.
稅吏就一定是罪人.
但是我們看到這句話.
有些稅吏.
他們很喜歡聽耶穌講道理.
可能耶穌喜歡用比喻.
講故事.
聽故事.
比喻其實是一個短故事.
但是藉著那個故事.
講出一些真理.
使聽到的人.
容易記得.
也容易明白.
我相信這些稅吏和罪人.
都一樣會請耶穌吃飯的.
而耶穌也不介意跟他們去吃飯的.
於是十五章二節就告訴我們.
「法利賽人和文士私下已論說.
這個人接納罪人.
又跟他們吃飯」.
這句是批評的說話.
其實這句批評說話.
表明那些人心裡不太開心.
不太滿意.
他們可能認為.
耶穌也是一個懂得律法的人.
壞壞的也是一個律法教師.
他們認為耶穌和他們是同等的.
起碼是同階層的人.
但是覺得耶穌令他們失望.
在他們的想法當中.
他們認為耶穌應該像他們這樣.
要鄙視罪人.

$^{201}$怎麼可以跟罪人去吃飯呢?.
這樣的事情怎會發生呢?.
耶穌才說了三個比喻.
來指責他們.
或者教導他們.
使他們要去反省自己的價值觀.
也要認識上帝的心.
在這三個比喻裡.
頭兩個比較短.
也比較簡單.
我簡單地說一說.
第一個比喻是針對男性的.
在希臘文裡面的名詞,代名詞,形容詞,動詞.
其實都有陽性,陰性的分別.
從字句裡面我們可以看到.
全部都是用陽性.
就是說那些男人.
這個比喻的主角.
就是一個牧羊人做主角.
主耶穌直接問在場的聽眾.
如果你們中間有人.
有一個人.
有一百隻羊.
走失了一隻.
會不把那九十九隻留在牧場.
去尋找那隻迷路的羊呢?.
這個反應大家要小心一點.
跟《和合本》有一點點差別.
它不是一個問題.
就是不是一個.
是一個問題.
但是其實他們有一個假設在後面.
就是假設他們一定會.
放下九十九隻去尋找迷失的那一隻.
那我就想一下這個問題.
如果是我.
我是不是一定會去尋找迷失的那一隻呢?.
我告訴大家,我不一定會去尋找的.
不一定.
因為.

$^{241}$一隻相對於九十九隻來說.
這個數量其實很小.
只有1\%.
好了,如果我放下九十九隻去尋找那一隻.
然後出現了一隻狼來吃了我十隻羊.
那我豈不是得不償失呢?.
我為甚麼這麼傻要去尋找那隻迷路的羊呢?.
除非我很有把握那九十九隻羊.
是絕對安全的.
我走開了,不要緊的.
牠們沒事的.
那我就會留下那九十九隻.
去尋找那隻迷路的羊.
但是.
這個牧羊人不是生長在我們這個時代.
猶太人的牧羊人.
對牠們的羊是很有情感的.
牠們每隻羊都會為牠們起個名字.
而羊也認識牠們主人的聲音.
所以.
牠們一般都一定會去尋找那隻迷路的羊.
並且當他們尋找到迷路的羊之後.
都一定會歡喜快樂.
猶太人真是聰明.
他們知道.
快樂一個人是不夠的.
一定要分享的.
於是就請了那些左鄰右舍.
一起來分享快樂.
第七節就是結論了.
迷失的羊.
就是代表罪人.
罪人迷失了.
或者好像夏娃那樣被人騙了.
但是假如他悔改.
悔轉.
不僅是尋找他的人.
尋找到他的人會快樂.
原來天助天使也一樣會快樂.
第八至十節是第二個比喻.

$^{281}$這次的主角換了.
路加醫生很注重男女平衡.
這次講的是女人.
第一個比喻是男人.
這個比喻是一個女人.
他說這個女人.
掉了.
不見了一個銀幣.
這個價值是比較高的.
因為一個銀幣.
是當時的工人要做一天工.
才賺得到的工資.
在這個比喻裡面.
失去的數量.
和第一個比喻比較.
高很多.
第一個比喻失去了百分之一.
但是這裡是失去了十分之一.
更加要賺錢.
這個人人都會賺.
這個女人也很正常.
快點燈.
周圍慢慢仔細地找.
找到了.
找到自然很歡喜.
正常人也會歡喜.
但是他們厲害的地方.
不是自己歡喜就算了.
請鄰居一起來.
一起歡喜快樂.
結論也提到.
天使也同樣會為一個罪人悔改而歡喜.
第三個比喻.
就是兩兄弟的比喻.
其中一個.
失落了.
如果從比例上來說.
這個是最高.
失落了二分之一.
兩個失去了一個.

$^{321}$這個比喻也是聖經裡面最長的比喻.
不過我相信大家都很熟悉這個.
浪子回頭的故事.
讓我們再一次.
詳細一點來思考這個比喻.
在這個比喻裡面有三個主角.
他們是一家人.
有一位慈愛的父親和兩個兒子.
大兒子是一個很乖很勤奮工作的人.
他在父親的爸爸的農場.
為爸爸工作.
時時都自己動手動腳.
也會打點工人.
去搞好農場.
小兒子就比較任性貪玩.
好食懶做.
不想做事的.
這個小兒子可能大概怕辛苦.
就不太想做農場工.
不過住在家裡.
不多不少.
都一定要幫幫忙.
因此這個小兒子就想離家出走.
不如我自己去外面闖闖天.
不要經常這麼辛苦做農場工.
於是有一天他就找爸爸.
「爸爸,不如這樣吧」.
「你將我應得的家業分給我吧」.
這個請求其實不太正確.
事實上不單是猶太人.
我相信華人都是一樣的.
產業應該在父親死了之後才分.
爸爸還活著的時候.
有什麼理由叫爸爸分家產.
難道想望爸爸早點死嗎?.
我相信我們中間有兒女弟兄姊妹.
如果你的兒女走來跟你說.
「不如你將你的家產分給我吧」.
我不知道你有什麼感覺.
會不會答應他的要求呢?.

$^{361}$我也有一個女兒.
如果我女兒要這樣跟我要求.
我說「我可以答應的」.
「不過房子不能給你,我要住的」.
「車不能給你,我要開的」.
「錢我沒有多少」.
「結果沒有東西可以給你」.
在耶穌的比喻裡.
這位慈愛的爸爸一口答應.
結果我相信除了房屋和田產之外.
他會將三分一的財產.
都分給小兒子.
為什麼分得三分一呢?.
因為猶太人的規矩.
大兒子是應該得雙份的.
所以他的財產要分成三份.
雙份給大兒子.
一份給小兒子.
在第十二節說.
「他父親就把財產分給他們」.
不是分給他,我們要留意.
是分給他們.
就是說不只是小兒子有財產分得到.
大兒子也分得到.
大兒子也得到他的那份.
小兒子拿到錢之後.
就離鄉別井.
去外地過一些醉生夢死的生活.
但是這樣他很快就變得一無所有.
當他走投無路.
機場碌碌的時候.
要吃東西怎麼辦呢?.
你打工吧.
你不打工哪有東西吃呢?.
那就要找工作了.
結果由於他在農場長大.
唯一的選擇就是找些農場工工作.
有一個農場老闆肯聘請他.
不過叫他去看豬.
如果我們了解猶太人的法律.

$^{401}$上帝定了豬是屬於不潔淨的動物.
不潔淨的家屬.
因此猶太人絕對不會養豬.
這個老闆應該不是猶太人.
怎麼可能養豬呢?.
天天與豬為伍.
全部都被他弄得不潔淨.
是不可以的.
這個小孩這個時候就覺得很醜.
他也回想到自己的人生.
為什麼會弄得這麼悽慘呢?.
年輕人可能有夢想.
要做一個比較特別的人.
要做一些比較特別的事情.
要出人頭地.
要與別不同.
這個小孩可能也曾經是這樣想過.
可能他也是一個相當有理想的人.
他不想再受父親和兄長的囉嗦.
於是就自己出去闖天下.
想幹一番事業.
想解放自己.
不過可能他社會經驗不足.
又或者遇上一些不務正業的朋友.
於是就開始過放蕩的生活.
令自己現在可以說是山窮水盡.
才懂得坐下來反省.
在這個時候.
上帝憐憫他.
讓他怎樣呢?.
「醒悟過來」.
他醒悟到些什麼呢?.
這位浪子想起自己的父親.
是有田有地的.
他還可以請傭人.
他最重要的一個「醒悟」.
就是認識到自己是一個罪人.
他也認識他爸爸.
他深信他爸爸是愛他的.
他就準備回家.

$^{441}$向爸爸認錯.
他心裡準備好台詞.
見到爸爸就要對他說.
「父親,我得罪了天,又得罪了你」.
「從今以後,我不配稱為你的兒子」.
「你把我當作一個故宮吧」.
於是他就開始回家的路程.
當他回到家鄉的時候.
他爸爸的行動.
竟然完全是出乎意料之外.
他準備好台詞.
最後一句「把我當作一個故宮吧」.
都沒有說出口.
爸爸已經忙著疼愛他.
然後吩咐僕人.
預備一個盛大的宴會.
歡迎他兒子回家.
這位父親.
就是代表我們的天父.
我們可以從他的行為上.
來認識上帝.
可能沒有人的眼睛.
比這位父親的眼睛更加明亮.
當他看到他的兒子離家還遠.
父親已經看到他.
認出這個人是他的兒子.
沒有人比這位父親的心更加充滿愛戀.
沒有人的步伐跑得比這位父親更快.
沒有人的擁抱比這位父親擁抱更加實.
也沒有人的親吻比這位父親的親吻更加真.
他沒有隱藏他對自己兒子的愛.
他也沒有責備這個兒子.
他不介意他身上可能還有一股豬味.
當一個人承認自己是罪人.
願意悔改.
願意回歸天家的時候.
上帝是一定會接納我們.
我們都是曾經得罪過天父的罪人.
但天父卻是會因主耶穌的死.
借免我們的罪.

$^{481}$而當我們願意歸家的時候.
天父也都一定會很喜樂.
天上的天使也都會一同喜樂.
這個比喻來到這裡.
如果只是說浪子回頭.
就可以在二十四節收場.
但這個比喻還沒有結束.
還有二十五節到三十二節.
大兒子那一幕.
這位大兒子和他父親對弟弟回來的反應.
是完全不同的.
他一知道好吃懶的弟弟回來.
他父親竟然為他開派對.
大事慶祝.
他的反應是怎樣?.
很開心地一起進去嗎?.
不是的.
他很不高興.
不肯進去.
二十八節說.
大兒子就生氣.
很生氣.
其實誰應該生氣?.
應該生氣的應該是父親.
但父親不生氣.
反而是這個哥哥非常生氣.
為什麼呢?.
為什麼他這麼生氣呢?.
因為他覺得父親偏心.
他很偏心.
愛弟弟多於愛自己.
當父親出去勸他的時候.
他跟父親說.
你看看你多麼偏心.
我服侍你這麼多年.
從來沒有違背過你的命令.
而你從來沒有給我一隻小山羊.
叫我和我的朋友開心.
但你這個臭小子回來.
他跟娼妓吃光了你的財產.

$^{521}$他一回來你都為他宰了肥牛毒.
這個大兒子.
這個大兒子滿心怨氣.
跟父親說出他的心聲.
他自覺他是一個孝順兒子.
父親我有什麼做得不好?.
你叫我做什麼我經常幫你.
從來都不違背你的說話.
但你都不欣賞我.
連小羊也不給我一隻.
讓我和朋友開心.
你這個臭小子回來.
你宰了一隻小牛毒慶祝.
當時牛的價值比羊高很多.
說來說去就是父親偏心.
其實他似乎忘記了一件事.
他忘記了父親已經把他的產業分了給他.
並且他得到雙份.
他自己不去享用.
為什麼要怪父親?.
他自己不享用是他的問題.
父親沒有禁止他享用.
你喜歡開巴士就開吧.
但他的說話.
反映出他是一個自我中心缺乏親情的人.
他不覺得能夠時時陪伴父親是他的福氣.
他也不紀念兄弟之間的親情.
認為這個臭小子回來是為他爭寵.
分分鐘還用了他的財產.
不過父親對他說.
兒子,兒啊.
我時時跟你同在.
我所有的都是你的.
不過你這個弟弟.
是死而復活.
失而復得.
我們是理當.
應該的.
我們是應該來歡喜慶祝的.
大家說這個父親說的話合不合理呢?.

$^{561}$為什麼耶穌要說這個比喻呢?.
我們回頭看十五章一節.
我們看到稅吏和罪人.
這些法利賽人和一般猶太人.
對稅吏的看法就認為稅吏就是罪人.
小孩子就代表這些罪人.
但十五章二節我們又看到.
法利賽人和經學教師是埋怨.
這個人接納罪人.
又跟他們吃飯.
大兒子就是指這些法利賽人和經學教師.
他們雖然是熟讀聖經.
經常禁食土國與天父接近.
留在聖殿當中.
在上帝的家中來事奉.
但是他們其實不是很了解天父的心.
他們太過聚焦於自己.
怎樣去事奉上帝.
自己怎樣去為上帝工作.
就好像大兒子很努力地在父親的田間工作.
不過他們缺乏憐憫之心.
對弟弟回家一點也不開心.
還認為父親這樣為這個罪人弟弟開派對是不對的.
所以耶穌一連說三個比喻.
去指責和教導法利賽人和問事.
首先是他們的價值觀.
頭兩個比喻都是說人失去了財物.
會努力去尋找.
這些就是一般人的價值觀.
錢,財物.
很緊張.
但是對失落了的罪人.
是他的事.
莫不關心.
甚至罪人迴轉也不會很歡迎.
但是耶穌不同.
耶穌是去尋找,去接觸罪人.
幫助他們認識上帝.
有些法利賽人還批評耶穌.
可見他們的價值觀有很大的偏差.

$^{601}$因為教導人遵行律法.
沒有教導人.
要用心來事奉上帝.
要用心來愛倫社.
第二方面.
法利賽人和問事.
要去認識天父慈憐的心.
上帝當然也愛那些法利賽人和問事.
因為他們真的很尊敬上帝.
認真查考聖經.
也教導人遵行律法.
只不過他們缺乏憐憫.
不愛惜那些失落的同胞.
只懂得歧視他們.
罪人來了.
不同他們交往還批評耶穌.
好去接納罪人.
雖然主要是針對問事和法利賽人.
但是在場也一樣會有遂利.
所以這些比喻.
我相信耶穌也會盼望他們.
懂得「醒悟」.
如果他們「醒悟」.
他們懂得回轉悔改.
上帝會歡迎他們.
和他們一起歡喜快樂.
我們已經得到天父的恩典.
但是我們不要自得其樂.
我們要體貼天父的心意.
天父歡喜在罪人悔改.
我們是否願意用盡各樣方法.
去幫助罪人悔改呢?.
求主幫助我們享受天父的愛.
也體貼天父的慈心.
不要再站在家外.
快點進去.
和天父和眾天使.
為了罪人的悔改回轉一起歡喜快樂.
理當如此.
應該這樣做.

$^{641}$讓我們一起祈禱.
因為你讓我們成為你的兒女.
可以享受你的同在.
與你同行的生活.
求你幫助我們有你的愛心,智慧.
與我們的朋友和鄰居.
一起分享信主的喜樂.
奉耶穌基督的名祈禱.
阿們.
\newpage



\section{歌羅西書 2:16-23}
\label{sec:032055}
\textbf{教會尋求社會認同的危機}~~~~黃紹權牧師
\newline
\newline
連結: \href{http://www.church.com.hk/acms/content.asp?site=cdc\&op=show\&type=product\&code=032055\&layout=sermon}{\texttt{032055}} ~~~~ 語音日期: 2026-04-26
\newline
\newline
\hyperref[sec:032054]{\small{< < < PREV SERMON < < <}}
~
\hyperref[sec:index_chronic]{\small{[返順時目]}}
~
\hyperref[sec:032056]{\small{> > > NEXT SERMON > > >}}
\newline
\newline
歌羅西書 2:16-23
\newline
\begin{longtable}{cl}
\hline
\hline
章節 & 經文 (和合本修訂版)\\
\hline
2:16 & \begin{tabularx}{0.7\textwidth}{X} 所以，不要讓任何人在飲食上，或節期、初一、安息日等事上評斷你們。 \end{tabularx} \\ \\ \relax
2:17 & \begin{tabularx}{0.7\textwidth}{X} 這些原是未來的事的影子，真體卻是屬基督的。 \end{tabularx} \\ \\ \relax
2:18 & \begin{tabularx}{0.7\textwidth}{X} 不要讓人藉著故作謙虛和敬拜天使奪去你們的獎賞。這等人拘泥在所見過的幻象，隨著自己的慾望無故地自高自大， \end{tabularx} \\ \\ \relax
2:19 & \begin{tabularx}{0.7\textwidth}{X} 不緊隨元首；其實，由於他全身藉著關節筋絡才得到滋養，互相聯絡，靠神所賜的成長而成長。 \end{tabularx} \\ \\ \relax
2:20 & \begin{tabularx}{0.7\textwidth}{X} 既然你們與基督同死而脫離了世上粗淺的學說，為甚麼仍像生活在世俗中一樣，去服從那「不可拿、不可嘗、不可摸」等類的規條呢？ \end{tabularx} \\ \\ \relax
2:21 & \begin{tabularx}{0.7\textwidth}{X} 見上節 \end{tabularx} \\ \\ \relax
2:22 & \begin{tabularx}{0.7\textwidth}{X} 這些都是根據人的命令和教導，論到這一切都是一經使用就都敗壞了。 \end{tabularx} \\ \\ \relax
2:23 & \begin{tabularx}{0.7\textwidth}{X} 這些規條使人徒有智慧之名，用私意崇拜，自表謙卑，苦待己身，其實在克制肉體的情慾上毫無功效。 \end{tabularx} \\ \\
[1ex]
\hline
\hline
\end{longtable}
$^{1}$好 等一會讓我們在聆聽主的說話之前.
我們一起同心低頭禱告.
把我們的心思意念交給我們主愛的神.
父神多謝你讓我們再次打開你的啟示.
讓我們能夠在其中明白上主你的心思.
也都讓我們在明白過後.
曉得怎樣真是按照你的規劃.
來去建立一間合神心意的教會.
也都成為一間有真正與神理關係的教會.
能夠彼此建立在真理的裡面.
我們為此再次獻上禱告.
祈求次智慧的聖靈.
讓我們能夠在每一天的生命生活當中.
不斷地去反思和觀察.
好使我們能夠更貼近.
上帝你要我們在每個時刻裡面.
當作的回應.
讓恰當的回應能夠成為這個世界.
更多屬主你的人.
能夠可以知道自己的罪.
明白上主的創始成中的身份.
重新將生命投靠在上帝你當中.
將人的主權.
再次與上帝你的主權.
能夠在其中蒙受主的引導.
禱告奉靠主耶穌基督得勝的明智教.
阿們.
在這幾種思想裡面.
其實我可以大膽些按照聖經的說法.
都是屬於異端教會的思維.
它比靈恩派的神學更加毒害.
主義的思維.
更加激烈.
比靈恩派的神學更加毒害.
主要的一個原因.
就是它對於這個世界.
和神的啟示.
作了一個很大的妥協的活動.
這一點我相信我今天講完這番話.
網上應該極多人會來找我找我晦氣的了.

$^{41}$不過不要緊.
但是怎樣啟示呢.
讓我們今天透過歌羅西書的第二章.
來去更加的進一步去思想.
不過.
透過在過去由2019年到2023年.
或者2022年.
所容許的這個所謂新冠狀病毒的廣泛流行.
上帝將這一類型的異端教會的真面目.
將它揭開了.
所以在近這一年半載之間.
我們會留意到.
有極多的這類型的教會.
把持著所謂使命人生.
使命教會.
或者所謂的社區教會.
這些種種的思維逐一爆破.
怎樣爆破.
無論在神學上的爆破.
以及財政上的爆破.
你可以看到這類型的教會逐一的倒閉.
而離開教會的信徒也都越來越多.
因為他看到一個真正的現象.
就是追求上主.
不是只是要在這些運動裡面參與其中.
才是合神的心意.
我剛剛所講的這番話.
其實在聖經裡面.
有一個同樣有同樣觀點的人物.
將他很清楚的否白.
這個就是保羅.
他亦都點出了給當時他寫這封書信.
俄羅斯書給的俄羅斯教會一個很重要的信息.
就是倘若教會如果要和這個社會有所結合的話.
其實是對於教會一個最大的危機.
亦都倒轉過來去看.
倘若教會是要尋求得到社會的認同的話.
這個亦都將教會帶進去一個教會離開神的危險.
所以我相信上帝是祝福加拿大的教會.
是透過這個所謂C9法案.

$^{81}$來將聖經裡面的內容宣講定性為歧視的項目之一.
是對今天教會一個很重要的警號.
告訴信徒聽究竟我們對上主的說話有多大的尊重.
所以上帝沒有在這個過程裡面去將C9法案廢除.
或者不讓它通過.
就好像早二十年前左右將同性戀的婚姻關係.
定為合法同樣的情況一樣.
其實對於教會是一次又一次的提醒信徒.
你們究竟對上帝的說話有多尊重呢?.
是不是你只要怕法律會說擔心被人指控為歧視.
你就對聖經裡面的主張或者原則.
將它摒除或者忍而不談呢?.
這一點正正是我看到上帝愛今天加拿大教會的原因.
更加愛教會的原因就是使加拿大的左膠政府.
更加成為一個majority government.
它喜歡做什麼就做什麼的時候.
加拿大教會所面對的事情應該會一浪接一浪來更多.
你怎樣去承接呢?.
是不是要和這個世界和這個社會再次尋求他們的認同.
而我們要去遷就它而使教會存在.
這樣就能夠見證上主呢?.
這一連串我的問題都在今天這篇聖經裡面.
保羅上希望逐一的去處理.
尤其是在當時哥羅西這個城市.
是正正在這個小亞細亞或者我們真正說亞洲的地方.
他們在裡面是否要存在當中的教會也要和社會尋求認同呢?.
而這個認同會帶來什麼影響呢?.
保羅就在這段聖經介紹給我們知道.
所以今天我講到的題目就是教會尋求社會認同的危機.
The Trap of Seeking Recognition from the World.
我們一起打開這段聖經.
歌羅西書第二章十六至二十三節.
我們再次的一起去讀一下這段聖經.
很想我們更加熟悉經文的內容之後.
能夠更加準確的去知道聖經裡面所介紹的是什麼內容.
好,我們一起打開之後一起讀.
故意謙虛和敬拜天使,就奪去你們的獎賞.
且等人驅離在所見過的,隨著自己的慾心.
無故地自高自大,不辭定元首.
全身既然靠著他,跟著得以相助聯絡.

$^{121}$就因神大德長進,基督裡的新生活.
你們若是與基督同死,脫離了世上的小學.
為什麼形象在世俗中活著.
服從那不可拿,不可嘗,不可摸等類的規條呢?.
這都是照人所吩咐,所教導的,說道這一切.
正用的時候就都敗壞了.
這些規條使人徒有自慰之名.
用施義崇拜,自表謙卑,苦待己身.
其實在克制肉體的程慾上是毫無功效.
在讀經時你會留意到.
我讀了和合本的一個標題.
《基督裡的新生活》.
正常我不讀的,為什麼我會讀呢?.
因為很多信徒當看見這段聖經的時候.
見到《基督裡的新生活》.
他就將這段聖經完全抽了出來.
覺得只要模仿,我就等於一個基督徒.
但是這段聖經裡面所講的究竟是一個什麼樣的新生活呢?.
是不是一個靠模仿裡面的事情?.
譬如有些信徒讀經不仔細去讀的.
讀那些服從不可,拿不可,嘗不可摸等類的規條呢?.
他以為真的要這樣做.
我們仔細去看,究竟裡面真的講《基督裡的新生活》是什麼?.
不過提起前文,究竟這段聖經講些什麼呢?.
哥羅西教會正如我先前所講.
是一間處於古近東亞洲地區的正正中心點的交通要塞.
因此它的商業發展是非常之蓬勃的.
保羅提醒一間處於一個大都會交通要塞的教會的重信徒.
你們可能會享受到最漂亮,最先進的一些角色,形式的社會,世界所得到的東西.
不過請留意,有很多東西在背後可能只是造假去吸引你們.
吸引你們最主要的原因是什麼?.
要使你們覺得不需要太過提高警覺.
因為你要求的可能就是世界上的一些享受.
所以今天的人很多時候看到的人生最重要的是開心.
這個觀念不是現在的人有的.
在古希臘的時候所有哲學家都是圍繞著這個問題去討論的.
究竟什麼叫開心?.
怎樣尋求到開心?.
現在的人沒什麼的.
滿足個人意願就是開心.

$^{161}$飲飲食食,有得玩就叫開心.
但是做人是不是那個原則,甚至作為信徒要想多一層.
人的意義是不是就是這樣?.
保羅所以在先前要他們提醒到留心他們所處的世界.
沒錯,可能是一個非常讓人很開心的環境.
但是你們要更加格外留神.
免得在一些事情上很容易將我們的心去被擄去.
這些全部在前文講過的,我不再重複了.
在這些種種的環境裡面.
真正相信耶穌基督的人.
最大的與不信主的人的不同之處就是.
人透過相信耶穌基督的死與及復活.
我們可以透過一件事.
可以將罪的債得到耶穌基督.
去將他付清,人再沒有機會來指控我們.
或者追溯我們在罪裡面的責任.
因為那個由耶穌基督去承擔了.
這個指的是在屬靈的角度裡面.
所以人只要他相信耶穌基督之後.
沒有人會去能夠再追我們以前是一個罪人身份.
而要陷我們於不義的狀態.
但是來到今天這段聖經的時候.
我們就會發現到一個現象.
使得我們墮落跌倒的人.
他使用的方式可以說是無所不用其技.
他竭力的去嘗試用很多方式捆綁我們走下去.
可以指控我們的條件當中.
而這個正正就是保羅對一間.
其實信耶穌上面已經很不錯的.
哥羅西教會作出的提醒.
我分兩部份和大家一起看.
第一部份就是從第十六節至第十九節.
第二章十六至十九節.
保羅要真的提醒哥羅西教會.
你們真的要留心.
一些事情越使你快樂的.
就越使你容易墮落.
所以我的附題叫做.
越快樂就越快樂.
那個樂就是跌下來的那個樂.

$^{201}$我們看看這段聖經說什麼.
十六節至十九節.
所以不拘在飲食上或節期越縮.
按昔日都不可讓人論斷你們.
這些原是後事的影兒.
那形體卻是基督.
不可讓人因著故意謙虛和敬拜天使.
就奪去你們的獎賞.
這等人驅離在所見過的.
隨著自己的慾心.
無故地自高自大.
不辭定元首.
全身既然靠著他.
跟接得以相助聯絡.
就因神大德長進.
這裡說什麼呢?.
這裡說的所有都是在教會裡的人.
所會聽到的主張.
尤其是基督的信仰.
是源自一位猶太人的耶穌基督.
很多時候在這個世界上.
有很多叫你快樂的東西.
但越不健康的不是我說的.
很多人都有認同的.
最簡單的例子就是.
很多人說如果你不開心.
吃點甜食吧.
我不知道你喜不喜歡吃甜食.
我爸爸剛剛過來參加我女兒的婚禮.
之後我跟他去旅行.
我怕他肚子餓.
我就問他你要不要吃東西.
我看到有東西吃.
我就問他你要不要吃東西.
他無端端說了一句話.
你經常吃東西嗎?.
然後我看到有些.
跟他逛超級市場.
他們要巡視業務.
即是這邊的業務.

$^{241}$不是他真正的業務.
是別人的業務.
不過巡視一下.
但他去超級市場.
跟他逛超級市場的時候.
有很多甜品.
我就問他吃不吃?.
你經常吃甜食嗎?.
其實這個世界.
越叫你去以為開心的東西.
就越要留心.
因為很多時候是一些.
很不健康的東西.
不要說造假的那些.
很多假的東西.
現在雞蛋也是假的.
你看到一隻煞蛋也是假的.
即是有蛋黃有蛋白包著的.
有沒有吃過?.
在酒店.
你不要一定以為是真的.
是小啤出來的.
如何怎樣?.
很多時候很多人以為.
只要越容易得到的東西.
越以為很好吃的東西.
其實那些東西.
使你更健康更滿足.
其實那些東西可能對你身體.
是最大破壞的東西.
最近有一位姐妹告訴我.
她什麼都不能吃.
我真的為她難過.
我說為什麼?.
因為她有病.
那些病.
吃多一點又會超標.
吃少一點又會暈.
真的很慘.
我自己也知道有些情況.

$^{281}$譬如糖尿病就是一個典型例子.
你不吃糖.
你又會暈.
吃少了糖.
你又會暈.
吃多了又會超標.
又會暈.
在這樣的情況之下.
很多時候我們.
在這個世界吃的東西.
其實都是一些味精糖精.
和一些香精的東西.
似乎在麻醉我們那個.
即是叫食覺.
味覺.
在屬靈上也有很多這些東西.
我曾經聽過一位老人家說.
他說在人生當中.
他走過了七十多年之後.
他說我最享受吃的一件事.
就是一碗由我媽媽煲的白粥.
白粥是最好吃的東西?.
他說因為這個世界.
我吃過這麼多東西之後.
我發覺我開始不知道什麼是其他味道.
只知道那碗白粥是最好吃的味道.
西人墮落的天使.
與及他的爪牙.
可以用各種宗教活動形式.
來去透過二次創作.
所搞出來的花神信仰來欺騙信徒.
有什麼花神信仰.
他們可以搞得出來.
就在第十六節那裡說了.
所以不拘在飲食上.
節期 月朔 安息日.
都不可讓人論斷你們.
這句話說的是什麼?.
如果我們仔細去看這四樣東西.
飲食 節期 月朔 安息日.

$^{321}$其實是包括了所有猶太人傳統.
甚至古近東文化裡面一連串的活動.
涉及宗教裡面的活動.
飲食 猶太人很多飲食規矩.
不用我說了.
在《申命記利未記》裡面已經詳細記載.
我不再重複了.
哪些可以吃 哪些不可以吃.
很多的規矩.
節期 節期是說什麼?.
指的是整個農耕的系統.
你不按照這個節期去活.
你根本上.
其實應該說得直一點.
整個古近東文化.
在節期的氣候裡面牽動著 帶動著去運作.
我們現在現代社會不是.
不是的 現在一樣.
現在入春了.
入春就要種一下東西了.
你不種東西.
至少你要搞一塊草地.
如果有一塊草地.
如果你家裡有種盆栽的.
你就要在這個時候換泥了.
每一樣東西.
節期在猶太人裡面.
和世界的創造有連結的關係.
月朔每一日.
每一個月裡面的循環.
安息日就是每一週裡面的系統.
這一連串所講的東西.
其實是回應在前文所講到的.
一連串特別二章第八節裡面所講的.
世上的小學.
那些小學是什麼東西.
好像是一些很基礎的東西.
沒有人走得掉的了.
這些東西全部個個都要做的了.
他就在這些的範疇裡面.

$^{361}$來去二次創作.
去重新用這些東西來去包裝.
使得你要去服從他們的宗教活動.
所以這段聖經所講的對象.
以及所講話的人群.
都在教會裡面.
保羅所談論到的就是.
在教會的當中.
當教會要去繼續走前面.
信徒的成長的道路的過程裡面.
一定有這些的人在教會裡面.
去宣揚這一類型的東西.
要回到一個他不是講傳統.
這些不關傳統事.
個個民族都要守的了.
你外邦人也好不是外邦人也好.
你都要按照這個循環去做.
農耕這個節期.
個個民族都要做的了.
今天都要做.
雖然你怎麼聰明都好.
什麼基因改造都好.
節期仍然是.
掌控著這個世界的農作物.
那個運作怎樣去生產.
在這些種種的活動裡面.
其實都有一個元素是共通的.
就是每一次.
這些的活動裡面.
都會和慶祝這件事情有關.
所以我們無論在舊約聖經裡面.
或者甚至今時今日的裡面.
每逢這些節期也好.
什麼飲食也好.
不用說了.
飲食現在什麼都慶祝的了.
節期也好.
所以為什麼有感恩節.
你有沒有過感恩節.
感恩節本身原本.

$^{401}$為什麼加拿大的會比美國早.
整個月的.
為什麼.
因為我們比他北.
我們的農耕早收工.
感恩節就是去慶祝.
給他能夠在一年裡面的.
耕種的成果.
月朔你不用說了.
每一天你的工作.
每一週你的工作.
種種的都和慶典有關.
當人走入到慶典的範疇的時候.
就很容易會給人有機會.
去作出指控.
是嗎牧師.
為什麼.
每逢有慶典的時候.
就會很容易被人看到.
人的放蕩的一面.
人就很容易犯錯.
不用說哪一個例子.
說約伯記.
如果大家熟悉約伯記的時候.
約伯記的第一章.
你有沒有留意約伯.
是天天向上帝.
為他的兒女去禱告獻上祭物.
為什麼.
留意他的兒女不是窮人.
而且關係很好.
哥哥慶祝也叫一些妹妹來.
因為慶典很容易使人墮入一個放蕩.
或者講得直述一點.
當你在慶典的時候.
是很容易犯罪的.
很容易去講一些狂妄的話.
所以弟兄姊妹.
保羅留心到在教會當中.
如果有一些人仍然在這樣的情況之下.

$^{441}$利用這些時間.
你們就要留意.
如果你參與其中的.
你一定要很小心.
不要讓人有機會來攻擊你.
論斷這個陳述.
就是指到你給了一個機會.
那些人來攻擊你.
會有犯錯的機會.
所以保羅在這個地方.
提醒哥羅西教會的人.
這些人其實他們想找的.
是你們的錯處.
不過請留心.
十七節.
他說:這些原是後事的影兒.
後事的影兒是指什麼後事.
就是那形體卻是基督.
中文在這裡比較難了解.
英文的聖經是這樣解讀的.
There are only the shadow of the things to come.
But the reality is Christ.
原來他在這裡所講到的.
所有這些活動原本的焦點.
應該是放在一個人物的上面.
這個人物叫做基督.
所有這些事情.
在猶太人的社群當中.
本來是平凡不過的事情.
但是如果我們的焦點.
只是放在自己的慶典裡面.
舊約也有講.
每逢什麼節期又有什麼慶典.
有時候什麼時候會期.
最後又有個慶典.
我們把焦點放在我們的慶典裡面.
但是我們失卻了.
原來這一切的事情.
不是指那位創造的上主的提供.
所以這裡所講到的.

$^{481}$後世的影兒是指什麼?.
指到其實所有成就這些事情.
都是由基督所締造出來的.
而這個講法也都和之後所談論到.
或者之前也有講到.
要聯手於耶穌基督這個身體.
是有直接的關係.
所以當我們去看回這段聖經.
保羅要提醒哥羅西教會.
縱使這些飲食,節期,月朔,安息日等等的事.
其實是讓我們看到.
上帝透過主耶穌基督.
怎樣在我們當中創造所締造出來的.
是讓我們更加把焦點放在基督的身上.
而不是把焦點放在我的慶典上.
因為慶典很容易使我們墮入.
這個試探引誘甚至犯錯的上面.
而這些人就在這些時候.
怎樣二次創作呢?.
十八節.
他們就會故意謙虛.
搞敬拜天使.
什麼來的?.
故意謙虛敬拜天使.
故意謙虛是指到一些人.
在那個過程當中.
表現得很敬虔.
每逢慶典的時候.
我們就會留意到一些宗教主持儀式的人.
特別穿得隆重.
而且表現出來的表情或者形態.
就好像我是最謙遜的人的代表.
而且他們搞敬拜天使.
天使是上帝的僕人.
差來傳話的.
這個做法其實就指到.
上帝不再成為敬拜的焦點.
而是指到這些天使來做在我們當中的事情.
那份事情是值得我們高舉的.
用一個簡單的例子.

$^{521}$很多人將某些人封為聖徒.
然後將他成為我們要學校膜拜的對象.
你說這些天主教才有的.
姓什麼姓什麼.
基督教沒有的.
我們沒有封聖這件事.
你看看基督教是不是真的沒有.
基督教將多少的名人名目名講員.
拿出來封為大家的學校模仿的對象.
我曾經在香港認識過一個商人.
他說他聲稱自己是基督徒.
不要緊 你說自己是基督徒無所謂的.
但每一次我跟他談的時候.
他就會說香港某一個著名的基督徒的商家.
是誰我不說了.
快來了 多倫多有說話的.
他怎樣怎樣怎樣.
所以我也要怎樣怎樣.
我要照著他怎樣怎樣.
你可能說牧師可能是他的錯誤.
你留意一下基督教教會裡面.
多少信徒他們經常口邊掛著.
哪位牧者做過什麼 講過什麼.
他們怎樣做.
所以他很喜歡看那些牧者的見證錄.
然後模仿那些牧者的做法.
曾經我聽過一個這樣的事情.
有一位都幾出名的牧者.
我跟他沒有相交的 私底下.
所以我沒有什麼特別.
他說他們教會建堂的時候.
教會一直都遇著很多困難.
不過教會看定了一塊地一定要建堂.
於是乎那個牧者提出了一件事.
我們不如圓了崇拜之後一起拉隊.
然後去我們看中的地盤 那塊地.
然後繞著那塊地一直大聲呼喊.
然後說我們要在這裡建堂 榮耀神.
這種做法我相信是模仿約書亞公艾成的.
因為他自己親口說是這樣的.

$^{561}$後來我認識了一位牧者在多倫多也是這樣做.
吸不吸引呢.
姐妹你不去就是壞人 你不去就是壞女.
你不愛教會 你不愛上帝.
一起攻取這塊地 成為那個地方的見證.
地就繞過了.
不過最後什麼牧者都跑掉了.
地又買不成 教會更加建不成.
你認為這是見證還是羞辱呢.
保羅在這裡所說的這些東西.
全部在大城市很容易出現的.
你不要被這些人用這些手段牽引著你走這條路.
所謂敬拜天使其實就是你去追隨.
或者我們叫做什麼 你做了某個牧者的粉絲.
所以在這裡我奉勸所有在網上聽的實體崇拜的弟兄姊妹都是.
你不是我的粉絲.
我沒有粉絲的 你不要做我的粉絲.
你不需要追隨我 你追隨上帝的話就行了.
我說的是上帝的話 你追隨上帝的話就行了.
但是這些人就用這些方法.
假謙遜在你的面前.
穿得很漂亮的 很虔誠的.
出來祈禱好像很有力量的.
然後怎樣的.
其實是轉過來叫你去敬拜那些神的代理人.
不敬拜上帝的.
保羅在這裡說.
不可讓這些人能夠在這個情況之下.
擄掠你們去失去了你們的獎賞.
這裡所指的獎賞是什麼.
是指到在上帝的神國裡面.
神對於你的一個永恆的肯定.
今天很多類型的這些所謂的謙遜的表現.
來到各式形式來表現.
除了在大聚會 在慶典裡面穿得很虔誠之外.
他們還講到很多 譬如所謂的大愛包容.
樣樣都去接受.
弱勢社群我們要去照顧.
是 我認同弱勢的社群是需要照顧.
但是不是拿出來成為一種炫耀.

$^{601}$而去將自己包裝為好像特別給人.
去值得尊重跟隨的人.
這種的情況就是變成一種逆向的剝削.
將很多我們今天的人.
一些的價值觀將它拆毀.
在這種種的活動裡面.
保羅就要拆開他們的西洋鏡.
免得哥羅西教會的信徒.
被他們這些種種的言論.
來去包裝之後誤導.
所以保羅在這裡再次跟他們講.
這等人驅離在所見的.
隨著自己的肉心無辜地自高自大.
他們就開始怎樣.
用這種種的說話來將自己去抬起.
還有一些空洞的說話.
自高自大的說話.
來將自己抬起.
所以去到最後第十九節.
保羅在這裡這樣說.
他說不辭令元首.
全身既然靠著他.
根折得以相助聯絡.
就因神大德長進.
保羅在這裡強調的一樣東西.
就是用了一個比喻.
好像一個身體.
隨著好像先前所談論到基督的身體.
每一個信徒應該是連繫於這個身體之上.
這些人好像都將自己.
連繫在基督的身上.
但是他們卻沒有將這個頭.
即是基督.
來作為馬首是瞻這個原則.
於是乎怎樣.
在這樣的無緣無故的情況之下.
不是帶你們去長進.
而是帶你們離開了上帝.
哥羅西教會所遇見的危機.
其實在今天每一個信徒同樣是面對.

$^{641}$很多時候我們在這個世界裡面.
因為我們對聖經不熟悉.
被人挫兩挫之後.
我們就變成無力招架.
為什麼會這樣呢.
保羅其實帶出一個很重要的信息.
就是沒有在神的裡面有成長.
我們只是浮在一個表面的裡面.
哥羅西教會不是不認識這些東西.
什麼飲食啊 節期啊 月朔啊 安息日啊.
什麼天使啊 全部這些他們都知道.
但是只是留在一個表面的裡面.
沒有將你深化理解.
我用一個簡單的例子.
譬如說禱告.
每個信徒都說禱告很重要.
但是你有沒有參加祈禱會.
一起在禱告上面操練.
另外說傳福音.
今天很多信徒都說要傳福音.
但是對於福音的整傳.
是怎樣的介紹.
有沒有全面的理解.
因為對於聖經的認識是支離破碎.
有沒有系統來參加主日學.
我看這裡不是推動要回聚會才可以成長.
怎樣祈禱怎樣才會認識聖經.
但是你有沒有靈修.
如果當我們沒有系統地去認識上帝的話.
教會這些主日學也好 祈禱會也好.
其實都是希望提供一個有系統的教導.
讓你知道能夠擦透這些人所說的東西.
是一些片段 是錯誤的.
怎麼錯你會知道.
今天的信徒有多少是真的有這樣的位身.
是呀 你知道什麼叫飲食.
你知道什麼節期.
什麼月朔 安息日 敬拜 怎樣謙遜.
全部你都知道.
你看不看得透 通看得透.

$^{681}$究竟他們在做的事情.
其實是帶你去上帝.
還是帶你去他們自己的權力當中.
要你馴服在他們裡面.
還顧今天的教會.
很多時候信徒熟悉的卻不是一本聖經.
裡面的內容.
而是熟悉他每一天回到家裡.
按著Netflix.
你看Netflix的時間長.
還是看聖經的時間長.
很簡單問這條問題就可以了.
或者你告訴我.
牧師我不看Netflix的.
我看YouTube的.
那你看YouTube的時間長.
還是你看聖經的時間長.
很多時候我們屬靈有問題.
當有測試來到的時候.
我們就見到這些情況.
就拿手不成勢.
於是怎樣回應.
就用回這個社會的套.
今天教會就是這樣.
因為教會根本沒有基礎.
我們站穩在生理的上面.
當這個世界來到.
有些挑戰的時候.
就用回這個世界的套來回應.
社區教會就是這些出現.
社會喜歡這些.
喜歡搞這些活動.
社區支援.
社區喜歡我們就搞多一點.
人進來就等同於信了主.
不是的.
人進來可能是因為那些利益.
那些什麼利益.
飲食 節期 月索 安息日.
所帶來的種種慶祝.

$^{721}$所以保羅在這裡其實.
想提出一點的.
這些東西其實本身.
應該不會左右信徒的生命成長.
但是卻是因為.
在這些事情的上面.
信徒沒有整全的理解.
而在片段之下.
很容易墮入這些困難.
或者這些的人有.
而自己也不知道.
而教會也都順應了這個情況.
去遷就這些信徒.
不理會這些等等的問題.
於是教會就更加遷就世界.
就更加進一步走樣.
沒錯 教會裡面可能很開心.
靈性上面你又參加事奉.
又覺得自己好像很成熟.
上主很喜歡你.
上帝的天堂肯定給你位置.
但結果是怎麼樣.
上帝說你所作的事.
你都不明白.
於是怎麼樣呢.
第二部份.
二十至二十三節.
在這裡造成了一個現象.
人的說話比上主的說話.
更加有力量 更加受歡迎.
這個就是教會更加遷就.
要尋求世界妥協的一個現象.
我們看看這裡二十至二十三節說什麼.
這樣說.
你們若是與基督同死.
脫離了世上的小學.
為什麼形象在世俗中活著.
服從那不可拿.
不可嘗 不可摸等類的規條.
這都是照人所吩咐所教導的.

$^{761}$說到這一切.
徵用的時候.
就都敗壞了.
這些規條使人徒有智慧之名.
卻用私義崇拜.
自表謙卑.
苦待己身.
其實是克制肉體情慾上是毫無功效.
保羅這段說話是很強烈.
用我神學院的教授說.
因為是Hardcore Christian.
什麼叫Hardcore.
真是很硬朗.
企硬的.
保羅企得這麼硬.
對一間跟他素未謀面的.
保羅西教會會不會說得嚴重了一點.
我們看看為什麼保羅要說得這麼嚴重.
在第20節這裡一開始說到.
和2章第8節剛才我們也說過.
世上的小學.
保羅再次提出.
這些人很多時候用一些好像很原始的.
討論或者題目來引發信徒.
來踩入他們的陷阱當中.
小學這個解釋.
可以說是說到一些很微細很原始.
沒有人可以再有爭論的一些題目.
其實也是滿足了很多信徒.
他們不想毀身花時間.
去投入在這個信上面.
去加以努力.
加以去成長.
那份的動力.
很討好人的.
今天的信徒.
你知不知道很多返教會是不想做這麼多事情的.
最好我有錢.
所以你聽有錢出錢有力出力.
其實事奉主不是有錢出錢有力出力.

$^{801}$以前我聽過一位牧者說.
事奉主是全人投入的.
有錢出錢.
不單止沒錢你有多少可以出的就出.
你有多少力就出多少力.
不是用錢或者用力來計算.
是全個人.
因為人有一種懶惰的傾向.
最好就不要太多投入.
懶得做最少就得到最大的益處.
最緊要上到天堂這個基本要求.
進到天堂可以見到很多.
我自己很想見的人.
和一些很老友的人.
可以一起生活.
和上帝生活.
我最想的就是這樣.
所以我做最少的事情.
怎樣可以做最少的事情.
這個就說了.
最少的事情.
什麼事情.
好像世俗上活著的態度.
猶太人最簡單的事情.
有幾樣東西.
當朝律法的規條.
不可拿不可嘗不可摸.
三類規條.
做了就等於搞定.
不過當你看到這個所謂.
談論到不可拿不可嘗不可摸.
這些規條的時候.
就讓我想起.
聖經其實真正的教導是什麼.
詩篇第一篇第一節.
我們一起去看看.
詩篇第一章第一節.
找到之後我們一起讀好嗎.
我們讀第二節.
好,找到我們一起讀.

$^{841}$不從惡人的計謀.
不暫聚人的道路.
不助洩萬人的座位.
唯喜愛耶和華的律法.
晝夜思想者人便為有福.
有沒有留意.
這裡所說的.
不是不可拿不可嘗不可摸.
那些東西是律法規條做了.
就好像很虔誠.
聖經不是這樣教的.
聖經教的是不從惡人的計謀.
不暫聚人的道路.
對不起.
不從惡人的計謀.
不暫聚人的道路.
不助洩萬人的座位.
是一個一連串持續的行動.
不是一次過.
而是常常面對著.
惡人,罪人,洩萬人.
你不要跟他們有任何的交往.
但是你應該有什麼交往.
你應該跟上帝的律法.
晝夜思想.
在聖經裡面所教導的東西.
跟西族人教導的東西.
是截然不同的.
不同在哪裡.
一輩子做了一些規條.
就以為滿足了上帝的要求.
但是上帝對信祂的人的教導就是.
你要持續不要在罪惡,罪人,洩萬人,惡人.
等等他們走在一起.
即是聖經所講到的.
聖潔這個題目.
聖潔這個字眼.
英文Holy.
是指到分開的意思.
分開什麼.

$^{881}$跟什麼分開.
跟罪惡分開.
你跟罪惡分開.
就是成聖的人.
信徒叫聖徒.
為什麼.
你和我都是聖徒.
為什麼.
因為我們跟罪惡竭力分開.
這個是持續的事情.
不是一次的事情.
你看看這段聖經所講的.
不可常,不可摸.
像是對一些飲食題目.
慶典上的題目.
所以這些慶典有很多規矩.
你不做那些就聖見了.
No.
正正很多做這類不可拿,不可常,不可摸的人.
這班人其實是假冒偽善在那裡.
保羅在這裡講得很清楚.
這班人怎樣假冒偽善活在他們當中.
他用很多人的說話在這裡作出教導.
表面上好像在講一些有智慧的說話.
其實是全部帶著私心來去參與.
我們看一下.
尤其是第二十二,二十三節裡面講.
這都是照人所吩咐的.
他們講就是律法了.
所以為什麼猶太人的整個猶太教的宗教.
會建立在他們自己個人所顯示出來的六百多條的律例呢.
就是這個原因.
他們所講的這一切.
都其實是要在這裡敗壞你.
使你怎樣呢.
離開了上帝的規矩.
跳進他們的規矩裡面.
這些人徒有規條.
而講出來的規條好像很有智慧.
好像很有智慧.

$^{921}$但其實目的是要施以崇拜.
按他的意思來崇拜.
最好你崇拜他.
講到這裡我順帶講一句.
如果弟兄姊妹你聽我說到有什麼學習呢.
因為我的工作就是將他的說話告訴你.
我不需要也不應該成為你們敬拜的對象.
你尊重牧師沒問題.
尊重沒問題的.
但你應該聽到神的說話.
讓你明白到上帝的心思意念.
但這些人是怎樣.
他敬拜都是希望帶你進入他.
用他的方式來敬拜.
因為他敬拜是希望你的人心去到他身上.
表面好像很謙卑.
甚至出苦肉計.
用苦代己生.
但保羅一句就說下去.
就算你所謂克制肉體的情慾.
其實對於你都是無功效.
因為這些都是包裝.
保羅在這裡所講到這一連串的事情.
就給我們看見.
教會今天所做的事.
教會做很多的事.
其實究竟是帶著信徒去事奉上帝.
還是事奉你自己的教會團體.
將你的團體做大.
我做牧者這麼多年.
常常聽見牧者之間常常講.
怎樣去牧養教會.
使教會更加興旺.
但結局很多時候他們的討論都是.
教會做了什麼有多少人來.
有很多人來就等同於教會成功.
萬事的人來只不過是因為其他原因.
我做牧者做了信徒三十多年.
我自己很明白一點.
很多人去教會其實是別有用心.

$^{961}$我希望我講這番說話.
有別有用心的信徒在這個時候.
可以回到正規的道路上.
我不是嚇你不是叫你離開你的教會.
是希望你回歸到上帝裡面.
有些人回教會可能只是為了找配偶.
有些人回教會可能是為了做生意.
有些人回教會只是為了找位置.
有些人回教會可能只是為了.
有些自己的宗教滿足感.
在這個世界上的工作失敗感太多.
教會怎麼做.
不要緊你在社會失敗你在教會成功.
上帝看見你都是成功.
用這些說法哄著你回教會.
所以很多牧者以為教會興旺等.
等同於人多這個說法.
我想你應該要去重新審視一下.
很多人回教會要找大教會多人教會.
到時就會多些見證多些資源.
就可以多些去事奉.
我甚至有牧者這樣說.
我不怕老實說.
甚至直接說以馬類不如執了他.
因為人少資源缺乏.
不如你留一些弟兄姊妹.
愛他們更多.
使他們回去大教會資源豐富.
你又舒服了Alan.
不用每個星期講道.
我最近才發覺原來做大教會的主任牧師.
一個月講一堂道.
我發覺他們.
我跟師母說了一句.
他們都幾好幾好賺.
那份工資又高.
七萬八萬九萬.
嘩還有數之不盡的.
那些叫做fringe benefit.
醫療保險什麼都有.

$^{1001}$幾好是吧.
好在什麼.
他就好.
騙了你掉進他的陷阱裡面.
這些就是保羅在這裡說.
他用很多方法來吸引你.
好像很有智慧帶你去到他那裡.
但是怎麼樣.
他說一切功效都沒有.
曾經有一個神學院教授.
指著那些神學生說.
他說你們這班神學生.
將來要出來事奉真理的主.
傳講真理的道.
不過我真的很不明白.
為什麼你們傳真理傳到這麼壞.
那些戲子佬在戲院裡面上演的戲.
虛情假意卻賺得你們的流淚.
但是你們做不到.
聽清楚這句話說什麼.
你聽不聽到那個教授.
神學院教授在說什麼意思.
你們坤人的手法不夠.
或者他的意思就是.
你們將上帝的說話說得很好.
都變成很壞.
哪一種我不敢說.
不過突然之間.
在聖經裡面.
以賽亞書第29章13節.
他這樣說.
主說因為這百姓親近我.
用嘴唇尊敬我.
深卻遠離我.
他們敬畏我.
不過是領受人的吩咐.
又在馬可福音第七章第七節裡面這樣說.
他們將人的吩咐.
當作道理教導人.
所以拜我也是枉然.

$^{1041}$這兩段聖經無獨有偶都提出.
在教會的圈子裡面.
有很多虛情假意的假敬虔的人.
用這種假敬虔的事奉.
來對吸引人去到他們的面前.
不過他訛稱.
你是去到上帝的當中.
哥羅西教會活在一個這樣的世界裡面.
保羅一語道破.
因為他們當中有很多人為的教導.
什麼時候知道?第三章會說的.
人在教會的當中.
很多時候因為要面子.
所以很容易在這些人的虛情假意裡面中了招.
以為有時妥協一點點在這個世界沒什麼問題.
反正就可以傳主的道.
例如社區教會運動.
Community church movement.
將教會塑造成為一個社區中心支援點.
聽起來挺好的這個世界.
尤其是在這個所謂大愛包容公平無差別.
這種普世價值的觀念之下.
你如果告訴別人你不做這些事.
反而人們會覺得你真的不是教會.
你看看今天很多人做了這些之後.
教會進去的是什麼人.
成為了牧者的是什麼人你就會明白.
全部都是在這個世界裡面進行妥協.
例如最簡單.
飲食上.
說起來已經二十年前了.
以前我送我女兒去一間教會的幼稚園.
叫她托兒.
每一次大約六五點我去接放學的時候.
都沒什麼特別.
但是有一次我回去的時候.
我就聞到很香.
很香什麼?.
煎牛扒的味道.
我就覺得很奇怪.

$^{1081}$我跟那間教會個別有幾個人都認識.
我就問很香.
為什麼這麼香?.
他說你有所不知.
我們今晚就是Alpha Course的開學禮.
我說Alpha Course開學禮有什麼這麼香?.
他說Alpha Course要請一些未信主的人來參加.
最好要人來參加的就是請人吃飯.
請人吃飯跟這麼香有什麼關係?.
他說我們今晚就吃牛扒.
我說Alpha Course參加者可以吃牛扒?.
我說我能不能來?.
他說你是做牧者的當然不如你.
不過你要來我也歡迎.
不過你要帶一個不信主的人來.
飲食上很合邏輯.
不要說吃東西.
教會活動,電腦班,排舞班.
不可以你自己一個人來.
你要帶多一個未信主的人來.
來到做什麼?.
一起跳排舞.
我們跳一些福音舞.
福音操你有沒有見過?.
我有一間教會有一次教課程.
他教會請我去教會課程.
晚上六點多走來做福音健美操.
還叫牧師你一起來做.
我說我不做了我很累.
這些很吸引人的東西.
在這個世界很多教會.
會成為他們說所謂事奉主的藉口.
或者理由.
去方便更多人來接觸這位上主.
因為他們更加經常拿著那句話.
耶和華耶和華.
是有憐憫有恩典的神不輕易發怒.
並且有豐盛的慈愛和誠實.
這句話在聖經裡面出現過七次.
他們就說我們教會無論怎樣.

$^{1121}$在恩典之下有少少妥協吸引到人來.
就沒有問題.
他來到最後都會聽到福音.
有多少個吃完牛扒之後留在這裡.
我不知道.
不過有免費牛扒吃必然有很多人來.
人多是不是教會成功呢?.
還是很多不成功的人在教會裡面.
是上帝看不到.
是正如保羅所說毫無功效呢?.
等於每次教會如果要尋求社會認同.
這是一件很危險的事.
也同時反映出信徒對神的說話很沒骨氣.
信不過上帝的說話.
沒錯上帝的說話傳了出來很多人不喜歡聽.
但是這些卻是令人最受益的說話.
就是上帝讓人去面對他們的罪.
也都是上帝願意主動幫助人去處理人的罪.
但是今天的教會卻是擔心因為這樣的情況之下.
別人卻不聽 可不可以反過來想一想.
他不肯去聽.
正正可能就是他對上帝的福音根本是沒興趣.
你說我可以用很多方法扭轉他的興趣.
我告訴你 我傳福音傳了這麼多年.
這些興趣是不會改變的.
他是沒興趣就是沒興趣.
因為得救與否是上帝命定的.
我相信一件事.
當神的道傳開的時候.
肯聽上帝說話的人就會有正面的回應.
不肯聽上帝說話的人.
就會更加去反抗上帝的說話.
不用我說 最近那些航機上面那些.
航空公司上面的航機那些人在做錯事.
指證全機的人都說他錯的時候.
他仍然不認自己錯.
就是這樣 西人就是這樣.
那你就會在這裡問.
哇 保羅這樣對哥羅西教會講到這麼直.
保羅不怕他們不肯聽.

$^{1161}$而且更加去跟隨那些傳這些.
施以崇拜自我謙卑.
假扮苦肉計的人來他們講的說話.
他們會趕到那裡去聽嗎.
保羅不怕的嗎.
弟兄姊妹 作為一個牧者.
如果你要怕的.
是怕你傳上帝的說話不夠準確.
不是怕那句說話傳了之後.
那個結果你不理想.
我不理想不要緊.
最重要是上帝他的決定.
在那個人身上他得救可以回轉.
保羅我看到他的那個回應就是.
他既然出得來去事奉上帝.
他就要承擔這些事奉所.
同時帶來的挑戰和擔當.
只要是對教會有益處的說話和見證.
以及教導他無懼任何人的反對.
這個就是保羅的生命特色.
現在每一個基督徒應有的特色.
來到最後我想問.
你有沒有這樣的生命.
還是你還是說算了.
我是小信徒而已.
沒什麼影響力的.
你是哪一類信徒.
我們一起禱告.
父親多謝您讓我們在這條聖經裡面.
看到保羅對哥羅西教會的愛護.
可以說是冒著險來去.
虧欠在哥羅西裡面的每一個信徒.
他明知教會裡面.
有一些人傳別樣的教導語及福音.
以及以錯誤的神學.
牽引人走去一個遠離上帝的境敗.
但保羅卻是勇於將這些的陳明.
讓有耳當聽的人都能夠聽到.
並且作出恰當合神心意的回應.
今天我們有沒有這份勇氣.

$^{1201}$以及骨氣去事奉主理.
我們的屬靈的生命.
是不是在上主理裡面.
能夠真是按照你的說話.
站硬來見證主理的教訓.
很多時候我們以為.
我們徒有智慧.
來懂得如何在這個世界裡面周旋.
以為我們稍稍作一些妥協.
或者隱藏.
就能夠在上帝你的事工上.
在這個世界當中.
好像臥底一樣隱藏在其中.
去傳揚主理的教訓.
但誰不知.
其實我們自己也跌入了一個陷阱裡面.
求主理的幫助.
讓我們有一個警惕的生命.
常常留心到.
這些種種此是而非.
包裝得好像很敬虔的事情.
讓我們學會如何辨別之後.
能夠再次站在一個正確的信的基礎上.
求你的幫助.
我們今天活在這個地上.
行你的說話.
最終的目的是要活像主理.
而不是活得更加合我們自己的意願.
求你讓我們的信仰.
不是一個自我滿足的信仰.
甚至一個以基督教詞彙包裝的自滿信仰.
而是一個不折不扣合神心意的一個相信.
活出一個聖潔的生命.
不與惡人,罪人,亂人同為一黨.
並且學會常常在主理的說話上思想,實踐,回應這個世代.
求你讓我們每一個都有一個智慧的靈.
並且有一份願意毀身的動力.
身體力行來追隨上主理的教訓.
有系統地去學習主理的說話.
求你的大靈禱告奉靠主耶穌基督得勝的名字而求.

$^{1241}$阿們.
\newpage



\section{出埃及記 23:12}
\label{sec:032056}
\textbf{聖經人聖經事:第七日要安息}~~~~黃克勤牧師博士
\newline
\newline
連結: \href{http://www.church.com.hk/acms/content.asp?site=cdc\&op=show\&type=product\&code=032056\&layout=sermon}{\texttt{032056}} ~~~~ 語音日期: 2026-04-23
\newline
\newline
\hyperref[sec:032055]{\small{< < < PREV SERMON < < <}}
~
\hyperref[sec:index_chronic]{\small{[返順時目]}}
~
\hyperref[sec:032057]{\small{> > > NEXT SERMON > > >}}
\newline
\newline
出埃及記 23:12
\newline
\begin{longtable}{cl}
\hline
\hline
章節 & 經文 (和合本修訂版)\\
\hline
23:12 & \begin{tabularx}{0.7\textwidth}{X} 「六日你要做工，第七日要安息，使牛、驢可以歇息，也讓你使女的兒子和寄居的可以恢復精力。 \end{tabularx} \\ \\
[1ex]
\hline
\hline
\end{longtable}
$^{1}$各位弟兄姊妹,平安!.
又一齊讀聖經了.
繼續看摩西這個人.
當然其實已經不是摩西了.
是我們的主.
他怎樣將他的心願誡命告訴我們.
所以看聖經的事.
是講到一些條例,律例.
今天很順利地看第七天要安息.
不是休息,是安息.
大家看看.
「華的野獸可以吃」.
「你的葡萄園和橄欖園也要照樣辦理」.
第十二節.
「六日你要做工」.
「第七日要安息」.
「使牛奴可以歇息」.
「並使你婢女的兒子和寄居的都可以舒暢」.
讀完這三節聖經.
心裡真的很感恩,很舒服.
因為神是很看顧人的.
《昌埃及記》二十三章第十節.
「六年你要耕種田地」.
「收藏土產」.
意思就是頭六年.
第一點要分享的頭六年.
其實神安排我們工作.
是有一段時間工作.
然後有休息.
有一段時間努力工作.
不單是工作是努力工作.
然後有休息.
頭六年拼命耕種田地.
要耕種田地.
不是隨便耕種.
要耕種田地.
耕種完其實耕種也辛苦.
耕種這件事很辛苦.
然後還要收藏土產.
耕種完收藏土產.

$^{41}$神給你足夠豐富的收藏.
特別是農民就會知道.
因為他不是像我們現在這樣上班.
每個月出糧.
又用了他儲一點每個月出糧.
他是一耕種下去.
然後仰望上帝.
要雨水,要陽光,要好天氣.
要適合的溫度.
然後收割很辛苦.
要快速收割.
然後將它收藏起來.
然後分開有哪些可以賣.
有多少可以吃.
要保養住,保護住.
不要變爛變壞.
因此做農民是靠農作來謀生.
是艱難很多的.
由於其實不要以為耕田很容易.
很多事情要想清楚.
什麼時候要灌溉,什麼時候要施肥.
什麼時候要除草,什麼時候要殺種.
哪個時間才對.
要仰望上帝.
有陽光,有水.
什麼時候收藏,收藏在哪裡.
收藏的地方的濕度要怎樣.
大概我們一家人要吃多少.
剩多少,然後賣多少.
很多這些事情.
如果現在的價錢不好.
要不要種別的,種別的自己會不會.
其實這六年真是頭都大了.
又要付出很多努力去耕種.
又要收藏.
一點都不容易.
第一點要賣力.
第二點到第七年只是.
不過你要留意到第七年.
就要叫地歇息.

$^{81}$忘記了螢光這個字.
第七年要歇息.
特別叫地歇息.
是不耕不種.
這個中文和學本的翻譯挺有趣的.
我留意到有些字跟大家分享一下.
歇息是sharmat.
sharmat是讓它跌下來.
到第七年你要讓土地跌下來.
那是什麼意思.
跌下來是你不用出力.
拿著它.
其實譬如有些東西下來.
你要拿著,拿著六年.
拿著它很辛苦.
第七年放下它.
將土地放下它.
不搞.
這個不耕不種其實只有一個字.
就是letharge.
letharge是離開任憑拋棄.
第七年地離開它.
任憑拋棄丟它在這裡.
第七年其實.
如果你還有很多東西.
它自己會長出來.
或者你還有些種子掉在地上.
它會長出來.
第七年其實最好.
看來以以下的字來看.
第七年你應該殺些種.
讓它可以有彪牙.
仍然有這些作物.
不過第七年不要搞.
讓它不用拿著.
拿著那麼辛苦.
又要出力.
放手.
讓它跌下來.
不過那塊地不耕不種.

$^{121}$有四個中文字.
講了耕和種.
其實就是離開.
讓它丟在這裡.
如果離開任憑拋棄.
就說讓它丟在這裡.
使你的民.
使你民中的窮人.
使你這個字是沒有的.
所以我用藍色的.
民.
讓你的M.
M就是人民百姓.
讓people.
讓些人.
其實到第七年.
我就讓塊地丟下.
就不理它了.
就等些人民百姓.
中的窮人.
這個中的都沒有的.
是讓人民.
和.
窮人.
或者有需要的人.
那些needy.
就是有需要的人.
他們都可以有吃的.
就是說.
這段時間.
就讓他們進去田地.
去拿東西吃.
拿東西吃.
摘東西.
其實也包括了.
我的地.
第一我要讓他進來.
我圍住他.
要開門讓他進來.
就是說大家.

$^{161}$雖然他是窮人.
但我都很信任的.
他進來我的田地.
他拿東西吃.
讓他.
有吃的.
其實在另外的經文.
都講了這件事.
不單止去到第七年.
平時每一年的收割.
要留一些角落的地方.
譬如我這個.
電視石那塊田.
如果收割.
是收割到邊的角落要留一些.
邊要留一些.
角要留一些.
邊要留一些.
不要全部收割.
為什麼不要全部收割.
因為是讓窮人.
可以去撿墨水的.
那老鄧就是撿墨水的.
就是很窮的.
那他拿鵝米.
他婆婆都很窮的.
那他就是撿墨水.
自己和婆婆衛生的.
那撿什麼.
就是撿波亞斯田地裡面.
邊邊那些.
特意留在那裡的.
這個是神教儀式的人.
你要照顧你裡面的窮人.
你會發現到.
神是有這樣的心.
祂很愛那些人.
特別愛窮人要照顧他們.
就是讓那些窮人可以.
在第七年種植.

$^{201}$吃東西.
其實塊田不理.
你看到有幾個窮人進去.
我留一些墨水在那裡.
他又學我種.
又來看.
又來拔草.
多好啊.
給他的.
好了.
如果他撿到的.
他摘到的.
給他.
不要馬上跑出去.
為什麼摘我的東西.
為什麼摘我的水果.
不會的.
給他吧.
所以剩下的.
人吃不完.
剩下的.
就給野獸可以吃.
野獸可以吃.
野獸這個字.
野獸會來吃這些田的嗎?.
野獸會來吃這些果的嗎?.
不會的.
其實.
就譯中文.
野獸譯得不好.
是K.
K是動物.
是吃剩的那些.
好.
第七年扔下塊地.
任他做.
有些東西剩下給他.
那些窮人.
有需要的人可以來吃.
吃剩了.

$^{241}$神就說.
我會給那些野獸.
就是那些動物吃.
其實動物大部分.
如果這樣多數吃農作的動物.
是什麼呢?.
如果是吃農作多數是飛鳥.
來吃那些種子.
飛鳥.
或者會有一些.
野牛野羊.
進來吃那些.
禾稈草.
吃那些草.
又或者是田鼠.
就是這些.
他才給野獸吃的.
如果你的葡萄園和橄欖園.
都要照樣扮你.
說得這麼嚴重.
其實照樣做.
照樣這樣做.
第二點我覺得神就很好.
有兩件事的.
第一就是.
那塊地都要休息.
給那塊地休息.
那塊地他都辛苦了.
六年不停地供應那麼多植物.
第七年要休息.
是很對的.
原來地裡面很多肥料.
是慢慢給你吸收.
六年很多肥料沒有了.
第七年我繼續種.
繼續種都沒有用.
因為他已經不夠肥了.
他要有一年時間.
等地下再重新恢復呼吸.
這樣才好.

$^{281}$再種才好.
所以神很厲害.
他當然了.
他創造的那塊地.
他很清楚.
六年又第七年.
給他恢復.
那塊地都要歇息.
而且歇息的時候.
就給窮人.
有需要的人.
甚至動物.
有東西可以吃了.
記得不要只是說田.
田他們多數都是種大麥小麥的.
田不要只是田.
葡萄園,橄欖園都要這樣.
講得很清楚.
第一六年努力耕種.
第七年休息.
那塊地休息.
我們很多時候忘記了那塊地要休息.
就算你做生意也是.
六年.
生意的產品.
很多時候都差不多了.
如果你不再出新東西.
那個產品就已經老了.
沒力了.
你要想一下新的東西.
舊的產品可能是要休息了.
是要丟下放下它了.
放下它.
然後給其他人可以做.
而你可能有新的要給歇息.
那塊地歇息.
可能你的產品也要歇息.
第三點.
六天你要做工.
又講我勞力做工.

$^{321}$第三點要講.
第七天要安息.
不是休息.
大家要留意.
是安息.
安息這個字.
suburb.
我想大家都聽過了.
suburb是休息.
很多時候不是.
是安息.
其實suburb是停止.
cease.
要停止.
停止.
到第七天要停止.
頭六天做工.
第七天是停止.
我們很多時候.
由於因為中文聖經.
將這個停止.
cease.
suburb.
是翻譯成安息.
所以很多時候.
我們有傳到人會.
第七年安息年.
很好.
安息年.
其實是停止.
cease.
做什麼呢.
第七天停止.
是我可以停止.
另外最重要要留意這個字.
令到牛和驢.
幫你耕田做事的那個.
可以歇息.
留意.
安息和歇息原來是不同的字.

$^{361}$suburb是安息.
其實是停止.
cease.
好了.
歇息.
就是new work.
CH讀K.
new work.
是休息.
休息.
牛和驢都要休息嗎.
牠晚上已經休息了.
是的.
有時你編得緊要.
趕得牠緊要.
牠坐開力.
嘩.
牠累了.
休息一夜就可以了.
不行的.
人一樣的.
體力過度虛耗了.
久了牠會瘦的.
會病的.
不能吃東西的.
不開心的.
太累了.
休息不夠.
所以.
到第七天要給牛和驢.
歇息就是休息了.
安息就是停止.
歇息就是休息.
還有給你的婢女.
這個字挺有趣的.
我們以前.
叫那些做工的女士.
叫做阿麻.
阿麻.
找些阿麻幫忙吧.

$^{401}$嘩.
這個字希伯來文.
婢女是阿麻.
阿麻.
真的很有趣.
其實.
我們很多時候.
一直這樣.
現在回看.
很有可能有幾個重要的字.
都很相似的.
在不同的語文裡面很相似.
於是大家想.
是不是有一個比較相同接近的.
祖先呢.
約你的婢女阿麻.
女僕女奴.
或者婢女.
可能是奴隸.
你知道.
男的.
就是做一些事情.
趕牛勞耕作什麼的.
女的呢.
女的不用做.
女的奴隸女僕呢.
女的奴隸.
或者給錢買的僕人呢.
都是要做事的.
打水啦.
煮飯啦.
洗衣服啦.
清潔啦.
照顧小孩啦.
一樣這麼做的.
第七天你要給那個阿麻.
休息的.
這裡還有有趣的.
讓阿嬤的兒子也休息.
那是什麼呢.

$^{441}$你要知道.
其實那些阿嬤是很辛苦的.
因為.
她要幫你煮飯.
洗衣服.
煮飯.
吃.
清潔地方.
各樣各樣各樣.
服侍很周到.
很貼身.
如果當她有個小孩的時候.
那小孩是怎樣的呢.
是不是她就照顧小孩.
這個小孩就很享受.
不會的.
其實那個小孩.
都會跟著她媽媽.
來幫忙做事的.
如果有個小孩.
那個阿嬤是更加辛苦的.
而那個小孩呢.
由於那個小孩.
都是女僕和奴隸來的.
因此.
她都會很賣力地跟媽媽做事.
即是她都服侍你的.
那小孩拿一點東西就說很重.
又累.
當然是小孩嘛.
正常是這樣的嘛.
所以都要給她休息的.
給她玩一下的.
寄居的.
聖經裡面很注重.
神很注重照顧那些寄居的人.
寄居呢.
這個字剛剛說到話了.
為什麼呢.
這個音叫GIRL.

$^{481}$新來的.
寄居又怎樣.
GIRL她啦.
是不是這樣呢.
GIRL.
外人.
新來的人.
就叫做GIRL.
就是寄居的.
聖經在中文就叫做寄居.
其實是婢女的兒子和外人.
新來的.
那麼這個.
婢女兒子和寄居.
怎樣放在一起呢.
寄居的就是說.
有些人.
來了你們那裡.
來了你的城市.
來了你的社會.
那裡做什麼呢.
他想移民來這裡.
或者他.
哎呀.
我家那邊很饑荒的.
我就想在這裡找一份工作.
哦.
新來的.
外人.
幫忙打工.
或者甚至.
很餓.
那當然你幫我.
那就幫我做事.
不過呢.
工錢就沒有了.
在這裡吃.
在這裡住.
有食無功.
我小時候.

$^{521}$是認識很多人.
甚至我們在柬埔寨.
二十年前去.
都是很窮困的時候.
很多人都是有食無功的.
來吧來吧.
在這裡幫我照顧孤兒院小朋友.
不過有飯吃.
有地方住.
不過沒有錢的了.
好啊好啊好啊.
都搶到來的.
因為他.
外人會走來你那裡住.
就是說其實他已經很困難.
是照顧他的.
第七天要給他停止他的工作.
歇息.
休息.
第七天是誰做的.
自己做.
或者留給他明天做.
第七天給他休息.
都可以舒暢.
這個字特別要留意了.
可能當時沒有我們用這些字.
給他舒服一點.
輕鬆一下.
舒暢.
今天不用做.
我們小時候就說.
今天不用上班嗎.
星期啊.
星期這個字.
以前的人都不認識.
不過原來是基督教的字.
不用做.
今天不用做.
今天休息.
星期啊.

$^{561}$休息星期啊.
休息星期而已.
星期那天是敬拜主.
休息一下.
給他舒暢.
都可以休息一下.
以為是星期不用做.
星期這個字.
Lavash.
是指恢復和振作.
英文是什麼.
今天懂了.
Refresh.
讓他重新得力.
怎麼可以重新得力呢.
一定要給他有一天休息.
停止休息.
地要休息.
這裡第一點.
六年落地工作.
六日落地工作.
第七年休息.
第七天休息.
這個是神特別指明要做的.
是讓我們休息.
Refresh.
休息了.
可以恢復重新振作.
地也是這樣.
人也是這樣.
牛也是這樣.
奴也是這樣.
外來的人也是這樣.
是需要的.
而且給些空間.
讓那些窮人.
就在這個時間.
可以出來撿到食物.
這個也是重要的.
很感恩的.

$^{601}$看到上帝.
是很照顧我們身體的需要.
我自己特別歡喜什麼呢.
有晚上有早晨.
是第一天.
猶太人以色列人.
是用晚上做一天的開始.
大概晚上六點鐘.
或者五點鐘.
當是一天的開始.
黃昏了.
吃一餐飯.
然後大家在家裡吃一餐飯.
很高興.
然後休息.
先儲備能量.
再開始工作.
就是這個意思.
就是要讓他休息.
然後再做得更加好.
走更遠的路.
我們一起祈禱.
多謝我們的主.
讓我們有休息.
叫我們休息.
主,今天我們很多時候.
不懂得休息.
或者在休息的那一天.
忙了其他的事情.
令到我們更加困倦.
以致我們星期一.
新開始的時間.
累幾天才到一個星期完.
然後我們又在weekend.
在週末再忙.
然後令到我們.
上班又很疲憊.
主,其實你叫我們在.
有一個安息的日子.
安息日就是安息的.

$^{641}$休息好.
以致我們有新的工作.
主,你自己是將安息日.
剛好放在星期五晚.
至到星期六晚.
這段時間叫我們星期六.
可以多休息.
星期日敬拜你.
又可以做其他的事情.
然後又開始一天新的工作.
你是叫我們那六天做得好.
主,讓我們更加努力.
在頭六天.
然後到那一天.
我們休息.
安息日.
你是叫我們守衛聖日.
主,我們尊敬你.
除了尊敬你自己的創造.
我們也尊敬你想我們休息.
垂灘祈禱.
奉耶穌基督命求.
好.
\newpage



\section{創世記 5:21-29}
\label{sec:032057}
\textbf{與神同行}~~~~黃若娟傳道
\newline
\newline
連結: \href{http://www.church.com.hk/acms/content.asp?site=cdc\&op=show\&type=product\&code=032057\&layout=sermon}{\texttt{032057}} ~~~~ 語音日期: 2026-04-26
\newline
\newline
\hyperref[sec:032056]{\small{< < < PREV SERMON < < <}}
~
\hyperref[sec:index_chronic]{\small{[返順時目]}}
~
\hyperref[sec:032058]{\small{> > > NEXT SERMON > > >}}
\newline
\newline
創世記 5:21-29
\newline
\begin{longtable}{cl}
\hline
\hline
章節 & 經文 (和合本修訂版)\\
\hline
5:21 & \begin{tabularx}{0.7\textwidth}{X} 以諾活到六十五歲，生了瑪土撒拉。 \end{tabularx} \\ \\ \relax
5:22 & \begin{tabularx}{0.7\textwidth}{X} 以諾生瑪土撒拉之後，與神同行三百年，並且生兒育女。 \end{tabularx} \\ \\ \relax
5:23 & \begin{tabularx}{0.7\textwidth}{X} 以諾共活了三百六十五年。 \end{tabularx} \\ \\ \relax
5:24 & \begin{tabularx}{0.7\textwidth}{X} 以諾與神同行，神把他接去，他就不在了。 \end{tabularx} \\ \\ \relax
5:25 & \begin{tabularx}{0.7\textwidth}{X} 瑪土撒拉活到一百八十七歲，生了拉麥。 \end{tabularx} \\ \\ \relax
5:26 & \begin{tabularx}{0.7\textwidth}{X} 瑪土撒拉生拉麥之後，又活了七百八十二年，並且生兒育女。 \end{tabularx} \\ \\ \relax
5:27 & \begin{tabularx}{0.7\textwidth}{X} 瑪土撒拉共活了九百六十九年，就死了。 \end{tabularx} \\ \\ \relax
5:28 & \begin{tabularx}{0.7\textwidth}{X} 拉麥活到一百八十二歲，生了一個兒子， \end{tabularx} \\ \\ \relax
5:29 & \begin{tabularx}{0.7\textwidth}{X} 給他起名叫挪亞，說：「在耶和華所詛咒的地上，這個兒子必使我們從工作和手中的勞苦得到安慰。」 \end{tabularx} \\ \\
[1ex]
\hline
\hline
\end{longtable}
$^{1}$丁姊妹早晨.
很高興又回到自己熟齡的家.
來分享訊息.
今天和大家分享的一個題目.
是很平淡.
但亦是我們一生一個很大的福氣.
就是與神同行.
為什麼會說這個題目呢?.
我發現過去七八年的時間.
許多不同的人都說同行.
特別我記得.
我2019年受傷的時候.
在休養期間.
當時很特別.
忽然之間有一天.
主就臨到我和我說.
你知不知道未來這幾年.
我會將一個詞彙.
放在香港整個城市.
這個詞語就是「同行」.
有時人很遲鈍.
後知後覺.
我自己也是.
所以當時我說.
主你這樣說我當然相信.
但其實是什麼意思呢?.
其實不是很清楚.
但七年之後.
到現在七年過去.
我發現是喔.
為什麼我過去七年.
我開電視.
看新聞.
看政府施政報告.
看商界朋友所推舉的東西.
很多時候都用「同行」這個詞彙.
特別教會在這幾年.
我講的不是「沙關浸信會」這麼簡單.
我講的是香港教會.
其實「同行」這個詞.

$^{41}$是我們經常說.
直至幾個月之前.
我心裡真的覺得.
真的要鑽研一下.
在聖經由舊約到新約.
其實「同行」是什麼呢?.
特別當很多人都說.
要與貧窮人同行.
商界夥伴要同行.
人與人之間.
經常說「同行」.
但其實一切的來源.
一切的根基.
其實在神的裡面.
因為祂真是眾光之父.
所以當如果一切的祝福.
都是從祂而來的話.
「同行」這個詞.
首先真的要問.
其實神怎樣與我們同行.
我們亦都怎樣與神同行呢?.
我很想邀請.
我們帶著一個謙卑和敬畏的心.
我們首先來到主的面前.
來躺開我們的心給祂.
我們一起祈禱.
天父在《瑪拉基書》三章.
你這樣說.
「萬君的耳和話如此說.
你要差遣你的使者.
為你預備道路.
你要向尋求你的人.
你說你會忽然臨到你的殿與」.
所以天父.
這個忽然真是好.
因為忽然.
是意味在我們沒有預設.
沒有計劃.
甚至乎可能我們.
常常都好似慣常.

$^{81}$來到去翻崇拜.
是我們一切的聚會.
但我們沒有辦法可以預計.
永生神.
亦都是這位萬君的耳和話.
你會忽然臨到你的殿與.
所以萬君的耳和話.
我們的神.
我們今日作為你的子民.
亦都是你地上的身體.
亦都是象徵你萬君的耳和話.
軍隊元帥.
所率領的軍隊的時候.
天父我們真是說主.
我們切務你.
我們的心尋求你.
當我們去繼續尋求你的時候.
主你忽然臨到.
是我們好想為了這個忽然.
在我們沒有辦法預計的情況下.
主我們真是帶著一個.
來到歡迎你這個忽然臨到.
所以聖父聖子聖靈三位一體的神.
主我們好想歡迎你.
主在今日四月二十六日.
今日的這個日子.
就是這個主日.
主我們說我們的心渴想你.
我們好想歡迎你忽然臨到我們的中間.
所以主我們就願意打開我們的心.
主來接受主你忽然之間的來臨.
多謝你的主.
讚美你.
賜納我們在你的面前.
所呈現的一切敬拜和禱告.
奉我主耶穌基督得勝的名.
Amen.
與神同行.
我們今天不看經文字.
我們先看看這個字.

$^{121}$與神同行其實是什麼意思呢?.
原來如果用希臘白萊文的原文來說.
這個字是說來回行走.
走來走去.
是什麼意思呢?.
是說多於一.
就是兩個三個甚至幾十個人.
是來回的行走.
來回的行走.
就是我們上年那個值班.
不知道你記不記得.
你看著今年.
不知道你記不記得上年呢?.
上年就是走在一起.
預備未來.
所以走在一起的時候.
原來今年再說.
原來這是我們同行的關係.
這個同行的字.
只是舊約已經出現了千多次.
很多.
來回的行走.
走來走去.
也是指經常慣性的來回行走.
如果用新約來說.
就是親密的團契關係.
你這樣看.
有幾個關鍵字你要留心.
就是形容詞.
形容是經常的.
形容是慣性的.
所以如果有些人.
你跟他同行的關係.
是你一年見他一次.
你覺得算不算是聖經說的同行呢?.
當然不算.
如果你說你半年出去吃飯一次.
這些是不是同行的關係呢?.
我也覺得不算.
因為這個同行.

$^{161}$你覺得你同行即是怎樣呢?.
你是不是同行一起去計劃一些東西?.
你同行是不是一起開會?.
你是不是裡面想籌謀一些計劃?.
原來其實他沒有說.
那個同行他之後是怎樣的.
他只是說有兩個或以上的人.
你經常慣性的來回行走.
這個是一個目的.
你不是來回行走為了甚麼.
而是你可能很無聊.
可能你會一起聊天.
聊心事.
或者大家一起聊天.
我們香港人說.
你無論是甚麼都好.
其實關鍵是我們是一起.
關鍵是我們是經常.
關鍵是我們是慣性.
這個是很重要的.
所以我發現原來在香港.
我們這個城市.
在過去這幾年我們說大了.
有些人其實不是經常.
不是慣性和他一起行走.
但你常常都會用同行這個字.
不過如果再繼續說聖經的話.
其實同行這個字.
是和經常慣性來回行走有關.
而如果今天說與神同行.
其實是說我們是否經常.
很慣性的去與神同行.
當我們用一個字去看的時候.
我們開始用經文.
去試試入獄去看看.
今天特意選取一段金句.
是用另外一個我們不太慣的譯本.
就是和合本2010.
這句譯本我覺得他譯得很好.
他這樣說.

$^{201}$神是信實的.
他呼召你們.
即是我們.
好與他的兒子.
我們的主耶穌基督.
共享團契.
舊約同行這個字.
在新約其實是.
很近團契這個字的意思.
如果用一句說話.
其實首先你要知道.
我們所信的神是信實的.
什麼叫信實呢?.
我們經常說我們的主是信實的.
信實的意思就是.
祂說了就算了.
一言九鼎.
祂不會說謊.
祂和你說得出就是了.
祂聖經記載得就是了.
所以我們信的神是信實的.
而祂呼召我們.
好與祂的兒子共享團契.
如果用繼續說其他的譯本.
甚至原文.
其實呼召後面有一個詞.
就是一個介詞.
就是我們英文的proposition.
就是進入.
所以是我們永生神.
祂是信實的.
祂呼召我們.
進入去.
與祂的獨生兒子.
耶穌基督.
共享團契.
所以當我們說.
是祂呼召我們進入.
而這個是呼召.
弟兄姊妹你們知道嗎?.

$^{241}$這是呼召.
所以當我們有時很好奇.
很想知道.
神啊.
我現在活了幾十歲了.
究竟你對我們生命.
有什麼獨特的呼召.
又真的對.
如果根據詩篇139篇.
其實神對我們每一個人.
祂不單做我們獨特.
祂對我們每一個人.
都有一個很獨特的計劃.
今天你宣召的經文.
你已經看到了.
當我們尋求主對我們的心意.
其實每一個人的呼召是很獨特的.
但原來當我們尋求主.
對我們獨特的呼召之前.
原來在聖經裡面.
其實主.
好像一個建築師一樣.
他建樓.
他不會一建就建十六樓.
他建.
他會由根基.
打樁.
根基開始建上去.
而那個根基.
打樁.
你要留心.
在聖經裡面.
有很多是普世教會.
是我們信耶穌的人.
我們一起擁有的呼召.
是祂呼召我們全部人.
去進入的呼召.
這些呢.
其實我們會稱為.
是一些根基.

$^{281}$你一定要進入這些.
呼召的根基.
接著你慢慢進入.
一些個別.
很獨特.
主是對你的.
祂沒有對你旁邊那個.
祂是對你的.
所以當在那個根基.
那個基礎.
那個普世教會裡面.
我們一起被呼召.
擁有的呼召的時候.
原來其中一個.
就是神.
祂呼召我們進入去.
與祂的獨生兒子.
耶穌基督.
共享團契.
我只是說這段經文.
說團契或者同行.
你最起碼呢.
有兩個的意思.
你看得出來.
一 就是團契或者同行.
祂不是一個.
達到目標的途徑.
祂是目標.
好像我們.
射標靶.
就是射個紅心.
所以.
祂呼召我們進入團契.
其實.
那個是一個目標.
如果我們這樣說.
我們如果用現在我們生活的例子.
信仰的生活來說.
很值得我們.
去深思.

$^{321}$我們有時人越大.
可能現在.
年青人會不會是這樣.
就是他們覺得.
團契很無聊.
我們經常坐在那裡.
四個W.
過去如此 現今如此.
經常都是W.
其實重點不是W不W.
重點就是我們要走在一起.
是我們一起.
與神同行.
所以我們個別.
與神同行之餘.
我們說要多過兩個人.
所以我們原來是彼此需要的情況下.
我們是一起.
同行 一起團契.
而.
哪怕你覺得無無聊聊.
有時甚至你說.
其實我回團契.
是因為我很想我們快點出去傳福音.
只是坐在那裡.
不行的 我們快點出去傳福音.
所以我回團契.
純粹是想激勵弟兄姊妹.
出去傳福音.
傳福音一定是好的.
但原來你不是.
透過團契.
為著達成.
另一個目的.
而是團契和同行.
本身就是一個目的.
所以我們被主.
呼召我們與.
他的獨生愛子.
耶穌基督共享團契的時候.

$^{361}$原來我們要.
享受我們走在一起.
我們享受我們常常.
經常慣性的.
與神同行.
我們又彼此的同行.
這是一個目的.
第二.
剛才說共享團契.
一定有一個.
很關鍵的動詞.
就是共享.
如果用新約聖經.
在紹興傳.
當時耶路撒冷的教會.
在他們經歷.
聖靈之後.
聖經記載他們有一個.
很特別的東西發生了.
在他們中間.
就是他們一起共享.
彼此有的資源.
這個共享.
是不分彼此.
所以當年有錢人.
他會賣地.
供養弟妹.
供養貧苦的.
孤兒寡婦.
但這個共享.
其實不是單單我們人與人之間.
因為如果.
一切萬物.
都是從我們的神而來.
原來祂是首先.
白白的.
與我們共享了.
一樣東西.
就是耶穌基督.
祂是三一神.

$^{401}$祂是耶穌.
祂本身是同祂的父.
並且同聖靈.
同榮同尊.
祂是神.
祂竟然來到這個世界.
祂是與我們分享.
祂的生命.
所以耶穌是生命.
祂是生命的娘.
祂是生命.
祂為我們寫十字架.
祂為我們進入.
救贖的進程.
祂要帶我們.
進入救贖的時候.
原來祂.
首先來與我們分享.
祂自己.
所以弟兄姊妹.
原來當我們.
白白的得來.
所以路加福音和.
馬太福音第十章說.
我們都白白的寫去.
所以主與我們分享.
我們便與人分享.
所以當你.
與人分享的時候.
你望著自己.
我都很好.
我很樂意與人分享.
我所有的一切的時候.
其實你的信任是給自己的.
即是你裡面是對自己的肯定.
你說好不好.
不是不好的.
不過原來如果你真的.
有一個這樣的看見.
你之所以有.

$^{441}$是因為神給你.
所以當神給你.
你付出去.
這真是很理所當然的事.
你說共享甚麼呢.
可以共享.
笑聲.
笑聲滿載溫馨.
不知道你有沒有聽過.
所以彼此分享我們的喜樂.
分享我們的憂愁.
還要分享我們的時間.
分享我們的金錢.
分享我們在神裡面的智慧.
分享神給我們的恩賜.
所以有時你覺得.
為何他唱歌這麼動聽.
為何他彈琴彈得這麼好.
我真的會欣賞不同的人.
因為他手上所有的恩賜.
是主給他.
他樂意與我們分享的時候.
其實不是想製造一樣.
你看看你.
這個又不懂那個又不懂.
其實不是想製造一個這樣的果效.
而是.
我們每個人.
主白白的給我們.
他給了我們甚麼呢.
而我們白白的分享的時候.
我們又和其他人.
共享了甚麼呢.
這個原來是同行.
和我們新約講的.
團契裡面的精髓.
在聖經裡.
有兩個人在舊約.
講得很清楚.
他形容這兩個人是與神同行.

$^{481}$我們今天很快.
看看這兩個人.
然後我們一齊進入下一步.
第一個.
就是爾諾.
剛才大家讀過經文就不再讀了.
現在不要看著我.
看著這段經文問大家三個問題.
看看你今天有多聰明.
爾諾.
他一生有多少年日.
在地上.
多少.
365年.
他幾多歲.
生小孩.
65歲.
這裡記載.
與神同行.
多少年.
38年.
所以用查經的角度.
你可不可以告訴我.
聖經並沒有記載.
他有多少年.
與神同行.
沒有說他是與神同行.
有多少年.
應該是減數.
大聲一點.
對了.
聖經是聰明的.
九點鐘崇拜也很聰明.
你不會回答我300.
365.
有寫嗎.
而是聖經明明記載.
爾諾一生365年.
地上的生命.
但是.

$^{521}$他說他是與神同行.
300年.
即是說一個減數的情況下.
有65年.
並沒有記載他是與神同行.
很清楚.
那65年發生什麼事.
不是很仔細.
有寫.
但你只知道一件事.
就是他65歲那年.
有一件事發生.
跟你旁邊的人說.
是什麼事.
就是他65歲的時候生了一個兒子.
那即是怎樣.
生了一個兒子的時候.
就開始.
與神同行.
嘩.
厲害.
生兒子 生女兒.
開始與神同行.
請你.
現在跟你旁邊的人交流.
或者跟你後面的人交流.
你覺得他發生什麼事.
為什麼.
生了兒子之後.
會開始與神同行呢.
請你交流.
為什麼生了兒子之後.
會突然之間.
與神同行呢.
還要300年.
沒有斷過.
生兒子也很重要.
生女兒也很重要.
發生什麼事呢.
聖經沒有說.

$^{561}$我們只知道一個事實.
就是.
他生了兒子之後.
這個兒子叫馬土撒拉.
跟著他就.
與神同行了.
我們只能猜測.
聖經真的很好看.
真的要看聖經.
從前文後理.
可以參透很多東西.
從脈絡可以參透很多東西.
由舊約到新約.
真的很好看.
還有最特別的是.
看來看去都覺得.
原來神.
神的靈.
想給我們一個空間.
默想是否很棒.
喜歡看聖經.
每一次都覺得歷久彌新.
看來看去.
真的要有些東西可以看.
若涓你快點說.
五歲已經發生什麼事.
我都說我猜測.
你也是猜測.
猜測是什麼呢.
可能他求有子女.
很多年都沒有.
突然之間六十五歲.
終於有兒子了.
心存感恩.
是猜測到的.
所以就開始.
你真是信實的.
你一言九鼎.
神是信實的.
所以我今天開始與你同行.

$^{601}$可能是這樣.
又或者是什麼呢.
我人生第一個兒子.
我怎樣教他.
所以主啊.
究竟怎樣教子女呢.
所以你開始怎樣.
每天好像翹著主的手一樣.
主啊你可不可以告訴我.
其實你怎樣教人.
你可不可以和我分享.
究竟怎樣教子女.
其實關鍵是什麼呢.
可能是次教.
所以開始了同行的關係.
亦都可以是怎樣呢.
天父啊.
我女兒真的很難教.
所以其實.
我不知道怎樣教她.
我只是叫她吃一碗飯.
吃一碗飯.
你可不可以坐定吃一碗飯.
我都教不到她.
所以主啊.
其實我怎樣教.
你快點告訴我奧秘是什麼.
是不是很多.
做父母.
我覺得大家今天早上都醒了.
我只是猜測.
但你就是知道.
原來.
與神同行.
是她生子之後.
開始進入與神同行.
繼續看著這段經文.
回答我最後這個問題.
當她開始65歲與神同行之後.
聖經記載她.

$^{641}$發生什麼事.
有記載的.
沒錯啦.
就是和你隔離關關.
繼續生子女.
是真的.
你快點跟她說.
與神同行原來繼續生子女.
我們是不是繼續生子女呢.
無論幾多歲.
當然可以.
不是肉身也可以生屬靈的下一代.
但原來.
與神同行.
你不要以為如果我開始與神同行.
我就要讀神學.
我就要做傳督人.
我就要做一個很厲害的領袖.
站在港台.
與神同行.
原來是很生活的.
與神同行.
你開始之後.
你不是要很刻意.
去做一些.
你認為很偉大的事.
因為.
偉大奇妙神.
你今天選歌真的很對.
勁敗隊你好像知道我想說什麼.
第一首歌就是.
偉大奇妙神.
宇宙創造主.
所以他已經很偉大.
你不需要刻意.
要做一些很偉大的事.
他明明沒有叫你做.
你已經做了.
找你搞.
你真的需要繼續與神同行的時候.

$^{681}$進入這種關係.
我繼續說.
同行是一個目的.
你不是透過同行.
你不是透過團契.
刻意要做一些事.
不是,而是這個就是一個目的.
弟兄姊妹.
說到這裡.
你講得這麼平淡.
你好像還沒有入獄.
其實我已經入獄了.
為什麼對於香港的教會來說.
其實我們很慘.
有一樣東西很慘.
慘在什麼.
我們經常都覺得.
要快點做一些很偉大的事.
我們經常都覺得.
做完休息一下.
不要休息,繼續做.
很辛苦.
不行,繼續做.
好像一定要用雙手做一些事.
那些才是.
其實香港教會.
最難的就是.
停.
最難就是停.
最難就是休息.
要知道.
耶和華是神.
很難的是.
我們享受.
關係.
我們不是.
利用關係.
達成一些目的.
我們不是利用.
一些領袖.

$^{721}$來很想他.
成就我.
你說領袖.
一定是成就其他人.
是.
但這個是他作為領袖的學習.
但作為.
我們神家裡面.
我們每一個人都是.
我們可不可以一起團契.
一起同行.
這就是一個偉大的呼召.
是神偉大的呼召.
因為神是信實的.
他呼召我們進入去.
與他的獨生兒子.
我們的主耶穌基督.
共享團契的時候.
同行.
團契英文寫得好.
Fellowship.
是一班夥伴.
一班信耶穌的人.
同一條船.
我們一起經常.
慣性的.
來回的行走.
所以你以為.
好像.
沒什麼好說的.
但其實我就是想說.
這個是我們最難.
突破的.
各種關係的狀態.
就是我們不是要.
利用關係.
不是在關係裡期盼.
快點下一個.
而是我們在Here and Now.
在關係裡.

$^{761}$我們要享受彼此.
我們享受彼此分享.
我們享受我們一起.
有同一位神.
他的名字叫耶穌.
我們一起進入去.
這個很深的親密的.
團契同行的關係.
原來.
以律與神同行.
三百年.
另一個人.
聖經記載,就是羅亞.
我們一說羅亞.
我們就一定會知道.
他是做了一隻方舟.
神叫他做一隻方舟.
但我們真的很少去默想.
他為什麼無端端.
做一隻方舟.
他是個起始點.
那個孵化這件事出來.
其實那一刻是發生什麼事呢.
我們試試看.
首先.
這個方舟真的有點難看.
但拉到最闊的.
我想說的瑟特是誰.
瑟特就是亞當的.
第三個兒子.
該隱亞伯之後.
第三個就是瑟特.
由瑟特去到.
閃含瓦弗中間就牽涉.
十代.
你說有十嗎.
因為有四代拉不到那麼闊.
方舟.
先刪卻這四代.
所以瑟特之後.

$^{801}$還有四代才到以洛.
以洛就生了瑪祖撒拉.
開始了與神同行.
瑪祖撒拉就生了拉默.
拉默就生了羅亞.
你就會看到.
原來以洛這個與神同行.
三百年的人.
原來他是羅亞的.
太爺.
對嗎.
曾祖父.
這個括號是什麼呢.
如果你的眼力很強的話.
這個是歲數.
請你們看看.
你們看到哪一個最短.
看到嗎.
知道.
看到.
原來十代人.
他的壽數.
平均來說都八百五十.
八百五十歲.
但就是有一個人.
差那麼遠.
這個就是.
就是他.
以洛.
他只有三百六十五歲.
其實我們說他多少歲.
有些誤導.
為什麼.
你覺得與神同行最短命.
與神同行為什麼最短命.
其實不是.
而是這十代人裡面.
其他人平均壽數八百五十歲.
原來聖經都是記載.
他們的終結都是這樣寫的.

$^{841}$就是這個人共活了.
多少歲就死了.
唯獨一個人.
他沒有這樣描述.
那個人就是誰.
就是以洛.
因為你讀經的時候有讀過.
就是三百六十五歲.
神就將他取去.
原來是聖經裡面.
與神同行的人.
以洛.
他是聖經記載.
這兩個人沒有經歷過死亡.
一個是以利亞.
另一個就是以洛.
所以以洛三百六十五歲.
他不是死了.
他是在地上三百六十五年之後.
神就.
接走了他.
很蒙福.
與神同行.
原來可以這麼蒙福.
第二個就是.
以洛的時候.
當我們說.
以洛.
羅亞.
羅亞方舟.
其實他是什麼呢.
我們要說一下.
原來同行和團契.
有第三樣東西.
是很重要.
原來同行關係.
想像一下同行關係.
現在我們好像與神同行.
好像神牽著你.
你牽著神.

$^{881}$甚至有其他弟兄姊妹.
經常慣性行走的時候.
原來.
同行的關係.
是主要透過同行的關係.
或者團契的關係.
來成為.
他誕生.
屬靈生命的所在.
那個所在.
不是一個地方.
那個是在關係裡.
他要誕生.
屬靈的生命.
他要在同行關係裡.
成為神的.
產房.
incubator.
原來.
孵化箱產房不是一個事工.
產房.
孵化.
原來是一個.
關係.
在一個關係裡孵化生命.
在一個關係.
同行關係裡.
孵化屬靈的事.
嘩.
糟糕了.
說到這裡.
如果我今天.
很冒昧地跟你說.
在神的設計裡.
祂的原創.
並不是.
用會議.
來決定.
重要的事.
如果我跟你說.

$^{921}$神原先的設計.
祂不是要成立委員會.
我們在桌上.
在桌上.
大家舉手投票.
祂本身的原創.
真的不是這樣.
祂本身的原創與神同行.
同行.
我要再說這個字的本身.
其實就是一個團契.
很親密經常的關係.
當你在同行裡.
好像無無聊聊.
好像要享受.
我們走在一起.
的日子.
去到祂忽然之間.
永生神忽然之間.
臨到這個關係.
祂突然之間.
要透過這種同行的關係.
來生畜.
祂要屬靈的事.
這個嬰孩.
這個BB.
原來神祂的旨意.
祂要透過同行的關係.
團契的關係.
生畜.
屬靈的事.
當我們.
看羅亞的時候.
其實很清楚.
羅亞.
聖經記載.
沒有說祂與神同行多少年.
因為你有理由相信.
祂一出生就已經與神同行.
原因是什麼呢?.

$^{961}$當然這裡寫.
下面那個.
拉瑪就是祂的爸爸.
祂活到182歲的時候.
祂生了一個兒子.
祂將他起名叫羅亞.
接著祂對著這個兒子說什麼呢?.
祂說這個兒子必為.
我們的操作和手中的奴婦.
安慰我們.
這操作奴婦是因為耶和華就坐地.
原來希伯來人.
為祂的子女改名.
很多時候.
都有兩個原因.
一.
就是神向他顯現.
告訴他.
你的兒子要改什麼名.
你不陌生.
如果用.
施洗約翰.
他也是主差派天使.
跟約翰的爸爸.
你一年之後要有個兒子.
你要將他改名叫約翰.
所以這個你不陌生.
第二個希伯來人.
爸爸改名的時候.
很多時候都會.
有一個期盼.
有一個寄語.
有一個願望.
甚至他很想向兒子.
發出一個祝福.
所以他就將兒子.
改為羅亞.
而羅亞的名字.
是什麼意思呢?.
竟然就是休息.

$^{1001}$安息.
rest的意思.
為什麼呢?.
原來聖經.
明明記載了.
他說這個兒子必為我們的操作.
和手中的奴婦安慰我們.
為什麼呢?.
什麼意思呢?.
原來翻查原文聖經.
奴婦這個字.
一看微讀聖經.
那些APP.
你不用.
覺得約翰.
為什麼你會認識這麼多原文.
其實沒有的.
你手機有的APP我也有.
我看到你一定看到.
我不是開玩笑的.
你一打開APP.
你就知道我說的是不是真的.
奴婦這個字.
在舊約的聖經只出現過三次.
第三次在這裡.
第一二次在哪裡呢?.
原來在創世記第三章.
就是.
人吃了神.
叫他不要吃的.
分辨善惡樹的果子的時候.
耶和華說.
要像你們一樣.
我要就坐地.
所以當時第一個奴婦記載.
他和夏娃說.
由那天開始.
你要經歷奴婦.
才可以生產.
第二次的奴婦.

$^{1041}$就是耶和華和亞當說.
你從今天開始.
你要終身奴婦.
才有得吃.
所以兩次的奴婦.
真的說得對.
是因為耶和華就坐地.
為什麼耶和華就坐地呢?.
是因為亞當夏娃得罪神.
想不到拉麥.
羅亞的爸爸.
當時.
為什麼他會這樣呢?.
很有理由.
他一直活到182歲.
他活了182年.
他經歷了很多人的奴婦.
你知不知道.
馬土撒拉.
的名字是什麼意思呢?.
你真是猜也猜不到.
如果我不是說到.
我也不會這樣查.
原來馬土撒拉的名字就叫.
猛衝.
即是怎樣呢?.
衝啊!我們衝啊!.
衝完之後我們再衝啊!.
衝啊!衝啊!.
做啊!做啊!.
原來馬土撒拉的名字.
就是猛衝的人的意思.
所以他的兒子拉麥.
他經歷了182年的.
這個奴婦.
他突然之間不知為什麼有一種啟發.
這個啟發就是.
耶和華做人.
第六日做人.
第七日是安息日.

$^{1081}$他定義其實做完人之後.
他要設立.
安息日.
你都要將.
人和他自己.
進入安息.
這個安息也是希伯來書.
所說.
我們無要竭力.
進入去.
安息.
所以弟兄姊妹你們知道嗎?.
原來你信了耶穌.
你說.
你願意一生以祂為主.
當你宣認了之後.
第一件事原來你要學的.
就是與神同行.
與神團契.
並且如何安息.
在祂的裡面.
原來第一件事要學的.
是同行.
或者團契.
和我們如何享受.
這份安息.
當我們在這個安息.
裡面的時候.
原來羅亞他從小到大就是這樣.
他一出生.
來來去去.
經常都在這個安息裡面.
他跟主說了很多無聊的事.
說了很多年.
直到有一天主忽然之間臨到.
跟羅亞說.
羅亞.
因為這個世代盡都事惡.
所以我要用洪水.
毀滅這個世代.

$^{1121}$但我要救你一家.
所以你要為我的名.
來做一隻方舟.
原來永生神.
要透過他和羅亞.
這種同行的關係.
與神同行的關係.
incubate.
birth out.
屬靈的東西.
這個就是.
方舟.
當時羅亞的反應是什麼呢?.
聖經記載.
告訴我們.
就由那一天開始.
他就預備這隻方舟.
砌一隻船.
好像這裡.
還沒完成.
他每天都在砌.
你想像一下.
他旁邊的鄰舍.
他和他一起長大的人.
他的親戚朋友.
你猜有什麼反應呢?.
有些情況很好猜.
我會猜.
你做什麼?.
我做一隻方舟.
什麼叫方舟?.
什麼叫船?.
沒見過.
你是不是懵了?.
你做不做?.
來和我們玩.
主說了.
這個是baby.
他要生產一個baby.
這個baby叫方舟.

$^{1161}$洪水會來.
什麼叫洪水?.
沒見過.
什麼叫洪水?.
什麼叫大船?.
來和我們玩.
一天這樣.
兩天這樣.
一個月又這樣.
一年又這樣.
五年十年.
日日如是.
在希伯來書.
他這樣形容.
他會稱他是一個信心的偉人.
我都覺得是.
但如果.
你想深一層.
原來信心是怎樣來的?.
原來信心.
我們之所以對神有信心.
對神說的話有信心.
是因為.
神是信實的.
而我們怎樣去經歷.
我們不是頭腦知.
我不是祖長告訴我.
神是信實的.
祂說的我都不知道是不是.
而是祖長告訴你.
聖經告訴你.
天天經歷.
和你同行的時候.
祂說了一句話.
天天這樣經歷.
經歷到突然之間.
主要透過這種同行關係.
來成為產房的時候.
你裡面其實都不知道.
是不是你自己對神有信心.

$^{1201}$因為這種關係.
這種關係.
踏實到一個地步.
你根本不可能.
你去懷疑.
是不是啊?.
主,是不是啊?.
因為這一種.
很踏實.
可信的關係.
是令到你不再會.
懷疑.
你旁邊和你同行的那個人.
祂說的話.
是不是真的.
弟兄姊妹.
信心.
是從神的道而來.
神的道.
是從永生神.
而來.
這一種是.
關係.
約翰福音.
曾經有一段經文.
耶穌說.
你們讀聖經.
以為聖經.
裡面有永生.
是的.
但是.
其實整本聖經指向的.
就是我.
耶穌說.
但是你們竟然不肯.
到我來得.
生命.
道已經在這裡.
神的話已經在聖經裡.
但是原來.

$^{1241}$聖經之所以.
很重要.
就是一個話語.
令我們認識神.
祂是怎樣.
但是聖經的重要之處.
就是原來祂指向我們的神.
祂是指向.
這個永生的神.
我們和祂的關係的建立.
是什麼一回事.
祂是鼓勵我們.
勸勉我們.
督責我們.
是要和永生的神.
我們進入這個親密.
慣性的同行的關係.
經常慣性的.
團契的關係.
我今天沒辦法.
再講這裡.
但是.
我想說在同行的關係裡.
弟兄姊妹你們知道嗎.
原來.
神現在.
要在香港的教會.
也都包括.
在列國.
很多的國家.
主現在.
在做一場很大的改革.
這個改革.
就是祂要帶領.
祂的身體.
是祂的身體.
你必須要知道不是我們的身體.
我們走在一起.
我們是在代表耶穌基督的身體.
是祂的身體.

$^{1281}$所以當耶穌現在.
在永活的神寶座的右邊.
為我們祈求的時候.
在地上我們是祂的身體.
所以當現在全球.
在列國我們是.
耶穌基督的身體的時候.
祂現在在做一個很偉大的工程.
是要帶領我們.
進入一個很親密的.
同行的關係.
這個是呼召來的.
是神事信實的.
祂呼召我們進入去.
與耶穌.
作為我們的頭.
和耶穌進入去.
這個很深的.
同行的關係.
很深的團契的關係.
這個是整個的根基.
是整個教會的根基.
祂不是means來的.
祂不是一個途徑來的.
祂不是一個目標來的.
所以其實我們不是透過團契.
達成我們某些的NO.
而是這個.
就是我們的目標.
當我們這樣一路一路.
去同行與神同行.
我們一齊與神同行的時候.
是常常我們都要迎向.
主忽然.
臨到祂的殿.
忽然臨到.
祂的身體.
忽然臨到我們的關係.
這個忽然.
是我們.

$^{1321}$沒得預計.
是我們只能夠帶著一個極大的.
longing.
極大的渴慕.
渴慕.
極大的渴慕.
於是到.
祂忽然的靈火.
來到.
Hallelujah.
我們不再.
說主究竟我們有沒有.
方言.
我不是說說方那些.
我們不再計較.
其實我有沒有方言.
我們不再覺得.
聖靈的.
表達.
應該有人躺下的,怎麼沒有人躺下.
這些其實是.
幾十年之前很多爭論.
很多討論,我們不再活在.
這麼低的層次.
因為我們重點的.
就是我們與神之間的關係.
是否進入一種很親密.
很經常很慣性很渴慕.
祂臨到.
我們在乎.
我們是否真的如實.
進入這一種的關係.
就正如耶穌說.
我上上.
應該說.
祂邀請我們上上在祂的裡面.
祂就上上.
在我們的裡面.
我最後.
我要說這個故事.

$^{1361}$本身.
這個故事我是打算.
十一點才說.
但是想不到.
今早.
這個故事我們很熟悉.
是說.
亞伯拉罕和以撒.
說一下這個故事.
這個故事.
很熟悉.
亞伯拉罕.
我們後人稱他為信心之父.
原來這個人.
你覺得.
他很有信心.
原來我剛才說信心是怎樣來的.
是聽道而來.
所以當神的道.
是誰說的.
因為有神差遣某些人去講道.
所以.
本身那個source.
那個來源是不是那個人.
不是 是神.
所以信心的建造.
是因為那個人或那兩個人.
和神之間的關係.
話說當年.
去到亞伯拉罕有一天.
他說.
你要將你當時.
只有一個兒子.
sorry 兩個.
不過當年這個是神的應許.
是以撒.
神叫亞伯拉罕.
將以撒獻給他.
當時聖經.
記載什麼呢.

$^{1401}$讀給大家聽.
亞伯拉罕把繁祭的柴放在他兒子.
以撒身上.
自己手裡拿著火如刀.
於是二人同行.
以撒對他父親.
亞伯拉罕說.
父親啊.
亞伯拉罕說 我兒我在這裡.
以撒說 請看.
火與柴都有了.
但繁祭的羊羔在哪裡呢.
亞伯拉罕說.
我兒.
神必自己預備作繁祭的羊羔.
於是.
二人同行.
同行.
當時.
聖經記載.
亞伯拉罕和他兒子.
和他兩個家僕.
一起上摩尼亞山.
上上下下的時候.
亞伯拉罕叫他的家僕.
你走啊.
我和兒子會帶著柴.
火和刀我們會上山.
所以只剩下他和兒子.
好像這幅圖一樣.
當他一路上的時候.
突然之間.
兒子和他父親說.
爸爸.
亞伯拉罕說.
兒子啊 我在這裡.
以撒說.
爸爸 你看.
火與柴都有了.
但繁祭的羊羔在哪裡呢.

$^{1441}$為什麼他會這樣問呢.
很簡單.
這個兒子.
他常常看到.
爸爸獻祭.
獻繁祭.
繁祭是什麼呢 燒了.
你獻在壇上.
所有都燒了 這些叫繁祭.
所以對於以撒來說.
他一點都不陌生.
因為他常常看到爸爸獻祭.
所以他會說.
爸爸 如果繁祭需要有羊.
但為什麼沒有羊呢.
爸爸就說.
兒子.
神必會自己.
預備作繁祭的羊羔.
我想問問你.
如果你是以撒.
你當時會想什麼呢.
我真的有理由相信.
以他這麼認識.
爸爸和神的關係.
他從爸爸的口中.
聽過這麼多.
耶和華和他爸爸的故事.
他見證爸爸.
他常常為自己的家.
獻繁祭給永生神.
我有理由相信.
他心裡開始盤算.
為什麼沒有羊呢.
那個羊羔.
那是什麼呢.
來到這裡.
可能小時候.
我們上主教學的時候.
你不知道有沒有這樣想過.

$^{1481}$我小時候上主教學.
我真的有想過.
我一直以為是他爸爸亞伯拉罕.
把兒子綁上祭壇.
但是人大了.
你真的會查經.
原來聖經學家.
在一個研究聖經裡.
告訴我們.
當時以撒最少有.
十五歲.
有些聖經學者說.
他應該大概二十歲.
你想想一個十五歲的.
年輕人.
你嘗試去青少年崇拜.
跟一個十五歲的年輕人說.
爸爸要把你放在祭壇裡燒你.
我現在要綁你.
你猜你這樣說.
他會不會讓你綁.
我肯定昨天那群人不會.
所以是怎樣.
我有理由相信.
以撒第一.
那些柴他不會叫年老的父親揹.
他自己揹柴.
所以我特意選這一張.
自己揹柴.
他作為十五歲的兒子.
他日日見到爸爸.
上上獻祭給神.
我有理由相信.
他心裡開始盤算.
我是不是自己揹住柴.
是之後.
上去祭壇.
我是隻羊羔呢.
哇.
於是.

$^{1521}$二人同行.
弟兄姊妹.
我們一路都以為.
獻祭的是.
亞伯拉罕.
但我今天.
想告訴你.
原來是兩代的人.
一齊的同行.
一齊的團契.
一齊上.
摩利亞山.
獻祭給神.
而這個是活祭.
這個故事的精彩.
是因為這個故事.
他預表了.
一個很偉大的關係.
這個就是我們的父神.
和他的獨生兒子耶穌.
你知道嗎.
你和我都知道.
耶穌是為我們.
寫十字架.
但你知不知道.
這個獻祭.
其實神的愛.
真的很荒謬.
神的愛的荒謬.
為什麼我會這樣形容.
是我們很難.
去用我們人有限的愛情.
去想像.
究竟神的愛.
無條件的愛是怎樣.
因為這位永生神.
他竟然因為我們.
作為一個受造的人.
去得罪他的時候.
他知道在靈界的法則.

$^{1561}$是必須要死.
我們必須要死.
因為我們得罪神.
罪的功架乃是死.
所以.
他不想.
人和他在一個.
完全無銀的裡面.
一個永恆的.
一個分隔.
所以他要猜探他的獨生兒子.
來到這個世界.
替我們去死.
但其實天父.
好像我們用舊約看亞伯拉罕.
如果問你.
你覺得如果他獻祭.
獻耳殺.
你覺得耳殺會痛一點.
還是亞伯拉罕會痛一點呢?.
作為現場.
如果你是父母.
你一定會覺得.
如果耳殺會痛的話.
亞伯拉罕會更痛.
耶穌為我們寫十字架上.
當.
亞伯拉罕和耳殺.
於是二人同行.
當耶穌.
他揹著十字架.
背負著十字架.
上去被釘.
原來我們永生的神.
他同樣和他的獨生兒子.
這個同行.
獻上自己.
為著我們.
的時候.
其實這真是一個很荒謬.

$^{1601}$很難.
明白的一種愛情.
弟兄姊妹.
你知道嗎?.
原來神.
他今天.
由舊約創世記一開始.
到你一路看.
亞伯拉罕和耳殺.
到你再看.
永生神和耶穌.
再到現在.
我們在一個末後的日子.
弟兄姊妹神.
現在呼籲我們.
你願意.
和三一神.
進入這個.
經常慣性.
的親密關係嗎?.
你想嗎?.
你渴望嗎?.
你不要再想著.
我渴望靈火.
我又很想傳福音.
不要再說這些.
傳福音一定是好的.
你先渴望它.
不要為了主.
渴望你想達成什麼.
我渴望你.
我很想名成利就.
你不要說主.
我很渴望你.
因為我很想向上爬.
去到某些高位.
我很想被你祝福.
好.
你可以全部放下嗎?.
你的未來是怎樣?.

$^{1641}$是在神的手上.
但你要的.
是要問你自己現在.
你是不是真的說主.
I'm longing.
for you.
我渴望的是你.
我渴望和你進入這個.
親密的關係.
並且主我要.
和你說.
原來在你的設計裡.
不是透過會議.
甚至乎.
我真的很膽敢告訴你.
不是透過月會.
不是透過委員會.
那些好不好?.
團契就好了.
你最好在義工職員會裡.
和義工職員一起團契就好.
如果你敢.
我就告訴你.
你得著神了.
我忘記了那個叫什麼會.
報告的時候.
我會告訴你.
你最好不要再經事工.
大家同在團契.
你覺得好像很無聊.
你要越過你裡面的無聊.
你要越過覺得.
原來這就是一個目的.
主我要得著你.
並且得著我們的彼此.
當你忽然淋到你的電雨的時候.
主.
我們說.
這裡.
我們就在這裡.

$^{1681}$你生任何東西.
都可以.
你生任何東西.
哪怕是一隻方舟.
主你必然得著.
我們.
給你的這個祭物.
因為.
你配得所有.
\newpage



\section{腓立比書 2:1-5}
\label{sec:032058}
\textbf{建立健康的使命教會 (13) 同心事奉的秘訣}~~~~林耀楠牧師
\newline
\newline
連結: \href{http://www.church.com.hk/acms/content.asp?site=cdc\&op=show\&type=product\&code=032058\&layout=sermon}{\texttt{032058}} ~~~~ 語音日期: 2026-04-26
\newline
\newline
\hyperref[sec:032057]{\small{< < < PREV SERMON < < <}}
~
\hyperref[sec:index_chronic]{\small{[返順時目]}}
~
\hyperref[sec:032071]{\small{> > > NEXT SERMON > > >}}
\newline
\newline
腓立比書 2:1-5
\newline
\begin{longtable}{cl}
\hline
\hline
章節 & 經文 (和合本修訂版)\\
\hline
2:1 & \begin{tabularx}{0.7\textwidth}{X} 所以，在基督裡若有任何勸勉，若有任何愛心的安慰，若有任何聖靈的團契，若有任何慈悲憐憫， \end{tabularx} \\ \\ \relax
2:2 & \begin{tabularx}{0.7\textwidth}{X} 你們就要意志相同，愛心相同，有一致的心思，一致的想法，使我的喜樂得以滿足。 \end{tabularx} \\ \\ \relax
2:3 & \begin{tabularx}{0.7\textwidth}{X} 凡事不可自私自利，不可貪圖虛榮；只要心存謙卑，各人看別人比自己強。 \end{tabularx} \\ \\ \relax
2:4 & \begin{tabularx}{0.7\textwidth}{X} 各人不要單顧自己的事，也要顧別人的事。 \end{tabularx} \\ \\ \relax
2:5 & \begin{tabularx}{0.7\textwidth}{X} 你們當以基督耶穌的心為心： \end{tabularx} \\ \\
[1ex]
\hline
\hline
\end{longtable}
$^{1}$兩位姊妹朋友早晨.
我們很快就會在《屬靈成長》裡面.
五個豬完成了多少個?.
又有人不知道,三個.
那開始第四個了.
第四個《屬靈成長》的題目是甚麼呢?.
紅色那裡是甚麼字?一起讀吧.
好,發揮恩賜,講恩賜有關的東西.
今天這一講是講關於.
同心似鳳的秘訣.
講秘訣都是不容易的.
有些好像很神秘的東西.
很重要的東西,很關鍵的東西.
不是我厲害,是我們一起學.
在今天的經文裡面,希望有一個好的學習.
這裡有多少位加拿大公民?.
差不多一半,可能多一半.
加拿大公民,你做了公民.
要上公民班的,是嗎?.
考公民試的.
你知道公民應有的概念.
有甚麼責任,有甚麼思想,行為.
有甚麼態度都很重要,是嗎?.
今天不是在推廣成為加拿大公民.
你們各自的努力.
你美國公民也好,哪裡公民也好.
都有些公民需要知道的東西.
我們今天不是講加拿大公民的問題.
但「公民」這一詞很特別.
「公民」這一詞在聖經裡面.
特別是今天的經文裡面.
翻譯成為「行事為人」.
我們要成為一個公民的話.
我們的「行事為人」是怎樣的呢?.
是一種公民需要付出的代價.
有沒有人知道這是哪個國家的旗幟呢?.
有沒有人沒見過的?.
多數沒有舉手,即是多數都知道,是嗎?.
這支叫做「基督徒的旗幟」.
你不要以為我拍了這條作出來,不是我作的.

$^{41}$130年前,大約1897年的時候.
在美國紐約市一個主學的老師.
教主學的時候.
發明了這個「基督徒的旗幟」.
之後的十年八年裡面.
其他的基督教裡面的牧者,領導.
就將它更加完美化.
解釋裡面白色代表甚麼?.
藍色代表甚麼?.
紅色甚至是代表甚麼?.
將這個成為「基督徒的旗幟」.
今天就用這個旗幟.
成為我們今天這個講題的年冠的旗幟.
就是天國公民.
我們行事為人應該有甚麼呢?.
有甚麼概念責任?.
跟剛才一樣,有甚麼意念?.
有甚麼行為?.
有甚麼態度?.
就嘗試用今天的經文.
跟著基督所需要的範圍.
來一起學習.
天國公民應該有甚麼概念,態度呢?.
我就用這兩個句子來包圍.
在聖經裡面第一章.
十七節裡說「行事為人」是要跟基督的福音相稱.
「行事為人」是要以基督的心為心.
就在這種概念和態度的包圍裡面.
中間有三樣東西.
總共有五樣東西.
我們要一起學習.
有甚麼責任?.
有甚麼意念?.
有甚麼行為?.
今天所講的在經文裡面.
一起慢慢地去看.
概念「行事為人」.
要跟基督的福音相稱.
當然我們要先明白.
甚麼叫做基督的福音相稱.

$^{81}$羅馬書第六章第三至五節.
古來跟我們說.
我們受洗歸入了基督耶穌的人.
然是叫我們一舉一動.
都有新生的陽息.
是不是很熟悉?.
幾個星期前.
父子有受葬的兄弟姊妹.
他們上浸浴班就會學過這個經文.
你曾經上過浸浴班都會學過.
然而我們有新生的陽息.
新生的陽息令我想起甚麼?.
最近教會有很大的恩典.
領受了神的恩典.
有了新的嬰兒出世.
大家知不知道教會最近多了多少個新的嬰兒?.
很多.
有四,五個,六個.
甚至我認識的人身邊都有新的嬰兒出世.
新的嬰兒出世有新生的陽息.
新的嬰兒出世有甚麼特別的地方呢?.
眼睛碌碌很可愛.
皮膚又白又滑.
手掌又暖又軟.
又軟又舒服.
大家想不想像嬰兒一樣變回皮膚?.
有多少個想的?.
其實你們每個都想.
不好意思舉手.
心裡知道只有想著都沒用.
但這個經文不是表面的新生陽息.
是說我們一舉一動.
不是我們信了耶穌.
得到福音之後.
我們的樣子變了.
我們全部的樣子都沒有變過.
雖然可能寬容了.
但沒有漂亮.
也沒有醜.
是說我們一舉一動.

$^{121}$會有新生的陽息.
藉著洗禮我們歸入耶穌基督的死.
和祂同埋葬.
和耶穌同復活.
在形狀上與耶穌聯合.
我們很珍惜這個福音帶來的福氣.
我們每一個人.
受洗後加入耶穌基督的教會.
和我們的主聯合.
是不是很珍貴?.
我們每一個人都珍惜這件事.
這是福音帶來的好處和結果.
但要留意一件事.
不單止我們個別.
每一個人和我們的主聯合.
因著和主聯合的緣故.
其實我們互相之間.
現在在主的聯合裡.
聯合在一起.
所以我們不單止和主聯合.
我們互相之間.
其實有一個聯合的關係.
我們有什麼責任呢?.
「勸勉,安慰,相交,慈憐」.
所以在基督裡.
既然我們和主聯合的時候.
保羅在這裡和我們說了一連串的事.
那個約詞不是假如有.
而是已經肯定有.
而是在那些情況裡.
在這些「勸勉,安慰,交通,慈悲,憐憫」的情況裡.
是有任何.
是所有的情況裡我們都要做.
我們不會回頭再說有質素的相交.
我們說了幾堂.
這個時候不是回頭再說.
有質素的相交的內容是怎樣去做.
有沒有留意這次保羅在這段經文裡.
說了一件很特別的事.
說了什麼呢?.

$^{161}$說了我們心裡.
有沒有這件事.
今天的主題是「事奉」.
在事奉裡面.
我們有沒有注意到.
觀察到和我們一起事奉的弟兄姊妹.
一些情況是怎樣呢?.
我們和弟兄姊妹一起事奉的時候.
有沒有留意到.
有些弟兄姊妹在掙扎裡面.
在時間的運用裡面.
甚至乎在過度疲勞的裡面.
我們有沒有用我們的心.
用我們的眼去留意.
觀察得到呢?.
有沒有和他們分擔.
他們的勞苦.
我們有沒有留意觀察到.
弟兄姊妹在事奉裡面.
其實有很多的喜樂的.
有很多的學習的.
和我們一起.
我們有沒有和他們一起喜樂.
有沒有觀察到.
嘗試到.
體會到.
感受到他們心裡面那份的喜樂.
有沒有留意到.
弟兄姊妹在事奉裡面遇到困難.
甚至乎受到傷害.
我們的眼.
我們的耳.
我們的心.
有沒有留意到呢?.
我們有沒有想起.
弟兄姊妹的需要呢?.
我們有沒有記掛在裡面呢?.
在事奉的秘訣的第一樣.
用我們的心.
用我們的靈.

$^{201}$提醒我們一件事.
事奉裡面.
我們要關懷弟兄姊妹的手.
我們有沒有上心.
第一個秘訣是什麼呢?.
聽不到.
再來一次.
用了很多心.
很多靈裡面.
我們要做剛才所講的那些事.
第一樣提醒我們的秘訣.
我們要.
上.
心.
原來在事奉裡面.
我們很多時候將我們的目標.
集中在我們做的事情上.
我們肯定我們以後的時間要準時到.
準時完.
我們肯定可以討論的內容討論完.
但是我們有沒有留意到.
弟兄姊妹一起走這條路的時候.
很多的困難.
同樣有很多的喜樂.
有很多的勞苦.
甚至有很多的無奈.
但是有很多的盼望.
有很多的付出.
我們的心.
有沒有留意到.
弟兄姊妹在各方面的需要呢?.
如果我們事奉要做得好的時候.
單是做得好都沒有用.
要上心.
要對弟兄姊妹的情況.
她的處境要上心.
我預備講這個分題的時候.
想起一段的故事.
這段故事是講什麼故事呢?.
耶穌和十二門徒.

$^{241}$在耶穌傳道的早期.
有一次耶穌跟門徒說.
我們去到.
外人的另一邊.
耶穌說我們去到.
外人的另一邊.
路上突然狂風大作.
波浪打入船的裡面.
連船都灌滿了水.
你留心看這張照片.
我覺得這張照片很有意思.
沒有拍到那些.
爬船的那些像.
沒有拍到那些.
船的那些旗桿.
你有沒有看到那些十二門徒的眼神呢?.
他們在看哪裡呢?.
不是看著耶穌.
他們全部向上看著天.
他們是不是很一致呢?.
其實我們在事奉裡面.
好像這十二門徒一樣.
當時其實坐在同一條的船.
我們在教會的事奉.
在我們心靈裡面的事奉.
我們其實是坐在同一條的船.
面對的很多時候是同一個困難.
當風浪來到的時候.
他們仰望天上的神.
但是他們一起講一句話.
那個老師.
我們就快沒命了.
你不理的嗎?.
那個時候還是耶穌傳道的早期.
門徒跟著耶穌.
還有一個很短的時間.
在他們的心理中.
耶穌只是一個老師.
是一個夫子.
他見到耶穌行那些神蹟奇事.

$^{281}$曾經醫病趕鬼.
有講到.
但是去到這個大風大浪.
這麼大的天的.
不要說天災.
在大自然這麼澎湃的情況裡面.
可能他們都沒有想過.
耶穌能夠做什麼.
但是他們知道這就是他們的夫子.
當我們事奉裡面.
無論遇到什麼情況都好.
我們是要同心.
來面對同一條船的那個境況.
我們要向我們的夫子.
來去呼求.
我們都知道故事的結束.
耶穌平定了風浪.
接著門徒說.
這個究竟是誰來的?.
為什麼他連風浪都要聽他的話?.
經歷裡面讓我們更加認識我們的主更加多.
就好像門徒當日認識他們的主更加多一樣.
他們認識當時的主只是一個老師的角色.
老師的身份.
老師的關係.
藉著他們一同事奉的經歷.
一起走這條人生路.
一條新的人生路.
跟隨主的人生路.
他們有機會認識他們的主多一點的時候.
他們就可以同心.
向他們的夫子來去呼求.
老師我們快要沒命了.
我們怎麼辦?.
你不理我們嗎?.
之後他們就開始認識.
這個不單止是一個夫子.
還有一個大能.
大力.
能夠呼風喚雨.

$^{321}$平定風浪的神.
好了,再去到第三個題目.
講到我們的意念.
和思想愛心和心思目的.
保羅說我們同心事奉的時候.
會帶給他喜樂.
讓他的喜樂可以滿足.
我們就要.
我們需要.
我們需要刻意去做的.
是用心機去做的.
不是隨隨便便就做到的.
是要用心機去做的.
做些什麼呢?.
保羅說你們要意念相同.
講思想上的相同.
愛心要相同.
有一樣的心思.
有一樣的意念.
一樣的目的.
那我問你.
老實說.
我們這裡一二百人.
我們的意念和不同嗎?.
和不同嗎?.
不同.
我們心思和不同嗎?.
不同.
我們思維方法和不同嗎?.
不同.
那保羅是不是在要求我們做一樣.
我們不可能做到的事呢?.
這裡有一個關鍵的地方.
就是我們第二個的秘訣.
剛才羅馬書院的保羅跟我們說.
我們有新生的樣式.
我們好像嬰兒一樣.
有新生的樣式.
不過我們一舉一動.
是有新生的樣式.

$^{361}$是和以前不同的.
為什麼我們能夠這樣做呢?.
是因為我們與主聯合.
我們每一個人都與主聯合.
我們每一個人是越來越像我們的主耶穌.
在主裡面.
我們能夠聯合在一起.
是因為我們與主先聯合了.
與頂智慧之間就可以聯合.
在聯合裡面.
我們有一樣的思想.
不是思想的方法一樣.
而我們的思想是靠近主的思想.
我們的愛心是靠近主耶穌.
其愛心相似.
我們的心思.
我們的意念.
我們的目的.
是越來越像我們的主.
我們不僅是互相彼此的影響.
而是我們一起在主的裡面.
被主去模造.
去改變.
雖然我們到現在為止.
仍然各有不同的思維方法.
各有不同的心思意念.
但隨著主給我們的改變.
我們是慢慢越來越像我們的主.
這是我們的目標.
如果我們願意的心.
才能夠達到這個目標的時候.
這就是我們使奉裡面的第二個秘訣.
是什麼呢?.
是憐心.
剛才第一個是要賞心.
第二個秘訣是要去憐心.
繼續看第三個.
古老說我們凡事要怎樣做呢?.
凡事不可結黨.
什麼叫做結黨?.

$^{401}$在政治的環境裡面.
很容易明白什麼叫做結黨.
結黨和偶像崇拜有些相似的地方.
偶像崇拜是怎樣的呢?.
就是覺得偶像所做的一切.
都是對的.
偶像崇拜就看不到偶像所做的錯事.
很多時候你在政治環境裡面.
看到一些領袖.
沒有一個政治領袖是完美的.
就好像我們沒有一個人是完美的一樣.
既然沒有一個政治領袖是完美的時候.
他所做的事就有機會做錯.
但崇拜偶像.
或者偶像崇拜的情況裡面.
就會看不到偶像做錯的事情.
即使偶像做錯了.
都會覺得他做得對.
那結黨又是怎樣的呢?.
結黨就是同在一個黨裡面.
一個黨派裡面.
有不同的心思意念.
本來是有不同的對錯.
不過因為同在一個黨派的時候.
就會隨著黨的意向.
而將他自己本來的意願收起.
偶像崇拜就可能分辨不到對錯.
結黨就其實知道對錯.
不過不願意站在自己對錯的位置.
而是隨著黨派的路向來行走.
我就說你們不要結黨了.
結黨的結果就會將教會分裂.
黨派各有自己的立場.
堅持自己的黨派對對方的黨派是錯的.
很自然就形成分裂.
如果我們在教會裡面.
做著同一樣的事情的時候.
我們就會形成教會的分裂.
我們知道對錯不是在於屬於哪個黨派.
不是在於崇拜哪個偶像.

$^{441}$而是在於聖經告訴我們的對錯標準在哪裡.
如果我們的對錯歸回聖經的標準.
歸回主人蘇蘇教導的真理來作為標準的時候.
我們就不會形成到結黨分派的情況.
除了不可結黨.
第二個不可就是.
不可以貪圖虛浮的榮耀.
我相信大家都明白.
我們事奉的人不單止是台上事奉的.
台下事奉的.
在門口進來的事奉的.
在幕後在家裡面.
在電腦上面的事奉的.
各方面的事奉.
都會去到一個情況.
所有的事奉人員.
都有機會被第一機會去欣賞.
甚至去讚賞.
去多謝.
去鼓勵的時候.
那件事本身沒有錯.
本來是一件好的事情.
人所讚賞鼓勵.
多謝的說話.
我們仍然去接受.
但是我們要知道.
這些都是人的說話.
始終都屬於地上的虛浮的事情.
很容易我們人會走去一個軟弱的地方.
就是被讚賞.
被欣賞多了.
就越來越對於這個欣賞.
成為一個追求.
本來是客觀地去接受.
成為主觀地去追求.
甚至去貪圖.
貪圖這個字很有意思.
不單止去圖.
不單止主動去找.
去拿.

$^{481}$還要去貪圖.
要求多過本來當得的欣賞鼓勵.
不要貪圖.
不要貪心.
多過你應該得的.
不要主動去追求那些人的欣賞鼓勵.
不要貪地上虛浮的榮耀.
這是第二個我們要知道的不可的事情.
第三個.
我們要全心謙卑.
個人看別人比自己強.
剛才那個和本聖經是看別人比自己強.
最新的和本聖經是看別人比自己更加重要.
其實是一個更加適合的翻譯方法.
不是看別人比自己強.
不是看別人比自己厲害的意思.
不是覺得別人比我們厲害的意思.
不是覺得我們要謙卑.
覺得自己不是那麼厲害.
別人比我們厲害.
不是這樣的意思.
是要謙卑自己看別人比我們更加重要.
我們要謙卑自己看別人比我們更加重要.
上個星期我們一班的教會領袖.
叫做Intergen一個跨文化的會議裡面有兩天的學習.
我記得其中一位講員講了一個很真實的地方.
我們事奉小朋友的時候.
大人我們和小朋友說話是怎樣說話的.
很自然地我們和小朋友說話當然從高而下.
還叫我們記得和小朋友說話應該怎樣.
你猜猜.
明白了嗎.
我們要苦就低微的人.
苦就低微的人不單止對小朋友是這樣的.
對我們的老人家也是一樣.
年輕人要聽著.
老人家坐在那裡.
不是說隨時挖你身體起來.
腳不是那麼聽話的.
老人家站起來和你聊天你就應該苦就.

$^{521}$怎樣苦就.
心悠姐坐在那裡的時候你就走過去.
苦就和心悠姐聊天.
全心謙卑看別人比自己更加重要.
我覺得年紀很大的時候有很多好處.
教會聚餐的時候通常是排隊出去吃東西.
有時候不是每次都是.
有時候是怎樣.
年紀大的先出去拿.
我們尊重老人家.
讓老人家先出去拿.
為什麼要這樣做.
老人家走得慢.
筆就筆得慢.
食就食得慢.
尊重老人家.
讓他們先拿.
好事還是不好事.
好事.
但是老人家出去拿的時候要留心一件事.
我還記得有一次.
我去拿披薩.
我很想拿兩塊披薩.
但結果拿了幾塊.
一塊.
為什麼呢.
不要單顧自己的事.
也要顧別人的事.
發覺以前在排隊拿東西吃的時候排得很後.
慢慢每年一退一退.
現在排得很前.
長者先出去拿東西吃.
我就不客氣.
出去拿.
但是長者也要知道自己.
我們跟長者說.
我們不單顧自己的事.
雖然很想拿兩塊披薩吃.
但是要顧別人的事.
免得不夠分.

$^{561}$留一些給下一個拿.
情願拿完之後.
還有剩下的時候.
才出去拿第二塊.
這個就是商戶之間那種.
不驕傲不自私的效果.
我剛才在地上.
撿了兩塊垃圾.
順主耶穌的時候.
三四十年前順主耶穌的時候.
聽那些牧師說.
有兩件事我記得很認真的.
一件是.
教會裡面有一個很特別的情況.
教會裡面.
什麼人都有.
什麼職業都有.
牧師說他最欣賞的就是那些.
在社會上很有地位.
很有身份的人.
來到教會就謙卑自己.
放下自己的身份.
來事奉.
無論你是醫生也好.
工程師也好.
什麼都好.
來到教會就很樂意去服侍.
我坐信的時候就記住這件事.
在教會裡面成長的時候記住這件事.
事奉裡面就要放下.
無論你在世界上有多厲害.
多突出.
多本事都好.
回到教會事奉就要放下自己.
那份驕傲.
全心謙卑.
沒有謙卑心來到事奉.
第二個學習就是.
他說見到有些人.
望見教會地上有垃圾的時候.

$^{601}$會掙扎去撿還是不撿.
你們有沒有人掙扎去撿垃圾還是不撿垃圾.
牧師我們很開心.
我們教會沒有人掙扎這件事.
你沒有人掙扎是撿還是不撿.
你沒有人掙扎是不撿還是不撿.
因為與我無關.
一會兒都會有人掃地.
我們走了之後.
都會有護堂的清潔的吸塵.
要不要撿呢.
因為你沒有掙扎就撿了.
剛才見到有兩件垃圾.
你猜我有沒有撿呢.
牧師有沒有垃圾.
前面地上沒有.
幾十年來我就抓住這件事.
很認真的.
教會初期初信的時候.
見到垃圾當然要撿.
自己初信.
甚至尾信的時候.
見到垃圾都不撿.
即是新兵.
初初去教會.
去下去下開始有事奉.
去下去下開始有地位.
開始有身份.
這個教導就不斷提醒我.
不要因為自己的身份事奉改變了.
見到垃圾就覺得.
這個與自己無關.
地上有垃圾.
是因為神的家.
又需要.
有人.
放下自己.
的身份.
去將地上那些污糟的東西.
撿起來.

$^{641}$扔掉.
這個其實是很簡單的動作.
剛才問大家有沒有掙扎.
我告訴你.
其實有難度在裡面.
剛才那兩個垃圾.
不是同一時間撿的.
第一個垃圾.
是一張白色的碎紙.
很明顯的.
我就想都不想就撿了.
第二個垃圾.
是一支牙籤.
不是很明顯的.
就在碎紙裡面.
我就想了幾分鐘.
原來我撿垃圾.
撿了幾十年.
都還要想幾分鐘.
撿剩不撿呢.
是否需要撿呢.
是需要撿的.
有證據在盒子裡.
在事奉的小事情裡面.
反映了我們的心態是怎樣.
我們的心情是怎樣.
不單顧自己的事.
也要顧別人的事.
最近有一個很好的領袖.
令我很感動.
很多時候我和Karen.
都有機會去不同的餐會.
去完餐會.
聽了很多東西.
通常都九點到十點.
甚至乎晚一點.
平均都去到十點.
最近有兩次機會.
有兩個很好的經歷.
是甚麼呢.

$^{681}$有一次如常的.
我們通常去餐會.
有一位姊妹.
她去Downtown.
坐巴士.
上到去北面的聚會地方.
參加聚會.
完了.
九點到十點.
沒理由叫自己坐巴士回去.
我們很多時候.
都負責載她回到Downtown.
然後再回家.
你可以想像.
是用很多時間.
和經過一天這麼長的時間.
Banquet完了.
還要載人回家.
心裡面或多或少.
都是她的擔子在裡面.
最近有一位弟兄和我說.
Vasilejo 我不知道你們.
原來要載人回家.
如果早知是這樣.
你出聲.
下次記得.
如果有需要.
你告訴我們.
我們可以幫你.
載人回家.
我心裡面很有感動.
弟兄姊妹.
不是單顧自己的事.
也要顧牧者的事.
不只是單顧牧者的事.
也要顧牧者所關顧的人.
其他人的事.
因為他們不認識的人.
都願意放出時間.
來關顧.

$^{721}$顧別人的事.
最近一次.
在姻緣的聚餐裡面.
同樣的情況.
我們已經在計劃裡面.
預備載一個姊妹.
同一個姊妹.
回Downtown的家.
吃著吃著.
這個姊妹跟我們說.
不用你們了.
她這麼奇怪.
為甚麼不用呢.
原來坐她旁邊的一個弟兄.
和在桌上吃飯.
大家都認識的一個弟兄.
主動的offer.
載她回Downtown的家.
全部都不是順路的.
都是特意做這件事.
這位弟兄我相信他.
因為知道我們.
預了要載這個姊妹回家.
他就自己主動地.
去offer載她回家.
讓我和Karen早點回家休息.
我們心裡面又很有感動.
感受到主裡面那份愛.
感受到主裡面那份捨己.
感受到主裡面那份.
不單只顧自己的事.
誰不想早點回家休息.
也顧別人的事.
擺出他們的時間.
有可能說.
要做到這些部落教導.
一點都不容易.
有沒有板可以看.
有沒有榜樣可以學習.
有沒有板可以給你看.

$^{761}$如果我告訴你.
其實我們有板可以看.
你知不知道是誰.
有板可以看是誰.
第三個秘訣.
就是虛心.
什麼叫虛心呢.
看我們的板就知道了.
保羅在講完今天的經文之後.
第二章六至八節的經文.
就講了這段話.
他本來是耶穌基督.
是本有神的形象.
但沒有堅持自己與神同等.
反倒虛己取了奴僕的形象.
成為人的樣式.
既有人的樣子.
就謙卑自己.
全心馴服.
以至於死.
且死在十字架上.
這段經文我不詳細解釋.
因為大家都很熟.
大家留意到那四個字.
反倒虛己這四個字.
是什麼意思呢.
反倒虛己不是主耶穌.
謙卑自己的意思.
不是.
祂不是謙卑自己的意思.
反倒虛己是祂將自己倒空.
將祂自己原本有的東西倒空.
神的形象.
祂就是神自己.
祂本來與神同等.
將祂與神同等的位份.
那份榮耀.
那份身份.
倒空了出來之後.
讓祂自己成為奴僕的樣式.

$^{801}$倒空了祂本來是神的身份.
那份榮耀之後.
祂就可以成為人的樣式.
有奴僕 有謙卑 有馴服.
神本來是主人.
耶穌基督是我們的主.
倒空了主人這個角色之後.
祂就成為奴僕.
倒空了神的樣式.
神的形象之後.
祂就成為人的形象.
祂本來是我們馴服的對象.
倒空了之後.
祂就可以去馴服.
去謙卑.
對我們有什麼學習的地方呢.
很明顯我們是想學.
像耶穌那樣虛己.
主耶穌是倒空了.
祂自己神的位份 榮耀.
我們要倒空我們裡面不好的東西.
我們人裡面就是有不好的東西.
我們有分裂的傾向.
有虛榮的傾向.
有驕傲的傾向.
有自私的傾向.
我們要將這些不好的垃圾倒空.
我們要虛己就倒空了它.
所以第三個秘訣就是要虛心.
倒空我們自己.
倒空我們裡面不好的東西.
讓一些屬於主耶穌的樣式.
好像主耶穌的心思意念.
來到我們裡面.
代替了我們原本不好的東西.
在這個過程裡面.
我們可能都覺得.
我們努力學習的時候.
努力祀奉的時候.
身邊有些人和我們一起祀奉.

$^{841}$他的行事為人.
可能令我們.
廣東話不好說.
不要說俗一點的一句.
很激氣.
可能令我們很激氣.
甚至乎很嘔心.
很不舒服 很不自在.
我們接著給兩段經文.
都是保羅講的兩段經文.
給予我們互相的鼓勵.
保羅說.
在羅馬書.
你是誰啊.
竟然聽見別人的僕人.
全部祀奉的人員.
都是主耶穌的僕人.
我們是誰呢.
我們有什麼資格評論主的僕人呢.
主的僕人站立或跌倒.
在任何情況裡面.
都有他的主人在.
而且這個僕人.
必會站立得穩.
是因為他的主.
能夠使他站穩.
我們見到身邊的.
祀奉的同工.
有跌倒的時候.
又可以令我們激心的時候.
又可以讓我們有困難.
挑戰的時候.
我們可能很想去評論他.
很有傾向去批評他.
很有衝動去.
說他.
去責備他.
保羅說.
我們是誰呢.
我們有什麼資格去評論.

$^{881}$主人的僕人.
他們都是屬於主的.
讓他的主.
來作他要的評斷.
主必要令他站得穩.
有沒有想過.
我們其實.
本身.
都是一個僕人.
同樣的眼光.
反過來看都是一樣的.
當我們覺得身邊的人.
看不順眼.
令我們憂心憂肺的時候.
同一道理.
別人可能看我們不順眼.
覺得我們憂心憂肺.
我們全部都是主耶穌的僕人.
就讓我們的主.
來使我們站得穩.
不要以為我們.
能夠事奉這個身邊的人.
我們就要扶起他.
批判他.
責備他.
雖然我們有責任去鼓勵他.
去提醒他.
但我們不是去論斷.
不是去批判.
讓他的主人.
做他的主人.
當做的.
我們在僕人的時候.
就讓我們的主人.
扶起我們自己.
扶起身邊的弟兄姊妹.
保羅在《哥林多前書》.
同樣有另一次的經文.
保羅在《哥林多前書》第十章.
提到以色列的時代.

$^{921}$這班以色列民多麼反叛.
做一些很不討神許的事情.
後來的人都覺得.
當時的以色列人很失敗.
很沒用.
很沒信心.
很叛逆神.
保羅在《哥林多前書》.
就用這段話來提醒.
當代的信徒.
也提醒今代的信徒.
保羅說這些事.
就是以前的以色列人.
發生在以色列人身上.
經歷神的事.
背叛神的事.
違抗神命令的事.
發生在他們身上.
是要作為尷尬.
寫下來是要警戒.
我們這一代的人.
當代的人.
今代的人.
末世的人.
所以自以為站得穩的人.
必須去謹慎.
免得跌倒.
不要以為我們好得過以色列人很多.
我們和以色列人.
其實是同一樣的人.
我們都會叛逆神的時候.
我們都有違抗神的時候.
我們都有跌倒的時候.
我們都有沒有信心的時候.
不要以為我們自己能夠站得穩.
保羅提醒我們.
要謹慎.
讓以色列人的學習.
經歷成為我們的尷尬提醒.
今天我們同樣的學習.

$^{961}$當代的門徒.
學習提醒我們自己.
最後我們講完了.
天國公民事奉的秘訣.
在概念和態度的包圍裡面.
我們行事為人.
是要與基督的福音相稱.
我們的態度是要.
以基督的心為心.
我們講了三個的秘訣.
在責任上.
我們要做各種美好相交的事情的時候.
我們要賞心.
不是為做而做.
而是賞心地去做.
掛心地去做.
認真地去做.
我們的意念是要相同.
不是絕對一模一樣的相同.
是要與我們主相連.
好讓我們能夠連心.
第三樣就是我們的行為.
要虛心.
不是謙虛的心.
是倒空自己的心.
將一些不好的東西.
一些罪性的東西.
放開.
讓主耶穌的形象.
進入我們的心裡面.
從我們的生命裡面流露出來.
今天我想教給大家的三個秘訣.
就是賞心.
連心.
和虛心.
作為天國公民.
我們行事為人.
其實有一個很簡單的目的.
是朝主的目標來進行.
我們要使主耶穌.

$^{1001}$我們是主的心.
起立.
要讓主的福音.
能夠廣傳.
我們剛才所說的.
那麼多的事奉的事情.
目的不是為了要做什麼.
而是讓主的心起立.
讓主的福音.
能夠廣傳.
可以啟動心祈禱.
主耶穌多謝你.
藉著你的使徒保羅.
他們有好的學習提醒.
我們所看到的.
不僅是要做的事情.
不僅是事奉的崗位.
而看到在一切一切裡面.
有一群我們的弟兄姊妹.
一群我們在主耶穌裡面.
能夠相連的弟兄姊妹.
一群你主你自己所愛的人.
包括我們自己每一個.
我們連在一起的時候.
在你們裡面連在一起的時候.
我們就成為一體.
專心的事奉.
成為一個身體.
一間教會.
要完全的擺相.
同心的事奉.
這就是你給我們的使命.
你給我們的目標.
求你給我們有真正的虛心.
真正能夠倒空自己的心.
願意倒空自己的心.
好像你的倒空一樣.
你倒空你的神聖.
你的榮耀.
我們要倒空我們的罪性.

$^{1041}$我們的軟弱.
讓我們盼望我們.
期盼我們.
越來越像你.
讓你的樣式在我們心靈裡面流露出來.
因為我們心靈裡面有你在.
多謝你給我們恩典.
多謝你聽我們的祈禱.
耶穌基督你得逞的名堂.
Amen.
\newpage



\section{彼得前書 2:2-10}
\label{sec:032071}
\textbf{活石}~~~~李均熊牧師博士
\newline
\newline
連結: \href{http://www.church.com.hk/acms/content.asp?site=cdc\&op=show\&type=product\&code=032071\&layout=sermon}{\texttt{032071}} ~~~~ 語音日期: 2026-05-03
\newline
\newline
\hyperref[sec:032058]{\small{< < < PREV SERMON < < <}}
~
\hyperref[sec:index_chronic]{\small{[返順時目]}}
~
\hyperref[sec:032072]{\small{> > > NEXT SERMON > > >}}
\newline
\newline
彼得前書 2:2-10
\newline
\begin{longtable}{cl}
\hline
\hline
章節 & 經文 (和合本修訂版)\\
\hline
2:2 & \begin{tabularx}{0.7\textwidth}{X} 要愛慕那純淨的靈奶，像初生的嬰孩愛慕奶一樣，好使你們藉著它成長，以致得救， \end{tabularx} \\ \\ \relax
2:3 & \begin{tabularx}{0.7\textwidth}{X} 因為你們已經嘗過主恩的滋味。 \end{tabularx} \\ \\ \relax
2:4 & \begin{tabularx}{0.7\textwidth}{X} 要親近主，他是活石，雖然被人所丟棄，卻是神所揀選、所珍貴的。 \end{tabularx} \\ \\ \relax
2:5 & \begin{tabularx}{0.7\textwidth}{X} 你們作為活石，要被建造成屬靈的殿，成為聖潔的祭司，藉著耶穌基督獻上蒙神悅納的屬靈祭物。 \end{tabularx} \\ \\ \relax
2:6 & \begin{tabularx}{0.7\textwidth}{X} 因為經上說：「看哪，我把一塊石頭放在錫安—一塊蒙揀選、珍貴的房角石；信靠他的人必不蒙羞。」 \end{tabularx} \\ \\ \relax
2:7 & \begin{tabularx}{0.7\textwidth}{X} 所以，這石頭在你們信的人是珍貴的；在那不信的人卻有話說：「匠人所丟棄的石頭已作了房角的頭塊石頭。」 \end{tabularx} \\ \\ \relax
2:8 & \begin{tabularx}{0.7\textwidth}{X} 又說：「作了絆腳的石頭，使人跌倒的磐石。」他們絆跌，因為不順從這道，這也是預定的。 \end{tabularx} \\ \\ \relax
2:9 & \begin{tabularx}{0.7\textwidth}{X} 不過，你們是被揀選的一族，是君尊的祭司，是神聖的國度，是屬神的子民，要使你們宣揚那召你們出黑暗入奇妙光明者的美德。 \end{tabularx} \\ \\ \relax
2:10 & \begin{tabularx}{0.7\textwidth}{X} 「你們從前不是子民，現在卻成了神的子民；從前未曾蒙憐憫，現在卻蒙了憐憫。」 \end{tabularx} \\ \\
[1ex]
\hline
\hline
\end{longtable}
$^{1}$弟兄姊妹平安.
讓我們首先低頭祈禱.
預備我們的心領受聖道.
主要求你就在我們中間.
與我們同在 幫助我們.
明白你的聖言 你的真理.
讓我們在世間造你的門徒的時候.
懂得如何去面對種種的挑戰和困難.
又以我口中的言語 心裡的思念.
成就你的話語 幫助弟兄姊妹認識你更深.
我這樣禱告是奉主耶穌的明智祈求 阿們.
神應該一步步將我們推向喜樂和盼望.
好像一切已經翻過新的一頁.
當今早我戴上這條肩帶的時候.
它是白色的 反轉了.
它是白色的 代表光明.
代表基督的復活.
但是在復活節的第五主日.
復活期的第五主日.
我們帶回一個極度不安的場景.
史提反在石頭底下倒下來.
今天我們開始去想想石頭.
但是這個石頭不是活石.
是殺人的石頭 是處決人的石頭.
這個故事放在活石期的中段.
本身就構成一個問題.
如果基督真是復活的話.
為什麼這個世界會用石頭來回應那些見證復活的人呢?.
經文說 史提反看見上主耶穌站在上帝的右邊.
但他承受的卻不是理解或對話.
而是用來處決他的石頭.
復活並沒有令這個世界立刻變得安全可靠.
耶穌基督在今天的經課裡也說.
你們心裡不要憂愁.
不過這句話不是所有問題解決之後才說的.
而是門徒仍然很困惑 仍然很害怕.
仍然不知道前面的路應該怎樣走.
所以多瑪和肥力都對上主說的話很有疑惑.
復活的信仰並不是一條沒有問題的道路.
而是問題仍然存在.

$^{41}$我們卻仍然在世界裡繼續前行.
其實從史提反開始.
教會已經知道見證復活並不保證免於受傷.
直到我們看今天的書信經課《彼得前書》.
在《彼得前書》成書的時候.
那個傷害形式就不再是公開的傷害.
而是轉化為日常地方性的排擠.
基督徒沒有再每天面對刀劍.
不過可能有些基督徒會面對刀劍,刀劍迫害.
但大部分的時候.
它是日復一日地在生活裡.
被慢慢推出社會的核心.
這種對基督徒的排擠.
在第一世紀是有一個很具體的歷史原因的.
今天就不詳細地跟大家說了.
簡單地說.
在古代的羅馬世界裡.
宗教並不是像我們今天那樣.
是一種私人的信念.
你相信什麼自己選擇.
不是的.
宗教是一個包含著經濟,社會,網絡的行為.
以前他們行業的行會.
可能用今天說的工會 不過他們不是叫工會.
他們每一個行業.
譬如你是裁縫,養魚,捕魚的.
他們有自己的行會.
在這些聚會裡.
他們往往是以制治和共善開始的.
這個既是宗教的行為.
也是一個經濟和社會網絡的關鍵.
當基督徒因為信仰而拒絕參加這些行為的時候.
他們並不是說我不拜就算了.
而是被理解為拒絕.
被視為拒絕一套理所當然的公共責任.
你可以再想想剛過去的清明節.
你的家族可能很多人.
去到墳墓面前上香,鞠躬,一起去吃飯.
但是你說我不去.
那代表什麼呢?.

$^{81}$今天我們可能有些父母開明.
說你信什麼無所謂了 你自己選擇吧.
但是在古代的羅馬世界是不行的.
那是代表你拒絕家族,拒絕行會.
拒絕埋堆,拒絕進入社會.
成為社會的一份子.
所以基督徒並不是犯法.
但是他們漸漸被視為不合群,不可靠.
對這個城市的秩序構成潛在的威脅.
再加上我們獨特的生活方式和倫理選擇.
於是謠言和誤解包圍著他們.
你可以看彼得前書接下來的經文.
他說他們是被謠言毀謗他們的.
古代我們今天有聖餐.
古代他們的聖餐是關門不讓其他人參加的.
接著他們又說吃肉,喝血.
於是古代有人就謠傳他們吃人肉.
甚至拿嬰兒來吃.
這些基督徒很奇怪的.
一方面不參加我們行會的祭祀飲食.
然後自己關上門又不知道吃什麼.
不知道搞什麼.
究竟這些基督徒是什麼人呢?.
他們會不會對社會構成危害呢?.
他們還說相信一個叛國賊叫耶穌.
他是用十字架釘死的.
他對羅馬世界,羅馬帝國會不會構成什麼危害呢?.
正正就是這些誤解,謠言,揣測.
使基督徒不能在當時的城市,社會裡面.
有很正常的生活.
於是彼得前書的信徒.
或閱讀彼得前書的信徒.
其實生活在一個極具磨蝕性的張力裡面.
他們仍然在羅馬帝國的城市裡面居住.
生活好像如常.
但其實越來越不像身處在城市裡面必然的部分.
不是劇烈的衝突.
而是長期的耗損.
迫使人反覆地去自問.
特別是那些基督徒會問.

$^{121}$我活在這裡是否仍然在見證上帝呢?.
我是否真的在建立天國呢?.
我是否真的在建立天國的工程裡面呢?.
在這樣的處境裡面.
彼得選擇用石頭這個意象.
但他刻意不把石頭的焦點放在信徒身上.
他沒有說教會要成為甚麼建築.
也沒有鼓勵信徒要將自己磨成有用的材料.
他首先說的是主基督那位須被人所棄.
卻是神所揀選所寶貴的活石.
在古代的建築裡面.
避棄的石頭意味著甚麼呢?.
就是已經被工匠檢驗過.
然後被判斷,被拒絕.
認為沒有用的,失去任何進入工程的可能.
彼得並沒有否定這個判定的真實性.
反而將之視為理解基督的關鍵.
他指出耶穌正正就是被世界判定為沒有用的那一塊.
然而這個避棄的石頭.
是成了防角石.
決定了整座建築的方向和穩定.
這不是成功的保證.
而是一個顛覆性的重新理解.
這句很AI.
不瞞大家我是用AI幫我寫的.
避棄未必等於走錯路.
不被需要未必代表不被揀選.
更加關鍵的是甚麼呢?.
並不是石頭本身能否帶來改變.
而是在整個建築裡面.
它的取向轉移到誰的手裡.
一塊避棄的石頭.
不會憑自身的力量.
令自己的價值反轉.
它不能為自己爭取甚麼位置.
它只是一個石頭.
它只是被動安放在一個位置.
但是在神的眼中它卻成為了防角石.
防角石這件事我稍後再解釋多些.
經文以一連串被動的動詞描寫這個過程.

$^{161}$避棄,被揀選,被立.
彷彿刻意淡化石頭本身的能動性.
Agency.
它不能自己決定自己做甚麼.
因為它是石頭.
但是我們要將注意力.
放在真正推動建築的主.
上帝的身上.
基督之所以成為防角石.
不是因為世界重新認可,肯定它.
而是因為神在自身的工程裡面.
重新確立了建築的方向和位置.
究竟建築的中心在哪裡.
這個曾被擱置的石頭,被棄的石頭.
還是這樣的轉移裡面.
放置在一個決定整體結構穩定的關鍵位置.
於是防角石的意義.
不在於證明它能夠成功或者可以翻身.
而是揭示一種秩序的轉換.
建築的根基,它的指向.
並不是再由人來決定.
而是交付那位被棄之處行使主權的神.
將基督放在防角石.
在這樣的邏輯下.
彼得進一步就稱信徒是活石.
不是自己去活.
不是自己去移動.
自己去砌成的石頭.
而是被建造成為靈功的石頭.
教會不是一群.
透過群體共識慢慢建造出來的建築物.
而是一個被建造的工程.
信徒的身份不是靠自我的肯定來獲得的.
而是靠基督作為防角石的前提之下.
被放進去,被承托,被安排.
當我們將這段經文放回香港.
今日的社會來讀,不知道大家會不會看到那種類比.
那種類比不是歷史事件的相似.
可能也有些相似.
那種類比是生活的節奏.

$^{201}$我們仍然可以敬拜.
我們仍然可以去教會.
教會仍然存在.
信仰好像也是合法的.
但可能不少信徒會覺得.
信仰的公共位置慢慢收縮.
有些價值不再容易被理解.
有些信念歸類為純粹個人選擇.
這個社會好像不需要基督徒.
這種轉變未必帶來恐懼.
它是慢慢塑造了一種退後修煉的生活姿態.
在這樣的處境裡.
彼得前書很適合我們今天去看.
彼得前書呈現的.
或者大家可以再看一次經文.
不是指向一個重新站回中心位置的道路.
信徒不是被催促去證明自己.
我在這個社會還有價值的.
經盟約我的目光放在防國石那裡.
是主 是回到和主的關係.
是那個石頭 活石依附的地方.
一個尚未完整的工程.
等候被放置的地方.
這樣不是消極的.
這是很具體的生活操練.
在小而實在的信仰群體裡.
我們可以彼此守望.
我們在不確定當中持續的禱告.
在我們仍然可以存在的空間裡.
誠實又不張揚地活出我們的信念.
但是這樣我們可能慢慢會覺得.
自己只是一塊孤立的石頭.
我在今天的教會更像一個施工的工地.
而不是一個完成了的電宇.
大家可以.
不知道有沒有機會.
你看一下崇基禮拜堂興建的時候.
之前他拍了一些照片.
不知道在哪裡看到.
圖書館的黨史的照片會不會看到.

$^{241}$是被興建當中的.
結構未明 方向未清.
但是在這樣的狀態裡.
我們被召繼續前行.
靠近那位被棄.
卻是復活的防角石.
我們等候被建造.
當我們在城市裡生活.
如常地工作.
發現自己越來越不像.
是其中不可或缺的一部份.
不知道大家會不會有一種孤立的感覺.
回到公司 在社會裡面.
感覺到基督徒好像.
你可以想一下在殖民地時期.
教會可以建學校 建醫院.
做很多辭位的工作.
政府和教會有很多聯合的工作.
但是今天.
我不敢說太多.
大家可以去教會學校看看就知道.
我不知道大家會不會覺得.
大家在現在的社會裡面.
有一種被孤立的感覺.
甚至可能因為你覺得.
這個社會好像不太需要基督徒.
你會更加越來越覺得自己.
不是這個城市必然不可或缺的一部份.
這種孤立的感覺就會慢慢形成.
原本理所當然的結連開始鬆懈.
一些關係不再穩固.
信徒亦很容易在這個過程裡.
自己理解為一個孤零零的石頭.
需要自己去證明自己的價值.
自己去承擔所有的壓力.
自己的決定是否應該留在這裡.
這個其實只是彼得前書讀者的處境.
他們不是突然被人趕離開城市.
而是在日常生活裡面慢慢推向邊緣.
這種變化往往令人將結構性裡面的排擠.

$^{281}$化為個人的責任.
當支持的網絡一點一點地鬆開.
人不自覺地以為.
撐住是自己的責任.
撐不住你就失敗.
你就小信.
你沒有信心.
於是本應由群體共同承擔的重量.
漸漸變成個人默默硬撐的重負.
自己擔著.
因為你被孤立.
你覺得和這個社會的連結慢慢被鬆開.
這種孤立感可能又慢慢演變成一種焦慮的感覺.
大家的生活是否充滿著焦慮.
大家的生活是否覺得很孤單.
大家作為一個基督徒.
在不回教會的時間.
你會覺得自己是否好像城市的一部份.
正在貢獻這個世界.
一旦你自己理解為一個孤立的石頭.
信徒就會以為.
如果我撐不住,站不穩,派不上用場.
最終只是被丟棄的命運.
於是那些原本一起承擔的壓力.
就慢慢變成各自面對的責任.
覺得應該是彼此分擔的重量.
不知不覺間推翻到個人的肩頭上.
你自己的選擇.
信耶穌是你自己選擇的.
人在這樣的處境裡.
很容易以為撐得住是本份.
撐不住就是自己的問題.
然而彼得前書所描繪的信徒.
從來沒有被召去單獨去成就什麼.
當然不是被留下來自己去處理自身的命運.
他們所面對的那一個其實已經開始了.
就不在他們掌控的建造過程裡面.
他們是被召進來的.
不是獨力承擔.
而在既有的結構裡面是被放置的.

$^{321}$被承托的.
正因為這樣.
活石這個稱呼本身不是孤立存在的可能性.
石頭之所以能夠承受重量.
不是因為它夠硬淨.
是因為它是被放置在一個結構裡面.
它和其他的石頭彼此的貼合.
彼此的支撐.
彼此的分擔承重.
離開了這種彼此的關係.
石頭在工程上就失去了存在的意義.
你們是活石.
被建造成為靈工.
彼得所描繪的教會不是一群各自站立.
各自硬撐的信徒.
而是基督裡面逐漸被建造.
被承托的整體.
關鍵就在這裡.
我講了那麼久不知道大家是否明白.
在整個結構的承重焦點.
是放在一塊活石身上.
就是那塊防割石.
彼得不斷把大家的視線拉到防割石.
就是基督.
防割石是甚麼呢.
它不可以取代其他的石頭.
它是承托定位較準整座建築方向而穩定的.
大家明白甚麼是防割石嗎.
在古代的建築裡面.
它不會像今天那樣要打地基.
被一個我們看不到的底層結構承托.
然後再向上建築.
在以前的建築裡面.
他們會畫一個坑.
就是正方形的坑.
然後找一塊防割石.
是夠正方的.
夠硬朗的.
放在坑的角落.
防割石.

$^{361}$因為接下來的磚石頭.
就跟著石頭的準則繼續建造下去.
然後力量不是將所有力量集中在石頭上.
而是那些力量如何分散在不同的石頭上.
是靠著這個防割石的排列.
來看到它的分散是怎樣的.
大家如果不知道甚麼.
回家google一下.
防割石或古代的建築是怎樣的.
在古代的建築世界裡面.
它是依賴一塊被準確放置.
份量極重的防割石.
這個石頭不是承擔壓力這麼簡單.
最重要的是它決定了整座建築的石頭.
是怎樣排列.
怎樣被貼合.
怎樣共同分擔重量.
這個防割石放得正確.
這個其他的石頭才能穩固.
這個結構才能穩固.
這個結構才能承擔一切的壓力.
彼得清基督是防割石所描繪的.
並不是一個將重量集中在中心一點.
而是想告訴我們這個建築可以成型的起點.
其餘的石頭是因為對齊.
找對位置分擔重量.
沒有了防割石.
各自站立的石頭就是一堆不相干.
不能夠承重的碎石.
有了防割石.
它就開始成為一座能夠承托壓力的建築.
正因為這樣.
活石才能夠承重.
但卻不是獨自的承重.
我們彼此撐住.
但不是彼此去救贖.
因為基督已經救贖了我們.
活石之所以能夠承重.
不是因為自己能夠承重.
而是因為靠著身邊其他的石頭.

$^{401}$一起靠著防割石彼此的承擔.
所有的一位避棄的防割石.
原本避棄的防割石.
現在它是托著整個工程的基督.
我不知道這樣說完.
大家能否想像到.
在古代建築.
當彼得用防割石這個比喻.
其實是在說什麼呢?.
彼得前書沒有留下給我們一個.
可以做什麼的答案.
而是我們不斷重複重讀.
想像的一個圖像.
我們不是被留下來各自站立的石頭.
也不是要求證明自己能夠堅強.
可以撐得住整座建築的石頭.
我們只是被放在一座仍在建築中的結構裡面.
那個結構能不能夠站立.
從來不是取決於我們每一塊石頭的本身.
而是防割石.
基督已經放定了.
其他的石頭只要被安排.
被承托.
被彼此紀念就可以了.
在這樣的結束裡面.
我們未必知道自己的位置有多重要.
我們未必看到最終完成的成態是怎樣的.
但我們只要仍然留在活石所定的關係裡面.
那個建造就沒有停止.
信心並不是確定自己能不能夠撐得住.
而是容許自己不再單獨承重.
不會急於看見結果.
仍然願意在未完成的狀態裡面繼續被安置.
或者用今日的說法.
可能我們說你去馴服等候.
其實是甚麼意思呢.
你不是沒用.
你不是沒意思的在這個社會裡面.
只不過究竟基督安放在你哪裡.
是很在乎你和基督的關係.

$^{441}$在結束的時候.
我們不是在說一個結論.
而是在說一個動作.
彼得前書描繪的建築.
不是一個項已經完成的工程.
而是一個仍然在進行的現實.
在今日我們所要做的.
不是為這個建築加上最後一句話.
而是走向桌前領袖.
一會兒的聖餐裡面.
讓我們每一個人都被提醒.
我們不是靠自己去站立而成為教會.
不是要硬撐自己.
我們其實要連結在一起.
我們所依靠的始終是那個被棄.
它是被立的基督.
餅被抹開.
杯被遞上.
並不是因為我們找到自己的位置.
而是因為我們願意再一次.
將我們交回到那位基督的生命當中.
放在這個防割石的次序裡面.
一會兒我們來到聖餐的前面.
請大家留意一下站在你身邊的弟兄姊妹.
他們的存在.
記得自己不是孤立的石頭.
我們是一群被呼召,被連結.
仍然在建造的活石.
被建造成為靈宮.
我們來到卓子面前一同的領袖.
同一個身體.
這個身體分給眾人.
我們再依次的按字在彼此之間.
放在基督的裡面.
或許那個建築還沒完成.
但是我們能夠領受這份恩典.
我們要被提醒.
真正托著我們的並不是自己.
而是那個成為我們生命之糧的基督.
一會兒聖餐.

$^{481}$大家留意一下.
我們有彼此.
然後一起禱告.
主啊 你表面上是一個被人所撇棄.
人不喜歡的一塊石頭.
但是神卻是揀選你成為那個防割石.
而我們就像活石一樣要靠近你.
依著你的方向來擺放自己.
但是我們知道不是我們能夠自己決定什麼.
而是我們能夠放手被你使用.
被你安置.
在這個過程裡面.
讓我們明白到我們不是孤單.
自己一個人撐著.
當社會慢慢將基督徒邊緣化的時候.
讓我們明白到我們仍有彼此.
我們需要在基督裡面一起結連.
我們都成為活石.
靠近你這個活石.
以致我們一起被建造成為靈宮.
亦都藉著你去奉獻那個靈祭.
令到上主得到榮耀.
讓上主在世間被人看見.
求主你就幫助我們每一個.
當我們生活感到孤單.
感到壓力 感到焦慮.
叫我們回去和你結連.
亦都回去教會當中.
去感受彼此的承擔和幫助.
讓我們知道我們不是單獨去面對.
這些種種的壓力和挑戰.
主啊求你真的幫助我們.
提醒我們 帶領我們.
這樣我們禱告交託.
是奉主耶穌基督的名字.
阿們.
\newpage



\section{約翰福音 4:1-42}
\label{sec:032072}
\textbf{天涯何處無芳草}~~~~石偉漢牧師
\newline
\newline
連結: \href{http://www.church.com.hk/acms/content.asp?site=cdc\&op=show\&type=product\&code=032072\&layout=sermon}{\texttt{032072}} ~~~~ 語音日期: 2026-05-03
\newline
\newline
\hyperref[sec:032071]{\small{< < < PREV SERMON < < <}}
~
\hyperref[sec:index_chronic]{\small{[返順時目]}}
~
\hyperref[sec:032073]{\small{> > > NEXT SERMON > > >}}
\newline
\newline
約翰福音 4:1-42
\newline
\begin{longtable}{cl}
\hline
\hline
章節 & 經文 (和合本修訂版)\\
\hline
4:1 & \begin{tabularx}{0.7\textwidth}{X} 耶穌知道法利賽人聽見他收門徒和施洗比約翰還多， （ \end{tabularx} \\ \\ \relax
4:2 & \begin{tabularx}{0.7\textwidth}{X} 其實不是耶穌親自施洗，而是他的門徒施洗，） \end{tabularx} \\ \\ \relax
4:3 & \begin{tabularx}{0.7\textwidth}{X} 他就離開猶太，又回加利利去。 \end{tabularx} \\ \\ \relax
4:4 & \begin{tabularx}{0.7\textwidth}{X} 他必須經過撒瑪利亞， \end{tabularx} \\ \\ \relax
4:5 & \begin{tabularx}{0.7\textwidth}{X} 於是到了撒瑪利亞的一座城，名叫敘加，靠近雅各給他兒子約瑟的那塊地。 \end{tabularx} \\ \\ \relax
4:6 & \begin{tabularx}{0.7\textwidth}{X} 雅各井就在那裡；耶穌因旅途疲乏，坐在井旁。那時約是正午。 \end{tabularx} \\ \\ \relax
4:7 & \begin{tabularx}{0.7\textwidth}{X} 有一個撒瑪利亞婦人來打水。耶穌對她說：「請給我水喝。」 \end{tabularx} \\ \\ \relax
4:8 & \begin{tabularx}{0.7\textwidth}{X} 因為那時門徒進城買食物去了。 \end{tabularx} \\ \\ \relax
4:9 & \begin{tabularx}{0.7\textwidth}{X} 撒瑪利亞婦人對他說：「你是猶太人，怎麼向我一個撒瑪利亞女人要水喝呢？」因為猶太人和撒瑪利亞人沒有來往。 \end{tabularx} \\ \\ \relax
4:10 & \begin{tabularx}{0.7\textwidth}{X} 耶穌回答她說：「你若知道神的恩賜，和對你說『請給我水喝』的是誰，你早就會求他，他也早就會給了你活水。」 \end{tabularx} \\ \\ \relax
4:11 & \begin{tabularx}{0.7\textwidth}{X} 婦人對耶穌說：「先生，你沒有打水的器具，井又深，哪裡去取活水呢？ \end{tabularx} \\ \\ \relax
4:12 & \begin{tabularx}{0.7\textwidth}{X} 我們的祖宗雅各把這井留給我們，他自己和兒女以及牲畜都喝這井裡的水，難道你比他還大嗎？」 \end{tabularx} \\ \\ \relax
4:13 & \begin{tabularx}{0.7\textwidth}{X} 耶穌回答，對她說：「凡喝這水的，還要再渴； \end{tabularx} \\ \\ \relax
4:14 & \begin{tabularx}{0.7\textwidth}{X} 誰喝我所賜的水，就永遠不渴。我所賜的水要在他裡面成為泉源，直湧到永生。」 \end{tabularx} \\ \\ \relax
4:15 & \begin{tabularx}{0.7\textwidth}{X} 婦人對他說：「先生，請把這水賜給我，使我不渴，也不用到這裡來打水。」 \end{tabularx} \\ \\ \relax
4:16 & \begin{tabularx}{0.7\textwidth}{X} 耶穌對她說：「你去，叫你的丈夫，再到這裡來。」 \end{tabularx} \\ \\ \relax
4:17 & \begin{tabularx}{0.7\textwidth}{X} 婦人回答，對耶穌說：「我沒有丈夫。」耶穌說：「你說沒有丈夫是對的。 \end{tabularx} \\ \\ \relax
4:18 & \begin{tabularx}{0.7\textwidth}{X} 你已經有過五個丈夫，你現在有的並不是你的丈夫。你這話是真的。」 \end{tabularx} \\ \\ \relax
4:19 & \begin{tabularx}{0.7\textwidth}{X} 婦人對他說：「先生，我看你是一位先知。 \end{tabularx} \\ \\ \relax
4:20 & \begin{tabularx}{0.7\textwidth}{X} 我們的祖宗在這山上敬拜神，你們倒說，應當敬拜的地方是在耶路撒冷。」 \end{tabularx} \\ \\ \relax
4:21 & \begin{tabularx}{0.7\textwidth}{X} 耶穌對她說：「婦人，你要信我。時候將到，你們敬拜父，既不在這山上，也不在耶路撒冷。 \end{tabularx} \\ \\ \relax
4:22 & \begin{tabularx}{0.7\textwidth}{X} 你們所敬拜的，你們不知道；我們所敬拜的，我們知道，因為救恩是從猶太人出來的。 \end{tabularx} \\ \\ \relax
4:23 & \begin{tabularx}{0.7\textwidth}{X} 時候將到，現在就是了，那真正敬拜父的，要用心靈和誠實敬拜他，因為父要這樣的人敬拜他。 \end{tabularx} \\ \\ \relax
4:24 & \begin{tabularx}{0.7\textwidth}{X} 神是靈，所以敬拜他的必須用心靈和誠實敬拜他。」 \end{tabularx} \\ \\ \relax
4:25 & \begin{tabularx}{0.7\textwidth}{X} 婦人對他說：「我知道彌賽亞—就是那稱為基督的—要來；他來了，會把一切的事都告訴我們。」 \end{tabularx} \\ \\ \relax
4:26 & \begin{tabularx}{0.7\textwidth}{X} 耶穌對她說：「我就是，正在跟你說話呢！」 \end{tabularx} \\ \\ \relax
4:27 & \begin{tabularx}{0.7\textwidth}{X} 正在這時，門徒回來了。他們對耶穌正在和一個婦人說話感到驚訝，可是沒有人說：「你要甚麼？」或說：「你為甚麼和她說話？」 \end{tabularx} \\ \\ \relax
4:28 & \begin{tabularx}{0.7\textwidth}{X} 那婦人留下水罐，往城裡去，對眾人說： \end{tabularx} \\ \\ \relax
4:29 & \begin{tabularx}{0.7\textwidth}{X} 「你們來看！有一個人把我素來所做的一切事都說了出來，難道這個人就是基督嗎？」 \end{tabularx} \\ \\ \relax
4:30 & \begin{tabularx}{0.7\textwidth}{X} 他們就出城，來到耶穌那裡。 \end{tabularx} \\ \\ \relax
4:31 & \begin{tabularx}{0.7\textwidth}{X} 就在這個時候，門徒求耶穌說：「拉比，請吃吧。」 \end{tabularx} \\ \\ \relax
4:32 & \begin{tabularx}{0.7\textwidth}{X} 耶穌對他們說：「我有食物吃，是你們不知道的。」 \end{tabularx} \\ \\ \relax
4:33 & \begin{tabularx}{0.7\textwidth}{X} 門徒就彼此說：「難道有人拿甚麼給他吃了嗎？」 \end{tabularx} \\ \\ \relax
4:34 & \begin{tabularx}{0.7\textwidth}{X} 耶穌對他們說：「我的食物就是要遵行差我來那位的旨意，完成他的工作。 \end{tabularx} \\ \\ \relax
4:35 & \begin{tabularx}{0.7\textwidth}{X} 你們不是說『到收割的時候還有四個月』嗎？我告訴你們，舉目向田觀看，莊稼熟了，可以收割了。 \end{tabularx} \\ \\ \relax
4:36 & \begin{tabularx}{0.7\textwidth}{X} 收割的人已經得工錢，為永生儲存五穀，使撒種的和收割的一同快樂。 \end{tabularx} \\ \\ \relax
4:37 & \begin{tabularx}{0.7\textwidth}{X} 『那人撒種，這人收割』，這話可見是真的。 \end{tabularx} \\ \\ \relax
4:38 & \begin{tabularx}{0.7\textwidth}{X} 我差你們去收你們所沒有辛勞的；別人辛勞，你們享受他們辛勞的成果。」 \end{tabularx} \\ \\ \relax
4:39 & \begin{tabularx}{0.7\textwidth}{X} 那城裡有好些撒瑪利亞人信了耶穌，因為那婦人作見證，說：「他把我素來所做的一切事都說了出來。」 \end{tabularx} \\ \\ \relax
4:40 & \begin{tabularx}{0.7\textwidth}{X} 於是撒瑪利亞人來見耶穌，求他在他們那裡住下，他就在那裡住了兩天。 \end{tabularx} \\ \\ \relax
4:41 & \begin{tabularx}{0.7\textwidth}{X} 因為耶穌的話，信的人就更多了。 \end{tabularx} \\ \\ \relax
4:42 & \begin{tabularx}{0.7\textwidth}{X} 他們對那婦人說：「現在我們信，不再是因為你的話，而是我們親自聽見了，知道這人真是世界的救主。」 \end{tabularx} \\ \\
[1ex]
\hline
\hline
\end{longtable}
$^{1}$一起我們有個祈禱.
讓我們能夠聽你的話.
讓我們能夠見識到耶穌.
如何將福音傳給未信你的人.
求主你幫助我們.
讓我們去學習.
我們去思考.
以致我們都能夠承擔使命.
我們願意讓更多人聽到福音.
求主你就是這樣幫助我們.
聽我們祈禱奉耶穌的名字而求.
阿們.
今天和大家分享的題目是.
「天涯何處無芳草」.
聽過這句說話沒有?.
先舉手.
剛剛才聽過這句說話?.
不認識這句說話的?.
這麼多人都不認識?.
其他人未聽過這句說話的?.
原來沒有人聽過這句說話的?.
我很驚訝.
這句說話不是很常見的嗎?.
是呀,很多人都認識這句說話的.
所以我不知道你們是怎樣認識的.
所以我還想問這條問題.
就是通常你聽到別人說.
「天涯何處無芳草」的時候.
通常別人會告訴你下一句是甚麼呢?.
你應該不知道.
你聽都沒有聽過.
坊間有些人說.
這一句的下一句是甚麼呢?.
「何必單戀一枝花?」.
坊間有些人聽過這樣的說法.
有些人為甚麼會這樣說呢?.
是因為有些人在說「失戀」.
「天涯何處無芳草,安慰你,唔使擔心」.
「失戀」只是小事.
正所謂「何必單戀一枝花?」.

$^{41}$總有其他花可以選擇的.
不過我想說清楚.
這一句的下一句其實不是這一句.
「何必單戀一枝花?」.
是人們後來作上去.
為了安慰別人而失戀.
事實上這一句說話不是這樣的處境.
或者這麼說.
「天涯何處無芳草」這句說話.
出處本身都不是情歌.
Marko一聽就知道很熟悉文學的東西.
所以要考驗一下.
或者講解一下.
大家既然沒有聽過.
「天涯何處無芳草」這句說話.
其實是出自大作家的.
在唐宋.
大家知道唐宋.
唐朝和宋朝.
中國有不同的朝代.
其中唐朝和宋朝有很多出名的大文豪.
即是一些文人.
叫唐宋八大家.
唐宋八大家裡面.
其中唐朝有兩個很出名的詩人.
韓愈和柳宗元.
另外有六個是屬於宋朝的.
而今天這個經文裡面.
我們講到的這一句說話.
其實是來自一個很出名的詩人.
就是蘇軾.
大家聽過蘇軾嗎?.
蘇軾就是蘇家有三父子.
大家有聽過嗎?.
這些很文學的東西.
希望大家有些文化氣息.
聽完崇拜.
忘記了牧師講道.
都聽到有些很出名的人出現過.
其中就是蘇軾.

$^{81}$蘇軾就寫了這個詞.
就叫「蝶戀花春景」.
我讀一讀這個「蝶戀花春景」.
其實是說什麼呢?.
他說.
花褪殘紅 青杏少.
燕子飛時 綠水人家擾.
枝上柳棉 吹又少.
天涯何處無芳草.
長女秋千 長愛渡.
長愛行人 長女皆人笑.
笑漸不聞 聲漸悄.
多情卻被無情擄.
大約這首歌.
這首詞就是.
正正就是蘇軾嘗試去按著.
春天時節裡面.
對著那些景色.
那些環境.
嘗試去寫出那個意境出來.
包括了周圍的景色.
還有當中一些人的聲音.
一些狀況 一些心境.
當然有些人就說.
這個詞除了寫春天的景色之外.
會不會也寫一下蘇軾.
可能他自己當時是被貶.
被貶一個地方.
所以他自己也有些記憶.
在正途裡不愉快.
無論怎樣也好.
他就寫了這個「蝶戀花」.
成為一個很出名的詞.
當中剛才我們讀到的.
「天涯何處無芳草」.
就是來自這個詞裡.
說到這個詞裡.
為什麼今天會特別去.
用這個詞作為引言呢.
因為當我看到蘇軾這個詞.

$^{121}$他當然很有意境.
更加重要的就是.
他當中提到有四個景象.
或者四個特別的問題.
這四個問題裡.
我覺得很配合今天.
我們所讀到的經文.
我們今天讀到的經文.
就是來自《韓福金四章》.
裡面關於薩瑪利亞婦人的處境.
我覺得蘇軾裡面用到的.
四個引起我想起的問題.
就很配合我們今天去想.
這個薩瑪利亞人的段落.
特別這個薩瑪利亞人的段落.
我想大家都很熟悉.
經常聽過不少很多次.
關於這個俗色.
簡單來說.
這個薩瑪利亞婦人的段落.
就是很多人.
其實有人信耶穌.
所以有時候.
我們傳福音都會提到這段經文.
但你記不記得這段經文裡.
說到有多少人信耶穌.
三個嗎?.
很多人吧.
很多人.
所以其實這段聖經.
不是只說一個人信耶穌.
我們經常只記住.
有一個女人信耶穌.
就是薩瑪利亞的婦人.
其實你看到最後.
你會發現.
這段聖經就說.
原來有很多在薩瑪利亞城.
這裡的人都信耶穌.
所以今天藉著這段聖經.

$^{161}$幫我們想一想.
特別下個星期福音主日.
我們要傳福音.
如何幫我們想一想傳福音的事情呢?.
所以我很希望藉著.
蘇軾的詞.
我抽了四個他的詞裡.
一些字.
一些概念.
去引申我們今天看.
約翰福音這個段落裡.
所以.
記住我今天不是解釋.
蘇軾這首詞裡的字的意思.
我只是引用裡面的概念.
去幫我們思考.
關於傳福音的事情.
所以今天有四個重點.
第一個重點.
就是「花退殘紅青杏少」.
「燕子飛時綠水人家擾」.
想起什麼呢?.
當我想起這個詞的時候.
想起蘇軾說.
花退殘紅.
是指花瓣落盡.
是指一種植物裡的衰亡狀況.
花的樹葉和花瓣都掉下來了.
似乎是什麼呢?.
時間已經過去了.
好像時間飛逝.
好像時間已經不行了.
但蘇軾又說青杏少.
青杏少是什麼呢?.
就是一些未熟悉的杏.
一些杏樹裡的杏.
正好準備開花.
在春天的景況裡.
慢慢有一個生長的機會.
所以似乎是雙春之情.

$^{201}$但其實是春意盛現.
思考到什麼呢?.
第一個重點的問題是什麼呢?.
就是我們有時候都說時機.
我們說時間.
有時候我們說季節.
我聽到有時候年輕人說.
談戀愛都要說戀愛季節.
戀愛季節就是適合談戀愛的時間.
所以戀愛都要等機會.
但你又會看到年輕人.
經常都有機會.
轉頭又談戀愛.
分手又談戀愛.
真的何必單戀一枝花呢?.
所以有時候我們說戀愛要等機會.
等一個戀愛的季節.
等一個機會.
但我們有時候要問一個真實的問題.
就是機會是否真的成為了過去呢?.
機會是否真的成為了過去呢?.
花瓣可能這一刻落盡.
但同時間杏樹裡的杏.
都慢慢成熟起來.
所以今天第一個重要的問題.
就是全福音是否都要等全福音的季節呢?.
全福音是否需要等全福音的季節.
我們才去全福音呢?.
我們有時候都要問機會何時才到呢?.
有沒有機會可以讓我們去全福音呢?.
今天我們見到耶穌.
在經文裡見到.
耶穌遇上一個什麼呢?.
全福音最差的時機.
耶穌遇上一個全福音最差的時機.
為什麼呢?.
因為你看回經文告訴我們.
在四章六誌說.
「雅各井就在那裡,耶穌因女徒疲乏坐在井旁,那時是正午,而那時耶穌的門徒就進了城去買食物.」.
講什麼呢?.

$^{241}$很多時候我們講耶穌的比喻.
半夜求餅的比喻.
半夜求餅的時候.
那些旅客通常等日落沒那麼曬的時候.
才開始行程.
走一些路途.
以致不用那麼辛苦.
所以那個旅客去到很晚的時間.
才去到一個地方問人拿餅吃.
就不會日光日白走這些路.
但是耶穌要趕著回到加里尼的時候.
耶穌就是日光日白.
白天的時候就是要離開耶路撒冷.
而是上回到上面的加里尼.
走了一段很長的路途.
但是這段路途的時間是什麼呢?.
是日光日白.
是早上上午的時間.
太陽剛剛開始越來越熱的時候.
通常是九點或十點後.
就越來越熱.
到了中午十二點.
真的熱得很厲害.
所以耶穌經歷了一個時間.
由下面耶路撒冷走上去.
途經撒瑪利亞的時候.
其實已經是旅途疲乏.
很累了.
人都很疲累了.
才坐在這個井旁.
而且已經到了中午.
中午代表什麼呢?.
是白天太陽很猛烈的時候.
所以你想一下.
人的心情是怎麼樣呢?.
都是差的.
人的狀態呢.
可能都是暖暖的.
不是一個好的狀態.
而且門徒去了城裡面買一些食物.

$^{281}$代表什麼呢?.
就是還沒吃東西.
走得又累.
食物又還沒回來.
所以你想一下.
你的狀態不是一個OK的狀態.
所以你想一下.
如果你自己有時候身體都不OK.
或者不舒服.
或者精神不好的時候.
你想一下.
你都會想.
這個時候不要傳福音了.
這個時候不要跟別人說耶穌.
因為你自己都不想說耶穌.
時機不對嘛.
我們經常這樣說.
時機不對 遲一點.
會不會這樣呢?.
有時候我們都發現.
有些狀況就是.
我們覺得.
這個時候我都不想跟你說了.
遲一點 有機會.
我再跟你說.
你的機會一定會回來.
有時候又不一定.
但這裡很有趣.
當耶穌遇上一個傳福音最差的時間.
但是耶穌正正在這段經文告訴我們.
祂把握了傳福音最佳的時間.
耶穌把握了傳福音最好的時間.
因為這個時間正正好就是.
撒瑪利亞人這個女人.
一個人來到的時間.
這個時間她來到這個地方.
不需要被人特別去看待.
特別去歧視.
或者特別去看不起她.
而來到這裡.

$^{321}$正正她也沒有四下無人.
沒有人騷擾著耶穌.
就是跟她聊天.
沒有人在旁邊.
分散注意力.
讓耶穌去聊天.
耶穌把握著這個中午熱辣辣的時間.
流了汗的時間.
也要跟這個女人說福音的事情.
今天我們很多時候都說.
可以遲一點 可以遲一點.
等一下 等一下.
但是耶穌怎麼說呢?.
當我們看回後面去到第34節.
去到第38節的時候.
耶穌就說.
我的食物就是要遵行猜我來的那位的旨意.
完成他的工作.
剛才說到耶穌可能很餓.
還沒有食物.
但是耶穌說.
如果對於他來說.
什麼是他的食物呢?.
最重要的食物就是.
遵行那位猜他來的旨意的.
那位神的心意.
耶穌在第35節說.
你們不是說到修國的時候.
還有四個月嗎?.
就是到修國還有四個月.
等一下 慢慢來.
不用急.
還有很多時間.
耶穌說門徒就是這樣說.
到修國的時候還有四個月.
但是耶穌的答案是什麼呢?.
我告訴你們.
莊稼熟了.
可以修國了.
當你以為還要再等的時候.

$^{361}$耶穌說.
不是啊.
修國的時間到了.
莊稼已經熟透.
而莊稼熟透的話.
代表什麼呢?.
你可以工作了.
不用再等了.
修國的人已經得到工錢.
為永生儲存五穀.
然後說.
殺種的和修國的一同快樂.
乃人殺種者人修國.
因為這是一個美好的時刻.
原來這個不是最壞的時刻.
這個反而是一個最美好的時刻.
能夠把握著這個機會.
和撒瑪利亞婦人談道.
讓她能夠信耶穌.
是何等美好.
原來今天神給了我們一個任務.
就是我們要去.
大人信耶穌.
你留意聖經裡面說.
大人信耶穌的時候.
有兩批人.
一批是殺種的.
一批是修國的.
哪一批早一點呢?.
真的想想哪一批早一點呢?.
殺種的.
殺完之後其實不知道有沒有料到.
殺了做了些事.
又不知道自己做得成不成功.
什麼都不知道.
好像沒有那麼大的盼望.
但是做修國的那個.
看到有東西了.
看到有得撿了.
那些和菌已經熟了.

$^{401}$大麥小麥已經來了.
你只要肯做.
就一定有.
只要肯做.
一定有.
要多少有多少.
問題就是你做不做.
所以這裡很有趣.
經文很清楚說.
第三十八字說.
我差你們去修你們所沒有新奴的.
他說神差派人去做修國的那個.
和剛才那個殺種的有什麼分別.
他說我差你們去修你們所沒有新奴的.
其實殺種的那個才是辛苦的那個.
但是今天原來上帝給我們.
做修國的那個不用那麼辛苦.
就是做修國的工作.
就一起分享那種喜樂.
分享一起的快樂.
分享那種新奴的成果.
原來今天上帝吩咐我們去傳福音.
不是要我們去做很辛苦的事情.
我們傳福音其實很多時候.
我們做修國的那個.
我們輕鬆的那個.
問題有時候我們輕鬆的那個都不肯做.
看到都不肯去修國.
就說等多四個月.
等多一會.
等多點時間.
因為我們的理由是什麼.
這個不是最好的時間.
花退殘紅.
但是我們可以換個角度.
從上帝的角度看.
現在什麼稱行小.
或者這麼說.
現在上帝的角度是什麼.
莊稼熟了.

$^{441}$莊稼熟了我們要怎麼樣.
修國了弟兄姊妹.
今天我們返教會返了很久.
我們坐了很久.
但是我們是要回應上帝.
我們要去修國上帝.
要收取的莊稼.
輕輕鬆鬆.
就讓上帝的神家裡面.
豐豐富富.
我們可以做一份子.
所以這裡說.
有時候會這樣想.
父親節母親節.
就是怎樣.
只是那天才請父母喝茶.
有364天.
都對父母一般般.
是不是這樣.
我們拍拖結婚那些.
會不會都是這樣的.
情人節先送花.
你問Tony就不會這樣.
4月12日.
都有送花的.
如果你喜歡的人.
天天都是情人節.
如果你對你的父母孝順.
天天都是父親節母親節.
如果你天天都是愛主的話.
是不是天天都是全福音的季節.
弟兄姊妹.
今天第一個很重要的問題就是.
全福音何愁沒時機呢.
我們不要讓自己找一個藉口.
時機未到.
事後未到.
就放下.
上帝交託給我們.
一個這麼好的收割的時間.

$^{481}$這是第一個.
我們今天要思考的問題.
我們全福音.
會不會經常說.
沒時機沒時機呢.
還是你從上帝的眼光見到.
莊稼熟了呢.
輕輕鬆鬆.
進入上帝的禍場裡面.
第二個重點.
就叫知上柳綿.
知上柳綿.
不是知上.
知上.
讀這個音好一點.
知上柳綿吹又少.
天涯何處無芳草.
我想起那個畫面.
就是怎麼樣呢.
就是知上的柳碎.
已經吹得越來越少了.
吹完柳碎.
但是他說什麼呢.
其實你再看一下.
天涯何處無芳草.
就是什麼呢.
周圍都見到青草.
就是說什麼意思呢.
你以為很難才找得到.
其實呢.
可能周圍都是.
周圍都是.
周圍都會見到青草.
不知道大家有沒有聽過一個.
非洲賣鞋的故事呢.
聽過了吧.
非洲賣鞋的故事裡面.
就是說.
有一次.
很多很多年前.

$^{521}$他說有一間鞋廠.
他就想想.
怎麼去增加公司的生意額.
其中呢.
他就派了兩個員工.
就去了非洲.
看看在非洲裡面呢.
那個.
就是.
有沒有機會呢.
可以做到生意.
其中第一名的銷售員呢.
他去到非洲之後.
他看完整個非洲之後.
他覺得.
非洲的人又沒錢.
又不喜歡穿鞋.
買鞋這門生意呢.
注定什麼.
失敗.
非洲的人都不穿鞋的.
你怎麼買鞋呢.
接著第二個銷售員呢.
就這樣跟他的老闆匯報.
他說.
哇.
非洲是一個未開發的市場.
非洲個個都不穿鞋的.
只要每人買一雙鞋.
買兩雙鞋給他呢.
你都發到盆滿缽滿.
你贊成誰的意見呢.
第二個意見.
是啊.
有些人說.
見到個個都不穿鞋.
他不需要鞋的了.
有些人說.
個個都.
他這個沒鞋.

$^{561}$個個都沒鞋.
他們就最好.
就賣鞋給他.
這個就是.
非洲賣鞋故事的1.0.
今天我講過.
非洲賣鞋故事的2.0.
你就沒聽過了.
2.0講什麼呢.
話說呢.
隔了很多很多年之後.
開始呢.
非洲的人會穿鞋.
是不是.
於是呢.
就.
很多人去非洲做賣鞋的生意.
其中有個年輕人.
就說.
我適不適合.
去非洲做銷售員賣鞋呢.
我適不適合.
去非洲做銷售員賣鞋呢.
其中呢.
他有個老行尊.
老行尊就跟他說.
他說你不適合去.
為什麼.
非洲那裡賣鞋.
已經很多人.
在那裡賣鞋的了.
你賣鞋.
你吃菊種了.
對不對這句話.
很多人都賣鞋的了.
你還賣鞋.
你吃菊種了.
你不要去賣鞋了.
你想一想.
你覺得對還是不對呢.

$^{601}$那個年輕人就說.
不是啊.
雖然.
已經賣鞋的.
越來越蓬勃.
但是呢.
仍然有.
賣鞋的需要的嘛.
仍然有人需要.
換新的鞋的嘛.
仍然有些地方.
是有人還沒穿鞋的嘛.
為什麼我還沒有機會呢.
我應該還有機會.
買到鞋的.
因為我還有對象.
我還有對象.
這個就是2.0的故事.
2.0的故事.
全福音是不是都是這樣呢.
你想想有些地方其實.
那些人沒有聽過福音的.
在二百多年前.
有一個地方.
那裡很少人.
聽過福音的.
聽過耶穌的事情.
這個地方是什麼地方呢.
這個地方叫中國.
是中國.
因為呢.
在那個地方來說.
其實當時還是清朝的時間.
其實雖然有.
有少少福音都叫進去.
不過.
但是真是.
在中國的地方.
能夠傳福音呢.
其實還是很荒蕪的事.

$^{641}$就因為有一些的宣教士.
他看到中國.
正好福音還沒有到.
他不是說.
哇.
那應該沒有人信耶穌.
我不要去這個地方了.
不是啊.
當中有兩位.
其中一個最出名的宣教士.
包括了.
戴德生.
戴德生就是.
屬於英國來到的.
中國內地會的創辦人.
由他來到中國裡面.
去傳福音.
他甘願.
好像中國人一樣的生活.
就來到這個.
福音還沒聽聞的地方裡面.
去告訴別人.
信耶穌是什麼一回事.
就因為他看到.
原來有很多的對象.
另外一個人物就是.
馬禮信.
馬禮信來了.
去了中國200年多一點的時間.
同樣.
也因為他進來中國裡面.
去大大去傳福音.
令到整個的中國裡面.
帶來一個很大的改變.
引致到更加多的宣教士.
去到中國傳福音.
所以中國就開始.
慢慢成為一個福音.
都想不到的地方.
今天.

$^{681}$起碼你知道.
如果你回到國內.
很多的地方其實.
只要沒有人禁止.
他們福音的門都很大.
其實很多人都信耶穌.
很多人都願意信耶穌.
究竟是不是天涯何處無芳草.
你是不是看到.
還有很多的草.
你是不是看到.
還有很多的對象.
都可以傳福音.
200年前沒什麼了.
到現在今天.
已經多了人信耶穌.
是不是就不用傳福音了.
還要不要傳福音.
回大陸回國內.
要嘛.
所以呢.
一個很簡單的問題.
你看這個是什麼地方.
佛羅特達.
幾十年前.
我聽廣哥那些.
經常說幾十年前.
佛羅特達教會都沒什麼了.
可能就是邁阿密.
比較荒蕪一點.
最初沒有教會的時候.
你們說沒有教會.
後來就開了邁阿密的教會.
慢慢又少了.
沒有唐人.
有機會聽到福音.
慢慢開始聽到一點.
聽到一點聽到一點.
現在多了人信耶穌了.
那對於你來說.

$^{721}$你現在覺得.
你還有機會傳福音.
還是已經沒有機會傳福音了.
已經沒有對象了.
你認為現在在佛羅特達.
在你身邊沒有福音對象的.
請舉手.
沒有啊.
那你身邊總有一些人.
是還沒認識的.
總有一些人可以聽福音的.
那我們應該是.
做賣鞋的那個.
還是不賣鞋的那個呢.
我們應該要傳福音的那個.
還是沒人信耶穌的.
不傳福音的那個呢.
是啊.
當你認定了.
周圍都是福音對象.
其實個個都可以成為.
你福音對象的時候.
很不一樣.
耶穌希不是這樣呢.
你看看耶穌怎樣去.
四字真一字說.
耶穌知道法利在人聽見耶穌.
收門徒和施洗比約翰還多.
施洗約翰大家都認識.
施洗約翰在當時.
也有很多人認識.
你記得在福音書裡面說.
那些人特意去河邊.
找施洗約翰去施洗.
很多人爭著去找他.
但這裡告訴我們什麼.
耶穌收的門徒和施洗.
比約翰還多.
為什麼會這樣呢.
為什麼會這麼多人.

$^{761}$能夠跟著耶穌呢.
成為他的門徒呢.
因為耶穌歡迎每一個人.
能夠認識他.
耶穌沒有選擇.
耶穌沒有說.
你不是我的福音對象.
我不和你談.
你不是我的福音對象.
我不和你談.
你不是我的福音對象.
我不和你談.
耶穌和什麼人傳福音呢.
和有錢人.
少年財主.
是吧.
有錢的.
他和他談.
年輕人.
耶穌和他談福音.
談信仰.
耶穌和什麼人.
稅吏.
做官的.
撒該.
甚至那些人不喜歡他.
覺得他貪錢.
猶太人很不喜歡他的.
撒該.
這麼多人不喜歡他的.
撒該.
耶穌也和他談信仰.
傳福音給他.
還有什麼.
愚夫彼得.
史達林傳說.
他們是無知學問的小民.
讀書不多.
彼得 約翰那些做愚夫.
讀書不多.

$^{801}$讀書不多的.
是不是沒有機會聽福音呢.
又不是.
耶穌又選擇他們.
甚至成為重要的門徒.
有些宗教分子.
甚至對耶穌有些人很抗拒.
特別那些法利賽人和文士.
其中法利賽人裡面第三章.
約翰福音第三章.
尼哥底姆或者尼哥德姆.
他就是半夜和耶穌談道.
他對信仰很有興趣.
耶穌不會說.
你都是那批人.
那些當權派.
在宗教上當權人.
我懶得跟你說那麼多.
不是.
耶穌半夜一來就和他談.
他想談就和他談信仰.
和尼哥底姆.
還有.
耶穌會傳福音給小孩子.
傳福音給女人.
女人和小孩子在當時來說.
地位相對比男人低.
但耶穌不介意.
所以耶穌身邊有很多女人.
都願意跟隨耶穌.
甚至耶穌連一些患重病.
被人抗拒的病.
麻風病的人.
耶穌甚至碰他們.
很希望他們得到拯救.
耶穌在傳福音沒有選擇.
他沒有說.
少吧.
沒有什麼對象.
這些又不行.

$^{841}$不是.
在耶穌眼裡.
這些是.
所以今天這段經文裡.
在說耶穌和哪個人傳福音.
和撒瑪利亞的婦人.
撒瑪利亞的婦人是什麼.
她要面對.
她是一個女人.
因為你知道.
耶穌在當時是什麼.
一男一女在井旁.
在當時文化來說.
其實有很多秘技.
不可以輕易.
一男一女單獨相處.
重要是什麼呢.
她是撒瑪利亞的婦人.
耶穌是猶太人.
他們之間有種族分歧.
因為撒瑪利亞人被認定了.
不是純種的猶太人.
所以他們很抗拒.
一些撒瑪利亞人.
還要撒瑪利亞的婦人.
所以你要知道.
就是很大的抗拒.
我看到地圖裡.
指著的地方就是.
剛才我們經文所說的.
耶穌來到罪加這裡.
在罪加這裡.
耶穌和這個人打開話題.
你留意.
是耶穌主動打開話題.
耶穌用一個很簡單的問題.
請給我水喝.
就開始一個話題.
可以嗎?弟兄姊妹.
有時候跟別人.

$^{881}$傳福音也是這樣.
你的東西很好吃.
是什麼來的?.
有時候只是搭搭訕訕.
你很有趣.
你的背包有個公仔.
什麼公仔來的?.
聊聊天.
搭搭訕訕.
你主動.
如果他善良的話.
他願意跟你聯繫的話.
他會給你回應.
原來他就是你的對象.
所以我們教會的Grace很厲害.
Grace專門去熱身.
她肯收她的單張.
就能聊天了.
我們真的要跟她學習.
要找設備.
讓她帶領我們.
就是這樣.
開始這樣的對話.
即使.
在當時來說.
第九節就說.
因為猶太人和撒瑪利亞人沒有往來.
即使好像大家.
搭訕訕訕.
你覺得法蘭翻譯.
溝通不了.
或者不想大家接觸的人.
就是猶太人和撒瑪利亞人.
原來耶穌都願意跟他們往來.
跟他們開一個說話.
以致到後來.
當門徒買完東西回來.
見到耶穌的時候.
二十七節說什麼?.
這班門徒對耶穌.

$^{921}$正在和撒瑪利亞婦人說話.
感到驚訝.
沒想到耶穌竟然和撒瑪利亞婦人.
去聊天.
弟兄姊妹.
今天要問一個問題.
誰是你的福音對象?.
你會不會覺得.
周圍都沒有福音對象?.
我記得在我讀神學院的時候.
在其中有一個科目.
這個科目叫CWT.
叫持續報道訓練.
在這個課程裡.
我們每一個的神學生.
都要求我們要上一些訓練.
陪談.
即是報道的課程.
怎樣和別人開話題.
怎樣傳福音.
講核心的內容.
怎樣跟進.
一個訓練.
而其中這個報道訓練裡.
有一份功課.
你認為有什麼功課?.
要做什麼功課?.
報道訓練的功課是什麼?.
就是要傳福音.
他說你們同學.
幾個人一隊.
但每一個人.
都要找一個福音對象.
找一個福音對象.
和他傳福音.
容易找嗎?.
其實有難度的.
因為當你去讀神學的時候.
你身邊的朋友.
都是信耶穌.

$^{961}$開始讀神學.
做傳道.
那些對象已經沒有了.
但是.
其實當我要做這份功課的時候.
用心想.
其實不是.
我們身邊真的會有對象.
其中我們一班同學.
就找了.
誰去成為我的對象.
就是傳福音.
就找了我祖母.
就是找了我祖母.
這裡可以看到.
師母的樣子.
師母年輕的時候.
就是亞太人.
我也不錯.
也很OK.
現在叛了一點.
但是當時就是.
就是找身邊的人.
總有一個.
就是找了我祖母.
我祖母當時住老人院.
我們一班同學就去.
按著CWT的課程.
就跟著說.
那個課程很簡單.
你不需要說新的東西.
你跟著讀都可以.
你跟著步驟.
跟著說.
跟著讀.
跟著就邀請問她.
你信不信耶穌.
我們就是這樣去到這裡.
找我祖母去做這份功課.
說完之後.

$^{1001}$問她信不信耶穌.
她的答案是什麼呢?.
她說信.
我都很驚訝.
原來她說信.
原來莊稼已經熟透.
是差我們那一步.
就是這樣.
我就帶她去信耶穌.
信了耶穌之後.
每一次探望她.
都和她一起祈禱.
兼顧她.
幫助她.
起碼知道.
有一天她離開世界.
她跟誰回去.
弟兄姊妹.
你有沒有想起.
現在這個畫面.
你身邊有哪一個對象.
可以帶她去教會聽福音呢?.
所以第二條問題就是.
傳福音何愁沒對象.
我們傳福音.
要不要再擔心.
沒有福音的對象.
讓我們和她分享呢?.
哪怕那些.
你覺得沒有可能的人.
其實都可以成為有可能.
所以告訴你.
其中香港黑社會大佬洪漢義.
他就是信了耶穌.
當然現在已經回天了.
但是一些我們認為.
不可能信耶穌的人.
其實在上帝眼裡.
他都是對象.
所以不要讓我們自己的眼光.

$^{1041}$去合擇這些人.
這個人不信.
這個人不會.
這個人沒有機會.
是你去限制了上帝.
不是上帝限制了你.
所以很重要.
今天第二個問題.
不要擔心沒有對象.
認真想想.
找一個對象.
將福音帶給她.
C.
第三個重點.
就是笑漸不聞聲漸悄.
多情卻被無情擄.
笑漸不聞聲漸悄.
就是說那個畫面.
慢慢聲音不同了.
當然在蘇軾這個詞裡.
就是說在牆裡有一個少女.
玩得很開心.
玩千秋玩到差不多.
慢慢聽不到少女的笑聲.
不知道去了哪裡.
但這裡我想帶出來.
我們想一想.
原來我們傳福音的時候.
有時候都會什麼呢.
笑漸不聞聲漸悄.
慢慢我們會聽不到一些笑聲.
聽不到笑聲的原因.
不是因為什麼.
而是如果你能夠傳福音.
傳到她最需要的位置的話.
她不會再在你面前扮笑了.
很多時候我們見到.
面下人與人之間.
很多嬉皮笑臉.
很多笑話等等.

$^{1081}$但我們總是.
藏起了我們的內心.
藏起了我們內心的心情.
真正我們心底裡的事情.
但傳福音我們其中一個策略.
就是我們幫助那個人.
漸漸去探討.
他心裡的底蘊.
多情卻被無情擄.
因為原來有很多人.
其實內心深處.
都有一份需要.
有一份真正的情.
而很多時候.
他表面做出來.
與內心不一樣.
表面做出來.
與內心的實際處境.
其實很不一樣.
所以第三個思考的重點.
就是我們傳福音.
要怎樣呢.
要一層一層.
由外而內進入.
人真正的需要.
由外而內進入.
人真正的需要.
說什麼呢.
就是說.
因為我們很多時候.
反過來我們有時候.
很心急.
很心急要傳福音.
傳到哪一個人.
一見到那個人.
Marco.
你有罪啊.
你一講.
第一句已經這樣說.
你就什麼.

$^{1121}$別人都降起了.
就談不到下一句.
有時候我們要想一下.
我們怎樣去.
能夠談到福音呢.
你能不能夠建立到.
一份信任的關係.
如果他都不信你.
你說什麼都沒用.
反過來他信你.
你說什麼都有用.
真是這樣的.
我們能不能夠建立一份關係.
讓他能夠信得過你.
將他一層一層.
在你面前撕開呢.
這裡就說的.
好像一個洋蔥.
表面OK.
但如果你能夠.
幫他一層一層.
揭開洋蔥的話.
你就會什麼.
看到淚水.
是吧.
切洋蔥會有淚水.
當你能夠幫助一個人.
去到生命裡的深處的時候.
你能夠探到他的淚水.
探到他的掙扎.
你才能夠真正.
幫到那個人信耶穌.
第一次我們要留意.
有時候我們傳福音太急.
什麼都第一時間.
要逼別人.
第一次認識你.
我已經要你信耶穌.
我怎樣能夠信耶穌呢.
我都不知道你是誰.

$^{1161}$但我能不能夠.
和他建立一份關係.
讓他能夠信任我這個人.
我能夠認識他這個人.
我知道什麼是他需要.
也能夠幫他.
鼓勵他能夠.
撕開一層又一層的洋蔥皮.
這也是我個人.
傳福音的一個做法.
很多時候.
我都不會第一時間.
和他說你要信耶穌.
但我會聽他說什麼.
我會花半小時.
一個小時.
一個半小時.
聽他說他的生命故事.
你不知道他的生命故事.
你不知道他掙扎什麼.
你也不會知道.
他什麼原因不肯信耶穌.
第一次我們要探討.
這個人的生命.
就正如我們今天讀到.
這個撒瑪利亞的婦人.
一開始就關心什麼.
我和你都不是什麼人.
一開始就說.
我是撒瑪利亞人.
你是猶太人.
不要搞那麼多.
不要談那麼多.
她是什麼.
撒瑪利亞的女人.
但耶穌一步一步.
幫她開展她的話題.
由最初一個最基本的.
就是耶穌問.
你給我水喝.

$^{1201}$從一個生活的問題.
就是水的問題.
生活我們需要喝水.
我們需要食物.
一些生活的細節開始.
耶穌慢慢幫她探討.
她的信仰.
所以你發現.
這個女人開始有些轉變.
當耶穌願意一步一步和她聊天.
她想聊什麼話題.
耶穌都和她聊.
聊到差不多.
耶穌引入她信仰的問題.
耶穌由生活的問題.
慢慢引導她去到生存.
甚至生命的問題.
這個女人會有轉變.
有什麼轉變呢?.
一開始撒瑪利亞的婦人.
會叫耶穌.
你是猶太人.
為什麼要向撒瑪利亞的婦人要水呢?.
分得很清楚.
你是猶太人.
我是撒瑪利亞的婦人.
沒有關係.
分得很清楚.
但到了後來.
聊到差不多.
撒瑪利亞的婦人說.
先生.
是不是禮貌了很多?.
先生.
難道你比她還大?.
她開始對你感到好奇.
你好像知道一些我不知道的東西.
你好像能夠給我一些.
我想尋求的答案.
當她有一份渴慕的時候.

$^{1241}$你開始和她聊天.
就越來越接近了.
先生.
你是不是比她還大呢?.
原來她知道你真的有那些水.
不是那些水.
不是井的水.
原來你還有另外一些好的水.
好的牌子的水.
是活水.
這牌子的水.
可不可以分一些給我?.
於是她說.
先生.
請把那水賜給我.
善我不可.
她很想在你生命裡.
拿到這個水.
這個時候.
你給她這個水.
就為她送上橄欖枝.
她甘心樂意拿.
之前.
你勉強給她一份禮物.
你要了.
不斷給她.
她很抗拒.
現在你給我吧.
你給我吧.
我真的很想要.
是她求你.
所以有時傳福音.
是她求你.
要你給她福音.
只要你能夠談到.
她裡面需要的位置.
你要花時間和她.
逐層逐層剝開.
讓她有信任去揭開.
她的裡面.

$^{1281}$所以去到後面.
不單止是水的問題.
還有她的婚姻的問題.
她向耶穌承認.
她沒有丈夫.
她的婚姻出現了問題.
其實她沒有丈夫.
接著.
後來她說.
先生.
我看你是一位仙子.
又不同了.
看耶穌不單止是一位先生.
看耶穌是一位仙子.
你是從神而來的.
神派你來到我當中.
給我一些提示.
給我一些幫助.
是吧.
我們知道我們有一個身份.
就是從神而來的.
去祝福她.
接著最後.
她說.
我知道彌賽亞就是那位.
稱為基督的要來.
他來了.
會把一切事都告訴我.
到了後來.
她問一個問題.
難道你就是基督?.
她開始認出耶穌.
她知道耶穌就是基督.
就是真正能夠幫她處理.
她生命裡.
底牌裡的問題那位.
就是耶穌.
等一下.
撒瑪利亞婦人.
在以往.

$^{1321}$她一直都在逃避.
她生命的問題.
她生活的問題.
她嘗試去掩飾她的瘡疤.
包括了她的婚姻的問題.
她生活的問題.
很多被人排斥的問題.
甚至找不到敬拜的那位問題.
耶穌就是這樣幫助她.
由最初的渴求水.
帶她去認識到自己.
其實去渴求一份安全的婚姻關係.
但原來這個婚姻.
都不能夠帶來給她生命裡的安穩.
所以她有五個丈夫.
經歷了五段的婚姻.
都找不到那份安穩.
耶穌再幫她帶到.
原來她需要一份信仰.
一份信仰幫助她.
而這份信仰裡.
她要認識到.
哪一個才是真正能夠幫助她的神.
你是否就是基督呢?.
第一姐妹.
你知不知道自己的生命.
只是渴求些什麼呢?.
你知不知道自己的生命.
只是渴求些什麼呢?.
第一姐妹.
有時候我們回教會.
回了很久.
我們稱為Church Goer.
會離教會的人.
我們是一個會離教會的人.
但有時候我們又不等同.
我們是信耶穌的人.
因為有時候我們來到教會裡.
我們仍然帶著很多的盔甲.
戰衣去保護自己.

$^{1361}$一層層的皮去包著自己.
但有時候我們坐了很久.
聽了很久.
我們都不肯在耶穌面前.
打開我們一層一層的皮.
第一姐妹真的要認真想想.
究竟我們真正的渴求是什麼呢?.
如果當你能夠認識到耶穌.
滿督你真正的渴求的話.
你知道耶穌你真的很真.
我真的願意將生命交給你.
求主幫助我們.
不只是來到教會這麼簡單.
也讓我們能夠認定.
耶穌是我們生命的主.
幫助我們.
也幫助我們去幫助人的時候.
帶他去認識真正的主.
不再是生活的問題.
你可以幫助他很多生活的問題.
但你要幫助他轉化.
去處理他生命的問題.
這就是我們全福音最需要做的事.
所以第三個重點的總結是什麼呢?.
全福音何愁沒需要呢?.
每個人生命一定有需要.
每個人生命的底牌裡.
一定有一種需要.
可能會有一點點不同.
但我自己的經驗就是.
每個人的生命裡.
需要上帝的份的愛.
你將他帶到上帝的愛裡.
捉到他的痛點.
他願意在神的面前去降服.
我們不用擔心.
這個人沒有需要信耶穌.
不要跟他說.
不是的 他有.
不過他包住了.

$^{1401}$不讓你知道.
你給他一點時間.
陪他走 幫他轉化.
幫他去揭露在上帝面前.
最後第四點.
長女秋千 長外杜.
長外行人 長女皆人少.
這裡重複了一個字.
就是「長」.
就是一長之隔.
當然這首詞裡的意境.
就是在場內有個少女.
玩秋千 笑得很開心.
後來慢慢聽不到她的聲音.
不知道去了哪裡.
很擔心.
在外面偷聽.
聽不到她的聲音.
很混亂 究竟怎樣呢?.
這個女孩子在裡面怎樣呢?.
有時候我們信仰都會這樣.
什麼意思呢?.
有很多人在場外.
不知道你覺不覺得.
什麼意思呢?.
有很多人見到信耶穌的人.
很開心.
首先你是不是信耶穌很開心呢?.
有很多人見到信耶穌的人.
特別開心 特別喜悅.
很有喜樂.
就好像這樣.
見到場內的人充滿佳人笑.
而很多人其實在場外.
聽到你笑.
他又想知道你為什麼笑.
但他總是覺得欠了一面牆.
又想知道.
又進不來.
又不知怎樣.

$^{1441}$有些時候就是這樣.
所以未信耶穌的時候.
在外面.
他又想探求多些.
但他又不明白.
唯有怎樣呢?.
你幫他打破這堵牆.
讓他進來.
他就知道為什麼這麼好笑.
你讓他進來.
教會裡面 進神的裡面.
接受耶穌.
他就知道為什麼信耶穌的人.
經常這麼好笑.
這麼有喜樂.
要打破這堵牆.
這堵牆是怎樣來的呢?.
是我們的罪.
是我們的罪引致這堵牆.
但誰能夠打破這堵牆呢?.
是不是我呢?.
今天很多人傳福音.
為什麼不肯去傳福音呢?.
就是因為覺得.
我都不會說話.
我又什麼都做不到.
我又不是讀書多.
我又不是有錢.
我又不是什麼.
那你就搞錯了.
很多人都搞錯了.
傳福音不是在於你是誰.
弟兄姊妹.
是在於神是誰.
我們經常都以為.
我都不夠厲害.
不是要你厲害.
弟兄姊妹.
如果他厲害.
他不是信耶穌 而是信你.

$^{1481}$搞錯了.
有些人說信牧師.
信牧師會死的.
如果你在美國.
信牧師介紹你買樓.
你應該死定了.
一定選錯樓.
我怎樣會呢?.
我有很多東西都不會的.
不是靠牧師的.
靠上帝的.
這間教會能夠維持到現在.
不是靠我的.
我沒有來的時候.
你們都在.
是靠神去建立這個教會的.
有沒有搞清楚?.
弟兄姊妹.
打破這道牆的那個.
不是你.
你是流通管道.
你只是一個中介.
以弗所書怎樣說.
從前你們是遠離上帝的人.
如今在基督耶穌裡面.
靠著耶穌基督的血.
已經得以親近.
因為他自己.
耶穌基督自己.
是我們的和平.
使雙方合二為一.
是耶穌基督拆毀了.
中間隔絕的牆.
他自己.
而且以自己的身體.
終止了冤仇.
是他廢掉那個.
寄在律法上的規條.
要使我們兩方.
藉著自己.

$^{1521}$造成一個新人.
促成了和平.
弟兄姊妹說得很清楚.
原來打斷這道牆.
不是靠我們.
是耶穌基督的寶血.
耶穌基督的救贖.
耶穌基督的拯救.
將那個萬子打開.
是他將人與人之間的牆.
人與神之間的牆拆開.
使人可以進到神裡面.
是耶穌自己.
千萬不要將自己.
抬高到等同神.
傳福音當你說.
我不行的.
那你搞錯了.
是上帝.
是神自己.
神要藉著你.
不是靠你行不行.
但當你願意成為一個渠道.
成為一個媒介的話.
那又很不同.
當你很清楚.
自己是一個媒介的話.
這個打破的過程裡面.
就很不一樣.
我們試試看.
那個撒瑪利亞婦人.
後來信耶穌是怎樣呢.
她說.
她經歷了和耶穌的對話的時候.
她說我知道彌賽亞.
那個稱為基督的要來.
他來了.
會將這一切的事.
都會告訴我.
耶穌就告訴她.

$^{1561}$我就是.
和你說話的.
我就是那個的神.
當然這裡補充一點.
我就是這個字.
其實就是我是.
我是這個字.
是很有意思的.
因為在聖經讀到.
出埃及記的時候.
神就說什麼.
我是.
I am that I am.
我就是我.
所以當耶穌說.
我是的時候.
正好回想到.
以色列裡面.
說出埃及記的神.
就是我是.
耶穌說其實我就是那個.
所以去到後來.
這個婦人就怎樣呢.
她能夠去.
經歷了耶穌.
打開那個牆.
打開她的困擾.
解鎖她的一切的時候.
第28節就說.
這個婦人就留下了一個水罐.
去到城裡面.
對眾人說.
你們看看.
有一個人將我.
數來所做的事說出來.
難道這個人就是基督嗎.
她知道.
這個就是基督.
她知道她就是基督.
是耶穌基督幫她.

$^{1601}$打破了那個牆.
她不需要再想.
是不是在撒瑪利亞.
還是耶路撒冷.
去敬拜神.
她知道.
她現在就能夠去敬拜神.
所以撒瑪利亞婦人.
能夠打破那個阻隔的牆.
但是你再留意這段聖經.
還有一個層面.
一個層面就是什麼呢.
就是什麼能夠打破.
撒瑪利亞人阻隔的牆呢.
因為除了撒瑪利亞.
這個女人.
原來還有一群撒瑪利亞人.
這群撒瑪利亞人.
後來都能夠去相信耶穌.
為什麼呢.
我們看看第39節.
那城裡有好些撒瑪利亞人.
信了耶穌.
為什麼呢.
耶穌都沒有接觸過這群人.
因為那婦人作見證說.
她把我所做的一切都說了出來.
原來這個牆.
這群撒瑪利亞人打破的牆.
就是用這個撒瑪利亞.
這個信了耶穌的女人.
藉著她的見證.
她的分享.
她怎麼見到耶穌.
就由她去告訴你.
就是這樣.
這群人.
第40節就說.
這群撒瑪利亞人.
就特意再要來見耶穌.

$^{1641}$然後求她住在那裡.
就在那裡住了兩天.
第41節說得很清楚.
讀清楚第1節.
因為耶穌的話.
信的人就更多了.
第42節說.
現在我們信.
不是因為你的話.
不是因為那個婦女.
那個女人的說話.
而是我們親自聽見了.
知道這人真是世界的救主.
這個婦人成為一個渠道.
吸引這群人.
親自去見到耶穌.
我們做見證.
就是吸引人.
親自去尋找耶穌.
藉著這群人.
親自去見到.
親自聽到耶穌所說的話.
他就知道耶穌是那個救主.
認定他是救主.
他們願意信耶穌.
我們做我們的責任.
出來講見證.
我們告訴別人.
耶穌怎樣也好.
但你記住.
最後拆毀牆的.
是耶穌自己.
經過我們之後.
他最後仍然要去到.
耶穌面前.
他才能真正得到.
那種生命的釋放.
所以不要再看自己.
行不行了.
大人上教會.

$^{1681}$大人信耶穌.
不要再看自己行不行了.
我們只是其中一個.
流通管子的渠道.
所以今天第四條的問題.
就是.
全福音何愁沒真道呢?.
我們全福音經常說.
我不懂說.
我不懂說什麼道.
什麼道理.
全福音不用擔心你說什麼道理.
你說你怎樣經歷神.
告訴他.
讓他能夠透過你.
被引導去到耶穌面前.
讓他親自聽到耶穌.
所以盼望.
我們下個星期.
就是母親節的福音主日.
我希望你認真想想.
究竟四個問題.
第一個問題.
是不是沒有時機呢?.
如果你說沒有時機的話.
下個星期是福音主日.
那這個怎樣也是時機吧.
應該是時候你邀請人來吧.
是你帶你的街坊.
帶你的朋友.
帶你的親人.
聽一次福音.
是不是擔心沒有對象.
除非你住在荒島.
不然你總有認識的人.
你總有朋友的.
特別是華人.
總有一些街坊.
一些華人認識的.
總有一些人.

$^{1721}$你邀請他們來.
第三就是.
你是不是擔心他不會相信耶穌.
他也不是這樣的人.
不是你決定的.
只不過你看不到.
他深層的需要.
耶穌看到.
你有耶穌親自跟他說.
第四.
不用擔心.
沒有道理.
不懂怎樣說.
因為最後不是你說.
是牧師說.
牧師說聖經給你聽.
所以不用擔心.
你們做什麼.
收割而已.
撒種已經有人做了.
現在一起做收割的工作.
弟兄姊妹.
肯不肯做.
做收割而已.
不用撒種了.
我們一起珍惜.
上帝給我們的時間.
讓更多人能夠聽到福音.
我們一起祈禱.
上帝我們.
求你幫助我們.
藉著今天.
撒瑪利亞婦人.
能夠信耶穌.
撒瑪利亞城裡的人.
能夠信耶穌.
讓我們能夠知道.
主進來福音的門已經大開.
幫助我們.
回應你.

$^{1761}$讓我們不只是.
做一個churchgoer.
只是坐在教會裡.
讓我們成為一個.
Christian.
一個Believer.
一個願意行道的人.
以致我們能夠.
讓這裡.
坐滿這個地方.
坐滿這個禮堂.
讓更多人聽到福音.
完成我們在地上.
你交託給我們的使命.
主啊 莊稼已經熟透了.
讓我們也開我們的眼睛.
看到熟透的莊稼.
我們一起去收割.
我們願意參與主你的角度.
我們多謝主.
聽我們祈禱.
奉耶穌的名字而求.
阿們.
\newpage



\section{彼得前書 4:7-11}
\label{sec:032073}
\textbf{建立健康的使命教會 (14) 作恩賜的管家}~~~~江山偉牧師
\newline
\newline
連結: \href{http://www.church.com.hk/acms/content.asp?site=cdc\&op=show\&type=product\&code=032073\&layout=sermon}{\texttt{032073}} ~~~~ 語音日期: 2026-05-03
\newline
\newline
\hyperref[sec:032072]{\small{< < < PREV SERMON < < <}}
~
\hyperref[sec:index_chronic]{\small{[返順時目]}}
~
\hyperref[sec:032074]{\small{> > > NEXT SERMON > > >}}
\newline
\newline
彼得前書 4:7-11
\newline
\begin{longtable}{cl}
\hline
\hline
章節 & 經文 (和合本修訂版)\\
\hline
4:7 & \begin{tabularx}{0.7\textwidth}{X} 萬物的結局近了。所以你們要謹慎自守，要警醒禱告。 \end{tabularx} \\ \\ \relax
4:8 & \begin{tabularx}{0.7\textwidth}{X} 最要緊的是彼此切實相愛，因為愛能遮掩許多的罪。 \end{tabularx} \\ \\ \relax
4:9 & \begin{tabularx}{0.7\textwidth}{X} 你們要互相款待，不發怨言。 \end{tabularx} \\ \\ \relax
4:10 & \begin{tabularx}{0.7\textwidth}{X} 人人要照自己所得的恩賜彼此服事，作神各種恩賜的好管家。 \end{tabularx} \\ \\ \relax
4:11 & \begin{tabularx}{0.7\textwidth}{X} 若有人講道，他要按著神的聖言講；若有人服事，他要按著神所賜的力量服事，好讓神在凡事上因耶穌基督得榮耀。願榮耀和權能都歸給他，直到永永遠遠。阿們！ \end{tabularx} \\ \\
[1ex]
\hline
\hline
\end{longtable}
$^{1}$各位弟兄姊妹 早安.
每逢受浸的時候.
一定是教會一件很大的事情.
眾弟兄姊妹都會為受浸的媳婦.
為她祝賀.
通常教會都會預備一份受浸的禮物.
有一個弟兄當日也收到很多禮物.
很多人為他祝賀 拍照.
他將這些禮物帶回家之後.
就把教會那份放在一旁.
一個星期又過去 幾個月又過去.
這份禮物一直都沒有動過.
直至有一天 牧者去探望他.
問他這份禮物是什麼.
就是教會送的受浸禮物.
牧者就問他.
你有沒有機會用過.
那個弟兄就很不好意思.
其實我一直都沒有打開過.
不知道你受浸的那份禮物.
是什麼呢?還記不記得呢?.
有沒有用過呢?.
有沒有很用心地去看看呢?.
通常都是聖經 或者是靈修書籍.
其實這個事件都讓我們想起.
我們的屬靈生命其實都是一樣的.
由我們信主那一刻.
其實上帝已經將恩賜.
這段時間我們去到第三個Gift.
恩賜其實早就已經放在你和我的生命裡面.
問題不是神有沒有給我們.
而是我們有沒有發現.
我們是否願意拿出來用呢?.
神給我們恩賜.
不是要我們收藏.
而是要我們去使用.
所以今天我們會透過.
彼得前書四章.
剛才主席說我們讀出的經文.
一起去看.

$^{41}$究竟我們如何去作神百般恩賜的好管家.
彼得寫這封書信的時候.
其實基督徒.
即是當時的弟兄姊妹.
都面對著一個很不容易的處境.
當時的皇帝尼祿.
去加諸很多的迫害在當時的信徒.
所以其實都面對著一個很動盪.
即是信徒受苦.
生命是受著很嚴峻的威脅.
在這個處境裡面.
彼得提醒當時的弟兄姊妹.
我們一起去讀出第一節.
第七節才開始.
預備 1 2 3.
萬物的結局近了.
所以你們要謹慎自守.
叫警醒禱告.
萬物的結局近了.
這裡不單止說這個世界很快終結.
耶穌基督很快再回來.
同時之間也在提醒我們.
今天我們如何能夠用一個永恆的目光.
去看此刻我們的生命和生活.
在當時彼得提醒弟兄姊妹.
要保持一個清醒.
和一個屬靈的生命.
不要被外在環境的影響.
那些逼迫和很多的憂患.
而需要上上的禱告.
和上帝保持一個連結.
敏銳神的心意.
接下來在第十節.
彼得也這樣說.
「人人要照自己所得的恩賜彼此服事」.
「作神百般恩賜的好管家」.
其實這兩節在原文是彼此連結出來的.
而不是兩個分開的命令.
這裡說到警醒禱告.
其實也是恩賜事奉的屬靈的基礎.

$^{81}$當一個人上上去警醒.
上上去禱告的時候.
他的心就能夠和神保持一個連結.
很敏銳上帝要他如何去活化他的人生.
他在一個這樣的生命屬靈的狀態.
就會很容易掌握到神在自己生命裡的帶領.
也很能夠察覺到神放在自己身上的恩賜.
也能夠發現神在教會和身邊的人的需要.
所以彼得在這裡不是先說事奉.
而是說屬靈生命的狀態.
當人和神的關係較正的時候.
其實他們就會很忠心地使用神所託付給他們的恩賜.
彼得在這裡沒有說等環境好一點才去事奉神.
等壓力小一點才去用你的恩賜.
反而彼得在這裡在提醒我們.
在困難的裡面就去使用你的恩賜.
因為神的關係是正常的.
因此因為神的教會很多時候都在困難的當中被建立.
所以很有意思的是在說恩賜的時候.
彼得也在第八和第九節裡面提及到.
我們一起讀出來好不好?.
預備 1 2 3.
「最要緊的是彼此切實相愛.
因為愛能遮掩許多的罪」.
彼得在這裡提醒我們.
在一個充滿動蕩,壓力不確定的時代裡面.
教會不單只是講事工.
而是我們如何去建立一個彼此相愛的群體.
有時候你有沒有想過為何聖經講來講去.
都是三幅,不是屁,是真理.
為何講來講去都是不斷強調呢?.
正正是我們不斷地提醒,再提醒,再提醒我們.
這個講就容易.
實踐的時候其實就很艱難.
在這裡所說的就是.
當經文說「愛能遮掩許多的罪」.
不是說我們假裝看不見罪.
好像問題不存在.
真正的愛其實是包含兩件事情.
一方面,愛是會誠實地指出罪的問題.

$^{121}$另一方面,愛也會溫柔地挽回一個人.
就好像在加拉太書六章一節裡面所說.
「弟兄們,若有人偶然被過犯所成.
你們屬靈的人」.
要怎樣做呢?.
就要用溫柔的心把他挽回過來.
所以彼此相愛不是忽略問題.
而是當我們看到彼此軟弱的時候.
不是要彼此定罪.
而是要彼此寬恕,彼此挽回.
如果我們只是有指責而沒有寬恕.
很容易就會帶來經文裡面所說的.
那個恨意,那個仇恨和爭端.
當愛和饒恕兩樣都存在的時候.
教會就會成為一個能夠醫治,重建人的群體.
所以彼得接下來也說到.
「你們要互相款待,不發怨言」.
換句話說,愛不是情感.
而是一個具體的行動.
當我們這樣實踐,彼此接納,彼此款待.
因此就會很自然地被使用,被建立.
其實在過去幾年前.
我們都經歷過類似的處境.
大家還記不記得疫情中的教會.
幾年前在COVID的時候.
其實很多教會都面對著很多挑戰.
和大家分享一下.
很快我們又不會忘記了.
因為當時我們常常都強調一個new normal.
一個新常態.
這三個字可能你都不陌生.
到這一刻你就已經back to normal了.
已經不再新了.
回想一下在當時最嚴重的時候.
教會又要關閉.
又不能夠實體聚會.
我們盼望多想能夠在當中見到大家.
不過當時有種種的原因.
又限聚令,又要close的時候.
不知道當時你的心情是怎樣.

$^{161}$會不會有孤單呢?.
會不會有焦慮呢?.
甚至乎很害怕.
究竟下一個是不是輪到我呢?.
所以當時也會有人想.
在這樣的狀態.
教會還能不能夠正常運作呢?.
很奇妙.
在那個處境裡面.
其實很多教會,包括我們.
反而在疫情裡面.
見到一些新的曙光.
並不是因為疫情的籠罩.
我們就好像做不到事.
我們就好像失去了教會的功能.
邱太英姐妹就開始去想.
怎樣透過科技去繼續去舉行一些崇拜.
也有人開始學用Zoom.
在當中繼續有小組.
也有弟兄開始學不同的AV器材.
以致希望令到崇拜能夠更加流暢.
當時講求怎樣收聲.
不是叫大家收聲.
是收音.
收音能夠更加清晰.
每件樂器都希望能夠出聲.
在困難的裡面.
教會沒有因為疫情而停止運作.
反而是更多的愛去活出來.
還記得當時有不少弟兄姊妹.
都嘗試去打電話.
彼此關心問候.
我也知道有些弟兄姊妹負責煮菜.
送到門口.
讓中招的弟兄姊妹.
能夠有些溫飽.
有人支持著.
當時的教目.
怎樣去面對主餐呢?.
將二合一的主餐送到.

$^{201}$當時分了不同的小隊.
送到每一個會友的門口.
我們在線上一起領主餐.
所以種種回憶回想這些片段.
豈不是引證了彼得的提醒.
當教會彼此相愛的時候.
因此就能夠被使用.
當因此被使用的時候.
教會就被建立.
剛才的小組.
牧師和師母當時的大小組.
在短短的期間.
一分為三.
在後來的時候.
整個教會因全家.
在短短這兩三年的時候.
神的祝福帶領了在座來到我們當中.
所以問題就來了.
如果神已經將恩賜放在我們的生命裡面.
究竟今天你有沒有發現祂呢?.
接下來我們要講四個D.
第一個D 我們怎樣去發現.
法國神已經給你的恩賜呢?.
彼得在這裡說 人人要照自己所得的恩賜彼此服侍.
留意這句說話.
彼得沒有說如果你有.
他是說照所得的恩賜.
意思是什麼?.
其實你和我已經領受了.
當我們說萬物的結局近了.
即是說時間有限的時候.
我們就不能夠將神的恩賜放在一邊.
不是我沒有恩賜.
而是神給我的恩賜.
我又沒有好好地去使用.
因為我們知道有一天.
我們都要在神的面前交仗.
我們怎樣去管理祂交託給我們的生命和恩賜呢?.
弟兄姊妹 很多時候我們都會有一個錯覺.
以為時間很多.

$^{241}$我不知道你有沒有講過以下的說話.
讓我忙完這件工作.
讓我的小朋友長大一點.
讓我的生活穩定一點.
讓我先準備好.
其實到死那天都準備不好.
讓我準備好一點才能夠更好.
很多時候等一等就變成再等等.
再等等就等了很多年.
時間就是這樣不知不覺地過去.
我們以為還有很多個明日.
不過聖經提醒我們.
人生其實很短暫.
《牙國書》四章十四節.
我們一起讀出來.
預備 1 2 3.
其實明天如何?.
你們還不知道.
你們的生命是什麼呢?.
你們原來是一片雲霧.
出現片刻就不見了.
特別在困難的時候.
這次經文我們很有共鳴.
當現在好像回復正常.
很安逸的時候.
我們又好像忘記了.
所以彼得在這裡不斷提醒我們.
萬物的結局近.
不是要恐嚇我們.
而是提醒我們一件事情.
神給我們的時間是寶貴.
而且是有限.
既然是這麼有限.
我們就更加不應該將神給我們的寶貝.
埋沒在一邊.
神給我們的恩賜不是到某一天才用.
而是神很想很想很想.
我們在此時此刻.
今日在現在.
在教會和在神擺放在我們的處境裡面.

$^{281}$能夠活現出來.
在實際裡面.
Pastor Jesse.
究竟如何能夠找到神給我們的恩賜呢?.
我們可以從四個向度找到一些足跡.
第一就是熱情.
有些福士當你拿在手上的時候.
你就會感覺到很累.
你就會很沮喪.
但有些福士當你越做的時候.
你越起勁.
起勁到一個點你睡不著.
就好像我今日要講國語堂和粵語堂.
昨晚我起了三次.
以為七點了.
誰知道三點,五點到最後終於七點了.
所以有些人.
當上帝給你一個熱情的時候.
你睡覺也不會忘記他.
經常會想著.
當上帝給你見到有需要的人.
你就會很自然.
不用提醒你.
你就會很自然去關心他,去探訪.
也有弟兄姊妹.
他們很樂於享受在幕後做很多的安排,組織.
總之在人看不見的地方做了很多事情.
令事情很順利地完成.
我相信這些都不是偶然.
很多時候.
上帝就是給我們每一個弟兄姊妹.
放了一種熱情在我們的心裡面.
引導我們去預備為上帝服侍.
第二方面就是果子.
在這件事情裡面.
有時候我們都知道.
當你用盡不知多少百分比的努力.
得出來的效果都是很有限的.
有些弟兄姊妹同樣做一件事.
他很快地三刨兩撥.

$^{321}$就已經做得很好了.
所以在我們在當中也會有時候說.
不知為何跟你聊完天之後.
心裡面很得安慰,舒暢.
有些弟兄姊妹可能他分享神的話語.
其實就很容易幫助到人的明白和提醒.
也有我們當中的領袖.
當帶領團隊的時候.
特別是吳牧師.
他很有方向感.
帶領著我們的時候.
我們很知道前面下一步應該怎麼走.
這些果子.
其實很多時候就是神在提醒我們.
神就是透過你這個領域來使用你.
第三方面就是別人對你的肯定.
很多時候我們都不知道自己有什麼恩賜.
當身邊的人看到的時候.
可能他們會說你真的很棒.
你解釋得很好.
你今天真的唱得很好.
你的行政方面很有能力.
這些可能有很多時候我們聽過就算了.
但其實可能是神都是藉著身邊的弟兄姊妹.
提醒我們成為一個引證.
最後一方面就是需求.
有時候其實神不是要找一些有能力的人.
他是先給我們有一個負擔的心.
每個人的負擔都不一樣.
特別是我們未來新堂的時候.
有很多不同的需要.
不是要你每樣東西都要參與.
你就是要找一樣東西.
上帝就放在你的心裡面.
這一方面可能是你特別關心.
可能你看到這個社區有很多人困苦流離.
或者看到年輕人很迷失.
你很想去陪伴他們.
又或者當你看到教會某一個崗位.
很需要童工的時候.

$^{361}$可能有時候你都會心想.
神究竟我可不可以幫到忙呢?.
這些其實都是神呼召的一個負擔的開始.
所以這四個元素.
當我們一起去禱告,去諮商.
去尋求神的心意的時候.
其實就幫助我們有一個更加清晰的方向.
神可能是帶領著我們.
去透過我們的恩賜來服侍.
當我們發掘完之後.
就千萬不要停在那裡.
第二個D就是我們怎樣去培育.
神所託付給我們的恩賜.
發現恩賜只是開始.
培育恩賜才是我們的責任.
恩賜其實就好像一粒種子.
剛剛三月的時候我們拿那些種子回家.
如果你不放在泥土.
不加水,不曬陽光.
其實它仍然是一粒.
可能已經死了也說不定.
我們需要去給它有養分.
給它一個適當的平台.
讓它不斷地成長.
所以很多時候我們都會有兩個誤解.
第一個誤解就是.
我們有時候都會帶著一個競爭的心態.
特別在現在這個社會裡面.
從小我們就被教導.
我們要累積很多的履歷.
建立成就.
讀書要比其他人好.
工作要比其他人成功.
能力要比其他人突出.
慢慢地我們都會很容易.
將這種比較,這種思維.
帶到教會裡面.
有時候我們都會問.
究竟誰的恩賜比較多.
誰的領袖有魅力.

$^{401}$誰是被神使用比較多.
甚至有些人可能會想.
如果我的恩賜沒有別人那麼明顯.
好像我就沒有那麼重要.
不過聖經在這裡再給我們一個很不同的角度.
人人要照自己所得的恩賜.
彼此服侍.
這裡有兩個很重要的提醒.
人人無一倖免.
是在說你和我.
你看看你旁邊那個.
是你.
是你有份.
第二個提醒就是彼此服侍.
恩賜不是要用來比較.
而是要用來服侍.
換句話說.
恩賜不是要為了證明.
自己有多厲害.
而是為了要祝福別人的生命.
就好像一個身體一樣.
每一個部份有不同的功能.
當我們看聖經的時候.
很多時候用身體來作為基督的肢體.
缺少不可.
眼不能說你的手指尾沒有用.
腳趾尾沒有用.
我有用一點.
你試試摸手指尾和腳趾尾會怎樣.
是很笨拙的.
同樣在教會裡面.
固然有人很擅長去牧羊.
有人會很擅長去教導.
有人很擅長去關懷.
有人會很擅長去做接待.
神看重的.
不是誰被誰突出.
而是看我們有沒有去忠心地.
去使用祂所託付給我們的恩賜.
在愛裡面被建立.

$^{441}$第二個誤解就是.
我們很多時候都會過分謙卑.
可能我們會說我不行的.
我不行的.
你去找別人吧.
我還不夠班.
還沒準備好.
表面上好像很謙卑.
我們先祈禱.
看看神有沒有什麼感動.
但其實很多時候.
我們都是為了自己來盤算.
彼得這裡沒有說.
等你準備好再服侍.
整本聖經裡面.
其實都是呼召我們每一個.
去成為神的用人.
其實恩賜從來都不是要等到成熟才使用.
而是我們在過程裡面.
要去被培育.
就好像學樂器一樣.
我唯一一件會彈的樂器.
就是結他.
當時我彈結他.
不是因為興趣.
而是因為在小組裡.
經常都要清唱.
很乾.
所以因為這樣的緣故.
我才去學結他.
當時我的年代沒現在那麼方便.
播YouTube就搞定.
經常清唱.
音又不準.
所以我是因為看到那個需要.
而去學的.
我還記得當時彈的時候.
很辛苦.
我也很欣賞我父母.
一方面他們給了我2000元買了一支結他.

$^{481}$另一方面天天聽我.
很吵耳的.
一年過去.
還是那麼吵耳.
到第二年開始.
彈到一首歌了.
可以開始在小組了.
當繼續有一個客人來服侍的時候.
神開始給我機會.
可以在台上.
透過敬拜隊來服侍.
所以我相信.
在整個過程裡面.
恩賜其實也是這樣.
很多時候神會.
為我們預備一個.
黑房.
Matthew 知道我在說什麼.
拍照.
很久以前.
在黑房裡面.
透過菲林.
菲林沖出來要曬.
用人手擦一擦.
在黑房裡面.
固然很暗.
也沒觀眾.
在一個看不見的地方.
那幅照片就會慢慢顯影.
神塑造人的方式.
往往都像這樣.
神很多時候.
不是馬上把我們拉到台上.
祂是先帶我們.
去一個隱藏的地方.
在那裡.
祂就塑造我們的生命.
磨練我們的信心.
建立我們屬靈的品格.
聖經裡面.

$^{521}$有不少的人物.
都經歷過這樣的過程.
大衛就是一個很好的例子.
當撒母耳.
來高沒他為王的時候.
大家還記得當時.
他只不過是一個.
牧羊的少年人.
連他爸爸.
我都看不見.
他.
都是找一個帥一點的.
一個大隻一點的.
最後才輪到他.
當時撒母耳.
高沒他的時候.
其實大衛還沒有馬上.
來到作王.
中間其實等了大約十五年.
在這段期間.
大衛在做什麼呢.
其實就是繼續在曠野裡面.
牧羊.
只有羊群.
在曠野裡面.
當然還有他和上帝之間.
的交往.
在那個地方.
他面對獅子.
學習什麼呢.
勇敢.
在那個地方他學彈琴.
去敬拜神.
學習如何去親近神.
在那個階段.
他經歷等待.
學習如何去信靠神.
當後來他面對.
哥利亞的時候.
他能夠說.

$^{561}$耶和華救我脫離獅子和紅的爪.
也必救我脫離.
這非利士人的手.
因為.
大衛已經在.
那十五年的黑房.
的經歷裡面.
學會如何去依靠神.
信靠神.
亦磨練他對上帝的依賴.
領事們很多時候.
神不是要我們.
很快地去獨當一面.
神反而在不同的人生.
的階段裡面.
神說此刻的我.
久不久我也要去.
黑房裡面.
繼續被上帝練成.
透過生活.
或是在事奉中的點滴.
一些挫敗.
再次磨練.
我對上帝的.
昔日的初心.
昔日上帝給我的託付.
我們如何能夠.
重新去.
re-calibrate.
去校正.
使我們能夠繼續.
為上帝走前面.
一個更遠更合上帝心意.
甚至乎是超乎.
我們所想到的.
一個里程.
如果沒有了.
黑房的操練.
可能我們的能力.
會超過我們的品格.

$^{601}$我們的位置.
可能都會超過我們的生命.
當我們的生命.
盛載不到這個呼召的時候.
呼召可能反而.
會帶來重擔.
所以神在.
隱藏的地方塑造我們.
是要重新再建立我們的.
信心,內心和馴服.
只有.
經過黑房的人.
他的生命才能夠.
盛載到神給他的呼召.
但培育.
因此並不是最終的目的.
因此最終的目的.
是要被使用.
要去祝福更多人建立教會.
完成神的工作.
這就是第三個D.
彼得.
在前書四章.
首一節我們一起再讀這節經文.
預備一二三.
若有人講道.
他要按著神的聖言講.
若有人服事.
他要按著神.
所賜的力量服事.
彼得在這裡很簡單.
神的恩賜分成兩大類.
包括.
話語的.
恩賜.
和服事的恩賜.
話語的恩賜包括教導,勸勉,講道.
分享神的話,帶領查經.
鼓勵別人,關心人等等.
而服事的恩賜.

$^{641}$可以是一些接待,關懷,行政.
安排一些事情.
幫助其他人.
默默地在背後支持.
教會的團隊運作.
所以在教會裡.
有些人是用口服事.
的確是.
有些人.
是用雙手和生命.
去事奉.
但彼得在這裡提醒我們.
一件非常重要的事.
真正服事背後的力量.
不是靠自己.
不是靠自己預備多少.
而是禮至神.
無論你話語的服事怎樣也好.
我們是需要在聖經裡講.
要按著神的聖言講.
或者我們在服事的過程裡.
我們要按著.
神所賜的力量服事.
當我們服事了一段時間.
可能你開始覺得.
得心應手.
我覺得我們.
經常要培養一種.
好像這節經文給我們的提醒.
我們需要去依靠神.
當我們服事的時候.
我覺得我們也是很好的.
我們會走在一起去祈禱.
每個星期當我們崇拜之前.
吾穆斯也會帶領我們一起.
圍著一群服事人員去禱告.
其實在這個過程裡.
再一次調教我們.
不是靠我們練了多少.
不是靠我們準備了多少.

$^{681}$而是再一次提醒我們.
我們怎樣去依靠神.
當上帝給了你一個.
好像剛才所講.
給了你一個負擔.
給了你一個體見的時候.
沒錯我們需要在過程裡.
去培育.
同時間當我們在那個場景裡.
提醒.
記得提醒提醒自己.
不是靠你的經驗.
不是靠你的才幹.
不是靠你的努力.
唯有靠著神的靈.
方能成事.
神從來都不是找.
最有能力的那個.
神在找.
那些是.
有一個謙卑受教.
願意的心.
當一個人他願意.
將自己交在神的手裡.
神就能夠去使用他.
我認識.
有一位姊妹她不喜歡在台上.
去作任何分享.
因為她說自己不擅長.
說話.
所以每次有人叫她.
分享見證她都很緊張.
可以的話就不要了.
不過她有一個很特別的恩賜.
就是很懂得關心人.
每個星期她都會打給.
不同的教會裡的弟兄姊妹.
尤其是長者.
問候一下他們的近況.
也都知道.

$^{721}$當誰誰不舒服的時候.
她又會虛寒問暖又會煮菜.
給他們.
她也都很記得人的生日.
都記下了.
到那個人生日的時候.
她就會發一個.
短的訊息.
或者預備一些小禮物.
她所做的可能都很微小.
很普通.
有一天有另一位.
接受過她服侍的姊妹.
就跟我說了一番話.
牧師牧師.
全賴她一直關心我.
不然的話我都不會有.
信耶穌的一天.
有些神的恩賜.
其實有很多不同的方式.
有些恩賜是被人看到的.
有些恩賜.
是在幕後.
默默的來發揮.
有些人是用聲音.
來服侍,有些人是用雙手.
有些人是透過一個關懷的心.
在神的眼中.
每一個恩賜.
都是同樣重要.
而當我們來到.
彼此配搭的時候.
我們事奉最終的目的.
不是要被人看到.
而是要讓神.
得著榮耀.
第四個D.
讓神.
在四章十一節裡面.
我們一起讀出來.

$^{761}$好讓神在凡事上.
因耶穌基督.
得榮耀.
這個其實就是.
一切恩賜.
最終的目的.
恩賜不是要建立.
自己的名聲.
也不是要展現自己的能力.
也不是要讓其他人看到我多厲害.
恩賜存在的目的.
是要讓.
神得著榮耀.
在神的國度裡面.
有些服侍.
是沒有什麼人會留意.
但是.
只要我們知道.
我們的心.
我們有沒有來調教.
而去做.
上帝是看見的.
所以當我們忠心使用神.
給我們的負擔.
我們好好發揮的時候.
最終的結果.
就是要讓神.
得著當得的榮耀.
可能在我們生命裡面.
正有一個禮物盒.
不知道你打開了沒有.
可能還有些人從未打開.
裡面其實是.
藏了神專屬.
為你預備的恩賜.
今天神再一次.
邀請你去打開.
使用它.
讓神的家.
讓神的國度.

$^{801}$讓這個世界可以.
讓神得著榮耀.
不要讓你的人生.
將它白白地.
繼續收藏埋沒.
請你打開它.
讓神去使用你.
成為神恩典流動的.
管道.
在這一刻.
我想大家去禱告.
我們一起閉上眼睛.
很想你們去想三個問題.
主啊.
你給我的恩賜是什麼呢?.
主啊.
你給我的恩賜是什麼呢?.
主啊.
在哪一個地方.
我需要繼續去成長呢?.
主啊.
在哪一個地方我需要繼續去成長呢?.
最後主啊.
你呼喚我.
要在哪一個地方.
去服侍你呢?.
最後主啊.
你呼喚我.
要在哪一個地方.
去服侍你呢?.
可能在座當中.
仍然有些姐妹不知道.
自己的恩賜是什麼.
或許你都覺得自己.
未夠成熟.
心裡仍然有些猶豫.
或許害怕.
神今天不是要.
去找那些.
很有能力的人.

$^{841}$神正在尋找.
願意為他.
擺上的人.
特別可能此刻你仍然.
沒有任何事奉的崗位.
神今天再一次輕輕地提醒我們.
我今天在你的生命裡.
就放了.
恩賜.
就是要你去祝福教會.
祝福其他人.
很想在這裡發出一個邀請.
如果這也是.
你在神面前的願意.
主啊,我願意讓你帶領我.
去發掘去培育.
去使用.
你放在我生命裡的恩賜的話.
我想邀請你.
在我的座位裡起立.
你去回應神.
你繼續向神.
立志.
擺上.
為自己去禱告.
如果這也是你的心.
很想讓你去.
餘生.
獻情給神去使用.
縱使這一刻.
你未清楚,未掌握.
但若然你有一個願意的心.
我邀請你.
你在座位裡起立.
主啊,我在這裡.
請差遣我.
今天在神的國度裡.
仍然很需要工人.
在神家裡.
我們仍然很需要頂子妹.

$^{881}$群策群力.
不只是落在.
那群.
頂子妹的身上.
神是呼喚著我們每一個.
求主帶領我們.
求主使用我們.
還有沒有頂子妹呢?.
讓我們一起.
去禱告.
請我們每一個都起立.
我們一起去.
接受上帝.
給我們的差遣.
接受上帝.
給我們的帶領.
藉著你的話語再一次提醒我們.
今天.
我們的生命.
是你救贖過來的.
我們的生命是屬於你的.
你將每一個恩賜.
放在我們每一個頂子妹的心裡.
在我們當中.
我們很感恩.
我們能夠在神的國度有份.
我們能夠成為你的兒女.
神也帶領我們來到恩泉家.
我們很感受到.
體會到我們在這個屬靈的家裡.
真的每一個頂子妹.
都是很用心竭力的.
在不同的崗位裡面.
去與神同工.
與神同行.
主也求你繼續向我們說話.
在前行的裡面.
當教會一路的發展.
當世代主耶穌很快再回來的時候.
主也求你.

$^{921}$叫我們不要再耽延時間.
讓我們每一個的生命.
就是按著.
今天這段經文.
你的話語.
你的提醒.
我們都能夠好好地.
善用我們的餘生.
善用我們的時間.
怎樣能夠在神給過我們的恩賜.
在不同的崗位角色裡面.
繼續去到中心的擺上.
去與神同工.
與神同行.
主也就願剛才我們在.
禱告的裡面.
你繼續向我們說話.
你繼續扭轉我們的生命.
我深信.
神的精神要使用我們的時候.
不是只是停留在.
我們所能夠思想得到的.
而是上帝你要.
讓我們去超越.
一個可能我們從未想過.
就是人心裡面.
未曾想過.
眼睛未曾看過.
耳朵未曾見過的一個境地.
主啊 就願意你去.
拆穿我們.
在這個時代的裡面.
我們要成為精兵.
主你的工人.
成為你合用的僕人侍女.
繼續在神的家.
神的國度.
在我們的職場.
繼續去榮耀主你的名.
多謝你.

$^{961}$求你就是這樣去.
聽我們的禱告 獻上我們的謝恩.
奉耶穌基督得勝的名求.
阿們.
\newpage



\section{帖撒羅尼迦後書 1:1-3:18}
\label{sec:032074}
\textbf{帖撒羅尼迦書 (2) 主快再來}~~~~劉少康牧師
\newline
\newline
連結: \href{http://www.church.com.hk/acms/content.asp?site=cdc\&op=show\&type=product\&code=032074\&layout=sermon}{\texttt{032074}} ~~~~ 語音日期: 2026-04-28
\newline
\newline
\hyperref[sec:032073]{\small{< < < PREV SERMON < < <}}
~
\hyperref[sec:index_chronic]{\small{[返順時目]}}
~
\hyperref[sec:032075]{\small{> > > NEXT SERMON > > >}}
\newline
\newline
帖撒羅尼迦後書 1:1-3:18
\newline
\begin{longtable}{cl}
\hline
\hline
章節 & 經文 (和合本修訂版)\\
\hline
1:1 & \begin{tabularx}{0.7\textwidth}{X} 保羅、西拉和提摩太寫信給帖撒羅尼迦、在我們的父神與主耶穌基督裡的教會。 \end{tabularx} \\ \\ \relax
1:2 & \begin{tabularx}{0.7\textwidth}{X} 願恩惠、平安從我們的父神和主耶穌基督歸給你們！ \end{tabularx} \\ \\ \relax
1:3 & \begin{tabularx}{0.7\textwidth}{X} 弟兄們，我們該常常為你們感謝神，這本是合宜的；因為你們的信心格外增長，你們眾人彼此相愛的心也都增加。 \end{tabularx} \\ \\ \relax
1:4 & \begin{tabularx}{0.7\textwidth}{X} 所以，我們在神的各教會裡為你們誇耀，因為你們在所受的一切壓迫患難中仍牢守著耐心和信心。 \end{tabularx} \\ \\ \relax
1:5 & \begin{tabularx}{0.7\textwidth}{X} 這正是神公義判斷的明證，使你們配得上他的國，你們就是為這國受苦。 \end{tabularx} \\ \\ \relax
1:6 & \begin{tabularx}{0.7\textwidth}{X} 既然神是公義的，他必以患難報復那加患難給你們的人， \end{tabularx} \\ \\ \relax
1:7 & \begin{tabularx}{0.7\textwidth}{X} 也必使你們這受患難的人與我們同得平安。那時，主耶穌同他有權能的天使從天上在火焰中顯現， \end{tabularx} \\ \\ \relax
1:8 & \begin{tabularx}{0.7\textwidth}{X} 要報應那些不認識神和不聽從我們的主耶穌福音的人。 \end{tabularx} \\ \\ \relax
1:9 & \begin{tabularx}{0.7\textwidth}{X} 他們要受懲罰，永遠沉淪，與主的面和他權能的榮光隔絕。 \end{tabularx} \\ \\ \relax
1:10 & \begin{tabularx}{0.7\textwidth}{X} 這正是主再來，要在他聖徒的身上得榮耀，就是要使一切信的人感到驚訝的那日子，因為你們信了我們對你們作的見證。 \end{tabularx} \\ \\ \relax
1:11 & \begin{tabularx}{0.7\textwidth}{X} 為此，我們常為你們禱告，願我們的神看你們與他的呼召相配，又用大能成就你們一切良善的美意和因信心所做的工作， \end{tabularx} \\ \\ \relax
1:12 & \begin{tabularx}{0.7\textwidth}{X} 使我們主耶穌的名，照著我們的神和主耶穌基督的恩，在你們身上得榮耀，你們也在他身上得榮耀。 \end{tabularx} \\ \\ \relax
2:1 & \begin{tabularx}{0.7\textwidth}{X} 弟兄們，關於我們主耶穌基督的來臨和我們到他那裡聚集，我勸你們： \end{tabularx} \\ \\ \relax
2:2 & \begin{tabularx}{0.7\textwidth}{X} 無論藉著靈，藉著言語，藉著冒我的名寫的書信，說主的日子已經到了，不要輕易動心，也不要驚慌。 \end{tabularx} \\ \\ \relax
2:3 & \begin{tabularx}{0.7\textwidth}{X} 不要讓任何人用甚麼法子欺騙你們，因為那日子以前必有叛教的事，並有那不法的人，那沉淪之子出現。 \end{tabularx} \\ \\ \relax
2:4 & \begin{tabularx}{0.7\textwidth}{X} 那抵擋者高抬自己超過一切稱為神明的，和一切受人敬拜的，甚至坐在神的殿裡，自稱為神。 \end{tabularx} \\ \\ \relax
2:5 & \begin{tabularx}{0.7\textwidth}{X} 我還在你們那裡的時候曾把這些事告訴你們，你們不記得嗎？ \end{tabularx} \\ \\ \relax
2:6 & \begin{tabularx}{0.7\textwidth}{X} 現在你們也知道那攔阻他的是甚麼，為要使他到了時機才出現。 \end{tabularx} \\ \\ \relax
2:7 & \begin{tabularx}{0.7\textwidth}{X} 因為那不法的隱祕已經運作，只是現在有一個阻擋的，要等到那阻擋的被除去才會發作， \end{tabularx} \\ \\ \relax
2:8 & \begin{tabularx}{0.7\textwidth}{X} 那時這不法的人必出現，主耶穌要用口中的氣滅絕他，以自己來臨的光輝摧毀他。 \end{tabularx} \\ \\ \relax
2:9 & \begin{tabularx}{0.7\textwidth}{X} 這不法的人來，是靠撒但的運作，行各樣的異能、神蹟和一切虛假的奇事， \end{tabularx} \\ \\ \relax
2:10 & \begin{tabularx}{0.7\textwidth}{X} 並且在那沉淪的人身上行各樣不義的詭詐，因為他們不領受愛真理的心，好讓他們得救。 \end{tabularx} \\ \\ \relax
2:11 & \begin{tabularx}{0.7\textwidth}{X} 故此，神就給他們一個引發錯誤的心，叫他們信從虛謊， \end{tabularx} \\ \\ \relax
2:12 & \begin{tabularx}{0.7\textwidth}{X} 使一切不信真理、倒喜愛不義的人都被定罪。 \end{tabularx} \\ \\ \relax
2:13 & \begin{tabularx}{0.7\textwidth}{X} 主所愛的弟兄們哪，我們本該常為你們感謝神，因為他揀選你們為初熟的果子，使你們因信真道，又蒙聖靈感化成聖，得到拯救。 \end{tabularx} \\ \\ \relax
2:14 & \begin{tabularx}{0.7\textwidth}{X} 為此，神藉著我們所傳的福音呼召你們，好得著我們主耶穌基督的榮光。 \end{tabularx} \\ \\ \relax
2:15 & \begin{tabularx}{0.7\textwidth}{X} 所以，弟兄們，你們要站立得穩，凡所領受的教導，無論是我們口傳的，是信上寫的，都要堅守。 \end{tabularx} \\ \\ \relax
2:16 & \begin{tabularx}{0.7\textwidth}{X} 願我們主耶穌基督自己，和那愛我們、開恩將永遠的安慰及美好的盼望賜給我們的父神， \end{tabularx} \\ \\ \relax
2:17 & \begin{tabularx}{0.7\textwidth}{X} 安慰你們的心，並且在一切善行善言上堅固你們！ \end{tabularx} \\ \\ \relax
3:1 & \begin{tabularx}{0.7\textwidth}{X} 末了，弟兄們，請你們為我們禱告，好讓主的道快快傳開，得著榮耀，正如在你們中間一樣， \end{tabularx} \\ \\ \relax
3:2 & \begin{tabularx}{0.7\textwidth}{X} 也讓我們能脫離無理和邪惡人的手，因為不是人人都有信仰。 \end{tabularx} \\ \\ \relax
3:3 & \begin{tabularx}{0.7\textwidth}{X} 但主是信實的，他要堅固你們，保護你們脫離那邪惡者。 \end{tabularx} \\ \\ \relax
3:4 & \begin{tabularx}{0.7\textwidth}{X} 我們靠主對你們有信心，你們現在遵行，以後也必遵行我們所吩咐的。 \end{tabularx} \\ \\ \relax
3:5 & \begin{tabularx}{0.7\textwidth}{X} 願主引導你們的心去愛神，並學基督的忍耐！ \end{tabularx} \\ \\ \relax
3:6 & \begin{tabularx}{0.7\textwidth}{X} 弟兄們，我們奉主耶穌基督的名吩咐你們，凡有弟兄懶散，不遵守我們所傳授的教導，要遠離他。 \end{tabularx} \\ \\ \relax
3:7 & \begin{tabularx}{0.7\textwidth}{X} 你們自己知道該怎樣效法我們。因為我們在你們當中從未懶散過， \end{tabularx} \\ \\ \relax
3:8 & \begin{tabularx}{0.7\textwidth}{X} 也從未白吃人的飯，倒是辛苦勞碌，晝夜做工，免得使你們中間有人受累。 \end{tabularx} \\ \\ \relax
3:9 & \begin{tabularx}{0.7\textwidth}{X} 這並不是因我們沒有權柄，而是要給你們作榜樣，好讓你們效法我們。 \end{tabularx} \\ \\ \relax
3:10 & \begin{tabularx}{0.7\textwidth}{X} 我們在你們那裡的時候曾吩咐你們，說若有人不肯做工，就不可吃飯。 \end{tabularx} \\ \\ \relax
3:11 & \begin{tabularx}{0.7\textwidth}{X} 因為我們聽說，在你們中間有人懶散，甚麼工都不做，反倒專管閒事。 \end{tabularx} \\ \\ \relax
3:12 & \begin{tabularx}{0.7\textwidth}{X} 我們靠主耶穌基督吩咐並勸戒這樣的人，要安分做工，自食其力。 \end{tabularx} \\ \\ \relax
3:13 & \begin{tabularx}{0.7\textwidth}{X} 弟兄們，你們行善不可喪志。 \end{tabularx} \\ \\ \relax
3:14 & \begin{tabularx}{0.7\textwidth}{X} 若有人不聽從我們這信上的話，要把他記下，不和他交往，使他自覺羞愧； \end{tabularx} \\ \\ \relax
3:15 & \begin{tabularx}{0.7\textwidth}{X} 但不要把他當仇人，要勸他如勸弟兄。 \end{tabularx} \\ \\ \relax
3:16 & \begin{tabularx}{0.7\textwidth}{X} 願賜平安的主隨時隨事親自賜給你們平安！願主與你們眾人同在！ \end{tabularx} \\ \\ \relax
3:17 & \begin{tabularx}{0.7\textwidth}{X} 我—保羅親筆向你們問安。凡我的信都以此為記，我的筆跡就是這樣。 \end{tabularx} \\ \\ \relax
3:18 & \begin{tabularx}{0.7\textwidth}{X} 願我們主耶穌基督的恩惠與你們眾人同在！ \end{tabularx} \\ \\
[1ex]
\hline
\hline
\end{longtable}
$^{1}$各位弟兄姊妹平安.
我們今天要講《貼撒羅尼加後書》.
這本書是保羅收到提摩他人給他的訊息之後.
他就有些釋懷.
因為原來這間教會很掛念他.
而且弟兄姊妹仍然信心堅固.
但是很可惜.
他們當中也有些人傳了一些不應該傳的.
對保羅的批評.
或者一些不正確的教導.
所以保羅很快就寫了《貼撒羅尼加後書》給這間教會.
這封信一開頭的時候.
是用保羅,西拉,提摩他三個人的名字發出去的.
三個人一起向這間教會的弟兄姊妹問安.
三個人一起說.
這間教會是主耶穌基督.
我們的父神和主耶穌基督的教會.
各位你要知道今天的教會.
是我們的神和我們的主所設立的.
而聖靈就在我們當中工作.
所以三位一體的神都與我們同在.
所以教會是神的家.
教會是神的殿.
教會是神的所在.
而我們這群弟兄姊妹就是教會.
所以教會不是只是一間禮拜堂.
一個地點.
因為地點可以改.
有些人是在家裡聚會.
有些人是在其他地方聚會.
或者租一些酒店.
租一些地方去聚會.
所以聚會的地方不要緊.
那個不是教會.
我們這群弟兄姊妹才是教會.
帖撒羅尼加後書說什麼呢?.
就是那些弟兄姊妹覺得主耶穌回來的日子很近了.
就有些錯誤的行動.
有些人說主耶穌就回來了.
那不用工作了.

$^{41}$不要工作了.
那又怎樣?.
那就靠其他弟兄姊妹幫他了.
聖保羅聽了之後心裡有些很不安心.
這樣也行?.
因為主耶穌回來了你就不工作?.
其實你是很不負責任.
於是他就責備他們.
那些不做工的不准吃飯.
哇!真的很嚴厲.
你怎麼可以這樣叫其他人供應你.
你自己又不去工作又不去賺錢.
那就靠其他人供養.
於是保羅就提出很多的警誡.
也是教他們怎樣要小心.
不要誤信一些這樣的假的道理.
帖撒羅尼加後書只有三章聖經.
但這三章聖經充滿了很多的真理.
第一章帖撒羅尼加後書.
是關於文安和感恩.
我發覺保羅是一個很懂得感恩的使徒.
每一封書述他都找到很多東西.
甚至是一些很細微的東西.
所以在帖撒羅尼加後書.
保羅就提到.
我們在患難的時候.
祂和我們一起得平安.
所以主耶穌基督祂的能力.
祂的大能在我們當中向我們顯現.
所以帖撒羅尼加後書有很寶貴的信息.
還有提到主在我們當中.
祂是得著榮耀的名.
主要得著榮耀.
教會有一個使命.
教會有一個目標.
就是榮耀我們的父神.
然後去到第二章.
提到我們很需要明白主在內的真理.
特別今天打仗的風聲.
在我們四周圍響起.

$^{81}$中東那種緊張牽涉到很多國家.
所以我們知道主回來的日子真是很近.
求神叫我們睜大眼睛看清楚.
所以聽新聞.
分析這些新聞同樣是很重要的.
我不是叫大家要聽一些閒言閒語的評論.
而是要聽一些很中肯的評論員.
他們很深刻的分析.
主降臨是一定有各樣不同的情況出現.
首先主降臨的時候.
不法者.
就是仇敵撒旦控制的不法者.
會顯露出來.
他的尾巴就被我們看見了.
他要做很多不好的勾當.
因此他就說門徒應該要守保羅的教訓.
你要遵守我給你們的教導.
要切切的遵行.
要好好的聆聽.
要堅守我教你們的真道.
然後去到第三章.
保羅有一個禱告.
主的道要快快走開.
就是主的福音要越傳越快.
越傳越遠.
越傳越多.
越來越多人信主.
今天是一個末世的時候.
特別是一個媒體的時候.
我的YouTube可以傳到很遠.
很多人的信息可以傳到很遠.
而且是很快.
當每一次我請頂尖幫我放上網的時候.
馬上那天已經有幾百人看了.
接著千人了.
接著千多人看了.
越來越多.
這個資訊爆炸的年代真的很厲害.
求主幫助我們要努力趕快做傳福音的工作.
讓主的道快快傳開.

$^{121}$然後就是剛才我提到的.
保羅警戒那些不做工的人.
你不做工就不吃飯.
不可吃飯是一個吩咐不准吃飯.
而且我們的行善是不可以喪志.
大家都知道做善事很重要.
特別在這個末世的時代.
有些弟兄姊妹或者有些有需要的人.
真的很可憐.
沒有人去幫他們.
如果你知道有弟兄姊妹有需要.
讓你解開你的手伸出你的援手.
去幫助去支援那些有需要的人.
然後到了十六節.
保羅就說.
次平安的主.
我們的主是一位次下平安的主.
祂會隨時任何時間.
祂會隨時每一件事都賜你平安.
我很喜歡這句金句.
次平安的主.
我們的主是平安的源頭.
祂隨時.
祂隨時每一件事都賜我們平安.
而且是親自給我們平安.
帶給我們境大安慰.
主親自來給我們平安.
到了十七節.
保羅就說.
我親筆問你們安.
這封信是保羅手寫的.
你看我的筆跡就是這樣.
我寫的筆跡就是這樣.
認住認住.
這些就是我的筆跡.
有一次教會更新運動.
因為他們四十週年紀念.
請我寫一些感恩的說話回去.
我就說.
我用手寫.

$^{161}$寫了一些我.
因為和這個運動一同成長.
我參與這個運動.
差不多有三十多年的時間.
所以就寫給他們.
那位總幹事收到之後很開心.
他說劉先生你親筆寫的.
很寶貴.
這個稿件很有價值.
因為我說回.
當初香港教會更新運動是怎樣成立的.
因為當時面對九七.
我們有一班的教務同工.
差不多一月一至.
去到粉嶺的宣道園那裡.
大家一起開會.
大家一起商量.
大家一起同心祈禱.
有一班的教務同工開始就很熟了.
後來就產生了信念書.
這個影響很深遠的.
一個信徒在面對時代轉變.
我們基本上十個信念是甚麼.
結果於是一個教會更新運動.
就這樣成立了.
求主都幫助我們.
當我們回顧歷史的時候.
神施恩的手怎樣報受我們.
神施恩的手怎樣帶領我們.
讓我們肯定.
在這個末世多變的世代裡面.
我們仍然需要依靠這位賜我們平安的主.
祂是親自將平安給過我們.
給過我 給過你 願主的平安與各位同在.
天父賜福給各位.
\newpage



\section{出埃及記 25:10}
\label{sec:032075}
\textbf{聖經人聖經事:用皂莢木做約櫃}~~~~黃克勤牧師博士
\newline
\newline
連結: \href{http://www.church.com.hk/acms/content.asp?site=cdc\&op=show\&type=product\&code=032075\&layout=sermon}{\texttt{032075}} ~~~~ 語音日期: 2026-04-25
\newline
\newline
\hyperref[sec:032074]{\small{< < < PREV SERMON < < <}}
~
\hyperref[sec:index_chronic]{\small{[返順時目]}}
~
\hyperref[sec:032076]{\small{> > > NEXT SERMON > > >}}
\newline
\newline
出埃及記 25:10
\newline
\begin{longtable}{cl}
\hline
\hline
章節 & 經文 (和合本修訂版)\\
\hline
25:10 & \begin{tabularx}{0.7\textwidth}{X} 「他們要用金合歡木做一個櫃子，長二肘半，寬一肘半，高一肘半。 \end{tabularx} \\ \\
[1ex]
\hline
\hline
\end{longtable}
$^{1}$大家好,生命說書的網絡和大家繼續看聖經人聖經史.
今次是用造甲木造藥櫃.
即是神在會幕裡面說要給我造一個櫃.
開頭是說要造一個櫃.
不過後面就知道是放法板,所以叫藥櫃.
用造甲木,特別說一下.
在以色列曠野的地方有這種樹.
這種樹叫造甲樹,造甲木.
你看到它在曠野周圍都是乾旱.
以色列地一年下雨大概由十二月中至二月中三個月.
是冬雨,平時一曬下去五十度C.
五十度攝氏或以上在曠野很熱.
全草不生,這些樹就在曠野裡面生存.
平時很少見到它,在曠野就見到.
差不多在曠野經常見到這種樹,其他樹都沒有.
還有這麼乾旱都能夠生存的就是棕樹,Palm,或者是Date,或者是棗樹.
但是神是說要用造甲木來做聖殿或是會幕裡面的器材.
我都覺得很難做,原因是它的木不是很粗.
又不是很高,又不是很長,又不夠幼.
於是如果要做一塊板,一條大一點的就很難找到.
板可能要甲一條條的木,要很多手工.
不過我覺得它都有一個好處,好處是什麼呢?.
就是它很實,重又不重,很堅硬.
它生長得超級慢,如果這棵這樣的樹就是很多年了.
我想一百幾十年才能夠生長到這樣,生長得很慢.
是好像象徵以色列人,或者象徵我們的信仰.
在一個很艱難曠野的地方都可以這樣.
頂得住這些艱難慢慢生長的.
好了,做藥櫃,又要先介紹一下藥櫃.
藥櫃我都在谷歌找的,在Google找的,你會看到.
當做藥櫃的時候,多數是這樣的.
就是一個四方形的,高度其實是好像人的高度.
有兩支很長的柱子可以拿來抬.
照著聖經這樣說的,照著這裡這樣做.
圍著它要有牙的,要有牙的,要有尖的.
然後有兩個基魯柏雕刻,它的翼有翼的,翼伸出來.
兩隻翼要連著的,遮蓋著藥櫃的頂.
藥櫃的頂這個地方,藥櫃的頂這個地方就叫做私恩座.
叫私恩座,神就說我在這裡和你相會,神就坐在這裡.
私恩座,大概是這樣.

$^{41}$好了,其實這個樣子是不是真的這樣?不知道.
因為無數樣子,你去聖地會看到很多種賣的模型都不同的.
這個基魯柏就不是站在這裡的.
是小小一個小了的櫃子,然後翼要重疊,遮蓋私恩座.
這裡有尖的牙,也是這樣抬的,有四條柱.
柱上面有一個圈,這邊伸一支棍去抬.
這個就是藥櫃的樣子.
你見到,剛才有些基魯柏這種天使站著,這個櫃子小一點.
這個又有另一種,這個表達就是裡面擺了兩塊法板.
和一個罐,擺了一些嗎哪和阿倫開花的那些像.
採了一些赤綠.
通常這些像已經是死了的木,竟然連枯木都開花重生.
象徵著上帝的生命的能力.
究竟哪一個才是比較對呢?.
不知道.
這個比較特別,我拍一下,這個就不對了.
原因就是這兩隻形黏在一起.
另外,你見到這兩隻是鳥仔.
今天去到以色列,在一些猶太人的書室.
猶太教書室,裡面是見到有聖殿的模型,有會幕的模型.
有藥櫃的模型,很多種樣子都不同.
有些基魯柏是兩隻鳥仔.
兩隻鳥類,這麼有趣,基魯柏是鳥類嗎?.
其實應該這樣說,到今天離開摩西做的第一個藥櫃.
已經是很久了.
其實也失傳了,現在藥櫃也不見了.
究竟藥櫃是什麼樣子呢?.
是差不多現在沒有人看過,見過的.
原因就是,第一,不是普通人可以看到和摸到.
一定是利未人,雲柱火柱升起了,我們起行了.
他就開始找一些很厚的布,皮蓋著藥櫃.
然後利未人就開始抬.
所以他們也不會說,順手看看才行.
不會的,不行的,不可以的.
很恭敬地,拿著一些布走過去,蓋著蓋著.
蓋的時候也能看到一點.
蓋著才抬出來.
所以百姓看藥櫃是不會看到樣子的.
因此今天藥櫃是怎樣呢?已經沒有人知道了.
大家知道了,不過知道有一樣東西是聖經說的.

$^{81}$好,很快和大家看看藥櫃.
怎樣做的呢?也是用三點和大家分享的.
《初一及記》二十五章第十節.
要用造甲木造一季.
只是說造一個季.
長兩爪半,寬一爪半,高一爪半.
一爪是多少呢?.
就是由手尖到這裡,這樣叫一爪.
一爪大概是尺半.
尺半,即是長三尺.
兩爪半,即是尺半尺半.
長三尺,藥櫃不是很大,長三尺.
寬一爪半,高一爪半.
所以其實也不是很高.
當然包著基督教的翅膀可能高一點.
大概的高度可能不會像一個人那麼高.
要裡外包上精金.
裡外包上精金,即是說譬如那個季的木板.
外面裡面都是包了精金.
被包的是什麼木?造甲木.
四圍雙上金牙邊,所以那些是有牙的.
要鑄四個金環安在季的四腳上.
借哪兩環用來做什麼?.
要用造甲木做兩根弓.
用金包裹,連那支棍也是用金包的.
要把弓穿在季旁的環內以便抬季.
即是說要常在季的環內不可抽出來.
一放了支棍進去就不拔出來了.
不會說要搬了才拿支棍放進去.
這個「降」即是將它抬出來.
必將我所要賜給你的法板放在季裡.
要用精金做私恩座,長兩爪半寬一爪半.
即是說這個藥櫃的頂,那塊板.
長方形,有牙,金的尖角做牙.
然後頂那塊也是精金做的.
原因就是那塊叫做私恩座.
上帝坐在藥櫃上面.
那裡叫私恩座,常說「阿神啊,你在私恩的寶座聽我祈禱」.
其實是指藥櫃的頂,不是天上的寶座.
是藥櫃的頂,叫私恩座.

$^{121}$要用金紙除出兩個基督帕.
基督帕是天使的種類而不是名.
這個天使很多時候在聖經裡面出現.
都是做保安工作的,保衛工作的.
譬如神將犯罪之後的亞當夏娃趕出伊甸園.
不想她回頭又去吃生命樹的果,那怎麼辦?.
於是就派了兩個基督帕把守生命樹的路.
去生命樹的路守著.
並且做一個機關,有四面轉動發火焰的劍.
派兩個基督帕守著生命的路.
另外你會見到在藥櫃的兩邊就是兩個基督帕的像.
名字是Cherubim,就是護衛的天使.
安在私恩座的兩頭.
基督帕是怎樣的?.
借頭做一個,拿頭又做一個.
四方人的櫃,這邊一個,這邊一個,像剛才所見到的.
兩個基督帕要連接一塊在私恩座的兩頭.
這個做法是非常重要的.
所以我把它用紅色字寫出來.
而基督帕要高張翅膀.
所以櫃不是很高,但基督帕要高張翅膀.
遮掩私恩座,保護著.
而基督帕要臉對臉朝著私恩座.
所以這邊一個,這邊一個,私恩座在中間的櫃下面.
它要翅膀伸出來,高高伸出翅膀.
翅膀我估計是一對.
究竟是左邊的翅膀,還是右邊的翅膀.
還是兩隻翅膀呢?.
沒有說,總之翅膀是要貼著的.
要將私恩座安在櫃的上方.
必然要將我所賜給你的法板放在櫃內.
當時特別是將法板放在櫃內,後來才放其他東西.
而我要在那裡與你相會.
又要從法櫃私恩座上.
而基督帕中間和你說我所要吩咐你.
而全給以色列人的一切使.
這個就是今天要和大家讀的經文.
有三點,藍色螢光是我要講的.
我要和大家提醒一下那三點.
第一,要做一個櫃.

$^{161}$櫃是用來擺法板的.
令你記起與神的一種關係.
其實神基本上與以色列的關係是因為藥.
與亞伯拉罕的藥.
與以色列的藥,與雅各的藥.
與西乃山的藥.
亞伯拉罕的藥自然是連帶在以色列國.
所以在講列祖就是亞伯拉罕在以色列國.
神是一個與人有來往的關係.
在聖經裡面是用一個字,就是藥.
藥不是合約,不是契約.
這個藥意思就是,我向你應許了那些.
你沒有守,我都守,起碼我都會守住.
舉個例子,現在分藥,你怨不怨,我都怨.
於是大家就納藥,好像買賣藥.
其實不是這個意思,藥的意思是.
我向你納了藥,然後我要令你的子孫.
很多天上的星,海邊的沙,我會賜福你的後裔.
當時亞伯拉罕是睡了覺.
是神自己在兩個納藥的祭身中間走過.
有火把和香爐,是像神拿著一樣走過.
是單獨單方面納藥.
所以如果分藥,兩個納藥不是的.
是兩個向對方納藥.
所以就算他不健康,他傻傻地,他沒錢.
我都一樣守這個藥.
就不會說你不愛我,我不愛你.
不是的,仍然要守這個藥.
現在我納了這個法板,你們要跟著這樣做.
我就會對這個法板有態度.
因此延伸到你們有態度,很清楚.
你要用造甲木象徵著很堅毅,很艱難.
慢慢都要生長到,來做一個櫃.
是要來承載著這個藥,顯示著很重要.
第二樣就是法板是放在櫃裡面.
於是用法板來做藥,所以開始叫藥櫃.
叫法櫃,以前叫法櫃,就放了法的十言.
上帝在西雅圖山上給他們講了十樣東西.
現在就叫十戒,就放了戒令.
其實是神最重要的十句說話.

$^{201}$跟以色列人的關係,叫他們這樣.
跟神的關係有五戒,跟神的關係有五戒.
跟人的關係,叫他們這樣做.
要來做法的,要去守的,在一個櫃裡面.
這個櫃很重要,以前叫法櫃,現在叫藥櫃.
其實叫法櫃是比較適合的.
有基督巴守著,即是說這件事是神聖不可侵犯.
是重要的.
最要留意的是藥櫃的頂叫私因座.
藥櫃的頂叫私因座,上帝就坐在那裡.
要安一個私因座在那裡.
我要在那裡,第三點.
我要在那裡與你相會.
見到這個叫法櫃.
是將法給你,不是叫藥的當時,最初期是法.
在那裡,我在這裡與你相遇,相會.
相會不是相遇,是相會.
兩個基督巴中間,我要將我吩咐以色列的事告訴你.
這裡有個很寶貴,特別是這裡很寶貴.
原來上帝是歡喜坐在私因寶座面前.
不是坐在一個審判的座位那裡.
坐在一個私因座位.
然後在那裡,我與你相會.
與你就是摩西,但其實摩西是代表全以色列人.
所以上帝說,我在那裡與以色列人相會.
在那裡我會將我的心事告訴他們.
這件事真的很寶貴,原來上帝是想與我們相會.
即是我們要給些時間給上帝.
靜聽祂與我們說甚麼,聽祂在私因座與我們說甚麼.
每一天給些時間與上帝在法季私因座前與祂相遇.
在那裡見祂朝見,讓祂與我們說些說話.
朝見,我用「相遇」這個詞是用得不太好.
太過存在主義想法,相遇,是神在那裡,我要去朝見祂.
大家記得,祂會將一切事情吩咐我們.
我們要樂意去到私因座前,請大家記得,我們祈禱.
天護上帝,我們多謝你,因為你是做私因座,然後坐在上面.
是很歡喜我們去到法季面前,記起你的律法.
你的心意,然後我們在那裡聆聽,與你相會.
幫助我們歡喜與你相會,因為我們很忙,很多時忘記與神相會.
沒有去到那個地方,主我們求你給我們歡喜去到這個安靜的地方.

$^{241}$在私因座前與你相會,求你對我們說話.
我們等候你,禱告奉耶穌名求,阿們.
\newpage



\section{腓立比書 3:12-16}
\label{sec:032076}
\textbf{努力!繼續進擊}~~~~莫潤輝牧師
\newline
\newline
連結: \href{http://www.church.com.hk/acms/content.asp?site=cdc\&op=show\&type=product\&code=032076\&layout=sermon}{\texttt{032076}} ~~~~ 語音日期: 2026-05-03
\newline
\newline
\hyperref[sec:032075]{\small{< < < PREV SERMON < < <}}
~
\hyperref[sec:index_chronic]{\small{[返順時目]}}
~
\hyperref[sec:032077]{\small{> > > NEXT SERMON > > >}}
\newline
\newline
腓立比書 3:12-16
\newline
\begin{longtable}{cl}
\hline
\hline
章節 & 經文 (和合本修訂版)\\
\hline
3:12 & \begin{tabularx}{0.7\textwidth}{X} 這不是說我已經得著了，已經完全了；而是竭力追求，或許可以得著基督耶穌所要我得著的。 \end{tabularx} \\ \\ \relax
3:13 & \begin{tabularx}{0.7\textwidth}{X} 弟兄們，我不是以為自己已經得著了；我只有一件事，就是忘記背後，努力面前的， \end{tabularx} \\ \\ \relax
3:14 & \begin{tabularx}{0.7\textwidth}{X} 向著標竿直跑，要得神在基督耶穌裡從上面召我來得的獎賞。 \end{tabularx} \\ \\ \relax
3:15 & \begin{tabularx}{0.7\textwidth}{X} 所以，我們中間凡是成熟的人，總要存這樣的心；若在甚麼事上存別樣的心，神也會把這些事指示你們。 \end{tabularx} \\ \\ \relax
3:16 & \begin{tabularx}{0.7\textwidth}{X} 然而，我們達到甚麼地步，就當照這個地步行。 \end{tabularx} \\ \\
[1ex]
\hline
\hline
\end{longtable}
$^{1}$各位記者 早安.
(記者)早安.
全因堂九周年.
上次我來應該是你獻堂的時候.
今天來你已經堂慶了.
是開心的日子.
也賣廣告了.
你們九周年了.
我們德勝堂.
今年也自立十周年了.
大家一起準備要做堂慶.
我們也會做堂慶.
小小的廣告.
我們應該六月九日會有感恩崇拜.
邀請大家.
不過是晚上的 不好意思.
我們的崇拜是晚上的.
晚上八時 星期二晚.
是播道會我們借了樂全堂的禮堂.
我們會做感恩崇拜.
如果大家有時間.
可以跟我們一起感恩.
我們也會有一個崇拜之前.
我們也有感恩聚餐.
如果你有興趣跟我們一起飲飲食食的話.
我們也歡迎你.
吃完之後我們就敬拜神.
你也可以跟我們說.
我們也是預備人數.
既然全恩堂也是自立十年了.
今年九周年.
德聖堂也是自立十年了.
我們應該在神的面前.
我們繼續去求什麼呢.
或者我們應該抱著什麼態度.
再向前望呢.
有些事情是要的.
特別我們大家都去到一個時間.
已經不短了.
全恩堂更加有不同的地方來到這裡.

$^{41}$我們更加要去想.
我們繼續向前的時候.
究竟我們是抱著什麼態度呢.
所以今天的題目就是.
努力繼續進擊.
這個是我用了.
我不知道大家有沒有聽過.
日本有一個漫畫.
叫進擊的巨人.
如果你是年輕的會知道.
很喜歡看的.
裡面就說.
很多巨人在外面.
你要去進擊他們.
你要打贏他們.
我把它套用到我們教會.
或者跟全恩堂去看的時候.
其實可能我們在外面.
仍然有很多巨人.
我們怎樣去進擊他們.
進入他們的範圍裡面呢.
所以我希望大家可以一起去看.
但是在這個過程裡面.
我們需要有這個態度.
或者有一些東西我們要去.
自己要留意的.
在這個過程裡面.
所以今天我用了.
剛才菲勒比斯.
大家都很熟悉的一段經文.
剛才已經讀了.
這個是保羅跟菲勒比去說的.
為什麼會說到.
第十二節那裡說到.
我不是說我已經得著了.
就是這件事.
保羅很厲害.
為什麼說了這句話.
因為之前.
如果你看回前面三章頭.

$^{81}$一至十一節裡面.
保羅是在說.
提醒菲勒比教會那些信徒.
有些人你要提防.
賈師傅那些東西.
而當中其中一個就說.
他們很喜歡說.
你要做什麼的禮儀.
割禮啊.
守律法啊那些東西.
這些才是最重要的.
保羅就在前面那裡.
講及這件事.
在那裡說.
如果講守那些.
我比他們厲害.
我是真以色列人.
我又受過割禮.
然後我又守了十戒.
然後我又在哪一個老師門下.
所有的東西.
人家以色列人.
或者猶太人.
他所要求的.
保羅說.
我都完全達標準.
但是.
他最後在前面那裡說.
他說.
我都因.
我信了耶穌基督.
他竟然預了基督之後.
他說.
我願將前面所有的東西.
全部放下.
就是所有覺得榮耀.
人家覺得榮耀.
重要的東西.
全部放下.
而是為了.

$^{121}$得著耶穌為至寶.
他講完這裡.
接著.
他就講第十二節.
但是保羅為什麼在這裡這麼說呢.
就是因為.
信號很自然.
一樣東西.
保羅.
就是你行.
你又可以.
你本來這麼高成就.
接著你又全部放下.
接著你可以做到.
我不是這樣.
那怎麼辦.
就是我覺得.
我們都是第一件事.
我們可以想.
教會.
很多時候我們去到某一個地步.
我們覺得.
行了.
就是.
獻完堂了.
有個新的地方了.
接著就覺得.
可以靜一靜.
休息一休息.
或者.
有一個叫做.
已經得著了的感覺.
所以保羅在這裡.
很明白信濤的心態.
所以呢.
他說.
我不是說我已經得著了.
而是已經完全了.
他說.
我仍然都是竭力的去追求.

$^{161}$他很想告訴信濤.
不要停啊.
我雖然是這樣.
雖然已經好像.
比人家覺得.
是特別一點.
我已經做得到了.
不是等同於信濤做不到.
今天很多時候.
我們都有這樣的心態.
你做得到.
江牧師你做得到.
我不是啊.
或者誰執事.
你可以啊.
我不行啊.
是不是.
就是如果.
當我們在教會的使命裡面.
傳人堂有你們的使命的時候.
我們看到是向前的.
但我們很容易就像.
那些肥辣比信濤一樣.
保羅你做得到.
江牧師你做得到.
執事會主席你做得到.
我未必做得到.
保羅就開始在這裡提醒.
提醒我們第一件事.
我們要.
就是不要.
自滿地停止信濟.
在這裡說.
我不是已經得著.
已經完全.
得著的意思就是一個領袖.
達到目標.
或者成功.
又可以很多個不同的翻譯裡面.
都是講及這個事.

$^{201}$就是保羅就說.
我不是已經成功了.
我不是已經完全了.
已經達到目標了.
同樣地去提醒我們.
今天我們不是達到目標了.
傳人堂已經由柴灣那邊.
過到來北角了.
現在很順利了.
也獻了堂了.
都過了這麼多個月了.
我們覺得OK了.
都不錯了.
保羅說.
不是已經得著.
不是達到目標了.
未是達到目標.
或者甚至我們已經成全了.
已經完全的意思.
他說我是.
還要竭力的追求.
竭力的追求.
竭力的意思.
是逼迫的意思.
逼迫的意思.
追逐獵物的意思.
你追逐獵物.
你不會停下來的.
如果你看動物片.
就知道獅子,豹那些去追逐.
去追逐獵物的時候.
那種快的速度一直去.
保羅就說這件事.
你發覺在第十二節和十三節.
保羅都重複說已經得著了.
我不是已經得著了.
所以不要自滿的停止去進擊.
或者今天有時候我們都會看到.
獻了堂了.
來了一個新的地方.

$^{241}$神都給了我們一個意象.
我們覺得有好的地方.
我們就坐在這裡.
我們自己去自滿地去崇拜.
而是在這裡.
保羅提醒我們.
不要.
我們還有進擊的.
我們還要去的.
當我們來到北角的時候.
當大家來到北角的時候.
我們有沒有眼光再去看.
我們有什麼巨人要進擊.
我們有什麼地方.
我們可以繼續進行.
所以第一件事就是說.
當我們繼續在這個地方.
今天神擺了在我們裡面.
在這個地方的時候.
給了我們這個地方的時候.
我們不要自滿的.
夠了我們已經得了.
不是.
我們還要去進擊.
接著第二件事.
保羅正在提醒我們.
不要自以為.
剛才是自滿.
但現在覺得自己都可以了.
他說已經得著了.
他說只有一件事.
只有一件事.
可以保羅說是一個目標.
但這個目標.
是有兩個的態度.
保羅正在告訴我們.
第一個態度就是不要自以為.
就是覺得我已經.
第一已經拿了很多東西.
拿了很多東西.

$^{281}$他說我們第一個.
我們要忘記.
忘記背後.
忘記這個字是一個現在式.
現在式就是現在進行之外.
還要怎樣?.
繼續進行的.
在原文的文法裡面.
用到這個字.
這個時式的時候.
就不是現在這麼簡單.
所以不要自以為.
我已經可以了.
我好了,樣樣都可以了.
樣樣都有了的感覺.
而是我們要.
忘記一些東西.
或者忘記這個字可以是一種.
就是.
忘懷.
或者置諸腦後.
那個含義,置諸腦後.
將我們之前得到的東西.
已經OK了,置諸腦後.
而保羅在這裡提出一個.
如果忘記背後.
這個態度是一種.
比較消極一點.
好像忘記一些東西.
這個就說.
保羅想提醒信徒一件事.
不要.
給過去.
給過去的成就.
或者所達到.
做得到的東西.
我們就影響.
現在我們的追求.
人很容易的一件事.
就是自以為是.

$^{321}$已經OK了,我過去的成就.
我過去所做的東西.
接著就停步了.
如果你看整個國家.
都是這樣.
整個大群體都是這樣.
何況教會.
當我們一停低.
自以為我很厲害的時候.
就會停了.
或者傳真堂.
可能在過去.
由上年我們已經拿到.
新的地方.
我們找了新的地方.
亦都能夠很滿意地.
大家可以在這裡一起敬拜.
有時我們就會.
很容易去到這個地步.
我們甚麼都行了.
有甚麼事工.
我們再慢慢來吧.
不是的,保羅說.
不要用這些東西.
來影響我們.
現在所做的事.
保羅正在提醒我們.
這個情況.
正在提醒我們.
然後呢.
保羅就說.
你努力的前面.
努力的前面.
努力又是一個.
現在時態.
所以在這裡.
你發覺保羅說的.
消極那邊.
就是你不斷地要將你過去的成就.
你覺得很厲害的東西.

$^{361}$不斷地磨忌.
拋諸路後.
但是另一個積極的態度就是.
努力前面.
好像賽跑一樣.
向前.
又是這個字.
他在描述.
努力這個字.
在描述賽跑者.
身體向前傾.
如果你是跑步.
你就知道你不會這樣跑.
你不會向後跑.
人的身體是不會向後的.
你一定向前的.
你看所有跑者都是這樣.
身體是向前傾的.
這個努力的意思就是.
你好像向前一樣.
一路向前去走.
他還要是在跑的.
好像賽跑一樣.
所以他接著就說.
再重複說一個.
是竭力的.
這個努力是去到什麼.
追求那個目標.
有一個目標.
就是剛才說的.
去追獵物.
在前面有一個獵物.
你在那裡向他追.
追他.
今天.
全銀行.
我們有.
即是你有你的使命.
或者你有.
你自己服事的福金的對象.

$^{401}$那些對象.
或者我們要做的事.
我們有沒有好像當他.
獵物那樣去追.
所以保羅說.
竭力我努力前面.
向前去追.
一個積極的.
今天.
全銀行九年了.
九周年了已經.
過去可能我們已經做了很多事.
在柴灣那邊都很成功.
然後來到北角的時候.
我們應該放下所有.
柴灣的東西.
置諸路後.
而今天我們就要向前.
去望.
我們要望前面.
好像賽跑者我只可以.
向前的你不可以向後.
向後望你知道結果會怎樣.
即是如果在賽跑.
你看人家在比賽.
如果那個考手.
突然向後望一望.
你知道後果會怎樣.
一是跌倒.
一是被人標過.
自己會慢.
所以在這裡.
全英堂我們一齊.
去看我們的使命.
我們一齊去.
是努力前面.
努力前面.
就是向著這個念脈.
我不知道有甚麼.
form隊長.

$^{441}$都做了很多form隊長.
這些form隊長是不是大英姐妹.
我們大家都成為一個好像當他們是念脈.
我們遵從它.
我們努力.
向著前面去跑.
而不要.
自我放棄.
哎呀 不行的.
有些信徒就.
做不到了.
我只有這麼多.
我就停下來.
自我放棄.
或者覺得.
哎呀 江牧師你做到了.
我不行的 我不是你.
我寶蘿就最怕這件事.
有肥立比的人說.
寶蘿你是而已 我不是你.
你可以繼續去.
你繼續傳福音.
你在監獄裡面繼續去做事.
我不行的.
就停下來.
很容易自我放棄.
覺得到這裡了.
停下來.
所以在這兩個.
大家一起的勸勉.
其實我都會勸勉.
我的姐妹也是.
不要停下來覺得現在很好.
一直在說某些情況.
某些聚會.
我們有很多人.
或者我們已經進入到某一個圈子.
我們不可以在這裡.
就行了.
或者有些人覺得你們可以做到.

$^{481}$我做不到.
你們可以的.
其實不是的.
當你這樣說.
就是自我放棄.
我自我放棄.
所以在德聖堂.
我用德聖堂的例子.
因為我自己在德聖堂.
他們很多時候都說.
你去做社團人士.
你們可以的.
我都不知道怎麼談.
我們不行的了.
是不是呢.
我們有沒有努力去做呢.
我們有沒有努力去嘗試呢.
所以我鼓勵很多弟兄.
你來吧.
你不要自我放棄.
坐在這裡.
同樣地今天.
全銀行一樣.
你要知道你的獵物在哪裡.
你就不要想著.
只有江牧師能做到.
只有執事們他們能做到.
我做不到.
然後就坐在這裡.
牧師你去吧.
牧師你做吧.
不是.
是大家一起的.
去做這件事.
所以保羅在這裡學得很清楚.
很清楚.
有個目標.
保羅就希望我們一起.
最終極目標就是耶穌基督.
所得著我們.

$^{521}$去到耶穌基督那裡.
但是我們可以將這個.
這個目標大家都一起.
但是我們去看現實.
現在的情況.
我們的目標.
是不是要去努力.
去努力.
忘記過去.
無論柴灣那裡是怎樣都好.
放下.
我們再向目標.
現在是北角.
既然是北角.
我們有什麼獵物呢.
江牧師已經定了很多.
給我們看.
但是我們是不是將這些.
覺得成為我們追的目標.
跟江牧師一起.
去追這個目標.
我在這裡不是叫大家.
去吧去吧.
進擊那裡去吧.
不是 保羅也不是這樣.
保羅也很小心.
也明白信號的處境.
每個人都一樣.
所以他最重要.
就是十五和十六這兩節.
他說我們中間.
凡是成熟的人.
成熟的我們都是成人.
都應當這樣.
思想就是說前面那裡.
說的那個思想.
思想又是一個.
現在式的.
都是這樣去想.
不斷進行的.

$^{561}$不斷的去想.
他說成熟人.
你一定要這樣想.
就是保羅前面說的那樣東西.
你要努力 你要前面.
你不要自我放棄.
你不要自以為是 也不要自以滿.
你不要這樣做.
我們應該這樣.
共同的有一個目標.
大家一起去的.
他說即使.
你們不是這樣思想.
神話會把這事指示你們.
不過.
我們到了什麼程度.
就要照著.
什麼程度去走.
保羅在這裡很小心的.
是明白信號的不同.
可能.
定了一個目標.
可能我們走.
大家都有這樣的思想.
但是在走的過程裡.
在進行的過程裡.
可能在程度上.
大家有所不同.
大家的步伐都不同.
大家思想上都會不同.
保羅就說.
不要緊.
最重要的是.
我們明白我們自己去到哪裡.
我們可以.
能夠做到什麼.
我們就做什麼.
而在這個能做到什麼的程度裡.
是互相的配合.
這個配合是什麼.

$^{601}$配合著.
我們要去那個目標那裡.
不是那個程度.
我覺得我不行了就坐在那裡.
而是我們.
背後怎樣能夠.
彼此的.
互相配合在當中.
慈隱堂有一個使命.
大家很清楚在前面.
但在這個過程裡面.
我相信大家有不同的程度.
有不同的限制.
我們明白.
但在這個限制裡面.
不會成為我們停低了在那裡.
我不再去進擊.
不再在一起.
如果我們在全隱堂裡面.
我們大家都是一個團隊的話.
我們就要發揮我們團隊精神.
怎樣互相配合.
保羅都提醒過我們.
萬事都是互相考慮.
你在哪個程度.
你就做那個程度.
做一個支援也好.
或者做一個怎樣協助也好.
我們一起的向著這個目標.
如果你有機會去看一些.
動物片.
特別有些是.
追逐獵物.
有一些是會團隊出發.
用團隊.
不是一隻.
炮通常都是一隻.
獅子也是會團隊.
獵狗也是.
胡狼也是.

$^{641}$他們很懂得怎樣一起.
去進擊那個獵物.
有人向前.
有人試咬他.
然後又走回去.
然後第二個又上去.
我們純粹看他們追逐獵物.
如果你看回這個情況.
在一個團隊的時候.
我們都是一個團隊.
我們在哪個崗位.
我們就在那個崗位.
我知道我自己是做到這件事.
保羅提過.
我們就好像一個身體.
每個人都有一個功用.
所以我們知道我們自己的位置的時候.
我們怎樣配合.
一起的.
向著這個使命.
去做這件事.
所以希望黎英姊妹.
我們現在要學的一件事.
不是.
只有一個人去做.
而是整體一切.
互相的配合.
我知道戈牧師.
他都很努力.
但有時需要人的配合.
我相信大家都知道.
我們做牧者.
很希望弟兄姊妹一齊配合.
德聖堂也是這樣.
我們鼓勵到弟兄姊妹.
他們可能明白自己的限制.
我們也明白他們的限制.
但我們怎樣可以幫到他們.
去發掘到.
自己的恩賜.

$^{681}$知道自己的位置.
而去做到應該做的事.
我們是需要這樣去做.
所以在這裡裡面.
保羅的提醒.
叫我們繼續.
去進擊的時候.
我們要有這個態度.
我們要先撇除它.
不要自以為是.
不要自我放棄.
更加要清楚.
自己的能力.
限制.
而是彼此.
配合的.
彼此配合.
以致教會.
可以向著標竿直跑.
向著目標.
而去做.
向著這個獵物.
去進擊.
哪怕外面還有很多巨人.
或者我們看到.
很多不同的範疇.
有很多巨人.
很多可能令到你卻步的.
但因為我們有同心合意.
我們這個團隊.
可以能夠去進擊.
繼續的.
是進擊.
所以希望傳統教會.
在未來的日子裡面.
更加的一起合一.
團結.
可以去進擊.
這些巨人.
我們追擊我們的獵物.

$^{721}$我們的福音對象.
我們一起祈禱.
推廣你.
帶領傳統教會.
一段時間裡面.
很多的領域.
讓他們能夠.
看到神利的恩典.
但讓我們不要停留.
在這個恩典當中.
只是這個地步.
那是我們竭力的一切.
向著前面.
我們的標桿.
神利所為我們所預備的.
繼續努力奔跑.
最後求你恩代傳學.
繼續保守前面的路程.
保守每個弟兄姊妹.
我們忘記背後.
我們努力前面.
得到神為我們所定的.
一切目標的掌上.
求你保守我們.
我們同心祈禱.
奉主耶穌命求.
\newpage



\section{創世記 35:1-15}
\label{sec:032077}
\textbf{把偶像埋在橡樹下}~~~~羅泰然牧師
\newline
\newline
連結: \href{http://www.church.com.hk/acms/content.asp?site=cdc\&op=show\&type=product\&code=032077\&layout=sermon}{\texttt{032077}} ~~~~ 語音日期: 2026-05-03
\newline
\newline
\hyperref[sec:032076]{\small{< < < PREV SERMON < < <}}
~
\hyperref[sec:index_chronic]{\small{[返順時目]}}
~
\hyperref[sec:032078]{\small{> > > NEXT SERMON > > >}}
\newline
\newline
創世記 35:1-15
\newline
\begin{longtable}{cl}
\hline
\hline
章節 & 經文 (和合本修訂版)\\
\hline
35:1 & \begin{tabularx}{0.7\textwidth}{X} 神對雅各說：「起來！上伯特利去，住在那裡。在那裡築一座壇給神，就是你逃避你哥哥以掃的時候向你顯現的神。」 \end{tabularx} \\ \\ \relax
35:2 & \begin{tabularx}{0.7\textwidth}{X} 雅各就對他家中的人，以及所有和他一起的人說：「除掉你們中間外邦的神明，要自潔，更換衣服。 \end{tabularx} \\ \\ \relax
35:3 & \begin{tabularx}{0.7\textwidth}{X} 我們要起來，上伯特利去，在那裡我要築一座壇給神，就是在我遭難的日子應允我，在我行走的路上與我同在的神。」 \end{tabularx} \\ \\ \relax
35:4 & \begin{tabularx}{0.7\textwidth}{X} 他們就把手中所有外邦的神明和自己耳朵上的環子交給雅各；雅各把它們埋在示劍那裡的橡樹下。 \end{tabularx} \\ \\ \relax
35:5 & \begin{tabularx}{0.7\textwidth}{X} 他們起行。神使周圍城鎮的人都驚恐，就不追趕雅各的兒子們了。 \end{tabularx} \\ \\ \relax
35:6 & \begin{tabularx}{0.7\textwidth}{X} 於是雅各和所有與他一起的人到了迦南地的路斯，就是伯特利。 \end{tabularx} \\ \\ \relax
35:7 & \begin{tabularx}{0.7\textwidth}{X} 他在那裡築了一座壇，給那地方起名叫伊勒‧伯特利，因為他逃避他哥哥的時候，神曾在那裡向他顯現。 \end{tabularx} \\ \\ \relax
35:8 & \begin{tabularx}{0.7\textwidth}{X} 利百加的奶媽底波拉死了，葬在伯特利下邊的橡樹下；那棵樹名叫亞倫‧巴古。 \end{tabularx} \\ \\ \relax
35:9 & \begin{tabularx}{0.7\textwidth}{X} 雅各從巴旦‧亞蘭回來，神又向他顯現，賜福給他。 \end{tabularx} \\ \\ \relax
35:10 & \begin{tabularx}{0.7\textwidth}{X} 神對他說：「你的名原是雅各，從今以後不要再叫雅各，你的名要叫以色列。」於是，神就叫他的名為以色列。 \end{tabularx} \\ \\ \relax
35:11 & \begin{tabularx}{0.7\textwidth}{X} 神又對他說：「我是全能的神；你要生養眾多，將來有一國和許多的國從你而來，又有許多君王從你生出。 \end{tabularx} \\ \\ \relax
35:12 & \begin{tabularx}{0.7\textwidth}{X} 至於我賜給亞伯拉罕和以撒的地，我必賜給你；我必賜這地給你的後裔。」 \end{tabularx} \\ \\ \relax
35:13 & \begin{tabularx}{0.7\textwidth}{X} 神就從與雅各說話的那地方升上去了。 \end{tabularx} \\ \\ \relax
35:14 & \begin{tabularx}{0.7\textwidth}{X} 雅各就在神與他說話的地方立了一根柱子，就是石柱，在它上面獻澆酒祭，又澆油。 \end{tabularx} \\ \\ \relax
35:15 & \begin{tabularx}{0.7\textwidth}{X} 雅各就給神與他說話的那地方起名叫伯特利。 \end{tabularx} \\ \\
[1ex]
\hline
\hline
\end{longtable}
$^{1}$說到偶像,香港每年都會有一個選舉.
就是人氣偶像榜,或者叫做風雲榜.
你們有沒有參加選舉?.
你們都過了這個年紀了.
年青人有沒有參加內地的偶像選舉?.
你們應該是這個年紀了.
原來是每年都有的.
在樂壇裡面特別是點名,點播,投票.
選出自己心目中最喜愛的歌手或者偶像.
當這個榜一出,應該是十個人.
媒體大肆報導,粉絲就說我那個偶像贏了.
大家知不知道十個裡面男多還是女多?.
男多,十個有九個是男的,一個是女的.
大家想到了,十個有九個是男的.
可能很多是來自同一隊的組合.
MIRROR,大家都聽過,聽過名字MIRROR.
我太太跟我說,4月30號,前幾天,千萬不要去銅鑼灣.
大家知不知道為什麼?.
原來其中有一個成員生日.
那天銅鑼灣就不再叫銅鑼灣了.
叫什麼灣?.
你們有沒有參加選舉,大家就知道了.
就是姜濤灣,就是這個姜濤.
十二個裡面,我只能說出他的名字,姜濤.
其他的就要想一下.
但實際上,偶像這種風氣,這種情況.
都是在一個現代的社會,商業的社會被人無可厚非.
這些都是商業化了娛樂的事業.
當然,剛才在聖經裡面說的.
偶像不是單單說舞台上的追星的偶像.
也不是單單說挑剔一些神像.
最簡單來說,任何取代了神.
在我們生命中首要位置的事物,人物.
就會成為我們的偶像.
所以今天這個題目是把偶像埋在橡樹下.
這是一個什麼事件呢?.
因為當時是雅各和一家人.
是經歷一個危機的時候.
這是什麼事情呢?.
可以說,大家很記得雅各和伯達利.

$^{41}$會記得是哪件事呢?.
雅各逃避,逃亡他哥哥追殺二嫂.
去到夜間有個地方,見到和神交替.
見到夢中有天使上上落落在這個地方.
他就說,如果神救我,我就會在這個地方祝壇.
他覺得這個地方是神的殿.
所以叫伯達利,伯達利是這個意思.
希望在當中敬拜神.
但是和這件事相差了多少年呢?.
大家知不知道?.
是已經三十年後的事了.
三十年後的事已經解決了和哥哥的復和.
已經去到安居在這個事件的地方.
但是發生了什麼事呢?.
是一件悲劇.
就是他的女兒叫做底娜被事件強暴了.
引致了家族很大的危機.
當然他哥哥會看到他們用一個想法.
就是很殘暴的,要去報復這些事件人.
所以開頭就欺騙他們.
他說,我都願意把妹妹嫁給你,嫁給強暴的人.
這樣就完成了這件事情,不再追究了.
但是表面上是這樣說.
但實際上他哥哥用了一些計謀.
就是說,好,我們這樣做都可以.
你們的男丁一定要受割禮.
當他們願意了,受割禮了.
但是男人最痛的時候,他們就會殺過一個不留.
就讓他們屠殺了全城的男丁.
這個這麼殘暴的做法.
當然令到士劍一帶的那些外族.
那些城鎮美人.
對於雅各一家很憎恨.
覺得他們很殘暴,明星敗壞.
所以他們就在一個危機裡面.
隨時被人結合一起去報復,去滅族.
在這個危機的一刻裡面.
就看到三十五章,一到十五節來到這段經文.
神就跟雅各說,起來,上伯達利去那裡祝壇.
這個就開始了今天這段的內容.

$^{81}$但事實上我們看回這個事件裡面.
當然我們會去很多的學習提醒.
特別我們今年教會的主題,結連傳承復興.
特別在復興的主題裡面.
我們會很多的分享,很多的實踐.
無論我們剛剛昨晚有一個敬拜祈禱會裡面.
我們都是尋求復興.
梅利牧師給我們的分享都是要回轉背改,認罪的內容.
所以我們都是繼續在這方面去幫助弟兄姊妹.
在這個主題裡面,是希望我們能夠完完全全地回到神那裡.
我們今天來到神的面前,我們一定要坦然無懼.
我們不會有些遮遮掩掩在神的面前.
就算我們是有罪惡.
我們能夠去打開我們的心.
求主去潔淨我們,認罪背改.
以致我們能夠完全跟我們的神相通.
在這段經文裡面就看到.
事實上當時雅各的一家裡面.
宿命的光景卻不是這樣.
今天的題目是把偶像埋在橡樹下.
分開三方面.
第一方面是離罪自結,埋掉偶像.
第二方面是復興的保障.
第三方面是從雅各到以色列.
第一方面是離罪自結,埋掉偶像.
首先說回那件事.
如果神跟雅各去追究.
是一件很嚴重的事情.
整件事裡面,如果用十誡.
可以犯了六誡,最少.
尤其是我們過去一年開始.
上年五月開始有個認罪背改祈禱會.
我們用十誡開始.
我們有參加都明白.
十誡有很多地方.
神要我們去遵守.
有很多原因.
很多原因都是希望我們跟祂不會有隔絕.
但是你看到這群屬於神的人.
雅各的家庭竟然犯了十誡.

$^{121}$最少有六誡.
這件事件裡面.
不可有別的神.
你也知道他們有很多敬拜偶像.
敬拜外邦的神.
不可挑剔偶像.
他們還把偶像藏起來.
放在身上.
不可殺人.
殺了這些事件人.
不可姦淫.
這件事是姦淫的事件.
不可偷渡.
他們殺人之後.
偷取他們的財物.
不可作假見證.
因為底那兩個哥哥.
西面和里面.
欺騙他們.
說他們會和他們聯親.
但其實是假的.
被迫殺了他們.
起碼有六個十誡.
犯了十誡.
這麼厲害的時候.
神跟雅各說.
雅各一定很害怕.
會否追討他們的罪.
追討他們家人的罪.
但這裡看到神的管教.
有慈悲和愛臨到當中.
這場血腥的屠殺.
足以令雅各意識到自己的家族.
會陷入很大的絕境.
會被人追殺滅族.
是走投無路的感覺.
但聽到神的聲音.
究竟是不是這樣呢?.
但神給他聽到的是這樣.
第一節.

$^{161}$上伯特利住在那裡,在那裡築一座壇給神.
就是當日你逃避哥哥耳掃的時候.
向你顯現的神.
他想到這個時候.
這個神又會再開恩給他.
又會在走投無路裡.
又會再一次給他有恩典.
所以雅各體會到.
是時候要回轉了.
是時候要回轉.
當人在困境裡.
我們會看到甚麼地方呢?.
如果我們有一個鏡頭.
我們有沒有想到.
我們是不是要回轉?.
是不是要回轉回神那裡呢?.
實際上他們當日.
在事件居住.
是安於現狀.
安於現狀最重要就是.
無風無浪.
能夠像上班一樣.
無事無幹.
又到五點鐘放工.
最好是這樣.
但是在屬靈生命卻不可以這樣.
無風無浪.
我們就很好的生活.
我們沒有甚麼打擾.
神又沒有甚麼特別的使命給我做.
我都會準時回到教會裡面崇拜.
有基本上的祈禱.
這樣就過我們屬靈的生活.
但是靈性慢慢麻木.
這樣我們慢慢聽不到神給我們微小的聲音.
慢慢我們感受不到神的催迫.
他們就在這樣的生活底下.
但是神有時真的跟我們開玩笑.
藉著一些生命的難題.
給我們一些甚至悲劇.

$^{201}$來到我們當中.
要引起我們去仰望我們的神.
專注我們的神.
所以這個復興往往不會在一個安逸裡面.
會臨到我們當中.
而是一個危機出現的時候.
我們要去醒覺過來.
就像今天.
你有沒有危機呢?.
你自己屬靈的生命你有沒有危機呢?.
生命走了一個方向.
你會不會越來越疏遠神.
或者你只不過是一個掛名的基督徒.
我們沒有仔細去體會.
我們那個屬靈的生命.
亞國家庭的生活.
那個新光景就是這樣.
即使在這個事件裡面有竹壇.
表面上仍然敬拜神.
但心裡面卻是遠離了神.
這裡看到神怎樣去讓他體會到自己的問題.
又讓他看到神自己的管教.
和讓他回轉.
神要他回到他曾經在那裡許願的地方.
伯特利.
這是多少年前的事呢?.
是三十年前的事了.
即是他逃避哥哥的時候.
是三十年前.
如果再計算其他時間.
他又去了那裡呢?.
舅父拉班那裡.
做了一年跟他打工.
娶了兩姐妹.
生兒育女.
那裡又過了二十年.
即是說他來到事件.
也有十年的時間.
十年時間裡面.
剛才我們說他們的光景.

$^{241}$無風無浪就可以了.
誰知搞出這樣的事.
好了.
但是.
神是一個提醒.
你要回到那裡.
你要去還願.
你能夠立誓.
說要祝壇.
但是你走了就走了.
就再沒有去記得這件事.
而誤了三十年.
神是記得的.
提醒他.
於是雅各就知道.
其實自己對神的一個關係.
對神的敬畏.
其實已經放在一邊.
就是所謂過橋的秋板.
忘記了.
誓言延了幾十年.
悲劇來到.
才去悔改.
但是這些情景.
仍然是神的一個管教.
仍然是帶著一個恩典.
愛在當中.
不過今天這些經文提醒我們.
不要再延遲去認罪悔改.
不要再延遲去信服跟隨神的大靈.
其實回到伯達利.
就由這個使劍過去.
其實有多遠呢?.
大概四十八公里.
大家知不知道馬拉松有多長?.
四十二公里.
我問了馬拉松的弟兄.
如果平均一小時走四公里.
即是十二個小時.
你當然不會走完十個小時.

$^{281}$三天也好,兩天也好.
你都走完了.
大家知不知道今天很多人來到香港.
特別內地一些人來到香港.
他們不打算住在一個酒店.
他們全部帶著重裝備.
背囊,螢幕,睡袋等等.
知不知道他們去哪裡露營?.
他們走麥理浩徑.
麥理浩徑是一百公里.
通常五日四夜.
當然不是比賽那種.
很多人,今天一時也很多.
你們有沒有帶背囊來?.
你們有地方住.
我看到一些視訊.
走麥理浩徑都要排隊.
看某個風景某個位置都要排隊.
多厲害.
但都是一百公里.
都那麼多人去.
是不是那麼難?.
現在四十八公里少一半.
為什麼都不去那個地方?.
所以不可以說太遠.
尤其是古人行慣路.
有什麼難處?.
所以看到這個屬靈的疏難.
一樣的.
回來教會很近.
但又不回來.
就是疏難.
你說住遠了.
香港遠極都有限.
半小時去元朗.
半小時去東涌.
在這裡計算.
不會很遠的地方.
但我們有很多藉口.
講這些情況.

$^{321}$但要回到起點.
回到最美好和最感動的一刻.
神就要他們回到那裡.
回轉的方向.
通常是回到最初和神相遇的那種感動.
我們每個人都有這種感動.
一刻裡面的感動.
是什麼感動?.
初戀時的感覺.
這樣的情況.
當然是美好的.
我最近不約而同都和一些頂展妹去接觸分享.
有時我看到頂展妹的那個屬靈的光景.
是有出現了問題.
或者生活上都出現了問題.
很多時我都覺得.
你找回你當初和神的納弱.
當初你那種的納智.
真是一講到這裡的時候.
大家都會黯然.
是很慚愧的.
曾經年青.
或者是早年的時候.
納智事奉.
納智親近神.
怎會好像其他人那樣.
回不回呢?.
怎知現在自己就是這樣的情況.
為何好像年輕人那樣.
真的不會親近神呢?.
現在自己慢慢都變了這樣的情況.
所以我們要回到最初和神相遇的感動.
這個是回轉的方向.
你要去重新建立這個關係.
然後你看到他們.
雅各就和他們講全家人一個行動的方案.
行動的指南.
就是第三節那裡.
除掉你們中間外邦人的神明.
要自己更換衣服.

$^{361}$這個是要重新連結的時候.
必須去解決那些難阻.
剛才講很多的難阻.
難阻可能是自己心裡仍有罪.
我們來到神的面前.
都是遮遮掩掩.
都是做個樣子.
有時慢慢覺得.
不來好過來了.
不去敬拜好過去敬拜了.
自己都不敢坦然無懼.
我們未解決罪的問題.
會是這樣的情況.
所以罪就繼續隔絕了我們和神的關係.
很多的難阻都是在講這個罪的問題.
他們除掉了外邦的神像.
當時雅各的太太拉傑.
曾經偷了拉班家裡的神像.
他們去到事件裡面.
殺了那些人之後.
又搶了那些所謂敬拜外邦人.
拿到那些附身符.
掛在身上裡面.
亂了他們的情況.
當然在今天來說.
我們不會.
我們家裡不會拜偶像.
我們信了主.
我們沒有這些.
但剛才我說.
當我們生活上任何和神爭奪第一位的事物.
那些已經成為我們的偶像.
即是我們心裡很重要的依靠.
真的沒有它是不行的.
我不是說我們信了神.
我們甚麼都不要.
而是那個態度是以神為首.
神自然會給我們.
但我們不是.
我們不是看神.

$^{401}$我們會去尋求金錢.
明星.
我們的工作.
我們的崗位.
我們的家庭.
我們的目標.
甚至自我.
我們如果這些擺在最首要的時候.
這個就成為了敬拜偶像.
敬拜偶像是一個罪.
所以罪.
你就會明白.
不是單單說我們德行上的一些過犯.
即是看得到.
聽得到.
一些很明顯的罪.
但是直到核心裡面.
將那些東西代替了神.
已經是罪了.
將那些東西排在神的面前.
已經是罪了.
這個我們一定要在神的面前.
將這些難阻去拆除.
才能夠去到神的面前.
所以雅各就說.
你要自潔更換衣裳.
如果你在字面上說.
就說洗個澡.
換件衣服.
乾乾淨淨.
當然在屬靈的生命裡面.
就是潔淨我們自己.
我們要心裡面的.
內心裡面的潔淨.
我們才能夠重新可以.
享受和神的同在和祝福.
我們才能夠坦然無僅神的面前.
能夠去站立.
因為神是聖潔的.
來到神的面前的人.

$^{441}$都要聖潔.
守潔心清.
才能夠去登到耶和華的聖山上.
新教裡面提醒我們.
好像換衣服一樣.
脫去舊人的行為.
穿上新人那種品格.
所以聖潔的定義.
就是我們要成為完全.
但是如果在一個字節裡面.
怎樣完全呢?.
我們不要做表面功夫.
做個樣子就不完全了.
我們局部清潔.
如果在家裡打掃.
垃圾掃完就放在地底裡面.
掃到櫃底裡面.
見不到就算了.
這些不是真正的潔淨.
真正的要完全的潔淨.
我們每天都洗澡.
洗乾淨.
我們不會堆在一起.
一個禮拜洗一次.
到時候都很乾淨.
我們不會這樣.
我們會覺得不舒服.
但是同樣的.
內心裡面有污穢的時候.
越來越多.
根本不敢去清洗.
不敢在神的面前去認罪.
反而你每天去清洗.
少少地你都覺得OK.
那怎樣清洗我們內心呢?.
自從去年我們有五月一個認罪悔改的祈禱會裡面.
我發覺你們開始就祈禱悔改.
認罪.
我發覺每樣東西.
十戒裡面.

$^{481}$一路提醒.
多多少少都有些得罪神的地方.
有些得罪人的地方.
所以發覺人都是很容易.
沾染了一些罪污.
我們就去祈禱去認罪.
清洗之後整個人真的爽快了.
無罪一身輕的感覺.
我們不會覺得來到神的面前.
好像不敢去爭取更多.
有時我們一個清心的人來到神的面前得見神.
能夠一個清心的人去祈求神扣門.
為你開門.
但有時我們不敢.
但認罪悔改的禱告.
慢慢就成為了我其實每一日.
一開口跟神禱告就是認罪悔改.
昨日我有甚麼地方做得不好.
想想有甚麼地方得罪了神.
認罪悔改.
我們能夠看到每日潔淨的自己.
求主幫助我們.
今天能夠聽到你的心意.
今天能夠行到你自己照我所做的.
認罪悔改就是一個自潔的實踐.
復興都是要從認罪悔改開始.
所以雅各就要家人自潔更換衣裳.
這個不是表面的.
是內心裡面做好.
還有一個行動就是.
雅各自己將收集到的偶像.
我不知道雅各是否一起收集.
起碼拉傑也收集了30年的偶像.
然後拿了很多護身符.
還說了耳朵上的環子.
簡單來說是耳環.
那些是有宗教的味道.
迦南宗教.
於是埋在事件的橡樹底下.
35章第4節.

$^{521}$吩咐他們全部拿出來.
雅各就開始做好這個行動.
這個代表徹底的斷絕.
再去追求偶像.
追求神以外的一些期望.
埋掉代表了死亡和終結.
雅各沒有說這些不要緊.
這些像不如我用掉.
很多是金.
變了其他金都叫做保值.
都可以這樣做.
變了不是那些偶像.
不過雅各沒有這樣做.
也沒有轉手賣給別人.
賣給你們.
沒有這樣做.
完全埋在橡樹底下.
這種徹底的從生命裡面去砍掉它.
這個象徵了這些東西.
在生命裡面不再有價值.
以前捍衛自保.
現在捍衛糞土.
自保的是我們的主耶穌基督.
不再有影響我們的生命.
成為攔阻我們去神的面前.
今天很多人不能夠長期.
或者慢慢疏離了神.
最簡單就是慢慢一步一步的攔阻了.
在神的面前.
一定是讀經祈禱少了.
慢慢慢慢生命裡面已經和神隔絕了.
但是雅各很清楚一件事情.
這個危機來到的時候.
很大的提醒.
他不再給自己後路的決心.
完全是去清理了它.
所以埋在橡樹底下.
對於我們今天的人的意義是.
找出這些東西.
這些是偶像.

$^{561}$銅像.
耳環.
附身符.
是什麼?是令你分心.
令你不再依靠神的事物.
不要再依靠那些東西.
以一種不再回頭的姿態去處理它.
當然你說我以前很重視錢.
走去賺錢.
很緊張怕沒錢.
現在把錢埋在肚子裡面.
一元都沒有.
不是這樣的意思.
最重要的是.
你自己不再看那些是最重要的.
將神視為首位.
反而先求神的國和神的義.
其他的東西.
其他的是什麼?.
其實我們所關心的東西.
憂心的東西.
其實都不用憂心.
神會賜給你的.
神賜給你出人意外的平安.
甚至一些我們意想不到的.
神賜給我們的好東西.
所以這裡提醒我們.
究竟今天有什麼是成為了偶像.
去難了我們和神的關係.
有什麼是令我們日思夢想.
但是卻離開了神.
有些什麼是罪.
是我們尋求罪終之樂.
使我們難了我們親近神.
亦成為了我們心中的偶像.
你願不願意立志埋掉在肚子裡面呢?.
第二第三點我會很快講完.
可以講這些行動之後.
帶來的一些結果或好處.
第二方面是復興的保障.

$^{601}$第五節說他們起行了.
神使周圍城鎮的人都驚恐.
不追趕雅各的兒子們了.
其實他們當時是驚弓之鳥.
已經殺戮了這些事件.
但是周圍的人都想滅了他們.
他們都不敢走出來.
但是這次神說你要出來.
要去百達利.
雖然不是很遠的路程.
但是你會面對到這些人追殺.
但是神主動令到周圍的人驚恐.
為什麼呢?.
這是很重要的.
一個復興的生命.
看到不信的人都會害怕.
神為什麼這麼厲害.
改變人心.
神為什麼這麼殘暴的人.
改變了他們的生命.
可能變得溫柔了.
變成他們完全不同的面目.
這是一個復興的現象.
外面的人會看到.
這個神很厲害.
這個神是值得尊崇.
值得畏懼.
所以城外的人都驚恐.
什麼神來的.
這麼有能力去改變這些暴戾的人物.
當然他們看到神親自保護.
神成為他們最好的防禦.
所以真正的安全感.
不是來自我們的武力.
或者我們金錢有什麼勢力.
而是我們和神的關係.
當我們除掉了這些偶像.
神會親自成為我們的盾牌.
成為我們的保障.
所以一個雅各再稱呼伯達利.

$^{641}$再加一個字.
正經為伊勒伯達利.
當然這個伊勒.
伯達利本身是希伯來文的神的意思.
全能者的意思.
全能者再貫了一個名稱在上高.
全能者的神的殿.
這裡很簡單看到.
悔改的人必定得到神的保護.
縱使過往很多地方得罪人.
很多人不喜悅.
但神會保護他.
神會保護一個人.
今天我們離棄了神的人.
我們當然很怕面對神.
但我想告訴你.
神仍然愛你.
神是期望你找回當日的初心.
你的立志.
神仍然期望你成就.
仍然期望三十年前的雅各.
他的起誓.
他要去祝壇敬畏神.
所以神要趕絕你.
神要繼續愛你.
繼續讓你回轉過來.
第三方面.
從雅各到以色列.
這裡神顯現給雅各聽.
你的名字不再叫第十節了.
雅各.
現在叫做以色列.
大家記得在雅各準備見哥哥的時候.
他曾經在一個雅博道口跟神較力.
神人.
神人說你叫做以色列.
這裡又過了十年的時間.
但是神再次跟他說.
你不要再叫雅各.
你叫做以色列.

$^{681}$這裡給我們看到.
不是單單名字改換.
因為雅各不是再一個.
捉住腳跟的人的意思.
而是跟神較力得勝的人.
即是你是滿有神的力量.
所以復興的生命是帶來本質的改變.
帶來生命的轉化.
雅各以前代表自己靠自己努力.
靠自己那些小聰明去掠奪.
但是以色列代表靠神的力量去得勝.
以靠神能夠去贏取.
就是以一個自我中心.
轉向為以神中心的人.
雅各到以色列.
我們需不需要由雅各去到以色列.
我們這個名字要改換.
由自我中心轉向以神為中心的生活.
所以雅各最終在第十四節.
在跟神說話的地方.
都是伯達利.
搭了一根柱子.
石柱在上面.
敲酒掣.
又敲油.
這是一個納藥.
重新的納藥.
重新的奉獻.
這個復興的曙光就是建立在一個地方裡面.
當我們能夠回轉到神那裡.
認罪悔改.
接著我們重新立志.
來到神的面前.
這個復興就臨到我們身上.
所以這個是永久的紀念.
當我們埋藏在一些偶像.
神的臨在和我們中間.
神的福氣.
神的恩典是更新了我們.
我們的生命,我們的家庭,我們的教會.

$^{721}$都會迎來復興.
所以在這段經文裡面.
這是一個復興的方向.
我們可以從中體會.
如果弟兄姊妹遇上一個艱難的時候.
好像雅各一家是無路可退的時候.
光景已經是麻木不仁的時候.
我們是不是這個時候要去回轉我們神.
是不是這個時候我們要起來.
重尋我們當日跟神相遇的那種感動.
能夠起來再一次的重新立志.
去建立我們跟神的美好關係.
求主幫助我們.
今天我們教會需要.
我相信由我們個人開始.
就會看到這個復興的來臨.
我一同的祈禱.
主啊,我們的機會不是多.
但我們要把握能夠去你的面前.
因為我們不知道你的懲罰.
或災難是何時臨到.
但我們更加知道.
今天是神你自己給我們的恩典臨到當中.
只要我們能夠去回轉回來的時候.
即使我們過去有很多得罪神的地方.
得罪人的地方.
我們願意去自見.
我們願意去認罪悔改.
神你是公義.
你是聽我們禱告.
你是願意去赦免我們的罪.
以致我們能夠成為一個守潔心清的人.
以致我們無罪.
能夠可以在你的面前站立得住.
更加的要使用我們.
成為了你自己的光明的兒女.
一個明亮的燈台.
在一個地方去照耀一個社區.
在我們的家庭裡面成為一個見證.
在我們的工作崗位裡面成為美好的榜樣.

$^{761}$更加的在我們的生命裡面.
能夠時時刻刻反映主耶穌基督來馨香的味道.
求主幫助我們.
今天再一次藉著經文裡面.
我們的自我審查.
能夠去每天自見.
每天尋求你的面.
願你禱告.
感恩交託.
奉主耶穌名求.
阿們.
\newpage



\section{馬太福音 28:1-7}
\label{sec:032078}
\textbf{When Hope Returns (英/粵)}~~~~趙崇健牧師
\newline
\newline
連結: \href{http://www.church.com.hk/acms/content.asp?site=cdc\&op=show\&type=product\&code=032078\&layout=sermon}{\texttt{032078}} ~~~~ 語音日期: 2026-05-03
\newline
\newline
\hyperref[sec:032077]{\small{< < < PREV SERMON < < <}}
~
\hyperref[sec:index_chronic]{\small{[返順時目]}}
~
\hyperref[sec:032079]{\small{> > > NEXT SERMON > > >}}
\newline
\newline
馬太福音 28:1-7
\newline
\begin{longtable}{cl}
\hline
\hline
章節 & 經文 (和合本修訂版)\\
\hline
28:1 & \begin{tabularx}{0.7\textwidth}{X} 安息日過後，七日的第一日，天快亮的時候，抹大拉的馬利亞和另一個馬利亞來看墳墓。 \end{tabularx} \\ \\ \relax
28:2 & \begin{tabularx}{0.7\textwidth}{X} 忽然，地大震動；因為有主的一個使者從天上下來，把石頭滾開，坐在上面。 \end{tabularx} \\ \\ \relax
28:3 & \begin{tabularx}{0.7\textwidth}{X} 他的相貌如同閃電，衣服潔白如雪。 \end{tabularx} \\ \\ \relax
28:4 & \begin{tabularx}{0.7\textwidth}{X} 看守的人嚇得渾身顫抖，甚至和死人一樣。 \end{tabularx} \\ \\ \relax
28:5 & \begin{tabularx}{0.7\textwidth}{X} 天使回應婦女說：「不要害怕！我知道你們是尋找那釘十字架的耶穌。 \end{tabularx} \\ \\ \relax
28:6 & \begin{tabularx}{0.7\textwidth}{X} 他不在這裡，照他所說的，他已經復活了。你們來！看看安放他的地方。 \end{tabularx} \\ \\ \relax
28:7 & \begin{tabularx}{0.7\textwidth}{X} 快去告訴他的門徒，說他已從死人中復活了，並且要比你們先到加利利去，在那裡你們會看見他。看哪！我已經告訴你們了。」 \end{tabularx} \\ \\
[1ex]
\hline
\hline
\end{longtable}
$^{1}$Let me begin with a question..
我地先從一個問題開始.
What do you think the world needs most right now?.
你認為呢個世界最需要ge 係D 咩呢?.
Some of you might say peace, because we keep hearing about wars and conflict..
有D 人可能會話和平,因為而家我地聽到好多世界有戰爭有衝突ge 消息.
Some may say food and resources, because so many people still go hungry..
有D 人可能會話食物同埋資源,因為呢個世界有好多人捱緊餓.
But let me make it more personal. What do you need most right now?.
Maybe it is something practical, like a better job, a better situation..
Maybe it is something even deeper, like love or a sense of purpose..
We all have our own needs..
大約我都將呢個問題講得更加切身一D ,你而家最需要ge 係D 咩呢?.
可能係一D 實際ge 東西,例如係一份更好ge 工作啦,更好ge 處境啦,又或者係一D .
更加深沉ge ye .
就好似係愛或者生命ge 意義.
我地每個其實我地都有自己ge 需要.
Have you heard of the term "Fur Babies"? Let me share something with you..
大家有冇聽過毛孩呢個字呢?.
我同你分享一件事呀.
I have one fur baby at my place. According to this picture,.
Is a picture from the World Bank data. According to World Bank data,.
In 1963, the global average birth rate.
was 5.3 children per woman..
趙牧師都有一個毛孩ge ,咁佢就想分享下呢個ge 圖表,就係世界銀行ge 數據.
由1963年全球平均ge 生育率係每位女性生5.3個小孩.
By 2023, that number had dropped to just 2.2.
去到2023年呢個數字已經降到2.2個.
And the countries with the lowest birth rate in the world, what are they?.
全球生育率最低ge 地區同埋國家係邊D 呢?.
China,.
中國,香港,澳門,南韓,新加坡,.
Many with birth rate below 1.0.
That means less than one child per woman.
好多地方ge 生育率係低於1.0,即係每個女性生唔到一個小朋友.
Fur babies is what happens when people.
choose not to have children and instead.
raise pets as if they were their own.
毛孩就係一個咁ge 背景底下誕生,係人選擇唔去生育,轉移去到養育自己ge 孩子.
好似養育個寵物一樣.

$^{41}$See, I am not here to say whether that is right or wrong.
我唔係係呢度判斷咁樣做arm 定錯.
My point is this, the trend reveals something deeper.
People have lost hope in the future.
我想講ge 係呢個趨勢揭示jor 一個更加深層ge 問題.
人已經失去jor 對未來ge 盼望.
People nowadays said, even if a war ends today, another war might start tomorrow.
即使今日一場ge 戰爭結束,明日另一場ge 戰爭可能又再開始.
Even if the economy recovers, it may not last.
即使經濟復甦都好啦,未必持久ge .
People are asking themselves, is there any reason to build a future at all?.
人個個都係度問自己,仲有冇任何理由去到建造一個未來呢?.
So let me ask you honestly.
所以我直接ge 去到問你.
Do you have hope for this world, for yourself, for your own future?.
你對呢個世界,對你ge 未來,對你ge 人生仲有冇盼望呢?.
I believe what the world lacks most today is not resources, it is hope.
我深信呢個世界今日最缺乏ge 其實唔係資源,而係盼望.
And perhaps what you lack today is the same.
或許你今日所缺乏ge 都正正係呢份盼望.
Let's bow and pray.
我地一齊禱告.
Lord, we come today carrying different burdens.
Disappointment and uncertainty.
主啊,我地今日帶住各自ge 重擔黎到你跟前,可能我地有失望有迷惘.
And Father, I pray that you give us hope.
我求主你比我地有盼望.
We ask that you meet us here.
求主你與我地相遇.
Speak to us, open our heart, and remind us that Father, you are alive and you go before us.
向我地說話,闖開我地ge 心.
提醒我地你係活著ge ,你行係我地ge 前面.
In Jesus name, Amen.
禱告,奉主名求,Amen.
See, today's passage is set in a very narrow window of time.
The days after Jesus died and before anyone knew he would come back to life.
英文係發生係一個好短暫ge 時間裡面,係耶穌死jor 之後幾日,gwo 陣仲未有人知道.
佢會復活.
The main characters are women who had followed Jesus, and in just a few days, one scene after another.
They had let them down, a row of great disappointment.

$^{81}$主角係一班跟隨耶穌ge 婦女,短短幾日內呢佢地一幕接一幕ge 將佢地帶入深深ge .
失望之路.
First of all, they watched Jesus died.
Scripture says they stood at a distance and saw everything.
首先呢佢地眼白白ge 望住耶穌死去,聖經話佢地企係遠處看見了一切.
See, watching any person take their last breath is deeply shaking.
Whether it is someone you love or even just a stranger.
This woman saw it all.
望住任何一個人死,其實都好深深震撼我地,無論gwo 個係你所愛ge 定係一個陌生人.
但係婦女係全都親眼目睹.
And that same evening, they followed and watched, as his body was laid in the tomb.
當日旁晚,佢地跟隨住耶穌ge 屍體,亦都目睹jor 耶穌ge 身體被安放係墳墓裡.
They noticed that the burial had been done without the traditional anointing of spices.
佢地注意到安葬ge 時候,並無按照傳統用香料去到高門呢個遺體.
So they went home and prepared the spices themselves, a final act of love.
That suddenly feels so close to them.
於是佢地返屋企自己準備jor 香料,呢個係佢地最後一個愛ge 行動.
死亡突然變得如此ge 真實,如此ge 靠近.
Then early the next morning, before dawn, because they are so desperate.
They really want to come to the tomb and find the stone roll away.
第二日一早,佢地天未亮ge 時候,佢地就好迫切黎到墓穴,發現石頭已經滾開jor .
The body was gone, someone had taken him, they said, and they didn't know where he is.
個遺體唔見jor ,有人將佢帶走jor ,而佢地唔知道去jor 邊度.
There is such helplessness in those works.
呢個話語裡面藏住好大ge 無力感.
They had come to offer one last act of love and couldn't even do that.
佢地本身係黎獻上最後一個愛ge 行動.
但係連呢樣ye 都做唔到.
They couldn't say goodbye, this moment they felt utterly helpless, disappointed, maybe even hopeless.
佢地感到完全ge 無助,失望,甚至可能係絕望.
Mary stood outside the tomb and wept, so overcome.
He bent over.
瑪利亞.
企係墓外面去到哭泣,悲痛到佢彎低佢ge 腰.
向裡面去到張望.
好重要我地去理解佢呢個動作.
She is like this.
佢會係咁樣.
佢係好似趙牧師咁樣.
What I mean is, she was so overcome, she couldn't even stand straight.

$^{121}$佢悲痛到彎低身落黎,企唔直gwo 日.
Looking into the tomb.
佢向個墓穴裡面望.
And that image is so clear.
呢幅圖畫好清晰, Someone who still cares, still tries, still shows up.
But has nothing left to do.
係一個仍然好在乎,仍然好嘗試,仍然出現,但係已經一無所有.
Bent over with grief.
被悲傷壓彎jor 佢ge 腰.
So I want to ask, when did you last cry?.
好想問你上一次喊係幾時呢?.
When did you last feel that kind of disappointment or heartbreak?.
你上一次感受到gwo 種ge 失望心碎係幾時呢?.
Maybe you didn't literally cry, but inside you felt it.
可能你冇真係喊出黎,但係心裡面你係確實感受到.
I also want to be honest with you.
好想同你坦誠ge 分享.
My younger son is 10 years old.
我ge 細仔今年10歲.
And he was born, he had several health challenges.
佢出世ge 時候有幾個健康方面ge 問題.
One of his eyes had significant farsightedness.
其中一隻眼佢有嚴重ge 軟視.
The doctor told us that if it didn't improve, he might never be able to drive.
醫生話比我地聽如果冇改善呢佢可能永遠都唔可以zar 車.
And in a place like Australia, not being able to drive makes life so complicated.
係澳洲呢個地方唔識zar 車呢會令到生活非常ge 困難.
He also had a curved spine.
佢仲有脊柱側彎.
The doctor warned that if it worsened.
He would need to wear a rigid.
Back brace for 20 hours a day.
如果醫生警告話如果惡化ge 話呢佢.
要每日去到穿戴一個硬式ge 背架20個鐘.
Even while sleeping.
訓覺都唔例外.
So my wife Winnie had a friend who is also a pastor's wife.
趙牧師ge 太太Winnie有一個朋友都係師母黎ge .
We met her family in 2014.
佢地2014認識.

$^{161}$And at that time she had just had her first child.
gwo 個時候佢ge 朋友arm arm 生jor 第一個小朋友.
Today in 2026.
而家係2026年.
She has 10 children.
佢有10個小朋友.
I saw their family photo on Facebook.
我見到係Facebook見到佢地ge 全家福.
When they go out, they need a minivan.
佢地出去旅行ge 時候要租架van仔.
I saw how they are going to have their own dinner.
見到佢地食飯要點樣食.
They have so many pots of pasta on the desk.
佢地煮意粉要煮好多盤.
See I admire her.
好欣賞佢.
Not because she had 10 children.
唔係因為佢呀10個小朋友.
But because all 10 are healthy.
因為10個小朋友都健健康康.
See I know what it is to sit in a doctor's office.
With that fear.
我好清楚知道坐係真係診室裡面gwo 個心懷恐懼ge 感覺.
Every appointment.
Every diagnosis carry weight.
每次ge 預約每次ge 診斷都有佢ge 重量.
So when we look at this woman's experience over these few days.
所以當我地望見呢幾位ge 婦女短短幾日裡面所經歷ge .
Don't we see something similar.
我地會唔會見到一D 我地都覺得好熟悉ge ye 呢.
Losing someone you love.
失去你所愛ge 人.
You are trying your best.
你盡力而為啦.
Things always go unexpectedly wrong.
但係事情總係出乎意料ge 出差錯.
And finally you are still standing.
最後你仲企係度ge .
You are still trying.
你仲係努力緊ge .

$^{201}$But you are barely holding it all together.
但係你係好勉強ge 撐住自己.
There is a pattern in the text worth noticing.
呢段經文裡面有一個好值得留意ge 規律.
The word soul keeps appearing.
看見呢個字不斷ge 出現.
They saw Jesus died.
They saw how he was buried.
They saw the tomb.
They saw the stone had been moved.
They saw the body was gone.
佢地見到耶穌死去見到佢地點樣被安葬呀見到呢個墓穴.
見到石頭被移開佢見到遺體唔見jor .
But they could not see the future they could not see hope.
但係佢地見唔到ge 係未來.
見唔到ge 係盼望.
Disappointment.
失望.
Is.
Not being able.
To see what you hope to see.
係見唔到你好盼望見到ge .
Despair.
Is not being able to see.
Anything at all.
絕望就係乜ye 都見唔到.
And only God.
Can open our eyes.
To see hope.
唯有上帝可以開啟我地ge 眼讓我地見到盼望.
See we all know that God is love.
我地都好明白上帝係愛ge .
And God does not leave.
The people he loves stuck in disappointment and despair without doing something about it.
但上帝都唔會讓佢所愛ge 人陷係失望絕望裡面去到袖手旁觀.
But the trouble is we so often place our hope in things of this world.
但問題在於其實我地係時常將盼望寄託係呢個世界ge 事物上.
Thinking they will secure a better future.
但我地以為呢D 事物可以保障一個更好ge 明天.
But real hope can only be found in Jesus who is leaving and return to us.

$^{241}$但係真正ge 盼望只能夠係耶穌ge 身上搵到.
佢活著ge 並且係我地中間.
Look at what happens.
我地睇下跟住發生咩事.
The woman stand at the tomb weeping helpless and then.
The angel speaks.
婦女們企係ge 墓旁.
哭泣無助.
就係呢個時候一位天使開口講ye .
See I want to highlight a few words in this passage.
我想係呢段經文裡面標註幾個ge 詞語.
If you have had the text in front of you, if you have a Bible, you want to highlight this as well.
如果你手上有經文聖經ge 不妨都將佢地畫出黎.
The first word is afraid.
第一個字係懼怕.
The angel says do not be afraid.
天使話唔好懼怕.
Why are the women afraid?.
佢地點解驚呢?.
Jesus coming back to life brings.
Peace, shouldn't them be overjoyed?.
係耶穌復活帶黎平安架嘛,佢地起唔係應該欣喜若狂咩?.
If.
Out of all the sudden, an angel suddenly appeared to anyone.
當然啦,突然出現D 天使可能會嚇親人ge .
But at a deeper level, they were afraid because they never expected God to work in this way.
但更深ge 層面,佢地驚係因為佢地從未預料到上帝會用呢一種ge 方式.
去到行事.
God's hope.
Doesn't always arrive in the way we expect.
上帝ge 盼望.
並唔總係以我地所預期ge 方式去到降臨.
Sometimes it.
Surpasses what we imagined.
有時佢遠超我地所能夠想像.
That kind of surprise can be unsettling.
But it is always a good kind.
呢種ge 驚喜可能係令人措手不及.
但佢總係令到呢一種ge .
呢種ge 美好.

$^{281}$係美好的措手不及.
Second word, the angel said.
I know.
第二個字係我知道.
I know you are looking for Jesus.
天使話我知道你地係想搵耶穌呀.
When we are afraid, when we are lost.
When we don't know what to do, God said.
I know.
當我地懼怕,當我地迷失,當我地不知所措ge 時候.
上帝話我.
知道.
I know what you have been searching.
我知道你一路都搵緊D 乜.
I know you couldn't find what you were looking for.
但我知道你搵唔到你想搵ge .
I know your feelings.
我知道你ge 感受.
Even the ones you can't put into work.
即使你講唔出ge 感受我都知.
So you don't need to be afraid.
所以你唔使懼怕.
What a comfort because God knows.
呢個係何等ge 安慰.
The third phrase.
第三個詞.
He is going ahead of you.
他在你門以先往那裡去.
Before.
The disciples.
Knew Jesus was alive.
Before.
The woman.
Found the tomb empty.
係門徒知道耶穌復活之前.
係婦女們搵到呢個空墳墓之前.
God had already.
Acted.
He was already ahead.
Of them.

$^{321}$He goes.
Before us.
原來上帝已經行動.
佢已經走係呢班婦女ge 前面.
神係走係我地ge 前面.
The fourth word.
Joy.
第四個詞.
大喜樂.
The woman ran from the tomb with fear and great joy.
婦女們帶住恐懼同埋大大ge 喜樂.
Just a moment ago.
They were weeping.
Not knowing which way.
To turn, now they were running with great joy.
前一刻佢地先係係度哭泣,唔知去邊好,但而家佢地就係帶住大大ge 喜樂去到奔跑.
See, hope doesn't.
Doesn't make you say, oh well, maybe.
Perhaps things.
Will work out.
Probably it may work.
I'm not so sure.
Hope isn't a wishful thinking.
Hope brings certainty, hope doesn't make you joyful because you might receive something.
Hope makes you joyful because you have already received it.
盼望帶黎確信.
盼望使你喜樂唔係因為你可能得到D 咩,而係因為你已經得到D 乜ye .
Hope makes you joyful because you have already received it.
盼望帶黎確信,盼望使你喜樂唔係因為你可能得到D 咩,而係因為你已經.
得到jor .
Weeping is the expression.
Of disappointment, joy is the expression of hope.
哭泣係失望ge 表達.
而喜樂就係盼望ge 表達.
When did you last feel that kind of joy.
你上次感受到呢種ge 喜樂係幾時呢.
I want to share a miracle that happened recently.
我想分享一個我最近發生ge 神蹟.
A month ago.
At a CNMA pastor's gathering.

$^{361}$I had to leave early to take my son to see the eye doctor.
係一個月前係宣道會一次ge 牧者ge 聚會上面呢我需要提前離開帶我個仔去睇眼科.
And before I left, someone suggested the pastor pray for him.
係我離開之前呢有人就建議牧者一班ge 牧者為我個仔去到禱告.
Since he wasn't there, I became his stand-in convoy.
Someone.
Lay hands on me and pray for him.
因為我個仔唔在場啦我就成為jor 佢ge 替代品.
大家都將手按係我ge 身上面但係為佢祈禱.
I am usually the one laying hands and pray for others.
我通常係gwo 個按手為其他人禱告ge 人.
I can't remember the last time anyone lay hands on me.
Apart from my ordination.
我已經唔係好記得上一次有人為我按手禱告係幾時呢除jor 我個案立典禮.
所以大家為佢祈禱多D .
所以有人為我祈禱ge 時候我好感動.
My wife took my son to the doctor that afternoon.
當日下午我ge 太太帶仔去睇醫生.
Afterward she told me.
之後佢就同我講.
The doctor said a lot of medical terms I didn't fully understand.
醫生講jor 我好多唔明ge 醫學術語.
But the conclusion was.
最重要ge 結論係咩呢.
He can try going without glasses entirely.
佢可以嘗試完全唔戴眼鏡.
Treat him like a normal person.
就好似對jor 一個正常人一樣.
See, even as a pastor, I feel helpless sometimes.
即使身為牧師其實我有時都感到無助.
Every doctor's wishes carries fear.
每次看診都帶住恐懼.
But God already knew.
但上帝早已經知道.
And he had already gone ahead to prepare the way.
佢已經係前面預備好道路.
That is why we can place our hope.
In Jesus.
And find joy.
In that hope.

$^{401}$呢個就係點解我地可以將盼望寄託係耶穌身上.
並且係個盼望裡面.
搵到喜樂.
Hope.
Is never.
Something you can hold in your hand.
Or see with your own eyes.
盼望唔係我地能夠握係手裡面或者用眼睛真係睇到ge .
And it is not.
Found in anything in this world that has no life in it.
佢唔存在於呢個世界上面任何冇生命ge 事物當中.
And because Jesus is living.
Hope is not just a feeling.
It is a fact.
It came back to life.
因為耶穌活著盼望唔係一種ge 感覺.
佢係一個ge 事實.
佢係從死裡面復活jor .
He has overcome death.
And nothing will overcome it.
Again.
And Jesus is alive forever.
佢已經勝過jor 死亡本身.
冇任何ye 能夠再勝過佢.
耶穌係永遠活著.
No matter how disappointed you are today no matter how hopeless this.
Things you feel.
You are not left with nothing.
無論今日你有幾咁ge 失望.
無論事情睇起黎有幾咁ge 絕望你唔係一無所有.
Remember this.
Jesus Christ is alive.
You have him and you have hope.
請記住一樣ye 耶穌基督係活著ge .
你有佢.
你就有盼望.
And the last section of our passage reveals the heart of Jesus toward people.
最後呢段經文ge 一部分揭示jor 耶穌對人ge 心意.
The text says.
Statini.

$^{441}$Jesus.
Meet them.
經文話忽然耶穌遇見他們.
It reads like a coincidence.
讀起上黎好似係一個即係偶然咁樣.
But it wasn't.
但係唔係偶然.
Statini is the woman's experience.
They didn't see it coming.
忽然係講緊婦女們ge 感受呀佢地估唔到呀.
But for Jesus.
There was nothing sudden about it.
He had intentionally gone to meet them.
但係對耶穌黎講根本唔係乜ye ge 忽然.
佢係有意去到搵佢地ge .
Think about this when did you last.
Visit someone.
Without telling that person.
諗下你上一次特登拜訪某一個人而對方唔知情ge 時候係幾時呢.
And Jesus did it the one who created the heaven and earth.
Who holds the universe together.
And he deliberately went to meet a small group of.
Grieving women.
但耶穌就係咁樣做gwo 位創天造地托住宇宙ge 主.
佢特登去一班傷心ge 婦女搵佢地去到相遇.
And that is love he does it for you.
呢個就係愛.
佢都同樣去到搵我地.
Then the text said Jesus spoke to them.
然後經文話耶穌對佢地說話.
We read God said and Jesus said so often in the Bible.
但我地時常係經文聖經裡面讀到下上帝話或者耶穌話.
We always want to skip them.
Sounds like nothing.
Particularly special.
我地讀聖經ge 時候讀到係好容易輕輕帶過.
But we better to pause here.
但係我地停一停諗一諗.
If Jesus personally.
Speak to you right now would you still take that for granted.

$^{481}$如果耶穌親自而家對你說話你仲會唔會當呢樣ye 係理所當然.
See Jesus could have sent an other angel.
Or speak to them in a dream.
耶穌可以差一位天使或者夢裡面同佢地講呀.
But he chose to come in person and speak directly and to be pleasant.
但係佢選擇親身ge 黎到直接ge 講ye 親身ge 同在.
This is how much.
He values us.
呢個就係佢幾咁ge 看中你.
This is the depth of his love.
呢個就係佢愛ge 深度.
And what is his very first word.
而耶穌講ge 第一個字係咩字呀.
It is an original Greek.
呢個係一個希臘ge 原文字.
Can we say it together.
我地一齊讀好嗎.
Do you know what is Karate.
大家知唔知咩叫Karate呀.
You know Karate right? It sounds like Karate.
聽落去好似跆拳道呀.
It is Karate.
係Karate呀.
So if you remember Karate you remember Karate sounds like the first word Jesus speak to this woman.
所以你記得空手道就記得呢個Karate.
所以希望大家都記得呢個Karate.
It was a common greeting of that time like hello or how are you.
呢個就係昔日ge 時候常用ge 問候語好似你好呀你點呀咁樣.
May you be well, may you be happy, may you be blessed.
大意思係願你平安願你喜樂願你蒙福.
It is what you say when you genuinely want someone to be okay.
呢個係你真心希望一個人一切安好ge 時候所講ge 一句說話.
The very first word.
The very first thing Jesus said.
To this woman.
After returning was.
耶穌復活之後對婦女們講ge 第一句ge 說話就係.
I want you to be well.
I want you to be happy.
I want you to be blessed.

$^{521}$我希望你平安.
我希望你喜樂.
我希望你蒙福.
呢個就係耶穌對你對我ge 心意.
當盼望返黎ge 時候而佢一定確實會返黎佢唔係淨淨你返黎.
佢係完全ge 帶住生命ge 返黎就係gwo 個活著ge 耶穌基督ge 身上返黎.
And if you need hope today I invite you to turn forward.
Turn to Jesus.
無論你今日係咪已經認識jor 耶穌或者摸索緊佢係邊個都邀請你轉向耶穌.
Not a concept, not a religion, but a living person who goes ahead of you.
唔係一個概念唔係一個宗教而係一個活生生ge 人.
He knows your pain.
佢知道你面對ge 痛苦.
And he wants to meet you where you are right now.
佢想係你所在ge 地方同你去到相遇.
有一首詩歌講得好靚ge .
因他活著我能面對明天.
因他活著不再懼怕.
我心知道他掌管明天.
生命充滿了盼望只因他活著.
當盼望返黎ge 時候.
It has a name.
佢有個名架.
His name is Jesus Christ.
佢個名係耶穌基督.
Let's bow and pray.
我地一齊低頭.
Jesus thank you for reminding us that hope is not a wishful thinking.
It is a living certainty.
耶穌多謝你提醒我地.
盼望原來唔係一廂情願ge 諗法.
而係一個活生生ge 確信.
Because you are living, as we return to our daily lives, go before us.
因為你係活著.
當我地返返去我地日常ge 生活裡面求你走在我地ge 前面.
In that moment we feel afraid or helpless, help us remember that you already know.
And you are already there.
係我地感到恐懼無助ge 時刻求你幫助我地記得.
你早已經知道並且你已經係gwo 度.
May the joy and hope.

$^{561}$of Jesus Christ carry us through.
In Jesus name.
願耶穌基督ge 喜樂盼望.
保持我們渡過一切.
我們禱告奉主名求.
Amen.
\newpage



\section{哥林多前書 12:27}
\label{sec:032079}
\textbf{同心合意,建造教會 (華/英)}~~~~劉彤牧師
\newline
\newline
連結: \href{http://www.church.com.hk/acms/content.asp?site=cdc\&op=show\&type=product\&code=032079\&layout=sermon}{\texttt{032079}} ~~~~ 語音日期: 2026-05-03
\newline
\newline
\hyperref[sec:032078]{\small{< < < PREV SERMON < < <}}
~
\hyperref[sec:index_chronic]{\small{[返順時目]}}
~
\hyperref[sec:032080]{\small{> > > NEXT SERMON > > >}}
\newline
\newline
哥林多前書 12:27
\newline
\begin{longtable}{cl}
\hline
\hline
章節 & 經文 (和合本修訂版)\\
\hline
12:27 & \begin{tabularx}{0.7\textwidth}{X} 你們是基督的身體，並且各自都是肢體。 \end{tabularx} \\ \\
[1ex]
\hline
\hline
\end{longtable}
$^{1}$讀一節聖經,在格列多前書十二章二十七節..
保羅說,你們就是基督的生死,並且各自做肢體..
今天是傳承主日..
什麼是傳承主日呢?.
傳承主日就是一個屬神的大家庭,一同來紀念神工作的日子..
我們一方面來追溯神的恩典,也來數算神這些年在我們中間的作為..
特別在過去三十年,我們經歷過多少神的神蹟奇事..
我們也經歷神多少偉大奇妙的作為在我們弟兄姊妹身上..
另一方面也需要傳承,讓我們下一代能夠明白神的法則,明白神所給我們的使命跟託付..
我們的世世代代一同來建造神榮耀的教會..
弟兄姊妹,建造教會不是只有外表的建築而已..
更重要的是裡面的建造,也不只是個別生命的建造..
更重要的是生命要跟生命連結在一起,就好像你蓋一個房子一樣..
最重要的倒不是這塊磚塊有多好,就是再好的磚塊,如果不能夠連在一起,那個房子還是蓋不起來..
同樣,屬靈的建造也是這樣..
我們不能只看重自己個人生命的建造,我們卻忽略了生命跟生命的連結..
保羅在這裡教導我們說,教會好像是基督的身體,我們就是這個身體裡面一個小小的肢體..
這個肢體每一個都要配搭得好,整個身體才能發揮功能..
同樣,今天在教會裡面,我們也需要看重肢體的連結跟團隊的事奉..
所以今天早上從三方面,我們特別要來思想在教會裡面的團隊事奉..
首先第一個我們看見團隊事奉是神的心意..
在聖經裡面我們看見我們所信的神是一位三位一體的真神..
三位一體是什麼意思?.
意思是我們所信的只有獨一的一位神,但他卻有不同的位格..
是聖父,聖子,聖靈..
什麼叫位格呢?.
位格就表明他們三個有獨立的思想,意志,情感..
所以今天我們所信的是一位獨一的神,但在神裡面有三個不同的位格..
他們三位是同等的,同尊同榮,但他們卻有不同的功能..
所以我們看見我們所信的神,祂的本身就是一個團隊..
所以今天神對教會的心意也是如此..
神自己自然是一個團隊,神在設立教會的時候也要我們學習團隊的事奉..
所以保羅在經文裡面告訴我們,我們好像一個身體一樣..
我們每一個肢體都有不同的功能,但只有當我們配搭,連結在一起,整個身體才能夠發揮它的功能..
換句話說,我們沒有一個人可以獨立作業..
就算你說你十八般武藝樣樣都會,你還是需要你的弟兄姊妹..
跟你媽說,你很需要你的弟兄姊妹,也跟她說..
就像聖經裡面的保羅,你還記得當保羅在往大馬塞路上走的時候,他遇見了主..
大光照向他,他撲倒在地.在那一剎那,他知道用光光照他,是他生命的主宰..
但他不知道他是誰.保羅在那個時刻就問了他一生中最重要的兩個問題..

$^{41}$第一個問題,主啊,你是誰?第二個問題,主啊,我當作什麼?請問你主耶穌回答他的問題沒有?.
其實我們仔細看聖經,耶穌只回答了一半.第一個問題耶穌回答了,耶穌說我就是你所逼迫的耶穌..
第二個問題耶穌沒有直接回答,耶穌沒有直接回答他,耶穌只對他說,你進城裡,有人會告訴你..
這麼重要的一個問題,主啊,我當作什麼?為什麼主耶穌沒有直接告訴他?.
而是說有人會告訴你.而且你要知道,當他在那裡被大光光照的時候,聖經說保羅的眼睛就有個鱗片擋住他..
聖經裡面告訴我們說,他其實那三天,他的鱗片一直都沒有掉,好像個瞎眼的人..
我相信在那三天當中,保羅不斷地呼求主啊,讓這個鱗片掉下來,掉下來..
可是好像神也沒有垂聽他的禱告,一直到神差遣亞拿尼亞來到他的面前..
這個亞拿尼亞是誰呢?其實是一個不大出名的弟兄..
我常想如果神要向保羅顯現,要向他說什麼話,起碼一個像彼得這樣分量的人來跟他講..
可是沒有神差遣亞拿尼亞,亞拿尼亞在聖經只出現這一次而已..
亞拿尼亞不是很出名的一個人,神就藉著亞拿尼亞到保羅面前,當亞拿尼亞按手在保羅身上的時候,他的鱗片就掉下來..
然後亞拿尼亞對他說,神的心意是要他做外邦的使者..
我常常在想,為什麼當保羅來尋求神的心意的時候,神不直接回答?.
當保羅為他自己求醫治的時候,神不直接醫治他?.
為什麼神要藉著這樣一個小弟兄來說明這件事?.
我告訴你為什麼.如果說每一次保羅來尋求神,神都直接回答他,每次保羅求醫治,神就醫治他..
那保羅可能以為他一輩子不需要別的弟兄姊妹了..
可是在這個時刻,在他事奉開始,神就很清楚地教導他說,嘿,你需要你的弟兄姊妹..
也許如果說神從一開始就決定他禱告的話,他可能一輩子只要靠他自己就可以了..
但現在他謙卑下來,他學習與人同工..
如果你仔細看保羅的事奉,保羅是一個非常懂得與人同工的人..
你看保羅書信,保羅書信有一個特點,就是裡面充滿了名字..
平常我們讀保羅書信都趕快過去,趕快掃過去,哎呀,怎麼那麼多名字,過去,過去..
保羅為什麼寫那麼多名字啊?不只是隨便寫寫,你發現保羅跟這些人的關係都蠻深厚的..
你發現保羅是個非常懂得與人同工的人,而且無論他到哪裡,都帶著同工一起去事奉..
甚至下到監獄,也帶著同工到監獄去..
所以你知道做保羅的同工不容易..
他下監獄,你也跟他下監獄..
你就看見保羅從第一天開始事奉,他就學到非常重要的一個功課..
他需要他的弟兄姊妹,他需要一個團隊..
跟你旁邊的人在說,你需要你的弟兄姊妹,你跟他說..
所以團隊事奉是神的心意..
第二,團隊事奉是一種教會生活的文化..
今天在教會裡面,我們若想發揮團隊事奉的功能,我們就必須先建立一個團隊事奉的文化..
意思就是,無論我們做什麼,我們都不能只想到我們自己,我們需要想到別人..
那什麼是團隊事奉的文化呢?.
第一個是互相尊重的文化..
是一個互相尊重的文化..
在教會讓你們事奉的時候,我們需要彼此的尊重..

$^{81}$為什麼我們要學習彼此尊重呢?.
有兩個原因..
第一個,因為他是按照神的形象,神所造的..
他有神所給他的自由意志,以及神所賦予他作為一個人的價值..
第二個,因為耶穌愛他..
耶穌的寶血也是為他而流..
所以我告訴你,無論我們是否同意他的想法跟決定,.
我們都要尊重他..
也要尊重神在他個人身上的帶領..
第二,你真的要知道一件事情,我們要明白神在每一個人身上的工作都是不一樣的..
我們千萬不要用我們的精力來勉強別人來跟隨..
因為神在每一個人身上的帶領是不同的..
無論我們明不明白,我告訴你,神的道路還是高過我們的道路..
所以我常常說,千萬不要太早蓋棺論定..
有時候我們看到一些人,這個人就是這個樣子..
你有沒有講過這句話?.
那個人就是那個樣子..
意思就是,好像連神都沒辦法..
好像連神都沒辦法..
千萬不要太早對別人失望..
因為連神都還沒有失望,你失望有什麼價值?.
神永遠有他自己的時候跟方法來成就他自己的工作..
我們就是要來學習來尊重神的工作..
你知道這些年當我在幕會當中,我學習到一個非常重要的功課..
就是絕對不要代替神來工作..
我們永遠只是站在祭司的地位為別人代求..
不是去責備他們,而是為他們求恩典..
因為只有在神的恩典之中,人才能夠重新找回自己的意義跟價值..
所以第一個很重要的是互相尊重..
第二個是互相接納的文化..
今天在教會裡面我們也需要學習彼此的接納..
你說為什麼需要學習彼此接納?.
最重要的原因是主耶穌也這樣接納了我們..
事實上耶穌不但接納了你,耶穌也接納了他..
所以在教會裡面我們不要太自以為是..
有一次有一個傳道人來到神面前禱告..
他看見另外一個傳道人有很多軟弱的地方..
但好像什麼神還在使用他呢?.
他心中就很不滿..
他就對神說,神啊,那樣子的人你怎麼能用啊?.

$^{121}$你怎麼可以使用他?.
神怎麼回答呢?.
神對他說,就像我怎麼使用你一樣啊!.
為什麼使用你呢?.
一樣啊!.
很多時候我們自以為是,覺得我們比別人好..
其實神使用我們不是因為我們好啊!.
我們都不配!.
是因為神接納了我們!.
神既然如此接納我們,我們同樣要接納別人..
因為主也接納了他..
我們就是要在主耶穌裡面彼此接納..
第三個,我們需要互相祝福的文化..
在教會裡面我們要學習彼此祝福啊!.
千萬不要看見神祝福別人就心裡嫉妒啊!.
特別是當你到人家家做客的時候,.
當你看見別人的廚房比你家客廳還大!.
你不必酸溜溜,哼!有什麼了不起!.
人活著又不是單靠食物!.
那嫉妒會使我們的心變得苦澀!.
以至於我們自己得不到神的祝福了!.
想想看當年約瑟的兄弟是因著嫉妒把他賣到埃及啊!.
事實上如果約瑟的兄弟如果早知道,其實神在他們身上也有奇妙的計畫啊!.
如果立威早知道他的子孫是唯一能夠在聖殿事奉的人,.
如果猶大早知道連救主耶穌基督都從他的子孫出來,.
如果他們早知道,可能整個故事就不會這樣發展下去啊!.
弟兄姐當你看見神祝福別人的時候千萬不要嫉妒啊!.
反而你要需要相信有一個屬於你的時候會來到的啊!.
跟你媽說你的時候會來到的啊!.
你要學習去祝福別人啊!.
你越多祝福別人,你越多神就會祝福你啊!.
所以我們看見團隊事奉是一種在教會裡面的生活的文化..
第三,團隊事奉也是一種生命的操練..
那你說我們怎麼樣在教會裡面這樣操練團隊事奉?.
在腓立比書第二章第二到第五節裡面保羅給我們四個很重要的原則..
第一個我們要建立共識..
腓立比書二章第二節這裡說你們要意念相同,愛心相同,有一樣的心思,有一樣的意念..
這個團隊若是想有效的服侍神,首先第一個就要有同樣的心思意念..
換句話說就要建立一個共識..
那你說怎麼樣建立這個共識呢?.

$^{161}$就是要透過不斷的禱告跟溝通..
禱告能夠使我們的心謙卑下來..
否則我們常常堅持己見..
禱告讓我們謙卑來到神的面前..
禱告也能夠讓聖靈把神的心意放在我們的裡面..
好讓我們能夠在神的心意裡面達到共識..
另外溝通也很重要..
溝通能夠幫助我們了解別人的想法跟立場..
其實你知道這些年我在教會發現,大家在教會服侍,大家都非常有熱心..
而且很樂意擺上自己..
但很多時候就是因為缺乏溝通,帶來不必要的誤解跟傷害..
丁鳳姐,既然是團隊事奉,就表明不是只有一個人的意見..
我們就必須謙卑下來學習聆聽..
不但聆聽神的聲音,也聆聽彼此的聲音..
而且我們必須存著寬容的心..
因為有些時候意見的成熟是需要花時間的..
每次當我們遇見衝突跟困難的時候,.
我們需要重新回到神的面前..
我們求神的聖靈來引導我們..
重新在我們生命當中來掌權..
直到我們在耶穌基督裡面有同樣的心思意念為止..
第二個,我們要互相保護..
我們要互相配搭..
腓立比書二章三節說,只要存心謙卑,個人看別人比自己強..
丁鳳姐,今天當我們在教會事奉,我們一定要明白一件事情..
我們是彼此互補的..
以前我還有機會做婚前輔導的事..
總是會提醒這些夫妻們,夫妻互補的觀念..
因為我們的神真的很幽默..
就像中國話所說的,急驚風一定碰到慢狼蟲..
你看夫妻就是這樣,一個愛講話,一個怎麼擠都擠不出一句話來..
這個就是會在一起的,很奇怪,就會在一起..
一個愛花錢,一個愛省錢,就是會在一起的..
就是互補嘛,你要明白互補的觀念..
不但在婚姻上是如此,在教會生活也是如此..
我們不要只去看別人的缺點,.
卻沒有看別人的缺點是需要我們去補足的..
所以我曾經讀到一句話,很觸動我的心..
他說要在別人的缺點上看見自己的責任..
我從來沒這樣想過,要在別人的缺點上能夠看見自己的責任..

$^{201}$所以你就不要隨便批評別人..
其實當你看見別人的缺點的時候,是提醒你,我們有這個責任去幫補別人..
另一方面我們需要謙卑下來,能夠看別人比自己強..
意思就是欣賞別人的優點吧..
有一首詩歌我很喜歡的,裡面講到說我在你的裡面看見君王神的容美..
哇,太好了,這句話..
換句話說,你看見別人的優點,看見別人的真的是被神所使用的,.
然後真正看見,透過他神的容美,能夠彰顯..
第二點,今天神把我們放在這裡,就是我們彼此配搭..
看見別人的優點要欣賞..
看見別人的缺點,是有責任去補足..
第三,要彼此相顧..
菲律賓是二章四節,這裡說不要單顧自己的事,要顧別人的事..
現在團隊事奉裡面最常見的一個錯誤,就是我們常常只顧自己的事情,只顧自己的部門..
我們常常覺得只要我這裡做好就好了,別人如何我不管,與我無關..
所以我們在教會裡面常常分門別類,甚至彼此責怪..
這不是我的錯,而是別的部門的事,不是我這裡出問題,是別的部門出問題..
我們常常以為只要我這裡顧好就好了,但是團隊事奉的時候,你發現只要一個地方有破口,仇敵就會進來攻擊..
不管這個破口是在你這裡,還在別人那裡,都一樣..
只要任何地方有破口,仇敵就會進來攻擊..
所以保羅才在這裡說,嘿,不要單顧自己的事..
想想看,如果主耶穌在天上只顧祂自己的事,祂就不會到這個世界上來完成救恩..
今天一個事奉神的人,我們要擴大我們生命的格局,不但顧念自己,也要顧念別人的事情..
因為我們是同一個身體..
既然是同一個身體,就沒有說誰成功誰失敗..
如果是同一個身體,那就一起失敗一起成功..
沒有說誰成功誰失敗..
一起成功一起失敗..
跟你朋友說一起成功一起失敗,你跟他說啊..
最後第四在於主理要合一..
基督徒書20章第五節說你們當以耶穌基督的心為心..
團隊事奉最重要的一件事情是在主理我們要合一..
今天無論我們前面的道路是如何,.
無論這個旁邊的風浪有多大,.
我們一定要決定在主耶穌的裡面一起禱告,一起掙扎..
一起成長一起得勝..
今天在神的大家庭裡我們不但彼此成為團隊,.
更重要與主成為一個團隊..
這是支屬神的大軍隊..
耶穌基督是軍隊的元帥..

$^{241}$我們尊他為主,同心合一向前進..
無論我們仇敵的攻擊有多厲害,.
無論我們的環境多麼艱難,.
主必與我們同在,.
帶領我們進入那至終的得勝..
事實上今年七月我們有一個非常精彩的活動要推動..
這是一個我們過去從來沒有做過的一件事情..
我們把這個行動叫做耶利哥行動..
什麼叫做耶利哥行動?.
因為當時你還記得神教導以色列人攻打耶利哥城的時候,神給他們那個策略..
其實是神的百姓團隊事奉最好的呈現..
你還記得神要求以色列百姓圍繞耶利哥城嗎?.
一天繞一圈..
第七天七圈..
最後大聲呼喊城牆就倒..
我們也盼望今年七月份推動這個耶利哥行動來為主贏得我們所居住的城市..
你說怎麼做法呢?.
我們將會從七月五號到七月十一號連續七個晚上舉行七天的敬拜禱告之夜..
每個晚上有不同墓區來帶領..
一方面我們要為每一個墓區所在的社區跟城市祝福禱告..
另一方面我們要輪流為影響世界的七個領域的其中的領域來祝福禱告..
最後一個晚上也是禮拜六的晚上,我們會聯合中英文青少年的所有人在一起我們一個復興之夜的聯合聚會..
我們要一同為我們的教會,城市,國家來求復興..
然後七月十二號主日的下午..
我們會發動全教會以小組為單位進入社區做禱告行徑..
我們要藉著禱告呼求復興..
我們要為主贏得我們所居住的城市..
弟兄姊,你聽見這行動興奮嗎?.
從七月五號到七月十二號,千萬不要錯過..
讓我們一同禱告,一同征戰,一同經歷神的得勝..
弟兄姊,今天當你來到這個教會,你一定要明白這個教會是非常看重團隊事奉..
你看見我們教會每一件事情都是一個團隊..
團隊事奉是一個突破跟擴展不可少的關鍵..
不是靠牧師們,不是靠少數的童工們,而是每一位信徒,每一個團隊,.
大家同一個心靈,同一個腳蹤,我們就同心建造神榮耀的教會..
我們一起禱告..
謝謝你給我們這麼美好的屬靈大家庭..
讓我們在這裡學習彼此的相顧,彼此的配搭..
讓我們同心合意與你一同來建造教會..
求你再次賜下謙卑的心在我們每個人心中..

$^{281}$讓我們學習彼此的連結在一起..
求你自己親自來祝福每一位弟兄姊妹..
讓我們在事奉當中能夠被你建造..
讓我們生命也因此成長..
讓我們在這個地方能夠一同經歷你的同在,.
經歷你的大能..
當我們在這裡為你建造榮耀的教會的時候,.
當世人看見,就要歸榮耀給你..
將一些榮耀歸給你..
奉耶穌得逞的名祈求..
阿們..
\newpage



\section{以弗所書 5:21-33}
\label{sec:032080}
\textbf{共同的使命:婚姻}~~~~陳志輝牧師
\newline
\newline
連結: \href{http://www.church.com.hk/acms/content.asp?site=cdc\&op=show\&type=product\&code=032080\&layout=sermon}{\texttt{032080}} ~~~~ 語音日期: 2026-05-03
\newline
\newline
\hyperref[sec:032079]{\small{< < < PREV SERMON < < <}}
~
\hyperref[sec:index_chronic]{\small{[返順時目]}}
~
\hyperref[sec:032081]{\small{> > > NEXT SERMON > > >}}
\newline
\newline
以弗所書 5:21-33
\newline
\begin{longtable}{cl}
\hline
\hline
章節 & 經文 (和合本修訂版)\\
\hline
5:21 & \begin{tabularx}{0.7\textwidth}{X} 要存敬畏基督的心彼此順服。 \end{tabularx} \\ \\ \relax
5:22 & \begin{tabularx}{0.7\textwidth}{X} 作妻子的，你們要順服自己的丈夫，如同順服主。 \end{tabularx} \\ \\ \relax
5:23 & \begin{tabularx}{0.7\textwidth}{X} 因為丈夫是妻子的頭，如同基督是教會的頭；他又是這身體的救主。 \end{tabularx} \\ \\ \relax
5:24 & \begin{tabularx}{0.7\textwidth}{X} 教會怎樣順服基督，妻子也要怎樣凡事順服丈夫。 \end{tabularx} \\ \\ \relax
5:25 & \begin{tabularx}{0.7\textwidth}{X} 作丈夫的，你們要愛自己的妻子，正如基督愛教會，為教會捨己， \end{tabularx} \\ \\ \relax
5:26 & \begin{tabularx}{0.7\textwidth}{X} 以水藉著道把教會洗淨，使她成為聖潔， \end{tabularx} \\ \\ \relax
5:27 & \begin{tabularx}{0.7\textwidth}{X} 好獻給自己，作榮耀的教會，毫無玷污、皺紋等類的缺陷，而是聖潔沒有瑕疵的。 \end{tabularx} \\ \\ \relax
5:28 & \begin{tabularx}{0.7\textwidth}{X} 丈夫也應當照樣愛妻子，如同愛自己的身體；愛妻子就是愛自己了。 \end{tabularx} \\ \\ \relax
5:29 & \begin{tabularx}{0.7\textwidth}{X} 從來沒有人恨惡自己的身體，總是保養愛惜，正像基督待教會一樣， \end{tabularx} \\ \\ \relax
5:30 & \begin{tabularx}{0.7\textwidth}{X} 因我們是他身體的肢體。 \end{tabularx} \\ \\ \relax
5:31 & \begin{tabularx}{0.7\textwidth}{X} 「為這個緣故，人要離開父母，與妻子結合，二人成為一體。」 \end{tabularx} \\ \\ \relax
5:32 & \begin{tabularx}{0.7\textwidth}{X} 這是極大的奧祕，而我是指基督和教會說的。 \end{tabularx} \\ \\ \relax
5:33 & \begin{tabularx}{0.7\textwidth}{X} 然而，你們每個人都要愛妻子，如同愛自己一樣；妻子也要敬重她的丈夫。 \end{tabularx} \\ \\
[1ex]
\hline
\hline
\end{longtable}
$^{1}$今天和大家分享一下講題.
我想 怎麼說呢.
讓我想起幾十年前結婚.
我不知道你們有沒有想過這些事情.
就是你結婚之後 過去那十幾二十年.
三十年 你有沒有想起你當時結婚的事情.
共同的使命及婚姻.
我回想起來 咦 不是的 我結婚的時候.
沒有想過共同的使命 我不知道你們有沒有想過.
結婚的時候想什麼呢.
第一 有沒有錢.
我們那時候沒有錢也結不了婚.
因為要擺酒的 對不對.
除了想錢 還想有沒有地方.
你不要以為有錢有地方.
我女兒找了一年也找不到地方結婚.
他們要找一些特別好的地方.
要到處去找.
那時候我們找的不是那麼浪漫.
我們是找酒樓.
你擺在哪裡 酒樓還行嗎.
你又要想什麼呢.
想請誰.
請酒樓那個就糟糕了.
請得多你就要負擔一下.
你會想很多這些事情.
但是沒有想過共同的使命 婚姻.
在這個講題.
婚姻有什麼共同的使命呢.
有兩段經文.
一段是在二弗所書.
五章二十一節至二十三節.
看過的 幾個月前我和你們分享過.
第二個傳道書 四章九節至十二節.
也都在以前和你們分享過的.
三個月前.
這次我們再看這兩段經文.
然後藉著這兩段經文.
對我們的婚姻有什麼提醒呢.
是不是我們有些東西真的看漏了呢.

$^{41}$因為那時候我們很簡單.
我和我師母拍拖那麼久了.
OK的 我們可以結婚了 有錢.
兩邊父母贊同.
找了酒樓 找了時間.
誰是什麼時候 請什麼人 搞定了.
那就結婚了.
共同的使命.
當時完全沒有討論過.
也都沒有想過.
在這個題目的時候.
我就在想.
有什麼共同的使命呢.
我們藉著這兩段經文.
我們去深入的去看一看.
我們先祈禱.
天父都是多謝你的.
你將你的話語留給我們.
讓我們藉著你的話語.
再一次的去反思我們的生命.
再一次反思你的話語.
給我們的提醒.
叫我們的生命能夠在地上.
活出你叫我們所走的一切.
讓我們的生命能夠去榮耀你.
也都得忘你的喜悅.
禱告為交託 奉旨耶穌基督得請命求.
Amen.
好.
再次重溫.
我會讀出來.
那你們可能會質疑.
OK的 那我讀出來.
在二弗所書五章二十一至六十三.
至十一節經文裡面.
第一個要講到.
要全敬畏基督的心.
彼此馴服.
作妻子的.
你們要馴服自己的丈夫.

$^{81}$如同馴服主.
因為丈夫是妻子的頭.
如同基督是教會的頭.
他又是這身體的救主.
教會怎樣馴服基督.
妻子也要怎樣凡事馴服丈夫.
丈夫的你們要愛自己的妻子.
正如基督愛教會.
為教會捨己.
以水浴著渡把教會洗淨.
使他成為聖潔.
好獻給自己.
作榮耀的教會毫無煎污.
就民等類的缺憾.
而是聖潔沒有瑕疵的.
丈夫也應當.
這裡幾節.
丈夫也應當照樣愛妻子.
如同愛自己的身體.
愛妻子就是愛自己了.
從來沒有人.
恨惡自己的身體.
總是保養外色.
正像基督待教會一樣.
因我們是他身體的肢體.
為這個緣故.
人要離開父母.
與妻子結合.
二人成為一體.
這是極大的奧秘.
而我是指基督和教會說的.
然而你們每個人.
都要愛妻子.
尊重她的丈夫.
考你們的記憶力.
三個月前.
我講這段經文時.
我提及過.
我提及過在二十二節.
作妻子的.

$^{121}$你們要信服.
自己的丈夫.
原文是沒有「信服」這個字.
在二十一節.
它有這個字.
有「信服」這個字.
但二十二節是沒有的.
但為何現在會放進去呢.
因為那時在原文.
很多時前面的意思.
會放進去.
但是不需要出這個字.
有其他例證.
在聖經裡面.
提一提這件事.
在這裡講.
很多時講到妻子怎樣.
丈夫怎樣.
十節的經文都是這樣.
妻子怎樣.
丈夫怎樣.
但是.
有沒有留意十節的經文.
講妻子怎樣做.
講丈夫怎樣做.
首先.
是講了第二十一節.
有一個.
最重要的那一節.
保羅.
他不講妻子.
他不講丈夫.
他先講擺了甚麼呢.
這些東西.
是對基督徒講的.
第一.
為何在二十一節.
他就說.
要傳敬畏耶和華的心.
這句說話.

$^{161}$已經是對基督徒講.
即是我和我太太.
結婚的時候.
和師母結婚的時候.
他這句說話是和我們講.
這個是.
最重要的一個根基.
你不要講.
老婆怎樣.
先生要怎樣做.
先講.
用在基督徒.
第二樣.
最重要的根基在哪裡呢.
就在這一節裡.
他說.
要傳敬畏耶和華.
然後彼此.
信服.
他說其他東西不要講.
不要講丈夫.
要做的事.
妻子要怎樣做.
先不講.
先講.
第一.
你是否基督徒.
是否神的兒女.
是不代表你敬畏.
聽清楚.
不代表一定敬畏神.
因為我們很多時候講的.
是恩典時代.
我們講很多恩典.
我們講愛.
但是我們未必會著重.
敬畏基督的心.
而保羅呢.
他先在.
這個根基之上.

$^{201}$就說你要做這件事.
然後他開始下經文.
就說你太太.
是怎樣做.
所以這個是基石.
是那個根基.
我們很多時候.
完全沒有想過.
我當時信了主.
大概十年.
我們.
打算結婚.
兩個人一起生活.
有自己的生活.
離開父母.
完全沒有想過.
原來最重要.
我們如果有.
敬畏基督的心.
我們就開始學習.
彼此信服.
這是一個很重要的.
共同態度.
就是說兩個人.
兩個基督徒.
合在一起的時候.
他們有一個敬畏神的心.
就很不同了.
為什麼呢.
我處理過幾次.
那個家庭的紛爭.
我想問你.
如果你們有些朋友.
或者你有些經歷.
他們是兩個爭的話.
就是兩個有紛爭.
家庭紛爭.
那個太太跟你說的時候.
會不會說.
他錯呢.

$^{241}$不會的.
太太當然說他對.
先生不要.
因為兩個都認識.
我就要聽老師說.
老師不會說我錯.
不會的.
他就會說太太.
不對的地方.
出現了什麼問題.
我不知道.
如果你試一下去調停一些東西.
一定是這樣的.
少了什麼呢.
少了這個.
這個是我們家裡.
掛了十幾年的一個牌.
基督是我家之主.
就是剛才保羅.
所說的.
你是基督徒嗎.
有沒有敬畏.
基督的心.
敬畏基督的心.
就是說.
不是你先生說的是對的.
不是你太太說的是對的.
基督說的是對的.
他是主.
但是很多時候.
在家庭裡不是他主.
我是主.
我是覺得這樣做.
太太可能覺得這樣做.
沒有了這個.
所以保羅為什麼這麼重要.
在根基就說.
不用說先生不用說太太.
先說你是不是基督徒.
第二樣.

$^{281}$就是說.
你有沒有敬畏基督的心.
這是第二樣.
有開始說.
做其他東西.
我不知道你們有沒有.
任何人當中.
有沒有學過功夫.
我小時候是學.
空手道的.
在中學的時候學空手道.
學其他功夫.
學空手道最重要的是什麼呢.
在國術來說.
就是紮馬.
我們學空手道.
他不會教你.
自由搏擊.
他教你一些根基.
在哪裡.
我們的根基.
就是練被人打.
學空手道.
頭手練是被人打.
被人打的意思就是說.
你躺在那裡.
不是升起上來.
是這樣的形狀.
用支棍戳你的肚子.
把肌肉拿出來.
那些是根基.
所以保羅也是.
保羅他不說教你們怎樣做.
先說.
你既然是基督徒.
你就要有敬畏耶穌基督的心.
好了.
你有了這東西.
就開始說你要怎樣做.
這個根基是很重要的.

$^{321}$在第21節那裡.
當你們兩個人.
態度都一樣的話.
那.
這個婚姻.
就會成為了什麼呢.
成為了一個教會的.
一個縮影.
你有一個共同態度.
是敬畏耶和華的.
OK.
家庭就成為教會的縮影.
你想想家庭怎樣會成為教會的縮影呢.
很簡單.
教會是敬拜的場所.
但是教會裡面.
很多的事奉.
實在在座當中都有參與.
在事奉當中.
也都成為你那個.
屬靈生命成長的一個場所.
一個.
去服侍的場所.
而在家庭那裡.
也都完全一樣.
你的婚姻.
有時我覺得婚姻比教會.
還適合成為.
一個訓練的場所.
很簡單.
有沒有聽過.
禮拜日的基督徒.
回到家裡.
就不是的.
教會容易做的.
為什麼呢.
一個禮拜.
你聰明的.
我當你回兩天教會.
一個禮拜.

$^{361}$你多少天在家庭.
七天都在家庭.
所以.
家庭反而.
還成為一個可以.
敬拜訓練的一個場所.
所以家庭.
這個婚姻在家庭裡面.
就成為一個教會的縮影.
就在那裡去訓練.
因為剛才保羅講過.
你們兩個.
是敬畏基督.
就開始去服侍.
開始互相.
彼此服侍.
在家庭.
這個基礎.
就是他提到的.
我們看看第二段經文.
給我們一些學習的地方.
在傳道書裡.
他講到兩個人總比一個人好.
他們勞碌同得美好的報償.
若是跌倒.
這人可以扶起他的同伴.
倘若孤身跌倒.
沒有別人扶起他來.
一兩人同睡就得暖和.
一人獨睡怎能暖和呢.
若遇敵攻擊.
孤身難擋.
兩人就能抵擋他.
三股合成的繩子.
不易截斷.
第二段經文講到.
那個關係.
他講到.
勞碌同得美好的報償.
你可能會問.

$^{401}$勞碌.
我去勞碌不就行了嗎.
我去服侍不就行了嗎.
有些是完全不行的.
很簡單.
我在事奉那裡.
我可不可以說.
我不理家裡.
我去讀神學.
不行.
如果太太.
當日她是反對的話.
我可不可以就這樣.
不上班.
就去讀神學.
不可以.
因為家庭裡有兩個人.
我有那麼多.
的服侍.
那麼多服侍甚麼呢.
帶團去以色列.
一個人不行.
我看得.
前面.
看不到後面.
我太太走前我就走尾.
我太太走尾我就走前.
我參加那麼多的.
事奉裡面.
開教會.
我太太都要默默.
她不想她都要去做.
我不知道你們參與.
事奉當中如果一對夫婦的話.
你是有些事.
她不參與.
是沒有辦法成就.
我開教會.
她要做關顧.
她在這裡.

$^{441}$她做金物.
所以很多時候.
在事奉當中.
為甚麼兩個人.
是更好效果呢.
因為我一個人做不到.
顧得頭就顧不到尾.
沒有可能.
但如果我和太太.
一起的話.
那就差很遠了.
因為姐妹很多時候做關顧.
很簡單 去探望姐妹.
我一個人去又不方便.
那我太太就要去.
我去事奉拖了我太太.
很多時間.
很多.
但在神的角度來看.
是更好的效果.
為甚麼呢.
神安排了你這個婚姻.
神安排了你兩個在這裡.
然後兩個一起去事奉.
從神的角度.
他會不會覺得.
好啊 夫婦一起.
在事奉當中.
所以效果會好很多.
因為有些東西.
一個是做不到.
跌倒.
我相信.
每個人在你人生當中.
都有跌倒的時候.
不論是屬靈方面.
不論是屬世方面.
我們屬世的.
老闆罵過我的.
甚至.

$^{481}$很尖酸刻薄的事.
都有的.
我太太的老闆.
也對我發脾氣.
回到家裡.
會怎樣.
那你就可以有一個安慰.
有一個安樂家.
可以彼此鼓勵.
事奉有跌倒的時候.
因為事奉.
有時候我見一些.
做media的弟兄.
我都見過.
事奉不能夠.
覺得永遠.
一切順利.
一切開心.
因為有很多時候我遇到的事奉.
是被人冤枉你的.
但是你不能夠.
為這些冤枉而放棄.
因為主耶穌基督.
都是被人冤枉的.
他是怎樣做呢.
他沒有出聲.
這些就是我們的操練.
但是沒有出聲.
我們有家庭有婚姻的時候.
互相能夠扶持.
所以我和太太.
很多時候.
早上走路的時候.
我們都討論教會的一些事情.
團契的一些事情.
我們會不同的角度去看.
有時候是說.
聖經裡面的東西.
彼此的扶持.
對於我們在事奉.

$^{521}$和走前面的人生路是很大的幫助.
因為我們有跌倒過.
我不相信.
你們一生都說.
我沒有跌倒的.
你有失落的時候.
而失落的時候.
你就要有人扶持.
你可以說.
不是 我有弟兄姊妹.
我請問你的弟兄姊妹.
是不是24小時呢.
是不是一個星期7天呢.
不是喔.
但是你的婚姻是.
因為你的婚姻.
是一生一世.
所以.
為什麼是這麼重要.
在婚姻裡面.
是成就了很多的磨練.
操練.
第三點.
就比較困難去解釋.
三股繩子的力量.
怎麼會有三股.
繩子的力量呢.
若遇敵攻擊.
孤身難擋.
不能抵擋.
為什麼會有三股呢.
不同的.
我給大家看一張照片.
這張照片.
你看到是兩條繩子.
當然是強過一條繩子.
強過一條繩子.
兩條.
為什麼要說到三條呢.
給你看看另一張照片.

$^{561}$有沒有人看得到這條是什麼繩子.
大船那些.
因為我行船.
我已經不知拿過多少次了.
這條繩子.
這條繩子和剛才那條繩子.
分別在哪裡呢.
一個婚姻.
先生.
太太.
兩條繩子.
很好.
剛才那兩條繩子很好.
但在這些船裡.
絕對不能夠用那兩條繩子.
就算有多粗也不行.
為什麼.
你留意這兩條繩子.
哪一條比較容易滑.
你留意一下.
這條是大船.
我們叫Morning Rope.
這條船.
你看看繩子的設計.
是不容易滑的.
當你拿的時候有多粗呢.
一隻手是拿不完的.
那條繩子來說.
是不容易滑的繩子.
你看看它破了.
斷了.
那裡已經磨了.
它照樣可以用.
因為這裡不是兩條繩子.
是三條.
婚姻.
不是兩個人.
是三條繩子.
太太.
先生.

$^{601}$還有第三條繩子是什麼.
神.
基督.
是一家之主.
當這三條繩子.
交叉地捲在一起.
這條.
是可以拖到十萬噸的船.
都可以Hold住這條繩子.
你不要小看這條繩子.
如果你留意.
大船停泊在碼頭的時候.
沒有這條繩子.
是停泊不了的.
你會見到有一等船.
推船過去.
但是沒有辦法.
拉住船.
這條繩子.
當時我們一個站在船頭.
一個站在船尾.
不是將繩子拋去岸上.
你拋不動.
因為太粗.
拿另一條小的繩子.
吊上去.
拉上去.
放在椅子上.
然後用船的機器.
拉船過去.
這條繩子你知道有多強.
幾萬噸.
十萬噸的船.
都可以拉過去.
就是靠這條繩子.
你去碼頭留意一下.
船頭兩條至三條.
船尾兩條至三條.
就是將大船拖到岸邊.
分一日吊上去.

$^{641}$當你能夠將.
基督的心.
敬畏基督的心在的時候.
再加上.
太太,先生.
加上基督.
三條.
你可以在.
打暈天下.
幾多宣教士.
整個家庭去.
一生都在.
荒蕪之地去宣教.
三樣.
這就是.
三股城的力量.
在事奉當中.
如果你.
夫婦一起去參與.
再加上.
敬畏基督的心.
在裡面.
你整個事奉就很不同.
你會有喜樂.
你會有驚喜.
你會經歷到神.
實在很多次的事奉.
是我太太.
是被動的.
我讀神學.
神感動的是我.
沒有感動她.
她是被動者.
但當她如果看到.
先生所做的.
是為神.
她就算沒有受感動.
她都會支持.
我們開展.
另外一個事工.

$^{681}$在斯多頓.
有斯多頓弟兄姊妹.
她都是一個被動者.
我們開車去.
一個多小時.
她都是一個被動者.
但她一直會支持的.
這個退休社區.
開教會.
她都是一個被動者.
因為我提出來.
但當我們兩個.
一起的時候.
為神的緣故.
而神也在我們中間.
教會開了.
我們的Christian Village.
在這個退休社區.
也開了團.
我們要開.
Non-profit Organization.
SOW 開了.
我要帶團 她要跟.
她是一個.
信服者.
師母是一個信服者.
她沒有試過一次.
反對我的事工.
沒有試過.
但每次都是.
跟隨者.
一起去.
接著會看到什麼呢?.
她就見到神.
她看到神.
看到神所做的事.
沒有可能的事.
就發生了.
婚姻.
是神所祝福的.

$^{721}$是神所設的.
但在婚姻當中.
你的根基.
在不在.
你有沒有敬畏.
基督的心.
當你夫婦.
都是敬畏基督的心的時候.
太太.
可以信服丈夫.
不是所有的事.
是信服.
是屬靈的事.
是神的事.
三古誠旨.
你的婚姻就會成為這三古誠旨.
誠旨.
在神的事工.
在神的安排之下.
你會看到神蹟.
你生命.
可以被造就.
第三.
第一.
保羅說到根基重要.
你要有敬畏的心.
第二.
當你在整個家庭.
有婚姻的時候.
你不只是兩古誠旨.
是要三戶.
因為神要在你的家庭當中.
你就有三樣.
而第三.
就是使命.
我以前結婚.
不會想過有什麼使命.
但你如果看回.
二弗所書.
他經常說有根基就是.

$^{761}$丈夫做什麼什麼.
妻子要怎樣怎樣.
丈夫作為寫記的僕人.
正如基督為教會寫記.
他將丈夫.
就代替了.
好像基督為教會寫記.
你即是說.
丈夫為妻子而放棄很多東西.
犧牲很多東西.
好像基督為了教會一樣.
一個比喻.
妻子呢.
就成為一個樂意的回應者.
是馴服主.
即是說馴服丈夫.
但現在說的是屬靈的事.
你不要說.
丈夫叫太太.
跟我一起去喝酒.
去賭錢.
這是屬靈的事.
不要說.
只是取了字面的意思.
那個馴服.
是馴服基督.
即是說當.
丈夫.
做一些事工.
做一些神所喜悅安排的事.
是太太去馴服那件事.
所以就不是.
斷章取義.
那你會留意到.
在二弗所書.
基本上講到教會.
和基督的關係.
那有沒有留意.
如果我們.
將丈夫.

$^{801}$譬如好像.
為教會寫基.
妻子又要馴服基督的話.
那婚姻.
的使命是什麼呢.
就成為了.
好像.
教會基督一起.
婚姻完成的.
大使命訓練的.
實踐的場所.
就好像教會可以有門徒訓練.
我相信大部分你們都可能.
參加過教會的門徒訓練.
完了.
你婚姻的門徒訓練.
沒得用.
一生一世.
你的生命成長.
是一生一世在婚姻裡面.
我相信你們會經歷.
很多的爭吵.
不同意見.
看法不同.
但是這些.
是成為你生命的成長.
婚姻並不是.
講理由的地方.
很多時候沒得講.
我明明.
覺得牙膏.
是凍膏的.
我太太.
不是放下.
對還是錯呢.
沒得講.
但是如果我贊同.
太太的看法.
我的是什麼呢.
我是接受.

$^{841}$不是看自己為重.
是看對方為重.
實際上門徒訓練.
最重要的東西.
就是你看到.
別人的需要.
在婚姻裡面.
實際上最大的操練.
不是對還是錯.
是你能不能夠看到.
你另外一邊的需要.
而用你的愛.
去包容他.
不是他對.
你的生命成長.
就是這樣.
教會有沒有什麼機會呢.
有,在事奉當中也有不同意見.
但是這個短暫.
教會也是七天.
兩天就足夠.
在小組裡面也可以.
討論一下.
不同意見也可以.
但是家庭呢.
一生一世.
婚姻神給你就是一生一世.
一生一世.
你的使命就要去.
互相的操練.
去成就對方的.
生命的成長.
你這樣做.
影響不只是你.
是影響你的下一代.
就好像我和我太太一樣.
我們影響到.
是我們的兒女.
兩個女兒.
我和我太太.

$^{881}$尼赫曼斯姆.
我可以說如果記得的話.
是沒吵過架的.
可以有不同意見.
但是沒吵過架的.
最多不開心.
就走了.
所以在我的兒女面前.
她沒見過我和師母.
吵架的.
意見見到我們有不同.
神要我們這段婚姻.
能夠互相建立.
生命成長.
也都因為這樣.
賜福給下一代.
我們的兒女.
見到我們.
是影響了他們.
所以就見到.
我和我太太沒吵過架.
這裡就影響下一代.
也都賜福給下一代.
第三樣.
公開的見證.
教會是成為一個燈台.
我們有敬拜.
有很多事公在這裡.
但是.
婚姻.
是一個活生生的.
見證.
人家看見你.
你活出來.
經文我們可以研讀.
我們可以研讀經文.
聽到容易.
行到能.
但是當你的婚姻.
活出來的時候.

$^{921}$人家就看到.
這兩夫妻.
從早上叫到晚上.
不知道你們有沒有.
聽過這些.
特別是婚姻破裂那些.
一定吵得很厲害.
婚姻要行一個大使命.
是怎樣呢?.
就是能夠將.
彼此的人生.
能夠建立去成就.
而這一樣.
亦都彰顯到在人的面前.
而婚姻.
既然你說是神設立的.
神設立一個甚麼的婚姻呢?.
活出甚麼的婚姻呢?.
這個使命.
是一生一世.
教會門徒訓練.
可以一個月兩個月三個月.
但是婚姻.
門徒訓練.
是一生一世.
而且這個訓練.
不只是在經文裡.
是要你行出.
經文要求的東西.
要行出神要你行的東西.
我們在座當中.
弟兄姊妹.
每個人的婚姻.
有每個人婚姻的問題.
但是記住.
婚姻裡面一定有問題.
即是沒有問題.
是沒有辦法磨到你的.
沒有問題.
你是沒有機會去改變的.

$^{961}$而你這個是一生的使命.
你的婚姻.
是成為你一個試圖.
操練的地方.
讓你的生命能夠.
繼續去改變.
朝著神要求你的改變.
你不要用理由.
不要用你個人.
你是將.
基督擺在一家之主.
當你有問題的時候.
請你將基督.
擺在你們兩個人的中間.
這樣.
你就能夠學到.
在生命當中的成長.
神要你的婚姻.
是怎樣彰顯到.
祂的心意出來.
我希望今次的講道.
能夠再一次思想.
幫助我們思想.
我們的婚姻原來就是這樣.
原來在神的心願裡面.
婚姻是要你.
彰顯到.
神要你所做的事.
在《馬太福音》講道.
向萬民.
作他的門徒.
要使萬民作他的門徒.
而婚姻就是要你.
一生的操練.
成為神的門徒.
你就不能.
與神所設立的人.
不難分開.
一生裡面去磨練.
將基督擺在一家之主.

$^{1001}$我們一同祈禱.
將婚姻賜在我們當中.
你揀選了.
我們的配偶.
幫我們成立家庭.
我們知道.
我們每個人有自己.
我們知道每個人都有自我.
但是.
你將主基督.
擺在我們的家庭當中.
讓我們這個婚姻.
成為一個門徒訓練的場所.
讓我們這個婚姻.
能夠活出.
你的見證.
能夠將這個婚姻.
展示在人的面前.
能夠將榮耀.
歸給你.
求你賜福我們每一段的婚姻.
讓每一段的婚姻.
都成為.
你.
彰顯.
主耶穌基督的.
生命的場所.
讓每一個婚姻.
都能夠.
去榮神.
去益人.
我們禱告交託.
奉主耶穌基督得勝命.
阿們.
\newpage



\section{以賽亞書 49:1-6}
\label{sec:032081}
\textbf{外邦人的光}~~~~李應新牧師
\newline
\newline
連結: \href{http://www.church.com.hk/acms/content.asp?site=cdc\&op=show\&type=product\&code=032081\&layout=sermon}{\texttt{032081}} ~~~~ 語音日期: 2026-05-03
\newline
\newline
\hyperref[sec:032080]{\small{< < < PREV SERMON < < <}}
~
\hyperref[sec:index_chronic]{\small{[返順時目]}}
~
\hyperref[sec:032082]{\small{> > > NEXT SERMON > > >}}
\newline
\newline
以賽亞書 49:1-6
\newline
\begin{longtable}{cl}
\hline
\hline
章節 & 經文 (和合本修訂版)\\
\hline
49:1 & \begin{tabularx}{0.7\textwidth}{X} 眾海島啊，當聽從我！遠方的眾民哪，要留心聽！自出母胎，耶和華就選召我；自出母腹，他就稱呼我的名。 \end{tabularx} \\ \\ \relax
49:2 & \begin{tabularx}{0.7\textwidth}{X} 他使我的口如快刀，把我藏在他手蔭之下；又使我成為磨利的箭，把我藏在他箭袋之中； \end{tabularx} \\ \\ \relax
49:3 & \begin{tabularx}{0.7\textwidth}{X} 對我說：「你是我的僕人以色列；我必因你得榮耀。」 \end{tabularx} \\ \\ \relax
49:4 & \begin{tabularx}{0.7\textwidth}{X} 我卻說：「我勞碌是徒然，我盡力是虛無虛空。耶和華誠然以公平待我，我的賞賜在我的神那裡。」 \end{tabularx} \\ \\ \relax
49:5 & \begin{tabularx}{0.7\textwidth}{X} 現在耶和華說話，他從我出母胎，就造我作他的僕人，要使雅各歸向他，使以色列聚集在他那裡。耶和華看我為尊貴，我的神是我的力量。 \end{tabularx} \\ \\ \relax
49:6 & \begin{tabularx}{0.7\textwidth}{X} 他說：「你作我的僕人，使雅各眾支派復興，使以色列中蒙保存的人歸回；然而此事尚小，我還要使你作萬邦之光，使你施行我的救恩，直到地極。」 \end{tabularx} \\ \\
[1ex]
\hline
\hline
\end{longtable}
$^{1}$叫我們這段經文.
剛才讀的這段經文.
叫我們作外邦人的光.
這裡是第二首「搏人之歌」.
我們又翻開聖經.
我今天騙了你去聖經.
快點翻出來.
我們翻到聖經裡說.
在哪裡呢?.
那裡是這段經文的第六節.
是要成為外邦人的光.
看見嗎?.
第六節.
自己手指指下去.
我還要使你們作外邦人的光.
這是「搏人之歌」第二首.
上一次跟你揭曉時揭了四首.
記不記得?.
有沒有印象?.
樣子好像沒有人記得.
四首.
這是第二首.
上次說了第四首.
這次說第二首.
為何跳下去不是順著呢?.
因為上一次貼著壽苦節,復活節.
被迫要先說第四首.
這次我回頭說第二首.
我們成為外邦人的光.
「搏人」是誰呢?.
記不記得我上次說過.
「搏人」最認是主耶穌.
但沒有理由只有主耶穌才對.
有可能是所羅巴伯.
有可能是以賽亞先知.
這裡像以賽亞先知比較多.
不過無論如何.
就算我們完全適合是主耶穌.
其實沒有理由也不關我的事.
同樣.

$^{41}$說主耶穌不關我的事.
不是.
主耶穌要做的事你也要做.
主耶穌是那個受苦的「搏人」.
難道你想自己卸駁.
裝模作樣不做「搏人」嗎?.
你是不是也是主的「搏人」.
各照自己地方本事應當.
行不行?.
你也是要成為主的「搏人」.
成為外邦人的光.
行不行?.
我今天要告訴你重點.
記住你們每個人都要做主的「搏人」.
首「搏人」之歌.
既預告是主耶穌.
同樣.
其實和你也有關.
你們也要成為主的「搏人」.
我們從頭又看下去.
這首歌的視野很寬廣.
「縱海島當聽我言遠方的眾民留心而聽」.
這裡說了兩件事情.
說的不是在附近.
「縱海島」和「遠方的眾民」.
都是在說超越附近的視野.
不只是說給自己的子民.
我們在這裡真好.
彼得登山變相的時候.
不是我們要下山光照在眾人前.
所以這個視野很廣闊.
而且後面那裡做了Echo.
剛才你還要作第六節.
說作外邦人的光.
這樣感受到.
後面那裡又再提醒你.
這裡不只是你附近.
這個視野是非常廣闊的.
還有叫你施行我的教恩第六節.
說直到地極.

$^{81}$你看著聖經寫著.
直到地極.
不要給太多你看.
你自己看你那本.
你自己抓來抓去.
你看你那本聖經.
要直到地極.
所以視野非常廣闊.
各位電視節目人.
我們都是的.
其實我們做基督徒.
不只是看自己教會好不好.
只是看我們身邊的事情.
我們關心我們身邊.
固然本事應當.
我們要開放日.
招呼我們的街坊.
本事應當.
但是基督徒的視野.
縱海島和去到地極.
去到外邦.
我們都需要關心.
而事實上我們都見到.
很多神的僕人.
他忠心的僕人.
去到遠方.
你想都想不到的地方.
上期.
上一期時論壇.
你自己找找.
上一期時論壇訪問了.
有個叫施達基金.
他們有兩位姊妹.
去哪裡呢?.
去了阿富汗.
做人道救援工作.
原來前一陣子.
阿富汗地震.
他們去做人道的工作.
阿富汗都去?.

$^{121}$美軍都走?.
不是.
美軍走.
但是神的僕人是會去的.
其實還有很多.
你沒有告訴你的地方.
他們那麼大膽登出來.
其實按我所知.
很多主崩忠心的僕人.
做了你都不知道.
不敢告訴你.
免得你害怕.
所以其實神的僕人.
視野.
作為忠心的僕人.
視野是去到縱海島.
去到地極.
作為忠心的僕人.
視野也不放得那麼窄.
只在我們身邊當中.
神感動你.
我們自己去.
神也可能感動你.
支持你身邊的弟兄姊妹.
或者一些他的僕人去.
都是參與在其中.
我們的視野擺在很遠的縱海島.
而且這裡所講的宣告.
是向著縱海島和遠方的眾民.
他也講解他自己.
這個僕人自我出胎.
耶和華就選照我.
自出母腹.
他就提我的名.
即是他未出世的時候.
上帝就認識他.
而且是一個Echo.
第五節.
耶和華從我出胎.
看著讀著.

$^{161}$指著來讀.
就是造就我.
作他的僕人.
要使我各歸向他.
使他哪裡聚集.
即是說這裡講到什麼呢.
是神揀選的.
我們成為神的選民.
這位僕人也好.
我們也好.
神認識的.
我們未出世.
神已經認識我們.
這個也非常之.
詩篇139.
即是怎樣啊.
牧師.
讀吧.
揭下去.
詩篇139.
我們揭下來讀一讀.
詩篇139篇.
揭下吧.
詩篇139.
13到17.
我們一起讀.
詩篇139.
也是我很喜歡.
很多電影節目我很喜歡的經文.
就是在回應以賽亞書49章.
詩篇139.
我們這本聖經820.
詩篇139.
13到16.
或者17.
17到17.
我們預備.
找到沒有.
旁邊那個幫幫他.
好 OK.

$^{201}$預備 1 2 3 一起讀.
我的肺腑是你所造的.
我在母腹中.
你已浮避我.
我要稱謝你.
因我受造奇妙可謂.
你的作為奇妙.
這是我深深知道的.
我在暗中受造.
在地的深處被聯絡.
那時我的形體並不向你隱藏.
我未成形的體質.
你的眼早已看見了.
你所定的日子.
我尚未到一日.
你都寫在你的冊上了.
神啊.
你的意念向我何等寶貴.
其數何其眾多.
這裡說到.
未曾出世的時候.
天父已經認識 揀選.
不知你有沒有這個感受.
放在我的家庭當中.
我不詳細說了.
但這亦是我常常的感受.
我能夠在香港成長.
在我的家庭長大.
甚至我讀的小學幼稚園.
我當然成長時有埋怨過.
為何我媽媽那麼窮.
或者這麼說.
後來想著想著.
都不錯.
神放在我身上最合適.
我尚未成形.
體質未曾出現的時候.
神都監察.
甚至我都埋怨到.
那間幼稚園這麼不好.

$^{241}$那間小學.
不是的 想著想著.
我只說一件事.
我那時都很埋怨.
為何我不帶眼識人.
我亂報學校.
報了那間中學.
其實只有一樣東西.
我就是貪它夠近.
可以回家吃飯.
但我後來發覺.
這不是一個好處.
這是一個壞處.
這麼近家.
去逛逛都不行.
但我後來.
最主要是甚麼.
我很埋怨甚麼.
埋怨為何會選藍校來讀.
那間藍校真是壞事.
養隻貓都是功.
為何會選這些學校來讀.
不過後來才發覺.
其實那間學校.
這麼齋沒問題.
因為那間學校.
下面就是我後來成長的教會.
神真的認識我.
見到我最好的東西.
擺在我面前.
所以在那間教會.
我認識了全教會.
最好的女孩子.
後來就做了我老婆.
所以神揀選為我預備的道路.
最好.
收不收到.
這裡說到.
我們都深深知道神揀選.
這位僕人也是.

$^{281}$自我出胎.
耶和華就選擇我.
自我出母峰.
祂就提我的名.
原來我們的名字.
神都是閱立和知道的.
我想多提一件事.
我們的名字很矜貴.
知道嗎.
我們的名字記載.
生命冊當中.
主在點名的時候.
會叫你的名字.
你應新.
到.
我在這裡.
當點名的時候.
我和你都在其中.
收不收到.
不過今天不唱這首歌.
但是我們繼續看下去.
這裡第一件事.
神知道我認識我.
不過我想跳下去看下去.
第三節.
對我說.
你是我的僕人以色列.
我必因你得榮耀.
這裡就說.
原來神叫這個僕人.
叫做以色列.
即是神的子民.
這是個體單數.
這是說那個人.
是神的子民.
好像雅各,以色列.
我必因你得榮耀.
神透過我.
可以得到榮耀.
即是神的心意.

$^{321}$我們發光.
就有榮耀.
所以我們發光.
主得著榮耀.
但是我卻說.
即是他領不領情.
一般般領.
我卻說.
我勞碌是徒然.
我盡力是虛無虛空.
然而.
我先停一下.
我剛才說.
這個僕人之歌.
這一段.
像是以塞瓦書才知道.
有什麼證據.
我證明給你看.
這裡是非常 Echo.
什麼呢?.
以塞瓦書第六章.
第六章.
所以又揭了.
我們又揭了.
旁邊左右.
揭前面就可以了.
第六章.
第六章是怎樣說的呢?.
第六章.
我們揭去第六章.
我們就看第六章.
這裡.
我們看第九節.
先看八節.
八節.
我們一起.
八節.
就看到第十節.
八到十.
預備.

$^{361}$我們一起讀.
我又聽見主的聲音說.
我可以差遣誰呢?.
誰肯為我們去呢?.
我說.
我在這裡.
請差遣我.
他說.
你去告訴這百姓說.
你們聽時要聽見.
卻不明白.
看時要看見.
卻不曉得.
要使這百姓.
心懵之猶.
耳朵發沉.
眼睛昏迷.
恐怕眼睛看見.
耳朵聽見.
心裡明白.
回轉過來.
便得醫治.
你通常都不太明白.
你只是讀到.
通常第八節讀完就算了.
請差遣我.
這樣就夠了.
我們都記得了.
這個爾賽亞才被天父呼召.
我可以差遣誰.
請差遣我.
然後你停在這裡.
原來後面還有一段.
這段是什麼呢?.
耶和華告訴爾賽亞.
你去跟百姓說什麼.
你說些什麼.
不過他們.
聽時聽見.
卻不明白.

$^{401}$看時看見.
卻不曉得.
那是什麼呢?.
白說.
齋做.
沒有效果.
明白嗎?.
這樣就知道.
Echo到我們剛才讀的那個僕人.
Echo著.
我們看回去.
第四十九章.
我勞碌是徒然.
我盡力是虛無虛空.
感受到嗎?.
所以這個.
我說似爾賽亞.
都似似地是.
講明了.
你做得來.
只是做而已.
其實你講沒有人聽.
聽都不明白.
是白做的.
不過.
原來去到第四十九章的時候.
其實是有突破位.
突破位在哪裡呢?.
突破位在後面.
繼續讀下去.
然而.
這些然而.
即是突破.
之前不明白.
做是白做.
不過然而.
我當得的理.
必在耶和華那裡.
我的賞賜.
必在我神那裡.

$^{441}$突破了想法.
雖然做是沒有效果.
不過原來.
我的道理.
我的新冤情.
我的道理.
在哪裡呢?.
新冤的權利.
道理.
在神那裡.
還有.
雖然是白做.
但還有什麼呢?.
賞賜.
在下面寫著.
賞賜.
我有賞賜.
不過.
你看不到.
在耶和華那裡.
在心裡.
祂是突破了.
白做.
白做就是白做.
白做不是白做.
再說一次.
白做不是白做.
因為原來在耶和華那裡.
所以祂是突破的.
祂知道在地上是這樣的.
但是.
我知道.
神會有祂賞賜的.
我突破了.
祂突破了想法.
又不知道你有沒有這個感受呢?.
有些時候.
做的事情都是白做的.
傳福音.
好像是白傳.

$^{481}$我聽不少鄧英之妹說.
跟我爸爸傳福音傳了幾十年都不行.
有時候這樣說.
我白做.
鄧英之妹又做團契導師.
只有「搞」字.
拿來搞.
大家都不服氣.
有時間中有這種感覺.
但是.
只要你忠心.
祂突破了什麼呢?.
我知道我的賞賜.
必在我神那裡.
你忠心做就夠了.
OK的.
祂也收穫.
以賽亞收穫.
希望你也收穫.
耶和華這裡說到.
從我出母科.
就造就我作.
他的僕人要使亞國歸向他.
使以色列到他那裡聚集.
這裡呼召.
這個是眾數.
這裡的亞國和以色列是一個群體.
其實後面都在回應.
在echo.
原來耶和華看我為尊貴.
我的神也成為我的力量.
這裡也提醒我們.
其實神是我們的力量.
不是靠別人的讚賞.
靠別人的讚賞.
也沒有什麼盼望.
剛才也說了.
說了又沒有人聽明白.
又不明白.
又不看見.

$^{521}$看見又不明白.
那些都是白搞.
不過.
我知道.
我是尊貴的.
我是上帝的僕人.
我們是神的子民.
我們個個都很尊貴.
好.
試試讓你看看旁邊那個.
那個尊不尊貴.
OK的.
是才好說.
是的.
真的.
你們都是上帝尊貴.
上帝所揀選.
上帝是僕人.
不過待會我還會說.
不是.
旁邊那個有時也很討人厭.
不會的.
我待會再說.
我們每個是上帝揀選.
看為尊貴.
好.
我們說到這裡的時候.
我們繼續讀下去.
第六節.
現在他說你作我的僕人.
使雅各眾志派復興.
使以色列中得保存的歸回.
好了.
這裡就說原來他要做的事.
剛才前面那裡也說了.
是一對的.
因為這是C的體裁.
所以很多時候叫平衡體.
使雅各歸向他.
使以色列到他那裡聚集.

$^{561}$還有這裡.
使雅各眾志派復興.
使以色列中得保存的歸回.
這句話很甩咳.
即是什麼呢.
被擄了一些可以回到家裡.
民族偉大復興.
是不是很厲害.
厲害嗎.
沒看聖經.
聖經明白嗎.
尚為小事.
其實就不厲害.
你以為厲害.
但其實一點都不厲害.
讀了幾次這段經文.
尚為小事.
不厲害.
你經常想著很厲害的其實都不厲害.
但不知為何為了不厲害的事.
總是想來想去.
令你睡不著.
所以目標搞錯了.
不知為何令你睡不著.
所以那些很偉大的事情.
我舉例子.
我以前認識一位.
我不認識他.
我聽別人的分享會.
有一次97後的立法會議員.
厲害嗎.
立法會議員.
全香港只有幾十個.
有些見證說.
欠他的票.
於是通過了政改法案.
其實我忘記了他說什麼.
他很自豪.
我那時候有份投票.
我投票後.

$^{601}$陳就很使用我.
通過了那個標.
那時候爭持很激烈.
政府游說.
他說得很眉飛色舞.
欠他的票.
最後通過了.
現在那些法案.
去了哪裡.
他就覺得很厲害.
我真的不太記得他說什麼法案.
總之他就覺得自己很厲害.
有不少人說他不厲害.
聽起來也很厲害.
但厲害之餘.
其實尚為小事.
我不懂說.
但不知為何.
近這三五七年.
有時候我們為了一些事情.
爭拗.
家裡也好 教會也好.
爭拗.
爭拗到很仇恨對方.
尚為小事.
同意我說也好.
不同意也好.
其實我們有時爭拗來爭拗去.
家裡也好.
那些是小事.
希望你自己慢慢答真一點.
哪些是大事 哪些是小事.
其實說到尚為小事.
那些是小事.
哪些是大事 愛幫人信主.
傳救恩到地極.
這些才是大事.
這是神的看法.
你說不同意 我沒有辦法.
不過真的這樣寫.

$^{641}$希望你答真一點.
回過頭來.
平時你做的事情究竟是大事還是小事.
不是那些大事小事.
而是那些大事小事.
我們經常.
將很多小事搞到很大件事.
搞到你睡不著.
吃不安不樂.
搞到你很吵.
就因為將一些不是很大的事.
搞到很大件事.
我想分享多一點.
不過有時小事.
我們教會有時用很多精力來做.
其實我有時也不懂分辨.
究竟哪些是大事.
你知道牧師有時.
其實我一天到晚做什麼.
我不知道為什麼我忍不住.
回想起來我做了一件事.
你幫我分享究竟是大事還是小事.
有一次我去.
有一個姐妹家庭家裡探訪.
不是太古城堂.
不要想是誰.
不是太古城堂.
她說經常睡不著.
我說你為什麼睡不著.
牧師你看看那個鍋.
有個瓦鍋.
那個瓦鍋她說借了另外一位姐妹.
另外一個人的.
瓦鍋來煮薑醋給她的孫子.
我說你的孫子會走路.
很久了.
不就是.
那個瓦鍋我借了另外一個人.
另外一個人想還給她.
她說不要.

$^{681}$放在我那裡.
我孫子也會走路.
她說你為什麼.
為這些睡不著.
我很怕孫子會走路.
打爛他就不好了.
要打爛別人的東西.
弄傷他也不太好.
原來她擔心這些.
你覺得是大事還是小事.
小事.
牧師聽到這些小事.
是不是應該不用理他.
小事你不理.
我真的不懂分辨.
我那天就想.
這下子.
搞到你睡不著.
好像很大件事.
不如這樣.
你有沒有叫她拿回那個鍋.
叫了很多次 叫了很多次都不肯拿.
這樣吧.
我拿去我辦公室.
我放在我辦公室正中央.
跟著告訴那位姊妹.
你的鍋在我那裡.
跟著怎麼樣.
跟著將來.
拿回去.
我解決了這位姊妹.
我也不知道我做了大事還是小事.
不過anyway.
我也想你回答一下.
究竟什麼是大事什麼是小事.
有些事尚為小事.
不知你搞到什麼.
大事就是沒時間做.
不過你被小事卡住.
你做不了大事.

$^{721}$我又不是很安樂.
你整天都這麼多小事搞到我.
我又做不了大事.
你感受到嗎.
再說一件事.
最近發生的.
不要生我的氣.
說完.
最近有件事.
好 說完.
這件事是什麼.
我們的廁所有個縮.
有個縮在縮.
就會鎖了門.
但有些人跟我說.
牧師你應該害死人.
害死人有什麼重要.
你如果縮了外面就開不到進去.
那你拆了它.
那你有鎖匙.
如果它在裡面暈了.
可以開鎖來救裡面的人.
那就不會害死人.
我就跟幹事說.
那就拆了它.
我跟幹事說.
牧師 幾年前有些人投訴.
很危險.
沒有那個縮的話.
又怕去廁所會鎖.
還有一元是可以開的.
一元這麼厲害可以開的嗎.
你們都懂得開.
你們這麼厲害.
所以才加一個縮.
是不是要拆.
所以我現在有沒有拆.
我現在不敢拆.
你究竟整天我們搞大事小事.
搞到我們忘記了傳福音.

$^{761}$我們應該要來傳福音.
你搞這麼多小事.
大事我們應該要做什麼.
求主幫助我們.
親愛的天父.
我求你賜福給每一個弟兄姊妹.
我們都做你忠心的僕人.
向世界發出光芒.
叫遠方的.
直到地極的.
都得到救恩.
使用我們.
以耶穌的名義祈福!.
阿們!.
\newpage



\section{出埃及記 32:4}
\label{sec:032082}
\textbf{聖經人聖經事:阿倫的妥協}~~~~黃克勤牧師博士
\newline
\newline
連結: \href{http://www.church.com.hk/acms/content.asp?site=cdc\&op=show\&type=product\&code=032082\&layout=sermon}{\texttt{032082}} ~~~~ 語音日期: 2026-04-30
\newline
\newline
\hyperref[sec:032081]{\small{< < < PREV SERMON < < <}}
~
\hyperref[sec:index_chronic]{\small{[返順時目]}}
~
\hyperref[sec:032083]{\small{> > > NEXT SERMON > > >}}
\newline
\newline
出埃及記 32:4
\newline
\begin{longtable}{cl}
\hline
\hline
章節 & 經文 (和合本修訂版)\\
\hline
32:4 & \begin{tabularx}{0.7\textwidth}{X} 亞倫從他們手裡接過來，用模子塑造它，把它鑄成一頭牛犢。他們就說：「以色列啊，這是領你出埃及地的神明！」 \end{tabularx} \\ \\
[1ex]
\hline
\hline
\end{longtable}
$^{1}$大家好,眨眼又一日,又來到《聖經》.
《生命說書》網絡和大家繼續在《聖經》人和《聖經》的事上.
來認識主耶穌,認識我們的天父.
今次仍然在曠野.
今日我給的題目不是《聖經》的字,是我的字.
就是叫《阿倫的妥協》.
《阿倫的妥協》,阿倫是妥協了一次.
今次事情真是很大了.
有時說妥協這件事,我們經常都要做的.
原因就是人與人之間,周時都是你遷就我,我遷就你.
很多時在某些事上是要妥協的.
是要雙贏的妥協,多贏的妥協才好.
有時有些妥協,如果你說是對於一些罪惡的事情.
對於不對的事情妥協,那不行了.
特別是在一些神學問題上,真理問題上.
我們是很不想妥協的.
今次阿倫就遭遇到一件事,壓力很大.
他做了一個妥協,我們看看這個妥協吧.
《阿倫的妥協》,就是我今日的題目.
《初級記》第32章第一節.
這段是出現過在聖經裡面好幾次的.
跟大家再看這一段,以前我記得我說過的.
都是這件事件.
32章第一節.
「百姓見摩西馳言不下山」.
「就大家聚集到阿倫那裡,對他說」.
「起來,為我們做神像,可以在我們前面引路」.
「因為領我們出埃及地的那個摩西」.
「我們不知道他做了什麼事」.
好,繼續讀完.
「阿倫對他們說」.
「你們去摘下你們妻子兒女以上的金環」.
「拿來給我,百姓就都摘下他們以上的金環,拿來給阿倫」.
「阿倫從他們手裡接過來很多金環,金戒指等等」.
「注了一隻牛毒,用雕刻的器具造成的」.
即是雕刻一隻金牛毒出來的像.
「他們就說:以色列,這是領你們出埃及地的神」.
「阿倫開堅,就在牛毒面前捉牌,且宣告說」.
「明日要向耶和華守節」.
「阿倫宣佈,明日向耶和華守節」.

$^{41}$「次日清早百姓起來,很認真,就獻梵祭和平安祭」.
對著那隻牛獻梵祭和平安祭.
「就坐下吃喝,起來還撒,起來換撒,還撒」.
故事就是這樣.
聖經人,聖經故事,今次聖經人就是人物阿倫.
故事就是跟那些以色列人做了金牛.
這個時間摩西和約書亞上了山.
很久都沒有下來.
而且上帝吩咐「不可以上來,不要跟上來」.
「其他人不可以,只有你們兩個上」.
其實約書亞去到半山都停了.
摩西一個人自己上到去跟神相會的地方.
交了兩塊法板給他.
很久都沒有下來.
你知道,馳言不下山,那些百姓在想什麼呢?.
於是他們就走去,心虛了,虛怯.
因為主要是沒有在地上眼見到可以信賴的領袖,摩西.
他們就聚集去找第二個領袖,次級的領袖,就是阿倫.
去到那裡就跟他說「阿倫,起來不如這樣,做點事吧」.
這個「起來」是指做點事,不是坐著要起床.
「做點事吧,為我們做神像」.
這是第一點我要說的,紅色字,藍色螢光.
「為我們做一個神像」,這是我第一點要說的.
「為我們做一個神像」.
其實這幫人怎樣?拜偶像嗎?是呀,很明顯.
其實又可以說是還是不是,是又是對還是不對.
因為他們是要用他來在我們前面引路.
沒有人帶我們路,摩西又不見了.
其實我們不能上山,不然我們上山去找他.
「摩西,你何時下來?我們在下面等你」.
這樣又不能上去,因為當時的人都不敢上.
照計那山是冒煙,震動,聾聲,真的不敢去.
等了很久,我們都不知道他怎樣.
「因為領我們去埃及地的那個摩西,我們不知道他遭了什麼事」.
會遭了什麼事呢?當然會,上了山那麼久,我們又看不到.
他跌了下山坑,死了都可以.
受了傷跌倒跛了,不能走路都可以.
慢慢餓死,傷死都可以.
我們怎樣?等多久?.
其實都是等了不會太久,不過心虛怯.

$^{81}$「唉呀,不見了,那怎麼辦呢?摩西」.
都挺有趣的,不過.
這個「摩西不在,但我們沒有人帶」.
這件事其實不完全對.
原因就是上帝帶領以色列人在抗疫的時候是怎樣帶的?.
是有雲柱和火柱的.
日間有雲柱,夜間有火柱.
雲柱,火柱未升起,他們是不會前行的.
如果日間雲柱是停在那裡,一直有雲柱在那裡.
他們是不會收拾東西,不會拆帳篷,不會搬的.
到了一定時間,雲柱收起.
於是大家見到「要起行了」.
於是他們收拾好,然後預備起行了.
所以其實是否沒有人帶我們前面的路.
於是我們做個神像?不是的.
其實是有雲柱,火柱帶的.
而且還有亞倫在那裡,繼續等雲柱,火柱.
就算摩西在那裡.
如果雲柱,火柱不動,大家都是不會搬遷的.
都不會收拾東西,不會拆帳篷的.
不過這個時候,百姓出現問題.
就是心裡面有些虛怯,又或者想找些事情做.
那些事情是摩西不在才適合做的.
他們看得很清楚了,於是就去找亞倫.
「做些事情吧,摩西不在,都不知道他上面是否死了」.
「不如我們要做一個神像,做神像」.
怎樣跟大家說呢?很多時候我們人是這樣的.
是會歡喜自己做一個神像,自己做一個神像的.
自己做一個神像有什麼好處呢?.
自己做一個神像,你會覺得容易溝通一點.
自己做一個神像,你會覺得他會比較聽你的.
自己做一個神像,是比較適合自己的口型.
很有趣的,就是這樣,因為我自己做的.
於是我覺得他是會聽我的.
是我自己的心理作用,就是自己投射一個神出來.
這個神是我很服的,因為他跟我很相似.
又接受到,很知道是這樣了.
所以你看到本來中國沒有什麼鬼故事.
就是夜晚的鬼,中國的鬼故事.
魏晉南北朝之後開始就有鬼故事出現了.

$^{121}$當日你看回魏晉南北朝時候的鬼故事.
或者到今天我們心裡想的鬼故事.
那些鬼都跟我們人很相似的.
你可以用到人的個性去解釋他的.
用到人的感情去解釋他的.
因為是我心裡的投射來的.
這些人是跟摩西和帶他們出埃及的耶和華神.
哇,厲害嗎?超厲害!.
但是呢,是厲害的,不過很高啊.
有少少惡啊,有少少就是很厲害啊.
心裡有時候覺得很奇怪,是很難跟下的.
是令到埃及軍兵死在海裡面.
但是他帶我們走的路都不是朝著迦南地去的.
屈了下來南邊來曠野的,這麼有趣的.
又不敢問,現在摩西不在就好了.
就去找二哥,二哥就是第二位的領導者.
當然阿倫是摩西的哥.
找第二位領袖,你可不可以跟我們做點事啊.
其實有少少逼他,摩西不見了,那怎麼辦?.
起來做點事吧,其實是壓力來的.
我們有時候容易妥協,就是因為有壓力.
或者不習慣有壓力.
那誰習慣有壓力呢?摩西就習慣有壓力.
是不是帶我們出來死在海邊啊?.
埃及沒有墳墓啊?.
是不是想我們餓死在曠野啊?.
是不是想我們渴死在曠野啊?.
你這麼陰毒啊,帶我們出來這裡死啊.
埃及沒有嗎?.
埃及沒有墳墓啊,我們這麼餓,吃完都沒東西吃.
我們在埃及地的時候,坐在浴窩旁邊,多厲害啊.
那班人是很喜歡發怨言給壓力的.
摩西呢,頂得住習慣有壓力.
阿倫呢,可能在這方面就沒那麼強力了.
不是很頂得住.
為我們做一個神像,為我們做一個神像.
大家怎麼要留意呢?有很多時候我們都會為自己做一個神像.
會為自己做一個領袖.
見到某一個人在做領袖的時候.
無論教會的牧師,教會的導師,教會的長老.

$^{161}$或者是我們公司的老闆,或者我們的組長等等.
甚至我爸爸.
我們都會將一個偶像,將一個領袖的.
我心目中的領袖呢,是投射在他身上.
希望在他身上見到的.
沒有啊,他不是這樣啊,很失望.
做領袖,做爸爸,做組長,做經理,應該是要這樣的嘛.
為什麼你不是這樣啊?這樣就會生氣的.
其實是自己的要求投射.
是好像我們為自己做了一個神像.
心目中虛擬的,就要打在別人身上.
希望是這樣的.
喂,我們做一個神像吧,那個魔性我們不想跟.
那阿倫竟然間妥協.
就說這樣啊,做一個神像.
那好吧,你們回去就找你們妻子兒女.
以上的金環拿來給我.
這個呢,表示那些人其實又不是說.
真的離棄耶和華的.
只不過他們是不懂的.
不過認真,你試一下回去.
問你太太,拿光那些黃金首飾給我.
我要拿出去用了它,大家做一個金牛像.
你試一下你太太肯不肯啊.
不是吧,拿光金耳環,拿兒子的耳環.
拿兒子的戒指,不是吧,拿女兒的.
那你就拿你自己的好了.
其實不是很容易的.
那你回去認真地拿光那些金環.
那些金環有些是穿耳,有些是手.
有些是穿鼻的.
百姓都摘下他們以上的金環給阿倫.
阿倫就將那些金器摺過來.
鑄了一隻牛獨.
問題就在這裡了.
牛獨第二點呢.
鑄了牛獨,牛獨是什麼呢.
就是他們在埃及地最多人拜的神.
就是牛獨,而是用金做的.
嘩,牛啊.

$^{201}$多數做出來的牛會不會小小一隻呢.
不會的,牛做出來當然要大隻.
恐怖有力,那個像都很大.
要很多金的.
做一隻金牛.
因為在阿倫來說.
都在埃及生活了幾十年.
一想起要做神像.
他馬上的形象是什麼呢.
金牛,是金做的牛.
就是神.
叫大家做一隻金,拿走金來.
鑄一隻金,一隻牛獨.
然後叫那些工匠雕刻成為漂亮的.
不是隨便一塊一塊的.
不是的.
然後他就說.
以色列,這是領你們出埃及地的神.
那個神見不到.
摩西山上跟他聊天.
又見不到.
他都是出現在雲柱火柱.
有些聲音跟摩西講的時候.
我們聽到有時是摩西透過.
摩西跟我們講這樣.
雲柱火柱我見到.
嗎哪妹又見到.
鵪鶉又食過.
埃及人死在海裡面都見到.
其實那個神從來沒有顯現過給他看什麼樣.
但由於他們要具體實在.
所以心裡沒有想到別的東西.
你要知道.
人是不能夠創作.
是想到一樣他從來沒有見過的東西.
因此他們想起神是什麼呢.
金牛毒.
一隻金牛.
於是就做了.
然後就是.

$^{241}$這個就是領你們出埃及地的神.
這樣都講得出.
這裡阿倫是知道他們的心意.
是做什麼好.
做一隻金駱駝.
金的馬行不行.
金的鷹.
金牛.
因為金牛就是他們見得最多的神.
是有力量的.
他就說這隻金牛就是領他們出埃及的神.
然後阿倫看見之後就在牛毒面前捉壇.
第五節就做了一件事.
令到我們都挺有趣的.
阿倫接過來.
第四節說.
那些金然後注尿一隻牛毒.
注尿成了一隻雕刻號.
即是說阿倫一直都知道.
這個就是領你們出埃及的神.
阿倫看見就在牛毒面前捉壇.
好像是阿倫現在才看到是一隻牛毒.
不是的.
其實他一直看.
做好了.
然後就在牛面前捉壇.
捉壇做什麼.
要來獻祭的.
那獻祭怎樣獻.
向耶和華守節.
然後在阿倫宣告.
明日終於現在做了一隻牛.
明日要向耶和華守節.
因此阿倫和百姓都覺得這隻金牛就是耶和華.
最大的問題就死在這裡.
這隻金牛就是耶和華.
所以向耶和華守節.
而且他們很認真的.
因為拿了自己的真金白銀出來.
明日向耶和華守節.

$^{281}$於是他們清晨起床.
昨晚都不知道有沒有睡.
今天起床就向著耶和華.
這隻金牛獻繁祭和平安祭.
很認真.
還要獻祭.
即是說還要拿出牛羊出來.
他們獻完繁祭和平安祭.
就坐下吃殼起來還傻.
大家要知道.
整個以色列這樣獻祭.
那個繁祭和平安祭不是一次.
可能很多人都拿.
譬如十二支派.
十二支派每一個支派就拿些牛羊出來.
這個可能拿五隻.
那個拿十隻.
我自己都有錢.
我很想都獻一隻.
其實很有可能那個繁祭.
獻的繁祭是一百幾十隻.
不奇怪的.
如果你看後面這個散步耳機上下的時候.
獻那些牛羊就七十隻.
什麼幾十隻.
很厲害.
很大堆頭.
獻平安祭,繁祭.
然後他們是坐下吃殼.
吃殼製生的肉和喝酒.
為什麼不喝水喝酒呢?.
他們是很少喝水的.
除非在很乾旱的時候.
如果在以色列地.
他們的水就更加不是很認真乾淨.
因為他們儲了很久的味道.
他們要混一些酒進去.
喝酒是他們很平常會做的.
不過正是說不要喝水.
他們坐下吃殼.

$^{321}$很有可能.
就這樣看.
關聯性.
他們吃殼那些祭物.
即是燒了.
燒了多少是吃的.
起來橫灑.
這個橫灑是玩.
其實有多少是慶祝.
很開心的慶祝.
這次如果再讀下去.
就是什麼呢?.
神就跟摩西說.
你下去吧.
那班百姓已經敗壞了.
做什麼事?.
他們在拜金牛.
不是吧?.
於是摩西下去.
就將兩塊法板.
扔在那些人身上.
將金牛磨了粉.
然後沖水叫大家喝.
讓你記得耶和華神是忌邪的神.
不可以用任何東西代替耶和華.
這個故事是一個領袖妥協的故事.
妥協什麼呢?.
是在屬靈的事上妥協了.
令到以色列陷在罪內.
我覺得如果亞倫在那時候站著.
不行,不可以這樣獻.
我們要等雲柱火柱動.
就繼續帶我們走.
就算摩西不下來.
我們都可以走.
跟著雲柱火柱就可以了.
跟著神的帶領就可以了.
如果這樣他堅持的話.
那就好了.
摩西也快要下來了.

$^{361}$也不會整個以色列陷在罪內.
神是發了很大的怒氣.
差不多想以色列人.
不是以色列人.
當時叫做我的百姓.
全部殺了.
然後我叫摩西的後裔成為大族.
不過摩西就說.
神啊,這樣是不可以的.
你要紀念你和阿巴拉罕.
以撒雅各的立約.
特別是阿巴拉罕的立約.
所以你一定要照顧這班人.
算了,你赦免他們.
不要找我來搞.
摩西挺厲害的.
今天就講到這裡.
其實就發覺一件事.
第一,人是會將自己心目中.
對於領袖甚至神的投射.
都會投射出來.
然後要求對方.
都好像這樣跟著我的心意去成就.
是我們人的弱點.
我們做基督徒千萬不要這樣.
神應該是這樣的.
應該是祝福我的.
這些時間應該是這樣的.
要治好我的病.
要讓我兒子進到那間學校.
要叫那個人嫁我.
要這樣的.
要給我多少錢.
要給我一份好工.
要這樣的.
教會要好的.
投射就會出來.
有時候我們個人.
教會的領袖.
特別教會的領袖.

$^{401}$千萬不要像這個負面的案例.
就是亞倫妥協了.
用了埃及的神.
就說這個就是耶和華.
這個是令到神很憤怒的.
不妥協也對.
我們的神是忌邪的.
只有一個神.
正如忌邪的意思就是.
你的爸爸永遠不會容許你.
叫別人做爸爸.
不要緊啦.
我們中國人說.
有鬍子就是爸爸.
不可以的.
爸爸只有一個.
就是我.
你叫別人做爸爸我就生氣了.
你叫別人做媽媽.
真正的媽媽就生氣了.
不一樣的.
你叫別人做真神.
那個真神就會生氣.
你會知道是會這樣的.
而且很嚴重.
雖然他們的內心虔誠.
要向耶和華守節.
不過走錯了路.
由於亞倫的妥協.
而導致百姓走錯路.
如果亞倫死守.
不理他.
不肯講清楚那個道理.
結果那些人連亞倫都打了.
然後自己做金牛毒.
那也沒辦法.
神是會對這班百姓做事的.
而作為領袖.
是要守得很穩的.
求主幫助我們.

$^{441}$我們一起祈禱吧.
天護上帝.
我們在這個世界一直生活.
學習了很多妥協的技巧.
不過原來有些範圍.
在真理的範圍.
不能妥協.
我們是要用生命去守的.
主求我們有勇氣與膽量.
主我們看到亞倫那一次.
守不住妥協了.
維吾爾人開了一條路.
他們拜了一隻金牛來代替耶和華.
雖然多麼虔誠.
多麼賣力.
用了真金耳環來容祂去做.
何時還雕刻好.
然後獻上梵祭和平安祭.
但是是錯的就是錯的.
主求你讓我們很小心你的心意.
讓我們很懂得.
要跟著你的道理.
你的心意去走.
不要將自己心裡的東西投射出來.
要求別人這樣做.
要求上帝這樣做.
讓我們謙卑敬拜你.
除非我們祈禱奉耶穌基督命求.
阿們.
\newpage



\section{約翰福音 7:1-13}
\label{sec:032083}
\textbf{榮耀君王:兄弟情深}~~~~林鍵牧師
\newline
\newline
連結: \href{http://www.church.com.hk/acms/content.asp?site=cdc\&op=show\&type=product\&code=032083\&layout=sermon}{\texttt{032083}} ~~~~ 語音日期: 2026-05-02
\newline
\newline
\hyperref[sec:032082]{\small{< < < PREV SERMON < < <}}
~
\hyperref[sec:index_chronic]{\small{[返順時目]}}
~
\hyperref[sec:032084]{\small{> > > NEXT SERMON > > >}}
\newline
\newline
約翰福音 7:1-13
\newline
\begin{longtable}{cl}
\hline
\hline
章節 & 經文 (和合本修訂版)\\
\hline
7:1 & \begin{tabularx}{0.7\textwidth}{X} 這些事以後，耶穌周遊加利利，不願在猶太往來，因為猶太人想要殺他。 \end{tabularx} \\ \\ \relax
7:2 & \begin{tabularx}{0.7\textwidth}{X} 這時猶太人的住棚節近了。 \end{tabularx} \\ \\ \relax
7:3 & \begin{tabularx}{0.7\textwidth}{X} 耶穌的兄弟們對他說：「你離開這裡上猶太去吧，好讓你的門徒也看見你所做的事。 \end{tabularx} \\ \\ \relax
7:4 & \begin{tabularx}{0.7\textwidth}{X} 因為人要揚名，沒有在隱祕的地方行事的，如果你要做這些事，該把自己顯明給世人看。」 \end{tabularx} \\ \\ \relax
7:5 & \begin{tabularx}{0.7\textwidth}{X} 原來連他的兄弟們也不信他。 \end{tabularx} \\ \\ \relax
7:6 & \begin{tabularx}{0.7\textwidth}{X} 於是耶穌對他們說：「我的時機還沒有到，你們的時機卻隨時都有。 \end{tabularx} \\ \\ \relax
7:7 & \begin{tabularx}{0.7\textwidth}{X} 世人不會恨你們，卻是恨我，因為我指證他們的行為是惡的。 \end{tabularx} \\ \\ \relax
7:8 & \begin{tabularx}{0.7\textwidth}{X} 你們上去過節吧！我現在不上去過這節，因為我的時機還沒有成熟。」 \end{tabularx} \\ \\ \relax
7:9 & \begin{tabularx}{0.7\textwidth}{X} 耶穌說了這些話，仍然留在加利利。 \end{tabularx} \\ \\ \relax
7:10 & \begin{tabularx}{0.7\textwidth}{X} 但他的兄弟們上去過節以後，他也上去，不是公開去，卻似乎是祕密地去的。 \end{tabularx} \\ \\ \relax
7:11 & \begin{tabularx}{0.7\textwidth}{X} 節期間，猶太人尋找耶穌，說：「他在哪裡？」 \end{tabularx} \\ \\ \relax
7:12 & \begin{tabularx}{0.7\textwidth}{X} 人群中有許多人對他議論紛紛，另有的說：「他是好人。」有的說：「不，他是迷惑群眾的。」 \end{tabularx} \\ \\ \relax
7:13 & \begin{tabularx}{0.7\textwidth}{X} 可是沒有人公開談論他，因為他們怕猶太人。 \end{tabularx} \\ \\
[1ex]
\hline
\hline
\end{longtable}
$^{1}$歡迎大家來到我們的周到崇拜.
上年七月我們開始研讀新約聖經中的約翰福音.
我們說到系列的題目是「榮耀君王」.
我們花了半年的時間去完成了約翰福音第一章至第四章的宣講.
接著今年的一月份我們開始研讀約翰福音第五章.
約翰福音第五章開始的時候.
耶穌來到耶路撒冷在靠近羊門的不是大池那裡.
醫治了一個騎腿了三十八年的人.
然而你必須要知道這個不是一個普通的神蹟.
而是一個十字架救贖的記號.
故事當中一開始就會有上帝的羔羊.
耶穌是上帝的羔羊來到聖殿的羊門.
將要被獻上通過羊門的羊其實是要送去聖殿那裡作為贖罪祭.
圖畫裡面又有一個永生的活傳在那裡出現.
這個永生的活傳要拯救要等待著這個將要被人救贖的人.
在眾多的病人裡面耶穌選了一個騎腿了正正好三十八年的人.
不是因為這個人有信心.
又不是因為他特別可憐而是因為三十八年.
剛剛好就是以色列人出埃及之後.
叛逆不信上帝流浪曠野.
沒辦法進入上帝安息的年數.
於是耶穌其實是在畫一幅屬靈的圖畫.
祂說我來是要呼喚救贖復活帶領靈裡面叛逆癱瘓罪孽深重的以色列人.
因為三十八代表了以色列人.
這群以色列人他們需要上帝的拯救.
他們騎腿走不入安息.
我要使他們復活.
我要使他們得一個新生命.
接著我要帶他們進入上帝永恆的安息.
耶穌說當年.
大家都記得以色列人癱瘓不能進入上帝迦南地的安息.
是因為埃及的偶像.
於是耶穌對這個人說不要再犯罪.
其實意思是不要再回頭.
不要再貪戀埃及的偶像.
你卻要愛慕堅守上帝的安息.
要奔跑朝向永恆的天國.
不要貪戀埃及的偶像.
要向前奔跑朝向上帝永恆的天國.
接著耶穌又說.

$^{41}$這是聖父和聖子共同努力持續工作要達成的救贖計劃.
你開始明白耶穌先做了神蹟.
接著他要解釋神蹟.
其實後面的信蹟比前面的神蹟重要得太多.
然而耶穌的神蹟和解說並沒有受到猶太人的接納.
猶太人就是猶太人的領袖.
反而忍耐他們憤怒的逼迫.
並且要殺害他.
我們說耶穌的神蹟做神蹟容易.
但是傳福音很艱難.
一開口不是他夾著你.
一開口就被人罰他.
那些人要殺他.
接著耶穌就有機會自變.
他開始一個很長很長的獨白.
很多時候我們誤會以為耶穌的獨白.
其實是在獨罵其他要去審判他的人.
其實不是.
其實在獨白裡面.
耶穌給了很多提示讓他們知道救贖的道路在哪裡.
因為上帝愛世人.
因為上帝愛三八.
你記得那時候把約翰福音第三章數字改了.
三八代表了以色列因為上帝愛三八.
甚至將他的獨生子次級丹們.
叫一切相信的不致滅亡.
反德永生.
因為上帝差遣他的兒子到三八.
以色列人那裡來.
不是要定三八的罪.
而是要使三八因信他而得救.
於是在那個很長的審判當中.
耶穌基本上是被他們抓住了.
接著要受審判.
耶穌給了一個很長的獨白.
由第五章的十九節一直到第五章的四十七節.
總共二十九節的獨白.
當我們仔細去讀那個獨白的時候.
我們驚訝發現每一段都是慰著以色列人的好處.
其實是在給他們很多的引導.

$^{81}$叫他們一直跟著他的腳蹤.
有一天他們一定可以得到拯救.
這個是約翰福音第五章的訊息.
我怎樣去總結它呢.
不是地上的神蹟.
不要將故事變成地上的神蹟.
乃是天國的記號.
你一想地上的神蹟.
我也有差不多的病.
耶穌可不可以都搞定他呢.
我其實沒有他那麼嚴重.
沒有三十八年的.
他可以搞定的.
你不可以將他讀做一個地上的神蹟.
這個不是一個地上的神蹟.
是天國的記號.
耶穌所做的一切.
約翰福音所寫的一切.
不是叫我們眼看地上的神蹟.
而是要耳聽天國的教導.
不是用眼的.
天國是用耳朵.
去聆聽天國的教導.
這是第五章的訊息.
接著是第六章的訊息.
一進入第六章.
我們就遇到.
最為人所知.
我們最喜歡的五餅二魚.
餵飽五千人的天國記號.
不是神蹟.
是記號.
整個五餅二魚的中心.
就是耶穌分餅那一幕.
起餅來足謝了.
默害那一幕.
那一幕其實是最後晚餐.
十字架寫命的寓言.
又不是一個普通的神蹟.
是一個天國的記號.

$^{121}$是一個.
救贖的記號.
但是很可悲.
也很可惡.
當天有五千個男丁.
加起來可能有二萬個男女老幼.
全部都是在那裡看神蹟.
不要聽耶穌的話.
大家就去吃餅.
然後說這個人很厲害.
現在可以做皇帝.
於是他們要耶穌作王.
他們只是看神蹟.
他們不聽耶穌的話.
猶太的群眾是這樣.
耶穌的門徒都是這樣.
在故事裡面.
如果你記得菲力和安德烈.
合力對抗不聽耶穌的話.
耶穌要行神蹟.
去餵餅.
他們堅決不要.
正如我之前所說.
五餅二魚很可能是耶穌.
一生裡面眾多神蹟裡面.
最熱鬧但又最孤單的一個.
因為只有祂知道祂在做什麼.
其他人沒有一個人.
接到祂的招.
沒有人接到祂的球.
每個人都以為祂在行神蹟.
但其實祂在給我們.
一個記號.
回到楊寶欽第六章.
五餅二魚的.
天國記號之下.
又是另外一個天國的記號.
短一點的.
小一點的.
但其實兩個是雙胞胎.

$^{161}$下面有一個天國的記號叫.
黑夜狂風海上行走.
耶穌在海上行走.
驟眼看來.
這兩個天國的記號.
九五八搭的.
上下.
好像沒有什麼關係.
但當你仔細研讀的時候.
我們發現原來五餅二魚.
和海上行走.
是一個龍鳳胎.
正所謂好事成雙.
五餅二魚的結尾是十二個男子.
相信是十二個門徒.
拿著十二個男子.
海上行走的起頭.
也是十二個門徒.
海上行走在楊寶欽裡面.
唯一一件行在.
十二門徒身上的神蹟.
沒有其他人的.
只有十二門徒見到經歷過這個神蹟.
他們是耶穌的門徒.
是耶穌為他們度身訂造的神蹟記號.
為要成為他們將來的鼓勵和應許.
為什麼?.
因為這班十二門徒.
他們很快會被耶穌.
差遣into darkness.
進入一個黑暗的世界裡面.
將他們的餅碎.
餅碎代表了天國的福音.
餅已經餵飽了五千人.
餵了二萬人.
但是還有餅碎.
餅碎可以繼續去救人.
於是將餅碎生命的糧.
傳到以色列的十二支派裡面.
再傳到地極.

$^{201}$而當你去做這個行動的時候.
當你做這個使命的時候.
是into darkness.
光照在黑暗的當中.
黑暗會拒絕光.
於是十二門徒極其需要耶穌的同在.
而通過海上行走的神蹟.
或者記號.
耶穌告訴他們.
他必聽天語.
他們同在直到世代的終結.
但是有一個條件.
他們要忠心耶穌天國的使命.
如果大家記得.
在這兩個天國記號.
做完之後.
約翰帶我們進入耶穌.
在加伯農會堂裡面的講道.
那篇道有一個名堂叫.
我是生命量.
那篇道其實是為了要解明.
五病二瘀究竟是什麼事.
你明白嗎?.
五病二瘀是小事.
那篇道是大大大大大事.
耶穌做了五病二瘀.
然後他要去將那篇道.
去解釋五病二瘀.
那篇道如果大家記得.
裡面的內容非常奇特.
講一下又講什麼吃我的肉.
喝我的血.
在當日在場的聽眾裡面.
我相信大部分.
沒有人聽得懂耶穌的信息.
然而.
只要他們繼續跟隨.
只需要一年之後.
你繼續留意耶穌.
繼續追蹤他的時候.

$^{241}$你會明白那篇道裡面.
所有的一切.
都是指向耶穌基督的.
十字架.
在那天耶穌講那篇道.
教導眾人的時候.
他也鼓勵大家.
要全人回應.
十字架的救贖.
這個是耶穌給我們的.
形容.
六章第二十六節.
耶穌回答他們說.
我實實在在地告訴你們.
不要為那毀壞的.
食物操勞.
卻要為那.
傳到永生的食物操勞.
這食物是人子.
要賜給你們的.
那食物是誰賜的?是人子.
人子是什麼料來的?.
人子就是耶穌.
但以理書第七章.
十三至十四節.
人子是說舊約聖經裡面的人子.
人子在舊約聖經裡面.
做什麼的?領受一個永遠的國度.
所以人子.
要賜給你的不是麵包.
這麼簡單.
人子要賜給你的是天上.
永恆的財寶.
永恆的天國.
所以耶穌的形容是不要為那毀壞.
地上的國度.
今生的榮耀而操勞.
耶穌叫我們要為天上的財寶.
要為永恆的榮耀.
永恆的國度.

$^{281}$而操勞.
然後耶穌告訴我們.
我就是人子.
我會為你上十字架.
通過十字架我會為你贏取.
天國的財寶.
永恆的國度.
你既然將會從我手上.
得著永恆的國度.
你就應該.
一無掛慮.
奉上所有.
單求我的國.
將你的餘生.
全部獻給我.
與我一起追求.
那個天國天上的財寶.
這樣就對了.
這是《讓福音》第六章的訊息.
接下來到今天.
今天我要帶你進入.
《讓福音》第七章.
我們只是看頭十三節.
有關耶穌.
與他的兄弟的故事.
這是四福音書.
獨有的耶穌與.
兄弟幫的故事.
在我們進入第七章之前.
我想帶你回顧.
《讓福音》第五章至第十章.
的大圖畫.
聖學者告訴我們.
在《讓福音》第五章.
到《讓福音》.
第十章.
我們進入了一個.
節慶的週期.
Festival Cycle.
第五章有一個節日.

$^{321}$叫做安息日.
第六章有一個節日.
背景是愉悅節.
第七至第九章.
是鑄棚節.
第十章是獻殿節.
最後到十一章.
又回到愉悅節.
在這個節慶週期當中.
耶穌通過神蹟.
和教導以節日為背景.
節日是一個背景.
在節日當中.
耶穌通過神蹟和教導.
神蹟和教導是啟示自己.
耶穌希望大家知道.
祂是誰.
祂的身份.
祂的使命.
所以不要亂來.
耶穌救救我.
你搞錯了.
你聽聽耶穌說話.
你聽聽耶穌的身份.
祂的使命.
安靜聆聽.
究竟你需要問耶穌.
拿些什麼.
耶穌用這些天國的記號.
祂的教導啟示.
祂的身份和使命.
很可惜很有趣.
就是祂越是清楚啟示自己.
猶太領袖就越是.
恨惡去逼迫他.
越清就越不清.
不是不清.
就是清了才越來越恨惡耶穌.
於是去到第十一章第五十七節.
當然我們沒有去到那裡.

$^{361}$但是你去到就知道了.
祭司長和法利賽人早已下令.
若有人知道耶穌的下落.
就要報告.
他們好去捉拿他.
要將他送上十字架.
這個是十一章最後的時候.
告訴你他們已經立定志向.
要殺死耶穌.
但是所有的一切其實從第五章已經開始了.
在第五章的.
黑色暗色人.
Black Sabbath.
為什麼呢?因為他們想殺人.
第五章的黑色暗色人裡面.
猶太人已經.
是立定志向要殺耶穌.
第五章第十八節.
猶太人越想殺他.
但因為他.
因為他不但犯了.
黑色暗色人而且稱上帝.
為他的父.
把自己和上帝.
看為同等.
於是由第五章開始.
每一章每一節.
每一個故事每一個段落.
都加深了猶太領袖.
對耶穌的仇恨和妒忌.
以致最後.
引發最後的一個爆炸.
就是在耶路撒冷裡面.
憤怒的呼喊.
Crucify Him! Crucify Him!.
大家都有一些印象.
有時候教會裡面去到受難節會演一個話劇.
當然要有一群人在那裡說.
Crucify Him! Crucify Him!.
所有一切帶我們去到最後.

$^{401}$憤怒的爆炸.
耶穌必須要死.
我想帶你回去看.
第五第六第七章的起頭.
你會發覺.
他們的起頭是完全一樣的.
這是以後.
第五章第一節.
到遙猶太人的一個節期.
耶穌上耶路撒冷去.
沒有告訴我們是什麼節期.
後來我們就知道他是一個暗色的.
但我們不知道那個節期是什麼.
第六章第一節.
這是以後耶穌渡過加利利海.
就是提比利亞海.
有一個節期叫逾越節.
第七章.
這是以後耶穌周遊加利利.
不願在猶太往來.
因為猶太人想要殺他.
第七章第二節.
告訴我們當時.
是鑄棚節.
近了耶穌會再次.
回到耶路撒冷.
第六章和第七章之間.
大概是半年的時間.
由春天到秋天.
春天的逾越節.
到秋天的鑄棚節.
所以我們知道第六節第七章.
中間是隔了大概半年的時間.
第五章和第七章之間.
就不是很肯定.
但一般人相信要不就一年.
要不就一年半的時間.
都比較長的時候.
然而第五章和第七章.
是前後非常連貫的.

$^{441}$你會發現.
第五章的數字這樣說.
所以猶太人迫害耶穌.
他做了一個神蹟.
三十八年騎去的人.
他們迫害耶穌.
因為耶穌再安息的做了這事.
猶太人想要殺他.
因為他稱上帝.
為他的父.
把自己和上帝看為同等.
去到第七章第一節.
我們在說的是.
一年半之後的事.
以後耶穌周遊列國不願再猶太來往.
他在加里尼亞在北部.
不去南邊.
一年半之後.
仍然想殺他.
這個所說的猶太人不是所有猶太人.
是猶太人的領袖.
他們對耶穌念念不忘.
寵愛有加.
晝夜思想.
不住等候.
到他一年半之後.
他再回到聖城耶路撒冷的時候.
他們仍然想要殺他.
都殺不到.
要再等多一年.
然後才殺到.
或者等多半年.
到第二年的逾越節.
他們才可以殺到.
如果你去看《讓福音》第七章.
從故事的時序裡面.
你會發現故事可以分為三部分.
第一個部分是.
第七章第二節.
猶太人的絕行接近了.

$^{481}$第二部分是第十四節.
節期已過了一半.
第三部分是第七章第三十七節.
至其最後的一天.
也是最隆重的一天.
那才是耶穌.
做天國教導自我啟示的高峰.
在那一天.
耶穌說了兩件事.
我就是生命的活水.
我就是世界的光.
大家有印象嗎?.
第一個我是生命的活水.
其實是我自己加插的.
因為是我幫他做的總結.
他沒有說過那句說話.
但是他的意思是在說那句說話.
我讀給你聽.
《約翰福音》第七章第三十七節.
就在最後的一天.
七天裡的最後一天.
節期裡的最後一天.
最隆重的一天.
耶穌站起來.
喊著說:人若渴了到我這裡來.
記不記得?.
渴了到我這裡來,他是誰?.
生命的活水,信我的人可以來.
就如經上所說.
從他的福中會流出活水的江河.
耶穌說我是誰?.
我是生命的活泉.
耶穌說:是指信他的人將要受聖靈說的.
那時還沒有賜下聖靈.
因為耶穌還沒有得榮耀.
即是未上十字架.
這是將來的事.
以後我們會有機會回來看這段聖經.
接著在同一天.
雖然已經去到第八章第九章.

$^{521}$耶穌再說兩次.
我是世界的光.
接著他做了一個神蹟.
他醫治了一個從一開始就已經黑眼的人.
將黑眼人.
大盡光明.
約翰好像在這個鑄棚節.
所有這一切,第七章,第八章,第九章.
其實都是鑄棚節的背景裡面.
發生.
約翰好像在這個鑄棚節的背景裡面.
在建築兩座.
耶穌自我啟示的金字塔.
你不要想歪了.
金字塔一定是壞事.
我現在只是打一個比喻.
耶穌在建兩個金字塔.
一個叫做我是生命的活水.
另一個叫做我是世界的光.
他.
將他.
高峰的時候將他自己的身份.
他的使命.
展示出來.
不過我們在說的是將來的事.
因為今天我們沒有拿到.
是以後的講道.
今天我們在做的是打地基.
今天我們在做的是畫背景.
因為我們還是第一節.
至第十三節.
今天的經文是第七章.
一章一指十三節.
十三節的聖經.
非常有趣精心編排.
呈現一個ABCD.
CBA前後呼應的結構.
我們看看最外層.
外層的A.
上面的A.

$^{561}$猶太人.
耶穌恩著他們.
恩著他們.
不願往來.
下去猶太的地方.
下面猶太人.
恩著他們不公開談論耶穌.
兩個都是和猶太人.
猶太領袖有關係.
中間的BB.
兄弟們想猶太去顯明世人.
B.
下面許多人.
他是好人有人說他是好人.
也有人說他是迷惑眾人.
是其他人的不是那些主角.
接著.
CC有一個敘述者.
Narrator就是一個講故事的人.
第一句他講的是.
耶穌的兄弟並不信他.
第二句他講的是.
耶穌的兄弟也不和耶穌同路.
他們在走兩條不同的路.
我們讀到聖經.
我解釋給你聽你就明白.
在中間整個故事的核心.
耶穌講了一句說話.
我的時候.
我的時機.
還會有到.
你熟悉仰福音的時候.
不是時候嗎.
不是我的時候嗎.
不是時機嗎.
原來這裡的原文.
Kairos是時機.
只是這一節出現.
他就和另外一個.
約翰很喜歡用的時候.

$^{601}$Aura.
The hour.
是不同的在原文上是不一樣的.
但是這兩個名詞.
其實都和時間有關係.
我會說他們兩個名詞.
其實都是相連.
都是貫通的.
所以這裡的Kairos.
他的時機.
最後會引領他前往.
來到耶穌的Aura.
耶穌的時候.
The time will become the hour.
這裡講的是The time時機.
後來講的是The hour.
時候.
這個時機會引他去到.
那個時候.
那個時候就是他被.
高舉在十字架上的.
那一刻.
無論你怎麼看.
所有故事都是指向.
耶穌基督的十字架.
我想你想想.
在這一切發生的同時.
耶穌正在上十字架.
耶穌的兄弟們.
是不相信他的.
耶穌的兄弟們.
是完全不知道發生了什麼事.
耶穌的兄弟們和耶穌.
是完全不同路的.
但是耶穌的兄弟們.
最後有沒有被耶穌拯救.
我先.
馬上告訴你.
就是有的.
這一切令我想起.

$^{641}$羅馬書第五章.
羅馬書第五章第六至八節.
這麼說.
我們還軟弱無力的時候.
基督就在特定的時刻.
為不敬虔的人死.
耶穌上十字架的時候.
不是因為已經有一群人在招手.
我相信,我會轉.
是當我們還軟弱無力.
所謂軟弱無力就是.
當我們還不相信的時候.
當我們還死在最終的時候.
當我們不是跟隨耶穌的時候.
當我們不順服上帝的時候.
當我們還不敬虔的時候.
耶穌正在上十字架.
如果耶穌是需要.
有人在招手.
歡迎他救我們.
耶穌是不會上十字架的.
耶穌的弟兄沒有一個向他招手.
稍後看下去.
他不但不向他招手.
還單單打壓他.
然而.
基督就在這特定的時刻.
上帝的時刻.
不是人的時刻.
為罪人死.
而我說的不是普通的罪人.
我說的是罪大惡極.
叛逆不悔的大罪人.
我們這裡每個人都是.
我們很難想像.
在我們的世界裡面.
有些是大逆不道的人.
有些是罪大惡極的人.
在上帝眼中每個人都是一樣的.
每個人都是吃蛋.

$^{681}$每個人都是負無限分.
沒有大沒有小.
每個人都是小.
每個人都是無限的負分.
保羅告訴我們.
當我們沒有人迴轉.
當我們沒有人追求上帝的時候.
上帝為我們預備了一個無價的救贖.
在第五章.
第七節裡面.
他插了一個人世間的比喻.
他很強調人世間的比喻.
是為要.
對比顯大耶穌基督.
這個無法想像.
不可理解的大愛.
保羅這樣說.
為異人死的是少有的.
為人人死的.
或者有感做.
什麼叫異人呢?異人就是.
正直公義的人.
道德高尚的人.
什麼叫人人呢?人人就是.
他自己奉上了自己.
造福了百姓的人.
一個是自己.
非常之道德高尚.
另一個不單止自己道德高尚.
還獻上自己.
去造福其他人.
保羅說.
對於那些造福眾人.
犧牲自己的人.
或者都有人敢為他們死.
為一個道德高尚的人.
是少有會為他死的.
當然我想說的就是.
這兩種人.
在上帝眼光是不存在的.

$^{721}$你不要想太多.
他是在說地上的比喻.
他說我們在地上我們要知道.
這些人難得有.
也很難得有人會為他死.
但我想說的是.
下面那節唯有基督在我們.
還作罪人的時候為我們死.
上帝的愛在此就向我們顯明了.
耶穌不是.
等我們做了異人.
做了人人.
才為我們死.
是我們全部都在罪人.
在上帝眼中其實是沒有異人.
沒有人人.
上帝獻上他自己.
獨生的兒子為要拯救罪人.
而且是大大大大.
大大的.
罪人.
我想你記得這個故事.
我也想你記得.
耶穌今天.
要去拯救的弟弟.
耶穌自己的弟弟.
全部都是這些人.
都是不去敬畏上帝.
軟弱無力.
死在最終的時候.
但是耶穌為他們.
上十字架.
但我們進入到楊福音第七章.
楊福音第七章.
開始說這個故事.
這時猶太人的.
柱彭節近了.
柱彭節是舊約聖經裡面猶太人.
三大的節期之一.
在耶穌的時代裡面.

$^{761}$特別的南丁.
每年要求.
理論上是要去耶路撒冷三次.
去紀念.
慶祝三個節期.
第一個是逾越節.
如月的羔羊.
第二個是逾越節之後.
四十九日的五旬節.
跟著還有另外一個.
不是春天.
而是秋天裡面的柱彭節.
所以柱彭節是三大節期之一.
甚至有人說柱彭節.
特別多事要做.
特別多事要看.
柱彭節紀念的.
是以色人出埃及之後.
上帝在曠野四十年裡面.
對他們的看顧.
和保守.
這個節期維持七日.
通常在.
我們的年曆裡面.
九月十月之間.
秋高氣爽非常舒服的時候.
除了柱彭之外.
每日還有水禮和火禮.
於是水禮和火禮就成為了.
耶穌.
我是生命活水.
我是世界之光的宣告背景.
你明白嗎.
耶穌在那裡做的時候.
其實是有特別效果的.
已經有水有光.
於是耶穌說我是世界的光.
我是生命的.
活泉.
第七章第三節.

$^{801}$耶穌的弟兄.
對耶穌說.
他們在那裡.
耶穌的弟兄全部都成人了.
耶穌要.
他們都要去過節.
他們過節之前就這樣跟耶穌說.
他說你離開.
這裡想由大去萬.
所謂想其實是下的.
因為他由北邊去南邊.
但是因為.
耶路撒冷是.
整個地方的最高峰.
屬靈高峰.
所以你個個去耶路撒冷.
無論向南向北都是上的.
他說.
你耶穌.
你離開這裡.
想耶路撒冷去吧.
好讓你的門徒也看見.
你所做的事.
你行的我們信你的.
你會做事的.
因為人要揚名.
沒有在隱蔽的地方行事.
你為甚麼在我們家鄉小地方.
做事.
如果你要做這些事.
該把自己顯明給世人.
現在是最多人.
在耶路撒冷的時候.
因為是轉行節.
你和我們一起去吧.
這裡說.
耶穌的.
兄弟們對他說.
這些兄弟是誰來的.
在演繹的歷史裡.

$^{841}$告訴我們.
有三個可能.
第一.
他們說的是耶穌的表兄弟.
不是在說.
他的弟弟.
第二.
是他的哥哥.
約瑟有前一段的婚姻.
前一段婚姻.
有一群兒子.
於是再和瑪利亞結婚.
於是.
耶穌成為弟弟.
耶穌有幾個哥哥.
你很不明白.
明明說是耶穌的兄弟.
就是瑪利亞生的.
為什麼要說是.
約瑟的前一段婚姻.
原因就是.
在教會歷史裡.
無論是羅馬天主教.
東正教或某些基督教的支排.
都有一個信念.
信念叫.
Perpetual Virginity of Mary.
聖母終身同征.
說的是他們有一個信念.
這是神學的觀念.
是教會的傳統但沒有聖經的證明.
因為瑪利亞在耶穌出世之前.
出世之時.
出世之後.
都是一個同女.
從來沒有性關係.
和任何一個人.
從來沒有和任何人有性關係.
從頭到尾.
她就是一個同女.

$^{881}$於是當你說耶穌有兄弟的時候.
他們一定要想辦法.
這些兄弟從哪裡來.
如果我們將這個.
拋開教會的傳統.
即是聖經沒有講.
來講.
最直接了當的解說.
這些耶穌的兄弟.
就是耶穌自己出世.
他是大哥.
之後約瑟和瑪利亞繼續生的弟弟.
你一會兒看.
他不單止有弟弟.
他還有妹妹.
這是我們所想的.
耶穌的一班弟弟.
我想你想想一個很奇怪.
不是很奇怪.
我自己想的東西.
在一個人的成長過程裡.
如果你有耶穌.
這樣的大哥會是怎樣的呢.
有沒有想過.
好還是不好呢.
耶穌是世界的光.
早上還在睡覺的時候.
世界的光照你了.
照醒你了.
起床擦牙.
穿衣服回教會了.
你有沒有想過.
有耶穌這樣的完美哥哥.
會是怎樣的.
天天照顧你.
照顧你.
講到.
耶穌和他的弟弟的關係.
這個非常微妙的關係.
其實令我想起.

$^{921}$聖經裡面舊約的一個故事.
在創世記裡面.
有一個故事是關於約瑟的.
大家記得.
我讀幾次給你聽.
你就明白了.
約瑟十七歲的時候與哥哥一同牧羊.
約瑟有一個專長.
約瑟把他們的.
惡行報告給父親.
有沒有想過.
每晚.
耶穌都在打小報告.
我那兩個弟弟.
今天做錯了什麼.
然後發電郵給上帝.
天天在打報告.
於是以色列人.
他的爸爸又特別愛約瑟.
為約瑟做了一件.
長袍.
約瑟又發了一個夢.
我的哥哥和父母.
以後都會拜我的.
令到哥哥很恨他.
接著他哥哥.
對他說.
難道你真的要作我們的皇嗎.
難道你真的要統治我們嗎.
因為他的夢.
他所講的話就非常恨他.
有沒有想過.
耶穌是不是這麼討人厭.
你知道我是誰嗎.
你以後就在這裡.
再想.
耶穌比他們年紀大.
就算他們比耶穌小幾年.
你可以想像.
耶穌每天.

$^{961}$長輕為富.
有人說約瑟早死.
那就是說耶穌不是做了爸爸.
管教弟弟的角色.
每天都在吵.
每天耶穌拿著聖經.
在吵他們.
我們什麼時候.
盡心盡性.
盡義愛上帝.
其實就是愛我.
你們豈有此理.
剛才吃飯又什麼.
很嚴厲的警告.
而且.
大型事後.
做爸爸.
你這臭小子.
打你兩巴掌.
打你屁股也好.
有沒有想過.
耶穌和弟弟的關係.
是不是這樣.
其實不是.
我猜不是.
弟弟.
有沒有留意弟弟對耶穌.
這句說話.
很不客氣.
弟弟是在教哥哥.
怎樣做人.
弟弟.
全部都是小朋友.
他是猶太人.
年歲很重要.
弟弟對他完全不客氣.
哥哥你不會做人.
我現在教你怎樣做人.
令我想起耶穌很可能.
有點像電影.

$^{1001}$Rainman.
有沒有看過很舊的電影.
Rainman裡面的Dustin Hoffman.
Dustin Hoffman扮演一個.
智商超高.
耶穌的情況除了智商超高.
也是能力超強.
但是患有自閉症.
有沒有想過.
弟弟這樣對耶穌.
會不會耶穌從小到大.
像一個困在自己的世界.
裡面的自閉症兒童.
他很厲害.
他很有能力.
但是在弟弟的眼中.
耶穌很奇怪.
經常活在自己的世界裡面.
耶穌也明白.
他們的世界和他的世界.
是完全兩個世界.
在他的世界裡面只想著天上的事.
在他們的世界裡面只想著地上的事.
事.
我又想你想一下.
如果你是耶穌的話.
作為一個.
完美的人.
作為一個完美的上帝.
作為一個完美的哥哥.
活在一群很不完美的.
弟弟的世界裡面.
是什麼感受.
我覺得.
耶穌是一個非常孤單的哥哥.
因為他是唯一一個.
不停不斷思想.
天上事的人.
但是他面對的.
起碼這群弟弟.

$^{1041}$是半逆不順的弟弟.
住在一起.
根本對天國事.
完全沒有任何的興趣.
耶穌是一個非常孤單的上帝.
內裡不孤單.
因為他和天父一起.
但是外面.
是完全被孤立的.
面對著罪人.
弟弟.
的半逆和不順.
耶穌彰顯著.
無限的忍耐.
和恩慈.
無盡的憐憫和大愛.
如果你問我.
他的成長就是.
他任由他們繼續做他們喜歡做的事.
但是他偶然.
會跟他們說.
你要小心將來的審判.
但是他又不會常常盯著他們.
這是沒有用的.
但是他的存在就代表了.
他們的前景.
因為他所做出的一切.
都是完美的.
弟弟知道他是完美的.
弟弟知道他們不對.
然而他有他的世界.
他們有他的世界.
耶穌只是在等待.
希望他們有一天會回轉.
但是耶穌要經過那三十年.
是需要無限的.
忍耐和恩慈.
無盡的憐憫和大愛.
我這裡說.
耶穌一對四.

$^{1081}$原因是因為在馬可福音和.
馬太福音裡面.
也給了我們耶穌幾個弟弟的名字.
馬可福音第六章第三節.
馬太福音第六章第三節.
有雅各,約瑟,西門和猶大.
馬太福音第六章第三節.
有雅各,約瑟,西門和猶大.
馬太福音第六章第三節.
有雅各,約瑟,西門和猶大.
馬太福音第六章第三節.
有雅各,約瑟,西門和猶大.
馬太福音第六章第三節.
有雅各,約瑟,西門和猶大.
馬太福音第六章第三節.
有雅各,約瑟,西門和猶大.
馬太福音第六章第三節.
有雅各,約瑟,西門和猶大.
馬太福音第六章第三節.
有雅各,約瑟,西門和猶大.
馬太福音第六章第三節.
有雅各,約瑟,西門和猶大.
馬太福音第六章第三節.
有雅各,約瑟,西門和猶大.
馬太福音第六章第三節.
有雅各,約瑟,西門和猶大.
馬太福音第六章第三節.
有雅各,約瑟,西門和猶大.
馬太福音第六章第三節.
有雅各,約瑟,西門和猶大.
馬太福音第六章第三節.
有雅各,約瑟,西門和猶大.
馬太福音第六章第三節.
有雅各,約瑟,西門和猶大.
馬太福音第六章第三節.
有雅各,約瑟,西門和猶大.
馬太福音第六章第三節.
有雅各,約瑟,西門和猶大.
馬太福音第六章第三節.
有雅各,約瑟,西門和猶大.

$^{1121}$馬太福音第六章第三節.
有雅各,約瑟,西門和猶大.
馬太福音第六章第三節.
有雅各,約瑟,西門和猶大.
馬太福音第六章第三節.
有雅各,約瑟,西門和猶大.
馬太福音第六章第三節.
有雅各,約瑟,西門和猶大.
馬太福音第六章第三節.
有雅各,約瑟,西門和猶大.
馬太福音第六章第三節.
有雅各,約瑟,西門和猶大.
馬太福音第六章第三節.
有雅各,約瑟,西門和猶大.
馬太福音第六章第三節.
有雅各,約瑟,西門和猶大.
馬太福音第六章第三節.
有雅各,約瑟,西門和猶大.
馬太福音第六章第三節.
有雅各,約瑟,西門和猶大.
馬太福音第六章第三節.
有雅各,約瑟,西門和猶大.
馬太福音第六章第三節.
有雅各,約瑟,西門和猶大.
馬太福音第六章第三節.
有雅各,約瑟,西門和猶大.
馬太福音第六章第三節.
有雅各,約瑟,西門和猶大.
馬太福音第六章第三節.
有雅各,約瑟,西門和猶大.
馬太福音第六章第三節.
有雅各,約瑟,西門和猶大.
馬太福音第六章第三節.
有雅各,約瑟,西門和猶大.
馬太福音第六章第三節.
有雅各,約瑟,西門和猶大.
馬太福音第六章第三節.
有雅各,約瑟,西門和猶大.
馬太福音第六章第三節.
有雅各,約瑟,西門和猶大.

$^{1161}$馬太福音第六章第三節.
有雅各,約瑟,西門和猶大.
馬太福音第六章第三節.
有雅各,約瑟,西門和猶大.
馬太福音第六章第三節.
有雅各,約瑟,西門和猶大.
馬太福音第六章第三節.
有雅各,約瑟,西門和猶大.
馬太福音第六章第三節.
有雅各,約瑟,西門和猶大.
馬太福音第六章第三節.
有雅各,約瑟,西門和猶大.
馬太福音第六章第三節.
有雅各,約瑟,西門和猶大.
馬太福音第六章第三節.
有雅各,約瑟,西門和猶大.
馬太福音第六章第三節.
有雅各,約瑟,西門和猶大.
馬太福音第六章第三節.
有雅各,約瑟,西門和猶大.
馬太福音第六章第三節.
有雅各,約瑟,西門和猶大.
馬太福音第六章第三節.
有雅各,約瑟,西門和猶大.
馬太福音第六章第三節.
有雅各,約瑟,西門和猶大.
馬太福音第六章第三節.
有雅各,約瑟,西門和猶大.
馬太福音第六章第三節.
有雅各,約瑟,西門和猶大.
馬太福音第六章第三節.
有雅各,約瑟,西門和猶大.
馬太福音第六章第三節.
有雅各,約瑟,西門和猶大.
馬太福音第六章第三節.
有雅各,約瑟,西門和猶大.
馬太福音第六章第三節.
有雅各,約瑟,西門和猶大.
馬太福音第六章第三節.
有雅各,約瑟,西門和猶大.

$^{1201}$馬太福音第六章第三節.
有雅各,約瑟,西門和猶大.
馬太福音第六章第三節.
有雅各,約瑟,西門和猶大.
馬太福音第六章第三節.
有雅各,約瑟,西門和猶大.
馬太福音第六章第三節.
有雅各,約瑟,西門和猶大.
馬太福音第六章第三節.
有雅各,約瑟,西門和猶大.
馬太福音第六章第四節.
有雅各,約瑟,西門和猶大.
馬太福音第六章第三節.
有雅各,約瑟,西門和猶大.
馬太福音第六章第三節.
有雅各,約瑟,西門和猶大.
馬太福音第六章第三節.
有雅各,約瑟,西門和猶大.
馬太福音第六章第三節.
有雅各,約瑟,西門和猶大.
馬太福音第六章第三節.
有雅各,約瑟,西門和猶大.
馬太福音第六章第三節.
有雅各,約瑟,西門和猶大.
馬太福音第六章第四節.
有雅各,約瑟,西門和猶大.
馬太福音第六章第三節.
有雅各,約瑟,西門和猶大.
馬太福音第六章第三節.
有雅各,約瑟,西門和猶大.
馬太福音第六章第三節.
有雅各,約瑟,西門和猶大.
馬太福音第六章第三節.
有雅各,約瑟,西門和猶大.
馬太福音第六章第三節.
有雅各,約瑟,西門和猶大.
馬太福音第六章第三節.
有雅各,約瑟,西門和猶大.
馬太福音第六章第三節.
有雅各,約瑟,西門和猶大.

$^{1241}$馬太福音第六章第三節.
有雅各,約瑟,西門和猶大.
馬太福音第六章第三節.
有雅各,約瑟,西門和猶大.
馬太福音第六章第三節.
有雅各,約瑟,西門和猶大.
馬太福音第六章第三節.
有雅各,約瑟,西門和猶大.
馬太福音第六章第三節.
有雅各,約瑟,西門和猶大.
馬太福音第六章第三節.
有雅各,約瑟,西門和猶大.
馬太福音第六章第三節.
有雅各,約瑟,西門和猶大.
馬太福音第六章第三節.
有雅各,約瑟,西門和猶大.
馬太福音第六章第三節.
有雅各,約瑟,西門和猶大.
馬太福音第六章第三節.
有雅各,約瑟,西門和猶大.
馬太福音第六章第三節.
有雅各,約瑟,西門和猶大.
馬太福音第六章第三節.
有雅各,約瑟,西門和猶大.
馬太福音第六章第三節.
有雅各,約瑟,西門和猶大.
馬太福音第六章第三節.
有雅各,約瑟,西門和猶大.
馬太福音第六章第三節.
有雅各,約瑟,西門和猶大.
馬太福音第六章第三節.
有雅各,約瑟,西門和猶大.
馬太福音第六章第三節.
有雅各,約瑟,西門和猶大.
馬太福音第六章第三節.
有雅各,約瑟,西門和猶大.
馬太福音第六章第三節.
有雅各,約瑟,西門和猶大.
馬太福音第六章第三節.
有雅各,約瑟,西門和猶大.

$^{1281}$馬太福音第六章第三節.
有雅各,約瑟,西門和猶大.
馬太福音第六章第三節.
有雅各,約瑟,西門和猶大.
馬太福音第六章第三節.
有雅各,約瑟,西門和猶大.
馬太福音第六章第三節.
有雅各,約瑟,西門和猶大.
馬太福音第六章第三節.
有雅各,約瑟,西門和猶大.
馬太福音第六章第三節.
有雅各,約瑟,西門和猶大.
馬太福音第六章第三節.
有雅各,約瑟,西門和猶大.
馬太福音第六章第三節.
有雅各,約瑟,西門和猶大.
馬太福音第六章第三節.
有雅各,約瑟,西門和猶大.
馬太福音第六章第三節.
有雅各,約瑟,西門和猶大.
馬太福音第六章第三節.
有雅各,約瑟,西門和猶大.
馬太福音第六章第三節.
有雅各,約瑟,西門和猶大.
馬太福音第六章第三節.
有雅各,約瑟,西門和猶大.
馬太福音第六章第三節.
有雅各,約瑟,西門和猶大.
馬太福音第六章第三節.
有雅各,約瑟,西門和猶大.
馬太福音第六章第三節.
有雅各,約瑟,西門和猶大.
馬太福音第六章第三節.
有雅各,約瑟,西門和猶大.
馬太福音第六章第三節.
有雅各,約瑟,西門和猶大.
馬太福音第六章第三節.
有雅各,約瑟,西門和猶大.
馬太福音第六章第三節.
有雅各,約瑟,西門和猶大.

$^{1321}$馬太福音第六章第三節.
有雅各,約瑟,西門和猶大.
馬太福音第六章第三節.
有雅各,約瑟,西門和猶大.
馬太福音第六章第三節.
有雅各,約瑟,西門和猶大.
馬太福音第六章第三節.
有雅各,約瑟,西門和猶大.
馬太福音第六章第三節.
有雅各,約瑟,西門和猶大.
馬太福音第六章第三節.
有雅各,約瑟,西門和猶大.
馬太福音第六章第三節.
有雅各,約瑟,西門和猶大.
馬太福音第六章第三節.
有雅各,約瑟,西門和猶大.
馬太福音第六章第三節.
有雅各,約瑟,西門和猶大.
馬太福音第六章第三節.
有雅各,約瑟,西門和猶大.
馬太福音第六章第三節.
有雅各,約瑟,西門和猶大.
馬太福音第六章第三節.
有雅各,約瑟,西門和猶大.
馬太福音第六章第三節.
有雅各,約瑟,西門和猶大.
馬太福音第六章第三節.
有雅各,約瑟,西門和猶大.
馬太福音第六章第三節.
有雅各,約瑟,西門和猶大.
馬太福音第六章第三節.
有雅各,約瑟,西門和猶大.
馬太福音第六章第三節.
有雅各,約瑟,西門和猶大.
馬太福音第六章第三節.
有雅各,約瑟,西門和猶大.
馬太福音第六章第三節.
有雅各,約瑟,西門和猶大.
馬太福音第六章第三節.
有雅各,約瑟,西門和猶大.

$^{1361}$馬太福音第六章第三節.
有雅各,約瑟,西門和猶大.
馬太福音第六章第三節.
有雅各,約瑟,西門和猶大.
馬太福音第六章第三節.
有雅各,約瑟,西門和猶大.
馬太福音第六章第三節.
有雅各,約瑟,西門和猶大.
馬太福音第六章第三節.
有雅各,約瑟,西門和猶大.
馬太福音第六章第三節.
有雅各,約瑟,西門和猶大.
馬太福音第六章第三節.
有雅各,約瑟,西門和猶大.
馬太福音第六章第三節.
有雅各,約瑟,西門和猶大.
馬太福音第六章第三節.
有雅各,約瑟,西門和猶大.
馬太福音第六章第三節.
有雅各,約瑟,西門和猶大.
馬太福音第六章第三節.
有雅各,約瑟,西門和猶大.
馬太福音第六章第三節.
有雅各,約瑟,西門和猶大.
馬太福音第六章第三節.
有雅各,約瑟,西門和猶大.
馬太福音第六章第三節.
有雅各,約瑟,西門和猶大.
馬太福音第六章第三節.
有雅各,約瑟,西門和猶大.
馬太福音第六章第三節.
有雅各,約瑟,西門和猶大.
馬太福音第六章第三節.
有雅各,約瑟,西門和猶大.
馬太福音第六章第三節.
有雅各,約瑟,西門和猶大.
馬太福音第六章第三節.
有雅各,約瑟,西門和猶大.
馬太福音第六章第三節.
有雅各,約瑟,西門和猶大.

$^{1401}$馬太福音第六章第三節.
有雅各,約瑟,西門和猶大.
馬太福音第六章第三節.
有雅各,約瑟,西門和猶大.
馬太福音第六章第三節.
有雅各,約瑟,西門和猶大.
馬太福音第六章第三節.
有雅各,約瑟,西門和猶大.
馬太福音第六章第三節.
有雅各,約瑟,西門和猶大.
馬太福音第六章第三節.
有雅各,約瑟,西門和猶大.
馬太福音第六章第三節.
有雅各,約瑟,西門和猶大.
馬太福音第六章第三節.
有雅各,約瑟,西門和猶大.
馬太福音第六章第三節.
有雅各,約瑟,西門和猶大.
馬太福音第六章第三節.
有雅各,約瑟,西門和猶大.
馬太福音第六章第三節.
有雅各,約瑟,西門和猶大.
馬太福音第六章第三節.
有雅各,約瑟,西門和猶大.
馬太福音第六章第三節.
有雅各,約瑟,西門和猶大.
馬太福音第六章第三節.
有雅各,約瑟,西門和猶大.
馬太福音第六章第三節.
有雅各,約瑟,西門和猶大.
馬太福音第六章第三節.
有雅各,約瑟,西門和猶大.
馬太福音第六章第三節.
有雅各,約瑟,西門和猶大.
馬太福音第六章第三節.
有雅各,約瑟,西門和猶大.
馬太福音第六章第三節.
有雅各,約瑟,西門和猶大.
馬太福音第六章第三節.
有雅各,約瑟,西門和猶大.

$^{1441}$馬太福音第六章第三節.
有雅各,約瑟,西門和猶大.
馬太福音第六章第三節.
有雅各,約瑟,西門和猶大.
馬太福音第六章第三節.
有雅各,約瑟,西門和猶大.
馬太福音第六章第三節.
有雅各,約瑟,西門和猶大.
馬太福音第六章第三節.
有雅各,約瑟,西門和猶大.
馬太福音第六章第三節.
有雅各,約瑟,西門和猶大.
馬太福音第六章第三節.
有雅各,約瑟,西門和猶大.
馬太福音第六章第三節.
有雅各,約瑟,西門和猶大.
馬太福音第六章第三節.
有雅各,約瑟,西門和猶大.
馬太福音第六章第三節.
有雅各,約瑟,西門和猶大.
馬太福音第六章第三節.
有雅各,約瑟,西門和猶大.
馬太福音第六章第三節.
有雅各,約瑟,西門和猶大.
馬太福音第六章第三節.
有雅各,約瑟,西門和猶大.
馬太福音第六章第三節.
有雅各,約瑟,西門和猶大.
馬太福音第六章第三節.
有雅各,約瑟,西門和猶大.
馬太福音第六章第三節.
有雅各,約瑟,西門和猶大.
馬太福音第六章第三節.
有雅各,約瑟,西門和猶大.
馬太福音第六章第三節.
有雅各,約瑟,西門和猶大.
馬太福音第六章第三節.
有雅各,約瑟,西門和猶大.
馬太福音第六章第三節.
有雅各,約瑟,西門和猶大.

$^{1481}$馬太福音第六章第三節.
有雅各,約瑟,西門和猶大.
馬太福音第六章第三節.
有雅各,約瑟,西門和猶大.
馬太福音第六章第三節.
有雅各,約瑟,西門和猶大.
馬太福音第六章第三節.
有雅各,約瑟,西門和猶大.
馬太福音第六章第三節.
有雅各,約瑟,西門和猶大.
馬太福音第六章第三節.
有雅各,約瑟,西門和猶大.
馬太福音第六章第三節.
有雅各,約瑟,西門和猶大.
馬太福音第六章第三節.
有雅各,約瑟,西門和猶大.
馬太福音第六章第三節.
有雅各,約瑟,西門和猶大.
馬太福音第六章第三節.
有雅各,約瑟,西門和猶大.
馬太福音第六章第三節.
有雅各,約瑟,西門和猶大.
馬太福音第六章第三節.
有雅各,約瑟,西門和猶大.
馬太福音第六章第三節.
有雅各,約瑟,西門和猶大.
馬太福音第六章第三節.
有雅各,約瑟,西門和猶大.
馬太福音第六章第三節.
有雅各,約瑟,西門和猶大.
馬太福音第六章第三節.
有雅各,約瑟,西門和猶大.
馬太福音第六章第三節.
有雅各,約瑟,西門和猶大.
馬太福音第六章第三節.
有雅各,約瑟,西門和猶大.
馬太福音第六章第三節.
有雅各,約瑟,西門和猶大.
馬太福音第六章第三節.
有雅各,約瑟,西門和猶大.

$^{1521}$馬太福音第六章第三節.
有雅各,約瑟,西門和猶大.
馬太福音第六章第三節.
有雅各,約瑟,西門和猶大.
馬太福音第六章第三節.
有雅各,約瑟,西門和猶大.
馬太福音第六章第三節.
有雅各,約瑟,西門和猶大.
馬太福音第六章第三節.
有雅各,約瑟,西門和猶大.
馬太福音第六章第三節.
有雅各,約瑟,西門和猶大.
馬太福音第六章第三節.
有雅各,約瑟,西門和猶大.
馬太福音第六章第三節.
有雅各,約瑟,西門和猶大.
馬太福音第六章第三節.
有雅各,約瑟,西門和猶大.
馬太福音第六章第三節.
有雅各,約瑟,西門和猶大.
馬太福音第六章第三節.
有雅各,約瑟,西門和猶大.
馬太福音第六章第三節.
有雅各,約瑟,西門和猶大.
馬太福音第六章第三節.
有雅各,約瑟,西門和猶大.
馬太福音第六章第三節.
有雅各,約瑟,西門和猶大.
馬太福音第六章第三節.
有雅各,約瑟,西門和猶大.
馬太福音第六章第三節.
有雅各,約瑟,西門和猶大.
馬太福音第六章第三節.
有雅各,約瑟,西門和猶大.
馬太福音第六章第三節.
有雅各,約瑟,西門和猶大.
馬太福音第六章第三節.
有雅各,約瑟,西門和猶大.
馬太福音第六章第三節.
有雅各,約瑟,西門和猶大.

$^{1561}$馬太福音第六章第三節.
有雅各,約瑟,西門和猶大.
馬太福音第六章第三節.
有雅各,約瑟,西門和猶大.
馬太福音第六章第三節.
有雅各,約瑟,西門和猶大.
馬太福音第六章第三節.
有雅各,約瑟,西門和猶大.
馬太福音第六章第三節.
有雅各,約瑟,西門和猶大.
馬太福音第六章第三節.
有雅各,約瑟,西門和猶大.
馬太福音第六章第三節.
有雅各,約瑟,西門和猶大.
馬太福音第六章第三節.
有雅各,約瑟,西門和猶大.
馬太福音第六章第三節.
有雅各,約瑟,西門和猶大.
馬太福音第六章第三節.
有雅各,約瑟,西門和猶大.
馬太福音第六章第三節.
有雅各,約瑟,西門和猶大.
馬太福音第六章第三節.
有雅各,約瑟,西門和猶大.
馬太福音第六章第三節.
有雅各,約瑟,西門和猶大.
馬太福音第六章第三節.
有雅各,約瑟,西門和猶大.
馬太福音第六章第三節.
有雅各,約瑟,西門和猶大.
馬太福音第六章第三節.
有雅各,約瑟,西門和猶大.
馬太福音第六章第三節.
有雅各,約瑟,西門和猶大.
馬太福音第六章第三節.
有雅各,約瑟,西門和猶大.
馬太福音第六章第三節.
有雅各,約瑟,西門和猶大.
馬太福音第六章第三節.
有雅各,約瑟,西門和猶大.

$^{1601}$馬太福音第六章第三節.
有雅各,約瑟,西門和猶大.
馬太福音第六章第三節.
有雅各,約瑟,西門和猶大.
馬太福音第六章第三節.
有雅各,約瑟,西門和猶大.
馬太福音第六章第三節.
有雅各,約瑟,西門和猶大.
馬太福音第六章第三節.
有雅各,約瑟,西門和猶大.
馬太福音第六章第三節.
有雅各,約瑟,西門和猶大.
馬太福音第六章第三節.
有雅各,約瑟,西門和猶大.
馬太福音第六章第三節.
有雅各,約瑟,西門和猶大.
馬太福音第六章第三節.
有雅各,約瑟,西門和猶大.
馬太福音第六章第三節.
有雅各,約瑟,西門和猶大.
馬太福音第六章第三節.
有雅各,約瑟,西門和猶大.
馬太福音第六章第三節.
有雅各,約瑟,西門和猶大.
馬太福音第六章第三節.
有雅各,約瑟,西門和猶大.
馬太福音第六章第三節.
有雅各,約瑟,西門和猶大.
馬太福音第六章第三節.
有雅各,約瑟,西門和猶大.
馬太福音第六章第三節.
有雅各,約瑟,西門和猶大.
馬太福音第六章第三節.
有雅各,約瑟,西門和猶大.
馬太福音第六章第三節.
有雅各,約瑟,西門和猶大.
馬太福音第六章第三節.
有雅各,約瑟,西門和猶大.
馬太福音第六章第三節.
有雅各,約瑟,西門和猶大.

$^{1641}$馬太福音第六章第三節.
有雅各,約瑟,西門和猶大.
馬太福音第六章第三節.
有雅各,約瑟,西門和猶大.
馬太福音第六章第三節.
有雅各,約瑟,西門和猶大.
馬太福音第六章第三節.
有雅各,約瑟,西門和猶大.
馬太福音第六章第三節.
有雅各,約瑟,西門和猶大.
馬太福音第六章第三節.
有雅各,約瑟,西門和猶大.
馬太福音第六章第三節.
有雅各,約瑟,西門和猶大.
馬太福音第六章第三節.
有雅各,約瑟,西門和猶大.
馬太福音第六章第三節.
有雅各,約瑟,西門和猶大.
馬太福音第六章第三節.
有雅各,約瑟,西門和猶大.
馬太福音第六章第三節.
有雅各,約瑟,西門和猶大.
馬太福音第六章第三節.
有雅各,約瑟,西門和猶大.
馬太福音第六章第三節.
有雅各,約瑟,西門和猶大.
馬太福音第六章第三節.
有雅各,約瑟,西門和猶大.
馬太福音第六章第三節.
有雅各,約瑟,西門和猶大.
馬太福音第六章第三節.
有雅各,約瑟,西門和猶大.
馬太福音第六章第三節.
有雅各,約瑟,西門和猶大.
馬太福音第六章第三節.
有雅各,約瑟,西門和猶大.
馬太福音第六章第三節.
有雅各,約瑟,西門和猶大.
馬太福音第六章第三節.
有雅各,約瑟,西門和猶大.

$^{1681}$馬太福音第六章第三節.
有雅各,約瑟,西門和猶大.
馬太福音第六章第三節.
有雅各,約瑟,西門和猶大.
馬太福音第六章第三節.
有雅各,約瑟,西門和猶大.
馬太福音第六章第三節.
有雅各,約瑟,西門和猶大.
馬太福音第六章第三節.
有雅各,約瑟,西門和猶大.
馬太福音第六章第三節.
有雅各,約瑟,西門和猶大.
馬太福音第六章第三節.
有雅各,約瑟,西門和猶大.
馬太福音第六章第三節.
有雅各,約瑟,西門和猶大.
馬太福音第六章第三節.
有雅各,約瑟,西門和猶大.
馬太福音第六章第三節.
有雅各,約瑟,西門和猶大.
馬太福音第六章第三節.
有雅各,約瑟,西門和猶大.
馬太福音第六章第三節.
有雅各,約瑟,西門和猶大.
馬太福音第六章第三節.
有雅各,約瑟,西門和猶大.
馬太福音第六章第三節.
有雅各,約瑟,西門和猶大.
馬太福音第六章第三節.
有雅各,約瑟,西門和猶大.
馬太福音第六章第三節.
有雅各,約瑟,西門和猶大.
馬太福音第六章第三節.
有雅各,約瑟,西門和猶大.
馬太福音第六章第三節.
有雅各,約瑟,西門和猶大.
馬太福音第六章第三節.
有雅各,約瑟,西門和猶大.
馬太福音第六章第三節.
有雅各,約瑟,西門和猶大.

$^{1721}$馬太福音第六章第三節.
有雅各,約瑟,西門和猶大.
馬太福音第六章第三節.
有雅各,約瑟,西門和猶大.
馬太福音第六章第三節.
有雅各,約瑟,西門和猶大.
馬太福音第六章第三節.
有雅各,約瑟,西門和猶大.
馬太福音第六章第三節.
有雅各,約瑟,西門和猶大.
馬太福音第六章第三節.
有雅各,約瑟,西門和猶大.
馬太福音第六章第三節.
有雅各,約瑟,西門和猶大.
馬太福音第六章第三節.
有雅各,約瑟,西門和猶大.
馬太福音第六章第三節.
有雅各,約瑟,西門和猶大.
馬太福音第六章第三節.
有雅各,約瑟,西門和猶大.
馬太福音第六章第三節.
有雅各,約瑟,西門和猶大.
馬太福音第六章第三節.
有雅各,約瑟,西門和猶大.
馬太福音第六章第三節.
有雅各,約瑟,西門和猶大.
馬太福音第六章第三節.
有雅各,約瑟,西門和猶大.
馬太福音第六章第三節.
有雅各,約瑟,西門和猶大.
馬太福音第六章第三節.
有雅各,約瑟,西門和猶大.
馬太福音第六章第三節.
有雅各,約瑟,西門和猶大.
馬太福音第六章第三節.
有雅各,約瑟,西門和猶大.
馬太福音第六章第三節.
有雅各,約瑟,西門和猶大.
馬太福音第六章第三節.
有雅各,約瑟,西門和猶大.

$^{1761}$馬太福音第六章第三節.
有雅各,約瑟,西門和猶大.
馬太福音第六章第三節.
有雅各,約瑟,西門和猶大.
馬太福音第六章第三節.
有雅各,約瑟,西門和猶大.
馬太福音第六章第三節.
有雅各,約瑟,西門和猶大.
馬太福音第六章第三節.
有雅各,約瑟,西門和猶大.
馬太福音第六章第三節.
有雅各,約瑟,西門和猶大.
馬太福音第六章第三節.
有雅各,約瑟,西門和猶大.
馬太福音第六章第三節.
有雅各,約瑟,西門和猶大.
馬太福音第六章第三節.
有雅各,約瑟,西門和猶大.
馬太福音第六章第三節.
有雅各,約瑟,西門和猶大.
馬太福音第六章第三節.
有雅各,約瑟,西門和猶大.
馬太福音第六章第三節.
有雅各,約瑟,西門和猶大.
馬太福音第六章第三節.
有雅各,約瑟,西門和猶大.
馬太福音第六章第三節.
有雅各,約瑟,西門和猶大.
馬太福音第六章第三節.
有雅各,約瑟,西門和猶大.
馬太福音第六章第三節.
有雅各,約瑟,西門和猶大.
馬太福音第六章第三節.
有雅各,約瑟,西門和猶大.
馬太福音第六章第三節.
有雅各,約瑟,西門和猶大.
馬太福音第六章第三節.
有雅各,約瑟,西門和猶大.
馬太福音第六章第三節.
有雅各,約瑟,西門和猶大.

$^{1801}$馬太福音第六章第三節.
有雅各,約瑟,西門和猶大.
馬太福音第六章第三節.
有雅各,約瑟,西門和猶大.
馬太福音第六章第三節.
有雅各,約瑟,西門和猶大.
馬太福音第六章第三節.
有雅各,約瑟,西門和猶大.
馬太福音第六章第三節.
有雅各,約瑟,西門和猶大.
馬太福音第六章第三節.
有雅各,約瑟,西門和猶大.
馬太福音第六章第三節.
有雅各,約瑟,西門和猶大.
馬太福音第六章第三節.
有雅各,約瑟,西門和猶大.
馬太福音第六章第三節.
有雅各,約瑟,西門和猶大.
馬太福音第六章第三節.
有雅各,約瑟,西門和猶大.
馬太福音第六章第三節.
有雅各,約瑟,西門和猶大.
馬太福音第六章第三節.
有雅各,約瑟,西門和猶大.
馬太福音第六章第三節.
有雅各,約瑟,西門和猶大.
馬太福音第六章第三節.
有雅各,約瑟,西門和猶大.
馬太福音第六章第三節.
有雅各,約瑟,西門和猶大.
馬太福音第六章第三節.
有雅各,約瑟,西門和猶大.
馬太福音第六章第三節.
有雅各,約瑟,西門和猶大.
馬太福音第六章第三節.
有雅各,約瑟,西門和猶大.
馬太福音第六章第三節.
有雅各,約瑟,西門和猶大.
馬太福音第六章第三節.
有雅各,約瑟,西門和猶大.

$^{1841}$馬太福音第六章第三節.
有雅各,約瑟,西門和猶大.
馬太福音第六章第三節.
有雅各,約瑟,西門和猶大.
馬太福音第六章第三節.
有雅各,約瑟,西門和猶大.
馬太福音第六章第三節.
有雅各,約瑟,西門和猶大.
馬太福音第六章第三節.
有雅各,約瑟,西門和猶大.
馬太福音第六章第三節.
有雅各,約瑟,西門和猶大.
馬太福音第六章第三節.
有雅各,約瑟,西門和猶大.
馬太福音第六章第三節.
有雅各,約瑟,西門和猶大.
馬太福音第六章第三節.
有雅各,約瑟,西門和猶大.
馬太福音第六章第三節.
有雅各,約瑟,西門和猶大.
馬太福音第六章第三節.
有雅各,約瑟,西門和猶大.
馬太福音第六章第三節.
有雅各,約瑟,西門和猶大.
馬太福音第六章第三節.
有雅各,約瑟,西門和猶大.
馬太福音第六章第三節.
有雅各,約瑟,西門和猶大.
馬太福音第六章第三節.
有雅各,約瑟,西門和猶大.
馬太福音第六章第三節.
有雅各,約瑟,西門和猶大.
馬太福音第六章第三節.
有雅各,約瑟,西門和猶大.
馬太福音第六章第三節.
有雅各,約瑟,西門和猶大.
馬太福音第六章第三節.
有雅各,約瑟,西門和猶大.
馬太福音第六章第三節.
有雅各,約瑟,西門和猶大.

$^{1881}$馬太福音第六章第三節.
有雅各,約瑟,西門和猶大.
馬太福音第六章第三節.
有雅各,約瑟,西門和猶大.
馬太福音第六章第三節.
有雅各,約瑟,西門和猶大.
馬太福音第六章第三節.
有雅各,約瑟,西門和猶大.
馬太福音第六章第三節.
有雅各,約瑟,西門和猶大.
馬太福音第六章第三節.
有雅各,約瑟,西門和猶大.
馬太福音第六章第三節.
有雅各,約瑟,西門和猶大.
馬太福音第六章第三節.
有雅各,約瑟,西門和猶大.
馬太福音第六章第三節.
有雅各,約瑟,西門和猶大.
馬太福音第六章第三節.
有雅各,約瑟,西門和猶大.
馬太福音第六章第三節.
有雅各,約瑟,西門和猶大.
馬太福音第六章第三節.
有雅各,約瑟,西門和猶大.
馬太福音第六章第三節.
有雅各,約瑟,西門和猶大.
馬太福音第六章第三節.
有雅各,約瑟,西門和猶大.
馬太福音第六章第三節.
有雅各,約瑟,西門和猶大.
馬太福音第六章第三節.
有雅各,約瑟,西門和猶大.
馬太福音第六章第三節.
有雅各,約瑟,西門和猶大.
馬太福音第六章第三節.
有雅各,約瑟,西門和猶大.
馬太福音第六章第三節.
有雅各,約瑟,西門和猶大.
馬太福音第六章第三節.
有雅各,約瑟,西門和猶大.

$^{1921}$馬太福音第六章第三節.
有雅各,約瑟,西門和猶大.
馬太福音第六章第三節.
有雅各,約瑟,西門和猶大.
馬太福音第六章第三節.
有雅各,約瑟,西門和猶大.
馬太福音第六章第三節.
有雅各,約瑟,西門和猶大.
馬太福音第六章第三節.
有雅各,約瑟,西門和猶大.
馬太福音第六章第三節.
有雅各,約瑟,西門和猶大.
馬太福音第六章第三節.
有雅各,約瑟,西門和猶大.
馬太福音第六章第三節.
有雅各,約瑟,西門和猶大.
馬太福音第六章第三節.
有雅各,約瑟,西門和猶大.
馬太福音第六章第三節.
有雅各,約瑟,西門和猶大.
馬太福音第六章第三節.
有雅各,約瑟,西門和猶大.
馬太福音第六章第三節.
有雅各,約瑟,西門和猶大.
馬太福音第六章第三節.
有雅各,約瑟,西門和猶大.
馬太福音第六章第三節.
有雅各,約瑟,西門和猶大.
馬太福音第六章第三節.
有雅各,約瑟,西門和猶大.
馬太福音第六章第三節.
有雅各,約瑟,西門和猶大.
馬太福音第六章第三節.
有雅各,約瑟,西門和猶大.
馬太福音第六章第三節.
有雅各,約瑟,西門和猶大.
馬太福音第六章第三節.
有雅各,約瑟,西門和猶大.
馬太福音第六章第三節.
有雅各,約瑟,西門和猶大.

$^{1961}$馬太福音第六章第三節.
有雅各,約瑟,西門和猶大.
馬太福音第六章第三節.
有雅各,約瑟,西門和猶大.
馬太福音第六章第三節.
有雅各,約瑟,西門和猶大.
馬太福音第六章第三節.
有雅各,約瑟,西門和猶大.
馬太福音第六章第三節.
有雅各,約瑟,西門和猶大.
馬太福音第六章第三節.
有雅各,約瑟,西門和猶大.
馬太福音第六章第三節.
有雅各,約瑟,西門和猶大.
馬太福音第六章第三節.
有雅各,約瑟,西門和猶大.
馬太福音第六章第三節.
有雅各,約瑟,西門和猶大.
馬太福音第六章第三節.
有雅各,約瑟,西門和猶大.
馬太福音第六章第三節.
有雅各,約瑟,西門和猶大.
馬太福音第六章第三節.
有雅各,約瑟,西門和猶大.
馬太福音第六章第三節.
有雅各,約瑟,西門和猶大.
馬太福音第六章第三節.
有雅各,約瑟,西門和猶大.
馬太福音第六章第三節.
有雅各,約瑟,西門和猶大.
馬太福音第六章第三節.
有雅各,約瑟,西門和猶大.
馬太福音第六章第三節.
有雅各,約瑟,西門和猶大.
馬太福音第六章第三節.
有雅各,約瑟,西門和猶大.
馬太福音第六章第三節.
有雅各,約瑟,西門和猶大.
馬太福音第六章第三節.
有雅各,約瑟,西門和猶大.

$^{2001}$馬太福音第六章第三節.
有雅各,約瑟,西門和猶大.
馬太福音第六章第三節.
有雅各,約瑟,西門和猶大.
馬太福音第六章第三節.
有雅各,約瑟,西門和猶大.
馬太福音第六章第三節.
有雅各,約瑟,西門和猶大.
馬太福音第六章第三節.
有雅各,約瑟,西門和猶大.
馬太福音第六章第三節.
有雅各,約瑟,西門和猶大.
馬太福音第六章第三節.
有雅各,約瑟,西門和猶大.
馬太福音第六章第三節.
有雅各,約瑟,西門和猶大.
馬太福音第六章第三節.
有雅各,約瑟,西門和猶大.
馬太福音第六章第三節.
有雅各,約瑟,西門和猶大.
馬太福音第六章第三節.
有雅各,約瑟,西門和猶大.
馬太福音第六章第三節.
有雅各,約瑟,西門和猶大.
馬太福音第六章第三節.
有雅各,約瑟,西門和猶大.
馬太福音第六章第三節.
有雅各,約瑟,西門和猶大.
馬太福音第六章第三節.
有雅各,約瑟,西門和猶大.
馬太福音第六章第三節.
有雅各,約瑟,西門和猶大.
馬太福音第六章第三節.
有雅各,約瑟,西門和猶大.
馬太福音第六章第三節.
有雅各,約瑟,西門和猶大.
馬太福音第六章第三節.
有雅各,約瑟,西門和猶大.
馬太福音第六章第三節.
有雅各,約瑟,西門和猶大.

$^{2041}$馬太福音第六章第三節.
有雅各,約瑟,西門和猶大.
馬太福音第六章第三節.
有雅各,約瑟,西門和猶大.
馬太福音第六章第三節.
有雅各,約瑟,西門和猶大.
馬太福音第六章第三節.
有雅各,約瑟,西門和猶大.
馬太福音第六章第三節.
有雅各,約瑟,西門和猶大.
馬太福音第六章第三節.
有雅各,約瑟,西門和猶大.
馬太福音第六章第三節.
有雅各,約瑟,西門和猶大.
馬太福音第六章第三節.
有雅各,約瑟,西門和猶大.
馬太福音第六章第三節.
有雅各,約瑟,西門和猶大.
馬太福音第六章第三節.
有雅各,約瑟,西門和猶大.
馬太福音第六章第三節.
有雅各,約瑟,西門和猶大.
馬太福音第六章第三節.
有雅各,約瑟,西門和猶大.
馬太福音第六章第三節.
有雅各,約瑟,西門和猶大.
馬太福音第六章第三節.
有雅各,約瑟,西門和猶大.
馬太福音第六章第三節.
有雅各,約瑟,西門和猶大.
\newpage



\section{腓立比書 4:4}
\label{sec:032084}
\textbf{喜樂,喜樂:生命的顛覆 (粵/華)}~~~~陳淳牧師
\newline
\newline
連結: \href{http://www.church.com.hk/acms/content.asp?site=cdc\&op=show\&type=product\&code=032084\&layout=sermon}{\texttt{032084}} ~~~~ 語音日期: 2026-05-03
\newline
\newline
\hyperref[sec:032083]{\small{< < < PREV SERMON < < <}}
~
\hyperref[sec:index_chronic]{\small{[返順時目]}}
~
\hyperref[sec:032085]{\small{> > > NEXT SERMON > > >}}
\newline
\newline
腓立比書 4:4
\newline
\begin{longtable}{cl}
\hline
\hline
章節 & 經文 (和合本修訂版)\\
\hline
4:4 & \begin{tabularx}{0.7\textwidth}{X} 你們要靠主常常喜樂。我再說，你們要喜樂。 \end{tabularx} \\ \\
[1ex]
\hline
\hline
\end{longtable}
$^{1}$我們一起來祈禱.
這一份是遲來的講章.
本來這篇講章在半年前就寫好了.
預備在將臨期宣講.
如果你記性好的話.
十一月二十六日.
就算你記性不好你也一定會記得.
洪福苑大火震驚整個全香港整個世界.
奪走了很多寶貴的生命.
燒碎了我們很多人的心.
到今天我們看到相關的新聞.
仍然心有戚焉.
所以當天我擱置了這篇講章.
因為作為一個牧者.
我不能夠在那個哀傷的過程裡面.
漠視大家的情緒.
雖然宣講真理.
但可能成為我們傷口不能承受的那種傷痛.
所以我當時用了羅馬書十五章十三節.
宣講盼望盼望.
因為在那個時候我們最需要.
仰望那位在黑暗中仍然發光的救主.
但我有一個祈禱.
我想找一個機會.
面對當時我不敢講或者不輕易講的題目.
半年前的盼望.
今天喜樂喜樂.
主要許可.
三個星期後輪到我講到.
我會講平安.
這三個的信息.
希望成為我們生命裡面的激勵.
說回今天的題目.
喜樂喜樂.
為什麼我要這麼講呢.
我是來平衡盼望盼望嗎.
不是的.
因為保羅在腓立比書裡面是這樣講.
我再說你們要喜樂.
這個重複是想告訴我們.

$^{41}$我們無論生活有多麼身不由己.
我們仍然能夠做出決定.
無論環境怎樣讓我們沮喪.
我們仍然能夠成為我們心情的主人.
無論人生有多少的低谷.
我們仍然能夠重新奪回生命的安穩和權利.
今天我們要透過保羅所講的喜樂.
從三個層次顛覆和反轉我們一些的人生邏輯.
在講到之前我們一起來禱告.
親愛的主我們來到您面前.
求您安靜我們的心.
除去一切的雜念和宣揚.
求您開啟我們的耳朵和心靈.
使我們今天能夠專注聆聽.
又能明白和順服您的話語.
我們如此祈禱.
奉我主耶穌基督的名求.
阿們.
保羅是這樣講的.
他說你們要靠主常常喜樂.
我再說.
你們要喜樂.
喜樂喜樂生命的顛覆.
始於從環境到根源的顛覆.
如果我們讀到你們要喜樂這個命令的時候.
不知道你心裡面有沒有疑惑呢?.
你有或沒有.
我就有了.
情緒難道是可以用命令蓋過的嗎?.
因為我們世界有個邏輯.
喜樂是結果.
是好事發生環境順暢的時候.
我們被動接收的禮物.
是好事發生我們被動接收的禮物.
真是這樣嗎?.
保羅在監獄裡面寫這個命令.
豈不是諷刺的語句?.
所以今天透過今天的經文.
想顛覆我們偏差的概念.
想要顛覆我們偏差的一個概念.

$^{81}$這個命令全部的奧秘.
開始於兩個字.
在哪兩個字呢?.
你看一下螢幕就知道了.
靠主.
中文兩個字.
原文也是兩個字.
在主前面有一個介詞.
這個介詞是說.
身處於哪一個位置.
或者身處於哪一個生存的空間.
直接一點翻譯.
你們要留在主裡面.
住在主裡面.
這個就是在主裡常常喜樂的秘訣.
你們要留在主裡面.
你們要住在主裡面.
這個就是你們常常喜樂的秘訣.
喜樂首要的條件.
不是你使用了什麼有效的方式.
或者那些方式很屬靈.
或者那些方式很有效.
而是你居住在哪裡.
你的所在.
決定了你的所有.
你的所在.
決定了你的擁有.
住在基督裡.
是整本新約聖經的基石.
這意味著.
你不是在風暴裡面的孤獨人.
你已經被遷到神的國度裡面.
被安置在萬古磐石的洞穴當中.
這是一個畏憤性的喜樂.
它不是從你擁有什麼開始.
而是從你在哪裡開始.
你出生在哪裡.
你湧流出來的喜樂.
是基督徒的喜樂.
根源不在環境的順逆.

$^{121}$而是那位已然來臨.
並且永不離開的耶穌基督.
祂就是我們生命裡面最穩妥.
最能遮風擋雨的保障.
所以請我們記得.
喜樂不是你看環境.
不是你看天氣.
甚至不是你看人的臉色來做人.
只要你留在主裡面.
外面的風雨再大.
都奪不走你內心的安穩.
保羅說你們要靠主上上喜樂.
我再說.
你們要喜樂.
題目說喜樂喜樂生命的顛覆.
是從情緒到狀態的顛覆.
我們大部分的人都經驗過.
喜樂總是跟某一些好的事情掛鉤.
升職.
加薪.
大病初癒.
家庭和睦.
心想事成.
好不好呢?.
這些好事發生的時候.
我們總是覺得心中喜樂.
但是這些事情結束了以後呢?.
或者好像半年前那樣發生悲劇呢?.
甚至今天你和我都有一些不好的事情發生.
甚至今天你和我也會有一些不好的事情發生.
這些喜樂隨之消失.
這些的喜樂隨之消失.
找都找不回.
找也找不回來.
這樣導致我們的生命.
很容易被分割兩部分.
喜樂的時段.
不喜樂的時段.
然後你會聽過說.
人生不如意的事情時常八九.

$^{161}$所以我們大部分的時間都是不喜樂的.
是不是這樣呢?.
請聽神的話.
祂說你們要靠主常常喜樂.
神的命令不是邀請你參加幾場喜樂盛宴.
神的命令不是邀請你參加幾場喜樂的盛宴.
祂是邀請我們進入.
一種持續居住在祂裡面的喜樂狀態.
一種持續的進入祂裡面的喜樂的狀態.
所以我們要從環境好,心情就好的被動反應.
所以我們要從環境好,心情就好的這種被動的反應.
轉變為無論日子苦或甜.
都靠主守護喜樂.
這個是主動的選擇.
為什麼這麼說呢?.
因為常常這個字.
不是一個字,是兩個字.
因為原文是一個字,所以我說錯了.
它裡面是說到在任何時候始終.
無例外的絕對的情況.
要喜樂.
這三個字.
原文都是一個字.
雖然我們翻譯的時候會把它放在不同的地方.
但是原文是一個字.
它用的是命令的語氣.
現在時態.
原文的現在時態是怎樣的呢?.
從現在開始做.
然後慢慢逐漸延續至死方休.
是一種永恆常規的狀態.
所以它描繪的是一種持續進行.
習慣性的生命常態.
所以這兩個詞語放在一起.
是想說喜樂不應該是人生樂章.
偶然出現的幾個興奮高音.
它是貫穿整個樂譜.
穩定而悠揚的愉快音調.
不知道你聽的時候有沒有這些想法.
你會不會想問我們.

$^{201}$牧師.
面對人間的悲劇.
人間的冤屈.
我心裡面的那種不忿.
或者我心裡面受的委屈實在太多了.
我怎樣能夠常常喜樂呢?.
所以答案不是靠喜樂的延續性.
也不是它的持續有多久.
而是取決於它的源頭根源的穩定性.
如果你的喜樂源頭是.
變幻的環境.
會貶值的金錢.
會改變的人事.
或者是某些人的心情.
這樣的喜樂就好像海上的波浪一樣.
隨波起伏.
如果你的喜樂是那位今日,昨日.
直到永遠的至高主.
你的喜樂就會得到穩定性.
得到堅固性.
我邀請你檢視一下你自己.
你自己內心跟主的連結狀態.
你生命裡面的那種喜樂.
是不是時強時弱呢?.
是不是一進入人生的低谷.
或者人生的隧道裡面.
你就立即斷線.
環境稍微差一點.
心情稍微差一點.
立即就跌落沒有訊號的狀態.
靠著你的意志力苦苦支撐.
或者某一些的記憶支持你自己.
還是說你是那種能夠穩定隨時.
接得通的天國的WIFI.
我這麼講.
不是代表要否定你的心情.
人的情緒會起伏不定.
是最真實而且最自然的人性.
是上帝設計來給我們守護自己的.
但是這個守護.

$^{241}$是透過一個信差的方式來守護.
它是透過你自己對你自己價值觀,人生觀的了解.
然後接觸外面的現實環境.
它跟你裡面的人生觀,價值觀相遇.
當它們產生衝突的時候.
它就會有不愉快的情緒出現.
它就會有不愉快的這種的心情出現.
當它合乎你的心意.
你就心裡面愉快了.
輕鬆了.
所以請記住它是一個信差來的.
信差送信息.
要不要留在家裡過夜呢.
還是留在你家裡做主人呢.
信差送完信就請他離開了.
所以但願主來幫助我們.
情緒只是一個信差.
而神的心意是要求我們持續住在主裡面.
無論情緒如何.
保持連線暢通.
每天透過祈禱和讀經.
每周透過穩定的團契成長班.
或者今天的崇拜.
我們一起來增強我們的信號.
無論日子多艱難.
依然能夠保持在線.
喜樂靠主.
永不中斷.
保羅說.
你們要靠主常常喜樂.
我再說.
你們要喜樂.
喜樂喜樂生命的顛覆.
最關鍵是從被動到主動的顛覆.
這一切是不是會自然出現呢.
當然不會.
如果你等.
你就等一等就好了.
不是等一等,是等到不知道什麼時候都等不到.
不是等一等,是等到不知道什麼時候才等得到.

$^{281}$因為我們面對所有的環境都有不同的心情.
你沒有心情就糟糕了.
你沒有心情就慘了.
要看醫生了.
或者你來跟我聊聊天.
但是無論怎麼樣也好.
一般的人.
面對災難.
面對無辜受苦.
面對冤屈.
面對和自己不同的設定.
你就會覺得憤怒.
你就會覺得沮喪.
你會覺得無奈.
你會覺得抑鬱.
當你心碎到無法自拔的時候.
我們應該怎樣做呢.
保羅用一個重複的命令.
將責任和能力.
全部交回給我們.
他說我再說,你們要喜樂.
這意味著喜樂不只是清晰的道理.
不是你明白就算了.
不是你知道就可以了.
而是需要發動你的意志.
去執行的軍令.
是從被動的情感反應.
邁向主動的意志行動.
這一步非常關鍵.
關鍵到什麼程度呢.
關鍵到你能不能請那位郵差離開.
會不會客人佔了你的主人的房間.
但願我們清楚的能夠知道.
保羅所說的我再說.
我們用黃色字告訴他們.
這個字在這裡非常重要.
因為在原文的語法裡面.
它不是三個字,它是兩個字.
其中一個字是極具衝擊力.
它的宣講是要衝擊你.

$^{321}$保羅很清楚的知道.
我們身在困境不開心的裡面.
是聽不進這些真理的.
或者覺得這些真理是不合理的.
聽不進耳.
或者聽完就忘記了.
因此他這麼說.
他說我知道很難.
所以你要很小心的聽我講.
聽清楚了.
這個不是一個情感的選項.
這是一個必須執行的屬靈動作.
在神的國度裡面.
我們很清楚的知道.
每一個的命令都是承載恩典的管道.
他必然為他的命令負責到底.
命令是神給我們的信號.
指引著我們接通那份屬天的恩典.
意志行動就是信心具體的表現.
所以經文不是叫你強顏歡笑.
不是叫你皮笑肉不笑.
不是叫你在人面前假裝歡笑.
它是信心具體的操練.
當你選擇喜樂.
不是代表你抹去你內心的傷痛.
不是你強行收著淚水不哭.
我們說情感會令我們流眼淚.
是健康的.
是一個很好的信號.
代表著我們心裡面非常無奈和傷感.
但是不應該停在這裡.
信心讓我們靠著得安穩.
當環境一片黑暗.
前途未卜.
信心對我們的意志這樣講.
現在你雖然感覺不到什麼喜樂.
但是應該選擇相信神的話語.
多過感覺.
所以你要憑著意志作出一個動作.
唱詩讚美.

$^{361}$回教會來聚會.
向人宣講.
見證神過去對你的信實.
你的意志就像一個開關.
拍一聲.
接通了神幫助你常常喜樂.
這份按開關的勇氣和力量.
不是靠意志.
是靠聖靈在我們心裡面來動弓.
但是聖靈在你心裡面動弓之後.
要做什麼?.
靠你的意志來採取行動.
所以我們不是等有感覺了才去信服.
而是透過舌戟背十字架跟隨主去信服.
而是透過舌戟背十字架順服去順服.
去重尋那份跟隨主的良好感覺.
過去幸福與主同行的感覺.
我們的主耶穌基督.
祂就是意志順服的最佳榜樣.
我們熟悉的克西馬尼園.
祂在那裡感覺極其憂傷.
甚至汗如雪滴.
但是祂向父神禱告說.
然而不要照我的意識.
只要照你的意識.
同樣保羅寫這封信的時候.
他的現場是監獄.
他的鎖鏈是他的現實.
但是他證明了.
喜樂可以跟外在的境況脫鉤.
他在枷鎖裡面.
可以選擇禱告.
可以選擇唱詩.
可以用讚美的行動.
去扣開他內心那個因情緒而關閉的門.
打開了天國的門.
當然那道門是耶穌基督開了的.
他邀請你走進去.
但是你我因著情緒整天都果步不前.
保羅示範了.

$^{401}$意志是信心的腳步.
喜樂沒有意志的選擇.
就是虛空和虛幻.
弟兄姊妹.
每一天我們都面臨選擇.
我們是順從感覺的引導.
讓情緒被環境牽著走.
做你所有決定的主人嗎?.
還是順服神的話語.
讓神的真理居住在你和我的內心.
做每一件事決定的標準.
神已經將得勝的方法交給我們了.
所以今天當你憂愁的時候.
我要請你開口讚美.
當你心裡面有怨言的時候.
你可以選擇數算神在你生命當中的恩典.
當你絕望的時候.
當你灰心的時候.
你應該選擇宣告神的應許.
這個不是心靈雞湯.
這個不是阿Q精神.
這是建立在基督已經得勝的事實上.
憑信靠耶穌基督的信心.
踐踏世界的謊言.
奪回神應許之地的實戰.
弟兄姊妹,你願意嗎?.
你的意志就是指向聖靈的箭頭.
當你憑著你的意志做出選擇順服.
你就停留在神聖的循環裡面.
什麼循環呢?.
順服,經歷神,信心增長.
更深的順服,更深的經歷神.
更多的信心增長.
這個循環.
幫助我們從喜樂一開始我們抗拒.
喜樂這個命令一開始我們抗拒覺得很艱難.
但是只要你開始慢慢地學習.
進入這個循環.
你的內在就會慢慢生命變得自然.
自然順服上帝.

$^{441}$敢於順服上帝.
因為你嘗試過甘甜.
最終你能夠成為榮耀上帝的見證.
弟兄姊妹.
我們來到結束的時候.
可能你心裡面覺得很開心.
有飯吃了.
其實我們最重要的是回顧一下我們今天的經文.
我們所有的艱難.
需要用盼望來渡過.
我們在惶恐和淚水當中宣告盼望.
還是在黑暗和幽閉裡面.
我們要站穩我們的腳步.
守護住內心聖靈微弱的聲音.
不是聖靈的聲音很微弱.
是我們總是不聽所以微弱.
今天我們在基督裡面宣告.
喜樂是在我們軟弱的情緒裡面.
我們能夠宣告重新來建立.
這個生命力來自至高的上主.
這兩者是互相不衝突的.
什麼不衝突?.
盼望和喜樂.
因為在基督徒的喜樂.
是有過去,現在,未來的三重奏.
我們為過去主耶穌基督.
已經成就的救恩而喜樂.
我們現在又為主耶穌基督.
真實的與我們同在同行.
而喜樂.
然後我們展望未來.
耶穌基督要榮耀的再來.
終結一切的不公平.
罪惡.
一切人間的醜陋.
那一份喜樂.
超過一切.
所以我們要記住這份常常喜樂.
我將今天所有的信息.
濃縮成四句的詩詞.

$^{481}$希望我們一起來讀.
無論是廣東話.
或者普通話.
我們一起來讀.
「反轉人生」.
可能你覺得很老套.
但卻是我最習慣來寄東西的方式.
可能你不記得不要緊.
你用你的方式來記住吧.
這篇信息是上帝對我說的.
今天你聽的信息是上帝對你說的.
這篇講稿遲了半年講給大家聽.
因為我有我的能力的限制.
我有我的諧諾的限制.
但是神的恩典從來不遲到.
讓我們記得.
我們不應該再被環境牽著走.
我們不應該讓情緒來做主人.
讓它來做信差就好了.
讓我們願意放下被動的情緒.
可以主動地選擇「我選擇喜樂」.
你就已經開啟了一份生命的顛覆.
願你每一次再說一次的喜樂.
成為你每天早上醒來的禮物.
你的心靈跟神接通了.
你就有勇氣來做出決定.
你就能夠持續地走下去.
願主祝福大家.
我們一起來禱告.
親愛的阿爸父.
我們來到你的面前來承認.
我們的喜樂常常跟著環境起起伏伏.
我們太容易被生活裡面的不如意.
牽著我們的鼻子走.
主啊.
求你給我們有勇氣.
讓我們學會即使在困難裡面.
也能夠憑著信心選擇喜樂.
我不想再靠自己的心情來過日子了.
而是我們將我們自己的心.

$^{521}$重新連上跟你永遠連線的WIFI.
求你幫助我們.
讓這份在主裡的喜樂.
成為我們每天生活真實的力量.
我們能夠如此祈禱.
來奉我主耶穌基督的名而求.
\newpage



\section{士師記 1:1-4}
\label{sec:032085}
\textbf{士師記析讀 (4) 互相扶助 迎挑戰}~~~~梁永善牧師
\newline
\newline
連結: \href{http://www.church.com.hk/acms/content.asp?site=cdc\&op=show\&type=product\&code=032085\&layout=sermon}{\texttt{032085}} ~~~~ 語音日期: 2026-05-06
\newline
\newline
\hyperref[sec:032084]{\small{< < < PREV SERMON < < <}}
~
\hyperref[sec:index_chronic]{\small{[返順時目]}}
~
\hyperref[sec:032086]{\small{> > > NEXT SERMON > > >}}
\newline
\newline
士師記 1:1-4
\newline
\begin{longtable}{cl}
\hline
\hline
章節 & 經文 (和合本修訂版)\\
\hline
1:1 & \begin{tabularx}{0.7\textwidth}{X} 約書亞死後，以色列人求問耶和華說：「我們中間誰當首先上去攻打迦南人，與他們爭戰呢？」 \end{tabularx} \\ \\ \relax
1:2 & \begin{tabularx}{0.7\textwidth}{X} 耶和華說：「猶大要先上去。看哪，我已將那地交在他手中。」 \end{tabularx} \\ \\ \relax
1:3 & \begin{tabularx}{0.7\textwidth}{X} 猶大對他哥哥西緬說：「請你同我上到我抽籤所得之地，與迦南人爭戰；我也同你去你抽籤所得之地。」於是西緬與他同去。 \end{tabularx} \\ \\ \relax
1:4 & \begin{tabularx}{0.7\textwidth}{X} 猶大就上去，耶和華把迦南人和比利洗人交在他們手中。他們在比色擊殺了一萬人。 \end{tabularx} \\ \\
[1ex]
\hline
\hline
\end{longtable}
$^{1}$我仍然很感受上個星期和大家分享加勒的生平.
今天我們會看回題目是「互相輔助迎挑戰」.
我盡量去到七個字.
我們主要集中在《士師記》第一章第一到第四節和第十九節.
我們先看到一件事是什麼呢?.
就是大張利虛.
但是他們更加倚靠神.
以色列人出埃及入迦南地的時候.
兩個很卓越的領袖都帶領著他們.
摩西就帶領以色列人對抗法路.
然後離開了法路管轄.
終於得到自由.
大地在曠野漂流四十年.
進入迦南地的時候.
就由他的接棒人約書亞來帶領.
但是在《士師記》一開始我們看到.
約書亞都離開世界了.
但是似乎沒有一個很突出的領袖帶領著他們.
我們可能都會想迦納都可以.
迦納和約書亞同年代.
唯有他們兩個可以入迦南地.
但是神沒有揀選迦納做他們的領袖.
上次我們看到迦納我很欣賞他.
他一官半職都沒有.
沒有做賦稅.
但是他沒有埋怨.
他仍然堅守他自己的立場.
還有他自己的使命.
我們看到沒有一個很突出的領袖.
他就陷入一個危機裡面.
正所謂蛇無頭而不行.
會不會大象一去.
馬上就風雨飄搖呢?.
但是危機往往會成為另外一種祝福.
更加他們這班人沒有一個很明顯的領袖.
帶領著他們的時候.
他們會一些人懂得更加倚靠神.
他們不會將焦點集在一個人的身上.
甚至將焦點放在神的身上.
我們看看《士師記》第一章第一節.

$^{41}$約書亞死後.
以色列人求問耶和華說.
我們中間誰當首先想去攻擊迦南人.
與他們爭戰?.
《士師記》一開始就講到約書亞死了.
所以他們沒有一個突出的領袖.
在一個完全沒有領袖興起的年代裡面.
神的指紋就一定要永遠認定.
真正帶領著我們的不是人.
是我們的神.
所以他們同心去求問神.
怎樣面對前面的挑戰.
是由哪一個支派的人先去打仗呢?.
你想想如果有約書亞在的時候.
他們就會聽約書亞的命令.
由約書亞去祈禱.
然後看看神怎樣帶領約書亞.
然後他們就跟著約書亞的指示去.
但是不要忘記.
人始終會有他的不足.
神始終會有他的錯.
約書亞都曾經被人欺騙的.
如果我們熟悉的時候.
欺騙人.
他們知道約書亞打贏很多的族.
他們很怕自己被滅.
去致死.
於是怎樣呢?.
他就去扮演很遠的地方.
風塵僕僕的走來.
約書亞一時不察.
因為他沒有求問神.
於是就跟他們簽了約.
以後都不可以殺他.
只可以令到欺騙人成為他們的奴僕.
及後就因為這樣的緣故.
而引起人對約書亞的不滿.
我們就不看了.
我只告訴大家.
這段經文記載約書亞第九章.

$^{81}$第一到第十八節.
這個就是很重要的.
我自己也很深體會.
特別我在COVID之前的二十年.
我是經常來英國.
每一年都有機會來英國.
就是港島.
我就見到當時的英國的華人教會.
不是這麼多的傳道人.
甚至可能在一些較為偏遠的地方裡面.
他們自己就由查經班開始.
然後有聚會,有崇拜.
沒有目視他們當中.
但是他們就互相之間的配搭.
他們一起同心協力去建立教會.
反而可能有一個很強的領袖在那裡的時候.
他就全部指望那個領袖.
所以你看到.
譬如我去過委內瑞拉.
那時候委內瑞拉也是沒有牧者.
那時候就由信徒或是執事.
來去承擔那些牧羊的工作.
探訪,去輔導,去探病.
所以有時候未必不是一件好事.
這次我們看到當約書亞死的時候.
他就一起去求問神的心意.
接著我們就看到剛才我說的.
因為約書亞他來大學去失策.
沒有好好的來大學去了解到.
被機票人騙了.
所以後來那些人都埋怨他.
我們看看約書亞記第九章第十八節.
因為會眾的首領已經指著耶和華.
以色列的神向他們起誓.
所以以色列不擊殺他們.
全會眾就向首領發怨言.
所以你記住.
他們就會向那些首領.
特別是向約書亞來發怨言.
記住每一個牧者包括我在內.

$^{121}$一定有不足的地方.
有缺陷的地方.
所以我們千萬不要將我們的信心.
放在人的身上.
叫仰望神.
特別是一些重要的事情.
就要尋求神的心意.
就要去信服神.
我們再看看這一章的聖經.
第一章第一節就說到.
當約書亞死了之後.
他們就要去求問神.
應該是哪一個支派來打這場仗.
當他們向神禱告的時候祈求.
即表示他們願意尋求神的心意.
一個自以為是,獨斷獨行的人是不會祈禱的.
除非出了問題.
他才會求神來幫助他.
他們認為自己能夠做的事情.
何必要祈禱呢?.
但我想和大家搞清楚.
尋求神的心意的人.
並不一定代表他很樂意順從神的旨意.
有沒有想過.
可能我們自己也有犯這個錯.
我們心裡有時候有自己的意願.
如果神的旨意和我們心裡所想的事情.
是相符的,吻合的.
我們當然願意遵從.
主你清楚地帶領.
我一定會樂意去遵從的.
但如果他們的心意不相符的時候.
甚至神很清楚地告訴你.
祂的意思和你所想的是背道而馳的時候.
我們會不會帶著它去否定.
這個不是神的心意.
就不去遵從呢?.
我們看看.
當猶大國滅亡的時候.
那些精英都老走了.

$^{161}$剩下的那些都是基層的人.
或者老弱,貧窮的人.
留在耶路撒冷.
其實是頹垣敗瓦 很慘.
當時有一個軍長.
他來求問耶和華.
究竟應不應該留在那裡.
有時候去埃及呢?.
這個軍長就去找耶利米.
跟耶利米說.
你可不可以來為我們禱告?.
在耶利米書42章第3節.
「願耶和華您益臣子使我們所當走的路.
所當做的事」.
這些軍長就叫耶利米為他祈禱.
耶利米很樂意的去幫助他們.
所以後來就為他們祈禱.
十天之後.
神就有那個答案給他們了.
耶利米就跟他們說.
神很清楚告訴他們.
你們不要去埃及.
要留在這個地方.
神是會祝福你的.
如果你去埃及的時候.
就會遇到這些災禍.
但是這幫人.
求神來指引我們走這條路.
誰知耶利米給他們答案.
「你講大話」.
很熟悉的一句話.
他就說到.
「你講大話 神沒有向你離帶去啟示.
你想害死我們而已」.
他就馬上否定了神的旨意.
我們看看耶利米書43章第2,第3兩節.
我們一起來看 好不好?.
「菩薩雅的兒子阿撒利亞和阿加利亞的兒子約哈蘭.
並一切旁傲的人.
就對耶利米說 你說謊言.

$^{201}$也會說我們的神並沒有差遣你來說.
你們不可進入埃及在那裡居住.
只是尼利亞的兒子巴魯教唆你害我們.
要將我們交在加勒底人的手中.
使我們有被殺的.
有被擄到巴比倫去的」.
這幫人表面上好像很願意順從神的心意.
但是當神的意思和他們所想的是相違背的時候.
他就拒絕了 甚至說.
「耶利米 你講大話.
其實神沒有向你離帶去啟示.
你想害我們而已」.
那會不會有時候我們也是這樣呢?.
我們要尋求神的心意.
但原來我們有自己既定的立場.
就好像這幫人這樣.
結果這幫人怎麼樣?.
就下埃及 還捉了耶利米下去.
結果他們就被殺了.
所以我們願意有一個順從神的心.
我們當然每個人都有自己的前設.
但是如果很肯定神的心意.
我們就要將自己的前設放下.
所以我們看回事實.
但今次我們看到有美好的一面.
他們樂於遵行神的心意.
當神指示猶大支派去作戰的時候.
他們沒有異議.
他們就會遵從.
特別是猶大支派沒有異議.
其他的支派就反而簡單.
「幸好沒有叫我」.
這時候就叫猶大.
他們當然樂於遵從.
但是猶大支派呢?.
它沒有離帶去抗議.
「為什麼是我?為什麼不是留便?」.
因為留便才是大兒子嘛 是不是?.
看回上次再上次.
我就有解釋過為什麼會找猶大支派.

$^{241}$一會兒都會再提一提的.
好了 他們要打仗.
但是不要忘記.
原來他們還沒打已經先贏了.
雖然還沒打先贏 但他們仍然要作戰.
當他們還沒作戰的時候.
神就已經給他們有應許.
他們一定會得勝的.
第一次就說了輸就死了.
他們求問神.
哪個支派上去打仗?.
神就說猶大支派.
好了 我們看看第二和第四節.
「耶和華說:猶大當先上去.
我已將那地交在他手中.
猶大就上去.
耶和華將加林人和比利西人交在他們手中.
他們在比利西逼殺了一萬人」.
看不看到?.
他們還沒打這場仗.
神就說:我會將他們交在你手中.
神已經給他們有保證.
一定會得勝.
神是會親自為他們作戰的.
將敵人殲滅.
但事實上過去的日子.
他們都經歷過這些情況.
面對龐大的敵人.
他們身處弱勢.
他們沒有好的軍事訓練.
他們沒有精良的武器.
但是雞蛋可以扔爛夠長的.
正如亞伯拉罕 大家記不記得.
他帶著318個僕人 所以亞伯拉罕多有錢.
僅僅是僕人都318個 打贏了四王聯盟.
有四個宗族 或者四個不同的民族.
他們聯手 但是亞伯拉罕帶著318個家丁.
就贏了他們 救回了羅德.
還有 埃及當時是強盛的國家.
但是以色列人卻是可以打贏他們.

$^{281}$紅海的水將他們淹沒 當然不是靠他們自己.
因此他們的勝利是什麼呢?.
是神的恩典 是神的作為.
但是不要忘記 他們絕對不能夠袖手旁觀.
他們必定要很積極的投入.
他們得勝 不是因為他們人強馬壯.
有精良的配備 而是因為是神的恩典.
雖然在恩典裡面 他們仍然要努力去作戰.
今天 神豈不是應許了我們.
我們是可以戰勝撒旦的嗎?.
當然要依靠神了.
難道我們會懷疑神的能力是低於撒旦的嗎?.
不會的 我想我們每一個機會都相信.
神的能力一定是超越了魔鬼.
但是今天我們能夠得勝撒旦的仇敵.
罪惡的試探 忍有 仍然是神的恩典.
但是我們仍然很需要努力追求.
好好的學習 多親近神.
以致我們內裡的那種屬靈生命是堅強的.
那種抗疫能力都大了.
所以我們在恩典之中得勝 但是要努力.
個人 教會 包括我們的人生.
不進則退 不勝即敗.
很簡單 你踩單車.
我現在也會踩單車的 我今天也會踩單車的.
你單車停下來的時候.
很自然會怎樣呢?.
就不動了 甚至會跌倒.
如果上斜路就會倒下來.
所以求神給我們體會到.
我們是活在神的恩典之中.
但是我們仍然是要努力地作戰.
好了 第三方面就是很重要的.
今晚我們所看的.
就是要互相輔助 然後去面對挑戰.
第一章的第三節 我們會看到.
當以色列人透過禱告.
污淩土明向神祈求的時候.
神就指出猶大之派應該先去作戰.
雖然猶大不是長子.

$^{321}$我上次也講過了.
但是因為長子劉弁犯了錯.
跟自己的父親的妾侍發生了關係.
在道德上失敗.
於是他失去了這個祝福.
雅各聰最大的祝福就是給了猶大.
而且這是皇族的支派.
後來耶穌基督也是從猶大支派而生的.
所以神就安排了第一個支派去作戰的.
絕對不是劉弁.
而是猶大.
這個也引證了 我經常這樣講.
恩典不單單是美麗 也是沉重的.
很簡單 我舉一個例子.
很有意思的 原本我應該是繼隊去以色列的.
如果真是去旅行的 真是很可惜.
我想我只可以去多一次以色列.
就是我結束了 因為體力都應該再去不了了.
這時候 誰知道當我二月一停了.
決定了不去以色列.
眼睜睜收到一個北愛爾蘭的姐妹.
我看回那個記錄.
她在去年11月跟我說了一個很簡單的hi.
她跟我說 多了香港人的團體.
牧師你可不可以過來跟我們分享一下.
我說好啊 其實我只是說兩堂而已.
告訴大家 剛剛過去的星期六.
我是非常累的 那個累是甚麼呢?.
就是喉嚨很累 原來有人很想見我.
早上已經約了八點多吃早餐.
八點多見完她之後.
有第二批人來 其實我不知道是哪些.
有六,七個人來 還有一些是未信的人.
要跟我們分享信仰 澄清他們對信仰的誤解.
我由八點多一直說了一點多鐘.
就很累 喉嚨很累 然後去到教堂.
三點多鐘就是聚會 我坐在那裡想休息一下.
誰知道都不夠五分鐘 又有些人說.
牧師啊 他很開心 想不到.
現在正在來的那對夫婦 他們是未信的.

$^{361}$你跟他們聊聊天吧 於是我就跟他們聊天了.
總之就是除了說了一個多小時.
我由早上八點多去到六點十分 為什麼呢?.
因為六點十分我就坐渡輪 坐旅遊巴.
從北岸去南岸 那個時候可以終於休息一下.
那你說這個是恩典還是什麼呢?.
是啊 我很開心 可以跟一些未信的人聊天.
甚至有些很特別的 一些自閉的人.
都跟我聊得很好 連媽媽都覺得很奇怪.
怎麼可以跟牧師這麼好的溝通呢?.
我覺得這個是一個祝福 但都是沉重的.
你想一下 如果我很冷酷地對人 不理睬人.
我來港島 何必要見人呢?.
那我就舒服很多了 是不是?.
所以你記住 原來恩典是沉重的.
正如聖經有一句話 路加方第十二章第四十八節.
「因為多給誰 就向誰多取」.
沒錯 因為多給誰 就向誰多取.
所以我仍然要感恩的 猶帶順命去打仗.
但他知道自己力量不足.
於是他就向另外一個支派 就是西面.
這個支派去求助了.
好 我們看看《史詩紀》第一章第三節.
猶帶對他哥哥西面說.
請你跟我到尖溝所得之地去.
好與加林人爭戰.
以後我也跟你到尖溝所得之地去.
於是西面與他同去.
為什麼他只是問西面而不問其它支派呢?.
你記不記得我上次提到.
原來西面猶帶都是利亞所生的兒子.
即是同父同母的兄弟.
因為他有四個妻子.
12個兒子是由不同的母親生的.
但是猶帶和西面是同一個媽媽生的.
即是親兄弟.
而上一次沒有提到的 留意一下.
這次我們會再提一點 就是什麼呢?.
原來他們的地業.
西面的地業.

$^{401}$是在猶帶支派的地域之內.
再聽聽我說.
西面的地域是在猶帶支派的地域之內.
我們看看《約書記》第19章的第一節 好不好?.
為西面支派的人按著宗族.
占出第二級 他們所得的地業.
是在猶大人地業中間.
我想香港人 特別是一定年齡的香港人.
就會明白這一段的說話.
怎樣西面的地業就在猶大支派中間呢?.
記不記得在東頭村 九龍城中間.
以前有個叫九龍城寨.
九龍城寨的特別的地方在哪裡呢?.
原來當時英國沒有割讓這個地方.
所以後來這個地方就成為三不管的地方了.
某個程度英國的法例不能夠在城寨裡面運作.
所以後來九七回歸時 一定要解決這個歷史問題.
於是後來就拆了這個城寨.
原來西面支派.
他們的地業就在猶大支派裡面.
所以除了有血緣的關係 還有.
上次沒有提的 原來我們有地緣的密切關係.
所以他就向西面求助 他們答應有一天.
西面當你有需要 我們一樣會拔刀相助.
一樣會手足扶持 一起來面對這個挑戰.
所以這幅是很美麗的圖畫.
因為集合了兩個支派的力量.
耶和華就使得他們戰勝了迦南人和比利西人.
而且在比色裡面擊殺了一萬人.
我們看看《史詩》的第一章第四節.
猶大就上去 耶和華將迦南人和比利西人.
交在他們手中 他們在比色擊殺了一萬人.
大家有沒有留意到 猶大他們願意信服.
但是神沒有側避他們 要找人來幫忙的.
他們兄弟之間輔助.
所以神仍然是要耶和華將迦南人和比利西人.
交在他們手中 是不是?.
他們因此而得到這個勝利.
所以不要忘記 今天神擺我們在教會裡面.
都是讓我們彼此建立互相的扶持.

$^{441}$好了 我們就看到 不單單這樣.
我們武打出來 更加將非利士人的加薩.
亞瑟基倫和以甲倫這些這麼重要的菲利城市.
都交在他們手中 使他們大獲全勝.
所以一個很重要的教導.
我們是要很努力去做事 要盡力.
承擔自己的責任 但是也要承認.
有時候我們都需要別人來幫助我們 輔助我們.
如果已經盡力 但是還沒有成功的時候.
我們請別人來幫助我們並不為過.
我舉兩個例子 以前我覺得不要搞著人.
第一位的弟兄很多年之前幫助我學習.
如何找人來幫我.
就是當年我募會的時候 後來有了我的女兒.
師妹要上班 所以我們請工人.
我們請菲傭的時候 每兩年就有一個合約.
我們就會換一個工人 因為我們已經熟悉.
這些古靈精怪的了 所以我家裡有門鎖的.
每一次換一個工人的時候.
我就會拿一個鎖膽出來 換另外一把鎖.
為什麼呢?因為工人他有我們的鎖匙.
你真是以小人之心度君子之腹.
以防有時候有些人會有些古古怪怪的事.
所以我一定會換門鎖的.
我記得有一次我想換門鎖.
我搞了四十多分鐘 鎖上了就鎖不上了.
或者是鎖不上了 搞了很久 我終於很無奈.
我打電話給一位弟兄 這個弟兄很好.
就是我教會的元老 如果你們聽我說.
就是上次我去日本 然後特意有幫人飛到台灣見我.
是其中一位來的 他就在我家旁邊.
當時我住在丁雅縣 他就在太平工業中心那邊工作.
我說 啊 炳哥 我搞不定 換不到門鎖.
如果你放工 你過來幫我搞一搞.
他說 行啊 沒問題 他真是五點多放工的時候.
就來我家幫我換那個門鎖.
不用七,八分鐘 他就換了.
我說 你這麼厲害 我搞了七,八個字都搞不定.
他就說了一句話 他說 牧師 這些是我們的專長.
你讓我們做這些 沒所謂的 我們很快就做得到.

$^{481}$但是如果你叫我擦一支筆寫東西 我寫給你看.
他說 你不如將那七,八個字 這些時間.
你寫廣宗也好 你關心別人也好.
不要浪費在這些事情上.
有時候我們覺得 哎呀 不是挺好意思的.
只要我們不是自己懶 去找人家笨.
於是我們就怎麼樣呢?我們真是有需要.
你可以找人幫忙 減輕了你的負擔.
但是你就不是翹著手 你是做其它的事情.
所以這個是很重要的.
好了 還有 就是你幫了人家 人家幫了你.
但是第二天有人有需要 你照樣去幫人.
未必幫在那個人的身上.
但是你要去幫助別人.
在美恩堂 我那個時候 牧會的時候都很感恩.
我建立了教會的發展.
但是我最感到驕傲 不是招職 而是驕傲的是什麼呢?.
就是很開心 我們有三十多,四十個傳道人出來了.
當我們自己會有讀神學的時候 我就很清晰.
那個時候我就說 學費 十多億的學費我們給.
但是我也不會要求你一定要回來美恩堂事奉.
但是我會說 你要忠心的在你的工場工作.
還有第二天你有機會 你都要去幫助別人.
去輔助別人 這個就已經是足夠了.
我幫了你 你不需要回報在我身上.
只要你繼續去幫其他人.
我當年讀神學 第二年就已經沒有錢了.
為什麼呢?其實我有錢的.
但是當年我媽媽不喜歡我讀神學.
我做事的時候又說我亂七八糟.
叫我把錢給她儲著 將來可以儲老婆有錢.
那我就給錢她儲 誰知道我媽媽很早離開世界.
我還沒有畢業 她已經離開了世界.
但是無證無據 那些錢就白白拿走了.
沒有計算的 是不是?.
第二年的時候沒有錢.
但是我第一間的實習教會.
有一個弟兄 他只是小學程度.
做傭人 但是他很有愛心.
那時候他每個月給我100元.

$^{521}$來支持我讀神學 讓我買書.
還有很紀念他 我沒有還錢給他.
因為他奉獻給我 但是現在我都會怎麼樣?.
去奉獻支持有需要的人.
就是美金堂現在有人在讀神學.
他前一段時間有需要 我也有錢支持他.
是我經歷了這份的愛.
於是我去幫助同樣有需要的人.
所以這個是互相輔助.
所以猶太就說 西面之派你們來幫我.
當有一天你有需要的時候 我都會幫你.
所以今天我們要學習彼此扶持.
互相的來幫助.
但是實行的工作是要破壞神的教會.
破壞肢體之間的團結生活.
所以我們一定要同心合意 興旺福音.
還有就是什麼呢?.
招納賢能齊作戰.
就是《士師記》第一章11到第15節.
這裡就說到在戰爭過程當中.
不單單是猶太之派請求西面之派來協助.
當中記載了猶太之派.
一位很重要的人物就是迦勒.
就是上個星期我們所看的.
他都請人來協助他們的支派.
去攻打激烈西忽的地方.
好 我們看看《士師記》第一章第12節.
迦勒說:誰能攻打激烈西忽 將城奪取.
我就把我女兒革 扼殺 及她為妻.
很重要啊 誰可以幫我.
贏了將我的女兒給他 做妻子.
我就一定不會開這些條件的了.
我都很疼愛女兒的.
我怎麼知道那個是誰.
怎麼知道是不是真的好的 是不是?.
但是迦勒真的開出了高架.
及後就是俄陀列回應了這個呼召.
下個星期我們就會看俄陀列了.
下個星期我們會看三個《史詩》.
好了 他幫助攻取了迦勒 即是攻取了地方.

$^{561}$後來就成為迦勒的女婿.
而俄陀列也都成為《士師記》裡面.
第一個的《史詩》 第三章第九節.
以色列人呼求耶和華的時候.
耶和華就為他們興起一位拯救者.
救他們就是迦勒兄弟基納斯的兒子 俄陀列.
看不看到 其實他們是親上加親.
即俄陀列就是迦勒的兄弟的兒子.
即其實是堂兄妹來的 是不是?.
好了 迦勒願意將女兒許配一個有能之士.
可以說是高架招聘賢能 真的要找專才的.
這個給我們看到 原來在事奉上面.
我們很多時候是需要請一些專才來幫助我們的.
我這次在北愛的時候 很驚訝.
就是一位基督教機構 一位聖經研究的學者.
原來他都移民去了北愛爾蘭.
我們都有短暫時間的相聚.
當年我冥恩堂除了我教主學之外.
即教聖經之外 我都邀請不同的人.
特別是在那方面有專長的 就過來幫我們.
很簡單 我還很感動.
張雲開老師就是基督神學院的老師.
他晚上去大埔幫我們教希伯來書.
我在想 遲些完事時我會不會教希伯來書.
希伯來書是很不容易去解釋的.
你想一下 他在長洲住.
他七點半來到教書 九點半完.
九點半由大埔回到長洲 還要上山頂的宿舍.
是很漫長的路 他都願意.
我們真是禮賢下士.
當日我們有一個好的發展就是怎樣?.
可以有不同的人來幫助我們.
所以我們當然也要給應有的報酬.
而不能夠單單講友情的.
所以這些事奉者未必以報酬為重點.
但是他要有這份的愛.
正如我這次去北愛爾蘭.
他們說我們這個團體剛剛成立.
牧師 我們是沒有收居費給你.
甚至你來的機票都要你自己出的.

$^{601}$我自己出機票 但是很不幸.
我要找Alex幫忙 但這次真是太複雜了.
我買那些廉航 但是check in不了.
結果去到counter要check in 要給55鎊.
但是他說可以cam 今晚完了.
明天我就試一下去cam.
昨天我在北愛爾蘭等了一個小時.
但是都沒有人去回覆我.
等一等我們會有真人 不是那些電腦的人.
要和你對話 但是等了一個多小時都不行.
我不是為了錢過去.
但是都感恩 也有些肢體原來預備了一些愛心給我.
所以你留意到 如果他們是沒有經濟能力.
我們事奉者絕對不是在乎錢 可以去幫忙.
但是如果有能力的時候 而不去給.
這個就需要真好地面對 這是一個很大的可悲.
好了 今天教會有好的發展.
其實需要不同的人去配合.
當然這個專才一定不只是技巧上面高.
是要靈明高 這個很重要.
好了 另一方面就是他們未克全功 需要再聘.
在《士師記》第一章第四到第十八節.
記載了猶太人攻取了不少地方.
包括我們熟悉的希伯倫.
為什麼是希伯倫呢?.
要再看回上兩次跟大家的分享.
解釋了為什麼他們要去希伯倫.
但是他們並沒有完全攻佔順位他們預備的地方.
《士師記》第一章第十九節.
耶和華語猶大同在 猶大就講出山地的居民.
小事不能講出平原的居民 因為他們有鐵車.
作者指出正面的是什麼呢? 神和猶大同在.
猶大就講出山地的居民.
事實上你要攻佔山地其實不容易的.
你知道人家在山上 你要攻上山就易守難攻.
但是他們成功了 但是很奇怪.
反而在平原那個地方.
他們卻是不能夠趕走那些人 為什麼?.
因為平原一定是適合更多人居住.
所以一定會團重兵在當中.

$^{641}$但是最大的問題是什麼呢?.
就是這些敵人有戰車 當時是最厲害的武器.
我們會認同同情他們沒有力量去克勝強敵.
但根本問題不是敵人的強 是自己弱.
弱是什麼呢? 就是信心的問題.
因為神一早已經應許他和子民同在.
我們看回《申命記》第七章第一到第二節.
都比你們強大 這七個的人民.
神已經說 是啊 他們是強啊 你們 但不用怕.
我一早已經應許過 你們會贏.
所以神一早已經告訴你們 你們一定會得到勝利.
事實上 在弱勢的帶領下 你記不記得.
他們不是贏了耶利哥城嗎?.
這個城是很難去攻陷的.
但是在神奇的帶領下 他們就贏了耶利哥城.
而且有很多馬匹車輛的下鎖王.
和其他的民族聯合對抗弱勢.
可能大家不太記得 或者不知道有這些戰爭.
《約書》第十一章第四到第五節.
「鼎西王和他們的眾軍都出來.
人數多如海邊的沙 並有許多馬匹車輛.
這諸王會合 來到米倫水邊.
一同安營 要與以色列人登戰」.
看不看到 當時有許多的外族聯手.
他們的人數多到好像海邊的沙.
而且同樣有許多的馬兵,車輛.
但是神一早宣告 你們必定會得勝.
事實上就是這樣 他們得勝是有經驗的.
他們不是未曾贏過這些強敵的.
《約書》也留下了寶貴的指令給他們.
「加拿人雖有鐵車 雖是強盛.
你也能把他們趕出去」.
記不記得《約書》的遺言跟我們說.
雖然加拿人有鐵車 雖然他們很強大.
但是放心 你們都能夠趕走他們的.
鐵車雖然厲害.
但是並不等於他們沒有機會得勝.
稍後我們就會看到《士師記》記載了.
耶賓王 他有鐵車900輛.
大大欺壓以色列人二十年.

$^{681}$我們看看《士師記》第四章第三節.
耶賓王 有鐵車900輛.
他大大欺壓以色列人二十年.
以色列人就呼求耶和華.
但是神後來又興起底波拉.
最終神用奇妙的方法.
使得耶賓王他們潰敗.
特別是元帥叫西西拉 一樣的死.
好 我們看看《士師記》第四章第十五節.
耶和華使西西拉和他一切車輛全軍遣亂.
在巴拉面前被刀殺敗.
西西拉下車步行同跑.
看不看到900輛戰車又如何.
一樣可以輸.
所以他們過去曾經打過幾次勝仗.
但是已經慢慢安定下來了.
因為安定了 老實說.
就不想再打下去.
這些所謂的滿足 安穩.
留在安全網之下的心態.
往往就令人裹足不前 不再冒險.
但是你不去前進 就會有下滑的危機.
所以求神幫助我們.
有時候我們在安全網裡面.
我們有時候會鬆懈.
所以我離開香港 都是離開了我的安全網.
特別是頭一年多來到這裡的時候.
都很辛苦 要重新適應過.
雖然我以往經常都來英國港島.
但是和你居住是不同的.
特別租地方的時候 家徒四壁.
一個一件的傢俬都沒有.
要自己買傢俬來哄.
以前是有人幫我打字 幫我做powerpoint.
但是現在什麼都沒有 那怎麼辦?.
所以是有困難的.
但是正正離開了那個舒適圈.
我是經歷到神更多.
某個程度我靈明會更大的一些突破.
很多時候我們見到敵人很龐大的勢力.

$^{721}$於是很難 根本不可能贏的.
於是這些說話就會出現在我們口中.
或者心裡面.
我們更加為自己不想前進來找藉口.
猶大之派不願意為征服迦南.
付上更大的代價.
故此缺乏了能力.
並不是因為他們失敗的因素.
反而是他們沒有信心.
不想付太大的代價.
而成為他們不能有得勝的原因.
神的能力 記住.
不會因為環境,困難,時間有所改變.
問題是我們對神的應許.
有多少信心和毅力.
我們都要有這個警覺.
我們絕對不能夠滿足於現狀.
仍然要努力的去前往.
仍然有很多人未聞福音.
我們仍然要繼續努力.
為神的角度去打仗.
拓展神的角度.
我們都見到前面好像.
有些無法克服的惡劣的環境和挑戰.
但我們不能夠得勝.
取決於除了我們對環境的觀察和和應之外.
我們更加需要將焦點放在神的手上.
我們每一個人將危機視為轉機.
並且以非凡的方式.
和竭盡所能的努力來事奉神.
這個就是很重要.
所以求神幫助我們.
透過這個很簡單的分享.
讓我們記住.
我們要互相守望.
今天神擺在你的群體裡面.
就是讓你成為別人的天使.
也相信別人會成為你的天使.
你有能力去幫助有需要的人.
但當你經歷到別人給你的幫助.

$^{761}$你也需要去幫助別人.
所以由大西面去彼此配合.
是一幅很美麗的圖畫.
希望我們能夠有這個體會.
\newpage



\section{馬太福音 6:25-33}
\label{sec:032086}
\textbf{先求祂的國}~~~~張健庭牧師博士
\newline
\newline
連結: \href{http://www.church.com.hk/acms/content.asp?site=cdc\&op=show\&type=product\&code=032086\&layout=sermon}{\texttt{032086}} ~~~~ 語音日期: 2026-05-03
\newline
\newline
\hyperref[sec:032085]{\small{< < < PREV SERMON < < <}}
~
\hyperref[sec:index_chronic]{\small{[返順時目]}}
~
\hyperref[sec:032087]{\small{> > > NEXT SERMON > > >}}
\newline
\newline
馬太福音 6:25-33
\newline
\begin{longtable}{cl}
\hline
\hline
章節 & 經文 (和合本修訂版)\\
\hline
6:25 & \begin{tabularx}{0.7\textwidth}{X} 「所以，我告訴你們，不要為你們的生命憂慮吃甚麼喝甚麼，或為你們的身體憂慮穿甚麼。生命不勝於飲食嗎？身體不勝於衣裳嗎？ \end{tabularx} \\ \\ \relax
6:26 & \begin{tabularx}{0.7\textwidth}{X} 你們看一看那天上的飛鳥，也不種也不收，也不在倉裡存糧，你們的天父尚且養活牠們。你們不比飛鳥貴重得多嗎？ \end{tabularx} \\ \\ \relax
6:27 & \begin{tabularx}{0.7\textwidth}{X} 你們哪一個能藉著憂慮使壽數多加一刻呢？ \end{tabularx} \\ \\ \relax
6:28 & \begin{tabularx}{0.7\textwidth}{X} 何必為衣裳憂慮呢？你們想一想野地裡的百合花是怎麼長起來的：它也不勞動也不紡線。 \end{tabularx} \\ \\ \relax
6:29 & \begin{tabularx}{0.7\textwidth}{X} 然而我告訴你們，就是所羅門極榮華的時候，他所穿戴的還不如這些花的一朵呢！ \end{tabularx} \\ \\ \relax
6:30 & \begin{tabularx}{0.7\textwidth}{X} 你們這小信的人哪！野地裡的草今天還在，明天就丟在爐裡，神還給它這樣的妝飾，何況你們呢？ \end{tabularx} \\ \\ \relax
6:31 & \begin{tabularx}{0.7\textwidth}{X} 所以，不要憂慮，說：『我們吃甚麼？喝甚麼？穿甚麼？』 \end{tabularx} \\ \\ \relax
6:32 & \begin{tabularx}{0.7\textwidth}{X} 這都是外邦人所求的。你們需要這一切東西，你們的天父都知道。 \end{tabularx} \\ \\ \relax
6:33 & \begin{tabularx}{0.7\textwidth}{X} 你們要先求神的國和他的義，這些東西都要加給你們了。 \end{tabularx} \\ \\
[1ex]
\hline
\hline
\end{longtable}
$^{1}$弟兄姊妹平安.
我們一起低頭祈禱.
天父我們願你的話語再一次的.
向我們講說話.
主讓我們明白你的心意.
也讓我們聽到你的聲音的時候.
我們愛慕.
我們也是跟從.
我們這樣祈禱.
奉耶穌基督的名祈求.
阿們.
今天就和大家講這個題目.
就是先求他的國.
為甚麼會選這個題目呢.
其實大家都知道.
其實今天開始.
或者明天開始.
我們會有一個.
靈修讀一本書的運動.
希望藉著有30天.
我們一起來思想一本書.
怎樣去親近神.
那本書的中文譯名.
就叫做移動的聖所.
完了之後.
在後面的10元一本.
希望大家盡量去參與.
去買一本書.
每天會讀一篇.
有30篇.
在30天裡.
我們一起來思想.
怎樣去追求神的國.
其實那本書的英文名叫.
Seeking the Kingdom.
所以就是今天的題目.
去追求神的國.
因為我們活在世上的時候.
我們要知道.
主人想我們怎樣來生活.

$^{41}$我想這個也是一個.
值得停下來思想的題目.
所以我們完了之後.
就在後面的Foyer.
每人去買一本書.
10元.
如果你是需要看英文的.
我們都是.
其實那本書原本是英文.
有英文的.
但是.
英文售賣的書.
就在英文堂那邊.
所以如果你要買英文的.
或者完了之後.
我再和你談談.
如果你要電子版.
其實都有的.
那本書可以買電子版.
如果你覺得看電子版.
方便一些的話.
剛才讀那一段經文.
是《馬太福音》第六章.
很熟悉的一段經文.
就是說在結束的時候.
登山寶訓當中.
其中一段很人記得的經文.
「你們要先求他的國和他的義」.
「這些東西都要加給你們了」.
但是.
說一說究竟.
關於主耶穌為什麼會說.
前面這一段經文.
因為其實主耶穌是知道.
神的國.
人是很想去追求.
但是好像有很多的難阻.
究竟是怎樣去追求神的國.
以致到.
所以他就說出剛才這一段經文.

$^{81}$或者.
未正式進入神的國的題目的時候.
我想問一下大家.
其實當你看到這本書的時候.
你會發覺有三樣東西.
主耶穌經常都做的.
那是哪三樣東西呢?.
大家知不知道呢?.
有三樣東西主耶穌經常都做的.
第一樣我先告訴你.
第一個送給你的答案.
就是主耶穌經常都講道.
那你說當然了.
但是主耶穌經常講道的主題.
其實都是關於神的國的.
天國,神的國.
這就是第一樣東西.
在福音書那裡.
你讀福音書必然會有一個這樣的印象.
主耶穌經常都講神的國.
一開始出來傳道的時候.
就已經講天國近了.
你們應該悔改.
就是神的國.
第二樣東西呢.
你看福音書的時候有什麼印象呢?.
主耶穌是做什麼的呢?.
做什麼呢?.
很多時候你看四福音的時候.
你會有一個印象.
就是神蹟.
神蹟當中最多的.
就是醫病和講鬼.
很多的.
醫病,講鬼.
那第三樣是什麼呢?.
經常都做的.
是什麼呢?.
可能你很少想.
但是其實你經常看到主耶穌經常做的.

$^{121}$禱告都有.
但是不是那麼多.
比我講的第三樣是多一點的.
就是什麼呢?.
教導.
教導到算是講道.
就是吃飯.
你有沒有發覺.
經常都吃飯.
然後那些人又叫主耶穌.
五餅二魚其實都是關於吃飯的.
五餅二魚都是關於吃飯的.
然後跟衰李又吃飯.
然後吃飯的時候又有些討論.
經常吃飯的.
是不是?.
其實這三個福音書當中.
你看到主耶穌很多時間.
很長篇幅來講的.
其實你有沒有想過.
這三個題目.
或者這三樣東西.
有沒有什麼共通呢?.
有什麼共通呢?.
其實就是剛才所講的那樣.
其實共通就是神的國.
在講道那裡是自然的.
因為主耶穌已經講明.
神的國是怎樣怎樣.
然後那些天國的比喻.
各樣都是講神的國.
但是你說醫病講鬼呢?.
醫病講鬼其實都是講神的國的.
因為他就是要講到神的國.
現在已經彌臨了.
當神的國彌臨的時候.
黑暗的權勢是被驅逐的.
黑暗的權勢被驅逐的時候.
有什麼實際的.
被人看得到的發生呢?.

$^{161}$就是魔鬼的勢力開始被驅逐.
被打敗.
這個就是講鬼.
而罪惡的勢力.
被驅逐的時候.
特別是疾病那些.
這個世界.
是因為罪進入了世界才會這樣.
所以就醫病講鬼.
這個世界的苦難各方面.
其實就是神的國彌臨的時候.
這些的.
無論是苦難.
無論是魔鬼的勢力.
都是被驅逐的.
所以第二個.
關於醫病講鬼那些的神跡本身.
其實是告訴人知道.
是實體看得到.
神的國現在已經好像一步一步進入人間.
而黑暗的勢力被驅逐了.
第三個你說.
吃飯呢?.
吃飯其實是關於神的國的.
你有沒有想過?.
因為這個.
在舊約的先知那裡.
其實多次的講到.
神的國彌臨.
就好像一個大延籍那樣.
一個很大的延籍.
當神的國彌臨的時候.
最終.
魔鬼或者是仇敵的勢力被打敗的時候.
最終神的子民.
就一起慶祝.
或者進入神的國.
好像進入一個大延籍.
所以主耶穌都不止一次用延籍的比喻.
講了一些教導.

$^{201}$而他實際來說.
他自己又不斷地跟別人吃飯.
當然在猶太人來說.
他們跟別人吃飯.
也是代表著一種.
生命的一種交流.
所以猶太人基本上.
不太跟外邦人吃飯.
因為他們的身份是.
為神分別為聖.
跟外邦人在生命上.
不想有那麼多交流.
所以他們一般來說.
不太願意跟外邦人吃飯.
所以你去《仕途行傳》讀到.
為什麼關於吃飯那裡.
彼得要去高利樓那裡吃飯.
那裡有很大的心理抗拒.
甚至乎去到加拉太書那裡.
講到保羅罵彼得.
怎麼跟外邦人吃飯.
然後慢慢退後.
諸如此類的.
就是因為吃飯的問題.
本身就是.
也是講到神的國的預表.
生命在神的國裡面.
屬神的子民.
是一個歡樂的情景.
好像一個羔羊婚宴的延籍.
當中一起來分享生命.
三方面其實都是關於神的國.
所以神的國本身就是.
整個福音書的主題.
其實也是延伸到新約的主題.
而這個主題本身.
很值得我們去關注.
究竟我們現在屬神的人.
其實我們就是天國的子民.
我們屬於天國的子民.

$^{241}$我們是怎樣來活出天國的樣式呢.
來成為神國的一份子呢.
是很值得我們去關注.
所以一直在三十日看那本書的時候.
其實就分了三方面.
第一就是向內.
第二就是向上.
然後就是向外.
去追求神的國的三個向度.
所以一直希望大家.
一直跟著讀的時候.
也是一直有機會去反省.
我們怎樣去追求神的國.
回到這裡的時候.
主耶穌其實是很明白的.
我們追求神的國的時候.
是很重要的.
所以祂說你們要先求祂的國和祂的義.
但是剛才你讀那段聖經.
我相信你之前也一定會讀過.
也是很有印象深刻的.
主耶穌就說叫人不要憂慮.
祂說要先求祂的國和祂的義之前.
就叫人不要憂慮.
為什麼呢.
因為似乎人對於活在世上的需求.
需要本身.
會產生一種對神的國可能的障礙.
所以祂就要先講了.
你不要這樣.
然後才追求神的國.
那什麼會對於追求神的國有障礙呢.
祂就有前面所說.
祂說就是憂慮吃什麼穿什麼.
是不是.
在這段經文這裡會這樣說.
憂慮吃什麼穿什麼.
那你說憂慮吃什麼穿什麼.
為什麼會阻礙人去追求神的國呢.
事實上是很真實的.

$^{281}$就是我們活在世上的時候.
有我們肉體的需要.
有我們世間各方面的需求.
是不是.
那這個世界的需要本身.
都是我們活在一個物質層面的世界的需要.
你要吃飯要喝水要吸氣.
是不是.
所以主耶穌就不是說.
你不用吃了.
或者不用穿衣服了.
不是.
但是祂的意思就是說.
你不要讓世間的憂慮本身.
讓世間的需求本身成為了一種憂慮.
來阻礙你追求神的國.
在這裡.
我相信大家在這裡都會看得明白.
不要將世間的需要本身.
成為了一種阻礙.
去追求神國的一個障礙.
我記得那時候.
在香港年輕的時候.
那時候還在讀書.
有一次在街上走.
在街上走的時候.
看到路邊有一本聖經.
那你知道了.
其實我那個時代.
現在的聖經可能都.
可能大家很少見實物聖經.
因為很少人買.
或者怎麼樣.
但是我那個時代.
讀中學那個時代.
那些聖經是很明顯的.
你離遠已經知道那本是聖經.
為什麼.
因為黑色的皮面.
不是皮面.

$^{321}$是黑色的面.
然後紅色的邊.
其實沒有這樣的書.
所以很離遠.
在街上看到街角有一本這樣的東西.
你就已經知道那本是聖經.
我是基督徒.
看到在街上有人扔了一本聖經.
就不太想聖經被人糟蹋.
就走過去.
就撿起來.
撿起來.
看到都OK.
沒有被人撕了幾頁.
就拿回家.
能用就用.
聖經多一本都無所謂.
就不用糟蹋那本聖經.
其實都OK的.
整本聖經很完整.
但是後來讀的時候.
看到前頭第一頁.
有個人寫了一些東西.
應該是一個基督徒.
送這本聖經給一個非基督徒.
他就寫了一篇聖經在臉上.
他就說.
人活著不是靠食物.
乃是靠耶和華口裡所出的一切話.
有沒有什麼問題.
有些人聰明就聽得出.
我再讀一次.
人活著不是靠食物.
乃是靠耶和華口裡所出的一切話.
有沒有什麼問題.
漏了一個單字.
人活著不是單靠食物.
乃是靠耶和華口裡所出的一切話.
這本聖經他送給別人.
然後抄漏了一個字.

$^{361}$抄漏了一個單字.
我有一次跟朋友說笑.
你看看這本聖經.
難怪被人扔了.
因為寫漏了一個單字.
這是開玩笑的說.
但其實是很重要的.
聖經的單字是說出了一件事.
其實聖經不要說.
我們不是追求世間的需要.
我們也需要的.
我們活在世上的時候.
我們世上很多的事情.
我們實在是需要操心的.
就好像我們人活著不是單靠食物.
也要靠食物的.
如果變成人活著不是靠食物.
那就錯了.
所以我們雖然不是單靠食物.
我們也要靠食物.
但問題就是.
我們所操心的事情.
那些世上的需要.
很容易就成為了霸佔我們心的憂慮.
成為了追求神國的一種障礙.
所以在《馬太福音》第十三章.
主耶穌就用了一個撒種的比喻.
來說出了那個困境.
那個困難.
撒種的比喻其實是主耶穌.
自己親自去解釋的比喻.
沒有什麼解錯的空間.
因為主耶穌自己已經解了話.
很多比喻本身主耶穌沒有解話.
就有一些解釋的空間.
但這個撒種的比喻是沒有的.
主耶穌自己解話解到怎樣呢.
他就說那些種子.
就是一個人出去撒種.
種子就是神的道.

$^{401}$泥土就是人的心.
然後結出來的果子就是.
神的道在他的生命當中.
會不會有發芽生長的結果.
這個就是主耶穌自己解釋的.
就沒得質疑的.
主耶穌就說到.
有一種種子是發芽了.
但是生長不好的.
那種種子是怎樣呢.
那種種子是落在一個荊棘的土地上.
那種種子落在荊棘的土壤上.
它是會發芽的.
它會長出來的.
但是它長到差不多.
就沒得生了.
因為那些荊棘就擠住了它.
就看不到果子.
主耶穌自己解釋.
自己解釋說什麼呢.
撒在荊棘裡的就是人聽了道.
後來有世上的思慮.
錢財的迷惑把道擠住了.
不能結實.
主耶穌自己解釋.
主耶穌就解釋到.
世上的思慮.
錢財的迷惑.
將人心裡面所接受到.
神的道的種子擠住了.
人的心不能再結果子.
這是一個很圖像性的圖畫.
你想像一下.
如果你自己的心.
或者一個人的心.
就好像長了一些荊棘擠住了一樣.
無論你聽了多少道.
神的話是不會有果效的.
結不了果.
那些種子能不能下到心裡呢?.

$^{441}$下了.
但是生兩生就沒有了.
生兩生停了就沒有結果了.
主耶穌自己去形容一個這樣的情況.
所以在《馬太福音》第六章.
祂之前為世上的食物.
衣裳去憂慮.
主耶穌是很明白的.
我們活在世上.
有這個世上的需要.
我們廣東話說得很貼切.
我們上班叫做「搵食」.
是不是?.
我們上班叫做「搵食」.
因為吃東西就每個人都要.
所以主耶穌沒有說不用吃的.
滿足人的需要.
其中一個很偉大的神蹟.
唯一的神蹟在四福音有記載的.
就是五餅二魚.
都是為了滿足人吃東西的需要.
主耶穌沒有說不用吃的.
但是就是去到一個地步.
世上的需要.
但是可能會變成了一種.
擠著人心的一種荊棘.
所以祂說.
你們要先求他的國和他的兒.
這些東西都要加給你們了.
這裡祂說了一個字.
是很有意思的.
祂說你們要先求他的國.
和他的兒.
那裡是說了一個先後的次序.
有先就是有後.
所以祂不是說不求其他的.
但是在各樣的需求都需要滿足之餘.
有一個次序是要先擺好.
那個次序就是要先擺好神的國和神的兒.
你擺好了這個次序.

$^{481}$你就不用怕.
後面的東西就會加給你了.
你沒有擺好這個次序.
你憂慮的東西.
你未必有的.
正如前面所說.
誰能使你的思慮.
使你的受數多加一克.
沒有人有的.
你思慮都未必有你所要的東西.
但是你的思慮本身.
你沒有把神的國先擺好的時候.
那些就會變成了荊棘.
會擠於你的心.
主耶穌就說了一個這樣的情況.
我記得那時候.
未出來讀神學之前.
聽到神的呼召.
就覺得神都呼召我要去做傳道.
我原本在香港有我固定的職業.
有我的工作.
如果要出來全時間讀神學.
然後去事奉.
第一步當然是要辭職.
先離開工作.
那時候讀神學.
現在讀神學的模式都有點變了.
很多人讀神學都不是全職的.
那時候我年代很多讀神學的人都是全職.
就是要辭了工作.
當中都有一些憂慮.
就是都很沒有把握的.
因為你辭了工作.
首先你沒有收入.
如果你去美國讀神學.
第一有多貴.
第二來到這裡讀F1 visa.
又不能工作.
又沒有收入.
又要交租.

$^{521}$又要吃飯.
又要交學費.
雖然有少少獎學金.
但都要交學費.
那時候都有少少擔憂.
其實兩邊都有些拉力.
但我都有些積蓄.
因為我知道要讀神學都要儲錢.
但錢好像永遠都覺得不夠.
心裡面都有少少.
究竟怎樣呢.
不是怎樣.
其實都覺得是這樣.
但心裡面好像很不踏實的感覺.
都是會去的.
但都有一種很不踏實的感覺.
有一次聽到有一個姐妹的見證.
那個姐妹是神學生.
她差不多讀完.
差不多畢業.
差不多類似畢業見證.
諸如此類.
她講一個見證.
她見證說.
她初初出來讀神學的時候.
又有類似我剛才所講的憂慮.
究竟讀神學會怎樣呢.
生活會怎樣打算呢.
會不會沒錢交學費呢.
事實上我讀呼臘神學院.
很多人是沒錢交學費的.
可能Kara都知道.
但那個姐妹就說.
她說有一次.
我忘了她講靈修還是聽到.
她聽到神跟她講.
問她一個問題.
哪裡有小朋友讀書.
要想有沒有錢交學費的.
給爸爸想就行了.

$^{561}$這些是爸爸想的.
不是小朋友想的.
她這樣講完之後.
她就把憂慮交給神.
然後就很平安.
然後就讀神學.
她這樣講完之後.
對我自己都很大的提醒.
也給了我很多的安慰.
然後就沒了.
那時候就沒什麼掙扎.
你交給神.
有些事是要想的.
但不是你想.
是神想的.
好像說交學費.
讀書交學費那些.
給爸爸想就行了.
讀幼稚園你還想沒錢交學費.
那些不是你想的.
有時候我們的信仰.
很多事我們在世上活著的時候.
事實上要你想.
你想到想不完.
但是你可以交給神.
這裡就是說先求他的國.
和他的義.
不是說那些東西不重要.
吃飯當然重要.
穿衣服當然重要.
那時候的人真的吃一餐就沒一餐.
耶穌時代真的是這樣的.
衣服破了就沒衣服穿.
她的憂慮是合理的.
但是問題就是.
在合理的憂慮之先.
你可以先將神的東西放在前面.
其他的東西.
他說這些東西都要交給你.
這些東西都要交給你.

$^{601}$我們是可以追求神的國和神的義.
不是說我們活在世上那些東西.
追求的東西不合理.
但是你可以選擇的.
你可以選擇放什麼先.
當你先將神的東西放在前面的時候.
其他的東西你可以放心的.
你可以放得低.
因為那些東西根本上也不會成為.
好像你先將神的東西放在前面.
天就會掉下來那麼大的一件事.
因為很多東西.
一步一步走的時候.
就好像擠麻辣那樣.
以色列人那樣.
又不會餓死.
你說會不會每一餐都豐富到吃九大鬼呢.
肯定不是的.
但是我不能死我不死的.
但是因為神每天就供應他們.
他們就有另外一種的經歷.
說到神的國.
我們是可以追求先義.
不要被世上的思慮攔阻的時候.
有一次有人問.
究竟神的國是怎樣追求呢.
有一次在路加福音那裡.
有人就這樣問耶穌.
伐利賽人就問.
神的國何時來到呢.
耶穌回答說.
神的國來到不是眼所能見的.
人也不得說.
看啦在這裡.
看啦在那裡.
因為神的國就在你們心裡.
這裡就說到.
一個追求神的國的另外一個提醒.
就是很多時候我們想.
神的國如果來臨的時候.

$^{641}$會在環境上面.
在外在方面.
有些什麼改變.
有些什麼.
譬如當時的以色列人.
生活那裡怎樣可以改善呢.
怎樣可以得到豐盛.
怎樣他們國家的尊嚴.
可以得到恢復.
推翻羅馬帝國.
諸如此類很多東西.
不知道就是一大堆.
你所想象的豐盛美好的東西.
如果神的國來臨就會來臨.
但是主題所在這裡.
好像要校正一個重點.
那個重點是什麼呢.
就是神的國.
你不只是想外在那些發生的事情.
你知道現在又打仗.
很多如果你看那些youtube.
世界末日.
好像說到.
神的國.
天國預言.
那些很過癮.
坦白說.
因為人真的有好奇心.
好奇心看這些都比看科幻小說好.
因為真的有一點像.
科幻小說那些天馬行空.
你好像很好奇.
但是問題就是.
傑尼蘇就提到.
神的國.
如果你要去追求神的國的時候.
你的注意會不會是.
首先要看一下自己的內心呢.
是看一下自己的內心呢.
就是.

$^{681}$神的國的關注點.
當然最終極.
神國的來臨.
我們都不只是說那個屬靈化.
我們說那個內心就有神的國.
就是最終極的神的國.
當然也是新天新地.
就是一切都得到恢復.
罪惡得到清除.
魔鬼撒旦被打敗.
那些都是神的國的終極顯現.
但是問題就是現在還沒到.
現在是already but not yet.
就是現在神的國已經開始進入人間.
但是又未完全成就的時候.
那要留意的呢.
他就提到.
是留意我們的內心.
因為神的國.
來臨的時候呢.
屬神的人就成為了神國裡面的子民.
而神國的子民.
有什麼標記呢.
是不是你的額頭那裡.
畫了一個十字架.
或者是那些不屬神的人.
額頭那裡就畫了666.
是不是這樣看得見的呢.
不是這樣看得見的.
是看不見.
但是內心那裡是有神的國.
在屬神的子民的內心裡面.
換句話說.
就是說神國那個reality.
神國那個現實.
是首先進入人間.
但是進入人間的時候.
是進入神子民的心裡面去實現的.
所以神國的子民.
是第一批人.

$^{721}$來去嘗到.
神國的真實的一班人.
再說一次.
神國的子民就是第一批人.
可以嘗到神國的真實的一班人.
但是怎麼嘗呢.
因為神的國首先進入他們的生命裡面.
所以主耶穌就說.
神的國就在你們中間.
或者神的國就在你們心裡.
在你們心裡這個字也可以翻譯成.
在你們中間.
所以.
首先我們就要監察我們的內心.
從我們自己的生命那裡做起點.
永遠都不是.
信耶穌那個實際那個reality.
那個真實.
永遠都不是只是一種外在的行動.
或者一種外在的行為.
或者是一種教會活動的一種儀式.
而是人心靈裡面.
經歷到一個神聖的實在.
是彌琳.
簡單一點來說就是.
接觸到神.
經歷到神是千真萬確的真實.
而這種的真實不是道聽途說.
是他已經經歷到這個.
神國的reality.
神國的真實.
已經確實在我們的生命裡面.
這就是這裡所說的意思.
所以追求神國的人.
當然.
他心靈一定是有所見.
他一定是有一種內在的conviction.
有內在的一種的確具.
不只是一種外在.
有很多活動.

$^{761}$我有份參與就代表.
我就是教會天國的一份子.
而是人是.
真的好像在裡面.
經歷到神就在他的心裡面.
在他的中間.
而接著下一個問題.
你就會問.
那我怎樣才可以經歷到.
這種內在的轉化呢.
是不是.
我怎樣才可以經歷到.
這個內在的轉化呢.
他就說.
主耶穌有一次有人問我.
實在告訴你們.
你們若不回轉變成小孩子的樣式.
斷不能進天國.
這裡又是說進天國.
這裡就說進天國的人就是.
需要有一個轉變的.
這裡回轉的意思就是說.
有一個轉變.
這個就不是變大.
是變小.
是不是.
在我們的經歷方面.
永遠都是變大的.
尼哥底姆也說.
人已經老了.
如何能夠重生呢.
如何能夠再中武福.
如何能夠再生一次呢.
是不是.
因為所有的經歷都是變大.
不是變小的.
但是主耶穌這裡就說變小.
是不是.
這裡說怎樣進天國.
是一個變小的過程呢.

$^{801}$我相信有幾方面.
第一方面就是.
你想想小孩子有一個什麼模式.
是跟成年人不同的呢.
小孩子有一個跟成年人不同的模式.
就是他是無助的.
他是無助的.
在動物界裡.
其實由一個小的.
譬如狗仔也好.
貓仔也好.
或者老虎獅子諸如此類.
所有的動物.
他都有一個成長的歷程.
小鳥也是.
是不是.
但是一般來說.
是短的.
不想這麼長的.
最多兩年也好.
就算是一些很複雜的.
要他傳授如何去捕獵的動物.
他也是兩年.
差不多要自己賺錢.
要不然.
他生存不了.
因為那個成長期是很短的.
為什麼.
為什麼要短.
為什麼不可以長.
為什麼.
因為在他還沒有完全成長的時期.
是一個非常軟弱.
非常脆弱.
是一個非常無助的階段.
譬如你想想.
當是鷹.
在天空之王.
但是我也看過一些紀錄片.
那些鷹的媽媽不在巢裡.

$^{841}$那些烏鴉也在搞那些小的鷹.
讓他游輪也有的.
為什麼.
即使你是鷹又怎麼樣.
你是嬰兒的階段就是很無助的.
很軟弱的.
很脆弱的.
我想這個就是第一點.
你要進入天國的時候.
第一件你要學的.
你要轉換的是什麼呢.
你要進入一種無助的模式.
要進入一種無助的模式.
什麼叫進入一種無助的模式呢.
就是你要承認在神的面前.
在進入天國的歷程當中.
你是絕對是無能為力的.
靠自己.
靠自己絕對是無能為力的.
你需要承認.
時時刻刻都要依靠神.
當我們做基督徒越來越久的時間.
越來越多答案.
很多事情你已經基本上有了一套答案.
因為你聽了很多.
當你聽了一切答案的時候.
你就越來越不感覺到自己無助.
很多事情你已經自動播.
你都懂得怎麼做.
教會生活是怎樣.
怎樣和大家團友相處.
事奉.
這件事你已經很懂得了.
你無助的感覺是越來越少的.
但是呢.
這裡正正就說出了.
不是的.
到真正去.
剛才說如果天國在你心裡.
就是觸摸到屬靈的reality.

$^{881}$屬靈的實際的時候.
就和你在教會.
怎樣去operate.
怎樣去運作那種生活模式.
完全是兩碼子的事情.
怎樣去愛神.
怎樣去對付肉體.
怎樣去勝過世界.
怎樣去將自己的私慾去隔離.
怎樣去明明是很軟弱都無力的時候.
連禱告都不知道.
怎樣去有什麼話要說.
這一種所有的無助感.
是真正去接觸屬靈的reality的起步.
你有沒有想過.
或者你知不知道.
你連怎樣去.
你說我也知道要愛神.
每個人都知道要愛神.
牧師經常都說要愛神.
但是問題就是.
當我真正要愛的時候我愛不到.
我明明都知道要對付肉體.
要捨己.
人家踩我一腳.
但是我真的生氣到.
我經常都寄予他的仇.
當我都知道.
每個人都知道.
常常都要依靠的.
不住禱告.
但是說禱告多過真正禱告.
很多時候.
做牧師更多問題.
因為牧師經常都要教人禱告.
是不是比大家禱告多很多呢.
可能多你一點.
是不是多很多呢.
我都不敢說.
我都要問自己說.

$^{921}$說禱告多過禱告.
靠自己都很不行.
所有屬靈的實際.
我們進入不到.
因為我們都不承認.
靠自己是不行的.
我們仍然都死撐住.
雖然我都不是很行.
坦白說.
但是.
接著看旁邊.
都是差不多.
就是我都不行.
他都不行.
大家都不行.
一起依靠神.
就是全世界都是罪人.
我都是全世界當中的一個人.
所以我都是罪人.
那種這樣的.
那種的承認.
我是一個神開恩可憐.
我是一個罪人.
那種看著聖殿哀哭.
即是那個隨哭頹空.
那個神開恩可憐.
我是一個罪人.
那種根本是完全兩回事.
大家都知道有罪.
神讓我看到我的罪.
是多麼的可惡.
以致我願意離開.
大家都知道人都是自私的.
神幫我來到去.
將我的自我中心.
來到去學習每日的社稷.
所有這一切.
都是從無助開始的.
你是沒有任何的能力.
去靠自己走多一步.

$^{961}$再說多一層的時候.
不單是從無助開始.
甚至到我今天做牧師的時候.
我都是一樣的.
剛才所說的一切的東西.
我都一樣沒辦法靠自己來多走一步.
對付自己.
對付世界.
對付愛神.
對付私慾各樣的東西.
都一樣我都沒辦法靠自己來多走一步.
都是一樣.
好像以前眼前一絲光.
我每一刻需要主.
這樣來到去.
而這個就是做基督徒的起點.
也是繼續走下去的.
所以主耶穌就說.
你不變為小孩子的樣式.
斷不能進天國.
這個斷不能是一個很決絕的.
不用旨意上天國的意思.
斷不能的意思就是.
不用旨意.
想都不用想.
就是這個意思.
否則就斷不能.
所以一切其實都是從基礎開始的.
一切都是從基礎開始的.
這個基礎在哪裡呢?.
基礎就在羅馬書第六章那裡.
我們是浸信會.
浸信會你應該很熟悉這段經文.
羅馬書第六章那裡說什麼呢?.
怎樣怎麼說呢?.
我們可以凝在最終.
叫因天獻當.
莫斷乎不可我們在上死了的人.
豈可凝在最終活著呢?.
豈不知我們這受洗歸入基督耶穌的人士.

$^{1001}$受洗歸入他的死嗎?.
所以我們藉著洗禮歸入死.
和他一同崖狀.
原是叫我們一舉一動.
有新生的樣式.
像基督藉著父的榮耀.
和死去復活一樣.
這段經文.
羅馬書第六章沒有時間讀.
後面其實還有.
很熟悉的.
受浸的時候.
很多時候都說這段經文.
就是說一個人做基督徒的開始.
就是受浸歸入基督.
加入教會.
受浸同主同死同埋葬同復活.
這個一定是開始.
是基礎.
浸信會一定聽過這段經文.
但是.
不知道你有沒有留意到.
這裡.
如果你讀下去.
他一開始就說到.
那個是一個生命的轉化.
就是一個人.
舊有的人.
怎樣藉著洗禮.
藉著同主同死同復活.
來進入一個新生命的模式的狀況.
這段經文就是這樣說的.
如果你讀下去的時候.
去到結尾的時候.
其實中間一直都是一段說這些的.
因為你們從前雖然做罪的奴僕.
如今卻從心裡信服了.
所傳給你們的模範.
你們既從罪裡得了釋放.
就做了義的奴僕.

$^{1041}$這樣.
一路你讀整段經文的時候.
他就不斷重複一個字.
就是奴僕這個字.
如果你不信你回家.
你自己看一下這段經文.
怎樣從罪的奴僕.
轉為耶穌基督釋放了我們.
然後成為義的奴僕.
這裡說奴僕其實是很客氣的.
意思是什麼呢.
意思就是奴隸.
奴隸現在就不喜歡別人說了.
現在隨時說奴隸那些.
被人覺得是很貶義的.
又說以前什麼黑奴那些.
當人都不是人才叫奴隸.
是不是.
但是聖經很厲害.
就是用奴隸這個字.
你說奴隸這個字不太好聽.
保羅當時可能沒什麼文化感.
當時沒有那些文化.
你以為保羅不知道嗎.
保羅當然知道.
保羅那個時代.
買賣奴隸的.
他天天看到的.
他保羅絕對知道奴隸這個字.
是多麼的將人貶低的.
但是他將奴隸這個字.
特意用在他自己.
和用在基督徒的經歷上.
為什麼呢.
因為奴隸有一樣東西.
和一般人是不同的.
奴隸和一般人有什麼不同.
奴隸沒有自由.
為什麼沒有自由.
奴隸為什麼沒有自由.

$^{1081}$一般人有自由.
為什麼有自由.
因為奴隸不是自己的.
他的主權不是屬於自己的.
所以奴隸可以買賣的.
為什麼.
因為如果主人不喜歡他.
不要他.
那就賣了他.
因為他的主權也不是自己的.
不要他你不要我.
那我自己做回自己行不行.
不要他賣了他就轉給別人.
因為奴隸的主權.
是不在自己那裡.
他無論怎麼轉.
他都是一個奴隸的時候.
他的主權只是不斷的易手.
這一點大家清楚.
但是保羅為什麼說這件事.
為什麼用奴隸這個字說這件事.
其實保羅很明白這一點.
因為人就是奴隸.
我們自己以為自己就是自己的主人.
但是誰不知我們以為自己是自己的主人.
其實我們就是罪的奴隸.
我們就是魔鬼的奴隸.
根本的主權我們以為自己在自己那裡.
根本就不在自己那裡.
根本就是騙在魔鬼那裡.
騙在罪那裡.
所以我們其實就是罪的奴隸.
這個可能你都聽過了.
你都明白了.
但是接著說那一點.
可能你不知道有沒有聽過.
保羅說.
你想得到自由嗎.
不行的.
不用想.

$^{1121}$為什麼.
你以為信了耶穌.
耶穌釋放了你.
然後你就自己做自己的主人嗎.
是不是這樣.
保羅說.
不是的.
你不能做自己的主人.
你只能選擇做罪的奴隸.
如果你要轉出來.
另外一個選擇.
就是做義的奴僕.
這裡所說.
就是做神的奴隸.
沒有第三個選擇.
信耶穌就是將我們的主權.
由自己那裡.
其實就是魔鬼那裡.
或者是罪那裡.
轉到另外一個主權那裡.
這個主權就是神那裡.
你要麼轉要麼不轉.
你以為轉到自己那裡.
就得到釋放.
其實你一轉頭來.
其實也是回到魔鬼那裡.
你唯一就是要轉到神那裡.
所以你也是奴隸.
但是你現在就是義的奴隸.
就是神的奴隸.
聽明白了嗎.
這個很難聽.
做奴隸.
信耶穌的人得自由.
現在牧師跟我說.
信耶穌做奴隸.
是很難聽的.
但是聖經是這麼說的.
我們信耶穌的人.
其實基督釋放我們.

$^{1161}$是脫離罪.
他沒有說釋放我們脫離神.
沒有說釋放我們自己.
說現在自己做主.
他只是說釋放我們脫離罪.
所以我們信耶穌.
所以在整段經文.
在羅馬書這段經文說到.
受浸同主同死同埋葬.
怎樣脫離一種舊有的生活模式.
進入一種新的生活模式的時候.
基本上就是說.
一個信耶穌的人.
他的主權的轉換.
如果你的生命的主權不轉換的時候.
你的生命的形態是不會改變的.
你的主權是需要交給神.
而交給神這個主權是一個很conscious.
是一個你要自己做的.
不會自動的.
我們交給罪的主權.
就是自動的.
交給神的主權.
是要你自己願意去給神.
我講個比喻也是很難聽的.
就是你知道了.
鬼上身是怎樣的.
可能你沒見過.
鬼上身就是不受自主的.
魔鬼控制人.
那個人是沒得選擇.
做什麼怎樣的.
就是好像被他控制了.
去魔鬼.
但是被聖靈充滿是怎樣的.
就是你自主的.
不要醉酒,酒能使人荒唐,乃要被聖靈充滿.
是你自主的.
同樣都是一樣.
給神的主權是要你自己給.

$^{1201}$神不會像鬼一樣要搶你的主權.
然後佔你的身體.
沒有的,永遠沒有這樣的一回事.
不會這樣發生,不會自動發生的.
你的生命不獻給神.
你一天舊有的想法,舊有的價值觀.
舊有的思維,舊有的模式.
是不會走的.
主耶穌講的比喻.
一隻鬼趕走了,七隻走了來.
是不會走的.
我們的主權是需要交給神.
我不知道你信了耶穌這麼久.
你的生命的主權有沒有交給神.
有時候你可能做了基督徒很久.
你那個屬靈的經歷.
好像去到一個地步已經去到瓶頸了.
已經好像不能再進前了.
因為好像各樣東西已經有答案了.
各樣東西已經自動波了.
教會生活也是自動波了.
然後屬靈的經歷有時候才進有時候才退了.
但是都是差不多.
都突破不了.
因為我們的主權是需要不斷的交給神.
從小孩子開始做起.
小孩子對於他信任.
就真的是完全的信任.
小孩子不會說一半.
一表面就說信你.
然後背後就猜度你.
或者背後就留一把.
他不會這樣.
小孩子的信任就真的是信任.
而很多時候.
人生.
我們將生命的主權交給主.
為什麼為之追求神.
先求神的國和神的義呢.
一開始那裡說.

$^{1241}$其實就是將神的事先擺.
就是將神做我的主.
我的主權交給他.
很多時候.
有很多基督徒信了耶穌很久.
他們的生命都沒有將主權交給神.
不代表你不是做基督徒的.
你都是做基督徒的.
但是很多事情你很掙扎.
因為你大大小小的事情都成為了掙扎.
我舉個例子.
譬如以前一些老一輩的爸爸.
就是我們中國人傳統的爸爸.
他們不會想.
今天要不要上班.
今天要不要賺錢養家.
今天很累了.
不如今天放一天假放一天工.
不會的.
你明白我說什麼嗎.
因為他們.
我講老一輩那些父親的形象.
因為他們已經有一個決定.
他成家之後.
他就要負責養家.
這個決定已經成為了他們人生最底層的決定.
擺在前面.
所以他不會一天上班.
一天就想今天怎樣蛇王.
怎樣偷懶.
沒有的.
但是相反來說.
現在可能一個大學生畢業.
找不找工作.
今天找了一份工作.
上了兩天班.
今天很累.
不如不要上班.
又看了一個電視節目旅遊.
這麼好.

$^{1281}$不如請假.
請不到假.
辭了工作去旅遊.
為什麼.
因為他根本沒有決定過他人生上班這件事.
不是一個很重要的.
他也沒有什麼負擔.
這個比喻可能不是這麼好.
但是你明白我說什麼嗎.
因為有一個基層的決定.
基本上成為了一些很小的決定的指導.
就好像一個運動員.
他決定要成為一個運動員的時候.
他怎樣吃東西.
每天生活的規律.
他不需要每天掙扎.
因為他已經知道自己是一個運動員的身份.
他有這樣的決定.
說回來的時候.
做基督徒.
生命主權交給主就是這樣的決定.
你沒有這樣的決定.
今天要不要早點起床.
還是要靈修.
很累要多睡一會.
沒有這樣的決定.
今天很晚睡了.
祈禱還是明天吧.
沒有這樣的決定.
又要團契.
不過看了劇.
追了兩集還想看完.
所有的每一樣東西.
我不是說你一定會選神或者不選神.
而是說你每一樣東西.
都要很費力的去打這些仗.
有時候要選贏.
有時候要選輸.
那時候就不斷在瓶頸裡來來回回.
但是如果你的生命交給神的時候.

$^{1321}$你沒有這些很大的掙扎的.
坦白說.
這個是我自己的經歷.
我想很多基督徒都有這樣的經歷.
生命主權交給神之後.
今天還要不要靈修.
還要不要問.
就等於你今天星期天要不要回教會.
還要不要問.
因為你根本沒有這些掙扎.
你已經在最底層.
你要做一個怎樣的人.
你的生命主權已經決定了.
因為你已經是神的奴隸.
合理嗎.
絕對是.
聖經是這樣教導.
我們是這樣來做基督徒的.
不是你是老闆.
然後每一樣東西再算數.
這件事要不要做.
你每一樣東西都這樣算數.
你做人都做到很累.
因為你跟從神.
所以就變成一個很大的負累.
變成一個很大的重擔.
變成每一樣東西都要算過量過.
或者每一樣東西都要你自己去掙扎完.
然後這次終於打贏了.
今天我這個星期持續零修了七天.
求神幫助我們.
我們要先求神的國和神的義.
在結束的時候.
保羅在哥林多後書有一句這樣的說話.
他說你們總要自己審查.
有信心沒有也要自己試驗.
豈不知你們若不是可棄絕的.
就有基督.
耶穌基督在你們心裡嗎.
哥林多教會面對著很多的問題.

$^{1361}$哥林多前書.
但從來沒有一個問題是關於.
耶穌基督是不是神的兒子.
他們信神而耶穌基督是神的兒子.
絕對沒有問題的.
他們是不是基督徒.
他們絕對覺得自己是基督徒.
哥林多前書絕對沒有這樣的問題.
但是有各樣的問題.
有紛爭啊.
有嫉妒啊.
然後有情慾啊.
各樣都有啊.
去到最後的時候.
保羅我不知道是不是受不了.
但是可能最後他也要問他.
喂 你想一下自己啊.
你審查自己有信心沒有啊.
這裡所說的信心.
我相信不是說信不信耶穌.
是不是神的兒子.
信不信自己是基督徒.
那些信心.
而是最底層的信心.
對神那種信靠的信心.
就是剛才所說主權交出的信心.
有沒有啊.
你是不是屬於主啊.
就是最後哥林多後書結束的時候.
保羅這樣問的時候.
我相信我們都需要問同樣的問題.
\newpage



\section{民數記 20:12}
\label{sec:032087}
\textbf{聖經人聖經事:摩西發脾氣出事}~~~~黃克勤牧師博士
\newline
\newline
連結: \href{http://www.church.com.hk/acms/content.asp?site=cdc\&op=show\&type=product\&code=032087\&layout=sermon}{\texttt{032087}} ~~~~ 語音日期: 2026-05-02
\newline
\newline
\hyperref[sec:032086]{\small{< < < PREV SERMON < < <}}
~
\hyperref[sec:index_chronic]{\small{[返順時目]}}
~
\hyperref[sec:032088]{\small{> > > NEXT SERMON > > >}}
\newline
\newline
民數記 20:12
\newline
\begin{longtable}{cl}
\hline
\hline
章節 & 經文 (和合本修訂版)\\
\hline
20:12 & \begin{tabularx}{0.7\textwidth}{X} 但是耶和華對摩西、亞倫說：「因為你們不信我，沒有在以色列人眼前尊我為聖，所以你們必不能領這會眾進入我所要賜給他們的地去。」 \end{tabularx} \\ \\
[1ex]
\hline
\hline
\end{longtable}
$^{1}$各位電視節目的觀眾朋友大家好.
又跟大家一起讀聖經了.
《生命傳書》的網絡.
鼓勵大家多看聖經.
了解裡面的人物.
裡面的故事.
而看到背後的上帝.
是有什麼心意的.
今天來到這件事就令人很遺憾.
唉,真是低頭了.
覺得上帝好像很….
唉呀,用不著這麼緊張.
用不著這麼尖銳.
這樣很Harsh.
摩西又是一個很慘的經歷.
這個當然是有我們的投射在摩西那裡.
摩西可能覺得沒有我們這麼慘的.
好.
摩西發脾氣出事.
這個不是聖經的字.
看就知道是我的字.
這個句子.
摩西發脾氣出事了.
這次看看要做什麼.
今天聖經很長的.
民數記二十章第一節.
「正月間,以色列傳會中,到了沉的曠野」.
這個「沉」呢.
Jin.
Jin好像普通話的「沉」.
所以就譯作「沉」.
Jin.
是平面的意思.
就是說這個曠野比較平.
去到一個沉的曠野.
他們就紮營停了在迦底斯.
這個Kadesh.
Kadesh就是神聖的意思.
這個Kadesh很多時候都會在其他經文出現的.
這個摩西的姐姐.

$^{41}$這個重要的人物.
不過一度都曾經攻擊過摩西的.
都是有權力的問題.
一定是你說了算的.
我堂堂姐姐也要說一下.
這些不要緊了.
人之常情.
米利庵就是Mia.
所以Mia就是米利庵.
其實這個名字很有趣.
反抗的意思.
以前也說過了.
很有可能.
由於當時他們在埃及地被埃及人迫害.
要殺殺男孩.
要做奴隸.
苦工.
所以很有可能這個女孩子出生的時候.
爸爸改了名字就是要反抗.
於是就把自己的心意.
給了兒女的名字.
米利庵就死在迦底斯的抗爺那裡.
撞在那裡.
匯縱去到迦底斯沒有水殼.
你知道了.
以色列人一有不妥當.
又來聚集攻擊摩西亞倫.
攻擊.
怎麼攻擊呢?.
下面看到了.
那些說話難聽到.
肉酸到.
如果我是摩西人.
一巴掌打過去.
真的很難說.
真的很難聽.
百姓在摩西爭吵.
說什麼呢?.
我們的弟兄曾死在耶和華面前.
我們恨不得與他們同死.

$^{81}$這句話怎麼說的?.
我們的弟兄死在耶和華面前.
我們恨不得與他們同死.
意思是什麼?.
就是耶和華害死我們的弟兄.
我們真的很慘.
我也想和他一起死.
不如害死我吧.
如果我的翻譯是廣東話的口語.
耶和華害死我們的弟兄.
我也不如害死我吧.
我也想和他一起死.
說這種話沒有水喝.
大家都沒有.
摩西亞倫你們為何把耶和華的匯眾領到這曠野.
使我們和牲畜都死在這裡.
為什麼你們要帶我們來這裡.
是不是想渴死我們和我們的牲畜.
你們為何逼著我們出埃及.
領我們到這壞地方.
逼著我們出埃及.
這樣也可以說的.
你不要來.
現在你回去吧.
我沒有逼著你.
是大家一起走.
你還拿著埃及人的金,銀器,衣服,牲口一起走.
你自願一腳走出來.
我怎會逼著你出埃及呢.
領我們來這個壞地方.
這裡壞到怎樣.
這個地方不好殺種.
因此他們可能在這裡停留了一段時間.
這裡又不能種田.
又沒有無花果樹.
又沒有葡萄樹.
又沒有石榴樹.
又沒有水飲.
你以為這裡是油果園嗎.
這裡都沒有無花果.

$^{121}$這裡是曠野.
沒有無花果樹.
沒有葡萄樹.
沒有石頭樹.
又沒有水井.
那怎樣生活呢.
這個鬼地方.
為什麼你逼我們出埃及.
為什麼你們兩兄弟帶我們來這裡死.
第一句是最壞的.
第一句是我們的弟兄曾死在耶華面前.
我恨不得與他們同死.
這句更壞.
耶華說害死我們的兄弟.
現在我都想與他們同死.
害死我吧.
如果我是上帝.
就一個雷劈下來.
摩西亞倫呢.
我想這只是一個特別的例子.
我想平日這群人的怨言已經多到不得了.
都有些聽到利利地.
頓頓地.
沒有費事反應.
摩西亞倫留意.
第一點要學的.
忘記了寫號碼.
螢光字第一點.
遇到這些問題.
都會莫門口.
俯伏在地.
哇,真是好.
摩西亞倫呢.
他真的拿到一個竅門.
一個秘訣.
一個武功心法是什麼呢.
別人攻擊到這樣.
如果我呢.
我就會立刻站好.
去反駁他.

$^{161}$是你自己跟出來的.
你這麼厲害.
現在回去.
明知道來抗議的.
是走一段路的.
怎麼會有無花果樹種好在等你.
怎麼會有石榴樹.
怎麼會有葡萄在種好.
等你讓走.
你不是吧.
你沒事吧.
今天吃了藥沒有.
那你就.
一定會大家頂過,生氣的.
你還敢說耶和華帶你.
殺死你的兄弟.
不是吧.
哇,我想如果我是摩西.
真的頂不住.
我想那支匠.
就飛過去插死幾個.
真的很生氣.
但是呢.
摩西亞倫很沉靜.
離開匯種.
留意一下.
離開匯種.
離開.
當去到這些火爆場面.
很厲害的時候.
其中一個.
今天上課也教一下.
就是.
怎樣去做呢.
就是如果只有你離開這個場面.
不要站在那裡.
這樣越來越大聲.
是有機會就是.
靜一靜.
當然有人幫你做梅就好一點.

$^{201}$哎呀,摩西亞.
不如你過去那邊休息一下.
冷靜一下.
你們又等一等.
我們一會再談.
現在先停一停.
暫停.
去個洗手間先.
離開那裡.
斷了它切斷了.
當然我們今天教人家.
去洗手間.
洗個臉那些.
其實是切斷了思路.
不要再繼續爭辯.
但是摩西亞倫厲害的就是.
離開匯種.
是去到匯莫門口.
俯伏在地求問上帝.
這一點大家要學.
我也要學.
你也要學.
就是去到這些最難的位置.
心情最壞的位置.
情緒真的氣到震的位置.
憤怒的位置.
那怎麼辦?.
去到神面前.
無論你是找廁所其中一格.
關上門.
自己靜一靜.
在那裡禱告.
去回自己的房間.
走一走出去也好.
看看窗口.
你要心靈來到屬靈匯莫的門口.
心俯伏在地上.
神啊.
怎麼辦?.
你教我吧.

$^{241}$這是第一點要學的.
於是耶和華的榮光向他們顯現.
於是神就曉喻摩西了.
你拿著像.
留意一下.
你拿著你的像.
那支像是象精神的同在與能力.
你拿著像.
和你的哥哥亞倫.
焦醉匯種在他們眼前.
分腹留意一下.
第二個榮光的地方.
第二點就是.
分腹盤石發出水來.
記得紅色字.
分腹盤石發出.
神就會叫你做一些事情.
你拿著一支像.
其實是表示你重拾你的信心.
重拾神給你的能力.
重拾以前這支像走過很多神蹟.
你是有信心的.
從丟到地上變成蛇.
撿起來變成一支像.
由一指向海水.
就分開百姓過海.
好似行乾地.
但都是這支像.
你拿著這支像.
就是信心的一個象徵.
重拾信心.
然後你分腹.
分腹是不用做事的.
不用動的.
是叫盤石出水就可以了.
於是就分腹盤石發出水來.
水就從盤石流出來.
吸會眾和他們的生促汗.
這樣就分腹盤石出水.
這樣就有水流了.

$^{281}$於是摩西照耶和華所吩咐的.
就是從耶和華面前取了像出去.
拿著一支像.
摩西和阿倫就招聚會眾到盤石前.
我想留意大家留意一下這個字.
盤石.
盤石我也和我的團友去過很多次.
以色列去到哥薩利亞菲勒比.
有很多神像的小小的小洞.
以前在那裡拜很多樣神.
那裡除了我講解之外.
我就叫大家找一下盤石在哪裡.
指給我看.
大家就低頭周圍看.
這塊石這麼大是不是盤石.
那塊石十幾二十呎真的挺大的.
那塊也不是盤石.
只是一塊石.
大家都找不到.
然後我叫大家抬起頭看看.
嘩!原來那個盤石是一座山.
一座山.
在山腳下我們幾個人仔在這裡.
嘩!那個盤石.
這個才叫盤石.
不是一塊的.
是一座山.
所以你會明白什麼叫做.
耶和華是我的高台.
是我的盤石.
是我的拯救.
因為如果你站在上面.
我要攻打你.
是沒有可能的.
你扔石頭下來也扔死我.
這樣就很高台.
是很穩陣的.
好了.
去到一塊盤石面前.
摩西就說.

$^{321}$你們這些.
你看這個摩西的說話.
他雖然沒怎麼出聲.
不過這個說話顯示他有點脾氣.
心裡有點生氣.
你們這些背叛的人.
聽我說.
我為你們使水.
從這盤石中流出來嗎?.
這裡又變成一個問號.
變成一個問句.
你們這班五逆仔.
聽著.
看看我可不可以在這塊石上.
為你們變些水流出來.
如果我做今天廣東話就是.
你們這班五逆仔聽著.
聽著我說.
我為你們.
看看行不行.
我要在這塊石上.
流些水出來.
五逆仔聽著.
摩西就舉手.
第三點.
螢光字.
用將擊打盤石兩下.
嘩!很多水流出來.
砰砰聲.
所以那些水.
你千萬不要以一塊石.
這麼小塊.
然後一打.
有一條水線流出來.
二百萬人喝.
嘩!不是吧.
要喝到什麼時候.
還要拿兩桶回去.
其實盤石是指這樣.
嘩!一座山.

$^{361}$然後摩西打了兩下.
整塊大盤石很多地方.
砰砰聲水流出來.
嘩!這麼厲害.
於是他們可以.
二百萬人連生畜.
都可以有水喝.
所以不是一條水線.
打下去水流出來.
不是那些.
哪裡夠.
喝到明年都未喝完.
未喝完又喝過了.
是那一塊巨大的山岩.
砰!很多水流出來.
就是這樣.
於是很多水.
生畜都喝了.
問題在哪裡.
神叫留意看上面.
神是你拿著象.
然後去第八節說.
你是吩咐盤石出水.
紅了兩個字.
是吩咐神給他的答案就是.
你行的.
你拿著象.
拿著你從實你的信心.
不要怕.
不要怕那班人.
你拿著你的象.
法老王你都不怕.
拿著這些不好的心亂.
吩咐盤石發出水.
你去吩咐就行.
誰知摩西舉手用象擊.
打盤石兩下.
是有水.
不過耶和華對摩西說.
因為你們不信我.

$^{401}$不在以色列人眼前尊我為聖.
所以你們必不得領者會眾.
盡我所賜給他們的地去.
嘩!這個要緊.
你們神是說.
我覺得這裡很嚴厲.
你們兩個不信我.
其實他們不信.
又怎會去那塊石擊打呢?.
其實摩西是做多了.
就是擊打而已.
他應該吩咐.
不過他擊打了.
神說你不信我.
不在以色列人眼前尊我為聖.
其實以色列人有沒有聽到.
神跟他們說什麼呢?.
沒有聽到.
他自己在會幕前禱告.
神教他.
不過神說.
我都當你在以色列人面前不尊敬我.
所以你不可以領他們進入.
我賜給他們的地方.
即是迦南.
兩件事.
打了兩下就沒有了.
沒有得去.
這水名叫米利巴.
米利巴.
是爭拗,紛爭的意思.
所以這裡的水.
自此就叫做紛爭水.
爭拗水.
因為以色列人向耶和華爭鬧.
嘩!你們夠膽.
向耶和華爭鬧.
好了.
你看到嗎?.
有沒有看得出這些經文的字.

$^{441}$是告訴我們.
神知不知道發生什麼事?.
知道的.
是百姓向耶和華爭鬧.
不過剛剛摩西亞倫站在中間.
一夾.
夾中.
就說你帶我們出來死嗎?.
其實他們是先鬧耶和華.
那個耶和華害死我們的兄弟.
是不是想害死我?.
不如我也死了.
嘩!如果我是耶和華.
真的一下子弄死他.
但神很忍住.
不過神反而是透過摩西亞倫去做一件神蹟.
於是就這樣發生了.
那些人沒有理由.
這裡也不好.
這塊地也沒有無花果樹.
沒有葡萄樹.
沒有石榴樹.
唉!有果樹等你二百萬人嗎?.
不會抗議.
好了.
因為他們向爭鬧.
神是知道是他們那幫人向耶和華爭鬧.
不過神在這裡就將他覺得不尊敬.
給了摩西亞倫身上.
好了.
看看另一段相關經文.
《新命紀》三十二章.
「當日耶和華吩咐摩西說:.
『你想借阿巴林山中的尼泊山.
在摩亞地與耶利哥相對.
這裡是摩西就快死了.
神就說你不可以進入迦南.
現在我帶你去哪裡呢?.
去尼泊山.
尼泊山剛剛望到迦南地.

$^{481}$隔著約旦河望到迦南地.
望到迦南地是甚麼?.
望到迦南地的耶利哥.
其實就是他們最先過了約旦河.
進入迦南就是進入耶利哥.
大家記得.
你看那裡就是耶利哥.
是相對的.
那裡就是迦南地.
你去那裡看看吧.
如果你是摩西.
你會有甚麼心情呢?.
然後神就說你去那裡看看吧.
有第五十節.
這句真的要命.
跟他講.
「你必死在你所登的山上.
歸你列祖去」.
好像哥哥亞倫死在荷爾山.
荷爾山即是荷列山即是西乃山.
死在那裡.
歸他列祖一樣.
其實只有這裡.
沒有記載亞倫死.
原來亞倫死在西乃山.
一個頒佈律法的地方.
可能亞倫是大祭司.
所以不在那裡死.
其實這個都是在說.
你死在尼泊爾山.
你哥哥死在荷爾山.
其實都是因為這件事.
就是我吩咐你.
吩咐盤石出水.
誰知你擊打兩下.
我知道你是發火的.
所以我很不開心.
你們兩個都不能入迦納.
哥哥死在荷爾山.
你死在尼泊爾山.

$^{521}$看看吧.
亞倫看都看不到.
歸他列祖.
因為你在沉的曠野.
迦底斯的米尼巴水.
憎鬧之水.
在以色列人中.
沒有尊我為聖.
得罪了我.
這個得罪是mild.
mild是行為不忠.
什麼叫行為不忠呢.
不是說那些貞潔等等問題.
行為不忠是指.
我叫你擔張椅子.
你又不擔.
我叫你擔兩張椅子來坐.
你又擔一張.
給我兩個朋友倒一杯水.
行為不忠.
就是叫你做的事你不做.
這麼簡單.
我叫你做的事你不做.
你沒有尊敬我.
所以我賜給以色列的地.
卻不得進去.
摩西從摩亞平原登尼泊山.
上了拉越耶利哥相對的比斯加山頂.
這是尼泊山的一個尖峰.
耶和華對他說.
這就是你看看對面耶利哥那塊地.
這就是我向亞伯拉向以色亞國.
起誓應許之地.
我必將這地賜給你的後裔.
現在我使你的眼睛看見了.
你卻不得過到那裡去.
於是耶和華的僕人摩西.
死在摩亞地.
摩亞地就是約旦河的西邊.
就是以色列地.

$^{561}$即是說是應許之地加勒.
他的東邊.
約旦河的東邊那邊的岸.
就是今天叫約旦.
以前就是摩亞阿曼.
你知道摩亞阿曼是羅德兩個女兒.
那些地方.
你會知道今天約旦首都阿曼.
和阿曼這個名字有點相似.
低一點藍一點就是摩亞.
北一點就是阿曼.
這個死在摩亞地.
正如耶和華所說的.
第六節.
這節寶貴了.
耶和華將他埋葬在摩亞地.
伯比爾對面的谷中.
只是到今天沒有人知道他的墳墓.
這個伯比爾是甚麼呢.
我寫了一個希伯來文出來.
但譯音出來不知道是甚麼.
只知道Beth.
Beth就是伯.
這就是屋.
屋.
Pet O.
Pet O間屋的對面.
那個伯比爾就是Pet O的家.
對面的山谷.
只是今天沒有人知道他葬在哪裡.
是第六節寶貴.
耶和華將他埋葬在摩亞地.
雖然是外邦之地.
不過都是屬於神的.
葬在那裡.
簡約的地方.
在這裡留意.
會不會一聲嘆息呢.
摩西一生人很努力.
神一直預備他.

$^{601}$到最後只是打了兩下那塊盆石.
出水就可以了.
你吩咐就可以.
打就不可以.
打了兩下就不能進去.
應許之地.
懲罰會不會重了一點呢.
老實說我心裡真的有這樣想.
神啊你有沒有重了一點.
這個刑罰.
那就算了.
是那班人激怒他.
讓你也生氣.
我也生氣.
聽到我也走來劈死他.
最生氣的那句是什麼呢.
這個耶和華害死我的兄弟.
不如害死我吧.
這樣說話.
是不是想渴死我們.
帶我們出來.
逼我們出埃及幹什麼.
為什麼來到這裡呢.
這塊地這麼爛.
有沒有葡萄樹.
讓我們沒有酒喝.
真是飛腳打過去.
摩西學到的是什麼.
安靜祈禱.
然後第二.
神給他解決方案.
我給你一個辦法.
給你一個神職的辦法.
不過摩西仍然沉不住他的氣.
其實整個摩西的生平裡.
只有這一次他沉不住氣.
後面聖經有說.
摩西是世上最謙和的人.
極謙和的人.
沒有人比摩西更謙和.

$^{641}$他真的厲害.
不過在這裡他也時失手.
這次發脾氣.
罵那些以色列人.
另外擊打兩次盆石.
其實神是有責任的.
如果神不叫摩西拿一支杖去.
你就這樣走過去盆石那裡.
吩咐吧.
他沒有杖在手.
他不會打.
譬如你拿一支杖.
信心的標記.
拿著去.
老老實實.
拿著一支杖的人.
發火會怎樣.
扮人.
他就擊打盆石兩次.
可能摩西那一刻是不小心的.
只是覺得盆石.
動兩次.
打兩次算什麼.
但神是講得很清楚.
你吩咐不可以擊打.
沒有說不可以擊打.
不過是吩咐.
令我們學到東西.
原來聽神的命令.
要小心一字一句.
聽得清楚.
有可能他就是叫我們吩咐.
我們做這件事.
不用擊打.
有時候吩咐我們做這件事.
跟別人說就可以了.
不用動它.
不用做這麼多動作.
所以在這裡雖然覺得很不幸.
很不開心.

$^{681}$不過信服神.
他是有最美好的旨意.
不過看到這裡.
神並沒有不喜歡摩西.
一樣的歡喜他.
不過他進不來.
都讓他看看.
摩西表現得很好.
一句怨言都沒有.
有沒有求情呢.
如果我是神.
不是吧.
我這麼辛苦.
你都讓我進去站一站.
這樣算怎樣.
沒有求情.
沒有求求.
沒有摩西.
都是我做得不對.
不能進去.
不要緊.
就上我山看.
看啊.
就在這裡.
然後他就死了.
他知道.
神是告訴你會死的.
他很接受.
第六節最寶貴.
耶和華將他埋葬在摩亞地.
在哪裡不要緊.
是耶和華將他埋葬.
很厲害.
如果說無辜.
有沒有事是無辜的呢.
這個世界.
當然有.
非常之多.
如果比摩西無辜的是誰.
阿倫.

$^{721}$他哥哥是幫他.
做助手跟著他.
都死了.
是不是.
在這裡提及.
所以跟這件事有關.
這裡有些事情.
我們是不明白的.
或者覺得上帝太苛刻.
不要緊.
我們交給他.
信靠他.
這樣就最好了.
神是有恩典的.
他會為我們擦乾眼淚.
他會埋葬.
一起祈禱.
到這件事看到的時候.
真是內心無語.
不知道.
不懂得說什麼.
一個.
這麼賣力.
不生事奉你.
受盡不少委屈.
走了不少.
風塵的路.
風塵撲撲.
多辛苦.
當日要帶著200萬.
不停埋怨的.
以色列人和埃及人.
走在抗日裡.
一天都嫌多.
他就帶領他們走了這麼久.
犯了一次錯誤.
神就覺得.
你既然跟我.
跟得這麼貼.
這樣都不懂.

$^{761}$就不要讓他入迦納.
主這個我們尊重.
我們敬畏你.
知道你有你的心意.
我們信服.
讓我們遇到事情.
我們冷靜下來.
去求問之後.
我們就小心小心地.
去跟從.
每一個細節求你教我們.
又讓我們樂意去跟從.
讓我們.
拿起那支杖.
信心的標記.
我們是更加緊的.
跟隨你.
縱然可能我們有時候做錯了.
有些事我們有接受過後果.
不過我們都願意信服.
願意順從.
讓我們像摩西一樣.
這麼謙和地跟隨你.
隨聽祈禱奉耶穌基督命求.
阿們.
\newpage



\section{以弗所書 5:15-20}
\label{sec:032088}
\textbf{要明白主的旨意:七要三不要}~~~~杜崇基牧師
\newline
\newline
連結: \href{http://www.church.com.hk/acms/content.asp?site=cdc\&op=show\&type=product\&code=032088\&layout=sermon}{\texttt{032088}} ~~~~ 語音日期: 2026-05-03
\newline
\newline
\hyperref[sec:032087]{\small{< < < PREV SERMON < < <}}
~
\hyperref[sec:index_chronic]{\small{[返順時目]}}
~
\hyperref[sec:032089]{\small{> > > NEXT SERMON > > >}}
\newline
\newline
以弗所書 5:15-20
\newline
\begin{longtable}{cl}
\hline
\hline
章節 & 經文 (和合本修訂版)\\
\hline
5:15 & \begin{tabularx}{0.7\textwidth}{X} 你們要謹慎行事，不要像無知的人，要像智慧的人。 \end{tabularx} \\ \\ \relax
5:16 & \begin{tabularx}{0.7\textwidth}{X} 要把握時機，因為現今的世代邪惡。 \end{tabularx} \\ \\ \relax
5:17 & \begin{tabularx}{0.7\textwidth}{X} 不要作糊塗人，要明白主的旨意如何。 \end{tabularx} \\ \\ \relax
5:18 & \begin{tabularx}{0.7\textwidth}{X} 不要醉酒，酒能使人放蕩；要被聖靈充滿。 \end{tabularx} \\ \\ \relax
5:19 & \begin{tabularx}{0.7\textwidth}{X} 要用詩篇、讚美詩、靈歌彼此對說，口唱心和地讚美主。 \end{tabularx} \\ \\ \relax
5:20 & \begin{tabularx}{0.7\textwidth}{X} 凡事要奉我們主耶穌基督的名常常感謝父神。 \end{tabularx} \\ \\
[1ex]
\hline
\hline
\end{longtable}
$^{1}$可能今天大家覺得這個題目比較有趣.
要明白主的旨意.
基本上大家都很熟悉.
七要三不要.
為什麼會有七要三不要呢.
一會兒聽我繼續分享的時候.
大家就會明白的了.
其實今天我們身處一個很複雜.
多變和急促.
而且令人容易迷失的時代.
香港人對這種感受是特別深的.
社會環境在變.
經濟環境在變.
家庭處境也在變.
教會的生態也在變.
人面對很多不確定.
前路應該怎樣走呢.
教會應該怎樣走呢.
家庭應該怎樣走呢.
我個人的信仰也要怎樣活下去呢.
今天很多人不是完全沒有方向.
而是方向太多.
不是完全沒有聲音.
而是聲音太多.
不是完全沒有選擇.
而是選擇太多.
但在選擇越來越多的時代.
人卻未必越來越清楚.
相反我們可能會越來越焦慮.
越來越疲憊.
越來越迷茫.
所以今天這段經文非常切合我們.
保羅不是向一群.
活在安穩年代的人說話.
而是向一群身處異教文化.
價值混亂.
靈性爭戰中的門徒說話.
他提醒當時的人.
也提醒我們.
你們不要活得糊塗.

$^{41}$不要活得無知.
不要被世界拖著走.
你們要作智慧人.
要明白主的旨意.
今天我會用一個很清楚的結構.
來傳講這段經文.
這個七要三不要.
總共有十個提醒.
這個不是零碎的道德命令.
而是一幅完整的屬靈生命的圖畫.
教導我們個人如何活.
也教導教會群體.
如何尋求神的方向.
而今天我想特別強調的一點.
就是講到明白主的旨意.
不單止是個人層面的問題.
也是關於教會群體辨識的問題.
我們先看這個圖.
我們看到藍色字.
有三個不要.
第一個不要.
就是不要仗無知的人.
相對就是要作智慧的人.
第二個不要.
就是不要作糊塗人.
相對來說.
就是要明白主的旨意如何.
第三個不要.
就是不要醉酒.
相對就是要被聖靈充滿.
這個就是三不要.
那麼七個要是什麼呢.
七個要在紅色字那裡.
要謹慎行事.
要仗智慧的人.
要把握時機.
要明白主的旨意.
要被聖靈充滿.
要口唱深和地讚美主.
大家很清楚了.

$^{81}$題目是七要三不要.
大家就知道什麼是不要.
什麼是要.
我們首先去看一下.
什麼叫三不要.
就是要離開錯誤的生命狀態.
經文說我們不要仗無知的人.
無知的人不是說他的智力不足.
而是缺乏屬靈智慧和分辨.
經文提醒我們.
不可以只懂得世界的方法.
卻不懂得神的道路.
在我個人來說.
我們作決定的時候.
是不是按真理.
或者是按我們的情緒呢.
我們是不是隻是得過忘字.
而不清醒呢.
作為教會教會是否只是重策劃.
輕禱告呢.
其實無論是策劃或者禱告.
兩者都是重要的.
但是我們要分清優次.
就是在教會行事.
我們要先著重禱告.
然後之後我們就去策劃.
我們是否只是學外面的模式.
卻少民主的心意呢.
傳道人其中一個最有趣的話題.
我們經常都講的.
就是究竟一間成功教會的策略.
應該是怎樣的呢.
有幾個傳道人就在一起聊天.
大家就互相問.
杜牧師.
你間教會有多少人呀.
我說某某有多少人呀.
接著.
麥牧師.
你間教會有多少人呀.

$^{121}$哇一千人呀.
哇一千人好厲害哦.
接著呢.
接著問第三個牧師.
那個牧師小小聲說.
我教會呀.
一百人而已.
我們經常都會有個意見就是.
一千人的教會當然是成功啦.
那你.
究竟你教會有什麼事工的呀.
你有什麼策略的呀.
為什麼你可以有一千人那麼多呀.
那我們這些教會呢.
教務呢經常很有興趣.
我記得有一年.
我就參加了華里克牧師的講座.
當時整個城鎮都坐滿了教務同工.
聽華里克牧師講.
那個直奔標杆.
就是可能大家早年都聽過.
哇就是怎麼樣呢.
是以慕道者為主的一個教會啦.
那他其中一個例子呢.
就是整個人像出來的.
他就說馬鞍風的老馬.
就是說他吸引呢.
在教會里面呢.
專吸引一些呢.
中上教育的人士啦.
可能是財經界的人士啦.
可能是專業的人士.
做電腦的人士啦.
專吸引這幫人呢.
來回教會的.
所以他的教會呢就增長得很大啦.
還有呢.
財務呢也是很豐盛的.
大家聽到呢津津有味.
哇那我們要學一下.

$^{161}$究竟在香港.
哪一個呢.
是我們的馬鞍風老馬呢.
那其實.
究竟.
我們這樣做是不是真的有智慧呢.
我們這樣做是不是真的合乎神的心意呢.
在《不在一樣》這本書裡面提及到.
作計劃是很有價值的操練.
但永遠不能取代聆聽神的話語.
我們與神的關係遠比人所能作的任何計劃重要.
其實最好的計劃會議.
就是禱告.
我們願意花時間到天父的面前.
發掘他想要我們做什麼.
倘若教會真是很認真要進行神的旨意.
就會花更多的時間.
來尋求神的旨意.
更少爭論或者辯論.
每個會友對於教會應該怎樣做的意見.
第二個不要呢.
就是我們要不要作糊塗人.
糊塗的意思就是指.
不願意分辨.
不願意理解神旨意的人.
這裡的糊塗與第十五節的無知.
不是完全同一個字.
在第十七節的糊塗.
更有愚昧.
遲鈍.
行事愚蠢的意思.
特別是指生活行為上的不明智.
與這種糊塗相反的.
就是明白主的旨意.
很多時候不是神沒有說.
而是人不願意信服那些已知的真理.
我問下自己.
我是不是真的尋求信服.
還是隻是想神支持我的決定呢.
我有沒有在已知的真理上忠心呢.

$^{201}$作為教會.
教會尋求方向是靠神的引導.
還是靠人的規劃呢.
是否我們只是追求成功.
而不追求忠心.
讓我們檢討我們開會的心態.
在我年輕的時候.
開會經常有個意思.
就是我經常想我有一個.
很好的想法.
很好的意見.
我希望在開會里面.
清楚地表達.
以至開會的人都聽我的意見.
聽我講.
但是隨著年日我發覺.
原來開會不是這樣的.
開會就是讓我們.
聆聽不同弟兄姊妹的意見.
以至我們集合一個最好的共識.
我們能夠在會議裡面.
我們仍然能夠開懷.
我們仍然能夠很慫容地.
面對不同的意見.
這個是我們要學習的.
以往我是要爭勝.
但是現在我卻是尋求一個和睦的聲音.
我們一起在神裡面尋求主的共識.
教會要有群體的辨識.
不是隻是靠少數人的督定.
也不是模仿成功的模式.
我們是要借神的話.
藉著禱告.
藉著聆聽.
藉著等候.
我們一起尋求神.
在過去5月1號.
過去的星期二.
星期三.
我們有教務的檢視營.

$^{241}$其實這個檢視營.
我們有一個心願.
就是我們由去年9月開始.
將教會今年的主題鋪排的時候.
差不多如果由去年9月到今年4月.
已經有七個月的時間了.
究竟我們在這七個月裡面.
我們是否能夠做到.
我們很想帶出給弟兄姊妹所知道的東西呢.
大家知道我們今年的主題.
上卷就是「渴慕靈火」.
究竟在今年裡面.
我們做了些什麼.
讓弟兄姊妹渴慕這個靈火呢.
我就數一數.
在今年裡面我們辦了幾個聖靈講座.
在聖靈講座裡面.
我們去聽到什麼叫做「被聖靈充滿」.
我們去處理一些外顯的恩賜.
我們也有機會.
有一群弟兄姊妹.
去了臺灣的烈火特會.
去領受神給他們的訊息.
我們也都願意在教會里面.
我們去多教導多分享多認識.
當我們今天看到我們的星期三晚的月土會.
我們越來越多弟兄姊妹都能夠領受感動來參與的時候.
我們知道聖靈在教會當中來動功.
但我們不單止去回顧.
我們更加去前瞻.
我們下半年我們的主題是什麼.
「點燃生命」.
我們如何去點燃弟兄姊妹的生命呢.
我們就發現我們需要重拾一個很基本的真理.
就是教會需要有跨代的同行.
教會需要有門徒的訓練.
可能你覺得這些都是很老生常談.
沒有什麼新意.
其實你想想.
教會的屬靈原則.

$^{281}$二千年以來都是這樣的.
我們只是因為時代不同的轉變.
我們在形式上做出調節.
但是那個核心思想仍然是沒有變的.
在5月1號我們三股合成醒止.
不是一個事工的會議.
而是我們一起祈禱.
我們一起彼此服侍.
我們一起安靜.
我們一起尋求主.
我們一起去思想到.
在神心意的教會應該是怎樣.
而在每一點裡面我們再發揮.
我們再分享.
其實大家知不知道每次做完這些聚會.
我們教會同工繼續要忙.
因為我們要整理.
我們要統合弟兄姊妹的體見.
我發現整個的聚會所著重的.
反而不是我們怎樣做事工.
怎樣去牧養.
反而是我們作為領袖.
生命是需要更新成長.
弟兄姊妹.
這個就是我們不要作糊塗人.
我們要明白主的旨意.
第三我們不要醉酒.
當然醉酒在希羅時代已經是一個問題.
但是到今天醉酒仍然是一個問題.
迂酒的問題在社會里仍然很嚴重.
迂酒令到一個人失去了常性.
很多家庭暴力的問題都是從迂酒引起.
而這裡所說的我們放蕩.
放蕩的意思就是說我們失控.
我們耗散.
我們敗壞.
其實如果我們引申不要醉酒這個原則的時候.
就不只是說喝酒的問題.
而是我們不要被其他的力量來控制我們.
人總會被某些東西影響和掌管.

$^{321}$問題是我們在什麼影響之下生活呢.
被酒精影響會帶來放蕩.
被情慾影響會帶來敗壞.
被恐懼影響會帶來退縮.
被世界影響會帶來妥協.
被性靈充滿才能帶來聖潔喜樂智慧和感恩.
我就問一下自己.
我現在正被什麼牽引麻醉和支配呢.
我發覺有兩個機器在影響我們.
第一個就是滑手機.
在我們今天的生活裡面.
我們自然都用很多時間滑手機.
我最近發現.
中國裡面有很多手機短片.
不是很長的.
但是它一滑的時候你就覺得挺好看的.
看著看著就看了一些時間去了.
這個其中一個.
還有滑手機還有什麼好處呢.
就是可以放鬆心情.
經常看看有什麼網上的產品.
有什麼買.
現在我發覺去淘寶買很多東西都很划算.
買喫的東西又划算.
買那些器材又很划算.
想著想著滑著滑著真的滑了一些時間.
不過通常不是我下單的.
通常在家裡說買什麼.
現在我穿著西褲都可以淘寶.
現在真的很便宜.
八十多塊錢.
但是你出去買二百塊錢.
好像很划算.
淘下淘下真的很划算.
另外一個機器是什麼呢.
打機.
打機真的很影響我們.
有些人覺得打機只是小朋友.
其實不是的.
打機大人小孩都一起玩的.

$^{361}$打機有什麼問題呢.
通常打得很晚.
第二天早上很晚起床.
這些都是影響著我們的.
雖然我們沒有喝醉酒.
但是我們在手機和打機裡面.
都是很影響我們的.
究竟我們有沒有一種熟練的敏銳呢.
我們不是說不可以.
不是說不可以打機.
而是我們去想一想.
這樣東西會不會合制了我們的生命.
令到我們有一些更重要的東西.
我們忽略了沒有去做呢.
有時我們有時想一想.
有沒有些人其實我們需要關心一下他呢.
有沒有些人我們心裡面.
應該是為他祈禱的呢.
如果我們經常在兩個機器裡面.
困擾我們的時候.
我們自然就少了時間去做這些事情.
今天杜牧師其實買了很多書.
我太太經常說你不要再買書了.
那我現在就不買實體書了.
其實都買的.
剛剛《浸神75週年》又送了兩本書給我了.
現在我就買電子書.
現在我今天看一看我的電子書庫.
我有一萬一千多本書.
我真的想一想.
其實我有沒有機會看這麼多書呢.
這次講這篇《道》的時候.
我大概看了四十本書左右.
我覺得挺有趣的.
現在真的不同了.
我將這四十本書的內容.
fit入我的AI裡面.
然後它就總結了.
當然那些重點是我自己的.
但是裡面很多內容都是.

$^{401}$總合了不同的書籍的內容.
其實不知道大家有沒有時間玩一下AI呢.
現在這樣東西都是沒辦法的.
這樣東西不可以說它是好還是不好.
而是說你不用.
你就一直不敢去用.
還有你慢慢就會落後了.
所以我決定幾十歲的人都去玩一下AI.
但是你知道玩這個東西是花你的時間.
但是你其實要看到你的目的是什麼.
你的目的就是看你的PowerPoint是不是變漂亮了.
以前的PowerPoint你記不記得.
很小的字.
現在我跟它說一定要18點字以上的.
現在它做到了.
字變大了.
其實你要慢慢跟它玩.
其實AI這東西就好像跟它聊天一樣.
你慢慢要求它.
這樣東西最浪費時間.
但是如果你能夠對準這樣東西的好處的話.
那你就可以使用這個工具.
作為教會.
教會會不會有時候是沉醉於一種形式.
數字 效率 傳統 或者品牌呢.
我們會不會事工很多卻是缺少性命的工作呢.
所以我們要進到七個要.
我們要進入神心意的生命.
你留意這七個要可以分為兩方面.
一個是行動方面.
一個是動力方面.
七個要我們最中心的.
我們就是要明白主的旨意.
行動我們就要謹慎行事.
我們要仗志為人.
我們要把握時機.
而動力方面.
為什麼我們能夠有動力去明白主的旨意呢.
因為我們能夠被聖靈充滿.
我們能夠藉著詩歌去讚美.

$^{441}$我們就逐個來看看.
第一我們要謹慎行事.
意思就是要仔細留意自己是怎樣生活的.
我們要到處去查看.
我們留意信仰不是隨隨便便地生活的.
而是我們要有警醒和有查驗的.
所以作為教會我們也要想想.
教會的事工是不是一種慣性的延續.
就是說今年做的事和明年做的事.
和後年做的事是完全一樣的.
而是我們要問.
這件事是不是出於神.
我經常都問自己.
究竟我們做每一個事工的時候.
最終的目的.
是不是隻是這個事工的成功.
還是這個事工能夠讓人的生命更新和改變呢.
大家剛纔留意到.
三福訓練已經是進入第四屆.
我自己看到神是很使用三福的工作.
讓很多弟兄姊妹能夠操練.
他如何和人分享整全的福音.
也讓他們能夠繼續燃點和影響身邊的人.
在這幾年裡面.
我們除了和未信者分享之外.
我們也和一些牧區的領袖去分享.
所以我們都發現在區裡面.
漸漸多一些的人去認識三福是什麼.
但是我不滿足於這個只是一個事工的推動.
我更加要看到.
神要興起領袖在這個工作上.
我們有兩位的姐妹成為三福的信徒教師.
今年我們有另一位的姐妹.
邀請她也成為信徒的教師.
我發現神讓這個工作.
不只是道牧師個人的感受或者感動.
而是能夠讓弟兄姊妹都能夠同心.
能夠一起去做.
教會的事工並不是一個個人的喜好.
教會的事工是我們尋求眾人同心.

$^{481}$我又看到我們今年的《暑期聖經班》.
大家初初覺得那個報名很少人.
現在的報名人數和去年都差不多.
我更加開心聽到有些弟兄姊妹.
他們是第一次參加的.
無論是三福的片傳.
或者無論是《聖經班》的.
都是第一次參加.
我們感受到弟兄姊妹那種熱切事奉的心.
我經常說今天作事工不是為事工而事工.
而是為了服侍人.
為了裝備人.
為了人的生命.
得著神自己的豐富.
第二方面我們要像智慧的人.
這個是對比那些無知的人.
智慧就是按神的真理而活.
不是說圓滑世故.
智慧也不是問什麼是最有效.
而是問什麼纔是最合神的心意.
教會的重點不是成功.
而是忠心.
最近我們聽到一堂講道.
就是新加坡神之心教會.
陳兆豪牧師.
他講那個信息的主題.
就是世上無人知曉的強壯教會.
他強調教會真正的強壯.
知名度或者華麗的事工.
而是來自一群願意安靜忠心.
簡樸地愛神愛人的牧者和門徒.
簡樸牧養的重點.
就是以單純的心去愛耶穌.
也都以真誠的心去愛人.
教會真正的力量.
不在於領袖的魅力或者個人的風格.
而是在於可以被傳承被複制的品格.
魅力可能會吸引人一時.
卻不一定能夠建立人的一生.
個性和才華或者會令人羨慕.

$^{521}$卻不是每一個人都能擁有.
但是品格不同.
忠心 謙卑 真誠 忍耐 愛心和責任感.
是每一個願意跟隨基督的人.
都可以培養的生命素質.
建立在品格上的教會.
才能夠有能力跨越世代.
持續的健康成長.
我最近教教牧書信的時候.
就講到教會領袖的資格.
我就用一個冰山作為比例.
大家知道冰山的上層.
其實只是佔整個冰山的很小部分.
或者容許我這麼說.
可能我們在冰山水面上看到的冰山.
只是整個冰山大概10\%.
90\%的冰山都是在水平線以下的.
我就特別提出一個領袖的條件.
有90\%是在乎他的品格.
20\%是在乎他的才能.
這個是一個比方的說法.
今天我們選擇領袖.
我們很多時候都是看他在人的裡面有什麼條件.
可能我們會選擇年輕的.
我們會選擇有專業的.
學歷高的.
有能力的.
但是聖經選領袖是看品格為先.
這個就是顛覆了我們和世界觀念一個很大的不同.
簡樸的牧養關係提醒我們.
教會不要只是被節目策略或者平臺建立起來.
而是要透過一段又一段真實的關係被建造.
回到愛耶穌,愛人,建立穩格,忠心陪伴的初心的時候.
教會就能夠成為一個不一定眾人知曉.
卻在神面前真正強壯的群體.
第三我們要把握時機.
你留意和合館是要愛惜光陰的.
其實有什麼分別呢.
光陰其實就好像說到一些時間.
可以數的時間.

$^{561}$但是原來聖經原文裡面.
這個字不只是說流逝的時間.
它更加是說一個時機.
一個關鍵的時機.
所以我們有一個科叫Kairos.
Kairos就是關鍵的時機.
不是指一個流逝的時間.
就是將神給我們贖靈的機會贖回來.
把握時機就是我們要贖回這些時間.
贖回這些機會.
邪惡的世代會吞噬機會.
所以我們要主動把握.
保羅使用的是用市場的用語.
好彩人在市場上將有價值的東西買下來.
基督徒要擁有永恆價值的事.
將時間從浪費,罪惡,懶散和無益中買回來.
作為教會要分辨此時此地的照明.
不要錯過神給我們的機會.
鼓勵大家一會兒參加月會.
月會除了收章一讀之外.
我們還會報告一下.
我們接下來要宣傳的下令會.
因為這次的下令會是和之前不同的.
我們之前是在營地舉行的.
我們這次的下令會是在教會舉行的.
而且我們這次的下令會.
我們特別邀請臺灣林慶宗牧師的團隊.
一起去帶領我們.
回想起去年九月的時候.
其實我們初步想.
其實神在烈火特會里祝福了我們很多弟兄姊妹.
其實我兩三年前已經是參加香港的烈火特會.
我們的同工在去年也都參加了臺灣的烈火特會.
其實大家回來的時候都感受到.
在信息當中給到我們靈明.
有很多的眺望和激勵.
我們就想究竟這一個火.
可不可以在我們弟兄姊妹身上.
更多的被燃點呢.
在過程裡面其實我們都經過很多的溝通.

$^{601}$我們經過很多和大家分享異象.
初初我們想著.
如果可以在營地舉行.
我們都真的在營地舉行.
但是結果就出乎我們意料.
就是剛剛營地那個星期.
我們就沒有辦法舉行.
因為營地不是給我們時間.
所以結果我們都在我們自己的教會進行.
如果大家想有興趣瞭解多一點.
可以在月會里面去聆聽去了解.
不過我們很快就會進入宣傳.
透露一點點.
牧師曾經很壞.
喂 這麼好的聚會.
要收錢的.
收錢.
其實我不是收錢.
我收了錢之後.
如果你回到錄像我會退錢給你的.
但是那些童工就說.
不要啦 牧師收錢.
我是鼓勵他出來的.
但是結果怎麼樣呢.
結果反過來.
你來齊錄像.
吸滿印子.
送份禮物給你.
大家是不是開心一點.
是不是 你是不是開心一點.
收錢就個個都不開心了.
所以送禮物給你.
豪給你.
等你得到這個屬靈的滿足.
我們也要學習明白神的心意.
對比我們就不要做糊塗人了.
明白的意思就是綜合領會和辨識.
我們的中心不是問我們要做些什麼.
而是問我要成為怎樣的人.
對教會來說.

$^{641}$不是先問哪一套模式有效.
而是先問神要我們成為怎樣的群體.
所以群體辨識的重點.
是在乎我們要回到聖經.
我們要去.
我們要去信服已知的真理.
我們要持續的禱告.
我們不只是說去策劃.
我們要等候.
我們要去彼此的聆聽.
我們要有群體共同的沉澱.
我們要信服聖靈.
我們不要迷信一些成功的模式.
其實我想想建立教會.
強調我們應該是吸引那些非基督徒.
而非僅僅是轉會的信徒.
不知道你明不明白.
教會人數多有兩種的可能.
一種就是在第二間教會轉來.
第二間教會轉來.
在神的角度裡面其實是沒有分別的.
就是說你都是吸引那些已信的人來這間教會.
但是在整個神的角度是沒有多到的.
但是如果教會重視那些未信的人.
他在這間教會信耶穌成長的時候.
真的對神的抗戰是有增長的.
所以其實一間教會.
哇 為什麼出名 為什麼出名.
其實可能都是那些已信的信徒去說的.
就是未信的人不會關心你什麼教會.
未信的人最重要是認識耶穌.
所以我也提醒自己.
教會的成長是要讓那些真的未信主的人去認識主.
這個纔是我們的焦點.
我們是要有更新生命的屬靈動力.
這個動力第一我們是要被聖靈充滿.
對比我們不要醉酒.
你留意聖靈充滿這個時式是一個被動的命令式.
就是說我們持續讓自己被聖靈的充滿.
留意 追求聖靈的充滿不是追求神祕的經驗.

$^{681}$而是整個生命被聖靈的管理.
包括在敬拜的生活和人際關係之上.
在教會來說 教會不是靠運作來維持.
而是靠聖靈的更新.
不知道有人是說笑還是說真的.
他說今天即使聖靈離開了教會.
教會仍然像一部機器一樣如常運作.
去聽的時候好像很隱移.
沒有聖靈教會怎麼運作呢.
但是是真的.
如果教會只是按照規劃 按照計劃.
按照我們每年做的事重複去做的時候.
那就真的好像一部機器一樣走.
所以沒有聖靈我們也沒有了辨識.
我們也要用詩篇讚美詩 靈歌去讚美.
敬拜不只是音樂的安排.
而是聖靈充滿的群體生活.
留意 經文講到彼此對說.
表明敬拜是具有一個群體建立的功能.
口唱心和表明敬拜要真誠.
我記得我上一科教會音樂的悼論的時候.
老師就告訴我們.
教會音樂有什麼作用.
教會音樂就是讓弟兄姊妹對神.
對信徒 對未信的世界.
去分享神的信息.
我覺得這是一個很重要的定義.
教會音樂的目的.
不是讓弟兄姊妹感到舒暢.
教會音樂的目的是指向人去尊榮上帝.
教會音樂的目的是彼此對唱.
讓弟兄姊妹能夠在音樂裡彼此得激勵.
教會音樂的目的是讓未信的世界.
他們能夠得聞福音.
所以真的敬拜是會幫助群體.
重新調教它的方向.
一間沒有敬拜深度的教會.
很難有辨識的深度.
其中一個指標就是會眾能夠欣賞不同風格的音樂.
你今天如果上午課.

$^{721}$你就會和早課不同.
因為午課我們有師班.
我們都有不同的方式去敬拜.
有些人告訴我.
哎呀 牧師很不習慣.
敬拜的詩歌很不習慣.
其實你要學一下.
擴闊你對教會音樂的欣賞纔行.
什麼音樂其實不重要.
甚至不是合我口味最重要.
最重要就是這個音樂能夠幫助我們對準神.
今天神要復興稍鎮的敬拜律師班聖惡的團隊.
要能夠去除自我.
讓我們能夠同心合一.
感恩不是否認困難.
而是在任何的處境中仍然承認神掌權.
凡事上上顯示感恩是一種生活的方式.
不是因為每件事本身都容易成為美好.
而是因為神知道什麼是對我們最好.
祂帶領我們的道路.
雖然現在有時看起來好似奧祕.
但是將來卻必顯明為良善和憐憫.
感恩的文化使教會更能分辨神的作為.
5月1號我聽到有領袖分享.
他說在2017年聽了一篇要感恩的信息之後.
他們的小組持續至今.
都是在實踐感恩.
反過來抱怨是會叫人只看環境.
而感恩是會叫人重新看見神.
接下來星期三我們用累代這個主題.
去分享一個牧養關係.
特別常常提醒我們.
神上上揀選那些在社會中被忽略.
被輕看甚至被邊緣化的年輕人.
親自賦予他們力量身份和使命.
使他們能夠在這個世代中為主發光發熱.
所以我們要重視我們的每一代.
在我來說我的上一代是一群.
在教會里面忠心服侍的長輩.
在我來說下一代.

$^{761}$就是我不太明白他們挑戰的年輕人.
但無論如何真正的牧養.
不是追求數字,成果或者外在的成功.
而是出於愛耶穌.
也學習真實的去愛人.
牧養的核心是願意看見每一位信徒的生命.
關心他們的處境.
陪伴他們走過軟弱和掙扎.
耶穌教導我們.
牧者的心不是隻滿足於九十九隻羊仍在羊圈當中.
而是願意為那一隻迷失的羊走出去.
尋找它,擁抱它.
並且將它帶回群體之中.
所以教會群體的本質.
就是要以愛心,責任感與真誠的關係.
來對待每一位的會友.
而不是將人視為完成事工.
達成目標或者建立個人成就的工具.
今天說起牧養.
我經常都有個感受.
我覺得自己真的不是那些很貼心的傳道人.
很多時候都不是一些主動去虛寒問暖的傳道人.
所以我不是那些很父親型的傳道人.
但是不重要.
因為在我們的童工隊裡面.
有不同的傳道人.
他們都能夠負起這個責任.
最近有一個建議.
就是叫我們的區域去不同年齡的牧區.
不要只是去自己的牧區.
去人家的牧區看看.
明白人家的挑戰困難.
我覺得這個挺好的.
找一次我去青少年區.
找一次我去長青區.
這樣我們就能夠瞭解到他們的需要.
是的.
傳道人如果你都不瞭解.
你那些弟兄姊妹的挑戰.
有時候我們留意什麼呢.

$^{801}$閉門造車.
我們只會罵那些弟兄姊妹不愛主.
不熱心.
不祈禱.
不靈脩.
但是我們有沒有一個牧者心呢.
我們都求主幫助.
保羅在《爾弗所書》五章十五到二十節.
開始轉入一段積極和具體的渲染.
說明信徒在邪惡充斥的世代中是怎樣生活.
信徒不能夠對現實的處境調以輕心.
也不能消極地期待環境自然改善.
相反.
我們必須警醒謹慎地.
省察自己的行事為人.
敏銳洞察周遭世界的實況.
並且竭力明白主的旨意.
對於當時文化中普遍存在的問題.
好像醉酒放蕩.
信徒不可以隨波逐流.
反而應當活在聖靈的引導和掌管之下.
在基督裡面不斷邁向成熟.
因此信徒群體會在聚會中.
以詩歌彼此造就.
因為上帝藉著他們與主耶穌基督的聯合.
已經豐豐富富地賜下恩典和福分.
所以在複雜多變的世代裡.
上帝呼召我們.
不僅僅是個人去作智慧人.
也呼召教會成為一個會群眾辨識.
會禱告等候.
會敬拜感恩.
會信服聖靈的群體.
教會前面的路不是靠人的規劃.
不是靠模仿成功的模式.
而是靠上帝的話語與聖靈的帶領.
一同明白主的旨意.
讓我們一起禱告.
三不要七要.
就是這段經文給到我們的信息.

$^{841}$很豐富很多.
可能今天大家未必全部可以領受.
但是可能神就給你一個點.
就給你一個提醒.
聖靈就藉著裡面的內容來幫助你.
所以我們願意弟兄姊妹.
在你的位置上你自己去禱告.
究竟哪一點今天是激動你的.
提醒你的.
你求神幫助你.
去克服我們生命裡面的軟弱.
靠主去得勝.
願意今天的信息.
能夠提醒我們.
如何去活出一個.
遵行神旨意.
遵行神心意的生活.
也都求主的聖靈繼續的催促督導.
我們能夠以回應禰的道.
在生活中實踐.
為我們的中心.
願天父幫助我們.
去帶領我們.
預納隨聽我們的禱告.
奉基督寶貴名字祈求.
\newpage



\section{羅馬書 1:1-15}
\label{sec:032089}
\textbf{持定福音使命的人生}~~~~黃志剛牧師
\newline
\newline
連結: \href{http://www.church.com.hk/acms/content.asp?site=cdc\&op=show\&type=product\&code=032089\&layout=sermon}{\texttt{032089}} ~~~~ 語音日期: 2025-10-04
\newline
\newline
\hyperref[sec:032088]{\small{< < < PREV SERMON < < <}}
~
\hyperref[sec:index_chronic]{\small{[返順時目]}}
~
\hyperref[sec:032090]{\small{> > > NEXT SERMON > > >}}
\newline
\newline
羅馬書 1:1-15
\newline
\begin{longtable}{cl}
\hline
\hline
章節 & 經文 (和合本修訂版)\\
\hline
1:1 & \begin{tabularx}{0.7\textwidth}{X} 基督耶穌的僕人保羅，蒙召為使徒，奉派傳神的福音。 \end{tabularx} \\ \\ \relax
1:2 & \begin{tabularx}{0.7\textwidth}{X} 這福音是神從前藉眾先知，在聖經上所應許的。 \end{tabularx} \\ \\ \relax
1:3 & \begin{tabularx}{0.7\textwidth}{X} 論到他兒子－我主耶穌基督，按肉體說，是從大衛後裔生的； \end{tabularx} \\ \\ \relax
1:4 & \begin{tabularx}{0.7\textwidth}{X} 按神聖的靈說，因從死人中復活，用大能顯明他是神的兒子。 \end{tabularx} \\ \\ \relax
1:5 & \begin{tabularx}{0.7\textwidth}{X} 我們從他蒙恩受了使徒的職分，為他的名在萬國中使人因信而順服， \end{tabularx} \\ \\ \relax
1:6 & \begin{tabularx}{0.7\textwidth}{X} 其中也有你們這蒙召屬耶穌基督的人。 \end{tabularx} \\ \\ \relax
1:7 & \begin{tabularx}{0.7\textwidth}{X} 我寫信給你們在羅馬、為神所愛、蒙召作聖徒的眾人。願恩惠、平安從我們的父神和主耶穌基督歸給你們！ \end{tabularx} \\ \\ \relax
1:8 & \begin{tabularx}{0.7\textwidth}{X} 首先，我靠著耶穌基督，為你們眾人感謝我的神，因你們的信德傳遍了天下。 \end{tabularx} \\ \\ \relax
1:9 & \begin{tabularx}{0.7\textwidth}{X} 我在他兒子的福音上，用心靈所事奉的神可以見證，我怎樣不住地提到你們， \end{tabularx} \\ \\ \relax
1:10 & \begin{tabularx}{0.7\textwidth}{X} 在我的禱告中常常懇求，或許照神的旨意，最終我能毫無阻礙地往你們那裡去。 \end{tabularx} \\ \\ \relax
1:11 & \begin{tabularx}{0.7\textwidth}{X} 因為我迫切地想見你們，要把一些屬靈的恩賜分給你們，使你們得以堅固， \end{tabularx} \\ \\ \relax
1:12 & \begin{tabularx}{0.7\textwidth}{X} 也可以說，我在你們中間，因你我彼此的信心而同得安慰。 \end{tabularx} \\ \\ \relax
1:13 & \begin{tabularx}{0.7\textwidth}{X} 弟兄們，我不願意你們不知道，我屢次計劃往你們那裡去，要在你們中間得些果子，如同在其餘的外邦人中一樣，只是到如今仍有攔阻。 \end{tabularx} \\ \\ \relax
1:14 & \begin{tabularx}{0.7\textwidth}{X} 無論是希臘人、未開化的人、聰明人、愚拙人，我都欠他們的債， \end{tabularx} \\ \\ \relax
1:15 & \begin{tabularx}{0.7\textwidth}{X} 所以願意盡我的力量把福音也傳給你們在羅馬的人。 \end{tabularx} \\ \\
[1ex]
\hline
\hline
\end{longtable}
$^{1}$『錄音或錄影』.
此錄音或錄影的版權是屬於基督教宣導會北角堂所有.
我們的屬靈傳統之一是以真理教導為本.
歡迎你將此錄音或錄影的連結傳給親友.
懇請注意未經授權請勿轉錄或剪輯本錄音或錄影.
甚願好像《仕途行傳》20章32節下所說.
這道能建立你們.
叫你們和一切誠誠的人同德基業.
剛才我們看了周牧師.
他為這一季的講題《羅馬書》做了一個介紹.
大家應該都會記得他說.
第一章去到第十一章.
如果你有看《羅馬書》.
你知道《羅馬書》總共有十六章.
一至十一章一般是講教義比較多.
剛才我們聽到周牧師講了很多不同的部份.
十二章開始我們很熟悉.
所以弟兄們我以神的慈悲勸你們.
一路背下去.
講到教義的實踐.
我們這個季度是選擇先看教義這個部份.
剛才我聽主席讀經的時候.
我又想起一個片段.
我想起我信主初初的時候.
你知道開始看聖經.
我不知道你信主初初的時候.
你是從哪一本書開始.
我猜一般是看福音書.
我想我應該也是.
但很快我就發覺.
我很喜歡看《羅馬書》.
那時候覺得他講的東西很清楚.
自己又是讀理科的人.
又很理性.
所以覺得《羅馬書》講救恩.
講福音.
講得很清晰.
當然慢慢信主來的時候.
又會看其他的新約舊約的書卷.
所以當剛才看介紹的短片.

$^{41}$聽主席讀的時候.
又想起原來曾幾何時初信的時候.
我喜歡看《羅馬書》.
很盼望這個季度.
當我們透過《羅馬書》一至十一章.
去領受訊息的時候.
或者你都可以花多點時間.
這一季你可以多看幾次.
《羅馬書》裡面的內容.
我深信對你.
特別對於福音.
好像剛才周牧師說.
不是單單停留在我們信主才聽.
其實福音,好消息,上帝的救恩.
歷久常新.
都會對我們生命有幫助.
我們一起禱告.
把這段時間讓給上帝.
神下示的如果不是.
主耶穌基督為我們的罪釘死.
並且從死裡面復活.
我們今天.
是一個沒有盼望的人.
我們的人生.
也談不上有什麼使命.
因為我們的人生就是沉淪.
就是沒有希望.
所以主今天我們能夠成為.
何等樣的人.
是蒙你的恩才成.
我們帶著期盼的心.
我們恭敬的仰望上帝.
願聖靈帶領我們進入真理.
不單只是今天.
在往後的日子當中.
讓我們從羅馬書裡面.
我們再一次得到上帝的提醒.
教導,安慰和鼓勵.
願你帶領我們以下.
一起聽道,講道的時間.

$^{81}$願我們都一同蒙恩.
我們禱告奉耶穌基督.
成名其琦.
阿們.
羅馬書.
給人的感覺是一封.
好像講教義.
的一個講章.
但其實羅馬書像剛才大家所聽.
它是一封書信.
新約的書信.
或者希臘的書信.
它有一個基本的格式.
你看到這裡有1,2,3,4,5,6.
有六個部份.
當然頭幾個部份.
其實不是佔很多篇幅.
在羅馬書裡面.
我們剛才讀的是.
1至15節.
一章1至15節正正就是.
源於正文之前.
所以.
今天不會講.
我們最熟悉的福音本是神的大能.
一章16節.
這個留給下個星期的同工.
跟大家分享.
我們今天會看的是頭15節.
我們說是序言.
在序言裡面我們看到的劃分.
包括了講這封信是誰寫的.
寫給誰.
有文案.
然後在新約的書信.
通常都會有.
文案或者感恩.
或者禱告這個部份.
這個部份有個特別的地方.
就是如果相比起.

$^{121}$接下來那些比較較易的.
硬邦邦的.
屬靈的道理.
這個我們會感受到.
保羅那種的情在裡面.
一會兒我們會一起慢慢再看.
我們就不如進入.
這個部份.
我就先將.
羅馬書第一章1和7節.
將寫信人.
和收信人這兩個.
放在一起看.
羅馬書第一章第一節.
基督耶穌的僕人保羅.
蒙召作使徒.
奉派去傳神的福音.
這個第一節.
是作者的自我介紹.
你看到他介紹.
他提到有三方面.
他是耶穌基督的僕人.
他是蒙召.
做使徒.
奉派去傳神的福音.
某程度上是解釋給你聽.
他作為使徒.
他的使命是什麼?.
他的任務是什麼?.
就是奉派去傳神的福音.
奉派去傳.
他原本的意思是.
特別是預留了.
分別出來.
是專做這件事.
所以某程度上.
我們可以說.
保羅是一個福音的特派專員.
他領受了上帝.
特別交託給他的一個託付.

$^{161}$一個使命.
他蒙召去做使徒.
要去傳福音.
而當他做這件事的時候.
其實他是耶穌基督的僕人.
僕人我們現在聽得很多.
我們覺得都沒有什麼特別.
但僕人一方面.
聖經學者說在舊約來說.
僕人其實是一個很尊貴的身份.
另一方面.
你知道在羅馬時代.
仍然都有奴隸.
奴隸或者僕人.
這個字更準確.
有些翻譯是奴僕.
所以當羅馬人.
一見到這個字.
一聽到他就知道.
他會想起的不只是基督耶穌的奴僕.
也會同時想起.
在他社會裡面.
可能每個家庭裡面.
都有奴僕.
在這個背景裡面.
保羅這樣去自稱.
他有這三個的身份.
很清楚.
保羅一開始講很清楚.
他自己的身份是怎樣.
他有什麼要去做.
是很清晰的.
如果用我們今天的說話講.
他的人生方向是很清晰的.
他的目標也都很明確.
他知道自己.
起碼在他順處得救之後.
事實上我們知道.
無論羅馬書,無論《使徒行傳》.
無論保羅書信.

$^{201}$保羅正正就是在他最後的一刻.
他仍然去到.
凱撒的面前去見證福音.
他很知道自己.
在做什麼.
不是他去尋找一些什麼的夢想,理想.
而是上帝.
將他的召命.
將他的使命.
交了給保羅.
所以今天的講題說.
持定福音使命的人生.
使命某程度上.
不是你去尋找.
是上帝呼召你的.
蒙召這個字.
在下面講到修信人.
也是保羅口中的表達.
在第七節說.
給你們所有在羅馬.
為神所愛蒙召作聖徒的人.
又是三樣.
第一.
是給在羅馬.
羅馬書.
很簡單很清晰.
但他如何形容羅馬這群信徒.
這群修信人.
第一他們說.
他們是神所愛的.
他沒有提到他們有什麼社會地位.
有什麼背景.
他只是很快就說.
你們是神所愛的.
而且不只是神所愛.
神也是對你有.
期盼.
有使命交託給你.
要做一個.
分別為聖的門徒.

$^{241}$聖徒的意思是.
最基本.
你要和未信的人.
和當時當代的外邦人.
有所不同.
你看到PowerPoint.
都highlight了兩個字.
一模一樣.
無論是對保羅.
還是對保羅所寫信的對象.
他都很看重.
是蒙召.
這個身份.
是被動式被召喚.
是上帝召喚你的.
是上帝呼召你.
讓你成為使徒.
讓你們成為.
聖徒.
有時.
當我們想想.
這麼多東西可以描述.
但保羅正正一開始的時候.
他只是強調.
羅馬教會.
你們是神所愛的.
神對你們.
有呼召交給你.
對你們有盼望.
對你們有期望.
你可能會說.
不要吧,很大壓力.
是的,有時對於.
一些對你有期望的人.
因為正正期望你可能做一些事.
或者要說一些說話.
等等,你會覺得好像有些無形的壓力.
不過你又明白.
如果反過來.
對你沒有期望.

$^{281}$我們用今天的表達.
你的存在感會很低.
我想我們保羅不是要.
將一塊大石壓下去羅馬的信徒.
只不過.
這個就是我們作為基督徒.
蒙神救贖一個客觀.
也是主觀的身份.
我們是被召喚的.
所以在第一章第六節.
其實剛剛上面.
他也這樣說.
其中也包括了你們這些蒙召.
後面多說一點.
蒙召是什麼?.
蒙召歸屬於耶穌基督的人.
所以簡單來說.
原來這三次這麼短的裡面.
這個序言的裡面.
保羅刻意三次用到.
蒙召這個表達.
他很強調的是.
這個蒙召是上帝做主動的.
是上帝主動去.
拯救人,呼喚人.
你回應我.
你信主,你來到我這裡.
歸屬於耶穌基督.
你去領受這個使命,去實踐.
去傳福音.
其實你想想,是很神聖的.
就等於.
我喜歡用一個例子.
你有沒有參加過.
運動會?.
你知道參加運動會有時候有什麼.
初賽,福賽那些.
我不知道現在是不是一樣.
我記得我以前做老師的時候.
運動會都要幫忙.

$^{321}$有時候其中一個崗位.
就要宣佈.
以下同學.
誰寫的,將會進入決賽.
接著就會讀名.
讀到名字是什麼意思?.
你有份.
你聽不到你的名字即是什麼?.
你進不了決賽.
那個被呼喚的名字.
不是單單說你的名字.
三個字,兩個字.
而是呼喚你,就表示.
你OK,你被選上.
你可以進入下一個比賽.
被召喚.
被選召.
是很重要的.
特別對於我們基督徒來說.
這是上帝定義.
認得我們,呼喚我們的名.
叫我們歸向祂.
所以這個第一個.
在那個簡單.
只不過是說.
修信人和寫信的人.
已經有一個.
隱含著本來要說的話.
所以第一方面.
我想做一個總結.
我和你都是蒙召.
屬於耶穌基督.
我們是分別出來.
甚至我們分別出來.
不是純粹自己為威位.
我們是有使命的.
這個在後面的篇幅我們會再看.
然後當說了.
修信的人.
也說了寫信的人.

$^{361}$修信的人之後,第七節的下半節.
就是一個書信格式的文案.
願我們的父神.
和主耶穌基督.
賜予你們恩惠和平安.
這裡我不特別.
花時間要去講解.
很多,只是在保羅的書信.
或者在福音書信.
在新約的年代.
很多時候的文案.
他們只有恩惠.
因為比較多是當時.
外邦人的書信都會講恩惠.
但是恩惠不單只是.
所謂的Grace.
也是有舊約所講的.
平安.
所以有學者說.
其實恩惠平安就是在講.
新約舊約互相輝映.
現在保羅說.
從神,從耶穌基督而來.
那個祝福既有平安.
也有恩惠.
李深水清.
你見到.
第一節講寫信人.
第七節講收信人和平安.
那麼中間那五節.
吃了些什麼進去呢?.
根據那個格式.
其實應該完結了.
保羅就是將了很多.
他很想講.
或者在這裡提綱揭領的地方.
放在了這幾節.
我們一起讀.
這兩節.
預備起.

$^{401}$基督耶穌的僕人保羅.
夢照作聖徒.
奉派去傳神的福音.
這福音.
是神直著眾先知.
在聖經上.
預先應許的.
這裡你特別highlight的是.
神的福音.
保羅一開始說他是奉派.
他是傳福音的特派宣言.
他沒有停下來.
他接著就解說.
其實這個福音是什麼.
首先他說.
第二節說.
這個福音原來不是新的東西.
或者說不是全新.
與舊約無關.
因為這裡說是神直著眾先知.
而且在聖經上.
預先.
宣告或者應許.
如果你說舊約.
聖經中先知.
你馬上想起的書卷就是.
先知書.
先知書不是很多.
在希伯來的正典劃分.
其實很簡單.
摩西五經就是律法.
接著有很大的篇幅.
我們說前先知書.
後先知書.
是包括了很多.
我們除了先知書以外.
歷史書等等的篇幅.
最後一個部分就是詩篇.
其他這些著作.
所以舊約聖經或者新約的時候.

$^{441}$作者都是這樣提的.
耶穌基督在他復活.
升天之前的顯現.
在路加福音24章都是這樣說的.
他說舊約聖經對他的預言.
其中說過的包括了.
律法,先知.
就是這裡.
所以這裡說.
這個神的福音原來.
保羅第一樣東西就是.
他是有前科.
還是有一個.
隱藏了尼賽亞會來的訊息.
舊約聖經說了.
接著在第三至四節.
他就用了一個.
比較的方式.
去表達.
第三節就說.
當中論道他的兒子.
按肉體說.
他生為大衛的後裔.
按聖潔的靈說.
他因從死人中復活.
立為神大有能力的兒子.
他就是我們的主.
耶穌基督.
你看到這裡就是說.
那個福音.
究竟內容是什麼.
你說福音內容很多的.
某程度上整本羅馬書一至十一章都在說.
不過在引言裡面.
保羅當然.
他是偉大的使徒.
你也知道他背景很熟悉律法.
他的學識又很淵博.
所以他那些.
顯示造句應該是很精煉的.

$^{481}$他就簡單用了兩節聖經.
他就將福音.
一個最核心的內容.
撮要了下來.
他說首先這個福音是關於誰的.
關於.
他的兒子.
他是在說上帝的兒子.
這位上帝的兒子.
可以用兩個.
階段去看.
按肉體說.
這位上帝的兒子.
其實他生的時候.
是大衛的後裔.
是生出來的.
但是接著作者就說.
按聖潔的靈說.
這位兒子.
是從死人中復活.
換句話說.
他經歷過出生.
經歷過死亡.
經歷過復活.
是不是很厲害?.
他只是用了.
兩句幾個字.
他就已經告訴你.
這位兒子.
這個福音的內容.
其實就是.
生死復活.
都經歷過的這位.
既是大衛的後裔.
人的兒子.
也是上帝的兒子.
而且這裡說.
是上帝大有能力的兒子.
聖經學者在這裡研究.
什麼意思?為什麼大有能力?.

$^{521}$相信是和.
上面聖潔的靈.
有些關係.
又是聖經學者在這些字眼.
用了很多的分析.
簡單來說.
比較多學者會接受.
第一,聖潔的靈其實簡單就是聖靈.
如果和下面要拉上一個關係.
更加是聖靈的能力.
耶穌基督.
成為人的兒子.
他死了,被釘在十字架上.
但是因著聖靈的能力.
或者藉著聖靈的能力.
神使他從死人中.
復活.
而上帝.
也定立他.
或者委立他.
成為一個.
都是神的兒子,不過已經不同了.
因為他勝過死亡,他復活.
所以他大有能力.
他基本上將.
福音的V字.
出生,死亡.
降在陰間.
然後復活,升天.
很簡潔地說了出來.
然後保羅最後做了一個總結.
這位從死人中復活.
又是大衛後裔.
有神性又有人性的.
這位兒子,他是誰呢?.
他是我們的主.
耶穌基督.
我們剛才讀.
《使徒信經》好像有句是這樣的.
主.

$^{561}$耶穌.
基督.
還有.
我們的主.
所以保羅.
不單止奉派去傳.
福音.
他某程度上.
好像故意要將.
福音用一個.
很簡單的方式再說出來.
讓這班羅馬的信徒.
可以聽到.
讓他們知道為什麼保羅.
會故意跟他們說多一點.
或是表達那種很想見他們的心情.
保羅.
他作為福音的特派專員.
他要很熟悉.
他傳福音的內容.
就好像.
你一會兒走的時候.
你要買一件電器.
你不是很懂.
你走進去.
通常有個人就會過來.
問你先生或女士.
有什麼可以幫到你.
你說我想買一部洗衣機.
他就會跟你.
推銷或介紹.
有些.
我不知道你們有沒有遇過一些.
就是滔滔不絕.
口若鉛河的.
他講得很詳盡.
我問什麼問題他都能回答.
但如果下一個.
sales.
你跟他談的時候.

$^{601}$他說不知道.
你看說明書說什麼.
那你可能心裡想.
你怎麼做sales的.
你還要是這個產品.
起碼你也要讀熟他.
你才肯跟我說.
你明白我這個比喻想表達的東西嗎.
保羅.
他是.
特別安排去傳福音.
他一定要很清楚福音.
是什麼.
應該怎樣理解.
傳讓人聽得進去.
今天我們.
我不知道你信了住多少年.
你對福音的理解.
有多少呢.
你能不能講到福音.
一個最基本的內容.
可以用這個幫幫你.
後面也看到.
五色珠.
這是其中一個.
姐妹做的.
十字架.
既有形狀也有顏色.
每一個顏色都是代表著.
福音的訊息.
今天不花時間講解.
其實我們.
福音使命.
我不懂得講.
我怕.
弟兄姊妹.
我們試試學.
你不一定要做保羅.
你不一定要做.
牧者傳道.

$^{641}$但如果你還是那位.
剛才說的.
蒙召作聖徒.
有使命人生的基督徒.
你大體上也要.
掌握一下.
福音的基本內容.
甚至學一下.
怎樣去表達.
所以當我們.
看到保羅表達了.
福音的內容之後.
他接著就說.
這個福音的內容其實和.
他使徒的積分是很相關的.
我們一起讀這兩節.
預備起.
「直著他,我們領受了.
恩典和使徒的積分.
為了他的名.
在萬族中使人.
因信而信服.
其中也包括.
你們這些.
蒙召歸屬於耶穌.
基督的人」.
由第三,第四節.
說到福音.
一個很核心的訊息.
然後他就說到.
這個核心的訊息其實和他.
或者和一班信主的人.
這裡用「我們」.
有什麼關係呢?.
當我們說我們領受福音.
我們信福音.
原來我們還有一個.
全福音的積分.
這個也是保羅的託付.
這裡你看到綠色那裡說.

$^{681}$「直著他,為了.
他的名」.
這個「他」是在說.
基督.
這裡很清楚說了很多方面.
第一,這個福音.
其實它的本質.
或者這個.
全福音積分的本質.
是什麼呢?.
不是我去賺取回來.
不是我很聰明.
我很聰明,我背今句.
背了一百句,我聖經考試.
很高分,不是.
完全是說.
直著基督.
也在我們去全福音的時候.
這個是領受了.
恩典和使徒的積分.
意思不是有兩樣獨立的東西.
而是想強調.
使徒這個積分本身.
是恩典.
不是每個人.
都有作為使徒.
特別是當時來說.
十二使徒,保羅是有人說.
第十三位.
對於保羅來說,這個積分.
是上帝恩代他的.
特別是我們知道保羅的歷史.
或者說知道保羅的.
黑歷史.
你更加會想.
我哪有資格.
去全福音.
我連得到福音都可能沒有資格.
我是迫不及待的.
你記得嗎?.

$^{721}$拉著人去耶路撒冷.
保羅就是在大馬士革.
的路上被光照.
當時他不是在讀聖經,祈禱.
他當時是在壓榨.
信基督的基督徒.
去耶路撒冷.
保羅是蒙上帝的憐憫.
能夠被光照.
被呼照.
蒙照.
接受耶穌基督.
然後蒙照成為使徒.
去將他之前所鄙視的福音.
去廣傳.
這個.
使徒的積分.
是恩典.
弟兄姊妹,我們得救都是恩典.
你同意嗎?.
我們有幾.
體會.
這件事.
有時我們做基督徒做得久.
我們有時很多事都可能是.
習以為常.
讀經,祈禱,靈修,崇拜,主約學.
回團契.
好像都是我們的生活日常.
但生活日常.
有時最怕就是.
很多事都是指定動作.
我們唱詩.
會不會是念口旺.
我們讀使徒信經.
會不會都是照背如儀.
但裡面說什麼.
或者你讀的時候.
那一次讀的時候.
你覺得責心.

$^{761}$你覺得.
聖徒相通,原來很美好.
你會不會有這些感受.
弟兄姊妹.
我們不要take it for granted.
保羅.
他很重視.
使徒的積分.
亦因為這個緣故,他知道不是為了他自己.
不是他出名.
不是他是大使徒.
其他的神特別呼召我.
他說我去做這個福音的使者.
是完全為了.
耶穌基督的名.
而他很清楚.
他領受這個積分.
他有個目標去做.
就是在萬族裡面.
相對於可能.
彼得,雅各,約翰.
他們在耶路撒冷對著猶太人.
保羅是猶太以外的.
各邦各族.
剛才我們讀經會聽到希臘人.
未開化的人.
等等.
這個是神託付給他的負擔.
等於今天有弟兄姊妹.
我知道你參與很多不同的事奉.
有人在施班.
施班聖餓.
是神給你們的負擔.
有人可能是相見群體的服侍.
這個是他的負擔.
每個人的負擔都有不同.
但都是很尊貴.
這個是上帝給.
我們知道.
我們是為了.

$^{801}$上帝去服侍.
保羅的負擔是說.
在萬族裡面使人信而信服.
不只是做信徒.
不只是得個頭腦.
不是拿到天堂入場券.
幫人申請了.
給了你了,你有個護照了.
他是要人因著相信耶穌.
而信服神.
有上英文詩歌的歌名.
Trust and Obey.
信靠.
信服.
你見到嗎?.
保羅很言簡意賅.
但他亦都將.
使徒的職份.
本質是甚麼.
為了誰而去做.
對著誰去做.
等等都講了出來.
所以我稱呼這一章.
一章五至六節.
我稱呼他做一個使徒.
或者一個使命的宣言.
是保羅作為一個蒙召的使徒.
他去傳福音.
一個使命的宣言.
我不需再重複.
這個使命宣言包括了一些甚麼的內容.
剛才都提過.
使命有個特別的地方.
在整個過程裡面.
在整個描述裡面.
使命.
是不是為自己的.
我們不會說.
我為自己做某件事.
我要完成這個使命.

$^{841}$使命是託付給你.
是為別人.
是為上帝.
你為自己的那個不叫使命.
最多給你的那個是你的夢想.
你的志業.
但當講到使命.
講到呼召.
那個對象.
為誰的一定不是自己.
一定是為了別人.
你回到北村.
如果多過十年的話.
你應該會記得十年之前.
或者再早一點.
2010至2015.
當時教會的主題.
成為使命.
教會.
你可不可以讀下一句出來.
我聽到了.
使命.
人生.
所以某程度上.
弟兄姊妹.
你不可以只是做信徒.
你是要做門徒.
你是要做使徒.
使徒的意思是.
你要帶著福音的使命.
無論你的對象是誰.
無論你有什麼方式.
不在於你有沒有口才.
不在於你是否聰明.
那些不重要.
只要告訴你.
你的使徒.
你的傳福音的職份.
是很尊貴的.
是為了上帝的緣故.

$^{881}$然後你就祈禱神.
你給我的負擔在哪裡.
這段時間我們有一批弟兄姊妹.
去了韓國.
在之前有不同的弟兄姊妹.
去不同的地方去做短宣.
甚至有些在猜拳聯會.
獻身去做宣教士.
去做傳道人.
你是一個家庭主婦.
你是一個在公司裡工作的人.
你都可以有福音使命的.
我聽過有些弟兄姊妹.
他們在職場裡.
他們嘗試聯絡基督徒開一個賣場.
開頭就是彼此打氣.
在一個微信的環境裡做基督徒是很難的.
慢慢他們發覺.
很想真的做點事.
然後就在他們的公司.
在他們的部門裡.
他們試試有些福音的工作.
可能有些BSF.
還有EBS的查經班.
或者搞一些聖誕的聚會.
然後邀請一些同事去參與.
可以的.
你想到很多不同的方式.
我很簡單的.
我做後援.
我做下手工都可以的.
我就做這些五色珠.
做完之後成為福音工具給別人去傳福音.
這個可以是上帝給你託付你的使命.
每一個人都可以有.
所以弟兄姊妹.
你想使命是什麼?.
這個是你要去回答的問題.
無論你是一個什麼的年紀.
你仍然都是帶著使命.

$^{921}$去活你的一生.
所以第二方面.
我想跟弟兄姊妹分享的是.
你要掌握那個福音.
你要活出上帝給你的使命.
這部份就完結了.
在那個序言裡面的一至七節.
當我們轉到第八節.
其實第八節.
如果你記得應該是進入那封書信.
講到感恩的部份.
所以你看到第八節.
首先我藉著耶穌基督.
因為你們的信心.
傳遍了天下.
剛才那句有點類似.
又是藉著誰呢?.
藉著耶穌基督.
為了誰呢?.
上次是為了神.
這次是為了你們.
這群修信人.
去獻上感恩.
但這裡他加了一個字.
我這麼自私的自肥.
神只是你的,你霸佔著.
不是的.
他提到個人的.
這個上帝.
是我的上帝.
當你說我的.
某程度上就是你很重視.
你很渴慕和上帝的關係.
你試想想.
如果不是你想和他拉什麼關係.
你多大都不敢說.
彈開,什麼我的朋友.
什麼我的拍檔.
我們已經很久了.
你會這樣.

$^{961}$保羅這裡說.
他要為這群信徒.
如果我們看下面.
你會看到原來他和上帝.
都有一個很親密的關係.
他就在神這裡這樣用.
然後他說這個原因.
因為你們的信心.
傳遍了天下.
天下,你就不要說.
全世界,中國聽不聽到.
印度聽不聽到.
你知道這個有一點誇張的表達.
其實就是說.
你們的信心傳遍到一個地步.
廣為人知.
我們有時候說話都是這樣.
全世界都對不起嗎.
你只是誇張表達.
你覺得很多人都令你很辛苦.
這裡的信心也不是.
有一個意思就是說.
羅馬書,羅馬教會.
信心比較厲害.
你們的信心傳遍天下.
哥林多教會的信心.
傳遍只是哥林多.
不是.
信心這裡不是說多還是少.
他這裡說的信心.
或者說因為你們的信.
你們那種對上帝的信靠.
決意去相信.
是眾所周知的.
你再感受一下這句說話.
保羅說的時候的心情.
我相信他真的很雀躍.
去為羅馬的信徒去感恩.
因為他竟然知道.
他聽到的是.

$^{1001}$羅馬信徒.
羅馬教會.
北宣家.
你們對上帝的信.
整個香港都聽到.
你馬上就明白了.
不是說我們特別厲害.
而是如果作為一個.
我儘管這樣借用.
如果蕭牧師.
周牧師.
當他聽到有人跟他說.
北宣家的弟兄姊妹.
個個都很愛主的.
很委身的.
你作為堂主任.
你開心嗎?.
你感不感恩?.
現在在家裡.
光宗耀祖.
保羅帶著一份這樣的心情.
首先去感恩.
接著他就禱告.
你見到「神」這個字出現嗎?.
首先這兩節.
我想先和你集中看對於人.
「神就是我在他兒子的福音上.
用心靈服侍的那一位.
他可以給我作見證.
作甚麼見證?.
我怎樣常常在禱告中.
不住地提到你們.
並且祈求.
或許按著神的旨意.
終有一天.
我能夠順利到你們那裡去.
到你們那裡去」.
我們首先聚焦看的就是.
他為信徒感恩.
很開心.

$^{1041}$因為他們的信心全遍天下.
然後保羅就說.
他就說.
其實是.
「我常常在禱告中.
不住地提到你們.
並且祈求」.
其實他原本的意思更濃烈.
他原本的意思是.
常常記得.
和提起你們的名字.
在禱告中.
不住地祈求.
懇求.
如果你發覺.
你在你的代禱名單.
你的群組裡.
你常常都被.
不住地提到.
懇切地為你祈禱.
你應該是開心的.
雖然你說不好,通常那些當然不是病.
甚麼才會懇切地為我祈求.
但也好,你的病有需要的時候.
別人記得你,常常為你祈禱.
情況就是這樣.
保羅就是.
將羅馬教會的弟兄姊妹.
放在心上.
常常都是如此.
而且後面.
我們又明白多一點.
為甚麼他會這麼肉緊.
為甚麼他要講福音的內涵.
為甚麼要講使命的宣言.
原來.
保羅是沒有去過羅馬教會的.
羅馬教會不是他建立的.
所以他說.
希望總有一天.

$^{1081}$我能夠順利.
去到你們那裡.
當時沒有.
視像.
當時沒有錄影.
當時沒有直播.
如果你真的要.
跟弟兄姊妹一起.
想知道他們的狀況.
只有一個方法.
就是親身去.
保羅.
他很想見這班弟兄姊妹.
雖然聽到他們的信心.
傳遍天下.
但聽到.
他很想親身去.
可能見到他們.
打個招呼.
一起坐下喝一口茶.
吃個包,聊聊天,一起祈禱.
當你轉看這個.
和一張的PowerPoint.
不過Highlight的東西不同的時候.
你發覺原來保羅不單止對人.
對這班信徒很有心.
他對上帝.
都是很敬重.
他說.
這位神是怎樣的神.
就是我在他兒子的福音上.
用心靈服侍的那一位.
他可以給我作見證.
他原本的表達就是.
神就是我的見證.
你記不記得上面那句.
我的神.
然後他說.
神是我的見證.
見證我怎樣為羅馬的信徒.

$^{1121}$晝夜不住地禱告.
而且他說.
這位上帝.
我是很用心.
心靈其實意思是.
他用一個屬靈上的.
那種想法.
那種心態去服侍.
他就在上帝.
兒子的福音上.
他用心去服侍.
這位.
我這樣服侍的上帝.
我可以很大膽的說.
你可以為我作見證.
我怎樣愛這班人.
他在下面說.
他去不了.
然後他很謙卑的說.
或者按著神的旨意.
有一天我去得到.
他沒有說.
神為什麼不讓我去.
神為什麼這間教會不是我建立的.
沒有.
保羅說或許.
他也不敢說一定.
或許如果按神的旨意.
主若願意.
你可能會說.
為什麼這麼想去.
其實剛才也說了一點.
這裡他的表達是.
我很想見你.
我想和你分享一些東西.
我想帶一些手信給你.
那個手信是什麼.
屬靈的恩賜.
屬靈的恩賜這個字.
雖然我們會即時想.

$^{1161}$你有恩賜講道.
你有恩賜教導.
你有恩賜關顧.
不是這些.
因為這是一個單數的字.
因此這個字在羅馬書其他地方.
很多時候是和什麼相對呢.
很多時候是和.
過犯罪的報應相對.
因此這個字.
其實本身的意思是gift.
是禮物.
所以為什麼要因此.
正正不是你賺回來.
不是你抽獎抽回來.
不是你做得好而可以籠絡回來.
而是白白上帝.
送給你的禮物.
而這個禮物.
是和救恩.
和福音有關.
所以我估計大抵就是.
保羅沒有去過.
但聽到他們的信德.
也和他們講了福音的概要.
如果可以的話.
我很想親身來你們當中.
和你們開十天的培靈研經會.
還是和你們講一個.
貴道的羅馬書講道.
就是一個這樣面對面的接觸.
將屬靈的恩賜.
將屬靈的救恩和福音.
去和你們宣講.
目的是.
令到你們得到堅固.
你作為老師.
你作為教練.
你一定會想是這樣的.
不過我再醒多兩招給你.

$^{1201}$讓你再做得好一點.
讓你可以更上一層樓.
作為教睦當然是說.
你已經去到這個level.
你要有門份.
你要做領袖等等.
保羅的心態是一樣的.
甚至保羅在下面再補充.
我們在你中間的時候.
我們彼此那種信心.
我們彼此那種對信仰的認識.
那種的經歷.
我們可以彼此豐富.
一同去鼓勵對方.
保羅正正是很愛這班弟兄姊妹.
很緊張他們.
既很想快點去.
很想堅固他們.
但也不是一個高高在上的心態.
現在君臨城下了.
堂堂大名的保羅使徒.
現在要來羅馬教會.
不是.
他只是很想快點見到這班弟兄姊妹.
和他們分享.
屬靈的好處.
這是保羅的心腸.
所以保羅在後面就說.
其實我.
我告訴你們.
你們不要不知道.
我很多次都想去你們那裡.
不過去不到.
去不到.
你看到那個至今仍受到攔阻.
在原文是放在中間.
我多次地說到你們那裡去.
然後豁出去小聲地說.
不過到現在都不行.
那你說給誰攔阻呢?.

$^{1241}$是不是撒旦呢?他沒有說.
不過這個攔阻一定有神的容許.
因為剛才不是說了.
說許神的旨意.
如果是的話他就去.
但他為甚麼要去呢?.
咦,這次說多了一點.
不單止為了去牧養.
去兼顧,去開培靈會.
他說要在你們當中.
都要得一些果子.
屬靈的果子?.
然後說好像我在其餘.
外族人當中一樣.
保羅你想說甚麼呢?.
我見到時間,我要快點看完這兩節.
他接著就自爆了.
其實我剛才閃過幾次要說這個term.
他就說.
無論希臘人.
還是未開化的人.
簡單來說就是.
講希臘語,受過教育.
有文化水平的人.
所以相對那些就是未開化.
英文的翻譯就是barbarian.
野蠻人.
聰明人或者愚笨人.
大家就在想.
上一個季度講箴言那些.
馮永坤老師就說.
可能這個是一個比較個人的表達.
無論聰明.
還是愚蠢.
其實你知道保羅很多時候寫東西都列很多對比.
所以這個可能只是一個例子.
不一定是說兩類很特別的人.
但是無論如何都好.
不同的族群.
不同的情況.

$^{1281}$保羅說甚麼呢?.
我都欠了他們.
保羅欠他們錢嗎?.
街訴這麼厲害嗎?.
你知道嗎?.
這個就是跟著羅馬書說的.
福音的債.
所以他最後說.
我恩切期望.
將福音傳給你們.
在羅馬的人.
但在這裡就不只是傳給羅馬已經.
信主的弟兄姊妹.
在羅馬裡面還有很多人還沒信主.
在香港還有很多人還沒信主.
在北角區還有很多人還沒信主.
保羅說他很想.
很切切的期盼.
希望上帝願意.
讓我可以來到你的當中.
所以這個果子.
我相信是說福音的果子.
所以你看到保羅.
他既是牧羊弟兄姊妹.
也是鼓勵弟兄姊妹.
要去到當中去傳福音.
因為他的人生.
就是一個福音.
使命的人生.
傳福音,建立教會,栽培信徒.
訓練工人.
這個是保羅的生命.
所以最後.
我們見到保羅.
他尊敬上帝.
他愛他的福音對象.
甚至愛世人.
愛愛邦的族群.
這個是他的和祥.
他忠於上帝.

$^{1321}$忠於上主.
被他的託付.
我們一起祈禱.
如果當中在神的話裡.
你領受感動.
你就將這個感動和上帝表達.
甚至你可以呼求上帝.
你給我的使命是甚麼.
你告訴我讓我去實踐.
你自己禱告.
主耶穌基督.
我們的主耶穌基督.
當我們只是剛剛開始.
看羅馬書第一章的上半部分.
我們已經體會到.
保羅的一種.
神福音的心智.
對弟兄姊妹.
對信徒那種的愛.
關懷.
對他們靈命的緊張.
也都對上帝你自己.
那種的尊崇.
對你所交付給他的職份.
那種的執著.
那種的 Passion.
神下今天求你都幫助我們.
我們不是要和保羅.
做甚麼的比較.
既然我們都是同樣蒙照.
作使徒.
作歸屬耶穌基督的人.
我們就深信我們每一個.
都有理所要呼召我們.
去回應的福音使命.
或者一些其他方面的使命.
求你讓我們知道.
當我們知道的時候.
幫助我們.
有從你而來的力量.

$^{1361}$我們可以去實踐這個使命.
以致我們的人生.
真是一個實踐使命.
甚至實踐福音使命的人生.
求你這樣去幫助我們.
聽我們暗中有聲無聲的禱告.
奉耶穌基督的聖名祈求.
阿們.
\newpage



\section{約翰福音 4:1-18}
\label{sec:032090}
\textbf{主耶穌傳福音的智慧}~~~~駱柱成牧師
\newline
\newline
連結: \href{http://www.church.com.hk/acms/content.asp?site=cdc\&op=show\&type=product\&code=032090\&layout=sermon}{\texttt{032090}} ~~~~ 語音日期: 2026-05-03
\newline
\newline
\hyperref[sec:032089]{\small{< < < PREV SERMON < < <}}
~
\hyperref[sec:index_chronic]{\small{[返順時目]}}
~
\hyperref[sec:032091]{\small{> > > NEXT SERMON > > >}}
\newline
\newline
約翰福音 4:1-18
\newline
\begin{longtable}{cl}
\hline
\hline
章節 & 經文 (和合本修訂版)\\
\hline
4:1 & \begin{tabularx}{0.7\textwidth}{X} 耶穌知道法利賽人聽見他收門徒和施洗比約翰還多， （ \end{tabularx} \\ \\ \relax
4:2 & \begin{tabularx}{0.7\textwidth}{X} 其實不是耶穌親自施洗，而是他的門徒施洗，） \end{tabularx} \\ \\ \relax
4:3 & \begin{tabularx}{0.7\textwidth}{X} 他就離開猶太，又回加利利去。 \end{tabularx} \\ \\ \relax
4:4 & \begin{tabularx}{0.7\textwidth}{X} 他必須經過撒瑪利亞， \end{tabularx} \\ \\ \relax
4:5 & \begin{tabularx}{0.7\textwidth}{X} 於是到了撒瑪利亞的一座城，名叫敘加，靠近雅各給他兒子約瑟的那塊地。 \end{tabularx} \\ \\ \relax
4:6 & \begin{tabularx}{0.7\textwidth}{X} 雅各井就在那裡；耶穌因旅途疲乏，坐在井旁。那時約是正午。 \end{tabularx} \\ \\ \relax
4:7 & \begin{tabularx}{0.7\textwidth}{X} 有一個撒瑪利亞婦人來打水。耶穌對她說：「請給我水喝。」 \end{tabularx} \\ \\ \relax
4:8 & \begin{tabularx}{0.7\textwidth}{X} 因為那時門徒進城買食物去了。 \end{tabularx} \\ \\ \relax
4:9 & \begin{tabularx}{0.7\textwidth}{X} 撒瑪利亞婦人對他說：「你是猶太人，怎麼向我一個撒瑪利亞女人要水喝呢？」因為猶太人和撒瑪利亞人沒有來往。 \end{tabularx} \\ \\ \relax
4:10 & \begin{tabularx}{0.7\textwidth}{X} 耶穌回答她說：「你若知道神的恩賜，和對你說『請給我水喝』的是誰，你早就會求他，他也早就會給了你活水。」 \end{tabularx} \\ \\ \relax
4:11 & \begin{tabularx}{0.7\textwidth}{X} 婦人對耶穌說：「先生，你沒有打水的器具，井又深，哪裡去取活水呢？ \end{tabularx} \\ \\ \relax
4:12 & \begin{tabularx}{0.7\textwidth}{X} 我們的祖宗雅各把這井留給我們，他自己和兒女以及牲畜都喝這井裡的水，難道你比他還大嗎？」 \end{tabularx} \\ \\ \relax
4:13 & \begin{tabularx}{0.7\textwidth}{X} 耶穌回答，對她說：「凡喝這水的，還要再渴； \end{tabularx} \\ \\ \relax
4:14 & \begin{tabularx}{0.7\textwidth}{X} 誰喝我所賜的水，就永遠不渴。我所賜的水要在他裡面成為泉源，直湧到永生。」 \end{tabularx} \\ \\ \relax
4:15 & \begin{tabularx}{0.7\textwidth}{X} 婦人對他說：「先生，請把這水賜給我，使我不渴，也不用到這裡來打水。」 \end{tabularx} \\ \\ \relax
4:16 & \begin{tabularx}{0.7\textwidth}{X} 耶穌對她說：「你去，叫你的丈夫，再到這裡來。」 \end{tabularx} \\ \\ \relax
4:17 & \begin{tabularx}{0.7\textwidth}{X} 婦人回答，對耶穌說：「我沒有丈夫。」耶穌說：「你說沒有丈夫是對的。 \end{tabularx} \\ \\ \relax
4:18 & \begin{tabularx}{0.7\textwidth}{X} 你已經有過五個丈夫，你現在有的並不是你的丈夫。你這話是真的。」 \end{tabularx} \\ \\
[1ex]
\hline
\hline
\end{longtable}
$^{1}$弟兄姊妹 住內平安 早安.
今天我是講傳福音的 題目是.
我們Elly姐有傳福音的活動 昨天.
盼望主多多地將人帶到我們月月堂.
我們平時都喜歡聊天 聊天就是聊聊天.
談話 通常我們會約人聊天.
喝一口茶 坐下來慢慢聊.
聊天聊天 其實聊天是什麼意思呢.
我們普通話就說談話.
聊天我們語語真是很豐富的意思.
聊天原來的意思是先清喉嚨.
後來就引申為我們交談.
清喉嚨和交談 即是我要先做好準備.
和你去有這個對話.
所以聊天是和我們的交談.
和我們平時有空喝一口茶聊聊天.
這個聊天是很悠閒的.
這個字其實和佛教的字是有關的.
他就說聊天就是智慧的話.
他不是單是一般的閒聊.
你聊的時候就帶出那個智慧.
所以聊就是吐出來.
我們倒出來.
將心裡那些智慧想表達的.
講出來在一個很輕鬆的環境當中.
如果開會呢 開會就麻煩了 很嚴肅.
大家都很緊張 因為很正式.
聊天就不同了 聊天就很容易建立關係.
世人就會聊天 帶出那個世界的道理.
你在當中聊的時候 你就知道很多的訊息.
或者甚至是學到很多世界的道理.
基督徒呢 我們聊天 聊的時候.
如果是和講福音有關的.
我們在這個過程就要有那種智慧.
不單止我們講的話 裡面有真理.
我們講的過程就需要有智慧.
因為你有智慧的時候.
你的關係就可以建立得到.
傳福音就比較容易入手.
主耶穌就是幫我們做了一個很好的榜樣.

$^{41}$在這段經文中.
經文就說到.
約翰是傳福音.
天國近了 你們應當悔改.
耶穌基督呢 祂同樣是傳福音.
大家都是做福音工作的.
但是當時約翰的門徒就說了一句話.
意思就是大家都去了耶穌那裡.
我們這裡人越來越少了.
所以就在試探 在挑起爭端.
這句話其實很容易.
所以如果處理不好 約翰處理不好.
他就和耶穌有爭執.
所以很容易同宮就會有分裂.
所以如果這兩個最突出的領袖.
一個是施洗約翰 一個是耶穌.
有矛盾 在這個關鍵的時刻.
不知道怎麼走下去.
是沒有人預料得到的.
因為這兩個人太重要了.
我記得以前看電視聽過一個故事.
康熙年間有個做官的 叫張瑛.
他的家人住在他出生的那裡.
他做官 在京城.
誰知道他想收屋.
收屋 但是就從隔離屋發生爭執.
他家人就寫信去京城.
就告訴張瑛 叫他出面動用人力.
但是張瑛沒有這樣做.
他就寫了一封很短的信回來.
千里送書子為長 讓他三隻又何妨.
萬里長城金猶在 不見當年秦始皇.
我們不需要爭的.
所以他家人就讓出三尺.
隔離屋的人就讓出三尺.
又讓出三尺 就變成六尺.
所以有一條港叫做六尺港.
就是這樣來的 當時康熙讚張瑛做得好.
所以因為他處理這些關係的時候.
做得特別好.

$^{81}$約翰就沒有將分裂擴大.
他沒有帶給人更大的隔閡.
他的態度是這樣 他不爭.
人若不是從天上得的 就不能得著什麼.
一切都是從神所賜的.
他也不搶 我又不是基督.
他很清楚自己的位置.
他也不至於 當人跟了耶穌.
他不覺得失落 為什麼呢.
因為他本來就是做僕人的.
主興旺了 我就高興了.
人信主 福音就興旺了.
這個才是最重要的.
至於他去哪間教會這班人.
以後這些人的去向 主帶領他們.
我們今天教會 傳福音是最重要的.
去哪間教會他們 去哪個堂 神帶領他們.
耶穌做僕 他也不爭 也不宣揚.
因為我們在其他福音書會看到.
耶穌基督是很安靜去做僕的.
他不爭 他也不跟約翰爭.
所以很多人以為.
耶穌比約翰施洗的人更多.
聖經告訴我們 他的門徒在這裡施洗.
約翰是用水施洗 做預備的.
引人到耶穌基督.
他施洗有一個儒表.
耶穌基督是真正的施洗者.
他就是用聖靈去讓人施洗.
讓人去重生 使人得到生命.
這件事只有耶穌基督才可以做.
約翰不可以做 其他人也不可以做.
約翰也很清楚.
他只不過是一個預備人心 引向耶穌基督的那一位.
耶穌基督才是真正的拯救者.
所以人真的不可以做赦罪的工作.
只有耶穌基督可以賜給人生命.
可以使人重生 可以赦罪.
所以約翰的工作是不可以跟耶穌基督相比的.
當約翰很興旺的時候.

$^{121}$很多人以為這個就是最重要的人物.
神就在這裡干預.
耶穌基督才是最重要的.
約翰不可以在這個歷史這麼重要的時期.
跟耶穌基督相比.
否則有人就以為約翰是最重要.
但是神就在這裡干預.
耶穌基督才是獨一的救主.
不容許其他人跟神相比.
所以很多人去到耶穌基督那裡.
當那麼多人去到耶穌基督那裡.
問題就來了.
法利塞人知道這件事.
就覺得有危機感.
就開始關注耶穌基督.
而且那種衝突.
本來你結證聖殿.
本來你做公國的時候已經是很多衝突.
現在那麼多人去到你那裡.
那種衝突就更加激烈了.
槍打出頭鳥.
所以耶穌基督就是準備被對付的那種.
但是耶穌基督在這個時候.
他就離開猶大.
離開最敏感的地方.
回到加里尼.
耶穌基督這個退讓,退開,離開和避開.
是不是表現他很軟弱的那個位置呢?.
不是的.
因為有一個原因就是事後未到.
聖吉卜人不是跟著情況去行事的.
是跟著神的時間表.
不是說有益處就退後.
有危險的時候.
應該是有益處的時候.
不是向前.
有危險的時候反而要站出來.
這個是聖吉卜人應該有的態度.
因為如果按我的說法.
很容易縮沙有危險的時候.

$^{161}$但耶穌不是這樣.
我們處事很多時候是這樣的.
無事不生事.
大事化小事.
小事化無事.
如是呢.
如果前面的三句還不夠.
如是的時候你就不要怕事.
耶穌基督是不怕事的.
我記得在國內的時候.
當教會興旺的時候.
很多人過來聚會.
宗教部門也開始關注.
那個老僕人.
平時是什麼都背了.
一有公安局找他聊天,喝茶.
就是他去的.
自己什麼都背了.
下面的同工他們壓力沒那麼大.
但是這位僕人已經回天家了.
那個時候.
那些新起來的領袖.
就要承擔那個責任.
好像我當時.
我還很年輕.
比現在更年輕.
是不夠資格的.
通常去背票的是多少歲呢.
四十多歲以上.
他們就說.
我來.
不是說他沒有家人.
沒有責任.
而是如是不怕事.
他願意去背票.
因為這些情況是要承擔的.
老僕人離開了.
有個追思會.
批文下來了.
你們這個教會.

$^{201}$要禁止在那天去辦追思會.
教會就信服了.
沒有抗衡.
因為大事化小事.
小事化無事.
但是很奇妙.
神的工作是沒有人難做得到的.
那天就很多人來了.
我那個時候在唱詩.
我一唱就唱了幾個小時.
我知道大概.
所以是過萬人.
幾萬人我估計是.
過來參加追思會.
所以神有祂的工作和計劃.
和祂的時間.
所以事後未到.
所以不用衝上去.
你衝上去的時候.
這個是蠢.
事後到了.
你該去的你不去.
這個就是膽怯.
所以進退不是.
關鍵不是我們想做什麼.
而是按著神的時間表.
時間未到好好預備.
耶穌基督一直預備.
預備了三十年.
三十三年半至真正上十字架.
信服神的時間.
這個不是退避.
不是逃跑.
耶穌基督離開猶大.
要往加尼尼去.
路線是從猶大的南方.
繞約旦河那條水路去.
要去到加尼尼.
中間有一個地方.
叫做撒瑪利亞.

$^{241}$撒瑪利亞看見粉紅色那塊就知道了.
黃色那塊就是猶大.
要繞過隔離那塊.
隔離那塊.
才上去加尼尼.
他們不去踏足猶太地.
以色列原本有十二個支派.
其中一個支派是猶大支派.
猶大支派他們很蒙福.
因為君王是他們的支派出來的.
所以猶大支派就有那種自尊和優越感.
真正的王在我們中間出來.
所以覺得很威風.
南國猶大是被巴比倫擄去的.
之前北國.
因為我們知道以色列分為南北國.
北國是被亞述擄去的.
北國被擄去之後.
他們沒有嚴格遵守那個律法.
在婚姻上的事他們就跟那些外族通婚了.
他們沒有遵守律法的要求.
南國被擄之後他們是堅守律法的.
所以他們的血統相對是比較純正的.
猶太人認為.
這些撒瑪利亞人是不夠純正的.
因為他們已經跟外族通婚了.
所以他們的血統是不純正的.
如果你看尼希米基就知道了.
他們的血統是不夠純正的.
你看尼希米基對付那些外族通婚的.
他是要撿出來去對付的.
猶太人認為我只是純正的以色列人.
我只是阿伯拉海真正的後裔.
所以他們是很自豪的.
所以他們就不會看得起撒瑪利亞人.
因為撒瑪利亞人跟外族通婚.
你們已經被污染了.
不守律法了.
如果你再看公元前八世紀.
亞述攻陷以色列之後.

$^{281}$這班以色列人就被遷出去了.
那些外族人就在那裡定居了.
那些外族人,外邦人就在那裡拜假神.
所以這塊撒瑪利亞地的人.
跟宗教有種混合.
所以當猶太人看撒瑪利亞人的時候.
他們就看到你們這班人.
跟這些拜偶像有混合在一起的時候.
是對神的一個冒犯.
所以你們的信仰已經是混合的一種信仰.
不夠順了.
所以猶太人很憎恨這些撒瑪利亞人.
拉比當中有一句話是這麼說的.
吃撒瑪利亞人的餅就等於吃豬肉.
我們知道猶太人因為豬肉是很骯髒的.
所以他們的問題不單止是文化的差異.
是信仰,血統和他們身份的問題.
猶太人已經否定撒瑪利亞人了.
所以他們要求耶穌繞路回加利利.
但經文很特別.
他說他必須經過撒瑪利亞.
這個必須不是地理上的必須.
因為你可以繞過去.
是神旨意的必須.
當耶穌這樣做的時候.
他就打破了傳統的界限.
民族的界限.
就是猶太人和撒瑪利亞人是沒有來往的.
現在耶穌主動進入這塊地.
打破了他們民族的界限.
還有打破了他們宗教的界限.
因為猶太人和撒瑪利亞人.
他們的信仰已經不是那麼純了.
很可能已經是異端了.
耶穌基督向他們傳福音.
所以打破了宗教的界限.
還有打破了人心的界限.
因為猶太人看撒瑪利亞人.
就是看你不順眼.
我就要指著你去罵.

$^{321}$所以耶穌主動去找他們的時候.
打破了心靈的界限.
所以我們今天看.
人會拒絕耶穌基督.
但是耶穌基督願意跨越那種界限.
跨越人的隔膜去找人.
真正的福音工作.
不是我們想去的地方.
而是要順服神的帶領.
必須去的地方.
所以這個必須不是地理上的必須.
是那種使命.
神既然有命令給你.
你就要出發.
所以這件事才是最重要的.
去撒瑪利亞那裡.
是有一個婦人.
一個不是那麼得體的婦人.
但是耶穌基督看重的是靈魂.
不是別人怎麼看這個人那種眼光去看.
我記得我在機場工作的時候.
跟同事一起.
我們是國內的.
有時候就要去支援國際檢查.
我們都是做安檢的.
他們是國際的.
我們是國內的.
我們就要支援他們.
人太多了.
國際區那裡.
有一個開車的同事.
開電車在機場主機大樓開電車.
他每次都走到登機辦理處那裡.
去繞個圈.
我就不明白為什麼他經常都繞個圈.
本來都不需要這樣做的.
後來隔壁那個就說.
那個啊那個啊.
那個是誰啊.
不是他未來太太.

$^{361}$所以他是故意去那裡繞個圈.
假裝一下他.
因為愛情的力量.
那耶穌基督呢.
耶穌基督的那種愛是性愛.
他是為了救人那種愛.
他不是看這個人不潔淨.
他不是看別人眼光怎麼看他.
他是看看中這個人的靈魂.
所以我們很多肢體.
如果你知道的.
很多是他故意去接人.
或者是接一些舞蹈者.
本來都不用走那條路的.
故意過去接他.
就是看中他的靈魂.
不是我們順便就去.
順路就去.
而是當我們敬畏神的時候.
看到神的心意.
看中神的心意的時候.
有這個負擔的時候.
對靈魂有愛的時候.
就願意這樣去做.
所以耶穌來到撒瑪利亞.
不是為了方便.
是順服神的計劃.
那耶穌基督來到一座城.
名叫罪加.
這座城本來都沒有什麼起眼的地方.
似乎是微不足道的地方.
耶穌就來到這裡.
因為這裡有一個女性.
一個撒瑪利亞的婦人.
按當時來說.
男女如果談話.
不是像我們現在這樣隨便.
如果在一個公眾地方去談話.
是很不合宜的.
是很不得體的.

$^{401}$他們猶太人的宗教人士.
他們的禱告是怎樣呢.
不要讓我生為撒瑪利亞人.
所以你會看到.
這些禱告是看到他們.
不喜歡撒瑪利亞人的.
甚至他們會禱告.
做我不是女人.
所以他們也有男女的差別.
不是我們現在這樣平等.
他們是有差別的.
所以如果一個拉比.
你做老師的.
跟女性在公開的場合去談話的話.
可能是你的妻子也好.
你的女兒也好.
是不太合體統的.
所以耶穌基督.
來找這個女性的時候.
他已經是跨越了性別的界限.
所以你會看到.
耶穌基督.
他沒有因為人類的傳統.
而難了去看待這個人.
去拯救這個人.
雖然這個人生活很混亂.
也被人看不起.
甚至是被社會所排斥.
但是耶穌基督.
就是因為看中人的靈魂的寶貴.
去尋找他.
當耶穌基督去找他的時候.
那些門徒就被耶穌基督叫出去了.
你去買餅買東西吃.
如果這班門徒在的時候.
問題就多了.
他們的對話就不是那麼容易.
他們可能在聊天的時候.
就會用不是那麼好的眼光去看待這個婦女.
甚至插口的時候.

$^{441}$這個婦女十三對一.
你怎麼爭辯呢.
如果是聊天去辯論的話.
是不能爭辯的.
所以這是一個干擾.
耶穌基督需要一個沒有干擾的.
可以真實談話的地方.
那個空間.
要有一個空間給他們去談.
耶穌基督很有智慧處理事情.
他去尋找尼哥底姆的時候.
尼哥底姆去尋找他的時候.
他允許他在夜晚.
在私下.
避免太多法利塞人對他的指責.
對撒瑪利亞婦人.
他是在白天.
光天化日之下.
是公開的.
突破那種.
而且他是主動去尋找的.
所以你會看到耶穌基督.
他傳福音的時候.
真的很有智慧.
他不會讓你拿著把柄去指責你.
你夜晚去尋找婦女.
這樣就不對了.
但白天呢.
白天就可以避免那種嫌疑.
所以耶穌基督每一個行動.
他的時間,地點,對象,方式.
都是有一個很有智慧的選擇.
所以傳福音不是隨便的.
你需要有從神來的那種智慧.
否則就會引起很多麻煩.
那個時候就是五政.
一般婦女都是選擇那種.
沒有那麼熱的時間去打水.
但在這裡呢.
這個婦女很特別.

$^{481}$她是在人少的情況才出來的.
因為她自己的.
我們讀一下哪個經文你會知道.
她自己的背景不是那麼好.
很容易被人指責和輕看.
所以她不願意去見人.
避免這些人去指責她.
所以她是一個避世的人.
不敢見人.
所以在這個時刻.
耶穌基督就在那裡了.
這個不是偶然.
是神祂一個很精心的安排.
耶穌基督不是照那些異人.
是照罪人悔改的.
所以我們要知道.
我們看人很多時候.
這個人不是那麼好.
但在神眼光不一定是這樣.
神的心很大.
很可能那些不潔淨的.
可能我們現在看不是那麼好的.
以後成為神重用的器皿.
神的目光是很遠大的.
我們知道莫大拉的瑪利亞.
她是那種邊緣化的人物.
因為她是被鬼附的.
但是她信主之後.
她是打一個為主做見證的.
主復活之後.
基里拉合她是外邦人.
也是被神拯救的.
而且她的名字寫在了.
梅塞亞的家譜裡.
神在不可能的當中.
祂會成就一些很奇妙的事.
很偉大的事.
因為神不願意放棄任何一個人.
祂是願意帶人回到神那裡.
耶穌一開口就說.

$^{521}$請將水給我.
但是這裡.
如果你看背景的時候.
耶穌基督沒有直接揭露這個人的背景.
祂有五個老公.
可能是五個半老公.
但是耶穌基督沒有揭露祂很多醜事.
祂也沒有說過一些刺激人的話.
太陽這麼曬.
你現在才出來.
你生活也挺特別的.
如果這樣說的時候.
這個女人的戒備心就很容易建立起來.
但是耶穌基督.
祂是要打開人的心門.
不是要關閉人的心門.
耶穌基督知道人的問題.
但是祂沒有一下子就指出來了.
祂是要建立那種關係.
建立了關係.
祂才帶出那個真理出來.
我們知道我們當中有不少弟兄姊妹.
他是有能力去幫人的.
但是為什麼你幫不到人呢.
很多時候就是因為我和他的關係還不夠好.
他還沒有真正信得過我.
如果他真正信得過我的話.
他會找我.
但是我們向人傳福音的時候.
如果我們的關係還沒有建立好的時候.
我們是沒有那個能力和機會去幫到人的.
不是因為你特別有知識.
人家就找你.
而是那種關係.
那個心有沒有真的拉近了.
我們知道很多問題.
不是因為你的問題而出來的.
很多時候就是那個關係不好而帶出來的.
中國人說一個詞很好說.
就是人事.

$^{561}$你是先有人再有事.
人放在前面.
事放在後面.
先有那種關係.
你才有那個影響力.
如果你關係不好的時候.
你很容易.
本來沒有問題的.
你都很容易變成有問題.
所以耶穌基督是要幫助這個人的.
這個人是被排斥的.
耶穌基督沒有給他機會去懷疑耶穌.
所以他就說了一句很合理的話.
請你給我水喝.
給我水喝.
什麼意思呢.
如果按這位婦女的身份來說.
她是一個在人前可以說是沒什麼面子的.
但是耶穌基督說.
你給我水喝的時候.
他就其實在幫助這個人.
耶穌基督很有智慧.
當你一個人的地位很低微.
看自己很自卑的時候.
他就很容易失去他的價值.
很容易.
但是耶穌基督對他的關心的時候.
他的價值觀就很容易建立起來.
耶穌基督打破了那種傳統的那種規條.
雖然很多的限制.
可以說你男女.
向耶穌基督男.
向這個婦女拿水是不合理的.
但是如果是在烈日之下.
走路困倦之下.
我又沒有拿水的工具之下.
向一個人去求助.
是很合理的.
所以耶穌基督沒有讓他去.
關閉這些條條框框.

$^{601}$這些宗教.
這些身份去關閉了洞門.
耶穌基督說請你把水給我喝.
很合理因為太熱了.
而且這句話有另外一個訊息.
你是一個可以幫助人的人.
你不是一個很低下的人.
沒有價值的人.
你不是一個被遺棄的人.
你現在可以幫助我.
你是有價值的.
我是沒有水具去打水的.
但是你可以.
所以當這樣說的時候.
這句話出來的時候.
耶穌表面來看是向他拿水.
但是這個婦女意識到.
原來我是不是沒用的人.
我是可以幫助人的人.
所以耶穌一開始.
祂不是指出人的問題是什麼.
祂是喚醒人的價值是什麼.
你是有價值的.
所以這句話真正表達的就是.
我尊重你.
我看待你是一個人.
可能其他人看你.
你是很低下的.
你可以說你是很低賤的.
你是不被人尊重的.
但是我現在尊重你.
所以這個婦女.
當被人尊重的時候.
她的心就打開了.
當有人幫助其他人的時候.
她的價值就出來了.
請你給我水喝.
你是可以付出的.
所以耶穌基督這句話.
其實是很觸動婦女的心.

$^{641}$祂是幫助這個人.
走出這個困惑當中.
這個處境當中.
祂是給人一個新的盼望.
你原先是沒有價值的.
你讓他變成一個有價值.
有盼望的人.
是非常有智慧的.
所以耶穌基督不是一開始就揭露人的罪.
也不是冷淡地對待這個人.
祂是要打開人的心.
我們很多時候.
我們傳福音的時候.
我們很愛我們的朋友.
我們很想他脫離那種罪.
回到神的面前.
但是我們很多時候太急.
急中無智.
人家急中生智.
很多時候急就沒有那種智慧.
我記得我小時候.
小時候只有幾歲.
很喜歡走出去玩.
是一個人走出去.
和其他小朋友玩.
我哥哥大我八歲.
他很多時候叫我回家.
你還玩.
你都不知道是時候了.
他想叫我回來.
但是他就是.
這句話我聽.
就是他在罵我.
他在責備我.
但是我媽媽叫我回去.
就不一樣了.
她說回來了.
家裡有東西吃了.
你快點回來幫我數筷子.
我數筷子分給其他家人.

$^{681}$我就快點回去.
為什麼?.
因為我有價值.
我就快點回去.
我感受到原來這個家是需要我的.
我是感受到媽媽是很愛我的.
如果不是我哥哥抓我回去.
你以為我真的被嚇回去.
我是被家人的愛吸引回去.
所以耶穌基督.
祂是沒有一下子就將人推出去.
神愛我們.
祂愛那些不聽話的小孩.
我們原先都是不配的.
都是不聽話的.
但是神祂是很有道量.
祂可以赦免我們的罪.
所以我們傳福音是需要調整的.
不是用審判人的態度.
去向人傳福音.
我們要尊重人性.
人性是很重要的.
你要知道它是有價值的.
它是在神的恩典當中有份的.
這件事是很重要的.
神愛他.
差在我們來到他的面前.
傳講這個福音.
讓他知道他是有價值.
被愛的.
當人知道這一點的時候.
他才有可能迴轉的機會和意向.
我們才有可能將人帶到神的面前.
另外一方面.
你看婦女的回應就知道了.
現在你尊重了她.
她本來是很低下的.
被人輕看的.
現在她一下子就和耶穌基督平起平坐了.
耶穌基督很尊重她.

$^{721}$但是她另一個反應就出來了.
你尊重我.
我現在和你一樣.
所以她的思維.
又是另一個思維.
我現在是和你平起平坐的.
你不了解我.
也不知道我原先的背景是怎樣一個人.
你是猶太人.
我是撒瑪利亞人.
本來我們之間是有爭執的.
你現在向我下水.
你是什麼想法.
你都不知仇的.
我們本來是敵對的.
你應該是不向我下水的.
你向我下水.
你到底是什麼意思.
所以她就帶出另外一件事.
我現在比你更好.
因為你不知道我的背景是怎樣的.
我現在是可以給你水的.
你現在是向我下水的.
你是撒瑪利亞人.
我們本來就是有敵對的.
你向我下水.
你是男人.
我是女人.
你這種動機是不好的.
所以她反而輕看了耶穌基督.
所以她又將那種.
當她平等的時候.
她又拉回到民族的鬥爭那裡.
那種隔閡那裡.
又下回原先那種地位.
民族敵對的地位.
人通常都是這樣.
人的本性就是.
你當不知道我犯過什麼罪的時候.
我就和你平等了.

$^{761}$我就和你是一樣的.
你原先那麼尊重我.
我現在被你尊重.
我原先是被人看不起的.
我已經忘記了我原先的背景.
我原先的思維又不一樣了.
所以很容易.
人的本性就很容易回到身份.
背景和種族的對立那裡去.
我們肉體很容易受地域和地區文化的不同.
而對人有不同的看法.
但是在靈裡面.
其實我們基督徒來看是沒有分別的.
因為人的靈魂.
不是按你的口音和地區.
去將人看低或者將自己看高的.
耶穌基督來降低了自己.
沒有用那種高高的姿態去向人傳福音.
如果一開口就很高調的話.
這個傳福音的過程就失敗了.
因為如果你是很高傲地指責人的時候.
別人很可能會說我不需要你的福音.
你可以說你沒有什麼優越感.
你是從天上來的.
你是天上的君王.
你就回到天上去吧.
我不需要你.
因為你太驕傲了.
所以和人傳福音那種態度是很重要的.
我記得有一個故事叫做.
借泥之息.
有一個人是在齊國的,他很有錢.
叫做禽奧.
他分食物給人的時候.
他說你過來吃吧.
你過來吃吧.
那個時候那些人是寧願不吃的.
我寧願餓死都不吃.
為什麼?.
因為你將我看低了.

$^{801}$我是一個人,我是有尊嚴的.
你這樣施捨給我,我是不接受的.
所以我們今天可能在我們教會有不同的地方.
我在廣州的時候也是有不同的地方.
山東,山西,湖北,湖南.
那個打扮不一樣.
衣著不一樣.
背景文化也不一樣.
在社會當中.
有些是在大城市的.
有些是在鄉下的.
但是一般按社會來說.
你是在鄉下來的.
我是在大城市來的.
我比你高一級.
如果在教會當中有這種聲音.
是不可以的.
因為人的分別不是關乎口音.
不是你的口音高一點.
你的調高一點.
你的靈魂就高一點.
也不是因為你的地方富裕一點.
你的生命就豐富一點.
在神的面前,人人是平等的.
所以無論你是什麼時候來到教會.
都是一視同仁的.
除非你第一次來,第二次來.
我們出於禮貌.
好好地去幫助你,引導你.
為你安排好位置.
在我們現在.
這裡有很多不同的弟兄姊妹.
因為洛杉磯.
就是匯聚了世界各地的人口.
我們聖谷華人宣導會儒堂.
每次來,有新人來.
我們都是主動去歡迎他.
問候他和陪伴他.
而且每個主日都是這樣做.
這個就是因為.

$^{841}$我們不是因為那種地區地域.
而將人分開.
而是因為每一個人在神的面前.
都是平等的.
我們不可以將人當作另類.
不是因為你講粵語講得好.
我就尊重你.
或者你的口音不好,我就不尊重你.
而是感到.
我記得有個姐妹.
她講粵語不是那麼流暢.
但是她說.
我來到這裡.
你們當中的長者對我很好.
我是特別感到關懷和尊重的.
所以我留下在這個教會.
所以一個人能夠留下.
是因為他感受到那個愛在這裡.
反過來.
我記得有一次去吃飯.
和一個朋友.
他說他只會講國語.
我們去到粵語餐館.
他去到那裡的時候.
店員會講粵語.
國語.
英語.
所以他點餐的時候.
就被對待.
就不覺得那麼受尊重.
所以我吃完之後.
都很好吃.
下次我們再來.
好吃是好吃.
但是以後都不來了.
為什麼.
他覺得未受尊重.
所以一個人不受尊重的時候.
他很容易離開那裡.
很可能這個人就是.

$^{881}$弟兄姊妹花了很多時間.
在意將他叫回來.
所以很可能.
因為我們一句不合宜的說話.
就將人搬走了.
很可能他就永遠都不來了.
所以真正需要的.
不是我們有什麼學問.
而是說我們真正需要的是.
我們要真正去愛人.
去尊重人.
不要傷害人.
如果不是的時候.
就很容易難了福音的工作.
我們要接納不同的人.
無論是廣東人也好.
廣西人也好.
因為保羅講過.
向什麼人就做什麼.
廣東人當中就做廣東人.
在外邦人當中也要做外邦人.
因為神愛世人.
不是分那種背景,語言,文化.
神是要我們在各族各方.
各國的人中間成為一個橋樑.
將人帶到神的面前去.
你既是猶太人.
為什麼向我一個撒瑪利亞人要水.
耶穌基督對他好.
他沒有對耶穌基督好.
耶穌基督沒有因此而退後.
你對我不好.
但是我不會對你不好.
因為耶穌基督要的是.
要用智慧和忍耐.
將引導這個婦人.
他要從井水.
讓這個婦人看到活水.
從她的生活當中.
看到她生活的需要.

$^{921}$看到她靈魂的需要.
最後才認識耶穌基督.
不是因為你對我好一點.
我就對你好一點.
你對我不好.
我就對你不好.
可能在戀愛當中的情侶是這樣的.
心計,攻心計.
但是當我們一種戰術的時候.
是不合聖經的.
因為我們對待一個人好不好.
是因為我的品格.
我的教育.
我的背景.
這些是有關的.
但是我怎樣去對待一個人.
不是關乎他怎樣對待我.
原則是.
你要怎樣待人待你.
你是要先怎樣待人.
我們要先怎樣去待人.
我們是屬主的人.
我們要先怎樣去待人.
不是因為他那種不好.
而我對他不好.
所以耶穌基督.
雖然知道這個婦人好不好.
知道她一切的過去.
但是在這個時候.
她沒有追究她的責任.
也沒有執行審判.
耶穌基督來.
祂不是去審判人的.
祂是要拯救人的.
祂降卑到.
到聖肉身來到我們中間.
是要了解.
祂了解人那種軟弱.
也提出人那種軟弱.
所以祂用了一種很自然的方式.

$^{961}$請你將水給我喝.
這個婦人提出一個問題.
文化不同 宗教不同 民族不同.
你為什麼要向我拿水.
耶穌基督沒有因為.
她這些質疑和對抗.
而反對她.
雖然祂說.
你向我男人 向我女人拿水是不對的.
猶太人看不起撒瑪利亞人.
雖然祂提出很多問題.
但耶穌基督沒有因為她這些問題.
跟著她走.
沒有因為她這些問題.
被她帶著走.
祂要將人帶回到.
祂要講的核心.
福音的核心那裡.
你如果知道神的恩賜.
又知道對你說.
請你給我水喝.
這個是誰.
你就早已經求祂了.
祂必給你活水.
所以耶穌基督再次將這些婦人.
兜兜轉轉.
在民族對抗當中.
又帶回到神的恩典和神那裡.
所以耶穌基督傳福音的時候.
不是跟著人的思路走.
我們傳福音.
很多時候是跟著人的思路走.
祂沒有.
祂要將人這些猜疑 質疑 對抗.
帶回到神的恩典那裡.
帶回到福音的核心那裡.
因為每一個人.
當我們去傳福音的時候.
他很可能都是很擅長去對抗的.
為什麼?.

$^{1001}$因為這個是撒旦的一個.
對抗上帝一個計劃的一部分.
人都落在罪當中.
都受撒旦的克制.
所以我們傳福音的時候.
如果你是中國文化的你就知道.
我以前做生意.
我做慈善的.
所以我對得住天地良心.
好人有好報.
不需要其他東西.
所以你覺得好人有好報.
很合理.
你說得對.
所以你跟著祂思路的時候.
很多時候話題就結束了.
你沒有辦法再去深入福音那裡.
因為祂已經將那樣東西停止了.
耶穌基督沒有.
耶穌基督沒有放棄祂.
也沒有順著祂那種民族和宗教的衝突.
去思路去走.
祂是將人那種焦點.
帶回到神的恩典那裡去.
一切都是神的恩典.
這口井是恩典.
你能夠打水也是恩典.
你能夠給人供應也是一個恩典.
你以為你在施恩.
其實你自己都還在這個恩典當中.
耶穌基督將人帶回到恩典那裡去.
接著進一步說.
你若知道向你要水的是誰.
你早已經求祂了.
祂要給你的是活水.
這句話就很有張力了.
為什麼?.
當他聽到一個活水的時候.
就擊中了他生命的核心.
你給我一個活水.

$^{1041}$我現在給你.
就要給你一個死水.
他每天都來打水.
每天都頸開.
每天都還要來.
而且每天都來.
重重複複都要來這種狀況.
也是他生命的一種狀況.
他就是一個追求慾望.
不斷地追求.
但一直在虛空當中的一個人.
他每天都在打水.
每天都要避開閒言閒語.
那些婦女天天都在論斷我.
在冷言冷語.
看不起我.
所以這是他生活的寫照.
他每天都要做這個重複的動作.
我本來是要靠臉吃飯的.
太陽這麼曬.
我都要出來打水.
為什麼?.
因為我就是需要這個水.
我要需要這個水.
但是我的臉一出來曬的時候.
很容易有斑.
我是靠臉吃飯的.
如果我靠這個臉吃飯的時候.
我以後還怎麼靠呢?.
但是我願意冒著這個險.
來到這裡當中.
冒著這個險.
讓別人看不起我都好.
我都還要來.
我每天都在打水.
罪惡帶給他很快樂.
他背後的事帶給他很快樂.
但是他面對這些罪的後果.
人的論斷.
他是很痛苦的.

$^{1081}$每天都要來打.
不是像我們現在這樣.
有自來水.
有快遞.
他就是每天都來.
人人都看我的眼光.
都在交仇接義.
我就是在這個困境當中.
很渴望一個不用打水.
或者一個令我心靈真正滿足的東西.
我不想被人笑.
我不想活在這種羞辱的壓力當中.
但是我就是為了這個水.
每天都要來.
冒著別人的指責.
我都要來.
所以當耶穌說活水的時候.
擊中了他的一個要害.
死水是停滯的.
是會變壞的.
活水是流動的.
帶給人生命的.
如果我們看舊約就知道了.
講到活水的時候.
耶利米蘇講到兩件事.
就是離棄我這個活水的泉源.
詩篇也都講到.
因為你那裡有生命的泉源.
我要將水囂冠.
口渴的人將靈囂冠你的後裔.
以塞拉書也都繼續講到.
你們一切乾渴的.
都當來就近這個水.
所以一早舊約就提到這個活水.
就是耶穌基督帶給人的救恩.
聖靈要賜給人一個生命.
耶穌基督說.
你若知道神的恩賜.
又知道對你講話的是誰.
你就祖求祂 祂必給你活水.

$^{1121}$這是極大的一個震撼.
因為祂就活在.
喝來喝去還頸渴這個水當中.
每天都要對著這班被指責祂的人.
但現在有人說可以活水給祂.
活水給祂代表著.
我還是不是不需要來這裡打水呢.
所以他就有一個質疑.
他說你是誰.
難道你比我的祖宗還要大嗎.
耶穌基督.
你是一個三十多歲出頭的人.
我們喝的一直是這個水.
你說你有活水.
你說你有水.
我們這個是不是活水.
是我們祖宗流下來的.
難道你比我的祖宗還要大.
所以他又從雅各的時代.
這口井是從雅各的時代.
到耶穌的時代.
大概是兩千多年.
很深.
一百多英呎.
在一些學者認為.
是在巴基斯坦最深的一個口井之一.
所以這位婦女說.
你憑什麼說.
你能夠賜這些水.
比我們這個口井還要好.
我們雅各.
我們的祖先.
留給我們的已經很好了.
你說還要比它好好.
所以當她聽到一個真理的時候.
不是她所熟悉的.
也不是她現在明白的時候.
她是懷疑這件事.
她也看到一個.
不是那麼有權威的人.

$^{1161}$說一個比祂祖先.
說得還要有權威的話.
她也不是那麼相信.
她也習慣了.
我喝的是雅各.
我們祖先所留下來的.
我們的傳統.
我們的信仰.
這個才是最好的.
所以人又回到自己.
習慣那裡才是最好的.
對那些稍微不是那麼明白.
不是那麼熟悉的事.
也是很容易去懷疑.
耶穌基督沒有跟祂爭辯.
沒有說.
我比雅各阿爺還要大.
我是亞伯拉罕.
我比亞伯拉罕還要大.
祂也沒有跟祂解釋.
我就是那個太初之道.
我是那個大二位格的.
今天我們也一樣.
不需要跟無道者說得很深的東西.
因為這個太難明白了.
耶穌基督沒有跟祂爭論.
只是說了一句.
犯喝者水的還要再喝.
人若喝我所賜的水.
就永遠不喝.
很簡單的一句話.
你喝了你的水還要再喝.
喝了我這些水你就不再喝了.
這個女的一直在追求一些她慾望的事.
她可能有很多慾望.
很多不同的慾望.
最明顯的就是她背後的.
這些不見得光的那種慾望.
世上的人一直都在追求.
明星,金錢和情慾.

$^{1201}$但這些都不是真正可以使我們滿足的.
我們得不到的時候.
覺得那樣東西就是我一直想要的.
如果我實現到這個目標的時候.
拿到這種生活的時候.
就是我最大的幸福.
但當你得到那種生活的時候.
你才發覺.
都好像沒有什麼特別.
不過如此.
傳統認為所羅門就是全途輸掉的作者.
他是擁有財富,智慧,權勢的.
他人生高峰的時候.
他什麼都有了.
但是他寫全途輸掉.
竟然寫了一句.
虛空的虛空.
凡事都是虛空.
我有這一切的.
都依然是虛空.
這個婦女.
她有個不斷地去喝這些.
每天要喝水.
每天去享受那種罪惡.
但她都沒有得到真正的滿足.
因為有限的事物.
是沒有辦法填滿人的那種心靈.
那種無限的渴望.
日光之下一切的事物都不可以.
奧古斯丁說過.
如果我的心.
不可以安息在神裡面.
我就永遠得不到那種安息.
因為人是神所做的.
他有一個向往永恆的渴望.
這個渴望只可以神.
他無限的一位神.
才可以使我們滿足.
人要進入到這位神的裡面.
才可以真正得到這份滿足.

$^{1241}$不是那些每天你去喝水.
不是那些每天享受的罪惡.
能夠讓你享受那種滿足.
是耶穌基督的活水.
這個活水的源泉來到的時候.
你進入這個活水的源泉的時候.
你的心靈才可以得到滿足.
所以耶穌基督就說.
人若喝我所賜的水.
這個水在他裡面.
成為泉源.
湧到永生.
所以井水可以供應幾千年.
但是你喝了之後.
每天擔每天飲.
你依然後悔.
唯有耶穌基督.
祂是活水.
在你接受祂的時候.
你心接受祂為救主的時候.
祂是活水的源泉.
祂可以不斷的湧流.
讓你的心靈真正得到滿足.
不是你每天去和其他人去搞那些事.
讓你得到滿足.
而是得到這個活水的時候.
你才可以.
所以在這裡.
又有一個很關鍵的轉折點.
先生請將這個水給我喝.
福音傳講的時候.
很關鍵的一個東西.
不是我逼著人去信.
是讓他去知道.
這個人自己需要什麼.
開始的時候.
耶穌基督說請你給我水喝.
現在呢.
祂說請你給我水喝.
原先.

$^{1281}$耶穌基督.
他以為耶穌基督口渴.
耶穌基督很精心來到這裡.
刻意來到這裡.
我們看後面知道.
耶穌基督沒有喝祂的水.
沒有寫到喝祂的水.
但是現在呢.
反而是這位婦人要喝這個活水.
所以是從主動又變成被動.
被動又變回主動.
原先是耶穌基督主動.
似乎耶穌基督是很被動的.
現在呢.
是祂主動.
耶穌基督.
讓祂能夠察覺到.
祂的需要是什麼.
我們向人傳福音.
不是我們逼著人去推給人.
而是讓祂知道.
祂真正需要的是什麼.
不是靠那種技巧和能力.
是出於神的能力.
所以耶穌基督沒有強迫祂.
而是一步一步引導祂.
知道讓祂看到.
祂真正那種饑渴是什麼.
祂沒有推祂.
反而引導祂.
讓祂變為一個主動.
主啊.
請你給這個水給我喝.
我們教會呢.
其實我們今天也是這樣.
不是我們去逼人.
而是說.
是讓祂察覺到.
祂真正的需要.
很多時候我們叫弟兄姊妹.

$^{1321}$回來團契或者是茶經.
結果呢.
每個人都來一兩次就覺得.
好像聽又不是聽得很明白.
勉強呢.
又好像.
勉強一會兒就可以.
但是久了又睏了.
所以接受又不是.
不接受又不是.
組長打電話.
團契領袖打電話.
來又不是.
不來又不是.
所以.
祂是沒有那種主動那種渴望的.
但是耶穌基督就很不一樣.
祂讓人察覺到那種主動那種.
祂真正渴望的那種東西是什麼.
祂真正饑渴的那樣東西是什麼.
讓祂自己去追求.
所以不是被逼的.
是喚醒人的那種渴望.
是被吸引.
內在那種渴望.
跟尼哥帝母說.
是你直接說.
你必須重生.
跟撒瑪利亞婦人說.
先建立關係.
引導她進入那種真理.
所以全福音.
沒有固定的一個形式.
兵無常形.
水無常態.
是變幻莫測的可以說.
但是我們需要的是.
要有那種肅靈的那種敏銳.
對人的那種了解.
從神而來的那種智慧.

$^{1361}$你才可以真正去打開人的那種心.
耶穌基督就是.
一直去捉住.
祂那種需要是什麼.
讓祂能夠.
察覺祂真正的需要是什麼.
祂沒有一下子就揭穿祂.
你就是一個罪人.
你就是一個污穢的人.
祂先建立關係.
引導祂那種興趣.
喚醒祂那種需要.
引導祂進入真理.
最後才捉住這個罪的問題.
所以耶穌基督說了一句.
很突然的.
你去叫你的丈夫來.
很關鍵.
耶穌基督已經不可以再逃走了.
祂只可以說.
我沒有丈夫.
祂沒有辦法去遮掩.
祂知道耶穌基督知道一切.
很多時候我們.
會為了自己的面子.
而保護那種羞恥之心.
耶穌基督沒有責備他.
雖然這句話很不老實.
你沒有丈夫是不錯的.
耶穌基督沒有立刻反駁他.
反而肯定他.
然後才揭露那種全部的事實.
你已經有五個丈夫.
現在有的.
這個都不是你的.
耶穌基督沒有一下子推翻他.
反而很自然地說出這個事實.
你之前五個都不是你的.
你已經有過五個了.
現在這個五個半不是你的.

$^{1401}$這半個好像是你.
也好像不是你.
真理就出來了.
沒有傷害人.
所以耶穌基督沒有騙人.
祂先是體諒人.
也保護人.
尊重人.
祂沒有一下子撕裂那種關係.
而是建立關係之後.
才帶他進入那個真理.
很多時候我們很容易受那種毀謗.
人家說得不好的時候.
或者做得不是那麼好的時候.
或者做得好也好.
人家會毀謗.
我們需不需要回應呢?.
耶穌基督.
人家對他不好.
他沒有立刻去回應.
因為人的口在人家那裡.
這是我們需要操練的那種功課.
不是急著去辯護.
為自己辯護什麼.
反而是願意去捨己.
我們不是為了自己的名譽去辯護.
很多時候我們知道.
其他弟兄姊妹受了委屈.
我們就會說主必伸冤.
但是我們自己受那種委屈.
人家無理的指責的時候.
我們很容易自己做主.
自己去伸冤.
所以我們不需要這樣做.
不是急著糾正人家.
也不是急著去辯護.
所以婦人說我沒有丈夫.
這是一種遮掩.
是為了保護自己.
耶穌基督沒有羞辱他.

$^{1441}$也沒有刺痛他.
反而說你沒有丈夫是對的.
現在先肯定他.
再說出那種事實.
如果一下子就指出的時候.
就很危險了.
你看關公刮骨療傷.
中了毒之後要刮骨療傷.
我聽到這個故事很愕然.
當然相傳是華佗幫他刮骨的.
但是一些其他事實不是.
如果那時候真的是華佗幫他刮的時候.
先用麻醉藥就很不一樣了.
為什麼?因為沒那麼痛.
你將人家的骨刮出來的時候.
將肉撕開的時候再刮骨的時候.
這樣就很難去承受了.
所以用麻藥的時候就很不一樣了.
我們傳福音呢?.
傳福音也是我們勸人去悔改.
不是直接指出人家的罪是什麼.
強烈去指出和拆穿.
受不了的也會逃避那種指責.
耶穌基督不是馬上就對付那個罪.
而是帶著恩典帶著智慧去處理這件事.
所以當用麻藥的時候.
人家就很容易去接受了.
所以三個原則.
你要真心去了解人.
人家的痛楚是什麼.
你要全面去了解人.
不要片面.
還有建立那種關係.
你才真正可以去幫人.
那個時候你才有機會去幫人.
才可以去輔導人.
我記得有一個同學.
他在一個教會服事.
教會的同工說.
你講話太粗俗了.

$^{1481}$另外一個人說.
你的笑聲太大了.
我們很難忍受.
有些人直接跟他說.
你這麼粗魯.
你也不像傳道人.
他那時候就說.
第一個反應就說.
我很想離開這間教會.
當他有這個想法的時候.
有一位很有經驗的牧師.
很懂得體諒人的牧師.
他出來說話說.
你們知不知道這位願意做傳道的.
他之前是在這個社會打滾的.
撈偏門的.
做過一些不見光的生意.
但是今天他能夠服事人.
是神很大的一個恩典.
他雖然說話還是這麼直接.
這麼大聲.
甚至這麼刺耳.
但是這個人很誠實.
沒有盔甲.
所以我們當要幫助人的時候.
向人傳福音.
去輔導人的時候.
去托人一把的時候.
關乎我們的知識有多少.
是你要有真正愛人的那種心.
真正體恤人的那種軟弱.
你建立那種關係.
你才可以真正去幫助人.
當人感受到你的真誠和關心的時候.
你明白他的時候.
他才可以向你闖開.
所以耶穌基督就是最好的榜樣.
祂對罪不妥協.
對人充滿了恩典.
祂說出真理又不是傷人.

$^{1521}$祂醫治人.
帶給人盼望.
祂是真正的一種模樣.
是最有能力的一種福音.
最後親愛的朋友.
如果你的心靈是乾渴的.
來到主耶穌基督那裡.
得著這份活髓.
親愛的弟兄姊妹.
願主賜給我們傳福音的智慧.
多為主得人.
也求主保守我們.
報受我們「教會合一 彼此相愛」.
\newpage



\section{馬可福音 9:30-37}
\label{sec:032091}
\textbf{耶穌如何回應門徒的爭論}~~~~羅杰才牧師
\newline
\newline
連結: \href{http://www.church.com.hk/acms/content.asp?site=cdc\&op=show\&type=product\&code=032091\&layout=sermon}{\texttt{032091}} ~~~~ 語音日期: 2026-04-26
\newline
\newline
\hyperref[sec:032090]{\small{< < < PREV SERMON < < <}}
~
\hyperref[sec:index_chronic]{\small{[返順時目]}}
~
\hyperref[sec:032092]{\small{> > > NEXT SERMON > > >}}
\newline
\newline
馬可福音 9:30-37
\newline
\begin{longtable}{cl}
\hline
\hline
章節 & 經文 (和合本修訂版)\\
\hline
9:30 & \begin{tabularx}{0.7\textwidth}{X} 他們離開那地方，經過加利利；耶穌不願意人知道， \end{tabularx} \\ \\ \relax
9:31 & \begin{tabularx}{0.7\textwidth}{X} 因為他正教導門徒說：「人子將要被交在人手裡，他們要殺害他；被殺以後，三天後他要復活。」 \end{tabularx} \\ \\ \relax
9:32 & \begin{tabularx}{0.7\textwidth}{X} 門徒卻不明白這話，又不敢問他。 \end{tabularx} \\ \\ \relax
9:33 & \begin{tabularx}{0.7\textwidth}{X} 他們來到迦百農。耶穌在屋裡問門徒說：「你們在路上議論的是甚麼？」 \end{tabularx} \\ \\ \relax
9:34 & \begin{tabularx}{0.7\textwidth}{X} 門徒不作聲，因為他們在路上彼此爭論誰最大。 \end{tabularx} \\ \\ \relax
9:35 & \begin{tabularx}{0.7\textwidth}{X} 耶穌坐下，叫十二個使徒來，說：「若有人願意為首，他要作眾人之後，作眾人的用人。」 \end{tabularx} \\ \\ \relax
9:36 & \begin{tabularx}{0.7\textwidth}{X} 於是耶穌領一個小孩過來，讓他站在門徒當中，又抱起他來，對他們說： \end{tabularx} \\ \\ \relax
9:37 & \begin{tabularx}{0.7\textwidth}{X} 「凡為我的名接納一個像這小孩子的，就是接納我；凡接納我的，不是接納我，而是接納那差我來的。」 \end{tabularx} \\ \\
[1ex]
\hline
\hline
\end{longtable}
$^{1}$各位弟兄姊妹,早安.
早安.
很高興在這個地方跟大家見面.
一齊來.
屬拜,一齊來練習上帝的說話.
這樣.
我今天是選了一段經文.
電話就讀了出來,是不是?.
可以打出來給大家看到的.
而我的題目就是經文裡面其中的一個段落.
就是耶穌如何去回應門徒的憎恨.
我的講道是有一個特色.
因為我比較多用這種方式.
就是一段經文.
我多數講道都是一段經文.
我就在這段經文裡面用來講道.
為什麼呢?.
我都是很想透過講道.
給弟兄姊妹多一個明白.
怎樣去讀一段經文.
講道的時候,我希望大家.
原來一段經文.
我們這樣讀是可以這樣.
或者這麼豐富的.
我想帶出這個年代的效果.
當我們讀聖經的時候.
原來牧師一段經文.
可以講到這麼多東西.
看到這麼多東西.
我們讀那段經文的時候.
是不是可以看到這麼多東西呢?.
希望我的講道可以令人對聖經.
是更加有一種興趣.
或者想在那段經文裡面.
明白更加多.
所以我的講道多數是這樣的.
當然不是絕對的.
有時候提高是不行的.
今天我都是用這個方式.
這段經文是講門徒是增大.

$^{41}$在這裡說是誰做大的.
都是一個比較多成為講道素材的經文.
所以如果回教會久了.
應該不會太陌生.
都聽過這段經文.
其實這段經文在這件事件.
在三短的福音書都有記載的.
我今天讀的是《馬河福音》.
其實《馬太》有記載.
《路加》也有記載.
不過短一點.
有一點點不同的.
因為重點內容是多一點的.
但我今天都注意了.
三段經文所講的都很接近的.
經文記載的事件是發生在何時呢?.
大概是在耶穌傳道的中後期.
司祭約翰被殺之後.
耶穌自己的處境也是有些危險的.
當時耶穌出來傳道.
就招到當時建制的法利塞人.
祭祀.
整個傳道就有些敵視.
門徒都感覺到壓力.
亦都被當時的政府因為有一群群眾在制度以外.
又這麼輕盛.
而且大家又覺得好像要將他推出來做彌賽.
然後帶動一個社會政治的改革.
門徒都有這些不同的心情在.
耶穌是知道的.
門徒都知道的.
所以到那時候司祭約翰被殺之後.
其實耶穌開始透露他將會有甚麼遭遇.
所以因為他是教導門徒.
人子將要交在人的手裡面.
他們要殺害他.
但耶穌就講他被殺之後.
三天後他會復活.
但起碼在前面聽到耶穌要遭難.
會被殺.

$^{81}$即是你如果跟著一個領袖.
這次暫時先講中華興.
會被殺.
即是他要被除掉.
門徒跟著這個領袖.
他就想講.
耶穌的說話門徒都是聽得很認真.
很震驚.
不敢問他.
這個問題就是不敢.
又不敢問他.
門徒不明白又不敢問他.
即是耶穌跟他們這樣講.
即是說在這個時候.
跟著他們又爭論.
跟著他們就在那裡討論一些事.
時秋就問他們.
你們在那裡爭論些甚麼.
要接回前面.
你想才行.
耶穌跟他們講了.
我會遇難.
又不知何時會遇難.
隨時都會遇難.
誰知整個團隊裡面.
是有一個很大的危機.
突然耶穌被人拿走了.
出事了.
那怎麼辦.
他們都很震驚.
又不敢問耶穌.
我繼續行.
路上繼續行.
這時他們開始討論.
爭論一件事.
哪個做大呢.
一定跟前面有些關係.
兩件事放在一起.
三段都是這樣的處境.
門徒明白.

$^{121}$他們就有一種擔心.
有一種擔心.
開始想怎樣去應對.
沒有了耶穌這個團隊.
誰去帶呢.
沒有了這個團隊.
誰去帶呢.
爭論誰為大.
我們聽到就很不開心.
說耶穌說誰來爭大.
爸爸都未死還爭身家.
但我們嘗試用另一個去想.
其實當時又沒有權.
又沒有資產.
甚麼都沒有.
這個團隊很初型.
甚麼都沒有.
爭甚麼呢.
其實沒有爭的.
是不是.
是沒有爭的.
其實爭甚麼呢.
爭責任.
這個團隊大家都做來珍惜.
大家都是重視.
即使出事都不可以散.
如果用這樣角度去解釋.
其實門徒不是爭誰拿多些東西.
誰拿多些威.
他們其實是想爭取.
想想如果沒有了耶穌的時候.
誰去做領導呢.
他們爭我們嘗試才代入的場景.
當然是我做.
約翰說何時輪到你.
我比你還聰明.
然後約翰又說甚麼.
爭一定是這樣.
而且多數都不是這樣.
他就說.

$^{161}$你都不如你吧.
我記得說肯定是我.
我不知.
你都不如你吧.
另外一個說約翰比較適合.
我們用誰帶呢.
我們回到一個小群體裡面.
像這裡這樣.
找誰來帶領教會呢.
就會有不同的堆.
這個這個.
未必是那個人自己就想站出來說我做.
誰就推誰.
幾種可能都有.
我想說的不要說他們爭大.
這些人就斬亂權力等等.
而是他們在一起.
就想議論耶穌不在的時候.
這個團隊是怎樣走下去呢.
哪些人.
哪個人或者哪些人.
是可以去擔當這個.
維持這個團隊的角色呢.
我覺得這樣的理解.
是比較合理的.
他們的出發點就是這樣.
這種情況一定都會發生在現在的教會.
我們就會捧這個捧那個捧這個捧那個.
不過現在我覺得有少少不妥.
一出事就不要與我玩.
這個比較多.
叫誰做呢.
有能之時反而就.
不知為何教會這麼小.
或者這麼多.
沒什麼好爭的.
威又不到他威.
著數又不到他著數.
你不要以為現在的教會.
我的理解就是這樣.

$^{201}$很多人一選入位.
選什麼的時候.
最好不要選我.
不是爭是縮的.
我們說門徒爭大.
我們說他不對.
其實我們想想都避去不如.
他們爭這個捧那個捧那個.
或者再是.
實際上他們不是擁有了很多.
權力財富在那裡準備分.
爭什麼呢.
他們在爭什麼呢.
爭怎樣去持續.
當耶穌不響動的時候.
這個團隊.
有人去幫他們.
耶穌就叫門徒來到加巴倫的時候.
走完那段路.
在路上議論說.
如果耶穌真的出事的時候.
誰作罷.
在屋裡就問.
你在路上爭論的.
議論的是什麼.
門徒就不作聲.
因為他們在路上爭論誰最大.
排到誰.
耶穌不在那裡.
於是就坐下.
就叫十二個使徒過來.
都是叫齊了.
是不是.
你知道嗎.
特別這樣寫.
叫齊了.
個個都到齊才說.
普通我們開會.
個個都齊了.
因為這件事都是重要的.

$^{241}$耶穌其實沒有否定他們.
也沒有責備他們.
不是說我都未死.
你們在爭.
不是這樣.
也沒有責備他們.
沒有否定他們.
這也是值得我們留意的.
聽完他們說話之後.
就沒有出聲.
沒有說.
那是彼得.
我早就說了是彼得.
約翰才對.
沒有耶穌也沒有判斷.
沒有責備他們.
也沒有說應該誰做.
你留意這個經文.
那個脈絡.
是不是.
接著就告訴他們.
你們誰願意做大那個.
就要做眾人之後.
做眾人的用人或者僕人.
講了一個原則.
我不管你們誰做大.
你們誰做大都不重要.
也不是誰做大之後問題就解決.
耶穌沒有完全越過了這一層.
因為他們在爭論誰做大.
耶穌是跳過了這一層.
說你們誰做大.
你們就要做世人.
你們要做僕人.
是不是.
意思就是說.
誰願意做領袖.
誰就要這樣做.
或者反過來說.
誰是這樣做.

$^{281}$誰就是領袖.
都可以.
誰願意做領袖.
誰就要這樣做.
所以因為帶著眾人的模樣.
去服侍其他人.
誰願意去服侍其他人.
誰願意做就是.
那個就是最大.
耶穌主耶穌的邏輯.
是和我們的邏輯不同.
但這是很重要很重要的一個屬性原則.
換言之.
你們不用爭話事權.
你們要爭的是服事權.
服事就不用權.
服事是甚麼權呢.
是應該爭服事的機會.
或者服事的參與.
我在螢幕會的時候.
都有這樣的經驗.
有些人.
特別有一個弟兄.
他經常想做團長.
做組長.
做隊長.
很喜歡做長的.
選不到他就發脾氣.
這樣的人.
不一定有些人喜歡他.
他本身有能力又有熱誠.
所以他有時爭到有時爭不到.
爭不到會怎樣呢.
他們都不給我機會服侍.
永遠投訴.
他想做.
我服侍都要給人機會嗎.
服侍不用人給機會.
職位要人.
崗位要人給機會.

$^{321}$但服事就不用.
沒有人阻到你.
鄧小平都阻到你.
我是要潛伏人.
我是要愛人.
我要關心人.
你怎樣做.
誰阻不到你.
所以你要去事奉.
是不能阻的.
所以你不是去爭一個位置.
所以越過去.
你要做眾人的僕人.
你去服侍.
網上教會甚麼都沒有.
沒有人認識我.
你這樣看自己是宿迷.
你有很多機會.
不會沒有機會服侍.
你想想.
有沒有機會服侍.
如果你說沒有.
我馬上打電話給你.
有很多機會.
你放一個小紅紙幫一個人.
做一些關心的事.
人人皆濟私.
怎會沒有呢.
所以你要做僕人.
大家一齊去服侍.
當然一個團隊會有人出來.
做決定.
做帶領.
但當大家這樣做.
這個人自然會被浮現出來.
因為大家都不爭.
他就自己浮現出來.
但不是.
你又撥開我我又撥開你.
就鬥了.

$^{361}$這是另一個場景.
耶穌說了一個很重要的事.
你們要做僕人.
如果我們有機會站在一個位置.
譬如做長.
我們不比其他人大.
我們是一個服侍的人.
有難處我們應該擔當多些.
應該這樣去想.
我們不是看我們自己身份.
而是看我們自己的責任.
然後耶穌說完之後.
他再做了一件事.
他叫一個小孩子來.
同一個場景來的.
叫一個小孩子來.
站在當中.
就配對話來.
凡是我的名接待一個小孩子.
就是接待我.
在這個馬達福音是這樣說的.
我實在告訴你們.
你們若不迴轉變成像小孩子一樣.
絕不能進天國.
可以自己謙卑像小孩子的.
他就是天國裡面最大的.
凡是我的名接待一個像小孩子的.
就是接待我.
馬太福音還多得少少大.
你要像一個小孩子.
這裡說你要接納一個小孩子.
是兩重意思.
我們讀的馬達福音.
他寫了要接納.
但如果你看這個馬達福音.
他還有多一重.
在這個禱告裡.
你要像一個小孩子.
做一個領導人.
要像一個小孩子.

$^{401}$領導人當然是成熟的.
不成熟怎做領導人呢.
但這個成熟的領導人.
就不是一個老氣橫秋.
而是他好像一個小孩子.
我們會以為小孩子是很幼稚.
領導的人.
不是像一個小孩子那麼幼稚.
而是他要保留一個小孩子的那種誠實.
那種良善.
那種好奇.
那種謙遜.
那種合法.
小孩子是很多好的物質.
耶穌是很特別的.
你們要接待小孩子.
你們要回傳小孩子.
因為天國裡面的人就是這樣.
我們長大了.
做大人了.
我們不做小孩子了.
但耶穌反過來說你們要做小孩子.
真的要記住.
因為天國裡面的人都是這樣的.
小孩子的特質就是天國人的特質.
如果你見到一個人.
好像小孩子那樣.
你會覺得跟他相處很舒服.
你同意嗎.
其實我們每個人都可以做小孩子.
為什麼一定要扮大人呢.
我七十多歲.
我都喜歡做小孩子.
我有一個孫子.
我跟孫子一起住.
我們住在新界的村屋.
三層一棟.
我女兒一家住三樓和天台.
我就住地下二樓.
一個複式單位.

$^{441}$新界的屋比較便宜.
感恩有這樣的環境.
住下去都舒服.
中間有一條來自樓梯的複式單位.
因為我孫子和我一起住.
所以他一日六十.
現在他都八歲了.
有一次我孫子給了我一個稱號.
公公是最寶貴的玩具.
我聽到心都跳了.
我是最寶貴的玩具.
我和他一起玩.
他很喜歡和我玩.
證明我有同質.
他出生時有很多感受.
我們大人教小孩子怎樣過生活.
要學懂技巧.
識字,識過馬路.
要教他們生活性的東西.
但小孩子向我們呈現甚麼是生命.
不要討厭小孩子.
耶穌有一次發怒.
有些人帶小孩子和嬰兒來.
英文叫禁止他們.
禁止貪奴.
聖經怎樣說.
耶穌受生氣.
你不覺得重要.
即是嚴重.
在整本福音書見到耶穌生氣了多少次.
為何你會生氣.
如果一個很成熟.
情緒又很穩定的人.
老師或校長.
亂生氣.
怎會這樣.
一定是那件事給他很大刺激.
或同事做了絕對不應該做的事.
你不應該這樣做.
而且做得很離譜.

$^{481}$帶來很大的傷害.
你才會生氣.
如果你做一個高級的領導人.
你都不會隨便發脾氣.
而且耶穌是一個超穩定的人.
他生氣了.
又為小孩子.
然後又說進天國的要是這樣.
因為門道當時的反應.
犯了極之嚴重的錯誤.
他把好人當賊犯.
當進天國的人拒絕了出去.
小孩子立即要受懲罰.
小孩子.
我們為何今天生活得那麼淪盡.
或做得那麼.
一個個都起了狂.
又不懂得變通.
誰知硬著頭皮.
小孩子是很靈活的.
你看他們玩是很靈活的.
很懂得集散.
玩什麼都搞得出來.
我孫女特別對勝.
經常想些東西出來.
誰知我一弄東西.
她說讓我教你.
她經常教我.
王老師.
我們就是這樣.
然後她又會問一些東西.
你覺得她無厘頭.
她很認真的.
她就不用這樣做.
小孩子就是這樣.
她是很挑戰你的.
她沒有不可以.
沒有說不可以.
有一次她在家裡想.
在二樓.

$^{521}$她拿一條繩子扔下來.
二樓扔下來.
於是就給她.
她就綁起繩子.
扔下來.
很開心.
然後又扔下來.
有大肌肉運動.
她不用煩我自己玩.
我很開心.
她扔了一天.
扔了兩天.
扔了整個星期.
不同的玩偶在扔.
小孩子可以重複做一些.
我們覺得很單調的東西.
她做得到.
用那種能力.
有些東西支持她.
然後一家人就去旅行.
十多八天.
那條繩子扔在那裡.
我已經不方便.
不過她在那裡玩.
沒辦法.
阿森.
你說我這樣拆她的好不好.
她說不要拆.
不經過她拆得.
很大權.
我就呆呆.
然後她就回應我.
外公.
她那時準備出來.
外公.
你有空.
跟他們玩滑梯.
他們會很開心.
我聽到呆呆.
她跟我上了一課.

$^{561}$原來她跟那些玩偶玩.
很開心.
是那些玩偶很開心.
她很開心.
是因為那些玩偶很開心.
這是甚麼人.
那些沒生命的東西.
她可以注入生命給她.
我們大人.
因為生的當死.
我們那些姑娘.
去到病房.
不知去到病房.
不明是多少號.
那些不是病人.
是多少號號碼.
管監獄那些.
她看不到人.
看到號碼.
我們見不到人.
很多時都是.
見到的是那件工作.
任務做完就算了.
我們看不到人.
這個小朋友.
原來她不是.
跟那個玩偶玩.
是那個玩偶開心.
她很開心.
那些沒生命的東西.
她都將它注入生命.
給她.
她那種想像力.
那種創意.
她有多大能量.
你想想如果一個領袖.
這不是麻煩.
這是機會.
是完全不同.
她有創意.

$^{601}$有注入.
她的動力.
在裡面出來.
小孩子.
我們沒有這些很久了.
如果你還有一些.
你覺得.
這個不是悶.
這個是靜.
不同.
悶和靜怎樣不同.
靜是很享受.
所以.
真是不同.
小孩子是很.
良善的.
很好奇的.
很感恩的.
她跟別人有衝突.
很快就忘記.
她很有生命力.
你做領袖的人要這樣.
耶穌說.
你們要好像一個小孩子.
小孩子其實.
是沒有權力的觀念.
怎樣得到權力.
怎樣鬥你不懂.
永遠都是.
她喜歡的時候.
她就很投入.
她能夠收.
能夠放.
她沒有謀生.
她只會吸收成長進步.
所以耶穌經常教.
我們要好像一個小孩子.
我們真是需要.
在小孩子身上學習.
我們失去了.

$^{641}$那些特質.
生命的特質.
小孩子.
亦都可以.
有另外一個.
解釋或看法.
耶穌說你做領袖.
你要接納小孩子.
另外一個解釋和看法就是.
小孩子代表幼小的.
細小的.
微小的.
弱小的.
耶穌為何會生氣呢.
因為門徒.
他不接納這些.
弱小的,幼小的,微小的.
人有一種傾向.
喜歡漂亮的.
強壯的.
重要的.
我們叫吹捧闊小的.
有些領袖都是.
哪裡可以爬得更高.
你要看那些.
弱小的,微小的,細小的.
幼小的.
接納小孩子可以說.
你要包容這些人.
你要愛護這些人.
教會有好多.
走也走得慢,說也說得慢.
論論盡盡.
年紀未必是小的.
剛好他的年紀很大.
或者他的健康,身體等等.
你要照顧到這些人.
教會怎樣去呈現.
他的角色呢.
他就要照顧這些人.

$^{681}$實際的.
操作上.
有很多東西我們要考慮.
但是.
好像我們要照顧一個小孩子.
我們照顧一個小孫女.
你要保護她.
愛護她.
你要教她,訓練她.
如果她有能力學的話.
你就跟她玩.
有很多事要做.
要幫助她.
那些年長的.
或者大的.
如果她是弱小的.
她是微小的,她是細小的.
她就會說.
沒有用的.
其實不是,可以說是為了轉化.
提升她.
當大家都看到.
感受到.
原來教會這個群體.
是這樣的氣氛.
去做的.
去凝聚.
她就有安全感.
她就這樣去尊重.
那個.
那個主輪就大得到起來.
所以.
如果一個團隊,一個群體.
一個社會,一個教會.
它是增大的.
向權力.
向能力.
去靠攏.
大家就爭哪個大力的.
大聲的.

$^{721}$出去行事.
這個是社會.
或者團隊的模式.
教會不應該這樣.
教會是甚麼呢?.
教會是俗世的.
俗世也是.
是服事的.
是純真的.
所以.
教會應該是一個服事的群體.
一個愛的群體.
一個生命的群體.
生命的群體.
真是有一個特色.
那些年長的是犧牲自己.
貢獻自己.
但那些年幼的.
幼小的可以成長.
你看過嘉玲就是這樣.
是不是.
讓她可以成長.
她不是說.
你小的你走一邊.
大人坐一邊.
不是這樣的.
你是照顧她.
這個道理是很顯示.
很清楚.
回到教會其實都一樣.
你和哪些人最有需要呢?.
哪些人要幫助呢?.
有能力的那些.
有領導的那些.
就照顧他們.
是不是.
有時你在這方面弱.
在另一方面強.
你都要貢獻你的所長.
才可以.

$^{761}$有牧師這樣說.
我個個都扮給人照顧.
不行的.
這個方向一調整之後.
就很不同.
所以耶穌怎樣回應那些門徒.
那些門徒真命的時候.
其實他們就處於一種思想模式.
如果耶穌講到我們誰大呢?.
耶穌就告訴他們.
如果我講到.
你們怎樣大呢?.
怎樣走呢?.
去做牧人.
去服侍人.
亦都要.
回轉那些孩子.
你要有一種.
生命.
一種純真.
你能夠看到.
看到.
人看到生命.
那種真正的價值.
其實都始因我們.
我們一起祈禱.
求主始因求主幫助.
你的教導.
是最好的.
求主讓我們都能夠.
聽得到你的教導.
學習.
去做一個牧人.
又回轉一個小孩子.
又接納.
那些小孩子.
就是那些幼小的.
弱小的,微小的.
讓我們都成為.
一個.

$^{801}$屬主的團體.
祈禱奉主耶穌基督的名.
\newpage



\section{路加福音 16:1-13}
\label{sec:032092}
\textbf{門徒的錢財觀}~~~~俞大鳴牧師
\newline
\newline
連結: \href{http://www.church.com.hk/acms/content.asp?site=cdc\&op=show\&type=product\&code=032092\&layout=sermon}{\texttt{032092}} ~~~~ 語音日期: 2026-05-03
\newline
\newline
\hyperref[sec:032091]{\small{< < < PREV SERMON < < <}}
~
\hyperref[sec:index_chronic]{\small{[返順時目]}}
~
\hyperref[sec:032093]{\small{> > > NEXT SERMON > > >}}
\newline
\newline
路加福音 16:1-13
\newline
\begin{longtable}{cl}
\hline
\hline
章節 & 經文 (和合本修訂版)\\
\hline
16:1 & \begin{tabularx}{0.7\textwidth}{X} 耶穌又對門徒說：「某財主有一個管家，有人向主人告管家浪費他的財物。 \end{tabularx} \\ \\ \relax
16:2 & \begin{tabularx}{0.7\textwidth}{X} 主人叫他來，對他說：『我聽到了，你做的是甚麼事？把你所經管的交代清楚，你不能再作我的管家了。』 \end{tabularx} \\ \\ \relax
16:3 & \begin{tabularx}{0.7\textwidth}{X} 那管家心裡說：『主人辭我，不用我再作管家，我將來做甚麼呢？鋤地嘛，沒有力氣；討飯嘛，怕羞。 \end{tabularx} \\ \\ \relax
16:4 & \begin{tabularx}{0.7\textwidth}{X} 我知道怎麼做，好叫人們在我不作管家之後，接我到他們家裡去。』 \end{tabularx} \\ \\ \relax
16:5 & \begin{tabularx}{0.7\textwidth}{X} 於是他把欠他主人債的，一個一個地叫了來，問頭一個說：『你欠我主人多少？』 \end{tabularx} \\ \\ \relax
16:6 & \begin{tabularx}{0.7\textwidth}{X} 他說：『一百簍油。』管家對他說：『拿你的賬，快坐下，寫五十。』 \end{tabularx} \\ \\ \relax
16:7 & \begin{tabularx}{0.7\textwidth}{X} 他問另一個說：『你欠多少？』他說：『一百石麥子。』管家對他說：『拿你的賬，寫八十。』 \end{tabularx} \\ \\ \relax
16:8 & \begin{tabularx}{0.7\textwidth}{X} 主人就誇獎這不義的管家做事精明，因為今世之子應付自己的世代比光明之子更加精明。 \end{tabularx} \\ \\ \relax
16:9 & \begin{tabularx}{0.7\textwidth}{X} 我又告訴你們，要藉著那不義的錢財結交朋友，到了錢財無用的時候，他們可以接你們到永遠的住處去。 \end{tabularx} \\ \\ \relax
16:10 & \begin{tabularx}{0.7\textwidth}{X} 人在最小的事上忠心，在大事上也忠心；在最小的事上不義，在大事上也不義。 \end{tabularx} \\ \\ \relax
16:11 & \begin{tabularx}{0.7\textwidth}{X} 若是你們在不義的錢財上不忠心，誰還把那真實的錢財託付你們呢？ \end{tabularx} \\ \\ \relax
16:12 & \begin{tabularx}{0.7\textwidth}{X} 如果你們在別人的東西上不忠心，誰還把你們自己的東西給你們呢？ \end{tabularx} \\ \\ \relax
16:13 & \begin{tabularx}{0.7\textwidth}{X} 一個僕人不能服侍兩個主；他不是恨這個愛那個，就是重這個輕那個。你們不能又服侍神，又服侍瑪門。」 \end{tabularx} \\ \\
[1ex]
\hline
\hline
\end{longtable}
$^{1}$主席.
今日的經文又是另外一個比喻.
實際上盧家醫生記載的大部份比喻.
都是在他從加利利去耶路撒冷的路途中講的.
即是說從《路加福音》九章五十一節以後就有很多的比喻.
而《路加福音》記載的比喻中.
其中有十七個都是其他福音書沒有記載的.
今日所記的比喻也是其他福音書找不到的.
這個比喻也是令很多人誤解或者覺得有些困惑.
比較難明白.
所以我想先和大家思考一下耶穌講比喻的目的和比喻的特色.
耶穌用比喻教導群眾和門徒.
很多時候都是有不同的目的.
有時候他用比喻是教導門徒.
但是由於有其他人不是門徒的人在場.
於是他就用比喻來說到是希望門徒明白.
而其他人就不明白.
但是我們又會想門徒都可能不明白.
門徒不明白不要緊的.
因為他們經常跟著耶穌說.
你白天聽完講到不明白.
晚上可以問嘛.
在聖經裡面的確是有這樣的記載的.
他們在晚上的時候去問耶穌.
在福音書中最有名的比喻.
撒種比喻就是這樣.
其實門徒聽完都不明白撒種比喻的意思.
後來問耶穌.
耶穌反而引用了以賽亞斯所說的話.
或者人問聽了又聽卻不明白.
看了又看卻看不見.
意思就是說有些外人是不會明白他講比喻的意思.
比喻的特色一般來說是一個故事的形式.
故事的形式有什麼好處呢?.
就是在聽的人.
你即使不是很明白比喻是想講什麼.
真正的涵義是什麼.
你都會記得故事的.
例如上星期講的第十五章那三個比喻.
其實是責備法利賽人和民事的價值觀.

$^{41}$但是他們是不是真的能夠明白和即時領會到那些比喻的涵義呢?.
我不敢說.
但即使他們不明白.
那幾個故事很容易他們會記得那三個比喻.
如果當中有些人比較聰明或者願意思考的話.
會回去再想想那三個比喻.
或者他們會明白的.
三個比喻中第一個是講牧羊人失了一隻羊.
只是失了1\%的羊.
他都要放下99\%的羊去找那隻迷失了的羊.
找到之後就歡喜快樂.
還要和鄰居一起分享.
然後講一個婦人失了一個銀幣.
是她財富的十分之一.
她一定要仔細尋找.
找到之後當然也是歡喜快樂的.
這些其實我相信都是當時猶太人的一般正常反應.
即使我們今天也會失了東西找當然是開心的.
但最後那個浪子的比喻.
那個爸爸有兩個兒子.
失去了一個.
那他失去的比例是非常之高的.
他失去了一半.
失去了二分之一.
而當那個兒子回來的時候.
他是非常歡喜快樂.
但是那個大哥就不開心.
爸爸就對大兒子說.
你這個弟弟是死而復活失而復得的.
所以我們理當歡喜才對.
理當歡喜慶祝.
即是說歡喜是正常的.
合理的.
慶祝是合理的.
只有思想不正常的人.
價值觀不好的.
不對的人才會不進去開派對.
要自己不參加.
就是說那些法理財人和文士.
他們的價值觀是完全錯誤的.

$^{81}$對羊對錢很緊張.
印他就不緊張的.
當然主耶穌都會用比喻教導門徒.
對財富是有正確的價值觀.
他曾經說過一個叫做無知財主的比喻.
記載在路加方十二章.
那個比喻說到一個財主.
那年大豐收.
於是他看到自己有那麼多的穀物.
就想我一於可以享受世界.
不如這樣.
我把小倉庫拆掉.
建一個大倉庫來儲糧.
從此以後就可以安安逸逸.
吃喝快樂.
比喻結束的時候.
是上帝對那位財主說.
無知的人.
今夜就要你的性命.
你所預備的.
要歸誰呢?.
這個比喻主要教導門徒.
調教好對錢財的價值觀.
不要為今生憂慮.
吃什麼穿什麼.
這一段經文是耶穌對門徒說的.
不過從十四節那裡我們知道.
其實都有法理錯人在場.
主要的目的仍然是教導門徒.
要有正確看待錢財的觀念.
《禾下本聖經》就比喻一個標題.
叫做「不義的管家」.
這個標題其實可能會給讀者一個感覺.
比喻一個管家是不義的.
我們不要學他.
但是從比喻本身的內容來說.
是無可厚非的.
但是從比喻的目的來說.
這個標題可能會給誤導讀者.
很可惜《還聖》的譯本也用了這個標題.

$^{121}$我相信他們用這個標題也有理由的.
因為第八節特別提到那個不義的管家.
但是新漢語就用了另一個標題.
新漢語的標題叫「善用金錢」.
這個是從比喻的目的來比喻一個標題.
「善用金錢」我覺得比較好的標題.
是現代中文譯本.
在這幾個版本中比較好的.
叫做「機警管家的比喻」.
很明顯講出機警.
管家是有機警性.
為什麼好呢?.
因為這個比喻不是一個反面教材.
而是要門徒學習這位管家的機警行為.
強調他是懂得把握時機運用錢財.
而不是說這個不義的管家.
雖然今生的生活有保障.
但是死後他死定的 會落地獄的.
所以我們千萬不要惡他.
大家要知道聖經出的標題.
不是上帝的默示.
不用緊張 改一改是沒有問題的.
這些標題只不過是翻譯聖經的人加上去.
方便讀者找經文的時候容易一點.
所以我們覺得標題不好.
刪掉改另一個沒有問題.
希望今天能夠幫到大家.
正確了解這個比喻.
使我們會成為錢財的主人.
而不要做金錢的奴隸.
很明顯這個比喻跟錢有關.
在福音書裡面.
主耶穌講錢比祂講天堂地獄更多.
我相信由於錢的確與我們的生活息息相關.
很多人都會說錢不是萬能.
但是沒有錢就萬萬不能.
故事主耶穌不能夠不教導不同.
怎樣對待錢.
我們也要學習正確對待錢.
有十三節的經文可以分為兩段.

$^{161}$一段八節是一個比喻.
其實九節到十三節就是主耶穌教導門徒.
從這個比喻裡面要怎樣去應用.
或者得到的教訓.
一開始在比喻當中提到一位主人和一位管家.
發覺以前的有錢人是很會賺的.
他們有錢之後自己就會開始去享受人生.
不想工作了.
怎樣處理日常的事務呢.
有錢嘛.
請人做就可以了.
請個管家和他做其他的事.
在《創世記》裡面我們已經看到.
那裡提到那個埃及的侍衛長波提弗.
他真是一個典型的例子.
他發覺這個藥膳很厲害.
做人又忠心.
交給他辦的事又一定順利.
藥膳辦事他真是完全放心.
於是他就很放心將一切的事都交給藥膳去做.
除了自己飲食之外.
什麼都不管了.
從今天比喻中這位主人.
我們看到他都是差不多.
他似乎請了管家打點一切之後.
自己就去享受生活了.
很可惜他這個管家就沒那麼忠心.
這位管家就被人控告.
他告訴主人.
你那個管家不行.
他浪費你的財物.
於是主人就要管家.
你快點搞清楚賬目.
你要交賬.
我要炒你魷魚.
夾些茶飯.
我們留意到這位管家收到主人的通知之後.
留意他的反應.
他沒有去找那些管家訴苦.
他也不是馬上去求職訊.

$^{201}$而是他坐下來思考.
他在想什麼呢.
第一他在想一個最關鍵的問題.
我為什麼要打工呢.
打工只不過是為了生活.
所以沒工作不是最關鍵的問題.
最重要的問題是.
沒工作之後我能不能夠生活下去.
我靠什麼繼續生活下去.
這位管家他也想到.
他其實很難再找一份管家做.
沒有其他人肯請他.
因為在那些小的地方.
你這個人被人解僱了.
因為不忠心.
一定臭名遠播.
不用想了.
耶穌時代的社會是農業社會.
在十五章也看到.
浪子回頭的比喻裡面.
大兒子在農場工作.
小兒子走投無路的時候.
也是找到一個農場工去放豬.
這位管家就在想.
去農場打工.
可能這位管家年紀也有點大.
他不夠力去挫地.
農場不只是放豬.
放豬其實也很舒服.
但在農場種植的時候就要挫地.
覺得自己已經沒什麼力去挫地.
知道這個不行.
不是好辦法.
什麼都不做.
行乞好不好呢.
利用別人的善心行不行呢.
行不是不行.
不過老年那些掛不住.
面子放不下.
所以也不是好辦法.

$^{241}$第三節提到他思考這兩點之後.
表明這個管家也挺聰明.
他起碼有自知之明.
知道自己的限制.
他不是一個愚蠢的人.
他知道自己有什麼限制.
也很坦白地知道自己也喜歡面子.
於是他小心評估了當前的情況之後.
他同時知道自己仍然有的資源和機會.
是什麼呢.
他仍然有的資源和機會就是.
他主人要他交帳.
但不是馬上交.
不是明天要交.
因此在沒交帳之前.
他仍然是管家.
他仍然有權可以處理主人的事務.
我想中國人大概現在也說起來.
可能十幾二十年前.
有句流行的說話.
有權不用過期無效.
意思其實是鼓勵貪官.
你還是做官的時候.
你有得貪你就去貪.
不然你一下台就沒得貪了.
這個管家也是想到這個.
他要在他的權力未曾過期之前.
就要運用他短暫的機會.
用主人的錢財去結交朋友.
以便那些朋友日後會回報他的恩惠.
會照顧他.
我們留意到第八節.
那位主人誇獎這位不誠實的管家.
但他並不是誇獎他不誠實這件事.
任何人都知道他不誠實的行為是錯的.
應當受責備.
但主人卻誇獎他的精明和機警.
事實上這個比喻奇怪的地方不是這個局.
主人的誇獎也是很奇怪的.
如果是主人可能會罵他一頓.

$^{281}$立刻抓他去見官有份.
還誇獎他做什麼呢.
正常人可能也會很生氣.
但這位主人卻稱讚他.
因為這位主人是代表上帝.
上帝是擁有一切的.
錢財對上帝來說並不重要.
所以主人反而看到管家精明和機警的地方.
比喻講完了.
我們看看主耶穌想我們怎樣應用這個比喻.
第九節是要我們明白.
錢只是在今世有用的.
當我們有一天離開世界的時候.
錢又不是廢紙.
但錢會變成其他人.
不再是我們.
所以我說你有財不用過期無效.
你有財好用.
否則你過期無效.
我們要把握機會善用金錢.
怎樣才算是善用呢.
第九節說要藉著不義的善財.
要交朋友.
這些朋友不是今世之子.
這些朋友可以接你們到永遠的住處.
最重要當然是主耶穌.
其次弟兄姊妹之間要有好的友誼.
要彼此扶持.
可以多些和弟兄姊妹相約去飲茶吃飯.
第十節.
所謂耶穌的門徒.
我們要忠心.
所以很明顯這個故事不是教我們學不義的管家.
不是忠心.
耶穌在這裡說明我們應該要忠心.
並且在小事上忠心.
在大事上也忠心.
在小事上不義.
在大事上也不義.
這個比喻不是教我們要學管家的不義行為.

$^{321}$要學習他機智的行為.
但我們要做忠心的門徒.
耶穌所說的比喻當中.
也有一個講到主人將錢財交給三個僕人.
然後出去遠行.
遠行之後回來要他們交帳.
這個比喻中的管家提早要交帳.
但其實每一個人將來都一定要向上帝交帳.
而錢財的確有相當大的吸引力.
但我們要記得錢財不用只是一個數字.
放在銀行裡面你看到的數字就開心.
不要這麼開心.
其實別人正在用你的錢.
我們自己要學習.
自己去用不要給別人去用.
第十一到十二節提到.
若是你們在不義的錢財上不忠心.
誰還把那真實的錢財託付你們呢?.
如果你們在別人的東西上不忠心.
誰還把你們自己的東西給你們呢?.
錢財很多時候都是用不義的想法賺取.
已經給了定義不義的錢財.
為什麼這麼說呢?.
因為其實在耶穌的門徒當中都有位稅吏.
我們知道稅吏上個星期講過.
稅吏賺錢都是不義的方法.
他們是用不義的方法去賺取錢財.
其實在我們現今這個世代.
都有很多人是用不義的方法去賺錢.
例如詐騙集團.
聽聞最近有位船東.
他想把船駛出去離開荷武之海峽.
於是付錢給伊朗政府.
他去了一個網站交了錢.
誰知道他們搞錯了.
原來那個網站是詐騙集團搞的.
付了錢給詐騙集團.
船仍然不能出去.
除了這些詐騙很明顯之外.
有些黑黑的不是很明顯的.

$^{361}$大家小心查就查到.
例如川普家族.
每逢川普有什麼特別重大的宣佈.
有關戰爭停戰和平之前.
股票市場石油市場一定有大手筆的上落.
這些很多時候牽涉到內幕消息.
用不義的手法賺錢.
十二節有些不是很容易了解.
十一節是講錢財.
十二節是講別人的東西.
都應該是講錢財.
從管家的比喻來看.
別人的東西應該是講主人的錢財.
不是自己的錢財.
不是管家自己的錢財.
如果管家在管理主人的錢財上不忠心.
他就會被炒魷魚.
可能工資也會被主人扣一些.
十三節是一個結論.
我們要做上帝的僕人.
做錢財的主人.
耶穌說.
「一個僕人不能事奉兩個主.
要不是恨者的愛那個.
就是重者的輕那個.
你們不能服侍上帝.
又服侍馬門」.
馬門就是錢的意思.
在《二汗福音》十二章二十六節.
主耶穌對門徒說.
「若有人服侍我.
我父必尊重他」.
但如果我們選擇做錢的奴隸.
錢是不會怎樣尊重我們的.
錢只會奴役我們的.
錢只會叫我們一個星期做七天工.
每天做十幾個小時.
你做八個小時.
快點去找個人兼職.
打兩份工.

$^{401}$賺多點錢.
你已經存了點錢.
我教你.
你快點去投資股市.
半夜教定鬧鐘三分鐘起床.
做什麼.
你要看歐洲美國的股市.
你不能只看澳洲.
要看其他市場的匯市,期貨市場.
昨天有以前教會的會友請我去喝茶.
其他會友的女兒也挺厲害.
去了美國.
讀完書就在美國做事.
說起來她就轉了工作.
以前她好像在投資銀行工作.
那份工作很辛苦.
半夜三點鐘起床看歐洲那邊的股市.
你給錢做你的主人.
你就預備好做錢的奴隸.
而且我提醒大家.
有一天如果你是這樣做錢的奴隸.
可能會讓股市大跌.
令到你心臟病或者精神病.
這些事件曾經在香港發生過.
但因為耶穌的講論.
是要我們明白.
錢只不過是生活上的工具.
而且只是今世有用的.
既然錢是工具.
我們就要用.
不要把錢當成愛的對象.
錢錢我愛你死了.
我們應該愛的是上帝.
是我們的親人和其他人.
不是錢.
我給大家一個問題自己反省.
錢究竟是你的主.
還是耶穌基督是你的主呢?.
我們一起祈禱.
將天上各樣屬靈的福氣賜給我們.

$^{441}$求你幫助我們帶著旅客的心態在世上活著.
幫助我們有機會和別人分享這份.
從天上而來的喜樂和平安.
也堅決做錢財的主人.
做你的僕人.
奉耶穌基督的名祈禱.
\newpage



\section{約翰福音 16:16-24}
\label{sec:032093}
\textbf{憂愁變喜樂}~~~~謝喬遠弟兄
\newline
\newline
連結: \href{http://www.church.com.hk/acms/content.asp?site=cdc\&op=show\&type=product\&code=032093\&layout=sermon}{\texttt{032093}} ~~~~ 語音日期: 2026-05-03
\newline
\newline
\hyperref[sec:032092]{\small{< < < PREV SERMON < < <}}
~
\hyperref[sec:index_chronic]{\small{[返順時目]}}
~
\hyperref[sec:032094]{\small{> > > NEXT SERMON > > >}}
\newline
\newline
約翰福音 16:16-24
\newline
\begin{longtable}{cl}
\hline
\hline
章節 & 經文 (和合本修訂版)\\
\hline
16:16 & \begin{tabularx}{0.7\textwidth}{X} 「不久，你們將不再見到我；再過不久，你們還要見到我。」 \end{tabularx} \\ \\ \relax
16:17 & \begin{tabularx}{0.7\textwidth}{X} 有幾個門徒彼此說：「他對我們說『不久，你們將不再見到我；再過不久，你們還要見到我』；又說『因我到父那裡去』。這是甚麼意思呢？」 \end{tabularx} \\ \\ \relax
16:18 & \begin{tabularx}{0.7\textwidth}{X} 於是門徒說：「他說『不久』到底是甚麼意思呢？我們不明白他說甚麼。」 \end{tabularx} \\ \\ \relax
16:19 & \begin{tabularx}{0.7\textwidth}{X} 耶穌看出他們要問他，就對他們說：「我說『不久，你們將不再見到我；再過不久，你們還要見到我』，你們為這話彼此詢問嗎？ \end{tabularx} \\ \\ \relax
16:20 & \begin{tabularx}{0.7\textwidth}{X} 我實實在在地告訴你們，你們將要痛哭，哀號，世人反要歡喜。你們將要憂愁，然而你們的憂愁要變成喜樂。 \end{tabularx} \\ \\ \relax
16:21 & \begin{tabularx}{0.7\textwidth}{X} 婦人生產的時候會憂愁，因為她的時候到了；但孩子一生出來，就不再記得那痛苦了，因為歡喜有一個人生在世上了。 \end{tabularx} \\ \\ \relax
16:22 & \begin{tabularx}{0.7\textwidth}{X} 你們現在也是憂愁，但我要再見到你們，你們的心就會有喜樂了；這喜樂沒有人能奪去。 \end{tabularx} \\ \\ \relax
16:23 & \begin{tabularx}{0.7\textwidth}{X} 到那日，你們甚麼也不會問我了。我實實在在地告訴你們，你們奉我的名無論向父求甚麼，他會賜給你們。 \end{tabularx} \\ \\ \relax
16:24 & \begin{tabularx}{0.7\textwidth}{X} 直到現在，你們沒有奉我的名求甚麼，如今你們求就必得著，使你們的喜樂得以滿足。」 \end{tabularx} \\ \\
[1ex]
\hline
\hline
\end{longtable}
$^{1}$主席.
非常之清楚.
下一節我們一起讀好不好.
我實實在在地告訴你們.
你們將要痛哭哀號.
世人都要喜樂.
你們將要憂愁.
然而你們的憂愁要變為喜樂.
所以當耶穌每一次說到實實在在的時候.
就是非常之重要.
英文是veryly veryly.
或者是truly truly.
原文是Amen Amen.
所以當耶穌說實實在在的時候.
我們就要注意聽.
所以告訴門徒們聽.
他們會痛哭會哀號.
因為很快門徒就會目睹他們的主耶穌.
被釘十字架.
是非常之血腥的.
所以門徒會痛哭會哀號會憂愁.
後面耶穌告訴他們.
世人都要喜樂.
所以這個世界這個邪惡.
不屬神的世界都要喜樂.
我預備的時候就想起在啟示錄十一章十節裡面說到.
當兩個見證人在末世的時候為神作見證之後.
當他們的使命完結了之後.
這個世界的撒旦就將他們殺害了.
世人不止是歡喜.
他們還互相去送禮物.
所以這個世界.
這麼屬主的世界.
其實就是這樣的邪惡.
而耶穌告訴門徒.
很快要發生什麼事情.
他告訴他們.
門徒會憂愁.
但門徒們的憂愁.
會變為喜樂.

$^{41}$所以耶穌告訴他們.
雖然未來他們不知道會發生什麼事.
但耶穌告訴他們會發生什麼事.
他的意思是.
當門徒有疑問.
他們問會發生什麼事.
會是什麼時候.
等不多時.
耶穌沒有將這個等不多時的答案告訴他們.
耶穌告訴他們後面會發生什麼事.
但耶穌沒有告訴他們那個時間.
我們經常說timing.
耶穌沒有告訴他們這個timing.
因為其實這個timing的事情.
我們也很想知道.
我們找工作.
我們等待一位牧者.
我們很想知道這個timing.
但神沒有給我們這個答案.
因為這個時間.
知道這個時候.
知道這個日期.
對我們來說沒有一個重要性.
甚至對我們來說可能不好.
因為我們不需要知道.
其實一個application.
一個應用.
是我們信心的一個考驗.
那個時候等待.
等待各樣.
等待找工作.
等待另一半.
其實都是一個信心的考驗.
後面耶穌也很清楚地告訴門徒.
他們會憂愁.
但憂愁會變為喜樂.
那會怎樣變呢.
我們一起讀下面的經文.
說了這話.
就把手和肋旁.

$^{81}$致給他們看.
門徒看見主就喜樂了.
所以前面我們今天的經文.
在約翰福音十六章.
下面是在約翰福音二十章二十節.
其實耶穌當他跟門徒說這番話的時候.
過了可能不到四十個小時.
耶穌就會怎樣.
他就會復活.
他從死裡面復活.
所以門徒約翰.
他在二十章二十節記載這段經文.
其實從前面.
是有完全一個印證.
印證在前面.
前面是說憂愁會變為喜樂.
約翰在主所愛的門徒.
他在二十章二十節就告訴大家.
發生了什麼事.
令到他們喜樂.
在路加福音裡面也說到.
他們就是門徒正喜得不敢信.
並且唏噓.
當他們見到復活的主.
他們喜得不敢信.
並且唏噓.
所以當門徒們.
他們見證了十字架的痛苦之後.
經過了兩天.
他們親眼見到復活的主.
他們就喜樂了.
有個牧師叫Chuck Smith.
他說人生有多少的事情.
我們在經歷的當下.
為之憂傷.
只因我們尚未看見最終的結果.
我們很多時候都這樣.
我們在當中.
真的很困難,真的很難過.
但因為我們不知道後面的結果.

$^{121}$而主也將結果告訴了門徒.
一起讀二十一節.
婦人生產的時候就憂愁.
因為她的時候到了.
既生了孩子.
就不再記念那苦楚.
因為歡喜世上生了一個人.
所以耶穌繼續跟門徒解釋.
什麼叫憂愁變喜樂.
耶穌給了一個日常的例子.
就是婦人生產.
我們當中很多姊妹一定都經歷過.
我自己可能比平均的弟兄.
也多了一點點的經驗.
因為我五個小孩.
我進了產房五次.
這件事情我都問過太太.
王靖她坐在後面.
她說確實是這樣.
她進了產房.
尤其是生第二個之後.
她就進去了.
突然問自己.
Why am I here again?.
為什麼我又回來了.
因為生產確實是很痛苦.
我們知道因為在創世記律三章.
當我們刺倒犯罪之後.
生產就變成一個非常痛苦的事情.
確實是當婦人在產房經歷完之後.
當孩子出來了之後.
她就完全不記得前面的苦楚.
不再記念那苦楚.
因為什麼呢.
因為歡喜世上生了一個人.
因為多了一個Baby.
其實在《爾賽雅書》的舊約當中.
很多很多的神都用生產和難產來形容.
以色列人在罪裡面的痛苦.
我不會一節一節去讀.

$^{161}$所以耶穌在這段話的時候.
再一次離開他的被捕.
離開他的釘十字架.
已經剩下幾個小時.
我們的門徒會經歷.
會見證.
會看到十字架上的痛苦和傷害.
但耶穌也知道.
當他經歷了十字架之後.
後面會發生什麼事情.
我們一起來讀下面那段經文.
在希伯來書十二章十二節一起讀.
仰望為我們信心創始成中的耶穌.
他因那擺在前面的喜樂.
就輕看羞辱.
忍受了十字架的腐爛.
便坐在神寶座的右邊.
所以今天我們主餐的主日.
我們一起領受了杯和餅.
來紀念耶穌.
祂的身體為我們破碎.
祂的寶血為我們流出來.
耶穌在《約翰福音》十六章的時候.
祂很清楚擺在祂前面的是什麼.
祂也求父.
如果可以的話.
把這個福杯可以拿去.
但父神因為愛我們.
因為他救贖的計劃.
他沒有把這個福杯拿去.
耶穌自己也很清楚.
當他忍受了這個痛苦之後.
後面會發生什麼事情.
他忍受了十字架的腐爛.
他之後因為擺在前面的喜樂.
再一次講到這個喜樂.
他就坐在神寶座的右邊.
他就好像確實是婦人生產幽壽.
過了之後.
他就不紀念她的苦處.

$^{201}$因為歡喜世上生了一個人.
耶穌當時也一樣.
祂因為擺在前面的喜樂.
祂知道祂通過十字架的腐爛.
祂會成就什麼事情.
The birth of the church.
祂會將教會生下來.
所以祂忍受羞辱.
忍受十字架的腐爛.
再讀二十二節.
你們現在也是幽壽.
但我要再見你們.
你們的心就喜樂了.
這喜樂也沒有人能奪去.
耶穌說你們現在.
十字架之前在最後的晚餐.
耶穌馬上要被捕.
門徒們也是幽壽.
但過後呢.
祂說.
我要再見你們.
很有意思.
在前面祂說.
等不多時你們不再見我.
門徒不會見到耶穌.
過不多時他們會再見到耶穌.
到了二十二節這裡.
耶穌說我要再見你們.
耶穌自己知道他復活之後.
他要再見到門徒.
在彼得前書.
彼得將經過很多年之後.
他說你們雖然沒有見過他.
見過耶穌.
但是愛他.
如今雖不能看見.
卻因信他又說不出來.
滿有榮光的大喜樂.
所以彼得說到.
我們其實每一個都一樣.

$^{241}$我們雖然沒有親身見過耶穌.
但我們愛他.
我們雖然現在見不到他.
但因為信他又說不出來.
滿有榮光的大喜樂.
這個喜樂是什麼呢?.
耶穌也說這個喜樂是.
沒有人能夠奪去的喜樂.
當然這個世界上有很多什麼.
能夠奪去的喜樂.
在應用的時候會說到.
什麼是不能奪去的喜樂.
這個喜樂是怎樣來的.
亦都是門徒們.
因為他們會經過非常悲慘.
非常憂愁的時間.
但後來他們會喜樂.
很有意思的是當耶穌復活了之後.
耶穌復活了之後.
到福音書的最後.
到使徒行傳.
其實很少門徒.
很少再提到十字架的痛苦.
為什麼呢?.
因為已經過去了.
施佩珍(Charles Burdett)說.
No grief mentioned after by the disciples.
因為耶穌完全復活了之後.
過去的事情已經過去了.
門徒不會再.
Dwell on.
不會再沉醉在痛苦當中.
因為耶穌已經復活了.
所以復活當然.
是我們基督徒信仰的核心.
非常之重要的一個.
我們相信不只是我們靈的復活.
我們身體的復活.
就好像耶穌一樣.
耶穌復活了之後.

$^{281}$祂可以吃東西,吃魚.
我們都會期待一個完全身體的復活.
其實這是一個見證.
我剛剛信主的時候.
對我來說也是一個很大的激勵.
我其實是一個非常理智的人.
在福音書裡面看到.
當門徒.
當耶穌被捕的時候.
門徒全部都走了.
全部都走了.
馬可福音應該是他自己形容自己.
他連衣服都不要.
丟下了就跑了.
因為他們很害怕.
所以我當時信主之後也在想.
如果耶穌沒有復活的話.
會怎樣呢?.
門徒多半都走了.
自己回去打魚.
做各種各樣的事情.
沒有後來的事情.
也不會有基督教.
就是因為耶穌真的復活了.
所以十二個門徒.
所有的門徒他們之後是怎樣的呢?.
他們每一個都基本上是殉道.
什麼樣的力量.
可以令他們為了.
可以擺上他們的生命.
就是因為他們每一個都看到復活的主.
他們很清楚主真的復活了.
也不只是耶穌的使徒.
他們的門徒.
之後歷世歷代的教會.
還有這麼多的基督徒和宣教士.
可以為主可以為這個信仰.
擺上他們的生命.
因為我們相信是復活了的主.
二十三節一起讀.

$^{321}$對不起.
二十三節.
你們什麼也不問我了.
我實實在在地告訴你們.
你們若向父求什麼.
他也必因我的名賜給你們.
所以耶穌繼續告訴他們.
到那一天.
再過兩三天.
你們什麼也不再問我了.
我實實在在地告訴你們.
所以再一次實實在在地告訴門徒.
你們若向父求什麼.
他也必因我的名賜給你們.
所以記得門徒.
之前有很多的疑問.
他們彼此的相問.
我們剛才說到耶穌是歡迎他們的問題.
但是耶穌告訴他們.
到那一天.
門徒們就什麼也不再問他了.
為什麼不問呢?.
因為他們已經有答案了.
他們見到復活的主.
他們一下子明白了很多事情.
耶穌復活了之後.
門徒們也沒有那麼多的問題.
到了使徒行傳之後.
當聖靈降臨之後.
因為有神持續的引領.
又有聖靈降臨在信徒的心裡面.
所以門徒們也沒有那麼多的問題.
在《瑪竇福音》七章七字都說到.
你們祈求就給你們.
尋求就尋見.
叩門就給你們開門.
所以門徒們.
他們有很多的問題都得到回答.
所以門徒們他們有很多的問題都得到回答.
二十四字.

$^{361}$向來你們沒有奉我的名求什麼.
如今你們求.
就必得著.
叫你們的喜樂.
可以滿足.
耶穌繼續向我們解釋.
他說奉我們的名求.
我們其實禱告.
我們都會奉耶穌的名.
我們每次禱告.
都會奉耶穌的名.
所以這不是一個套話.
也不是一個傳統.
我們本地的牧師John Piper很有名.
他說如果我們使徒.
不靠著耶穌來到神面前.
你會被焚燒.
不要嘗試.
英文大概的原話是.
If you try to come to God.
without Jesus you will be incinerated.
Don't try.
所以我們來到耶穌.
來到神的面前來到父神的面前.
是完全因為耶穌所成就.
其實今天的主題是.
耶穌因為十字架所成就.
我們可以坦然無懼的來到神的面前.
不要嘗試.
靠著自己的義行.
來到神的面前.
是沒有可能的.
大陸也學我們.
功氣發財身體健康.
四個字這樣說.
我小時候就聽.
後面還有一句叫萬事如意.
這個這麼厲害.
如果萬事如意的話.
我前面其實也不用再祝福了.

$^{401}$還有四個更厲害的.
叫做心想事成.
如果我們心想事成.
那真是不得了.
其實我自己在想.
心想事成.
所以耶穌說我們求.
就必得著.
是不是我們世俗說萬事如意.
心想事成.
當然不是.
我們後面會解釋究竟是什麼意思.
其實這個世界我們熟悉的都知道.
如果我們有罪的人.
我們是萬事如意.
心想事成的話.
這個世界會變成怎樣.
非常之可怕.
世界可能已經沒有了.
有罪的人如果心想事成的話.
我們就更恐怖.
耶穌也說叫我們的喜樂.
可以滿足.
joy be full.
joy be complete.
後面有另外.
我們今天讀的是約翰福音十六章.
其實十六十七章.
都有很類似的.
耶穌的吩咐.
我們一起讀中間那段經文.
十五章.
你們若遵守我的命令.
就常在我的外裡.
正如我遵守了我父的命令.
常在他的外裡.
這些事我已經對你們說了.
是要叫我的喜樂.
存在你們心裡.
並叫你們的喜樂可以滿足.

$^{441}$所以耶穌告訴門徒.
告訴我們.
怎樣我們的喜樂可以滿足.
怎樣我們的joy可以complete.
可以full.
就是我們要遵守神的命令.
遵守神的命令.
遵守神的話.
下面最後一段一起讀.
現在我往你那裡去.
我還在西上說這話.
是要叫他們.
心裡充滿我的喜樂.
就是最終我們不能奪去的喜樂.
是誰的喜樂.
是耶穌的喜樂.
是主的喜樂.
是神的喜樂.
所以下面我們有兩個.
的應用.
亞國書說我們.
看完神的話.
不要好像照鏡子然後就轉身.
忘記了.
當我們明白了神的話.
我們將會應用在我們自己身上.
就是在二十四節說.
向來你們沒有奉我的命求什麼.
如今你們求.
就必得著.
叫你們的喜樂可以滿足.
下面一節十四章.
十三節我們一起讀.
你們奉我的命.
無論求什麼.
我必成就.
叫父因而子得榮耀.
所以我們怎樣奉耶穌的命.
怎樣求怎樣就必得著.
這個禱告是怎樣的禱告呢.

$^{481}$在約翰福音十四章.
十三節耶穌告訴我們.
我們奉祂的命無論求什麼.
耶穌會成就.
叫父因而子.
得榮耀.
所以當我們禱告是以神為中心.
是叫神的兒子耶穌.
得榮耀的話.
這樣的禱告神必定會答應.
下面最後一節經文一起讀.
你們若常在我裡面.
我的壞事常在你們裡面.
凡你們所願意的祈求.
就給你們成就.
所以我們怎樣禱告呢.
我們怎樣求我們怎樣禱告神.
我們求神就必定會答應.
我們必定會得著.
十五章這裡七節告訴我們.
我們要常在神的裡面.
常在主裡面.
常在主的話裡面.
當我們常常在祂的話.
常常被神的話充滿的話.
神必定會通過祂的話來工作.
當我們禱告是充滿.
被神的話飽和的話.
神必定會回應我們的禱告.
當然這一段我們今天到十六章這裡.
因為耶穌馬上要上十字架.
門徒直著耶穌.
馬上要成就的事情.
就是十字架上的死.
可以無限地.
沒有任何東西可以有權阻擋.
直接來到神面前.
當耶穌在十字架上.
祂死的時候.
祂說什麼.

$^{521}$祂說成了完成了.
救恩完成了.
當時萬子怎樣.
從上到下裂開.
其實神告訴我們.
再一次人可以直接到.
神的面前.
因為耶穌所成就.
第二個應用.
我們的憂愁如何變成喜樂.
我們憂不憂愁.
憂愁.
有姊妹很大聲說憂愁.
她跟我說.
一定要在台上笑.
帶大家喜樂.
所以在這段短短的經文裡.
憂愁出現了四次.
喜樂這個詞出現了六次.
所以我們查經讀經都知道.
如果這個詞不停地重複的話.
在這一段神話裡.
非常重要的關鍵字.
所以憂愁.
怎樣變為喜樂呢.
其實生活中.
非常多的憂愁.
說到憂愁.
說到喜樂.
不能不提苦難.
人生中大大小小都經歷過.
各種各樣的苦難.
當我們說到憂愁變成喜樂的時候.
我們不是縮小任何的苦難.
不是減少任何的苦難.
在苦難當中.
我們很多時候有沒有答案.
其實我們沒有答案.
幾個月前香港大暴一場大火.
我們每個人都沒有答案.

$^{561}$完全沒有答案.
當我們自己沒有答案.
亦都不可能幫別人找答案.
更加不可以幫神回答.
我很多年前看約伯記的時候.
約伯記提到苦難的時候.
當約伯在苦難的時候.
他的朋友來找他.
我們常說不要做約伯的朋友.
其實你看約伯的三個朋友.
他們開始的時候.
他們做得非常之好.
他們從很遠很遠的地方來到.
因為知道約伯的苦難.
他們從很遠的地方看到約伯.
他們就撕裂他們的衣服.
將灰塞在頭上.
因為他們和約伯在一起.
亦都非常哀痛.
因為看到約伯的光景.
他們都非常之難過.
他們都七日七夜.
在約伯的旁邊陪著他.
連一句話都沒有說.
問題其實出在後面.
後面他們忍不住.
他們要說話.
其實對我們都一樣.
我們弟兄姊妹朋友有苦難的時候.
其實很多時候他們不需要我們.
給他們什麼答案.
我們都不可能給他們答案.
我們可以陪他們.
我相信約伯的朋友都會為他們禱告.
我們可以為弟兄姊妹.
為苦難中的肢體.
為苦難中的朋友 鄰居.
家人來禱告.
聖經裡面.
神亦都告訴我們.

$^{601}$我們的苦難.
我幾次流離.
你都記數.
求你擺我眼淚.
將在你的皮袋裡.
這筆都記在你的冊子上.
所以主知道我們的苦難.
祂特別有一個皮袋.
是裝我們的眼淚.
祂特別有一個冊子.
在天上是裝我們的哀傷.
讓我們知道當我們在苦難當中.
神是知道.
神是和我們同行.
當然.
我們憂愁怎麼變喜樂.
我們剛才說到不能奪去的喜樂.
這個世界上很多很多的喜樂.
只要不是在主裡面.
都會被奪去.
在我們生活中.
就算我們信了主之後.
追求的很多東西.
這個世界上的學業.
工作 綠卡 身份.
升級 加薪 錢財.
男朋友 女朋友 婚姻.
孩子 這些東西好不好.
這些東西都可以好好.
都可以有神的祝福.
但如果是變成我們心思意念.
全部時間精力所思想的話.
這些都會變成.
我們能奪去的喜樂.
很多時候我們很想得一樣東西.
最空虛就是得了之後.
一直想著要升級加薪.
升級加薪你開心多久.
開心一陣.
所以都是這個世界上.

$^{641}$能夠奪去的喜樂.
那什麼喜樂是不能夠奪去的呢.
什麼喜樂是永恆的呢.
其實基督徒都知道.
面對生死的時候.
我們要知道.
什麼的喜樂是不能奪去的.
我預備的時候.
看到一段說話.
一個茶經家叫 Matthew Henry.
他說 因此對於我們在主內的弟兄姊妹.
我們可以如此說.
等不多時.
我們將不再見到他們.
不時他們離開我們.
就是我們離開他們.
這不久必須分別.
但可以確定的是.
這不是永久的分離.
這不過是暫時的廣曼安.
對那些我們盼望在清晨的喜樂當中.
再見的人而言.
相信我們都經歷過.
我們的親人.
離世的話我們都很難過.
但我們知道.
等不多時.
只要他們在主內的話.
我們必然會再見到他們.
最後我再讀一次這一節.
你們現在也是憂愁.
但我要再見你們.
你們的心就喜樂了.
這喜樂也沒有人能奪去.
這喜樂是沒有人能奪去的喜樂.
是神裡面的喜樂.
因為我們神已經是一位復活的主.
所以我們可以非常喜樂.
就好像門徒一樣.
當我們見到復活的主.

$^{681}$他們的心就喜樂.
我們都一樣.
我們在生活當中.
有各樣各樣大小的難處.
但我們知道主因他活著.
所以我們可以喜樂.
我們一起禱告.
主對我們說.
復活在你 生命也在你.
信你的人雖然死了.
也必復活.
凡活著 信你的人必永遠不死.
主我們信你的話.
主將我們每一位.
將孩子放在你的手中.
求你將我們在每一天生活的時候.
意識到.
明白到你是復活.
和我們同行的主.
我們將你身邊的朋友.
在我們當中.
我們相信你的朋友.
交到你的手中.
求你的靈打開他們的眼睛 打開他們的心.
讓他們明白.
復活在於你 生命也在於你.
是奉我主耶穌得勝的名義祈求.
Amen.
\newpage



\section{撒母耳記上 1:1-2:36}
\label{sec:032094}
\textbf{一位母親的奉獻}~~~~梁廷益牧師
\newline
\newline
連結: \href{http://www.church.com.hk/acms/content.asp?site=cdc\&op=show\&type=product\&code=032094\&layout=sermon}{\texttt{032094}} ~~~~ 語音日期: 2026-05-03
\newline
\newline
\hyperref[sec:032093]{\small{< < < PREV SERMON < < <}}
~
\hyperref[sec:index_chronic]{\small{[返順時目]}}
~
\hyperref[sec:032095]{\small{> > > NEXT SERMON > > >}}
\newline
\newline
撒母耳記上 1:1-2:36
\newline
\begin{longtable}{cl}
\hline
\hline
章節 & 經文 (和合本修訂版)\\
\hline
1:1 & \begin{tabularx}{0.7\textwidth}{X} 以法蓮山區有一個拉瑪的瑣非人，名叫以利加拿，他是蘇弗的玄孫，託戶的曾孫，以利戶的孫子，耶羅罕的兒子，是以法蓮人。 \end{tabularx} \\ \\ \relax
1:2 & \begin{tabularx}{0.7\textwidth}{X} 他有兩個妻子：一個名叫哈拿，另一個名叫毗尼拿。毗尼拿有孩子，哈拿卻沒有孩子。 \end{tabularx} \\ \\ \relax
1:3 & \begin{tabularx}{0.7\textwidth}{X} 這人每年從本城上到示羅，敬拜萬軍之耶和華，向他獻祭。在那裡有以利的兩個兒子何弗尼和非尼哈當耶和華的祭司。 \end{tabularx} \\ \\ \relax
1:4 & \begin{tabularx}{0.7\textwidth}{X} 每逢獻祭的日子，以利加拿把祭肉分給他的妻子毗尼拿和毗尼拿所生的兒女。 \end{tabularx} \\ \\ \relax
1:5 & \begin{tabularx}{0.7\textwidth}{X} 他給哈拿的卻是雙分，因為他愛哈拿。耶和華卻不使哈拿生育。 \end{tabularx} \\ \\ \relax
1:6 & \begin{tabularx}{0.7\textwidth}{X} 她的對頭毗尼拿因耶和華不使哈拿生育，就常常惹她發怒，要使她生氣。 \end{tabularx} \\ \\ \relax
1:7 & \begin{tabularx}{0.7\textwidth}{X} 年年都是如此。每當她上到耶和華殿的時候，毗尼拿就這樣惹她發怒，以致她哭泣不吃飯。 \end{tabularx} \\ \\ \relax
1:8 & \begin{tabularx}{0.7\textwidth}{X} 她丈夫以利加拿對她說：「哈拿，你為何哭泣？為何不吃飯？為何傷心難過呢？有我不比有十個兒子更好嗎？」 \end{tabularx} \\ \\ \relax
1:9 & \begin{tabularx}{0.7\textwidth}{X} 他們在示羅吃喝完了，哈拿就站起來。祭司以利坐在耶和華殿門框旁邊的位子上。 \end{tabularx} \\ \\ \relax
1:10 & \begin{tabularx}{0.7\textwidth}{X} 哈拿心裡愁苦，就痛痛哭泣，向耶和華祈禱。 \end{tabularx} \\ \\ \relax
1:11 & \begin{tabularx}{0.7\textwidth}{X} 她許願說：「萬軍之耶和華啊，你若垂顧你使女的苦情，眷念不忘你的使女，賜你的使女一個子嗣，我必使他終生歸給耶和華，不用剃刀剃他的頭。」 \end{tabularx} \\ \\ \relax
1:12 & \begin{tabularx}{0.7\textwidth}{X} 哈拿在耶和華面前不住地祈禱，以利注意她的嘴。 \end{tabularx} \\ \\ \relax
1:13 & \begin{tabularx}{0.7\textwidth}{X} 哈拿心中默禱，只動嘴唇，聽不到她的聲音，因此以利以為她喝醉了。 \end{tabularx} \\ \\ \relax
1:14 & \begin{tabularx}{0.7\textwidth}{X} 以利對她說：「你要醉到幾時呢？不要再喝酒了！」 \end{tabularx} \\ \\ \relax
1:15 & \begin{tabularx}{0.7\textwidth}{X} 哈拿回答說：「我主啊，不是這樣。我是心裡愁苦的婦人，清酒烈酒都沒有喝，只在耶和華面前傾心吐意。 \end{tabularx} \\ \\ \relax
1:16 & \begin{tabularx}{0.7\textwidth}{X} 不要將你的使女看作不正經的女子。我因極其難過和生氣，所以一直禱告到如今。」 \end{tabularx} \\ \\ \relax
1:17 & \begin{tabularx}{0.7\textwidth}{X} 以利回答說：「平平安安地回去吧。願以色列的神允准你向他所求的！」 \end{tabularx} \\ \\ \relax
1:18 & \begin{tabularx}{0.7\textwidth}{X} 哈拿說：「願你的婢女在你眼前蒙恩。」於是婦人上路，去吃飯，臉上不再帶愁容了。 \end{tabularx} \\ \\ \relax
1:19 & \begin{tabularx}{0.7\textwidth}{X} 他們清早起來，在耶和華面前敬拜，就回去，往拉瑪自己的家裡。以利加拿和妻子哈拿同房，耶和華顧念哈拿。 \end{tabularx} \\ \\ \relax
1:20 & \begin{tabularx}{0.7\textwidth}{X} 時候到了，哈拿懷孕生了一個兒子，給他起名叫撒母耳，說：「這是我從耶和華那裡求來的。」 \end{tabularx} \\ \\ \relax
1:21 & \begin{tabularx}{0.7\textwidth}{X} 以利加拿和他全家都上去，要向耶和華獻年祭和還願祭。 \end{tabularx} \\ \\ \relax
1:22 & \begin{tabularx}{0.7\textwidth}{X} 哈拿卻沒有上去，因為她對丈夫說：「等孩子斷了奶，我就帶他上去朝見耶和華，讓他永遠住在那裡。」 \end{tabularx} \\ \\ \relax
1:23 & \begin{tabularx}{0.7\textwidth}{X} 她丈夫以利加拿對她說：「就照你看為好的去做吧！可以留到兒子斷了奶，願耶和華應驗他的話。」於是婦人留在家裡乳養兒子，直到他斷了奶。 \end{tabularx} \\ \\ \relax
1:24 & \begin{tabularx}{0.7\textwidth}{X} 斷奶之後，她就帶著孩子，連同一頭三歲的公牛，一伊法細麵，一皮袋酒，上示羅耶和華的殿去。那時，孩子還小。 \end{tabularx} \\ \\ \relax
1:25 & \begin{tabularx}{0.7\textwidth}{X} 他們宰了公牛，就領孩子到以利面前。 \end{tabularx} \\ \\ \relax
1:26 & \begin{tabularx}{0.7\textwidth}{X} 婦人說：「我主啊，請容許我說，我向你，我的主起誓，從前在你這裡站著祈求耶和華的那婦人就是我。 \end{tabularx} \\ \\ \relax
1:27 & \begin{tabularx}{0.7\textwidth}{X} 我祈求為要得這孩子，耶和華已將我向他所求的賜給我了。 \end{tabularx} \\ \\ \relax
1:28 & \begin{tabularx}{0.7\textwidth}{X} 所以，我將這孩子獻給耶和華，使他終生歸給耶和華。」他就在那裡敬拜耶和華。 \end{tabularx} \\ \\ \relax
2:1 & \begin{tabularx}{0.7\textwidth}{X} 哈拿禱告說：「我的心因耶和華快樂，我的角因耶和華高舉。我的口向仇敵張開；我因你的救恩歡欣。 \end{tabularx} \\ \\ \relax
2:2 & \begin{tabularx}{0.7\textwidth}{X} 「沒有一位聖者像耶和華，除你以外沒有別的了，也沒有磐石像我們的神。 \end{tabularx} \\ \\ \relax
2:3 & \begin{tabularx}{0.7\textwidth}{X} 不要誇口說驕傲的話，也不要口出狂妄的言語，因耶和華是有知識的神，人的行為被他衡量。 \end{tabularx} \\ \\ \relax
2:4 & \begin{tabularx}{0.7\textwidth}{X} 勇士的弓折斷，跌倒的人以力量束腰。 \end{tabularx} \\ \\ \relax
2:5 & \begin{tabularx}{0.7\textwidth}{X} 飽足的人作雇工求食；飢餓的人也不再飢餓。不生育的生了七個；兒女多的反倒孤獨。 \end{tabularx} \\ \\ \relax
2:6 & \begin{tabularx}{0.7\textwidth}{X} 耶和華使人死，也使人活，使人下陰間，也使人往上升。 \end{tabularx} \\ \\ \relax
2:7 & \begin{tabularx}{0.7\textwidth}{X} 耶和華使人貧窮，也使人富足；使人降卑，也使人升高。 \end{tabularx} \\ \\ \relax
2:8 & \begin{tabularx}{0.7\textwidth}{X} 他從灰塵裡抬舉貧寒人，從糞堆中提拔貧窮人，使他們與貴族同坐，繼承榮耀的座位。地的柱子屬耶和華，他將世界立在其上。 \end{tabularx} \\ \\ \relax
2:9 & \begin{tabularx}{0.7\textwidth}{X} 「他必保護他聖民的腳步，但惡人卻在黑暗中毀滅，因為人不是靠力量得勝。 \end{tabularx} \\ \\ \relax
2:10 & \begin{tabularx}{0.7\textwidth}{X} 與耶和華相爭的，必被打碎；他必從天上打雷攻擊他們。耶和華審判地極的人，將力量賜給所立的王，高舉受膏者的角。」 \end{tabularx} \\ \\ \relax
2:11 & \begin{tabularx}{0.7\textwidth}{X} 以利加拿往拉瑪回自己的家去了。那孩子在以利祭司面前事奉耶和華。 \end{tabularx} \\ \\ \relax
2:12 & \begin{tabularx}{0.7\textwidth}{X} 以利的兩個兒子是無賴，不認識耶和華。 \end{tabularx} \\ \\ \relax
2:13 & \begin{tabularx}{0.7\textwidth}{X} 這二祭司對待百姓的規矩是這樣：凡有人獻祭，正煮肉的時候，祭司的僕人就手拿三齒的叉子來， \end{tabularx} \\ \\ \relax
2:14 & \begin{tabularx}{0.7\textwidth}{X} 將叉子往盆裡，或鍋裡，或釜裡，或壺裡一插，插上來的肉，祭司都拿了去。他們對所有上到示羅的以色列人都這樣做。 \end{tabularx} \\ \\ \relax
2:15 & \begin{tabularx}{0.7\textwidth}{X} 甚至在未燒脂肪之前，祭司的僕人就來對獻祭的人說：「把肉給祭司，讓他烤吧。他不要拿你煮過的肉，要生的。」 \end{tabularx} \\ \\ \relax
2:16 & \begin{tabularx}{0.7\textwidth}{X} 獻祭的人若說：「他們必須先燒脂肪，然後你才可以隨意拿。」僕人就說：「不，你立刻給我，不然我就要搶了。」 \end{tabularx} \\ \\ \relax
2:17 & \begin{tabularx}{0.7\textwidth}{X} 這些年輕人的罪在耶和華面前非常嚴重，因為這些人藐視耶和華的祭物。 \end{tabularx} \\ \\ \relax
2:18 & \begin{tabularx}{0.7\textwidth}{X} 那時，撒母耳還是孩子，穿著細麻布的以弗得，侍立在耶和華面前。 \end{tabularx} \\ \\ \relax
2:19 & \begin{tabularx}{0.7\textwidth}{X} 他母親每年為他做一件小外袍，同丈夫上來獻年祭的時候帶來給他。 \end{tabularx} \\ \\ \relax
2:20 & \begin{tabularx}{0.7\textwidth}{X} 以利為以利加拿和他妻子祝福，說：「願耶和華由這婦人再賜你後裔，代替從耶和華求來的孩子。」他們就回自己的地方去了。 \end{tabularx} \\ \\ \relax
2:21 & \begin{tabularx}{0.7\textwidth}{X} 耶和華眷顧哈拿，她就懷孕生了三個兒子，兩個女兒。那孩子撒母耳在耶和華面前漸漸長大。 \end{tabularx} \\ \\ \relax
2:22 & \begin{tabularx}{0.7\textwidth}{X} 以利年紀老邁，聽見他兩個兒子對以色列眾人所做一切的事，又聽見他們與會幕門前伺候的婦人同寢， \end{tabularx} \\ \\ \relax
2:23 & \begin{tabularx}{0.7\textwidth}{X} 就對他們說：「你們為何做這樣的事呢？我從這眾百姓聽見了你們的惡行。 \end{tabularx} \\ \\ \relax
2:24 & \begin{tabularx}{0.7\textwidth}{X} 我兒啊，不可這樣！我聽到耶和華的百姓傳出你們不好的名聲。 \end{tabularx} \\ \\ \relax
2:25 & \begin{tabularx}{0.7\textwidth}{X} 人若得罪人，有神可以裁決；人若得罪耶和華，誰能為他代求呢？」然而他們還是不聽父親的話，因為耶和華想要他們死。 \end{tabularx} \\ \\ \relax
2:26 & \begin{tabularx}{0.7\textwidth}{X} 撒母耳這孩子漸漸長大，耶和華與人越發喜愛他。 \end{tabularx} \\ \\ \relax
2:27 & \begin{tabularx}{0.7\textwidth}{X} 有神人來見以利，對他說：「耶和華如此說：『你祖宗的家在埃及法老家的時候，我不是向他們顯現嗎？ \end{tabularx} \\ \\ \relax
2:28 & \begin{tabularx}{0.7\textwidth}{X} 在以色列眾支派中，我不是揀選他作我的祭司，上我的祭壇，燒香，在我面前穿以弗得嗎？我不是將以色列人所獻的火祭都賜給你祖宗的家嗎？ \end{tabularx} \\ \\ \relax
2:29 & \begin{tabularx}{0.7\textwidth}{X} 你們為何踐踏我所吩咐獻在我居所的祭物和供物呢？你為何尊重你的兒子過於尊重我，將我百姓以色列所獻美好的祭物都拿去養肥你們自己呢？』 \end{tabularx} \\ \\ \relax
2:30 & \begin{tabularx}{0.7\textwidth}{X} 因此，耶和華－以色列的神說：『我確實說過，你和你祖宗的家必永遠行在我面前，但現在耶和華卻說，我絕不會這樣做。因為尊重我的，我必尊重他；藐視我的，他必被輕視。 \end{tabularx} \\ \\ \relax
2:31 & \begin{tabularx}{0.7\textwidth}{X} 看哪，日子將到，我要折斷你的膀臂和你祖宗家的膀臂，使你家中沒有一個老年人。 \end{tabularx} \\ \\ \relax
2:32 & \begin{tabularx}{0.7\textwidth}{X} 你在以色列人享福的時候必看見我居所的衰敗，你家中必永遠沒有一個老年人。 \end{tabularx} \\ \\ \relax
2:33 & \begin{tabularx}{0.7\textwidth}{X} 你家中的人，我沒有從我壇前剪除的，必使你眼睛失明，心中憂傷。你家中所添的人口都必夭折。 \end{tabularx} \\ \\ \relax
2:34 & \begin{tabularx}{0.7\textwidth}{X} 你的兩個兒子何弗尼、非尼哈所遭遇的事是給你的預兆：他們二人必在同一日死亡。 \end{tabularx} \\ \\ \relax
2:35 & \begin{tabularx}{0.7\textwidth}{X} 我要為自己立一個忠心的祭司，他行事必照我的心、如我的意。我要為他建立堅固的家，他必天天行走在我受膏者的面前。 \end{tabularx} \\ \\ \relax
2:36 & \begin{tabularx}{0.7\textwidth}{X} 你家所剩下的人都必來叩拜他，求一塊銀子，一個餅，說：求你給我一個祭司的職分，好使我得點餅吃。』」 \end{tabularx} \\ \\
[1ex]
\hline
\hline
\end{longtable}
$^{1}$今日的講道題目是一位母親的奉獻.
快到母親節了.
我特別選了這段經文.
跟大家分享一位母親的奉獻.
今天是講奉獻.
但我用一位母親的奉獻.
來作為一個見證.
我們來看看聖經裡的經文.
耶和華卻不使哈拿生育 他的對頭比利拿因.
耶和華不使哈拿生育 就常常惹他發怒.
大家知道當時的文化.
這個人有兩個太太.
哈拿沒有嬰兒.
那個所謂的二奶.
常常罵他.
常常使他發怒.
使他生氣 每一年都是如此.
每當他上到耶和華殿的時候.
比利拿就惹他發怒.
以致他哭泣不吃飯.
令他連飯都吃不下.
他丈夫伊利加拿對他說.
哈拿 你為何哭泣不吃飯.
為何傷心難過呢.
有我不比有十個兒子更好嗎.
幾千年前的男人都不懂得取悅.
不懂得安慰.
老公 你的頭髮掉了.
沒關係 你有我就行了.
都是這樣.
老公不懂得說話 不懂得安慰.
哈拿心裡就愁苦.
聽完老公更加沒有安慰.
於是痛痛哭泣 向耶和華祈禱.
他許願說 萬君子耶和華.
你若垂過你侍女的苦情 眷念不忘你的侍女.
賜你的侍女一個子嗣.
若我必使她終生歸給耶和華.
不用剃刀剃她的頭.
他們稱道起來的時候.

$^{41}$耶和華面前敬拜就回去.
往拉瑪自己的家裡.
伊利加拿和妻子哈拿同房.
耶和華顧念哈拿 事後到了哈拿還生了一個兒子.
給他起名叫薩姆爾.
說 這是我從耶和華那裡求來的.
伊利加拿和他的全家都想去.
要向耶和華獻念祭和還願祭.
我祈求為要得著孩子.
耶和華已將我向他所求的賜給我了.
所以我將這孩子獻給耶和華.
使他終生歸給耶和華.
他就留在那裡敬拜耶和華.
那時候薩姆爾還生孩子.
穿著細麻衣的伊伏德.
侍納在耶和華面前.
耶和華眷顧哈拿.
他就還人生了三個兒子.
兩個女兒 因為不是沒有生.
是一生一鬥.
那孩子薩姆爾在耶和華面前漸漸長大.
這就是我們今天讀的一段經文.
大家知道我讀聖經都是希望從聖經裡的華語裡.
學習一下神要我們學習一下功課.
每逢講到關於奉獻的題目.
我都會用這個例子.
如果你以前聽過我講關於奉獻.
我都喜歡用這個例子.
就是一些前塵的往事.
這對於我來說是一個很深刻的.
有很多事跡都是和我的家庭一起的時候.
你知道現在的講道 講章 很多的牧者都可以找AI做.
但是我可以告訴你.
這些是你自己的生命的經歷.
不能由AI做出來.
一些生命裡面的經歷.
我現在有兩個女兒.
一個三十歲 一個二十六歲.
但是以前小時候.
我們一家人都有很多的活動.

$^{81}$那個是我.
女兒還很小的時候.
我們回Texas探望自己的母校.
這個就是講到麥當勞的故事.
有些人可能聽過.
再聽一下.
小時候在我家裡住.
我住在香港的馬鞍山.
馬鞍山新港城樓下.
有一家麥當勞.
所以我很多時候都會帶女兒去逛街.
或者回來.
我們沒有時間煮東西吃.
所以就去吃麥當勞的兒童餐.
因為兒童餐有玩具送.
小朋友一定喜歡的.
喜歡玩具.
但是也喜歡吃麥當勞的食物.
他們吃麥當勞.
我就每人買了一個餐.
我當然是吃Big Mac.
他們兩個就吃了兒童餐.
兒童餐很快就吃完那個麵包.
很小的.
就慢慢地吃薯條.
很快就吃完我的餐.
我看著他們兩個慢慢地吃薯條.
於是我就向大女兒打了注意.
我就看著她的薯條.
於是我就伸了手.
想拿她Tray上的薯條來吃.
於是大女兒也一貫作風.
大女兒就馬上用兩隻手掩著薯條.
她說「喂,是我的」.
「是我的」.
「Don't touch」.
於是我就說.
「你不給我就向小女兒打了注意」.
小女兒還小嘛.
想著她不太懂.

$^{121}$於是我就拿了她的薯條.
於是小女兒就說.
「不要碰,Hands off」.
就像她姐姐那樣.
用兩隻手掩著薯條.
不讓我去拿.
這件事對我來說有很深刻的印象.
因為這件事其實就是.
講解我們每一個人對上帝的公認的看法.
兩個小女孩吃薯條.
我只是拿了一條而已.
她們兩個的薯條都是我給她的.
主權是我的.
是我給她的.
沒錯我送了給你.
但我送給你之後.
我只是叫你給我一條而已.
上帝給了我們這麼多東西.
上帝說給十分之一來.
你說是我的.
我不想給你.
不明白這件事背後.
是你爸爸給你的.
不是你自己賺錢去買的.
第一,第二.
我心想.
兩個死美女.
你知不知道我比你大.
我一隻手把你們兩個推開.
我就拿了你的薯條.
我強而有力過你.
如果我要奪取你的薯條.
我絕對有能力這樣做.
我不這樣去做.
是因為我給你們恩典.
我想你們享受.
我想和你們一起享受.
為什麼這麼斤斤計較.
這麼吝嗇呢.
如果我一發火.

$^{161}$全部都沒得吃.
我自己吃完.
我絕對可以這樣做.
第三.
很重要的一件事.
兩個根本不明白.
我在口袋裡當時有的信用卡.
是可以買得完.
全香港麥當勞的薯條.
是不是.
只是欠錢而已.
我有無限的資源.
為什麼你要這麼吝嗇.
這麼斤斤計較.
只想保存著.
你自己口袋裡的薯條.
這就是講到奉獻一個很重要的.
三點很重要的訊息.
第一件事.
其實我們現在享有所有的東西.
其實基本上不是屬於你的.
是屬於上帝的.
我們要明白主權在哪裡.
我們今生所有的東西.
是屬於誰的.
不是屬於你的.
如果我們是基督徒.
我們應該看一看.
第二.
我們的上帝絕對有權.
拿走我們所有的一切.
就像路加夫人說的無知的財主.
那個財主說.
我多少的錢.
多少的收入.
收入多少的產物.
我可以建幾個倉庫.
明天慢慢坐著看.
有多少錢剩下呢.
我隨時可以拿走你的生命.

$^{201}$你今晚就要你的靈魂.
這些東西你可以什麼都沒有.
上個星期.
我一個好童工.
我現在在一間哨院教一課.
剛剛跟他教完書.
他開一班我開一班.
同一時間.
一班哨生.
李勝林牧師.
阿Mo的爸爸.
突然間心痛.
嘔吐.
進醫院.
當晚危急.
然後去世.
不足十二小時.
有時人就是這樣.
上帝接他回天家.
但是人的生命.
就像你做了皇帝的希律.
上帝都可以用蟲咬死你.
我們每一個人都要學習.
一個感恩的心去領受上帝給我們的.
因為生命氣息都是神給的.
如果神不給我們根本就沒有可能.
享受到現在有的東西.
最後.
亦都很重要.
我們要看到的是.
神有無限的資源.
神不是稀罕我們那十分之一.
稀罕我們給他有的東西.
而是上帝想我們學習一樣東西.
你想不想像我這樣.
我對兩個女兒.
我心裡想沒有下次.
你不給啊.
沒有下次.
你願意學習分享的時候.

$^{241}$上帝不會這麼吝嗇.
我們想的就是.
如果你願意在小事上忠心.
上帝亦都會在大事上.
對你有更多的賜予.
這是我們學習展望.
投資未來.
在我們用我們的資源的裡面.
所以今天我們講奉獻的時候.
當然我們不只是講錢.
我們講時間.
因為我們上次講到敬拜的時候都說到.
你來上帝面前敬拜.
不要說是付出.
是一個願意你擺上.
是一個你有的時間.
精神,體力,恩賜.
你有的物品.
你可以有的事奉能力.
你可以有的心智.
你可以有的金錢.
但為什麼金錢這麼重要呢?.
聖經裡面講了這麼多.
是因為金錢最實際.
貪念.
或者想保存金錢.
財富.
是人要突破的.
進天國裡面.
不要像下台穿針眼.
一樣的困難.
聖經裡面.
耶穌遇到很多不同的人.
包括少年的官也好.
很多富有的人.
也有些富有的人願意擺上.
像薩戈一樣.
但是我們一開始就看到.
哈娜這個母親.
她開始奉獻她的兒子之前.

$^{281}$都是從一個缺乏的開始.
很多聖經的故事都是.
我們覺得.
你找一些有錢人.
講奉獻金錢.
應該找伊隆·馬克斯.
不應該找我們.
這麼窮的人.
跟他比不上.
但是在聖經裡面.
耶穌遇到的和挑戰的.
不只是富有的人.
反而他覺得富有的人挑戰了之後.
他知道他不願意比.
就算了.
但是他的教導.
耶穌三十八個的比喻裡面.
有人說三十六個,三十七個.
基本上在三十八個的比喻裡面.
有超過一半.
是跟.
錢,財,財富有關.
在新約裡面.
所有的經文裡面.
經節裡面.
有超過五分之一.
是跟金錢有關.
講錢.
是比起講.
祈禱.
比起講信心.
是更加的多.
耶穌基督講的.
登山寶訓.
或者周圍跟人講的教導.
大部分的聽眾.
都是.
普羅大眾.
貧窮的人.
去到荒野裡面.

$^{321}$聽他講到連吃都沒有.
這個小朋友只有一頓午餐.
五餅二元.
是從一個貧窮的開始.
所以耶穌讚賞.
給養生兩個小時前的窮寡婦.
就是.
這是一個很窮的人.
他都懂得怎樣將他有的東西.
奉獻給上帝.
我們看到哈娜.
母仔生.
當時母仔生.
被人看為一個咒詛.
因為以色列人覺得.
生兒育女是神的祝福.
所以反過來的邏輯.
就是你不生兒育女.
就不是神的祝福.
所以就被人鄙視.
就不生育.
並且被二奶長期刺激.
在經文裡.
很生動地描述.
他經常刻意去刺激她.
不知刺激她什麼.
可能就是.
母仔生.
單單打打說話.
不單止這樣.
每年都是這樣.
去到聖殿敬拜.
我有個兒子你沒有.
每年都這樣刺激她.
於是她就.
當然很不開心.
老公剛才都說了.
完全不理解.
不懂得安慰.
有我就行了.

$^{361}$何必要有兒子.
你不明白她.
所以她就很不開心.
她就去.
求問耶和華.
她就.
將主權交給他.
因為沒有生兒育女.
她只明白.
主權在耶和華的手裡.
我想第一件事.
我想我們.
面對奉獻的時候.
無論是錢還是精神.
無論是時間體力.
還是恩賜.
你願意為神擺上的話.
我們.
是要明白.
主權是屬於上帝.
子女的主權.
也是屬於上帝.
萬君之耀.
求你顧念.
如果你給的話.
我就會有.
這就是第一件事.
我們要明白的心態.
就是我們每個人都赤身出院無胎.
赤身歸回.
這就是.
《約伯記》裡所說的.
試想想.
我們每個人都赤身出院無胎.
我看著兩個女兒出生.
在醫院.
有很多人都看過.
自己的小朋友出生.
沒有拿著信用卡.
沒有拿著iphone.

$^{401}$沒有帶著鑽石戒指.
沒有的.
我也有機會看著.
我四個老人家.
我岳父,岳母,我爸爸,我媽媽.
四個都是.
我主持.
安息禮.
我爸爸,我媽媽.
我去殮房認屍.
這個望著.
什麼都沒有.
帶給他的.
就是你放進去.
像岳父那樣.
放了一瓶白蘭地給他.
放在棺材裡.
因為他喜歡喝酒.
但是那些.
你帶不走.
岳父經常說.
兩個手伸出來.
空空如也.
什麼都帶不走.
我自己很感觸.
一年來美國讀書的時候.
讀大學.
去了West Virginia.
讀大學.
我當時坐飛機.
是.
兩個旅行行李.
來讀書.
接著我十幾年前.
從香港.
回到美國.
四十呎的貨櫃.
搬回來.
在貨櫃上.
幾百箱東西.

$^{441}$當貨櫃來到.
上海西的時候.
一看.
二百多箱.
幾百箱東西.
不斷搬進家裡.
因為.
我們省錢.
我們連香港的筲箕,筷子.
搬回來全部掃清.
整間屋掃清.
因為全部東西.
連摺椅子.
四十呎貨櫃.
當時很便宜.
運輸不是很貴.
但我看著那個貨櫃.
我自己心裡有一個感覺.
我們豈不是赤身出於母胎.
赤身歸會.
由我們出生到死.
中間所有的東西.
都是上帝.
豐盛的賜予.
我根本不是因為.
我這麼厲害.
我可以拿到這麼多東西.
你要賺錢.
生命是上帝賜予的.
你有能力.
你有智慧.
你要讀書.
都是上帝賜予的.
我們首先要明白.
要認識我們的主權.
我們的奉獻不是.
失去.
不是一個損失.
又傳奉獻代價.
奉獻又少了十分之一.

$^{481}$而是.
我們去施予.
我們去分享.
本來這些東西就是屬於神的.
神給你十個蘋果.
提醒你一件事.
你給花果出來.
一會兒我們先說說.
為什麼我們要給花果出來.
為什麼上帝要我們這樣做.
我們是為上帝去做.
我們不是損失了自己屬於的東西.
這個就是我們.
後書九章九節.
他施捨周濟貧窮.
他的義全道榮.
上帝提醒我們.
只要我們願意施予.
願意去比喻去奉獻的話.
我們就是一個義人.
很厲害.
為什麼我們會是義人.
因為耶穌說了這個比喻.
我每次都說.
這個比喻在.
馬太福音二十五章裡面.
有很多的標準.
評估的標準.
這個不是唯一一個.
其中一個評估的標準.
就是你如何運用上帝給你的東西.
去幫助.
去為神去做一些事情.
以前有分開.
綿羊和山羊.
大家記得二十五章.
三十一至四十六節.
上帝於是.
和右邊的綿羊.
即是神稱讚的人說.

$^{521}$你是有福的.
你要享受.
創世以來神給你的.
預備的國.
因為我餓了你給我吃.
我渴了你給我喝.
我流浪你給我住.
我赤身露體你給衣服我穿.
我病了你來看我.
我在監獄裡你來探望我.
聖經裡面.
耶穌所說.
義人就會踏船.
你自己回去看聖經.
這個義.
不是關義哥那種義.
不是黑社會那種義.
這個義人.
是神稱你為義.
即是你.
神的標準你已經達到了.
不是真的因為你達到了.
而是神覺得你已經達到了.
為何我能成為一個義人.
這群義人就回答.
主啊.
我什麼時候見你痛啊.
什麼時候見你渴啊.
什麼時候見你坐牢啊.
什麼時候見你病啊.
來.
回應你啊.
王就會踏船.
我實在告訴你們.
這些事你們做在.
我的弟兄當中一個最小的身上.
就是做在.
我的身上.
真的這麼簡單.
你今天.

$^{561}$好像Terence.
和一班crew去搞音響.
你是.
做在我們這班.
小的弟兄姊妹身上.
我們不認識.
你可以幫我們.
你有你的恩賜你要付出.
Terence 你就是義人.
一會兒完了崇拜.
恭喜他 他就是義人.
是上帝說的.
你願意為這班弟兄姊妹.
做的話 你就是義人.
這個就是.
上帝看我們的.
其中一個標準.
弟兄姊妹.
第一件事我們要明白.
這個主權.
我們就不會那麼多的斤斤計較.
有時候我們.
給了出來一些東西.
我們把時間放在弟兄姊妹身上.
這些人吃愛心.
你怎麼愛她 怎麼關心她.
她都走了 不理你.
拿了你很多好處.
又或者.
見到乞丐你給錢她.
回到家 媽媽告訴你.
乞丐有幾棟樓.
不要給她.
你覺得.
這個蝕了.
被人騙了.
你心裡覺得.
不開心.
因為給了錢出去.
原來人家騙你的.

$^{601}$但弟兄姊妹.
如果你給.
是因為給上帝.
神說你每樣東西.
你的目的.
當然我們都有很多不同的目的.
但其中一個很重要的是.
我們是給上帝.
神告訴你.
祂會眷顧你.
因為主權是屬於祂.
因為你有東西都是屬於祂.
你為神去做 為神去給.
時間 金錢 恩賜.
能力等等.
都是在給上帝.
如果在給上帝.
你就不要再斤斤計較.
因為你已經給了上帝.
當然你要有.
聰明智慧去.
分辨.
就算你被人騙了.
你也是在給上帝.
這是神提醒我們第一件事.
第二件事也很重要.
回應上帝的恩典.
哈娜跟上帝說.
我祈求.
得了這個兒子.
將祂所求的賜給我.
所以.
我將這個孩子獻給耶和華.
使祂終生.
歸給耶和華.
要將投生的.
獻給上帝.
為什麼要將投生的獻給上帝.
為什麼聖經裡經常說.
投生的初熟的果子.

$^{641}$我們上次說敬拜.
語節要獻.
投生的牛或羊.
五旬節要獻.
初熟的果子.
豐收時要獻.
感恩的祭.
為什麼經常說.
投生的要獻給上帝.
因為.
第一是感恩.
那隻牛不一定.
可以生.
那些果子.
不一定.
最近我在家中種苦瓜.
種了四五棵.
現在又再種.
不一定可以生.
但如果真的可以生.
首先要感恩.
沒有了.
所以第二件事很重要.
你的信心在哪裡.
你信不信上帝.
真的會給你.
有時候我們很吝嗇.
我們不願意給.
我們不想給.
但是.
上帝說.
你以此試試我.
我真的會上尖下流.
延年搖大安.
賜福給你嗎.
你試試吧.
上帝就是這樣挑戰我們.
你沒有了.
那隻小牛.
你奉獻給上帝.

$^{681}$你信不信上帝會賜福給.
這隻母牛再繼續生.
多些小牛出來.
就好像哈娜.
生了一個之後.
再生幾個子女.
成為一個.
莫大的祝福和美好的見證.
這就是.
上帝要我們學習受一奉獻的原因.
大家知道聖經裡面.
追溯源頭.
受一奉獻就是.
由阿不蘭.
四王五王的時候.
他打敗了.
打敗了.
他的敵人.
然後麥羈洗德.
出來迎接他.
然後.
獻酒.
餅.
為他祝福.
為阿不蘭祝福.
這個就是.
阿不蘭將他戰利品的.
十分之一獻給.
麥羈洗德.
大家要記住.
麥羈洗德是怎樣做的呢.
如果你看.
創世記裡面.
麥羈洗德不是.
那場戰爭裡面的.
其中一份子.
他是耶路撒冷的王.
但他不是五王四王裡面其中一個.
但他封聞到.
這個阿不蘭.

$^{721}$他的封聖.
他是主動出來迎接.
和獻上.
他的禮物.
而阿不蘭也願意將.
十分之一獻給他.
因為他是代表.
他是祭司.
耶路撒冷的王也是祭司.
他是神的代表.
所以阿不蘭就願意將.
他有的十分之一獻給.
上帝.
接著.
神設立了會幕.
然後有聖殿.
在會幕裡面.
你奉獻錢.
給會幕裡面.
有三個很重要的.
一個祭祀的.
制度.
如果你回去看.
文宿紀十八章.
裡面講到.
第一件事就是供養祭司.
因為亞倫和.
利美的後裔.
他們不做其他事情.
只是做獻祭.
每天為.
以色列人獻祭.
大節的時候以色列人來到.
他們幫忙獻祭.
專心去做.
神的工作服事.
所以要去.
供養他們.
好像帕西傑奇那樣.
沒有兼職.

$^{761}$全時間沒有兼職.
不出聲.
我們要供養他.
支持他每天.
為教會所做的事.
第二件事.
這個申命記十四章.
裡面講到.
聖殿裡面的需要.
當時以色列人起初的時候.
沒有匯膜沒有聖殿.
只有利未人.
有匯膜的時候.
匯膜裡面要有.
運作的費用.
殺一隻牛.
磨一把刀都要錢.
磨一把石頭都要.
清潔那些.
聖殿.
匯膜.
What a bloody mess.
一隻牛殺羊.
要清潔.
很多東西.
要有資源去運作.
所以這個申命記十四章.
告訴我們.
錢也是在聖殿裡面運作.
最後幫助當中.
以色列人的需要.
就是社區的需要.
這個.
以色列一些貧窮的人的需要.
所以教會.
做辭位的工作是必須的.
有需要的人.
我們要去幫助他們.
所以.
奉獻在.

$^{801}$殿裡面去到馬拉基斯的時候.
已經是指示.
人們一定要將.
奉獻給.
上帝的殿 直至今天.
所以要說.
多一點.
經常人們都跟我說.
為什麼要奉獻給教會.
為什麼不可以奉獻給.
地球.
我經常說.
在整個十一奉獻的發展裡面.
大家都知道.
起先的時候.
阿伯和該人的奉獻不是十分之一.
是將最好的獻上.
羅馬書12章.
一字二字.
將所有的獻上.
生命獻上.
但十分之一是一個.
學習的機會.
是神讓你學習感恩.
所以每逢你.
做十一奉獻的時候.
神定了一個規矩.
給我們學一個ABC.
要將生命獻上.
但十分之一是要讓你學去感恩.
所以每逢你.
給奉獻的時候.
是一個提醒你.
你有多少恩典是上帝給你的.
有十個蘋果給了你.
你只要給一個出來.
是提醒你這十個蘋果.
另外九個蘋果都是上帝給你的.
你要學習信心.
你將第一樣最好的東西.

$^{841}$給了上帝之後.
你信不信得過.
聖經所說的祂會賜福給你.
在實際的裡面賜福給你.
第三.
神的群體.
我們經常說這是屬靈的家.
是你的家.
我真的不明白.
實際上去想.
這個家裡面.
你不去供養家裡的需要.
你走出去.
人家的家裡面.
去供養他們.
家裡面也有很多的需要.
所以我覺得財政.
是需要去透明.
讓大家知道.
我們的錢會怎樣用.
我們的預算會怎樣用.
我們會明白.
我們去奉獻是給上帝.
也是照顧家裡的需要.
因為一件事都是錢.
你今天坐下來.
拿杯咖啡.
拿杯子.
都要錢買的.
我和我太太.
負責茶水部.
我經常和她.
去Costco.
Business Costco.
買那些杯.
你知道那些人.
喝一半就丟掉.
很快就用完.
我和太太說不要買那麼大杯.
像自助餐.

$^{881}$那些碟子弄得小小的.
那你就煩她.
讓她拿少一點就行.
太太說不要.
給吧.
那就給吧.
我想.
我們要明白.
很多東西都要用錢.
去運作.
你今天.
可以回來.
停車場不用錢.
下來喝咖啡不用錢.
吃茶葉蛋不用錢.
去廁所不用錢.
為什麼廁所要錢.
大家歐洲都知道.
是美國才不用錢.
歐洲.
沒錢也不要去廁所.
接著.
你坐下來.
燈給電.
你用這些PowerPoint.
Projector.
這些音響.
所有東西都要錢.
你一個線都不用出.
是因為有人幫你付了.
是教會幫你付了.
有什麼是免費的.
是不是.
Nothing.
是因為有人付了.
救恩是免費的.
是因為有人幫你付了.
就是這麼簡單.
今天你可以.
免費來這裡聚會.

$^{921}$是因為有人幫你付了.
所以今天我們.
學習就是.
受恩奉獻就是一個.
基本學習感恩和信心的.
功課.
奉獻不是一個損失.
是一個提醒你.
恩典和信心的.
功課.
而恩典和信心的功課不是感覺.
而是實際的行動.
你說我沒有錢.
你給時間.
給精神出來.
如果你有的話.
學一下去做受恩奉獻.
最後也很重要的.
投資未來.
好像我對我兩個女兒.
心裡面想.
你今天學去願意分享的話.
我以後再給多一點的東西.
你給都給得很少.
是不是.
那時三妹仍是孩子.
依依不得.
在耶和華面前.
給她一生生了一竇.
那孩子.
在耶和華面前.
漸漸長大.
這個就是聖經的一個真理.
這個.
就是.
哥林多後書.
第九章六章說.
小眾是小受,多眾是多受.
各人要隨心所願.
不要為難,不要勉強.

$^{961}$愛那些樂捐的人.
上帝能將各樣的恩惠.
多多加給你們.
使你們能有.
充足去做各樣善事.
這個就是一個.
所謂的投資.
捆存你有的東西.
就是.
一片的田.
你種東西.
你有十粒種子.
你說不要那麼吝嗇.
只給一粒.
那你就得一粒.
出來的苦瓜.
是不是.
但如果你說你要種多一點.
希望收成多一點.
理論就是.
你要殺多一點種.
多種是多收.
少種是少收.
肯定的.
就像銀行利息.
銀行有十元有十元利息.
你怎麼可以多一點利息.
就多放點錢進去.
這樣的理論.
但不要勉強.
因為上帝是愛樂捐的人.
但上帝告訴你.
當你願意.
捆存.
不是集中.
不是積存.
是一直投資出去.
上帝會給你越來越多.
為了什麼呢?.
就是你能夠給更多.

$^{1001}$這就是聖經裡面.
上次提醒我們的真理.
就是要將我們有的東西.
殺出去.
做貿易.
到時候.
你就會有收成.
這是未來的事.
因為你不知道.
生命何時會結束.
沒有機會這樣做.
結束的時候.
我想講這個比喻.
耶穌的比喻.
都是正面的.
有些是負面的.
用負面的東西.
來提醒我們正面的理念.
他說.
路加福音第十六章.
一個負面的比喻.
這個負面的比喻就是.
有一個不好的僕人.
做得不好.
老闆又要炒他魷魚.
這個僕人就說.
快要炒魷魚了,怎麼辦呢?.
他管錢財的.
於是他就走出去.
抓住主人的債.
就跟債人說.
大家把他拆了.
你欠十塊.
我給你八塊算了.
沒有債.
免你的債項.
很多個都這樣.
你收他四個半算了.
你簽了名.
沒有債.

$^{1041}$當然.
錢就是主人虧了.
但耶穌就反而讚這個人.
為什麼要讚他呢?.
因為耶穌說.
今世之子.
應付自己.
的世代.
比光明之子更加正.
他說.
這個.
狡猾的.
這個僕人.
他也懂得.
運用錢.
為他未來.
被人炒魷魚之後.
去找工作.
他心裡想.
這群人,我有恩於他.
將來他們就會請我做工.
我不用失業.
於是對他好.
耶穌說.
這個狡猾不義的僕人.
也懂得用錢.
為他未來.
打算去投資.
懂得用主人給他的錢.
令他未來不用失業.
你們這些所謂光明之子.
我們反而不懂得這樣做.
或者不會這樣做.
這不是很搞笑嗎?.
耶穌說我們要學習.
要投資.
在未來.
下面就是耶穌這樣說.
在小事上忠心.
在十一奉獻中忠心.

$^{1081}$大事上也忠心.
小事上不義,大事上也不義.
所以.
在不義的錢財中不忠心.
誰還會將.
真實的錢財.
托付給你們呢?.
如果你們在別人的東西上不忠心.
誰還會把你自己的東西.
給你們呢?.
主人不給你的東西.
你也不忠心的話.
你怎樣去學去忠心呢?.
所以一個僕人不能夠屬於兩個主.
大家很熟悉這句話.
不是恨這個就愛那個.
是重者過,輕那過.
你們不能既事奉.
或服侍上帝,又服侍.
錢財,服侍.
馬門.
等下哈娜給我們看到.
她將兒子奉獻給耶和華.
如果她把兒子藏在家裡.
只是一個.
數學很厲害.
STEM很厲害.
讀醫學碩士.
哈佛,史丹佛.
是一個名校的學生.
但是.
如果她將兒子.
奉獻給耶和華.
她就成為了一個.
能夠改變時代的人.
神就用撒母耳.
帶領以色列人.
走過一段艱難的日子.
領子永遠記住.
奉獻不是失去.

$^{1121}$而是.
投資在神的國度裡.
你說投資在哪裡?.
其中一件事.
我覺得勝經說得很清楚.
神的國裡.
是有人.
一定有人.
一定有.
信了主的人.
一定有基督徒.
在當中.
所以我鼓勵大家.
用你的錢財.
投資在福音的工作.
投資在.
基督徒訓練的工作.
因為這些投資.
是會去到永恆的裡面.
求主幫助我們.
我們低頭,我們禱告.
因為從哈娜母親的奉獻裡面.
你讓我們學習到.
我們應該怎樣.
過我們這個豐盛的人生.
因為豐盛的人生.
不是積楚.
而是要將你給我們的.
我們願意去施如.
願意去奉獻,願意去擺上.
願意去付出.
正因為我們給.
是因為天父你先給我們.
你將你的兒子.
耶穌基督.
去賜下給我們.
以致我們能夠白白得罪這個舊恩.
我們亦都好像耶穌基督提醒我們.
我們白白得來.
亦都學習怎樣白白去.

$^{1161}$捨去.
主,遂聽我們禱告,奉耶穌聖名.
阿們.
\newpage



\section{約書亞記 2:11}
\label{sec:032095}
\textbf{勇敢機智信靠神的妓女喇合}~~~~黃克勤牧師博士
\newline
\newline
連結: \href{http://www.church.com.hk/acms/content.asp?site=cdc\&op=show\&type=product\&code=032095\&layout=sermon}{\texttt{032095}} ~~~~ 語音日期: 2026-05-07
\newline
\newline
\hyperref[sec:032094]{\small{< < < PREV SERMON < < <}}
~
\hyperref[sec:index_chronic]{\small{[返順時目]}}
~
\hyperref[sec:032096]{\small{> > > NEXT SERMON > > >}}
\newline
\newline
約書亞記 2:11
\newline
\begin{longtable}{cl}
\hline
\hline
章節 & 經文 (和合本修訂版)\\
\hline
2:11 & \begin{tabularx}{0.7\textwidth}{X} 我們一聽見就膽戰心驚，人人因你們的緣故勇氣全失。耶和華－你們的神是天上地下的神。 \end{tabularx} \\ \\
[1ex]
\hline
\hline
\end{longtable}
$^{1}$各位大家好,眨眼又有一日讀聖經了.
生命說書的網絡和大家再看聖經人,聖經和事.
看看上帝背後在做些什麼.
最重要不是看到前面的人和事.
而是背後上帝是怎樣的.
今次我給了一個這麼長的題目.
就是講妓女拉俠.
我將它形容做勇敢,機智,而且信靠神的.
這個不是聖經這樣說的.
是我覺得讀完之後.
覺得拉俠真的是勇敢,機智和信靠神.
我們就會妓女加拉俠講兩次的.
先看一下這個地圖.
希望你看得清楚.
因為我都是在Google那裡找.
就蒙蒙茶茶放到這麼大之後.
所以我加了一些字的.
這個沒有加都知道了.
在這個位置就是埃及.
這是他們做奴隸的時候建的其中一個織貨城.
就是貨倉城.
就叫蘭塞.
蘭塞是一個埃及法老的名字.
他們就離開了.
雨至之後出來了.
撒掌以災.
然後就走了去曠野.
然後繞回上去.
繞回上去.
埃及人就說哈哈哈他們在曠野迷路了.
我們去追他們.
於是就看到他們這樣繞了.
繞到這裡有少少藍色是海.
這個地中海有部分的水來到這裡.
可能是斷或者是有一個湖.
或者是水漲水退的時候會通.
就是這裡.
這裡就是他們過紅海的地方.
不是紅海這裡過的嗎?.
不是這裡過的嗎?.

$^{41}$這裡才是紅海.
紅海的頂.
紅海很大.
直到下面.
有兩個很長的海灣.
其實原來紅海這個字很有可能是譯錯了.
Souf.
Souf是希伯來文.
希伯來文譯做英文字母.
Souf.
希伯來文Souf的意思就是奴隸.
可能那裡很多奴隸.
奴隸海一定是離開了海邊遠一點.
因為如果是太淡水就沒有奴隸.
可能是淡水的地方.
這裡現在通了.
就是蘇伊士運河.
淡水交界才有奴隸或者全淡水.
我聽到他們這樣說.
現在通了.
打通了變成蘇伊士運河.
這樣進去紅海.
然後出來了.
蘇伊士運河就是這個位置.
奴隸海.
然後進去尋的曠野.
然後來到順的曠野.
就是米利巴的地方.
前幾次說.
沒有水喝.
發火了.
於是神就叫摩西.
拿著你的杖子.
去吩咐盆石出水.
誰知道摩西那一刻.
真的沒見過他內心這麼脾氣這麼憤怒.
擊打盆石兩下.
於是出了很多水.
所以盆石我上次說過.
是一個很大的一塊巨型.

$^{81}$好像山一樣的岩石.
不是一塊石板.
有一條水柱.
200萬人能喝到什麼時候.
而吩咐或者擊打盆石.
哇!整個山出水,這麼大.
就是米利巴.
然後去到河列山就有十界了.
然後要向上走.
就是加里斯曠野.
來到這裡有個山叫河爾山.
我記得我曾經說錯了.
就是河列山.
其實是河爾山記錯了.
亞倫死在河爾山.
死的時候123歲.
當時摩西還沒死.
繼續帶著他們走這段路.
差不多要進入迦南.
死海.
然後這裡上面就是約旦河.
他們要在這裡過約旦河才進入耶利哥.
過了約旦河就叫進入應許之地.
當然神給他們不單止是約旦河的西邊.
東邊這一區都給了他們.
摩亞,阿滿,以東這些都給了他們.
一直上到很遠.
去到尼巴利遠那麼遠.
去到這裡有個山叫尼波山.
這個重要了.
實際上跟摩西說.
你曾經發脾氣.
所以你是不能夠進入應許之地了.
不過我帶你去尼波山.
看一看應許之地吧.
如果我是摩西真的慘.
真的第二次發火.
喂 你不是吧.
這麼辛苦走這麼遠的路四十年.
來到這個位置.

$^{121}$百姓埋怨得不得了.
去到這裡你不讓我進入.
就因為打了兩下.
好像比較敏感吧.
是不是.
好像比較苛刻.
但是摩西很信服.
沒有爭拗.
好啊好啊.
他就上山了.
看完之後他就死在那裡了.
誰撞在哪裡不知道.
誰撞他.
他怎樣被撞不知道.
就在尼波山看耶利哥.
然後他們繼續走上去.
高一點繞一圈.
來到要過河的地點.
這個紅色點.
下面這個紅色點就是尼波山.
上面這個紅色點.
就叫做十亭.
過河就是耶利哥了.
大家知道了.
大家先有些概念.
約書亞記第二節.
當下亂的日子.
約書亞從十亭.
剛剛打發兩個人做探子.
在這裡要過河.
那還沒去過.
於是就找兩個人做探子.
你過河先去耶利哥那裡.
去看看.
打聽一下.
窺探一下.
和我們這樣.
這件事就是這樣的.
打發兩個人做探子.
就吩咐他了.

$^{161}$你們去窺探那地.
和耶利哥.
因為耶利哥是我們進去的第一個城.
於是二人就去了.
就去了一個妓女名叫拉合的家女.
就在那裡躺臥.
這個躺臥就只是在那裡預備過夜.
怎麼會搞到去一個妓女家裡.
一去就去這麼巧.
你可以想像.
他們會不會有家男人的服裝在身上呢.
沒有的.
兩個人都穿回自己的衣服.
就去扮遊客.
這樣走就去.
這些外人一看就看到了.
如果今天.
如果你要.
我們這樣說.
著草.
著草去哪裡.
找條村住吧.
你千萬不要.
找條村住.
一走進去.
每個村民一眼就知道你不是本地人.
大隱隱於市.
你去大城市.
大都會那裡住.
國家戰鬥門前說.
誰住隔壁問都沒有問.
我隔壁跟屋主已經住了十幾年.
我都不知道他吃什麼.
真的不知道.
大家不打招呼又不同一個時間出門口.
我只知道他吸煙吸得很多.
什麼都不知道.
如果現在在街上碰到面.
我是不知道的.
同一輛電梯我都不知道.

$^{201}$因為我都沒有看他的樣子.
他是日本人.
這樣.
去到那裡.
又天黑了.
快天黑了.
你知道妓女會怎樣.
會不會在屋裡.
她就在門口.
見到兩個外人.
先生.
要不要進來喝杯茶.
聊聊天.
走到很累了.
她說好啊.
找這個跟我們聽一下.
收一下風.
雖然她進去了.
可能她去之前.
她不會說.
兩位先生.
我是妓女來的.
來我這裡.
她不會的.
是不是走得很累.
不如進來喝杯茶.
喝點水.
大家聊聊天.
你們是不是外面來的.
亦都想快點進屋.
不要被人在街上.
個個眼望望.
所以她找到一個人.
所以妓女在門口招呼到她.
亦都不出奇.
亦都最有可能.
有人告訴耶利哥王.
今夜有意識的人.
來到這裡窺探此地.
因為被人見到.

$^{241}$又知道他走進去.
拉合那間屋那裡.
於是王就說.
去找他們回來.
於是人們就去拉合家裡.
跟拉合說.
來到你這裡.
進了你家的人.
要交出來.
因為他們來窺探我們.
其實這個時候.
拉合知不知道.
這兩個意識裂人的來意呢.
都大概知道.
或者神也有感動她.
或者她的聰明智慧都知道.
眉頭眼額識看.
聞起來已經知道.
她知道.
她就來窺探我們.
這個第一點要看的.
女人將二人隱藏.
這個拉合一個被欣賞的地方.
就是她很夠膽.
勇敢兼機智.
因為如果被人找到.
我藏起來.
我死人的.
她都將他藏起來隱藏.
即是說藏起來.
而且藏起來.
當時你的屋子那麼小.
一找就找到.
一定死人的.
所以她要用機智的方法.
是能夠打發這些士兵走.
所以我覺得拉合是勇敢機智.
又勇敢又機智.
女人將拉合隱藏.
回答說.

$^{281}$答什麼呢.
真的.
那些人果然來到我這裡.
他們從哪裡來.
我就不知道.
天黑要關城門了.
當時天黑的時候.
那些城都會關門的.
要保護.
天黑要關城門的時候.
他們就走了.
他們出去了.
往哪裡去.
我卻不知道.
不如你們快點去追吧.
這樣就會追到的.
沒多久而已.
走了沒多久.
你快點去追吧.
那些人就沒有說.
那樣嗎.
先搜搜你的屋子.
那樣也沒有.
因為她說得似曾相識.
啊.
厲害.
又勇敢又機智.
你快點走吧.
他們剛剛要關門.
所以你快點離城出去了.
就應該走回到道口那邊.
約旦河的道口.
過去河那邊.
於是.
好啊.
我們快點追他.
追他.
先是女人領了二人上了房頂.
多聰明.
一早將他們.

$^{321}$哇.
你們來的.
那就快點收起來.
先躲在那裡.
就擺在那些麻吉那裡.
即是.
在那些麻的.
碎的麻.
麻的.
收那些麻繩.
麻的質的東西.
就擺在那裡.
躲起來.
就沒那麼容易找到.
帶過去.
那些人就往約旦河的道口追趕去了.
追趕他們的人.
一出去.
城門就關了.
是對的.
那些人為什麼會追趕呢.
第一要快.
第二呢.
城門快要關了.
那我關了門.
那些城門關了門.
就不是那麼容易開了.
那麼麻煩一點了.
於是.
哇.
快點快點追出去.
這樣就.
追了.
二人還會躺臥.
即是雖然晚上天黑了.
那些人還沒睡.
兩個探子還沒睡.
女人馬上就上房頂告訴他們.
就這些東西要留意了.
我知道耶和華.

$^{361}$已經把這地賜給你們.
並且因你們的緣故.
我們都驚慌了.
這地的一切居民在你們面前.
心都消化了.
喂.
他自己給情報啊.
就是我知道.
那麼你們呢.
是耶和華已經把這個地賜給你的了.
他怎麼這麼懂事呢.
其實呢.
他可能以前聽過了.
又或者這次.
他們來之前聊天已經知道了.
那就哦.
這樣.
那其實呢.
我們知道耶和華.
是做過這些事情.
我們很害怕.
誰知道.
因為.
耶和華的名字都知道.
是啊.
因為我們聽到你們出埃及的時候.
在遠方出埃及.
耶和華怎樣在你們面前.
令到洪水的海乾了.
並且你們是怎樣.
是等待約旦河的兩張.
約旦河東的兩個王.
阿摩利王西王.
和岳將他們進行毀滅.
哇.
我知道耶和華幫你們.
耶和華是真的.
是厲害的.
我們已經很害怕了.
因為我們聽到什麼.

$^{401}$耶和華幫你們.
令到這麼強大的埃及軍隊.
完全葬身海底.
洪海乾了.
你們走過.
然後埃及的軍兵全沒了.
這個洪海這個洪字.
我剛才講過了.
洪海這個字在這裡的希伯來文.
不是red.
不是洪.
是soof.
s-o-o-f.
希伯來文的譯.
譯做英文字.
soof.
soof就是蘆葦.
你們蘆葦海乾了.
你們這樣走過的.
這樣.
於是就.
哦.
哎呀.
你們那個神耶和華真厲害.
我們的心全部消化了.
消化就是融化.
唉.
站不住了.
哎呀.
不行了.
沒有了.
因為你的緣故.
沒有一個人有膽氣.
我們這裡.
就是沒有膽量.
沒有勇氣.
那麼膽氣.
蠻有趣的膽氣.
中文膽氣.
那麼膽量勇氣.

$^{441}$沒有膽量勇氣.
又說你們的神本是上天下地的神.
因此我看到你的神.
真是天上地下都管了.
是真神來的.
很厲害.
那麼.
我都信靠他.
喏.
拉合做這些事情.
保護他們.
其實他已經心裡面.
是哇.
羨慕了.
以色列族那個神.
知道他叫耶和華的.
傳說都知道的了.
通了天了.
這個神真是真的.
我也要信祂的.
那麼現在你們兩個來到.
就是那裡派來的.
哎呀.
太好了.
他當然保護了.
出盡力.
又有機智又有勇敢.
好了.
現在我是恩代了你們.
求你指著耶和華向我起誓.
也要恩代我父家.
並且給我一個實在的證據.
救回我父母弟兄姊妹.
和一切屬於他們的.
拯救我們的生命不死.
在說什麼呢.
他們兩位就說.
那麼.
嗯.
拉合就說.

$^{481}$是啊.
現在我就恩代救了你們.
你們可不可以指著耶和華神起誓.
也要恩代我父家.
他不只是自己的.
不如你們兩個娶了我吧.
隨便一個娶我.
那不是.
是我的父家.
是我全家人.
全個家族.
而且不單止你們兩個口說.
還想你們有一個確實的.
實在的證據.
就是說.
誓言啊.
信物啊.
那些這樣給我zar 手.
那麼其實呢.
我是想.
我恩代你.
也想你能夠恩代我父家.
救回我們全家人.
有人就對他說.
如果你不洩露我們來過.
我們情願替你死都可以.
因為你剛才也替我們死了.
已經是很危險的了.
不過你這麼夠膽又替我們死.
為我們擋住.
好,我們的命其實是屬於你的.
如果不洩露到這個軍情.
我們必以慈愛誠實待你.
我們也會一樣這樣恩代你的.
於是女人用繩子將二人從窗戶裡垂下去.
因為她的房子在城牆邊上.
她也住在城牆上面.
住在耶利哥城牆上面.
因此她把繩子扯下去.
夜晚讓她快點走.

$^{521}$不過她更好.
做人做到底就送什麼.
她對他們說.
你們不要馬上走去渡口.
恐怕你撞到他們回來.
那你要做什麼呢.
恐怕追趕的人會碰見你.
你們先在那裡隱藏三天.
你們上山躲三天.
然後等追趕的人回來了.
他們不單止去到渡口.
何的渡口追不到你呢.
他們會到處找.
然後等他們幾天你們才走.
交了路.
二人對她說.
你這樣行不然.
你要這樣做.
否則我們起的誓無竿.
就是什麼呢.
將來我們來到的時候.
你要第三.
借一條珠紅線繩.
繫在誰我門下去的窗戶上.
這樣就會令到你的父母兄弟.
全家聚集在家裡.
這樣就會成功.
如果走出家門.
對不起.
走在街上.
我們打起來.
殺起來.
不知道是誰.
他就自己歸到自己身上.
這個責任.
在你家裡.
是沒有人向他下手.
如果不是有人向你家人下手.
罪在我們那裡.
這樣就歸到我們頭上.

$^{561}$這個就是他們說的.
第三點.
這條珠紅線繫在窗戶上.
是會救你們的.
信律就是一條珠紅的繩.
用珠紅的線織成的繩.
你們掛在窗口.
於是我們進城的時候.
見到這條珠紅色的線在這個窗戶.
就是這一家人我們不會殺.
這樣.
其實到下一集.
明天你就會看到.
那兩個探子是去這間屋裡救他們的.
帶他們走的.
他們全部都活命的.
這條珠紅色的繩是一個信實的標記.
在這裡用屬靈的.
沒有特別的根據的.
只不過我會多說的就是.
這條珠紅色的繩是救了他一家人.
正如在雨至的時候.
門尾門腔上的血.
是救了那一家人的.
救了那一家人的長子.
這條珠紅線繩是會救了他的一家人.
珠紅色血的顏色.
我會將它解釋為.
其實都是會救他一家人的.
最重要就是你這一家人.
是在這條珠紅線織成的繩上面.
的底下.
就可以了.
都會被救.
就是說凡是這個救贖裡面都可以的.
珠紅是象徵血的顏色.
這個可能是我自己多餘的估計.
不過我都喜歡這樣估計的.
所以用珠紅色的繩.
做這條救命的證據.

$^{601}$這樣放在窗上.
都挺好的.
好了.
這個我是將它稱呼為.
勇敢機智又信靠神的基里拉俠.
他就做了這件事.
保存了自己家裡的性命.
也都保存了這兩個探子的性命.
而將來神是比基里拉俠有很高的身份的.
好.
我們今天先講到這裡.
這件事是很值得我們去感恩讚美和思考的.
禱告.
天護受弟多謝你.
你無所不能.
又無所不在.
又無所不知.
你是在茫茫人海裡面.
你是預備了這個基里拉俠.
在耶利哥城等這兩個探子.
他們進入了基里拉俠的家.
真是進得正好.
不然恐怕他們已經沒辦法活命了.
已經死了.
因為他們進城就已經被人看到了.
已經有人通風報信要去抓他們了.
就拉俠勇敢將他們隱藏了.
救了他們.
而他們都知恩忘報.
於是又給一個記號他.
一條珠紅的繩掛在窗口.
你的部人覺得珠紅的繩子象徵著血.
亦都想到令我聯想到救贖.
你自己可能都給了一個隱喻在當中.
幫助我們記得.
其實在這條紅繩的底下.
就得救.
我們多謝你.
為了拉俠的勇敢機智和對你的信靠.
我們感恩我們學習.

$^{641}$還有下一集更多看到你自己怎樣使用他.
怎樣尊重這個人.
怎樣尊敬他.
都是你的恩典.
幫助我們.
祈禱奉耶穌基督名求.
阿們.
\newpage



\section{申命記 7:1-26}
\label{sec:032096}
\textbf{上帝也講KPI?(3) 超出預計的愛}~~~~李觀南傳道
\newline
\newline
連結: \href{http://www.church.com.hk/acms/content.asp?site=cdc\&op=show\&type=product\&code=032096\&layout=sermon}{\texttt{032096}} ~~~~ 語音日期: 2026-05-02
\newline
\newline
\hyperref[sec:032095]{\small{< < < PREV SERMON < < <}}
~
\hyperref[sec:index_chronic]{\small{[返順時目]}}
~
\hyperref[sec:032097]{\small{> > > NEXT SERMON > > >}}
\newline
\newline
申命記 7:1-26
\newline
\begin{longtable}{cl}
\hline
\hline
章節 & 經文 (和合本修訂版)\\
\hline
7:1 & \begin{tabularx}{0.7\textwidth}{X} 「耶和華－你的神領你進入你要得為業之地，從你面前趕出許多國家，就是比你更強大的七個國家：赫人、革迦撒人、亞摩利人、迦南人、比利洗人、希未人、耶布斯人。 \end{tabularx} \\ \\ \relax
7:2 & \begin{tabularx}{0.7\textwidth}{X} 當耶和華－你的神把他們交給你，你擊殺他們的時候，你要完全消滅他們，不可與他們立約，也不可憐惜他們。 \end{tabularx} \\ \\ \relax
7:3 & \begin{tabularx}{0.7\textwidth}{X} 不可與他們結親；不可將你的女兒嫁給他的兒子，也不可叫你的兒子娶他的女兒。 \end{tabularx} \\ \\ \relax
7:4 & \begin{tabularx}{0.7\textwidth}{X} 因為他必使你的兒女離棄我，去事奉別神，以致耶和華的怒氣向你們發作，迅速將你除滅。 \end{tabularx} \\ \\ \relax
7:5 & \begin{tabularx}{0.7\textwidth}{X} 你們卻要這樣處置他們：拆毀他們的祭壇，打碎他們的柱像，砍斷他們的亞舍拉，用火焚燒他們雕刻的偶像。 \end{tabularx} \\ \\ \relax
7:6 & \begin{tabularx}{0.7\textwidth}{X} 「因為你是屬於耶和華－你神神聖的子民；耶和華－你的神從地面上的萬民中揀選你，作自己寶貴的子民。 \end{tabularx} \\ \\ \relax
7:7 & \begin{tabularx}{0.7\textwidth}{X} 耶和華專愛你們，揀選你們，並非因你們人數比任何民族多，其實你們的人數在各民族中是最少的。 \end{tabularx} \\ \\ \relax
7:8 & \begin{tabularx}{0.7\textwidth}{X} 因為耶和華愛你們，又因要遵守他向你們列祖所起的誓，耶和華就用大能的手領你們出來，救贖你脫離為奴之家，脫離埃及王法老的手。 \end{tabularx} \\ \\ \relax
7:9 & \begin{tabularx}{0.7\textwidth}{X} 所以，你知道耶和華－你的神，他是神，是信實的神。他向愛他、守他誡命的人守約施慈愛，直到千代； \end{tabularx} \\ \\ \relax
7:10 & \begin{tabularx}{0.7\textwidth}{X} 向恨他的人，他必當面報應，消滅他們。凡恨他的，他必當面報應，絕不遲延。 \end{tabularx} \\ \\ \relax
7:11 & \begin{tabularx}{0.7\textwidth}{X} 所以，你要謹守我今日所吩咐你的誡命、律例、典章，遵行它們。」 \end{tabularx} \\ \\ \relax
7:12 & \begin{tabularx}{0.7\textwidth}{X} 「你們若聽從這些典章，謹守遵行，耶和華－你的神必照他向你列祖所起的誓，對你守約，施慈愛。 \end{tabularx} \\ \\ \relax
7:13 & \begin{tabularx}{0.7\textwidth}{X} 他必愛你，賜福給你，使你人數增多，也必在他向你列祖起誓要給你的地上賜福給你身所生的，你地所產的，你的五穀、新酒和新的油，以及你的牛犢、羔羊。 \end{tabularx} \\ \\ \relax
7:14 & \begin{tabularx}{0.7\textwidth}{X} 你必蒙福勝過萬民；你沒有不育的男人和不孕的女人，牲畜也沒有不生育的。 \end{tabularx} \\ \\ \relax
7:15 & \begin{tabularx}{0.7\textwidth}{X} 耶和華必使一切的疾病遠離你；你所知道埃及各樣的惡疾，他不加在你身上，反要加在所有恨你的人身上。 \end{tabularx} \\ \\ \relax
7:16 & \begin{tabularx}{0.7\textwidth}{X} 你要吞滅耶和華－你的神交給你的各民族，你的眼目不可顧惜他們。你也不可事奉他們的神明，因為這必成為你的圈套。 \end{tabularx} \\ \\ \relax
7:17 & \begin{tabularx}{0.7\textwidth}{X} 「你若心裡說，這些國的人數比我多，我怎能趕出他們呢？ \end{tabularx} \\ \\ \relax
7:18 & \begin{tabularx}{0.7\textwidth}{X} 你不必怕他們，要牢牢記住耶和華－你神向法老和埃及全地所行的事， \end{tabularx} \\ \\ \relax
7:19 & \begin{tabularx}{0.7\textwidth}{X} 你親眼見過的大考驗、神蹟、奇事、大能的手和伸出來的膀臂，都是耶和華－你神領你出來所施行的。耶和華－你的神也必照樣處置你所懼怕的各民族， \end{tabularx} \\ \\ \relax
7:20 & \begin{tabularx}{0.7\textwidth}{X} 並且耶和華－你的神必派瘟疫 攻擊他們，直到那剩下而躲起來的人都從你面前滅亡。 \end{tabularx} \\ \\ \relax
7:21 & \begin{tabularx}{0.7\textwidth}{X} 不要因他們驚恐，因為耶和華－你的神在你中間是大而可畏的神。 \end{tabularx} \\ \\ \relax
7:22 & \begin{tabularx}{0.7\textwidth}{X} 耶和華－你的神必將這些國從你面前漸漸趕出；你不可迅速把他們消滅，免得野地的走獸多起來危害你。 \end{tabularx} \\ \\ \relax
7:23 & \begin{tabularx}{0.7\textwidth}{X} 耶和華－你的神必將他們交給你，大大擾亂他們，直到他們被除滅。 \end{tabularx} \\ \\ \relax
7:24 & \begin{tabularx}{0.7\textwidth}{X} 他又要將他們的君王交在你手中，你必從天下除去他們的名；必無一人能在你面前站立得住，直到你把他們除滅了。 \end{tabularx} \\ \\ \relax
7:25 & \begin{tabularx}{0.7\textwidth}{X} 你們要用火焚燒他們神明的雕刻偶像；不可貪愛偶像上的金銀，也不可私自收起來，免得你因此陷入圈套，因為這是耶和華－你神所憎惡的。 \end{tabularx} \\ \\ \relax
7:26 & \begin{tabularx}{0.7\textwidth}{X} 你不可把可憎之物帶進你的家，否則，你就像它一樣成為當毀滅的。你要徹底憎恨它，極其厭惡它，因為這是當毀滅的。」 \end{tabularx} \\ \\
[1ex]
\hline
\hline
\end{longtable}
$^{1}$好的 姐妹們安安.
今天有點忙 一腳踢.
歡迎大家再一次回來.
我們繼續講我們這個申命記的訊息.
去到這個系列.
我們叫做上帝也講KPI問號的最後一講.
整個系列其實我們一直集中在講什麼呢.
就是在講摩西向新一代的以色列人.
去重新去提問.
究竟昔日耶和華和這群以色列人.
在荷列山上怎樣立約.
亦都說怎樣去頒佈一些誡命給他們.
叫他們記得要遵守.
當我去想這件事的時候.
其實我經常都會想.
我們這群作為崇尚自由的現代人.
我們很多時候一聽到別人告訴我們要守規矩.
有些什麼規條.
例如一些校規的時候.
我們會有什麼感覺呢.
那個感覺通常都是覺得.
我們的自由被剝奪.
我們有些東西好像在束縛我們一樣.
會覺得我們個人自由的發揮受到限制.
而這種想法其實在我們的信仰裡面.
我們都會這樣想的.
我們都很容易會有同一個的思想.
就是在講.
我們經常看到聖經很多時候都告訴我們.
做得什麼不做得什麼.
有很多的規條在裡面.
我們會無形間想.
其實上帝是不是都在用這些規條.
去規範我們的人生呢.
我們會想.
究竟是不是上帝在用這些準則.
就是他設定了這些規條之後.
成為了一個準則.
去衡量我們是不是一個所謂好的信徒呢.
所以這個系列為什麼會叫做.

$^{41}$「上帝也講KPI」的問號呢.
就是因為這個原因.
如果為免了我們大家當中有人不知道什麼叫做KPI.
稍微簡單解釋.
KPI的全名就是.
Key Performance Indicator.
就是用來做什麼的呢.
就是用來評核一個公司.
或者那個員工的表現的一些指標.
所以這個系列.
其實我們最想做的就是什麼呢.
就是很想探討一下.
其實我們這個信的上帝.
是不是都是這樣的呢.
他頒布十誡給我們的時候.
其實是不是就是成為了那個評核我們的基準.
去批判我們的人生做得好與不好呢.
那如果你上兩個星期有回答的話呢.
那理論上你應該已經有答案了.
就很明顯不是了.
那這個系列的頭兩講講什麼呢.
那就是Ivan和阿松.
其實透過《申命記》的第五和第六章的經文.
去提醒我們很多東西.
就是在講.
上帝是那個主動與人立約的神.
他又提到上帝是願意先祝福我們.
繼而先管教我們的上帝.
就好像在講的.
他當初拯救以色列人的時候.
同樣他是先帶他出埃及的.
之後才頒布律例的.
那他也沒有用剛剛所說的指標.
就是十誡那些他本身還沒頒的.
都沒有這些指標在裡面.
他計算一下究竟這群以色列人.
其實是不是值得上帝去救呢.
他沒有想這件事.
他反而做的就是選擇先跟這群以色列人去立約.
然後才將律法賜給他們.

$^{81}$然後Ivan又提到.
他頒布的那些十誡.
其實不是一些憑空想像出來的定下的規則.
他不是要設一些這樣的規則.
去奴役那群以色列人.
而是針對當時本身住在迦南地的那群迦南人.
他們一直以來的很多不同的風俗.
很多文化.
十誡是針對這些東西.
從而去設計出來的.
為了要保護那群以色列人.
預備他們即將要進入迦南地的時候.
他們將會面對的種種的挑戰.
種種的文化衝擊.
是在講這件事.
而上個星期阿松又提了.
在講的是.
摩西為什麼要跟那群以色列人去解釋.
究竟他們為什麼要遵守律法這件事.
他說的是希望新一代的以色列人.
不要很單純地按照上一代那群人.
那個第一代的以色列人.
他之前傳了什麼.
叫我做什麼.
我就單純跟著來做.
成為了那個所謂的舊有的人傳下來的傳統.
為了守而守.
他說的是.
摩西要講的原因是.
很希望這群新一代的以色列人.
能夠很真實地去理解到.
要守律法背後的意義.
以致他們可以有真正的能力.
或者起碼有願意的心.
去嘗試去遵守.
去嘗試用這個能力去做到.
最終能夠敬畏耶和華這件事.
所以在講的也是為了讓那群以色列人.
可以一直在祝福上帝對他們的這種祝福.
在地上能夠一代一代地.

$^{121}$到最終.
到他們有一天.
真的在迦南地得地為業.
能夠定居之後.
他們可以在那個地區.
改變那個地方的文化.
建立到一個真真正正.
敬畏上帝的一個文化,一個環境.
當我聽完他們兩個的分享之後.
一方面我當然覺得很好.
因為在講的是.
上帝頒佈十誡的原因.
他們講了.
而摩西為什麼要傳.
以色列人為什麼要守.
他們又解釋了.
很完整吧.
那我就想.
糟了 那我講什麼.
我還有什麼好講呢.
好像整段都講完了.
不過.
我覺得很好.
上帝在我預備的過程裡面提醒我.
申命記的特點是什麼呢.
申命記是在講摩西重新將.
由上帝拿到的那個律法.
再次解釋給這群新世代的以色列人去聽.
去理解.
那我又會想.
其實摩西當時也是這樣.
那我重複又有什麼所謂呢.
也是重複摩西做的事情而已 不是很重要.
最重要是什麼呢.
就是忠於上帝的說話.
繼續講解給大家聽.
至於上帝到最終.
怎麼用.
怎麼讓你們聽到什麼.
都是上帝的事情.

$^{161}$我也不可以猜太多.
所以盼望.
今天上帝跟大家親自說話.
讓我們一起祈禱開始好嗎.
上帝我們也多謝你.
多謝你因為你是親自與人對話的上帝.
所以上帝我求你今天祝福我們.
祝福我們今天去到你的信息裡面.
去到你的話語當中的時候.
上帝我們不是很單純地看文字.
我們不是很單純地看聖經以前發生的故事.
而我們仍然相信上帝到今天.
你仍然會跟我們說話.
所以主我求你大大你祝福我們.
當我們今天回到這裡裡面的時候.
願上帝你親自打開我們每個人的耳朵.
讓我們有耳當聽的就應當聽.
聽我們禱告教託.
奉耶穌的特性明智祈求.
Amen.
好.
雖然我剛剛說.
這個系列所說的經文.
都是環繞著上帝頒佈十誡的這件事.
內容雖然很近.
但其實過去甚至到今天.
我們是沒有重複過經文的.
阿松和Ivan分別在用第四章的尾.
一直說到第六章.
今天我和大家看的就是第七章.
我們會看的經文就是七章的一至二十六節.
不過由於經文都比較多.
如果一次過讀.
我想大家到最後又會忘記這些經文是什麼.
所以我決定和大家分段落去看.
和大家一起讀.
然後在每個段落再和大家講解一下.
才再去下一個段落.
我們先看第一個段落.
第一個段落就是七章一至六節.

$^{201}$好,我想邀請大家看到.
1,2,3.
耶和華,你臣領你進入要德為帝.
從你面前講出許多國文.
就是黑人,蟑螂,殺人.
阿摩利人,迦南人,比利西人.
希美人,耶布斯人.
共七國之民.
都比你強大.
耶和華,你臣將他們交給你擊殺.
那時你要把他們滅絕淨盡.
不可與他們立約.
也不可連戌他們.
不可與他們結親.
不可將你的女兒嫁他們的兒子.
也不可叫你的兒子娶他們的女兒.
因為他必使你兒子轉離不更從世.
去事奉別臣.
以至耶和華的怒氣向你們發作.
就速速地向你們滅絕.
你們卻要這樣待他們.
拆毀他們的祭壇.
打碎他們的柱像.
砍下他們的木偶.
用火焚燒他們雕刻的偶像.
因為你歸耶和華你臣為聖傑的民.
耶和華你臣從地上的萬民中揀選你.
特作自己的子民.
謝謝大家.
這個段落說什麼呢.
其實摩西很直白.
一開始就跟這班新一代的以色列人說.
就是說你們將來.
進到迦南地之後得地為業.
要完完全全地將剛才那一撮.
迦南地的七個不同國家裡面.
即是不同民族的居民趕走.
不可以連戳他們.
不可以跟他們立約.
不可以跟他們結親.

$^{241}$摩西很清楚地說.
他說如果你們這樣做.
你們的後代就會受著他們的文化.
他們這班人的影響.
走去事奉其他的神明.
以致最終令到耶和華會發怒.
去滅絕這班以色列人.
所以這段其實正正就是.
之前Ivan如果你有聽的話.
他所分享的就是.
摩西正在預備這班以色列人.
要他們getting ready.
他特別地祝福他們的是甚麼呢.
就是他們將來要進到迦南地之後.
他們一定要做的就是要除滅一切的偶像.
除去當地一切叫.
耶和華所厭惡的那些拜偶像的一些風土習俗.
摩西正在預備以色列人的心.
他正在預備他們將來.
很快就會面對的那些不同的文化的衝擊.
他正在提醒的是這班以色列人.
你們要小心.
你們會慢慢被同化.
你們會慢慢被影響.
你寧願將他們趕盡殺絕.
都不要讓這些人和事.
影響到你們自己和上帝之間的關係.
摩西將這件事看得非常重要.
不過.
我們生活在香港.
是不是.
其實我們有時候很難明白.
為甚麼要這麼盡.
因為我們生活的環境都是.
香港是一個多元文化的.
多元文化交融了的一個城市.
很多不同的文化都在這裡.
我們會想其實多元文化有甚麼不好呢.
是不是.
你不會覺得旁邊.

$^{281}$可能在拜祭其他東西.
走過一條街已經有.
有甚麼影響對我們來說.
是不是 我們會這樣想.
不過如果真的這樣想的話.
或許我們是放輕了一些東西.
我們本身生存在這裡.
可能覺得不是一件甚麼事.
但是你不可以否認的就是.
其實文化,習俗.
對一個人的影響.
是非常之大的.
給大家一個例子.
就是大家有沒有.
看Instagram的views.
嘩 為甚麼來的.
對了 IG.
它這個views的功能.
大家知不知道是甚麼時候出的.
其實不是很久的.
其實是 我看網上資料.
是2020年8月5日.
是其實來了這個世界.
不過是五年多的事.
但是.
你看到你現在的生活.
你現在看到你所有.
大家打開的網上平台.
這個功能.
已經是侵佔了你的生活.
有的 我上網找AI.
我老婆教我 要這麼多AI.
那我就AI一下.
那時候它有解釋的.
為甚麼views會去得這麼快呢.
就是會增長得這麼快呢.
其實它有一個好處.
就是它的娛樂性.
和那個新奇感.
它是會令到大家看的人.

$^{321}$其實那個滿意度.
是會大大提升.
為甚麼呢.
因為你以前煲劇的時候.
你是要用很多時間了解很多東西的.
但是現在那些.
短視的頻.
短視頻.
其實你不需要理解整個故事的人物.
主角的架構.
他的背景是甚麼.
你很快就知道.
他一句說話交代.
甚至大家IG裡面看的東西.
你不需要理解很多東西.
你純粹看見他旅行.
帶一些東西出來.
你就知道他去了那裡旅行.
他在吃甚麼.
全部都很快就知道.
但是呢.
這個Instagram的views呢.
對了 今天不知道發生甚麼事.
對了 OK.
原來過度使用是會有很大影響的.
根據美國的研究.
如果過度使用views.
其實是會刺激這個大腦.
過度釋放多巴胺.
而導致成了癮.
而亦都會因為這樣.
會影響到我們長期看的話.
會降低我們對於事情的注意力.
以及我們的控制.
它叫做抑制控制.
簡單來說就是注意力集中的問題.
他又說的就是.
如果每天你看多於一個小時的短的影片.
我們看的人的專注力明顯是會下降的.
而這件事是不關.

$^{361}$你是否成人 青年 還是小朋友的事.
所有人都會的.
而另外亦都會有心理健康的影響.
他就說你看多了.
就會增加這個焦慮和壓力.
尤其你睡之前如果使用的話.
是會影響你睡覺.
因為你會放大你的情緒.
因為看的東西很快.
所以整個人都會load得很快.
而亦都有另一個心理健康的影響.
就是說那些用戶如果要看得太多這些東西.
因為剛才說他們.
share他們的生活share得太過豐富了.
你會常常看得多的時候.
你就會比較自己和他人的生活.
哇 從而產生自卑.
從而會令到你對所有事情都不感興趣.
因為裡面的東西你會覺得自己.
怎樣都追不到的了.
所以呢.
亦都影響了你日常的工作.
會拖延了你日常的一些要做的事情.
而現實裡面.
有些人說他聽過我說.
我都有說的就是我有個姪女.
同樣都很喜歡看一些視頻.
亦都影響到她.
其實連坐定定看一套戲.
都是有困難的.
都是一個專注力不足的問題.
所以我們在說的是.
其實在這個例子見到.
一個文化的出現.
其實很容易就會對一個人.
產生一個很大的影響.
甚或他的日常生活有很多的改變.
所以剛剛說.
我們生活在香港.
雖然我們會覺得多元文化沒有問題.

$^{401}$就正如其實迦南人本身都不會覺得.
迦南的那些文化會有問題的.
但是摩西卻在這個位置.
他要提醒那班以色列人.
我要提防這件事發生.
而在第六節.
剛才看的經文第六節.
其實摩西是有解釋的.
解釋為什麼以色列民要做那一籌.
要趕他們走.
要趕盡殺絕那些.
他的原因是解釋了.
原來以色列民的身份尊貴.
是屬神的一個聖潔的子民.
在下一個段落摩西繼續說.
想邀請大家跟我一起讀.
第七至十六節.
一二三.
耶和華尊愛你們.
簡選你們.
並非因你們的人數多於別民.
原來你們的人數在萬民中是最少的.
只因耶和華愛你們.
又因要守他向你們列祖所起的誓.
就用大能的手領你們出來.
從遺牢之家救贖你們脫離埃及王法老的手.
所以你要知道耶和華你的神.
他是神.
是信實的神.
向愛他守他誡命的人守約.
思池外.
直到千代.
恨恨他的人當面報應他們.
將他們滅絕.
凡恨他的人必報應他們.
決不辭而.
所以你要遵守遵行我今日所吩咐.
你的誡命律例典章.
你們果然聽從這些典章謹守遵行.
耶和華你神就必照他向你列祖所起的誓守約.

$^{441}$思池外.
他必愛你.
賜福於你.
使你人數增多.
也必再向你列祖起誓.
應許給你地上.
賜福於你.
地所殘的.
並你的五穀生酒和油以及牛毒.
楊高.
你必蒙福勝過萬民.
你的男女沒有不能生養的.
生畜也沒有不能生育的.
耶和華必使一切的病症離開你.
你所知道埃及各樣的惡疾.
他不加在你身上.
只加在一切恨你的人身上.
耶和華你神所要交給你的一切人民.
你要將他們除滅.
你眼不可顧惜他們.
你也不可事奉他們的神.
因這必成為你的忘羅.
由我們剛剛看的一至六章.
其實摩西本身在一至六章的時候用了很多「他們」.
都是集中講述那些迦南人他們拜偶像等等.
去到這個段落我們看到的是.
摩西講了很多「你們」.
摩西將那些以色列人的目光.
放在他們自己身上.
他說的是你們這些以色列民.
其實是上帝所揀選.
是一群蒙福的子民.
是上帝對他們這一群人的一個肯定.
對他們身份的一個肯定.
上帝為什麼要揀選以色列人.
上帝為什麼要刻意說我尊愛這一群以色列人.
經文沒有說.
摩西沒有說.
他只是在說他們愛蒙上帝的恩寵.
被上帝揀選.

$^{481}$不是因為他們族大人多.
也不是因為他們本身有什麼厲害.
純粹是出於上帝的一份恩典.
而這份恩典.
不單止是過去耶和華.
在埃及地救他們出來.
更是在說將來當他們進入這個應許之地的時候.
他們會得到更大的祝福.
他們後裔會有很多.
五谷有豐盛.
新畜,綠畜興亡.
還要是在埃及地的時候.
他看到一切的災病.
全部都與這群人無關.
但有一個條件.
條件就是他要求這一群以色列民.
要用行動去回應上帝.
他要求的是這群以色列民.
去遵守耶和華頒給他們的律例典章.
他們要用這個方式.
去回應上帝對他們的愛和揀選.
從而突出這個被揀選的身份.
突出這個上帝揀選這群聖潔子民的身份.
這個段落摩西其實是在提醒一群新一代的以色列人.
你要看重你的身份.
他提到以色列人這個身份是不同的.
這個身份是上帝給他們的一個肯定.
是尊貴,是聖潔的.
亦都是這群以色列人.
為甚麼會得到很大的福氣的一個很重要的原因.
摩西為甚麼要特意去.
這群以色列人去肯定自己的身份呢.
其實不難明的.
因為一個人的身份不同.
其實是很容易影響那個人.
如何看待一些事情.
而他對那些事情的理解.
亦會因著你的身份不同.
而會有分別.
例子是甚麼呢.

$^{521}$就用上星期我們為這群雲福縣居民禱告為例.
可能對於某些官員來說.
其實給三個小時回去收拾東西.
還要一家幾口加起來幾十個小時.
他覺得很合理的 其實是足夠的.
特別是火災後其實你有甚麼東西可以收拾呢.
但如果你是真的住在裡面的居民.
那個地方是我住了大半生的地方.
你給我三十天我都收拾不夠.
又以另一個例子 我之前都用過.
就是我們當成為父親之後.
你會對上帝那個無條件的愛的理解.
是會完全不一樣.
可能我們以前.
即是我啦 這些部分未有小朋友.
是我.
有部分做我爸爸的可以一同理解.
以前你會經常聽到別人說.
上帝怎樣無條件愛你 你會想.
到你今天 你的兒子.
那些甚麼屎尿屁.
你有份要收拾的時候.
你就會明白甚麼叫無條件的愛.
你不會收拾隔壁屋那個.
你不會去隔壁不小心的地方幫他.
你不會的 你只會避開.
但是那個是你小朋友的時候.
你沒得選擇.
當然有部份原因是沒得選擇.
但是你的付出.
你的兒子不會回報你的.
起碼現在這十年八載都不會.
我們看淡一點.
不過.
你仍然會付出.
天父都是這樣.
你充滿罪惡.
你犯罪得罪他.
你只要回到他裡面的時候.
他仍然都會愛你.

$^{561}$所以.
這段經文摩西用了很大的篇幅.
去解釋.
以色列民這個身份.
究竟有多麼的尊貴.
有多麼的特別.
有多麼的好處 有多麼的祝福.
他正正想做的就是想這班以色列民.
去擁抱這個身份.
因為他們擁不擁抱這個身份.
其實就是會影響.
他們對於上帝要頒佈律例給他們的時候.
的一個看法.
他可能跟我們一樣會覺得上帝設一些條件來.
規範我們.
他可以是這樣理解.
但是當他知道原來他這個身份.
是多麼的尊貴 是多麼的聖潔.
是上帝特意要揀選的時候.
他理解到這件事.
也都會成為他們會不會肯繼續堅持.
去遵行耶和華頒給他們一切的律命.
戒命律例典章.
就好像我們今天坐在這裡的.
這裡我看到的應該全部都是.
都是基督徒 是不是?.
你是不是基督徒?.
你擁不擁抱這個身份?.
正正就是你會不會理聖經這件事?.
如果你不擁抱基督徒這個身份.
就等於一個不信的人.
聖經說的話其實與我無關.
但是就是因為你擁抱這個身份.
你就會明白聖經其實為什麼是這樣.
身份的不同.
帶來眼光的不同.
同樣帶來我們對事情理解的不同.
最後一段.
摩西又回到這一章的主題.
我想繼續讀 我們一二三.

$^{601}$「你若心裡說.
這些國的民比我更多.
我怎能講出他們呢?.
你不要懼怕他們.
要牢牢紀念耶和華.
你臣向法老和埃及全地所行的事.
就是你親眼所看見的大事.
嚴神職 歧視和大能的仇.
並伸出來的棒臂 都是耶和華你臣領你出來所用的.
耶和華你臣必照樣待你所懼怕的一切人民.
並且耶和華你臣必打發黃蜂飛到他們中間.
直到那剩下而躲藏躲的人.
從你面前滅亡.
你不要因他們驚恐.
因為耶和華你臣在你們中間是大而可畏的神.
耶和華你臣必將這些國的民從你面前漸漸趕出.
你不可把他們速速滅盡.
恐怕他們的熱受多起來害你.
耶和華你臣必將他們交給你.
大大地擾亂你他們.
直到他們滅絕了.
又要將他們的君王交在你手中.
你就使他們的名從天下消滅.
必無一人能在你面前站立得住.
直到你將他們滅絕了.
他們雕刻的神像你們要用火焚燒.
其上的金銀你不可貪圖也不可收取.
免得你因此陷入囊瓤.
這原是耶和華你所珍惚的.
可珍的物你不可大進加去.
不然你就成了當毀滅的軟硬物一樣.
你要十分厭惡十分珍厭.
因為這是當毀滅的物.
第三個段落.
摩西用了一個假設性的問題開始.
他就跟那班以色列民說.
如果你心裡想著.
那些國內的民比我多.
我怎能趕走他們呢.
我不知道大家聽這句說話的時候.

$^{641}$有沒有覺得摩西有種酸溜溜的感覺呢.
如果大家有印象的話.
以色列人為何會無端端在曠野漂流四十年.
就是當日摩西派了十二探子去窺探迦南泉地的時候.
其中十個探子回報回來的時候.
這十個探子就說.
我們不可以去攻擊那裡的民.
因為他們比我們強壯.
那十個探子見到迦南地的人比他們強.
他們想勸摩西不如我們走吧.
而這十個探子的說話.
也都令到當時全個以色列的所有回眾都埋怨.
他們還說不如我們另立一個領袖.
帶我們回埃及吧.
摩西今天在用這個假設性的問題.
在提醒新一代的以色列人.
你不要重複當年第一代以色列人所犯的錯誤.
你不要再單純地看別人有多厲害.
別人有多多人.
你要記得那個大而可畏的神在你們裡面.
所以你們不需要懼怕.
他再次提到當日耶和華神.
究竟如何幫助以色列人脫離埃及法老的手.
今天也一樣.
他會將那些迦南人的領袖交到你們手上.
摩西在這個段落.
他再一次將那班以色列民的目光.
從很開初怎樣看迦南人裡面的風俗拜偶像.
到了第二個段落.
講到很多以色列民.
你要擁抱你們自己的尊貴身份.
然後再從這班以色列民的身上.
叫他們看看耶和華神.
今天在做什麼.
他今天在你們裡面.
他一步一步地將以色列人所看見的東西.
慢慢地校正.
慢慢地教他們從別人有多厲害.
到看你們的身份有多尊貴.
到你看耶和華神昔日是怎樣.

$^{681}$他今天未來會怎樣.
這裡講的是.
他在讓這班新一代的以色列人看到.
你們未來是要打仗的.
但是打的不是你們.
不是你們自己.
是耶和華神他也同樣在參與.
他和這班以色列人一同去爭戰.
所以.
當這班以色列人.
要知道這場仗.
未來的這麼多場仗.
要依靠耶和華去打的時候.
摩西再一次警告這班以色列人.
在這段經文的尾段.
就是說你將來進到迦南地的時候.
你一定要消滅所有耶和華所珍惠的一些市民.
連用來整偶像的那些金銀.
無論多貴重你們都不要貪.
要完完全全地除去.
看到這裡.
大家有沒有覺得摩西其實也有些.
創傷後遺症.
他開頭說完.
現在又說.
你們不要擺偶像了.
你們不要讓這件事發生在你們的群體裡面.
不要有機會啊 完全不要有機會.
他提完又提.
不過.
如果我們想真一點.
我們是摩西的話.
其實會有這個所謂的.
創傷後遺症其實是挺正常的.
大家還記得.
金牛獨步這個事件嗎?.
長話短說.
就是當年.
摩西上了這個.
荷里山.

$^{721}$是不是荷里山?.
去領十屆的時候.
而山下這一班人因為很久都沒看到摩西回來.
他們就.
碎阿倫.
做了一隻金牛出來.
而金牛代表什麼呢?.
有《釋經書》解釋的就是.
過去如果大家對埃及有一點點認識的話.
他們的神明很多時候都有動物在頭上.
其中一個神明就是用.
牛神.
就是這樣的.
簡單來說就是這班以色列人.
他們.
拿回以前在埃及的一些偶像.
他們盼望.
摩西不見了.
不如我們找回以前的神明去救我們.
去帶領我們.
他們盼望這件事.
結果怎麼樣呢?.
就是激怒了耶和華.
耶和華其實在經文裡面有寫.
就是想殺光這班以色列人.
不過慶幸的是因為.
摩西不斷地.
力挽狂瀾地勸回上帝.
上帝最終肯放過這班以色列人.
沒有當下就殺了他們.
但是無可避免的就是.
摩西在這些事裡面.
在第一代以色列人的身上.
其實他是很深感歷略到一些.
錯誤.
或者不討神喜悅的文化.
其實是會對他們這一群人.
帶來一個很深遠的影響.
如果大家有印象.
我們講申命記的初期.

$^{761}$我們有講的.
申命記是在講給一班新一代以色列人聽.
那代表舊一代是怎麼樣呢?.
舊一代全死了.
除了另外那兩個探子.
講的不是幾個人.
不是幾十人.
不是幾百人.
如果看聖經講的.
只是計算男人也有六十萬.
所以有些學者去推斷.
當時跟摩西離開埃及的人口.
其實有機會是高達二百萬人.
摩西見證著這一班.
一起生活了四十年的第一代以色列人.
一個又一個在曠野裡死.
到最終只剩下迦納和約書亞兩個.
如果我們代入去摩西的心情.
其實我們會有創傷後遺症也很正常.
所以他為什麼要不斷地講完又講.
是因為他在用一個過來人的身份.
他要去千叮萬囑這一代新的以色列人.
不要再跟那一班第一代一樣.
只是用一些世人的方法.
去看一些事情去決定一些事情.
也不要被最初.
那些不同的國家不同的民族的文化習俗影響.
第一代以色列人可能被.
以前他們住在埃及地的時候.
那些文化影響了.
以致他們會有這樣的選擇.
所以摩西知道接下來你們會見到的就是.
迦南地會有另一種的文化出現.
你們要小心.
你們要認真地去遵循.
我頒給你們的一切誡命,律例,典章.
因為你要抓緊.
我耶和華是你們的神.
我們要確信的是.
耶和華神今天仍然在這班以色列人裡面.

$^{801}$去到這裡.
同樣我想問大家.
你又信不信.
耶和華神在我們裡面.
這個系列我們一直在說的就是.
摩西和新一代以色列人重新講十誡.
但他不是在重溫一件事.
他不是單純讀多一次十誡給你聽.
他是在按照對於新一代的以色列人.
他們這班人當時的情況.
他們將要面對的種種事情,種種問題.
他用最初上帝頒布的十誡.
去對著這班人.
按照當時的情況.
加以描述,加以解釋.
加以應用這些律例典章在他們的生命裡面.
他希望的是.
你們真的要在第一代的以色列人身上.
吸取當時的教訓.
以至你們將來走的路.
會有一個不一樣的結局.
這件事令我想起的是什麼呢.
近期很多電影或者電視劇.
過去一些經典的.
過了N年後或者10N年後.
開始拍電影版,開始拍第二集.
對於看過第一集的人來說.
其實會選擇去看或者去戲院看什麼都好.
當然有部分是因著追回那份情懷.
可能為了找回我以前看過這套戲的一些集體回憶.
但是我相信的是同樣有部分的原因.
我們之所以會選擇繼續去看那些電影的話.
其實是很盼望我們故事裡面.
即是我們回憶裡面的那個故事.
可以繼續延續下去.
甚至很盼望那個故事走下去.
會不會有一個更美好的結局呢.
可能我們想看的是這些東西.
在兩代的以色列人的故事裡面.
第一代的以色列人.

$^{841}$到最終其實大部分.
他們都不可以進加拿大.
而第二代的以色列人.
我們看聖經裡面我們知道.
他們雖然最終是進了加拿大得地為業.
但是當時代過去的時候.
他們的後代並沒有繼續按耶和華的吩咐.
他們繼續去做.
而最終搞到以色列國.
到了國破家亡的地步.
不過我想說.
這個故事.
到今天其實還沒有完.
只不過故事的結局.
故事的主角.
不再只是以色列人.
而是變了每一個相信耶穌的人.
今天我們看到上帝的一切說話.
聖經寫下了.
跟當年一樣.
頒佈給我們這群愛主愛神的人.
究竟我們今天去看這些故事的時候.
究竟我們可不可以在第一代.
甚至第二代的以色列人身上.
我們同樣學習如何吸取教訓.
在兩代的以色列人身上.
這個故事的裡面告訴我們.
人不同的回應.
同樣在上帝裡面帶來一個不同的結局.
面對著上帝給我們的愛.
面對著上帝將這些.
很多我們經常都覺得是.
do's and don'ts的保護.
究竟你怎麼看?.
究竟是不是上帝.
將這一切的事情.
這些所謂的枷鎖.
放在我們身上?.
上帝是不是在用這些律例.
典章,誡命.

$^{881}$去平衡我們呢?.
而更重要的是.
這一切的事.
寫在聖經裡面.
我們已經知道.
最重要的是我們究竟如何回應.
我們會不會.
單純地繼續覺得是一條條例?.
抑或我們可以像摩西那樣.
將上帝頒出來的那些律例.
典章,他的心意更好地應用呢?.
他在告訴後世的一代以色列人.
我們同樣是否可以做到這件事呢?.
在今天的社會裡面.
在我們今天的生活當中.
究竟我們是單純地覺得上帝在管我們呢?.
抑或當你明白上帝頒佈這一切的律例.
本身是為了保護我們.
本身是為了愛我們的時候.
我們願意去回應.
願意去跟隨.
甚至就像摩西一樣.
將這些事情傳下去.
傳開去.
讓「敬畏神」這個故事.
能夠繼續延續下去.
所以這星期我設計了一個小功課給大家.
小幼地.
我們也不是什麼祈禱.
自己反省那麼簡單.
我覺得我們太多那些了.
我覺得要實際一點開始.
就是這星期我想大家做一件事.
就是活用這個十誡.
你在十誡裡面選擇其中一個教導.
你嘗試發揮你的創意.
在今天這個社會裡面走出來.
很空泛吧?.
我想了一些例子給大家.
當孝敬父母.

$^{921}$你嘗試在這個星期花一點時間陪伴他們.
當守安息日.
你給自己一點時間安靜在上帝裡面.
與上帝親近.
不可以貪戀別人的房屋.
你嘗試在這個星期除了三餐吃飯.
和必需品之外不要買任何東西.
控制自己.
不要貪戀一些外在的東西.
再有其他的交給你們自己想.
不過.
上帝的說話的確.
到今天其實我們是可以應用的.
只不過是.
究竟我們選擇想與不想.
還是純粹單純擺了條規矩在這裡.
「就這樣吧」.
就算了.
盼望我們聽到.
行到.
活出上帝給我們的一切.
這麼久以來的教導.
我們一起禱告.
向上帝多謝你.
多謝你愛我們每一個.
以致我們在你裡面的時候.
我們不再只是單純追求自己.
想要的東西,想做的事.
願上帝你親自去幫助我們.
讓我們明白你.
讓我們知道.
你縱然在聖經裡面有很多不同的教導.
或許我們今天仍然都當這些東西是一些規矩.
但主要求你幫助我們.
你轉換我們的目光.
讓我們知道這一切的事情.
完完全全是出於你愛我們.
是出於你想保護我們.
所以神我求你你繼續的教導我們.
使用我們.

$^{961}$讓我們繼續學習怎樣將我們的生命獻上給你.
成為活祭.
為你所用.
求主你這樣去帶領我們更多.
聽我們禱告教堂.
奉耶穌基督的聖名致祈求.
Amen.
\newpage



\section{士師記 3:12-30}
\label{sec:032097}
\textbf{士師記系列 (2) 委身那意想不到的神}~~~~謝又生牧師
\newline
\newline
連結: \href{http://www.church.com.hk/acms/content.asp?site=cdc\&op=show\&type=product\&code=032097\&layout=sermon}{\texttt{032097}} ~~~~ 語音日期: 2026-05-03
\newline
\newline
\hyperref[sec:032096]{\small{< < < PREV SERMON < < <}}
~
\hyperref[sec:index_chronic]{\small{[返順時目]}}
~
\hyperref[sec:032098]{\small{> > > NEXT SERMON > > >}}
\newline
\newline
士師記 3:12-30
\newline
\begin{longtable}{cl}
\hline
\hline
章節 & 經文 (和合本修訂版)\\
\hline
3:12 & \begin{tabularx}{0.7\textwidth}{X} 以色列人又行耶和華眼中看為惡的事。耶和華使摩押王伊磯倫強大，攻擊以色列，因為他們行耶和華眼中看為惡的事。 \end{tabularx} \\ \\ \relax
3:13 & \begin{tabularx}{0.7\textwidth}{X} 伊磯倫召集亞捫人和亞瑪力人到他那裡，他就去攻打以色列，佔據了棕樹城。 \end{tabularx} \\ \\ \relax
3:14 & \begin{tabularx}{0.7\textwidth}{X} 於是以色列人服事摩押王伊磯倫十八年。 \end{tabularx} \\ \\ \relax
3:15 & \begin{tabularx}{0.7\textwidth}{X} 以色列人呼求耶和華，耶和華就為他們興起一位拯救者，便雅憫人基拉的兒子以笏，他是個慣用左手的人。以色列人託他送禮物給摩押王伊磯倫。 \end{tabularx} \\ \\ \relax
3:16 & \begin{tabularx}{0.7\textwidth}{X} 以笏打造了一把兩刃的劍，長一短肘，綁在右腿上衣服裡面。 \end{tabularx} \\ \\ \relax
3:17 & \begin{tabularx}{0.7\textwidth}{X} 他把禮物獻給摩押王伊磯倫。伊磯倫是個很肥胖的人。 \end{tabularx} \\ \\ \relax
3:18 & \begin{tabularx}{0.7\textwidth}{X} 以笏獻完禮物的時候，就把抬禮物的人送走。 \end{tabularx} \\ \\ \relax
3:19 & \begin{tabularx}{0.7\textwidth}{X} 但他自己卻從靠近吉甲的雕像那裡轉回來，說：「王啊，我有一件機密的事要奏告你。」王說：「迴避吧！」於是所有侍立在他左右的人都退去了。 \end{tabularx} \\ \\ \relax
3:20 & \begin{tabularx}{0.7\textwidth}{X} 以笏來到王那裡，那時他獨自一人坐在陰涼的頂樓。以笏說：「我有神的話向你報告。」王就從座位上站起來。 \end{tabularx} \\ \\ \relax
3:21 & \begin{tabularx}{0.7\textwidth}{X} 以笏伸出左手，從右腿上拔出劍來，刺入王的肚腹。 \end{tabularx} \\ \\ \relax
3:22 & \begin{tabularx}{0.7\textwidth}{X} 劍柄連同劍刃都刺進去了，肥肉夾住了劍刃。他沒有把劍從王的肚腹拔出來，糞便就流出來了。 \end{tabularx} \\ \\ \relax
3:23 & \begin{tabularx}{0.7\textwidth}{X} 以笏出到門廊，把王關在樓門裡面，就上了鎖。 \end{tabularx} \\ \\ \relax
3:24 & \begin{tabularx}{0.7\textwidth}{X} 以笏出來之後，王的僕人就來了。他們觀看，看哪，樓門鎖住，就說：「他必是在陰涼的房間裡大解。」 \end{tabularx} \\ \\ \relax
3:25 & \begin{tabularx}{0.7\textwidth}{X} 他們等得不耐煩，看哪，樓門仍然不開，就拿鑰匙打開樓門，看哪，他們的主人已經倒在地上死了。 \end{tabularx} \\ \\ \relax
3:26 & \begin{tabularx}{0.7\textwidth}{X} 他們耽延的時候，以笏就逃跑了。他經過雕像那裡，逃到西伊拉。 \end{tabularx} \\ \\ \relax
3:27 & \begin{tabularx}{0.7\textwidth}{X} 他到了那裡，就在以法蓮山區吹角。以色列人跟隨他從山區下來，他在他們前面引路， \end{tabularx} \\ \\ \relax
3:28 & \begin{tabularx}{0.7\textwidth}{X} 對他們說：「緊跟著我！因為耶和華已經把你們的仇敵摩押交在你們手中。」於是他們跟著他下去，佔據了摩押對面約旦河的渡口，不准一人過去。 \end{tabularx} \\ \\ \relax
3:29 & \begin{tabularx}{0.7\textwidth}{X} 那時，他們擊殺了約一萬摩押人，都是強壯的勇士，連一個也沒有逃脫。 \end{tabularx} \\ \\ \relax
3:30 & \begin{tabularx}{0.7\textwidth}{X} 那日，摩押在以色列手下制伏了。於是這地太平八十年。 \end{tabularx} \\ \\
[1ex]
\hline
\hline
\end{longtable}
$^{1}$世界上有很多事是意想不到的.
即是你覺得它是這樣.
其實它不是這樣.
我舉個例子.
你以為這張照片是一個文件夾.
看上去好像萬字夾.
但其實不是.
它是一個公事包.
第二張.
你以為它是一枝樹枝.
但其實不是 它是一座台燈.
第三個.
為何沒有第三呢 因為意想不到.
第四.
你以為它是一本書.
但其實不是 它不是一本書.
它是一個枕頭.
一打開.
為何看不到.
一打開就睡著了.
第五.
你以為它是一塊豬肉.
但其實不是 它是一個手提行李箱.
你想要一個這樣的手提行李箱嗎.
我很想要一個.
拿著豬肉在街上走.
第六個.
很恐怖.
你以為它是手指.
但其實不是 它是肥皂.
拿著手指洗澡.
甚麼是意想不到呢.
因為你覺得事情一定會發生.
那些人一定是這樣.
其實結果是出乎你意料之外.
不知道大家有沒有看過一套戲.
叫做《食神》 周星馳的《食神》.
谷德超說.
真的猜不到.
就是這個意思.

$^{41}$不過可能你覺得《食神》太歐了.
現代有一個國家的領袖.
你真的猜不到.
意想不到有兩種意想不到.
第一是意想不到的壞事.
可能大家知道現在美國和伊朗打仗.
開頭你以為強弱完事.
很快就打完.
相距那麼遠.
誰知一打就打了兩個多月.
還不知道何時完.
搞到世界大亂.
真是意想不到的壞事.
也有意想不到的好事.
David今天生日.
但所有身邊的人.
無論是家人還是朋友.
全部都超忙.
他想著今晚自己過生日.
誰知回到家.
所有人都在.
還有一個驚喜的生日派對.
是意想不到的好事.
在《士師記》第三章.
迦南地督發生了意想不到的大事.
話說迦南日報登頭條.
裡面說什麼呢.
摩訶王遇刺震撼全國.
以色列恐怖分子已乏成功逃脫.
是一個刺殺的行動.
是一個暗殺的行動.
這是頭條.
但當你看詳細的新聞報道.
你會發覺有三件都是意想不到的地方.
第一個意想不到的地方.
就是欺壓以色列人的摩訶人.
竟然是這一班民族.
為什麼我這樣說呢.
因為從以色列人的角度來說.
這個摩訶人是弱小的民族.

$^{81}$是微不足道的.
無論你看《出埃及記》.
或者《民數記》.
裡面分別都記載.
摩訶人因為以色列人的強大人數眾多.
大大懼怕.
所以我就像想當年那樣.
假設這班民族不屑一顧.
但事實上卻是意料之外.
因為這班弱小民族.
竟然聯同阿摩人和阿瑪力人.
打敗了以色列人.
還有以色列人佔領迦南地的象徵是什麼呢.
耶利哥城是一個勝利征服迦南地的象徵.
但竟然被這些弱小的摩訶人佔領.
還將耶利哥城轉變成為欺壓以色列人的基地.
實在意想不到第一件事.
第二件事意想不到的地方.
就是拯救以色列民族.
這個以弗的作為.
並不是我們一般想像的.
傳統的事先拯救方式.
是什麼呢.
我們舉兩個例子.
第一個事先是俄陀列.
那個拯救方式是神的靈.
充滿俄陀列.
於是俄陀列就帶領以色列人打贏仗.
第二個例子.
基典有一樣.
神的使者在基典顯現.
祂說你是大能的勇士.
我差遣你.
於是基典就帶領以色列人打贏仗.
於是我們可能想.
這個以弗.
這個第二個事先.
一定是用這樣的模式拯救以色列.
結果出乎意料.
不是.

$^{121}$拯救以色列的模式並不是打贏仗.
而是暗殺.
出乎意料.
第三個出乎意料是什麼呢.
這個新的事先.
第二個事先.
是一個左撇子.
為什麼我這樣說呢.
因為在猶太人的傳統來說.
猶太人的傳統是用左.
是用右手的.
你看回聖經裡面.
不斷地.
以色列民族不斷地用右手.
神用右手去拯救以色列.
詩篇十八篇三十五節說.
你的右手扶持我.
神也說.
我用我的右手去消滅仇敵.
還有.
復活了的耶穌基督.
坐在傳行者的右手邊.
所以傳統猶太人的文化.
右手比較厲害.
右手代表能力.
右手代表榮耀.
偏偏我們假設.
神的事先一定用右手.
不是.
結果是出乎意料.
這個事先竟然是左右的.
左手仔.
出乎意料.
可能在我們人生裡面.
在我們很多時候.
當我們禱告的時候.
可能我們的潛意識.
都有這種想當然.
覺得一定會發生的心態.
我舉個例子.

$^{161}$神啊.
我祈禱.
你一定要照這樣做.
我祈禱發生的這件事.
一定要在那個時間發生.
我祈禱.
這樣的形式.
一定要這樣的形式出現.
結果一定是這樣.
我們的祈禱有時候.
潛意識就會發覺.
我們的禱告就像過年的揮春一樣.
心想事成.
等一下我想問問大家.
我們同福堂.
灣景樓上面有個大牌.
寫著四個什麼字.
耶穌是主.
就不只是說耶穌是主.
就是說耶穌真是主宰.
什麼叫做主宰.
是祂決定的事.
不是我們吩咐耶穌做什麼.
所以耶穌是主的意思就是說.
耶穌決定一切的事.
為什麼是出人意料呢.
因為祂決定.
不是我決定.
如果是我決定.
就不是出人意表了.
是不是.
耶穌是主的意思是說.
祂真是主宰.
祂不是讓我們.
想當然的方程式去限制.
等一下.
這句話不是我說的.
是聖經這樣說的.
以賽夏書55章8-9節.
讓我們一起讀.

$^{201}$123.
耶穌說.
我的意念非同你們的意念.
我的道路非同你們的道路.
天怎樣高過地.
照樣我的道路高過你們的道路.
我的意念高過你們的意念.
等一下.
神帶領人的方式.
就好像剛才聖經這樣說.
神所做的事.
是超乎我們意念所想.
所以神帶領人.
幫助人的方式.
往往都是超過我們理性.
所能夠.
可以推測.
可以計算.
所以我們就要憑信心來接受.
我再說一次.
神幫助人.
帶領人的方式.
往往超過我們理性所想像.
所以我們要憑信心來接受.
我舉例在聖經裡面.
以色列人攻打耶利哥城.
用我們理性的方式.
就是我們戰略是怎樣.
我們人多不多.
我們有攻城器械行不行.
是不是天時地利人和.
我們會這樣想.
但是神偏偏吩咐以色列人.
你們這群人就圍繞過城13次.
你怎樣.
什麼意思.
當時帶領的約斯蘭完全不明白他在做什麼.
理性上不知道他在做什麼.
唯有用信心去接受.
我舉多一個例子.

$^{241}$另一個是斯基丁.
他帶領一萬人.
去對抗十三萬五千人.
的米地人軍隊.
其實已經強弱懸殊.
但是神險.
一萬人太多了.
減到三百人.
嘩 搞事啊.
完全不明白.
但是基丁唯有用信心去接受.
我們所信的耶穌基督.
是天地萬物的神.
祂說一句話就毀滅這個世界.
但是偏偏.
祂變成你和我一樣.
這麼軟弱的人.
並且羞辱痛苦的.
寫在十字架上.
我們完全理性上完全不明白.
我們唯有用信心來接受.
在我們人生裡.
有很多時候.
我們所想的東西.
神已經超越這一切.
我們所想的東西.
因為神的安排.
我們未必在這個時候.
完全明白.
但是當整件事完成.
整個圖畫完成後.
你才明白神在做什麼.
所以耶穌就是主.
我以前曾經.
也分享過我自己的經歷.
我媽媽離開我.
回到天家主人的安息.
是很短很短的時間.
我和我媽媽的關係非常好.
所以很短的時間.

$^{281}$我完全沒辦法接受這個事實.
所以很長的時間.
我很不開心.
除了這個不開心.
是我媽媽回到天家.
很短的時間離開我之外.
另外一個很大的難處.
就是我要面對.
從小到大.
和我的關係不好的爸爸.
為什麼不好呢.
因為我覺得我爸爸從小到大.
都不是在家.
我覺得他不關心我.
不愛我.
然後因為星賀的問題.
大家搞得很不愉快.
大家沒什麼兩句.
沒什麼聊天.
一聊天大家語氣都不好.
所以很多時間都不歡而散.
我曾經有一段時間.
那段時間很不明白.
不知道為什麼.
神在搞什麼.
但是神往往都是出乎意料的.
我不明白祂的安排.
於是我自己有些不好的想法.
我強調.
我以下說的話.
我已經完全在神的面前認罪了.
為什麼不好的想法呢.
我說神啊.
我媽媽也信耶穌.
我爸爸也不信耶穌.
不如拿走我爸爸吧.
留下我媽媽在這裡.
我們一起回去教會敬拜你.
那不就是最美的事了.
是你自己得察為妖.

$^{321}$那麼多人敬拜你.
不是很好嗎.
但是當然了.
神絕對不會聽我這個奇怪的想法.
因為神是超過我所想像的.
慢慢我知道神的安排.
最後是祝福我的.
包括有一晚.
又是很不開心的時候.
在一個夢裡.
我發了一個夢.
這個夢裡.
我媽媽很開心.
穿了一件很白的光的衣服.
和一群天使去到很光的地方.
在這個夢裡.
我從小到大從來沒有見過.
我媽媽笑得這麼開心.
就好像神藉著這個夢.
跟我說.
你放心吧.
你所愛的媽媽.
今天就在我的平安裡.
在我的懷抱裡.
在我的安息裡.
你一無其他.
所以你放下這個重擔.
我剛才說.
我媽媽回到天家了.
我就要和我爸爸朝夕相對.
我最不開心的就是這件事.
但是為什麼呢.
有這件事為什麼會發生呢.
就是我以前.
我和爸爸真的沒有兩句.
我剛才所說.
但是事實上都是要溝通的.
那怎麼溝通呢.
我爸爸有什麼要說.
就和我媽媽說.

$^{361}$我媽媽就和我說.
我有什麼要和爸爸說.
就和我媽媽說.
我媽媽就和爸爸說.
結果大家可以想像.
我媽媽在家裡的功能.
除了煮飯和做家務之外.
有一個新的功能.
就是人肉傳聲筒.
但是大家可以想像.
如果是這樣的話.
怎麼和好呢.
這個關係怎麼再次去改變呢.
再次去進步呢.
結果呢.
神就將我媽媽.
拿在她的同在.
她的懷抱裡.
結果現實上.
我就每天都要見到我爸爸.
但是這個我覺得不好的東西.
到最後是變成好的.
為什麼呢.
因為你越見得多.
你就開始聊天.
你開始聊天.
就開始溝通.
開始溝通.
就開始了解.
開始了解.
你就開始明白.
原來接觸得多.
我就知道原來.
我上一代的男性.
是不習慣說話的.
不習慣表達的.
他心裡的那些感情.
那些愛.
那些關懷.
完全不懂得說.

$^{401}$但是我小時候.
年輕的時候.
他不說就當他沒有.
但是當你了解多.
明白多的時候.
你就知道原來他這個人是這樣.
你開始了解.
開始明白.
開始接納.
於是我和我爸爸.
因為神的幫助和好.
因為這樣.
我爸爸就信了耶穌.
等於今天.
我爸爸媽媽.
都已經是天價了.
主為人安息了.
但是.
神給我的祝福是什麼呢?.
出人意外的祝福.
就是我今天沒有遺憾.
什麼叫遺憾?.
是一些對的事.
你應該做的事.
但是你不想做.
但是有一天.
你發覺對的.
應該做的.
但是你沒有機會做.
這個叫遺憾.
但是等於神給我的祝福是什麼呢?.
祂出人意外.
祂出人意表.
祂是意料之外的神.
但是祂最後給我的祝福.
就是我再沒有遺憾.
所以我今天想和弟兄姊妹說.
我們的神出人意料.
並不是因出人意料而出人意料.
神出人意料.

$^{441}$最終的目的.
是祝福你和我.
我再說一次.
神出人意料的.
但是祂出人意料.
不是因為出人意料.
為出人意料出人意料.
神出人意料的作為.
祂最終的目的.
是祝福你和我.
為什麼我這樣說呢?.
我們繼續看《釋詩紀》.
在《釋詩紀》裡面.
我剛才說.
以弗拯救以色列的行動.
是一個暗殺.
這個表面上是一個暗殺.
但是在屬靈上.
卻是一個憲制.
為什麼這樣說呢?.
我們看看.
《釋詩紀》裡面的經文.
有幾個字 有幾個原因.
第一個原因.
我們剛才說.
魔壓王伊基倫.
是一個大肥佬.
但是伊基倫的.
希伯來原文.
是肥牛毒.
用回香港的說話.
小肥牛.
現在還有沒有這家.
吃火鍋的餐廳呢?.
小肥羊小肥牛.
簡單來說.
伊基倫就是小肥牛.
肥牛毒是當時.
憲制經常用的牲口.
用羊,白鴿,牛.

$^{481}$肥牛毒是憲制的牲口.
我們再看.
《釋詩紀》三章十七至十八節.
裡面說.
以弗釋詩獻賞禮物.
希伯來原文.
是將憲制的牲口.
帶到祭壇的面前.
這個叫獻賞.
那個禮物是什麼呢?.
禮物的希伯來原文是.
獻給耶和華的貢物.
獻祭給耶和華.
繁祭的禮物.
所以越來越像獻祭.
《釋詩紀》三章二十二節.
裡面形容.
以弗行釋的時候.
將劍插入伊基倫的身體.
連劍都看不到.
插了進去.
這把劍的希伯來原文是火.
整本聖經只是在這個地方.
用希伯來原文的文字.
大家看這幅圖畫.
這件事表面上.
是一個行釋.
刺殺的行動.
暗殺的行動.
但在屬靈上.
卻是一種獻祭.
可能大家會問.
獻祭和今天的禮儀有什麼關係呢?.
這個獻祭.
這個擺上.
這個刺殺的行動.
就是一件獻祭.
和我們今天有兩個訊息說.
第一個訊息.
神是超越一切.

$^{521}$掌管一切的神.
祂不但出人意料.
祂並且是超越一切.
掌管一切的神.
為什麼這樣說呢?.
我們一起看看.
《釋詩紀》.
我們剛才所說的聖經.
是《釋詩紀》三章十九至二十六節.
其實這段聖經裡面.
有一個結構.
頂智慧不單止在聖經的文字.
和我們說話.
聖經的結構也和我們說話.
我們一起看看這個結構.
我解釋一下.
這個結構叫倒影結構.
什麼叫倒影結構呢?.
可能有些頂智慧.
你去旅行.
比如去西湖.
你見到這個湖真的漂亮.
水平如鏡.
好像一塊鏡一樣.
中間有一條橋.
有些頂智慧拿手機.
這麼漂亮.
拍了它.
就見到一條橋.
有一個倒影在下面.
就是這個意思.
所以你見到.
ABC中間是交叉.
下面就是CBA.
就是ABC的倒影.
換句話說.
其實這兩個說話.
有頭一尾和中間.
都是在說一些東西.
你看看A和A的倒影.

$^{561}$分別都有作石之地.
作石之地的原文是什麼呢?.
是偶像的石雕.
就是說你去那個地方.
整個地方都是偶像.
換句話說.
那是拜偶像的地方.
那是供奉偶像的地方.
好了.
兩邊呢.
拜偶像的地方.
中間呢.
BC.
CB.
X.
全個是一個整個.
是一個暗殺的行動.
是拯救以色列的行動.
但是在屬靈上面.
卻是憲制的行動.
它在說什麼呢?.
以佛.
在所有那些.
好像很厲害的那些偶像的面前.
在他們的眼前.
用拜偶像的人.
向神獻祭.
向耶和華獻祭.
這個行動.
而那些偶像只是眼光光地看著.
沒有任何的能力去干涉.
沒有任何的能力去阻撓.
就是說什麼呢?.
神是最大的神.
神是勝過一切的神.
神是超越一切.
勝過一切的神.
他勝過偶像.
勝過那些邪神.
勝過這個混亂的世界.

$^{601}$神已經勝過這個世界.
鄧小平今天跟我們說什麼呢?.
鄧小平說這個世界動蕩不安.
你同不同意嗎?.
俄烏戰爭打了幾年.
中東戰爭之後繼續打下去.
所引發的結果不僅是那些地區的死亡破壞.
並且大家都知道油價高成怎樣.
大家入一缸油.
平時一千元可以入一缸油.
現在一千元可能入一半都不夠.
不僅如此.
不僅只是開車的底子會受影響.
拿來開電車是沒問題的.
所有的成本都上升.
物價上升.
航運上升.
所有的價錢上升.
引發的是什麼呢?.
經濟混亂.
經濟低迷.
地震,天災不斷發生.
AI的急速發展可能令很多人擔心.
我讀完書的工作會不會沒了.
我將來的工作會不會被AI代替.
簡單來說是今天不知道明天的事.
這個世界是充滿動蕩不安的.
在這個動蕩不安的世界.
誰才是最大的依靠.
就是那個勝過一切.
那個超越一切的神.
那個成為我們保障的神.
我們的神為什麼出人意料呢?.
為什麼是我們意料之外呢?.
因為祂掌管一切.
祂明白的是超過我們明白的東西.
所以我們的神掌管一切.
但更重要的是什麼呢?.
我們的神超越一切掌管一切.
但祂這些意料之外的行動.

$^{641}$最終是祝福你和我.
所以這個是非常重要的.
亭姐妹.
我也和大家分享過一個.
我自己的經驗.
就在2002年.
我覺得神父召我全職事奉.
所以我在2002年辭了那份工作.
2003年1月.
打算在同福就全職事奉.
但辭了那份工作之後.
有個跨國公司跟我說.
我給你一個更好的發展機會.
我有些動搖.
但我心裡很清楚.
神父召我用神的方法去事奉祂.
並不是我想的方法去事奉祂.
我記得那一晚.
當我跪在神的面前去認罪.
我和神說我願意信服你的心意.
不是我自己的心意.
我記得那一刻.
我心裡有一個很大的平安.
我心裡也感受到.
神給了我一節聖經.
就是耶利米茲33章3節.
祂說你求供我.
我就應運你.
並且將你所不知道.
又大又難的事指示你.
亭姐妹.
現在可能我們覺得動蕩不安.
但那時候是2002年.
世界一切OK.
周圍一切都順利.
老老實實.
我完全不明白神究竟說了什麼.
你完全不明白這段聖經說了什麼.
但是很奇怪.
這段聖經不用背.

$^{681}$一直都藏在我心裡.
一直都不拿出來.
我完全不明白.
為什麼這段聖經一直記住.
當時是2002年.
我完全不明白祂說了什麼.
但是17年後.
當我接任同福堂主任牧師的時候.
我就完全明白.
什麼叫做又大又難的事.
大家想像2019年之後.
發生了什麼事.
一鑊大過一鑊.
一鑊金過一鑊.
一鑊結過一鑊.
沒辦法用任何以往的經驗.
人的智慧能力.
可以面對所有的事.
但是那時候我就知道.
什麼叫做又大又難的事.
但是最奇妙的是什麼呢.
竟然17年前.
神告訴我.
在17年後.
你將會面對又大又難的事.
最奇妙就是這樣.
並且應祂說.
你將會面對又大又難的事.
但是你放心.
你求過我.
我就應運你.
並且會指示你.
今天是2026年.
我們之前那五六年的挑戰難處.
神很感恩.
因為神已經帶領我們同福堂.
去完全渡過這些危機.
我們的神是超越一切的神.
是不是.
所以我在這裡想和各位姐妹見證.

$^{721}$我們的神是超越一切.
掌握一切.
你沒辦法想像祂.
但是我們的神是越是好的.
我們的神在一切的超乎意料之外.
不是為了超乎意料之外.
出人意料.
而是祂出人意料的作為.
祂最終的目的.
是祂想祝福你和我.
這是第一個訊息.
我們的神是超越一切.
掌管一切的神.
第二個訊息.
當我們看神的作為.
我們看神的說話的時候.
我們要明白我們的神是一個什麼神.
祂的心究竟是怎樣.
祂的心態究竟是怎樣.
究竟祂想和我們說什麼.
第二個訊息.
是神何德渴慕去閱納我們.
我們剛才說獻祭.
獻祭的目的是什麼呢.
獻祭的目的是仰望神赦免我們的罪.
以致我們可以和神復和.
和神建立親密的關係.
這是獻祭的目的.
在獻祭裡面.
以弗所做的行動是一個獻祭.
但以弗是不知道的.
他所做的就是一個完全的毀身.
什麼叫毀身呢.
傾罪他生命一切擺上彼實.
包括他這個暗殺行動.
可能無論成不成功.
可能他沒有命.
他冒著這個風險.
擺上他一切的生命.
並且他做了很周詳的部署.

$^{761}$伊基倫是一個很貪心的人.
以弗為了討他的喜悅.
送了大量的禮物給他.
博取伊基倫的好感.
並且送禮物.
當禮物擺到整個店滿的時候.
他叫人搬行李.
搬禮物的人全部走.
剩下我一個在這裡.
他一個人面對伊基倫.
在聖經裡面說.
伊基倫是左手便利.
不只是他左撇子.
其實是代表他右手殘障.
大家可以想像一下伊基倫.
他如何體而弗.
我是一個大胖子.
你是一個人.
並且右手殘障.
我一隻手指尾都打贏你.
砸都砸死你.
於是伊基倫就非常輕視他.
藐視他.
疏於防範.
結果如何呢.
結果.
以弗一擊即中.
但這件事.
以弗怎樣都猜不到.
他正在做一個傳言的毀身.
但神將這個傳言的毀身.
變成了一個憲制.
自己閱立這個憲制.
所以後來以色列人大獲全勝.
在這裡大家砌一砌這幅圖畫.
這個傳言的毀身.
是以弗自己做的.
但神將這個傳言的毀身.
變成了一個憲制.
然後這個憲制獻給自己.

$^{801}$自己又閱立這個憲制.
而祝福他的子民.
並且與這些子民磨好.
大家看不看到這幅圖畫.
神自己將這個毀身.
變成憲制.
自己又閱立.
換句話說.
神的心是怎樣.
大家看不看到.
神多麼想閱立我們.
多麼想和我們復和.
簡單來說.
阿彼斯你申請銀行貸款.
你肯申請就批了.
就是這個意思.
神多麼想我們與他復和.
神多麼想我們毀身.
你一毀身.
我就和你復和.
我就閱立你.
閱立你一切的擺相.
就是這個意思.
看不看得到神的心.
這個說話.
原來.
以弗的毀身.
在神的眼裡是一個憲制.
是蒙神閱立的.
這個說話又不是我說的.
是羅馬書說的.
我們看看羅馬書十二章一節.
我們一起讀.
一二三.
所以弟兄們.
我以神的慈悲勸你們.
將身體獻上.
當作活祭.
是聖潔的.
是神所喜悅的.

$^{841}$你們如此事奉.
乃是理所當然.
毀身等於憲制.
什麼是毀身呢?.
一個基督徒毀身.
如果一個基督徒向神毀身.
他不只是口中說一些高言大字的東西.
也不是說我認識很多聖經.
認識很多神學.
而是用實際具體的行動.
去回應神.
信服神.
遵守神的吩咐.
以弗就是一個榜樣.
擺上他一切.
全然的毀身.
等到今天神都挑戰我們每一位.
你願不願意將你生命的主權交給神.
你願不願意將你的時間交給神.
你願不願意將你的金錢交給神.
你願不願意將你的財幹.
擺上交給神.
完全擺在神的手裡.
你過一個以神的旨意為中心的生活.
我再說一次.
今天神都挑戰我們.
你願不願意將你的生命主權交給神.
你願不願意將你的時間交給神.
你願不願意將你的金錢交給神.
你願不願意將你的財學.
才華.
才能.
交給神.
而過一個以神的旨意為中心的生活.
這是全然的毀身.
也是獻祭.
也是蒙神喜悅.
在去年的七月.
我講道的時候.
曾經播過我們同福堂一個姐妹.

$^{881}$卓菲娜的見證短片.
當時卓菲娜去年患上末期的肺癌.
她的癌細胞擴散到她全身.
當時她這樣形容自己.
醫生跟她照X光.
她全身的癌細胞佈滿全身.
就好像你晚上看到天上的星.
天上滿天星.
就是這個意思.
她日常生活有什麼影響呢?.
無論幾輕微的活動.
包括躺下,起身,拿筷子.
她背部都極痛.
痛到什麼程度呢?.
她要三種強烈的止痛藥才能止痛.
大家要明白痛法.
卓菲娜這麼多年跟神結了親密關係.
卻成為她面對人生極大挑戰.
一個重要的信仰基礎.
在她的口中.
神是超越一切的神.
神是掌管一切的神.
神是出人意料的神.
但卻神的出人意料的最終目的.
是想祝福我們.
所以God is good God.
God is good God.
所以在她心裡.
無論她跟什麼人說.
她心裡流露出一種平安和穩妥.
不是做出來.
沒有人可以在這樣的處境裡.
做到一些平安穩妥.
我記得那一次.
我跟師母上去探望她.
她一直在形容自己的狀況的時候.
我跟師母眼有淚光.
但我看到菲娜那種信心.
那種積極去面對當年的處境.
我心裡想起兩個字.

$^{921}$偉神.
全然的偉神.
將她生命的主權完全交在神的手裡.
這全然的偉神.
在神的眼裡卻成為一個憲制.
蒙神越立.
過了大半年.
她在醫療的進展是怎麼樣呢.
今天的Trina.
她本來的癌指數是48.
但今天已經跌到4.5.
一個正常人的指數.
我剛才說她照X光.
全身的癌細胞就像滿天的星一樣.
但今天只剩下肺部6mm的小腫瘤.
我們繼續為她祈禱.
所以Trina上星期六.
很平安喜樂地出來說話.
你看到她肥肥白白嗎.
Hallelujah.
鄧梓慧 我在這裡也想和大家說說聽見證的心態.
可能有些人會說.
「卓天蘭 我不是每個人都是這樣的」.
「她就是這樣 不關我的事」.
可能很多香港人聽見證的時候.
第一句的反應是「不關我的事」.
但是鄧梓慧.
我們聽見證不是聽人的作為.
是聽神的作為.
所以如果你聽見證.
除了你自己發生之外.
這個世界的所有事都不關你的事.
你就會少聽很多事.
所以我們是聽神的作為.
我們未必每個人都像Trina那樣.
但是你相不相信.
我們每一個的際遇處境都不同.
但是神知道的.
既然神那麼厲害.
祂必定可以在你的處境.

$^{961}$用祂的方法去祝福你.
你明明神是掌管一切 超越一切的神.
祂明白你.
祂必定可以在你的處境祝福你.
況且神出乎意料.
最終的目的是祝福你和我.
所以我們每個人未必像Trina的處境.
大家的際遇都不同.
但是神都給大家一個挑戰.
都是一樣的挑戰.
你願不願意將你的生命主權交給神呢?.
你願不願意將你的時間交給神呢?.
你願不願意將你的金錢交給神呢?.
你願不願意將你的財幹交給神呢?.
過一個以神的心意為中心的生活.
最後.
都有兩種人生給大家選擇.
第一種人生.
就是一個沒有神的.
或者就算你相信耶穌.
都不以神為中心的人生.
不知道大家有沒有看過一部電影.
叫《東方不敗》.
《少傲江湖》是《東方不敗》.
這邊王瞻作了一首詩.
天下風雲出我背 一入江湖歲月吹.
妄圖霸業談笑中 不勝人生一場醉.
有一次小朋友說.
你鼓勵人家喝醉酒.
錯.
其實這個意思不是這個意思.
妄圖霸業談笑中 不勝人生一場醉.
是什麼意思呢?.
他說當人我行我素.
追逐自己的野心.
去建立很多很多不屬於神的東西.
在你生命裡跟神的心意去做的時候.
到最後的結果.
就好像我們《聖經》的傳道書說.
人生就是空虛.

$^{1001}$不如喝醉就好了.
就是這個意思.
不是教人喝醉酒.
這是一種人生.
另外一種人生.
就是《羅馬書八章》二十八字說的.
頂智和我一起讀.
一二三 我們曉得萬事互相效力.
叫愛神的人得益處.
就是按他旨意備詔的人.
用回我們的說話來講是什麼意思呢?.
我們人生就像拼圖一樣.
你未拼之前可能大家都拼過圖.
你拼圖的時候.
你未拼的時候.
喂 什麼來的?.
你拿著一隻 一塊 是什麼來的?.
換句話說 我們今天可能在人生裡面.
這一刻好像沒什麼意義.
這一刻很慘.
這一刻可能很苦毒.
這一刻可能很痛苦.
這一刻可能很迷惘.
很迷失.
不知怎樣.
但是我們的神.
是會將這些拼圖.
在你的人生裡面一塊一塊地砌下去.
到最後會成為美麗的圖畫.
就好像這幅圖畫一樣.
重點是什麼呢?.
這幅圖畫是給予什麼人呢?.
就是給予愛神的人.
就像剛才以輝的見證.
全然毀身的人.
弟兄姊妹 你願意選擇哪種人生呢?.
我們一起祈禱.
當我們禱告的時候.
讓我們來到神的面前.
今天神的心意在我們生命裡顯明.

$^{1041}$當我們安靜禱告的時候.
聖靈要帶領我們.
要跟我們說話.
將神的心意顯明在我們的生命裡.
今天聖靈對大家有兩個呼召.
第一個呼召是.
對已經信於耶穌的弟兄姊妹.
今天無論你的生命光景是怎樣.
你信了多久.
可能你人生裡都走過不同的階段.
不同的道路.
今天有人可能很生氣.
亦都可能有人正在面對難處.
亦都有人可能有一些迷失.
有很多人亦都有其他不同的際遇.
但今天神很呼召你.
你是否好像以輝一樣.
將你人生的主權.
生命的主權交給神.
跟從神的旨意.
不是你自己的旨意去走.
雨後的弟兄姊妹.
我邀請你在你的座位張開你的手.
你願意將你生命的主權.
你的未來.
你心裡面最記掛的.
你覺得最重要的.
這一切都在神的計劃裡面.
你願意將這一切主權.
交在神的計劃裡面.
今天神呼召你.
如果你願意這樣做.
你願意回應祂.
我邀請你在你的座位張開你的手.
請我主祀我父耶穌的名.
去祝福這些弟兄姊妹.
這些弟兄姊妹是願意回應你.
親近你.
他們來到你的私人寶座面前.
他們祝福你張開雙手.

$^{1081}$主要我們來到你面前是無保留的.
我們只是將這一切生命的主權交給你.
我們深信我們出人意外的神.
我們出乎意料的神.
是一個對我們最好的神.
你為我們安排的永遠都是最好的.
所以今天我們仰望神你的帶領.
我們仰望神你的同在.
我們很深信你為我們編織的圖畫是最好的.
當然我們每一個人的圖畫都不同.
但我們知道這個圖畫是對我們最好的.
所以今天我們回應你的時候.
我們都仰望你.
我們全心信靠你.
我們不靠我們自己的才智能力去做事.
我們願意將這一切交在你的手裡面.
主人你會使用.
你會帶領.
我們會與你同工.
你亦都會與我們同工.
主人我們將這一切交上的時候.
都成為一個美好的獻祭.
是蒙你血脈.
主你為這些弟兄姊妹開路.
讓他們見到耶穌基督再不一樣.
他們的生命是比你更新.
給他們新的能力.
無論他未來的前面是怎樣.
神你給他們能力衝破一切的難關.
讓他們見到神在你的結果的時候.
祂要將榮耀歸給你.
當我們繼續閉上眼睛祈禱的時候.
我不知道在座有沒有一些朋友.
可能你回了同福的聚會一段時間.
但你從來都沒有告訴別人你信耶穌.
今天神都呼召你.
耶穌要為你編織.
拼埋一幅美麗的圖畫.
如果你願意領受這幅圖畫的時候.
你必須要做一件事.

$^{1121}$就是與耶穌建立關係.
你怎樣與耶穌建立關係呢.
就是打開你的心.
去接受耶穌成為你的主.
所以如果在座當中有一些朋友.
今天你願意打開你的心.
去承認耶穌是你的主.
我邀請你在座為舉手.
有沒有這些朋友呢?.
今天你願意打開你的心.
接受耶穌.
去接受耶穌為你編織的這幅圖畫.
砌成這幅圖畫是美麗的圖畫.
你願意接受這幅圖畫.
你必須要打開你的心.
承認耶穌是主.
如果有的,有朋友.
你願意打開你的心.
接受耶穌成為你的主.
我邀請你在座為舉手.
有沒有這些朋友呢?.
有沒有這些朋友在我們當中呢?.
今天你願意打開你的心.
去接受耶穌成為你的主.
有沒有這些朋友在我們當中呢?.
好,我們祈禱.
因為神下你的靈就在我們當中.
帶領我們.
幫助我們更深的進度你的同在.
讓我們知道.
神下你的恩典也永遠來臨.
多謝你聽我們禱告.
奉耶穌基督得勝阿們.
阿們.
\newpage



\section{創世記 12:1-50:26}
\label{sec:032098}
\textbf{碎片}~~~~陳崇基牧師
\newline
\newline
連結: \href{http://www.church.com.hk/acms/content.asp?site=cdc\&op=show\&type=product\&code=032098\&layout=sermon}{\texttt{032098}} ~~~~ 語音日期: 2026-05-03
\newline
\newline
\hyperref[sec:032097]{\small{< < < PREV SERMON < < <}}
~
\hyperref[sec:index_chronic]{\small{[返順時目]}}
~
\hyperref[sec:032099]{\small{> > > NEXT SERMON > > >}}
\newline
\newline
創世記 12:1-50:26
\newline
\begin{longtable}{cl}
\hline
\hline
章節 & 經文 (和合本修訂版)\\
\hline
12:1 & \begin{tabularx}{0.7\textwidth}{X} 耶和華對亞伯蘭說：「你要離開本地、本族、父家，往我所要指示你的地去。 \end{tabularx} \\ \\ \relax
12:2 & \begin{tabularx}{0.7\textwidth}{X} 我必使你成為大國，我必賜福給你，使你的名為大；你要使別人得福。 \end{tabularx} \\ \\ \relax
12:3 & \begin{tabularx}{0.7\textwidth}{X} 為你祝福的，我必賜福給他；詛咒你的，我必詛咒他。地上的萬族都必因你得福。」 \end{tabularx} \\ \\ \relax
12:4 & \begin{tabularx}{0.7\textwidth}{X} 亞伯蘭就遵照耶和華的吩咐去了；羅得也和他同去。亞伯蘭離開哈蘭的時候年七十五歲。 \end{tabularx} \\ \\ \relax
12:5 & \begin{tabularx}{0.7\textwidth}{X} 亞伯蘭帶著他妻子撒萊和姪兒羅得，以及他們在哈蘭積蓄的財物、獲得的人口，往迦南地去。他們就來到了迦南地。 \end{tabularx} \\ \\ \relax
12:6 & \begin{tabularx}{0.7\textwidth}{X} 亞伯蘭經過那地，直到示劍地方，摩利橡樹那裡；當時迦南人住在那地。 \end{tabularx} \\ \\ \relax
12:7 & \begin{tabularx}{0.7\textwidth}{X} 耶和華向亞伯蘭顯現，說：「我要把這地賜給你的後裔。」亞伯蘭就在那裡為向他顯現的耶和華築了一座壇。 \end{tabularx} \\ \\ \relax
12:8 & \begin{tabularx}{0.7\textwidth}{X} 從那裡他又遷到伯特利東邊的山，支搭帳棚；西邊是伯特利，東邊是艾。他在那裡又為耶和華築了一座壇，求告耶和華的名。 \end{tabularx} \\ \\ \relax
12:9 & \begin{tabularx}{0.7\textwidth}{X} 後來亞伯蘭漸漸遷往尼革夫去。 \end{tabularx} \\ \\ \relax
12:10 & \begin{tabularx}{0.7\textwidth}{X} 那地遭遇饑荒。亞伯蘭因那地的饑荒嚴重，就下到埃及，要在那裡寄居。 \end{tabularx} \\ \\ \relax
12:11 & \begin{tabularx}{0.7\textwidth}{X} 將近埃及，他對妻子撒萊說：「看哪，我知道你是美貌的女人。 \end{tabularx} \\ \\ \relax
12:12 & \begin{tabularx}{0.7\textwidth}{X} 埃及人看見你會說：『這是他的妻子』，他們就會殺我，卻讓你活著。 \end{tabularx} \\ \\ \relax
12:13 & \begin{tabularx}{0.7\textwidth}{X} 所以，請你說你是我的妹妹，使我可以因你得平安，我的性命也因你存活。」 \end{tabularx} \\ \\ \relax
12:14 & \begin{tabularx}{0.7\textwidth}{X} 亞伯蘭到達埃及時，埃及人看見那女人極其美貌。 \end{tabularx} \\ \\ \relax
12:15 & \begin{tabularx}{0.7\textwidth}{X} 法老的臣僕看見了她，就在法老面前稱讚她。那女人就被帶進法老的宮中。 \end{tabularx} \\ \\ \relax
12:16 & \begin{tabularx}{0.7\textwidth}{X} 法老就因她厚待亞伯蘭，給了亞伯蘭許多牛、羊、公驢、奴僕、婢女、母驢、駱駝。 \end{tabularx} \\ \\ \relax
12:17 & \begin{tabularx}{0.7\textwidth}{X} 耶和華因亞伯蘭妻子撒萊的緣故，降大災擊打法老和他的全家。 \end{tabularx} \\ \\ \relax
12:18 & \begin{tabularx}{0.7\textwidth}{X} 法老召了亞伯蘭來，說：「你向我做的是甚麼事呢？為甚麼沒有告訴我她是你的妻子？ \end{tabularx} \\ \\ \relax
12:19 & \begin{tabularx}{0.7\textwidth}{X} 為甚麼說『她是我的妹妹』，以致我把她接來要作我的妻子呢？現在 ，看哪，你的妻子在這裡，帶她走吧！」 \end{tabularx} \\ \\ \relax
12:20 & \begin{tabularx}{0.7\textwidth}{X} 於是法老吩咐人把亞伯蘭和他妻子，以及他一切所有的都送走了。 \end{tabularx} \\ \\ \relax
13:1 & \begin{tabularx}{0.7\textwidth}{X} 亞伯蘭帶著他的妻子與羅得，以及一切所有的，從埃及上尼革夫去。 \end{tabularx} \\ \\ \relax
13:2 & \begin{tabularx}{0.7\textwidth}{X} 亞伯蘭的牲畜和金銀極多。 \end{tabularx} \\ \\ \relax
13:3 & \begin{tabularx}{0.7\textwidth}{X} 他從尼革夫漸漸往伯特利去，到了伯特利和艾的中間，當初他支搭帳棚的地方， \end{tabularx} \\ \\ \relax
13:4 & \begin{tabularx}{0.7\textwidth}{X} 也是他起先築壇的地方。亞伯蘭在那裡求告耶和華的名。 \end{tabularx} \\ \\ \relax
13:5 & \begin{tabularx}{0.7\textwidth}{X} 與亞伯蘭同行的羅得也有牛群、羊群、帳棚。 \end{tabularx} \\ \\ \relax
13:6 & \begin{tabularx}{0.7\textwidth}{X} 那地容不下他們住在一起；因為他們的財物非常多，使他們不能同住一起。 \end{tabularx} \\ \\ \relax
13:7 & \begin{tabularx}{0.7\textwidth}{X} 當時，迦南人與比利洗人在那地居住。亞伯蘭的牧人和羅得的牧人之間起了爭端。 \end{tabularx} \\ \\ \relax
13:8 & \begin{tabularx}{0.7\textwidth}{X} 亞伯蘭就對羅得說：「你我不可以相爭，你的牧人和我的牧人也不可以相爭，因為我們是一家人。 \end{tabularx} \\ \\ \relax
13:9 & \begin{tabularx}{0.7\textwidth}{X} 遍地不都在你眼前嗎？請你離開我吧！你向左，我就向右；你向右，我就向左。」 \end{tabularx} \\ \\ \relax
13:10 & \begin{tabularx}{0.7\textwidth}{X} 羅得舉目，看見約旦河整個平原，直到瑣珥，都是水源充足之地。在耶和華未毀滅所多瑪、蛾摩拉以前，那地好像耶和華的園子，又像埃及地。 \end{tabularx} \\ \\ \relax
13:11 & \begin{tabularx}{0.7\textwidth}{X} 於是羅得選擇了約旦河整個平原。羅得往東遷移，他們就彼此分開了。 \end{tabularx} \\ \\ \relax
13:12 & \begin{tabularx}{0.7\textwidth}{X} 亞伯蘭住在迦南地；羅得住在平原的城鎮，他漸漸遷移帳棚，直到所多瑪。 \end{tabularx} \\ \\ \relax
13:13 & \begin{tabularx}{0.7\textwidth}{X} 所多瑪人在耶和華面前罪大惡極。 \end{tabularx} \\ \\ \relax
13:14 & \begin{tabularx}{0.7\textwidth}{X} 羅得離開亞伯蘭以後，耶和華對亞伯蘭說：「你要從你所在的地方，舉目向東西南北觀看； \end{tabularx} \\ \\ \relax
13:15 & \begin{tabularx}{0.7\textwidth}{X} 你所看見一切的地，我都要把它賜給你和你的後裔，直到永遠。 \end{tabularx} \\ \\ \relax
13:16 & \begin{tabularx}{0.7\textwidth}{X} 我要使你的後裔好像地上的塵沙，人若能數地上的塵沙，才能數你的後裔。 \end{tabularx} \\ \\ \relax
13:17 & \begin{tabularx}{0.7\textwidth}{X} 你起來，縱橫走遍這地，因為我必把這地賜給你。」 \end{tabularx} \\ \\ \relax
13:18 & \begin{tabularx}{0.7\textwidth}{X} 亞伯蘭就遷移帳棚，來到希伯崙，幔利的橡樹那裡居住，在那裡為耶和華築了一座壇。 \end{tabularx} \\ \\ \relax
14:1 & \begin{tabularx}{0.7\textwidth}{X} 當暗拉非作示拿王，亞略作以拉撒王，基大老瑪作以攔王，提達作戈印王的時候， \end{tabularx} \\ \\ \relax
14:2 & \begin{tabularx}{0.7\textwidth}{X} 他們攻打所多瑪王比拉、蛾摩拉王比沙、押瑪王示納、洗扁王善以別和比拉王，比拉就是瑣珥。 \end{tabularx} \\ \\ \relax
14:3 & \begin{tabularx}{0.7\textwidth}{X} 這些王都會合在西訂谷，西訂谷就是鹽海。 \end{tabularx} \\ \\ \relax
14:4 & \begin{tabularx}{0.7\textwidth}{X} 他們已經服事基大老瑪十二年，第十三年就背叛了。 \end{tabularx} \\ \\ \relax
14:5 & \begin{tabularx}{0.7\textwidth}{X} 第十四年，基大老瑪和與他結盟的王都來了，在亞特律‧加寧擊敗利乏音人，在哈麥擊敗蘇西人，在沙微‧基列亭擊敗以米人， \end{tabularx} \\ \\ \relax
14:6 & \begin{tabularx}{0.7\textwidth}{X} 在何利人的西珥山擊敗何利人，一直到靠近曠野的伊勒‧巴蘭。 \end{tabularx} \\ \\ \relax
14:7 & \begin{tabularx}{0.7\textwidth}{X} 他們轉回，來到安‧密巴，就是加低斯，擊敗了亞瑪力全地的人，以及住在哈洗遜‧他瑪的亞摩利人。 \end{tabularx} \\ \\ \relax
14:8 & \begin{tabularx}{0.7\textwidth}{X} 於是所多瑪王、蛾摩拉王、押瑪王、洗扁王和比拉王，比拉就是瑣珥，都出來，在西訂谷擺陣，與他們交戰， \end{tabularx} \\ \\ \relax
14:9 & \begin{tabularx}{0.7\textwidth}{X} 就是與以攔王基大老瑪、戈印王提達、示拿王暗拉非、以拉撒王亞略交戰；這就是四王對五王之戰。 \end{tabularx} \\ \\ \relax
14:10 & \begin{tabularx}{0.7\textwidth}{X} 西訂谷有許多柏油坑。所多瑪王和蛾摩拉王逃跑，掉在坑裡，其餘的人都往山上逃跑。 \end{tabularx} \\ \\ \relax
14:11 & \begin{tabularx}{0.7\textwidth}{X} 四王就把所多瑪和蛾摩拉所有的財物和所有的糧食都擄掠去了； \end{tabularx} \\ \\ \relax
14:12 & \begin{tabularx}{0.7\textwidth}{X} 他們也把亞伯蘭的姪兒羅得和羅得的財物都擄掠去了。當時羅得住在所多瑪。 \end{tabularx} \\ \\ \relax
14:13 & \begin{tabularx}{0.7\textwidth}{X} 有一個逃脫的人來告訴希伯來人亞伯蘭；亞伯蘭正住在亞摩利人幔利的橡樹那裡。幔利、以實各和亞乃都是弟兄，曾與亞伯蘭結盟。 \end{tabularx} \\ \\ \relax
14:14 & \begin{tabularx}{0.7\textwidth}{X} 亞伯蘭聽見他姪兒被擄去，就把三百一十八個生在他家中、受過訓練的壯丁全都出動去追，一直到但。 \end{tabularx} \\ \\ \relax
14:15 & \begin{tabularx}{0.7\textwidth}{X} 在夜間，他和他的僕人分隊擊敗了敵人，並且追殺他們，直到大馬士革北邊的何把。 \end{tabularx} \\ \\ \relax
14:16 & \begin{tabularx}{0.7\textwidth}{X} 他把一切被擄掠的財物奪回，也把他姪兒羅得和他的財物，以及人和婦女都奪回來。 \end{tabularx} \\ \\ \relax
14:17 & \begin{tabularx}{0.7\textwidth}{X} 亞伯蘭擊敗基大老瑪和與他結盟的王回來的時候，所多瑪王出來，在沙微谷迎接他，沙微谷就是王的谷。 \end{tabularx} \\ \\ \relax
14:18 & \begin{tabularx}{0.7\textwidth}{X} 又有撒冷王麥基洗德帶著餅和酒出來；他是至高神的祭司。 \end{tabularx} \\ \\ \relax
14:19 & \begin{tabularx}{0.7\textwidth}{X} 他為亞伯蘭祝福，說：「願至高的神、天地的主賜福給亞伯蘭！ \end{tabularx} \\ \\ \relax
14:20 & \begin{tabularx}{0.7\textwidth}{X} 至高的神把敵人交在你手裡，他是應當稱頌的！」亞伯蘭就把所有的拿出十分之一給他。 \end{tabularx} \\ \\ \relax
14:21 & \begin{tabularx}{0.7\textwidth}{X} 所多瑪王對亞伯蘭說：「你把人還給我，財物你自己拿去吧！」 \end{tabularx} \\ \\ \relax
14:22 & \begin{tabularx}{0.7\textwidth}{X} 亞伯蘭對所多瑪王說：「我指著耶和華—至高的神、天地的主起誓： \end{tabularx} \\ \\ \relax
14:23 & \begin{tabularx}{0.7\textwidth}{X} 凡是你的東西，就是一根線、一條鞋帶，我都不拿，免得你說：『是我使亞伯蘭富足！』 \end{tabularx} \\ \\ \relax
14:24 & \begin{tabularx}{0.7\textwidth}{X} 我甚麼都不要，只是僕人所吃的，以及與我同去的亞乃、以實各、幔利所應得的份，讓他們拿去吧！」 \end{tabularx} \\ \\ \relax
15:1 & \begin{tabularx}{0.7\textwidth}{X} 這些事以後，耶和華的話在異象中臨到亞伯蘭，說：「亞伯蘭哪，不要懼怕！我是你的盾牌，你必得豐富的賞賜。」 \end{tabularx} \\ \\ \relax
15:2 & \begin{tabularx}{0.7\textwidth}{X} 亞伯蘭說：「主耶和華啊，我還沒有兒子，你能賜我甚麼呢？承受我家業的是大馬士革人以利以謝。」 \end{tabularx} \\ \\ \relax
15:3 & \begin{tabularx}{0.7\textwidth}{X} 亞伯蘭又說：「看哪，你沒有給我後嗣。你看，那生在我家中的人要繼承我。」 \end{tabularx} \\ \\ \relax
15:4 & \begin{tabularx}{0.7\textwidth}{X} 看哪，耶和華的話又臨到他，說：「這人不會繼承你，你本身所生的才會繼承你。」 \end{tabularx} \\ \\ \relax
15:5 & \begin{tabularx}{0.7\textwidth}{X} 於是耶和華帶他到外面，說：「你向天觀看，去數星星，你能數得清嗎？」又對他說：「你的後裔將要如此。」 \end{tabularx} \\ \\ \relax
15:6 & \begin{tabularx}{0.7\textwidth}{X} 亞伯蘭信耶和華，耶和華就以此算他為義。 \end{tabularx} \\ \\ \relax
15:7 & \begin{tabularx}{0.7\textwidth}{X} 耶和華又對他說：「我是耶和華，曾領你出迦勒底的吾珥，為要把這地賜你為業。」 \end{tabularx} \\ \\ \relax
15:8 & \begin{tabularx}{0.7\textwidth}{X} 亞伯蘭說：「主耶和華啊，我怎能知道我必得這地為業呢？」 \end{tabularx} \\ \\ \relax
15:9 & \begin{tabularx}{0.7\textwidth}{X} 耶和華對他說：「你為我取一頭三歲的母牛犢，一隻三歲的母山羊，一隻三歲的公綿羊，一隻斑鳩和一隻雛鴿。」 \end{tabularx} \\ \\ \relax
15:10 & \begin{tabularx}{0.7\textwidth}{X} 亞伯蘭就把這些都取來，每樣從中間劈成兩半，一半對著另一半排列，只有鳥沒有劈開。 \end{tabularx} \\ \\ \relax
15:11 & \begin{tabularx}{0.7\textwidth}{X} 當鷙鳥下來，落在這些屍體上時，亞伯蘭就把牠們趕走了。 \end{tabularx} \\ \\ \relax
15:12 & \begin{tabularx}{0.7\textwidth}{X} 日落的時候，亞伯蘭沉睡了。看哪，有大而可怕的黑暗落在他身上。 \end{tabularx} \\ \\ \relax
15:13 & \begin{tabularx}{0.7\textwidth}{X} 耶和華對亞伯蘭說：「你要確實知道，你的後裔必寄居在別人的地，服事那地的人；那地的人要虐待他們四百年。 \end{tabularx} \\ \\ \relax
15:14 & \begin{tabularx}{0.7\textwidth}{X} 但我要懲罰他們所服事的那國，以後他們必帶著許多財物從那裡出來。 \end{tabularx} \\ \\ \relax
15:15 & \begin{tabularx}{0.7\textwidth}{X} 至於你，你要平平安安歸到你祖先那裡，必享長壽，被人埋葬。 \end{tabularx} \\ \\ \relax
15:16 & \begin{tabularx}{0.7\textwidth}{X} 到了第四代，他們必回到這裡，因為亞摩利人的罪惡到現在還沒有滿盈。」 \end{tabularx} \\ \\ \relax
15:17 & \begin{tabularx}{0.7\textwidth}{X} 日落天黑的時候，看哪，有冒煙的爐和燒著的火把從那些肉塊中經過。 \end{tabularx} \\ \\ \relax
15:18 & \begin{tabularx}{0.7\textwidth}{X} 在那日，耶和華與亞伯蘭立約，說：「我已賜給你的後裔這一片地，從埃及河直到大河，幼發拉底河， \end{tabularx} \\ \\ \relax
15:19 & \begin{tabularx}{0.7\textwidth}{X} 就是基尼人、基尼洗人、甲摩尼人、 \end{tabularx} \\ \\ \relax
15:20 & \begin{tabularx}{0.7\textwidth}{X} 赫人、比利洗人、利乏音人、 \end{tabularx} \\ \\ \relax
15:21 & \begin{tabularx}{0.7\textwidth}{X} 亞摩利人、迦南人、革迦撒人、耶布斯人的地。」 \end{tabularx} \\ \\ \relax
16:1 & \begin{tabularx}{0.7\textwidth}{X} 亞伯蘭的妻子撒萊沒有為他生孩子。撒萊有一個婢女，是埃及人，名叫夏甲。 \end{tabularx} \\ \\ \relax
16:2 & \begin{tabularx}{0.7\textwidth}{X} 撒萊對亞伯蘭說：「看哪，耶和華使我不能生育。你來和我的婢女同房，也許我可以從她得孩子 。」亞伯蘭聽從了撒萊的話。 \end{tabularx} \\ \\ \relax
16:3 & \begin{tabularx}{0.7\textwidth}{X} 於是亞伯蘭的妻子撒萊把她的婢女，埃及人夏甲，給了丈夫為妾；那時亞伯蘭在迦南已經住了十年。 \end{tabularx} \\ \\ \relax
16:4 & \begin{tabularx}{0.7\textwidth}{X} 亞伯蘭與夏甲同房，夏甲就懷了孕。她看見自己有孕，就輕視她的女主人。 \end{tabularx} \\ \\ \relax
16:5 & \begin{tabularx}{0.7\textwidth}{X} 撒萊對亞伯蘭說：「我因你受了委屈。我把我的婢女放在你懷中，她見自己懷了孕，就輕視我。願耶和華在你我之間判斷。」 \end{tabularx} \\ \\ \relax
16:6 & \begin{tabularx}{0.7\textwidth}{X} 亞伯蘭對撒萊說：「看哪，婢女在你手裡，你可以照你看為好的對待她。」於是，撒萊虐待她，她就從撒萊面前逃走了。 \end{tabularx} \\ \\ \relax
16:7 & \begin{tabularx}{0.7\textwidth}{X} 耶和華的使者在曠野的水泉旁，在書珥路上的水泉旁遇見夏甲， \end{tabularx} \\ \\ \relax
16:8 & \begin{tabularx}{0.7\textwidth}{X} 對她說：「撒萊的婢女夏甲，你從哪裡來？要到哪裡去？」她說：「我從我的女主人撒萊面前逃出來。」 \end{tabularx} \\ \\ \relax
16:9 & \begin{tabularx}{0.7\textwidth}{X} 耶和華的使者對她說：「你要回到你的女主人那裡，屈服在她手下。」 \end{tabularx} \\ \\ \relax
16:10 & \begin{tabularx}{0.7\textwidth}{X} 耶和華的使者對她說：「我必使你的後裔極其繁多，多到不可勝數。」 \end{tabularx} \\ \\ \relax
16:11 & \begin{tabularx}{0.7\textwidth}{X} 耶和華的使者又對她說：「看哪，你已懷孕，要生一個兒子。你要給他起名叫以實瑪利，因為耶和華聽見了你的苦楚。 \end{tabularx} \\ \\ \relax
16:12 & \begin{tabularx}{0.7\textwidth}{X} 他為人必像野驢。他的手要攻打人，人的手也要攻打他；他必常與他的眾弟兄作對。」 \end{tabularx} \\ \\ \relax
16:13 & \begin{tabularx}{0.7\textwidth}{X} 夏甲就稱那向她說話的耶和華為「你是看見的神」，因為她說：「他看見了我之後，我還能在這裡看見他嗎？」 \end{tabularx} \\ \\ \relax
16:14 & \begin{tabularx}{0.7\textwidth}{X} 所以這井名叫庇耳‧拉海‧萊，看哪，它位於加低斯和巴列的中間。 \end{tabularx} \\ \\ \relax
16:15 & \begin{tabularx}{0.7\textwidth}{X} 後來夏甲為亞伯蘭生了一個兒子；亞伯蘭給夏甲生的兒子起名叫以實瑪利。 \end{tabularx} \\ \\ \relax
16:16 & \begin{tabularx}{0.7\textwidth}{X} 夏甲為亞伯蘭生以實瑪利的時候，亞伯蘭年八十六歲。 \end{tabularx} \\ \\ \relax
17:1 & \begin{tabularx}{0.7\textwidth}{X} 亞伯蘭九十九歲時，耶和華向他顯現，對他說：「我是全能的神。你當在我面前行走，作完全的人， \end{tabularx} \\ \\ \relax
17:2 & \begin{tabularx}{0.7\textwidth}{X} 我要與你立約，使你的後裔極其繁多。」 \end{tabularx} \\ \\ \relax
17:3 & \begin{tabularx}{0.7\textwidth}{X} 亞伯蘭臉伏於地；神又對他說： \end{tabularx} \\ \\ \relax
17:4 & \begin{tabularx}{0.7\textwidth}{X} 「看哪，這就是我與你立的約，你要成為多國的父。 \end{tabularx} \\ \\ \relax
17:5 & \begin{tabularx}{0.7\textwidth}{X} 從今以後，你的名字不再叫亞伯蘭，要叫亞伯拉罕，因為我已經立你作多國之父。 \end{tabularx} \\ \\ \relax
17:6 & \begin{tabularx}{0.7\textwidth}{X} 我必使你生養極其繁多；國度要從你而立，君王要從你而出。 \end{tabularx} \\ \\ \relax
17:7 & \begin{tabularx}{0.7\textwidth}{X} 我要與你，以及你世世代代的後裔堅立我的約，成為永遠的約，是要作你和你後裔的神。 \end{tabularx} \\ \\ \relax
17:8 & \begin{tabularx}{0.7\textwidth}{X} 我要把你現在寄居的地，就是迦南全地，賜給你和你的後裔永遠為業；我也必作他們的神。」 \end{tabularx} \\ \\ \relax
17:9 & \begin{tabularx}{0.7\textwidth}{X} 神又對亞伯拉罕說：「你和你的後裔一定要世世代代遵守我的約。 \end{tabularx} \\ \\ \relax
17:10 & \begin{tabularx}{0.7\textwidth}{X} 這就是我與你，以及你的後裔所立的約，是你們所當遵守的，你們所有的男子都要受割禮。 \end{tabularx} \\ \\ \relax
17:11 & \begin{tabularx}{0.7\textwidth}{X} 你們要割去肉體的包皮，這是我與你們立約的記號。 \end{tabularx} \\ \\ \relax
17:12 & \begin{tabularx}{0.7\textwidth}{X} 你們世世代代的男子，無論是在家裡生的，或是用銀子從外人買來而不是你後裔生的，都要在生下來的第八日受割禮。 \end{tabularx} \\ \\ \relax
17:13 & \begin{tabularx}{0.7\textwidth}{X} 你家裡生的和你用銀子買的，都必須受割禮。這樣，我的約就在你們肉體上成為永遠的約。 \end{tabularx} \\ \\ \relax
17:14 & \begin{tabularx}{0.7\textwidth}{X} 不受割禮的男子都必從民中剪除，因他違背了我的約。」 \end{tabularx} \\ \\ \relax
17:15 & \begin{tabularx}{0.7\textwidth}{X} 神又對亞伯拉罕說：「至於你的妻子撒萊，不可再叫她撒萊，她的名要叫撒拉。 \end{tabularx} \\ \\ \relax
17:16 & \begin{tabularx}{0.7\textwidth}{X} 我必賜福給她，也要從她賜一個兒子給你。我必賜福給撒拉，她要興起多國；必有百姓的君王從她而出。」 \end{tabularx} \\ \\ \relax
17:17 & \begin{tabularx}{0.7\textwidth}{X} 亞伯拉罕就臉伏於地竊笑，心裡想：「一百歲的人還能有孩子嗎？撒拉已經九十歲了，還能生育嗎？」 \end{tabularx} \\ \\ \relax
17:18 & \begin{tabularx}{0.7\textwidth}{X} 亞伯拉罕對神說：「但願以實瑪利活在你面前。」 \end{tabularx} \\ \\ \relax
17:19 & \begin{tabularx}{0.7\textwidth}{X} 神說：「不！你妻子撒拉必為你生一個兒子，你要給他起名叫以撒。我要與他堅立我的約，成為他後裔永遠的約。 \end{tabularx} \\ \\ \relax
17:20 & \begin{tabularx}{0.7\textwidth}{X} 至於以實瑪利，我已聽見你了：看哪，我必賜福給他，使他興旺，極其繁多。他必生十二個族長，我要使他成為大國。 \end{tabularx} \\ \\ \relax
17:21 & \begin{tabularx}{0.7\textwidth}{X} 到明年所定的時候，撒拉必為你生以撒，我要與他堅立我的約。」 \end{tabularx} \\ \\ \relax
17:22 & \begin{tabularx}{0.7\textwidth}{X} 神和亞伯拉罕說完了話，就離開他上升去了。 \end{tabularx} \\ \\ \relax
17:23 & \begin{tabularx}{0.7\textwidth}{X} 在那一天，亞伯拉罕遵照神所說的，給他的兒子以實瑪利和家裡所有的男丁，無論是在家裡生的，或是用銀子買來的，都行了割禮。 \end{tabularx} \\ \\ \relax
17:24 & \begin{tabularx}{0.7\textwidth}{X} 亞伯拉罕受割禮時，年九十九歲。 \end{tabularx} \\ \\ \relax
17:25 & \begin{tabularx}{0.7\textwidth}{X} 他兒子以實瑪利受割禮時，年十三歲。 \end{tabularx} \\ \\ \relax
17:26 & \begin{tabularx}{0.7\textwidth}{X} 在那一天，亞伯拉罕和他兒子以實瑪利一同受了割禮。 \end{tabularx} \\ \\ \relax
17:27 & \begin{tabularx}{0.7\textwidth}{X} 家裡所有的男人，無論是在家裡生的，或是用銀子從外人買來的，也都一同受了割禮。 \end{tabularx} \\ \\ \relax
18:1 & \begin{tabularx}{0.7\textwidth}{X} 耶和華在幔利橡樹那裡向亞伯拉罕顯現。天正熱的時候，亞伯拉罕坐在帳棚門口。 \end{tabularx} \\ \\ \relax
18:2 & \begin{tabularx}{0.7\textwidth}{X} 他舉目觀看，看哪，有三個人站在他附近。他一看見，就從帳棚門口跑去迎接他們，俯伏在地， \end{tabularx} \\ \\ \relax
18:3 & \begin{tabularx}{0.7\textwidth}{X} 說：「我主，我若在你眼前蒙恩，請不要離開你的僕人走過去。 \end{tabularx} \\ \\ \relax
18:4 & \begin{tabularx}{0.7\textwidth}{X} 容我拿點水來，請你們洗腳，在樹下休息。 \end{tabularx} \\ \\ \relax
18:5 & \begin{tabularx}{0.7\textwidth}{X} 既然你們來到僕人這裡了，我再拿點餅來，讓你們恢復心力，然後再走。」他們說：「就照你說的去做吧。」 \end{tabularx} \\ \\ \relax
18:6 & \begin{tabularx}{0.7\textwidth}{X} 亞伯拉罕急忙進帳棚到撒拉那裡，說：「你趕快拿三細亞細麵，揉麵做餅。」 \end{tabularx} \\ \\ \relax
18:7 & \begin{tabularx}{0.7\textwidth}{X} 亞伯拉罕又跑到牛群裡，牽了一頭又嫩又好的牛犢來，交給僕人，僕人就急忙去預備。 \end{tabularx} \\ \\ \relax
18:8 & \begin{tabularx}{0.7\textwidth}{X} 亞伯拉罕取了乳酪和奶，以及預備好了的牛犢來，擺在他們面前，自己在樹下站在旁邊，他們就吃了。 \end{tabularx} \\ \\ \relax
18:9 & \begin{tabularx}{0.7\textwidth}{X} 他們對亞伯拉罕說：「你妻子撒拉在哪裡？」他說：「看哪，在帳棚裡。」 \end{tabularx} \\ \\ \relax
18:10 & \begin{tabularx}{0.7\textwidth}{X} 有一位說：「明年這時候 ，我一定會回到你這裡。看哪，你的妻子撒拉會生一個兒子。」撒拉在那人後面的帳棚門口也聽見了。 \end{tabularx} \\ \\ \relax
18:11 & \begin{tabularx}{0.7\textwidth}{X} 亞伯拉罕和撒拉都年紀老邁，撒拉的月經已停了。 \end{tabularx} \\ \\ \relax
18:12 & \begin{tabularx}{0.7\textwidth}{X} 撒拉心裡竊笑，說：「我已衰老，我的主也老了，怎能有這喜事呢？」 \end{tabularx} \\ \\ \relax
18:13 & \begin{tabularx}{0.7\textwidth}{X} 耶和華對亞伯拉罕說：「撒拉為甚麼竊笑，說：『我已年老，果真能生育嗎？』 \end{tabularx} \\ \\ \relax
18:14 & \begin{tabularx}{0.7\textwidth}{X} 耶和華豈有難成的事嗎？到了所定的時候，我必回到你這裡。明年這時候，撒拉會生一個兒子。」 \end{tabularx} \\ \\ \relax
18:15 & \begin{tabularx}{0.7\textwidth}{X} 撒拉因為害怕，就不承認，說：「我沒有笑。」那人說：「不，你的確笑了。」 \end{tabularx} \\ \\ \relax
18:16 & \begin{tabularx}{0.7\textwidth}{X} 三人從那裡起程，面向所多瑪觀望，亞伯拉罕與他們同行，要送他們一程。 \end{tabularx} \\ \\ \relax
18:17 & \begin{tabularx}{0.7\textwidth}{X} 耶和華說：「我所要做的事豈可瞞著亞伯拉罕呢？ \end{tabularx} \\ \\ \relax
18:18 & \begin{tabularx}{0.7\textwidth}{X} 亞伯拉罕必要成為強大的國；地上的萬國都必因他得福。 \end{tabularx} \\ \\ \relax
18:19 & \begin{tabularx}{0.7\textwidth}{X} 我揀選他，為要叫他命令他的子孫和後代家屬遵行耶和華的道，秉公行義，使耶和華所應許亞伯拉罕的話都實現了。」 \end{tabularx} \\ \\ \relax
18:20 & \begin{tabularx}{0.7\textwidth}{X} 耶和華說：「所多瑪和蛾摩拉罪惡極其嚴重，控告他們的聲音很大。 \end{tabularx} \\ \\ \relax
18:21 & \begin{tabularx}{0.7\textwidth}{X} 我要下去察看他們所做的，是否真的像那達到我這裡的聲音一樣；如果不是，我也要知道。」 \end{tabularx} \\ \\ \relax
18:22 & \begin{tabularx}{0.7\textwidth}{X} 二人轉身離開那裡，往所多瑪去；但亞伯拉罕仍然站在耶和華面前。 \end{tabularx} \\ \\ \relax
18:23 & \begin{tabularx}{0.7\textwidth}{X} 亞伯拉罕近前來，說：「你真的要把義人和惡人一同剿滅嗎？ \end{tabularx} \\ \\ \relax
18:24 & \begin{tabularx}{0.7\textwidth}{X} 假若那城裡有五十個義人，你真的還要剿滅，不因城裡這五十個義人饒了那地方嗎？ \end{tabularx} \\ \\ \relax
18:25 & \begin{tabularx}{0.7\textwidth}{X} 你絕不會做這樣的事，把義人與惡人一同殺了，使義人與惡人一樣。你絕不會這樣！審判全地的主豈不做公平的事嗎？」 \end{tabularx} \\ \\ \relax
18:26 & \begin{tabularx}{0.7\textwidth}{X} 耶和華說：「我若在所多瑪城裡找到五十個義人，我就為他們的緣故饒恕那整個地方。」 \end{tabularx} \\ \\ \relax
18:27 & \begin{tabularx}{0.7\textwidth}{X} 亞伯拉罕回答說：「看哪，我雖只是塵土灰燼，還敢向主說話。 \end{tabularx} \\ \\ \relax
18:28 & \begin{tabularx}{0.7\textwidth}{X} 假若這五十個義人少了五個，你就因為少了五個而毀滅全城嗎？」他說：「我在那裡若找到四十五個，就不毀滅。」 \end{tabularx} \\ \\ \relax
18:29 & \begin{tabularx}{0.7\textwidth}{X} 亞伯拉罕又對他說：「假若在那裡找到四十個呢？」他說：「為這四十個的緣故，我也不做。」 \end{tabularx} \\ \\ \relax
18:30 & \begin{tabularx}{0.7\textwidth}{X} 亞伯拉罕說：「求主不要生氣，容我說，假若在那裡找到三十個呢？」他說：「我在那裡若找到三十個，我也不做。」 \end{tabularx} \\ \\ \relax
18:31 & \begin{tabularx}{0.7\textwidth}{X} 亞伯拉罕說：「看哪，我還敢向主說，假若在那裡找到二十個呢？」他說：「為這二十個的緣故，我也不毀滅。」 \end{tabularx} \\ \\ \relax
18:32 & \begin{tabularx}{0.7\textwidth}{X} 亞伯拉罕說：「求主不要生氣，我再說一次，假若在那裡找到十個呢？」他說：「為這十個的緣故，我也不毀滅。」 \end{tabularx} \\ \\ \relax
18:33 & \begin{tabularx}{0.7\textwidth}{X} 耶和華與亞伯拉罕說完了話就走了；亞伯拉罕也回到自己的地方去了。 \end{tabularx} \\ \\ \relax
19:1 & \begin{tabularx}{0.7\textwidth}{X} 兩個天使在傍晚到了所多瑪，羅得正坐在所多瑪的城門口。羅得一看見，就起身迎接他們，臉伏於地下拜， \end{tabularx} \\ \\ \relax
19:2 & \begin{tabularx}{0.7\textwidth}{X} 說：「看哪，我主，請你們轉到僕人家裡過夜，洗你們的腳，清早起來再上路。」他們說：「不！我們要在廣場上過夜。」 \end{tabularx} \\ \\ \relax
19:3 & \begin{tabularx}{0.7\textwidth}{X} 羅得懇切地請他們，他們就轉向他，進到他屋裡。羅得為他們預備宴席，烤無酵餅，他們就吃了。 \end{tabularx} \\ \\ \relax
19:4 & \begin{tabularx}{0.7\textwidth}{X} 他們還沒有躺下，所多瑪城的人，連老帶少所有的人，個個都來圍住那屋子。 \end{tabularx} \\ \\ \relax
19:5 & \begin{tabularx}{0.7\textwidth}{X} 他們呼叫羅得，對他說：「今天晚上到你這裡來的人在哪裡？把他們帶出來，讓我們親近他們。」 \end{tabularx} \\ \\ \relax
19:6 & \begin{tabularx}{0.7\textwidth}{X} 羅得出了門，把身後的門關上，到眾人那裡， \end{tabularx} \\ \\ \relax
19:7 & \begin{tabularx}{0.7\textwidth}{X} 說：「我的弟兄們，請你們不要做這惡事。 \end{tabularx} \\ \\ \relax
19:8 & \begin{tabularx}{0.7\textwidth}{X} 看哪，我有兩個女兒，還沒有親近過男人，讓我領她們出來給你們，就照你們看為好的對待她們吧！只是這兩個人既然到我舍下，請不要向他們做這事。」 \end{tabularx} \\ \\ \relax
19:9 & \begin{tabularx}{0.7\textwidth}{X} 眾人說：「站到一邊去吧！」又說：「這個人來寄居，還想扮審判官呢！現在我們要害你比害他們更厲害。」眾人就往前衝向羅得，要攻破大門。 \end{tabularx} \\ \\ \relax
19:10 & \begin{tabularx}{0.7\textwidth}{X} 那兩個人伸出手來，把羅得拉進屋子他們那裡，就關上門。 \end{tabularx} \\ \\ \relax
19:11 & \begin{tabularx}{0.7\textwidth}{X} 他們擊打門外的人，無論老少，都眼睛迷糊，找門找得很煩躁。 \end{tabularx} \\ \\ \relax
19:12 & \begin{tabularx}{0.7\textwidth}{X} 那兩個人對羅得說：「你這裡還有甚麼人嗎？無論是女婿，是兒女，這城中所有屬你的人，你都要把他們從這地方帶出去。 \end{tabularx} \\ \\ \relax
19:13 & \begin{tabularx}{0.7\textwidth}{X} 我們要毀滅這地方，因為控告城內百姓的聲音在耶和華面前非常大，耶和華派我們來毀滅這城。」 \end{tabularx} \\ \\ \relax
19:14 & \begin{tabularx}{0.7\textwidth}{X} 羅得出去，告訴娶了他女兒的女婿們說：「起來，離開這地方，因為耶和華要毀滅這城。」他的女婿們卻以為他說的是笑話。 \end{tabularx} \\ \\ \relax
19:15 & \begin{tabularx}{0.7\textwidth}{X} 天亮了，天使催逼羅得說：「起來！帶著你的妻子和你這裡的兩個女兒出去，免得你因這城的罪孽同被剿滅。」 \end{tabularx} \\ \\ \relax
19:16 & \begin{tabularx}{0.7\textwidth}{X} 但羅得遲延不走。二人因為耶和華憐憫羅得，就拉著他的手和他妻子的手，以及他兩個女兒的手，把他們領出來，安置在城外； \end{tabularx} \\ \\ \relax
19:17 & \begin{tabularx}{0.7\textwidth}{X} 領他們出來以後，就說：「逃命吧！不可回頭看，也不可在平原站住。要往山上逃跑，免得你被剿滅。」 \end{tabularx} \\ \\ \relax
19:18 & \begin{tabularx}{0.7\textwidth}{X} 羅得對他們說：「我主啊，不要這樣！ \end{tabularx} \\ \\ \relax
19:19 & \begin{tabularx}{0.7\textwidth}{X} 看哪，你僕人已經在你眼前蒙恩，你又向我大施慈愛，救我的性命。但是我不能逃到山上去，恐怕這災禍追上我，我就死了。 \end{tabularx} \\ \\ \relax
19:20 & \begin{tabularx}{0.7\textwidth}{X} 看哪，這城又近又小，比較容易逃到那裡。這不是一座小城嗎？求你容我逃到那裡，使我的性命可以存活。」 \end{tabularx} \\ \\ \relax
19:21 & \begin{tabularx}{0.7\textwidth}{X} 天使對他說：「看哪，這事我也應允你，不傾覆你所說的這城。 \end{tabularx} \\ \\ \relax
19:22 & \begin{tabularx}{0.7\textwidth}{X} 你要趕快逃到那城，因為你還沒有到那裡，我不能做甚麼。」因此那城名叫瑣珥 。 \end{tabularx} \\ \\ \relax
19:23 & \begin{tabularx}{0.7\textwidth}{X} 羅得到了瑣珥，太陽已經升出地面。 \end{tabularx} \\ \\ \relax
19:24 & \begin{tabularx}{0.7\textwidth}{X} 當時，耶和華把硫磺與火，從天上耶和華那裡降與所多瑪和蛾摩拉， \end{tabularx} \\ \\ \relax
19:25 & \begin{tabularx}{0.7\textwidth}{X} 把那些城和全平原，城裡所有的居民和土地上生長的，都毀滅了。 \end{tabularx} \\ \\ \relax
19:26 & \begin{tabularx}{0.7\textwidth}{X} 羅得的妻子在他後邊回頭一看，就變成了一根鹽柱。 \end{tabularx} \\ \\ \relax
19:27 & \begin{tabularx}{0.7\textwidth}{X} 亞伯拉罕清早起來，到了他先前站在耶和華面前的地方， \end{tabularx} \\ \\ \relax
19:28 & \begin{tabularx}{0.7\textwidth}{X} 面向所多瑪和蛾摩拉，以及平原全地觀望。他觀看，看哪，那地有濃煙上騰，好像燒窯的濃煙。 \end{tabularx} \\ \\ \relax
19:29 & \begin{tabularx}{0.7\textwidth}{X} 當神毀滅平原諸城的時候，他記念亞伯拉罕；在傾覆羅得所住之城的時候，就把羅得從傾覆中帶出來。 \end{tabularx} \\ \\ \relax
19:30 & \begin{tabularx}{0.7\textwidth}{X} 羅得因為怕住在瑣珥，就同他兩個女兒從瑣珥上去，住在山上。他和兩個女兒住在一個洞裡。 \end{tabularx} \\ \\ \relax
19:31 & \begin{tabularx}{0.7\textwidth}{X} 大女兒對小女兒說：「我們的父親老了，這地又沒有男人可以照世上的禮俗來與我們結合。 \end{tabularx} \\ \\ \relax
19:32 & \begin{tabularx}{0.7\textwidth}{X} 來！我們叫父親喝酒，然後與他同寢。這樣，我們可以從我們的父親存留後裔。」 \end{tabularx} \\ \\ \relax
19:33 & \begin{tabularx}{0.7\textwidth}{X} 於是，那晚她們叫父親喝酒，大女兒就進去和她父親同寢；她幾時躺下，幾時起來，父親都不知道。 \end{tabularx} \\ \\ \relax
19:34 & \begin{tabularx}{0.7\textwidth}{X} 第二天，大女兒對小女兒說：「看哪，我昨夜與父親同寢。今晚我們再叫他喝酒，你進去與他同寢。這樣，我們可以從父親存留後裔。」 \end{tabularx} \\ \\ \relax
19:35 & \begin{tabularx}{0.7\textwidth}{X} 於是，那晚她們又叫父親喝酒，小女兒起來與她父親同寢；她幾時躺下，幾時起來，父親都不知道。 \end{tabularx} \\ \\ \relax
19:36 & \begin{tabularx}{0.7\textwidth}{X} 這樣，羅得的兩個女兒都從她們的父親懷了孕。 \end{tabularx} \\ \\ \relax
19:37 & \begin{tabularx}{0.7\textwidth}{X} 大女兒生了兒子，給他起名叫摩押，就是現今摩押人的始祖。 \end{tabularx} \\ \\ \relax
19:38 & \begin{tabularx}{0.7\textwidth}{X} 小女兒也生了兒子，給他起名叫便‧亞米，就是現今亞捫人的始祖。 \end{tabularx} \\ \\ \relax
20:1 & \begin{tabularx}{0.7\textwidth}{X} 亞伯拉罕從那裡往尼革夫遷移，寄居在加低斯和書珥之間的基拉耳。 \end{tabularx} \\ \\ \relax
20:2 & \begin{tabularx}{0.7\textwidth}{X} 亞伯拉罕稱他的妻子撒拉為妹妹。基拉耳王亞比米勒派人把撒拉帶走。 \end{tabularx} \\ \\ \relax
20:3 & \begin{tabularx}{0.7\textwidth}{X} 夜間，神在夢中來到亞比米勒那裡，對他說：「看哪，你要死了，因為你帶來的女人，她是有丈夫的女子！」 \end{tabularx} \\ \\ \relax
20:4 & \begin{tabularx}{0.7\textwidth}{X} 亞比米勒還未親近撒拉；他說：「主啊，連公義的國，你也要毀滅嗎？ \end{tabularx} \\ \\ \relax
20:5 & \begin{tabularx}{0.7\textwidth}{X} 那人豈不是自己對我說『她是我妹妹』嗎？連這女人自己也說：『他是我哥哥。』我做這事是心正手潔的。」 \end{tabularx} \\ \\ \relax
20:6 & \begin{tabularx}{0.7\textwidth}{X} 神在夢中對他說：「我也知道你做這事是心中正直的；是我攔阻了你，免得你得罪我。所以我不讓你侵犯她。 \end{tabularx} \\ \\ \relax
20:7 & \begin{tabularx}{0.7\textwidth}{X} 現在你當把這人的妻子歸還給他；因為他是先知，他要為你禱告，使你存活。你若不歸還，你當知道，你和你所有的人都必定死。」 \end{tabularx} \\ \\ \relax
20:8 & \begin{tabularx}{0.7\textwidth}{X} 亞比米勒清早起來，叫了他的眾臣僕來，把這一切事說給他們聽，他們就很害怕。 \end{tabularx} \\ \\ \relax
20:9 & \begin{tabularx}{0.7\textwidth}{X} 亞比米勒召了亞伯拉罕來，對他說：「你怎麼向我這樣做呢？我甚麼事得罪你，你竟使我和我的國陷在大罪中呢？你對我做了不該做的事了！」 \end{tabularx} \\ \\ \relax
20:10 & \begin{tabularx}{0.7\textwidth}{X} 亞比米勒對亞伯拉罕說：「你看見甚麼才做這事呢？」 \end{tabularx} \\ \\ \relax
20:11 & \begin{tabularx}{0.7\textwidth}{X} 亞伯拉罕說：「我以為這地方的人根本不敬畏神，必為我妻子的緣故殺我。 \end{tabularx} \\ \\ \relax
20:12 & \begin{tabularx}{0.7\textwidth}{X} 況且她也真是我的妹妹；她與我是同父異母的，後來作了我的妻子。 \end{tabularx} \\ \\ \relax
20:13 & \begin{tabularx}{0.7\textwidth}{X} 當神叫我離開父家、飄流在外的時候，我對她說：我們無論走到甚麼地方，你要對人說：『他是我哥哥』，這就是你以慈愛待我了。」 \end{tabularx} \\ \\ \relax
20:14 & \begin{tabularx}{0.7\textwidth}{X} 亞比米勒把牛、羊、奴僕、婢女送給亞伯拉罕，也把他的妻子撒拉歸還給他。 \end{tabularx} \\ \\ \relax
20:15 & \begin{tabularx}{0.7\textwidth}{X} 亞比米勒說：「看哪，我的地都在你面前，你看為好的地方就居住吧。」 \end{tabularx} \\ \\ \relax
20:16 & \begin{tabularx}{0.7\textwidth}{X} 他對撒拉說：「看哪，我給你哥哥一千銀子。看哪，這要在你全家人面前遮羞，向眾人證實你是清白的。」 \end{tabularx} \\ \\ \relax
20:17 & \begin{tabularx}{0.7\textwidth}{X} 亞伯拉罕向神禱告，神就醫好亞比米勒和他的妻子，以及他的使女們，他們就能生育。 \end{tabularx} \\ \\ \relax
20:18 & \begin{tabularx}{0.7\textwidth}{X} 因耶和華為亞伯拉罕的妻子撒拉的緣故，已經使亞比米勒家中的婦人不能懷孕。 \end{tabularx} \\ \\ \relax
21:1 & \begin{tabularx}{0.7\textwidth}{X} 耶和華照著他所說的眷顧撒拉，耶和華實現了他對撒拉的應許。 \end{tabularx} \\ \\ \relax
21:2 & \begin{tabularx}{0.7\textwidth}{X} 亞伯拉罕年老，到神對他說的那所定的時候，撒拉懷了孕，給他生了一個兒子。 \end{tabularx} \\ \\ \relax
21:3 & \begin{tabularx}{0.7\textwidth}{X} 亞伯拉罕給撒拉所生的兒子起名叫以撒。 \end{tabularx} \\ \\ \relax
21:4 & \begin{tabularx}{0.7\textwidth}{X} 以撒出生後第八日，亞伯拉罕遵照神所吩咐的，為以撒行割禮。 \end{tabularx} \\ \\ \relax
21:5 & \begin{tabularx}{0.7\textwidth}{X} 他兒子以撒出生的時候，亞伯拉罕年一百歲。 \end{tabularx} \\ \\ \relax
21:6 & \begin{tabularx}{0.7\textwidth}{X} 撒拉說：「神使我歡笑，凡聽見的人必與我一同歡笑」， \end{tabularx} \\ \\ \relax
21:7 & \begin{tabularx}{0.7\textwidth}{X} 又說：「誰能預先對亞伯拉罕說，撒拉要乳養孩子呢？因為在他年老的時候，我為他生了一個兒子。」 \end{tabularx} \\ \\ \relax
21:8 & \begin{tabularx}{0.7\textwidth}{X} 孩子漸漸長大，就斷了奶。以撒斷奶的那一天，亞伯拉罕擺設豐盛的宴席。 \end{tabularx} \\ \\ \relax
21:9 & \begin{tabularx}{0.7\textwidth}{X} 那時，撒拉看見埃及人夏甲為亞伯拉罕所生的兒子戲笑， \end{tabularx} \\ \\ \relax
21:10 & \begin{tabularx}{0.7\textwidth}{X} 就對亞伯拉罕說：「你把這使女和她兒子趕出去！因為這使女的兒子不可與我的兒子以撒一同承受產業。」 \end{tabularx} \\ \\ \relax
21:11 & \begin{tabularx}{0.7\textwidth}{X} 亞伯拉罕為這事非常憂愁，因為關乎他的兒子。 \end{tabularx} \\ \\ \relax
21:12 & \begin{tabularx}{0.7\textwidth}{X} 神對亞伯拉罕說：「你不必為這孩子和你的使女憂愁。撒拉對你說的話，你都要聽從；因為從以撒生的，才要稱為你的後裔。 \end{tabularx} \\ \\ \relax
21:13 & \begin{tabularx}{0.7\textwidth}{X} 至於使女的兒子，我也必使他成為一國，因為他是你的後裔。」 \end{tabularx} \\ \\ \relax
21:14 & \begin{tabularx}{0.7\textwidth}{X} 亞伯拉罕清早起來，拿餅和一皮袋水，給了夏甲，搭在她肩上，把她和孩子一起送走。夏甲就走了，但她卻在別是巴的曠野流浪。 \end{tabularx} \\ \\ \relax
21:15 & \begin{tabularx}{0.7\textwidth}{X} 皮袋的水用完了，夏甲就把孩子放在一棵小樹下， \end{tabularx} \\ \\ \relax
21:16 & \begin{tabularx}{0.7\textwidth}{X} 自己走開約有一箭之遠，相對而坐，說：「我不忍心看見孩子死」。她就坐在對面，放聲大哭。 \end{tabularx} \\ \\ \relax
21:17 & \begin{tabularx}{0.7\textwidth}{X} 神聽見孩子的聲音，神的使者就從天上呼叫夏甲說：「夏甲，你為何這樣呢？不要害怕，神已經聽見孩子在那裡的聲音了。 \end{tabularx} \\ \\ \relax
21:18 & \begin{tabularx}{0.7\textwidth}{X} 起來！把孩子扶起來，用你的手握住他，因我必使他成為大國。」 \end{tabularx} \\ \\ \relax
21:19 & \begin{tabularx}{0.7\textwidth}{X} 神開了夏甲的眼睛，她就看見一口水井。她就去，把皮袋裝滿了水，給孩子喝。 \end{tabularx} \\ \\ \relax
21:20 & \begin{tabularx}{0.7\textwidth}{X} 神與這孩子同在，他就漸漸長大，住在曠野，成了一個弓箭手。 \end{tabularx} \\ \\ \relax
21:21 & \begin{tabularx}{0.7\textwidth}{X} 他住在巴蘭的曠野；他母親從埃及地為他娶了一個妻子。 \end{tabularx} \\ \\ \relax
21:22 & \begin{tabularx}{0.7\textwidth}{X} 那時候，亞比米勒和他的將軍非各對亞伯拉罕說：「凡你所做的事，神都與你同在。 \end{tabularx} \\ \\ \relax
21:23 & \begin{tabularx}{0.7\textwidth}{X} 我願你如今在這裡指著神對我起誓，不要虧待我和我的兒子，以及我的子孫。我怎樣忠誠待你，你也要照樣忠誠待我和你所寄居的這地。」 \end{tabularx} \\ \\ \relax
21:24 & \begin{tabularx}{0.7\textwidth}{X} 亞伯拉罕說：「我願意起誓。」 \end{tabularx} \\ \\ \relax
21:25 & \begin{tabularx}{0.7\textwidth}{X} 先前，亞比米勒的僕人霸佔了一口水井，亞伯拉罕為這事責備亞比米勒。 \end{tabularx} \\ \\ \relax
21:26 & \begin{tabularx}{0.7\textwidth}{X} 亞比米勒說：「我不知道誰做了這事，你也沒有告訴我，我到今日才聽到。」 \end{tabularx} \\ \\ \relax
21:27 & \begin{tabularx}{0.7\textwidth}{X} 亞伯拉罕把羊和牛給了亞比米勒，二人就彼此立約。 \end{tabularx} \\ \\ \relax
21:28 & \begin{tabularx}{0.7\textwidth}{X} 亞伯拉罕把七隻小母羊另放在一處。 \end{tabularx} \\ \\ \relax
21:29 & \begin{tabularx}{0.7\textwidth}{X} 亞比米勒對亞伯拉罕說：「你把這七隻小母羊另放一處是甚麼意思呢？」 \end{tabularx} \\ \\ \relax
21:30 & \begin{tabularx}{0.7\textwidth}{X} 他說：「你要從我手裡接受這七隻小母羊，作我挖了這口井的證據。」 \end{tabularx} \\ \\ \relax
21:31 & \begin{tabularx}{0.7\textwidth}{X} 所以他給那地方起名叫別是巴，因為他們二人在那裡起了誓。 \end{tabularx} \\ \\ \relax
21:32 & \begin{tabularx}{0.7\textwidth}{X} 他們在別是巴立了約，亞比米勒就和他的將軍非各起身回非利士人的地去了。 \end{tabularx} \\ \\ \relax
21:33 & \begin{tabularx}{0.7\textwidth}{X} 亞伯拉罕就在別是巴種了一棵柳樹，在那裡求告耶和華—永恆神的名。 \end{tabularx} \\ \\ \relax
21:34 & \begin{tabularx}{0.7\textwidth}{X} 亞伯拉罕在非利士人的地寄居了許多日子。 \end{tabularx} \\ \\ \relax
22:1 & \begin{tabularx}{0.7\textwidth}{X} 這些事以後，神考驗亞伯拉罕，對他說：「亞伯拉罕！」他說：「我在這裡。」 \end{tabularx} \\ \\ \relax
22:2 & \begin{tabularx}{0.7\textwidth}{X} 神說：「你要帶你的兒子，就是你所愛的獨子以撒，往摩利亞地去，在我指示你的一座山上，把他獻為燔祭。」 \end{tabularx} \\ \\ \relax
22:3 & \begin{tabularx}{0.7\textwidth}{X} 亞伯拉罕清早起來，預備了驢，帶著跟他一起的兩個僕人和他兒子以撒，劈好了燔祭的柴，就起身往神指示他的地方去了。 \end{tabularx} \\ \\ \relax
22:4 & \begin{tabularx}{0.7\textwidth}{X} 到了第三日，亞伯拉罕舉目遙望那地方。 \end{tabularx} \\ \\ \relax
22:5 & \begin{tabularx}{0.7\textwidth}{X} 亞伯拉罕對他的僕人說：「你們和驢留在這裡，我和孩子要去那裡敬拜，然後回到你們這裡來。」 \end{tabularx} \\ \\ \relax
22:6 & \begin{tabularx}{0.7\textwidth}{X} 亞伯拉罕把燔祭的柴放在他兒子以撒身上，自己手裡拿著火與刀；於是二人同行。 \end{tabularx} \\ \\ \relax
22:7 & \begin{tabularx}{0.7\textwidth}{X} 以撒對他父親亞伯拉罕說：「我父啊！」亞伯拉罕說：「我兒，我在這裡。」以撒說：「看哪，火與柴都有了，但燔祭的羔羊在哪裡呢？」 \end{tabularx} \\ \\ \relax
22:8 & \begin{tabularx}{0.7\textwidth}{X} 亞伯拉罕說：「我兒，神必自己預備燔祭的羔羊。」於是二人同行。 \end{tabularx} \\ \\ \relax
22:9 & \begin{tabularx}{0.7\textwidth}{X} 他們到了神指示他的地方，亞伯拉罕在那裡築壇，把柴擺好，綁了他兒子以撒，放在壇的柴上。 \end{tabularx} \\ \\ \relax
22:10 & \begin{tabularx}{0.7\textwidth}{X} 亞伯拉罕就伸手拿刀，要殺他的兒子。 \end{tabularx} \\ \\ \relax
22:11 & \begin{tabularx}{0.7\textwidth}{X} 耶和華的使者從天上呼喚他說：「亞伯拉罕！亞伯拉罕！」他說：「我在這裡。」 \end{tabularx} \\ \\ \relax
22:12 & \begin{tabularx}{0.7\textwidth}{X} 天使說：「不可在這孩子身上下手！一點也不可傷害他！現在我知道你是敬畏神的人了，因為你沒有把你的兒子，就是你的獨子，留下不給我。」 \end{tabularx} \\ \\ \relax
22:13 & \begin{tabularx}{0.7\textwidth}{X} 亞伯拉罕舉目觀看，看哪，一隻公綿羊兩角纏在灌木叢中。亞伯拉罕就去牽了那隻公綿羊，獻為燔祭，代替他的兒子。 \end{tabularx} \\ \\ \relax
22:14 & \begin{tabularx}{0.7\textwidth}{X} 亞伯拉罕給那地方起名叫「耶和華以勒」 。直到今日人還說：「在耶和華的山上必有預備。」 \end{tabularx} \\ \\ \relax
22:15 & \begin{tabularx}{0.7\textwidth}{X} 耶和華的使者第二次從天上呼喚亞伯拉罕， \end{tabularx} \\ \\ \relax
22:16 & \begin{tabularx}{0.7\textwidth}{X} 說：「耶和華說：『你既行了這事，沒有留下你的兒子，就是你的獨子，我指著自己起誓： \end{tabularx} \\ \\ \relax
22:17 & \begin{tabularx}{0.7\textwidth}{X} 我必多多賜福給你，我必使你的後裔大大增多，如同天上的星、海邊的沙。你的後裔必得仇敵的城門， \end{tabularx} \\ \\ \relax
22:18 & \begin{tabularx}{0.7\textwidth}{X} 並且地上的萬國都必因你的後裔得福，因為你聽從了我的話。』」 \end{tabularx} \\ \\ \relax
22:19 & \begin{tabularx}{0.7\textwidth}{X} 於是亞伯拉罕回到他僕人那裡。他們一同起身，往別是巴去，亞伯拉罕就住在別是巴。 \end{tabularx} \\ \\ \relax
22:20 & \begin{tabularx}{0.7\textwidth}{X} 這些事以後，有人告訴亞伯拉罕說：「看哪，密迦也為你兄弟拿鶴生了幾個兒子： \end{tabularx} \\ \\ \relax
22:21 & \begin{tabularx}{0.7\textwidth}{X} 長子烏斯、他的兄弟布斯、亞蘭的父親基摩利、 \end{tabularx} \\ \\ \relax
22:22 & \begin{tabularx}{0.7\textwidth}{X} 基薛、哈瑣、必達、益拉和彼土利。」 \end{tabularx} \\ \\ \relax
22:23 & \begin{tabularx}{0.7\textwidth}{X} 彼土利生利百加。這八個人都是密迦為亞伯拉罕的兄弟拿鶴生的。 \end{tabularx} \\ \\ \relax
22:24 & \begin{tabularx}{0.7\textwidth}{X} 拿鶴的妾名叫流瑪，她也生了提八、迦含、他轄和瑪迦。 \end{tabularx} \\ \\ \relax
23:1 & \begin{tabularx}{0.7\textwidth}{X} 撒拉享壽一百二十七歲，這是撒拉一生的歲數。 \end{tabularx} \\ \\ \relax
23:2 & \begin{tabularx}{0.7\textwidth}{X} 撒拉死在迦南地的基列‧亞巴，就是希伯崙。亞伯拉罕來哀悼撒拉，為她哭泣。 \end{tabularx} \\ \\ \relax
23:3 & \begin{tabularx}{0.7\textwidth}{X} 然後，亞伯拉罕起來，離開死人面前，對赫人說： \end{tabularx} \\ \\ \relax
23:4 & \begin{tabularx}{0.7\textwidth}{X} 「我在你們中間是外人，是寄居的。請給我你們那裡的一塊墳地，我好埋葬我的亡妻，使她不在我的面前。」 \end{tabularx} \\ \\ \relax
23:5 & \begin{tabularx}{0.7\textwidth}{X} 赫人回答亞伯拉罕說： \end{tabularx} \\ \\ \relax
23:6 & \begin{tabularx}{0.7\textwidth}{X} 「我主請聽。你在我們中間是一位尊貴的王子，只管在我們最好的墳地裡埋葬你的死人；我們沒有一人會拒絕你在他的墳地裡埋葬你的死人。」 \end{tabularx} \\ \\ \relax
23:7 & \begin{tabularx}{0.7\textwidth}{X} 於是，亞伯拉罕起來，向當地的百姓赫人下拜， \end{tabularx} \\ \\ \relax
23:8 & \begin{tabularx}{0.7\textwidth}{X} 對他們說：「你們若願意讓我埋葬我的亡妻，使她不在我面前，就請聽我，為我求瑣轄的兒子以弗崙， \end{tabularx} \\ \\ \relax
23:9 & \begin{tabularx}{0.7\textwidth}{X} 把他田地盡頭的麥比拉洞賣給我。他可以按照足價賣給我，作為我在你們中間的墳地。」 \end{tabularx} \\ \\ \relax
23:10 & \begin{tabularx}{0.7\textwidth}{X} 那時，以弗崙正坐在赫人中間。赫人以弗崙就回答亞伯拉罕，說給所有出入城門的赫人聽： \end{tabularx} \\ \\ \relax
23:11 & \begin{tabularx}{0.7\textwidth}{X} 「不，我主請聽。我要把這塊田送給你，連田間的洞也送給你，在我同族的人眼前都給你，讓你埋葬你的死人。」 \end{tabularx} \\ \\ \relax
23:12 & \begin{tabularx}{0.7\textwidth}{X} 亞伯拉罕就在當地的百姓面前下拜， \end{tabularx} \\ \\ \relax
23:13 & \begin{tabularx}{0.7\textwidth}{X} 對以弗崙說，也給當地百姓聽：「你若應允，請你聽我。我要把田的價錢給你，請你收下，我就在那裡埋葬我的死人。」 \end{tabularx} \\ \\ \relax
23:14 & \begin{tabularx}{0.7\textwidth}{X} 以弗崙回答亞伯拉罕說： \end{tabularx} \\ \\ \relax
23:15 & \begin{tabularx}{0.7\textwidth}{X} 「我主請聽。四百舍客勒銀子的地，在你我中間算甚麼呢？只管埋葬你的死人吧！」 \end{tabularx} \\ \\ \relax
23:16 & \begin{tabularx}{0.7\textwidth}{X} 亞伯拉罕聽從了以弗崙。亞伯拉罕就照著他說給赫人聽的，把買賣通用的銀子，秤了四百舍客勒銀子給以弗崙。 \end{tabularx} \\ \\ \relax
23:17 & \begin{tabularx}{0.7\textwidth}{X} 於是，以弗崙把那塊位於幔利對面的麥比拉田，和其中的洞，以及田間周圍的樹木都成交了， \end{tabularx} \\ \\ \relax
23:18 & \begin{tabularx}{0.7\textwidth}{X} 在所有出入城門的赫人眼前，賣給亞伯拉罕作為他的產業。 \end{tabularx} \\ \\ \relax
23:19 & \begin{tabularx}{0.7\textwidth}{X} 後來，亞伯拉罕把他妻子撒拉安葬在迦南地幔利對面的麥比拉田間的洞裡，幔利就是希伯崙。 \end{tabularx} \\ \\ \relax
23:20 & \begin{tabularx}{0.7\textwidth}{X} 從此，那塊田和田間的洞就從赫人移交給亞伯拉罕作墳地的產業。 \end{tabularx} \\ \\ \relax
24:1 & \begin{tabularx}{0.7\textwidth}{X} 亞伯拉罕年紀老邁，耶和華在一切事上都賜福給他。 \end{tabularx} \\ \\ \relax
24:2 & \begin{tabularx}{0.7\textwidth}{X} 亞伯拉罕對他家中管理他一切產業最老的僕人說：「把你的手放在我大腿底下。 \end{tabularx} \\ \\ \relax
24:3 & \begin{tabularx}{0.7\textwidth}{X} 我要叫你指著耶和華—天和地的神起誓，不要為我兒子娶我所居住的迦南地的女子為妻。 \end{tabularx} \\ \\ \relax
24:4 & \begin{tabularx}{0.7\textwidth}{X} 你要往我的本地本族去，為我的兒子以撒娶妻。」 \end{tabularx} \\ \\ \relax
24:5 & \begin{tabularx}{0.7\textwidth}{X} 僕人對他說：「如果那女子不肯跟我來到這地，我必須把你的兒子帶回到你出來的地方嗎？」 \end{tabularx} \\ \\ \relax
24:6 & \begin{tabularx}{0.7\textwidth}{X} 亞伯拉罕對他說：「你要謹慎，不可帶我兒子回那裡去。 \end{tabularx} \\ \\ \relax
24:7 & \begin{tabularx}{0.7\textwidth}{X} 耶和華—天上的神曾帶領我離開父家和本族的地，對我說話，向我起誓說：『我要將這地賜給你的後裔。』他要差遣使者在你面前，你就可以從那裡為我兒子娶妻。 \end{tabularx} \\ \\ \relax
24:8 & \begin{tabularx}{0.7\textwidth}{X} 倘若那女子不肯跟你來，我叫你起的誓就與你無關了，只是你不可帶我的兒子回到那裡去。」 \end{tabularx} \\ \\ \relax
24:9 & \begin{tabularx}{0.7\textwidth}{X} 僕人就把手放在他主人亞伯拉罕的大腿底下，為這事向他起誓。 \end{tabularx} \\ \\ \relax
24:10 & \begin{tabularx}{0.7\textwidth}{X} 那僕人從他主人的駱駝中取了十匹駱駝，他手中也帶著他主人各樣的貴重物品離開，起身往美索不達米亞去，到了拿鶴的城。 \end{tabularx} \\ \\ \relax
24:11 & \begin{tabularx}{0.7\textwidth}{X} 傍晚時，眾女子出來打水，他就讓駱駝跪在城外的水井旁。 \end{tabularx} \\ \\ \relax
24:12 & \begin{tabularx}{0.7\textwidth}{X} 他說：「耶和華—我主人亞伯拉罕的神啊，求你施恩給我的主人亞伯拉罕，讓我今日就遇見吧！ \end{tabularx} \\ \\ \relax
24:13 & \begin{tabularx}{0.7\textwidth}{X} 看哪，我站在井旁，城內居民的女子們正出來打水。 \end{tabularx} \\ \\ \relax
24:14 & \begin{tabularx}{0.7\textwidth}{X} 我向哪一個少女說：『請你放下水瓶來，給我水喝』，她若說：『請喝！我也給你的駱駝喝』，願她作你所選定給你僕人以撒的妻。這樣，我就知道你施恩給我的主人了。」 \end{tabularx} \\ \\ \relax
24:15 & \begin{tabularx}{0.7\textwidth}{X} 話還沒說完，看哪，利百加肩頭上扛著水瓶出來。利百加是彼土利所生的；彼土利是亞伯拉罕的兄弟拿鶴妻子密迦的兒子。 \end{tabularx} \\ \\ \relax
24:16 & \begin{tabularx}{0.7\textwidth}{X} 那少女容貌極其美麗，是未曾與人親近的童女。她下到井旁，打滿了瓶子的水，就上來。 \end{tabularx} \\ \\ \relax
24:17 & \begin{tabularx}{0.7\textwidth}{X} 僕人跑上前去迎著她，說：「請你讓我喝你瓶子裡的一點水。」 \end{tabularx} \\ \\ \relax
24:18 & \begin{tabularx}{0.7\textwidth}{X} 少女說：「我主請喝！」就急忙拿下瓶子托在手上，給他喝水。 \end{tabularx} \\ \\ \relax
24:19 & \begin{tabularx}{0.7\textwidth}{X} 那少女給他喝足了，又說：「我也為你的駱駝打水，直到駱駝喝足了。」 \end{tabularx} \\ \\ \relax
24:20 & \begin{tabularx}{0.7\textwidth}{X} 她就急忙把瓶子裡的水倒在槽裡，又跑到井旁打水，為所有的駱駝打了水。 \end{tabularx} \\ \\ \relax
24:21 & \begin{tabularx}{0.7\textwidth}{X} 那人定睛看著少女，一句話也不說，要知道耶和華是否使他的道路亨通。 \end{tabularx} \\ \\ \relax
24:22 & \begin{tabularx}{0.7\textwidth}{X} 駱駝喝足了，那人就拿出一個比加重的金環，一對十舍客勒重的金手鐲， \end{tabularx} \\ \\ \relax
24:23 & \begin{tabularx}{0.7\textwidth}{X} 說：「請告訴我，你是誰的女兒？你父親家裡有沒有地方可以讓我們過夜？」 \end{tabularx} \\ \\ \relax
24:24 & \begin{tabularx}{0.7\textwidth}{X} 少女說：「我是密迦為拿鶴生的兒子彼土利的女兒。」 \end{tabularx} \\ \\ \relax
24:25 & \begin{tabularx}{0.7\textwidth}{X} 又說：「我們家裡有充足的乾草和飼料，也有住宿的地方。」 \end{tabularx} \\ \\ \relax
24:26 & \begin{tabularx}{0.7\textwidth}{X} 那人就低頭向耶和華敬拜， \end{tabularx} \\ \\ \relax
24:27 & \begin{tabularx}{0.7\textwidth}{X} 說：「耶和華—我主人亞伯拉罕的神是應當稱頌的，因他不斷以慈愛信實待我主人。至於我，耶和華一路引領我，直到我主人的兄弟家裡。」 \end{tabularx} \\ \\ \relax
24:28 & \begin{tabularx}{0.7\textwidth}{X} 那少女跑去，把這些話告訴她母親家裡的人。 \end{tabularx} \\ \\ \relax
24:29 & \begin{tabularx}{0.7\textwidth}{X} 利百加有一個哥哥，名叫拉班，拉班就跑到外面井旁那人那裡。 \end{tabularx} \\ \\ \relax
24:30 & \begin{tabularx}{0.7\textwidth}{X} 當他看見金環和戴在他妹妹手上的金鐲，又聽見他妹妹利百加說的話：「那人如此對我說」，他就來到那人面前，看哪，他還站在井旁的駱駝旁邊， \end{tabularx} \\ \\ \relax
24:31 & \begin{tabularx}{0.7\textwidth}{X} 就對他說：「你這蒙耶和華賜福的人，請進來吧！為甚麼站在外面？我已經收拾了房屋，也為駱駝預備了地方。」 \end{tabularx} \\ \\ \relax
24:32 & \begin{tabularx}{0.7\textwidth}{X} 那人就進了拉班的家。拉班卸了駱駝，用飼料餵牠們，拿水給那人和隨從他的人洗腳， \end{tabularx} \\ \\ \relax
24:33 & \begin{tabularx}{0.7\textwidth}{X} 把食物擺在他面前，請他吃。他卻說：「我不吃，等我把我的事情說完了再吃。」拉班說：「請說。」 \end{tabularx} \\ \\ \relax
24:34 & \begin{tabularx}{0.7\textwidth}{X} 他說：「我是亞伯拉罕的僕人。 \end{tabularx} \\ \\ \relax
24:35 & \begin{tabularx}{0.7\textwidth}{X} 耶和華大大地賜福給我主人，使他發達，賜給他羊群、牛群、金銀、奴僕、婢女、駱駝和驢。 \end{tabularx} \\ \\ \relax
24:36 & \begin{tabularx}{0.7\textwidth}{X} 我主人的妻子撒拉年老的時候為我主人生了一個兒子；我主人把他一切所有的都給了他。 \end{tabularx} \\ \\ \relax
24:37 & \begin{tabularx}{0.7\textwidth}{X} 我主人叫我起誓說：『不要為我兒子娶我所居住的迦南地的女子為妻。 \end{tabularx} \\ \\ \relax
24:38 & \begin{tabularx}{0.7\textwidth}{X} 你要往我父家、我本族那裡去，為我的兒子娶妻。』 \end{tabularx} \\ \\ \relax
24:39 & \begin{tabularx}{0.7\textwidth}{X} 我對我主人說：『恐怕那女子不肯跟我來。』 \end{tabularx} \\ \\ \relax
24:40 & \begin{tabularx}{0.7\textwidth}{X} 他就說：『我所事奉的耶和華必要差遣他的使者與你同去，使你的道路亨通，你就可以在我父家、我本族那裡，為我的兒子娶妻。 \end{tabularx} \\ \\ \relax
24:41 & \begin{tabularx}{0.7\textwidth}{X} 只要你到了我本族那裡，我叫你起的誓就與你無關。他們若不把女子交給你，我叫你起的誓也與你無關。』 \end{tabularx} \\ \\ \relax
24:42 & \begin{tabularx}{0.7\textwidth}{X} 「我今日到了井旁，就說：『耶和華—我主人亞伯拉罕的神啊，願你使我所行的道路亨通。 \end{tabularx} \\ \\ \relax
24:43 & \begin{tabularx}{0.7\textwidth}{X} 看哪，我站在井旁，對哪一個出來打水的女子說：請你讓我喝你瓶子裡的一點水， \end{tabularx} \\ \\ \relax
24:44 & \begin{tabularx}{0.7\textwidth}{X} 她若說：你只管喝，我也為你的駱駝打水；願那女子作耶和華給我主人兒子所選定的妻子。』 \end{tabularx} \\ \\ \relax
24:45 & \begin{tabularx}{0.7\textwidth}{X} 「我心裡的話還沒有說完，看哪，利百加肩頭上扛著水瓶出來，下到井旁打水。我對她說：『請你給我水喝。』 \end{tabularx} \\ \\ \relax
24:46 & \begin{tabularx}{0.7\textwidth}{X} 她就急忙從肩頭上拿下瓶子來，說：『請喝！我也給你的駱駝喝。』我就喝了；她也給我的駱駝喝了。 \end{tabularx} \\ \\ \relax
24:47 & \begin{tabularx}{0.7\textwidth}{X} 我問她說：『你是誰的女兒？』她說：『我是彼土利的女兒，彼土利是密迦和拿鶴生的兒子。』我就把環子戴在她鼻子上，把鐲子戴在她雙手上。 \end{tabularx} \\ \\ \relax
24:48 & \begin{tabularx}{0.7\textwidth}{X} 然後我低頭向耶和華敬拜，稱頌耶和華—我主人亞伯拉罕的神，因為他引導我走合適的道路，使我得著我主人兄弟的孫女，給我主人的兒子為妻。 \end{tabularx} \\ \\ \relax
24:49 & \begin{tabularx}{0.7\textwidth}{X} 現在你們若願以慈愛誠信待我主人，就告訴我；若不然，也告訴我，使我可以或向左，或向右。」 \end{tabularx} \\ \\ \relax
24:50 & \begin{tabularx}{0.7\textwidth}{X} 拉班和彼土利回答說：「這事既然出於耶和華，我們不能向你說好說歹。 \end{tabularx} \\ \\ \relax
24:51 & \begin{tabularx}{0.7\textwidth}{X} 看哪，利百加就在你面前，可以將她帶去，遵照耶和華所說的，給你主人的兒子為妻。」 \end{tabularx} \\ \\ \relax
24:52 & \begin{tabularx}{0.7\textwidth}{X} 亞伯拉罕的僕人聽見他們這些話，就向耶和華俯伏在地。 \end{tabularx} \\ \\ \relax
24:53 & \begin{tabularx}{0.7\textwidth}{X} 僕人拿出金器、銀器和衣服送給利百加，又將貴重的物品送給她哥哥和她母親。 \end{tabularx} \\ \\ \relax
24:54 & \begin{tabularx}{0.7\textwidth}{X} 然後，僕人和隨從的人才吃喝，並且住了一夜。早晨起來，僕人說：「請讓我回我主人那裡去吧。」 \end{tabularx} \\ \\ \relax
24:55 & \begin{tabularx}{0.7\textwidth}{X} 利百加的哥哥和母親說：「讓她同我們再住幾天，也許十天，然後她可以去。」 \end{tabularx} \\ \\ \relax
24:56 & \begin{tabularx}{0.7\textwidth}{X} 僕人對他們說：「耶和華既然使我道路亨通，你們就不要耽誤我，請讓我走，回我主人那裡去吧！」 \end{tabularx} \\ \\ \relax
24:57 & \begin{tabularx}{0.7\textwidth}{X} 他們說：「我們把她叫來問問她。」 \end{tabularx} \\ \\ \relax
24:58 & \begin{tabularx}{0.7\textwidth}{X} 他們就叫了利百加來，對她說：「你和這人同去嗎？」她說：「我去。」 \end{tabularx} \\ \\ \relax
24:59 & \begin{tabularx}{0.7\textwidth}{X} 於是他們送他們的妹妹利百加和她的奶媽，同亞伯拉罕的僕人，以及隨從他的人走了。 \end{tabularx} \\ \\ \relax
24:60 & \begin{tabularx}{0.7\textwidth}{X} 他們就為利百加祝福，對她說：「我們的妹妹啊，願你作千萬人的母親！願你的後裔得著仇敵的城門！」 \end{tabularx} \\ \\ \relax
24:61 & \begin{tabularx}{0.7\textwidth}{X} 利百加和她的女僕們起來，騎上駱駝，跟著那人去。僕人就帶著利百加走了。 \end{tabularx} \\ \\ \relax
24:62 & \begin{tabularx}{0.7\textwidth}{X} 那時，以撒住在尼革夫。他剛從庇耳‧拉海‧萊回來。 \end{tabularx} \\ \\ \relax
24:63 & \begin{tabularx}{0.7\textwidth}{X} 傍晚時，以撒出來，到田間默想。他舉目一看，看哪，來了一隊駱駝。 \end{tabularx} \\ \\ \relax
24:64 & \begin{tabularx}{0.7\textwidth}{X} 利百加舉目看見以撒，就急忙下了駱駝， \end{tabularx} \\ \\ \relax
24:65 & \begin{tabularx}{0.7\textwidth}{X} 對那僕人說：「這從田間走來迎接我們的人是誰？」僕人說：「他是我的主人。」利百加就拿面紗蓋住自己。 \end{tabularx} \\ \\ \relax
24:66 & \begin{tabularx}{0.7\textwidth}{X} 僕人把他所做的一切事都告訴以撒。 \end{tabularx} \\ \\ \relax
24:67 & \begin{tabularx}{0.7\textwidth}{X} 以撒就領利百加進了母親撒拉的帳棚，娶了她為妻，並且愛她。以撒自從母親離世以後，這才得了安慰。 \end{tabularx} \\ \\ \relax
25:1 & \begin{tabularx}{0.7\textwidth}{X} 亞伯拉罕再娶了一個妻子，名叫基土拉。 \end{tabularx} \\ \\ \relax
25:2 & \begin{tabularx}{0.7\textwidth}{X} 她為他生了心蘭、約珊、米但、米甸、伊施巴和書亞。 \end{tabularx} \\ \\ \relax
25:3 & \begin{tabularx}{0.7\textwidth}{X} 約珊生了示巴和底但。底但的子孫是亞書利族、利都是族和利烏米族。 \end{tabularx} \\ \\ \relax
25:4 & \begin{tabularx}{0.7\textwidth}{X} 米甸的兒子是以法、以弗、哈諾、亞比大和以勒大。這些都是基土拉的子孫。 \end{tabularx} \\ \\ \relax
25:5 & \begin{tabularx}{0.7\textwidth}{X} 亞伯拉罕把他一切所有的都給了以撒。 \end{tabularx} \\ \\ \relax
25:6 & \begin{tabularx}{0.7\textwidth}{X} 至於亞伯拉罕妾的兒子，亞伯拉罕趁著自己還活著的時候把財物分給他們，打發他們離開他的兒子以撒，往東方去，直到東方之地。 \end{tabularx} \\ \\ \relax
25:7 & \begin{tabularx}{0.7\textwidth}{X} 這是亞伯拉罕一生的年日，他活了一百七十五年。 \end{tabularx} \\ \\ \relax
25:8 & \begin{tabularx}{0.7\textwidth}{X} 亞伯拉罕壽高年邁，安享天年，息勞而終，歸到他祖先 那裡。 \end{tabularx} \\ \\ \relax
25:9 & \begin{tabularx}{0.7\textwidth}{X} 他兩個兒子以撒、以實瑪利把他安葬在麥比拉洞裡。這洞在幔利的對面、赫人瑣轄的兒子以弗崙的田中， \end{tabularx} \\ \\ \relax
25:10 & \begin{tabularx}{0.7\textwidth}{X} 就是亞伯拉罕向赫人買的那塊田。亞伯拉罕和他妻子撒拉都葬在那裡。 \end{tabularx} \\ \\ \relax
25:11 & \begin{tabularx}{0.7\textwidth}{X} 亞伯拉罕死了以後，神賜福給他的兒子以撒。以撒住在庇耳‧拉海‧萊附近。 \end{tabularx} \\ \\ \relax
25:12 & \begin{tabularx}{0.7\textwidth}{X} 這是撒拉的婢女、埃及人夏甲為亞伯拉罕生的兒子以實瑪利的後代。 \end{tabularx} \\ \\ \relax
25:13 & \begin{tabularx}{0.7\textwidth}{X} 以實瑪利兒子們的名字，按著他們後代的名字如下：以實瑪利的長子尼拜約，又有基達、亞德別、米比衫、 \end{tabularx} \\ \\ \relax
25:14 & \begin{tabularx}{0.7\textwidth}{X} 米施瑪、度瑪、瑪撒、 \end{tabularx} \\ \\ \relax
25:15 & \begin{tabularx}{0.7\textwidth}{X} 哈大、提瑪、伊突、拿非施和基底瑪。 \end{tabularx} \\ \\ \relax
25:16 & \begin{tabularx}{0.7\textwidth}{X} 這些都是以實瑪利的兒子們。他們的村莊和營寨按著他們命名；他們作了十二族的族長。 \end{tabularx} \\ \\ \relax
25:17 & \begin{tabularx}{0.7\textwidth}{X} 以實瑪利一生的歲數是一百三十七歲，斷氣而死，歸到他祖先那裡。 \end{tabularx} \\ \\ \relax
25:18 & \begin{tabularx}{0.7\textwidth}{X} 他的子孫住在哈腓拉，直到埃及東邊的書珥，向著亞述，在他眾弟兄的對面安頓下來。 \end{tabularx} \\ \\ \relax
25:19 & \begin{tabularx}{0.7\textwidth}{X} 這是亞伯拉罕的兒子以撒的後代。亞伯拉罕生以撒。 \end{tabularx} \\ \\ \relax
25:20 & \begin{tabularx}{0.7\textwidth}{X} 以撒四十歲時娶利百加為妻。利百加是巴旦‧亞蘭地的亞蘭人彼土利的女兒，是亞蘭人拉班的妹妹。 \end{tabularx} \\ \\ \relax
25:21 & \begin{tabularx}{0.7\textwidth}{X} 以撒因他妻子不生育，就為她祈求耶和華。耶和華應允他的祈求，他的妻子利百加就懷了孕。 \end{tabularx} \\ \\ \relax
25:22 & \begin{tabularx}{0.7\textwidth}{X} 胎兒們在她腹中彼此相爭，她就說：「若是如此，我為甚麼會這樣呢？」她就去求問耶和華。 \end{tabularx} \\ \\ \relax
25:23 & \begin{tabularx}{0.7\textwidth}{X} 耶和華對她說：兩國在你腹中；兩族要從你身上分立。這族必強於那族；將來大的要服侍小的。 \end{tabularx} \\ \\ \relax
25:24 & \begin{tabularx}{0.7\textwidth}{X} 到了生產的日期，看哪，腹中是對雙胞胎。 \end{tabularx} \\ \\ \relax
25:25 & \begin{tabularx}{0.7\textwidth}{X} 先出生的身體帶紅，渾身有毛，好像皮衣；他們就給他起名叫以掃。 \end{tabularx} \\ \\ \relax
25:26 & \begin{tabularx}{0.7\textwidth}{X} 隨後，以掃的弟弟也出生，他的手抓住以掃的腳跟，因此給他起名叫雅各。兩個兒子出生時，以撒六十歲。 \end{tabularx} \\ \\ \relax
25:27 & \begin{tabularx}{0.7\textwidth}{X} 兩個孩子漸漸長大，以掃善於打獵，常在田野；雅各為人安靜，常住在帳棚裡。 \end{tabularx} \\ \\ \relax
25:28 & \begin{tabularx}{0.7\textwidth}{X} 以撒愛以掃，因為常吃他的野味；利百加卻愛雅各。 \end{tabularx} \\ \\ \relax
25:29 & \begin{tabularx}{0.7\textwidth}{X} 有一天，雅各熬了湯，以掃從田野回來，疲憊不堪。 \end{tabularx} \\ \\ \relax
25:30 & \begin{tabularx}{0.7\textwidth}{X} 以掃對雅各說：「我累死了，請你讓我吃這紅的，這紅的湯吧！」因此以掃又叫以東。 \end{tabularx} \\ \\ \relax
25:31 & \begin{tabularx}{0.7\textwidth}{X} 雅各說：「你今日把長子的名分賣給我吧。」 \end{tabularx} \\ \\ \relax
25:32 & \begin{tabularx}{0.7\textwidth}{X} 以掃說：「看哪，我快要死了，這長子的名分對我有甚麼用呢？」 \end{tabularx} \\ \\ \relax
25:33 & \begin{tabularx}{0.7\textwidth}{X} 雅各說：「你今日對我起誓吧。」以掃就向他起誓，把長子的名分賣給了雅各。 \end{tabularx} \\ \\ \relax
25:34 & \begin{tabularx}{0.7\textwidth}{X} 於是雅各把餅和豆湯給了以掃，以掃吃喝以後，起來走了。這樣，以掃輕看他長子的名分。 \end{tabularx} \\ \\ \relax
26:1 & \begin{tabularx}{0.7\textwidth}{X} 那地有了饑荒，不是亞伯拉罕的時候曾有過的那次饑荒，以撒就到基拉耳，非利士人的王亞比米勒那裡去。 \end{tabularx} \\ \\ \relax
26:2 & \begin{tabularx}{0.7\textwidth}{X} 耶和華向以撒顯現，說：「你不要下埃及去，要住在我所指示你的地。 \end{tabularx} \\ \\ \relax
26:3 & \begin{tabularx}{0.7\textwidth}{X} 你要寄居在這地，我必與你同在，賜福給你，因為我要將這一切的地都賜給你和你的後裔。我必堅定我向你父親亞伯拉罕所起的誓。 \end{tabularx} \\ \\ \relax
26:4 & \begin{tabularx}{0.7\textwidth}{X} 我要使你的後裔增多，好像天上的星，又要將這一切的地賜給你的後裔，並且地上的萬國都必因你的後裔得福， \end{tabularx} \\ \\ \relax
26:5 & \begin{tabularx}{0.7\textwidth}{X} 因為亞伯拉罕聽從我的話，遵守我的吩咐、誡令、律例和教導。」 \end{tabularx} \\ \\ \relax
26:6 & \begin{tabularx}{0.7\textwidth}{X} 於是，以撒住在基拉耳。 \end{tabularx} \\ \\ \relax
26:7 & \begin{tabularx}{0.7\textwidth}{X} 那地方的人問起他的妻子，他就說：「她是我的妹妹。」原來他害怕說「我的妻子」。他想：「或許這地方的人會因利百加殺我，因為她容貌美麗。」 \end{tabularx} \\ \\ \relax
26:8 & \begin{tabularx}{0.7\textwidth}{X} 他在那裡住了一段很長的日子。有一天，非利士人的王亞比米勒從窗戶往外觀看，看哪，以撒在撫愛他的妻子利百加。 \end{tabularx} \\ \\ \relax
26:9 & \begin{tabularx}{0.7\textwidth}{X} 亞比米勒召以撒來，說：「看哪，她實在是你的妻子，你怎麼說『她是我的妹妹』呢？」以撒對他說：「因為我想，恐怕我會因她而死。」 \end{tabularx} \\ \\ \relax
26:10 & \begin{tabularx}{0.7\textwidth}{X} 亞比米勒說：「你向我們做的是甚麼事呢？百姓中有一個人幾乎要和你的妻子同寢，你就把我們陷在罪中了。」 \end{tabularx} \\ \\ \relax
26:11 & \begin{tabularx}{0.7\textwidth}{X} 於是亞比米勒命令眾百姓說：「凡侵犯這個人，或他妻子的，必要把他處死。」 \end{tabularx} \\ \\ \relax
26:12 & \begin{tabularx}{0.7\textwidth}{X} 以撒在那地耕種，那一年有百倍的收成。耶和華賜福給他， \end{tabularx} \\ \\ \relax
26:13 & \begin{tabularx}{0.7\textwidth}{X} 他就發達，日漸昌盛，成了大富翁。 \end{tabularx} \\ \\ \relax
26:14 & \begin{tabularx}{0.7\textwidth}{X} 他有羊群牛群，又有許多僕人，非利士人就嫉妒他。 \end{tabularx} \\ \\ \relax
26:15 & \begin{tabularx}{0.7\textwidth}{X} 他父親亞伯拉罕在世的時候，他父親的僕人所挖的井，非利士人全都塞住，填滿了土。 \end{tabularx} \\ \\ \relax
26:16 & \begin{tabularx}{0.7\textwidth}{X} 亞比米勒對以撒說：「你離開我們去吧，因為你比我們強盛得多。」 \end{tabularx} \\ \\ \relax
26:17 & \begin{tabularx}{0.7\textwidth}{X} 以撒就離開那裡，在基拉耳谷支搭帳棚，住在那裡。 \end{tabularx} \\ \\ \relax
26:18 & \begin{tabularx}{0.7\textwidth}{X} 他父親亞伯拉罕在世的時候所挖的水井，在亞伯拉罕死後，都被非利士人塞住了，以撒就重新把井挖出來，仍照他父親所取的名為它們命名。 \end{tabularx} \\ \\ \relax
26:19 & \begin{tabularx}{0.7\textwidth}{X} 以撒的僕人在谷中挖井，就在那裡得了一口活水井。 \end{tabularx} \\ \\ \relax
26:20 & \begin{tabularx}{0.7\textwidth}{X} 基拉耳的牧人與以撒的牧人相爭，說：「這水是我們的。」以撒就給那井起名叫埃色，因為他們和他相爭。 \end{tabularx} \\ \\ \relax
26:21 & \begin{tabularx}{0.7\textwidth}{X} 以撒的僕人又挖了一口井，他們又為這井相爭，以撒就給這井起名叫西提拿。 \end{tabularx} \\ \\ \relax
26:22 & \begin{tabularx}{0.7\textwidth}{X} 以撒離開那裡，又挖了一口井，他們不再為這井相爭了，他就給那井起名叫利河伯。他說：「耶和華現在給我們寬闊之地，我們必在這地興旺。」 \end{tabularx} \\ \\ \relax
26:23 & \begin{tabularx}{0.7\textwidth}{X} 以撒從那裡上別是巴去。 \end{tabularx} \\ \\ \relax
26:24 & \begin{tabularx}{0.7\textwidth}{X} 當夜耶和華向他顯現，說：「我是你父親亞伯拉罕的神。不要懼怕，因為我與你同在，要賜福給你，也要為我僕人亞伯拉罕的緣故，使你的後裔增多。」 \end{tabularx} \\ \\ \relax
26:25 & \begin{tabularx}{0.7\textwidth}{X} 以撒就在那裡築了一座壇，求告耶和華的名，並且在那裡支搭帳棚；他的僕人就在那裡挖了一口井。 \end{tabularx} \\ \\ \relax
26:26 & \begin{tabularx}{0.7\textwidth}{X} 亞比米勒同他的顧問亞戶撒和他軍隊的元帥非各，從基拉耳來到以撒那裡。 \end{tabularx} \\ \\ \relax
26:27 & \begin{tabularx}{0.7\textwidth}{X} 以撒對他們說：「你們既然恨我，趕我離開你們，為甚麼又到我這裡來呢？」 \end{tabularx} \\ \\ \relax
26:28 & \begin{tabularx}{0.7\textwidth}{X} 他們說：「我們明明看見耶和華與你同在；因此就說，讓我們雙方彼此起誓，我們跟你立約， \end{tabularx} \\ \\ \relax
26:29 & \begin{tabularx}{0.7\textwidth}{X} 使你不加害我們，正如我們未曾侵犯你，素來善待你，並且送你平平安安地走。你是蒙耶和華賜福的！」 \end{tabularx} \\ \\ \relax
26:30 & \begin{tabularx}{0.7\textwidth}{X} 以撒為他們擺設宴席，他們就一起吃喝。 \end{tabularx} \\ \\ \relax
26:31 & \begin{tabularx}{0.7\textwidth}{X} 他們清早起來，彼此起誓。以撒送他們走，他們就平平安安地離開他去了。 \end{tabularx} \\ \\ \relax
26:32 & \begin{tabularx}{0.7\textwidth}{X} 那一天，以撒的僕人來，把挖井的消息告訴他，說：「我們得到水了。」 \end{tabularx} \\ \\ \relax
26:33 & \begin{tabularx}{0.7\textwidth}{X} 他就給那井起名叫示巴，因此那城名叫別是巴，直到今日。 \end{tabularx} \\ \\ \relax
26:34 & \begin{tabularx}{0.7\textwidth}{X} 以掃四十歲的時候娶了赫人比利的女兒猶滴，和赫人以倫的女兒巴實抹為妻。 \end{tabularx} \\ \\ \relax
26:35 & \begin{tabularx}{0.7\textwidth}{X} 她們使以撒和利百加心裡愁煩。 \end{tabularx} \\ \\ \relax
27:1 & \begin{tabularx}{0.7\textwidth}{X} 以撒年老，眼睛昏花，不能看見，就叫他大兒子以掃來，對他說：「我兒。」以掃對他說：「我在這裡。」 \end{tabularx} \\ \\ \relax
27:2 & \begin{tabularx}{0.7\textwidth}{X} 他說：「看哪，我老了，不知道哪一天死。 \end{tabularx} \\ \\ \relax
27:3 & \begin{tabularx}{0.7\textwidth}{X} 現在拿你打獵的工具，就是箭囊和弓，到田野去為我打獵， \end{tabularx} \\ \\ \relax
27:4 & \begin{tabularx}{0.7\textwidth}{X} 照我所愛的做成美味，拿來給我吃，好讓我在未死之前為你祝福。」 \end{tabularx} \\ \\ \relax
27:5 & \begin{tabularx}{0.7\textwidth}{X} 以撒對他兒子以掃說話的時候，利百加聽見了。以掃往田野去打獵，要把獵物帶回來。 \end{tabularx} \\ \\ \relax
27:6 & \begin{tabularx}{0.7\textwidth}{X} 利百加就對她兒子雅各說：「看哪，我聽見你父親對你哥哥以掃說： \end{tabularx} \\ \\ \relax
27:7 & \begin{tabularx}{0.7\textwidth}{X} 『你去把獵物帶回來，做成美味給我吃，讓我在未死之前，在耶和華面前為你祝福。』 \end{tabularx} \\ \\ \relax
27:8 & \begin{tabularx}{0.7\textwidth}{X} 現在，我兒，你要聽我的話，照我所吩咐你的， \end{tabularx} \\ \\ \relax
27:9 & \begin{tabularx}{0.7\textwidth}{X} 到羊群裡去，從那裡牽兩隻肥美的小山羊來給我，我就照你父親所愛的，把牠們做成美味給他。 \end{tabularx} \\ \\ \relax
27:10 & \begin{tabularx}{0.7\textwidth}{X} 然後，你拿到你父親那裡給他吃，好讓他在未死之前為你祝福。」 \end{tabularx} \\ \\ \relax
27:11 & \begin{tabularx}{0.7\textwidth}{X} 雅各對他母親利百加說：「看哪，我哥哥以掃渾身都有毛，我身上卻是光滑的； \end{tabularx} \\ \\ \relax
27:12 & \begin{tabularx}{0.7\textwidth}{X} 倘若父親摸著我，我在他眼中就是騙子了。這樣，我就自招詛咒，而不是祝福。」 \end{tabularx} \\ \\ \relax
27:13 & \begin{tabularx}{0.7\textwidth}{X} 他母親對他說：「我兒，你所受的詛咒臨到我身上吧！你只管聽我的話，去牽小山羊來給我。」 \end{tabularx} \\ \\ \relax
27:14 & \begin{tabularx}{0.7\textwidth}{X} 他就去牽來，交給他母親。他母親就照他父親所愛的，做成美味。 \end{tabularx} \\ \\ \relax
27:15 & \begin{tabularx}{0.7\textwidth}{X} 利百加把大兒子以掃在家裡最好的衣服給她小兒子雅各穿， \end{tabularx} \\ \\ \relax
27:16 & \begin{tabularx}{0.7\textwidth}{X} 又用小山羊的皮包在雅各的手上和頸項光滑的地方， \end{tabularx} \\ \\ \relax
27:17 & \begin{tabularx}{0.7\textwidth}{X} 就把所做的美味和餅交在她兒子雅各的手裡。 \end{tabularx} \\ \\ \relax
27:18 & \begin{tabularx}{0.7\textwidth}{X} 雅各來到他父親那裡，說：「我的父親！」他說：「我在這裡。我兒，你是誰？」 \end{tabularx} \\ \\ \relax
27:19 & \begin{tabularx}{0.7\textwidth}{X} 雅各對他父親說：「我是你的長子以掃。我已照你吩咐我的做了。請起來坐著，吃我的野味，你好為我祝福。」 \end{tabularx} \\ \\ \relax
27:20 & \begin{tabularx}{0.7\textwidth}{X} 以撒對他兒子說：「我兒，你怎麼這樣快就找到了呢？」他說：「因為這是耶和華—你的神使我遇見的。」 \end{tabularx} \\ \\ \relax
27:21 & \begin{tabularx}{0.7\textwidth}{X} 以撒對雅各說：「我兒，靠近一點，讓我摸摸你，你真的是我的兒子以掃嗎？」 \end{tabularx} \\ \\ \relax
27:22 & \begin{tabularx}{0.7\textwidth}{X} 雅各就靠近他父親以撒。以撒摸著他，說：「聲音是雅各的聲音，手卻是以掃的手。」 \end{tabularx} \\ \\ \relax
27:23 & \begin{tabularx}{0.7\textwidth}{X} 以撒認不出他來，因為他手上有毛，像他哥哥以掃的手一樣。於是，以撒就為他祝福。 \end{tabularx} \\ \\ \relax
27:24 & \begin{tabularx}{0.7\textwidth}{X} 以撒說：「你真的是我兒子以掃嗎？」他說：「我是。」 \end{tabularx} \\ \\ \relax
27:25 & \begin{tabularx}{0.7\textwidth}{X} 以撒說：「拿給我，讓我吃我兒子的野味，我好為你祝福。」雅各拿給他，他就吃了，又拿酒給他，他也喝了。 \end{tabularx} \\ \\ \relax
27:26 & \begin{tabularx}{0.7\textwidth}{X} 他父親以撒對他說：「我兒，靠近一點來親我！」 \end{tabularx} \\ \\ \relax
27:27 & \begin{tabularx}{0.7\textwidth}{X} 他就近前親吻父親。他父親一聞他衣服上的香氣，就為他祝福，說：「看，我兒的香氣好像耶和華賜福之田地的香氣。 \end{tabularx} \\ \\ \relax
27:28 & \begin{tabularx}{0.7\textwidth}{X} 願神賜你天上的甘露，地上的肥土，和豐富的五穀新酒。 \end{tabularx} \\ \\ \relax
27:29 & \begin{tabularx}{0.7\textwidth}{X} 願萬民事奉你，萬族向你下拜。願你作你弟兄的主，你母親的兒子向你下拜。詛咒你的，願他受詛咒；祝福你的，願他蒙祝福。」 \end{tabularx} \\ \\ \relax
27:30 & \begin{tabularx}{0.7\textwidth}{X} 以撒為雅各祝福完畢，雅各才從他父親那裡出來，他哥哥以掃正打獵回來。 \end{tabularx} \\ \\ \relax
27:31 & \begin{tabularx}{0.7\textwidth}{X} 以掃也做了美味，拿來給他父親，對他父親說：「父親，請起來，吃你兒子的野味，你好為我祝福。」 \end{tabularx} \\ \\ \relax
27:32 & \begin{tabularx}{0.7\textwidth}{X} 他父親以撒對他說：「你是誰？」他說：「我是你的兒子，你的長子以掃。」 \end{tabularx} \\ \\ \relax
27:33 & \begin{tabularx}{0.7\textwidth}{X} 以撒就大大戰兢，說：「那麼，是誰打了獵物拿來給我呢？你未來之前我已經吃了，也為他祝福了，他將來就必蒙福。」 \end{tabularx} \\ \\ \relax
27:34 & \begin{tabularx}{0.7\textwidth}{X} 以掃聽了他父親的話，就大聲痛哭，對他父親說：「我父啊，求你也為我祝福！」 \end{tabularx} \\ \\ \relax
27:35 & \begin{tabularx}{0.7\textwidth}{X} 以撒說：「你弟弟已經用詭計來把你的福分奪去了。」 \end{tabularx} \\ \\ \relax
27:36 & \begin{tabularx}{0.7\textwidth}{X} 以掃說：「他名叫雅各，豈不是這樣嗎？他欺騙了我兩次：他先前奪了我長子的名分，看哪，他現在又奪了我的福分。」以掃又說：「你沒有留下給我的祝福嗎？」 \end{tabularx} \\ \\ \relax
27:37 & \begin{tabularx}{0.7\textwidth}{X} 以撒回答以掃說：「看哪，我已立他作你的主，使他的弟兄都給他作僕人，並賜他五穀新酒可以養生。我兒，那麼，現在我還能為你做甚麼呢？」 \end{tabularx} \\ \\ \relax
27:38 & \begin{tabularx}{0.7\textwidth}{X} 以掃對他父親說：「我父啊，你只有一個祝福嗎？我父啊，求你也為我祝福！」以掃就放聲而哭。 \end{tabularx} \\ \\ \relax
27:39 & \begin{tabularx}{0.7\textwidth}{X} 他父親以撒回答說：「看哪，你所住的地方必缺乏肥沃的土地，缺乏天上的甘露。 \end{tabularx} \\ \\ \relax
27:40 & \begin{tabularx}{0.7\textwidth}{X} 你必倚靠刀劍度日，又必服侍你的兄弟；到你強盛的時候，必從你頸項上掙開他的軛。 \end{tabularx} \\ \\ \relax
27:41 & \begin{tabularx}{0.7\textwidth}{X} 以掃因他父親給雅各的祝福，就怨恨雅各，心裡說：「為我父親居喪的時候近了，到那時候，我要殺我的弟弟雅各。」 \end{tabularx} \\ \\ \relax
27:42 & \begin{tabularx}{0.7\textwidth}{X} 有人把利百加大兒子以掃的話告訴利百加，她就派人去，叫了她小兒子雅各來，對他說：「看哪，你哥哥以掃想要殺你來洩恨。 \end{tabularx} \\ \\ \relax
27:43 & \begin{tabularx}{0.7\textwidth}{X} 現在，我兒，聽我的話，起來，逃往哈蘭，到我哥哥拉班那裡去， \end{tabularx} \\ \\ \relax
27:44 & \begin{tabularx}{0.7\textwidth}{X} 同他住一段日子，直等到你哥哥的怒氣消了。 \end{tabularx} \\ \\ \relax
27:45 & \begin{tabularx}{0.7\textwidth}{X} 等到你哥哥向你消了怒氣，忘了你向他所做的事，我就派人去，把你從那裡帶回來。我何必在一天之內喪失你們二人呢？」 \end{tabularx} \\ \\ \relax
27:46 & \begin{tabularx}{0.7\textwidth}{X} 利百加對以撒說：「我因這赫人的女子活得不耐煩了；倘若雅各也從本地女子中娶像這樣的赫人女子為妻，我為甚麼要活著呢？」 \end{tabularx} \\ \\ \relax
28:1 & \begin{tabularx}{0.7\textwidth}{X} 以撒叫了雅各來，為他祝福，並吩咐他說：「你不要娶迦南的女子為妻。 \end{tabularx} \\ \\ \relax
28:2 & \begin{tabularx}{0.7\textwidth}{X} 你起身往巴旦‧亞蘭去，到你外祖父彼土利的家，從你舅父拉班的女兒中娶一位作你的妻子。 \end{tabularx} \\ \\ \relax
28:3 & \begin{tabularx}{0.7\textwidth}{X} 願全能的神賜福給你，使你生養眾多，成為許多民族， \end{tabularx} \\ \\ \relax
28:4 & \begin{tabularx}{0.7\textwidth}{X} 將應許亞伯拉罕的福賜給你和你的後裔，使你承受你所寄居的地為業，就是神賜給亞伯拉罕的地。」 \end{tabularx} \\ \\ \relax
28:5 & \begin{tabularx}{0.7\textwidth}{X} 以撒送雅各走了，雅各就往巴旦‧亞蘭去，到亞蘭人彼土利的兒子拉班那裡，拉班是利百加的哥哥，利百加是雅各和以掃的母親。 \end{tabularx} \\ \\ \relax
28:6 & \begin{tabularx}{0.7\textwidth}{X} 以掃見以撒已經為雅各祝福，而且送他往巴旦‧亞蘭去，在那裡娶妻，並且見以撒祝福雅各的時候吩咐他說：「不要娶迦南的女子為妻」， \end{tabularx} \\ \\ \relax
28:7 & \begin{tabularx}{0.7\textwidth}{X} 又見雅各聽從父母的話往巴旦‧亞蘭去了， \end{tabularx} \\ \\ \relax
28:8 & \begin{tabularx}{0.7\textwidth}{X} 以掃就看出他父親以撒看迦南女子不順眼。 \end{tabularx} \\ \\ \relax
28:9 & \begin{tabularx}{0.7\textwidth}{X} 於是他往以實瑪利那裡去，在兩個妻子之外， 又娶了瑪哈拉為妻，她是亞伯拉罕兒子以實瑪利的女兒，是尼拜約的妹妹。 \end{tabularx} \\ \\ \relax
28:10 & \begin{tabularx}{0.7\textwidth}{X} 雅各離開別是巴，往哈蘭去。 \end{tabularx} \\ \\ \relax
28:11 & \begin{tabularx}{0.7\textwidth}{X} 到了一個地方，因為已經日落，就在那裡過夜。他拾起那地方的一塊石頭枕在頭下，就躺在那地方。 \end{tabularx} \\ \\ \relax
28:12 & \begin{tabularx}{0.7\textwidth}{X} 他做夢，看哪，一個梯子立在地上，梯子的頂端直伸到天；看哪，神的使者在梯子上，上去下來。 \end{tabularx} \\ \\ \relax
28:13 & \begin{tabularx}{0.7\textwidth}{X} 看哪，耶和華站在梯子上面，說：「我是耶和華—你祖父亞伯拉罕的神，以撒的神。你現在躺臥之地，我要將它賜給你和你的後裔。 \end{tabularx} \\ \\ \relax
28:14 & \begin{tabularx}{0.7\textwidth}{X} 你的後裔必像地上的塵沙，必向東西南北開展；地上萬族必因你和你的後裔得福。 \end{tabularx} \\ \\ \relax
28:15 & \begin{tabularx}{0.7\textwidth}{X} 看哪，我必與你同在，無論你往哪裡去，我必保佑你，領你歸回這地。我總不離棄你，直到我實現了對你所說的話。」 \end{tabularx} \\ \\ \relax
28:16 & \begin{tabularx}{0.7\textwidth}{X} 雅各睡醒了，說：「耶和華真的在這裡，我竟不知道！」 \end{tabularx} \\ \\ \relax
28:17 & \begin{tabularx}{0.7\textwidth}{X} 他就懼怕，說：「這地方何等可畏！這不是別的，是神的殿，是天的門。」 \end{tabularx} \\ \\ \relax
28:18 & \begin{tabularx}{0.7\textwidth}{X} 雅各清早起來，拿起枕在頭下的石頭，立作柱子，澆油在上面。 \end{tabularx} \\ \\ \relax
28:19 & \begin{tabularx}{0.7\textwidth}{X} 他給那地方起名叫伯特利；那地方原先名叫路斯。 \end{tabularx} \\ \\ \relax
28:20 & \begin{tabularx}{0.7\textwidth}{X} 雅各許願說：「神若與我同在，在我所行的路上保佑我，給我食物吃，衣服穿， \end{tabularx} \\ \\ \relax
28:21 & \begin{tabularx}{0.7\textwidth}{X} 使我平平安安回到我父親的家，我就必以耶和華為我的神。 \end{tabularx} \\ \\ \relax
28:22 & \begin{tabularx}{0.7\textwidth}{X} 我所立為柱子的這塊石頭必作神的殿；凡你所賜給我的，我必將十分之一獻給你。」 \end{tabularx} \\ \\ \relax
29:1 & \begin{tabularx}{0.7\textwidth}{X} 雅各起行，到了東方人之地。 \end{tabularx} \\ \\ \relax
29:2 & \begin{tabularx}{0.7\textwidth}{X} 他觀看，看哪，田間有一口井，看哪，有三群羊臥在井旁；因為人都取那井裡的水給羊喝。井口上的那塊石頭很大。 \end{tabularx} \\ \\ \relax
29:3 & \begin{tabularx}{0.7\textwidth}{X} 羊群都在那裡聚集，人就把石頭移開井口，取水給羊喝，然後又把石頭放回井口原處。 \end{tabularx} \\ \\ \relax
29:4 & \begin{tabularx}{0.7\textwidth}{X} 雅各對他們說：「弟兄們，你們從哪裡來？」他們說：「我們是從哈蘭來的。」 \end{tabularx} \\ \\ \relax
29:5 & \begin{tabularx}{0.7\textwidth}{X} 他對他們說：「你們認識拿鶴的孫子拉班嗎？」他們說：「我們認識。」 \end{tabularx} \\ \\ \relax
29:6 & \begin{tabularx}{0.7\textwidth}{X} 雅各對他們說：「他平安嗎？」他們說：「平安。看哪，他女兒拉結和羊一起來了。」 \end{tabularx} \\ \\ \relax
29:7 & \begin{tabularx}{0.7\textwidth}{X} 雅各說：「看哪，日正當中，不是牲畜聚集的時候。你們取水給羊喝，再去牧放吧！」 \end{tabularx} \\ \\ \relax
29:8 & \begin{tabularx}{0.7\textwidth}{X} 他們說：「我們不能這樣，必須等所有的羊群聚集，人把石頭移開井口，我們才可以取水給羊喝。」 \end{tabularx} \\ \\ \relax
29:9 & \begin{tabularx}{0.7\textwidth}{X} 雅各正和他們說話的時候，拉結和她父親的羊來了，因為她是牧羊的。 \end{tabularx} \\ \\ \relax
29:10 & \begin{tabularx}{0.7\textwidth}{X} 雅各看見他舅父拉班的女兒拉結和舅父拉班的羊群，就上前把石頭移開井口，取水給舅父拉班的羊喝。 \end{tabularx} \\ \\ \relax
29:11 & \begin{tabularx}{0.7\textwidth}{X} 雅各親了拉結，就放聲大哭。 \end{tabularx} \\ \\ \relax
29:12 & \begin{tabularx}{0.7\textwidth}{X} 雅各告訴拉結，自己是她父親的親戚，是利百加的兒子。拉結就跑去告訴她父親。 \end{tabularx} \\ \\ \relax
29:13 & \begin{tabularx}{0.7\textwidth}{X} 拉班聽見外甥雅各的消息，就跑去迎接他，抱著他，親他，帶他到自己的家。雅各把這一切的事告訴拉班。 \end{tabularx} \\ \\ \relax
29:14 & \begin{tabularx}{0.7\textwidth}{X} 拉班對他說：「你實在是我的骨肉。」雅各就和他同住了一個月。 \end{tabularx} \\ \\ \relax
29:15 & \begin{tabularx}{0.7\textwidth}{X} 拉班對雅各說：「雖然你是我的親戚，怎麼可以讓你白白服事我呢？告訴我，你要甚麼作工資呢？」 \end{tabularx} \\ \\ \relax
29:16 & \begin{tabularx}{0.7\textwidth}{X} 拉班有兩個女兒，大的名叫利亞，小的名叫拉結。 \end{tabularx} \\ \\ \relax
29:17 & \begin{tabularx}{0.7\textwidth}{X} 利亞的雙眼無神，拉結卻長得美貌秀麗。 \end{tabularx} \\ \\ \relax
29:18 & \begin{tabularx}{0.7\textwidth}{X} 雅各愛拉結，就說：「我願為你的小女兒拉結服事你七年。」 \end{tabularx} \\ \\ \relax
29:19 & \begin{tabularx}{0.7\textwidth}{X} 拉班說：「我把她給你，勝過給別人，你與我同住吧！」 \end{tabularx} \\ \\ \relax
29:20 & \begin{tabularx}{0.7\textwidth}{X} 雅各就為拉結服事了七年；他因為愛拉結，就看這七年如同幾天。 \end{tabularx} \\ \\ \relax
29:21 & \begin{tabularx}{0.7\textwidth}{X} 雅各對拉班說：「日期已經滿了，請把我的妻子給我，我好與她同房。」 \end{tabularx} \\ \\ \relax
29:22 & \begin{tabularx}{0.7\textwidth}{X} 拉班就擺設宴席，請了當地所有的人。 \end{tabularx} \\ \\ \relax
29:23 & \begin{tabularx}{0.7\textwidth}{X} 到了晚上，拉班帶女兒利亞來送給雅各，雅各就與她同房。 \end{tabularx} \\ \\ \relax
29:24 & \begin{tabularx}{0.7\textwidth}{X} 拉班也把自己的婢女悉帕給女兒利亞作婢女。 \end{tabularx} \\ \\ \relax
29:25 & \begin{tabularx}{0.7\textwidth}{X} 到了早晨，看哪，她是利亞，雅各對拉班說：「你向我做的是甚麼事呢？我服事你，不是為拉結嗎？你為甚麼欺騙我呢？」 \end{tabularx} \\ \\ \relax
29:26 & \begin{tabularx}{0.7\textwidth}{X} 拉班說：「大女兒還沒有給人就先把小女兒給人，我們這地方沒有這樣的規矩。 \end{tabularx} \\ \\ \relax
29:27 & \begin{tabularx}{0.7\textwidth}{X} 你先為這個滿了七日，我們就把那個也給你，不過你要另外再服事我七年。」 \end{tabularx} \\ \\ \relax
29:28 & \begin{tabularx}{0.7\textwidth}{X} 雅各就這樣做了。滿了利亞的七日，拉班就把女兒拉結給雅各為妻。 \end{tabularx} \\ \\ \relax
29:29 & \begin{tabularx}{0.7\textwidth}{X} 拉班又把自己的婢女辟拉給女兒拉結作婢女。 \end{tabularx} \\ \\ \relax
29:30 & \begin{tabularx}{0.7\textwidth}{X} 雅各也與拉結同房，並且愛拉結勝過愛利亞，於是他又服事了拉班七年。 \end{tabularx} \\ \\ \relax
29:31 & \begin{tabularx}{0.7\textwidth}{X} 耶和華見利亞失寵，就使她生育，拉結卻不生育。 \end{tabularx} \\ \\ \relax
29:32 & \begin{tabularx}{0.7\textwidth}{X} 利亞懷孕生子，給他起名叫呂便，因為她說：「耶和華看見我的苦情，如今我的丈夫必愛我。」 \end{tabularx} \\ \\ \relax
29:33 & \begin{tabularx}{0.7\textwidth}{X} 她又懷孕生子，給他起名叫西緬，說：「耶和華因為聽見我失寵，所以又賜給我這個兒子。」 \end{tabularx} \\ \\ \relax
29:34 & \begin{tabularx}{0.7\textwidth}{X} 她又懷孕生子，說：「我給丈夫生了三個兒子，現在，這次他必親近我了。」因此，雅各給他起名叫利未。 \end{tabularx} \\ \\ \relax
29:35 & \begin{tabularx}{0.7\textwidth}{X} 她又懷孕生子，說：「這次我要讚美耶和華。」因此給他起名叫猶大。於是她停了生育。 \end{tabularx} \\ \\ \relax
30:1 & \begin{tabularx}{0.7\textwidth}{X} 拉結見自己不給雅各生孩子，就嫉妒她姊姊，對雅各說：「你給我孩子，不然，讓我死了吧。」 \end{tabularx} \\ \\ \relax
30:2 & \begin{tabularx}{0.7\textwidth}{X} 雅各對拉結生氣，說：「是我代替神使你生不出孩子的嗎？」 \end{tabularx} \\ \\ \relax
30:3 & \begin{tabularx}{0.7\textwidth}{X} 拉結說：「看哪，我的使女辟拉在這裡，你可以與她同房，使她生子歸在我膝下，我也可以藉著她得孩子。」 \end{tabularx} \\ \\ \relax
30:4 & \begin{tabularx}{0.7\textwidth}{X} 拉結就把她的婢女辟拉給丈夫為妾，雅各與她同房。 \end{tabularx} \\ \\ \relax
30:5 & \begin{tabularx}{0.7\textwidth}{X} 辟拉懷孕，為雅各生了一個兒子。 \end{tabularx} \\ \\ \relax
30:6 & \begin{tabularx}{0.7\textwidth}{X} 拉結給他起名叫但，說：「神為我伸冤，也聽了我的聲音，賜給我一個兒子。」 \end{tabularx} \\ \\ \relax
30:7 & \begin{tabularx}{0.7\textwidth}{X} 拉結的婢女辟拉又懷孕，為雅各生了第二個兒子。 \end{tabularx} \\ \\ \relax
30:8 & \begin{tabularx}{0.7\textwidth}{X} 拉結給他起名叫拿弗他利，說：「我與我姊姊大大較力，並且得勝了。」 \end{tabularx} \\ \\ \relax
30:9 & \begin{tabularx}{0.7\textwidth}{X} 利亞見自己停了生育，就把她的婢女悉帕給雅各為妾。 \end{tabularx} \\ \\ \relax
30:10 & \begin{tabularx}{0.7\textwidth}{X} 利亞的婢女悉帕為雅各生了一個兒子。 \end{tabularx} \\ \\ \relax
30:11 & \begin{tabularx}{0.7\textwidth}{X} 利亞給他起名叫迦得，說：「真是幸運！」 \end{tabularx} \\ \\ \relax
30:12 & \begin{tabularx}{0.7\textwidth}{X} 利亞的婢女悉帕又為雅各生了第二個兒子。 \end{tabularx} \\ \\ \relax
30:13 & \begin{tabularx}{0.7\textwidth}{X} 利亞給他起名叫亞設，說：「我真有福啊，眾女子都要稱我有福。」 \end{tabularx} \\ \\ \relax
30:14 & \begin{tabularx}{0.7\textwidth}{X} 收割麥子的時候，呂便到田裡去，找到曼陀羅草，就拿給他的母親利亞。拉結對利亞說：「請你給我一些你兒子的曼陀羅草吧。」 \end{tabularx} \\ \\ \relax
30:15 & \begin{tabularx}{0.7\textwidth}{X} 利亞對她說：「你奪走了我的丈夫還是小事嗎？你還要奪取我兒子的曼陀羅草嗎？」拉結說：「今夜他可以與你同寢，來交換你兒子的曼陀羅草。」 \end{tabularx} \\ \\ \relax
30:16 & \begin{tabularx}{0.7\textwidth}{X} 到了晚上，雅各從田裡回來，利亞出來迎接他，說：「你要與我同寢，因為我真的用我兒子的曼陀羅草把你雇下了。」那一夜，雅各就與她同寢。 \end{tabularx} \\ \\ \relax
30:17 & \begin{tabularx}{0.7\textwidth}{X} 神應允了利亞，她就懷孕，為雅各生了第五個兒子。 \end{tabularx} \\ \\ \relax
30:18 & \begin{tabularx}{0.7\textwidth}{X} 利亞給他起名叫以薩迦，說：「神給了我工價，因為我把婢女給了我的丈夫。」 \end{tabularx} \\ \\ \relax
30:19 & \begin{tabularx}{0.7\textwidth}{X} 利亞又懷孕，為雅各生了第六個兒子。 \end{tabularx} \\ \\ \relax
30:20 & \begin{tabularx}{0.7\textwidth}{X} 利亞給他起名叫西布倫，說：「神賜給我厚禮了；這次，我丈夫必看重我，因為我為他生了六個兒子。」 \end{tabularx} \\ \\ \relax
30:21 & \begin{tabularx}{0.7\textwidth}{X} 後來她又生了一個女兒，給她起名叫底拿。 \end{tabularx} \\ \\ \relax
30:22 & \begin{tabularx}{0.7\textwidth}{X} 神顧念拉結，應允她，使她能生育。 \end{tabularx} \\ \\ \relax
30:23 & \begin{tabularx}{0.7\textwidth}{X} 拉結懷孕生子，說：「神除去了我的羞恥。」 \end{tabularx} \\ \\ \relax
30:24 & \begin{tabularx}{0.7\textwidth}{X} 拉結就給他起名叫約瑟，說：「願耶和華再增添一個兒子給我。」 \end{tabularx} \\ \\ \relax
30:25 & \begin{tabularx}{0.7\textwidth}{X} 拉結生約瑟之後，雅各對拉班說：「請讓我走，回到我的本鄉本土去。 \end{tabularx} \\ \\ \relax
30:26 & \begin{tabularx}{0.7\textwidth}{X} 請你把我服事你所得的妻子和孩子給我，讓我走吧！我怎樣服事你，你都知道。」 \end{tabularx} \\ \\ \relax
30:27 & \begin{tabularx}{0.7\textwidth}{X} 拉班對他說：「願你看得起我，因我占卜得知，耶和華賜福給我是因你的緣故。」 \end{tabularx} \\ \\ \relax
30:28 & \begin{tabularx}{0.7\textwidth}{X} 又說：「請為我定你的工資，我就給你。」 \end{tabularx} \\ \\ \relax
30:29 & \begin{tabularx}{0.7\textwidth}{X} 雅各對他說：「我怎樣服事你，你的牲畜在我這裡變得怎樣，你都知道。 \end{tabularx} \\ \\ \relax
30:30 & \begin{tabularx}{0.7\textwidth}{X} 我未來以前，你擁有的很少，現在卻已大量增加，因為耶和華隨著我的腳步賜福給你。現在，我到甚麼時候才可以成家立業呢？」 \end{tabularx} \\ \\ \relax
30:31 & \begin{tabularx}{0.7\textwidth}{X} 拉班說：「我該給你甚麼呢？」雅各說：「你甚麼也不必給我，只要你為我做這件事，我就繼續牧放你的羊群。 \end{tabularx} \\ \\ \relax
30:32 & \begin{tabularx}{0.7\textwidth}{X} 今天我要走遍你的羊群，把綿羊中凡有點的、有斑的，和小綿羊中凡是黑色的羊；以及山羊中凡有斑的、有點的，都從那裡挑出來，作為我的工資。 \end{tabularx} \\ \\ \relax
30:33 & \begin{tabularx}{0.7\textwidth}{X} 以後你來當面查看我的工資，任何我這裡的山羊不是有點有斑的，小綿羊不是黑色的，就算是我偷的。這就可以證明我是正直的。」 \end{tabularx} \\ \\ \relax
30:34 & \begin{tabularx}{0.7\textwidth}{X} 拉班說：「看哪，就照你所說的做吧。」 \end{tabularx} \\ \\ \relax
30:35 & \begin{tabularx}{0.7\textwidth}{X} 當日，拉班把有紋的、有斑的公山羊，一切有點的、有斑的、有少許白色的母山羊，以及小綿羊中所有黑色的，都挑出來，交在他兒子們的手裡， \end{tabularx} \\ \\ \relax
30:36 & \begin{tabularx}{0.7\textwidth}{X} 又使自己和雅各相隔三天的路程。雅各就牧放拉班其餘的羊。 \end{tabularx} \\ \\ \relax
30:37 & \begin{tabularx}{0.7\textwidth}{X} 雅各拿楊樹、杏樹、楓樹的嫩枝，把皮剝出白色的條紋，使枝子露出白色來。 \end{tabularx} \\ \\ \relax
30:38 & \begin{tabularx}{0.7\textwidth}{X} 他把剝了皮的枝子對著羊群，插在羊喝水的水溝和水槽裡。羊來喝水的時候，牠們彼此交配。 \end{tabularx} \\ \\ \relax
30:39 & \begin{tabularx}{0.7\textwidth}{X} 羊對著枝子交配，就生下有紋的、有點的、有斑的來。 \end{tabularx} \\ \\ \relax
30:40 & \begin{tabularx}{0.7\textwidth}{X} 雅各把小綿羊分出來，讓羊對著拉班羊群中有紋的和所有黑色的。於是他把自己的羊群分開，不叫牠們和拉班的羊混在一起。 \end{tabularx} \\ \\ \relax
30:41 & \begin{tabularx}{0.7\textwidth}{X} 當肥壯的羊交配的時候，雅各就把枝子插在水溝裡，使羊對著枝子交配。 \end{tabularx} \\ \\ \relax
30:42 & \begin{tabularx}{0.7\textwidth}{X} 可是當瘦弱的羊交配的時候，他就不插枝子。這樣，瘦弱的就歸拉班，肥壯的就歸雅各。 \end{tabularx} \\ \\ \relax
30:43 & \begin{tabularx}{0.7\textwidth}{X} 於是這人極其發達，擁有許多的羊群、奴僕、婢女、駱駝和驢。 \end{tabularx} \\ \\ \relax
31:1 & \begin{tabularx}{0.7\textwidth}{X} 雅各聽見拉班兒子們的話，說：「雅各把我們父親所有的都奪去了！他從我們父親所擁有的獲得這一切的財富。」 \end{tabularx} \\ \\ \relax
31:2 & \begin{tabularx}{0.7\textwidth}{X} 雅各見拉班的臉色，看哪，待他不如從前了。 \end{tabularx} \\ \\ \relax
31:3 & \begin{tabularx}{0.7\textwidth}{X} 耶和華對雅各說：「你要回你祖先之地，到你本族那裡去，我必與你同在。」 \end{tabularx} \\ \\ \relax
31:4 & \begin{tabularx}{0.7\textwidth}{X} 雅各就派人叫拉結和利亞到田野他的羊群那裡去， \end{tabularx} \\ \\ \relax
31:5 & \begin{tabularx}{0.7\textwidth}{X} 對她們說：「我看你們父親待我的臉色不如從前了，但我父親的神向來與我同在。 \end{tabularx} \\ \\ \relax
31:6 & \begin{tabularx}{0.7\textwidth}{X} 你們也知道，我盡了全力服事你們的父親。 \end{tabularx} \\ \\ \relax
31:7 & \begin{tabularx}{0.7\textwidth}{X} 可是你們的父親欺騙我，十次更改我的工資，但神不容許他害我。 \end{tabularx} \\ \\ \relax
31:8 & \begin{tabularx}{0.7\textwidth}{X} 他若說：『有點的歸給你作工資』，羊群所生的都是有點的；他若說：『有紋的歸給你作工資』，羊群所生的都是有紋的。 \end{tabularx} \\ \\ \relax
31:9 & \begin{tabularx}{0.7\textwidth}{X} 這樣，神把你們父親的牲畜拿來賜給我了。 \end{tabularx} \\ \\ \relax
31:10 & \begin{tabularx}{0.7\textwidth}{X} 「羊群交配的時候，我在夢中舉目一看，看哪，跳母羊的公羊都是有紋的、有點的、有花斑的。 \end{tabularx} \\ \\ \relax
31:11 & \begin{tabularx}{0.7\textwidth}{X} 神的使者在夢中呼叫我說：『雅各。』我說：『我在這裡。』 \end{tabularx} \\ \\ \relax
31:12 & \begin{tabularx}{0.7\textwidth}{X} 他說：『你舉目觀看，跳母羊的公羊都是有紋的、有點的、有花斑的。拉班向你所做的一切，我都看見了。 \end{tabularx} \\ \\ \relax
31:13 & \begin{tabularx}{0.7\textwidth}{X} 我是伯特利的神；你曾在那裡用油膏過柱子，向我許過願。現在你起來，離開這地，回你本族之地去吧！』」 \end{tabularx} \\ \\ \relax
31:14 & \begin{tabularx}{0.7\textwidth}{X} 拉結和利亞回答雅各說：「在我們父親家裡還有我們可分得的產業嗎？ \end{tabularx} \\ \\ \relax
31:15 & \begin{tabularx}{0.7\textwidth}{X} 我們不是被他看作外人嗎？因為他賣了我們，還吞吃了我們的銀錢。 \end{tabularx} \\ \\ \relax
31:16 & \begin{tabularx}{0.7\textwidth}{X} 神從我們父親所拿走的一切財物，都是我們和我們孩子的。現在，凡神所吩咐你的，你只管去做吧！」 \end{tabularx} \\ \\ \relax
31:17 & \begin{tabularx}{0.7\textwidth}{X} 雅各起來，叫他的孩子和妻子都騎上駱駝， \end{tabularx} \\ \\ \relax
31:18 & \begin{tabularx}{0.7\textwidth}{X} 又趕著他一切的牲畜和他所得的一切財物，就是他在巴旦‧亞蘭所得的，他擁有的牲畜，往迦南地他父親以撒那裡去了。 \end{tabularx} \\ \\ \relax
31:19 & \begin{tabularx}{0.7\textwidth}{X} 當時拉班去剪羊毛，拉結偷了他父親家中的神像。 \end{tabularx} \\ \\ \relax
31:20 & \begin{tabularx}{0.7\textwidth}{X} 雅各瞞住亞蘭人拉班，不通知他就逃走了。 \end{tabularx} \\ \\ \relax
31:21 & \begin{tabularx}{0.7\textwidth}{X} 雅各帶著他所有的逃走了；他起程，渡過大河，面向著基列山。 \end{tabularx} \\ \\ \relax
31:22 & \begin{tabularx}{0.7\textwidth}{X} 到第三天，有人告訴拉班，雅各逃跑了。 \end{tabularx} \\ \\ \relax
31:23 & \begin{tabularx}{0.7\textwidth}{X} 拉班帶著他的弟兄們去追他，追了七天，就在基列山追上了。 \end{tabularx} \\ \\ \relax
31:24 & \begin{tabularx}{0.7\textwidth}{X} 夜間，神來到亞蘭人拉班那裡，在夢中對他說：「你要小心，不可對雅各說好說歹。」 \end{tabularx} \\ \\ \relax
31:25 & \begin{tabularx}{0.7\textwidth}{X} 拉班追上雅各。雅各在山上支搭帳棚；拉班和他的弟兄們也在基列山上支搭帳棚。 \end{tabularx} \\ \\ \relax
31:26 & \begin{tabularx}{0.7\textwidth}{X} 拉班對雅各說：「你做的是甚麼事呢？你瞞著我把我的女兒們帶走，好像用刀劍擄去一般。 \end{tabularx} \\ \\ \relax
31:27 & \begin{tabularx}{0.7\textwidth}{X} 你為甚麼暗暗地逃跑，瞞著我，不通知我一聲，叫我可以歡樂、唱歌、擊鼓、彈琴送你回去呢？ \end{tabularx} \\ \\ \relax
31:28 & \begin{tabularx}{0.7\textwidth}{X} 為甚麼不容許我與外孫和女兒吻別呢？你現在所做的真是愚蠢！ \end{tabularx} \\ \\ \relax
31:29 & \begin{tabularx}{0.7\textwidth}{X} 我的手本有能力害你，只是你父親的神昨夜對我說：『你要小心，不可對雅各說好說歹。』 \end{tabularx} \\ \\ \relax
31:30 & \begin{tabularx}{0.7\textwidth}{X} 現在你既然這麼想念你的父家，不得不去，為甚麼又偷了我的神明呢？」 \end{tabularx} \\ \\ \relax
31:31 & \begin{tabularx}{0.7\textwidth}{X} 雅各回答拉班說：「因為我害怕，我想，恐怕你把你的女兒從我這裡奪走。 \end{tabularx} \\ \\ \relax
31:32 & \begin{tabularx}{0.7\textwidth}{X} 至於你的神明，你若在誰那裡搜出來，就不讓誰活。當著我們弟兄面前，你認一認在我這裡有甚麼東西是你的，你就拿去吧。」原來雅各並不知道拉結偷了神明。 \end{tabularx} \\ \\ \relax
31:33 & \begin{tabularx}{0.7\textwidth}{X} 拉班進了雅各、利亞，以及兩個使女的帳棚，卻沒有找到，就從利亞的帳棚出來，進入拉結的帳棚。 \end{tabularx} \\ \\ \relax
31:34 & \begin{tabularx}{0.7\textwidth}{X} 拉結拿了神像，藏在駱駝的鞍子裡，自己坐在上面。拉班搜遍了那帳棚，並沒有找到。 \end{tabularx} \\ \\ \relax
31:35 & \begin{tabularx}{0.7\textwidth}{X} 拉結對她父親說：「請我主不要生氣，因為我恰有月事，不能在你面前起來。」拉班搜尋，卻找不到神像。 \end{tabularx} \\ \\ \relax
31:36 & \begin{tabularx}{0.7\textwidth}{X} 於是雅各發怒，斥責拉班。雅各對拉班說：「我有甚麼過犯，有甚麼罪惡，你竟這樣火速地追我？ \end{tabularx} \\ \\ \relax
31:37 & \begin{tabularx}{0.7\textwidth}{X} 你搜遍了我一切的物件，你找到甚麼呢？可以放在你我弟兄面前，叫他們在我們兩個之間評評理。 \end{tabularx} \\ \\ \relax
31:38 & \begin{tabularx}{0.7\textwidth}{X} 我在你那裡這二十年，你的母綿羊、母山羊沒有掉過胎。你羊群中的公綿羊，我沒有吃過； \end{tabularx} \\ \\ \relax
31:39 & \begin{tabularx}{0.7\textwidth}{X} 被野獸撕裂的，我沒有帶來給你，是我自己賠償的。無論是白日被偷的，或是黑夜被偷的，你都從我手中索取。 \end{tabularx} \\ \\ \relax
31:40 & \begin{tabularx}{0.7\textwidth}{X} 我常常白日受盡炎熱，黑夜受盡寒霜，不得合眼入睡。 \end{tabularx} \\ \\ \relax
31:41 & \begin{tabularx}{0.7\textwidth}{X} 我這二十年在你家裡，為你兩個女兒服事了你十四年，為你的羊群服事了你六年，你卻十次更改我的工資。 \end{tabularx} \\ \\ \relax
31:42 & \begin{tabularx}{0.7\textwidth}{X} 若不是我父親以撒所敬畏的神，就是亞伯拉罕的神與我同在，你如今必定打發我空手而去。神看見我的苦情和我手的辛勞，就在昨夜責備了你。」 \end{tabularx} \\ \\ \relax
31:43 & \begin{tabularx}{0.7\textwidth}{X} 拉班回答雅各說：「這兩個女兒是我的女兒，這些孩子是我的孩子，這些羊群也都是我的羊群；凡你所看見的都是我的。我的女兒和她們所生的孩子，我今日還能對他們做甚麼呢？ \end{tabularx} \\ \\ \relax
31:44 & \begin{tabularx}{0.7\textwidth}{X} 現在，來吧！讓我和你立約，作你我之間的證據。」 \end{tabularx} \\ \\ \relax
31:45 & \begin{tabularx}{0.7\textwidth}{X} 雅各就拿一塊石頭立作柱子， \end{tabularx} \\ \\ \relax
31:46 & \begin{tabularx}{0.7\textwidth}{X} 對弟兄們說：「大家來堆積石頭。」他們拿石頭堆成一堆，於是在那裡，在石堆旁邊吃喝。 \end{tabularx} \\ \\ \relax
31:47 & \begin{tabularx}{0.7\textwidth}{X} 拉班稱那石堆為伊迦爾‧撒哈杜他，雅各卻稱那石堆為迦累得。 \end{tabularx} \\ \\ \relax
31:48 & \begin{tabularx}{0.7\textwidth}{X} 拉班說：「今日這石堆成為你我之間的證據。」因此這地方名叫迦累得， \end{tabularx} \\ \\ \relax
31:49 & \begin{tabularx}{0.7\textwidth}{X} 又叫米斯巴，因為他說：「我們彼此離別以後，願耶和華在你我中間鑒察 。 \end{tabularx} \\ \\ \relax
31:50 & \begin{tabularx}{0.7\textwidth}{X} 你若苦待我的女兒，或在我的女兒以外另娶妻，雖沒有人在場，你看，有神在你我中間作證。」 \end{tabularx} \\ \\ \relax
31:51 & \begin{tabularx}{0.7\textwidth}{X} 拉班又對雅各說：「看哪，這石堆，看哪，這柱子，是我在你我中間所立的。 \end{tabularx} \\ \\ \relax
31:52 & \begin{tabularx}{0.7\textwidth}{X} 這石堆是證據，這柱子也是證據。我必不越過這石堆去害你；你也不可越過這石堆和柱子來害我。 \end{tabularx} \\ \\ \relax
31:53 & \begin{tabularx}{0.7\textwidth}{X} 願亞伯拉罕的神和拿鶴的神，就是他們父親的神，在你我中間判斷。」雅各就指著他父親以撒所敬畏的神起誓， \end{tabularx} \\ \\ \relax
31:54 & \begin{tabularx}{0.7\textwidth}{X} 又在山上獻祭，請弟兄們來吃飯。他們吃了飯，就在山上過夜。 \end{tabularx} \\ \\ \relax
31:55 & \begin{tabularx}{0.7\textwidth}{X} 拉班清早起來，與他外孫和女兒親吻，為他們祝福，就回到自己的地方去了。 \end{tabularx} \\ \\ \relax
32:1 & \begin{tabularx}{0.7\textwidth}{X} 雅各繼續行路，神的使者遇見他。 \end{tabularx} \\ \\ \relax
32:2 & \begin{tabularx}{0.7\textwidth}{X} 雅各看見他們就說：「這是神的軍營。」於是給那地方起名叫瑪哈念 。 \end{tabularx} \\ \\ \relax
32:3 & \begin{tabularx}{0.7\textwidth}{X} 雅各派使者在他前面到西珥地，就是以東地他哥哥以掃那裡。 \end{tabularx} \\ \\ \relax
32:4 & \begin{tabularx}{0.7\textwidth}{X} 他吩咐他們說：「你們要對我主以掃說：『你的僕人雅各這樣說：我在拉班那裡寄居，延遲到如今。 \end{tabularx} \\ \\ \relax
32:5 & \begin{tabularx}{0.7\textwidth}{X} 我有牛、驢、羊群、奴僕、婢女，現在派人來報告我主，為了要在你眼前蒙恩。』」 \end{tabularx} \\ \\ \relax
32:6 & \begin{tabularx}{0.7\textwidth}{X} 使者回到雅各那裡，說：「我們到了你哥哥以掃那裡。他正迎著你來，並且有四百人和他一起。」 \end{tabularx} \\ \\ \relax
32:7 & \begin{tabularx}{0.7\textwidth}{X} 雅各就很懼怕，而且愁煩。他把跟他同行的人和羊群、牛群、駱駝分成兩隊， \end{tabularx} \\ \\ \relax
32:8 & \begin{tabularx}{0.7\textwidth}{X} 說：「以掃若來擊殺其中一隊，剩下的另一隊還可以逃脫。」 \end{tabularx} \\ \\ \relax
32:9 & \begin{tabularx}{0.7\textwidth}{X} 雅各說：「耶和華—我祖父亞伯拉罕的神，我父親以撒的神啊，你曾對我說：『回你本地本族去，我要厚待你。』 \end{tabularx} \\ \\ \relax
32:10 & \begin{tabularx}{0.7\textwidth}{X} 你向僕人所施的一切慈愛和信實，我一點也不配得。我先前只用我的一根杖過這約旦河，如今我卻成了兩隊。 \end{tabularx} \\ \\ \relax
32:11 & \begin{tabularx}{0.7\textwidth}{X} 求你救我脫離我哥哥的手，脫離以掃的手，因為我怕他來殺我，連母親和兒女都不放過。 \end{tabularx} \\ \\ \relax
32:12 & \begin{tabularx}{0.7\textwidth}{X} 你曾說：『我必定厚待你，使你的後裔如同海邊的沙，多得不可勝數。』」 \end{tabularx} \\ \\ \relax
32:13 & \begin{tabularx}{0.7\textwidth}{X} 當夜，雅各在那裡住宿，就從他手中所擁有的拿禮物要送給他哥哥以掃， \end{tabularx} \\ \\ \relax
32:14 & \begin{tabularx}{0.7\textwidth}{X} 就是二百隻母山羊、二十隻公山羊、二百隻母綿羊、二十隻公綿羊、 \end{tabularx} \\ \\ \relax
32:15 & \begin{tabularx}{0.7\textwidth}{X} 三十匹哺乳的母駱駝和牠們的小駱駝、四十頭母牛、十頭公牛、二十匹母驢和十匹公驢。 \end{tabularx} \\ \\ \relax
32:16 & \begin{tabularx}{0.7\textwidth}{X} 他把每種牲畜各分一群，交在僕人手中，對僕人說：「你們要在我的前頭過去，使群和群之間保持一段距離」。 \end{tabularx} \\ \\ \relax
32:17 & \begin{tabularx}{0.7\textwidth}{X} 他又吩咐領頭的人說：「我哥哥以掃遇見你的時候，問你說：『你是誰的人？要往哪裡去？你前面這些是誰的？』 \end{tabularx} \\ \\ \relax
32:18 & \begin{tabularx}{0.7\textwidth}{X} 你就說：『是你僕人雅各的，是送給我主以掃的禮物。看哪，他自己也在我們後面。』」 \end{tabularx} \\ \\ \relax
32:19 & \begin{tabularx}{0.7\textwidth}{X} 他又吩咐第二、第三和所有趕畜群的人說：「你們遇見以掃的時候要照這樣的話對他說， \end{tabularx} \\ \\ \relax
32:20 & \begin{tabularx}{0.7\textwidth}{X} 你們還要說：『看哪，你僕人雅各在我們後面。』」因雅各說：「我藉著在我前面送去的禮物給他面子，然後再見他的面，或許他會寬容我。」 \end{tabularx} \\ \\ \relax
32:21 & \begin{tabularx}{0.7\textwidth}{X} 於是禮物在他前面過去了；那夜，雅各在營中住宿。 \end{tabularx} \\ \\ \relax
32:22 & \begin{tabularx}{0.7\textwidth}{X} 他夜間起來，帶著兩個妻子，兩個婢女和十一個孩子，過了雅博渡口。 \end{tabularx} \\ \\ \relax
32:23 & \begin{tabularx}{0.7\textwidth}{X} 他帶著他們，送他們過河，他所有的一切也都過去， \end{tabularx} \\ \\ \relax
32:24 & \begin{tabularx}{0.7\textwidth}{X} 只剩下雅各一人。有一個人來和他摔跤，直到黎明。 \end{tabularx} \\ \\ \relax
32:25 & \begin{tabularx}{0.7\textwidth}{X} 那人見自己勝不過他，就摸了他的大腿窩一下。雅各的大腿窩就在和那人摔跤的時候扭了。 \end{tabularx} \\ \\ \relax
32:26 & \begin{tabularx}{0.7\textwidth}{X} 那人說：「天快亮了，讓我走吧！」雅各說：「你不給我祝福，我就不讓你走。」 \end{tabularx} \\ \\ \relax
32:27 & \begin{tabularx}{0.7\textwidth}{X} 那人說：「你叫甚麼名字？」他說：「雅各。」 \end{tabularx} \\ \\ \relax
32:28 & \begin{tabularx}{0.7\textwidth}{X} 那人說：「你的名字不要再叫雅各，要叫以色列，因為你與神和人較力，都得勝了。」 \end{tabularx} \\ \\ \relax
32:29 & \begin{tabularx}{0.7\textwidth}{X} 雅各問他說：「請告訴我你的名字。」那人說：「何必問我的名字呢？」於是他在那裡為雅各祝福。 \end{tabularx} \\ \\ \relax
32:30 & \begin{tabularx}{0.7\textwidth}{X} 雅各就給那地方起名叫毗努伊勒 ，說：「我面對面見了神，我的性命仍得保全。」 \end{tabularx} \\ \\ \relax
32:31 & \begin{tabularx}{0.7\textwidth}{X} 太陽剛出來的時候，雅各經過毗努伊勒，他的大腿就瘸了。 \end{tabularx} \\ \\ \relax
32:32 & \begin{tabularx}{0.7\textwidth}{X} 因此，以色列人不吃大腿窩的筋，直到今日，因為那人摸了雅各大腿窩的筋。 \end{tabularx} \\ \\ \relax
33:1 & \begin{tabularx}{0.7\textwidth}{X} 雅各舉目觀看，看哪，以掃來了，有四百人和他一起。雅各就把孩子們分開交給利亞、拉結和兩個婢女。 \end{tabularx} \\ \\ \relax
33:2 & \begin{tabularx}{0.7\textwidth}{X} 他叫兩個婢女和她們的孩子走在前頭，利亞和她的孩子跟在後面，而拉結和約瑟在最後。 \end{tabularx} \\ \\ \relax
33:3 & \begin{tabularx}{0.7\textwidth}{X} 他自己卻走到他們前面，一連七次俯伏在地才挨近他哥哥。 \end{tabularx} \\ \\ \relax
33:4 & \begin{tabularx}{0.7\textwidth}{X} 以掃跑來迎接他，將他抱住，伏在他的頸項上親他，他們都哭了。 \end{tabularx} \\ \\ \relax
33:5 & \begin{tabularx}{0.7\textwidth}{X} 以掃舉目看見婦人和孩子，就說：「這些和你一起的是誰呢？」雅各說：「這些孩子是神施恩給你僕人的。」 \end{tabularx} \\ \\ \relax
33:6 & \begin{tabularx}{0.7\textwidth}{X} 於是兩個婢女和她們的孩子前來下拜， \end{tabularx} \\ \\ \relax
33:7 & \begin{tabularx}{0.7\textwidth}{X} 利亞和她的孩子也前來下拜，隨後約瑟和拉結也前來下拜。 \end{tabularx} \\ \\ \relax
33:8 & \begin{tabularx}{0.7\textwidth}{X} 以掃說：「我所遇見的這些畜群是甚麼意思呢？」雅各說：「是為了要在我主眼前蒙恩。」 \end{tabularx} \\ \\ \relax
33:9 & \begin{tabularx}{0.7\textwidth}{X} 以掃說：「弟弟啊，我的已經夠了，你的你自己留著吧！」 \end{tabularx} \\ \\ \relax
33:10 & \begin{tabularx}{0.7\textwidth}{X} 雅各說：「不，我若在你眼前蒙恩，就請你從我手裡收下這禮物；因為我見了你的面，如同見了神的面，並且你也寬容了我。 \end{tabularx} \\ \\ \relax
33:11 & \begin{tabularx}{0.7\textwidth}{X} 請你收下我帶來給你的禮物，因為神恩待我，使我一切都充足。」雅各再三求他，他才收下。 \end{tabularx} \\ \\ \relax
33:12 & \begin{tabularx}{0.7\textwidth}{X} 以掃說：「讓我們起身前行，我和你一起走吧。」 \end{tabularx} \\ \\ \relax
33:13 & \begin{tabularx}{0.7\textwidth}{X} 雅各對他說：「我主知道孩子們還年幼嬌嫩，我的牛羊也正在哺乳中，只要催趕一天，群羊都會死了。 \end{tabularx} \\ \\ \relax
33:14 & \begin{tabularx}{0.7\textwidth}{X} 請我主在僕人前面先走，我要按著在我面前的牲畜和孩子的步伐慢慢前進，直走到西珥我主那裡。」 \end{tabularx} \\ \\ \relax
33:15 & \begin{tabularx}{0.7\textwidth}{X} 以掃說：「讓我把跟隨我的人留幾個在你這裡。」雅各說：「何必這樣呢？只要能在我主眼前蒙恩就夠了。」 \end{tabularx} \\ \\ \relax
33:16 & \begin{tabularx}{0.7\textwidth}{X} 於是，以掃當日起行，回西珥去了。 \end{tabularx} \\ \\ \relax
33:17 & \begin{tabularx}{0.7\textwidth}{X} 雅各就往疏割去，在那裡為自己蓋房屋，又為牲畜搭棚，因此那地方叫疏割。 \end{tabularx} \\ \\ \relax
33:18 & \begin{tabularx}{0.7\textwidth}{X} 雅各從巴旦‧亞蘭平安地回到迦南地的示劍城，他在城的前面支搭帳棚。 \end{tabularx} \\ \\ \relax
33:19 & \begin{tabularx}{0.7\textwidth}{X} 他用一百可錫塔從示劍的父親哈抹的眾子手中買了搭帳棚的那塊地。 \end{tabularx} \\ \\ \relax
33:20 & \begin{tabularx}{0.7\textwidth}{X} 雅各在那裡築了一座壇，起名叫伊利‧伊羅伊‧以色列。 \end{tabularx} \\ \\ \relax
34:1 & \begin{tabularx}{0.7\textwidth}{X} 利亞給雅各所生的女兒底拿出去，要探望那地的女子們。 \end{tabularx} \\ \\ \relax
34:2 & \begin{tabularx}{0.7\textwidth}{X} 那地的族長希未人哈抹的兒子示劍看見她，就拉住她，與她同寢，玷辱了她。 \end{tabularx} \\ \\ \relax
34:3 & \begin{tabularx}{0.7\textwidth}{X} 示劍的心喜歡雅各的女兒底拿，愛上這少女，甜言蜜語地安慰她。 \end{tabularx} \\ \\ \relax
34:4 & \begin{tabularx}{0.7\textwidth}{X} 示劍對他父親哈抹說：「求你為我聘這女孩為妻。」 \end{tabularx} \\ \\ \relax
34:5 & \begin{tabularx}{0.7\textwidth}{X} 雅各聽見示劍污辱了他的女兒底拿。那時他的兒子們正和牲畜在田野，雅各就沉默，等他們回來。 \end{tabularx} \\ \\ \relax
34:6 & \begin{tabularx}{0.7\textwidth}{X} 示劍的父親哈抹出來，到雅各那裡，要和他講話。 \end{tabularx} \\ \\ \relax
34:7 & \begin{tabularx}{0.7\textwidth}{X} 雅各的兒子們聽見這事，就從田野回來，人人悲憤，十分惱怒，因示劍在以色列中做了醜事，與雅各的女兒同寢，這本是不該做的事。 \end{tabularx} \\ \\ \relax
34:8 & \begin{tabularx}{0.7\textwidth}{X} 哈抹和他們談話，說：「我兒子示劍的心喜歡你們家的女兒，請你們把她嫁給我的兒子。 \end{tabularx} \\ \\ \relax
34:9 & \begin{tabularx}{0.7\textwidth}{X} 你們與我們彼此結親；你們可以把你們家的女兒嫁給我們，也可以娶我們家的女兒。 \end{tabularx} \\ \\ \relax
34:10 & \begin{tabularx}{0.7\textwidth}{X} 你們與我們同住吧！這地都在你們面前，只管在這裡居住，做買賣，置產業。」 \end{tabularx} \\ \\ \relax
34:11 & \begin{tabularx}{0.7\textwidth}{X} 示劍對女子的父親和兄弟們說：「願你們看得起我，你們向我要甚麼，我必給你們， \end{tabularx} \\ \\ \relax
34:12 & \begin{tabularx}{0.7\textwidth}{X} 無論向我要多貴重的聘金和禮物，我必照你們所說的給你們，只要你們將這少女嫁給我。」 \end{tabularx} \\ \\ \relax
34:13 & \begin{tabularx}{0.7\textwidth}{X} 雅各的兒子們因示劍污辱了他們的妹妹底拿，就用詭詐的話回答示劍和他父親哈抹， \end{tabularx} \\ \\ \relax
34:14 & \begin{tabularx}{0.7\textwidth}{X} 對他們說：「我們不能做這樣的事，把我們的妹妹嫁給沒有受割禮的人為妻，因為那是我們的羞恥。 \end{tabularx} \\ \\ \relax
34:15 & \begin{tabularx}{0.7\textwidth}{X} 惟有一個條件，我們才答應你們，就是你們所有的男丁都要受割禮，和我們一樣， \end{tabularx} \\ \\ \relax
34:16 & \begin{tabularx}{0.7\textwidth}{X} 我們就把我們家的女兒嫁給你們，也娶你們家的女兒；我們就與你們同住，大家成為一族人。 \end{tabularx} \\ \\ \relax
34:17 & \begin{tabularx}{0.7\textwidth}{X} 倘若你們不聽從我們受割禮，我們就帶我們家的女兒走了。」 \end{tabularx} \\ \\ \relax
34:18 & \begin{tabularx}{0.7\textwidth}{X} 這些話在哈抹和他兒子示劍的眼中看為美。 \end{tabularx} \\ \\ \relax
34:19 & \begin{tabularx}{0.7\textwidth}{X} 那年輕人毫不遲延做這事，因為他愛上了雅各的女兒；他在他父親家中也是最受人尊重的。 \end{tabularx} \\ \\ \relax
34:20 & \begin{tabularx}{0.7\textwidth}{X} 哈抹和他兒子示劍到他們的城門口，對城裡的人講說： \end{tabularx} \\ \\ \relax
34:21 & \begin{tabularx}{0.7\textwidth}{X} 「這些人對我們友善，不如允許他們在這地居住，做買賣；看哪，這地寬闊，足以容納他們。我們可以娶他們家的女兒，也可以把我們家的女兒嫁給他們。 \end{tabularx} \\ \\ \relax
34:22 & \begin{tabularx}{0.7\textwidth}{X} 惟有一個條件，這些人才答應和我們同住，成為一族人，就是我們中間所有的男丁都要受割禮，和他們一樣。 \end{tabularx} \\ \\ \relax
34:23 & \begin{tabularx}{0.7\textwidth}{X} 他們的牲畜、財物和一切的牲口豈不都歸給我們嗎？只要答應他們，他們就與我們同住。」 \end{tabularx} \\ \\ \relax
34:24 & \begin{tabularx}{0.7\textwidth}{X} 凡從城門出入的人都聽從了哈抹和他兒子示劍的話。於是，凡從城門出入的男丁都受了割禮。 \end{tabularx} \\ \\ \relax
34:25 & \begin{tabularx}{0.7\textwidth}{X} 到第三天，他們正疼痛的時候，雅各的兩個兒子，就是底拿的哥哥西緬和利未，各拿刀劍，不動聲色地來到城中，把所有的男丁都殺了， \end{tabularx} \\ \\ \relax
34:26 & \begin{tabularx}{0.7\textwidth}{X} 又用刀殺了哈抹和他兒子示劍，把底拿從示劍家裡帶走，就離開了。 \end{tabularx} \\ \\ \relax
34:27 & \begin{tabularx}{0.7\textwidth}{X} 雅各的兒子們因為他們的妹妹受污辱，就來到被殺的人那裡，洗劫那城， \end{tabularx} \\ \\ \relax
34:28 & \begin{tabularx}{0.7\textwidth}{X} 奪走了他們的羊群、牛群和驢，以及城裡和田間所有的； \end{tabularx} \\ \\ \relax
34:29 & \begin{tabularx}{0.7\textwidth}{X} 又俘擄搶劫他們一切的財物、孩童、婦女，以及房屋中所有的。 \end{tabularx} \\ \\ \relax
34:30 & \begin{tabularx}{0.7\textwidth}{X} 雅各對西緬和利未說：「你們連累了我，使我在這地的居民中，就是在迦南人和比利洗人中壞了名聲。我的人丁稀少，他們必聚集來擊殺我，我和全家的人都要被滅絕。」 \end{tabularx} \\ \\ \relax
34:31 & \begin{tabularx}{0.7\textwidth}{X} 他們卻說：「他豈可待我們的妹妹如同妓女呢？」 \end{tabularx} \\ \\ \relax
35:1 & \begin{tabularx}{0.7\textwidth}{X} 神對雅各說：「起來！上伯特利去，住在那裡。在那裡築一座壇給神，就是你逃避你哥哥以掃的時候向你顯現的神。」 \end{tabularx} \\ \\ \relax
35:2 & \begin{tabularx}{0.7\textwidth}{X} 雅各就對他家中的人，以及所有和他一起的人說：「除掉你們中間外邦的神明，要自潔，更換衣服。 \end{tabularx} \\ \\ \relax
35:3 & \begin{tabularx}{0.7\textwidth}{X} 我們要起來，上伯特利去，在那裡我要築一座壇給神，就是在我遭難的日子應允我，在我行走的路上與我同在的神。」 \end{tabularx} \\ \\ \relax
35:4 & \begin{tabularx}{0.7\textwidth}{X} 他們就把手中所有外邦的神明和自己耳朵上的環子交給雅各；雅各把它們埋在示劍那裡的橡樹下。 \end{tabularx} \\ \\ \relax
35:5 & \begin{tabularx}{0.7\textwidth}{X} 他們起行。神使周圍城鎮的人都驚恐，就不追趕雅各的兒子們了。 \end{tabularx} \\ \\ \relax
35:6 & \begin{tabularx}{0.7\textwidth}{X} 於是雅各和所有與他一起的人到了迦南地的路斯，就是伯特利。 \end{tabularx} \\ \\ \relax
35:7 & \begin{tabularx}{0.7\textwidth}{X} 他在那裡築了一座壇，給那地方起名叫伊勒‧伯特利，因為他逃避他哥哥的時候，神曾在那裡向他顯現。 \end{tabularx} \\ \\ \relax
35:8 & \begin{tabularx}{0.7\textwidth}{X} 利百加的奶媽底波拉死了，葬在伯特利下邊的橡樹下；那棵樹名叫亞倫‧巴古。 \end{tabularx} \\ \\ \relax
35:9 & \begin{tabularx}{0.7\textwidth}{X} 雅各從巴旦‧亞蘭回來，神又向他顯現，賜福給他。 \end{tabularx} \\ \\ \relax
35:10 & \begin{tabularx}{0.7\textwidth}{X} 神對他說：「你的名原是雅各，從今以後不要再叫雅各，你的名要叫以色列。」於是，神就叫他的名為以色列。 \end{tabularx} \\ \\ \relax
35:11 & \begin{tabularx}{0.7\textwidth}{X} 神又對他說：「我是全能的神；你要生養眾多，將來有一國和許多的國從你而來，又有許多君王從你生出。 \end{tabularx} \\ \\ \relax
35:12 & \begin{tabularx}{0.7\textwidth}{X} 至於我賜給亞伯拉罕和以撒的地，我必賜給你；我必賜這地給你的後裔。」 \end{tabularx} \\ \\ \relax
35:13 & \begin{tabularx}{0.7\textwidth}{X} 神就從與雅各說話的那地方升上去了。 \end{tabularx} \\ \\ \relax
35:14 & \begin{tabularx}{0.7\textwidth}{X} 雅各就在神與他說話的地方立了一根柱子，就是石柱，在它上面獻澆酒祭，又澆油。 \end{tabularx} \\ \\ \relax
35:15 & \begin{tabularx}{0.7\textwidth}{X} 雅各就給神與他說話的那地方起名叫伯特利。 \end{tabularx} \\ \\ \relax
35:16 & \begin{tabularx}{0.7\textwidth}{X} 他們從伯特利起行，到以法他還有一段路程，拉結生產，生得十分艱難。 \end{tabularx} \\ \\ \relax
35:17 & \begin{tabularx}{0.7\textwidth}{X} 她生得十分艱難的時候，接生婆對她說：「不要怕，你又要有一個兒子了。」 \end{tabularx} \\ \\ \relax
35:18 & \begin{tabularx}{0.7\textwidth}{X} 她快要死，還有一口氣的時候，就給她兒子起名叫便‧俄尼；他父親卻給他起名叫便雅憫。 \end{tabularx} \\ \\ \relax
35:19 & \begin{tabularx}{0.7\textwidth}{X} 拉結死了，葬在往以法他的路旁；以法他就是伯利恆。 \end{tabularx} \\ \\ \relax
35:20 & \begin{tabularx}{0.7\textwidth}{X} 雅各在她的墳上立了一塊碑，就是拉結的墓碑，到今日還在。 \end{tabularx} \\ \\ \relax
35:21 & \begin{tabularx}{0.7\textwidth}{X} 以色列起行，在以得臺的那一邊支搭帳棚。 \end{tabularx} \\ \\ \relax
35:22 & \begin{tabularx}{0.7\textwidth}{X} 以色列住在那地的時候，呂便去與他父親的妾辟拉同寢，以色列也聽見了這件事。雅各共有十二個兒子。 \end{tabularx} \\ \\ \relax
35:23 & \begin{tabularx}{0.7\textwidth}{X} 利亞的兒子是雅各的長子呂便，還有西緬、利未、猶大、以薩迦、西布倫。 \end{tabularx} \\ \\ \relax
35:24 & \begin{tabularx}{0.7\textwidth}{X} 拉結的兒子是約瑟、便雅憫。 \end{tabularx} \\ \\ \relax
35:25 & \begin{tabularx}{0.7\textwidth}{X} 拉結的婢女辟拉的兒子是但、拿弗他利。 \end{tabularx} \\ \\ \relax
35:26 & \begin{tabularx}{0.7\textwidth}{X} 利亞的婢女悉帕的兒子是迦得、亞設。這是雅各在巴旦‧亞蘭所生的兒子。 \end{tabularx} \\ \\ \relax
35:27 & \begin{tabularx}{0.7\textwidth}{X} 雅各來到他父親以撒那裡，到了幔利，基列‧亞巴，就是希伯崙，是亞伯拉罕和以撒寄居的地方。 \end{tabularx} \\ \\ \relax
35:28 & \begin{tabularx}{0.7\textwidth}{X} 以撒共活了一百八十年。 \end{tabularx} \\ \\ \relax
35:29 & \begin{tabularx}{0.7\textwidth}{X} 以撒年紀老邁，安享天年，息勞而終，歸到他祖先那裡。他兩個兒子以掃和雅各把他安葬了。 \end{tabularx} \\ \\ \relax
36:1 & \begin{tabularx}{0.7\textwidth}{X} 這是以掃的後代，以掃就是以東。 \end{tabularx} \\ \\ \relax
36:2 & \begin{tabularx}{0.7\textwidth}{X} 以掃娶迦南的女子為妻，就是赫人以倫的女兒亞大和希未人祭便的孫女，亞拿的女兒阿何利巴瑪， \end{tabularx} \\ \\ \relax
36:3 & \begin{tabularx}{0.7\textwidth}{X} 又娶了以實瑪利的女兒，尼拜約的妹妹巴實抹。 \end{tabularx} \\ \\ \relax
36:4 & \begin{tabularx}{0.7\textwidth}{X} 亞大為以掃生了以利法；巴實抹生了流珥； \end{tabularx} \\ \\ \relax
36:5 & \begin{tabularx}{0.7\textwidth}{X} 阿何利巴瑪生了耶烏施、雅蘭、可拉。這些都是以掃的兒子，是在迦南地生的。 \end{tabularx} \\ \\ \relax
36:6 & \begin{tabularx}{0.7\textwidth}{X} 以掃帶著他的妻子、兒女和家中所有的人，以及他的牛羊、牲畜和一切財物，就是他在迦南地所得的，往別處去，離開了他的兄弟雅各。 \end{tabularx} \\ \\ \relax
36:7 & \begin{tabularx}{0.7\textwidth}{X} 因為他們擁有的很多，不能住在一起。因為牲畜的緣故，寄居的地方容不下他們。 \end{tabularx} \\ \\ \relax
36:8 & \begin{tabularx}{0.7\textwidth}{X} 於是以掃住在西珥山；以掃就是以東。 \end{tabularx} \\ \\ \relax
36:9 & \begin{tabularx}{0.7\textwidth}{X} 這是以掃的後代，他是西珥山裡以東人的始祖。 \end{tabularx} \\ \\ \relax
36:10 & \begin{tabularx}{0.7\textwidth}{X} 以掃子孫的名字如下：以掃的妻子亞大生以利法；以掃的妻子巴實抹生流珥。 \end{tabularx} \\ \\ \relax
36:11 & \begin{tabularx}{0.7\textwidth}{X} 以利法的兒子是提幔、阿抹、洗玻、迦坦、基納斯。 \end{tabularx} \\ \\ \relax
36:12 & \begin{tabularx}{0.7\textwidth}{X} 亭納是以掃兒子以利法的妾，她為以利法生了亞瑪力。這是以掃的妻子亞大的子孫。 \end{tabularx} \\ \\ \relax
36:13 & \begin{tabularx}{0.7\textwidth}{X} 流珥的兒子是拿哈、謝拉、沙瑪、米撒。這是以掃妻子巴實抹的子孫。 \end{tabularx} \\ \\ \relax
36:14 & \begin{tabularx}{0.7\textwidth}{X} 以掃的妻子阿何利巴瑪是祭便的孫女，亞拿的女兒。她為以掃生了耶烏施、雅蘭、可拉。 \end{tabularx} \\ \\ \relax
36:15 & \begin{tabularx}{0.7\textwidth}{X} 這是以掃子孫中作族長的：以掃的長子以利法的子孫中，有提幔族長、阿抹族長、洗玻族長、基納斯族長、 \end{tabularx} \\ \\ \relax
36:16 & \begin{tabularx}{0.7\textwidth}{X} 可拉族長、迦坦族長、亞瑪力族長。這是在以東地，從以利法所出的族長，是亞大的子孫。 \end{tabularx} \\ \\ \relax
36:17 & \begin{tabularx}{0.7\textwidth}{X} 以掃的兒子流珥的子孫中，有拿哈族長、謝拉族長、沙瑪族長、米撒族長。這是在以東地，從流珥所出的族長，是以掃妻子巴實抹的子孫。 \end{tabularx} \\ \\ \relax
36:18 & \begin{tabularx}{0.7\textwidth}{X} 以掃的妻子阿何利巴瑪的子孫中，有耶烏施族長、雅蘭族長、可拉族長。這是從以掃的妻子，亞拿的女兒阿何利巴瑪的子孫中所出的族長。 \end{tabularx} \\ \\ \relax
36:19 & \begin{tabularx}{0.7\textwidth}{X} 以上的族長都是以掃的子孫；以掃就是以東。 \end{tabularx} \\ \\ \relax
36:20 & \begin{tabularx}{0.7\textwidth}{X} 這是那地原來的居民，何利人西珥的子孫：羅坍、朔巴、祭便、亞拿、 \end{tabularx} \\ \\ \relax
36:21 & \begin{tabularx}{0.7\textwidth}{X} 底順、以察、底珊。這是在以東地，從何利人西珥子孫中所出的族長。 \end{tabularx} \\ \\ \relax
36:22 & \begin{tabularx}{0.7\textwidth}{X} 羅坍的兒子是何利、希幔，羅坍的妹妹是亭納。 \end{tabularx} \\ \\ \relax
36:23 & \begin{tabularx}{0.7\textwidth}{X} 朔巴的兒子是亞勒文、瑪拿轄、以巴錄、示玻、阿南。 \end{tabularx} \\ \\ \relax
36:24 & \begin{tabularx}{0.7\textwidth}{X} 祭便的兒子是愛亞、亞拿，當時在曠野牧放他父親祭便的驢，發現溫泉的就是這亞拿。 \end{tabularx} \\ \\ \relax
36:25 & \begin{tabularx}{0.7\textwidth}{X} 亞拿的兒子是底順，亞拿的女兒是阿何利巴瑪。 \end{tabularx} \\ \\ \relax
36:26 & \begin{tabularx}{0.7\textwidth}{X} 底順的兒子是欣但、伊是班、益蘭、基蘭。 \end{tabularx} \\ \\ \relax
36:27 & \begin{tabularx}{0.7\textwidth}{X} 以察的兒子是辟罕、撒番、亞干。 \end{tabularx} \\ \\ \relax
36:28 & \begin{tabularx}{0.7\textwidth}{X} 底珊的兒子是烏斯、亞蘭。 \end{tabularx} \\ \\ \relax
36:29 & \begin{tabularx}{0.7\textwidth}{X} 這是從何利人所出的族長：羅坍族長、朔巴族長、祭便族長、亞拿族長、 \end{tabularx} \\ \\ \relax
36:30 & \begin{tabularx}{0.7\textwidth}{X} 底順族長、以察族長、底珊族長。這是從何利人所出的族長，都在西珥地，按著族長來分。 \end{tabularx} \\ \\ \relax
36:31 & \begin{tabularx}{0.7\textwidth}{X} 以色列未有君王治理之前，這些是在以東地作王的。 \end{tabularx} \\ \\ \relax
36:32 & \begin{tabularx}{0.7\textwidth}{X} 比珥的兒子比拉在以東作王，他的城名叫亭哈巴。 \end{tabularx} \\ \\ \relax
36:33 & \begin{tabularx}{0.7\textwidth}{X} 比拉死了，波斯拉人謝拉的兒子約巴接續他作王。 \end{tabularx} \\ \\ \relax
36:34 & \begin{tabularx}{0.7\textwidth}{X} 約巴死了，提幔人之地的戶珊接續他作王。 \end{tabularx} \\ \\ \relax
36:35 & \begin{tabularx}{0.7\textwidth}{X} 戶珊死了，比達的兒子哈達接續他作王，哈達曾在摩押地擊敗米甸人，他的城名叫亞未得。 \end{tabularx} \\ \\ \relax
36:36 & \begin{tabularx}{0.7\textwidth}{X} 哈達死了，瑪士利加人桑拉接續他作王。 \end{tabularx} \\ \\ \relax
36:37 & \begin{tabularx}{0.7\textwidth}{X} 桑拉死了，大河邊的利河伯人掃羅接續他作王。 \end{tabularx} \\ \\ \relax
36:38 & \begin{tabularx}{0.7\textwidth}{X} 掃羅死了，亞革波的兒子巴勒‧哈南接續他作王。 \end{tabularx} \\ \\ \relax
36:39 & \begin{tabularx}{0.7\textwidth}{X} 亞革波的兒子巴勒‧哈南死了，哈達爾接續他作王，他的城名叫巴烏。他的妻子名叫米希她別，是米‧薩合的孫女，瑪特列的女兒。 \end{tabularx} \\ \\ \relax
36:40 & \begin{tabularx}{0.7\textwidth}{X} 這是以掃的族長，按著他們的宗族、住處和名字：亭納族長、亞勒瓦族長、耶帖族長、 \end{tabularx} \\ \\ \relax
36:41 & \begin{tabularx}{0.7\textwidth}{X} 阿何利巴瑪族長、以拉族長、比嫩族長、 \end{tabularx} \\ \\ \relax
36:42 & \begin{tabularx}{0.7\textwidth}{X} 基納斯族長、提幔族長、米比薩族長、 \end{tabularx} \\ \\ \relax
36:43 & \begin{tabularx}{0.7\textwidth}{X} 瑪基疊族長、以蘭族長。這些是以東人在所得為業的地上，按著他們住處的族長。以掃是以東人的始祖。 \end{tabularx} \\ \\ \relax
37:1 & \begin{tabularx}{0.7\textwidth}{X} 雅各住在迦南地，就是他父親寄居的地。 \end{tabularx} \\ \\ \relax
37:2 & \begin{tabularx}{0.7\textwidth}{X} 這是雅各的事蹟。約瑟十七歲與他哥哥們一同牧羊。他是個少年，與他父親的妾辟拉和悉帕的兒子們常在一起。約瑟把他們的惡行報給父親。 \end{tabularx} \\ \\ \relax
37:3 & \begin{tabularx}{0.7\textwidth}{X} 以色列愛約瑟過於其他的兒子，因為約瑟是他年老生的；他給約瑟做了一件長袍。 \end{tabularx} \\ \\ \relax
37:4 & \begin{tabularx}{0.7\textwidth}{X} 哥哥們見父親愛約瑟過於他們，就恨約瑟，不與他說友善的話。 \end{tabularx} \\ \\ \relax
37:5 & \begin{tabularx}{0.7\textwidth}{X} 約瑟做了一個夢，告訴他哥哥們，他們就更加恨他。 \end{tabularx} \\ \\ \relax
37:6 & \begin{tabularx}{0.7\textwidth}{X} 約瑟對他們說：「請聽我做的這個夢： \end{tabularx} \\ \\ \relax
37:7 & \begin{tabularx}{0.7\textwidth}{X} 看哪，我們在田裡捆禾稼；看哪，我的捆起來站著；看哪，你們的捆圍著我的捆下拜。」 \end{tabularx} \\ \\ \relax
37:8 & \begin{tabularx}{0.7\textwidth}{X} 他的哥哥們對他說：「難道你真的要作我們的王嗎？難道你真的要統治我們嗎？」他們就因他的夢和他的話更加恨他。 \end{tabularx} \\ \\ \relax
37:9 & \begin{tabularx}{0.7\textwidth}{X} 後來他又做了另一個夢，告訴他哥哥們說：「看哪，我又做了一個夢；看哪，太陽、月亮和十一顆星都向我下拜。」 \end{tabularx} \\ \\ \relax
37:10 & \begin{tabularx}{0.7\textwidth}{X} 約瑟告訴他父親和哥哥們，他父親就責備他說：「你做的這是甚麼夢！難道我和你母親、你的兄弟真的要俯伏在地，來向你下拜嗎？」 \end{tabularx} \\ \\ \relax
37:11 & \begin{tabularx}{0.7\textwidth}{X} 他的哥哥們都嫉妒他，他父親卻把這事存在心裡。 \end{tabularx} \\ \\ \relax
37:12 & \begin{tabularx}{0.7\textwidth}{X} 約瑟的哥哥們到示劍去放他們父親的羊。 \end{tabularx} \\ \\ \relax
37:13 & \begin{tabularx}{0.7\textwidth}{X} 以色列對約瑟說：「你哥哥們不是在示劍放羊嗎？來，我派你到他們那裡去。」約瑟對他說：「我在這裡。」 \end{tabularx} \\ \\ \relax
37:14 & \begin{tabularx}{0.7\textwidth}{X} 以色列對他說：「你去看看你哥哥們是否平安，羊群是否平安，再回來告訴我。」於是他派約瑟出希伯崙谷，約瑟就往示劍去了。 \end{tabularx} \\ \\ \relax
37:15 & \begin{tabularx}{0.7\textwidth}{X} 有人遇見他，看哪，他在田野走迷了路。那人問他說：「你找甚麼？」 \end{tabularx} \\ \\ \relax
37:16 & \begin{tabularx}{0.7\textwidth}{X} 他說：「我找我的哥哥們，請告訴我，他們在哪裡放羊。」 \end{tabularx} \\ \\ \relax
37:17 & \begin{tabularx}{0.7\textwidth}{X} 那人說：「他們已經離開這裡走了，我聽見他們說：『我們往多坍去。』」約瑟就去追哥哥們，在多坍找到了他們。 \end{tabularx} \\ \\ \relax
37:18 & \begin{tabularx}{0.7\textwidth}{X} 他們遠遠看見他，趁他還沒有走近他們，就圖謀要殺死他。 \end{tabularx} \\ \\ \relax
37:19 & \begin{tabularx}{0.7\textwidth}{X} 他們彼此說：「看哪！那做夢的來了。 \end{tabularx} \\ \\ \relax
37:20 & \begin{tabularx}{0.7\textwidth}{X} 現在，來吧！我們把他殺了，丟在一個坑裡，就說有惡獸把他吃了。我們且看他的夢將來怎麼樣。」 \end{tabularx} \\ \\ \relax
37:21 & \begin{tabularx}{0.7\textwidth}{X} 呂便聽見了，要救約瑟脫離他們的手，說：「我們不可害他的性命。」 \end{tabularx} \\ \\ \relax
37:22 & \begin{tabularx}{0.7\textwidth}{X} 呂便又對他們說：「不可流他的血，可以把他丟在這曠野的坑裡，不可下手害他。」呂便要救他脫離他們的手，把他還給他父親。 \end{tabularx} \\ \\ \relax
37:23 & \begin{tabularx}{0.7\textwidth}{X} 約瑟到了他哥哥們那裡，他們就剝去他的外衣，就是他身上那件長袍。 \end{tabularx} \\ \\ \relax
37:24 & \begin{tabularx}{0.7\textwidth}{X} 他們抓住他，把他丟在坑裡。那坑是空的，裡頭沒有水。 \end{tabularx} \\ \\ \relax
37:25 & \begin{tabularx}{0.7\textwidth}{X} 他們坐下吃飯，舉目觀看，看哪，有一群以實瑪利人從基列來，用駱駝馱著香料、乳香、沒藥，要帶下埃及去。 \end{tabularx} \\ \\ \relax
37:26 & \begin{tabularx}{0.7\textwidth}{X} 猶大對他的兄弟們說：「我們殺我們的弟弟，遮掩他的血有甚麼好處呢？ \end{tabularx} \\ \\ \relax
37:27 & \begin{tabularx}{0.7\textwidth}{X} 來，我們把他賣給以實瑪利人，不要下手害他，因為他是我們的弟弟，我們的骨肉。」他的兄弟們就聽從了他。 \end{tabularx} \\ \\ \relax
37:28 & \begin{tabularx}{0.7\textwidth}{X} 那時，有些米甸的商人從那裡經過，就把約瑟從坑裡拉上來。他們以二十塊銀子把約瑟賣給以實瑪利人，他們就把約瑟帶到埃及去了。 \end{tabularx} \\ \\ \relax
37:29 & \begin{tabularx}{0.7\textwidth}{X} 呂便回到坑旁，看哪，約瑟不在坑裡，就撕裂自己的衣服， \end{tabularx} \\ \\ \relax
37:30 & \begin{tabularx}{0.7\textwidth}{X} 回到他兄弟們那裡，說：「孩子不在了。我往哪裡去才好呢？」 \end{tabularx} \\ \\ \relax
37:31 & \begin{tabularx}{0.7\textwidth}{X} 於是，他們宰了一隻公山羊，拿了約瑟的那件外衣染上了血， \end{tabularx} \\ \\ \relax
37:32 & \begin{tabularx}{0.7\textwidth}{X} 派人把長袍送到他們的父親那裡，說：「我們發現這個， 請認一認，是不是你兒子的外衣？」 \end{tabularx} \\ \\ \relax
37:33 & \begin{tabularx}{0.7\textwidth}{X} 他認出來，就說：「這是我兒子的外衣，惡獸把他吃了，約瑟一定被撕碎了！」 \end{tabularx} \\ \\ \relax
37:34 & \begin{tabularx}{0.7\textwidth}{X} 雅各就撕裂衣服，腰間圍上麻布，為他兒子哀傷了多日。 \end{tabularx} \\ \\ \relax
37:35 & \begin{tabularx}{0.7\textwidth}{X} 他的兒女都起來安慰他，他卻不肯受安慰，說：「我必哀傷著下陰間，到我兒子那裡。」約瑟的父親就為他哀哭。 \end{tabularx} \\ \\ \relax
37:36 & \begin{tabularx}{0.7\textwidth}{X} 米甸人把約瑟賣到埃及，給法老的官員，就是護衛長波提乏。 \end{tabularx} \\ \\ \relax
38:1 & \begin{tabularx}{0.7\textwidth}{X} 那時，猶大離開他兄弟們下去，到一個名叫希拉的亞杜蘭人的家附近支搭帳棚。 \end{tabularx} \\ \\ \relax
38:2 & \begin{tabularx}{0.7\textwidth}{X} 猶大在那裡看見一個名叫拔‧書亞的迦南女子，就娶她為妻，與她同房， \end{tabularx} \\ \\ \relax
38:3 & \begin{tabularx}{0.7\textwidth}{X} 她就懷孕生了兒子，猶大給他起名叫珥。 \end{tabularx} \\ \\ \relax
38:4 & \begin{tabularx}{0.7\textwidth}{X} 她又懷孕生了兒子，給他起名叫俄南。 \end{tabularx} \\ \\ \relax
38:5 & \begin{tabularx}{0.7\textwidth}{X} 她又再生了兒子，給他起名叫示拉。她生示拉的時候，猶大正在基悉。 \end{tabularx} \\ \\ \relax
38:6 & \begin{tabularx}{0.7\textwidth}{X} 猶大為長子珥娶妻，名叫她瑪。 \end{tabularx} \\ \\ \relax
38:7 & \begin{tabularx}{0.7\textwidth}{X} 猶大的長子珥在耶和華眼中看為惡，耶和華就殺死了他。 \end{tabularx} \\ \\ \relax
38:8 & \begin{tabularx}{0.7\textwidth}{X} 猶大對俄南說：「你當與你哥哥的妻子同房，向她盡你的本分，為你哥哥生子立後。」 \end{tabularx} \\ \\ \relax
38:9 & \begin{tabularx}{0.7\textwidth}{X} 俄南知道如果與嫂嫂同房，所生的孩子不屬於自己，就洩在地上，不為哥哥生子立後。 \end{tabularx} \\ \\ \relax
38:10 & \begin{tabularx}{0.7\textwidth}{X} 俄南所做的在耶和華眼中看為惡，耶和華也殺死了他。 \end{tabularx} \\ \\ \relax
38:11 & \begin{tabularx}{0.7\textwidth}{X} 猶大對他媳婦她瑪說：「你去住在你父親家裡守寡，等我兒子示拉長大。」因為他說：「恐怕示拉也像兩個哥哥一樣死去。」她瑪就去，住在她父親家裡。 \end{tabularx} \\ \\ \relax
38:12 & \begin{tabularx}{0.7\textwidth}{X} 過了一段很長的日子，猶大的妻子，書亞的女兒死了。猶大受到了安慰，就和他朋友亞杜蘭人希拉上亭拿去，到他的剪羊毛的人那裡。 \end{tabularx} \\ \\ \relax
38:13 & \begin{tabularx}{0.7\textwidth}{X} 有人告訴她瑪說：「看哪，你的公公上亭拿剪羊毛去了。」 \end{tabularx} \\ \\ \relax
38:14 & \begin{tabularx}{0.7\textwidth}{X} 她瑪見示拉已經長大，卻還沒有娶她為妻，就脫去她寡婦的衣裳，用面紗蒙著，蓋住自己，坐在往亭拿的路上，伊拿印城門口。 \end{tabularx} \\ \\ \relax
38:15 & \begin{tabularx}{0.7\textwidth}{X} 猶大看見她，以為是妓女，因為她蒙著臉。 \end{tabularx} \\ \\ \relax
38:16 & \begin{tabularx}{0.7\textwidth}{X} 猶大就轉到路邊她那裡，說：「來吧！讓我與你同寢。」他並不知道她就是他的媳婦。她瑪說：「你要與我同寢，把甚麼給我呢？」 \end{tabularx} \\ \\ \relax
38:17 & \begin{tabularx}{0.7\textwidth}{X} 猶大說：「我從羊群裡取一隻小山羊，派人送來給你。」她瑪說：「在未送之前，你能給我一個信物嗎？」 \end{tabularx} \\ \\ \relax
38:18 & \begin{tabularx}{0.7\textwidth}{X} 他說：「我給你甚麼信物呢？」她瑪說：「你的印、你的帶子和你手裡的杖。」於是猶大給了她，與她同寢，她就從猶大懷了孕。 \end{tabularx} \\ \\ \relax
38:19 & \begin{tabularx}{0.7\textwidth}{X} 她瑪起來走了，除去面紗，照常穿上寡婦的衣裳。 \end{tabularx} \\ \\ \relax
38:20 & \begin{tabularx}{0.7\textwidth}{X} 猶大託他朋友亞杜蘭人送一隻小山羊去，要從那女人手裡取回信物，卻找不到她。 \end{tabularx} \\ \\ \relax
38:21 & \begin{tabularx}{0.7\textwidth}{X} 他問那地方的人說：「伊拿印路旁的神廟娼妓在哪裡？」他們說：「這裡沒有神廟娼妓。」 \end{tabularx} \\ \\ \relax
38:22 & \begin{tabularx}{0.7\textwidth}{X} 他回到猶大那裡說：「我找不到她，並且那地方的人說：『這裡沒有神廟娼妓。』」 \end{tabularx} \\ \\ \relax
38:23 & \begin{tabularx}{0.7\textwidth}{X} 猶大說：「讓她拿去吧，免得我們被人譏笑。看哪，我把這小山羊送去了，可是你找不到她。」 \end{tabularx} \\ \\ \relax
38:24 & \begin{tabularx}{0.7\textwidth}{X} 大約過了三個月，有人告訴猶大說：「你的媳婦她瑪行淫，並且，看哪，她因行淫而懷了孕。」猶大說：「拉她出來，把她燒了！」 \end{tabularx} \\ \\ \relax
38:25 & \begin{tabularx}{0.7\textwidth}{X} 她瑪被拉出來的時候，就派人到她公公那裡，對他說：「這些東西是誰的，我就是從誰懷了孕。」她又說：「請你認一認，這印、這帶子和這杖是誰的？」 \end{tabularx} \\ \\ \relax
38:26 & \begin{tabularx}{0.7\textwidth}{X} 猶大承認說：「她比我更有理，因為我沒有把她給我的兒子示拉。」猶大再也不跟她同寢。 \end{tabularx} \\ \\ \relax
38:27 & \begin{tabularx}{0.7\textwidth}{X} 她瑪生產的時候到了，看哪，腹裡懷的是雙胞胎。 \end{tabularx} \\ \\ \relax
38:28 & \begin{tabularx}{0.7\textwidth}{X} 生產的時候，一個孩子伸出手來；接生婆拿紅線綁在他手上，說：「這是頭生的。」 \end{tabularx} \\ \\ \relax
38:29 & \begin{tabularx}{0.7\textwidth}{X} 這孩子把手收回去，看哪，他哥哥生出來了；接生婆說：「你竟然為自己衝出一個裂縫！」於是，他的名字叫法勒斯。 \end{tabularx} \\ \\ \relax
38:30 & \begin{tabularx}{0.7\textwidth}{X} 後來，那手上有紅線的兄弟也生出來，他的名字叫謝拉。 \end{tabularx} \\ \\ \relax
39:1 & \begin{tabularx}{0.7\textwidth}{X} 約瑟被帶下埃及去。有一個埃及人波提乏，是法老的官員，是護衛長，他從那些帶約瑟下來的以實瑪利人手中把約瑟買了去。 \end{tabularx} \\ \\ \relax
39:2 & \begin{tabularx}{0.7\textwidth}{X} 約瑟在他埃及主人的家中，耶和華與他同在，他是一個通達的人。 \end{tabularx} \\ \\ \relax
39:3 & \begin{tabularx}{0.7\textwidth}{X} 他主人見耶和華與他同在，又見耶和華使他手裡所辦的事都順利， \end{tabularx} \\ \\ \relax
39:4 & \begin{tabularx}{0.7\textwidth}{X} 約瑟就在主人眼前蒙恩，伺候他主人，主人派他管理家務，把一切所有的都交在他手裡。 \end{tabularx} \\ \\ \relax
39:5 & \begin{tabularx}{0.7\textwidth}{X} 自從主人派約瑟管理家務和他一切所有的，耶和華就因約瑟的緣故賜福給那埃及人的家；凡家裡和田間一切所有的，都蒙耶和華賜福。 \end{tabularx} \\ \\ \relax
39:6 & \begin{tabularx}{0.7\textwidth}{X} 波提乏把他一切所有的都交在約瑟手中，除了自己所吃的食物，其他的事一概不知。約瑟英俊健美。 \end{tabularx} \\ \\ \relax
39:7 & \begin{tabularx}{0.7\textwidth}{X} 這些事以後，約瑟主人的妻子以目送情給約瑟，說：「你與我同寢吧！」 \end{tabularx} \\ \\ \relax
39:8 & \begin{tabularx}{0.7\textwidth}{X} 約瑟拒絕，對他主人的妻子說：「看哪，一切家務我主人一概不知，他把所有的都交在我手裡。 \end{tabularx} \\ \\ \relax
39:9 & \begin{tabularx}{0.7\textwidth}{X} 在這家裡沒有人比我更大，除你以外，他也沒有留下一樣不交給我，因為你是他的妻子。我怎能行這麼大的惡，得罪神呢？」 \end{tabularx} \\ \\ \relax
39:10 & \begin{tabularx}{0.7\textwidth}{X} 她天天這樣對約瑟說，約瑟卻不聽從她，不與她同寢，也不和她在一起。 \end{tabularx} \\ \\ \relax
39:11 & \begin{tabularx}{0.7\textwidth}{X} 有一天，約瑟進屋裡去辦事，家裡沒有一個人在那屋子裡， \end{tabularx} \\ \\ \relax
39:12 & \begin{tabularx}{0.7\textwidth}{X} 婦人就拉住他的衣服，說：「你與我同寢吧！」約瑟把衣服留在她手裡，逃出外面去了。 \end{tabularx} \\ \\ \relax
39:13 & \begin{tabularx}{0.7\textwidth}{X} 婦人看見約瑟把衣服留在她手裡逃到外面， \end{tabularx} \\ \\ \relax
39:14 & \begin{tabularx}{0.7\textwidth}{X} 就叫了家裡的人來，對他們說：「看，他帶了一個希伯來人到我們這裡戲弄我們。他到我這裡來，要與我同寢，我就大聲喊叫。 \end{tabularx} \\ \\ \relax
39:15 & \begin{tabularx}{0.7\textwidth}{X} 他聽見我放聲大喊，就把他的衣服留在我這裡，逃出外面去了。」 \end{tabularx} \\ \\ \relax
39:16 & \begin{tabularx}{0.7\textwidth}{X} 婦人把約瑟的衣服放在身邊，直到他主人回家， \end{tabularx} \\ \\ \relax
39:17 & \begin{tabularx}{0.7\textwidth}{X} 就用這樣的話對他說：「你帶到我們這裡來的那希伯來僕人進來要調戲我， \end{tabularx} \\ \\ \relax
39:18 & \begin{tabularx}{0.7\textwidth}{X} 我放聲大喊，他就把衣服留在我身邊，逃到外面。」 \end{tabularx} \\ \\ \relax
39:19 & \begin{tabularx}{0.7\textwidth}{X} 主人聽見他妻子對他說的話，說：「你的僕人就是這樣對待我」，就非常生氣。 \end{tabularx} \\ \\ \relax
39:20 & \begin{tabularx}{0.7\textwidth}{X} 約瑟的主人把他抓起來，關在監獄裡，就是王的囚犯被關的地方。於是約瑟在那裡坐牢。 \end{tabularx} \\ \\ \relax
39:21 & \begin{tabularx}{0.7\textwidth}{X} 但耶和華與約瑟同在，向他施恩，使他在監獄長的眼前蒙恩。 \end{tabularx} \\ \\ \relax
39:22 & \begin{tabularx}{0.7\textwidth}{X} 監獄長就把監獄裡所有的囚犯都交在約瑟手下；在那裡的一切事都由他處理。 \end{tabularx} \\ \\ \relax
39:23 & \begin{tabularx}{0.7\textwidth}{X} 任何交在約瑟手中的事，監獄長一概不察，因為耶和華與約瑟同在，耶和華使他所做的都順利。 \end{tabularx} \\ \\ \relax
40:1 & \begin{tabularx}{0.7\textwidth}{X} 這些事以後，埃及王的司酒長和司膳長得罪了他們的主埃及王。 \end{tabularx} \\ \\ \relax
40:2 & \begin{tabularx}{0.7\textwidth}{X} 法老就對司酒長和司膳長兩個官員發怒， \end{tabularx} \\ \\ \relax
40:3 & \begin{tabularx}{0.7\textwidth}{X} 把他們關在護衛長府內的監獄裡，就是約瑟被囚的地方。 \end{tabularx} \\ \\ \relax
40:4 & \begin{tabularx}{0.7\textwidth}{X} 護衛長把他們交給約瑟，約瑟就伺候他們。他們被關了一段日子。 \end{tabularx} \\ \\ \relax
40:5 & \begin{tabularx}{0.7\textwidth}{X} 關在監獄裡的這兩個人，就是埃及王的司酒長和司膳長，在同一個晚上各自做了一個夢，每個夢都有自己的解釋。 \end{tabularx} \\ \\ \relax
40:6 & \begin{tabularx}{0.7\textwidth}{X} 到了早晨，約瑟來到他們那裡看他們，看哪，他們很憂愁。 \end{tabularx} \\ \\ \relax
40:7 & \begin{tabularx}{0.7\textwidth}{X} 他就問一同關在他主人府內法老的官員，說：「你們今日為甚麼面帶愁容呢？」 \end{tabularx} \\ \\ \relax
40:8 & \begin{tabularx}{0.7\textwidth}{X} 他們對他說：「我們各自做了一個夢，卻沒有人能講解。」約瑟對他們說：「解夢不是出於神嗎？請你們把夢告訴我。」 \end{tabularx} \\ \\ \relax
40:9 & \begin{tabularx}{0.7\textwidth}{X} 司酒長就把夢告訴約瑟，對他說：「在我的夢中，看哪，有一棵葡萄樹在我面前， \end{tabularx} \\ \\ \relax
40:10 & \begin{tabularx}{0.7\textwidth}{X} 樹上有三根枝子。枝子發了芽，開了花，結出串串成熟的葡萄。 \end{tabularx} \\ \\ \relax
40:11 & \begin{tabularx}{0.7\textwidth}{X} 法老的杯在我手中，我就拿葡萄擠在法老的杯裡，把杯遞到他手中。」 \end{tabularx} \\ \\ \relax
40:12 & \begin{tabularx}{0.7\textwidth}{X} 約瑟對他說：「夢的解釋是這樣：三根枝子就是三天； \end{tabularx} \\ \\ \relax
40:13 & \begin{tabularx}{0.7\textwidth}{X} 三天之內，法老要讓你抬起頭來，叫你官復原職。你仍要遞杯在法老的手中，像先前作他的司酒長一樣。 \end{tabularx} \\ \\ \relax
40:14 & \begin{tabularx}{0.7\textwidth}{X} 但你得福的時候，請你記得我，向我施慈愛，在法老面前提起我，救我出這監牢。 \end{tabularx} \\ \\ \relax
40:15 & \begin{tabularx}{0.7\textwidth}{X} 我實在是從希伯來人之地被拐來的，我在這裡也沒有做過甚麼，好叫他們把我關在牢裡。」 \end{tabularx} \\ \\ \relax
40:16 & \begin{tabularx}{0.7\textwidth}{X} 司膳長見夢解得好，就對約瑟說：「在我夢中，看哪，我頭上頂著三個裝餅的籃子； \end{tabularx} \\ \\ \relax
40:17 & \begin{tabularx}{0.7\textwidth}{X} 最上面的籃子裡有為法老烤的各樣食物，有飛鳥來吃我頭上籃子裡的食物。」 \end{tabularx} \\ \\ \relax
40:18 & \begin{tabularx}{0.7\textwidth}{X} 約瑟說：「夢的解釋是這樣：三個籃子就是三天； \end{tabularx} \\ \\ \relax
40:19 & \begin{tabularx}{0.7\textwidth}{X} 三天之內，法老要讓你抬起頭來，身首異處，把你掛在木架上，必有飛鳥來吃你身上的肉。」 \end{tabularx} \\ \\ \relax
40:20 & \begin{tabularx}{0.7\textwidth}{X} 到了第三天，正是法老的生日，他為眾臣僕擺設宴席，使司酒長和司膳長從眾臣僕中抬起頭來， \end{tabularx} \\ \\ \relax
40:21 & \begin{tabularx}{0.7\textwidth}{X} 讓司酒長官復原職，仍舊遞杯在法老手中， \end{tabularx} \\ \\ \relax
40:22 & \begin{tabularx}{0.7\textwidth}{X} 卻把司膳長掛起來，正如約瑟向他們所講解的。 \end{tabularx} \\ \\ \relax
40:23 & \begin{tabularx}{0.7\textwidth}{X} 然而，司酒長不記得約瑟，竟忘了他。 \end{tabularx} \\ \\ \relax
41:1 & \begin{tabularx}{0.7\textwidth}{X} 過了兩年，法老做夢，看哪，自己站在尼羅河邊， \end{tabularx} \\ \\ \relax
41:2 & \begin{tabularx}{0.7\textwidth}{X} 看哪，有七頭母牛從尼羅河裡上來，長相俊美，肌肉肥壯，在蘆葦中吃草。 \end{tabularx} \\ \\ \relax
41:3 & \begin{tabularx}{0.7\textwidth}{X} 看哪，隨後又有七頭母牛從尼羅河裡上來，長相醜陋，肌肉乾瘦，與那七頭母牛一同站在河邊。 \end{tabularx} \\ \\ \relax
41:4 & \begin{tabularx}{0.7\textwidth}{X} 這長相醜陋，肌肉乾瘦的七頭母牛吃了那長相俊美又肥壯的七頭母牛。法老就醒了。 \end{tabularx} \\ \\ \relax
41:5 & \begin{tabularx}{0.7\textwidth}{X} 他又睡著，第二次做夢，看哪，一株麥桿長了七個穗子，又肥大又佳美， \end{tabularx} \\ \\ \relax
41:6 & \begin{tabularx}{0.7\textwidth}{X} 看哪，隨後又長出七個穗子，又細弱又被東風吹焦了。 \end{tabularx} \\ \\ \relax
41:7 & \begin{tabularx}{0.7\textwidth}{X} 這細弱的穗子吞了那七個又肥大又飽滿的穗子。法老醒了，看哪，是個夢。 \end{tabularx} \\ \\ \relax
41:8 & \begin{tabularx}{0.7\textwidth}{X} 到了早晨，法老心裡不安，就派人把埃及所有的術士和智慧人都召來。法老把所做的夢告訴他們，但是沒有人能為法老解夢。 \end{tabularx} \\ \\ \relax
41:9 & \begin{tabularx}{0.7\textwidth}{X} 那時司酒長對法老說：「我今日想起我的罪來。 \end{tabularx} \\ \\ \relax
41:10 & \begin{tabularx}{0.7\textwidth}{X} 從前法老對臣僕發怒，把我和司膳長關在護衛長府內的監牢裡。 \end{tabularx} \\ \\ \relax
41:11 & \begin{tabularx}{0.7\textwidth}{X} 我們兩人在同一晚上各做一夢，每個夢都有各自的解釋。 \end{tabularx} \\ \\ \relax
41:12 & \begin{tabularx}{0.7\textwidth}{X} 同我們在一起有一個希伯來的年輕人，是護衛長的僕人。我們告訴他，他就為我們解夢，照著各人的夢講解。 \end{tabularx} \\ \\ \relax
41:13 & \begin{tabularx}{0.7\textwidth}{X} 後來事情正如他給我們講解的實現了，我官復原職，司膳長被掛起來了。」 \end{tabularx} \\ \\ \relax
41:14 & \begin{tabularx}{0.7\textwidth}{X} 於是法老派人去召約瑟，他們就急忙把他從牢裡提出來。他就剃頭刮臉，換衣服，進到法老面前。 \end{tabularx} \\ \\ \relax
41:15 & \begin{tabularx}{0.7\textwidth}{X} 法老對約瑟說：「我做了一個夢，沒有人能講解。我聽人說，你聽了夢就能講解。」 \end{tabularx} \\ \\ \relax
41:16 & \begin{tabularx}{0.7\textwidth}{X} 約瑟回答法老說：「這不在乎我。神必應允法老平安。」 \end{tabularx} \\ \\ \relax
41:17 & \begin{tabularx}{0.7\textwidth}{X} 法老對約瑟說：「在我的夢中，看哪，我站在尼羅河邊， \end{tabularx} \\ \\ \relax
41:18 & \begin{tabularx}{0.7\textwidth}{X} 看哪，有七頭母牛從尼羅河裡上來，肌肉肥壯，外形俊美，在蘆葦中吃草。 \end{tabularx} \\ \\ \relax
41:19 & \begin{tabularx}{0.7\textwidth}{X} 看哪，隨後又有七頭母牛上來，虛弱，外形很醜陋，肌肉又乾瘦，在埃及全地，我沒有見過這樣醜陋的牛。 \end{tabularx} \\ \\ \relax
41:20 & \begin{tabularx}{0.7\textwidth}{X} 這乾瘦又醜陋的母牛吃了那先前的七頭肥母牛， \end{tabularx} \\ \\ \relax
41:21 & \begin{tabularx}{0.7\textwidth}{X} 進了肚子以後卻看不出已經進了肚子，那醜陋的長相仍舊和先前一樣。我就醒了。 \end{tabularx} \\ \\ \relax
41:22 & \begin{tabularx}{0.7\textwidth}{X} 我又在夢中觀看，看哪，一株麥桿長了七個穗子，又飽滿又佳美， \end{tabularx} \\ \\ \relax
41:23 & \begin{tabularx}{0.7\textwidth}{X} 看哪，隨後又長出七個穗子，枯槁，細弱，又被東風吹焦了。 \end{tabularx} \\ \\ \relax
41:24 & \begin{tabularx}{0.7\textwidth}{X} 這些細弱的穗子吞了那七個佳美的穗子。我告訴術士，卻沒有人能為我講解。」 \end{tabularx} \\ \\ \relax
41:25 & \begin{tabularx}{0.7\textwidth}{X} 約瑟對法老說：「法老的夢是同一個。神已把要做的事指示法老了。 \end{tabularx} \\ \\ \relax
41:26 & \begin{tabularx}{0.7\textwidth}{X} 七頭好母牛是七年，七個佳美的穗子也是七年，這是同一個夢。 \end{tabularx} \\ \\ \relax
41:27 & \begin{tabularx}{0.7\textwidth}{X} 那隨後上來的七頭乾瘦又醜陋的母牛是七年；那七個空心，被東風吹焦的穗子也一樣，都是七個荒年。 \end{tabularx} \\ \\ \relax
41:28 & \begin{tabularx}{0.7\textwidth}{X} 這就是我對法老所說，神已把要做的事顯明給法老了。 \end{tabularx} \\ \\ \relax
41:29 & \begin{tabularx}{0.7\textwidth}{X} 看哪，必有七個大豐年來到埃及全地， \end{tabularx} \\ \\ \relax
41:30 & \begin{tabularx}{0.7\textwidth}{X} 隨後又有七個荒年，甚至埃及地的人都忘了先前的豐收，這地必被饑荒所滅。 \end{tabularx} \\ \\ \relax
41:31 & \begin{tabularx}{0.7\textwidth}{X} 因為那後來的饑荒非常嚴重，就不覺得這地先前有豐收。 \end{tabularx} \\ \\ \relax
41:32 & \begin{tabularx}{0.7\textwidth}{X} 至於法老兩次做夢，是因為神已經確定這事，神必速速成就。 \end{tabularx} \\ \\ \relax
41:33 & \begin{tabularx}{0.7\textwidth}{X} 現在，請法老選一個聰明又有智慧的人，委派他治理埃及地。 \end{tabularx} \\ \\ \relax
41:34 & \begin{tabularx}{0.7\textwidth}{X} 請法老這樣做，委派官員治理這地，在七個豐年的期間，徵收埃及地出產的五分之一， \end{tabularx} \\ \\ \relax
41:35 & \begin{tabularx}{0.7\textwidth}{X} 叫他們聚集未來豐年一切的糧食，積存五穀歸在法老的手下作糧食，儲藏在各城裡。 \end{tabularx} \\ \\ \relax
41:36 & \begin{tabularx}{0.7\textwidth}{X} 這糧食可以為這地作儲備，為了埃及地要來的七個荒年，免得這地被饑荒所滅。」 \end{tabularx} \\ \\ \relax
41:37 & \begin{tabularx}{0.7\textwidth}{X} 這事在法老和他眾臣僕眼中都覺得好。 \end{tabularx} \\ \\ \relax
41:38 & \begin{tabularx}{0.7\textwidth}{X} 法老對臣僕說：「像這樣的人，有神的靈在他裡面，我們豈能找得著呢？」 \end{tabularx} \\ \\ \relax
41:39 & \begin{tabularx}{0.7\textwidth}{X} 法老對約瑟說：「神既指示你這一切事，就沒有人像你這樣聰明又有智慧。 \end{tabularx} \\ \\ \relax
41:40 & \begin{tabularx}{0.7\textwidth}{X} 你可以治理我的家；我的百姓都必服從你口中的命令。惟獨在寶座上，我比你大。」 \end{tabularx} \\ \\ \relax
41:41 & \begin{tabularx}{0.7\textwidth}{X} 法老又對約瑟說：「看，我委派你治理埃及全地。」 \end{tabularx} \\ \\ \relax
41:42 & \begin{tabularx}{0.7\textwidth}{X} 法老就脫下手上帶印的戒指，戴在約瑟的手上，給他穿上細麻衣，把金鏈戴在他的頸項上， \end{tabularx} \\ \\ \relax
41:43 & \begin{tabularx}{0.7\textwidth}{X} 又給約瑟坐他的副座車，在他前面有人呼叫說：「跪下。」於是，法老委派他治理埃及全地。 \end{tabularx} \\ \\ \relax
41:44 & \begin{tabularx}{0.7\textwidth}{X} 法老對約瑟說：「我是法老，若沒有你的命令，埃及全地的人都不可擅自辦事。」 \end{tabularx} \\ \\ \relax
41:45 & \begin{tabularx}{0.7\textwidth}{X} 法老給約瑟起名叫撒發那特‧巴內亞，又將安城的祭司波提‧非拉的女兒亞西納給他為妻。約瑟就出去治理埃及地。 \end{tabularx} \\ \\ \relax
41:46 & \begin{tabularx}{0.7\textwidth}{X} 約瑟在埃及王法老面前侍立的時候年三十歲。約瑟從法老面前出去，巡行埃及全地。 \end{tabularx} \\ \\ \relax
41:47 & \begin{tabularx}{0.7\textwidth}{X} 七個豐年之內，地的出產極其豐盛， \end{tabularx} \\ \\ \relax
41:48 & \begin{tabularx}{0.7\textwidth}{X} 約瑟聚集埃及地七年一切的糧食，把糧食積存在各城裡，就是把各城周圍田地的糧食都積存在該城裡。 \end{tabularx} \\ \\ \relax
41:49 & \begin{tabularx}{0.7\textwidth}{X} 約瑟積存的五穀很多，如同海邊的沙，無法計算，數也數不清。 \end{tabularx} \\ \\ \relax
41:50 & \begin{tabularx}{0.7\textwidth}{X} 荒年未到以前，安城的祭司波提‧非拉的女兒亞西納為約瑟生了兩個兒子。 \end{tabularx} \\ \\ \relax
41:51 & \begin{tabularx}{0.7\textwidth}{X} 約瑟給長子起名叫瑪拿西，因為他說：「神使我忘了一切的困苦和我父的全家。」 \end{tabularx} \\ \\ \relax
41:52 & \begin{tabularx}{0.7\textwidth}{X} 他給次子起名叫以法蓮，因為他說：「神使我在受苦的地方興盛。」 \end{tabularx} \\ \\ \relax
41:53 & \begin{tabularx}{0.7\textwidth}{X} 埃及地的七個豐年一過， \end{tabularx} \\ \\ \relax
41:54 & \begin{tabularx}{0.7\textwidth}{X} 七個荒年就來了，正如約瑟所說的。各地都有饑荒，惟獨埃及全地有糧食。 \end{tabularx} \\ \\ \relax
41:55 & \begin{tabularx}{0.7\textwidth}{X} 等到埃及全地也有了饑荒，眾百姓就向法老哀求糧食。法老對所有的埃及人說：「你們到約瑟那裡去，凡他所說的，你們都要做。」 \end{tabularx} \\ \\ \relax
41:56 & \begin{tabularx}{0.7\textwidth}{X} 當時饑荒遍滿了全地，約瑟就開了各處的糧倉，賣糧食給埃及人。埃及地的饑荒非常嚴重。 \end{tabularx} \\ \\ \relax
41:57 & \begin{tabularx}{0.7\textwidth}{X} 各地的人都去埃及，到約瑟那裡買糧食，因為全地的饑荒非常嚴重。 \end{tabularx} \\ \\ \relax
42:1 & \begin{tabularx}{0.7\textwidth}{X} 雅各見埃及有糧，就對兒子們說：「你們為甚麼彼此對看呢？」 \end{tabularx} \\ \\ \relax
42:2 & \begin{tabularx}{0.7\textwidth}{X} 他又說：「看哪，我聽見埃及有糧，你們可以下到那裡，從那裡為我們買些糧來，我們就可以存活，不至於死。」 \end{tabularx} \\ \\ \relax
42:3 & \begin{tabularx}{0.7\textwidth}{X} 於是，約瑟的十個哥哥都下去，到埃及買糧食。 \end{tabularx} \\ \\ \relax
42:4 & \begin{tabularx}{0.7\textwidth}{X} 至於約瑟的弟弟便雅憫，雅各沒有派他和哥哥們同去，因為雅各說：「恐怕他遭難。」 \end{tabularx} \\ \\ \relax
42:5 & \begin{tabularx}{0.7\textwidth}{X} 以色列的兒子們來了，在前來的人當中，為要買糧食，因為迦南地也有饑荒。 \end{tabularx} \\ \\ \relax
42:6 & \begin{tabularx}{0.7\textwidth}{X} 當時在埃及地掌權的人是約瑟，賣糧給各地眾百姓的就是他。約瑟的哥哥們來了，臉伏於地，向他下拜。 \end{tabularx} \\ \\ \relax
42:7 & \begin{tabularx}{0.7\textwidth}{X} 約瑟看見他哥哥們，就認出他們，卻對他們裝作陌生人，向他們說嚴厲的話，對他們說：「你們從哪裡來？」他們說：「我們從迦南地來買糧。」 \end{tabularx} \\ \\ \relax
42:8 & \begin{tabularx}{0.7\textwidth}{X} 約瑟認得他哥哥們，他們卻不認得他。 \end{tabularx} \\ \\ \relax
42:9 & \begin{tabularx}{0.7\textwidth}{X} 約瑟想起從前所做的那兩個夢，就對他們說：「你們是奸細，你們來是要窺探這地的虛實。」 \end{tabularx} \\ \\ \relax
42:10 & \begin{tabularx}{0.7\textwidth}{X} 他們對他說：「我主啊，不是的，僕人們是來買糧的。 \end{tabularx} \\ \\ \relax
42:11 & \begin{tabularx}{0.7\textwidth}{X} 我們都是同一個人的兒子，我們是誠實的人。僕人們並不是奸細。」 \end{tabularx} \\ \\ \relax
42:12 & \begin{tabularx}{0.7\textwidth}{X} 約瑟對他們說：「不，你們一定是窺探這地的虛實來的。」 \end{tabularx} \\ \\ \relax
42:13 & \begin{tabularx}{0.7\textwidth}{X} 他們說：「僕人們本是兄弟十二人，我們都是迦南地同一個人的兒子。看哪，最小的今日在我們父親那裡，有一個不在了。」 \end{tabularx} \\ \\ \relax
42:14 & \begin{tabularx}{0.7\textwidth}{X} 約瑟對他們說：「我剛才對你們說過了，你們是奸細！ \end{tabularx} \\ \\ \relax
42:15 & \begin{tabularx}{0.7\textwidth}{X} 我指著法老的性命起誓，若是你們最小的弟弟不到這裡來，你們就不可以離開這裡；這樣你們就可以證實自己了。 \end{tabularx} \\ \\ \relax
42:16 & \begin{tabularx}{0.7\textwidth}{X} 要派你們當中的一個人去，把你們的弟弟帶來。至於你們，都要關在這裡，好證實你們的話是不是真的。若不是，我指著法老的性命起誓，你們一定是奸細。」 \end{tabularx} \\ \\ \relax
42:17 & \begin{tabularx}{0.7\textwidth}{X} 於是約瑟把他們一起都關在監裡三天。 \end{tabularx} \\ \\ \relax
42:18 & \begin{tabularx}{0.7\textwidth}{X} 第三天，約瑟對他們說：「我是敬畏神的，你們這麼做就可以活。 \end{tabularx} \\ \\ \relax
42:19 & \begin{tabularx}{0.7\textwidth}{X} 如果你們是誠實的人，留你們兄弟中的一個關在監牢裡，你們帶糧食回去，救你們家的饑荒， \end{tabularx} \\ \\ \relax
42:20 & \begin{tabularx}{0.7\textwidth}{X} 再把你們最小的弟弟帶到我這裡來。如此，你們的話就是真的了，你們也不至於死。」他們就照樣做了。 \end{tabularx} \\ \\ \relax
42:21 & \begin{tabularx}{0.7\textwidth}{X} 他們彼此說：「我們在弟弟身上實在犯了罪。他哀求我們的時候，我們看見他的痛苦，卻不肯聽，所以這場苦難臨到我們。」 \end{tabularx} \\ \\ \relax
42:22 & \begin{tabularx}{0.7\textwidth}{X} 呂便回答他們說：「我不是對你們說過，不可傷害那孩子嗎？只是你們不肯聽，看哪，他的血在追討了。」 \end{tabularx} \\ \\ \relax
42:23 & \begin{tabularx}{0.7\textwidth}{X} 他們不知道約瑟在聽，因為在他們之間有傳譯官。 \end{tabularx} \\ \\ \relax
42:24 & \begin{tabularx}{0.7\textwidth}{X} 約瑟轉身離開他們，哭了一場，又回來對他們說話，就從他們中間抓了西緬，在他們眼前捆綁他。 \end{tabularx} \\ \\ \relax
42:25 & \begin{tabularx}{0.7\textwidth}{X} 約瑟吩咐人把他們的器皿裝滿糧食，把各人的銀子退還在各人的袋裡，又給他們路上需用的食物。人就為他們這樣做了。 \end{tabularx} \\ \\ \relax
42:26 & \begin{tabularx}{0.7\textwidth}{X} 他們把糧食馱在驢上，離開那裡去了。 \end{tabularx} \\ \\ \relax
42:27 & \begin{tabularx}{0.7\textwidth}{X} 到了住宿的地方，有一個人打開袋子，要拿飼料餵驢，就看見自己的銀子，看哪，仍在袋口上。 \end{tabularx} \\ \\ \relax
42:28 & \begin{tabularx}{0.7\textwidth}{X} 他對兄弟們說：「我的銀子退回來了，看哪，還在我袋子裡！」他們戰戰兢兢，心都快跳出來了，彼此說：「神向我們做的是甚麼呢？」 \end{tabularx} \\ \\ \relax
42:29 & \begin{tabularx}{0.7\textwidth}{X} 他們來到迦南地他們的父親雅各那裡，把所遭遇的事都告訴他，說： \end{tabularx} \\ \\ \relax
42:30 & \begin{tabularx}{0.7\textwidth}{X} 「那地的主對我們說嚴厲的話，把我們當作窺探那地的奸細。 \end{tabularx} \\ \\ \relax
42:31 & \begin{tabularx}{0.7\textwidth}{X} 我們對他說：『我們是誠實的人，並不是奸細。 \end{tabularx} \\ \\ \relax
42:32 & \begin{tabularx}{0.7\textwidth}{X} 我們本是兄弟十二人，都是同一個父親的兒子，有一個不在了，最小的今日和我們父親在迦南地。』 \end{tabularx} \\ \\ \relax
42:33 & \begin{tabularx}{0.7\textwidth}{X} 那地的主對我們說：『只有這樣我才知道你們是誠實的人：留你們兄弟中的一個在我這裡，你們帶糧食回去，救你們家的饑荒， \end{tabularx} \\ \\ \relax
42:34 & \begin{tabularx}{0.7\textwidth}{X} 再把你們最小的弟弟帶到我這裡來，我就知道你們不是奸細，是誠實的人。然後，我就把你們的兄弟交還你們，你們也可以在此地做買賣。』」 \end{tabularx} \\ \\ \relax
42:35 & \begin{tabularx}{0.7\textwidth}{X} 後來他們倒空袋子，看哪，各人的銀囊都在袋子裡。他們和父親看見銀囊就都害怕。 \end{tabularx} \\ \\ \relax
42:36 & \begin{tabularx}{0.7\textwidth}{X} 他們的父親雅各對他們說：「你們害我喪失了我的兒子：約瑟不在了，西緬也不在了，你們還要帶走便雅憫！這些事都臨到我身上了。」 \end{tabularx} \\ \\ \relax
42:37 & \begin{tabularx}{0.7\textwidth}{X} 呂便對他父親說：「我若不帶他回來給你，你可以殺我的兩個兒子。只管把他交在我手裡，我必帶他回來給你。」 \end{tabularx} \\ \\ \relax
42:38 & \begin{tabularx}{0.7\textwidth}{X} 雅各說：「我的兒子不可與你們一同下去。他哥哥死了，只剩下他。他若在你們行走的路上遭難，你們就害我白髮蒼蒼、悲悲慘慘下陰間去了。」 \end{tabularx} \\ \\ \relax
43:1 & \begin{tabularx}{0.7\textwidth}{X} 那地的饑荒非常嚴重。 \end{tabularx} \\ \\ \relax
43:2 & \begin{tabularx}{0.7\textwidth}{X} 他們從埃及帶來的糧食吃完了，父親對他們說：「你們再去給我們買些糧來。」 \end{tabularx} \\ \\ \relax
43:3 & \begin{tabularx}{0.7\textwidth}{X} 猶大對他說：「那人嚴厲地警告我們說：『你們的弟弟若不和你們同來，你們就不要來見我的面。』 \end{tabularx} \\ \\ \relax
43:4 & \begin{tabularx}{0.7\textwidth}{X} 你若派我們的弟弟跟我們同去，我們就下去給你買糧； \end{tabularx} \\ \\ \relax
43:5 & \begin{tabularx}{0.7\textwidth}{X} 你若不派他去，我們就不下去，因為那人對我們說：『你們的弟弟若不和你們同來，你們就不要來見我的面。』」 \end{tabularx} \\ \\ \relax
43:6 & \begin{tabularx}{0.7\textwidth}{X} 以色列說：「你們為甚麼這樣害我，告訴那人你們還有弟弟呢？」 \end{tabularx} \\ \\ \relax
43:7 & \begin{tabularx}{0.7\textwidth}{X} 他們說：「那人詳細問到我們和我們的家人，說：『你們的父親還在嗎？你們還有兄弟嗎？』我們就按著他的這些話告訴他，我們怎麼知道他會說：『把你們的弟弟帶下來』呢？」 \end{tabularx} \\ \\ \relax
43:8 & \begin{tabularx}{0.7\textwidth}{X} 猶大又對他父親以色列說：「請派這年輕人和我同去，我們就動身前去，好叫我們和你，以及我們的孩子都得存活，不至於死。 \end{tabularx} \\ \\ \relax
43:9 & \begin{tabularx}{0.7\textwidth}{X} 我為他擔保，你可以從我手中要人，我若不帶他回來交在你面前，我就對你永遠擔當這罪。 \end{tabularx} \\ \\ \relax
43:10 & \begin{tabularx}{0.7\textwidth}{X} 我們若沒有耽擱，現在第二趟都回來了。」 \end{tabularx} \\ \\ \relax
43:11 & \begin{tabularx}{0.7\textwidth}{X} 父親以色列對他們說：「如果必須如此，你們要這樣做：把本地土產中最好的乳香、蜂蜜、香料、沒藥、堅果、杏仁各取一點，放在器皿裡，帶下去送給那人作禮物。 \end{tabularx} \\ \\ \relax
43:12 & \begin{tabularx}{0.7\textwidth}{X} 手裡要帶雙倍的銀子，把退還在你們袋口的銀子親手帶回去；或許那是個失誤。 \end{tabularx} \\ \\ \relax
43:13 & \begin{tabularx}{0.7\textwidth}{X} 帶著你們的弟弟，動身再去見那人。 \end{tabularx} \\ \\ \relax
43:14 & \begin{tabularx}{0.7\textwidth}{X} 願全能的神使你們在那人面前蒙憐憫，放你們另一個兄弟和便雅憫回來。我若要失喪兒子，就喪了吧！」 \end{tabularx} \\ \\ \relax
43:15 & \begin{tabularx}{0.7\textwidth}{X} 於是，他們拿著那些禮物，手裡也帶雙倍的銀子，並且帶著便雅憫，動身下到埃及，站在約瑟面前。 \end{tabularx} \\ \\ \relax
43:16 & \begin{tabularx}{0.7\textwidth}{X} 約瑟見便雅憫和他們同來，就對管家說：「把這些人領到屋裡。要宰殺牲畜，預備宴席，因為中午這些人要跟我吃飯。」 \end{tabularx} \\ \\ \relax
43:17 & \begin{tabularx}{0.7\textwidth}{X} 那人就照約瑟所說的去做，領他們進約瑟的屋裡。 \end{tabularx} \\ \\ \relax
43:18 & \begin{tabularx}{0.7\textwidth}{X} 這些人因為被領到約瑟的屋裡，就害怕，說：「領我們到這裡來，必是因為當初退還在我們袋裡的銀子，要設計害我們，抓我們去當奴隸，搶奪我們的驢。」 \end{tabularx} \\ \\ \relax
43:19 & \begin{tabularx}{0.7\textwidth}{X} 他們就挨近約瑟的管家，在屋子門口和他說話， \end{tabularx} \\ \\ \relax
43:20 & \begin{tabularx}{0.7\textwidth}{X} 說：「我主啊，求求你，我們當初下來，真的是要買糧食。 \end{tabularx} \\ \\ \relax
43:21 & \begin{tabularx}{0.7\textwidth}{X} 後來到了住宿的地方，我們打開袋子，看哪，各人的銀子還在自己的袋口上，銀子的分量一點不少。現在我們親手把它帶回來， \end{tabularx} \\ \\ \relax
43:22 & \begin{tabularx}{0.7\textwidth}{X} 我們手裡又帶了另外的銀子來買糧食。我們不知道是誰把銀子放在我們袋裡的。」 \end{tabularx} \\ \\ \relax
43:23 & \begin{tabularx}{0.7\textwidth}{X} 他說：「你們平安！不要害怕，是你們的神和你們父親的神把財寶放在你們的袋裡。你們的銀子，我已經收了。」他就把西緬帶出來，交給他們。 \end{tabularx} \\ \\ \relax
43:24 & \begin{tabularx}{0.7\textwidth}{X} 那人領這些人進約瑟的屋裡，給他們水洗腳，又給他們飼料餵驢。 \end{tabularx} \\ \\ \relax
43:25 & \begin{tabularx}{0.7\textwidth}{X} 他們預備好禮物，等候約瑟中午來，因為他們聽說他們要在那裡吃飯。 \end{tabularx} \\ \\ \relax
43:26 & \begin{tabularx}{0.7\textwidth}{X} 約瑟來到家裡，他們就把手中的禮物拿進屋裡給他，俯伏在地，向他下拜。 \end{tabularx} \\ \\ \relax
43:27 & \begin{tabularx}{0.7\textwidth}{X} 約瑟問他們安，又說：「你們的父親，就是你們所說的那位老人家平安嗎？他還在嗎？」 \end{tabularx} \\ \\ \relax
43:28 & \begin{tabularx}{0.7\textwidth}{X} 他們說：「你僕人，我們的父親平安，他還在。」於是他們低頭下拜。 \end{tabularx} \\ \\ \relax
43:29 & \begin{tabularx}{0.7\textwidth}{X} 約瑟舉目看見他同母的弟弟便雅憫，就說：「你們向我所說那最小的弟弟就是這位嗎？」又說：「我兒啊，願神賜恩給你！」 \end{tabularx} \\ \\ \relax
43:30 & \begin{tabularx}{0.7\textwidth}{X} 約瑟愛弟之情激動，就急忙找個地方去哭。他進入自己的房間，哭了一場。 \end{tabularx} \\ \\ \relax
43:31 & \begin{tabularx}{0.7\textwidth}{X} 他洗了臉出來，勉強忍住，就說：「開飯吧！」 \end{tabularx} \\ \\ \relax
43:32 & \begin{tabularx}{0.7\textwidth}{X} 他們為約瑟單獨擺了一席，為那些人又擺了一席，也為和約瑟同吃飯的埃及人另擺了一席，因為埃及人不和希伯來人一同吃飯；那是埃及人所厭惡的。 \end{tabularx} \\ \\ \relax
43:33 & \begin{tabularx}{0.7\textwidth}{X} 兄弟們被安排在約瑟面前坐席，都按著長幼的次序，這些人彼此感到詫異。 \end{tabularx} \\ \\ \relax
43:34 & \begin{tabularx}{0.7\textwidth}{X} 約瑟把他面前的食物分給他們，但便雅憫所得的比別人多五倍。他們就喝酒，和約瑟一同暢飲。 \end{tabularx} \\ \\ \relax
44:1 & \begin{tabularx}{0.7\textwidth}{X} 約瑟吩咐管家說：「按照他們的驢子所能馱的，把這些人的袋子裝滿糧食，再把各人的銀子放在各人的袋口上， \end{tabularx} \\ \\ \relax
44:2 & \begin{tabularx}{0.7\textwidth}{X} 我的杯，就是那個銀杯，要和買糧的銀子一同放在最年輕的那個人的袋口上。」管家就照約瑟所說的話去做了。 \end{tabularx} \\ \\ \relax
44:3 & \begin{tabularx}{0.7\textwidth}{X} 天一亮，這些人和他們的驢子就被送走了。 \end{tabularx} \\ \\ \relax
44:4 & \begin{tabularx}{0.7\textwidth}{X} 他們出城走了不遠，約瑟對管家說：「起來，去追那些人，追上了就對他們說：『你們為甚麼以惡報善呢？ \end{tabularx} \\ \\ \relax
44:5 & \begin{tabularx}{0.7\textwidth}{X} 這不是我主人用來飲酒，確實用它來占卜的嗎？你們這麼做是不對的！』」 \end{tabularx} \\ \\ \relax
44:6 & \begin{tabularx}{0.7\textwidth}{X} 管家追上他們，把這些話對他們說了。 \end{tabularx} \\ \\ \relax
44:7 & \begin{tabularx}{0.7\textwidth}{X} 他們對他說：「我主為甚麼說這樣的話呢？你僕人們絕不會做這樣的事。 \end{tabularx} \\ \\ \relax
44:8 & \begin{tabularx}{0.7\textwidth}{X} 看哪，我們從前在袋口上發現的銀子，尚且從迦南地帶來還你，我們又怎麼會從你主人家裡偷竊金銀呢？ \end{tabularx} \\ \\ \relax
44:9 & \begin{tabularx}{0.7\textwidth}{X} 你僕人中無論在誰那裡找到杯子，就叫他死，我們也要作我主的奴隸。」 \end{tabularx} \\ \\ \relax
44:10 & \begin{tabularx}{0.7\textwidth}{X} 管家說：「現在就照你們的話做吧！在誰那裡找到杯子，誰就作我的奴隸，其餘的人都沒有罪。」 \end{tabularx} \\ \\ \relax
44:11 & \begin{tabularx}{0.7\textwidth}{X} 於是他們各人急忙把袋子卸在地上，各人打開自己的袋子。 \end{tabularx} \\ \\ \relax
44:12 & \begin{tabularx}{0.7\textwidth}{X} 管家就搜查，從年長的開始到年幼的為止，那杯竟在便雅憫的袋子裡找到了。 \end{tabularx} \\ \\ \relax
44:13 & \begin{tabularx}{0.7\textwidth}{X} 他們就撕裂衣服，各人把馱子抬在驢上，回城去了。 \end{tabularx} \\ \\ \relax
44:14 & \begin{tabularx}{0.7\textwidth}{X} 猶大和他兄弟們來到約瑟的屋裡，約瑟還在那裡，他們就在他面前俯伏於地。 \end{tabularx} \\ \\ \relax
44:15 & \begin{tabularx}{0.7\textwidth}{X} 約瑟對他們說：「你們做的是甚麼事呢？你們豈不知像我這樣的人必懂得占卜嗎？」 \end{tabularx} \\ \\ \relax
44:16 & \begin{tabularx}{0.7\textwidth}{X} 猶大說：「我們對我主能說甚麼呢？還有甚麼話可說呢？我們還能為自己表白嗎？神已經查出你僕人的罪孽了。看哪，我們與那在他手中找到杯子的人都是我主的奴隸。」 \end{tabularx} \\ \\ \relax
44:17 & \begin{tabularx}{0.7\textwidth}{X} 約瑟說：「我絕不能做這樣的事！誰的手中找到杯子，誰就作我的奴隸。至於你們，可以平平安安上到你們父親那裡去。」 \end{tabularx} \\ \\ \relax
44:18 & \begin{tabularx}{0.7\textwidth}{X} 猶大挨近他，說：「我主啊，求求你，讓僕人說一句話給我主聽，不要向僕人發烈怒，因為你如同法老一樣。 \end{tabularx} \\ \\ \relax
44:19 & \begin{tabularx}{0.7\textwidth}{X} 我主曾問僕人們說：『你們有父親、兄弟沒有？』 \end{tabularx} \\ \\ \relax
44:20 & \begin{tabularx}{0.7\textwidth}{X} 我們對我主說：『我們有父親，他已經年老，還有他老年所生的一個小兒子。他哥哥死了，他的母親只剩下他一個孩子，父親也疼愛他。』 \end{tabularx} \\ \\ \relax
44:21 & \begin{tabularx}{0.7\textwidth}{X} 你對僕人說：『把他帶下到我這裡來，讓我親眼看看他。』 \end{tabularx} \\ \\ \relax
44:22 & \begin{tabularx}{0.7\textwidth}{X} 我們對我主說：『這年輕人不能離開他父親，若是離開，父親就會死。』 \end{tabularx} \\ \\ \relax
44:23 & \begin{tabularx}{0.7\textwidth}{X} 你對僕人說：『你們最小的弟弟若不和你們一同下來，你們就不要來見我的面。』 \end{tabularx} \\ \\ \relax
44:24 & \begin{tabularx}{0.7\textwidth}{X} 我們上到你僕人，我們父親那裡，就把我主的話告訴了他。 \end{tabularx} \\ \\ \relax
44:25 & \begin{tabularx}{0.7\textwidth}{X} 後來，我們的父親說：『你們再去給我買些糧來。』 \end{tabularx} \\ \\ \relax
44:26 & \begin{tabularx}{0.7\textwidth}{X} 我們說：『我們不能下去。最小的弟弟若和我們同去，我們就可以下去。因為，最小的弟弟若不和我們同去，我們必不能見那人的面。』 \end{tabularx} \\ \\ \relax
44:27 & \begin{tabularx}{0.7\textwidth}{X} 你僕人，我父親對我們說：『你們知道我的妻子給我生了兩個兒子。 \end{tabularx} \\ \\ \relax
44:28 & \begin{tabularx}{0.7\textwidth}{X} 一個離開我走了，我說他必是被野獸撕碎了，直到如今我再沒有見過他； \end{tabularx} \\ \\ \relax
44:29 & \begin{tabularx}{0.7\textwidth}{X} 現在你們又要把這個從我面前帶走。倘若他遭難，那麼你們就害我白髮蒼蒼、悲悲慘慘下陰間去了。』 \end{tabularx} \\ \\ \relax
44:30 & \begin{tabularx}{0.7\textwidth}{X} 如今我回到你僕人，我父親那裡，若沒有這年輕人和我們同去，我父親的命是與這年輕人的命相連的， \end{tabularx} \\ \\ \relax
44:31 & \begin{tabularx}{0.7\textwidth}{X} 當我們的父親看見沒有了這年輕人，他就會死。這樣，我們就害你僕人，我們的父親白髮蒼蒼、悲悲慘慘下陰間去了。 \end{tabularx} \\ \\ \relax
44:32 & \begin{tabularx}{0.7\textwidth}{X} 僕人曾向我父親為這年輕人擔保，說：『我若不帶他回來交給父親，我就在父親面前永遠擔當這罪。』 \end{tabularx} \\ \\ \relax
44:33 & \begin{tabularx}{0.7\textwidth}{X} 現在，求你把僕人留下，代替這年輕人作我主的奴隸，讓這年輕人和他哥哥們一同上去。 \end{tabularx} \\ \\ \relax
44:34 & \begin{tabularx}{0.7\textwidth}{X} 若這年輕人不和我一起，我怎能上到我父親那裡呢？恐怕我要看到災禍臨到我父親了。」 \end{tabularx} \\ \\ \relax
45:1 & \begin{tabularx}{0.7\textwidth}{X} 約瑟在所有侍立在他旁邊的人面前情不自禁，就喊叫說：「每一個人都離開我，出去吧！」約瑟和兄弟相認的時候沒有一人站在他那裡。 \end{tabularx} \\ \\ \relax
45:2 & \begin{tabularx}{0.7\textwidth}{X} 他放聲大哭，埃及人聽見了，法老家中的人也聽見了。 \end{tabularx} \\ \\ \relax
45:3 & \begin{tabularx}{0.7\textwidth}{X} 約瑟對他兄弟們說：「我就是約瑟。我的父親還在嗎？」他兄弟們不敢回答他，因為他們在他面前都很驚惶。 \end{tabularx} \\ \\ \relax
45:4 & \begin{tabularx}{0.7\textwidth}{X} 約瑟又對他兄弟們說：「靠近我一點。」他們就近前來。他說：「我是被你們賣到埃及的兄弟約瑟。 \end{tabularx} \\ \\ \relax
45:5 & \begin{tabularx}{0.7\textwidth}{X} 現在，不要因為把我賣到這裡而憂傷，對自己生氣，因為神差我在你們以先來，為要保全性命。 \end{tabularx} \\ \\ \relax
45:6 & \begin{tabularx}{0.7\textwidth}{X} 現在這地的饑荒已經二年了，還有五年不能耕種，沒有收成。 \end{tabularx} \\ \\ \relax
45:7 & \begin{tabularx}{0.7\textwidth}{X} 神差我在你們以先來，為要給你們在世上存留餘種，大施拯救，保全你們的性命。 \end{tabularx} \\ \\ \relax
45:8 & \begin{tabularx}{0.7\textwidth}{X} 這樣看來，差我到這裡來的不是你們，而是神。他又使我如同法老之父，作他全家之主，和埃及全地掌權的人。 \end{tabularx} \\ \\ \relax
45:9 & \begin{tabularx}{0.7\textwidth}{X} 你們要趕緊上到我父親那裡，對他說：『你兒子約瑟這樣說：神已立我作全埃及之主，請你下到我這裡來，不要耽擱。 \end{tabularx} \\ \\ \relax
45:10 & \begin{tabularx}{0.7\textwidth}{X} 你和你的兒子孫子，羊群牛群，以及一切所有的，都可以住在歌珊地，與我相近。 \end{tabularx} \\ \\ \relax
45:11 & \begin{tabularx}{0.7\textwidth}{X} 我要在那裡奉養你，因為還有五年的饑荒，免得你和你的家屬，以及一切所有的，都陷入窮困中。』 \end{tabularx} \\ \\ \relax
45:12 & \begin{tabularx}{0.7\textwidth}{X} 看哪，你們的眼睛和我弟弟便雅憫的眼睛都看見，是我親口對你們說話。 \end{tabularx} \\ \\ \relax
45:13 & \begin{tabularx}{0.7\textwidth}{X} 你們要把我在埃及一切的尊榮和你們所有看見的事情都告訴我父親，也要趕緊請我父親下到這裡來。」 \end{tabularx} \\ \\ \relax
45:14 & \begin{tabularx}{0.7\textwidth}{X} 於是約瑟伏在他弟弟便雅憫的頸項上哭，便雅憫也在他的頸項上哭。 \end{tabularx} \\ \\ \relax
45:15 & \begin{tabularx}{0.7\textwidth}{X} 他又親眾兄弟，伏著他們哭。過後，他的兄弟就和他說話。 \end{tabularx} \\ \\ \relax
45:16 & \begin{tabularx}{0.7\textwidth}{X} 這消息傳到法老的宮裡，說：「約瑟的兄弟們來了。」法老和他的臣僕眼中都看為好。 \end{tabularx} \\ \\ \relax
45:17 & \begin{tabularx}{0.7\textwidth}{X} 法老對約瑟說：「你要吩咐你的兄弟們說：『你們要這樣做：把馱子抬在牲口上，動身到迦南地去， \end{tabularx} \\ \\ \relax
45:18 & \begin{tabularx}{0.7\textwidth}{X} 請你們的父親和你們的家屬都到我這裡來，我要把埃及地的美物賜給你們，你們也要吃這地肥美的出產。』 \end{tabularx} \\ \\ \relax
45:19 & \begin{tabularx}{0.7\textwidth}{X} 你要吩咐他們：『要這樣做：從埃及地帶著車輛去，把你們的孩子和妻子，以及你們的父親都接來。 \end{tabularx} \\ \\ \relax
45:20 & \begin{tabularx}{0.7\textwidth}{X} 你們的眼不要顧惜你們的家具，因為埃及全地的美物都是你們的。』」 \end{tabularx} \\ \\ \relax
45:21 & \begin{tabularx}{0.7\textwidth}{X} 以色列的兒子們就照樣做了。約瑟遵照法老的吩咐，給他們車輛和路上需用的食物。 \end{tabularx} \\ \\ \relax
45:22 & \begin{tabularx}{0.7\textwidth}{X} 他又給所有哥哥每人一套衣服， 卻給便雅憫三百銀子，五套衣服。 \end{tabularx} \\ \\ \relax
45:23 & \begin{tabularx}{0.7\textwidth}{X} 他也送給父親十匹公驢，馱著埃及的美物，以及十匹母驢，馱著給他父親在路上需用的穀物、餅和糧食。 \end{tabularx} \\ \\ \relax
45:24 & \begin{tabularx}{0.7\textwidth}{X} 於是約瑟送他的兄弟們回去，對他們說：「你們不要在路上爭吵。」 \end{tabularx} \\ \\ \relax
45:25 & \begin{tabularx}{0.7\textwidth}{X} 他們從埃及上去，來到迦南地他們的父親雅各那裡， \end{tabularx} \\ \\ \relax
45:26 & \begin{tabularx}{0.7\textwidth}{X} 告訴他說：「約瑟還活著，並且作了埃及全地掌權的人。」雅各心裡冰涼，因為不信他們。 \end{tabularx} \\ \\ \relax
45:27 & \begin{tabularx}{0.7\textwidth}{X} 他們就把約瑟對他們所說一切的話都告訴了他。他看見約瑟派來接他的車輛，他們父親雅各的靈就甦醒了。 \end{tabularx} \\ \\ \relax
45:28 & \begin{tabularx}{0.7\textwidth}{X} 以色列說：「夠了！我的兒子約瑟還活著，我要趁我未死之前去見他。」 \end{tabularx} \\ \\ \relax
46:1 & \begin{tabularx}{0.7\textwidth}{X} 以色列帶著一切所有的，起程到別是巴去，獻祭給他父親以撒的神。 \end{tabularx} \\ \\ \relax
46:2 & \begin{tabularx}{0.7\textwidth}{X} 夜間，神在異象中對以色列說：「雅各！雅各！」他說：「我在這裡。」 \end{tabularx} \\ \\ \relax
46:3 & \begin{tabularx}{0.7\textwidth}{X} 神說：「我是神，你父親的神。不要害怕下埃及去，因為我必使你在那裡成為大國。 \end{tabularx} \\ \\ \relax
46:4 & \begin{tabularx}{0.7\textwidth}{X} 我要和你同下埃及去，也必定帶你上來；約瑟要親手合上你的眼睛。」 \end{tabularx} \\ \\ \relax
46:5 & \begin{tabularx}{0.7\textwidth}{X} 雅各就從別是巴起行。以色列的兒子讓他們的父親雅各和他們的孩子、妻子都坐在法老為雅各派來的車上。 \end{tabularx} \\ \\ \relax
46:6 & \begin{tabularx}{0.7\textwidth}{X} 他們也帶著迦南地所得的牲畜和財物來到埃及。雅各和他所有的子孫都一同來了。 \end{tabularx} \\ \\ \relax
46:7 & \begin{tabularx}{0.7\textwidth}{X} 他把他的兒子、孫子、女兒、孫女，他所有的子孫一同帶到埃及。 \end{tabularx} \\ \\ \relax
46:8 & \begin{tabularx}{0.7\textwidth}{X} 這些是來到埃及的以色列人，雅各和他子孫的名字：雅各的長子是呂便。 \end{tabularx} \\ \\ \relax
46:9 & \begin{tabularx}{0.7\textwidth}{X} 呂便的兒子是哈諾、法路、希斯倫、迦米。 \end{tabularx} \\ \\ \relax
46:10 & \begin{tabularx}{0.7\textwidth}{X} 西緬的兒子是耶母利、雅憫、阿轄、雅斤、瑣轄，還有迦南女子生的兒子掃羅。 \end{tabularx} \\ \\ \relax
46:11 & \begin{tabularx}{0.7\textwidth}{X} 利未的兒子是革順、哥轄、米拉利。 \end{tabularx} \\ \\ \relax
46:12 & \begin{tabularx}{0.7\textwidth}{X} 猶大的兒子是珥、俄南、示拉、法勒斯、謝拉；珥與俄南死在迦南地。法勒斯的兒子是希斯崙、哈母勒。 \end{tabularx} \\ \\ \relax
46:13 & \begin{tabularx}{0.7\textwidth}{X} 以薩迦的兒子是陀拉、普瓦、約伯、伸崙。 \end{tabularx} \\ \\ \relax
46:14 & \begin{tabularx}{0.7\textwidth}{X} 西布倫的兒子是西烈、以倫、雅利。 \end{tabularx} \\ \\ \relax
46:15 & \begin{tabularx}{0.7\textwidth}{X} 這是利亞在巴旦‧亞蘭為雅各所生的兒孫，還有女兒底拿，兒孫共三十三人。 \end{tabularx} \\ \\ \relax
46:16 & \begin{tabularx}{0.7\textwidth}{X} 迦得的兒子是洗非芸、哈基、書尼、以斯本、以利、亞羅底、亞列利。 \end{tabularx} \\ \\ \relax
46:17 & \begin{tabularx}{0.7\textwidth}{X} 亞設的兒子是音拿、亦施瓦、亦施韋、比利亞，還有他們的妹妹西拉。比利亞的兒子是希別、瑪結。 \end{tabularx} \\ \\ \relax
46:18 & \begin{tabularx}{0.7\textwidth}{X} 這是拉班給他女兒利亞的婢女悉帕的兒孫，她為雅各所生的共有十六人。 \end{tabularx} \\ \\ \relax
46:19 & \begin{tabularx}{0.7\textwidth}{X} 雅各之妻拉結的兒子是約瑟和便雅憫。 \end{tabularx} \\ \\ \relax
46:20 & \begin{tabularx}{0.7\textwidth}{X} 約瑟在埃及地生了瑪拿西和以法蓮，是安城的祭司波提‧非拉的女兒亞西納為約瑟生的。 \end{tabularx} \\ \\ \relax
46:21 & \begin{tabularx}{0.7\textwidth}{X} 便雅憫的兒子是比拉、比結、亞實別、基拉、乃幔、以希、羅實、母平、戶平、亞勒。 \end{tabularx} \\ \\ \relax
46:22 & \begin{tabularx}{0.7\textwidth}{X} 這是拉結為雅各所生的兒孫，共有十四人。 \end{tabularx} \\ \\ \relax
46:23 & \begin{tabularx}{0.7\textwidth}{X} 但的兒子是戶伸。 \end{tabularx} \\ \\ \relax
46:24 & \begin{tabularx}{0.7\textwidth}{X} 拿弗他利的兒子是雅薛、沽尼、耶色、示冷。 \end{tabularx} \\ \\ \relax
46:25 & \begin{tabularx}{0.7\textwidth}{X} 這是拉班給他女兒拉結的婢女辟拉的兒孫，她為雅各所生的共有七人。 \end{tabularx} \\ \\ \relax
46:26 & \begin{tabularx}{0.7\textwidth}{X} 那與雅各同到埃及的，除了他媳婦之外，凡從他生的共有六十六人。 \end{tabularx} \\ \\ \relax
46:27 & \begin{tabularx}{0.7\textwidth}{X} 還有約瑟在埃及所生的兩個兒子。到埃及的雅各全家共有七十人。 \end{tabularx} \\ \\ \relax
46:28 & \begin{tabularx}{0.7\textwidth}{X} 雅各派猶大先到約瑟那裡，請他先指示到歌珊去的路；於是他們來到了歌珊地。 \end{tabularx} \\ \\ \relax
46:29 & \begin{tabularx}{0.7\textwidth}{X} 約瑟備好座車，上歌珊去迎接他的父親以色列。他見到父親，就伏在父親的頸項上，在父親的頸項上哭了許久。 \end{tabularx} \\ \\ \relax
46:30 & \begin{tabularx}{0.7\textwidth}{X} 以色列對約瑟說：「我見了你的面，知道你還活著，現在我可以死了。」 \end{tabularx} \\ \\ \relax
46:31 & \begin{tabularx}{0.7\textwidth}{X} 約瑟對他兄弟和他父親的全家說：「我要上去告訴法老，對他說：『我在迦南地的兄弟和我父親的全家，都到我這裡來了。 \end{tabularx} \\ \\ \relax
46:32 & \begin{tabularx}{0.7\textwidth}{X} 他們是牧羊人，是牧放牲畜的人；他們把羊群牛群和一切所有的都帶來了。』 \end{tabularx} \\ \\ \relax
46:33 & \begin{tabularx}{0.7\textwidth}{X} 等到法老召見你們，說：『你們是做甚麼的？』 \end{tabularx} \\ \\ \relax
46:34 & \begin{tabularx}{0.7\textwidth}{X} 你們就說：『你的僕人，從幼年直到現在，都是牧放牲畜的人，我們和我們的祖宗都是這樣。』如此，你們就可以住在歌珊地，因為凡牧羊的都被埃及人厭惡。」 \end{tabularx} \\ \\ \relax
47:1 & \begin{tabularx}{0.7\textwidth}{X} 約瑟進去告訴法老說：「我的父親和我的兄弟帶著羊群牛群，以及他們一切所有的，從迦南地來了。看哪，他們正在歌珊地。」 \end{tabularx} \\ \\ \relax
47:2 & \begin{tabularx}{0.7\textwidth}{X} 約瑟從他所有兄弟中挑選五個人，引他們到法老面前。 \end{tabularx} \\ \\ \relax
47:3 & \begin{tabularx}{0.7\textwidth}{X} 法老對約瑟的兄弟說：「你們是做甚麼的？」他們對法老說：「你僕人是牧羊的，我們和我們的祖宗都是這樣。」 \end{tabularx} \\ \\ \relax
47:4 & \begin{tabularx}{0.7\textwidth}{X} 他們又對法老說：「迦南地的饑荒非常嚴重，僕人的羊群沒有牧草，所以我們來到這地寄居。現在求你准許僕人住在歌珊地。」 \end{tabularx} \\ \\ \relax
47:5 & \begin{tabularx}{0.7\textwidth}{X} 法老對約瑟說：「你的父親和你的兄弟到你這裡來了， \end{tabularx} \\ \\ \relax
47:6 & \begin{tabularx}{0.7\textwidth}{X} 埃及地都在你面前，只管讓你父親和你兄弟住在最好的地，他們可以住在歌珊地。你若知道他們中間有能幹的人，就派他們看管我的牲畜。」 \end{tabularx} \\ \\ \relax
47:7 & \begin{tabularx}{0.7\textwidth}{X} 約瑟帶他父親雅各來，站在法老面前，雅各就為法老祝福。 \end{tabularx} \\ \\ \relax
47:8 & \begin{tabularx}{0.7\textwidth}{X} 法老對雅各說：「你平生的年日是多少呢？」 \end{tabularx} \\ \\ \relax
47:9 & \begin{tabularx}{0.7\textwidth}{X} 雅各對法老說：「我在世寄居的年日是一百三十年，我一生的歲月又短又苦，比不上我祖先在世寄居的年日。」 \end{tabularx} \\ \\ \relax
47:10 & \begin{tabularx}{0.7\textwidth}{X} 雅各又為法老祝福，就從法老面前退出去了。 \end{tabularx} \\ \\ \relax
47:11 & \begin{tabularx}{0.7\textwidth}{X} 約瑟安頓他的父親和兄弟，遵照法老的命令，把埃及境內最好的地，就是蘭塞地，給他們作為產業。 \end{tabularx} \\ \\ \relax
47:12 & \begin{tabularx}{0.7\textwidth}{X} 約瑟用糧食供給他父親和兄弟們，以及他父親全家的人，照扶養親屬的人口供給。 \end{tabularx} \\ \\ \relax
47:13 & \begin{tabularx}{0.7\textwidth}{X} 饑荒非常嚴重，全地都絕了糧，埃及地和迦南地都因饑荒耗損了。 \end{tabularx} \\ \\ \relax
47:14 & \begin{tabularx}{0.7\textwidth}{X} 約瑟收集了埃及地和迦南地所有的銀子，就是眾人買糧的銀子，約瑟就把那些銀子都帶到法老的宮裡。 \end{tabularx} \\ \\ \relax
47:15 & \begin{tabularx}{0.7\textwidth}{X} 埃及地和迦南地的銀子都花光了，埃及眾人到約瑟那裡，說：「我們的銀子都用完了，求你給我們糧食吧！我們為甚麼要死在你面前呢？」 \end{tabularx} \\ \\ \relax
47:16 & \begin{tabularx}{0.7\textwidth}{X} 約瑟說：「銀子若是用完了，可以把你們的牲畜賣給我，我就以你們的牲畜換糧食給你們。」 \end{tabularx} \\ \\ \relax
47:17 & \begin{tabularx}{0.7\textwidth}{X} 於是他們把牲畜帶到約瑟那裡，約瑟就拿糧食換了他們的馬、羊、牛、驢；那一年他因換他們一切的牲畜，用糧食養活他們。 \end{tabularx} \\ \\ \relax
47:18 & \begin{tabularx}{0.7\textwidth}{X} 那一年過去，第二年他們又來到約瑟那裡，對他說：「不瞞我主，我們的銀子都花光了，牲畜也都歸於我主了。我們在我主面前，除了自己的身體和土地以外，一無所剩。 \end{tabularx} \\ \\ \relax
47:19 & \begin{tabularx}{0.7\textwidth}{X} 你為甚麼要眼看著我們人死地荒呢？求你用糧食買我們和我們的地，我們和我們的地就要為法老效力。求你給我們種子，使我們可以存活，不致死亡，土地也不致荒蕪。」 \end{tabularx} \\ \\ \relax
47:20 & \begin{tabularx}{0.7\textwidth}{X} 於是，約瑟為法老買了埃及所有的土地，埃及人因饑荒所迫，都賣了自己的田地；那些地都歸給法老了。 \end{tabularx} \\ \\ \relax
47:21 & \begin{tabularx}{0.7\textwidth}{X} 至於百姓，從埃及邊界的一端到另一端，約瑟使他們作奴隸。 \end{tabularx} \\ \\ \relax
47:22 & \begin{tabularx}{0.7\textwidth}{X} 只有祭司的土地，約瑟沒有買，因為祭司從法老領取薪俸，靠法老的薪俸過活，所以沒有賣自己的土地。 \end{tabularx} \\ \\ \relax
47:23 & \begin{tabularx}{0.7\textwidth}{X} 約瑟對百姓說：「看哪，我今日為法老買了你們和你們的土地。看哪，這些種子是給你們的，你們可以耕種土地。 \end{tabularx} \\ \\ \relax
47:24 & \begin{tabularx}{0.7\textwidth}{X} 將來收割的時候，你們要把五分之一納給法老，另外四分可以給你們作田地的種子，作你們和你們全家大小的食物。」 \end{tabularx} \\ \\ \relax
47:25 & \begin{tabularx}{0.7\textwidth}{X} 他們說：「你救了我們的性命，願我們在我主眼前蒙恩，我們情願作法老的奴隸。」 \end{tabularx} \\ \\ \relax
47:26 & \begin{tabularx}{0.7\textwidth}{X} 於是約瑟為埃及的土地立下定例，直到今日，就是收成的五分之一要歸法老。惟獨祭司的土地例外，不歸於法老。 \end{tabularx} \\ \\ \relax
47:27 & \begin{tabularx}{0.7\textwidth}{X} 以色列人住在埃及境內的歌珊地。他們在那裡得了產業，並且生養眾多。 \end{tabularx} \\ \\ \relax
47:28 & \begin{tabularx}{0.7\textwidth}{X} 雅各住在埃及地十七年。雅各一生的年日是一百四十七年。 \end{tabularx} \\ \\ \relax
47:29 & \begin{tabularx}{0.7\textwidth}{X} 以色列的死期快到了，就叫了他兒子約瑟來，對他說：「我若在你眼前蒙恩，把你的手放在我大腿底下，以慈愛和誠實向我承諾，必不將我葬在埃及。 \end{tabularx} \\ \\ \relax
47:30 & \begin{tabularx}{0.7\textwidth}{X} 我與我祖先同睡的時候，你要將我帶出埃及，把我葬在他們所葬的地方。」約瑟說：「我必遵照你的吩咐去做。」 \end{tabularx} \\ \\ \relax
47:31 & \begin{tabularx}{0.7\textwidth}{X} 雅各說：「你向我起誓吧！」約瑟就向他起了誓。於是以色列在床頭敬拜。 \end{tabularx} \\ \\ \relax
48:1 & \begin{tabularx}{0.7\textwidth}{X} 這些事以後，有人告訴約瑟說：「看哪，你的父親病了。」他就帶著兩個兒子瑪拿西和以法蓮同去。 \end{tabularx} \\ \\ \relax
48:2 & \begin{tabularx}{0.7\textwidth}{X} 有人告訴雅各說：「看哪，你的兒子約瑟到你這裡來了。」以色列就勉強在床上坐起來。 \end{tabularx} \\ \\ \relax
48:3 & \begin{tabularx}{0.7\textwidth}{X} 雅各對約瑟說：「全能的神曾在迦南地的路斯向我顯現，賜福給我， \end{tabularx} \\ \\ \relax
48:4 & \begin{tabularx}{0.7\textwidth}{X} 對我說：『看哪，我必使你生養眾多，成為許多民族，又要將這地賜給你的後裔，永遠為業。』 \end{tabularx} \\ \\ \relax
48:5 & \begin{tabularx}{0.7\textwidth}{X} 我未到埃及你那裡之前，你在埃及地所生的以法蓮和瑪拿西這兩個兒子，現在他們是我的，正如呂便和西緬是我的一樣。 \end{tabularx} \\ \\ \relax
48:6 & \begin{tabularx}{0.7\textwidth}{X} 你在他們以後所生的後裔就是你的，這些後裔可以在自己兄弟的名下得產業。 \end{tabularx} \\ \\ \relax
48:7 & \begin{tabularx}{0.7\textwidth}{X} 至於我，我從巴旦回來的時候，拉結在我身旁死了，就是在往迦南地的路上，離以法他還有一段路程。我就把她葬在往以法他的路旁；以法他就是伯利恆。」 \end{tabularx} \\ \\ \relax
48:8 & \begin{tabularx}{0.7\textwidth}{X} 以色列看見約瑟的兒子，就說：「這些是誰？」 \end{tabularx} \\ \\ \relax
48:9 & \begin{tabularx}{0.7\textwidth}{X} 約瑟對他父親說：「這是神在這裡賜給我的兒子。」以色列說：「領他們到我跟前，我要為他們祝福。」 \end{tabularx} \\ \\ \relax
48:10 & \begin{tabularx}{0.7\textwidth}{X} 以色列年紀老邁，眼睛昏花，不能看見。約瑟領他們到他跟前，他就和他們親吻，抱著他們。 \end{tabularx} \\ \\ \relax
48:11 & \begin{tabularx}{0.7\textwidth}{X} 以色列對約瑟說：「我沒有想到能夠見你的面。看哪，神還讓我看見你的兒子。」 \end{tabularx} \\ \\ \relax
48:12 & \begin{tabularx}{0.7\textwidth}{X} 約瑟把他們從以色列兩膝中間領出來，自己臉伏於地下拜。 \end{tabularx} \\ \\ \relax
48:13 & \begin{tabularx}{0.7\textwidth}{X} 然後，約瑟牽著他們兩個，帶到父親跟前，右手牽以法蓮到以色列的左邊，左手牽瑪拿西到以色列的右邊。 \end{tabularx} \\ \\ \relax
48:14 & \begin{tabularx}{0.7\textwidth}{X} 以色列卻伸出右手來，按在次子以法蓮的頭上，又交叉伸出左手來，按在長子瑪拿西的頭上。 \end{tabularx} \\ \\ \relax
48:15 & \begin{tabularx}{0.7\textwidth}{X} 他就為約瑟祝福說：「願我祖父亞伯拉罕和我父親以撒所事奉的神，就是一生牧養我直到今日的神， \end{tabularx} \\ \\ \relax
48:16 & \begin{tabularx}{0.7\textwidth}{X} 救贖我脫離一切患難的那位使者，賜福給這兩個孩子。願我的名，我祖父亞伯拉罕和我父親以撒的名藉著他們得以流傳。又願他們在全地上多多繁衍。」 \end{tabularx} \\ \\ \relax
48:17 & \begin{tabularx}{0.7\textwidth}{X} 約瑟見父親把右手按在以法蓮的頭上，他看為不好，就提起他父親的手，要從以法蓮的頭上移到瑪拿西的頭上。 \end{tabularx} \\ \\ \relax
48:18 & \begin{tabularx}{0.7\textwidth}{X} 約瑟對父親說：「我父，不是這樣。這個才是長子，請你把右手按在他頭上。」 \end{tabularx} \\ \\ \relax
48:19 & \begin{tabularx}{0.7\textwidth}{X} 他父親卻不肯，說：「我知道，我兒，我知道。他也要成為一族，也要強大。可是他的弟弟將來比他還要強大；他弟弟的後裔要成為許多國家。」 \end{tabularx} \\ \\ \relax
48:20 & \begin{tabularx}{0.7\textwidth}{X} 以色列就在當日為他們祝福，說：「以色列人要指著你們祝福，說：『願神使你如以法蓮、瑪拿西一樣。』」於是他立以法蓮在瑪拿西之上。 \end{tabularx} \\ \\ \relax
48:21 & \begin{tabularx}{0.7\textwidth}{X} 以色列又對約瑟說：「看哪，我快要死了，但神必與你們同在，領你們回到你們祖先之地。 \end{tabularx} \\ \\ \relax
48:22 & \begin{tabularx}{0.7\textwidth}{X} 從前我用刀用弓從亞摩利人手下奪取的那一份，我要把它賜給你，使你比你的兄弟多得一份。」 \end{tabularx} \\ \\ \relax
49:1 & \begin{tabularx}{0.7\textwidth}{X} 雅各叫了他的兒子來，說：「你們都來聚集，讓我把你們日後要遇到的事告訴你們。 \end{tabularx} \\ \\ \relax
49:2 & \begin{tabularx}{0.7\textwidth}{X} 雅各的兒子們，你們要聚集，要聆聽，聽你們父親以色列的話。 \end{tabularx} \\ \\ \relax
49:3 & \begin{tabularx}{0.7\textwidth}{X} 呂便啊，你是我的長子，我的力量，我壯年頭生之子，極有尊榮，權力超群。 \end{tabularx} \\ \\ \relax
49:4 & \begin{tabularx}{0.7\textwidth}{X} 你卻放縱如水，必不得居首位；因為你上了你父親的床，你上了我的榻，污辱了它！ \end{tabularx} \\ \\ \relax
49:5 & \begin{tabularx}{0.7\textwidth}{X} 西緬和利未是兄弟；他們的刀劍是殘暴的兵器。 \end{tabularx} \\ \\ \relax
49:6 & \begin{tabularx}{0.7\textwidth}{X} 願我的心不與他們同謀，願我的靈不與他們合夥；因為他們在烈怒中殺人，任意割斷牛腿的筋。 \end{tabularx} \\ \\ \relax
49:7 & \begin{tabularx}{0.7\textwidth}{X} 他們火爆的烈怒可詛，他們兇殘的憤恨可咒！我要把他們分散在雅各中，使他們散居在以色列。 \end{tabularx} \\ \\ \relax
49:8 & \begin{tabularx}{0.7\textwidth}{X} 猶大啊，你的兄弟必讚美你，你的手必掐住仇敵的頸項，你父親的兒子要向你下拜。 \end{tabularx} \\ \\ \relax
49:9 & \begin{tabularx}{0.7\textwidth}{X} 猶大是隻小獅子；我兒啊，你捕獲了獵物就上去。他蹲伏，他躺臥，如公獅，又如母獅，誰敢惹他呢？ \end{tabularx} \\ \\ \relax
49:10 & \begin{tabularx}{0.7\textwidth}{X} 權杖必不離猶大，統治者的杖必不離他兩腳之間，直等細羅來到，萬民都要歸順他。 \end{tabularx} \\ \\ \relax
49:11 & \begin{tabularx}{0.7\textwidth}{X} 猶大把小驢拴在葡萄樹上，把驢駒拴在佳美的葡萄樹上。他在葡萄酒中洗衣服，在葡萄汁中洗長袍。 \end{tabularx} \\ \\ \relax
49:12 & \begin{tabularx}{0.7\textwidth}{X} 他的眼睛比酒紅潤，他的牙齒比奶潔白。 \end{tabularx} \\ \\ \relax
49:13 & \begin{tabularx}{0.7\textwidth}{X} 西布倫必住在海邊，必成為停船的港口；他的疆界必延到西頓。 \end{tabularx} \\ \\ \relax
49:14 & \begin{tabularx}{0.7\textwidth}{X} 以薩迦是匹強壯的驢，臥在羊圈之中。 \end{tabularx} \\ \\ \relax
49:15 & \begin{tabularx}{0.7\textwidth}{X} 他看見居所安舒，土地肥美，就屈肩負重，成為服勞役的僕人。 \end{tabularx} \\ \\ \relax
49:16 & \begin{tabularx}{0.7\textwidth}{X} 但必為他的百姓伸冤，作為以色列支派之一。 \end{tabularx} \\ \\ \relax
49:17 & \begin{tabularx}{0.7\textwidth}{X} 但必作道旁的蛇，路邊的毒蛇，咬傷馬蹄，使騎馬的人向後墜落。 \end{tabularx} \\ \\ \relax
49:18 & \begin{tabularx}{0.7\textwidth}{X} 耶和華啊，我等候你的救恩。 \end{tabularx} \\ \\ \relax
49:19 & \begin{tabularx}{0.7\textwidth}{X} 迦得必被襲擊者襲擊，他卻要襲擊他們的腳跟。 \end{tabularx} \\ \\ \relax
49:20 & \begin{tabularx}{0.7\textwidth}{X} 亞設必出豐盛的糧食，要供應君王的佳肴。 \end{tabularx} \\ \\ \relax
49:21 & \begin{tabularx}{0.7\textwidth}{X} 拿弗他利是被釋放的母鹿，他要生出可愛的小鹿。 \end{tabularx} \\ \\ \relax
49:22 & \begin{tabularx}{0.7\textwidth}{X} 約瑟是多結果子的樹枝，是泉旁多結果的枝子；他的枝條伸出牆外。 \end{tabularx} \\ \\ \relax
49:23 & \begin{tabularx}{0.7\textwidth}{X} 弓箭手惡意攻擊他，敵對他，向他射箭。 \end{tabularx} \\ \\ \relax
49:24 & \begin{tabularx}{0.7\textwidth}{X} 但他的弓仍舊堅硬，他的手臂靈活敏捷，這是因雅各的大能者的手，從那裡，他是以色列的牧者，以色列的磐石。 \end{tabularx} \\ \\ \relax
49:25 & \begin{tabularx}{0.7\textwidth}{X} 你父親的神必幫助你；全能者必賜福給你：天上的福，深淵下面蘊藏的福，以及生育哺養的福。 \end{tabularx} \\ \\ \relax
49:26 & \begin{tabularx}{0.7\textwidth}{X} 你父親的福勝過我祖先的福，直到永世山嶺的極限。這些福必降在約瑟的頭上，臨到那與兄弟有分別之人的頭頂上。 \end{tabularx} \\ \\ \relax
49:27 & \begin{tabularx}{0.7\textwidth}{X} 便雅憫是隻抓撕掠物的狼，早晨要吃他的獵物，晚上要分他的擄物。」 \end{tabularx} \\ \\ \relax
49:28 & \begin{tabularx}{0.7\textwidth}{X} 這一切是以色列的十二個支派。這是他們的父親對他們所說的話，他按照各人的福分為他們祝福。 \end{tabularx} \\ \\ \relax
49:29 & \begin{tabularx}{0.7\textwidth}{X} 他又吩咐他們說：「我快要歸到我祖先那裡。你們要將我葬在赫人以弗崙田間的洞裡，與我的祖先在一處， \end{tabularx} \\ \\ \relax
49:30 & \begin{tabularx}{0.7\textwidth}{X} 就是在迦南地幔利對面的麥比拉田間的洞裡，那田是亞伯拉罕向赫人以弗崙買來作墳地的產業。 \end{tabularx} \\ \\ \relax
49:31 & \begin{tabularx}{0.7\textwidth}{X} 亞伯拉罕和他的妻子撒拉葬在那裡；以撒和他的妻子利百加也葬在那裡。我也在那裡葬了利亞。 \end{tabularx} \\ \\ \relax
49:32 & \begin{tabularx}{0.7\textwidth}{X} 那塊田和田間的洞是向赫人買的。」 \end{tabularx} \\ \\ \relax
49:33 & \begin{tabularx}{0.7\textwidth}{X} 雅各囑咐眾子完畢後，就把腳收在床上斷了氣，歸到他祖先那裡去了。 \end{tabularx} \\ \\ \relax
50:1 & \begin{tabularx}{0.7\textwidth}{X} 約瑟伏在他父親的臉上，在他臉上哭，又親他。 \end{tabularx} \\ \\ \relax
50:2 & \begin{tabularx}{0.7\textwidth}{X} 約瑟吩咐伺候他的醫生們用香料塗他父親，醫生就用香料塗了以色列。 \end{tabularx} \\ \\ \relax
50:3 & \begin{tabularx}{0.7\textwidth}{X} 四十天滿了，就是塗香料所規定的日子滿了。埃及人為他哀哭了七十天。 \end{tabularx} \\ \\ \relax
50:4 & \begin{tabularx}{0.7\textwidth}{X} 過了哀悼的日子，約瑟對法老家中的人說：「我若在你們眼前蒙恩，請你們對法老說： \end{tabularx} \\ \\ \relax
50:5 & \begin{tabularx}{0.7\textwidth}{X} 『我父親曾叫我起誓說：看哪，我快要死了，你要將我葬在迦南地，在我為自己所掘的墳墓裡。』現在求你准我上去葬我父親，然後我必回來。」 \end{tabularx} \\ \\ \relax
50:6 & \begin{tabularx}{0.7\textwidth}{X} 法老說：「你可以上去，照你父親叫你起的誓，將他安葬。」 \end{tabularx} \\ \\ \relax
50:7 & \begin{tabularx}{0.7\textwidth}{X} 於是約瑟上去葬他父親。與他一同上去的有法老的眾臣僕和法老家中的長老，以及埃及地所有的長老， \end{tabularx} \\ \\ \relax
50:8 & \begin{tabularx}{0.7\textwidth}{X} 還有約瑟的全家和他的兄弟們，以及他父親的家屬；只留下他們的孩子和羊群牛群在歌珊地。 \end{tabularx} \\ \\ \relax
50:9 & \begin{tabularx}{0.7\textwidth}{X} 又有車輛和駕駛兵和他一同上去，隊伍非常龐大。 \end{tabularx} \\ \\ \relax
50:10 & \begin{tabularx}{0.7\textwidth}{X} 他們到了約旦河東亞達的禾場，就在那裡大大地號咷痛哭。約瑟為他父親哀哭了七天。 \end{tabularx} \\ \\ \relax
50:11 & \begin{tabularx}{0.7\textwidth}{X} 迦南的居民看見亞達禾場上的哀哭，就說：「這是埃及人一場極大的哀哭。」因此那地方名叫亞伯‧麥西，是在約旦河東。 \end{tabularx} \\ \\ \relax
50:12 & \begin{tabularx}{0.7\textwidth}{X} 雅各的兒子們遵照父親的吩咐去辦了， \end{tabularx} \\ \\ \relax
50:13 & \begin{tabularx}{0.7\textwidth}{X} 他們把他送到迦南地，葬在幔利對面的麥比拉田間的洞裡；那田是亞伯拉罕向赫人以弗崙買來作墳地的產業。 \end{tabularx} \\ \\ \relax
50:14 & \begin{tabularx}{0.7\textwidth}{X} 約瑟葬了他父親以後，就和他的兄弟，以及所有同他上去葬他父親的人，都回埃及去了。 \end{tabularx} \\ \\ \relax
50:15 & \begin{tabularx}{0.7\textwidth}{X} 約瑟的哥哥們見父親死了，就說：「也許約瑟仍然懷恨我們，會照我們從前待他一切的惡，重重報復我們。」 \end{tabularx} \\ \\ \relax
50:16 & \begin{tabularx}{0.7\textwidth}{X} 他們就傳口信給約瑟說：「你父親未死之前曾吩咐說： \end{tabularx} \\ \\ \relax
50:17 & \begin{tabularx}{0.7\textwidth}{X} 『你們要對約瑟這樣說：從前你哥哥們惡待你，你要饒恕他們的過犯和罪惡。』現在求你饒恕你父親的神之僕人們的過犯。」他們對約瑟說了這話，約瑟就哭了。 \end{tabularx} \\ \\ \relax
50:18 & \begin{tabularx}{0.7\textwidth}{X} 他的哥哥們又來俯伏在他面前，說：「看哪，我們是你的奴隸。」 \end{tabularx} \\ \\ \relax
50:19 & \begin{tabularx}{0.7\textwidth}{X} 約瑟對他們說：「不要怕，我豈能代替神呢？ \end{tabularx} \\ \\ \relax
50:20 & \begin{tabularx}{0.7\textwidth}{X} 從前你們的意思是要害我，但神的意思原是好的，要使許多百姓得以存活，成就今日的光景。 \end{tabularx} \\ \\ \relax
50:21 & \begin{tabularx}{0.7\textwidth}{X} 現在你們不要害怕，我必養活你們和你們的孩子。」於是約瑟安慰他們，講了使他們安心的話。 \end{tabularx} \\ \\ \relax
50:22 & \begin{tabularx}{0.7\textwidth}{X} 約瑟和他父親的家屬都住在埃及。約瑟活了一百一十年。 \end{tabularx} \\ \\ \relax
50:23 & \begin{tabularx}{0.7\textwidth}{X} 約瑟看到以法蓮第三代的子孫。瑪拿西的孫子，瑪吉的兒子，出生時都放在約瑟的膝上。 \end{tabularx} \\ \\ \relax
50:24 & \begin{tabularx}{0.7\textwidth}{X} 約瑟對他的兄弟說：「我快要死了，但神必定看顧你們，領你們從這地上去，到他起誓應許給亞伯拉罕、以撒、雅各之地。」 \end{tabularx} \\ \\ \relax
50:25 & \begin{tabularx}{0.7\textwidth}{X} 約瑟叫以色列的子孫起誓：「神必定眷顧你們，你們要把我的骸骨從這裡帶上去。」 \end{tabularx} \\ \\ \relax
50:26 & \begin{tabularx}{0.7\textwidth}{X} 約瑟死了，那時他一百一十歲。人用香料塗了他，把他收殮在棺材裡，停放在埃及。 \end{tabularx} \\ \\
[1ex]
\hline
\hline
\end{longtable}
$^{1}$好,在十週年裡,我們有一個講道的系列,就叫做「NEXT」.
我們這次是第二講.
上一次我們講到阿布拉罕的故事.
我們說整本聖經,由《創世記》到《啟示錄》.
去到介紹一位上帝是怎樣的上帝呢?.
這位上帝是「Guardian Mission」.
是「Missionary God」.
是一個宣教宗的上帝.
上帝沒有放棄過這個世界.
雖然這個世界是亂七八糟.
我們生命亂七八糟.
但是上帝一直沒有放棄過我們.
祂的「Mission」就是去拯救我們.
拯救世人和拯救整個的地.
我們這次開始的時候,用這張圖片來開始.
這個是在巴黎聖母院的天主教教堂裡面拍攝的.
大家有沒有去過?.
或者見過這樣的影像.
就是「Stained Glass」.
教堂裡面有一些特別傳統的教堂.
他們的裝修或者他們的裝置.
都會有一些彩色的玻璃窗.
「Stained Glass」.
陽光照下來的時候,就會很漂亮的.
有沒有印象?.
你知不知道我們教會有的?.
其實我們沒有的.
我們沒有「Stained Glass」.
因為我們不屬於那個世代.
但是這個世代,就是以前的教堂,古代教堂的那個世代.
他們就有一種這樣的裝置.
話說這個裝置就是十二世紀開始的.
當時有個神學家他認為教會的牆.
其實應該要有光照進來.
因為光就代表什麼?.
代表上帝,是不是?.
代表上帝.
所以他就將這些牆拆了,然後弄一些窗上去.
當時也有一些彩色的玻璃開始技術成熟.
所以就用這些新的技術來做一些「Stained Glass」.

$^{41}$很多時候就弄一些標誌,或者弄一些聖經的故事.
或者一些圖畫等等.
就給信眾去到教堂的時候,他能夠看到聖經的一些信息等等.
「Stained Glass」告訴我們每一個碎片.
是一起放在一起才會組成整幅圖畫.
那個「Stained Glass」每一塊的碎片.
它自己只是一個很小的角落.
它需要其他的碎片來和它一起組織.
並且背後有一個工匠.
然後工匠很巧妙地將「Stained Glass」不同的顏色,圖案等等.
放在一起才會有一幅圖畫.
為什麼用這個圖畫來開始呢?.
因為我們繼續看創世記的故事的時候.
我們發現創世記的故事,特別是那些所謂族長的故事.
亞伯拉罕,以撒,雅各,約瑟,猶大等等.
由亞伯拉罕開始,直到創世記最後的時候.
有十二眾子,一個大的家族出現.
原因是因為上帝要藉著亞伯拉罕來做一個任務.
使萬族能夠因祂得福.
這個上帝,永遠是在任務上的上帝.
就是將這些破碎的生命組成在一起.
使世界和祂身邊的人能夠得到祝福.
我們上一次說到在這個系列開始的時候.
聖經介紹這個上帝是一個宣教宗的上帝.
面對著一個被罪破壞的世界上帝.
沒有放棄,並且願意去祝福.
所以呼召亞伯拉罕出來.
這一段的經文就是創世記第十二章.
最後那句就是最重要的.
他說地上的萬族都要因你亞伯拉罕或者亞伯蘭來得福.
就是上帝要藉著他來施行.
他自己在寶世上祝福萬族的救恩計劃.
使萬族能夠去認識他.
上一次我們聚焦亞伯拉罕的故事.
亞伯拉罕的故事.
我們說亞伯拉罕並非打不死的英雄.
其實是一個很真實和很脆弱的人物.
今天我們的鏡頭就較闊一點.
我們就看亞伯拉罕之後的後代.
到底發生了什麼事.

$^{81}$如果你有一點印象的話.
你記得亞伯拉罕的兒子是誰嗎?.
當然有一個以實瑪利.
有以實和埃塞.
埃塞是不同的一代.
記載了亞伯拉罕的故事之後.
就講到以薩的故事.
以薩的妻子就是尼伯格.
他們也是不辱的.
人到盡頭的時候.
他們不能夠繼續有後代.
不能夠成為一個大國的時候.
上帝就插手去介入.
有一個所謂的Miracle Child出現.
這個Miracle Child其實是雙倍的.
就是以素和雅各.
你記得嗎?.
他們是雙胞胎.
從小到大就是在爭奪東西.
爭奪什麼呢?.
就是爭奪產權.
你看這些族長的故事.
幾乎是在看何鴻燊家族的故事.
到底誰爭奪了多少.
爭奪了多大塊等等.
這個家族的故事.
其實就是在講這些後代.
怎樣去處理他們有的所謂Inheritance.
他們的基業.
以薩和尼伯格.
他們有的雅各,以素和雅各.
就是兄弟之間相爭的.
你記得在雅各和以素的故事.
是不是有兩次.
這個雅各去欺騙和奪取.
這個所謂長子的名分.
你記得嗎?.
記得嗎?.
記得吧?.
如果你沒看過聖經.

$^{121}$當然你不會有印象.
但是如果你稍微有查過經.
或者稍微聽過.
你就知道他們經常爭奪祝福.
在近東以前的社會裡面.
長子.
永遠都是得到的產業是雙倍的.
就是有個蛋糕.
家父有的遺產.
好像一個蛋糕一樣.
就切開三塊.
如果有兩個兒子就切開三塊.
怎樣分呢?.
小兒子就分三分之一.
長子就分三分之二.
OK嗎?.
66\%和33\%.
長子永遠都是雙倍的.
所以他們就很緊張那個所謂長子的名分.
First born right of inheritance.
我在香港長大的.
雖然我沒有那種香港長大的背景.
但是偶爾聽到一些這樣的故事.
譬如香港有一些叫丁屋的故事.
丁權.
什麼叫丁權?.
就是那個產業是按誰來分的?.
按丁來分的.
丁就是兒子.
就按兒子.
如果你是女兒的話.
不好意思.
你就沒份的.
因為女兒是嫁出去的.
不屬於你的.
兒子留下來.
所以那些丁權的遺產的爭霸.
就有一點像這個族長的故事裡面的背景.
長子名分既然是收到雙倍的話.
產業權是雙倍的話.

$^{161}$這個一定要爭.
所以雅各是雙胞胎小的那個.
他就經常想維持什麼?.
大那個所謂的長子的名分.
你記得他欺騙他兩次.
我不知道你從小到大聽這些故事的時候.
為什麼會欺騙他兩次呢?.
你欺騙他一次.
你記得紅豆湯的故事.
這個耳素出外打獵.
然後肚子餓回來的時候.
餓得不得了.
然後雅各就弄了一碗粥給他.
是不是?.
很便宜的.
然後說如果你肚子餓.
你吃吧.
你吃的話就跟我簽一簽那個 share.
剛才說股票.
就是那個股份.
你簽給我.
OK.
我就拿66.
你就拿33.
這樣好不好?.
這個耳素就很流行看這些.
可能是世外高人打獵.
他就很流行看這些.
算了我拿33.
你拿66.
可能中間有這樣的交易.
所以雅各第一次騙了這個耳素.
就拿了66\%.
拿了三分之二的股份.
然後家父死之前.
就是以薩死之前.
他要祝福他的兒子.
那裡又有另外一個所謂 deception(欺騙)的 story(故事).
就是騙祝福的故事.
你記得應該是家父以薩去祝福那個大兒子耳素.

$^{201}$結果雅各就假扮耳素.
假扮他的哥哥.
然後拿那個祝福.
是不是?.
奇怪了.
應該是臨終的時候.
家父臨終的時候.
還要騙他一次.
之前已經騙了.
你拿了66\%的股份.
那個應該是你的.
是不是?.
原來是這樣的.
在舊約裡面的背景.
那個所謂的長子名分.
有分兩個層次的.
第一個層次可以叫做長子名分.
就是66\%.
你看到受益人按道理應該是長子.
結果小兒子雅各就騙了.
拿了雙倍的股份.
就是產業權的雙份他拿了.
性質就是他得到了.
用什麼方法?.
就是用紅豆湯的交易.
但是到最後的時候.
家父的祝福.
那個就不是產業權的問題了.
那個是管理權的問題.
管理權的問題.
本身理所當然有管理權.
就是所謂的Trustee.
之前那個就是你的股份.
你的Portion.
然後最後最後決定的就是Trustee.
那個管理權.
Administration.
那個行政權.
就是執行權.
就是在家父那裡.

$^{241}$理所當然應該是.
給Firstborn長子的.
但是雅各就用欺騙的方法.
去爭取了.
二數的管理權.
管理權裡面有什麼呢?.
為什麼這麼重要要爭呢?.
原來管理權.
Administration.
就是那個Execution.
就是執行上.
什麼時候分紅.
什麼時候分Share.
什麼時候分利息.
是那個管理權去負責的.
所以就算你得到多少的產業也好.
多少部份也好.
都是要靠那個管理權.
來去施行的.
大家記不記得?.
類似這樣.
你可以想像.
譬如一個大家族的企業.
就是阿巴罕爾薩雅各.
他們有很多產業的.
你可以想像一個大的企業.
就是父親分給你.
當你這邊是長子.
分給你66\%.
然後這邊分多少?.
33\%多.
是不是?.
三分二在這邊.
三分一在這裡.
但是呢.
你分到多少也好.
你分到三分二也好.
但是執行權CEO在這邊.
CEO就負責什麼時候分.
他可以拖到你最後.

$^{281}$才分給你.
所以你拿著的股份.
其實就白紙而已.
記不記得?.
記得吧?.
或者我這樣舉個例子.
我昨天想到這個例子.
因為吃飯的時候.
就跟子女聊天.
假設我的退休金.
那些湊一湊.
那我有一億.
OK.
有一億.
如果我有一億就好了.
當我沒有一億.
當我有一億.
那我女兒有個兒子.
你不要分男女.
總之就是.
兒子就分三分之二.
多不多?.
都挺多.
女兒就分三分之一.
女兒當然不服氣.
她就騙了她第一次.
騙了她弟弟第一次.
就拿了三分二.
但是始終執行權在誰那裡?.
就在兒子那裡.
OK.
就在兒子那裡.
就是不要讓他長子賤子.
賤子.
總之執行權就是.
因為重男輕女.
是不是?.
所以執行權始終都在我兒子那裡.
這個是例子.
就是這樣分.

$^{321}$一個分三分二.
一個分三分一.
臨死之前.
開個你講句.
如果我在ICU那裡.
上氣不接下氣的時候.
那應該兒子呢.
他有三分之二.
他應該.
他就馬上來探望我.
他說老爸.
你那個Online account.
那個password在哪裡?.
password是什麼?.
沒有人知道.
只有你知道.
你現在還有膽子.
你告訴我那個password.
因為那個password就是執行權.
你就可以login.
然後可以.
將那些股份來cash了它.
是不是?.
怎樣cash你才會有我的password.
好壞不好聽.
如果是.
女兒截股的話.
女兒早一個小時來探望我.
然後假扮我的兒子.
然後騙了我.
拿個password來.
老爸.
你快死了.
你沒氣了.
快要拔喉了.
我就糊裡糊塗.
半夢半醒.
我就半昏迷狀態.
就告訴我女兒.
那個password.

$^{361}$然後之後一個小時.
我兒子才真正趕來.
老爸.
你那個password是什麼?.
我就跟他說.
我剛才不是給了你嗎?.
我說那不是你.
原來是你姐姐.
你快點給我password.
我給了他password.
給password兩次.
他馬上login.
login不了.
他為什麼login不了?.
當然了.
他姐姐拿了之後.
那三分鐘馬上改了password.
就是執行權的問題.
所以為什麼雅各一定要騙他兩次呢?.
因為你不拿到執行權的時候.
你實際上是拿不到那些所謂的股份的.
你看到這個故事裡面就很搞笑.
所以.
以素發現自己兩次被雅各欺騙他的祝福的時候.
他大哭然後說什麼呢?.
二十七章三十六節.
以素就說.
他名叫雅各.
雅各的意思就是說.
就是拿著不放的意思.
OK?.
拿著不放的意思.
他叫雅各豈不是這樣.
是真的.
他的名字真的叫Jacob.
是這麼壞的.
是這麼狡猾的.
他欺騙了我多少次呢?.
兩次.
先前拿了我的股份就是長子的名分.

$^{401}$後來他又奪了我的福分.
拿了我的執行權.
於是你發現發生什麼事?.
就起殺機是嗎?.
是.
就等於如果我兒子發現了這件事的時候.
剛才那個例子.
他一定會殺他姐姐的.
他一定會拿他的命.
這個故事就是講到他們之間的爭奪.
這些是非常人性的黑話.
但是他在聖經的族長的故事.
就是神的子民的.
Founding Fathers.
聖經是這樣記載的.
聖經在這裡告訴我們這些所謂的黑歷史.
他們之間這樣去爭奪上帝的祝福.
好像上帝的祝福是要爭回來一樣.
用自己的手段來拿回來的.
這些絕對不是一些很可愛的故事.
我經常都講聖經的故事絕對不是給小孩子看的.
是真的.
是一些聖人的故事.
所以結果的後果是怎樣呢?.
以素就凍殺基.
然後亞國就流亡了在米索不達米亞二十年的時間.
然後才有機會再回去.
這個是第二個所謂族長的故事.
現在我們再看看.
以索看完了.
我們再看看亞國的妻子.
亞國在米索不達米亞有兩個老婆.
第一個叫做莉亞.
第二個叫做拉傑.
他比較愛誰呢?.
他比較愛拉傑,Rachel.
為什麼還有另外兩個女人的名字呢?.
就是因為這兩個老婆.
他們爭什麼呢?.
丁權嘛.

$^{441}$所以就爭誰的生子生得多.
你生得越多你的share就越多.
因為按丁來分的.
所以大家在不育的當中.
他們就想去生.
然後生不了的時候.
以前古代的習俗就是叫屎女來跟你生.
因為工人生的都是屬於你的產業.
所以莉亞就找了西柏.
西柏.
另一個他的工人.
去生兒子.
然後拉傑就見到自己不育.
然後就找了碧拉去幫他生.
結果生了多少個出來呢?.
十二個.
再加上女兒.
女兒不是很漂亮.
她叫底娜.
十三個.
總之生了十三個.
這個創世記二十九章.
尾的時候.
去到交代.
但是你看到那個故事是在說什麼呢?.
就是一幫很破碎的人.
他們去到爭奪一些世上的東西.
爭奪所謂上帝給他們的產權.
在質度裡面去爭奪.
在計算裡面去爭奪.
然後才有十二支派出現.
然後第三個故事就是大家很清楚了.
這個就是約瑟和眾兄弟.
大家的糾纏.
因為雅各慢慢地繼續下去.
由以薩去到雅各.
阿巴罕,以薩,雅各.
雅各就有十二個兒子.
他最愛的是誰?.
他偏愛的是誰?.

$^{481}$偏愛的就是約瑟.
所謂的拉仔.
最後的時候生的約瑟.
所以因為偏愛他的緣故.
那些兄弟就不喜歡他了.
這個不純粹是偏愛的問題.
如果你想想產業的問題.
就是為什麼呢?.
因為那個 password(密碼)在哪裡?.
那個 password(密碼)在他最愛的兒子那裡.
是不是?.
也將會給他的.
所以你說為什麼那些兄弟這麼討厭他呢?.
討厭他嗎?.
如果你是其中一個兄長.
特別你是長子的話.
我當然不妥他.
我分到多少都好.
那個 password(密碼)在他那裡.
是不是?.
是他負責.
是他有權.
他什麼時候分給我.
到我死那天.
可能他都還沒分給我.
所以他們的兄弟就很想殺約瑟.
這裡有幾段說話.
這個故事一開始的時候.
就講到這些兄長怎樣對約瑟.
他就說哥哥們.
三十七章第四節.
見父親愛約瑟過於他們的時候.
就惱恨約瑟.
就不與他說友善的話.
其實原文就是 shalom.
為什麼這麼重要呢?.
因為希伯來人到今天.
以色列人.
他們大家去到問安的時候.
他們講什麼呢?.

$^{521}$shalom.
講 shalom.
所以其實這裡不是說不講友善的話這麼簡單.
這裡簡直是hello都不會講.
就是你每天見到約瑟的時候.
如果你是兄長的話.
你就問他什麼.
where is the password(密碼).
I want the password(密碼).
OK.
那個 password(密碼)應該在我這裡.
但是他就不講.
是不是?.
你說吹什麼我不講給你聽.
是不是?.
然後你看到每一個的故事.
每一個兄長.
第一個.
最大的長子.
就是Rubin.
裡面.
你看到他們在爭家產.
為什麼呢?.
因為當拉傑.
就是所謂第二個老婆.
死了之後.
然後他就和碧拉.
同寢.
就是拉傑的侍女同寢.
按照當時近東的社會.
去到和那個媽媽的侍女同寢.
是什麼呢?.
我們現代人就會想.
一定是性的問題.
其實不是的.
是爭家產的問題.
就是如果和他侍女同寢的話.
他所生的那兩三個小孩.
那些的丁權.
就是誰的?.

$^{561}$當然是他.
現在我是你爸爸.
是不是?.
你明白我的意思嗎?.
三十五章第二十二節.
然後到第二第三仔.
就是Simeon和Levi.
西敏和利美.
他們就是滿有殺機的人.
那三十四章講了一個故事.
就是他們殺了另外一個外邦人的城市.
這個事件.
裡面的人.
來將他們通殺.
就用了割禮這件事.
來設一個局.
那個詳細我們不講了.
就是主要是想交代的一個方面.
就是這兩個人.
是滿有殺機的.
是violent people.
第四個就是猶大Judah.
猶大這個故事.
我們就更加清楚了.
三十七章.
他就是想賣掉藥室的.
其實為什麼要將他賣掉呢?.
不是只是憎恨他那麼簡單.
而是因為他有password.
賣了你看看還有沒有password.
對嗎?.
當你死了.
然後password應該理所當然給誰?.
給誰爭到那個所謂.
長子丁權的那個人.
將那個競爭者剷除.
那你就可以拿到.
其他的兄弟做什麼?.
就丟藥室在坑裡面.
丟在坑裡面.

$^{601}$然後做什麼?.
就坐下吃飯.
坐下吃飯.
你就見到這個就是.
以色列.
神的子民.
那個祖先的故事.
兩個字.
嘉辯.
為了什麼?.
為了利益.
是吧?.
為了利益.
這些我們看完的故事.
我們會發現.
一點都不光榮.
對嗎?.
一點都不光榮.
跟我們平時看聖經的印象.
不是這樣的.
我們覺得放在聖經裡面的.
都是一些比較神聖的東西.
是吧?.
神聖的故事.
但卻是.
聖經的作者.
不寫一些光榮的東西.
然後偏偏選一些.
所謂黑歷史.
它中間有欺騙.
是一些破碎的人.
是一些.
是一些會出賣大家.
以及會彼此相憎的一些人.
為什麼會這樣呢?.
你就可以問一個問題.
到底聖經是不是一些普通的.
道德的教訓的書呢?.
但很多人.
特別一些可能.

$^{641}$以前在香港長大的一些朋友.
都是在教會的學校.
長大的.
是吧?.
小學啊,中學啊.
為什麼我們經常喜歡將我們的小朋友.
發送到教會的學校.
天主教,基督教的學校.
為什麼呢?.
你們也是.
是吧?.
因為他們教得好.
是吧?教得好.
因為德育啊.
discipline(教育)這些東西.
通常在教會的學校是比較好的.
但是你們看到.
這個當然是對的.
但很多時候教會學校裡面.
你學的信仰或者看聖經的故事.
是不會看這些的.
是吧?.
這些太黑了.
其實很難看的.
是吧?.
沒有人解釋給你聽.
你怎麼知道發生什麼事呢?.
是吧?.
所以我們通常看的都是比較溫馨的東西.
是吧?愛人如己啊.
怎樣愛人.
怎樣憐恤人啊等等.
但是你看到聖經的作者.
很多時候就是將這些的部分.
不提這麼多.
然後提一些比較黑的地方.
他的焦點不是教我們做好人.
聖經的焦點不是做這件事.
聖經的焦點是藉著這些黑歷史.
甚至是神的子民的祖先.

$^{681}$就是我們的創造父親.
去講述他們的黑歷史的時候.
就是藉著這些黑歷史,黑故事.
去告訴我們.
這位上帝是一個怎樣的上帝.
你看完這些故事.
你真的是搖頭了.
Just shake my head.
這群人如果你是神的話.
如果你說要藉著這個家族.
來使萬國得福的話.
你真的是.
搖頭了.
你一定會很生氣的.
確實在近東的宗教的概念裡面.
神明呢.
看到那些人不聽話的時候.
他們會做什麼呢.
神明都會有殺機的.
神明就會殲滅這些人.
就是禁個reset.
全部東西重新來過.
是不是?因為這群人太吵了.
太麻煩了.
就是在近東的那些宗教的概念裡面.
他們的神學裡面.
暴力其實是世界的宇宙的一個規律的其中一樣東西.
由始至終都會有暴力的.
由始至終都會有一些violence出現的.
這是一個正常的工具.
你不喜歡那些人.
你就將他全部消滅.
然後再重新來過.
是不是?.
你看到聖經所表達的上帝.
就不是一個這樣的上帝.
是截然不同的.
上帝看到這些人搞到亂七八糟的時候.
看到他們很多破碎的時候.
上帝他選擇了一條.

$^{721}$not the easy way.
not the easiest way.
the easiest way is to reset.
但是上帝去救贖.
上帝救贖一些破碎的人.
破碎的人生.
然後要將他們重新再組合.
使得這幅圖畫能夠呈現出來.
這樣做上帝絕對不是洗白那些哥哥.
或者其他族長的惡事.
惡人始終做的事情是evil.
但是上帝是可以克服人的軟弱.
克服人的惡意.
最終是成就他自己的計劃.
這個就是故事的目的.
就是想告訴我們theology.
告訴我們上帝是一個怎樣的上帝.
所以這個故事創世記最後41章至50章.
說到約瑟被賣到埃及的時候.
然後在高高低低的當中.
他就成為了埃及的宰相.
最後發生什麼事呢.
在饑荒的裡面.
約瑟藉著他的苦難.
藉著哥哥在他身上行的惡事.
他就成為了宰相.
然後41章第57節.
各地的人或者甚至全地的人在故事裡面.
全地的人都去了埃及.
向約瑟買糧.
因為天下的饑荒甚重.
弟兄姊妹之間他們再相認之後.
一個很感動的故事.
那些哥哥終於所謂認錯了.
大家有reconciliation(認同).
45章第8節約瑟說什麼.
猜我來的不是你們.
你們原本是惡意使我被賣到埃及的.
我受了很多苦.
不過我現在知道上帝.

$^{761}$is on a mission(有一個任務).
他的mission是猜我來這裡.
猜我來埃及.
然後保存著雅各家.
並且祝福各地的人.
是不是.
50章第20節了.
最後一章.
他說從前你們興獎你們的義士.
是要害我.
但是上帝的義士.
原是好的.
要使許多人.
萬族能夠得以存活.
成就今天的光景.
這裡你就看到.
老祖耶和華對阿巴喊的呼召.
那個的應許就是萬國要因你得福.
這裡就慢慢地實現出來.
這個就是.
族長的故事.
或者可以這麼說.
這個是教會的前輩.
上一代,第一代的故事.
數他們的黑歷史出來.
不是要描黑他們.
而是想表達清楚這位上帝.
我們所信的上帝.
我們所讚美的這位上帝.
是一個怎樣的神.
而我們是一群怎樣的人.
我們是需要祂的恩典.
我們是一群很脆弱.
很不行的人.
族長是虧缺上帝的榮耀.
我們每一個都是一樣.
神才是最厲害的.
在歷史歷代裡面.
教會都會有不同的會議會開.
在新約的教會.

$^{801}$到現在大概有二千年的歷史.
久不久有一些所謂的church council.
或者大公的教會.
普世的教會.
他們大家坐在一起.
去解決一些神學問題.
或者想寫出.
或者總結出.
到底我們的信仰.
應該怎樣了解呢.
這些所謂叫信條或者信經.
所以在過去二千年裡面.
有很多這樣的信經是產生的.
其中一個信經叫做Belgic Confession.
比利時信條.
在第二十七條那裡.
這個十六世紀的時候有的信條.
我覺得他描述教會是描述得很好.
他講什麼呢.
他說教會.
自世界之初就已經存在了.
在創世記的時候已經存在.
教會的初葉.
也都將會存在世界末了的時候.
因為基督是永恆的君王.
神的子民不可或缺的.
只要有基督就要有基督的子民.
就會有教會.
這個所謂神聖的教會.
是蒙上帝的保守.
免受世上一切的侵擾的.
即便在世人的眼中.
教會或許顯得渺小.
甚至瀕臨消亡.
extinct.
即便在世人的眼中.
或許是顯得渺小.
或者瀕臨消亡.
我覺得這句話.
點出教會的神學.

$^{841}$是點得很通的.
教會是很脆弱的.
神的子民都是很脆弱.
由第一代開始.
到現在.
為什麼呢.
不是因為我們可以.
而是因為上帝可以.
不是因為我們沒有黑歷史.
不是因為我們可以.
我們都是規矩上帝的榮耀.
但唯獨是基督.
基督是我們主的.
基督去拯救我們.
族長的社潮.
教會的社潮.
就是這樣.
我記得以前二十世紀的時候.
有一個很出名的天主教神學家.
Hans Kuhn.
他曾經說過類似這樣的說話.
他說不是教會守住真理.
而是真理守住了教會.
是福音的真理守住了教會.
你回去看回這張照片.
所謂Stained Glass.
每一個碎片.
不是單獨的.
是最後有一個工匠.
將這些碎片放在一起.
然後成為一個很偉大的圖畫.
看族長的故事.
甚至看我們自己的故事.
我們都會覺得.
我們真的很差勁.
我們真的很不行.
蓄意的破壞.
或者我們會不在乎的心態.
來破壞上帝的工作.
但上帝是可以剋聖我們所有人.

$^{881}$我們每一個人.
我們的軟弱和脆弱的地方.
今天有三個邀請.
我們說過這個系列.
我們會有一些東西給大家.
而這個實物.
這個小的物件.
是給你一個提醒.
提醒你上帝的恩典和上帝的工作.
上一次開始我們有一個瓶子.
裡面有沙.
你記得我們說過亞伯拉罕的沙的所謂應許.
多如海沙.
你的後裔會多如海沙.
這次你出去的時候.
你可以拿一塊拼圖.
我們當然不會給你玻璃.
不會給你玻璃碎片.
但這個拼圖.
好像碎片一樣.
它會告訴我們.
原來我們每一個人都是一塊碎片.
是不完整的.
我們需要上帝的恩典.
將我們拼在一起.
第一個邀請就是一個謙卑的邀請.
我們放下一些有的錯覺.
或者那種錯覺就是我夠好的話上帝就會用我.
原來上帝一直都不需要等我們好.
才能用到我們.
由第一代開始到現在.
上帝不需要等完美的人.
來繼續他的計劃.
上帝有的能力是高過我們的不善.
我們的不足.
第二個邀請.
是一個醒覺的邀請.
醒覺的邀請.
問一問你自己.
我有沒有用一些旁觀者的眼光去看上帝.

$^{921}$看教會.
上帝既然有他的使命.
去宣教.
教會就自然有同一個使命.
我們未必會彼此爭兵權.
我們未必會彼此爭password.
但我們可能都會有一種.
關我什麼事的心態.
等於十二個兒子.
不是每個人都出主意出賣弱勢.
有很多都是旁觀者.
他們就是關我什麼事的心態.
我想這才是最根本的問題.
就是上帝的子民很多時候說.
與我何干.
這種心態.
這種應該是一種醒覺的邀請.
你告訴我上帝繼續在做他的工作的時候.
為何我們還會說是關我什麼事.
第三個邀請就是參與.
我們就是一個碎片.
我們可不可以教我們自己.
奉獻我們自己.
在上帝公像的手中.
不是因為我夠好.
而是我自己不完美也好.
上帝的能力是高過我的.
整幅圖畫我們就見到上帝的恩典.
上帝自己的大能.
等我出去的時候拿一塊.
拿回你的瓶子.
如果你今天是第一次來的話.
你不妨也拿一個瓶子.
放在當下的地方.
放這塊碎片進去.
去告訴你自己.
我們需要大家.
需要大家繼續參與上帝自己的工作.
上帝你拼我自己.
不完美的一塊.

$^{961}$放到你自己的圖畫裡面.
使你宣教的工作.
使更多人認識上帝.
使這個工作能夠繼續延伸.
好嗎?.
我們同心協力地祈禱.
因為你是.
那位完美的神.
你是去不斷地宣教.
不斷地去救贖人.
我們能夠成為你.
被救贖的子民的一部分.
赦免我們的罪.
我們的不是.
我們的脆弱.
我們的不好的地方.
讓我們同心協力地來到你面前.
主要是我們願意將我們的生命拼進去.
你一個救贖普世的計劃當中.
叫人藉著我們的生命.
我們的生活.
更加認識耶穌基督愛我們的主.
今天我們都是聖餐的主人.
這個聖餐.
這個餅.
這個杯.
就是告訴我們.
完全因為你自己的恩典.
我們能夠成為你自己的門徒.
成為你的教會.
主要是這個恩典要提醒我們.
活在恩典裡面.
站穩在恩典裡面.
讓你的名.
能夠得以被稱頌.
我們祈禱奉祖彌陀佛.
\newpage



\section{詩篇 36:1-9}
\label{sec:032099}
\textbf{人的邪惡,神的美善 (華)}~~~~張磊牧師
\newline
\newline
連結: \href{http://www.church.com.hk/acms/content.asp?site=cdc\&op=show\&type=product\&code=032099\&layout=sermon}{\texttt{032099}} ~~~~ 語音日期: 2026-05-03
\newline
\newline
\hyperref[sec:032098]{\small{< < < PREV SERMON < < <}}
~
\hyperref[sec:index_chronic]{\small{[返順時目]}}
~
\hyperref[sec:032100]{\small{> > > NEXT SERMON > > >}}
\newline
\newline
詩篇 36:1-9
\newline
\begin{longtable}{cl}
\hline
\hline
章節 & 經文 (和合本修訂版)\\
\hline
36:1 & \begin{tabularx}{0.7\textwidth}{X} 過犯在惡人的心底向他說話，他的眼中不怕神。 \end{tabularx} \\ \\ \relax
36:2 & \begin{tabularx}{0.7\textwidth}{X} 他自誇自媚，以致罪孽無法察覺，不被恨惡。 \end{tabularx} \\ \\ \relax
36:3 & \begin{tabularx}{0.7\textwidth}{X} 他口中的言語盡是罪孽詭詐，他不再有智慧，也不再行善。 \end{tabularx} \\ \\ \relax
36:4 & \begin{tabularx}{0.7\textwidth}{X} 他在床上圖謀罪孽，定意行不善的道，不憎惡惡事。 \end{tabularx} \\ \\ \relax
36:5 & \begin{tabularx}{0.7\textwidth}{X} 耶和華啊，你的慈愛上及諸天，你的信實達到穹蒼， \end{tabularx} \\ \\ \relax
36:6 & \begin{tabularx}{0.7\textwidth}{X} 你的公義好像高山，你的判斷如同深淵；耶和華啊，人民、牲畜，你都救護。 \end{tabularx} \\ \\ \relax
36:7 & \begin{tabularx}{0.7\textwidth}{X} 神啊，你的慈愛何其寶貴！世人投靠在你翅膀的蔭下。 \end{tabularx} \\ \\ \relax
36:8 & \begin{tabularx}{0.7\textwidth}{X} 他們必因你殿裡的豐盛得以飽足，你也必叫他們喝你那喜樂的泉水。 \end{tabularx} \\ \\ \relax
36:9 & \begin{tabularx}{0.7\textwidth}{X} 因為在你那裡有生命的泉源，在你的光中，我們必得見光。 \end{tabularx} \\ \\
[1ex]
\hline
\hline
\end{longtable}
$^{1}$弟兄姊妹平安.
感恩.
我們能夠每個周在一起聚集.
分享神的話.
照著我們今年教會的.
每天的晨更導讀.
我們仍舊停留在詩篇的分享上.
如果.
如果我.
讓我介紹給弟兄姊妹.
十本書的話.
這一定是其中的一本.
世界觀的故事.
這是我個人認為寫得非常好的一本書.
作者是.
大家知道水門事件吧.
尼克松總統的那位特別顧問.
就是寇爾森.
那他在獄中.
因為這個事件下台.
他在獄中倖主之後呢.
他寫了這本書.
非常有份量的.
一本基督教的刊物.
那在這本書當中.
他提到一個觀念.
就是今天.
他看到今天在整個西方.
基督教信仰的.
示威衰弱.
是發生在價值觀層面的.
一個失敗.
在價值觀層面上.
我們沒有抵禦.
這個世界的價值觀.
這個世俗的價值觀.
他認為價值觀的失敗.
就是今天基督教信仰的失敗.
在.
十多年前.

$^{41}$第三屆在南非開的.
洛桑會議當中他們有一個調查.
你知道當時全球.
大約三千多個基督教界的.
領袖與會.
那有一個機構在他們三千個.
領袖當中問他們一個問題.
今天你們在帶領.
教會當中帶領弟兄姊妹當中.
面對的最大的壓力和挑戰.
是什麼.
有百分之二十多的人他們認為是.
伊斯蘭教在全球的.
擴張.
當然伊斯蘭教是我們在宣教路上.
最大的一個藍組但是你知道.
這是number two number one是什麼嗎.
number one不是一個層面的.
有百分之七十的.
這些的教會領袖.
回答是世俗主義.
或者是.
世俗世界的世界觀.
世俗主義.
不是我們宣教的時候.
給我們的衝突.
而是我們這些.
特別是這些西方.
據為基督教國家的.
這些基督徒們.
基督徒的團體當中.
已經開始從內部.
被這個世界的思想.
所滲透.
我曾經聽過一個宣教士的.
素質見證讓我真的.
非常的驚訝.
他是在一個典型的伊斯蘭教國家.
當中長時間的宣教.
他說各位.

$^{81}$你千萬不要以為在伊斯蘭教國家.
阿薩特阿拉伯.
土耳其.
聯合酋長國.
你以為他們都是伊斯蘭教.
每個人都是很堅定的.
伊斯蘭教徒不是的.
其實他們的年輕人很大程度上.
也受到世俗主義的衝擊.
他們的信仰也都.
很膚淺很表面甚至.
幾乎要垮塌.
他們的伊斯蘭教信仰.
但是你知道他說什麼嗎.
你知道那些國家的人很多.
他想要移民到歐洲.
或者移民到北美.
這些第一世界國家.
Once他們來到.
這些所謂的.
文明的基督教國家.
這些原本.
伊斯蘭教信仰已經很.
動搖的這些的年輕人.
他們的伊斯蘭教信仰會再.
一次堅定起來.
你說奇怪了.
不是吧應該是在那邊.
很難信主啊應該是.
來到美國之後更容易信耶穌.
不是哦他說.
來到美國來到加拿大來到.
這些西方所謂的.
基督教國家他們原本.
很膚淺的伊斯蘭教信仰.
再一次.
再次堅固起來了.
你知道為什麼嗎.
因為他們發現.
這些所謂的基督徒.

$^{121}$他們的道德他們的價值觀.
是如此的敗壞.
是如此的荒謬.
他不羨慕基督徒的生活.
基督徒沒有一個見證.
讓他們羨慕反而讓他們.
激發他們哦原來他們.
老祖宗所信的伊斯蘭教.
才是真理才是真道.
今天北美的教會有多少.
已經和同性戀妥協的.
今天北美的教會.
有多少已經世俗化的.
當他們來到這樣的一個.
西方世界各位.
他們不是來被福音感動.
他們是被我們絆倒了.
世俗主義.
一個最基本的特徵.
用兩個字來表達叫做去昧.
一個去昧的時代.
一個不再相信超自然的.
一個不再相信.
形而上的.
一個不再追求永恆的.
一切都是在日光之下.
用趙天恩牧師.
他的一句話我非常喜歡這句話.
他說基督教信仰是一個.
統攝性的文化.
什麼叫做是統攝性呢.
我們的信仰什麼叫做.
統攝性呢就是基督教的.
信仰是一個整全的.
世界觀.
統攝性的文化就意味著.
這個信仰不是星期天.
不是在教堂裡.
而是無時無刻.
充斥在我們.

$^{161}$所有的日常生活.
一切的層面和領域當中.
曾經有一個基督徒.
我在大陸的時候遇到他.
年紀輕輕.
信了耶穌而且很堅定.
他告訴我說.
他很高興的有一天告訴我說.
趙牧師我告訴你.
基督教信仰已經成了我生活當中.
不可少的一部分.
你知道他什麼意思嗎.
他每個星期天都去教會.
每個星期五都去傳契.
不可少了已經成了固定模式了.
我稍微校正了他一下.
我說我們這句話.
可以反過來講.
不要說信仰成為.
你生活不可少的一部分.
而是你的生活.
要成為信仰的一部分.
星期一是不是你信仰的一部分.
旅遊是不是.
你信仰的一部分.
購物是不是你信仰的一部分.
工作考學.
是不是你信仰的一部分.
基督教的信仰有沒有成為一個.
整全的世界觀.
統設在統治在.
你整個的生活領域.
還是只在我們聚會的時候.
若是不然.
其實我們的.
世界觀已經.
世俗化了.
已經失去了.
那個最初的見證.
你我的重生得救.

$^{201}$可能是一瞬間的一個扭轉.
但是我們的信仰.
要經過一個漫長的.
世界觀的更新.
承認吧.
我們需要每一天甚至一生之久.
我們需要不斷的被更新.
Otherwise.
真的是逆水行舟.
我們就一定趨向世俗.
但是世界觀不是基督徒才有的.
每一個信仰.
甚至沒有信仰的人他們有他們所謂的世界觀.
我給大家念一下.
中國的黨報.
發表了一篇評論員文章.
這篇文章怎麼說呢.
這篇文章說.
公民有信仰宗教的自由.
但共產黨員.
要做堅定的.
馬克思主義無神論者.
嚴守黨章黨綱.
堅定信仰理念.
牢記黨的宗旨.
不信仰宗教.
不參加宗教活動.
這是由黨的性質和.
共產黨的世界觀.
所決定的.
人家的世界觀.
是決定了他能做某些事情.
不能做某些事情.
如果你是一個堅定的共產黨員的話.
那麼如果你是一個堅定的基督徒的話.
請問.
有哪些事情人家可以做.
你不能做.
有哪些事情你跟這個世界分別出來了.
但是呢.

$^{241}$我也給大家.
其實這是在很多的紀錄片影片當中.
我們都聽過的.
就是這位毛主席.
毛澤東他晚年的時候.
他年長晚年的時候.
他接受了很多外國記者採訪的時候.
他非常喜歡說一句話.
我不知道你們有沒有聽過.
我不會講那個湖南口音.
他撇著他的湖南腔.
他說我快要去見馬克思了.
他說我不久之後就要去馬克思了.
我突然覺得.
這是一個幽默呢.
一個唯物主義的幽默呢.
馬克思已經成了一具屍骨了.
他也成為一具屍骨了.
一具屍骨怎麼去見另外一個屍骨呢.
你怎麼知道這具屍骨.
是革命同志還是反革命的屍骨呢.
各位.
一個無神論的人.
他們在某些時候.
他們沒有嚴守黨紀黨綱.
毛澤東沒有嚴守黨紀黨綱.
在他的世界觀有一個盡頭.
在他生命的.
面對死亡.
面對生命的重大的這種衝擊.
關卡的時候.
他的世界觀是如此的狹窄.
弟兄姊妹今天我們要回到詩篇.
我要藉著詩篇36篇.
特別是藉著一至九節.
我們看到兩種世界觀.
我們看到了一至四節.
講到一個罪人的世界.
一個被無神.
或者是一個偶像.

$^{281}$所籠罩的一個世界.
沒有福音的一個世界.
那五至九節呢.
我們看到了一個讚歌.
大衛把他的眼睛.
從這個無神的世界.
沒有神的世界觀.
所籠罩的世界轉向這一位.
被耶和華獨一上帝.
所掌管所護理的世界.
當然我們今天只讀到第九節.
如果你讀第十一節十二節.
他就獻上了他的禱告.
首先我們看到.
當大衛在開始寫.
做這首詩篇的時候.
他自稱他自己是什麼.
我忘記給大家看了.
他自稱自己是普人.
對不對.
他自稱他自己是耶和華的普人.
弟兄姊妹.
這個稱謂在詩篇當中.
十八篇七十八篇九十六篇.
反反復復的出現.
請問大衛是誰.
在以色列.
在當時世代大衛的身份是什麼.
是君王對不對.
一個做王的人.
很容易飄飄然.
很容易忘了他是誰.
很容易找不著北.
因為他上一個王就是這樣子.
大衛上一個王是誰.
掃羅對吧.
掃羅一旦做王.
掃羅做王前很謙卑的.
但是掃羅一旦做王就飄了.
但是大衛.

$^{321}$之所以是和珅心意的王.
他反而更加的清醒.
他寫這首詩篇.
他關注這個世界.
他看到這個世界的罪惡.
和一切難以解開的問題的時候.
弟兄姊妹.
他是以耶和華的普人.
他是耶和華的奴僕的身份.
而不是以君王的身份在這裡看這個世界.
他是以大衛的哀嘆.
大衛的讚美.
是以一個敬拜者的身份.
而不是以一個哲學家的身份.
在這裡尋求答案.
弟兄姊妹.
我們也要為我們的國家來禱告.
為我們的領袖來禱告.
無論他們是誰.
巴不得他們不是一個.
以政治政客的身份.
以一個.
學者的身份.
而是一個敬拜者.
一個耶和華普人的身份.
如此一來.
我們才有方向.
弟兄姊妹我們看到今天的經文.
在開始的時候.
他就說.
惡人的罪過在他心裡說.
我眼中不怕神.
你如果讀西班牙原文.
西班牙原文的第一個字.
其實是一個動詞.
就是這裡說.
惡人的罪過在他心裡說.
的那個說.
其實是整個原文的第一個字.
說.

$^{361}$你說這有什麼特別嗎.
在整個舊約當中.
是非常非常強烈的一個字眼.
這個字眼用在什麼呢.
基本上是用在耶華神的啟示.
用在耶華親口說的.
或者是透過先知說的.
耶華如此說.
這個字.
不是普通的你說我說.
我們落家常.
而是.
而是一道神諭.
而是一道聖旨.
你知道.
而是一個超自然的啟示.
用這個字.
大衛用這個字在表達.
用一個非常強烈的反諷.
說.
我們的罪惡好像是一個先知.
好像是一個假先知.
在那裡發聖旨.
在那裡發神諭奉天承運.
罪人照約.
約什麼沒有神.
不怕神.
這是這個罪惡的世界.
反反復復.
在強調的一種敘事.
不是嗎.
各位.
這是一個假傳聖旨的世代.
這是一個篡奪王位的世代.
有一個聲音.
哪怕今天在座很多人我知道你們都是基督徒.
在你做重大的決定的時候.
你有去求問神嗎.
甚至在你軟弱犯罪.
沒有人看見的時候.

$^{401}$你有在神面前有一種深深的敬畏和恐懼嗎.
沒有.
有一個聲音.
有一道神諭.
在我們的心中告訴我們說不怕不怕.
別怕別怕.
他不管的.
他不會審判的.
他不會刑罰的.
他看不見.
沒事.
這些聖經不是說沒有神.
這些聖經是說什麼.
不怕神.
這些聖經還不是在講無神論.
讓我告訴大家.
無神論是在這200年左右才出現的.
在人類歷史當中所出現的一個怪胎.
在古今中外.
過去從來沒有過無神論這個詞.
請問亞當夏娃.
人類的始祖.
在創世記第三章.
他們犯罪被逐出伊甸園的時候.
怎麼樣.
他們是無神論了嗎.
亞當夏娃是無神論嗎.
我告訴你.
亞當和夏娃他們的心裡.
壓根就沒有無神論.
這麼邪惡的一個概念.
其實墮落之後的人類.
仍然是有宗教的人類.
對吧.
從某種意義上.
墮落不是除去了宗教.
而是使我們的生命.
和生活的一部分.
不再被上帝統治了.
不再被上帝管理了.

$^{441}$各位.
不是完全不被上帝統治.
而是不完全被上帝統治.
聽得懂嗎.
就是剛才趙天安說的.
那個統攝性.
已經在我們的生命當中.
大打折扣了.
不是.
不完全.
不是完全不被上帝統治.
而是不完全被上帝統治.
丁慈美.
這是不是你我的生活.
有的時候也要停一下身.
有些時候呢.
有些時候最好別來管我.
讓我自己也瀟灑一把.
各位.
即便.
即便一個無神論者.
你知道嗎.
無神論不只是口頭上否認神.
其實它是用在.
一種生活方式.
一種不怕神的生活方式.
去生活.
有一句話.
這一位.
切特斯頓.
他說的.
這是他跟.
大衛休默在辯論的時候.
大衛休默.
這是二十世紀非常厲害的一個.
無神論的哲學家.
他是被稱為無神論的鬥犬.
他們辯論的時候.
最後他發出了一個結論.
這個切特.

$^{481}$不是切斯特頓.
他說什麼呢.
他說當人不再相信上帝的時候.
你不要以為他什麼都不信了.
他不是從此以後.
他什麼都不信了.
而是從此以後他什麼都會去信.
你打開了潘多拉的盒子.
人什麼都會去信.
我借用他這句話.
我說.
當人不再敬畏上帝.
你知道敬畏是fear god.
對不對.
當人不再敬畏上帝.
懼怕上帝的時候.
人變得什麼都怕.
人變得前怕狼後怕虎.
人變得怕這個世界上.
一切比他自己更強大的力量.
我記得有一個人.
他說.
在日本的時候他跟我說.
我發現只要是帶老字的.
好像我一輩子都在怕.
小的時候怕老師.
長大了怕老闆.
結婚了呢.
然後他甚至怕老鼠.
不光怕老虎.
也怕老鼠.
什麼都怕.
為什麼.
因為他不怕上帝.
一個不怕上帝的人什麼都開始怕.
我記得在某種場合.
跟大家分享過.
大家知道嗎.
深藏之間的這位叫做霍金.
史迪文霍金.

$^{521}$他也說過一句話.
你知道他也是一個無神論者.
他說.
there is no heaven or afterlife.
沒有來生沒有天堂.
這是一個fairytale.
for people.
afraid of dark.
darkness.
他說.
宗教信仰.
你們這些信耶穌的你們怕黑你們才信耶穌.
英國這一位叫約翰.
德穆斯他怎麼回答呢.
他回應說.
atheism 無神論.
你們這些無神論者.
無神論是你們的一個.
童話故事.
是給你們這些怕光的人.
各位我想.
這位霍金已經過世了.
記得嗎前一段時間.
鬧得瘋瘋火火的那個愛波斯坦案.
披露出來的文件當中.
我很難想像.
這一位坐輪椅的人.
竟然在那個蘿莉島上.
所以他說我不怕黑.
但是他怕光.
他不怕犯罪.
但是他怕神.
刁子妹.
在我們生命當中.
到底什麼讓你懼怕.
你怕這個世界嗎.
你怕股市崩盤嗎.
你怕戰爭嗎.
你怕失業嗎.
你怕這個世界比怕耶和華還怕嗎.

$^{561}$你心中的恐懼.
似乎是怕到一個地步.
比怕上帝還怕.
這就是我們的軟弱.
這就是我們的世俗化的開始.
這就是我們基督教信仰.
世界觀的崩塌.
和那個統攝性的缺失.
講到這裡.
我常常想到一個人.
我想到一個.
敬畏上帝的一個典範.
這是在.
他應該是156幾年生人.
他是蘇格蘭長老會之父.
被稱為.
叫做諾克斯.
這個人他敬畏到上帝.
敬畏到一個地步.
他是不怕任何人.
你知道在156年以後.
這一位叫諾克斯的牧師.
這是一位蘇格蘭長老會的牧師.
他與當時的一位統治者.
被稱為血腥瑪麗.
如果你讀過基督教教會歷史的話.
你知道這個人.
這位女王有四次會晤.
這位血腥瑪麗.
她為什麼被稱為血腥瑪麗.
她殘害新教徒.
她是天主教徒.
她是殘害一切改教的新教徒.
最後一次會晤的時候.
瑪麗女王傲慢地責備.
這個曾經做過奴隸.
現今做牧師的諾克斯.
她說在這個國家你以為你是誰.
你敢如此對你的女王說話.
我不知道你會怎麼回應.

$^{601}$這位諾克斯這樣回答.
他說殿下.
我是一個在蘇格蘭土生土長的百姓.
儘管我不是伯爵也不是領主.
但是上帝使我成為國中的一個公民.
無論在您的眼中我多麼的卑賤.
是的女士.
權力不是來源於你.
而是來源於上帝.
因為你對真信仰比較無知.
哎呦喂.
你敢這麼對川普說嗎.
我覺得比較切合.
他說因為你對真信仰比較無知.
所以我在信仰問題上.
沒有服從君王的義務.
難道古代的希伯來人.
應當皈依法老的宗教嗎.
難道但依理應該信仰彌布賈尼薩和大流士的宗教嗎.
哇哦.
各位.
這是一位敬畏上帝的人.
因為他敬畏上帝.
這個世界上最有權力的人.
手握他生死大權的人.
他可以在這個人面前如此的勇敢.
一個得勝的基督徒.
就是除了神以外.
再也不怕的基督徒.
一無所懼的.
因為我們怕那位上帝.
為什麼.
因為他是聖潔的.
他是公義的.
我們害怕違背他的律法.
我們不遵循他的道路.
我們害怕這件事情.
勝過害怕其他所有的事情.
甚至因為害怕這件事情.
其他所有的事都不怕了.

$^{641}$這就是我們的生命.
罪人.
就是說詩篇當中.
大衛所敘述的世人.
這些邪惡的人.
惡人為什麼不怕上帝呢.
因為這裡說.
在第二節說.
他以為他的罪孽終不顯露.
我的罪孽被我包裝起來了.
被我粉飾起來了.
被我藏起來了.
終不顯露.
你們有沒有看過這部電影.
上世紀末的.
我忘了九幾年了.
蠻老的一部電影.
很有名.
楚門世界.
看過嗎.
你們都在讀聖經.
所以不看電影.
這部楚門世界.
講到一個人.
他從出生.
其實他的整個生活.
他活了大半輩子.
他發現他的生活.
他房間每一條街道.
甚至天空和太陽都是假的.
全裝滿了攝像頭.
當他發現的時候.
怎麼樣.
崩潰了.
原來我沒有任何隱私可言.
原來我的生命.
就是這樣赤裸敞開的.
他想方設法要捍衛他的隱私權.
他想方設法要逃離這個世界.
各位啊.

$^{681}$你我在上帝的眼中.
就如同楚門一樣.
你一個心思和意念.
你以為你不顯山不露水.
你以為說我的罪惡永不顯露.
其實呢.
其實你早被他看穿了.
你所有的誤會的淫穢的念頭.
慾望動機.
讓我說.
都暴露在日光之下了.
各位.
罪人會.
不惜一切的代價.
來捍衛這樣的一種隱私權.
所以.
我們會說審不審判.
甚至我們會說沒有審.
弟兄姊妹.
這是這個世界.
是我們生命當中最大的謊言.
彌天大謊.
沒有審.
這個謊言是關於.
你對你整個的生命.
整個的世界.
最基本的一個認知.
換句話說就是構成你世界觀的.
最基本的認知.
如果這個最基本的認知.
是虛謊的.
那麼你的一生無論說什麼.
無論做什麼都籠罩在這個謊言之下.
因為這個謊言.
統治了你的一生.
我們很多都是從中國來的.
特別我們.
最近大陸來的蠻多的.
請問中國大陸是一種唯物主義的.
無神論的教育對不對.

$^{721}$你猜中國大陸的信仰狀況如何.
這是我引用的.
這是北京的一個調查公司.
叫做林鼎調查公司.
他調查了一個報告.
中國人的精神狀況全國抽樣調查.
經過大數據的.
統計調查整理推算.
他得出這樣一個結論念給大家.
16歲以上的中國人.
85\%都持有某種.
超自然信仰或從事.
某種宗教活動.
雖然有58\%的人.
聲稱自己是無神論者.
但其中44\%的人.
在這12個月中.
有過某種形式的宗教實踐.
有什麼樣呢.
1.2億人自稱相信神靈.
鬼佛.
有1.4億人.
相信財神.
這個信得多.
有1.45億.
相信風水.
3.62億.
過去在12個月.
算過命.
看過相.
然後有.
多達7.54億人.
在過去一年.
清明節的時候.
去上墳去燒香.
去祭祖.
其中2.06億人.
承認祖先神靈的存在.
而且有1.23億人.
在家裡供奉著.

$^{761}$祖先的牌位.
中國人.
我們不是無神論嗎.
我們不是唯物主義者嗎.
我們不是都要去馬克思嗎.
去見馬克思嗎.
各位.
沒有的.
不是的.
我常常提到這個名字.
楊鳳岡教授.
如果大家聽過我帶領導高輝和講道的話.
他是普渡大學宗教系的一個教授.
也是白宮的特別的宗教顧問.
他是一個長老.
他寫了一篇文章叫做.
當代的中國宗教復興.
與宗教短缺.
什麼意思啊.
他說中國人的精神信仰的整體狀況.
他用了八個字.
宗教不彰.
巫術盛行.
他說.
中國人的民間信仰.
或巫術在中國市場.
如此的迅猛發展.
關鍵是人們對超自然.
對超自然.
和神秘力量的需求.
在市場化的過程當中被喚醒了.
如果沒有辦法.
在制度化的宗教當中.
他所謂的制度化的宗教.
就是佛教 基督教 伊斯蘭教當中.
得著滿足.
那麼人們呢就去灰世.
就去黑世裡去.
找滿足.
你說好了.

$^{801}$我不是那85\%.
我既不嫉俗也不犯命.
我崇尚科學.
我高舉理性.
中世紀有一位神學家.
他說理性就是無神論的至高神.
感覺.
就是這位神明的祭司.
理性和感覺構成了.
所有無神論者的幻想.
理性是被造的.
理性也是墮落的.
弟兄姊妹 聖經明明地告訴我們.
無論你是信仰什麼宗教.
或者你什麼宗教不信.
聖經明明地告訴我們一件事情.
就是今天沒有一個人的生命.
中間是沒有敬拜的.
沒有一個人是真正意義上的無神論者.
在上帝以外.
如果你把單單屬於上帝的屬性.
單單只有上帝才能夠帶來的意義.
歸給了其他的受造物.
無論是人 無論是思想.
無論是主義.
無論是什麼物件.
弟兄姊妹 那個就是你的上帝.
從這個角度上來說.
聖經控訴所有的人.
如果你不信那位獨一的真實人.
你就是一個偶像崇拜者.
你知道現代人最大的偶像是什麼嗎.
當我們講到世俗化的時候.
在一個世俗化的社會.
在一個宗教不彰.
巫術橫行的社會當中.
你知道最大的偶像是什麼嗎.
John Piper.
約翰·派伯牧師他說.
最大的偶像就是人的自我.

$^{841}$每一個中國人.
你知道嗎.
我們裡邊都有個小秦始皇.
當我們推翻了紫禁城裡邊的帝王.
我們就在自己的心中登基.
我們的身體就是我們的紫禁城.
我的身體就是我的王土.
我要在這個身體裡縱育.
我要在這個身體裡墮胎.
我要在這個身體裡做我一切的事情.
只有我能做主.
我的人生我要做主.
約翰·派伯講到這裡的時候.
他突然說.
你怎麼來證明.
這個是我們世俗世代最大的偶像呢.
自我是最大的偶像呢.
他說一件事情.
你就看吧.
整個西方世界婚姻制度的一種崩潰.
婚姻的一種崩塌.
就是最好的解釋.
你知道為什麼嗎.
很簡單.
婚姻是什麼.
婚姻是二人合為一體.
婚姻是兩個人的聯合.
結果突然發現有兩個自主權.
有兩個巨大的自我.
當這兩個人走到一起的時候.
就要崩塌.
就要吵架.
女王姊妹.
無論我們崇尚自我.
無論我們追求這個世界上的人事物.
我們都是成了世俗世界的一份子.
女王姊妹.
普蘭丁格他說.
這個世界上根本就沒有無神論.
只有反神論.

$^{881}$同意嗎.
這個世界上根本壓根就沒有無神論.
只有反神論.
反神論的世界觀是一個得了妄想症的世界觀.
整個世界都陷入在這種的反神論當中.
世人一個不認識神的惡人.
這裡說他不顧一切的.
他竭盡全力的去反對神的勇能和神性.
反對神的檢查和審判.
甚至反對神的慈愛.
我再用毛主席語錄的一句話.
叫做什麼呢.
與天鬥.
大陸來的就知道.
與天鬥 其樂無窮.
這個世界就是這樣的一個世界.
女王姊妹.
當他與天鬥其樂無窮的時候.
你會發現.
他的罪孽.
他的詭詐.
他圖謀罪孽.
他們不行善道.
他們不憎惡事.
這些道德上倫理上就產生了一個崩塌.
我覺得.
我們中國人.
我們說我們是禮儀之邦.
我們是泱泱大國.
我們是文明古國.
對吧.
古代的中國人雖然不認識耶和華上帝.
但是人家敬天.
人家說人在做.
天在看.
對吧.
但是今天呢.
今天.
人越發越過了對神對天的敬畏.
我們開始扯彌天大謊.

$^{921}$我們開始膽大包天.
我們開始想要做齊天大聖.
我們失去了對上帝的敬畏.
我們失去了一切美善的道德基礎.
有人說.
你知道中國為什麼在這些年間.
有這麼多道德.
就一次一次地突破道德的底線.
你知道在中國傳統社會是一個農業社會.
是一個儒家思想.
或者儒教所管理的這種傳統的農業社會.
但是呢.
今天中國已經快速轉型成了一個現代社會.
兩個現象.
第一工業化.
第二都市化.
你知道西方美國從農業社會轉型成工業社會.
花了多少年.
整個西方花了兩百年.
你知道中國花了多少年.
有人說二十年.
有人說四十年.
總而言之.
我們用二十年完成了人家兩百年的壯舉.
你知道.
你知道這樣的一個改變.
徹底顛覆了中國人的這種道德約束力和人際關係.
你知道在一個農業社會當中.
你世世代代.
你三代同堂.
你世居故里.
你一輩子.
老張家就一輩子在這個屯.
在這個屯.
對不對.
而村裡的張老爺李大媽都是世交.
如果你犯了錯會怎麼辦.
如果你做了不道德的事怎麼辦.
門口村口就有幾個大媽坐在板凳上.
老張家他那個三兒子.

$^{961}$你知道嗎.
這形成了一種道德約束力.
千夫所指.
無疾而死.
但是今天中國快速的城市化了.
請問在一座摩天大樓當中.
你跟你的靈社有多深的這種的連接.
請問你跟你的上下樓之間有多深的交往.
沒有.
每一個人成了一個獨立的個體.
我再見毛澤東.
今天高舉他了.
的一句話叫做什麼呢.
中國就是光合常打散.
無法無天.
為什麼無法.
就是沒有天.
沒有上帝.
所以沒有法治.
沒有道德.
這段經文再次讓我們看到.
一個不懼怕上帝.
一個把上帝想要排擠在外的.
他就是無法無天.
不僅無法無天.
我今天會多停留在前四節.
第二節這裡告訴我們什麼呢.
他說他自誇自媚.
其實是一個希伯來文.
其實這是一個希伯來文.
這個希伯來文.
這個希伯來文在原文當中.
其實就是一種自吹自擂自誇.
愛自己.
體貼自己.
奉承自己.
自戀.
現代人的語言.
罪人會自戀的.
罪人成功地說服了自己.

$^{1001}$賄賂了自己的良心.
他不斷地告訴自己說.
你還不錯.
你是個好人.
你比上不足比下有餘.
我曾經看一個人一篇文章.
他說.
God does not grade on curve.
你知道現在大學生考試.
他是有一個curve.
如果你只得了一分.
但是全班都是零分.
你還是A.
對吧.
他說God不會curve你的.
God的標準不會根據.
你比這間教會300人.
你勝過了299人.
你就是一個藝人了.
不會的.
弟兄姊妹.
我們用這樣的方式自吹自擂.
我們用這樣的方式自誇自媚.
弟兄姊妹.
我們以為我很光鮮.
我很體面.
我們把罪藏得很好.
很美.
我們不斷地往自己身上化妝塗粉.
以至於我以為我的罪惡終不顯露.
話說回來.
我有的是刷短視頻.
我刷到一個化妝博主.
我的天哪.
化妝前化妝後幾乎判若兩人.
對吧.
你們笑.
你們也刷過那種視頻嗎.
各位.
當人不敬為神.

$^{1041}$就開始自己向自己獻媚.
給自己化妝.
給自己修圖.
你說.
不.
我從來不化妝.
也不修圖.
各位.
我說的不只是修你的照片.
我們都在修你的自我形象.
我們都在修我們自己的良心.
我們都在修我們人.
人.
罪人的那個本性.
我們不願意面對神光照下那個自己.
我們不願意在神的光中怎麼樣.
得以見光.
為什麼見光死.
我不想知道真相.
真相很傷人.
我不想面對我自己的本性.
因為我的本性很不堪.
讀到過一篇文章.
總結了三個問題.
自我檢查一下.
你自戀嗎.
基督徒.
第一.
你是否很容易承認自己的錯誤.
還是說.
你是一個非常.
英文怎麼說.
defensive.
只要被說的時候.
你馬上就開始變得defensive.
自我辯護.
第二個.
你是不是很容易發現別人的缺點.
我千萬不要把這個當成聖靈的恩賜.
這基本上是法利賽人的特長.

$^{1081}$不是聖靈的恩賜.
馬上就看到人眼中的什麼.
刺.
小小的.
火眼金睛.
我一眼就看出來.
第三個.
你的事奉.
或者你在教會教的事奉.
你在家裡幹活.
你在上班的時候工作.
你是不是非常的在乎別人的肯定.
你知道嗎.
心理學上有一個詞叫做假想觀眾.
這個心理從青少年開始出現.
從teenager開始.
我不知道你們記得嗎.
當你們teenager的時候.
你走進教室.
你覺得所有的班上的人都在看我.
看我每一根髮型有沒有到位.
看我的衣著.
這種觀眾.
這種心態其實蔓延在我們裡頭.
這三個問題你都是是的話.
各位啊.
小心.
你可能是很自戀的.
時間的原因不講它了.
這是一個希臘神話當中自戀的美男子.
愛上了自己在湖中的導引.
給大家念一段雞湯.
這段雞湯就是現代人不斷美化自己的.
對不起.
我甚至在教會當中也聽過類似的雞湯.
給他念幾句話.
一切從自己出發.
努力做最美的自己.
不要等別人來愛你.
你要學會先愛自己.

$^{1121}$不要在意別人的看法.
做最真實的自己.
為自己而活.
別人怎麼想想你不重要.
重要的是你自己怎麼看自己.
幸福取決於你的心態.
與別人無關.
你每天要對著鏡子喊.
我是最優秀的.
我是最棒的.
不要為你的人生設限.
你的未來在你自己手中.
你不覺得這是這個世界最主要的一種論調和敘事嗎.
今天的年輕人不就是這樣子嗎.
自我饞媚.
自誇自媚.
自我膨脹.
讓我這樣說.
這不光是今天年輕人的錯.
這是在整個人類歷史當中.
特別是近代歷史當中.
一個強烈的現象.
我把它稱作一個分水嶺.
這個分水嶺.
第一是中世紀的宗教改革.
第二宗教改革之後的叫做啟蒙運動.
大家知道宗教改革吧.
但是啟蒙運動是在那之後.
17到18世紀的時候.
一種思想文化的變遷.
啟蒙運動高舉.
以人的理性為皈依.
高舉人性的美.
人性的善.
他們樂觀地認為.
在人類美善的加持下.
我們的未來更美好.
更和諧.
更進步.
但是宗教改革是什麼.

$^{1161}$對於宗教改革的這些改教家.
對於馬丁路德.
對於約爾加爾文.
對於慈韻禮而言呢.
你知道嗎.
就好像大衛.
他觀察到這個世界所有的問題.
這個世界上所有的罪惡的原因是什麼.
每一個人都犯罪的違背了上帝.
因此宗教改革的核心就是離棄自我歸向上帝.
那他強調了每一個人.
所以每一個人歸向上帝之後呢.
叫做人人皆祭司.
但是你知道啟蒙運動就把這樣的一個人人皆祭司的思想.
在一個趣味的世界觀裡邊推向極致.
叫做什麼呢.
人人皆上帝.
對於啟蒙運動來說.
是人人皆上帝的一個世界觀.
我請大家記住啟蒙運動這個妖孽.
我真的啟蒙運動是近代所有思想當中.
你可以說是最大的妖孽.
因為其他所有接下來違背基督教文化.
基督教理念的都是以啟蒙運動為母體的.
他們的一個最強烈的slogan說什麼呢.
天地之間個體為尊.
他把人拔高到世界宇宙的中心.
把人自誇自媚發揮到了一種極致.
都覺得自己是somebody.
各位對於這群人而言.
你知道嗎.
我就說接下來我們所處的歐美世界了.
他們的盼望就是打破啟蒙的意思.
就是打破舊世界建立新世界.
聽上去很高大上對吧.
但是打破舊世界對於歐洲人而言.
你知道第一件事情就是要把上帝從他們的世界當中趕出去.
因為整個歐洲的世界其實就是建立在有神論.
甚至就是建立在一神論.
獨一神論的基礎上所建立起來的.

$^{1201}$他們說以理性來代替信仰.
或者說理性就是他們的信仰.
以科學代替宗教.
或者說科學就是他們的宗教.
以自我代替上帝.
以自由代替權威.
以進步代替天國.
不需要等待天國了.
我們在這裡進步建立一個美好的國家.
多好啊.
中子妹.
這道分水嶺.
這形成了人類歷史上的兩個敘事.
你是活在宗教改革的敘事當中.
承認我是個徹頭徹尾的罪人.
我需要救贖呢.
你還是活在啟蒙運動的敘事當中.
剛才那些歐耶的人.
天天照鏡子說魔鏡魔鏡.
告訴我.
這個教會當中.
最進錢的人是誰.
是我嗎.
這個教會當中.
最愛主的人是誰.
中子妹.
陷我於不義.
各位啊.
所以整個人類歷史就分成了兩種的敘事.
有一本書叫做《耶穌在北京》.
是一位宣教士叫做大衛·艾克曼.
他寫的.
他在這本書當中講過一個故事.
很有意思.
有一天馬克思.
就是卡爾·馬克思.
他的女兒問恩格斯.
他說恩格斯啊.
世界上你最恨的人是誰.
你知道恩格斯怎麼說嗎.

$^{1241}$他說有一個叫斯布真的英國牧師.
就因為他.
我們在英國的革命沒成功.
同樣.
大家知道《雨果》嗎.
就是寫《悲慘世界》的那個.
他還寫了一本書叫做《九三年》.
他也有一個很類似的描述.
他這樣說.
他說.
每一個村莊都有一位點火的啟蒙教師.
但是呢.
每一個村莊都有一位滅火的牧師.
啟蒙運動就是想把火點起來.
今天美國的白左們就是要點起火來.
但是呢.
把人類的邪惡的火澆滅下去.
把聖靈聖潔的火點起來.
這是神的兒女肩負的使命.
丁子維.
在這場全人類的狂歡當中.
自戀的狂歡當中.
只有一副藥方.
在無神論和絕對的一神論當中.
你找不到任何的立足點.
換句話說.
除非你活在一個絕對的一神論當中.
基督徒就是絕對的一神論.
什麼意思啊.
我再次用那句話.
就是統攝性的信仰.
你是徹底的絕對的一神論嗎.
耶和華神在你生命當中.
一切的領域當中都是絕對的嗎.
徹底的嗎.
還是有的時候你會悄悄地自告自答.
有位牧師他說的很有意思.
他很幽默.
他說我每天早晨起來的時候都很自負.
晚上回家的時候都很自卑.

$^{1281}$今天他又失敗了.
各位.
是不是我們當中還會這樣呢.
求主用他那個絕對的一神論來改變我們.
不僅如此.
弟兄姊妹.
這裡說我們乃是與良善怎麼樣.
智慧良善.
第三節說已經斷絕了.
各位啊.
斷絕了.
連藕斷絲連都不是哦.
你的全部的生命跟美善跟聖潔跟公益斷得一乾二淨.
好傷人哦.
為什麼呢.
他接著說他在床上圖謀罪孽.
各位在床上是幹嘛的.
睡覺休息的.
對吧.
但是他還睡不著.
睡不著.
他就在那裡想.
他就在那裡打算盤.
打什麼算盤.
罪孽的算盤.
弟兄姊妹.
然後呢.
定義不行不善的道.
他走路的時候也是在那裡走不善的道.
弟兄姊妹.
他沒有想到在在這個這個申命記的六章那裡講到說這是教育的大憲章.
無論你坐在家裡行在路上.
躺下起來都要談論.
談論什麼.
談論神的道.
談論耶和華的話.
各位罪人是躺下起來.
怎麼樣.
罪孽罪孽.
基督徒應當是什麼.

$^{1321}$躺下起來.
神的話來指引我們.
躺下起來就是指向全部的生活.
你是淪陷在罪惡當中.
你還是活在神的道當中.
我覺得基督徒有一個.
我們中間也有剛受洗的.
有一個成長的路徑.
就是你剛信耶穌的時候.
當然你信耶穌就表達你已經認罪.
承認自己是個罪人.
你知道你是怎麼認罪的嗎.
你做那個認罪表也好.
或者認罪清單也好.
或者你思想也好.
我犯了這個罪那個罪.
所以呢.
所以我是個罪人.
但是呢.
希望這首詩篇讓你看到一件事情.
不光是因為你犯了罪.
所以你是個罪人.
而是說.
你已經與良善斷絕了.
本質上你已經壞透了.
你是個罪人.
所以你會犯這個罪那個罪.
因為你信了耶穌.
因為你受了洗.
甚至因為你參加了門徒培訓.
你還會犯這個罪那個罪.
這是我們對罪的一個認知.
弟兄姊妹.
你對罪的認知越深切.
你在接下來.
第五節到第九節.
那個世界對你來說就越吸引你.
越發的讓你覺得.
Amazing.
Amazing Grace.

$^{1361}$各位.
為什麼.
因為第一節到第四節和第五節到第九節之間.
那個落差好大喔.
各位.
如果在你的信仰當中那個落差很小.
你覺得這個世界還蠻美好的.
上帝的世界也沒有那麼好的時候.
你的信仰不會Amazing.
不會有Amazing Grace.
但是.
這是為什麼我花了這麼多的時間.
弟兄姊妹.
我們停留在前四節.
但是好.
我們用最後的時間進到後五節.
進到一個天賦的世界.
一個美善的世界.
一個上帝掌權的世界.
寫得好好.
寫得好美.
我們門訓要背.
對不對.
這節聖經.
他說什麼.
耶華.
你的慈愛上及諸天.
你的信實達到窮倉.
他的公義好像高山.
他的判斷如同深淵.
耶華.
人民生畜.
你都救護.
弟兄姊妹.
你知道這是對應著上邊.
上邊說什麼.
我的罪惡不顯露的.
他看不見的.
這裡說什麼呢.
窮倉.

$^{1401}$諸天.
高山.
深淵.
神都在那裡.
神的慈愛和信實.
神的判斷都在那裡.
那個判斷也可以翻成審判.
都在那裡.
你藏.
你往哪裡藏.
弟兄姊妹.
無論是宏觀的宇宙.
微觀的分子.
哪裡沒有耶華的慈愛.
哪裡沒有耶華的信實公義判斷呢.
他創造一切.
他關注一切.
他的慈愛公義充滿整個宇宙.
遍滿整個受造界.
兄姊妹.
這就是為什麼我們不能拜偶像.
十誡當中的第二誡教導我們.
不可以拜偶像.
不可以為自己雕刻偶像.
做什麼形象.
服什麼上天下地.
地底下和水中的百物.
不可跪拜那些不可事奉.
為什麼.
因為版權所有.
他代表的是耶華的慈愛.
你拿他來幹什麼.
肖像權都已經歸屬於耶華了.
這四個字.
諸天.
瓊蒼.
高山深淵.
在我們和合本的翻譯當中.
比較注重自然界地理空間層面的一種描繪.
對不對.

$^{1441}$無所不在.
但是我很喜歡另外一個譯本.
是司高本.
他把諸天翻譯成青天.
青天大老爺那個青天.
他把那個高山翻譯成聖山.
你知道這種翻譯更傾向一個道德層面.
一個價值觀層面的一種敘事.
對不對.
神的慈愛在青天.
青天就不是那個sky.
青天是講到一個道德層面的聖潔公益.
所有的中國人在沒有遇見上帝之前.
我們都喜歡包青天不是嗎.
但是只有上帝.
他的慈愛上達青天.
弟兄姊妹.
我們看到他的慈愛.
其實如果你讀到第十節.
反反復復.
其實講到葉劃的慈愛.
葉劃的慈愛在哪裡呢.
今天的經文.
第一在諸天之上.
第二出現在哪裡.
出現在翅膀的陰下.
第七節.
你的慈愛何其寶貴.
神投靠在他的翅膀的陰下.
各位其實是一個聚焦.
你知道嗎.
首先是個廣角.
然後呢.
祝天.
瓊山.
瓊倉.
然後聚焦在哪裡.
人.
聚焦在你我.
聚焦在他的翅膀的陰下.

$^{1481}$弟兄姊妹.
基督徒才有隱形的翅膀.
你們都聽不懂我的梗.
所以你們也不聽張韶涵的歌.
各位.
這個翅膀的陰下.
在聖經當中有兩個比喻.
第一個就是.
今天我聽到儒學講路德紀.
路德紀裡邊他把波拉斯和路德.
用了一個翅膀的陰下來比喻.
在波拉斯那種救贖的愛中.
這位落難的女子得到了呵護和救贖.
另外一個呢.
就是在馬太福音當中.
主耶穌說的.
他說.
以指著耶路撒冷.
他說我要如同母雞.
想把你們聚集在翅膀的陰下.
但是他們不願意.
弟兄姊妹.
無論是哪一個比喻.
都是在講到那個細膩的救贖的遮蓋的愛.
我們都玩過老鷹捉小雞吧.
玩過嗎.
老鷹捉小雞的時候.
有一隻母雞張開他的手臂.
張開他的翅膀.
小雞就在後邊.
就是安全的.
哪一個最危險.
越往後越危險.
對不對.
弟兄姊妹.
我們的信仰也是這樣子.
給大家講一個總統.
約翰遜總統.
他雖然是基督徒.
但是他的信仰又不太堅定.

$^{1521}$他最怕什麼.
他恐高.
恐高最怕幹什麼.
坐飛機.
但是作為國家元首.
他不得不常常坐飛機.
他坐飛機的時候.
他又害怕.
有一天天氣不好.
他又有行程.
天氣預報說.
這個判斷說.
可以起飛.
空軍一號沒問題的.
絕對可以起飛.
他心裡戰戰兢兢.
你知道人家幹什麼.
你們怎麼辦.
他打電話給格裡漢牧師.
他說格裡漢牧師.
他不是說請你為我禱告.
請你和我一同登基好不好.
我有可能掉下來.
但上帝應該不會讓你掉下來.
他說在格裡漢的翅膀的陰下.
他比較安全.
坦白說比較危險.
因為他是排隊排得比較好.
弟兄姊妹.
每個人.
上帝說他的慈愛何其寶貴.
他要遮蓋你.
他用他的慈愛來遮蓋你.
弟兄姊妹.
求主幫助我們.
最後一點.
我講到這裡.
這裡除了慈愛這個詞反復出現.
還有一個反復出現的字.
就是公義.

$^{1561}$弟兄姊妹.
慈愛和公義.
在今天這個世界當中.
這個墮落的世界當中.
這兩者是斷裂的.
是撕裂的.
是彼此衝突的.
是沒有辦法同時成立的.
若要有慈愛.
那就要犧牲一點點公義.
若要有公義.
那你就不能太慈愛了.
不是嗎.
有半黑臉的.
有半白臉的.
但是耶和華到底是黑臉還是白臉.
在詩篇第八十.
五篇第十節這裡說.
我們一起讀一下好吧.
慈愛和誠實彼此相遇.
公義和平安彼此相親.
相親.
這中文不就結婚嗎.
結婚才相親.
相親才結婚.
對吧.
它的原文是親嘴的意思.
其實很親密.
公義與平安.
慈愛與誠實.
其實在現代的時代當中.
是一個狹路相逢.
格外眼紅的敘事.
你知道嗎.
怎麼可能又慈愛.
又公義又誠實呢.
弟兄姊妹.
今天是八一病.
所以我必須結束到這個點上.
當我們來到耶穌基督和他的十字架前.

$^{1601}$我們發現公義與慈愛完美的契合.
契合在耶穌基督和他的十字架上.
請問你在他的十字架上有沒有看到公義.
上帝是不是妥協了公義來赦免我們.
沒有.
我們看到了上帝的公義.
上帝的審判.
請問你在十字架上有沒有看見慈愛.
同樣的.
我們也看見了慈愛.
在耶穌基督的十字架上.
公義.
慈愛.
匯聚在那裡.
他們彼此相和.
他們彼此親吻.
那裡是我們的救贖.
是我們的遮蓋.
所以.
我用一個字好嗎.
來表達.
如果你真的認識上帝的慈愛.
你就會撒嬌.
你如果真的認識上帝的公義.
你就會敬畏.
一個只有敬畏.
不會撒嬌的信徒.
是很累的.
一個只會撒嬌.
沒有敬畏的信徒.
是不會成長.
不會破碎的.
求主幫助我們.
在中國的電視裡邊有一個字.
在古裝劇裡邊.
我覺得特別好.
用我們基督徒用來.
我們怎麼稱呼這位上帝呢.
父王.
祂是我們的父.

$^{1641}$所以我們可以撒嬌.
祂是我們的王.
所以我們要服服.
我們要順服.
我們要敬畏.
求主幫助我們.
願我們對神的認識.
重新塑造我們的世界觀.
我們的價值觀.
讓我們在祂的光中得以見光.
阿們.
\newpage



\section{撒母耳記上 25:1-44}
\label{sec:032100}
\textbf{作個好男人,作個好女人——大衛與亞比該 (1)}~~~~李思聰牧師
\newline
\newline
連結: \href{http://www.church.com.hk/acms/content.asp?site=cdc\&op=show\&type=product\&code=032100\&layout=sermon}{\texttt{032100}} ~~~~ 語音日期: 2026-05-03
\newline
\newline
\hyperref[sec:032099]{\small{< < < PREV SERMON < < <}}
~
\hyperref[sec:index_chronic]{\small{[返順時目]}}
~
\hyperref[sec:032101]{\small{> > > NEXT SERMON > > >}}
\newline
\newline
撒母耳記上 25:1-44
\newline
\begin{longtable}{cl}
\hline
\hline
章節 & 經文 (和合本修訂版)\\
\hline
25:1 & \begin{tabularx}{0.7\textwidth}{X} 撒母耳死了，以色列眾人聚集，為他哀哭，把他葬在拉瑪他的家裡。大衛動身，下到巴蘭的曠野。 \end{tabularx} \\ \\ \relax
25:2 & \begin{tabularx}{0.7\textwidth}{X} 在瑪雲有一個人，他的產業在迦密。這人是一個大富翁，有三千隻綿羊，一千隻山羊；他正在迦密剪羊毛。 \end{tabularx} \\ \\ \relax
25:3 & \begin{tabularx}{0.7\textwidth}{X} 這人名叫拿八，他的妻子名叫亞比該。拿八的妻子有美好的見識，又有美麗的容貌，但拿八為人剛愎兇惡，是迦勒族的人。 \end{tabularx} \\ \\ \relax
25:4 & \begin{tabularx}{0.7\textwidth}{X} 大衛在曠野聽見拿八正在剪羊毛， \end{tabularx} \\ \\ \relax
25:5 & \begin{tabularx}{0.7\textwidth}{X} 就派十個僕人，對他們說：「你們上迦密到拿八那裡，提我的名向他問安。 \end{tabularx} \\ \\ \relax
25:6 & \begin{tabularx}{0.7\textwidth}{X} 你們要如此說：『願你來年平安，願你家平安，願你一切所有的都平安。 \end{tabularx} \\ \\ \relax
25:7 & \begin{tabularx}{0.7\textwidth}{X} 現在我聽說你有剪羊毛的人，你的牧人和我們在一起，他們在迦密一切的日子，我們沒有欺負過他們，他們也未曾失去甚麼。 \end{tabularx} \\ \\ \relax
25:8 & \begin{tabularx}{0.7\textwidth}{X} 你問你的僕人，他們會告訴你。願我的僕人在你眼前得歡心，因為我們是在好日子來的。請你隨手分點食物給僕人和你兒子大衛。』」 \end{tabularx} \\ \\ \relax
25:9 & \begin{tabularx}{0.7\textwidth}{X} 大衛的僕人到了，就提大衛的名，把這一切話告訴拿八，他們就停頓下來。 \end{tabularx} \\ \\ \relax
25:10 & \begin{tabularx}{0.7\textwidth}{X} 拿八回答大衛的僕人說：「大衛是誰？耶西的兒子是誰？今日悖逆主人奔逃的僕人很多。 \end{tabularx} \\ \\ \relax
25:11 & \begin{tabularx}{0.7\textwidth}{X} 我豈可把飲食，以及我為剪羊毛的人所宰的肉給那些我不知道從哪裡來的人呢？」 \end{tabularx} \\ \\ \relax
25:12 & \begin{tabularx}{0.7\textwidth}{X} 大衛的僕人轉身從原路回去，照這一切的話告訴大衛。 \end{tabularx} \\ \\ \relax
25:13 & \begin{tabularx}{0.7\textwidth}{X} 大衛對他的人說：「你們各人都要佩上刀！」各人就都佩上刀，大衛也佩上刀。跟隨大衛上去的約有四百人，留下二百人看守物件。 \end{tabularx} \\ \\ \relax
25:14 & \begin{tabularx}{0.7\textwidth}{X} 拿八的一個僕人告訴拿八的妻子亞比該說：「看哪，大衛從曠野派使者來向我主人問安，主人卻辱罵他們。 \end{tabularx} \\ \\ \relax
25:15 & \begin{tabularx}{0.7\textwidth}{X} 但是那些人待我們真好；我們在田野與他們一切來往的日子，沒有受他們欺負，也未曾失去甚麼。 \end{tabularx} \\ \\ \relax
25:16 & \begin{tabularx}{0.7\textwidth}{X} 我們在他們那裡牧羊，一切的日子他們晝夜保護我們。 \end{tabularx} \\ \\ \relax
25:17 & \begin{tabularx}{0.7\textwidth}{X} 現在，你當知道，看怎樣做才好。不然，禍患必定臨到我主人和他全家。他性情兇暴，無人敢與他說話。」 \end{tabularx} \\ \\ \relax
25:18 & \begin{tabularx}{0.7\textwidth}{X} 亞比該急忙將二百個餅，兩皮袋酒，五隻宰好的羊，五細亞烘熟的穗子，一百個葡萄乾餅，二百個無花果餅，都馱在驢上， \end{tabularx} \\ \\ \relax
25:19 & \begin{tabularx}{0.7\textwidth}{X} 對僕人說：「你們在我前面走，看哪，我跟著你們去。」她卻沒有告訴丈夫拿八。 \end{tabularx} \\ \\ \relax
25:20 & \begin{tabularx}{0.7\textwidth}{X} 亞比該騎著驢，正下山坡，看哪，大衛和他的人正迎著亞比該下來，她就去迎接他們。 \end{tabularx} \\ \\ \relax
25:21 & \begin{tabularx}{0.7\textwidth}{X} 大衛曾說：「我在曠野為那人看守他一切所有的，以致他未失去任何一樣東西，實在是徒然了！他竟然向我以惡報善。 \end{tabularx} \\ \\ \relax
25:22 & \begin{tabularx}{0.7\textwidth}{X} 凡屬拿八的男丁，我若留一個到明日早晨，願神重重懲罰大衛！」 \end{tabularx} \\ \\ \relax
25:23 & \begin{tabularx}{0.7\textwidth}{X} 亞比該看見大衛，就急忙下驢，在大衛面前臉伏於地叩拜。 \end{tabularx} \\ \\ \relax
25:24 & \begin{tabularx}{0.7\textwidth}{X} 她俯伏在大衛的腳前，說：「我主啊，願這罪歸於我！求你容許使女向你進言，更求你聽使女的話。 \end{tabularx} \\ \\ \relax
25:25 & \begin{tabularx}{0.7\textwidth}{X} 我主不必理會拿八這性情兇暴的人，他就像他的名字一樣；他名叫拿八，為人也真是愚頑。至於我，你的使女並沒有看見我主所派來的僕人。 \end{tabularx} \\ \\ \relax
25:26 & \begin{tabularx}{0.7\textwidth}{X} 現在，我主啊，耶和華既然阻止你親手報仇，避免流人的血，我指著永生的耶和華起誓，又指著你的性命起誓：『現在，願你的仇敵和謀害我主的人都像拿八一樣。』 \end{tabularx} \\ \\ \relax
25:27 & \begin{tabularx}{0.7\textwidth}{X} 現在求我主把婢女送來的禮物給跟隨我主的僕人。 \end{tabularx} \\ \\ \relax
25:28 & \begin{tabularx}{0.7\textwidth}{X} 求你原諒使女的冒犯。耶和華必為我主建立堅固的家，因為我主為耶和華爭戰，並且你一生的日子查不出有甚麼惡來。 \end{tabularx} \\ \\ \relax
25:29 & \begin{tabularx}{0.7\textwidth}{X} 雖有人起來追逼你，要尋索你的性命，我主的性命在耶和華－你的神那裡，如同藏在生命的寶藏中。至於你仇敵的性命，耶和華必甩去，如用機弦甩石一樣。 \end{tabularx} \\ \\ \relax
25:30 & \begin{tabularx}{0.7\textwidth}{X} 耶和華照所應許你的福氣賜給我主，立你作以色列王的時候， \end{tabularx} \\ \\ \relax
25:31 & \begin{tabularx}{0.7\textwidth}{X} 我主就不至於因為親手報仇，流了無辜人的血，而心裡不安，良心有虧了。耶和華賜福給我主的時候，求你記得你的使女。」 \end{tabularx} \\ \\ \relax
25:32 & \begin{tabularx}{0.7\textwidth}{X} 大衛對亞比該說：「耶和華－以色列的神是應當稱頌的，因為他今日派你來迎接我。 \end{tabularx} \\ \\ \relax
25:33 & \begin{tabularx}{0.7\textwidth}{X} 你和你的見識也配得稱讚，因為你今日攔阻我親手報仇、流人的血。 \end{tabularx} \\ \\ \relax
25:34 & \begin{tabularx}{0.7\textwidth}{X} 我指著阻止我加害於你的耶和華－以色列永生的神起誓，若不是你很快地來迎接我，到早晨天亮的時候，凡屬拿八的男丁，必定一個也不留。」 \end{tabularx} \\ \\ \relax
25:35 & \begin{tabularx}{0.7\textwidth}{X} 大衛從亞比該手中收了她帶來的禮物，對她說：「平平安安上你的家去吧！你看，我看了你的情面，聽了你的話。」 \end{tabularx} \\ \\ \relax
25:36 & \begin{tabularx}{0.7\textwidth}{X} 亞比該到拿八那裡，看哪，他在家裡擺設宴席，如同王的宴席。拿八心情舒暢，酩酊大醉。所以亞比該大小事都沒有告訴他，直等到早晨天亮的時候。 \end{tabularx} \\ \\ \relax
25:37 & \begin{tabularx}{0.7\textwidth}{X} 到了早晨，拿八酒醒了，他的妻子把這些事都告訴他，他就發心臟病快死了，僵如石頭。 \end{tabularx} \\ \\ \relax
25:38 & \begin{tabularx}{0.7\textwidth}{X} 過了十天，耶和華擊打拿八，他就死了。 \end{tabularx} \\ \\ \relax
25:39 & \begin{tabularx}{0.7\textwidth}{X} 大衛聽見拿八死了，就說：「耶和華是應當稱頌的，因為我從拿八手中受了羞辱，他為我伸冤，又阻止他的僕人行惡；耶和華使拿八的惡歸到他自己頭上。」於是大衛派人去向亞比該說，要娶她為妻。 \end{tabularx} \\ \\ \relax
25:40 & \begin{tabularx}{0.7\textwidth}{X} 大衛的僕人來到迦密，到亞比該那裡，對她說：「大衛派我們到你這裡，要娶你作他的妻子。」 \end{tabularx} \\ \\ \relax
25:41 & \begin{tabularx}{0.7\textwidth}{X} 亞比該起來叩拜，俯伏在地，說：「看哪，你的使女情願作婢女，為我主的僕人洗腳。」 \end{tabularx} \\ \\ \relax
25:42 & \begin{tabularx}{0.7\textwidth}{X} 亞比該立刻起身，騎上驢，五個女僕跟著她走。她跟從大衛的使者去，就作了大衛的妻子。 \end{tabularx} \\ \\ \relax
25:43 & \begin{tabularx}{0.7\textwidth}{X} 大衛先前娶了耶斯列人亞希暖，她們二人都作了他的妻子。 \end{tabularx} \\ \\ \relax
25:44 & \begin{tabularx}{0.7\textwidth}{X} 掃羅已把他的女兒米甲，就是大衛的妻子，給了迦琳人拉億的兒子帕提為妻。 \end{tabularx} \\ \\
[1ex]
\hline
\hline
\end{longtable}
$^{1}$今天我們繼續看三位機上下的經文.
亦是當中的訊息.
今天跟大家講的題目叫做甚麼呢.
做個好男人.
做個好女人.
就是大衛與亞比蓋的事.
這個就是上,下何時講呢.
下個禮拜,下個禮拜是母親節.
所以今天我們先講做個好男人.
下個禮拜呢.
做個好女人.
不過其實不是的,因為聖經的真理.
是不只是為男.
或者不只是為女的,當然有些是.
但是特別是今天這個.
大衛與亞比蓋他們的經歷.
事實上我們兩個都看過有很多的特質.
我們可以學習,無論是男還是女.
所以我們一起做個,甚麼.
好男人,又一起做個好女人.
三位機上二十五章.
剛剛是在二十四和二十六章之間.
如果大家記得上一次.
我們就講二十四和二十六章.
就講到大衛被誰追殺.
被蘇勞追殺.
但是他兩次,在二十四章,二十六章.
兩次有機會,都可以.
殺了蘇勞.
好像天賜給他一個良機.
其他人都這樣看.
但是大衛他不是用普通人的眼光去看.
他是以神的眼光去看.
他認為這個不是神的心意.
所以雖然我們說.
大衛是面對.
蘇勞這些無理的追殺和迫害.
但是他沒有以惡報惡.
相反是以甚麼.
以善報惡,這就是上一次.

$^{41}$我們跟大家分享的訊息,他以善報惡.
那麼.
這個學者Berlin,他亦都將.
二十五章,在二十四和二十六章之間.
他認為有一個神學的連貫性.
因為這裡同樣的.
大衛是遭受一種無理的對待.
但是,按人.
這裡很明顯將他人性表達出來.
按人性,他就是要報仇報復的.
按人性,他就是要.
你有半斤,那又如何.
我也有八兩,所以他本來.
他人性那方面,要來殺敵人.
但是很奇妙,同時.
同樣的,有神的提醒.
讓大衛轉變他的心意.
和二十四章,二十六章那種.
我們發覺人的靈命,的確有高有低.
他回復回二十四和二十六章.
那種心態,就是他不以.
武力來除掉敵人.
相反,以神的方法和幫助.
以善去對待這些對他不好的人.
所以這裡.
好像有一種連貫的意義.
和神學意義給我們看見.
所以這次我們看二十五章.
我們看到他仍然在逃避掃羅追殺之間.
既然被逃亡,物質上.
是有需要的,所以他就向.
其中一個配角,很重要的配角.
就是這個大富戶,名字叫.
拿伯,拿伯希望請求.
他的幫助,但是.
遭遇到整個故事的發展.
我們看見一會兒,就顯出一個.
壞的惡的男人和一個.
好的男人的分別,所以我們看見.
在這裡,其實當時這個背景就是.

$^{81}$以色列遭受很多外敵,時常有.
非利士人,或者亞馬力人.
的攻擊,所以很多時候.
國家的領袖都派遣一些將領.
或者兵士,去保護自己的屬地.
很明顯的,駐軍在那裡,或者.
去確保那裡的人民.
平安,保護他們.
所以被派到各處去守護這些牧場.
特別是以畜牧為主,免受.
侵略者,敵人和匪徒.
的迫害,或者侵害.
所以這個故事的背景就在這裡.
我們讀一讀這段經文,我們由.
第一節讀到很長的,其實.
不過我們今天只選到十七節,有時間.
等一下節目可以回來看完全章,我們看一看.
第一節的下半節開始,我們用啟營的.
方式,我讀一節,請會眾讀一節.
我們一起讀完這段經文,第一節.
大衛就起身,下到巴蘭的.
曠野,請讀.
人名叫拿伯,是.
加勒族的人,加勒族是其中一個.
以色列當時覆蓋,管理.
地方,其中一個民族.
他們是以色列人.
他的妻名叫阿比蓋,這個就是主角.
阿比蓋,聖經形容她.
是聰明俊美的婦人.
相反,拿伯就是.
她的丈夫,為人剛.
迫凶惡,剛剛相反.
的形容,一個就是聰明.
俊美,一個就是剛迫.
凶惡,所以剛剛你看見.
非常不相配.
第四節,請讀.
願你家平安.
願你一切所有的都平安.

$^{121}$就問他們安,問了幾次.
三次,很少問三次.
表示很恭敬,一種祝福語.
願你平安,已經很好了.
願你的家平安,已經更好了.
還不止,願你一切所有的.
都平安,這個真是一個很祝福.
很有禮貌,很謙卑.
帶著一種和平的心.
來祝福他們.
相反,我們見到第七節.
第七節,請聽我讀.
現在我聽說.
有人為你剪羊毛.
這個就是拿伯所說的.
你的牧人在加密的時候.
和我們在一處.
我們沒有欺負他們.
他們也未曾失落什麼.
好,第八節,請讀.
我可以問你們牧人.
他們要告訴你.
所以願望你們牧人.
在你們眼前同甘.
願你們身受好日子,多謝您.
求您順手順意.
主佑我們每一天.
所以這裡.
大衛的牧人跟拿伯說.
你問問你的牧人.
就是我們沒有欺負他們.
使他們失去什麼.
反而我們時常去保護他們.
這個就是他們一向以來.
幫助拿伯和他的產業.
但是第九節.
大衛的牧人到了.
將這話題大衛的名都告訴了拿伯.
就住了口,所以這個就是大衛吩咐.
他的牧人去到拿伯的時候.

$^{161}$就照這樣說.
第十節就是拿伯的回應.
我們一起請你們讀.
所以這裡.
十一節.
這裡就是拿伯的回應.
請問他用了多少句說話回應他們.
大衛思緒.
然後耶西的兒子.
這裡有少少侮辱性.
說了他爸爸的名字出來.
耶西的兒子是誰.
連你爸爸的兒子是誰.
然後近來博奕主人逃跑的瀑人.
甚多.
博奕主人的瀑人.
我怎可以將我的飲食和子母所載的肉.
給這些我不知道的人.
請問這幾句說話帶著什麼語氣.
藐視侮辱.
一種語氣.
所以本來他們帶著三個問安語.
現在回覆就是三個.
侮辱和.
輕視的說話.
這個問號其實不是問.
是一種挑戰和鄙視的態度.
十二節.
大衛的瀑人就轉身從原路回去.
招者告訴大衛.
大衛的反應是怎樣.
十三節請你們讀.
有喇叭的瀑人告訴.
喇叭的妻子給該說.
大衛從曠野打發.
使者來問我主人的安.
主人卻辱罵他們.
結果大衛就說.
豈有此理.
將他殺了.

$^{201}$各人拿刀.
留下二百人.
二百人守住大本營.
四百人要將他殺了.
這個消息一到.
那些瀑人馬上就不告訴喇叭.
告訴誰呢.
告訴喇叭的妻子阿比蓋.
阿比蓋聽到這些瀑人.
這個壞丈夫竟然辱罵他們.
這瀑人再加上一個註腳.
第十五節.
但是那些人待我們甚好.
這就是事實.
我們在田野和他們來往的時候.
沒有受他們的欺負不止.
也未失落什麼.
還有十六節請讀.
好不好啊.
待我們好不好.
不但沒有欺負我們還成為我們的保護.
這個喇叭這樣對待他們.
真是以惡報善.
以惡報善.
十七節就是最後一節.
我們一齊讀.
看怎樣行才好,不然和萬定要臨到我主人和他全家,他性情凶暴,無人敢與他說話..
這個就是喇叭.
所以按當時的背景.
本來在慶祝的時候.
因為他們在剪羊毛的時候.
慶祝的時候一起吃喝,所以有肉分.
這個不明文的做法.
就像今天我們這樣.
那些人服侍我們,幫助我們.
在這個時候既然分了.
分一份給他們.
就好像我們今天給什麼.
給一些貼士.
雖然不是明文的做法.

$^{241}$但是卻是很普遍的一種做法.
但是我們見到這個喇叭.
真是非常的.
吝嗇.
所以吝嗇.
我們廣東話是什麼.
所以稍後節目不要吝嗇.
有時候我發覺我們不知為何.
去吃東西的時候.
對於服務員.
很吝嗇.
特別去到華人餐館.
有幾次見到.
西人餐館沒空位,拿張單來,現在幾個百分比.
十五,十八.
現在十五,十八.
但是不知為何.
有時候我們中國人.
你去華人餐館.
這裡不用給,十個百分比.
十二個百分比.
為什麼西人餐館給十八,中國人餐館給十二呢.
請問.
吝嗇.
我覺得.
我們應該學習這些.
不要吝嗇.
其實這是一個美德.
所以我們看到.
吝嗇是一個以惡報善的情況.
大衛.
很有禮貌,用祝福.
問安.
所以我們看到這是一種不應該有的.
侮辱在他身上.
吝嗇其實是以色列人.
所以在整個加勒族.
主要城市,加勒族在希伯倫.
希伯倫其中一個主要的城市.
一個小的地方,不過他是一個大富戶.

$^{281}$所以他在那裡可以說是稱王稱霸.
在氣派甚大.
不過很可惜.
他的人剛愎兇惡,不單止忘恩.
又很淪色.
他拒絕了大衛的要求.
和將他們侮辱一番.
印章.
印章仇報.
生氣嗎?.
所以大衛說豈有此理.
你知道大衛是何許人.
帶著的是何許人.
果然令他怒火中燒.
四百人準備將他全家滅絕.
當然這是他們應得有的結果.
也是帶著怒氣而發出的一種回應.
幸好有一個阿比蓋.
阿比蓋.
他知道環境,明白事理.
知道誰是對誰是錯.
他沒有爭著他的丈夫.
拿伯的惡夫.
真的不應該這樣對待別人.
所以他臨急置身化解危機.
逃過全家被殺的大能.
他自己在途中與大衛相遇.
這就是他們相遇之間的記載.
突出了大衛和拿伯之間的分別.
也突出了大衛與阿比蓋相遇.
大家彼此欣賞.
以致後來當神擊殺了拿伯.
拿伯罪有應得.
他很惡.
不過一知道有禍的時候.
怕得全身僵硬.
被神擊殺.
所以惡有惡報.
大衛以善報惡.
結果善有善報.

$^{321}$結果他成功娶了一個非常美德.
非常好的女人.
這個人叫什麼名字.
就是阿比蓋.
所以我們今天來看一看.
聖經的記載.
的確在拿伯和大委身上有很多特點和形容.
我自己做研究做了一個對比.
在我的書裡面.
如果大家有的話可以去看第三章.
講到大衛與阿比蓋.
在那裡我列了一個表.
就是表310.
如果大家有我的書可以去看表310.
講到大衛和阿比蓋.
拿伯和大衛的一個比較.
書裡面很詳細.
這裡我只是簡單和大家分享.
第一我們看到拿伯有什麼特點.
聖經形容他.
可見他是一個壞男人.
真的不要學他.
第一就是對人侮辱.
我剛才講過.
人家好好的問你安多少次.
三次.
你回答他也是三次.
不過不是問安是侮辱的說話.
所以後來大衛都自己這樣講.
連拿伯的服人都說他辱罵大衛.
拿伯死後大衛回應.
如果大家回去看25章39節.
當稱頌耶和華.
因他伸了拿伯辱罵我的冤.
所以他深深感受到.
拿伯對他的辱罵.
然後神幫他伸冤.
這真是大衛的感受.
又顯出對人侮辱的惡性情.
但相反大衛對人祝福.

$^{361}$所以他是三層祝福.
吩咐他的服人.
到拿伯的時候要祝福他.
祝福他的家 祝福他所有的一切.
三層的祝福.
後來他對阿彼戈.
如果大家回去自己詳細看那章經文.
他也對阿彼戈有三層的稱讚.
還稱讚上帝.
稱讚阿彼戈的明智和才德.
也確認阿彼戈採取這個行動.
是大利益處.
所以大衛也是時常對人祝福.
弟兄姊妹.
無論男人女人.
特別是男人.
做個好男人.
我們應該對人祝福.
多些對人祝福.
不要對人時常辱罵貶低.
或者侮辱.
所以這個實在好與壞再次提出.
我們看到拿伯的特點是什麼.
以惡報善.
這個我們剛才說過.
但是相反我們看到.
大衛卻以善待人.
他不單止在這個衝突的時候.
更加之前.
之前已經很友善的去保護拿伯他們的家產.
保護他們的姓名.
以致他們沒有失落.
不單止仍然得到他們的保障.
這個就是以善待人.
你想做好人還是壞人.
我們當然不要以惡報善.
更加需要以善待人.
所以這個真是好.
第三我們看到.
這個惡男人.

$^{401}$就是目空一切.
正經形容他叫剛愎兇惡.
剛愎兇惡.
即是非常性情兇惡.
他不單止粗暴.
沒有禮貌的回答別人.
更加天不霸地不霸.
目空一切.
真是.
又不敬畏神.
又對人不禮貌.
所以這種就是目空一切.
相反我們看到.
大衛卻是敬神敬人.
因為大衛.
他是敬畏神.
願意放下自己的面子.
衝動接受勸告.
等一下我們會說.
他更加不斷的稱讚耶和華.
稱讚神.
難阻他殺無辜人的血的罪.
所以在這裡.
見到他對神的敬畏.
而且我們見到.
他回答亞比該的說話.
大家如果回去看看.
你會發覺他在思想和言辭當中.
不斷以耶和華為中心.
他事實上想到.
耶和華怎樣看神的心意.
是甚麼.
他回答剛愎的那八.
那些侮辱的說話.
他用甚麼來說呢.
他自稱.
你兒子大衛.
所以他謙卑他自己.
你兒子大衛現在向你請求.
所以我們見到.

$^{441}$他對神也好對人也好.
都是一種尊敬的心.
一種敬愛敬畏的心.
這是很值得我們學習的.
你想做好還是壞.
壞的人目空一切.
好的男人.
甚麼.
敬神敬人.
這個就是分辨.
再看我們見到.
那八是甚麼.
神憎鬼厭.
為甚麼.
無人敢與他親近.
即是沒有人敢勸他.
沒有人敢走近他身邊.
只會被他罵.
所以.
你想做這些人嗎.
無人敢走近跟你說話.
因為個個都怕被你罵.
你想不想做一個好的男人.
相反我們見到.
這個就是人緣親近.
我們記得在《撒慕爾記》上22章.
當時有很多很多人.
犯受困逼的欠債的心裡.
勞苦的都聚集到大衛那裡.
當時他就做了.
這些眾人的頭目.
跟隨的人有多少人.
有400人.
後來再加上200人.
總共600人.
全部這些人來到大衛那裡.
願意親近他.
以他為他們的頭目.
所以你看到這個分別.
縱使大富豪那八.

$^{481}$他有很多僕人.
但無人敢與他講話.
因為他的性情.
剛愎兇惡.
但相反我們見到大衛.
有智有勇.
有大能的作為.
大能的勇士.
有很多人願意來到他旁邊.
接著我認為.
這個人叫做假皇帝.
因為他過著.
好像皇帝般的生活.
他奢華艷樂.
真的.
他臨死的時候.
還是奢華艷樂.
真的.
見到他這種.
好像皇帝的生活.
卻不是真正皇帝.
相反大衛是真天子.
他是和平的使者.
他是以色列的受高者.
他是真是君王.
弟兄姊妹.
我們不要做假基督徒.
要做真天子.
因為你是神的兒女.
你是天國的子民.
讓我們都能夠.
好像大衛活出這種.
尊貴的生命的品格.
能夠顯出.
他的天國子民的身份.
當然他就是那位受高者.
然後我們見到.
最後他不懂得欣賞.
真的不懂得欣賞.
有一個這麼好的太太在他身邊.

$^{521}$竭力維護他.
私下來為他講說話.
但是他自己.
只知吃喝爛醉.
艷樂.
來過他的人生.
他沒有一些來欣賞.
太太在他身邊的勞苦.
只是.
自己顧自己.
自己自私自利的.
過他的人生.
相反我們見到.
大衛善於欣賞.
他稱讚阿彼該.
有不同方面對他稱讚.
所以當藍靶死.
大衛立刻.
這麼好的女人.
怎麼可以浪費呢.
不要讓鮮花插在甚麼.
問鼎姐妹你想做哪一邊.
你想像拿八一樣.
還是像大衛一樣.
求主幫助我們.
在大衛的身上.
特別講到男人.
其實女人都是一樣.
我們可以應用在幾方面.
和大家.
諮詢一下.
第一.
大衛的特點.
快適路.
三十二到三十五節.
大衛對阿彼該說.
以色列的上帝是應當稱讚的.
你看看內容.
一樣是.
神做他的中心.

$^{561}$因為神今天.
使你來迎接我.
你和你的見識.
應當稱讚因為你今天攔住我.
親手報仇流人的血.
我只著.
阻止我加害於你的耶和華.
又要提耶和華.
以色列永生的上帝.
速速地來迎接我.
到明日早晨.
凡屬拿八的男丁必定不留一個.
三十五節.
大衛受阿彼該送來的禮物.
就對他說.
我聽了你的話 準了你的情面.
你可以平平安安地回家去了.
嘩 本來有個危機 本來有個殺機.
現在轉成平平安安的回家去.
這裡反映一件事.
大衛.
快適路.
所以本來很生氣.
生氣得怎樣.
不單殺拿八 要殺什麼.
殺他所有人.
很生氣.
不單丟一個碗.
全國.
全國貴的碗都丟了.
很生氣.
所以你看到.
這裡就是.
當時大衛那種.
情緒表現.
突然間我們很明顯看到.
準了那個.
很明顯剛才讀到這幾節.
特別是最後這一節.
他說你平平安安地回家去了.

$^{601}$請問大衛還殺不殺拿八.
不殺.
轉變快不快.
很快.
情緒本來很激昂.
要動殺機.
不單是意念.
拉隊正在走.
在半路上.
但是在半路上有截住.
半路上收回來.
這個真的難.
未出手之前.
收回來容易一點.
出到一半收回來很難.
大衛的特點就在這裡.
他一以遙和為中心.
一想到神的心意.
一明白是理.
一冷靜下來就收.
所以這個.
顯出大衛.
大衛.
這種值得我們學習的地方.
有沒有發過憤啊.
大英姐妹.
沒有人敢說.
沒有人敢說即是個個都是.
誰會不發過憤呢.
有沒有試過很憤啊.
有沒有試過憤得想動手啊.
有沒有試過憤得想動手.
中間被聖靈提醒.
收回啊.
又有啊.
這個就是我們不斷人生去學的.
激嬲的經常都有.
聖經這樣說.
詩篇三十七篇.
請讀.

$^{641}$當止住怒氣.
離棄憤怒.
不要心懷不平.
以致作惡.
因為憤怒.
很容易就會變成.
變成作惡.
因為過份.
去回應的時候就作惡.
以惡報惡.
所以我們見到.
他克服了他的怒氣.
避免了.
因為自尊心的緣故.
當時戴維帶著四百人.
個個看著他.
說收就收.
既然出口.
說殺就要殺.
但是他沒有因為.
他的自尊.
這麼壞.
就走回頭路.
他能夠收回.
他的怒氣.
這個不至於殺無辜人的血.
其實那個人都不是無辜的.
不過拿破崙的親屬是無辜的.
所以.
他避免了這件事.
所以他在這裡.
真的沒有作惡.
因為他回復了上帝的心意.
的思想.
做他的決定.
所以我們見到.
這個就是特別的地方.
說起情緒控制.
男女有時會有些不同.
當然.

$^{681}$我們有些性格背景.
家裡是不是.
以前的原生家庭.
經常都吵吵鬧鬧的.
出刀的.
我以前的外婆.
我記得我還有少少印象.
她真的拿菜刀去砍我的阿姨.
關了門.
砍在她的門裡.
砍下去.
不過也有男女之分.
上幾個禮拜.
李師母帶領了一個聚會.
叫做.
還記不記得.
叫做男女之別.
我們分開男和女分組.
來聊天.
男就盡吐心中言.
講女人怎樣怎樣.
女人又盡吐心中言.
講什麼.
講男人男人男人.
當然就不只是.
扁的.
都有煲.
男人的好處.
男人有什麼好處.
其中一樣.
他夠力.
粗重的事他做.
所以.
粗枝大葉.
要有力的.
男人就搞定.
大膽一點.
師母前兩天才大聲叫.
蜘蛛呀.
誰出動.

$^{721}$當然我出動.
蜘蛛俠出動.
那麼.
但是男人因為.
比較粗暴一點.
蜘蛛就死了.
有時候也容易脾氣.
容易爆火.
當然我不是絕對.
有些弟兄很好脾氣.
有些姐妹也不簡單.
但是無論如何.
我們都知道.
能夠有一個溫柔的心.
有怒氣的時候.
能夠可以止住.
能夠可以.
得到.
以神的心去思想.
免除一些.
更嚴重的後果.
這個都是明智的.
在一本書叫.
Love and Anger in Marriage.
這是David Mace所寫的.
裡面他就用.
動物天然本性.
他說面對衝突的時候.
動物就會.
第一就是fight.
這個是本能.
你生氣就罵我.
我又很生氣就罵你.
這個fight.
作戰狀態.
有時候動物也有.
走 明知不夠距離.
走人.
也有的 有時候走.
走其實就是avoidance.

$^{761}$有時候未必一定是最好的.
應該找些方法去調解.
會比較好.
動物也會freeze.
freeze即是明知不夠距離.
裝作扮懵.
一如另外一張avoidance.
免得跟你.
裝作無事 全不招架.
其實最好的方法.
就是.
David Mace說.
我們應該學習relax.
relax的意思就是說.
這個時候可能真的需要祈禱.
深呼吸 數由1.
數到100.
然後就.
不要馬上回應.
他說有幾種.
對於anger的態度.
第一叫做expressing anger.
expressing anger就是說.
你很憤怒 怎樣爆出來.
很憤怒爆出來.
可以用言語 可以用你的表情.
可以用你的動作 可以用擲門.
可以用爆門.
這些都是expressing.
就是讓對方知道你很憤怒.
但是他說.
其實另外一種處理anger.
就是venting.
venting即是什麼叫vent.
就是讓氣.
壓力煲.
壓力煲 到很大壓力的時候.
怎樣.
放氣.
放出氣.

$^{801}$那怎樣放氣呢.
其實這個時候就是說.
我們需要祈禱.
主下幫助我.
這個時候就是跟神說.
你跟人說很憤怒.
即是有限.
你跟神說 神是很憤怒.
真的很憤怒.
所以這個venting.
不是爆出來.
所以就跟上帝說.
第三.
processing.
processing即是說.
你要想一想 為什麼我會這麼憤怒.
那時候就需要祈禱.
聖靈幫助我們.
理智回復.
想一想為什麼我會這麼憤怒.
你踩到我哪一樣東西 按到我哪個按鈕.
令到我這麼憤怒.
究竟這件事我值不值得這麼憤怒.
我如果回應的時候.
後果帶出來是什麼.
所以processing.
大衛也是這樣.
大衛就processing.
他就想一想.
如果我這樣做會殺無關人的血.
如果我這樣做.
將來登基的時候.
真的做皇帝的時候.
就會有很不好的記錄.
所以他就process了之後.
覺得是.
所以當堂的氣就順了.
所以這個.
真的是.
很好的提醒.

$^{841}$快適路.
真的是一種.
值得我們學習的表現.
轉個角度想一想.
當堂.
就是海闊天空.
因為process了之後.
其實很小事.
不要緊 誤會而已.
有機會澄清就可以了.
所以這個就是.
聖經給我們提醒.
這是新約的.
雅各書一章智慧書.
我親愛的弟兄們.
這是你們所知道的.
但你們各人要.
快快的聽.
慢慢的說.
慢慢的生氣.
因為.
人的怒氣.
並不成就神的義.
因為可能會令事件更差.
也不能夠帶來.
神的.
作為和見證.
因為這種怒氣.
只會傷害.
所以在這裡說.
秘訣就是快快的聽.
快快的聽就是聽清楚.
快快的意思不是隨便的意思.
快快的意思就是.
聽得很清楚.
掌握聽清楚每一句說話.
快快就是不要漏了每一句說話.
所以聽清楚.
未聽清楚就爆.
越誤會加深.

$^{881}$衝突就更大.
所以聽清楚.
原來不是這個意思.
問清楚.
所以就不要這麼生氣.
然後慢慢的動腦.
讓他生氣的時候.
都應該生氣.
拖緩時間.
讓自己想清楚.
多些機會溝通.
然後以弗所恕.
四章二十六二十七節請讀.
一起讀.
生氣卻不要犯罪.
不可含怒到日落.
也不可給魔鬼留地步.
請問可不可以生氣.
這句你答得這麼大聲.
生氣是可以的.
因為我們在說這個人的.
情緒的表達.
上帝賜給我們的.
有喜怒哀樂.
沒有feeling就糟了.
所以生氣是可以的.
不過不要犯罪.
不要過了火.
的意思.
更加不要含怒到日落.
所以如果你生氣.
就麻煩你早上生氣.
起碼要八個小時.
不要晚上生氣.
因為沒有時間.
這裡當然不是這個意思.
盡快取消它.
平息它.
因為魔鬼最適用這些機會.
來破壞家人與家人的關係.

$^{921}$人與人之間的關係.
破壞弟兄姊妹之間的關係.
所以不要讓魔鬼留地步.
然後回到舊約的智慧書.
一起讀.
不輕而發怒的.
勝過勇士.
自服己心的.
強如取勝.
這句說話真好.
當你慢慢操練.
每個人都有怒氣.
包括我.
但是我們去操練.
讀聖經的智慧.
快快的聽.
慢慢的說.
慢慢的動腦.
然後生氣.
不要讓它成為犯罪.
慢慢操練聖靈幫助我們.
好像大衛的例子.
幫助我們的話.
能夠process.
正確方法.
能夠process正確方法.
這些怒氣的時候.
聖經裡面說.
不輕而發怒.
不只是發怒.
不輕而發怒.
這個就好了.
所以不是天生的好脾氣.
而是它已經操練到.
不是那麼容易.
就會很快很生氣.
因為他懂得去想一想.
懂得process.
很快就process.
很快就釋怒.

$^{961}$不讓他那麼快用情緒.
控制了我們整個人.
所以他不輕而發怒.
勝過勇士.
自服己心.
強如取勝.
這個就是.
非常好.
我們基督徒在神的話語當中.
得到提醒.
恨的傷患是什麼.
愛.
所以其實怒氣.
能夠克服.
愛是一個很大的原因.
當你想到.
那麼生氣幹什麼.
我仍然很愛他.
那就不要生氣了.
所以這個就是幫助我們.
好像大衛一樣.
快釋怒.
第二.
就是肯聽勸.
大衛肯聽勸.
二十五章的三十五節.
我們再讀一次.
大衛受了亞彼該送來的禮物.
就對他說.
我聽了你的話.
準了你的情面.
你可以平平安安地.
回家去吧.
那次我們.
前幾個禮拜的喜樂團日聚會.
有些姐妹就讚.
讚她先生.
他很聽我話的.
不過有很多.
有很多回應就說.

$^{1001}$他不聽人說的.
那些.
很多.
原來我們也有些.
丈夫是很聽話的.
不過也有很多的丈夫.
是不聽人說的.
不單止.
有幾個評語說.
他永遠都不肯認錯的.
嘩你這麼大聲的回應.
因為有很多人共鳴.
大衛真好.
大衛自己說出口.
跟亞彼該說.
我聽了你的話.
準了你的情面.
這裡給你看見.
他不是一種大男人的情況.
蔡元雲醫生.
我們知道她是一位輔導.
很有.
專業的.
這方面的.
寫了很多書關乎這方面.
做了很多個案.
她其中一本書.
講到男人的時候.
她就提到男人很多時候.
以強人的姿態出現.
男性的尊嚴不容易放下.
很容易洗刷.
她說她自己是男性.
所以她這樣說.
很公平.
但是我們又發覺.
女性比較重感性.
比較細心.
有時候她的直覺.
感情.

$^{1041}$有時候我們男性發生了什麼事.
都不太清楚.
因為不太敏銳於別人的感覺.
男性.
相反很能夠分析.
很理智.
有時候姊妹的事.
要分析.
要理智一點.
男性比較容易理智分析.
看大圖比較多.
但是其實真的互相補足.
所以.
真的有時候女要聽什麼.
要聽男的說話.
因為她有時候看得遠一點.
看得大一點.
看得理智一點.
明白事理分析得好一點.
但是相反男要聽什麼.
你們不出聲.
男要聽.
女的說話.
因為有時候她的敏銳情況.
她的感情的了解.
可能真的成為我們.
很大的幫助.
所以.
有時候特別說到一些.
人的表現.
性格那些東西的時候.
可能真的女性的感受.
真實的感受.
是帶來一些.
很重要的參考.
不過.
女性也有一個問題.
說完一次她會再.
又再說多一次.
重複.

$^{1081}$所以男性很怕她.
什麼話.
你是自己說的.
所以男性很怕她囉嗦.
所以姊妹請你留意.
勸導丈夫的時候.
不要太.
囉嗦.
不過弟兄.
請你記得.
原來姊妹囉嗦的時候.
是因為我們時常都.
不肯聽她說.
所以.
這裡我們又覺得.
都很公平的.
彼此學習.
大衛我覺得.
他能夠這樣聽.
勸告.
是因為他有一個受教的心.
這個是一個關鍵.
弟兄姊妹.
你有沒有受教的心.
有.
因為他習慣了在神面前.
有一個受教的心.
所以他願意當人有.
一些東西提出來的時候.
他不會馬上說不行.
或者馬上就會.
認為.
自己對.
我們每個人都有.
有些比較嚴重.
沒有這麼嚴重.
不過我們每個人都學習.
很多時候是因為他.
那種受教謙卑的心不夠.
所以他在.

$^{1121}$人神面前.
不是很習慣受教.
去到人面前.
都不容易受教.
大衛相反.
我們見到他有一個受教的心.
他第一句回應亞比哥.
就是耶和華意識你的上帝.
是應當稱重的.
這是他第一句說話.
顯明他的內心.
反射出來,他的心態.
所以弟兄姊妹.
那些技巧,那些道理.
我們都明白的.
跟丈夫男女相處.
但是這個是關鍵.
就是我們有沒有在神面前.
有一個謙卑受教的心.
第三我們見到.
大衛他就是懂欣賞.
很明顯.
在top step.
上25章33節.
我們一起讀33節.
你和你的見識.
也當稱讚.
因為你今天難了.
我親手報仇.
留人的緣.
所以這裡我只是說一句.
因為時間關係.
你回去自己看一下.
他有三重稱讚.
稱讚亞比哥.
這件事,稱讚他的才德.
稱讚他的做法.
能夠幫助到他.
沒有犯下大的錯.
所以原來這個是懂稱讚.

$^{1161}$這個真的重要.
但又不容易學的.
特別我們中國人很含蓄.
不容易稱讚別人.
明明自己很崇高.
得到幫助也好.
不容易稱讚.
特別丈夫.
不容易稱讚太太.
太太一片苦心.
換來.
有片.
沒有煲.
煲有,煲湯就有.
煲飯有.
不過我每一次.
師母煮完飯.
我一定多謝她.
多謝.
事實上.
下次有好東西吃.
哈哈哈哈.
其實.
不是為了這樣.
從心裡稱讚.
真的好,多謝.
我前兩天先多謝師母.
稱讚她.
那天你在大祈禱的時候.
我很有感動.
你的祈禱.
很有感動,我相信會對頂尖很有幫助.
所以.
我就稱讚她的祈禱.
是的,的確.
很多東西都可以稱讚.
我們可以.
去留意的時候.
就有了.
嗯.

$^{1201}$其實阿彼該.
也是,如果你回去讀.
所以不只是男人,女人也是.
他很懂得稱讚大衛.
他就說,你知不知道.
你是神的受告者.
他說你.
不應該和這些惡人.
以惡報惡.
等等.
他稱讚大衛.
你有遠大的眼光.
所以我們看到.
這個大衛和阿彼該.
真是相配.
在這方面,我們看到他們互相稱讚.
大衛懂得.
欣賞阿彼該.
所以一下子,這個男人死了之後.
第一時間.
這個不可以放過的好太太.
就馬上將他取為.
妻.
在這個「愛就是彼此珍惜」.
這本書裡面.
提出了七項.
在我的書裡面有引用的.
「愛就是彼此珍惜」.
其中這七項.
我們很快看看.
第一,就是愛是什麼.
愛就是懂得在生活的細節上.
欣賞對方的一舉一動.
一鵬一笑.
所以就多點欣賞.
欣賞一下他.
第二,就是愛.
是能細心體會對方.
為你所做的一切.
並感激他,有沒有.

$^{1241}$這麼小聲.
學一下,我們應該學一下.
你想男女之間的關係.
特別夫婦之間的關係.
好啊,多點感激.
第三.
愛是願意.
花時間與對方在一起.
也不吝惜於.
告訴對方.
你愛他.
真好啊,你告訴他.
我很珍惜你,你真好.
我愛你,這個真是一個.
很好的言語.
還有用時間.
陪伴上去表達出的愛.
第四,愛是承認.
兩人都不完全,是不是.
但是能夠.
接納彼此的缺點.
讓對方做真正的自己.
給他一點空間.
然後讓他.
繼續成長,接納.
他的缺點,不要經常針對他.
所以反而.
多點稱讚他,是好的地方.
第五,我們看到.
愛是熱情.
加上友情.
培養友情,大家有些共同的嗜好.
傾談,分享.
一起祈禱.
激情加上同情.
結婚之前.
就說很熱情,結婚之後.
就冷冰冰.
無情.
但是我們仍然可以.

$^{1281}$回復溫馨之情.
這是真愛的表現,Amen.
然後第六.
愛是.
愛是懂得傾聽.
就是聽,不要太快反駁.
不要太快發表.
愛是懂得傾聽.
了解對方的需要.
聽對方.
心裡面的說話,不要太快.
就馬上論斷他.
所以這個也是.
我們不斷要學習.
不斷要學習.
第七,愛.
是能夠與對方享受一起.
歡愉的時刻.
如果到達一個境地,兩人必須.
經常練習正面的互動.
就好像.
黃維仁博士說的.
愛就是.
要存款.
什麼叫存款.
就是這段爭議.
稱讚多一點,感激多一點.
因為為什麼你經常稱讚他,感激多一點.
當你有時提起他.
罵他一句的時候.
就算了.
因為存款夠.
存款夠,提款的時候.
就不會.
虧空.
供款,透支.
所以求神幫助我們多一點存款,表達愛意.
所以例如早上有個微笑.
言語多一點,謝謝.
體貼對方,辛苦了.

$^{1321}$一天,這麼辛苦煮一餐.
謝謝.
特意幫我買東西.
從來都不用.
入油,都是你搞定.
很開心.
這種是愛的表達.
用言語表達.
用你的感激.
稱讚表達.
讓我們.
一起溫習.
讓我們像大衛這次.
讓我們看見,不單是好男人.
也是好女人.
對人祝福,而不是.
對人侮辱.
見待人,不要以惡報惡.
敬神敬人.
不要.
目空一切,不看一世.
接著那三點.
三個字,第一.
快適路,第二.
肯聽勸.
第三.
動人上.
求主幫助我們.
讓我們一起.
拍照拍掌.
不用這麼含蓄.
因為我們鼓勵.
跟旁邊的弟兄姊妹說.
做個好男人.
我們一起祈禱.
我們讚美天上的父神.
耀華的名是應當稱頌的.
多謝你.
在聖經裡給我們很多智慧.
透過.

$^{1361}$撒米爾記二十五章給我們很多提醒.
讓我們在當中學習.
一些美好的榜樣.
一生追求成為別人的祝福.
奉耶穌基督的名.
\newpage



\section{約書亞記 5:2-12}
\label{sec:032101}
\textbf{聖經人聖經事:入門三件事}~~~~黃克勤牧師博士
\newline
\newline
連結: \href{http://www.church.com.hk/acms/content.asp?site=cdc\&op=show\&type=product\&code=032101\&layout=sermon}{\texttt{032101}} ~~~~ 語音日期: 2026-05-08
\newline
\newline
\hyperref[sec:032100]{\small{< < < PREV SERMON < < <}}
~
\hyperref[sec:index_chronic]{\small{[返順時目]}}
~
\hyperref[sec:032102]{\small{> > > NEXT SERMON > > >}}
\newline
\newline
約書亞記 5:2-12
\newline
\begin{longtable}{cl}
\hline
\hline
章節 & 經文 (和合本修訂版)\\
\hline
5:2 & \begin{tabularx}{0.7\textwidth}{X} 那時，耶和華對約書亞說：「你要造火石刀，第二次為以色列人行割禮。」 \end{tabularx} \\ \\ \relax
5:3 & \begin{tabularx}{0.7\textwidth}{X} 約書亞就造了火石刀，在哈爾拉勒山為以色列人行割禮。 \end{tabularx} \\ \\ \relax
5:4 & \begin{tabularx}{0.7\textwidth}{X} 約書亞行割禮的原因是這樣：從埃及出來的眾百姓，所有能打仗的男丁，出了埃及以後，都死在曠野的路上。 \end{tabularx} \\ \\ \relax
5:5 & \begin{tabularx}{0.7\textwidth}{X} 這些從埃及出來的眾百姓都受過割禮；但是那些出埃及以後，在曠野的路上所生的眾百姓卻沒有受過割禮。 \end{tabularx} \\ \\ \relax
5:6 & \begin{tabularx}{0.7\textwidth}{X} 以色列人在曠野走了四十年，直到那從埃及出來，全國能打仗的人都消滅了，因為他們沒有聽從耶和華的話。耶和華曾向他們起誓，必不容許他們看見耶和華向他們列祖起誓要給我們的地，就是流奶與蜜之地。 \end{tabularx} \\ \\ \relax
5:7 & \begin{tabularx}{0.7\textwidth}{X} 他們的子孫，就是耶和華興起接續他們的，都沒有受過割禮；因為在路上他們沒有受割禮，約書亞就為他們行割禮。 \end{tabularx} \\ \\ \relax
5:8 & \begin{tabularx}{0.7\textwidth}{X} 全國的人都受了割禮，留在營中自己的地方，直到痊癒。 \end{tabularx} \\ \\ \relax
5:9 & \begin{tabularx}{0.7\textwidth}{X} 耶和華對約書亞說：「我今日將埃及的羞辱從你們身上除掉了。」因此，那地方名叫吉甲，直到今日。 \end{tabularx} \\ \\ \relax
5:10 & \begin{tabularx}{0.7\textwidth}{X} 以色列人在吉甲安營。正月十四日晚上，他們在耶利哥的平原守逾越節。 \end{tabularx} \\ \\ \relax
5:11 & \begin{tabularx}{0.7\textwidth}{X} 逾越節的第二日，他們吃了當地的出產，就在那一天，吃了無酵餅和烘過的穀物。 \end{tabularx} \\ \\ \relax
5:12 & \begin{tabularx}{0.7\textwidth}{X} 他們吃了當地出產的第二日，嗎哪就停止了。以色列人不再有嗎哪了。那一年，他們就吃迦南地的出產。 \end{tabularx} \\ \\
[1ex]
\hline
\hline
\end{longtable}
$^{1}$各位弟兄姊妹大家好.
今天跟大家談談入門三件事.
在生命樹宣教網絡.
跟大家一起讀聖經.
看看聖經人和聖經的故事.
看到上帝背後的工作.
入門三件事.
這些不是聖經的字.
是我作的字.
入門其實是指他們去到迦南地.
進到應許之地.
一進去就做了三件事.
那三件事是什麼呢?.
就跟大家看看.
也有意思的.
也是神叫他們做的.
入門三件事.
就好像開門三件事.
入門三件事.
約書亞記第五章第二節.
那時耶和華吩咐約書亞說.
你製造火石刀.
第二次給以色列人行割禮.
約書亞製造了火石刀.
再脫皮生那裡.
給以色列人行割禮.
大家可能留意到這次.
我用大隻的字.
如果我做powerpoint的時候.
這隻字是60級.
很大隻的字.
想大家透過螢光幕.
手機.
有些頂尖姐妹說.
她是用手機看的.
這隻字大隻一點好一點.
那時他們進去了.
耶和華叫約書亞.
吩咐他要做.
製造火石刀.

$^{41}$什麼叫火石刀呢.
其實應該叫岩石刀.
不是用打火那樣.
有火花閃出來.
用來點燃木材.
用來起火的那些刀.
那些火石不是.
其實是用石磨一把刀.
用石磨一把刀.
為什麼不用鐵的刀.
似乎神不是很喜歡用鐵的刀.
或者用鐵器.
或者譬如在會幕裡面.
有很多種器具.
譬如那些要堆起一個很大的祭壇.
祭壇可能是五米乘五米.
可能是不止的.
很大的祭壇.
柴放在地上燒祭身放在上面.
有斜路上去.
因此你知道底下有很多塊石頭堆下去.
才能堆起這麼大.
幾十平方尺的祭壇.
那怎麼做呢.
說明那些石頭在海裡拿回來.
或者在哪裡拿回來.
是不可以用鐵器碰過,釘過的.
不可以用鐵器接觸到的.
否則那塊石不能用.
可能上帝不是很喜歡鐵器,金屬.
撞過的那些東西.
叫他們用石來做.
火石這個字不算譯得很好.
叫他做石刀.
或者盤石做出來.
做一把刀磨出來.
是第二次以色列人行割禮.
這個割禮我想大家都知道.
割禮其實今天是一個男性.
性器官包皮割了他.

$^{81}$於是做這個割包皮的手術.
當時用石磨起一把鋒利的刀割了前面的包皮.
為什麼要這樣做呢.
因為上帝和亞伯拉罕說.
你每一個世代我和你立約.
你以後世世代代每一個以色列人.
你的後裔都要行割禮.
是我的標記.
是和我立約的標記.
所以這個都要做的.
什麼時候做呢.
通常是出世第八天行割禮.
但是這是第二次做.
他們這次是集體做的.
以色列人行割禮第二次.
原因是他們出埃及的時候.
走曠野的路.
走了四十年.
那一代人都做了才出來.
但是在路程上不方便之類.
他們就沒有行割禮.
來到加拿大第二次行割禮.
是集體全國人一起做的.
全部民族一起做的.
你知道當時他們上一代四十年前.
出埃及入曠野來加拿大.
那四十年裡面出埃及的人.
都倒閉在曠野了.
因此新一代是沒有行過割禮的.
現在就要做了.
第二次為以色列人行割禮.
入門第一件事就是行割禮.
保持那個記號與印記.
這樣的意思.
於是約書亞就做了「火石刀」.
其實就是石刀.
為什麼要中文譯作「火」字呢.
我都不算太清楚.
在這個徐皮山那旅級以色列人行割禮.
你找來找去聖經都沒有第二個徐皮山.

$^{121}$也不知道徐皮山在哪裡.
大概他們找一個地方.
那裡是全部百萬人行割禮.
是不是一把刀由約書亞做的呢.
應該不是的.
應該是每一個族裡面.
由族長長老.
他們會從所有自己支派的人做的.
如果一個人做就要很長時間了.
徐皮山由於集體.
所以那個地方就叫徐皮山.
徐皮山應該叫做Foreskin.
Foreskin徐皮是前面的皮.
橫形的意思.
就行了割禮.
你知道上帝很厲害.
叫他們第八天要行割禮.
出世所有男的第八天行割禮.
第八天行割禮.
那些嬰兒就不痛了.
很有趣.
大人行割禮痛得不得了.
大人就要痛很久了.
他們的子孫就是耶和華.
所興起來接續他們的.
都沒有受過割禮.
因為在路上沒有給他們行割禮.
約書亞借財給他們行了.
在路上沒有給他們行割禮.
聖經自己說的.
好了.
國民都受完了割禮.
住在營中自己的地方等到痊癒了.
這個痊癒起碼是傷口要痊癒.
而且痛.
我想沒有兩個星期.
一個星期不是很痊癒.
一兩個星期都不是很痊癒.
特別頭一個星期都很痛.
又流血.

$^{161}$而且我不知道會不會有些人有感染.
當時不清潔沒那麼好.
無菌技術差不多在抗疫不會有.
所以等到痊癒.
痊癒了之後.
在這裡有什麼意義呢?.
受完了割禮.
就是說他們這一代人.
進來迦南了.
這一代新的一代.
上帝是重新確認和他的關係.
要行割禮確認關係.
然後耶和華對約書亞說.
我今日將埃及的羞辱.
從你們中間令你們身上滾去了.
因此那地方叫吉甲.
吉甲的意思就是車輪滾動.
滾走直到今天叫吉甲.
其實在抗疫經過很多日子.
現在總算來到迦南的應許之地落腳了.
來到這裡第一件事要行割禮.
於是重新恢復.
以色列人和神的印記關係.
神要他們這樣做.
做完了然後神就說.
這些全部我認是我的子民.
一直都認的.
由現在開始.
他已經離開埃及來到迦南.
因此所有當日在埃及圍爐之間的羞辱.
都滾去了離開了.
重新確認和上主的關係.
因此他們受割禮的地方叫吉甲.
吉甲的意思就是滾動.
滾走就像車輪滾動.
直到今天.
所以開門第一件事.
入門迦南第一件事就是受割禮.
第二件事就是守禹節.
以色列人在吉甲安營.

$^{201}$正月十四日晚上在耶利哥的平原.
耶利哥城外面的平原.
第二件事就是守禹節.
入門第二件事就是守禹節.
禹節是這麼巧合的?.
是的.
就是他們出埃及的時候.
就是守禹節出來的.
禹節羔羊的血在門尾門腔上.
保守了一些家人.
Passover就是滅命天使不進去殺他們的長子.
他們就走出來.
紀念守禹節就是他們這樣出來的.
現在離開了遺老之地.
走了四十年來到這裡.
再守禹節做世代的定例.
再剛剛正月十四日又守.
在這個時候再記起我們出埃及.
現在來到這裡紀念神的工作.
禹節的次日他們開始吃了那地的土產.
正當那日吃無教餅和烤烘的菊.
什麼意思呢?.
去到禹節是要吃無教餅的.
他們以前在抗疫的時候的無教餅是怎樣的呢?.
就是有時停下來就要種一下地.
種一下麥.
或者到處找野生的麥.
又或者去到哪裡跟當地的民族買.
各樣都會有.
抗疫裡面有的.
這次是吃了那地的出產.
入門第二件事禹節是開始吃那地的出產.
出埃及的時候是吃抗疫裡面有的東西.
現在去到應許之地.
是吃那地的土產.
吃當地的麥子做的無教餅.
和菊都是烤熟的.
可以吃當地的東西.
留意一下首禹節第二件事.
他們第十二節.

$^{241}$他們吃了那地的土產的第二日.
即是說在他們首禹節入門第二件事.
首禹節的第二日.
第三點.
嗎哪就止住了.
嗎哪什麼時候停的?.
他們是在整個抗疫都吃嗎哪.
一直吃到入應許之地都在吃.
吃到什麼時候呢?.
吃到正月十四日禹節.
禹節他們就吃當地的土產做出來的無教餅.
吃完之後第二日.
嗎哪就止住了.
即是說沒有嗎哪了.
他們就開始吃當地的土產.
其實即是說神的供應是無微不至.
由整個抗疫四十年都吃嗎哪.
來到當地有了.
自己可以種到,買到.
好,就停止了嗎哪.
以色列人再也不再有嗎哪了.
那一年他們開始吃加拿地的出產.
開始吃當地的出產.
所以三件事.
第一.
是國禮入門第一件事.
第二.
守禹節.
第三件事.
就是吃當地的土產.
而嗎哪就停止了.
嗎哪停止了,這又是一個神蹟.
一直吃,天天都有.
還是多拿了,明天會爛.
還是安息日要多拿一天.
第二天又不爛.
經過四十年的神的神蹟供應.
很厲害.
在這裡他們就「嗎哪止住了,開始自食其力」.
神已經帶他出來.

$^{281}$走了四十年供應了.
現在進入了應許之地.
大家「自食其力」.
種東西,買東西.
開始「嗎哪止住了,自食其力」.
吃自己耕的土產.
這就是入門.
入加拿,入門之後三件事.
這三件事記得.
守國禮,守禹節.
然後「嗎哪停止,自食其力」.
我們一起禱告.
主我們見到你實在是.
做事很有步驟的.
在抗疫不方便.
於是你沒有行國禮.
其實你當天是說.
每一個人在第八日南午就要做國禮.
不過這次你讓他們保持到四十年.
進入了加拿大耶利哥外面的平原.
才開始做這件事.
另外我們再次提醒他們.
禹節是記起他們出埃及.
在埃及那種奴隸羞辱.
現在已經過去了.
完全過去了.
經過了一個抗疫的階段.
亦都好像一個人.
由維勞之家最愛的奴隸出來.
行人生的路.
經過抗疫很艱難.
然後入到應許之地.
好像我們經過死門.
然後去到永生當中.
在永生當中亦都不再有地上那些生活方式.
沒有嗎哪了.
在那裡是上帝無比的供應.
主我們讚美你.
你就是這樣恩代以色列人.
同樣又恩代我們.

$^{321}$讓我們更加感恩.
更加看你自己偉大的作為.
願你得著榮耀.
祈禱奉耶穌名求.
阿們.
\newpage



\section{提摩太前書 1:1-6:21}
\label{sec:032102}
\textbf{牧者手冊:提摩太前後書 (1)}~~~~劉少康牧師
\newline
\newline
連結: \href{http://www.church.com.hk/acms/content.asp?site=cdc\&op=show\&type=product\&code=032102\&layout=sermon}{\texttt{032102}} ~~~~ 語音日期: 2026-05-06
\newline
\newline
\hyperref[sec:032101]{\small{< < < PREV SERMON < < <}}
~
\hyperref[sec:index_chronic]{\small{[返順時目]}}
~
\hyperref[sec:032103]{\small{> > > NEXT SERMON > > >}}
\newline
\newline
提摩太前書 1:1-6:21
\newline
\begin{longtable}{cl}
\hline
\hline
章節 & 經文 (和合本修訂版)\\
\hline
1:1 & \begin{tabularx}{0.7\textwidth}{X} 奉我們的救主神，和我們的盼望基督耶穌的命令，作基督耶穌使徒的保羅， \end{tabularx} \\ \\ \relax
1:2 & \begin{tabularx}{0.7\textwidth}{X} 寫信給那因信主作我真兒子的提摩太。願恩惠、憐憫、平安從父神和我們主基督耶穌歸給你！ \end{tabularx} \\ \\ \relax
1:3 & \begin{tabularx}{0.7\textwidth}{X} 我往馬其頓去的時候，曾勸你留在以弗所，好囑咐某些人不可傳別的教義， \end{tabularx} \\ \\ \relax
1:4 & \begin{tabularx}{0.7\textwidth}{X} 也不要聽從無稽的傳說和冗長的家譜；這樣的事只會引起爭論，無助於神的計劃，這計劃是憑著信才能了解的。 \end{tabularx} \\ \\ \relax
1:5 & \begin{tabularx}{0.7\textwidth}{X} 但命令的目的就是愛；這愛是出於清潔的心、無愧的良心和無偽的信心。 \end{tabularx} \\ \\ \relax
1:6 & \begin{tabularx}{0.7\textwidth}{X} 有人偏離了這些而轉向空談， \end{tabularx} \\ \\ \relax
1:7 & \begin{tabularx}{0.7\textwidth}{X} 想要作律法教師，卻不明白自己所講的是甚麼，也不知道所主張的是甚麼。 \end{tabularx} \\ \\ \relax
1:8 & \begin{tabularx}{0.7\textwidth}{X} 我們知道，只要人善用律法，律法是好的； \end{tabularx} \\ \\ \relax
1:9 & \begin{tabularx}{0.7\textwidth}{X} 因為知道律法不是為義人訂立的，而是為不法和叛逆的，不虔誠和犯罪的，不聖潔和戀世俗的，弒父母和殺人的， \end{tabularx} \\ \\ \relax
1:10 & \begin{tabularx}{0.7\textwidth}{X} 犯淫亂和親男色的，拐賣人口和說謊話的，並起假誓的，或是為任何違背健全教義的事訂立的。 \end{tabularx} \\ \\ \relax
1:11 & \begin{tabularx}{0.7\textwidth}{X} 這是按照可稱頌、榮耀之神交託我的福音說的。 \end{tabularx} \\ \\ \relax
1:12 & \begin{tabularx}{0.7\textwidth}{X} 我感謝那賜給我力量的我們的主基督耶穌，因為他認為我可信任，派我服事他。 \end{tabularx} \\ \\ \relax
1:13 & \begin{tabularx}{0.7\textwidth}{X} 我從前是褻瀆、迫害、侮慢神的人；然而我還蒙了憐憫，因為我是在不信、不明白的時候做的。 \end{tabularx} \\ \\ \relax
1:14 & \begin{tabularx}{0.7\textwidth}{X} 而且我們的主的恩典格外豐盛，使我在基督耶穌裡有信心和愛心。 \end{tabularx} \\ \\ \relax
1:15 & \begin{tabularx}{0.7\textwidth}{X} 這話可信，值得完全接受：「基督耶穌到世上來是要拯救罪人」，而在罪人中我是個罪魁。 \end{tabularx} \\ \\ \relax
1:16 & \begin{tabularx}{0.7\textwidth}{X} 然而，我蒙了憐憫，好讓基督耶穌在我這罪魁身上顯明他完全的忍耐，給後來信他得永生的人做榜樣。 \end{tabularx} \\ \\ \relax
1:17 & \begin{tabularx}{0.7\textwidth}{X} 願尊貴、榮耀歸給永世的君王，那不朽壞、看不見、獨一的神，直到永永遠遠。阿們！ \end{tabularx} \\ \\ \relax
1:18 & \begin{tabularx}{0.7\textwidth}{X} 我兒提摩太啊，我照從前指著你的預言把這命令交託你，使你能藉著這些預言打那美好的仗， \end{tabularx} \\ \\ \relax
1:19 & \begin{tabularx}{0.7\textwidth}{X} 常存信心和無愧的良心。有些人丟棄良心，在信仰上觸了礁； \end{tabularx} \\ \\ \relax
1:20 & \begin{tabularx}{0.7\textwidth}{X} 其中有許米乃和亞歷山大，我已經把他們交給撒但，讓他們學會不再褻瀆。 \end{tabularx} \\ \\ \relax
2:1 & \begin{tabularx}{0.7\textwidth}{X} 所以，我勸你，首先要為人人祈求、禱告、代求、感謝； \end{tabularx} \\ \\ \relax
2:2 & \begin{tabularx}{0.7\textwidth}{X} 為君王和一切在位的，也要如此，使我們能夠敬虔端正地過平穩寧靜的生活。 \end{tabularx} \\ \\ \relax
2:3 & \begin{tabularx}{0.7\textwidth}{X} 這是好的，在我們的救主神面前可蒙悅納。 \end{tabularx} \\ \\ \relax
2:4 & \begin{tabularx}{0.7\textwidth}{X} 他願意人人得救，並得以認識真理。 \end{tabularx} \\ \\ \relax
2:5 & \begin{tabularx}{0.7\textwidth}{X} 因為只有一位神，在神和人之間也只有一位中保，是成為人的基督耶穌。 \end{tabularx} \\ \\ \relax
2:6 & \begin{tabularx}{0.7\textwidth}{X} 他獻上自己作人人的贖價；在適當的時候這事已經證實了。 \end{tabularx} \\ \\ \relax
2:7 & \begin{tabularx}{0.7\textwidth}{X} 我為此奉派作傳道，作使徒，在信仰和真理上作外邦人的教師。我說的是真話，不是說謊。 \end{tabularx} \\ \\ \relax
2:8 & \begin{tabularx}{0.7\textwidth}{X} 我希望男人舉起聖潔的手隨處禱告，不發怒，不爭論。 \end{tabularx} \\ \\ \relax
2:9 & \begin{tabularx}{0.7\textwidth}{X} 我也希望女人以端正、克制和合乎體統的服裝打扮自己，不以編髮、金飾、珍珠和名貴衣裳來打扮。 \end{tabularx} \\ \\ \relax
2:10 & \begin{tabularx}{0.7\textwidth}{X} 要有善行，這才與自稱為敬畏神的女人相稱。 \end{tabularx} \\ \\ \relax
2:11 & \begin{tabularx}{0.7\textwidth}{X} 女人要事事順服地安靜學習。 \end{tabularx} \\ \\ \relax
2:12 & \begin{tabularx}{0.7\textwidth}{X} 我不許女人教導，也不許她管轄男人，只要安靜。 \end{tabularx} \\ \\ \relax
2:13 & \begin{tabularx}{0.7\textwidth}{X} 因為亞當先被造，然後才是夏娃； \end{tabularx} \\ \\ \relax
2:14 & \begin{tabularx}{0.7\textwidth}{X} 亞當並沒有受騙，而是女人受騙，陷在過犯裡。 \end{tabularx} \\ \\ \relax
2:15 & \begin{tabularx}{0.7\textwidth}{X} 然而，女人若持守信心、愛心，又聖潔克制，就必藉著生產而得救。 \end{tabularx} \\ \\ \relax
3:1 & \begin{tabularx}{0.7\textwidth}{X} 「若有人想望監督的職分，他是在羨慕一件好事」，這話是可信的。 \end{tabularx} \\ \\ \relax
3:2 & \begin{tabularx}{0.7\textwidth}{X} 監督必須無可指責，只作一個婦人的丈夫，有節制、克己、端正，樂意接待外人，善於教導， \end{tabularx} \\ \\ \relax
3:3 & \begin{tabularx}{0.7\textwidth}{X} 不酗酒，不打人；要溫和，不好鬥，不貪財。 \end{tabularx} \\ \\ \relax
3:4 & \begin{tabularx}{0.7\textwidth}{X} 要好好管理自己的家，使兒女順服，凡事莊重。 \end{tabularx} \\ \\ \relax
3:5 & \begin{tabularx}{0.7\textwidth}{X} 人若不知道管理自己的家，怎能照管神的教會呢？ \end{tabularx} \\ \\ \relax
3:6 & \begin{tabularx}{0.7\textwidth}{X} 剛信主的，不可作監督，恐怕他自高自大，落在魔鬼所受的懲罰裡。 \end{tabularx} \\ \\ \relax
3:7 & \begin{tabularx}{0.7\textwidth}{X} 監督也必須在教外有好名聲，免得被人毀謗，落在魔鬼的羅網裡。 \end{tabularx} \\ \\ \relax
3:8 & \begin{tabularx}{0.7\textwidth}{X} 同樣，執事也必須莊重，不一口兩舌，不好酒，不貪不義之財； \end{tabularx} \\ \\ \relax
3:9 & \begin{tabularx}{0.7\textwidth}{X} 要存清白的良心固守信仰的奧祕。 \end{tabularx} \\ \\ \relax
3:10 & \begin{tabularx}{0.7\textwidth}{X} 這些人也要先受考驗，若沒有可責之處，才讓他們作執事。 \end{tabularx} \\ \\ \relax
3:11 & \begin{tabularx}{0.7\textwidth}{X} 同樣，女執事 也必須莊重，不說閒話，有節制，凡事忠心。 \end{tabularx} \\ \\ \relax
3:12 & \begin{tabularx}{0.7\textwidth}{X} 執事只作一個婦人的丈夫，要好好管兒女和自己的家。 \end{tabularx} \\ \\ \relax
3:13 & \begin{tabularx}{0.7\textwidth}{X} 因為善於作執事的，為自己得到美好的地位，並且無懼地堅信在基督耶穌裡的信仰。 \end{tabularx} \\ \\ \relax
3:14 & \begin{tabularx}{0.7\textwidth}{X} 我希望盡快到你那裡去，所以先把這些事寫給你； \end{tabularx} \\ \\ \relax
3:15 & \begin{tabularx}{0.7\textwidth}{X} 倘若我延誤了，你也可以知道在神的家中該怎樣做。這家就是永生神的教會，真理的柱石和根基。 \end{tabularx} \\ \\ \relax
3:16 & \begin{tabularx}{0.7\textwidth}{X} 敬虔的奧祕是公認為偉大的：神在肉身顯現，被聖靈稱義，被天使看見，被傳於外邦，被世人信服，被接在榮耀裡。 \end{tabularx} \\ \\ \relax
4:1 & \begin{tabularx}{0.7\textwidth}{X} 聖靈明說，在末後的時期必有人離棄信仰，去聽信那誘惑人的邪靈和鬼魔的教訓。 \end{tabularx} \\ \\ \relax
4:2 & \begin{tabularx}{0.7\textwidth}{X} 這是出於撒謊者的假冒；這些人的良心如同被熱鐵烙了一般。 \end{tabularx} \\ \\ \relax
4:3 & \begin{tabularx}{0.7\textwidth}{X} 他們禁止嫁娶，又禁戒食物—就是神所造、讓那信而明白真理的人存感謝的心領受的。 \end{tabularx} \\ \\ \relax
4:4 & \begin{tabularx}{0.7\textwidth}{X} 神所造之物樣樣都是好的，若存感謝的心領受，沒有一樣是不可吃的， \end{tabularx} \\ \\ \relax
4:5 & \begin{tabularx}{0.7\textwidth}{X} 都因神的話和人的祈禱而成為聖潔了。 \end{tabularx} \\ \\ \relax
4:6 & \begin{tabularx}{0.7\textwidth}{X} 你若把這些事提醒弟兄們，就是基督耶穌的好執事，在信仰的話語和你向來所服從的正確教義上得到了栽培。 \end{tabularx} \\ \\ \relax
4:7 & \begin{tabularx}{0.7\textwidth}{X} 要棄絕那世俗的言語和老婦的無稽傳說。要在敬虔上操練自己： \end{tabularx} \\ \\ \relax
4:8 & \begin{tabularx}{0.7\textwidth}{X} 因操練身體有些益處；但敬虔在各方面都有益，它有現今和未來的生命的應許。 \end{tabularx} \\ \\ \relax
4:9 & \begin{tabularx}{0.7\textwidth}{X} 這話可信，值得完全接受。 \end{tabularx} \\ \\ \relax
4:10 & \begin{tabularx}{0.7\textwidth}{X} 我們勞苦，努力正是為此，因為我們的指望在乎永生的神。他是人人的救主，更是信徒的救主。 \end{tabularx} \\ \\ \relax
4:11 & \begin{tabularx}{0.7\textwidth}{X} 你要囑咐和教導這些事。 \end{tabularx} \\ \\ \relax
4:12 & \begin{tabularx}{0.7\textwidth}{X} 不可叫人小看你年輕，總要在言語、行為、愛心、信心、清潔上，都作信徒的榜樣。 \end{tabularx} \\ \\ \relax
4:13 & \begin{tabularx}{0.7\textwidth}{X} 要以宣讀聖經，勸勉，教導為念，直等到我來。 \end{tabularx} \\ \\ \relax
4:14 & \begin{tabularx}{0.7\textwidth}{X} 不要忽略你所得的恩賜，就是從前藉著預言、在眾長老按手的時候賜給你的。 \end{tabularx} \\ \\ \relax
4:15 & \begin{tabularx}{0.7\textwidth}{X} 這些事你要殷勤去做，並要在這些事上專心，讓眾人看出你的長進來。 \end{tabularx} \\ \\ \relax
4:16 & \begin{tabularx}{0.7\textwidth}{X} 要謹慎自己和自己的教導，要在這些事上恆心，因為這樣做，既能救自己，又能救聽你的人。 \end{tabularx} \\ \\ \relax
5:1 & \begin{tabularx}{0.7\textwidth}{X} 不可嚴責老年人，要勸他如同父親。要待年輕人如同弟兄， \end{tabularx} \\ \\ \relax
5:2 & \begin{tabularx}{0.7\textwidth}{X} 年老婦女如同母親。要清清潔潔地待年輕婦女如同姊妹。 \end{tabularx} \\ \\ \relax
5:3 & \begin{tabularx}{0.7\textwidth}{X} 要尊敬真正守寡的婦人。 \end{tabularx} \\ \\ \relax
5:4 & \begin{tabularx}{0.7\textwidth}{X} 寡婦若有兒女，或有孫兒女，要讓兒孫先在自己家中學習行孝，報答親恩，因為這在神面前是可蒙悅納的。 \end{tabularx} \\ \\ \relax
5:5 & \begin{tabularx}{0.7\textwidth}{X} 獨居無靠的真寡婦只仰賴神，晝夜不住地祈求禱告。 \end{tabularx} \\ \\ \relax
5:6 & \begin{tabularx}{0.7\textwidth}{X} 但好宴樂的寡婦活著也算是死了。 \end{tabularx} \\ \\ \relax
5:7 & \begin{tabularx}{0.7\textwidth}{X} 這些事，你要囑咐她們，讓她們無可指責。 \end{tabularx} \\ \\ \relax
5:8 & \begin{tabularx}{0.7\textwidth}{X} 若有人不照顧親屬，尤其是自己家裡的人，就是背棄信仰，還不如不信的人。 \end{tabularx} \\ \\ \relax
5:9 & \begin{tabularx}{0.7\textwidth}{X} 寡婦登記，年齡必須在六十歲以上，只作一個丈夫的妻子， \end{tabularx} \\ \\ \relax
5:10 & \begin{tabularx}{0.7\textwidth}{X} 又有行善的名聲，就如養育兒女，收留外人，洗聖徒的腳，救濟遭難的人，竭力行各樣善事。 \end{tabularx} \\ \\ \relax
5:11 & \begin{tabularx}{0.7\textwidth}{X} 至於年輕的寡婦，你要拒絕登記，因為她們情慾衝動、背棄基督的時候，就想嫁人， \end{tabularx} \\ \\ \relax
5:12 & \begin{tabularx}{0.7\textwidth}{X} 她們因廢棄了當初所許的願而被定罪。 \end{tabularx} \\ \\ \relax
5:13 & \begin{tabularx}{0.7\textwidth}{X} 同時，她們又學了懶惰，習慣於挨家閒逛；不但懶惰，而且說長道短，好管閒事，說些不該說的話。 \end{tabularx} \\ \\ \relax
5:14 & \begin{tabularx}{0.7\textwidth}{X} 所以，我希望年輕的寡婦嫁人，生養兒女，治理家務，不讓敵人有辱罵的把柄， \end{tabularx} \\ \\ \relax
5:15 & \begin{tabularx}{0.7\textwidth}{X} 因為已經有一些人轉去隨從撒但了。 \end{tabularx} \\ \\ \relax
5:16 & \begin{tabularx}{0.7\textwidth}{X} 信主的婦女若有親戚是寡婦，要救濟她們，不可拖累教會，好使教會能救濟真正無助的寡婦。 \end{tabularx} \\ \\ \relax
5:17 & \begin{tabularx}{0.7\textwidth}{X} 善於督導教會的長老，尤其是勤勞講道教導人的，應該得到加倍的敬奉。 \end{tabularx} \\ \\ \relax
5:18 & \begin{tabularx}{0.7\textwidth}{X} 因為經上說：「牛在踹穀的時候，不可籠住牠的嘴」；又說：「工人得工資是應當的。」 \end{tabularx} \\ \\ \relax
5:19 & \begin{tabularx}{0.7\textwidth}{X} 有控告長老的案件，非有兩三個證人就不要受理。 \end{tabularx} \\ \\ \relax
5:20 & \begin{tabularx}{0.7\textwidth}{X} 繼續犯罪的人，要在眾人面前責備他，使其餘的人也有所懼怕。 \end{tabularx} \\ \\ \relax
5:21 & \begin{tabularx}{0.7\textwidth}{X} 我在神、基督耶穌和蒙揀選的天使面前囑咐你要遵守這些話，不可存成見，做事也不可偏心。 \end{tabularx} \\ \\ \relax
5:22 & \begin{tabularx}{0.7\textwidth}{X} 不可急於給人行按手禮；也不可在別人的罪上有份，要保守自己純潔。 \end{tabularx} \\ \\ \relax
5:23 & \begin{tabularx}{0.7\textwidth}{X} 為了你的胃，又常患病，不要只喝水，要稍微喝點酒。 \end{tabularx} \\ \\ \relax
5:24 & \begin{tabularx}{0.7\textwidth}{X} 有些人的罪是明顯的，已先受審判了；有些人的罪是隨後跟著來。 \end{tabularx} \\ \\ \relax
5:25 & \begin{tabularx}{0.7\textwidth}{X} 同樣，善行也有明顯的，就是那不明顯的也不能隱藏。 \end{tabularx} \\ \\ \relax
6:1 & \begin{tabularx}{0.7\textwidth}{X} 凡負軛作奴隸的，要認為自己的主人配受各樣的尊敬，免得神的名和教導被人褻瀆。 \end{tabularx} \\ \\ \relax
6:2 & \begin{tabularx}{0.7\textwidth}{X} 奴隸若有信主的主人，不可因他是主內弟兄就輕看他們，更要越發服侍他們，因為得到服侍的益處的正是信徒，是蒙愛的人。 \end{tabularx} \\ \\ \relax
6:3 & \begin{tabularx}{0.7\textwidth}{X} 若有人傳別的教義，不符合我們主耶穌基督純正的話語與合乎敬虔的教導， \end{tabularx} \\ \\ \relax
6:4 & \begin{tabularx}{0.7\textwidth}{X} 他是自高自大，一無所知，專好爭辯，擅於舌戰，因而生出嫉妒、紛爭、毀謗、惡意猜疑， \end{tabularx} \\ \\ \relax
6:5 & \begin{tabularx}{0.7\textwidth}{X} 和心術不正與喪失真理的人不停地爭吵，以敬虔為得利的門路。 \end{tabularx} \\ \\ \relax
6:6 & \begin{tabularx}{0.7\textwidth}{X} 其實，敬虔加上知足就是大利。 \end{tabularx} \\ \\ \relax
6:7 & \begin{tabularx}{0.7\textwidth}{X} 因為我們沒有帶甚麼到世上來，也不能帶甚麼去； \end{tabularx} \\ \\ \relax
6:8 & \begin{tabularx}{0.7\textwidth}{X} 只要有衣有食，我們就該知足。 \end{tabularx} \\ \\ \relax
6:9 & \begin{tabularx}{0.7\textwidth}{X} 但那些想要發財的人就陷在誘惑、羅網和許多無知有害的慾望中，使人沉淪，以致敗壞和滅亡。 \end{tabularx} \\ \\ \relax
6:10 & \begin{tabularx}{0.7\textwidth}{X} 貪財是萬惡之根。有人因貪戀錢財而背離信仰，用許多愁苦把自己刺透了。 \end{tabularx} \\ \\ \relax
6:11 & \begin{tabularx}{0.7\textwidth}{X} 但你這屬神的人哪，要逃避這些事；要追求公義、敬虔、信心、愛心、忍耐、溫柔。 \end{tabularx} \\ \\ \relax
6:12 & \begin{tabularx}{0.7\textwidth}{X} 你要為信仰打那美好的仗；要持定永生，你為此被召，也已經在許多見證人面前作了那美好的見證。 \end{tabularx} \\ \\ \relax
6:13 & \begin{tabularx}{0.7\textwidth}{X} 我在那賜生命給萬物的神面前，並在向本丟‧彼拉多作過那美好見證的基督耶穌面前囑咐你： \end{tabularx} \\ \\ \relax
6:14 & \begin{tabularx}{0.7\textwidth}{X} 要守這命令，毫不玷污，無可指責，直到我們的主耶穌基督顯現。 \end{tabularx} \\ \\ \relax
6:15 & \begin{tabularx}{0.7\textwidth}{X} 到了適當的時候都要顯明出來：他是那可稱頌、獨一的權能者，萬王之王，萬主之主， \end{tabularx} \\ \\ \relax
6:16 & \begin{tabularx}{0.7\textwidth}{X} 就是那獨一不死、住在人不能靠近的光裡，是人未曾看見，也是不能看見的。願尊貴和永遠的權能都歸給他。阿們！ \end{tabularx} \\ \\ \relax
6:17 & \begin{tabularx}{0.7\textwidth}{X} 至於那些今世富足的人，你要囑咐他們不要自高，也不要倚賴靠不住的錢財；要倚靠那厚賜萬物給我們享受的神。 \end{tabularx} \\ \\ \relax
6:18 & \begin{tabularx}{0.7\textwidth}{X} 又要囑咐他們行善，在好事上富足，甘心施捨，樂意分享， \end{tabularx} \\ \\ \relax
6:19 & \begin{tabularx}{0.7\textwidth}{X} 為自己積存財富，而為將來打美好的根基，好使他們能把握那真正的生命。 \end{tabularx} \\ \\ \relax
6:20 & \begin{tabularx}{0.7\textwidth}{X} 提摩太啊，要持守所給你的託付。要躲避世俗的空談和那假冒知識的矛盾言論。 \end{tabularx} \\ \\ \relax
6:21 & \begin{tabularx}{0.7\textwidth}{X} 有人自稱有這知識而偏離了信仰。願恩惠與你們同在！ \end{tabularx} \\ \\
[1ex]
\hline
\hline
\end{longtable}
$^{1}$又能夠和大家見面.
這次我要用一堂的時間.
來介紹很重要的一卷書.
是保羅寫給他.
作恩信主.
做了他真兒子的年輕的提摩太.
大家知不知道我的英文名字.
叫做Timothy.
即是提摩太.
我記得在中學的時候.
有一個西方來的journeyman.
即是來香港兩年.
就起遣你中學教英文.
他說他幫我們每個學生.
要走出去外面.
就給一個英文名給我們.
我不記得了.
當時他好像給了一個Henry.
我覺得這麼隨便.
我不喜歡.
於是我就沒有用那個名字.
後來我自己就選了.
Timothy這個名字.
因為碰巧我爸爸也是牧師.
後來我又獻身做牧師.
所以我覺得保羅怎樣傳承給提摩太.
等於我爸爸怎樣傳承給我.
我又傳承給年輕一代的傳道人.
讓我能夠去接派.
讓我能夠繼續為神事奉.
所以提摩太這個名字.
真是很適合我.
保羅一共寫了三卷書.
通常被稱為教牧書信.
就是提摩太前書,提摩太後書.
和提多書.
這三卷書可以說是牧師的手冊.
當然不是牧師先看.
信徒更加應該要看.
因為你會知道神對教會.

$^{41}$有一個很高的期望.
對教會有一個很高的要求.
然後神期望教會可以在地上發展.
榮耀祂的名.
今天就讓我來介紹提摩太前書這卷書.
首先我介紹一下提摩太這個年輕的信徒.
他原來是在路斯德這個地方成長的.
他爸爸是一個希臘人.
而他媽媽是一個猶太人.
所以可以說是異族通婚的.
可能他就在他媽媽.
特別是他媽媽的媽媽.
愛祖母那裡.
就承傳了這個信仰.
根據《使徒行傳》十六章的記載.
提摩太的媽媽就叫做有妮基.
而她的婆婆.
即是愛祖母就叫做羅爾.
所以在提摩太後書第一章.
就有說到你母親有妮基.
你愛祖母羅爾那裡.
傳承了這個信仰.
而他的爸爸這個希臘人.
可能是後期才信主的.
我們很感恩.
結果他爸爸也信了耶穌.
保羅在他第一次宣教行程的時候.
就去過路斯德這些地方.
可能就在那時候.
提摩太就聽到保羅傳的道.
就信主了.
所以可以說是保羅帶提摩太信主的.
這個年輕人信了主之後.
就開始很追求.
他心裡面火熱.
所以保羅就說.
當時路斯德那個地方的.
個個都讚這個年輕人.
所以後來保羅第二次.
他要去宣教行程的時候.

$^{81}$提摩太就申請.
加入保羅這個宣教隊.
好像做一個助手.
或者做一個年輕的傳道.
保羅就收了他.
於是保羅和提摩太之間.
就建立了一個很緊密的關係.
保羅不單止一路在宣教行程裡面.
一路傳福音.
一路去訓練這個年輕人.
他看到這個年輕人可以造就.
他將來一定是可以成為一個.
很有能力事奉神的人.
結果後來保羅就派提摩太.
很多很重要的使命.
讓他跟他去一些地方.
然後後來更加將牧羊一間大教會.
就是以弗所教會的重任.
交了給提摩太.
所以提摩太.
他可能是一個生性比較膽怯.
不是很有膽的.
有一點害羞的.
其實我都告訴大家.
以前我都很害羞的.
很怕在別人面前說話.
但是神真的很奇妙.
不知為何每一次站上講台.
我講話的時候我的聲音都不同了.
所以有些弟兄姊妹經常說我.
劉牧師跟你講電話你的聲音不同.
但是你一講到.
你那種洪亮的聲音就出來了.
以往我很害羞.
不喜歡講話.
到現在我都不是那麼喜歡講話.
所以師母經常罵我.
你講多點東西.
講多點東西給我聽.
但是我站出來.

$^{121}$我全講神的道的時候.
不知為何神的靈就好像充滿我.
給我有力量.
讓我將裡面我所領受的道.
好像很流暢那樣傾倒出來.
我想提摩太.
就是這樣一個比較膽怯的人.
所以保羅就鼓勵他.
你不要膽怯.
劉剛強.
因為神賜給我們的.
不是一個膽怯的心.
是一個勇敢.
一個山脈.
一個有頭腦的頭腦.
然後我們可以為主.
來辯論真道.
分辨是非.
特別今天這個時代.
真是一個黑白不分的時代.
你都不知道什麼是真什麼是假的.
到處那些傳媒.
很多消息傳來.
所以各位.
我們必須要靠著神.
給我們的幫助.
我們有判斷.
懂得去判別.
你所聽的東西是真的還是假.
保羅在提摩太前書.
最主要跟提摩太講的是什麼呢?.
就是你管理一間教會.
怎樣設立不同的一些職份.
譬如好像有監督.
有執事.
所以保羅就提到.
對監督有什麼要求?.
做執事有什麼要求?.
然後教會.
教會就是永生神的家.

$^{161}$是真理的柱石.
和根基.
又勸提摩太.
你不要讓人小看你年輕.
你年輕.
他說你總要在言語行為.
信心,愛心,清潔.
這五方面上面.
你要做信徒的榜樣.
保羅實在是一個很有心思的人.
他知道提摩太太年輕.
所以可能教會裡面會有些長者.
有些信了幾十年的信徒.
覺得你這個年輕人做牧師.
我懂事也比你多.
我吃鹽比你吃米多.
所以有些長者就會很看不起.
這個年輕的牧師.
保羅說.
你不要小看你年輕.
你要在很多方面.
你做他們的榜樣.
這個也成為我一個很大的激勵.
當我在教會被安納做牧師的時候.
我剛剛三十歲.
那一年也是剛剛我新婚那一年.
就被安納做牧師.
真的什麼都不懂.
結婚之後連怎樣和太太相處都不懂.
所以撞了很多板.
但是教會有些長者很好.
他們教我很多東西.
他們在各方面來鼓勵我.
特別我講完道之後.
有些長者過來握手.
他說「紹康,你講得很好」.
「你努力講」.
所以神給我一個很大的負擔.
就要在神身邊事奉他.
全港神活潑上傳的道.

$^{201}$所以保羅特別提醒提摩太太.
他說「聖經都是神所默示的」.
「越教訓督責使人歸正」.
「教導人學義都是有益的」.
所以提摩太太前後書.
充滿了很有智慧的華語.
所以今天成為了一本牧師的手冊.
如果你要在教會裡面事奉.
多些讀提摩太太前後書.
你會在提摩太太前後書裡面.
得著很多的啟迪.
得著很多的幫助.
讓我繼續來介紹提摩太太前書.
提摩太太前書一共有六章.
第一章就講到對教義的一些指示.
特別我很欣賞的就是.
保羅在第一章下半章那裡.
他就用自己來做例子.
他說「在罪人裡面我是一個罪魁」.
「因為我迫不及待教會」.
於是有首歌唱過.
「在罪人中我是個罪魁」.
「然而我卻蒙了神的憐憫」.
「在罪人中我是個罪魁」.
「然而我蒙了憐憫」.
我雖然是一個罪魁.
但是我在神面前.
祂是蒙了憐憫.
因為神是要救贖世人.
神連一個迫不及待教會的人.
祂都能夠將他360度轉變過來.
成為很偉大的使徒.
保羅寫的書信.
佔了新約聖經很大的部分.
所以保羅呆呆被神使用.
雖然他以往的歷史很不光彩.
神不是看我們過去的歷史.
神要用你,祂就可以用你.
第二章就講到我們在公共崇拜裡面.
有什麼指示.

$^{241}$對那些婦女我們應該有什麼態度.
然後第三章就講到監督有什麼要求.
執事有什麼要求.
然後第四章就對那些假教師指示他們.
保羅特別在提摩泰的年代裡面.
很多假教師起來.
他們混亂神的道.
所以保羅叫提摩泰要提醒防備假教師.
然後第五章就教會是什麼.
就剛才我提到的.
這家就是永生神的教會.
真理的注釋和根基.
還有你怎樣待人.
所以不要叫人小看你年輕.
你總要在言語行為信心開深.
清潔神明都做信徒的榜樣.
然後到第六章.
就是你要恭敬人.
你要做一個僕人.
僕人是要聽從主人.
如果你打一份工.
你要聽你老闆的話.
老闆叫你做什麼.
你要做什麼.
當然不是什麼都討好他.
而是有原則地去順從主人.
而且保羅在提摩泰前書.
幾次提到要打美好的仗.
這個好像是他一個很特別的語調一樣.
在第一章的十八節.
在第六章十二節.
和最後在提摩泰後書第四章.
都提到打美好的仗.
各位今天我們都是打一場仗.
一場屬靈的爭戰.
我們要和惡者爭戰.
所以我們需要求主幫助我們.
作我們的力量.
我們都要打贏這場仗.
今天撒旦用很多方法.

$^{281}$要來唆擺我們.
要來攻擊我們.
要來打擊我們.
但是我們不要聽撒旦的話.
我們依靠我們的主.
我們穿起全副軍裝.
我們就能夠抵擋惡者一切的攻擊.
後主幫助你.
幫助我.
幫助我們每一個.
我們都上上打美好的仗.
我們站立得穩直到再見主的面.
願神賜福給各位.
\newpage



\section{哥林多前書 9:1-6}
\label{sec:032103}
\textbf{意想不到的自由}~~~~陳嘉信牧師
\newline
\newline
連結: \href{http://www.church.com.hk/acms/content.asp?site=cdc\&op=show\&type=product\&code=032103\&layout=sermon}{\texttt{032103}} ~~~~ 語音日期: 2026-05-03
\newline
\newline
\hyperref[sec:032102]{\small{< < < PREV SERMON < < <}}
~
\hyperref[sec:index_chronic]{\small{[返順時目]}}
~
\hyperref[sec:032104]{\small{> > > NEXT SERMON > > >}}
\newline
\newline
哥林多前書 9:1-6
\newline
\begin{longtable}{cl}
\hline
\hline
章節 & 經文 (和合本修訂版)\\
\hline
9:1 & \begin{tabularx}{0.7\textwidth}{X} 我不是自由的嗎？我不是使徒嗎？我不是見過我們的主耶穌嗎？你們不是我在主裡面工作的成果嗎？ \end{tabularx} \\ \\ \relax
9:2 & \begin{tabularx}{0.7\textwidth}{X} 假若對別人來說，我不是使徒，對你們來說，我總是使徒；因為你們在主裡正是我作使徒的印證。 \end{tabularx} \\ \\ \relax
9:3 & \begin{tabularx}{0.7\textwidth}{X} 對那些質問我的人，這就是我的答辯。 \end{tabularx} \\ \\ \relax
9:4 & \begin{tabularx}{0.7\textwidth}{X} 難道我們沒有權利靠著傳福音吃喝嗎？ \end{tabularx} \\ \\ \relax
9:5 & \begin{tabularx}{0.7\textwidth}{X} 難道我們沒有權利帶著信主的妻子一起出入，如同其餘的使徒，和主的兄弟們，和磯法一樣嗎？ \end{tabularx} \\ \\ \relax
9:6 & \begin{tabularx}{0.7\textwidth}{X} 只有我和巴拿巴沒有權利不做工嗎？ \end{tabularx} \\ \\
[1ex]
\hline
\hline
\end{longtable}
$^{1}$今年的連題是「走出去傳遞愛 活出來見證神」.
盼望透過哥林多前書第九章今天的訊息.
可以讓我們跟思想怎樣走出去傳遞愛.
我們一起禱告.
多謝你給我們有你的說話.
求因主祝福保守我們.
讓我們透過你的說話.
彼此的激勵.
在你的愛裡.
能夠活出這份真正的自由.
我們這樣禱告 奉主耶穌基督的聖名自祈求.
阿們.
在香港基督教對微信的朋友來說.
除了給我們有兩天的假期.
讓我們可以放鬆一下之外.
有人說基督教給人的印象一點都不輕鬆.
就好像在頭影片看到孫悟空的緊箍咒.
是沒有自由的.
很多規矩的.
又要回教會.
又要服侍.
又要奉獻.
綁手綁腳.
我不知道大家認識一些微信的朋友.
知道你是基督徒的時候.
他們會抱著什麼樣的眼光看著你.
會不會是同情的眼光.
哎呀 真是淒涼 你信了耶穌.
這麼看不開.
但是不過這樣.
很特別的是基督教又給人一種釋放的感覺.
脫離枷鎖.
有趣的是福音解讀.
福音解禱.
這些力量是很浩大的.
福音帶給人釋放的能量.
福音解讀在現時來說.
可以說是最有效的解讀方法.
比起沒有宗教.
甚至其他的宗教.

$^{41}$更加具威力.
上面我看到有其他宗教的學者.
居然是其他宗教的學者.
在問一個問題.
研究為什麼基督教這麼大的威力.
我就想你回來教會就知道了.
原來福音的解讀.
可以讓人認識福音.
而福音帶來對罪人的接納.
讓人脫離毒癮 捆綁.
有真正的自由.
今天的經文臨前第九章.
本來是告訴我們.
如何活出耶穌的真自由.
在臨前第九章第一節.
本來一句反問.
「我不是自由的嗎?」.
講出他是自由的.
第一至十四節.
提到神派他去傳福音.
他有自由結婚.
有自由收入.
有人供養他的生活所需.
讓他可以專心去傳福音.
在第四和第五節.
保羅提出兩個反問.
難道我無權靠傳福音有收入嗎?.
難道我無權結婚嗎?.
保羅帶出他有權有自由.
其實這一點可以推遠一點.
每一個人不只是神交托職務的人.
不只是仕途.
不只是傳道人 牧者.
其實每一個人都可以享受神的創造.
正正是因為保羅知道這一切都是好東西.
只要懷著感恩的心就可以去享受.
另一個地方有人不讓人結婚.
另一個書卷要人禁戒飲食.
保羅在提摩太前書裡說.
凡神所做的物都是好的.

$^{81}$能夠享受神的祝福.
去結婚有收入 有吃喝 有生活.
令我們有滿足.
其實不用覺得不好意思.
試想像一下.
你去吃自助餐 吃放題.
你有沒有見到有人一邊吃一邊唉聲嘆氣.
唉呀 醉過了 唉呀 折墮了 唉呀 浪費了.
我想說的是.
不需要這樣想.
不需要不好意思.
但要留意一件事.
這些祝福不可以令我們看為絕對.
不會取代神的位置.
這些祝福 這些自助餐 這些美食.
不可以取代神.
否則神的祝福會變成取代他的偶像.
所以第一個結論是什麼呢.
第一個結論就是自由是什麼呢.
是自由地去感恩 享受神的賜予.
在一至十四節裡面讓我們看得見.
保羅知道神賜予他有權利.
他可以自由地去享受.
他更加有自由去不用這些權利.
保羅居然叫人不要給他錢.
在十二節和十五節.
十二節說 然而我們並沒有用過這權利.
十五節 但這權利我傳言沒有用過.
指的就是他暗示自己沒有結婚.
重點其實都不是結不結婚.
而是金錢.
在第九章裡面他強調.
他不靠哥林多人的奉獻去支持他.
他自己工作賺錢.
自費支持自己去傳福音.
這個正正就是記載在《士堂傳》第十八章.
如果大家回家有時間可以看看這張聖經.
在《士堂傳》第十八章裡面說.
保羅在哥林多.
就是哥林多這個地方.

$^{121}$他做了什麼職業.
大家知不知道.
沒錯 織帳棚.
他織帳棚做工去賺錢.
在安息日他才在會堂那裡傳道 傳福音.
到第十八章第五節.
當西拉和提摩泰從馬其頓到哥林多去會合保羅的時候.
他們其實是帶著馬其頓那些東西奉獻給他.
經文就沒有再提起保羅織帳棚.
保羅就為道迫切全程投入去傳福音.
所以收人工去傳道不是必然的.
是處境性的.
有些宣教士去到某些國家.
不能夠以宣教士的身份或者警察會的支持.
給他去傳福音.
可能他就要怎樣做呢.
他就要帶職或者做營商去經濟支持自己去宣教.
香港普遍的教會都在經濟上能夠支持同工去做聖經.
去教聖經 去傳福音.
去做就門徒的工作.
假使有一天環境不容許的時候.
我其實都要自己去想.
我都要帶職去工作.
去支持自己去傳道.
有人有自由這樣去運用權利.
更加有人自由這樣不用權利.
保羅正正就是行使他的仕途.
行使他的自由.
他不用他的權利.
不收那麼多人的錢去傳福音.
目的是什麼呢.
在第九章第十八節那裡.
保羅就說.
就算我傳福音的時候.
洗人不花錢得福音.
原來保羅餵著那麼多人去節制自己.
他不用他的權利.
其實他都是在用他的自由.
他有權利都不用.
有句話叫做什麼呢.

$^{161}$有風 使盡力.
保羅沒有這樣做.
他沒有給錢去操控他自己.
他沒有讓這些東西去操控.
他反而為了不妨礙那麼多人傳福音.
他能夠節制.
其實這樣說.
很多沒有自由節制的人.
其實都很羨慕保羅這種知足的人生態度.
有節制能力才有真正的自由.
所以使人不花錢得福音.
這個就是保羅的原則.
但這句話聽起來好像很豪情壯志.
保羅很瀟灑 很大方.
但是他在第九章第十六節前兩節.
他都有些事情去擔心.
他說什麼呢.
傳福音是不得已的.
若不傳福音我就有禍了.
傳福音他沒有自由不傳.
是神給他的一個責任.
是沒得選擇的.
保羅提過人應該看他為基督的僕人.
那麼僕人有沒有自由呢.
僕人是沒有自由的.
試想像一下.
可能家裡有工人姐姐.
會不會有一天工人姐姐跟你說.
這些事我不做的.
昨晚我在港島之後.
我媽媽就跟我說.
現在的工人可能不是這樣想的.
當然那時候保羅的時代.
他用回僕人這個概念.
那時候的時代是沒得選擇的.
沒得選擇你做什麼還是不做什麼.
傳福音他是沒得選擇的.
不過態度怎樣去傳是有得選擇.
他不是無奈地去接這個工作.
不是無奈地去履行這個責任.

$^{201}$他是積極去面對.
在經文裡面他沒有提到自己傳福音有多辛酸.
遇到多少的困難.
有多可憐.
第九章裡面沒有提到.
相反他說如果沒得傳福音他就慘了.
他很緊張他的職責.
保羅還提到他履行職責的時候會得到賞賜.
他的賞賜是什麼呢.
賞賜是原文賞賜是有人工.
有出糧的含義.
那他的人工是什麼呢.
在第十八節裡說.
不如我們一起讀.
這讓我的賞賜是什麼呢.
就是我傳福音的時候使人不花錢得福音.
你會看到原來保羅對他人工是什麼.
就是可以傳一個免費的福音.
你想想這樣的人工也有.
簡單來說他收的人工就是不收錢.
那就是收還是不收.
不收就是收了.
收了就是沒得收.
那就是收還是不收.
不收就是收了.
你會覺得很奇怪.
你會問這是什麼想法呢.
他把別人的益處當成自己賺了.
你有沒有認識這樣的人.
相信這就是基督徒的生命.
他努力賺取的就是自由地.
很自由地去傳免費的福音.
即是free of charge的福音.
令到人能夠得自由.
可以讓人free.
所以是free.
三個free.
他的工作態度真是令人意想不到.
對於我們今天用什麼態度去打一份工很有啟發.
大家有沒有見過這樣的同事.

$^{241}$他以自由的心態去做一份看起來不自由的工作.
保羅傳不傳福音他沒有自由.
但是他有自由去用什麼態度去傳.
這裡保羅教我們怎樣運用自由.
我們有自由去顧及別人的情況或者感受.
不做某些事.
其實這就是更大的自由.
有人說回教會要奉獻好妻良.
但是有人會問.
為什麼教會不收錢.
為什麼我剛剛拿茶點的時候沒有人收我錢.
為什麼我喝東西的時候沒有人收我錢.
我們教會恩福晚生堂會開一個自修室.
在疫情的時候看到很多家庭困獸鬥搞不定.
我們就開始了自修室.
讓基層家庭可以帶小朋友來做功課.
可以來溫書.
有些叔叔姨姨陪著他們.
有些家長會走來問我們.
為什麼教會不收錢.
我們曾經試過和社署一起合作.
在教會裡搞嘉年華會.
邀請到社署接觸的家庭來到.
家長都會問為什麼不收錢.
其實很簡單我們這樣回答.
因為我們根本不想賺錢.
相反我們很想和晚生有困難的家庭去同行.
當然我們都很想他們能夠信耶穌.
我們會找機會和他們分享.
但是信耶穌不是一個條件.
而是我們很想將最好的東西介紹給他們.
不只是做功課這麼簡單.
而是將主耶穌這份這麼寶貴的禮物和他們分享.
教會有時候有人會問.
你們怎麼圍著被子.
有沒有問過這個問題.
教會怎麼圍著被子.
怎麼能夠維持到.
因為神感動馬其頓.
我們的弟兄姊妹回應主耶穌的愛.

$^{281}$神感動我們奉獻去支持自修室的服侍.
支持童宮前線的牧養.
推動福音的工作和其他教會的一切開支.
我看到倭福堂代討事項裡.
很多宣教士為他們禱告.
為關愛的地方禱告.
為一些生命去祈禱.
原來做了這麼多的服侍.
教會也有支持一些不同的服侍.
福音的工作.
我看到報告有很多福音的活動.
可以用場地打乒乓球.
又可以有電影的欣賞.
又請到一些導演來分享福音見證.
其實這些是很美的事情.
但是這些怎麼可以圍著被子呢.
這些就是我們的弟兄姊妹.
在主的感動下.
行使他們運用金錢的自由.
來奉獻給神.
所以第二個結論是什麼呢.
就是自由是自由裡面的節制.
自由有時候我們會覺得是擺脫一些東西.
脫離一些枷鎖.
脫離金錢的控制.
就好像保羅一樣.
保羅他不收錢.
但是他不單止不收錢.
他為傳福音更進一步採取主動.
去做一件事.
對於自由來說.
聖經有兩個概念.
第一個概念就是freedom.
有自由去擺脫.
幫助我們脫離罪.
解除壞習慣.
有自由去不做某些事情.
不做這個不做那個.
原來聖經還有一個概念.
是積極的.

$^{321}$是free to.
有自由去做某些事情.
而在第十九節.
第九章第十九節.
這兩個概念都有出現.
我們一起讀一讀這一節.
我雖然是自由的.
不受人管轄.
但我甘心作了眾人的僕人.
為贏得更多的人.
保羅有自由.
他free from all men.
他有自由去擺脫所有人的管轄.
不受人管轄.
他free from all men.
但是下一句.
有另一個概念.
還有就是.
他很特別.
相反.
他free to serve all men.
他有自由去服侍所有人.
他甘心作眾人的僕人.
保羅他不單止是基督耶穌的僕人.
這裡還說到他是所有人的僕人.
這個就真是為人民服務.
你們要問誰為人民服務.
原來不是誰.
原來是保羅.
保羅說.
我雖然有自由.
但是我寧願做僕人.
做僕人本身一方面是放棄自由.
但另一方面是他用自己的自由去選擇做僕人.
那是什麼意思.
那就是自投羅網.
他甘心受束縛.
去做所有人的僕人.
他的目的是什麼.
在經文裡面說.

$^{361}$他為要贏得更多的人.
可能你會問.
保羅怎樣去做眾人的僕人呢.
在下兩節.
下幾節有說.
一起讀一讀.
20至22節.
預備起.
對猶太人.
我就作猶太人.
為要贏得猶太人.
對律法以下的人.
作律法以下的人.
為要贏得律法以下的人.
對沒有律法的人.
我就作沒有律法的人.
為要贏得沒有律法的人.
對軟弱的人.
我就作軟弱的人.
為要贏得軟弱的人.
在這裡.
保羅在20至21節.
他提到律法.
律法給人的感覺是怎樣的.
是沒有自由.
可以令到我沒有自由.
想怎樣就怎樣.
而律法的制裁是什麼呢.
就是令到你沒有自由.
要坐牢.
但是第20至21節.
保羅對律法有些看法.
在第20節說.
對於律法以下的人.
就是說對於有神律法的猶太人.
他們受規限於舊約摩西律法.
保羅說他不在律法以下.
為什麼呢.
耶穌來到已經承傳了舊約律法的精神.
耶穌救贖了.

$^{401}$釋放了律法底下的人.
身為基督徒的保羅已經有自由了.
他不再受到舊約規條所束縛.
他不再需要為贖罪去獻祭.
雖然如此.
他說他還是作律法以下的人.
次唐人傳記載到保羅陪猶太人一起行結正禮.
他雖然有自由.
但是他甘願為他的猶太同胞.
去維持守一些律法規條.
目的是什麼.
為要贏得律法以下的人.
而下一節對於沒有律法的人.
就是說所有非猶太人.
就是外邦人.
我就作沒有律法的人.
又是什麼意思呢.
他怎樣作沒有律法的人呢.
他突破了舊約的限制.
猶太律法是怎麼樣的呢.
猶太律法將聖結和污穢分得很清楚.
猶太人不會和外邦人吃飯.
大家有沒有一些猶太教徒的朋友.
你們有沒有和他們吃飯.
如果你和他們吃飯是很麻煩的.
你不可以帶他們去吃叉燒.
你去喝茶的時候很多點心都不能吃.
為什麼呢.
因為豬肉是不潔的.
所以和外邦人吃飯很多東西都不能吃.
因為是不潔.
對於猶太人來說.
但是保羅這個猶太人.
他會和外邦人吃飯.
他為要贏得沒有律法的人.
但是保羅補充一句.
其實我在神面前不是沒有律法.
保羅說不要看我freestyle.
不要看我像無法無天那樣.
其實我有律法的.

$^{441}$他說他在基督的律法之下.
他想說他自己不在摩西律法底下.
而在基督底下.
在基督裡面按照加拉太書.
所有律法全律法都包含什麼.
在愛人如己裡面.
他受這一份愛去驅使.
去愛外邦人.
向外邦人傳福音.
他就成就了律法的精神.
所以有人說基督教很沒有自由.
欠缺彈性.
我們看到保羅這個很有自由度的去服事.
是相當具彈性的.
有自由去服事.
負彈性地融入社群不封閉.
這個其實是基督教的特色.
基督教不是僵化的文化.
一點都不是緊箍咒.
事實上基督教能夠適應多元的文化.
基督教的擴展其實不再是在發源地.
最興旺的教會是不是在耶路撒冷.
當然不是.
很多的地方.
但其實很多傳統宗教來來去去都是在發源地.
印度教最多人信在哪裡.
就是印度.
佛教 孔教都是在發源地裡面多.
但基督信仰現在最多人信在哪裡.
人們說在南半球.
南半球有什麼地方.
有南美 亦有非洲.
當然亞洲都有.
但我們多說一點非洲.
以非洲為例.
1900年非洲的基督徒數目是9\%的人口.
都不少.
但現在大約有多少基督徒在非洲.
有44\%的非洲人是基督徒.
當然今天中國 南韓都有很多基督徒.

$^{481}$我們再深入了解一下非洲這片土地.
為什麼非洲這麼多基督徒.
有非洲學者研究.
非洲人相信超自然有好的靈.
亦有壞的靈.
當他們讀到聖經.
很多人在耶穌身上看到自己歷史上的盼望.
非洲人信了耶穌之後.
能否保存自己的習俗.
保存自己的傳統.
他們會為神去唱歌.
非洲人很喜歡唱歌.
他們會為神演奏樂器.
他們最擅長的是什麼樂器.
非洲鼓.
要付錢學的很貴的.
他們最擅長的.
用他們的非洲鼓能夠去讚美神.
所以基督教是幫助非洲人成為新造的非洲人.
而不是被複製出來的歐洲人.
大家明白我說什麼嗎.
我們經常以為基督教信仰在歐洲開展.
所以歐洲人去到全部地方.
要人去信耶穌之餘.
亦要保留一些歐洲的文化.
其實不是的.
在這裡來說.
是要讓他們成為新造的非洲人.
非洲學者強調.
西方的世俗主義.
在他們反超自然的主義和個人主義的主張底下.
對於非洲文化來說.
基督教更加適合他們.
更加適合非洲人.
在聖經裡.
當非洲人讀到耶穌勝過邪靈.
透過基督教.
非洲人可以批判他們的傳統.
讓非洲人可以脫離.
脫離巫術.

$^{521}$脫離拜精靈.
但仍然可以讓非洲人持守某些傳統.
所以就是這樣的緣故.
基督教會在世界各地遍地開花.
在不同的地方孕育成長.
能夠適應融入不同的文化.
卻沒有對自身的宗旨妥協.
為什麼他們要融入不同的社群.
保羅很簡單在第22節裡說出一個原因.
對什麼樣的人.
我就作什麼樣的人.
給的原因是什麼.
無論如何.
我總要救一些人.
我們要了解一個真相.
就是每一個人都需要救贖.
就是這個世界都會終結.
我們今天所享受的一切.
就是蒲妃,放緹都會過去.
我們每一個人都需要獲得拯救.
所以當天介紹你來認識耶穌的朋友.
他們都很希望你分享到這份自由.
所以今天大家坐在這裡的時間.
我們都盼望能夠走出去.
讓別人都能夠.
我們的親友,崗位裡所認識的人.
都能夠享受到這一份自由.
所以保羅接著說.
凡我所做的都是為福音的緣故.
為要與人共享這福音的好處.
可能看到這裡的時間.
你或者會問.
對著不同的人.
成為不同的人.
有什麼奇怪的想法呢?.
這樣不就很假了嗎?.
對什麼樣的人.
我就作什麼樣的人.
原來香港人有一個版本.
叫做什麼?.

$^{561}$大家口型我都看到了.
見人講人話.
這樣不就很假了嗎?.
沒錯.
保羅可以很假.
除非一件事.
除非什麼呢?.
除非是出於愛.
來華的宣教士戴德生.
他曾經說過一句話.
假如我有千磅英金.
不留下一磅不給中國.
這份愛其實不是選擇性的愛.
是愛與我們不同的人.
你說愛非洲人?.
我們都不敢走過去.
不同的語言.
不同的文化.
不同的飲食習慣.
不同的身體氣味.
但是一位英國人.
來到中國人當中.
他更加做些什麼呢?.
他留邊.
他穿唐裝.
去融入當地的社群.
目的就是贏得中國人.
我們有弟兄.
有弟兄見到教會有樂器的需要.
他之前曾經去學打鼓.
為了想幫手.
我們有些導師.
為了青少年.
為了跟兒童一起.
他們很努力地去學習補課.
學習.
他們常常看見矇著雙眼在掃.
學一些社交平台.
什麼Fresh.
IG.

$^{601}$學這些東西.
跟青少年溝通.
跟孩子溝通.
他們更加重要的是什麼呢?.
他們冒著生命的危險.
他們做了什麼呢?.
他們居然去.
平時不會做的.
去超級市場買一些零食.
看見成分.
反式脂肪.
人造色素.
他冒著生命危險.
去跟青少年吃.
捨命陪青少.
很大的心理關口.
又或者有些.
終於等到機會可以吃.
他們要為了什麼呢?.
不是為了自己.
是為了融入他們.
所以第三個結論是什麼呢?.
原來自由是有自由去服侍.
去負擔性地融入社群.
愛是超越界限.
也是超越自我中心.
我們見到一群弟兄姊妹.
很愛我們的年輕人.
很愛我們的小朋友.
本來有動力去融入不同的文化.
但是他自己都要堅定.
能夠克制自己.
他要克制己身.
我們一起讀一讀第24至27節.
預備起.
你們不知道在運動場上賽跑的.
大家都跑.
但得獎賞的只有一人.
你們也要這樣跑.
好使你們得著獎賞.

$^{641}$凡參加競賽的.
在各方面都要有節制.
他們不過是要得會擾壞的冠冕.
我們卻是要得不會擾壞的冠冕.
所以我奔跑.
不像無目標的.
我鬥拳.
不像打空氣的.
我克制己身.
完全順服.
免得我傳福音給別人.
自己反而被淘汰了.
本來用運動員的概念.
去比喻基督徒的人生.
他不只要求自己.
不只要求自己做運動員.
他還叫大家一起去報渣馬.
去做運動員.
鼓勵所有基督徒.
第24節問好之後那句是什麼.
你們也要這樣跑.
好使你們得著獎賞.
而25節那裡本來要我們.
很清楚自己的人生目標.
是為不擾的冠冕.
是為耶穌而跑.
是為神的國而跑.
為什麼要這麼強調呢.
因為無私的付出.
有時候做著做著都會迷失.
可能大家在事奉裡都有一段日子.
起初一片單心.
一片初心去服侍神.
但做著做著.
有時候不知道自己為什麼要這樣做.
我為了什麼呢.
他忘記了自己的初心.
所以我們要認清目標.
所以戴德生所說.
假如我有千磅英金.

$^{681}$不留下一磅不給中國.
原來後面還有一句說話.
大家知道是什麼嗎.
「不是中國,乃是基督」.
戴德生學後保羅.
教清自己的目標.
他知道自己在服侍的是基督.
27節保羅強調.
他要克制己身.
使他完全順服免得我全福音級別人.
自己反而被淘汰.
他要努力放下自我中心.
否則動機不好的時候.
可能別人一句批評.
別人一句說話.
又或者對方一些反應.
可能我們就會迷失.
全福音給其他人.
自己反而會迷失被淘汰.
其實在第9章24至27節.
其實是第9章和第10章中間的轉接說話.
所以保羅除了提醒我們.
要為其他人負擔性運用自己的自由之外.
也要為自己小心.
不要跌倒.
提醒我們不要參與.
原來第8至10章是說.
是否可以吃祭國偶像的食物.
保羅說為了別人的緣故可以吃.
你為了他也可以陪他一起吃.
或者你知道凡物都是潔淨的.
所以可以吃.
但是你要徹底知道.
如果別人不是這樣想的時候.
你又要不要吃.
你要節制.
但為人的時候你要節制不要吃.
但你為自己也要小心.
因為你參與的是什麼.
你要知道.

$^{721}$你不要參與在偶像崇拜的裡面.
不要參與在偶像的祭祀底下.
因為第10章裡面說那些是祭鬼.
是得罪神的.
所以我們要克制自己.
讓整個人都信服在主基督裡面.
免得我們傳福音給其他人.
自己反而被淘汰.
你可能接著又會問.
剛才我們不是一直在說自由嗎.
要克制自己.
這樣就很沒有自由.
我們有時候將自由這個定義.
定得太簡化.
定得太窄.
以為自由就是負面地排除所有限制.
沒有界線.
可以想怎樣就怎樣.
這個就是自由.
但我想說的是界線和限制.
其實是通往自由的一個路徑.
我想問大家有沒有人喜歡彈琴.
有沒有人喜歡玩樂器.
也有.
我問其他崇拜弟兄姊妹的時候.
也有人回應說喜歡.
我想問練樂器辛苦嗎.
練樂器辛苦.
練琴辛苦.
可能大家在家裡有兒女在練琴.
催他們辛苦嗎.
他辛苦之餘你也辛苦.
我想問練琴的時候有沒有自由.
練琴的時候同一時間做其他事.
沒有自由的.
你不會練琴一邊練一邊回覆手機.
你做不到這件事.
這樣就很沒有自由.
但我想說的是.
當你練琴練到很熟的時候.

$^{761}$你不需要再看樂譜去彈琴.
我想問一下感覺你猜爽不爽.
很爽的.
在電視上看到朗朗扭腰揮灑自如.
看到他很陶醉很high很自由.
原來這就是釋放出來的自由.
獲得更大的自由.
愛能夠推動自己.
我們剛才說最有效的解讀方法是福音解讀.
原來福音解讀裡面很有紀律.
要很有紀律.
我曾經訪問過.
和福音解讀機構裡的一位弟兄聊天.
他說其實福音解讀是很有紀律.
很辛苦很嚴格.
比坐牢還要嚴格.
但我又問為何這些學員會聽你們的要求.
因為他們知道導師愛他們.
因為他們知道耶穌愛他們.
因為體會到神的愛的時候.
讓他們有力量去克制自己.
我們經常聽到青少年說.
我要打機的自由.
有沒有聽過?家裡有沒有人說過?.
可能你會有這樣說過.
我要打機的自由.
但有時候我會去問.
其實打機這個操作本身是很沒有自由的.
遊戲有這麼多遊戲規則.
招式有時候可以出有時候不可以出.
然後又要走到哪裡去儲哪一樣東西.
很多規限很沒有自由.
其實關鍵不是在於自由.
關鍵是在於你喜不喜歡.
試想像一下一個男生在街上等女朋友.
買著束花在等她.
他還要提早.
甚至乎天寒地凍.
凍冰冰在等.
女生本身早就到了.

$^{801}$不過她不出現.
她埋伏在附近.
去看看男朋友在做什麼.
看他的表情.
看他的熱切的時候.
很開心.
看到他這麼凍冰冰的時間.
我想告訴大家.
其實這些是違反了男生的人物設定的.
這不是他的本性.
試想像他約其他朋友會不會這樣做.
拿著束花好好的.
會這樣做.
我想問問大家.
大兄姐妹.
這些叫什麼?.
我怕大家會說這些就是傻子.
我今早聽到說這些就是蠢子.
有沒有人這樣說你傻子過?.
我想說的是.
大兄姐妹.
這個是愛.
明白嗎?.
這個是愛.
愛是釋放.
是會令人釋放很大的力量.
是可以令一個人違反他人物設定.
去做一些你意想不到的行為.
但有人說.
愛是喜歡做什麼就做什麼.
我想告訴大家.
其實.
愛不是最大的.
大家有沒有聽過這句話.
經常聽的.
愛不是你想怎樣就怎樣.
保羅在再多幾章的時間.
他提到愛.
哥林多前書哪裡最出名是愛?.
十三章.

$^{841}$沒錯.
保羅說愛是恆久忍耐.
愛是不作害羞的事.
害羞的事是什麼?.
不作羞恥的事.
是怎樣?.
是甘願自限約束自己.
但不是每一樣限制都是釋放自由.
紀律和限制只是合適於我們本質上的性情.
才能夠釋放出來.
好像你在湖邊看到一條魚.
在湖裡游泳.
你真是辛苦.
真是可憐.
我釋放你出來.
我將你安全地放在草地上.
你猜這條魚有什麼情況?.
會反肚.
會死.
所以要適合的本質.
自由不是除去所有的限制.
而是找到適合的限制.
所以找到合適的職業.
找到合適的興趣.
運動對象是好的.
但更何況是信仰呢?.
是關係到整個人生的自由.
所以要找到適合的信仰.
不是什麼信仰都要找.
所以大家來到的時候.
我看到大家找到適合的信仰.
在臨前第十一章第一節裡.
保羅說你們該效法我.
像我效法基督一樣.
保羅叫我們學他.
學他學基督.
那究竟基督做了什麼.
讓我們可以學他呢?.
在約翰一書裡第四章第九節裡說.
神猜他獨生子到世間來.

$^{881}$使我們直著他得生.
神愛我們的心就在此顯明了.
你看到神的愛.
就是他來到世間上.
這就是聖誕的真諦.
聖誕就是紀念神的愛子到世間來.
目的是令我們重獲新生.
就顯明了神愛我們的心.
愛是最大的自由.
但亦是最有限制.
一般的宗教信仰是單向的.
即是一位神要求人.
而那個人要自限.
要適應這位神.
但基督信仰是雙向的.
耶穌道成肉身.
神成為人.
我們強調一位救主.
他甘願與當時被視為罪人的人一起吃飯.
甚至愛當時迫害他的人.
耶穌甘願自限.
受制於痛苦與死亡.
為我們犧牲他自己.
這就是看到他的一份愛.
而去服侍我們.
我們怎樣能夠回應他.
我們就是領略到這份愛.
原來基督的愛激勵我們.
當我們能夠體會領略到這份愛的時候.
我們就實踐使命.
所以看起來好像放棄自由.
實踐來說就是得到真正的自由.
怎樣才有如魚得水的自由呢?.
希望你能夠領略到和正在經歷.
就是降伏在耶穌裡面.
被他的愛激勵你去愛家人,去愛鄰舍.
愛是會喪失自由.
也都最具釋放力的.
放棄個人的自主性.
去苦就服侍對象.

$^{921}$卻是釋放出最大的力量和滿足感.
但是我要先旨聲明.
愛是有風險的.
大家可能在經驗裡面都得到了.
愛可以是有心碎的風險.
保羅為什麼要寫哥林多前書第九章呢?.
因為第三節裡面說.
有人質問他.
有人誤會他.
有人攻擊他.
所以他要答辯.
保羅甘心命底去愛哥林多人.
不收錢向他們傳福音.
建立了這些哥林多的信徒.
建立了哥林多教會.
但是他竟然受著哥林多人的攻擊.
這就是愛的風險.
但是保羅沒有放棄.
保羅依然耐性地向哥林多人解釋.
依然甘心去愛攻擊他的人.
這就是仕途的生命是愛的動力.
相反如果你很安全地將這份愛.
不要有風險了.
將它鎖在保險箱裡面的話.
這份人生是了無生氣.
是會很孤單的.
我們陰風晚山堂.
有位很奇怪的堂委會主席.
他提早退休.
他和一班舊同事一起去吃飯聊天.
去問一些退休的舊同事一起聊天.
關心一下他們.
發現到一個事情.
原來每個人都弄得身體差了.
通常退休大家想怎樣.
環遊世界.
他們就會說去旅行.
暴飲暴食.
這樣高那樣高.
就會問那位弟兄.

$^{961}$那位堂委會主席.
為什麼你這麼fit.
他就會回答.
我也是暴.
但我不是暴飲暴食.
我是天天一晚暴.
他每天走路回來教會.
來到自修室去開檔.
去設檯去設牆去服侍.
他當做運動.
他身體還好了.
更加重要的是什麼呢.
他覺得他遲了信耶穌.
他要好好珍惜他的生命.
去認識主和服侍主.
兼且我們很多的導師.
付出了愛.
你以為做導師這麼容易嗎.
很勞氣的.
吐血的.
你看到小孩子和年青人的成長.
他們都很開心.
所以最後一個結論是什麼呢.
自由是.
愛是在自由裡面.
甘願受限.
自律的生活.
這份愛是釋放很大的能量.
所以今天的結論是什麼呢.
這份愛是能夠釋放很大的力量.
再重複一下.
自由是什麼呢.
是我們去感恩享受神的賜予.
在自由裡面要有節制.
基督信仰令我們自由地去服侍.
負擔性融入社群.
不封閉.
愛是在自由裡面.
甘願受限.
自律的生活.

$^{1001}$這份愛釋放出很大的力量.
讓我們一起禱告.
給我們領略主耶穌的愛.
亦體會到天賦你.
犧牲你的愛.
只為我們成就.
讓我們一同在這份愛裡面被融化.
我們昔日都是因為這樣.
而來到你的根前.
我祈求主你帶領著我們的人生.
讓我們能夠被你的愛激勵.
讓身邊的人.
我們的鄰舍.
我們的家人.
我們在崗位裡面所服侍.
所接觸到的人.
都能夠有機會體會到.
這份莫大的自由.
垂聽禱告.
奉主明求.
阿們.
\newpage



\section{創世記 25:24-28}
\label{sec:032104}
\textbf{利百加:愛與掌控的母親}~~~~譚綺敏傳道
\newline
\newline
連結: \href{http://www.church.com.hk/acms/content.asp?site=cdc\&op=show\&type=product\&code=032104\&layout=sermon}{\texttt{032104}} ~~~~ 語音日期: 2026-05-10
\newline
\newline
\hyperref[sec:032103]{\small{< < < PREV SERMON < < <}}
~
\hyperref[sec:index_chronic]{\small{[返順時目]}}
~
\hyperref[sec:032105]{\small{> > > NEXT SERMON > > >}}
\newline
\newline
創世記 25:24-28
\newline
\begin{longtable}{cl}
\hline
\hline
章節 & 經文 (和合本修訂版)\\
\hline
25:24 & \begin{tabularx}{0.7\textwidth}{X} 到了生產的日期，看哪，腹中是對雙胞胎。 \end{tabularx} \\ \\ \relax
25:25 & \begin{tabularx}{0.7\textwidth}{X} 先出生的身體帶紅，渾身有毛，好像皮衣；他們就給他起名叫以掃。 \end{tabularx} \\ \\ \relax
25:26 & \begin{tabularx}{0.7\textwidth}{X} 隨後，以掃的弟弟也出生，他的手抓住以掃的腳跟，因此給他起名叫雅各。兩個兒子出生時，以撒六十歲。 \end{tabularx} \\ \\ \relax
25:27 & \begin{tabularx}{0.7\textwidth}{X} 兩個孩子漸漸長大，以掃善於打獵，常在田野；雅各為人安靜，常住在帳棚裡。 \end{tabularx} \\ \\ \relax
25:28 & \begin{tabularx}{0.7\textwidth}{X} 以撒愛以掃，因為常吃他的野味；利百加卻愛雅各。 \end{tabularx} \\ \\
[1ex]
\hline
\hline
\end{longtable}
$^{1}$靚姊妹平安 母親節快樂 祝姊妹們.
今日講一個記載舊約的母親.
經文剛才都講了 25章24至28節 請聽我讀.
生產的日子到了 福中果然是雙子.
先產的身體發紅 渾身有毛 如同皮衣.
他們就給他起名叫耳掃 就是有毛的意思.
隨後又生了耳掃的兄弟 手抓住耳掃的腳筋.
因此給他起名叫雅各 就是抓住的意思.
尼伯加生了兩個兒子的時候 以撒年正六十歲.
兩個孩子漸漸長大 耳掃善於打獵 常在田野.
雅各為人安靜 常住在帳幕裡.
以撒愛耳掃 因為常吃他的野味.
尼伯加卻愛雅各.
我們先一起來祈禱.
設立家庭 設立父母 這個又尊貴又重要的角色.
求主幫助我們每年的母親節.
而且我們做母親的 我們要學習.
怎樣用神的話語去教導我們的孩子.
求主幫助我們 多謝主.
祈禱奉主耶穌基督 得勝命祈求.
阿們.
今天是母親節 所以問問大家的問題.
你覺得你媽媽對你的人生影響大不大呢?.
如果你覺得媽媽對你的人生影響有一半以上.
舉手來看看.
不是很多 我以為每個人都會舉手.
其實媽媽對孩子的影響不止是孩子.
一直對孩子的影響都是非常深遠的.
可以是小時候 也可以是一生之久.
現代心理學家亦都研究.
他說父母的教養方式 特別是母親.
很大程度影響了孩子將來怎樣看自己.
怎樣跟人相處.
有沒有安全感.
有沒有自信.
甚至怎樣面對神.
所以我們看到.
父母的影響其實很大.
現在的研究.
大約將父母的教養方式分為四種.

$^{41}$這四種都是建立在兩個核心上.
一個是你給孩子的愛.
另一個是界線 就是原則.
甚麼是愛呢?.
包括你接納他 你支持他.
你怎樣陪伴他.
你怎樣關心他.
這些都是愛的表現.
當然你供養他.
飲食各樣東西都是愛的表現.
第二是界線.
就是怎樣有原則 有規矩去管教他.
或者對他有要求.
這些都是在這兩個界線上.
其實你發覺很有趣.
現代的概念在這兩個界線上.
其實神創造亞當夏娃的時候.
創世記第二章.
其實神就是做了這件事.
神做了亞當夏娃之後.
他供應他一切的所需.
然後就設立了一個界線.
就是分辨善惡術的果子.
你們不可以吃.
這個就是一個界線.
所以你發覺很特別地.
原來神也是這樣教導我們.
真正健康的愛不是完全放任的.
不是沒有規矩的愛.
一定是包含了有界限的.
所以我們現在來看看這四類母親.
第一個就是權威型.
一聽名字就以為權威不是很好.
但其實這個才是最健康的母親.
她又有愛又有高的界線.
她會去抱孩子 聽孩子說.
同時也會跟孩子說.
這件事不可以做.
也會跟孩子解釋原因.
她不只是發命令.

$^{81}$她更加懂得怎樣去教孩子.
所以研究發現這種媽媽.
帶出來的孩子會經常跟他們溝通.
所以比較有安全感.
有自律 有責任感.
人際關係也會比較健康.
這是最接近神教導我們的方向.
第二個就是專制型的母親.
她有很多命令 但是很嚴格.
但她比較少關懷.
可能媽媽有句口頭禪.
「我說了是就是」.
不能讓孩子去反駁.
這種媽媽很重視規矩.
但很少一些情感的交流.
可能跟我們上一輩有點相似.
孩子表面上很乖.
但其實裡面有很多壓抑.
很容易去焦慮.
自己想怎樣做.
就會做的自主性.
長大之後.
孩子很容易反叛.
不喜歡聽媽媽說話.
或者很討好人.
因為從小覺得她未被真正接納.
所以做得很討好人.
第三種就是縱容型的母親.
很疼愛她們.
非常愛孩子.
她們不會跟孩子說不.
因為不想令孩子失望.
亦不想令孩子哭.
慢慢地孩子說怎樣就怎樣.
這種孩子長大之後.
就會缺乏自律 比較自我中心.
很難去接受一些挫折.
因為從小都沒有界限去規範孩子.
所以有挫折就會有很大問題.
第四種就是忽視型的母親.

$^{121}$意思就是甚麼都不理會.
放任她們.
這種母親對孩子最大的傷害.
她們不會打孩子.
未必是打孩子的傷害.
因為她們不理會 冷漠.
但可能她們太忙 太累.
或者自己有很多壓力.
所以令孩子覺得自己不重要.
所以這種孩子很缺乏安全感.
或者很自卑 很渴望被愛.
我曾經聽過一個姐妹分享.
不是姐妹 是一個企業主管.
有人問她人生最大的影響是甚麼.
她說不是學校 不是工作.
不是環境 是她媽媽.
她媽媽每天送她上學的時候.
都會對她說.
你是神所愛的孩子.
你可以做一些正確的事情.
你是做得到的.
她說小時候聽這些說話.
她沒有甚麼感覺.
但長大後遇到困難.
或者試探的時候.
或者很多壓力的時候.
她說這些說話就會浮現在她心裡.
她知道她是神所愛的孩子.
有甚麼壓力 神都會幫她擔當.
但也知道她可以選擇做正確的事.
當試探來到的時候.
她自己可以選擇 也做得到.
所以母親的聲音.
其實是影響孩子的一生.
所以我們做媽媽的不是很容易.
特別是現代的媽媽.
又要工作 因為房價高 租金高.
其實現在的媽媽多數都要工作.
很少不工作.
又要照顧家庭 又要煮飯.

$^{161}$又要陪讀 陪讀書.
真的很辛苦.
但小孩子的時間是一個黃金時期.
又會過得很快.
要影響他們 其實小孩子很重要.
所以又要陪他們.
所以做媽媽真的不容易.
求主給我們智慧.
知道怎樣去做一個神所喜悅的媽媽.
今天要說一個很特別的母親.
是李百家.
李百家其實很聰明 很能幹.
也很有愛心.
在《創世記》24章的時候.
他在井旁幫助一個陌生人.
打水給他的駱駝和陌生人喝.
其實不是一個普通人可以做到.
當他做了母親的時候.
他在教養孩子方面出了一些問題.
這些問題不單影響了兩個孩子.
甚至影響了整個家族 整個以色列的民族.
所以我們剛才讀的經文.
在《創世記》25章.
說到李百家生了雙胞胎.
以素和雅各.
當孩子漸漸長大的時候.
兩個孩子的個性都不太不同.
以素喜歡打獵.
經常在田野裡.
雅各比較安靜.
在帳篷裡.
孩子不同其實沒什麼問題.
但聖經記載了一個很重要的一句話.
以撒愛以素.
李百家卻愛雅各.
問題就在這裡開始.
我們看到父母對孩子有偏心.
其實李百家和以撒都不是不愛孩子.
他們都很愛孩子.
他們不是冷漠的父母.

$^{201}$他們都有愛 有關心.
還有情感.
但他們的問題就出在哪裡呢.
就是出在他們偏愛 偏心.
爸爸很愛以素.
因為經常吃她的食物.
打獵回來 已經長大了.
媽媽就愛雅各.
可能因為她比較安靜.
經常可以陪在她身邊.
所以各有不同.
但不是因為孩子不同.
我們就會有偏愛.
但其實我們做父母都有軟弱.
有時候不知不覺裡.
我們比較喜歡哪個小孩.
對於乖 經常聽你說.
或者很像你的小孩.
自然愛他多一點.
但如果比較難搞.
經常走來走去.
好像過度活躍.
可能少一點耐性去對待他.
所以有時候都會偏心對待他們.
偏心造成什麼呢.
偏心令孩子容易嫉妒 比較.
或者破壞關係.
偏心會留在孩子心裡.
覺得很不公平.
或者被拒絕.
有些苦毒在心裡.
我聽過一位姊妹分享她的成長.
她說一直覺得媽媽比較偏心哥哥.
可能還男輕女.
有一次她在班裡考到第二名.
她很開心 馬上回家跟媽媽說.
我考到第二名了.
以為媽媽都讚她厲害.
她很期待.
但媽媽不是這樣說.

$^{241}$媽媽說 你看哥哥考第一名.
你也沒有 努力一點吧.
這一刻她覺得.
媽媽經常疼愛哥哥.
當然 哥哥讀書又厲害.
我總是不夠他好的.
所以她很不開心.
苦毒留在她心裡.
雖然她仍然很努力.
很感恩她仍然很努力.
但她一直有苦毒.
也沒有安全感.
覺得自己經常不夠好.
跟媽媽和哥哥的關係也不太好.
直到她信了耶穌的時候.
信了主.
她知道神愛她.
也接納她.
肯定她做所有的一切.
她慢慢信主的時候.
慢慢體會其實神很疼她.
也發覺媽媽不是不疼她.
所以她開始願意.
跟家人建立好一點的關係.
所以我看到.
父母偏心其實對孩子有影響.
或者一句說話.
留在孩子心裡.
令他們有一些苦毒.
在《創世記》25章和27章.
我們可以看到.
耶穌和雅各的媽媽關係其實不太好.
以撒和耶穌.
他們無法從小培養孩子.
互相照顧.
互相彼此合作.
可能他們沒有做到.
培養孩子這方面的事情.
反而因為他們偏心.
而令孩子中間有一些裂痕.

$^{281}$在27章35到36節.
當以撒要為兒子祝福的時候.
祝福錯了.
我們跳到後面去看.
以撒說了一句話.
「你兄弟已經用詭計.
將你的福分奪去了」.
耶穌就說.
「他是叫雅各」.
「不是抓住機會騙了我兩次」.
所以因為這樣的時候.
耶穌就因為他的父親為雅各祝福了之後.
他就怨恨雅各.
心裡就會說.
「當我父親居喪的日子近了」.
「死了的時候」.
「我就要殺我的兄弟雅各」.
弄得這麼厲害.
你會發覺聖經很特別.
記載了父子的對話.
我覺得很有意思.
為何以撒很疼愛耶穌.
但祝福錯了的時候.
他就說.
「你兄弟已經用詭計.
將你的福分奪去了」.
我覺得這句很韌義.
不知道你覺不覺得.
其實以撒自己做父親.
祝福錯了也有負責任.
他自己是不清楚了.
糊里糊塗了.
祝福錯了.
但因為他很疼愛以撒.
覺得自己祝福錯了很震驚.
所以他就說了這句話.
我覺得這句話.
其實也很傷心.
以撒一聽到就很激怒.
覺得「我的弟弟有沒有搞錯」.

$^{321}$「用詭計奪了我的福分」.
「是什麼弟弟來的」.
可以這樣想.
本身他們兩個已經不和.
你還說了這句話.
所以他覺得自己要殺了弟弟.
所以有時候一些說話.
挑起了兄弟之間的矛盾.
其實是很不應該的.
所以父母的偏心.
其實在孩子裡.
留下了很深的傷痕.
也造成了家庭的糾紛.
有些小朋友.
其實你不說一些話.
他也會有競爭性.
彼此之間也會有競爭.
但當父母再加上偏心的時候.
肯定情況更加不好.
更加不能和諧.
這個傷害可能一直傳給下一代.
你會發覺雅各從小就在偏心的家庭長大.
當他做爸爸的時候.
他也偏愛弱勢.
在《創世記三十七章》之後的段落.
結果兄弟之間也有彼此的嫉妒.
因為雅各愛弱勢過於其他兒子.
弱勢在他年老的時候生的.
他還為弱勢做了一件彩衣.
當他哥哥見到的時候.
他也很恨弱勢.
覺得為什麼爸爸要殺小孩呢.
所以就和他說一些不友善的話.
最後把他賣掉.
還想殺小孩.
後來沒有殺到.
就賣了去埃及.
所以聖經教導我們.
《以弗所書》六章四節.
你們作父親的.

$^{361}$不要激怒兒女.
卻要照著主的教訓和勸誡.
養育他們.
這裡不只是和父親.
我覺得是和父親和母親說的.
不要激起兄弟姐妹之間的競爭.
反而是要培養他們彼此相愛.
在神的話語中去養育他們.
所以這是第一個.
我們看到父母偏心造成了一些問題.
第二件事我們看到.
就是母親利伯嘉的掌控.
到了二十七章.
即是推翻剛才二十七章那裡.
當以撒想為耳掃祝福的時候.
利伯嘉偷偷地聽見了.
而以撒沒有和妻子溝通.
偷偷地想為耳掃祝福.
利伯嘉聽見的時候.
他沒有真正和丈夫溝通這件事.
他選擇了操縱整件事.
什麼意思呢.
利伯嘉生產之前.
他記得神跟他說過.
在《創世記》的二十五章二十三節.
他說過兩國在你之內.
兩族要從你的身上出來.
這族必強於那族.
將來大的要服侍小的.
利伯嘉其實有聽神的話.
他聽見神說.
小的服侍大的.
神的心意是小的服侍大的.
祝福是大的.
不是照傳統的長子.
接受家族的祝福.
是分兩倍.
也是一路承繼下去.
他知道神的心意.
但問題是他不是想用神的方法.

$^{401}$他沒有等候神.
也沒有去禱告.
他反而是用自己的方法.
用操控的方法.
去成就一些事情.
很多時候操控是因為.
你不相信神會去做這件事.
你想自己去操縱畫面的局勢.
所以利伯嘉利用了丈夫年老的時候.
眼睛看不清楚.
他安排了雅各去冒充耳掃.
這是一個很不尊重丈夫的表現.
夫妻之間如果失去了尊重的表現.
整個家庭都會受影響.
如果要去停止一些偏心或操縱.
其實家庭的婚姻本身.
是一個很關鍵的因素.
所以第一步.
婚姻本身我們要盡力去維護.
保持和諧和尊重.
這可能是送給孩子第一個最好的禮物.
如果夫妻之間失去了鞏固的關係.
其實小小的裂痕.
孩子都會看到.
也會受到影響.
所以小心去維護婚姻的合一.
是很重要的.
尤其是在一些孩子的事上.
要不要去讀哪間學校.
要不要去補習.
或者將來孩子的選科.
或者你給予意見上.
小孩子都會看見.
如果大人不合一.
小孩子很聰明.
他們會懂得看你們的反應.
如果有偏心有偏愛.
小孩子更加知道站在哪裡.
他會贏的.
所以父母應該先去多溝通.

$^{441}$在一切孩子的事上.
要團結要溝通.
只要互相尊重.
彼此信任.
亦都知道怎樣去放手.
才會與孩子建立最好的關係.
所以願我們都去謹記.
夫妻關係是很重要的.
第三我們看見.
李百嘉母親的錯誤教導.
李百嘉因為偏愛了雅各.
他直接策劃了一個騙局.
運用母親的特權.
要兒子聽從她的吩咐.
幫助雅各得到一個祝福.
他教雅各怎樣做.
要雅各去順從他.
去田間裡面找些野味.
造成一些美味給他的父親.
他就幫雅各穿上一些衣服.
包著他的身體.
穿上雅各的衣服.
扮得像耳掃一樣.
因為耳掃有很多毛.
他教雅各說謊.
說我是他的長子耳掃.
當他去到父親的前面的時候.
雅各說如果父親識破了.
以為他欺騙我.
我就會有咒詛.
得不到祝福.
怎麼辦.
他很害怕.
李百嘉更加說.
我兒你不用怕.
你的咒詛歸在我身上.
你只要聽我的話就行.
其實那時候雅各已經是成人.
他應該知道這樣是不對的.
所以才跟母親說.

$^{481}$但當母親這樣說的時候.
雅各沒有違背母親的意願.
而母親李百嘉.
亦在教一件很可怕的事情給兒子.
為了得到神的祝福.
可以不擇手段.
為了得到你要的結果.
欺騙是可以做的.
當兒子害怕的時候.
他說兒子你不用害怕.
不是你要去負責任.
而是母親來擔當.
有什麼事都有媽媽來負責.
其實你教孩子的是什麼呢.
心理學家告訴我們.
孩子最早的是非觀念.
是從父母來的.
當母親親自去示範欺騙.
亦告訴小孩.
沒關係的.
你聽著我說.
你不用負責任.
小孩學到什麼呢.
手段比誠信更重要.
不擇手段是可以的.
自己不用承擔選擇回來的後果.
這其實是一個非常大的錯誤.
亦比我們今天要小心.
如果過度重視小孩.
或者保護小孩.
做錯事都不用自己負責任.
這是一個很嚴重的後果.
所以當我們過度去掌控孩子.
或者教導他們的時候.
得出來的結果是怎樣呢.
一是好像雅各那樣.
表面很馴服媽媽.
很馴服父母.
但內裡就學了一些隱藏.
不懂做決定.

$^{521}$或者避開去承擔責任.
懂得逃走.
或者就像二嫂那樣.
二嫂很反叛.
突然憤怒就爆發.
徹底地切斷了關係.
小孩都不要了.
將來要殺小孩.
無論是哪一種的結果.
哪一種的孩子.
最後父母都會失去了孩子.
和失去了真正和孩子的關係.
所以利伯家的偏心和掌控去欺騙.
最後負上了沉重的代價.
我們看《創世記》27章41到46節.
二嫂因為怨恨雅各.
心裡想殺雅各.
而利伯家知道了.
他再次給雅各一個主意.
就是你走到舅父拉班那裡.
和他住一些日子.
住一些日子在原文的意思是住幾天.
等哥哥的怒氣消失了.
我就叫你回來.
然後他又和丈夫說.
騙他的丈夫.
不是騙的.
是用好的說話和丈夫說.
我看見我的長子二嫂.
娶了這個黑人女子為妻.
一天到晚都不安寧.
我們應該打發雅各.
要去舅父那裡住幾天.
娶一個希伯來人的婦人.
這樣才可以.
所以他就瞞住了爸爸.
就打發了雅各去.
這一去不是幾天.
而是二十年.
聖經再沒有記載.

$^{561}$利伯家和小兒子雅各的重逢.
利伯家可能在兒子雅各回來之前.
他都去世了.
他用盡心機將祝福.
留給最愛的兒子雅各.
但結果就是永遠的分離和失去了他.
或者更加諷刺的說.
這位美麗能幹.
影響了全家庭的母親.
最後因為她的偏心.
因為她的掌控.
因為她的欺騙.
落得了家庭破碎.
丈夫二嫂蒙在鼓裡.
最後長子二嫂怨恨雅各.
雅各也需要逃亡.
自己孤獨終老.
這就是她愛小兒子雅各的人生.
而她愛的小兒子雅各的人生又是怎樣呢.
充滿了欺騙和被欺騙的循環.
她欺騙了爸爸.
後來逃到舅父拉班.
被騙了二十年.
吃盡苦頭.
所以我們發覺.
對於小孩子的管教.
對父母的影響是很深遠的.
每一天.
小孩子都是看著你怎樣去做人.
怎樣去處理衝突.
怎樣去面對不公平.
怎樣不如意發生的事情.
在你身上的時候.
你是否真的要祈禱去交託神呢.
還是你要去掌控一切呢.
你的口裡說.
我們什麼都交託給神了.
但你走出來就是.
倚靠媽媽的聰明.
所以在黎伯家的身上.

$^{601}$我看到這位母親.
帶給孩子的問題.
令我們做媽媽很值得反思.
就是我們不要去偏愛.
而是要公平對待我們的孩子.
小孩子很容易說.
It's not fair.
你發覺小孩子.
我經常都聽到.
是不公平的.
有時候我都很喜歡問小孩子.
你覺得媽媽是否公平對待你.
有沒有偏心.
很怕我不知不覺.
處理一些事情有偏心.
所以盼望我們知道.
每個小孩子都是獨特的.
都有不同的.
我們都要知道小孩子是被愛的.
第二就是我們不要去操縱.
要去交託給神.
我們很容易靠自己的聰明.
很容易想家人都聽我們的說話.
但我們要學習去交託.
將孩子帶到神面前.
用神的話語去管教他們.
也要設立一些規則.
一些界限.
為他們多點去祈禱.
第三就是不要去虛假.
只要去真實跟隨神.
今天我們教導小孩子.
其實對於他們是很重要的.
今天我們說我們去崇拜很重要.
我們跟隨神很重要.
但當有生日會 踢球 補習.
這些一切來到的時候.
不用回禮拜堂.
不用去兒童主義學校.
我們先去這些.

$^{641}$當我們做這些的時候.
其實我們可以告訴孩子.
跟隨神不是很重要.
去少一天 教會沒有所謂.
最重要是你讀書好.
最重要是你健康 做運動.
最重要是一切都理想.
其實你是在灌輸一些價值觀給孩子.
跟隨神是否真的那麼重要呢.
盼望我們都去反思.
其實做母親很不容易.
做父母也不容易.
做父親也不容易.
所以當孩子小的時候.
我們又擔心他們長不大.
但當他們長大的時候.
我們會專心去學業 朋友 工作 結婚 信仰.
這些我們都會擔心.
有時候比擔心自己更加緊張.
所以今天提醒我們.
我們愛孩子.
我們要怎樣去愛呢.
要靠著神的恩典.
不是要靠自己的聰明.
聖經裡面都沒有完美的母親.
剛才那個都不是完美的母親.
也是以色列的始祖.
更加不是完美的母親.
我經常都有軟弱.
但我們要靠著神的恩典.
相信神的恩典是夠用的.
或者過去我做得不好.
神總是有憐憫的.
你為孩子去做過的事.
或者流過的眼淚.
神是看見的.
小孩子也知道的.
將來小孩子也不會忘記的.
所以盼望我們懂得怎樣陪孩子.
去走完我們的人生路.

$^{681}$為他們祈禱 陪伴他們.
最後用一句說話結束.
我們要愛孩子.
但不要失去界線.
要設立界線.
但不要失去關係.
要信靠神.
不是操控孩子.
盼望我們一起去祈禱.
設立母親這個尊貴的角色.
我們不多不少都受著.
我們自己母親的影響.
在過去可能母親的一些說話.
也都對我們成為傷害.
今天如果我們當中有受傷害的弟兄姊妹.
求主您特別去醫治他們.
願他們相信耶穌之後.
可以在愛裡得到自由.
求主特別用恩手托著我們做母親.
我們都有軟弱的時候.
讓我們學懂怎樣去等候主.
學懂怎樣去教導我們的孩子.
用神的話語去教導我們的孩子.
也都求主賜智慧.
讓我們看待每個孩子都是公平公正的.
以致我們的孩子可以得著愛.
他們都需要被關心被愛.
也都幫助我們知道怎樣設立一個界限.
當他們需要去學習的時候.
我們都不怕有個界限要令他們明白.
我們不是去操控孩子.
而是去教導孩子在神的手裡.
用神的話語去教導他們.
求主親自去引導我們.
多謝主聽我們祈禱.
奉主耶穌基督得勝一命祈求.
\newpage



\section{約翰福音 7:11-36}
\label{sec:032105}
\textbf{榮耀君王:終極問答}~~~~林鍵牧師
\newline
\newline
連結: \href{http://www.church.com.hk/acms/content.asp?site=cdc\&op=show\&type=product\&code=032105\&layout=sermon}{\texttt{032105}} ~~~~ 語音日期: 2026-05-09
\newline
\newline
\hyperref[sec:032104]{\small{< < < PREV SERMON < < <}}
~
\hyperref[sec:index_chronic]{\small{[返順時目]}}
~
\hyperref[sec:032107]{\small{> > > NEXT SERMON > > >}}
\newline
\newline
約翰福音 7:11-36
\newline
\begin{longtable}{cl}
\hline
\hline
章節 & 經文 (和合本修訂版)\\
\hline
7:11 & \begin{tabularx}{0.7\textwidth}{X} 節期間，猶太人尋找耶穌，說：「他在哪裡？」 \end{tabularx} \\ \\ \relax
7:12 & \begin{tabularx}{0.7\textwidth}{X} 人群中有許多人對他議論紛紛，另有的說：「他是好人。」有的說：「不，他是迷惑群眾的。」 \end{tabularx} \\ \\ \relax
7:13 & \begin{tabularx}{0.7\textwidth}{X} 可是沒有人公開談論他，因為他們怕猶太人。 \end{tabularx} \\ \\ \relax
7:14 & \begin{tabularx}{0.7\textwidth}{X} 節期已過了一半，耶穌上聖殿去教導人。 \end{tabularx} \\ \\ \relax
7:15 & \begin{tabularx}{0.7\textwidth}{X} 猶太人驚訝地說：「這個人沒有學過，怎麼那樣熟悉經典呢？」 \end{tabularx} \\ \\ \relax
7:16 & \begin{tabularx}{0.7\textwidth}{X} 於是耶穌回答他們，說：「我的教導不是我自己的，而是差我來那位的。 \end{tabularx} \\ \\ \relax
7:17 & \begin{tabularx}{0.7\textwidth}{X} 人若立志要遵行神的旨意，就會知道這教導究竟是出於神，還是我憑著自己說的。 \end{tabularx} \\ \\ \relax
7:18 & \begin{tabularx}{0.7\textwidth}{X} 憑著自己說的人是尋求自己的榮耀；但那尋求差他來那位的榮耀的人，他是真誠的，在他心裡沒有不義。 \end{tabularx} \\ \\ \relax
7:19 & \begin{tabularx}{0.7\textwidth}{X} 摩西不是傳了律法給你們嗎？你們卻沒有一個人守律法。為甚麼想要殺我呢？」 \end{tabularx} \\ \\ \relax
7:20 & \begin{tabularx}{0.7\textwidth}{X} 眾人回答：「你是被鬼附了！誰想要殺你呢？」 \end{tabularx} \\ \\ \relax
7:21 & \begin{tabularx}{0.7\textwidth}{X} 耶穌回答，對他們說：「我做了一件事，你們都驚訝。 \end{tabularx} \\ \\ \relax
7:22 & \begin{tabularx}{0.7\textwidth}{X} 摩西傳割禮給你們（其實割禮不是從摩西開始，而是從列祖開始的），你們就在安息日給人行割禮。 \end{tabularx} \\ \\ \relax
7:23 & \begin{tabularx}{0.7\textwidth}{X} 人若在安息日受割禮，是為了不違背摩西的律法，我在安息日使一個人痊癒了，你們就向我發怒嗎？ \end{tabularx} \\ \\ \relax
7:24 & \begin{tabularx}{0.7\textwidth}{X} 不要憑外表斷定是非，總要按公平斷定是非。」 \end{tabularx} \\ \\ \relax
7:25 & \begin{tabularx}{0.7\textwidth}{X} 於是耶路撒冷人中有的說：「這個人不是他們想要殺的嗎？ \end{tabularx} \\ \\ \relax
7:26 & \begin{tabularx}{0.7\textwidth}{X} 你看，他還公開講道，他們也不對他說甚麼。難道官長真的認為這是基督嗎？ \end{tabularx} \\ \\ \relax
7:27 & \begin{tabularx}{0.7\textwidth}{X} 然而，我們知道這個人從哪裡來；可是基督來的時候，沒有人知道他從哪裡來。」 \end{tabularx} \\ \\ \relax
7:28 & \begin{tabularx}{0.7\textwidth}{X} 那時，耶穌在聖殿裡教導人，喊著說：「你們認識我，也知道我從哪裡來；我並不是憑著自己來的。但差我來的那位是真實的，你們不認識他。 \end{tabularx} \\ \\ \relax
7:29 & \begin{tabularx}{0.7\textwidth}{X} 我卻認識他，因為我從他那裡來，是他差遣了我。」 \end{tabularx} \\ \\ \relax
7:30 & \begin{tabularx}{0.7\textwidth}{X} 於是他們想要捉拿耶穌，只是沒有人下手，因為他的時候還沒有到。 \end{tabularx} \\ \\ \relax
7:31 & \begin{tabularx}{0.7\textwidth}{X} 但人群中有好些人信他，他們說：「基督來的時候，他所行的神蹟難道會比這人行的更多嗎？」 \end{tabularx} \\ \\ \relax
7:32 & \begin{tabularx}{0.7\textwidth}{X} 法利賽人聽見群眾對耶穌這樣議論紛紛，祭司長和法利賽人就打發聖殿警衛去捉拿他。 \end{tabularx} \\ \\ \relax
7:33 & \begin{tabularx}{0.7\textwidth}{X} 於是耶穌說：「我跟你們在一起的時候不會太久了，我要回到那差我來的那裡去。 \end{tabularx} \\ \\ \relax
7:34 & \begin{tabularx}{0.7\textwidth}{X} 你們要找我，卻找不到；我所在的地方，你們不能去。」 \end{tabularx} \\ \\ \relax
7:35 & \begin{tabularx}{0.7\textwidth}{X} 於是猶太人彼此問：「這人要往哪裡去，使我們找不到他呢？難道他要往散居在希臘的猶太人那裡去教導希臘人嗎？ \end{tabularx} \\ \\ \relax
7:36 & \begin{tabularx}{0.7\textwidth}{X} 他說『你們要找我，卻找不到；我所在的地方，你們不能去』，這話是甚麼意思呢？」 \end{tabularx} \\ \\
[1ex]
\hline
\hline
\end{longtable}
$^{1}$歡迎大家來到我們的周六崇拜.
很多年前當我還在神學院學習的時候.
有一天我互教學的老師.
和我們分享有關人生的終極問題.
他說幾千年以來西方的哲學家.
都是不斷探討人生裡面的四個終極問題.
這四個終極問題分別是我從哪裡來.
我要到哪裡去.
就是你死了之後.
為什麼這個世界會有苦難.
又是另一個文化.
為什麼這個世界有罪惡.
人生的意義在哪裡.
罪的救贖在哪裡.
當時我第一個反應.
聽完這四個終極問題的時候.
就是原來這個才是人生的終極問題.
我一直還以為人生的終極問題就是衣食住行.
我原來真的太膚淺.
太過無知了.
信主當時已經有幾年.
還在思想地上的事.
連不認識上帝的哲學家都會探求天上的事.
我們作為基督徒.
是不是更加應該要去思想人生的終極問題呢.
正如C.S. Lewis說.
All that is not eternal is eternally out of date.
所有不屬於永恆的事情.
從永恆的角度來說.
其實都已經全部過期了.
有趣的是.
當我們打開這本聖經.
這本非常厚厚的聖經的時候.
1189章的新舊約聖經.
這四個終極問題.
其實在頭三章裡面已經開始浮現.
而且在頭三章裡面.
你已經開始對這四個問題有所解答.
所以其餘的1186章的聖經.
都是終極問題和答案的延伸和解說.

$^{41}$我想和你開始看這四個問題.
第一我從《雅履來》.
《創世記》第一章第27節.
「上帝照著他的形象創造人,照著上帝的形象創造他們.
他創造了他們,有男有女」.
我要往哪裡去.
我死了之後要往哪裡去.
這個問題比較複雜.
我們需要先回答第三條.
然後再回來看第二條.
第三條.
「為什麼這個世界有苦難,為什麼這個世界有罪惡」.
答案是因為撒旦的引誘.
因為亞當夏娃的叛逆.
《創世記》第二章第15節.
「若和上帝把人安置在伊甸園」.
伊甸的意思是喜樂.
將我們放在伊甸喜樂生命園裡面.
「若和上帝對那人說園中各樣樹上所出的.
你都可以隨意吃」.
隨便吧不要客氣.
「只是知善惡的樹所出的你不能吃.
因為你吃的日子你必定死」.
園中有多少樹我們真的不知道.
我們知道很多很多多不勝數.
「然而只有一棵」.
你不要以為只有十棵裡面這一棵不能吃.
是千千萬萬幾百萬幾千萬幾億棵樹裡面.
只有一棵是不可以吃的.
「然而就是那棵」.
《創世記》第三章第四節.
「吃撒旦的化身堆女人說你們不一定死」.
就是吃了那棵東西都不一定死.
因為上帝知道你們吃的日子眼睛就開了.
你們就像上帝一樣知道善惡.
「於是女人見那棵樹好作食物又悅人的眼目.
那樹令人喜愛能使人有智慧」.
在古代的時候誰需要智慧.
做皇帝需要智慧.
這個女人想做皇帝.

$^{81}$她就摘下那樹吃了又吸了.
她的丈夫一起吃她也吃了.
在那天人類的始祖亞當夏娃犯了滔天的大罪.
你可能會說.
他們沒有殺人放火.
也沒有虐待動物.
也沒有火燒伊甸園.
只是吃了一個他們不應該吃的果子.
這樣就是滔天大罪嗎?.
這是什麼樣的法律道理?.
所以我告訴你.
那是因為罪的定義和本質就是不聽上帝的話.
和我們現在的法律系統是完全不一樣的.
現在的法律系統你自己喜歡做什麼都可以.
只要你不要干犯別人.
但是他們干犯的是上帝.
上帝的話.
當他們不聽上帝的話.
這個就是滔天的大罪.
他們不敬畏上帝.
他們不敬拜上帝.
他們以自己為重心.
他們叛逆上帝.
他們去做自己眼中看為對的事情.
這個也成為了亞當叛逆以後.
所有墮落人類的通病.
你從小到大都知道的.
This is mine.
那個是我的.
兩個小孩在搶玩遊戲的時候.
搶玩具的時候.
This is mine.
是我的.
我們都是在做我們自己眼中看為對的事情.
於是我們回到第二條問題.
我要到哪裡去.
因為我們現在犯了罪.
因為我們墮落.
那麼我們以後會去哪裡呢?.
《創世記》第三章第23節.

$^{121}$若說上帝就驅逐他們出伊甸園.
再沒有喜樂再沒有生命.
若說上帝把那些人趕出去.
就在伊甸園東邊.
安設了幾路柏和.
發出火焰轉動的劍把生命樹.
守住生命樹.
就是說你不會有機會得享永生.
世人都犯了罪.
規缺了上帝的榮耀.
罪的公家乃是死死後.
且有審判.
如果你看完這本書三章之後.
你可以再看最後的三章.
你要看書的時候.
你會說這本書太過厚了.
明天要考試那怎麼辦呢?.
看頭三章.
然後看尾三章.
那你就差不多知道這本書寫了些什麼了.
啟示錄第20章第15節.
凡明知沒有記載生命冊上的人.
全部要被應盡火壺裡.
但我們去到第四條問題.
人生的意義何在?.
或者罪的救贖在哪裡?.
答案是罪的救贖在耶穌基督裡面.
創世記第三章第15節.
第三章而已.
已經開始給我們這個答案了.
上帝說我要使你和女人彼此為仇.
你的後裔和女人的後裔也彼此為仇.
她要傷你的頭.
你要傷她的腳筋.
這個是一個有關耶穌基督的預言.
指向耶穌基督在十字架上要承傳的救贖.
祂要被傷.
但祂會復活.
而耶穌的十字架會勝過撒旦的權勢.
令我們罪得赦免.

$^{161}$有機會重新進入永生.
我想帶你回到《約翰福音》第七章今天的經文.
很奇妙的是.
今天我們剛才讀的20多節.
短短20多節的聖經裡面.
耶穌其實一口氣回答了千年以來.
西方哲學家所問的四個終極問題.
耶穌都答完了.
耶穌的終極答案也是我們的終極出路.
你聽到耶穌的答案.
你明白耶穌的答案.
你跟著耶穌的答案.
你就有可能出路.
耶穌回答我從哪裡來.
耶穌說我從上帝哪裡來.
耶穌回答我要去哪裡去.
耶穌說我要回到上帝哪裡去.
耶穌也回答了為什麼這個世界有苦難.
為什麼這個世界有罪惡.
因為在故事裡面.
耶穌做了醫治的神蹟.
這是解決苦難的問題.
他也說過罪惡的根源是在撒旦.
在故事當中.
人們跟耶穌說了一句話.
他說耶穌你被鬼附了.
其實真正被鬼附的不是耶穌.
而是所有不信耶穌的人.
所有叛逆他不願相信他的人.
於是耶穌告訴我們.
這個世界為什麼有苦難.
為什麼會有罪惡.
是因為撒旦的人有.
因為人的墮落.
最後人的人生的意義何在.
罪的救贖在哪裡.
答案是再跟隨耶穌.
敬拜真神回歸.
因為我們起初罪就是不敬拜真神.
現在我們要跟隨耶穌回去敬拜真神.

$^{201}$當然今天的故事裡面.
有另外一個subtitle.
裡面有一個小節.
耶穌告訴我們.
我們怎麼可以得到罪的救贖.
就是在基督裡面領受割禮.
所以我今天如果將講道的題目重寫的時候.
我可以說榮耀福音之領受割禮都可以.
因為你會發覺你跟著我去到最後.
你會發現耶穌最後講的割禮.
其實是和我們的永生息息相關的.
於是回到現在這個圖畫.
你會發現左邊其實是黑暗.
我們其實在黑暗的當中.
右邊是光明.
耶穌才是光明.
左邊是永遠的審判.
右邊是永恆的天國.
唯有耶穌.
耶穌的十字架空墳墓.
可以拯救我們從黑暗的審判逃跑.
直到上帝光明永恆的天國.
這個是我們所相信的福音.
接著我們現在要回到楊福音第七章裡面.
剛才問的最後一個問題的回答.
耶穌說我怎麼可以得到罪的救贖.
耶穌給我們的答案是在基督裡面領受割禮.
這八個字帶我們回到約翰福音.
上年七月我們開始研讀新約聖經裡面的約翰福音.
我們提到的題目是耶穌是上帝榮耀的君王.
我們花了半年的時間完成了約翰福音第一章至第四章.
接著今年一月我們開始閱讀約翰福音第五章.
你會發現我們今天正在讀的是約翰福音第七章.
但原來和第五章是息息相關的.
所以我們必須要帶你回去看第五章.
聖經學者告訴我們自從約翰福音第五章開始.
一直到第十章.
約翰誕我們進入了一個節慶的周期.
第五章有一個背景是安息日.
第六章有一個背景是逾越節.

$^{241}$第七章有一個背景是駐彭節.
接著第七,八,九都是駐彭節.
第十章有一個背景是獻殿節.
最後在第十一章裡面我們回到逾越節.
在這個節慶的周期裡面.
耶穌通過了祂的神蹟和教導.
以節期為一個背景啟示祂自己的身份和使命.
正如我之前不斷強調.
祂越是清楚啟示祂自己.
猶太的領袖越是恨惡逼迫祂.
所有的一切其實是從約翰福音第五章.
安息日風雲那裡開始的.
我們今天要講這段聖經.
我們已經來到約翰福音第七章.
約翰福音第七章其實.
聖經學者相信大概已經是約翰福音五章之後的.
一年半十八個月之後的事情.
但是猶太領袖對於耶穌在安息日所行的神蹟記號.
仍然是耿耿於懷.
所以你一定要記得那個故事.
我不知道你還記不記得他.
十八個月前在安息日做了些什麼事.
在約翰福音第五章裡面.
耶穌來到耶路撒冷.
在靠近陽門的不是大池.
醫治了一個騎腿三十八年的人.
這是故事的大綱.
然而我們說這個不是一個普通的神蹟.
這個是十字架的記號.
是一個救贖的記號.
在故事當中.
上帝的羔羊.
耶穌有另外一個小名.
他叫做上帝的羔羊來到聖殿的陽門.
經過陽門.
羔羊經過陽門.
就是要被獻上作為一個贖罪祭.
然後耶穌又是永生的活水.
因為這個醫治是在哪裡發生的.
是在那個池裡面發生的.

$^{281}$耶穌是永生的活水.
祂出現在等待救贖的人群當中.
那時候的人群可能有幾百幾千人.
而在眾多的病人裡面.
可能上千的人裡面.
耶穌選了一個騎腿三十八年的人去醫治.
不是因為那個人獨有信心.
也不是因為他特別可憐.
而是因為三十八年剛剛好.
就是以色列人出了埃及之後.
因為他們的叛逆不順.
流浪,曠野.
沒辦法進入上帝應許之地的年數.
就是三十八年.
於是耶穌藉著這個天國的記號告訴我們.
祂來是要呼喚救贖,復活.
帶領靈裡面叛逆癱瘓.
罪孽深重的以色列人.
三八.
進入上帝所應許的永恆安息.
接著耶穌再加了一句.
祂說這是聖父和聖子共同努力.
持續工作為要達成的救贖計劃.
那句話就真的爆燈了.
因為祂不是說單單祂自己.
在安息日做了這件事.
祂說其實是我天父.
上帝和我一起做的.
因為那句話.
耶穌的神職和解說.
不單止沒有得到猶太領袖的接納.
而且還引來他們憤怒的逼迫.
想要殺他.
第五章第十八節這樣說.
為了這緣故.
猶太人越發想要殺他.
因為他不但犯了安息日.
安息日是小事.
犯完安息日之後還說.
是上帝和他一起做的.

$^{321}$他還稱上帝是他的父.
把自己和上帝捍衛同等.
這是一年半之前所發生的事.
這是很重要的.
我和你說了這麼大的故事.
然後我想你記得一件事.
這是一年半之前發生的事.
信不信由你.
這件事仍然餘波未了.
將要成為今天故事的重要背景.
今天我們要研讀的經文是.
《約翰福音》第十第七章.
第十一節開始一直到三十六節.
故事的篇幅是比較長的.
而且故事裡面有很多人物.
很多對話很多事情.
幸好聖經學者告訴我們.
這段聖經前後呼應.
形成一個ABCDE DCBA的文學架構.
為什麼呢.
因為我們可以不看其他的雜項.
我們可以專注在它的中心.
你這樣有了這個架構.
我們就可以清晰看到整段經文的重心和主旨.
我隨便和你看看外圍發生了什麼事.
你會發現外面AA.
兩次提到有人對耶穌議論紛紛.
耶穌到底是誰.
很明顯他告訴你這是前後腳.
然後中間有BB有兩條問題.
他沒上過學為什麼可以這麼熟悉聖經.
第二個問題是他行的神蹟這麼多.
如果基督真的來的時候.
會不會比耶穌行的神蹟更多.
兩條問題.
然後CC.
耶穌舉目望天.
第一個他說我的教導是從那差我來的那位來的.
然後下面那個C.
我是從上面他上面來的.

$^{361}$是他差派我.
兩個都是和耶穌被差派來到這個世界有關係.
然後這些又是兩個問題.
第一個問題是你被鬼附誰想要殺你呢.
你是一個被鬼附的人.
誰有人想要殺你.
下面有另外一個問題.
你這不就是他們想要殺的人嗎.
於是整個架構的中心是E.
在E裡面我們有故事的中心和主旨.
耶穌回答說.
在安息日做兩件事是可以的.
第一件事是行割禮.
第二件事是使人存如.
第一件事是按照猶太傳統是可以的.
一會兒回來再說你就會明白了.
後面那件事就是耶穌所做的.
我想你這樣去想.
一個是Part A一個是Part B.
耶穌將這兩個故事連體變成.
連成一體.
安息日行割禮是Part A.
安息日使人存如是Part B.
因為時間關係.
我們今天是不會看完所有經文的.
也不需要的.
因為很多經文是襯托的.
我會集中在頭一半的經文.
接著我想和你zoom in.
去看割禮的意義.
究竟什麼是割禮.
為什麼耶穌要這麼強調.
安息日行割禮這件事情.
但我們進入今天的經文.
楊鳳蔭第七章第十一節.
主彭節的節期間.
猶太人尋找耶穌.
說他在那裡.
為什麼要尋找耶穌.
因為想要殺他.

$^{401}$整件事是發生在十八個月前的衝突.
而這十八個月前的衝突.
安息日的衝突.
仍然對於猶太人領袖來說.
歷歷在目記憶猶新.
很記仇.
於是第十二節.
人群中有很多人對他議論紛紛.
另外有人說他是好人.
有人對耶穌念念不忘.
覺得他是一個好人.
也有人說他是迷惑群眾.
可惜沒有人敢公開的評論.
因為他們知道猶太人或者猶太領袖.
對耶穌恨之入骨.
他們恨惡耶穌想要將他殺死.
但是整件事是基於十八個月前的一件事.
我們現在想的是十八個月前的一件事.
這樣也行.
第十四節.
節期已經過了一半.
耶穌上耶路撒冷的聖殿去教導人.
節期是七天.
所以七天去了一半.
即是三天半左右.
於是在第四天.
耶穌終於上聖殿去教導人.
我要你想一想.
耶穌終於再次在耶路撒冷公開露面.
已經是十八個月前的事情.
你記不記得.
其實耶穌十八個月沒有見過他.
有幾次有大節來.
耶穌都沒有在聖殿講道.
突然間十八個月之後.
耶穌終於再次公開露面.
如果你知道當時的情景.
就是很多有名的拉比.
每個學院學派.
一早已經霸佔了最有利的位置.

$^{441}$在聖殿裡面.
聖殿很大.
大家找了自己的位置.
擺了一個攤位.
去收納他的粉絲.
收納他的followers.
而耶穌遲來了三天.
我們問他還有沒有可能脫穎而出呢.
我們大家都在這個時候搶followers.
搶fans.
第十五節.
猶太人猶太的領袖驚訝地說.
這個人沒有學過怎麼那樣熟悉經典呢.
猶太人驚訝.
連他的敵人.
這個猶太人不是普通的猶太人.
是猶太人的領袖.
都很欣賞很肯定耶穌的才華.
十八個月不見.
耶穌風采依然.
十八個月不見.
耶穌功力大增.
你明白嗎.
以至他們雖然是他的敵人.
死敵仍然驚訝耶穌的才華.
然而school matters.
學校是很重要的.
猶太人和中國人是一樣的.
都很著重名校名師名門.
耶穌什麼都不是.
於是第七章第十六節.
耶穌這樣回答他們.
耶穌回答他們說.
我的教導不是我自己的.
而是差我來的那一位.
耶穌什麼都不是.
其實一什麼都不是都無所謂的.
家教自學都是可以被認可.
被肯定的.
但是你有沒有留意.

$^{481}$這個不是耶穌的回答.
我想幫你去解讀.
耶穌究竟在回答什麼.
耶穌說.
我所屬的學校.
你去不到.
你的學校在半山.
我的學校在天上.
你知不知.
根本無法比較.
天父是我的老師.
耶穌不是說我是家教出來的.
所以他說一句.
我的教導不出於我自己.
我不是homeschooler.
我不是自己在家裡.
剁磨出來的.
其實我是有學校的.
不過我的學校你去不到.
你頂多去半山.
我的學校在天上.
天父是我的老師.
不是我沒有名校.
不是我沒有名師.
是你有眼無珠.
核眼耳聾.
井底之蛙.
被困於黑洞裡面.
你根本不知道我的底細.
這個耶穌真的在耍他們.
我不是在家自學的.
我的教導不出於我自己.
我的教導是差我來的那位.
差我來的那位.
如果你熟習楊富英就知道.
在第五章裡面.
耶穌重複又重複說.
我是天父差下來的.
差我來的是天父.
耶穌分明就是要重提舊事.

$^{521}$耶穌分明就是要重揭舊有的傷疤.
reopen the old wounds.
你知不知道.
祂專題.
你記得十八個月前.
我一直跟你說.
我是有人差下來的.
而差我的是天父.
現在耶穌又再說那件事.
我是差我來那位教我的.
而祂很明顯是重提十八個月前的事情.
原來對於十八個月前的安息日事件.
僅僅於懷的不僅僅是猶太領袖.
連耶穌也是僅僅於懷.
於是十八個月之後.
耶穌要重提十八個月前的安息日的風雲.
接下來.
耶穌還要在他們的傷口上灑鹽加辣.
第七章第十七節十八節.
人若立志要進行上帝的旨意.
就會知道這教道究竟是出於上帝.
還是我憑著自己說的.
憑著自己說的人是尋求自己的榮耀.
但那尋求差他來的那位榮耀的人.
他是真誠的.
在他裡面沒有不然.
我想幫你翻譯一下耶穌究竟說了什麼.
耶穌是這樣說的.
你們如果真的從心裡面.
愛慕,跟隨,惹禍,上帝.
舊約的那個上帝.
你們所有都是以色列人.
你們熟讀舊約的.
如果你們真的從心裡面.
愛慕,跟隨,惹禍,上帝.
那你就必定會知道我的教導出於上帝.
好搞笑.
你們就會很自然地.
如果你真的敬拜,跟隨,愛慕,舊約的上帝的話.
你就會很自然地認出我.

$^{561}$認識我.
愛慕我.
跟隨我.
肯定我就是上帝所差來的那一位.
其實如果你正在讀《讓福音》.
你知道其實在《讓福音》裡面.
已經有一群人.
很自然地認出耶穌的.
包括施洗約翰.
包括拿彈業.
包括那個神奇女俠.
撒把戲阿普仁.
包括撒瑪利亞人當中的信徒.
最後說他是世界的救主.
其實是有些這樣的人.
不過就是不是以色列人.
於是耶穌在遞禮他們.
耶穌告訴他們.
但是你們這群人有眼無珠.
黑眼異龍.
井底之蛙.
被困黑洞.
接著耶穌再說.
還有一件事.
你們尋求的是自己的榮耀.
是世界的榮耀.
之前耶穌的親弟弟們也是一樣.
為什麼你們不信我.
為什麼你們不愛我.
為什麼你們不跟從我.
其實很簡單.
因為你心有所屬.
因為你一早就愛這個世界.
因為你一早就愛這個地上的榮耀.
去到這裡我們都要問問我們自己.
親愛的朋友親愛的弟兄姊妹.
或者你都可以想一下.
你來Crossway聽到都可能有一段日子.
你聽到的是真道.
你聽到的是真福音.

$^{601}$但是為什麼你不信呢.
為什麼你經常在猶豫呢.
為什麼你又說信又不信呢.
為什麼你不全心全意.
All in的去跟隨上帝呢.
這是不是福音的問題.
這是不是上帝的問題.
這是不是耶穌的問題.
還是你自己的問題.
是因為你貪愛世界.
因為你追求不低地上的榮耀.
因為你不願捨棄.
因為你不願背負十家.
可能去到某個位有人會說.
其實是林牧師的問題.
說真的.
是他將福音講得太複雜.
所以我信不到.
是他將跟隨耶穌講得太艱難.
所以我沒辦法跟隨.
我的回應是.
你不用信我.
你也千萬不要信我.
但是你必須要相信聖經.
你一定要好好思想.
究竟你所聽的道.
我所講的道.
是不是貼近聖經.
還是我偏離真道.
如果我是貼近聖經.
而你不信.
你知道你的問題.
你的問題不是我.
你的問題是上帝.
我們回到《讓富欽》第七章.
我們繼續讀第七章第十九節.
第二十節.
「我誓不是傳了律法給你們嗎?」.
「你們卻沒有一個人守律法」.
「為什麼不守律法?」.

$^{641}$「為什麼你們都想要殺我呢?」.
眾人回答說.
「你是被鬼附了」.
「誰想要殺你呢?」.
在《讓富欽》裡面.
耶穌至少有五次被人稱為被鬼附.
其實事實剛剛好相反.
被鬼附的不是耶穌.
被鬼附的是不相信耶穌.
不愛慕耶穌.
不跟隨耶穌的人.
但是我們去到第七章第二十一節.
進入我們今天福音故事的.
戲玉和高潮.
第七章第二十一節這麼說.
耶穌回答說.
對他們說.
「我做了一件事」.
「你們都驚訝」.
你會發現耶穌又重提舊事.
這個舊事不是兩章之前.
這個舊事是十八個月前所發生的事件.
正如我再次強調一樣東西.
對於這件安息日的風雲.
耶穌耿耿於懷.
耿耿於懷不單只是猶太的領袖.
耶穌一樣是耿耿於懷.
於是耶穌開始了.
第七章第二十二節.
「摩西傳國禮給你們」.
其實國禮不是從摩西開始的.
而是從列祖阿伯漢那裡開始.
一會兒再說.
列祖開始的.
「你們就在安息日給人行國禮」.
我想請大家留意.
這是一個全新的教導.
全新的.
是之前沒有聽過的.
也是國禮會成為整段對話裡面的核心.

$^{681}$他說「摩西傳國禮給你們」.
這是出於利未記第十二章一至三節.
我讀一讀給大家聽.
又說「吩咐摩西說.
你要吩咐以色列人說.
婦人若懷孕生男孩.
第八天要給孩子行國禮」.
問題來了.
如果他生出來的第八天.
剛剛好是安息日.
那怎麼辦呢.
猶太人的傳統說「照割」.
所謂猶太人傳統.
就是猶太的律法師.
對於舊約律法的演繹.
變成了他的傳統.
他演繹的傳統告訴你.
舊約沒有說.
現在去到安息日.
他又去第八天.
那割不割.
他說割.
他說「你們的傳統說照割的」.
其實你想想這個小朋友.
真的是幸運童子.
如果他是第八天是安息.
就是說他出生的那天也是安息日.
他是生於安息.
割於安息.
很棒吧.
因為你數第八天.
剛好又去到安息日.
在這裡耶穌又加了一句.
其實不是出於摩西開始的.
是從列祖開始的.
但是我們回到創世記第十七章.
待會回來才跟你讀.
我們繼續讀.
第七章第23節.
「人若在安息日受國禮.

$^{721}$是為了不違背摩西的律法」.
因為摩西說第八天要割.
就算安息日割不割.
照割.
因為要遵守摩西.
當然這個是他們演繹.
他們演繹也可以說.
不如第九天才割.
也可以的.
不過當時的演繹就是.
照割.
耶穌說.
「人在安息日受國禮.
是為了不違背摩西的律法.
我在安息日使一個人痊癒了.
你們就向我發怒.
何必憑外表斷定是非.
總要按公平斷定是非」.
我們問耶穌在說什麼.
耶穌究竟在說什麼.
一般聖經學者的解釋.
明白是這樣的.
耶穌是為自己辯護.
解釋他為什麼在安息日治病.
他在安息日治病.
為38年的人治病.
是合理和合法的.
但你想想這件事.
是發生在18個月前.
你覺得.
耶穌是不是這麼小器呢.
這麼記仇.
為什麼要18年.
18.
不是18年.
18個月前的事情.
現在都要拿回來說.
一會兒我再跟你開估.
但是這裡.
這裡說的是.

$^{761}$大部分的聖經神學家的理解.
大概的邏輯是這樣的.
如果按照現在猶太人的傳統.
為了不違背摩西的律法.
就是要在第八天為嬰兒行割禮.
容許人在安息日行割禮.
使以色列的男孩得潔淨.
這是割禮的其中一個意義.
一會兒我會回來再跟你深切思想.
割禮的意義是什麼.
其中一個意義就是.
為這個男孩子使他得潔淨.
於是耶穌.
從小的潔淨到大潔淨.
耶穌說.
我現在在安息日使一個殘缺不潔的人.
完全康復治療潔淨.
為什麼這個人已經38年沒辦法去敬拜上帝.
38年是因為他不乾淨.
因為他是一個殘缺的人.
殘缺的人在當時的傳統裡面.
他就是大罪.
於是他沒辦法見證.
現在耶穌在安息日裡面.
使這個殘缺不潔的人.
康復,痊癒,潔淨.
耶穌說.
豈不是一個更美的事情.
更加貼近安息的意志.
是不是應該得到你們的接納和鼓掌呢.
我剛才告訴你這個邏輯.
是大部分聖經學者和牧師們的看見和理解.
耶穌是在抵抗他自己.
他是告訴他為什麼他這麼做是OK的.
但我認為這個解說.
是有些太尋常,太簡單.
沒有考慮到故事背景的不尋常和複雜性.
讓我解釋一下.
在整個故事裡面最關鍵的地方.
就是第五章和第七章.

$^{801}$其實相隔一年半.
已經相隔了一年半.
為什麼還要舊事重提.
耶穌不說就算了.
重新來過.
為什麼要重新引爆炸彈.
為什麼要舊事重提.
耶穌是不是真的這麼記仇呢.
這是第一.
第二.
耶穌其實不一定要用割禮.
來做安息日可以做事的例子.
例如.
猶太的傳統是容許安息日打聖戰的.
耶穌大可以用這個角度來大寫文章.
你知道我現在在做什麼.
天國降臨.
我現在勝過撒旦的權勢.
耶穌可以說其實他在打聖戰.
但他不用聖戰.
他用割禮.
第三.
在讓富欽第五章結束的時候.
耶穌開始提及摩西.
然後他的對話就中斷了.
現在去到讓富欽第七章.
耶穌又主動提起摩西.
因為他說割禮是摩西給你割禮.
耶穌之前第五章裡面提起摩西.
接著就斷了纜.
接著第七章.
他自己再次回到耶路撒冷.
他主動提起摩西.
有沒有想過為什麼.
或者我們可以這樣想像一下.
一年半之前的那天.
耶穌本來是預備繼續說割禮的.
有沒有想過.
有沒有這樣的可能.
一年半前耶穌其實還沒說完.

$^{841}$他想繼續在他的神蹟記號前面.
再加上列祖割禮.
接著將神蹟.
醫治的神蹟和割禮是連在一起的.
不過很可惜.
那天猶太的領袖就已經爆燈了.
他沒辦法再聽.
沒辦法再聽耶穌下去.
於是情況轉壞.
急轉直下.
耶穌指好.
耶穌以前習慣了會做這件事.
神蹟般的在人群中消失.
於是第五章就完結了.
因為耶穌突然間神蹟般的離開.
這些全部都是我的想像.
但是是有理由的想像.
猶太領袖心猶不禁對此事耿耿於懷.
因為事情還沒完.
他們想抓耶穌殺祂.
對於耶穌來說一樣都是耿耿於懷.
因為祂還有最後那句還沒說出來.
故事還沒完結就被人趕走了.
於是雙方面.
無論是猶太的領袖耶穌.
都是在等待時機再續前緣.
於是到了第七章.
他們終於可以再續前緣.
如果我這個猜想是合理的話.
列祖的割禮其實就是.
耶穌依次38年騎腿男人的前傳.
他拍了第二套.
然後再拍前面那一套.
故事部分A是列祖的割禮.
故事部分B是安息日復活的依次.
這樣就是一個完整的福音故事.
我們再看多三次聖經.
完結後我再帶你回去看列祖割禮的多重意義.
我們跳到第七章第32,33和34節.
法利賽人聽見群眾對耶穌這樣議論紛紛.

$^{881}$祭司長和法利賽人就打發聖殿的警衛去捉拿他.
於是耶穌說.
我跟你們在一起的時刻不會太久了.
我要回去那差我來的那裡去.
你們找不到我.
你們也不能去我去的地方.
我想用五點來歸納總結今天的故事.
其實幫大家去refresh一下.
讓你記得今天講的故事講了些什麼.
五點是這樣.
人們議論紛紛.
耶穌到底是誰.
第二.
耶穌的教導不是出於自己.
是出於那個差我來的那位.
第三.
猶太領袖打發聖殿的警衛.
要去捉拿耶穌.
要將耶穌置諸死地.
第四.
耶穌說我的時候不久了.
因為我要回到那差我來的那裡.
那個人.
接著在我回去之前.
我要承傳一件事.
是什麼呢.
使一個聖約中受割禮的人.
全月走進按食.
我想你留意三件事.
我叫他做ABC.
因為我這裡有12345.
A就是人們議論紛紛.
耶穌到底是誰.
這個是人生最終極的問題.
耶穌是誰.
也是人生最終極的答案.
因為如果你答到耶穌是誰.
你就可以得到人生終極的拯救.
接著Part B.
我要你考慮是234.

$^{921}$在234這三點裡面.
顯示了耶穌的Round Trip Journey.
耶穌是買了來回機票的.
耶穌一下來的時候.
從差他那裡的人來.
為要有一天他會回去.
但是他必須經過第三.
他必須經過十字架.
他必須要被捉拿.
他必須要被殺.
而他被捉拿被殺.
不是一個意外.
也都不是一個悲劇.
是他降世的使命和目的.
第三點C點.
耶穌降世的目的是什麼呢.
我用一句話總結.
使一個在聖約裡面受割禮的人.
痊癒.
接著可以走進安城.
我將兩個故事加起來.
事隔一年半.
約翰福音第五章的Part B.
再加上約翰福音第七章的Part A.
終於諸聯璧合.
合而為一.
你可以這樣想.
有兩塊肉失落了.
於是一年半之後.
終於可以合而為一.
我想你想一下.
這個騎腿三十八年的人.
其實是一個以色列的男人.
他有沒有受割禮.
他有受割禮.
他是不是第八天.
是否安息日受割禮.
我就不知道.
但他是一個有受割禮的以色列人.
耶穌來不是單單去拯救世人.

$^{961}$耶穌來是要拯救.
在約中受了割禮的以色列人.
這個騎腿三十八年的人.
於是Part A和Part B.
割禮和騎腿三十八年的人.
兩個故事是息息相關.
不可分割的.
當你來到割禮.
在舊約聖經裡面的重要意義.
我想帶你去看.
會帶你去看《創世記》的第十七章.
我們只看兩節的聖經.
《創世記》第十七章第十節.
「這就是我與你.
以及你的後裔所納的藥.
是你們所當遵守的.
你們所有的男子都要受割禮」.
這是上帝與亞伯拉罕所說的話.
第十九節.
「你妻子撒拉必為你生一個兒子.
你要給他起名叫以撒.
我要與他建立我的藥.
成為他後裔永遠的藥」.
這裡所說的一開始.
我們都要交代一下什麼叫割禮.
簡單來說割禮就是割掉.
男性生殖器官的包皮.
你知這麼多就可以了.
不用去Google它.
這個割禮將要成為耶和華上帝.
和亞伯拉罕與他所有後裔所納藥的記號.
這個記號在舊約裡面有多重的意義.
我想說在舊約裡面.
因為這個割禮是舊約的東西.
你去到新約我們不再行割禮.
但是在新約裡面他講一個割禮.
不是用人手所做的割禮.
所以割禮和基督徒有沒有關係.
仍然有關係.
所以我們現在說的割禮.

$^{1001}$在舊約裡面有多重意義.
但是對於基督徒而言.
儘管我們不是猶太人.
儘管我們身上未必有割禮.
儘管我們不一定是男性.
但是割禮的屬靈意義.
對於我們而言.
仍然是最真實最寶貴的東西.
所以我們必須要去明白什麼叫割禮.
我先說第一重的意義.
割禮的第一層意義是切割分離.
切割分離.
割禮就是切割分離.
切割分離是什麼呢.
我想你想想.
亞當是上帝所做的第一個人.
亞當是人類的代表.
是人類的始祖.
是人類的首領.
但是亞當犯下滔天的大罪.
你繼續讀《創世記》的時候.
你會這樣去想像.
通過亞當.
通過亞當生殖器的器官.
通過亞當的生殖器官.
人類的罪性從一代傳到另外一代.
在人類的歷史中繼續蔓延.
是不是.
你明白嗎.
因為所有下面的罪人.
全部都是通過亞當的生殖器而來的.
所謂割禮.
就是代表這個出生的新生命.
需要和起初的亞當有個切割分離.
你明白嗎.
為什麼要切割分離.
因為告訴你.
亞當的罪一定要被切割.
唯有切割才可以使他進入上帝的新創造.
所以聖上學家牧師經常在討論.

$^{1041}$為什麼要在第八天切割.
我不知道你有沒有聽過.
有些人說第八天是特別不痛的.
是有這樣的醫學理由.
但是加爾文有這樣的想法.
他說因為第八天是創造之後的一個新創造.
上帝用了七天.
就是創造了七天.
於是第八天是一個新創造.
所以第一個割禮的意義.
就是與亞當切割使這個孩子.
可以進入一個新的創造.
這是第一個.
割禮的第二層的意義是什麼呢.
就是流血逃難.
流血逃難是什麼意思.
有聖經學者指出.
割禮是在創世記第十七章裡面頒佈的.
接著創世記第十八章裡面.
記載著一個天降的審判.
創世記第十八章裡面記載了什麼.
所多瑪俄摩拉被上帝的審判.
於是割禮中所流出的血.
彷彿就好像逾越節羔羊的血.
庇護了上帝的子民逃過審判.
大家明白嗎.
創世記第十七章裡面這個割禮的血.
成為了.
The Passover before the Passover.
逾越節之前的逾越節.
是一個逾越節的影子.
所以這個是割禮的第二重意思.
割禮的第三重意思是什麼呢.
割禮的第三重意思是分別為聖.
原來當時在以色列人沒有割禮之前.
古時的埃及已經早有割禮這個傳統.
在埃及.
大家記得亞伯拉罕其實去過埃及.
所以亞伯拉罕應該是很熟悉.
埃及人的割禮的傳統.

$^{1081}$在埃及裡面的割禮.
是預留給埃及的君王和他的王室成員.
再加上埃及的祭司.
所以埃及的割禮的意思是什麼呢.
就是將埃及的君王和他的王室成員.
再加上埃及的祭司分別為聖.
使他們忠心竭力.
服侍埃及的神明.
OK.
現在上帝用一個已有的所謂的割禮.
但是和埃及不一樣的是.
亞伯拉罕和他的後裔全民皆祭司.
所以你去到出埃及記第十九章第六節的時候.
他會說你們要歸我作祭司的國度.
所有以色列人都成為全民皆祭司.
不是只是某一些人.
才是奉獻給上帝.
所有的人都是奉獻給上帝.
分別為聖.
第四.
割禮的第四層意義是.
盡心愛主.
割禮不單只是身體上的記號.
割禮代表了上帝的子民盡心.
在心裡面受割禮.
心裡面盡心愛主的納節.
我們知道你繼續讀舊約聖經.
你會意識到有一個很大的問題.
就是他們受了身體的割禮.
但是他心裡面沒有受割禮.
他不是一個真意識的人.
聖經不斷告訴我們.
唯有深受割禮盡心愛主.
才可以成為真正上帝的子民.
最後.
割禮的第五層的意思是什麼呢.
等待救主.
原來割禮的設立不單只是為亞伯拉罕.
你知道在《創世記》第十七章裡面.
亞伯拉罕多少歲.

$^{1121}$99歲.
亞伯拉罕當時有沒有兒子.
有的 叫以撒瑪利.
但是這個割禮不是為以撒瑪利去預備.
這個割禮是為以撒來說的.
割禮要等到以撒出來.
所以割禮不單只是為99歲的亞伯拉罕.
更是為明年將要出生的應許之子.
以撒所生.
而且這個割禮不單只是給以撒.
還要成為他後裔.
他有一個他的種子.
這個後裔是誰呢.
女人的後裔.
在《創世記》裡面一直說的.
後裔等待一個後裔的來臨.
就是你的後裔和女人的後裔要彼此為敵.
他要傷他的頭.
他要傷他的腳筋.
我們在等待這個後裔的出現.
這個後裔是誰呢.
這個後裔就是人類的救主耶穌基督.
你有沒有想過.
耶穌出生之後在第八天受割禮.
是耶穌一生裡面第一次流血.
你當他好命.
一生裡面耶穌流過兩次血.
第一次流血在哪裡.
在他第八天的時候.
第二次他死那天.
流盡所有的寶血.
所以第一次的流血.
已經在預表他將來在十字架上要捨身流血.
是不是很像耶穌的婚宴.
在婚宴裡面隨便走.
耶穌只是兩次喝過酒.
一次在婚宴裡面.
另一次就是最後.
盲顫所有的東西.
彷彿告訴我們.

$^{1161}$耶穌將來要為我們所流的血.
我想用五點來總結.
幫大家總結一下割禮的意義.
五個非常有意思的意義.
再說一次.
這些意義不單只是給意識的人.
也是給所有基督徒裡面.
所有的屬靈的意義.
除掉始祖亞當·聚根進入生生.
第二天降烈火.
索多瑪·阿摩拉先血庇護逃過審判.
第三分別為聖.
這個背景是埃及人的割禮.
分別為聖君王祭司服侍真神.
第四聖約中的子民深受割禮.
盡心愛主.
最後應許之子永約的後裔.
等待耶穌救主的來臨.
當我們仔細去閱讀舊約聖經的時候.
你會發現.
上帝重複又重複呼籲意識的人.
不要單受身體上的割禮.
身體上的割禮是沒有意義的.
一定要深受割禮才可以盡心愛主.
才是上帝真正的子民.
我想跟你讀幾段舊約聖經.
你就會明白.
你還沒記第26章第40節.
他們要承認自己的罪孽和祖先的罪孽.
就是背叛我與我作對.
我也行事與他們作對.
把他們遣送到敵人之地.
這是在說巴比倫的事情.
那時他們未受割禮的心.
身體受了割禮.
但是心靈不受割禮.
他們的心若肯謙卑.
這就是深受割禮.
也服了罪孽的懲罰.
我就要紀念我和雅各的若.

$^{1201}$紀念我和以撒的若.
紀念我和亞伯拉罕的若.
我也要紀念這地.
使他們從被擄中回歸.
《申命記》第6章第5節.
你要盡心盡性盡力愛耶和華.
你的上帝.
你怎麼可以盡心呢.
除非你深受割禮.
第10章第16節.
心要受割禮.
不可再硬著頸看.
你唯有深受割禮.
才可以盡心愛耶和華.
《耶利米蘇》第4章第4節.
猶太人和耶路撒冷的居民.
你們當自行割禮.
心裡面的割禮.
歸耶和華.
將你們心裡的污穢除掉.
這些所有的呼召.
也是向基督徒所說的.
我們一定要深受割禮.
真正的迴轉.
才可以愛上帝.
在結束的時候.
我想帶你看多一段.
新約的聖經.
我想帶你看.
歌羅西書第2章第11節.
上面那個是你中文的版本.
下面那個是我的翻譯.
我們讀下面那個版本.
歌羅西書第2章第11節.
在他.
這個他就是耶穌基督裡面.
你們也受了不是人手所行的割禮.
我們不是在說身體上的割禮.
而是心裡面的割禮.
就是在基督的割禮中.

$^{1241}$你們脫去了肉體的情慾.
我們問什麼是基督的割禮.
簡單來說.
基督的割禮就是他在十字架上的.
流血,寫命,完全的贖罪.
本來的意思是什麼呢.
本來的意思是當我們迴轉悔改.
當我們專注仰望耶穌的十字架.
在耶穌基督的割禮當中.
我們領受了不是人手所行的割禮.
那個割禮是什麼呢.
有什麼對我們影響呢.
就是脫去了肉體的情慾.
當我們迴轉悔改.
我又想帶你去十字架.
又想帶你去看那段聖經.
當我們迴轉悔改.
當我們專注仰望耶穌十架的時候.
我們就會因真信與他聯合.
與耶穌基督聯合.
在耶穌的割禮裡面.
我們會脫去肉體的情慾.
什麼意思呢.
就是十字架的耶穌會將肉體的情慾.
從生命的寶座上擊倒下來.
你明白嗎.
當我們仰望耶穌的時候.
彷彿耶穌進入了我們的人生裡面.
在我們人生的寶座上.
將從前的老我.
肉體的情慾擊倒下來.
是通過仰望耶穌.
耶穌進入我們的人生.
耶穌擊倒我們從前的老我.
使我們不再被肉體或撒旦所佔據.
而是開始一個新的人生.
保羅繼續說.
還有兩節.
歌羅西書第二章第十二節.
你們既受洗.

$^{1281}$現在從割禮到洗禮了.
與他一同埋葬.
也就在此禮上.
洗禮上面.
因信那使他從死人中復活的上帝.
的作為跟他一同復活.
你們從前在過飯和.
未受割禮的肉身中死了.
上帝使你們與耶穌一同活過來.
一個是割禮.
一個是洗禮.
洗禮是活過來.
赦免你們的過飯.
我想你看看保羅是怎樣將割禮和洗禮連上.
那個割禮是一個真實的割禮.
是一個不是人手行的割禮.
那個洗禮也是一個真實.
一個真信徒的洗禮.
我們可以這樣說.
如果割禮強調的是脫去肉身的情慾.
這個肉身的情慾是在說老我.
如果割禮是代表我們脫去舊有的生命.
包括我們舊有錯誤的人生方向.
舊有的人生方向是這裡.
這裡這裡這裡這裡.
現在你要脫去這個舊有的人生方向.
以前追求地上的榮耀.
現在要追求上帝.
是脫去舊有錯誤的方向.
脫去對世俗榮耀的追求.
洗禮強調的是復活的生命.
一個是負面的一個是正面的.
洗禮強調的是我們被聖靈所重生.
我們現在注目著天上永恆的國度.
天上的榮耀天上的財寶.
於是我們割禮我們離開.
我們受洗與基督同死同復活.
我們復活我們開始全人全心的追求.
上帝的國和上帝的義.
耶穌的割禮.

$^{1321}$耶穌的洗禮會成為我們信心仰望的方向和焦點.
你明白嗎?.
你怎麼可以得勝這個世界.
你怎麼可以得勝你的老我.
注目耶穌的割禮.
耶穌的洗禮.
你發現耶穌的割禮和耶穌的洗禮都很相似.
耶穌的割禮是第八天的割禮.
耶穌的割禮是十字架上面的割禮.
耶穌的洗禮是約旦河的洗禮.
約旦河的洗禮會帶領祂到十字架上面.
耶穌被洗.
耶穌有說過祂十字架上.
我所受的洗你能受嗎?.
那個洗的意思是什麼?.
羅亞洪水的洗禮.
耶穌為我們承擔所有東西.
當我們注目將我們的信心仰望.
恩信看著耶穌的時候.
我們就可以除掉我們的老我.
我們可以開始一個新的人生.
上個星期二.
剛剛過了一個星期二.
在福音班裡面.
我用了兩張大支票來解釋.
表明什麼叫信耶穌.
什麼是基督徒的人生.
第一張大支票是上帝給我們的支票.
我以前有個同事叫做Bonnie Sin.
我說可不可以借你的名字來.
因為你的名字起得很好.
Bonnie Sin 好像Bonnie Sin.
一個最佳的基督徒的名字.
祂生於最終.
但是你想想在耶穌基督裡面.
上帝給我們一個何等偉大的應許.
上帝給我們清義的禮物.
給我們領養的禮物.
給我們誠信的禮物.
給我們榮耀的禮物.

$^{1361}$上帝自己簽了那張支票.
聖父,聖子,聖靈的名字.
現在我們需要領受這張支票.
然後在後面簽自己的名字.
這個叫信心.
這張支票叫求他的義.
但是我告訴你.
基督徒的人生真信心.
是簽兩張支票.
不是一張是兩張.
兩張都要有.
不是二選一.
而是兩張要有.
第二張支票叫求他的國.
第一張叫求他的義.
第二張叫求他的國.
當你簽完這張求他的義之外.
你要簽完求他的國.
才是真正的信心.
這次這張支票是Bonnie Sin給上帝的.
收錢的人是聖父,聖子,聖靈.
是那三億的上帝.
慈愛的救主.
寶貴的十價.
而Bonnie Sin簽下去的是一個無價支票.
將她的人生是一個空頭支票.
沒有號碼的上面寫著的是.
毫無保留,永不後退.
沒有遺憾的獻上自己.
兩張支票加起來.
就是一個基督徒的人生.
就是一個真正的信心.
求他的國,求他的義.
在耶穌基督裡面我們收了割禮.
除去舊的生命.
在耶穌基督裡面我們接受了水禮.
與基督同死同復活.
開始一個全全暫生的人生.
求主幫助我們.
讓我們得著這樣的生命.

$^{1401}$我們祈禱.
去到《養父音》第七章裡面這個故事.
真的很奇怪.
但是後來當我們明白的時候是很奇妙.
為什麼一件十八個月之前所發生的事.
仍然耶穌耿耿於懷.
因為祂想幫我們完成這個救贖的圖畫.
祂想告訴我們.
所有一切是從割禮開始.
割禮是恩典聖若當中的記號.
割禮告訴我們上帝會拯救.
這一群受割禮的人.
上帝會差遣他的獨生兒子.
成為亞伯拉罕唯一獨有的後裔.
這個後裔要在地上接受割禮.
這個後裔要在十字架上為我們承傳救贖.
幫我們深思割禮和我們的關係.
幫我們深思不單止是割禮也是洗禮.
去到新約的時候.
我們再不需要有割禮.
但同時在靈裡的割禮仍然是真實.
是自保.
於是我們在耶穌基督裡面.
我們接受祂的割禮.
又領受祂的水禮.
一個是負面的一個是正面的.
一個是除去舊時的人生.
一個是進入耶穌復活的生命.
幫助我們求他的兒.
幫助我們求他的國.
幫助我們窮一生去榮耀.
祂將我們的生命獻上.
多謝你.
我們感激祈禱.
奚奉主耶穌基督.
\newpage



\section{以賽亞書 49:1-6}
\label{sec:032107}
\textbf{未竟全功?}~~~~周弘毅牧師
\newline
\newline
連結: \href{http://www.church.com.hk/acms/content.asp?site=cdc\&op=show\&type=product\&code=032107\&layout=sermon}{\texttt{032107}} ~~~~ 語音日期: 2026-05-09
\newline
\newline
\hyperref[sec:032105]{\small{< < < PREV SERMON < < <}}
~
\hyperref[sec:index_chronic]{\small{[返順時目]}}
~
\hyperref[sec:032108]{\small{> > > NEXT SERMON > > >}}
\newline
\newline
以賽亞書 49:1-6
\newline
\begin{longtable}{cl}
\hline
\hline
章節 & 經文 (和合本修訂版)\\
\hline
49:1 & \begin{tabularx}{0.7\textwidth}{X} 眾海島啊，當聽從我！遠方的眾民哪，要留心聽！自出母胎，耶和華就選召我；自出母腹，他就稱呼我的名。 \end{tabularx} \\ \\ \relax
49:2 & \begin{tabularx}{0.7\textwidth}{X} 他使我的口如快刀，把我藏在他手蔭之下；又使我成為磨利的箭，把我藏在他箭袋之中； \end{tabularx} \\ \\ \relax
49:3 & \begin{tabularx}{0.7\textwidth}{X} 對我說：「你是我的僕人以色列；我必因你得榮耀。」 \end{tabularx} \\ \\ \relax
49:4 & \begin{tabularx}{0.7\textwidth}{X} 我卻說：「我勞碌是徒然，我盡力是虛無虛空。耶和華誠然以公平待我，我的賞賜在我的神那裡。」 \end{tabularx} \\ \\ \relax
49:5 & \begin{tabularx}{0.7\textwidth}{X} 現在耶和華說話，他從我出母胎，就造我作他的僕人，要使雅各歸向他，使以色列聚集在他那裡。耶和華看我為尊貴，我的神是我的力量。 \end{tabularx} \\ \\ \relax
49:6 & \begin{tabularx}{0.7\textwidth}{X} 他說：「你作我的僕人，使雅各眾支派復興，使以色列中蒙保存的人歸回；然而此事尚小，我還要使你作萬邦之光，使你施行我的救恩，直到地極。」 \end{tabularx} \\ \\
[1ex]
\hline
\hline
\end{longtable}
$^{1}$頂至每午安.
這個題目稍後會解釋.
因為這四個字每個字都認識.
加起來是甚麼呢.
好.
來到五月份.
大概這個時間.
我相信大家可能都會留意.
一個日子.
因為周邊都會有不同的廣告.
來提醒你.
明天.
是甚麼日子.
想起來了.
母親節.
所以在這裡.
先代表教會.
恭祝在座的每一位母親們.
身體健康.
母親節快樂.
好.
上個主日.
我去了外面.
港島.
星期六的時候.
我有機會.
負責向幾位.
姐妹頒發.
成長教練.
栽培員證書.
我在台上.
這樣說.
其實要拿這張證書.
都不容易.
因為除了.
要有課程.
需要學之外.
其實是要有實踐.
要幫助.
弟姐妹成長.

$^{41}$這個.
屬靈成長.
好像照顧孩子.
真是艱難.
作為一位.
地上的母親.
我相信.
都差不多.
甚至可能更加艱難.
因為每個兒女的.
成長.
不同的階段.
都有不同的挑戰.
有一天.
有三個媽媽.
三十多歲和四十多歲.
這三個媽媽.
約在一起.
吃午飯.
聊聊天.
他們討論到一個問題.
究竟在.
小朋友的成長階段.
哪個階段.
最困難.
在座當中.
周圍的媽媽.
可以考慮一下.
二十多歲的媽媽.
當然是小朋友.
兩歲多的時候.
最困難.
英文有句說話.
terrible two.
每逢兩歲至四歲.
基本上.
小朋友開始有自我意識.
然後.
她喜歡唱反調.
很容易情緒失控.

$^{81}$然後無理取鬧.
很惡搞.
另外一個.
三十多歲的媽媽.
說.
兩歲多的小朋友.
到小朋友.
開始讀書.
不知道為什麼.
這麼多墨書.
測驗考試.
很辛苦.
她指著自己.
不是小朋友很辛苦.
然後.
四十多歲的媽媽.
又說.
你們有沒有想像過.
當孩子.
進入青春期.
的時候.
究竟會面對什麼.
她當然覺得很艱難.
另外.
兩個媽媽因為還沒到這個階段.
所以她們就.
答不上嘴.
突然間旁邊桌子.
有一位婆婆就搭訕.
她很嚴肅地說.
她說你們的子女.
到了五十多歲你就知道.
五十多歲的子女.
覺得他們很難搞.
其實.
在不同的.
生命階段.
的確要生命成長.
也是艱難的.
每個階段.

$^{121}$都這麼艱難.
今天我要和大家分享.
到耳塞書的四十九章.
是用.
第二首的僕人詩歌.
作為開始.
揭示神的僕人.
他怎樣扭轉.
上帝.
在這個過程.
看到他們面對的.
困局.
幫助這些.
以色列民復興.
然後.
還可以將舊恩傳到地極.
大家應該記得.
講耳塞書的.
背景的時候.
跟大家提過.
其實當時.
先知耳塞生活在一個.
非常動蕩.
不安的時間.
其實當時.
他們是要面對.
巴比倫老弱的威脅.
這是地上的威脅.
屬靈雙邊.
其實他們同樣充滿疑惑.
因為他們會問.
其實那位.
由他們的祖宗.
承傳下來的.
耶和華上帝.
為何會看著.
他們死亡.
不出手去幫他們.
其實.
神已經回答過.

$^{161}$正正是因為.
他們的列祖.
不斷犯罪離棄神.
所以就帶來了.
今天的困局.
你解到.
不代表解決到.
分得好.
你解釋得到.
不代表解決得到.
可以怎樣.
神在這裡.
藉著.
祂所說的我的僕人.
來扭轉這個困局.
並且成就他們.
從來沒有想過.
更加遠大的目標.
在開始分享.
這一段經文.
先和大家.
讀一讀這六節.
的聖經,預備起.
仲海道.
當聽我言,遠方的民眾.
留心耳聽.
自我出胎.
耶和華就選召我.
自出母腹.
祂就提我的名.
祂使我的口語快到.
將我藏在祂手任之下.
又使我成為.
聖上的箭.
將我藏在.
祂箭袋之中.
對我說.
你是我的僕人以色列.
我必因你得榮耀.
我卻說.

$^{201}$我勞碌是徒然.
我盡力是虛無虛空.
然而我當得的利.
必在耶和華那裡.
我的賞賜.
必在我神那裡.
耶和華從我出胎.
我歸向祂.
使以色列到祂那裡聚集.
原來耶和華看我為尊貴.
我的神也成為.
我的力量.
然在祂說.
你作我的僕人.
使我各眾之派復興.
使以色列中得.
保存的歸回尚為小事.
我還要使你作.
外邦人的光.
叫你施行我的救恩.
直到地極.
經文在第一節.
就是這樣說.
縱海島.
當聽我言.
遠方的眾民留心耳聽.
自我出差.
耶和華就選召我.
自出武福.
他就提我的名.
在解釋的時候.
我相信我們第一樣要搞清楚.
就是在開頭那裡.
說到那個縱海島.
縱海島.
不僅是以色列.
還包括所有.
外邦的國家.
當中的人在其中.
這些人都要留心.

$^{241}$來到去聽.
那個我說話.
這個我究竟是.
誰來的呢.
第一節下半節.
就說自我出差.
耶和華就選召我.
自出武福.
他就提我的名.
你想清楚.
因為他是.
耶和華所選召的.
所以不會是.
耶和華選召.
耶和華的.
不合邏輯.
所以呢.
這個我就不是耶和華.
如果不是.
耶和華.
你就再想.
那不就是主耶穌.
我相信大家聽了這麼久.
的牧者去說.
這個僕人是誰.
由年頭開始.
我們都應該.
掌握得到.
這個所謂的僕人.
其實未必.
和主耶穌畫上等號.
又不是耶和華.
又不是主耶穌.
究竟是誰.
我們如果.
看清楚去到第三節.
上半節.
我們跳去看的時候.
答案就來了.
對我說你是我的僕人.

$^{281}$以色列.
所以呢.
這個你就是以色列.
即是今天的教會.
即是今天的你和我.
就是那個僕人.
跟著.
自我出差.
出自武福.
又是什麼意思呢.
如果我們看完.
其他聖經的翻譯.
我們就會更加清楚.
呂鎮忠譯本的翻譯.
就是在母胎中.
從我武福女.
這裡是說.
那群以色列人.
今天的你和我.
不只是從出世那天開始.
而是在媽媽的肚子裡.
已經被神揀選.
我們沒有辦法.
改變我們的遺傳基因.
生出來爸爸媽媽帶下來的.
其實就是這樣.
如果我們讀書.
我們回想當年.
如果選科選得不對.
最多用了一點點時間.
就放棄它.
就讀別的科.
但是我們不是.
好像買錯一件東西.
送給別人.
第二天就再買一件.
有機會看中.
不是的.
神對於他所揀選的以色列.
不是這樣.

$^{321}$第二節的上半節說.
他使我的口如快刀.
將我藏在他手刃之下.
又使我成為無量的箭.
將我藏在他箭袋之中.
這裡所說的快.
不是速度.
不是快和慢.
他的字義.
的解釋是鋒利.
所以.
詩高譯本就譯得好一點.
就將他譯作利劍.
無量的意思.
也不是說.
磨到他的箭閃閃亮亮.
重點在磨的字.
打磨一支箭.
我相信.
大家也明白.
只有一個目的.
就是令箭頭.
像剛才所說的那把刀.
能夠尖一點.
鋒利一點.
以致造成更大的殺傷力.
所以英文的聖經.
NRSV就將他譯成.
Sharp Arrow.
一支很尖很鋒利的箭.
在神的眼中.
以色列民.
在媽媽的肚子裡.
其實已經.
這麼厲害.
這是上帝的看法.
如果是這樣.
為什麼接下來.
神又說要將他.
藏在他手刃之下.

$^{361}$然後.
藏在箭袋之中.
這麼厲害.
為什麼要暗器地收起來呢.
我也在台上.
說過不止一次.
我很喜歡看金庸武俠小說.
其中.
他寫的一套.
叫做神雕俠侶.
如果你.
喜歡看金庸.
你應該記得.
裡面有一個角色.
叫做裘千尺.
有沒有印象.
裘千尺.
其實是全身殘廢.
不過他擅長.
將一個.
含在嘴裡.
紅棗或任何一種棗.
的核套出來.
好像暗器一樣.
我圈住.
那裡你也看不到.
這是金庸小說的插圖.
其實那顆核.
是小到.
手指甲那麼大.
不過它的力度.
是強勁到.
令到主角.
楊過.
這裡是這樣描寫的.
感覺手臂劇陣.
五指竟然拿捏不住.
當的一聲.
匕首掉到了地下.
這句話在金庸那裡.

$^{401}$引出來的.
一顆核打過去.
武功高強的人.
拿著一支.
短的匕首都掉到地上.
經文.
這裡不是說.
要把刀和箭.
收藏好.
為了要.
保護他們被人偷走.
不是的.
而是因為刀和箭.
本身都是具有.
殺傷力的武器.
根本不需要被保護.
正確的原因是因為.
他已經準備好.
這一班的以色列民.
往後會被上帝.
大大使用.
但是.
無可否認.
這個塑造的過程.
其實很不容易.
因為當我們繼續看.
第四節一直讀下去的時候.
我們就見到.
筆鋒一轉.
經文是有這樣的描述.
第四節開始的時候.
就說我卻說.
即是說.
就算神.
想到以色列民.
多厲害多好.
怎麼積極.
鼓勵他們.
他們都不同意.
以色列民覺得.

$^{441}$感受不到.
其實.
以賽亞先知.
他的呼召早在.
第六章出現.
記不記得六章一節.
當烏西瓦王.
崩的那年.
我聽見主在.
座在高高的寶座上.
他的衣裳下擺.
遮滿聖殿.
這個很震撼的場面.
出現了一個撒拉弗.
然後拿著.
一個燒紅了的炭.
然後點先知的口.
跟著那句很熟.
六章八節.
我聽見主的聲音說.
我可以猜險誰.
誰肯為我們去.
我說我在這裡.
請猜險我.
多有信心.
為上帝做大事.
豁出去了.
是不是這樣.
我們不要停在第八節.
如果你看完.
第九和第十節.
就精彩了.
他說你去告訴.
這百姓說你們聽了又聽.
卻不明白.
看了又看卻不曉得.
要使這百姓深夢猶知.
耳朵發沉.
眼睛昏花.
恐怕他們眼睛看見.

$^{481}$耳朵聽見.
心裡明白迴轉過來.
就得醫治.
原來神召先知.
去做的大事.
就是煲糯米粥.
即是你做完.
一餐.
其實是浪費力氣.
過了.
四十多章之後.
那裡是第六章.
來到今天.
這裡四十九章.
都是這個樣子.
說了.
四十九章這麼多.
仍然沒有變.
神你說過.
以色列今天的宗教會.
我們每一個.
其實.
都是你的快刀.
剛才說無量的箭.
不過橫看.
碰看都不像.
我們很誠實.
不是這樣.
更普遍的.
感覺就是.
覺得很沮喪.
做來都浪費力氣.
做來做什麼呢.
就好像.
第四節.
上半節那裡說.
我勞碌是徒然.
我盡力是虛無虛空.
由於.
所做的欠缺果效.

$^{521}$結果其實是.
帶來了一個.
很沉重的.
無力感.
在這個情況下.
其實我們會問.
這麼艱難.
做來都浪費力氣.
為什麼我還要堅持下去.
為什麼.
先知的答案.
在第四節的下半節.
他這樣說.
我當得的利.
必在耶和華那裡.
我的賞賜必在我神那裡.
你看到.
兩個必字.
我專登將它寫成.
紅色的.
他認定自己所做的.
無論成功也好.
無論失敗也好.
都不是由.
其他人.
甚至不是由自己.
去判斷.
只是由神去判斷.
並且去到第五節.
上半節又說.
耶和華說.
從我出胎.
造就我作他的僕人.
這麼相似.
第一節說過.
他再次認定.
他自己是被揀選的僕人.
他記得.
神應許過他.
要使雅各歸向他.

$^{561}$使以色列到他那裡.
聚集.
神說.
要將這些以色列民.
再次焦聚回來.
一起回家.
縱然.
目前的境況.
很惡劣.
巴比倫在門口.
但是神的僕人.
有一個很深的體會.
原來.
耶和華看我為尊貴.
我的神也成為.
我的力量.
他知道.
神怎樣看他.
這裡說.
是尊貴的.
就是這一份和上帝的關係.
賜給他.
一個很沉穩的力量.
行了.
就能夠有力量.
去面對生命的挑戰.
回到.
我們今天的處境.
我知道大家覺得.
你今天的生活過得怎樣.
不過不失.
好像我開始的時候.
所說.
昔日的以色列民.
很艱難.
兩個多星期前.
我收到一個消息.
知道我認識的.
當我讀神學的時候.
有一位女同學.

$^{601}$因為膽管癌.
擴散了.
去了舒緩治療的病房.
於是.
我們幾位中神的同學.
約在一起去探望她.
她讀的科目是.
基督教研究碩士.
MCS.
和我讀的科目不同.
她的目標.
不是牧會.
不是做牧者.
因為.
這個學位的課程.
不是做牧者.
所以讀的比較短.
會比我早一年畢業.
因此.
我不算和她.
很熟悉的一位女同學.
但我隱約記得.
在我一年班的時候.
神學院一年班.
她好像和我.
同一個靈命.
培育小組.
我也有個印象.
由於她回天家的日子.
應該不遠了.
所以我也和一班同學.
一起去探望她.
那天.
去到病房.
除了分享了很多.
很美麗.
當時讀神學的時候.
一些回憶之外.
我還聽到一件很感恩的事.
這位生病的女同學.

$^{641}$其實.
她素來在教會.
都是做青少年工作的.
導師.
已經很久了.
那天我們去探望她的時候.
我其中一個.
中神的同班同學.
其實她現在.
是中神的教務長.
她說.
在她們老師的祈禱會裡.
聽到一個年輕老師.
分享了.
她說.
她自己有一位.
昔日帶領著她.
成長的青少年導師.
患上癌症.
情況相當嚴峻.
於是.
請老師們.
在他們的祈禱會裡.
為她昔日.
這位導師.
來祈禱.
這個世界很小.
原來這個青少年的導師.
就是那天.
我們去探望她.
睡在病床上.
和我同屆的.
中神的女同學.
當天.
這位患上癌症的.
女同學.
昔日.
悉心帶領成長的.
這位青少年.
今天.

$^{681}$成為了神學院的.
老師.
我很感動.
因為神使用了.
這位女同學.
比我年輕了一截.
她不是牧者.
我剛才說.
她讀的課程不是牧會.
但神仍然.
很使用她.
能夠.
做好.
一個很美麗.
模樣的.
果校.
我想說.
神怎樣用人.
其實是遠遠.
超過我們的想像.
你想也想不到.
這位女同學.
因為她不是牧者.
她讀完.
基督教研究碩士之後.
回到自己的母會.
做行政幹事.
一做就.
做了十五年.
然後.
因為她聽聞.
知道有另外一個.
教會.
很需要人.
行政的角色.
沒有人可以承擔.
於是.
她就拍心口.
去頂住三個月.
誰知一做就做了.

$^{721}$五年.
三個月變成五年.
當日.
我們幾個中神同學.
去探望她的時候.
她說她自己預備好了.
等上帝.
要接她.
她很多年.
一直來到忠心.
去服侍.
她當日.
雖然睡在病床.
很病.
不過她的眼睛.
所散發出來的光彩.
和二十多年前.
我認識她的時候.
一樣.
你從她的眼睛.
你是完全看不到.
一個人可以病成這樣.
但是.
仍然散發出.
那種光彩.
不過.
今早我收到一個訊息.
她回天家了.
我寫這篇講章的時候.
還沒有.
不過今天我和大家講的時候.
她走了.
昔日的以色列民.
面對的處境.
很艱難.
今日.
我們.
是否只是.
當我們的癌症.
得著醫治的時候.

$^{761}$我們才.
可以為神作見證.
事實.
告訴我們.
癌症得醫治.
總有下次.
總會再病.
終於.
有一天.
得醫治.
終有一天.
是不得醫治.
好了.
雖然不得醫治.
但是在過程裡.
我們仍然可以.
數算上帝的恩典.
讓上帝.
得著榮耀.
看看第三節下半節.
怎麼說.
我等神.
你.
祂的子民以色列.
透過祂.
能夠得著榮耀.
神很多時候.
都不是.
在我們得勝.
而是在我們.
很困難的時候.
甚至一仆一碌的時候.
這些經歷.
讓上帝.
得著榮耀.
因為在得勝的時候.
我們很容易.
有意無意.
將榮耀.
歸了給自己.

$^{801}$其實.
教會不是一樣.
求主幫助.
我們能夠認清.
不要只是去求.
單單去求.
事公發展順利.
一個.
兩個.
三個四個的困難.
是會有的.
上帝.
想見到的.
是更多的罪人.
可以悔改.
讓上帝.
這樣.
來得著榮耀.
接著.
神繼續透過.
先知在第六節.
這樣說.
現在他說.
你作我的僕人.
眾知派復興.
使以色列中.
得保存的歸回尚為小事.
我還要.
使你作外邦人的光.
叫你施行我的救恩.
直到地極.
相信這一節經文.
當時.
星日的以色列文.
聽起來.
真的非常之啞耳.
他們面對.
亡國.
自老到巴比倫.
歷盡滄桑.

$^{841}$詩篇也有這樣說.
一想起.
錫安就哭了.
但是神竟然.
告訴他們.
使雅各.
眾知派復興.
使以色列中.
得保存的歸回.
尚為小事.
不是吧.
如果是這樣.
這些是小事.
什麼才是大事.
第六節下半節就說.
我還要.
使你作外邦人的光.
叫你施行我的救恩.
直到地極.
神最終的心意.
是要使用.
以色列文.
他揀選的教會.
揀選的你和我.
每一個人.
是為要.
使到救恩.
可以臨到地極.
那你問.
地極.
即是什麼.
保羅在.
使徒幸傳十三章.
四十七節.
曾經引用過.
這一段經文.
直接引用.
我用呂鎮忠.
的翻譯.
給大家看.

$^{881}$看清楚一點.
因為主曾這樣.
囑咐我們說.
我立了你作外國人的光.
使你作救星.
直到地極.
看不看到.
雙引號裡面.
藍色那兩句.
是平行句.
平行句是什麼意思.
就是上面和下面的意思.
一樣.
地極.
地域.
去到那麼遠.
北極南極.
當時也不知道.
不是.
而是等於.
外國人.
即是和合本所說的.
外邦人.
有一個近代的美國神學家.
叫做雷安德森.
Ray Anderson.
他很嚴正地提醒.
我們教會.
最重要的.
服事和使命.
是什麼.
他說.
耶穌基督的事工.
是植根於天父的旨意.
而不是由世界的需要.
或者要求.
所定義.
或者作為起點.
英文那段就在下面.
打出來給你.

$^{921}$by the needs or demands.
of the world.
這一位神學家.
用了.
耶穌.
使拉撒路復活.
作為一個例子.
如果.
耶穌按照.
人的想法.
去做的話.
他聽見拉撒路病成這樣.
好了.
會怎樣做.
去治療.
很困難.
是不是這樣.
其實.
神兜口兜面.
要跟.
這些伊斯利文說.
雖然.
你們是.
在經歷巴比倫的.
郵輪.
很期待.
復興.
希望可以歸會.
不過.
這些尚為小事.
這個.
不是焦點.
也不是.
上帝的心意.
所在.
所以.
很值得我們想一想.
究竟我們的.
生命焦點在哪裡.
究竟你的生命.

$^{961}$焦點在哪裡.
因為.
我年紀大.
所以.
我應該退了.
留給年輕人.
我聽到太多.
我正在.
面對一千件事.
搞不定.
神.
你不要來搞我.
哪裡搞得定.
應付不來.
但是神說.
這些尚為.
小事.
Sorry.
不是.
原因.
不是藉口.
我們看得到.
上帝的心意.
又願意.
放下自己的喜好.
覺得那些東西很重要.
甚至是困難.
去校正我們的焦點.
去作外邦人.
的光.
和我的救恩.
直到地極.
我很感恩.
上帝把你放在官宣.
就算你低頭.
上帝也會把你放在官宣.
上帝.
把你放在官宣.
其實.
接下來.

$^{1001}$幾個星期.
我們有.
六月份的會員大會.
在當中.
我們眾位.
同工和執事.
會跟大家分享.
神在福音的.
向導.
期盼.
怎樣用官宣.
我需要.
大家去配合.
你又願意.
參與.
和回應.
上帝給教會.
這個福音機遇嗎.
船造出來.
不是用來.
停泊在碼頭.
車造出來.
也不是放在停車場.
當神.
去造你和我.
呼召我們.
作為門徒.
其實也是一樣.
不是只坐在牆邊.
而是.
請你下牆.
每一個人.
都下牆.
當我們.
願意下牆.
參與.
我們就能夠一起.
彰顯上帝的榮耀.
即使是一場球賽.
我告訴你.

$^{1041}$可能你只是.
參與這場球賽.
第一節的比賽.
接著教練.
就按照他的策略.
將你換走了.
你就出來了.
下一個日.
球賽還沒結束.
但是不代表.
你下場.
那段時間.
是沒有意義的.
因為我們已經.
配合了上帝.
要我們.
作為他的僕人.
作為他的大使.
這是其中一部分.
最後我想借鏡.
伊斯蘭信仰的.
發展.
來幫助我們.
做一個反思.
可能大家都會知道.
以國家來計算.
印尼是.
世界上穆斯林.
人口最多的國家.
以國家來計算.
大約有.
2300萬.
對不起.
2300萬人.
信奉伊斯蘭教.
佔了.
他們整體的人口.
比例87\%.
都很高.
我不知道大家有沒有想過.

$^{1081}$為什麼伊斯蘭信仰.
可以在印尼.
發展得這麼蓬勃.
有沒有想過.
可能你的答案.
是.
因為他們戰鬥格.
好打.
又或者.
他們信奉伊斯蘭教.
的男人.
一個娶四個老婆.
於是生小朋友.
都比你多.
我想告訴你.
其實.
印尼人.
信奉伊斯蘭教.
其實是和軍事.
征服又或者.
國家的強權.
法令.
甚至有很多.
信奉伊斯蘭的宣教士.
去那裡傳教.
其實.
完全不關事.
印尼原本.
是一個印度教.
佛教和.
萬靈主義.
為主的.
混合的國家.
原來這個改變.
是因為.
貿易的路線.
是因為文化的融合.
由主後.
七世紀已經開始.
所以.

$^{1121}$其實是經歷了.
幾個世紀的轉化.
幾百年.
伊斯蘭教最初.
是經海路.
由阿拉伯 由波斯.
透過商人.
來傳入.
由於當時.
走陸路.
經商.
有很多賊.
盜賊猖獗.
所以.
經海路過去.
是比較安全的路線.
所以.
做生意的人.
會選擇坐船.
不過當然.
有利也有弊.
坐船的問題是什麼.
被天氣影響.
其實.
東南亞.
有季候風.
大家都知道.
這個風向.
基本上是半年一次.
所以.
如果風向不對.
他們去到印尼.
沒辦法回家.
那怎麼辦.
坐在那裡等.
在他們不同的島嶼.
生活.
等到風向轉.
對了.
坐船回家.

$^{1161}$在他們逗留期間.
他們不單是做生意.
他們更加會.
跟當地人做朋友.
甚至.
有些人開始.
跟當地人通婚.
當地人看到.
這些穆斯林.
信仰的朋友.
怎樣默默實踐.
他們的信仰.
就對他們的.
伊斯蘭信仰產生了興趣.
融入了他們的文化.
結果就越來越多.
印尼島嶼的.
當地人.
信奉伊斯蘭教.
是這樣的.
這個過程.
很長.
也很不容易.
剛才說了幾百年.
即是你幾百年.
都要做得這麼好.
老友.
就好像一開始的時候.
我說到那三位.
二十幾歲.
三十幾歲 四十幾歲的媽媽一樣.
其實.
生命成長.
見證.
其實是要時間的.
真金不怕.
紅爐火.
過程很慢的.
不過在這個過程裡.
我們要認定.

$^{1201}$兩樣東西.
第一 上帝很厲害.
很信實.
第二 我們很不行.
我們很需要福音.
你想一下.
上帝拯救我.
拯救我.
拯救你的時候.
不是按一個按鈕.
是經歷了.
一段的過程.
這個正正.
就是上帝.
拯救的方法.
作為一個.
簡單總結.
我想和大家.
分享一位.
丹麥哲學家.
他叫琪克果.
他的說話很正經.
他這樣說.
如果用中文翻譯.
基督教的真理.
本身就有一雙眼去看.
事實上它就是一雙眼.
想像一下.
當你去.
看一幅畫的時候.
如果你發覺.
那幅畫的人像.
他的眼睛都看著你.
你又看著他.
他又看著你.
其實不是很舒服.
神的道.
是很尖銳的.
我想告訴你.
就是有這個銳利的目光.

$^{1241}$來看著你的靈魂.
你另一面也沒有用.
他也是看著你的.
所以我們.
更加要.
認真去聽.
我們不能將.
他的說話聽完.
覺得不錯.
然後走出教會門口.
就放下.
置之不理.
神的話.
我們是必須要回應的.
有一天.
有一位男士來找牧師.
他和牧師說.
我想了解一下.
基督教信仰的要點.
簡單說給我聽.
事源.
因為他太太.
剛剛信了耶穌.
他希望到馬克.
和太太吃飯的時候.
聊天可以搭上嘴.
不然太太說話.
不知道會發生什麼事.
他告訴牧師.
所以.
實在沒有時間.
去了解.
看清楚.
他太太的信仰.
究竟是怎樣.
於是.
牧師就回答他.
我看得出你真的很忙.
所以我覺得這樣做.
其實不太適合你.

$^{1281}$這位男士.
他有點尷尬.
他問牧師.
因為.
如果我將.
基督信仰的重點.
來跟你說.
你很聰明.
你又聽得明白.
神的真理.
一定會深深.
影響你的生活.
甚至會徹底.
反轉.
打亂你的生活.
你會.
需要重新.
編配你的時間.
重新釐定.
你的工作目標.
甚至是人生意義.
我猜你不太想這樣做.
這位牧師.
跟他說.
我不知道.
大家會怎樣回應.
我今天.
透過這六節.
僕人之歌.
問你的問題.
你還記得.
昔日.
神揀選了你.
成為鋒利的刀.
和箭嗎.
無論你的生命.
面對著.
什麼挑戰.
祂說要透過.
這些經歷來塑造你.

$^{1321}$目的是讓.
上帝得著榮耀.
還有.
要求我們去.
校正我們生命的.
焦點.
從我們的小事.
回到.
上帝福音的托付.
神交托.
給官孫.
和我們每一個人.
這個福音托付.
還未完.
要解那個題目.
這個就是.
未竟全功.
的意思.
未做完.
我在後面.
專登加上一個.
文畫符號.
因為我想.
代上帝問下你.
問下你們.
每一個.
究竟你.
會怎樣回應.
\newpage



\section{以弗所書 2:1-10}
\label{sec:032108}
\textbf{超越七大奇蹟的奇蹟}~~~~禤凱榮牧師
\newline
\newline
連結: \href{http://www.church.com.hk/acms/content.asp?site=cdc\&op=show\&type=product\&code=032108\&layout=sermon}{\texttt{032108}} ~~~~ 語音日期: 2026-05-10
\newline
\newline
\hyperref[sec:032107]{\small{< < < PREV SERMON < < <}}
~
\hyperref[sec:index_chronic]{\small{[返順時目]}}
~
\hyperref[sec:032109]{\small{> > > NEXT SERMON > > >}}
\newline
\newline
以弗所書 2:1-10
\newline
\begin{longtable}{cl}
\hline
\hline
章節 & 經文 (和合本修訂版)\\
\hline
2:1 & \begin{tabularx}{0.7\textwidth}{X} 從前，你們因著自己的過犯罪惡而死了。 \end{tabularx} \\ \\ \relax
2:2 & \begin{tabularx}{0.7\textwidth}{X} 那時，你們在過犯罪惡中生活，隨從今世的風俗，順服空中掌權者的領袖，就是現今在悖逆的人心中運行的邪靈。 \end{tabularx} \\ \\ \relax
2:3 & \begin{tabularx}{0.7\textwidth}{X} 我們從前也都生活在他們當中，放縱肉體的私慾，隨著肉體和心中的意念去做，和別人一樣，生來就是該受懲罰的人。 \end{tabularx} \\ \\ \relax
2:4 & \begin{tabularx}{0.7\textwidth}{X} 然而，神有豐富的憐憫，因著他愛我們的大愛， \end{tabularx} \\ \\ \relax
2:5 & \begin{tabularx}{0.7\textwidth}{X} 竟在我們因過犯而死了的時候，使我們與基督一同活過來—可見你們得救是本乎恩— \end{tabularx} \\ \\ \relax
2:6 & \begin{tabularx}{0.7\textwidth}{X} 他又使我們在基督耶穌裡與他一同復活，一同坐在天上， \end{tabularx} \\ \\ \relax
2:7 & \begin{tabularx}{0.7\textwidth}{X} 為要把他極豐富的恩典，就是他在基督耶穌裡向我們所施的恩慈，顯明給後來的世代。 \end{tabularx} \\ \\ \relax
2:8 & \begin{tabularx}{0.7\textwidth}{X} 你們得救是本乎恩，也因著信；這並不是出於自己，而是神所賜的； \end{tabularx} \\ \\ \relax
2:9 & \begin{tabularx}{0.7\textwidth}{X} 也不是出於行為，免得有人自誇。 \end{tabularx} \\ \\ \relax
2:10 & \begin{tabularx}{0.7\textwidth}{X} 我們是他所造之物，在基督耶穌裡創造的，為要使我們行善，就是神早已預備好要我們做的。 \end{tabularx} \\ \\
[1ex]
\hline
\hline
\end{longtable}
$^{1}$小時候在學校已經學到.
這個世界七大奇蹟.
七大奇蹟.
當然了.
大家知道我記性差.
很快就忘記了.
我就上網查查看.
發覺原來.
現在的七大奇蹟.
和我小時候學的七大奇蹟已經不同了.
不知道大家知不知道.
現在版本的七大奇蹟.
是哪七個呢?.
或者讓我們很快的看看.
當然有我們祖國的萬里長城.
有約旦的佩特拉古城.
有巴西的里約熱內盧的基督像.
有被擄的馬丘比丘.
有墨西哥的金字塔.
有意大利的荊棘牆.
還有印度的泰基瑪哈尼.
不知道大家去過多少個呢?.
我只去過萬里長城.
別人說北上長城非好看.
我也做過兩次好看的.
不過今天我不是說七大奇蹟.
我想跟大家說的就是.
比這七大奇蹟裡面.
更奇蹟的一個神蹟.
一個奇蹟.
而這個奇蹟是不會受到任何時間的改變的.
你說這是什麼奇蹟呢?.
如果剛才有留心聽.
我們所讀的經文的時候.
這個奇蹟.
其實就是你和我.
嘩!為什麼這樣說?這麼厲害?.
讓我們通過剛才讀的經文.
看看究竟這一段的經文.
教導的我們是什麼.

$^{41}$我想大家想一想.
我們決定一樣的東西.
是不是一個奇蹟呢?.
有什麼條件呢?.
或者我們需要考慮什麼呢?.
我相信很重要的.
我們看看究竟這個東西所涉及的人.
當時的科技.
資源.
很多東西都是考慮的.
我舉個例子.
如果我跟大家說.
大家有沒有看到三層高的.
這麼漂亮的豪宅呢?.
如果我跟大家說.
這是一個很厲害的建築師.
一個人做出來的.
我們會怎樣呢?.
我們就會讚嘆這個建築師的辛勞.
他的技巧.
他的手工.
不過我們不會覺得這是一個神蹟.
這是一個奇蹟.
是不是?.
我們不會覺得.
不過如果我跟大家說.
建築這個三層樓高的.
是一個十歲大的小孩呢?.
這樣就不同了.
沒有可能的.
一個十歲大的小孩.
怎麼可以做出這樣的建築物呢?.
不過大家有沒有想過.
正正是這個不可能.
就是決定那個東西是不是一個奇蹟.
是不是一個神蹟?.
所以當我們去檢討去思考.
神在人類裡面.
所行過的神蹟奇事的時候.
我首先一個很重要要考慮的因素是什麼?.

$^{81}$就是當時在神未介入人類歷史之前.
人類的光景.
是什麼呢?.
這是一個很重要我們需要考慮的.
我們看看經文.
經文一開始的時候.
保羅告訴我們.
你們都死在過犯罪惡之中.
他叫你們活過來.
這個他當然是指到神.
那麼你們呢?.
你們是指到誰呢?.
很明顯是指到當時寫信給以弗所書的讀者.
是什麼呢?.
就是說在以弗所教會裡面的信徒.
是重生得救的基督徒.
在這群人未成為基督徒之前.
保羅就說.
那個時候.
你們隨從今世的風俗.
信服空中掌權者的首領.
放縱肉體的私慾.
隨從肉體和心中所喜好的去行.
如果大家留意短短這兩節的經文.
保羅是形容給我們看.
一個沒有基督的人生.
其實是被三種很厲害的力量所困鎖住的.
我們一起看看這三個力量是什麼力量.
第一個是世界.
保羅說你們在其中行事為人.
隨從今世的風俗.
很多時候大家要想一想.
人總是在某一個時空底下.
受著當時當刻的處境.
當時的文化.
你和什麼人在一起.
當時社會的意識形態等等.
所有這些事情是影響著的.
所有剛才我們所說的東西.
其實都是形成了一些外在的壓力.

$^{121}$去讓人認同同化當時世代的價值觀.
舉個例子.
今天如果有人對你很沒禮貌.
有人對你不給你面子.
不給你面子,不尊重你.
你會怎樣?.
中國人最喜歡說你對我怎樣我就對你怎樣.
你錯了我十五.
當然要報復.
曾經有一個剛剛讀完神學的牧師出來.
去到一個新的教會.
去牧羊的教會.
去到的時候很熱心.
就想著不如我們.
有些未信的人在當中.
不如我們搞一個報道會.
就想著崇拜的題目是我定的.
我的責任範圍.
所以做了決定而沒有和當時的執事們說.
不過他這樣做的時候.
就令當中的一個執事很不滿意.
你這樣真的不給我面子.
這樣的時候就很不開心.
我想問大家當你不開心的時候.
如果按照聖經的原則應該怎樣處理.
當然要跟聖經說.
我們要私下跟他說.
跟他說的時候不要用對錯的立場.
首先要承認他是一個很有熱心傳福音的人.
也要提醒他對公的重要性.
為了要讓事情得到和解.
不過如果不是按照聖經的原則.
而是按照世界的原則.
會怎樣處理?.
你不給我面子.
我不看你台戲.
這個受到所謂傷害的執事.
最後就沒有推廣這個福音主人.
當日也沒有出現.
不單止沒有出現.

$^{161}$還要叫其他人不要出席.
所以當然了.
這件事情搞得很不開心.
丙姐妹.
很重要的.
我們千萬不要隨從今世的風俗.
我們按照聖經的原則做事.
第二.
保羅告訴我們.
我們都要順從空中的掌權者.
誰是空中的掌權者?.
聖經告訴我們.
很清楚就是邪靈.
魔鬼.
撒旦.
使徒彼得在他的經文裡也說過.
我們的仇敵魔鬼和嚎叫的獅子.
四處遍地遊行.
做什麼?不是觀光.
是找一些可以吞吃的人.
什麼意思?.
就是說看看有哪一個人.
信心沒那麼堅固的時候.
就令他跌倒.
在真道上跌倒.
我想問大家.
你要一個人跌倒.
最有效的方法是什麼?.
當然是攻擊他最弱的地方.
最弱的地方.
我們今天所說的攻擊是短板.
撒旦知道我們每一個人的短板.
你知不知道你自己的短板是什麼?.
對於那些可能自我形象比較低.
自信不是很高的人.
撒旦可能會在耳邊跟我說.
你也是沒用的.
神怎麼會理你?.
你永遠都不會改的.
如果你是一些.

$^{201}$很多壞習慣.
改也改不了.
撒旦就會在你耳邊說.
個個都是這樣做的.
做完這次再說吧.
有什麼大不了?.
對於那些.
經常很容易受到私貪.
撒旦可能在你耳邊說.
多看一眼吧.
做完這次再說吧.
沒有人知道的.
大家有沒有這種經驗?.
沒有,我沒有這種經驗.
如果你真的沒有的話.
大概你不是被人強.
大概你是患了屬靈的弱聽症.
聽不到東西.
所以大概你需要什麼呢?.
你需要安一個屬靈的助聽器.
所以大家留心.
彼得給過我們的提醒.
撒旦真的是遍地遊行的獅子.
看看誰可以一口吞了他.
他快要吞了我們.
第三個.
保羅告訴我們.
我們都是放縱肉體的私慾.
中國人說人總有七情六慾.
我相信沒有人反對.
神是讓人有身體上.
情緒上的需要.
這是人性.
神從來沒有否認這些人的需要.
不過神提醒我們.
我們千萬不要放縱這些情慾.
身體上的需要.
免得反過來這些需要.
控制我們.
所有我們今天所認知的上癮.

$^{241}$有什麼上癮?.
毒品上癮.
酒精會上癮.
賭博都會上癮.
性會上癮.
種種都是當人放縱他的私慾的時候.
就會發生的事情.
過去舉個例子.
過去我們對一些吸毒的人.
上癮的人.
我們怎樣來稱呼他們?.
我們香港叫他們導友.
我們不是排斥這些人.
是貶低這些人.
我們是在社會的意識形態裡面.
我們這樣認定這些人.
不過今天.
今天我們所吸的毒不是像以前那樣.
聽說現在那些是冰毒.
那些毒的毒性是厲害很多的.
今天叫那些人是什麼人?.
不是叫導友.
叫喪屍.
很多的國家的城市.
都有這些喪屍出現.
保羅告訴我們.
一個很可怕的事實.
你們豈不是跟這些人.
以前就是跟那些人一樣.
豈不是你們都是和他們和別人一樣.
這個別人是指什麼?.
就是指那些沒有耶穌在生命裡面.
對生命失去了應有的掌控.
為什麼會這樣?.
保羅一針見血就說出來了.
為什麼?.
因為我們都死在過犯罪惡之中.
當然大家要留意.
保羅這裡所說的死.
並不是指身體上的死亡.

$^{281}$我們今天個個都活著.
我們沒有死.
保羅這裡所說的是靈性上的死.
什麼叫靈性上的死?.
聖經說得很清楚.
靈性的死是指當人和神分開.
被阻隔的時候.
這個就是靈性上的死.
神是生命的源頭.
一個人可以生龍活虎.
成功一族.
但是他仍然可以在靈性上.
是無知無覺.
一無所知.
對神是毫無反應的.
人為什麼會和神隔絕呢?.
先知以賽亞曾經這樣說過.
但你們的罪孽使你們與神隔絕.
你們的罪惡使神掩面不聽你們.
當我們和神分開的時候.
我們就和神完全失去了那個關係.
我們不會有任何的興趣.
去關心一些屬靈的事情.
所以你不要和一些沒有神在心中的人.
去覺得他們會祈禱.
或者會敬拜.
或者是追尋真理.
他們沒有這些需要.
當人失去了和神的關聯之後.
我們每一個人都活在一個.
聖經說是被罪所控制的世界.
既然是被罪所控制.
而我們活在這個世界裡的時候.
當然人自然地.
就會被罪所控制.
這個就是.
當我們死在過犯罪惡之中的時候.
人的光景就是這樣.
不過感恩我.
感恩就是.

$^{321}$這個並不是人類歷史的終結.
如果我們看第四節的時候.
第四節.
這個保羅說.
然而.
這個然而是一個很寶貴的連接詞.
因為神並沒有容許到死亡.
成為人類歷史的句號.
並不是.
保羅說.
由於神的憐憫與愛.
他猜險他的獨生子.
進入人類的歷史.
向罪.
向死亡.
來宣戰.
有一個行動.
這個行動就是他的獨生兒子耶穌基督.
替罪人死在十字架上面.
為什麼要死在十字架上面.
以至他可以流出來的血.
洗淨人所犯過的罪.
以至人可以在罪惡裡面被釋放出來.
這個.
就是基督徒傳了二千年的福音.
保羅告訴我們.
當我們面對福音.
接受福音的時候.
會有三個很清楚的結果.
就是與基督一同活過來.
與基督一同復活.
與基督一同坐在天上.
今天讓我們一起看看.
這三個領受了福音之後的結果.
第一個.
與基督一同活過來.
當我們死在.
我們的罪惡過犯當中的時候.
由於神的恩典.
和我們在耶穌基督裡面的信心.

$^{361}$這裡保羅就是說.
我們就會被赦免.
以至我們不再與神隔絕.
意思就是說我們靈命得到復甦.
我們能夠與神恢復溝通.
雖然我們肉眼看不到這個改變.
但是這是屬靈的事實.
所以保羅在羅馬書的六章裡.
說到我們.
信的人向罪.
應當看自己是死的.
向神在基督耶穌裡.
我們應該看自己是活的.
用一個例子.
好像一個燈掣.
電一直都在那裡.
你不開燈就沒電了.
你開燈就有電了.
就是這麼簡單.
意思就是說耶穌基督.
我們在耶穌基督裡面的信心.
其實就好像開這個電掣一樣.
按著的時候.
救恩就進入我們的生命.
讓我們的靈魂得著甦醒.
讓神人之間的阻隔.
可以被除去.
這個就是與基督一同活過來的意思.
第二個.
第二個就是與耶穌基督一同復活.
不單止神.
潔淨了我們的罪.
讓我們重新得著屬靈的生命.
他們也說.
叫我們與基督耶穌一同復活.
一同再起來.
什麼意思.
保羅在羅馬書的六章裡.
也有一個較詳細的解釋.
保羅在羅馬書這樣說.

$^{401}$他說我們藉著進入耶穌的死.
和祂一同埋葬.
其實就是叫我們一舉一動.
有新生的樣式.
以致我們可以與祂同活.
就好像從死裡復活的人.
將自己獻給神.
並將肢體作義的器具獻給神.
換言之.
神不單止賜我們新的屬靈生命.
更加重要的.
給我們一個新的生命的方向.
一個新的生活的方式.
換言之.
神不單止讓我們與主同活.
也讓我們有足夠的能力與恩典.
為主而活.
一個新的方向.
新的生活的內容.
第三.
與基督一同坐在天上.
保羅進一步告訴我們.
不單止給我們新的生命.
給我們新的生命的方向.
現在更加告訴我們.
我們與基督一同坐在天上.
什麼意思?.
你會說牧師坐在天上?.
不是吧?.
我相信除非你現在是混淆障礙.
夢鄉正甜.
否則我們都會坐在禮堂.
怎麼會與耶穌坐在天上呢?.
很明顯這句話並不是字面的解釋.
如果我們看回羅馬書的六章.
那裡有更深的解釋.
「我們若在祂死的形狀上與祂聯合.
也要在祂復活的形狀上與祂聯合」.
什麼意思呢?.
大家記不記得.

$^{441}$在耶穌基督復活之後.
祂向門徒顯現了幾次.
在最後一次要跟門徒說的話.
在馬可芬的記載.
當耶穌跟門徒說完要說的話之後.
被接升到天上.
然後坐在神的右邊.
由於我們與耶穌在復活上與祂聯合.
我們今天屬靈的意義上.
我們都是與祂一同坐在天上.
這件事很空泛和抽象.
是什麼意思呢?.
雖然我們基督徒今天仍然坐在地上.
我們今天不是身處天堂.
但我們要明白一件事.
我們屬靈的地位已經在神的國度裡.
我們的身份有一個很基本的改變.
我們不再屬於地上.
我們屬於天上.
簡而言之我們雖然活在地上.
但我們屬於天上.
有天國子民的身份.
神也給過我們一個身份.
對於一個這麼大的恩典.
我們應該怎樣作出回應呢?.
保羅在這裡又給了我們三種.
人應該有的回應.
我們現在跟大家看看.
第一個回應.
我們需要接受神的恩典.
「彭遜接受」.
如果我們看八節到九節.
救恩是神給人的一個禮物.
那裡說到你們得救是本乎恩.
也是因著信.
所以我們一定要搞清楚.
我們信耶穌.
我們不是靠著我們一生以來所做的好事.
我們經常說.
我們要行善積德.

$^{481}$來得到神明的眷由.
對不起.
基督教裡面不是在說這件事.
基督教說.
聖經保羅已經說得很清楚.
我們得救是本乎恩.
也是因著信.
我們能夠在神面前得到接納.
完全因為神的恩典.
容讓我們將我們的信心.
投放在耶穌基督裡面.
以致我們可以得到罪的赦免.
得到神寶貴的禮物.
想像一下.
如果我今天有一份禮物.
剛才大家進來的時候也有.
希望大家如果還沒有拿到母親節的禮物.
等一下記得臨走的時候拿.
大家想像一下.
如果我要給你一個禮物.
你雙手都拿著東西的時候.
拿著自己所堅持的東西.
你拿不到這個禮物.
拿不到的.
要怎麼樣.
你要將你手上的東西放下.
然後才可以接受新的禮物.
所以當我們仍然持守著.
一些對將來.
這個永生.
一些執著的時候.
我們沒有辦法領受.
這個神所為我們預備的恩典.
所以我們第一樣.
我們需要憑信接受.
第二樣是信靠信服.
信靠信服.
聖經告訴我們.
罪為什麼會出現在這個世界.
聖經告訴我們.

$^{521}$當我人類的始祖犯罪.
不信服神的時候.
罪因此而進入世界.
影響控制這個世界.
所以才知道.
撒慕爾是這麼說的.
他說.
神看人的信服多過人所給過他的獻祭.
這是一個很出名的說話.
聽命勝於獻祭.
順從勝於公羊的自由.
無論你獻多少牛羊.
都及不上你在神面前的信服.
在信服裡面.
神很多時候必然會做一些神所喜悅的事情.
這裡就是保羅所說.
當我們在神的救恩裡面的時候.
我們每一個人就是神的工作.
我們每一個人.
神都盼望我們得著這個恩典之後.
叫我們行善.
所以大家看到分別了.
我們行善不是為了得著神的救恩.
是因為我們得著神的救恩之後.
我們有能力.
有動力去行善.
這個行善大家要搞清楚.
這個行善並不是指社會上公義的事情.
做公益就可以得到神的眷有.
不是這個意思.
這裡所說的行善.
是指行善捍衛好的事情.
所以大家要很小心.
我們不是靠善行得救.
但是我們得救以後.
我們需要靠著神的恩典與能力.
行出善行.
第三.
我們需要傳揚福音.
這裡說.

$^{561}$第三個就是.
神在第七節那裡這樣說.
要將他極豐富的恩典.
就是他在基督耶穌裡面.
向我們所施的恩慈顯明給後來的世代.
意思就是說.
神一直期盼著他的恩典.
他的救恩.
他所為我們預備一切的美善.
賜予給每一個人.
無論你是在什麼時代.
無論你是什麼國家.
什麼民族.
每一個人都歡迎來到神面前.
領受這個神為我們預備的救恩.
故此.
在耶穌升天之前.
聖經有這樣的記載.
祝福他的門徒說.
將來聖靈會降臨到你們的身上.
你們就必得著能力.
在不同的地方直到地極.
一路傳傳傳開來.
作我的見證.
這個就是.
教會傳了二千年的福音.
弟兄姊妹.
我不知道大家怎麼認同.
對我來說.
救恩是整個世界裡面.
超越每一個.
這個世界我們所定義的奇蹟與神蹟.
而這個神蹟.
你和我.
只要是真心相信的基督徒.
就是這個神蹟的例子.
我們是顯示上帝.
無比豐盛恩典的見證.
不過留意.
我們在整個救恩裡面.

$^{601}$我們不是主角.
我們只是配角.
神才是整個救恩的主角.
如果我們用一個藝術作品.
來思考的時候.
我們只不過是一個畫布.
神才是將美麗的藝術.
畫上去的藝術家.
如果我們用一個獎杯來形容的時候.
我們是一個獎杯.
神才是真正得獎的冠軍.
一切都是在於神.
本於神.
所以弟兄姊妹.
願意我們都能夠更加明白聖經.
讓我們都能夠用正確.
應當的方式.
回應神給我們的恩典.
如果大家仍然是.
未認識耶穌基督的時候.
盼望你有機會.
或者等下你和我.
來詢問信耶穌的事情.
如果大家已經是基督徒的時候.
我們都希望大家.
能夠繼續與主聯合.
服侍祂.
傳揚祂.
和跟隨祂.
請我們低頭祈禱.
因為使徒寶萊告訴我們.
是你的恩典.
讓人可以通過信耶穌.
得著救恩.
可以得著永遠的生命.
可以得著你在天上.
為我們存留的基業.
如果我們在座當中.
仍然有未認識你.
未相信你.

$^{641}$未投靠你的朋友.
求你幫助.
讓他們能夠更深的認識你.
以至他們能夠將心交你.
成為一個滿有你豐盛連運的信徒.
一生跟隨你.
為我們已經信你的人.
求你幫助.
使我們更加的每天跟隨你.
親近你.
多謝你給我們一切的祈禱.
奉主名求.
阿們.
\newpage



\section{彼得前書 3:13-22}
\label{sec:032109}
\textbf{耶穌不在的時候}~~~~林豪恩先生
\newline
\newline
連結: \href{http://www.church.com.hk/acms/content.asp?site=cdc\&op=show\&type=product\&code=032109\&layout=sermon}{\texttt{032109}} ~~~~ 語音日期: 2026-05-10
\newline
\newline
\hyperref[sec:032108]{\small{< < < PREV SERMON < < <}}
~
\hyperref[sec:index_chronic]{\small{[返順時目]}}
~
\hyperref[sec:032110]{\small{> > > NEXT SERMON > > >}}
\newline
\newline
彼得前書 3:13-22
\newline
\begin{longtable}{cl}
\hline
\hline
章節 & 經文 (和合本修訂版)\\
\hline
3:13 & \begin{tabularx}{0.7\textwidth}{X} 你們若熱心行善，有誰會害你們呢？ \end{tabularx} \\ \\ \relax
3:14 & \begin{tabularx}{0.7\textwidth}{X} 即使你們為義受苦，也是有福的。不要怕人的威嚇，也不要驚慌； \end{tabularx} \\ \\ \relax
3:15 & \begin{tabularx}{0.7\textwidth}{X} 只要心裡奉主基督為聖，尊他為主。有人問你們心中盼望的理由，要隨時準備答覆； \end{tabularx} \\ \\ \relax
3:16 & \begin{tabularx}{0.7\textwidth}{X} 不過，要以溫柔、敬畏的心回答。要存無虧的良心，使你們在何事上被毀謗，就在何事上使那些凌辱你們在基督裡有好品行的人自覺羞愧。 \end{tabularx} \\ \\ \relax
3:17 & \begin{tabularx}{0.7\textwidth}{X} 神的旨意若是要你們因行善受苦，這總比因行惡受苦好。 \end{tabularx} \\ \\ \relax
3:18 & \begin{tabularx}{0.7\textwidth}{X} 因為基督也曾一次為罪受苦 ，就是義的代替不義的，為要引領你們到神面前。在肉體裡，他被治死；但在靈裡，他復活了。 \end{tabularx} \\ \\ \relax
3:19 & \begin{tabularx}{0.7\textwidth}{X} 他藉這靈也曾去向那些在監獄裡的靈傳道， \end{tabularx} \\ \\ \relax
3:20 & \begin{tabularx}{0.7\textwidth}{X} 就是那些從前在挪亞預備方舟、神容忍等待的時候不信從的人。當時進入方舟，藉著水得救的不多，只有八個人。 \end{tabularx} \\ \\ \relax
3:21 & \begin{tabularx}{0.7\textwidth}{X} 這水所預表的洗禮，現在藉著耶穌基督的復活拯救你們，不是除掉肉體的污穢，而是向神懇求有無虧的良心。 \end{tabularx} \\ \\ \relax
3:22 & \begin{tabularx}{0.7\textwidth}{X} 耶穌已經到天上去，在神的右邊，眾天使、有權柄的、有權能的都服從了他。 \end{tabularx} \\ \\
[1ex]
\hline
\hline
\end{longtable}
$^{1}$今日是母親節.
讓我們祝願天下母親身心靈健康.
願上主的恩典.
常與每一位母親同在.
說起母親.
想起一個與母親有關的心理學實驗.
實驗是這樣進行的.
母親和幼兒.
被帶到一個不熟悉的房間.
房間裡擺滿了很多玩具.
當母親在場的時候.
幼兒會被鼓勵去探索周圍的環境.
放心去玩.
過了幾分鐘之後.
有一位陌生人會進來.
然後母親就會離開.
經過這次短暫分離之後.
母親又會回來.
研究人員要觀察的是.
母親離開房間以及母親重新回到房間的時候.
幼兒有什麼反應呢.
這個心理學實驗和今天的經課有什麼關係呢.
剛才大家聽福音經課的時候.
有沒有聽到.
耶穌說自己要離開.
我們試試看.
《要人福音》第十四章十八節.
耶穌說:我不撇下你們為孤兒.
我必到你們這裡來.
還有不多的時候世人不再看見我.
你們卻看見我.
因為我活著你們也要活著.
耶穌要離開.
在耶穌不在的時候.
你猜猜跟隨耶穌的人會有什麼反應呢.
然後有一天耶穌又再回來的時候.
那些跟隨者又會有什麼反應呢.
我們回到剛才的心理學實驗.
研究人員將嬰兒對母親離開和重回的時候的反應.
在最初分為三種.

$^{41}$第一種就稱之為安全型依附.
安全型依附的嬰兒在母親在場的時候.
他會將母親視為探索世界的安全基地.
當母親離開之後.
嬰兒就會變得很沮喪.
但是當母親回來的時候.
嬰兒就會去接近母親.
並且感覺到很快樂.
第二種嬰兒的反應稱為迴避型不安全依附.
那種迴避型不安全依附的嬰兒.
他們跟母親顯得疏離.
母親在場的時候都當他沒有到.
母親離開之後其實又不會特別失望.
又沒有什麼反應.
到母親回來了都沒有什麼所謂.
第三種就是稱為焦慮或矛盾型的不安全依附.
矛盾焦慮型的不安全依附的小孩.
他們在母親離開的時候顯得相當不安和焦躁.
但當母親回來的時候.
他們不會特別開心.
他們沒有辦法輕易被安撫.
如果他哭泣的時候.
你不容易安撫他.
同時間他對母親保持著憤怒和抗拒.
但不僅如此.
因為他又很渴望可以跟母親接觸.
這個心理學的實驗後來被廣泛擴充.
和應用在親子關係以及很多人際關係裡面.
今天我們不是講心理學.
我們嘗試用這個角度來閱讀.
我們剛才閱讀過的幾段經課.
我們要問的問題是.
耶穌離開之後.
他的跟隨者有什麼反應呢.
他們和耶穌維持著一個怎樣的關係呢.
雖然當年耶穌的跟隨者處境.
和我們今天的處境不一樣.
但類似的問題也值得我們問一問自己.
我們和那一位不在的耶穌的關係是怎樣的呢.
我們感覺到很安全嗎.

$^{81}$還是當他沒有到呢.
還是很矛盾地徘徊在刑拒之間呢.
其實耶穌離開的情況和心理學實驗並不相同.
因為在實驗裡面.
母親離開之前沒有預告.
但在耶穌離開之前是有充分說明和交代的.
我們一起看看福音書是怎樣交代的.
我們試著想一想 現實上就算耶穌這麼有交代提供允許.
並且設下聖靈作保衛師.
是不是就等於他沒有離開呢.
有了這些是否和他沒有離開的狀況一樣呢.
耶穌在和不在 跟隨者的經驗是不同的.
耶穌在的時候是看得見的 摸得到的.
耶穌離開之後是看不見 摸不到的.
耶穌去了哪裡呢.
我們剛才讀彼得前書的時候.
你會見到作者是需要向他的跟隨者解釋.
耶穌去了哪裡 是不是.
你試著看看彼得前書怎樣說.
彼得前書第三章二十二節.
彼得說 耶穌已經到天上去.
在上帝的右邊.
讓天使有權柄的有權能的都服從了他.
什麼意思呢 耶穌不在了.
就像小朋友問媽媽去了哪裡.
爸爸去了哪裡 去上班了.
下班會回來的 還會買東西給你.
為什麼要這樣說 祂現在不在了.
跟隨者被告知耶穌在天上有權柄.
但你試著想想 對於跟隨者來說.
耶穌在天上的權柄 跟隨者在地上看不看得到.
他經不經驗得到呢.
在彼得前書當時的字裡行間.
我們看得到的是跟隨者似乎在受苦.
第十四節.
耶穌說 你們就是畏而受苦也是有福的.
不要怕人的威嚇 也不要驚惶.
第十六節說 存著無愧的良心.
叫你們在何事上被毀謗.
就在何事上可以叫無賴你們在基督裡.

$^{121}$有很品行的人自覺羞愧.
你看到書信裡出現受苦 威嚇 驚惶 被毀謗這些字眼.
就顯示出 耶穌不在的時候.
跟隨者直接面對著攻擊和困苦.
耶穌在的時候不同嗎.
是不同的.
耶穌在的時候真的不同.
耶穌在的時候 耶穌會出來為他的門徒遮風擋雨.
你記不記得耶穌被捕那一幕是多麼英勇.
我和大家一起看 在《約翰福音》裡有這樣的片段.
有一大隊士兵就拿著燈籠 火把 武器來圍捕.
然後他們來的時候 耶穌跟他們說.
你們找誰 我們來找拿撒勒人耶穌.
耶穌怎麼回應.
我就是.
然後 這句我就是 經文說那些兵兵退後倒在地上.
是不是很戲劇性.
然後耶穌又問他 你們找誰 找拿撒勒人耶穌.
我已經告訴你 我就是.
你們如果是找我的 你就放這些人走.
有沒有想像這個場面.
這個場面 如果你是跟隨者 你會覺得很威猛.
他可以保護你 他在的時候是會為你遮風擋雨的.
但是他不在的時候 跟隨者就直接受到攻擊.
所以耶穌在的時候跟隨者和不在的時候跟隨者的狀況是不同的.
對於耶穌不在的時候 跟隨者所受的困苦.
彼得是要向他們解釋的 為什麼會這樣.
彼得怎麼說 17-18 第三章.
彼得說 上帝的旨意 如果是叫你們因行善受苦.
總強如因行惡受苦.
因基督也曾一次為罪受苦 就是易的代替不易的.
為要引我們到上帝面前 按著肉體說.
他被致死 按著聖靈說 他復活了.
彼得勸勉跟隨者學習耶穌的榜樣 為而受苦.
但是無論有什麼恩許 將來有什麼獎賞.
無論多有意義 受苦是當下的狀況.
是此刻的煎熬.
曾經在痛苦裡的人 必定感受到那一種情況.
無論是身體上的痛苦 心靈上的痛苦 或者是關係上的痛苦.
都是好當下 甚至是不能為外人道的.

$^{161}$雖然有人告訴你 耶穌在你心中 聖靈在你心中.
但是正是因為在心中 就摸不到 捉不實.
可以好像當年耶穌在的時候一樣 隨時出來為你遮風擋雨.
耶穌不在 會不會令人疑惑.
其實那些所說的 聖靈在心中 耶穌在心中.
會不會只不過是一廂情願的想像.
對於正在這裡受苦的人 有什麼幫助呢.
如果去到連聖靈在心中 耶穌在心中都不能肯定 那又如何呢.
如果有人出現這些疑問 是不是等於他們不屬靈呢.
是不是等於他們放棄了神呢.
不一定的 因為在歷史裡面的屬靈先賢都有類似的經驗.
德蘭修女就經歷了40年的心靈黑夜.
在靈魂的深處 她感受到耶穌經常缺席.
心靈的黑夜不是德蘭修女的專利.
聖十加約翰談及黑夜的時候.
他說這是人類和基督徒的典型經驗.
越認真的信徒 就越有可能經驗到這種狀況.
耶穌不在的時候 跟隨耶穌的人都有可能經歷這種掙扎.
改編自阮騰周作的同名小說《沉默》.
探討信仰苦難和神的沉默 其中有一句對白這樣說.
你沉重的沉默如此可怕 我祈禱 但我迷失了.
或者我只是在向虛無祈禱嗎 古時有人挑戰受苦的約伯跟他說.
你仍然持守你的純正嗎 你背棄上帝 死了去吧.
現代有人宣告上帝已死.
可能這些挑戰和口號都反映著人的處境和經驗.
但是人類是否可以這麼容易否定了神呢.
人類是否真的可以當神死了呢.
《哨道行傳》記載了保羅在雅典遇到的狀況.
讓我們對人的認識有多一個面向.
我們一起看看經過裡面的《哨道行傳》的片段.
第十七章二十二節.
保羅站在阿略巴谷當中說.
眾位雅典人 我看你們凡事肯敬畏鬼神.
遇見一座壇上面寫著「微色之神」.
雅典是希臘世界的文化之都 也是學術重鎮.
但是雅典人心裡面仍然有神.
只是對神的認識有限 於是稱為「微色之神」.
耶穌離開了 人類是不能夠在物質世界裡見到耶穌.
摸到耶穌 捉住耶穌.
是不是等於人類可以否定了耶穌呢.

$^{201}$或者有人可以認為.
是的 隨著人類的發展和隨著科技的進步.
其實人類是可以不理神 可以沒有了神的.
宗教是會消失的 神都可以被消滅的.
人類的世界的確在發展.
科技的確在進步 但是宗教是不是消失了呢.
神是否被消滅了呢.
以下三個發生在我們身邊的事例.
是在顯示人類的需要.
第一個狀況發生在五六十年前的某個地方.
當時宗教遭受到毀滅性的打擊.
寺廟,教堂,道觀被破壞 經典被焚燒.
僧侶和神職人員被批鬥和強迫還俗.
縱使如此強烈的狀況.
人的宗教信仰是不是被跟隨了呢.
暴風雨過後 鄉間的小寺廟重新被建造起來.
香火重新燃點.
在大城和名山 那些被封閉的寺廟,教堂,道觀都重新開放.
善男信女不斷增加.
他們不單是星斗市民 還有知識份子.
透過求神問卜,求神進香,冥想,讀經,參加宗教聚會.
安頓一己的心靈.
更加有些是致力於社會倫理的昇華和改造.
重建人民的世界.
第二個例子發生在我們附近.
在很近的地方.
有一間叫應用科學研究院的公司被揭發.
豪花十八萬公帑看風水.
引起了政府和社會強烈不滿.
有議員批評應科院作為科研機構.
不應該花十八萬來看風水.
反問風水重要還是科研重要.
第三個例子是在哪裡體驗到在人類社會中宗教氣氛原來如此強烈.
就是殯儀館.
當我們去殯儀館的時候.
我們會發現原來這麼多人在用宗教儀式.
無論那些是道教的破地獄,燒紮寺.
還是佛教的誦經,超道或者是基督教的安息禮拜.
你都會發現那個地方原來人去到某個時間是充滿著宗教感的.
耶穌不在.

$^{241}$但是人類就不輕易因此而當神舞蹈.
如果在耶穌不在的時候.
怎樣可以言說耶穌呢.
耶穌不在.
你怎樣去言說耶穌呢.
怎樣可能在拿不到神出來的情況下.
來去言說神呢.
你叫你的耶穌下來給我看看.
是叫不到.
那怎樣還說耶穌呢.
我們看看保羅怎樣說.
保羅做了一個示範.
我們回到《哲學行傳》第十七章二十四節.
保羅對雅典人說.
創造宇宙和其中萬物的上帝是天地的主.
不住在人手所造的殿.
也不用人手服侍.
好像缺少什麼.
自己倒在生命氣息,萬物,次級,萬人.
祂從一般做出萬俗的人住在全地上.
並且預先定準他們的年限和所住的疆界.
要叫他們尋求上帝或者可以揣摩而得.
其實距離我們各人不遠.
我們生活動作全流都在乎祂.
就如你們作詩的有人說.
我們也是祂所生的.
我們既是上帝所生的.
就不當以上帝的神聖.
人用手藝心思所雕刻的金銀石.
世人愚昧無知的時候.
上帝並不監察.
如今卻吩咐各處的人都要悔改.
因為祂既定了日子.
要藉著祂所設立的人.
按公義審判天下.
並且叫祂從死裡復活.
給萬人作可信的憑據.
古代希臘的哲人認為人類有三個大問題.
其中兩題就是.
我從哪裡來的呢?.

$^{281}$我會去哪裡呢?.
大家都可能看得出.
古代向雅典人介紹神.
大致上都是請他們思考這兩個問題.
我是從哪裡來的呢?.
我會去哪裡呢?.
歷世歷代不少先賢先哲.
在思考這兩個問題的時候.
他們透過思考這兩個問題.
來思想到神.
透過思考這兩個問題來認識神.
並且與神建立關係.
保羅示範了.
我拿不到一個可以摸得到.
可以捉得實的神出來的狀況底下.
來去言說神.
保羅示範了.
耶穌離開了.
耶穌還沒回來的情況底下.
來去言說耶穌.
其實詩篇的作者提供了另外一個方法.
給我們參考.
如何在那位摸不到的神不在的時候.
我們來去言說祂呢?.
第66章第10節.
在詩篇那裡.
作者說:上帝啊!.
你曾試煉我們.
熬煉我們.
如熬煉銀子一樣.
你使我們進入網羅.
把重擔放在我們身上.
你使人坐車.
砸我們的頭.
我們經過水火.
你卻使我們到豐富之地.
凡敬畏上帝的人.
你們都來聽.
我要述說.
祂為我所行的事.

$^{321}$我曾用口求告祂.
我的舌頭也稱祂為高.
我若心裡注重罪孽.
主必不聽.
但上帝實在聽見了.
祂則已聽了我禱告的聲音.
上帝是應當稱頌的.
祂便沒有推卻我的禱告.
也沒有叫祂的慈愛離開我.
詩人在做什麼呢?.
詩人如何去說那位.
「摸不見看不見」的神呢?.
他正在用的方法.
是分享他與神交往的經驗.
就類似我們今天所說的.
「講見證」.
希伯來書也是這樣勸勉信徒.
希伯來書作者這樣說.
「所以既然我們有這許多見證人.
如同雲彩圍繞著我們.
就該卸下各樣重擔和謹慎的罪.
以堅忍的心.
奔那擺在我們前頭的路程.
仰望我們信心的創始成中者耶穌」.
耶穌不在這裡.
希伯來書作者說.
「但有許多見證人在這裡.
有人認為信仰是對真理的一個本質性.
的肯定和認同.
而不是直接的找實.
回想起那個心理學實驗.
它會不會也在啟發我們.
來檢視我們與神的關係呢?.
耶穌離開了 耶穌答應有一天會回來.
在這一段耶穌不在的時間裡.
我們與神的關係是怎樣的呢?.
究竟是安全還是其實祂不在也無所謂.
當祂不在還是會否很矛盾很焦慮.
祂不在 要我受那麼多苦.
不來幫我 不來救我.

$^{361}$我恨祂 但我又不可以沒有祂.
究竟我們與這位不在的耶穌的關係是怎樣的呢?.
或許也會引發我們很多自我的檢視和思考.
不過無論是哪一種狀況也不太重要.
為什麼呢?.
因為我們與神的關係不是決定於我們.
而是決定於神的恩典和憐憫.
以及我們對祂的信實和信賴.
願上主幫助我們.
\newpage



\section{約書亞記 5:14}
\label{sec:032110}
\textbf{聖經人聖經事:我來是要作元帥}~~~~黃克勤牧師博士
\newline
\newline
連結: \href{http://www.church.com.hk/acms/content.asp?site=cdc\&op=show\&type=product\&code=032110\&layout=sermon}{\texttt{032110}} ~~~~ 語音日期: 2026-05-09
\newline
\newline
\hyperref[sec:032109]{\small{< < < PREV SERMON < < <}}
~
\hyperref[sec:index_chronic]{\small{[返順時目]}}
~
\hyperref[sec:032111]{\small{> > > NEXT SERMON > > >}}
\newline
\newline
約書亞記 5:14
\newline
\begin{longtable}{cl}
\hline
\hline
章節 & 經文 (和合本修訂版)\\
\hline
5:14 & \begin{tabularx}{0.7\textwidth}{X} 他說：「不，我現在來是要作耶和華軍隊的元帥。」約書亞就臉伏於地下拜，說：「我主有甚麼話，請吩咐僕人吧！」 \end{tabularx} \\ \\
[1ex]
\hline
\hline
\end{longtable}
$^{1}$今日的題目是「我來是要作元帥」.
大家留意,「我來是要作元帥」.
這句說話是上帝說的.
「我來是要作你的元帥」.
「你們的元帥」.
神是否你的元帥呢?.
還是我們以神做我們的一師.
做我們的跟班.
我叫祂,點祂做事.
大神說「我來到你那裡是要作元帥」.
大家留意.
《約書記》第五章第十三節.
「約書亞靠近耶利哥的時候.
第一點.
有一個人手裡有拔出來的刀.
對面站立.
約書亞到他那裡.
問他說.
你是幫助我們呢?.
是幫助我們敵人呢?.
問號」.
多有趣.
約書亞會問這樣的問題.
你是幫我們還是幫我們敵人呢?.
讓我想起當年我們小時候.
幾十年前看我們香港的粵語殘片.
曹達華多數是做探長的.
石堅多數是做奸人做賊的.
我記得一幕很清楚.
就是曹達華探長問石堅.
究竟你是忠還是奸呢?.
是差不多的語氣.
你是忠還是奸呢?.
多有趣.
探長會問賊你是忠還是奸呢?.
你不是忠還是奸呢?.
他會回答你.
你是幫助我們還是幫助我們敵人呢?.
多有趣.
有原因的.

$^{41}$「約書亞靠近耶利哥的時候.
其實這是說什麼呢?.
就是他們預備要進攻耶利哥.
奪取第一個迦南地的城.
在這個時候你作為元帥.
作為領軍的人.
那你要怎樣呢?.
當然是出去.
到處逛逛.
觀察一下耶利哥城.
去看看.
所以「靠近耶利哥的時候.
看看耶利哥這個城.
是怎樣的城.
城牆估計有多厚.
能不能頂住.
有多少個城門.
這個「約書亞靠近耶利哥的時候」.
會不會只是他自己一個人出去呢?.
不知道.
可能他一出去就有十個八個人跟著.
不知道.
不過不要緊.
「約書亞靠近耶利哥.
想不到不料.
第一點.
是有一個人手裡拔出來的刀.
對面站立」.
在逛逛的時候看到有一個人站立.
手裡拔出來的刀拿著刀.
對面對著他站立.
「約書亞如果你是約書亞你會怎樣呢?」.
大家害怕.
立刻把手放在刀柄上.
隨時預備要拔刀.
這個不知是否耶利哥走出來的人.
但是他應該又不是這樣.
他會問這樣的問題.
如果是一個人.
他會問「你哪一邊的?」.

$^{81}$似乎他一看就知道.
他不是一個普通的人.
看來他看到的是一個.
天使或者神靈的人物.
可能是有靈氣的.
一看他就知道他不是普通人.
是上帝派來的人.
或者是靈.
不過他又有一個人站在那裡.
問他「請問你是幫我們還是幫助敵人的?」.
「你是敵是友?」.
很直接地問「你是敵是友?」.
他回答說「不是的」.
即是說「你問錯了」.
你問我是幫你的還是幫你的敵人?.
「是敵是友?」.
這個人回答「不是的」.
「你問錯了」.
「我來是要作耶和華軍隊的元帥」.
「我來是做元帥,現在上任」.
「我是來做耶和華的元帥」.
然後約書亞就俯伏地下拜.
其實約書亞見他的時候.
就知道他應該是一個很有能力.
不單是一個普通的人.
看來他可能是天使之類.
他就覺得有一個這樣的人.
究竟是幫我們神派來.
還是幫對面派出來的?.
他就問了.
他說「你問錯了」.
「其實我來是要做你們的總司令」.
「我來是要做耶和華軍隊的元帥」.
這個人是誰呢?.
不樂乎只有兩種.
一個是天使,神派來的天使.
而另一個就是上帝自己.
很多時候在聖經裡面.
特別是舊約聖經.
很多時候神的使者.

$^{121}$可能是使者,又可能是神自己.
不知道,也無需猜測.
約書亞一見到他就說.
「我是來做耶和華軍隊」.
「就是你全軍的元帥」.
於是約書亞就說.
「喂,你是元帥還是我是元帥?」.
「我是元帥,現在出來看看軍情」.
「現在你是元帥,那我是什麼?」.
「那你欠我的位子嗎?」.
約書亞馬上俯伏地下拜.
「我主要什麼吩咐僕人的話」.
「有什麼吩咐我?」.
「我是你的僕人」.
他馬上俯伏.
「原來是上帝參拜你的元帥」.
「那我就是你的僕人」.
即是說,他馬上知道這是元帥.
他就放下自己元帥的位子.
臣服在上帝的手下.
這個很重要.
老老實實.
我們基督徒有些信主幾十年.
他都是自己要作爺和華的元帥.
他怎樣爺和華做事.
「臣啊,我要什麼?我要什麼?」.
「我什麼時候要第一桶金?」.
「臣啊,我要這個女孩做我的朋友」.
「哎呀,我兒子要讀什麼學校?」.
「哎呀,我想怎樣?我想怎樣?」.
「你幫助我吧」.
「我每個星期去教會」.
「風雨不改,還捐了錢」.
「那你一定要幫我」.
「哎呀,你幫我,你幫我」.
自己出了主意.
然後在一個最關鍵的轉折點.
自己不能夠控制.
先求上帝你幫我這邊.
但其實.

$^{161}$爺和華君子的元帥.
就是他來出現統領我整個人.
而不是去到某個.
我覺得不搞得定的位.
才叫他幫我.
變了其實.
那個神,我們的神.
其實變得很小.
在我們心目中很小.
所以有人寫一本書.
基督徒寫的.
「Your God is too small」.
「你的神是太小了」.
「是阿四來的」.
「是阿姐姐來的」.
「就是幫手做事的」.
「不是帶領我們」.
「不是我們的元帥來的」.
神說「我是來做你的元帥」.
「今天神就是跟你說」.
「我要做你生命的元帥」.
「願不願意?」.
「我主有什麼吩咐僕人的說話嗎?」.
這個態度真好.
「我主有什麼吩咐僕人的說話呢?」.
「當你早上起床」.
「最好的那一招的神禱是什麼呢?」.
「主啊」.
「多謝你給我新的一日」.
「又有氣息又有呼吸」.
「你有什麼要我做的呢?」.
「今天讓我知道」.
就是孩童薩姆爾.
跑到他的以利濟斯面前.
「有什麼要我做的呢?」.
濟斯以利說.
「我沒有叫過你」.
「如果你再聽到叫你的名字」.
「就是神叫你的了」.
「請說僕人敬聽」.

$^{201}$「早上起床最好的禱告就是」.
「主啊」.
「今天有什麼要我做的呢?」.
「請說僕人敬聽」.
「我是他的僕人」.
馬上呼服.
「我主有什麼吩咐僕人的話呢?」.
這個很好.
跟隨主更好.
然後耶和華軍隊的元帥對約書亞說.
「將你腳上的鞋脫下來」.
「因為你所站的地方是聖的」.
約書亞就自在行了.
由這一句.
你就會知道.
這個出現的並非別人.
是不是一個天使呢?.
由這裡到第三點所說.
「喂,你脫下鞋子」.
「因為這裡是聖地」.
「是聖的」.
「分別的」.
這樣.
就是說.
這個是上帝向約書亞顯現.
一個武將這樣出現.
拔刀預備為他打仗.
帶領他們打仗.
為什麼這樣呢?.
因為你所站的地方是聖的.
這個通常都是上帝說的.
「摩西啟礦日見到荊棘被火燒著」.
「而是神在荊棘被火燒著裡面」.
「呼召他過來」.
「然後跟他說你脫下鞋子」.
「因為這裡是聖地」.
這樣.
他們的習慣就是.
如果面見上帝的時候.
要脫下鞋子.

$^{241}$就是這個鞋子.
是表示著我們的行為.
我們的行徑.
是污穢的.
脫下鞋子.
用自己的腳踩下去.
才是自己.
是真正的尊敬的.
因為鞋脫下來.
你站的地方是聖的.
有時候是朝著陷廟.
在這裡知道.
每一件事.
我們的元帥叫我們做的.
無論我去到他面前祈禱靈修.
默想聽他的說話.
又或者他叫我做的事.
全部都是.
聖的.
聖的意思不是乾淨.
不單只是聖潔清潔.
是指不同.
聖哉聖哉聖哉不是乾淨乾淨乾淨.
而是.
是不同不同不同.
神人之間有一個很大的分別.
我拍馬追不到上帝的神聖.
見到他的時候.
這裡是分別的.
要知道我是分別為.
聖站在神面前.
所以當我們的元帥.
上帝叫我們做事的時候.
你要知道這一件是.
聖.
神聖的工作.
神聖的事奉.
這樣就可能是神了.
今天有什麼要做?.
請說服人聽聽.

$^{281}$神就說.
你上班做好它.
今天的工作做好它.
上班不是有國家大事嗎?.
這樣要我做嗎?.
不是的.
你上班.
將你的工作做好它.
這個也是一個神聖的工作.
因為是神叫你做的.
就不要上班.
隨便一點.
只是發工資而已.
不是的.
將工作做好它.
耶穌所說的叫我們這樣.
不要只在人的眼前事奉.
像是討人喜歡.
而是要好像服侍神一樣.
所以每一個工作.
都是聖的.
都是神聖的.
做好它.
因為是服侍著上帝的.
都是懈脫下來的.
記得你站所走的地方.
都是神聖的.
都是神選給你要跟隨的.
所以每一次.
神有什麼給我做?.
請說服人聽聽.
然後神叫你做的.
全部都是聖的.
其實我們每一天的生活.
都是屬於神的.
都是要脫鞋.
都是尊敬為聖的.
今天說的三點.
在這裡和大家看到了.
再溫習一下.

$^{321}$上帝向你顯現.
拔出來的刀.
我是來做元帥.
我是準備幫你的.
你不用擔心神不幫你.
神說我已經來了.
拔出的刀.
我是來做你的元帥.
然後你所做的每一件事.
都是神聖的.
祈禱.
主,這段聖經.
給我們心裡很震驚很惶恐.
因為你是我們生命的元帥.
主,很多時候我們會指指點點.
哎呀,叫你做媽.
希望你做媽.
為什麼你不做?.
為什麼你不在那個時間做這件事?.
不過其實原來我是要跟隨.
正如若說一見到你.
他馬上俯伏下拜.
請問有什麼吩咐給僕人呢?.
這個態度真是好.
然後神就叫他.
你脫下鞋子.
這裡是聖地.
主,幫助我們記得.
每做的每一件事.
都是屬於你的.
是歸主為你而作的.
讓我們這樣去生活.
這樣去跟從.
順聽祈禱.
奉靠耶穌基督的名.
阿們.
\newpage



\section{箴言 22:6}
\label{sec:032111}
\textbf{為母則剛}~~~~林羅麗斯牧師
\newline
\newline
連結: \href{http://www.church.com.hk/acms/content.asp?site=cdc\&op=show\&type=product\&code=032111\&layout=sermon}{\texttt{032111}} ~~~~ 語音日期: 2026-05-10
\newline
\newline
\hyperref[sec:032110]{\small{< < < PREV SERMON < < <}}
~
\hyperref[sec:index_chronic]{\small{[返順時目]}}
~
\hyperref[sec:032112]{\small{> > > NEXT SERMON > > >}}
\newline
\newline
箴言 22:6
\newline
\begin{longtable}{cl}
\hline
\hline
章節 & 經文 (和合本修訂版)\\
\hline
22:6 & \begin{tabularx}{0.7\textwidth}{X} 教養孩童走當行的道，就是到老他也不偏離。 \end{tabularx} \\ \\
[1ex]
\hline
\hline
\end{longtable}
$^{1}$我想問一下其實我們當中的影子會有多少自己的媽媽仍然健在的 揚下手 都很多 和隔壁那個說你很有福氣.
還可以和媽媽慶祝母親節 我又想問多少其實自己已經是媽媽的 揚下手 更加多 老公更加要和他們說 祝他們母親節快樂.
這個主日是母親節 所以選了一個題目叫做為母則剛 為什麼會想起這個字呢 當然之前有一條片都經常提及過.
這句說話其實是來自梁啟超的新聞說 原文是說婦人若也 而為母則強.
你知道梁啟超這些名字 是上一代的名字 文初的時間 在當時的環境很不同 婦女大部分都是在家中相夫教子.
你知道那段時間還在說女權開始在中國興起 爭取讀書 甚至做工的機會都是很少的 因此一般就會覺得女人是弱一點的.
除了身體機能弱之外 在社會地位賺不到錢 當然是會低微一點的.
我亦相信因為要靠著男人生活 你一定會溫婉柔順 但到今時今日 二零二六年.
看看你旁邊的那個 可不可以用溫婉柔順來形容自己呢 都有點打冷顫.
但無論如何 起碼一個很柔弱的 可能在家中吞聲忍氣的女士 往往在甚麼時間會變得剛強呢.
就是母愛的偉大了 通常就是為了孩子 為了家中的事認仇 在因無可因之下 當時的女人都會變得比較剛強.
但時間過得很快 我們的社會也適應和轉變得很快 由開始可以讓女士讀書 其實只是幾十年之間.
現在整個社會的態度都已經很不同了 甚至香港被視為女權非常高漲的一個地方.
但是否真的是一個很幸福的狀態呢 平心而論 我們也不是這樣想的 原來平權平到又要上班 加上又要照顧孩子.
早知道就平一點好了 不要平那麼多 因此你看到現在一直是這二三年 應該說十多二十年 做媽媽的一定是這樣的姿態.
我們被稱為是雙職的母親 即是她不單要上班 她也會打另一份工 就是要繼續做媽媽 做家庭主婦 其實是同時來去兼備的.
而且還有一個很大的變化 我覺得在我們九十年代做媽媽的 其實我是很幸福的那一代.
我們開始要出來工作 但是在那個年代就開始引入工人的時代.
但是你問我那個氣氛還是在考慮請不請工人 覺得工人進了屋會有什麼影響.
還有大部分的男人在當時來說 他們仍然可以擔得起頭家 養到老婆不用出去做工是引以為榮的.
所以我經常說我老公一生的一件事就是他沒有試過養我 因為我們是一起出來事奉的.
但是是不是一起出來工作的經濟壓力好像少了之後 夫妻之間的壓力就少了呢.
其實我覺得也不是的 雖然今天是母親節 但是平心而論 今天母親的角色已經有很多爸爸在當中去分擔.
我們今年搞的啟發的親子課程是第一次是爸爸參加多過媽媽的.
這個是新景象來的 可能將來也會是這樣的 親子房是爸爸坐多過媽媽坐的.
但是為母則綱是好聽的 好像很讚揚媽媽那樣 我趁母親節也弄一句讚揚我們的父親.
因為我覺得現在這一代的爸爸 我稱他們為為父則勇 為什麼呢.
因為這一代年輕的爸爸其實他們都是媽寶來的 同意嗎.
他們都是被媽媽疼大的 特別始終中國人還是有少少重男輕女 覺得男孩子不應該做家務的.
所以就算包括老公那些 是完全不需要理家裡的事和進廚房的.
但是突然之間做了人家老公之後 就發覺很糟糕了 原因是真的發覺自己低能.
好像你老婆以為你什麼都懂 但你發覺原來你什麼都不懂.
所以洗碗也要教三次才洗得完整的一隻碗.
你想一下在他們的心目當中他會有多怯呢 外表就好像很強壯的樣子.
但是進到家裡 進到廚房就發覺自己好像弱智一樣 又不知道怎麼去面對.
還有好像我老公說的 你以為是男人就不怕老鼠 不怕蟑螂嗎.
我沒有想 因為我給蟑螂有個耳朵 我見到蟑螂一定瘋的.
但是他說 問題是我們都還沒叫得出口 你們還瘋過我們 還瘋狂過我們.
還有例如老公老公什麼的 你以為我們真的可以有這個勇氣去打死蟑螂 或者追著老鼠嗎.
他說很坦白 大大聲聲叫在那裡拍 其實是想說老友你走快一點.

$^{41}$我們搞你不行 連想執屍的能力都不是拍夠的.
因此事實上我覺得現代的家庭 我們被父母那一代愛得太緊要了.
所以無論我們做媽媽做爸爸都是很大的壓力 很多的憂慮.
所以你明白為什麼現在很多年輕的一代都不是很敢去生小朋友.
我們兒童教會是缺幼稚的 你明白嗎 很想多一點有嬰兒能夠進入當中.
在信仰裡面 我們知道上帝創造人是要我們一男一女去成家立室.
然後生小朋友 一代一代的傳下去.
但是很有趣的 聖經教如何去做父母的經文 說真的不是太多.
講來講去一定例牌 不可以離開這一節我們要做金句的.
就是在《箴言》的22章第6節 教養孩童 使他走當行的路.
就算到老都不偏離 是不是個個都知道.
個個在教會做過媽媽也好 做子女也好.
我們都聽過 都看過 但是我不知道大家怎樣去看待神的話語.
既然這一句已經是我們兒童教會做父母一講座必講的時候.
你問我真的做了幾十年媽媽 都是很困難.
因為事實上聖經講就很容易 但是你問我什麼是當行的路.
我的心裡面已經想了很多年 哪一條路是當行的呢.
我們自己都不知道前面的路是怎樣走.
我們不同階段 現在退休之後怎樣呢 無限循環.
但是還要教他們走他們當行的路.
上帝不如清清楚楚告訴我 他走東走東 走西走西.
我都容易跟 我就叫他走東 上帝叫你走東 你就走東.
但是聖經沒有寫得很清楚 而最重要還有一個很考我們深思.
就是怎樣令到他可以到他老都不偏離.
這個事實上真的是一個很難很難的功課.
其實我們在人生成長的過程裡面.
做老婆為人夫的 或者做媽媽的 做爸爸的.
每一個階段 現在做祖父母的我們已經開始.
好像每一個階段我們都不知道應該怎樣去帶領我們的下一代.
或者搞不定的時候某程度來說我們不是放手.
很多是放棄的 甚至祖長也是 祖長你們都是做父母的心腸.
用盡各樣的板斧 瘦也好不瘦那些 搞到頭髮都白了都搞不定的時候.
到頭來不是真的放手 是真的心裡面說 我都放棄了 做不了多少.
很多男人就說 唉 不知道為什麼 娶個老婆的時候一定是喜歡她這樣的樣子.
不是漂亮 是少了伊人的樣子.
幾個女人她能夠結婚她總會有少了伊人的樣子出現.
才能夠令到有個男士想要娶她.
但是不知道為什麼當人一做了媽媽的時候.
我們總會發現你身邊的老婆就變樣了.
就會變成一隻鵝雞那樣出現.

$^{81}$多麼溫柔多麼小聲也好 不知道為什麼做了媽媽的聲音就突然擴大.
不用拿咪高峰 音頻也非常高.
頂智慕其實也是的 我們在香港的壓力很大.
情緒已經是成為我們一個很主要的課題.
但是想深一層 其實真的不知道為什麼.
往往情緒出現了 你才是不知道怎麼樣來得及去修課.
所以為什麼我們需要在我們行事的時候.
好好地去認識自己 好好地去學習倚靠神.
因為不要說要控制你的下一代.
其實有時候人控制自己都已經覺得很難.
究竟我們的情緒出現背後是因為我們在想些什麼.
很多時候連我們都不知道.
但是行為已經說出來了.
我們明明心裡是想愛的.
但是我們說出來是刻薄的 串串共的.
我們總是詞不達意 全情達意上是表達不了的.
難怪其實每一個年代的父母和子女之間的關係都是很複雜.
在普羅在羅馬書的十二章的第二節.
我們平時經常讀第一節.
但是我覺得第二節單獨去看的時候是一個很好的提醒.
我自己也很喜歡在這段經文裡面去想.
普羅說第一件事 不要效法這個世界.
要心意更新而變化.
目的是什麼呢.
是叫我們察現什麼是上帝的旨意.
背今句中間的形容詞是凹凸不平的.
但是當你這樣掩著它 你就很明白.
普羅告訴我們其實這個世界的變化很大 是真的.
魔鬼會用盡各樣的方法讓人的心來偏離上帝的路.
但是我們卻很喜歡跟從這個世界.
一個人說的話 你聽起來覺得怪怪的.
但是十個人都敢說 每個人都敢說的時候.
我們就覺得是真理.
普羅提醒我們不要效法這個世界.
反而要心意更新.
在幾個禮拜的討論裡 我們的姐妹就說.
為什麼經常說要更新 不要說要更新.
我現在都很好 何必更新.
普羅說不是.
因為我們活在這個世界裡 都一直在改變.

$^{121}$魔鬼會用各樣的方法來迷惑人.
以前那一代 拜樹 拜山 拜海.
接著就拜偶像 有不同的人物出現.
現在下一代都沒有形象了.
你很少看到他們特別搭擺什麼.
除了喜歡玩偶之外.
但是他們是一樣迷信的.
你就知道這個世界有哪些惡者不斷用各樣似是非的道理.
來引導人離開我們的上帝.
因此我們需要心意更新的原因.
是要測驗什麼是上帝的旨意.
如果我們沒有安靜的心 謙虛的心.
願意去測驗的心.
其實我們人云亦云 說著說著就亂了.
我記得我們90年代做爹地媽咪.
那時候很流行有個小盆的枕頭.
因為它套上去 讓寶寶的頭圓一點.
我媽那一代是用來搓的.
搓完寶寶的頭 但我不會搓.
有個枕頭都怪怪的 你都願意買的.
睡了不到一歲 原來小朋友趴著睡.
就會好一點 於是很多人都在那裡趴著睡.
但我心裡就覺得怪怪的.
剛剛才說那個盆令到頭圓一點.
沒理由臉會扁的 是不是.
感覺很奇怪 但我沒有跟著.
不理會它 頭圓也合理.
如果趴在裡面 焗死或者會扁.
沒理由臉會扁的 是不是.
我就沒有理會.
其實去到我女兒七歲至八歲.
就已經說 因為美國先用了.
就發現趴著睡的寶寶.
患心臟病的機率是高一點的.
於是一下子之間.
每個媽媽就把子女反過來.
你明不明白.
如果你留心去想想你身邊所說的話.
無論吃東西也好 養生也好.
我們那一代的三角形.

$^{161}$現在吃了這麼多年.
你才跟我說要調轉它.
怎麼調轉 人是站直的.
難道調轉它嗎 是不是.
很辛苦的 要改變吃東西的習慣.
各位兄弟姐妹 這個世界就是這樣.
問題是我們怎樣真的與神同行.
走在上帝的任避裡面.
以致我們的安全.
是我們需要有一個願意.
謙卑 謹守的心.
不是說將所有的東西放下.
但是你都願意被人改變.
因為事實上 時代真的會進步.
但是又不可以一窩蜂.
人說什麼 我們又說什麼.
的確是一件艱難的事.
不過我覺得彼此要勉勵.
捉到保羅的好意.
我們所謂改變的目的.
不是因為時代興.
不是因為人人都是這樣做.
不是因為怕倒閉.
是因為我們要留心.
上帝在這個時代的心意究竟是什麼.
其實上帝很愛我們.
你有沒有想過.
上帝創造人的時候.
他已經明白我們的軟弱.
你有沒有想過.
為什麼那些嬰兒這麼可愛.
當你回去看照片的時候.
我們這裡比較好.
我們有照片.
年輕人很久沒有照片.
我們看照片的時候.
往往你很喜歡看.
就是他小時候的樣子.
坐到歪了 諸如此類.
但是你覺得真的很可愛.

$^{201}$神不只是做人.
嬰兒的時候這麼可愛.
很坦白 貓貓狗狗.
甚至是熊貓.
他小時候就很可愛.
但不知為什麼很快.
一下子就變大了.
為什麼呢.
你有沒有想過.
上帝的創造都有他的智慧.
如果我們的嬰兒不是這麼可愛.
我相信大部分過不了兩歲.
就被我們悶死了.
你有沒有想過.
無端端又哭.
然後床又濕了.
然後學吃東西更恐怖.
手又拿起來.
頭二分在頂上擋頭髮.
你見到他.
你真是火都來了.
但是沒辦法.
他的樣子很可愛.
傻傻地笑著看著你.
就算你忍不住罵了他兩句.
然後他就很可憐地.
含著一堆眼淚對著你.
那你怎麼辦.
罵不罵得下.
大部分.
其實在他們很小的時候.
他們都不聽你的話.
因為他們根本聽不懂你說什麼.
然後他們都任意妄為.
衣服喜歡怎樣搞.
杯水喜歡怎樣搖.
都是他們自己說了算.
完全不聽話.
但是我們都能夠愛與寬恕.
我們都能夠包容他們.

$^{241}$是不是.
你知不知道.
你有沒有想過.
那種感覺是什麼.
這種叫做愛.
而且是由上帝而來的愛.
我們普遍來說.
是一個叫做惻隱之心.
一般你想一下.
你突然間看電影也好.
或者在街上.
或者人家需要也好.
你總是有那種觸動.
就是那一刻.
就是感動來.
明明是不應該幫的.
或者明明你真的想罵.
但是你都會收得很快.
就是那份上帝的愛.
觸動來的那一刻.
聖經很喜歡用耶穌.
去行神蹟奇事的時候.
都有這四個字.
就叫做動了慈心.
就是上帝將這份愛.
特別加強在做父母的身上.
當我們按理來說.
就應該打到抽筋的時候.
總是那一刻.
我們的心就被這份愛觸動了.
我們就願意放下自己的要求.
放下嘴巴.
我們就可以放下脾氣.
變得很溫柔地跟他們說話.
其實聖靈的果子.
某程度來說.
當你常常連繫在神的愛裡.
我們個個其實都做得到.
當你做過父母.
當你做過媽媽,爺爺.

$^{281}$或者你幫人家在家裡.
照顧小孩的時候.
你就總會領略到.
當這份愛與觸動臨在的時候.
你就不知道為什麼.
多麼辛苦.
你都會很開心.
當我們看到.
應該想捏死他的時候.
突然間就很溫柔.
又可以忍耐.
又可以控制到自己.
這個就是.
當我們常存在上帝的愛裡.
那種威力.
做父母我們知道.
這條路是不容易的.
子女越大的時候.
他們未必願意聽話.
當在爭拗裡.
你看到大部分.
都是父母引養.
我們的父母.
去到最後.
我們自己做了成人.
我們回頭想都是一樣.
當我們賺到錢.
我們惡.
我們自己作主的時候.
父母越來越弱的時候.
你看到其實每一次.
自己衝口而出的.
都知道自己不對的時候.
但是你看到做父母的.
到頭來.
都是他們引養.
這一份就是.
最值得我們去感恩.
和多謝的.
就是爸爸媽媽.

$^{321}$對我們這一份的愛.
但是很有趣.
這一份愛不知道為什麼.
在嬰兒開始長大的時候.
就開始減退.
原因是.
可能我們受了世界的教導.
我們要教養孩童.
我們要讓他為自己負責任.
將來的路是他走的.
於是我們對他們的期望.
變了.
我們對他們的要求.
也因為這樣.
不斷地去爭奪.
以致往往我們在愛.
和要求裡面.
常常自己也會打架.
究竟我們何時要很嚴厲.
何時要很恩慈.
其實真的很困難.
但我很喜歡用恩威並重.
成為我們做父母學習的目標.
我自己也這樣想.
究竟何時要完全的寬容.
去接納.
在什麼時候.
要站穩一些界線.
讓他自己站起來.
真的很難拿捏.
很需要在神的面前.
不斷地禱告和祈求.
因為我們真的需要.
神給我們的智慧.
去傳遞愛.
其實真的很需要智慧.
不過我們每一個人的背後.
的取向.
的衝動的能力.
的快慢.

$^{361}$其實很代表我們身在不同的時代.
因為不同的時代.
世界的價值觀.
我們如何受父母教導.
都會影響我們.
我自己當自己是中間.
我們上一代的.
可能現在是八十歲的媽媽.
九十歲的媽媽.
都有很多健在.
在他們的身上.
我們自小就知道.
打者愛也是正常的.
媽媽打到你抽筋的時候.
都說愛你和為你好.
我們在對於傳統父母的權威.
是很強的時候.
我們明知是不合理.
但是我們都不敢反抗.
所以那所謂的聽話.
某程度來說.
在一個很大的壓力之下.
令我們做好自己的本分.
但是這個壓力.
是因為他們是我們的父母.
是講長幼有序的觀念.
而最重要的一件事.
你慢慢會明白.
當我們自己生了小朋友.
父母的壓力就來了.
我們通常已經學了讚賞.
對小朋友是好一點的.
容易愛表達.
但是父母的理解.
就是你不要整天這樣讚賞.
會讚壞的.
所以一生出來.
就說你女兒醜.
你想想孕婦已經很辛苦了.
剛剛看到女兒多可愛的時候.

$^{401}$每個親戚就說.
我女兒很醜.
其實都挺大傷害的.
但是沒辦法.
他們的價值觀就是這樣.
他們的背後目的.
是想女兒越來越漂亮.
是人格分裂的.
但是我們這一代表達的愛.
父母表達的愛就是這樣.
永遠都是很嚴厲.
然後就是不罵你.
就等於讚你.
所以你明白.
我們這一代.
都不是這麼容易用說話讚人的.
要學的.
因為我們整個觀念就是.
不可以亂讚的.
讚是會壞的.
而大部份在我們的年代.
很多男生仍然很有責任感.
覺得有使命要養起頭家.
他們認為養起太太.
才是他們的光榮.
但是很坦白.
現在這一代沒有人這樣想的了.
結婚的時候.
你都找不到誰想養你.
大家想孕婦一份.
最重要是大家還要有自己的錢.
所以其實我們現在在香港.
時代走得很快.
我們自己做阿婆阿媽.
我們這一代開始.
你發覺如何祝福下一代呢.
大部份都想不要理那麼多了.
不是放手.
其實是放棄.
覺得理不了.

$^{441}$但是我想頂替我們去明白.
現在這一代長大的.
新一代的父母.
他們都不是這樣想的.
因為他們都受了下一代的氛圍影響.
近十幾二十年.
你知道情緒比天大.
加上現在越來越多的測試.
每個子女都總有一個問題.
所以發放的學生.
現在已經過半數.
你明白嗎.
所以跟隔壁的弟兄姊妹說.
發放是正常的.
大部份是發放.
發放就是正常.
不發放的才是不正常.
是很難理解的.
我們最近有兩個個案.
有一個中一二的學生.
很多中一二不上學.
他去看心理醫生.
拿了一張心理醫生的紙證明.
他是低能量.
不是低能.
他不是低能.
他是低能量.
低能量是什麼意思.
低能量就是他的能量.
不能夠滿足一天這麼長的學習.
所以他可能一個多小時.
或者兩個小時.
拿了一張紙就大了.
現在學校很有趣.
因為我們知道有聲音的學生.
每一個學校的資助.
就可以多拿一份.
所以你真的搞錯了.
以前是很難發放的.
現在是很容易的.

$^{481}$所以你不要受這個世界影響.
你應該從神的眼光去看你的弟弟.
不要拿著那張紙.
有張紙就大了.
回家給老師.
於是小朋友晚回來也可以.
早走也可以.
因為你現在有張紙.
大家不敢鬥你.
免得你投訴.
然後有個中三中四的女生.
評定她是不語症.
小時候拼了拼到飛起來.
現在叫不語症.
她不說話.
她在病中不能夠和你們溝通.
不能夠和你們說話.
於是老師要接受她.
她不需要和同學玩.
因為你不應該逼她生病.
不語.
然後老師舉手問問題.
問學生.
你可以問她.
因為她生病.
不語.
你強行要她說.
就對她的病情增加.
然後影響她的學習.
終於很坦白.
一個小孩一個小時就回家.
兩個小時就回家.
一個小孩坐在那裡.
一天在那裡發呆也好.
沒有人理會.
兄弟姐妹.
這樣可以幫到他們嗎.
不知道.
太多壓力了.
太多困難了.

$^{521}$太多新鮮的事情沒想過.
家長有時候覺得.
有張紙就好了.
免了自己很大的責任和困難.
總之她有病就算了.
在新的高舉體貼情緒.
令到很多年輕的父母.
放棄了權威型的教導.
因為說真的.
我們自己都不開心.
我們被父母打大.
罵到大.
我們都知道不開心.
己所不欲勿施於人.
所以我們不想我們的子女是這樣.
所以明明想罵的時候.
都收回運火.
又哭又哄.
因為我們希望我們沒有快樂的童年.
兒子你可以有.
但是這樣是可以的嗎.
我們很怕小朋友哭.
很怕小朋友不高興.
但是不斷遷就.
他們是否高興呢.
當他們沒有正面的教導.
不能懂得去做人.
又不懂得社交.
你保得到他在家開心.
保不到他上學開心.
因此那個困難.
兒子應該受的裝備與訓練.
是越來越少.
所以無論多大都不懂得說話.
疫情之後.
說話和社交是一個很大的課題.
十多年前.
都不知道誰那麼成功.
就推出了一個.
拍孩子輸在起跑線的廣告.

$^{561}$於是又是這樣.
突然間全世界就在這個氛圍裡面.
到底我們的理性是從何而來.
我們怕他輸在終點.
你為何怕他輸在起跑線.
他贏在起跑線.
不過他輸在終點.
是我們想要的結局嗎.
但是沒有人提出.
你問的時候都不對.
你問的時候都會被人說.
我們這二十年.
就是在這樣的氛圍.
疼愛小朋友.
很疼愛.
但是為他將來設想.
於是不斷加班.
兼職傾家蕩產.
都是為了他們有最好的教育.
他們能夠追得上隔離的家庭.
可以看到這個世界.
小時候我就可以去哈佛交流一下.
然後去不同的地方旅遊.
開闊他們的眼界.
說真的.
六歲以下去了哪裡.
你問問你的子女.
他真的不知道去了哪裡.
聽姐妹,不要怪你的媳婦.
不要怪你的女兒.
不跟你的套.
因為你想想.
這兩代的氛圍.
真的是兩個不同的世界.
我們沒有辦法說.
哪個叫好與不好.
各有各好.
但是如果我們不能夠.
好像經文這樣說.
我們要審查.

$^{601}$甚麼是上帝的旨意的時候.
我們這個世界.
永遠都是令我們一頭煙.
現在你想吃那個維他命甚麼.
對你有幫助.
其實可能不夠兩年後.
又說吃了會有事.
現在不斷叫你吃蛋.
一人吃四隻蛋一日.
你還相信嗎.
以前吃四碗飯.
現在吃四隻蛋.
你還相信.
聽姐妹,真的要拆業.
我們有時也會想.
我們也要看看結果.
我很喜歡跟小朋友聊天.
很喜歡知道他們的事.
為甚麼.
因為我有幸在這裡三十多年.
我看著幾代在這裡.
所有的事可以驗證.
因為我們所有事都要看結果.
你問我.
我們這代的父母很疼愛小朋友.
是真的.
個個捏到自己大眼泡.
所以今天的媽媽.
有很多眼淚.
那些眼淚是有苦自己知.
說不出口.
但為甚麼我們的小朋友.
全部都是玻璃.
你明不明.
我們小時候.
爸爸媽媽哪裡用講道理.
打來就打來.
打完你還要恭恭敬敬.
但問題是.
我們的心非常強壯.

$^{641}$我們感恩就是.
上一代可能學識少.
我們的父母都比較嚴.
我都說我媽媽只看我的成績表.
最重要是有藍色字.
沒有紅字.
這樣就過關了.
因為他們不知道你是甚麼科.
最重要是沒有紅字.
不要不合格.
因此我們經常說.
我有個同學三十六分.
他回去就拿藍色研究筆.
很輕輕地填.
填到變成藍色.
媽媽就過關了.
他就用砂膠輕輕地擦.
擦完之後.
寫一個紅色.
他就可以交給老師.
瞞天過海.
但弟兄姊妹你們太厲害了.
你們真是厲害.
厲害的父母.
怎麼辦.
你的小朋友有沒有一條罅.
可以靜靜地.
說話.
讓他們自己.
舒服一點.
你明白他們為甚麼這麼大壓力.
他們不是想精神病.
但因為他們逃不過.
你們的法眼.
攝錄機.
你明白嗎.
現在個個都被監控.
所有小朋友.
有手機.
任何時間檢查到.

$^{681}$做好定位.
真的想找時間休息都不行.
而最重要.
有一樣東西.
但卻沒有機會.
讓他們溝通.
因為說.
你不會過.
你說的話對不對.
對的.
但我做不到.
我會偶爾失手.
你問我為甚麼.
解不開.
因為我不是刻意失手.
因此在很多.
父母認為自己很用心.
去教導他們.
但同樣期望都很高.
因此有很多小朋友.
的確在很多情緒裡.
原因就是.
面對新一代的父母.
你們學識高.
你們能力高.
但你們要求更高.
但你們很忙.
你們不是一天到晚在家.
他需要幫助需要求救的時候.
你沒有時間理他.
你問為甚麼.
解不開.
不懂做.
往往我們要求高.
但沒有執手教.
這樣永遠.
都是做錯.
這樣做又錯.
這樣做都錯.
怎會有自信呢.

$^{721}$怎會覺得.
自己叻呢.
只會覺得自己蠢.
然後最慘是.
家庭本來是一個.
我們可以休息.
可以休息的地方.
但因為太重視成績的時候.
所有時間都.
趕趕趕趕.
不夠不夠不夠.
於是我們家裡從來都沒有歡樂.
很多人連一起.
吃飯的時間都沒有.
一家人都是輪流.
誰回來就誰吃.
小朋友就一天到晚趕功課.
補習補習補習.
我們下面的小朋友.
一個星期真的.
課外活動加補習.
八課七課差不多.
還是上全日.
早知道上半日.
上完全日還這樣.
為什麼沒有人喜歡讀書.
家裡是一個.
令大家開心的地方.
可以放鬆的地方.
以前我們教會.
星期日下午.
通常都帶小朋友出去玩.
甚至我們搞親子活動.
都是星期日下午.
一起玩.
現在的家長.
還拜完吃飯.
全家回去睡.
起床4點多.
大家看看手機.

$^{761}$就可以吃飯.
其實真的歡樂的時間不是那麼多.
如果我們一直.
被世界這樣恐嚇我們.
我們家庭一直這樣.
我們夫妻生活.
連拍拖的機會都沒有.
完全不能夠有任何多餘的時間.
浪費機會.
你想想.
哪裡講愛呢.
所有.
沒有時間去愛.
沒有空間去愛.
因為.
連感動的機會都沒有.
聽不到 看不到.
明明覺得很慘.
做到這麼晚.
很快就說明天要交功課.
不可以讓老師覺得他不好.
接著我們又繼續逼迫.
鼎姐妹.
你有沒有想過.
你的小朋友.
或者現在這一代的小朋友.
為什麼他們那麼容易輕生.
你試一下.
你和他一樣.
的時間表.
看看你捱不捱到一個月.
可能你會死得快過他們.
鼎姐妹.
現在這一代真的很慘.
我們照顧女兒的年代.
在青衣的年代.
也不是可以去旅行的年代.
其實我們都很遲才懂得去旅行.
但是我們通常和女兒做什麼.
你問我和大家一樣.

$^{801}$我們父母相織.
我們兩個還是做教務同工.
有事上來24小時會打給你.
但是我們很珍惜.
很守一個位置.
就是睡覺之前的半小時.
很重要.
小時候講故事.
上一代拿著書不會聽你講.
大家一起作故事.
最有名的就是.
肥狗的故事.
一隻肥狗由創世記看到.
就用肥狗來講.
一直都是這樣.
當然你不需要講故事.
但是那段期間已經習慣了.
講一下生活.
遇到什麼事.
最重要的一件事.
為什麼我們搬到教會樓上.
我們就是為了那頓飯.
六點鐘我們就開飯.
那頓飯很重要.
上去吃飯.
很忙聊聊天.
說家裡的事.
沒什麼特別的議題.
輕輕鬆鬆地過.
但是每一天我們都知道.
能夠有時間.
我們沒什麼錢可以花.
所以一定是在樓下.
那時候樓下有很多公園.
或者是鋼架.
我很喜歡捧著女兒去爬鋼架.
因為女孩子都要鍛鍊一下.
當我現在閒談.
和我們年輕人分享.
為什麼回家睡.

$^{841}$不如去公園玩.
接著回應就看到.
這四個字是常常出現的.
師母 時代不同.
講來講去不夠我說.
師母 時代不同.
這些千秋已經拆毀了.
弟兄姊妹.
我今天很想和大家說.
我們真的是時代不同.
所以不等於我們過往用的.
可以是適合現在這個年代.
但是我很想警告我們每一位.
其實什麼年紀都走不去.
你以為你可以走得去嗎.
接下來的時代.
是真的不同.
而且是快得很厲害.
又有一家有電視.
到家家有電視.
都十幾年了.
但是你想一下.
用手機.
跟著智能的手機.
我們還想著.
自己不需要新媒體.
接著我們學新媒體.
我們預先學的教會.
我們才在疫情的時候.
突然間大家可以上網.
家樂木區那些.
都可以上網.
你說算了只是回教會.
我不管了.
我有個八達通到處走就算了.
一下開放上回家玩.
我已經懂得任何的電子支付.
為了去玩.
你明不明白.
短短幾年之間.

$^{881}$我們上去玩都是一兩年的事.
馬上就學會了.
等於姐妹.
這件事很快會出現.
科技是你控制不到.
這樣去的.
你以為我們還可以用父母的威嚴.
用父母的知識.
來鎮壓下一代.
沒有了.
你想想無論講知識.
你不夠AI厲害.
一按就知道.
你都未必可以即答.
第二個問題.
我們可以和他玩得很開心.
由他拿著奶嘴.
按ipad的時候.
他們已經知道和他玩.
比和你們玩好.
何況現在.
你想和他打機.
我們和他一起打機.
不好意思.
你不會打他的機.
語言你都不知道他在說什麼.
你說不是的.
我們還是擁抱和親吻.
和他聊天諸如此類.
不好意思.
現在他的什麼水.
24小時on call.
還有一樣他不會發瘋.
他不會罵人.
他說話絕對不會傷人.
你有沒有看過他們的回答.
多溫柔.
接待員都不夠他.
所以被解僱.
他們開口不懂回答.

$^{921}$總之回答你.
氣氛很順.
很受尊重.
你想想.
我們還在想老套的.
和他講故事.
和他去公園玩.
說得粗俗一點.
你省一點.
其實你的小朋友已經不需要你.
往後的日子我們真的不留意.
你想想.
夫妻關係.
在一間屋他不需要有小三.
他的小三在這裡.
你明白嗎.
他們多溫柔.
對老公說話.
然後我們朋友之間都會被蒙騙.
我們下一代的小組.
他說他們只需要一個朋友.
其實很坦白將來他們都是這樣說.
但他只需要一個在這裡的朋友.
他連眼鏡都不需要.
不用約出來那麼麻煩.
都不用理會他吃什麼.
弟兄姊妹.
事實上我們現在講的.
父母的問題.
真的很嚴峻.
不要還在舊有的框框裡面.
無論是哪個時代.
我們都承認.
我們是有框框的.
我們不是將所有東西都扔掉.
但是要心意更新.
而變化.
你做阿婆做阿媽.
你怎樣才可以和孫子.
去接觸一個家庭.

$^{961}$一個家族怎樣維繫.
以前我們媽媽那一代.
八十歲都煮十幾人飯.
現在搞不定.
怎樣可以維繫一個家庭.
這些我們真的要去想.
或者不斷地.
將人心奪去.
不單止.
是因為去外國讀書.
去其他地方移民.
就算坐在你旁邊.
他的靈魂都.
進入了這裡.
你會發現現在很多小朋友.
都說.
你是我最好的朋友.
但他的朋友.
是在這裡.
是否我們很害怕.
還在想要否拒絕AI.
AI是否666.
那些問題.
親密的弟兄姊妹.
我們雖然走得很前.
但我們要繼續向前.
我都明白.
我都和牧詩歌說了.
連PPT都沒有.
連PowerPoint都沒有.
還有膠片的時候.
你怎樣和牧詩歌說AI.
弟兄姊妹 時代是很快.
我們怎樣都要去.
在這個潮流裡保住.
我們自己的家庭.
讓我們的小朋友.
每個靈魂攝了上去.
真的在這個地上.
成為主可以.

$^{1001}$因待和使用的人.
所有新的都會令我們驚恐.
但我很想說.
人類發明出來的工具.
是我們可以善用的.
但這個課題是需要我們願意更新.
我們願意變化.
我們需要去學習.
如果不是的時候.
我們就會和下一代.
完全失聯.
即是說.
他和你是完全無話可說.
無話可說.
因為.
媽媽.
解一輪.
你都是不明白的.
其實是他們鄙視我們.
低能.
不是我們鄙視他們低能.
弟兄姊妹 是難的.
因為一代又一代的路.
這條當走的路.
做了人越老的我們就明白.
他那條路我沒有走過.
我們都不知道.
靠什麼去給他們意見.
不過你要知道.
我們不需要害怕.
他們也是上帝創造.
我們是同樣的.
有一個肉體.
我們有一個情緒的魂.
我們會有一個靈魂在.
無論一切.
做出來的機械人.
就算他有溫度.
暖暖地也好.
他始終不同.

$^{1041}$抱著父母那種感覺.
所以我們珍惜我們的子女.
能夠親近的時候.
去努力.
建立好的關係.
以致我們的全程達意.
不是要查機.
而是我們本身的一個眼神.
一個擁抱.
就已經勝過萬語千言.
第二 我想鼓勵.
AI可以給你很多知識.
但是他給不到榜樣.
我們生在他前面的幾十年.
成為他的父母.
成為他的婆婆 奶奶.
我們做了長輩那些.
其實我們的生命.
就是機械人取代不到.
我們走過的路.
我們怎樣成功.
我們怎樣失敗.
這些對小朋友來說.
其實是很重要的.
所以爸爸媽媽我鼓勵你.
不要經常說你小時候多乖多聰明.
你這樣是得不到他的心.
你不如說說你小時候多頑皮.
讀書多差.
他馬上覺得你很可愛.
有你才明白我.
我們要給他一個榜樣.
我們小時候都很頑皮.
可能我們小時候都很差.
讀書都差.
但是為什麼我們今天會有一個好的工作.
為什麼我們可以養起一個家.
這些是機械人給不到他的.
因為這個是生命.
而最重要的一件事.

$^{1081}$就是機械人永遠都給不到他的.
就是我有父母撐腰.
那種骨氣.
這個世界有很多衝擊.
他有他的路.
他永遠知道.
他有一個家可以貴.
他有兩個很愛他.
無條件愛他的父母.
在他後面撐腰.
當這樣想的時候.
你發覺在人生路上.
除了有父母.
誰可以給他無條件的陪伴.
守護和鼓勵.
我們當中.
60歲以上的弟兄姊妹.
我鼓勵你.
我們真的努力了一生去經營我們的家庭.
是時候我們真的要.
好好珍惜.
我們還有機會在後面.
做啦啦隊的時間.
他們自己向前走.
追不到的了.
但是我們能夠在旁邊.
不斷的鼓勵他們.
為他們禱告.
這個是我們今天仍然可以做得到.
現在是聖靈的時代.
我們每一個都能夠聽到神的聲音.
神是那個智慧者.
祂會幫助我們.
能夠醒覺.
看到世代的邪惡.
我們怎麼做呢.
不斷的去尋求.
有一個謙卑的心.
知道我們那一套.
未必受用.

$^{1121}$但是我們敬其愛主的.
Foundation.
是千古不變的.
求上帝指引我們.
懂得怎樣來成為.
下一代的祝福.
在應用上其實我沒什麼要說的.
都是陳腔舊調.
因為我只是做這些.
進入你的內室.
記得貼上你孩子的名字.
我們要行切的.
提名你去禱告.
行切是很重要的.
有誰是可以.
每一天都做到這件事.
你的子女就無福了.
聖經裡面說.
只有兩三個人.
一起同心合意的祈求.
無論祈什麼.
都會為他們成全.
所以War Room是不夠的.
War Room是我們自己和上帝.
Fighting的地方.
但是我們要抓緊這個應許.
我們有兩夫妻.
都是信主的.
你明白我的意思嗎.
只要你們夫婦同心.
提名去為你的孩子禱告.
你不要管他有回教會沒回教會.
他做過什麼事都好.
你們兩個不放棄他的時候.
你所求的事.
是會成全的.
我們單身的.
或是沒有子女的.
家庭很尷尬.
我們做守望小組.

$^{1161}$12個人一起祈禱.
說話太浪費時間.
我們找兩三個一起就夠了.
這個抓緊上帝的應許.
當我們兩三個.
彼此同心的為對方的家庭.
對方的需要去禱告的時候.
神說他就在我們當中.
他就會應允我們的禱告.
弟兄姊妹當我們禱告.
不是無奈的禱告.
我們的禱告是我們的能力的源頭.
而且事情是會成就的.
今天我們在母親節的時候.
我也想鼓勵我們教的弟兄姊妹.
我們很多一家一主.
一個家族都在.
我們年老的.
跑不了多少.
但是在這個禱告的守望.
是我們可以做的.
很多在做爸爸媽媽.
的路上.
真的捱了很多年.
看到現在的狀況.
你還在想.
自己當初做得好不好.
但是弟兄姊妹們.
爸爸媽媽們.
我鼓勵你不要再被那惡者.
牽引你走.
事情已經過去.
我們已經將我們最好的擺上.
我們應該來到上帝的面前.
來領受.
上帝給我們最後的使命.
就是.
好好地為你的家庭.
去禱告.
讓上帝親自來介入.

$^{1201}$我們努力做.
都做得不好的時候.
上帝來幫我們收尾.
只可以這樣說.
你信不信得過我們上帝會愛你.
你信不信得過上帝會聽你的禱告.
當我們信的時候.
我們就行動.
讓我們一起起來.
我們來到上帝的面前.
安靜自己.
天父我們懇求你.
親自來祝福我們.
讓我們能夠聽到你的話語.
讓我們來到你的面前.
去靠近你.
可能我們做爸爸媽媽的日子.
已經過去.
我們的子女已經成人.
求你親自來.
帶領我們整個家庭.
我們做得不好的時候.
父阿求你親自.
為我們去修補.
所以我們奉耶穌的名.
做一些禱告的服事.
弟兄姊妹如果你真的很想為你的家.
為你的孩子去禱告.
無論你是男是女.
我都鼓勵你來到台前.
舉起你的手去接受禱告的祝福.
讓我們每一個的弟兄姊妹.
都能夠向主呼求.
求上帝的恩典臨在我們裡面.
所以爸爸媽媽們.
可能做媽媽的,爺爺的.
婆婆的,公公的.
我鼓勵你為你的孩子.
為你的下一代來禱告.
「神,求你在我們當中,臨在我們裡面,與我們同在.」.

\newpage



\section{使徒行傳 2:42-47}
\label{sec:032112}
\textbf{教會增長的迷思 The Growing Church}~~~~鄺偉志牧師
\newline
\newline
連結: \href{http://www.church.com.hk/acms/content.asp?site=cdc\&op=show\&type=product\&code=032112\&layout=sermon}{\texttt{032112}} ~~~~ 語音日期: 2026-05-10
\newline
\newline
\hyperref[sec:032111]{\small{< < < PREV SERMON < < <}}
~
\hyperref[sec:index_chronic]{\small{[返順時目]}}
~
\hyperref[sec:032113]{\small{> > > NEXT SERMON > > >}}
\newline
\newline
使徒行傳 2:42-47
\newline
\begin{longtable}{cl}
\hline
\hline
章節 & 經文 (和合本修訂版)\\
\hline
2:42 & \begin{tabularx}{0.7\textwidth}{X} 他們都專注於使徒的教導和彼此的團契，擘餅和祈禱。 \end{tabularx} \\ \\ \relax
2:43 & \begin{tabularx}{0.7\textwidth}{X} 眾人都心存敬畏；使徒們又行了許多奇事神蹟。 \end{tabularx} \\ \\ \relax
2:44 & \begin{tabularx}{0.7\textwidth}{X} 信的人都聚在一處，凡物公用， \end{tabularx} \\ \\ \relax
2:45 & \begin{tabularx}{0.7\textwidth}{X} 又賣了田產和家業，照每一個人所需要的分給他們。 \end{tabularx} \\ \\ \relax
2:46 & \begin{tabularx}{0.7\textwidth}{X} 他們天天同心合意恆切地在聖殿裡敬拜，且在家中擘餅，存著歡喜坦誠的心用飯， \end{tabularx} \\ \\ \relax
2:47 & \begin{tabularx}{0.7\textwidth}{X} 讚美神，得全體百姓的喜愛。主將得救的人天天加給他們。 \end{tabularx} \\ \\
[1ex]
\hline
\hline
\end{longtable}
$^{1}$各位兄弟姐妹早晨.
有一個簡單的祈禱.
天父將兄弟姐妹交給主.
將孩子交給主.
讓我們有機會一起崇拜.
透過你的聖言的講解.
以及聆聽.
我們一同思想你的話語.
不單理性的分析.
用你的聖靈來引導我們進入一切的真理.
幫助我們不是頭腦明白.
更加我們從心裡面領受.
主及我們願意緊守遵行我們領受的道.
所以這個工作絕對不是只能說的.
主啊.
你自己的聖靈去幫助我們.
就好像初期教會一樣.
灌輔在我們當中.
以致我們一起聆聽.
包括孩子去準備這個講道.
都是一起聽這個道的時候.
來領受你的話語.
以致我們心裡面得著激勵.
得著安慰.
得著幫助.
也特別慰問我們今天母親節.
有很多母親我們一起慶祝.
主但願你的愛在我們當中.
也欣賞母親節.
再一次提醒我們.
主你賜給我們偉大的母親.
我們有很多人自己都是母親.
讓我們看見一切都是源於主你自己.
主我們將榮耀都歸給你.
奉耶穌基督的名字.
阿們.
主要是以教員身體的語言.
或者身體的表達.
來砌一個你理想中的教會是怎樣的.
不知道大家有沒有玩過這個遊戲.

$^{41}$或者可能你玩過一種簡單版.
或者小朋友版.
就是你拍張紙.
你畫出一個你心目中的教會.
一個理想的教會是怎樣的.
你可以想像到.
其實那些人砌了些什麼出來呢.
比較深刻有幾款.
其中一款就是.
他們當有十個人.
每個人有不同的崗位.
當然有一個人是講道的.
有一個人是在祈禱的.
有一個人是在關懷人的.
全部都是硬招來的.
砌了身體語言.
有些人是傳呼音的.
這個是不是一個理想的教會呢.
另外一種的款式就是.
有一個講台.
很深刻的印象.
有一個人是講道的.
然後有很多人像你們這樣.
但你們不是看的時候.
有時候有點呆板.
不是的.
起碼扮到.
希望大家聽到的時候也是這樣.
扮到很留心的樣子.
有些人在寫筆記.
有些人讚美主人.
有一個這樣的硬招.
做了出來.
我很深刻的印象.
這間是不是理想的教會呢.
有一款很深刻的.
就是它是一種抽象式的.
怎樣抽象式呢.
它不是做一些不同的工作.
在這裡一個硬招表達出來.

$^{81}$而是通常為一個什麼形狀的東西.
圓形.
然後有很多款式.
有些是手牽手.
有些手牽手也不夠.
是怎樣呢.
手扣手.
就是很緊密的.
一個圓形.
做了一個很統一的.
好像我們常常唱歌.
有一個機構叫做同心圓.
我不是說要推銷它.
但是是一種這樣的表達.
又是很好.
就是說理想的教會.
不是只是說要做些什麼.
而是說是一種合一.
同一的一種方向和感覺.
從這個開場白.
的確和我們今天的經文很有關係.
我首先說一說.
這張圖看了兩秒.
準備了.
其實為什麼有一個英文的呢.
我給大家看一下.
其實我們很感恩的.
謝牧師帶著我們同工隊伍.
三語.
上次Pastor Candy也說過.
三語我們都是跟這個題目.
有五次的小徒行傳的講解.
我是第二次的.
所以其實這個是謝牧師給的題目.
The Growing Church.
就是這兩段的經文.
上次就說了上面的經文.
接著那幾個星期就有不同的經文.
是一起去在這個舅堂.
在這個經歷裡面.

$^{121}$我們去思想從小徒行傳.
從爾忽所書和菲勒比書.
怎樣在教會裡面.
同心合一.
去走.
我是負責The Growing Church.
但我給的題目是教會增長的迷思.
究竟教會增長是一回事呢.
所以我們從這裡去看.
今天的經文.
多謝主席讀出.
有兩段的.
我集中都是說第一段的.
當然兩段是絕對有關.
尤其是第二段給大家看一下.
第二段.
你會看到是第二章和第四章.
如果我們看回前文後理.
當然第一段.
剛剛才說的那一段.
一會兒會說多一點的那一段.
是第二章.
第二章就是小行傳.
聖靈降臨之後.
後來彼得就講道.
接著就有三千人信主.
所以我特別要寫多一節出來.
本身就沒有包括這一節.
但是你看到背景就是.
四十一節領袖的人.
彼得的說話就受洗.
接著一天就門徒多了三千人.
是這樣的背景.
接著就形容教會是一個怎樣的狀態.
接著第四章其實是很類似的.
你會看到.
也是有一個場景聚集.
繼續是教會的群體.
禱告完了.
聚會的地方震動.

$^{161}$接著被聖靈充滿.
講神的道.
很多人就一心一意.
還有凡物公用.
接著就說了.
主要是描述怎樣公用.
怎樣凡物公用.
各人將所需分給各人.
特別highlight了一個使徒.
叫做巴拿巴.
是有錢人.
但是他做了信徒之後.
做了使徒.
後來他就成為了一個使徒.
應該說是做了使徒.
他的名字是勸慰主.
但是他特別出名的.
就是他將田地賣了.
接著就給教會.
也是一個很大的見證.
所以這一段其實就是.
你可以說在史倫傳裡面.
是不斷的.
尤其是頭八九章.
不斷會形容教會的情況.
是怎樣的情況.
是很合一的.
甚至如果說回前面那一段.
第四十七節.
主將得救的人數天天加給他們.
第四章那個就是.
又是有一群人信主.
第二章又是一群人信主.
接著其實有一個描述.
在史倫傳有第五章.
第六章.
第九章都有出現的.
主將得救的人數天天加給他們.
用這樣的表達.
就是說教會是不斷增長.

$^{201}$所以說教會增長是對的.
也都是史倫傳說到初期的.
教會是不斷的增長.
增長是什麼意思呢.
很清楚你看英文.
你看回原文也好.
就是人數很明顯.
就是天天增加.
不是說你身量肥了.
或者大家的愛心增長.
不是.
是人數增長的.
是很明顯的.
所以我們回去.
看回這一段的經文.
你會發覺到.
究竟他想說什麼.
當然我們知道.
這是一個歷史的事實.
是發生過的.
甚至有很多人討論.
究竟這段經文.
是想表達什麼的時候.
這個是不是共產主義.
你有沒有聽過有人讀這段經文.
去說這些.
的確有一些解釋.
是會將這件事連上去.
其實根源在哪裡呢.
當然不要時空錯置.
去解釋聖經.
共產主義正式在近代.
共產主義其實就是.
在這段經文很久之後.
一直發展.
所以倒過來應該說.
其實會不會.
這段經文.
這些歷史的事實.
是在影響後來.

$^{241}$共產主義的思潮.
你可以這麼說.
我們就不可以說.
這個其實就是共產主義.
不是後來的事.
這是我們解釋聖經.
或者去解釋一些事的時候.
要留意的地方.
究竟.
這個經文究竟想表達什麼呢.
你看到.
我highlight了幾個紅色的字眼.
是彼此交接.
彼此交接就是團契的意思.
交接其實就是團契.
我們很熟悉團契的意思.
凡物公用.
不用特別解釋了.
同心合意.
很特別的.
特別說到你的mind.
第二個就是一心一意.
heart和那個字眼是soul.
就是靈魂.
一心一意.
大家公用.
你發覺這幾個字眼.
我們首先說背景.
其實我們會發覺.
有一個影響.
其實應該.
其中很大的影響就是.
當時的希臘哲學家.
的一些說法.
在這段經文裡面.
我們看一看.
我曾經說過很多次.
有些經文.
如果你cross reference.
去看.

$^{281}$cross reference怎麼說.
你去參照其他的經文.
或者參照其他當代的文學.
或者是一些他們的思想.
如果有一段經文密集地.
有兩次都密密地.
但如果有三次以上.
我們會解釋.
他應該是受那個影響.
或者他是想透過.
這些表達方法.
去.
經文聖經的作者.
去連結當代的一些思想.
或者更新他.
或者要印證他.
所以我們不是亂說的.
我要證實這件事.
剛才highlight了幾個字眼.
一會兒會再出現.
其實.
一句到尾就是小行傳的經文.
這一段.
是歷史的事實.
但是他也是.
路加作者.
去表達一個教會的典範.
理想的教會是一個怎樣的教會.
和哲學家有什麼關係呢.
看看你認識多少個.
這個你猜猜是誰.
拿著三角形.
大家讀小學已經讀過了.
畢氏定理.
小學中學.
畢達哥拉斯.
大家記不記得.
畢氏定理.
三角形的.
這個數學家很厲害.

$^{321}$在主前.
五百多年的一個人物.
其實除了數學.
其實他.
是有一個做法.
他沒有寫過書.
但是他有一班門徒.
跟他一起生活.
你可以說很神秘.
強調怎樣去.
透過思想.
哲學.
數學.
一起去.
找一些精英一起去生活.
這個是他.
共同生活的一種.
model.
一種表達.
早在這麼久之前.
主前五百多六百年.
即是說.
距今二千多年.
就有這種.
畢達哥拉斯其實影響後世的哲學家.
跟著這個.
是誰.
帕拉圖.
二百年後.
The Republic.
這本書中文譯名叫理想國.
可能大家都聽過.
跟著.
他的弟子.
阿里士多德.
他又寫過很多不同的書.
其中又講道德論.
倫理等等.
這三個哲學家有什麼關係.
我們看看.

$^{361}$這些詞.
剛才說的那幾個詞.
其實你看史來傳.
第二章第四十二節.
凡物公用.
一來就是原來.
衰朔到好像.
共同生活的那群人.
到後來.
其實影響.
剛才說帕拉圖的理想國.
其實就直接用這個字眼.
來形容.
他的理想國.
當然.
你深入去研究去看.
有很多的東西.
但這個字眼.
凡物公用這個字眼.
當然我們不讀希臘文.
你會發覺到.
很特別.
常去形容這個理想的社會.
的一種情況.
就是說我們凡物公用.
當然那個表達方法.
如果看回帕拉圖.
應該就是.
我不是專家.
應該讀他的著作.
說出來就是.
他建議.
或者他鼓勵一些人.
當然他是在說精英階層.
就拿出你自己的財物.
然後共同使用.
你說有什麼特別呢.
後來.
亞里士多德就說.
他就不同意他老師.

$^{401}$帕拉圖的說法.
他說.
那些自己的物件.
應該有.
那個所謂.
私有的一個權利.
譬如我們現在都是這樣的.
我們的產業.
是我們自己私有的.
你說那你共用我的產業.
我全拿出來.
為什麼呢.
我做事這麼辛苦.
我賺這麼多錢.
不是說我不給出來.
但是我這樣共用出來那怎麼辦.
所以他不同意老師的.
他就建議一件事.
產權是私有的.
但是使用是共用的.
明白嗎.
就是說大家都可以保留你的東西.
但是.
我們一起用吧.
一個這樣的想法.
所以凡物共用.
你想想.
在士林傳寫的時候.
早在三四百年.
已經在說這個概念.
而當時.
仍然是希臘的哲學主導.
在耶穌的時代.
我們所謂希羅文化.
一早已經有這樣的想法.
而這種的說法.
就是說一個理想社會應該怎樣.
我們都是這樣將.
我們應該所有的.
就算阿里士多德的修正版.

$^{441}$你是擁有那些東西.
但是你都要一起去使用.
這個就是最理想的世界.
接著交接.
特別是說到.
阿里士多德.
是說到朋友之間.
那種團契.
交接.
所謂團契的意思.
原來不只是我們在聖經寫的.
聖經寫團契這個字.
交接這個字.
其實就是源於這班哲學家說的.
就是說怎樣才算團契.
就是我們不單止.
是一種朋友.
聊天.
而是心靈心意都是相通的.
這樣就是團契.
當然到了聖經的團契.
你看約翰一書也好.
或者你再研究這個字.
士林傳也有用過不少次.
這種的情況.
是有神.
和我們一起團契.
所以這個就是基督徒.
或者你可以說聖經的作者.
將團契這個字.
不只是朋友之間那種心靈.
而是說我們是有神的靈在我們當中.
所以我們一起.
這樣的團契.
同心合意也是.
他本身的字是一體的.
什麼叫同心合意.
就是一個身體.
甚至下面那個一心一意.
我不詳細說.

$^{481}$他在用什麼著作.
這兩個字.
尤其是後面那兩個字.
就是一個身體.
兩個靈魂.
什麼意思呢.
倒轉了.
兩個身體一個靈魂.
這個叫一心一意.
就是兩個方法.
第三個同心合意.
就是.
大家是一個身體.
但是一心一意.
就是兩個身體但是一個意志.
一個靈魂.
都是在表達一種的同一.
尤其是亞里士多德.
我讀到真的很好.
他說什麼叫朋友.
你想一下什麼叫朋友.
兩個朋友.
雖然我們兩個身體.
但是一個靈魂.
你和你老婆.
或者你最好的朋友.
你最愛的那個人.
但是一個的想法和靈魂.
當然有很多的解釋.
但是我想說.
也不是要拋書包.
事實上.
在當時社會有希臘哲學家.
這種理想國的.
這種概念在當中.
所以路加就用.
這些字眼來形容.
三個以上.
形容著.
士同人傳的教會是一個.

$^{521}$好像理想國的.
一個狀態.
一個群體.
原來是這樣.
當時當然也有猶太人的影響.
如果你聽過.
很快講一講.
但是不可以詳細解釋.
譬如說死海古卷.
大家知道聽過.
我們發現聖經有些古卷.
是在印證聖經的真實性.
而死海古卷.
那個群體.
多數我們會看到.
愛錫尼派.
而那些愛錫尼派也是共同生活的.
一起.
其實是有這些影響.
在士同人傳的背景.
講回士同人傳.
很colorful.
我想說什麼呢.
所以我們去讀這段聖經.
我們.
好像剛才講的那個例子.
我們其實是.
想一間心目中.
最理想的教會是怎樣.
或者我們.
看回原來士同人傳.
路加已經是用這個方式.
他就說其實一個教會.
理想的典範應該是怎樣呢.
就好像士同人傳.
那個第二章的教會.
所以有些教會.
曾經我也聽過.
我們要回到士同人傳.
第二章那個教會的模式.

$^{561}$如果這樣.
我們怎樣做.
一起販物公用.
是不是這樣.
是不是教會要走這條路呢.
這樣就增長了.
有些教會真的想這樣做.
你發覺其實很有難度.
當然其實它是一個歷史的事件.
也都我們.
如果從一個常理.
去推斷那個情況.
為什麼那裡有這麼大班的.
信徒.
一下三千人信主.
你想回頭.
為什麼會有五旬節那個聖靈降臨.
這麼多人.
是因為五旬節的時間.
就是二月節的五十天之後.
其實就是那些猶太人.
去當時的聖殿那裡.
去敬拜.
所以為什麼這麼多人聚集.
其實很多萬人.
因為如果看回資料.
其實不是每個人.
在耶路撒冷.
或者一個聖地之外.
那些猶太人很多散居在羅馬帝國.
不同地方.
有些人可能一生只去幾次.
一家大小.
或者是.
籌備了很久.
這樣才要很多錢.
去朝聖.
當時也是更加困難.
全部人都走在一起.
所以其實.

$^{601}$是很奇妙的一個事件.
就是.
神就在那個時間.
所有猶太人就聚集在那裡.
神在那個.
時間就在聖靈降臨.
以致有一個.
很大的復興.
以致有很多人就信主.
但是很多時候我們想到這裡.
就停在這裡.
看回經文繁物公用.
很開心.
繼續在.
教會那裡發展.
但其實真正的情況.
你猜維持多久.
你有沒有想過這個問題.
這班人可以聚集多久.
應該不是超大型教會.
這個超大型教會.
為什麼不存在.
你想想最初有三千人信.
接著有五千人信.
起碼有八千人.
但其實算南丁.
一開始是萬多人的教會.
為什麼沒有了.
正正因為它是過節.
之後它就會散開.
它又回到不同地方.
就是說在歷史上.
其實這個教會.
這樣的繁物公用情況.
是一段很短的時間.
最多幾個月.
為什麼會繁物公用.
你想想很合理.
因為那些過節的人.
最初他沒有想這麼久.

$^{641}$但有復興.
信主很開心.
當他幾個月.
最初想來兩個月.
誰知道三四個月.
慢慢不是每個人這麼多旅費.
又不是每個人住.
在那裡住.
那就要分享繁物公用.
其實合理的推斷.
它是一段時間.
都會沒有了.
但是又回來了.
路加就是說.
雖然是這樣.
但這個歷史事件.
這個初代教會.
就是讓我們看到一個教會的典範.
所以我們不是說學它的方式.
要繁物公用.
當然我一會再解釋.
我們學什麼呢.
就是那種質素.
或者教會.
如果回到《少林傳》第二章.
最初成立的時候.
應該有什麼元素.
我們去看看.
這裡.
輕輕一說.
你見到經文.
有一個平衡.
所謂交叉平衡.
第41節和47節.
有個藍色就是教會增長.
就是平衡.
你發覺平衡一直收下去.
有說到.
第42節.
去回應第46節.

$^{681}$說到他們的聚會.
的情況.
第43節,44節,45節.
就是說他們.
繁物公用等等.
綠色那裡.
43節和47節是說到.
其他人的反應.
其實聖經有很多.
平衡的題材.
他想強調一些東西.
他就會用猶太方式的題材.
有些研究說.
究竟.
這間教會有什麼特質.
會做什麼.
大家數一下.
第一.
一定是有教訓.
不是現在睡覺.
是有試圖的教訓.
教會的元素.
一定是有人.
來教導.
並且有人來學習.
這是很基本的東西.
但這是理想的教會.
一定要有的.
而且是先行的.
以神的道先行.
你想想我們教會.
彼此交接.
其實就是團契.
彼此在神裡面.
有團契.
你也要留意.
那些人.
在哪裡呢.
第46節.
他們一起聚集的時候.

$^{721}$聖靈降臨在聖殿.
但是他們還有一個地方.
叫做且在家中.
抹餅.
所以他們在家裡.
如果你說團契.
當然不只是一個形式.
是一種人與人之間.
和神之間.
團契.
一起共同的.
一個經歷.
有時候我們.
開始要想想.
你想想你的教會生活.
或者我們的教會生活.
是不是有這些呢.
有時候有些信徒.
真的只是回大聚會.
你發覺這裡有大聚會.
敬拜.
但是也在家中抹餅.
另外一件事就是抹餅.
但是抹餅.
不只是在現在.
我們五堂上星期也有聖餐.
不是這種形式.
是怎樣的.
我們叫做愛言.
他們其實很簡單.
他們也要吃飯.
但是他們一起吃飯的時候.
就有一個元素.
就像耶穌的吩咐.
我們要抹餅.
要透過餅和杯來紀念主.
所以他們吃飯是一種交往.
也是一種愛的紀念.
每次都有一個抹餅.
紀念耶穌基督.

$^{761}$在他們身上做的.
提醒我們是一個群體.
我們一起吃飯.
吃飯在我們中國人來說.
在很多民族來說.
是很重要的.
我們在吃飯裡面.
就代表我們的友誼.
你一會兒想一想.
什麼時候在哪裡吃飯.
不要漏了.
當然不是抹餅.
但是他們吃飯.
抹餅是一件事.
代表那種共同的生活.
是家人才一起吃飯.
所以這一種.
是教會.
當然有一種紀念主的心.
第四樣東西.
就是祈禱.
一定是有祈禱的.
先說完.
46節.
吃飯是怎麼吃的.
開開心心.
誠實的心用飯.
吃飯的心態也不是神迷苦面.
是一種愛言.
喜樂的一種吃飯.
我們檢討一下吃飯的經歷.
接著就是祈禱.
我們有沒有祈禱呢.
這不是個人祈禱.
這是一個教會群體的祈禱.
也不是說祈禱會.
而是一種祈禱的生活.
我們弟兄姊妹.
也會需要有幫助.
一起經歷.

$^{801}$所以小組一定會一起祈禱.
再說不是一個表面的行動.
而是說.
這是一個很重要的元素.
因為很多時候.
我最近和一個弟兄聊天.
他也說.
如果不是靠祈禱真的過不了.
這個是教會的特質.
一定有祈禱的.
好像上次我曾經提過.
講過主禱文.
祈禱生活是不好的.
或者和弟兄姊妹.
祈禱的交往.
那種團契.
如果我們想想.
很少和弟兄姊妹祈禱.
也很少為弟兄姊妹祈禱.
我們教會的生活.
就差了一點.
接著.
第43節.
使徒有很多神蹟奇事.
是不是我們教會也有神蹟奇事呢.
我這樣演繹.
當然看後面.
尤其初代教會有很多神蹟奇事.
但是我們不是靠神蹟奇事吃飯的.
但是有些教會.
就是這個問題.
拿了這個做一個範本.
有神蹟奇事所以我們要追求神蹟奇事.
不是的.
反而就是說一種熟練的經歷.
或者那些歧視.
其實我們走在一起.
共同去傳福音.
已經是一個神蹟.
我這樣演繹.

$^{841}$會不會我們一起的時候.
我們繼續去經歷神呢.
講見證.
不是講30年前你怎樣經歷神.
是我們一起的.
不斷神是活著的.
有很多的神蹟奇事.
不是說醫人.
接著他馬上.
不是那種.
也不是排除.
而是說我們在生活.
我們教會的群體裡面.
有沒有神的工作.
在我們當中.
如果不是也是有些問題.
另外就是.
凡物公用.
剛才補充了.
不是說我們要拿錢出來.
一起公用.
其實後來發展到.
在哥林多前書第八第九章.
保羅講了一個.
講法出來.
他叫做.
愛心的均平.
就是說我們每個人都有不同的.
賺錢能力或者有不同的財富.
但是我們信徒.
走在一起.
有人有需要.
有人可能突然之間有需要.
或者有些窮人在我們周邊.
我們是不是應該.
用一個愛心.
去供應不同的人.
有的人.
夠實力的.
經濟上或者能力上.

$^{881}$你能夠去給多一點.
你就給多一點.
一些不夠的人.
缺乏的時候就得到供應.
是一種這樣的情況.
多過硬要做.
所以也有人指出.
教會那種.
如果說共產跟其他共產.
不同就是.
那是主動性不是強制性.
不是說教會拿錢出來.
不是這樣的.
是怎麼走的.
是在講主動是在講愛心.
你沒有愛心.
我沒有愛心.
我不肯給錢出來.
我看到有需要也不給.
這是你的問題.
所以前書教奉獻.
你個人按著.
你的心.
去做.
你不用作難.
不用勉強.
你按你自己的心.
理想是怎樣.
理想就是這樣去共產.
神感動你.
我看到有需要的人.
我將我有餘的或有能力的.
去幫其他人.
其實我想我們教會都是走這一套.
是不是人人走呢.
當然不知道.
但這就是所謂繁脈共用.
也在當時在使團轉的情況.
特別有個情況.
所以大家一起分享.

$^{921}$但這個原則.
就是一種背後的愛心.
一種的愛心.
好了.
有一個很特別的地方.
有兩個.
他要描述.
當教會是這樣的時候.
我們在想教會是這樣的時候.
第四十七節.
應該說是第四十三.
和第四十七節.
其實有一個元素.
是很特別.
很忽略的.
就是教會不只是四面牆.
或者只是我們一群群體.
你怎樣屬靈也好.
你是有一個向外的.
所以他要說.
綠色那裡眾人都懼怕.
第四十七節.
眾民都喜愛.
你們有沒有留意到那個矛盾.
教會的存在.
其他的人是害怕.
還是開心.
你有沒有看到這個矛盾.
同一節經文這麼短的出現.
其實.
我們要問一個問題.
那些人害怕什麼.
出現的時候.
那些人害怕那些人.
可能是害怕神.
原來神這麼厲害.
有神之騎士.
那些人就害怕.
也可能是一種敬畏.
敬佩.

$^{961}$你們這個群體很厲害.
其實.
在四十七節有一個譯法.
在英文有譯作respect.
就是說尊重.
得眾民的喜愛.
是尊敬.
就是說其實我們教會的存在.
一個理想的教會.
對於你身邊的人.
是有一個影響力在的.
你會使得那些人.
不是不喜歡你.
如果我們教會在Cheshire.
在Lincoln那裡.
這班人真是瘟神.
搞到我們泊車又搞不定.
又吵吵鬧鬧.
當然有個別人.
有什麼感覺沒辦法.
不一定是關教會的事.
但是我們教會.
對於其他.
非教會的人.
或者暫時不是信徒的人.
有什麼影響呢.
他們會不會因為我們.
respect我們尊重我們.
尊敬我們.
敬畏我們背後的神呢.
這件事我們要去想.
除了我們這樣想.
我們自己自身.
在我們的工作崗位.
別人會不會尊敬你.
因為你是基督徒.
你根本沒有告訴別人你是基督徒.
會不會尊敬你.
會不會好像這個逆法.
是喜愛你.

$^{1001}$或者我們作為一個信徒.
之於其他的非信徒.
是一個怎樣的感覺呢.
一個理想的教會.
是會使得.
起碼你對眾人有影響.
當然後來有迫害.
會有很多反應.
但如果沒有事發生.
我們教會在這個地方.
在我們的工作崗位.
你是基督徒.
你存在與不存在是沒有分別的.
這個教會是廢的.
基督徒也是廢的.
最後要處理一件事.
第四節.
一次過有很多人信徒.
你可以說是大型報道會.
有些人就說.
大型報道會的根源就是這樣.
不過他是沒有.
一個很深入的計劃.
是神的工作.
有人來信主因為聚集了.
這個後話.
但第四十七節.
第五章第六章第八章.
都特別是說.
這個教會.
當這樣的狀態.
就會主將天天得救的人.
加給他們.
這個很有意思.
你有沒有想過一個問題.
有些教會.
他會這樣看.
教會不用做全副音.
為什麼?因為經文這樣說.
你教會如果是這樣.

$^{1041}$那麼神.
你看主動是誰?主.
神就會將得救的人.
加給你.
你不用做神會加人給你.
這個說法對不對?.
為什麼不對?.
很有意思.
甚至有些進化到.
不用做了.
我們祈禱就行了.
彼此交接就行了.
我們做得對.
主就天天將得救人加給你.
是不對的.
當然我們看到.
我們有時候.
有些教會增長.
最近我們.
都經常看這些書.
去研究教會怎麼增長.
有很多說法.
其中一樣就是說.
剛才那些去量度.
你有沒有.
這個教會是不是很健康.
將神的教導.
是不是很健康.
去祈禱或者交往等等.
讚美就是在說敬拜神.
教會健康自然就會增長.
這個是的或者有吸引力.
等於說.
為什麼有些人買apple.
其實我這部也是apple.
不是sell apple.
那他用不用sell.
他要.
不過有些人排隊買.
因為它的quality好.

$^{1081}$因為有它的吸引力.
當然那些反apple的.
sorry.
那些不買apple的最討厭apple.
但是為什麼那麼多人.
排隊去買.
因為它的質素好.
很多東西都是這樣.
那樣東西好.
就自然會來.
但是.
秘密就是.
主將得救的人.
加給他們.
其實你看英文就清楚一點.
主是將.
人數加給他們.
是講人數.
但得救的人就是be safe.
being safe.
原來是不斷有人被拯救.
第一方面就是說.
所有東西都是因為主的緣故.
我們做多少傳呼音.
都是因為主.
設立這個救恩.
我們能夠順主得救.
不是我們的工作.
也是主去保守這個教會.
使得教會不斷多人.
但是be safe的意思.
被拯救的意思.
被主拯救.
但是被主拯救.
當然這個強調點不是.
說你去傳呼音.
但是人是怎樣能夠得救.
你沒有人來傳呼音.
他怎樣能夠得救.
你說怎樣吸引他來.

$^{1121}$他來到也要傳呼音.
等於新人來.
我們教會.
總有一個機會知道他沒有信主.
帶他去信主.
就要傳呼音.
甚至我以前認識的教會.
每個弟兄姊妹都要傳呼音.
不是交數.
而是說.
每個人應該要傳呼音.
沒有什麼恩賜.
但不是.
因為教會的model.
如果有一間教會.
你會去讀聖經.
你會學習聖經.
聽教訓.
你會祈禱.
其中一件事傳呼音.
你是一定要做的.
無論你怎樣做.
你都要學.
這是我自己的經歷.
我是否有恩賜傳呼音.
我現在有恩賜傳呼音.
因為做基督徒有幾十年.
由零開始.
由很害怕開始.
到慢慢學.
但動力在哪裡.
因為我是信了主.
最後結束.
我用上次.
上星期.
有手針.
Nicky弟兄.
是不是在當中.
我用這個分享.
他分享.

$^{1161}$信主很有趣.
以前.
教會的傳道人.
那間教會很傳呼音.
被那件事吸引.
為什麼.
一個基督徒.
為什麼你們這麼喜歡傳呼音的教會.
傳道人說了一個.
很簡單的比喻.
基督徒傳呼音.
就好像魚在水裡游.
聽不聽得懂.
如果你是一個基督徒.
你自然就會傳呼音.
好像魚的生命.
就在水裡游.
我用這個放在這裡.
如果一個理想的教會.
是這樣.
怎樣增長.
怎樣去見證主.
其實你是一個基督徒的時候.
我們一個群體.
很自然就會同心合意.
不是說完美.
什麼都沒有不同的意見.
但是我們的生命本身.
有流露一些東西.
好像魚在水裡.
所以當然有很多時候.
尤其現代.
很多教會的問題.
或者有些人很不幸.
回到教會有很多不好的經歷.
以後不回教會了.
我信主但我不信教會.
但是.
這段經文告訴我們.
其實教會美好的生活.

$^{1201}$或者一個好像理想國的生活.
是不是沒有可能.
這是我們追求的目標.
一個典範.
我就會挑戰大家.
一起去想一想.
這個小人傳教會的表現.
我想回我自己.
作為參與教會的一個門徒.
我是一個基督徒.
我有什麼缺陷呢.
如果我不傳呼音.
如果我不祈禱.
如果我不喜歡聽聖經教訓.
我出了問題.
我的生命應該是這樣.
不是這樣.
我自己不是這樣.
我們教會不是這樣.
出了什麼問題.
我們就要一起去追求.
一起去求神的恩典去幫助我們.
我們有一個想像.
究竟我們怎樣可以在神的心意裡面.
做一個合神心意的教會.
如果是這樣.
教會自然會增長.
主張得救的人數加給他們.
當然我們會傳呼音.
包括在內.
我們一起祈禱.
讓我們一起思考這段經文.
願你幫助我們.
一起去檢視.
我們自己的屬靈生命和教會的生活.
讓我們體會到.
主你的恩典.
你的愛在我們當中.
以致我們能夠.
不是勉強做出來.

$^{1241}$而是從心裡面活出主你給我們的生命.
多謝主聽我禱告.
逢耶穌基督明智.
\newpage



\section{腓立比書 4:4-7}
\label{sec:032113}
\textbf{在主裡得滿足}~~~~張健庭牧師博士
\newline
\newline
連結: \href{http://www.church.com.hk/acms/content.asp?site=cdc\&op=show\&type=product\&code=032113\&layout=sermon}{\texttt{032113}} ~~~~ 語音日期: 2026-05-10
\newline
\newline
\hyperref[sec:032112]{\small{< < < PREV SERMON < < <}}
~
\hyperref[sec:index_chronic]{\small{[返順時目]}}
~
\hyperref[sec:032114]{\small{> > > NEXT SERMON > > >}}
\newline
\newline
腓立比書 4:4-7
\newline
\begin{longtable}{cl}
\hline
\hline
章節 & 經文 (和合本修訂版)\\
\hline
4:4 & \begin{tabularx}{0.7\textwidth}{X} 你們要靠主常常喜樂。我再說，你們要喜樂。 \end{tabularx} \\ \\ \relax
4:5 & \begin{tabularx}{0.7\textwidth}{X} 要讓眾人知道你們謙讓的心。主已經近了。 \end{tabularx} \\ \\ \relax
4:6 & \begin{tabularx}{0.7\textwidth}{X} 應當一無掛慮，只要凡事藉著禱告、祈求和感謝，將你們所要的告訴神。 \end{tabularx} \\ \\ \relax
4:7 & \begin{tabularx}{0.7\textwidth}{X} 神所賜那超越人所能了解的平安，必在基督耶穌裡，保守你們的心懷意念。 \end{tabularx} \\ \\
[1ex]
\hline
\hline
\end{longtable}
$^{1}$母親節快樂.
梁君彥議員.
陳偉業議員.
湯家驊議員.
陳偉業議員.
湯家驊議員.
陳偉業議員.
湯家驊議員.
陳偉業議員.
湯家驊議員.
陳偉業議員.
湯家驊議員.
陳偉業議員.
湯家驊議員.
陳偉業議員.
湯家驊議員.
\newpage



\section{以弗所書 6:1-4}
\label{sec:032114}
\textbf{當孝敬父母}~~~~王正傑牧師
\newline
\newline
連結: \href{http://www.church.com.hk/acms/content.asp?site=cdc\&op=show\&type=product\&code=032114\&layout=sermon}{\texttt{032114}} ~~~~ 語音日期: 2026-05-10
\newline
\newline
\hyperref[sec:032113]{\small{< < < PREV SERMON < < <}}
~
\hyperref[sec:index_chronic]{\small{[返順時目]}}
~
\hyperref[sec:032115]{\small{> > > NEXT SERMON > > >}}
\newline
\newline
以弗所書 6:1-4
\newline
\begin{longtable}{cl}
\hline
\hline
章節 & 經文 (和合本修訂版)\\
\hline
6:1 & \begin{tabularx}{0.7\textwidth}{X} 作兒女的，你們要在主裡聽從父母，這是理所當然的。 \end{tabularx} \\ \\ \relax
6:2 & \begin{tabularx}{0.7\textwidth}{X} 當孝敬父母，使你得福，在世長壽。這是第一條帶應許的誡命。 \end{tabularx} \\ \\ \relax
6:3 & \begin{tabularx}{0.7\textwidth}{X} 見上節 \end{tabularx} \\ \\ \relax
6:4 & \begin{tabularx}{0.7\textwidth}{X} 作父親的，你們不要激怒兒女，但要照著主的教導和勸戒養育他們。 \end{tabularx} \\ \\
[1ex]
\hline
\hline
\end{longtable}
$^{1}$各位姐妹平安.
今日是我們的敬親主日.
但亦都是母親節.
如果你覺得你旁邊的是母親的話.
你也跟她說一聲祝你母親節快樂.
今日我們看的一段經文是一段.
我們返教會.
如果已經有一段日子的話.
其實你會很熟悉的一段經文.
就在《以弗所書》的六章一至四節.
剛才大家都已經讀過了.
特別今日因為是我們的敬親主日.
所以我們要一起去學習的是.
當孝敬父母.
但今日我很想和大家去看.
再次去看這段經文的時候.
不是單單只是看這幾節的經文.
去理解究竟經文講什麼.
而是我們真的可以明白.
前面《以弗所書》.
保羅他已經告訴我們.
當我們信徒的生命改變了之後.
我們在神的裡面.
我們脫去舊人穿上新人.
已經是一個更新的生命的時候.
究竟我們做父母的.
做兒女的.
怎樣可以經歷到.
和得到孝敬父母的祝福.
所以今日我希望和大家一起.
這樣去看這段經文.
對於我們華人來說.
講孝順父母這個課題.
傳統的父母去想.
子女孝順父母.
其實都是在想什麼呢.
就是養兒防老.
等自己老了的時候.
子女孝順父母.
他們可以照顧自己.

$^{41}$傳統的父母.
就會這樣去想.
就是養兒防老.
但是近代的父母.
已經開始不會再想.
子女在自己老了的時候.
會照顧自己.
而是想子女養得起自己.
已經很好了.
已經很孝順了.
這是近代的父母才會這樣想.
但是去到新一代.
特別如果你有留意.
有關於娛樂新聞的時候.
近幾天比我們更加學會一件事.
就是新一代的父母.
最要擔心最要防備的.
就是那些子女.
變成一個狠老族.
大家知道什麼叫狠老族嗎.
狠老族就是那些.
二十到三十歲以上的人.
他們選擇由守候下.
不上班.
然後靠父母去供養.
那些就是狠老族.
所以新一代的父母可能就要想.
只要你不是一個狠老族.
已經是孝順了.
究竟當我們去看這段經文.
今天我們再次去思想.
在說孝敬父母的時候.
究竟即以法所書所說.
保羅究竟要告訴我們.
什麼是孝敬父母呢.
當我們去看這段經文.
我們知道是來自十誡的.
是十誡中的第五誡.
前面第一至到第四誡.
都是在說神和人的關係.

$^{81}$然後去到第五誡.
先說到人和人的關係.
但當說人和人的關係.
第一個誡命就已經是說父母和子女.
所以我們明白.
這個是代表著.
父母和子女是一個很重要的誡命.
對我們來說是一個很重要的提醒.
而這段誡命更加重要的是什麼呢.
就是說他們帶著應許.
所以這裡不單止是說一個道德的規範.
更加是說一條通往祝福和生命延續的道路.
所以我們有時候要小心的是.
很多時候我們看聖經.
我們看到很多的聖經有些教導.
我們都會覺得他只是教導我們.
如何在道德上有一些規範.
好像當我們看今天的這段經文.
我們就會想.
他就是在說父母和子女.
是這麼簡單.
我們會這樣去理解.
其實這裡告訴我們的是.
一條帶著應許的誡命的時候.
就是說我們能夠去遵守這條誡命的時候.
我們是能夠得到祝福.
所以這段誡命是關乎著.
神的子民整體的一個延續.
當人能夠可以活出孝敬父母的時候.
神的祝福就臨到那個家庭.
並且信仰就能夠得以一代傳承一代.
所以當我們要去理解.
今天這四節的經文.
我們要先回到前面去看.
如果大家還有印象.
上個禮拜我們說了什麼.
全部還給我了.
上個禮拜我們說了合一.
合一的建基是什麼.
建基就是我們認識了耶穌基督.

$^{121}$我們脫去舊人穿上新人.
我們有一個更新的生命.
以致我們能夠可以脫去我們很多的.
壞習慣或者壞的行為.
以致我們能夠成為一個新做的人.
一同去建立教會.
這個是我們上個禮拜大概說的東西.
我正在考慮一件事.
我正在考慮就是.
每個禮拜將我的講道.
放進去AI.
然後弄一幅圖.
去整合了我那一個禮拜的講道.
然後發給大家.
不過雖然我心裡面也有個答案.
我知道發了給大家.
大家都可能未必看.
不過我都希望我可以這樣做.
當我們說回我們要去孝敬父母.
我們看這四節的經文.
我們要看前面的.
前面就是在說.
保羅說我們行事為人.
要與蒙召的恩相稱.
保羅在這裡提醒我們.
家庭的關係不單止是我們個人的.
而是建立在整個的屬靈的生命上.
因為我們每一個都是蒙恩的罪人.
我們已經得救了.
我們可以脫去我們的舊我.
脫去我們的老我.
我們不再有我們舊的行為.
不再與神的生命隔絕.
以致我們可以穿上新人.
在一個更新的生命.
我們能夠活出神的形象.
以致我們在關係上.
在家庭的關係上.
我們是有恩慈有憐憫.
能夠可以饒恕彼此相待.

$^{161}$所以我們明白.
如果當我們今天看回.
所說這一部分的時候.
我們必須要知道.
我們也是要建基於什麼.
建基於我們脫去舊人穿上新人.
我們是一個更新的生命.
以致我們的家庭才能夠活出.
恩慈憐憫饒恕.
以致這樣彼此相待.
所以因此孝敬父母.
不單只是一個外在的行為.
而是說我們一群人.
我們一個新的生命.
一種自然的流露.
當一個人真正被福音更新的時候.
他在家庭的態度就不同了.
就會改變了.
就會從自我中心轉向.
有愛有尊重.
就與他繼教的一些權利.
轉向甘心去付出.
是樂意的去付出.
並且基督更加成為我們一個榜樣.
祂提醒著我們.
祂本有神的形象.
祂甘願降卑自己.
所以使我們的孝敬是怎樣.
是不會被迫的.
是因著我們敬畏基督而去甘心樂意.
因為我們生命更新.
而去甘心的去孝敬我們的父母.
而當一個家庭裡面.
我們每一個人都以基督為中心.
我們以神為我們的元帥.
為我們生命的主的時候.
我們去彼此馴服的時候.
就不再是一種被迫.
亦都沒有任何的壓力.
而是只是一種愛的表達.

$^{201}$在家庭裡面我們很需要怎樣.
有時是你又遷就我.
我又遷就你.
對不對?弟兄姊妹.
如果你又不遷就我.
我又不遷就你.
那個家庭會怎樣?.
哪裡過得了?對不對?.
所以我們明白.
當我們每一個人.
我們生命更新.
我們脫去舊人.
穿上新人.
我們以基督為我們中心的時候.
我們的生命.
在我們的家庭裡面.
我們就能夠彼此去馴服.
這是一種愛的表達.
我認識一個姐妹.
她在一個雙親離異的家庭長大.
父母在她很小的時候.
就已經分開了.
她是由阿婆去養大的.
她每天都很不開心地過她的生活.
所以她一直都很想快點長大.
就可以離開這個家.
當她真的差不多夠年紀的時候.
她就真的離開了這個家.
離家出走.
但她形容就是過著一些很糜爛的生活.
吸毒.
亂搞男女關係等等.
所以她那段是一段很糜爛的生活.
但很感恩的是.
在她人生她真的去認識了主耶穌.
她回到教會.
她從神的話語.
從弟兄姊妹身上.
她很感受到神那份愛.
所以她生命有很大的改變.

$^{241}$她也很認真地去聽教會的牧者.
去講道.
上主一學.
查經等等.
但她說.
在聖經裡面每一個教導.
她都接受.
但唯獨是有一個教導.
她是接受不到的.
那個就是饒恕.
因為她心裡想起的就是.
她從來都饒恕不到她的爸爸和媽媽.
而正在這個時間.
突然之間她爸爸再一次出現.
並且她爸爸已經患上癌症.
這個姐妹當她見到她爸爸的時候.
她很掙扎.
很困難.
她不知道怎樣和她爸爸相處.
但當她真的認真地去祈禱.
和牧師去聊天之後.
牧師去鼓勵她.
讓她憑著愛和她爸爸建立關係.
去饒恕她爸爸的時候.
她鼓起勇氣再次去找她爸爸.
並且照顧他.
她在這段時間真的經歷了一段.
饒恕的功課.
她陪伴了她爸爸走了最後一段路.
更加她在這段時間帶著她爸爸.
決絕相信耶穌.
這個姐妹最後的總結她說.
她很感恩.
她很感恩的是因為她的信耶穌.
她認識了主.
她的生命真的改變.
以至她真的有能力去饒恕.
所以姐妹.
今天當我們說家庭的關係.
我們要能夠可以有饒恕.

$^{281}$能夠有彼此的接納.
彼此的馴服.
是因為我們都相信了這位耶穌基督.
祂成為我們生命的救主.
我們脫去舊人.
穿上新人.
以致我們真的有能力去愛和被愛.
然後接著進一步.
今天看的這段經文.
保羅他說了.
兒女要在主裡面去聽從父母.
在主裡其實在這裡是一個很重要的關鍵.
這裡很講給我們聽.
就是在不違背信仰的前提下.
來去聽從我們的父母.
而這個也是合乎神的心意.
也是誠心錄中的一個意義.
其實在說什麼呢.
是因為當時對於那些爸爸.
他們去教養他們的兒女.
他們是很嚴苛的.
他們是可以去責打他們的兒女.
甚至他們是可以操他們兒女的生死大權.
所以在這裡保羅說.
當我們脫去舊人穿上新人的時候.
保羅就提醒當時做兒女的要點.
就是在不違背信仰的前提下.
來去馴服父母.
在這裡我們非常容易就可以理解.
就算旁邊的人叫你做一些壞事.
你都不會去做.
更何況就算是爸爸.
他叫我們去做一些不合乎聖經的事.
他教我們去做.
我們都不會去聽.
這個是我們很容易可以理解.
而聽從父母在這裡.
其實保羅更加要讓我們知道.
其實我們是在回應神的權柄.
特別是猶太人來說.

$^{321}$爸爸是負責將信仰傳承的那一位.
他們要教導他們的兒女去聖經.
所以在這裡保羅是要讓我們知道.
聽從父母其實是我們在回應神的權柄.
但更重要的就是.
這個不是一個消極的服從.
而是一個積極的學習.
在當中在整個過程.
我們就要去學習尊重.
學習馴服.
學習愛.
弟兄姊妹.
如果你的爸爸媽媽是很愛你.
我們要去學習尊重.
學習馴服.
學習愛是很容易的.
但相反.
如果我們的父母是很嚴苛的.
甚至可能對我們有過分的要求.
在養育的過程帶來很多傷害.
我們要去學習尊重.
學習馴服.
學習愛.
相對就很困難.
所以保羅在這裡讓我們看到.
當我們去救人串上新人.
我們要去聽從父母.
我們要去回應神的權柄.
縱使我們的父母.
可能在他的教導會有嚴苛.
有時他的教導可能是一個過分的要求.
我們在當中都要去學習尊重.
學習馴服.
學習愛.
這個品格正是信仰能否傳承的關鍵.
等於我們在說孝敬父母這件事.
很重要的一件事.
就是我們的信仰.
能否傳承給下一代.
甚至我們的子子孫孫.

$^{361}$特別當保羅說孝敬父母.
當我們去理解孝敬的時候.
孝敬在希伯來文有一個很重的意思.
所以如果你看英文的翻譯.
就是尊敬父母.
尊榮父母的意思.
但如果我們看聖經裡面.
未有很具體的提及.
究竟我們怎樣做才是孝敬.
但聖經就有提到.
父母和子女之間的關係.
有一些要做的事.
有一些不可以去做的事情.
包括下面的父母有幾節經文.
都有特別提到.
究竟我們要怎樣和父母相處.
子女的.
首先今天看的經文就說.
我們要聽從父母.
我們都知道.
因為我們已經脫去舊人.
穿上新人.
我們要從主裡面去聽從我們的父母.
第二個就是說.
不可以打父母.
我相信這個應該一般都沒有.
但如果你有打父母.
你不用現在舉手給我們看.
咒罵父母是不可以的.
我沒有經歷過被父母咒罵.
但我聽過.
這麼多人死不見你死.
不知道被父母咒罵的多少.
還有就是.
不可以輕慢父母.
我們不會隨便輕慢我們的父母.
還有一樣東西是供養我們的父母.
我相信無論多與少.
我們都要供養我們的父母.
特別主要就是說.

$^{401}$如果你只是拿錢去奉獻.
就說你沒有能力去供養父母.
這個是不合理的.
這是不正確的.
所以我們在聖經裡面.
找到一些經文告訴我們.
我們有什麼可以做.
有什麼不可以做.
然後進一步補下去.
更加告訴我們.
這條是帶著應許的一條誡命.
它告訴我們能夠使我們得福在世長壽.
甚至有另一個經文是在地長久.
突然之間看到這裡.
不知道你有沒有同感.
我一聽到.
使你得福在世長壽.
我就很害怕.
我其實不太想這麼長壽.
所以我很害怕.
但是其實這條經文不是在說.
是不是一個長壽.
意思不是在說.
你孝敬父母.
你會得到福.
福氣是什麼呢.
你可以長命一點.
我看到這裡還想起.
我以前年三十晚.
家裡有一個傳統.
或者是習俗.
或者是我的姐姐給我的要求.
就是說如果想爸爸媽媽長命一點.
年三十晚就不准睡.
你越夜睡.
爸爸媽媽就越長命.
我不知道有沒有突然之間.
我家裡有這樣的傳統習俗.
可能都是我的姐姐騙我的.
這條經文說的得福.

$^{441}$不是在世長壽.
是說你孝敬父母的時候.
文藥能夠傳承.
當你孝敬父母.
你去尊重你的父母.
你去聽從你的父母的時候.
你的信仰就可以傳承給下一代.
傳承了文藥.
帶來納藥的應許.
就是他們能夠以色列人.
世世代代都可以做神的子民.
能夠在神所賜的應許之地.
可以長久居住.
所以你明白為什麼保羅.
他不單說長壽.
更加說在地長久.
所以他說的是.
我們孝敬父母怎樣得到祝福.
就是當我們去謹守.
遵行神的誡命.
我們去孝敬父母的時候.
我們能夠將我們和神所納的藥.
我們的信仰.
能夠傳承.
我們可以得到神的祝福.
更加讓我們的子孫都怎樣.
可以蒙神的祝福.
各位姐妹.
今日信義代的信仰是我們一定要.
亦都很必須.
我們要去關心和緊張的一件事.
各位姐妹.
我們沒有人期望.
我們的下一代是不信主的.
我們更加不會想.
我們的下一代是不回教會的.
所以我們.
能夠可以去謹守.
遵行神的誡命.
我們去孝敬父母.

$^{481}$我們就能夠將.
這個信仰.
可以傳承下去.
讓我們的下一代.
我們一代一代.
都能夠可以經歷到神的祝福.
而這裡經文的核心.
不是說一個條文.
而是讓我們的父母.
能夠得到尊榮.
當父母得到尊榮.
他能夠感受到.
被尊重.
被愛.
家庭就怎樣.
穩固.
當家庭穩固.
信仰就有根基.
當信仰有根基.
下一代就能夠承接.
在主裡面的祝福.
所以這個應許.
它是帶有一個群體性.
就是當我們一個群體.
我們都能夠可以.
活出孝敬父母的時候.
整個信仰的群體.
我們的信仰都能夠.
可以得以延續.
所以各位姐妹.
今天我們只是在說.
孝敬父母.
不是只是對我們父母好.
而是對我們整個家庭好.
不單止我們的父母.
他們得尊榮.
更加讓到我們的家庭裡面.
能夠讓這個信仰.
可以傳承下去.
我不知道有多少弟兄姊妹知道.

$^{521}$我的女兒是有病的.
她曾經有一段時間.
進了醫院.
住了大概二十天的時候.
現在都一直保持有覆診.
在整個過程.
對我來說是很痛苦的.
我在那段時間.
學會一件事是什麼.
就叫做以淚洗面.
我住美孚.
美孚後面有一個公園.
我很喜歡晚上走進去.
你知道為什麼嗎.
因為走進公園.
你哭是沒有人看到的.
在整個過程裡面.
我學了一個很重要的功課.
我經常都問的就是.
我不能夠一生一世.
都去看著我的女兒.
但是我發現一件事.
如果我的女兒.
在她人生.
她能夠認識耶穌基督.
她有耶穌看著她.
就算有一天.
當我回到天家.
我走都走得放心.
所以弟兄姊妹.
當我們能夠.
可以在家庭裡面.
我們去實踐.
孝敬父母的時候.
我們能夠得著.
神給我們的祝福.
讓我們的信仰.
能夠一代傳承到下一代.
大濟祖牧師.
是第五代的宣教士.

$^{561}$如果你對宣教的事情.
是有關心有認識的話.
你都知道他是.
戴德新第五代的孫.
戴德新最經典的一句名句.
就是他說.
我若有千磅英金.
中國可以全數之取.
我若有千條性命.
絕對不留下一條不及中國.
不是為中國是為基督.
這個就是戴德新牧師.
他當年來華.
張福音傳給中國人的時間.
他所立下的誓言.
我們看到的就是.
在戴德新牧師他的家族裡面.
他不單止傳承了信仰.
他更加傳承了什麼.
傳承了他們的儀仗.
他們的呼召.
所以他們一代一代.
繼續做宣教士.
來到第五代大濟祖牧師.
他的兒子幫他改了個名.
叫做戴承曰.
傳承他們和神所立的約.
戴德新牧師他說.
他對聖經裡面的族譜.
他感受很深.
在他人生或者在他信仰路程裡面.
很快他見到自己沒有力量的時候.
他就和神說.
神是他的阿拉罕的神.
也是以色的神.
更加是亞國的神.
他說的這一段經文.
他說他最大的感受就是.
神不單止是這幾位.
阿拉罕以色亞國的神.

$^{601}$更加是他的家族都一樣.
是他的真真真祖父.
真真祖父.
真祖父.
他們的神.
所以他心裡面的聲音.
告訴他.
過去神在戴德新身上.
以及內地會裡面.
行的許多神蹟.
他說今天.
神依然會彰顯在他這個時代.
因為.
我們心盼我們能夠在家庭裡面.
去實踐孝敬父母.
最重要的是.
我們能夠得著.
這條大應許的誡命.
就是我們能夠得著祝福.
我們的信仰能夠一代.
傳承一代.
然後接著.
看回這段經文.
不單止只是對.
做子女說.
對做父親的.
我們相信也是對父母說.
都有一些提醒.
他這裡說.
不要惹兒女的氣.
避免.
那個表達就是.
避免去壓傷.
去否定.
或者過度要求他們.
然後要照著主的教導.
去勸戒去養育.
這個其實.
保羅是特別要提醒.
當時做父親的.

$^{641}$因為他們已經去救人穿上新人.
從前他們在家庭的舊約裡面.
父親是可以去扎打.
甚至他有一個權.
做父親是怎樣呢.
他是可以承認.
或者去接納.
那個孩子.
是不是合法.
可以生存.
才可以.
那個嬰兒.
最終可以生存.
即是說.
如果你生了一個兒子.
是有病的.
有問題的.
那個爸爸是可以決定.
如果他有問題.
我可以殺了他.
當時他們就是有這個權利.
但保羅就要去提醒他們.
當那些爸爸.
他們認識了耶穌基督.
去救人穿上新人.
他們的信仰就已經不同.
在教育.
養育兒女的時候.
都已經變了.
不同.
因為真的好像在前面.
我們說的是.
他們要活出.
憐憫,接納,彼此馴服.
所以在這裡.
保羅要去提醒這些做父親的.
不要惹兒女的氣.
不要壓傷他們.
不要否定他們.
不要過度的要求.

$^{681}$以致信仰能夠可以傳承.
但更重要的就是.
信仰的傳承.
不單只是教導.
更加是一個榜樣.
當我們說到.
做父母要教導兒女的時候.
新教這件事.
我們從來都不能夠逃避.
你說得多好聽.
你怎樣做.
是更加重要.
所以我們的新教.
就成為我們兒女的榜樣.
所以兒女是否真的能夠去孝敬.
往往是來自甚麼.
他們是否在家庭裡.
在父母當中.
他們經歷到被愛,被尊重.
我在網上找了.
兒女最不喜歡父母說的話.
頭五句是這樣的.
第一句是說.
你學一下.
你看看別人吧.
我不知道大家有沒有說過.
你看看旁邊那個強仔.
他多好.
第二句就是.
你沒用的.
這個可能是很生氣的時候.
就會這樣罵子女.
第三句是粗口.
不可以說的.
第四句就是.
你沒得救了.
最後那句.
這麼老套還說.
生一塊叉燒比生你好.
我不知道大家在成長中有沒有聽過.

$^{721}$做父母的.
當我們去教育我們的兒女.
縱使我們在教育的過程中.
是很不容易.
我們都要謹慎我們的說話.
因為我們的說話是帶著權柄.
亦都很有破壞力.
所以我們要小心我們的說話.
不要惹兒女的氣.
然後保羅說不要惹兒女的氣.
他在說不要觸動子女的怒氣.
這裡不是在說.
因為我們不可以激怒他們.
不要惹他們的氣.
所以現在有些家長.
變成另一種.
什麼都迎合子女.
遷就你了.
你想要什麼就給你.
這裡保羅說.
他要知道.
父親不要使用過份的權力.
引致兒女動怒.
甚至他們會犯罪.
所以當父母.
他能夠用主的方式去對待他的兒女.
能夠有真理.
有恩典的時候.
他的兒女.
就能夠得到祝福.
我們要小心的是.
我們在教育子女的時候.
很容易有真理.
但沒有恩典.
很多規矩.
很多原則.
我們對他很多的評估.
究竟有什麼做得好.
有什麼做得不好的事.
我們有很多真理.

$^{761}$但我們可能沒有恩典.
一做錯就鬧到他狗血淋頭.
但又相反.
有些是沒有真理 只有恩典.
沒有真理 只有恩典.
這些是什麼呢.
是溺愛.
什麼都是好的.
千錯萬錯都不是你的錯.
是別人的錯.
為什麼你這麼壞.
是因為別人做得不好.
為什麼你讀書讀得這麼差.
是因為你的老師教得不好.
我們只有恩典 沒有真理.
我們也幫不了我們的子女.
所以當我們要用主的方式.
來對待我們的孩子.
我們要有真理 也要有恩典.
有原則 也要有一個溫度.
我們的子女.
他們這樣的成長.
才能夠更容易地認識神.
也才更願意去尊榮他們的父母.
所以相反.
我們除了知道有什麼不要說.
兒女也有一些很喜歡父母說的話.
頭五位是這樣的.
第一句是.
「中了力就可以了」.
「我覺得你已經很厲害了」.
這個應該也說得出.
然後就說.
「有進步 繼續努力」.
第三句說.
「失敗不要緊」.
「最重要不要放棄」.
這個很像「你敢不敢試」的口號.
第四句是.
「你做什麼都好 我們一定支持你」.

$^{801}$第五句是.
「不要給自己太大壓力」.
這些可能是我們子女都喜歡說的話.
可能對很多做父母的.
特別年輕一代.
我們要去說一些鼓勵他們的話.
可能不知道習不習慣.
不過我們生一代的父母可能會習慣一點.
但如果有做父母你覺得很困難的時候.
其實我會鼓勵你可以練習.
可以練習.
我們有空看著快鏡.
我們練習一下去鼓勵一下我們的子女.
「盡了力了」.
「盡了力就可以了」.
「我覺得你已經很厲害了」.
跟你分享.
前兩天有一頓晚飯的時候.
我女兒分享了一件事.
分享了一件事是什麼呢.
那天有一堂課.
老師是代課的.
整堂的同學.
有些人在後面唱卡拉OK.
有些人在那邊跳舞.
有些人在那邊捉棋.
玩啤牌之類的東西.
他說唯獨是他和他的同學.
兩個人坐在那裡做數學.
講完之後.
我老婆和我兒子都問.
「你講這些幹什麼」.
就問他.
然後我女兒就在那邊等我說.
我唯有馬上說.
「為什麼整班同學有人在那邊唱卡拉OK」.
「有人在那邊跳舞」.
「有人在那邊玩」.
「但是只有你這麼乖」.
「和你的同學一起做數學」.

$^{841}$我女兒就不停地在那邊笑.
反而我是鼓勵你.
我們的兒女都需要鼓勵.
需要認同.
不要以為你的子女已經五六十歲.
你覺得不需要.
肯定不需要認同.
我可以很坦白地和大家分享.
我做過一些輔導的個案.
曾經有一個男士.
他去到五十多歲.
他在社會上很有成就.
但是在他人生裡.
他很不開心.
因為他的成長.
唯一缺乏的一句.
就是他爸爸對他的肯定和認同.
所以鄧小妹.
儘管子女有多大都好.
你仍然可以肯定和認同他們.
在猶太人的古典.
這個塔勒木.
很具體地教導了父母怎樣做一個榜樣.
有五方面.
第一樣就是要教導托拉.
即是律法書.
其實父母都要去學律法書.
因為父母的屬靈生命都要去突破.
都要去成長.
第二樣做父母的.
就是要帶兒女進入一個信仰的群體.
堅持.
如果我們現在形容.
就是堅持要帶我們子女回教會.
他多不想都拖他回來.
另外亦都幫助我們的子女.
能夠經濟獨立.
這是要教導我們的子女.
能夠對自己的生活負責.
幫助了子女.

$^{881}$他們能夠生活之後.
亦都要幫助他們生存的技巧.
不單止讓他們能夠生活.
他們更加懂得怎樣去生存.
最後就要教導他們人際關係.
尊重別人.
亦都要去接納別人.
這個在猶太古典去教導做父母.
我們怎樣用一個榜樣去教導我們的兒女.
所以做父母的弟子姐妹.
我鼓勵大家.
我們不單止在賢教上面.
我們做好我們的榜樣.
我們更加在新教上面.
都能夠走出來去影響.
去祝福我們的下一代.
所以今天當我們去看這段經文.
這段經文說了很重要的訊息.
就是什麼呢.
當我們做父母的.
我們要以主的方式去養育我們的兒女.
兒女就可以學習到尊重馴服.
元女就可以去孝敬父母.
讓家庭成為見證.
然後我們的信仰可以傳承下一代.
下一代又再成為父母.
他們又再以主的方式.
能夠可以養育下一代.
以致我們的信仰就能夠可以怎樣.
可以傳承.
當我們的家庭關係健全.
福音就有土壤可以生耕.
可以茁壯去延續.
弟子姐妹.
當我們能夠在家庭裡面.
我們活出孝敬父母.
我們的信仰就能夠傳承.
可以祝福我們的下一代.
所以讓孝敬父母成為我們一個祝福的起點.
不過最後提醒我們.

$^{921}$在孝敬父母在家庭相處的關係裡面.
有一個很重要的重點就是.
好像服侍主.
不像服侍人.
這個理解很簡單.
就是當你看到你的子女.
很難受的時候.
你就跟自己說.
我在這裡服侍主.
不是服侍人.
當你看到你的父母很惡搞的時候.
你又跟自己說.
好像服侍主.
不是服侍人.
你就能夠可以愛你的父母.
孝敬你的父母.
從主裡去養育你的兒女.
所以孝敬父母.
不再只是一個家庭的責任.
而是一種崇尼的敬拜.
當我們能夠這樣活的時候.
兒女就學會尊明.
父母也都學會引導.
我們的家庭就成為見證.
信仰也都不會只會停在我們這一代.
而是能夠可以流動.
延續傳承到下一代.
最後我有少少的鼓勵.
在我們當中.
可能今天有些弟兄姊妹.
你的父母可能已經離開這個世界.
或者可能他回了天家.
可能你現在很掛念他,很捨不得他.
弟兄姊妹我會鼓勵你.
讓我們有一個值得我們懷念.
值得我們思念的父母.
我們獻上感恩.
又或者我們當中的弟兄姊妹.
可能在地上的父母.
他帶給你的可能是很多的傷害.

$^{961}$或者很多的困難.
我預祝親自去醫治你的心.
讓你能夠得到釋放.
不再被過去的經歷傷害困擾你.
而最後一個鼓勵.
我想分享一個見證.
弟兄姊妹你知不知道這個人是誰.
在文化界應該認識的.
似乎我們都可能不在文化界.
應該都不是沒有什麼文化的.
這個是林語堂.
國學大師.
和你分享他的見證.
林語堂他是生於一個三代基督徒的家庭.
他爸爸是牧師.
他原本十八歲的時候預備去讀大學.
已經預備去讀神學.
可能將來都是做一個牧師之類.
但當他進入大學讀神學的時候.
他看到很多的規條.
所以他覺得他接受不了.
他當時也被西方的哲學吸引了他.
所以他不單止離開了讀神學.
他更加離開了信仰.
他去到外國留學.
他繼續去學習.
他更加成為一代的國學大師.
很多人認識他.
我們土宅的弟兄姊妹不認識他.
因為我們不是這方面.
所以他是享負盛名的一個國學大師.
但當他在他人生最後的一個階段.
他遇到有兩個人.
一個是他太太.
他太太每個星期日都捉了他回教會去崇拜.
每次在崇拜的時候.
他都覺得牧師講了一些教條規則規矩.
他都覺得很接受不到.
但因為太太捉住他.
所以他唯有陪著太太一起回崇拜.

$^{1001}$另一個就是他的兒甥女.
兒甥女在他患病的時候.
對他有很好的照顧.
亦都不斷跟他說.
你相信耶穌.
你回去天父的身邊等等.
去陪伴著這個兒甥女.
而林語堂最終.
他就是因為他太太不停地帶他回教會.
他有一次遇到一個牧師.
這個牧師他去講道的時候.
他講了一篇生命的道.
他不是在講一些教條規則規矩的時候.
他覺得他的生命得著釋放.
所以他真的決定重新回到信仰裡面.
林語堂他講了一句說話.
他說當他在晚年的時候.
他人生最大的得著的安息.
是因為他回到信仰裡面.
所以他在晚年的時候.
唯一一本信仰的書就是.
你們的左上角的信仰之旅.
分享他的信仰和見證.
弟兄姊妹為什麼我臨尾要去.
當我說要做最後的鼓勵的時候.
會去和你分享林語堂的見證呢.
弟兄姊妹我很想和大家分享的是.
弟兄姊妹可能在你的家庭裡面.
可能是你父母也好.
可能是你的兄弟姐妹也好.
可能是你的兒女也好.
可能他從前有回教會.
可能他有認識主.
但可能因為很多不同的原因.
他今天已經不在信仰裡面.
他沒有回教會.
更加沒有和神去結合關係.
我鼓勵你不要放棄.
仍然的去捉住他們.
哪怕只是一句簡單的慰問.

$^{1041}$一個簡單的問候.
哪怕只是為他去祈禱.
哪怕只是威逼利誘.
逼他回來教會也好.
讓他們仍然有機會.
可以回轉回到神的身邊.
以至我們的家庭.
今天當我們說孝敬父母.
我們都能夠經歷到.
這條大應許的誡命的祝福.
我們的信仰能夠一代傳承給下一代.
讓我們的下一代.
甚至我們的後裔.
都能夠可以一直在主裡面.
蒙神的祝福.
弟兄姊妹,你願不願意呢?.
\newpage



\section{歌羅西書 3:5-11}
\label{sec:032115}
\textbf{我們怎樣彰顯神的榮耀?}~~~~黃紹權牧師
\newline
\newline
連結: \href{http://www.church.com.hk/acms/content.asp?site=cdc\&op=show\&type=product\&code=032115\&layout=sermon}{\texttt{032115}} ~~~~ 語音日期: 2026-05-10
\newline
\newline
\hyperref[sec:032114]{\small{< < < PREV SERMON < < <}}
~
\hyperref[sec:index_chronic]{\small{[返順時目]}}
~
\hyperref[sec:032116]{\small{> > > NEXT SERMON > > >}}
\newline
\newline
歌羅西書 3:5-11
\newline
\begin{longtable}{cl}
\hline
\hline
章節 & 經文 (和合本修訂版)\\
\hline
3:5 & \begin{tabularx}{0.7\textwidth}{X} 所以，要治死你們在地上的肢體；就如淫亂、污穢、邪情、惡慾和貪婪—貪婪就是拜偶像。 \end{tabularx} \\ \\ \relax
3:6 & \begin{tabularx}{0.7\textwidth}{X} 因這些事，神的憤怒必臨到那些悖逆的人。 \end{tabularx} \\ \\ \relax
3:7 & \begin{tabularx}{0.7\textwidth}{X} 當你們在這些事中活著的時候，你們的行為也曾是這樣的。 \end{tabularx} \\ \\ \relax
3:8 & \begin{tabularx}{0.7\textwidth}{X} 但現在你們要棄絕這一切的事，就是惱恨、憤怒、惡毒、毀謗和口中污穢的言語。 \end{tabularx} \\ \\ \relax
3:9 & \begin{tabularx}{0.7\textwidth}{X} 不要彼此說謊，因為你們已經脫去舊人和舊人的行為， \end{tabularx} \\ \\ \relax
3:10 & \begin{tabularx}{0.7\textwidth}{X} 穿上了新人，這新人照著造他的主的形像在知識上不斷地更新。 \end{tabularx} \\ \\ \relax
3:11 & \begin{tabularx}{0.7\textwidth}{X} 在這事上並不分希臘人和猶太人，受割禮的和未受割禮的，未開化的人、西古提人、為奴的、自主的；惟獨基督是一切，又在一切之內。 \end{tabularx} \\ \\
[1ex]
\hline
\hline
\end{longtable}
$^{1}$今天我們會分兩個部分去看這段聖經.
去嘗試去理解榮耀神這件事.
和我們作為神的子民.
甚至作為一個人的關係是怎樣.
先看看提及前文講些甚麼.
在前文我們已經講過.
作為一間不似其他眾多的教會.
只是依賴著去獲得這個世界的認同.
作為一種好像傳遞的藉口.
或者踏腳石的哥羅西教會.
它不是尋求這些東西的.
雖然在一個大城市裡面.
它不是去尋求人的肯定.
或者社會的肯定.
而去嘗試藉此成為教會.
影響這個世界的原因.
它不會妥協.
它不會去因為這個世界潮流流行甚麼.
而跟著去做的時候.
我相信這個更加合乎上帝對於教會的要求.
因此今天即或我們處在這個.
所謂二十一世紀的年代.
我們作為教會也都應該學會.
不要以一種只是去巴結這個世界.
意圖使這個世界人喜歡教會而吸引人教會這種做法.
因為這種做法很容易將教會拖入一個現象.
就是教會會變成向上帝的說話.
作出各方面的.
我們可以這樣說.
去刪除.
很多東西不提的.
所以在教會裡面.
今天你會留意到.
你上網看一下.
很多港獨.
牧者很多都不講的.
很多書卷從來不會有人碰的.
為甚麼?.
因為那些講了出來會得罪人.
講完歌羅西書之後.

$^{41}$我會回去講舊約.
我會去講俄巴底亞書.
一卷很少人會去講的經文.
為甚麼?.
因為這些書很多信徒覺得.
或者教會覺得只是給神學院來研究.
不是用來給信徒讀.
甚至給信徒看.
因為那些東西講了出來之後會得罪很多的人.
教會應有的存在的基礎.
是不應該為了擔心因為傳主的道而得罪人.
反而我們要明白到.
我們的角色就在這個世界裡面.
當傳講主的道的時候.
會興起了一場屬靈的爭戰.
而這場的爭戰會非常之激烈.
我們會被這個世界的人.
甚至去討厭或者憎惡.
但是教會仍然把守著這個對真理的堅持.
你說牧師這樣會受傷的.
是呀教會會受傷的.
打仗怎麼會不受傷呢?.
但是在上文我們看到的就是.
原來我們的上帝會把我們放在一個隱蔽之下.
我們已經講過那個隱蔽的字眼.
可以把它翻譯為掩護.
好像打一場戰爭一樣.
你的同伴為了你去履行.
可以進攻的目的或者拯救的目的.
於是乎他會為了因此而作掩護.
保守著我們去完成這個工作.
這些就真正是基督徒的事奉.
也是我們在聖經裡面看到的.
上帝給我們作為信徒.
當我們在這個世界裡面.
願意去按真理而活的時候.
基督耶穌成為我們的掩護.
使到在這個世界上面.
縱使有多少惡勢力.
怎麼討厭我們都好.

$^{81}$我們卻在這個過程當中得蒙保守.
但更重要的是.
可以因此而吸引或者拯救一些.
真正上帝要拯救的人.
而我們今天作為信徒的.
宣講神的說話其實不是我們去拯救那個人.
而是我們將上帝的說話.
讓那些要上帝被拯救的人.
他被吸引回到上帝的面前.
有回一份同樣對上帝堅持不會妥協的生命.
這個是上文我們所看的.
來到接著這段的聖經.
進入了一個題目是更加令人費解.
因為他談論到有關於神的榮耀的問題.
而這個榮耀的問題是包括了.
很多有關於創世的種種的關係.
我們今天會從這個角度.
重新揣摩歌羅西書想講的主題是什麼.
讓我們逐步去嘗試.
看看怎樣好像在聖經裡面.
我們可以去介紹去參與在.
這份彰顯神榮耀的工作當中.
什麼叫做榮耀神?.
今天很多時候我們回教會很多牧者可能就說給你聽.
事奉主.
什麼叫事奉主?.
事奉主就是在教會擔當一官半職.
不要只是做平信徒.
不要只是坐在教會.
只是坐在那裡去崇拜齋聽.
要做一些事情的.
所以有些信徒就認為.
那個榮耀神的情況就是先由平信徒發展.
接著就去做一下教會的一些事奉參與.
做一下做義工.
我強調教會的工作.
不應該用義工這個觀念去看的.
因為義工是只不過好像一份興趣.
你推動你去做的.
教會的事奉是一種上帝的委託.

$^{121}$你是有職責去委任.
並且去到履行的.
不是義工.
做又可以不做又可以.
有空就做一下.
沒空就不做.
不是.
你沒空你都需要做.
這個就是事奉和做義工的不同.
但是很多教會將教會的事奉變稱義工.
然後怎樣?.
你再熟悉一點的做教會的一些團契的職員.
再熟悉一點的你就做教會的掌執.
再熟悉一點的讀神學院.
做牧師.
很多信徒是這樣看的.
聖經是不是這樣看所謂榮耀神.
怎樣彰顯到神呢?.
很有趣.
我們一起去看一下.
第一部分我們要看的.
是從第五節到第八節.
講及到榮耀神原來是參與一份與罪爭戰的.
你死我活的爭鬥.
榮耀神不是坐在那裡唱「嘻嘻哈哈」.
或者「你好我好大家好」.
或者教會出去有空有什麼活動.
大家可以覺得很光榮.
不是.
榮耀神是參與著一場.
你死我活的爭戰裡面.
而這個爭戰是針對著罪來去參與的.
所以教會無論作任何的事.
任何的事.
無論你是教會唱詩也好.
你是崇拜也好.
你是參加授浸也好.
你參加團契也好.
其實你都是在參與著這份與罪有一個.
直接的衝擊和爭戰的活動.

$^{161}$榮耀神不是一個口號.
不是一個嘴巴在那裡說.
更加不是只是純粹做一下教會時工.
榮耀神是整個人由心而發.
在整個的心裡面.
承諾毀身在這份全面與罪.
進行衝突的爭戰的當中.
所以當我們去看回第三章的第五節裡面.
它就用了一連串.
其實不只是第三章的五節.
一直下去有一連串的命令語.
在希臘文裡面用了一個叫做.
imperative的語氣.
imperative可以只是一個指示.
一個命令.
但是也都可以說是一個.
充滿導向性的一個勸誡.
向哥羅西教會指出.
我們必須要自死在地上的肢體.
到這裡來算.
很多信徒就說.
哦.
原來要自死地上的肢體.
指的是什麼呢?.
指的是不是好像天主教那樣.
當有什麼信徒.
他與他們的教義不同的時候.
就將他驅逐出教會這樣的行為.
英文叫做excommunications.
知不知道什麼叫excommunications.
communications你懂的.
交通嘛.
excommunications就.
驅除他不讓他與你交通.
昔日天主教會就是這樣運作.
今天很多信徒都以為是這樣的.
最近我們福音浸聯會曾經.
發出了一個這樣的宣佈.
其實剛剛在開年會.
這個年會裡面就講及到.

$^{201}$在草原省份的我們福音浸聯會的分會.
就將當中有五間的教會.
將他除了他們的教會聯會的會籍.
有些人就說.
哇.
是不是又是走天主教那一套.
將他驅逐出去.
請大家留意.
如果我們在真理的上面.
是作了多次的規勸的時候.
如果我們再將他視為.
好像和我們一模一樣的信徒.
但是他是仍然沒有改變的.
沒有改過的.
在這方面.
就不是和你同心.
我們教會.
我們聯會叫做fellowship baptist.
你知不知道為什麼叫fellowship baptist.
福音浸聯會.
為什麼叫做福音浸聯會.
因為我們是有共信才叫fellowship.
如果你和我的信不同的時候.
我和你沒有聯絡的.
那浸聯會是不是現在在做的就是.
驅除他出來.
不是的.
我們聯會做了一樣東西.
雖然他的會籍不再是福音浸聯會.
但是仍然是在報告裡面說.
我們希望他們的教會.
是可以改變他們的信仰.
回到聖經裡面的標準.
但是這個會籍的問題.
其實只不過是顯示出.
他和我們的信有分別.
但是不代表我們驅除他.
不和他溝通.
天主教就不同了.
驅除他出去.

$^{241}$所以為什麼會有馬丁路德改革.
這個宗教改革的事情.
驅除他出去.
大家要記住.
馬丁路德是一個天主教的教士.
他不是一個我們以為基督徒.
就是基督教會裡面的信徒.
不是的.
他其實是想天主教會歸回真理裡面.
但是天主教會拒絕了.
不再和他溝通.
於是乎他才搞到要釘出.
那幾十條的教條出來.
去展示出正確的信仰是什麼.
和我們浸聯會做的不一樣.
但是在這裡這段聖經所說.
要治死你們地上的肢體.
是不是就指.
和這些人斷絕來往呢.
甚至有些信徒解經的上面就說.
哎是啊牧師我記得你先前曾經說過.
上文談論到原來有一個基督的身體.
這個基督的身體在這個時候這樣的情況.
要將他治死的肢體.
是不是就指到那裡.
將那個肢體上面的一些軀幹.
將它斬除它呢.
是不是這個意思呢.
我們看看.
這段聖經說到.
治死地上的肢體有清楚的界定.
英文叫quantifier.
或者qualifier.
應該講得準確.
qualifier.
有界定的.
界定什麼呢.
我們讀一下這裡幾樣東西是什麼.
第五節.
咒語什麼.

$^{281}$淫亂.
污衛.
邪情.
惡辱.
和貪婪.
加多句括號.
貪婪咒語拜偶像一樣.
這裡有沒有指任何一個人.
在教會裡如果做得不對.
信仰不同.
所以我們將它斬了斬斷它.
excommunicate它呢.
沒有.
相反在這裡說到的是.
一連串由人內心裡面.
所誘發的一連串的邪惡.
所以保羅在這裡向哥羅西教會所展示出的.
其實不是好像天主教那一套.
將他們的人或者怎樣.
斬除跟他再沒有關係.
相反.
保羅在這裡.
作出一個這麼不客氣的說話.
是要讓教會針對人裡面的問題.
這些什麼問題.
歸咎於同一個字.
罪.
保羅在這裡向哥羅西教會告訴他們知道.
我們要一同參與這場屬靈的爭戰.
不是去針對哪個人.
不是去針對哪個群體.
或者不是去針對哪個教會.
而是去針對人裡面的罪.
和它相關所誘發種種的後遺症.
這個才是核心所在.
去到第六節.
因這些事.
上帝的憤怒必臨到那撥逆之子.
在這裡我要先做個註腳.
在原文聖經的抄本裡面.

$^{321}$有一些很重要的抄本.
可靠的抄本.
是沒有這個臨到那撥逆之子這句話的.
沒有撥逆之子這句話的.
這句話在和合本有的.
我們看看有沒有大分別呢.
我們看看.
為什麼要致死這一連串在人裡面的罪呢.
因為這些事是引發神的憤怒.
這個才是主題.
臨到撥逆之子.
這個其實可有可無的記載.
因為世人都犯了罪.
這個說話是沒有錯的.
所以全部人都是撥逆之子.
包括在哥羅西教會後面會說的.
你們都是撥逆之子的一部分來的.
但是只要你願意參與.
在致死這一連串的罪惡.
我們就可以解除上帝的憤怒.
上文說到要去榮耀神.
要怎樣榮耀神.
所以第五節說.
所以你要怎樣做.
要致死這一連串的罪.
為什麼要致死這一連串的罪呢.
因為最重要的目的.
是免了上帝的憤怒.
要去擴展在他創造的人當中.
人是神創造的時候的巔峰.
人是上帝創造裡面.
那個屬於反映上帝的重要的角色.
如果人能夠歸回神的樣式的話.
換句話來說.
神就將他創造的榮耀.
能夠繼續存留在他的創造當中.
因此.
如果我們能夠參與在致死人裡面的罪.
換言之.
我們是將人帶回原本上帝神造人的樣式.

$^{361}$剛才有沒有聽清楚.
人是神創造裡面.
反映神的角色.
如果有真正的人的質素的時候.
沒有罪在他裡面的時候.
神的榮耀就彰顯.
顯露在他的創造當中.
所以保羅在這裡所談論到的.
這個所謂榮耀神.
是一件很嚴肅的事情.
不是在教會搞一下聖功.
搞一下什麼活動.
從而使大家覺得.
自滿地覺得好像神聖了.
我們在教會裡面.
一同去學會去榮耀神.
是我們一起去致死我們的罪.
所以教會的群體.
我們做每一件事.
其實是幫助弟兄姊妹.
去面對著我們自己的罪.
所以保羅很精準地去把握到.
你必須要處理好.
這個行為或者這個問題.
不過我們的現況是怎麼樣.
第七節說.
當你們在這些事中活著的時候.
也曾這樣行過.
但現在你們要欺絕這一切的事.
以及惱恨,憤怒,惡毒,毀謗.
並口中污穢的言語.
保羅在這裡.
說了給哥羅西教會知道.
應該怎麼樣可以去榮耀神.
榮耀神的目的是怎麼樣.
原來要收復人的形象.
是上帝做人的原本的形象.
就可以將人放回一個上帝的創造裡面.
去顯出上帝的榮美.
不過你們是怎麼樣呢.

$^{401}$保羅好像在數他們以前的不是.
當你們在這些事中活著的時候.
這句話中文讀出來.
可能你未必太清楚時態是什麼.
英文聖經New English Translation.
他將原文翻譯得很準確.
他這樣說.
如果要準確些翻譯成中文.
就是說其實你們以前.
曾幾何時都是活在.
同樣這種充滿淫亂,污穢,邪情,惡辱和貪婪的裡面.
哇!保羅你寫一封信給一間教會.
其實你跟他們沒有直接掃尾謀妙.
沒有傳道去過那間教會.
雖然你認識他們裡面一些人.
你寫些這樣的東西.
數他們的不善.
你會不會聽啊.
換句話來說.
保羅寫封信來.
以馬來尼亞人浸信會.
然後找一個人讀.
我們相信在第一章裡面.
講了那位人士.
在第七節裡面.
以巴弗.
將他讀給以馬來尼亞浸信會的人聽.
然後指著和睦思.
你知不知道你以前.
其實你曾幾何時.
是曾經活在淫亂,污穢,邪情,惡辱,貪婪.
哪有人聽啊.
你聽不聽啊.
什麼人不聽呢.
尤其是做過這些事的人.
你講我啊.
所以有些弟兄姊妹.
以前我講道.
他經常說.
牧師你在台上講我.

$^{441}$我怎知道你發生過什麼事.
我怎知道你淫邪過.
污穢過.
我怎知道.
我怎知道你臭爛口過.
我不知道啊.
我以前都講粗口的.
是不是啊.
我沒理由寫一堂道.
特意針對你的.
是不是.
本來在這裡.
講什麼啊.
他不單止講我.
你曾幾何時.
活在這樣的情況.
也曾這樣走過.
好像你知道我們.
做過什麼事.
哪有人聽啊.
你聽不聽啊.
如果你讀經讀到.
遍體鱗傷的讀法.
那當然是.
當然不明白.
讀完第八節.
但現在.
你們要棄絕這一切的事.
英文.
New English Translation.
是這樣譯的.
All such things.
接著那堆東西.
本來在這裡.
給第二個的命令.
或者勸喻.
但是現在是怎樣.
不同了.
你們肯接受了.
我現在給你的命令.

$^{481}$是怎樣.
將一連串.
這些不好的事情.
開始將它扭轉過來.
不再做.
是哪一連串的事啊.
惱恨憤怒.
惡毒誹謗.
並口中的污穢的言語.
寫文呢.
精彩.
有深度的人就這樣寫東西.
怎樣深度啊.
你看不看到.
頭第五節.
講了五樣他的罪.
第八節又講了五樣的罪.
頭第五節就講.
你以前做過這堆事.
第八節講.
你現在就決意用這一連串.
不再做這堆事.
來改善.
這種在寫法裡面.
就叫做inclusio.
前後對應回應的說法.
的陳述.
有什麼作用呢.
是要讓收信的人.
看見自己前後的改變.
因此如果我們在這個基礎裡面.
去看的時候.
第七節保羅講.
哥羅西教會曾經做這些事.
不單止活在其中.
並且做這些事的時候.
就不是去挑釁他們的.
壞處來講.
而是欣賞他們願意.
肯勇敢地去面對.

$^{521}$昔日的自我.
自我是怎樣.
那個舊我是充滿罪惡的人.
真正信耶穌的人.
被人點出他的罪.
他的回應不是撐.
他的回應是承認.
所以怎樣知道一個基督徒.
是真基督徒和假基督徒.
是非常之容易的.
就是當你.
一個人一個基督徒.
被人看到他錯誤的地方的時候.
他第一件事不是.
不是呀.
你誤會我了.
不是這樣講的.
他說對呀.
我真的做了這些事.
這個其實是很簡單的神學來的.
叫做認罪.
基督第一件事就是認罪.
保羅在這裡講的是.
你如果要參與.
在榮耀神的事情上的時候.
不是走去.
在教會做些什麼.
驅除哪些人出會呀.
誰讓你出來公審呀.
No 全部不是.
而是一起去面對自己的罪.
去承認自己的罪.
而承認過後.
是有行動付諸去改變的.
所以.
網上有個弟兄.
應該是弟兄.
我猜吧.
他寫了好幾次電郵來給我.
隨便寫.

$^{561}$因為寫電郵詳細些.
寫短信又.
寫短信通常寫得很短.
還有會錯的.
他問我多次同一條問題.
他說牧師呀.
我現在在讀聖經.
什麼叫做.
這個所謂Inclusive.
首尾呼應的寫作方法.
我順便回答了.
現在聖經任你問.
回答不了那麼多.
暫時回答兩句吧.
首尾呼應其實這種寫法.
一方面是使得讀者更深刻.
知道的內容.
有一個回應.
Echoing.
前句和後句.
會互相呼應的說法.
使得你更深刻.
最重要的是Inclusive裡面的核心.
想說些什麼.
核心寫在中間.
治罪.
是基督徒第一樣東西.
這個網上的弟兄.
我當弟兄.
他問了我兩次.
我後來在中間問他.
你是不是在讀神學.
因為很少弟兄姊妹會走來問.
去哪裡聖經找這些Inclusive去研究.
我說你是不是在讀神學院.
有些什麼paper.
有一段經文給我.
叫我分析給他.
然後寫完給他.
我很仔細寫給他.

$^{601}$寫完之後我給他.
他回覆我.
他說多謝你的回覆.
他說牧師.
不過我不是讀神學的.
我只是在讀一個讀經的查經班.
查經班.
弟兄姊妹.
有一連串的問題.
有發覺很想知真理.
我們以馬利雲浸信會的.
網上的弟兄姊妹.
我們網上會有.
是這樣的質素.
不是問垃圾問題.
是問一些很有質素的問題.
是希望自己在靈性上成長.
保羅在這裡是不是想.
哥羅西教會成長.
他是不是不好.
不是 保羅知道.
所以保羅說.
你們現在就會棄絕.
去行以下一連串的惡事.
保羅這個陳述.
是讓哥羅西教會的信徒知道.
如果你要榮耀神.
你必須要參與這份.
承認自己曾經在罪裡面的參與.
但是也願意參與在.
撕掉這個罪的增薦裡面.
原來真正榮耀神.
不是在教會搞事工.
而是在怎樣幫助信徒.
處理面對人的罪.
甚至我們自己參與.
在處理我們的罪裡面.
原來處理罪.
這個就是榮耀神的平台.
The platform to exemplify the lost glory.

$^{641}$弟兄姊妹.
我們有沒有.
參與在這一場.
你死我活的爭戰裡面.
不是我說的 聖經說的.
致死.
罪你要死 我才可以活.
當我們能夠致死罪.
我們活的時候.
活出的就是上帝的生命.
活出的就是上帝原本在創造的時候.
人的屬性是怎樣.
這樣就能夠彰顯榮耀上帝.
因此從這一點.
帶到我們進入了.
一個更加深層次的反思.
伊甸園裡面那條蛇.
很多信徒說牧師.
為什麼上帝要放這條蛇在裡面.
我現在首先說一個很重要的信息.
上帝放這條蛇其實是.
成就整個創造.
完美的創造的一個很重要的元素.
因為上帝容許在伊甸園當中蛇的出現.
陰柔他進行引誘.
甚至好像讓他摧毀了神的創造.
包括神的最高的設計.
人.
但是神就是藉著蛇所做到的工作.
將繼就繼.
將蛇所種下來死的這個種子.
將它轉化成為.
將創造變得完美的一個重要環節.
怎樣做?.
你和我的參與.
當我們願意去致死罪的時候.
這個創造就變得完全.
換句話來說.
簡單點說.
人能夠認罪悔改歸回神.

$^{681}$走回從以前的昔日的罪的生活.
走出來走回上帝的道路.
這個就是使上帝榮耀彰顯的一刻.
怎樣將它進一步全面發揮出來呢?.
第九至十一節就告訴我們.
我們透過穿上一個有基督樣式的新人.
來將它全面顯露出來.
我們一起讀讀這段聖經好不好?.
第九至十一節.
不要彼此說謊.
因你們已經脫去舊人和舊人的行為.
穿上了新人.
這新人在知識上漸漸更新.
正如做他主的形象.
在此並不分希臘人,猶太人.
受過禮的,未受過禮的.
花外人,西古提人,遺奴的,自主的.
唯有基督是包括一切.
又住在各人之內.
在這個部分一連串地談論到一堆.
讓我們有一個截然不同的生命的改變.
承接上文第三章第八節.
說了一連串在口舌上的問題.
得到一一處理之後.
保羅就向一個屬靈的群體.
即是哥羅西教會.
其實不單是哥羅西教會.
去到第四章我們知道.
是向整個亞洲教會宣佈.
建立成為一個屬靈的群體.
我們遲些去看.
當他去說到這個屬靈的群體得勝的時候.
就有一個截然不同的發展出現了.
不過在這個時候.
我們首先就要處理第九節裡面.
有一句很難明白的說話.
不知道為什麼會出現了出來.
第九節這樣說.
不要彼此說謊.
失驚無神說什麼.

$^{721}$不要彼此說謊.
弟兄姊妹.
如果你讀經.
我敢肯定百分之九十九的信徒.
一定會把這句說話當作沒有讀過一樣.
因為和上文下理好像不是很接的地方.
你知不知道讀聖經.
你的讀經能力.
可以提升它的進步空間.
或者我這樣說.
你想了解神說話上的功力能夠盡心.
這些經文你就一定要處理.
因為當你處理到.
你就驚啊.
無端端說什麼.
不要彼此說謊和上文下理有什麼關係呢.
剛才我已經引導了一點點大家去看.
在第八節裡面說了一連串.
在口舌上面的問題.
惱恨.
憤怒.
惡毒.
毀謗.
並口中污穢的言語.
逐步去引伸到言語這個範疇的上面.
保羅是要告訴.
哥羅西的教會的信徒知道.
你們要對自己誠實.
什麼誠實啊.
你根本上在這些一連串的事情上面.
一直是有犯的.
不過你肯去願意撕掉它.
你就會變成一個完全脫胎換骨的人.
你知不知道為什麼很多人經常說.
牧師我知不知道我是信主啊.
我什麼時候是信主.
我怎樣去印證我自己是信主.
其實很簡單的.
當你肯承認你自己的罪的時候.
誠實面對你自己.

$^{761}$我是一個罪人.
當然這句說話不是說出來.
那怎樣啊.
不是啊.
你知道很多人很爛的.
你跟他傳福音.
譬如很簡單.
我用一個簡單的例子.
不是假的.
如有類同實屬巧合.
有一天我在Queen街走著走著.
或者香港的砵蘭街走著走著.
經過的時候.
看到那個弟兄.
不要說Alan了.
我怕誤會.
原來是Peter.
我們這裡沒有Peter.
網上有.
對不起.
純粹虛構.
叫Peter在那裡跟一個妓女在那裡交談.
然後就上了樓上.
我後來回教會見到他.
我就跟他說.
弟兄Peter.
你是不是犯奸淫啊.
如果你是一個.
不說謊的人.
請留意.
第九次要說.
不要.
有兩個字.
彼此.
不是說你自己.
我自己可以了.
我自己承認自己的罪.
彼此是一個公開的關係.
在這樣的情況之下.
你肯承認你自己的罪.

$^{801}$誠實面對.
而且不是.
當然我沒有叫你公開來說我去召妓.
或者我犯了什麼罪.
沒有我沒有叫你這樣做.
不過你肯承認自己是一個真正的罪人.
這個就是能夠讓你脫胎換骨的時間.
你想知道你是不是真的信耶穌.
你有沒有真正承認你自己.
真相是怎樣墮落呢.
曾經有一個弟兄.
有一次在主日崇拜的時候.
聖餐主日.
在聖餐的時候.
往往很多時候我會這樣說.
請大家檢省大家的罪.
不過那一次.
不知道為什麼.
突然間聖靈或者怎麼樣.
催逼我就多說了一句話.
我說如果你是犯罪.
犯罪.
就請你不要領聖餐.
我說了一句這樣的說話.
其實.
牧師為什麼你會說.
我想都有一些原因的.
因為我收到線報.
你說牧師你很多線人.
不是很多線人.
我收到人是見到有人去犯奸淫.
不過我又不敢相信是不是.
我只能夠.
我不是警察.
我不是靈性警察.
我來檢控你的.
我沒有來進行調查的.
是不是.
不過希望你自己審查.
我那個時候沒有留意.

$^{841}$其實你們領聖餐.
有些我們就.
通常是我派的.
有些教會是懶鬼.
自己每人.
入場的時候.
每人派一個餅杯一套.
到時候撕開一層一層那些.
不要說他懶鬼了.
算了.
收回我這一句.
過了一天.
我收到一個電話.
有個弟兄打電話來跟我說.
他說牧師我想跟你聊聊天.
好 聊天沒所謂.
首先.
你找我聊天.
不一定是你犯罪的.
請大家留意一下.
不是說找親牧師.
當然是我犯罪才找他.
我說好啊.
出來吃東西.
吃個午餐.
吃個麵包.
一人一個sandwich.
他說他請我.
我說不要不要.
他說要的要的.
我覺得為什麼.
算了.
下次大家有來有往.
下次我請.
接著他跟我說.
他說牧師我昨天沒領聖餐.
我說為什麼.
為什麼你沒領.
他說我就是聽到你說了一句.
其實你不說那句.

$^{881}$說如果有犯罪的.
就不要領聖餐.
我都照領的.
不過你說了這句話的時候.
再回答了最後.
就說不要乾飯煮的身體煮的血.
我就很宅心.
我跟這位弟兄.
你知不知道怎樣回答他.
我說你肯承認自己的罪.
就已經證明你是一個真的基督徒.
你肯面對已經是一個.
基督徒一樣犯罪.
我首先重複再說.
基督徒一樣犯罪.
但是犯了罪會懊悔的.
會竭力求上帝.
用各用方式來幫助你.
撕掉這些罪.
哪裡說.
就是第三節說.
致死你的罪.
你肯這樣做的.
為什麼你肯做.
因為你想棄絕這一切.
這樣惡毒的言行.
而最重要的言行.
最大的言行.
就是自己欺騙自己.
你自己錯你自己不認.
弟兄姊妹.
上帝差一個人來你身邊.
指出你的罪.
而不是.
死撐.
死撐的我知道.
我永遠不會和你撐的.
你不用擔心.
不會和你撐.
因為我只知道你不信主.

$^{921}$當然我有別的方法.
希望幫你去認識.
你自己是怎樣犯罪.
希望你認罪悔改.
但如果你仍然是不肯的.
你是自己在欺騙自己.
所以保羅在這裡很有智慧.
不要彼此說謊.
突然間爆了一句這樣的說話出來.
是要告訴你和我知道.
大家不要去以.
自己欺騙自己.
作為一種.
活在教會裡面的生活.
你知不知道教會有很多這樣的人.
不肯面對自己的罪.
甚至被人指出他有罪的時候.
他都不肯認的.
保羅告訴哥羅西教會.
沒錯你們的靈性的確很好.
你充滿信心.
但是在這個過程當中.
其實你們仍然需要有一個更新的過程.
你肯承認.
肯看到你自己的問題.
願意去致死他.
你就能夠在這個榮耀的.
上帝的事情上有參與.
當你肯正確面對.
不會再自己欺騙自己的時候.
你的生命就開始進入一個.
截然不同的改變.
保羅就用了好像穿衣服的比喻.
這裡說因你們已經脫去救人.
和救人的行為.
然後穿上新人.
而且這個新人.
不是一下子就全部變成新人.
是怎樣呢?.
這新人在知識上漸漸更新.

$^{961}$正如做他主的形象.
所以基督徒榮耀神的過程.
是在穿上這個新人的時候.
不斷在改變.
改變改變改變.
不是走向墮落.
不是走向放棄.
不是走向墮落.
不是走向犯罪.
淫亂污衛.
邪情惡毒貪婪.
不是向老恨憤怒.
惡毒誹謗口中污衛語言.
不是.
他是漸漸將這些一路一路改變.
越來越像上帝原本做他的形象.
創世記的題目.
上帝做人.
是要讓人在上帝的創造裡面.
顯明上帝.
這些就叫榮耀神.
今天作為基督徒.
要上榮耀神.
就請參與在這個面對罪的時候.
逐步在這個新人身上.
不斷成長起來.
越來越似越來越似越來越似.
越來越似上帝原本做人的形象.
如果你的形象.
越來越不像上帝.
越來越像這個世界.
越來越和這個世界妥協.
你就做不到榮耀神.
所以你明白為什麼上文.
保羅說得這麼清楚.
哥羅西教會你不要在這個世界裡面.
去拜那些小學.
說一下那些什麼DNA.
說一下那些傳統.
說一下那些什麼中國人.

$^{1001}$什麼華人.
西人.
那些什麼的習俗.
然後怎樣將基督教吸下去.
對不起.
吸下去是不是英文.
還是中文我不知道.
吐下去.
我在神學院讀過一科.
叫做如何面對.
教會與這個世界的關係.
在這個課程裡面.
其中一個題目.
就是說將教會本色化.
本色化英文.
你如果想查叫contextualization.
不是將聖經在context上解釋的意思.
是要將聖經.
去fit進去哪一部分.
適合這個世界.
而去使教會就用那一套.
好像一個過濾器.
這個世界成為過濾的那塊網.
聖經穿得過這塊世界過濾網.
就用那些.
穿不過那些就不要用了.
在神學院裡面.
在我那班裡面.
充滿著名牧者.
著名啊.
我告訴你.
有神學院院長.
有號稱在多倫多建立.
八間浸信會的主任牧師.
他們在課程裡面.
很大事在課程裡面張揚.
要將教會怎樣本色化.
套進去華人教會裡面.
按照什麼方法.
華人文化.

$^{1041}$因為我們是華人教會.
所以我們要用華人的文化.
去過濾聖經裡面.
哪些適合華人教會用的.
我們就用.
這樣就使人容易進入教會.
鄧師傅.
如果是這樣的話.
華人很喜歡娶妻妾.
那可不可以華人教會都容許娶妻妾呢.
你知不知道他怎樣說.
亞伯拉罕有幾個老婆.
阿當有幾個老婆.
聖經裡面那些.
我們叫做.
Patria.
即是那些.
在創世記裡面著名人物.
雅各.
那些有幾個老婆.
他說那些可以用.
在座那些母親.
你想不想.
以馬來華陣信會實踐這一套.
想不想.
牧師多娶兩個老婆.
師母你肯不肯.
Daniel你多娶兩個老婆.
問問你老婆肯不肯.
但很多人會喜歡的.
突然之間.
聖經保羅寫得很清楚.
脫虛舊人和舊人的行為.
很多時候.
教會為了妥協.
就把很多舊人的東西.
套入教會.
用聖經洗白了它.
然後以為可以用.
其實是一種很危險的現象.

$^{1081}$要做回一個原本有神形象的人.
在這個世界而活.
而這個世界也是充滿罪惡.
一定有一場爭戰在裡面.
你會使到你的存在.
如果你按照神的話而活.
你會使到這個世界.
面對你的存在.
會使到你渾身不自在.
但這個正正反映出.
教會是否真的在上帝裡面.
能夠反映出神榮耀的樣式.
我用了一個說法.
這是一個怎樣成為榮耀神的.
benchmark.
就是你和我在真理上面.
那個接受實踐的程度.
直接成為有一個正比的關係.
你越肯面對罪.
你越肯對罪有個譴責.
你越肯對罪的去執行.
改變過來.
不再犯罪.
你就越有上帝的形象.
當你有上帝的形象.
你就會顯出上帝的榮美.
原來保羅在這裡用這麼多.
在創世記裡面的神學.
來介紹給.
哥羅西教會知道.
當你肯穿上.
漸漸更新的時候.
上帝就會將原本的.
因為罪所污染了的創造.
或者破壞了的創造.
直接將你和我扭轉過來.
堵塞了罪的繼續的發展.
罪沒有在我們生命上的任何影響.
不但如此.
並且上帝將破損了的創造.

$^{1121}$就是你和我.
換言一新,所以這個是新人.
換言一新.
重新包裝過.
底底裡裡都是一個新人的時候.
我們就能夠變成一個完美的創造.
大家知不知道.
在約翰福音第十九章.
第二十八至三十節裡面.
講了一句話.
有些人不明白.
不如這樣吧.
我們揭開來讀讀.
約翰福音第十九章.
二十八至三十節.
好,找到之後.
我們一起讀,好嗎?.
好,我們一起讀.
「這事已後,耶穌知道各樣的事已經成了.
圍繞史經上的話應驗.
就說『我學了』.
有一個器皿盛滿了醋.
放在那裡.
他們就拿凱翁也滿了醋.
綁在牛膝草上.
送到他口.
耶穌上了那醋.
就說『成了』.
便低下頭.
將靈魂交付上了」.
這裡連續講了兩次.
同一個字眼.
成了.
第一次講.
耶穌知道各樣的事已經成了.
什麼叫做各樣的事成了呢?.
第二次講是三十節.
耶穌上了那醋.
就說『成了』.
成了.

$^{1161}$成了什麼呢?.
不懂解聖經的人很簡單.
成了就是成了救恩.
死在十字架上成了救恩.
不過又未復活.
成什麼呢?.
所以大家聽解聖經.
別人講解聖經要很小心.
如果他很輕易講成了.
講救恩已經成了.
耶穌死了是不是救恩成了?.
耶穌不復活就不成.
耶穌復活了就成了.
對不對?.
那現在耶穌復活了沒有?.
剛剛才.
蓋國旗.
剛剛才死了.
成了些什麼呢?.
成了些什麼呀?.
這樣的事就成了.
指的是什麼事就是這樣各樣的事.
讓我告訴你.
上帝整個創造完美地完結.
這個完成了.
什麼?.
創造完.
這一刻才完成?.
牧師你說什麼?.
大家有沒有留意.
創世記的創造.
你看看.
直接做那些.
光呀.
雨水呀.
草草木木呀.
動物呀.
你留意一下那裡有些什麼.
在創造.
是父神和聖靈直接的工作.

$^{1201}$你有沒有留意.
創世記第一章第一節說什麼?.
起初.
神創造天地.
地是.
什麼?.
空虛混動.
冤冤黑暗.
神的靈在水面上運行.
是不是?.
你什麼時候聽到耶穌做過什麼?.
很多牧師說創造是聖父聖子聖靈一起做的.
是.
開會做的沒錯.
對.
但是在創世記裡.
從頭到尾我們看見.
在創世記裡創造的時候.
是父神和聖靈直接參與做的.
聖子做什麼?.
不是吃那些聖子.
聖子耶穌基督在做什麼?.
聖子耶穌是complete the whole creation.
怎麼complete?.
怎麼將它完成?.
蛇.
在伊甸園裡.
蛇引誘女人.
去違抗了.
神透過亞當所給的命令.
除了那棵樹上的果子不可吃.
第一個.
兒子.
亞當.
男人.
陷在罪當中.
直到什麼時候.
是完結了.
由上帝的兒子降生.
去到新約的時候.

$^{1241}$死在十字架上.
復活成了.
上帝的創造.
的修復點.
就在耶穌死在十字架上.
所以耶穌在這裡說.
各樣的事已經成了.
這裡是指到.
上帝整個的創造.
在這一刻complete.
跟著他的復活.
就是打開.
神和人在運行當中的.
另一個chapter.
另一個生活.
新天新地裡面的生活.
所以耶穌.
和亞當的關係是什麼.
亞當.
是首先的.
耶穌.
是在幕後的.
Omega.
Alpha and Omega.
耶穌包了其中.
不過神.
做人.
就是你和我.
可以有機會參與這個創造裡面.
耶穌說成了.
是告訴他.
他已經完完全全.
將這個創造.
完美地去.
砌造好.
砌造好出來.
能夠經歷.
罪的.
抵抗能力.
因此經過主耶穌基督.

$^{1281}$的拯救的人.
他能夠在創造裡面.
有超過一百級的.
抵抗能力.
去抵抗罪.
而罪也都不再沒有任何的影響.
你和我就可以參與在這個事情上.
怎樣參與?.
相信主耶穌.
就可以了.
所以.
十一節這樣說.
在此並不分希臘人.
猶太人.
受割禮的.
未受割禮的.
化外人.
西古提人.
為奴的.
自主的.
唯有基督是包括一切.
又住在各人之內.
說什麼?.
你會覺得很奇怪這句話.
說什麼?.
我不知道他說什麼.
尤其是我不知道什麼西古提人.
西古提人是什麼人來的?.
順帶說一下.
按照.
初期教會歷史學家.
做師弗所說.
西古提人這種人.
這種民族的人.
雖然我們今天不知道他實質是什麼.
在哪裡.
但是根據約瑟夫他所說.
西古提人.
他們的質素.
只是略比野獸好一點.

$^{1321}$即是用今天的說法.
比畜生好一點點.
這一大堆的人.
有分種族的.
有分在宗教上有沒有行過割禮的.
有沒有分他是哪一個類別的文化的人.
他的文化是野獸的文化.
還是有些真是學曉有文化的人呢?.
即是用今天說的.
他讀過書好還是沒讀過書好.
識字好還是不識字好.
他圍爐的好還是自由的好.
所有這些的人.
其實都歸入在基督的裡面.
什麼意思?.
不論你是是什麼的人.
只要你是能夠.
治死你身體地上的肢體.
那些邪情私慾.
能夠脫虛救我.
穿上身我.
肯承認自己罪過的人.
你就會包攬在基督的拯救創造裡面的一份子.
總之你能夠成為創造的其中一份子.
我成為創造的其中一份子.
只有一個因由.
就是因為我們肯脫虛救我.
穿上身我.
而這個過程就是一個誠實面對罪的人.
而哥羅西教會.
你們就是應該走在這條路上.
聖經告訴我們耶穌基督創造的角色.
是將整個創造完美地完成了.
當這個創造完美地完成的時候.
就只有一個結果.
榮耀神.
能夠反映上帝創造的目的.
彰顯我們神的榮耀.
不只是純粹一種心靈改變.
或者口舌上說得更加漂亮.

$^{1361}$是一個整全一個人的改變.
而我們要脫去一種膚淺的信仰.
以為在教會裡面做多些聖功.
你就等同於榮耀神.
你做福音工作你就榮耀神.
弟兄姊妹.
原來所做的所有工作.
要向自己和你身邊所有弟兄姊妹的罪去處理.
這樣才是榮耀神.
教會這麼多年來.
很膚淺地去理解榮耀神.
就是做多些聖功.
其實是對上帝的創造的一種侮辱.
以為沒有你和我做的事.
上帝就得不到榮耀.
對不起.
你和我做了多少事.
不是榮耀神的途徑.
你和我能夠榮耀神的是.
因為你和我都願意一起誠實地處理罪.
當我們肯面對我們的罪的時候.
歸回上帝原有的樣式.
就能夠充分反映出.
上帝在創世記裡面創世的原意.
最近大家都知道.
有輛車燒了引擎.
燒了引擎.
燒那一下其實都挺嚇人的.
全輛車燈都亮了.
又不想開車 又有很大陣燒的味道.
後來很心急地花了一筆錢去修理.
順帶說一句.
日後如果你的車有這麼大的事要修理.
千萬不要第一時間拿去修理.
首先要聯絡那輛車的車廠.
跟他討價還價.
因為後來我跟他討價還價.
他給我offer我非常好的.
不過我都給了人修理.
算了 做人要有口齒給人修理.

$^{1401}$那個人是貴的.
貴了幾千元.
沒辦法.
答應了人給人修理就給人修理.
車修好了.
我發覺有一樣東西.
扭力變強了.
有些人會跟我說.
我鄰居也這樣說.
賣了它吧.
反正它燒過修過.
可能未必那麼好.
不過我心裡知道那輛車.
我因為有駕駛過.
用得講得足一點.
那輛車很鏈了.
上高速公路是變強了.
可能因為它燒的位置.
本身是設計上有缺陷的.
現在換了很多.
弟兄姊妹.
有時有些東西.
錯了.
我們會覺得很內疚.
會覺得很羞愧.
但是如果有些東西是錯了.
而上帝容許的.
一定有個原因.
就是讓你在錯誤裡學習.
使自己不會再犯錯的話.
你會變得更加可貴.
尤其是在上帝眼中.
有人跟我說賣了它吧.
那輛車.
修理完之後.
可能七爐八箱有第二樣可以修理.
我明白的.
車可能真的是這樣.
但屬靈的生命.
上帝欣賞的.

$^{1441}$是你真的肯改變.
雖然我不太喜歡用俗語來講解聖經.
但是中國人有一句說話.
一個肯回頭的浪子.
是金不換的.
你說這是世人的說話.
不如我引用路加福音浪子的比喻.
浪子的比喻.
當量子散盡父親的家財.
可以說是.
阿爸有損失.
阿仔有損失的時候.
但是當兒子肯回來.
阿爸說了一句這樣的話.
浪子的爸爸說.
這是失而復得的兒子.
上帝最欣賞的是.
失而復得而變得更完美的兒子.
他的完美不是因為他犯過錯之後改過.
他的完美就是因為他犯了錯.
他以後不會再犯錯.
這個才是上帝最欣賞的生命.
弟兄姊妹.
俄羅斯教會.
保羅將一個這麼重要的訊息交了給他們.
今天我們有幸看到.
你有沒有改變你的生命.
活得越來越像上主的形象呢.
我們一起同心低頭禱告.
父神多謝你讓我們在這段聖經裡的學習.
看到保羅在遠處為了俄羅斯教會的成長而費心神.
他所寫的每一字每一句.
讓他們都可以細味.
屬靈的生命的栽培.
是需要上主你怎樣一步一步去掩護我們.
雖然我們或有閃失或有犯錯.
但是我們的生命卻從來沒有離開主你.
這個正是你最珍而重之最欣賞的地方.
求你給我們有這份毅力.
有這份屬靈的智慧.

$^{1481}$洞悉我們自己在哪一方面仍然去得罪上帝你的.
是在罪裡面受捆綁的.
讓我們去將祂致死在這個過程當中讓我們榮耀上帝你.
禱告感恩奉靠主耶穌基督得勝名至求.
阿們.
\newpage



\section{箴言 22:6}
\label{sec:032116}
\textbf{如何把兒女帶到神面前}~~~~劉彤牧師
\newline
\newline
連結: \href{http://www.church.com.hk/acms/content.asp?site=cdc\&op=show\&type=product\&code=032116\&layout=sermon}{\texttt{032116}} ~~~~ 語音日期: 2026-05-10
\newline
\newline
\hyperref[sec:032115]{\small{< < < PREV SERMON < < <}}
~
\hyperref[sec:index_chronic]{\small{[返順時目]}}
~
\hyperref[sec:032117]{\small{> > > NEXT SERMON > > >}}
\newline
\newline
箴言 22:6
\newline
\begin{longtable}{cl}
\hline
\hline
章節 & 經文 (和合本修訂版)\\
\hline
22:6 & \begin{tabularx}{0.7\textwidth}{X} 教養孩童走當行的道，就是到老他也不偏離。 \end{tabularx} \\ \\
[1ex]
\hline
\hline
\end{longtable}
$^{1}$很多時候我們領受了母親的愛.
常常習以為常.
不覺得有什麼特別.
每年我們習慣性的就是慶祝一下母親節.
寫張卡片.
喫頓飯就算了.
但事實上母親在我們一生當中所付出的愛.
是遠超過我們所能領會的.
我就想到我的母親.
從我小時候開始.
我就記得她每天總是在床前為我不禱告.
一直到我長大.
甚至做了牧者.
她還是每天為我在禱告.
甚至我相信現在她已經被接到天家去.
現在在天上.
恐怕她也是天天在為我禱告.
Amen.
讓我們從小長大.
母親的愛跟掛念.
總是永不止息的.
所以讓我再一次真的.
再一次向我們中間每一位偉大的母親說一聲.
謝謝你們.
是因為你們的愛.
讓我們能明白上帝的愛是何等強闊高深.
Amen.
有人說在現今的世代要教養兒女.
是一件叫人謙卑的功課.
我相信每一位現代的父母.
我們都深有同感.
時代真的不同了.
教養兒女越來越不容易.
雖然隨著科技的發達.
我們的生活越來越方便.
但另一方面我們面對的世界.
卻是越來越複雜的世界.
充滿了各種的選擇.
以前我們做父母的.
在管教兒女的事情上.

$^{41}$我們是無上的權威.
可是現在只不過是眾多選擇中的一個而已.
難怪很多的父母.
在教養兒女的事情上.
都深深感覺一種無力感.
我相信我們做基督徒的父母.
我們都很渴望把兒女帶到神的面前不是嗎.
那怎麼樣能夠在現今的世代.
把我們的兒女帶到神的面前來呢.
今天早上從四方面跟大家來分享.
首先第一個.
不要輕忽你對兒女的影響力.
跟你媽說不要輕忽你對兒女的影響力.
雖然今天我們的世代是越變越複雜.
雖然在兒女成長的過程當中.
我們感覺我們的影響力越來越小.
但我還是要提醒你.
你在你兒女身上的影響力還是很大的.
其實隨著年日的過去.
回頭看一下你的兒女.
他們不是充滿你的影子嗎.
他們的長相.
他們的個性.
他們的喜好.
他們的習慣.
不是就是你的樣子嗎.
你也是一樣啊.
如果你仔細看看你自己.
無論你察覺與否.
喜歡與否.
你多麼的像你的父親跟你的母親不是嗎.
你知道我的父親已經過世很多年了.
可是現在我回頭看.
我是越來越像我的父親.
你知道嗎.
他的長相.
他的表情.
他的個性.
他的舉動.
太像了.

$^{81}$連頭髮都一模一樣.
有一次我看到一張在非洲的照片.
忽然間看到我怎麼爸爸也在其中.
嚇我一跳.
仔細一看是我自己啊.
再看看我們自己的兒女.
不也是這樣嗎.
其實很多時候.
我看見我的女兒.
她們真的跟我跟師母都很像.
雖然她們自己不見得喜歡如此.
但也沒有辦法.
也沒有人故意安排她們一定要這樣.
但就是自然而然成長的結果嘛.
這些年我真的回頭來看的時候.
發現我兩個女兒在成長的過程當中.
無論我們中間過去有多少的爭吵.
無論有多少的歧見.
不能溝通.
但至終回頭一看.
她們真的是踩著我們的腳印在向前走.
所以作為父母的人.
當我們面對兒女成長的時候.
我真的要說.
不必先太擔心兒女.
你先要看重的是.
你自己先要好好為主而活.
Amen.
跟你媽說你好好為主而活.
跟她說.
因為有人說孩子在12歲之前.
他會照你所說的去做.
12歲以後.
他會照你所做的去做.
換句話說.
你的榜樣還是太重要了.
如果你希望你的兒女愛神.
那你自己最好先來愛神.
如果你不希望你的兒女.
只是一個一個禮拜做.

$^{121}$半天的基督徒.
是主日半天的基督徒的話.
那你自己最好做一個全時間的基督徒.
Amen.
若是你希望你的兒女.
心中有主的火來燃燒.
而不是世界上的火.
那你自己最好先讓主的火.
在你生命先燃燒起來.
不是嗎.
多少時候我們在教會里面.
很容易批評論斷別人說.
這個人很愛主.
那個人不愛主.
但真正的問題是.
請問你愛主嗎.
所以Oscar Smith.
他曾經說過一句話.
他說福音書.
其實一共有五本.
你說不是隻有四本嗎.
沒錯是四本.
馬太,馬可,路加,約翰.
第五本就是你了.
你就是第五本福音書.
一般人看的是前四本.
你身旁周圍人看的福音書.
就是你這本福音書了.
許多時候我們常常急於改變我們的兒女.
但事實上今天.
我們的兒女正在看著我們.
看著我們的生命的榜樣.
你可知道今天在我們生命當中.
你的兒女是怎麼看你嗎.
請問你在兒女眼中是怎麼樣一個父母呢.
更重要的.
請問你在你兒女眼中.
是怎麼樣一個基督徒呢.
不要忘記.
你在你兒女身上的影響力是很大的.

$^{161}$你的兒女還是要從你身上.
來認識到神的真實.
所以第一個.
不要忽略你對兒女的影響力.
第二.
要按照主的教訓教導兒女.
跟你說按照主的教訓教導兒女.
跟他說.
你說要教導些什麼呢.
我們看見一個人在成長的過程當中.
有三方面的需要.
第一個他需要認識他的神.
第二他需要建立一個正確的自我形象.
第三個他需要明白他的社會責任.
同樣今天我們需要從這三方面.
來教導我們的兒女.
首先第一個.
我們要教導他們來認識神.
跟你們說教導他們認識神.
作為一個父母.
我們需要明白神所託付給我們.
教導兒女的最主要的目標是什麼.
我相信神讓我們教導兒女最主要的目標.
不是隻是讓他們出人頭地而已.
神是要我們教導他們.
成為一個屬靈的人.
跟你說前書第二章十五節.
保羅講得很清楚.
什麼是屬靈的人.
一個屬靈的人能夠看透萬事.
卻沒有一個人能夠看透了他.
這裡說因為屬靈的人.
是能夠看透萬事.
他能夠不被黑暗的權勢.
黑暗的世代所迷惑.
要注意在這裡沒有說看透萬事.
好像看透看破萬事一樣.
好像看破紅塵一樣.
是看透萬事.
他能夠不被黑暗的世代所迷惑.

$^{201}$反而他能夠明白神的心意.
他知道什麼是對.
什麼是錯.
在這樣一個錯綜複雜的一個世代當中.
這不是今天我們所最需要的嗎.
因為有一天孩子會長大的.
當他長大的時候.
他們所需要面對這個世界上的挑戰.
他們所要面對這個世界上的選擇.
很多都是我們做父母愛莫能助的.
無論今天你送他到多好的學校.
你送他到多好的學區.
無論你送他到私立學校.
到基督教學校.
甚至自己叫homeschooling.
我告訴你沒有一件事情.
能夠保證他們在生命當中做正確的選擇.
而只有在他生命裡面有神的同在的時候.
神的靈就要在他每天生活當中.
親自來引導他們.
所以在我們中間做父母的.
一定要鼓勵你的小孩子來到教會.
不要因為課業來攔阻他.
你讓他在教會當中經歷到神的同在.
不只是有屬神的知識而已.
而是能夠經歷神的真實.
因為一旦他經歷神的真實.
他的價值觀就會被神所調整.
他就不會轉向世界.
你也讓他參加教會的服飾.
不要嫌花太多時間.
想想看你平常要去學多少東西.
又要學鋼琴 又學游泳 學足球.
最終還學中文.
花多少時間.
你說我明白服飾也是很重要的.
只有服飾的過程當中.
他的屬靈生命纔可能成長.
他纔可能體會神的心意.
一個親至成熟的孩子.

$^{241}$他就不容易被這個世界所影響.
更重要你要幫助他參加教會的小組.
因為在孩子成長的過程當中.
同伴是很重要的.
朋友是很重要.
若是有好的屬靈同伴.
就能夠在主的道路上快跑來跟隨神.
所以第一個要教導他認識神.
第二個教導他有建立正確的自我形象.
跟你們說要幫助他建立正確的自我形象.
跟他說.
你知道在美國我看見華人的孩子.
一向自我形象都很低.
我們真的需要幫助我們的孩子.
讓他能夠建立一個正確的自我形象.
你說怎麼樣幫助他們.
首先第一個.
一定要關心他們所關心的.
跟你們說關心他們所關心的.
有時候我們做父母的太忙碌.
不常關心孩子們正在做的事.
或者我應該說.
我們常只關心我們所關心的事.
我們只看重他們上學,讀書,放學,做功課.
有的喫有的睡就好了.
我們很少關心他們所關心的事情是什麼.
不關心就表示不看重.
以至於他們也不覺得自己的生命是有意義有價值的.
其次要常常鼓勵他們.
跟你們說要常常鼓勵他們.
因為人都是需要鼓勵的嘛.
不但是做得好的時候需要鼓勵.
做得不好的時候更需要鼓勵.
就像打球的啦啦隊一樣嘛.
打到的時候說棒極了.
打得不好的時候說It's ok.
下次再來.
千萬不要一天到晚拿你的孩子跟別人來比較.
因為中國話說人比人怎麼樣.
氣死人.

$^{281}$而且非常傷他們的自尊心.
我常常說如果你們家有老二.
特別要鼓勵做老二的.
因為老大跟老二成長過程不大一樣.
做老大的無論做什麼.
長大過程當中都是驚喜萬分.
老二無論做什麼都是理所當然.
感覺不被重視.
要常常鼓勵做老二.
第三要給予適當的責任.
給你們說要給予他們適當的責任.
責任是很重要.
當一個人有責任的時候.
他就覺得自己很重要.
就算很簡單的事.
無論是洗個碗啊.
或者是倒個垃圾啊.
都會讓孩子覺得他在家裡面是一個不可少的一份子.
事實上當一個孩子懂得負責任的時候.
就代表他正在成長.
第三我們要教導他明白自己的社會責任.
跟他們說要讓他明白自己的社會責任.
在傳統上啊.
中國人一向是有個人至少門前雪.
莫管他人瓦上霜的這種概念.
在這種觀念之下.
我們很少去教育我們的孩子.
去關心社區的需要.
在美國第二代的華人更是如此.
我們看見我們美國第二代的華人有很多個人的成就.
有得到科學獎的很多優秀的一些博士跟醫生都有.
但卻很少有影響力的社區領袖.
根據統計在美國.
他們把每一個族群的一些社區領袖跟他們的人數相比.
在各種族群當中.
華人的比例是最低的.
我相信這是因為我們缺乏教導.
我們下一代關於我們的社區跟我們對社區的責任的緣故.
但事實上今天我們的信仰是最注重社區的.
主耶穌為什麼道成肉身來到這個世界.

$^{321}$耶穌基督道成肉身來到這個世界.
就是我們進入社區.
進入人群最好的榜樣.
我們應當以身作則.
帶領我們的孩子走進學校.
走進社區一同參與.
一同關懷好讓神的恩典跟愛.
從我們生命當中能夠流到社會里面.
所以我們開始求神來幫助我們.
能夠真的在每天的生活當中.
按照主的教訓教導我們的兒女.
第三也是非常重要.
要相信神今天仍然能夠工作.
跟你們要相信神今天仍然能夠工作.
跟他說.
以佛手術第三章二十節怎麼說呢.
所以說神能造成運行在我們心裡的大力.
充充足足成就一切超乎我們所求所想.
哈利路亞.
換句話說.
神今天仍然能夠工作.
神今天仍然在我們生命當中運行.
Amen.
神今天仍然能夠將活潑的信心.
放在我們兒女的心中.
神今天仍然能夠將屬靈的夢想.
放在我們兒女的心中.
Amen.
第二點.
你一點不要被現今黑暗的世代所欺騙.
叫你覺得心裡害怕.
充滿憂慮.
我們的神今天仍坐在寶座上.
今天他仍然能夠工作.
他仍然能夠突破黑暗的權勢.
在我們心靈深處來運行.
Amen.
在後書第一章十二節怎麼說了我們一同來唸.
為這緣故我也受這些苦難.
然而不以為恥.

$^{361}$因為知道我所信的是誰.
也深信他能保全我所交付他的直到那日.
哈利路亞.
這裡講的進入我們的神.
不但是全能的神.
而且他能保全我們所交付他的直到那日.
這裡不是說我們.
而是說我的託付.
是我所交託他.
我所託付給神的.
你可知道神多麼看重你嗎.
你可知道神多麼看重你所交託給他的事物嗎.
而且神在應許我們說.
他要保全我們所交付他的直到那日.
這個保全的意思.
就原文裡面是護衛.
神要護衛你所交託給他的直到那日.
換句話說直到主再來的日子.
第二點一旦你明白神的應許這樣的寶貴.
你就要抓住神的應許.
為你的家人來禱告.
為你的兒女來禱告.
也許你說穆斯只好禱告了吧.
不是最後沒有辦法的辦法.
禱告是最好的辦法.
跟你媽說禱告是最好的辦法.
你說怎麼禱告法呢.
我鼓勵你每一天.
要為你的孩子做三個禱告.
第一個求神賜下盼望的靈.
跟你媽求神賜下盼望的靈.
因為多少時候你的孩子的眼睛是很短淺.
他們只能看見現在.
他們看不見將來的願景.
他們只看見眼前的挫折.
卻看不見神在他們生命當中美好的計劃.
所以我們一定要求神.
不斷將盼望的靈充滿他們.
讓他們在人生的道路當中不至於絕望.
也許人會叫他們失望.

$^{401}$但神永遠不會叫他們失望的.
我們要求神賜下盼望在他們的心中.
讓他們確信他們生命是有明天的.
而且明天在神的手中.
第二要求神賜下信心的靈.
跟你媽求神賜下信心的靈.
多少時候當孩子們面對周圍的環境.
跟同學的壓力.
常常會將他們心中害怕失去信心.
所以他們覺得沒有力量.
去勝過罪惡跟試探跟引誘.
所以我們一定要不斷求神.
把信心的靈充滿在他們心中.
好讓他們能夠持守從神那裡來的信心.
能夠相信無論環境如何.
靠著那加給他們力量的.
他們凡事都能做.
求神賜給他們耶穌那樣的信心.
我們常說現今的世代年輕是耶穌的世代.
就是要有耶穌的信心.
面對約旦河.
相信神在大水當中為他們開道路.
面對耶利哥.
相信神能夠為他們拆毀一切堅固的隱悅.
我們要求神把信心的靈充滿在他們的心靈中.
第三要求神賜下智慧的靈.
跟你媽求神賜下智慧的靈.
多少時候當孩子們在面對各樣擺在他面前的選擇.
他們不知道如何選擇的時候.
我們要為孩子們來禱告.
不但是為他們的功課禱告.
為他們成績禱告.
為他們申請大學禱告.
跟著求神賜下智慧的靈.
為什麼因為真言四章七節講得很清楚.
真言四章七節說智慧為首.
所以要得智慧在你一切所得之內就必得聰明.
你要求神不斷的用智慧的靈來充滿他們.
因為你不可能每時每刻與孩子在一起.
但神的靈卻能夠與他們同在.

$^{441}$神賜下智慧.
讓他們每天每時每刻能夠正確的時機.
做正確的決定.
所以我鼓勵你每天當你孩子出門的時候.
你孩子出門的時候.
每一天為他們祝福禱告.
求神賜下盼望的靈.
賜下信心的靈.
賜下智慧的靈.
甚至當你的兒女遠離神悖逆神的時候.
你仍然要繼續禱告.
因為深信神能夠保全你所交付他的.
直到那日.
你們.
我就想起奇異恩典那首詩歌.
它的作者是John Newton.
他出生在英國.
早年的時候他喫喝嫖賭無惡不作.
讓他父母非常的傷心.
但他的母親卻持續不斷的為他禱告.
後來他到非洲去做販賣黑奴的工作.
結果沒想到他被非洲人抓住.
後來反而成為非洲人的奴隸.
他的父親聽見就非常的著急.
想盡各樣的辦法.
最後終於把他從非洲人手中救出來.
沒想到在他回到英國的船上的時候.
忽然遇見大風暴.
船幾乎要沉到海里去.
就在這個時刻.
他在甲板上就想起他母親的禱告.
他想起他母親流淚滿面為他祈求.
他就跪在甲板上求神赦免他的罪.
並且將他的生命奉獻在神的手中.
結果神就隨聽他的禱告.
神搭救他.
讓他回到英國.
當他回到英國.
他就全時間奉獻來服侍神.
後來成為神非常重要的傳道人.

$^{481}$寫下這首奇異恩典的詩歌.
是流傳千古的詩歌.
只因為他的母親相信神.
今天仍然能夠工作.
Amen.
所以同樣.
今天你不要灰心.
無論你的光景如何.
你的家庭光景如何.
你與兒女的關係是如何.
你要相信神今天仍然能夠工作.
Amen.
最後第四.
要陪伴你的兒女成長.
跟你們說要陪伴你的兒女成長.
成長是一個過程.
在這個過程當中.
你知道常常做父母跟做兒女.
看同一件事情角度是不同.
對一個成長的兒女來講.
他們所看見的永遠是過程的本身.
做父母的.
永遠看見將來美好的結果.
所以這為什麼要叫小孩子去讀書.
做功課那麼困難的緣故.
因為他們所看見的只是現在的辛苦.
做父母看見的是將來美好的願景.
所以對父母來講.
明天這麼重要.
為什麼今天不能夠忍耐一下呢.
對小孩來講.
今天都過不去了.
還談什麼明天呢.
難怪怎麼講都講不通了.
其實都是成長的過程.
在這過程當中.
學習忍耐是很重要的.
跟你媽說要學習忍耐.
特別在兒女成長的過程當中.
你總是覺得他們成長好緩慢.

$^{521}$怎麼講怎麼還是不懂事.
但事實上.
現在也許你覺得他成長很緩慢.
但其實他們一生的過程當中.
一生的歲月當中.
做孩童的這些時期.
還是非常短暫的.
而且時間真的非常快速會過去.
不是嗎.
昨天他還在你大腿上玩耍.
一轉眼今天就要大學畢業了.
我知道很多父母.
你的孩子恐怕下禮拜就要畢業典禮了.
要學習忍耐.
陪伴他們成長是很重要的.
我常說就像聖經提醒我們.
像一個農夫忍耐等候地理的出產一樣.
千萬不要壓苗助長.
什麼叫壓苗助長.
就是農夫看見這個秧苗長得那麼慢.
就想幫他一下.
幫他拔一拔幫他成長.
一拔就死掉了不是嗎.
同樣的.
今天做父母的求神幫助我們.
不要急於用自己的方法.
勉強你兒女成長.
聖經上說要各按其時叫他們成長.
太重要.
各按其時這個句子太重要.
表明每個孩子都有他自己成長的時刻.
跟你們說每個孩子他自己成長的時刻跟他說.
不要著急.
忍耐等候.
只要我們存著接納的心.
忍耐的心去陪伴他們.
自然而然他們必然在神的時候長大成熟了.
今天無論你的家庭光景是如何.
無論你與孩子的關係是如何.
今天你都要再一次憑著信心.

$^{561}$把你的家庭把你的兒女交託給神.
因為今天我們所信的神.
仍然能夠使父母的心轉向兒女.
使兒女的心轉向父母的神.
只要你來信靠他.
順從他的引導.
神必然能夠在你家庭當中.
重新帶來恢復的恩典的.
Amen.
我們一起禱告.
親愛的主我們再次來到你面前.
謝謝主真的是紀念我們一切的需求.
主我們也謝謝你.
你真的保全我們所交付你的.
能夠直到那日.
何等寶貴的應許.
今天就把你的應許深印在我們心中.
求你祝福在我們中間的每一個家庭.
祝福在我們中間的每一個母親們.
主常常我們為家庭兒女憂慮的時候.
求主親自成為我們的安慰.
成為我們的力量.
成為我們的幫助.
也讓我們真的懂得交託仰望在你手中.
讓我們依靠你.
在每天生活當中.
就能夠經歷你的恩典.
經歷你的大能.
謝謝主.
Amen.
\newpage



\section{以弗所書 5:21}
\label{sec:032117}
\textbf{與子偕老的見證}~~~~招世超牧師
\newline
\newline
連結: \href{http://www.church.com.hk/acms/content.asp?site=cdc\&op=show\&type=product\&code=032117\&layout=sermon}{\texttt{032117}} ~~~~ 語音日期: 2026-05-10
\newline
\newline
\hyperref[sec:032116]{\small{< < < PREV SERMON < < <}}
~
\hyperref[sec:index_chronic]{\small{[返順時目]}}
~
\hyperref[sec:032118]{\small{> > > NEXT SERMON > > >}}
\newline
\newline
以弗所書 5:21
\newline
\begin{longtable}{cl}
\hline
\hline
章節 & 經文 (和合本修訂版)\\
\hline
5:21 & \begin{tabularx}{0.7\textwidth}{X} 要存敬畏基督的心彼此順服。 \end{tabularx} \\ \\
[1ex]
\hline
\hline
\end{longtable}
$^{1}$我最近看到一個Barnett Research.
大家知道Barnett Research是一個很受歡迎的調查研究機構.
這個研究說的是什麼呢?.
他告訴我們有47\%的基督徒母親.
47\% 差不多一半的基督徒母親表示.
他們的教會從來沒有提供過幫助他們做媽媽.
都可以嗎? 預備吧.
從來沒有提供過怎樣幫助做媽媽的一些材料或者訓練.
當我看到這個報告時.
我自己有些反思.
想想我們教會又如何呢?.
我記得在我們教會開創的期間.
我們曾經搞過一個叫做醒目父母系列.
我們請一些專家來教我們父母怎樣做父母親.
我們又搞過代產的母親班.
師母和姐妹比較有經驗的.
有很多快要生的父母.
我們怎樣教他們做呢?.
很緊張 怎樣做一個代產的母親.
準備去迎接嬰孩的出生呢?.
我記得大概在2010年.
我們推行一個全教會的門徒訓練時.
我一個人帶了三班.
一個是弟兄班 一個是姐妹班 一個是混合班.
姐妹班很特別 為什麼呢?.
全部都是媽媽.
那時候很挑戰的.
媽媽要說服老公.
星期三晚留在家裡帶孩子睡覺.
我還記得爸爸說要穿著老婆的睡衣.
躺在兒子旁邊.
兒子才能睡得著.
每個星期三晚.
有一個媽媽有三個小朋友.
放下小孩 放下老公.
回來接受門徒訓練.
我們過得很好.
除此之外.
我們還有什麼為了我們當中的父母或媽媽做的呢?.
其實沒有多少.

$^{41}$我最能夠想到的可能是十幾二十年前.
我們曾經搞過一個叫做BonBonFamilyCamp.
家庭營.
自此之後就好像沒什麼了.
所以當我看到這個報告的時候.
我真的不好意思.
我要跟我們的媽媽說聲對不起.
因為似乎是忽略了你們的需要.
事實上我心裡面其實都是記掛了這件事.
特別剛才這個姐妹也提到.
我們說Intergenerational.
其實我們很希望整個家庭都能夠去認識上帝.
敬拜上帝.
事奉上帝.
現在我們這個講道系列是說什麼呢?.
就是怎樣能夠幫助你的家庭.
變成一個訓練門徒的地方.
怎樣能夠培育我們的下一代.
成為耶穌基督的門徒.
而我們做父母的.
做夫妻.
在子女面前的表現.
就至關重要.
不知道我們當中.
做父母.
有沒有想過.
在家裡.
我們夫妻之間的關係.
對我們的子女.
或者在家裡.
其他人的信仰.
會有什麼影響呢?.
不知道大家有沒有想過這個問題.
今天早上我先掉了書包.
中國有一個書叫《詩篇》.
大家有沒有聽過.
《詩篇》裡面有一句話.
要打出來.
因為讀不清楚字眼.
這個是《詩篇》說.

$^{81}$死生契闊.
很立刻.
死生契闊.
與子成說.
後面那兩句大家熟悉一點.
執子之手.
與子皆老.
其實我們只記得.
開頭最後那兩句.
執子之手.
與子皆老.
給我們的感覺.
很浪漫.
拖著老婆老公的手.
一路走到白頭.
似乎給我們的印象.
我們愛情的長久.
永結同心.
一個很浪漫的期盼.
我們一般的感覺.
看法就是這樣.
但其實前面那兩句.
你發覺原來.
這句話的意思不是這樣.
意思的重點是說.
開口要立志.
是一個志向.
在人生漫長的旅途當中.
兩夫妻不離不棄.
彼此守護一個承諾.
其實是一個承諾.
一個promise.
不只是浪漫的感覺這麼簡單.
在這個過程當中.
大家知道.
如果結婚不是一兩年.
或者未過七年.
結婚之後.
你會發覺.
那種愛的感覺會昇華.

$^{121}$還有沒有那種feel.
你看見老婆不會心跳.
但對配偶的承諾和付出.
是可以不斷提升的.
大家明白我的意思吧.
但問題是什麼呢.
我們很現實的.
看回現實環境.
看回我們所認識的朋友.
我們會發覺.
越來越多的婚姻.
是功能失調的.
功能失調的婚姻.
舉目皆視.
甚至基督徒夫妻.
似乎都不可以倖免.
大家都知道.
我做了牧師幾十年.
有很多經驗.
什麼經驗呢.
我看過不少夫妻之間的吵架.
這些是最輕微的.
冷戰.
還有家暴.
離婚有沒有.
當然有.
甚至離家出走的案件.
他們都是這樣.
我在想.
如果有兒女在當中.
看到我們父母是這樣的時候.
在這樣的環境裡長大.
他們會怎樣看待教會呢.
他們會有一天回來教堂嗎.
他們會怎樣看待信仰.
怎樣去解讀.
他們想不想去做一個.
好像父母那樣的基督徒呢.
或許我這幾個問題.
比較沉重一點.

$^{161}$比較嚴肅一點.
但其實這是一個很嚴肅的問題.
如果我們還有配偶在身邊.
當我拖著我的太太.
拖著我的老公.
邁向白頭的這條路上.
不知道我們有沒有嘗試過.
刻意地去互相建立.
讓我們能夠一起成長.
抑或我們只是見步行步.
結了婚.
似乎是這樣.
隨遇而安.
這是我最喜歡說的話.
今天我想和大家看看.
我們夫妻之間應該怎樣相處.
我們應該怎樣去對待對方呢.
先看一節.
第五章二十一節.
可不可以一起讀一讀.
五章二十一節.
一起讀一讀.
你們要彼此信服.
因為你們是敬畏基督的.
再來多一次.
你們要彼此信服.
因為你們是敬畏基督的.
當我聽到信服這個字的時候.
在今天的世界社會當中.
會覺得很不舒服.
其實不只是今天.
在耶穌的時代更加嚴重.
為什麼呢.
耶穌的時代在羅馬帝國統治下.
不只是羅馬帝國統治.
其實那時是階級制度盛行的時代.
所以信服是單向的.
什麼叫單向呢.
僕人要信服主人.
下屬要聽上司的話.

$^{201}$子女要信服父母.
妻子都要信服丈夫.
所以彼此信服這個字.
這個說話.
對當時的人來說.
是一個很新穎,窄身的.
和難以接受的觀念.
我剛才提過.
不要說耶穌的時代.
就算今天.
我相信很多人都不能夠接受.
我們不明白.
為什麼聖經會有這樣的教導.
為什麼聖經會教基督徒.
其實如果我們看信服這個字.
在聖經裡面.
其實不是說要盲從.
或者覺得自己很弱.
所以要聽別人的話.
希臘文這個字.
信服這個字的原意.
其實就是排列在某某人的底下.
是一個排列.
是自願的.
將自己放在某個人的領導和照顧的下面.
所以重點就是即是和秩序.
而不是地位身份高低價值的大與小.
或者是誰比較厲害.
聖經強調的是什麼.
我再說一次.
是即是和秩序.
就好像當軍的.
現在打仗.
裡面有不同的階級.
有將軍有上校有少校.
將軍和士兵.
在人格上完全平等.
同不同意嗎.
將軍有些很好品格的人.
有些不好品格的人.

$^{241}$士兵一樣.
將軍和士兵在人格上.
但是為了要完成任務.
所以有不同的崗位.
有不同的秩序.
去安排.
所以其實馴服這個字.
在聖經裡面.
其實是一個有深刻神學意義的字眼.
保羅怎樣解釋馴服.
特別是彼此馴服.
他說因為我們敬畏忌妒.
他有解釋.
為什麼.
因為老婆也好.
老公也好.
我們都敬畏忌妒.
所以你看到.
馴服彼此馴服.
不是對方比我們強.
也不是對方比我們完美.
是因為對方.
無論是老婆還是老公.
和我們一樣.
有同一個主.
就是基督的耶穌.
所以夫妻可不可以比較.
可以.
不過不應該.
我知道很多夫妻.
在暗地裡.
自己在競爭.
看誰賺錢多一點.
誰在公司職位高一點.
在家裡誰當家作主.
誰應該聽誰的話.
但如果我們明白.
聖經說到彼此馴服這個字的意思.
我們就不應該這樣想.
聖經所教導我們.

$^{281}$就是要心甘情願地.
去互相扶持彼此.
幫助大家互相尊重.
將對方的需求.
看作是自己的需求.
一樣重要.
如果是這樣的話.
所以我們會盡量去滿足.
我們配偶.
我們的妻子.
我們丈夫的需要.
為了什麼?.
要建立對方.
讓對方能夠成為.
神所預立的人.
讓神因為這樣的緣故.
得到榮耀.
我剛才提過.
現在我們說什麼?.
說家庭門訓.
在家裡的門訓.
家裡的人.
今天我想跟大家說的是.
我們做夫妻.
包括父母和長輩.
在家裡的.
首先要學.
怎樣幫助自己的配偶.
成為耶穌基督的門徒.
就是說.
我們做丈夫的.
要學怎樣幫妻子成為耶穌基督的門徒.
而我們做妻子的.
同樣要幫助丈夫.
成為耶穌基督的門徒.
這就是互相建立的意思.
如果我們在座有單身的.
無論什麼情況.
單身也好.
我鼓勵大家.

$^{321}$找一個同性的屬靈夥伴.
和他建立一個.
互信互助的關係.
於是我們可以彼此幫助.
學習成為一個.
土主門徒.
剛才說到夫妻之間.
要學習互相建立.
聖經是怎樣說的呢.
我請在座的姊妹.
和我一起讀.
姊妹先讀.
第五章二十二節開始.
一起讀.
你們作妻子的.
當信服自己的丈夫.
如同信服主.
因為丈夫是妻子的頭.
為何沒有人讀.
如同基督的頭.
他又是教會的教主.
教會怎樣信服基督.
妻子也要怎樣.
凡事信服自己的丈夫.
讀完之後.
不知道有什麼感覺.
2011年的時候.
耶魯大學有一個.
法學博士教授.
叫做蔡美儀.
Amy Chow.
她出版了一本書.
這本書叫什麼名字呢.
《虎媽的戰歌》.
Bed to him of the tiger mother.
聽到虎媽這個字.
對我們今天的人來說.
是不陌生的.
但那時候是第一次.
出現在她那本書裡面.

$^{361}$現在為何我們不陌生呢.
因為已經成為一個文化的標籤.
聽到這個字的時候.
你怎樣去解讀.
什麼叫虎媽呢.
通常虎媽是我們形容什麼呢.
例如亞洲家庭裡面的媽媽.
對小朋友的管教很嚴格.
為什麼呢.
望子成龍.
這些媽媽通常會.
用一些比較高壓.
嚴厲的手段或者態度.
要求孩子.
在學業,在才藝.
比如鋼琴,小提琴,運動.
各方面都要有卓越的表現.
我剛才提過了.
現在虎媽這個字.
已經是被廣泛去形容.
對子女嚴重操控的媽媽.
我告訴大家.
通常這些媽媽.
都是很能幹.
都是做事很厲害.
很有自己的主見.
所以她才覺得.
我們的孩子應該怎樣怎樣.
不知道大家知不知道.
跟虎媽相對.
用來形容爸爸的.
是叫什麼.
有沒有人知道.
虎媽我們聽過了.
那Daddy呢.
怎樣形容Daddy呢.
你猜不到吧.
貓爸.
喵喵喵.
貓爸.

$^{401}$為什麼呢.
很多爸爸原來在虎媽面前.
變成貓.
一樣溫馴可愛.
所以你有沒有發覺.
很多孩子都很喜歡.
伏在Daddy的肚腩上.
舒服嘛.
事實上.
我知道很多家庭.
包括Galto的家庭.
當家作主的真的是那個媽媽.
是姊妹來的.
我剛才說過了.
其實他們都很厲害.
爸爸回到家.
很多時候什麼都不做.
要做什麼呢.
要不就躺著休息.
要不就上網.
要不就打個電視機.
有時間中都要抱著孩子去看球賽.
所以當時的經說.
你要做妻子的.
要信服自己的丈夫.
如同信服主.
如果老公回到家.
好像爛泥.
一坨東西.
躺在沙發上.
什麼都不肯幫忙.
可能真的令到.
做妻子的嘔血.
有沒有試過?.
姐妹們.
希望沒有吧.
但聖經說.
因為丈夫是妻子的頭.
如同基督是教會的頭.
如果大家記得.

$^{441}$剛才我解釋給大家知道.
聖經裡面說到.
彼此信服這個字的時候.
意思不是說到地位的高低.
而是說到夫妻都同有一個主.
就是基督耶穌基督.
所以我應該明白.
其實這裡的教導.
保羅不是叫我們做妻子的.
讓丈夫怎樣做都好.
不要理他.
任意妄為.
不是.
相反.
做妻子要積極想方法.
怎樣幫助你的丈夫.
成為一個好丈夫.
一個好爸爸.
一個好的基督徒.
妻子對丈夫的信服.
是基於對丈夫的愛和尊重.
而更重要的是什麼呢?.
當妻子明白到.
原來神創立家庭.
是不是神創立家庭?.
神設立婚姻,神創立家庭.
整個設計就是這樣.
我們是自願.
我希望我的姊妹.
能夠自願將自己.
放在丈夫的領導下.
但是我們要有意識.
有方法地去幫助自己的丈夫.
按神所設計的秩序去.
能夠鼓勵丈夫學習平等.
做好爸爸.
做一家之主.
做屬靈領袖的責任.
妻子能夠這樣做.
我希望你知道.

$^{481}$其實你是在建立自己的丈夫.
讓丈夫成為一個.
神所預立基督的門徒.
我記得有一次.
有一個姊妹請我去家裡吃飯.
當然開心.
特別是師母喜歡吃東西.
她和她的丈夫都很尊重牧師.
當然請牧師來吃飯.
很開心的.
吃飯前.
「牧師,借飯祈禱」.
「吃飯了,吃飯了」.
我給她一個措手不及.
我轉頭望向她的丈夫.
「誰是一家之主?」.
「一家之主」.
「The man of the house」.
「出來借飯祈禱」.
她的丈夫聽到震了一下.
「不用了,牧師,你來吧」.
我笑著對他說.
「你搞定了,你是一家之主」.
我告訴你.
多做兩次.
以後他再請我和師母去吃飯的時候.
不用我叫.
「The man of the house」.
自動就會向你借飯祈禱.
親愛的姊妹.
我對姊妹說.
親愛的姊妹們.
我鼓勵大家.
回家試一下.
幫你的丈夫成為一家之主.
他的祈禱可能沒你那麼順口.
嘴巴沒你那麼聰明.
他處理家裡的事.
可能沒你那麼有潮不敏.
他安排家裡的活動.

$^{521}$可能沒你那麼細心.
但是可不可以.
給了你的意見和看法之後.
讓他去做主思純呢?.
因為聖經說.
丈夫是妻子的頭.
如同基督是教會的頭.
基督也是教會全體的救主.
目的是什麼呢?.
要幫助你的先生.
成為神所閱立.
基督所喜愛的門徒.
聖經說.
教會怎樣信奉基督.
妻子又要怎樣凡事信奉丈夫.
什麼意思呢?.
教會是怎樣凡事信奉基督呢?.
我講三件事.
第一.
教會對基督信奉.
是基於耶穌基督對待愛.
基於愛.
要聽話.
因為愛他.
所以我們要聽他的話.
第二件事.
教會的決策.
導向和做事的方式.
應該以耶穌基督的教導為依歸.
而不是憑.
牧師夠老,夠經驗,夠智慧.
NO.
教我們的傳統.
NO.
經驗.
NO.
按照聖經的教導.
第三.
教會要讓神在耶穌基督裡面得到榮耀.
要讓神得榮耀.

$^{561}$所以做妻子.
要信奉丈夫.
像教會信奉基督一樣.
所以妻子對丈夫的信奉.
是什麼呢?.
出於你對丈夫的愛.
因為你愛他.
你願意聽他的話.
而不是.
但是以此同時.
不可以違背聖經的教導.
還有一樣東西很重要.
妻子要讓丈夫得榮耀.
什麼意思?.
千萬不要讓丈夫在人面前覺得尷尬.
覺得難堪.
如果你丈夫做了什麼.
你不喜歡,不高興.
或者他做錯了什麼事.
說錯了什麼話.
但是有人在身邊.
那怎麼辦?.
求神給你有智慧和力量.
回家才跟他算帳.
我師母懂得做這件事.
我希望你們懂得做.
讓我再說.
做妻子的.
有責任去建立你的丈夫.
而同樣的.
做丈夫的.
要怎麼樣呢?.
也應該刻意去建立你的妻子.
幫助你的妻子成為神所閱立.
基督所喜愛的門徒.
那麼先生又應該怎麼做呢?.
弟兄們一起讀讀今次.
弟兄們.
大聲一點.
你們做丈夫的.

$^{601}$要愛你們的妻子.
正如基督愛教會.
為教會捨己.
要用水接著倒.
把教會洗淨.
成為聖潔.
可以獻給自己.
作個榮耀的教會.
毫無瞻污,咒文等類的病.
乃是聖潔.
絕有瑕疵的.
剛才我說過.
妻子對丈夫的信服.
是基於愛.
這是第一原因.
因為愛丈夫要聽話.
除非你的妻子覺得被愛.
如果妻子不覺得被愛.
我們很難要求妻子對我們信服.
聽得明白嗎?.
不知道我們做丈夫的.
有沒有讓妻子覺得被愛.
或者讓妻子覺得.
他們是免費勞工.
替丈夫生孩子的工具.
甚至是褻瀆的工具.
所以.
我跟弟兄說.
還是我自己說.
有人說.
男人不善表達感情.
有沒有聽過?.
我們自己沒有說過.
男人似乎不善於表達我們的感情.
不過我越來越發覺.
這其實是我們做男人的藉口.
不懂可以學.
不可以說不懂.
可以學的.
在我成長的年代.

$^{641}$我小時候在哪裡學的.
我都不知道.
我們就知道.
做丈夫要照顧自己的妻子.
照顧自己的子女.
養妻活兒.
在我們做男人.
做丈夫.
做爸爸.
應該有責任.
我不知道.
爸爸也不喜歡跟我說話.
不過小時候.
我們那個年代.
那個文化文化.
傳統的中國男人就是這樣.
我以前也是很傳統的.
現在也很傳統.
信耶穌之後.
神就改變了.
為什麼呢?.
信耶穌之後.
聖經教導我們要愛護妻子和子女.
將最好的給他們.
要服侍他們.
兩樣加起來的時候.
在我學習做丈夫.
做爸爸的過程當中.
最容易做什麼呢?.
用心做工.
大家有沒有?.
用心做工.
賺多點錢.
嘗試送子女去最好的學校讀書.
太太喜歡買什麼.
買吧.
喜歡吃什麼.
吃吧.
喜歡去旅行.
去吧.

$^{681}$我們覺得.
做了我們要做的事.
I try my best.
我們會這樣跟自己說.
我做了我做丈夫.
應該做的事.
去 entertain 我的妻子.
大家覺得是不是?.
我們這樣做夠不夠呢?.
弟兄姊妹.
或者是弟兄們.
我跟弟兄說.
我們做丈夫.
如果想讓我們的妻子.
感受到丈夫的愛.
單單做我剛才做的那些事.
是不夠的.
不是叫你不做.
單單做那些事是不夠的.
因為重點.
不是為我們的妻子.
做了多少事.
而是我們能否了解.
體會妻子的需要.
讓她覺得.
跟我們可以連線.
可以相通.
問題是什麼呢?.
就是每個女性.
每個姐妹.
對愛的那種感覺.
和連線的方式.
都不一樣.
我的一般性.
一般的女性.
往往是透過言語.
和情感的交流.
來確認大家是否有連線.
大家是否相通.
如果我們明白這一點.

$^{721}$我們做弟兄做丈夫.
第一樣要學的是什麼呢?.
專心聽老婆說.
目前還在學.
因為我是坐不定的人.
所以師父經常抱怨.
我不專心聽他說話.
專心聆聽太太說話.
用言語去表達.
對妻子的欣賞和感激.
很多時候.
有時候老婆回家.
如果她去做工的話.
會和丈夫抱怨.
今天做工很辛苦.
那班人都不合作.
回到家裡.
太多瑣碎事了.
有時候會和丈夫抱怨.
無論是工作,家庭,生活.
不知道我們弟兄覺得.
妻子需要什麼呢?.
我們男人.
第一個反應.
1,2,3,4已經解決了.
我們是問題解決師.
給問題的答案.
告訴她如何解決問題和方法.
所以我們就失敗了.
其實老婆很厲害.
她跟你說這個問題.
她心中有數.
她知道怎麼做.
她不是要你拿答案.
她希望你明白.
How she feel.
她感覺什麼.
她知道怎麼解決.
但她想丈夫明白她的感受.
所以.

$^{761}$其實只要丈夫說.
是嗎?.
你現在在那裡很辛苦.
如果換作是我.
我也好像你這麼累.
我告訴你.
你這樣說.
她馬上發脾氣.
你不用再說什麼.
她覺得你很體貼.
她覺得你很愛她.
明白我說什麼嗎?.
要學的.
因為這件事不容易做.
除此之外.
可以主動幫太太.
做一些家裡的事.
用言語.
剛才說過.
去表達對妻子的欣賞.
白天要上班.
你又上班.
她又上班.
回到家裡.
還有很多家務.
又照顧小孩.
又吃東西.
說一句多謝的話.
不需要長氣.
是不是?.
時不時.
如果可以的話.
應該可以的.
帶些小禮物回家.
不要等生日.
週年紀念才做那些事.
有時不是送小禮物.
一起出去吃飯.
今晚不用煮了.
可以嗎?.

$^{801}$還有一件事.
很重要.
我剛才教過.
妻子要怎麼做.
是什麼?.
必須在朋友家人面前.
守護你妻子的尊嚴.
要不要我再說一次?.
在朋友家人面前.
守護你太太的尊嚴.
千萬不要在人家面前.
批評你老婆.
這樣最好.
千萬不要讓人覺得.
你在踩老婆.
有時男子是不經意的.
求神憐憫幫助我們.
諸多弟兄.
我做了牧師幾十年.
現在我教你們.
不過我自己也在學.
我們彼此鼓勵.
聖經說.
你們作丈夫的.
要愛你們的妻子.
正如基督愛教會.
為教會寫記.
我們要為妻子放下自我.
按聖經的教導.
讓妻子能夠全心全意.
毫無玷污咒民等類的病.
是聖潔無瑕疵的.
我們要做什麼?.
獻給我們自己.
你回去再看一次.
獻給我們自己.
做一個可以專一愛我們.
信服我們的妻子.
妻子對丈夫的信服.
是要我們贖回來的.

$^{841}$是要老公賺取回來的.
諸多弟兄姊妹.
我們應該怎樣去賺取.
妻子對我們的愛和信服呢?.
弟兄們.
再讀二十八節到三十節.
聖經是兩次重複說的.
姊妹說一次.
弟兄要留心聽多一次.
一齊讀二十八節.
丈夫也當照樣愛妻子.
如同愛自己的身子.
愛妻子便是愛自己了.
從來沒有人恨惡自己的身子.
總是保養姑息.
正像基督代教會一樣.
因我們是他身上的主體.
諸多弟兄姊妹.
聖經教導我們做夫妻.
是要互相建立.
互相建立.
幫助對方成為神所悅立.
耶穌所喜愛的門徒.
幸福快樂的婚姻.
能夠讓夫妻兩個人.
不單是白頭到老.
是一起成長.
當我們能夠彼此相待的時候.
我們就能夠見證基督的信仰.
帶給我們的祝福.
時間關係,我只剩一分鐘多一點.
如何見證?我沒時間說了.
不過我很希望和大家.
讀一段聖經.
傳道書第四章.
傳道書第四章九到十二節.
請大家留心聽我讀.
這段聖經很多時候在婚禮中會聽到.
當然不是說婚禮.
但無論如何,先讀一讀.

$^{881}$第九節說.
兩個人總比一個人好.
因為二人勞碌,同得美好的果效.
為什麼呢?.
因為有人和你同工.
效果會好一點.
成果會多一點.
讓你的喜樂加倍.
第十節.
若是跌倒,這人可以扶他的同伴.
若是孤身跌倒,沒有別人扶起他來.
這人就有禍了.
我們都會遇到危難.
人生當中怎麼可能沒有危難呢?.
但如果有人能夠扶我們一把的時候.
我們就能夠撐過去.
十一節.
在者,二人同睡,就都戀和.
一人獨睡,怎能戀和呢?.
大家放在一張床.
意思是我們能夠互相安慰.
互相鼓勵.
我們就能夠一起捱過那個嚴冬.
十二節.
有人功勝孤身一人.
若有二人便能抵擋他.
大家都聽過.
二人同心,其利斷金.
我們就能夠克敵制勝.
三股合成的誠旨,不容易截斷.
因為大家願意通力合作.
我們可以加強.
抵受外界衝擊的韌力.
能夠面對,克服一切的磨難.
諸多兄弟姐妹.
聖經對夫妻應該怎樣一起生活.
互相建立.
而不是互相批評.
互相指責.
有很清楚的教導.

$^{921}$我看大家的時候.
有些人結婚十年.
不止,二十年.
三十年,四十年.
甚至有五十年,六十年.
執子之手,與子偕老.
現在越來越多人可以實現.
現在人們越來越長命.
但不代表我們夫妻的關係.
能夠榮耀上帝.
因為很多夫妻.
還未學會彼此信服.
互相建立.
我們不能讓基督的信仰.
改善我們夫妻的關係.
這就是我們今天需要去面對.
和處理的問題.
如果你是單身.
我剛才說過.
我鼓勵你找個同性伴侶.
與他建立深厚的友誼.
然後你可以有機會.
時時互相提醒,鼓勵.
追求靈性上的成長.
今天結束的時候.
我鼓勵大家.
當你的配偶.
你的丈夫,你的妻子.
是你屬靈的夥伴.
找些機會一起祈禱,讀經.
或者參與教會的事.
一起做些事情.
然後一起事奉.
一起去讀經,祈禱.
可以彼此提醒,互相幫助.
今天我經常重複這兩個字.
目的是什麼?.
讓我們都能夠成為.
上帝所喜愛,耶穌所閱立的.
基督門徒.

$^{961}$這就成為我們建立一個.
屬靈家庭的基礎.
為建立門徒.
在家中奠定了一個.
很堅固的基石.
所以求神幫助我們弟兄姊妹.
你是丈夫,你是妻子.
回去後該怎麼做呢?.
求聖靈指引你們.
賜給你們.
我們一起祈禱.
\newpage



\section{撒母耳記上 25:1-44}
\label{sec:032118}
\textbf{作個好男人,作個好女人—大衛與亞比該 (下)}~~~~李思聰牧師
\newline
\newline
連結: \href{http://www.church.com.hk/acms/content.asp?site=cdc\&op=show\&type=product\&code=032118\&layout=sermon}{\texttt{032118}} ~~~~ 語音日期: 2026-05-10
\newline
\newline
\hyperref[sec:032117]{\small{< < < PREV SERMON < < <}}
~
\hyperref[sec:index_chronic]{\small{[返順時目]}}
~
\hyperref[sec:032119]{\small{> > > NEXT SERMON > > >}}
\newline
\newline
撒母耳記上 25:1-44
\newline
\begin{longtable}{cl}
\hline
\hline
章節 & 經文 (和合本修訂版)\\
\hline
25:1 & \begin{tabularx}{0.7\textwidth}{X} 撒母耳死了，以色列眾人聚集，為他哀哭，把他葬在拉瑪他的家裡。大衛動身，下到巴蘭的曠野。 \end{tabularx} \\ \\ \relax
25:2 & \begin{tabularx}{0.7\textwidth}{X} 在瑪雲有一個人，他的產業在迦密。這人是一個大富翁，有三千隻綿羊，一千隻山羊；他正在迦密剪羊毛。 \end{tabularx} \\ \\ \relax
25:3 & \begin{tabularx}{0.7\textwidth}{X} 這人名叫拿八，他的妻子名叫亞比該。拿八的妻子有美好的見識，又有美麗的容貌，但拿八為人剛愎兇惡，是迦勒族的人。 \end{tabularx} \\ \\ \relax
25:4 & \begin{tabularx}{0.7\textwidth}{X} 大衛在曠野聽見拿八正在剪羊毛， \end{tabularx} \\ \\ \relax
25:5 & \begin{tabularx}{0.7\textwidth}{X} 就派十個僕人，對他們說：「你們上迦密到拿八那裡，提我的名向他問安。 \end{tabularx} \\ \\ \relax
25:6 & \begin{tabularx}{0.7\textwidth}{X} 你們要如此說：『願你來年平安，願你家平安，願你一切所有的都平安。 \end{tabularx} \\ \\ \relax
25:7 & \begin{tabularx}{0.7\textwidth}{X} 現在我聽說你有剪羊毛的人，你的牧人和我們在一起，他們在迦密一切的日子，我們沒有欺負過他們，他們也未曾失去甚麼。 \end{tabularx} \\ \\ \relax
25:8 & \begin{tabularx}{0.7\textwidth}{X} 你問你的僕人，他們會告訴你。願我的僕人在你眼前得歡心，因為我們是在好日子來的。請你隨手分點食物給僕人和你兒子大衛。』」 \end{tabularx} \\ \\ \relax
25:9 & \begin{tabularx}{0.7\textwidth}{X} 大衛的僕人到了，就提大衛的名，把這一切話告訴拿八，他們就停頓下來。 \end{tabularx} \\ \\ \relax
25:10 & \begin{tabularx}{0.7\textwidth}{X} 拿八回答大衛的僕人說：「大衛是誰？耶西的兒子是誰？今日悖逆主人奔逃的僕人很多。 \end{tabularx} \\ \\ \relax
25:11 & \begin{tabularx}{0.7\textwidth}{X} 我豈可把飲食，以及我為剪羊毛的人所宰的肉給那些我不知道從哪裡來的人呢？」 \end{tabularx} \\ \\ \relax
25:12 & \begin{tabularx}{0.7\textwidth}{X} 大衛的僕人轉身從原路回去，照這一切的話告訴大衛。 \end{tabularx} \\ \\ \relax
25:13 & \begin{tabularx}{0.7\textwidth}{X} 大衛對他的人說：「你們各人都要佩上刀！」各人就都佩上刀，大衛也佩上刀。跟隨大衛上去的約有四百人，留下二百人看守物件。 \end{tabularx} \\ \\ \relax
25:14 & \begin{tabularx}{0.7\textwidth}{X} 拿八的一個僕人告訴拿八的妻子亞比該說：「看哪，大衛從曠野派使者來向我主人問安，主人卻辱罵他們。 \end{tabularx} \\ \\ \relax
25:15 & \begin{tabularx}{0.7\textwidth}{X} 但是那些人待我們真好；我們在田野與他們一切來往的日子，沒有受他們欺負，也未曾失去甚麼。 \end{tabularx} \\ \\ \relax
25:16 & \begin{tabularx}{0.7\textwidth}{X} 我們在他們那裡牧羊，一切的日子他們晝夜保護我們。 \end{tabularx} \\ \\ \relax
25:17 & \begin{tabularx}{0.7\textwidth}{X} 現在，你當知道，看怎樣做才好。不然，禍患必定臨到我主人和他全家。他性情兇暴，無人敢與他說話。」 \end{tabularx} \\ \\ \relax
25:18 & \begin{tabularx}{0.7\textwidth}{X} 亞比該急忙將二百個餅，兩皮袋酒，五隻宰好的羊，五細亞烘熟的穗子，一百個葡萄乾餅，二百個無花果餅，都馱在驢上， \end{tabularx} \\ \\ \relax
25:19 & \begin{tabularx}{0.7\textwidth}{X} 對僕人說：「你們在我前面走，看哪，我跟著你們去。」她卻沒有告訴丈夫拿八。 \end{tabularx} \\ \\ \relax
25:20 & \begin{tabularx}{0.7\textwidth}{X} 亞比該騎著驢，正下山坡，看哪，大衛和他的人正迎著亞比該下來，她就去迎接他們。 \end{tabularx} \\ \\ \relax
25:21 & \begin{tabularx}{0.7\textwidth}{X} 大衛曾說：「我在曠野為那人看守他一切所有的，以致他未失去任何一樣東西，實在是徒然了！他竟然向我以惡報善。 \end{tabularx} \\ \\ \relax
25:22 & \begin{tabularx}{0.7\textwidth}{X} 凡屬拿八的男丁，我若留一個到明日早晨，願神重重懲罰大衛！」 \end{tabularx} \\ \\ \relax
25:23 & \begin{tabularx}{0.7\textwidth}{X} 亞比該看見大衛，就急忙下驢，在大衛面前臉伏於地叩拜。 \end{tabularx} \\ \\ \relax
25:24 & \begin{tabularx}{0.7\textwidth}{X} 她俯伏在大衛的腳前，說：「我主啊，願這罪歸於我！求你容許使女向你進言，更求你聽使女的話。 \end{tabularx} \\ \\ \relax
25:25 & \begin{tabularx}{0.7\textwidth}{X} 我主不必理會拿八這性情兇暴的人，他就像他的名字一樣；他名叫拿八，為人也真是愚頑。至於我，你的使女並沒有看見我主所派來的僕人。 \end{tabularx} \\ \\ \relax
25:26 & \begin{tabularx}{0.7\textwidth}{X} 現在，我主啊，耶和華既然阻止你親手報仇，避免流人的血，我指著永生的耶和華起誓，又指著你的性命起誓：『現在，願你的仇敵和謀害我主的人都像拿八一樣。』 \end{tabularx} \\ \\ \relax
25:27 & \begin{tabularx}{0.7\textwidth}{X} 現在求我主把婢女送來的禮物給跟隨我主的僕人。 \end{tabularx} \\ \\ \relax
25:28 & \begin{tabularx}{0.7\textwidth}{X} 求你原諒使女的冒犯。耶和華必為我主建立堅固的家，因為我主為耶和華爭戰，並且你一生的日子查不出有甚麼惡來。 \end{tabularx} \\ \\ \relax
25:29 & \begin{tabularx}{0.7\textwidth}{X} 雖有人起來追逼你，要尋索你的性命，我主的性命在耶和華－你的神那裡，如同藏在生命的寶藏中。至於你仇敵的性命，耶和華必甩去，如用機弦甩石一樣。 \end{tabularx} \\ \\ \relax
25:30 & \begin{tabularx}{0.7\textwidth}{X} 耶和華照所應許你的福氣賜給我主，立你作以色列王的時候， \end{tabularx} \\ \\ \relax
25:31 & \begin{tabularx}{0.7\textwidth}{X} 我主就不至於因為親手報仇，流了無辜人的血，而心裡不安，良心有虧了。耶和華賜福給我主的時候，求你記得你的使女。」 \end{tabularx} \\ \\ \relax
25:32 & \begin{tabularx}{0.7\textwidth}{X} 大衛對亞比該說：「耶和華－以色列的神是應當稱頌的，因為他今日派你來迎接我。 \end{tabularx} \\ \\ \relax
25:33 & \begin{tabularx}{0.7\textwidth}{X} 你和你的見識也配得稱讚，因為你今日攔阻我親手報仇、流人的血。 \end{tabularx} \\ \\ \relax
25:34 & \begin{tabularx}{0.7\textwidth}{X} 我指著阻止我加害於你的耶和華－以色列永生的神起誓，若不是你很快地來迎接我，到早晨天亮的時候，凡屬拿八的男丁，必定一個也不留。」 \end{tabularx} \\ \\ \relax
25:35 & \begin{tabularx}{0.7\textwidth}{X} 大衛從亞比該手中收了她帶來的禮物，對她說：「平平安安上你的家去吧！你看，我看了你的情面，聽了你的話。」 \end{tabularx} \\ \\ \relax
25:36 & \begin{tabularx}{0.7\textwidth}{X} 亞比該到拿八那裡，看哪，他在家裡擺設宴席，如同王的宴席。拿八心情舒暢，酩酊大醉。所以亞比該大小事都沒有告訴他，直等到早晨天亮的時候。 \end{tabularx} \\ \\ \relax
25:37 & \begin{tabularx}{0.7\textwidth}{X} 到了早晨，拿八酒醒了，他的妻子把這些事都告訴他，他就發心臟病快死了，僵如石頭。 \end{tabularx} \\ \\ \relax
25:38 & \begin{tabularx}{0.7\textwidth}{X} 過了十天，耶和華擊打拿八，他就死了。 \end{tabularx} \\ \\ \relax
25:39 & \begin{tabularx}{0.7\textwidth}{X} 大衛聽見拿八死了，就說：「耶和華是應當稱頌的，因為我從拿八手中受了羞辱，他為我伸冤，又阻止他的僕人行惡；耶和華使拿八的惡歸到他自己頭上。」於是大衛派人去向亞比該說，要娶她為妻。 \end{tabularx} \\ \\ \relax
25:40 & \begin{tabularx}{0.7\textwidth}{X} 大衛的僕人來到迦密，到亞比該那裡，對她說：「大衛派我們到你這裡，要娶你作他的妻子。」 \end{tabularx} \\ \\ \relax
25:41 & \begin{tabularx}{0.7\textwidth}{X} 亞比該起來叩拜，俯伏在地，說：「看哪，你的使女情願作婢女，為我主的僕人洗腳。」 \end{tabularx} \\ \\ \relax
25:42 & \begin{tabularx}{0.7\textwidth}{X} 亞比該立刻起身，騎上驢，五個女僕跟著她走。她跟從大衛的使者去，就作了大衛的妻子。 \end{tabularx} \\ \\ \relax
25:43 & \begin{tabularx}{0.7\textwidth}{X} 大衛先前娶了耶斯列人亞希暖，她們二人都作了他的妻子。 \end{tabularx} \\ \\ \relax
25:44 & \begin{tabularx}{0.7\textwidth}{X} 掃羅已把他的女兒米甲，就是大衛的妻子，給了迦琳人拉億的兒子帕提為妻。 \end{tabularx} \\ \\
[1ex]
\hline
\hline
\end{longtable}
$^{1}$她說作過好男人.
今天我們繼續下去下講.
作過什麼?.
好女人.
大衛和阿彼哥的故事.
這個就是我們接續上一星期.
我們看的經文.
就是《三國義記》上的第25章.
因為今天母親節.
我都講了一點點.
大家知不知道母親節的由來?.
就是由這位Anna Jarvis發起的.
因為她的母親在120年前.
120年前的昨天.
或者今天昨天.
就是5月9日.
就是1906年.
剛剛120年.
之前她母親去世.
她母親是一位敬虔的基督徒.
是主一學的老師.
好像剛才的新聞仔.
他們的老師.
所以當媽媽過去之後.
這位Anna非常懷念母親.
不單止對其他的兒童.
也特別對她的撫養.
和在信仰上的幫助.
所以其實母親節和基督教有關.
所以我們看到這位基督徒.
就開始倡導.
紀念普世母親們的偉大.
直至表揚她的媽媽.
和上帝的愛.
直至到1913年 隔了七年之後.
在5月10日 就是今天.
正式簽署了由兩個上下院.
美國的Congress.
由威爾遜總統.
President Wilson.

$^{41}$我也可以偷一點光.
由President Wilson簽署.
成為美國的一個紀念日.
每逢五月第二個星期日.
在Anna離世之前.
他已經推廣到.
有43個國家設立了母親節.
到今天為止.
我相信差不多全世界.
大部分的國家.
都會以五月第二個星期日為母親節.
元氣其實都是由一位基督徒的母親.
她的見證引起的.
在母親節的時候.
讓我們可以紀念我們的母親們.
Focus on Family.
曾經做過一個調查.
調查三千個小孩子.
由6歲至14歲.
所以比較小的那些.
剛剛轉做teenager.
大一點的那些可能不同看法.
不過小的那些都很衷心表達出來.
這三千個從6歲至14歲的小孩子.
他們有76個.
當人問到他媽媽對他們多重要.
有76個表示給他們很好的意見.
所以有大約四分之三的小孩子.
都覺得媽媽給他們好的意見.
對他們有很大影響.
又差不多一半.
當然有些不懂得表達.
但是他表達出來.
如果問媽媽具備有什麼條件.
應該具備什麼條件.
他們認為希望媽媽.
有四分之三的人認為他們是支持他們.
這個給我們的母親們一些鼓勵.
也是了解.
然後也差不多的人數.

$^{81}$他們希望媽媽對他們的了解.
所以支持和了解.
有大約三分之一的認為.
希望他們能夠去鼓勵他們.
去發展他們的長處.
他們的潛能.
Inspire他們.
Inspiring.
所以這些都是一個調查.
給我們反映.
小孩子對母親的一些感受.
一些感恩.
也有一些期望.
所以我們今天來看.
做個好男人做個好女人.
當然我特別是說到.
好的母親.
好的婦女.
好的女人.
在《三妹記二十五章》這段經文裡面.
有學者研究.
Berlin他認為這個敘事文體.
這個敘事文體的分析.
就是去表達出.
有一個主角出來.
一直在說大衛的.
當然是大衛的故事.
但是表達了有一個很重要的角色.
就是阿比蓋.
阿比蓋他是一個.
可以說是一個敘事文體裡面.
一個主角.
他是一個type.
如果敘事文體是一個type.
type就是一個預表角色.
他要帶出一個故事.
和後面一個神學意義出來.
所以以敘事文體來分析.
這個type這個角色.
Berlin認為是可以表達出.

$^{121}$一個perfect wife.
或者一個完美的妻子的形象.
所以透過阿比蓋的故事.
其實將一些美好的婦女的德行.
或者女子應該有的.
在神眼中認為寶貴的東西.
向我們展示.
向我們啟示出來.
所以今天我們透過.
《三文魚記》第25章.
我們的思想.
集中在好女人那裡.
什麼是好女人呢?.
我們可以看到從經文.
第一.
讓我們看到的就是.
有什麼?.
有美貌.
為什麼呢?.
因為《三文魚記》上第25章的第三節.
開始形容這個故事.
敘述整件事的時候.
介紹阿比蓋.
怎麼介紹她呢?.
我們一起讀一讀這一節.
好嗎?第三節請讀.
那人名叫拿伯.
是加勒族的人.
他的妻名叫阿比蓋.
是聰明盡美的婦人.
拿伯為人剛愎兇惡.
所以當他來展示.
介紹一個故事.
一個配角的時候.
就是拿伯.
上一次我們將拿伯和大衛比較.
一個是壞男人.
一個是好男人.
相比較的時候.
就突出了大衛的優勝之處.

$^{161}$也是我們學校的地方.
但是拿伯這個壞男人.
他是屬於加勒族.
是以色列人.
我們知道他們有一個小的部落.
是加勒的後人.
就是入以色列的加勒他們的後人.
就是和約書亞時期.
一直到後來延續下去.
他的妻就叫做阿比蓋.
這是第一個介紹阿比蓋.
第一個形容就是說.
她是聰明盡美的婦人.
相比拿伯一樣.
一個這樣的形容.
就和他剛剛成為一個強烈的對比.
就是剛愎兇惡.
因為剛愎兇惡表示他的心性.
聰明盡美.
就有心性.
也有外表.
因為說到盡美.
而事實上在原文.
這個盡美這個字.
也有善解人意.
或者是有智能的意思.
有智慧有智能.
所以其實也比較偏向他的內心.
所以我們看到拿伯.
這個拿伯就是一個天壤之別.
相反如果我們查考.
《撒慕爾傑》上十六章.
形容大衛的時候.
當大衛出場.
聖經形容他什麼呢.
聖經形容大衛.
很相似.
他說他是容貌盡美.
所以我們發覺.
原來這個容貌盡美.

$^{201}$和這個聰明盡美.
原文是同一個字.
所以奇妙的地方就在這裡.
原來亞比蓋稱他為聰明盡美.
其實和拿伯不相配.
和誰相配.
和大衛相配.
因為大衛形容他的時候.
就是容貌盡美.
所以我們看到.
這個實在是一個很奇妙的安排.
結果後來他就成為了大衛的妻子.
在一本書.
由諾奇瓦所寫的.
《王與侯》一本書裡面.
他提到亞比蓋.
亞比蓋後來成為了大衛的妻子.
成為了皇后.
所以他的形容.
形容這個皇后是受愛為宗.
受愛為宗.
就是指他外面很盡美.
裡面很有智慧.
和聖經的形容很相似.
所以強調是他內在的一些真美.
而實際上在書裡也列出其他的皇后.
例如誰呢.
例如剛剛我們在.
春令福音退修海上遊.
我們的主題.
就是講誰呢.
伊斯帖他有一個皇后.
所以伊斯帖記形容.
這個皇后也是非常美貌盡美.
所以在當中他表達出.
去研究的時候.
這兩位皇后非常有相似的地方.
相似也是因為伊斯帖.
我們每個人去過春令會就知道.
伊斯帖講到他的優點的時候.

$^{241}$他不只是外表.
更加著重是他的美德.
就是他不貪心.
他不忘恩.
他願犧牲.
他夠勇敢.
也有智慧.
所以這些都是伊斯帖的特長.
也都是聖經裡面.
特別是這本書.
《皇與后》裡面所表達出來的.
當然我們講到姐妹們.
女性們.
都是愛美的.
是很自然的事.
而事實上美麗是上帝所賜的.
我們知道神是一位.
不單只是美善的神.
也是一位創造美好.
創造美麗的上帝.
我們看見周圍大戰的環境就知道.
上帝是愛美的.
看所有所做的一切.
祂都畫得很甚好.
今天我們看到各式各樣的花朵.
各式各樣的動物.
在天地間.
不同的顏色.
表揚出來.
在整個宇宙界.
所給我們看見的神設.
的月亮星宿.
都看到神是一位.
非常美麗的上帝.
祂愛美的.
這種美是神所賜的.
所以姐妹們愛美.
是很自然的.
扮扮美麗是應該的.
我提過在上個月.

$^{281}$我們喜樂團契.
李小妹帶了一個聚會.
叫做男女之別.
我們分了男和女.
當男士去說女士有什麼特點.
包括優點或者缺點.
優點裡面.
其中一樣東西.
表達是什麼呢.
大家都同意的.
就是她們是有美感.
比起男士有美感.
你們同意嗎.
同意.
所以我們很多美藝組.
全部都是姐妹們.
沒有弟兄們.
因為他們的美感比較好.
我們男士就隨便.
什麼顏色亂碰.
有時候碰到不清不楚.
但是姐妹們就有美感多一點.
事實上美是姐妹特別可以有的.
當然我們不希望.
我們的姐妹們不修邊幅.
個個都端莊美好.
也不希望我們的弟兄們.
衣冠不整.
所以我們都是.
其實這件事都是.
應用在無論男和女身上.
不修邊幅不好.
衣冠不整都不好.
因為美麗齊整.
端莊是上帝所喜愛的.
不過我們又發覺.
在聖經裡面講到美麗的時候.
很多時候.
聖經的著重點.
都不只是外表的美貌.

$^{321}$所以有美貌是好的.
但是更加.
原來這種美.
是真的好像我們剛才所說的.
那本書裡面所著重的.
一種內在的真的美.
在聖經裡面.
時常都這樣提到.
就是好像.
在《箴言》31章30節.
我們一起讀讀這篇經文.
請讀.
艷麗是虛假的.
美容是虛浮的.
為敬畏搖滑的婦女.
必得稱讚.
這裡我們見到.
她著重的是內在.
這個內在就是指到.
她有敬畏搖滑的心.
而事實上.
艷麗也好.
或者美容也好.
都是會很快過去的.
亦都是在表面很虛浮的.
當我們去看見一個女士.
一個好女人.
其實很多時候.
我們看見她美麗.
是她內心表達出來.
那種美麗.
多於外表的裝飾.
今天其實外表裝飾很多.
因為很多選美也好.
電視明星也好.
逛街.
見到各式各樣的.
在外面吸收的美麗.
都可以做到.
但這些外表的美麗.

$^{361}$卻是虛浮的.
虛假的.
因為當她卸妝.
就可能變成另一個人.
但我們見到.
真正的美麗.
就是在乎這些內在.
所謂敬畏搖滑的婦女.
是必得稱讚.
這個指裡面的心.
在新約.
無論在保羅的書信.
或者彼得的書信.
都有講.
我們又讀讀這節好嗎.
是提摩大前書.
是保羅的書信.
請讀.
我又願女人廉恥.
自首以正派衣裳為裝飾.
不以編發黃金珍珠.
和貴價的衣裳為裝飾.
只要有善行.
且才與自稱是.
敬神的女人相宜.
所以你見到這裡.
是指外表和內在.
外表當然是.
指編發黃金等等.
都不是沒有問題的.
最重要是有正派的衣裳.
但內在更加重要.
只有善行.
這才是與自稱為.
敬神的女人相宜.
所以這裡我們見到.
又是有外表.
又有裡面.
外表的美固然好.
但內在的美更加重要.

$^{401}$在彼得也是同樣相似的教導.
彼得的書信.
彼得前書三章.
三至四節我們再讀.
請讀.
你們不要以外面的扁頭髮.
戴金色穿美衣為裝飾.
只要以裡面長存久溫柔.
安靜的心為裝飾.
這在神面前是極寶貴的.
所以親愛的姐妹們.
我們母親們.
我們看見聖經教導我們.
外面固然有好的裝飾.
但不要以這些為.
真正的美麗評估.
真正的美麗評估是什麼.
真正的美麗評估是裡面.
裡面究竟我們有沒有這種溫柔.
這種長存溫柔.
安靜的心為裝飾.
這個在神的面前.
是極寶貴的.
所以這個不是在人的面前.
更重要的是在神的面前.
是極寶貴的.
姐妹們.
巴不得我們都這樣追求.
姐妹應該都是有這種溫柔.
有這種長久安靜.
一個平靜.
安靜的心.
為我們真正的裝飾.
所以我們看見.
我們很多美麗的姐妹們.
美麗的母親們.
我們看見她們不是以外表.
一些艷麗的打扮.
而是內在敬愛神.
有善行.

$^{441}$和這些長久溫柔.
安靜的心為裝飾.
這個實在是更好的.
阿彼戈在寶貴地方.
就是他善解人意.
有智慧.
當大衛稱讚他的時候.
他說他聰明.
而且有見識.
所以這個正是他.
吸引人的地方.
正是他美麗的地方.
在多年前有一篇導向雜誌.
他將一些女士.
外在一些化妝的材料.
賦予屬靈的意義.
所以我們可以看看.
這個有什麼屬靈意義.
粉底霜.
對不起.
我又不太清楚是什麼.
粉底霜.
粉膏是什麼.
可能幾十年前的化妝品.
粉餅.
我想是有人用的.
胭脂.
有沒有人用胭脂.
也有的.
口紅.
一定有的.
眼蓋膏.
所以.
這些都是那些女士的化妝品.
不過他賦予寶貴的屬靈意義.
他說粉底霜好像看不到的.
Foundation.
但是他的善行.
是可以表達出來的.
粉膏.

$^{481}$很多時候我們瘡黃的皮膚.
就代表屬靈的貧乏.
不過以一些敬畏神的心.
和精緻的品行.
就成為我們的粉膏.
粉餅就是以端正的品行.
和正派的衣裳.
胭脂就是以溫柔的心.
雖然看不到.
不過可以感受到.
所以這個很好.
因為溫柔的心.
因為胭脂不是很明顯的.
不過感受到.
口紅更加.
你的嘴巴怎麼了.
他說隨時說造就人的說話.
在《箴言三十一章》也說.
他開口就發自慰.
他舌上有仁慈的法則.
不要罵人那麼多.
有仁慈的法則.
然後眼蓋高.
眼蓋高是什麼.
他說能夠看到別人好的地方.
所以他的心裡面很美.
看的東西是什麼.
看的東西都是美善的.
所以這個希望.
不知道你有哪幾樣正在調查.
任何姐妹.
口紅也有.
眼蓋高有沒有.
粉底有沒有.
希望我們不單止有這些化妝品.
更加有右臉這些.
就是我們有的.
這些真正的.
去表達出來的.
一些內在而發出來的美好.

$^{521}$有一本書叫做《傳言美麗》.
《傳言美麗》其實是用了《箴言三十一章》.
我們很熟悉的第十至三十一節.
知不知道《箴言三十一章》是誰寫的.
就是尼姆伊勒王的媽媽寫的.
所以在聖經裡面特別說出.
這個是從母親的角度而看.
什麼叫做才德的婦人.
在這個《箴言三十一章》.
第十至三十一節.
這位母親去表達出來的時候.
才德的婦人.
今天我們沒有時間說了.
我們看到有很多方面.
包括了什麼.
包括了她有好的品格.
有正確的人生的價值觀.
有精操.
有對家庭的貢獻.
有使命感.
她能夠供應別人.
她自律.
而且有遠見.
就是有異象的人生.
她勤勞的工作.
勤奮的去生活.
而且對於自己有自信.
這些都是《箴言三十一章》裡面.
特別透過這本書.
《傳言美麗》去解釋出來.
也是她有憐憫人的心.
她有預備的.
不是沒有計劃的.
就很有預備的.
然後她很會做手工.
我們很多姐妹們都很厲害.
成為丈夫得益的好處.
丈夫的幫助.
她懂得處理事業.
就好像今天我們很多的職業婦女一樣.

$^{561}$工作很有成就.
她懂得做衣服.
她的華語很能表達出智慧和幫助人.
她能夠管理好自己的僕人,侍女.
也很好地照顧她的家庭.
她是被人稱讚的.
稱讚她為一位才德的婦人.
她是對神有信心的.
因為她是敬畏上帝的.
而且她是得到回報的.
因為她所勞碌.
所擺上她的信心.
她的服務.
她的貢獻.
是帶來果後.
在城門人,城牆.
當時以色列人在城門口.
那些人都稱讚她.
說她是幼福的.
不得我們都能夠開始.
《真言三十一章》.
才德的婦人裡面.
不斷地成為我們這些好女人的學習.
事實上姐妹們.
看見這些真的不簡單.
真的要很大忍耐力.
才能夠兼顧這麼多東西.
所以在我們上一次的男女有別裡面.
其中一個特點也是.
每個弟兄都贊成的.
就是姐妹是很有忍耐力.
比起我們男士.
忍耐力更大.
只是生孩子都知道.
所以真的我們發覺.
我們的姐妹們.
能夠繼續忍耐.
繼續堅持.
繼續去追求.
能夠表達出這些好女人的特色.

$^{601}$第一有美貌.
真的.
這個美貌不單是外在美.
更加是什麼?是內在美.
第二就是有美德.
我們發覺在《真言十一章》第二十二節這樣說.
他說.
一起讀好不好?.
他說.
「婦女有美貌而無見色.
如同金環戴在珠臂上」.
真的有大諷刺.
原來我們發覺.
有美貌是源於有美德.
因為我們發覺.
婦女有美貌.
但是沒有見色.
這個見色就是指她的內涵.
即是裡面.
這個內涵.
如果沒有內涵的.
只有樣貌外表.
如同金環很漂亮.
不過不配合美好的身份.
因為戴在珠臂上.
有一次有一個太太.
走來丈夫問她.
她說.
「你覺得我美不美?」.
丈夫跟她說.
「太太,人沒有十全十美.
不過我都覺得你有十全八美了」.
十全八美.
太太聽到都很開心.
你們開心嗎?.
如果80分都不錯.
即是80分都算高分.
太太開心完之後.
忽然間想一想.
還有那兩個是什麼呢?.

$^{641}$十全八美.
還有那兩個美是什麼呢?.
丈夫跟她說.
「一個就是外在美.
一個就是內在美」.
弟兄們千萬不要回去.
跟你太太開這個玩笑.
會死得很慘的.
不過事實上是一個很好的提醒.
究竟我們有沒有外在美.
更加有沒有內在美.
有美貌有美德.
所以我們看到.
這個就讓我們看見.
亞比該他兼備這種外在美.
更加他靚在他的美德.
有什麼美德.
我來思想聖經裡面的記載.
第一我給它叫做持家有道.
持家有道.
因為為什麼呢?.
如果大家還記得上個星期講的故事.
當拿伯來侮辱大衛.
大衛有什麼反應?.
有什麼反應?.
他就帶著多少人?.
400個隨從.
準備去剷除拿伯.
殺光他全家.
幸好有一個好女人.
他的妻子就是亞比該.
亞比該聽到那些僕人向他的提示.
說大能臨頭.
你看見亞比該臨危不亂.
這個就是他的地方.
他臨危不亂.
馬上叫人預備什麼?.
25章的第18節.
亞比該急忙將200餅.
2皮袋酒.

$^{681}$5隻收拾好料的羊.
5細瓦控好料的穗子.
100葡萄餅.
200無花果餅.
都帶在路上.
帶去哪裡?.
帶去在路上截住大衛.
希望可以平息他的怒氣.
希望可以除去他的殺機.
所以這個就是亞比該的急志.
在這裡我們看見什麼呢?.
形容他急忙.
急忙不是指他亂七八糟.
急忙是指他很有效率.
很有效率.
馬上處理.
請問這裡有多少樣東西?.
又有餅.
又有酒.
又有羊.
又有控好的穗子.
知道在哪裡.
有些還沒控好.
他知道控好在哪裡.
100個餅.
無花果全部.
這麼多東西.
在短短的時候.
馬上整齊好.
馬上收拾好.
當然他有人幫他.
但是他指示得很清楚.
馬上知道家裡有什麼.
知道東西放在哪裡.
知道要準備多少.
所以你會發覺.
他在很短的時間.
能夠有迅速的行動.
證明他平時非常持家有道.
非常明白到.

$^{721}$誰負責什麼.
哪裡在哪裡拿.
他有多少東西.
然後預備好這些東西.
在這麼短的時間裡面.
馬上出發.
我們看到這個實在.
看到他持家有道的一種表現.
親力親為.
不單止他預備好.
更加在這裡.
第25節.
這裡用了幾個形容詞.
在這個敘事文體.
特別用這些來形容他.
第23節.
他被該見大衛便急忙下路.
在大衛面前.
俯伏於地求拜.
我們看到他又是迅速的行動.
證明他這個人是很敏捷的.
是很有計劃的.
而且能夠在短時間之內.
做好這些一切.
第34節.
大衛稱讚他.
大衛稱讚他什麼.
我們一起讀一讀.
請讀.
我指著難阻止我加害於你的.
耶和華以色列永生的神起世.
你若不速速地來迎接我.
到明日早晨.
凡屬那伯的南丁必定不留一個.
大衛稱讚他.
你來得快就好了.
一過了今晚就太遲了.
所以你看到.
阿彼該真的把握住時間.
把握住這個機會.

$^{761}$而且他有這樣的能力.
能夠做到.
證明他是一個很能幹的婦人.
一個持家有道.
臨危不亂.
卻能夠冷靜的敏捷去應變一切.
這個真是一個非常值得.
我們看到的優點.
今天我們的婦女們.
我們的母親們.
都是持家有道.
很多家庭所需的事務.
很多事情是井井有條.
安排有理.
所以很多我們這些能幹.
又要上班.
又要顧家.
的姐妹.
實在是值得我們的欣賞.
她們精打細算.
勤奮勤勞.
持家有道.
這個實在是一種美德.
使丈夫有益無損.
所以我們很開心.
很感恩.
有我們很多很好的太太們.
很好的母親.
實在對於我們的家庭.
非常重要的一個貢獻.
我們上一次的聚會.
也說到母親或者婦女.
有什麼優點呢.
有兩件事我們也提出.
第一就是他們愛護兒女.
他們愛護兒女.
比我們父親愛護子女更多.
更加細微.
這個我深深感受到.
也是他們為家庭的犧牲.

$^{801}$為家庭的犧牲也很大.
真的很多時候.
眠乾睡濕.
丈夫就大被蓋頭.
但是他們很多時候.
眠乾睡濕.
你要照顧兒女.
我們的母親.
對於我們的兒女.
也對於家庭的貢獻.
對於我的家庭.
和我的女兒.
很大的貢獻.
持家有道.
這個就是第一個美德.
第二我們看到.
情理兼備.
就是《十美記》上25章27節.
我們一看.
讀一讀這節經文.
請讀.
如今求你將婢女.
送來的禮物.
給跟隨你的毒人.
所以我們剛才看到.
他預備了很多這些禮物.
來送給大衛.
和所有跟隨他的人.
這就顯出.
阿彼戈的智慧.
他不是空手去.
他帶著禮物去.
其實很明顯.
因為大衛來求.
就是希望求這些幫助.
因為他們現在這麼多人.
接著幫了拿百你們這麼多.
現在你們分羊毛了.
給我們一點我們需要的.
但是這個拿百拒絕了他.

$^{841}$還侮辱他三次.
記得大衛祝福了他一次.
祝福了他三次.
他侮辱他三次.
現在阿彼戈送來禮物.
這個禮物有原文的學者分析.
禮物這個字.
其實就是代表祝福的意思.
所以他就將祝福還給大衛.
大衛祝福我們.
他就將禮物成為祝福.
去祝福大衛和他所跟隨的人.
所以這種禮物表達善意.
表達善意.
因為本來拿百這麼兇.
現在他表達善意.
這個禮物是表達敬意.
因為代表尊敬他們.
多謝你們來保護我們.
這個禮物要代表歉意.
因為拿百這麼壞.
現在我代表我的丈夫.
表達我的歉意.
希望你原諒.
所以這個禮物有三重意義.
善意,敬意,歉意.
讓我們看見充滿著智慧.
而且我們看見他一見到大衛.
馬上做什麼行動.
立刻下來.
然後叩拜在他面前.
所以這是一種真誠謙卑的態度.
甘望下路.
俯伏於地.
而且替她的丈夫擔當他的罪名.
開口就認錯了.
就是說我丈夫實在不對.
實在太不對.
所以希望你可以原諒我們.
所以這種真是他所表達一種.

$^{881}$真誠謙卑的心.
他用了九次說「我主」這個字.
我主.
多少次?.
九次.
代表他很謙卑.
不斷稱呼大衛「我主」.
而且稱呼他自己做「婢女」.
稱呼自己「婢女」有六次.
婢女.
多少次?.
六次.
這實在表明他那種真誠謙卑的心.
所以我稱他為動之以情.
就是用感情,用情感.
輸掉大衛的怒氣.
所以這種動之以情.
實在是很有智慧.
不單止動之以情.
我說他情理兼備.
他素知於理.
因為他和大衛分析.
他說大衛你想想.
如果你殺了那八他該死.
殺了他那麼多人.
流了那麼多無辜人的血.
將來你作為皇帝就有個瑕疵.
你也不想有這個瑕疵.
瑤華會祝福你.
將來你必定作為皇帝.
他給予一個預言和應許.
所以在這個陳述裡面.
他不單止整個態度是動之以情.
更加素知於理.
他的分析層次分明.
首美呼應.
大家回去讀一下.
勸諫有力.
所以我們看到這個就是.
非常合情合理的一個勸誡.

$^{921}$所以看到這種充滿著美德.
是因為他情理兼備.
而事實上子妹們說話比較聰明.
因為她說話多.
比弟兄多得多.
所以就練到口齒.
口齒伶俐.
不過口齒伶俐就不要說閒話太多.
所以我們要小心一點.
姐妹們利用我們的口多說造就人的說話.
口出恩言.
這個就像主耶穌.
所以其實這種性格謙卑誠懇.
做事合情合理.
正是主耶穌基督給我們追求的美德.
所以不只是姐妹.
弟兄也是.
弟兄能夠學會.
如何控制我們的情緒.
管理好情緒.
學會合情合理去做事.
當然弟兄比較合理.
姐妹有時情感操縱多些.
不過大家都要學的.
如何能夠合情又合理.
這樣就減少衝突.
增加溝通.
這個實在是基督徒追求的美德.
其實聖靈幫助我們就是有這些.
能夠有這些結出來美好的果子.
另一本書叫做《屬神的女人》.
《屬神的女人》裡面講到.
姐妹們其實有很多方面可以追求.
包括知識,身體,靈性,事工,社交.
家政,照顧家庭,休閒,財政.
如何學得浪漫一些.
有很多方面.
不過綜合來說.
可以有幾方面說出來.
就是身,心,智,靈,社交.

$^{961}$這五樣東西.
其實五樣東西都能夠好好平衡.
身,心,智,靈,社交.
就是說身體懂得照顧自己.
心境懂得如何培養美好的情緒.
好的心境.
智能夠有智慧去處理事情.
靈和神的關係.
社交和人的關係.
所以我們見到這個正是耶穌基督.
他長大的時候.
聖經形容他.
他說耶穌基督長大的時候.
神人在我的心.
他身體,心智,神人在我的心.
都一同增長.
求主幫助我們.
我們不單止姐妹.
弟兄也是.
我們能夠繼續在這些成長.
這本書特別說到.
屬神的女人有的一些特性.
第三我們看到.
這個有美德第三方面就是丈夫的冠冕.
在《箴言》這樣說.
《箴言》十二章第四節.
一起讀好嗎.
請讀.
「才德的婦人是丈夫的冠冕」.
哎呀.
亞比蓋嫁給了那個叫什麼.
我們知道當時是沒有選擇的.
因為他不是自由戀愛.
很多時候是父母來訂婚姻.
或者男的他有錢.
他就去娶很多個老婆.
所以我們看到.
拿伯娶了亞比蓋.
相不相配.
哎呀.

$^{1001}$真的不相配.
我們說好像什麼.
鮮花插上牛糞上.
非常不配.
但是亞比蓋沒有選擇.
慘了馬上就被這個壞老公.
那怎麼辦.
機會來了.
大衛說現在我要來殺了他.
好了好了.
是不是.
不是.
就算這個丈夫這麼不好.
亞比蓋仍然去做他太太應該做的事.
這個就是寶貴的地方.
這個就是有美德.
因為她對於丈夫仍然去保護他.
所以其實他殺了殺男丁.
他自己沒事的.
所以女的沒事的.
反正殺了殺了他就對了.
不是.
他大衛希望止息了這個大能臨頭.
在他的丈夫身上.
而且留意回去讀.
他將他的罪歸在他自己身上.
所以他自己為丈夫擔罪.
他說他真的不對.
不過請你饒恕我們.
得罪你的地方.
所以我們見到這個他既合理.
但是也公義.
因為他沒有說拿白很對.
是不是.
就是不對就不對.
幫理不幫親.
真的.
不過他不是指責他.
這個丈夫真壞.
而是為他擔罪.

$^{1041}$站在他一邊.
來承擔這個錯失.
固然他不偏袒.
但是他一同擔當.
這個就是好太太這種美德.
聖經裡面說.
財德的婦人誰能得罪.
她的價值遠勝過真主.
真是我們發覺.
如果我們太太.
都能夠成為丈夫們的榮耀.
多多保護丈夫.
記得不要輕視他.
不要奚落他.
這個是作為一個.
有美德好女人的一個很重要的.
以丈夫成為丈夫的冠冕.
丈夫覺得有這個太太.
真是一個冠冕.
事實上我們知道.
子女尊不尊敬爸爸.
很多時候是看太太.
尊不尊敬丈夫.
所以聖經裡面特別說.
你們要敬重你們的丈夫.
在伊弗所書第五章.
不過又不是單方面.
在彼得書信又說.
丈夫你要敬重你的妻子.
所以我們看到.
原來這種互相尊重.
互相敬重.
是神的心意.
因為其實聖經裡面也說.
要對眾人尊敬.
即是表示基督徒的美德.
所以當我們對人尊敬.
對眾人表示尊敬.
更加何況在我們家庭.
好的太太.

$^{1081}$不要奚落你的丈夫.
不要藐視你的丈夫.
好的丈夫.
不要藐視.
不要看低你的太太.
相反進一步是應該敬重他們.
所以求神幫助我們.
希望我們能夠在這方面.
能夠更加追求有好美德.
我們第一.
阿彼該有好的美貌.
因為他源於有什麼話.
有好的美德.
第三我們看到.
他有好的靈性.
原來好的美德源於什麼.
源於好的靈性.
所以在這裡我們看到.
貼在三國義基上.
二十五章.
三十至三十一節.
我們一起讀一讀這兩節經文.
請讀.
我主現在若不親手報仇.
流無辜人的血.
到了耶和華.
照所應許你的話.
賜福與你.
納你作以色利的王.
那時我主必不至心裡不安.
覺得良心有愧.
耶和華賜福與我主的時候.
求你紀念疲累.
所以我們看到.
這個就是代表阿彼該.
他流露出對神的信心.
和敬畏的心.
留意這裡提到兩次.
到了耶和華.
照所應許你的話.

$^{1121}$賜福與你.
他知道福氣的源頭在於神.
納誰作王在於上帝.
所以他也說.
耶和華賜福與我主.
即是大衛的時候.
求你紀念我這個疲累.
所以我們看到.
這裡他兩次提到耶和華.
其實在整段的敘事文體裡.
在二十五章裡.
二十六節.
二十八節.
二十九節.
三十一節等等.
提到八次.
八次提到耶和華.
大家自己讀讀.
所以我們看到.
在阿彼該整個陳述裡.
他不斷提到上帝 耶和華.
耶和華的福.
信靠耶和華.
耶和華會賞賜你.
不要得罪耶和華.
耶和華會紀念你.
這就代表了他對神的關係.
所表達出來的信心.
在這裡有一個更加預言.
大衛必能作王.
將來你必作王.
你記得大衛當時正在流離.
他當時正在逃難.
我們記得剛剛二十四章.
他有殺掃羅的機會.
他沒有殺.
還在逃難.
來到這二十五章還未完.
還有二十六章.
二十六章仍然被掃羅追殺.

$^{1161}$但我們發覺.
他來賦予一個預言.
相信神的旨意會成就.
所以這個有神學的意義.
在二十四,二十五,二十六章.
其實有一個預言的作用.
除了表達出大衛以善勝惡之外.
這幾章裡都表達出.
原來透過竟然掃羅.
都說大衛將來會作王.
然後阿彼該又說.
大衛將來會作王.
為什麼呢?.
因為這已經是在第十六章.
神高了大衛之前.
在之前發生的事.
他們知道大衛會作王.
但是很多人不相信.
大衛固然都在掙扎.
但他再成為了大衛的鼓勵.
告訴他.
你將來必定照著神的應去作王.
所以你看到.
阿彼該有遠見.
有看見神的應去.
他有信心去知取.
他明白到神的旨意.
這個就是他屬靈的地方.
留意,當大衛後來.
拿伯死了,那怎麼辦呢?.
哇,這麼好的女人.
不要浪費了.
多大都要娶她做太太.
就娶了她做太太.
請問,如果你是阿彼該.
跟不跟大衛?.
跟?.
想清楚.
老公死了.
你就怎樣?.

$^{1201}$你發達了.
因為他是大富戶.
老公死了.
他承繼了所有的產業.
大衛沒有殺死他們的人.
他有這麼多產業.
打屁腳,享受生活都可以.
跟大衛?.
大衛當時在做什麼?.
逃難中.
沒錢.
窮淡.
跟不跟?.
再問一次.
跟.
證明你是屬靈.
因為他有屬靈的眼光.
他知道將來大衛必作王.
做王就在那裡做那八個寡婦.
所以如果大家記得上次.
我分析他一個叫假皇帝.
一個叫真天子.
大衛是真天子.
所以我們知道.
阿彼該屬靈的眼光.
他的信心.
表明了他的靈性的美好.
所以親愛的姐妹們.
我們說回來.
有美貌源於有美德.
有美德源於有好的靈性.
這是我們的源頭.
我們不得不有一個正確的價值觀.
合乎聖經倫理的道德觀.
以神的話語為我們人生的中心.
以神的應許為我們信心的依據.
來將這些應用在我們的人生.
和我們的家庭.
我們的丈夫.
我們的兒女身上.

$^{1241}$Amen.
這就是蒙福.
這就是好的女人.
當我們有一套以聖經健全的價值觀.
和道德觀念.
例如敬畏神.
有永恆的價值.
追求永恆.
有人愛.
誠實.
信心.
德行.
聖潔.
忍耐.
信實.
愛.
捨己.
要恕.
等等.
這是聖經所強調的靈性表現.
面對我們人生大小的抉擇.
面對不同的挑戰.
成為我們家庭的指南.
這就是蒙福的地方.
有一位姐妹叫做Joyce Landreth.
她說到關乎一個美好的女性.
有幾方面.
第一.
她說到信心.
這就是她的靈性.
第二.
她說到祈禱.
所以當姐妹們都能夠追求.
當然包括弟兄.
我們有信心.
我們能夠透過祈禱.
信心勝過內心的恐懼.
祈禱勝過憂慮.
她懂得接納.
她接納自己.

$^{1281}$接納自己不會自卑.
接納別人不會驕傲.
她懂得饒恕.
能夠勝過憤怒.
有一個和平的心.
有一個美善的表現.
她有人愛勝過仇恨.
勝過惱怒.
這些性情是很寶貴的.
這些性情都是靈性的表現.
所以在這位姐妹裡.
她說到這些信心.
祈禱.
接納.
饒恕.
人愛.
聖經的教導.
她說這樣就可以達到.
包括了.
有一種凡事主的眼光.
看事情.
和活出屬靈的美善.
一種難以形容的甜美的表現.
她的內心散發出來的那種美.
一種甜美.
一種勝過試探引誘的力量.
一種關懷別人的神聖舉止.
一種關懷別人的靈性表現.
一種安靜而有屬靈影響的能力.
一種生命散發出各樣的香氣.
有幽默感.
有友情.
有副組織力.
懂得款待別人.
有知彩.
有誠實.
不斷學習.
不斷上進.
人生有方向和優先.
懂得欣賞別人.

$^{1321}$你好不好?.
真的很好.
所以我們一生不斷努力.
這個我們看見.
在聖經裡面給我們很多提醒指引.
從亞比蓋的描述裡面.
我們看見他信靠耶和華.
相信大衛在神的手中蒙母手.
他敬畏耶和華.
所以知道神揀選大衛.
他願意跟隨.
他知道大衛是神所使用的.
他有遠見.
所以願意成為他的妻子.
我們發覺後來.
當大衛繼續經歷他流浪的人生.
有很多很多.
遲些我會在這裡說.
亞比蓋相信.
這位有才德.
有智勇.
有遠見.
有能力.
一個太太.
是成為大衛一個美好的幫助.
他提醒了大衛.
以致大衛能夠恢復和神美好的靈性.
敬畏耶和華.
聽他的勸導.
不用血氣.
讓神自己親自成就.
所以最終他們能夠.
大家合一成為一對.
真的敬畏上帝.
信靠神的應許.
同心的一對美好的夫婦.
一個美好的靈性.
就是他們的根基.
我覺得我們作為夫婦.
作為男士女士.

$^{1361}$都能夠學習這些美好的榜樣.
上星期記不記得.
做一個好男人.
有三件事.
除了和娜伯比較之外.
大衛顯出來.
第一.
有沒有人記得.
三件事.
每個都是三個字.
記不記得.
第一好男人就是他快適路.
大衛很快地調整他的怒氣.
快適路.
在聖經裡說到.
我們能夠止住怒氣.
實在是美好的.
不輕易發怒.
這就是愛.
所以愛不輕易發怒.
所以我們看到.
真的值得我們學習.
男女都學.
第一個快適路.
第二個就是肯聽勸.
肯聽勸.
所以大衛聽到亞彼該的勸導.
馬上就聽勸.
這代表他有一個受教的心.
我們每個人都追求這個受教的心.
第三.
快適路.
肯聽勸.
第三就是懂得欣賞.
所以我們能夠懂得欣賞.
欣賞別人.
欣賞妻子.
欣賞.
令我們更加有一個美麗的人生.
好男人.

$^{1401}$好女人,今天我們說什麼?.
有美貌.
源於.
有美德.
有美德源於.
好的靈性.
不得不說.
我們每個人都做過.
做過好男人.
做過好女人.
和你旁邊的大型節目說一聲,做過好男人.
我們低頭祈禱,請祈禱.
我們很希望.
我們不單止彼此鼓勵.
在人的面前.
能夠成為一個好男人,好女人.
更重要的是在天婦女的面前.
我們都能夠被你稱讚.
能夠被你稱讚為.
好.
又良善.
又忠心的僕人.
願耶主你自己都親自幫助我們.
繼續在我們的教會.
在我們的社區.
更加在我們的家庭當中.
你要實踐美善.
追求合乎你的心意.
被你稱讚.
因為一切愛的源頭.
全因為你.
聽我們祈禱.
奉耶穌基督的名.
阿們.
\newpage



\section{箴言 31:1-31}
\label{sec:032119}
\textbf{智慧與恩慈:生命的織布者}~~~~袁惠鈞牧師博士
\newline
\newline
連結: \href{http://www.church.com.hk/acms/content.asp?site=cdc\&op=show\&type=product\&code=032119\&layout=sermon}{\texttt{032119}} ~~~~ 語音日期: 2026-05-10
\newline
\newline
\hyperref[sec:032118]{\small{< < < PREV SERMON < < <}}
~
\hyperref[sec:index_chronic]{\small{[返順時目]}}
~
\hyperref[sec:032120]{\small{> > > NEXT SERMON > > >}}
\newline
\newline
箴言 31:1-31
\newline
\begin{longtable}{cl}
\hline
\hline
章節 & 經文 (和合本修訂版)\\
\hline
31:1 & \begin{tabularx}{0.7\textwidth}{X} 瑪撒王利慕伊勒的言語，就是他母親教導他的。 \end{tabularx} \\ \\ \relax
31:2 & \begin{tabularx}{0.7\textwidth}{X} 我兒，怎麼了？我腹中生的兒，怎麼了？我許願而得的兒，怎麼了？ \end{tabularx} \\ \\ \relax
31:3 & \begin{tabularx}{0.7\textwidth}{X} 不要將你的精力給婦女，也不要有敗壞君王的行為。 \end{tabularx} \\ \\ \relax
31:4 & \begin{tabularx}{0.7\textwidth}{X} 利慕伊勒啊，君王不宜，君王不宜喝酒，王子尋找烈酒也不相宜； \end{tabularx} \\ \\ \relax
31:5 & \begin{tabularx}{0.7\textwidth}{X} 恐怕喝了就忘記所頒的法令，顛倒所有困苦人的是非。 \end{tabularx} \\ \\ \relax
31:6 & \begin{tabularx}{0.7\textwidth}{X} 可以把烈酒給將亡的人喝，把酒給心裡愁苦的人喝， \end{tabularx} \\ \\ \relax
31:7 & \begin{tabularx}{0.7\textwidth}{X} 讓他喝了，就忘記他的貧窮，不再記得他的苦楚。 \end{tabularx} \\ \\ \relax
31:8 & \begin{tabularx}{0.7\textwidth}{X} 你當為不能自辯的人開口，為所有孤獨無助者伸冤。 \end{tabularx} \\ \\ \relax
31:9 & \begin{tabularx}{0.7\textwidth}{X} 你當開口按公義判斷，當為困苦和貧窮的人辯護。 \end{tabularx} \\ \\ \relax
31:10 & \begin{tabularx}{0.7\textwidth}{X} 才德的婦人誰能得著呢？她的價值遠勝過寶石。 \end{tabularx} \\ \\ \relax
31:11 & \begin{tabularx}{0.7\textwidth}{X} 她丈夫心裡信賴她，必不缺少利益； \end{tabularx} \\ \\ \relax
31:12 & \begin{tabularx}{0.7\textwidth}{X} 她終其一生，使丈夫有益無損。 \end{tabularx} \\ \\ \relax
31:13 & \begin{tabularx}{0.7\textwidth}{X} 她尋找羊毛和麻，歡喜用手做工。 \end{tabularx} \\ \\ \relax
31:14 & \begin{tabularx}{0.7\textwidth}{X} 她好像商船，從遠方運來糧食， \end{tabularx} \\ \\ \relax
31:15 & \begin{tabularx}{0.7\textwidth}{X} 未到黎明她就起來，把食物分給家中的人，將當做的工分派女僕。 \end{tabularx} \\ \\ \relax
31:16 & \begin{tabularx}{0.7\textwidth}{X} 她想得田地，就去買來，用手中的成果栽葡萄園。 \end{tabularx} \\ \\ \relax
31:17 & \begin{tabularx}{0.7\textwidth}{X} 她以能力束腰，使膀臂有力。 \end{tabularx} \\ \\ \relax
31:18 & \begin{tabularx}{0.7\textwidth}{X} 她覺得自己獲利不錯，她的燈終夜不滅。 \end{tabularx} \\ \\ \relax
31:19 & \begin{tabularx}{0.7\textwidth}{X} 她伸手拿捲線桿，她的手掌把住紡車。 \end{tabularx} \\ \\ \relax
31:20 & \begin{tabularx}{0.7\textwidth}{X} 她張手賙濟困苦人，伸手幫助貧窮人。 \end{tabularx} \\ \\ \relax
31:21 & \begin{tabularx}{0.7\textwidth}{X} 她不因下雪為家裡的人擔心，因為全家都穿上朱紅衣服。 \end{tabularx} \\ \\ \relax
31:22 & \begin{tabularx}{0.7\textwidth}{X} 她為自己製作被單，她的衣服是細麻和紫色布做的。 \end{tabularx} \\ \\ \relax
31:23 & \begin{tabularx}{0.7\textwidth}{X} 她丈夫在城門口與本地的長老同坐，為人所認識。 \end{tabularx} \\ \\ \relax
31:24 & \begin{tabularx}{0.7\textwidth}{X} 她做細麻布衣裳來賣，又將腰帶賣給商家。 \end{tabularx} \\ \\ \relax
31:25 & \begin{tabularx}{0.7\textwidth}{X} 能力和威儀是她的衣服，她想到日後的景況就喜笑。 \end{tabularx} \\ \\ \relax
31:26 & \begin{tabularx}{0.7\textwidth}{X} 她開口就發智慧，她舌上有仁慈的教誨。 \end{tabularx} \\ \\ \relax
31:27 & \begin{tabularx}{0.7\textwidth}{X} 她管理家務，並不吃閒飯。 \end{tabularx} \\ \\ \relax
31:28 & \begin{tabularx}{0.7\textwidth}{X} 她的兒女起來稱她有福，她的丈夫也稱讚她： \end{tabularx} \\ \\ \relax
31:29 & \begin{tabularx}{0.7\textwidth}{X} 「才德的女子很多，惟獨你超過一切。」 \end{tabularx} \\ \\ \relax
31:30 & \begin{tabularx}{0.7\textwidth}{X} 魅力是虛假的，美貌是虛浮的；惟敬畏耶和華的婦女必得稱讚。 \end{tabularx} \\ \\ \relax
31:31 & \begin{tabularx}{0.7\textwidth}{X} 她手中的成果你們要賞給她，願她的工作在城門口榮耀她。 \end{tabularx} \\ \\
[1ex]
\hline
\hline
\end{longtable}
$^{1}$今天我要和大家說的是.
《箴言三十一篇》的財德的婦人.
剛才彭長老跟我們說的經文.
我們應該很熟悉.
我給她的題目叫做.
「智慧與恩慈 生命的織布者」.
財德的婦人.
這個好的媽媽.
就是一個生命的織布者.
她不僅是編織布料.
她也是生命的編織者.
今時今日.
我們知道婦女的地位不斷在社會提高.
無論是家庭,社會,國家.
婦女現在都有很大的影響力.
我們也知道.
在很多國家.
有婦女成為元首,成為總統.
影響力是越來越大.
我們今天看一看這個財德的婦人.
她在家庭裡面有很大的影響力.
這個屬靈的母親.
聖經裡面的教導又是怎樣的呢?.
很多時候當牧師在樹台上.
說《箴言三十一章》的財德的婦人的時候.
很多的姐妹就會很怒,很害怕.
為什麼會這麼害怕?.
他們會說這個《箴言三十一章》的財德的婦人.
都不知道是真是假.
是不是真的有一個這樣的人呢?.
這麼厲害的?.
聖經說得這麼好.
是不是一個真正的女人呢?.
牧師,你不要將我和一個虛構的女人來比較.
還要我們來學習這個虛構的女人.
其實這個財德的婦人.
跟我們在座當中很多的母親已經很相似.
有什麼是這麼相似的呢?.
首先都是一個妻子.
都是一個媽媽.

$^{41}$都是一個煮飯婆.
都是一個洗衣婆.
都是一個清潔的女工.
營養的專家,家庭的老師.
還有是司機,會計師,顧問.
在家裡排欄糾紛的又是她.
還做了很多其他的事.
所以這個女人跟財德的婦人已經有些接近了.
我們再看看在聖經裡面告訴我們.
這個財德的婦人是怎麼樣的.
首先我們要知道.
聖經告訴我們一個賢妻是什麼.
在《箴言》的十二章第四節說.
財德的婦人是丈夫的冠冕.
是老公的冠冕.
每個人看到老公戴著冠冕漂亮嗎?.
財德的婦人就好像你老公的冠冕一樣.
還有《箴言》的十九章十四節.
房屋財產是祖宗遺留的.
惟有賢慧的妻是爺和華所賜的.
一個好的妻子不單是丈夫的冠冕.
還是臣賜給丈夫.
是爺和華所賜的.
今天的經文一開始《箴言》的三十一篇第一節.
就說利姆伊勒王的言語.
是他母親教訓他的真言.
我們知道《箴言》的三十一篇是誰作的.
是利姆伊勒王.
是他母親給他的教訓.
他就將它寫在《箴言》的第三十一篇.
你說這個利姆伊勒王是誰?.
有人說可能是所羅門王.
因為《箴言》是所羅門王所寫的.
有人說利姆伊勒王就是他的別名.
也有人說是希西加王.
無論這個利姆伊勒王是誰.
我們看到他的母親給了他很好的教導.
有什麼教導呢?.
第一.
在《箴言》的第三十一篇第三節.

$^{81}$叫他不要犯奸淫.
然後在第四至第七節.
叫他不要醉酒.
「兒子,你不要喝那麼多酒」.
還有在八至九節.
說他要去周祭那些窮苦貧乏的人.
還要替他們伸冤.
你說為什麼可以為他們伸冤?.
這個利姆伊勒王是個王子.
稍後我們會看到.
他會坐在城門口.
和長老一起審判案件.
所以他絕對可以替人伸冤.
然後在今天的經文.
就是《箴言》的第三十一篇十至三十一節.
利姆伊勒王的母親教他.
如何揀選一個賢慧的太太.
賢慧的妻子.
在《箴言》的三十一篇第十節.
說「才德的婦人誰能得著呢?.
她的價值遠勝過珍珠」.
「才德」這個字代表力量有能力的意思.
我們稍後就會看到.
這個才德的婦人.
不單非常屬靈.
是有道德的觀念.
她是有智慧.
她也真是實體上很有力量,很有能力.
稍後我們會回來看.
找一個這樣的太太不是容易的.
《和合本》說.
比找珍珠更難.
珍珠在哪裡?.
在蠔裡面.
你能否看到珍珠在蠔裡面有多大?.
是看不到的.
你可能要打開很多個蠔.
你才能找到一粒真正名貴的珍珠.
找一個好的太太也不是一件容易的事.
在七十字譯本叫她做「寶石」.

$^{121}$在其他的經文叫她做「紅寶石」.
總而言之很難找到這麼好的太太.
我剛才說我很公平.
說了這麼多堂父親節的訊息.
我又要說這麼多母親節的訊息.
聖經也一樣這麼公平.
不單止《箴言》的三十一篇告訴我們.
一個好的女人很難找到.
在《箴言》的二十篇也告訴我們.
一個好男人一樣這麼難找到.
在《箴言》的二十篇第六節說.
「人多說說自己的人詞.
但忠信的人誰能遇見呢?」.
一個忠信的好男人.
一個好男人不是這麼容易找到.
我們說在《箴言》的三十一篇.
十至三十一節.
就是尼穆伊勒王給他的兒子的教訓.
教他如何找到一個好媳婦.
有沒有發覺.
十至三十一節是多少節的經文?.
二十二節的經文.
希伯來文有多少個字母?.
又是二十二個.
如果你留意的話.
其實在《箴言》三十一篇的十至三十一節.
是一首字母詩.
每一節的經文.
都是用一個希伯來文的字母來開始.
由Aleph到Beth到Gimel到Dalet.
然後去到Taf.
一共二十二個希伯來文的字母.
每一節就是一個字母開始.
為甚麼尼穆伊勒王要這樣來寫這個詩歌呢?.
容易記得.
他媽媽教他的東西.
他經常要放在腦裡記得.
所以就用一首字母詩來敘述.
在《箴言》三十一篇.
讓我們看到一個才德兼備的女人是怎樣的.

$^{161}$一共是由六個角度來看待這個女人.
第一個角度她是一個好的妻子.
第二個角度她是一個好的主婦.
第三個角度她是一個好的鄰居.
不單止對她的家人好.
她也對鄰居對外人好.
還有第四點她是一個好的老師.
第五點她是一個好的媽媽.
第六點就是最重要了.
你前面那五點你都做不到的.
如果你不是一個好的信徒是一個好的聖徒.
每一樣都是基於這一點.
當你是一個好的信徒.
你就能夠做到前面那五點.
首先我們來看這個女人是一個好的妻子.
《箴言》三十一篇第十一節.
她丈夫心裡倚靠她必不缺少利益.
這個女人值得她丈夫來信賴.
她是可靠的.
她丈夫在外面工作.
甚至有時可能會離開家庭一段時間.
但她知道她太太在家裡.
一切都做得很好.
她絕對很放心.
在三十一篇我們看到.
是講利慕伊勒王.
是一個王子.
在皇宮裡面有很多事務需要做.
有很多資源需要來管理.
不是一件容易的事.
但是這個利慕伊勒王的王妃.
她得到她的丈夫來相信她.
來管理一切的資源.
這個其實和猶太人的傳統是很不同的.
因為猶太人是大男人主義.
有很多時她把太太當成傭人一樣.
好像一個僕人一樣.
所以當她要遠遊的時候.
她會怎樣做?.
把家裡的名貴的東西全部鎖起來.

$^{201}$怕太太拿去變賣.
這個利慕伊勒王就不同.
他絕對信得過太太.
把一切的東西交在太太的手中來管理.
在《尊嚴》三十一章第十二節.
「她一生使丈夫有益無損」.
她永遠都是為她的丈夫來設想.
她追求的丈夫最大的利益.
她的丈夫在外面賺錢.
她就在旁邊來支持她的丈夫.
來鼓勵她的丈夫.
她每做一件事.
她每做所有的事.
都是為了她的丈夫.
為了她的家.
經文說「她一生」.
其實原文就是她生命中的每一天.
她生命中的每一天.
都是出於對她丈夫的愛.
不會因為環境改變.
而她對她的丈夫的愛會改變.
無論是好的時光.
是壞的時光.
無論是健康的時候.
或者是生病的時候.
無論是開心快樂或者是悲傷的時候.
無論是富足.
無論是虛空.
我說這些像什麼?.
結婚時的誓言.
很多時候我們會說同一樣的誓言.
這個才德的婦人也是一樣.
她一心向著她的丈夫.
在彼得前書的三章六節.
也給我們看到一個榜樣.
就是撒拉.
撒拉她以純樸的心.
事奉亞伯拉罕為主.
這個就是聖經裡給我們的教導.
也好像剛才我說的那一節經文.

$^{241}$彼得前書.
也好像提多書的二章第四節.
一個好的婦人.
會指教比較年輕的婦人.
會愛她的丈夫.
會愛她的兒女.
也好像今天的經文.
在《箴言》的第31篇第23節所說.
她的丈夫在城門口與本地的長老同坐.
為眾人所認識.
你說他是一個王子.
他當然是為眾人所認識.
為什麼他在城門口?.
我剛才說過了.
城門口就是長老來辦案.
判斷是非的地方.
他與長老坐在一起.
他之所以能夠安心地坐在城門口.
天天來工作.
是因為他知道他的太太.
張頭家打理得井井有條.
這個妻子願意為她的丈夫犧牲.
為她丈夫的事業為重.
安心地在城門口工作.
不單止他是一個好的妻子.
第二點,他也是一個好的家庭主婦.
現在有很多人說.
尤其是女權運動的人.
他們會說.
叫女人做家庭主婦的職責.
浪費了,貶低女性.
他們會這樣說.
但在聖經裡.
在神的眼裡.
家庭主婦是非常重要的職業.
一個家庭主婦.
就是在家裡管理家庭的一個.
在《新約》有一個字.
叫做「統治者」.
「愛寇despotel」這個字.

$^{281}$就是「統治者」.
這個家庭主婦是家裡的統治者.
正如提摩太前書的五章十四節.
說是治理家庭的.
又例如在提多書的二章五節.
這個字稍微有少許不同.
但出於同一個字根.
就是料理家庭的那一位.
把家打理得頭頭是道.
她的丈夫絕對不會缺少任何利益.
這個就是聖經裡面的家庭主婦.
在《箴言》的三十一篇.
一共講出了十二點.
告訴我們這個主婦是怎樣的.
第一點就是.
她是有手藝的主婦.
在《箴言》的三十一篇第十三節.
她尋找羊絨和麻.
甘心用手作工.
她是有好的手藝.
她是把羊毛和麻做成衣服.
炭紙.
炭紙這個字比較好的翻譯可能是窗簾布.
羊毛和麻有甚麼分別.
羊毛就是冬天穿的衣服.
能夠保暖.
麻呢.
麻是比較熱的時候可以穿.
是比較涼快的.
但是無論是羊毛或者麻.
這個婦人都會找最好的質料.
來為自己的家人做衣服.
所以無論甚麼季節.
她都已經做好很好的衣服.
給她的家人穿.
而且是最好的材料.
這個主婦絕對不是遊手好閒.
她是落手落腳.
她是充滿活力的一個主婦.
她是非常有手藝.

$^{321}$第二點.
她還是新勞的主婦.
森賢的三十一篇十四節.
她好像雙船從遠方運糧來.
好像一隻雙船把食物從很遠的地方運回家.
有人跟我說.
牧師我也是這樣.
我時時去Costco.
買很多食物回來給我的家人來吃.
你是怎麼去的?.
你是不是開車去的?.
開車就容易.
這個森賢的三十一篇裡面.
這個婦人沒有車.
分分鐘她要走路.
走到很遠的地方去買這些食物回來.
為甚麼走到這麼遠?.
身邊沒有嗎?.
因為她是選擇最好的食物.
也是選擇最合理的價錢.
所以她要去到很遠的地方.
不是說她家附近馬上就可以買得到.
她為她的家人去找食物.
不單是最好的價錢也是最合理.
為甚麼她願意走這麼遠?.
因為她愛她的家人.
她願意表達她的愛心.
現在我都知道有很多媽媽喜歡在網上查.
哪一樣食物有營養價值.
為自己的家人都會這樣做.
這一樣有沒有營養?那一樣有沒有營養?.
古時就沒有這些東西.
她只可以走到很遠的地方.
找一些最好的食物來給家人.
還有這個家庭的主婦.
不單是手藝好,她是辛勞.
她還是一個願意犧牲的主婦.
在《箴言》的三十一篇十五節頭.
說「未到黎明他就起來」.
在東方的社會.

$^{361}$其實他們的傳統是.
在家裡點一盞油燈.
這個油燈是一個很小的陶瓷做的油燈.
裡面是裝著油.
而燈芯是飄在油的上方.
那個燈要有油才可以點燈.
沒有油怎麼不可以點燈?.
這個就是在《馬太福音》二十五章裡面的原因.
十個童女記不記得.
晚上睡覺睡過頭.
不記得起來點燈.
那會怎樣?.
錯過了.
新郎回來他們都不知道.
是一個家庭主婦猶太人.
在東方的社會.
是家庭主婦晚上爬起來點燈.
你說搞甚麼鬼.
晚上睡覺你還點燈來做甚麼?.
我們現在家裡都有夜燈.
有盞小燈.
你每次都起來上廁所.
黑漆漆的撲死你.
所以猶太人和東方人都是一樣.
點一盞油燈.
是太太的責任.
是妻子的責任.
晚上起來點燈.
而這個女人呢.
一到黎明又爬起來.
為她的家人準備一天的食物.
準備粟米食物給她的家人.
其實在中東.
大部份的時間.
白天都是很炎熱.
晚上其實是最好做工的時候.
又安靜.
所以這個婦人.
就是一早爬起來做工.
她是願意為她的家庭犧牲.

$^{401}$還有第四點.
她也是一個管理的主婦.
在第十五節的尾.
說把食物分給家中的人.
將當作的工分派疲累.
她首先將食物分配好給自己的家人.
然後分派工作.
不要忘記這個是王妃.
她用不著這麼早起床.
不用這麼早起床.
但她不是.
一早就爬起來.
晚上還要點燈.
你想想.
她不是說太陽曬到八月十五她才起床.
她一早就起床了.
「工作」這個字.
原文的希伯來文就是「Hog」這個字.
就是在七十四日本《易成任務》這個希伯來字.
亦在《出埃及記》的五章十四節翻譯成「工作」.
她一早起床就將工作分配好疲累.
每個人應該做些甚麼.
她絕對是一個很懂得管理家庭的主婦.
還有她懂得做買賣.
在《箴言》的三十一篇第十六節.
說她想得的田地就買來.
用手所得之利.
栽種葡萄園.
她見到她家附近有塊好的田地.
她就買了回來.
買回來做甚麼呢?.
做葡萄園.
種葡萄.
她覺得這個價錢合理.
她就去買了.
你說她哪裡有錢買?.
你看看經文說甚麼.
用手所得之利.
是她自己賺回來的錢.
不是她老公的錢.

$^{441}$是她的私己.
這個女人她居然懂得做甚麼?.
經營小生意.
這就帶領我們去到下一點.
她亦是一個創業的主婦.
在《箴言》的三十一篇十九節.
她手拿纜線桿.
手把紡線車.
她將羊毛和馬紡成線.
然後將這些線織成布.
然後在這些布上畫了衣服的圖案.
剪了出來之後.
將它縫起來.
成為一件一件的衣服.
給她的家人.
你看看第二十四節說甚麼.
她作細麻布衣裳去出售.
又將腰帶賣予商家.
她不單為自己的家人製作衣服.
她製作了衣服做甚麼?.
去賣.
來賺錢.
她才有錢可以買葡萄園.
經文說商家.
其實如果你看原文的《希伯來文》.
商家這個字就是家男人.
就是那些斐理基人.
斐理基人是做甚麼的?.
我們知道斐理基人是行船的.
是水手來的.
她就是將這些衣服給了這些斐理基人.
去幫她去不同的地方來賣.
這個女人非常之明智.
她懂得管家.
她懂得理財.
她還懂得做生意來幫補家用.
來幫助她的丈夫.
她有創業的精神.
還有第七點.
她是有力的主婦.

$^{481}$我們剛才一開始的時候.
我說財德的婦人.
財德這個字是甚麼意思?.
就是有力量有能力.
你說為甚麼這麼有能力這麼有力量?.
你看看在《箴言》的三十一篇十七節.
她以能力促腰使膀臂有力.
照字面的解釋.
就是她有能力圍繞著她自己.
這些能力她怎樣得來?.
很多婦女現在去做甚麼?.
去健身室去健身.
來鍛鍊身體.
這個婦人她不是去健身室.
她日常的家務她都親力親為.
以致她有膀臂有力.
膀臂有力是否要來描述一個女人?.
當然不是.
女人就要苗條的身材.
苗條淑女.
膀臂有力,孔武有力.
是形容甚麼人?.
男人.
但她就是.
每件事她都自己做.
所以她有膀臂有力.
她為了她的家庭.
她不僅出錢.
她還出力.
還有第八點.
她是一個無私的主婦.
在主第十八節.
她覺得所經營的有利.
她的燈終夜不滅.
她在田裡種葡萄.
葡萄有修成.
她感覺良好.
她很開心.
不是為了自己.
你也看到了.

$^{521}$賺了錢要來做甚麼?.
是幫補家用.
所以她絕對不是為了自己的利益.
有很多人.
拿私房錢來做甚麼?.
用在自己身上.
但她用在家人身上.
所以她絕對是一個很虔誠的女人.
她亦是神所設計的女人.
為何我說是神所設計?.
記不記得在《創世記》的第二章.
神看到亞當獨居不好.
做了夏娃來做甚麼?.
經文說是要幫助亞當.
而這個婦人就是幫助她的丈夫.
她真的是神所設計的女人.
因為她覺得自己的工作有果效.
所以她開心到.
不是開心就馬上去購物.
很多女人開心到.
不開心又去購物.
開心又去購物.
但這個女人不是.
她開心的時候.
夜晚的燈終夜不滅.
開心的時候更加去工作.
更加為她的家.
那盞燈終夜不滅.
她完全是為了丈夫和家人的利益.
她絕對是一個無私的女人.
還有第九點.
她是一個不單是內奴的女人.
她還是一個內涵的女人.
在《真言》的三十一篇二十一節頭.
她不因下雪為家裡的人擔心.
你說耶路撒冷有雪下的嗎?.
是啊.
隔一年就會下雪了.
耶路撒冷.
我記得我女兒住在達拉斯的時候.

$^{561}$她跟我說.
「爸爸,達拉斯也會下雪的」.
「隔一年就會下雪了」.
「但近年好像每年都會下雪」.
耶路撒冷也是一樣.
你說為什麼耶路撒冷會這麼冷?.
耶路撒冷在山上.
摩利亞山和錫安山長高.
所以海拔很高.
所以非常寒冷.
古代的家庭都沒有暖氣.
我們真是舒適.
住在這裡,尤其是加州.
冬天又說有暖氣.
夏天又說有冷氣.
其實相差多少度?.
沒有多少度.
但我們都有暖氣和冷氣.
你想想.
那些猶太人.
冬天冷得鼻子都脫了.
他們沒有暖氣.
他們怎樣?.
用炭爐來取暖.
全家人圍著炭爐.
真的要很小心.
你在家裡燒炭.
空氣不流通是很危險的.
我們剛才看到.
為什麼這個婦人不會擔心?.
她已經為她的家人做好了羊毛衣.
所以她知道她不用怕.
她的家人一定不會著涼.
第十點.
她也是一個籌備的主婦.
在《箴言》第31篇第21節.
因為全家都穿著珠紅衣服.
我相信這些就是羊毛衣服.
是冬天來寒的.
很保暖的.

$^{601}$她還染上了珠紅色.
這個女人真的有品味.
我們知道古時那些人很喜歡珠紅色.
記不記得《史徒行傳》第十六章.
那個婦人叫做呂底亞.
她是做什麼生意的?.
紫色布.
高貴大方.
個個都喜歡.
這個女人把自己的羊毛衣服染成了珠紅色.
中國人喜不喜歡大紅大紫?.
當然喜歡.
很意頭.
紫色好像還是皇上才穿的.
紫色和黃色.
她還用了深色.
做冬天的衣服.
用深色真的很有智慧.
我現在記起.
我現在是牧師.
我以前是熱學的專家.
深色的衣服是怎樣的?.
深色的衣服的吸收能力特別高.
吸光,吸熱.
所以這個婦人真的很有智慧.
她不單做的衣服是美觀的.
還是實用的.
她也不怕天氣寒冷.
因為她已經為她的家人準備好一切.
她也是一個很愛護的主婦.
在第31篇的第22節.
她為自己製作繡花毯子.
她的衣服是細麻和紫色布做的.
字面上這個毯子比較好的翻譯.
是床單,枕頭套,被子.
她親手為她的家人做的.
她要她的家人睡得舒服.
有很漂亮,很實用.
又美觀的床,床單,被子.
她為她的家人準備好一切.

$^{641}$這出於她犧牲她的牧師.
也是出於她對她家人的愛心.
最後一點.
這也是一個美麗的主婦.
你說牧師?你有沒有搞錯?.
還說她是美麗的主婦?.
這個女人一整晚都不睡覺.
她的眼睛不是像熊貓一樣有黑眼圈.
這樣就奇怪了.
她還去做什麼?.
織布,耕田.
兩隻手髒得蓬頭垢面的樣子.
就是她了.
還怎樣?.
旁臂有力,像個男人.
你不要說我聽她美麗.
她簡直就是骯髒,骯髒,污穢.
非常殘忍.
不,經文怎麼說?.
她在丈夫面前表現她自己的美麗.
她希望丈夫看到自己有多美.
在《箴言》三十一篇二十二節.
她說她的衣服是紫紅色染的細麻布.
雖然不是玲瓏綢緞.
雖然她也沒有珠光寶氣的首飾來襯托.
她就是穿著細麻布.
但顯得她很高貴.
我們聽到麻布就會覺得.
披麻戴孝,不妨好事.
不,細麻布染成紫色也能夠很高貴.
她完全在丈夫面前表現出.
她自己有高貴美麗的氣質.
待會我們會見到這個女人.
不單有外在的美.
她也有內在的美才是最重要.
而她在丈夫面前非常樸素.
不是很妖艷.
很樸素,剛剛好,不會過火.
非常適當.
真正的美麗是甚麼?.

$^{681}$這個婦人也明白到.
真正的美麗是內在的,不是外在的.
就好像保羅在提摩太前書的二章九節.
或者好像彼得在彼得前書的三節.
一個婦人真正的裝飾是甚麼?.
是她的純潔,端莊,美德.
是她的敬虔.
全部都是內在,不是外在.
雖然我們見到這個婦人.
在她的丈夫面前表現自己的美麗.
但更重要的是她的內在美.
而不是她的外在美.
講完了十二點.
這個才德的婦人,她是一個好的主婦.
第三點,她是一個好的鄰居.
在《箴言》的第三十一篇第二十節.
說她將守周濟困苦的人.
又伸手幫補窮乏的人.
她不單止對她的家人充滿了愛心.
她也對家外的人一樣充滿了愛心.
我們見到她為她的家人做甚麼?.
就是做衣服,預備食物給他們.
所以我相信她對家外的人.
對鄰居有需要的人都是一樣.
一樣做衣服,一樣做食物給他們.
令我想起《使徒行傳》第九章.
有一個叫多加的姊妹.
不知道大家是否記得.
她生前是廣行善事.
她多思周濟.
她幫了很多人.
以致她死了之後.
那些她幫過的寡婦.
站在使徒彼得旁邊哭泣.
她把多加生前所做過的女衣,外衣.
都給彼得看.
很懷念這個多德的姊妹.
這個多德的姊妹.
就是周濟窮人寡婦.
她是一個虔誠的女人.

$^{721}$她是一個很有美德的女人.
就好像我們今天《箴言三十一篇》所說.
在《箴言三十一篇》說這個婦人.
說她向有需要的人伸出雙手.
她對鄰居充滿了憐憫,充滿了慷慨.
她不單止對她的家人好.
她對外面的人一樣都是那麼關心.
一樣都是對他們那麼好.
第四點.
我們說一共要說這個才德婦人有六點.
第四點就是她是一個很好的老師.
在第二十五節.
說能力和威而是她的衣服.
一個好的老師不是紙上談兵.
尤其是說到一個媽媽.
真的要以身作則.
你說一套做一套是不行的.
一個好的.
如果你在神學院教神學.
你可以說一套做一套.
你在學生面前可以表現得自己很淑齡.
你有很好的教導.
但是下課後你會怎樣呢?.
沒有人知道.
但是一個媽媽呢.
你教完你的子女能不能拍拍屁股離開?.
不能夠的.
你的子女跟你住一間屋.
你這裡教完他.
你那裡要讓他們見到.
你要活出你自己的教導.
這個就是以身作則.
你要實踐你的教導.
如果你實踐不到的話.
你的子女會不會信服你?.
當然不會.
我們見到經文說這個婦人.
她就是以能力和威而就是她的衣服.
她是有屬神的品格.
她在靈命上也是非常強.

$^{761}$也是非常淑齡的一個女人.
經文說威而是甚麼?.
威而是一個境界.
這個婦人提高了她自己的境界.
因為她自己所教的她是要做出來的.
所以她將自己的境界提高了.
她不是外在的.
她教的東西是內在的.
讓她的子女見到.
她的確能夠實踐她自己的教導.
你問牧師.
牧師在《箴言》的31章25節.
只說能力和威而是她的衣服.
哪裡有說到以身作則.
她是一個好的老師.
你要看下面那一節經文26節.
說「她開口就發智慧.
她舌上有仁慈和法則」.
這個婦人是怎樣?.
她是以仁慈來引領她的家人.
我相信在她的家裡.
真正教子女的是她的丈夫.
但是這個婦人.
她仍然是以身作則.
而且她是怎樣?.
當她開口教她子女的時候.
是充滿仁慈和智慧.
我們見到她的態度是怎樣?.
是仁慈的.
是充滿了仁慈.
而她教的東西是神的法則.
是神的智慧.
其實以法則.
以仁慈來教子女.
容易嗎?.
絕對不容易.
很多時候被子女氣得怎樣?.
生蝦咁跳.
氣得你生蝦咁跳.
你還要怎樣?.

$^{801}$還要仁慈的態度,親切的.
還要善意,柔和地來教你的子女.
真的不容易.
但是這個女人.
這個才德的婦人.
她就做得到.
而她的口流露出來的是甚麼?.
是從神而來的智慧.
從神而來的法則.
我們剛才說二十五節的時候.
似乎還說漏了半節的經文.
哪半節?.
就是二十五節的尾.
「她想到日後的景況就喜笑」.
她為了家人一切安排妥當.
所以她絕對不會有甚麼憂慮.
因為她知道自己所信靠的是誰.
她只需要做好自己的本份.
她知道一切掌握在神的手裡.
這些她不需要擔心.
她所要做的是自己的本份.
她安排好一切為她的丈夫.
以至她的丈夫得到所有的利益.
她亦為她的家人以身作則.
教好她自己的子女.
她就不怕.
她知道她的子女長大之後.
一定會做得好.
因為她自己是真正以身作則的教導.
她的丈夫亦會安然.
完全有益無損.
這是實質上她知道這一切都會發生.
就是因為她做好自己的本份.
但在屬靈方面.
她亦知道全家人都有永恆的保障.
她全家人都是信靠神.
她是教出神的智慧,神的法則.
以至她知道她的家人都有永遠的盼望.
所以這裡就讓她看到.
無論實質上一切是否妥當.

$^{841}$屬靈上.
她的家庭在屬神的角度裡.
永遠會得到神的祝福.
所以當她想到日後的景況.
她就會喜笑.
她就會很安樂.
不單止她是一個好的妻子,主婦,鄰居,老師.
她亦是一個好的媽媽.
在第27節.
「她觀察家務並不吃閒飯」.
觀察家務是甚麼意思?.
即是管理好自己的子女.
管理好她家裡的資源.
不吃閒飯是甚麼意思?.
就是她不會去享用別人的成果.
她得到一切都是靠她自己的努力.
是她自己付出的代價.
所以她不需要丈夫來養她.
而她的子女亦都是尊重她,尊敬她.
這就是聖經裡面給我們的教導.
一個媽媽.
她是人生的上半場.
是將她的投資放在家庭和子女身上.
撫養她的子女.
然後下半場是收割.
當她的子女長大之後.
她的子女就會用餘下的生命.
來報答這個媽媽.
來照顧她.
為她們獻出生命的女人.
這正是神給每一個家庭的設計.
當子女長大之後.
亦都會有自己的子女.
而他們的子女亦都會明白媽媽的苦心.
記得媽媽溫柔的教導.
智慧的建議.
有愛心的管教.
以身作則.
記得媽媽無私的付出.
這些永遠都會在子女的心裡記得.

$^{881}$亦都會傳去下一代.
就是這樣一代傳給一代.
這就是神為家庭所設計.
在《箴言》的第31篇第28節.
說她的兒女起來稱她有福.
她的丈夫也稱讚她.
當我們投資在我們子女的身上.
他們會回報.
子女和丈夫都會稱讚這個婦人.
這個稱讚還在哪裡呢?.
《箴言》的第31篇第31節告訴我們.
她的丈夫會在城門口來稱讚自己的太太.
我們知道城門口是非常多眼的地方.
他就在那裡來稱讚他自己的太太.
最後一點.
不單止她是一個好的媽媽.
我們說她要做一個好的主婦.
一個好的妻子.
一個好的鄰居.
好的老師.
好的母親.
都基於甚麼?.
基於這個人屬靈的質素.
就是她要首先是一個聖徒.
她才可以做得到這一切.
真正能夠做得到這一切.
很可惜這個俗世.
當我們選擇太太的時候.
我們選擇甚麼?.
都是選美貌,學歷,個性,嗜好.
或者是有沒有成就.
我們都選擇這些.
但真正應該選擇的是甚麼?.
是內在的美德,外在的品格.
是屬靈的質素.
而不是外在的質素.
在《箴言》的第31篇第30節.
說「艷麗是虛假的,美容是虛浮的」.
「艷麗」這個字源民的希伯來文.
就是「charming」的意思.

$^{921}$這個女人很「charming」.
我們就是看外表.
看這個女人夠不夠「charming」.
夠不夠美貌.
其實這個字源民的意思.
是身體的形態.
身體的形態重要嗎?.
有人說重要.
但形態會怎樣?.
會漸漸衰老.
這是聖經教導我們.
你可以去健身室.
來訓練你的身體.
訓練得很好.
很「fit」,很有型.
但始終都會衰老,衰殘.
但內裡的質素是永遠不會衰殘.
這就是聖經教導我們.
但很可惜.
一般俗世的人都追求甚麼?.
就是追求艷麗,追求美容.
這一切都是虛假.
這一切都是沒用.
在《書與真言》的第31篇第30節.
「或為敬畏耶和華的婦女必得稱讚」.
這個才德的婦人就是一個愛神.
一個真正的敬拜者.
是一個真正敬畏神的人.
能夠娶到一個敬畏神的太太.
這是所有男人夢寐以求.
應該尋求的女人.
就是一個敬畏神的人.
娶到一個敬畏神的人.
才是一個男人一生中最大的祝福.
因為一年過一年.
我們身體的形態,外在都會衰敗.
但內裡的質素會越來越美麗.
如果一個女人是敬畏神的話.
外在是衰敗.
但她的內裡會越來越淑齡.

$^{961}$所以這才是最重要.
Matthew Henry說.
《真言》的第31篇.
其實是每一位基督徒的婦女.
都應該用來照自己的鏡子.
用《真言》的第31篇來自省.
看我做成怎樣.
是否一個好的媽媽.
是否一個好的母親.
在結束的時候.
我也想跟大家總結一下.
一個才德的婦人究竟是怎樣.
一個才德的婦人.
是每一個男人都希望成為妻子.
亦應該是每一個婦人都渴望能夠成為.
她是一個對伴侶忠誠.
她懂得管理她的家.
她甚至對家外的人也有憐憫的心.
願意來照顧他們.
她有善良,慈愛和優雅的態度.
來教導她的子女.
把神的智慧,神的法則.
來教導她身邊的人.
她亦是一個盡了媽媽責任的人.
她得到她子女的尊敬.
她不是去尋求子女的讚美.
但當她這樣做的時候.
她的子女自然就會來讚美她.
我希望每一個聽我這篇道的婦女.
不一定是在座.
可能是網上.
都能夠以《真言》的第31篇.
來成為自己的座右銘.
來看看自己做得怎樣.
而在我們當中.
若果是未結婚的男士.
怎樣找太太?.
是照著《真言》的第31篇來找.
另外那些已經結婚的男士.
可以時時以《真言》的第31篇.

$^{1001}$來鼓勵他們身邊的配偶.
願神祝福他們自己的懷孕.
讓我們一起低頭禱告.
讓我們見到一個才德的婦人.
究竟是怎樣.
她是一個好的妻子.
是一個好的主婦.
是一個有憐憫的心.
來幫助家外的人.
還有她是一個好的老師.
一個好的母親.
最重要的她就是一個好的信徒.
因為我們知道她要做到前面那五點.
她真的很需要在屬靈方面.
是一個真的敬畏神的人.
是一個蒙神喜悅的人.
願意在我們當中或者網上的朋友.
所有的婦女都能夠以《真言》的第31篇.
成為她們的鏡子.
成為她們的座右銘.
願意未婚的男士.
都能夠以《真言》的第31篇.
來尋找她們將來的配偶.
而已經結婚的弟兄.
願意都能夠用《真言》的第31篇.
來鼓勵她們身邊的太太.
願意我們能夠在你的話語上繼續追求.
在你的話語上繼續能夠長進.
潔淨自己.
以致我們能夠成為一個蒙利喜悅的人.
將一切的榮耀還賺得貴給你.
\newpage



\section{瑪拉基書 1:1-5}
\label{sec:032120}
\textbf{誰曾愛我?}~~~~石偉漢牧師
\newline
\newline
連結: \href{http://www.church.com.hk/acms/content.asp?site=cdc\&op=show\&type=product\&code=032120\&layout=sermon}{\texttt{032120}} ~~~~ 語音日期: 2026-05-10
\newline
\newline
\hyperref[sec:032119]{\small{< < < PREV SERMON < < <}}
~
\hyperref[sec:index_chronic]{\small{[返順時目]}}
~
\hyperref[sec:032121]{\small{> > > NEXT SERMON > > >}}
\newline
\newline
瑪拉基書 1:1-5
\newline
\begin{longtable}{cl}
\hline
\hline
章節 & 經文 (和合本修訂版)\\
\hline
1:1 & \begin{tabularx}{0.7\textwidth}{X} 耶和華的話，藉瑪拉基傳給以色列的默示。 \end{tabularx} \\ \\ \relax
1:2 & \begin{tabularx}{0.7\textwidth}{X} 耶和華說：「我曾愛你們。」你們卻說：「你在何事上愛我們呢？」耶和華說：「以掃不是雅各的哥哥嗎？我卻愛雅各， \end{tabularx} \\ \\ \relax
1:3 & \begin{tabularx}{0.7\textwidth}{X} 惡以掃，使他的山嶺荒涼，把他的地業交給曠野的野狗。」 \end{tabularx} \\ \\ \relax
1:4 & \begin{tabularx}{0.7\textwidth}{X} 以東若說：「我們雖被毀壞，卻要重建荒廢之處。」萬軍之耶和華如此說：「任他們建造，我必拆毀；人必稱他們為『邪惡之境』，為『耶和華永遠惱怒之民』。」 \end{tabularx} \\ \\ \relax
1:5 & \begin{tabularx}{0.7\textwidth}{X} 你們必親眼看見，你們要說：「耶和華在以色列疆界之外必尊為大！」 \end{tabularx} \\ \\
[1ex]
\hline
\hline
\end{longtable}
$^{1}$今天和大家分享一個題目叫「誰曾愛我」.
誰曾經愛我們呢?.
或者我們會一直看這個題目.
誰曾經愛我們呢?.
特別在母親節的日子.
希望我們都想一想.
回想一下.
我們都是蒙愛的.
特別是我們是被母親生出來.
讓我們母親受了很多的痛.
很多的苦.
讓我們能夠出生.
但當中我們都要問一問.
我們有時候能不能夠珍惜.
母親對我們的愛.
又或者我們能不能珍惜.
神給我們的愛呢?.
事實上我上網去找.
一個這樣的情況.
就是找一下.
2026年十大搜尋的一些題目.
我發現原來上網一找的時候.
他告訴我.
原來我們今天上網找得最多的.
搜尋的次數是什麼呢?.
原來排第一的是YouTube.
原來都不是搜尋GBT.
原來是YouTube.
你發現原來YouTube是排第一.
搜尋的次數最多的.
接著就是搜尋GBT.
接著就是Amazon.
看看買什麼.
接著就是Google.
接著就是Gmail,Facebook.
接著就是天氣.
Google Maps等等.
但是你有沒有發現.
原來我們今天.
這十大搜尋的題目.

$^{41}$都不是我們要找的東西.
你聽得明白嗎?.
原來今天我們.
最多人搜尋的.
就是那些可以幫我們搜尋的搜尋.
聽得明白了吧?.
就是說我們希望找到一些東西的時候.
原來我們第一時間找的.
是找一些幫我們找到的東西.
但是最後這件事.
都未能夠告訴我們那個答案.
始終都還不是我們最終要找的一件事情.
原來有時候我們都是這樣的.
就是說我們努力去找.
原來我們找的那些.
都不是我們所要的.
今天我們在我們的生命裡面.
我們尋找了很多東西.
我們努力去建造了很多東西.
但是那些東西.
有時候往往會不會只是其中一個搜尋器.
並不能夠帶我們去到終極呢?.
所以今天我們去問一個問題.
就是.
這是不是你一生都在尋找的一個答案呢?.
就是哪一個愛我們呢?.
「誰曾愛我」就是今天這個題目.
而今天我選了一個比較不尋常的書卷.
就是選了馬拉基書.
因為通常你看到一些福音主日.
或者一些報道聚會.
通常都不會找馬拉基書的.
那我就是另類的牧師來的.
所以我找的都是另類的段落.
不過我想藉著馬拉基書裡面.
說的是一個什麼事情呢?.
其實馬拉基書是一封情書.
這封情書裡面就是在說神向以色列.
所發的一份愛情的傾訴.
或者是一個討論或者是一個爭辯.

$^{81}$因為原來神面對著祂所愛的百姓以色列.
但是這群百姓似乎不覺得神愛他們.
他們體會不到上帝的愛.
以致他們迷失在當中.
一直去尋找他們以為尋找的.
但是他們一直尋不到他們應該要尋找的.
所以今天我們藉著這個馬拉基書裡面.
有三個重點我們去看一看.
怎樣去藉著上帝與以色列人的對話.
從現在想一想我們與神之間的關係.
這段聖經裡面在一章二節.
一開始的時候已經說了一個很重要的事情.
就是耶和華說我曾愛你們.
耶和華說我曾愛你們.
一開始的時候你發現原來神就在做什麼呢?.
在示愛.
我不知道你試過追女孩沒有?.
有沒有試過被男孩追?.
對方會向你示愛去表達他對你的愛.
原來這裡聖經裡面一開始就說.
耶和華我們的神.
都是一個會向我們示愛的神.
而今天我們也要問.
我們能不能夠在愛的路途裡面找到祂自己呢?.
今天也說了一個見證就是.
你知不知道牧師怎樣信耶穌的呢?.
剛才你剛剛聽到霞姐怎樣去經歷信耶穌的.
深一點的一些片段.
有沒有想過牧師怎樣信耶穌的呢?.
有沒有人聽過說我信耶穌的見證?.
沒有人聽過嗎?.
真的沒有人聽過嗎?.
你們聽過嗎?.
不知道.
可能大家都沒有印象.
原來真的很少聽過.
原來我們上一個禮拜有團契裡面分享.
分享你怎樣信耶穌,有沒有其他的信仰經歷.
其中上個禮拜我也說了一點.
原來在我年輕的時候.

$^{121}$當我還十多歲的時候.
其實我也很想找到我的信仰.
找到我的神.
原來我小時候十幾歲的時候.
原來已經很想知道生命裡面需要有一個神.
在我生命裡面去塗造我.
去建立我.
或者幫我去認識我的生命應該怎樣去走.
所以原來在我年輕的時候.
我自己就會怎樣呢?.
原來我就是背這個經.
大家背過這個經沒有?.
叫《般若波羅密多心經》.
有沒有背過?.
沒有背過.
這個其實是一個佛教的經文.
那時候原來我自己對佛教有很大的期待.
會不會在當中找到我的信仰呢?.
因為我很喜歡那種有即是沒有,沒有即是有.
很虛幻飄渺的那種感覺.
很神秘的一幕.
所以我自己覺得這個經文是一個很有趣的經文.
所以我就背了它.
不過背了之後也不太明白它說什麼.
已經很努力了.
背了之後就很熟悉了.
然後隔了兩天又忘記背什麼.
因為原來我根本就不明白裡面說什麼.
以致我一直在認真去看.
它能不能夠滿足我去尋找我的信仰的路呢?.
因為我在我的路裡面.
其中一個很重要的就是尋找一個愛我的上帝.
一個愛我的神.
當時我不知道叫什麼上帝.
一個愛我的神.
於是我就再花點時間去了解一下.
佛教裡面怎麼看愛呢?.
怎麼看愛呢?.
原來佛教裡面認為痛苦的根源就試探愛.
原來我們追求愛反而是痛苦的.

$^{161}$追求愛反而是痛苦的.
所以我們應該不要追求太多或是貪得太多.
這個令我覺得很奇怪.
我覺得心裡面有一種渴求.
我是人來的.
我覺得很需要一份的愛.
去學懂愛別人.
也希望能夠得到別人的愛.
但是當我一直看佛教的書籍或者內容的時候.
我發現我找不到.
甚至乎我再看多一點內容的時候.
我發現究竟佛祖有沒有說自己是神呢?.
佛祖有沒有說自己是神呢?.
答案是沒有的.
他沒有自稱自己是一個創造宇宙的神.
他只是說他在當中理解佛的意思是什麼呢?.
他對人生一些生活裡面.
一些人的痛苦裡面有一些覺悟.
他只是認為自己是一個覺悟者.
原來佛祖沒有說自己是神.
他不敢說自己是神.
也不會自稱自己是神.
他只是覺得人生裡面很多反省之後.
他成為一個覺悟者.
對於我來說第一個問題就是.
既然他不是神.
究竟我現在看的這個不是我要找的神.
究竟在哪裡找到我所要找的那位上帝呢?.
就是這樣.
我就開始了一直去看一些不同的書籍.
直到有一天.
我看了一本書.
一本叫二人行的書.
不過你應該也是沒有看過的.
一本好像散文的書籍.
這本二人行的書裡面就說到.
當中有不同的作者.
有不同的作品裡面說到.
有些人有相愛的經驗.
有些人有失戀的經驗.

$^{201}$有些人面對親人離開的經驗.
但是這本書裡面很有趣的地方就是.
每一個在當中的人.
都說到他們的心靈都感受到愛.
在愛裡面得到安慰.
於是我就很奇怪.
為什麼這本書裡面的每一個人.
都有這種感覺呢?.
於是我直到我了解多一點這本書的時候.
就知道原來這本書.
講基督教的書籍.
是講上帝的愛.
就是這樣我就開始去追求.
去回到教會裡面.
有機會之後再回到教會裡面去認識神.
原來今天我們在說的是.
教會裡面我們在說的是什麼呢?.
就是基督教裡面說的是一個愛的信仰.
一個愛的信仰.
愛的本體就是神.
因為聖經很清楚說神就是愛.
神就是愛.
當我們說想追求一份的愛的時候.
原來在神裡面就找到.
神就是那個的愛.
神本身就是一個愛的本體.
我們怎樣能夠愛.
是因為由神去定義.
告訴我們怎樣去愛.
神也吩咐我們去愛.
以致我們能夠去明白什麼是愛.
所以聖經裡面有一句話.
就是說神愛世人.
甚至將他的獨生子賜給他們.
叫一切信他的.
不致滅亡.
反得永生.
耶穌自己就是親自去示範了.
什麼叫做愛.
祂顯示了祂對人的愛.

$^{241}$顯示了祂願意來到這個世上.
去顯明了祂這份的愛.
也讓我們每一個認識祂的人.
能夠學懂這一份的愛.
所以原來我們在人的愛裡面.
我們有幾種不同的愛.
當我們說人很想追求一份的愛.
和被愛的時候.
究竟人有什麼的愛在當中呢?.
其中一位作者叫C.S. Lewis.
路易斯.
他寫了一本書叫《四種愛》.
他正正好就是說.
愛有四種不同的程度.
或者不同概念的愛.
看看你有沒有這四種的愛.
第一種的愛叫做親愛.
親愛就是一種.
人最基本最普遍的.
通常就是父母和子女之間的愛.
兄弟姐妹之間的愛.
一種比較多的是不具選擇性的愛.
沒有選擇的.
因為血緣關係我們就成為了一家人.
我們就在這個環境裡面一起生活.
一起認識.
以致我們學習當中的愛.
這是第一種我們常常都會有的.
我們親人的愛.
所以今天我們的母親節.
就是說我們要去表達.
我們對母親的那份的愛.
這是其中一種.
第二種是什麼呢?.
第二種就是一種叫做友愛.
友愛就是說特別一些朋友之間的愛.
一種友情之間的關懷.
兩個人本來可能沒有關係的.
但是因為相識.
因為一些共同的價值.

$^{281}$一些共同的嗜好.
以致能夠彼此去欣賞.
以致一個人能夠有愛.
第三種愛是什麼呢?.
就是一種叫做情愛.
這種的情愛.
甚至乎可能被表達成為一種情慾.
一種性.
一種之間很想擁有對方.
很想和對方有一種深入關係的愛.
就是一種的情愛.
所以通常夫婦之間就是一種情的愛.
兩個人不單可以傾心事.
傾密談話.
更加可以進入一個親密的性的關係裡面.
但是魯易斯也說到.
原來人世間裡面還有一種的愛.
就叫做大愛.
這種的大愛裡面就是說一種無私的愛.
這種無私的愛就不是取決於對方的特質.
而是取決於你愛者對於被愛者的憐憫.
和那種的愛的關懷.
特別是原來我們都從上帝而來得到這份的愛.
這份的大愛裡面就是說.
一份他願意為你所愛的對象無條件作出犧牲.
不介意對方的樣貌.
不介意對方有錢還是貧窮.
不介意對方究竟有多少學歷.
不介意對方是否身體患病.
仍然願意在對方身上發動那份的愛.
去表達那種的關懷,問候,溫暖.
等於原來這裡說的四種的愛裡面.
其中去到最後這種的愛.
就是說一種從神而來的愛.
他願意愛世間屬他的人.
他就愛他們到底.
所以今天在這段經文裡面我們見到.
一開始的時候神就說什麼.
我曾愛你們.
神沒有介意以色列是一個什麼人.

$^{321}$或者一個什麼的民族.
不過上帝就確定了他揀選了以色列.
成為他所愛的對象.
可能你會說這個的神.
神說我曾愛你們是否是一段逝去的愛情呢.
不是的弟兄姊妹.
在我曾愛你們這句說話裡面.
原文解釋時間的概念是什麼呢.
就是我從過去到現在一直愛你們.
原來這一刻神就對以色列這個民族.
去發出一個愛的宣告.
就是說我從過去到現在.
原來一直都愛你們這班的以色列人.
哪怕你們本來是不值得愛.
甚至乎在《馬拉基書》裡面不斷說.
神願意在一個約的限制裡面去愛.
在一個約的限制裡面去愛.
神的愛是怎樣在約的限制裡面去愛呢.
因為原來神經常說同人立約.
make a covenant.
同人立約是代表什麼意思呢.
就是跟你簽一個合同簽一個合約.
以致雙方堅定守這個合約的關係.
而神本來從來都不需要同人有任何的合約.
有任何的立約.
不過他甘願降低自己同人納一份約.
去表明這個愛.
以致他在這個限制裡面的時候.
即使你不守約.
你不守規矩你違反了約.
他仍然堅持要守住這份約.
這個就是神對人的愛的宣言.
所以第一點我想大家要記住的就是.
原來神是那位主動向我們示愛的神.
祂不是一個虛無的偶像.
祂不是一個妄自稱為神的那些人.
而祂是真真實實的神.
甘願來到降下與我們人親近.
以致祂表明了祂很愛我們.
但原來不是人人都懂得愛的.

$^{361}$不是人人都覺得被愛的.
所以以後以色列怎樣去反駁上帝呢.
以色列人問祂說.
你作何事上愛我們呢.
即是說你何時愛過我.
不知道你有沒有跟媽媽說過這些說話.
你媽媽很辛苦.
你預備一個早餐讓你上學吃.
然後你說不好吃的.
你都不明白我的.
你都不愛我的.
有時候會不會這樣去回應你的母親呢.
現在這裡以色列人正正好也是這樣.
他們向上帝反駁.
你在何事上愛我們呢.
你不愛我的.
神你不是真的愛我的.
你說得很愛我你是假的.
以致在整個馬拉基書裡面.
我們見到一群百姓.
就好像一些反駁隊似的.
提出很多很多的反駁.
說神不是愛我的.
當中在整個馬拉基書裡面.
有很多段人對神的反駁.
包括你在何事上愛我們呢.
即是說你何時愛過我.
我不覺得.
我們在何事上藐視你呢.
我何時有不尊重過你.
我在何事上使你不潔淨呢.
我何時沒有聽你的說話.
後面又說.
主的公卓是不潔淨的.
公卓上的果子和食物是可藐視的.
看哪 這是何等煩鎖.
即是說神你做的那些事情.
不值得尊重的.
有時候我們會不會這樣跟我們的父母說.
接著右邊又說.

$^{401}$我們在何事上使厭煩呢.
凡所行的亦或說.
看為善並且喜愛他們.
公平的上帝在哪裡呢.
我不開心的時候.
你是否在我身邊.
我哭的時候.
你是否在我旁邊安慰我.
有時候我們對父母發出這些質疑.
我們如何轉向你呢.
我們在何事上搶奪你呢.
我何時用背部得罪你.
我何時有不理睬你.
是你不理睬我.
我們會不會這樣對父母說話呢.
我何時頂撞你呢.
事奉上帝是枉然的.
我們會不會.
用些說話頂撞上帝.
又說上帝的不是呢.
這裡說的是.
以色列百姓裡面.
正正好就是他們.
向上帝作出一連串的反駁.
否定 抗辯.
以致到.
我們說上帝原來不愛他們.
但我們剛才知道.
原來那位神是願意.
降低自己身份去示愛的神.
神又怎會不理睬我們呢.
但你有沒有想過.
原來我們的關係.
會這麼奇怪呢.
原來可能很多時候.
都是我們人的問題.
是我們人的問題.
是我們甘於拒絕那位上帝.
所以.
有本書裡面.

$^{441}$他就說到.
末日四騎士.
不知道你有沒有聽過.
這裡正正好.
說的是一種愛的盲點.
說的是一種愛的關係裡面.
原來特別是人與人.
夫妻之間等等.
原來本來關係好.
本來大家是相愛的時候.
原來我們日漸容許.
我們的關係在扭轉.
初初你還覺得對方很好.
但我們容讓.
某些習慣出現.
以致我們與身邊的人.
關係越來越差.
而當中.
葛文就說到.
其中.
末日四騎士.
在關係最差.
差到成為.
冰點的時候.
很多時候就會有這四種現象.
有這四種現象.
你想一想.
你與人的關係裡面有沒有這些現象.
第一個是什麼.
第一個就是批評.
批評的意思是什麼.
就是原來.
不只是對一些.
事物的事情的批評.
而是.
對於人.
對方的人格.
性格.
作出一些批評.
很多時候開始慢慢.

$^{481}$你不是在說這件事.
做得對與不對.
而是一開始就說是你.
如果不是你.
這件事就不會搞我.
就是你這個人這麼自私.
以致從來都不考慮.
我的感受.
一直都是你.
我們就會說這些批評的說話.
以致.
將大家的關係.
再一次拉遠距離.
以色列也有這樣.
以色列也有說.
拉遠上帝的距離.
第二種是什麼.
稱為輕描.
或者藐視.
藐視就是.
你不會給我優勝.
就是說你與人相處.
別人是劣等.
你是優等.
所以你對別人的眼光.
總是反白眼.
冷嘲熱諷.
我不用你教我怎麼做.
你這種學歷.
你這種背景.
不懂的.
你有沒有說過這些說話.
有時候我們與人說這些說話.
所以說你不用理我.
你沒有資格理我.
有時候我們都是這樣.
與父母不知不覺.
說了這些傷害性的說話.
有時候令到我們與親人的關係.
都不知不覺疏離.

$^{521}$第三種.
就是辯解.
當別人說了一些說話.
令你不舒服的時候.
你就下意識去發出一個攻擊.
發出一個攻擊.
以致希望將責任推卸.
反過來指責對方.
是不是這樣.
當有問題發生的時候.
當然第一時間保護自己.
別人是最差的.
別人是最不講道理的.
為了什麼呢.
為了保護自己.
是完全不會犯錯的立場.
這一種.
是一種阻斷溝通的形式.
是不能夠真正解決問題.
所以通常我們都說.
這不是我的錯.
是你遲到.
你才令到我做不到那件事.
是你不對.
最後一樣.
就是築牆或者叫冷牆.
冷牆是什麼呢.
如果你與人的關係.
是完全隔絕.
就是說你不願意.
再與他有任何的回應.
任何的對話.
兩夫婦.
見到面.
一句話都不說.
由早到晚.
一句話都不說.
建立了一個高牆.
就會怎樣.
關係就會疏離.

$^{561}$特別今天母親節.
我們怎樣對待我們的父母.
我們檢視一下.
我們會不會對他們作出一些批評.
藐視.
防禦築高牆.
以致有時我們不懂得珍惜.
母親.
或者父親對我們的關愛.
這裡說的一個情況.
就是原來以色列.
就是這樣.
以色列的百姓不斷說了很多抱怨.
投訴.
神不愛他 神不體諒他.
拒絕神 等等.
原來是他們拒絕上帝的愛.
各位朋友.
今天有時想想.
我們是否真的沒有聽過.
神很愛你.
我們是否真的沒有體會過.
神很愛你.
我們是否真的不知道.
我們有一位很愛我們的神.
在我們身邊.
我們需要有一個方法.
提醒我們.
我們不要身在福中不知福.
身在愛中不知愛.
我們要做一件事.
就是有一個愛的回憶.
愛的回憶裡面說的是.
你回想一下.
神說.
以素不是雅各的哥哥嗎.
我卻愛雅各.
護以素 使他的山嶺放良.
把他的地業交給曠野的野狗.
以東若說.

$^{601}$我們雖被毀壞.
卻要重建荒廢之處.
萬君子耶和華如此說.
任他們建造.
我必拆毀.
人必稱他們為邪惡之警.
為耶和華永遠牢牢之民.
原來神是在說.
當以色列投訴.
神不愛他們的時候.
神就說.
你回想一下.
我怎樣對你.
神就說我是愛雅各.
護以素.
愛以色列.
愛雅各 愛以色列.
神願意.
祝福他們.
這個民族好不好.
好不好.
如果你看歷史.
這些民族.
說不上有多好.
有多優勝.
一點都說不出.
但是神就說.
你要回想一下.
神怎樣愛你.
第三點.
當我們回望.
過去與歷史.
就知道神堅定的愛.
當我們回望.
就知道神一直.
維持著.
祂的愛.
即使我們不懂愛祂的時候.
就像我們媽媽.
跟媽媽吵架.

$^{641}$第二天.
仍然有早餐放在桌上.
準備給你吃.
當你.
還在發燒的時候.
媽媽.
都沒有理會.
跟你吵架.
不要跟你聊太多.
她仍然會進來.
定時定候摸摸你的溫度.
拿藥給你.
倒杯熱水給你.
想一想這些片段.
就好像以色列.
在雅各的時候.
雅各.
其中一個最大的問題是什麼.
就是不斷說很多謊言.
欺騙長子名分.
結果要流落.
去舅父那裡.
但在舅父那裡.
神有沒有說.
你這麼壞說謊騙長子名分.
不要理你了.
沒有啊弟兄姊妹.
他去舅父拉班那裡.
二十年的時間.
上帝祝福他讓他有好的能力.
養好好的羊.
到最後祝福他.
要帶著很多羊群.
很多牲畜.
離開了他舅父的地方.
神並沒有.
因為他的不是.
而離開了他.
以色列人出埃及的時候.
是不是有很多投訴.

$^{681}$投訴什麼.
沒東西喝.
沒東西吃.
又有追兵.
生活不好.
不斷投訴.
但上帝沒有說既然如此.
不理你們了.
不是啊.
給你吃足四十年馬騮.
有鵪鶉.
有水.
二百萬人口.
在礦業裡養活.
不是上帝的公認嗎.
回望這些歷史.
你就知道.
上帝一直看顧他們.
當韓國國家滅亡的時候.
他們被擄.
到巴比倫帝國之下.
但上帝沒有說不理他們.
沒有.
上帝仍要.
藉著他的手幫助以色列的百姓.
所以我們有.
但以理書.
在尼保格維薩之下.
也有什麼.
有以斯拉尼希米等等.
回來建造聖殿.
回來建造城牆.
甚至今天.
讓這個亡了的國.
能夠再次復興.
回望歷史何等重要.
所以.
剛才我們看到.
霞姐自己寫了很多見證.
正正就是說.

$^{721}$為什麼要出來講見證.
因為我們想.
你經歷過上帝.
你多點時間想一想.
上帝為你做了什麼.
上帝幫助過你什麼.
讓你能夠知道.
原來上帝真的很愛你.
剛才聽到霞姐的分享.
說到上帝很愛你.
所以今天我們也要想一想.
特別母親節的日子.
有沒有一些片段.
有沒有回憶想起.
你的母親很愛錫你.
有的請你舉手.
有沒有母親很愛錫你的片段.
有啊,每個人都有.
當你想起這些片段.
你會知道母親是愛你的.
我經常分享一個片段.
我母親很愛我的片段.
我自己覺得.
看你記不記得.
我記得.
我們香港經常住屋村.
我們的屋村.
一棟樓有十三層樓.
十三層的屋村.
屋村裡面.
我們住九樓.
我們住九樓的時候.
很簡單.
一個門口.
鐵枝窗.
窗花.
一間這樣的屋.
一家七口.
在那裡生活.
當我還很小的時候.

$^{761}$貪得意.
就搞家裡的東西.
就搞家裡的門.
就鎖住了家裡的門.
當時我母親在街裡.
所以回來的時候.
連鎖匙都開不到.
當時因為太小的時候.
又不明白原來開門可能很簡單.
但就不懂.
那怎麼辦呢.
我母親就想了一個方法.
就是.
上樓上那一間.
去到十樓.
走上去叫人開門讓他進去.
他在十樓的窗口.
爬下來做蜘蛛俠.
因為我們剛好.
那個鐵窗.
那些沒有鎖.
他爬下來的時候.
就可以開窗口進來.
就可以開門.
每一次我想起這個畫面.
就想起.
我母親連命都不管.
可以爬下來.
因為你知道跌下來是很大件事.
她不是一個.
手長腳長的人.
如果像Ivan那樣六尺高.
加上手長八尺.
可能這邊一爬下去已經碰到下面.
但不是.
我母親是一個五尺高的女人.
她要爬一個很遠的距離.
她就爬到當中.
就是為了要開窗.
讓人進去.

$^{801}$照顧我們.
在你裡面.
是哪一個回憶.
母親節值得我們想一想.
回憶想起那份的愛.
當我們每一次願意回憶.
就想起每一次那份的愛.
所以今天.
我們要想一想.
你有沒有在你的回憶裡.
想起一個上帝對你的愛和拯救.
剛才霞姐說到.
他的女兒被醫治的過程.
但你呢.
可能你還沒有信耶穌.
但你有沒有想想.
在你的人生路上.
這麼多年的歲月裡.
你有沒有試過因為一個祈禱.
經歷過上帝的幫助.
經歷過上帝的醫治.
剛才.
霞姐說找醫生.
拿第二次建議.
我記得Stella也說.
她的眼睛本來.
是去做現在的醫生.
之前的醫生.
但上帝有感動她.
找第二次建議.
去找另一個醫生.
有一個更好的出路.
有一個更好的做法.
她經歷神給她的提醒.
但你呢.
有什麼是神感動你.
讓你去經歷過祂的幫助呢.
有些人說.
他不信耶穌是因為.
他不覺得神很愛他.

$^{841}$會不會就像.
這群以色列人.
但你有沒有想想.
原來神很多次拯救你.
幫助你.
以致到.
其實神一直對你不離不棄.
其中一個.
神對我們最大的愛.
就是祂在十字架上.
這個片段是真實的.
神的愛就在此.
向我們顯明.
還記得剛才我說.
當我去找哪一個神愛我的時候.
直至到我去到.
教會裡聽到.
原來神為我犧牲.
為我捨棄我的生命.
原來神一早已經愛我.
我望回上帝的足跡.
原來我讀小學的時候.
在一間天主教的學校.
去到中學的時候.
在一間天主教的學校.
有機會聽多些關於耶穌的事情.
原來神很早已經為我.
有機會聽到.
關於祂為我捨棄生命的事情.
直至有一天.
我終於可以找回.
一個愛我的神的時候.
我就找到.
原來我一早已經有.
我一早已經聽過祂的名字.
只是我不懂得去找祂.
但是今天.
我很想大家去想一想.
特別如果你還沒信耶穌.
還沒決定信耶穌的話.

$^{881}$我很希望你認真想一想.
其實在我們生命裡.
有一個愛我們的主耶穌.
祂很真實.
在瑕疵的身上出現.
是真實在我生命裡出現.
是真實在剛才Stella的生命裡出現.
亦都真實在很多人的生命裡出現.
其實祂都很真實.
在你生命裡出現過.
可能你會說.
神啊.
你何時愛過我.
你可能不知道.
但神一直向你示愛.
一次又一次.
只不過我們.
建了一個高牆.
去阻隔見到上帝的愛.
但耶穌.
死在十字架上.
就是拆斷人與人之間.
的牆.
祂願意打通了.
你和上帝的關係.
讓你恢復在一個愛裡.
願意你能夠回望過去.
看到神一直與你同在.
我想我們一起.
有個祈禱.
在祈禱之前.
我們都希望在今天的福音主日裡.
我們都邀請.
如果你聽到剛才的訊息.
或者你自己.
在你過去的日子裡曾經聽過.
很多關於耶穌的事情.
你想一想.
你有沒有聽過耶穌對你的愛.
你想一想.

$^{921}$你是不是曾經一次又一次.
拒絕耶穌的愛.
或許你見過一些很奇怪的異象.
見到一些有神有佛跟你聊天.
但這些你見到.
那些神佛跟你聊天.
他有沒有說過他愛你.
我相信沒有.
但你有沒有聽過耶穌說他愛你.
他一次又一次說他愛你.
又或者在你的生命裡.
你曾經拜過很多偶像.
你幫我被幫我.
但這些偶像從來都沒有幫過你.
但你有沒有想起.
你有一次向神求.
他幫你的時候.
他真真實實幫助你.
他醫好你.
他拯救你在交通意外裡的命.
他拯救你在.
婚姻裡的困難.
他拯救你在家庭裡的糾結.
他拯救你.
你不知道怎樣幫助你的子女成長的時候.
他一次又一次幫你.
我希望你想一想.
你有沒有親自經歷過.
這位的神.
如果你未曾.
去決定去信耶穌.
我邀請你這一刻.
我邀請你閉上雙眼.
你問一問.
我想問邀請在待當中.
有沒有人未信耶穌.
你都願意想去信耶穌.
求耶穌去幫你.
求耶穌的愛去幫助你.
如果有請你輕輕舉起你的手.

$^{961}$你願不願意.
得到耶穌對你的幫助.
對你的請求.
如果你願意的話.
我想你輕輕舉起你的手.
告訴耶穌.
其實我都願意被你所愛.
我都很想更深的經歷你的愛.
我們一切祈禱.
親愛的天父.
我們去仰望你.
我們知道原來很多人.
都不知道原來神.
很愛我們.
求你讓我們多謝你.
讓我們知道你很愛我們.
求你讓我們.
常常睜開雙眼.
回望你的愛.
讓我們有一天.
能夠去捉緊你.
能夠相信你.
讓我們已經信你的人.
更加堅定的愛你.
願意對你的愛.
作出回應.
求你幫助我們.
藉著今天母親節的日子.
我們想起.
母親對我們的點點滴滴.
這份愛的時候.
我們願意去表達.
我們這份愛的感動.
無論打電話給她.
無論是.
打個視像給她.
又或者親自跟她說一些話.
或者請她吃一頓飯.
讓我們去真正珍惜.
我們的母親.

$^{1001}$因為她同樣是上帝所差派來.
去愛我們的.
願我們能夠回應.
這份上帝給我們的愛.
和這份給我們母親的愛.
我們願意將我們的祈禱.
都交在天父的手裡.
聽我們祈禱.
奉耶穌的名字而求.
\newpage



\section{約拿書 1:1-4:11}
\label{sec:032121}
\textbf{出走的先知}~~~~陳雅莉傳道
\newline
\newline
連結: \href{http://www.church.com.hk/acms/content.asp?site=cdc\&op=show\&type=product\&code=032121\&layout=sermon}{\texttt{032121}} ~~~~ 語音日期: 2026-05-10
\newline
\newline
\hyperref[sec:032120]{\small{< < < PREV SERMON < < <}}
~
\hyperref[sec:index_chronic]{\small{[返順時目]}}
~
\hyperref[sec:032122]{\small{> > > NEXT SERMON > > >}}
\newline
\newline
約拿書 1:1-4:11
\newline
\begin{longtable}{cl}
\hline
\hline
章節 & 經文 (和合本修訂版)\\
\hline
1:1 & \begin{tabularx}{0.7\textwidth}{X} 耶和華的話臨到亞米太的兒子約拿，說： \end{tabularx} \\ \\ \relax
1:2 & \begin{tabularx}{0.7\textwidth}{X} 「起來，到尼尼微大城去，向其中的居民宣告，因為他們的惡已達到我面前。」 \end{tabularx} \\ \\ \relax
1:3 & \begin{tabularx}{0.7\textwidth}{X} 約拿卻起身，逃往他施去躲避耶和華。他下到約帕，遇見一條船要往他施去。約拿付了船費，就上船，與船上的人同往他施，為要躲避耶和華。 \end{tabularx} \\ \\ \relax
1:4 & \begin{tabularx}{0.7\textwidth}{X} 耶和華在海上颳起大風，海就狂風大作，船幾乎破裂。 \end{tabularx} \\ \\ \relax
1:5 & \begin{tabularx}{0.7\textwidth}{X} 水手都懼怕，各人哀求自己的神明。他們把船上的貨物拋進海裡，為要減輕載重。約拿卻下到艙底，躺臥沉睡。 \end{tabularx} \\ \\ \relax
1:6 & \begin{tabularx}{0.7\textwidth}{X} 船長到他那裡，對他說：「你怎麼還在沉睡呢？起來，求告你的神明，或者神明顧念我們，使我們不致滅亡。」 \end{tabularx} \\ \\ \relax
1:7 & \begin{tabularx}{0.7\textwidth}{X} 船上的人彼此說：「來吧，我們來抽籤，看看這災難臨到我們是因誰的緣故。」於是他們就抽籤，抽出約拿來。 \end{tabularx} \\ \\ \relax
1:8 & \begin{tabularx}{0.7\textwidth}{X} 他們對約拿說：「請你告訴我們，這災難臨到我們是因誰的緣故呢？你做甚麼行業？你從哪裡來？你是哪一國的人？屬哪一族？」 \end{tabularx} \\ \\ \relax
1:9 & \begin{tabularx}{0.7\textwidth}{X} 他說：「我是希伯來人，我敬畏耶和華，天上的神，他創造了滄海和陸地。」 \end{tabularx} \\ \\ \relax
1:10 & \begin{tabularx}{0.7\textwidth}{X} 那些人就大大懼怕，對他說：「你做的是甚麼事呢？」原來他們已經知道他在躲避耶和華，因為他告訴了他們。 \end{tabularx} \\ \\ \relax
1:11 & \begin{tabularx}{0.7\textwidth}{X} 海浪越來越洶湧，他們就問他說：「我們當向你做甚麼，才能使海浪平靜呢？」 \end{tabularx} \\ \\ \relax
1:12 & \begin{tabularx}{0.7\textwidth}{X} 他對他們說：「你們把我抬起來，拋進海裡，海就會平靜了；我知道你們遭遇這大風浪是因我的緣故。」 \end{tabularx} \\ \\ \relax
1:13 & \begin{tabularx}{0.7\textwidth}{X} 然而那些人竭力划槳，想要把船靠回陸地，卻是不能；因風浪越來越大，撲向他們。 \end{tabularx} \\ \\ \relax
1:14 & \begin{tabularx}{0.7\textwidth}{X} 於是他們求告耶和華說：「耶和華啊，求求你不要因這人的性命使我們滅亡，不要使流無辜人血的罪歸給我們；因為你－耶和華隨自己的旨意行事。」 \end{tabularx} \\ \\ \relax
1:15 & \begin{tabularx}{0.7\textwidth}{X} 他們把約拿抬起來，拋進海裡，海的狂浪就平息了。 \end{tabularx} \\ \\ \relax
1:16 & \begin{tabularx}{0.7\textwidth}{X} 那些人就大大懼怕耶和華，向耶和華獻祭許願。 \end{tabularx} \\ \\ \relax
1:17 & \begin{tabularx}{0.7\textwidth}{X} 耶和華安排一條大魚吞下約拿，約拿在魚腹中三日三夜。 \end{tabularx} \\ \\ \relax
2:1 & \begin{tabularx}{0.7\textwidth}{X} 約拿在魚腹中向耶和華－他的神禱告， \end{tabularx} \\ \\ \relax
2:2 & \begin{tabularx}{0.7\textwidth}{X} 說：「我在患難中求告耶和華，他就應允我；我從陰間的深處呼求，你就俯聽我的聲音。 \end{tabularx} \\ \\ \relax
2:3 & \begin{tabularx}{0.7\textwidth}{X} 你將我投下深淵，直到海心；大水環繞我，你的波浪洪濤漫過我身。 \end{tabularx} \\ \\ \relax
2:4 & \begin{tabularx}{0.7\textwidth}{X} 我說：『我從你眼前被驅逐，然而我仍要仰望你的聖殿。』 \end{tabularx} \\ \\ \relax
2:5 & \begin{tabularx}{0.7\textwidth}{X} 眾水環繞我，幾乎淹沒我；深淵圍住我；海草纏繞我的頭。 \end{tabularx} \\ \\ \relax
2:6 & \begin{tabularx}{0.7\textwidth}{X} 我下沉到山的根基，地的門閂將我永遠關住。耶和華－我的神啊，你卻將我的性命從地府裡救出來。 \end{tabularx} \\ \\ \relax
2:7 & \begin{tabularx}{0.7\textwidth}{X} 我心靈發昏時，就想起耶和華。我的禱告進入你的聖殿，達到你面前。 \end{tabularx} \\ \\ \relax
2:8 & \begin{tabularx}{0.7\textwidth}{X} 那信奉虛無神明的人，丟棄自己的慈愛； \end{tabularx} \\ \\ \relax
2:9 & \begin{tabularx}{0.7\textwidth}{X} 但我要以感謝的聲音向你獻祭。我所許的願，我必償還。救恩出於耶和華。」 \end{tabularx} \\ \\ \relax
2:10 & \begin{tabularx}{0.7\textwidth}{X} 耶和華吩咐那魚，魚就把約拿吐在陸地上。 \end{tabularx} \\ \\ \relax
3:1 & \begin{tabularx}{0.7\textwidth}{X} 耶和華的話第二次臨到約拿，說： \end{tabularx} \\ \\ \relax
3:2 & \begin{tabularx}{0.7\textwidth}{X} 「起來，到尼尼微大城去，把我告訴你的信息向其中的居民宣告。」 \end{tabularx} \\ \\ \relax
3:3 & \begin{tabularx}{0.7\textwidth}{X} 約拿就照耶和華的話起來，到尼尼微去。尼尼微是一座極大的城，約有三天的路程。 \end{tabularx} \\ \\ \relax
3:4 & \begin{tabularx}{0.7\textwidth}{X} 約拿進城，走了一天，宣告說：「再過四十天，尼尼微要傾覆了！」 \end{tabularx} \\ \\ \relax
3:5 & \begin{tabularx}{0.7\textwidth}{X} 尼尼微人就信服神，宣告禁食，從最大的到最小的都穿上麻衣。 \end{tabularx} \\ \\ \relax
3:6 & \begin{tabularx}{0.7\textwidth}{X} 這消息傳到尼尼微王那裡，他就從寶座起來，脫下朝服，披上麻布，坐在灰中。 \end{tabularx} \\ \\ \relax
3:7 & \begin{tabularx}{0.7\textwidth}{X} 他叫人通告尼尼微全城，說：「王和大臣有令，人、畜、牛、羊都不可嘗任何東西，不可吃，也不可喝水。 \end{tabularx} \\ \\ \relax
3:8 & \begin{tabularx}{0.7\textwidth}{X} 人與牲畜都要披上麻布，切切求告神。各人要回轉離開惡道，離棄自己掌中的殘暴。 \end{tabularx} \\ \\ \relax
3:9 & \begin{tabularx}{0.7\textwidth}{X} 誰知道神也許會回心轉意，不發烈怒，使我們不致滅亡。」 \end{tabularx} \\ \\ \relax
3:10 & \begin{tabularx}{0.7\textwidth}{X} 神察看他們的行為，見他們離開惡道，神就改變心意，原先所說要降與他們的災難，他不降了。 \end{tabularx} \\ \\ \relax
4:1 & \begin{tabularx}{0.7\textwidth}{X} 這事令約拿大大不悅，甚至發怒。 \end{tabularx} \\ \\ \relax
4:2 & \begin{tabularx}{0.7\textwidth}{X} 他就向耶和華禱告，說：「耶和華啊，這不就是我仍在本國的時候所說的嗎？我知道你是有恩惠，有憐憫的神，不輕易發怒，有豐盛的慈愛，並且會改變心意，不降那災難。我就是因為這樣，才急速逃往他施去的呀！ \end{tabularx} \\ \\ \relax
4:3 & \begin{tabularx}{0.7\textwidth}{X} 耶和華啊，現在求你取走我的性命吧！因為我死了比活著更好。」 \end{tabularx} \\ \\ \relax
4:4 & \begin{tabularx}{0.7\textwidth}{X} 耶和華說：「你這樣發怒，對嗎？」 \end{tabularx} \\ \\ \relax
4:5 & \begin{tabularx}{0.7\textwidth}{X} 約拿出城，坐在城的東邊，在那裡為自己搭了一座棚。他坐在棚子的蔭下，要看看城裡會發生甚麼事。 \end{tabularx} \\ \\ \relax
4:6 & \begin{tabularx}{0.7\textwidth}{X} 耶和華神安排了一棵蓖麻，使它生長高過約拿，影子遮蓋他的頭，使他免受苦難；約拿因這棵蓖麻大大歡喜。 \end{tabularx} \\ \\ \relax
4:7 & \begin{tabularx}{0.7\textwidth}{X} 次日黎明，神卻安排一條蟲來咬這蓖麻，以致枯乾。 \end{tabularx} \\ \\ \relax
4:8 & \begin{tabularx}{0.7\textwidth}{X} 太陽出來的時候，神安排炎熱的東風，太陽曝曬約拿的頭，使他發昏，他就為自己求死，說：「我死了比活著更好！」 \end{tabularx} \\ \\ \relax
4:9 & \begin{tabularx}{0.7\textwidth}{X} 神對約拿說：「你因這棵蓖麻這樣發怒，對嗎？」他說：「我發怒以至於死，都是對的！」 \end{tabularx} \\ \\ \relax
4:10 & \begin{tabularx}{0.7\textwidth}{X} 耶和華說：「這棵蓖麻你沒有為它操勞，也不是你使它長大的；它一夜生長，一夜枯死，你尚且愛惜； \end{tabularx} \\ \\ \relax
4:11 & \begin{tabularx}{0.7\textwidth}{X} 何況這尼尼微大城，其中不能分辨左右手的就有十二萬多人，還有許多牲畜，我豈能不愛惜呢？」 \end{tabularx} \\ \\
[1ex]
\hline
\hline
\end{longtable}
$^{1}$今天我們一起來看約拿書的時候.
約拿書很短的 四章 不是很長.
要讀很快讀完 是不是.
甚至乎我們主學的老師.
應該每個人都可能都教過也說不定.
總是會教有條魚吞了若拿嘛 是不是.
是 若拿可以說是其中一個.
我們都很熟悉的舊約人物.
他是一位很出名的以色列先知.
第一章他講述的時候就講述到.
神就呼召他 吩咐他要去東邊走.
往東邊 你留意聖經是往東邊.
去亞述的尼尼尾城.
但是他偏偏往西走 走了去約柏.
然後在碼頭上船.
想著要去他斯 上了他斯的船那裡.
接著發生什麼事呢.
船就起大風浪.
大家就在想誰連累我們起大風浪的.
誰連累我們這樣的.
若拿就知道自己了 他就叫船員.
不要煩了 你們扔我下海吧.
扔我下海就沒事了.
接著在第二章 上帝就安排一條大魚.
我們很熟悉的 這條魚叫吞拿魚.
吞了若拿的魚 在漁福裡面.
當他非常接近死亡之際.
我們應該每個人都沒有試過.
被一條魚吞了進肚子裡.
我想你都覺得自己死定了.
在死定的時候 覺得他呼求上帝.
他呼求上帝.
他就向上帝禱告.
願意再重新一次進行上帝的旨意.
上帝就叫那條魚吐出來.
吐出來之後 上帝就要叫他做事了.
原本他去亞述的尼尼尾城宣講信息.
他去了 的確去了.
還要一講 全城人信主.
整個尼尼尾城馬上悔改.

$^{41}$上帝於是就決定不毀滅尼尼尾城.
於是呢 上帝叫你做的事做好了.
你完成了.
但是剛才主席帶我們讀的時候.
第四章怎麼說呢?.
他宣講 全城悔改.
他心裡很憤怒.
然後離開那個城 搭了座棚.
我就要來坐下來看看尼尼尾城的命運.
於是上帝安排一棵比馬樹.
一夜之間長大.
很高興 很開心 很舒服.
一陣涼風吹過來.
但是接著 神又把一條蟲蟲咬死了.
一夜又枯乾了.
接著熱不熱的太陽曬若拿.
曬到他又很生氣 生氣得要死.
上帝用這麼多的經歷.
要若拿去明白他的心腸.
這個大概就是約拿書所說的話.
當我們去看約拿書的時候.
我們不難發現.
其實若拿是大情大性的人.
其實約拿書形容若拿的性格很鮮明.
好像今天看話劇看電影.
那個主角的性格描寫得很鮮明.
我們能夠想像若拿經歷了很多的情緒.
甚至他真的感覺到很受傷害.
他很難過 他很生氣.
上帝讓他走到絕路.
是不是真的想逼死他呢.
其實不是.
上帝要他體會上帝的心腸.
所以今天我們希望透過約拿書.
我們再一次一起來體會我們上帝的心腸.
哪怕你已經讀過這本書多少次.
你很熟悉這個故事.
但我們一起來體會.
今天用體會的心.
用一個要去明白我們上帝的心.

$^{81}$這樣去聽這堂道行嗎.
其實約拿書好像一齣話劇.
一幕又一幕.
我們會看到若拿一開始不是不小心迷路.
我們在座的太太就知道.
初初來學駕駛的那些.
特別是沒有什麼方向感的那些.
經常都迷路的 是不是.
但是若拿不是迷路的.
他是故意的.
他故意不去東邊去西邊.
他不是不小心走錯方向.
聖經一次又一次提到.
他是下渡.
他先是下渡約柏.
接著下渡往他斯的船.
接著再下渡船艙底睡覺.
然後最終他被拋下海.
去到下渡漁福的深處.
這個一連串的下渡.
不是剛好這麼巧.
聖經的作者是刻意的描述.
若拿其實是走一條往下的路.
這條也是一條一步一步一步遠離上帝的路.
他一路下渡下渡下渡.
距離上帝越來越遠.
這個一路往下走.
不只是他人身體移動走了.
而是他的心,他的屬靈狀況.
都落到低處.
一路往下再往下.
若拿可能以為他能躲避耶和華的臉.
可以躲避開耶和華.
他卻不知道這個躲避.
其實是讓自己走向黑暗.
他抗拒上帝的吩咐.
他逃避上帝.
上帝就用狂風巨浪.
令他走投無路.
上帝用狂風巨浪.

$^{121}$用大雨吞噬.
不是真的想他死.
是要喚醒他沉睡麻木的心靈.
很感恩若拿真的喚醒了.
不過要落到餘福臨死那一刻.
才喚得醒.
確實是.
有時我們真的要走到山窮水盡.
走到覺得走投無路.
我們才想起上帝.
若拿在餘福中說.
我心靈發昏的時候.
我就想起耶和華.
我的禱告就進入你的殿.
達到你的面前.
他當時不在殿裡面.
他在雨的肚裡面.
但當他想起耶和華的時候.
他想起神.
在絕境絕望裡面想起神的時候.
他知道他的禱告.
能夠達到耶和華的耳中.
確實有時是這樣.
上帝會利用絕境.
利用說出.
我們才聽得見.
上帝有話對我們說.
當我們繼續被自己的情緒.
自己的悲傷.
自己的不開心.
自己的困擾困住我們的時候.
我們聽不到上帝的聲音.
當我們繼續在自己很多的籌算.
很多的計劃.
覺得我可以預備很多東西.
我們被這一切困住的時候.
我們忽視上帝的聲音.
反而當你落到雨幅裡面.
我不知道雨的肚子是怎樣的.
可能真的很安靜.

$^{161}$可能真的要去到絕望的境地.
你才停下來.
聽見上帝的聲音.
想起上帝.
本應是混亂翻騰的海.
卻是成為了若拿.
重新認識上帝的地方.
本身是象徵死亡的魚肚.
成為了若拿禱告之處.
其實有時上帝.
不是要我們面對困難的時候.
要我們有多偉大多聰明.
給我們多少智慧.
想出很多的解決方法來.
甚至上帝沒有要我們有超乎常人的技能.
然後幫助我們自己走出一些絕境.
上帝只是要我們在絕境裡面.
第一時間想起上帝.
在絕望裡面.
叫我們「你回來呼求我吧」.
這個往往是我們人生.
是我們遇上很多事情裡面的轉捩點.
在詩篇十八篇.
大衛也曾經這樣求告耶和華.
他說當死亡的繩索纏繞他.
死亡的網羅淋到他的時候.
他在急懶中求告耶和華.
向神呼求.
神就從電裡面聽他的聲音.
他也和若拿一樣.
相信他的聲音達到神的耳中.
上帝其實是盼望我們呼求他的名字.
和他對話.
確實今天我們有很多方法.
很多的智慧.
很多的計謀.
甚至你今天有AI.
無限地幫助你.
但我盼望.
那些是一些工具.

$^{201}$那些都是我們的謀算.
在困境裡面.
你才能動得呼求神.
弟兄姊妹.
我不知道你現在在什麼困境裡面.
可能是工作.
可能是經濟需要.
可能你和人際關係出現了問題.
身體軟弱.
和家人.
夫妻之間.
教養兒女的時候.
遇到困難.
還是事奉裡面都料理灰心絕望.
還是可能只是你單單在自己的靈性上.
落到低處.
你覺得這是一個絕境.
無論你困在什麼情況.
在餘福中與主對話.
在絕望裡面和神說話.
去尋求祂.
去呼求祂.
上帝會垂聽我們的禱告.
當我一直看這次在預備這堂道的時候.
我都想到在世界各地的宣教士.
很多的宣教士.
他們都好像若拿一樣.
被神父召去.
你去宣講悔改的信息.
你去傳悔改的福音.
其實他們都落在很多不容易的困境.
很多宣教士在前線都很疲乏.
很孤單.
很多的不明白.
亦都難以向人傾訴.
雖然今天上帝未必呼召我們在座.
每一個都出去做宣教士.
但我盼望我們不只是坐在這裡.
我們都能夠多一份關懷和關心.
對那些在前線宣教士.

$^{241}$我們沒有辦法出去.
但有人出去了.
我們都在背後.
以禱告,以關懷.
讓我們的宣教士感覺不被遺忘.
所以我們教會將會成立.
宣教關懷小組.
去關心在前線的宣教士.
為他們禱告.
當我們願意同心合意地一起來呼求神.
神就願意再一次拯救.
將我們.
我們都看到好似祂將若拿一樣.
從餘福裡面提升上來.
在海的深處再一次拯救祂上岸.
叫大魚吐祂出來.
不過吐出來之後.
不是完的.
有事要做的.
上帝再一次呼召祂.
「你去尼尼微聖宣講我要你講的信息」.
雖然上帝再一次呼召若拿.
再給祂一次機會.
祂有沒有去呢?.
祂去了.
我們熟悉的.
剛才也有讀過的.
祂最終去了.
祂人是去了.
不過祂似乎心沒有去.
或者我們會問.
其實若拿生什麼氣呢?.
祂為什麼這麼生氣呢?.
《聖經》記載.
當神察看尼尼微聖的人的行為.
看到他們離開惡道之後.
神就改變心意.
原先說了他們不悔綱.
就要降禍給他們.
那個災禍不降了.

$^{281}$現在上帝決定不降.
若拿大大不悅.
發怒.
他生氣到怎樣呢?.
跟神說.
「你取我的性命吧」.
「我死了比活著還好」.
其實真的很生氣.
真的很生氣.
你有沒有試過?.
你有沒有試過生氣到想死?.
我不知道你有沒有試過.
真的生氣到想死.
生氣到覺得死了比活著還好.
真的很生氣.
他覺得世上不應該有公義嗎?.
尼尼微聖的人要抵擋上帝.
為什麼可以這樣接受福音.
願意悔改就沒有懲罰.
公義在哪裡?.
可能我們也有試過.
見到世上很多不公不義的事.
不要說世界那麼遙遠.
每一天在我們身邊.
也有很多不公不義的事.
我們有時候會不會真的很想.
惡人應該要有惡報.
萬一他臨死之前做了那麼多壞事.
他臨死之前信耶穌.
那就真的糟糕了.
他得赦免怎麼辦呢?.
跟你有什麼關係呢?.
有關係.
因為將來我要在天下見到他.
我真的不想在天下見到.
一些這樣一輩子在害人的人.
他臨死前悔改了.
我真的不想見到他.
公義在哪裡呢?.
所以我們也能夠理解.

$^{321}$為什麼若拿那麼生氣.
若拿深知道.
若拿深知道上帝是一個怎樣的上帝.
他跟上帝說.
我知你的.
你就是這樣.
我仍在本國的時候已經說過.
我知道你是有恩惠有憐憫的神.
你不輕易發怒.
你有豐盛的慈愛.
你會改變心意.
你會不降災.
我就是因為這樣.
我就逃去他私.
我知道你是這樣的神.
所以若拿這一段他很熟悉,他是先知.
他引用《出埃及記》的經文.
是的,的確.
《出埃及記》記述過神的慈愛.
但其實神也有公義的原則.
《出埃及記》說過.
我萬不以有罪為無罪.
甚至必追討他的罪.
至父及子,直到三,四代.
若拿很想說的是.
你不是公義的神嗎?.
為什麼你會後悔不降災?.
大大嬲,嬲得要死.
若拿本來嬲的是什麼人呢?.
是利利未成的人,是不是?.
但現在他嬲的是誰呢?.
上帝.
他到頭來竟然.
由嬲人變成嬲神.
當你向上帝發怒的時候.
你有沒有留心.
你和上帝的關係就受影響了.
你和上帝的關係開始出問題了.
其實若拿的憤怒.
其實我們也不是不能理解.

$^{361}$或者我們明白多一點.
當一個人發出求死的信號的時候.
其實他在呼求什麼呢?.
我不知道你在生命裡有沒有遇過.
有人跟你說.
「我想死」.
「我現在就要去死」.
「我要想死」.
我不知道當你聽到身邊有人.
這樣跟你說的時候.
你聽到的是什麼.
黃家良博士有一本書.
他透過約拿書有一個新的息息.
是一個新讀.
其實若拿所承受的是什麼呢?.
他說約拿書裡面一再提到.
大,大,大,大.
大風,大雨,大水.
其實我們讀的時候.
可能沒有留意這些「大」字.
但是這些「大」是顯示著.
若拿覺得他生命面對著.
大到他不能面對的事情.
而第四章 若拿的大大不悅.
他有一棵菲麻樹遮著他.
他又大大歡喜.
這些情緒的大上大落是不尋常的.
黃博士提出.
他說其實若拿是一個.
正在面對他無法面對,承受的困境.
無法面對的情緒裡面.
他發出尋死的訊號.
也提醒我們.
當如果我們身邊真的有人.
受到極大的創傷.
或者我們外人.
看他有一些難以理解的表現的時候.
甚至乎覺得他情緒這麼誇張.
用不著這麼誇張.
其實我們看不看到.

$^{401}$有些人落在夾縫裡面.
落在絕望裡面.
甚至乎他有混亂.
他有思緒混亂.
他有用言語表達不到的感受.
如果我們多一點同理.
多一點明白.
我們會不會讓對方得意志呢?.
不過,另一邊廂.
也不是說那些發脾氣的人.
全對全對.
確實.
當別人做了一些.
令我們很憤怒,很生氣的事的時候.
我們特別覺得.
當道理在我自己那一邊的時候.
我們會覺得我們發怒.
都是合情合理的.
上帝也沒有說不可以發怒.
不過.
不要一直讓憤怒佔據你的思想.
我曾經聽過一位姐妹這樣分享.
她說可以生氣.
不過不要越想越生氣.
聖經也教我們.
生氣卻不要犯罪.
不可含怒度日落.
所以不是說我們不可以生氣.
甚至遇到不公義的事.
我們也要沒有感覺嗎?.
不是.
我們可以有情緒.
我們也可以發怒.
但不要讓這個思想佔據了你.
甚至成為你的主人.
你成為了他的奴隸.
你不要被憤怒主宰了你.
當我們遇見不公平,不公義的事.
甚至我們有憤怒的情緒.
可以說是非常正常,非常合理.

$^{441}$但關鍵是不要容忍這個情緒.
一直侵佔你的思想.
甚至操控你的行為.
靈卿姐妹.
憤怒的時候.
你回來聆聽神.
其實上帝是很喜歡和我們對話的.
如果我們仔細看.
其實約拿書的第四章.
剛才主席一直帶我們讀.
你有留意嗎?.
整章幾乎都是對話.
是他和上帝你一句,我一句,你一句,我一句.
上帝是會對我們說話的.
當你憤怒,不公義.
過不了自己.
你來和上帝對話吧.
就好像約拿那樣.
他不斷和上帝對話.
以至上帝足之讓他在對話裡明白.
上帝的心腸.
上帝要我們明白他的心腸.
如果我們讀約拿書.
我們不得不面對的一個課題.
就是公義與憐憫的課題.
這是我們每一次讀約拿書.
都不能逃避的課題.
我們一定會問.
公義和憐憫如何並存?.
公義和憐憫哪邊應該多些?.
約拿書一開始讓我們看到一個出走的先知.
他逃避上帝可能有很多原因.
我相信他最主要的原因就是抵不住尼尼尾城的人犯罪.
他的罪惡不應該得救.
也有一個可能.
因為尼尼尾城是亞述國.
是以色列的敵國.
他也可能因為愛國之情.
不盼望敵人可以得救.
當時他很害怕.

$^{481}$害怕他宣講成功.
他不是害怕宣講失敗.
你留心.
其實所有宣教士,傳道人,傳科.
當然是渴望宣講成功.
很怕別人悔改.
信主我不得不坐在這裡.
都舉手決志相信耶穌.
悔改 應該是這樣的.
但約拿剛剛相反.
最好他們全部都不要悔改.
因為他很憎恨他們.
很憤怒他們.
逃跑是因為他很害怕.
宣講會成功.
他不覺得這個城應該值得被赦免.
其實約拿去尼尼尾城.
他最終雖然去了.
但他心裡不舒服.
約拿其實是一位很公義的先知.
你想清楚一點.
我們雖然說今天的題目是出走的先知.
但其實他是很公義的先知.
他心裡有一套.
很公義的原則在裡面.
非常公義.
但他不明白恩典的道理.
他沒有明白恩典的道理.
正如我剛才所說.
他告訴耶和華.
你選擇去尼尼尾城.
我不去.
是因為我知道你會後悔.
你會不降災.
約拿在裡面.
他覺得一個人.
甚至這個世界.
做得好的人就應該有獎賞.
做得不好.
就應該要罰.

$^{521}$好就應該有好報.
惡就應該有惡報.
當船翻起風浪的時候.
船員問誰錯.
約拿沒有說.
沒有說其他.
他只說你們抬我下海.
將我拋下海.
我就平靜了.
我知道你們遭遇大風是因為我.
他不是對人這樣.
他要求公義在人身上.
他要求在自己身上也一樣.
他對自己有同樣要求.
是我引起的.
是我連累你們.
所以你看到他裡面.
有一套很公義的準則.
但是.
當我們去想.
約拿這樣敵我分明.
他是不是敵我分明.
他很公平.
很公正.
很公義.
不對嗎.
我們不是應該這樣做人嗎.
甚至我們做基督徒.
不是更加應該這樣嗎.
更加應該要敵我分明.
更加應該要公平.
公正公義.
不是這樣嗎.
紐約時報有一個暢銷作家.
提摩泰.可納.
我很喜歡看他的書.
他是一位神學家.
他出了一本書叫《出走的先知》.
他裡面說到.
公平,公正,公義.

$^{561}$不是不對.
不過.
不能夠高過上帝的旨意.
好像約拿.
他將自己自以為對的.
擺高過上帝的旨意的時候.
他就跌入了一個.
自以為是的陷阱裡.
提摩泰.可納提到.
其實在舊約.
有約拿,這位出走的先知.
新約.
其實也有一個人出走了.
不知你想不想到呢.
是浪子的比喻.
浪子的比喻裡.
也有一個出走的兒子.
如果我們記得這個故事.
這個比喻.
耶穌說的比喻.
一個爸爸有兩個兒子.
有大兒子和小兒子.
小兒子很不長進.
在爸爸還在生的時候就要分家產.
分爸爸的家產.
離家出走.
拿著爸爸的家產揮霍無度.
放蕩生活.
生窮水盡的時候就回來找爸爸.
醒不過來.
回來爸爸又接納他.
還宰了一隻肥牛.
跟他慶祝一番.
叫親朋戚友來.
原本應該是一件很歡喜快樂的事.
我們讀這個比喻的時候.
浪子回頭.
根不換.
不過我們知道聖經裡形容.
有一個人很不高興.

$^{601}$那個是大兒子.
大兒子很生氣.
大兒子生氣.
為什麼你會原諒他.
你為什麼不追究.
甚至大兒子說.
我天天在這裡服侍你.
我天天來服侍你.
你都沒有為我慶祝過.
你沒有宰過一隻羊給我.
叫我朋友來和我一起吃飯.
你沒有這樣對我.
大兒子很生氣.
作者提摩泰海納.
他提到.
其實.
出走的基督徒.
有兩種.
一種好像小孩一樣.
真的出走了.
但還有一種基督徒.
好像大兒子一樣.
人好像沒有出走.
但他的心.
出走了.
表面看來.
他的人還在.
但他的心.
沒有辦法體會.
他父親的心腸.
沒有辦法.
明白.
父親的愛和憐憫.
他內心有一套.
很強烈的準則.
他感到心心不忿.
不公平.
他開始不是生氣弟弟.
他最終是生氣爸爸.
生氣這個父親.

$^{641}$他沒辦法體會上帝的心腸.
沒辦法,他心懷不平.
沒有慈悲憐憫.
其實.
他同樣是出走了.
若拿這一刻都一樣.
好像大兒子一樣.
他的人是去了尼尼眉城.
你要我做的事.
我都做了.
好像表面上很信服上帝的吩咐.
是不是?.
但他心不服.
他沒辦法體會上帝的慈愛.
理解不到.
上帝恩典的道理.
在這本書裡作者有一句話.
當我們變得非常虔誠.
和有道德的時候.
都一樣可以.
用來逃離上帝.
當我們很多自以為是的時候.
我以為我們站在上帝那邊.
我們為上帝.
幫上帝審判這個世界.
但其實可能令你.
非常遠離上帝.
我們有時候.
總是執著我們所信的.
我們所堅持的很多道理.
我們可能都會這樣.
也許.
我們每個星期日都回來教會.
我們都乖乖坐在這裡.
我們好像大兒子一樣.
沒有離開,是不是?.
但我們心裡面.
可能有很多的審判.
很多的心懷不平.
很多的分黨分派.

$^{681}$我們只是沒有做出面.
我們將人.
在我裡面敵我分明.
你不是我的黨.
我們可能在家庭裡面.
都是這樣.
與我們的丈夫,與我們的太太對立.
與我們的孩子對立.
我們在職場上.
都可能是這樣.
敵我分明.
公平公正.
不是不對.
不過似乎約拿書重新幫助我們.
重新修正.
我們的觀念.
當我們恃著.
這些公平公正.
拿著很多準則的時候.
我們不單止可能破壞了.
與人的關係.
要記得.
你是破壞了你與父母的關係.
如果今天.
我們抓住這些.
我們心目中.
覺得最重要的.
去審判我們的弟兄姊妹.
審判外面的人.
最終可能.
是你與上帝的關係.
受到障礙.
我們不可以.
一方面求上帝以恩典待我.
但一方面卻用.
公義去審判人.
當我們讀下去的時候.
其實我們會發現很奇妙的.
其實上帝有沒有.
責怪約拿呢?.

$^{721}$他看不到上帝的恩典.
他不體會上帝的心情.
上帝有沒有罵他呢?.
「臭小子,醒來吧!」.
「你不體會我的心情?」.
沒有.
上帝問他.
「你這樣.
發怒,對不對?」.
上帝於是用一棵樹.
一夜長大.
一夜枯死.
來開導他.
上帝安排一棵樹.
一夜之間上帝來幫約拿遮太陽.
令他很舒服.
他就大大喜悅了.
但上帝一晚又用一條蟲子咬死他.
白天出來就曬到約拿死死.
他又生氣了.
當他一路生氣.
到要死的時候.
一切跟上帝說.
上帝.
仍然問他.
「你因這棵披麻樹.
這麼生氣,對不對?」.
但他很可愛.
他繼續回答上帝.
「我發怒以致死也合理.
我也是對的」.
其實約拿真的很可愛.
我看他.
難免有一點我們自己的影子.
於是上帝問他.
上帝問他.
「這棵披麻樹你沒有操勞過.
你沒有叫它生長過.
它一夜生長.
一夜枯死.

$^{761}$你也愛惜它.
何況利利尚尚.
其中不能分辨左右手的.
有十二萬多人.
不能分辨左右手的意思是什麼?.
是鷹蟹.
沒有辦法.
最無辜,最弱小的那些.
連左右手也不懂分辨.
有十二萬多.
還有很多生育.
上帝問約拿.
「我豈能不愛惜呢?」.
約拿書很特別.
停在一個.
問題上.
如果你回去看約拿書.
四章十一節就是全卷書最後的一句.
這麼有趣的一個話劇.
上帝沒有.
向他大大發怒.
上帝沒有激打他.
上帝是在教導約拿.
你在公義,懲罰,審判之外.
你看不看到.
還有一樣東西.
叫做「慈愛與憐憫」.
還有很多無辜的.
是很值得珍惜的.
你看不看到.
上帝沒有說約拿這些執著是錯誤的.
他沒有.
他沒有說你這樣是錯的.
上帝只是想讓約拿知道.
在守住公義,公正,公平.
審判的同時.
不能夠.
沒有「慈愛與憐憫」.
上帝是要約拿去體會.
用一棵樹.

$^{801}$要他體會我的心情.
今天一樣.
上帝要我們去憐憫.
向世人.
顯慈愛的同時.
他沒有要我們放棄公義的準則.
上帝沒有要我們放棄.
雖然世界上有很多東西.
我們沒有辦法改變.
但上帝仍然.
將公義放在我們心裡.
否則世界各地.
不需要有那麼多戰爭歷史紀念博物館.
如果我們要忘記.
上帝沒有要我們忘記.
我們就不能夠讓這個不忘記.
不能夠讓這個.
高過上帝的旨意.
我們不是要當什麼事都沒有發生過.
我們不是要不去紀念這些不公義的事.
但我們怎樣去紀念.
不忘記的同時.
我們仍然持續.
為這些不公義的事守望.
仍然堅持繼續仰望上帝的公義.
終有一天會降臨.
這是我們要做的事.
約拿書是讓我們.
很真實地去思考.
很真實地去思考公義和憐憫的課題.
可能你看的時候覺得是一個故事.
但其實不是.
它是說很深奧的道理.
剛才我說.
約拿書很特別.
它停在一個開放式結尾的問題裡.
一個開放式結尾.
聖經最後是沒有交代的.
你們猜若拿最後有沒有悔改.
有沒有明白上帝的深層.

$^{841}$你們猜.
有吧.
上帝沒有.
最後他沒有寫.
有一句有第十二節.
若拿就明白上帝的深層了.
沒有這一節.
我們怎樣知道.
若拿最後是否明白.
如果我們認同.
約拿書.
是若拿所寫的.
若拿能夠寫得出.
你猜他明白了沒有.
他領悟到的是什麼.
他寫下了在這本書裡.
約拿書其實是一本很獨特的書.
它當中記載了.
神不只是愛意識的人.
神竟然愛.
外邦人.
愛犯罪的人.
愛不是信奉他,敬拜他的人.
是不是.
神愛世人.
神雖然揀選以色列人.
成為他的子民.
但是神.
都會用外邦人審判以色列人.
讓他們亡國.
同樣.
當外邦人不認識耶和華神的時候.
神都可以對他們施憐憫.
施慈愛.
讓他們有悔改的機會.
我們有時會不會.
將上帝.
限制在我們的框框裡.
覺得上帝.
應該要怎樣做.

$^{881}$我們會不會都覺得.
上帝是在我們的想像裡.
上帝應該.
做我們覺得應該要他.
這樣做的神.
有時是我們限制了我們的上帝.
雖然其實你限制不了.
我們的上帝.
他是全地的神.
他是有辦法.
他是有主權.
去拯救任何人.
所以約拿書最後留下的問題.
我豈能不愛惜嗎.
仍然是對我們.
這一代的信徒.
對今天採到的你和我.
發出的呼喚.
我們能不能體會上帝.
他豈能不愛惜嗎.
我們願不願意放下.
我們心目中的框框.
我們所謂的界線.
我們要承認上帝的道路高於我們的道路.
上帝的旨意高於我們的旨意.
我們願不願意放下我們固有的看法.
願意幫助我們.
再一次看到.
我們人沒有遠離.
但我們會不會很多時候.
心已經遠離了.
請姐妹.
當你去尋思.
去思想的時候.
真的覺得.
原來有時候我已經遠離上帝了.
我盼望你記得.
無論在任何地方.
任何時候.
哪怕在餘福中.

$^{921}$絕境裡.
只要你一轉念.
上帝仍然與你親近.
最後我要用一個小小的語道故事.
來結束我今天的講道.
這個故事是說.
你的杯子裡面裝著什麼呢.
What's in your cup.
試想像一下.
現在你手上正拿著.
一杯很熱的咖啡.
有人過來撞到你.
搖晃到你.
令到你將咖啡倒到周圍.
如果有人走過來問你.
為什麼你會倒滿咖啡呢.
你會怎樣回答.
他撞到我.
因為有人撞到我.
是,確實是.
是因為有人撞到你.
所以你倒到周圍都是咖啡.
因為你的杯子裡面.
裝了的是咖啡.
如果你杯子裡面裝的是茶.
倒出來的就是茶.
你杯子裡面裝的是清水.
倒出來的就是水.
其實在我們的生活裡.
每一天.
每時每刻.
每一個環境.
都會撞到我們.
你是沒有辦法.
避免有人撞到你.
你沒有辦法阻止環境.
衝擊你.
撞你.
甚至乎.
撞你的是我們的上帝.

$^{961}$是我們的上帝.
要衝擊你.
當生活的搖晃.
環境的衝擊.
來到我們的時候.
我們寫出來的是什麼呢.
如果.
我們裡面裝載的.
全部都是憤怒.
不平.
如果我們裡面裝載的.
全部都是惡毒.
我們寫出來的就會是這些.
但如果.
我們裡面裝載的都是神的話.
裝載的是感恩.
喜悅.
和平.
謙遜,忍耐.
你寫出來的就會是這些.
衝擊你的時候生活很平靜.
所有事情都可以很美好.
但往往就是有人撞你的時候.
你才發現原來你裡面裝載著什麼.
我盼望.
我們所裝載的.
都是神的道.
神的話.
我們裝載的都是神的屬性.
有人說.
生命.
生存本來給了你一個杯.
但裝載什麼在裡面呢.
是你選擇的.
是你選擇用什麼去填滿你的生命.
而且有一點我認為非常重要的就是.
當你裝載著熱的咖啡的時候.
寫出來的時候.
不單止會傷到人.
其實最傷那個.

$^{1001}$可能是你自己.
往往錄到的那個.
是你自己.
確保你的杯裡面.
裝載的是那些.
不會燙傷別人.
也不會燙傷你自己的東西.
我們常說.
願世界變得更美好.
其實是從我們開始.
是你決定.
你的杯裡面裝載什麼.
是不是有東西撞到你的時候.
一定要寫出燙人.
和燙到自己的東西呢.
願主的話.
願主的道.
願主的公義.
願主的憐憫與愛.
都在我們裡面.
讓主填滿我們.
願我們願意用感恩.
寬恕.
喜悅.
去對待我們身邊的人.
用良善.
溫柔,恩慈,忍耐.
去輔助我們身邊的人.
好不好,我們一起禱告.
親愛的天父上帝.
多謝你.
你從來沒有嫌棄我們.
哪怕我們真的出走了.
還是說.
我們心已經遠離你.
多謝主你仍然讓我們.
可以在絕境裡.
去尋求你.
在憤怒裡我們仍然可以.
與你對話.

$^{1041}$多謝你.
教導我們公義與憐憫的課題.
願主你幫助我們.
讓我們在這個世代裡.
滿滿裝載.
是你的聖情.
你的話語.
滿滿裝載的.
是你的教導.
直到所有的衝擊來臨的時候.
我們所寫出來的.
都是主的恩.
主的愛.
與主的道.
願我們能夠彼此.
相愛.
讓我們在教會中.
彼此以恩慈相待.
讓我們多一點.
多一點憐憫.
多一點體恤.
我們雖然心裡不是.
沒有公義,沒有底線,沒有規矩.
不是.
但是主啊.
當這些要出來的時候.
我們願你的旨意.
高過我們的旨意.
我們願你親自來教導我們.
我們怎樣秉行公義.
怎樣行出真理.
願求主幫助世代.
幫助這個世代所有的信徒.
幫助我們教會.
以耶穌基督名字祈求.
\newpage



\section{歌羅西書 3:12-21}
\label{sec:032122}
\textbf{最好的禮物}~~~~林瑞貞傳道
\newline
\newline
連結: \href{http://www.church.com.hk/acms/content.asp?site=cdc\&op=show\&type=product\&code=032122\&layout=sermon}{\texttt{032122}} ~~~~ 語音日期: 2026-05-10
\newline
\newline
\hyperref[sec:032121]{\small{< < < PREV SERMON < < <}}
~
\hyperref[sec:index_chronic]{\small{[返順時目]}}
~
\hyperref[sec:032123]{\small{> > > NEXT SERMON > > >}}
\newline
\newline
歌羅西書 3:12-21
\newline
\begin{longtable}{cl}
\hline
\hline
章節 & 經文 (和合本修訂版)\\
\hline
3:12 & \begin{tabularx}{0.7\textwidth}{X} 所以，你們既是神的選民，聖潔、蒙愛的人，要穿上憐憫、恩慈、謙虛、溫柔和忍耐。 \end{tabularx} \\ \\ \relax
3:13 & \begin{tabularx}{0.7\textwidth}{X} 倘若這人與那人有嫌隙，總要彼此容忍，彼此饒恕；主怎樣饒恕了你們，你們也要怎樣饒恕人。 \end{tabularx} \\ \\ \relax
3:14 & \begin{tabularx}{0.7\textwidth}{X} 除此以外，還要穿上愛心，因為愛是貫通全德的。 \end{tabularx} \\ \\ \relax
3:15 & \begin{tabularx}{0.7\textwidth}{X} 你們要讓基督所賜的和平在你們心裡作主，也為此蒙召，歸為一體。你們還要存感謝的心。 \end{tabularx} \\ \\ \relax
3:16 & \begin{tabularx}{0.7\textwidth}{X} 當用各樣的智慧，把基督的道豐豐富富的存在心裡，用詩篇、讚美詩、靈歌，彼此教導，互相勸戒，以感恩的心歌頌神。 \end{tabularx} \\ \\ \relax
3:17 & \begin{tabularx}{0.7\textwidth}{X} 你們無論做甚麼，或說話或行事，都要奉主耶穌的名，藉著他感謝父神。 \end{tabularx} \\ \\ \relax
3:18 & \begin{tabularx}{0.7\textwidth}{X} 你們作妻子的，要順服自己的丈夫，這在主裡面是合宜的。 \end{tabularx} \\ \\ \relax
3:19 & \begin{tabularx}{0.7\textwidth}{X} 你們作丈夫的，要愛你們的妻子，不可虐待她們。 \end{tabularx} \\ \\ \relax
3:20 & \begin{tabularx}{0.7\textwidth}{X} 你們作兒女的，要凡事聽從父母，因為這是主所喜悅的。 \end{tabularx} \\ \\ \relax
3:21 & \begin{tabularx}{0.7\textwidth}{X} 你們作父親的，不要惹兒女生氣，恐怕他們會灰心。 \end{tabularx} \\ \\
[1ex]
\hline
\hline
\end{longtable}
$^{1}$各位兄弟姊妹,平安!.
今日是母親節,首先慶祝各位母親,母親節蒙印!.
今日是母親節,不知大家會怎樣慶祝?.
會不會一會兒完了崇拜之後,一齊去飲茶?.
跟你的媽媽,家人一齊飲茶?.
還是今晚會吃一餐大餐呢?.
又或者買一份媽媽很喜歡的神秘禮物?.
不知會不會是這樣呢?.
其實買禮物都真的很講深思的,.
我經常很羨慕那些人買禮物買得很好,.
我就沒有這一方面的專才,.
我又很擔心自己買了一份禮物,.
別人是不合用的,.
即是拿著我送給他的禮物,.
他都不知道怎樣處理好,.
所以我都真的很怕買禮物,.
但是在這些特別的日子上,.
我們又真的需要送禮物去慶祝,.
不知大家有沒有聽到有些說話,.
買不買禮物不重要,.
同不同我食飯不重要,.
你同我食飯買禮物給我,.
不如平日對我好些,.
不如平日陪多些我,.
不知有沒有聽過這樣的說話呢?.
當我們平日能夠好好善待我們的家人,.
好好跟他們相處,.
比起跟他們吃一餐大餐,.
送一份他們喜歡的禮物反而更加重要,.
將我們自己成為一份禮物,.
送給我們的家人,.
可能對他們來說,.
他們會更加開心..
今天我和大家一起去看.
《歌羅西書》三章十二至二十一節,.
希望我們透過那段經文,.
提醒我們怎樣成為一個好的家人,.
作為一份禮物,.
送給我們的家人,.
不單是母親節,父親節,生日,.

$^{41}$而是每一天讓我們的家人.
都能夠收到這份最好的禮物..
首先我們去看.
我們去看《歌羅西書》三章十八至二十一節,.
三章十八至二十一節是關於家庭角色的本份,.
這四節經文裡面我們看到.
有兩種家庭基本的關係,.
第一種就是夫妻關係,.
在聖經裡面人倫是由夫妻開始的,.
亞當和夏娃,.
第二種關係就是父母和子女,.
我相信在座當中有些是身為別人的丈夫,.
有些是身為別人的妻子,.
有些是身為別人的父母,.
但可能有些三種都不是,.
沒有這三種身份,.
但有一種身份我們每一個人都有,.
就是別人的子女,.
我們每一個人都是別人的子女,.
這四節經文提到基督徒在他們家庭角色應有的本份,.
首先保羅提醒妻子應該是怎樣的,.
經文提到妻子要信服丈夫,.
只要丈夫的心意是合乎神的心意的話,.
妻子就應當經常信服丈夫,.
因為這是基督徒最恰當的態度,.
和信仰的身份相稱,.
可能你們有些疑問,.
現在的社會裡,.
女性的教育程度,工作能力,社會地位都不斷提升,.
妻子的收入,職位,能力都不比丈夫低,.
既然妻子已經不是活在一個父權的社會中,.
不是一定要依靠丈夫,.
有人就會疑問,.
為什麼妻子還要信服丈夫呢?.
我都不一定要靠丈夫的,.
這個吩咐似乎已經不合時代了,.
不知道你們有沒有這個想法呢?.
首先我們要認清一件事,.
就是我們的才華,能力都是神所賜的,.
沒有什麼可以誇,.

$^{81}$另外神吩咐妻子要信服丈夫,.
不是因為妻子的能力比丈夫低,.
不是能力的問題,.
也不是男比女優勝,.
也不是想貶低妻子的地位或尊嚴,.
而是因為丈夫是妻子的頭,.
好像基督是教會的頭,.
基督又是教會全體的救主,.
教會怎樣信服基督,.
妻子又要照樣凡事信服丈夫,.
歌羅西書二不鎖書五章二十四節這樣教導妻子,.
當妻子去信服丈夫,.
就是認同神設立的秩序,.
妻子信服丈夫,.
亦是代表著信服神的主權,.
然後神不偏心,.
祂不只是吩咐妻子要怎樣做,.
祂亦吩咐丈夫要愛妻子,.
不可以當妻子是傭人,.
妻子不是讓你娶回來做家務服侍你,.
不是一個性伴侶,.
或者可以隨意拋棄的人,.
丈夫愛妻子,.
不只是情感上的喜愛,.
不只是男女之間的喜愛,.
而是要真心顧念對方的好處的愛,.
就好像愛自己身體一樣,.
我們有沒有人不愛自己?.
神吩咐作為丈夫要愛妻子,.
好像愛自己一樣,.
亦吩咐丈夫不可以苦待妻子,.
意思是丈夫不可以帶著任何苦毒怨恨對待妻子,.
作為妻子要主動信服丈夫,.
丈夫無權吩咐妻子信服他,.
而丈夫亦要主動愛妻子,.
丈夫無權吩咐或命令妻子愛他,.
丈夫要自己主動愛妻子..
講完妻子丈夫應有的本分之後,.
寶來又吩咐作兒女和作父母的應該是怎樣?.
經文講到作兒女的要凡事聽從父母,.

$^{121}$兒女不只是小朋友的兒女,.
即不是小小的中學生小學生的兒女,.
就算是成年人,像我們這麼大年紀,.
都是要信服,都是要聽從父母,.
而且是凡事,不是你選你喜歡的就聽,.
不是你有空的時候才聽,.
是凡事都要聽從父母,.
當然凡事的意思不是父母的吩咐,.
父母的吩咐,如果是違背神都要聽,.
父母的吩咐是不可以違背神的,.
兒女凡事聽從父母,這是主所喜悅的..
在不同的聚會裡,.
我自己在教會除了負責成人群體,.
我自己也有負責兒童的,.
也會和家長接觸,.
有時候也會聽到家長說,.
現在的兒女真的很難教,無論大小也好,.
手永遠都拿著手機,機不離手,.
吃頓飯也拿著手機,.
想聊聊天也不行,.
兒女就坐著看著手機,.
小時候也乖乖聽話,.
到差不多升中,進入中學的時候,.
聽課後很想自主,很想獨立,.
不想被父母太多的管束,.
當父母有吩咐時,他們會怎樣呢?.
頂嘴,一言九鼎..
所以對父母來說,.
如果兒女能夠凡事聽從他們的話,.
真是一件非常好的禮物..
神又吩咐作父母的應該是怎樣?.
作父親的,不要激怒兒女,免得他們灰心喪志..
雖然這裡只是說父親,.
但當然同樣包括母親..
這裡的意思不是說不要做一些事,.
令到兒女厭煩反感,.
但不是說我們不可以去管教兒女,.
不是說我們不可以去責備兒女,.
只不過不要做一些事令他們灰心喪志,.
例如不斷去挑剔,不斷去說他們,.

$^{161}$令到兒女覺得自己沒價值,沒用,.
甚至會氣餒,不想再去嘗試,.
令到自己沒心機..
剛才我也說過,我是做小朋友工作的,.
就算他們升了中,.
我都會在長假期裡,約他們出來吃飯,聊天..
有時他們也會分享家裡的事給我聽,.
分享他們學習的困難,同學之間的關係,.
有時也有校園欺凌,杯葛等等..
我問他們,你們有沒有將這些事告訴父母,.
讓父母幫你?他們很快說,沒有..
告訴他們是沒用的,他們也不理,.
他們也不聽,他們不相信我,.
只說我不對,只說我不夠努力,.
但其實我已經很努力..
有個朋友跟我分享,.
他說在學校有些同學對他不太好,.
他回去跟父母說,.
父母說,你不要理他們就行了,.
為什麼叫你經常黏在一起?.
是不是自己拿來不好?.
所以這個年輕人就很不開心,.
明明是我不開心,明明是他們杯葛我,.
為什麼現在好像說是我不對?.
所以他不想跟父母分享一些不開心的事..
有個五年級的女孩子也跟我分享,.
她不只是跟我分享在主學的時候說,.
她不是一次,是經常在不同時候主學會說,.
突然又說,成績不夠好,.
但其實她的中文是全級第一,.
她知道數學沒有人那麼好,.
但沒有人那麼好不代表不合格,.
合格的,即是沒有人那麼好..
然後又說自己畫畫不夠漂亮,.
她畫畫不是差的,.
只不過她家裡的姐姐特別有藝術天分,.
畫畫很漂亮..
然後沒多久,隔兩星期她又說,.
我彈琴不夠好,其實她也八級了,.
小學,八級,可能她跟其他人比是演奏級..

$^{201}$然後又說自己身材不夠高,.
她也不矮,不過她不是模特兒的高度..
然後又說自己畫畫不夠漂亮,.
她畫畫也不差,只不過她畫得不夠漂亮..
然後我真的忍不住了,.
隔幾個星期她又說自己不好,.
這樣不夠好,那樣不夠好..
然後我就跟她說,.
你是神所做的,你是獨特的,.
你有你的好處的,你有你的強項的,.
你不是垃圾,神不做垃圾的,.
你有你的強項的,看到嗎?.
有時候,我們很多時候都會看到.
我們的子女不好的地方,.
很容易看到的,.
當我們看到子女不好的地方時,.
我們就會去提點..
我們有好的意圖,.
我們是有好的意圖,.
希望他們更好,更出色,.
但是有時候我們這個好的意圖,.
反而影響了他們的成長,.
影響了他們的自信心,.
影響了他們的自我形象,.
而這個影響不只是他們小時候,.
這個影響會帶到他們成人,.
影響了他們跟其他人的關係..
妻子要主動信服丈夫,.
丈夫要主動愛妻子,.
兒女要在各方面聽從父母,.
父母不要惹怒兒女,.
這些都是保羅教導,.
信徒在家庭角色的本分,.
但是要盡上這些本分,.
一點都不容易,.
剛才我講得很快,.
講就容易,.
但是要盡上這些本分,.
其實一點都不容易,.
是需要有美好的性情,.

$^{241}$和生命來配合..
歌羅西書只有四章經文,.
直接講到家庭關係的,.
只是這四節,.
而且是放在這本書較為後的部分,.
在第三章後半節,.
保羅不是用很多經文來講家庭角色本分,.
反而他用了較多的經文來講信徒的生命,.
家庭是信仰生活的一部分,.
信徒的生命會影響他跟家人的相處,.
影響他的家庭生活,.
信徒和主聯合在身份上和生命上會有改變,.
而這些改變也同時影響了跟家人的關係,.
所以保羅在歌羅西書很強調「棄絕舊人」,.
歌羅西書三章一至二節這樣講,.
「所以你們既然與基督一同復活,.
就應當尋求天上的事,.
哪裡有基督坐在神的右邊,.
你們要思念的是天上的事,.
不是地上的事.」.
當我們與耶穌基督聯合,.
我們相信了祂,.
我們要思念的是天上的事,.
要求的也是天上的事,.
代表著我們的生命是願意更生改變,.
在思想和行為之上也是願意信服神的旨意..
保羅在第三章命令信徒要棄絕舊人的行為,.
在三章五至八節這樣講,.
「所以要致死你們在地上的肢體,.
就如淫亂污穢邪情,.
惡辱和貪心,.
貪心就是敗牛將.」.
因為當我們貪一樣東西,.
可能每個人都不同,.
可能是貪權,貪錢,貪靈,.
這些我們所貪的東西,.
就取代了神在我們心目中的地位,.
就如牛將一樣..
在第八節亦講,.
「但現在你們要除去憤怒,.

$^{281}$惱怒,惡毒,毀謗,.
以及粗言穢語這一切事..
棄絕這一切的事.」.
第八節說要棄絕這一切,.
即是說原先是我們有的,.
現在我們不要了..
這些舊的行為,.
不單止不能幫助我們成為一個好的家人,.
反而會破壞我們和家人的關係..
保羅提醒我們,.
不單止要脫去舊人,.
同時要穿上新人..
保羅在第十二節提醒我們,.
我們的身份,.
根據這個身份我們應該有所改變..
我們既是神所揀選,.
是聖潔蒙愛的人,.
這就是我們的身份,.
我們被分別為聖出來,.
神的子民..
我們有這個身份的話,.
我們的言語行為就要和我們這個身份相稱..
我們有這個身份應該有什麼行為?.
他繼續說,.
我們要有憐憫的心腸,.
恩慈,謙卑,溫柔和忍耐..
因為我們是蒙揀選,.
我們有這些美德和人相處..
我們的身份,.
這個身份不是我們自己拿回來的,.
是神給我們才有的..
所以我們的生活行為,.
是要有這些美德去配合神,.
給我們這個身份..
我教會就在九龍城,.
有一次我走去九龍城廣場的時候,.
剛剛就去到過馬路前面,.
就是紅公仔,.
他自然就站在那裡等,.
在我旁邊有其他人在等過馬路,.

$^{321}$其中有兩個是穿著名校校服的學生,.
其中一個同學就和他旁邊的同學說,.
衝呀!意思是什麼?.
不要看燈了,衝過馬路..
然後他的同學就說,.
你不要命呀,.
有些車在行駛,.
然後又等了一會兒,.
沒有那麼多車了,.
那個同學,.
即是說衝的那個,.
又說衝呀,.
沒有什麼車了..
然後第二個同學就說,.
校風呀,.
大家明白什麼叫校風吧?.
他穿著校服的,.
代表著他的學校,.
有他國家學校學生的身份,.
如果他亂衝燈,.
是名校來的,.
會影響他什麼?.
影響他的校譽..
當第二個同學說完校風之後,.
第一個同學就沒有再出聲說衝,.
乖乖地站在那裡等過馬路,.
等轉燈..
很特別吧?.
這個同學,.
沒有因為怕被車撞到,.
而不選衝燈,.
反而怕影響校風,.
而選擇乖乖地等轉綠燈..
一個人的身份,.
當一個人在意他的身份的時候,.
都會影響他的行為..
我們有沒有在意我們的身份?.
這個蒙簡選,.
正潔蒙愛的身份,.
是神分別出來成為子民的身份..

$^{361}$當我們在意這個身份的時候,.
我們的言語,行為,.
都要和這個身份相稱..
憐憫的心腸,.
因此謙卑,溫柔和忍耐..
還有的就是,.
如果有人對別人有遙規,.
總要彼此饒恕,.
互相饒恕,.
主怎樣饒恕了你們,.
也要饒恕別人..
這些美德,.
生命的特質,.
是可以促進和諧群體的生活,.
可以幫助我們,.
怎樣好好對待我們的家人..
讓我們一起去看看這些美德..
身為基督徒,.
我們都經歷過神給我們的憐憫和恩慈,.
我們要嘗試去學習,.
以恩慈來對待我們的家人..
恩慈和嚴厲是相對的..
我們很多時候,.
比較容易以恩慈來對待,.
不是我們的家人,.
其他人我們會有恩慈..
有時候,.
我兒子小的時候,.
都會招呼同學來我們家玩,.
當然有小食,.
如果我兒子的朋友,.
吃多了,.
吃了別人的那一份,.
你都不會很大責怪他們,.
只會跟他們說,.
「這是誰的,.
我們吃自己的,.
如果你想再要,.
你告訴我.」.
但如果我們的小朋友,.

$^{401}$去到別人的家,.
吃了不是他的東西,.
我們會怎樣呢?.
就會說,.
「你做什麼?.
你這麼沒規矩的,.
為什麼吃別人的東西?.
要問別人的嘛.」.
總是會這樣..
可能我們會說,.
「那是別人的兒子,.
不可以說的.」.
我們的兒子就不同了,.
可以說的..
都是本著這個心,.
都是好意圖來的,.
都是想教育自己的兒女,.
但有時候,.
是否要用這樣的語氣呢?.
是否要這麼嚴厲呢?.
可不可以用第二個語氣,.
來教導我們的兒女呢?.
當我們回想,.
神在我們身上,.
我們是蒙神的憐憫,.
恩待我們,.
我們要學習上帝的恩賜,.
來待我們的家人..
聖經有提到謙卑,.
謙卑是一種態度,.
是認清自己的才能和能力,.
是神所賜,.
所以沒有可誇的,.
不需要增大,.
不需要比較,.
反而是願意降低自己的地位,.
願意信服,.
有信服和服侍的態度..
而謙卑帶出溫柔,.
溫柔不是那些陰聲細氣,.

$^{441}$說話好像很輕聲,.
不是說這些的溫柔,.
這裡說的溫柔,.
是存著對別人的好處,.
為別人著想和需要,.
願意放棄自己的權利,.
不斤斤計較..
主耶穌就是一個很好的榜樣,.
祂放棄了自己的權利,.
為了我們的好處,.
為了我們的著想,.
走上十誡的路..
聖經說的忍耐,.
不是那種忍你,.
現在我先忍你,.
然後你就知道,.
不是這種忍耐,.
因為聖經說的忍耐,.
是明白別人的限制,.
明白對方需要時間和空間,.
來改變和成長,.
對對方有正面的態度和支持,.
這些忍耐最好在哪裡訓練,.
你知道嗎?.
就是小朋友,.
自己的朋友,.
我剛才說我有小朋友,.
我有兩個小朋友,.
我小的小朋友,.
兩個都是兒子,.
我小的有特殊需要,.
專注力不足過度活躍,.
這些小朋友有什麼特質呢?.
就是不專注,.
不專注力不足,.
不專注,.
而且很容易忘記東西,.
很善忘的東西,.
又不記得帶這個,.
又不記得帶那個,.

$^{481}$他們小時候,.
我已經開始教他們煮飯,.
他們懂得煮就不用煮了,.
他們很小的時候,.
我就教他們煮飯,.
我小兒子煮了兩次炒菜,.
我就放手,.
快回不到家的時候,.
媽媽快回來了,.
你現在開始炒菜,.
回來就可以吃飯了,.
我回到家,.
進門口換衣服的時候,.
怎麼聞到糯米,.
我就問我兒子,.
你炒焦菜嗎?.
他說不是,.
出來沒多久,.
他就跟我說,.
媽媽,.
我忘了放油,.
然後我就說,.
那你有沒有放水?.
他說有,.
那可以了,.
那可以吃嗎?.
可以吃,.
好不好吃,.
熟就可以了,.
那就開飯吃吧,.
我真的吃了一半,.
然後我就教他,.
要先放油,.
就算你不放油,.
你只是煮菜,.
就要放多點水,.
我沒有罵他,.
只是教他,.
因為我明白,.
他不是一次兩次就會懂,.

$^{521}$有一次,.
我給他一個食譜,.
就是糖醋骨,.
高難度挑戰,.
但之前已經教過,.
我也煮過給他看,.
煮過,.
然後回到家,.
回到家的廚房,.
看一看,.
有一個湯煲,.
我剛才已經心知不妙,.
不過沒有糯米,.
吃飯的時候,.
打開來看,.
你猜我看到什麼,.
蚊骨,.
是糖醋排骨湯,.
我就問他,.
姐姐你為什麼用這個煲,.
他就說,.
你的食譜說要用心鍋,.
是的,.
這個湯煲真的夠心了,.
我就說,.
這樣吧,.
因為我平時煮不是用那個煲,.
他見過,.
一起煮過,.
我說,.
心鍋,.
這個真的是心鍋,.
心鍋的意思,.
就是不要用平底鍋,.
那麼薄的,.
用這個就可以了,.
因為太多水,.
然後我兒子又問同一個問題,.
你猜他問,.
吃不吃得,.

$^{561}$他就說,.
吃得,.
看好不好吃,.
吃吧,.
這次只是淡了一點,.
還是有點味道,.
不要喝湯了,.
都可以吃的,.
我又和他一起吃,.
又真的吃了一半,.
我明白他,.
還是未上手,.
還是未純熟,.
即是說,.
我還未教懂他,.
即是需要繼續練習,.
而不是拍打他,.
我不是教過你,.
一次兩次,.
你為甚麼不記住他,.
你是甚麼人,.
你有沒有聽我說話,.
我選擇的是,.
忍他,.
我忍不是我忍你,.
下次你再煮多一天,.
整鍋倒在你身上,.
你才吃完,.
不是這樣,.
是和他一起吃,.
正面的鼓勵他,.
希望他下次可以更好,.
這是剛才說的忍耐,.
明白對方的限制,.
給時間,.
給機會,.
給空間去成長,.
不是小朋友是這樣,.
成人也是這樣,.
都是需要時間空間,.

$^{601}$所以對我們的配偶,.
都一樣要這樣,.
需要忍耐,.
需要溝通之外,.
需要等候,.
這些美德特質,.
都能夠幫助我們,.
和我們家人,.
好好和諧地相處,.
好好善待他們,.
第十四節說,.
在這一切之上,.
還要有愛心,.
愛心是聯繫全得的,.
在這些美德之上,.
我們更加需要一樣,.
就是愛心,.
有愛,.
愛就是我們能夠,.
活出這些美德的動力,.
我不愛他,.
我不會忍他,.
我不愛他,.
我不會對他恩賜,.
我不愛他,.
我不會憐憫他,.
就是因為我愛,.
有這一份的愛,.
主耶穌也是因為愛我們,.
所以才可以為我們犧牲,.
穿上新人,.
在第十三章十五節上半部分說,.
又要讓基督的平安,.
在你們心裡作主,.
以基督的平安在我們心裡,.
作為和我們家人相處的時候,.
一些決定的原則,.
什麼意思呢?.
當我們和家人有些張力,.
有些困難的時候,.

$^{641}$我們怎樣去決定,.
怎樣去對他們,.
基督耶穌在我們心裡的平安,.
成為一些原則,.
我結婚兩年多之後,.
就和我奶奶一住,.
一住就住了十年,.
我奶奶是一個高度極度整潔的人,.
而我剛剛相反,.
是一個高度極度自由擺放的人,.
你想像,.
不是我很辛苦,.
我奶奶也很辛苦,.
她要遷就我,.
在她看不到的範圍,.
即是我和我先生的房間,.
我就隨意擺放,.
在她看到的範圍,.
即廳,廚房,廁所,.
我就按照她的規則,.
怎樣放怎樣擺放,.
盡量,盡量,.
我奶奶其實很疼我,.
對我很好,.
她也和我先生說,.
我們叫她做奶奶,.
奶奶說,.
水晶什麼都好,.
就是收拾東西就一般般好,.
但是她也明白我,.
所以我們很和平地相處了十年,.
後來中間也有衝突,.
但真的不多,.
但有一次,.
我很急趕,.
把東西放進廳裡,.
然後就出去,.
很急趕,.
因為我很急趕放進廳裡,.
沒有放好,.

$^{681}$然後我回到家,.
又很急趕,.
我想拿回那東西就走,.
怎料,.
回到家,.
拿走那東西,.
我又很急趕時間,.
我就找找找,.
後來我奶奶才告訴我,.
放在哪裡不好,.
阻礙了,.
然後我就很生氣,.
因為很急趕時間,.
我就很生氣,.
就罵我奶奶,.
你為什麼拿走我的東西,.
我放了你一會兒,.
一會兒也不行,.
我等了很久,.
你知不知道,.
然後我拿完就走,.
我很兇地對我奶奶說,.
然後我去辦事的時候,.
其實我辦事的時候,.
我心裡不平安,.
我心裡真的不平安,.
我覺得,.
其實何必這樣罵我奶奶呢,.
她一向都是這樣收拾東西的,.
但是我另一個聲音就來了,.
在我心裡又說,.
你平時都沒有這樣對你奶奶說話,.
一次而已,.
而且她真的拿走你的東西,.
但是我心裡仍然不平安,.
是,.
我同意,.
我這麼久都沒有很兇地對我奶奶說話,.
其實是她亂放了我的東西,.
但是我心裡仍然不平安,.

$^{721}$然後我覺得是聖靈責備我,.
然後我回到家裡就對我奶奶道歉,.
我跟她說,.
媽媽,.
不好意思,.
剛才我真的很趕,.
很生氣,.
不好意思,.
好像罵到你一樣,.
我奶奶的反應是什麼呢,.
知道你的了,.
趕時間,.
沒事的,.
我們又很開心地繼續相處,.
有時候我們信主,.
我們聖靈在我們心裡,.
基督在我們心裡會告訴我們,.
究竟應該是怎樣去做,.
有什麼是不應該去做,.
當我們心裡沒有這個基督的平安的時候,.
那樣東西不要去做,.
如果做了是有平安的話,.
我們就去做,.
接著,.
保羅又提醒我們,.
你們要讓基督的道,.
豐豐富富的住在你們心裡,.
以各樣的智慧彼此教導,.
互相勸誡,.
用詩章,.
聖詩,.
靈歌,.
懷著感恩的心歌頌神,.
我們要將神豐豐富富的道放在我們心裡,.
要將神的話語放在我們心裡,.
透過無論是個人的,.
或者是群體性的學習,.
群體性學習,.
即是無論是個人去學習神的話語,.
還是群體性一起學習神的話語,.

$^{761}$彼此的提醒,.
我們去敬拜神的時候,.
用詩歌敬拜神,.
不只是唱詩那麼簡單,.
詩歌裡面的歌詞都是在講神,.
詩歌的歌詞都是在教導我們,.
所以無論是透過自己個人學習神的話,.
集體聽到主日學,.
一起去看神的話語,.
透過敬拜的時候,.
都將神的話豐豐富富,.
充滿我們的心,.
讓神的話去改變我們,.
有神的話,.
作為我們怎樣去和家人相處..
試想想,.
雖然是一家人,.
一起住,.
但是仍然是背景不同,.
性格不同,.
想法,.
看東西的方法都不同,.
一起生活一定會有差異,.
這些都會影響我們的相處,.
但是假如我們能夠因著愛的緣故,.
明白對方的限制,.
願意謙卑放下自己的意見和權利,.
為著對方的好處的話,.
我相信我們能夠和家人好好相處..
保羅去到第十七節,.
給信徒一個概括全部生活的原則,.
你們無論做什麼,.
或說話,.
或行事,.
都要奉主耶穌基督的名,.
保羅提醒我們,.
當我們認主,.
基督為我們的主,.
我們的主人,.
和主聯合,.

$^{801}$有新的身份的時候,.
我們在日常生活裡,.
在各方面的層面,.
都應該展現出配合這個身份的行為,.
所以在我們的平日生活裡,.
不只是在教會裡,.
是在日常生活裡,.
全部的生活裡,.
無論信主的人對著信主的家人,.
或未信主的家人,.
我們和他們相處的時候,.
無論說話,.
行事,.
都應該傳著敬畏主的心,.
為主而做,.
不是為了自己,.
凡事按主的吩咐,.
遵主的命令,.
憑著主的能力,.
令我們在日常生活裡,.
在每一件事上,.
去經歷神,.
敬拜神,.
不只是在教會裡,.
不只是在詩歌裡敬拜,.
是在我們日常生活裡,.
去活出主自己,.
這就是敬拜..
這段經文對我們有什麼提醒呢?.
首先,.
我們靠主追求合神心意的美德,.
作為一個好家人,.
要謙卑,.
要對人有恩慈,.
要溫柔忍耐,.
彼此寬容饒恕,.
這些我們都知道的,.
但也不容易做到,.
提到的這些美德,.
不容易建立的,.

$^{841}$不過,.
不要忘記我們的身份,.
記不記得我們是什麼身份?.
我們是被揀選,.
蒙愛聖潔的,.
這個身份不是我們自己可以有,.
是神給的,.
既然是神給的時候,.
神會負責幫我們,.
去追求這些喜悅的美德,.
我們可以向神求,.
去倚靠祂,.
給我們能力去愛我們的家人,.
單靠自己一定做不到,.
不是單單意志可以幫助自己,.
是需要靠神,.
給這個身份我們的神,.
來建立這些美德,.
當然,.
我們都有我們的本份,.
就是我們願意放開自己,.
讓神的令,.
讓神的話,.
來塗做我們,.
去愛我們的家人,.
有一位姐妹,.
她很怕見她媽媽,.
很怕接她媽媽的電話,.
為什麼呢?.
她媽媽不知道為什麼,.
退休之後,.
人有很多負面的情緒,.
很大的負能量,.
每次這位姐妹去探望她媽媽的時候,.
她媽媽就會數這個不好,數那個不好,.
數那位姐妹也不好,.
那位姐妹聽得多就生氣,.
有時候不歡而散,.
跟她媽媽吵架,.
久而久之,.

$^{881}$她就很怕,.
很怕見她媽媽,.
很怕再聽這些負面的事情,.
好像很大壓力,.
好像承載不了她媽媽,.
所以她就過時過節,.
回去做節,.
才去探望她媽媽,.
她媽媽打電話來的時候,.
一看到內電顯示,.
媽媽馬上找她兒子,.
你婆婆來的,.
你跟婆婆聊幾句,.
婆婆不會罵孫子,.
叫「洗個仔」,.
有時候媽媽就「洗個仔」,.
打給婆婆,.
你打給婆婆,.
跟婆婆說,.
你訂了哪裡吃飯,.
叫她幾點去,.
記住跟婆婆說,.
問候她幾句就可以了,.
就「洗個仔」她,.
但她心裡不平安,.
她覺得這樣對媽媽總不好,.
心裡不平安,.
後來這位姐妹就跟她教會的牧者說,.
她教會的牧者建議她,.
「你生先辭卒」,.
即是甚麼?.
「你去探望你媽媽,.
由得你媽媽說,.
你甚麼都不要說,.
你只要聽,.
由得你媽媽說」,.
然後她說「很難」,.
但她又說「好,儘管試試」,.
所以每一次她要去探望她媽媽,.
她當然不只是過時過節,.

$^{921}$最起碼一個月都要去一次,.
當她要去見她媽媽,.
或者打一碼給她媽媽的時候,.
她就找大英姐妹,.
找團契的大英姐妹為她祈禱,.
「我何時要去探望我媽媽,.
我待會要打給我媽媽,.
你為我祈禱」,.
祈甚麼?.
「你為我只可以努力去聽,.
不要亂說話,.
有智慧地去回應,.
不要和我媽媽吵架,.
給我有忍耐,.
給我有愛心,.
你為我祈禱吧」,.
即是叫團契的人在群組立即WhatsApp出來,.
這個姐妹她祈禱,.
不是求她媽媽不要發那麼多牢騷,.
不是求她媽媽不要說那麼多不好的話,.
她在求甚麼?.
求自己去改變,.
求自己不要和媽媽吵架,.
求自己上帝給她有智慧去回應,.
求她有忍耐,.
很多時候我們會容易看到人需要改變的地方,.
我們就希望對方去改變,.
但想想當每一個人只是期望對方去改變,.
會發生甚麼事?.
關係是否會好?.
當我只是希望對方改變,.
關係很難改變,.
各位弟兄姊妹,.
聖經是神給我們去看,.
去教導我們,.
首先是教導自己去改變去成長,.
然後才去提醒,.
而不是一看到,.
聖經這樣說,.
你有沒有做?.

$^{961}$你應該怎樣?.
神的話是首先給我們自己看,.
給自己改變,.
這個姐妹去求的不是媽媽減少發牢騷,.
她去求的是改變自己,.
這個姐妹就這樣去求,.
然後就去探望媽媽多了,.
打電話給媽媽多了,.
她媽媽繼續在發牢騷,.
是少了一點點,.
不過這個姐妹的確有改變,.
當她聽她媽媽多了,.
媽媽說她自己的話的時候,.
她明白她媽媽多了,.
她明白她媽媽以往受了很多苦,.
她不是不知道的,.
不過沒有那麼詳細,.
她體諒了她媽媽裡面的苦毒是從何而來,.
她就更加可以聽得久了,.
沒有出聲,.
她不是在對著空氣說話,.
她真的用心聽,.
她也給媽媽有回應,.
讓她媽媽知道有人聽到她說話,.
有人見到她,.
終於有人見到她,.
有人聽到她,.
後來這個姐妹還有機會跟她媽媽傳福音,.
雖然她媽媽當時還未信主,.
還未願意,.
不過她媽媽開始願意團聚這個聚會,.
也願意接受其他人為她祈禱,.
所以我們可以去求神,.
去追求這個美德,.
不是靠自己,.
神是會幫助我們,.
但我們自己也需要開放自己的生命,.
我們有我們的本分,.
然後就是需要有同行支援,.
剛才說我們最大的支援是神,.

$^{1001}$不過地上我們仍然需要支援,.
我們不能靠自己,.
剛才那個姐妹也說要去探望她媽媽的時候,.
馬上WhatsApp請人類去祈禱,.
我們需要有同行者,.
所以我們需要團契,.
需要群體,.
因為我們需要彼此提醒,.
我們每個人都有盲點,.
有時候看不到自己的盲點,.
很需要別人去提點我們,.
當然我們要找一個成熟的弟兄姊妹來提點我們,.
在神學院第二年的時候,.
神學院要求我們每一個神學生,.
要找一個生命師傅,.
神為我預備了一個女牧師,.
在另一間教會,.
她每個月都會見我一次,.
我會跟她聊天,.
什麼都說,可能在家裡,.
可能學習,可能在實習,各樣東西,.
她說的時候,.
她會聽到一些我自己沒有留心,.
我自己生命上的一些問題,需要,.
她聽到我的一些源頭,.
而這些生命上的一些問題的源頭,.
如果我不去處理,.
不單會影響我日後的事況,.
也會影響我和人際的關係,.
甚至影響我自己,.
對自己的看法,.
所以她引領我去看回我自己的需要,.
幫助我,.
一起為愛祈禱,.
所以我們需要同行者,.
都盼望你們在教會裡面,.
能夠找到一些祈禱的伴侶,.
可以和你一起去祈禱,.
為你守望,.
不要小看這些祈禱,.

$^{1041}$實在是可以幫助我們,.
幫助我們去和人建立和諧的關係..
這段經文,.
讓我們知道家庭的角色,.
我們應該在這個角色應盡什麼本份,.
作為夫妻,父母和子女,.
在家庭裡面,.
有什麼本份應該去擔當,.
我們要盡上這個家庭角色的本份,.
需要除去舊人,穿上新人,.
來建立美好的關係,.
促進和我們家人和諧的關係..
但是無論是脫去舊人,穿上新人,.
建立這些美德,.
我們都需要依靠神,.
揀選我們那位的神,.
單靠自己,絕不能去完成,.
絕不能去好好的改變自己,.
願我們在日常生活裡面,.
都能夠全著敬畏主的心,.
遵著主的吩咐,.
靠著主的能力,.
和未信主和信主的家人,.
好好地去相處,.
藉著這樣去稱頌神,.
我們一齊去祈禱..
摯愛的阿爸父,.
你揀選我們,.
成為你的兒女,.
你也讓我們各人有自己的家庭,.
求主你也在當中幫助我們,.
能夠靠著主你,.
去建立一個美好的生命,.
用這個美好的生命,.
和我們的家人,.
無論是肉身的家人,.
又或者是屬靈家裡面的家人,.
能夠好好地相處,.
彼此去建立,.
讓我們在我們的生活當中,.

$^{1081}$一齊去見證你的名,.
榮耀主你自己,.
多謝你,.
祈禱奉主耶穌得勝,.
明治祈求,.
阿們..
\newpage



\section{瑪拉基書 3:7-18}
\label{sec:032123}
\textbf{失感情——十一奉獻的探討}~~~~梁永善牧師
\newline
\newline
連結: \href{http://www.church.com.hk/acms/content.asp?site=cdc\&op=show\&type=product\&code=032123\&layout=sermon}{\texttt{032123}} ~~~~ 語音日期: 2026-05-10
\newline
\newline
\hyperref[sec:032122]{\small{< < < PREV SERMON < < <}}
~
\hyperref[sec:index_chronic]{\small{[返順時目]}}
~
\hyperref[sec:032124]{\small{> > > NEXT SERMON > > >}}
\newline
\newline
瑪拉基書 3:7-18
\newline
\begin{longtable}{cl}
\hline
\hline
章節 & 經文 (和合本修訂版)\\
\hline
3:7 & \begin{tabularx}{0.7\textwidth}{X} 從你們祖先的日子以來，你們就偏離我的律例而不遵守。現在你們要轉向我，我就轉向你們。這是萬軍之耶和華說的。你們卻說：『我們如何轉向呢？』 \end{tabularx} \\ \\ \relax
3:8 & \begin{tabularx}{0.7\textwidth}{X} 人豈可搶奪神呢？你們竟搶奪我！你們卻說：『我們在何事上搶奪你呢？』其實就是在你們當納的十分之一奉獻和當獻的供物上。 \end{tabularx} \\ \\ \relax
3:9 & \begin{tabularx}{0.7\textwidth}{X} 因你們全國上下都搶奪我的供物，詛咒就臨到你們身上。 \end{tabularx} \\ \\ \relax
3:10 & \begin{tabularx}{0.7\textwidth}{X} 你們要將當納的十分之一全然送入倉庫，使我家有糧，以此試試我，是否為你們敞開天上的窗戶，傾福與你們，甚至無處可容。這是萬軍之耶和華說的。 \end{tabularx} \\ \\ \relax
3:11 & \begin{tabularx}{0.7\textwidth}{X} 我必為你們斥責蝗蟲，不容牠毀壞你們的土產。你們田間的葡萄樹，果實未熟以先也不會掉落。這是萬軍之耶和華說的。 \end{tabularx} \\ \\ \relax
3:12 & \begin{tabularx}{0.7\textwidth}{X} 萬國必稱你們為有福的，因你們必成為喜樂之地。這是萬軍之耶和華說的。」 \end{tabularx} \\ \\ \relax
3:13 & \begin{tabularx}{0.7\textwidth}{X} 耶和華說：「你們用話頂撞我。」你們卻說：「我們說了甚麼話頂撞你呢？」 \end{tabularx} \\ \\ \relax
3:14 & \begin{tabularx}{0.7\textwidth}{X} 你們說：「事奉神是枉然，我們遵守神所吩咐的，在萬軍之耶和華面前哀痛而行，有甚麼益處呢？ \end{tabularx} \\ \\ \relax
3:15 & \begin{tabularx}{0.7\textwidth}{X} 現在，我們稱狂傲的人為有福，並且行惡的人得以建立；他們雖然試探神，卻得以逃脫。」 \end{tabularx} \\ \\ \relax
3:16 & \begin{tabularx}{0.7\textwidth}{X} 那時，敬畏耶和華的人彼此談論，耶和華側耳而聽，且有紀念冊在他面前，記錄那敬畏耶和華、思念他名的人。 \end{tabularx} \\ \\ \relax
3:17 & \begin{tabularx}{0.7\textwidth}{X} 萬軍之耶和華說：「在我所定的日子，他們必屬我，是我寶貴的產業。我必憐憫他們，如同人憐憫那服侍他的兒子。 \end{tabularx} \\ \\ \relax
3:18 & \begin{tabularx}{0.7\textwidth}{X} 那時你們必再一次看出義人和惡人，事奉神和不事奉神的人有何差別。」 \end{tabularx} \\ \\
[1ex]
\hline
\hline
\end{longtable}
$^{1}$各位電影節目你們好.
我今天很開心跟大家分享.
但可能大家不是很開心.
今天聽我講到.
今天在香港是母親節.
我已經說了.
來了英國的婦女.
你們真的特別尊貴.
因為你們有兩個母親節要過.
英國的母親節不是今天.
是3月.
所以一定要和你們慶祝.
我們要日香隨俗.
但今天可不可以不慶祝呢.
是不可以的.
因為你會說這麼忘本.
我們香港是母親節.
我都提醒青少年.
我有個青少年的聚會.
都叫他們一定要.
今天特別要做一點事.
就是去給媽媽.
讓媽媽可以開心一點.
所以都鼓勵.
今天在我們當中參加聚會的.
如果你的配偶.
即是你太太在你身邊.
你很小聲跟她說.
母親節快樂.
我多謝你.
沒有你我都沒有今天.
真的要說的.
我原本是想說母親節的課題.
我講一篇道德叫做.
「偉人母親的條件」.
但牧師跟我說.
梁牧師.
你下次來我們教會.
講奉獻的道理.
OK.

$^{41}$他叫我講.
我說沒問題的.
那時候教會都是叫我講奉獻的道理.
我都問一些牧者.
我不是問羅牧師.
我問其他牧者.
為什麼?.
我說不是.
因為你們是講員.
你講完就講完.
這裡主要是你說.
我是沒問題的.
所以我就照樣講.
今天我會用一個馬戲班的故事作為開始.
也用一個馬戲班的故事.
作為今天講道的結束.
在這個馬戲班裡面.
有一個節目.
有一個大力士.
很厲害的.
他說.
有一個很厲害的西柚.
然後叫人拿一個瓶子來.
我拿.
我拿西柚.
西柚拿了之後.
真的滴了很多汁出來.
這個被我拿乾了.
有誰可以出來.
再拿.
拿到汁出來.
我們馬戲班.
是會有獎賞給他的.
年輕人出去.
不斷捏.
都沒有.
到最後有一個伯伯出來.
震一震.
讓我試一下.
不要玩了 伯伯.

$^{81}$你都可以的.
讓我試一下.
於是他就慢慢出汁.
真的有幾滴汁出來.
那個人就很驚訝.
他這麼厲害.
伯伯.
你做聖巷.
伯伯說.
我是牧師.
你會行神蹟.
不是.
我習慣會有奉獻.
所以.
他們沒有.
我都可以捉到他.
這個就是西方的笑話.
挖苦我們牧師.
說我們經常叫人去奉獻.
所以今天的題目就叫做「失感情」.
但是.
我想說的.
不只是為了鼓勵大家奉獻.
是很想大家了解奉獻的意義.
很多人就說.
現在新約時代裡面.
沒怎麼說到守一奉獻.
所以我們就可以.
不需要遵守守一奉獻了.
我們特別很多教會用錯了一段經文.
哥林多後書第八第九章.
特別一句.
「捐得樂意是神所喜悅的」.
不要作難.
不要心裡面離題.
總之你樂意就可以了.
如果你以為這段經文是講奉獻.
你是解錯聖經.
你看回這段經文是講.
慈惠.

$^{121}$因為耶路撒冷的教會有大饑荒.
於是他們有一班人籌錢.
去幫助他們.
所以這個「捐得樂意」是指什麼呢?.
是指你去幫助人.
就好像當年橫幅會園.
大火.
你認識一些人.
他們的家園被毀了.
你去幫助他們的時候.
這就是「捐得樂意」.
不是勉強的.
奉獻.
很現實地去說.
你不願意的話.
你都要去做的.
所以剛才我讀那個馬拉基書的時候.
是很清晰地講到.
是當立的十分之一.
有哪個人喜歡交稅的?.
我很現實地說.
我也不太喜歡交稅.
但是來到英國.
那個稅項很重的.
但是你不交不行.
如果你瞞稅的時候.
很大件事的.
而神也很現實的.
你看回馬拉基書第三章第七到第八節說到.
萬軍耶華說.
你們列祖幾日子以來.
你們時常偏離我典章而不遵守.
現在你們要轉向我.
我就轉向你們.
就說你們要歸回.
你們經常都犯錯.
人們就問.
那我怎樣去轉回?.
不是叫他去認罪祈禱.
不是叫他讀聖經.

$^{161}$不是叫他參加崇拜.
你看看神真的很現實.
竟然叫他們去奉獻.
而且這個奉獻是什麼?.
引起可奪取神的物.
你是搶了神的東西.
原來你不奉獻是搶了神的物品.
我們有沒有想過.
原來不奉獻是一件很嚴重的罪.
接著就說到.
你們在當立的十分之一上.
你們來來去去.
沒有做到.
是一定要做的.
你們沒有做到.
所以今天我們就會探索.
十一奉獻.
究竟什麼時候開始有十一奉獻呢?.
如果你留意到.
摩西眉斑部律法之先.
就已經有十一奉獻了.
亞伯拉罕很值得我們欣賞.
他以德報怨.
記不記得那個羅德.
他選擇了一個平原.
後來漸漸來移地.
去到索德瑪俄摩拉.
後來四王與五王爭戰的時候.
他就在夾縫之中.
他被人擄走了.
有些人走出來告訴亞伯拉罕.
亞伯拉罕說該死了.
選擇了一個這樣的地方.
搞到我要上山.
你記不記得.
什麼叫先輸後贏.
亞伯拉罕忍讓便走上山.
那一晚神就跟他說.
你出來看看.
你見到的地方.

$^{201}$都是屬於你的子孫了.
你不要忘記.
亞伯拉罕那時候已經很大年紀了.
大年紀別人有怨事.
所以他看得很遠.
所以那些地方很多是屬於他的.
但是羅德只是看著眼前.
結果什麼都沒有了.
後來亞伯拉罕帶著三百多家丁.
去與四王爭戰.
結果贏了.
但是你不要忘記.
那三百多家丁從來沒有試過打仗.
沒有精良的武器.
沒有軍事訓練.
所以他贏了這場仗.
他絕對肯定是神給他祝福.
當他回去的時候.
他就做了什麼呢?.
我們一起讀這段經文.
就是在《創世記》的二十八章.
二十到二十二節.
我們一起讀出來好不好?.
第一次出現十分之一.
就在這段經文裡面.
亞伯蘭(亞伯拉罕).
他很深感受到神給他的眷顧.
他懷著感恩的心.
為了獻十分之一給默基使德.
希伯來書裡面有提到.
默基使德是誰呢?.
不是說.
我使的是默基嘛.
很能使的.
他經常扮飛仔.
不關事的.
那個默基使德是撒冷王.
這是另一個課題.
他不是被逼.
他是因為感恩.

$^{241}$於是他獻上十分之一.
第二個說十分之一的是誰呢?.
是雅各.
為什麼雅各欺騙了哥哥的長子名分.
因為耳熟說要殺害他的時候.
他要逃亡.
這個是裙腳仔.
平時經常在家裡煮紅豆湯.
現在要走的時候.
他很孤單.
他躺在這裡的時候.
但是神就向他顯現.
我會祝福你.
雅各就說.
如果你真的祝福我.
我就會做些什麼呢?.
我們一起來讀這段經文.
在《創世記》二十八章二十到二十二節.
神向雅各說.
我會祝福你.
他就跟神說.
如果你真的祝福我.
我會獻上十分之一給你.
後來他有沒有獻上我不知道.
但至少他有這個應許.
所以不是被逼的.
所以你留意到.
原來十分之一不是在摩西的年代才定下來.
四百年之前.
在摩西以前已經有這件事了.
後來當他們成為國度的時候.
律法的制定就要獻上十分之一.
原來十分之一是有三個十分之一的.
我有筆記已經給了教會了.
你可以來看.
第一個十分之一是什麼呢?.
在利未記二十七章三十到三十三節.
以色列人要獻上十分之一.
就是讓聖殿或者當時會望.
可以維持著運作.

$^{281}$因為是很大的需要.
第二方面是什麼呢?.
就是給利未人在民數記十八章二十一到三十二節.
因為利未人就是服侍神的人.
他是沒有自己的地業.
他是要分散在不同的支派裡面.
去服侍人們的.
所以要讓他們有生活上無後顧之憂.
來供養他們.
正如保羅所說.
今天傳福音的就要靠福音來養生.
所以這個第二個十分之一.
就是要來支持那些服侍神的人.
第三個十分之一就很特別.
是三年要獻上一次.
這是什麼呢?.
放在基金那裡.
是幫助那些孤兒寡婦.
因為那個年代裡面.
都有福利制度的.
如果他們的丈夫死了.
一家之主已經沒有那個體度的時候.
他們是很貧窮.
於是就需要關心他們.
但是如果我們看回整體的子民.
去遵守十分之一.
是很少的.
原來這是人性.
在舊約年代裡面.
都不是這麼多人遵守十分之一.
什麼時候要珍惜全國.
要遵守十分之一呢?.
就是欺世嫁王.
欺世嫁王去尼德雅親近神復興.
去尼德雅看到神的祝福.
於是他就下命令.
要所有的子民都要來呈獻十分之一.
我們一起去看這一段經文.
在歷代志下第31章的第3,第4兩節.
我們一起去讀出.

$^{321}$今次是欺世嫁王下命令.
要全國的人都要來去獻上十分之一.
當全國人獻上十分之一的時候.
是令到當時的祭司非常非常驚訝.
在歷代志下第3章第10節.
原來當每個人都肯來奉獻的時候.
是令到他們非常的驚訝.
第二次全國十分之一的奉獻.
就是約西亞王又是宗教上的復興.
去親近神.
於是有第二次的全國十分之一.
最後一次是什麼呢?.
就是尼希米當他們回歸的時候.
尼希米責備他們.
我們已經回到自己的故土了.
但不要忘記當時仍然是頹垣敗壓.
但原來有些人.
你們沒有去行十分之一.
所以你們有問題.
所以你留意到神要求子民的十分之一奉獻.
不是神自己拿走的.
神絕對不會伸手下來.
拿走所有的東西.
原來就是讓聖殿可以維持運作.
讓一些孤兒寡婦得到援助.
讓服侍神的人可以得到支援.
這個很重要.
所以奉獻是神律法的規定.
也是子民應該盡的責任.
但很多時候人性是怎樣呢?.
我們很想有更多的權利.
但我們不肯承擔我們應有的義務.
這是很重要的.
正如我剛才所說.
我們可能很要求有什麼福利.
但當我們納稅的時候.
我們就會牙痛了.
又要繳稅了.
是很痛苦的一件事.
十分之一的延續是什麼呢?.

$^{361}$教會的歷史讓我們知道.
不是很多人有十一奉獻的.
原來十一奉獻的原則是沒有時間性的.
剛才我們已經看到了.
在律法之前已經有十分之一了.
在亞伯拉罕雅國的時候已經有十分之一了.
但也不是恩典時代的終結.
正如今天很多基督徒說.
《新約聖經》裡沒有提到十分之一.
如果你留意到耶穌基督也有提到十分之一.
但這個十分之一是什麼呢?.
祂是在罵法利賽人的.
我們帶大家去看《馬太福音》23章23節.
我們一起來讀.
你們這假冒為善的民事.
和法利賽人有禍了.
因為你們將薄荷,茴香,芹菜獻上十分之一.
那律法上更重的是.
就是公義,憐憫,信實,反倒不行了.
這更重的是你們當行的.
那也是不可不行的.
你別以為我們買的西芹,芹菜是那時的調味料.
即是你的胡椒粉.
也要分十分之一出來.
你連這些細微的東西也奉獻.
但最重要的是什麼?.
你的公義,憐憫,信實.
你沒有做到.
你只是表面做這些.
所以你是錯的.
所以你留意到原來沒有否定十分之一的奉獻.
而今天我們在恩典之下是不是不用守律法呢?.
我想這段經文很熟悉.
《羅馬書》第六章十四十五兩節.
我們一起來讀出,好嗎?.
「罪必不能作你們的主.
因你們不在律法之下.
乃在恩典之下.
這卻怎麼樣呢?.
我們在恩典之下.

$^{401}$不在律法之下.
就可以犯罪嗎?.
斷乎不可.」.
很清楚.
今天我們不是在律法之下.
我們不是因為很害怕.
死顧著守律法.
我們是要守律法.
不過不是再因為害怕.
更加加上因為如果我犯錯.
我的神是會失望的.
是因為愛之下而守律法.
我們不可以因為在恩典之下.
於是就放棄這個律法.
其實在恩典時代裡面.
律法是應該更加被提升的.
很簡單舉個例子.
十誡裡面.
不可殺人,是不是?.
拿把刀捅你.
這就叫做殺人.
但耶穌怎麼說呢?.
當你生一個人的氣的時候.
你恨一個人.
你就已經怎樣呢?.
殺人了.
我想問大家.
耶穌將律法降低了還是提升了呢?.
不敢回答我嗎?.
提升還是降低?.
提升了,是不是?.
好了.
十誡裡面不可姦淫.
你是和一個婚姻以外的人.
有性行為.
這就叫姦淫.
但耶穌怎麼說呢?.
我看著弟兄而已.
當你看到一個姐妹.
一個女性.

$^{441}$然後你心裡動淫念的時候.
你就已經怎樣呢?.
犯了姦淫.
我想問大家.
他是將律法降低了還是提升了呢?.
提升了.
但很奇怪的.
一到奉獻.
舊約是十分之一的.
如果照這個道理.
我們應該是超十分之一.
還是少於十分之一呢?.
理論是應該超十分之一的.
但很多時候我們都會隨心而.
我記得我自己小時候.
怎麼會奉獻呢?.
其實這幾天特別深感受.
因為我當年最好的朋友.
也是北宣的.
我們是同年的.
他在加拿大當牧師.
剛好這幾天來英國探望我.
我昨天會見他.
我明天也會送他上路.
明天就回加拿大了.
我們在這兩天裡勾起了很多童年回憶.
我想想我當年怎麼會奉獻呢?.
我是按照聖經的吩咐去奉獻的.
聖經說到你奉獻的時候.
你捐的時候.
你左手不要叫右手支.
所以我不知道我捐了多少錢.
每次奉獻袋傳過來的時候.
我們那個年代不流行袋錢包.
所有錢都放在袋子裡.
我馬上從袋子裡慢慢按.
紙幣甩開.
大的紙幣甩開.
拿幾個小錢.
就放進去.

$^{481}$傳到多少錢.
我們很聰明的.
聽聲音是會知道他扔了多少錢.
那些五元,一元和五毫子.
是不同的聲音.
所以我們有兩種方法.
一種是輕輕按下去.
讓他沒聲音.
一種是插進去.
嘩 好像很多.
這就是我小時候.
回到少年團體去奉獻的方法.
我們經常說不要緊.
我們被一句話誤導了.
捐得樂意.
所以我沒有計算.
但肯定少過十分之一.
是少很多的.
但是後來慢慢長大了.
才認識到這件事.
剛才我說.
在恩典時代裡面.
律法我們仍然要遵守.
但基於愛.
不是基於懼怕.
第二就是.
律法不是降低了.
而是提升了.
所以《羅馬書十二章》.
大家很熟悉.
將身體獻上當作活祭.
不只是錢.
是我們整個人來奉獻給神.
所以奉獻的心態是什麼呢?.
剛才我已經說了.
奉獻是我們必定要做的一個原則.
所以即使你自己很不願意.
但仍然要盡力去做.
但是如果你樂意的時候.
你會很不同.

$^{521}$我不知道有沒有跟大家分享過一個經歷.
是我太太教會我.
什麼叫做樂意.
就是.
我太太是大女兒.
譬如外母生日或祭祀.
很不幸地她們是星期六晚吃飯的.
因為小朋友星期日要上學.
所以星期六晚祭祀.
但永遠沒有我的份.
因為我在香港的時候.
星期六晚.
無論是粉嶺堂還是大埔堂.
我是有晚堂崇拜的.
我一定要出席.
所以我是不能去的.
雖然不能去.
但我是要付錢的.
即是外母生日.
做節日當然要付錢.
那當然是我付錢.
每次我覺得.
飯我又不能吃.
錢我又要付.
又或者可以吃.
都是夾回來凍一凍給我吃.
所以每一次我都是.
很不開心地付.
不拿去.
每次都苦不堪言.
有一次我太太跟我說.
你不開心又要付.
你開心都要付.
為什麼你不開開心心給我.
讓我可以覺得你疼我.
那我就會開心些.
當然她的語氣不是這樣.
我是很溫柔的.
她就教曉我這個道理.
所以.

$^{561}$後來我終於付錢的時候.
我是強顏歡笑.
拿去吧,開心心吧.
然後我心裡在唱陳百強的歌.
什麼歌?.
「眼淚在心裡流」.
拿去吧.
但當我笑著付錢的時候.
那就是了.
明白嗎?開心.
明白這個道理嗎?.
原來當我們知道.
這是我們不可以逃避的責任.
你開不開心都要.
但如果你能夠開心去的時候.
你會差很遠.
所以其實我們奉獻的是什麼?.
原來是回報神給我們的恩典.
今天神不單單拯救我們.
你留意到今天神給我們有多少的祝福.
我經常這樣說.
我到今天都很熱心去事奉.
雖然我生命都有一定的缺陷.
但我不會懶.
為什麼?.
因為我是一個敏感的人.
我是敏銳於神和人給我的恩典.
我去感恩.
如果我要和你分享我的見證.
是多到不得可了.
神就是這樣.
很簡單.
原來神是什麼呢?.
我們一起來讀《已忽所思》第一章第三節.
神除了救贖我們.
不是救了我們就算了.
祂一直都照顧我們.
眷顧我們.
祝福我們.
《已忽所思》第一章第三節.

$^{601}$我們一起來讀出.
願頌讚歸於我們主耶穌基督的父神.
祂在基督裡曾賜給我們天上各樣屬靈的福氣.
今天神賜給我們各樣屬靈的福氣.
去眷顧我們.
保守我們.
不單單這樣.
我們再讀《使徒行主願》第十七章第二十八節.
我們生活,動作,存留都在乎祂.
就如你們作詩的.
有人說:「我們也是他所生的」.
我們要上格拉斯哥講道.
我不知道大家知不知道有位叫楊健強牧師.
那星期我應該要見他.
誰知我去愛丁堡那次.
原來楊牧師那天沒有出席.
就好像今天羅牧師不在.
他去了格拉斯哥開新的教會.
我們約了一星期見面.
但他說他有點不舒服.
我說:「好了,不要見了」.
「有些會友仍然堅持叫我們去探望他」.
我說:「不要了」.
誰知星期三去醫院.
我還在跟他WhatsApp.
但兩天之後他就離開世界.
因為COVID-19很嚴重.
他年輕過我,健碩過我.
我昨天收到分領堂的會友跟我說.
「牧師,有一個移民過來英國的會友」.
「在倫敦去世了」.
「不要想著活著是必然的」.
「今天我們生活,動作,存留都在乎他」.
「你有沒有去感恩」.
「我今年都七十二歲了」.
「我仍然可以活著是神的靈魂」.
我們有時忽略了這些.
還有時我們有很多壓力.
很多困難.
很多重擔.

$^{641}$但聖經怎說呢?.
我們一起來讀司篇六十八篇第十九節.
「天天背負我們重擔的主」.
「就是拯救我們的神」.
「是應當稱重的」.
今天.
原來我們的神為我們背負很多重擔.
我們將我們的問題交託給祂.
我們的神這麼愛護我們.
為什麼我們不去進行祂自己的指令?.
今天原來奉獻是愛的回應.
是感恩的回報.
這是很重要的.
而且我們更加要不要忘記.
這是我們的特權.
你想想.
不信的人.
他們是不會奉獻的.
當不信的人去奉獻.
神也不會越立他們.
是我們的特權.
我們有沒有來重視.
原來這是我們少有的.
我們來的是去珍惜這個機會.
而且主耶穌用了自己的生命.
經歷.
來成為我們的鼓勵.
這段經文其實很寶貴.
但很多人都不明白.
我們一起讀哥林多號書第八章第九節.
「你們知道我們主耶穌基督的恩典」.
「祂本來富足」.
「卻為你們成了貧窮」.
「叫你們因祂的貧窮」.
「可以成為富足」.
是不是很複雜?.
本來我們主耶穌基督所謂的富足.
祂是因為我們成為貧窮.
就因為祂貧窮我們富足.
其實很容易明白.

$^{681}$做父母一講你就明白.
你跟子女說.
我原本是富足的.
不過生了你之後.
我就成為貧窮.
你一毫子都沒有帶來這個世界.
你想想你現在多少東西都是我賣給你的.
就因為我的貧窮你就富足.
是不是很簡單?.
《率性經》不是很難明白.
你一講就明白了.
這裡就說到耶穌基督是為了我們的緣故.
而放下了祂自己.
在天上的榮耀尊貴.
來救贖我們.
就因為祂的捨棄.
於是我們成為富足.
所以今天原來奉獻.
是將我們和基督的犧牲放在一起的.
這個奉獻.
就是反映我們自我犧牲的重要.
所以求神來幫助我們.
而當我們去奉獻的時候.
記得我剛才所說的不是神奉獻了祂.
而是會成為別人的幫助.
我們一起來讀.
哥林多後書第九章十一到十三節.
這一段經文.
不是只奉獻.
是講到慈慰.
但是也可以讓我們看到背後的精神.
我們一起來讀.
哥林多後書第九章十一到十三節.
「叫你們凡事富足.
可以多多施捨.
他們從這供給的事上得了憑據.
知道你們承認基督順服他的福音.
多多的捐錢給他們和眾人.
便將榮耀歸於神」.
這裡他帶出兩個道理.

$^{721}$當你捐錢給一些有需要的人.
不單單實質幫到他.
是更加令他們感受到神的恩典.
因為我剛才所說.
耶路撒冷大饑荒.
基督徒被迫害.
他們已經被人拿走所有財產.
所以當他們在困難時.
一定向神呼求.
後來神透過米其頓的教會.
尼泰去怎樣呢?.
將捐獻透過保羅拿給他們.
當他們拿到捐獻後會怎樣?.
神聽了我的祈禱.
神尼泰去藉著一些人來幫助我們.
他便會感受到他們的禱告上.
神的憐憫.
也不單幫助在物質上解決困難.
也讓我們更深經歷神.
所以我透過你們的捐獻.
讓他們尼泰去在生命上.
可以有更大的信仰上的突破.
但對你們來說是甚麼呢?.
是你們信服他的福音.
因為這是神的教導.
所以今天我們有沒有做的時候.
說是容易信服.
但我們有沒有實踐?.
如果我們說那些是信服神.
但連最基本的事情我們都沒有做的時候.
那很悲哀.
你這個信服是說謊的.
所以求神尼泰去幫助我們.
一句說話很好.
很寶貴.
我們所花用的.
都會失去.
即是你的錢用了便沒有.
好了.
你不尼泰去用.

$^{761}$一味留在這裡.
但一日你死了.
那些叫遺產.
你是帶不走的.
但今天如果我們奉獻的話.
某個程度.
因為神一定會記得.
你是投資一個回報率很高的基金.
一定不會虧損的.
即是如果你確信私比壽更為有福.
你奉獻了的時候.
其實不是一個損失.
而是你這個祝福.
是存留到永恆的.
好了.
我們想一想.
究竟我們應該奉獻多少.
才夠呢?.
你留意到耶穌基督所看的.
不是金錢的數額的多寡.
是你的心意.
大家很熟悉的.
這個不是比喻.
是耶穌留意人們去奉獻.
然後作出一個回應.
我們一起來讀.
路加福音21章第一到第四節.
我們一起來讀.
有一個有錢人.
特意在眾人面前捐獻.
有很多錢.
他真的很好.
他捐了那麼多錢.
一個窮寡婦.
只是捐幾個小錢.
你留意到耶穌沒有否定.
那個有錢人的捐獻.
耶穌說其實他的捐獻也好.
但他已經得到他的獎賞.
人們拍了他的掌.

$^{801}$但這個寡婦這幾個小錢.
是更加值得讚賞.
如果用金錢的幣值來看.
這個寡婦的那兩個小錢.
是很少的.
但耶穌為何大力讚賞這個寡婦.
因為她犧牲度很高.
所以你留意到最重要是心態.
而不是只看那個幣值.
但有時候.
我特別說父母那些.
有時候怎樣去教導你的兒女信仰.
很強調一定要好好的.
要承傳這個信仰.
但有時候.
特別是母親節.
我們便會看看.
你甚麼樣的樣本可以給你的兒女.
第一.
你的兒女是否看到你.
每天都很著緊零售.
親近神.
第二就是.
回到主一學.
我不知道你們主一學有否收奉獻.
有的.
我當年在牧會的時候.
我們主一學也有收奉獻.
當我們收奉獻的時候.
我看一看.
我不會碰的.
我看一看.
為何我們教會有很多寡婦.
那麼多小錢.
已經很少見到那些兩毫或一毫.
五毫在香港.
竟然在我們的奉獻中出現了不少.
其實我們教會不是那麼多寡婦.
只是那麼多小錢出現.
喂.

$^{841}$你出去給小貼士也不止這個數.
但可能你給五毫一元.
主一學奉獻.
如果你的兒女懂事的話.
這個神原來很低賤.
原來就是給這些錢.
你知不知道可能我們這些行為.
某個程度是在貶低我們的信仰.
你在教你的兒女甚麼.
所以我自己.
譬如我現在也會.
我不可以有錢給我的孫子.
但我跟他們說一定要收奉獻.
我很壞.
你們知道我在明恩堂的時候.
牧會是怎樣的.
我們每個月都會出.
有多少錢的奉獻.
但接著我就會出.
我們平均這個月有多少人出席崇拜.
但接著我就會出.
用封奉獻的比例是多少.
我不是看數額.
某個程度好像也很好.
我們第一個星期有時候已經過了一秒.
即是一個的崇拜.
但當然我們的支出也很大.
我們有社會服務.
我不是看數額.
我是看有多少人去奉獻.
計一計數.
原來有三分之一以上的人沒有奉獻.
所以今天最重要不是多少個.
你留意到的時候.
今天奉獻是怎樣的.
剛才我說到.
馬其頓的教會是很值得我們欣賞的.
原來他自己也在很窮困.
在逼迫,在患難當中.
但他們見到人有需要.

$^{881}$他竟然在困難當中去捐錢.
我們一起來看哥林多號書第八章第一到第四節.
我們一起來讀出.
弟兄們.
我把臣賜給馬其頓宗教會的恩告訴你們.
就是他們在患難中受大試煉的時候.
仍有滿足的快樂.
在極窮之間還格外顯出他們樂捐的口恩.
我可以證明.
他們是按著力量.
而且也過了力量.
自己甘心樂意地捐助.
在三地求我們.
准他們在這供給聖徒的恩情上有分.
你看到嗎?.
原來馬其頓的教會同樣面對困難.
面對打壓.
面對迫害.
他們是極窮困的.
這是赤貧.
但我們看到耶路撒冷的弟兄姊妹有需要.
他們就甘心樂意地捐助.
有一句說話你要留意.
他們不但按著力量過了力量.
保羅見到他們的困難都說.
「行了,你不如留下吧」.
「你給自己用吧」.
保羅也不肯要這份捐獻.
但是他求保羅.
在三地求我們.
准他們在供給聖徒的恩情上有分.
「你收吧」.
只有別人逼人去捐錢.
但保羅卻求保羅收他的捐獻.
在我生命當中.
我經歷過特別有兩個人.
一個是我自己的會友.
他已經離開了世界.
因為血癌,他很年輕.
我記得我去探望他.

$^{921}$當時他在中軍澳醫院.
我去探望他.
他已經沒有回到外面.
因為以前在粉嶺.
他叫Tracy.
後來就搬到中軍澳.
當我知道他有血癌.
我就去探望他.
當時他還未造化.
我和他聊天時,他說「牧師,我知道教會有需要」.
「最近我做事都沒有奉獻」.
不知道他帶著一本檢查書.
去醫院.
他馬上把支票.
我記得是二千一還是二千四.
我又做八一幫人.
不過今天沒有帶出來.
叫我去奉獻給教會.
我說「不要了」.
「現在你的病這麼嚴重」.
「你要留一些錢給自己用」.
「不行,這是我的責任」.
「你給我捐吧」.
我收了兩支支票.
我就把支票拍出來給弟兄姊妹看.
我說「Tracy,我們自己的會友」.
「現在有血癌,但他仍然願意奉獻」.
接著我說.
「不過我不會把這張支票」.
「放進教會的奉獻箱」.
「因為當我不放進去」.
「就不會提取那些錢」.
「那些錢就留在那裡」.
幾個月後我去探望他的時候.
原來他有上網聽我講道.
我不知道.
你有沒有搞錯.
他這句說話我都不會忘記.
「你是剝奪了我的祝福」.
「你為什麼不把我的錢奉獻」.

$^{961}$接著我再開一張.
二千五百元多了幾百元.
那次我不敢不奉獻.
我真的奉獻了.
這是我自己永遠都不會忘記的.
一位很年輕的姐妹.
血癌過身.
但另一個弟兄.
九龍城潮語浸信會.
不是外面的殷堂.
他有腦癌.
如果你看他.
他寄了一張卡給我.
上面講道他接受了MRI.
做了手術都沒有果效.
而且壓著神經都很辛苦.
他很厲害.
他唱過歌劇.
但終於早幾年他離開了世界.
但感恩.
醫生說他只有幾個月命.
他多了超過十年的命.
但他寄給我的時候.
還有一張支票.
我們剛剛開粉嶺堂.
他奉獻了五千元給我們.
我跟他說不要了.
你自己留下吧.
因為那時候他太太沒有工作.
他說他們不是很有錢的人.
他住慈雲山的公屋.
這些錢你留給自己用吧.
他這樣說.
我們知道你很樂意傳福音.
你們開新的教會.
所以我們願意奉獻給你.
我永遠都不會忘記這些人.
但今天英國的教會.
人數特別是香港來了不少的人.
但有些牧者就慨嘆人數有增加.

$^{1001}$要求也多了.
但奉獻是不會增加的.
為什麼呢?.
我們會來到這裡.
第一就是.
我們可能找不到工作.
我們找不到工作.
不是我們以前香港的收入.
我們這裡的水費電費.
Council Tax都很貴.
我們看著我們的存款.
可能一路一路減低.
但其實你減低多少.
如果你早幾年過來.
你賣了樓.
都很有一筆收入.
當年你賣了樓.
然後你用了一半的錢來買.
所以為什麼現在英國這麼多關門.
是真的.
你的錢在銀行修飾.
沒事做就關門.
但是我們奉獻是不行的.
先保住.
因為怕坐食山崩.
這是很可悲的.
請你記住.
神不是要求我們沒有.
然後要獻上.
是不會這樣的.
是你有能力.
在你的能力裡面.
你要做一些事.
就正如好像捐獻一樣.
我們一起來讀哥林多號書.
第八章第十四節.
我們一起來讀.
乃要均平.
就是要你們的富裕.
現在可以保他們的不足.

$^{1041}$使他們的富裕.
將來也可以保你們的不足.
這就均平了.
如果你能力承擔不了不要緊.
兩家或者三家的窮人.
一起去買這一隻的班鳩.
呈獻給神.
即是不可以用窮.
或沒有能力.
作為一個藉口.
我自己這樣說.
如果我真的沒有能力去奉獻.
那就跟神說.
神不是一個盲不講理的神.
神知道我現在來到這裡.
我又要給很貴的租金.
或者租金.
我的支出又很多.
現在我又未曾賺到錢.
即是未有一份工作.
所以神說我沒有能力.
但是可以守一奉獻.
你接納吧.
我相信如果我們能夠坦然地跟神說.
神是會接納的.
當然你又要檢討一下.
究竟你平時會有很大的消費.
用在一些無意識上.
或者又買衣服.
又吃得很厲害.
那你可不可以吃少一餐?.
我經常都這樣說.
是呀你沒有能力不要緊.
那就奉獻二十.
或者三十.
或者奉獻二百三百.
如果你找到工作.
下個月都是奉獻二百三百.
第三個月可不可以加少少?.
加多二十元.

$^{1081}$會不會奉獻多二十元.
那個月你就立刻崩潰.
即是經濟全部陷入一個很拮据的情況.
如果沒有的時候.
那你就多加二十元.
一直加上去.
你有很多事情都要求進步的.
我如果住在寮屋.
我就希望上公屋.
上公屋我希望住居屋.
住居屋我希望住私樓.
但在奉獻事上.
我們有沒有為自己設一個目標.
我試下多些.
所以都求神來幫助我們.
最後原來奉獻絕對不是一個損失.
神是會有祝福的.
大家很熟悉這段經文.
《瑪麗基書》第三章第十到第十一節.
聖經裡面很清楚說.
不可試探處理的神.
唯一一段經文可以試下神的.
是神挑戰我們.
你儘管將當立的十分之一.
全然送入倉庫.
不是神拿了.
使我家有糧.
以此試試我.
試下.
看看我會不會傾福於你們.
為你們躺在天上.
那股福氣淋到你們無處可用.
有時你為了保留一些.
但很可能在另一處已經損失.
很感恩我經歷了很多.
我剛剛開始明恩堂的時候.
薪金是非常非常低.
我由七千元在信義會的事奉.
去開荒的時候變成三千元.
但很感恩.

$^{1121}$我的小孩都很健康.
很少病痛.
但我見到有些人.
他們的小孩原來又說哮喘.
又說年紀很小.
不可以看普通醫生.
要看專科醫生.
那時看專科醫生.
幾十年前都要千多元一次.
一個月看兩三次.
就已經沒有幾千元了.
原來我看了原來我的心情很健康.
就已經很足夠.
後來還有一個基督徒醫生.
如果真的有事去看他的時候.
他收百多元.
這麼多年都收百多元.
已經省下很多了.
你明白這個道理嗎.
原來有時你為了自己保留.
但你會有另外的損失.
哈該書也是這樣.
我們自己享受天花板的房屋.
即是最好的木材.
我們自己用.
但神的加方良力不顧的時候.
你就會怎樣?.
身體有問題.
穿不暖.
顫顫震動.
吃了都沒有.
有蝗蟲.
咬了你的農作物.
你就已經有損失.
難道你不會確信神會祝福你嗎?.
所以求神來幫助我們.
魏國倫就是這樣說.
神不是呼召我們追求成功.
是呼召我們有忠心.
因為我們對神的期待太少.

$^{1161}$嘗試為神的事也太少.
這是很大的可悲.
以下這裡絕對不是我來帶你去.
想表現自己.
當我昨天再去看今天的講章.
神都提醒我.
最近有沒有奉獻?.
因為我每個星期去不同的教會.
我沒有固定的奉獻.
但我會將我的奉獻.
給香港有需要的教會.
如果你看回上兩章.
這兩章和這兩章.
昨天我奉獻給了葵涌堂和市民教會.
數目加起來應該是五千元.
但我預了十分之一.
應該有五萬元.
我講道是不會收到這麼多錢的.
我想告訴大家.
我不是只講理論給大家.
我今天講的這篇道.
也激勵了我自己.
提醒了我自己.
我要奉獻.
不是講道理.
另外牧師就很擔心.
問我是否教會友捐錢給黃福縣.
我說不是.
黃福縣至少有十一個家庭.
我看到他們有需要.
我每一個家庭.
都會給他們二千元的援助.
這是超過我的收入.
但我不想講理論給大家聽.
但我很豐富.
某程度上.
譬如下兩個星期.
我要去格拉斯哥.
我只講了一堂道.
車費不便宜.

$^{1201}$本來過去一個星期.
我應該去以色列.
但取消了.
但神祝福.
我二月取消以色列.
北愛爾蘭一位姐妹聯絡我.
我看回.
原來過去.
她對上一次是去年十月聯絡過我.
我們北愛爾蘭多了一個香港人團契.
牧師你可不可以來港島.
我們很想你來.
我二十七年都滿了.
我想最快也要二百年.
只有幾個星期.
但我想不是.
我取消了以色列.
我有個星期天.
不就行了嗎.
我上個星期已經去了北愛爾蘭.
然後去南愛爾蘭.
他們說我們這個團契剛新成立.
所以我們不會給你機票錢.
是你自己出的.
我說我沒有問題.
他甚至說我們連港元費都沒有.
但後來他有.
但我不是在乎你一點.
這都是我的奉獻.
但我是絕對沒有損失.
神給我很多祝福.
所以我想跟大家說.
絕對不是一個理論.
是一個真實.
有一天你不在這間教會.
你就離開.
奉獻也不用在這間教會.
牧師可能會罵我.
但你離開了就不用了.
就去另一間教會.

$^{1241}$奉獻給另一間教會.
就不可以說我還未穩定.
好了 完了.
記不記得我說.
一個馬戲班的故事.
快要結束了.
好了 另外一個節目是什麼呢.
就是走鋼線.
大概六層樓那麼高.
還沒有安全帶.
沒有安全帶很危險.
你見到那個人嗎.
在大陸.
掉下去了.
他拿著平行木.
即是棍子.
走鋼線.
下面是沒有安全網的.
來走.
於是每個人都很緊張.
他走完.
不只是這麼簡單.
我會再有一個.
我請我的拍檔出來.
我背著他來走這條鋼線.
嘩 不是吧.
嘩.
那個拍檔一出來的時候.
好像還未減肥的欣怡.
不是吧.
背著他走.
跌下來.
不跌死也宅死.
一步一步一驚心.
那些人走.
那些人很留心看.
不要看了.
終於走完了.
所有的人很自然地站起來.
大力鼓掌.

$^{1281}$然後這個人就說.
現在給大家一個挑戰.
有誰願意上來.
讓我背著走這條鋼線.
只要他願意的話.
我們就會有一些現金送給他.
不過仍然是沒有安全帶.
仍然是沒有安全網的.
接著那些人說 喂 你去吧.
神經病 要我去.
那你自己不去嗎.
原來全場的人.
沒有人願意走上去讓他背.
他們見到他曾經很成功地走過了.
很成功地背著那個較為健碩的人.
都過了.
但到最後你上來吧.
我背著你.
不要客氣.
明明我說這個故事.
是甚麼意思.
今天你聽了這篇道.
是對的.
這是聖經的教導.
但你去不去實踐呢.
由你自己決定.
求神幫助我們.
我再說.
不只是要大家去奉獻.
而是我想讓大家看到.
聖經給我們的教導.
我們應該怎樣實踐.
真理不是聽 是去行.
當你行真理.
你就會明白真理的真實.
我的生命 我的經歷.
就可以成為大家的提醒.
我非常豐富.
因為神給我很多祝福.
願意我們學習到.

$^{1321}$不只是金錢的奉獻.
是放上你的時間 精神 力量.
將身體現象當作活祭.
是神所喜悅的.
願以這篇信息成為大家的提醒.
祝福大家.
\newpage



\section{約書亞記 9:27}
\label{sec:032124}
\textbf{聖經人聖經事:自己吃虧也不改誓}~~~~黃克勤牧師博士
\newline
\newline
連結: \href{http://www.church.com.hk/acms/content.asp?site=cdc\&op=show\&type=product\&code=032124\&layout=sermon}{\texttt{032124}} ~~~~ 語音日期: 2026-05-14
\newline
\newline
\hyperref[sec:032123]{\small{< < < PREV SERMON < < <}}
~
\hyperref[sec:index_chronic]{\small{[返順時目]}}
~
\hyperref[sec:032125]{\small{> > > NEXT SERMON > > >}}
\newline
\newline
約書亞記 9:27
\newline
\begin{longtable}{cl}
\hline
\hline
章節 & 經文 (和合本修訂版)\\
\hline
9:27 & \begin{tabularx}{0.7\textwidth}{X} 那日，約書亞分派他們到耶和華選擇的地方，為會眾和耶和華的壇劈柴挑水，直到今日。 \end{tabularx} \\ \\
[1ex]
\hline
\hline
\end{longtable}
$^{1}$大家好,又來讀聖經.
生命樹為你準備的聖經人物.
聖經故事又來了.
這次仍然是約書亞記.
約書亞背景之下.
這次以色列人被人陷害.
被騙了.
吃了西瓜皮,香蕉皮,被香蕉皮陷害.
入世未深的以色列人與領袖陷害.
這是什麼呢?.
就是自己吃虧也不改世.
哇,厲害.
看看約書亞記第九章.
約書亞記第九章第一節.
約旦河西住山地高原.
並對住黎巴嫩山原大海一帶的諸王.
包括了黑人,阿摩利人,迦南人,比利西人.
希美人,耶布斯人的諸王.
聽見什麼呢?.
哇,聽見以色列人厲害.
百萬民族搖擁善戰.
經過了曠野打敗巴山王岳等等.
接著進來一叫,耶利哥城倒塌了.
接著叫城倒塌了.
前後夾擊,又有辦法,又好打.
糟糕了,現在繼續來了.
約旦河由北部黑門山一直流下來.
去到加里里湖,然後再下去死海.
然後這個約旦河這樣下來.
然後一片土地.
山地高,然後平地,然後去到海邊.
海邊就是地中海.
在迦南地來說,最北的就是黎巴嫩.
最南的就是接近埃及.
其實你說大不大,又說小不小.
因為由最北,如果是今天的以色列.
高速公路,高速公路很漂亮的.
高速公路,一百一十公里.
由北部一直踩車去到南部.
兩個邊界,我想是十五小時左右.

$^{41}$十多小時.
在車行,都幾遠,人行固然遠.
但又不是很大,不過很多民族住在那裡.
所以約旦河這邊是東,約旦河這邊是西.
所以約旦河西是一個山地,然後是海,地中海.
北部是黎巴嫩的山區,一直直南面.
約旦河西住在山地高原.
並且對著黎巴嫩山沿著大海一帶的珠黃.
因為這一片以色列土地,以前叫迦南地.
住了黑人,阿摩利人,迦南人,比利時人.
希美人,耶布斯人等等.
聽見以色列人來,那怎麼辦?.
他們進入我們的地區,不如我們一起圍堵他們.
因為沒有一個單獨自己救打.
於是珠黃聽見這事,就都聚集.
大家同心合意地要與約書亞和以色列人爭戰.
同心合意一起開集團的峰會.
欺騙的居民,其中有希美人.
希美人居住的地方,其中一個地方叫做基片.
基片的居民聽見約書亞.
向耶利哥和艾聖所行的事,嘩不得了.
他們就設詭計,為何要設詭計?.
因為上面珠黃就想,不是的,我們要夾在一起對抗.
否則我們單打獨鬥,我們沒有一個人救他鬥.
你看看耶利哥這麼大的城,艾成都被他搞成這樣.
百姓被他搶走,佔領了耶利哥和艾成的王都被殺.
艾成的王還被掛在樹上,被人看見,嘩死了.
其實很簡單,將王殺死掛在樹上,全世界都是這樣做.
就是示眾被人看見害怕,最害怕是誰呢?.
就是在做其他珠黃的王就害怕,因為怕自己有這樣一天.
於是不行,要跟他打.
希美人欺騙整個希美的王就去了開會.
國際峰會就要跟他打,跟以色列打.
但是這個小城市的欺騙這個地方,這個地區的居民.
我們希美王怎麼做,我們不管了,我們先保留性命,那怎麼做呢?.
這麼多人都不敢跟他打,又怕要一起.
難道我們在一個小地福音派兵嗎?更加不行.
但是我們有辦法,這裡有辦法,就設詭計.
那個詭計是什麼呢?就派幾個戲子出去演戲.
冒假充事者拿著舊口袋,舊口袋就是背著一個袋子.

$^{81}$拿著舊爛袋和破裂蓬佈的舊皮酒袋.
酒袋呢?新酒就要裝新皮袋,為什麼呢?.
因為新酒未發酵完畢,容易繼續有氣體膨脹.
所以要用新皮袋,新皮袋是新的皮,還有韌性.
可以頂得住那些氣,脹了的氣.
舊皮袋因為已經是進了,大了,脆了,老化了.
如果你放新酒進去,它再發酵一下,多一點氣就會推爆舊皮袋.
所以新酒放在新皮袋,舊酒放在舊皮袋.
舊皮袋就會穿,磨損,各樣硬化,老化.
所以一穿了怎麼辦?就要補了,把皮製成舊的,然後縫起來.
蓬蓬補補的舊皮袋,放在爐上面.
然後把穿過的舊鞋,舊鞋補過的舊鞋,穿了補的舊鞋.
穿在腳上,穿著舊衣服在身上,戴的餅乾的,還要發露.
以色列地很乾,所以如果你的餅放久了就會乾.
但是應該不會發露,應該不會發露,然後就會很乾.
如果放久了也發露,就已經很舊了,超超超舊了.
已經很久不值得了,長了露,長了霉.
他們到了吉甲營,就直接去找約書亞,跟以色列人說.
我們從很遠來的,現在求你跟我們納藥.
很遠來的嗎?納藥?納什麼藥?大家想有什麼藥?和藥.
以色列對這些希味人說,我又不認識你,又說來跟我們納藥.
我怕你原來就是住在附近,你不要扮遠行離道,我不認識你.
只怕你門是住在我們中間的,即是附近.
如果是這樣,我怎樣跟你納和藥?.
他們就跟約書亞說,其實我們是你的僕人,這句紅色字特別.
我們是你的僕人,我們來是要投靠你,是想做你的僕人.
我們想做你的僕人,所以和藥就是你不要殺我們.
我們想做你的僕人,而且我住到很遠,很有誠意來的.
你看我們破破爛爛,這樣就不行了.
所以我們來是想跟你搵啖茶飯,想跟你執條命仔,羞辱我們.
約書亞就問他們,你們是什麼人?在哪裡來的?.
他們就回答,我們是從極遠之地而來.
你想想,如果你是約書亞,你是那些領袖,你會怎樣想呢?.
你在哪裡來的?哦,我在黑龍江走來的,衣服都爛了,鞋都爛了.
哦,那你有個地方,你在哪裡來的?.
哎,我們很遠來的,都不講出地方名的,都恐防有詐.
不過當日的約書亞和領袖都說得很老實,又不懷疑的,很單純的.
哎呀,我們是從極遠之地來的,難道我在北極嗎?.
沒有問他,總之我們很遠,是因為聽到耶和華,你神的明星.
他就在埃及地,已經這樣過蘆葦海了,跟著曠野了.

$^{121}$然後耶利哥艾城,我們聽到這些.
然後聽到他在約旦河東,在阿摩利王,希什本王,西王.
還有在阿斯他祿他們,巴山王,鄂一切所做的事.
這些全部都是在未入迦南地,曠野擊敗西王和鄂.
哇,這些很出名的皇帝,之前的更出名的就是埃及法老.
都搞他不行,還淹死了,還在蘆葦海.
然後還在曠野殺死了這兩個王.
然後再來毀滅了耶利哥,毀滅了艾城.
我們聽到你們的神真厲害,我們從極遠之地來.
就是想做你的僕人,可能這樣說,我真的很崇拜你們.
所以我們一聽到就跑來,路遠跑來,遠到哪裡?很遠.
可能約書亞和他的眾守領同伴聽到之後.
哇,眉飛色舞,可能一下子就沒有小心了.
是我們的長老,就是那幫做戲子的那幫,使者那幫.
我們的長老,我們的領袖和我們的居民就說.
你們手裡拿著路上的食物去迎接以色列人.
帶些禮物送給他們,對他們說.
我們是要預備做你的僕人的了,現在求你與我們立約.
收我們做僕人,收我們幫你劈柴遞水,打水.
我們做你的僕人就可以了,不要殺我們就可以了.
我們出來是要你來這裡的日子,在家裡面.
你看看我們真的很誠意的,很苦心,為什麼呢?.
我在家裡面出來的時候,這些餅是新鮮的,熱的,你看看.
現在我們走到腳都瘸了,餅乾了,而且長了霉,你看看.
餅真的發霉了,這些皮袋,我們出來的時候是新皮袋來的.
脹鼓鼓的,現在走到我們這麼久,皮袋破裂.
衣服和鞋子都是滲了,軟路穿舊了.
唉,很可憐,很辛苦,很有誠意的.
第十四節,真的搞事了,以色列人竟然受了.
這個就叫做那些騙子,不怕舊,最重要是受.
以色列人受了他們的古惑,他們這些食物受了.
即是拿著那些發霉的餅,不知道有沒有牛肉乾.
唉,這隻我們帶來送給你的,小小禮物,不見得人很收家的.
不過都是我們一點心意來的,誠意,誠意,於是送給他們.
以色列人的領袖約書等等和領袖同伴竟然收了,信了.
就收了,所以下面第十四節的一點是什麼呢?.
就是我要講的第一點,並沒有求問耶和華.
就想一想,談一談,我們好不好接受他們呢?.
唉,這麼有誠意,不過說不出自己從何而來,有疑問.
唉,算了,看他們的樣子,出什麼古惑.

$^{161}$來做我們僕人也好,無端端有姐姐,有哥哥用也不錯.
不如收他吧,算了,算了,那些食物不算什麼.
我們一場誠意收吧,自己決定了,並沒有求問耶和華.
唉,這裡出事了,第十五節,於是第二點.
約書跟我們講和,跟他立約,會眾的首領也向他們起誓.
用他們活著,做我們的僕人,那就算你撿回一條命.
哇,多厲害啊,不用死,跟我們,做我們僧,做僕人就行了.
多厲害啊,行了,那些爛鬼食物放在一邊.
總之回去預備,我們什麼時候到你那裡,就行了,不會殺你們.
不過挺有趣的,最絕的一點在哪裡呢?.
是沒有問人家從何而來,來幾個代表做我們僕人.
那成族人在哪裡?沒有說,那行,去到你那裡,做我們僕人.
去哪裡?沒有說,好了,他們都走了.
以色列人與他們立約之後,過了三天才聽到他們是近鄰.
住在以色列人中間的,是近鄰,很近的,原來很近.
差不多跟我們一齊,哇,過了幾天才知道,傳回來.
哎呀,厲害了,幾個棄子哄到我們就死了.
以色列人繼續起行的時候,第三天就到了他們的城一.
就是基片,基斐拉比祿,基列耶林,主要是基片.
在這些地方,原來走三天就到了,又說極遠,裝模作樣.
哇,不是啊,去到了,歡迎,歡迎,不要殺我們,不要殺我們.
我們不會打,因為我們都是你的僕人.
因為會眾首領已經指著耶和華以色列的神向他們起誓.
所以以色列人就骨擊殺他們.
傳會眾就向首領發願言,喂,神叫我們進來是要擊殺他們的.
為什麼你跟他們立約的,要不殺他們還起了誓.
我們發誓不會殺你這樣,納了和約的.
搞什麼的,你們怎麼做領袖的,有沒有腦的.
哎呀,慘了,眾首領對傳會眾說,真是的,這句我加一句.
不好意思,我們已經指著耶和華以色列的神向他們起誓.
現象我們不能害他們,因為我們要起了誓,發了誓.
在神面前發誓,所以我們不能夠不做.
我們要如此待他們,就是不殺他們,用他們活著.
第三點要講的就是,免得有憤怒因我們所起的誓會臨到我們.
我們是做錯了,就是你們罵我們,我們是對的,我們做錯了.
其實當時沒有求問耶和華,我們見到這樣就說好吧.
當時都是一片好心的,見他們這樣.
但是我們雖然幾十歲人,一時就入世未深,哎呀,擅親了,踩錯腳了這次.
但是我們已經向耶和華起誓,要容他們傳活.
因此我們不能夠殺他們了,就免得我們破壞自己向神起的誓言.

$^{201}$反而招致神的憤怒.
於是首領又對會眾講,我們唯有要用他們傳活了.
不過他們是答應做我們的傭工的,做我們僕人.
於是他們就要為我們全會眾做劈柴挑水的人.
正如領袖對他們所說的話,我們給他們生活.
不過他們要幫我們做事,做牛做馬,劈柴挑水.
約書就召了他們來,對他們說,為什麼欺控我們?.
我們來你們甚遠?為什麼這樣講?.
其實你們住在我們附近旁邊,現在你們是被詛咒了.
你們中間的人必斷不了做勞僕,為我們神的殿堂做劈柴挑水的人.
為什麼你欺控我們?約書就叫他們來,為什麼你欺控我們?.
又說很遠,其實跟我們一起住的,三天就來到這裡,又說很遠.
走到鞋都爛,又說走到皮袋都爛了,為什麼會這樣?.
不過,照你所說的,你要永遠不斷做我們劈柴挑水的僕人.
特別為我們神的殿堂做服侍.
他們就回帝洛書說,因為有人實在告訴你們的僕人.
耶和華你們的神曾經吩咐僕人摩西要把全地賜給你們.
並且在你們面前滅絕借地的一切居民.
所以我們作為這件事的緣故,我們很怕死,很怕喪命.
就行了這件事,我們都是想撿回一條命,讓一條生路走走,大哥,不要殺我們.
我們都是想活著,我們已經想了,交換就是我們做你一世的奴隸.
我們成族人做你的奴隸,你們以後那些劈柴挑水的粗工我們做.
那就留我們一條命,好過死,因為我們知道你的神多厲害.
而且他又叫你們要殺我們,甚至現在我們在你的手中.
你們怎樣待我為善為正,你們就怎樣做.
我都沒得說了,因為我們一定不夠你打,你要怎樣做就怎樣做.
你認為好就做,我們就想不用死,就是為你打工.
於是約書亞就這樣待他們,好了,我就讓你們不被以色列人我們的手所殺.
那就沒殺他們,那做什麼呢?.
當日約書亞就令他們在耶和華所需要選擇的地方.
就是為會眾和耶和華的譚做劈柴挑水的人,直到今天.
於是欺騙人等等附近的人,來扮個棄子所代表的百姓.
為以色列人打工,做奴隸,點他,特別為神的譚做好功夫.
那即是怎樣?即是被人騙了,你騙我我就反口,你騙我的就殺.
約書亞說,因為我們是信神的,我們向他們起誓.
所以我們要守誓言,講信用守誓言.
請大家看詩篇15篇第4節,第4節講什麼呢?.
敬畏神的人是眼中藐視匪類,卻尊重敬畏耶和華的人.
留意這兩句,他發了誓,雖然自己竟然被騙是吃虧,也不更改.
自己吃虧也不更改,無論做生意做人.

$^{241}$我插錯腳,答應錯了價錢,就這個價錢,我在一些外面的地方.
即不是香港,在一些地方,那些地方很公道.
就是在超市裡面,在香港就沒有了.
在超市裡面見到這件貨品,放在這裡,寫著標誌多少錢.
我就拿這件貨品拿去付錢,很便宜,我就拿這件貨品拿去付錢.
現在我們知道,這麼便宜,真的沒有理由.
那我快點拍一張相,拿這件貨品去付錢.
他一掃過機,就說這件貨品20元.
我就說不是,那裡標誌著5元,不是吧,機是說20元.
那可能你標誌錯了,不能說擺錯,是真的那件貨品.
標誌著,5元而已,我拍了,是喔,那他就快點走去看看.
是喔,5元,那他就真的用5元來賣給我的.
他用5元來賣給我的,不可以說不行,我擺錯了,標誌錯了.
我們同時標誌錯了,不行,20元,你5元,我不賣了.
不行的,他一定要賣,這樣的,我們也要.
那香港就沒有的,一句標誌錯了,機是20元的.
你不給20元,我就不賣給你了,不可以給你了.
標誌著5元而已,標錯了都說,哎呀.
這裡如果我賣東西標錯價,我都要照那個價格.
自己吃虧也不更改,如果我是貴買貴了,哎呀,價錢開得差了.
用這個錢給我們買,哎呀,這次虧了,自己吃虧了,行不行.
這個就是說,自己吃虧也不更改,這個是敬畏神的人會做的.
因此,我在這裡和大家說,回到今天的題目.
很長的,今天說了,這個題目是什麼呢?.
自己吃虧也不改世,為什麼呢?.
因為當我向神納了願,以前就起世了.
現在又說不用起世了,新的又說不要起世了.
講真說話就可以了,不用用起世來證實你所說的話.
是另一個真理來的,如果我承諾過的.
就算我吃虧,我都要做的,要履行.
這個是一個正直敬畏神的表現,幫助我們做幾件事.
第一,不要自高,人家捧一捧就暈了.
第二,就是小心凡事求問耶和華.
這些人信不信得過呢?求你給我心中的平安.
另外,小心考察,這些事是否真的,先證實一下.
原來住在我們當中三日路程又說極遠,我就知道了.
而且起了世,不可以修改,要跟著承諾來履行的.
是聖經教導來的,有祝福的,我們一起祈禱.
我們見到約書亞和銀色列的領袖,這次是踏錯腳.
被人家的棄子騙到,還發誓要保護他們的性命.

$^{281}$但是後來發覺到自己傻了.
雖然被百姓罵,不過他們說我們都沒有辦法了.
既然在神的名下起世,我們要遵守.
我相信這個傻傻的,繼續要遵守,不會反口也是神所寬喜的.
而且詩篇也說,發了世,不要更改.
幫助我們守著自己的信用與承諾.
我們知道這是蒙福的,是敬畏你的人的表現.
幫助我們有信心與胸襟,隨聽禱告,奉耶穌基督名求,阿們.
\newpage



\section{路加福音 15:11-32}
\label{sec:032125}
\textbf{當層出不窮的愛,遇上深不見底的愛}~~~~冼嘉琪傳道
\newline
\newline
連結: \href{http://www.church.com.hk/acms/content.asp?site=cdc\&op=show\&type=product\&code=032125\&layout=sermon}{\texttt{032125}} ~~~~ 語音日期: 2026-05-10
\newline
\newline
\hyperref[sec:032124]{\small{< < < PREV SERMON < < <}}
~
\hyperref[sec:index_chronic]{\small{[返順時目]}}
~
\hyperref[sec:032126]{\small{> > > NEXT SERMON > > >}}
\newline
\newline
路加福音 15:11-32
\newline
\begin{longtable}{cl}
\hline
\hline
章節 & 經文 (和合本修訂版)\\
\hline
15:11 & \begin{tabularx}{0.7\textwidth}{X} 耶穌又說：「一個人有兩個兒子。 \end{tabularx} \\ \\ \relax
15:12 & \begin{tabularx}{0.7\textwidth}{X} 小兒子對父親說：『父親，請你把我應得的家業分給我。』他父親就把財產分給他們。 \end{tabularx} \\ \\ \relax
15:13 & \begin{tabularx}{0.7\textwidth}{X} 過了不多幾天，小兒子把他一切所有的都收拾起來，往遠方去了。在那裡，他任意放蕩，浪費錢財。 \end{tabularx} \\ \\ \relax
15:14 & \begin{tabularx}{0.7\textwidth}{X} 他耗盡了一切所有的，又恰逢那地方有大饑荒，就窮困起來。 \end{tabularx} \\ \\ \relax
15:15 & \begin{tabularx}{0.7\textwidth}{X} 於是他去投靠當地的一個居民，那人打發他到田裡去放豬。 \end{tabularx} \\ \\ \relax
15:16 & \begin{tabularx}{0.7\textwidth}{X} 他恨不得拿豬所吃的豆莢充飢，也沒有人給他甚麼吃的。 \end{tabularx} \\ \\ \relax
15:17 & \begin{tabularx}{0.7\textwidth}{X} 他醒悟過來，就說：『我父親有多少雇工，糧食有餘，我倒在這裡餓死嗎？ \end{tabularx} \\ \\ \relax
15:18 & \begin{tabularx}{0.7\textwidth}{X} 我要起來，到我父親那裡去，對他說：父親！我得罪了天，又得罪了你， \end{tabularx} \\ \\ \relax
15:19 & \begin{tabularx}{0.7\textwidth}{X} 從今以後，我不配稱為你的兒子，把我當作一個雇工吧。』 \end{tabularx} \\ \\ \relax
15:20 & \begin{tabularx}{0.7\textwidth}{X} 於是他起來，往他父親那裡去。相離還遠，他父親看見，就動了慈心，跑去擁抱著他，連連親他。 \end{tabularx} \\ \\ \relax
15:21 & \begin{tabularx}{0.7\textwidth}{X} 兒子對他說：『父親！我得罪了天，又得罪了你，從今以後，我不配稱為你的兒子。』 \end{tabularx} \\ \\ \relax
15:22 & \begin{tabularx}{0.7\textwidth}{X} 父親卻吩咐僕人：『快把那上好的袍子拿出來給他穿，把戒指戴在他指頭上，把鞋穿在他腳上， \end{tabularx} \\ \\ \relax
15:23 & \begin{tabularx}{0.7\textwidth}{X} 把那肥牛犢牽來宰了，我們來吃喝慶祝； \end{tabularx} \\ \\ \relax
15:24 & \begin{tabularx}{0.7\textwidth}{X} 因為我這個兒子是死而復活，失而復得的。』他們就開始慶祝。 \end{tabularx} \\ \\ \relax
15:25 & \begin{tabularx}{0.7\textwidth}{X} 「那時，大兒子正在田裡。他回來，離家不遠時，聽見奏樂跳舞的聲音， \end{tabularx} \\ \\ \relax
15:26 & \begin{tabularx}{0.7\textwidth}{X} 就叫一個僮僕來，問是甚麼事。 \end{tabularx} \\ \\ \relax
15:27 & \begin{tabularx}{0.7\textwidth}{X} 僮僕對他說：『你弟弟回來了，你父親因為他無災無病地回來，把肥牛犢宰了。』 \end{tabularx} \\ \\ \relax
15:28 & \begin{tabularx}{0.7\textwidth}{X} 大兒子就生氣，不肯進去，他父親出來勸他。 \end{tabularx} \\ \\ \relax
15:29 & \begin{tabularx}{0.7\textwidth}{X} 他對父親說：『你看，我服侍你這麼多年，從來沒有違背過你的命令，而你從來沒有給我一隻小山羊，叫我和朋友們一同快樂。 \end{tabularx} \\ \\ \relax
15:30 & \begin{tabularx}{0.7\textwidth}{X} 但你這個兒子和娼妓吃光了你的財產，他一回來，你倒為他宰了肥牛犢。』 \end{tabularx} \\ \\ \relax
15:31 & \begin{tabularx}{0.7\textwidth}{X} 父親對他說：『兒啊！你常和我同在，我所有的一切都是你的； \end{tabularx} \\ \\ \relax
15:32 & \begin{tabularx}{0.7\textwidth}{X} 可是你這個弟弟是死而復活，失而復得的，所以我們理當歡喜慶祝。』」 \end{tabularx} \\ \\
[1ex]
\hline
\hline
\end{longtable}
$^{1}$母親節快樂,大家.
我們先開始,再一次祈禱.
願神的話語在我們心裡,預備在我們心裡.
親愛的主人,多謝你.
再一次很感恩.
你讓我們可以一齊去敬拜.
天父,我們知道.
每一次可以敬拜,都不是必然的事.
多謝你,讓我們可以有實體的聚會.
可以一齊有這個時間.
分別為聖,一齊去敬拜你.
預備在我們心裡,領受你的話語.
天父,多謝你.
我們同心禱告,奉主明求,Amen.
大家好.
在去年年底的時候.
有一個新的研究.
在南韓有一個新的發明.
即是今年年初.
那裡是在韓國的科學研究院.
宣佈它這個新的發明.
是一個止血的粉末.
可以一秒內止血.
這個研究是很實際的.
用途是它預備.
譬如說打仗,一些災難性.
嚴重交通意外的時候.
是可以立刻在大動物流血的時候.
噴下去,一秒可以止血.
停止了血液.
止血粉可以一秒內止血.
但如果一個人愛到心痛的時候.
可不可以立刻停止呢?.
可不可以呢?.
相信我們都曾經試過.
愛子女愛到心痛.
我們都試過愛父母愛到心痛.
如果有一支感情的止血粉.
讓你噴.
就可以在你心痛的時候.

$^{41}$立刻讓你噴一噴.
你就停止去愛.
你會不會噴呢?.
為何會呢?.
為何不會呢?.
因為我們都知道.
愛是要經歷的.
雖然愛有時會痛.
但我們要經歷.
那份愛是很真實的.
人的愛很複雜.
但很真實.
所以我形容人的愛是.
層出不窮的愛.
層出不窮.
即是很多樣化.
一層一層彼此交織.
但有時會是矛盾.
人的愛有時很矛盾.
人的愛有時很溫柔.
有時很傷人.
那句說話叫一啖砂糖.
一啖甚麼?.
不出得街.
有時人的愛是怎樣矛盾呢?.
有時很理性.
很成熟.
但有時又很傻氣.
做些很傻的事.
早陣子上一年.
好像在美國有一群男士.
因為想哄女生開心.
他們學甚麼呢?.
一起學.
他們學幫人紮邊.
讓女生可以開心一點上學.
滿足一點.
開心一下.
這個熱潮不是只在美國.
接著風波到澳洲.

$^{81}$歐洲也有.
人的愛有時很認真.
很成熟.
有時很傻氣.
人的愛很複雜.
加上不同的個性.
不同的背景.
不同的經歷.
令人的愛更加複雜.
所以層出不窮的愛.
在一個家庭裡.
形成了張力.
但亦令一個家庭很豐富.
我們今天一起看看.
《聖經》裡面的故事.
這個故事是家庭故事.
某程度上是我們每個人發生的故事.
我們一起來看看.
《路加福音》十五章.
十一到三十二節.
今天我們不會一次過讀完經文.
但我會大概講個總結出來.
裡面講到一個浪子的比喻.
可能大家都很熟悉.
如果看《聖經》.
它的標題是浪子的比喻.
為何會說比喻呢?.
因為我們看到十五章.
耶穌一連棄講了三個比喻.
我們講浪子的比喻.
裡面一個故事是講甚麼呢?.
耶穌說有一個父親.
有兩個兒子.
大兒子和小兒子.
小兒子要求爸爸提早將家產分給他.
他拿了家產後幾天.
帶了所有東西離開.
遠走高飛.
去到遠方花光所有錢.
揮霍的時候.

$^{121}$誰知遇上饑荒的地方.
他沒辦法.
逼於無奈去打工.
打甚麼呢?.
去幫人看珠.
猶太人很討厭珠.
要去看珠.
很餓的時候.
他看著珠就想起.
真的想連豬哨都吃.
誰知連豬哨都不能吃.
那一刻他突然想起.
我爸爸連故宮都不用吃這些東西.
為何我不回爸爸那裡呢?.
於是他就想起.
我寧願回去求爸爸請我做工人.
好過在這裡.
於是他就滴起心肝回去.
回去的時候.
他爸爸怎樣呢?.
他爸爸離遠看見他.
已經衝過去.
抱著他 親他.
原諒了他小兒子.
衝過去.
然後大兒子怎樣呢?.
當時他大兒子見到他的時候.
就很生氣.
因為他覺得.
為何弟弟敬拜父親家.
父親也可以這樣原諒他呢?.
但他自己一直陪著父親.
就沒有這些事.
他爸爸跟他說了兩次.
解釋同一句說話.
因為你這個弟弟.
是死而復生.
失而復得.
因為你弟弟是死而復活.
失而復得.

$^{161}$這個故事.
有沒有像我們家庭裡的故事呢?.
如果我們有兄弟姊妹.
我們不多不少.
都會有少少這些淩亂.
我們一起看下去.
看看小兒子的愛是甚麼.
小兒子的愛是一個自我型的愛.
小兒子跟父親說.
父親請把我應得的那份家產給我.
父親未過世就要分家產.
無論在當時.
無論在現在.
都是一個大逆不道的要求.
是一個非常非常大逆不道.
就好像我們現在的兒子.
跟我們說.
你把房子寫在我的名字.
反正你將來死了也是我的.
大家接受嗎?.
好像有少少道理.
但這句說話是很無情.
很無情.
這個小兒子不是跟父親說.
父親我很想去玩.
你給我一點錢.
他不是這樣說.
他是說那些家產是我的.
我應得的.
現在就分了它.
在猶太人.
大兒子.
第一個兒子.
是會分到雙倍的財產.
他家裡有兩個兒子.
一大一小.
即是大兒子分到的.
是三分二.
小兒子是三分一.
小兒子這一刻說.

$^{201}$將我應得的那份家產.
給我.
這個小兒子已經計算過.
三分一的家產.
他都足夠可以.
去風流快活.
足夠可以過他自己想過的生活.
他計算過了.
即是說.
這個小兒子愛不愛爸爸呢?.
很難說.
但很肯定的是.
他愛自己.
他愛他自己想過的生活.
多過愛他爸爸.
應得.
這個字很深層的意思.
即是說.
小兒子不是說.
只是要錢那麼簡單.
他是說.
這些是我應得的.
這句說話裡面.
即是說.
這些我認為是我配得擁有.
我配得擁有.
I deserve it.
這個關係裡面.
很多時候破裂.
不是因為.
有缺乏.
有不夠開始.
很多時候.
這個關係裡面的破裂.
有人覺得.
我配得.
我應該擁有更多.
小兒子其實.
在這個家庭裡面.
一早已經有.

$^{241}$很多的愛.
被供應.
但小兒子他覺得.
我應該擁有更多.
我有權利.
應該得到更多.
要自己決定.
他的人生要怎樣走.
今天這個世代.
其實我們很容易培養到.
這東西我稱為「應得感」.
雖然這不是科學詞語.
我會形容這東西.
叫「應得感」.
「應得感」是怎樣呢?.
有一個年輕人.
他剛剛考完車牌.
很開心的跟爸媽說.
爸媽我終於考到車牌了.
我有架健身車你可不可以買給我.
媽媽就說.
你剛剛考完車牌.
不要買這麼漂亮的車.
然後小兒子就說.
但是.
你猜他怎樣回答.
但是.
為什麼你跟爸媽可以有這麼漂亮的車.
爸媽就不明白.
為什麼我這麼辛苦.
這麼辛苦上班.
開一部好一點的車.
都不行.
兒子又不明白.
為什麼家裡有的.
又不可以給我.
這就是「應得感」.
當有的時候.
理所當然.
當沒有的時候.

$^{281}$覺得很委屈.
為什麼沒有.
有兒童發展的研究.
指出.
在物質豐富.
科技發達.
和可以有即時滿足感.
的環境裡.
小孩子很難.
培養到感恩的心.
物質豐富.
科技發達.
容易滿足.
他的環境.
這三樣東西是不是很熟悉.
是不是很熟悉.
好像是不是我們這裡的環境.
很難.
有感恩的心.
所以教會.
為什麼兒童試工會叫.
Grateful Kid.
因為我們祈禱.
求神幫助我們.
可以讓孩子.
認識到神.
可以幫他.
自小培養到感恩的心.
在神裡有信心.
感恩去長大.
因為我們相信.
能夠有感恩的心.
是回應神一個.
恩典最自然.
最自然的流露.
我們很想從小教小朋友.
教他們.
什麼是恩典,什麼是感恩.
我們相信神的.
恩典是豐富的.

$^{321}$但是神的恩典.
不是理所當然.
我們很想很想教小孩子.
所以我們需要.
每一個大人.
我們需要全教會.
為小朋友一起去祈禱.
祈禱托著這個試工.
在早前.
我看到一個新聞.
我發現.
應得感最深最深.
的問題.
不是對物質的追求.
不是不停不停.
想要買某些物質.
這麼簡單.
最深層的問題.
是應得感培養到.
一個很深很深.
很深很深的自我中心.
一個滿有.
應得感的人.
慢慢將自己看得最大.
他會怎樣呢?.
我覺得什麼對就是對.
我覺得我有權.
走我自己想選擇的路.
我覺得我有權.
過我怎樣的生活.
覺得自己有權.
非常自我中心.
那個新聞.
看到什麼呢?.
說到有一個.
同性戀的.
情侶.
有同性戀的男士.
抱著一個嬰兒.
問嬰兒.

$^{361}$你想要.
爸爸.
還是爹爹.
一樣.
你想要爸爸還是爹爹.
嬰兒.
開始哭.
開始嘮叨.
媽媽.
接著那個.
男士還覺得.
很好笑.
你想要爸爸.
還是爹爹.
No媽媽.
他覺得好像很好笑.
但其實他在教育.
那個嬰兒什麼呢?.
他在教育那個嬰兒什麼呢?.
這個世代就是.
顛倒是非的世代.
一個顛倒是非的世代.
人的自我中心扭曲了正常的愛.
扭曲了正常的愛.
所以我們預計.
越來越多.
新一代的小孩.
新一代的嬰兒.
會在一些扭曲了的愛裡成長.
扭曲了的舊約裡成長.
扭曲了正常的愛.
小孩回來的時候.
他裡面的比較.
他裡面的苦讀.
他裡面的委屈.
就一次過爆出來.
一次過爆出來.
他跟他爸爸怎麼說.
他說:你看,我服侍了你這麼多年.
從來都沒有違背過你的命令.

$^{401}$我們有很多人.
一直照顧著爸爸媽媽.
一直照顧著子女.
好像有些夾心餅.
一直撐著家庭,一直服侍很多人.
外面看起來很堅固,很忠心.
但是裡面慢慢枯乾了.
都沒有人知道.
有一些成長後的子女.
很疼爸爸媽媽.
陪伴他們看醫生.
幫他們看文件.
整理家具.
在電話裡.
經常盯著電話.
因為很害怕兩老打電話來.
看著電話.
有時不小心錯過電話.
就很內疚.
心裡面很掛念父母.
很掛念年長的父母.
很想做得更好.
但是有時無奈工作.
無奈要照顧自己的家庭.
所以有時可能對兩老發脾氣.
發完之後又會很內疚.
弟兄姊妹.
有時愛令人覺得有壓力.
有壓力的時候又覺得累.
有時累的時候又覺得內疚.
為何照顧家庭都覺得累.
一方面知道父母年長.
另外一方面.
又覺得自己做得不夠多.
很多虧欠.
明明很累.
但又要繼續.
弟兄姊妹有時最累.
不是因為不愛.
是因為想愛得更加好.

$^{441}$更愛他.
所以外表撐著.
裡面就枯乾了.
這是人的愛.
人一直愛的時候.
就算最真最真的愛.
人的愛會怎樣呢.
會疲累.
會受傷.
慢慢會.
burn out.
慢慢會形成.
苦毒.
人不是不愛.
很想愛,是真愛.
不過.
人的愛有限.
所以人的愛會.
層出不窮扭曲.
在香港.
早前有一個女孩.
自從她11歲開始.
被她繼父性侵犯.
一路如是者好幾年.
頻繁性侵犯她.
到了這個女孩長大成人.
離開了這個家.
過了一段時間.
才去報案.
別人問她.
為何當初不早出聲.
你猜這個女孩怎樣回答.
這個女孩說.
她清楚記得當年.
她親生父母離婚.
她清楚記得當年.
她親生父母離婚.
那種傷心.
她不想在她母親.
第二段婚姻裡.

$^{481}$再次看到她母親那麼傷心.
所以她一路不出聲.
這個女孩用了自己的方式.
背起了一個錯誤.
而且不應該是她背負的責任.
用錯了方式去愛她母親.
這個世代就是這種.
層出不窮.
扭曲了的愛.
一層一層的交雜.
一層一層的變形.
靈姐妹.
我們可以做什麼呢?.
我們要做什麼呢?.
所以我們不是要從一個不完美的人身上去找完美的愛.
我們知道一個污染了的水龍頭.
我們不可能有乾淨的水走出來.
我們不可能.
我們真正需要的.
是需要找一個愛的源頭.
最乾淨的源頭.
那個就是神的愛.
一份深不見底的愛.
讓神的愛去囂觀我們.
神的愛.
這個父親在這個比喻裡.
如何去回應他的兒子呢?.
聖經裡面說.
「於是他起來往自己父親那裡去」.
「當他還在遠處的時候,父親看見他,就動了憐憫之心,奔跑進去」.
在當時猶太人的文化裡.
父親是不會跑的.
代表了父親放棄了父親的尊嚴和尊貴.
跑去擁抱他這個不值得愛的兒子.
弟兄姊妹.
各位朋友.
神都是這樣跑向每一個需要祂的人.
兒子未解釋.
未跟父親道歉.
由他之前是怎樣的?.

$^{521}$他之前是.
I deserve it.
我應得的甚麼都.
現在去到他看到.
他自己不配.
由他覺得.
甚麼都是應得.
現在看到他自己是不配.
這個就是.
屬靈裡面的醒悟.
人承認不配的時候.
就可以重新看到神的恩典.
人承認自己不配的時候.
就可以重新體會到神的恩典.
在一個家庭裡面.
剛剛說.
是層出不窮.
而又扭曲了的愛.
因為一個家庭裡面.
的愛的狀況.
不只是.
爸爸媽媽之間的關係.
不只是.
子女之間和父母之間.
不只是.
爸爸媽媽和他爸爸媽媽之間的關係.
是全部.
炒在一起.
這是人扭曲了的愛.
所以.
在座的父母.
如果我們當中.
發現我們和子女.
好像有些狀況.
如果我們發現和上一代.
好像有些狀況.
都想鼓勵大家.
我們嘗試先從.
我們的婚姻裡面.
去鞏固我們的關係.

$^{561}$因為在婚姻裡面.
一起走向神.
這個關係是很重要.
是我們家庭裡面的一個奠基石.
小小的廣告.
如果我們沒有參加過.
一些夫婦的迎會.
鼓勵夫婦.
要去參加一下.
因為夫婦的迎會.
不只是.
說一個夫婦的關係.
我們是.
教導大家.
在夫婦的關係裡.
去重新發現.
你們婚姻裡面.
神的恩典.
這個很重要.
由最核心的去恢復.
所以想知道更多.
有關夫婦的迎會.
可以來找我們.
回到這裡.
剛剛說到.
細仔的狀況.
很明顯.
大仔有什麼錯?.
大仔真正的問題.
是他說話的時候.
已經暴露了.
大仔說.
我服侍你這麼多年.
你從來都沒有給過我什麼.
這句說話就是說.
我這麼好.
我應該得到更多.
我這麼好.
我比我弟弟好.
還有.

$^{601}$在聖經裡面.
大兒子還說了一句話.
他特意.
跟父親說.
他說.
你這個兒子和昌記一起.
花光你的錢.
他特意.
要數弟弟的黑歷史.
在哥哥和他爸爸說的話裡.
他說.
你看.
中文和合本.
看不到這個字.
李漢.
所以環球聖經的本.
譯了出來.
英文有些也譯了出來.
就是李漢.
你看.
他心情是在發晦氣.
他很生氣.
他和他爸爸說.
你看.
這個大兒子.
他認為.
恩典.
臨到一個.
他看來.
不配擁有恩典的人身上.
這個大兒子.
覺得.
爸爸的恩典臨到一個.
不配擁有.
恩典的人身上.
這是最大的問題.
覺得別人不配.
擁有這個恩典.
其實他自己.
也是不配.

$^{641}$小兒子.
用他活躍.
應得感.
離開了.
其實大兒子.
他也是自義做好人.
他的優越感.
他的心.
其實都已經離開了他爸爸.
今天我們像.
小兒子多還是大兒子多呢?.
我們像小兒子多還是大兒子多呢?.
福音.
最奇妙的地方.
是神從來都不是.
給我們應得的.
神從來.
不是給我們應得的.
神給我們的是我們不配得的.
不配得的.
一切都是恩典.
我們本來.
就是不配得給神愛.
本來.
就不配得給神救贖.
真正的愛.
是什麼呢?.
真正的愛.
是有真理的.
真正的愛.
不是放縱的.
有是非黑白.
真正的愛.
是沒有懼怕的.
不害怕的.
真正的愛.
是沒有綑綁的.
是讓人得自由.
所以 活躍的人需要悔改.
有優越感的人.

$^{681}$都要悔改.
因為.
我們覺得自己.
如果我們越有應得感.
越有優越感的時候.
我們就看不到.
聽不到.
感受不到神的愛.
弟兄姊妹 我們不是因為做得好.
所以才被愛.
我們是因為.
有神愛我們.
所以我們才懂得愛.
人的愛.
是層出不窮.
而且.
不盡地扭曲的愛.
神的愛是心不見底.
是順傳的愛.
所以我們不要在.
扭曲了的愛裡.
去尋求愛.
我們要努力.
回到這個愛的源頭.
我們要回到.
神的愛裡.
我們一起低頭祈禱.
在這個彎曲薄茂的世代裡.
你彰顯.
你真正的愛.
給我們知道.
父啊.
我們知道.
這個世代有很多.
很多破碎了的.
關係.
很多破碎了的家庭.
我們知道.
我們很多人覺得.
很無奈.

$^{721}$很傷心.
有時很乏力.
很疲累.
很無奈地.
接受這些破碎.
父啊.
多謝你再一次.
提醒我們.
你的愛.
是可以修復.
可以醫治.
可以讓我們重新.
在你裡面.
得到被愛.
父啊.
如果我們當中.
有一些家庭.
有一些弟兄姊妹.
他們很需要.
很需要更深的明白.
你的愛是甚麼的愛.
父啊我求你.
特別去跟他們講說話.
父啊.
我們相信.
你的愛.
是豐豐富富.
是高過我們.
所想像.
深過我們.
可以理解.
父啊.
我求你.
讓你的愛囂觀.
讓你的愛.
醫治.
讓你的愛充滿我們每一個.
天父多謝你.
多謝你聽我們禱告.
我們同心禱告.

$^{761}$主明求.
Amen.
\newpage



\section{以弗所書 1:15-23}
\label{sec:032126}
\textbf{保羅為教會禱告 (華/粵)}~~~~姚臣泰長老
\newline
\newline
連結: \href{http://www.church.com.hk/acms/content.asp?site=cdc\&op=show\&type=product\&code=032126\&layout=sermon}{\texttt{032126}} ~~~~ 語音日期: 2026-05-10
\newline
\newline
\hyperref[sec:032125]{\small{< < < PREV SERMON < < <}}
~
\hyperref[sec:index_chronic]{\small{[返順時目]}}
~
\hyperref[sec:032127]{\small{> > > NEXT SERMON > > >}}
\newline
\newline
以弗所書 1:15-23
\newline
\begin{longtable}{cl}
\hline
\hline
章節 & 經文 (和合本修訂版)\\
\hline
1:15 & \begin{tabularx}{0.7\textwidth}{X} 因此，我既然聽見你們對主耶穌有信心，對眾聖徒有愛心， \end{tabularx} \\ \\ \relax
1:16 & \begin{tabularx}{0.7\textwidth}{X} 就不住地為你們感謝神，禱告的時候常常提到你們， \end{tabularx} \\ \\ \relax
1:17 & \begin{tabularx}{0.7\textwidth}{X} 求我們主耶穌基督的神，榮耀的父，把那賜人智慧和啟示的靈賜給你們，使你們真正認識他， \end{tabularx} \\ \\ \relax
1:18 & \begin{tabularx}{0.7\textwidth}{X} 照亮你們心中的眼睛，使你們知道他呼召你們來得的指望是甚麼，他在聖徒中所得榮耀的基業是何等豐盛， \end{tabularx} \\ \\ \relax
1:19 & \begin{tabularx}{0.7\textwidth}{X} 並知道他向我們這些信的人所顯的能力是何等浩大，這是照他的大能大力運行的。 \end{tabularx} \\ \\ \relax
1:20 & \begin{tabularx}{0.7\textwidth}{X} 這大能曾運行在基督身上，使他從死人中復活，又使他在天上坐在自己的右邊， \end{tabularx} \\ \\ \relax
1:21 & \begin{tabularx}{0.7\textwidth}{X} 遠超越一切執政的、掌權的、有權能的、統治的和一切有名號的；不但是今世的，連來世的也都超越了。 \end{tabularx} \\ \\ \relax
1:22 & \begin{tabularx}{0.7\textwidth}{X} 神使萬有服在他的腳下，又使他為了教會作萬有之首； \end{tabularx} \\ \\ \relax
1:23 & \begin{tabularx}{0.7\textwidth}{X} 教會是他的身體，是那充滿萬有者所充滿的。 \end{tabularx} \\ \\
[1ex]
\hline
\hline
\end{longtable}
$^{1}$我們又再次見面了,母親節快樂,記得崇拜之後,你們做母親的,有東西吃,做父親的要幫忙下..
我們再回到義父所述的時候,我們看到神自己對教會的心意..
我們再看二弗所書上,我們再次看到神對教會的心意..
那我們在禱告上面,保羅也有一些給我們的提醒..
特別在禱告方面,保羅也會對我們有一些提醒..
我不知道你的禱告如何?.
我不知道你自己的禱告是怎樣的?.
如果現在我們叫你們閉目禱告的時候,給你一分鐘你禱告,是不是能夠講?.
請你自己帶領我們,因為這裡好像有什麼東西還沒有做..
你知道我的意思嗎?.
我希望大家知道我的意思是什麼..
我們常常禱告不專心,容易分心..
我們常常禱告的時候很容易不能專心,常常分心..
那我相信禱告是需要操練..
我覺得禱告是需要操練的..
那聖經有給我們很多的模範..
在聖經裡面我們看到有很多的模範..
那包括今天的這段信息也是如此..
也包括今天我們所讀的經文也是如此..
所以我提醒一下我們第一章三節告訴我們說,.
神自己在基督裡已經賜給我們天上各樣屬靈的福氣..
我們在第一章的時候看到神已經在基督耶穌裡面賜給我們各樣的屬靈的恩賜..
所謂的三福..
所以在這裡是已經告訴我們說,神不是那個很吝嗇,不給我們恩典的神..
我們知道我們的神不是一個很吝嗇的神,不給我們任何東西的神..
使徒保羅給我們看見說,他已經賜福了..
經過保羅,神已經賜福了給我們..
天上各樣或可以說是一切全部的福氣都在基督裡賞給我們了..
我們已經看到天上各樣包括一切和全部已經在主耶穌基督裡面賜給我們了..
所以當我們看到神永恆的計畫的時候,我們發現說神自己在萬世之前,.
祂就簡選,祂就預定一些人..
我們看見在創世以前,神已經是簡選,預定了一些人..
那祂的救贖就從萬人中呼召一些人,特地歸給祂自己..
祂的救贖就是在所有這麼多的世人裡面來招贅我們..
所以我們發現說在伊弗所書很關鍵的一個中心思想是在基督裡..
我們看到伊弗所書裡面很中心的思想就是在基督裡..
所以聽說我們有了什麼東西,那在這裡看見說有新生命,有成為聖傑,有成為兒女的名分,也成為後世..
所以我們看到在基督裡面我們得到一個新的生命,成為聖潔德兒女的名分,並且成為他的後知..
那知道說從自己神對於萬事的進展的時候,天上地下各樣的權柄都要歸在他的..
我們看到在永禽裡面,我們看見天上地下已經是將福氣給了我們..

$^{41}$所以在我們如今預測不了將來,對嗎?.
我們本身不能夠預測將來..
所以但是聖經給我們看見以後的將來的事情將會是怎麼樣子..
但是聖經給我們看到,預先告訴我們將來的事情應該會變成怎樣..
所以我們發現我們的視野要有一個改變..
我們看見我們的視野是不會改變..
如果我們信了主之後,我們的視野人就是說我能夠看見的這個層面的話..
如果我們信了主,我們所看見的視野是怎樣的話..
只是我們能夠看見的這個層面..
我們只能夠看見我們所看見的層面..
所以我們的價值觀,我們的世界觀也要有一個轉向..
我們的價值觀,我們的世界觀都應該有一個轉變..
因為神自己將來要帶進的是他實現的這個天國..
將來神會帶我們進入一個神的天國..
那天國也現在在我們中間也正在越來越敞開的時候,我們發現說到底我們能看見的是多少的成分?.
如果我們今天當這個天國慢慢展開的時候,我們可以看見的是一個什麼樣的層面?.
你們知道這位是誰嗎?.
大家知不知道這位人是哪位呢?.
知道的是幾首?.
我知道他是幾首..
不知道有沒有所謂?.
不知道也不要緊..
他有沒有錢呢?.
發達..
哇,這個大富翁..
所以請問他的兒女需不需要愁吃愁穿?.
他的子女?.
他的兒女需不需要愁吃愁穿?.
大家都明白他說什麼..
謝謝..
需不需要?不需要..
為什麼呢?因為他父親簡直是非常非常富裕..
他們的兒女不需要愁吃愁穿..
當然有人稱他為投資傳奇或股神..
有人稱他為股神..
那我再回顧一下,你們的兒女如果,他的兒女如果,所以千萬不要說我只是借用一些圖案..
如果他的女兒是有這樣子的奢華生活的話..
如果他的女兒,這個不是他本人的仗義,如果他的女兒有過奢華的生活..
我們會覺得說不足為奇怪..
大家不會覺得很奇怪..

$^{81}$他有得享受的話盡量的享受..
他有機會享受就盡量去享受..
當然,或許他有什麼跑車,有自己的私人飛機都好..
無論他有自己的私人飛機或者跑車都可以..
我的意思是說,當一個人知道他自己的父親有多富有的時候,他的生活,他整個社交的層面很不一樣..
當然如果他知道他父親是這麼有錢的話,他的朋友,他的交際圈子會很不一樣..
當然做父母親的心可能會很擔心,因為這個字一看就是一個寵壞的樣子..
但是他的父母可能會很擔心,因為一看就知道他們是一個寵壞的孩子..
從我們的視野,我們發現說,伊夫索告訴我們神自己是那位萬有的神..
從我們的視野,我們看到我們的神是一個萬有的神..
所以今天我們讀到的這個保羅為伊夫索教會的禱告,在伊夫索書中有兩個禱告..
在伊夫索書中有兩個不同的禱告..
這個禱告一開始就說,因此..
這個禱告一開始就用兩個字,因此..
所以如果是因此的話,就把前面的全部要拉進來,對嗎?.
如果用因此的話,我們就知道是說以前的一些東西..
那也很顯然的..
很顯然的..
落實在你我身上的就是一個信心,愛心的一個表現..
就是在你我裡面就是一個信心和愛心的表現..
那如果你讀剛才我們四位幫我們讀的這個荷荷本..
剛才我們主席幫我們讀的是用和諧本..
所以你會留意的話,他不太一樣..
他說我聽見你們的什麼?.
剛才是讀到我聽見你們..
然後親愛眾信徒,對吧?.
親愛眾信徒..
所以可能別的翻譯給我們一個比較平衡的一個看到..
這裡我是用這個新標點的荷荷本..
姚長老用的是用新標點的譯本..
所以他告訴我們說有信心,而且是對聖徒的愛心..
不單是有信心,而是對眾聖徒都有愛心..
金主妹,依夫所教會如果你知道它的背景的話,是一個非常非常黑暗的社會..
如果你看到當時的依夫所的時候,那個背景是一個非常黑暗的背景..
那有敬拜這個Diana女神..
是敬拜一個Diana的女神..
所以記不記得他們的販賣這個偶像大到一個程度是非常非常大的一個生意..
對販賣一些偶像的時候是一個很大很大的生意..
然後保羅讓金主妹教會看見說讚美,praise and prayer..
在這裡我可以講到讚美是一個禱告..

$^{121}$為此我們的讚美先然後進入禱告..
先讚美然後進入禱告..
我不知道您的禱告如何..
我實在不知道你們的禱告是怎樣..
我們有求的東西先放在前面..
我們很多時候將有求的東西放在我們的禱告的前面..
所以它告訴我們求的一件事情是什麼東西..
它就是求神就是那位榮耀的父將那四人智慧和啟示的靈賜給你們..
這個禱告是什麼呢?.
它說求我們主耶穌基督的神榮耀的父將賜人智慧和啟示的靈賜給我們..
保羅很清楚在禱告中祂的一個很清楚的祈求就是金主妹要真正認識神..
保羅在一個禱告裡面是很懇切的希望信徒能夠真正認識神..
我們的禱告通常是祝你什麼時候你自己聽我的禱告..
很多時候我們在祈禱的時候我們會想神你要怎樣聽我的禱告..
認識在我們人間有不同的層面..
認識我們有不同的層面..
那我這裡舉的例子是一般的用法..
我所舉的是一般的用法..
所以你可以說我知道神太是誰..
你可以說我知道姚神太是誰..
看到照片,我知道,但是知道跟你認識還是差一段距離,對吧?.
但是知道跟認識有一段距離..
那如果認識到你了解這位人,這位弟兄或這位姐妹..
從認識到了解弟兄或姐妹..
又是一個過程..
那如果進一步到說他是我的知心的人..
如果要達到那個是我的知心的人..
還是一個過程..
仍然是一個過程..
那使徒保羅這裡講到的認識是一個夫妻之間這樣的認識,是親密的認識的程度..
保羅在這裡講認識是一個很親密的認識,就好像夫妻之間的認識..
這肯定不只是知識上的一個認識..
肯定不是一個知識上層面的認識..
並且是隨著時間能夠跟主同行,越來越明白他的心意的這樣的一個認識..
也就是說在與神同行的過程裡面,是一個越來越深的一個認識..
那我們有不同層面的認識,應該帶出不同層面的一個效果..
有不同程度的認識應該會帶來不同的效果..
因為如果我單單知道巴菲特是誰..
如果我單單知道巴菲特是誰..
那對我沒有什麼影響嗎?.

$^{161}$對我自己是沒有什麼影響..
他的股票暴跌或者是整個暴漲,對我可能如果投資的話可能會心痛..
不過基本上沒有什麼影響..
就算他的股票猛跌猛起,對我基本上是沒有什麼影響..
但是如果說我認識一個人,那有可能我需要的時候我再跟他交接..
如果我認識一個人的話,我會有機會有時候會跟他交接..
那我們再深入一點,如果我們了解一個人的話,我們會在意他的反應..
如果我更加深入一點了解一個人的時候,我會很在意他的想法..
那如果你說你是個知心的人呢,我相信應該有一個相對的回應,.
就是應該會跟他更同步..
所以如果你是知心的話,可能更深入的就是好像與他同步一樣..
第一次會真認識他..
我們一直在尋求合成心意的領袖..
我們一直在尋找一個合成心意的領袖..
當然在這個人的身上,我們的素質是非常少的,不能夠少的..
也通常需要時間來驗證..
那我們看到說,真言給我們一句很好的提醒,他說:敬畏耶和華是智慧的開端,認識自聖者便是聰明..
在《箴言第九章第十節》提到,提醒我們敬畏耶和華是智慧的開端,認識自聖者便是聰明..
你發現我們在教會跟我們外面的父母的處理有沒有不一樣,我不是很肯定..
我不知道我們當中是跟外面比較我們怎樣處理父母的,是不是一樣,我們不知道..
如果你要你的兒女聰明的話,.
你會照聖經給我們的提示,還是說我們只是指在世俗上的一些東西讓她學習?.
不是說,如果照聖經的話,不是只是叫他們認識世俗上的一些知識..
那史徒保羅在十八節告訴我們說:.
求神不單只是要我們是智慧的,能夠認識祂,並且要我們照明你們心中的眼睛..
所以在《耶和華所書》一章十八節,他特別提到希望神照明你們心中的眼睛..
你們的心有沒有眼睛?.
我們內心有沒有眼睛?.
沒有,對吧?.
沒有,沒有眼睛..
他的意思是說讓我們不單只是在頭腦上知道..
就是說我們不單只是在頭腦上知道..
而且在我們的心我們能夠有一個相對的領受..
也就是說我們的心裡面有一個相對的領受..
所以知跟體會還是一個過程..
所以知跟體會也是一個過程..
所以開我們的心眼其實在聖經當中很清楚就是我們所謂的光照或者是illumination..
所以開我們的眼睛意思就是光照,英文就是illumination..
那意思是說如果神沒有光照的話,有些東西是看不見的..
如果不是因為神的光照,有些東西我們是看不清楚的..

$^{201}$所以如果神若不開恩,開恩讓你讀聖經的時候,.
我們還是看不到的..
如果不是因為神的恩能夠開你的眼睛的時候,.
我們都有很多東西是看不透的..
所以在我們幫助我們的孩子成長的時候,.
當我們幫助我們的孩子成長的時候,.
什麼時候他能夠摸到你的心呢?.
很可能是他哪一天不是從你那邊想要更多錢..
不是因為你賺了很多錢..
不是哪一天他想從你那邊拿到更多的錢..
不是因為他想到某一天能夠從你身上拿到很多錢..
他會自己預備一些東西送給你..
所以有些兒女和父母之間的過程,.
或許可以幫助我們去體驗一下,.
什麼是神要我們開我們的心眼,.
能夠知道的事情..
從父母和兒女的關係裡面,.
可以幫助我們明白神怎樣藉著我們來光照我們心裡面的眼睛..
他要使徒保羅讓我們知道幾件事情..
第一,就是我們的恩昭有何等之望..
因為神曾經在你我信徒的人的身上,.
神曾經呼召..
在我們這時已經是信眾的人,.
神在我們身上有呼召..
我們的呼召我們通常寫見證的時候..
我們的呼召是當我們寫見證的時候..
可能某一天,神感動我,讓我知道某些事情..
我們關起門看到背後寫著:我早已認識你..
神在創立世界以前已經揀選,已經預定..
神在創世已經揀選了我們..
包括在這個邀請呼召裡面,是一個很大的套餐..
因召是一個很大的禮物..
當你想的時候是大得不得了..
當我們想的時候,這個因是何等的大?.
你的恩召有何等的指望?.
你的因召有何等的指望?.
第二,他們看見他在聖徒中得的積業有何等豐盛的榮耀?.
我可以這樣講,基督或三一神在他的永世計畫裡面,.
已經把教會成為一個他與他分不開的,他也不可少的一樣事情..
所以三位一體的神,已經把教會成為他與他分不開的一樣東西..

$^{241}$所以他說他在聖體中得的積業有何等豐盛的意思?.
他得了我們,也成為他自己的榮耀,他的光彩..
所以在聖徒中所得的積業,在我們當中是神的榮耀..
反面,我們也在他裡面得到一切的豐盛..
而且反面來說,我也在他們裡面得到一切的豐盛..
兩面都是合乎聖經的..
第三,他給我們看見他向我們這信的人所顯的能力是何等浩大?.
你注意這三點的祈求..
你們看看這三樣的祈求..
不是人不知道..
不是說人不知道..
但是保羅要他們知道是how big,何等的..
就是說,保羅希望他們知道這是何等,是how big..
所以有何等的指望,有何等的榮耀,有何等的浩大?.
何等的指望,何等的榮耀,何等的浩大?.
世人都在追求有財富,也有權利..
很多人都追求這個世上,追求的是財富和權力..
弟兄姊妹,我們回到神,看到神祂自己的萬能..
但是我們一直看到神自己的萬能..
我們知不知道祂的能力是何等的浩大?.
我們是不是真的知道祂的能力有何等的浩大?.
如果我們看以下的經文的話,你會發現我們,或許給我們一些提示..
我們看以下的經文的話,你會發現我們,或許給我們一些提示..
祂在二十節告訴我們,就是照祂在基督身上所運行的大能大力..
在一章十九節講到,我們這些信了的人所顯現的能力是何等的浩大,.
是照著祂大能大力所運行的..
我們發現說,在這裡都是照著祂大能大力運行的..
是照著祂自己的大能大力來運行的..
所以如果你看原文的話更有意思..
如果你看原文的話更加有意思..
因為四個關於能力的詞彙都用了..
關於能力這四個詞都會用..
所以我不知道廣東話有沒有多少能力,有沒有多少不同的形容詞..
我不知道廣東話裡面能力有什麼特別的形容詞..
全部都用了..
是不是全部的形容詞都用了呢?.
所以能力就是dynamic,炸藥,dynamite..
能力在原文來說就是dynamite..
那你發現說,而且祂也有其他的名詞,祂的大能,祂的大力,.
而且祂不單只是有這些能力,而且祂要energize,.

$^{281}$祂使一切的運作能夠照著祂的意思來運行..
所以這個不單是大能大力,而是用energize,用大力來運行的..
因為當我們想到神做事情的時候,絕對不是說,.
主啊,我求你,好像是說要拉祂的手,請你做你的事情..
我們看見神做事情,不是像我們拉著神的手,你要做這些事情..
使徒保羅要我們看見的是,神的能力有何等的浩大?.
保羅希望我們知道神的能力有何等的浩大..
跟你我的禱告相當的不一樣..
所以鄧哲維,這個跟我們的禱告有很大的不同..
從一個公道的層面來講,.
那對於我也在作為一個長老的一個身份,.
也可以有同樣的一個對弟兄姊妹有這樣的一個盼望..
所以我作為長老對鄧哲維有一個同等的盼望..
我們一同來看見神對我們的恩昭,對我們的基業,對我們的能力是多大..
我希望大家能夠看見神的恩昭,神的基業對我們的能力是何等的大..
你說我現在軟弱..
有時候可能你說我現在很軟弱..
以馬內利的神所能夠成就的完全不受任何權勢的限制..
以馬內利的神所做的事不受一切的限制..
鄧哲維,我不知道你信主多少年?.
我不知道你們信了耶穌有多少年?.
你隨著你的年日,你對主的盼望是不是與日俱增?.
我不知道你對神的盼望是不是與日俱增?.
那還是你每天的生活裡面你還是覺得很無奈?.
或者你覺得你自己生活裡面很無奈?.
我們都好像被別人擺佈..
好像我們都被別人擺佈..
所以我們每天的怨言只能說都是他..
我們每天的怨言都是他..
我老闆就是這樣子苦待我..
每天擠壓我..
但是我們有時候想起來,.
我們活著其實不再是只活在這個能看見的層面..
所以我們過去神對我們呼召,.
其實帶給我們的是一個無限的盼望..
神叫我們能夠得求我們有一個無限的盼望..
我希望我們看自己是我們的老豆,.
我們的天父,.
是一個非常非常富有的一個神..
我應該讓我們明白我們的天父是非常非常的富有..

$^{321}$所以當他不給你的時候,.
絕對是有他的智慧..
絕對是有他自己的智慧..
要按他的方式,.
按照他的方式,.
按照他的時間,.
來成就他的旨意..
來成就他自己的旨意..
弟兄姊妹,你留給你的兒女是什麼東西?.
我想問問弟兄姊妹,.
你給你的兒女留下什麼東西?.
我想問問弟兄姊妹,.
你給你的兒女留下什麼東西?.
我們對一些多年在教會事奉的弟兄姊妹..
如果人來到我們教會看見你的腳鐘,.
他們看見的是什麼東西?.
是你的厲害還是神的偉大?.
二十二,二十三節,我們請弟兄姊妹一起來讀..
我們一起讀二十二和二十三節..
神使萬有服在他的腳下,.
由使他為了教會作萬有之首,.
教會是他的身體,.
是那充滿萬有者所充滿的..
弟兄姊妹,你讀這兩節的時候,.
什麼詞是忠心?.
你讀的時候,哪個詞是忠心?.
有沒有誰想提?.
當神,這個是絕對對的..
當神使萬有復齊在他的腳下的時候,.
就令到他成為教會作為萬有之首..
你發現我們對教會的認知通常很渺小?.
我一般對教會的了解是覺得教會很渺小..
通常有時間才來?.
有時候,自己有時間才來..
我們要來不來?.
但是從聖經的角度來看的話,教會絕不是這樣子的..
但是從聖經來說,這個教會絕對不是這樣..
如果神極其地顧惜他的教會,.
他絕不會輕易地,隨隨便便地,.
讓你我隨隨便便地進出..

$^{361}$他絕不會隨隨便便地,讓你我隨隨便便地進出..
所以,如果你看基督或神的事情為中心的話,.
那我也可以說教會在他心意也是絕對的中心..
但我們也看到教會在神的眼裡絕對是一個中心..
那我想請問你們的重心放在哪裡?.
我想請問你們的重心放在哪裡?.
禱告要操練..
禱告需要操練..
我的禱告也很需要操練..
我自己本人的禱告也是需要操練..
那我相信可以幫助我們的有幾方面?.
第一是聖經的模範..
我們可以學習聖經的模範..
就是一種的模式..
你可以看丹尼里的禱告或其他的就約聖徒的禱告也好..
或者我們可以看丹尼里的禱告或其他的就約聖徒的禱告也好..
你會發現他們都是先明白神自己對於當時,對於永世的一個旨意..
所以不一定要聽教會某某人禱告的華麗..
而不是按照教會某某弟兄的禱告..
這個人很會禱告,一定很熟齡..
如果聽到這個人的禱告,他這麼華麗,一定是很熟齡的人..
我們可以回到聖經看這些聖徒他們的禱告..
第二,我們一定要認識神..
第三,我們一定要認識神..
第四,我們一定要認識神..
第五,我們一定要認識神..
第六,我們一定要認識神..
第七,我們一定要認識神..
第八,我們一定要認識神..
第九,我們一定要認識神..
第十,我們一定要認識神..
張子銘,神在我們的禱告中一定要是中心,一定要是重要..
我希望我們的禱告以認識神為重點..
那我們心眼要看見..
我們希望我們的心眼能夠看見..
我們一定要看見..
我們一定需要看見..
如果我們多年來聚會的時候,我們的團契查清,我們個人的成長,都為了這樣一個宗旨,.
不是說神你要動功,是神已經在做功了,是要我們心眼看見..
希望看到神在我們心靈的動功,是神讓我們的心眼能夠看見..

$^{401}$所以你發現我們在帶領一個弟兄或姐妹的時候,.
我們發覺當我們帶領一個弟兄或姐妹的時候,.
當他心眼開的時候,.
當他心靈開的時候,.
我告訴你,那是不得了..
實在是不得了..
因為你發現說,忽然間,他對神的事情忽然間就感到非常有興趣..
他突然間對神的事情非常有興趣..
那他看的視野每次讀經的話,不一樣..
所以他看聖經的視野完全不一樣..
所看的字句沒有改變..
所看到的字句完全沒有改變..
人改變了..
只有人改變了..
那最後,神的旨意的中心是教會..
最後就是講到神的旨意的中心就是教會..
如果教會掌管,教會是屬神的..
教會是屬神的..
那祂為教會之首..
祂是教會之首..
那教會和這個世界你們怎麼去衡量呢?.
那這個是變成很切身的一個問題..
那這個是很切身的一個問題..
也要問你,也要問我..
也是需要問你和問我自己..
如果神看教會是什麼?.
如果神看教會是什麼的話?.
多大?.
有多大?.
何等的大?.
何等的豪大?.
那祂看這個敗壞的這個世界怎麼看?.
那祂看這個敗壞的世界怎麼看?.
丟子妹,我可以問你嗎?.
教會在神的心目中是多大,多重要,多重?.
我可以問你,教會在神的心裡面是多重?.
那我們幫助弟兄姊妹成長的,.
難道只是幫助她,只是個人的需要,.
這些東西是一直需要的..
我們幫助弟兄姊妹的時候,.

$^{441}$不就是單單幫助她個人的需要,.
這些是需要的..
但是能夠提升她,.
讓她看見說神自己看教會是多大?.
我可不可以將她提升,.
來看到神看教會是多大?.
我們鼓勵說丁永哲妹去找去別的地方,.
我們要深造,或者要找工作都好..
或者看到有些弟兄姊妹去深造或者是找工作..
他們就問我,你覺得我去這邊好不好?.
或者他們問我去這裡好不好?.
我說我不知道..
我只能夠答,我不知道..
我們可以看看禱告..
我們可以看看,然後禱告..
那我通常會問他一句..
或者我們問他一句..
你對教會怎麼看?.
你對教會有什麼看法?.
如果你去那邊,你喪失了你自己的生命..
如果你去到那個地方,你喪失了自己的生命..
你再也不愛教會..
你再也不愛聚會..
你也不愛聚會..
也沒有跟你能夠同心的丁永哲妹一起來服事..
那你真的要去嗎?.
你是不是真的要去呢?.
教會是不是我們在很多計畫,在很多想法當中,是不是我們的重心?.
是不是教會在我們眾多的選擇裡面,是不是我們的重心?.
還是其實我們萌新自問的時候,我們是把它放在有時間再說..
而不是我們有時間再說..
丁永哲妹,我們的禱告需要有一個轉變..
我們的禱告需要有一個改變..
先看到神的心意..
然後我們再禱告..
然後我們再去禱告..
我們在粵語部當中有很多的禱告的戰士..
我們在粵語部裡面有很多禱告的戰士..
他們也甚至為牧長禱告..
丁永哲妹,我可以跟你講,我們很需要你們的禱告..

$^{481}$我跟你們很認真地說,我們很需要你們的禱告..
因為我們絕對不是憑人意來行事情..
我們盼望我們所做的,是能夠按照神自己的旨意..
我們希望我們所做的,都是按照神自己的旨意..
所以當我們繼續尋求一個合神心意的牧者的時候,.
所以當我們尋找一個合神心意的牧者的時候,.
我盼望我們自己也在我們禱告當中繼續地成長..
以致你所求的東西,是神所喜悅的東西..
也是神所喜悅的東西..
我們禱告..
我們一起禱告..
因為你是那位大有能力的神..
並且你按照你的旨意,已經將一切都展開..
因為你自己按照自己的旨意,將一切展開..
求你自己打開我們的心眼,讓我們看見..
求你打開我們的心眼,讓我們能夠看見..
我們如果是沒有指望的,你賜給我們真正的盼望..
如果我們沒有指望的,求你給我們真正的盼望..
如果我們自己覺得我們沒有能力的話,.
我們讓我們看見你自己的大能..
當我們覺得自己沒有能力的時候,.
讓我們看見你自己的大能..
當我們看見我們什麼都沒有的,.
讓我們看見在你裡面樣樣都有..
當我們覺得自己一無所有的時候,.
我們看見你自己裡面樣樣都有..
你幫助我們..
\newpage



\section{以賽亞書 49:14-18}
\label{sec:032127}
\textbf{必不忘記}~~~~馮潔儀傳道
\newline
\newline
連結: \href{http://www.church.com.hk/acms/content.asp?site=cdc\&op=show\&type=product\&code=032127\&layout=sermon}{\texttt{032127}} ~~~~ 語音日期: 2026-05-10
\newline
\newline
\hyperref[sec:032126]{\small{< < < PREV SERMON < < <}}
~
\hyperref[sec:index_chronic]{\small{[返順時目]}}
~
\hyperref[sec:032128]{\small{> > > NEXT SERMON > > >}}
\newline
\newline
以賽亞書 49:14-18
\newline
\begin{longtable}{cl}
\hline
\hline
章節 & 經文 (和合本修訂版)\\
\hline
49:14 & \begin{tabularx}{0.7\textwidth}{X} 錫安說：「耶和華離棄了我，主忘記了我。」 \end{tabularx} \\ \\ \relax
49:15 & \begin{tabularx}{0.7\textwidth}{X} 婦人焉能忘記她吃奶的嬰孩，不憐憫她所生的兒子？即或有忘記的，我卻不忘記你。 \end{tabularx} \\ \\ \relax
49:16 & \begin{tabularx}{0.7\textwidth}{X} 看哪，我將你銘刻在我掌上，你的城牆常在我眼前。 \end{tabularx} \\ \\ \relax
49:17 & \begin{tabularx}{0.7\textwidth}{X} 建立你的勝過毀壞你的，使你荒廢的必都離你而去。 \end{tabularx} \\ \\ \relax
49:18 & \begin{tabularx}{0.7\textwidth}{X} 你舉目向四圍觀看，他們都聚集來到你這裡。我指著我的永生起誓，你定要以他們為妝飾佩戴，帶著他們，像新娘一樣。這是耶和華說的。 \end{tabularx} \\ \\
[1ex]
\hline
\hline
\end{longtable}
$^{1}$今日係母親節.
特別選jor 一段同母親有關係ge 經文.
同大家分享同埋思想.
經文係選至二賽亞書四十九章十四至十八節.
聖經咁講.
石安說 耶和華離棄了我 主忘記了我.
婦人焉能忘記她吃奶的嬰孩.
不憐恤她所生的兒子.
即或有忘記的 我卻不忘記你.
看呀 我將你銘刻在我掌上.
你的長緩尚在我眼前.
你的兒女必急速歸回.
毀壞你的使你荒廢的.
必都離你出去.
他們都聚集來到你這裡.
耶和華說 我指著我的永生起誓.
你必要以他們為裝飾佩戴.
以他們為華戴束腰.
像新婦一樣.
聖經中其實.
神好少用母親的向度去表達自己.
多數用父親的向度.
所以在聖經中比較少這個記載.
今天看的這段經文是其中一段.
為什麼在這個地方.
在《以賽亞書》49章.
神要用一個母親的向度來表達自己呢?.
原來由《以賽亞書》的40章到55章.
一個這麼大的篇幅中.
神其實是透過先知和他的百姓石安.
石安不單只是一個地方名.
石安更加是名字和身份.
神要透過先知以賽亞和石安說.
你們被擄的子民.
我將來要給你一個新的一天.
新的局面.
耶路撒冷可以重被建造.
你們要回歸在神的荒土地圖裡面.
當神這樣去表達自己的時候.
去表出祂的應許的時候.

$^{41}$為什麼去到第49章.
神的百姓石安會發出一個這樣的呼喊.
耶和華利希了我 主忘記了我.
明明前面神很斬釘截鐵地.
將祂的應許說了出來.
為何他們會出現一個這樣的呼喊.
就因著他們的呼喊.
神就以一個母親的向度去表達.
去回應他們的呼喊.
神要告訴他們 我沒有忘記你.
今天神都透過這番說話.
跟我們說 神沒有忘記我們.
究竟當時.
石安 神的百姓面對著什麼樣的情況.
何時沒有忘記我們呢.
神在我們見不到前路的時候.
祂沒有忘記我們.
當時的石安.
他們面對的是淪亡之後.
淪亡之後 因為他們被擄到巴比倫.
他們現在面對的是一個很艱難的環境.
很奇怪 剛才我們讀是49章第14節開始.
49章的第1至13節說什麼呢.
第1至13節的經文裡面.
其實是包含著其中一篇的「獨人之歌」.
在以賽亞書裡面總共有四首「獨人之歌」.
四首「獨人之歌」是什麼呢.
「獨人」是特別說到指向這位尼塞雅.
以賽亞先知去說「獨人之歌」的時候.
是在說給百姓聽.
「尼塞雅會來拯救你們」.
「尼塞雅會來帶來一個新的局面」.
明明在49章1至13節才說完「獨人之歌」.
為什麼去到第14節.
石安竟然發出一個這樣的呼喊.
「耶和華離棄了我 忘記了我」.
「太遠了 尼塞雅何時來」.
「太遙遠了 這個局面」.
「我看不到前面有什麼去路」.
「困在現在這個困局裡面」.

$^{81}$「我就是一個被擄的人」.
其實我是一個長期病患者.
我有什麼病呢?我有濕疹.
你現在看我的樣子好像沒事.
但在我最嚴重濕疹的時候.
那個情況是怎樣呢?.
除了身體很多地方都有濕疹出現.
連我的頭皮這些地方都有.
它是會自己流一些黃色的水出來.
如果你有濕疹你很明白我在說什麼.
濕疹有一個最大的特點就是什麼呢?.
癢 要嘔吐.
很多時候我晚上是癢醒的.
我就要坐著嘔吐.
嘔吐到什麼時候才可以再睡呢?.
嘔吐到流血,嘔吐到皮破.
所以很多時候我起床的時候.
我的床單和我的枕頭袋.
除了有些黃色的水印之外.
還是血漬斑斑.
在那個時候有沒有看醫生呢?.
當然有了.
我看了很多醫生.
中醫,西醫,教授,專科都看了.
我想如果我沒有記錯.
我看了六至七個醫生.
我那個時候的狀況是怎樣呢?.
好像在新約有一個血漏婦人.
她說什麼?.
《聖經》記載她.
在許多醫生手裡受了許多的苦.
花了很多的錢.
病情一直都沒有變好.
我就是這樣.
我還可以怎樣呢?.
見不到前面有路可走.
所以那個時候.
我都簽過.
神啊!你是不是忘記了我?.
我的檔案已經放上來了.

$^{121}$是不是疊起來了?.
你忘記了我.
當我們見不到前路.
在困局裡的時候.
我們就會很容易發出錫安的呼喊.
耶和華說你棄了我.
忘記了我.
在現實很差的環境.
同樣都會有一個這樣的呼喊.
當時的錫安.
面對耶路撒冷.
本來是一個大城.
是一個很繁榮的城市.
但是當時耶路撒冷被淪陷.
一片頹垣敗瓦.
方直一片.
看著自己的家園被毀.
自己的身份被磨滅.
可以怎樣收科?.
那個現實很難處理.
不知道怎樣處理.
當時錫安說.
神啊,你忘記了我.
有一個故事.
沙灘上的竹印.
可能我們中間有不少的電影姐妹.
都聽過這個故事.
我簡單複述一下.
這個是怎樣的故事.
是有一個主人翁.
他有一晚做夢的時候.
他做夢夢見自己在沙灘上行走.
在沙灘上行走的時候就有腳印.
他就在行走.
其實他見到.
就是他自己人生一個回顧的片段.
他在行走的時候發現.
不單是他那雙腳印在那裡.
旁邊有另一雙腳印.
原來旁邊那雙腳印是神.

$^{161}$他很快樂.
見到沙灘上兩雙腳印並排而走.
很開心有神與我同在.
不過當他一度回顧的時候.
他發現其中一個片段.
沙灘上只剩下一雙竹印.
而這個時刻.
就是他人生最落魄.
最坎坷.
最痛苦的時間.
只剩下一雙竹印.
他就向神問.
神啊為什麼在那段最艱難的時間.
你離開我.
有時我們都可能像主人翁一樣.
向神發出這樣的呼喊.
在我們最痛苦的時候.
神啊你離開我.
那我怎麼辦.
我認識一個家庭有幾姐妹.
我認識其中一個比較親的姐妹.
這個家庭在兩年時間內.
喪失了四位家人.
分別是她的爸爸,媽媽,媽媽和婆婆.
四個家人相繼安息主懷.
有一次跟這個姐妹聊天的時候.
她叫我做導師.
她說導師.
我快撐不住了.
我現在還可以怎樣.
很難過.
有時我們人生歷程中.
都有這樣的時刻.
我們撐不住了.
神啊現在這個環境很差.
何時才可以改變.
當神的百姓.
向神發出這樣的呼喊.
她不單止懷疑神.
她有機會懷疑自己.

$^{201}$為什麼這樣說呢.
在以賽亞書前的篇章.
神透過以賽亞才點出.
神的百姓錫安做錯了什麼.
他們投靠外邦.
離棄神.
這個時候她也會想.
最壞是我做錯了.
做錯了神就不理我.
神就懲罰我.
我發現很多弟兄姊妹.
都有這樣的概念.
一遭遇緊患難.
就會想是不是我做錯了.
得罪神了.
所以神就透過這件事懲罰我.
原來這種想法.
在跟從耶穌的門徒心裡.
也是這樣.
在新約裡有行福音.
記載了一個生來黑暗的人.
你見耶穌.
當門徒見到這個生來黑暗的人.
他就跟耶穌說.
他說夫子.
這個生來黑暗的人.
是他犯了罪.
還是他的父母犯了罪.
後面其實是隱藏了.
蘊含了一些還沒說出來的說話.
是他犯罪.
是他父母犯罪.
所以才會得到這樣的下場.
這樣的報應.
才會跌到一個這樣.
這麼慘的結果.
有時候我們.
在一些軟弱跌倒的時候.
我們心裡都會有想.
是不是我不對.

$^{241}$所以神忘記我.
任由我在困難裡面.
弟兄姊妹不要誤會神.
當神聽見石安這樣的呼喊.
神馬上回應.
因為錯不是這樣.
沒有忘記.
神就馬上回應.
神說.
「婦人焉能忘記她吃奶的嬰孩.
不連恤她所生的兒子.
即或有忘記的.
我卻不忘記你」.
神以一個普遍母親的特質.
去回應石安的呼喊.
這裡特別說.
媽媽怎麼能夠忘記自己的寶寶呢.
這個寶寶是什麼時段的寶寶呢.
是正在吃奶的嬰孩.
正在吃奶的嬰孩有一個特點.
就是什麼都做不到.
她自己肚餓沒得找東西吃.
她不舒服.
不知道怎樣.
不懂得處理自己.
她奶濕了.
她也不懂得幫自己處理.
她整個人依賴著媽媽的照顧和看顧.
她有什麼需要都是靠媽媽.
既然如此.
媽媽很清楚吃奶的嬰孩很需要自己.
媽媽就會很緊張這個嬰孩.
所以一哭.
很多媽媽就很緊張.
「是否冷呢?是否熱呢?.
是否濕呢?是否乾呢?」.
以前我的教會.
小小一間.
我自己的教會.
有一個地方是崇拜的.

$^{281}$旁邊隔著一個玻璃房.
嬰孩放在那裡.
小朋友放在那裡.
媽媽爸爸就去外面崇拜.
但是在玻璃房.
你知道小朋友的哭聲是很響亮的.
還有很高頻率的.
當一有哭聲的時候.
不會全部媽媽都衝出去.
只有一個媽媽衝出去.
不會衝錯的.
一定是自己的小朋友哭.
媽媽怎會忘記自己的孩子.
人在崇拜.
心還在連繫著嬰孩.
怎會忘記呢?.
有些媽媽這樣說.
吃奶的時候.
當餵哺母乳的時候.
有些母親這樣說.
那個時刻.
餵哺母乳的時刻.
是媽媽和嬰孩的交融.
那個交流是很快樂的交流.
我不是媽媽.
我不明白.
是一個很微妙的事.
我們中間的弟兄都不明白.
但如果我們中間的姐妹.
你曾經是餵哺母乳的媽媽.
你會體會.
這是一種很親密的關係.
很親密.
彼此在交流中.
有一種很微妙的連結.
神跟石安說.
我就是這樣不忘記你.
我沒有忘記你.
然後神說.
即或有忘記.

$^{321}$我覺得這個補一補非常好.
真的有這麼大頭蝦的媽媽嗎?.
真的有嗎?.
我都不會忘記你.
其實神這樣說.
是表示萬萬不可能.
石安你知不知道.
要我忘記你嗎?.
沒有可能.
我怎麼可能忘記你.
神其實是在表達一個這樣的意思.
但當我們人落在困苦的時候.
的確是真的.
我們有時都會好像石安那樣.
發出一個懷疑.
神忘記了我.
神離棄了我.
在新約.
當耶穌和門徒.
坐船渡過加利利海的時候.
刮起大風浪.
耶穌在船上睡覺.
門徒走來找耶穌.
有沒有印象門徒說什麼?.
他說夫子.
我們喪命了.
然後呢.
理不顧摸.
即是說.
耶穌.
我們快要沒命了.
很大風浪.
你不理我們.
難道任由我們死嗎?.
原來人落在困難當中.
真的會這樣.
為什麼呢?.
為什麼會出現一個.
你不理我們.
你忘記了我們.

$^{361}$你離棄我們的想法呢?.
是因為我們習慣了.
用看得到的東西去衡量我們的神.
我們看到的是環境.
我們看到環境不順利.
環境出現困難.
我們就覺得.
神啊,你肯定是忘記了我.
你肯定是不理我.
當環境很好.
很順利的時候.
神啊,多謝你看顧我.
賜我恩典.
我們有時候會有這樣的弊病.
用看得到的東西.
去衡量我們的神.
是否真實.
不好.
神啊,不是的.
不是這樣.
我沒有忘記你.
不僅如此.
神不僅用母親的向度.
去回應錫安的呼喊.
還怕錫安不明白.
然後祂說.
我怎樣不忘記.
我告訴你.
我不忘記你的程度.
漢啊.
我將你.
這個你就是錫安.
我將你銘刻在我掌上.
你的長緩常在我眼前.
神怎樣不忘記錫安呢.
他用了兩個意象去表達.
第一個意象.
我將你.
這個就是錫安.
銘刻在我掌上.

$^{401}$很寶貴.
什麼是我們整個身體.
要拿著的東西.
最重要的東西要拿著.
就是這樣.
就是掌心.
銘將錫安刻的名字在他的掌上.
拿著.
拿著是很重要的東西.
大家有沒有見過.
應該有見過.
有些人很喜歡紋身.
紋什麼多.
有些弟兄姊妹今早說.
紋條龍.
都有的.
紋圖畫.
大家有沒有見過一些.
是紋一些人名下去的.
有吧.
那些是什麼人.
紋下去的人名.
通常是什麼人.
一定不會是仇人.
不會的.
通常紋什麼.
女朋友.
男朋友.
先生太太.
兒女的名字.
刻得紋得下身的那些.
是我們認為.
他很重要的人.
刻骨銘心的人.
神現在說.
他將我們紋在.
刻在他的掌上.
很重要.
你和我對於神來說.
是很重要的.

$^{441}$弟兄姊妹.
還有.
原來這種將名刻在一個人的身上.
這種做法在舊約上有的.
通常是出現在什麼時候.
譬如我是一個主人.
我有一個奴僕.
我有一個奴隸.
我是主人.
我就會將我自己的名.
刻在我奴隸的身上.
證明是什麼呢.
這個奴隸是屬於我的.
不過神現在.
調轉了.
現在神是將我們的名.
將錫安的名.
刻在他的掌上.
那就是說.
他不是要表達.
我們屬於神.
不是的.
神是在說.
我是屬於你的.
我怎會忘記你.
我是屬於你的.
我們的名在神的掌上.
我覺得這是一個很親密.
很鞏固的關係.
一個屬於自己的.
互屬的關係.
我怎會忘記呢.
聖經用這個刻字.
不是寫.
寫可以磨滅.
刻下去.
是永不磨滅.
神說.
我要告訴你.
錫安.

$^{481}$我沒有忘記你.
到一個這樣的地步.
我將你的名.
寫在我的手掌上.
第二個意象.
你的長門.
你是錫安.
你的長門常在我眼前.
剛才我說過.
耶路撒冷是一個大城.
一個長門.
對於一個大城來說.
是非常重要的設施.
如果這個長門倒塌.
裡面的人就不得安寧.
失去保障.
失去安全.
但是在現實當時的環境.
耶路撒冷的長門倒塌了.
但是神說.
不是啊.
你們的長門的情況.
在我眼前.
我看在眼裡.
我知道.
常在我眼前.
我很喜歡這個常字.
常在我眼前.
這個常含有.
繼續延續下去的意思.
以前記得.
那個長門.
好好的時候.
我記得.
但是現在倒塌了.
不是.
我仍然記得.
神說我仍然記得.
繼續記得沒有忘記過.
由何時開始.

$^{521}$一直都沒有忘記.
以上書46章.
已經去表達過.
這個想法.
這個信息.
「雅各家」.
「以色列家一切如誠的」.
「要聽我言」.
「你們自從生下」.
「就蒙我保譜」.
「自從出胎」.
「便蒙我懷差」.
「直到你們年老」.
「我仍這樣」.
「直到你們發白」.
「我仍懷差」.
「我已造作」.
「也必保譜」.
「我必懷抱」.
「也必拯救」.
看到嗎.
從生下來的時候.
直到年老發白.
神都沒有忘記.
神仍然去記得以色列.
仍然記得石安.
記得他的子民.
這樣的事對於以色列.
對於石安.
是天方夜譚.
因為他們當時的環境.
不是傳遞一個這樣的信息.
他們覺得自己沒有盼望.
如果弟兄姊妹今天.
有一些現況.
令到你產生懷疑.
出現擔心.
希望你記得.
神跟你說.
神沒有忘記你.

$^{561}$神一直都記得你.
終久不忘記.
這個終久的意思.
是直到這個事情有轉機.
神將轉機.
發生出來的時候.
第十七十八節.
「你的兒女必急速歸回」.
「毀壞你的使你荒廢的」.
「必都離你出去」.
「他們都聚集」.
「來到你這裡」.
「耶和華說」.
「我指著我的永生起誓」.
「你必要以他們為裝飾佩戴」.
「以他們為華戴束腰」.
「像新婦一樣」.
其實由第四十九章十七節打後.
去到五十二章.
這麼多的經文.
神是在說將來要做的.
要做成一個新的局面.
這個局面是一個怎樣的局面.
是一個怎樣華麗.
怎樣繁榮.
怎樣復興的局面.
當時的石安.
他認為自己.
好像一個寡婦喪子一樣.
這麼慘淡.
但是神說.
「我要讓你成為像新婦一樣」.
「這麼華麗」.
這是神要做的工作.
我們的神.
不只是「我知道」.
「我知道你辛苦」.
「我記得你困難」.
「在困局中」.
「我知道」.

$^{601}$「我明白」.
我們的神不只是這樣.
神說「我要做」.
「我記得之餘」.
「我明白之餘」.
「我還要做一個新的現象」.
「新的事情要發生」.
其實這種做法.
這種神的計劃.
在出埃及的時候已經有.
當以色列人在埃及地圍爐的時候.
神要拯救他們.
他就要派摩西來救他們.
當時出埃及的神.
跟摩西這樣說.
要派摩西去的時候.
說了以下的說話.
「耶和華說」.
「我的百姓在埃及所受的困苦」.
「我實在看見了」.
「我看到了」.
「他們因受督公的合制所發的哀聲」.
「我也聽見了」.
「我聽見了」.
「我原知道他們的痛苦」.
「你辛苦我知道」.
如果聖經只停在這裡.
沒有了第八節.
就慘了.
有時我們都會有一個這樣的情況.
「哎呀哎呀」.
「都是替你辛苦了」.
「哎呀你真是慘了」.
「真是慘了」.
這樣就沒有了.
「不是啊」.
「神說我下來」.
「什麼叫我下來」.
「就是我來」.
「我要救」.

$^{641}$「是要救他們脫離埃及人的手」.
「令他們出了那地」.
「到美好寬闊流乃與物之地」.
「神說我要來幫你」.
「我要來拯救你」.
「我不單止知道你辛苦」.
「我要來拯救你」.
弟兄姊妹,困難會過的.
真的會過的.
因為神記得,神明白,神知道.
而且神會工作.
會過去的.
神會帶給我們一個新的景象.
這是神給我們的應許.
給當時的錫安,也給我們.
神必不忘記她的指紋.
好像媽媽怎麼也不會忘記.
她吃奶的鷹蟹那樣.
不會.
神必不忘記我們.
那我們呢?.
我們也要不忘記神.
讓我們在任何時刻.
我們都能夠紀念我們的主.
因為我們知道.
我們的主紀念我們.
請祈禱.
你昔日向錫安發出的應許.
鐵一般的應許.
今天向我們展現.
主多謝你從來沒有忘記我們.
的確在人生裡面.
不是天天都花香常蔓.
不是天天都陽光普照.
在我們人生裡面的確有風雨交加的時候.
但是神讓我們不忘記你.
讓我們清楚肯定.
神你是那位不忘記我們的神.
在我們痛苦,辛淫的時候.
主啊,求你去幫助我們.

$^{681}$能夠抓住你.
讓我們能夠繼續堅心信靠你.
直至我們見到神為我們展開新的一頁的時刻.
主啊,願你去幫助我們.
神你透過母親的特性.
去告訴我們.
你萬萬不可能忘記我們.
我們知道,我們明白,我們相信.
誠心祈禱,奉主耶穌的命求.
阿們.
\newpage



\section{列王記下 4:8-37}
\label{sec:032128}
\textbf{婦人的愛與信心}~~~~溫艾衛傳道
\newline
\newline
連結: \href{http://www.church.com.hk/acms/content.asp?site=cdc\&op=show\&type=product\&code=032128\&layout=sermon}{\texttt{032128}} ~~~~ 語音日期: 2026-05-10
\newline
\newline
\hyperref[sec:032127]{\small{< < < PREV SERMON < < <}}
~
\hyperref[sec:index_chronic]{\small{[返順時目]}}
~
\hyperref[sec:032129]{\small{> > > NEXT SERMON > > >}}
\newline
\newline
列王記下 4:8-37
\newline
\begin{longtable}{cl}
\hline
\hline
章節 & 經文 (和合本修訂版)\\
\hline
4:8 & \begin{tabularx}{0.7\textwidth}{X} 一日，以利沙經過書念，在那裡有一個富有的婦人強留他吃飯。此後，以利沙每次經過就轉到那裡去吃飯。 \end{tabularx} \\ \\ \relax
4:9 & \begin{tabularx}{0.7\textwidth}{X} 婦人對丈夫說：「看哪，我知道那常從我們這裡經過的是神聖的神人。 \end{tabularx} \\ \\ \relax
4:10 & \begin{tabularx}{0.7\textwidth}{X} 我們可以為他蓋一間有牆的小閣樓，裡面安放床榻、桌子、椅子、燈臺。每當他來到我們這裡，就可以住在那裡。」 \end{tabularx} \\ \\ \relax
4:11 & \begin{tabularx}{0.7\textwidth}{X} 一日，以利沙來到那裡，轉進那閣樓，躺臥在那裡。 \end{tabularx} \\ \\ \relax
4:12 & \begin{tabularx}{0.7\textwidth}{X} 以利沙吩咐僕人基哈西說：「你叫這書念婦人來。」他把婦人叫了來，婦人就站在以利沙面前。 \end{tabularx} \\ \\ \relax
4:13 & \begin{tabularx}{0.7\textwidth}{X} 以利沙吩咐僕人說：「你對她說：『看哪，你為我們費了許多心思，我可以為你做甚麼呢？我可以為你向王或元帥求甚麼呢？』」她說：「我已住在自己百姓之中。」 \end{tabularx} \\ \\ \relax
4:14 & \begin{tabularx}{0.7\textwidth}{X} 以利沙說：「究竟可以為她做甚麼呢？」基哈西說：「她真的沒有兒子，她丈夫也老了。」 \end{tabularx} \\ \\ \relax
4:15 & \begin{tabularx}{0.7\textwidth}{X} 以利沙說：「叫她回來。」於是他叫了她來，她就站在門口。 \end{tabularx} \\ \\ \relax
4:16 & \begin{tabularx}{0.7\textwidth}{X} 以利沙說：「明年這時候，你必抱一個兒子。」她說：「神人，我主啊，不要這樣欺哄婢女。」 \end{tabularx} \\ \\ \relax
4:17 & \begin{tabularx}{0.7\textwidth}{X} 婦人果然懷孕，到了明年那時候，生了一個兒子，正如以利沙向她所說的。 \end{tabularx} \\ \\ \relax
4:18 & \begin{tabularx}{0.7\textwidth}{X} 孩子長大，一日出去到他父親和收割的人那裡。 \end{tabularx} \\ \\ \relax
4:19 & \begin{tabularx}{0.7\textwidth}{X} 他對父親說：「我的頭啊，我的頭啊！」他父親對僕人說：「把他抱到他母親那裡。」 \end{tabularx} \\ \\ \relax
4:20 & \begin{tabularx}{0.7\textwidth}{X} 僕人抱去，交給他母親。孩子坐在母親的膝上，到中午就死了。 \end{tabularx} \\ \\ \relax
4:21 & \begin{tabularx}{0.7\textwidth}{X} 他母親上去，把他放在神人的床上，關了門出來， \end{tabularx} \\ \\ \relax
4:22 & \begin{tabularx}{0.7\textwidth}{X} 呼叫她丈夫說：「你叫一個僕人給我牽一匹驢來，我要趕去見神人，然後回來。」 \end{tabularx} \\ \\ \relax
4:23 & \begin{tabularx}{0.7\textwidth}{X} 丈夫說：「今日不是初一，也不是安息日，你為何要到他那裡去呢？」婦人說：「平安無事。」 \end{tabularx} \\ \\ \relax
4:24 & \begin{tabularx}{0.7\textwidth}{X} 於是她備上驢，對僕人說：「走，趕緊走，除非我吩咐你，不要為了我而慢下來。」 \end{tabularx} \\ \\ \relax
4:25 & \begin{tabularx}{0.7\textwidth}{X} 婦人往迦密山去，到了神人那裡。神人遠遠看見她，對僕人基哈西說：「看哪，書念的婦人來了！ \end{tabularx} \\ \\ \relax
4:26 & \begin{tabularx}{0.7\textwidth}{X} 現在你跑去迎接她，對她說，你平安嗎？你丈夫平安嗎？孩子平安嗎？」她說：「平安。」 \end{tabularx} \\ \\ \relax
4:27 & \begin{tabularx}{0.7\textwidth}{X} 婦人上了山，到神人那裡，就抱住神人的腳。基哈西前來要推開她，神人說：「由她吧！因為她心裡愁苦。但耶和華向我隱瞞這事，沒有告訴我。」 \end{tabularx} \\ \\ \relax
4:28 & \begin{tabularx}{0.7\textwidth}{X} 婦人說：「我何嘗向我主求過兒子呢？我豈不是說過，不要欺哄我嗎？」 \end{tabularx} \\ \\ \relax
4:29 & \begin{tabularx}{0.7\textwidth}{X} 以利沙吩咐基哈西說：「你束上腰，手拿我的杖前去。若遇見人，不要向他問安，人若向你問安，也不要回答。要把我的杖放在孩子臉上。」 \end{tabularx} \\ \\ \relax
4:30 & \begin{tabularx}{0.7\textwidth}{X} 孩子的母親說：「我指著永生的耶和華，又指著你的性命起誓，我必不離開你。」於是以利沙起身，隨著她去了。 \end{tabularx} \\ \\ \relax
4:31 & \begin{tabularx}{0.7\textwidth}{X} 基哈西在他們以先去了，把杖放在孩子臉上，卻沒有聲音，也沒有動靜。基哈西回去，迎見以利沙，告訴他說：「孩子還沒有醒過來。」 \end{tabularx} \\ \\ \relax
4:32 & \begin{tabularx}{0.7\textwidth}{X} 以利沙進了屋子，看哪，孩子死了，放在自己的床上。 \end{tabularx} \\ \\ \relax
4:33 & \begin{tabularx}{0.7\textwidth}{X} 他進去，關上門，只有他們兩個人，他就向耶和華祈禱。 \end{tabularx} \\ \\ \relax
4:34 & \begin{tabularx}{0.7\textwidth}{X} 他上去伏在孩子身上，口對口，眼對眼，手對手。他伏在孩子身上，孩子的身體就漸漸暖和了。 \end{tabularx} \\ \\ \relax
4:35 & \begin{tabularx}{0.7\textwidth}{X} 然後他下來，在屋裡來回走了一趟，又上去伏在孩子身上。孩子打了七個噴嚏，眼睛就睜開了。 \end{tabularx} \\ \\ \relax
4:36 & \begin{tabularx}{0.7\textwidth}{X} 以利沙叫基哈西說：「你叫這書念婦人來。」於是他叫了她來。婦人來到以利沙那裡，以利沙說：「把你兒子抱起來。」 \end{tabularx} \\ \\ \relax
4:37 & \begin{tabularx}{0.7\textwidth}{X} 婦人就進來，在以利沙腳前俯伏於地，向他下拜，然後抱起她兒子出去了。 \end{tabularx} \\ \\
[1ex]
\hline
\hline
\end{longtable}
$^{1}$各位兄弟姐妹 早安.
剛才主席也有提及.
今天是母親節.
容讓我在此和大家致敬一下.
求神讓我們每一位母親都有健康的身體.
平安和安康.
今天想和大家分享一位母親的故事.
她是一位善良而且滿有信心的母親.
先問問大家.
你相不相信一個很簡單的舉動.
就可以改變人的生命呢?.
今天我們在列王記夏四章八至三十七節.
我們將會認識一位婦人.
是一位母親的故事.
她敏感到以利沙在事奉的需要.
主動以愛接待以利沙.
她提供了以利沙所需的食物和她住的地方.
她這個小小的善行.
將她帶到可以有一個意想不到的回報.
但後來因為她遇到一個很大的危機.
所以她又怎樣去面對呢?.
她怎樣去應對這件事呢?.
各位兄弟姐妹.
在我們生活中遇到困難.
遇到挑戰或危機的時候.
我們會怎樣去反應或面對這件事呢?.
希望今天我們分享這個事件的時候.
我們可以有一點點的學習.
我們一起同心禱告.
我們主我們的神.
讓我們在母親節的主日.
能夠來到您的殿裡.
我們一起學習.
求聖靈這一刻與我們同在.
帶領我們打開我們的眼睛.
打開我們的心.
來接受上帝的說話.
讓我們今天能夠聽得明白.
並且可以走出來.
在上帝的旨意當中.

$^{41}$我們祈禱.
奉主耶穌基督的聖名自祈求.
阿們.
我們先看看第一個重點.
就是不求回報的愛.
就在皇下四章八至十三節.
事件發生在舊約的時代.
樞臬這個地方.
先知以利沙去周遊列國.
去宣講上帝的說話.
有一天他經過樞臬.
遇到一個很富有的婦人.
他很熱情地邀請以利沙到他家裡吃飯.
起初以利沙也是拒絕的.
但婦人盛意拳拳地邀請以利沙.
最終以利沙就接受了他的邀請.
婦人主動地去招待.
熱情地邀請以利沙到他家裡吃飯.
從這次的善行開始.
就成為一個持續的故事.
一個持續的伺候.
以利沙到婦人家裡吃飯.
他丈夫是知道的.
婦人對於跟丈夫說過這件事.
甚至她丈夫也有跟她一起去接待以利沙.
一起去吃飯.
婦人當時有感動地去款待以利沙.
但也沒有忘記要去尊重她的丈夫.
要問他的意見.
想問他之後才去接待這位神人以利沙.
婦人不單有良善的心.
並且對家庭也很忠誠.
以利沙每逢經過書唸這個地方的時候.
大概也會去探望一下這位夫家的太太.
起初以利沙也沒有準備去她家裡吃飯.
但經過好幾次認識了她.
彼此的交情也深了.
以利沙也明白到這位夫家的婦人.
她的真摯和誠意.
所以也放膽接納了她的邀請和款待.

$^{81}$所以在第八節最後說到.
以利沙每次經過就轉到那裡去吃飯.
於是以後每次路過.
他就會去探望一下這位太太.
然後在家裡坐一會兒.
然後再去吃飯.
這就是書唸婦人第一輪的接待.
我們看看第九節.
那裡描寫到婦人對她的丈夫說.
「康,我知道那常從我們這裡經過的是神聖的神人」.
婦人有著屬靈的眼光.
有著分辨能力和對神的熱心.
她的屬靈的分辨能力是她確信.
以利沙是神所分別為聖的一位先知.
我們也要向神多些禱告.
求神給我們有屬靈的分辨能力.
開啟我們的眼睛.
看見我們身邊的人的需要.
我們教會的需要.
我們福音的需要.
而去回應.
書唸婦人和以利沙相處了之後.
見到以利沙其實也是一位處事很穩重的成熟人.
所以她就樂意更加慷慨為以利沙的需要去設想.
想到如果她路過書唸這個地方的時候.
不單可以有一餐飽飯吃.
還有一個潔色的地方是多好呢?.
應該對於以利沙和她的事供.
是會有更大的方便和幫助的.
由於這件事不是很小的事.
所以她也徵得她丈夫的同意.
然後就和他商量.
第十節說到.
讓我們可以為她蓋一間有牆的小閣樓.
安放床塔,桌子,椅子,燈台.
每當她來到我們這裡.
就可以住在那裡.
意思就是說這個地方.
不是用一些帳幕或者一些棚去搭建成的.
而是用一些木頭,石頭.

$^{121}$或者一些瓦磚去砌成一個樓房.
小閣樓是什麼呢?.
其實就是蓋在屋頂上的一個閣樓.
屋頂的閣樓通常在外面有一個樓梯.
上到這個閣樓.
不需要拆開裡面的結構.
或者做一個洞.
不需要在外面就可以進到這個閣樓裡面.
書唸婦人也很想在裡面擺設一些日常用品.
譬如床,桌子,椅子,燈台.
讓以利沙可以到家的時候.
住在那裡也可以有一點點的休息.
還有過夜也可以.
短暫的住宿也可以.
這個婦人的好意.
更是對於她認識的神有熱心和信心的表現.
婦人說服了她的丈夫.
給一位陌生的男士來到家裡吃飯.
最後也留宿在自己家裡的小閣樓.
可以看到她不單認得出.
以利沙是神所分辨為性的人.
而且還將自己的確信付諸行動.
她對神有一個很大的熱心.
我們知道愛不單是一個心動.
並且是一種行動出來的.
婦人認得出以利沙是神聖的神人.
看見神服人的需要也看見神工作的價值.
她與丈夫商量接待安排.
在良善之中也有尊重和忠誠.
讓愛更加有根基.
愛是可以讓人看到的.
可以看到別人的需要.
愛是看到別人的需要.
婦人的服侍不單止停留了一次吃飯.
她為神人預備了一個可以休息的地方.
可以停留的地方.
也可以繼續在這裡工作的一個空間.
真正的愛是會從人的實際需要出發.
並且為神的工作作長遠的預備.
這就是《書斂婦人》第二輪的接待.

$^{161}$以利沙深深感受到這位婦人.
為她所做的這麼多事情.
花了很多心思.
於是以利沙想到.
可以怎樣為這位婦人做多一些事呢?.
有一天以利沙就叫她的僕人基哈西.
請這位婦人上來.
詢問有什麼需要她的幫助.
意思是說如果她生活上有什麼重要的需要.
困難或者是冤屈.
以利沙可以代她向有關的人等.
包括若蘭王或者是軍長去轉達.
《書斂婦人》是怎樣回答呢?.
她就說「我已住在自己百姓之中」.
意思就是說她跟官場的人其實沒有什麼關係.
沒有什麼重要的困難或者是冤屈要她申訴.
所以她不需要以利沙去為她在官庭.
或者在官員那裡打什麼交道去幫她.
婦人已經知足了.
在環境裡已經感到安樂滿足.
一無所缺了.
這位《書斂婦人》款待以利沙.
她用細心和用心去服侍.
但是她不求任何的回報.
她的愛不單止是心動.
並且是一種慷慨的行動.
我們很多時候都會邀請弟兄姊妹.
一起去回應上帝的愛.
不單止有經常服侍的弟兄姊妹.
還有見到一些新力軍.
好像我們敬拜隊多了很多新力軍.
好像兒童暑期聖經班.
前一陣子招募事奉人員.
發現有些肢體.
我們的弟兄姊妹以往報名了.
但是從來都沒有來過我們兒童服侍.
甚至是第一次來參加我們教會的服侍.
裡面有年輕人和退休了的兄長.
當他們看到崇拜報告的時候.
有這個招募.

$^{201}$我深信他們已經將這件事放在心裡.
因為當日崇拜了之後.
分別有兩位弟兄姊妹不約而同地心裡這樣說.
如果今天午堂見到Ivy傳道的時候.
我就去回應這個招募的邀請.
你知道午堂我是服侍兒童的.
一般在午堂我都會不見到午堂的弟兄姊妹.
我就會上去六樓或者在不同的樓層服侍一群兒童.
那天突然之間我心裡有感動.
在十二點半我就下去大堂.
想著和各位弟兄姊妹道別.
我深信這些都是聖靈提醒或者是感動我自己.
就因為這樣.
這兩位新服侍的弟兄姊妹.
她就和我說一番話.
如果今天早上見到我 我就會回應這個邀請.
我會報名來參加.
這些都是神的預備.
另外有一位弟兄他回了祖堂.
我看看他在不在.
OK 我看到他了.
他主動完了之後就上到六樓去找我.
他就問這服侍的內容是什麼.
當日他即時回應了.
就報名去參加填了Google form.
就回應了上帝的召喚.
他們三個人從來都沒有服侍過兒童.
但他們就是這樣憑著這個信心.
勇敢 都要勇敢的去回應上帝的呼召.
上帝的愛.
願意付出了時間 心力和意外去轉化為行動.
不求什麼的回報.
願主親自去報答所有回應服侍主的弟兄姊妹.
這位婦人有屬靈分辨的能力.
我們都要去求屬靈分辨的能力.
向神禱告 求祂開啟我們的眼睛.
看見身邊的人 教會和福音工作的需要.
我們快點去回應.
當心裡面有感動 不要耽誤了它.
哪怕去服侍 上裝備 傳福音 做見證.

$^{241}$都是需要的時候不要熄滅聖靈的提醒.
然後就是以愛去服侍.
無論是新歷君也好.
或者一直去服侍的弟兄姊妹.
你們所付出的時間 心力和忠心.
在神眼中都是自為補的.
雖然婦人不只是停留在她的感動之中.
正如我們剛才分享的那些弟兄姊妹.
他們將她看見的需要.
實際的需要 變成可以被觸摸的接待和供應.
今天我們同樣都是被神呼召.
希望可以將這份愛轉化成信服的行動.
愛從看見的需要開始.
並且化為具體慷慨持續的行動.
當神向我們發出呼召的時候.
有感動去配合神的工作的時候.
大家不要耽誤了.
聖經都教導我們不要熄滅聖靈的感動.
我們看看第二個重點.
就是對神有堅定的信心.
在聖經第四章14-26節.
書簾的婦人表示她一無所缺.
沒有什麼所需要.
她就回到樓下.
以尼沙仍然想起可以為她做多一點什麼.
所以她跟她的僕人基哈西商量.
究竟可以為她做什麼呢?.
這位僕人回答她.
「她真的沒有兒子,她丈夫也老了」.
經僕人一提.
以尼沙就知道可以為這位婦人做什麼.
就是想要求神賜一個兒子給她.
於是他就叫僕人回到樓下.
然後就叫這位婦人上來.
婦人沒有強求過任何什麼事.
對於有沒有兒子也沒有很大的強求.
自己跟她的丈夫和諧相處.
家庭也富裕.
但這件事這位婦人的招待仍然被神紀念.
書簾的婦人就上到樓下.

$^{281}$以尼沙就跟她說.
「明年這個時候你必抱一個兒子」.
神看見她的缺乏.
並且主動將這個恩典帶進她的生命裡.
婦人的反應是不相信以尼沙所說的話.
因為這件事已經遠遠超過她的常識和期待.
救恩是神主動發出的應許.
只要憑信心去領受神所賜的.
這件事就可以成就.
當書簾婦人還不知道自己需要什麼的時候.
神已經預備好這個恩典賜給她.
經文說「沒有兒子,她丈夫也老了」.
似乎這位婦人的年紀也不大.
但她的丈夫可能已經是年紀老邁.
所以無法生產.
但對於當時的希伯來婦人來說.
不生產是一件比較羞恥的事.
不像現在的現代.
所以對於她來說.
生孩子無疑是一個願望.
雖然如此.
但以尼沙問這個婦人有什麼要求的時候.
雖然她沒有說需要一個兒子保存家產的產權.
但她對於自己的境況已經很知足.
我們看到婦人對於以尼沙的服侍的確是全心全意.
而且沒有半點私心.
就算對於以尼沙可以為她爭取一些減稅,優惠.
她也不需要.
反而她最主要的憂慮是什麼呢?.
就是不願意被人欺控.
在第十六節尾已經說到.
後來婦人真的懷孕了.
到了明年某個時候生了一個兒子.
正如以尼沙向她所說的.
神先發出這個應許.
人只要憑著信心去領受.
神的發生不單是為了人的需要而施行.
同時也是為了神的事奉和榮耀而施行.
神喜悅婦人去接待以尼沙.
並且立即賜福給她.

$^{321}$當我們善待最小的一個.
即是善待我們的神.
在第十八節說到孩子長大.
一天出去到父親和修國人那裡.
有學者估計當時這位小孩子大概兩三歲.
孩子找到父親自然很高興.
誰不知不久對父親說.
在第十九節說「我的頭啊我的頭啊」.
顯然這位小孩子開始頭痛.
不舒服或者中暑.
父親因為有事無在身.
所以他吩咐僕人將這位小孩子抱到母親那裡.
母親接過他之後.
大概也沒想到兒子的病情這麼嚴重.
以為只是抱著他.
放在膝蓋上就可以休息一會.
但經文在第二十節最後說到.
到中午就死了.
人生最大的危機.
突然淋到這位家庭裡.
孩子的母親的反應是怎樣呢?.
她的反應是非常冷靜.
表現出很冷靜.
她沒有驚惶失措.
第二十一節說到.
她先不動聲色.
將孩子抱起來上去閣樓.
然後就將孩子放在神人的床上.
關了門出來.
這個時候孩子的母親大概也是忍著拋眼淚.
不讓她的家人知道.
夫人一直都沒有驚動她的丈夫.
或者其他家人.
孩子死的時候.
她只是一個人知道.
連她近身的僕人都不知道.
這位母親就下去樓下.
在第二十二節說.
夫叫她丈夫說.
你叫一個僕人給我牽一匹奴來.

$^{361}$我要趕去見神人然後回來.
二十三節.
丈夫說.
今日不是初一.
也不是安息日.
你為何要到他那裡去呢?.
夫人說.
平安無事.
其實大家都知道.
事件並不是那麼平安.
因為他們的兒子已經死了.
夫人之所以這樣說.
不是否認他的困難.
而是想將裡面最深的一種痛.
帶到神的面前.
相信這些都是神對這位婦人的一種考驗.
丈夫就派了一個僕人和一匹奴給她.
婦人就吩咐僕人.
在二十四節說.
走 趕緊走.
除非我吩咐你不要為了我而慢下來.
婦人吩咐僕人不要慢下來.
專心前往加密山去找以利沙.
她急切地尋求神的幫助.
神人以利沙剛剛也沒有出去.
在家裡沒有外遊.
二十五節說到母親還沒有去到加密山區的時候.
以利沙就很遠就看到這位婦人.
這位母親.
以利沙就對這位僕人這樣說.
在第二十五節說.
看啊 書臉的婦人來了.
她這樣說.
相信以利沙也沒有想過.
這位母親會過來找她.
有點意外.
也猜到她會來找她.
她就叫僕人在第二十六節說.
跑去迎接她並且問候她.
問她你平安嗎?你丈夫平安嗎?孩子平安嗎?.

$^{401}$這位婦人回答她說平安.
同樣她回答這位僕人的問安的時候.
就好像回答她丈夫的問題一樣.
意思其實是說.
現在暫時沒有時間跟你說這件事.
因為她有更重要的事要找神人幫助.
這位母親雖然失去了兒子.
雖然也在經歷一個很大的考驗.
但這位母親沒有失去她的信心.
面對一個這麼大的危機.
她沒有驚惶失措.
也沒有怨天尤人.
而是靜悄悄把她的兒子放在以利沙的床上.
關上門.
不讓人知道這件事.
避免她的丈夫或是僕人看到的時候.
都會害怕起來.
然後這位婦人就馬上去加密山找神人.
她知道孩子是耶和華所賜的.
她有信心只要找到以利沙的時候.
孩子就可以得救了.
這位婦人沒有把真相告訴她的丈夫和以利沙的僕人.
大概是因為她認清了她信心的正確對象是誰.
就是以利沙的神.
在危急的時候.
她不讓這些恐懼和解釋去耽誤她找神人的腳步.
而且要趕著去找他的求助.
她專心去尋求神.
這個行動顯示了她對神的信心和信靠.
即使在困難的時候.
她仍然保持一份信心.
這也提醒我們.
當我們面對危機的時候.
我們需要有堅定的信心.
將問題交回給神.
無論我們的環境是怎樣.
神總是與我們同在.
我認識一位母親.
她是主內的姊妹.
她很努力地求神.

$^{441}$求神賜她一個兒子給她.
她終於生了一個兒子.
但沒多久.
這個兒子很小.
他得到重病.
也很危急.
然後她馬上送他去急症室.
這位母親在外面等待消息.
但她的心情與書廉的母親不同.
她很焦急.
她哭著求救.
她不停地說一句話.
「賞賜的是耶和華.
收取的也是耶和華.
耶和華的名是應當稱頌的」.
她不斷地說.
「賞賜的是耶和華.
收取的是耶和華.
耶和華的名是應當稱頌的」.
這句話是聖經的約伯記.
一章二十一節.
她願意去信服神的旨意.
不論結果是怎樣.
她都甘心情願.
她不埋怨.
結果是怎樣呢.
結果她兒子長大了.
現在還成家立室.
職場裡.
我遇到一個比較大的危機.
大家知道我之前是做IT的行業.
最重要不是趕著限期.
交系統設置或上線.
最重要是做好基本網絡安全.
首先要定好網絡安全的保護程序.
和守則.
並且確保持之以恆.
有一天回到公司.
發現香港有一部伺服器.
中了電腦病毒.

$^{481}$這個電腦病毒.
我們知道是不會傳染給人的.
不要太擔心.
不過連著這些伺服器.
是會一直感染的.
當時我們有十多部伺服器.
連著運作中.
除了香港有一部伺服器出事.
連接網絡回國內的伺服器.
也感染了.
當時我要負責帶領這個團隊處理這件事.
當時我心情是鎮定的.
我自己也沒有打亂陣腳.
因為我知道當務之急是要找出源頭和原因.
並且要停止這個電腦病毒繼續擴散.
要發通告給我們的用戶.
告訴他們伺服器要停止運作.
我們就去找方法救回伺服器的資料.
當我找到這個感染的原因的時候.
我大可以發一個很大的脾氣.
去責罵一位粗心大意的IT的同事.
他沒有按程序去做網絡的安全.
但是我沒有這樣去做.
我選擇了跟團隊一起去解決這個問題.
鼓勵大家一起去面對這麼大的難關.
我們用了大約一個星期的時間去做.
這班團隊都做到很晚.
或者連續沒有睡到很多晚.
特別這位粗心大意的同事.
他差不多每晚都在通宵.
最後我們都能夠救回香港.
三部伺服器救回兩部回來.
那些數據回來.
記得團隊裡面除了我之外.
是基督徒.
還有一位同事是基督徒.
我們兩個都一起為這件事去禱告.
去交託給神.
求神幫助.
我們知道真的需要神蹟出現.

$^{521}$我們不是怕公司會問責.
而是不希望我們公司大部份的同事.
在一兩個禮拜裡面所入的所有資料.
都沒有了.
我記得當我禱告之後.
我心裡都是平安的.
有一份的安心.
知道神必定會幫助我們.
我也有和我公司的老闆去說這件事.
不過我沒有將這位粗心大意的同事的名字告訴公司.
因為我知道.
我希望看到我們這班團隊.
能夠完成一個近似不可能的任務.
這位婦人在面對突如其來的危機的時候.
她冷靜地處理.
沒有將恐慌擴散.
然後急切地尋求神的幫助.
很多人遇到危機的時候.
他們可能都會有.
不知道該怎麼辦.
這位婦人.
主要原因她能夠冷靜地處理.
因為她認識了以利沙的神.
她認識到這位神一定會給她平安.
平安不是環境沒有事.
而是信心有了對象.
她也不是否認痛苦.
也不是因為環境有什麼的發生.
她因為信心有對象.
在一個不能夠承擔的危機上.
她選擇將這些危機交回神的手中.
而且這位婦人時時都有知足的心境.
所以她認定了恩典的源頭.
所以她就交託了.
平靜地去面對.
我們已經是上帝的兒女.
更加要有交託的心.
我們要信得過神.
知道無論有什麼困難.
我們都可以將這個重擔帶到神的面前.

$^{561}$神必定與我們同在.
主是會包底的.
就像我們的sticker.
主包底.
讓我們安然地渡過這個逆境.
我們看第三個重點.
得神恩後不忘這恩.
在皇后的四章27-37節.
27節這樣說.
她的婦人上了山,到神人那裡就抱著神人的腳.
基克西前來要推開她.
神人說:由她吧,因為她心裡受苦.
但耶和華向我隱瞞這事,沒有告訴我.
受苦的母親帶著一個真實的痛去找神.
當這位母親見到以利山的時候.
她抱著她的腳.
一方面表示她的謙卑.
另一方面表示她自己很無助.
心裡很受苦.
需要對方的幫助.
她一直忍著拋淚.
相信這一刻已經掉下來了.
以利山的婦人.
不好意思.
以利山的僕人見到這位母親.
直覺立即想推開她.
但以利山的僕人.
由她吧,因為她心裡很受苦.
以利山見到這位母親舉動的時候.
他就知道這位母親心裡一定有很多.
很難說的出聲.
因為已經抱著她的腳.
如果他覺得這位婦人可以抱著她的腳.
令她心情好一點的話.
就讓她抱著自己的腳.
反正以利山在經文中說.
他不知道應該怎樣對待這件事.
因為耶和華並沒有告訴他發生了什麼事.
既然神沒有跟以利山說發生了什麼事.
那以利山大概會問這位婦人.

$^{601}$發生了什麼事.
那婦人就在28節說.
我何嘗向我主求過兒子呢?.
我豈不是說過不要欺控我嗎?.
意思就是說.
婦人當初既然沒有說.
以利山或神賜予她一個兒子.
神那時候又為什麼要講預言.
說她可以得到這個兒子呢?.
如果現在這個孩子已經長大了.
但這麼小又死了.
為什麼當初給她有這個兒子.
現在就空歡喜一場呢?.
甚至換來一個傷心欲絕的失望.
那又何必呢?.
意思就是說.
你現在叫我得而又失.
不是等於在欺控我或戲弄我嗎?.
婦人這個表達雖然有些抱怨.
也是對神一個傾心吐意的表達.
婦人這樣說.
也是告訴以利山.
她的兒子已經死了.
希望以利山可以做一些事.
可以救一救她的兒子.
神是著重我們單純信靠的心.
不會責怪我們.
那個禱文或禱告裡面的軟弱和無知.
然後以利山就派了婦人去樞廉.
在第29節說到.
要將以利山的像放在孩子的臉上.
在第29節最後說到.
於是婦人就照以利山的吩咐.
馬上回去樞廉.
一到樞廉就馬上去找孩子.
但是這位婦人顯然是不滿足這個安排.
所以在第30節那位母親就說.
我指著永生的耶和華.
又指著你的性命起誓.
我必不離開你.

$^{641}$於是以利山就起來跟著她回去.
她的意思是.
這位母親說我不會離開你的.
她要以利山跟著她一起回到她家裡.
那婦人就當然先回到婦人那裡.
她就馬上上了閣樓.
也是照以利山的吩咐.
要將這個像放在孩子的臉上.
但這個孩子在第31節說到.
他有沒有聲音有沒有動靜.
即是說以利山的手像對孩子來說有沒有功效.
沒有功效.
他不能讓孩子復活過來.
大概是神故意讓孩子暫時未能復活.
神之所以這樣做.
就是因為書斂婦人將這個孩子放下了.
出去找以利山的時候.
完全沒有第三個人知道她委罪證.
這個孩子真的死了.
甚至連書斂婦人這位母親本人.
也可能會懷疑到.
如果這個孩子回去見到她醒了.
懷疑到可能剛才孩子暈了.
或者睡著了.
現在由這位僕人來作一個見證.
知道這個孩子真的已經死了.
見到孩子死在自己的床上.
以利山在第32節.
去到家裡的時候.
見到她死在自己的床上.
以利山在第33節說.
她進去關上門.
只有他們兩個人.
向耶和華祈禱.
同時將書斂的婦人和這位僕人.
叫他們在外面關上門.
以利山要單獨施行神蹟.
他要直接去求問神.
經文沒有說到究竟討問是甚麼.
但相信以利山是求神向他顯明.

$^{681}$他應該怎樣做才能使這位孩子復活.
第34節說到.
他上去伏在孩子身上.
口對口 眼對眼 手對手.
他伏在孩子身上.
孩子的身體就漸漸暖和了.
以利山這樣做.
都是神親自告訴他.
親自指示給他.
而不是他隨便想出這個做法.
所以嚴格來說.
這個做法本身沒有甚麼神奇的魔力.
可以使孩子的身體暖和.
以利山這樣做.
所以有效.
是完全因為神指示他這樣做.
是神的權能的成績.
最後.
第35節說到.
孩子打了七個噴刺.
眼睛就睜開了.
這位兒子終於復活了.
神藉著以利山去施行這個拯救.
使得兒子重新得到生命.
我們看最後兩節.
或者邀請各位影視姊妹.
一起讀出36和37節.
預備起.
以利山叫希雪.
你叫這書念婦人來.
於是她叫了她來.
婦人來到以利山那裡.
以利山雪把你兒子抱起來.
婦人就進來.
在以利山腳前伏伏於地.
向她下拜.
然後抱起她兒子出去了.
書念的婦人.
就從她的門外就進來了.
她立即在以利山腳前伏伏在地上.

$^{721}$經歷了神的拯救之後.
這位母親就以敬拜和感恩.
去回應神這份恩典.
她就抱起她兒子出去了.
可以想像當時這位母親是多麼興奮.
當初未能有孩子的這位婦人.
是神主動賜下這份恩典.
同時也是賜給她這份信心.
當書念婦人的孩子死那一刻.
也是需要這位母親憑著這份信心.
及時去尋求神.
才能夠經歷神的拯救.
恩典不是人爭取回來的成果.
而是神主動賜下來給這位母親.
親自去成就,拯救.
人要憑著信心去領受.
來成就神的救贖.
當我們經歷神的恩典的時候.
也不要忘記我們要感恩.
並且以行動去回應神.
奇妙的事就會發生了.
神蹟其實每一天都會發生.
只要我們細心去留意.
去聆聽神的聲音和神的啟示.
聖靈時刻都會給我們指引.
做一些正確的事.
從而改變我們的生活.
我時常都會這樣想.
如果不是神的恩典.
給我們各位弟兄姊妹成為神的兒女.
我們人生將會是怎樣的呢?.
你想想當你未信主之前.
你的生活是怎樣的呢?.
可能我們經常做一些傷害人的事.
覺得自己都沒有問題.
或者我們都有機會去吸煙.
或者去酗酒.
或者無緣無故都會發脾氣.
可能現在都會有.
不過相對來說都應該收斂了很多.

$^{761}$以我自己為例.
我以前上班下班都很穩定.
下班後都會和朋友去飲飲食食.
學習玩樂.
那一刻都沒有想過.
人生有什麼大的方向.
可能今天因為我信了耶穌.
神也給了我一個恩典呼召.
去回應神.
成為神的僕人.
每個人身上都有一份神蹟.
也是神的恩典.
書斂婦人經歷了兩個神蹟.
由愛開始.
從母到有不能夠生育的婦人.
得到一個兒子.
顯明神主動去看顧她.
未開口說.
神已經應允了她.
然後從死到生.
這位兒子重新去睜開眼睛.
顯明神的權能能夠進入最絕望的時刻.
他作出了兩個回應.
信心和感恩.
信心是交託.
在危機中尋求神.
不會讓這份恐懼成為最後的聲音.
而是將生命交託給神.
榮耀歸神.
經歷過拯救之後.
他承認一切恩典都是從神而來的.
也將榮耀歸給他所相信的神.
信心是在人們尋求神的時候.
恩典就會大大地經歷.
信心是人們尋求神.
我們有信心去尋求之後.
恩典就會有更深的經歷.
當這份恩典臨到的時候.
我們最好的回應.
不是就這樣放下.

$^{801}$而是不要忘記.
並且我們要以生命去繼續去服事.
繼續去敬拜.
\newpage



\section{路得記 1:6-22}
\label{sec:032129}
\textbf{在神的季節中選擇}~~~~毛譚美虹傳道
\newline
\newline
連結: \href{http://www.church.com.hk/acms/content.asp?site=cdc\&op=show\&type=product\&code=032129\&layout=sermon}{\texttt{032129}} ~~~~ 語音日期: 2026-05-09
\newline
\newline
\hyperref[sec:032128]{\small{< < < PREV SERMON < < <}}
~
\hyperref[sec:index_chronic]{\small{[返順時目]}}
~
\hyperref[sec:032130]{\small{> > > NEXT SERMON > > >}}
\newline
\newline
路得記 1:6-22
\newline
\begin{longtable}{cl}
\hline
\hline
章節 & 經文 (和合本修訂版)\\
\hline
1:6 & \begin{tabularx}{0.7\textwidth}{X} 拿娥米與兩個媳婦起身，要從摩押地回去，因為她在摩押地聽見耶和華眷顧自己的百姓，賜糧食給他們。 \end{tabularx} \\ \\ \relax
1:7 & \begin{tabularx}{0.7\textwidth}{X} 她和兩個媳婦就起行，離開所住的地方，上路回猶大地去。 \end{tabularx} \\ \\ \relax
1:8 & \begin{tabularx}{0.7\textwidth}{X} 拿娥米對兩個媳婦說：「你們各自回娘家去吧！願耶和華恩待你們，像你們待已故的人和我一樣。 \end{tabularx} \\ \\ \relax
1:9 & \begin{tabularx}{0.7\textwidth}{X} 願耶和華使你們各自在新的丈夫家中得歸宿！」於是拿娥米與她們親吻，她們就放聲大哭， \end{tabularx} \\ \\ \relax
1:10 & \begin{tabularx}{0.7\textwidth}{X} 對她說：「不，我們要與你一同回你的百姓那裡去。」 \end{tabularx} \\ \\ \relax
1:11 & \begin{tabularx}{0.7\textwidth}{X} 拿娥米說：「我的女兒啊，回去吧！為何要跟我去呢？我還能生兒子作你們的丈夫嗎？ \end{tabularx} \\ \\ \relax
1:12 & \begin{tabularx}{0.7\textwidth}{X} 我的女兒啊，回去吧！我年紀老了，不能再有丈夫。就算我還有希望，今夜有丈夫，而且也生了兒子， \end{tabularx} \\ \\ \relax
1:13 & \begin{tabularx}{0.7\textwidth}{X} 你們豈能等著他們長大呢？你們能守住自己不嫁人嗎？我的女兒啊，不要這樣。我比你們更苦，因為耶和華伸手擊打我。」 \end{tabularx} \\ \\ \relax
1:14 & \begin{tabularx}{0.7\textwidth}{X} 兩個媳婦又放聲大哭，俄珥巴與婆婆吻別，但是路得卻緊跟著拿娥米。 \end{tabularx} \\ \\ \relax
1:15 & \begin{tabularx}{0.7\textwidth}{X} 拿娥米說：「看哪，你嫂嫂已經回她的百姓和她的神明那裡去了，你也跟你嫂嫂回去吧！」 \end{tabularx} \\ \\ \relax
1:16 & \begin{tabularx}{0.7\textwidth}{X} 路得說：「不要勸我離開你，轉去不跟隨你。你往哪裡去，我也往哪裡去；你在哪裡住，我也在哪裡住；你的百姓就是我的百姓；你的神就是我的神。 \end{tabularx} \\ \\ \relax
1:17 & \begin{tabularx}{0.7\textwidth}{X} 你死在哪裡，我也死在哪裡，葬在哪裡。只有死能使你我分離；不然，願耶和華重重懲罰我！」 \end{tabularx} \\ \\ \relax
1:18 & \begin{tabularx}{0.7\textwidth}{X} 拿娥米見路得決意要跟自己去，就不再對她說甚麼了。 \end{tabularx} \\ \\ \relax
1:19 & \begin{tabularx}{0.7\textwidth}{X} 於是二人同行，來到伯利恆。她們到了伯利恆，全城因她們騷動起來。婦女們說：「這是拿娥米嗎？」 \end{tabularx} \\ \\ \relax
1:20 & \begin{tabularx}{0.7\textwidth}{X} 拿娥米對她們說：「不要叫我拿娥米，要叫我瑪拉，因為全能者使我受盡了苦。 \end{tabularx} \\ \\ \relax
1:21 & \begin{tabularx}{0.7\textwidth}{X} 我滿滿地出去，耶和華使我空空地回來。耶和華使我受苦，全能者降禍於我。你們為何還叫我拿娥米呢？」 \end{tabularx} \\ \\ \relax
1:22 & \begin{tabularx}{0.7\textwidth}{X} 拿娥米從摩押地回來了，她的媳婦摩押女子路得跟她在一起。她們到了伯利恆，正是開始收割大麥的時候。 \end{tabularx} \\ \\
[1ex]
\hline
\hline
\end{longtable}
$^{1}$剛才我們唱了一首詩歌.
給我們提醒上帝真的有一個預備和一個恩典在我們每一個人的心中.
但是在這個時間.
我不知道大家弟兄姊妹怎麼看.
我們接觸很多的訊息都是可能是打仗.
或者是一些的霸權.
甚至乎是詐騙.
近一點我們香港有很多災難.
有火災 有失去.
甚至乎我們自己牧者都會傾.
原來我們很快就會主懷安息.
我們會突然間都會主懷安息.
所以在這個世界裡面.
我們不知道究竟我們還有什麼可以抓緊.
或者我們還有什麼可以抓緊是最重要的.
但是在這個這麼特別的時間.
這麼特別的季節.
我們就很想知道上帝提醒我們.
告訴我們.
究竟在這個世代裡面.
上帝叫我們看到這些事情.
我們身邊有的人.
發生的事情.
究竟在演繹什麼呢.
就是說其實發生在我們周圍的事情不是偶然.
是一定有意思的.
我們在這個世代是一定有一些特別的事情.
在我們這個世代裡面去經歷的.
那又跟我們有什麼關係呢.
所以這次我選擇了《路得記》.
那個焦點我剛才說了一次.
我們不是說苦難.
《路得記》裡面有很多苦難.
但是我們這次說選擇.
我們切入選擇裡面.
特別在這個不正常的時代裡面.
我們其實怎樣去尋求上帝.
因為我們在這裡是屬於上帝的.
如果當我們做每一個選擇.
我們都是自己來的話.

$^{41}$其實我們不需要上帝.
但是我們有上帝的時候.
我們就應該將每一個選擇.
都是跟向上帝那邊.
我們就要尋求祂.
當我們尋求的時候.
我們就需要有上帝的指引.
屬於祂的想法.
將我們自己的想法放下.
這個世界的想法拿走.
上帝的想法究竟是怎樣.
神是怎樣看待我這個決定.
我應該怎樣做.
或者我怎樣做才可以選擇到.
在今天這個時間是一個正確的選擇.
因為時間都是很重要.
所以我們看今天的經文.
但是經文很長.
是一章六至二十節.
我在這裡就不會跟大家讀.
大家可以回去看一看.
還有這個背景我會跟大家說一說.
其實整個《路得記》.
特別是第一章那裡說.
路德是一個魔鴨的婦人.
她是一個魔鴨的婦人.
她成為寡婦之前.
她就嫁入了以色列的一個家庭裡面.
路德是《路得記》的主角.
所以路德往往是帶出不同的人物.
在她身邊出現.
其中一個就是那奧米.
那奧米是她的奶奶.
這個奶奶也是一個寡婦.
但是路德對她那份忠心和忠誠的愛.
以至到路德的生命.
是由這個魔鴨走到去伯利恆.
回到以色列.
去到神預備的地方.
讓路德嫁了給那奧米.

$^{81}$她丈夫的親戚叫波亞斯.
如果大家都讀過《路得記》.
有這個印象.
那奧米因為藉著這個歸回.
他有一個後裔.
他可以有孫.
他本來沒有的.
但是因為這樣.
回到去.
不單止有一個後人.
他還可以擺脫貧困.
他們建立的.
不單止是他們自己的家族.
還是神國度的王國.
如果你記得勝經的話.
波亞斯一直數下去的後裔.
就是大衛的血脈.
大衛的血脈一直再數下去.
就是耶穌基督.
所以我們一起看.
《路得記》裡面這三個女人.
怎樣選擇.
而是在神的時間裡面選擇得對.
我們來看看.
第一個選擇是那奧米.
那奧米有很多苦.
其實這個是不可以否認的.
但是原來當我們一個人有苦的時候.
是會阻擋我們做任何決定的.
原來苦難是會影響我們去想.
怎樣做一個決定.
如果我們信徒.
我們信耶穌的.
我們看苦難觀.
其實有很多種.
有很多種就是不同派系.
而那奧米對自己那種的苦毒.
其實正正就是在他失去了所有之後.
他覺得他就是在苦難裡面.
他失去他丈夫.

$^{121}$然後他失去丈夫之前.
其實為何要去摩訶呢.
他是逃避饑荒.
然後他十年裡面.
在摩訶當中失去了丈夫.
然後兩個兒子都一起死了.
這個失去成為他生命裡面的苦.
這個苦是令他看這個苦難觀.
成為他家庭裡面一個被咒詛的一環.
我們來看看.
路德的苦.
剛才我們看那篇經文.
那奧米怎樣說呢.
他講了一番很經典的說話.
他說你們不要再叫我做那奧米了.
因為我很苦.
我很苦是因為你叫我做瑪拉.
這個全能者使我受大苦.
我滿滿的出去.
我空空的回來.
耶和華降禍給我.
全能者使我受苦.
這個大家看不看到.
好像是那奧米一個苦難觀.
那奧米原本的名字叫做愉悅.
那奧米的意思是愉悅.
他要回去以色列的時候.
那班婦女去歡迎他.
但他說我很苦.
你們不要再叫我做那奧米了.
你叫我做瑪拉吧.
因為瑪拉的意思就是叫苦.
這個苦就好像他認為上帝加添了給他.
滿滿的出去.
但我現在空空的回來.
那奧米就是失去之後.
他意識到這個苦影響到他看他自己.
我原來失去了我老公.
失去了我兩個兒子.
我就失去了我的身份.

$^{161}$甚至我應該叫做愉悅那奧米.
我都變為苦.
當原來遇到苦難的時候.
或者困境的時候.
我們很容易在這個位置失腳.
就是說我們會不會在苦難困境中.
我們都會失去了我們的價值.
和我們自己本身的身份.
那奧米的苦怎樣出來呢.
我們說要醫治.
都要找苦毒根.
這個苦毒根就是我們剛才讀的經文.
我們剛才讀到的經文.
那奧米說這個全能者.
希伯來文就是Shaddai.
他形容上帝用全能者.
其實上帝可以用很多的名字去形容他.
不過他正正就是要說.
這個全能者有絕對的權利.
去諷刺我們的神.
你有絕對的權利.
但是你給我苦.
你讓我的生命變苦.
那奧米的困境.
是令到他重新去看這個神.
就算那個全能是多大都好.
在他的生命裡面.
這個全能的上帝雖然偉大.
但是一點都不良善.
全能是不是就是這麼偉大呢.
在那奧米那段時間不是.
這是他扭曲了上帝真實的屬性.
因為他的經歷.
他的經歷是他拒絕上帝給他這個苦.
他甚至拒絕他想放棄這個信仰.
雖然他回去.
但是他裡面真的很苦毒.
去到很多時候在這個苦毒的位置.
有一天有一個姐妹.
我們上上去聊天的時候.

$^{201}$她就說道.
她說我可不可以問一個問題.
我說好啊 你問什麼問題呢.
她說為什麼我們信耶穌的時候.
弟兄姊妹或者我們自己.
都會有時遇上一些問題的時候.
我們會用一種好像拜偶像.
舉個例子拜黃大仙那種心態去信耶穌呢.
這個姐妹的意思是什麼呢.
這個姐妹的意思是.
為什麼.
為什麼信耶穌就是一定要成功.
為什麼信耶穌就是一定得醫治.
為什麼信了耶穌.
我們就不會遇到災難.
要一路平坦.
是不是這樣我們遇到困難的時候.
我們就會立刻說上帝不愛我.
是不是這樣.
這個會不會都是娜奧米當時的苦難觀.
但其實事實上我們在信仰的路途裡面.
我們都見到有另外一個極端.
另外一個極端就是.
我們信耶穌就是必然要受苦.
必然要受苦.
去到另外一個極端是.
我們要和耶穌基督一起受苦.
我們還要拿一個受苦的恩賜.
其實這一句說話我聽到我有點害怕.
其實當我聽到有人求受苦的恩賜.
我真的沒有膽量去求.
但是這個就是我們信仰的裡面.
有時在A點.
但是很多時候也有不同的弟兄姊妹.
在信仰的關口裡面去到Z點.
很大的極端.
但是當這位姊妹這樣發問的時候.
我自己就有感而發.
有感而發是提醒我.
怎樣去思考我現在信的神.

$^{241}$我信這個神和我的生命有什麼關係.
我反省我信神究竟的目的.
我想起我為什麼第一天我決志.
我想起我決志的那天我為什麼會信祂.
是不是因為祂醫治了我.
是不是因為我避了災難.
還是一些什麼原因.
但是在這裡思考的時候.
我就想起這段經文.
這段經文是我很喜歡.
我日常常去記住.
這段經文就是這樣提醒我們.
是耶利米書29章11節.
耶和華說:我知道我向你們所懷的意念.
是持平安的意念.
不是降災的意念.
是要叫你們末後有展望.
這段經文是告訴我們.
神的屬性其中一樣.
祂是持平安.
但是在《瑪太福音》裡面.
我也找到另一段經文.
這段經文告訴我們.
我們要進閘門.
你們要進閘門.
因為引導滅亡的那門是闊的.
路是大的.
但是進去的人是多的.
但是引導去永生的那門是閘門.
路是少的.
找著的人也少.
這個令我們立刻要去思考.
究竟我怎樣去看這個苦難觀.
怎樣去看我面對的難阻.
難處 困境 挑戰.
這是我們需要去選擇.
所以當我們這樣問的時候.
我們遇上生命裡面突然的苦難.
我們就應該去問.
去尋求神.

$^{281}$究竟這個苦是甚麼意思呢.
又或者我們應該怎樣選擇呢.
我們選擇的是進入閘門.
還是去大門呢.
如果說到生命突然而來的苦.
我想起一個很接近的例子.
很接近的例子說一下我自己.
我有兩個弟弟.
我大弟弟結了婚.
他有一個兒子.
他們兩夫婦都很投入工作.
五年前.
因為他們想預備兒子入外國讀書.
在暑假裡面.
三個提早到當地定居.
長話短說.
定居後.
再加上想像力.
定居後.
當時是沒有計劃.
沒有安排.
沒有帶家當.
沒有帶錢的情況下.
他們在不預期之下.
三個人要留在當地.
這個你想像一下.
我不說得那麼實.
想像一下那個困境.
他不能離開.
要留下.
這個會不會算是.
放入人遇上生命突然遇到的挑戰.
困難 甚至是苦難的時刻.
他們三個在外地的狀況.
有很多東西可以說.
不如我就說一下.
在留港的家人的感受.
留港的家人的感受是甚麼呢.
其實當時的情況.
我們收到消息.

$^{321}$他們不能回來香港.
我爸爸媽媽就馬上感受到.
突然間失去兒子.
是很短暫的.
不知道何時可以再見到他.
失去兒子 失去孫子 失去媳婦.
失去一個家人.
而我做姐姐的呢.
我做姐姐的.
其實我能夠了解的也很少 很有限.
我只可以有少許的關心問候.
其實我真的不太懂.
她也不說得多.
所以少許問候 少許關心.
甚麼都做不到.
但事實上.
其實他們兩夫婦當時.
都不知道應該怎樣走下一步.
只知道只可以停頓在那個地方.
這個決定最清楚方向的.
我相信是我侄子.
他當時入中一準備.
我侄子順理成章.
他們三個去是為了讀書.
他就真的留在那裡讀書.
這個突然的生命轉變.
是不預期的.
其實對我來說 我覺得是一個很大的trauma.
是不知道怎樣去面對前路.
三個生命的前路.
和我們家人的狀況.
一點都不容易 不簡單.
但又很真實.
甚至我會用人生的災難.
或者大事件簿去形容這件事.
這個故事未完的.
我一會兒再說下去.
所以當我們無論在苦難或者災難.
出現在我們人生的那條path的時候.
我們需要改路.

$^{361}$我們需要本來在走的路是要改的.
所以苦難困境也好.
我們在那個時候是需要尋求.
一定去尋求幫助.
尋求指引 尋求方向.
不過我們怎樣尋求.
去尋求什麼人 或者怎樣尋法.
這個是關鍵.
所以當我們去尋求的時候.
我們需要尋求上帝的指引.
這是一個很必然的 和是必須的.
弟兄姊妹.
我們不要在一個關鍵的時刻.
一個人背了一個重要的擔子.
這個可能會成為我們後悔的決定.
或者是一個影響很多很多人的決定.
所以當問題來到拿訛米的身上.
他失去了一切.
神將一些人放在他身邊.
跟他去說.
若果我們有這樣的狀況的時候.
我相信上帝都一樣.
將有他的天使.
他差派來的人 會放在我們身邊.
不單止聽我們的禱告.
他還會用不同的人和環境來幫助我們.
所以當時拿訛米是跟他身邊的人.
去說他的決定.
第二 我們看看另外一個決定.
就是俄爾巴.
俄爾巴是一個不顯眼的角色.
我們看路得記.
我們會記得路德和拿訛米.
但我們不會記得俄爾巴.
不過我因為預備.
我現在很記得他 俄爾巴.
俄爾巴 我說他用了一個最安全的選擇.
為什麼呢 當時俄爾巴 路德和拿訛米.
他們都是面對相同的處境.
他們相同失去了老公.

$^{401}$失去了一切.
貧窮.
接著下一步怎樣走.
但俄爾巴他做了一個決定.
他做了一個很安全又很合理的選擇.
俄爾巴他是摩訶人.
所以他回到摩訶是很合理.
而在他的身份裡面.
如果一個寡婦去到一個陌生的地方.
其實對寡婦來說是一個很不利.
也沒有保障的決定.
當時拿訛米就提議他們.
俄爾巴和路德回到你的家鄉.
他們兩個都說不要哭著 不捨得.
但最後俄爾巴是接受了拿訛米的提議.
其實很正常.
而拿訛米提出這個提議.
也是很正常.
他是出於善意.
他不想這兩個媳婦.
和他走這條漫長的苦路.
因為在當代的處境.
他們兩個這麼年輕的婦女.
回到當地的地方.
他們可以再婚再嫁.
以致他們有歸宿.
可以有一個家庭去養活她.
但是因為現實和這份同情心.
拿訛米就說不要跟我走.
你們回到家鄉.
跟誰走就要走在對的位置.
俄爾巴就選擇回到家鄉的安全的路.
這條路是每個人都會接受的.
因為這是當時當地處境的習俗.
如果有這樣的遭遇.
他們回到家鄉是很自然不過的事.
他相信人的期望.
習俗的規矩.
所以他做這個決定是很安全的.
而他也聽他的婆婆說.

$^{441}$聽他的奶奶說.
我叫你去.
他就跟.
他跟回去的時候.
俄爾巴其實做了這個決定.
他代表什麼呢.
他代表當我們面對一些挑戰難關困難的時候.
他選擇走在一個比較安殊或者穩陣的方向.
而這個安殊穩陣的方向是比較熟悉這個世界的.
而俄爾巴就選擇了這個穩陣的方向.
在那個安殊位裡面去決定.
繼續完成他的人生.
所以他的路是安全的.
是有保障的.
當他回去的時候.
俄爾巴選擇了這個安全的選擇.
我們有一個反思.
我們的反思是什麼呢.
如果當我們自己的信仰.
我們遇到一些選擇性的挑戰.
我們面臨要抉擇的時候.
我們會選擇對我們自己比較舒服的.
或者是安全的.
或者是世人都會說.
OK的你選擇這個沒有問題.
是大多數人都會這樣跟你說的.
我們到那個時候我們會怎麼選擇呢.
俄爾巴的選擇其實是.
Not evil.
他沒有問題的.
絕對沒有問題.
而我們看到俄爾巴這個選擇.
反而是給我們一個提醒.
我們是不是會做這個選擇.
就是說當我們有這些困難的時候.
我們是不是做一個選擇.
安舒 安全.
跟隨世界的選擇.
我們怎麼去看到呢.
這個反思令我們怎麼可以再這樣想下去呢.

$^{481}$因為當俄爾巴去決定的時候.
我不是告訴他.
你回到你的家族裡面.
你是沒有問題的.
但是他的生命呢.
就是停止了在那個位置裡面.
他的生命那個屬上帝.
耶和華神.
不是摩訶的神.
是耶和華神.
那個接觸就停止了在那個位置裡面.
所以當他選擇安舒的時候.
讓我們去看看我們自己的生命的狀況.
我們都會不會有一樣的處境在裡面.
我們有這樣的思考.
去到第三個選擇.
第三個選擇就是路德.
路德是這個港島裡面的主角.
也是他的選擇是我們要留心的.
路德裡面的十六至十七節是很經典的經文.
我讀一次給大家聽.
路德說不要催我回去不跟從你.
你是拿我米.
你往哪裡去 我也往哪裡去.
你在哪裡住宿 我也在哪裡住宿.
你的國就是我的國.
你的神就是我的神.
你在哪裡死 我就在哪裡死.
你葬在那裡 我也葬在那裡.
除非死才能夠使我和你相離.
不然耶和華就重重降罪於我.
這句是今天我們的焦點.
是神怎樣讓路德去藉著他的回應.
以致選擇一個走向命定的選擇.
生命裡面我們每一個人都有一個命定.
或者在不同季節的選擇裡面.
上帝是會教我們引導我們.
去到神預備給我們的生命的計劃當中.
但是我們會不會有苦難去阻礙我們.
我們沒有走到原本上帝預備的那條路呢.

$^{521}$或者第二個位置就是.
我們會不會選擇一個比較安全.
安恕的決定 以致上帝的那條路.
我們都還未可以進入呢.
但是當路德選擇這條他命定的路的時候.
我們看看他怎樣會選擇這條路.
和裡面當中怎樣去影響他.
第一個就是我們看到.
路德原來不只是跟著.
他是會行動 因為他承諾.
他承諾給他的婆婆.
那俄米 他亦都承諾給上帝.
這個是他的回應.
回應是終身和終成的回應.
這個是他對那俄米剛才說的那番說話.
這個承諾是很堅的.
不是說我跟你走就跟你走.
而是你去哪裡我去哪裡.
你住去哪裡我跟你住去哪裡.
這個很大的不確定性.
我們是會不知道的 但是他都願意.
你甚至你死 我才跟你分離.
這個是一個很大的毀身.
這個我相信亦都不是路德一時的衝動.
他是經過深思熟慮.
在這十年裡面他怎樣和那俄米相處.
他怎樣開始認識那俄米的神.
這個我想是一個重點 是我們要看到的.
但是對今日的信徒 今日我們的挑戰.
我們在信仰裡面.
當我們真的遇到需要選擇的時候.
其實我們很膚淺.
我們真的會很膚淺.
很多時候我們會將信仰.
視為可以選擇的生活方式其中一部份.
信仰只是一個信仰.
需要的時候我就跟.
去到某些位置的時候我就放下.
這個會不會都是反照了我們的情況呢.
都是給我們一個提醒.

$^{561}$順利順境的時候很熱情的.
但是一到困難 我們開始退縮.
但是路德的榜樣給我們去思考.
思考是甚麼呢.
我們對上帝的毀身 是不是真的無條件呢.
如果你說 嘩 用不著說得這麼厲害.
是不是真的要毀身到這麼committed.
無條件.
如果我再問.
或者我們現在是去到哪一個毀身的狀態.
好嗎.
我們究竟有多毀身呢.
我們會不會現在都可以思考.
我們跟上帝的關係呢.
如果我們說 當一個人進入人生的命定.
是在上帝的計劃裡面.
其中一個重點就是.
我們能夠毀身在上帝的當中.
這個是我們需要考慮的.
第二個是在逆境裡面我們怎樣去見證.
其實當在逆境裡面我們能夠成為見證.
是一個很勇敢的行動.
因為我在逆境裡面還敢認我這個神.
如果我是成功的時候.
我就說 嘩 我的神保護我.
祝福我 保守我.
豐盛的人生在我裡面.
但當我有苦難的時候.
我困苦的時候.
我逆境的時候 我迷失的時候.
我怎樣去認這個神.
說這個就是我的神呢.
如果在路德裡面.
他是一個寡婦.
他是一位外邦人.
他只是一個女人.
在當時的處境 他是很卑微.
卑微到很貧困.
他很窮.
他是窮到飯都沒得吃.

$^{601}$那種的悲傷情況其實是很困難的.
但當我們在困難的時候.
在逆境的時候.
我們會不會都是藉著這個位置.
就是更加去偏向親近我們的神.
以致到我們可以公開承認.
我們在這個逆境裡面.
我們都是靠著我們的主呢.
可能你們都會說 我知道的.
詩歌都唱過很多.
神是我們的避難所.
是我們的依靠 是我們的幫助.
不過當真是你很困難的時候.
那個苦毒根會不會先出來呢.
還是這個信靠在這邊強過那個苦毒根呢.
這個是一個拉扯來的.
弟兄姊妹 這個很真實.
因為在逆境的時候 記不記得.
我們是很容易身份就被拉走了.
我們會忘記了我們是天父的兒女.
甚至我們會說 上帝你是不是真的那麼愛我.
記不記得第一個point.
這個是一個關鍵的時刻.
也是在危機裡面.
我們最容易在逆境中影響到.
我們自己怎樣去繼續持守.
我們基督徒的價值觀和我們的行為.
好像撓米一樣.
撓米他叫路德.
喂 你走啦 不要跟我那麼苦啦.
他叫他離婚啦 即時.
再娶嘛 再嫁嘛.
其實這個都往往可能是我們在現代裡面.
其中一個選擇來的.
我們會不會選擇一個容易一點的呢.
會不會選擇一個.
甚至乎在道德上 價值觀上.
都不是合乎我們信徒.
神的那個心意呢.
所以當路德有這個勇氣去表明.

$^{641}$他說 No.
奶奶 婆婆 我要跟你.
你去哪裡 我就跟你去哪裡.
這份力量不是他一個人來的.
我相信是這十年裡面.
他怎樣看到這個神和撓米一起走.
哪怕他們是經歷那麼慘的遭遇.
所以當我們見證到.
在逆境裡面見證的力量就是.
當我們能夠成為一個信徒.
而我們是很軟弱 很脆弱的時候.
我們怎樣去回應上帝.
這個就是一個生命的見證.
其實是不容易.
但是我們是做到的.
是可以靠的.
因為我們靠的不是我們自己.
我們靠的是我們很堅那個全能的神.
所以接下來的那一點.
接下來的那一點除了是.
那個毀身.
逆境中我們有見證的能力之外.
路德也讓我們看到.
原來屬神的子民是一個家庭.
是一個家庭.
我們的教導天父是我們的阿爸.
你的阿爸父也是我的阿爸父.
教會是耶穌基督的身體.
即是說我們這裡每一個都是手手腳腳.
鼻子眼睛是肢體來的.
沒有了一個都不行.
聖徒是弟兄姊妹.
這就是我們叫大家是弟兄姊妹.
聖靈是我們一起合一的印記.
所以上帝三一的神很清楚.
在我們的生命 在我們的教會.
在我們的關係裡面.
是一個沒有血緣的關係.
但是我們可以合一地一起去建立.
是超越血緣的關係的親密.

$^{681}$這個怎樣可以看到呢.
我剛剛就和我們宣教部的主管.
Pastor Betsy 聊天.
他剛剛又去完一些宣教的活動.
他告訴我他當時在會議裡面.
他有三個陪伴者.
這三個陪伴者是他出工場去招師閣.
那麼久到現在即是十幾年了.
他們三位都是很忠心地陪著他媽媽.
因為他在香港有一個媽媽.
當時爸爸還沒有去世.
都會去探望爸爸 探望媽媽.
他還有一個弟弟.
他弟弟有病.
他三位是定時輪住去關心這三個家人.
到爸爸去世.
他們有很多的支援照顧.
以致到Pastor Betsy.
在這十幾年裡面的宣教工作.
他可以去到很遠.
去到上帝叫他去的地方.
他那天和我午餐.
他和我分享.
他說當日在會議裡面.
這三位支援者.
即是Prayer Team.
他們那天是第一次.
在一些聚會三個都出現了.
他說他很感動.
那一刻神替他.
叫他們出來去honour這三位姊妹.
為他們禱告.
其實這三個是我們不認識的人.
他們三個只是在背後支持著宣教士.
但是對Betsy來說.
是一個很重要的人.
Betsy告訴我.
她比我的親人更加重要.
因為我不在的時候.
他們去照顧我的父母.

$^{721}$我不在的時候.
他去支援.
他有什麼需要.
他即時去達到.
來讓Betsy可以放心.
繼續完成他面前的事才回來.
這個是很真實.
還有是一個.
在天父裡面我們才有這份關係.
就好像今天.
這裡的弟兄姊妹一樣.
我相信YA.
你們的群體.
有一班很熟的弟兄姊妹.
一起成長.
結婚.
生孩子.
大家共同渡過不同的階段.
但是這個屬靈的關係.
和這份的合一.
是很寶貴的.
為什麼你們會走在一起.
就是這位信得過的神.
我們每一次去宣講.
去宣認和去見證.
都是因為這個神在我們中間.
讓我們能夠連繫.
能夠去見證到.
我們每一個生命的時間.
每一個生命的點滴.
最後.
納瓦米她的選擇.
也彰顯了神的旨意.
其實路得記.
一定是講到耶穌那裡.
是一定的.
這個很簡單的故事.
但往往是最重要的故事.
因為路德最終是帶到他救贖.
大胃王的祖先.

$^{761}$耶穌基督.
是他的後裔.
所以當我們提到.
我們現代的信徒.
我們做一些決定.
或者我們做一些回應.
我們好像無關痛癢.
或者是.
沒有果效的犧牲.
有時甚至乎是.
但原來在上帝裡面.
是有一份計劃.
弟兄姊妹.
這個計劃是很深遠.
我們可能這一刻不知道.
但最後上帝是知道.
所以當我們說.
原來路德他選擇.
去跟從他婆婆.
回到以色列的時候.
這個救贖的故事就能夠圓滿.
因為這個計劃.
在上帝的裡面.
是預定了.
讓這些人有份.
但當然.
這個也是我們的選擇.
我們再總結去看看.
這三個女性在路得記裡面.
她們做的選擇.
三個女性揭示了她們的身份.
她們怎樣的身份認同.
和她們在屬靈裡面的命定.
那奧米是一個失喪的見證者.
那奧米的名字叫如月.
但他變為虎.
但他命定裡面.
就不是停在虎那裡.
因為他回到以色列的時候.
我們聖經裡面沒有記載很詳細他的故事.

$^{801}$但我們相信.
那奧米是一個敬畏上帝的人.
他愛人.
所以路德會跟著他.
去守律法.
他覺得我守了.
神也降災給我.
他雖然身苦獨根.
但他也有生命的好見證.
他堅定地相信.
他要跟隨這個神.
他最後在苦難裡面.
他還是回到上帝的命定裡面.
他成為了在這個關鍵.
這三個女人裡面.
一個橋樑.
帶領她們怎樣去認識這位真神.
也帶領她們怎樣去見證.
她們能夠在神裡面.
走出這個命定的計劃.
所以那奧米是一個失喪.
但他也是一個見證人.
第二就是.
俄爾巴.
剛才我們說到俄爾巴.
他選擇安全.
他選擇妥協.
但安全妥協也有代價.
俄爾巴其實有個.
希伯來文的翻譯.
意思是硬頸.
當我知道是硬頸之後.
我就心想死硬.
因為每次說頸看硬.
或者硬頸.
都是沒有好結果的.
這個象徵甚麼呢.
其實象徵俄爾巴.
他是一個不屈服的性格.
硬頸這個名字.

$^{841}$他選擇留在熟悉的地方.
安屍的地方.
所以他選擇.
到最後.
停留在回到摩格的時間.
他的生命和上帝的觸摸.
就在那裡停留.
之後就是他和摩格.
拜的神的關係.
所以俄爾巴去做這個決定.
在我們怎樣去看待他.
我們不可以指控.
我們沒有指控.
我們亦都沒有說.
你這樣選擇就是壞事.
我們不用這個角度.
我們會認為他可以選擇.
而因為這本聖經裡面.
都沒有責備過俄爾巴.
他這樣去選擇.
以致我們值得去說他不對.
但是他只是損失了他那份.
他那份是和上帝繼續的計劃.
但是去到路德.
路德就是一個超越身份的偽身.
路德的意思是朋友.
他選擇的是.
原來他正在走一條引領.
他去救贖計劃的道路.
那個命定.
路德為了什麼呢.
他沒有想到那麼遠.
他為了陪著這個值得去陪伴的奶奶.
這位敬畏上帝的婦人.
這位是就算他面對很多苦.
他都會回到神身邊的女子.
所以當他和拿米一起.
去經歷這個狀態的時候.
他願意和她一樣獻出他的生命.
毫無保留的.

$^{881}$捨棄的.
因為他可以和挪耳巴一樣.
回去嫁個人.
可以做安全的決定.
他沒有.
但是路德的象徵是什麼呢.
象徵就是.
好像我們一班外邦的基督徒.
我們選擇怎樣去.
偽身在神的裡面.
我們藉著不同的經歷.
去認識了上帝.
但是我們怎樣繼續去偽身.
而是學習怎樣去無條件.
在上帝給我們的這份愛裡面.
繼續去見證他.
所以影響這三位女士走的路.
是有很大的清晰.
被我們看到了.
拿奧米成為一個見證者.
挪耳巴成為一個警示.
一個警告.
對信徒的警告.
路德成為一個救贖的樞紐.
他們的選擇是讓他們自己個人.
和他們可以承受的事情.
作出一個很關鍵性的決定.
上帝其實我們看到路德的故事.
他是允許人面對苦難.
面對試煉.
但他不是給你.
他不是糟蹋你.
但是會有這些事情發生.
天災人禍.
各樣的事情都會在我們身邊發生.
但是我們是會.
因為這些事情顯明了我們內心.
我們對神的生命的信仰.
和我們自己生命的動機.
這是很值得我們自己去檢視.

$^{921}$我們自己內在的生命.
而這三位女士.
這三位女人.
她們做的選擇就實現了.
她們自己的決定.
最後.
剛才我講到我弟弟.
差不多結束了.
我弟弟五年之後.
現在五年後.
她們的故事是怎樣呢.
她們夫婦面對困境.
其實和五年前的困境.
她們選擇了留下.
她們的選擇.
其實是可以有不同的選擇的.
但她們選擇留下.
她們回不來.
但可以做不同的事情.
但她們留下所做的選擇.
我們是看到上帝的美善.
怎樣看到上帝的美善呢.
因為這個突然來的改變.
我弟弟我弟夫.
其實他們一直在香港.
很忙很忙很忙.
超級忙地工作.
也超級超級超級使命感.
為他們的工作去投身.
當她們決定留下之後.
她們兩個人有一個全新的生活方式.
其實這個是可以.
選擇另一個生活方式和態度.
但她們選擇了這個.
選擇了誰呢.
她們選擇了那個陪伴她的兒子.
一起進入這個青春期.
進入這個很重要的成長階段.
她們三個人一起去適應新環境.
因為其實她們兩個忙到.

$^{961}$以往在香港都很難看到兒子.
但在那一刻她們留下.
她們做了這個選擇.
一起去過新生活.
一起去面對這個挑戰.
一起去調節她們的生活節奏.
重新去學怎樣三個人一起相處.
更多時間他們一起生活.
這五年裡我自己去演繹.
我沒有問過她們.
是她們一家三口的關係.
一個建造的季節.
本來是一個災難.
本來是一個挑戰.
本來是一個苦難.
但我看到上帝的心意.
願是好敵.
讓她們三個進入一個關係.
可以趁這些時間.
原來是可以建造的季節.
她們也是進入一個彼此滋潤.
彼此認識的季節.
很重要.
如果這裡有爸爸.
這裡有媽媽.
我看到很多初生嬰兒.
你們的時間對你們的孩子是很重要.
她們也結出了一個家庭愛的果子的季節.
在裡面.
我侄子應該是成功過渡青春期.
他已經十六歲了.
他去到當地的時候是十一歲.
我看著他的相片.
看著他臉書的頭髮長.
然後長大.
這個就是他.
更出人意外的是這個苦難.
原來神的心意.
不只是他們三個親密關係.
而是帶他們進入救贖.

$^{1001}$我弟弟一出來工作.
就誓一禮是天主教徒.
但從來不回教會.
從來不會親近上帝.
我弟父因為很忙工作.
很投入.
他也沒有時間去認識這一神.
但有一次機會.
他離開香港之前認識了.
但也沒有時間去接觸.
但當他們三個一起在一個空間裡面.
在這個季節的時候.
他們選擇了在這個地方停留下來.
聽一下上帝的聲音.
他們每一天的星期日.
他們會去彌撒.
回來傳訊息給我.
讓我看到他們去彌撒的訊息.
和我分享.
我很感動的是.
上帝的心意原是好的.
因為人可以看到很多災難.
很多苦難.
但原來當我們走到神的美善.
神的心意.
神的命定裡面的時候.
我們不需要在那裡掙扎.
我們今天一起去聽.
一起去看到老德.
看到今天的訊息.
我想告訴大家.
我今天沒有一個方程式.
來教大家怎樣進入神國的命定.
或者我也沒有一個魔法.
可以告訴你.
在你選擇的時候.
這樣選你就一定是中了.
我反而今天是想提醒在這裡每一位.
我鼓勵你們更多思考.
你們和神的關係.

$^{1041}$你信神的目的是甚麼呢?.
然後.
神在你的生命裡面的心意其實是甚麼呢?.
你不如問一下祂.
你在我的生命裡面.
你想我成為一個甚麼樣的人?.
我邀請你更多的禱告.
和我們的神聊天.
你問祂.
如果你這一刻都有選擇的需要.
我應該怎樣選?.
我怎樣選擇才是合乎神.
你在我生命裡面的命定.
我能夠踏入這條路.
而不是走了下去.
我不能夠再有一個遇見.
\newpage



\section{約翰福音 19:25-27}
\label{sec:032130}
\textbf{看,你的母親}~~~~李頌華傳道
\newline
\newline
連結: \href{http://www.church.com.hk/acms/content.asp?site=cdc\&op=show\&type=product\&code=032130\&layout=sermon}{\texttt{032130}} ~~~~ 語音日期: 2026-05-10
\newline
\newline
\hyperref[sec:032129]{\small{< < < PREV SERMON < < <}}
~
\hyperref[sec:index_chronic]{\small{[返順時目]}}
~
\hyperref[sec:032131]{\small{> > > NEXT SERMON > > >}}
\newline
\newline
約翰福音 19:25-27
\newline
\begin{longtable}{cl}
\hline
\hline
章節 & 經文 (和合本修訂版)\\
\hline
19:25 & \begin{tabularx}{0.7\textwidth}{X} 站在耶穌十字架旁邊的，有他的母親、姨母、革羅罷的妻子馬利亞，和抹大拉的馬利亞。 \end{tabularx} \\ \\ \relax
19:26 & \begin{tabularx}{0.7\textwidth}{X} 耶穌見母親和他所愛的那門徒站在旁邊，就對母親說：「母親，看，你的兒子！」 \end{tabularx} \\ \\ \relax
19:27 & \begin{tabularx}{0.7\textwidth}{X} 又對那門徒說：「看，你的母親！」從那刻起，那門徒就接她到自己家裡去了。 \end{tabularx} \\ \\
[1ex]
\hline
\hline
\end{longtable}
$^{1}$弟兄姊妹平安.
今日首先是主日.
我們敬拜的對象是上帝.
所以讓我們轉眼去仰望.
這位三位一體的上帝.
讓我們同心祈禱.
瑪利亞的神.
你使用侍女瑪利亞.
去養育子上帝耶穌基督.
在祂三十多年的人生裡.
陪伴著祂.
更加跟隨祂成為門徒.
所以在這一刻.
求天父你差派聖靈幫助我們.
讓我們明白跟隨基督是甚麼一回事.
讓我們學習.
像瑪利亞一樣跟隨主.
我們這樣祈禱.
奉主耶穌基督得勝之名而求.
阿們.
假如你是一個正在上班的打工仔.
你每天的身份.
就是一個公司的員工.
不過突然有一天.
你心血來潮想報一個學位.
心血來潮會遞交申請.
然後會有面試.
然後會考筆試.
最後學校應該會寄一封信給你.
說恭喜你.
我們現在出offer給你.
你可以讀這個學位了.
就在那一刻你的身份就轉變了.
你要辭工去入學了.
你的身份不再是一個打工仔.
你的身份是一個學生.
你有教授的指導.
你可以入學校的圖書館.
你可以去餐廳.
有學生假的飯吃.

$^{41}$你可以有你身邊的同學.
你是在學校的群體當中生活.
你的身份不再是打工仔.
你是一個學生.
同樣地今日這段經文.
都是在說一個身份的問題.
耶穌的死亡.
成就了我們今日身份的轉變.
我們從前是罪人.
如今我們成為了門徒.
進入上帝所賜的新家庭.
在這個新家庭的蔭蔽底下去生活.
而這個新家庭是怎樣的呢?.
今日有一個榜樣可以給大家參考.
這個正正就是瑪利亞和主所愛那門徒.
去作為示範.
今日這段經文主要是講十九章二十五至二十七字.
可以聽我讀出這個短短地.
站在耶穌十字架旁邊的.
有他的母親姨母.
格羅伯的妻子瑪利亞.
和莫大拉的瑪利亞.
耶穌見母親和他所愛的那門徒站在旁邊.
就對母親說.
夫人看你的兒子.
又對那門徒說看你的母親.
從那刻起那門徒就接他到自己家裡去了.
首先簡介一下這卷福音書其中一個很重要的主題.
若行福音是一反舊社會的傳統.
女性的角色不再是一個被動.
不再是卑微的身份.
相反女性的角色經常與男性的角色相比.
事實上在福音書中記載.
女門徒在多處都表現得比男門徒更加優秀.
例如撒瑪利亞井旁的婦人.
對比起尼哥底姆.
又或者是拉薩路的姐姐.
瑪利亞用香膏高抹耶穌.
對比起綠色的猶太.
而今天第一位在福音書中第一個出場的女門徒.

$^{81}$就是耶穌母親瑪利亞.
她最先在若行福音二章.
水變酒的故事登場.
水變酒的故事大家很熟悉.
我就這樣讀一至五節給大家聽.
第三天在加里利德加拿有一個婚宴.
耶穌的母親在那裡.
耶穌和祂的門徒也被請去赴宴.
酒用完了.
耶穌的母親對她說.
他們沒有酒了.
耶穌說:母親,我與你何干呢?.
我的時候還沒有到.
她母親對傭人說.
她告訴你們什麼.
你們就做吧.
耶穌如何與瑪利亞這個女門徒互動呢?.
第二節她這樣說.
她說:他們去參加婚宴.
但是婚宴沒有酒供應.
瑪利亞對耶穌說:沒有酒了.
耶穌如何回答她的母親呢?.
你可能會發覺很奇怪.
因為更加原本的翻譯.
應該是:婦人或女人.
我與你何干.
我的時候還沒有到.
用不著這樣吧.
沒有酒喝而已.
你不需要這樣說.
關我什麼事.
婦人或女人.
為什麼耶穌這樣叫自己的母親呢?.
耶穌下半句給我們一個提示.
祂說:我的時候未到.
耶穌的母親上前來.
請求耶穌視耶穌是一個行神職的人.
多於是自己的兒子.
瑪利亞憑著信心上前請求耶穌出手.
去解決飲品短缺的問題.

$^{121}$瑪利亞憑著信心去找耶穌幫忙.
對比起耶穌的血緣兄弟.
在第七章裡這樣說.
連他的弟兄們也不信他.
而耶穌用類似的回應.
去回答他的親兄弟.
我的時機還沒有到.
你們的時機卻隨時都有.
世人不會恨你們.
卻是恨我.
因為我指正他們的行為是惡的.
你們上去過節吧.
我現在不上去過這節.
因為我的時機還沒有成熟.
耶穌多次這樣說.
「時候未到 時候未到」.
即是什麼時候才是時候到了呢?.
作者不斷重複「時機」這個字.
是為了什麼呢?.
什麼時候才是時候到了呢?.
如果我們看第十三章.
會有少許提示給我們.
預節以前耶穌知道自己離世歸附的.
「時候到了」.
耶穌的時候是上十字架死亡的時候.
這裡也解釋了為何在加拿大婚宴.
耶穌說自己「時候未到」.
因為當耶穌真的行神跡.
開始傳到的時候.
祂的死期就開始精準地倒數.
我們回到二章水變酒的故事.
耶穌對瑪利亞的回答.
其實是約翰福音的作者埋下了一個很重要的伏筆.
「母親,我與你何干呢?我的時候還沒有到」.
但第四節這一句.
如果再直譯地翻譯的話.
應該是這樣的次序.
「為什麼是我和你?女人?」.
約翰巧妙地用「女人」這個字而不是「母親」這個字.
而在這一卷福音書裡.

$^{161}$第二次約翰記載耶穌叫自己母親做「女人」.
正是在十九章二十五至二十六節.
我只讀二十六節.
「耶穌見母親和她所愛的那門徒站在旁邊就對母親說.
「母親」或者譯「婦人」,「漢」,「你的兒子」.
耶穌在第二章第一次叫母親做「婦人」或「女人」.
而在這裡十字架上被釘的耶穌第二次叫自己母親做「婦人」.
當然不是很明白為什麼耶穌要這樣叫自己母親做「女人」.
耶穌這樣叫似乎是刻意拉開了自己和瑪利亞的關係.
好像拉遠了一點.
不過可以肯定一點的是.
耶穌和瑪利亞的關係除了母子關係外.
更加是老師和門徒的關係.
因為一開始加拿大婚宴.
耶穌和瑪利亞說「時候未到 關我和你什麼事?」.
在此時此刻 耶穌的回答是回答瑪利亞.
「時候到了」 就是祂死的那一刻.
瑪利亞一直跟隨耶穌.
由祂出世到祂出道,傳道,被捕,被釘,死亡.
瑪利亞一直跟隨到尾.
瑪利亞成為了門徒之中一個非常重要的榜樣.
她一直默默地跟隨耶穌到底.
去見證著耶穌的「時候到了」 就是祂死的那一刻.
當耶穌臨終的時候.
祂才交託祂的母親給另一個跟隨到底的門徒.
就是主所愛那門徒去照顧.
在這裡 死亡中的耶穌.
祂不僅為瑪利亞提供一個保護.
即是在那門徒的照顧下的保護.
耶穌的死亡更加為瑪利亞提供拯救.
不過耶穌的死亡不僅為瑪利亞提供拯救.
祂更為全人類提供了一個拯救.
此時此刻 耶穌的身份.
祂首先是世界的拯救者.
然後才是瑪利亞的兒子的身份.
耶穌首先是世界的拯救者.
之後才是瑪利亞的兒子的身份.
如果相反站在瑪利亞的角度的話.
祂首先是耶穌的門徒.
之後才是耶穌的母親.

$^{201}$瑪利亞明白到祂首先是一個門徒.
然後才是母親的身份.
門徒的身份必須是善於優先過母親的身份.
門徒的身份必須優先過母親的身份.
作者提醒讀者或聽眾一件事.
瑪利亞是我們學習的對象.
如果我們不懂得做耶穌門徒的話.
我們就會像耶穌的其他兄弟一樣走散.
各位弟兄姊妹.
無論你是母親還是父親.
還是兒子還是女兒還是單身.
沒有東西是優先過你是門徒的身份.
學習做門徒.
我們才有我們隨之而來的家人.
看完女門徒 我們就講一下男的門徒.
在第27節.
「又對那門徒說:看你的母親.
從那刻起 那門徒就接她到自己家裡去了」.
當耶穌被捕的時候.
所有男性的門徒都雞飛狗走.
留下來的只有婦女和唯一一個男性.
就是那門徒.
耶穌看見追隨的人只有那門徒一個.
唯有把他母親瑪利亞交給那門徒照顧.
所以在最後那一句.
「從那刻起 那門徒就接她到自己家裡去了」.
在上一節我剛才說.
耶穌稱自己母親為女人.
似乎拉遠了與瑪利亞的關係.
但是耶穌終究是擔心自己的母親.
在自己過世之後.
或者在升天之後沒有人照顧.
在猶太人的習俗裡.
都會有一個習慣.
就是死囚在臨死的時候.
他可以在見證人的底下.
去交代屬於自己的女性親屬.
他的法定地位是什麼.
正常來說.
當耶穌應該把瑪利亞交給自己有血緣關係的兄弟.

$^{241}$但是當耶穌去環顧四周.
只有一個男性的門徒在那裡.
不過當耶穌把瑪利亞交託給那門徒的時候.
這位沒有血緣關係的門徒.
就成為了瑪利亞的兒子.
在這一刻.
這個新家庭就出現了.
究竟這個新家庭有什麼特別呢?.
是不是我們所理解的核心家庭.
即是父母,兒子,女兒這四個人就算了呢?.
事實上如果我們看希羅世界歷史的話.
「Familia」這個字不是指父母,兒子,女兒的核心家庭.
「Familia」這個字是指家主或家父.
在他的權威底下.
他控制著所有成員的身份和財產.
他的成員很廣闊.
可以包括他的妻兒,親屬.
甚至是奴隸或被釋放的奴隸.
或者自由人.
是包含有血緣關係和沒有血緣關係的成員的大家族.
家主會為自己旗下的成員提供庇護.
相反受保護的家族成員會服侍家主.
當耶穌離世的時候.
瑪利亞即將會失去一個自己的家主.
成為一個無依無靠的寡婦.
但耶穌為瑪利亞提供了一個新的家主.
正如舊約的律法或新約的書信都不斷重複.
耶穌示範一次甚麼是照顧寡婦.
照顧那些失去經濟之處的女性.
英文怎樣去呈現這個概念.
從那刻起,那門徒就接他到自己家裡去了.
在十字架底下成就了一個新的家庭出來.
這個家庭不是建基於血緣關係.
正正呼應約翰福音14:2-3.
在我父的家裡有許多住處.
我去完事為你們預備地方去.
我若去為你們預備了地方.
就必在來接你們到我那裡去.
我在那裡,叫你們也在那裡.
耶穌說,我父的家不是指另一個次元的天堂.

$^{281}$不是指天堂.
住處不是指天堂有多少間房讓你搬進去住.
而是指父子聖靈和信徒群體之間相互內住的狀態.
當信徒彼此相愛,遵從主的話語的時候.
神就在這個群體當中.
人就在這個家裡找到上帝.
耶穌所預備的家是因著他的死亡和復活.
建立了一個新的群體.
這個新的群體就是教會的群體.
是新的家庭.
信徒不需要再留在耶路撒冷的聖殿裡做敬拜.
因為耶穌基督的身體正正就是真正的殿.
所以耶穌說:我在那裡,叫你們也在那裡.
當耶穌復活升天的時候.
門徒再次相聚起來.
成為了一個新的群體,一個新的家.
約翰福音的讀者也不例外.
這群讀者非常可能是在公元九十年.
耶穌基督所生活的一群基督跟隨者,是一群門徒.
這一群門徒正正是面對二十年前,即是公元七十年.
聖殿被毀,沒有辦法在聖殿做敬拜的人.
另外這一群基督跟隨者,這一群門徒.
他們也要面對被猶太人逐出猶太會堂.
看似沒有了一個敬拜的地方.
正當門徒這一群基督跟隨者流離失所的時候.
他們很需要被領袖的鼓勵,繼續去做基督的門徒.
作者寫這本書正正就是要鼓勵這一群門徒.
要繼續去相信基督,在新家庭裡繼續彼此相愛.
彼此洗腳,彼此服侍地去生活.
門徒在這一刻有新的身份,有新的家人.
不再是血緣關係的家人.
到最後,我們作為基督徒,在我們的角度.
什麼叫做一家人?什麼叫做家庭?.
外面主流社會所說的家人,仍然停留在核心家庭.
父母,兒女的概念.
這個概念甚至被推崇至社會的核心價值.
不過這樣,我們基督徒是這樣看的嗎?.
經文那裡這樣說:婦人看你的兒子,看你的母親.
誰才是你的家人呢?.
我們常說窩福是我們屬靈的家.

$^{321}$但其實可能背後那句是說.
她的老公,老婆,兒女才是真正的家.
其實這樣是錯的.
窩福是基督所建立的新家庭.
窩福是由一群跟隨基督的門徒而被建立起來.
窩福不是十幾個核心家庭而組成的群體.
我們被建立不是因著我們的血緣關係.
我們被建立是因為我們認識上帝.
我們敬拜上帝而被建立起來的門徒群體.
在這裡,我們首先是門徒.
我們這個角色是先於我們其他所擁有.
我們背著的角色.
我們才會說窩福不是你屬靈的家.
窩福是你真正的家.
請你好好地照顧,提供庇護.
護佑著窩福裡面每一位弱勢的人.
給你年幼的人.
正如那門徒去接瑪利亞.
去自己的家裡去照顧她.
那門徒不是帶著瑪利亞回去自己的家.
然後叫她自生自滅.
不是叫她自己去查經.
那門徒是要照顧瑪利亞的起居,飲食,衣食,住行.
我們怎樣去實踐這個超越血緣關係的新家庭呢?.
上星期我和幾個主內的弟兄姊妹一起吃飯.
其中一個朋友M.
他分享自己信主做門徒的經歷.
他信主過程當中其中一個故事是這樣的.
當時在教會裡有一個工作要一起服事.
成立了一個工作小組.
在這個小組裡有不同的弟兄姊妹.
朋友M有機會去服事.
但那時候他還未信耶穌.
在這個小組裡有一個很友善的姨姨.
她的女兒和我們差不多大.
她很關照朋友M.
到朋友M生日那天.
姨姨去找朋友M請他吃飯.
慶祝他的生日.
這個舉動好像很普通.

$^{361}$對於朋友M來說.
這是他信主歷程的其中一個很重要的片段.
這位姨姨確實的輩分是做了我們媽媽.
但她沒有血緣關係.
縱然她沒有血緣關係.
但她仍然關顧著朋友M的生命.
無論是肉身還是屬靈的生命.
她仍然關照他.
當他是兒子般.
慷慨地去照顧朋友M.
姨姨看你的兒子.
朋友M看你的母親.
窩福的弟兄姊妹.
請你首先成為一個門徒.
跟隨耶穌基督的門徒.
我們學習成為一個門徒.
當我們學習成為一個門徒的時候.
我們才懂得怎樣去成為一個母親.
或者成為一個父親.
或者一個兒子.
或者一個女兒.
當我們去跟隨基督的時候.
我們就學習去接待在我們群體中.
有需要的弟兄姊妹.
給你年幼的弟兄姊妹.
有需要輔助的弟兄姊妹.
請你慷慨地去接待他們.
接待群體中貧乏的.
給你年幼的.
給你弱勢的.
正如那門徒.
他慷慨地去接瑪利亞到自己家裡去.
所以請你慷慨地去接待.
窩福家裡每一個人.
今日首先是主日.
我們首先去敬拜上帝.
因為我們首先是門徒.
不要搞錯秩序.
讓我們同心祈禱.
天父求你幫助我們.

$^{401}$學習去成為基督的門徒.
當我們成為門徒的時候.
我們學習愛神愛人.
當我們學習彼此相愛的時候.
眾人就認得出.
我們是你的門徒.
幫助我們去實踐.
求聖靈去加力給我們.
讓我們慷慨地去接待.
每一位有需要的弟兄姊妹.
求主去加力給我們.
我們這樣祈禱.
奉主耶穌基督得勝之名而求.
阿們.
\newpage



\section{撒母耳記上 1:1-28}
\label{sec:032131}
\textbf{偉人母親的質素}~~~~梁永善牧師
\newline
\newline
連結: \href{http://www.church.com.hk/acms/content.asp?site=cdc\&op=show\&type=product\&code=032131\&layout=sermon}{\texttt{032131}} ~~~~ 語音日期: 2026-05-10
\newline
\newline
\hyperref[sec:032130]{\small{< < < PREV SERMON < < <}}
~
\hyperref[sec:index_chronic]{\small{[返順時目]}}
~
\hyperref[sec:032132]{\small{> > > NEXT SERMON > > >}}
\newline
\newline
撒母耳記上 1:1-28
\newline
\begin{longtable}{cl}
\hline
\hline
章節 & 經文 (和合本修訂版)\\
\hline
1:1 & \begin{tabularx}{0.7\textwidth}{X} 以法蓮山區有一個拉瑪的瑣非人，名叫以利加拿，他是蘇弗的玄孫，託戶的曾孫，以利戶的孫子，耶羅罕的兒子，是以法蓮人。 \end{tabularx} \\ \\ \relax
1:2 & \begin{tabularx}{0.7\textwidth}{X} 他有兩個妻子：一個名叫哈拿，另一個名叫毗尼拿。毗尼拿有孩子，哈拿卻沒有孩子。 \end{tabularx} \\ \\ \relax
1:3 & \begin{tabularx}{0.7\textwidth}{X} 這人每年從本城上到示羅，敬拜萬軍之耶和華，向他獻祭。在那裡有以利的兩個兒子何弗尼和非尼哈當耶和華的祭司。 \end{tabularx} \\ \\ \relax
1:4 & \begin{tabularx}{0.7\textwidth}{X} 每逢獻祭的日子，以利加拿把祭肉分給他的妻子毗尼拿和毗尼拿所生的兒女。 \end{tabularx} \\ \\ \relax
1:5 & \begin{tabularx}{0.7\textwidth}{X} 他給哈拿的卻是雙分，因為他愛哈拿。耶和華卻不使哈拿生育。 \end{tabularx} \\ \\ \relax
1:6 & \begin{tabularx}{0.7\textwidth}{X} 她的對頭毗尼拿因耶和華不使哈拿生育，就常常惹她發怒，要使她生氣。 \end{tabularx} \\ \\ \relax
1:7 & \begin{tabularx}{0.7\textwidth}{X} 年年都是如此。每當她上到耶和華殿的時候，毗尼拿就這樣惹她發怒，以致她哭泣不吃飯。 \end{tabularx} \\ \\ \relax
1:8 & \begin{tabularx}{0.7\textwidth}{X} 她丈夫以利加拿對她說：「哈拿，你為何哭泣？為何不吃飯？為何傷心難過呢？有我不比有十個兒子更好嗎？」 \end{tabularx} \\ \\ \relax
1:9 & \begin{tabularx}{0.7\textwidth}{X} 他們在示羅吃喝完了，哈拿就站起來。祭司以利坐在耶和華殿門框旁邊的位子上。 \end{tabularx} \\ \\ \relax
1:10 & \begin{tabularx}{0.7\textwidth}{X} 哈拿心裡愁苦，就痛痛哭泣，向耶和華祈禱。 \end{tabularx} \\ \\ \relax
1:11 & \begin{tabularx}{0.7\textwidth}{X} 她許願說：「萬軍之耶和華啊，你若垂顧你使女的苦情，眷念不忘你的使女，賜你的使女一個子嗣，我必使他終生歸給耶和華，不用剃刀剃他的頭。」 \end{tabularx} \\ \\ \relax
1:12 & \begin{tabularx}{0.7\textwidth}{X} 哈拿在耶和華面前不住地祈禱，以利注意她的嘴。 \end{tabularx} \\ \\ \relax
1:13 & \begin{tabularx}{0.7\textwidth}{X} 哈拿心中默禱，只動嘴唇，聽不到她的聲音，因此以利以為她喝醉了。 \end{tabularx} \\ \\ \relax
1:14 & \begin{tabularx}{0.7\textwidth}{X} 以利對她說：「你要醉到幾時呢？不要再喝酒了！」 \end{tabularx} \\ \\ \relax
1:15 & \begin{tabularx}{0.7\textwidth}{X} 哈拿回答說：「我主啊，不是這樣。我是心裡愁苦的婦人，清酒烈酒都沒有喝，只在耶和華面前傾心吐意。 \end{tabularx} \\ \\ \relax
1:16 & \begin{tabularx}{0.7\textwidth}{X} 不要將你的使女看作不正經的女子。我因極其難過和生氣，所以一直禱告到如今。」 \end{tabularx} \\ \\ \relax
1:17 & \begin{tabularx}{0.7\textwidth}{X} 以利回答說：「平平安安地回去吧。願以色列的神允准你向他所求的！」 \end{tabularx} \\ \\ \relax
1:18 & \begin{tabularx}{0.7\textwidth}{X} 哈拿說：「願你的婢女在你眼前蒙恩。」於是婦人上路，去吃飯，臉上不再帶愁容了。 \end{tabularx} \\ \\ \relax
1:19 & \begin{tabularx}{0.7\textwidth}{X} 他們清早起來，在耶和華面前敬拜，就回去，往拉瑪自己的家裡。以利加拿和妻子哈拿同房，耶和華顧念哈拿。 \end{tabularx} \\ \\ \relax
1:20 & \begin{tabularx}{0.7\textwidth}{X} 時候到了，哈拿懷孕生了一個兒子，給他起名叫撒母耳，說：「這是我從耶和華那裡求來的。」 \end{tabularx} \\ \\ \relax
1:21 & \begin{tabularx}{0.7\textwidth}{X} 以利加拿和他全家都上去，要向耶和華獻年祭和還願祭。 \end{tabularx} \\ \\ \relax
1:22 & \begin{tabularx}{0.7\textwidth}{X} 哈拿卻沒有上去，因為她對丈夫說：「等孩子斷了奶，我就帶他上去朝見耶和華，讓他永遠住在那裡。」 \end{tabularx} \\ \\ \relax
1:23 & \begin{tabularx}{0.7\textwidth}{X} 她丈夫以利加拿對她說：「就照你看為好的去做吧！可以留到兒子斷了奶，願耶和華應驗他的話。」於是婦人留在家裡乳養兒子，直到他斷了奶。 \end{tabularx} \\ \\ \relax
1:24 & \begin{tabularx}{0.7\textwidth}{X} 斷奶之後，她就帶著孩子，連同一頭三歲的公牛，一伊法細麵，一皮袋酒，上示羅耶和華的殿去。那時，孩子還小。 \end{tabularx} \\ \\ \relax
1:25 & \begin{tabularx}{0.7\textwidth}{X} 他們宰了公牛，就領孩子到以利面前。 \end{tabularx} \\ \\ \relax
1:26 & \begin{tabularx}{0.7\textwidth}{X} 婦人說：「我主啊，請容許我說，我向你，我的主起誓，從前在你這裡站著祈求耶和華的那婦人就是我。 \end{tabularx} \\ \\ \relax
1:27 & \begin{tabularx}{0.7\textwidth}{X} 我祈求為要得這孩子，耶和華已將我向他所求的賜給我了。 \end{tabularx} \\ \\ \relax
1:28 & \begin{tabularx}{0.7\textwidth}{X} 所以，我將這孩子獻給耶和華，使他終生歸給耶和華。」他就在那裡敬拜耶和華。 \end{tabularx} \\ \\
[1ex]
\hline
\hline
\end{longtable}
$^{1}$各位弟兄姊妹,你們好.
今天很高興來到大家當中.
跟大家分享一句話語.
在英國的姊妹們是特別有福的.
為什麼呢?.
當你在香港過來的時候.
有兩個母親節.
在英國是三月.
一定要和你慶祝.
如果你家人不和你慶祝.
你就要入鄉隨俗.
你竟然忘記和我慶祝母親節.
罪大惡極.
慶祝就不等於今天.
五月第二個星期就是香港.
或其他地方的母親節.
不可以因為慶祝了.
英國的母親節.
今天不慶祝.
你的丈夫或同事說.
三月和你慶祝了.
這麼快就忘本了.
我們的根在香港.
所以怎麼可以沒有母親節的慶祝.
所以你們真的很好.
有兩個母親節.
所以我一直說這個世界男女不得平等.
父親節是什麼時候?.
記得也算好了.
有些人是不記得的.
六月第三個星期.
希望你們記得可憐你的丈夫.
或者可憐你的父親.
好不好?.
今天是母親節.
所以我特別講一篇訊息.
是為母親而設的.
如果在我們當中.
你的妻子或你媽媽在的時候.
可不可以你現在輕聲跟她說.

$^{41}$媽媽或太太.
沒有你,我真的沒有今天.
要多謝她.
是真的.
是真心的.
我今天也沒有機會.
跟我太太.
慶祝母親節.
但我明天下午會獻屬序制.
明天就會.
跟她慶祝母親節.
但我想在講之前.
先講幾句我自己的感受.
我自己從佳節陪侍親.
我今年.
你記不記得復活節.
和清明節是同一天.
我自己也有很大的感慨.
因為熟悉我,我也曾經講過.
我媽媽四十二歲離開世界.
這段時間.
她在醫院裡,是一九七八年.
她六月六日離開世界.
我媽媽很年輕.
四十幾歲,我還在讀神學.
第三年.
我到今天仍然有一個很大的抱恨.
特別是母親節.
和清明節.
因為我完全沒有機會.
去回報我媽媽對我的恩情.
稍後也會講我媽媽對我有什麼恩情.
我媽媽九十幾歲離開.
離開後.
我也感恩,為什麼呢?.
她看到我怎樣去事奉.
她由反對我去做傳道人.
到後來她接受.
甚至鼓勵我要做好一點.
她看到是怎樣.

$^{81}$但很可惜.
我媽媽沒有機會看到.
那時我神學還未畢業.
所以我經常鼓勵特別是年輕人.
一定要去.
珍惜你的父母.
要把握機會.
有花堪折.
直須折.
莫待無花空折枝.
即是今天.
有時我們很想去.
回應父母對我們的關心.
但樹欲靜時.
風不息.
時辰不在.
那種痛苦你們未必經歷到.
我經歷到.
所以我鼓勵大家.
要珍惜你的配偶.
珍惜你的家人.
這很重要.
母親節會講一個偉人的母親.
有什麼質素可以承擔呢?.
很多年輕人.
很想做名門之後.
爸爸有錢.
有地位,有名譽.
名門之後.
特別是咬著金鑰匙出世的.
但往往.
有時那些有錢人.
中國人有句說話.
「富不過三代」.
那條褲子是穿不過三代的.
不是.
有錢人是不會落到第三代.
因為他們太享受.
父母的庇蔭.
他們是不懂珍惜.

$^{121}$四十五歲人都可以不出去工作.
我現在怎樣工作?.
所以這是一個悲哀.
甚至可以一代不如一代.
如果你有聽我講永攝論壇.
我原本明天.
預計是講永攝論壇.
「一個也不能少」.
是講我們要著重生命.
但現在改了.
我今晚回去會講.
今晚會錄音.
星期二會上載.
「鼎爺悲哀對我們的反省」.
我會講.
鼎爺的問題.
是給我們很多提醒.
很簡單舉個例.
有些人喜歡做名門之後.
但可能未必是好的.
但更加有些人.
是很想做名門的父母.
我當年都拍過電影.
你知不知.
我拍過幾套電影.
你上去搜索一下.
極度智能.
賭神之神.
我沒有騙你.
你上去會看到.
我記得我一年去紐西蘭.
當年COVID之前.
每年三月都去紐西蘭廣道.
有個家庭.
我都不會忘記.
「牧師你還有沒有戲拍?」.
我說應該遲些都有一套.
「可不可以找個角色.
簡簡單單都可以」.
「叫我兒子過來.

$^{161}$做童星可以嗎?」.
我說「不當我話事」.
「不收錢的.
我會叫他過來」.
很多人都很想.
讓子女去表演.
所以你見到很多人.
可能叫子女去參加比賽.
一登龍門.
已經升價十倍.
如果子女讓人來賞識.
我是星媽.
星爸都可能很開心.
但現在你見到.
完全不是這樣.
中國人有句說話.
「陽明星」.
知道下面那句是甚麼嗎?.
但那些文化.
已經沒有了.
雖然我說話粗俗一點.
但其實我都有中文的.
「陽明星顯父母」.
即是說你厲害的時候.
你父母都得到甚麼呢?.
即是聲譽.
你記不記得.
你小學讀書的時候.
媽媽照顧你放學的時候.
你考第一.
你當然很開心.
家長問「你兒子考成怎樣?」.
「考第一」.
你都會覺得很開心.
如果他是考到最後.
你趕不趕快離開.
所謂「陽明星顯父母」.
當然那個人本身.
是要努力.
但問題是我仍然覺得.

$^{201}$很重要.
父母的生命是在影響自己的兒女.
你想兒女好的時候.
你又要怎樣呢?.
所以今天我們要留意.
父母的愛護.
教導.
和生命.
是影響下一代很重要的.
我已經很欣賞我爸爸.
我爸爸當時都未信耶穌.
但我爸爸是不吸煙.
不喝酒,不賭錢.
所以我到今天.
我都不會吸煙.
只是吃煙仔餅,沒有吸煙.
是不會喝酒.
更加不會賭錢.
這是爸爸在影響我.
當然我爸爸未信主.
但我爸爸是白手興家.
很努力,很勤力工作.
年初一都開工.
他不是做飲食業,他做帳戶.
這就影響我自己.
所以請你記住.
原來你要怎樣影響你的兒女.
是你的生命.
你的價值觀,這是很重要.
正如.
在中國裡面.
記不記得有一個.
萬母三千.
這個媽媽很好.
因為原來那一區怎樣影響她的兒女.
就是萬母三千.
我相信今天在我們當中.
都有很多萬母在.
萬父都在.
有時你搬一間屋.

$^{241}$就是那個名校區.
地價或屋價貴一點.
你都去.
我兒子終於在一年多前.
搬到另一區.
因為那裡有一間基督教女子中學.
他就很想.
兩個女兒.
就入這間學校,所以就搬去那裡.
那些樓價貴一點.
原來我才知道.
那間屋苑.
被稱為外府村.
因為十個有八個.
都是為了入那間基督教女子中學.
所以那個叫外府村.
就是這樣.
所以現在這個年代裡面都有很多萬母.
聖經裡面有一個很偉大的人物.
這個人是聖經有.
兩卷書用他的名字.
來命名的.
他亦都是以色列.
支派裡面,即是人民裡面.
第一個和第二個王.
的高立者.
他亦都是最後一個.
的事事,誰呢?.
撒母耳.
撒母耳是.
一代偉人.
但是撒母耳的成長.
除了他自己本身.
的努力之外.
他的母親是很重要.
撒母耳那個年代裡面.
是事司年代.
我現在星期三.
的茶經是在查事司記.
是一個非常混亂的.

$^{281}$世代.
如果你看事司記.
最後那幾章的時候.
你會看到甚至當時一些.
的祭司是為了賺錢.
而已.
母親正正的教導人去敬拜神.
甚至來拜偶像.
哪個地方出錢多一些.
我就去那個地方.
是很可悲的.
而當時的宗教領袖.
是誰呢?是以利.
但是以利兩個兒子.
他應該要繼承父業.
但是已經是宗教領袖.
但是以利都沒有好好地.
教導自己的孩子.
他來的是去搶.
制慾,人家是獻祭.
最好那些獻給神.
他竟然不理人的反對.
拿了那些慾私下自己要了.
甚至和.
管理會幕那些婦女.
發生淫亂的行為.
所以當時是一個.
非常非常.
道德敗壞.
各人沒有尊敬神的.
世代.
但是這個年代還要做一個.
好的母親,還要教自己的.
兒女敬拜神,絕對不容易.
我們看到.
哈拿.
有什麼值得我們欣賞的地方.
第一點就是.
他不畏艱辛,忠心去敬拜.
我們一起來讀.

$^{321}$《撒慕爾記》上第一章.
第一到第三節,我們一起來讀.
他們就描繪到.
有一個人有兩個妻子.
正室是什麼呢?是哈拿.
他們就是.
每一年都會來.
到上道神的地方.
來敬拜神.
大家記不記得我剛才說.
史詩時代是非常.
宗教混亂,道德很敗壞.
連宗教領袖的兒子.
都不願意去.
去敬拜神.
他們就是.
連宗教領袖的兒子.
都亂來.
所以其實可以有很多藉口.
還有我們不要想著.
他們去敬拜是很容易.
你看這張地圖.
原來他們要敬拜的時候.
至少要.
走一段很長的路.
他們從伯利行一路走的時候.
要走很遠的路.
這張地圖就顯示出來了.
不要忘記那個年代.
沒有私家車.
沒有柏油爐.
而且中東地方很熱.
我昨天來的時候.
都很漂亮.
因為我早上出門.
穿上衣服很熱.
下午簡直是焗桑拿.
所以我來到這裡.
就想找個T恤來穿.
我要換一件短袖T恤.

$^{361}$中東地方更加熱.
他們是不畏艱辛.
去敬拜神.
縱使當時的宗教.
很糟糕.
其實他們可以有很多藉口不去.
現在的宗教領袖.
都不知所謂.
去做什麼呢?.
路途很遙遠.
還有神是無所不在.
我們現在在家裡.
都可以敬拜神.
為什麼要去這麼遠的地方.
去敬拜呢?.
剛才我說.
未來的宗教領袖.
惡行超張.
人人共憤.
為什麼要去敬拜神呢?.
他有很多藉口.
而且這些藉口可能是正確的.
但他們沒有因為這個緣故.
而放棄敬拜神.
今天我們知道.
神所喜悅的是什麼呢?.
就是願意人.
去敬拜祂.
正如你的子女.
讚賞你.
你都會開心.
但記住.
崇拜不是只聽一篇道.
其實崇拜最重要的是什麼呢?.
是敬拜的環節.
你們唱詩的時候.
是歌頌神的.
神是很開心的.
正如你聽到你的子女讚賞你.
你不是很開心嗎?.

$^{401}$所以我們不可以輕忽.
這個敬拜的時間.
但很可悲地.
今天有不少的基督徒.
用不同的藉口.
不出席崇拜.
可能他們因為工作.
可能工作真的有些人.
星期天要上班.
我接受到的.
但你會不會用另一種方法.
去敬拜神呢?.
有時候.
我們就會想.
太大風了.
特別是冬天.
這裡很冷.
不好了.
又或者小朋友.
今天是誰說的呢?.
梁應善.
不好了.
不要回去了.
一定有教訓的.
你知道什麼叫教訓嗎?.
我們有很多的原因.
甚至可能.
我跟誰不和.
我不想見到他.
這間教會又比較淪盡.
一眼就見到了.
無遮無掩.
無謂回去了.
不可能我們的藉口.
其實可能都是事實.
但是.
不要忘記敬拜是什麼?.
敬拜是朝向我們的神.
我們將我們的焦點.
放在神那裡.

$^{441}$而不是因為人.
所以今天.
我們怎麼可以輕忽敬拜.
以致我們不出席崇拜.
這是很重要的.
所以你留意到.
哈拿.
他不辭勞苦.
一來一回.
去敬拜神.
這是值得神所欣賞.
你想想.
如果父母愛主.
很熱切地親近神.
對神有尊重.
的時候.
神怎麼會不祝福他呢?.
而且不是只祝福他.
聖經裡面很清楚說.
是會祝福到他的後裔.
大家應該很熟悉這段經文.
有三篇.
第十七到第十八節.
未結婚的不需要讀.
結婚的.
沒有身分證的都要讀.
因為可能你像哈拿.
將來會有的.
結婚的讀這段經文.
聖經裡面很清楚說.
如果我們敬畏神的時候.
神必定會.
不但祝福你.
祝福你的兒女你的孫兒.
我自己不斷地說.
你看你的子女多少年.
你看你的孫兒多少年.
這一天你都會離開.
但如果你.
能夠將他們.

$^{481}$帶到神面前認識神的時候.
你都可以安心離開.
因為我們懂得怎樣仰望神.
和你愛主的時候.
神必定會祝福他們.
所以你要做.
偉人母親.
先檢視你和神的關係.
第二個不是什麼呢?.
他就不理會挑剔.
身冤在主.
為什麼呢?.
原來哈拿是被別人.
挑剔著的.
那時候他做什麼事呢?.
就是因為.
他沒有小孩.
但耶穌就有一個.
妾侍.
但他又生了小孩.
那你知道.
三分顏色上大雄.
眼睛都沒有直生.
屁都沒有睡.
他天天去挑戰.
哈拿他很辛苦.
原來在當時.
在律法下.
是容許多於一個妻子.
但不等於.
這個真是神的心意.
律法上批准.
不等於正確.
正如英國.
LGBT沒有問題.
同性戀沒有問題.
但這是不是神的心意呢?.
在聖經裡面.
給我們看到一個很重要的教導.
我們一起來讀.

$^{521}$申命記第21章15.
到第17節.
我們一起來讀.
所愛的.
所護的都給他生了兒子.
但長子是.
所護之妻生的.
到了把產業分級.
兒子成熟的時候.
不可將所愛的妻生的兒子.
納為長子.
在所護之妻生的兒子.
以上.
其要認所護之妻生的兒子.
為長子.
將產業多給一分級.
因這兒子.
強壯的時候生的.
長子的名分.
本當歸他.
我們很清楚知道.
瑪拉基書講到.
原來是要一夫一妻.
一生一世.
我們一起看瑪拉基書第二章第15節.
雖然神有靈的餘力能造多人.
他不是單造一人嗎?.
為何只造一人呢?.
乃是他元人的虔誠的求裕.
所以當謹守你們的心.
誰也不可以詭詐.
待有靈所取的妻.
很清楚講到.
神造阿當的時候.
只是造一個妻子給夏娃.
神有能力造多一個.
但一定不會造給她.
是要一夫一妻.
然後期望有虔誠的後代.
為甚麼呢?.

$^{561}$如果兩夫婦關係好.
會很不同.
可以共同一致.
教導自己的兒女.
但如果關係不好就糟糕了.
譬如兩夫婦.
終於鬧離婚.
一個疼愛大人.
一個疼愛小人.
互相批評對方.
不單令兄弟不和.
更令他們有很多問題.
這是很現實的.
所以我們要很小心.
但難道.
以利加納不知道.
多妻帶來的痛苦嗎?.
他們的先祖.
他們應該很熟悉.
亞伯拉罕被稱為信心之父.
但亞伯拉罕有多少個妻子?.
多少個?.
三個.
不止下甲.
後來還造多一個.
真是離譜.
但他娶了下甲.
當然是撒拉叫他.
做男人當然沒問題.
是很樂意遵從.
但帶來很多衝突.
是帶給他很多煩惱.
亞伯拉罕曾經為這事煩惱.
好了.
真是.
一代新人勝舊人.
亞伯拉罕的孫子.
是誰呢?.
大家要想一想.
亞伯拉罕生了.

$^{601}$以撒.
以撒是生了一個孫子.
我是說下甲.
下甲有多少個妻子?.
多少個?.
四個.
但是很苦.
很煩惱.
以利加難道.
不知道他們的先祖.
多妻帶來痛苦嗎?.
但仍然.
他們都多妻.
所以很不幸地.
歷史是不斷不斷重複.
你想想當年阮國新.
何生他四個妻子.
生日很痛苦.
知道嗎?.
去完這一家吃飯.
為什麼不來我家?.
要去第二家.
去第三家.
要溝口.
不然死定了.
所以絕對不是一件好事.
但很可惜.
很多人.
不懂得吸收歷史的教訓.
不知道什麼叫.
前車可鑑.
歷史不斷重複.
文革是一件很大悲哀的事.
當時的人互相.
有告發甚至批鬥自己的父母.
但現在.
文革再一次顯現在.
我們所熟悉的地方.
很多人沒有吸取教訓.
哈娜雖然被丈夫所愛.

$^{641}$但她經常.
被妾侍.
激她.
挑戰她.
所以她很難受.
我看.
《撒慕爾記》第一章第七節.
我們一起讀.
比利娜仍是激動他.
以致他哭泣不喜歡.
雖多一份.
但我有兒子.
激死了.
但你留意到.
哈娜做什麼.
中國人一句說話是這樣說.
人爭一口氣.
接著.
佛爭一爐香.
原來佛也會爭.
那口氣是最重要.
我來也是激氣.
如果我是哈娜.
我就一定反神相向.
那怎麼辦呢?.
你又想怎樣?.
慘得我丈夫愛我.
你看,雙份的.
你那裡有多少?.
我一定反過來.
我沒理由.
不斷被她激.
但你看到.
哈娜有一份很美麗的德性.
就是學到.
什麼叫忍耐.
她知道潑婦麻街.
只會令事情更加惡化.
繼續和.
妾侍爭吵.

$^{681}$只會令愛她的丈夫.
很難做.
所以,哈娜.
她沒有離帶去.
反神相向.
也沒有向丈夫投訴.
縱使.
苦大仇深.
但她仍然不和敵人.
逼迫妾侍離帶爭拗.
她只是將問題.
帶到神的面前.
向神離帶去申冤.
向神離帶去傾訴.
我們一起來讀第一章第十節.
哈娜心裡受苦.
就痛痛哭泣.
祈禱耶和華.
她真的受很多苦.
所以.
她.
很受苦.
叫大家有一點.
視力都看得到.
這時候.
她就哭泣.
向神離帶去傾訴.
這是.
很重要的一件事.
今天.
你和我都會遇到一些不公平的事.
或者.
這些不公平的事.
不是發生在你身上.
而是發生在社會上.
你都會激氣.
最近,大家知不知道.
有一個個案.
一個消防隊伍.
擴大了.

$^{721}$自己的財富.
然後借錢.
就被判了八個月監禁.
他肯認罪.
但判了八個月監禁.
都是紀律部隊.
但有一個警察督察.
同樣犯了這個罪.
借的錢是差不多.
可以判幾個星期監禁.
但這個人是不肯認罪的.
大家都是紀律部隊.
怎麼會有這麼不公平的判決.
有很多不公平的事.
甚至可能你被人惡意來挑剔你.
甚或.
這些挑剔可能是來自.
你的家人.
你的好友.
但我們會不會.
缺乏了.
Hannah這種美德.
而引伸到一些慘劇呢.
所以.
我們應該要.
學習的是什麼呢.
就是要好像Hannah這樣.
去仰望神.
有冤屈.
向神來傾訴.
因為聖經裡面說.
「身冤在主,主必報應」.
向神來講.
甚至奏作這些人.
神可以接納.
透靠神.
這個看似是.
你不出聲,你不是弱者所為.
所以一定要反部隊.
但原來.

$^{761}$你忍耐才是強者.
因為苦難.
有時會令人質疑神.
神你是否愛我.
你是否眷顧我.
因為有些人遇到苦難的時候.
可能會離棄神,放棄這個信仰.
但同樣有時.
苦難是驅使人更加去親近神.
他知道自己無能為力.
去仰望神.
他知道神會垂聽禱告.
正如.
夏格.
夏格是誰呢.
他是亞伯拉罕的妾侍.
他是一個婢女.
是代母.
當他生了二十萬.
又是壞人.
又是招職.
去藐視撒拉.
撒拉很憤怒,要趕他走.
叫亞伯拉罕.
一個外籍傭工.
但帶著小孩離開.
很痛苦.
完全不知如何是好.
這樣的時候.
神會怎樣呢.
神會來祝福.
神也祝福哈拿.
當哈拿向神傾訴的時候.
神會怎樣幫助他.
但同樣.
夏格.
被藐視的時候.
但神又怎樣祝福.
我們一起讀《創世記》十六章第十一節.
這裡很清楚.

$^{801}$耳實馬利.
原本的意思是.
神聽見我的苦情.
所以我們留意到夏格是一個很卑微的人.
但神也眷顧他.
所以當我們遇到困難.
挑戰.
苦情.
我們向神傾訴.
正如詩篇69篇第20節和第33節.
我們一起讀出.
我指望有人提出.
卻沒有一個.
我指望有人安慰.
卻找不著一個.
因為耶和華聽了窮乏人.
不藐視被囚的人.
有時候你發覺好像沒有人了解你.
沒有人聽你的傾訴.
但不要忘記.
神是.
眷顧你.
誰聽你的苦情.
我們願不願意將我們的問題.
帶到神的裡面.
當他禱告的時候.
竟然.
被當時的宗教領袖.
誤會了他.
他含冤.
但他辯白.
他受了宗教領袖的冤枉.
但他是.
溫柔地回答.
我們一起來看.
第12節到第14節.
第一章我們一起讀.
原來哈拿心中蜜桃.
只動嘴唇不出聲音.
因此伊利.

$^{841}$以為他喝醉了.
伊利對他說.
你要醉到幾時.
你不應該喝酒.
哈拿一路去祈禱.
其實我都鼓勵大家.
有時候祈禱最好.
都是開聲.
但可以很小聲.
因為你沒有留意到你一祈禱.
你的思想不知道去了哪裡.
可能回來又再來.
但如果你小聲開聲祈禱.
是可以專注一點.
哈拿就在那裡.
可能他自己聽到.
但原來伊利.
就留意這個女人.
我不幹甚麼呢.
就以為他喝醉了.
我也有這個經歷.
這個經歷是甚麼呢.
我冥恩堂在仁興街.
火車博物館斜對面.
如果去過大埔就知道.
那裡有一間火車博物館.
那時每星期有一天.
我就禁食祈禱.
我就去了火車博物館.
我會找.
你知道有些舊的車卡.
我會找頭等車商坐.
因為我小時候.
從來沒有坐過頭等車卡.
所以拿回我童年的失落.
正如現在我來這裡.
很喜歡吃脆皮雪條.
因為我小時候沒有脆皮雪條.
我就會祈禱.
但我又.

$^{881}$我冥恩堂是低頭祈禱.
很容易會打瞌睡.
就像哈娜.
我會祈禱.
你知道他們有些保安走過.
他懷疑有個神經老.
坐在那裡看著外面.
他嘴是這樣.
眼睛是這樣.
他真的以為我是神經.
一直在留意我.
但我真的這樣.
但被人誤會.
被人誤會以為他是喝醉酒.
最慘誤會他.
是誰?.
全國宗教領袖.
如果你是這樣.
你會怎樣?.
被人誤會是一件很悲慘的事.
這時候.
特別誤會你的人.
可能是一些你應該受尊重的人.
喂!.
哈娜已經可以這樣禱告.
但以你這樣的時候.
他馬上可以說:我走!.
你都不了解我.
不體諒我.
被一些屬靈的人誤會.
是一件不幸的事.
甚至影響你.
我和你分享過.
我這幾天都很大感受.
因為有個牧師在加拿大來了.
是我十三歲回屯契的時候.
我的團友.
他也做了牧師.
回來了幾天.
我昨天也和他在一起.

$^{921}$他明天就會回加拿大.
我明天也會送他上樓.
這時候.
我們有很多童年回憶.
我們一直聊童年回憶.
聽到我太太和我女兒眼睛都亮了.
你小時候是這樣的嗎?.
這時候.
有一件事是什麼呢?.
當年我和大家分享過.
我喜歡唱歌.
我的聲音應該也不差.
我後來參加.
教會的詩歌班.
不要說詩班.
詩班是死了出來的.
詩歌班.
我練習了很久.
終於有機會去唱歌.
不是獨唱,是一群人一起唱.
唱完之後回到台上.
我們的詩琴就是我們的詩歌班長.
他還沒有轉過身.
叫我們坐下來的時候.
後面就有人這樣說.
「剛才林永善唱歌走音了」.
「被他害慘了」.
「他以為很小聲」.
「而且很靈敏」.
「而說我是誰呢?」.
「我現在都記得」.
「他的樣子不是很記得」.
「他的名字我一定記得」.
「他是團契的靈修」.
「姓容的」.
「不再說全名」.
「不知道他還在不在生」.
「我們在團契做總務」.
「就是靈性」.
「不要做總務搬來搬去」.

$^{961}$「靈修一定是靈明好的」.
「你說我唱歌走音」.
「害你?好!」.
「我就不連累你」.
「從此我就退出江湖」.
「不再參加詩歌班」.
「但是我有小小的心靈」.
「受著傷害」.
「到今天都受著這個傷害」.
「所以到今天我唱歌走音」.
「都是因為他」.
「你明不明白?」.
「有時候你被人冤枉」.
「你是很慘的」.
「甚至因為這個緣故」.
「可能你會放棄」.
「但是哈拿絕對不是這樣」.
「他只是尊重義理」.
「沒有去怪責他」.
「沒有怪責他」.
「你沒有體諒我」.
「你沒有動策力」.
「他沒有去責備他」.
「更加你沒有好好輔導我」.
「我現在這麼難過」.
「你只會罵我」.
「你沒有牧養我」.
「你可以有很多責備」.
「來面對義理」.
「但是他沒有這樣做」.
「他只是很柔和地回答」.
「他為自己辯白」.
「好,我們一起來看這段經文」.
「好,我們一起來讀出」.
「十五頁以上第一章的十五十六兩節」.
「哈拿回答說」.
「主啊,不是這樣」.
「我是心裡仇父的婦人」.
「清酒濃酒都沒有喝」.
「但在耶和華面前傾心吐怡」.

$^{1001}$「不要將彼女看作不正經的女子」.
「我因被人激動」.
「受苦太多」.
「所以祈求到如今」.
「你聽到嗎?哈拿」.
「她不是去罵義理」.
「不是怪責她」.
「沒有體諒她」.
「她很溫柔地說」.
「我現在受了很多苦」.
「所以我很痛苦」.
「我現在在神面前」.
「很激動地向神禱告」.
「結果得以澄清」.
「我想跟大家說」.
「女性最厲害的武器是什麼呢?」.
「不是惡」.
「是你的眼淚」.
「一個正常的男性」.
「如果看到你所愛的人流眼淚」.
「就整個人變軟了」.
「但你罵,我們一定會反駁你」.
「我們這些轉數快的」.
「反駁到你沒有聲音」.
「所以記住」.
「最厲害的武器是眼淚」.
「很希望每一個女士都有這份美德」.
「還有一點」.
「她不再保留恆切的禱告」.
「剛才我們看到」.
「第十到第十二節的時候」.
「哈拿心裡受苦」.
「在神面前哭泣」.
「祈禱」.
「她有一個應許」.
「我們一起來讀讀」.
「第一章第十到第十二節」.
「我們一起來讀」.
「上去」.
「我們一起來讀」.

$^{1041}$「哈拿心裡受苦」.
「就痛痛哭泣」.
「祈禱耶和華許願說」.
「萬君子耶和華」.
「你若誰顧疲累的苦情」.
「眷念不忘疲累」.
「賜我一個兒子」.
「我必使他終身歸於耶和華」.
「不用提頭倒提他的頭」.
「哈拿在耶和華面前不住的祈禱」.
「以利定睛看他的嘴」.
「哈拿是很難過」.
「去曲執」.
「但祂的禱告是很簡潔」.
「祂只求一件事很明確」.
「你給我一個兒子」.
「如果你給我一個兒子」.
「我會奉獻給你」.
「做拿撒爾人」.
「拿撒爾人是不會提頭的」.
「所以哈拿的禱告很值得我們學習」.
「祂是用感情」.
「熱切的心向神傾訴」.
「祂的禱告很簡潔」.
「很厄妖」.
「情詞迫切」.
「祂沒有很美麗的修辭」.
「修飾」.
「悅耳動聽」.
「但是目標明確」.
「賜我一個兒子」.
「這個禱告不單單是祈求」.
「更加傾訴內心的苦情」.
「就像和朋友傾訴一樣」.
「傾心吐意」.
「就是最好的祈禱」.
「所以求神了幫助我們」.
「學習怎樣將我們的需要」.
「很坦白地跟神了說」.
「但有時候我們隨便而已」.

$^{1081}$「就像交差一樣說了幾句」.
「神啊你是無所不知的神」.
「按你的雖明白之道去賜福」.
「Amen」.
「那就搞定了」.
「你的神學觀點一定是對的」.
「但這樣不是祈禱」.
「你有沒有跟神交往呢?」.
「而且哈拿的禱告」.
「不是一次的」.
「祂的禱告是什麼呢?」.
「是恆久的」.
「所以那段經文第十二節」.
「怎麼說呢?」.
「哈拿在耶和華面前不住地禱告」.
「所以今天」.
「你在祈禱怎樣呢?」.
「有時候禱告幾次沒有那一刻」.
「可能你會放棄」.
「有一個很深的感受」.
「有一個媽媽的兒子吸毒」.
「勸也勸不了」.
「罵也罵過」.
「後來這個媽媽用了八年的時間」.
「每天下午不吃飯」.
「那些叫濕漉嗎?」.
「我不知道」.
「減肥的那些叫什麼?」.
「168那些」.
「她早就學會了」.
「每天下午都不吃飯」.
「就禁食為自己的孩子來禱告」.
「終於這個浪子回頭」.
「還大學畢業了」.
「所以突然之間可能你的兒女」.
「你的配偶有時候令你失望」.
「你勸也勸過」.
「罵也罵過」.
「但是好像都沒有什麼改變」.
「你懂不懂」.

$^{1121}$「將它帶到主的面前」.
「還不是一次兩次」.
「而是繼續堅持這個禱告」.
「如果你相信神是願意人歸正的」.
「這個禱告你繼續堅持的時候」.
「必定會有一天答應」.
另一樣值得我們學習的就是「不」.
它就是沒有帶去看環境.
信心的倚靠神.
「以利」其實都是值得我們欣賞的.
「以利」是一個很簡單的招數.
以利是一個宗教領袖.
他誤會了人.
但是他肯認錯.
你明白嗎.
你現在見到黃福苑的聽證會.
錯到不得可料.
這個是制度的問題.
都是不太理想.
一個人可以講15次.
一天在一個聽證會裡.
這個不理想.
到現在都沒有一個人肯來.
但是他認自己是錯.
留意到「以利」.
他承認自己原來看錯了「哈拿」.
然後他怎樣來回應「哈拿」呢.
我們一起去看.
17,18兩節好不好.
所以你留意到「以利」他很好.
他來但是他.
馬上承認自己看錯了.
然後他就來祝福「哈拿」.
他說神必定會來.
但是他賜福給你.
而「哈拿」亦都來.
信服這個屬靈的長者.
神的僕人.
他對神有信心.
他相信.

$^{1161}$神是會賜福給他的.
但不是那一刻他馬上.
「哎呀,我不知道不是懷孕了嗎?」.
你記住.
「以利」講完之後.
他沒有即時懷孕.
更加他沒有回家的時候.
那個妾侍馬上啞了.
想激他都激不到.
可能他仍然在面對這個挑戰.
但縱使.
外在的環境.
沒有即時的改變.
但他內心.
有這個改變.
他將他的目光.
帶去轉向神的時候.
什麼叫信心呢.
信心就是接受那些未曾發生.
但你會相信.
會發生的事.
「哈拿」改變了.
之前是什麼呢,記不記得.
是哭泣不吃飯.
不開心你是吃不下飯的.
但現在他接受了.
「以利」那個祝福.
他都確信.
這個祝福是會成就的.
於是有什麼反應呢.
是很大的不同.
第十八節我們一起來讀.
立即吃飯.
面上再不帶愁容了.
立即吃飯.
而且最重要的是.
他不再難過.
面帶愁容.
這個很強烈的對比.
就是他對神有這個信心.

$^{1201}$今天我們有沒有同樣樂於.
信服神的應許呢.
我們有信心.
不是對人.
是對神有信心.
以致我們的生命可以改變.
當然更重要的是.
我們有沒有實際的行動.
去配合你的信心.
相信神是會祝福.
求神幫助我們.
最後兩點就是.
他對神的尊重,信守承諾.
我們要知道.
為什麼「哈拿」這麼緊要生孩子.
你有沒有留意到舊約的女士.
沒有孩子生.
是被視為.
神放棄了你的.
那是什麼原因呢.
生孩子不是為了要顯自己.
我可以生孩子.
不是的,原來是基於一點.
當人類犯罪之後.
始終亞當夏娃犯罪.
接著神就.
下了一個應許.
在《創世記》.
第三章第十五節.
我們一起來讀出.
好,我繼續.
女人的頭銳.
要傷你的頭.
你要傷她的腳筋.
當人類犯罪之後.
神就說到.
女人的後銳,這個很特別.
因為在希伯來的文字裡.
每一個人都是男人的後銳.
不是稱為女人的後銳.

$^{1241}$但是女人的後銳.
要傷蛇的頭.
就應許了救世主會出現.
去對付魔鬼撒旦.
那就是說.
如果你連孩子都不能生.
那是怎樣呢.
中的機會都沒有.
也就是說神不祝福你.
所以他自己就覺得.
我要得到神的祝福.
我要生這個孩子.
他更加跟神說.
我願意將我這個孩子獻給你.
他很清晰的.
去承諾.
我要將我的兒子獻給你.
做那世義人.
當真是.
撒慕爾出世之後.
他一直照顧著她.
到他.
斷奶的時候.
四五歲的時候.
他要將她奉獻.
每一個做母親的.
你想一想.
一個小孩子.
由懷孕到他出世.
你一直培育他.
小孩子最可愛是哪個年齡.
我有這個感受.
是三四歲那些.
我現在都很懷念.
我孫子三四歲的時候.
那些感受.
會講廣東話.
爺爺呀.
現在全部講英文.
見到我都不知怎樣講.

$^{1281}$三四歲的時候.
他講很多東西.
我現在還有影片去講.
我現在都很懷念.
但到你十幾歲的時候.
他都不跟你說話.
只會對著電話.
所以.
三四歲或四五歲.
小孩子真是最可愛.
但當他斷奶的時候.
他竟然要怎樣呢.
他就要.
將撒母耳.
離大去送走.
我們一齊離大去看.
這一段經文.
第23到28節.
我們一齊讀出.
如養兒子.
直到斷了奶.
既斷了奶就把孩子帶上士郎.
到了耶和華的殿.
又帶了三隻公牛.
一衣髮細面.
一皮袋酒.
那時孩子還小.
載了一隻公牛.
就領孩子到爾利面前.
婦人說.
主呀我敢在你面前起誓.
從前在你.
這裡站著祈求耶和華的那婦人.
就是我.
我祈求為了得這孩子.
耶和華已將我所求的賜給我了.
所以我將這孩子歸於耶和華.
使他終身歸於耶和華.
於是在那裡敬拜耶和華.
撒母耳.

$^{1321}$斷了奶.
會走會走.
也會說話.
可以和人溝通.
這時候是最好玩.
不要忘記.
那是你的獨生兒子.
絕對不容易.
但是.
哈娜願意這樣做.
她有很多藉口.
可以不去.
她還小.
這個老人家.
我兒子送了過去.
她沒有好好教他.
你看她教自己的兩個兒子.
都教得這麼差.
我怎會放心將兒子交給她.
這全部是合理的藉口.
可以不去.
但是.
哈娜沒有因為自己愛撒母耳.
而忘記了.
她之前在神面前的承諾.
她自己很愛自己的骨肉.
但她更加愛救贖她的主.
她自己.
看重自己的兒子.
但她更加看重.
她所相信的神.
就是賜兒子給她的神.
她重視自己的後裔.
但她更加重視.
神的角度.
和自己的承諾.
所以.
當孩子斷奶的時候.
她就將孩子送到耶和華的殿裡.
第27,28節.

$^{1361}$她再一次說.
我祈求為了得這孩子.
耶和華已將我所求的賜給我了.
所以我將這孩子歸於耶和華.
使他終身歸於耶和華.
於是在那裡敬拜耶和華.
我經常這樣說.
神從來沒有迫我們許下甚麼承諾.
如果神要幫我們.
即使我們不許下承諾.
神都會幫我們.
但如果我們許下承諾的時候.
我們就不可以.
不去承擔這個承諾.
不然你只是在騙神.
不可以不去承擔這個承諾.
在《申命記》裡提醒我們.
我們一起去讀《申命記》23章.
21至23節.
我們一起去讀.
你不償還就有罪.
你若不許願倒無罪.
你嘴裡所出的.
就是你口中應許.
甘心所獻的.
要照你向耶和華.
你臣所許的願.
謹守遵行.
所以記住,我再說.
神沒有迫我們許願.
但如果我們許願的時候.
就要來帶去償還.
今天我們要做一個.
信守承諾的人.
這樣我們的生命.
一定有美麗的影響力.
影響到我們的下一代.
再者我們要反思.
會不會我們的兒女.
我們很重視他們.

$^{1401}$過於神呢?.
我已經說過很多人.
可能當你有了孩子的時候.
將所有注意力集中在孩子身上.
特別是媽媽.
連自己的丈夫都輕忽了.
這樣就是一個危機.
我們怎可以來愛兒女.
多於愛神呢?.
所以我們不是不愛他們.
但仍然以神為首.
我們不要讓我們的兒女.
影響到我們與神的關係.
聖經裡提醒我們.
馬達方第十章二十三十七到三十八節.
我們一齊讀出.
愛父母過於愛我的.
不配作我的門徒.
愛兒女過於愛我的.
不配作我的門徒.
不配著他的十字架.
也不配作我的門徒.
這裡很重要.
就是今天我們以神為首.
最後一點就是.
他將自己的兒子.
奉獻了給神.
但絕對不是一個損失.
神會繼續祝福的.
我們先來看第二十一節.
伊利加拿和他全家都嘗試來去.
要向耶和華獻年祭.
並還所許的願.
他們要去獻年祭.
還神所許的願.
但他們仍然很年輕.
大家不要忘記.
在那時候.
伊利加拿大一定很擔心.
自己的兒子.

$^{1441}$那時候沒有智能電話.
沒有FaceTime.
他們也不是經常.
可以探望這個兒子.
一年可能去一兩次.
但每一次去的時候.
你看這個媽媽.
多麼愛自己的兒子.
我們一起來讀.
撒母耳記上第二章.
十八到第十九節.
那時撒母耳還是孩子.
穿著細麻布的耳筆.
是立在耶和華面前.
他母親每年.
為他作一件小外袍.
同著丈夫上來.
獻年祭的時候.
大來給他.
這個詞沒有多愛自己的孩子.
一年可能見一兩次.
每次他都會來帶去.
弄一件衣服給他.
讓他可以替換.
自然就想起一首.
唐詩.
慈母守中線.
由子身上衣.
臨行密密逢.
義勇馳馳歸.
誰言寸草心.
步得三春輝.
我親身經歷.
現在新一代.
太多享受.
太容易.
我第一次坐飛機是1975年.
我去菲律賓.
出席一個國際性的會議.
叫Love China 75.

$^{1481}$我住在菲律賓契妹的家.
我媽媽.
檢視我的行李.
她就帶了這衣服去.
沒有睡衣的嗎?.
我們那個年代.
怎會穿睡衣.
1975年的時候.
面衣就是睡衣.
一天在家裡住.
我媽媽說.
那一代的母親真的很厲害.
一晚之內.
拿一塊布.
馬上縫一套睡衣給我.
我到現在都記得那套睡衣是什麼顏色.
是怎樣的款式.
一晚之內就縫好了.
慈母守中線.
由子身上衣.
當我讀神學的時候.
我媽媽不是很高學識的人.
很少出街.
我見到神學院讀書.
我媽媽竟然走進長洲.
拿一件棉被給我.
記住不是棉被.
是壓下來的.
很重的那些.
她知道我喜歡吃香蕉.
還帶了一串香蕉給我.
我記得我去碼頭.
去接我媽媽.
因為我媽媽很少出街.
是她人生第一次自己坐船進長洲.
就是她對我的愛.
所以記不記得我開始講的時候.
我很遺憾.
我沒有看到她的愛.
你看到哈娜每一年.

$^{1521}$去探望自己的兒子的時候.
為自己的兒子縫一件的衣服.
是多麼美麗的一件事.
同樣你有沒有留意到.
她都沒有忽略.
對神的呈獻.
去到神面前她都有獻祭.
她不但獻上自己的孩子.
她更加獻上禮物.
所以今天.
我們絕對不可以以為.
我已經事奉神了.
神還代替了我們金錢的奉獻.
但是.
哈娜是沒有損失.
她好像.
沒有了撒慕兒.
但是神賜福給她.
結束了.
我們一起來讀撒慕兒寄上第二章.
第二章是下一節.
沒有了一個.
多了五個.
比生這麼多個更辛苦.
加上六個是很慘的.
我媽媽也有生過六個.
其實應該七個.
有一個夭折了.
生這麼多個不容易的.
但是神賜福給她.
她好像沒有了撒慕兒.
她呈獻了給神.
但是神仍然是她的兒子.
但神仍然多給她三個兒子和兩個女兒.
所以今天我們為神擺上的時候.
絕對不是一個損失.
是一個祝福.
施比受.
更為有福.
神自己親自宣告.

$^{1561}$我們一起來讀.
撒慕兒寄上第二章第三十節.
因為尊重我的.
我必看從他.
藐視我的.
他必被輕視.
如果我們尊重神.
神會重視我們.
我們輕忽了神.
我們也不會得到祝福.
所以聖經裡很清楚說.
這人是神所知道的.
所以母親節裡的一個信息.
提醒我們每一個人.
因為哈拿她的生命.
她對神的愛.
對神的尊重.
她用她的生命.
來教導自己的孩子.
就促使了撒慕兒.
成為.
在當世代裡.
一個屬靈偉人.
很願意今天.
我們每一個姊妹做母親.
你有哈拿.
這麼美麗的生命.
有這一份.
親近神的心.
有忍耐.
或許面對挑戰.
但你會向神傾訴.
更加重要的是.
重視自己的配偶和孩子.
影響他們.
這樣的時候.
你的孩子一定會蒙福.
而你更加是一個蒙神祝福的人.
很願意.
在母親節這一天裡.

$^{1601}$跟大家來思想.
這個屬靈偉人的母親.
那種質素.
讓我們去學習.
同樣我已經將筆記交給教會.
大家都可以來大日去再看.
我們一起來禱告.
天父我多謝你讓孩子有機會.
和弟兄姊妹來大日去分享.
神的願意你的話語.
都成為我們的提醒和幫助.
祝福今天在你面前每一個姊妹們.
特別做母親的.
或者做婆婆.
或者奶奶的.
他們都很辛苦.
求主你來大日去眷顧他們.
他們所努力的付出.
你來結納他們所做的工作.
但更加透過他們的生命.
可以成為孩子的祝福.
願神你親自幫助他們.
聽下在你面前的禱告交託.
來奉主特赦的名.
\newpage



\section{路加福音 1:46-55}
\label{sec:032132}
\textbf{生命的逆轉}~~~~梁錫強醫生
\newline
\newline
連結: \href{http://www.church.com.hk/acms/content.asp?site=cdc\&op=show\&type=product\&code=032132\&layout=sermon}{\texttt{032132}} ~~~~ 語音日期: 2026-05-17
\newline
\newline
\hyperref[sec:032131]{\small{< < < PREV SERMON < < <}}
~
\hyperref[sec:index_chronic]{\small{[返順時目]}}
~
\hyperref[sec:032133]{\small{> > > NEXT SERMON > > >}}
\newline
\newline
路加福音 1:46-55
\newline
\begin{longtable}{cl}
\hline
\hline
章節 & 經文 (和合本修訂版)\\
\hline
1:46 & \begin{tabularx}{0.7\textwidth}{X} 馬利亞說：「我心尊主為大； \end{tabularx} \\ \\ \relax
1:47 & \begin{tabularx}{0.7\textwidth}{X} 我靈以神我的救主為樂； \end{tabularx} \\ \\ \relax
1:48 & \begin{tabularx}{0.7\textwidth}{X} 因為他顧念他使女的卑微；從今以後，萬代要稱我有福。 \end{tabularx} \\ \\ \relax
1:49 & \begin{tabularx}{0.7\textwidth}{X} 因為那有權能的為我做了大事；他的名是聖的。 \end{tabularx} \\ \\ \relax
1:50 & \begin{tabularx}{0.7\textwidth}{X} 他憐憫敬畏他的人，直到世世代代。 \end{tabularx} \\ \\ \relax
1:51 & \begin{tabularx}{0.7\textwidth}{X} 他用膀臂施展大能；他趕散心裡妄想的狂傲人。 \end{tabularx} \\ \\ \relax
1:52 & \begin{tabularx}{0.7\textwidth}{X} 他叫有權柄的失位，叫卑賤的升高。 \end{tabularx} \\ \\ \relax
1:53 & \begin{tabularx}{0.7\textwidth}{X} 他叫飢餓的飽餐美食，叫富足的空手回去。 \end{tabularx} \\ \\ \relax
1:54 & \begin{tabularx}{0.7\textwidth}{X} 他扶助了他的僕人以色列，不忘記施憐憫， \end{tabularx} \\ \\ \relax
1:55 & \begin{tabularx}{0.7\textwidth}{X} 正如他對我們的列祖說過，『憐憫亞伯拉罕和他的後裔，直到永遠。』」 \end{tabularx} \\ \\
[1ex]
\hline
\hline
\end{longtable}
$^{1}$大英姐妹主內平安.
很高興再一次在港台和大家見面.
今天的題目是生命的逆轉.
所謂逆轉就是一個叫reversal.
其實我們很多時候.
我們想的逆轉就是我們的環境改變了.
可能以前住屋村 後來參加港姐.
就贏了半軍 嫁了豪門.
我不是說我.
這種的逆轉就是生命一個完全的改變.
我們世人看的逆轉就是.
一個生命的環境很大的改變.
可能你很多的生意做到很大.
但是可能金融風暴 令到你破產.
所以破產之後可能屋又要賣.
車又要賣 變成一無所有.
這種的逆轉都不一定.
今天我們說的題目是為什麼?.
因為聖經裡面有很多很多的逆轉.
也是一個上帝給我們從舊約到新約.
一個中心的神學意思在裡面.
之前也說過有不同的婦女.
聖經裡面其實不是說很多婦女.
可能都是比較上當時的時代.
都是重男輕女 就像我們中國古代一樣.
就變成聖經很少提及一些婦女.
但是當它提及一些婦女的時候.
很多時候都是影響很大的.
大家剛才看那段的經文.
是路加福音就是瑪利亞的歌頌.
大家都很清楚知道.
瑪利亞不是一個普通的婦女.
本來是一個很普通的.
但是逆轉了之後就真的不普通了.
到今時今日我們知道瑪利亞是誰.
就算你不信基督教 不信天主教.
你問她瑪利亞是誰 耶穌是媽媽.
其實是知道的.
因為她真是一個可以說是歷史上.
聖經上一個很重要的人物.

$^{41}$之前我其實也有說過.
上次我們說過路德 是不是?.
路德也是一個婦女.
也是一個受人歧視的.
因為她是一個摩訶.
一個外邦的女子.
她也是一個寡婦.
婦女已經被人歧視了.
跟著外邦人更加被人歧視.
你是寡婦就更加被人歧視.
當時的年代就是這樣的.
你是寡婦就不是怪你老公死了.
是怪你把你老公嚇死.
通常人們是這樣想的.
但是後來路德.
後來成為一個什麼人呢?.
大家都知道了.
因為他的忠誠 他對上帝的忠誠.
他甘願跟隨他的奶奶.
回到這個南俄米 回到這個猶大的地方.
結果他後來就是再嫁了.
嫁給了一個那個人.
那個人後來就生了一個兒子.
就成為了大衛的祖先.
也是耶穌基督的家譜.
大家如果看《馬太福音》第一章的時候.
看到耶穌的家譜是有這個路德的名字在裡面.
其實那個家譜通常都是說男人的.
特別是提到這個路德.
這個是他生命的逆轉.
還有其它很多例子的.
之前我們看《約書記》.
打耶利哥城的時候.
就派了一些探子去的.
探子幾乎沒命 全靠誰呢?.
全靠拉俠 拉俠是誰呢?.
我說拉俠 大家可能不記得了.
如果我說妓女拉俠 大家就記得了.
其實對她都很不公平的 知不知道?.
每次說拉俠一定要加上妓女拉俠.

$^{81}$我們就記得 就是那個.
但是因為她的忠誠.
她救回了探子.
間接令到以色列人打贏這場仗.
結果她是在家譜裡面有份.
這個也是一個很神奇的事情.
我們知道後來她嫁給了大衛的祖先.
剛才看回《路加福音》的時候.
很多時候我們就當故事看.
我們看《路加福音》 我們看耶穌基督出生.
我們會留意到一些故事.
就算我們在教主教學的時候.
都是講一些故事出來的.
《路加福音》一開始的時候.
祂是預言這個屍體約翰的誕生.
天使來到以利沙白那裡跟祂說.
說你會有一個屍體約翰.
這個約翰 當時不知道祂是屬於屍體的.
是會出生的.
接著就預言耶穌基督降生 大家都很熟悉.
每天聖誕節我們都會讀到這段經文.
《路加福音》一章30節那裡.
瑪利亞這個天使長 拍一個更加高的天使長.
加巴黎來到 基督來到就說.
瑪利亞不要怕 你在上帝面前已經蒙恩了.
你要懷孕生子 給他起名耶穌.
他將要為大 稱為至高者的兒子.
求主上帝要把祂祖先大位的王位給他.
這個就是大家都很熟悉的一段.
接著就發生什麼事呢?.
大家看下去就是瑪利亞去探望以利沙白.
原來他們有親戚關係的.
當她探望的時候 以利沙白懷著女兒的胎兒.
當時已經懷胎幾個月了 開始跳動.
以利沙白被聖靈充滿.
就說你跟瑪利亞說 你是婦女中最有福的.
你所懷的胎兒也是有福的.
接著就說先世約翰出生.
接著我們就跳到第二章看耶穌基督出世.
但原來不是的 有一個很重要的一段.

$^{121}$就是瑪利亞的專注中.
很多時候我們就奇怪 講開歷史 講開故事.
突然間瑪利亞在唱歌 但這些才是重點.
瑪利亞所講的東西是重點.
也是盧家專門放她進去這段.
就是因為這個是配合整個盧家福音的主題.
為什麼呢? 大家剛才看那段聖經就知道了.
瑪利亞開始就是我心尊主為大.
我靈以上帝 我救主為大.
大家都很熟悉這兩句話.
很多詩歌都是用這兩句話 瑪利亞專注中.
但當她看下去的時候 就好像有些離題.
她講些什麼呢? 好像有個預言.
就是她看到那段 她說.
她趕散心裡妄想的狂傲人.
她叫有權柄的失去 叫卑賤的人升高.
她叫飢餓的人飽吃美餐 叫富足的人空手回去.
她的名成為聖 這就是講耶穌基督.
她來到是做什麼呢? 很大的逆轉.
狂傲人就被趕散 有權位的人失去.
卑賤的人反而升高.
飢餓的人反而有飽餐的美食.
以為自己是富足的人 結果就空手回去.
這就是很大的逆轉 是不是?.
瑪利亞就預言了這件事出來.
放在聖經《路加福音》第一章最後那裡.
其實正正就是耶穌基督來到.
祂要顛覆這個世界 顛覆當時人的思維.
祂來一個大的逆轉 之前都講過了.
就是《路加福音》一開始第四章的時候.
耶穌出來傳道 第一件事是做什麼呢?.
第一件事不是說祂行了什麼神蹟.
也不是說祂醫了什麼病.
一開始第四章《路加福音》.
說到祂一出來的時候 受完曠野試探之後.
第三章第四章講祂出來傳道的時候.
祂去到安息日 去到猶太人的會堂.
拿起這個以賽亞書去讀 那段經文是什麼呢?.
祂說「被擄的人得釋放 失明的人得看見.
受壓迫的人得自由 宣告上帝閱立人的欺憐」.

$^{161}$之前我講欺憐都講過一段.
原來耶穌基督來到是顛覆當時人的思維.
被擄的人 被人抓了他反而可以釋放.
看不到的人可以看見 受壓迫的人可以得到自由.
這個就是聖經告訴我們 耶穌基督來的目的.
就是顛覆當時人的思維.
但是其實不是什麼大的顛覆.
因為在整個舊約的歷史當中.
我們看到的就是正正常常的作為就是這樣.
以色列人在埃及的地方做奴隸很多年.
結果是摩西帶領他們.
我英文叫Against all odds.
其實是不可能的.
法老王那麼厲害 他的軍隊那麼厲害.
但是居然摩西是可以有這樣的能力.
帶領一幫以色列人離開埃及地.
當然這個是上帝的大能.
失明的人可以看到.
沒錯 後來耶穌真的去醫治了很多失明的人.
但是這個更加有一個屬靈的意思.
就是以前看不到上帝的人.
現在看得很清楚了.
這個是一個逆轉.
聖經有很多的例子.
約伯就是其中一個例子.
約伯是什麼人?.
約伯是一個義人.
約伯是一個好人.
他是有錢的.
有錢不是犯法的.
有錢是好.
但是這個是上帝賜給他的.
但是他或者他身邊的人.
是以他的財物來衡量他有沒有福氣.
聖經說他是有很多牛很多羊.
有很多的僕人 兒孫滿堂的時候.
上帝就拿走.
結果他是一無所有.
經歷了一無所有之後.
上帝賜給他.

$^{201}$他就真的明白了.
收取的耶和華.
賞賜的也是耶和華.
他開始明白了.
以前的豐文有上帝.
今天他是親眼看到上帝.
以前他是盲的.
現在是看到.
失明的人可以看到.
之前不會看到上帝.
現在是看到上帝.
雅各也是一樣.
要說舊約的雅各.
那個頑皮仔雅各.
其實他不是好人.
約瑟的爸爸.
他12個兒子.
12支派這樣來的.
他是妒忌他哥哥.
他是很貪心.
很貪戀權利.
他媽媽又很偏心.
很多事情加起來.
全部都是人性的罪.
結果他真的是要受過苦.
被放逐 被他哥哥追殺.
去到那裡被他舅父又騙了很多年.
結果他後來回來的時候.
他仍然還是不明白.
很多事情看不到.
看不到上帝的旨意.
結果怎麼樣?.
結果有一晚他和那個天使摔跤.
鬥到天亮.
這個是神學的一個難處.
為什麼那個天使是誰?.
究竟是誰打贏?.
很多不同的解讀.
但最重要的是經過這個摔跤.
他看到上帝.

$^{241}$他成為12支派的族長.
他這個很清楚他的逆轉.
就是因為從他失明得到漢見.
明白了沒有?.
所以法利賽人耶穌形容他們是盲的.
是他讀很多書.
很熟悉那些禮儀.
很熟悉那些律法.
但是在屬靈方面他們是看不到的.
看不到的.
受壓迫的人得自由.
這個也是耶穌基督來的.
本來受壓迫的現在得自由.
今天我們知道.
這一期我們一直在說舊約.
用新約來引證舊約.
發現在《撒母耳記》上一章有一個哈拿的故事.
其實和瑪利亞的土種是很相似的.
哈拿是誰呢?.
Hannah是誰呢?.
《撒母耳記》上一開始不是說撒母耳.
是說撒母耳的媽媽.
這個哈拿聖經形容她是一個婦人.
當然她是以利加拿的太太.
這個以利加拿是猶太人.
她有兩個妻子.
一個叫做哈拿.
另外一個叫做比利娜.
我們重點不是要記住她的名字.
重點就是這個比利娜就有生了很多孩子.
就是有小孩的.
哈拿沒有小孩.
那你說沒有就沒有吧.
不要緊啦 她老公也行.
傳宗接代就行了.
要明白當時那個年代那個世代不是這樣想的.
你沒有小孩就是你沒有福氣.
你不可以為這個家開枝散葉就是你沒用.
沒得生一定是賴那個女人的.
所以她是被人歧視.

$^{281}$被很多人歧視.
甚至是被另外一個妻子歧視.
我生了這麼多個了 你一個蛋都沒有.
我以前看那些粵語片最流行這樣說.
所以她是被人歧視的.
所以她很沒趣的 你們不知道嗎?.
一到什麼大場合 喜慶啊.
那個小的老婆就出來講風涼話.
她整天都在哭 整天都在哭.
但是那個老公對她很好.
這個就不跟粵語片那個橋段 你知道嗎?.
她老公是很疼她的.
她為什麼哭呢? 她說沒有小孩.
她說有我就行.
其實是很好的老公.
但是她覺得她做不到妻子的責任.
有一次哈娜來到耶和華的殿外那裡去祈禱.
當時是沒有聖殿的 慈思時代.
所謂耶和華的殿就是會幕.
會幕在哪裡呢? 在士羅的地方.
大家可能不是很清楚.
但是士羅很重要 在聖經裡面的歷史.
是很早期的 可以說是以色列的首都.
就在士羅 不是耶路撒冷的.
這個士羅的地方 會幕在那裡幾百年.
是在慈思時代的會幕就在那裡了.
她就在那裡 女人不可以進去的.
就在外面那裡祈禱.
但是裡面的大祭司以利就聽到她的禱告.
一開始不知道她在做什麼.
沉沉的 以為她喝醉了 但是後來她說.
她說 以利說你在說什麼?.
她說我是許願.
我說萬軍的耶和華.
如果你誰顧你侍女的苦情.
眷戀不忘你的侍女.
侍女侍女一個子字.
我必使她終生歸給耶和華.
不用剃刀剃她的頭.
這個以前猶太人很流行許願.

$^{321}$就是今天說得不好聽.
就好像和上帝討價還價.
上帝你給我這個 我就給你這個.
我們今天不是鼓勵大家用這一種.
因為有更加深的意思的 猶太人的許願.
要還的 你不要亂說 要還願的.
不要說了當吃生菜這樣不行的 知道嗎?.
雖然我們英國經常吃生菜.
但是這些就不要當吃生菜.
她說如果我小時候有兒子的話.
你給我這個兒子.
將來我們是歸還給.
我會歸還給耶和華.
這個簡直是奉獻給耶和華.
結果撒母耳在會幕那裡大同意你一起.
她說不用剃刀剃她的頭.
這個也有個意思.
大家那次說參選說過了.
就是她是做一個納西爾人.
是一個很忠心的 是信服上帝的人.
是一個標記的.
以利就不可以代表上帝回答她.
OK 給你吧 不是的.
以利沒什麼能力.
她只是說平平安安回去.
她說願以色列的上帝允許你所求的.
她希望以色列的上帝.
就是耶和華可以答應你.
結果很快哈娜就生了一個兒子.
這個兒子就是撒母耳.
這個意思就是上帝聽到她的苦情.
今天的重點是說什麼呢.
是說哈娜在二章裡面有個禱告.
這個禱告和瑪利亞的重禱是很有相關的.
哈娜的禱告 開頭那些我先不說了.
說到中間那一段又是一個逆轉.
為什麼她這樣說呢.
因為這個禱告是上帝啟示她說的.
如果普通人不會突然說這種話的.
她說什麼呢.

$^{361}$她說勇士的弓會截斷.
跌倒的人可以有力量.
飽足的人要做故宮來求食.
飢餓的人就不再飢餓.
不生育的生了七個.
兒女多的反而孤獨.
然後耶和華使人死也使人活.
使人下陰間也使人往上升.
耶和華使人貧窮也使人富足.
使人降低也使人升高.
全部都是逆轉 是不是.
勇士 勇士靠什麼謀生呢.
不是拿槍謀生 那時沒有槍.
拿個弓箭.
但是如果弓箭斷了的話.
就什麼能力都沒有了.
你多勇都沒有用 是不是.
勇士的弓會截斷.
跌倒的人反而有力量.
吃得飽的人反而沒得吃.
要做宮人去求食.
反而飢餓的人不再飢餓.
這全部都是逆轉.
沒有生育的生了七個.
但是兒女多的反而孤獨.
很奇怪.
兒女多的反而孤獨就不是很奇怪.
不要以為生得多就一定好.
可能個個不孝順就真是浪費力氣.
李家鼎生兩個就已經頭癢 是不是.
所以不一定是好的.
但最重要的是主權在上帝那裡.
逆轉的權在上帝那裡.
耶和華使人死使人活.
使人下陰間使人上升.
使人貧窮使人富足.
使人降臂使人升高.
這個就是哈拿的禱告.
很奇怪 哈拿和瑪利亞是相差幾百近千年.
但是講的意思都是聖經的中心來的.

$^{401}$就是一個逆轉的福音.
生命的逆轉.
從中我們學到什麼呢?.
逆轉是從看到神開始.
如果你說生命逆轉.
我們很多時候就說環境改變就逆轉.
我種了六合彩就逆轉 是不是.
我們是這樣想的.
我們用這些東西來衡量我們的生命是否改變.
用環境來看.
但是聖經說不是 逆轉的開始是你看到神.
逆轉的開始就是我們信了耶穌.
上帝改變了我們的心.
而不是改變我們的環境.
環境會隨之慢慢改變.
但是最重要的是你的心態改變.
所以仍然是有苦難.
仍然是有很多不如意的時候.
仍然你的老闆都是這麼刻薄.
仍然都是你的子女這麼不乖.
不會即時改變.
但是你的心態,你的看法是改變的.
所以哈拿第一句的禱告就說.
我的心是因耶和華快樂.
而瑪利亞就說我的心尊主偉大.
她的禱告第一句就是尊主偉大.
心是因為耶和華而快樂.
這個改變是從上帝而來.
改變是上帝改變了她的心.
不是改變了環境.
沒錯 環境是改變的.
瑪利亞生了耶穌出來.
哈拿後來生了撒慕爾.
這個不是最重點.
不是說一開始就已經是這樣.
這個是隨著上帝有她的時間表.
慢慢去成就.
這個是逆轉.
這兩個女人本來是什麼都沒有的.
她是很普通的人來的.

$^{441}$就如今天我們所說的.
她們是沒有什麼社會的資源.
什麼叫社會資源.
社會資源我們現在經常說的就是.
我們有多少社會資源就是什麼呢?.
有多少人力,物力,財力.
原來我們在這個世界要生存.
要靠這幾樣東西.
有人力,物力,財力.
不然為什麼我們要培養小朋友.
一開始就要這樣呢?.
要名校,幼稚園呢?.
為什麼要和他做一個Portfolio.
去入小學呢?.
為什麼我們要擺最好呢?.
為什麼我們要去打一些關係,認識多一些人呢?.
是不是?.
為什麼那個幼稚園名校,幼稚園那個.
那個誰啊?.
那個老師在那裡做.
那就去認識一下他吧.
這樣就可以有些關係.
我們是這樣的.
他今天出來做事.
我們都是靠什麼呢?.
靠人脈 是不是啊?.
很多時候.
人們說不是你認識多少東西.
是認識多少人.
你找工作能找到嗎?.
不是你認識不認識東西.
不是what you know 是什麼呢?.
是who you know 嘛.
經常都說的.
實在又是.
這個世界就是這樣.
我兩個女兒她都在商界發展.
不過她很差勁.
香港人很了不起.
她不是像我做的那些.

$^{481}$有時候和一些人要出來聊聊天.
約一下別家公司的人.
吃一下飯,喝一下酒.
她說搞好social networking.
原來這個字很流行.
要搞好一些社交.
社交的那個關係.
網絡 社交網絡.
她說可以啊.
將來轉工比較容易.
認識了那邊的人.
我轉什麼工呢?.
好像我這樣做NHS.
做到退休就算了.
她說你不明白的.
現在不是這樣的年代.
她說你不轉工呢.
就不會升職的.
她說魚不過塘就不會肥的.
我女兒當然不會說這句.
這個我幫她加進去的.
那又是真的.
她真的要這樣去認識一下人.
搞一下這個搞一下那個.
這些全部是人.
這個世界造成他們要這樣做.
但上帝的逆轉不是這樣的.
上帝的逆轉不是靠你自己做了多少.
那哈娜和瑪麗亞是什麼人呢?.
很普通的人.
瑪麗亞大家熟悉多一點.
她是一個猶太這個地方.
一個很普通的婦女.
天使找她的時候.
她已經許配給了藥室.
她很奇怪.
你說她是童貞 瑪麗亞 童女.
其實都不是很奇怪.
那個年代到處都是童女.
你明白嗎?.

$^{521}$那個年代是不同現在.
也不是什麼特別個人之處.
上帝揀選她.
不是因為她有什麼賤職.
不是因為她的Portfolio有多好.
不是因為她做過什麼.
經常都說.
上帝是揀選她.
是有上帝的原因.
上帝認為她可以為上帝做多少事.
之前我講參宣,講機電都是一樣.
不是因為她之前做了多少事.
而是上帝認為她可以為上帝做多少事.
這個就是瑪麗亞.
這個就是哈娜.
普通人.
但她生命得到一個很大的逆轉.
但他們不要忘記.
她知道逆轉是從上帝而來.
所以他們的土文.
第一句就是尊主為大.
第一句就是以耶和華為樂.
他們兩個都是很痛苦,很卑微的年代出生.
瑪麗亞不是.
瑪麗亞普通童女而已.
不要忘記.
那個年代你婦女已經是卑微的了.
被人歧視的.
真的.
那個年代是這樣的.
但是上帝不是先逆轉他們的人生.
而是逆轉了他們的心.
讓他們明白.
今天我們如果反思一下.
今天我們面對很多很多的問題.
我們被很多問題困擾著.
有時候我們透不過氣.
很多壓力.
無論是工作,家庭.
或者教會的事奉.

$^{561}$都可能有壓力.
我們在想很多很多問題.
問題天天都有.
是不是天天都多?.
但是我們問我們自己.
我們今天想的東西.
我們正在想的是問題.
還是我們正在想我們有沒有看見上帝?.
如果我們逆轉一下.
顛覆一下我們的思維.
我們正在看上帝.
其他的問題就會解決.
你整晚去想這個問題.
你為你自己的健康憂慮.
為你的工作憂慮.
為你的子女憂慮.
你解決不了的.
很多事情不是我們可以解決的.
我們知道的.
但是我們知道解決不了.
我們都去想.
一定的.
就是交給上帝.
我們不要想著這個問題.
我們想著究竟我們是不是.
定睛的仰望上帝.
瑪利亞那首詩的討文.
就是那個眾討和哈拿的討文.
很清楚是說到這個逆轉.
就是逆轉是從上帝而來.
是上帝顛覆我們的一生.
甚至上帝是顛覆世界的思維.
為什麼?.
因為這個世界的人想的是想什麼呢?.
我們要強嘛.
我們要厲害嘛.
我們要比別人厲害.
要贏過別人嘛.
是不是?.
這個是從小到大已經有這樣的思維的了.

$^{601}$小朋友如果長得高大就比較有好處的.
實在這個世界又是這樣的.
有數據高大的找工作成功率是高一些的.
不是假的.
有人做過統計.
這個世界是這樣的.
人是看這些的.
看外表的.
香港的電視台報告天氣.
都是找個美女出來報告的.
是不是?.
這裡不是的.
要找個師奶阿婆的.
因為她有經驗.
不是看外表.
如果看外表呢.
如果香港找個師奶阿婆就沒有人看天氣報告.
層次不同的.
但是人就是這樣的.
我們要贏.
我們要高大.
小時候就已經和隔壁的B仔量高的.
我就很不開心的了.
B仔當然比我高一吋.
那又怎樣呢?.
但是要做這些事的.
要贏.
你兒子考第幾.
以前有考第幾的.
現在就政治不正確.
好像沒有考第幾的.
我哥就考第幾.
比較.
要面子的.
因為如果你兒子讀書學得好.
你爸爸也有面子的.
我們是看這些.
我們所謂的資源.
世界的資源就是要強,要贏,要面子.
但是上帝說不是.

$^{641}$你要謙卑.
好像剛才那首歌.
記不記得?.
最後一句要謙卑.
有些歌很有趣的.
可能是領事妹妹.
硬要將它推排到.
那首歌好像還沒唱完.
它不是回到主音那裡停.
有些人說那首歌好像還沒完.
但是有時候專門是這樣.
就記得了.
好像無線的劇集.
不講結局給你聽.
你整天想著.
究竟會怎樣呢?.
有時候專門是這樣安排的.
那首歌叫我們謙卑下來.
上帝看是要謙卑,要依靠它.
要有信心,要有信服.
就不是要強,要贏,要面子.
所以上帝逆轉你的生命的起點.
就是要你變得更加強.
首先你要謙卑下來.
這個就是逆轉.
謙卑你才可以強.
耶穌基督的十字架.
祂說死了之後三天.
祂是下到陰間再升上來.
這個是一定要這樣做的.
今天你要接受上帝.
你要倒空自己.
將不好的東西全部倒出來.
你才可以接受新的進來.
是不是?.
那些人說呼吸的字.
為什麼叫呼吸呢?.
不是吸呼呢?.
你要呼出來才可以吸進去.
想想也是.

$^{681}$裡面裝滿滿你還怎麼吸呢?.
吸不了.
你要先呼出來才可以吸.
這個就是上帝教導我們的.
我們很明白這件事.
但我們世人不是.
我們看東西是很簡單的.
我們以為逆轉就是.
剛才說嫁到豪門.
就像中了六合彩那樣.
是 你生命可能改變.
但是你是不是真是.
你是物質上改變.
那些呼籲.
但是不是你人生有一個很大的改變呢?.
英國有六合彩.
其實就是二十多三十年歷史.
都不是很久的.
這幾十年就很多人中過了猴獎.
當然也是.
生命是改變了.
也有些人做了一些研究.
去追蹤他們過去的十年,二十年.
你發覺什麼呢?.
生命是改變了.
但是不一定是好的.
有很多是好不好的收場.
其實是大半是不好的收場.
不開心的收場.
有些就是突然間暴富.
暴發戶.
他們不懂得怎樣去管理那些錢.
因為那些錢不是你一筆錢一筆錢掙出來的.
不是你這樣慢慢建立.
你就一直投資.
你有很多方法.
突然間給你一億,半億.
你怎麼用?.
你放在哪裡?.
你只要.

$^{721}$要麼放在這裡.
放在這裡又不值錢 是吧?.
那怎麼辦呢?.
那就聽一下投資顧問.
他們又抽水又找你笨.
接著你的親戚朋友知道你中了.
又走去問你借.
借一次不借一次又不行.
那個人借你又借.
那個人又不借說你不公平.
搞到你天天為這些擔心.
是吧?.
結果他們不開心.
結果怎麼樣?.
有些人選擇離開自己的地方.
去一個沒有人認識的地方最好了.
就是眾叛親離.
就是你不要找我.
我不認識你.
我從此失蹤.
有些選擇就是錢放在這裡不用.
不用沒人知道我發財了.
有什麼用呢?.
錢不用又沒人知道你發財.
你又帶不走.
那有什麼用呢?.
那些搬去那些地方住.
就是搬來這個孫臣衣的舊泥區.
就去到當是Chartingham.
那裡的人不接受你.
他們是貴族錢.
那些Old Money.
個個已經建立起來了.
跟你說這些.
你兒女讀書哪間學校好啊 什麼名牌.
你都不懂 什麼都不懂.
是吧?.
不懂.
做不到 交不到朋友給他們.
脫節.

$^{761}$舊的朋友沒有了.
新的朋友交不到.
他們小看你.
就是暴發戶.
不開心的.
如果物質上看生命的逆轉.
你未必一定是好的.
這個不是上帝的旨意.
在你身上的計劃.
告訴你.
知不知道?.
所以不要待會崇拜完.
出去買張彩票.
這個不是上帝的身上的計劃.
我肯定不是.
我告訴你.
上帝不是要我們這樣.
知不知道?.
上帝不是要我們這樣.
很多人就是因為發達.
就變成了.
七爺子山也有.
什麼都有.
上帝是怎樣呢?.
上帝是他.
他改變我們的一生.
是由改變我們的心.
讓我們面對這個世界.
看的是闊一點.
看的是深度一點.
更加去依靠上帝.
這個就是我們的生命的逆轉.
這個是我們信了主的生命的逆轉.
但上帝是聽人禱告.
尤其是哈拿的禱告.
上帝是有聽.
很多時候我們覺得.
祈禱是不是這麼重要.
我們會懷疑.
祈禱了很多年.

$^{801}$上帝好像都沒有聽.
還有些人甚至說.
我不祈禱都可以了.
上帝都知道.
怎麼會不知道上帝有什麼傳之?.
但上帝是要我們去祈禱.
哈拿真的走去聖殿那裡祈禱.
他不走去聖殿祈禱.
大祭司是聽不到的.
上帝是聽到的.
在你心裡面他已經知道你想要什麼.
但上帝要我們去祈禱.
要我們將所要的話給他聽.
將我們所要的話給上帝聽.
祂就會保守我們心懷的意念.
上帝是看到的.
祂是聽到的.
祂是紀念我們每一件的苦情.
在適當時候合乎祂的旨意.
祂會介入.
但就不是絕對.
所以我們祈禱不是用我們的意思.
說上帝不如你介入.
你幫我這樣這樣搞定.
不是.
如果是這樣.
你就比上帝厲害.
你是指揮上帝怎麼做.
我們的神不是這樣.
我們的神很差勁.
如果我們的神是有求必應.
祂不是厲害.
祂是差勁.
祂要聽我們的說話.
是很簡單的道理.
但上帝要我們祈禱要跟我們說.
昨天我們教會有一個祈禱日.
非常好的.
應該是需要多搞一些這樣的祈禱日.
很可惜不是很多人參加.

$^{841}$預了不會有很多人參加.
陳全道很早回來.
找了幾個弟兄姊妹幫他搬椅子.
他不是這樣擺的.
有一個特別的模式擺的.
昨天的祈禱會.
擺了一些椅子在中間坐的.
旁邊給我們擺一些行李.
手機,iPad,褸.
讓你可以專注在中間.
陳全道一來到就跟我說.
Andrew我擺了三十張椅子.
我的信心是去到三十個人.
結果只有二十多人.
都挺好的.
我不是說你們為什麼不參加祈禱會.
其實不錯的.
我們教會做這些已經是不錯.
不過可以更加好.
我見過很大的教會.
以前我回香港教會.
幾百人,成千人.
他的祈禱會都是二十多人.
所以其實我們不是差.
不過就一定要注重禱告.
上帝要我們禱告.
祂是聽禱告的神.
祂聽哈拿的禱告.
也會聽我們的禱告.
這是很重要的一件事.
提醒我們.
生命的逆轉.
是要我們跟神去求問.
是要神去介入的改變.
而且是改變了我們的心.
結果是什麼呢?.
結果就是瑪利亞成為了救主的母親.
厲害了.
本來她很卑微.
但是在歷史裡面.

$^{881}$今天每個人都記得她.
她是聖母.
她是耶穌基督的媽媽.
哈拿呢?.
哈拿的兒子是撒母耳.
撒母耳後來取代了以利的位置.
他成為了大祭司.
他也成為了以色列的領袖.
大衛王也是他高舉的.
安慰他出來做皇帝.
所以耶穌基督的族譜可以說是推順了.
有撒母耳的工作在裡面.
雖然他不是在那條線上.
但是他的工作.
因為祭司一定要是亞倫的後裔.
祭司就一定是葉美美人的後裔.
但是他有不同的工作.
他這個工作是很重要的.
撒母耳大家知道是一個很偉大的先知.
在教育那裡.
所以他們很清楚.
哈拿和瑪利亞都很清楚.
他們這個逆轉不只是逆轉他們自己.
他們的生命是改變了.
哈拿就受人尊敬了.
因為生了個兒子.
有些人不敢出聲.
那個小婆不敢說她了.
不敢說她生了個兒子.
她是有轉變的.
瑪利亞也是有轉變的.
受人敬仰.
但是他知道他逆轉不是個人的逆轉.
是整個上帝的角度來一個顛覆逆轉.
耶穌基督來到這個世界上.
就是作為全人類的救贖者.
撒母耳是帶給以色列一個屬靈上的更新.
因為那時候他們已經很一般了.
大家知道時思時代那幫人都是一般.
是撒母耳出來之後他們可以好一點.

$^{921}$結果後來就正式立國.
有皇帝,有大衛王,有所羅門王,有起聖殿.
輝煌的一段時間.
當然你看下去他們又衰了.
結果整個國家分裂.
但是撒母耳當時他的角色就是令到.
以色列整族人一個屬靈的更新.
所以這個是上帝角度的逆轉.
所以今天我們信了主就是一樣.
不要想著就是上帝逆轉了我的生命.
很高興了 我可以上天堂了.
OK 搞定.
不是 上帝說不是.
你只不過是你的逆轉.
也是我的角度來一個逆轉.
你是part of我這個大的計劃.
所以上帝給我們每一個人有一個使命.
就是要將福音講傳使萬民作主的門徒.
所以這個就是上帝藉著我們個人生命的改變逆轉.
就是要我們一起去為上帝去顛覆這個世界.
今天這個世界屬世的思維.
要大家去迴轉回到屬基督.
大家明白了嗎?.
我們是part of上帝的用人.
我們是他的僕人 我們要這樣做.
瑪利亞很清楚知道她的逆轉不是這樣改變她的生命.
哈拿也是一樣.
如果不是她就不用將兒子奉獻給聖殿.
給上帝 是不是?.
她生了兒子就可以了.
功成身退 搞定.
沒有人再歧視我.
但不是 這個只不過是開始.
這個只不過是開始.
瑪利亞生了耶穌基督 只不過是一個開始.
耶穌後來的作為 這個才是重點.
但他們因為他們個人的逆轉.
是令到上帝的角度來一個大逆轉.
今天上帝也是和我們一樣講一句話.
就不是說信了主 我躲在家裡就行了.

$^{961}$回一下崇拜 守一下聖餐.
團啟 很多八卦是非 不要來了.
查經 很悶的.
祈禱會 我怕出聲.
很多時候我們就怕這些 怕那些.
他們說不是啊 我逆轉你的生命.
就是要你為我做更大的事.
我不懂的 我不行的.
沒有啊 哈拿也不行的.
瑪利亞也不是厲害的.
但上帝改變了我們的生命.
今天上帝改變了生命.
上帝選人不是選我們有多厲害.
而是要我們看到上帝有多厲害.
不是說我們有多厲害.
而是要我們看到上帝有多厲害.
要我們親眼看到上帝的榮耀.
這就是上帝和我們講的說話.
大家低頭.
看到哈拿和瑪利亞雖然是不同時代.
但他們是同一個上帝.
就是創造天地自有永久的耶和華.
他們很清楚神會抬舉謙卑人.
也會降卑驕傲人.
生命的逆轉不是環境的改善.
而是一切由心裡面開始.
一切是由仰望神 從禱告那裡開始.
神看到我們的卑微 看到我們今天的軟弱.
神聽到我們的呼求.
在合乎上帝的旨意之下.
合乎上帝的時間表.
我們會親自介入.
是會翻轉我們的生命.
翻轉我們的家庭.
甚至是翻轉這個世代.
今天這個世代是一個沉淪的世代.
負責要我們去翻轉這個世代.
結合上帝的國.
直至主耶穌再回來的日子.
讓我們過去幾年來到英國的地方.

$^{1001}$來到瑞典的地方有一個橫的教會.
負責這個翻轉的開始.
負責要我們為你做更大的事.
要我們走得更加遠 走得更加深.
負責求你讓我們有信服謙卑的心來到你面前.
讓我們知道我們一定要依靠你.
憑著一個謙卑的心來到你面前.
就好像昔日瑪利亞和哈拿一樣.
我們有困難的時候.
我們要和上帝去哀求.
負責你必定會聽到我們的禱告.
求你帶領我們在未來的日子當中.
我們真的能夠改變我們的思維.
改變我們對人生的看法.
讓我們更加像耶穌基督.
我叫我們能夠每一天都有一個平安喜樂陪伴我們.
我們祈禱奉耶穌的名祈求 阿們.
\newpage



\section{歌羅西書 3:12-17}
\label{sec:032133}
\textbf{信徒,你們在攪甚麼?}~~~~黃紹權牧師
\newline
\newline
連結: \href{http://www.church.com.hk/acms/content.asp?site=cdc\&op=show\&type=product\&code=032133\&layout=sermon}{\texttt{032133}} ~~~~ 語音日期: 2026-05-17
\newline
\newline
\hyperref[sec:032132]{\small{< < < PREV SERMON < < <}}
~
\hyperref[sec:index_chronic]{\small{[返順時目]}}
~
\hyperref[sec:032134]{\small{> > > NEXT SERMON > > >}}
\newline
\newline
歌羅西書 3:12-17
\newline
\begin{longtable}{cl}
\hline
\hline
章節 & 經文 (和合本修訂版)\\
\hline
3:12 & \begin{tabularx}{0.7\textwidth}{X} 所以，你們既是神的選民，聖潔、蒙愛的人，要穿上憐憫、恩慈、謙虛、溫柔和忍耐。 \end{tabularx} \\ \\ \relax
3:13 & \begin{tabularx}{0.7\textwidth}{X} 倘若這人與那人有嫌隙，總要彼此容忍，彼此饒恕；主怎樣饒恕了你們，你們也要怎樣饒恕人。 \end{tabularx} \\ \\ \relax
3:14 & \begin{tabularx}{0.7\textwidth}{X} 除此以外，還要穿上愛心，因為愛是貫通全德的。 \end{tabularx} \\ \\ \relax
3:15 & \begin{tabularx}{0.7\textwidth}{X} 你們要讓基督所賜的和平在你們心裡作主，也為此蒙召，歸為一體。你們還要存感謝的心。 \end{tabularx} \\ \\ \relax
3:16 & \begin{tabularx}{0.7\textwidth}{X} 當用各樣的智慧，把基督的道豐豐富富的存在心裡，用詩篇、讚美詩、靈歌，彼此教導，互相勸戒，以感恩的心歌頌神。 \end{tabularx} \\ \\ \relax
3:17 & \begin{tabularx}{0.7\textwidth}{X} 你們無論做甚麼，或說話或行事，都要奉主耶穌的名，藉著他感謝父神。 \end{tabularx} \\ \\
[1ex]
\hline
\hline
\end{longtable}
$^{1}$同心圖各.
父神多謝你.
多謝你將你的寶貴兒子耶穌基督.
以一個祂體貼你的心思意念來到這個世界.
滿足父神你拯救這個世界眾多的罪人.
父神多謝你.
因為你能夠將我們可以再次納在你的身旁.
依傍著你在人生裡繼續走向永恆.
求你的引導.
求你的安慰觸摸我們每一個屬你的信徒的生命.
讓我們能夠在你的當中.
無論或順或逆.
我們都曉得.
歐歌頌讚.
或者能夠在主你裡緊緊的依傍.
多謝你.
你那份對我們的接納.
當我們走歪路的時候.
陷在對罪當中的時候.
你沒有離棄我們.
求你讓每一個在今天再次去接觸主你啟示的人.
心靈都受到上主你的說話.
以及你聖靈的感動.
好使我們能夠在你的當中.
重新得力.
多謝你.
禱告奉靠主耶穌基督得勝的名字而求.
阿們.
其實包括教牧和公.
我們知不知道教會裡面.
我們的靈性根本裡面發生了甚麼問題呢?.
事實上很多時候.
當我們仔細去留心看一下.
今天的信徒.
我所指的是真正追求上主的信徒.
他們對於在這條屬靈的天路所行走的時候.
其實那個方式和那個心智.
是發生了一個很根本上的改變.
而我相信這個改變.
是他們再次尋索聖經裡面所指的真正的教會的意義.

$^{41}$而作出的一連串的改變.
只不過是今天很多時候.
在辦教會的人.
他只是仍然困在他自己的一套.
唯一以為是對的方式.
以我所對教會.
尤其是信徒.
他們近年的發展.
我會看到他們的發展是向著一個.
正確而且更符合聖經的啟示的一個.
屬靈成長的路向而走.
因此今天.
當我們回來看到歌羅西書.
從第三章第十二節至十七節裡面.
我們會思想一個題目.
就是叫做.
信徒呀!你們究竟在搞些甚麼呢?.
這個信徒包括了教會的長輩.
他們究竟知不知道.
他們現在在教會裡面.
究竟在發生甚麼情況.
而自己的靈性狀態又是怎樣呢?.
是否仍然是困在傳統教會的模式.
而認為其他的都是不合法的呢?.
這個就是今天這段聖經.
讓我們一起思想的題目.
讓我們作一個合神心意.
或者配稱為神子民的屬靈生命.
能夠反映出或者折射出.
上帝的榮耀.
我們一起來看看這段聖經說些甚麼.
第三章第十二節至到第十七節.
如果大家剛才已經打開了聖經.
聽Sandy幫我們讀了一次.
不過人家讀始終可以雲遊障礙.
你自己讀我們自己讀.
就能夠將聖經的說話讀進心版上.
讓我們一起再次打開.
歌羅西書第三章十二至十七節.
我們一起去讀一次好不好?.

$^{81}$好,找到之後我們一起讀了.
所以你們既是神的選民.
聖潔蒙愛的人.
就要全憐憫.
恩慈.
謙虛.
溫柔.
忍耐的心.
倘若這人與那人有嫌隙.
總要彼此包容.
彼此饒恕.
主怎樣饒恕了你們.
你們也要怎樣饒恕人.
在這一切之外.
要存著愛心.
愛心就是聯絡全得的.
又要叫基督的平安在你們心裡作主.
你們也為此蒙召.
歸為一體.
當用各樣的智慧.
把基督的道理.
豐豐富富地存在心裡.
用詩章.
頌詞.
靈歌.
彼此教導.
互相勸誡.
心被感根.
歌頌神.
無論做什麼.
或說話.
或行事.
都要奉主耶穌的名.
上一段的經文.
我們所看到的就是談論到.
彰顯神的榮耀.
不是只是純粹一種心靈改變.
更加不是只是流於口舌上的宣告.
而是一個人整全的改變.
讓我們可以脫離.

$^{121}$一種膚淺的以為榮耀神的做法.
因為今天教會很多時候說到榮耀神.
就是做聖公.
做多些在聖經裡.
要我們教會建立的事情.
甚至建一個教堂.
都是榮耀神的.
但很多時候我們只是流於這種.
很膚淺的一種信仰.
而缺乏了真正榮耀神.
是一個人的生命裡裡外外.
在他每一時每一刻.
在生活每一個環節裡面.
每一舉手每一投足.
都能夠散發出基督的生命.
因此在上文我們說到的.
當我們在這個世界活著的時候.
我們會遇上很多的挑戰.
因為我們堅持行神的說話.
我們放棄這個世界的價值觀.
人就會拒絕甚至會來攻擊.
我們這些的取態.
但是在這個時候.
我們的上帝透過主耶穌基督.
來向我們作出掩護.
保護我們.
在這個屬靈的.
你死我活的戰場上.
讓我們能夠真真正正做到.
可以實踐上帝的說話.
從而顯明上帝.
這個是上文我們所看的內容.
你們經常都會應該有個疑問.
牧師你每次講到.
上個禮拜講的那段聖經.
只用了幾分鐘就講了.
上個禮拜一個多小時的講道.
你會不會覺得.
牧師你是不是講多了.
或者講少了.

$^{161}$不是.
我只是想給大家看到.
怎樣去將一段聖經.
很詳細地去.
將自己浸淫在裡面.
去給聖經的教導去學習.
但是亦都怎樣很有效率地.
去將它濃縮成為一個.
很重要的主題.
成為我們每個禮拜.
每一日.
作為基督徒.
走面前的路.
順帶在這裡講的.
我記得上個禮拜.
完了崇拜.
Ellen問我一條問題.
牧師你講的話.
其實當時的信徒明不明白.
是不是他們真的這樣明白呢.
這一點我在這裡順帶介紹.
一個很重要的.
看聖經或者對聖經的一個看法.
給大家知道.
這叫做聖經神學.
什麼叫做聖經神學.
聖經神學不是一個.
宗教典籍.
只是講給我聽.
然後行為手律法的做法.
聖經神學是將聖經的本體.
視為一個活體.
你說是嗎.
可以這樣的嗎.
可以.
為什麼不可以.
聖經講給我們知道.
約翰福音第一章.
第一節就講了給我們聽.
上帝起初是做什麼的.

$^{201}$道與神同在.
道就是神.
道是神.
在原文聖經講.
就是神的啟示.
Logos.
聖經就是耶穌.
耶穌就將他活活地去.
在你和我現在這個時代裡.
按照我們所能夠.
需要知的資料.
或者啟示.
用同一套說話.
來將他的主旨告訴他.
每個人在不同的年代裡.
包括在初期教會.
甚至舊約聖經的時候.
以至於今時今日.
同一段的聖經.
都在不斷地.
向我們講同一番說話.
但卻是將深度和闊度.
不同的情況.
按照你和我的情況.
介紹了出來.
我為什麼要提出這一點.
給大家知道呢.
因為這是對聖經的一個尊重.
也可以怎樣讓聖經.
能夠在我們的.
屬靈生命上生影響.
但更重要的就是.
在這個星期裡.
我曾經收到.
在網上的一個回應.
這個回應.
在上面就將.
我所講的話.
變成遍體鱗傷.
不過你們不用擔心.

$^{241}$你們會看不到的.
因為那個人.
留言之後.
我用準確.
用神的說話.
告訴他之後.
他馬上將它全部刪除.
所有刪除.
以為沒有事發生過.
不過你不用擔心.
我全部有記錄的.
信徒很多時候.
對聖經不理解.
我們作為牧者的角色.
就是要解說.
當你被解說之後.
你看到自己的缺陷的時候.
而勇敢面對去改變.
而不是將它塗抹.
塗抹,隱藏.
不是基督徒的行為.
但是我陳述到這一點.
給大家知道.
因為今天你去看.
這段聖經的時候.
有絕大部分的信徒.
包括牧者.
他一定會告訴你.
尤其是去到最後.
講到教會裡面.
用詩章,頌詞,靈歌.
這些就是教會要應做的事.
變成事公化.
但是他其實不明白.
這些陳述本身的意義.
是想說些什麼.
而聖經神學.
就將它去剖析出來給我們知道.
我們拭目以待.
逐步去看一看.

$^{281}$既然在這段聖經裡面.
上文講到榮耀神.
是將一個人生命上裡裡外外.
要去改變.
直到你信了主之後.
能夠有一個徹頭徹尾的改變之後.
怎樣參與其中呢?.
這個正正就是這段聖經.
要很實質地去告訴他.
怎樣去推行.
我們會分為兩個部分去看.
第一個部分.
就是從第十二節至十四節的上面.
保羅在這裡向哥羅西教會提出.
你們要穿上一件全新的衣服.
這件衣服是和你過去截然不同的.
講到衣服的時候.
今日祈禱會都講到衣服.
大家都去介紹一下.
我今日身上這件穿了多久.
聽說丹尼爾那件至少都三十年了.
根據小儀說.
我這件都不差的.
這件都二十年了.
不過都挺新的.
我就想起以前我年紀很小的時候.
每一年一到農曆年尾的時候.
我媽媽就一定會.
包括我爸爸.
不過我爸爸經常快一點.
就一起去買衣服.
過新年穿新衣服.
為什麼每逢過新年要穿新衣服呢?.
其實大家明白為什麼嗎?.
我其實不是很明白.
不過我猜測其中一個.
就是商家想做多點生意的意思.
但是後來當我去.
每一年臨過年買新衣服之後.
就發覺跟著就去做什麼呢?.

$^{321}$就去拜年.
你們有沒有拜年?.
以前年輕的時候有沒有拜年?.
你不想拜.
爸爸媽媽都夾著你去的.
我爸爸媽媽都是這樣.
由年初一去到年初四至年初五.
包括有些日子他們去跟家人娛樂.
去到不同的地方.
我們其實來來去去.
都是穿著那套新衣服的.
跟著就跟親戚說同一番說話.
恭喜發財.
新年快樂.
出入平安.
那些說話.
我就發覺每一次.
當我去跟他們一起去到拜年.
拜完年說完一番說來說去.
都是那一堆的說話之後.
親戚除了給利是我們.
他多數都會多說一句這樣的說話.
啊 長大了.
如果是女孩子.
長大了.
祝你快高長大.
終於我明白.
為什麼要每一年買新衣服了.
因為真的小時候.
年年長高年年變.
你不買新衣服.
下年就穿不了.
其實全年最出得大場面.
就是那套衣服了.
即是你那一年內.
無論去喝酒去哪裡.
都是穿那一套的了.
因為最散鏡就是那一套.
但是.
年年換.

$^{361}$因為年年身體都長高長大.
有變的.
男孩子多數是長高.
女孩子身形越來越標緻.
所以那些衣服年年都要換的.
人在屬靈的生命上.
其實當你信了主之後.
都會有不斷的改變.
而你這個改變.
是需要有穿上一件新的衣裳.
所以在這段聖經裡面.
從第十二節至第十四節裡面.
保羅就在這裡向.
哥羅西教會.
他作為一個屬靈的長輩.
一個過來人的身份.
他向這班的信徒.
信了主的哥羅西教會的信徒.
告訴他們你們往後的屬靈生命.
有一個很重要的大原則.
你們就要怎樣呢?.
要穿上一件新的衣服.
不過這件衣服是有什麼特色呢?.
從這段聖經裡面.
首先讓我們看見.
這件衣服不是每個人都有份穿.
每個人都可以穿的.
第十二節.
所以你們既是神的選民.
聖潔蒙愛的人.
這件衣服不是每個人都有份穿的.
好像我過年.
我爸我媽去買件新衣服給我.
為什麼他要買新衣服給我?.
又不買給隔壁樓下.
二樓的那兩兄弟?.
因為他不是他的兒子.
他們不是他的兒子.
我是我爸爸媽媽的兒子.
所以他會為我的新一年.

$^{401}$去和我一起拜年的時候.
他會為我去裝潢.
為我去添新衣服.
和他一起在一種光榮的裡面.
去顯露在眾親戚面前.
上帝同樣是這樣.
他同樣將一個他的選民蒙聖潔蒙愛的人.
他才選他可以穿上這件新的衣服.
所以換言之.
這是一個很重要的神學.
什麼人的生命會改變呢?.
特別是屬靈的生命會改變呢?.
是上帝選的人.
這個正正是歌羅西書在這裡.
保羅首次清楚地去指出.
哥羅西教會與及保羅他自己都是.
是神的選民.
大家沒有聽錯.
傳統上來計選民這個字眼.
只是用作在猶太人的身上.
但哥羅西教會是一間什麼教會?.
還記不記得?.
哥羅西教會是一間充滿外邦人的教會.
外邦人成為選民這個身份.
是可以說是一個史無前例的介紹.
等同於你和我.
我們不是猶太人.
我們是什麼人?.
我們是華人.
我們不是在中東由猶太人家族生出來.
沒有血緣關係.
當然有些人.
有些基督徒.
出盡辦法去拉關係.
又說我們中國人華人.
是從閃含亞佛到亞洲開枝散葉而生.
嘗試和猶太人拉上關係.
沒有需要這樣說.
因為這是上帝對人的選擇.
這個陳述就是在聖經裡經常提到的.

$^{441}$眼中的同仁.
就是這個意思.
如果大家有機會去打開.
如果我沒有記錯.
應該是詩篇第十七篇第八節.
或者箴言的第七篇.
裡面都談論到有關於.
上帝的選民是他眼中的同仁.
什麼叫眼中的同仁?.
你用中文好說.
情人眼裡出西施.
個個都不選擇看中你.
因為你是被上帝看中過一個.
於是乎你就會有一個生命上.
和上帝結下一個好不同的關係.
而這個關係就引發到.
你有以下連串的改變.
你會怎樣?.
你會傳以下這七個內容.
如果我沒有記錯.
七個的內容.
包括什麼?.
我讀一下.
就要傳什麼?.
憐憫.
恩慈.
謙虛.
溫柔.
忍耐的心.
未完.
倘若這人與人有嫌隙.
總要彼此包容.
彼此饒恕.
主怎樣饒恕了你們.
你們也要怎樣饒恕人.
總共七樣東西.
有沒有書哪七樣?.
憐憫一樣.
恩慈第二.
第三是謙虛.

$^{481}$溫柔是四.
忍耐的心是五.
另外哪兩樣東西?.
原來還有兩樣東西.
是包容和饒恕.
總共七樣東西.
不過這七樣東西的陳述.
又有它自己本身的特色.
這七樣東西是什麼呢?.
似乎都是說到一個人的心的狀態.
於是乎很多信徒甚至牧者去解釋.
這段聖經就說到.
其實信主都是講心的.
我經常聽見有些人.
尤其是未信主的人.
他這樣說.
基督徒都是心好.
就會覺得這句說話有點貶義.
即是你的東西全部收在心裡.
外表是怎樣呢?.
可能是另一回事.
事實上很多基督徒都是這樣的.
心裡相信.
不過沒有了口裡承認的行動.
如果大家還記得前文.
剛才講榮耀神不是只是講內心的.
對上帝尊重這麼簡單.
而是有行動配合的.
是參與一場你死我活的.
見證上帝的說話的戰爭.
是有行動的.
在這裡第一樣東西.
不只是講心.
而是講到有所行動.
特登上帝體重.
特登揀他的人.
他們的生命.
就有以下串上這一連串的改變.
除了憐憫恩慈.
謙虛溫柔忍耐的心.

$^{521}$還有一連串.
人與人之間的關係建立.
這個關係的建立.
負面的是包容.
為什麼我叫負面.
什麼時候你要去包容別人.
這不就講了.
這人與那人有嫌隙的時候.
就是說你去給人欺負的時候.
廣東話比較合理.
普通話怎麼講.
欺負你.
輪不到你選擇的.
負面的.
你只能夠作回應.
這個負面負面的意思不代表不好.
負面的意思就是輪不到你主動的.
你是被動者.
你的被動只能夠怎樣.
包容他.
牧師可不可以解釋一下.
包容是什麼.
如果要講.
我就要你給我十天八天去逐樣講.
要講很多例子.
將那個字再去解策.
簡單點講.
包容就是你沒有錯.
人家是錯誤地誤會你.
好像我剛剛剛才講.
網上的評語.
我猜他是女士.
不過不要緊.
男女不要緊.
其實我很接受到大家的批評.
因為你越多批評.
就可以幫助到我.
究竟我自己的靈性了解性是否有錯.
而且.
也有一個很好的機會.

$^{561}$我去解釋解說.
我記得那個評語裡面這樣講.
黃牧師你已經講到走火入魔.
我心想我成魔了.
真是厲害了.
沒有人敢提出意見給你去更改.
現在我就提出了.
你用了一套系統神學來解釋聖經.
錯誤.
首先我很清楚告訴大家.
我不會用系統神學解經的.
因為我心知系統神學有什麼問題.
系統神學是一個人的哲學.
誰告訴你他專研系統神學.
我可以告訴你.
他說的話你先打50\%折扣就可以了.
因為有一半可能是神說話.
但另一半多數是人說話.
我是沒有錯.
我解釋給你聽.
我寫了一段都不短的解釋.
還有引述聖經告訴你.
在哪裡可以幫助他去想.
他將它刪除了.
連我的回應都刪除了.
頂尖門.
我有什麼可以做得到.
我沒有什麼可以做得到.
不過我可以做到的其中一樣東西.
是從我的角度.
接受包容有一個這樣的信徒.
他是不了解聖經.
我為他禱告.
希望他能夠在靈性上更加成長.
這個就是我們作為牧者應有的思維.
保羅在這裡也鼓勵哥羅西教會.
是的,有時候有人誤會你.
誤會你有很多原因.
最主要的一個原因可能是他不了解.
或者他聽不清楚.

$^{601}$大家不要以為大家年紀還很年輕.
聽得很清楚.
有時候我也聽不清楚.
我已經叫出名耳利.
我是玩音樂的人.
耳聽是有些sensitivity高一點的.
但是頂尖門.
有時候很多原因.
人家會誤會你.
會甚至得罪你.
在他不明白的時候.
可能說了一些事情傷害你.
甚至可能他在錯誤的神學建立裡面.
他以為他自己對.
而去對一些在聖經有啟示的事情.
一知半解.
這些是很普遍的情況.
在教會裡面,我自己做了基督徒30年.
我很清楚知道在教會裡面.
很多人只是信人的說話的權威.
而沒有去看聖經裡面所說的權威.
當他誤會你的時候.
你不是跟他去喊殺.
你不是跟他去撐.
你是解釋給他知道.
但是在未解釋之前.
你首先要做到的一件事,包容.
所以在這條聖經說.
這人與人有嚴規的時候.
你總要包容.
這個是從負面或者我們叫做.
被動的角度來去看.
但是同時間我們要有一個積極的行動配搭.
就是要彼此饒恕.
彼此饒恕這個說話.
就是一個主動的行為.
你有選擇的.
怎樣饒恕?.
比如我講一個簡單的例子.
曾經在Alberta有一個年青人.

$^{641}$無端端被一個人去殺死了.
不是很久以前的事情.
兩三個月前的事.
被殺的家庭.
原來是一個基督徒.
他有理沒理第一句就出來說.
我要寬恕你了.
弟兄姊妹.
基督徒要寬恕人是事實.
但是不是盲目地去做的.
聖經又清清楚楚在這裡說.
怎樣去饒恕那個人.
請讀清楚這句說話.
第十三節尾.
要彼此饒恕主.
怎樣饒恕你們.
你們又要怎樣饒恕那些人.
我們去寬恕其他的人.
其實動機是想那個人回歸上主.
去饒恕我們在罪裡面的目的.
如果我們不是在這個基礎上.
就隨隨便便說.
我饒恕你了.
你做了一些事傷害了我.
我饒恕你了.
但是那個人毫無悔意的話.
這個就不是饒恕.
這個就叫做縱容.
我們讀到這裡為止.
我們會看到有一連串的事情.
是要在大家彼此的狀態之下.
才可以發揮出來.
所以從這一連串的介紹裡面.
讓我們看見.
這個不只是心態的問題.
不是有心的問題.
是整個信徒在一個實際的群體裡面.
有心也有行動.
互相連接.
配合.

$^{681}$而且是一致的回應.
保羅在這裡說到.
這一切都是那些神所揀選.
蒙蓮愛的人.
才有這樣的表現.
所以一個真正的基督徒.
不會是因為人家得罪你.
你就大發裂怒.
罵到人家像一隻狗一樣.
就以為伸張正義.
不是.
在聖經這裡說.
在行動的上面.
你可能會是作出.
被動的包容.
或者主動的饒恕.
但最終的目的.
都是能夠幫助人.
能夠歸回上主饒恕人脫離罪.
這個目的.
所以如果我們去包容.
行饒恕.
以及我們去憐憫恩慈.
謙虛溫柔忍耐的心.
等等這些事情.
其實目的都是讓人.
能夠在罪的上面.
歸回上主耶穌.
希望他們罪得赦免這個目的.
因此.
當保羅說完這一番之後.
十四節.
他這樣說.
在這一切之外.
要存著愛心.
愛心就是聯絡全得的.
一般教會一讀到這裡.
又來了.
愛字上頭了.
我都說教會要講愛了.

$^{721}$神就是愛.
是不是這個意思.
我們仔細看看這節聖經說什麼.
在這一切之外.
要存著愛心.
愛心就是聯絡全得的.
首先我們清楚知道.
這裡所說到要聯絡的全得.
是指什麼.
全得就是指上面那一連串.
那七個的內容.
你們要穿上這個.
有七個特性的.
新的生命.
在你們的屬靈生命當中.
但是原來這七個屬靈的美德的聯絡.
是靠愛來將它維繫的.
或者將它連繫起來.
在聖經裡面.
愛這個字眼.
是用著好像一個接著劑的意思.
所以就用了聯絡這個字眼.
聯絡這個動詞.
原文的意思就好像一個球一樣.
英文叫Fasten.
大家不知道喜不喜歡修理東西.
所以東西爛了去修理.
其實是一個屬靈的操練.
你不要以為爛了就扔掉它.
爛了就扔掉它.
爛了的東西你肯去修理的時候.
你會學會很多東西.
尤其是修理.
是不是有很多叫做Fasten的東西.
扣.
我想我們這樣說.
Alan應該最熟悉了.
他讀Engineering.
很多種扣的.
譬如我們有時候見到一些螺絲.

$^{761}$為什麼扭一粒螺絲.
螺絲頭之前要加一塊戒指圈.
大家知不知道為什麼.
原來那個戒指圈是幫助那粒螺絲.
扭下去的時候.
鎖緊它.
並且將那個螺絲所做出來的壓力.
將它有系統平均分佈.
在要接觸的接觸面上.
這些叫做Fastener.
聖經在這裡用了這個動詞.
去形容愛.
其實就是這個Fastener.
它是一個接著的重要環節.
它將這七個屬靈的傳德.
來將它緊緊的扣在上面.
和合本來說.
在這一切之外.
這個翻譯我覺得不是很好.
因為在接下來的經文.
其實是甚至指導於這個愛.
其實一直都存留在上帝的裡面.
向著他所揀選的人.
一路施行的.
所以換句話來說.
這個愛怎樣將這些傳德緊緊的扣呢.
不是外加的事物.
而是在神裡面的人.
上帝就用這件事.
早早將這些傳德黏在一起.
能夠使每一個認識主的人.
都能夠緊緊被這七個的傳德.
焊死了.
或者緊緊連接在他們生命上面.
這幅是什麼圖畫呢.
一個屬靈的生命.
這幅圖畫是什麼呢.
如果大家在加拿大活得久一點.
通常你都會去那些民族村.
即是那些什麼Creek.

$^{801}$什麼Village去探訪.
你們很多時候會見到.
在加拿大的一些民居裡面.
在傳統上面.
他們很多東西都自己做的.
其中一樣東西就是他們的被鋪.
他們的被鋪是怎樣做的.
是不是好像現在那樣.
整幅用機器車它織出來的.
不是的.
原來是用很多不同的布碎.
一片一片剪起了它.
一塊塊方格形.
然後用一針一線將它連起來.
如果你今天要去買一幅這些被鋪.
是非常之貴的.
我記得有一次.
我們去多倫多的皇家博物館.
Rome.
如果我沒有記錯應該是四樓.
那裡有一個專門介紹民族服裝與生活的環節.
來的了.
可以登記免費進去.
大家有機會去一下Rome.
叫Open Door.
可以去看一下.
很多這些布.
在這些被鋪上面.
都往往會貼上一個介紹.
是什麼家族裡面去編織給他的兒子.
或者他的女兒.
或者他的家人.
而每一幅的布雖然是不同樣.
好像這裡七樣的美德.
每一幅都不同.
但是能夠將它連繫起來的.
就是那個織布的人對他家人的愛心.
他一針一線將每一幅連接起來.
成為一張被子.
給他的家人享用.

$^{841}$上帝就是這樣.
為我們每一個屬祂的子民.
一針一線.
只不過是用愛這個題目.
來將這些東西連繫起來.
加在我們的身上.
為什麼上帝要這樣做.
因為你和我是上主的心上人.
是祂眼中的同人.
是祂特意揀選.
蒙聖潔去愛的人.
因此我們在這個過程裡面.
我們就在他的聖潔裡面去被塑造起來.
而這一連串七樣的內容.
其實就是幫助教會.
哥羅西教會以至你和我.
來到去成為聖潔的子民.
因此屬神的子民.
不單止能夠得到上帝的憐愛.
或者蒙祂的寵愛.
蒙祂的成為心上人這份心意.
而我們作為神的屬民.
也都很願意.
心甘情願地將我們的芳心.
委託在上帝.
藉著耶穌基督所成就的這份關係.
於是乎我們在我們的社群的互動裡面.
就能夠顯出包容.
顯出饒恕.
看到嗎?.
這種的互動關係.
現在在教會裡面在運作中.
為什麼要這樣運作呢?.
因為我們在永恆裡面.
和上帝同樣是這樣運作.
什麼運作?.
和上帝交心的運作.
因為我們是上帝很痛愛揀選.
我們應該的回應是.
同樣對上帝是毫無保留的付出.

$^{881}$而這一點也都驅使我們.
怎樣在我們實際榮耀神這件事上面去參與.
這是上帝的心願.
人要脫離罪是上帝的心願.
耶穌基督首先做了一個榜樣給我們知道.
現在活在這個地上.
教會的角色就是你和我去幫助人脫離罪.
滿足上帝的心願.
這就是對上帝有愛的人表現.
真正的教會就是這樣.
而不是由一連串的活動.
去包裝組織起來.
所以從這段聖經裡面.
我們可以簡單的說.
一個很重要的訊息.
因為上帝揀選的人.
其實他是願意用他人生往後的日子.
互動地和上帝互相教同.
大家都知道我剛剛嫁女兒.
嫁完女兒,女兒搬了出去和女婿一起住.
女婿這個星期病了.
我女兒馬上看到老公病了之後.
想著煮些什麼給他吃.
我知道他女婿病了.
我作為外父也會關心的.
有什麼可以幫助他.
我女兒說我想煮些什麼給他吃.
她打了這些短信出來.
我都不太懂得看.
Do you have P.Dan?.
P.Dan?.
Where is P.Dan?.
我不知道什麼叫P.Dan.
P.E和D.A.N.
我又沒有一個人叫P.
D.A.N.
又是皮蛋.
她想皮蛋瘦肉粥.
不過她沒有瘦肉.
她有魚肉.

$^{921}$她說我有魚肉.
你明不明白.
一個人和另一個人的關係親近.
你想體貼他的需要.
你就會想盡他最需要的去付出.
當然她問正一個對的人.
因為她爸爸常備這些東西.
她當然知道爸爸有.
所以她都問.
我們和上帝的關係是不是這樣.
是不是這樣.
上帝想的東西.
我們有沒有去幫.
竭力去滿足上帝想做的東西.
所以你看看兩夫婦.
她究竟恩不恩愛很簡單.
不是她常常有浪漫的晚餐.
燭光晚餐.
你又笑了.
你又笑什麼.
你沒有的.
是不是笑呢.
不是的.
你看看她兩夫婦之間.
她留不留意到她的需要.
你就會知道他們是不是恩愛.
所以師母常常都說.
你有沒有常常看我.
我的褲子有顆白色.
她會說你的褲子有白色一點.
你知道我對女性的意識.
我不是很敢碰.
又不敢留心.
見到女士.
我的眼睛都要看遠一點.
離開她.
我很少看她.
但是現在學了.
她有什麼需要.
你要留意她的需要.

$^{961}$譬如她說什麼.
你常常掛在心裡幫她做.
我不告訴你.
我幫師母做了什麼.
你有沒有想過.
和上帝喜歡什麼.
你去幫上帝.
去和上帝做.
未必做得到.
因為上帝的要求.
我們不能夠去赦免人的罪.
但是我可以去幫助.
將上帝赦免人的罪.
這個心意去介紹給人知道.
上帝欣賞這些.
因為你愛她.
她又愛你.
當這個正確的觀念定準了之後.
當我們穿上這件全新的.
淑齡的衣服之後.
是什麼帶動到我們去參與.
在這一連串的.
榮耀神的活動裡面呢.
十五至十七節.
我們一起讀.
又要叫基督的平安.
在你們心裡作主.
你們也為此夢照.
歸為一體.
當用各樣的智慧.
把基督的道理.
豐豐富富地傳在心裡.
用詩章.
頌詞.
靈歌.
彼此教導.
互相勸勉.
心被恩感.
歌頌神.
無論做什麼.

$^{1001}$或說話.
或行事.
都要奉主耶穌的名.
在這個部分.
我們看到一個很重要的訊息.
就是以感恩的心.
來帶動我們全人的生命的發展.
感恩上帝給我們這一連串的事情.
而我們也在感恩裡面.
去推動我們去實踐上帝的心思意念.
這段聖經很多的教會.
都用來講關於教會崇拜學.
我曾經聽過一個牧者.
他在香港專門搞崇拜學.
性惡.
他很喜歡引用這段聖經.
講什麼道他都講這段聖經.
因為他搞性惡.
所以這段聖經是講詩章.
要唱詩的.
頌詞.
有時唱歌.
有些獨白.
靈歌.
他們很深入研究這幾個字眼.
歌頌神.
全部都講崇拜學.
性惡的主題.
是不是講這些呢?.
我們仔細看看.
和剛剛我所講的這段聖經.
其實是講.
用感恩去帶動我們去傳言.
活出一個合神生命的主子.
是截然不同的.
當我們去看看這段聖經的時候.
我們會發覺到.
保羅在這裡提出了.
兩項很重要的勸喻.
其實這兩項的勸喻和上文.

$^{1041}$一會兒我們所講到.
有關於包容和饒恕.
是好像一個倒影的結構.
首先我們看看第十五節.
第十五節這裡說.
又要叫基督的平安.
在你們心裡作主.
這個所謂作主.
這個說法.
在原文聖經裡面就是說.
管理的意思.
要由基督的平安.
來去管理我們的心.
保羅在這裡說.
我們要以基督.
來去管理基督的哪一方面.
原來有清楚界定.
基督的平安.
來去管理.
很多時候我們覺得心煩意亂.
其實主要原因.
因為我們用了.
我們自己認為平安的東西.
取代了基督的平安.
說得白一點的.
我們就以自己個人的滿足.
來去斷定一切的事情.
取代了以基督的平安.
或者和平來去管理我們一切的事情.
當人將人成為一切事情的核心.
因為我們是一個有限的人.
當我們有限的時候.
換句話來說.
有很多事情是做不到的.
當有些事情是做不到的時候.
我們就很容易會有失落感.
你知不知道為什麼有些人.
經常在辦公室罵人.
特別是老闆.
不過現在我發覺.

$^{1081}$罵人的不是老闆.
老闆反而不敢罵人.
是下屬罵老闆反而多了.
因為他們覺得自己很重要.
我覺得.
俗套一點就是.
他們失去了賓主之間的關係的定調.
其實這個情況尤其嚴重.
就是這些所謂左傾的思想.
爭取權益.
你知道什麼叫左傾嗎?.
左傾思想就是說.
要下面的都定為弱勢一群.
Employee 僱工就是弱勢.
所以要向僱主來爭取.
左傾思想就是在這個基礎上去建立.
什麼都反.
大眾權益就是王道.
大眾認為對的東西是否一定對呢?.
我們加拿大政府的Majority Government.
你準備圖定一個trillion.
供我們加拿大政府用.
他說你們大家.
他會用一個trillion的錢.
來幫你減電費.
你信不信呢?.
錢還那裡來?.
左傾思想最懂得用這種掩眼法.
使你支持他.
當你以人為管理核心的時候.
失去神來管理的時候就出問題.
你會覺得很惆悵.
因為你達不到自己的目的.
有些事想得到但得不到.
有些事想達成但能力又不能及.
左傾思想就會發難.
但上帝在這裡怎麼說呢?.
作為基督的群體.
他們要讓基督的平安為我們作主.
而且說道你們有一次為此蒙召歸為一體.

$^{1121}$這一段聖經.
看起來就好像.
很踐踏我們個人.
怎麼踐踏呢?.
沒有什麼主導的.
因為你只能夠接受.
好像上面頒語給你.
但是看清楚.
聖經在這裡說道.
上帝是將你和我蒙召在一個這樣的狀態裡.
歸為一體.
使我們去學會欣賞上帝的創造的公認.
當你懂得欣賞的時候.
上帝多謝你.
這不是高高在上的神給我們的.
而是高高在上的神.
祂很體貼我們的需要.
而去按照我們的需要去公認.
而我們接受公認的人.
曉得回應.
這個所謂感恩的神學.
其實早在《創世記》裡面已經介紹了出來.
是嗎?.
上帝創造世界的時候.
有沒有狂修過問過你想不想?.
沒有.
全部上帝決定.
專權啊.
不是專權.
祂考慮創造的時候.
是考慮人每個人的需要去創造.
上帝將最好的.
去設計出來.
給人擺去使用.
享受.
分享神的榮耀.
上帝給人的東西.
是讓你懂得去欣賞感恩.
只不過人.
抗拒上帝.

$^{1161}$怎麼抗拒?.
就是被蛇線動一下.
就挑出他要抗拒上帝.
現在不要上帝你做事給我.
我做事.
我給我自己.
我要做回神.
所以這個就是《創世記》裡面.
所談論到賜予和感恩的關係.
當人一失去了這個平衡的時候.
就落到這段聖經裡面.
所談論到你失去了基督的平安.
來去管理.
所以在聖經裡面.
當我們留意到這一節的經文第十五節.
談論到你們為此蒙召歸為一體.
其實和在以弗所書的第四章十六節裡面.
有一個異曲同工的陳述.
在裡面這樣說.
全身都靠他聯絡得合適.
百節各按各式.
照著各體的功用彼此相助.
便叫身體漸漸增長在愛中建立自己.
這個陳述.
正正就是告訴我們知道.
上帝是要讓這個身體.
這個上帝的屬靈群體.
建立是合一.
而讓你在這個建立過程.
帶著感恩的心去回應.
所以保羅在這裡再次邀請.
哥羅西教會的人.
繼續去對上帝這條屬靈的道路.
怎樣去榮耀他而去努力.
第二個字眼.
在第十六節.
十六節這樣說.
當用各樣的智慧把基督的道理.
豐豐富富存在心裡.
存在心裡.

$^{1201}$這個存在這個字.
在原文就是居在其中.
英文叫Dwell.
同樣與前面講到那個作主.
都是以一個勸喻的語調.
去向哥羅西教會陳述.
告訴哥羅西教會.
要有基督的教訓.
充充足足.
豐豐富富這樣住在你們裡面.
為什麼要住在你們裡面.
因為神的教訓.
就成為一個最全面的智慧.
存在在你們的當中.
十六節所講.
用各樣的智慧就是指這樣東西.
各樣的智慧不是人的智慧.
而是指到.
以基於神的說話啟示.
來作基本的一切的思想.
取捨.
教導.
而這些的教導.
從聖經裡面告訴他.
他用了各式各樣的媒介.
來讓你和我學習和回應.
哪幾樣東西.
詩章,頌詞,靈歌.
在過程當中.
彼此教導.
互相勸勉.
心,鼻,恩,感.
恩感這麼熟悉.
這個字眼.
同一個字的字根.
聖經在這裡告訴他.
原來保羅在這裡講出.
教會用各樣有的活動.
不同的文體.
不同的方式.

$^{1241}$來促進我們對基督的道理.
能夠彼此教導.
又出彼此這個字眼.
馬上就回應在上文.
談論到我們在人與人之間.
有嫌隙的時候.
如何彼此包容的問題和寬恕.
兩者的角色同樣都是.
被動與及主動.
教導是一種積極面的指導.
是一個人有計劃地去策動的事情.
有系統地去宣講.
教會最重要的.
最重要的.
就是要將基督的道理.
有系統地去教導.
而在教導的過程當中.
一定是有一些問題出現的.
所以要怎樣?.
彼此互相勸誡.
所以基督徒的成長.
不是一個人在修練.
練到心很平靜.
就等同於是一個基督徒.
基督徒的生命是在一個群體裡面.
不是單打獨鬥.
是一個群體裡面.
一起彼此造就.
再談談網上那個人的回應.
弟兄姊妹.
雖然我做了包容.
做了饒恕.
並且我會作出教導.
我會將基督的道理.
清清楚楚講給他知道.
在這個過程當中.
我對於傳講.
亦有更加學會全面地做得更好的方法.
但是對方肯不肯接受呢?.
能不能夠因此而帶著恩感來回應呢?.

$^{1281}$如果一個人被另一個人指出他的錯誤.
並且告訴他正確是甚麼時候.
那個人肯接受回應的話.
我們稱這種思想為甚麼呢?.
先不要太熟悉.
在這個世界我們稱為.
「袋錢落你個袋」.
對嗎?.
別人肯告訴你你的錯.
是「袋錢落你個袋」.
在屬靈的角度是怎樣看呢?.
你當被人指出你的錯.
肯願意承認之後.
你就會「歌頌神」.
甚麼是「歌頌神」呢?.
向上帝發出感恩.
你在我的靈性上.
差一個人來提醒我.
所以我再次重複說.
一個真正的基督徒是不難辨認的.
當人指出他的錯誤的時候.
他不是「撐」而是「感恩」.
他會不會覺得難堪呢?.
會!.
如果我說出我錯.
我也會覺得尷尬.
但你要肯承認.
我記得有一次.
五,六年前.
有位姊妹.
請大家紀念她.
叫慧清姊妹.
慧清姊妹是一個追求上主的基督徒.
但她不是經常去教會.
她說她為何來我們.
以馬來橫浸信會和我們相識.
是上帝用YouTube的方法.
帶她回來.
有一次.
每次她聽完我說道之後.

$^{1321}$她都會寫一些問題給我.
有一次有一條很有趣的問題.
她說「牧師啊」.
「有一次你說道」.
「上帝在荷莉山上」.
「講道到他和」.
「以利亞和巴力先知」.
「然後去爭奪」.
「過了之後」.
「以利亞先知馬上跌倒」.
「投訴上帝」.
「現在只有我一個先知」.
「仍然肯跟隨你」.
聖經說甚麼呢?.
上帝跟她說.
我捱了多少個先知?.
多少個?.
是不是七千?.
我回答她的時候.
我答錯了五千.
她馬上回覆我.
馬上.
半夜三更.
牧師不是七千嗎?.
以一個傳道的牧者.
準確度都非常高的牧者.
給一個信徒.
點出這一點.
你怎樣回答?.
我的回答是.
對不起.
是我錯.
應該是七千.
真正的基督徒就是這樣.
我沒有寫其他東西.
不要說我答錯.
不需要再辯駁.
正確.
七千就是七千.
所以Daniel你對的.

$^{1361}$衛東你對的.
Ellen你也是對的.
七千.
牧師寫錯.
不要再辯駁.
是七千.
耶穌餵飽了多少個人吃?.
五餅二魚嗎?.
三千?.
五千?.
幾千呢?.
幾千個吧.
聽姐妹.
因為耶穌不是.
一次餵飽多少個人.
在聖經記載了有兩次.
是.
一個曉得.
在上帝面前.
認罪悔了的人.
當被人指出你的錯的時候.
有一位這樣的姐妹.
能夠這麼勇敢.
去向我講出.
糾正的說話.
請姐妹.
這個就是.
保羅耀的教會.
一個真正的教會.
是在十七節所講.
無論做什麼.
或說話或行事.
都是奉主耶穌的名.
父神.
一個基督徒能夠繼續走前面的路.
就是因為有基督耶穌的說話.
一直在指出.
我們生命.
需要進步的地方.
然後我們可以接受到.

$^{1401}$這個就是整個屬靈生命.
發展的動力.
我們怎樣去榮耀神.
最重要的目的就是要有.
神的名.
主耶穌的教訓.
準確地不會虧損上帝的名.
來走前面的路.
當我們走的時候.
無虧損上帝的說話的時候.
肯定見到上帝與我們同在.
怎樣去.
作出掩護.
我們就會帶著感恩的心.
一起走.
今天我們見到.
這段聖經從第十五節.
以至於去到.
第十七節.
以感恩開始.
以感恩作結.
我們作為基督徒.
其實整個人生.
就是一個感恩的裡面.
上帝的恩.
來一起走這條屬靈的道路.
去榮耀上主.
因此這份感恩的心.
同時也能夠.
使我們有諸於內.
仍諸於外的生命的改變.
而今天教會.
很多時候.
還在那裡呱呱嘈.
那些人不回教會.
實體崇拜.
這些人不行.
甚至引用聖經.
亂引用聖經.
說那些人不回教會.

$^{1441}$像慣常的一樣.
不回去聚會.
因為聖經說.
不可停止聚會.
像慣常的一樣.
問題是慣常.
不是不回教會.
實體崇拜.
以馬濑橫浸信會.
正正看見.
今天教會那種.
不務正業.
沒有系統地將上帝的說話.
清楚有系統地去傳講.
於是造成今天的教會.
你去細心留意.
無論他們唱詩.
不是看見歌詞.
能夠將歌詞的內容.
成為你生命.
向上帝發出的感恩.
也不能成為.
上帝向你說話的過程.
我認識一個朋友.
在香港唱福音流行曲.
他認為自己.
在事奉主 事奉得很好.
現在開了一間.
聖樂流行詩歌訓練學校.
他會說.
你看看上帝多祝福我.
現在我可以開一間學校.
一舉兩得.
上帝win 我又win.
教會有多少.
跌在這個陷阱裡.
借上帝為名.
順便賺一筆.
今天的教會很多都是這樣.
唱詩歌講有feel.

$^{1481}$但是沒有對聖經裡面.
或者詩歌裡面.
歌詞成為一個.
獻上給上帝的歌頌.
因為這裡所講的.
是一種詞 是靈歌.
是發揮 發表.
人對上帝的感恩態度.
今天的教會領袖.
經常說.
最壞就是covid.
搞到人們不回教會.
現在人們不回教會.
教會就少人.
教會少人 奉獻不夠.
教會奉獻不夠.
就倒閉.
多少教會倒閉.
弟兄姊妹.
假如真的沒有covid.
個個回教會.
那些人回教會.
是否真心.
來到上帝面前.
帶著感恩的心.
不要迷戀一個number.
人數.
而明知裡面.
那些很多人不信主的人.
依馬利華浸信會人少.
我不覺得人少.
但是依馬利.
至少在裡面的人.
是認真相信主的.
不是隨隨便便的.
他是追求的.
今時今日的世界.
不少的信徒.
不回實體教會聚會.
是因為實體教會.

$^{1521}$已經失去了.
那個真正原本有.
那個追求認識主的.
那個群體.
那個本質.
而我們教會.
依馬利華浸信會.
就在做一個見證.
我們有很多.
不回實體教會.
崇拜的弟兄姊妹.
但是他們同樣是追求主導的信徒.
他們願意.
將他們所學到的真理.
能夠在他們的生命.
生活裡面.
付諸於實行.
我們依馬利華浸信會.
也有一些雖然不是很多.
但仍然肯參加實體崇拜聚會.
同樣追求主導的信徒.
他們學了聖經裡面的真理之後.
能夠同樣.
以向上帝感恩的心.
能夠彼此.
在一個實體裡面.
鼓勵.
執行神的律例.
典章和教訓.
從我這樣的角度來看.
教會.
不是在乎那些人.
實體還是不實體.
教會.
是在乎那批.
參與者.
是否真正神所.
揀選.
蒙年愛.
會懂得穿上.

$^{1561}$這七項.
屬靈的美德.
能夠連絡起來.
以感恩.
來推動他們行人生的.
信徒.
今天的教會.
遺留下來的只是.
經常想著一個.
教會的結構.
有個教堂.
有一連串活動.
他們為了.
有的是保存教會.
傳統的所謂.
宗教觀.
但卻失去了.
真心追求.
主導的人.
那份被栽培.
被一起在感恩裡面.
行前面的屬靈的道路.
那份原本.
聖經的啟示.
因此.
我們會見到.
為什麼很多信徒.
特別是教會的.
教睦領袖.
嘩嘩嘈.
他都不知道.
原來很多.
屬神的子民.
回到聖經的基本真理.
他不想再只是為了.
維繫一小撮人.
為了要保存教會.
傳統的一種方式.
首先請大家留意.
我不是否定有教會實體存在.

$^{1601}$但那只是.
其中一種.
追求主導.
的一個.
做法.
有很深的歷史.
但那只是其中一個.
以馬來尼環陣社會.
要躺開歸回.
聖經裡面.
追求主導.
有真正生命更新的信徒.
能夠被栽培.
在上帝面前.
的一間教會.
不在乎是實體.
不在乎是否.
虛擬或者.
網絡形式.
或者還有其他的形式.
不過.
上帝讓我們在21世紀的時候.
能夠提供了.
一個這樣的空間.
讓教會可以有一個.
修正的機會.
我重複說多一次.
是一個讓教會修正.
他對於教會.
的意義的本身.
的發展.
你有實體.
繼續實體.
但你不要否定有很多的信徒.
是不在.
實體崇拜的裡面.
或者實體聚會的裡面.
不追求主.
他仍然沒有.
慣常不回教會.

$^{1641}$因為他的心仍然.
與眾信徒走在一起.
求主幫助我們.
讓我們真的成為一個.
神眼中的同仁.
知道我們自己在做什麼.
而不是只是為了.
維持一個人為的.
傳統.
的一個硬兵兵.
上帝的.
屬靈的家.
讓我們一起和心低頭禱告.
父神多謝你.
在這段聖經的啟示.
藉著.
你的僕人保羅.
對你的說話的.
透徹的掌握.
讓他向當時.
哥羅西的教會.
宣告出一個.
屬靈生命要成長的路向.
以及他的動力.
的來源.
就是透過主你的.
啟示.
為了滿足.
上帝你.
拯救人的心願.
我們甘願.
將我們生命委託在.
其中.
讓上帝你去掌管.
讓你的說話.
以及你的心思意念.
成為我們教會的宗旨.
每一個的信徒.
以感恩的心.
來體會到我們怎樣.

$^{1681}$體貼上帝你的心思.
與及同在.
為此我們獻上感恩.
讓以馬來成為.
一間上帝你所喜悅的.
教會.
是以不同形式.
將屬主.
你的揀選的.
選民焦聚在一起.
同心追求.
主你的道理.
在我們被你所安託的.
時間空間.
環境當中.
榮耀上主.
禱告奉靠.
主耶穌基督得生命至久.
阿們.
\newpage



\section{約書亞記 11:23}
\label{sec:032134}
\textbf{聖經人聖經事:靠主奪了那全地}~~~~黃克勤牧師博士
\newline
\newline
連結: \href{http://www.church.com.hk/acms/content.asp?site=cdc\&op=show\&type=product\&code=032134\&layout=sermon}{\texttt{032134}} ~~~~ 語音日期: 2026-05-16
\newline
\newline
\hyperref[sec:032133]{\small{< < < PREV SERMON < < <}}
~
\hyperref[sec:index_chronic]{\small{[返順時目]}}
~
\hyperref[sec:032135]{\small{> > > NEXT SERMON > > >}}
\newline
\newline
約書亞記 11:23
\newline
\begin{longtable}{cl}
\hline
\hline
章節 & 經文 (和合本修訂版)\\
\hline
11:23 & \begin{tabularx}{0.7\textwidth}{X} 這樣，約書亞照著耶和華所吩咐摩西的一切話奪了那全地，就按著以色列支派所得的份把地分給他們為業。於是國中太平，沒有戰爭了。 \end{tabularx} \\ \\
[1ex]
\hline
\hline
\end{longtable}
$^{1}$各位電影姐妹大家好.
眨眼又一天了.
我們一起讀聖經.
在生命樹宣教網絡為大家預備聖經人,聖經事.
在這些人物與故事的背後.
看看上帝的手有多大能和智慧.
今次仍然在說約書亞.
就是「靠主奪了那全地」.
一看就知道這句不是聖經.
「靠主」是我加進去的.
「奪了那全地」是經文.
是「依靠主奪得全個應許之地」.
跟大家看看聖經.
約書亞記第十一章.
很快的十一章.
故事很長.
夏所王耶賓聽見這事.
就打發人去馬頓王約巴.
新倫王壓殺王與北方山地.
基尼列南邊的阿拉巴高原.
並西邊多爾山江的諸王.
先看這兩節.
看這兩節你會發覺背景是什麼呢?.
是「聽見這事」.
即是說是聽見上一章.
約書亞打勝南方一群王.
由於基片人與以色列人約書亞講和.
基片是一個大城市.
那麼多勇士與以色列人結盟.
令其他王害怕.
於是組織成五個王.
一起進攻基片.
基片向約書亞求救.
約書亞帶領其人.
上帝指示他.
晚上出發.
忽然與五個王打.
終於將五個王殺死.
毀滅打退軍隊.
五個王又被殺.

$^{41}$很害怕.
不過他們都在死海附近.
在當時的耶路撒冷附近.
在耶利哥附近.
在吉甲這個地方駐守.
被南方王.
夏索王是北方王.
南邊一群王被以色列人攻佔贏了.
殺死了.
北方王害怕起來.
特別是夏索王最害怕.
可能是領袖.
耶賓.
聽到之後就去找馬頓王.
新倫王.
亞薩王這些名字真的聽都沒聽過.
北方山地.
北方山地就是以色列.
現在你看以色列國境內最北的山地.
基尼列.
大家聽聽這個名字.
基尼列就是加里尼海.
基尼列湖.
是在加里尼海南邊的阿拉巴高原.
還有西邊.
南邊又西邊又北邊.
一起的諸王要與以色列打仗.
不僅如此.
北邊南邊.
再加上東方西方的迦南人.
不管低地高地.
山地的阿摩利人.
黑人,比利,洗人,耶布斯人.
甚至去到以色列最最最北的地方.
就是黑門山.
黑門山那裡有一個叫根米斯巴.
有些希美人住在那裡.
都找了.
差不多夏所王要找到的人.
結盟的人都要一起.

$^{81}$我們去與以色列作對.
打仗.
為什麼那麼多人聽他的呢?.
很了不起嗎?.
其實不用很了不起.
因為只要是以色列了不起.
是以色列的耶和華了不起.
他們聽到耶和華神帶以色列人.
由埃及這樣出來.
真的很了不起.
埃及王法老都制他不住.
經過蘆葦海.
即是紅海過來.
經過了西乃曠野.
對巴山王那邊的西弘,鄂等皇帝又被他殺了.
然後入耶利戈.
艾城又倒塌了.
艾城又被他用妓女攻佔了.
然後五個王打他又輸了.
只要將這些事告訴他.
大家一起要合作.
如果不合作.
我們全部沒得玩.
因為都是一條死路.
一定要跟他打仗.
這個時候如果我們不能夠壓制他.
我們的地區與我們的國位都沒有.
王位都沒有.
於是大家一起來.
這裡的人真的多到數不完.
這些王和他們的眾軍都出來.
人數多如海邊的沙.
並有許多馬匹車輛.
很大陣容.
豬王會合來到米倫的水邊.
米倫距離剛才第一節下鎖王下鎖城很近.
在那裡附近.
一起安營.
今天要跟大家講的第一點.
就是這麼多人一起安營聚集.

$^{121}$數不出這麼多的兵馬.
這麼多的王要與以色列人爭戰.
再一次第一點跟大家分享的.
就是其實我們跟隨上主.
祂帶我們去哪裡.
如果我們興旺的話.
更多人會攻擊你.
如果你不興旺的話.
也會有很多人攻擊你.
因為祂一定踩扁你.
看到這麼厲害.
你會來霸佔我的地方.
我的主權我的霸權會失去.
於是祂又會攻擊你.
以色列進入新地方必定遭遇這些爭戰與攻擊.
不過是神帶領的.
你就不用擔心.
放心面對.
還記不記得主耶穌說.
這個世界會恨你們.
不過我將平安留給你們.
這個世界恨你.
為什麼?.
恨我們信主的人跟隨主的人.
為什麼?.
因為你們不屬於這個世界.
所以世界就恨你們.
有時候你會發覺無端端.
這麼多人不喜歡我們.
因為我們不屬於這個世界.
世界的王會恨我們.
很多人會無端端做些事去攻擊你.
或者不歡喜你.
或者背後刺你.
不知道.
最要知道.
是很多人要預備.
預了很多人會跟我們這些屬神的人爭戰.
留意不要說為什麼會這樣.
主耶穌一早說了.

$^{161}$準備好了.
看看地圖.
這個地圖昨天看過.
以色列人在吉甲安營.
在吉甲住的很大的地方.
平原來的.
一大河谷一個平地.
那裡住了二百萬人.
還有耶利哥這一帶.
以色列人住在那裡.
機騙跟他們結盟.
攻打了艾城.
於是耶末.
耶路撒冷,希伯倫,伊基倫,拉根王.
五個王一齊跟以色列人打仗.
進攻機騙.
結果神幫助他們打贏了.
而神親自出手.
就在冰雹很大一粒.
將很多士兵.
對方的士兵打死了.
被冰雹打死的人.
比死在以色列士兵刀下還多.
聖經這樣說.
很感恩.
上帝真的跟他們爭戰了.
看看今天這張簡單一點的地圖.
希望你能夠看得到.
或者你用手機看.
希望看得到真.
上次發生的事.
就是吉甲在這個地方.
機騙在這個地方.
耶路撒冷的王在這裡.
在這裡這幾個王一齊進攻機騙.
這是昨天的.
今天.
所以你看到是以色列地的南面.
接著這件事.
嘩!大事了.

$^{201}$就被上面.
加利利海.
剛才說的甲尼撒勒湖.
就是這個地方.
這個湖邊.
有一個城叫夏索.
就是夏索王.
就帶動了附近所有北邊.
無論約旦河.
東南西北的人.
都一齊.
嘩!數不盡這麼多的王.
就一齊進攻吉甲勒.
要想打.
誰知道.
就是輸了.
他們以色列人.
是去進攻他們的.
這件事就是這樣.
你知道昨天說的就是南邊.
今天是北邊.
今天的題目是全地.
那全地.
那全地是指不是南邊的.
是指北邊那全地.
不過北邊都拿了.
南邊又在手裡.
於是.
在這個時候開始.
說完今天的故事.
以色列人就得著整個英雄之地.
就由這個故事開始了.
得著整個英雄之地.
好.
葉國華對約書亞說.
第二點要和大家分享的.
不要因他們懼怕.
懼怕呢.
fear 恐懼.
恐懼呢.

$^{241}$是我們人類最大的敵人.
很多事情你會發覺到.
有事一發生.
有新的事情出現.
有事情改變.
我們第一個會發生的.
是恐懼的.
死了.
為什麼會這樣呢.
因為恐懼.
就是發生了一些事.
不是照著我們的計劃.
不是照著我們的想像出現.
一有新的事情出現.
於是你的心會害怕.
為什麼會這樣呢.
於是你馬上會想防備了.
如果以前有一些不好的經歷.
你馬上就會想走的了.
離開.
逃跑了.
未必敢面對它了.
不敢面對它自然就會逃跑了.
當你想得負面的時候.
你就會發現越來越害怕.
今天害怕這件事.
明天再想一想.
害怕多幾件事.
晚上睡不著.
再想一想.
還害怕多十件事.
很害怕.
如果我是約書亞.
或者我是以色列人.
聽到剛才讀都讀不完.
數都數不到.
這麼多王要預備跟我們打仗.
整個北部的所有民眾.
預備要跟我們打仗.
真的人多到.

$^{281}$海邊沙天上的星一樣多.
而且他們很富有.
有馬有車.
不夠他們來.
而且他們地頭蟲.
地形又熟.
那我們怎麼辦.
如果是這樣.
我想很多人都害怕.
如果百姓聽到真的想偷走.
如果約書亞聽到.
馬上就「哎呀 阿神 怎麼辦」.
不過你會看到.
第六節是講夜禍.
對約書亞說.
你不要因他們懼怕.
為什麼呢.
不怕就搞了.
因為留意這四個字.
所以我沒有highlight它.
明日這時.
明日這個時候.
明日就說不用久.
明日這個時候.
我必將他們交付以色列人.
將他們全然殺了.
我教你要做什麼.
你要叫所有以色列人做.
將他們得勝之後.
將他們全然殺了.
並且砍斷他們馬的蹄筋.
很有趣.
本來很多牲口拿來自己可以用.
是砍斷他們的蹄筋.
即是說這些馬不能再做戰馬了.
沒有用了.
可以tong 了.
馬最好就是要跑得快兼有力.
砍斷了蹄筋.
還有什麼用呢.

$^{321}$多數真的要死了.
用火焚燒了他們的車輛.
這麼有趣.
那麼多馬留給自己用.
那麼多車留給自己用.
我們又增加那麼多馬和車輛.
全世界都怕我們.
又可以用.
但神說我不要.
雖然看來很有用.
不過我不要.
如果要我會給你.
其實都不用的.
放心吧.
於是約書亞率領一切兵兵.
在米倫水邊突然向前攻打他們.
米倫近處下鎖.
即是說.
是否等他們下來.
不是.
大家不介意.
我給大家看一張圖.
他們在吉甲.
營地在吉甲.
神跟約書亞說.
你不用怕.
明天這個時候.
我會將他們全部北方的士兵都交給你.
你要勝過他們.
你要殺了他們.
還有你要將馬的蹄筋斬斷.
將車燒了.
於是約書亞真的很勇敢.
他又帶著以色列全部的士兵.
非常勇敢.
同心合意.
很堅定.
自己攻上去.
去到下鎖.
如果用我這支筆.

$^{361}$來指下下鎖在這裡.
如果剛才說米倫水邊.
米倫水邊大概在這個鎖字下面.
這個位置.
很近.
即是說他們自己攻上去.
這個攻其無避.
真的很厲害.
而且要留意.
在米倫水邊突然向前攻打他們.
又說將他們交在以色列人手裡.
那些下鎖的人.
所有北方的王.
以色列人就擊殺他們.
追趕他們去到西頓大城.
去到米斯利弗馬欽.
直到東邊米斯巴的平原.
將他們擊殺.
沒有留下一個.
這個是神的命令.
沒有留下一個.
入兵許之地.
你要將裡面的人殺了.
雖然今天我有機會跟大家說.
當天這個就是種族滅絕.
可以說是的.
這麼壞.
是我們今天的思想無法接受的.
有機會再跟大家解釋這個問題.
好了.
去到西頓大城.
不妨又再看回地圖.
他們在這裡上去米倫水邊.
攻打下鎖.
所有的王聚集在下鎖一起.
然後打到他們大敗而逃.
竟然有部分人逃到西頓.
走到這麼遠.
這麼北.
去到西頓這麼遠.

$^{401}$如果今天去到西頓這麼遠.
今天的以色列地圖.
西頓已經差不多去到黎巴嫩.
很遠.
他們就一直逃跑.
周圍逃跑.
以色列人追上他們.
只要吩咐他們行.
砍斷他們馬的蹄筋.
用火燒了他們的車.
約書亞就奪了那全地.
那全地不是整個英雄之地.
是指今天打仗的北方.
就是山地.
一帶的南地.
哥山全地.
高原阿拉巴以色列的山地.
和山下的高原.
由上西爾的哈拉山.
到黑門山下黎巴嫩平原的巴力加德.
並且擒獲那些地方的諸王.
將他們殺死.
那些王一打仗輸就逃跑得最快.
去到西頓跑到最北黑門山.
去到黎巴嫩平原.
不過都被以色列人追上.
你知道以色列士兵大部分都是步兵.
真的每天跑很遠.
當時約書亞就來到.
將山地希伯倫底壁阿拉巴.
猶大山地以色列山地.
所有阿納族人剪除了.
約書亞將他們和他們的城邑盡都毀滅.
由這一節第21節就留意了.
約書亞來到.
將山地希伯倫.
希伯倫是南的.
底壁是南的.
猶大山地比較南的.
其實這裡開始說.

$^{441}$北部的先是昨天打仗南部.
平定了.
贏了殺了.
然後今天說北部又贏了.
來到這裡就說全國.
北部南部都贏了.
並且將當中整個應許之地都獲得了.
而且當中有一種人叫阿納族人.
特別要提一提.
為什麼阿納族人要特別在聖經提一提阿納族人呢.
因為阿納族人.
如果照聖經這樣說.
如果是這個解釋的話.
阿納族人是比較高大.
聖經裡面叫他們做偉人.
偉人不是說很偉大.
是指很大很高大的人.
有些聖經是譯作巨人.
好像以色列和打仗的時候都發生過.
在這個非利士人裡面有哥利亞.
照看他差不多九尺多高.
很大隻的.
阿納族人就是這一類很高大的人.
非利士人裡面有阿納族人.
這樣的後裔也說不定.
但你會看到阿納族人.
就是那些很大很大隻很難對付的人.
為什麼這麼大隻呢.
有可能是像聖經另一種人的名字.
叫利弗欽人.
利弗欽就是利弗欽.
意思就是以前說看守黑門山的刺經.
看守黑門山的天使.
看到人的女子美貌.
然後化身成為人和他們交合.
於是生出很多上古英武有為的人.
就是巨人.
教他們很多新的技術.
懂得鑄金屬.
懂得木器,雕刻等等.

$^{481}$打仗等等.
那些就是利弗欽.
利弗欽就是利弗欽人.
或者叫巨人.
後來就被叫做阿納族人.
如果是這樣的話.
世界上剩餘少數很少.
神沒有讓他們生存很久.
把他們都剪除了.
即是最難對付的那些大隻的都剪除了.
以色列人的地方沒有留下一個阿納族人.
因為阿納族人如果這樣看.
是特別指那些違背上帝.
神不歡喜而由那些墮落天使.
壞了天使而生出來的人.
如果有這個解釋的話.
神不歡喜他們.
很有可能有些人認為.
其實洪水就是要毀滅地上這些人.
和當日很多罪惡的人.
這是另一個問題,有機會再說.
阿納族人他們沒有留下一個.
不過.
只在加薩,加特和阿薩特有些留了.
我不知道是不是.
加薩地帶,有沒有聽過今天在以色列加薩地帶.
即是以前的非利士地.
又和以色列不停地打仗.
可能有些大隻佬在裡面.
所以有些阿納族人可能在非利士裡面.
在這個地區加薩,加特,阿薩特.
海邊這個地方.
是非利士人住在那裡.
他經常和以色列對抗爭戰.
裡面出現了哥利亞很高大.
不知道是不是阿納族人的後代.
剩下他們.
他沒有留下這些人,其實是留下了一些的.
敵人的罪根.
那些根在這裡,他們日後再興旺和以色列人打仗的.

$^{521}$好了,23節.
這樣若書下趙爵耶和華所吩咐摩西的一切話.
奪了那傳遞.
這個那特別有個那就是指.
是神說要給以色列.
給亞伯拉罕的後裔那塊地.
所以是那傳遞.
傳遞就是整個地球.
那傳遞那塊地是奪了.
拿到.
這個奪了也不是隨便伸出手牌就拿到.
是要想辦法去爭奪回來的.
出力的.
之後就按著以色列支派的宗族.
就是他們的支派.
和他們每個支派的很多族人.
很多兄弟姐妹很多族.
將地分給他們為業.
於是留意了.
奪了那傳遞之後.
國中太平沒有爭戰了.
這裡有個字挺有趣的.
國中太平沒有爭戰.
當時都沒有以色列國.
但是說國中太平沒有爭戰.
這個是.
睡不著了,不懂得解釋.
為什麼有個國呢?.
約書亞都沒有立國.
去到塑羅斯立國.
去到薩姆爾.
塑羅斯開始立國.
約書亞記憶裡面已經說.
國中太平沒有爭戰.
這個會不會是一個很早的隱喻.
在四經裡面說將來會立國呢?.
很有可能.
不過我找不到其他的線索可以去解釋.
有時候你會看到奪了那傳遞.
神真的給他的.

$^{561}$你看到不簡單.
留意由以色列人.
叫一個摩西去帶他們從埃及出來.
接著被法老追兵.
然後經過蘆葦海,紅海.
一直進曠野沒有水沒有東西吃.
由瑪拉一直餵養他們.
然後他們在那裡得到律法.
然後擊敗沿途對抗他們的人.
然後開始入迦南.
過約旦河.
好像又走一次乾地.
然後對付耶利哥.
然後艾城.
然後贏了五個黃南方.
然後對抗北方下屬無數的盟友.
奪了那傳遞.
整個過程是上帝的手在後面工作.
雖然他進入這個世界.
這個世界已經被佔領了.
你進去的時間沒有人歡迎你會攻擊你.
不過你要知道.
有神是說不要懼怕.
就可以了.
勇敢地獲得.
神就帶他們連夜上去獲得了.
奪了那傳遞.
然後照著神的應許.
要獲得那塊地.
然後繼續繁衍下來.
讓人認識耶和華.
在這個時間.
整個迦南地的人都認識耶和華.
是最頂呱呱的神.
和最厲害的民族.
讚美主我們一起祈禱吧.
天父上帝我們看回整個.
由埃及出來到曠野.
再進入迦南地.
兩次在水裡走乾地.

$^{601}$然後呼喊一個城牆倒塌下來.
然後南方五個王.
北方無數的王.
一起跟他們作戰.
不過都是主帶領他們.
能夠突襲成功.
主對付這麼多的軍隊.
誰夠膽去突襲呢.
而且又不是自己地頭.
不過是上帝帶領.
主我們再次記起你的話.
你跟約書亞說.
跟你的百姓說.
不要懼怕.
懼怕是我們最大的敵人.
我們是剛強壯膽跟隨你.
你叫我們做什麼我們做.
你叫我們要砍斷那些馬的腳筋.
你叫我們燒了那些馬車.
這個地你賜給我們.
主又幫助我們.
其實我們在生命當中.
正如他們這樣說.
後面還有很多未得之地.
主我們一定要得到那些地方.
是討你的喜悅.
我們的靈性上.
我們的生活上.
有很多未得的地方.
給我們可以為主得著.
奪取而榮耀你.
誰聽祈禱.
奉靠耶穌基督命求.
阿們.
\newpage



\section{使徒行傳 1:6-14}
\label{sec:032135}
\textbf{天上天下 唯主獨尊}~~~~趙承恩宣教師
\newline
\newline
連結: \href{http://www.church.com.hk/acms/content.asp?site=cdc\&op=show\&type=product\&code=032135\&layout=sermon}{\texttt{032135}} ~~~~ 語音日期: 2026-05-17
\newline
\newline
\hyperref[sec:032134]{\small{< < < PREV SERMON < < <}}
~
\hyperref[sec:index_chronic]{\small{[返順時目]}}
~
\hyperref[sec:032136]{\small{> > > NEXT SERMON > > >}}
\newline
\newline
使徒行傳 1:6-14
\newline
\begin{longtable}{cl}
\hline
\hline
章節 & 經文 (和合本修訂版)\\
\hline
1:6 & \begin{tabularx}{0.7\textwidth}{X} 他們聚集的時候，問耶穌：「主啊，你就要在這時候復興以色列國嗎？」 \end{tabularx} \\ \\ \relax
1:7 & \begin{tabularx}{0.7\textwidth}{X} 耶穌對他們說：「父憑著自己的權柄所定的時候和日期，不是你們可以知道的。 \end{tabularx} \\ \\ \relax
1:8 & \begin{tabularx}{0.7\textwidth}{X} 但聖靈降臨在你們身上，你們就必得著能力，並要在耶路撒冷、猶太全地和撒瑪利亞，直到地極，作我的見證。」 \end{tabularx} \\ \\ \relax
1:9 & \begin{tabularx}{0.7\textwidth}{X} 說了這些話，他們正看的時候，他被接上升，有一朵雲彩從他們眼前把他接去。 \end{tabularx} \\ \\ \relax
1:10 & \begin{tabularx}{0.7\textwidth}{X} 他升上去的時候，他們定睛望天，看哪，有兩個人身穿白衣站在他們旁邊， \end{tabularx} \\ \\ \relax
1:11 & \begin{tabularx}{0.7\textwidth}{X} 說：「加利利人哪，你們為甚麼站著望天呢？這離開你們被接升天的耶穌，你們見他怎樣升上天去，他也要怎樣來臨。」 \end{tabularx} \\ \\ \relax
1:12 & \begin{tabularx}{0.7\textwidth}{X} 有一座山，名叫橄欖山，離耶路撒冷不遠，有安息日可行走的路程。那時，門徒從那裡回耶路撒冷去， \end{tabularx} \\ \\ \relax
1:13 & \begin{tabularx}{0.7\textwidth}{X} 他們一進城，就上了所住的樓房；在那裡有彼得、約翰、雅各、安得烈、腓力、多馬、巴多羅買、馬太、亞勒腓的兒子雅各、激進黨的西門，和雅各的兒子猶大。 \end{tabularx} \\ \\ \relax
1:14 & \begin{tabularx}{0.7\textwidth}{X} 這些人和幾個婦人，包括耶穌的母親馬利亞，和耶穌的兄弟，都同心合意地恆切禱告。 \end{tabularx} \\ \\
[1ex]
\hline
\hline
\end{longtable}
$^{1}$各位兄弟姐妹,早安.
那麼,剛剛過去的星期四呢.
就是這個耶穌基督的升天日.
可能大家不提起呢.
就未必記得,或者未必知道的.
原因是呢.
因為無假放嘛,是不是?.
大家星期四都可能要上班上學啦,這樣.
那麼.
所以呢,其實今天我們可以說是升天日之後的第一個主日.
那麼所以在說到的時候呢.
我就想跟大家看一段短片.
那麼其實這段片呢,就是剛剛過去的復活節呢.
就是南韓那邊的教會就舉行了一個聚會.
那麼就是想重現這個耶穌復活和升天的場景.
那麼所以他們就找了一個演員去扮演耶穌.
那麼扮演耶穌的時候當然是復活啦.
那怎麼知道呢,在升天的過程裡面呢.
就因為這個機件故障.
那麼他就被一個起重機一路吊一路吊一路吊.
然後吊到上半空.
然後呢,就在這個半空之中呢.
就因為天…就是天氣令到那個機件故障.
然後吊了上半空之後呢.
就下來不回來了.
就一直在那個半空那裡呢.
就…大概好像不知道說.
四個小時至半天左右才下來回來,這樣.
那麼.
那麼幸好感恩就是那位演員呢.
就扮演耶穌的演員平安無事回到地上.
那麼當我看這段片的時候呢.
那麼因為呢,其實我本身看那段片呢.
這段不是我原…看的原片來的.
就…就說呢,因為看那段原片呢.
就說…這位….
就說這位演員是想扮演耶穌而已.
不是想去見耶穌.
就是大家明白這個意思.
那麼所以我看完之後呢.

$^{41}$就一直在笑.
那麼然後呢.
看到笑的過程裡面呢,也都….
讓我去到重新思索一下.
其實耶穌升天.
對我們有什麼意義呢.
那麼為什麼又要升天呢.
那麼甚至乎呢,為什麼我們說.
每一個星期我們去….
以使徒信經去迎信的時候呢.
我們都會讀到…我們都會迎信.
耶穌基督升天的呢.
那麼我們今天可以從.
今天的經文去到發現呢.
基督升天的意義.
以致到呢.
對於今時今日的我們來說.
到底耶穌升天.
跟我們有什麼關係呢.
那麼我們想邀請大家打開聖經.
那就是這個《使徒行傳》的第一章.
《使徒行傳》的第一章.
好了,那現在差不多.
那麼當我們去看這個.
《使徒行傳》的第一章的時候呢.
其實就是記載到耶穌復活之後.
就吩咐門徒不要離開耶路撒冷.
要等候天父的應許.
那麼這個應許又是怎麼樣呢.
就是門徒要領受聖靈的洗.
接著呢.
門徒聚集在一起.
就問耶穌.
主啊.
你就要在這時候復興以色列國嗎.
可能我們第一次看這段經文的時候.
就會覺得.
為什麼門徒已經經歷了耶穌的死.
耶穌的復活.
但是他仍然還是這麼狹窄的眼光.

$^{81}$仍然只是將目光.
望向以色列.
不過如果我們去看這段經文的時候.
我們發現耶穌沒有否定他們.
就是耶穌是沒有否定門徒.
不過呢.
祂就講明.
就是指明天父憑著自己的權柄.
所定的時候和日期.
就不是門徒可以知道.
不過.
就有一件事可以很肯定.
就是聖靈會降臨在他們身上.
就會讓他們有能力.
在耶路撒冷.
猶太傳地.
撒瑪利亞.
直到地極.
作耶穌的見證.
那麼其實整件事的目光.
是闊了的.
就是耶穌就帶領了這群門徒.
將目光不只是放在以色列.
而是將整個目光.
帶到地上.
就是整個世界.
那麼這個問題就是.
門徒要將福音.
一步一步傳到地極.
那麼之後耶穌講了這番話.
就是吩咐了他們要做之後.
就是門徒要做這個傳福音的舉動之後.
耶穌就被接上聖.
就有一朵雲彩從門徒眼前.
就將接了耶穌到升天了.
那麼可能呢.
我們就會問.
耶穌升天有什麼意義呢?.
那麼或者其實.
我們其實在過去.

$^{121}$不同的基督徒.
甚至乎教會裡面有很多神學家.
都對耶穌的升天.
作了不同的解釋.
也都講到明就是.
耶穌升天其實有多重的意義.
不過在講這些意義之前.
可能我們都會有一個很現實的問題.
就是.
如果耶穌不升天離開.
一直讓我們看到祂.
就是說.
今時今日.
到2026年了.
我們仍然可以在一個地方那裡.
就找到我們救主耶穌.
看到祂有救主耶穌的時候.
其實不是更加好嗎?.
就是有時候我們都.
我們作為信徒.
都會被一些未信耶穌的人.
去挑戰我們.
挑戰我們是什麼呢?.
就是.
哇!你的耶穌這麼厲害.
又說永生.
又說是永恆的主宰.
不如祂下來讓我見一見.
讓我摸一摸.
讓我碰一碰.
可能我們都會有這些挑戰.
我們那時候.
如果我們真的跟祂說.
你去這個地方.
你就可以找到耶穌了.
可能整件事立刻很不同.
我們就會覺得.
真的帶人信耶穌.
好像會容易很多.
不過.

$^{161}$我們細心一點想一想.
是不是真的這樣呢?.
所以我們這次去解釋.
升天有什麼意義的時候.
其實就是逐步逐步.
去拆解這些問題.
首先我們要知道一件事.
就是在聖經裡面.
不只是只有耶穌從死裡復活.
例如.
耶穌曾經使拉撒路復活.
而在舊約裡面.
以利亞和以利沙兩個先知.
也曾經使死人復活.
不過這些復活了的人.
到最後都是死了的.
所以今時今日我們不會說.
我們會看到拉撒路.
就是活生生的拉撒路在我們面前.
但是.
如果我們說.
耶穌的復活跟這些死了的人復活.
有什麼不同呢?.
我們就會說.
因為耶穌復活了之後.
祂擁有了一個永恆的生命.
就是說耶穌不會再死.
而升天就象徵著.
祂擁有著統管一切的權柄.
表明了就連死亡.
都是在耶穌的統管之下.
所以我們說耶穌升天.
雖然說有很多個象徵.
但是其中一個最重要的象徵.
就是象徵著耶穌基督的升天.
就是祂已經得盡了一切的榮耀.
權柄.
超越也統管了一切的受眾.
所以.
耶穌的升天.

$^{201}$就跟那位扮耶穌的演員的升天.
是絕對不同的.
因為那位扮耶穌的演員.
一路升天的過程.
他仍然受制著什麼呢?.
首先受制著天氣.
因為天氣的影響.
導致機件故障.
然後.
導致他一路升上天的時候.
死亡仍然是威脅著他的.
是不是?.
就是如果你想一下.
一個不幸地一下子掉下來.
就真的是從天降下的那個狀態.
也不是開玩笑.
真的是搞出人命.
所以死亡.
空間.
時間.
仍然是威脅著那位扮耶穌的演員.
而耶穌基督的升天.
我們會說.
其實祂不是一個空間上的升天.
什麼意思呢?.
就是說.
耶穌不是升到上宇宙.
所以有些人就說猜測.
以前有些人猜測.
耶穌是不是外星人.
可能近期也多了一些人猜測這件事情.
因為美國公開了一些外星發現的文件.
我會說.
耶穌不是外星人.
耶穌就是一個人.
而祂的升天不是說祂升上宇宙.
而是在說祂重新回到天父上帝的角度.
和天父一同掌權.
所以耶穌的升天其實是一件超越的事蹟.
祂在說.

$^{241}$就是耶穌是一個統管一切萬有.
和上帝同等的存在.
所以我們今天的講題.
天上天下為主獨尊.
就正是要表達.
耶穌基督的升天事件就是指出.
世上沒有任何人和任何事物.
能夠高過耶穌基督.
所以耶穌基督的升天.
是表達了祂是至高的上帝.
剛剛也問到大家的問題.
如果耶穌沒有升天.
祂在我們這個地上一直和我們一起.
好像幫助到我們傳福音.
因為耶穌也叫我們.
臨升天之前也跟我們說我們要傳福音.
如果祂在傳福音就容易很多了.
我會說耶穌基督的升天.
其實不是要遠離我們.
也不是不顧這個世界而去.
因為歷世歷代的基督徒和神學家也指出了一件事.
就是祂的升天不會遠離了我們.
反而是拉近了和我們的距離.
如何解釋這件事呢?.
明明說祂升天.
怎麼會拉近了和我們的距離呢?.
好像很矛盾.
問一個問題就可以解答得到.
我問兩個問題.
我想問一下去過耶路撒冷的弟兄姊妹.
可不可以舉手?.
在座.
去過耶路撒冷的弟兄姊妹.
我有去過的.
我看到十個以內.
好,問第二個問題.
我想問一下.
有向耶穌祈禱的弟兄姊妹可不可以舉手?.
理論上應該全部都是.
如果沒有的話.

$^{281}$是不是有些東西阻礙著你.
我們可能會聊一下.
我為什麼會問這兩個問題呢?.
我相信一件事.
就是我們每一個基督徒都有向耶穌基督祈禱的.
而耶穌基督升天.
其實是什麼意思呢?.
就是因為耶穌基督升天.
就令到祂不再固定在耶路撒冷.
或者要我們去某一個地方.
我們才能找到祂.
就是說我們不用故意說.
要買一張機票.
然後就走過去耶路撒冷.
然後可能排長龍.
排完一段長龍.
看到祂.
然後親吻祂的右手.
然後就走.
或者好像多嗎哪樣.
要戳祂的肋旁.
然後才走.
不是這種意思.
耶穌基督的升天.
是讓祂可以藉著聖靈.
無處不在.
甚至乎可以藉著聖靈.
住在我們每一個人的心裡面.
令到祂和我們更加近.
讓我們可以藉著祈禱.
向祂一起和祂聊天.
和祂接觸.
但問題是.
我們會不會和祂有更加多的祈禱.
有更加多的接觸呢?.
有時候人是很有趣的.
我們很想去一個地方朝聖.
就會更加多地向耶穌祈禱.
所以.
實情就是.

$^{321}$耶穌和我們很近.
耶穌就住在我們的心裡面.
但是.
我們可以離上帝.
離耶穌很遠.
因為我們也不太想和祂接觸.
所以祈禱.
正正就是耶穌基督.
和我們和耶穌基督更加貼近的方法.
所以為什麼我們教牧之家會說.
我們很著重.
弟兄姊妹有沒有祈禱的生活.
有沒有靈修的生活.
其實是因為.
我們有沒有和這位主更加貼近.
我們回到經文.
當門徒看到耶穌升天.
被雲彩接去之後.
就有兩個白衣人告訴他們.
我們相信這兩個白衣人是天使.
這兩個白衣人告訴他們.
你們看到祂怎樣升上天.
祂也要怎樣彌臨.
耶穌基督的彌臨.
就是天國的臨到.
我們就說天國什麼時候來.
我們沒有人知道.
正如經常聽到年輕人引用高牧師說的那一句.
就是耶穌什麼時候來.
或者天國什麼時候來.
我們不知道的.
但是我們只會知道一件事.
就是每一天又近了.
又近一點了.
所以在這個居間的時間.
我們說中間的時間裡面.
我們其實就是靠著聖靈的恩賜.
去為耶穌做見證.
為耶穌做見證是什麼呢?.
就是讓更多的人.

$^{361}$去認識耶穌是主.
耶穌是主.
不只是讓其他人認識到.
耶穌愛他.
說耶穌愛他我們也很常說.
但是真正要傳福音給其他人.
其實是要讓其他人知道.
耶穌是主.
什麼意思呢?.
就是要讓其他未信的人.
到最後能夠承認.
耶穌是一個所有造物的主.
是一個救主.
祂是至高無上.
天上天下唯有祂是獨尊的主.
要讓其他未信的人.
到最終能夠信服於祂.
所以我們所傳的福音.
就不是一個幸福音.
不是單純就是說.
耶穌愛你就完了.
或者是說一個成功追求幸福的信仰.
不是這樣.
而是一種能夠否定.
地上有任何事物.
能夠高於耶穌的信仰.
難不難呢?.
我就覺得超難.
作為教母是超難.
我單純說一個「耶穌愛他」.
人們都不太理睬我們.
更何況我們還要是一個.
好像越了很多級的挑戰.
去說耶穌是一個至高無上的主的時候.
就更加難.
原因是當我們去宣揚這個信仰的時候.
就必然需要面對內外的各種挑戰.
這些挑戰可能是什麼呢?.
第一個挑戰就是.
我們自己是不是真的都承認.

$^{401}$耶穌是至高的主.
不只是口中承認.
而是我們的生命裡面.
有沒有真的這樣的表達.
可能我們怕我們的上司.
多於怕上帝.
有時候我們真的上班嘛.
上班是很實在的.
但是我們再推多幾步就是.
其實可能是我們怕那些錢.
怕沒錢.
多於怕我們的上帝.
又舉一個例子.
可能呢.
我有時候跟我太太聊天的時候.
她經常自己都會說她自己.
大概公開試的時候她的一些經歷.
她的經歷是什麼呢?.
就是她自己的經歷就是.
她兩次公開試都是一通過的.
就是都過了的.
那麼說兩次就不是DSE了.
大家都明白.
就是有一定的年齡了.
就是會考和高考.
那麼什麼意思呢?.
就是.
但是她就說.
她兩次的考試裡面.
她是沒有停止過.
她在教會上團契.
以及崇拜.
沒有停止過這件事情.
因為她知道信仰的生活.
其實是應該要保持著的.
所以在一個過程裡面的時候.
她都是覺得.
她沒有停止過這件事情.
所以她就繼續.
那麼她也都.

$^{441}$我再延伸下去.
一直聊天的過程就是.
其實.
不行的就找另一條路.
但是有些事情我們不應該放棄.
所以其實我太太就在這個過程裡面.
她都說.
其實她自己都覺得.
這個也是她的見證來的.
對於她來說.
就是.
公開試不能夠高於上帝.
其實你怎麼表達呢?.
就是.
你在公開試的過程裡面.
你有沒有帶她回去.
你有沒有堅持自己回教會.
又或者作為家長的.
有沒有放過子女回教會.
這些是很真實的例子.
到底公開試是不是高於我們上帝.
這些可能是一些內部的挑戰.
那麼對外的挑戰可能就是.
你說.
上帝耶穌基督是高於一切的時候.
有很簡單的例子.
就是.
你養.
是不是耶穌基督扔了餅乾下來養我.
可能是一些很.
這是以前我也被我父母問過的一個問題.
可能是這些問題.
這些挑戰.
不是那麼容易解答得到.
或者我們自己也不是那麼容易實踐得到.
但是.
耶穌基督正正在說一樣東西.
我們不是靠著自己的能力去面對.
而是靠著聖靈給我們的力量去面對.
當我們宣認耶穌是主的時候.

$^{481}$其實.
耶穌的升天.
正正就是在帶給我們聽.
我們不是靠自己.
而是我們靠著.
耶穌的能力.
耶穌藉著聖靈給我們的能力.
去面對這些挑戰.
一些很現實的挑戰.
我們不是要逃避.
而是我們靠著聖靈的力量去面對.
正正因為這個挑戰那麼艱鉅.
特別是第一世紀的教會面對的問題更加困難.
因為他們相信耶穌.
是會死的.
是會殉道的問題.
所以在一個生命的問題的時候.
門徒知道要面對這樣的挑戰.
他們第一件事要做的就是聚集在一起.
同心合意地懇切禱告.
因為他們門徒意識得到.
唯有貼近上帝的心意.
也遵照著基督的教導.
靠著聖靈的恩賜.
我們才能夠去見證耶穌是主.
在最後我想問幾個問題.
作為我今天講到的結束.
不用回答我的.
自己心裡面就可以了.
這幾個問題其實就是.
反映著我們梁法堂的弟兄姊妹.
是不是真的視耶穌是主.
如果我們有這些特質.
或者有齊這些特質的時候.
其實正正就是我們反映著我們的信仰.
真的將耶穌視為至高.
第一個問題.
你自己心裡面回答就可以了.
不用回答我.
你認為梁法堂的弟兄姊妹.

$^{521}$就是我們.
是不是一個願意視耶穌基督為最高的主.
完全信服於祂的群體呢?.
我再說一次.
你認為梁法堂的弟兄姊妹.
是不是一個願意視耶穌基督為最高的主.
完全信服於祂的群體呢?.
大家可以想一想.
第二個問題.
你認為我們是不是一個願意為福音的緣故.
同心合意懇切禱告的群體呢?.
第三個問題.
你認為我們是不是靠著聖靈的恩賜.
而不是靠自己的能力.
去用生命見證耶穌是主的群體呢?.
可能這三個問題對我們來說.
都不是那麼容易回答.
但是當我們宣告天上天下為主獨尊的時候.
其實我們都應該在這三個問題裡面.
回答「是」.
回答「Yes」.
我們都是有這些特質.
不過可能在一個實際的場景裡面.
我們可能離這些特質有些距離.
不要說很遠.
不要說完全沒有.
可能差一點點.
但是在這個差一點點的過程裡面.
我們有一個不二的法門.
就是我們能夠藉著祈禱.
願上帝幫助我們能夠成為合祂心意的見證人.
藉著聖靈的幫助.
讓我們每一個人能夠在我們生活,生命裡面.
去見證耶穌是主.
耶穌是我們最高的主宰.
讓我們一同禱告.
因為你藉著你的愛子耶穌基督.
祂的生平,祂的死亡,復活,升天.
你讓我們知道我們所信靠的主是那位最高的主宰.
大父上帝我們會有軟弱的時候.

$^{561}$因為當耶穌說要我們用生命去作祂的見證的時候.
我們發現現實有很多的挑戰.
有很多的困難.
我們要面對.
但是耶穌不是要我們靠著自己的力量去面對.
而是靠著聖靈的力量去面對.
因為我們知道聖靈的力量能夠讓我們的生命更加豐富.
以致我們有勇氣可以去到見證耶穌是主.
就願你的靈幫助我們.
讓我們每一個人可能我們還是有軟弱的時候.
能夠使我們剛強.
當我們有恐懼的時候.
讓我們有勇氣.
當我們不知道怎樣去選擇的時候.
願你的旨意臨到我們當中.
讓我們的生命能夠作耶穌基督.
我們這位升天的基督.
我們這位升天的救主的見證.
我們的禱告交託是靠我們救主耶穌基督的名.
阿們.
\newpage



\section{路加福音 7:11-17}
\label{sec:032136}
\textbf{夕陽西下 斷腸人在天涯}~~~~鄧瑞強牧師博士
\newline
\newline
連結: \href{http://www.church.com.hk/acms/content.asp?site=cdc\&op=show\&type=product\&code=032136\&layout=sermon}{\texttt{032136}} ~~~~ 語音日期: 2026-05-17
\newline
\newline
\hyperref[sec:032135]{\small{< < < PREV SERMON < < <}}
~
\hyperref[sec:index_chronic]{\small{[返順時目]}}
~
\hyperref[sec:032137]{\small{> > > NEXT SERMON > > >}}
\newline
\newline
路加福音 7:11-17
\newline
\begin{longtable}{cl}
\hline
\hline
章節 & 經文 (和合本修訂版)\\
\hline
7:11 & \begin{tabularx}{0.7\textwidth}{X} 過了不久，耶穌往一座城去，這城名叫拿因，他的門徒和一大群人與他同行。 \end{tabularx} \\ \\ \relax
7:12 & \begin{tabularx}{0.7\textwidth}{X} 當他走近城門時，有一個死人被抬出來。這人是他母親獨生的兒子，而他母親又是寡婦。城裡的許多人與她一同送殯。 \end{tabularx} \\ \\ \relax
7:13 & \begin{tabularx}{0.7\textwidth}{X} 主看見那寡婦就憐憫她，對她說：「不要哭。」 \end{tabularx} \\ \\ \relax
7:14 & \begin{tabularx}{0.7\textwidth}{X} 於是耶穌進前來，按著槓，抬的人就站住了。耶穌說：「年輕人，我吩咐你，起來！」 \end{tabularx} \\ \\ \relax
7:15 & \begin{tabularx}{0.7\textwidth}{X} 那死人就坐了起來，開始說話，耶穌就把他交給他的母親。 \end{tabularx} \\ \\ \relax
7:16 & \begin{tabularx}{0.7\textwidth}{X} 眾人都驚奇，歸榮耀給神，說：「有大先知在我們當中興起了！」又說：「神眷顧了他的百姓！」 \end{tabularx} \\ \\ \relax
7:17 & \begin{tabularx}{0.7\textwidth}{X} 關於耶穌的這事就傳遍了猶太和周圍地區。 \end{tabularx} \\ \\
[1ex]
\hline
\hline
\end{longtable}
$^{1}$各位弟兄姊妹早晨.
願神的恩惠常與你們同在.
大家有沒有試過絕望呢?.
親手經營幾十年.
以為可以養老的家園.
轉眼之間就灰飛煙滅了.
一起成長的玩伴.
幾十年來互相支持的老街坊.
忽然之間就告別人間.
摯愛的家人.
相依的老伴.
一下子就永別.
當中的淒慘.
真是無法言傳.
身在其中的人.
連絕望的眼淚.
都流不出來.
曾經有一個晚上.
這一刻.
我還和我太太聊天.
下一刻.
她就離開了這個塵世.
整個世界一剎那之間.
真是一剎那之間.
返回太古之前的蒼白和虛無.
絕望臨到的時候.
其實只是一片空白.
空白到連眼淚都不懂流出來.
路加福音筆下的耶穌.
就是上帝走入人間.
擁抱絕望的人.
不同的福音書描繪.
描繪出耶穌不同的畫像.
馬太福音筆下的耶穌.
好像摩西一樣.
走上山頒佈律法.
宣講登山寶訓.
塑造神的新指紋.
馬可筆下的耶穌.
背著沉重的十家.

$^{41}$呼喚門徒跟著祂.
而路加福音筆下的耶穌.
最富有人情味.
最喜歡和人吃飯.
喜歡飯局的弟兄姊妹.
應該多讀路加福音.
這個福音書裡有很多路加獨有的故事.
都是很有人性溫度的.
譬如好撒瑪利亞人故事.
浪子回頭故事.
瑞利薩該的故事.
這些都是路加獨有的.
瑞利薩該因為耶穌和他吃了一餐飯.
他就將他一半的財富分給窮人.
這餐飯貴到不得了.
但是他認為值得.
今天講到路加福音的一個小故事.
也是路加福音獨有的.
這個故事以一個悲劇為背景.
真是在夕陽西下的歲月.
一個斷腸人在天涯.
絕望至極,哀傷至極.
耶穌以祂神聖的人詞.
與這個絕望的人相遇.
人性的溫度讓冰冷的生命重拾希望.
我們看看這個小故事.
記載在路加福音第七章第十一節至第十七節.
經文在我們程序表較後的部分.
網上的朋友也可以打開聖經.
打開了就不要蓋上了.
路加福音第七章第十一節.
過了不多時耶穌往一座城去.
這座城名叫拿恩.
他的門徒和極多的人與他同行.
將近城門有一個死人被抬出來.
這人是他母親獨生的兒子.
他母親又是一個寡婦.
有城裡很多的人同著寡婦送殯.
主漢見那寡婦就憐憫她.
對她說不要哭.

$^{81}$於是進前按著缸.
抬的人就站住了.
耶穌說少年人我吩咐你起來.
那死人就坐起並且說話.
耶穌便把他交給他母親.
眾人都驚奇.
歸榮耀與上帝說.
有大先知在我們中間興起來了.
又說上帝眷顧了他的百姓.
他者是的風聲就傳遍了猶太和周圍地方.
這個故事很簡單.
只有七節經文.
沒有冗長的情節.
也沒有長篇大論.
卻包含了福音一個好的消息.
福音的多項元素.
例如死亡與生命.
失去與重獲.
絕望與希望.
哭泣與安慰.
天母眼與天友眼.
人間的苦難與上帝的憐憫.
獨生子的死亡與復活.
神的話語的創造的權能.
等等等等.
七節經文而已.
有很多內容.
有人說這是福音信息最濃縮的版本.
就是福音的微型.
福音不是一些很抽象的大道理.
福音一個好的消息.
就是神與人相遇的故事.
是神擁抱絕望的人.
我們慢慢分開幾點看這個故事.
第一點就是生與死的相遇.
這個小故事就開始在一幕.
就是生與死的相遇.
第十一節說到.
耶穌的門徒和極多的人與耶穌同行.
耶穌是生命的主.

$^{121}$這是散發生命的行列.
一大班人應該都是興高采烈.
另外一邊廂是死亡的行列.
是絕望的行列.
是世人必經的行列.
第十二節經文說.
將近城門有一個死人被抬出來.
這個人是他母親獨生的兒子.
他母親又是一個寡婦.
有城裡很多人與這個寡婦送殯.
幾個字就已經道盡當事人的淒涼.
這是一個寡婦.
她過去所依靠的丈夫已經死了.
她有一個獨生子.
這是她的將來,她的希望.
一切的人生都是寄託在這個獨生子的身上.
如今這個獨生子都死了.
沒有了過去,沒有了將來.
這個寡婦的當下只是一片空白.
在這個故事裡.
她一直都沒有出過一句聲.
她的無言是天地間最震天動地的哀哭.
就在這個生與死的相遇裡.
耶穌主動地終止了這個死亡行列的前行.
第二點,我們看到主看見,主憐憫.
經文第十三節.
主看見那寡婦就憐憫她.
耶穌看到這個寡婦.
耶穌看到在痛苦深淵裡的絕望的人.
耶穌看到我們.
然後經文說耶穌憐憫她.
憐憫就是我們所信的上帝的本性.
千萬不要經常想上帝的本性.
什麼全能,全善.
不要經常想那些哲學的概念.
我們信仰裡的上帝是很仁慈的.
祂憐憫我們.
這就是上帝的本性.
人間的憐憫.
常常只不過是一種同情的情感.

$^{161}$但神的憐憫其實就是一種行動.
當我們說神憐憫.
其實祂就有行動.
堅決的行動.
就因為這一種憐憫.
神進入了人間.
將人生的苦難.
擺在自己的肩頭上.
一力擔起.
這個故事裡刻意的說.
這是獨生子.
是獨生子的死亡.
其實也預示了天父的獨生子.
耶穌祂自己的死亡.
對神來說.
憐憫一個人.
就不只是同情祂.
而是行動上為這個人.
付出自己的一切.
神的憐憫就以祂救贖的行動.
對人類的苦難.
大聲的說不.
憐憫的上帝進入人間.
與絕望的我們相遇.
祂解救我們.
接著耶穌開始說話.
接著的一點就是不要哭.
第十三節.
主對寡婦說.
不要哭.
在絕望的處境裡叫人不要哭.
不要哭.
是不是說風涼話呢?.
能夠不哭.
根本就不是真正的絕望.
很多時候我們叫人不要哭.
或者某些場合我們說節哀順變.
這些都是廉價的安慰.
其實沒有辦法安慰到當事人.
很多時候顯出你根本不明白當事人的那種哀傷.

$^{201}$在這裡耶穌對寡婦說.
不要哭.
祂不是要做一些廉價的安慰.
而是讓寡婦體會到神的介入.
一些新的事情.
想不到的事情將要介入.
這個女士將要面對的.
不是單單苦難的本身.
而是要準備體會闖入她生命裡的生命的神.
在這個苦難的人間.
我們有千八個哭泣的理由.
但是在神的恩典底下.
我們同樣有千八個不去哭泣的理由.
生命充滿悲劇.
要哭總有理由.
而要不哭.
就必須要有一個克服悲劇的神.
這個神.
耶穌說要與寡婦相遇.
接著耶穌出手了.
並且祂發言.
第十四節.
於是耶穌進前來.
按住那個鋼.
台的人就站住了.
耶穌說少年人.
我吩咐你起來.
耶穌很主動的.
祂走來主動的.
經文說祂進前.
祂主動的多走一步.
用手按住棺材架.
抬著棺材的人就停在那裡.
棺材架是不潔的.
帶有死亡的氣息.
挖出了死亡的界線.
一般人.
猶太人.
是不應該觸碰這個界線的.
何況是尊貴的拉比耶穌.

$^{241}$祂更加不應該觸碰這個界線.
觸碰的人會沾染不潔.
耶穌是刻意的.
進前來.
祂就按住這個棺材架.
去接觸這個祂不應該接觸的界線.
真正的神聖.
不是遠離這些所謂不潔的領域.
而是觸碰祂.
瓦解祂背後導致絕望的根源.
耶穌正面的攔住死亡的行列的前進.
耶穌就站在死亡的界線的前面.
攔住祂.
阻止祂的前進.
耶穌闖入了死亡的界限.
要在絕望的深淵.
燃起生命的火.
然後耶穌就發言了.
祂發出了神聖主權的命令.
沒有任何的東西.
可以違抗這個命令.
耶穌說少年人.
我吩咐你起來.
耶穌不需要和死神討價還價.
不需要與人討論死人復活.
在科學上是不是可能的.
耶穌不需要考慮.
這個合不合乎人類理性的框架.
耶穌就直接發出祂的命令.
少年人.
我吩咐你起來.
這是創造主在創造萬有的時候.
發出的那種帶有權能的命令.
神說要有光.
就有了光.
神說有就有.
命立就立.
神從虛無裡面創造出萬有.
祂命令光明從黑暗裡照出來.
祂從死亡中呼喚生命.

$^{281}$上帝的話語帶著權能.
祂的命令一出.
事就這樣成了.
耶穌以神的命令.
創造出一個新的世界.
神的話語一出.
對這個寡婦來說.
舊事已過.
都變成新的了.
第十五節.
那死人就坐起並且說話.
耶穌便把他交給他母親.
就算死人.
也要聽從神的命令.
這個死人就坐起來並且說話.
在這裡我們看到.
死亡原來是沒有話事權的.
我們每個人都高估了它.
我們怕了它.
很害怕.
但是原來死亡根本沒有話事權.
在神的話語的面前.
死亡只能夠俯首聽命.
現在是復活期.
說這個故事真是好了.
主從死裡復活了.
生命雖然有限.
但是死亡對我們來說.
沒有終極的權柄.
它只能夠降服於神.
耶穌讓我們每一個人都看到.
在死亡的界限以外.
其實是有更廣闊的生命世界.
有神的話語所創造的明天.
有神開啟的希望.
耶穌把那個復生的人交回給他媽媽.
我們不知道他媽媽會怎樣想.
但是這個明顯地顛覆了他一切的想法.
明明是沒有可能的.
怎會變成有可能的呢?.

$^{321}$明明是死亡,黑暗,絕望.
為何會變成生命,光明,希望呢?.
他一生生活的邏輯.
就是死了就死了,不能挽回的.
殊不知在他的邏輯以外.
有一個永活的上帝.
這個上帝恩著憐憫.
他甘願闖入人間.
改變了他的一生.
我要講結論了,很快故事又完了.
由頭到尾,大家都看到.
那個寡婦沒有說過一句話.
她沒有求過耶穌任何的東西.
她其實是一個很沉默而痛苦的人.
全部事情都是由耶穌基督主動搞出來的.
是耶穌主動做的.
我們說這個就是白白的恩典.
各位弟兄姊妹.
無論你現在在什麼光景.
或者是憂慮,或者是絕望.
施恩的神都會介入你的生命.
他看到你,憐憫你.
他正面的面對虛無的力量.
以他權能的命令創造光明.
所有這些都是上帝主動送贈給我們的禮物.
都是恩典.
這個就是福音.
我們說傳福音.
是不是很好的一個好消息.
為什麼不好好傳給其他的人.
上帝闖入我們這個人生裡.
給我們一個新的生命.
恩典在這裡.
福音在這裡.
神在這裡.
故此各位弟兄姊妹.
親愛的弟兄姊妹.
請不要哭.
願神賜福你們.
願榮耀歸給聖父,聖子,聖靈的名.

\newpage



\section{瑪拉基書 1:6-14}
\label{sec:032137}
\textbf{「A貨」的敬拜}~~~~石偉漢牧師
\newline
\newline
連結: \href{http://www.church.com.hk/acms/content.asp?site=cdc\&op=show\&type=product\&code=032137\&layout=sermon}{\texttt{032137}} ~~~~ 語音日期: 2026-05-17
\newline
\newline
\hyperref[sec:032136]{\small{< < < PREV SERMON < < <}}
~
\hyperref[sec:index_chronic]{\small{[返順時目]}}
~
\hyperref[sec:032138]{\small{> > > NEXT SERMON > > >}}
\newline
\newline
瑪拉基書 1:6-14
\newline
\begin{longtable}{cl}
\hline
\hline
章節 & 經文 (和合本修訂版)\\
\hline
1:6 & \begin{tabularx}{0.7\textwidth}{X} 萬軍之耶和華對你們說：「兒子孝敬父親，僕人敬畏主人；我既為父親，孝敬我的在哪裡呢？我既為主人，敬畏我的在哪裡呢？你們這些藐視我名的祭司啊！」你們卻說：「我們在何事上藐視你的名呢？」 \end{tabularx} \\ \\ \relax
1:7 & \begin{tabularx}{0.7\textwidth}{X} 「你們將不潔淨的食物獻在我的祭壇上，卻說：『我們在何事上使你不潔淨呢？』你們說，耶和華的供桌是可藐視的。 \end{tabularx} \\ \\ \relax
1:8 & \begin{tabularx}{0.7\textwidth}{X} 你們將瞎眼的獻為祭物，這不算為惡嗎？將瘸腿的、有病的獻上，這不算為惡嗎？那麼，請把這些獻給你的省長，他豈會悅納你，豈會抬舉你呢？這是萬軍之耶和華說的。」 \end{tabularx} \\ \\ \relax
1:9 & \begin{tabularx}{0.7\textwidth}{X} 現在我勸你們要懇求神，好讓他施恩給我們。這事既出於你們的手，他豈會抬舉你們任何人呢？這是萬軍之耶和華說的。 \end{tabularx} \\ \\ \relax
1:10 & \begin{tabularx}{0.7\textwidth}{X} 萬軍之耶和華說：「甚願你們中間有人把殿的門關上，免得你們徒然在我壇上燒火。我不喜歡你們，也不從你們手中悅納供物。」 \end{tabularx} \\ \\ \relax
1:11 & \begin{tabularx}{0.7\textwidth}{X} 萬軍之耶和華說：「從日出之地到日落之處，我的名在列國中必尊為大。在各處，人必奉我的名燒香，獻潔淨的供物，因為我的名在列國中必尊為大。 \end{tabularx} \\ \\ \relax
1:12 & \begin{tabularx}{0.7\textwidth}{X} 你們卻褻瀆我的名，說：『主的供桌是不潔淨的，供桌上的果子和食物是可藐視的。』 \end{tabularx} \\ \\ \relax
1:13 & \begin{tabularx}{0.7\textwidth}{X} 你們又說：『看哪，這些事何等煩瑣！』並嗤之以鼻。這是萬軍之耶和華說的。你們把搶來的、瘸腿的、有病的拿來獻上為祭，我豈能從你們手中悅納它呢？這是耶和華說的。 \end{tabularx} \\ \\ \relax
1:14 & \begin{tabularx}{0.7\textwidth}{X} 行詭詐的人是可詛咒的！他的群畜中雖有公羊，他許了願，卻將有殘疾的獻給主。因我是大君王，我的名在列國中是可畏的。這是萬軍之耶和華說的。」 \end{tabularx} \\ \\
[1ex]
\hline
\hline
\end{longtable}
$^{1}$弟兄姊妹主內平安.
各位弟兄姊妹和家中的弟兄姊妹都平安.
特別是新朋友歡迎你們來到當中.
我們一起有個祈禱.
天父我們去將我們設置在你面前.
幫助我們去知道你是何等尊貴.
配受敬拜的那一位.
讓我們真的去將我們這個時刻.
恭敬的仰望你.
讓我們去打開我們的耳朵.
去聆聽你的說話.
幫助我們.
聽我們祈禱奉耶穌的名字而求.
阿們.
今天和大家分享的題目叫做A4的敬拜.
不知道大家知不知道是什麼意思.
稍後我會講解一下.
不過接下來有一件大事發生.
不知道大家知不知道有什麼大事發生.
不過可能對我來說是大一點的事.
可能對你來說是普通的事.
對強哥來說應該一說就知道是什麼事.
強哥接下來有什麼大事發生.
對了強哥一說就知道.
一定是世界盃.
四年一度的世界盃.
接下來很快就會開始比賽.
而今年世界盃的特別地方是什麼呢.
就是有幾十隊伍.
今年以往沒有這麼多隊伍.
現在是多到48隊去打世界盃.
打了很久的循環 初賽 復賽等等.
如果你是喜歡看球賽的人.
我想你都很期待著這個時刻.
甚至乎可能你都會花一點時間.
去了解一下究竟哪一隊會贏.
有很多網上的分析會說哪一隊會贏.
當然我不是專家.
我問問我們的專家.
強哥今年哪一隊會贏.

$^{41}$哪一隊.
中國.
英國.
大家都很期望著英國會贏.
我期待著的應該是阿根廷.
阿根廷有美斯應該會贏的.
但你有沒有想過.
既然你很期待著參加這個比賽的時候.
通常我們為了迎接這個比賽.
我們做一些預備.
做什麼預備呢.
其中就是.
作為支持他的球隊.
買一些球衣.
帽子.
一些設備.
買一些紀念的足球等等.
很多時候在這些世界盃的日子裡面.
都很期待很準備去做這件事.
但你有沒有想過.
當在這個期待著這麼興奮的節日裡面.
如果甚至乎你有機會.
進到球場看球賽的話.
你會穿一件怎樣的球衣呢.
但原來在這個坊間裡面.
亦都出現一個情況就是.
在世界盃的名牌球衣裡面.
同時間有一些叫做A4的球衣.
什麼叫做A4呢.
才可以嚇嚇大家.
沒有人做製衣業的.
明不明白什麼叫做A4呢.
假的.
假的這種說法都不是很準確.
怎樣叫做A4呢.
扮真.
仿真.
再多一點點字眼.
不是.
方向是對的.

$^{81}$但原來是這樣的.
原來所謂A4的概念.
就是在貿易的貨品裡面.
去分類裡面嘗試描述的一種說法.
原來貨品裡面都有A4 B4和C4.
其中所謂的A4是什麼呢.
就是叫高仿.
即是說.
有些帽牌.
你一看就知道是帽牌.
其實你都買.
你都不懂那件事.
但所謂的A4就是高仿到什麼呢.
基本上和真的有很少的差別.
少到一個地方就是.
你除非是專家.
或者是很熟悉的.
你就會看到那種的差別.
那些就叫做A貨.
但你又要留意一下.
A貨是不是真貨.
是不是那個正牌貨.
不是.
是不是.
但沒有人買A貨.
有.
是不是.
為什麼你會買A貨.
便宜.
是不是.
便宜很多.
比正版便宜很多.
所以你都會買A貨.
但我想問一下強哥.
又是時候問你了.
你作為這麼捧英格蘭的.
你會不會買A貨的英格蘭球衣.
進去球場支持英格蘭呢.
強哥很堅定的說.
不會.

$^{121}$只有一些始之以備的.
不是吧.
是不是.
正版就可以了.
為什麼呢.
你會發現原來.
真正捧英格蘭球衣的人.
如果他進得了球場.
去支持自己的球隊.
一定穿一件正版的.
正品的.
去支持他的隊伍.
原來是這樣的.
因為你會知道.
如果你穿一件翻版.
進去球場.
你站在穿著正版球衣的人之間.
那些人是歧視你的.
藐視你的.
你都會覺得自己很羞愧.
很噁心.
你立即遮住自己的胸口.
就是會這樣的.
說一個很有趣的問題.
一方面我們說.
我們支持我們的球隊.
我們要用很多方式去支持他.
但另一方面.
有些人就是買一些便宜的.
廉價的東西去頂.
扮高檔.
進去之後你會發現.
原來是帶來自己的羞恥.
自己的沒面子.
不過仍然有人會這樣做.
理由很簡單.
他們希望付最少的代價.
希望混入在正版之中.
原來有很多人都不願意付上代價.
去成為正版貨.

$^{161}$而寧願希望用高仿.
接近正版 扮正版.
最後他知道不是正版.
而混在當中.
他希望套一些光.
今天我們讀到的經文.
在《瑪麗基書》裡面.
正正是在說.
以色列裡面有一群百姓.
同樣都獻上一種A4的敬拜.
他們要去敬拜神.
去到神靈殿裡面去敬拜神.
不過他又不想付這麼多代價.
交一些A4出來.
所以今天我們又看看.
今天這段聖經裡面.
怎樣幫我們去明白更加多這一個段落.
當中我們都解釋少少一些背景.
因為可能有些人不太熟悉聖經.
或者根本上不知道歷史背景.
我們嘗試逐部份去看看這段聖經說什麼.
在我們今天讀的《瑪麗基書》裡面.
其實《瑪麗基書》我上次說到.
是一封情書.
是一封神向以色列所發出的情書.
說的是神對以色列民族的一份愛.
不過當神付出了很多的愛.
但卻收不到應有的愛.
所以神就向他們發出這封書信.
說的是這段愛情.
是神做了什麼呢?.
所以當我們去看少少背景的時候.
我們就要了解一下.
究竟這是一個什麼時間呢?《瑪麗基書》的時間.
因為原來我們知道在以色列裡面.
其中有兩段的歷史.
第一段就是講以色列人被擄.
就是說我們知道有所謂的南國和北國.
以色列去到北國滅亡了.
去到後來連南國僅有的猶大國都滅亡了.

$^{201}$國家滅亡.
所以那些人就不再有自己的國家.
甚至乎如果你記得Jenny講到.
但以理書就被當時打敗了南國猶大的尼布甲尼薩王.
就是巴比倫王.
就帶到利福圖的東面.
就是右手邊去到巴比倫.
當中就是這樣滅了這個國家的時候.
帶了一群知識份子.
有錢一些有能力多一些的人去到巴比倫王國裡面.
就成為了漁民.
漁民就是剩下的.
我們簡單的說就是死症種.
就留下一群死症種帶到巴比倫裡面.
去將他們同化培育.
甚至將他們化為巴比倫帝國的人.
當中那些中等低等窮的就留在猶大國裡面.
所以這段歷史正正好就是在講以色列.
就是被亡國甚至被擄走的一段日子.
但你要知道當他們能夠被擄走的時候.
上帝仍然是愛他們.
因為上帝為他們留下了漁民.
不是讓他們滅族.
而是留下了一些漁民一些死症種.
留在這裡仍然可以生根建造.
重新建立他們的家庭.
這個就是第一段的歷史.
但當去到第二段的歷史就是.
當過了一段時間.
原來巴比倫這個帝國以為很厲害的帝國.
之後都滅亡了.
滅亡的時候就興起了一個新的王朝.
就是馬代波斯的國度.
在這個波斯的國度裡面.
波斯王的政策就是一個比較懷柔的政策.
就容許你們這些在巴比倫的死症種.
可以回到自己的家鄉.
於是就有第二段的歷史.
就是以色列人的歸回.
回到自己的以色列那個地方裡面.

$^{241}$這個聖經很清楚地說都是上帝的恩典.
原來我們不能夠整個國家沒有滅亡.
或者說那些人沒有死.
是上帝的恩典.
而上帝也藉著70年的日子.
讓他們回到自己的國家.
又是上帝的恩典.
上帝容許波斯王讓他們回到他們的地方.
當中裡面.
回去的主要的歷史有兩大段.
第一段就是.
在公元前538年至515年的時候.
讓他們回去重建聖殿.
由索羅巴伯帶領回去重建聖殿.
後來去到大約八多年後.
再去到公元前457年至428年.
由以斯拉和尼希美再帶領他們重建聖殿.
所以你們見到很慶幸的.
他們之前有日子可以回去重建聖殿.
回去又重建聖殿.
當然他們很期望重建家園.
上帝恩典他們終於可以回去.
應該重新興起家園.
希望有一個不一樣的生活.
他們很期待著會發生.
但在這個歷史的發生出現的時候.
原來不是像他們那麼理想.
第一你要知道.
這班人離開了很久.
你想想本來留在以色列的人.
經歷了這八多年.
他們有什麼事發生.
就是那些本來是中產的.
因為走了一班富貴的.
他們就慢慢成為有錢人.
基層的人經歷了這麼多年.
就開始成為中產.
所以這班人就開始在那個地方.
去作主管理的人.
所以這班人回來後都很不一樣.

$^{281}$有很多矛盾衝突發生.
甚至乎他們即使回到以色列這個地方.
他們仍然被波斯去管理.
波斯就派了省長去管理著以色列.
所以他們期望著重新建立一個王國.
恢復以前以色列的國.
的日子不發生.
不能發生.
在這樣的情況下.
他們覺得好像很期待.
有好的生活有好的結果.
有好的環境有好的國家再一次興起.
但是都似乎面對很多的失望.
不行.
不可以.
就是這樣他們就說.
神啊你不愛我的.
神你不愛我的.
所以就出現了第二段聯絡.
就是原來我們上一次看到.
馬拉基書第一章一至到五節.
去到今天看到六至十四節的時候.
其實你會發現就是.
原來一開始的時候.
神就說我曾愛你們.
因為原來以色列.
回到耶路撒冷的這幫人裡面就覺得.
我們回來都見不到上帝.
有很多很多的祝福.
所以他們就說神你不愛我.
所以神一開始就宣告.
我曾愛你們.
神就向這幫以色列人示愛.
說其實我真的很愛你們.
由一開始到現在都這麼愛你們.
不過這幫人的反應是什麼.
你什麼時候愛過我.
我不覺得.
所以上一次我們看到的段落.
就是神向這幫人示愛.

$^{321}$接著我們今天看到的段落就相反了.
就是神就和這幫百姓說.
你不愛我.
你聽得懂嗎.
前一段一至五節.
神就說很愛以色列.
以色列就說我不覺得你愛我.
現在去到這裡六節.
去到十四節就是.
以色列人覺得自己很愛神.
我們有敬拜神.
我們有返教會等等.
今天的概念就是返教會.
但是神就說.
其實你們不愛我.
你們給我的敬拜.
其實是一個A符.
所以我們今天就看到這個段落.
看看究竟有什麼問題發生了.
神說你們的敬拜是有問題.
有什麼問題呢.
我們就先看第二節.
原來敬拜神有兩種態度.
第一種叫應有的態度.
我們看看經文怎麼說.
萬君子和華對你們說.
兒子孝敬父親.
僕人敬畏主人.
我就是神了.
神既為父親.
孝敬我的人在哪裡呢.
神既為主人敬畏神的人在哪裡呢.
原來神說的一個問題是.
原來兒子和僕人對父親和對主人.
都有一種態度.
什麼態度呢.
我們先看兒子對父親是什麼呢.
有一種孝敬父親.
孝敬父親在和學本的另一個翻譯.
用了尊敬父親.

$^{361}$剛剛英文的NIV版本用了honor.
什麼意思呢.
原來神就說.
兒子對父親的態度應該是什麼呢.
就是一種對他孝敬尊敬.
就是說究竟他有多少重量.
你有沒有給他應該有的重量.
不知道你明不明白意思.
就是說究竟那樣東西值多重.
你有沒有給他那種重量.
剛才說那件正牌波衫是多少錢.
你有沒有想過給他應該有的價錢.
去買這件正牌波衫呢.
最近你知道總統特朗普去了哪裡嗎.
去了中國.
中國誰去接見特朗普.
川普.
誰去啊.
國家主席習近平.
最後是說見面.
就是說哪一個級數.
如果他派一個二打六的外交官.
迎接川普.
跟著跟著跟他見了三天.
跟著拜拜.
那你想一下什麼事.
不給面子的了.
所以川普美國總統來到.
你要派中國的主席.
那叫什麼.
因為你的份量重.
不是說他們兩個重量.
他們的地位是很厲害的.
那你要派這樣的地位的人跟他見面.
才叫國家級的對話.
一樣的.
原來說我們對父親.
你覺得你的父親有多少的重量.
有多少的份量.
如果你能夠給一個他應該有的份量給他.

$^{401}$這種就叫做尊敬.
所以聖經教我們.
我們知道父母親對我們重要.
所以上個禮拜母親才說.
我們要答謝我們的母親.
其實父親一樣.
所以聖經裡面怎麼說.
以弗所說.
作兒女的你們要在主裡聽從父母.
這個是理所當然的.
而且說什麼.
要當孝敬父母.
使你得福.
在世長壽.
這是第一條大應許的誡明.
聖經很清楚說.
父親他對我們有養育之恩.
也讓我們能夠出生.
他一定有一定的份量.
甚至乎一個按著猶太的教導.
教育子女的父親.
是一個值得配受頌讚.
尊敬的那位父親.
所以我們做子女要孝敬父母.
做子女要這樣對待父親.
這樣你就明白了.
剛剛最近又有一單新聞.
我也有看這些新聞.
有一個不知道大家認不認識.
他叫李家鼎.
有沒有印象.
李家鼎說他有兩個兒子.
兩個兒子.
其中兩個兒子裡面.
有兩個兒子的媽媽死了.
死了之後原來有很多很多的報道.
搞了很久.
終於越來越清晰了.
那幅圖畫說什麼呢.
直到李家鼎.

$^{441}$我們認識的藝人李家鼎.
他公開了錄音說.
原來那個大兒子是什麼呢.
當他爸爸是什麼呢.
人育ATM.
就是人育提款機.
如果我當Tony是人育提款機.
我一見到Tony就.
嘟嘟嘟嘟嘟嘟嘟.
對著他的臉按他的密碼.
他就過錢給我了.
就是那些.
他每一次回去見爸爸.
就是問李家鼎拿錢.
拿錢拿錢.
爸爸就是這樣的角色.
所以報道說.
李家鼎的前妻死了.
原來十天後.
才知道他的前妻死了.
原來在不知道.
他的前妻死了十天之前.
這個大兒子還在問爸爸.
媽媽有些需要要拿錢.
其實他媽媽已經死了.
媽媽有些需要要拿錢.
你再給我一點錢.
現在還說.
他監視著他定期的200萬.
這個就是大兒子的態度.
小兒子又怎樣呢.
小兒子現在正在照顧他爸爸.
保護他爸爸.
幫他爸爸搞一些法律的事情.
甚至乎嘗試去處理.
他媽媽遺產的問題.
我想問你.
誰叫孝順兒子.
誰才叫孝敬父親呢.
我想都一目了然.

$^{481}$父親對兒子是這樣.
另外神也說到.
另外一個情況是什麼呢.
僕人呢.
僕人要怎樣對主人呢.
僕人要對主人敬畏.
什麼叫敬畏呢.
敬畏就是畏懼和害怕.
就是fear.
就是說你知道他高級過你.
他是大過你.
甚至乎他有很重要的角色.
我們今天不是很明白.
我們習慣上班的概念.
我們上班.
我來了這裡幾年.
我聽了一些新聞.
大家第一時間對他說.
現在那些人上班.
不做就什麼.
就劈牌就走了.
Frank 說什麼.
前門不走都後門走.
有兩道門.
總之找到門就走.
不做了.
他不會和你交代的.
今天我們香港有些新聞.
就說什麼呢.
辭工是怎樣的.
就WhatsApp跟你說一句.
就不上班了.
但是在古代不是這樣的.
古代不是這樣的.
古代呢.
僕人和主人是一個什麼關係呢.
因為在古代來說.
原來一個僕人.
能夠有工作.
已經是莫大的恩典.

$^{521}$因為為什麼呢.
原來有很多人都生活不了.
生活不了.
以致他真的叫賣身契.
去賣身.
就希望有人給錢.
買了他之後.
他就能夠依附在這個家庭裡面.
能夠生活.
能夠被照顧.
以致包你飲包你食.
你能夠還有生命.
還有生活氣息.
因為他會養活你.
所以在古代來說.
就能夠有人照顧你.
養活你.
照顧你.
甚至乎如果你看回.
我們剛剛在看利未記.
之前大家看到利未記.
等等.
出埃及記.
原來主人如果好的.
還會安排老婆給你.
可以這樣的.
所以僕人.
其實不像我們今天這樣想.
很慘.
依附在主人身上.
被人鞭打虐待.
其實不一定是這樣的.
所以僕人能夠有一個.
這麼愛惜他.
照顧他.
照顧他一輩子的主人.
是一個福氣.
所以他們反過來.
如果主人有一天說.
我不要你.

$^{561}$他還害怕.
就是說.
你從此就露宿街頭.
沒有地方住.
沒得吃.
生活在一個擔憂之下.
所以他說.
一個這樣的狀況.
就是什麼呢.
所以僕人會怎樣.
很怕主人.
所以你記得.
我們之前說過.
就是有一個耶穌的比喻.
有五千.
兩千和一千的比喻.
其中一千那個.
雖然他做得不好.
不過他說得一個.
很真實的狀況.
就是什麼呢.
他說主人我知道.
你是個嚴厲的人.
沒有種的地方要收穫.
沒有播的地方要收穫.
我就害怕.
因為他知道.
主人是有一個絕對的權.
去對待他.
可以不要他.
於是他就害怕.
將錢收起來.
所以我們見得到.
兒子對父親是一種孝敬.
僕人對主人是一種敬畏.
說的是什麼呢.
這種是一種對父親和主人.
一種有份量的態度去回應.
原來神都說.
很期待著我們要對他.

$^{601}$給一種有份量.
有份量是價值.
是對等於神的價值.
那種回應給他.
所以你就要問.
今天神在你眼裡值多少.
神在你眼裡值多少.
五毛錢 五元.
你一分鐘.
你和神一分鐘的相處都很困難.
和神有靈修.
十分鐘都覺得很奢侈.
你是不是覺得和神.
談十分鐘都很辛苦呢.
你有沒有給上帝應得的價值呢.
這種是一種應有的態度.
但是以色列人是什麼態度呢.
看看怎麼說.
以色列人的態度.
是一種錯誤的態度.
他說了兩句話.
其中第六和第七節.
我們在何事上藐視你的名呢.
他說什麼.
我什麼時候藐視過你啊.
你想想.
他給不了應該有的態度對神.
他還說.
我什麼時候藐視過你啊.
你想想.
如果你和你的子女這樣說話.
你說幫我做點事.
什麼時候又不聽你的話.
你這樣聽下去.
你都覺得馬上火都來了.
他又說什麼.
我何事上使你不潔淨啊.
我什麼時候使你不潔淨啊.
原來說這群以色列人.
他們死不悔改.

$^{641}$不願意去回應.
神對他們提出的疑問.
神指出他們的問題.
他們不願意去承認他們有問題.
到底哪一種態度.
是我們面對神應該有的態度呢.
是不是應該要敬畏和尊重他呢.
還是說.
我什麼時候不認真啊.
我什麼時候沒有好好對待你啊.
曾經有人嘗試去玩一個遊戲.
他說.
怎樣能夠讓生活得到滿分呢.
他就說.
如果可以將這些的概念.
能夠呢.
態度呢.
這些事情呢.
能夠去化為ABCD去到Z呢.
就是說.
A等於1分的話.
B等於2分.
C等於3分.
Z等於26分的話.
那究竟怎樣才能夠.
令到你的生活能夠滿分呢.
有人就說.
Luck.
幸運就能夠令你的生活滿分.
但如果用這個字.
你會見到有多少分呢.
47分.
有些人又嘗試說.
既然是這樣.
我們試試用Money這個字.
Money這個字.
如果M是13分.
O是15分的話.
加起來會有多少分呢.
72分.

$^{681}$即是都能夠100分嗎.
不能夠.
有人說.
我們怎樣能夠令到生活能夠滿意呢.
滿分呢.
他說.
是不是靠我們的Hard work.
勤奮工作.
勤奮工作只佔98分.
是不是靠我們讀書多.
我們的知識就能夠令我們生活得到滿分呢.
96分.
是不是靠我們的領導能力.
做事出色.
能夠帶領一些人.
就能夠有滿分呢.
這裡說有89分.
似乎都拿不到滿分.
所以哪一個字才能夠拿到滿分呢.
就是Attitude.
Attitude.
A-T-T-I-T-U-D-E.
加起來就是100分.
當然這個是一個數字遊戲.
有些人這樣玩.
不過它帶出一個很重要的問題就是.
原來我們改變我們的態度.
才能夠改變我們的生活.
我們態度錯誤.
我們的生活就糟糕.
我們糟糕不是因為我們沒有錢.
不是因為我們沒有工作.
不是因為我們沒有學歷.
不是因為我們沒有家庭.
而是你的態度.
令到你不能夠有好的工作.
不能夠有好的學業.
不能夠有好的家庭.
往往是因為我們的態度.
決定了我們的高度.

$^{721}$我們的態度決定了我們的結果.
我們對神的態度.
都決定了我們能夠.
怎樣去經歷神更多呢.
當你還覺得自己很OK很行.
靠著你的學歷.
靠著你的勤奮.
靠著你的知識.
靠著你的領導的時候.
原來我們都是徒然.
因為沒有辦法真正能夠改變到.
我們和神的那種關係.
我們敬拜神最重要的期望.
就是你用什麼態度去敬拜神.
用什麼態度去敬拜神.
剛才我們說到.
有一群人在說話.
我什麼時候藐視你.
我什麼時候弄髒你的地方.
令你不潔.
究竟是什麼人這麼大膽呢.
你不要想究竟什麼人這麼大膽.
這樣和神說.
我哪有藐視你.
接著我們看看原來有兩群人.
兩群人.
這兩群人是怎樣呢.
看一看.
應該遮住了.
上面猜一下下面那四個字.
你們猜中了.
就是死不悔改.
哪群人死不悔改呢.
原來有兩群人這樣去.
說愛神其實不愛神.
正在給A4的敬拜神.
就是上良和下良.
大家認識這句話.
上良不正下良壞.
上良是什麼人呢.

$^{761}$原來上良就是祭司.
原來《馬拉基書》裡面說.
那些上良就是祭司.
第六節說你們這些藐視我們的祭司.
原來《馬拉基書》裡面說.
那些敬拜神為什麼這麼憤怒.
這麼大反感.
因為原來現在帶領一些人.
給A4的敬拜上帝的人.
就是那群祭司.
大家之前讀到利未記.
都看到很多關於祭司的事情.
你要知道.
祭司很高要求的.
他們要進入聖殿會幕裡面.
去獻祭.
會幕外面獻祭會裡面.
換陳設品等等.
每一個過程都很高要求.
只是祭司穿的那件衣服都有要求.
如果你記得這裡.
我們見過.
特別是大祭司的衣服裡面.
有胸袋又有腰帶.
又有衣服.
又有外袍又有內袍.
內袍下面有些鈴鈴噹噹.
因為整套服裝裡面表明了.
他們要用一種正確的態度去面對神.
是有要求的.
弟兄姊妹.
原來在說.
神對這群祭司能夠進入會幕.
無論是獻祭的.
無論是進入聖所的.
或者至聖所的.
都很有要求.
所以我們見到.
為什麼要讀利未記呢.
因為利未記其實就是一本祭司的手冊.

$^{801}$這本祭司手冊裡面說了很多.
關於做祭司裡面的吩咐.
要求.
有些什麼要做.
怎樣叫做聖潔.
怎樣叫做不潔.
當中你看到第一至五章.
特別說到很多不同的獻祭.
有凡祭,素祭,平安祭,贖罪祭和贖天祭.
我們說基本上在以色列裡面的獻祭只有五類.
就是這五大類.
其他你聽到那些所謂的瑤祭.
其實動作是代表什麼動作.
但真的祭的名字就是這五種.
其中贖天祭就是一種叫做賠償的祭.
等於一種這樣的祭.
所以不要緊.
但是就是說.
原來祭司他自己很有要求.
他獻祭物的安排是有要求的.
所以是不容易的.
但他們要怎樣呢.
既然這麼有要求的人.
所以他們事奉的態度就有要求.
但我們看到這裡說什麼呢.
這群這樣的祭司他們的態度出了問題.
他說什麼.
我們在何事上藐視你的名呢.
我們在何事上使你不潔淨呢.
後面解釋得很清楚.
為什麼他們會這麼大膽說話呢.
什麼時候藐視你什麼時候不潔淨呢.
因為原來他們說神是可藐視的.
而胡為的貢卓是可藐視的.
貢卓的意思就是指他們神的祭壇是可以藐視的.
就是說屬神的東西原來可以藐視.
輕視.
什麼叫輕視.
輕視就是說不是這麼重要.
不用這麼嚴謹.

$^{841}$不用這麼有要求.
只是回教會而已.
坐在那裡就行了.
回到教會就算了.
用不著說準不準時.
遲到不行嗎.
反正你都開門給我了.
有時候做得不好.
有什麼問題呢.
司復忘記了不要緊一定有人會做的.
有時候我們擔任一些領導的角色.
一些領導的角色都會有這種情況.
就是不認真.
那種藐視就是那種態度令到出問題.
怎麼藐視呢.
我們看看他們怎麼做.
因為原來我們剛才說的就是.
有很多不同的獻祭.
原來他們獻的祭物是有要求的.
我們看到有時候獻牛,獻羊,獻斑溝等等都好.
是有要求的.
你還沒記清楚說什麼嗎.
要沒有殘疾的.
沒有殘疾的.
初生的,投生的等等.
但是瑪雅基書說這群人獻什麼呢.
什麼呢.
黑暗的祭物.
他們將黑暗的祭物獻為祭物.
不算是惡嗎.
他們都不覺得有什麼問題.
但神就說這些都不是有問題.
他就將瘸腿的有病的獻上.
即是說腳有點跛.
即是說還在生病.
即是說可能你的羊已經病了.
即是說其實快要死了.
反正我的羊也快要死了.
獻給神又可以交差.
又獻羊給神.

$^{881}$反正羊也快要死了.
我又不會有財物的損失.
這樣豈不是兩全其美嗎.
那些以色列是這樣想的.
所以瘸腿有病的賣不出去的.
不值錢的獻給神.
又可以交差.
又可以令自己沒有任何損失.
但神說這些豈不是惡嗎.
神還說.
就算你將這些獻給省長.
剛才我們說到他們在以色列.
現在由省長管理.
即是說以色列這個地方最高級的省長.
你打算送給省長.
我送一隻跛的羊給你.
送一隻有病的牛給你.
你猜省長要不要.
不要.
你看報紙報道.
中國送了很多東西給美國.
他們上飛機之前都扔掉了.
全部都不拿.
因為怕有問題.
他們怕有問題.
一樣東西都不會帶回美國.
今天我們去到神的面前.
獻東西給神.
那些東西全部都有問題.
神怎能夠粗粗地要你的東西呢.
這個行為是有問題的.
這班作為祭司.
就是一班憲制的領袖.
因為他們才能夠進到會幕.
進到聖所.
他們這些憲制領袖都說沒問題.
那你想想那些百姓怎樣呢.
那些百姓說.
我真的有一隻跛的羊.
行不行啊.

$^{921}$那些祭司就說.
OK 照拿來吧.
沒把關嘛.
下面的人就說原來這麼低要求.
如果在教會裡面也是這樣.
如果教會那些領袖.
那些理事會的人.
那些人員個個都沒有要求的.
對事奉都是隨隨便便的.
接著那些其他事奉的人就問.
那些回來教會的人就問.
這樣行不行啊.
OK 隨隨便便行不行.
行.
穿拖鞋回來行不行.
行.
穿件內衣回來行不行.
行.
是不是.
都行的.
我奉獻的那些錢.
真是煤臭的.
行了 不要緊.
反正銀行收的.
你想想我們.
在香港封利是.
有沒有封利是.
香港封利是知道要做些什麼嗎.
要去銀行排隊的.
拿新銀紙.
你知道嗎.
舊一點都不拿出來.
要排隊拿新銀紙.
過年要拿最漂亮的銀紙.
好的.
炒了都不行.
是不是.
我們封利是都這麼重要.
我們奉獻.
看看袋有多少.

$^{961}$給你.
是不是這樣.
當然我們一會兒奉獻的時候.
你炒的收不收.
我不知道.
問問那些理事.
我說不了事.
我不是這方面的領袖.
我是港獨.
是不是.
最好不要.
多謝你這麼說.
最好不要.
我們自己.
我們教子女都是這樣.
我們現在兒童崇拜.
有奉獻環節.
兒童崇拜有奉獻環節.
我們每一次都.
之前我們要.
一兩天前找找.
有沒有錢給他們可以做奉獻.
我們要之前找到.
到那天給他們.
你袋好.
讓他們知道.
這個是我們獻給神的.
要做足準備.
去迪士尼都做準備.
為什麼去獻祭.
去到神面前.
去敬拜他.
不做足準備.
是不是迪士尼.
比回教會重要.
真的.
要認真想想.
所以.
下面那群百姓是怎樣.
看這幅圖就知道.

$^{1001}$這幅圖就是剛才那隻羊.
是AI幫我做的.
那隻羊是瞎了.
你可能遠處看不到.
那隻羊是瘸了.
下面那群人是怎樣.
就是這樣.
這種叫什麼.
恥之以備.
你想想.
神就說.
十二至十三節.
就說那些下梁.
下梁那些百姓是怎樣.
見到上梁是這樣.
下梁是什麼.
原來是可以褻瀆的.
神的名是可以褻瀆的.
所以他們說什麼.
神就說他們說什麼.
主的公卓是不潔淨的.
公卓上的果子和食物.
是可藐視的.
就是說.
可以骯髒.
可以藐視.
可以隨便.
甚至說什麼.
這些事何等煩鎖.
用什麼這麼多要求.
真的很煩.
是不是.
不要妨礙我.
由得我做.
不要煩到我.
最後是什麼.
恥之以備.
我都做不到.
我都不會做這些.
誰最熟.

$^{1041}$做恥之以備.
馬上對著神.
是啊.
我們以為是這樣.
所以我們以前在香港.
我們做青少年的時候.
青少年崇拜.
就在學校禮堂裡面.
很多時候崇拜之前.
青少年就帶了一杯.
二杯飲料.
然後崇拜之前就隨便擺.
那些就擺在崇拜台前.
擺三杯.
差點以為可以鞠躬.
擺到整個崇拜.
我們敬拜神.
擺到崇拜台前這麼髒.
每一次我們都要說.
你快點收拾回去.
搬開現在崇拜.
我們都叫青少年這樣做.
我們做大人.
是不是應該有好的示範.
怎樣認真去面對我們的敬拜.
前兩天有機會去到.
另外一間教會.
去教會門口原來有這樣的字.
教會門口說.
Cell phones off.
No eating or drinking.
原來他們左右兩道門.
兩道門都寫這些大大的字.
就是什麼呢.
記得關掉手機.
因為崇拜要專心.
原來崇拜要有一個好的環境.
給自己敬拜.
給別人敬拜.
不要騷擾了別人.

$^{1081}$關掉手機不是歧視別人.
而是我們做好我們的準備.
做好我們敬拜的準備.
所以我提醒大家.
如果你有時候不懂關掉手機.
找些年輕人幫你.
有時候我們上了年紀.
不懂怎樣關掉手機.
讓它不要響.
有時候不懂.
你找些年輕人.
年輕人.
你是年輕人.
你也是.
年輕人幫他做.
做完崇拜後.
年輕人幫我弄好它.
這樣就可以了.
我們主動找人幫我們.
解決我們的問題.
令我們有好的敬拜.
有好的環境.
所以他們說.
不可以吃東西不可以喝東西.
因為什麼.
你想想.
如果在崇拜你又吃又喝.
開大餐也不是辦法.
因為我們在專注敬拜神.
所以說.
這班人就是這樣的狀態.
他們對神是褻瀆的.
認為神很藐視.
認為神很煩.
認為神可以恃之以備.
有這種心態的人.
你猜他面對神的時候.
會怎樣呢.
有一個圖案.
一個新的公仔.

$^{1121}$不知道大家見過沒有.
給大家看看.
你猜這個是什麼呢.
這個原來是.
在香港最受歡迎的.
政府的宣傳物和公仔.
原來最受歡迎.
我上網看.
知道這個是什麼嗎.
不知道嗎.
這麼大顆也不知道嗎.
就是葡萄.
你們沒見過葡萄嗎.
M是M字.
這是葡萄.
葡萄.
他叫葡萄.
提防騙子.
這個宣傳.
葡萄提防騙子.
因為原來有很多騙子.
這個世界有很多騙子.
很多騙子.
神說以色列的百姓裡.
也有一群騙子.
這群騙子有什麼行為呢.
十三十四節說.
他們把搶來的.
即是賊贓.
渠腿的.
即是爛的.
不完整的.
有病的.
即是快要死的.
都拿來獻上為祭.
這麼面熟.
是誰搞的.
就是上梁.
上梁嘛.
現在下梁又怎樣.

$^{1161}$再這樣.
搶來的也可以.
我講清楚.
我們教會是不收賊贓的.
不要拿賊贓來奉獻.
所以我們檢查不到.
搶來的.
渠腿的.
有病的.
都拿來獻祭.
然後神也問一個問題.
我怎樣能夠歡歡喜喜.
怎樣能夠月立.
神收禮物也要開心.
你拿這些來獻給神.
神說我怎樣能夠歡歡喜喜.
因為你們是什麼人呢.
十四節說.
是行鬼詐的人.
原來這群就是提子.
不是提子.
提子裡面的子.
就是這群騙子.
他們向神說謊.
欺騙神.
以為騙得了神.
神說這是可咒詛的.
他們的群促裡有公羊.
這群人不是沒有錢.
這群人是有羊的.
這群百姓有羊.
而且有公羊.
是符合要求的.
他有.
他不是窮到沒有那些.
神不會窮到沒有那些.
要強行要來.
但有些人是真的有.
但他收起來.
而且把殘疾的獻給主.

$^{1201}$這就是問題.
各位姐妹.
今天我們想想.
A4敬拜的時候.
究竟我們今天.
在敬拜崇拜的過程裡.
我們所獻上的.
有多少是A4.
有多少是頂級真課.
要很認真想.
因為當我們獻上的時候.
只會有兩個結局.
第一個結局是什麼呢.
人被拒絕.
第九第十節.
他說現在我要勸你們.
要懇求上帝.
好讓他施恩給我們.
這是既出於你們的手.
他豈會抬舉你們任何人呢.
神怎會抬舉你們呢.
本來神願意在糞堆裡.
提拔窮乏人.
但如果我們今天這樣對神.
神怎會能夠提拔我們呢.
第十節說.
甚願你們中間.
有人把電的門關上.
免得你們徒然在我壇上燒火.
神說如果我們在神面前獻上的.
這麼隨便.
這麼不認真.
拿些雜物.
拿些爛貨.
拿些不要的垃圾去獻給神的話.
神說我寧願你們教會關門.
教會關門.
大家就連獻的機會都沒有.
總比你羞辱主名好.
第十節說.

$^{1241}$今天很多教會在世界裡.
美國也有很多教會已經關門.
關門就是因為越來越少人.
認真去對待上帝.
以致一間教會越來越少人.
越來越不肯奉獻.
越來越不肯認真對待他們的事奉.
越來越不肯回到教會裡去敬拜神.
教會人凋零.
教會沒有財政.
沒有能力.
教會關門是遲早的事.
神說如果這個地方都沒有人願意敬拜祂.
神說寧願將聖殿的門關上.
總比你獻上些雜物好.
總比你將一些不應該獻上祭壇的東西獻上去好.
所以關門大吉是其中一個神說會有的出路.
會有的結果.
以色列人都要面對這種上帝的審判.
當他們不認真去對待神的殿.
就算他們已經重建了聖殿.
第二聖殿都要被毀.
都要關門大吉.
但同時還有另一個結果.
神會被尊大.
如果人認真去對待神.
甚至你今天不去珍惜神這個敬拜的地方.
總有人會珍惜.
神說今天這裡的人不欣賞上帝的福音.
福音要傳給其他人.
以色列的百姓拒絕上帝的憐憫和愛.
神就讓耶穌基督將福音臨到萬邦.
從日出之地到日落之處 我的名在列國中必被尊為大.
你不尊上帝為大.
其他地方會尊上帝的名為大.
在那裡認識上帝.
知道上帝是我們的父.
知道上帝是我們的主.
他們就會奉神的名去燒香.
去獻最好的貢物獻給神.

$^{1281}$你知道神在那些地方被尊為大.
剛才說了很多教會因為越來越少人.
越來越多人不認真去對待神而關閉.
但同時也有人認真去對待神的話.
在佛羅里達州能夠將戲院賣掉成為教會.
讓這個地方有多一間可以榮耀神的教會.
究竟結果是誰呢?.
是教會倒閉越來越多還是開越來越多教會呢?.
其實結果取決於我們.
如何去對待我們對神的敬拜.
對神的敬拜就是我們給多少重量上帝.
你得上帝是最珍貴的.
你獻上最寶貴的.
可能是你的時間 可能是你的金錢.
可能是你的知識 可能是你的才幹.
但不是懶得獻給神.
那就是給一個應該有的分量上帝.
神就閱立 神就使用.
神就讓更多人得到福氣.
當你輕輕忽忽.
神就寧願不要閱立你的東西.
所以福音能否傳到世界各地.
其實神會自己彰顯祂的命.
神會做的事.
但我們做錯事 神就會自己關了他的門.
最後總結.
如何去對待我們的敬拜.
什麼叫做不是A課的敬拜.
是真課的敬拜呢?.
就是一種以神為中心的敬拜.
當我們以神為我們生命的中心.
我們就能夠知道.
無論何事何物何人.
我們都以神為首.
不是看人去決定是否回教會.
而是看神.
我回教會不是固然見到人.
但我重要的是見到神.
我事奉的不是有人看到我做就做.
沒有人看到我就不做.

$^{1321}$因為我們對著的事奉就是對神事奉.
神是無處不在.
看到我們獻上什麼.
之前有個事奉人員的手冊.
崇拜的事奉人員.
要認真去對待我們在神面前的崇拜的事奉.
因為事奉這個字就是敬拜.
我們的事奉就是敬拜.
我們的敬拜就是在我們事奉中顯現出來.
所以我們不要想著時間到才埋位.
我們五分鐘十分鐘前已經有應該有的預備.
翻崇拜也是這樣.
所以讓我們重新檢視.
特別是每一個教會的事奉人員.
都想想我們的敬拜或事奉的態度.
因為教會已經出了一本事奉手冊.
但有時我們都不太根本手冊做.
我做這本手冊時講了很多每一個崗位都有它的要求.
它的角色其實有很多東西要留心.
無論是你的服飾.
無論是你的時間觀念.
無論你在位置有沒有做好你俠如其分的.
無論是接待 施事等等.
有些人說施事是否出來祈禱就行了.
每次都以為看到他出來就出來祈禱.
其實施事有很多事情要做.
一個認真的施事.
由你一進到門口.
他已經監視著每一個人.
招呼每一個人.
安頓好每一個人.
坐好.
是很重要的.
今日我盼望大家重新檢視我們整個崇拜.
能否有一個好的崇拜獻給神.
以神為中心.
由崇拜開始的預備.
講道 回應.
有主餐 有商交 差遣.
每一個都是以神為中心的聚會.

$^{1361}$求主幫助我們.
我們一起祈禱.
天父我們求你幫助我們.
神我知道你很不喜歡我們交A課給你.
甚至乎我們很虧欠地說.
有時我們連A課都不如.
我們交B課 C課.
甚至乎是劣質課給神.
求你真的讓我們去審查自己.
不要再自高自大.
不要以為已經做得最好.
而讓我們去檢視.
有沒有時刻以神為我們的中心.
以神為我們敬拜和事奉的中心.
以致我們知道.
我們唯有將最好的獻給你.
才能夠討你喜悅.
求主你幫助我們在家的.
現場的.
每一個都能夠去思考到我們.
怎樣有一種正確的敬拜的態度.
願主你幫助我們.
聽我們祈禱.
奉耶穌的名字而求.
阿們.
你輕輕舉起你的手告訴耶穌.
其實我都願意被你所愛.
我都想更深的經歷你的愛.
我們一起祈禱.
親愛的天父我們去仰望你.
我們知道原來很多人.
都不知道原來神很愛我們.
大主我們多謝你.
讓我們知道原來你很愛我們.
求你讓我們常常去睜開雙眼.
回望你的愛.
以致有一天我們能夠去捉緊你.
能夠相信你.
也讓我們已經相信你的人.
我們更加堅定的愛你.

$^{1401}$願意對你的愛去作出回應.
也求主你幫助我們.
藉著今天母親節的日子.
我們想起母親對我們的點點滴滴.
這份愛的時候.
我們願意去表達這份愛的感動.
無論打電話給她.
無論打個視像給她.
或者親自跟她說一些話.
或者請她吃一頓飯.
讓我們真正珍惜我們的母親.
因為她同樣是上帝所差派來愛我們的.
願意我們能夠回應這份上帝給我們的愛.
和這份你給我們母親的愛.
我們願意將我們的祈禱.
都交在天父的手裡.
聽我們祈禱奉耶穌的名字而求.
\newpage



\section{提摩太後書 1:1-4:22}
\label{sec:032138}
\textbf{牧者手冊:提摩太前後書 (2)}~~~~劉少康牧師
\newline
\newline
連結: \href{http://www.church.com.hk/acms/content.asp?site=cdc\&op=show\&type=product\&code=032138\&layout=sermon}{\texttt{032138}} ~~~~ 語音日期: 2026-05-13
\newline
\newline
\hyperref[sec:032137]{\small{< < < PREV SERMON < < <}}
~
\hyperref[sec:index_chronic]{\small{[返順時目]}}
~
\hyperref[sec:032139]{\small{> > > NEXT SERMON > > >}}
\newline
\newline
提摩太後書 1:1-4:22
\newline
\begin{longtable}{cl}
\hline
\hline
章節 & 經文 (和合本修訂版)\\
\hline
1:1 & \begin{tabularx}{0.7\textwidth}{X} 奉神旨意，按照基督耶穌裡所應許的生命，作基督耶穌使徒的保羅， \end{tabularx} \\ \\ \relax
1:2 & \begin{tabularx}{0.7\textwidth}{X} 寫信給我親愛的兒子提摩太。願恩惠、憐憫、平安從父神和我們的主基督耶穌歸給你！ \end{tabularx} \\ \\ \relax
1:3 & \begin{tabularx}{0.7\textwidth}{X} 我感謝神，就是我接續祖先用純潔的良心所事奉的神，在祈禱中晝夜不停地想念你。 \end{tabularx} \\ \\ \relax
1:4 & \begin{tabularx}{0.7\textwidth}{X} 我一想起你的眼淚，就急切想見你，好讓我滿心快樂。 \end{tabularx} \\ \\ \relax
1:5 & \begin{tabularx}{0.7\textwidth}{X} 我記得你無偽的信心，這信心先存在你外祖母羅以和你母親友妮基的心裡，我深信也存在你的心裡。 \end{tabularx} \\ \\ \relax
1:6 & \begin{tabularx}{0.7\textwidth}{X} 為這緣故，我提醒你要把神藉著我按手所給你的恩賜再如火挑旺起來。 \end{tabularx} \\ \\ \relax
1:7 & \begin{tabularx}{0.7\textwidth}{X} 因為神賜給我們的不是膽怯的心，而是剛強、仁愛、自制的心。 \end{tabularx} \\ \\ \relax
1:8 & \begin{tabularx}{0.7\textwidth}{X} 所以，不要以給我們的主作見證為恥，也不要以我這為主被囚的為恥；總要靠著神的大能，與我為福音同受苦難。 \end{tabularx} \\ \\ \relax
1:9 & \begin{tabularx}{0.7\textwidth}{X} 神救了我們，以聖召召我們，不是按我們的行為，而是按他的旨意和恩典；這恩典是萬古之先在基督耶穌裡賜給我們的， \end{tabularx} \\ \\ \relax
1:10 & \begin{tabularx}{0.7\textwidth}{X} 但如今藉著我們的救主基督耶穌的顯現已經表明出來；他把死廢去，藉著福音，將不朽的生命彰顯出來。 \end{tabularx} \\ \\ \relax
1:11 & \begin{tabularx}{0.7\textwidth}{X} 我為這福音奉派作傳道，作使徒，作教師。 \end{tabularx} \\ \\ \relax
1:12 & \begin{tabularx}{0.7\textwidth}{X} 為這緣故，我也受這些苦難。然而，我不以為恥，因為我知道我所信的是誰，也深信他能保全他所交託我的，直到那日。 \end{tabularx} \\ \\ \relax
1:13 & \begin{tabularx}{0.7\textwidth}{X} 你從我聽到那健全的言論，要用在基督耶穌裡的信心和愛心常常守著，作為規範。 \end{tabularx} \\ \\ \relax
1:14 & \begin{tabularx}{0.7\textwidth}{X} 你要靠著那住在我們裡面的聖靈，牢牢守住所交託給你那美好的事。 \end{tabularx} \\ \\ \relax
1:15 & \begin{tabularx}{0.7\textwidth}{X} 你知道，所有在亞細亞的人都離棄了我，其中有腓吉路和黑摩其尼。 \end{tabularx} \\ \\ \relax
1:16 & \begin{tabularx}{0.7\textwidth}{X} 願主憐憫阿尼色弗一家的人，因為他屢次令我欣慰。他不以我的鐵鏈為恥， \end{tabularx} \\ \\ \relax
1:17 & \begin{tabularx}{0.7\textwidth}{X} 反而一到羅馬就急切尋找我，並且找到了。 \end{tabularx} \\ \\ \relax
1:18 & \begin{tabularx}{0.7\textwidth}{X} 願主使他在那日能蒙主的憐憫。他在以弗所怎樣多服事我，你是清楚知道的。 \end{tabularx} \\ \\ \relax
2:1 & \begin{tabularx}{0.7\textwidth}{X} 我兒啊，你要在基督耶穌的恩典上剛強起來。 \end{tabularx} \\ \\ \relax
2:2 & \begin{tabularx}{0.7\textwidth}{X} 你在許多見證人面前聽見我所教導的，也要交託給那忠心而又能教導別人的人。 \end{tabularx} \\ \\ \relax
2:3 & \begin{tabularx}{0.7\textwidth}{X} 你要和我同受苦難，作基督耶穌的精兵。 \end{tabularx} \\ \\ \relax
2:4 & \begin{tabularx}{0.7\textwidth}{X} 凡當兵的，不讓世務纏身，好使那招他當兵的人喜悅。 \end{tabularx} \\ \\ \relax
2:5 & \begin{tabularx}{0.7\textwidth}{X} 運動員在比賽的時候，不按規則就不能得冠冕。 \end{tabularx} \\ \\ \relax
2:6 & \begin{tabularx}{0.7\textwidth}{X} 勤勞的農夫理當先得糧食。 \end{tabularx} \\ \\ \relax
2:7 & \begin{tabularx}{0.7\textwidth}{X} 我所說的話，你要考慮，因為主必在凡事上給你聰明。 \end{tabularx} \\ \\ \relax
2:8 & \begin{tabularx}{0.7\textwidth}{X} 要記得耶穌基督，他是大衛的後裔，從死人中復活；這就是我所傳的福音。 \end{tabularx} \\ \\ \relax
2:9 & \begin{tabularx}{0.7\textwidth}{X} 我為這福音受苦難，甚至像犯人一樣被捆綁，然而神的話沒有被捆綁。 \end{tabularx} \\ \\ \relax
2:10 & \begin{tabularx}{0.7\textwidth}{X} 所以，我為了選民事事忍耐，為使他們也能得到那在基督耶穌裡的救恩和永遠的榮耀。 \end{tabularx} \\ \\ \relax
2:11 & \begin{tabularx}{0.7\textwidth}{X} 這話是可信的：我們若與基督同死，也必與他同活； \end{tabularx} \\ \\ \relax
2:12 & \begin{tabularx}{0.7\textwidth}{X} 我們若忍耐到底，也必和他一同作王。我們若不認他，他也必不認我們； \end{tabularx} \\ \\ \relax
2:13 & \begin{tabularx}{0.7\textwidth}{X} 我們縱然失信，他仍是可信的，因為他不能否認自己。 \end{tabularx} \\ \\ \relax
2:14 & \begin{tabularx}{0.7\textwidth}{X} 你要向眾人提醒這些事，在神面前囑咐他們不可在言詞上爭辯；這是沒有益處的，只能傷害聽的人。 \end{tabularx} \\ \\ \relax
2:15 & \begin{tabularx}{0.7\textwidth}{X} 你當竭力在神面前作一個經得起考驗、無愧的工人，按著正意講解真理的話。 \end{tabularx} \\ \\ \relax
2:16 & \begin{tabularx}{0.7\textwidth}{X} 要遠避世俗的空談，因為這等空談會使人進到更不敬虔的地步。 \end{tabularx} \\ \\ \relax
2:17 & \begin{tabularx}{0.7\textwidth}{X} 他們的話如同毒瘡越爛越大；其中有許米乃和腓理徒， \end{tabularx} \\ \\ \relax
2:18 & \begin{tabularx}{0.7\textwidth}{X} 他們偏離了真理，說復活的事已過去，敗壞了好些人的信心。 \end{tabularx} \\ \\ \relax
2:19 & \begin{tabularx}{0.7\textwidth}{X} 然而，神堅固的根基屹立不移；上面有這印記說：「主認得他自己的人」，又說：「凡稱呼主名的人總要離開不義。」 \end{tabularx} \\ \\ \relax
2:20 & \begin{tabularx}{0.7\textwidth}{X} 大戶人家不但有金器銀器，也有木器瓦器；有作為貴重之用的，有作為卑賤之用的。 \end{tabularx} \\ \\ \relax
2:21 & \begin{tabularx}{0.7\textwidth}{X} 人若自潔，脫離卑賤的事，必成為貴重的器皿，成為聖潔，合乎主用，預備行各樣的善事。 \end{tabularx} \\ \\ \relax
2:22 & \begin{tabularx}{0.7\textwidth}{X} 你要逃避年輕人的私慾，同那以純潔的心求告主的人追求公義、信實、仁愛、和平。 \end{tabularx} \\ \\ \relax
2:23 & \begin{tabularx}{0.7\textwidth}{X} 但要棄絕那愚拙無知的辯論，因為你知道這等事只會引起爭辯。 \end{tabularx} \\ \\ \relax
2:24 & \begin{tabularx}{0.7\textwidth}{X} 主的僕人不可爭辯，只要溫和待人，善於教導，恆心忍耐， \end{tabularx} \\ \\ \relax
2:25 & \begin{tabularx}{0.7\textwidth}{X} 用溫柔勸導反對的人。也許神會給他們悔改的心能明白真理， \end{tabularx} \\ \\ \relax
2:26 & \begin{tabularx}{0.7\textwidth}{X} 讓他們這些已被魔鬼擄去順從他詭計的人能醒悟過來，脫離他的羅網。 \end{tabularx} \\ \\ \relax
3:1 & \begin{tabularx}{0.7\textwidth}{X} 你該知道，末世必有艱難的日子來到。 \end{tabularx} \\ \\ \relax
3:2 & \begin{tabularx}{0.7\textwidth}{X} 那時人會專愛自己，貪愛錢財，自誇，狂傲，毀謗，違背父母，忘恩負義，心不聖潔， \end{tabularx} \\ \\ \relax
3:3 & \begin{tabularx}{0.7\textwidth}{X} 沒有親情，抗拒和解，好說讒言，不能節制，性情兇暴，不愛良善， \end{tabularx} \\ \\ \relax
3:4 & \begin{tabularx}{0.7\textwidth}{X} 賣主賣友，任意妄為，自高自大，愛好宴樂，不愛神， \end{tabularx} \\ \\ \relax
3:5 & \begin{tabularx}{0.7\textwidth}{X} 有敬虔的外貌，卻背棄了敬虔的實質，這等人你要避開。 \end{tabularx} \\ \\ \relax
3:6 & \begin{tabularx}{0.7\textwidth}{X} 他們當中有人潛入別人家裡，操縱無知的婦女；這些婦女被罪惡壓制，被各樣的私慾引誘， \end{tabularx} \\ \\ \relax
3:7 & \begin{tabularx}{0.7\textwidth}{X} 雖然常常學習，終久無法達到明白真理的地步。 \end{tabularx} \\ \\ \relax
3:8 & \begin{tabularx}{0.7\textwidth}{X} 從前雅尼和佯庇怎樣反對摩西，這等人也怎樣抵擋真理；他們的心地敗壞，信仰經不起考驗。 \end{tabularx} \\ \\ \relax
3:9 & \begin{tabularx}{0.7\textwidth}{X} 然而，他們沒有進步，因為他們的愚昧必在眾人面前顯露出來，像那兩人一樣。 \end{tabularx} \\ \\ \relax
3:10 & \begin{tabularx}{0.7\textwidth}{X} 但你已經追隨了我的教導、行為、志向、信心、寬容、愛心、忍耐， \end{tabularx} \\ \\ \relax
3:11 & \begin{tabularx}{0.7\textwidth}{X} 以及我在安提阿、以哥念、路司得所遭遇的迫害和苦難。我忍受了何等的迫害！但從這一切苦難中，主都把我救了出來。 \end{tabularx} \\ \\ \relax
3:12 & \begin{tabularx}{0.7\textwidth}{X} 其實，凡立志在基督耶穌裡敬虔度日的，也都將受迫害。 \end{tabularx} \\ \\ \relax
3:13 & \begin{tabularx}{0.7\textwidth}{X} 只是作惡的和騙人的將變本加厲，迷惑人也被人迷惑。 \end{tabularx} \\ \\ \relax
3:14 & \begin{tabularx}{0.7\textwidth}{X} 至於你，你要持守所學習的和所確信的，因為你知道是跟誰學的， \end{tabularx} \\ \\ \relax
3:15 & \begin{tabularx}{0.7\textwidth}{X} 並且知道你從小明白聖經，這聖經能使你因在基督耶穌裡的信有得救的智慧。 \end{tabularx} \\ \\ \relax
3:16 & \begin{tabularx}{0.7\textwidth}{X} 聖經都是神所默示的，於教訓、督責、使人歸正、教導人學義都是有益的， \end{tabularx} \\ \\ \relax
3:17 & \begin{tabularx}{0.7\textwidth}{X} 叫屬神的人得以完全，預備行各樣的善事。 \end{tabularx} \\ \\ \relax
4:1 & \begin{tabularx}{0.7\textwidth}{X} 我在神面前，並在將來審判活人死人的基督耶穌面前，憑著他的顯現和他的國度鄭重地勸戒你： \end{tabularx} \\ \\ \relax
4:2 & \begin{tabularx}{0.7\textwidth}{X} 務要傳道；無論得時不得時總要專心，並以百般的忍耐和各樣的教導責備人，警戒人，勸勉人。 \end{tabularx} \\ \\ \relax
4:3 & \begin{tabularx}{0.7\textwidth}{X} 因為時候將到，那時人會厭煩健全的教導，耳朵發癢，就隨心所欲地增添好些教師， \end{tabularx} \\ \\ \relax
4:4 & \begin{tabularx}{0.7\textwidth}{X} 並且掩耳不聽真理，偏向無稽的傳說。 \end{tabularx} \\ \\ \relax
4:5 & \begin{tabularx}{0.7\textwidth}{X} 至於你，凡事要謹慎，忍受苦難，做傳福音的工作，盡你的職分。 \end{tabularx} \\ \\ \relax
4:6 & \begin{tabularx}{0.7\textwidth}{X} 至於我，我已經被澆獻，離世的時候到了。 \end{tabularx} \\ \\ \relax
4:7 & \begin{tabularx}{0.7\textwidth}{X} 那美好的仗我已經打過了，當跑的路我已經跑盡了，該信的道我已經守住了。 \end{tabularx} \\ \\ \relax
4:8 & \begin{tabularx}{0.7\textwidth}{X} 從此以後，有公義的冠冕為我存留，就是按著公義審判的主到了那日要賜給我的；不但賜給我，也賜給凡愛慕他顯現的人。 \end{tabularx} \\ \\ \relax
4:9 & \begin{tabularx}{0.7\textwidth}{X} 你要趕緊到我這裡來。 \end{tabularx} \\ \\ \relax
4:10 & \begin{tabularx}{0.7\textwidth}{X} 因為底馬貪愛現今的世界，已經離棄我，往帖撒羅尼迦去了；革勒士往加拉太去；提多往撻馬太去； \end{tabularx} \\ \\ \relax
4:11 & \begin{tabularx}{0.7\textwidth}{X} 只有路加在我這裡。你來的時候把馬可帶來，因為他在服事上於我有益。 \end{tabularx} \\ \\ \relax
4:12 & \begin{tabularx}{0.7\textwidth}{X} 我已經打發推基古往以弗所去。 \end{tabularx} \\ \\ \relax
4:13 & \begin{tabularx}{0.7\textwidth}{X} 我在特羅亞留給加布的那件外衣，你來的時候要帶來，那些書也帶來，特別是那幾卷羊皮的書。 \end{tabularx} \\ \\ \relax
4:14 & \begin{tabularx}{0.7\textwidth}{X} 銅匠亞歷山大多方害我；主必照他所行的報應他。 \end{tabularx} \\ \\ \relax
4:15 & \begin{tabularx}{0.7\textwidth}{X} 你也要防備他，因為他極力抗拒我們的話。 \end{tabularx} \\ \\ \relax
4:16 & \begin{tabularx}{0.7\textwidth}{X} 我初次上訴時，沒有人前來幫助，竟都離棄了我，但願這罪不歸在他們身上。 \end{tabularx} \\ \\ \relax
4:17 & \begin{tabularx}{0.7\textwidth}{X} 惟有主站在我身邊，加給我力量，使我能把福音完整地傳開，讓所有的外邦人都聽見；我也從獅子口裡被救出來。 \end{tabularx} \\ \\ \relax
4:18 & \begin{tabularx}{0.7\textwidth}{X} 主必救我脫離一切的兇惡，也必救我進他的天國。願榮耀歸給他，直到永永遠遠。阿們！ \end{tabularx} \\ \\ \relax
4:19 & \begin{tabularx}{0.7\textwidth}{X} 請向百基拉、亞居拉和阿尼色弗一家的人問安。 \end{tabularx} \\ \\ \relax
4:20 & \begin{tabularx}{0.7\textwidth}{X} 以拉都在哥林多住下了。特羅非摩病了，我把他留在米利都。 \end{tabularx} \\ \\ \relax
4:21 & \begin{tabularx}{0.7\textwidth}{X} 你要趕緊在冬天以前到我這裡來。友布羅、布田、利奴、革老底亞和眾弟兄都向你問安。 \end{tabularx} \\ \\ \relax
4:22 & \begin{tabularx}{0.7\textwidth}{X} 願主與你的靈同在！願恩惠與你們同在！ \end{tabularx} \\ \\
[1ex]
\hline
\hline
\end{longtable}
$^{1}$給姐妹平安.
我用一次的時間介紹提摩太后書.
提摩太前書大概是保羅在羅馬坐牢.
第一次在羅馬坐牢獲釋之後.
他就寫了提摩太前書.
然後很稍後的時間.
因此有很多研究聖經的人就猜.
提摩太后書.
差不多是保羅臨死之前的絕筆.
所以有人說提摩太后書.
好像是保羅的遺囑一樣.
大家都知道.
保羅是被一個很殘酷的羅馬皇帝.
尼羅尼祿王.
將他斬頭而死的.
所以他是畏主殉道的.
而在他被斬首之前.
他最後寫的一本書就是提摩太后書.
所以提摩太后書充滿了.
保羅對這個年輕牧者的一份叮嚀.
好像對他有很多話要說.
言之不盡.
「我有太多話要勸你」.
「我有太多話要提醒你」.
「我有太多話要鼓勵你」.
好像一個老牧人.
他知道自己的日子不多了.
所以就傾盡他的心.
全都倒出來.
告訴他的接班人.
他最愛的學生.
這個年輕的提摩太子.
所以提摩太后書充滿了.
很溫暖.
保羅一個很深的提醒教導.
在第一章提摩太后書.
一個月有四章.
第一章保羅提到.
他說「當我為你按手的時候.
我就將以恩賜賜給你」.

$^{41}$這個恩賜在你的心裡面.
你時時要眺望它.
不要等它熄滅了.
要點著它.
因為我按納牧師有四十多年.
有時都想不起當初按納的時候是怎樣.
因為當初我按納牧師的時候.
是張有光牧師.
九龍城浸信會做了很久的主任牧師.
他來訓練我的.
新牧師要站起來聽訓練.
張牧師的訓練對我很深刻印象.
他說你做牧師.
所有人都看著你.
所以你要很小心.
我記住這個牧者對我的勸勉.
所以保羅說.
我按手的時候給你的恩賜.
你要如火一樣再一次燒著.
不要熄火.
千萬不要熄著.
我們時時要火熱的愛主.
火熱的來去事奉.
各位弟子們.
你有沒有有些時候.
你覺得好像自己都凍了.
沒錯.
我們基督徒的旅程裡面.
我們有時會高峰.
有時會低谷.
有時會高峰.
有時會低谷.
我們說去參加培靈會.
就好像打了一支興奮劑.
哇!谷一谷興奮起來.
但很快又謝了.
求神給我們一個持久的愛.
耶穌在比啟示錄七間教會的信裡面.
第一封已忽所教會.
就是你不要熄了起初的愛.

$^{81}$你剛剛信耶穌的時候.
哇!你多火熱.
你剛剛信耶穌的時候.
哇!你衝得多厲害.
千萬不要失去了起初的愛心.
求神今日都時時提醒我們.
讓我們對主的愛是火熱的.
是拉著的.
是眺望的.
幾時都這麼旺盛的.
然後在提摩太后書第二章.
就提到要做光強的僕人.
要做一個光強的僕人.
就是我們要做耶穌基督的精兵.
神的道是不被捆綁的.
而且我們要做無愧的工人.
特別我在浸會神學院讀.
浸信會神學院的院訓.
就是你當竭力在神面前.
做無愧的工人按著正義去分解真理的讀.
我相信有不少神學院.
都是用這句來做校訓的.
因為讓我們知道.
我們要成為一個沒有慚愧的工人.
神給我有一個積分.
就按著正義,正確的意思.
按著那把切鑽石的那把刀.
有一把很鋒利的刀.
那些鑽石割下去.
你千萬不要割錯.
不然那麼大顆鑽石給你切成一半.
那就沒什麼價值了.
你要把表面的石頭砍掉.
然後一顆很大,很光芒的鑽石就露出來了.
所以我們要按正義去分解真理的讀.
那我們就是一個無愧的工人.
成為我們神學院的院訓了.
然後我們要做貴重的器皿.
因為在每一個人家裡面.
在大戶人家有金器,銀器,木器,瓦器.

$^{121}$所有的器皿有各類不同.
最重要的是什麼?乾淨.
如果你去別人家裡.
別人說我拿一隻杯來給你倒杯水,倒杯茶給你.
他就打開他的飾櫃.
他在巴黎那裡買了一隻杯.
原來放了很久都沒有洗.
那你敢不敢喝?不用了.
你隨便拿一隻玻璃杯,拿個瓦杯出來.
倒水給我喝就行了.
倒茶,倒水給別人喝.
最重要是那個器皿乾淨.
不是那個器皿貴不貴重.
一隻金杯來的,一隻銀杯來的.
我不在乎,我不要緊.
最重要是乾淨.
你找一隻普通的瓷器杯,玻璃杯給我喝水就行了.
所以我們要做貴重的器皿.
做一個潔淨的器皿.
被神來使用.
這個同樣是今天我們很適用的一個訊息.
求神上上潔淨我們.
讓我們成為神手中一個被神使用的器皿.
貴重的器皿.
好,然後第三章就說到面對末日的警方.
在末世的時候危險的日子來到了.
必有很多人會有不同的表現.
我曾經說過一篇道.
就叫做「三國愛,五國不愛」.
三國愛,愛什麼呢?.
愛艷樂.
愛自己.
愛錢財.
今天就是這個世代.
我掛念自己.
我最威,我最叻.
掛念著錢.
我要賺多一點,賺多一點.
第一桶金,第二桶金,第三桶金.
我要去玩.

$^{161}$我要享受.
我要去很多地方旅行.
我要吃盡天下的美食.
有五個不愛.
不愛神.
不愛公義.
不愛正直.
還有不愛什麼?.
不愛良善.
不愛親情,沒有親情的.
哇,子女可以殺父母.
父母可以殺子女.
現在你會聽見這樣的新聞.
令你毛骨悚然.
沒有親情.
不愛神,不饒恕人.
今天這是一個很凶惡的世代.
沒有饒恕的.
於是我們就看見.
保羅老早已經預告了.
在末後的日子.
有這樣危險的日子出現.
然後保羅就提醒我們.
聖經都是神所默示的.
與教訓,督責,使人歸正,教導人學而.
都是有益的.
因此他叫提摩泰要好好教導神.
然後到了最後第四章.
就保羅臨終的吩咐.
他說:我奉這位快來.
末後很快再來的主.
我祝福你.
務要傳道.
無論得時不得時.
你總要傳福音.
無論時機是對時機是不對.
你都要去傳.
因為這個是已經交託了給我們的責任.
各位弟兄姊妹.
我們今天需要起來.

$^{201}$努力的傳福音.
我又傳,你又傳,他又傳.
大家都傳.
而且用不同的方法.
你可以寫信給人.
你可以買一些小冊子送給人.
你可以用YouTube.
你可以用很多的媒介.
努力將福音傳開.
然後保羅就說.
因此美好的帳我已經打過.
所信的道我已經守住了.
當跑的路我已經跑盡了.
末後的時候我會見主.
他就提到.
好像一隻船.
本來泊在岸邊.
有條纜綁上岸上面.
你解纜這隻船做啟航.
你箍住牛的鼻子那裡.
有個洞有條繩子.
綁起地上有個泉塞.
你解開牛就可以去吃草了.
你搭了一個帳幕.
用釘子拉著帳幕的釘子在地上.
現在收營了.
就拉起釘子.
於是就走了.
原來保羅說.
我離世的日子到了.
我被囂墮的日子到了.
因此保羅預知自己很快離開世界.
因此他就勸提摩泰.
怎樣在主的面前要盡力.
當跑的路你跑盡了.
當打的仗你打完美好的仗.
所信的道你守著.
然後你就會見這位持公義的官員.
給我們的主.
然後最後結束的時候.

$^{241}$保羅就有很多的問安.
特別他說那些書卷最緊要帶來.
書卷在原文裡面就是.
那些羊皮卷.
很證明那些不是普通的書.
是聖經.
他說最緊要帶聖經給我.
帶些衣服給我.
因為那裡很冷.
然後說所有人都會不在我身邊.
唯有路加在我身邊.
但是主與我同在.
保羅很肯定一件事.
沒錯他坐牢的時候很孤單.
沒有人陪到他.
可能只剩下一個路加.
但是他說主與我同在.
今天都是一樣.
我們有時會覺得好像只剩下我自己.
但是主從來沒有離開我們.
祂與我們同在.
祂加力給我們.
祂與我們同行.
直到我們再見主的面.
所以願神賜福給各位.
讓您忠心事奉主.
直到我們再見祂的面.
阿們.
\newpage



\section{約翰福音 13:3-10}
\label{sec:032139}
\textbf{愛的醒覺}~~~~蔡允瑤傳道
\newline
\newline
連結: \href{http://www.church.com.hk/acms/content.asp?site=cdc\&op=show\&type=product\&code=032139\&layout=sermon}{\texttt{032139}} ~~~~ 語音日期: 2026-05-17
\newline
\newline
\hyperref[sec:032138]{\small{< < < PREV SERMON < < <}}
~
\hyperref[sec:index_chronic]{\small{[返順時目]}}
~
\hyperref[sec:032140]{\small{> > > NEXT SERMON > > >}}
\newline
\newline
約翰福音 13:3-10
\newline
\begin{longtable}{cl}
\hline
\hline
章節 & 經文 (和合本修訂版)\\
\hline
13:3 & \begin{tabularx}{0.7\textwidth}{X} 耶穌知道父已把萬有交在他手裡，且知道自己是從神出來的，又要回到神那裡去， \end{tabularx} \\ \\ \relax
13:4 & \begin{tabularx}{0.7\textwidth}{X} 就離席站起來，脫了衣服，拿一條手巾束腰， \end{tabularx} \\ \\ \relax
13:5 & \begin{tabularx}{0.7\textwidth}{X} 隨後把水倒在盆裡，開始洗門徒的腳，並用束腰的手巾擦乾。 \end{tabularx} \\ \\ \relax
13:6 & \begin{tabularx}{0.7\textwidth}{X} 到了西門‧彼得跟前，彼得對他說：「主啊，你洗我的腳嗎？」 \end{tabularx} \\ \\ \relax
13:7 & \begin{tabularx}{0.7\textwidth}{X} 耶穌回答他說：「我所做的，你現在不知道，但以後會明白。」 \end{tabularx} \\ \\ \relax
13:8 & \begin{tabularx}{0.7\textwidth}{X} 彼得對他說：「你絕對不可以洗我的腳！」耶穌回答他：「我若不洗你，你就與我無份了。」 \end{tabularx} \\ \\ \relax
13:9 & \begin{tabularx}{0.7\textwidth}{X} 西門‧彼得對他說：「主啊，不僅是我的腳，連手和頭也要洗！」 \end{tabularx} \\ \\ \relax
13:10 & \begin{tabularx}{0.7\textwidth}{X} 耶穌對他說：「凡洗過澡的人不需要再洗，只要把腳一洗，全身就乾淨了。你們是乾淨的，然而不都是乾淨的。」 \end{tabularx} \\ \\
[1ex]
\hline
\hline
\end{longtable}
$^{1}$各位姊妹早晨.
我代表西雅圖正道堂向大家問安.
很開心今天能夠和你們一起敬拜.
其實我和多倫多也有一份緣份.
我上一次在多倫多應該是三十年前.
我認識多倫多.
但其實我不是在這裡住過.
不過我有六年的時間在Kingston Queen's皇后大學讀大學.
所以當時一有節日或者是Reading Week的時候就會出土.
不知道現在是否還是用這個名詞「出土」.
我記得第一年Reading Week出到來多倫多.
當時我還未信耶穌.
但也有去過團契.
來到的時候有朋友邀請去一個聚會.
我說好吧去吧沒甚麼事要做.
原來是大多市的報道會.
當年講員是馮臻牧師.
我已經忘記他說甚麼了.
不過當他呼召的時候.
他教會眾一起來唱「他是主」.
我記得我其實是坐在這裡的後排.
在邊邊的.
當時他一直叫我們閉上眼.
一直聽著「他是主」.
嘩很感動的心.
好像有人在敲門.
「Wendy開門吧我想進來」.
很感動但同時很掙扎.
我可以告訴你.
在掙扎的當中.
我沒有那個勇氣走到台前去決志信主.
之後其實都覺得為甚麼我這麼沒用.
我是不是沒得救了.
真的沒得救那個救贖的救.
真的敲到我悶了.
但我竟然不去回應他.
覺得自己很失敗.
今天我們看的經文.
有另一個人彼得.
在當中可能他都是經歷一段很失敗.

$^{41}$他覺得很失敗的一個經歷.
我們從兩段的經文.
來看一看彼得的經歷.
和神在當中的心意是甚麼呢?.
多謝剛才Elisa幫我們.
和一班敬拜隊的同工.
來讀出這段經文.
我想帶大家回到三年多前.
耶穌呼召彼得.
耶穌呼召彼得在四個福音都有記載.
但在約翰福音第一章裡.
特別講到耶穌看著彼得說.
你是約翰的兒子西門.
你要稱為基發.
本來叫做西門.
現在耶穌幫他改名成為基發.
翻出來就是彼得.
我想問問這裡有沒有人叫Peter?.
這裡沒有人叫Peter.
Peter或者希臘文Petros.
的意思就是石頭.
嘩!這個本來叫西門的人.
耶穌幫他改名.
還要是改了成石頭.
如果我們看回聖經的時候.
嘩!改名是一件不得了的事.
是神的呼召將生命來到扭轉.
進入神的計劃當中.
亞伯蘭改成亞伯拉罕.
雅各改成以色列.
西門改成石頭.
當彼得來回應耶穌這個呼召的時候.
的而且確彼得很想很想.
耶穌呼召我做石頭.
我就要做到這一塊大石頭.
我們想一想.
究竟在彼得心目中那塊大石頭.
是一塊怎樣的大石頭呢?.
我相信在他心目中這一塊大石頭.
就是好像耶穌基督那樣.

$^{81}$祂行神蹟奇事.
大而有力.
被敬仰尊敬.
去到在《馬太福音》.
在《該撒尼亞肥勒比》.
耶穌問:你們說我是誰?.
彼得怎樣去回答?.
在《馬太福音》第十六章十六節.
西門彼得回答說:.
「你是基督,是永生神的兒子」.
啊!這個就是那塊大石頭.
「你是基督」.
基督是什麼?.
尼賽亞.
君王.
受高者.
永生神的兒子.
彼得都希望好像.
這一塊大而有力的石頭一樣.
我都要做一塊大石頭.
但是彼得絕對沒有想過.
這一塊大石頭.
或者他所敬仰的耶穌.
竟然會被人踐踏.
擊打.
受苦.
所以同樣在剛才那段經文.
十六章那裡.
後來耶穌對門號說.
我將會受苦.
要踏上十字架的道路.
彼得怎樣跟耶穌說?.
彼得說:.
「萬萬不可」.
不可能的.
你是這一塊大石頭.
怎可以這樣的.
主啊.
萬不可如此.
這事必不臨到你身上.

$^{121}$這一件事必不臨到一塊.
大而有力的石頭上.
諷刺的是什麼呢?.
彼得真的成為了一塊石頭.
在這個時刻.
但是一塊什麼石?.
耶穌怎樣跟彼得說?.
祂說:.
「你是伴我腳的」.
真的成了一塊石頭.
可惜是一塊伴腳石.
好了.
來到今天的經文.
三年過了.
去到耶穌真的要上十字架.
昨天休穆斯說到.
在十三章一開始的時候.
耶穌跟一班門徒一起.
跟他們說最後一份訓言.
耶穌愛他們愛到底.
而其中.
耶穌怎樣去表達.
祂愛他們愛到底呢?.
就是拿一條毛巾.
束上腰.
跪在這裡.
來跟他們洗腳.
但去到彼得的時候.
祂說:.
「你永不可以洗我的腳」.
為什麼彼得會這樣說呢?.
「你永不可以洗我的腳」.
彼得看到.
祂和耶穌兩者的差異.
「你」.
「我」.
「你是永生神的兒子」.
「你是基督啊」.
「你是我的老師」.
「你是我的拉比啊」.

$^{161}$「我」.
「只不過是彼得而已」.
「我是你的學生啊」.
「我怎可以讓你來洗我的腳呢?」.
「你」.
「和我」.
「差得太遠了」.
然後耶穌怎樣說呢?.
耶穌說:.
「我若不洗你」.
「你就與我無分了」.
耶穌也說到「你」和「我」.
祂說:.
「如果我不洗你的腳的話」.
「你就與我無分了」.
耶穌所著重的是什麼呢?.
都是「你」和「我」.
不過不是「你」,「我」.
或者在耶穌的角度來說.
就是彼得「你」,「我」的差別.
而是在說一個很重要的字在當中.
「與」.
「你」,「與」,「我」.
「就無分了」.
耶穌不是在說兩個的差異.
而是在說兩個的連接.
在英文《NIV》那裡翻譯.
祂說:.
「No」.
祂說:.
「如果我不洗你的腳的話」.
「Then you have no part with me」.
「with」.
這個「與」.
耶穌所著重的是那份關係.
那個聯繫.
但是你猜猜彼得他明白嗎?.
相信彼得不明白.
因為彼得他接著怎麼說呢?.
他說:.

$^{201}$「哎呀!」.
「如果是這樣,主啊!」.
「你不要單是我的腳」.
「手和頭你都要洗啊!」.
如果耶穌跟彼得已經有這個「with」的聯繫.
「你與我的時候」.
為什麼還要洗手又要洗頭呢?.
因為彼得還在想兩者的差異.
彼得還在想:.
「啊!是啊是啊!」.
「耶穌啊!」.
「如果你洗我的腳的話」.
「我就可以好像走出你所走出的路了」.
「不如你洗我的手吧!」.
「於是我的手就可以做出你的手做出的事了」.
「洗我的頭啊!」.
「我的思想我的智慧要好像你一樣」.
「我的言語要好像你一樣」.
彼得還在想的是什麼呢?.
不是「with」.
彼得還在想的是「from」.
我要在耶穌那裡拿到什麼呢?.
洗我的腳.
洗我的手.
洗我的頭.
「with」還是「from」?.
我想其實今天我們做了基督徒.
無論是多久我們都要想一下.
究竟我們信耶穌.
究竟我們說我們和耶穌同行.
我們做門徒.
究竟是「with」耶穌.
「與」耶穌.
還是「from」耶穌啊?.
在耶穌裡面拿些東西吧.
拿了之後.
我就成了這個大石頭了.
越來越像這個大石頭了.
我好像基督一樣.
如果是「from」的時候.

$^{241}$基督只不過是一個模範.
或者祂給了一些恩賜給我們.
給了一些能力給我們.
但是那個關係是否存在呢?.
有沒有那個聯繫呢?.
「from」還是「with」呢?.
彼得在這一刻.
其實他都可能覺得.
我都差不多了.
我都差不多是那個大石頭了.
以致後來彼得對耶穌說.
他說:主啊!我願意為你寫命啊!.
耶穌啊!我已經是一塊大石頭了.
你有什麼困難的話.
我幫你頂住啦!.
我願意為你寫命啊!.
但是彼得做不做到啊?.
當我們去看另一段經文的時候.
就發覺.
彼得連這一塊大石頭.
幫耶穌頂住.
都做不到啊!.
在《路加福音》.
彼得三次不認主.
昨天我們玩這個破冰遊戲.
大家一看到三隻雞.
就已經知道了.
是說彼得不認主啊!.
在這裡我就用《路加福音》這一段經文.
一開始的時候.
有人說:這個人不就是和那個人.
即是耶穌一起的嗎?.
彼得說:我不認得他.
不是說:我認不得他的樣子啊!.
太遠了,看不到啊!.
不是啊!.
英文是說:I do not know him.
我不認識他啊!.
我不認識這個人啊!.
接著,第二次了.

$^{281}$有人和彼得說:.
你不是和他們一黨的嗎?.
彼得說:不是啊!.
我不是和他們一黨的.
我還我,他們還他們啊!.
不要將我們扯上關係啊!.
接著,第三次了.
有人說:這個人.
即是說彼得啊!.
真是和那個人耶穌一起的啊!.
他也是加利利人啊!.
彼得再次說:我不認識他啊!.
我不知道你在說什麼啊!.
在這三次之中.
我不知道他是誰.
我不認得他.
我不認識他.
我不是和他們在一起的.
我不是那個人啊!.
彼得和耶穌劃清界線啊!.
莫說你與我啊!.
這個關係沒有了.
甚至乎連彼得心目中的你啊!.
那個耶穌啊!.
他也是否定啊!.
我不認識他啊!.
或者在當中.
彼得真是有一點點的.
不太認識這個耶穌啊!.
因為在彼得的心目中.
耶穌不是這樣的啊!.
怎麼會被人抓了不反抗啊!.
怎麼會被人打了又不出聲啊!.
這個不是彼得心目中.
耶穌那塊大石頭啊!.
所以在這裡.
彼得和耶穌.
你與我沒有了.
甚至乎你彼得也很懷疑.
為什麼會這樣啊!.

$^{321}$彼得的石頭夢是什麼啊!.
彼得的石頭夢是一塊大石頭啊!.
可惜經過這三次不認主的時候.
粉碎了他的夢.
究竟耶穌是誰啊!.
究竟我是誰啊!.
完全地被拆毀,被粉碎.
為什麼我要選《路加福音》這段經文呢?.
在這裡他特別講到.
當彼得三次不認耶穌之後.
雞就叫鳥.
這一句只是在《路加福音》裡出現的.
主轉過身來看彼得.
然後彼得就想起主對他的說話.
今天在雞叫耳先.
你要三次不認我.
他就出去痛哭了.
我想問問你們.
幻想一下.
如果你是彼得的話.
當時見到耶穌轉過身來看著你.
耶穌的眼神是一個什麼樣的眼神呢?.
想一想.
你見到耶穌轉過身來看著你.
耶穌用一個什麼樣的眼神來看著彼得.
看著你呢?.
會不會是一個很失望的眼神呢?.
唉!彼得你真是三次不認我.
會不會是一個責備的眼神呢?.
如果我們看回之前在《路加福音》.
耶穌說到彼得你要三次,你會三次不認我.
在那裡有一節的經文.
耶穌又對彼得說,又對西門說.
西門!西門!撒旦想要得著你們.
好西你們像西物質一樣.
但我已經為你祈求.
叫你不至失了信心.
你回頭以後要堅固你的弟兄.
我相信耶穌在那一刻轉身過來看著彼得.
不是一個失望責備的眼神.

$^{361}$而是一個愛和憐憫的眼神.
是一個告訴彼得.
彼得我明白的.
但彼得你不用怕.
彼得我已經為你祈求了.
並且我現在就是為你祈求了.
我為你祈求,你不會失去信心的.
你回頭吧!你回頭吧!.
在這一刻,彼得看見耶穌.
看見他自己.
這是一個彼得見到耶穌.
見到彼得.
你與我.
一個的連結.
我見到你,見到我.
不知道你們知不知道.
當你跟一個人很近的時候.
你走去看著他的眼睛,他看著你的時候.
你知不知道你會見到什麼?.
看看自己.
一個的連接.
一個愛的相遇.
彼得在這一刻.
跌到最低.
他的夢想.
他以為他有那個大石球的身份.
他以為神所呼召他,我要做一個大石球.
完全的粉碎.
跌到一個谷底那裡.
好像剛才那首歌.
去到無水的盡處.
沒得救了.
但是在一個盡處的當中.
耶穌仍在.
並且耶穌看著彼得.
所以當彼得見到耶穌.
看著他的時候.
其實這個就是彼得開始.
來到被重建的那一刻.
當然,這個重建的過程.

$^{401}$不是我看到耶穌看著我.
就馬上沒事了.
我信心回來了,我回轉了.
不是,可能是一段長的時間.
是不容易的一段時間.
我們今天有了整本聖經.
我們見到去到使徒倖轉的時候.
彼得站起來了.
他真的回頭了.
來到去見過他的弟兄.
是教會的領袖.
甚至我們去看彼得前書.
彼得前書那裡.
他明白到耶穌改他的名字.
做彼得石頭的意思是什麼.
原來他們是被神建造成為靈宮的石頭.
原來一塊很大的石頭.
是不能夠被神建造成為一個靈宮的.
一塊大的石頭就在這裡.
就是這麼多了.
你要被神來到去建造成為一個靈宮.
不是說你一塊石頭.
還有,你看這裡有多少弟兄姊妹.
每一個弟兄姊妹都是一塊石頭.
你試想一下.
這裡你們個個都是大大小小的.
怎麼建造呢.
我想建造一層,第二層就倒塌了.
需要被雕琢的.
甚至需要被拆毀,粉碎的.
這個過程不容易.
有時會很痛苦的.
但是在這個過程當中.
耶穌在這裡.
你與我.
並且耶穌說我已經為你祈求了.
我現在也是為你祈求了.
你回頭了.
留意的是主轉身過來看彼得.
不是因為彼得叫耶穌.

$^{441}$然後耶穌就轉過來看一看.
是耶穌主動的轉身過來看彼得.
在彼得去到這個無水的盡處裡.
沒得救了.
跌到谷底的時候.
其實他不懂得去呼叫耶穌.
但是耶穌愛他.
愛到底.
就算在他的失敗當中.
在他不認耶穌的當中.
耶穌仍然愛到底.
耶穌主動的轉身過來看彼得.
愛與憐憫的一個目光.
彼得,我已經為你祈求了.
我為你祈求了.
你不會失去信心的.
你回頭的時候.
你就可以堅固你的弟兄了.
彼得看見耶穌看著他.
在一個低谷的當中.
無水盡處的當中.
看見耶穌看著他.
這個就是一個愛的相遇.
愛的連結.
你與我在愛的裡面.
在這份關係裡面.
說回我.
剛才說到第一年的時候.
我不敢走下來來決志信主.
我心想,哎,失敗.
沒得救了.
我相信神都會.
哎,算了.
敲了你門.
你都不認了.
我去敲別人的門吧.
但是.
神仍然沒有放棄我.
當我自己放棄了自己的時候.
神沒有放棄我.

$^{481}$也不是一段短的日子.
直到我第三年讀幼稚園的時候.
我記得在開學的時候.
我們那時候.
還要是個人註冊.
有一些不同的學校的組織.
會擺攤子.
其中一個就是我們校園的團體.
KCCF.
我記得當時.
我看著這一班人.
然後我說.
他們有些東西我很想要.
我不知道是什麼東西.
不過我很想要他們有的東西.
所以在第三年.
我就回到團體裡面.
來參與團體.
之後不是在一個報道會.
沒有在一個很大的激動感動當中.
來信耶穌.
只不過是有一位同學.
有一天上課之前.
他就叫我過來.
Wendy.
你已經回到團體一段時間了.
你想不想信耶穌?.
我就問.
什麼叫信耶穌?.
他就把福音告訴我.
在一個很平淡很自然的境況裡.
我決志信主.
神沒有放棄我.
我以為在我的一個低谷.
在一個沒得救的境況.
原來神沒有放棄過我.
就好像彼得一樣.
彼得在一個低谷當中.
在一個無水的盡處當中.
耶穌沒有放棄他.

$^{521}$耶穌仍然在那裡.
主動地望著他.
而當彼得望著耶穌的時候.
就在這個關係當中.
來相遇.
彼得的石頭夢.
本來是要做一個大而有力.
很威水的一塊石頭.
可惜就真的很威水.
威水在什麼呢?.
成為了一個絆腳石.
以致神需要來拆毀彼得.
神需要來雕琢他.
這時候才能夠用彼得來建造成一個靈宮.
不容易啊.
這個過程實在是不容易.
彼得他對自己的身份認同.
他的成就.
甚至他的期望完全被粉碎.
但是在粉碎的當中.
耶穌仍然與他同在.
神的計劃仍然在彼得的身上進行.
因為主耶穌就好像在十三章跟彼得說.
如果我不洗你的腳,你與我無份.
耶穌所著重的不是差距.
而是你與我with.
十五章的葡萄樹與枝子.
常在基督裡面.
基督常在我們裡面.
with.
不是from.
是與耶穌.
而不是在耶穌那裡拿到東西.
而是我自己來去活.
我自己基督徒的人生.
在這裡我邀請大家一起來閉上眼.
去到神的面前.
看了這兩段經文.
看到彼得一開始的時候.
他的石頭夢.

$^{561}$後來被拆毀.
被粉碎.
在這個時候邀請你.
來與主耶穌有一段對話.
問問主耶穌.
在這段經文裡.
你想對我說什麼呢?.
不知道.
你心裡面又有沒有一個石頭夢呢?.
這個石頭夢是with耶穌.
還是from耶穌呢?.
你願意放下你自己的石頭夢.
一個from耶穌的石頭夢.
來進入一個with耶穌.
與耶穌一起的夢呢?.
又或者你們現在是否在一個破碎當中.
被神拆毀的當中.
去到一個無水的盡處.
不容易呀!.
但是耶穌在這裡.
耶穌看著你.
用祂愛和憐憫的眼光.
看著你.
耶穌已經為你祈求.
並且現在為你祈求.
你願意看著耶穌嗎?.
有時我們都不明白.
究竟你愛我們愛到底.
是一回什麼事.
有時甚至我們覺得我們不配你愛我們愛到底.
但是主呀!你卻是常常的來尋找我們.
你主動的轉身來看著我們.
主耶穌,多謝你!.
多謝你看著我們!.
讓我們能夠在見到你見到我的時候.
就是在這個時刻當中.
就算是在谷底的裡面.
就算是在一個無水的盡處當中.
只要我們願意看著耶穌.
我們就能夠看著耶穌看著我們.

$^{601}$在這個你與我的當中.
在這個愛的關係的當中.
與耶穌你相遇.
阿們!.
\newpage



\section{歌羅西書 3:1-4}
\label{sec:032140}
\textbf{基督徒的現在與未來}~~~~羅錫堅牧師
\newline
\newline
連結: \href{http://www.church.com.hk/acms/content.asp?site=cdc\&op=show\&type=product\&code=032140\&layout=sermon}{\texttt{032140}} ~~~~ 語音日期: 2026-05-17
\newline
\newline
\hyperref[sec:032139]{\small{< < < PREV SERMON < < <}}
~
\hyperref[sec:index_chronic]{\small{[返順時目]}}
~
\hyperref[sec:032141]{\small{> > > NEXT SERMON > > >}}
\newline
\newline
歌羅西書 3:1-4
\newline
\begin{longtable}{cl}
\hline
\hline
章節 & 經文 (和合本修訂版)\\
\hline
3:1 & \begin{tabularx}{0.7\textwidth}{X} 所以，既然你們已經與基督一同復活，就當求上面的事；那裡有基督，坐在神的右邊。 \end{tabularx} \\ \\ \relax
3:2 & \begin{tabularx}{0.7\textwidth}{X} 你們要思考上面的事，不要思考地上的事。 \end{tabularx} \\ \\ \relax
3:3 & \begin{tabularx}{0.7\textwidth}{X} 因為你們已經死了，你們的生命與基督一同藏在神裡面。 \end{tabularx} \\ \\ \relax
3:4 & \begin{tabularx}{0.7\textwidth}{X} 基督是你們的生命，他顯現的時候，你們也要與他一同在榮耀裡顯現。 \end{tabularx} \\ \\
[1ex]
\hline
\hline
\end{longtable}
$^{1}$我曾經同一位傳道人傾偈.
他大約五十歲.
溫文有禮.
常怒笑容.
他不講我都不知道原來他的身世這麼坎坷.
他告訴我.
曾經多次犯法.
打家劫舍.
坐監加埋都十幾廿年.
吸毒戒毒都過jor 好多次.
後來係戒毒所裡面.
信jor 耶穌生命就完全改變.
最後仲做jor 傳道人.
我有認識好多.
係屬於勝利人生組ge 基督徒.
點樣勝利人生組呢.
佢地出身良好.
名校畢業.
有學識.
就有工作.
有兒有女.
有樓有車.
咁呀.
退休仲有一大筆養老金.
但係佢地唔快樂唔開心.
佢地對自己ge 健康家庭經濟都唔滿意.
佢地覺得生活冇意義.
前途冇希望.
佢地係教會裡面ge 參與都唔積極.
聚會就循例ge .
唔去就唔好意思.
有時落雨就係網上聚會.
但過jor 兩三次唔黎.
前途人關心佢.
佢又覺得囉嗦.
咁結果又返返黎.
你地今日就唔應該係咁.
因為今日落雨你地都黎.
咁樣先好難得.
佢地覺得.

$^{41}$耶穌基督講ge 永生好似得唔到.
耶穌基督應許ge 豐盛生命.
佢又唔經歷過.
如果你覺得你係咁ge 基督徒.
或者話多或少似你咁.
咁今日ge 訊息就好arm 你用.
剛才gwo 四字經文.
係好寶貴.
呢四字經文係話比你知道.
點樣可以得到屬靈生命ge 最高境界.
咁又係有D 咁ge 事.
點解頭先奔波讀我好似唔覺咁.
咁我讀多一次.
係咪通常你好似.
即係好似食叉燒包咁樣.
食第一個就唔飽ge .
第二個就飽就唔係因為第二個好.
因為食jor 第一個先.
所以你們既然與基督一同復活就應該尋求天上的事.
哥羅西斯三章第一節.
哪裡有基督在在神的右邊.
你們思念的是天上的事不是地上的事.
因為你們已經死了.
你們的生命與基督一同隱藏在神裡面.
基督就是你們的生命.
但顯現的時候你們也要和他一同在榮耀裡顯現.
所以今日我要同大家分享一個屬靈秘笈.
點樣得到屬靈生命最高境界.
而且一定可以做到.
打功夫都未必一定得.
羅木斯說我一定可以做到.
咁我地首先睇第一節.
第一節講到基督在在神的右邊.
呢度講jor 兩樣ye .
第一就有基督坐係度.
第二基督坐係神ge 右邊.
舊約時代ge 祭司就站著獻祭.
講道通常都係站著獻祭.
舊時ge 祭司就站著獻祭.
但基督作為大祭司獻上一次ge 永遠贖罪祭.

$^{81}$就係佢自己ge 身體獻上就坐低.
即係話人類ge 救贖工作就徹底完成.
就唔需要再站一度坐低.
同埋係古代ge 傳統裡面.
中東坐低坐係右邊.
就係最高ge 榮耀同埋權益.
如果你睇希伯來書你就睇到.
希伯來書十章十一到十二節.
佢話所有ge 祭司都是天天站著事奉.
多次獻上同樣ge 祭品.
gwo 些祭品永遠不能把罪除去.
唯有基督獻上了一次永遠有效ge 贖罪祭.
就係佢拎住自己ge 血進入至聖所一次.
就係神ge 右邊坐下.
所以呢度話比佢聽.
基督ge 救贖工作完成.
而且佢已經坐低搞掂.
咁第二點呢.
就係係第三軸gwo 度我地已經死了.
我地中國人就唔鍾意早晨漏漏就講個死字.
今日我要講比你聽ge 死字就唔係唔好ge .
而係非常之好ge .
因為呢個死ge 意思就釋放.
我地對罪ge 權勢死jor .
即係罪ge 權勢原本係綁住我地.
但因為我地向佢死jor .
佢我地就得jor 釋放.
佢就唔能夠再管轄我地.
所以保羅話你地已經死了.
就係我地已經向世界死了.
呢個係同魔鬼爭建最有效ge 方法.
就係我釋放你管唔到我.
因為我與基督同死.
咁然後第三點呢就去返第一節gwo 度.
我地已經與基督一同復活.
但如果你死jor 唔復活.
咁就好弊架喇.
咁就好絕望.
佛教呢同基督教有好大ge 唔同.
佛教其中一個唔同佢最高境界就係空.

$^{121}$即係冇ye .
即係咩都冇.
但呢個呢就唔能夠滿足我地ge 渴望.
人死jor 之後話咩都冇.
咁我個生命係咩ye 呢.
但係基督唔係咁ge .
耶穌基督死後復活.
而聖經多次提醒我地.
信佢ge 人都會與基督一同復活.
呢個復活呢.
唔係靈魂飄呀飄呀上天堂唔知自己點.
呢個復活係身體ge 復活.
呢個好重要.
基督教ge 復活係你ge 身體復活.
有人話我做jor 白內障.
我做白內障係11年前ge 事.
咁醫生話你個鏡片可以.
保證50年.
50年gwo 陣時我已經超過100歲.
咁但係你復活gwo 陣時就冇白內障.
你個鏡片都唔係度.
因為上帝俾你個榮耀身體係perfect.
咩叫perfect呀.
我都唔知.
總之樣樣都掂晒.
咁唔使每日度血壓.
咁我地因為基督復活.
我地都會同佢一齊復活.
咁唔單止將來會復活有個榮耀身體.
你而家已經與基督一同復活.
呢個就勁.
因為我地已經脫離jor 罪惡.
脫離jor 魔鬼ge 管轄.
我地已經有個新ge 生命.
所以聖經話如果有人信耶穌.
你已經係一個新做ge 人.
你將來gwo 個榮耀身體係賜比你.
但你而家就算仍然戴住白內障眼鏡.
好似我咁.
你都已經有個新ge 生命.

$^{161}$舊事已過都變成新.
係上帝眼中你已經係一個新做ge 人.
所以有新ge 生命.
有新ge 生活.
有新ge 能力.
新ge 目標.
所以加拉提書第二章二十字.
保祿話我已經與基督同定十字架即死jor .
而家活著ge 唔再係我.
因為基督活在我裡面.
基督復活我又復活.
如今我係肉身復活.
係因為信神ge 兒子而活.
他愛我為我捨己.
第四.
大家記唔記得.
第一基督在在神的右邊.
第二我地已經死了.
第三我地已經與基督一同復活.
第四我地與基督一同隱藏在神裡面.
隱藏在神裡面就有幾樣ye .
第一象徵我地被保護.
冇人搵到我地.
捉迷藏.
捉迷藏通常都俾人搵到.
但係隱藏係神裡面就好隱蔽.
亦象徵著舊生命過去同新生命與基督一同聯合.
所以我地再睇返第三第四節.
呢個非常寶貴.
因為你們已經死了.
你們的生命與基督一同隱藏在神裡面.
基督就是你們的生命.
呢兩節經文係非常奇妙.
所以如果你今日下午冇ye 做.
我又鼓勵你返屋企諗下呢段經文.
乜ye 叫做與基督一同隱藏在神裡面.
一個人點樣隱藏係另一個人裡面.
我地用一個例子黎諗.
就係一個baby係媽媽個肚裡面.
你知唔知baby係媽媽個肚裡面.

$^{201}$gwo 個環境係人世間最完美.
你出jor 世之後都冇咁好.
點樣第一佢又安全又舒服.
係媽媽個肚裡面呢個地心吸力就減輕jor .
頭gwo 半年呢gwo 個babyge 地心吸力就減輕百分之七十.
咁你越黎越大個呢就冇咁減輕.
因為佢個胎裡面ge 羊水浮起佢.
你開唔開心呀.
減輕百分之七十喎你知唔知呀記唔記得呀.
唔記得.
仲有呢係佢ge 胎裡面ge 溫度就恆溫.
你個冷氣機最新出gwo 個penicillin都冇咁好ge .
恆溫.
咁仲有呢你所有ge 需要呢個baby所有需要都媽媽公認實.
咁係媽媽個肚裡面呢係又安全又滿足.
咁仲有一個基督徒ge 生命係神裡面藏住.
佢ge 意思亦都係咁樣.
係受到上帝ge 保護同埋供應.
係母胎裡面ge baby仲有一個特點ge 就係被隱藏.
外界D 人睇唔到唔知道佢係點樣ge .
唔單止外界隱藏佢自己都唔知道自己係點.
呢個baby都唔知道佢自己係點.
呢個隱藏呢連佢自己都隱藏只有上帝知道佢.
咁你話係咪特別奇妙呢.
呢個baby呢係媽媽ge 肚裡面就冇掙扎.
無欲無求冇掙扎.
媽媽唔會對個胎裡面ge baby話.
阿仔你勤力D 呀.
阿仔你精歸寧利呀有冇架.
我就唔記得到阿媽佢對我咁講ge .
呢個係媽媽ge 肚裡面ge baby係冇掙扎ge .
所以如果你聖經話你與基督一同藏在神裡面.
就係你安全你舒服你唔需要有掙扎.
你所有ge 供應完全ge 滿足.
我諗到個baby係媽媽ge 肚裡面.
媽媽有個心跳ge .
個心跳其實好有用ge .
佢俾個baby有安全感.
同埋佢隔絕jor 外面ge 噪音.
所以如果個爸爸等電車等巴士都唔影響個baby.

$^{241}$因為媽媽有心跳bang bang bang bang 對佢黎講就好好笑.
我地係神裡面都係咁樣.
仲有一點.
你地ge 生命與基督一同藏在神裡面.
咩叫與基督一同藏在神裡面.
點解與基督一同呢.
一個baby同媽媽ge 關係.
可以話人與人之間最親密就咁.
但係都仍然係兩個生命兩個心跳.
個babyge 心跳同媽媽心跳唔同.
媽媽心跳可能70 baby係120 130.
但係呢度話你地與基督一同藏在神裡面.
基督是我地ge 生命.
唔係兩個生命係一個生命.
你同基督係一個生命一個心跳.
頂至尾呢個就係屬靈生命ge 最高境界.
就係呢句話你地與基督一同藏在神裡面.
咁你點樣可以做到與基督一同藏在神裡面呢.
咁你要話我要做到咩.
我要問阿江牧師.
我係咪要經過屬靈十課.
或者要考埋個牌.
又要識彈琴又要有事房咁樣先得呢.
保羅話唔係喎.
你地已經死了.
你地已經與基督同復活.
你地與基督一同藏在神裡面.
所以頂至尾.
其實呢今日我要鼓勵大家.
如果你覺得自己好得.
如果你覺得自己好多困難.
覺得自己冇乜成就.
覺得自己比隔離gwo 個唔係好得.
你唔使擔心.
因為你與基督一同藏在神裡面.
上帝唔叫你掙扎.
上帝冇個quota比你.
信你冇話寶珠你每日要讀兩篇聖經.
如果唔係就唔夠屬靈.
你用三篇.

$^{281}$上帝冇要你今日要同幾多個人去講福音.
因為你與基督一同藏在神裡面.
你唔需要做ye .
屬靈ge 最高境界呢.
就係係主裡面享受.
聖經話非勢力非才能係靠我ge 靈黎到成事.
大衛對葉哥尼亞話.
爭戰ge 勝敗係在乎耶和華.
The battle belongs to the Lord.
唔係我打你 上帝打你.
非勢力非才能係靠我ge 靈成事.
你地得救係在乎皈依安息.
你地得力係在乎平靜安穩.
所以我做基督徒.
我係好舒服ge .
今日我地係房企祈禱.
肇方話穆斯林好似做好多ye .
點解唔gui 呀.
我每個月來回香港.
台灣其實我同我太太都累.
非勢力真係累.
我星期三晚返黎.
由夜晚八點訓到早上七點.
訓jor 十一個鐘.
但係我唔覺得好肉緊.
唔覺得好辛苦.
因為我以基督一同創作神明.
呢場仗係佢打.
我返黎之前三日.
星期一星期二星期三好忙.
返黎之後.
星期四星期五星期六都好忙.
好多ye 發生.
佢周圍好似好多隻猴子.
咁但係我就提醒自己.
我以基督一同創作神明.
基督係我生命.
一個生命一個心跳.
我已經以基督同釘十字架.
而家活著佢唔再係基督.

$^{321}$我以基督係上帝ge 肚裡面.
上帝唔會要我有鋪打.
上帝唔會要我努力.
呢個就屬靈生命ge 最高境界.
就係全部交比上帝.
咁容易唔容易做呀.
其實可以做到個個人都做到.
上帝好fair.
上帝唔會對你ge 要求.
嘩你ge IQ150都做唔到.
所以我今日ge 題目就係.
基督徒ge 現在與未來.
四個現在.
基督在我神的右邊.
我命已經死了.
我們的生命與基督一同藏在神裡面.
我與基督同復活.
咁基督ge 未來係點樣呢.
第一基督將會顯現.
基督會返來.
佢復活升天gwo 陣時你唔係度.
但佢返黎gwo 陣時你會睇到.
佢會以榮耀ge 身份出現.
王者ge 身份顯現.
佢第一次黎咁隱蔽ge .
佢黎ge 時候亦都係.
等於時代ge 終結.
所以你話係時代.
而家時代臨臨黎你唔使擔心.
因為佢會終結.
係馬大福音.
耶穌自己就話.
那時人子的徵兆.
要顯在天上.
地上的萬族都要哀號.
並且看見人子帶著能力.
滿有榮耀.
帶著天上的雲降臨.
當浩同發出響聲.
祂要差派使者.

$^{361}$把祂的選民從四方.
從天的這邊到天的那邊.
都召聚來.
你知唔知.
仲有喎 基督顯現ge 時候.
我地都會同佢一齊顯現.
呢四節都提到.
我地也與基督.
一同顯現在榮耀裡.
嘩.
所以你話我而家好up 得.
我又冇錢又冇面.
但係你與基督一同.
顯現在榮耀裡面就好緊要.
同響榮耀.
那時亦都係基督同教會.
同屬於佢ge 人完全的結合.
罪同埋死亡.
係被取締.
啟示錄二十二章.
一到七節就講到當時的情況.
天使又指示.
我一道明亮如水晶.
的生命水的河.
從神的.
和羔羊的寶座.
那邊流出來.
經過城裡的街道.
河的兩邊有生命樹.
結十二次果子.
每月都結果子.
事業可以醫治列國.
所有的咒詛都不再有了.
城裡有神和羔羊的寶座.
他的僕人都要事奉他.
你同我.
就係佢ge 僕人.
又要見他的面.
他的名字必寫在他們的額上.
不再有黑夜了.

$^{401}$不需要燈光或日光了.
因為主神的光照他們.
他們要作王.
你和我要作王.
直到永永遠遠.
天使對我說.
這的話是可信的真實的.
是真的.
主就是眾先知的靈的神.
差遣了他的天使.
把那快要發生的事.
指示他的僕人.
看那我必快來.
遵守書上預言的.
是有福利的.
所以上帝.
應許我們基督會.
榮耀他回來.
你和我會一同顯現.
好了講完這裡.
既然你知道屬靈生命.
最高境界是怎樣得到.
我們怎樣去有些應用呢.
第一.
就是看開一點.
你現在覺得自己很低.
或者別人傷害了你.
或者很生氣.
有些信徒.
小時候被媽媽罵.
不知為何.
我被媽媽罵就完全忘記了.
他被媽媽罵到現在五十歲.
還記得很清楚.
已經賺了錢已經提早退休.
仍然覺得放不下.
看心理醫生.
心理醫生就越看越麻煩.
因為他經常提起他以前忘記了的事.
令他越來越差.

$^{441}$我基督徒.
為何會這樣呢.
其實最好的心理醫生.
就是聖經.
看開一點.
你已經這麼好了.
你已經與基督一同創造神.
什麼慣例都沒有.
什麼恩怨都沒有.
什麼怒氣都沒有.
前途光明最好.
現在又找堂.
看開一點.
不要這麼肉緊.
第二.
既然人生短暫.
來勢無窮.
你就要投資在永恆.
有位弟兄.
那次我講到幾個月前.
有位弟兄在教會講到.
他完了之後跟我聊天.
他約我去他家.
他和他太太約我們去他家.
他說:羅牧師我要支持你台灣的士工.
每一年四萬元.
三年.
他說:為何我要說三年呢.
因為我現在八十幾歲.
三年之後.
我都不會在這個世上.
我說:我為你祈禱.
求神給你很多很多年.
不是要他每一年支持我四萬元.
而是.
他.
其實我們的生命都是在乎上帝.
他懂得投資永恆.
第一次.
其實我都是投資永恆.

$^{481}$為什麼我要這麼辛苦.
我今年七十三歲.
本來應該要交盤.
不是我不想交.
我要找接班人.
香港警察總幹事的位置.
我要找接班人.
台灣商與教會工廠主任.
我要找接班人.
還沒有找到.
但是.
上帝留我在這裡.
對我來說.
都是一個恩典.
因為你能夠事奉主.
你就不要很得意.
以為你給面子上帝.
給你做.
做這個歡迎大使.
這是恩典.
上帝給你彈琴.
給你.
做這個網絡.
全部都是上帝恩典.
我們每一樣東西.
大大小小都是投資永恆.
不單止這樣.
我們看到世界上很多的需要.
上帝都要我們去服務.
社稷.
關心我們周圍困苦的人.
其實最好的方法.
就是從自己教會開始.
全恩堂我覺得真的很美.
我們人數不是多.
不是大.
但我們接觸很多社區裡面的困難.
新移民.
有些少數民族.
還有一些.

$^{521}$貧窮困苦的人.
有很多需要.
每個人都需要福音.
這是很接觸他們的機會.
這是你值得投資永恆.
知道人生短暫.
第三.
檢視你的生命.
究竟是不是讓主掌管.
你有沒有放手.
讓耶穌基督帶領你.
你願不願意.
聽耶穌基督的話.
等於你開一部車.
一部車有兩個座位.
一部跑車.
耶穌基督跟你一起.
究竟是祂開車還是你開車.
祂是你的客人.
還是你跟著祂.
基督徒就是.
基督的跟隨者.
我們將生命交給主.
當我們將生命完全交給主的時候.
就是我們完全與基督.
一同藏在神裡.
那時候就再沒有掙扎.
再沒有quota.
再沒有自己自怨自危.
因為你.
與基督.
連結在一起.
一個生命一個心跳.
現在活著的不再是你.
而是基督在你裡面活著.
我們一起低頭祈禱.
這個奇妙的聖經.
讓我們知道你愛我們.
你為我們預備了新生命的寶貴.
是叫我們能夠在你裡面.

$^{561}$得著安息.
得著穩妥.
不是靠自己.
你又應許我們榮耀.
將來是多麼美妙.
我們實在讚美主.
我們沒有辦法.
主求你的靈.
今天感動我們在座每一個.
如果我們當中有怨恨.
有不滿有困難有掙扎的.
你讓我們知道.
我們得救是在乎.
皈依安息.
我們得力是在乎平靜安穩.
我們不是靠勢力.
是靠你的靈來成事.
祈禱感恩奉耶穌基督的名.
\newpage



\section{詩篇 127:1-5}
\label{sec:032141}
\textbf{建造的人不枉然勞力}~~~~李美儀傳道
\newline
\newline
連結: \href{http://www.church.com.hk/acms/content.asp?site=cdc\&op=show\&type=product\&code=032141\&layout=sermon}{\texttt{032141}} ~~~~ 語音日期: 2026-05-17
\newline
\newline
\hyperref[sec:032140]{\small{< < < PREV SERMON < < <}}
~
\hyperref[sec:index_chronic]{\small{[返順時目]}}
~
\hyperref[sec:032142]{\small{> > > NEXT SERMON > > >}}
\newline
\newline
詩篇 127:1-5
\newline
\begin{longtable}{cl}
\hline
\hline
章節 & 經文 (和合本修訂版)\\
\hline
127:1 & \begin{tabularx}{0.7\textwidth}{X} 若不是耶和華建造房屋，建造的人就枉然勞力；若不是耶和華看守城池，看守的人枉然警醒。 \end{tabularx} \\ \\ \relax
127:2 & \begin{tabularx}{0.7\textwidth}{X} 你們清晨早起，夜晚安歇，吃勞碌得來的飯，本是枉然；惟有耶和華所親愛的，必叫他安然睡覺。 \end{tabularx} \\ \\ \relax
127:3 & \begin{tabularx}{0.7\textwidth}{X} 看哪，兒女是耶和華所賜的產業，所懷的胎是他所給的賞賜。 \end{tabularx} \\ \\ \relax
127:4 & \begin{tabularx}{0.7\textwidth}{X} 人在年輕時生的兒女好像勇士手中的箭。 \end{tabularx} \\ \\ \relax
127:5 & \begin{tabularx}{0.7\textwidth}{X} 箭袋充滿的人有福了！他們在城門口和仇敵爭論時必不蒙羞。 \end{tabularx} \\ \\
[1ex]
\hline
\hline
\end{longtable}
$^{1}$今天是成年主日.
成年報今年的主題是.
「全程投入 共享豐收」.
這幅是我們的展板.
如果大家剛才在外面經過時.
都看到我們同工 部員一起拍照.
我們用十秒的時間.
去欣賞一下這幅展板 好嗎.
我們觀望一下裡面有甚麼元素.
這幅畫呈現了.
一個我們對於事工的展望和遠景.
是一個結果累累的場景.
大家都忙於收割.
共享原子裡面的豐收.
不知道大家在人生的經歷裡.
我們縱使知道豐收是來自上帝.
有沒有曾經有一個疑問.
究竟人的努力和神的供應.
到底誰才是真正的根基.
真正我們結局的來源呢.
相信標準答案.
一定我們都知道次福是來自神的.
但既然如此.
我們是否停留著等神去供應這個狀態.
我們怎樣才可以推動自己或信徒.
能夠甘心擺上 竭力而為.
在風雨裡面我們不言放棄.
知道賞賜是來自上帝的.
但同時我們明白建造的人.
不枉然努力的真理.
今天要分享一個經文.
就是詩篇127篇.
詩人正正反思和回應這個問題.
他用三個很具體的生活場景去切入.
思考著在神的主權底下.
怎樣去理解人努力的成果.
我們開始分享之前.
我們先一起低頭祈禱.
神啊!你是賞賜人的主.
求主你用著今天我們要分享的說話.

$^{41}$來激勵著我們.
當我們每一天努力建造的時候.
我們學習以你為根基.
享受你所給我們的豐盛.
以上的祈禱是奉教主耶穌基督的命求.
阿們.
我們一起讀出今天的經文.
127篇的一到五節.
詩篇120篇.
至到這一個詩篇的最後一節.
我們就來看看.
這首詩篇的結局.
我們就來看看.
這首詩篇的結局.
我們就來看看.
這首詩篇的結局.
我們就來看看.
這首詩篇的結局.
我們就來看看.
這首詩篇的結局.
我們就來看看.
這首詩篇的結局.
我們就來看看.
這首詩篇的結局.
我們就來看看.
這首詩篇的結局.
我們就來看看.
這首詩篇的結局.
我們就來看看.
這首詩篇的結局.
我們就來看看.
這首詩篇的結局.
我們就來看看.
這首詩篇的結局.
我們就來看看.
這首詩篇的結局.
我們就來看看.
這首詩篇的結局.
我們就來看看.
這首詩篇的結局.

$^{81}$我們就來看看.
這首詩篇的結局.
我們就來看看.
這首詩篇的結局.
我們就來看看.
這首詩篇的結局.
我們就來看看.
這首詩篇的結局.
我們就來看看.
這首詩篇的結局.
我們就來看看.
這首詩篇的結局.
我們就來看看.
這首詩篇的結局.
我們就來看看.
這首詩篇的結局.
我們就來看看.
這首詩篇的結局.
我們就來看看.
這首詩篇的結局.
我們就來看看.
這首詩篇的結局.
我們就來看看.
這首詩篇的結局.
我們就來看看.
這首詩篇的結局.
我們就來看看.
這首詩篇的結局.
我們就來看看.
這首詩篇的結局.
我們就來看看.
這首詩篇的結局.
我們就來看看.
這首詩篇的結局.
我們就來看看.
這首詩篇的結局.
我們就來看看.
這首詩篇的結局.
我們就來看看.
這首詩篇的結局.

$^{121}$我們就來看看.
這首詩篇的結局.
我們就來看看.
這首詩篇的結局.
我們就來看看.
這首詩篇的結局.
我們就來看看.
這首詩篇的結局.
我們就來看看.
這首詩篇的結局.
我們就來看看.
這首詩篇的結局.
我們就來看看.
這首詩篇的結局.
我們就來看看.
這首詩篇的結局.
我們就來看看.
這首詩篇的結局.
我們就來看看.
這首詩篇的結局.
我們就來看看.
這首詩篇的結局.
我們就來看看.
這首詩篇的結局.
我們就來看看.
這首詩篇的結局.
我們就來看看.
這首詩篇的結局.
我們就來看看.
這首詩篇的結局.
我們就來看看.
這首詩篇的結局.
我們就來看看.
這首詩篇的結局.
我們就來看看.
這首詩篇的結局.
我們就來看看.
這首詩篇的結局.
我們就來看看.
這首詩篇的結局.

$^{161}$我們就來看看.
這首詩篇的結局.
我們就來看看.
這首詩篇的結局.
我們就來看看.
這首詩篇的結局.
我們就來看看.
這首詩篇的結局.
我們就來看看.
這首詩篇的結局.
我們就來看看.
這首詩篇的結局.
我們就來看看.
這首詩篇的結局.
我們就來看看.
這首詩篇的結局.
我們就來看看.
這首詩篇的結局.
我們就來看看.
這首詩篇的結局.
我們就來看看.
這首詩篇的結局.
我們就來看看.
這首詩篇的結局.
我們就來看看.
這首詩篇的結局.
我們就來看看.
這首詩篇的結局.
我們就來看看.
這首詩篇的結局.
我們就來看看.
這首詩篇的結局.
我們就來看看.
這首詩篇的結局.
我們就來看看.
這首詩篇的結局.
我們就來看看.
這首詩篇的結局.
我們就來看看.
這首詩篇的結局.

$^{201}$我們就來看看.
這首詩篇的結局.
我們就來看看.
這首詩篇的結局.
我們就來看看.
這首詩篇的結局.
我們就來看看.
這首詩篇的結局.
我們就來看看.
這首詩篇的結局.
我們就來看看.
這首詩篇的結局.
我們就來看看.
這首詩篇的結局.
我們就來看看.
這首詩篇的結局.
我們就來看看.
這首詩篇的結局.
我們就來看看.
這首詩篇的結局.
我們就來看看.
這首詩篇的結局.
我們就來看看.
這首詩篇的結局.
我們就來看看.
這首詩篇的結局.
我們就來看看.
這首詩篇的結局.
我們就來看看.
這首詩篇的結局.
我們就來看看.
這首詩篇的結局.
我們就來看看.
這首詩篇的結局.
我們就來看看.
這首詩篇的結局.
我們就來看看.
這首詩篇的結局.
我們就來看看.
這首詩篇的結局.

$^{241}$我們就來看看.
這首詩篇的結局.
我們就來看看.
這首詩篇的結局.
我們就來看看.
這首詩篇的結局.
我們就來看看.
這首詩篇的結局.
我們就來看看.
這首詩篇的結局.
我們就來看看.
這首詩篇的結局.
我們就來看看.
這首詩篇的結局.
我們就來看看.
這首詩篇的結局.
我們就來看看.
這首詩篇的結局.
我們就來看看.
這首詩篇的結局.
我們就來看看.
這首詩篇的結局.
我們就來看看.
這首詩篇的結局.
我們就來看看.
這首詩篇的結局.
是取決於耶和華是否在當中.
耶和華是否參與在其中.
所以我們回望一百二十七篇的背景.
私人所處的是一個農業的社會.
他很深明耕種之道.
播種施肥都是農夫的責任.
所以耕種的人有耕種人的責任.
但豐收的收穫是在乎上帝.
人的建設固然是真實的.
我們的努力擺上是真實的.
不是我們所做的一切被上帝取代了.
而是我們的擺上是在上帝的主權底下.
神必須站在這一切的之前.
神的份,神的責任,神的崗位.

$^{281}$我們是不可以取代的.
這正正是當時以色列民屬神的子民.
一個很重要的信仰宣告.
是他們建造的根基.
因為當時迦南地充滿著巴力的敬拜.
巴力正正是掌管著雨水和豐收.
外邦人相信的神子.
所以這些外邦人以祭祀,獻祭.
甚至畢業的儀式去換取豐收.
但在以色列人心中他們很清楚這些神是假神.
所以為了要明證上帝的真實.
他們必須抓緊這個很重要的價值觀.
以這個價值觀去行事生活.
以致他們能夠在外邦人面前.
去見證耶和華是掌管天地的主.
我們不需要靠巴力.
我們的神才是真正的神.
所以在下半節裡面.
詩人繼續說.
若不是耶和華看守城池.
看守的人就枉然警醒.
警醒就是一個軍事的用詞.
很明顯代表小心翼翼.
隨時警戒.
因為要留意有沒有敵人入侵.
勞力就是代表著日間的工作建造.
警醒就代表夜間持續的看守.
日與夜就是代表一個圓滿的保障.
不過我們不可以不承認.
我們人是很有限的.
我們不能夠日夜去看守.
這個幾乎是不可能的.
我們不能夠靠自己去持守.
這個日夜的看守.
可能我們都曾經試過照顧家人.
你最多可能陪他不睡一天.
可不可以三天不睡.
五天不睡.
可能你只是坐在他旁邊休息.
但你不可以睜大眼睛24小時看著他.

$^{321}$照顧一個人都不能夠可以守護他日夜.
更何況我們日常這麼多的事務.
因為人是無法超越我們受造物的限制.
我們是被造的.
我們是受造物.
所以我們一定是有限制.
很正常.
所以面對著我們今天的身體.
工作,事業,服事.
我們不能夠完全靠自己是正常的.
因為還有一份在上帝那裡.
所以斯人說建造的人不枉然努力.
是因為他以基督,以神為根基.
先有神的恩典作為根基.
人的努力才有成果.
因為我們等待能夠遇見上帝給我們的結果.
去到第二部分.
去到第二節.
人不枉然努力.
是因為神給予人平靜安穩.
在第二節裡說到.
「你們清晨早起,夜晚安歇.
赤露碌德來的法本是枉然.
惟有耶和華所親愛的.
必叫他安然睡覺」.
古代的以色列人常常處於戰爭,備戰.
不夠糧食,望天打卦的狀態.
通常古代不夠資源都是這樣.
所以其實他們在各方面.
內心裡面存在很多不安.
例如他們耕田會擔心有天災.
天氣不好.
積蓄了之後怕被人掠奪.
城牆有可能被人攻破.
以致他們被侵略.
最後可能生活很多事情.
令他們落入荒謬的狀態.
「拼命努力,卻依然不安」.
這就是當時百姓的寫照.
幾千年過去了.

$^{361}$今天我們是不是都是這樣呢?.
現代人早出晚歸.
我們工作.
甚至現在很怕被AI取代.
有一天我們去到農部.
全部都是AI機械人給你點心.
可能我們很擔心.
有一天我們的職業被AI取代.
我們的積蓄.
因著經濟的不好,不穩定.
或者一些指標的不同.
政策的改變.
令我們所擁有的財富貶值,安保.
甚至我們內心的平穩安靜.
很多人幻想焦慮症,抑鬱症.
驚恐症.
甚至我們對著世間信息萬變的社會.
甚至一些我們不能夠防止.
靠自己去面對的突如其來的生病.
其實都是我們今天現代人.
拼命努力之後依然不安的狀態.
所以幾千年的過去.
人內心的限制都是一樣.
所以當時有人宣告一個很震撼的訊息.
就是他唯有耶和華所親愛的.
必叫他安然睡覺.
其實這兩行的字是四個希伯來文字.
中文很難用四個字直接翻譯.
最簡單就是神給一個安睡你.
神給一個安睡你.
神給安睡你.
這個安睡是一個名詞.
好像一份禮物.
詩篇三十六篇說.
他們必因你殿內的豐盛得以飽足.
在我們很熟悉的二十三篇中.
耶和華是我的牧者.
我必不自缺乏.
他使我躺臥在青草地上.
兩節的經文都說到.

$^{401}$當人當信徒進入神的婉語時.
他不單有飽足.
他能夠躺臥,能夠平靜.
這個景象就像剛才一百二十七篇第二節所說的.
跟耶和華給一個安睡給人的情景相似.
神給人心裡平靜安穩.
其實詩人不單是說一個作息的狀態.
不是說我們信了耶穌之後.
我們每晚都睡得著.
而是他將這個睡覺的模式.
借用這個形象去說明.
這個平安只有上帝耶和華能夠給我們.
只有神能夠去讓人內心安靜.
親愛的弟兄姊妹.
今天我們會不會都意識到.
我們有神的供應和保護.
你就可以安心無憂呢?.
在睡覺之前.
你會不會跟自己說.
我在耶和華的供應和保護下.
我現在可以無憂無慮地躺下.
我們不是在乎最終的結果.
而是我們知道這一刻.
能夠平靜我們內心波濤的.
就只有上帝的同在.
也很確信即使我們孤單一人.
神都會陪著我們.
縱使詩篇127篇的背景.
是進入神的聖殿的場景.
就像舊約當中說到進入神的殿宇.
進入錫安.
進入聖殿的時候.
人就有這份安靜.
但很可惜今天我們不是在聖殿時代.
難道我們今天要進入教會.
只有星期六,日才有平安的信徒.
或者我們要見到牧師,傳道人.
我們才能夠經歷那份進入神殿的平安.
其實作者背後所帶出的.
無論任何時代和背景.

$^{441}$詩人其實是很想說.
我們要進入神的裡面.
進入基督當中.
神已經為我們釘上十架.
我們每時每刻都在基督的裡面.
在約翰福音說到.
我將這些事告訴給你.
是要叫你在我裡面有平安.
臣所指出人以外的平安.
必在基督耶穌裡.
保守你們的心懷意念.
在耶穌基督的裡面.
我們回想一下.
回想一下耶穌的狀態.
很多時候耶穌的狀態是怎樣的呢?.
是不是都很平靜安穩呢?.
英國作家理查·卡特穆斯.
在他的著作其中一本書裡面說到.
這本書叫《在俗修道都史靈修萬部》.
這本書當中他分享了一個經歷.
他說他每當壓力難當的時候.
他會做一件事.
不知道你們壓力難當的時候會做什麼呢?.
是不是去吃東西呢?.
當他壓力難當的時候.
他很喜歡去英國國家的博物館.
英國國家的美術館.
看一幅畫.
是由荷蘭畫家格雷特去畫的.
這幅畫叫《大祭司前的基督》.
Christ before the High Priest.
理查穆斯去形容畫中耶穌的形象.
他說基督的雙手被綁.
不過相比下指控人的手.
指著他指控耶穌基督.
同時他注意到耶穌身體的那種靜態.
彷彿他用身體拒絕了這些惡意的攻擊.
並且拒絕作出回應.
甚至他覺得耶穌基督在這個畫像裡面.
他的態度是鬆弛的.

$^{481}$甚至有些慵懶.
跟這個看著他.
釘住他的人 指責他的審判者.
有一個很大的對比.
甚至他最後形容.
雖然耶穌是被審判.
但指控者才是真正惶恐 懼怕的人.
因為他拒絕了.
讓真理進入他的心裡.
安然是基督的記號.
是給我們信徒信仰的記號.
因為安然就是代表著.
反映我們信徒對神的信靠.
今天你安不安然?.
安不安心?.
以前聽媽媽說抱小朋友.
我不是媽媽.
我問姐妹們.
小朋友哭的時候.
用手拍一下他的胸口.
拍前面.
不然拍後面.
拍前面.
你摸摸自己.
不要摸旁邊那個.
摸自己.
搓自己的手 暖一點.
不要冷冰冰.
拍自己.
耶穌基督就是這樣.
祂讓你安心.
今天我們是否一邊交託給神.
一邊深夜反思 輾轉反側.
想著怎樣出路.
一邊我們知道上帝在建造.
但我們又不自覺將建造的成敗.
壓在自己身上.
忘記了有上帝的一份.
所以詩人這句溫柔的宣告.
唯有耶和華可以給你一個安睡.

$^{521}$這個安然.
就是來自於這個掌管房屋.
守護城牆.
賜下兒女福氣的耶和華.
我們不要再被這些焦慮吞噬.
記得更高之上有上帝.
有一份在上帝那裡.
當我們勞碌.
為著我們生活奔波的時候.
要記得不是明天睡醒這件事就解決.
而是我們知道明天仍然在上帝的掌管下.
祂沒有離開.
祂說你睡醒之後就離開了.
上帝仍然在你的身邊.
所以在第三個部分.
詩人繼續說.
他說建造的人不枉然勞力.
是因為神的豐盛傳自後代.
這裡說兒女是耶和華的產業.
產業很明顯.
就是耶和華所給信徒承諾的一些東西.
例如迦南地.
應許之地.
神對信徒的承諾.
這個承諾是包括.
祂會看管.
祂會守護.
祂甚至會使其豐盛.
這是祂承諾的責任.
所以當我們說.
兒女都是耶和華給我們的產業的時候.
就是說這些兒女是屬於耶和華.
由耶和華負責任.
負上全責.
祂是應許下.
這些生命可以成長.
最後成為豐盛的賞賜.
所以在經文裡面很強調.
耶和華所賜的產業.
很強調所懷的胎是祂給的.

$^{561}$所賜所給.
兩個重複都是強調.
生命是來自上帝.
祂會負責任.
所以對於以色列人來說.
兒女是很真實.
明明確確.
恩典的明證.
我們留意一下這個詩篇.
從第一節第二節.
說到神是根基.
豐盛的供應來源.
第二祂保守.
保守這些豐盛.
存留在你的生命裡面.
最後祂說到.
兒女成為豐盛的核心.
最後人成為豐盛的供應者.
從源頭經歷.
到最後你自己成為供應者.
勇士手中的箭.
其實就是提到這一班的兒女.
有些人很強調.
少年的時候所生的.
因為少年的時候.
家庭就有足夠的時間.
去培養塑造他們.
磨練成為有用的人.
箭袋裡面充滿人變為有福.
就是說.
當這個家庭被上帝祝福.
滿有兒女的時候.
是有一個可見的豐盛在裡頭.
城門口仇敵.
其實就代表一個戰爭的場景.
整體而言.
其實就是說.
當這些敵人來到.
我們要打仗的時候.
兒女就成為父親當中的裝備.

$^{601}$他們可以護衛著這個國家.
同時可以向著這些敵人進擊.
成為家園的保護者.
所以我們見到這個轉變.
從一個被動的接收者.
去到神的賜予.
將我們的後代模造成一個強壯的建造者.
我們可以保護其他人.
神的賜予最終的目的.
就是要我們信徒成為強壯的建造者.
要我們祝福其他人.
能夠有這個能力.
去將福氣傳遞開去.
在十七世紀.
有一位姊妹叫蘇珊娜·韋斯利.
可能大家都認識她.
她一生中.
生過十九個小朋友.
不過很不幸有九個夭折離開.
最後只剩下十個.
蘇珊娜的人生.
她如何去養育她的孩子呢?.
她每天早上三小時.
下午三小時.
還有每個小朋友.
每個星期.
有一個一個小時的單獨會面.
早上三小時.
下午三小時.
十個小朋友.
一個星期.
會有一個小時單獨見面.
會教她聖經.
唱詩歌.
會教她其他的學習.
一個很美好的圖畫.
其實這位姊妹都經歷很多.
她的丈夫曾經入獄.
她家庭極度貧窮.
曾經房屋被燒.

$^{641}$甚至她自己經歷過.
喪失小朋友的經歷.
後來她在廚房裡.
帶領著家庭崇拜.
她一家十幾口.
慢慢她的鄰居一起參與.
結果這個家庭聚會有二百多人.
在她養育孩子的過程裡.
她有兩個小朋友.
可能大家會認識.
一個是約翰.韋斯利.
成為秦道運動的創始人.
向超過百萬人傳福音.
第二個是查理.韋斯利.
他一生寫過超過九千首詩歌.
他一生完美去詮釋了.
一百二十七篇三至五節的內容.
兒女是耶和華所賜的產業.
他不是靠人的方式去預養他們.
而是他知道責任在上帝那裡.
他管理著這些兒女.
將神的話語進入他們的生命裡.
以致這些兒女成為真正豐盛的產業.
將來能夠祝福其他人.
這正是不枉然奴隸最豐盛的詮釋.
使後代成為有力的建造者.
蘇格蘭首府愛丁堡的城市盾輝很美.
有盔甲,城堡,有隻動物,有個人.
上面還有一句說話.
翻譯為英文就是.
昨天弟兄跟我說.
他說城門水塘也有一句.
你們行山時幫我看看.
這句說話正正來自詩篇一百二十七篇.
過往愛丁堡這個地方.
它的市民,它的領袖.
都是以這一句上帝的說話.
作為他們前進的動力.
成為他們建立城市的根基.
成為他們做人的價值觀.

$^{681}$不過在2022年一個統計調查出來.
有51\%的蘇格蘭人表示不再相信宗教.
在2001年開始.
蘇格蘭的教會遞減一半.
而16至24歲的人.
超過62\%的人.
他覺得自己是沒有宗教的.
而最終這句說話.
就淪為一句文化遺產.
仍然標誌在他們的盾輝上面.
但再不大有能力.
今天我們城鎮有沒有盾輝呢?.
我們有一個Logo.
代表我們有一個城鎮的教會.
不過我們有一個電機石.
大家知不知道在哪裡?.
這裡是哪裡?.
在外面.
在外面,就是我們的車進來那裡.
耕田外面那裡.
我拍不到那個電機牌.
那天也撞到幾個童工.
我最後爬了進去.
拍到了.
全都爬了進去就拍到了.
1963年11月17日.
我們一起讀這八個字好不好?.
1,2,3.
機電盤石,種播美田.
盤石就是基督.
機電盤石就是將根基紮在基督裡面.
種播美田.
就像以色列人一樣.
他們所有的耕種農作.
都是依靠著上帝這份供應.
當年是這樣的.
未有的時候是這樣的.
我相信這張相片.
已經是1963年之前.
在城鎮剛興建的時候.

$^{721}$有沒有認出是誰你認識的?.
看不到吧.
剛興建,又低頭.
正在禱告.
又擔了雨傘.
他們今天在哪裡呢?.
我相信這些的牧者.
長執.
弟兄姊妹.
可能大部分都不在這裡.
是天家.
不在我們在座當中.
在興建的時候.
他就將這個根基交給上帝.
今天.
這個堂子.
就是當年先賢憑信心.
相信有神的份.
有人的份建造而成.
雖然他們不在這裡.
但你就是他們的風聖延續者.
你就是他們風聖的儲備.
將神的風聖一代一代.
傳流在教會裡.
我們的努力是不會枉然的.
因為我們有上帝在裡面.
神的風聖不會只停留在紙張上.
不只是我們出面的展板上.
而是進入我們的生命裡面.
成為我們就是那個風聖.
我們延續下去.
將這份風聖祝福.
帶給其他人.
建造的人不枉然努力.
建造的人就是你.
但願我們今天教會.
在神的風聖供應下.
建造,增長,茁壯.
成為這個城市永久的祝福.
不會停留.

$^{761}$有一天成為一個標語.
毫無力量.
也是永遠帶著神的應許力量.
影響著這個地區.
以致我們身邊未認識主的人.
能夠認識神.
我們一起低頭祈禱.
親愛的天父.
求你幫助我們以你為首.
在工作和生活裡與你同工.
我們知道在基督裡.
你賜下安然.
使我們不被焦慮纏繞.
求你幫助我們.
讓我們忠心栽培你所託付的生命.
使下一代成為你手中的箭.
願我們全情投入.
能夠一同共享神所賜的豐收.
以上的祈禱.
是奉靠著耶穌基督的名求.
阿們.
\newpage



\section{約書亞記 4:8}
\label{sec:032142}
\textbf{聖經人聖經事:迦勤 專心跟從耶和華}~~~~黃克勤牧師博士
\newline
\newline
連結: \href{http://www.church.com.hk/acms/content.asp?site=cdc\&op=show\&type=product\&code=032142\&layout=sermon}{\texttt{032142}} ~~~~ 語音日期: 2026-05-19
\newline
\newline
\hyperref[sec:032141]{\small{< < < PREV SERMON < < <}}
~
\hyperref[sec:index_chronic]{\small{[返順時目]}}
~
\hyperref[sec:032143]{\small{> > > NEXT SERMON > > >}}
\newline
\newline
約書亞記 4:8
\newline
\begin{longtable}{cl}
\hline
\hline
章節 & 經文 (和合本修訂版)\\
\hline
4:8 & \begin{tabularx}{0.7\textwidth}{X} 以色列人就照約書亞所吩咐的做了。他們按以色列人支派的數目，從約旦河中取了十二塊石頭，正如耶和華所吩咐約書亞的。他們把石頭帶過去，到他們所住宿的地方，就放在那裡。 \end{tabularx} \\ \\
[1ex]
\hline
\hline
\end{longtable}
$^{1}$各位姐妹,讓我們一起讀聖經.
《生命書宣教網絡》.
和我們一起看聖經人物,聖經故事.
最重要的是看到背後上帝的手工作.
看到祂的心意.
今天的題目是.
迦勒這個人.
迦勒專心跟從耶和華.
「專心跟從耶和華」這一句.
真的是聖經裡面的字.
迦勒是「專心跟從耶和華」.
是很寶貴的.
這個人物是很值得我們去學習的.
《約書亞記》第十四章第六節.
「那時猶大人來到吉甲見約書亞」.
其中有基尼洗族耶夫尼的兒子迦勒.
就跟約書亞說.
其實是迦勒特意去找約書亞.
去吉甲找約書亞.
吉甲是他們駐守的安營地方.
總部在吉甲.
有二百萬人在那裡住.
因此他們分散比較遠.
今天來了一個猶大支派的人.
他是猶大裡面的基尼洗族耶夫尼.
基尼洗族有一個人叫做耶夫尼.
他的兒子叫迦勒.
迦勒和約書亞其實是認識到.
我相信是有感情的.
生死之交.
為什麼會這樣呢?.
就說給大家聽.
來到找約書亞就說.
「耶和華在迦底斯巴尼亞」.
「指著我與你對神人摩西所說的話」.
你也知道了.
你記得了.
「耶和華在迦底斯巴尼亞」.
迦底斯巴尼亞是一個地方名字.
不過應該是兩個字加起來.

$^{41}$最主要聖經出得多的是迦底斯.
「指著耶和華對摩西所說的話」.
「我與你」.
即是我和你.
聽到神這樣說的.
你也知道了.
當日是這樣說的.
忘記了.
迦勒這個字是專心跟隨的.
很厲害.
迦勒這個字.
比較少人改這個名字.
做英文名或者做名字.
為什麼呢?.
因為Khalab.
Khalab這個字.
Khalab是好名字.
但是Khalab這個字在希伯來文是狗.
是一隻狗.
為什麼會改名叫狗呢?.
都說了.
嬰兒出生不會自己改名.
爸爸改名.
或者媽媽改名.
兩個人談了改名.
很有可能未必改住.
後來見到「哎呀,這隻狗仔」.
於是改名叫狗.
狗仔.
「哎呀,我兒子好像狗仔」.
跑一跑很有趣.
可能叫狗.
Khalab是狗.
不過改得很好.
狗呢.
有時候狗真的在某些角度看.
狗真的比人好.
狗很忠心.
很忠誠.
對主人很忠誠.

$^{81}$狗就是迦勒.
Khalab就跟余約瑟談了.
當日在迦底斯的時候.
神對指著我和你.
對摩西說的話.
你記得嗎?.
記得.
耶和華的僕人從迦底斯帕尼亞.
打發我們去窺探迦南地.
那個時候我是四十歲.
我按著心意來回報摩西.
心意的意思.
就是喇叭.
喇叭是內心或者理解.
我是照著我真心內心和我想理解的.
來回報給摩西.
不過然而.
和我上去的眾弟兄.
就不是.
是使百姓的心消化.
這個消化的意思就是.
唉.
不行了.
怎麼去到.
做不到.
打不到.
那裡的人很大隻.
又強壯.
我們打不過他.
但我專心跟從耶和華的神.
他們這麼說.
不過我專心跟從耶和華的神.
我是知道.
我跟從神.
神說可以就可以.
儘管去吧.
為什麼?.
神叫我們得到這塊地.
神由我們的先祖亞伯拉罕就答應了.
所以我們一定可以.

$^{121}$很有信心.
我是專心跟從耶和華的神.
祂叫我去拿這塊地.
我就一定可以.
所以我窺探完就說.
不怕的.
可以去.
這塊地好.
這樣子.
其實當時他們經過了埃及出來曠野.
預備進入迦南.
當時他們是沒有過約旦河的.
就在迦南地的南部.
南地.
派人上去北部.
進入迦南.
打聽一下.
有些地方.
他一進去的地方.
馬上看到什麼人呢?.
就看到阿納族人.
非利士那邊.
阿納族比較近.
哇.
阿納族人原來個個都是.
高頭大馬.
八九尺的.
八尺都算矮仔的.
我走過去都算矮仔的.
阿納族人是很大隻的.
而.
譬如.
看到加特人.
加特人又很大隻的.
例如我昨天都講了.
加特呢.
就是非利士地.
出了個什麼呢?.
哥利亞.
跟大衛打.

$^{161}$哥利亞是綠爪.
九尺多高.
哇.
很大的.
我見到很多人.
摩西多納派十二個探子去.
有十個回來說.
不行了.
沒辦法了.
他們是偉人和偉人的後代.
看到阿納族人.
哇.
不行了.
太大隻了.
我們好像裝猛一樣.
他們像巨人.
不過剩下兩個人.
就是約書亞和加納華.
不怕的.
我們足能得勝.
因為有耶和華同在.
加納華專心跟從耶和華.
我將我內心所見的.
不是鳳凰那些弟兄.
害到每個人心裡面害怕.
當日摩西就起誓了.
跟我說.
你腳踏之地要歸你和你的子孫.
永遠為業.
就是你去過的地方.
我要代表上帝.
上帝會叫我將你去的地方.
給你和你的後裔永遠為業.
為什麼?.
因為你專心跟從耶和華我的神.
你這麼專心.
第二次說.
第一次是加納自己說的.
我專心跟從耶和華.
我就夠膽說我們可以上去.

$^{201}$這裡是摩西說.
你專心跟從耶和華.
好.
你腳踏之地.
我就特別賜給你和你的後裔.
當然賜是指代表神.
好了.
自從耶和華對摩西說這話的時候.
耶和華是照他所應許的使我存活.
這是加納的看法.
我專心跟隨神.
我跟約書亞里當日一起說.
我們可以去德里法的應許之地.
雖然我們十個兄弟說不行.
不過我們兩個照著心裡面的話實說可以.
其實他們十個去摘了剛剛葡萄季節的葡萄回來.
拿了其他的水果回來給他們看.
這裡看看水果多大.
那一串葡萄是要兩個人抬的.
就是說多豐富多多.
是表示這塊地很豐富.
我們去拿有神同在.
我們足能得勝.
十個兄弟說不行.
很大隻.
於是就說你專心跟隨我這塊地給你.
你去過窺探那塊地給你.
其實耶和華給我這塊地.
他一定會讓我存活.
不然我怎麼拿得到.
耶和華照他所應許的使我存活.
這四十五年.
就是說現在我八十五歲了.
我去窺探的時候四十歲.
現在四十五年了.
八十五歲.
期間以色列人在曠野行走.
可能現在我八十五歲了.
我們一直都在曠野生活.
很艱苦的.

$^{241}$八十五歲了.
但是很感恩.
我還是強壯.
像摩西四十五年前.
打發我去拿天一樣.
無論征戰或者出入.
我仍然有力量.
我的力量那時如何.
現在還是如何.
他在那裡見證和感恩.
就是神給了我.
我跟隨他沒有跟錯.
他按照他的應許.
給我地.
又令我很強壯.
直到今天.
其實同代的人.
只剩下我和李約書亞兩個人.
當代我們同代的人.
全部倒閉曠野.
因為不信神.
並沒有專心跟隨耶和華.
求你將耶和華那日應許我的山地給我.
那裡有阿納族人.
阿納族人.
阿納.
阿納奇人.
阿納奇.
這個奇是讀K的.
阿納奇.
阿納奇就是阿納.
阿納族人.
阿納奇是什麼意思.
就是頸的意思.
叫頸.
那裡有頸人.
頸人可能很有趣.
你見那些高大大隻的人.
特別那些摔角舉重.
那些人很大隻.

$^{281}$頸是很粗的.
你一看上去.
他的頸.
真的像面一樣粗.
很粗.
還要斜著兩肩.
那條頸很粗.
可能這樣.
很大隻.
那裡有頸人.
阿納族人.
並且他們有寬大堅固的誠.
你都曾聽見是這樣的.
現在求你把那塊地給我.
是按著耶和華當日的應許給我.
所以今天迦勒尼找約書亞.
是想告訴他.
你可不可以在分地的時候.
分給十二支派.
可不可以特別把那塊.
是上帝透過摩西應許給我的那塊地.
就是我們最初進去窺探的時候.
見到阿納族人那塊地給我.
這樣.
那裡其實有阿納族人.
現在還有.
還有寬大堅固的誠.
你給我吧.
給你又怎樣.
即是你要和阿納族人作戰.
要攻破那些寬大堅固的誠才獲得.
你見到.
迦勒仍然很有信心.
今天講這個故事是講迦勒.
主角是迦勒不是約書亞.
這句很有趣.
或者耶和華.
照他所應許的與我同在.
我就把他們趕出去獲得那個地方.
這個都明白.

$^{321}$不過最不明白的是「或者」是什麼意思呢.
「或者」即是「或者」.
意思是說不定.
那些人又高大.
個個九尺.
又大又堅固的誠.
給你都沒有用.
我都想過的.
不過說不定.
神幫助我.
照他的應許與我同在.
我把他們趕走.
很有信心.
不過這個「或者」都有個含義.
我會去攻到他們的誠.
把他們趕走.
不過如果不行.
或者神沒有幫助我.
那也算了.
不要緊.
不過我是想照著當日.
我們的耶和華神透過摩西.
在你面前你都知道的.
是應許了這一塊地給我的.
我想去拿.
我想去得.
喂.
老友.
大家兄弟一場.
我不告訴你.
那裡有阿納族人.
九尺高.
另外還有些.
很高大堅固的誠.
你去不到的.
你怎麼攻啊.
或者耶和華照他所應許與我同在.
我就把他們趕走.
我今天八十五歲.
和當日四十歲.

$^{361}$一樣強壯.
眼都沒有花.
老鼠仔還在.
這樣的意思.
你見到這樣說.
妹真的很有信心.
好了.
於是約書亞就為耶夫尼的兒子迦勒祝福.
好啊.
願上帝賜福你.
就把那塊地給你.
就把希伯倫這個地方給他.
希伯倫.
大家聽過希伯倫這個地方在舊有都很出名的.
有呢.
吸他為業.
我本身信主出身.
就用這個地方做名.
說在教堂.
希伯倫堂.
這樣.
好了.
吸他為業.
所以呢.
希伯倫就作了基尼洗族.
耶夫尼的兒子迦勒的產業.
直到今天.
因為第三次講.
用這個字.
他專心跟從耶和華.
是今天主題.
專心跟從耶和華以色列的神.
神是會祝福他的.
希伯倫的意思就是聯合.
United.
Join together.
希伯倫是一切.
我的族人.
我會將他們一齊聯合.
將他們得到.

$^{401}$不過這些希伯倫.
是他們得到之後的生命.
不過現在是回頭記述當日.
約書亞記述當日.
是不是他臨死前寫呢?.
不知道了.
總之他後來記述當日.
當時的地方已經叫做希伯倫.
好了.
希伯倫就是聯合.
Cabral.
Cabral就是聯合共壩.
從前以前的名字.
當他要去攻打的時候.
當說好了.
給你吧.
耶夫尼的兒子迦勒.
我祝福.
然後給你.
那你去吧.
拿去吧.
以前就叫做基列阿爸.
基列阿爸.
基列.
是兩個音來的.
Q讀K.
奇亞.
奇亞阿爸.
奇亞其實是城的意思.
你見到奇亞就是城.
這個字是城.
阿爸的城.
所以叫做阿爸之城.
其實阿爸是甚麼來的呢?.
阿爸是阿納族裡面.
最尊大的一個人物,一族人.
阿爸.
阿爸其實.
阿爸是他們阿納族人的始祖.
所以就叫做阿爸的城.

$^{441}$但是在這裡我用藍色字.
其實希伯倫是沒有這個字的.
只有Who.
Who就是他是阿納族中最偉大.
Gado.
Gado是最偉大的一個人種或人類.
或者一個人.
如果是眾數就是人種人類.
如果是單數就是一個人.
即是說.
阿爸是阿納族一個最偉大的人.
可能是阿爸是阿納族的始祖.
他一直下來的人種就叫做阿爸.
阿爸是阿納族的人.
於是國中就太平沒有爭戰.
就去拿了這塊地.
國中呢.
這個國是甚麼意思.
都未有國.
都未立國.
到蘇羅斯才立國.
國這個字其實是譯得不好.
可以直接說是譯錯了.
是Edwards.
Edwards就是土地.
於是那塊土地就太平了.
因為加納公取了那塊地方.
好了.
我今天跟大家分享的有三點.
改了分享方式.
第一.
你專心跟從耶和華.
是你記得要做的.
我們要專心跟從耶和華.
甚麼叫專心呢.
一心一意跟著看著.
很厲害的.
我們都帶過很多人.
十團去以色列.
沒有甚麼人有去以色列的經驗.

$^{481}$有些是很專心的.
我發覺有好幾個.
有時帶過幾十人的團.
或者是少數很相識認識的小團大團.
有些人是時時.
我都發覺他.
我一站好預備要介紹這個地方的時候.
那個人就站在我旁邊了.
就拿著一本簿.
這樣寫東西.
他很想聽.
專心跟著我.
有些是不見影.
只有這麼少人來到.
還未可以說.
他一回來又會.
哎呀剛才說甚麼.
都不知道何時開始講解好.
有些我都講過很多次.
有些是很不專心的.
他一見到有地方漂亮.
他就馬上拍照打卡.
你和我拍.
我和你拍.
又這樣姿勢又這樣姿勢.
爬上那裡蹲下那裡.
嘩.
很久才兩個人跑過來.
有些就是一見到有裝飾品.
一見到有紀念品.
紀念品店已經不見影.
鎖進去了.
有些是見到店鋪都已經進去了.
接著笑笑口地拿著兩袋東西跑出來了.
來到.
講甚麼呀剛才.
講完啦.
離開這麼遠不專心的.
拍照與購物最為重要.
那不專心的.

$^{521}$專心跟隨的話你才跟得貼的.
主耶穌的十二個門.
為甚麼後來那麼洗德.
因為他們專心貼著主耶穌.
衣食住行一起的.
所以跟得很貼.
他才能夠專心跟隨主.
記得專心跟隨.
跟貼些.
你們會知道整個心意的.
跟得久些更加知道了.
第二呢.
他有信心.
迦納是照著他應許的使我存活.
他是相信神是有應許在我身上.
是神使我存活.
如果神沒有使我存活.
叫我收回去.
我就回去了.
我就不會到今天八十五歲.
仍然跟當日四十歲去.
做探子窺探的時候.
一樣守乖起展.
一樣那麼健康.
你現在給亞納族的地方.
希伯倫或基烈阿巴給我.
我完全不怕.
又高大的亞納族人.
又寬又高的城.
我都不怕.
我趕他們出去.
因為是耶和華的應許.
令到我成功.
第三點.
或者耶和華照他所應許的.
與我同在.
這四個字最重要.
照他的應許與我同在.
有甚麼比神與你同在重要.
有甚麼強得過神與我同在.

$^{561}$就算我奄奄一息快死了.
但神與我同在.
我仍然那麼平安.
面對這樣的敵人.
神的應許與我同在.
我好平安.
好有膽量.
好有勇氣.
就因為耶和華照他的應許與我同在.
或者是否不信得救.
不是.
或者只是說.
不過我尊重上帝的心意.
如果他沒有救我,就沒有救我.
他救我就救我.
他沒有給力量,我就沒有力量.
給我就有.
但以你身為我的朋友.
將來我就火坑了.
喂!一扔你下去,一定死的.
明白嗎?.
他們說.
神會救我脫離這些火.
還有你的手.
即或不然.
如果不救.
我燒死了.
我都不會拜你的金像.
即或不然.
即是這裡的或者.
說不說神可以幫我呢?.
即是說.
他不敢強行說.
神一定幫我的.
不幫我就是他欠我.
他不是.
他沒有欠.
神幫我就是幫我.
就是給我恩典.
沒有.

$^{601}$就是我應得的沒有了.
這個很重要的.
將這三點送給大家.
約瑟拉的十四章說.
迦勒這個人非常好.
他很忠心.
很忠心.
當然.
迦勒這個字就是狗.
希伯來文.
但我不想說.
他忠心到像隻狗一樣.
其實.
狗我真的很欣賞.
也很喜歡.
也很想養狗.
不過香港不方便.
很忠心的.
很忠心.
沒得頂.
真的有時候覺得.
真的比養人容易.
好.
今天跟大家分享到這裡.
我們一起禱告.
多謝我們的天父.
在你帶領的人裡面.
在以色列裡面.
猶大之派裡面.
耶夫尼的兒子.
出現了一個叫迦勒的人.
他專心跟隨耶和華.
只是這一段聖經.
已經說了三次.
專心跟隨耶和華.
他自己說.
我專心跟隨耶和華.
又被神摩西說.
你專心跟隨耶和華.
是稱譽.

$^{641}$最後多謝你.
幫助我們學習迦勒.
他是專心跟隨耶和華.
他滿有信心.
是靠著你的應許而生活.
靠著你的應許而勇敢.
是靠著知道你的同在.
而他相信自己會有力量.
多謝我們的主.
你的同在.
讓我們有迦勒的忠心與勇敢.
這樣的信靠.
實在是很寶貴.
幫助我們這麼有信心.
聚聽祈禱奉耶穌基督的名.
阿們.
\newpage



\section{羅馬書 1:16-32}
\label{sec:032143}
\textbf{以信為始 以信為終}~~~~蕭壽華牧師
\newline
\newline
連結: \href{http://www.church.com.hk/acms/content.asp?site=cdc\&op=show\&type=product\&code=032143\&layout=sermon}{\texttt{032143}} ~~~~ 語音日期: 2025-10-11
\newline
\newline
\hyperref[sec:032142]{\small{< < < PREV SERMON < < <}}
~
\hyperref[sec:index_chronic]{\small{[返順時目]}}
~
\hyperref[sec:032144]{\small{> > > NEXT SERMON > > >}}
\newline
\newline
羅馬書 1:16-32
\newline
\begin{longtable}{cl}
\hline
\hline
章節 & 經文 (和合本修訂版)\\
\hline
1:16 & \begin{tabularx}{0.7\textwidth}{X} 我不以福音為恥；這福音本是神的大能，要救一切相信的，先是猶太人，後是希臘人。 \end{tabularx} \\ \\ \relax
1:17 & \begin{tabularx}{0.7\textwidth}{X} 因為神的義正在這福音上顯明出來；這義是本於信，以至於信。如經上所記：「義人必因信得生。」 \end{tabularx} \\ \\ \relax
1:18 & \begin{tabularx}{0.7\textwidth}{X} 原來，神的憤怒從天上顯明在一切不虔不義的人身上，就是那些行不義壓制真理的人。 \end{tabularx} \\ \\ \relax
1:19 & \begin{tabularx}{0.7\textwidth}{X} 神的事情，人所能知道的，原顯明在人心裡，因為神已經向他們顯明。 \end{tabularx} \\ \\ \relax
1:20 & \begin{tabularx}{0.7\textwidth}{X} 自從造天地以來，神的永能和神性是明明可知的，雖然眼不能見，但藉著所造之物就可以了解看見，叫人無可推諉。 \end{tabularx} \\ \\ \relax
1:21 & \begin{tabularx}{0.7\textwidth}{X} 因為，他們雖然知道神，卻不把他當作神榮耀他，也不感謝他。他們的思想變為虛妄，無知的心昏暗了。 \end{tabularx} \\ \\ \relax
1:22 & \begin{tabularx}{0.7\textwidth}{X} 他們自以為聰明，反成了愚昧， \end{tabularx} \\ \\ \relax
1:23 & \begin{tabularx}{0.7\textwidth}{X} 將不能朽壞之神的榮耀變為偶像，仿照必朽壞的人、飛禽、走獸、爬蟲的形像。 \end{tabularx} \\ \\ \relax
1:24 & \begin{tabularx}{0.7\textwidth}{X} 所以，神任憑他們隨著心裡的情慾行污穢的事，以致彼此羞辱自己的身體。 \end{tabularx} \\ \\ \relax
1:25 & \begin{tabularx}{0.7\textwidth}{X} 他們將神的真實變為虛謊，去敬拜事奉受造之物，不敬奉那造物的主—主是可稱頌的，直到永遠。阿們！ \end{tabularx} \\ \\ \relax
1:26 & \begin{tabularx}{0.7\textwidth}{X} 因此，神任憑他們放縱可羞恥的情慾。他們的女人把自然的關係變成違反自然的； \end{tabularx} \\ \\ \relax
1:27 & \begin{tabularx}{0.7\textwidth}{X} 男人也是如此，放棄了和女人自然的關係，慾火攻心，男的和男的彼此貪戀，行可恥的事，就在自己身上受這逆性行為當得的報應。 \end{tabularx} \\ \\ \relax
1:28 & \begin{tabularx}{0.7\textwidth}{X} 他們既然故意不認識神，神就任憑他們存扭曲的心，做那些不該做的事， \end{tabularx} \\ \\ \relax
1:29 & \begin{tabularx}{0.7\textwidth}{X} 裝滿了各樣不義、邪惡、貪婪、惡毒，滿心是嫉妒、兇殺、紛爭、詭詐、毒恨，又是毀謗的、 \end{tabularx} \\ \\ \relax
1:30 & \begin{tabularx}{0.7\textwidth}{X} 說人壞話的、怨恨神的、侮辱人的、狂傲的、自誇的、製造是非的、忤逆父母的、 \end{tabularx} \\ \\ \relax
1:31 & \begin{tabularx}{0.7\textwidth}{X} 頑梗不化的、言而無信的、無情無義的、不憐憫人的。 \end{tabularx} \\ \\ \relax
1:32 & \begin{tabularx}{0.7\textwidth}{X} 他們雖知道神判定做這樣事的人是該死的，然而他們不但自己去做，還贊同別人去做。 \end{tabularx} \\ \\
[1ex]
\hline
\hline
\end{longtable}
$^{1}$這錄音或錄影的版權是屬於基督教宣導會不各堂所有.
我們的屬靈傳統之一是以真理教導為本.
歡迎您將此錄音或錄影的連結傳給親友.
懇請注意,未經授權,請勿轉錄或剪輯本錄音或錄影.
甚願好像《仕途行傳》20章32節下所說.
這道能建立你們,叫你們和一切聖聖的人同得基業.
我們都會在這個時候想念到很多我們的先輩,前輩.
他們如何在艱辛當中開展了教會的服事.
我講一點歷史,大家知道73年之前的10月12日.
就是北宣的開堂日,明天的日子.
當時很多人都會問,為什麼會有人來這個地方傳福音呢?.
其實大家如果留心一些香港的歷史,你會知道.
在50年代的時候,因為中國內地政權的轉變.
所以大量的難民湧來香港.
而香港政府當時就將難民集中在一條鏡嶺的地方.
那裡變成一個難民集中營的地方.
那個地方大家不要以為現在很快地在地鐵兩步就去到了.
其實當時是一些很荒蕪的山嶺,是沒有交通工具去到那個地方.
我們稱之為一條鏡嶺,其實因為這個地方以前的名稱叫做吊頸嶺.
因為是一個荒蕪的山嶺,所以有些人自殺就躲在那裡吊頸.
當時的難民在那個地方,你可以想像是很艱難的生活.
但是當時的九龍塘宣導會的貝岡道牧師帶同婦女會的姐妹們.
他們就去那個營地去救濟,去關心,也在當地傳福音.
後來他們就邀請宣導會的西差會同工,就在那裡開始了教會的工作.
但是最多人信主的日子,短短兩年之間教會有超過二千人信主.
我想到今天都未必有這樣的情況,很多人信主.
但是在信主之後就發覺很困難,但是也有神所興起的牧者.
教會開始了一個叫做當時的宣導聖經學校.
他們訓練神學生來承擔很多的牧職的工作.
這些神學生他們在暑假期間本來應該馬上找機會去賺一些錢.
生活很艱難,他們都是當地的一些居民,一些逃離香港的難民.
但是他們全部都願意用那段時間在香港四處去走,說要能夠傳福音.
其中一班就來了北角,當時來北角都很艱難.
沒有交通,搭船,搭一些街道,來到這個地方,來到北角碼頭地區傳福音.
就是這樣,有很多的朋友是當時的信徒主.
經過一段時間之後,這班神學生他們要回到學校上課.
但是有很多信徒主的北角的街坊怎麼辦呢?.
神就安排一些宣教士和一些弟兄姊妹.
就在這個地方開始了宣導會北角堂.
他們當時租用了在北角道十四號的地方.

$^{41}$在現在北角道,麻布道交界的地方.
租用一個地下的單位來開始教會的崇拜.
在當時我們,我自己都未曾參與教會.
但是後來看到有關的消息的時候,心裡是很感動的.
是很多的前人,他們是真的付上了很大的代價.
很辛勞來帶領在這個地區的人認識住.
他們當時是每一天都在北角碼頭的地方.
是趁著很多船隻下船上船,很多的乘客.
也有些時候,傍晚的時候,很多人在乘涼的時候.
就會公開用一些手風琴來拉奏.
吸引一些人來聽福音.
然後過程當中,很多人信了住.
我自己是1969年來到北角宣導會的.
當時我來到的時候,心裡面開始知道這些歷史.
心裡就很感動,我知道如果不是有這一群人.
今天我也不會像今天那麼順暢.
不過到14號,不知道大家有沒有留意.
今天已經是完全不同了.
不過當你經過的時候,你會發現一座建築物.
外面寫了四個大英文字母,叫做NPAC.
什麼NPAC?北角宣導會.
North Point Alliance Church.
你說為什麼還有這間教會呢?.
當然已經不是以前的教會.
已經傳動建築物,已經拆卸重建.
但是剛好建了一座建築物.
這座建築物的英文名字就是NPAC.
大家可能經過的時候會看到.
我們也會問一個問題,為什麼上帝會這樣做呢?.
或者是上帝要我們不要忘記.
在那個地方,有一群人.
他們是付上生命的代價.
他們毀身在福音上面,來建立教會.
而我們今天不在那個地方聚會.
但是經過的時候,也要再思想上帝自己的恩惠.
今天是北川家的唐慶崇拜.
大家知道崇拜,當然主日崇拜我們常常都有參與.
但是唐慶是一個慶祝.
那你說慶祝什麼呢?.
慶祝這間教會在那個地方73年.

$^{81}$沒錯,是慶祝這間教會.
但是教會是什麼呢?.
教會就是神的子民,就是我們.
我們就是教會.
所以今天如果說是唐慶.
你可以說,今天也是我們的生日慶祝.
所以好不好,和你旁邊的弟兄姊妹說一聲.
生日快樂,好不好?.
各位北川家弟兄姊妹,生日快樂.
願神賜福我們教會.
我們大家一起祈禱,我們同心同意.
在人以為最艱難的日子.
常常是你向人施恩的日子.
你讓人經歷生活的缺乏.
但是你自己癡下豐盛的人生.
癡下奇妙的救恩.
主我們今天在你面前能夠蒙受你的拯救.
主我們知道是因為你曾經感動人.
負上代價來向我們的前人傳福音.
而將這份恩典一直延伸到我們身上.
如今我們到你的面前,同樣將自己交托給主.
主我求你自己幫助我們.
知道過一個福音的人生.
知道委身在你自己奇妙的福音裡面.
我們知道你也去見證你,去傳揚你.
求主施恩,帶領我們今天讀你聖經說話的時候.
更深明白你福音的奧秘.
更知道委身福音.
我們仰望祈禱,奉主耶穌成名祈求,阿們.
我們今天會繼續思想聖經羅馬書的經文.
大家知道我們在今季一起思想到回到福音核心的.
在羅馬書裡面的訊息.
首先我們思想一件事情.
就是聖經一開始在羅馬書裡才讀的經文.
你就見到它很清楚地提出一件事情.
就是福音是神的大能.
福音是神的大能.
在聖經的第十四節,十五節,十六節剛才讀過.
之前經文,保羅就表達一件事情.
就是他很深切地渴望羅馬人能夠得救.

$^{121}$寫說:我恩切期望把福音也傳給你們在羅馬的人.
當時保羅在哥林多這個地方寫信給羅馬的教會.
他其實很清楚知道羅馬地區很艱難.
如果他去的時候,其實很多人會反對他.
甚至他會被抓捕.
但是保羅的內心裡面有一個很強的心願.
他很盼望有機會去到.
能夠將福音帶給這一群在當地未曾認識主的朋友.
他很恩切地期望.
或者你問為什麼這麼恩切呢?.
為什麼這麼期待呢?.
接著十六節,保羅就解釋.
因為我不以福音為恥.
無論別人怎麼看.
無論別人對基督徒的看法怎麼樣.
我知道福音是什麼.
我不會以福音為恥.
為什麼不以福音為恥?.
因為這福音是神的大能.
福音不單單是一個好消息.
不單單是神拯救人的道理.
直接地說,福音就是神的大能.
神的福音能夠完全地,充足地拯救我們歸向上帝.
能夠拯救我們出死入生.
是有大能的神的道.
當然,福音不單單臨到一個人或者一個群體.
是臨到所有,每一個相信的人.
每一個相信的人都能夠經歷神給人一份生命改變的大能.
有一次,在一個很偏遠的地方.
有一群西方人士.
他們去探訪非洲裡面的一些部族.
一些比較落後的土族.
當他們到達一個鄉村的時候.
他們留意到有一個當地的土人.
就在一棵樹下坐在地上看書.
他們說,他們看書了.
這些未必有自己文字的民族.
他們開始看書了.
其中一個訪者就走過去.
想看看他們在看什麼書.

$^{161}$到達的時候,一看發現原來是聖經.
這位朋友就忍不住對著這個土族說.
他說,你看這些這麼落後的東西.
你知不知道聖經已經過時了?.
這位土族明白他的意思,然後就跟他說.
他說,是嗎?已經過時了嗎?.
你知不知道,如果不是因為這本聖經.
你現在應該在我的肚子裡面.
後來他們明白這個土族其實是什麼.
是食人族.
是一個很強悍,凶殘的民族.
過去曾經殺了很多很多的人.
但是今天他們全族迅速.
他們整個群體改變.
有人以為聖經的道理已經過時了.
這個時代這麼發展,怎會再讀這些書.
但是他們不知道,這福音本是神的大能.
讓生命經歷奇妙的改變.
能夠領受神給我們的完全的新生命.
當聖經提到,這福音是神的大能的時候.
接著進一步,聖經在第十七節繼續說.
其實神的義就是藉這福音顯明出來.
這個福音帶著能力.
同時福音是將神的義顯明出來.
這裡所說的神的義是包含了神的一切的美善.
祂的光明,祂的正義的屬性.
包括了這些一切神自己的美善的屬性.
這些一切都是藉著福音來顯明出來的.
但是人如果要藉著福音去領受神的義.
就必定是因為信.
是因為信靠上帝而領受神的義.
你見到聖經所說.
神的義就是藉著這福音顯明出來.
是以信為始,而信為終.
是藉著信而得到神給我們的義.
當聖經提到信的時候.
有些人會覺得很奇怪.
為什麼會由信始,由信為終.
究竟是什麼意思呢?.
法經裡面講述當我們領受神的道的時候.

$^{201}$是因為我們承認自己的無助.
我們沒有能力去救自己.
因此憑著相信神自己的工作.
接受祂的拯救而得到神的義.
當大家有機會去讀一讀聖經的時候.
在聖經跟著的羅馬書.
我也有跟大家列出來.
在第十章講及一件事情.
羅馬書第十章提到當時的猶太人.
他們很想憑著自己進行律法.
就可以成為義人.
保羅說他們這一群人.
第三節說.
他們不明白神的義.
又企圖建立自己的義.
不明白神的義是什麼意思呢?.
因此自己憑著自己的努力.
自己行善來建立自己的義.
就不信從神的義了.
不明白信從神要賜給人義的方式.
他們不知道原來基督就是律法的終結.
基督來到結束了舊約.
但是開展了新約.
在新約裡面人是憑著信心去領受救恩.
或者說基督其實就是律法的總結.
將整個律法得義的重要核心.
在他的生命裡面表達出來.
讓人可以得到義.
讓哪些人?.
聖經說:好讓所有信的人都得到神的義.
以信為開始.
是聖經清楚說的.
但是當在繼續跟隨主的路上.
你必定是堅定地信靠主.
不會在中間以為我都已經不同了.
我自己有自己的力量.
我可以憑著自己繼續闖一條光明的路.
不是.
由始到終都是以信靠的心.
以靠神的恩惠能力去領受神給你的完全的拯救.

$^{241}$《迦勒體書》第三章這段經文.
大家可能會留意過.
這段經文很特別說一件事情.
就是說保羅問當時的信徒.
為什麼你們不知道一件事情.
忘記一件事情.
就是你們早期領受了聖靈.
是因為遵行律法還是聽信福音呢?.
是什麼呢?.
是因為聽信福音.
不是因為他們憑自己遵守律法就得到聖靈的恩惠.
但是跟著他就說.
第三節就說.
你們敢無知嗎?.
你們既然以聖靈開始.
現在卻以肉身來完成嗎?.
你們竟然以為可以憑著自己去完成神在你身上的拯救?.
基督徒當我們信著之後.
有很多時候都會覺得.
我跟別人不同.
我已經有一個新的生命.
因此面對一些試探的時候.
就知道我自己可以忍得住.
我不會被他影響.
但是慢慢就發覺.
其實何時不是憑著信心依靠神的幫助.
他就何時重新跌倒.
唯有由此到終.
你都堅定地信靠神.
才會能夠領受聖靈.
繼續得著聖靈給你的幫助.
能夠領受神給你一個完全的拯救.
我記得有一件事情是很多年前.
我自己女兒還小的時候.
有一次我跟她去長洲.
大家知道長洲有一個地方叫張保仔洞.
有沒有去過?.
那件事情很深刻.
我記得當我去到那個洞口的時候.
就發覺原來去到那個洞.

$^{281}$在地面和洞底有一段距離.
我要自己很小心地爬下去.
站在那個洞底.
我女兒在上面.
我就問我怎樣下來.
我說不要緊.
爸爸抱你下來.
我就伸高手抱起她.
就把她放下在洞底.
然後我們就走進去那個洞.
洞裡面有一條比較窄的通道.
但是完全漆黑一片.
看不到裡面是怎樣的.
我大概可以摸一下就走進去.
我走了兩步的時候.
女兒就站在那裡不敢動.
然後我走回去.
她一句話都不出聲.
第一時間捉住我的手.
爸爸帶著我走.
我就很放心地拖著她向前走.
就這樣兩父女很順利地去到洞口.
但是過程裡面我突然想.
如果女兒下來洞底.
突然她就說.
我搞定了.
我下來平地了.
我自己會走.
你自己先走.
我跟著來.
她能不能走呢.
我相信當時她的狀況.
她是完全不能走的.
只會碰撞那些石頭.
她不能夠走過那條黑洞.
我想起這件事情的時候.
我就想起.
在很多人生跟隨主的路上.
各位弟兄姊妹.
記得你是憑著信領受救恩.

$^{321}$不要以為你可以將來憑著自己.
去過聖潔的生活.
若然不是神自己的恩惠和工作.
你很可能在中間會重新跌倒.
若然不是你繼續緊握著神的手.
讓祂拖帶著你.
憑著信心作為開始.
憑著信心作為整個過程的基礎.
直到終結.
否則你不會得著.
神要給你的完全的拯救.
各位弟兄姊妹.
記得福音本是神的大能.
這句話是我們最大的安慰.
你一定會經歷過.
在跟隨主的路上.
看到自己很多的軟弱.
甚至有時候跌倒之後.
心裡責怪自己.
甚至憤怒自己.
在這些時候.
你記得.
你放心再次將自己交回給神.
很清楚地告訴神.
雖然我軟弱.
但你剛強.
好像那個小女孩一樣.
伸出手來握著神的手.
你會經歷那種能力在你身上.
讓你能勝過生活很多的試探.
一直堅定地跟隨上帝.
得著祂要給你的完全的恩惠.
我們再看.
在聖經當我們再讀的時候.
聖經提到憤怒的問題.
是神的憤怒.
神的憤怒出現在這個世界上.
在聖經這一段經文18至25節這一段裡面.
聖經提到神的憤怒.
顯明在一些不乾不淨不義的人身上.

$^{361}$為什麼會顯明在不乾不義的人身上呢?.
是因為神已經顯明了祂自己.
祂在自然界裡面顯明了祂自己.
人不是不知道的.
他們是知道的,他們是明白的.
但是他們是拒絕去接受神.
因此神的憤怒顯明在這個世界上.
這一段經文.
我將18至20節拿出來.
大家可能會不會太小.
看不到.
或者你看中間.
中間在其中的第19節.
聖經提到這個自然界的情況.
其實神的事情人所能知道的.
在他們心裡都是明顯的.
因為神已經向他們顯明了.
自從創造世界以來.
神那些看不見的特質.
看不見的.
是屬神的那種特質.
就是祂的永恆的大能.
和祂的神聖.
是顯而易見的.
從這個世界的大自然.
和一切祂所創造的事物上.
是顯而易見的.
聖經說.
這從祂所造的萬物.
就可以領悟使人無可推諉.
我自己在一間中學讀書.
其中一位老師.
他是一個業餘的天文學家.
可能有些人會認識他 廖慶齊先生.
他後來成為了香港太空館的第一任館長.
是一個很著名的業餘.
研究天文的一位學者.
我後來有機會和他聊天.
當時已經是他壯年的時候.
他跟我說.

$^{401}$你知不知道我是信了耶穌.
我不知道 以前他做老師的時候.
他還沒有信主.
他跟我說.
我過去這麼多年一直研究天文.
常常用很多時間去觀察星空.
有很多不同的學習.
但經常有一種很驚訝的感受.
很多時候一個人在山頂觀星的時候.
他心裡有一種很大的驚詫.
感受到整個天際是很奇妙.
帶著智慧的創造.
那種設計是想不到的奇妙.
當他觀察到很多天文地理的時候.
發覺整個天際的秩序運行不息.
很穩定地穩定在天上的時候.
他就感覺這些秩序後面有一份愛.
有一份對人的慈愛和正直.
他說當我一直觀察的時候.
我深信不疑有一位創造主.
他說後來有一次.
他在自己的酒店裡.
一個人在那裡安靜.
他看到聖經.
一看到聖經的時候.
他清楚知道主耶穌就是那位生命的創造主.
然後他自己一個人.
在酒店裡面去自信耶穌.
當你真的安靜下來.
接觸這個大自然.
接觸你所看到的各種世上被創造的事物的時候.
你會知道你會領悟一位的上帝是何等的真神.
上帝讓人知道祂的存在.
但是聖經就說人雖然知道.
但是故意不去承認.
聖經十八節我們剛才看到的就是.
很特別的情況.
聖經說神的憤怒顯明在哪些人身上呢?.
就是那些以不義壓制真理的人.
是用不義來壓抑真理.

$^{441}$去排斥真理.
去否定真理.
因為聖經說.
神的憤怒顯明在他們身上.
大家都會留意到.
聖經很多時候都會提及神的憤怒.
我們有時候都覺得.
不要這樣說.
神怎麼會憤怒呢?.
我們才會憤怒.
因為我們想起人的憤怒.
很多時候是因為一種仇恨.
或者一種要傷害人的情緒.
所以才會發怒.
但當聖經說到上帝的憤怒的時候.
是用一種我們所能夠明白的語言.
去描述神的公義的一種屬性.
神有一種疑慮.
按照祂的公義去審判不義.
在審問的過程當中.
神以烈怒來描述祂的審判.
神會審判人的不義.
正如光明出現.
自然就會驅逐黑暗.
神的公義顯明的時候.
就會審判不義.
神的烈怒就會出現.
當然當你讀很多聖經歷史的時候.
你會發覺聖經提到神的憤怒出現的時候.
有些時候是立刻出現.
但有些時候是延後才出現.
但總會出現.
例如大家會記得.
聖經提到有一對夫婦.
阿拿尼亞薩菲拉士堂主所說.
當時他們欺騙聖靈.
欺哄聖靈.
這些是我們所買的錢財.
全部拿出來了.
因為他們欺哄聖靈.

$^{481}$當時他們怎樣呢?.
即時死在施道的腳前.
你會想起歷史裡面.
當羅亞的日子.
人類犯罪.
聖經就容許洪水.
將整個大地.
和其中一切人類消滅.
神很多時候會立刻讓祂的憤怒出現.
帶來審判.
但有些時候.
神好像是延後.
不會讓這些審判即時出現.
有些人以為可能是神都不理.
或者神已經睡了.
祂現在都沒有時間理.
我記得有一年我去印度的時候.
探訪當地的宣教士.
他帶我去當地的一些印度廟.
入廟之前.
我見到所有人都做一件事情.
就是在門口有一個鈴鐘.
啪響才走進去.
我就問宣教士.
為什麼他們打鐘才進去.
是不是不打鐘就不准進去呢?.
他說不是.
他們相信有些時候神明可能會睡了.
或者他們走開了.
所以響下鐘.
讓祂知道我來了.
請你留意.
人以為因為上帝有些時候會睡覺.
所以這個時候不理那麼多.
你犯什麼錯.
有你了.
但是他們不知道.
聖經所說.
有些時候延後審判.
是因為寬容.

$^{521}$我們下星期會讀到的.
羅莫書體二章裡面.
其中的第四第五節說到.
聖經說.
你知不知道神的恩慈.
寬容和忍耐.
為什麼會忍耐?.
為什麼會讓你繼續犯惡.
而沒有任何的虧損?.
仍然好像風生水起.
我們廣東人所說.
聖經說.
是要令你悔改.
但是由於你的心環境不肯悔改.
你就在自己在神發怒和彰顯公義審判的日子.
有公義審判的日子.
但是現在未曾讓審判臨到你.
當你繼續在罪當中的時候.
你是在積累神對你的憤怒.
那些憤怒不會消失.
而只是積累到那審判的日子.
因此聖經當找到審判的時候.
當找到神的怒氣的時候.
來到羅莫書剛才讀的第一章.
這裡發覺聖經提到.
神的怒氣顯明在這些不乾不易的人身上.
但是有些時候我們發覺.
神的怒氣的顯明不一定是即時.
也沒有延後.
但是竟然透過任憑來顯出來.
什麼是任憑呢?.
我們再看回聖經.
聖經提到一件事情.
聖經提到在當時的人.
他們故意不認識神.
因此神就讓他們心思變得虛妄.
他們的心就昏暗.
他們的心靈已經去到一個.
屬靈的黑暗的狀況.
就在這樣的情況.

$^{561}$接著二十二節就說.
他們覺得自己有智慧.
所以就製造了不同的偶像.
包括了心中一些角色的偶像.
然後第二十四節.
所以,大家一起讀,好嗎?.
所以神任憑他們隨從自己心裡的私慾.
愛的是自己的身體.
所以說,神任憑他們隨從自己心裡的私慾.
行污穢的事情.
任他們做.
這樣就行了.
他不用那麼麻煩.
我自己搞定自己.
任憑,很多人以為是不理會的意思.
其實聖經任憑的字是包含了一種交出來的意思.
不是被動的.
是主動地將人交了給人自己的情慾.
意味著情慾是可以控制著人.
人自己沒有力量去避開這種情慾的試探.
意味著神不再寬容.
所以你可以說,這已經是審判的開始.
一個家庭的兒子染上了賭癮.
終日去賭博.
爸爸嘗試幫他.
有些地方嘗試幫他解決一些債項.
一直勸導他.
而且也找一些人.
盡量幫他有些中心去戒賭的中心.
來盡量幫他.
這個兒子嘗試去改變.
但後來又再次再次重新賭博.
甚至有一次跟爸爸說.
爸爸你不用理我了.
我自己搞定.
在那個時候,爸爸就說.
好,你既然不肯放棄賭癮.
我就任你自己做.
各位聽眾,在這些生活上的事情就會明白.
一個人如果任憑你做你想做的事情.

$^{601}$其實是你最大的災難.
你再不得到幫助.
得不到別人給你的輔助.
而面對著那種罪惡對你的捆綁.
聖經給我們一個很清楚的方向.
各位靈智妹.
在人類所面對各樣的試探當中.
如果你領受很多的提醒.
你要去回應.
但如果你領受提醒之後你不去回應.
你發覺會去到一個狀況.
就是神任憑你.
將你交給你自己的私慾.
這時候就很悲哀.
聖經繼續說第二種的任憑.
我們再看.
在根據聖經其中的25到27節就說到.
神任憑人隨從他可恥的私慾.
可恥的情慾.
來做他想做的事情.
26到27節這段經文.
繼續說一些關於同性行為的事情.
聖經稱這些行為為可恥的情慾.
是很羞辱的.
很可恥的事情.
26節裡有人說.
這裡說到自然的性關係.
有些人就說.
我覺得我和同性的人的關係才是最自然的.
但是小心再讀27節.
聖經所說的是什麼是自然的關係.
男性也放棄了與女性自然的性關係.
是說神所創造的婚姻裡一男一女的性關係.
那是被稱為自然的.
出於神的創造設計裡的關係.
但是如果男和男.
有些人認為這是自然的.
但聖經說到男和男去行.
是可恥的事情.
而提到在他們身上受到這種怪貿行為所得到.

$^{641}$所產生對他們身體的傷害.
或者說報應.
有人曾經說.
其實同性戀問題是天生的.
很多研究證明了.
是天生的.
所以是自然的關係.
它自然就是這樣的.
為什麼你不接受呢?.
當我們細心去看這些研究的時候.
你會發覺絕大多數這方面的研究.
其實都沒有一個確定性的結論.
很多時候是因為他們所做研究的.
這些樣本的個案有偏頗.
因此所達到的結論是有疑問的.
或者是統計上不足的數目的個案.
因此研究也不足以成為固定的結論.
當你細心去看的時候.
你不要隨便相信研究證明.
其實研究沒有證明.
聖經反而在過程當中不斷提醒.
這種情況會帶來一種危險.
就是人隨從這種私慾來生活的時候.
產生很大的對人生命的傷害.
當然當我們這樣說的時候.
我們不是要去歧視.
或者輕視一些同性戀行為的朋友.
只是按照聖經所說的.
誠實地講述神自己創造的設計.
神自己對人生命的設計是怎樣.
那個真理是怎樣.
讓人知道如何朝向真理迴轉.
歸向上帝的一個重要方向.
因此我們看到聖經裡提到.
人有些時候陷在罪當中.
因此神就任憑他們隨著私慾行事.
去建立不同的偶像.
包括自己心中的偶像.
第二 神就有些時候任憑人.
隨著他自己的可恥的情慾.

$^{681}$來將自己的運用成一個危險.
在一個不自然的情況出現.
產生對自己身體的傷害.
有一位聖經學者曾經這樣說.
拜偶像和同性戀是兩大例證.
是證明人嚴重違反.
神藉自然界和人的天性所彰顯的真理.
神透過自然界讓人知道有神.
但人卻創造心中的偶像.
神藉著人生命身體的創造.
顯明出人的天性如何.
但人違背神的設計.
因此去破壞壓抑剛才所說的真理.
陷在罪當中.
第三 我們簡單看一看.
在接下來的28,29節.
很多人不肖形承認神.
不肖的意思是.
人覺得不值得去承認.
沒有價值.
故意不去承認有神.
所以神就將他們交給.
任憑他們全敗壞的心.
英文所說的depraved mind.
將他們交給敗壞的心思.
在聖經裡面.
不肖和敗壞原來是同字源.
他們文字的根源是一樣.
亦有一個含義表達出.
不肖承認神帶來敗壞.
神將他們交給敗壞的心思.
所以你小心讀經文.
書中不斷列出很多不同的敗壞行為.
很多惡事一路列出來.
正如我們中國人信眾所說.
如果我們心中無神.
接下來會怎樣?.
大家知道.
目中無人.
如果你心裡不承認有神.

$^{721}$你覺得我甚麼都可以.
我為甚麼不可以這樣做?.
他就是目中無人.
各種惡行都會做出來.
不肖承認神.
生命就越發敗壞.
有眾多的惡事.
正如下面的經文可以看到.
很多很多不同的惡事就會走出來.
而跌在罪的裡面.
甚至最後30節所說.
不單自己喜歡做.
明明知道這件事是惡事.
但都想別人跟你一起做.
或者一起死.
而惡到一個地步.
是惡上加惡.
請知道我們讀完這段經文.
我想大家都會很清楚體會這件事.
人是沒有能力自救的.
人被罪完全捆綁著.
唯有當人願意謙卑.
信靠上帝.
承認自己的不足.
承認我只能靠著神.
我才能夠遠離邪惡.
正如聖經所說.
我們得救是在乎歸回.
在乎信靠.
歸回向這位上帝.
承認祂的公義的審判.
以致當你信靠祂的時候.
你可以遠離神公義的陣路.
我們各位已經信了主的弟兄姊妹.
或者這段經文我們更加明白.
這個世界的情況.
但是何時你察覺.
你裡面的私慾.
剛才所說的那些可恥的情慾.
或者那些敗壞的心思.

$^{761}$若然你察覺在你裡面出現.
你要很小心.
在神面前悔改悔轉.
免得你進入了神的憤怒當中.
憑著得救的信心.
以信為始以信為終.
一生堅定的信靠神.
跟隨拯救你的主.
我們一起祈禱.
滿有拯救的恩惠.
我們不配,我們無能為力.
但是主啊,你自己賜下你奇妙的救恩.
叫我們恩著信.
接納你自己的拯救.
能夠得勝罪惡,得勝死惡.
主啊,幫助我們人生的每一步.
我們由始至終堅定的信靠你.
行在你自己的心意當中.
以致我們能夠歡歡喜喜.
有一天,回到你的庚前.
領受你自己的稱讚.
主啊,願你這樣帶領,幫助我們.
奉靠耶穌基督,盛名祈禱.
阿們.
\newpage



\section{以弗所書 2:11-22}
\label{sec:032144}
\textbf{近在咫尺,遠在天涯}~~~~禤凱榮牧師
\newline
\newline
連結: \href{http://www.church.com.hk/acms/content.asp?site=cdc\&op=show\&type=product\&code=032144\&layout=sermon}{\texttt{032144}} ~~~~ 語音日期: 2026-05-17
\newline
\newline
\hyperref[sec:032143]{\small{< < < PREV SERMON < < <}}
~
\hyperref[sec:index_chronic]{\small{[返順時目]}}
~
\hyperref[sec:032145]{\small{> > > NEXT SERMON > > >}}
\newline
\newline
以弗所書 2:11-22
\newline
\begin{longtable}{cl}
\hline
\hline
章節 & 經文 (和合本修訂版)\\
\hline
2:11 & \begin{tabularx}{0.7\textwidth}{X} 所以，你們要記得：從前你們按肉體是外邦人，是「沒受割禮的」；這名字是那些憑人手在肉身上「受割禮的人」所取的。 \end{tabularx} \\ \\ \relax
2:12 & \begin{tabularx}{0.7\textwidth}{X} 要記得那時候，你們與基督無關，與以色列選民團體隔絕，在所應許的約上是局外人，而且在世上沒有指望，沒有神。 \end{tabularx} \\ \\ \relax
2:13 & \begin{tabularx}{0.7\textwidth}{X} 從前你們是遠離神的人，如今卻在基督耶穌裡，靠著他的血，已經得以親近了。 \end{tabularx} \\ \\ \relax
2:14 & \begin{tabularx}{0.7\textwidth}{X} 因為他自己是我們的和平，使雙方合而為一，拆毀了中間隔絕的牆，而且以自己的身體終止了冤仇， \end{tabularx} \\ \\ \relax
2:15 & \begin{tabularx}{0.7\textwidth}{X} 廢掉那記在律法上的規條，為要使兩方藉著自己造成一個新人，促成了和平； \end{tabularx} \\ \\ \relax
2:16 & \begin{tabularx}{0.7\textwidth}{X} 既在十字架上消滅了冤仇，就藉這十字架使雙方歸為一體，與神和好， \end{tabularx} \\ \\ \relax
2:17 & \begin{tabularx}{0.7\textwidth}{X} 並且來傳和平的福音給你們遠處的人，也傳和平給那些近處的人， \end{tabularx} \\ \\ \relax
2:18 & \begin{tabularx}{0.7\textwidth}{X} 因為我們雙方藉著他，在同一位聖靈裡得以進到父面前。 \end{tabularx} \\ \\ \relax
2:19 & \begin{tabularx}{0.7\textwidth}{X} 這樣，你們不再是外人或客旅，是與聖徒同國，是神家裡的人了， \end{tabularx} \\ \\ \relax
2:20 & \begin{tabularx}{0.7\textwidth}{X} 被建造在使徒和先知的根基上，而基督耶穌自己為房角石， \end{tabularx} \\ \\ \relax
2:21 & \begin{tabularx}{0.7\textwidth}{X} 靠著他整座房子連接得緊湊，漸漸成為在主裡的聖殿。 \end{tabularx} \\ \\ \relax
2:22 & \begin{tabularx}{0.7\textwidth}{X} 你們也靠他同被建造，成為神藉著聖靈居住的所在。 \end{tabularx} \\ \\
[1ex]
\hline
\hline
\end{longtable}
$^{1}$不知道大家知不知道.
雪梨有一些餐廳.
在入座之前首先要鎖起手機.
大家有沒有去過這樣的餐館.
這麼可怕.
你說偶爾去一下.
每次去到餐館都要先鎖起手機.
搞不定搞不定.
大部分人都不會去.
有些人甚至說.
進去餐館.
手機沒了好像沒穿衣服的感覺.
不知道大家是不是這麼嚴重.
去到一個什麼地步.
中國人就說.
同桌吃飯就各自修行.
不過我相信現今的版本應該是說.
同桌吃飯就各自低頭.
爸爸就玩手機.
媽媽就玩iPad.
子女就玩Game Boy.
你可以說.
一家人坐在同一張餐桌上進餐的時候.
近在咫尺.
不過他們的心境卻是遠在天涯.
各自有各自忙的世界.
今天的科技.
人們都說.
可以讓人有更好的生活質素.
不過相反來說.
它似乎讓人不自覺地.
就好像建了一道牆一樣.
將自己困在裡面.
讓我們彼此的溝通慢慢被侵蝕.
不過今天我們所讀的經文.
保羅告訴我們.
人生有另外一種牆.
比這些人與人之間隔膜的牆.
來得更加破壞性.
保羅說這個牆就是將人與神.

$^{41}$或者是人與人之間相隔的.
經文也告訴我們.
耶穌基督已經將這道牆拆毀了.
我們要問的問題就是.
今天我們在地上的生活.
是不是活出一個沒有牆的人生呢?.
讓我一同去看一看今天的經文.
我想問一問大家.
你相不相信有人說.
每一個人都是平等的.
人人是平等的.
你相信這句話嗎?.
如果你相信的話.
我相信很多時候你會活在痛苦當中.
因為現實上.
你很難找到一個地方是有絕對的平等的.
不平等的事情.
你可以說比比皆是.
發生在我們的種族之間.
國與國之間.
財富,社會地位,教育.
你的樣貌生不如何.
你的體型如何.
這些都讓人感受到不平等.
而這些不平等.
就好像一個看不到的牆一樣.
將人分成不同的等級.
在保羅的時代.
有一種對他們來說很困擾的隔膜.
就是猶太人和外邦人.
似乎中間有一個深不可測的鴻溝.
猶太人就叫外邦人是.
未有受過國禮的.
當然這是一個貶詞.
而外邦人就在經常諷刺猶太人.
你們有這種飲食法真的麻煩得不得了.
所以.
也由於這樣的鴻溝.
當時在耶路撒冷的聖殿裡面.
確實是有一面牆.

$^{81}$將外邦人和猶太人分開的.
考古學家發現.
這面牆上面甚至刻了一些文字.
文字是這麼說的.
外邦人不可以進入聖殿圍牆的範圍.
圍者必救.
追究的是什麼呢?.
被處死在殿裡面.
雖然你可以說是一牆之隔.
但是牆的裡面.
是代表著神的同在.
而牆的外面.
就是代表神不與他們同在.
外邦人是被拒於牆外.
當時的以弗所教會.
是保羅寫信給他們的.
而以弗所教會.
大部分都是外邦人.
可能甚至全部都是外邦人.
所以保羅寫給他們的時候.
就說你們.
就是說當時的以弗所的信徒.
你們以前是與耶穌基督無關.
在以色列的國民以外.
在所應許的諸葛上面.
你們都是局外人.
你們沒有份的.
而且在這個世界上.
沒有指望.
沒有神.
這番話很嚴厲.
不僅是在說一些.
社會階層上的冷落或者諷刺.
聽保羅所說簡直是判死刑.
我想大家想像一下.
如果你是一個難民.
大家知道什麼是難民嗎?.
難民到了一個很豐盛很富庶的邊界.
到了邊境.
這個難民.

$^{121}$沒有護照.
沒有通行證.
沒有擔保人.
又不懂得那裡的語言.
在他背後.
是他破碎的家鄉.
在他前面.
很高的牆.
前面有一道閘門.
門上面寫著.
如非國民.
不得內進.
很慘.
而保羅形容的外邦人的光景.
其實就是這樣.
一幢牆的左隔.
兩邊是完全不同的世界.
一邊就是代表著神的豐盛.
神的同在.
而另一邊就是沒有應許.
沒有希望.
沒有神的一個地方.
很大的分別.
你可以說是光明與黑暗的分別.
是罪惡與良善.
是生命與死亡的分別.
而我相信更加的危險.
不是在於牆所代表的距離.
大家想一想.
我相信真正的危險在於.
人究竟知不知道有這幢牆的存在.
人究竟知不知道如果有這幢牆的存在.
他究竟是牆的哪一邊呢?.
大家都記得.
不是很久之前.
在香港發生的火災.
大家知道死了一百多人.
很多人罹難.
大家知不知道這麼多人死的原因是什麼?.
他們發現其中一個最重要的原因.

$^{161}$就是那些人不知道.
原來當事發的時候.
大廈的警報系統被那些人關掉了.
關掉了警報系統.
以致很多困在大廈裡面的時候.
都不知道有火警發生.
到他們知道的時候.
太遲了.
已經身陷火海.
沒有辦法走得出來.
所以朋友們.
這個牆是代表著.
剛才我們所說聖殿裡面那幢牆是.
分隔了外邦人.
以致外邦人.
和神好像是分開了一樣.
連這個牆就是代表著神與人之間的阻隔.
聖兵告訴我們這幢牆是罪.
罪將人和神分開了.
這是一個我們不能改變的事實.
但是重要的問題我們要問的.
我們究竟知不知道.
我們在這幢牆的哪一邊呢?.
如果我們不知道這幢牆.
或者我們不知道我們在這幢牆的哪一邊的時候.
我們需要像這個火災的悲劇.
告訴我們一件很重要的事情.
如果在火災的時候.
我們不知道有這個火災的時候.
這是致命的.
我們究竟是不是有火災的危險呢?.
今天我們是不是在火災的危險裡面.
在屬靈的意義上.
是的.
我們是在火災的危險裡面.
我們看一段經文.
耶穌基督曾經這樣說過.
信子的人有永生.
不信的人得不隨永生.
不信的人怎樣呢?.

$^{201}$神的陣路常在他身上.
這是什麼陣路呢?.
耶穌在馬太福音裡面這樣說.
他說將來那些不信耶穌的人.
會進入那位魔鬼和他的使者所預備的永火裡面去.
原來有一個火在等待將來最終極的審判.
如果我們真的不幸地.
在這個牆的不正確那邊.
我們怎樣才可以由不正確那邊去正確那邊呢?.
以弗撒河書剛才我們所讀的經文第十三節告訴我們.
你們.
那時候的信徒們.
以前你是遠離神的.
如今怎樣呢?.
在基督耶穌裡面.
靠著他的血.
已經得親近了.
親近什麼?.
親近永生的上帝.
保羅在這裡告訴以弗撒河的信徒.
你以前是跟神很遠的.
遠在天涯.
不過現在靠著耶穌的寶血.
你已經近在咫尺.
跟神很親近.
這就告訴我們.
如果我們今天在牆的錯誤那邊的時候.
又想去到另外一邊的時候.
我們唯一的途徑.
就是通過耶穌基督的寶血.
保羅在第十四節告訴我們.
耶穌基督能夠將罪人從遠處拉到他跟神很親近的原因.
是什麼?.
是因為.
耶穌基督已經拆毀了中間隔斷的牆.
讓人與神可以得著和諧.
那你說.
耶穌是怎樣將這個牆拆毀的呢?.
在第十五節告訴我們.
耶穌是以自己的身體廢掉冤仇.

$^{241}$就是律法上的規條.
如果單看這兩句話.
我們很容易有一個錯誤的感覺.
就是說什麼律法這麼邪惡.
竟然讓人有冤仇.
如果這樣當然是錯誤.
我們的意思並不是說.
律法是邪惡.
或者是律法是鼓動冤仇.
並不是這個意思.
保羅這裡所說的.
就是說在律法的禮儀和盟約的規條上.
將猶太人和外邦人分開了.
為什麼會這樣呢?.
神原初所定律法的時候.
是為了要什麼呢?.
是為了要將以色列民.
從萬民當中分別出來.
成為聖潔.
以致他們可以成為神的子民.
有什麼條例呢?.
例如國禮.
飲食的條例.
憲制.
節旗.
聖殿裡面的規例等等.
將這些眾多的律法的規條.
讓以色列人從萬民當中被分別出來.
不過很可惜.
由於人的軟弱.
以色列人得到這些規條之後.
反而將這些規條讓自己高高在上.
感覺超人一等.
成為看不起其他民族的原因.
因此引致很多的冤仇出來.
現在保羅告訴我們.
由於耶穌基督在十字架上的捨身犧牲.
已經將這些冤仇完全廢掉.
為什麼可以廢掉呢?.
我們看幾次經文.

$^{281}$大家記得耶穌在十字架上斷氣的時候.
祂最後一句話說什麼呢?.
「成了,然後低下頭.
將靈魂交付給父神」.
在《馬可福音》也告訴我們.
在耶穌斷氣之後.
馬上電裡的萬子從上到下.
列為兩半.
為什麼從上到下呢?.
很明顯.
是從天上撕開萬子.
是神的一個作為.
而《出埃及記》告訴我們.
這個被從上到下撕成兩半的萬子.
是什麼萬子呢?.
就是將聖所和致聖所分開的一個萬子.
什麼是致聖所呢?.
致聖所裡面裝著藥櫃.
代表著神的同在.
致聖所這個位置不是人人都可以去到的地方.
是每年一次贖罪日的時候.
大祭司才可以進去為百姓贖罪的地方.
而這個萬子很明顯就代表著.
聖潔的神和罪惡的子民.
中間的隔絕.
當這個隔絕的萬子被撕開的時候.
神就用一個很形象的畫面告訴我們每一個人.
現在每一個人都可以去到神的同在那裡.
而希伯來書的作者更加的很清晰.
很直接的告訴我們.
他說弟兄們.
你們既因耶穌的血.
得以坦然進入致聖所.
是藉著耶穌為我們開了一條又新又活的路.
從萬子經過.
這個萬子就是他的身體.
耶穌的身體在十字架上被撕裂.
就代表著這個分開聖所與致聖所之間的萬子被撕裂.
代表著人現在不會因為萬子的阻隔.
而可以直接去到致聖所.

$^{321}$直接去到神的同在.
而當耶穌基督說「成了」.
「成了」什麼?.
「成了」就是宣稱.
耶穌基督在地上為拯救世人的救贖工作.
已經得到完成.
包括了.
記不記得我說過.
罪的公加乃是死.
人的罪只能夠用血來洗淨.
耶穌基督就說這個律法的要求.
就是在他死在十字架上流出寶血的同時.
滿足了這些律法上的要求.
換言之.
剛才我們所說的律法上的禮儀和盟約的規條.
就在那一刻.
被完全耶穌基督的救恩完成.
而得到滿足.
所以.
今天的基督徒.
不再需要像以前那樣來獻祭.
或者受割禮.
或者要嚴格的遵守食物的規條.
聖殿的規條等等.
因為耶穌基督已經一次過在十字架上.
為所有的.
律法上對人的要求.
完全的滿足了.
所以耶穌曾經這樣說過.
祂說.
你們不要以為我來到這個世界是要廢掉律法.
相反我來到這個世界是要承傳律法.
所以很明顯.
耶穌基督所廢掉的.
千萬要弄得很清楚.
不是要廢掉律法.
而是廢掉.
因為律法而引起的人與人之間的冤仇.
也是挪開了因為律法而顯明讓我們知道的罪.
不單止耶穌沒有廢掉這些律法和規條.

$^{361}$相反.
祂更加將這些律法承傳了.
所以.
耶穌就說.
不要以為我是來廢掉律法.
你是承傳.
好了.
耶穌基督承傳了這個律法之後.
複清了罪的功架之後.
也為人開啟了一條路進入神的同在之後.
我們就需要明白.
我們在耶穌基督裡面.
是需要帶著信心.
讓神與人之間的牆可以得到拆毀.
留意.
神將和諧帶入猶太人與外邦人.
並不是將猶太人變成外邦人.
也不是將外邦人變成猶太人.
讓他們可以不要這麼多爭拗.
可以和平共處.
並不是這樣.
神如何將這兩個不同或很嚴拘的群體帶在一起呢?.
聖經這裡說.
「祭」是造成一個新人.
造成一個新人.
什麼意思?.
什麼新人?.
教會.
這個新人就是今天的教會.
教會從來不是建築在一個民族.
或者國家.
或者社會階層.
或者某一個文化上面.
所有不同地方不同時代的教會.
都是建築在耶穌基督的身上.
所以弟兄姊妹.
當這個牆.
或者人與人之間的仇恨.
被耶穌基督所除去的時候.
我們是需要活出在教會裡面.

$^{401}$基督所代表的合一.
就是說不管我們的膚色.
不管我們在社會裡面是否高階層還是低階層.
不管我們的教育程度.
所有的東西都在基督耶穌裡面.
在教會裡面不再重要.
如果我們真的有愛和合一的時候.
我告訴大家.
教會很難不吸引外面的人進來教會.
享受教會在以外的世界體驗不到的愛.
也因為這樣教會才能夠榮耀神的名.
咦 今天教會很難感受到這件事.
很多時候我們還是看到很多冤仇.
如果有的時候.
求神憐憫我們.
我們的責任是要活出神的合一.
不過聖經也告訴我們.
這些冤仇的事情必定會發生.
所以求神讓我們不要在這些冤仇的製造者.
而是在這些冤仇的調解者.
讓我們活出神的恩典.
也許因為.
也許因為人的軟弱.
你會發覺.
很多時候我們真的沒辦法活出.
教會裡面應該有的完美.
應該有的合一.
不過.
無論教會如何的不濟.
在神的心意裡面.
教會始終都是.
屬天的子民.
所歸屬的地方.
如果我們看第十九節.
保羅告訴當時的外邦人.
你們以前是外人.
是客女.
是沒有身份的.
不過現在你們已經成為神家裡的一份子.
什麼意思呢.

$^{441}$我想大家都去過旅遊.
當你做旅客的時候.
很多時候我們去旅遊就是觀光.
很享受那個地方的風景.
吃他們當地的美食.
或者是享受.
嚮往當地的文化等等.
不過無論如何我們怎樣的喜歡.
到我們行程完的時候.
我們都要回到老家.
回到我們以前所屬的地方.
而在神的恩典裡面.
通過耶穌基督的信心.
這裡告訴我們.
我們已經不再是神國裡面的遊客.
也不是神家裡面的租客.
我們是什麼.
我們是神家裡面的一份子.
我們是天國的子民.
我們是神的兒女.
旅客如果沒有身份的時候.
這裡告訴我們.
是可以成為神家裡面的一份子.
罪人可以成為神的兒女.
而這個新的身份我們大家要留意.
不是靜態的.
什麼意思呢.
不是說我們拿了這個身份之後.
哇 怎樣當.
我們翹起雙手等上天堂.
正.
不是這樣的.
聖經告訴我們.
我們拿了這個身份之後.
不是繼續我們平常的生活依舊.
等上天堂.
不是這麼簡單.
保羅在這裡告訴我們.
他用了一個比喻.
我們好像一塊石頭.

$^{481}$被神所採用.
搜集來建造神的聖殿.
而這個建造的過程.
是一個正在進行.
和會持續進行的過程.
今天神仍然在建造這個教會.
祂的聖殿.
或者耶穌基督的身體.
每一個真心相信耶穌基督的人.
其實就是一個建築的材料.
神是使用來建造祂屬靈的聖殿.
當然這是一個比喻告訴我們.
神並不住在人手所造的建築裡面.
也不是處在某一個地理的位置上.
這個比喻告訴我們.
神是住在什麼.
是住在真實信徒當中.
這一段經文也要大家留意.
祂不是說種族的和解.
而是說宇宙性.
整體創造性的一個更新工程.
神不是單是將猶太人和外邦人.
令他們可以和諧地相處.
不是這樣的.
經文告訴我們.
神在當中是在造一個新的群體.
這個新的群體.
剛才我們所說是教會.
是屬靈的聖殿.
這個也是神要居住在當中的地方.
留意這裡建築物的根基.
聖經所說.
這個聖殿是建造在使徒和先知的根基上.
也要留意這個所謂的防角石就是基督.
大家知道防角石是古代建築時作為標準的基石.
這個基石在某一個角度不準確的時候.
整個建築就會斜了.
好像意大利的比薩斜塔.
大家有沒有去過.
我去過,真是斜的.

$^{521}$你讓我上去,我都不敢上去.
因為如果那個防角石是斜了的時候.
整個建築就會斜了.
我想大家想像一下大教堂.
一個大教堂.
大家知道大教堂不是從一個大塊石造出來的.
大教堂是從很多數以萬計的小石.
砌合而成的.
每一塊小石都有它們各自的形狀,顏色.
或者是質地.
如果在建築裡面.
你拿走一塊石,兩塊石的時候.
你會發覺久而久之.
這個大教堂裡面的結構就會受到影響.
在整個建造裡面.
每一個真實信耶穌的人.
都是一塊容易作為建造的材料.
也許你可能是一塊粗糙的小石.
不過你會發覺旁邊的那塊石不是.
是打磨得很漂亮的.
為什麼一塊粗糙的石.
打磨得很漂亮的石.
可以拼合在一起成為神的聖殿呢?.
聖經就告訴我們.
全因耶穌基督的愛和聖靈的能力.
讓我們本來不一樣的信徒.
可以連合在一起來建造神的聖殿.
不過大家要留意.
當不同的信徒被連在一起的時候.
你會說:好吧.
既然神這樣編我在這個教會裡面連在一起的時候.
你旁邊那個我很難頂你.
不過我忍.
為了神的工作而已.
對不起.
當不同背景性格的信徒.
被連在一起的時候.
神不是要我們容忍.
神是要我們連在一起.
以致我們可以彼此學習.

$^{561}$共同成長.
你會發覺新約裡面.
很多的教導.
都有彼此互相的教導.
彼此認罪.
互相代求.
彼此學習.
彼此包容.
全部都是說人與人之間如何建立一個合一的群體.
所以今天大家要搞清楚.
教會從來不是一個建築物.
從來不是一個活動.
一個傳統.
而是神的家.
神的殿.
神居住的地方.
當保羅強調耶穌基督就是整個教會的防國石的時候.
給我們要認清一件事.
教會是需要在任何事情上.
我們所傳講的信息.
我們所做過的決定.
我們所籌辦的事工.
完全都是要以耶穌基督作為中心.
這個就是以耶穌基督作為防國石的意思.
而當保羅強調教會是神居住的地方的時候.
即是告訴我們甚麼.
即是告訴我們今天我們在教會裡面所做的一切.
需要活出彰顯到神的性情.
甚麼性情呢.
神的愛.
神的聖潔.
神的合一.
神的真理.
我們都要在教會裡面.
將這些所有屬神的品質.
完全地活出來.
所以基督徒應該.
甚至是必須要有教會的生活.
必須與其他信徒來相交.
不是啊牧師.

$^{601}$那個我真的受不了.
不可以和他一起合作做事的.
不是容忍.
我們是互相學習.
不是呀我沒有錯.
全部都是他錯.
對不起.
還是那一句.
互相學習.
彼此包容.
這是神給我們一個很重要.
在教會裡面我們需要知道的東西.
我們每一個人.
我們要知道.
都是教會.
當神要做這個教會的時候.
我們每一個真實的信徒.
都是一個不能或缺的一部份.
而當神將這些所有不同的個體.
聯合在一起的時候.
我們要依靠耶穌基督的愛.
祂的能力.
祂的帶領.
祂的保守.
讓每一個人.
在整個過程當中.
我們首先被建造.
首先被無造.
以致我們可以.
能夠成為一個.
彰顯神榮耀的群體.
成為一個堅固的聖殿.
讓神住在當中.
得逐榮耀.
請我們從低頭祈禱.
正如今天的講題一樣.
最後如果我們當中雖然.
近在咫尺的周圍的弟兄姊妹.
但是我們的心智.
我們的所有的一切.

$^{641}$都是遠在天涯的時候.
神求你首先原諒我們.
讓我們學習從你而來的愛.
去成為一個合一的群體.
互相學習.
互相包容.
互相成長.
以致這個地方.
能夠讓你居住在其中.
也讓這個地方.
能夠成為一個榮耀你的地方.
求你帶領.
求你賜力.
祈禱.
奉主名求.
阿們.
\newpage



\section{啟示錄 2:12-17}
\label{sec:032145}
\textbf{別迦摩:勇敢卻走偏的教會}~~~~林穎豪傳道
\newline
\newline
連結: \href{http://www.church.com.hk/acms/content.asp?site=cdc\&op=show\&type=product\&code=032145\&layout=sermon}{\texttt{032145}} ~~~~ 語音日期: 2026-05-17
\newline
\newline
\hyperref[sec:032144]{\small{< < < PREV SERMON < < <}}
~
\hyperref[sec:index_chronic]{\small{[返順時目]}}
~
\hyperref[sec:032146]{\small{> > > NEXT SERMON > > >}}
\newline
\newline
啟示錄 2:12-17
\newline
\begin{longtable}{cl}
\hline
\hline
章節 & 經文 (和合本修訂版)\\
\hline
2:12 & \begin{tabularx}{0.7\textwidth}{X} 「你要寫信給別迦摩教會的使者，說：『那有兩刃利劍的這樣說： \end{tabularx} \\ \\ \relax
2:13 & \begin{tabularx}{0.7\textwidth}{X} 我知道你的居所，就是有撒但座位之處；當我忠心的見證人安提帕在你們中間，在撒但所住的地方被殺之時，你還堅守我的名，沒有否認對我的信仰。 \end{tabularx} \\ \\ \relax
2:14 & \begin{tabularx}{0.7\textwidth}{X} 然而，有幾件事我要責備你，就是在你那裡有人服從了巴蘭的教訓；這巴蘭曾教唆巴勒將絆腳石放在以色列人面前，使他們吃祭過偶像之物，並且犯淫亂。 \end{tabularx} \\ \\ \relax
2:15 & \begin{tabularx}{0.7\textwidth}{X} 同樣，你那裡也有人服從了尼哥拉派的教訓。 \end{tabularx} \\ \\ \relax
2:16 & \begin{tabularx}{0.7\textwidth}{X} 所以，你當悔改；若不悔改，我很快就到你那裡來，用我口中的劍攻擊他們。 \end{tabularx} \\ \\ \relax
2:17 & \begin{tabularx}{0.7\textwidth}{X} 凡有耳朵的都應當聽聖靈向眾教會所說的話。得勝的，我必將那隱藏的嗎哪賜給他，並賜他一塊白石，石上寫著新的名字，除了那領受的以外，沒有人認識。』」 \end{tabularx} \\ \\
[1ex]
\hline
\hline
\end{longtable}
$^{1}$在《蝙蝠俠》裡有一個我很喜歡的角色.
這個角色曾經是這個市裡面一個很正義的檢察官.
代表著體制裡面一個很正義,希望的代表.
不過這個檢察官經歷了一次嚴重的意外.
他不單止留下了一個嚴重的創傷.
也令到這個人毀容.
不知道有沒有人知道我在說《蝙蝠俠》裡面是誰呢.
這個就是哈維丹特.
這個人昔日是《蝙蝠俠》裡面的光明騎士.
但是那次令他毀容的意外和創傷.
他是令道德最堅定,成就最卓越的人徹底的墮落.
最後他化身成為否認一切理性和道德準則的一個雙面人.
我說我很喜歡這個角色不是因為我很欣賞他或者我很認同他.
而是因為我覺得他很真實地道出了人性.
將內心的惡和他原本生命裡面有的善.
完全割裂成為兩個獨立運作的極端.
所以我覺得《蝙蝠俠》真的很吸引人.
因為他其實不只是一套給小朋友看的卡通.
而是他同樣講述很多人的深層思想和矛盾.
當我覺得《蝙蝠俠》很吸引人,很好看的時候.
其實我們手上的聖經同樣吸引甚至更加有魅力.
是因為聖經不單止與我們陳明人性的真實和軟弱.
聖經也在黑暗裡刺下盼望和答案.
時隔三個星期,今天我們再次回來啟示錄系列的《正道》.
來到耶穌基督第三封寫給教會的信.
我們見到別加摩這間教會.
雖然這間教會是很勇敢,很堅定.
但可能我們剛才讀了經都見到.
他們最終都是墮落,妥協.
心願耶穌基督寫給這間雙面教會的這封信.
都成為我們今天教會每一個信徒.
一個警告,一個提醒.
我們再次低頭祈禱.
親愛的主,願你大能的說話.
今天否開我們的生命.
成為我們生命裡很重要的提醒,很重要的警告.
主啊!今天這個世界.
沒錯有很多的患難,有很多的苦難.
但同樣這個世界都有很多的誘惑,很多的試探.
願主你堅固我們.

$^{41}$讓我們知道在世的日子怎樣作一個蒙你喜悅的人.
蒙你喜悅的信徒.
最後幫助我們今天讓講的能夠講得清晰.
也讓聽的有聖靈的導引讓我們去明白.
我們將以後的時間恭敬教堂.
奉主名求.
別加摩這間教會是怎樣的呢?.
大家可能透過剛剛讀經的時候.
已經對別加摩這個地方可能有一些初步的印象.
覺得別加摩對於基督徒來說一定是很凶險的地方.
因為耶穌第一句就說.
這是有撒旦座位之處.
這是撒旦所住的地方.
別加摩是一個這樣的地方.
是羅馬帝國當時在小亞細亞的一個行政中心.
很早期的時候這個地方已經很蓬勃.
懂得鑄造一些錢幣,發行一些錢幣.
並且這個行政中心還有一個很重要的標誌.
一個規模很大的圖書館.
我不知道你覺得圖書館有多大.
有多少書才可以叫做一個大的圖書館.
這間圖書館有二十萬卷的藏書.
如果是我們的趙牧師一定流連忘返.
除了別加摩是一個行政的機關之外.
別加摩也充斥著各種異教的神殿.
當中有三座神殿是用來敬拜當時希臘的宙斯.
女神雅典娜和酒神.
我們看看別加摩的神殿究竟是怎樣.
有一座神殿很出名.
出名到如今如果你去德國旅行.
德國有一個博物館將別加摩這個神殿重新造出來.
讓大家看看這個敬拜宙斯的神殿.
如果大家去德國今天也可以看到.
大家覺得如果這座神殿坐落在別加摩的山上.
居高臨下一座這樣的聖殿.
看起來像不像一個座位.
然而耶穌基督說別加摩是撒旦所住的地方.
不單止是因為有這些神殿.
更加是因為那裡有一座神殿.
故意不是用來敬拜神 而是敬拜人 敬拜羅馬皇帝的神殿.

$^{81}$昔日羅馬帝國推行帝王崇拜.
將羅馬皇帝尊為神一樣.
在人間彰顯.
就好像中國也將皇帝尊為天子一樣.
別加摩這座城市就是帝王崇拜的重鎮.
你不在那裡敬拜皇帝做神.
不敬拜他不單止會沒工作.
甚至會沒了這條命.
所以耶穌口中的安提柏就是其中一位殉道的使徒.
歷史相傳他因為拒絕敬拜羅馬皇帝.
被放在一個銅造的牛.
燒紅後將人擺在裡面 活活地燒死.
不知當你想像一下.
如果你親眼見到你的屬靈長輩.
可能見到趙牧師 可能見到你的弟兄.
因為堅持信仰而慘死.
你心裡有什麼感受呢?.
不知你心裡會不會害怕呢?.
但是別加摩的教會竟然沒有退縮.
他們可能流著眼淚 但抵住恐懼.
耶穌說他們堅守他的命.
堅守這個字在希臘文裡有「捉住」的意思.
意思是無論多艱難 有人死了.
他們仍然緊緊地捉住耶穌.
仍然宣告耶穌基督是他們唯一的主.
他們至死都沒有棄絕耶穌的福音.
而耶穌說「我看到你們的勇敢」.
今天耶穌同樣可能見到我們的勇敢.
今天我們身處澳洲.
雖然我們不需要面對帝王崇拜.
但或許我們只是處於一個唯利是圖的職場環境.
當公司的上司要求你欺騙你的客戶.
又或者不是欺騙的.
是暗示你要和你的客戶隱瞞事實的全部.
你不是騙他 你只是沒有說完.
是用不合神心意的手法在商場競爭的時候.
可能你覺得這樣不對 不是神想我做的事.
所以你就冒著一個失去升遷機會.
被同事孤立 甚至失去工作的風險.
你仍然選擇拒絕去做.

$^{121}$在這個以金錢 以業績為王的現代神殿裡.
你堅持宣告耶穌才是你唯一的主.
耶穌今天跟我們說 我知道你身處的環境是怎樣.
我也見到你所承受的壓力.
見到你仍然很努力捉住我.
耶穌見到我們的眼淚 見到我們的堅持.
也深深珍惜 紀念我們這份勇敢.
不過就正當我們為別加摩教會的勇敢喝彩的時候.
耶穌卻話鋒一轉.
然而有幾件事我要責備你.
因為在你那裡有人服從了巴蘭的教訓.
也有人照樣服從了尼哥拉一黨人的教訓.
很諷刺 服從這個字其實在這段經文裡.
和上一節堅守的字在希臘文裡是同一個字.
就是說 別加摩教會的信徒.
原來他們不單單只是捉住了耶穌.
原來他們另一隻手還捉住了其他東西.
捉住了另外兩種教訓 兩種不合神心意的教訓.
為什麼會有這麼大的反差呢?.
為什麼他們的信仰不是很堅定呢?.
他們連死都不怕這間教會.
為什麼一間死都不怕的教會在信仰裡會走偏呢?.
讓我們先回去舊約 看看巴蘭是誰.
民數記二十二章至二十五章去到記載.
巴蘭是一個很貪財的術士.
當他發現自己不能正面去救助以色列人的時候.
他為了得到摩押王巴勒的報酬.
他就教了巴勒王一個惡毒的詭計.
派一些年青貌美的女子去誘惑以色列人.
邀請他們一起食祭偶像的筵席.
一起行刑.
原來是魔鬼撒旦發現.
用恐怖的壁壩 用刀劍原來打敗不了別加摩教會.
不可以令別加摩教會離開神.
於是魔鬼撒旦換了另一個策略.
不再用刀劍去威脅.
他用世俗的利益.
用人際關係的溫暖.
用肉體的享受.
去同化教會的信徒.

$^{161}$他在我們信徒的耳邊 在昔日別加摩教會.
信徒的耳邊去講.
大家都這樣做的.
這些只是生存法則.
因為在別加摩這個商業行會的聚餐.
其實他們多數都是在異教的神殿裡進行.
他們聚會的時候.
這些商會聚會的時候.
就是吃那些祭過偶像的肉.
飯後往往伴隨著廟祭的淫亂狂歡.
然而 別加摩教會是一間怎樣的教會呢?.
剛才我們說是一間勇敢的教會.
緊緊抓住耶穌的教會.
為什麼一間這樣的教會.
在信仰上會軟弱跌倒呢?.
在這個時候.
我們就知道原來這就是尼哥拉一黨人的教訓.
在信仰的層面乘虛而入.
就好像先前在以弗所教會的信裡的正道提及過.
尼哥拉一黨是當時教會面對一班很有影響力的異端.
他們說肉體和靈魂是可以分開的.
他們教導會眾.
我們的靈魂已經因為我們信耶穌基督得救了.
靈魂已經得救了.
所以我們肉體做些什麼其實不重要.
神的恩典是無限的.
並且神都會理解我們.
今天我們在別加摩這裡住.
這麼多人信主.
神明白我們身不由己.
你不去參加這些聚會的應酬.
你不去做這些事情.
你生意怎樣做?.
你怎樣養你的家人?.
我們去這些宴教.
其實我們不是去拜神.
我們去參加聚會.
我們屬於應酬性質.
我們吃這些祭偶像的肉.
算得了什麼呢?.

$^{201}$這些是我們逢場作戲.
我們不是從心裡面想吃的.
我們出席這些聚會.
不會影響我們心裡面對神的信仰.
我們心裡面有主就好.
我們心裡面愛主就行了.
外面做些什麼.
手碰些什麼.
口碰些什麼.
不重要.
簡單來說.
尼哥拉一黨的人主張的就是.
唯一一樣重要的就是.
你心裡面愛主就行了.
外在的生活方式.
可以完全和世界一樣.
所以別加摩教會的悲哀.
就在於.
其實他們可以為耶穌死.
但是他們不懂得為耶穌而活.
他們可以勇敢面對外在的獅子.
但是他們不懂得處理內在的毒蛇.
他們視而不見.
他們包容了錯誤的教訓.
讓真理和世俗混雜在一起.
並且他們用恩典作為包裝.
去合理化他們所做的錯誤的事情.
這種就是妥協.
而妥協從來不是一瞬間的背叛.
而是一步一步慢慢慢慢跌入陷阱.
好像溫水煮蛙一樣.
更可怕的是.
其實或許尼哥拉一黨人的教訓.
到今天彷彿都成為我們信仰上的絆腳石.
令我們將信仰從我們的生活去到割裂.
我們仍然覺得肉身所做的事不影響我們的信仰.
可能在職場上 在財務上.
我們為了爭取高一點的業績.
或者我們為了避稅.
我們選擇對於做假訴.

$^{241}$對於簽一些不實的文件.
我們隻眼開隻眼閉.
但我們心裡就跟自己說.
商場如戰場.
真是個個都是這樣生存的.
不要緊 我賺多點錢.
我回來教會奉獻多點錢.
無壞的.
或者我們在感情的關係上.
可能我們只是處於一段誤討神喜悅的關係.
甚至我們妥協於世俗的婚姻觀.
我們每次當心裡的聖靈提醒我們.
責備我們的時候.
我們卻是用恩典麻醉我們自己.
神是充滿愛的.
祂一定會體諒我的孤單.
體諒我身不由己.
更何況身體的界線是愛在野.
不影響我心裡對神的愛.
對神的信心.
或者在私密的私生活裡.
可能我們手機裡.
藏著一些不該看的網頁.
或者我們毫無顧忌地.
沉浸在充滿扭曲價值觀的娛樂裡.
但我們跟自己說.
我平時工作太大壓力.
我放鬆一下自己.
反正神有無限的恩典.
祂過去赦免我.
今天 未來 祂也會繼續赦免我.
我們是否把信仰.
當成了一種屬靈的積分制.
我們上教會.
我們在小組裡服事.
我們有靈修的時候.
就好像在儲分.
儲完的那些分就可以用來抵消.
甚至合理化我們星期日到星期六.
在世俗裡妥協的那種罪疚感.

$^{281}$我讀了經就可以了.
不要緊.
星期一至五.
繼續做那些事.
繼續見那些人.
魔鬼最可怕的武器.
原來不是碧白.
而是糖衣的毒藥.
魔鬼欺騙我們每一個人.
讓我們以為自己可以同時擁有世界.
又同時跟隨神.
魔鬼欺騙我們.
讓我們以為可以一手抓住耶穌.
另一手繼續抓住世間的響樂.
耶穌在福音書警告過我們.
一點點的面教.
就能夠令到整團的面發起.
走偏的信仰.
有時可能比不信更加可怕.
因為它引導我們.
在不知不覺之中.
甚至在我們自以為自己無罪.
自以為自己很正常.
每個人都是這樣.
慢慢卻步向滅亡.
面對著走偏的教會.
耶穌的反應是什麼呢?.
十六節.
耶穌說.
所以你當悔改.
如果不悔改.
我就快臨到你那裡.
用我口中的劍去攻擊.
耶穌很嚴厲的警告教會.
他們不可以一手抓住耶穌.
另一手又抓住罪惡.
口中的劍所說的.
所代表的.
就是代表著耶穌帶著能力的話語.
代表著真理公義的審判.

$^{321}$耶穌不容忍並且要審判罪.
是因為罪毀滅了祂所愛的兒女.
但是請注意.
耶穌在審判以先.
祂給我們有悔改的機會.
是因為耶穌的目的.
從來都不是毀滅.
而是要還回.
祂的嚴厲是出於祂對我們深切的.
就好像一個爸爸.
看到兒子一路走.
越來越近懸崖邊的時候.
這個爸爸一定很心急.
很急切.
大聲跟兒子說.
停呀!不要動!.
回來我這裡.
但是我們今天怎樣停?.
怎樣回轉?.
讓我們將我們的目光.
轉向那位向教會出信的耶穌基督.
耶穌基督二千多年前來到這個世界.
祂的使命是什麼?.
沒錯.
耶穌基督終極的使命.
就是在十字架上成就救恩.
用死贖回我們.
以復活帶來盼望.
但是祂不只是做了這件事.
如果祂只是做這件事.
可能祂幾天已經可以來到這個世界.
但是耶穌在人世間生活了三十多年.
祂花時間體驗世間的甜酸苦辣.
並且耶穌親身進入人群.
親手服侍他們.
摸他們.
醫他們.
餵飽他們.
所以當我們看到福音書記載耶穌所做的事.
其實我們知道祂從來不是將人的需要分開.

$^{361}$分割為一些靈魂的需要.
一些肉身的需要.
耶穌所做的事.
讓我們看到祂在乎的.
從來都是我們全人整全的需要.
靈魂和肉體不是分開的.
所以當我們生病的時候.
我們不舒服的時候.
有痛的時候.
我們的靈魂,我們的情緒,我們的感受是受影響的.
當你知道你的朋友在病患,很痛的時候.
你不可以跟他說你樂觀一點.
你看開一點.
你想多一點開心的回憶就沒事了.
你這樣說是很不人性的.
人的靈魂和肉身是整全的.
所以保羅也教導我們.
恭克己身.
叫身服我.
就好像今天我們做運動的時候.
跑步的時候.
操練身體的時候.
其實我們同時也是在操練,鍛鍊自己的意志力.
以至鍛鍊我們的靈性.
以致我們的生命不是由我們肉體的需要和慾望去主導.
而是由我們人的靈性信仰帶動我們的生命.
既然信仰是全人,是整全的.
我們不單止要調整我們的心思和意念.
我們更加要在行為上悔過或改善.
宣道會有一個很重視的傳統.
不知道大家知不知道.
我們是說來華人宣道會.
宣道會有一個很重視的傳統就是信徒的成聖.
成聖不是靠個人的意志力.
而是藉著聖靈.
在人裡面的耶穌基督親自將聖潔的生命活出來.
所以我們不單止要承認.
我們自己靠肉體得不到聖.
我們要將生命的主權完全交還給神.
我們也要主動和主耶穌基督連結.

$^{401}$和神同行.
隨著時間在我們日常的生活裡.
被主去更新,被主去充滿.
叫我們活出一個積極的生活.
一個具體去對付罪的生活.
剛才說到不同的範疇.
可能有時我們也有軟弱的地方.
在職場裡,在財務裡.
如果你過去在職場裡.
習慣用不同的方言去爭取利益.
或者只是掩飾你的錯誤.
真正的悔改和成聖.
就是當明天你回到辦公室的時候.
你勇敢地.
好像別加摩教會的信徒一樣.
你勇敢地拒絕在那份不實的文件上去簽署.
在感情上,在關係上.
如果你只是處於一段不聖潔的關係.
或者你只是和不應該的人同居.
實質的悔改就是勇敢地搬出去.
離開那段關係.
設立一個聖潔的界線.
因為那段關係不斷地拉扯你離開神.
在你隱密的私生活裡.
倘若你習慣了沉溺於色情的網絡,暴力的影視.
或者你經常依賴算命,星座.
成為你對未來的安全感的時候.
真實的成聖就是今天回家之後.
你徹底刪除那些應用程式.
徹底取消那些會引誘你犯罪的訂閱.
並且將你生命的主權重新交回給神.
是很難的,很不舒服的.
因為悔改一定會帶來痛楚.
因為我們要將肉體裡對世界的安逸,享樂的依賴割捨.
並且讓基督在我們的生命裡掌權.
但是當你願意付上代價悔改的時候.
主耶穌在17節給了我們一個很美麗的應許.
祂說:得勝的,我必將隱藏的嗎哪賜給他.
並且要賜他一塊白石.
石上寫著新的名字.

$^{441}$隱藏的嗎哪.
那些昔日在信仰上妥協的基督徒.
為了肉體的滿足,為了社會的地位.
他們走去吃的是那些異教神殿裡祭偶像的食物.
但主耶穌說:如果你願意為我的緣故.
放棄世界那些骯髒的筵席.
放棄那些不義之財.
放棄那些最終之路.
我就會親自餵養你.
嗎哪是在舊約裡面,曠野裡面神為以色列人賜下的糧食.
代表著神親自的供應.
世俗的利益看起來是很美味.
但吃下去之後只會帶來心靈的枯乾.
帶來被罪捆綁的短暫快樂.
但耶穌說要賜給真正信徒的是隱藏的嗎哪.
哪怕這個嗎哪是俗世的眼光所不能看見.
哪怕這個嗎哪是世人所不能明白.
但他的生命,他的真理的道.
耶穌說他是生命的糧.
叫我們的靈魂得到真正的飽足和滿足.
我們心靈的空虛可以被神的滿足填滿.
還有那塊白石.
白石在羅馬時期有兩個最主要的意義.
第一個就是在法庭上.
當陪審團投下白石代表無罪的宣判.
另外還有一個白石的第二個意義.
就是在當時的聖宴裡.
你拿著白石代表你被邀請的入場券.
耶穌說:沒錯,這個世界可能會因為你不肯妥協.
去定你罪,排擠你,不讓你入場.
但只要你忠心到底.
我就會在天上,在終末的法庭宣判你無罪.
並且我賜給你一張刻著你的名字.
代表著你是被尊崇,被邀請進入黃賢寺的白石入場券.
那種親密,那種無人能夠取代的愛.
只有你和我才明白.
昔日別交摩架會的信徒.
他可以很勇敢去面對暴君的刀劍.
但他們抵擋不了世間世俗,燕洛的誘惑.
他們不怕為耶穌死.

$^{481}$但他們不知道如何在日常的柴米油鹽,工作,娛樂裡為主而活.
今天,我們是否也像一個別交摩架會的信徒?.
弟兄姊妹,讓我們此刻也誠實地.
去檢視一下自己的生命.
有沒有一些地方是我們已經向世界妥協了?.
是否我們為眼前的好處.
為了怕我們被人孤立.
或者單純為了肉體的享受.
我們慢慢有些聖經的原則.
過去我們聽過的,但我們慢慢放下了.
然後我們向世界的價值觀低頭.
主見到我們的掙扎.
主也知道要在這個彎曲薄謀的世代.
要活出聖潔其實是困難.
但是我們的主耶穌愛我們.
我們的主耶穌愛我們.
以致祂不願意看到我們繼續在罪裡面沉淪.
讓我們記得信仰和生活不是分割.
我們不可以將靈魂交給神.
但是我們的身體肉身的主權就是交給世界.
今天耶穌呼召我們要做一個得勝者.
不單是面對大逼迫的時候.
我們要堅守主名.
而是更加要在我們每一天.
每一天很多微不足道的選擇裡.
我們拒絕妥協.
活出聖潔.
願主幫助我們.
不貪戀世俗賢直.
願主幫助我們.
單單去到仰慕神的隱藏的萬能.
並且期盼在永恆裡.
由主的手中接過那塊刻著我們的名字的白石.
我們一起同心低頭祈禱.
主耶穌基督.
你是那位眼目如火焰.
口中套出兩刃利劍的主.
你的話語大有能.
你知道我們的居所.
你知道我們生活在一個充滿著試探.

$^{521}$甚至常常和你的真理去到背道而馳的世界裡.
主也求你赦免我們.
因為我們有時都好像別加摩教會的信徒一樣.
我們雖然宣告你的名.
但是我們生活裡隨從了很多世俗的標準.
我們有時將信仰和生活去到割裂.
濫用了你的恩典.
在工作,在感情,在日常的娛樂生活裡.
我們向罪惡妥協.
求聖靈光照我們.
賜我們有個真實悔改的勇氣.
又求聖靈兼顧我們.
幫助我們明白信仰是整全的.
讓基督的生命在我們裡面真實地活現出來.
也在每一個微小的選擇上.
幫助我們去追求聖潔.
追求一個德性的生命.
也求你賜給我們力量.
去拒絕世俗那些包裝美麗卻帶著毒素的現象.
單單渴望你允許我們隱藏的萬物.
讓你成為我們生命裡最大的滿足.
讓我們一生忠心到底.
不再用恩典成為我們放縱的藉口.
直到那一天.
我們可以從你的手中.
接過那塊刻著我們的名字的白石.
享受和你永遠同住的榮耀.
禱告奉主耶穌基督.
得勝的名字祈求.
\newpage



\section{使徒行傳 8:9-25}
\label{sec:032146}
\textbf{無法複製的聖靈工作}~~~~吳銘淇傳道
\newline
\newline
連結: \href{http://www.church.com.hk/acms/content.asp?site=cdc\&op=show\&type=product\&code=032146\&layout=sermon}{\texttt{032146}} ~~~~ 語音日期: 2026-05-17
\newline
\newline
\hyperref[sec:032145]{\small{< < < PREV SERMON < < <}}
~
\hyperref[sec:index_chronic]{\small{[返順時目]}}
~
\hyperref[sec:032147]{\small{> > > NEXT SERMON > > >}}
\newline
\newline
使徒行傳 8:9-25
\newline
\begin{longtable}{cl}
\hline
\hline
章節 & 經文 (和合本修訂版)\\
\hline
8:9 & \begin{tabularx}{0.7\textwidth}{X} 有一個人名叫西門，向來在那城裡行邪術，自命為大人物，使撒瑪利亞的居民驚奇。 \end{tabularx} \\ \\ \relax
8:10 & \begin{tabularx}{0.7\textwidth}{X} 所有的人，從小到大都聽從他，說：「這個人就是神的能力，那稱為大能者的。」 \end{tabularx} \\ \\ \relax
8:11 & \begin{tabularx}{0.7\textwidth}{X} 他們聽從他，因他很久以來用邪術使他們驚奇。 \end{tabularx} \\ \\ \relax
8:12 & \begin{tabularx}{0.7\textwidth}{X} 當他們信了腓利所傳神國的福音和耶穌基督的名，連男帶女都受了洗。 \end{tabularx} \\ \\ \relax
8:13 & \begin{tabularx}{0.7\textwidth}{X} 西門自己也信了；既受了洗，就常與腓利在一處，看見他所行的神蹟和大異能，就覺得很驚奇。 \end{tabularx} \\ \\ \relax
8:14 & \begin{tabularx}{0.7\textwidth}{X} 在耶路撒冷的使徒聽見撒瑪利亞人領受了神的道，就打發彼得和約翰到他們那裡去。 \end{tabularx} \\ \\ \relax
8:15 & \begin{tabularx}{0.7\textwidth}{X} 兩個人下去，就為他們禱告，要讓他們領受聖靈， \end{tabularx} \\ \\ \relax
8:16 & \begin{tabularx}{0.7\textwidth}{X} 因為聖靈還沒有降在他們任何一個人身上，他們只奉主耶穌的名受了洗。 \end{tabularx} \\ \\ \relax
8:17 & \begin{tabularx}{0.7\textwidth}{X} 於是使徒按手在他們頭上，他們就領受了聖靈。 \end{tabularx} \\ \\ \relax
8:18 & \begin{tabularx}{0.7\textwidth}{X} 西門看見使徒一按手，就有聖靈賜下，就拿錢給使徒， \end{tabularx} \\ \\ \relax
8:19 & \begin{tabularx}{0.7\textwidth}{X} 說：「請把這權柄也給我，使我手按著誰，誰就可以領受聖靈。」 \end{tabularx} \\ \\ \relax
8:20 & \begin{tabularx}{0.7\textwidth}{X} 彼得對他說：「你的銀子和你一同滅亡吧！因為你想神的恩賜是可以用錢買的。 \end{tabularx} \\ \\ \relax
8:21 & \begin{tabularx}{0.7\textwidth}{X} 你在這道上無份無關；因為你在神面前心懷不正。 \end{tabularx} \\ \\ \relax
8:22 & \begin{tabularx}{0.7\textwidth}{X} 你要為你這樣的惡而悔改，祈求主，或者你心裡的意念可得赦免。 \end{tabularx} \\ \\ \relax
8:23 & \begin{tabularx}{0.7\textwidth}{X} 我看出你正在苦膽之中，被不義捆綁著。」 \end{tabularx} \\ \\ \relax
8:24 & \begin{tabularx}{0.7\textwidth}{X} 西門回答說：「請你們為我求主，使你們所說的，沒有一樣臨到我身上。」 \end{tabularx} \\ \\ \relax
8:25 & \begin{tabularx}{0.7\textwidth}{X} 使徒既作了見證，並且宣講了主的道，就回耶路撒冷去，一路在撒瑪利亞好些村莊傳揚福音。 \end{tabularx} \\ \\
[1ex]
\hline
\hline
\end{longtable}
$^{1}$你是輔召我們的主 讓我們能夠相聚.
也讓我們能夠一同去聆聽上帝的話語.
願上帝今日藉著孩子的講道.
也讓我們一同去聽道.
叫聖靈今日在我們當中.
我們去被提醒 被鼓勵.
也要讓我們的生命被改變.
逢著耶穌的聖名祈求 阿們.
然後他很開心.
他覺得自己很厲害.
好像大人一樣 像媽媽一樣.
推著車子把家裡所有有需要的東西都買回去.
之後媽媽就停下了.
放下了車子 然後就買洗衣粉 買清潔用品.
還有小朋友 這些當然不是我們的興趣.
所以他當然想推著車子去買薯片.
突然兒子發現車子不動.
明明他也在用力 但車子是不會向前的.
為什麼他怎麼推都推不動呢.
弟兄姊妹你猜猜你知不知道為什麼.
當然是因為一直在推車子的不是小朋友自己.
而是背後的媽媽在推.
小朋友只是扶著車子 以為自己在推車子.
有時候基督徒也好像小朋友一樣.
我們以為教會的增長是因為我們有一個很好的方法.
我們以為某些事工的興旺是因為我們的活動有夠吸引力.
甚至我們以為人生命的改變是因為我們有某些技巧.
但慢慢我們就會好像小朋友一樣.
因為車子是我們推動的.
所以《士徒行傳》第八章的故事正正要拆毀我們這種想法.
《士徒行傳》第八章說什麼呢.
當中路加描述當時的基督徒被猶太人逼迫.
所以他們開始離開耶路撒冷.
同時福音傳到更多的地方.
在八章四節說那些分散的人要去到各處傳道.
而各處在哪裡呢 八至十一章裡路加告訴讀者.
究竟這些人去到各處傳道的時候發生了什麼事.
而說到要離開耶路撒冷開始傳道工作的時候.
當然要包括離耶路撒冷不遠.
只是隔六十公里的撒瑪利亞地區.

$^{41}$路加就說到肥利.
肥利是什麼人呢.
就是在初期教會負責管理飯食的其中一個執事.
就要去到撒瑪利亞傳講福音.
另外肥利也不斷去行一些神蹟.
包括人從鬼屋裡面去釋放.
也可以在癱瘓和瘸腿的人得著醫治.
由於這些神蹟令到長年受鬼和疾病困擾的人得著釋放.
所以肥利在當中的時候 大家都是大大的喜樂.
路加要講肥利的工作.
為了要帶出我們今天要講的西門.
今天要講的故事的主人翁.
我們會分三個部分去討論一下今天的故事.
首先是講西門信主的經歷.
這個可以讓我們去想一下信徒的動機.
之後我們會去看西門對彼得和約翰的請求.
這裡是說我們可能會有面對的陷阱.
最後就是彼得的回應.
也讓我們去想一下如何去重整.
事奉誰才有那個控制權.
我們首先進入第一個部分.
八章九到十三節.
講的是關於西門信主的一個經歷.
西門是一個什麼人呢?.
西門是一個行邪術的人.
行邪術是講什麼呢?.
就是當時用咒語.
為了影響或控制超自然力量的儀式.
雖然在摩西律法已經禁止一切形式的邪術.
但似乎還未阻止到猶太人對神秘主義的興趣.
他們仍然有很多人是行邪術.
而西門甚至是透過販賣咒語去賺取金錢.
當時行邪術的人很多.
而西門似乎是這個行頭的代表.
就是行頭的大佬一個.
經文形容八章九到十節.
他說他是妄自尊大洗撒瑪利亞的百姓都驚奇.
無論大小都會聽從他.
說這個人就是稱為神的大能者.
其實這句說話是有得被證明的.

$^{81}$尤斯汀他是一個殉道者.
但同樣他也是一個撒瑪利亞人.
所以他擁有著很多關於西門的一些資訊.
都幫助我們去明白當時的一些狀況.
尤斯汀說西門他會怎樣自稱.
他更加詳細.
他說他是超越一切權能 權柄 能力以上的一個神.
好像很厲害 說他和神是一樣的.
而撒瑪利亞人也承認他是守衛的神.
所以其實當時一個情況下.
他真的好像神一樣的存在.
甚至有些文獻說.
什麼驅使人們會這樣相信呢.
撒瑪利亞人都不會這樣相信.
有一個這樣的人會是神.
為什麼他會這樣相信呢.
原來是因為西門當時是有些異能的.
他有些東西是能人所不能.
就是他能夠懸浮.
就好像飛的一樣.
所以因此在羅馬都聽過他的名字.
覺得這個人很厲害.
所以我們可以想像一下.
當時西門自稱是神.
也擁有一些能人所不能的能力.
所以當撒瑪利亞人見到的時候.
就說他是神.
他又做了這些飛的行為.
所以他應該是真的.
所以不單覺得他是神.
聽他的話.
而西門不單有能力.
也有權力.
這個就是當時的情況.
不過肥利的出現.
好像一切都改變了.
在經文裡告訴我們.
人不再聽西門說話.
因為西門只是精通這些邪術.
而是反過來聽肥利說話.

$^{121}$他們也不再因為西門.
藉著這些邪術所做的事而感到驚奇.
反過來因為肥利所行的神蹟而感到驚奇.
他也不再承認西門對自己是神的宣稱.
因為他相信只有耶穌才是拯救他們的神.
也知道天國將會降臨的好消息.
因著肥利將福音傳給撒瑪利亞人.
上帝當中的工作.
所以無論男女都受洗.
他們都毀身在信奉基督的記號當中.
社會有這麼大的轉變.
其實已經很厲害.
你試想像一下.
如果今天在查理斯堡裡面.
所有的居民都信了主.
然後浸禮是排著隊來浸禮.
然後謝牧師就要一個一個浸.
他應該做到了.
因為剛剛放完假回來.
所以他就要一個一個浸.
我相信這件事不單止是轟動在靈福堂裡面.
應該是整個查理斯堡都會是一個很大的新聞.
作為一個看了這麼多聖經的故事.
可能你都會想.
通常去到這些場景的時候.
都不會有好事發生.
因為西門應該會感到受到威脅.
他應該會有些回應.
可能他會挑戰肥利.
然後看看誰的異能更厲害.
不過經文沒有這樣說.
路加直接說.
西門不單止沒有直視肥利是他的敵人.
甚至他還信了主.
也受了洗.
因為路加想藉著西門的信.
去進一步和他的讀者說.
你看看福音的力量.
你看看上帝的工作是多麼厲害.
可能我們和當時的讀者都一樣.

$^{161}$是不是真的.
怎麼可能是這樣呢.
會不會西門其實他都不是真心悔改呢.
其實可能他都想刺探一下君丞.
看看為什麼這個人這麼厲害.
為什麼他可以拉走我們的客人.
會不會其實畢竟故事繼續發展下去.
西門都會有很多問題.
但是路加透過三個行動.
很簡潔的去告訴他的讀者.
西門的信是很真實的.
首先他說西門是信了.
然後做了什麼呢?.
就是受了洗.
他作為他的悔改和承認耶穌是彌賽亞的一個認信.
而這個認信是帶著一個行為.
就是西門和肥利在一起.
一同的生活.
西門因著受了信.
所以就繼續和這些初期教會的信徒一起生活.
並且因著上帝的行為.
整天都感到很震撼.
西門就好像一個初信的信徒一樣.
他是真心悔改的.
他真的相信耶穌基督的拯救.
而在教會的生活裡面.
就讓他繼續認識上帝.
學習怎樣成為信徒在主的當中.
在上帝的裡面看見上帝的作為而感到驚歎.
路加沒有暗示說其實是肥利太心急了.
他經驗尚淺.
所以太早被西門受針.
他被西門欺騙了.
也沒有說西門怎樣心懷不軌.
他只是很簡單地說.
他信了也受了洗.
他真的好像我們每一個信徒一樣.
的確西門和我們每一個都好像沒有分別.
都是一樣的信徒.
但他和我們每一個都一樣.

$^{201}$是他在信主開始改變他的生命過程當中.
難免有時候都會有舊我.
一些舊有的觀念.
怎樣去理解這個信仰.
理解我們這個上帝.
八章十四到十六節那裡說.
使徒去到耶路撒冷.
見到撒瑪利亞人領受了神的道.
之後就打發了彼得和約翰去到他們那裡.
兩個人去到就為他們祈禱.
要叫他們受聖靈.
因為聖靈沒有降臨在他們一個人的身上.
他們只是奉了主耶穌的名去受了洗.
撒瑪利亞的信徒信了主也受了浸.
但聖靈都沒有降臨在他們的身上.
其實這是一件很不尋常的事.
因為在使徒行傳的其他經文裡.
通常都是說他們信了主的時候就領受聖靈.
但是路加沒有因著這個差異去做什麼解釋.
不過在這段經文裡.
我們似乎都很看見.
上帝是刻意要等彼得和約翰去到那裡.
然後才讓撒瑪利亞的信徒領受聖靈.
或者可能因為當時我們一直都知道.
猶太人和撒瑪利亞人是水火不容的.
關係非常差.
一早已經分了家 彼此的不信任.
上帝似乎藉著使徒的出現.
要說兩件事.
一方面要告訴耶路撒冷教會.
撒瑪利亞人是屬於上帝的.
但同時也要讓撒瑪利亞的信徒知道.
他們不是另外一個分開的群體.
他們和其他信徒都是同屬於基督的教會當中.
而這件事同樣有另一個效果.
就是讓西門暴露自己對上帝的誤解.
彼得和約翰按手在撒瑪利亞的信徒身上.
他們領受了聖靈.
在八章十七節裡說.
我們相信領受聖靈的過程.

$^{241}$在當時的過程裡.
應該是有些東西可以看到或聽到的東西.
去彰顯出來.
所以才引起了西門的注意.
印著西門以前一直對權力或控制.
都有很多的想法.
他以為上帝和他以前所認識的靈界力量是一樣.
只要掌握到一個正確的方法或技巧.
甚至你有相應的權柄的時候.
就能夠帶出同樣的果效.
因為西門本身也是用咒語.
就能夠覺得可以操控到超自然的力量.
所以結果他就提出了一個很大膽的想法.
就是一個操控上帝的想法.
他拿錢給彼得和約翰就說.
你把這個權柄給我吧.
叫我可以按手在誰的身上.
誰就可以領受聖靈.
其實你可以想一下.
西門竟然敢光明正大去說這件事.
其實就證明了他說的時候.
他真的不覺得這句說話有多有問題.
我們今天一定會嘩然覺得.
怎麼可能說這些東西.
你想都不要說.
但他當時可以這樣說.
就是說他真的覺得這樣是沒有問題.
你試想像一下有沒有幾年前.
你在那些社交平台裡面.
有見過有人分享自己和一個五歲的兒子.
坐在車的駕駛座那裡.
然後小朋友扶著tie 盤.
大人又扶著tie 盤.
然後他們就慢慢踩油門開車.
其實你可以想一下.
為什麼會分享這段影片呢.
當時我相信他一心一意.
是想分享一下我和兒子的歡樂時光.
又或者是和兒子一起分享對汽車的熱愛.
我想他應該在分享的時候.

$^{281}$他從來都沒有想過.
當其他人看到的時候.
會告訴他.
喂 你知不知道這樣是犯法的.
喂 你知不知道你在路上.
被小朋友這樣開著tie 盤開車.
其實對道路使用者有多大的危險.
他沒有想過.
因為他那一刻是想分享他的喜悅.
可能西門也一樣.
他不是一下子就去想要怎樣去抵擋上帝.
甚至他也不是懷著一個壞的心腸.
他怎樣去想自己的好處.
怎樣用這件事去賺更多的錢.
我們都沒有辦法去證實這件事.
可能他只是想買這些能力.
就好像使徒一樣.
去和不同的地方去傳道.
行神蹟.
可以幫助一下這些使徒.
然後將聖靈傳到當地的信徒當中.
他認為自己都可以用錢去買到這個權力.
他就好像可以像使徒一樣.
去一起參與在這個服事的裡面.
不過西門不明白的是.
聖靈不是一種被人操控的能力.
希望操控聖靈.
操控上帝是當時西門的一個陷阱.
但可能也是今天教會的一個陷阱.
我們經常說教會增長學.
甚麼是教會增長學呢.
增長學本身是宣教士去到一些陌生的國度.
他要開始建立教會.
所以他們就嘗試去制定一些策略.
一些方法.
希望可以更有效地進入人群當中去建立教會.
這個原本是好的.
因為宣教士彼此要分享一些經驗.
希望可以啟發到其他或者更新的宣教士.
在他們的工場裡面怎樣去服事.

$^{321}$但是再好的方法.
經過時間也會開始慢慢失去焦點.
有些教會領袖開始用世界怎樣營運的模式.
去理解教會的增長.
例如可能用管理的思維去理解教會.
好像有限公司一樣.
又或者是用量化的果效去衡量.
甚麼是成功 甚麼是失敗.
甚至用消費的文化去建構崇拜.
建構教會的聚會.
但是其實這樣的想法會產生很多的問題.
就是當教會以為人數多或者增長得夠快.
就能夠衡量到這間教會是否健康.
它的屬靈狀態是怎樣.
但是這樣的話.
最終教會只能夠看到上帝在人數上有沒有做甚麼工作.
但卻看不到上帝在信徒的生命裡面在做甚麼工作.
當教會以人的感受和需要為優先.
期望信息能夠怎樣輕鬆.
怎樣能夠有話題性.
甚至敬拜的音樂和畫面要怎樣去吸引到人.
但是這個讓我們在會中.
在崇拜裡面最終學習到的是怎樣去滿足人.
卻是不能夠去為了敬拜.
它不是為了要敬拜上帝.
由當教會以為藉著複製一套的方法.
就能夠帶來果效.
結果就是忽視了上帝一直以來的工作和會中.
有甚麼特質.
他們的文化是甚麼.
其實教會都很容易會跌入這些陷阱.
特別是面對著困難的時候.
當面對著流失.
面對著下一代不肯再回教會.
聚會人數減少的時候.
我們都很自然地會問.
怎樣做才會好一點呢.
我們未必會因為這樣就要操控上帝.
我們就會好一點.
未必會這樣想.

$^{361}$但很多時候當教會感到失控的時候.
我們就很自然地對某些曾經有效的方法.
或者是其他教會有效的方法.
我們會產生一種依賴.
久而久之.
甚至我們會將盼望都放在這些制度.
這些流程和技巧當中.
所以今天值得我們去反思的是.
到底我們是否在尋求上帝呢.
還是我們都慢慢將這個安全感.
本身在上帝裡面轉移到某些.
我們認為有效的方法上面呢.
我們理解了西門.
應對了今天在教會裡面.
講的盲點是什麼的時候.
我們又回到經文.
八章廿二十到二十五節裡面.
給了一個很嚴厲的回應.
去揭示西門的一個問題.
也揭示了一個屬靈的能力.
很重要的一個認知.
就是上帝的能力是恩賜.
這個不是可以透過技術.
又或者怎樣的代價.
就可以換來或者複製到的一個商品.
彼得的回應可以分為四個部分.
第一他說的是關於審判.
彼得說願你和你的銀子一同滅亡.
因為西門將上帝當成了一場交易.
可以操控的力量.
所以彼得要在這裡提醒他.
甚至警告他.
就是上帝的恩賜不是人可以用金錢.
或者用技巧.
甚至你付什麼代價就可以換回來的.
第二就是講了西門的信心的缺陷.
彼得說你在這個道上面是無分無關的.
因為你的心在神的面前是不正.
彼得在這裡要講的是.
西門雖然已經相信了耶穌.

$^{401}$但是他仍然很想擁有很多的能力.
甚至可以掌握到這些能力.
他的問題不是單單做錯了說錯了這件事.
而是彼得更加看見他的心.
就是他想將上帝的工作.
變成一些他能夠控制.
他都能夠去掌控的話事的一件事.
因此彼得就要呼籲他悔改.
彼得說你要去懊悔你的罪惡.
要去祈求主.
因為西門他的動機是不對的.
但是上帝仍然給他機會回轉.
所以彼得就去呼籲西門要去悔改.
最後就是要去揭露.
西門不單止犯罪得罪了上帝.
甚至他也被罪所傷.
彼得說我看出來正在苦膽之中.
被罪惡所捆綁.
彼得要講的是西門因著過去.
他是這樣理解權力的.
他是這樣理解能力.
這樣去理解怎樣去控制.
所以結果他就被這些控制的慾望.
去綁住他自己.
這個捆綁就限制他看不到上帝那種超越的能力.
而不是上帝的主動.
他也看不到上帝高於他的能力.
是來自於上帝自己的主動權.
他認為一切都要靠他的手去做.
才可以做得到.
他看不到上帝的那種超越性.
當我們覺得自己都要去擔起教會的職責.
擔起教會的重擔的時候.
我們去找一個方法.
希望可以解決教會一切的問題.
其實這樣不單止得罪了上帝.
也是因為我們過份依賴著某些方法的時候.
而不去再依靠我們的上帝.
更困難的問題是什麼呢.
就是我們知道這些東西都沒有辦法可以完完全全.

$^{441}$所有的問題都可以被解決.
都會被這些想法所傷.
很快我們就會撐不住.
然後因為這些根本就超越我們能力所及.
而使徒行傳要提醒我們的是什麼呢.
就是真正改變人生命的不是人的方法.
而是上帝自己.
無所教會的都是上帝.
信徒仍然要很努力地去準備.
彼此的學習 忠心的服侍.
我們永遠都不能夠以為.
我們只要掌握到正確的方法.
就能夠去控制聖靈怎樣去工作.
西門的回應都很特別的.
在八章二十四節這樣說.
他說他不再好像以前一樣.
他沒有提出一些他自己的方法或想法.
反過來他退下來.
然後他謙卑去請求彼得為他代求.
真正的信仰.
很多時候都不是我們掌握到什麼.
而是我們願意承認.
我們不能夠去掌握這個上帝.
我們要交回給上帝.
今年教會的主題是.
在基督裡同心前行.
其實最初主題只是定在同心前行的裡面.
不過同工會開會的時候.
有牧者就說應該要加回在基督裡.
其實坦白說一開始聽的時候.
心裡就想「不用了吧」.
「這不是理所當然的嗎?」.
「有什麼需要特別寫出來呢?」.
但今天這段經文對我有一個提醒.
就是原來真的需要.
因為在基督裡同心前行.
或者我們說同心前行.
都不等於一定是在基督裡.
今天我們準備進入五一八.
我們經常說三禹要同一屋禪客.

$^{481}$一起聚會,事奉和生活.
我們都會知道.
同心前行必然是教會一個很重要的大方向.
但問題是我們靠著什麼去讓我們同心呢?.
今天的經文提醒我們.
真正叫教會合一的.
從來不是人有什麼技巧.
我們都不是教會的主人.
甚至我們都沒有辦法掌控教會的方向.
我們想去合一他們就去合一.
我們也沒有辦法靠著自己的能力.
去保證教會一定能夠做到合一.
其實我們都知道.
三禹大家的文化不同.
一直很久都沒有聚在一起.
要合一很難.
因為這是上帝的工作.
不是人能夠掌握技巧帶來的後果.
所以我們只能夠懇求上帝的工作.
懇求上帝去改變我們.
懇求上帝讓我們更願意放下自己.
去看見別人的需要.
甚至我們懇求上帝讓我們去看見.
今天他想塑造我們成為一個怎樣的群體.
所以今天的挑戰是.
與其我們去想很多的方法.
可以的.
但是我們有沒有更迫切的.
去為我們的合一去一起禱告呢?.
如果你認為這是上帝的工作的時候.
最後 弟兄姊妹.
還記不記得一開始說的那個手推車的故事.
很多時候我們都好像那個小朋友一樣.
以為是我們推動了那個車.
但是我們以為教會只要我們能夠一起努力.
就能夠撐起的了.
教會的增長是我們努力就能夠帶來的.
合一是我們努力就可以做到.
但是今天我們要知道的是.
原來真正在背後推動教會前行的.

$^{521}$從來都不是我們.
是上帝自己.
所以弟兄姊妹.
在基督裡面同心前行.
不是我們一起努力去推動上帝叫他去做事.
而是我們一起去學習.
怎樣去跟隨主.
懇求上帝讓我們知道他的心意.
貼近他的心意.
再為主去工作.
讓我們同土共.
主我們去懇求你施恩和憐憫我們.
就是在我們的生命當中.
都有自大的時候.
都有想控制的時候.
或者在無力的裡面.
我們會找尋更多不是從你而來的.
我們希望用這些方法.
可以去讓教會成長.
但是我們知道.
這些都不是一個你喜悅的方法.
願上帝我們今天就謙卑在你的裡面.
看見上帝你的工作.
而感到驚奇.
求上帝你繼續祝福和帶領我們.
在合一的事上.
求你去改變我們的心.
求你改變我們.
讓我們成為一個願意看見別人需要的群體.
奉主耶穌基督的聖名祈求.
\newpage



\section{約書亞記 12:7}
\label{sec:032147}
\textbf{聖經人聖經事:從北到南的勝利,因應許而來}~~~~黃克勤牧師博士
\newline
\newline
連結: \href{http://www.church.com.hk/acms/content.asp?site=cdc\&op=show\&type=product\&code=032147\&layout=sermon}{\texttt{032147}} ~~~~ 語音日期: 2026-05-17
\newline
\newline
\hyperref[sec:032146]{\small{< < < PREV SERMON < < <}}
~
\hyperref[sec:index_chronic]{\small{[返順時目]}}
~
\hyperref[sec:032148]{\small{> > > NEXT SERMON > > >}}
\newline
\newline
約書亞記 12:7
\newline
\begin{longtable}{cl}
\hline
\hline
章節 & 經文 (和合本修訂版)\\
\hline
12:7 & \begin{tabularx}{0.7\textwidth}{X} 這些是約書亞和以色列人在約旦河西所擊殺的諸王，他們的地從黎巴嫩平原的巴力‧迦得，直上到西珥的哈拉山。約書亞按著以色列支派所得的份把這地分給他們為業， \end{tabularx} \\ \\
[1ex]
\hline
\hline
\end{longtable}
$^{1}$各位聽眾姐妹大家好 又和大家見面了.
今天和大家就是透過今天的經文去了解.
聖經人,聖經事.
看看背後上帝奇妙的工作.
這個是生命樹宣教網絡為大家去準備的.
都經過很長的日子了.
我們今天繼續看約書亞.
約書亞這個人物雖然我們都聽過很多次,認識他.
但其實你會發現比較少人去說約書亞怎樣怎樣.
都很少的.
當時如果說舊約就說亞伯拉罕,雅各,摩西,羅亞.
或者後來有大衛或者有些皇帝.
比較少提約書亞.
他似乎只不過是摩西的接班人,助手.
但其實他是一個非常重要的人物.
因為他是上帝指明他接替摩西.
帶以色列人進去加拿大.
而且是由於他的勝利.
而帶給整個的以色列人.
在他們進加拿大之後怎樣去分土地.
以及怎樣得勝的.
今天就來到一個總結了.
你記得昨天就說到北面的皇帝.
北面的很多地方很多不是皇帝叫王的.
很多北面的各個城的王.
特別是下所城的王就聽到.
以色列人在南方這麼厲害.
在出埃及,在曠野這麼厲害.
而且將當地的王都殺死掉在樹上面.
害怕了 整個北方的王就幾十個一起.
數不了這麼多 一起去進攻以色列.
預備打仗了 趕他們走 離開應許之地了.
前天就說南面五個王一起聯合.
要抗拒以色列人.
因為看到他們進來的時候.
耶利哥在艾城這麼厲害.
結果昨天看到北面的王.
全部都不是以色列的對手.
打臉他們是得勝了.
今天就總結了.

$^{41}$今天的題目是我自己找出來的.
自己想的.
是從北到南的勝利是因為應許而來.
其實就是說是因為神的應許而他們得勝.
因為傳而得勝.
由北到南的地方得勝.
而且在聖經裡面是記載了的.
今天和大家去看一下.
從北到南的勝利 因應許而來.
讀這段聖經 心裡面是非常感恩和讚美的.
今天和大家看的經文.
就是在約書亞記第十二章第七節.
經文是這樣說的.
約書亞和以色列人在約旦河西擊殺了諸王.
他們的地是從黎巴嫩平原的巴力加德.
直到上西爾的哈拉山.
約書亞就將那地按著以色列支派的宗族分給他們.
這節聖經大家都能讀到.
但是你會發現是很少人看.
很少人讀.
很少人去欣賞它的.
再看一次.
約書亞和以色列人在約旦河西擊殺了諸王.
是哪個人呢?.
是約書亞和以色列人一起.
在約旦河西.
約旦河西就是我們今天比較多看到.
就算去到今天的以色列.
大部分地方都是約旦河西的.
而河東那邊其實也有三個支派.
後來你看到.
其實約旦河東的一片土地.
都是給了以色列的.
但是後來在王國被擄之後.
被人霸佔了.
一直霸佔搶到現在.
約旦河東岸的地方比較多現象是屬於約旦的.
今天.
在這裡約旦河西和河東都擊殺了諸王.
他們的地是從黎巴嫩平原的.

$^{81}$記得這個highlighting.
就是巴黎加德和哈拉山.
就是黎巴嫩平原和西爾這兩個地方.
我們先介紹一下.
一會兒有地圖.
巴黎加德.
其實巴黎.
以色列人拜巴黎.
迦南地拜巴黎.
其實可能是指一個泛指當地所拜的每一個王.
每一個族.
每一個不同的人.
他們拜不同的神.
他們大部分都用相同的字叫做巴黎.
巴黎.
就是加德.
這是一個城的名字.
加德就是好運.
Bio就是主.
巴黎加德的意思就是好運的主.
不過這次不好運.
在黎巴嫩最北的地方.
西爾.
去到最西東南的地方.
按著以色列的支派.
約書亞這個人很重要的地方.
是他按著以色列支派.
來分地給他們為業.
這就是今天要說的地方.
看看下一張.
這張地圖在大家看.
特別是用手機就更加小了.
但是不要緊.
看到顏色和幾個重要的地方就行了.
我會放大它.
你看到這裡有一個湖.
基尼列湖就是這個位置.
其實就是今天的加利利海.
而這條河流一直下來.
這條河流就叫做約旦河.

$^{121}$而這個阿拉巴海就是死海或者嚴海以前的名字.
今天經過了幾千年.
地圖都改了.
以前阿拉巴海就是嚴海是一個海.
現在已經中間水不夠斷了.
就變成上下嚴海了.
黎巴嫩在哪裡呢?.
黎巴嫩是指這個地方.
你看到地圖寫著大馬的.
其實這個部份.
這個部份就是敘利亞.
敘利亞很大的.
敘利亞就是亞屬.
如果下去敘利亞就很遠了.
應該更寬一點.
敘利亞是這樣的.
這裡是黎巴嫩.
這個地方就是以色列.
但是在這個圖來說.
特別要標耶路撒冷.
大概耶路撒冷是在這個位置.
這裡是地中海.
有兩個顏色的.
黃色和粉紅色.
黃色是以德之地.
就是他們經過這麼多場打仗.
由北面的巴黎加達.
整個去到南面的西伊.
就是這個地方.
黃色是他們已經打仗回來獲得的.
在應許之地的土地.
如果昨天你讀.
就會發覺到.
這兩天讀了有些地方.
就連裡面的阿納族人都殺死了.
不過在加薩,阿特基倫,以甲倫這些地方.
沒有完全將它毀滅.
你看到原來這塊地方.
粉紅色這塊.
就是今天的叫做非利士地.

$^{161}$或者今天一直都是和以色列打仗.
打到今天的非利士人.
或者叫做加薩地帶.
是那些叫做巴勒斯坦自治政府裡面.
是很極端的人射的火箭出來.
這個他們沒有在約瑟克時代將它擊敗.
拿了它那塊地.
都叫做拿了的.
但是都沒有將它們全部殺死.
而留到今天成為一個.
由大陸開始.
你看到有這個巨人哥利亞.
就是出自非利士自治和以色列對抗.
這裡上面就是西頓,推羅,西頓這一帶地方.
其實黎巴嫩的地方.
就是以色列沒有拿到的.
還有大馬色這一段.
也都沒有拿到的.
就是一直和今天對抗.
中間有一點點加瑪和基塞都沒有拿到的.
就是屬於加里里海的對面.
而上面這個地方.
我現在線上指的地方.
就是今天叫做哥蘭高地.
對面這裡也都已經是屬於亞述的了.
如果今天來說.
亞述就是敘利亞了.
剩下的地方.
粉紅色就是未得之地.
黃色就是他們得了的.
好 看看下一張.
我先放大.
如果放大了.
放大了你會看到剛才我們說過的.
以色列在約書亞帶領之下.
而得到的土地就是黃色這些.
而紅色的粉紅色未得.
最北的地方就是這裡了.
最北的地方在這裡.
就是巴黎加德.

$^{201}$就是這個位置.
最尖的一個地方.
最北的了.
是他們得到的最北的地方.
就是屬於在黎巴嫩平原.
最北的地方.
下面下去.
就會沿著大河下去.
去到加里尼海了.
然後大河繼續下去.
很多個王在這裡.
每一個城都有個王.
你看到城的名字.
其實裡面都有個王在這裡.
是做了王的.
而黑門山大概就在巴黎加德這個字遮住.
就是黑門山.
去到北的黑門山.
看下一張就是剛才地圖的下半.
你會看到.
黃色的地方都是以色列得到的.
粉紅色整個菲律賓.
就是加薩,阿塞斯特,阿塞基倫,以加倫.
這邊的土地就是菲律賓地.
好了.
西耳的哈拉山.
剛才讀的經文.
在北邊就由巴黎加德去到西南邊的西耳.
哈拉山都是屬於以色列的.
剛才是看北.
今天是看南.
如果我們查考聖經.
或者看各樣的資料.
其實是找不到哈拉山在哪裡的.
不過有兩個字.
讓我們能夠知道它大概在那附近就是西耳.
西耳其實就是「毛」的意思.
很多毛.
就是以東或者叫以塑的另一個名字.
以塑就是叫做「毛」.

$^{241}$很多毛.
以東的意思就是紅色.
以塑後來離開了.
沒有住在巴勒斯坦.
是住在這邊這個山上.
這個山區上.
就是哈拉山.
哈拉山可能是很多山之一.
以色列得到這裡.
由於哈拉山不知道在哪裡.
所以找不回來.
所以畫地圖的人沒有將黃色都畫在這裡.
不知道在哪裡.
可能亂畫了.
所以沒有畫.
其實由北面的巴勒加德.
直到南面的哈拉山.
西耳哈拉山.
都是他們獲得的.
西耳哈拉山後來沒有再在以色列手中堆砌.
因為聖經是說.
我不將西耳以東那邊給以色列人.
因為神說我已經給了以塑.
和他的後裔住在以東地.
就是說西耳是給以塑的.
而這個黃色的英雄之地.
特別是約旦河東邊.
約旦河東部這些地區.
就是約旦河西的小地區.
是給以色列的.
其實開頭以色列有幾個支派.
在約旦河的東邊.
他們是沒有進入英雄之地的.
在外面.
他們都佔領了.
大家看到聖經是描述.
從巴勒加德直到南面的西耳哈拉山.
西耳哈拉山可能是向南面更低一點.
不過地圖所限表現不了.
好.

$^{281}$我想大家看看剛才整個地圖是這樣的.
西耳哈拉山可能是更低一點也說不定.
更低到很低也說不定.
低多一點.
我想要看的.
再來一次.
第一點我想和大家說的.
神賜給他們的第一.
是賜給他們有地的.
有塊地.
地就是由巴勒加德直到哈拉山給他們的.
於是他們就按著以色列支派.
分離做他們的地業.
他們就分了地了.
是誰做的?是約書亞帶領他們做的.
約書亞代表上帝將地給他們.
所以他們進入英雄之地的第一件事.
就是真的得到.
第一是地.
是上帝給的.
好.
然後第二.
聖經就說了.
第八節.
就是黑人,阿摩利人,迦南人,比利西人,希美人,耶布斯人的山地.
高原,阿拉伯,山坡,曠野和藍地.
這裡巴含的人數之不盡.
有幾個大的.
在這裡約書亞說了.
有黑人,有阿摩利人,迦南人,比利西人,希美人,耶布斯人.
耶布斯人其實就是約瑟蘭一帶地區的山地.
約瑟蘭是山地.
另外很多人沒有給名字.
就是他們住在高原,阿拉伯一帶.
有些是住在山坡的一些人.
有些是住在曠野平地.
沒有人去的人.
有些住在藍地.
更加不說了.
藍地是什麼呢?.

$^{321}$如果以色列由北向南走.
如果今天根據一條高速公路.
或者由舊約習黨藍地.
就是一直下了耶路撒冷.
然後再向南走.
走完死海,嚴海.
再向南走到很遠.
去到真正的紅海那麼遠.
大概是今天高速公路都有六個小時車.
或者多一些.
這樣的時候去到那裡.
以下那些全部都是荒地,曠野.
將它稱為藍邊那塊土地.
藍地,去哪裡?去藍地.
很荒蕪.
藍地.
整個以色列地.
英海之地有多少種人呢?.
很多.
第一就是我將地給你.
第二.
我將裡面的人都給你.
不過這些人.
全部地都不要留了.
雖然拿了很多人.
我都可以做什麼?.
做奴隸.
這樣就可以幫忙耕田.
做粗工.
男的可以幫忙種田.
其實神說不要稀罕這些.
因為我知道他們會帶領你.
會拜那些假神.
拜偶像.
其實神甚至在舊約聖經之前.
摩西已經說.
其實我知道你們.
進到迦南地會被那些假神.
來騷擾你.
然後你個個都拜了.

$^{361}$你們很迅速就會拜壞了.
這樣說.
所以.
神要將他們滅絕了.
有機會再說一下.
神為什麼會這樣滅絕他們呢?.
然後將地給以色列人.
話說回來.
這麼多種人在這裡住.
你猜在迦南地裡面有多少人口呢?.
我也不知道.
因為沒有什麼人統計過.
沒有人說過出來.
但是你知道以色列人出埃及的時候.
很有可能是近200萬人.
經過了這個抗疫40年.
當代的人都死了.
他們全部進迦南的人都是在抗疫出生的.
我相信他們都補回這個數字.
百多二百萬都可能有.
而且久了你繁衍得更加多.
整個地如果有二三百萬人在這裡住.
還不算很擠.
不過不會說很空.
因為你看到有這麼多種人在裡面住.
每個都幾萬幾萬幾萬.
那裡一百幾十萬一定有.
而且未必有以色列人這麼多.
所以他們對於以色列這麼多人要進來住.
他們是很惶恐,很驚慌.
而且想著如果你進來.
你們的耶和華神有這麼好打.
我們沒有,一定要和你對抗.
所以南邊的五個王,北邊的無數王.
都要和以色列打仗.
就是這個意思.
神真是按著應許帶領他們得勝了.
記得如果約瑟奇一開始就攻艾城.
在耶利哥之後.
如果沒有一個惹麻煩者叫阿甘.

$^{401}$如果沒有這個人.
就應該他們一直都勝利的.
是神特別帶領他們勝利的.
早兩次都說了.
神落駁來攻擊他們.
那些兵搏大鬥到打死的人比被以色列的士兵都殺的還多.
就是說神親自介入.
幫助他們獲得這一塊應許的地方.
這個就是第一,得到地.
第二,神真是讓他們得到人了.
第三,真是稀奇.
因為真的很少讀.
第九字,他們的王.
這個三字是指我要講的第三點.
他們的王有誰呢?.
一個是耶利哥王.
是艾城的王.
艾城近伯特利吉.
一個是耶路撒冷王.
一個是希伯倫王.
一個是耶末王.
拉吉王.
伊基倫王.
基瑟王.
底碧王.
基德王.
荷爾瑪王.
阿拉德王.
納娜王.
阿杜蘭王.
瑪基大王.
伯特利王.
哈普亞王.
希弗王.
阿弗王.
拉沙倫王.
瑪瑣王.
夏瑣王.
申倫米倫王.
阿薩王.

$^{441}$哈納王.
米吉多王.
基底斯王.
靠近加密的約連王.
多爾山當的多爾王.
吉甲的戈仁王.
德薩王.
三十二個王.
三十一個王.
總共.
這三十一個王.
數剛才的名字.
數到有三十一個王.
這三十一個王.
他們得到什麼呢.
以色列人第一得到地.
第二呢.
就得到很多人.
第三呢.
就得到很多王.
得到人可以補充一下.
他們就一路.
似乎他們逐點下去.
將男的戰士王.
男人士兵.
等等特地做.
留下比較多女的.
以來做婢女奴隸.
當時婢女不是只做.
煮飯金紙嗎.
不是的.
女的都耕田的.
都可以養牲口的.
什麼都做的.
都辛苦的.
第一得到地.
第二得到人.
第三得到三十一個王.
條命.
將他們都殺死了.

$^{481}$原來以色列人這一群.
烏合之眾.
你想像.
由埃及出來.
全部奴隸.
拿錘就叻.
砸石就叻.
拉麻繩.
拉木就叻.
個個都手瓜起繭.
但其實要打仗.
不是只講大隻.
是要講戰術.
講戰略的.
但神幫了他們鬥.
勝過了那麼多.
最先勝過了.
埃及法老的軍兵.
然後一路勝過.
曠野的王.
和他的軍兵.
然後勝過三十一個.
迦蘭地的王.
嘩.
看到這裡.
讀那些王的名字的時候.
心裡面真的步步跳.
因為很感恩.
很看到上帝的大能.
看聖經人物.
看聖經故事.
是看到上帝的手在背後工作.
上帝的手在天上降下冰雹.
上帝的感動.
令到那些聯軍.
那麼多不同的王.
不同的城出來的士兵.
那些叫戰士.
或者叫勇士.
他們是亂了.

$^{521}$於是大家賄亂不堪.
大家自己殺起來.
就不夠英職的人打.
第三你看到上帝帶約書亞.
是很有謀略.
很有謀略.
譬如夜間出動.
那些士兵還沒準備好.
敵人還沒準備好.
就馬上到.
突然攻擊他們.
艾成輸了一次.
第二次他又怎樣呢.
早一晚就叫那些士兵.
去到艾成的後面.
大家還記不記得了.
埋伏.
然後約書亞帶著另一幫人.
在正面攻擊城.
然後詐輸逃跑.
等艾成的人.
上次追你不夠.
這次繼續追你.
一追出來.
他們馬上轉身走了.
裝輸了.
當他們走到差不多.
引開了敵軍之後.
後面城後埋伏的士兵.
出來攻佔了城.
放火燒.
嚇到那些艾成軍隊.
回頭一看.
我們的城原來被中計了.
我們被燒了佔領了.
然後城裡面那些.
埋伏掠奪了城的士兵.
就出城.
於是令到艾成的勇士.
前後鬥敵.

$^{561}$於是就輸了.
沒有一個可以逃脫.
這麼厲害又有謀略.
又戰勝了.
三十一個王.
確實是上帝奇異的功臣.
讓他們得到地.
得到王.
得到人.
然後得到他們的王.
完全是得勝.
大家看到這裡.
心裡面實在非常感恩.
我自己看完之後.
上帝你真是厲害.
大家看到我用兩節聖經.
來結束今天講的.
創世記第十五章.
第十七十八節.
這一段.
用藍色字來表示是參考.
日落天黑.
不料.
有冒煙的爐.
並燒著的火.
火把從那些.
玉塊中經過.
是什麼呢.
原來神童亞伯蘭.
合立約.
耶和華與.
亞伯蘭立約.
就是這一段.
當那日耶和華與亞伯蘭立約.
創世記第十五章的約.
是講.
我要將地賜給你們.
這樣.
我以賜給你們的後裔.
從埃及地.

$^{601}$到又法拉底河.
大河之地.
都給你們.
他們沒有拿到埃及.
埃及仍然是法老.
但你看到他們二百萬人離開的時候.
是拿了埃及很多的.
財寶,食物.
與他們的金器.
就.
掠奪了埃及.
留在那裡.
都沒有用.
又法拉底河他們去不到.
他們是在約旦河.
西已經停了.
不過如果照這樣看.
其實他們再向西走.
再向西走去到又法拉底河.
即是兩河流域.
都說我已經給你了.
不過他們沒有拿到.
貴德之地實在很多.
從來都沒有什麼記載.
不過這一段.
創世記十五章是講了.
我是將埃及.
與又法拉底大河.
之地都給大家.
給以色列人.
這裡講什麼呢?.
不料有冒煙的爐.
燒著的火把.
當日.
立約人是這樣的.
將一隻牲口斬開一半.
斬開兩邊是對稱的.
擺在地上.
或者擺在.
木架上之類.

$^{641}$一半.
牲口死了流血.
象徵著.
這是一個很認真流血.
大家都承諾要做的藥.
然後大家就拿著.
自己祈禱用的.
爐.
一個圓形的爐.
有幾條鏈這樣吊著.
有些香料在中間.
放一些火在下面燒.
香氣出來.
有味道香味.
每個人都有一個爐.
象徵著祈禱.
每個人如果要立約.
譬如我和你立約.
我們就把牲口斬開.
譬如我兒子要娶你女兒.
又或者.
我和你合作.
我們要分家.
你賣給我之類.
大一點的藥.
把牲口斬開放在兩邊.
然後大家就拿著爐.
手牽著手.
拿著爐.
兩個人一起在兩邊.
兩個斬開的牲口當中.
走過.
象徵著.
上帝在我們中間造證.
在那些肉塊當中經過.
這裡說到肉塊.
就是把牲口斬開.
放在一半.
天黑.
所以天黑就開始用火把.

$^{681}$有冒煙的爐.
並燒著一個火把.
從那些肉塊中經過.
火把就是晚上用的.
冒煙的爐.
就是他們祈禱用的.
這裡有多少呢?.
有一個爐和一個火把.
肉塊在兩邊.
這是什麼意思呢?.
天黑,如果你看上面.
創世記十五章挺好看的.
神使亞伯拉罕沉睡.
亞伯拉罕就睡不著.
睡著了.
什麼都不知道.
然後去到日落天黑.
有冒煙的爐和燒著的火把.
在那些肉塊中經過.
這是象徵什麼呢?.
不是象徵,是表示什麼呢?.
是神自己單獨與亞伯拉罕立約.
我不用你.
你其實都給不了什麼我.
是我應去給你的.
我立這個約.
燒著火把.
我炸著爐.
我當時是人用來做.
神明在上.
你知道我們立約的意思.
神都沒有什麼.
再大過自己.
人用的方法我都用.
神自己單獨與亞伯拉罕立約.
就是我要把地持給你.
就是這一段聖經.
就是埃及自由法拉第學.
有些什麼人在這裡呢?.
有很多人.

$^{721}$有基尼人,基尼洗人.
甲摩利人,哈人.
比利洗人,利索欽人.
阿摩利人.
迦南人,甲迦薩人.
耶布斯人之地.
大家都說了.
耶布斯人之地就是耶路撒冷的地方.
迦南人,哈人.
這些我們經常都聽到.
就是張之外都讀到.
很多地的,很多日.
很多王,就是和這些有關係.
我特別想和大家說的.
神是單獨與亞伯拉罕立約.
是要把這塊地.
給他們.
當然.
在地圖上.
很多粉紅色,是以色列人.
利得之地.
得了地,一分了地中支派之後.
他們就沒有再.
進攻其他地方.
利得之地很多.
就沒有要.
這個是他們.
去到那裡就停步了.
我今天和大家.
分享到這裡,我很想大家.
記得我們的題目.
就是.
從北到南的勝利.
有兩個地方,不知道你記不記得.
就是在這個.
北面.
叫做最北的地方.
去到尼泊爾縣附近.
那個就叫做.
巴力加德.

$^{761}$那個地方開始.
就是去到南面.
在以東西義山那裡.
西義的地方.
的山都是以色列.
得勝的,看回地圖.
你見到他們獲得很大便宜.
其實是不是真的他們.
這麼好打呢?成班做奴隸.
出身,沒有練過.
兵,在曠野.
又未必練兵練得很好.
只是有個勇氣.
有個信心.
就進去了.
你看到過約旦河.
是誰讓他們走.
乾路過去的?上帝.
走七次.
第七天走.
每天走一次,第七天走七次.
耶利哥城島塔是誰做的?.
上帝做的.
有這個謀略去指證.
艾城的人.
成為他們一個的.
標誌,每個人都記得這件事.
是誰?.
艾城和耶利哥城的攻陷.
是成為他們所有.
後面的王都害怕的一個.
標記,很厲害.
接著南邊的王.
北邊的王,越來越多和他打仗.
都被勝利.
大家知道,純粹是.
因為應許.
而來的,所以我今天.
是要讀回.
創世記第十五章.

$^{801}$上帝賜地.
給他們的應許.
雖然上帝賜的地.
是更大的,他們拿的是很少.
這個都提醒了我們.
神是給我們.
得勝的,由北到.
南都得勝的.
是應許我們得勝的.
我們有沒有信奉這個應許呢?.
第二就是.
有很多地神給了我們.
不過我們沒有去拿,夠了.
我們已經好了,開始.
大家就要分地了.
他們就商量好.
將分了地之後.
就拿去業了.
有很多未得之地.
是沒有拿到的.
不過到現在.
都仍然亡國之後.
重新立國,都見到.
神是祝福他們的.
記得,我們的得勝.
完全因為應許而來.
無論你得勝的地位.
有多大,如果上帝沒有幫助你.
你是不會得勝的.
你是沒有這些謀略的.
但是說多謝我們的神.
教我們的指頭打仗.
教我們.
怎樣想.
去計謀.
用手指.
我們可以這樣這樣.
其實這些全部是上帝幫我們的.
因此不要邀功.
不要唱上帝的榮耀.

$^{841}$是忠心信服的.
跟隨他,有信心就好了.
讓我們一起祈禱吧.
天父上帝.
我們今天讀這個經文.
其實很激動,很感恩.
一個.
這樣的以色列民.
做開奴隸.
竟然可以.
出埃及.
過蘆葦紅海.
然後走完曠野.
然後過約旦河.
走乾地.
然後摧毀了.
耶利哥.
和艾城.
令到整個迦南地.
整個王都和百姓都震驚.
然後他們群起.
要圍毆.
攻擊以色列.
不過都被勝過.
是你自己的大流.
無論是你的兵搏.
無論是你的計謀.
無論是你的大力.
無論你在天上被那些敵軍的擾亂.
完全是你的工作.
這些.
其實不是因為以色列的好.
而是因為你.
他們的祖先.
阿伯納.
是要將地給他們.
你就按自己的應許.
將那些地給他們.
將裡面的人給他們.
然後將裡面的王都給他們得勝.

$^{881}$主我們讚美你.
就讀完這聖經.
讓我們更加勇敢地走前面的路程.
知道你是會祝福.
會帶領,會報仇.
我們整個一生人.
就仰望你.
我們走這條曠野.
直到入迦蘭的路都象徵著.
我們在世界上的日子.
走過很多曠野.
荒涼.
淒涼,缺乏.
驚恐的日子.
是你的瑪蘭每一日支持我們.
是你帶領我們勝過.
很多的禍.
而最終我們會入到天家.
走完路程.
讓我們這一世人更加抬頭仰望你.
跟隨你.
主能夠帶領我們.
經過曠野路.
讚美主.
多謝你顯示你自己站在大樓.
祈禱奉靠耶穌基督的名.
阿們.
\newpage



\section{利未記 6:8-13}
\label{sec:032148}
\textbf{常常燒著的火}~~~~周曉沂傳道
\newline
\newline
連結: \href{http://www.church.com.hk/acms/content.asp?site=cdc\&op=show\&type=product\&code=032148\&layout=sermon}{\texttt{032148}} ~~~~ 語音日期: 2026-05-17
\newline
\newline
\hyperref[sec:032147]{\small{< < < PREV SERMON < < <}}
~
\hyperref[sec:index_chronic]{\small{[返順時目]}}
~
\hyperref[sec:032149]{\small{> > > NEXT SERMON > > >}}
\newline
\newline
利未記 6:8-13
\newline
\begin{longtable}{cl}
\hline
\hline
章節 & 經文 (和合本修訂版)\\
\hline
6:8 & \begin{tabularx}{0.7\textwidth}{X} 耶和華吩咐摩西說： \end{tabularx} \\ \\ \relax
6:9 & \begin{tabularx}{0.7\textwidth}{X} 「你要吩咐亞倫和他的子孫說，燔祭的條例是這樣：燔祭要放在壇的底盤上，從晚上到天亮，壇上的火要不斷地燒著。 \end{tabularx} \\ \\ \relax
6:10 & \begin{tabularx}{0.7\textwidth}{X} 祭司要穿上細麻布衣服，又要把細麻布褲子穿在身上，把在壇上燒剩的燔祭灰收起來，放在壇的旁邊。 \end{tabularx} \\ \\ \relax
6:11 & \begin{tabularx}{0.7\textwidth}{X} 然後，他要脫去這衣服，穿上別的衣服，把灰拿到營外潔淨之處。 \end{tabularx} \\ \\ \relax
6:12 & \begin{tabularx}{0.7\textwidth}{X} 壇上的火要不斷地燒著，不可熄滅。每日早晨，祭司要在壇上燒柴，把燔祭擺在壇上，並燒平安祭牲的脂肪。 \end{tabularx} \\ \\ \relax
6:13 & \begin{tabularx}{0.7\textwidth}{X} 壇上的火要不斷地燒著，不可熄滅。」 \end{tabularx} \\ \\
[1ex]
\hline
\hline
\end{longtable}
$^{1}$各位鄧兄姊妹 早安 各位平安.
今天的講題是「常常燒著的火」.
我想起BBQ的話.
大家有沒有什麼方法可以保持BBQ的火是常常燒著的呢?.
有沒有人自認自己是燒烤的專家呢?.
這樣好像很驕傲.
大家平常燒烤的時候有沒有什麼方法呢?.
先問問 加壓 加燙.
還有沒有其他方法呢?.
噗火 這邊聽到有些噗火.
其實我不是一個專家.
通常都是等人幫忙設好爐.
但我上網搜尋.
如果BBQ要透爐的話.
有什麼秘訣呢?.
其實有很多方法.
簡單來說 因為火是有個火三角的東西.
如果大家有學過.
火三角就是要有足夠的燃料.
要有足夠的溫度 要有足夠的氧氣.
所以剛才大家講的方法.
其實都是在應對這些火三角不同的元素.
例如燃料就是加壓.
這邊的噗火就是要保持通風和氧氣.
有不同的方法.
其實都是在講我們要預備一個合適的環境.
去讓火去生存或者更加旺盛.
而踏入五月 教會都很積極地在預備夏令會.
同時不單止我們自己去預備一些事務性的東西.
其實我們都很想弟兄姊妹一同去預備我們一個環境.
我們心靈的環境去預備.
一起去迎接夏令會.
所以如果大家有留意貓子的話.
就應該見到其實貓子都邀請大家去預備自己的心靈.
去花多些時間去禱告.
花些時間去預備我們的.
可能是花些時間.
分別為性地去同時去禱告.
所以不如我們這一刻都一起去祈禱.
「差主 多謝你賜下你很寶貴的話語」.

$^{41}$「神啊 願意我們這一刻」.
「當我們回到這個禮堂裡面」.
「我們將這個時間分別為性給你」.
「並且我們都在預備我們的心靈」.
「願意我們真的能夠在今天去聆聽」.
「上主你自己的話語」.
「所以主 多謝你」.
「求你去開我們的耳 開我們的眼」.
「祈禱系奉主明求 阿們」.
今天的經文就在利未記.
我不知道大家對利未記這個書卷.
其實有甚麼印象.
可能大家一聽的時候就會覺得.
會不會是一個好像很累贅.
很冗長 篇幅很長.
又或者是很多憲制的條例.
或者對於祭司的條例.
於是就會想 祭司即是領袖.
是事奉人員.
好像不是很關我事.
只是想回來參與崇拜.
但相信大家都有聽過.
信徒皆祭司.
正如我們之前也有分享過.
彼得前書第二章也說.
我們是君尊的祭司.
但除了這個原因之外.
第二個原因.
為何這段經文與我們有關.
是因為其實在利未記.
由第一章到第六到七節.
是講述不同的祭.
然後到第六章第八節.
即是我們今天讀的經文.
又再開始講述五祭.
就好像在重複一樣.
但如果我們細心看的話.
其實會看到頭幾章.
特別是對會眾講.
讓這群要獻祭的人.

$^{81}$他們知道自己要怎樣去獻祭.
而第二次 即是第六章第八節.
這裡所謂重複的時候.
其實是對祭司講的.
於是你就會想.
是呀 那不就是只對祭司講.
但是有學者指出.
其實他們研究了發現.
其實利未記現在的場景.
其實是神當著一群會眾對祭司講.
所以不是一個閉門造車的狀況.
不是叫祭司來 我們只是這樣談.
而是當著所有會眾講.
而為什麼要這樣當著所有會眾講.
不是一個閉門造車的狀況.
就是因為這句說話是很重要.
並且很想大家都知道神的說話.
所以相對於其實是說.
祭司條例或者獻祭條例.
很一字一句.
其實利未記更想讓以色列人知道.
他們自己怎樣可以與神建立關係.
不只是自己.
不只是自己在信仰,行為.
或者個人衛生等各方面.
而是作為一個國民.
即是一個群體 也是以色列人.
一整個群體的分別為勝.
所以不只是自己.
更加是一整個群體的分別為勝.
所以這一刻可以跟你旁邊的人說.
今天這段經文是與我們有關的.
是的 是與我們有關的.
所以讓我們再一次讀出這段經文.
1,2,3起.
「通宵一刻 祭壇上的光 與水人的肖像」.
「祭司要穿上祭壇的外套」.
「又把祭壇的褲子穿在身上」.
「加要把祭壇上的藍色紫色的灰色的袋子」.
「放在祭壇旁邊」.

$^{121}$「然後針對這些衣服」.
「穿上這些衣服」.
「把灰拿出營外見識的地方」.
「祭壇上的火必須常常燒著」.
「不可適應 每天早晨」.
「祭司要在上面燒著」.
「把祭壇列在壇上」.
「又在上面燻燒地方祭祀之光」.
「火必須不斷在祭壇上燒著」.
「不可適應」.
剛才我們也提及過第六章第八節.
是第二次再講繁祭.
如果比對第一次講.
即是第一章 如果大家有印象的話.
或者我們剛才就這樣讀.
可能你也會發現.
為何講一個祭的條例.
好像完全沒有提及過動物.
大家有沒有發現.
平時講任何獻祭好像也會提及羊和象.
但剛才那裡一點動物相關的東西也沒有.
最多最多都是脂肪.
第十二節的這裡.
為何會這樣呢.
因為這一段剛才也提及過.
是一個祭司的實務.
讓祭司知道他們可以怎樣.
好好地去預備一個環境.
知道他們每一天的工作.
他們預備的功夫要做甚麼.
所以同樣地.
願意今天的訊息.
也是在幫助我們去預備我們生命的祭壇.
我們對神的學務是不惜滅的.
剛才在讀這段經文的時候.
不知道大家有沒有發現.
或者觀察到一些常常重複的地方呢.
差不多每一次都會問大家.
有甚麼重複的地方呢.
這一段經文.

$^{161}$其實這段經文也挺明顯的.
也是今天的講題.
就是燒著.
無論它是說一直燒著,常常燒著.
還是必須不斷燒著.
其實都是在說燒著.
而第一個去形容.
火一直燒著的狀態.
一個時態就是要通宵達旦.
通宵達旦在第九節是這麼說的.
它就是說要通宵達旦.
而另一個中文譯本.
就是說由晚上到天亮.
那麼通宵達旦就是由晚上一直.
可能黑眼圈都出現了.
然後到早上天亮.
所以是由晚上到天亮.
這個火都是一直燒著的.
這一段經文其實沒有很明確的說.
怎樣可以令到這個火.
通宵達旦地燒.
但是它告訴我們的是.
要令到這個火保持著通宵達旦地燒.
這是我們要去設法確保的事.
什麼意思呢?.
其實它是在暗示.
原來火是有機會會熄滅的.
原來火是有機會會熄滅.
大家有沒有想過.
原來火是有機會會熄滅的.
當我們不去聚精會神.
或者想一些方法.
或者不去檢視火是什麼狀態的時候.
原來火是有機會會熄滅的.
當時的以色列人.
只有火不熄滅的時候.
他們才可以繼續去獻祭.
如果火萬一熄滅了的話.
是不是就可以隨便拿一個.
可能打火機或者當時的蠟燭.

$^{201}$火柴.
就可以點燃呢?.
當然不是.
因為在《利未記》有一個概念.
就是「聖火」和「煩火」.
「聖火」和「煩火」.
在第九章之後的一些段落有提及.
提及到有火從耶和華面前出來.
燒盡了繁祭和自由.
緊接著到第十章的時候.
提及到亞倫的兒子拿達和阿比胡.
拿著自己在耶和華面前獻上了煩火.
於是就有火從耶和華面前燒滅了.
所以為什麼要這麼竭力去保持著火呢?.
因為其實由神而來的火.
我們是沒有辦法用自己的方法去點燃的.
而當火沒有了的話.
對於當時的以色列人來說.
就是他們不能夠去獻祭了.
於是我們要問的是.
對於我們而言.
如果火沒有了的時候.
我們會怎樣呢?.
大家可以有少許時間去想想.
如果我們生命對於神的火沒有了的時候.
會是怎樣的呢?.
其實我們來了神就不能夠作神魔了.
而且這個要保持著火.
其實重點不是一個宗教的儀式.
好像在《聖母意記》上第十五章這裡提及到.
「聽命聖於賢祭」.
其實比起一些宗教的儀式.
更加重要的是我們心中的馴服.
我們有沒有遵行神的話語呢?.
而不只是一個形式上的敬虔.
所以去保持著這個火.
說的是我們有沒有去聽從神的話.
我們有沒有在心中的馴服呢?.
說到火要不熄滅的時候.
我就想起奧運的傳火具.

$^{241}$奧運這件事其實是源於希臘的神話.
也是他們的宗教.
所以大家知不知道.
或者有沒有聽過.
如果他們要一開始收集這個所謂的聖火.
我們姑且都叫聖火.
他們要收集這個聖火是有什麼方法的呢?.
大家有沒有聽過?.
咦 聽到 是的.
就是要用太陽.
因為正如剛才所說.
他們都相信是聖火.
所以不能用打火機或者什麼點著.
所以他們首先要有太陽.
然後就一直射著這個凹面鏡.
然後就去收集這個火.
然後有了第一把火之後.
就一直傳下去.
大家都會在新聞可能見過.
有搭車 有搭船 有搭飛機.
要搭這些不同的交通工具的時候.
他們都需要有一些很特別的器皿去裝著.
以至於這個火是不熄滅的.
一直傳 傳 傳 傳.
傳到當屆那個國家的主場館的所謂聖火台.
當然我們想一想.
在奧運這麼漫長的歷史裡面.
有沒有發生過一個聖火熄滅的意外呢?.
也不難想像是有的.
這麼久的歷史 是不是?.
在Santa紀錄發現.
在1976年的加拿大的Montreux.
即是蒙特利爾的奧運.
他們傳到開幕之後.
就有暴雨淋熄了那個火.
因為這個聖火台是很重要的.
所以一直都有看守著的.
看守著的人一看到火熄滅了.
就馬上從褲袋裡拿出打火機.
就點燃了它.

$^{281}$於是 因為這個火不是來自太陽的.
不是神聖的起源.
所以奧委會就要把這個不純正的火熄滅了.
然後就要用一個叫做.
官方備用的火種去點燃.
是由希臘他們帶來的.
事實上 其實歷史上.
都不止這一次的熄滅的意外.
而為了應對這些意外.
他們就有一個方法.
就是當他在.
剛才說到太陽射在鏡子上的時候.
他們就不止收集一個火種.
而是會收集多過一團的火.
就會有一些叫做官方的備用燈.
火種燈.
就一直在整個全聖火的過程都會跟著.
回到今天的經文.
昔日的以色列人.
他們要 剛才說到要通宵達旦.
去保持著火去燒.
有什麼方法呢.
是不是要派一個夜班的漢奸也是這樣呢.
隨時在褲子上掩著打火機呢.
還是要做什麼方法呢.
一早計算好什麼.
經文似乎沒有很明確地說清楚.
怎樣 就是how.
可以不熄.
他沒有說清楚 沒有明確地指出.
但是有學者就指出.
當第九節 剛剛讀的第九節.
他說到要放在柴上.
雖然相對於第十二節.
他說直接是說在上面燒柴.
十二節是比較明確的.
第九節是好像暗示多一點.
但其實都已經在說.
要確保有柴.
所以我們要從經文學習的就是.

$^{321}$我們要去預備.
無論如何 我們要知道的是我們可以預備.
我們都要警醒.
就是我們要留意著.
究竟這個火的狀態是怎樣呢.
我們要檢視自己那個火的狀態是怎樣.
也都要為這個祭壇設置好一個環境.
而另一個經文說到.
可以令火繼續燒著.
然後不熄滅的方法.
就是處理一些灰.
其實這個方法都不難去理解.
大家有沒有聽過青少年會去宿一宵營會.
可能家長也有聽過.
其實這個宿一宵營會.
都很久歷史了.
我都數不到.
我想至少都有十幾年了.
每年青少年在農曆年假的時候.
就會去宿一宵營會.
很明顯宿一宵這個字是食字的.
是什麼意思呢 就是我們除了睡一晚之外.
我們會燒烤的.
無論多熱.
雖然農曆年有時候都會熱的.
雖然是冬天.
但是都會熱.
我們都要燒烤的.
是一個傳統.
我們每年燒烤都是第一晚的時候.
到最後大家吃飽了.
要關火了.
我們要把牆交給別人的時候.
我想問一下大家.
有什麼方法可以關掉這個炭火.
燒烤的火.
大家平時會怎樣關的.
淋水.
還有沒有其他方法.
用沙沙在哪裡.

$^{361}$燒完它我聽到.
其實我經常都要幫忙.
善後的時候留到最後.
就會說.
怎樣可以關掉這個火.
其實我最常聽到的一句話就是.
可以了 就這樣放在這裡.
它會關掉的.
可以了 就這樣放在這裡.
它會關掉的.
然後平時我都沒有怎麼想的.
因為我不是專家.
我剛才說了.
所以就說 好吧.
你們說就這樣放在這裡就關掉.
那就就這樣放在這裡.
但是今天我就在這裡.
不是這次預備的時候.
我就在這裡想.
其實為什麼.
就這樣放在這裡就會關掉的呢.
到底為什麼就這樣放在這裡就會關掉的呢.
於是.
大家好像很厲害.
大家是不是經常勾引去燒烤.
是的 就是找了上網查.
有兩個原因.
其實都是跟灰有關的.
第一就是.
她說到.
因為火 剛才說要有氧氣.
要通風.
所以如果一直燒.
很多灰的時候.
就會塞住爐的底部.
所以空氣就不流通了.
所以火就會窒息.
那就關掉了.
這是第一個原因.
而第二個原因就是.

$^{401}$當灰也是.
慢慢燒 越來越多灰.
就會包住了整個炭.
所以就會阻隔了.
炭和空氣的接觸.
就會令到火勢一直慢慢減弱.
然後就關掉了.
所以簡單來說.
這兩個原因都是關於灰的.
就是因為有灰.
阻隔了空氣的接觸.
所以剛才當經文讀到.
要祭司去清走灰.
以至於祭壇上的火.
可以常常燒著.
不會熄滅.
原來是一個很科學化的方法.
非常合邏輯.
非常合理.
而我們剛才一開始也說到.
要通宵達旦讓火燒著.
都應該想像到.
其實真的會有很多灰.
當要在早上清的時候.
如果我們繼續很有邏輯.
這樣去想.
我們要將這段經文.
套到自己生命的實際處境的時候.
我們就會想.
怎樣去應用這段經文呢?.
很自然就會想.
我們的生命的火.
就是一個祭壇.
讓火繼續燒的時候.
就要清走那些灰.
清走那些灰即是怎樣呢?.
即是我們要去處理.
可能我們生命的一些垃圾.
即是要去倒走一些.
可能生命中不屬於神的東西.

$^{441}$一些會阻礙到神的關係的東西.
都要拿走.
這是一個很自然.
很合邏輯的想法.
但是如果我們再看一段經文的話.
可能有不同的看法.
所以不如我們再讀.
剛才說除灰那部分的那兩節經文.
一二三起.
.
是的.
剛才我們讀的那段經文.
第十一節最後那一句是甚麼.
大家記得嗎?.
把灰拿出.
去哪裡?.
.
大家還記不記得這一卷.
這麼講求.
聖潔不聖潔.
乾淨不乾淨的經卷.
它特別強調要把灰.
拿出營外潔淨的地方.
其實是很值得我們思考和留意的.
拿出去潔淨的地方.
上帝不會叫我們把垃圾拿去乾淨的地方.
特別是剛才讀的那段經文裡面.
很清晰地說到要穿一些特別的衣服去收灰.
然後當它放在旁邊.
然後出去的時候要脫掉那些衣服.
穿回普通的衣服.
證明整件事是很清晰.
很具體.
很清楚地說整個流程.
是一件很認真去處理的事情.
所以如果叫那些灰做生命的垃圾.
似乎好像有點奇怪.
有點卡住.
不太講得通.
於是我們要問的就是.

$^{481}$為什麼這些灰要放在一個乾淨的地方呢?.
又或者我們要問的就是.
其實這些灰本身是什麼東西呢?.
如果看英文的譯本的話.
其實它有一個名字叫Fatty Ashes.
就是一些肥的灰.
那是什麼呢?.
其實就是燒盡了.
因為凡是燒盡的.
其實燒之前是什麼呢?.
就是一些祭物.
可能是柴.
所以這些灰是有燒盡了.
並且是被神所閱立的意思.
所以這些灰就不是我們平時燒的食物.
很忌諱的.
如果一粒魚蛋或者一塊麵包掉下去.
我們就不懂了.
就不理它了.
不是這樣的.
而是有神所閱立的意思.
是神指定了那些動物用來獻祭.
當它燒盡了.
就是神所閱立.
是屬於神的.
所以這些灰就要放在一個乾淨的地方.
但是為了讓火繼續燒的時候.
祭司都要在每個早上.
在擺柴之前.
就要將這些灰拿走.
當然我們生命如果有垃圾的時候.
一定是要清走它.
一些神不喜歡的地方.
阻礙了我們與神的關係的東西.
我們是要拿走的.
不過今天的經文.
好像也在提醒我們.
原來有一些我們以為是好的事.
可能是我們過往的經歷.
我們昔日的事奉.

$^{521}$過去的成就.
原來今天也會成為了灰.
阻礙了一些火繼續燃燒.
其實簡單來說.
就是不要吃老本.
不要覺得我仗著昨天的火燒得很旺盛.
今天就覺得不用做任何事去建立與神的關係.
其實每一天都是新的.
所以我們不可以仗著昨天.
我做過怎樣怎樣.
然後今天就覺得今天不用了.
我們需要持續地去更新我們的生命.
每一天都向神去支取力量.
每一天都與神去建立關係.
最後兩節都讓我們再一同讀出.
一二三起.
祭壇上的火.
應常常燒著.
不可熄滅.
每天早晨.
祭司要在上面招柴.
把凡祭牌列在壇上.
又在上面焚燒.
平安祭的之風.
火必須不斷在祭壇上燒著.
不可熄滅.
這兩節有兩個字詞.
想和大家再看得深入一點.
第一個就是第十二節所說的.
每天的早晨.
其實在原文的希伯來文.
它是早晨早晨地說.
兩個名詞重複.
其實是一個很常見的修辭手法.
有兩個功用.
第一是說很持續.
並且有規律性.
每天不間斷.
持之以恆去做.
而第二個意思或修辭的作用.

$^{561}$就是在說這個行為很重要.
所以剛才說到每天早晨.
即是早晨早晨.
要在這裡燒柴.
其實是在說燒柴這個行為十分重要.
並且要不間斷.
持之以恆地去做.
而聖經其實也有很多地方.
會提到早晨早晨這個修辭手法.
而另一個有用到這個修辭手法的地方.
就是在《出埃及記》.
提到嗎哪的時候.
提到嗎哪大家應該都熟悉.
就是要說每個早晨.
即是早晨早晨去按照個人的食量去收取.
即是在說每一天要去拿.
而且不能夠是今天餵到明天去存取.
不能夠的.
所以回到今天的經文.
其實在壇上燒柴的這個行動.
其實都是要每天去做的.
不能夠今天為了明天去燒柴.
又或者你覺得.
昨天好像燒多了很多.
那今天就不用了.
是不可以的.
是需要每一天都去建立和神的關係.
而第二個字去看的就是「不斷」.
其實剛才說了.
這段經文都很重複「火要一直燒著」.
即是有重複這個字.
但是只有第十三節這裡.
它是有一個額外的副詞.
就是要強調那個時間性「不斷」.
而這個「不斷」.
不是只是說很頻繁的意思.
其實它是一個永恆的意思.
什麼意思呢?.
即是它不是只是說.
我這一個星期,兩個星期才不斷.

$^{601}$而是說永遠一直都是這樣去不斷.
所以它說到要一直燒著的時候.
是永遠一直不斷地去燒著.
在《初一及記》第29章第42節就說到.
「者要在耶和華面前.
在會幕的門口.
作為你們世世代代經常獻的梵祭.
我要在那裡與你們相會.
和你說話」.
「梵祭」是一個完全的焚燒.
而「梵祭」其實它的意思是「向上升」.
所以「梵祭」就是一個向上燒去獻情給神.
一個輕香的祭去獻情給神.
而一個世世代代的梵祭.
它在每天的早晨和黃昏獻上.
衛爵在這裡說的就是.
衛爵去預備一個地方.
給神和人去相會.
有一個地方給神和人去說話.
同時間他也在提醒以色列的群體.
每一天一張開眼.
到閉上眼的時候.
都是在神的恩典之下.
所以這個世世代代.
所說的不只是某一代的人.
而是一代傳一代.
而要一代傳一代.
就必然需要上一輩.
或者父母的新教,延教.
每一天我們在生活上的實踐.
其實也會影響下一代.
我們怎樣在生活上實踐神的話.
其實是會影響身邊的人.
怎樣才可以做到生活上的實踐呢?.
其實回歸最基本.
就是我們要抽時間和神建立關係.
將一天的時間分別繪聖給神.
但是可能有弟兄姊妹會說.
其實我也有的.
我真的有花時間.

$^{641}$我真的有畫了一個界線.
或者一個時間就是每天定時定候去靈修.
但是我發現自己有個問題.
就是我好像每天靈修的那些.
可能神跟我說了的事或者經文.
我好像帶不到去日常生活.
又或者當我完了靈修一天.
很多工作或者要照顧家裡.
或者各樣忙碌的時候.
我就忘記了所有的經文.
或者忘記了神對我的提醒.
在這裡想和大家分享一個靈修的方法.
就是意識醒察的練習.
其實是一個靈修的方式.
它可以有不同的方法.
或者很細微的不同.
但是其實目的都是一樣的.
就是去察覺神在我們當一天的生活裡面.
祂是怎樣去臨在.
並且去辨識自己.
我內心其實是什麼在推動我去做一些事呢?.
究竟是神在推動我去做一些事.
還是我自己內心的慾望在推動.
是幫助我們去加深對神的連結.
我在神學院學的方法就是.
一天有四個時間.
每一個時間都很短.
就是早上刷牙的時候.
一直在刷牙.
然後你吃完午餐.
想站起來走走.
在消化的時候.
還有晚餐吃飯之前.
在預備,可能剛剛回到家裡.
最後就是晚上睡覺之前.
頭三個時段.
其實每一次都是一個很簡單的祈禱.
就是你去祈禱和神說.
「神啊,求你繼續幫助我去面對自己的生活」.
可能你會想起早上靈修的一些提醒.

$^{681}$或者你也可以去祈禱.
去檢視究竟我剛剛那幾個小時的生活.
有沒有得罪神.
有沒有得罪人呢.
而另外一個比較長一點的時段.
就是晚上睡覺之前.
就是這裡看到的一些問題.
有幾個問題其實是在幫助自己去檢視自己當日的全日.
首先去想一下.
今天有沒有什麼事情.
你覺得是很深感受的.
而去想這件事的時候.
不只是想著「哦,是的」.
「有人打我尖」就算了.
而是想「當有人打我尖的時候」.
「那個喜恩,我的感受是怎樣呢」.
「我的反應是怎樣呢」.
「究竟有人打我尖的時候」.
「我是會跟他有衝突嗎」.
「去問他」.
「還是我一下子就發脾氣」.
「還是我算了,我不想跟別人吵」.
「我就躲起來,我就看不到」.
「我就走去別的地方」.
就是有不同的喜恩,或者有不同的行為.
然後結果就是究竟是怎樣發生這件事.
究竟是可能我躲起來,我不想跟別人吵.
於是我可能就逃避.
但是我心裡有憂鬱.
所以憂鬱了一整天之後.
我就把這些情緒可能發洩在我的家人身上.
可以有不同的結果.
就是很仔細地去想.
究竟一件很深感受的事情.
究竟它的喜恩,行為和結果是怎樣.
第二就是發生這個事情的時候.
很仔細地去想所有的感受.
不是一個很籠統的感受.
即是不只是說開心,不開心.
而是很仔細地.

$^{721}$例如我覺得很委屈.
我覺得很尷尬.
我覺得很想逃避.
即是很仔細地去命名一些感受.
但更加重要的是.
其實不只是去想這些感受.
而是當想到這些感受的時候.
去想究竟.
這些感受是更加幫我們去緊了神.
還是遠離了神呢?.
即是當我們很想發脾氣的時候.
好像聽到神的聲音叫你停一停的時候.
其實是幫助你去緊神.
但有時可能有些感受是幫我們.
即是拉離開了神.
這是一個很重要的一步.
去感受或去想究竟.
這些事件是幫你緊了神.
還是遠離了神呢?.
而最後就是去想.
有沒有什麼是你想跟神說的呢?.
這是一天的流程.
你去檢視.
但其實更加重要或對我來說.
覺得更加有用的.
不是一天的檢視.
而是一個做法是持續.
可能一段長時間.
兩三個星期.
然後當你兩三個星期之後.
你會再拿每一天的紀錄出來看.
你去比對一下.
其實有沒有一些感受.
是我好像每一天都發生的.
而不是發生同樣的事.
即是不是都是被人打劫.
可能有不同的事件.
但你都是有一些感受.
而同時間也要留意.
究竟這兩三個星期.

$^{761}$當你去有這些不同的感受的時候.
究竟是遠了神還是近了神.
對於我來說.
我是一個很感性的人.
其實很多時候.
我會有很多很細微的感受.
出來之後.
我都不知道發生什麼事.
可能我要很難才意識到.
自己在發生什麼事.
但當我用這些方法.
去檢視自己每天有什麼感受.
甚至是感受在哪裡來的時候.
我發現.
即是之前的時候.
就發現經過了兩個星期.
原來真的有一些感受.
是重複地出現的.
而事件是真的不同的.
但激發我的點都是一樣的.
更加重要的是.
當我再鑽深一點的時候.
我就發現.
當這些激發我有這些感受的點.
原來一直都卡住我和神的關係.
原來卡的位置就是這裡.
所以我希望這個方法.
都可以幫到大家.
去檢視自己.
也是在建立和神的關係.
去梳理一些究竟自己在生命上.
卡住和神的關係是什麼事情.
或者那點是什麼呢.
是的.
最後剛才我們也說到.
其實利未記就是幫助我們.
自己和神建立關係之餘.
其實也是一個整個群體.
去獻情自己的生命給上帝.
所以今天我們要問的就是.

$^{801}$我可以怎樣獻上這個祭呢?.
到底這個祭是我們有沒有恆常的獻上呢?.
到底這個祭是不是每天都在生一個活祭.
而不是吃腦本的?.
到底我們有沒有去清理一些生命的灰燼呢?.
之前我們教務或童工也有討論過.
如果夏令會的主題有英文翻譯的話.
究竟會是用ING.
就是Igniting還是Ignited.
其實是在說動詞和名詞的分別.
我們最後是沒有定案的.
因為暫時也不需要用到英文名.
很隨便地聊的時候.
但是其實不是要討論英文的文法.
不過這一刻想問大家的就是.
如果在你生命裡面.
你的火的狀況是怎樣呢?.
到底是一個ING的形式.
就是進行中.
是還在燒的一個狀況.
還是一個過去式ED的狀況.
就是曾經燒著過.
但是現在好像沒有再去檢視.
沒有再去理會.
所以願意今天的訊息.
就是提醒我們要常常去被燒著.
使得我們的生命.
真的是一個活祭去獻情給神.
是一個我們每一天.
恆常獻情給神的祭.
\newpage



\section{撒迦利亞書 5:1-4}
\label{sec:032149}
\textbf{審判全地的神}~~~~黃梅玲牧師
\newline
\newline
連結: \href{http://www.church.com.hk/acms/content.asp?site=cdc\&op=show\&type=product\&code=032149\&layout=sermon}{\texttt{032149}} ~~~~ 語音日期: 2026-05-17
\newline
\newline
\hyperref[sec:032148]{\small{< < < PREV SERMON < < <}}
~
\hyperref[sec:index_chronic]{\small{[返順時目]}}
~
\hyperref[sec:032150]{\small{> > > NEXT SERMON > > >}}
\newline
\newline
撒迦利亞書 5:1-4
\newline
\begin{longtable}{cl}
\hline
\hline
章節 & 經文 (和合本修訂版)\\
\hline
5:1 & \begin{tabularx}{0.7\textwidth}{X} 我又舉目觀看，看哪，有一飛行的書卷。 \end{tabularx} \\ \\ \relax
5:2 & \begin{tabularx}{0.7\textwidth}{X} 他問我：「你看見甚麼？」我回答：「我看見一飛行的書卷，長二十肘，寬十肘。」 \end{tabularx} \\ \\ \relax
5:3 & \begin{tabularx}{0.7\textwidth}{X} 他對我說：「這就是向全地面發出的詛咒。凡偷竊的必按書卷這面的話除滅，凡起假誓的必按書卷那面的話除滅。 \end{tabularx} \\ \\ \relax
5:4 & \begin{tabularx}{0.7\textwidth}{X} 萬軍之耶和華說：我要把這書卷送出去，進入偷竊者的家和指著我名起假誓者的家，停留在他家裡，連房屋帶木頭和石頭都毀滅了。」 \end{tabularx} \\ \\
[1ex]
\hline
\hline
\end{longtable}
$^{1}$各位弟兄姊妹平安.
今天我們會繼續看撒加拉尼亞書裡面的異象.
不經不覺已經去到第六個了.
這個異象是關乎甚麼呢.
我以審判全地的神作為這四節經文的題目.
因為這個異象的重點就是審判.
甚麼是審判呢.
今天就讓我們透過第六個異象.
來思考審判這個重要又嚴肅的課題.
在進入經文之先.
我們一起同心圖圖告.
尋求阿們.
說到審判的時候.
我想我們都會想起懲罰.
不知道大家小時候有沒有做錯事呢.
一定有的.
有沒有試過因為做錯事而受罰呢.
我想很多人都會試過.
我記得我讀小學的時候.
有一次是測驗.
我忘記是甚麼科的測驗.
我發現班長竟然出貓.
班長都出貓 有沒有搞錯.
於是我就做了一個很大膽的行動.
就是我舉手向老師告發.
那個班長出貓.
結果這個班長就被老師罰.
大家知不知道罰甚麼呢.
就是不能做班長.
結果誰做班長呢.
就是我.
不過我舉手的時候沒有想過這個結局.
但當我做了班長之後.
雖然好像很威風.
夾了班長的牌.
但有一次我遇到一位老師.
很高要求 很嚴厲.
他要求我這個班長.
每一天小息的時候.
都要去教員室拿一疊作業.

$^{41}$放在課室裡面.
他就預備上課的時候可以派.
但我就很貪玩.
有一次在小息的時候.
我就想著玩 忘記了.
這個嚴厲的老師就馬上罰我.
不是罰我不能做班長.
而是罰我在教員室門口企.
夾著班長的牌罰企都很醜怪.
那個老師都看著.
這個班長為甚麼要給老師罰企呢.
於是我就拔了班長的牌.
藏在裙袋裡.
因為那一刻真的很想找洞鑽.
這次是一個很深刻的教訓.
我明白到即使我被老師揀選成為班長.
即使我有這個身份.
但如果我達不到老師的要求.
我都要承擔後果受罰.
因為我學了一個很重要的功課.
犯錯受罰是理所當然的.
特別是帶著被揀選的身份.
也是沒有例外.
因為被揀選並不代表有特權.
被揀選的神的指紋又是怎樣呢.
事實上神的指紋都是一樣的.
以色列人作為神所揀選的指紋.
他的身份極其的尊貴.
神更親自來跟他們立約.
在《申命記》十一章二十六至二十八節.
裡面提到神的約.
這也是《申命記》神學裡的重點.
就是講到神的祝福和咒詛.
根據神和以色列所納的約.
以色列如果他被揀選作為神的指紋的話.
他就必須要遵守神所定的誡命律例典章.
他們才會得到神的祝福.
相反如果神的指紋違背神的約.
違背了神的誡命神的道的話.
他們就會被神咒詛,要接受懲罰.

$^{81}$在這裡所講神的誡命是什麼呢.
如果按照《申命記》的脈絡來看的話.
神的誡命其實就是指這三個部分.
第一個部分是五至十一章裡的十誡.
這個十誡是摩西在這裡再次向神的指紋重申.
因為在倡埃及的時候.
神在西乃山頒佈了這個十誡.
第二個部分就是十二至二十六章裡.
講到的《申命記》的法典.
在當中看到的律例典章包含了很多的部分.
包括了宗教的義務.
包括了作為宗教領袖他們的治理.
也包括了在社會裡要行公義的事.
以及每一個神的指紋.
他的行為的規範.
當中涉及婚姻,財富,戰爭,公共衛生等等的法律.
而最後的部分就是二十七至三十章裡提到的.
祝福和咒詛.
所以當我們看到《申命記》的時候.
我們會很清楚的知道.
清楚明白為什麼北角以色列.
為什麼南角猶大作為神的指紋.
竟然會落得亡國被擄的下場.
因為這個就是他們違背神的約的懲罰.
是咒詛.
即使北角以色列人.
南角猶大人他們都是神所揀選的指紋.
但是當他們面對神的審判的時候.
他們都必須要降服在神的審判台前.
因為犯錯受罰是理所當然的.
尤其是帶著這個被揀選的身份.
就是有一個責任去履行這個約裡所定的規矩.
因為被揀選並不代表有特權.
所以當我們讀到今天的經文.
撒迦利亞書裡的第六個意象的時候.
如果我們說這個意象是關於神的審判.
我們就必須要回到以上所提到.
神的約和神的誡命的基礎上.
來理解神的審判.
這裡所提到神的審判對象.

$^{121}$並不是神指紋的敵人.
這裡所提到神的審判對象.
其實是神所揀選的指紋.
是神的家.
在彼得前書裡有一句經文.
說到審判要從神的家起手.
從神的家開始.
以前我讀到這節經文的時候.
我很不明白.
為什麼審判要從神的家開始呢?.
撒迦利亞書為我們提供了答案.
在我們進入第六個意象之先.
我想和大家一起重溫前面五個意象.
第一個意象就是有一個騎紅馬的人.
他在翻石榴的林中.
然後有紅色黃色白色的馬匹來跟隨.
這個意象代表著神差派天使.
來偵察這個世界裡所發生的事情.
也代表著神對神指紋的關心.
神不是冷漠的.
神是愛這些指紋.
所以祂會關心他們.
也在這個意象裡.
去應許會重建聖殿.
接著第二個意象就是四國和四將.
這四個國代表著攻打神指紋的敵人.
而這四位將人就是神差派來攻擊這些敵人.
也代表著神不會袖手旁觀.
祂一定會幫助祂所愛的指紋.
第三個意象是有一個人拿著準繩.
要去測度耶路撒冷城的長寬.
這個意象代表著耶路撒冷的城牆是無界限的.
是一個無牆之城.
而神就會親自成為這個城的火牆.
來保護神的指紋.
第四個意象是大祭司約書亞身穿污穢的衣服.
而神就親自為他來換上一套華美的禮服.
象徵著神會親自來赦免神指紋的罪.
神會親自恢復神指紋作祭司的職份.
也是預表尼賽亞的來臨.

$^{161}$第五個意象就是金燈台的意象.
有金燈台就有聖殿.
所以金燈台就是預告聖殿必被重建.
換言之我們發現前面這五個的意象.
其實全部都是對神指紋所提出的安慰的訊息.
我們可以說五個的意象都是對神指紋的祝福.
其實這一系列的意象就停在金燈台的意象.
豈不是更美好嗎?.
因為大家都喜歡Happy Ending.
會很開心的.
金燈台 聖殿可以被重建.
停在這裡就可以了.
但是這一系列的意象 神並沒有停在Happy Ending裡面.
神卻是繼續次下其餘三個的意象.
而下一個意象也就是第六個意象.
並不是好像之前那五個意象.
是會為人帶來歡喜快樂的情景.
反而是進入了一個很嚴肅的課題.
就是審判.
什麼是審判?.
在現實世界裡.
審判的場景就是一個法庭.
當中有法官和陪審團來做裁決者.
而原告或者檢察官就會和被告形成一個對立的關係.
當中就是由原告向被告提出指控.
而在過程裡面.
法官和陪審員就要去權衡雙方所提出的證據.
作出辯論商討.
然後對被告的行為的後果和法律責任進行最終的確認和判決.
而這一切判決的基準就是建基於人所制定的法律.
在信仰世界裡.
審判的場景不是這個法庭.
而是在神寶座前的審判台.
作出審判的是那位公義又全知的神.
祂所審判的不單只是人的行為.
而是包括人的內心.
因為神是察看人心的主.
神會按著祂的標準來作出審判.
神的律法 神的道 神的話.
就是這個審判裡面的基準.

$^{201}$事實上審判是一個非常重要的概念.
特別是神的審判.
試問 如果沒有神最後的審判.
沒有神在終末最後的審判的話.
今天我們在這個世界上做任何事情.
其實都不會有任何的後果.
如果這樣說的話.
個個都可以任意妄為.
個個都可以不計後果來做最差.
為自己帶來最大利益的事.
如果這個世界沒有審判的話.
後果是不堪設想.
這個世界會進入一個完全崩壞的境況當中.
所以審判是一個既重要又嚴肅的課題.
讓我們看看第六個意象是怎樣描述審判的訊息.
第六個意象其實都很簡單.
在先知的眼前出現的就是一個飛行的書卷.
這個飛行的書卷 先知從來都沒有見過.
今天我們有飛機我們坐過飛機.
我們都可以理解.
但當時的先知的時代沒有飛機.
只可以見到有小鳥在天上飛.
為何會有個書卷在天上飛.
而這個書卷也帶著好幾個特點.
第一個特點當然就是它會飛.
這個飛行的意思其實是指進行中的意思.
或者我們可以理解為現在進行式.
就好像在手術室做手術的時候.
在門外會有個燈寫著手術進行中.
又或者在工程裡面.
有個牌寫著前面工程進行中.
這個飛行的書卷會飛的意思就是.
這個訊息是進行中.
在這個書卷上所寫的內容是進行中的事.
而這個書卷上的訊息.
當然就是從神而來的訊息.
飛行的書卷代表這個訊息是進行中的.
是一個活潑的訊息.
是一個有功效的訊息.
是一個迅速的訊息.

$^{241}$正在飛.
不是慢慢的.
是很快的飛行.
是帶有速度.
是將神的訊息有效地傳播的一個意思.
而翼象的第二個特點就是它的尺寸.
這是一個很特別的尺寸.
長20肘 寬10肘.
如果用我們現在的理解.
大概是30尺乘15尺.
又或者是10米乘5米.
大家想像到.
好像我們在街上看到巴士上的廣告一樣.
又或者我們在地鐵上看到的燈箱.
這麼大.
或者是更大.
又或者我們在電視上看過.
或者在現場看過.
有些飛機要宣傳.
做一些宣傳的時候.
它會掛一個橫額.
飛在天空上就好像這樣.
如果你看到有一架飛機這樣拉著橫額.
你都會怎樣.
看看那個橫額上面有什麼內容.
宣傳著什麼.
這個意思就是這樣.
為什麼尺寸這麼重要.
是代表著這個書卷是打開的.
而且尺寸不是小的 是大的.
打開而且是大的一個書卷.
是要被人看見.
要得到人的注意.
要請每一個人來留意.
這個書卷上面所寫的內容.
所帶有的訊息.
這個就是飛行的書卷.
為什麼是這麼大.
打開.
而且是在飛的一個意思.

$^{281}$這個書卷就好像一個廣告牌一樣.
很矚目.
要引人注意.
讓人清晰地看到那個內容.
而且還很特別的.
這個尺寸其實是跟.
聖殿的門廊的尺寸是一模一樣的.
在《列王記》上的六章二至三節裡.
記載所羅門所建的這個聖殿.
是長二十爪 闊十爪.
跟剛才這個飛行的書卷的尺寸.
是一模一樣的.
這個門廊是在哪裡呢.
其實就是一個空間.
是要進到聖所和致聖所之前.
每一個要進入的時候.
都必定會經過的地方.
飛行的書卷跟聖殿的門廊的尺寸.
大小相同.
其實不是偶然.
而是有一個聯繫在當中.
是要讓神的子民明白到.
這個書卷的重要性.
是跟進入聖殿有莫大的關係.
因為每一個進入聖殿的人.
都必須要經過這個門廊.
就好像在這裡.
神所宣講的審判的訊息.
就是要告訴神子民知道.
他們要去重建聖殿.
成為聖殿裡的祭司.
去敬拜神 去事奉神.
他們就必須要掌握.
這個書卷上的訊息.
事實上在第三節當中.
天使的解釋.
也都支持了剛才的說法.
因為天使說出.
這個飛行的書卷.
其實是法出行在遍地上的咒詛.

$^{321}$咒詛在原文的意思就是.
違背誓言或者違背盟約所帶來的懲罰.
這個豈不是.
《申命記》神學裡所提到的咒詛嗎?.
天使宣告這個咒詛是從神而來的.
這是神的審判帶來的咒詛.
這個咒詛是很徹底的.
因為它遍及整個地上.
所有的神的子民都包括在其中.
神的審判會覆蓋全地.
覆蓋整個國家.
覆蓋整個民族.
覆蓋所有神所揀選的子民.
意思是沒有一個人能夠逃避這個審判.
神的審判的對象不是神子民的敵人.
而是審判他所愛所揀選的子民.
神審判的準則是什麼呢?.
天使繼續說.
凡偷竊的必按捲上者面的話除滅.
凡起假誓的必按捲上那面的話除滅.
天使宣告這個咒詛.
這個審判的對象就是這兩類的人.
偷竊者和起假誓的人.
事實上這兩種的罪都是在十誡裡面被禁止的.
不可妄稱耶和華你神的名就是起假誓者的罪名.
不可偷渡就是偷竊的人的罪名.
在十誡裡面偷竊是違反關於愛人愛鄰舍的誡命.
而起假誓是妄稱神的名.
就是涉及到違反關於愛神的誡命.
換言之這兩類的人其實是代表全部的律法.
當中的意思就是凡違背神的命令.
無論他是違反了對愛人的命令.
還是對愛神的命令他們都會受到神的咒詛.
沒有人可以逃避神的審判.
這個標準是寫了在這卷飛行書卷上面的兩邊.
就是兩邊都是寫上了神的審判的要求.
而這個審判的準則既然是寫在書卷的正面和背面.
就是代表著神的誡命.
就是代表著神在西乃山和祂子民所立的約.
所定下的誡命律例典章.

$^{361}$因此這個飛行的書卷是提醒我們.
神是會遵守這個和人所立的約的.
如果人願意遵守的話.
他就會得到祝福.
相反他就會審判這些罪人.
他們就會得到救助.
換言之.
在這個書卷上面所寫的一切內容.
是為了提醒神的子民要聖潔.
因為任何不悔改的罪人.
都沒有辦法進入神的殿.
沒有辦法進入神的同在.
因為每一個要進到聖殿當中.
去經歷去享受神同在的子民.
都必須達到聖潔的標準.
這個標準並不是由人所定立的.
而是由神親自定立.
親自頒佈.
神之所以要頒佈這一切的誡命.
就是要為了幫助神的子民.
來處理他們罪的問題.
因此這本書不只是陷入懲罰的威脅.
其實也帶有潔淨神子民的盼望.
這本書的出現.
出現在先知撒迦利亞的眼前.
那個時間點就正是.
這班猶太人面對重建聖殿的歷史處境.
因此這本書的出現.
其實是正確地出現在這個時間點.
因為重建聖殿.
從來都不只是講硬件.
而更重要的是軟件.
重建聖殿的軟件就是神子民的心.
是神子民的心.
神要以這個意象來提醒猶太人.
他們要預備好自己的心.
要謹守進行神的誡命.
這樣才會得到神的祝福.
這樣才可以重建聖殿.
否則就會招致神的就坐.

$^{401}$接著神親自去宣告這個審判.
祂說「我必使者書卷出去」.
「我必使者書卷出去」.
神的意思是.
祂的審判一定會按照祂的心意.
祂的計劃.
祂的應許會臨到成就.
這個事是決定性的.
是絕對會發生的.
神是絕對不會收回祂的審判.
神要親自潔淨祂的子民.
神繼續說.
這本書會進入偷竊人的家.
會進入假釋人的家.
會經常在這些人的家裡.
連房屋大木石都毀滅了.
神的意思是甚麼呢.
神的意思是.
這個代表著神審判的書卷.
進到偷竊.
進到發假誓的人家裡.
進去要做甚麼.
神要進行審判.
神要審判他們.
神要將這些罪.
從他們的子民當中除去.
神的審判就是要這樣.
進到這些罪人的家裡.
要去完全徹底毀滅他們的家.
這些罪人以為自己.
建立他們的家.
寫書的地方.
留在裡面就安全了.
但是神說錯了.
我的審判會進到你們.
以為安全的地方.
以為安全的家.
我要將你們以為安全的東西.
徹底的毀滅.
這個毀滅是.

$^{441}$沒有一塊木頭.
沒有一塊石頭.
在另一塊木頭 另一塊石頭上面.
徹徹底底.
完完全全的毀滅.
這個意思.
其實就是代表著罪惡.
會由內而外.
來摧毀神子民的生命.
神子民的屬靈生命.
他會被罪惡來毀滅.
侵蝕.
這個是要指出.
罪惡的破壞性和毀滅性.
神三次的重複家這個字.
事實上.
這裡的家.
其實和聖殿.
希伯來文裡的神的殿.
是同一個字.
神在這裡要責備神的子民.
當我們看同期的哈該書.
一章三至四節的時候.
裡面記載神透過先知哈該.
去責備這班猶大人.
神的責備是這樣說.
聖殿仍然荒涼.
你們自己還住天花板的房屋嗎.
神責備這班猶大人.
竟然放下聖殿重建工程的這件事.
但是自己就掛起自己的房屋.
建立自己的家.
他們自己的家有天花板.
但是聖殿呢 神的家連天花板都沒有.
為什麼這班猶大人竟然會忘記了神呢.
為什麼他們會本末倒置.
他們作為神所揀選的子民.
豈不是應該將神放在首位.
放在第一位嗎.
為什麼在這一刻.

$^{481}$竟然他們是將自己的安書放在第一位呢.
這個就是神子民的罪.
第六個異象.
這個飛行的書卷是一個.
令到神子民沒有辦法不看到的異象.
一個打開的書卷.
一個這麼大型的書卷.
裡面記載了這麼震撼的.
神的審判 神的咒詛的內容.
唯一的目的就是要讓人看見.
要讓神的子民看見.
才看見了.
猶大人看見了.
這本書的每一個特點.
都是指向唯一的目的.
是要被看見.
因為神要讓所有神的子民都知道.
都明白神的關注.
神的心意 神的計劃.
是要建立一個聖潔蒙愛的群體.
神的審判 神的懲罰 神的警告.
這一切一切從來都是出於神對神子民的愛.
昔日先知撒迦利亞看到.
昔日的猶大人他們看到.
他們看見神是那位審判全地的主.
他們看見神是那位愛他們到底的神.
他們看見神是祝福他們的那位神.
因著他們每一個人的看見.
他們懂得回應神的警告.
他們懂得回應神的審判.
他們懂得回應神的愛.
他們終於選擇踏上重建聖殿的路.
以至到聖殿最終是順利被重建起來.
這就是神對神子民的祝福.
在被擄之前 神子民因著他們犯罪.
違背神的約 以至到神的家聖殿被毀.
在回歸之後 神的子民領受神審判的訊息.
這個愛的訊息 神的子民選擇回轉.
歸向神 認罪 悔改 歸向神.
神親自賜下福興帶領他們來重建聖殿.

$^{521}$讓他們恢復在聖殿裡的敬拜.
最重要的是恢復他們的心.
對準上帝 成為神所喜悅的子民.
今天我們在這個異象裡我們看見什麼呢?.
我們看見的是不是昔日先知所看見的呢?.
我們看見的是不是昔日這群猶大人所看見的呢?.
今天神要向我們佔領這個屬靈的群體.
這個屬於祂的群體.
這個被祂所揀選的群體.
要我們怎樣來回應.
回應祂的愛 祂的審判呢?.
我們知道耶穌基督.
祂會在終末的時候再次來到.
祂會帶來最後的審判.
祂就是審判全地的主.
但是這個審判的訊息.
不是一個令我們惶恐的訊息.
不是單單令我們感到很害怕.
不想面對 不想聽.
令我們不開心的訊息.
在這個審判的訊息的背後.
其實是帶著這個福音的好訊息.
因為耶穌基督在十字架上為我們流血 捨身.
以祂的寶血去洗淨我們眾人的罪.
以致到聖靈內住在我們的生命裡面.
我們的身子成為聖靈的殿.
罪在昔日令到實體的聖殿被拆毀.
但是神親自賜下祝福.
臨在神的子民身上去重建這個實體的聖殿 重建這個群體.
今日我們都一樣經歷耶穌基督的拯救.
祂的救恩 祂的福音就是要我們每一個人成為聖靈的殿.
當我們今日要講復興的訊息.
我們必須要認真面對我們生命裡面罪的問題.
我們不回轉 我們不悔改 我們沒有辦法復興.
這個就是一個真實的現象.
而且是會面對最終神的審判.
沒有人想被罰的 是不是.
有沒有人想被罰的.
應該沒有人會舉手的.
我們每一個人都想得到寶貴被神揀選 作神子民的身份.

$^{561}$我們很想進到這個榮耀的國度裡面.
我們知道神已經差遣耶穌基督來拯救我們.
我們在這個國度裡面是有份的.
神今日要提醒我們要復興的時候我們必先回轉歸向祂.
這個就是飛行書卷這個異象的訊息.
要看見神的審判.
神的審判是兩面的 好像一個銀子的兩面.
有祝福亦有就坐.
但是我們會得到祝福還是就坐就是決定於我們的選擇.
我們是否願意回轉歸向神.
如果我們願意的話.
神是無條件的將這個祝福的窗戶打開.
傾福給我們每一個人.
因此今日我們處在耶穌基督還未再次來臨的時候.
我們是期待著祂第二次的降臨.
我們期待著新聖殿的降臨.
這個就是我們終末的盼望.
我們在信仰裡面神給我們的應許.
但是在這個已然未然already but not yet的時候.
我們要做什麼呢?我們是否什麼都不做呢?.
今日的訊息提醒我們.
這個書卷是飛行的 是全速的向前.
是要向全地宣講.
神是那位愛我們的主.
祂的審判亦會臨在全地.
因此我們要怎樣去回應這個訊息.
在個人的方面我們要再一次提醒自己.
我們有沒有認真去到神的面前去認罪去悔改呢?.
在教會的層面裡面.
我們有沒有圍著我們教會這個屬靈的群體來去禱告.
去祈求我們能夠得著神的憐憫.
將復興賜下給我們的教會呢?.
在世界的方面我們有沒有這個全福音的心智.
神是很迫切去向當時的先知.
當時的尤大人去宣告這個訊息.
用一個很大很大的書卷來向人宣告.
祂的訊息 祂的愛.
我們今天有沒有這個迫切性去將神的福音.
將神審判的訊息 將神愛的訊息.
向我們身邊未信的人分享呢?.

$^{601}$我們有沒有這個迫切性.
我們是不是行動中 我們是不是進行中呢?.
盼望今天的訊息能夠再次成為我們眾人的提醒.
我們祈求神的祝福的同時.
我們都要有一個行動來去回應.
希望神的話語潔淨我們眾人 潔淨沾浸.
以致我們能夠一起成為神所喜悅.
聖潔蒙愛的群體.
我們一起低頭同心禱告.
親愛的主耶穌基督.
讓我們可以看得見.
讀得到.
聽得明白.
放在心裡.
因為你從一開始.
就是因著你對我們的愛去揀選我們.
因著你對我們的愛.
去留下這麼寶貴的訊息.
我們明白到.
神父的訊息並不是只是為我們帶來恐懼.
而是讓我們帶著這一份感恩的心去領受.
因為你向我們所留下寶貴的話語.
為要提醒我們.
你是愛我們的.
而你的愛就是在主的十字架上向我們彰顯.
我們知道.
你甘願為我們背負我們的罪.
潔淨我們.
好叫我們能夠成為聖潔.
好叫我們能夠在你的審判台前站立.
主啊求你幫助我們.
讓我們真是天天來到你的面前去認罪去悔改.
我們要去操練自己的靈命.
我們要去潔淨我們自己.
因為你說過.
人若自潔脫離卑賤的事.
就必作貴重的器皿.
成為聖潔合乎主用.
我們就是帶著這一份的心.
來到你的庚前向你呼求.

$^{641}$願你潔淨我們.
讓我們成為你合用的器皿.
在這個世代的當中.
與你同行.
見證耶穌基督.
你就是世上的真光.
照亮這個黑暗的世界.
求你親自祝福我們每一個.
我們這樣禱告.
交託奉靠我主耶穌基督.
得勝有得勝的名字.
祈求.
阿們.
\newpage



\section{路加福音 2:41-52}
\label{sec:032150}
\textbf{成為一個充滿恩寵的人 (華)}~~~~曾興才牧師
\newline
\newline
連結: \href{http://www.church.com.hk/acms/content.asp?site=cdc\&op=show\&type=product\&code=032150\&layout=sermon}{\texttt{032150}} ~~~~ 語音日期: 2026-05-17
\newline
\newline
\hyperref[sec:032149]{\small{< < < PREV SERMON < < <}}
~
\hyperref[sec:index_chronic]{\small{[返順時目]}}
~
\hyperref[sec:032151]{\small{> > > NEXT SERMON > > >}}
\newline
\newline
路加福音 2:41-52
\newline
\begin{longtable}{cl}
\hline
\hline
章節 & 經文 (和合本修訂版)\\
\hline
2:41 & \begin{tabularx}{0.7\textwidth}{X} 每年逾越節，他父母都上耶路撒冷去。 \end{tabularx} \\ \\ \relax
2:42 & \begin{tabularx}{0.7\textwidth}{X} 當他十二歲的時候，他們按著過節的規矩上去。 \end{tabularx} \\ \\ \relax
2:43 & \begin{tabularx}{0.7\textwidth}{X} 守滿了節期，他們回去，孩童耶穌仍舊在耶路撒冷。他的父母並不知道， \end{tabularx} \\ \\ \relax
2:44 & \begin{tabularx}{0.7\textwidth}{X} 以為他在同行的人中間，走了一天的路程才在親屬和熟悉的人中找他， \end{tabularx} \\ \\ \relax
2:45 & \begin{tabularx}{0.7\textwidth}{X} 既找不著，就回耶路撒冷去找他。 \end{tabularx} \\ \\ \relax
2:46 & \begin{tabularx}{0.7\textwidth}{X} 過了三天，他們發現他在聖殿裡，坐在教師中間，一面聽，一面問。 \end{tabularx} \\ \\ \relax
2:47 & \begin{tabularx}{0.7\textwidth}{X} 凡聽見他的人都對他的聰明和應對感到驚奇。 \end{tabularx} \\ \\ \relax
2:48 & \begin{tabularx}{0.7\textwidth}{X} 他父母看見就很驚奇。他母親對他說：「我兒啊，為甚麼對我們這樣做呢？看哪，你父親和我很焦急，到處找你！」 \end{tabularx} \\ \\ \relax
2:49 & \begin{tabularx}{0.7\textwidth}{X} 耶穌對他們說：「為甚麼找我呢？難道你們不知道我應當在我父的家裡嗎？」 \end{tabularx} \\ \\ \relax
2:50 & \begin{tabularx}{0.7\textwidth}{X} 他所說的這話，他們不明白。 \end{tabularx} \\ \\ \relax
2:51 & \begin{tabularx}{0.7\textwidth}{X} 他就同他們下去，回到拿撒勒，並且順從他們。他母親把這一切的事都存在心裡。 \end{tabularx} \\ \\ \relax
2:52 & \begin{tabularx}{0.7\textwidth}{X} 耶穌的智慧和身量，並神和人喜愛他的心，都一齊增長。 \end{tabularx} \\ \\
[1ex]
\hline
\hline
\end{longtable}
$^{1}$弟兄姊妹平安..
也歡迎網絡上的弟兄姊妹跟我們一起來敬拜詔明文的主..
很高興能夠進到你的家中,一起來分享神的話語..
今年教會背景行動正在如火如荼地進行..
到目前為止,已經有1452位弟兄姊妹註冊..
請問你註冊舉手一下..
但是呢,到昨天為止,背完1月份到4月份的經文只有532位..
大家要加油,高總管要加油..
要繼續努力來思想背誦神的話語,因為神的話語太重要了..
不但我們要背誦神的話語,我們更是要在日常生活當中來應用神的話語,.
好讓神的話語真的可以成為我們每天祝善的光,矯情的燈..
阿們.大人小孩也是如此,要懂得如何應用神的話語..
在正經書第22章第6節告訴我們說,教養兒童,使他走當行的道,.
就算到老,他也不偏離..
有一天,老張三歲多的女兒就哭著來他面前跟他說,.
"弟弟坐在我的玩具上,怎麼叫他都不肯起來.".
老張一邊刷牙一邊跟他說,.
"你要想想辦法啊."他心裡想,.
其實我們常常做的事情就是用另外一個玩具把弟弟引開不就好了嗎?.
沒想到大概三十秒之後,聽到他對弟弟大聲地喊叫說,.
"我奉拿撒勒人耶穌基督的名叫你起來行走.".
高中旁邊說,讀聖經的孩子是很聰明的..
今天我們就來思想,拿撒勒人耶穌,拿撒勒人耶穌基督..
聖經裡面很少記載耶穌基督童年的事情,.
但是唯有《路加一生》很清楚地記載一段非常有趣耶穌基督童年的一個事件..
在馬太,馬可,約翰福音都沒有記載,唯一在《路加福音》第二章41節到52節我們可以來看這段非常有趣耶穌基督童年的一個記錄..
這裡說,每年到育月節,他母親父母就上耶路撒冷去..
當他十二歲的時候,他們按著節旗的規矩上去..
守滿了節旗,他們回去,孩童耶穌仍舊在耶路撒冷,他的父母並不知道..
42節說,以為他在同行的人中間走了一天的路程,就在親屬和熟悉的人中找他..
既找不著,就回耶路撒冷去找他..
過了三天,就遇見他在店裡,坐在教室中間,一面聽,一面問..
凡聽見他的,都稀奇他的聰明和他的應對..
他父母看見,就很稀奇..
他母親對他說,"我啊,為什麼像我們這樣行呢?看啊,你父親和我三行的來找你啊.".
49節,耶穌說,"為什麼找我呢?.
豈不知我應當以我父的事為念嗎?".
他說這話,他們不明白..
他就同他們下去,回到拿撒勒,並且順從他們..
他母親把這一切的事都存在心裡..

$^{41}$52節,我們一起來讀來..
耶穌的智慧和聖量,並神和人喜愛他的心,都一起增長..
再來一遍來.耶穌的智慧和聖量,並神和人喜愛他的心,都一起增長..
這是耶穌記錄所有事件當中,記載他童年最有趣的一段事情..
我們先來做個幾個重要的觀察..
首先說,耶穌第十二歲的時候,他父母親按照猶太人的規矩,.
就是根據《申命記》第十六章十六節裡面記載,.
所有猶太人的男人一年三次必須回到耶路撒冷去朝聖,敬拜神..
特別是一月節這麼大的慶典,.
他們會襲家帶眷地一起朝向耶路撒冷,到神的殿中來敬拜神..
成為非常敬虔的約瑟和瑪利亞,他們肯定也一樣,.
帶著十二歲的孩子耶穌一起到耶路撒冷去..
那你說,耶穌在人群當中怎麼會走丟了呢?.
其實按照猶太人的習俗,他們一般上都會借伴而行,.
跟他們的親戚或者鄰居一群人一起到耶路撒冷去來朝聖,來敬拜神..
因此這時候,可能約瑟在想,.
因為按照他們的習慣,通常婦女跟小孩子會走在前面,.
然後男人就在殿後..
因此父親可能約瑟可能想,可能耶穌就在他母親的一群人當中..
但是他母親也以為耶穌可能跑到後面去跟爸爸在一起了..
結果他們到第二天的時候,.
發現走了一天的路程,可能喫點心的時候,發現耶穌不見了..
在他們親戚當中,熟悉的人當中找不著,.
結果他們又回去耶路撒冷..
你可知道從拿撒勒到耶路撒冷大概需要一百二十公里的路程,.
需要三到四天或者四到五天..
他們只走了一天的路程,可能還到三分之一的路程,.
他們就發現耶穌不見了,又花了一天的路程走回到耶路撒冷,.
在那裡到處尋找耶穌..
甚至到了第三天,找了一整天以後,在那裡找到耶穌呢?.
到聖殿裡面,看見耶穌基督坐在老師中間..
在那裡面做什麼呢?一面坐,一面聽,一面講..
你可知道父母親孩子不見了多麼操心啊!.
看到耶穌以後,耶穌的母親瑪利亞有一點埋怨,責備的心:.
"你為什麼這樣待我們呢?讓我們做父母親的,到處傷心地來找你呀!".
耶穌怎麼回答呢?.
"你們為什麼找我呢?難道你不知道我應當,.
我必須,我一定要以父的事為念呢?".
"你們為什麼找我呢?這個很有趣的一個問題,對不對?.
難道孩子不見了,父母不會找你嗎?.

$^{81}$怎麼耶穌基督好像心中有一點失望,.
父母親難道不了解耶穌基督他應當做的事情嗎?.
他說我應當,我必須,我一定,.
英文聖經說I must,必須以神天父的事情為念..
結果耶穌基督沒有固執不走,硬要待在聖殿,.
結果他跟著他的父母親回去納撒勒,並且順從他父母親..
所以瑪利亞看到這些事情以後,他心中不明白,.
他就把這些事情藏在心裡,這是非常有智慧的女人所做的事情..
但你不明白的事情,不要下任何的定論,不要做任何的批評,.
他就把這些事情放在心裡,時候到的時候他肯定會明白的..
第52節很清楚地告訴我們,.
耶穌的智慧與生亮,並神和人喜愛他的心,都一起增長..
換句話說,耶穌基督在他成長的過程當中是充滿智慧的,.
身體是非常健壯的,而且他得神的喜愛,.
而且是有屬靈的恩寵,神非常喜悅他,這個恩寵靈到他身上..
不但是如此,因此也得眾人的喜愛,.
人也喜愛他,得人的喜愛,得人的恩寵..
換句話說,他是一個充滿恩寵的一個人..
所以你見耶穌基督後來出來服侍的時候,.
他所到之處前呼後擁的集滿了人,.
大家都想來糾盡耶穌,都想來親近耶穌,.
因為他是充滿神的恩寵,也充滿人的恩寵的人..
今天早上我們一起來思想,我們都要成為一個充滿恩寵的人..
不論是充滿神的恩寵,或者是我們人的恩寵,.
我們都需要有恩寵..
因為一個沒有恩寵的人,他所做的事情盡都不順利,到處碰壁..
一個有恩寵的人,是有神的恩典在他身上,.
有眾人也在他身上,看到神的柔美,也會特別對他有這樣的恩寵..
為什麼需要恩寵呢?你想想看,.
我們生活在充滿邪惡,而且彎曲背廟的時代,.
我們每天面對很多不同的假新聞,不同的挑戰,.
我們需要神給我們的恩寵,讓我們有個辨別是非的能力,.
這是恩寵是要從神而來的..
第二個理由為什麼要需要恩寵,.
因為我們生活在充滿挑戰的時代,IT迅速的發展,.
工作上,職業上受到很大的挑戰,.
婚姻,家庭,人際關係,服飾也受到很大的挑戰,.
我們需要神的恩寵,也需要人的恩寵,.
我們才能夠渡過或者勝過這一切的挑戰..
第三個理由,今天我們教會的主題是什麼?.

$^{121}$興起一支屬神的大軍..
一個屬神的大軍必須要有神的恩寵,.
還有我們需要人的恩寵在我們身上,.
我們才能夠繼續不斷的跟隨耶穌,.
活出耶穌基督柔美的生命,.
能夠成為神偉大,而且屬靈的大軍,.
能夠為神做美好的見證.阿們..
今天早上,你我是否很想成為充滿恩寵的人呢?.
你想不想?.
好像只有一些人..
你想不想成為充滿恩寵的人?.
我想成為滿有恩寵的人,.
因為當你有恩寵的時候,.
你知道你手中所做的盡都順利,.
神和人喜愛你的心都隨著年日來增長,.
這是好的無比的..
因此我們一起來思想,.
從最有恩寵的人身上,耶穌基督的身上,.
來學習如何成為一個充滿恩寵的人..
首先第一,我們要全心全意的來可牧神,.
告訴旁邊的人說要全心全意的來可牧神..
要全心全意的來可牧神,.
不是三心兩意的有時間再來可牧神,.
沒有時間就不來可牧神,.
而且是全心全意的來可牧神..
你知道耶穌基督當他走丟了以後,.
他父母就到處找他,找不到..
到了第三天,我們看看第46節跟47節,再看一下..
過了三天就遇見他在殿裡,坐在教師中間,.
動詞很重要,.
坐在教師中間一面聽一面問,.
凡聽見他的都稀奇他的聰明和他的應對..
特別這裡講到他在殿裡面坐在教師中間,.
不是站在他們當中,.
那你站在那裡的時候可能是你不太有興趣,.
隨便聽一下就走了..
一個站在旁邊的人是不是那麼專心,.
要來把時間投入在大家所談論的問題當中..
耶穌基督不但是坐在那裡,.
而且坐在那裡,坐在教師中間..

$^{161}$你可以想象,哇,12歲的孩子坐在教師中間,.
這怎麼講呢?.
特別是預鑰節的時候,.
當時在聖殿裡面許多猶太人的老師,.
包括什麼所謂的神學教授都會為眾人來講解聖經,.
耶穌基督是一面聽,一面問..
之所以能夠聽,能夠問,.
表示說他是非常專心地聽,非常渴慕地聽,.
他纔會想這些問題來問..
因此他被邀請坐在教師中間..
你可以想象,很可能他問了一些問題,.
連這些神學教授,猶太人老師都無法回答..
他們都稀奇,哇,天才兒童,12歲的孩子,.
怎麼問出這樣的一個問題來?.
12歲的孩子怎麼答得出這麼困難就業聖經背景的問題呢?.
他們想都想不到,.
12歲的孩子一面坐在那裡,一面聽,一面問..
這是耶穌基督全心全意地來渴慕..
所以他說,難道你不知道我應當,我必須,.
以我父的是為念嗎?.
所以首先你要做什麼?.
你要做,你要聽,你要問,你要做,你要聽,你要問..
你要做,你要聽,你要問..
基督你坐在這裡有福了,好好地做,好好地聽..
當你聽以後,你瞭解在你腦海當中彷彿思想,.
你纔會問問題嘛..
那你問問題的時候,你會更多地學習,更多地吸收..
這乃是一個全心全意渴慕神的人要做的事情,.
必須懂得坐在那裡,坐在那裡..
但是你知道,當耶穌基督的母親找來找去找不到他的時候,.
看到耶穌的時候,.
祂說,你們為什麼這樣待我們呢?.
難道你不知道你和父親上天來找你嗎?.
然後第49節耶穌說什麼呢?.
很有趣,耶穌說,為什麼找我呢?.
豈不知我應當以我父的事為念嗎?.
中文聖經翻譯成應當,其實更好的翻譯應該是,.
我必須,我一定要以我父的事情為念..
所以英文聖經翻譯成I had to,或者I must,.
我必須,我一定要以我父的事情為念..

$^{201}$為什麼找我呢?難道你不知道嗎?.
這裡耶穌基督帶點失望的口氣,.
難道你們還不知道我應當必須一定要以我父的事為念嗎?.
這就引起很多神學家的討論,.
到底耶穌基督什麼時候知道他是要來的彌賽亞,.
是人的舅子,什麼時候知道他真正的身份呢?.
神學家辯論來辯論去,他們最終的結果是,.
耶穌可能一早就知道他是要來的彌賽亞,.
他要來拯救世人的彌賽亞..
所以他才說,你們為什麼找我呢?.
這句話聽起來不是很有趣嗎?.
父母親找你,應該的,你不見掉,怎麼不找你呢?.
你是我的孩子嘛,怎麼不找你呢?.
你為什麼要找我呢?.
難道你不知道,豈不知道我應當必須以我父的事為念嗎?.
看到耶穌基督從小就知道他生命的優先順序,.
當我們全心全意的來尋求神,來渴慕神的時候,.
必須把我們的生命的優先次序擺好..
他怎麼還沒說,你要有正確的優先次序?.
正確的優先次序,一個有正確優先次序的人,.
他的生命纔是成功的,纔是會有影響力的,.
因為他懂得把最重要的事情,最首要的事情放在最重要,最首要的位置,.
不能本末倒置..
那麼把次要的事情,不關痛癢的事情,.
放在最重要的地位上,花很多時間去做這個事情的時候,.
我們就失去了我們生命的中心..
所以耶穌就非常清楚知道,我要以我父的事情為念,.
他懂得他生命中的優先順序,.
把優先順序擺正的時候,其他的事情都會按部就班地成就了..
這是非常重要的..
所以耶穌就在《登山寶訊》講了一大篇道理以後,.
在快要結論的時候,在《馬太文》第六章三十三節,.
大家非常熟悉的,都可以背起來的,一起來讀來..
你們要先求祂的國和祂的義,這些東西都要加給你們了..
所以說你們要先求,當你看到先這個字,.
see first,表示說這裡面有什麼?.
有優先的順序,什麼是第一,什麼是第二,什麼是first,什麼是second..
換句話說,耶穌基督講到說,我們生命當中每天都做出選擇,.
早上你要來崇拜,要穿什麼衣服,女生要打扮成什麼樣子,.
我早上要穿什麼西裝,配什麼領帶,都要做很多的決定,.

$^{241}$每天都要做出選擇..
所以你們要先求神的國和神的義..
換句話說,耶穌基督說在你生命中最優先,最重要,最首要的任務是什麼呢?.
不是求這些東西,前面所講外八人所求的一事駐行,.
而是先求神的國..
那你說,到底什麼是神的國?.
神的國代表神在哪裡?.
坐寶座,掌權的地方..
神的國代表神的權柄,代表神掌權的地方..
換句話說,你的婚姻家庭,夫妻關係,.
是否耶穌基督在那裡坐寶座掌完全呢?.
誰生病比較大,說了算,還是耶穌基督說了算?.
在你的工作服飾,人際關係上,.
是否耶穌基督也在那裡掌權,坐寶座掌完全呢?.
你是不是按照祂的心意來經營你的人際關係,來做你的工作,做你的服飾呢?.
這是神的國的意思,神在這個地方掌權,你叫做尊主偉大..
那你說先求神的國和神的義,什麼是神的義呢?.
神的義代表神的屬性,神的位格,祂的character,祂的attributes..
我們的神是有位格的,祂是有屬性的,.
祂是公義的,聖潔的,慈悲,憐憫,良善,慈愛..
因此當我們先求神的國,神的義的時候,.
那是在我們生命當中如何能夠更多彰顯.
神賦給我們的祂的屬性,聖潔的生活,良善,公義,憐憫,慈悲..
那我們要先求神的國,神啊,禰現在在哪裡掌權?.
是不是在我的家庭,我的生命中掌權?.
是不是要透過我的生命,透過我的夫妻關係,透過我的人際關係,我的工作服飾,.
能夠彰顯出主你的生命來?.
這是耶穌所講的,你先求神的國和神的義,結果是什麼?.
這些東西,這些東西都喊不喊,隆重所要..
不論是給你們的,都要加給你們的,哇,好不好啊!.
你已經有了嘛,還要加給你啊!.
你看,馬太福音第二三十三之前三十二節到前面二十節左右,.
都是講說外邊人每天想說要穿什麼,喫什麼,去哪裡,做什麼..
耶穌基督說,當你先求我的國和我的義的時候,.
這些東西全部都要加給你了,好不好啊!.
換句話說,當你把優先的次序擺正的時候,.
可以說你是一無所缺的人,你凡事所行的人是盡都順利的,.
因為神的恩高,神的恩寵在你身上嘛,阿們!.
聖經裡面有一個非常明顯的例子,.
就是常常到迦百農一個家庭裡面做客,就是馬大和瑪利亞的家..

$^{281}$非常熟悉的故事,在路加福音第十章,我們看第三十八節和三十九節,.
他們走路的時候,耶穌進了一個村莊,有個女人名叫馬大,.
接她到自己家裡來,她有一個妹子名叫瑪利亞,.
在耶穌腳前做什麼?很熟悉的字啊,坐坐,聽耶穌的道..
瑪利亞也是坐坐,不走了,選擇就坐在耶穌腳前做什麼?.
聽耶穌講道..
後來馬大忙得不可開交,就開始生氣埋怨,甚至有點責問耶穌,.
你看我的妹子把我一個人丟在廚房,請你讓她來幫助我..
耶穌說什麼?.
馬大,馬大,你為許多事情憂慮煩惱,但是不可少有的只有一件,.
瑪利亞已經什麼?選擇..
當你看到選擇這個字的時候,就表示說有修仙順序了..
你選擇馬大,瑪利亞選擇他上好的福分是不能奪去的..
馬大和瑪利亞都要做選擇嘛,.
馬大選擇要來把十八般無異的初遇,.
展現在耶穌基督的面前,好好地服侍耶穌,.
瑪利亞卻坐在耶穌基督的腳前聽祂的道..
親愛的兄弟姐妹,這裡不是說服侍不重要,.
服侍固然重要,.
但在服侍跟親近主的優先順序裡面,.
瑪利亞是做對的,.
先坐在耶穌基督的腳前,先明白耶穌基督到底要的是什麼,.
要你怎麼樣來服侍祂,在哪裡服侍祂,.
讓你明白耶穌的心意以後,你纔來服侍祂的話,.
你肯定是十半功倍的,阿們?.
所以瑪利亞纔得到耶穌基督的稱讚,.
祂選擇了上好的福分,世人不能奪去的..
愛的弟兄姊妹,你目前生命的優先順序正確嗎?.
神在你心中的地位是第一還是第二呢?.
是否常常忙碌,以你的工作服侍取代了神應當在你生命中的地位呢?.
一個全心全意可慕神的人,一個滿有神的恩寵人的恩寵的人,.
必須懂得如何全心全意的可慕神,阿們?.
第二,我們要竭力遵從神的話,要竭力遵從神的話..
換句話說,我們不但只是聽神的話,不但只是來背誦神的話,.
更重要的是我們要來順從聽從神的話,.
神的話才能夠影響我們的生命,我們的生命才能夠帶出影響力..
神的話是大有功效的,阿們?.
我們看到第五十一節,當耶穌的母親瑪利亞說以後,.
耶穌基督沒有強制自己,不讓祂母親帶祂回拿撒勒,.
祂就從他們下去,回到拿撒勒,並且順從他們,.

$^{321}$祂母親把這一切的事都存在心裡..
耶穌基督雖然是神的兒子,祂也知道祂是仁子,.
乃至透過在地上的母親瑪利亞聖靈感應懷胎生子,.
祂有做兒子的責任..
你們做兒女的要在主裡孝敬父母,這是唯一帶有應許的誡命..
記得嗎?一部手書,跟哥羅斯的書跟我們講的..
所以耶穌基督順從祂父母親,跟他們回到耶路撒冷去..
雖然在聖殿非常的好,祂可以坐在教師當中一面聽一面問,.
但是祂還是跟他們回去,並且順從他們..
並且順從他們的意思是什麼呢?.
回去,乖乖地做祂月社和瑪利亞的兒子,.
幫忙祂父親做木匠,幫忙祂父親做木工..
當然想說,耶穌基督幫祂父親做木匠的時候,.
祂所做出來的遺址,坐址是否是十分的完美,.
而且非常的耐用..
現在你到拿撒勒,到一路上,.
是不是還可以找到耶穌基督所做的遺址,坐址啊?.
你想想看,神的兒子做遺址,坐址啊,不得了..
耶穌就是如此地順從祂的父母親,.
盡祂兒子丹靖的本分..
這是為什麼耶穌基督身上有這樣一個恩寵..
記得嗎?十八年以後,耶穌開始坐在床上,.
到三十歲的時候,在施洗約翰維拉受洗之後,.
天上的天打開了,聖靈向各自降在祂的身上,.
天上有聖靈說什麼?.
這是我的愛子,我所洗約的..
那時候耶穌根本一個福音都沒有傳,.
一個天國的哲理都還沒有講,.
一個神理啟示都還沒有做,.
天父已經說這是我的愛子,我所洗約的..
因為耶穌基督順服父母,.
祂是一個懂得順服,而且是全然順服的人..
所以耶穌是全然順服,.
不但祂是順服地上的父母親,.
祂也順服在天上的阿爸父,.
祂也順服聖經裡面所講的這些話..
我非常喜歡使徒保羅當他講到基督論的時候,.
最精闢的基督論就在腓立比書第二章第五到第十一節,.
這裡說:你們當以基督的性為心,.
祂本有神的形象,卻不以自己與神同等為強奪者,.

$^{361}$反倒虛極虛弱,奴僕的形象,成為人的樣子..
第八節是什麼呢?既有人的樣子,就怎麼樣?.
自己卑微,全心順服,以至於死,且死在十字架上..
這是全然順服最好的一節經文..
他既有人的樣子,就自己卑微,.
沒有人強迫他要謙卑他自己..
他可以說我就是要在神的殿中,.
好好來渴慕天父的事情,為父的事情,.
因為我就是為這個事情出來的..
他沒有,他自己卑微,而且全心順服,.
選擇順服,而且順服到什麼地步呢?.
以至於死,如果到這裡結束的話,還不覺得稀奇..
且死在,保羅特別強調,even death unto the cross,.
且死在十字架上..
什麼人是可以定在十字架上的?.
外辦人,而且是那種重犯,殺人放火,大強盜,.
是要被釘死在十字架上的..
這個完全公義,完全聖潔,沒有罪的耶穌基督,.
他全心順服,以至於死,且死在十字架上..
難怪保羅說,因此神將他升為至高,.
是給他超過萬名的名..
萬名人叫一切在天上的地上的地底下的,.
無不屈膝,無不口稱耶穌基督為主..
耶穌全然的順服..
這是為什麼耶穌基督有這麼多的恩高在他身上,.
這麼多的恩寵在他身上,因為他是一個全然順服的人..
也是如此的順服神,順服神給你們的護照,.
給你們的帶領,全心全意來服侍主..
接下來我們看到,怎麼知道一個人他是完全全心的順服呢?.
看到他是否遵行神的話,是否是一個行道的人..
一起告訴萬名人,要行道不要單單聽道..
這是我們來逐日崇拜,讀聖經,背誦聖經最重要的一個目的,.
要把神的道行出來..
不是隻是單單聽道,聽道容易,.
行道需要一點功夫,真的需要付代價..
所以在雅各書第一章二十二節,雅各很清楚告訴我們說什麼?.
只是你們要行道,不要單單聽道..
欺侮就是欺騙自己..
換句話說,一個不行道的人,只是聽道的人,是欺騙你自己嗎?.
以為你聽懂了,明白了,但是你卻沒有消化行出來的話,.

$^{401}$你是自己在騙你自己..
接下來第二十三節,第一章雅各書,雅各繼續說什麼呢?.
因為聽道而不行道的人,就好像人對著鏡子看自己本來的面貌,.
還不錯,還蠻帥,還蠻漂亮的,.
看見走了以後,隨機忘了他的相貌如何..
是很可笑嗎?明明看得很清楚,.
特別是姐妹們照鏡子照好多次,照好久..
照完以後,離開以後,自己長什麼樣子都忘記了,.
剛纔有沒有化妝,有沒有放胭脂,都忘記了..
二十五節說什麼呢?.
唯有詳細查看,詳細查看那權貝使人自由之律法的,.
並且偶爾做一次嗎?禮拜天做一次嗎?.
時常如此,這人豈不是聽了就忘,乃是實在的行出來,.
就在他所行的事上必然得福..
你注意,行道的人是有福的..
行道人是有福的,因為他詳細考察以後,.
所以背誦聖經就讓你詳細考察,明白了神的話,.
然後詳細地考察,然後把神的道行出來..
時常如此行,你就成為一個有福的人..
就像一棵樹栽在溪水旁,按說結果只葉子不枯乾,.
反正所做的竟都順利..
一個行道人,耶穌記錄在《馬太福音》第七章告訴我們說,.
一個聽了我的話又去行的人,.
好像一個聰明人把房屋建造在磐石上,.
能夠經得起風吹雨打,房子也不會倒塌,.
你的生命是堅固的,你所行的事情都是順利的,.
因為你有神的恩寵,你有人的恩寵在他裡面,阿們..
然後最後第三,要跟隨耶穌的腳蹤行,.
告訴他們要跟隨耶穌的腳蹤行..
我們已經532個人已經背完1月份到4月份的經文,.
記得4月份的經文是什麼嗎?.
你們蒙昭原是為此,因為基督也為你們受了苦,.
給你們留下了榜樣,叫你們跟隨祂的腳蹤行,.
記得嗎?一起來看這段經文,會背的一起來背了,來..
你們蒙昭原是為此,因為基督也為你們受過苦,.
給你們留下榜樣,叫你們跟隨祂的腳蹤行,.
祂並沒有犯罪,口裡也沒有詭詐,.
祂被罵不還口,受害不受,為何的話,.
只將自己交託那按公義審判人的主,阿們..
這講得很清楚,我們蒙昭原是為此,為什麼?.

$^{441}$就是要來跟隨耶穌基督的腳蹤行..
耶穌在地上三年半,加上祂在父母親的家三十三年半的時間,.
祂曾經也為我們受過苦..
耶穌受的苦多不多?.
被人逼迫,被人誣賴,甚至被人出賣,被人誣告,.
結果還被釘死在十字架上,.
祂的確為我們受過苦..
那麼祂留下的榜樣,目的是什麼?.
叫我們可以跟隨祂的腳蹤行..
因此我們要成為充滿恩寵的一個人,一定要緊緊地跟隨耶穌..
效法這位最有恩寵德神的恩寵德人的恩寵的耶穌,.
他是怎麼得到這些恩寵的?.
就是效法祂..
很多方面我們可以效法耶穌基督,.
但是基督本身親自告訴我們,.
祂最想我們學習效法祂的模式是什麼呢?.
就是效法祂的柔和謙卑..
一起講一遍來,效法耶穌的柔和謙卑..
所有美德,所有生命裡面的美德里面,.
耶穌認為最重要的就是柔和謙卑..
怎麼知道呢?.
在《馬太福音》第十一章二十九節三十節,.
耶穌講得非常清楚,.
我們都非常熟悉之前的經文說,.
當眾丹可以到我這裡來,我就使你們得安息..
接下來二十九節說什麼呢?.
我心裡柔和謙卑,你們當負我的惡,學我的樣式,.
這樣你們心裡就必得享安息..
在原文的排版應該是這樣講的,.
你們當負我的惡,學我的樣式,我心裡柔和謙卑..
這樣你們的心裡就必得享安息,.
因為我的惡是容易的,我的丹子是清神的..
你們有沒有想過,耶穌很多方面我們可以學習的,.
結果耶穌基督的愛心,祂的憐憫,.
祂的門徒們的忍耐,包容,接納,都是我們可以學習的,.
但是最重要的是,耶穌說你要學我的樣式,.
Learn of me, I am meek and humble,我是柔和謙卑..
一個柔和謙卑的人也是完全全然順服的一個人..
一個心裡沒有謙卑的人根本不會順服,.
不要說順服權柄,順服小豬黨都困難..

$^{481}$一個心裡沒有柔和謙卑的人很難順服,.
因為他心裡面想說,他們根本不會順服神,.
我幹嘛要順服神,我學歷夠高,我的經驗夠豐富,.
我幹嘛要順服神,我自己有自己的辦法可以解決這個問題..
我有心裡柔和謙卑的說,主啊,我真的需要禰給我的智慧..
所以柔和謙卑的人真的是我們需要學習的..
所以子路保羅很清楚知道,.
所以他說,你們當效法我,如同像我效法基督一樣..
如果保羅只是說,你們當效法我就停止的話,.
哇,保羅你好驕傲啊,效法你,對不對?.
他說,像我效法基督一樣,.
你們效法我的時候不要只是停留在效法我保羅這個人,.
而是效法我是如何效法耶穌基督的柔和謙卑..
所以保羅可以說是聖經裡面最傑出的使徒,.
做了最多的事情,開拓教會,帶領門徒,寫了十三卷新約的書信,.
但是他也是最謙卑柔和的人,怎麼知道呢?.
當他寫信給很驕傲,充滿驕傲自大的哥林多教會的人的時候,.
他怎麼說?.
哥林多教會第二章,第一節,第二節..
他說,弟兄們,從前我到你們那裡去,.
並沒有用高嚴大智對你們宣傳神的奧祕,.
因為我曾定了主意,在你們中間不知道別的,只知道什麼?.
耶穌基督並他定時之駕..
嘿,保羅,保羅是猶太人中最傑出的法利賽人啊!.
他是法利賽人中的法利賽人,博學多聞,.
而且他是在法利賽人的老師加瑪略大師門下的大弟子..
他可以用高嚴大智為王的話,跟哥林多教會弟兄姊妹來說話,.
他說,我已經定了主意,我沒有用高嚴大智給你們宣講神的奧祕,.
而且他知道神的奧祕最多的一個使徒,因為他被帶到三層殿裡面,.
他明白耶穌基督對教會的信義,所以才寫出這麼多的書信,.
十三卷的書信都是保羅所寫的..
他裡面充滿了神給他的奧祕真理,多麼豐富的一個人!.
他怎麼說我已經定了主意,在你們當中不知道別的,只知道什麼?.
兩件事情,耶穌基督並他時之駕..
我只知道如何效法耶穌基督的柔和謙卑..
謙卑到一種地步,自己卑微,全境順服,以至於死,且死在十字架上..
一個柔和謙卑的人,是充滿神的favor的人,.
因為神阻擋驕傲的人,喫恩典給謙卑的人..
什麼時候你驕傲,你身上就少了神的恩典,少了神的恩寵,.
更是少了別人對你的恩寵..

$^{521}$什麼時候你謙卑柔和,神就駕恩典給你..
我們教會有很多柔和謙卑的弟兄姊妹,.
知道你們當中學歷都非常高,好幾個學位,.
但是你們還是謙卑柔和地在教會跟我們牧者同文一起佩戴來服侍主..
但是我知道很多的屈牧小組長,你們都是有PhD,.
不止一個,兩個PhD,你們還是如此謙卑地來服侍弟兄姊妹,來牧養他們..
給我們的小組長牧牧們鼓掌鼓勵一下..
這是耶穌要我們學習的,你要成為充滿恩寵的人嗎?.
所以耶穌啊,我要向你學習柔和謙卑..
柔和謙卑的人是人人都喜愛的人,而且是充滿影響力的人..
那你說我怎樣柔和謙卑呢?.
你要怎麼樣?常常跟那些柔和謙卑的人在一起..
最柔和謙卑的人就是耶穌基督嘛..
要怎麼樣?常常和耶穌在一起..
我們中國不是有句話說,敬朱者赤,敬墨者黑..
你想要成為什麼樣的人,就跟這樣的人在一起..
在久了以後呢,你就被他的生命所感染嘛..
所以有人對你說,你們很有夫妻相..
那些大妹子說,太太像我不好,我像太太比較好..
夫妻相就是說,你越來越像彼此了嘛..
因為你生活在一起啊,幾十年在一起了,你肯定彼此影響..
同樣,我們要常常和耶穌在一起..
我們的生命才能夠充滿耶穌基督的恩寵,也充滿別人對我們的恩寵..
這是耶穌基督訓練門徒的法則..
耶穌花了很多的時間,這麼多人跟隨他..
晝夜禱告以後,第二天選了十二個門徒..
他怎麼訓練他們呢?.
也是我們每位訓練門徒,做導師跟學徒的很好的一個例子..
在馬可福音第三章第十四節十五節裡頭,告訴我們,.
他設立了十二個人,耶穌要他們常和自己同在,.
要差他們出去傳道,並給他們全並趕回..
耶穌把十二個人招來,要他們日夜與他們同在,.
要他們常和,不是偶爾,有時週末有空的時候禮拜天,.
乃是常和自己同在,做什麼呢?.
好讓他們能夠看見耶穌基督如何服侍人,.
看耶穌基督如何醫治病人,如何趕回,.
耶穌基督如何傳講天國的真理,.
面對試探,面對挑戰,面對逼迫,面對誤告..
他們在耶穌基督身上看到,.
榮美的耶穌,三年半的時間,日夜跟耶穌基督在一起,.

$^{561}$賜賀跟耶穌基督在一起,.
所以他們的生命就被耶穌基督的影響,.
然後耶穌纔派他們出去傳道,並且給他們全並趕回..
同樣我們訓練門徒也是如此,.
做門訓也是如此,我們黎明加油站也是如此,.
我們把這些寄宋人員叫來,.
教練陪伴他們,雖然不能天天跟他們在一起,.
常跟他們在一起,為了他們能夠從彼此身上學習,.
能夠吸收,讓他們的生命能夠繼續不斷地成長,.
並不斷地能夠充滿神的同在.阿們..
你看這些門徒跟隨耶穌久了以後,.
耶穌復活生命以後,他們就是成為小耶穌了..
所以在使徒行傳充滿了耶穌基督的門徒們所做的神蹟啟示,.
特別是在第三章五星節之後,.
彼得跟約翰走到聖殿的美門那裡,.
看到一個伸來斜腿的人,就說:.
"淨害都沒有我奉那撒冷的耶穌的名,叫你起來行走.".
不是姐姐叫弟弟,是彼得叫斜腿的人起來行走..
他就起來,又唱又跳,趕緊讚美一組..
大家拍一首哈利路亞,沒想到官府跟長老跟這些祭司們,.
因為嫉妒的緣故把他們抓起來,一直審問他們,.
發現這兩個門徒背景是目不識丁的,.
連幼兒園都沒上過,不要說小學,.
都是打魚的漁夫..
怎麼有這樣的膽量,竟然針對我們這些官府長老,.
大膽地來宣講耶穌基督和耶穌基督復活的事情?.
他們沒有其他的想法,只有一個結論,什麼結論呢?.
使徒行傳第四章十三節,.
他們看見彼得和約翰的膽量,.
又看出他們源自沒有學問的小民,.
被人藐視,看不起的小民,.
就稀奇認明他們是什麼?.
跟過耶穌的..
哇!跟過耶穌的人,別人是可以看得出來的..
雖然你是沒有任何的學歷,或者有很多的學歷,.
或者是沒有文化,或者很有文化,.
只要你長期地跟過耶穌的話,.
你不必到處宣講,我是基督徒,我是生命領導人的會員,.
你不需要講,別人在你身上就能夠看見耶穌基督的樣子..
因為一個充滿耶穌基督生命的人,.

$^{601}$是一個充滿恩寵,充滿神的恩寵的人的恩寵的人,.
他們的生命有吸引力,他們的生命有感染力,.
他們的生命有影響力,.
別人喜歡跟他在一起,想從他身上得到一些智慧,.
想從他身上得到一些幫助,.
為什麼?他因為有神的恩寵在他身上,.
因為他是跟過耶穌的人..
你不必講太多,.
你的家人,特別是你的配偶就知道,.
你的同事,跟你同工的弟兄姊妹就知道,.
到底你有沒有跟過耶穌?騙不了人的,.
不是你的口中所說的,.
那是從你的生命當中,你的一舉一動,.
你的行事為人,你的處事待人,.
是否能夠流露出耶穌基督柔和謙卑的樣子..
所以你可以告訴他們說,.
我像耶穌基督一樣,.
智慧和聖量,並神和人喜愛的心一起增長,.
而且繼續不斷地增長,.
因為你是一個充滿恩寵的人..
我們一起禱告..
謝謝主禰給我們加那個恩典,.
加那個榮幸,.
讓我們有機會認識禰,.
不但認識禰,.
主啊,禰也留下榜樣,.
讓我們可以有樣可學,.
讓我們有耶穌可以跟隨,.
不要在黑暗中摸索,.
但是清楚地來跟隨禰,.
花時間和禰在一起,.
不能透過主的崇拜,.
小主的聚會,.
做美麗的靈脩禱告,.
我們就是願意花時間跟你在一起,.
在一起久了,.
主啊,.
別人看見我們,.
沒有太多的學歷,.
沒有太多的文化,.

$^{641}$沒有太多的背景,.
卻有這樣的一個膽量,.
卻有這樣的一個智慧,.
他們不得不承認說,.
這個人是跟過耶穌的,.
主幫助我們每一個人,.
都能夠成為跟過耶穌的人,.
充滿你的恩寵,.
也充滿人對我們的喜愛..
謝謝我們的主聽我們的禱告,.
\newpage



\section{創世記 18:17-19}
\label{sec:032151}
\textbf{何日再歸家?}~~~~李業基牧師
\newline
\newline
連結: \href{http://www.church.com.hk/acms/content.asp?site=cdc\&op=show\&type=product\&code=032151\&layout=sermon}{\texttt{032151}} ~~~~ 語音日期: 2026-05-17
\newline
\newline
\hyperref[sec:032150]{\small{< < < PREV SERMON < < <}}
~
\hyperref[sec:index_chronic]{\small{[返順時目]}}
~
\hyperref[sec:032152]{\small{> > > NEXT SERMON > > >}}
\newline
\newline
創世記 18:17-19
\newline
\begin{longtable}{cl}
\hline
\hline
章節 & 經文 (和合本修訂版)\\
\hline
18:17 & \begin{tabularx}{0.7\textwidth}{X} 耶和華說：「我所要做的事豈可瞞著亞伯拉罕呢？ \end{tabularx} \\ \\ \relax
18:18 & \begin{tabularx}{0.7\textwidth}{X} 亞伯拉罕必要成為強大的國；地上的萬國都必因他得福。 \end{tabularx} \\ \\ \relax
18:19 & \begin{tabularx}{0.7\textwidth}{X} 我揀選他，為要叫他命令他的子孫和後代家屬遵行耶和華的道，秉公行義，使耶和華所應許亞伯拉罕的話都實現了。」 \end{tabularx} \\ \\
[1ex]
\hline
\hline
\end{longtable}
$^{1}$我上星期母親節在聖地安索已經說過.
開始時玩一個有獎搶答遊戲.
今天想再玩一次.
但如果你在網上已經聽過這篇說到重溫.
就麻煩不要參與.
因為你聽過就知道答案.
但如果沒有聽過就可以參加.
這是一個搶答遊戲.
第一個舉手而且正確答案的人就獲得獎品.
這個Irene師母親手作一份.
不是,這杯是我的.
好,這個遊戲是這樣的.
這首詩大家很熟悉.
慈母手中線,柔子身上衣.
臨行脈蜂,以孔慈慈歸.
最後兩句是.
喂喂喂,阿Peter先舉手,不要吵.
你是否脫牙垂燕,到你了.
是垂燕,Peter先上來.
多謝,快與準.
串草心是甚麼意思呢.
小草指子女,心有相同含義.
即是指敬幹,亦是指子女的心意.
三春暉,農曆正月叫孟春.
二月叫重春,三月叫季春.
合稱三春.
暉即是陽光.
形容母愛好像春天的溫暖.
和海的陽光照耀子女.
這首詩是唐代詩人孟郊的作品.
唐代詩人孟郊早年漂泊無依.
一生可以說是貧困潦倒.
但五十歲才獲得一個卑微的職位.
於是接了母親一起居住.
詩人飽嘗世態炎涼.
因而更加體恤到親情的可貴.
我們說到系列的問題是.
如何讓你的家成為培育門徒的中心.
這是我們來到尾聲.
但其實對於很多我們的子女已經長大了.

$^{41}$在外地.
有些像孟郊的母親的人來說.
心裡面的問題可能反而是.
何日再歸家呢.
昨天我的兒子Sophus.
讀完大學一年就回來了.
最近幾個星期.
我發現很多親戚朋友不同人.
至少問了我十幾次.
Sophus何時回來.
何日再歸家.
聖經裡面.
浪子的故事.
有部份我想是到了某個年紀.
我開始領會.
但何日再歸家這句說話.
當然是家中父母所想.
亦都是那位柔子心中所問.
人成為柔子其實各有原因.
柔子去踏上這個旅途.
或者是要闖天下.
為自己幹一番事業.
找些家中找到的東西.
但可能亦都是被迫離家.
因為他家中再沒有容身之處.
甚至可能被人驅逐出境.
被要求離開.
以色列的歷史.
就是一個柔子的故事.
耶和華賜這個名.
以色列給雅各.
的時候雅各被迫離家.
因為他的孿生大佬要追殺他.
雅各這個故事.
以色列的故事.
其實是聖經裡面反覆重演的橋段.
當饑荒迫使雅各和他全家.
七十多人 我想不止.
到了埃及的時候.
以色列再次成為異鄉人.

$^{81}$雅各再次成為異鄉人.
摩西逃離埃及的時候.
又是重演了這個柔子的故事.
當聖經講到第一次聖殿被毀的時候.
以色列人被流放到巴比倫七十年.
第二個聖殿公元七十年.
聖殿被毀的時候.
以色列人散居各國.
所以柔子的故事.
其實就是以色列歷史的縮影.
何日再歸家.
就是貫徹始終的核心問題.
聖經一開始就跟我們講柔子的故事.
其實是全人類的故事.
人類始終亞當夏娃被趕出伊甸園.
成了柔子.
所以其實我們每一個人心底的骨子裡面.
都在問這條問題.
都在追問這條問題.
何日再歸家.
聖經告訴我們.
只是當另一位柔子.
離開天家的耶穌基督.
來到這個世上.
透過他的工程.
透過他的工作.
人才能夠回到我們真正屬於的家.
在來自中國大陸很多海外知識份子裡面.
其實有一本書是九零年代中.
就已經廣為流傳的.
一本信仰入門書叫做《柔子淫》.
作者的筆名是李程.
李程其實是馮炳程牧師的筆名.
如果大家有朋友.
同事.
其實可以屬於一個中國知識份子的背景.
我推介大家自己看看和送那本書.
現在還是第四版.
超過一百萬本.
非常受歡迎.

$^{121}$非常有同感的一本書.
我們國語事公其實會點是一個發展重點.
這就是為什麼我們邀請了.
盛曉梅姐妹和泥土音樂.
月底來到我們當中.
去搞這個幸福音樂的聚會.
但如果我們想接觸他們.
不只是搞一些聚會.
搞一些活動.
你真的要明白他們的思維模式.
用他們的心境.
用他們的言語和他們說話.
其實全福音萬變不離其宗.
第一步就是要建立一個共同的話題.
我在波士頓讀大學的時候.
組織過幾次為香港留學生而設計的報道會.
那時候香港留學生遠比現在多.
何日再歸家呢.
令到很多人能夠產生共鳴的一個話題.
何日再歸家呢.
從前那些人考得功名.
衣錦還鄉就回家了.
後來那些人來到錦山.
那些台山人賺到錢.
回鄉下就賣田賣地.
我記得我小時候70年代.
很多人剛剛出了香港.
賺到錢就會抬電視機.
飯煲回鄉下.
不過也不一定是回鄉下.
孟交考到進士將母親接回家一起住.
其實家人去到那裡.
家就在那裡.
家和家鄉就有這個分別.
其實當人問何日再歸家.
想了一段時候一定會想.
究竟何處是吾家呢.
摩西在米殿人的生活當中是40年.
當神在焚燒的荊棘叢果當中.
吩咐摩西回埃及的時候.

$^{161}$摩西不願意.
摩西一早已經視米殿為他的家.
出埃及的以色列人.
在曠野中激怒了和神無數次.
才口服心服.
他們才明白他家已經不是埃及.
而是在他前面的應許之地.
我想想我家在那裡.
我祖籍是台山.
但從來沒有去過.
我香港長大的.
18歲就離開了.
其實已經很久沒有說回香港了.
只是去香港.
但是在波士頓生活了30年.
那裡仍然有很多親戚.
因此雖然我的父母都和我一起搬到這邊.
但當Irene和我去波士頓的時候.
我們都是說我們回波士頓.
剛剛這個星期有些事.
要探望一個親戚.
就回了波士頓兩天.
當我寫《落霞》的時候.
《霧耳》當然心裡想著香港.
但現在我唱的時候.
我懷疑心裡想的是波士頓.
其實本地也是這樣.
大家記得去年四月.
我去Millpeters幾個月.
聖地安索的弟兄姊妹就開始問.
什麼時候回來.
現在剛好相反.
去到Millpeters才說.
很久沒見了.
回來啦.
經常一個遊子人生的轉捩點.
就是從何日再歸家.
當中醒覺.
轉移到問何處是吾家.
浪子回頭這個故事.

$^{201}$他醒覺那一刻的時候.
其實我想他不是在想何日再歸家.
哪有面子.
還有面子回去見爸爸嗎.
但是他反而可能是在想.
何處是吾家呢.
面前那個豬欄.
是不是真的要成為他人生的家呢.
亞伯拉罕和薩拉的老家在哪裡.
在一座叫做吾爾的泛華大城.
位於巴比倫東南140英里.
今天我們叫做有法拉底河的西邊.
大概十英里左右.
亞伯拉罕為什麼要離開他老家吾爾呢.
搬家去到迦南呢.
如果你看牛羊的數目.
和瀑人的數目來說.
亞伯拉罕當時是算為極端富有的.
我在哥倫康威爾神學院的院長.
也是舊約教授.
Dr.Walter Kaiser就說.
亞伯拉罕是當時的洛克菲勒.
亞伯拉罕離開吾爾.
是因為神呼召他去創立一個大國.
神應許亞伯蘭.
那時候還是.
他的新的家園將是留乃與物之地.
那裡亞伯蘭還會得到一個兒子.
而且從兒子的後裔當中.
建立一個民族出來.
亞伯拉罕和他的妻子薩拉.
有沒有懷念他們吾爾的舊居呢.
他們有沒有想過何日再歸家呢.
我猜也有的.
你想想將迦南地變成為一個家園.
其實不是這麼容易的事情.
如果他們不是經常想念家鄉.
就不用派一個年紀老邁的僕人.
回家鄉為自己的兒子找個老婆.
但我想他們也因此.

$^{241}$一定有問過何處是吾家這個問題.
是不是真的可以捱得住.
在迦南地繼續住下去.
但他們仍然願意留在這個新的應許之地.
建立他們的家園.
是因為他們相信神對他們的應許.
當以色列人從摩西帶領他們口中災.
出埃及聽到這個故事的時候.
他們其實都自身在一個相似的處境當中.
當然他們不像亞伯拉罕和易萬富翁.
他們其實背景就是埃及的奴隸.
但是神對亞伯拉罕所作的應許.
也是對那些出埃及的以色列人的應許.
你想想那些以色列人有沒有懷念過他們住了這麼多年的埃及呢.
他們有沒有心想著埃及的時候問何日再歸家呢.
我想也有的.
但他們願意繼續向前走.
因為他們知道埃及最終不是何處是吾家這個問題的答案.
因為神已經應許了他們一個新的永久的家居.
所以無論是亞伯拉罕或是離開埃及的迦南.
進入迦南的以色列人來說.
何日可成家比何日再歸家是一個更加重要的問題.
神曾經應許亞伯拉罕的子孫.
會好像天空的星和海上的沙這麼多.
但日復一日 年復一年.
這個應許還沒有兌現.
在創世記十八章.
神就在以撒出生之前的一年.
開始為亞伯拉罕解答這條問題.
他何日可成家.
亞伯拉罕其實是七十五歲離開哈蘭的.
亞伯拉罕九十九歲那年.
主和另外兩位天使路過他家.
可能亞伯拉罕就是想弄清楚.
究竟何日可成家.
所以亞伯拉罕勉強他們.
懇請他們留下來吃一頓飯.
果然他們願意.
所以在這個席中.
主向亞伯拉罕宣告.

$^{281}$一年之後.
即是亞伯拉罕離開哈蘭二十五週年.
或者說亞伯拉罕建國二十五週年.
亞伯拉罕將要在撒拉身上得到一個兒子.
因為預備這篇講道.
原本是母親節.
所以我特別留意故事當中的母親.
當撒拉聽到自己.
下年要生孩子的時候.
她有什麼反應.
亞伯拉罕和撒拉都年紀老邁.
撒拉的月經已停了.
撒拉心裡竊笑說.
我已衰老.
我的主也老了.
怎能有這事呢.
亞伯拉罕看到撒拉心裡竊笑.
對亞伯拉罕說.
撒拉為什麼竊笑.
說我已年老.
果真能生育嗎.
夜禍說.
豈有難成的事嗎.
到了鎖定的時候.
我必回到你這裡.
明年這時候.
撒拉會生一個兒子.
撒拉害怕了.
就不承認了.
沒有我沒有笑.
亞伯拉罕對撒拉說.
不是的.
你真的笑了.
當你一世人在等一件事發生.
等了二十五年.
都一直不斷思索.
何日可成家.
但始終得不到答案.
到這個事情開始實現的時候.
我想你都大概會和撒拉一樣反應.

$^{321}$我?.
撒拉?.
LOL.
這個就是信心的考驗.
很多時候相信.
神在他慎自己的時間.
勝就他的應他.
甚至是需要神職.
神本身是全能全知的.
但是他要回應人何日可成家這個問題.
首先未必給答案.
有可能是時候.
神亦未必在年曆表畫過日期.
我聽聞某間公司.
最近要離職.
已經說好.
未來星期三不知會有多少千人.
但神有時候不是這樣做的.
他可以透過一些先兆.
透過一些先決條件.
告訴他何時成就他的工程.
這個就是我們今天的經文.
創世記十八章十七至十九節.
耶和華說.
我所要做的事豈可瞞著亞伯拉罕呢.
亞伯拉罕必要成為強大的國.
地上的萬國都必因他得福.
我揀選他.
為要叫他命令他的子孫和後代家屬.
遵行耶和華的道.
秉公行義.
使耶和華所應許亞伯拉罕的話都實而了.
何日可成家.
總有一日.
總有一日亞伯拉罕會成為強大的國.
事實上地上的萬國都必因亞伯拉罕得福.
明年這個時候.
就在他立國二十五週年之際.
神就要賜一個兒子給亞伯拉罕.
以此作為未來應許的一個down payment.

$^{361}$一個保證.
但是神對亞伯拉罕的計劃.
不是只要建立一個小家庭.
只給他一個兒子.
要他建立一個國家.
所以何日可成家這個國家呢.
三件事情之後.
天火焚成之後.
圖理滿門之後.
應許實現之後.
何日可成家.
天火焚成之後.
十七次其實是一個反問句.
我要做的事情不可以瞞著亞伯拉罕.
即是說他要告訴他.
神要實現一件事情之前.
就要向亞伯拉罕和他的子孫.
首先講定他的計劃.
所以當事情真的發生了之後.
他們就知道他們相信這個主.
是一個全能全知的神.
擁有絕對的歷史主權.
所以之後耶和華就和亞伯拉罕解釋.
他將要審判索多瑪和俄摩拉.
我想亞伯拉罕也知道.
索多瑪和俄摩拉.
確實是邪惡到無可辯護.
但是他的侄子在那裡.
亞伯拉罕就用神的公義向他懇求.
倘若那裡有人其實是罪不至死的.
這樣將他們和該死的人一同毀滅.
神你好像有點不公義.
結果耶和華神聽亞伯拉罕的請求.
即使有十個不該死的人.
他都會饒恕那裡的人.
會不滅亡他們.
這個是神的本性.
即使有很少人能夠得救的時候.
神仍然給予時間.
但是神也會洞悉過去,現在,未來,歷史.

$^{401}$當一個民族已經到了一個無可救藥的時候.
神公義的聖情就出事.
他必須施行審判.
其實索多瑪和俄摩拉的故事.
是後來那些以色列人進入迦南地的時候.
一個預告.
迦南地其實有很多邪惡的成癮.
都會遭遇這件事.
這件事是一個歷史先例.
如果沒有人住在那裡.
迦南地很快就會變成荒涼.
因此當亞伯拉罕的子孫在埃及回來之前.
神容許迦南地成癮的罪惡.
容忍他們.
但是時機成熟的時候.
仍然需要潔淨這片土地.
所以當以色列人進入迦南地的時候.
就是神潔淨土地工程的一個部分.
是神潔淨的火.
神將亞伯拉罕的子孫從埃及領出來.
為他們建立一個嶄新而且聖潔的國度.
不是要他們進入迦南地.
融入一個異教文化當中.
所以可以這樣說.
以色列人就算進入了.
他們已經身處應許之地.
他們還沒有成為一個聖潔的國度.
在這個地方能夠再一次成為他們的家之前.
他們必須清除一切不屬於神的人.
不屬於神的東西.
保羅在羅馬書第九章二十七節.
引述以賽亞的話.
關於以色列人.
以賽亞喊著.
雖然以色列人多如海沙.
得救的將是剩下的餘數.
教會也是這樣的.
特別是那些夢想成為Mega Church的教會.
不要以為我們有一個宏偉的建築物.
人山人海.

$^{441}$我們就一定是一個被神閱立的教會.
何謂教會.
它是一個Ecclesia.
被召出來屬於神的群體.
在我們成為神的家之前.
我們必須成為君尊的祭司.
聖潔的國度.
如果我們不把我們自己分別為聖.
最終還是會變成.
所多瑪俄摩拉.
和迦南帝國的國的誠邑.
一樣下上.
果然.
所多瑪俄摩拉.
是可以存活的.
是配得存活的.
十個都不夠.
亞伯拉罕的侄子羅德是其中一位.
所以當兩位天使來到羅德面前.
祝福羅德快點離開.
不要再回頭看.
羅德確實沒有再回頭看.
保存了自己的姓名.
但是不幸的是.
羅德的妻子回頭一看.
就變成了嚴處.
我因為這篇講到母親節.
所以值得思想先前的撒拉一樣.
想想羅德的女兒.
即是他的妻子.
撒拉信心不足.
但至少她面對主的責備時.
都知道害怕.
但是.
所多瑪俄摩拉.
已經天火焚成了.
已經爆炸了.
背部可能已經很熱.
烈浪已經來到身.
但是為什麼羅德的妻子.

$^{481}$仍然不聽天使的話.
仍然要往後望呢.
其實就算你在羅德的一家裡.
都救不了她.
神是充滿恩典的.
但是如果有人嘗過天恩的滋味之後.
仍然是回頭望的時候.
等待的.
就只有審判.
或者羅德的妻子心裡真的在想.
何日再歸家.
或者對她來說.
所多瑪和俄摩拉.
就是一早已經是她的家了.
但是如果神已經向你顯明.
一個新的家.
永遠的家.
你還是眷戀著.
一個神已經在毀滅的地方.
不停地回頭張望的時候.
你就要小心.
恐怕你和羅德的妻子.
有同樣的結局.
這就是摩西說這個故事.
要提醒那些出埃及的以色列人的道理.
這個也是今天很多教會的記錄圖.
我們需要明白的真理.
亞伯拉罕的祖孫.
能夠成為一個國家之前.
他們必須在自己的家居裡面.
每一個人家中.
使家人成為一個遵行神的道的門徒.
何日可成家.
圖你滿門之後.
我選擇他.
為要叫他命令他的子孫和後代家屬.
行耶和華的道 秉公行義.
你想想這個其實不是一個建立.
強大國家的常規.
古代世界建立一個強大國家是怎樣做呢.

$^{521}$好像巴比倫那樣.
打人 組織一個強大的軍隊去征服別人.
同化別人.
或者至少好像埃及那樣.
將人民的奴役為國家的經濟引擎.
但是神對亞伯拉罕的子孫.
有一個不同的計劃.
我們可能叫做organic growth.
有機生長.
一個一個的門徒.
你家裡有多少人.
一個一個的門徒去建立這個國家.
毫無疑問.
我想今天有很多創業精神的founder.
看到這個計劃肯定會搖頭.
No.
多少時間才夠 我一輩子都不夠.
確實你沒有足夠時間.
但是神有.
其他那些搞工程 管理的人.
現在搞AI的人.
看到這個計劃可能會說.
這個計劃要成功.
除非有神蹟.
起步都不容易.
撒拉今年貴庚.
但是大家看看結果.
古代的巴比倫 古代的埃及現在在哪裡.
亞伯拉罕的後裔仍然在.
而且不只是一個民族.
而在萬國萬民當中都有亞伯拉罕的後裔.
為什麼.
因為亞伯拉罕後裔.
遵行耶和華的道.
秉公行義.
所以對一個建立一個國家來說.
沒有一樣東西比建立國家的品格.
國民的品格和道德標準是更加重要的.
正因為這個緣故.
神先將培育後裔成為亞伯拉罕成家的核心策略.

$^{561}$何日可成家.
應許實現之後.
雖然不是透過一個世界的常規來達成的.
但是神仍然使亞伯拉罕成為一個強大的國度.
神的計劃表面上是沒有什麼可能的.
但是在恰當時候.
必然是成就的.
甚至其他民族也都因為亞伯拉罕而蒙福.
正如在十九節裡面.
神對亞伯拉罕說.
耶和華必為亞伯拉罕成就他所應許的.
這個不是意味著亞伯拉罕的後人.
總是如美無缺地去遵守這個律法.
但是即使神的計劃是涉及到容易犯錯的人的時候.
神的計劃都不會失敗的.
在母親節說世上的母親是沒有完美的.
聽來可能是有點不禮貌.
但是是真事來的.
的確世上沒有一個完美的母親.
就好像我們說某個家庭是完美一樣.
或者某個地方教會是完美一樣.
但是母愛的天性仍然是值得被表揚的.
因為這個天性是天賦到來的.
是值得讚揚的.
母親也是人.
當母親因為無知或者各樣感情的預估.
做錯事的時候.
我們是值得用同理心去接納的.
當我們去想羅德的兩個女兒的古怪行徑的時候.
可能需要明白一個背景的態度.
當羅德的兩個女兒發現他媽媽變成了鹽柱的時候.
他們就切法灌醉了羅德的爸爸.
然後跟羅德同寢.
生了摩訝和便亞米.
摩訝人和阿曼人的祖先.
後來聖經經常提到這兩個民族.
大衛王的曾祖母羅德就是摩訝人.
倘若羅德的妻子沒有違抗天使的命令.
沒有回頭望的時候.
這個可能是可以避免的事情.

$^{601}$但是即使人會犯錯.
神的計劃都會有一個辦法去承受的.
許多民族的確是因為阿巴拉罕蒙福.
其中也包括一個近親.
即是羅德而出的民族.
總結時間到了.
如果我們明白神在耿古永遠當中.
已經為我們預備了一個家.
有一個新天新地的時候.
我們何日可以再歸家.
何日可成家.
其實是同一條問題的兩個不同的問法.
其實都有同樣的答案.
有三個的答案.
天火焚成之後.
圖你滿門之後.
應許實現之後.
那天是哪天我們不知道.
但是我們知道那天必然來臨.
知道的時候首要任務是什麼呢.
聖經中的預言聽來是很難以置信的.
但是我們千萬不要像撒拉那樣痴笑神的話語.
第一樣就是要信.
第二不要回頭望.
不要像羅德的妻子那樣.
神將以色列人從巴比倫和埃及領出來.
不是要他們做回迦南人.
不是他們順應世界的誘惑.
不單是在外面.
在家中.
在教會裡面.
我們都需要對應這個世界給予我們的試探.
我們需要忘記背後努力面前.
不要只是思想地上的事.
而是真的有思念天上的事.
其實你來自何處呢.
其實不重要的.
何處是吾家.
歸根究底.
就是你能夠將下一代造就成為門徒的地方.

$^{641}$那個就是你的家.
一個由此人生的轉捩點.
就是從問何日在貴家當中醒覺.
轉移到問何處是吾家.
我們在聖尼安卓分堂有一個口號.
Be a Saint Neanderthal.
Neanderthal其實是Leo加Angel兩個字.
即是Lion Man.
所以Be a Saint Neanderthal.
我們的意思就是.
做一個獅子心.
人子身的聖徒.
當然這是提醒.
神差派我們到那裡是要直跟這個社區當中.
我們的意象不是只是將我們的家居.
或者是我們的教會.
變成一個門徒訓練中心.
我們也是要將我們的Hometown.
我們的家鄉.
成為一個門徒的訓練中心.
所以下季我們就有一個講道系列.
就是按照這個題目.
是藉著約書亞記.
士師記和盧德記去思想.
以色列人進入應許之地的歷史.
反思我們的當地的使命.
差不多每一個星期日.
我禱告的時候.
當然神給我們有一個建築物.
我們是感恩的.
還沒有拿到這個的時候.
你拿到日本教會的時候.
我們就一直在等待這個所謂應許之地.
但其實在教堂裡面也好.
在教堂外面也好.
總之我們是宣揚神的福音.
我們是在教導門徒的.
那個就已經是神的家.
因為教會健康的治標.
最終不是有多少人進來的.

$^{681}$而是我們有多少人猜出去.
所以第三.
我們必須盡我們的所能去開枝散葉.
圖利滿門.
必須按照主的道去養育我們的子女.
透過門徒訓練建立神的國度.
其實是很緩慢的事情.
所謂Organic Growth.
和征服一個國家不同.
門徒訓練是一代傳一代的長遠計劃.
但當我們一代又一代去盤約神的指紋的時候.
好像伊斯蘭一樣.
最終會成為一個偉大的民族.
所以最後我們要記住的事情是什麼.
就是即使我們失敗.
即使我們有時候不是做到那麼完美.
實際上他的計劃.
他必然會成就的.
何日再歸家.
何處是吾家.
何日可成家.
其實都是同一條問題.
應許實現之後.
神要保守祂自己的信譽.
神要保守祂自己的榮耀.
所以祂的計劃必然成就.
讓我們低頭禱告.
你招聚我們成為一家人.
你招聚我們在這個世界當中走出來.
成為你的教會.
成為你被揀選的指紋.
願我們都是正宗我們的身份.
我們這個家人的身份.
我們懂得怎樣去互相去愛護我們弟兄姊妹.
亦懂得將我們的家園.
我們所住的地方.
我們就算在Mt.Peters在南灣這裡.
仍然有很多我們失散了.
我們出世就已經失散了的家人.
我們要找回他.

$^{721}$帶回來神的家.
求神你給我們能力.
在這裡成就實力的計劃.
禱告 奉主耶穌的聖名 阿們.
\newpage



\section{出埃及記 20:1-26}
\label{sec:032152}
\textbf{石板}~~~~陳崇基牧師
\newline
\newline
連結: \href{http://www.church.com.hk/acms/content.asp?site=cdc\&op=show\&type=product\&code=032152\&layout=sermon}{\texttt{032152}} ~~~~ 語音日期: 2026-05-17
\newline
\newline
\hyperref[sec:032151]{\small{< < < PREV SERMON < < <}}
~
\hyperref[sec:index_chronic]{\small{[返順時目]}}
~
\hyperref[sec:032153]{\small{> > > NEXT SERMON > > >}}
\newline
\newline
出埃及記 20:1-26
\newline
\begin{longtable}{cl}
\hline
\hline
章節 & 經文 (和合本修訂版)\\
\hline
20:1 & \begin{tabularx}{0.7\textwidth}{X} 神吩咐這一切的話，說： \end{tabularx} \\ \\ \relax
20:2 & \begin{tabularx}{0.7\textwidth}{X} 「我是耶和華－你的神，曾將你從埃及地為奴之家領出來。 \end{tabularx} \\ \\ \relax
20:3 & \begin{tabularx}{0.7\textwidth}{X} 「除了我以外，你不可有別的神。 \end{tabularx} \\ \\ \relax
20:4 & \begin{tabularx}{0.7\textwidth}{X} 「不可為自己雕刻偶像，也不可做甚麼形像，彷彿上天、下地和地底下水中的百物。 \end{tabularx} \\ \\ \relax
20:5 & \begin{tabularx}{0.7\textwidth}{X} 不可跪拜那些像，也不可事奉它們，因為我耶和華---你的神是忌邪的神。恨我的，我必懲罰他們的罪，自父及子，直到三、四代； \end{tabularx} \\ \\ \relax
20:6 & \begin{tabularx}{0.7\textwidth}{X} 愛我，守我誡命的，我必向他們施慈愛，直到千代。 \end{tabularx} \\ \\ \relax
20:7 & \begin{tabularx}{0.7\textwidth}{X} 「不可妄稱耶和華－你神的名，因為妄稱耶和華名的，耶和華必不以他為無罪。 \end{tabularx} \\ \\ \relax
20:8 & \begin{tabularx}{0.7\textwidth}{X} 「當記念安息日，守為聖日。 \end{tabularx} \\ \\ \relax
20:9 & \begin{tabularx}{0.7\textwidth}{X} 六日要勞碌做你一切的工， \end{tabularx} \\ \\ \relax
20:10 & \begin{tabularx}{0.7\textwidth}{X} 但第七日是向耶和華---你的神當守的安息日。這一日你和你的兒女、奴僕、婢女、牲畜，以及你城裡寄居的客旅，都不可做任何的工。 \end{tabularx} \\ \\ \relax
20:11 & \begin{tabularx}{0.7\textwidth}{X} 因為六日之內，耶和華造天、地、海和其中的萬物，第七日就安息了；所以耶和華賜福與安息日，定為聖日。 \end{tabularx} \\ \\ \relax
20:12 & \begin{tabularx}{0.7\textwidth}{X} 「當孝敬父母，使你的日子在耶和華－你神所賜你的地上得以長久。 \end{tabularx} \\ \\ \relax
20:13 & \begin{tabularx}{0.7\textwidth}{X} 「不可殺人。 \end{tabularx} \\ \\ \relax
20:14 & \begin{tabularx}{0.7\textwidth}{X} 「不可姦淫。 \end{tabularx} \\ \\ \relax
20:15 & \begin{tabularx}{0.7\textwidth}{X} 「不可偷盜。 \end{tabularx} \\ \\ \relax
20:16 & \begin{tabularx}{0.7\textwidth}{X} 「不可做假見證陷害你的鄰舍。 \end{tabularx} \\ \\ \relax
20:17 & \begin{tabularx}{0.7\textwidth}{X} 「不可貪戀你鄰舍的房屋；不可貪戀你鄰舍的妻子、奴僕、婢女、牛驢，以及他一切所有的。」 \end{tabularx} \\ \\ \relax
20:18 & \begin{tabularx}{0.7\textwidth}{X} 眾百姓見雷轟、閃電、角聲、山上冒煙，百姓看見就都戰抖，遠遠站著。 \end{tabularx} \\ \\ \relax
20:19 & \begin{tabularx}{0.7\textwidth}{X} 他們對摩西說：「請你向我們說話，我們必聽；不要讓神向我們說話，免得我們死亡。」 \end{tabularx} \\ \\ \relax
20:20 & \begin{tabularx}{0.7\textwidth}{X} 摩西對百姓說：「不要害怕；因為神降臨是要考驗你們，要你們敬畏他，不致犯罪。」 \end{tabularx} \\ \\ \relax
20:21 & \begin{tabularx}{0.7\textwidth}{X} 於是百姓遠遠站著，但摩西卻挨近神所在的幽暗中。 \end{tabularx} \\ \\ \relax
20:22 & \begin{tabularx}{0.7\textwidth}{X} 耶和華對摩西說：「你要向以色列人這樣說：『你們親自看見我從天上向你們說話了。 \end{tabularx} \\ \\ \relax
20:23 & \begin{tabularx}{0.7\textwidth}{X} 你們不可為我製造偶像，不可為自己造任何金銀的神像。 \end{tabularx} \\ \\ \relax
20:24 & \begin{tabularx}{0.7\textwidth}{X} 你要為我築一座土壇，在上面獻牛羊為燔祭和平安祭。凡在我叫你記念我名的地方，我必到那裡賜福給你。 \end{tabularx} \\ \\ \relax
20:25 & \begin{tabularx}{0.7\textwidth}{X} 你若為我築一座石壇，不可用鑿過的石頭，因為你在石頭上動了工具，就使壇污穢了。 \end{tabularx} \\ \\ \relax
20:26 & \begin{tabularx}{0.7\textwidth}{X} 你不可用臺階上我的壇，免得露出你的下體來。』」 \end{tabularx} \\ \\
[1ex]
\hline
\hline
\end{longtable}
$^{1}$好,今次我們繼續講這個罌.
如果你新來的朋友,你未有的話.
等下出去的時候,我們送一個給你,好嗎?.
一個玻璃罌,裡面有不同的物件.
就是我每一次講到新的一年.
我們十週年這個系列的時候.
我們就會從聖經由頭到尾.
我們看看這位的上帝是一個怎樣的上帝.
你記得我第一次講的時候.
今次是第三次了.
第一次講的時候,我說上帝是一個怎樣的上帝?.
God on a mission.
上帝是由頭到尾.
祂不單止創造了這個世界.
知道人類是叛逆祂也好.
但是祂沒有放棄過這個世界.
祂繼續是on a mission去拯救這個世界.
這個第一次講的時候就是說.
這個世界這麼瘋狂.
這麼不濟,這麼破碎.
上帝做了什麼?.
沒有放棄,祂就召喚了Abraham出來.
召喚了亞伯拉罕出來.
藉著亞伯拉罕,以撒雅各.
成為一個家,成為一個國.
成為一個子民.
這個上帝的子民就是在上帝的計劃裡面.
去做上帝的工具.
以至上帝的祝福能夠流向萬國.
這是第一次.
上一次我們講到族長的故事.
就是說這些族長,亞伯拉罕,以撒雅各.
他們不是什麼英雄人物.
我們不需要將他們美化.
他們都是在掙扎,跌又碰的裡面跟隨上帝.
去承接這個呼召.
拿起這個使命.
去使萬國能夠得福.
這個是整本《創世記》的結束.
《創世記》就是告訴我們.

$^{41}$上帝藉著這些弱小的人.
他要去成就他自己的心意.
所以我們每一個上帝的子民.
我們有一個identity.
一個名分.
這個名分就是上帝賜給我們.
作為他自己的people.
我們都是on the mission.
因為上帝就是on the mission.
這次我們繼續trace整個聖經的大故事.
完了《創世記》就到出《埃及紀》.
你記得《創世記》最後的時候發生什麼事嗎?.
雅各他們一家去了哪裡?.
去了埃及,是嗎?.
去了埃及.
然後在埃及他們一代一代受苦.
然後上帝聽見他們的聲音叫了摩西出來.
然後拯救他們.
用恩典拯救了他們.
他們就能夠出埃及.
經過曠野去到西大山.
他們真正成為了上帝的子民.
這個就是出《埃及紀》20章所記載的.
講到石板.
我的主題是石板.
你記得摩西拿著兩塊石板從山那裡下來,是嗎?.
記得嗎?.
那塊石板大概是雕著什麼?.
雕著十戒,是嗎?.
就是上帝和他的子民正式來到立約的document來的.
所以有兩塊,是嗎?.
一塊是上帝的.
一塊是上帝的子民的.
大家互相去到見證的.
所以講到石板的時候.
今天我們就拿一些小石子走,好嗎?.
上一次我們有沙.
提到阿布拉罕的沙土.
就是他的後代多如海沙的那個認許.
然後上一次我們講到族長的故事.

$^{81}$我們拿起這個item.
就是什麼?.
那顆puzzle.
我們每個都是有破碎的碎片.
是上帝幫我們砌在一起.
以致能夠成為一個美麗的圖畫.
今次我們拿什麼呢?就是小石子.
為什麼拿小石子呢?.
因為講起律法.
上帝的子民能夠成為他自己的people的時候.
上帝就賜給他們石板.
這個石板就記載著上帝.
去引導他的子民給他們的law.
給他們的律法.
現代人我們需要先明白什麼叫做律法.
舊約裡面的律法.
law of the old testament.
因為我們想起law的時候.
法律我們就想起一些不好的東西,是嗎?.
是嗎?.
我們通常經歷所謂的law,法律.
我們就覺得心裡面不是很爽的.
兩個禮拜前我就去某個地方.
我就在街上找了一個位置.
其實我是在high park的.
我就在那裡找了一條街有一個位置.
繞來繞去.
終於找到一個位置,多好.
但是街的兩旁就有一些橙色的牌.
掛在樹上.
大家在芝加哥本地的都知道,是嗎?.
city好陰濕.
我已經繞了幾次.
這邊就no parking on monday.
對面就很陰濕.
no parking on tuesday,沒錯.
no parking on monday, no parking on tuesday.
我找位置的那天就是monday.
我找到一個空位.
我已經看清楚了.

$^{121}$no parking on monday.
no parking on tuesday.
繞來繞去有一個U turn,然後就park了那個位置.
我說這次沒死,這個免費的位置.
因為no parking on monday.
我park的那個位置應該在哪裡?.
exactly.
我park了很安心,我說這次真的沒死.
幾個小時之後.
我回到去就拿了一張橙色的東西.
掛在玻璃窗那裡.
豈有此理,原來我出事了.
原來出事了,因為我繞來繞去的時候我亂了.
我就park了no parking on monday那邊.
我就以為我park了什麼?.
no parking on tuesday那邊.
但我找到一個位置明明是no parking on tuesday那邊.
為什麼我會park了no parking on monday那邊呢?.
所以罵人就60元.
其實不是CCTV陰濕的,是我自己大意.
我相信大家也受過這些罵,對吧?.
你心裡會覺得很大聲,經常受罵.
我心同感受.
你的反應可能就是真的很不公平.
我良民一個,我平時不犯法.
我park了位置你就罰我60元.
我對法律的感覺也不是很好.
因為通常我想起法律的時候就是懲罰.
或者法律就是陷阱.
所以有時候我們把這個concept.
我們就去了解舊約的律法.
我們回到舊約裡面有很多很長篇的律法.
我們的眼光就是覺得.
上帝當然好像那些交通警.
有沒有見過那些?那些city worker.
拿著部機,穿著衣服,然後走來走去.
他就是想抓你.
你心裡面很不滿.
但我們心裡面的律法就是這樣的.
就是懲罰.

$^{161}$但是我想今天告訴你.
在舊約裡面講的律法不是這樣的.
我們有那些舊的概念.
我們一定要把它delete的.
我們經常說我們不在律法之下.
我們現在在恩典之下.
你有沒有聽過這句話?.
有啊,從小到大在教會裡面都是聽一樣東西的.
we are not under the law.
we are under grace now.
以前就是律法的時代.
現在就是恩典的時代.
現在我們不講律法了,很free的.
no no,這個是舊的概念.
這個是不明白聖經,解錯聖經的一種概念.
在舊約的律法裡面講的不是一種懲罰.
它是講到上帝給我們的恩典.
所以有一個誤解.
我們就不要像畫面上那樣講.
舊約那個神就很像交通警.
很嚴的.
你一犯法,他就罰你六十塊.
罰你六百塊.
就是引頭落地.
新約那個就不是了.
新約那個講愛的,講自由的.
這個是錯誤的一個見解.
所以這裡講到.
今天想講舊約的律法.
當上帝和他的子民立約的時候.
給他們法板,給他們石板的時候.
其實那個律法是有哪幾方面需要留意的?.
第一方面.
舊約裡面的律法.
是以恩典來建立的.
是由恩典來開始的.
我們去土耳其.
我們很多時候講起土耳其的時候.
想起什麼?.
聖經的故事或者聖經的事蹟.

$^{201}$初期教會.
啟示六七教會.
全部都是新約的東西.
但很多時候我們忘記了舊約的土耳其.
土耳其舊約的時期.
其實是一個帝國在那裡存在的.
那個帝國叫做Hittite.
黑人帝國.
黑人,不是嚇死人,是黑人.
黑人帝國.
那些考古學家在黑人帝國.
即是今天的土耳其其中一個地方叫做Hattusa.
那裡找到當時黑人帝國的石板.
他們留下的一些文件.
是以前政府的文件.
黑人帝國大概是什麼時間呢?.
就是出埃及的時候.
大概主前1500至1200年的時候.
即是距離現在大概三千多年的時候.
當時黑人帝國Hittite Empire.
它是一個超大的superpower.
它和埃及或者其他小國有很多條約.
有很多外交的文件來往.
其中一種文件叫做Suzerain Treaty.
Suzerain Treaty.
什麼叫做Suzerain Treaty呢?.
我們中文大概也是一種不平等的條約.
即是說一個大的power和一個小的power.
大家之間有一個條約.
一個大國和一個小國之間有一個條約.
我們叫它宗主條約.
找到這些石板的時候.
翻譯出來的時候就發現.
這些石板很多時候.
一些外交文件都有一個format.
那個format就是條約開頭的時候.
那個大的王就會自稱或者介紹自己.
我就是什麼什麼Great King.
我就是Great King of Hittite.
Super Great King of Hittite.

$^{241}$I am Great.
我很厲害的.
很厲害的.
講了幾十句.
就是講他自己多厲害.
這個就是自我介紹.
序言preamble的地方.
然後那個格式落到下面.
那個大國大王就會和那個小王去介紹.
之前我恩代過你什麼.
你記不記得很多年前.
你被人差不多攻陷的時候.
我派兵去救你等等.
我曾經幫過你的.
我是一個大國.
我曾經恩代過你的小國.
現在我就告訴你.
What I have already done for you.
以前我曾經為你做過的好事.
所以到第三部分.
現在你應該怎樣和我建立關係.
你是一個小國.
你應該怎樣回應我對你的恩典.
然後就有很多條款列出來了.
然後到最後第四部分.
就是一些祝福和奏座的語句.
就是說如果你是faithful的話.
我就會繼續祝福你.
如果你不faithful的話.
你反背我的話.
那你就麻煩了.
我就會來打你.
就是祝福和奏座.
大概這樣的條約的形式.
我們找回的時候.
我們就發現.
《翻開聖經》就是舊約律法的格式.
雖然舊約五經裡面記載十誡.
或者不同的條約的時候.
用一個故事式來記載.

$^{281}$但是裡面有很多的元素.
其實是和《舊約律法》很有關係的.
並且它的格式是跟著這樣寫的.
你記得十誡第一誡是什麼.
「除了我以外,不可有別的神」.
我們小學開始.
如果是接觸教會的.
我們就會背這些東西.
第一誡.
「There is none other than God himself」.
「除了我以外,不可有別的神」.
然後第二誡就是.
「不可以雕刻偶像」等等.
我們就說背十誡.
但其實十誡那一段.
《出埃及傑》第二十章.
開始的時候不是.
「除了我以外,不可有別的神」.
你記得是怎樣開始的.
是「我是耶和華你的神」.
「我是耶和華你的神」.
曾將你從埃及地遺老之家領出來.
只是這兩句說話.
其實就是跟著Suzerain Trini的格式.
第一句是說什麼.
「我是」.
「I am the Lord your God」.
「I am Yahweh your Lord your God」.
「我就是耶和華你的神」.
「我就介紹我自己」.
「我是一個怎樣的神」.
當然這個就不像那些客人的帝王.
在自己煲自己.
或者讚自己.
「我是多麼偉大」.
「我是多麼厲害」.
好像今天那些國家領導人.
讚到自己「笨」.
「落花流水」.
其實他不是的.

$^{321}$他虛報自己.
是 propaganda 而已.
但聖經裡面就很簡單地說.
「我是耶和華你的神」.
然後第二句.
其實是按照Suzerain Trini的第二部分.
就是「我曾經為你做過什麼」.
「我曾經將你從埃及地領出來」.
大家記得那個故事.
到創世記最後的時候.
他們下了埃及.
是不是因為他們很好.
是不是那些希伯來人.
他們很乖,很聰明.
他們是世上最好的人.
最仁慈,最良善的人.
所以上帝救他們.
不是的.
上帝是用恩典去救他們.
不是因為你值得救你.
而是因為「我是耶和華」.
「我用恩手將你從埃及領出來」.
不是跟你講條件.
是因為我愛你.
所以我拯救你出來.
將你領出來.
這兩個說話就是跟這個格式的.
就是說我已經將你拯救了.
現在你已經有一個新的身份.
被我恩典拯救出來的時候.
你應該如何生活呢.
所以拯救恩典是在前的.
然後所謂的律法.
其實是指引.
是從恩典流出來告訴他們.
你們現在應該要如何生活.
所以我們不要再說.
現在不在律法的時代.
那是舊約的事.
現在新約我們不講律法.

$^{361}$我們只講恩典.
這是錯誤的觀念.
因為舊約裡面十誡開始的時候.
已經是從恩典開始.
上帝拯救你.
我們今天新約的信徒.
我們知道在基督裡面拯救了你.
你應該有一個新的身份.
所以你應該如何去生活.
這些就是律法的條款.
所以不是說我做好一點.
守好律法.
我就能夠成為神的子民.
應該相反.
我現在已經因著上帝的恩典.
我成為上帝的子民.
有一個新的身份.
所以我現在應該要如何生活.
等於我的兒子.
我兒子二十多歲了.
他百分之九十九點九都很乖.
很乖.
因為他乖所以成為我的兒子.
但因為他明白他是我的兒子.
所以他應該生活如此.
有我們的家規.
或者我們家的文化.
所以今天譬如我們飛去土耳其.
我們有幾個電影節目.
我們要去機場.
他就自己幫忙.
他說我載你.
嘩 多好.
多好.
當然如果他開車不好.
我就會覺得害怕.
但我知道他開車OK的.
他開車很好.
所以說真的好.
他免費載我.

$^{401}$不用搭Uber.
非常好.
因為他有這個身份.
所以他如此行.
不是說我載你.
如果我載你.
載得我乖的話.
我就做你的兒子.
如果不乖.
我就做你的五翼仔.
你趕我走.
不.
上帝的子民和上帝的關係都是一樣.
用恩典來開始.
第二方面.
這些律法.
雖然好像和我們文化沒什麼關係.
講那些東西大部分都是和古代的社會有關.
不過.
其實律法的精義是什麼呢.
是顯出上帝的心.
和上帝的心腸.
上帝的智慧.
上帝對他的子民的領導.
應該給他們一個怎樣的生活模式.
以至他們的生活模式在萬國的中間的時候.
萬國見到他們這樣生活.
見到他們這樣去處理他們家裡面的事.
生活裡面的事.
他們就會說.
他們是不同的.
我給兩個例子給大家.
第一個例子就是.
講到房屋的.
為什麼講房屋.
這個就是我們看舊的律法的時候.
我們覺得不太合理的.
為什麼講這些建築法呢.
這些建築法又不適用我們.
我們看看是不是適用我們.

$^{441}$起碼我們知道一個.
他說你建造新的房屋的時候.
他就吩咐那些以色人去到迦南地的時候.
他們建造房屋.
他說你屋頂上一定要有安欄杆.
免得有人從房上掉下來流血的罪.
就歸在你的家了.
ok.
為什麼.
你屋頂有沒有欄杆.
為什麼.
因為芝加哥通常是斜的.
是不是.
還有通常芝加哥屋頂是用來做什麼的.
很少用的我們.
防水當然是.
我們很少上去BBQ.
跳舞或者做運動.
是不是.
通常不會的.
因為我們的生活方式不同.
以前以色列人.
他們的生活方式是平頂的房屋.
一定是平頂.
因為他們要儲雨水.
然後將雨水儲在地下水池等等.
然後房屋的屋頂.
很多時候用來曬禾稈草.
或者曬稻米等等.
禾稈放在上面等等.
去處理一些事情.
所以他們的屋頂.
其實是工作間.
他們在上面不是曬太陽.
是去到農業的生產.
一個很重要的空間.
所以因為房頂經常要用.
他們很多時候在旁邊有樓梯上去.
爬兩層樓梯就上到上面的房頂.
你就可以估計到他們的生活會是怎樣.

$^{481}$有時候小朋友就會爬上去.
然後他們就會在上面.
不知道發生什麼事就會玩.
然後就會掉下來.
對不對.
所以為什麼以色列人去到他們的家暗地的時候.
他們要在房屋的頂上面建圍欄.
其實就是保護生命.
顯出了上帝給神的指紋要怎樣.
告訴他們生命是神聖的.
所以那個真正的律法的意義.
不是我們照辦主管今天下來這裡.
21世紀我們這樣用.
而是我們要知道背後的意義.
那個意義就是上帝對生命的看法.
上帝很珍惜每一個生命.
因為每一個生命都是祂所愛的.
所以你們建房屋的時候都有建建屋.
這些外邦的律法裡面.
不是時常有的.
但是以色列人的律法卻是有的.
還有什麼樣本可以給大家.
這一句應該大家很熟悉了.
以前可能曾經說過.
在利未記第十九章.
他說你收割莊稼的時候.
為什麼收割莊稼.
因為是農田的生活.
農業的社會.
你收割莊稼的時候.
不可以割盡田割.
也不可以收拾所遺落的.
要留給什麼人.
窮人和寄居的外人.
什麼意思.
就是你去收割田和的時候.
你可以什麼都收了.
是不是.
就是說你賺錢可以賺到盡.
所有割了頭的東西.

$^{521}$你全部都拿了.
所有大麥小麥的東西.
全部都要了.
你可以這樣的.
但是以色列人你不同.
上帝的子民你不能夠這樣.
因為你不可割盡田割.
那些剩下來的東西.
你不要走去再收拾.
不要扎乾扎盡它.
原因不是主要環保.
原因是留給一些窮人.
當那些比較有錢的人.
他們去到收割完之後.
那些窮人很多時候就會在傍晚的時候.
他們就會來到.
然後他們就可以收拾這些剩下來的東西.
這些窮人寄居者.
他們未必是有合法的居留權的那幫人.
你要留下給他們.
因為他們都是上帝所愛的.
這裡說到律法是什麼.
上帝想熱神的子民.
是有一個公義慷慨的生活.
我們怎樣看我們的錢財.
看我們的資源.
不要割盡它.
意思是不要算死草.
那些婆娘數不要講那麼多.
我們很喜歡講婆娘數的.
我們懷別人的格是這樣的.
小哥崩都不給你.
為什麼?因為公平.
Yes I know.
是應該這樣的.
這裡不是說你買菜的時候不要跟別人算死草.
你應該算到盡.
不是說要被人騙.
而是你要慷慨的生活.
因為上帝是一個慷慨的上帝.

$^{561}$祝福你的時候.
你就應該去慷慨的祝福其他人.
甚至一些你覺得不應該被祝福的人.
我用不著理他.
他又不是公民.
他又不是我們自己人.
外人來講清楚.
外人寄居者.
和我有什麼益處.
他不會回饋我的.
但是你應該有個慷慨的心態.
應該這樣去生活.
這個不是說要照辦煮碗.
你都不會收割莊稼.
你不會趕著走去收割莊稼.
大家在講什麼.
我們怎樣處理我們的資源.
上帝祝福我們的東西.
我們應該慷慨的去奉獻給別人.
祝福其他人.
所以立法會一直都是這樣講.
說顯出了上帝自己的心意.
和上帝對上帝的子民.
有什麼智慧.
如果有人這樣生活的時候.
其他人才會覺得好奇.
咦,我們每個人都是割盡田國的.
為什麼你不割盡田國.
What is the purpose?.
以色列人就會告訴其他人.
因為上帝是這樣.
上帝有沒有和你割盡田國.
我告訴你.
上帝沒有欠過你什麼.
上帝沒有欠過你一些東西.
但是上帝在基督裡面.
祂滿滿祝福給你.
你有沒有想過.
你不要常常覺得上帝欠你這個欠你那個.
不給這個給你.

$^{601}$不給那個給你.
上帝一點東西都沒有欠過你.
由你出生那天開始.
到你死那天.
上帝沒有欠過你一個斗靈.
但是上帝是一個Generous的教.
Generosity.
甚至祂的獨生子耶穌基督.
在十字架上為你死.
這個就是上帝.
如果你說你敬拜這個上帝.
如果你信這個上帝的話.
Please.
在你生活裡面.
給多一點人看一下.
看到你信這個上帝.
作為上帝子民的身份.
這個就是舊約的律法的本意.
第三方面就是講到Mission.
我們說這個上帝是一個God of Mission.
上帝永遠都是想拯救罪人的.
由始至終都是.
由創世記到啟示錄.
上帝是希望去藉著祂的教會.
藉著祂的子民.
去顯出祂自己的大名.
去告訴人聽.
上帝是好的.
去吸引人來到上帝的面前.
在《新冥紀》第四章.
摩西去吩咐那些以色列人.
他們還沒有進入迦南地之前.
就用一大堆東西來吩咐他們.
講了很多篇幅.
告訴他們你進入迦南地的時候.
應該怎樣生活.
剛才我們說莊惠蘭.
不要割盡田國.
這些都是剛才那個package來的.
但《新冥紀》第四章這裡說了一句話.

$^{641}$說什麼呢?.
他說你進入去生活的時候.
在上帝的同在下去生活的時候.
成為上帝的子民.
你們要謹守進行上帝的律法.
你們就在萬國的面前.
其他你的neighbor面前.
在你身邊的人面前.
去顯出智慧和聰明.
他們 即是說其他的人.
非以色列人.
聽見這一切的律例.
他們必定會說.
這大國就是以色列國.
這大國的人.
真是有智慧有聰明.
哪有一個大國.
有神明與他們相近.
好像耶和華我們的上帝.
在我們求告他的時候.
與我們相近呢.
他們會讚嘆.
他們會問.
這麼奇怪.
為什麼他們會這樣.
為什麼我們這麼好奇.
他們的生活的方式.
或者生活的心態是與我們不同的.
因為他們的神明是不同的.
是不同的.
我們的生活其實都是一樣.
《十誡》是上帝親手寫的法版來的.
在不同的經文都有提過我們.
說摩西拿著兩塊的石板.
兩塊記得.
因為當時立約的條款就是這樣的.
你一個copy我一個copy.
你就將你的copy放在你最重要的地方.
就是你最大的廟那裡.
你那個拿了一個exactly的copy.

$^{681}$你就放在你那個廟裡面.
然後每年就拿出來種稻等等等等.
記住這些律例這些約.
所以兩塊的石板是.
two tablets of the covenant.
two tablets.
exactly一樣.
所以以色列國.
他們有那兩塊石板就不同了.
很有趣他們兩塊石板都是放在藥櫃裡面.
在他們的致勝所裡面.
因為上帝耶和華大王.
他最大的最重要的地方在哪裡.
就是以色列中間那個temple.
回望.
而以色列人作為一個小國.
他們最重要的地方在哪裡.
就是敬拜耶和華的地方.
回望.
和將來在耶路撒冷的聖殿.
所以在聖殿裡面的致勝所.
有一個藥櫃有一個盒.
那個盒裡面有幾樣不同的東西.
很重要的東西.
其中一個最重要的就是石板.
所以兩塊都在裡面.
一塊代表耶和華.
大的王.
the great king.
一塊是代表以色列國.
上帝的子民.
the small kingdom.
大家建立了關係.
就是大家放在這裡.
所以說到這個石板的時候.
上帝用指頭親自寫上去.
親自去雕上去.
你一份我一份.
所以那份東西不只是備忘錄這麼簡單.
是一個永久的盟約.

$^{721}$所以為什麼作為上帝的子民.
他們有一個新的身份之後.
他們就有一堆新的.
去啟蒙他們.
怎樣做一個新的子民.
怎樣生活的一種方式.
這些就是每一個律法的條款.
律法是什麼呢?再說一次.
是先有恩典.
然後我們如此行.
不是因為我們做好一點.
我們乖一點.
所以我們能夠成為神的子民.
反過來.
是既然是上帝的子民.
我們應當如此行.
既然我們已經是被耶穌基督救贖的一班人.
我們就應當如此行.
上帝既然是一個慈愛的上帝.
我們應當在我們生活裡面流露慈愛.
說話的時候流露慈愛.
待人的時候流露慈愛.
既然上帝是一個忍耐我們的上帝.
我們就應當怎樣?.
恆久忍耐.
對不對?.
既然上帝是一個赦免人的上帝.
我們處事和其他人生活的時候.
我們就應當學校上帝的現象.
上帝是一個慷慨的上帝.
我們就應當是一個慷慨的人.
這就是我們是誰.
所以最後今天一個問題.
就是到底我的生活方式.
我的生活方式,我的心態.
我們的言語.
我們待人接物等等各方面.
我們有沒有引起別人.
讓別人覺得.
是有點不同.

$^{761}$I want to know.
What is going on?.
到底我們的生活是怎樣?.
講律法的時候不是講一些要去欺騙我們.
要去陰我們的東西.
上帝的律法原來就是一種邀請.
叫我們活得似上帝.
叫我們活得有智慧.
活得像上帝一樣慷慨.
以致我們的生命能夠on a mission.
我們教會在講抗疼.
如果我們的心態不擴張的話.
擴多大也沒有用.
或者就算我們不擴.
我們舒服服也沒有用.
今晚照一照鏡.
我自己的生活有沒有被上帝的恩典啟動過.
我經歷過上帝的恩典,上帝的愛.
我的生活有沒有像上帝多些.
有沒有渴求我更加像上帝.
我講一句說話的時候.
以前是這樣講.
現在我有沒有再想一想.
基督會怎樣講.
我以前這樣處理我的資源,我的時間.
現在我有沒有有些改變.
是因為我身份上已經有改變.
以致我的生活,我的lifestyle都會有改變.
讓這個問題.
這個星期你帶走.
我的生活方式有沒有讓其他人感到好奇.
好奇我這位這麼慷慨,這麼愛我的神.
是一個怎樣的神.
各位.
我心底裡要祈禱.
親愛上帝多謝你,因為你先拯救我們.
然後你設下你美好的律法.
主要是你自己的律法.
有如金紙一樣.
有如金銀一樣.

$^{801}$不是去威脅我們,不是去限制我們.
而是一份禮物.
叫我們這樣進行你的生活的時候.
使人能夠渴慕你.
使人能夠認識你.
主要是讓我們認識你更多.
讓我們的言談舉止.
讓我們作的決定.
讓我們處理我們生活裡面有的東西.
讓我們做一個好的管家.
讓我們的生命流露出.
上帝你自己更多美好的本性.
讓人能夠因此而被教會.
被你自己神的指紋來吸引.
去認識你.
祈禱奉著名.
\newpage



\section{約書亞記 17:14}
\label{sec:032153}
\textbf{聖經人聖經事:不可僅有一鬮之地}~~~~黃克勤牧師博士
\newline
\newline
連結: \href{http://www.church.com.hk/acms/content.asp?site=cdc\&op=show\&type=product\&code=032153\&layout=sermon}{\texttt{032153}} ~~~~ 語音日期: 2026-05-21
\newline
\newline
\hyperref[sec:032152]{\small{< < < PREV SERMON < < <}}
~
\hyperref[sec:index_chronic]{\small{[返順時目]}}
~
\hyperref[sec:032154]{\small{> > > NEXT SERMON > > >}}
\newline
\newline
約書亞記 17:14
\newline
\begin{longtable}{cl}
\hline
\hline
章節 & 經文 (和合本修訂版)\\
\hline
17:14 & \begin{tabularx}{0.7\textwidth}{X} 約瑟的子孫對約書亞說：「耶和華到如今這樣賜福給我，我百姓眾多，你為甚麼只給我抽一籤，分一份的土地為業呢？」 \end{tabularx} \\ \\
[1ex]
\hline
\hline
\end{longtable}
$^{1}$眨眼有一日,仍然是講以色列人.
不過今次是講弱勢的後裔.
你知道,以色列十二支派.
約書亞帶他們入了迦南地.
摩西帶他們離開埃及,一路行完抗日入迦南.
入迦南約書亞接手帶他們入去.
一路打低了當地人.
南部五個王,北部數不到那麼多個王.
都被他們平定了,吊在樹上.
將那些人趕走或殺死.
開始分地了,不過還有很多未得之地.
如果你早幾日看到地圖.
有些是粉紅色未得之地,黃色得了的地.
後來再講,他們都開始得了.
未得那怎麼辦?未得已經分了.
分了又未得怎麼辦?你自己去打仗獲得.
你自己去拿,昨日讀就是.
你迦南給你一塊地,你自己去打,自己去獲得.
你不要等,等你給我,天跌下來.
我們無論自己屬靈的經歷.
無論世上有什麼都是神給我們.
叫我們朝著目標去獲得.
你獲得給你,是應許的,你拿.
我不拿,不拿那個應許的禮物永遠給不到你.
起碼你要自己伸手去拿.
要做些事的,都要戰鬥的,都要獲得的.
好了,現在分地了.
這句是「不可僅有一溝之地」.
即是什麼呢?即是說.
「尖溝」的意思是抽籤「十二之派,長老來哂」.
這塊地我已經畫好地圖了.
其實沒有畫地圖.
約書就說這塊地在我們這裡.
大概這些山名,城市名都在這裡.
大家看看這個地圖,我們圍著看.
現在我們有十二之派.
利美之派,照神的說話是沒有地的.
不過他們住在我們其他十一之派的人當中.
帶領我們敬拜神.
為我們做祭司,為我們做這些敬拜靈修的事.

$^{41}$獻祭的事.
而十一個之派裡面.
有一個之派是會要多些地的.
為什麼?人多.
那怎麼分呢?.
我們大概是抽籤「尖溝」.
抽籤極有可能是將瓦片寫了名.
十一塊,十二塊寫了名.
然後十一塊寫了名.
放在袋子裡.
這裡先說了怎麼分.
這個之派大概這裡到這裡.
這個之派這裡到這裡.
沒有很準確的地界.
沒有拉繩子,沒有測量.
大概這樣.
不過兄弟去到哪裡不要爭.
行了,好,明白.
抽了.
抽到,現在我們抽這塊地.
抽這塊瓦片.
又是大之派.
這裡,這樣就尖溝了.
每一個之派都有一塊瓦片拿出來的.
一塊瓦片就是一塊地.
不可以只有一塊地.
一塊尖溝就代表一塊地.
這個我們不能只有一塊,我要兩塊.
什麼意思?我們看看.
約瑟是記十七章.
約瑟的兒子來到跟約瑟說.
這個可能其他人都聽著.
不過沒有人反對.
「也或說到如今既然賜福於我們」.
約瑟的子孫,約瑟的子孫是誰呢?.
以法連嗎哪西,就是約瑟的子孫.
約瑟的子孫有兩個.
你知道十二之派裡面少了利未之派.
十一之派分地.
十二之派,但是他們是優待約瑟的子孫.

$^{81}$所以約瑟是沒有地的.
他們兩個兒子,以法連嗎哪西.
就是由約瑟出來的兩個.
於是他們有個以法連之派的地.
有個嗎哪西之派的地.
約瑟的子孫對約瑟說.
「也或說到如今賜福於我們」.
約瑟的子孫.
「我們族大人多」.
其他人都覺得「是呀,是呀」.
約瑟也認為「是呀,你真是族大人多」.
「你為什麼緊張一溝一段之地,分級我們為業呢?」.
譬如大家都大概差不多一大.
大家有些山,有些平地.
差不多一大.
但是我們族大,族大人多.
那樣很擠迫,不夠住.
有什麼不夠?你看香港這麼小的豆腐.
可以容納八百萬人.
加上流動人口可能有千萬.
八,九百萬都不夠.
建高一些樓,一百層.
當時沒有的.
而且他們習慣住整個曠野.
很闊的.
又要有農地,又要圈住地.
又要有牲口,很闊很大.
我們族大人多.
一段一溝的地方不夠.
分給我們為業,不夠呀.
約瑟先生.
是一個很真誠,很真誠,很真確的需要.
約瑟先生說「是喲,你們如果族大人多.
嫌以法連山地,即給了你那塊地宅小.
不如你去比利洗人.
利佛欽人之地.
那裡是樹林中坎佛樹木.
那裡有些山地,你去斬樹.
住在那裡,你不夠.
你就去拿比利洗人,利佛欽人的地.

$^{121}$在那裡如果是樹林,你斬給他.
大家在這裡不要用我們今天的想法來想.
斬樹這麼不環保.
當時沒有這些想法.
如果要斬樹,初期的社會都是這樣想.
甚至現在還是這樣做,不讓做就偷偷做.
利佛欽人.
利佛欽人其實和阿蠟族有關係.
他們很大隻,很高大的利佛欽.
很高大.
如果不說那些典故創世記.
那些人很兇猛,很大隻.
你們人多勁,你拿他們的地,趕他們走.
殺了他們.
而且不夠住斬樹,你還要拿森林.
約瑟的子孫說.
那山地都容不下我們.
並且住在平原的迦南人.
住在白線和屬白線的鎮市.
並住耶斯利平的人都有鐵車.
那山地你叫我們上山斬樹住在那些地方.
我們都不夠.
那邊有平原.
不過那些平原的名字叫白線.
是一個城市,一個鎮市.
白線中間有一點,就是兩個字.
白線.
Beth.
Sheran.
Sheran就是善.
Beth就是白.
其實Beth就是地.
Sheran.
唯獨Sheran都可以.
就是靜,安靜之地.
或者住在安靜之地的城市.
和附近耶斯利平的人.
那邊我們都可以去拿.
不過那邊的人是很強壯.
而且有鐵車.

$^{161}$即是說他們很好打.
軍備很強.
我們用木頭車,我們用步兵,他們用鐵的戰車.
很厲害.
那我們怎樣去打.
那你拿去吧.
好了.
再看一些經文.
留意一下.
約瑟的子孫說.
我們族大人多.
只抽一個籤是不夠的.
所以我們要大些地方.
約瑟就說.
是喔,你族大人多.
如果只是比二百年的山地不夠.
那不如你們拿給你洗人.
你的陰人和森林的地.
如果都約瑟就說.
其實都不夠我們用.
那就是我們想要的.
就是可以拿到白善.
和附近的城鎮.
和耶斯利平原的地方.
那就好很多了.
那都有鐵車的.
我們有排打的.
可不可以給我們.
好啊,給你吧,你拿去吧.
好了,今天在這裡.
和大家將剛才那段說話.
看看有三點要和大家學習.
第一,約瑟的子孫.
記不記得.
創世記約瑟.
約瑟被賣到埃及.
然後被賣做奴隸.
一直做了宰相.
七隻肥牛,七隻瘦牛.
七顆肥的麥子.

$^{201}$七顆瘦的麥子等等.
然後他做了宰相.
大家記得了.
然後埃及法老王找了一個祭司.
將女兒嫁給約瑟.
約瑟就生了兩個兒子.
兩個兒子,一個大兒子叫以法連.
小兒子叫嗎哪西.
當時已經沒有說.
不過在這裡說一下.
以法連,嗎哪西.
剛才以法連做代表.
做約瑟的子孫.
來求約瑟傾.
我們族大人多不夠.
真的族大人多嗎?.
當然了.
我有以法連.
還有我弟弟嗎哪西.
我們的人多得不得了.
大家照著一溝一段地.
我們不夠住.
約瑟說.
你拿給你洗陰人的地和山地,斬樹.
可否要白善,安靜之地.
和附近耶穌列平原.
不過他們有鐵車.
我們都很頭痛.
不要緊,我們去打.
給了他.
我想提一提.
約瑟的兩個子孫厲害嗎?.
厲害.
昌不昌盛?.
昌盛.
族大人多.
以法連.
Efraim.
以法連很有趣.
Efraim.

$^{241}$意思是忘記.
你怎會改兒子名叫忘記.
忘記做了功課嗎?.
忘記昨晚睡得好嗎?.
忘記吃飯了嗎?.
很有趣叫兒子名叫忘記.
意思是神使我忘記我父家和環境給我的痛苦.
大家要想族大人多.
興旺你必須要先忘記.
不要想著回望手扶著泥向後看.
你做不到事.
有人在基督律就是新族有人舊事已過都變成新的.
你要向著前面看.
將事情做好.
努力忘記後面的仇虎.
先忘記後昌盛.
嗎哪西他改了第二個兒子名叫Manachet.
Manachet.
應該有個N字跟這裡都可以.
Manachet.
Manachet就是昌盛.
神使我在借呂昌盛.
約西用兩個兒子的名記起上帝的作為.
第一個神使我忘記以往的痛苦.
第二個神使我在借呂昌盛.
很有意思.
你要先忘記後來才會昌盛.
這兩點跟大家說.
確實是約西的後裔變成族大人多.
非常感恩.
連平原耶斯利平原有鐵車.
我們要敢去打他.
獲得.
第三點跟大家看公平.
咦.
大家都明白什麼是公平.
我想在這裡煽釋一下我對公平公正的想法.
公平等不等於公正呢.
是不等於的.
什麼叫公平呢.

$^{281}$現在我們有十個人.
我們就十個人每人有一碗飯.
公平了.
每個人都有一碗.
一個人頭一碗飯.
人頭飯.
每個人一碗.
公平了.
但我覺得這不是公平.
一人一碗飯.
這叫公正.
按照規矩去做一人一碗飯.
公正.
你多給他一碗就不公正了.
不公平他有兩碗.
不是.
公正.
公平是什麼意思呢.
公平和公正是不同的.
不過我們混亂了.
我覺得.
我對公平的看法是什麼呢.
你是族大人多.
並且強盛.
什麼叫強盛.
多東西.
很多.
人又多.
吃的東西又多.
不可僅有一溝之地.
如果你有以法連有.
馬來西族大人多.
弱勢的後代.
又強盛.
很多東西.
而且馬來西是很多生畜.
不可以僅有一溝之地.
即是說.
大家抽出來的地.
大家都這麼多.

$^{321}$地這麼多.
不過放在不同的位置.
人多不行.
一溝一地你不夠的.
弱勢是這麼說的.
你不公平.
不是.
是公正.
不過不是公平.
什麼叫公平呢.
十個人.
有一個像我這樣.
嘩這麼高的.
要吃兩碗飯才飽的.
一個人.
這麼矮.
又不喜歡吃飯的.
所以半碗飯他就夠了.
如果人頭一碗飯.
我就不夠飽.
他就太多.
人不理.
人人一碗飯.
是公正.
但並不公平.
因為我的需要是大一點.
他的需要是小一點.
所以我應該多一點飯.
他應該少一點飯.
不要一人一碗.
他半碗夠了.
我半碗夠了.
那剩下的半碗怎樣.
給我吧.
我需要多一點.
如果來一個加特人.
哥利亞.
他就會說我九尺.
沒有九碗飯不夠.
那怎樣.

$^{361}$所以其實公平的意思就是.
按著需要而分配.
那叫公平.
公正就是按著法理而公正對待.
沒有偏頗.
那叫公正.
所以這個辦法你族大人多.
又強盛.
所以你不要因為一溝一地.
你是不夠的.
那不如你拿多一點地.
那些人的地你拿去.
山上你拿去斬樹.
平原和安靜之地可以嗎.
可以.
給你.
那怎樣.
族大人多.
我覺得這個應用了公平的原則.
而不是公正的原則.
我想和大家分享我這個看法.
我們到這裡一起禱告.
天會上帝我們感恩多謝你.
因為你教導了我們.
去更多學習公平公正.
又見到主令到若失經歷了很多艱難.
祂將祂在神面前的經歷.
放在兩個兒子的名字上面.
使我忘記以前的痛苦.
然後我是現在這裡昌盛起來.
如果還是拿著以往的痛苦不肯放手.
相信今天很難昌盛.
所以就幫助我們學習若失.
又用祂兩個兒子的名字.
來警惕我們自己.
提醒我們自己.
放下以往的痛苦.
而去獲得今天的昌盛.
幫助我們因著不同的需要.
我們去彈性處理.

$^{401}$而可以獲取更多公平的安排.
雖然很難做的.
不過求主幫助我們做到最好.
讓我們又公平又公正又公開.
都是需要的.
求主自己祝福.
誰聽祈禱奉耶穌名求.
阿們.
多謝大家 明天再見.
\newpage



\section{但以理書 6:10}
\label{sec:032154}
\textbf{與素常一樣}~~~~容永祺先生
\newline
\newline
連結: \href{http://www.church.com.hk/acms/content.asp?site=cdc\&op=show\&type=product\&code=032154\&layout=sermon}{\texttt{032154}} ~~~~ 語音日期: 2026-05-17
\newline
\newline
\hyperref[sec:032153]{\small{< < < PREV SERMON < < <}}
~
\hyperref[sec:index_chronic]{\small{[返順時目]}}
~
\hyperref[sec:032155]{\small{> > > NEXT SERMON > > >}}
\newline
\newline
但以理書 6:10
\newline
\begin{longtable}{cl}
\hline
\hline
章節 & 經文 (和合本修訂版)\\
\hline
6:10 & \begin{tabularx}{0.7\textwidth}{X} 但以理知道這文件已經簽署，就進自己的家，他家樓上的窗戶開向耶路撒冷。他一天三次，雙膝跪著，在他的神面前禱告感謝，像平常一樣。 \end{tabularx} \\ \\
[1ex]
\hline
\hline
\end{longtable}
$^{1}$其實今天我想和弟兄姊妹分享的題目就是「與素常一樣」.
今天我講完之後.
弟兄姊妹什麼都忘記了也不要緊.
記住「與素常一樣」就行了.
很容易的「與素常一樣」.
是不是?.
我們親近耶穌是不是應該「與素常一樣」?.
我們祈禱是不是應該「與素常一樣」?.
我們不是應該臨急抱耶穌的腳嗎?.
應該有什麼事都是經常與神親近祈禱的,對不對?.
那我請弟兄姊妹和你旁邊的人說一聲.
我們親近神「與素常一樣」.
如果你沒有親近神的話.
就不要說「與素常一樣」了.
你就要說.
我們親近神要更多親近神「與素常一樣」.
好不好?.
學習「與素常一樣」.
1,2,3.
和你旁邊的人說一聲.
好了.
我們看一看今天我們要讀的經文.
其實是來自《大氣曆書》第六章第十節.
不如我們一起讀好不好呢?.
我們找到的弟兄姊妹一起讀好不好?.
1,2,3.
在這段聖經裡面我們看到.
丹爾尼知道王的禁令蓋了玉璽回到自己的家.
在樓上窗戶開向耶路撒冷.
然後一天三次雙膝跪在他神的面前.
這本聖經之前是在說什麼呢?.
其實是在說大利烏王的時間.
其實他的意思是立了120個總督和3個總長.
而丹爾尼是其中一個總長.
在120個總督之上.
因為丹爾尼聖經仍然有美好的靈性.
所以他很顯然.
他的智慧和能力在各方面.
是超過了這些總長和總督.
而這個王甚至很想納他去管治全國.

$^{41}$我們發覺丹爾尼其實在被擄之地.
大家知道他是被擄到巴比利的少年人.
他和他三個朋友其實很蒙恩.
因為他們很特別.
以求神和親近神.
丹爾尼是共歷四朝.
由巴比倫王尼布甲尼撒.
接著是白沙薩王.
接著就是馬代人大利烏王.
最後就是塞魯士.
我們經常讀舊的聖經.
就是古列王的年代.
歷時是六十載.
但是在六十載裡面.
仍然能夠在那些環境裡面.
好好地被神使用.
服侍神.
也被興起.
在被擄之地能夠大大的被神使用.
幾朝的王都很需要他.
因為他很會解夢.
因為他裡面有美好的靈性.
聖經在丹爾尼書中說過.
他有兩次美好的靈性.
一次是我們剛才所讀所看到的.
就是在第六章第三次.
因為這丹爾尼有美好的靈性.
所以他能夠顯然超乎其餘的總長和總督.
往日立區治理通過.
另一個就是白沙薩王的母親.
太后也承認他有美好的靈性.
所以一個與神親近的人.
其實是有美好的靈性.
我們回顧丹爾尼.
他為什麼能夠有美好的靈性.
他為什麼能夠這麼被神使用.
而且在危急的時間都能夠.
神在他身上打一條出路.
其實我們留意這段聖經.
這段聖經其實就是剛才所讀到的.

$^{81}$與素常一樣.
他親近神是與素常一樣.
他和大家的關係是與素常一樣.
就算他知道生命危急的時間.
是與素常一樣.
因為他身邊的人都妒忌他.
總長和總督都妒忌他.
所以就一定是想設計陷害他.
然後就獻了一條計給大尼烏王.
就是說總而言之.
在三十天之內.
不論什麼人.
他只可以向他的王求.
如果不是的話.
他向他的神或其他人求.
什麼都好.
就要扔他進獅子坑.
這個故事我相信大家都很清晰.
都讀過,聽過很清楚.
但是在這個故事的背景之後.
就是形容大義利.
就是什麼呢?.
就是說他是怎樣呢?.
他就是知道禁令玉璽.
蓋了玉璽之後.
就在自己的家裡.
在樓上的窗戶開向耶路撒冷.
一日三次.
雙膝跪在他神的面前.
與素常一樣.
就是說其實無論有禁令沒禁令.
有這些危急.
他都是與素常一樣.
他親近神就是與素常一樣.
他平時就是這樣親近神.
他生活就是親近神.
他生活就是與神一起.
與神同行.
學習與神一起.
每一天他的生命就是這樣的生命.

$^{121}$所以無論什麼環境.
在他身上你看到的就是.
神與他同在的生命.
所以他能夠面對這麼大的挑戰環境.
都能夠面對這個挑戰.
而能夠我們叫做逢凶化吉.
甚至能夠彰顯神的榮耀.
彰顯神的名.
那麼聖經後來就形容了.
大義利後來當然沒事了.
連家獅子坑.
神保佑了他.
他雖然不聽話.
不聽王的話.
然後人們誣告他.
他真的沒辦法了.
教改了玉璽.
要扔進獅子坑.
但是扔進去之後.
第二天他平安無事.
出來之後.
王都很開心.
因為王其實很不想放他進去.
但是禁了玉璽.
沒辦法了.
但是問題就是.
那些陷害他的人.
連家大小都扔進獅子坑.
一扔進去.
聖經又怎麼形容呢.
永續獅子坑.
他們每有一個坑.
坑底獅子就抓住他們.
神的保守就是大義利.
但是那些陷害他的人.
都沒有一個坑底獅子抓住他們.
大義利為什麼能夠這樣呢.
我覺得有幾點.
是我的領受.
第一他經常心是親近神.

$^{161}$心是於神.
我想大家都有手機.
我們有手機.
我們去外國旅行.
很多時候我們會帶什麼.
除了帶手機還要帶什麼.
充電寶.
為什麼要帶充電寶.
因為你去拍那些照片.
拍那些影片.
去旅行拍很多影片.
很快就沒有了.
電都沒有了.
你會帶什麼呢.
有充電寶.
我在香港就沒有用充電寶.
我去外國.
我想前天從武漢回來.
都是用充電寶.
沒有充電可能中途就不夠電.
在香港我不用的.
但是是不是等於我不用充電呢.
不是.
因為我在車上可以充電.
在車上可以充電.
同樣一樣.
其實我們親近神.
我們需要的是什麼呢.
需要能夠時刻親近神.
大義利就是一天三次.
親近神.
按照猶太人信仰的傳統.
一天三次.
甚至在危難的時候.
與照常一樣.
一天三次.
親近神.
他的生命就是親近神的生命.
所以他能夠有這樣的勇氣.
有這樣的信心.

$^{201}$能夠面對這麼大的危難之前.
他好像泰山崩於前掩.
不動.
因為他裡面有什麼呢.
有神.
當你時刻親近神的時候.
那種信心.
和神賢你那種信心.
那種能力.
那種平安.
才能夠在你身上表現出來.
所以我們為了.
彈劾你這個經歷.
親近神是值得我們學習的.
還有一件事.
他不會因為環境而改變他的信仰.
至死忠心.
但是你在當時的環境裡面.
其實只要他不繼續.
一天三次祈禱.
或者他關上窗.
滿祈禱.
沒有人知道他在祈禱.
他的命都沒事.
他這裡的選擇是什麼呢.
要麼就依靠神.
要麼就死亡.
其實沒有人會想到.
他可以在獅子坑被拯救出來.
是沒事的.
神與他同在.
平安處給他.
沒有人會想到.
在那時候沒有人會知道.
但是.
他卻選擇了什麼呢.
站在神的華語當中.
站在依靠神.
謹守神的華語當中.
使這個性命於道外.

$^{241}$這個也讓我們能夠.
學習到在我們平時的生活.
我們是不是真的能夠.
站穩在神的華語呢.
面對世界的衝擊.
世界的挑戰.
世界的引誘.
有時候我們會去引誘那邊.
挑戰那邊.
而沒有站在神的華語呢.
當我們經營做生意.
我們有沒有站在聖經的華語.
去做生意呢.
當我們待人接物.
有沒有站在聖經的華語.
去待人接物呢.
還是我們隨波逐流.
在社會上.
隨著社會的價值觀.
去飄來飄去呢.
這個我覺得我們都值得.
好好地學習.
當有重大挑戰的時候.
我們是不是仍然站在.
神的華語當中呢.
這個也是我們很需要.
值得我們去學習的.
第三點.
我覺得在大眼你的身上.
你在這裡是信仰生活化.
日日如是.
他一天三次祈禱.
什麼呢.
就是根本就是.
與素常一樣.
平時都是這樣的.
他的生活就是信仰.
信仰就是生活.
其實我想起我們小時候剛進主.
我想很多弟兄姊妹都可能會經歷過.

$^{281}$當你剛進了耶穌的時候.
好像你吃飯會怎麼樣.
祈禱.
有沒有人在吃飯的時候.
閃閃的這樣祈禱.
有沒有這樣祈禱.
我那時候小時候也有.
不好意思怕別人笑.
會不會有些人說笑.
吃完飯你忘記了一碗飯.
其餘都忘記了一碗飯.
忘記了一隻雞.
都有些說笑這樣說.
但是其實.
你在人面前認住.
神在天目面前.
主耶穌在天目面前也認你.
所以我們很重要.
我們要什麼呢.
我們要認.
我們的主耶穌基督.
我們要有膽量在人面前認.
認神.
認主耶穌基督.
認我們的信仰.
我記得有一次.
也有幾次跟朋友在內地.
交流的時間.
甚至有些官員在那裡.
有個朋友說祈禱.
好啊祈禱.
一樣有謝萬月圖.
別人也很尊重你的祈禱.
所以我覺得我們要能夠.
有很好的見證.
信仰生活化.
生活信仰化.
在我們生命裡面.
流露出基督.
讓基督的信仰.

$^{321}$成為我們的生活的傳播.
讓我們能夠養成阿華羽.
引領帶領我們活像基督.
另一件事就是對神的角度和心智.
心向耶路撒冷.
你會發覺到.
在聖經怎樣說.
就是說.
他樓上的窗戶是開向耶路撒冷.
開向耶路撒冷.
他的心是向耶路撒冷.
其實這裡有好幾個意思.
第一.
他是抓住了神的英雄.
其實待會我們有一段聖經會說到.
是所羅門獻殿的祈禱.
其實他向著耶路撒冷的祈禱.
是很回應到.
所羅門獻殿的祈禱的內容.
神和所羅門所納的藥的內容.
他在找神和以色列人所納的藥.
和所羅門所納的藥.
就算他被擄到外邦.
仍然是向著聖殿.
向著耶路撒冷的祈禱.
他的心是嚮往著聖殿.
同樣他的心是嚮往著.
他的信仰就在耶路撒冷.
在耶路撒冷那裡.
他的根就在耶路撒冷.
他的角度就在神的角度.
在耶路撒冷.
所以你看到.
他不是在巴比倫.
雖然他人在巴比倫.
但他的心是什麼呢?.
耶路撒冷.
所以他照樣是向著什麼呢?.
是耶路撒冷.
而且還有一個意義.

$^{361}$就是他期望什麼呢?.
有復興.
能夠為了耶路撒冷的祈禱.
盼望復興以色列民.
這個禱告是一個很大的信心的禱告.
盼望的禱告.
所以你看到他的心智.
是面向神的角度.
面向耶路撒冷.
所以在這裡.
這個是看到大禹利的特質.
或者我快一點.
先讀了這篇聖經.
如果大家可以的話.
跟我一起讀好不好?.
讀一讀這篇聖經好不好?.
一,二,三.
祂們若在被堵到之地想起罪來.
心意轉移.
懇求你說.
我們有罪了.
我們背逆了.
我們作惡了.
祂們若在被擄之地盡心盡性服你.
又向自己的地.
就是你賜給他們列祖之地.
和你所選擇的聖.
並我為你們所建造的殿禱告.
求你在天上的居所.
聚聽他們的禱告祈求.
為他們申冤.
饒恕得罪你的民.
赦免他們的一切過犯.
使他們在被擄他們的人面前蒙面述.
大家看到整段就是當時獻殿.
在最榮華,最風光的所羅門建了殿之後的一個祈禱.
大家也都.
你看到它都預言到將來.
在這裡也都預言到將來.
會有被擄的機會.

$^{401}$在這裡.
五百年後的但以來.
祂的祈禱是找緊了神的心意.
找緊了那個約.
是一個回應的.
面向耶路撒冷.
在聖經裡你看得很清楚.
求你在天上的居所.
聚聽他們的禱告祈求.
為他們申冤.
饒恕得罪你的民.
赦免他們的一切過犯.
使他們在被擄他們的人面前蒙面述.
看到整個.
但以來向耶路撒冷祈禱.
我相信他們都熟讀摩西的律法.
和當時他們可以學習到的教導.
所以你看到整個禱告.
是面向耶路撒冷.
回應得很收徹.
因為這次聖經說.
你的民若得罪你.
你向他們發怒.
將他們交給仇敵.
擄到仇敵之地.
或遠或近.
他們若在被擄到之地.
想起罪來.
回心轉意.
懇求你說.
我們有罪了.
我們背逆了.
我們中了.
他們若在被擄到之地.
盡心盡性歸服你.
又向自己的地.
就是你賜給他們列祖之地.
和你所選擇的城.
這裡說.
他們列祖之地是亞南美地.

$^{441}$城是耶路撒冷的城.
為你的明所建築的殿.
就是聖殿.
求你在天上的居所.
垂聽他們的禱告祈求.
為他們申冤.
饒恕他們.
得罪你的民.
赦免他們的一切過犯.
使他們在被擄他們的人面前.
蒙面術.
整個蘭爾里的祈禱的行動.
其實就回應了這一段的經文.
你看到他很明白神的心意.
很清晰向耶路撒冷祈禱.
他可以關上門.
但是他沒有關上門.
他可以不向耶路撒冷.
他向耶路撒冷.
他知道他的信心.
他知道他的信仰根基在哪裡.
在神那裡.
在神的國度那裡.
他的心是向著神的國度.
他有一個對神國度的伏木士心智.
所以這是很難得的.
另一件事就是.
我看回前面的一兩段.
第五點就是表明.
神主民的身份不被榮華富貴所腐化.
這個其實很不容易的.
特別是在現今的社會.
有很多不同的吸引力.
這次我聽說是職場的聚會.
所以先聽我們職場的弟兄姊妹來分享.
在職場上有很多的吸引.
很多的商業利益.
很多社會利益.
很多各方面的利益.
我們沒有被這個社會的富貴榮華.

$^{481}$腐化了我們自己.
我們單純對主的信仰.
會不會慢慢地付諸流水呢?.
這裡很值得我們一日三次向神祈禱.
我們去學習.
一日三次的意思就是.
你要全心全意向神祈禱.
你要向神祈禱.
親近神.
時刻親近神.
在這裡是很重要的.
其實主耶穌在世的成長都給了很好的榜樣.
祂雖然是神的兒子.
祂有神性.
但是祂來到世上.
降生為人.
道成肉身.
祂也有人性.
在西馬遼祈禱.
祂也有哀傷.
祂也有痛苦.
求主將你的腹部拿開.
但是求主按祂的意思.
不是按祂的意思.
但是你看到主耶穌在祂的生活裡面是怎樣呢?.
清晨祈禱.
晚上祈禱.
當很多人續擁祂的時間.
祂會祈禱.
當祂行神跡之前.
祂會祈禱.
差不多你看到祂.
還有五千人.
五餅二魚.
借飯祈禱.
四千人吃得飽.
祈禱.
耶穌的生活就是祈禱的生活.
其實祂有神性.
但是祂是謙卑在天父的面前.

$^{521}$在父的面前.
祂給了我們美好榜樣.
祂是時刻祈禱.
聽天父的旨意.
信服天父的旨意.
聖父聖子連在一起.
這是一個美好.
給我們一個很好的學習.
所以主耶穌教導我們祈禱的.
主通文第一句是怎樣說的呢?.
我們在天上的父.
願人都尊.
你的名為聖.
祂將榮耀歸給父神.
祂向父神祈禱.
所以祈禱是我們生活必需的.
我用養料營養.
或者是一種生命的泉源.
祈禱讀的是我們生命的泉源.
令我們更加能夠時刻親近神.
了解神.
認識神.
得到從神而來的諸腿力量.
讓我們能夠時刻提醒.
在這裡我想說一個見證.
就是我記得九九年的時間.
在年初我們雙人團祈有一次的除外會.
我記得是禮拜六.
說起一千年尾.
在土耳其以弗所有一個世界祈禱會.
不如我們也參加吧.
一起去吧.
為著接下來的千年國度來祈禱.
因為大家知道.
二千年就是一個新的一千年的國度.
那時候你有沒有發覺.
在社會上很流行討論的一個叫什麼.
千年叢.
你記得嗎?.
你有記得.

$^{561}$所以很有一個吸引去.
所以我說好啊去啊.
自己也承諾了去.
誰知道到了九月的時間.
其實那一年.
也是我們國家非常重要的一年.
就是我們的立國五十週年.
在九月份的時間.
我忘記了九月初還是九月中.
我收到一個邀請.
請我們北京參加.
參觀閱兵.
不是參加.
是參觀.
自己不能閱的.
自己不是兵.
參觀.
還有國慶的活動.
甚至可以國宴那些活動.
這是很榮華.
很尊貴的身份才能去到的.
也是很難得的機會.
人生有幾何呢?.
不是很多機會.
但是那時候是九月.
但是我三月已經承諾了.
我們要去耶路撒冷祈禱.
這次是主讓我知道.
要先求神的國和神的義.
一切都會交給你.
讓我知道.
我們要盡心盡性盡義.
盡力愛祖我們的神.
所以我選擇了去耶路撒冷祈禱.
去耶路撒冷.
先去土耳其的耶福特祈禱.
然後再去耶路撒冷.
但是那個旅程和我太太在一起.
那個旅程給我一個靈性上的更新.
特別是去到架欖山.

$^{601}$祈禱.
看著整個聖城.
祈禱的時候.
那個感動.
那些眼淚就流出來.
那個屬靈的經歷.
又是另一個新的領受.
所以其實這些是.
沒有其他東西可以換到的.
選擇了去那裡.
我還被人說.
你真是不識抬舉.
有機會讓你也不要去.
我真是有些不識抬舉.
但是我會選擇最好的.
就是什麼?.
我們的主耶穌基督.
我們的神.
結果就是.
去了祈禱.
去了耶路撒冷.
去了巴黎行城.
真的很開心.
但是呢.
接著有五年後.
就是國家的慶典55週年.
那時候我還沒有是政協委員.
但是那時候就有.
我心裡面.
這次有沒有機會呢?.
都還沒有正式的祈禱.
心裡面想一想.
很奇妙.
主又為我預備了一個邀請.
就已經有邀請我了.
從此之後.
每一次大的節慶.
我都有份.
所以其實呢.
我們將榮耀歸給主.

$^{641}$主知道你的需要.
但是最重要的是.
我們要先求神的國.
和神的義.
是不是?.
願我們一起都能夠先求神的國.
和神的義.
我們跟教友說一聲.
你要先求神的國和神的義.
這一切都會交給你.
基督徒每一天都有很多不同的選擇.
你是選擇站在神的華語當中.
神的應許當中.
神的價值觀當中.
還是選擇不同的價值觀呢?.
但是我們最蒙福的是選擇.
在神的華語當中.
最蒙福的是選擇.
在神的價值觀裡面.
好了.
全然的信靠.
毫無懼怕.
我們應該信得過神.
十多年前.
我覺得我們的身份其實是神給的.
我們信得過神.
十多年前.
在內地有一個城市.
其實有一些政策.
就是出了.
拆一些違規的建築.
叫做三國一策.
其實這本政策是很好的.
但是怎知道.
當中也牽涉到一些教會的建築.
甚至是十字架.
在當時那段時間.
都比較多人議論.
也都.
我相信政府.

$^{681}$當地的某些政府部門.
和教會之間.
可能都有大家改善的空間.
但是引申出來的是.
會令到.
可能.
在執行政策的時間.
可能會有機會.
令到一些信徒的.
信仰宗教感情.
受到傷害.
和國家的公共外交政策受到影響.
甚至可能有些人會誤解到.
國家的宗教政策.
有所改變.
在當時我都透過我自己的渠道.
有兩三次那兩年.
都反映了一些意見.
和一些有機會.
影響到宗教政策的官員.
是改善了一些,後來又好像還沒有處理到.
他們有一次在兩會期間.
我有一個感動.
豈不知你.
得著王浩的位份.
就是為了現今這個機會.
這個就是大家知道的.
以斯帖的經歷.
我覺得既然處給我這個位份.
我就憑信心祈禱.
希望能夠透過.
我的平台.
在兩會期間去反映意見.
當時我找了很多的牧者和我一起祈禱.
祈禱之後.
平安就上去表達意見.
結果是.
合規合法的方式.
去表達.
結果.

$^{721}$其實國家是同善如流.
亦都可能收到很多不同的意見.
結果很快,後來事件.
就慢慢處理了,但我得到一個訊息.
那個訊息就是.
國家宗教政策沒有改變,會處理.
但在這件事上.
那個學習是甚麼呢.
就是當時那個選擇.
你夠不夠膽去表演.
表達意見呢.
表達意見的時候會不會影響.
你將來的發展呢.
如果政策不同的話.
有沒有這個勇氣去做呢.
主讓我選擇了上好的福份.
因為.
豈不知你得到王雄的位份.
是為了現今的機會嗎.
我們經常祈禱.
求神讓我們能夠做一個.
不同的角色.
做一個流通管治的角色.
我們是.
將人帶到神的面前.
與神和好.
亦都藉著宗教基督.
我們是與人和好.
而且我覺得我們這個行動是愛神.
愛國愛港愛內地的教會.
在祈禱平安之後.
死就死吧.
進去吧.
結果後來.
更高層的.
反應是聽了.
意見後來這件事已經很快.
慢慢慢慢.
我相信不是我一個意見.
可能是整個很多方面的.

$^{761}$意見是.
處理了.
亦都算是.
平靜了.
其實主讓我們每一個人都有不同的崗位.
我們會不會學習.
在主的面前.
能夠站穩在整個話語裡面.
豈不知神.
豈不知我們今天得了皇后的位份.
是為了現今的.
機會嗎.
好了.
但以你的祈禱.
一日三次.
其實是一個恆常的祈禱.
這是一個.
很能夠令你感覺到.
就是一個.
要恆常的祈禱.
要恆切的.
我記得.
我早.
我記得我剛出來工作的時候.
當然我是中學的時候.
已經是被主的我的級別.
心裡很想奉獻給耶穌.
是呀 心裡很想奉獻給耶穌.
然後就.
學習守神庚.
早上祈禱 晚上祈禱.
結果我記得出來工作.
在銀行工作的第一份工作.
那兩三年 給我一個很好的學習.
就是打工的時候 你知道午餐被取消了.
結果吃完午餐之後.
有很少空間可以在自己的位置看聖經 祈禱.
一日三次.
一個學習.
後來我做了保險 做了行業就不行了.

$^{801}$但是在很多不同的環境.
有需要的時間 也都會在公司.
祈禱 還有.
甚至乎是回家祈禱.
特別是一些很需要 很大壓力.
或者有時候一些很需要.
決策的時間 決定重要的事情的時間.
來到聖經面前.
但我覺得一件事就是說 早上的祈禱.
是不可以缺少.
晚上的祈禱也都不可以缺少.
晚上的感恩 是不是?.
一日過了感恩.
主耶穌清神就去抗夜祈禱.
所以這幾句話就在我腦海裡面.
主耶穌清神就是抗夜祈禱.
但願我們都能夠.
每一天支取力量 從神而來支取力量.
但同樣我們需要什麼呢?.
充電寶一樣 要時刻.
所以你要時刻充電.
所以你時刻來到聖經面前.
我舉個例子 你要時刻充電.
即是支取從神而來的力量.
時刻來到主的面前.
一日三次行賽祈禱.
這就是在.
大義利的生活我們看到的.
第二就是祂雙膝跪在.
神的面前.
這是一個很特別的形容.
雙膝跪在神的面前.
你看到祂那種敬虔.
看到祂那種對神的謙卑.
看到祂那種.
在神的面前.
完全的一種撫摸.
這個我覺得是一個很值得我們學習.
我記得我中學的時間.
那幾年其實被我打下.

$^{841}$一個很好的屬靈的基礎.
就是放了學.
我上午班放了學 吃完飯.
還沒做功課前.
跪下祈禱 當時也有半小時至一小時.
祈禱,讀經等等的時間.
彼此建立了很好的一個屬靈的基礎.
我自己覺得.
跪下祈禱 到現在.
跪下祈禱其實是.
更加集中 更加能夠.
可以更長的時間 我的感覺.
所以這個也是一個很好的.
很重要的大義利.
讓我們看到.
不同的地方.
另外在祂的禱告裡面.
我們面對很多困難挑戰.
我們可能是埋怨神的 但是大義利是怎樣的?.
聖經經常教導我們 要常常喜樂不絕的禱告.
凡事在耶恩 因為你是神在.
基督耶穌裡 祂向我們鎖定的是什麼?.
主義.
所以其實我們要.
給我 設在我身上 我相信你的美意.
我們相信萬事都會互相效力 叫我神的人得益處.
就是我們主義被召的人.
萬能生隱 萬能生老年.
叫我們不至於羞愧.
我們真的很相信.
是能夠同樣的 讓我們能夠.
謝恩 這個是讓我們能夠.
更加在挑戰.
環境 困難環境的時候.
要常常喜樂 不住的禱告.
所以這個在新約的時候.
保羅也有這樣的教導.
你有用 英多姆掛慮.
將一切力量 所有的告訴 神所賜出的意願.
平安 必在基督耶穌.

$^{881}$保守裡面的心懷 意念.
從基督裡面的平安.
保守裡面的心懷 意念.
這個很重要.
所以就是 你看到.
但你也是 祈禱之後.
平安就臨到.
當然他的祈禱不是林吉抱耶穌的那種祈禱.
他不是林吉抱耶穌的.
他是什麼呢?.
與素常一樣.
他就是這樣的人 他就是這樣的生活.
所以這個.
是很值得我們大家.
學習.
另外他.
與素常一樣 剛才我說過.
但他的特點就是與素常一樣.
與素常一樣.
平時他就是這樣.
心理學家經常說.
心理學家就說.
要留意你的思想 思想影響你的言語.
要留意你的言語 言語影響你的行動.
你要留意你的行動.
你的行動會怎樣? 影響你的.
習慣 要留意你的習慣.
因為你的習慣會影響你的性格.
你的性格會怎樣? 影響你的命運.
同樣一樣.
其實習慣是很重要的.
人生很多時候 我們的習慣.
都是影響你自己 整個人生.
你的習慣好 你的人生就好.
沒有你的習慣不好 你的人生也會有很多問題.
習慣很重要.
同樣一個基督徒 同樣我們需要.
有好的習慣在神的面前.
去親近神.
親近神要成為我們生活的習慣.

$^{921}$你的生活裡面.
習慣親近神 其實都不需要說.
不需要問 不需要.
很特別.
其實就是與素常.
大聲一點說 與素常.
一樣.
與素常一樣.
我們親近神 與素常甚麼?.
一樣.
我們要學習.
與素常一樣.
好了 另外就是.
但以理.
其實我覺得他能夠被提升的原因.
其實有兩個 第一個.
我們剛才說了 他活在處女 與素常一樣.
第二個.
其實就是他分別為性的心慈.
他不以王的.
酒和錢.
去招呼自己.
但以理書一章八字 他是這樣說.
但以理卻立志不以王的善.
和王所賜的酒.
招呼自己 所以求太監長.
容他不.
招呼自己 因為.
當他們被擄到巴比倫的時候.
當時他和他三個朋友.
是被培養的對象 有聰明.
有智慧 這些年青人.
是.
被揀選出來的.
因為在他們進尾.
羊貓進尾 母蟬.
通達各種學問 知識 聰明具備.
王要訓練他們.
在王面前訓練 服侍王.
打算派他.

$^{961}$給他們.
三年的時間.
享受王的美食.
喝王的酒 和王的善一樣.
我想問你 今天你每天吃皇帝宴.
你吃不吃?.
可能你說 不是 我喜歡鹹魚白菜 都很好吃.
這個是真的 但問題是.
很多人說 美食當前.
可能你發覺.
有時候會忘記自己.
但是.
他們卻不容.
納至不以王的善.
和王所飲的酒 去折磨自己.
就像舊台監獎一樣.
他們不要折磨自己.
結果很奇妙的 十天之後.
他們比那些有吃王的善.
喝王的酒的少年.
更加進尾.
更加.
進尾肥胖.
然後那些委辦.
就撤去他們.
用的善 就給他們.
吃素菜 喝水.
這代表什麼呢? 就是他們有一個分別為勝的生活.
他不以王的酒.
和善去折磨自己 因為當時.
巴比倫王 他們吃的東西.
是經過祭祀的 經過祭祀的.
其實拜過偶像.
他要分別為勝 不吃拜偶像之物.
第二 他們所吃的東西裡面 有很多.
是違反了當時摩西的律法.
有馬 有豬 這些肉.
第三 更重要的是.
他不會和他們連結 因為當時的巴比倫.
王是很想將他們.

$^{1001}$能夠同化 學巴比倫的文字.
改他們巴比倫的名字.
喝王的善 將他們同化成為.
巴比倫人 但是他知道.
他是神的選民 神的兒女.
他站在神的話語當中 站在神裡面.
所以他不以王的善.
和王的酒.
來選擇.
水.
和素菜 結果.
更加勇貌俊美.
其實.
大禹李是給了我們一個很好的提醒.
也是一個很好的學習.
他只要稍微.
一說yes.
他不站在神的話語 他立刻榮華富貴.
分分鐘他可能還會.
在那三個總長之上.
因為找他自己傳過 本來那個王也有.
這樣的想法 但是他卻選擇.
在神的裡面 站在神的話語當中.
選擇他的.
根 他的信仰.
結果神更加興起他.
你看到他是四朝的元老.
直到六十年.
他在被擄之地.
能夠被大神使用.
最後我鼓勵我們.
一同學習主耶穌在.
哈西瑪利園和門徒祈禱的時候.
那些門徒睡覺.
所以主耶穌提醒他們.
總要警醒禱告 免得入了迷惑.
我們心靈固然願意.
但是肉體.
卻是軟弱.
守神金國也很願意.

$^{1041}$不過真的要五六點起床的時候.
或者七八點起床的時候 天氣又冷的時候.
會不會寒風陣陣的時候.
就不捨得呢.
這些我們要一同學習.
當你面對挑戰的時候.
究竟是不是站在神的話語當中.
還是站在世界價值觀當中呢.
這個也值得我們學習.
\newpage



\section{約翰福音 7:25-52}
\label{sec:032155}
\textbf{榮耀君王:活水靈風}~~~~林鍵牧師
\newline
\newline
連結: \href{http://www.church.com.hk/acms/content.asp?site=cdc\&op=show\&type=product\&code=032155\&layout=sermon}{\texttt{032155}} ~~~~ 語音日期: 2026-05-16
\newline
\newline
\hyperref[sec:032154]{\small{< < < PREV SERMON < < <}}
~
\hyperref[sec:index_chronic]{\small{[返順時目]}}
~
\hyperref[sec:032156]{\small{> > > NEXT SERMON > > >}}
\newline
\newline
約翰福音 7:25-52
\newline
\begin{longtable}{cl}
\hline
\hline
章節 & 經文 (和合本修訂版)\\
\hline
7:25 & \begin{tabularx}{0.7\textwidth}{X} 於是耶路撒冷人中有的說：「這個人不是他們想要殺的嗎？ \end{tabularx} \\ \\ \relax
7:26 & \begin{tabularx}{0.7\textwidth}{X} 你看，他還公開講道，他們也不對他說甚麼。難道官長真的認為這是基督嗎？ \end{tabularx} \\ \\ \relax
7:27 & \begin{tabularx}{0.7\textwidth}{X} 然而，我們知道這個人從哪裡來；可是基督來的時候，沒有人知道他從哪裡來。」 \end{tabularx} \\ \\ \relax
7:28 & \begin{tabularx}{0.7\textwidth}{X} 那時，耶穌在聖殿裡教導人，喊著說：「你們認識我，也知道我從哪裡來；我並不是憑著自己來的。但差我來的那位是真實的，你們不認識他。 \end{tabularx} \\ \\ \relax
7:29 & \begin{tabularx}{0.7\textwidth}{X} 我卻認識他，因為我從他那裡來，是他差遣了我。」 \end{tabularx} \\ \\ \relax
7:30 & \begin{tabularx}{0.7\textwidth}{X} 於是他們想要捉拿耶穌，只是沒有人下手，因為他的時候還沒有到。 \end{tabularx} \\ \\ \relax
7:31 & \begin{tabularx}{0.7\textwidth}{X} 但人群中有好些人信他，他們說：「基督來的時候，他所行的神蹟難道會比這人行的更多嗎？」 \end{tabularx} \\ \\ \relax
7:32 & \begin{tabularx}{0.7\textwidth}{X} 法利賽人聽見群眾對耶穌這樣議論紛紛，祭司長和法利賽人就打發聖殿警衛去捉拿他。 \end{tabularx} \\ \\ \relax
7:33 & \begin{tabularx}{0.7\textwidth}{X} 於是耶穌說：「我跟你們在一起的時候不會太久了，我要回到那差我來的那裡去。 \end{tabularx} \\ \\ \relax
7:34 & \begin{tabularx}{0.7\textwidth}{X} 你們要找我，卻找不到；我所在的地方，你們不能去。」 \end{tabularx} \\ \\ \relax
7:35 & \begin{tabularx}{0.7\textwidth}{X} 於是猶太人彼此問：「這人要往哪裡去，使我們找不到他呢？難道他要往散居在希臘的猶太人那裡去教導希臘人嗎？ \end{tabularx} \\ \\ \relax
7:36 & \begin{tabularx}{0.7\textwidth}{X} 他說『你們要找我，卻找不到；我所在的地方，你們不能去』，這話是甚麼意思呢？」 \end{tabularx} \\ \\ \relax
7:37 & \begin{tabularx}{0.7\textwidth}{X} 節期的最後一天，就是最隆重的一天，耶穌站著，喊著說：「人若渴了，到我這裡來喝！ \end{tabularx} \\ \\ \relax
7:38 & \begin{tabularx}{0.7\textwidth}{X} 信我的人，就如經上所說：『從他腹中將流出活水的江河來。』」 \end{tabularx} \\ \\ \relax
7:39 & \begin{tabularx}{0.7\textwidth}{X} 耶穌這話是指信他的人要受聖靈說的；那時還沒有賜下聖靈，因為耶穌還沒有得到榮耀。 \end{tabularx} \\ \\ \relax
7:40 & \begin{tabularx}{0.7\textwidth}{X} 眾人聽見這些話，有的說：「這真是那先知。」 \end{tabularx} \\ \\ \relax
7:41 & \begin{tabularx}{0.7\textwidth}{X} 另有的說：「這是基督。」但也有的說：「難道基督是出自加利利嗎？ \end{tabularx} \\ \\ \relax
7:42 & \begin{tabularx}{0.7\textwidth}{X} 經上不是說『基督是大衛的後裔，出自大衛的本鄉伯利恆』嗎？」 \end{tabularx} \\ \\ \relax
7:43 & \begin{tabularx}{0.7\textwidth}{X} 於是眾人因耶穌而分裂了。 \end{tabularx} \\ \\ \relax
7:44 & \begin{tabularx}{0.7\textwidth}{X} 其中有人要捉拿他，只是沒有人下手。 \end{tabularx} \\ \\ \relax
7:45 & \begin{tabularx}{0.7\textwidth}{X} 警衛們回到祭司長和法利賽人那裡。他們對警衛說：「你們為甚麼沒有帶他來呢？」 \end{tabularx} \\ \\ \relax
7:46 & \begin{tabularx}{0.7\textwidth}{X} 警衛回答：「從來沒有像他這樣說話的！」 \end{tabularx} \\ \\ \relax
7:47 & \begin{tabularx}{0.7\textwidth}{X} 於是法利賽人說：「你們也受了迷惑嗎？ \end{tabularx} \\ \\ \relax
7:48 & \begin{tabularx}{0.7\textwidth}{X} 難道官長或法利賽人中有信他的嗎？ \end{tabularx} \\ \\ \relax
7:49 & \begin{tabularx}{0.7\textwidth}{X} 但這些不明白律法的眾人是被詛咒的！」 \end{tabularx} \\ \\ \relax
7:50 & \begin{tabularx}{0.7\textwidth}{X} 其中有尼哥德慕，就是從前去見過耶穌的，對他們說： \end{tabularx} \\ \\ \relax
7:51 & \begin{tabularx}{0.7\textwidth}{X} 「不先聽本人的口供，查明他所做的事，難道我們的律法還定他的罪嗎？」 \end{tabularx} \\ \\ \relax
7:52 & \begin{tabularx}{0.7\textwidth}{X} 他們回答他說：「你也是出自加利利嗎？你去查考就知道，加利利是不出先知的。」 \end{tabularx} \\ \\
[1ex]
\hline
\hline
\end{longtable}
$^{1}$歡迎大家來到我們的周六崇拜.
上年七月我們開始研讀新約聖經中的約翰福音.
我們講到系列的主題是榮耀君王.
我們花了半年的時間完成了約翰福音一至四章的宣講.
今年一月我們開始研讀約翰福音第五章.
死亡學者告訴我們.
從約翰福音第五章開始到約翰福音第十章.
約翰帶我們進入一個節慶的週期.
第五章是安息日.
第六章是逾越節.
第七章至第九章是鑄棚節.
第十章有的背景是獻殿節.
在這個節慶週期裡面.
耶穌以節期為背景.
通過祂的神蹟和教導.
去啟示祂自己的身份和使命.
我們會說祂越是清楚啟示祂自己.
猶太的領袖就越是恨惡要逼迫祂.
所有的一切其實都是從約翰福音第五章.
安息日的風雲開始.
約翰福音第五章第十八節.
為了這個緣故耶穌行了神蹟.
但這個神蹟是在安息日行的.
猶太人越想殺祂已經想殺祂了.
第五章因為祂不但犯了安息日.
而且稱上帝為他的父.
把自己和上帝看為同等.
然而正是因為耶穌甘願替罪人死.
他們想殺耶穌.
但耶穌也甘願替罪人死.
想十字架.
所以叛逆的猶太人才可以有機會得著救贖.
醫治和最後上帝裡面的安息.
在這裡有兩點我希望你特別留意.
第一雖然在約翰是重複又重複強調.
耶穌的時候還沒到.
但是我們來到約翰福音第七章.
其實距離耶穌上十字架只剩下六個月的時間.
所以你去到第七章的時候.
你會感覺到耶穌是特別的焦急.

$^{41}$祂的呼喊.
希望那些人真的認識祂.
希望祂的人來祂的身邊.
因為唯有祂才可以拯救他們.
所以第一我想你看到的是.
我們現在距離十字架只剩下六個月.
第二約翰福音第五章至約翰福音第七章.
中間相隔了大概有十八個月的時間.
然而耶穌對約翰福音第五章的陳年舊事.
仍然是念念不忘的.
我們現在讀的是第七章.
但是耶穌對十八個月前發生的事情.
是念念不忘的.
在上個星期的講道裡面.
我們看到耶穌如何去舊事重提.
祂重提十八個月前祂在靠近陽門的不是大池.
醫治了一個騎腿三十八年的人.
使這個人可以走進安息.
是安息日行的神蹟.
那件事情.
祂再重提那件事情.
然後將那件事情連結在列祖的國禮的盟約當中.
你看看你記不記得上個星期.
我給你看那個文學架構.
前後呼應的文學架構.
整個文學架構的中心在哪裡.
在中間的「伊」.
耶穌回答說.
安息日祂做了兩件事.
一個行國禮使人存余.
另外一個是在安息日祂做了一個使一個人存余.
起來可以行走.
然後走進安息.
來到這裡其實你會發現一年半之前.
約翰福音第五章裡面.
那個故事其實是後來的故事.
在一年半之前的第二章.
終於在第七章裡面.
耶穌加插了一個第二章.
那個第二章講到安息日是可以行國禮的兩件事.

$^{81}$於是將兩件事株連璧合.
合而為一.
去表明耶穌的身分和使命.
耶穌是上帝認祂揀選的君王救主.
祂來是要使聖約國禮人.
存余的走進安息的當中.
如果你用猶太人的角度來講.
你開始記得.
那個騎腿三十八年.
三十八年代表以色列人在曠野三十八年.
沒辦法進入上帝認祂之地的三十八年.
這個騎腿的以色列人三十八年.
這個騎腿的人其實是受過國禮的.
一出生第八天已經受了國禮.
然後等到現在耶穌終於來拯救.
你會知道如果你用一個以色列人的眼光去看整件事.
Part A和Part B.
國禮是上帝恩典蒙約的認許.
安息是上帝恩典蒙約的承傳.
這個人受國禮.
這個人沒辦法活動.
這個人三十八年好像舊約的以色列人.
沒辦法進入上帝的安息.
而耶穌來要拯救這個三十八年.
騎腿不能走動的人.
終於可以起來復活.
進入上帝的安息.
所有的東西都是一個圖畫.
一個福音的圖畫.
上帝要來拯救以色列人.
讓他們都起來進入上帝的安息.
所有的一切出於耶和華上帝.
所有的一切要承傳在耶穌的十字架上.
然而還沒說完.
還有很多很多的東西.
今天我們繼續說完.
耶穌拯救完之後還有下文.
或者拯救有一個很大的延伸.
就是聖靈的降臨.
我想帶你回去看.

$^{121}$《讓步音》第七章裡面的大圖畫.
從故事的時序來說.
《讓步音》第七章可以分為三大段.
第一 猶太人駐彭捷近了.
第二 捷旗已經過了一半.
一開始的時候耶穌沒有去.
祂去到一半.
一半就是第四天左右.
然後祂才出現.
然後去到捷旗的最後一天.
今天我們已經來到捷旗的最後一天.
我們要研讀的是捷旗最後一天的Part A.
是什麼呢?.
耶穌做了一個自我的宣告.
我就是生命的活水.
當然祂沒有說這句話.
但是我們將祂那句話.
那三節聖經總結的時候.
耶穌其實在說祂是誰.
說祂是生命活水.
是兩個耶穌在駐彭捷裡面.
重要啟示的第一個.
第二個就是.
「我是世界的光」.
這個就是耶穌指名道姓.
說了「我是世界的光」.
但是我們今天要看的是.
「我是生命活水」.
在我們還沒有進入今天的經文之前.
我想帶你回去看舊約聖經裡面.
駐彭捷的設立.
我帶你回去看《利未記》第23章.
《利未記》第23章.
我們從33節看到36節.
耶和華吩咐摩西說.
你要吩咐以色列人說.
這7月15日是駐彭捷.
要向耶和華守這節7天.
接著第35節.
第一天當有聖會.

$^{161}$任何勞動的工作都不能做.
要將火祭獻給耶和華7天.
第8天當守聖會.
並要獻火祭給耶和華.
這是嚴肅會任何勞動的工都不可以做.
我想你留意幾件事.
第一.
猶太人的駐彭捷是在他們的月曆裡面.
他們的月曆裡面是7月15日.
大概是我們的月曆裡面的9月至10月之間.
第二.
無論節期的第一天落在那週的哪個日子.
都會變成做一個安息日.
如果剛好是安息日.
那就是安息日.
但是如果他是星期一,二,三.
那天只要他第一天落在那日子.
那日子突然間會變成做安息日.
那天會有聖會.
任何勞動的工都不可以做.
但是更有趣的.
不單止是第一天.
第八天也是一樣.
無論第八天落在哪一天.
都會變成做安息日.
當守聖會.
這是嚴肅會任何勞動的工都不可以做.
讀到這裡你會開始問.
究竟駐彭捷是七天還是八天呢?.
你橫看怎麼看.
第八天好像很重要.
第八天都變成做安息日.
但是他又說駐彭捷是七天.
究竟七天還是八天?.
這個也是聖經學者考量的其中一樣東西.
待會我們說讓福音的時候.
你就會明白七天八天的分別在哪裡.
我想帶你回到今天的經文.
剛才我們讀的經文都很長.
由第25節開始一直讀到第52節.

$^{201}$這篇經文你問我是覺得有點雜亂的.
組織起來是比較困難的.
我自己就嘗試勾畫了一個故事的輪廓給大家看.
這個前後呼應的架構是我做的.
學者不一定會接納我做的這個前後呼應.
但是這個前後呼應是幫助我自己去理解整個故事的.
他們未必同意這個架構.
但是所有的學者基本上都同意一樣東西.
就是故事的核心在哪裡呢?.
就是中間的D.
故事的核心就是第七章第37至39節.
我講一下我這個架構的想法.
上面最開面的層是A和A.
法利賽人聽群眾紛紛議論.
下面的是法利賽人聽尼哥德姆為耶穌辯護.
接著BB警衛打發去要捉拿耶穌.
B警衛回來說沒有人像他.
警衛不捉他.
接著說沒有人像他.
接著CC兩個問題的.
耶穌要往哪裡去他們彼此問.
接著下面的C是耶穌從哪裡來他們分裂了.
而在整個架構的中間就是D.
是最重要的耶穌站起來呼喊大聲說.
說的是活水.
實際接著約翰告訴我們.
其實他說的是勝利.
但是這幅圖畫裡面.
我們彷彿看見耶穌是站立在人群當中去吶喊.
接著旁邊有很多人.
可能分為兩個群組.
三個群組四個群組.
有人說他是一個好人.
有人說他是基督.
有人說他是先知.
但也有其他人說不是的.
他不是的.
於是耶穌自己獨立雞群.
他站起來大聲呼喊.
旁邊的人不斷的議論耶穌到底是誰.

$^{241}$接著他們問值不值得我們信靠他.
值不值得我們委身給這個叫耶穌的人.
或者我們當中也有些人.
也是這樣想的.
你在聽耶穌.
接著你在問.
究竟耶穌值不值得我去奉上一切.
All in 去跟隨他.
所以我邀請你今天也聽一下.
耶穌是怎樣宣告他自己的身份.
以及他的使命.
但是我們進入我們今天要看三節聖經.
其實是三節聖經的第一節.
《約翰福音》第七章第三十七節.
節期的最後一天.
就是最隆重的一天.
耶穌站起來.
接著大聲喊著說.
這個是耶穌宣告的背景.
背景是什麼呢.
背景是駐彭捷.
是駐彭捷的最後一天.
聖經學者告訴我們.
七日節期的每一天.
在駐彭捷的時候.
我們先當它是七日.
駐彭捷有七日.
每一天大祭司都會手拿著一個金水壺.
或者金酒壺.
他會帶領一群人遊行的隊伍.
從一個地方叫做西羅亞池子.
西羅亞池子是在聖殿的南邊.
你會發現西羅亞池子.
遲些在第九章又會再出現的.
那個池子和耶穌有關係的.
西羅亞的意思是奉差遣的.
他要去西羅亞的池子那裡拿一些水.
接著經過了水門.
來到聖殿裡面.
聖殿的祭壇.

$^{281}$接著倒在祭壇裡面.
這個是大祭司需要做的事情.
但是大祭司不是自己一個人去做的.
一群人一直跟著他.
大祭司拿著金水壺.
接著很多人在那裡依依掛掛.
跟著他一群的祭司.
接著其他的來耶路撒冷過節的人.
就更加圍著他們.
當他經過水門的時候.
要有三響的喜樂的號角.
砰砰砰.
就是我不知道怎麼吹.
但是他說是很開心的喜樂.
接著到他們來到聖殿.
在祭壇上面又要繞幾個圈的.
在當時聖殿裡面就有一群師班.
他們會高唱Hallelujah.
不只是說Hallelujah這麼簡單.
所謂的Hallelujah在舊約聖經裡面.
就是詩篇第113篇至118篇.
六篇的詩篇.
所以聖殿裡面的師班.
不只是唱Hallelujah那些.
是唱舊約的詩篇113至118.
他們唱到詩篇118的時候.
以色列的吳南丁.
右手就拿著樹枝在搖動.
接著左手拿著當地的水果.
接著他們要一起同聲吶喊三次.
要稱謝耶和華.
Give thanks to the Lord.
接著才倒水到聖殿的祭壇上面.
有聖經學者告訴我們.
當他們將水倒到祭壇的時候.
他們會加上酒祭.
原來聖殿裡面有酒祭的.
所以不只是倒水.
倒了酒.
你知道酒是什麼顏色的.

$^{321}$酒當然是紅紅的.
是葡萄酒.
於是有水有酒一起倒到聖殿的祭壇上.
所謂祭壇是什麼呢?.
祭壇就是用來殺動物的.
燒動物的地方.
所以有水有酒.
有聖殿祭壇裡面流進去.
也從聖殿祭壇流出來.
聖殿祭壇又有血.
你想想.
其實你熟悉楊漢福音.
你也知道這是一個十字架的圖畫.
擺明了以後耶穌.
上帝的羔羊在哪裡殺的呢?.
就在這個祭壇裡面.
又有水又有血.
就是十字架的圖畫.
所以你已經開始覺得.
你只是想一個背景.
所有東西已經將你引到耶穌的十字架裡面.
又有學者說.
來到第七天的時候.
同一樣的水禮將會重演七次.
所以你頭六天就是做一次的水禮.
到了第七天.
大祭司也很辛苦.
要走七次.
接著他們相信耶穌的宣告.
很有可能是七次的最後一次.
就是整個節期的高峰.
就是最後一次了.
七日最後.
七七.
就是最後的第七天的七次水禮.
最後最高峰的時候.
接著耶穌就呼喊出來.
就是下面那一句.
這麼說.
不過也有學者說.

$^{361}$如果你第七天呼喊.
哪裡有人聽到呢?.
所以有另一個學者說.
其實這裡所說的最後一天.
未必是第七天.
可能是第八天.
也有它的道理.
因為第八天是沒有水禮的.
接著在第八天已經沒有水禮的時候.
耶穌起來說.
我替代了水禮.
這個也有它的道理在這裡.
所以學者說第八天.
雖然不是.
可能已經包在節期裡面了.
你記得舊約聖經也是說第八天.
第八天也是安息日.
我們可以叫它做擦鵬日.
因為鑄鵬節第八天是用來擦鵬的.
雖然沒有水禮了.
但是其實大家仍然是興高采烈.
他說第八天.
所有的以色列藍丁.
一樣在高唱哈利路亞頌.
而且是重覆又重覆重覆.
不用choir了.
不用試班了.
我們自己去唱.
於是.
有些學者覺得.
他是在第八天.
當所有水禮已經結束的時候.
耶穌自己取替了水禮.
祂走出來.
今天沒有水禮.
但是我今天要成為那個水禮.
於是耶穌宣告.
我就是水禮的承傳.
我就是以色列救贖盼望的達成.
OK.

$^{401}$又有道理的.
無論你第七天也好.
無論是你第七天裡面七次水禮之後.
接著大家靜下來的時候.
耶穌再說出來.
如果當時祂立刻大聲說.
可能人聽不到.
還是第八天也有祂的道理.
無論是第七天.
或者第八天也好.
耶穌所宣告的不是一時一刻的水禮.
祂說我要做一個水禮.
但是這個水禮不是一時一刻.
不是好像你們的水禮一做完之後就沒有了.
我要做一個源源不絕.
永恆不息的.
活水江河的水禮.
OK.
這個帶我們.
去到耶穌所說那句很大的謊話.
不是大話精的大話.
是很大很大不可思議的立限.
耶穌說什麼呢?.
祂說人若喝了.
到我這裡來.
信我的人就如經上所說.
從他的福中將流出.
活水的江河來.
這個是一個很驚人的宣告.
我們問這個宣告究竟是說什麼呢?.
約翰.
作者約翰嘗試為我們解謎.
解謎在哪裡呢?.
下一節.
第七章第39節.
這個是約翰給我們的解謎.
耶穌沒有說的.
耶穌只是吶喊了.
人若喝了.
到我這裡來.

$^{441}$信我的人.
如經上所說的.
從他的福中將流出.
活水的江河來.
OK.
真的很大的謊話.
接著約翰就解釋了.
耶穌這話是指.
信他的人要領受聖靈所說的.
那時.
還沒有賜下聖靈.
因為耶穌.
還沒有.
得到榮耀.
約翰說.
所謂的活水江河.
其實指的是聖靈的降臨.
約翰說.
聖靈的時代.
新時代將要.
快將要展開.
OK.
我之所以這樣說的原因是.
因為.
如果你熟悉舊約.
你知道舊約聖經裡面.
都有聖靈.
舊約聖經裡面.
有聖靈的工作.
舊約聖經裡面.
信徒都有聖靈的經驗.
不過.
約翰現在所說的.
或者耶穌所說的.
約翰正在幫他翻譯出來.
是一個前所未有的.
新時代.
聖靈將要全面降臨.
在每一個屬於上帝的人身上.
再不是.

$^{481}$在舊約裡面.
是一小撮領袖級的人才有.
摩西.
約書.
大衛.
這樣.
他們領導級的人.
才有聖靈的經驗.
但是現在.
從最小的到最大.
從最體弱.
很殘的人.
去到最強健的人.
都可以領受同一樣的聖靈.
聖靈會同在.
聖靈會與你有一個聖靈的經驗.
OK.
而且這個經驗.
是沒有時限的.
打遊戲機也是.
五分鐘就停了.
又要再給錢了.
這個經驗現在是沒有時限的.
再不是做完工作.
然後聖靈就回歸上帝.
而是聖靈會降臨.
在這個人的身上.
一生一世.
無時無刻的與這個人同在.
你說你在說什麼.
令我想起.
《行傳》第二章裡面.
彼得那篇道裡面.
講到舊約.
先知約矣的預言.
我讀一讀給你聽.
你就明白了.
這正是.
藉著先知約矣所說的.
這個是聖靈降臨1.

$^{521}$彼得講那篇道的時候所講的.
上帝說在末後的日子.
我要將我的靈.
囂冠在凡血肉之軀的.
你們的兒女要說預言.
你們的少年要見異象.
你們的老人家要做異夢.
就是什麼都包了.
小的中年的.
年輕的老人家.
全部都一樣.
你說說預言.
見異象.
發異夢.
其實都是同一樣東西.
因為有上帝的福音.
上帝的天國臨在.
大家都在想同一樣東西.
在那些日子.
我要把我的靈囂冠.
甚至給我的僕人和婢女.
連阿姐也有.
全部人都有.
個個都有份.
不是某幾個人做領袖的才有.
他們都要說預言.
聖靈的時代將要展開.
聖靈的時代.
從耶穌第一次的彌臨展開.
一直延續到耶穌第二次再回來.
整個時代.
成為了一個叫做聖靈的時代.
然而.
但是.
你必須要知道.
那時還沒有刺客聖靈.
因為耶穌還沒有得到榮耀.
那時的意思是約翰福音第七章.
就快了.
還有六個月.

$^{561}$但是約翰福音第七章.
主彭節.
耶穌大叫的那一段.
真的是謊言.
其實是沒有實際的.
因為聖靈還沒有來得及.
因為耶穌還沒有得到榮耀.
在約翰福音裡面.
耶穌的榮耀是十字架.
很快很快就有了.
不過現在還沒有.
還有六個月.
耶穌就要替罪人死去贖罪.
還有六個月.
他受死之後.
他會復活得勝.
再過六個月.
再加四十天.
他會升天.
升天作王.
然後他會派遣他的聖靈.
你可以計算.
我們還剩下六個月.
七個月左右.
耶穌會得到榮耀.
他的榮耀包括了他的十字架.
他的榮耀包括了他的復活.
他的榮耀包括了他的升天.
但是所有的東西.
從十字架開始.
如果他上十字架.
聖靈就可以開始來到我們的中間.
《水學》寫給我們知道.
駐彭捷和其中的禮儀.
剛才跟你說了很多關於水禮的東西.
其實那些禮儀是裝載了.
以色列人對末世救贖的終極盼望.
那些禮儀是全部背後有神學的理念.
有些人會說.
水因為他們在中東很缺水.

$^{601}$所以水是求上帝降雨等等.
沒錯.
這也是對的.
對於以色列人來說.
熟習舊約聖經的以色列人來說.
那些禮儀背後帶著的是屬靈的意義.
耶穌自己深深明白.
駐彭捷對以色列人的意義.
因為他熟讀舊約聖經.
而這也是他在世上最後一次的駐彭捷.
因為還有六個月他就再見了.
他不會再有下一年的駐彭捷.
他不在了.
所以你可以想像耶穌當時.
內心深處那種激烈的情感.
於是耶穌起來說.
我來為要成全駐彭捷.
其實他說的是.
人要渴要來要什麼什麼福中什麼江河.
所有的東西.
其實他說的是.
我來是要成全駐彭捷中.
上帝子民所迫切渴望的活水.
駐彭捷是說水的.
我就是那個活水.
而聖靈降臨.
聖靈時代的開啟.
就是那個活水盼望的成全.
我想帶你去看兩段舊約聖經.
兩段舊約聖經.
這兩段舊約聖經.
是當時的猶太拉比.
不只是耶穌.
是當時的猶太拉比.
認為是和駐彭捷水利息息相關的末世寓言.
不是只有耶穌自己在想.
當時的猶太人都知道.
一想起駐彭捷裡面的水利.
就想起這兩段舊約的聖經.
第一段大家是比較熟悉的.

$^{641}$是以西結書第47章.
其實嚴謹來說.
我們應該讀完1至12節.
因為時間關係.
我們只可以讀一節.
那一節也已經是那個寓言的精髓.
以西結書第47章第一節.
看啊.
這是一個意象.
不是真實的.
沒有紀錄片的.
這個不是紀錄片.
看啊.
在一個意象裡面.
有水從聖殿的門檻下.
下面往東流出.
又是從祭壇的南邊流出.
所謂從祭壇南邊.
就是和祭壇很近.
不過在祭壇的旁邊流出.
這段聖經.
其實我們之前說約翰福音第4章的時候.
耶穌拜訪聖瑪利亞的講道裡面.
我已經有和大家仔細讀過.
這裡我只可以略略和你思想.
整個意象的中心是耶路撒冷的聖殿.
不是紀錄片.
是一個意象.
在末後的日子裡面.
聖殿會有活水流出.
從聖殿流出.
而在聖殿流出的時候.
是經過聖殿那個祭壇的旁邊.
接著往東流.
什麼叫做往東流呢?.
東是墮落世界的代表.
因為亞當夏娃被趕出伊甸園的時候.
是向東邊走的.
所以這個聖殿的活水.
是要將拯救帶給一個失落墮落的世界.

$^{681}$從聖殿裡面的祭壇旁邊流.
你會想會不會沾了一些血.
就是了.
That's the point, right?.
就是有水帶著血.
去到一個墮落的世界.
在這裡.
活水的江河.
是生命和豐盛的代號.
我讀兩節以西結書第47章給你聽.
第一個是.
看那河那邊的岸有極多的樹木.
就是它這個河.
生命的江河.
去到哪裡.
在河的兩岸.
很多的樹木就會繁衍.
這個代表了生命.
代表了豐盛.
接著更有趣的.
第47章第8節.
這水往東方流.
直到海裡面.
那個海是死海.
所流出來的水.
就是進了海之後.
一進海就變成淡水.
這樣也可以.
海裡面的水.
是有鹽的.
鹽是會弄死人的.
是沒有魚可以在那裡生存的.
然而這個江河去到海的時候.
是可以讓鹹水變淡水.
你覺得你講.
這個整個意象告訴我們.
那個生命的活存.
代表了一種生命的力量.
上帝的豐盛.
這個是上帝天國降臨的圖畫.

$^{721}$這幅圖畫剛剛好連接了.
如果你還記得.
駐彭哲的水禮.
有水和酒祭.
接著從聖殿祭壇流出來.
也連上了耶穌的十字架.
你已經想到.
又有水又有祭壇旁邊的血跡流出來.
彷彿就是十字架的前瞻.
這是一個何等榮耀美麗的圖畫.
這個圖畫會在耶穌的駐彭哲.
成全的裡面得到成全.
不過這個不是唯一的圖畫.
還有另外一個舊約的寓言.
和駐彭哲的水禮息息相關.
那個寓言我更加喜歡.
我帶你去看.
只有兩節.
撒加利的書第十四章第八節和第九節.
我讀第八節給你聽.
在那一天必有活水從耶路撒冷出來.
和以西結書很相似.
一半往東海流一半往西海流.
這個就不只是東方.
總之東西南北全部都去了.
東下都是如此.
但是一定要讀第九節.
第九節是什麼呢.
第九節是這樣說.
耶和華要作全地的主.
那一天耶和華必為獨一無二.
他的名也是獨一無二.
你讀這個寓言的時候.
你不要只是讀第八節.
你要第八節和第九節.
兩節一起連起來你才明白.
如果你就這樣讀第八節.
那就和以西結書很相似.
但是第九節.
給我們看到一些更深層的東西.

$^{761}$如果你是一個忠於聖約的以色列人.
在舊約的時代裡面.
這個是一個榮耀的意象.
榮耀的盼望.
撒迦利亞書第十四章第八節.
是駐彭哲水禮的成全.
聖殿裡面有水流出來.
越流越多.
流向全世界.
第十四章第九節.
都是水禮成全的一部分.
所以不要看漏第九節.
第九節說的是什麼呢.
這兩節聖經加起來意味著什麼呢.
聖靈會降臨.
因為水代表聖靈.
聖靈的時代要開展.
接著要做什麼呢.
耶穌基督的福音要廣傳.
接著全地東南西北所有的人.
聽了這個福音.
領受福音的人.
要俯伏勁大.
以至耶和華作王.
獨一無二.
什麼叫獨一無二呢.
如果你去讀英文.
The Lord is One.
一.
我們中文翻譯叫獨一無二.
其實你可以翻譯成獨一.
什麼意思呢.
它是帶我們回到生命的第六章.
Shema O Israel.
最大的誡命.
以色列人你們要聽.
耶和華.
我們的上帝是獨一的主.
你明白嗎.
當這個生命的江河流向全世界的時候.

$^{801}$人們要聽到這個福音.
要俯伏勁大.
而俯伏勁大出來的那句話是什麼呢.
耶和華是全地的王.
耶和華是獨一的.
這些人聽到福音的人.
真的俯伏勁大的人.
會成為耶和華的敬拜者.
這個是那個意象.
耶穌說我來.
我要猜我的聖靈來.
要承傳這個偉大的意象.
這個偉大的意象是什麼呢.
就是全地要充滿著耶和華的敬拜者.
你想想這個時候.
耶和華的敬拜者只是在以色列人當中有.
而且只有那幾個.
以色列人是很廢的.
整群人幾百萬人其實沒有幾個敬拜者.
但是現在這個意象.
這個末世意象說有福音要傳到地極.
全地都要充滿著耶和華的敬拜者.
不是該人那種隨意獻奢濟的敬拜者.
你明白嗎.
不是說那種.
那種不是Jemah Yisrael.
不是以色列人要聽那種.
不是盡心盡性盡意愛耶和華那種.
它這裡的意象是.
全地會呼喊耶和華是獨一的.
The Lord is One.
全地會盡心盡性盡意愛耶和華.
全地都是阿伯一般獻賞.
最好獻賞所有的敬拜者.
這個是駐彭治的意象.
想都想不到吧.
玩到這麼大.
當你有活水的時候.
活水的效應就是心受割禮心愛主.
活水代表了聖靈的降臨.

$^{841}$聖靈的工作.
聖靈的工作使人做什麼.
心受割禮心愛主.
我想帶你回到約翰福音.
在約翰福音裡面.
其實這個聖靈時代的來臨.
已經是多次被預告.
約翰福音第一章第三十二節.
施洗約翰又作見證說.
我曾看見聖靈彷彿甲子一樣降下.
停留在耶穌的身上.
先前我不知道耶穌是什麼水.
可是那差我來的上帝.
用水施洗的對我說.
叫我用水施洗的對我說.
如果你看到有聖靈降臨.
停留在誰身上.
那個人就是將來會用聖靈施洗的.
我看見了施洗約翰福音.
所以我現在作見證.
他這一位是上帝的兒子.
約翰福音在預表將來.
他說上帝的兒子耶穌來了.
他最大的使命要做什麼.
他要用聖靈替我們施洗.
當耶穌用聖靈替我們施洗的時候.
耶穌會開啟一個聖靈的新時代.
聖靈要賜我們力量.
將天國的福音.
第一聖靈會使我們心受割禮心愛主.
這是很重要的.
第二聖靈會給我們力量.
將天國的福音傳到地極.
因為是生命的光學.
這個光學要傳到地極.
拯救罪人.
復活死人.
更生上帝的子民.
讓全地充滿耶和華的敬拜者.
聖靈來最重要最偉大的意象是什麼.

$^{881}$不是幫你考試.
真的麻煩你.
不是你心驚膽跳的時候.
聖靈給我平安.
不要玩這些.
真的沒有關係.
聖靈來最大最大的意象.
是要全地充滿耶和華的敬拜者.
我再說一次.
不是該人一般隨意獻奢東西的敬拜者.
而是阿伯一般獻賞.
最好獻賞所有的敬拜者.
這個是聖靈的工作.
我們再看.
第三章第五節第六節.
接著三十四三十六節.
耶穌回答說.
我實實在在告訴你.
人若不是從水和聖靈生的.
就不能進上帝的國.
從肉生生的就是肉身.
從靈生的就是靈.
聖靈來要使我們重生.
聖靈來要我們深受國禮.
深愛主.
第三十四節.
上帝所差來的說上帝的話.
首先是說耶穌的.
所賜給他的聖靈是沒有限量的.
上帝給耶穌的聖靈是無限量的.
但我想問你.
這個聖靈拿來做什麼.
原來是和耶穌宣講上帝話有關係的.
接著說.
信旨的人有永生.
不信旨的得不到永生.
上帝的憤怒在這個人的身上.
在這裡.
其實背後的道理和邏輯是一樣的.
耶穌要將他無限量的聖靈和他的門徒分享.

$^{921}$這個門徒包括你和我.
如果你真信耶穌的時候.
耶穌會將他無限量的聖靈和我們分享.
分享我們做什麼.
分享以至我們都可以和他一起.
一樣忠心宣講上帝的話.
你明白嗎.
上帝給我們這個聖靈做什麼.
不是出入平安.
這些人真的很傻.
以為有聖靈保住保身.
出入平安身體健康.
聖靈護體.
傻的嗎.
完全不是這樣.
上帝要將聖靈與我們分享主要的目的.
以至我們可以忠心宣講他的話語.
我們明白他的話語.
我們宣講他的話語.
以至我們可以拯救罪人.
復活死人.
更生上帝的子民.
最後的意義是什麼.
全地要充滿耶和華的敬拜者.
再說一次.
不是該引一般隨意獻奢東西的敬拜者.
那些敬拜者不算數.
而是阿伯一樣.
獻賞最好.
獻賞所有的敬拜者.
讓我帶你回到我們那三節的聖經.
《約翰福音》第七章第37節這麼說.
耶穌說.
人學了道我者對來.
信我的人就如經上所說.
從他的福中將流出活水的江河來.
原來這個宣告在希臘原文裡面.
視乎你怎麼擺那個標點符號.
是有兩個不同的翻譯的可能.
我叫做一個A和B.

$^{961}$A就是和合本聖經的翻譯.
人學了道我者對來.
信我的人就如經上所說.
從他的福中將流出活水的江河來.
問題來了.
福中活水的江河是從誰身上流出來.
基於這個翻譯.
其實是基於信耶穌那個人.
因為他信我的人.
信我的人就如經上所說.
從信我的人的福中將要流出活水的江河.
所以我們現在讀的翻譯.
正常來講.
就要這樣讀.
就說哪個肚子裡面.
或者心裡面.
會有活水的江河.
是信主的人.
但是你現在看一下.
B有另外一個可能的翻譯.
看看你怎麼擺那個標點符號.
是這樣的.
人若學了.
讓他道我者對來.
我就是耶穌.
他若信我.
讓他學.
接著就是句號了.
接著下面那個就如經上所說.
從他福中將流出活水的江河來.
我們再問同一個問題.
他福中活水江河.
從誰的福中流出來.
你就這樣看.
開始不一定是那個信主的人.
他說人若學了.
道我者對來.
如果他信我.
讓他喝.
接著是句號.

$^{1001}$接著從他的福中.
這個福中是耶穌的福中.
所以你看那個演繹有點不同.
前面A那個翻譯的福中.
是信主的人福中會有活水.
這裡的福中說.
要讓他學.
學什麼.
學耶穌福中所流出來的活水.
你看到那個分別嗎.
所以有兩個不同的翻譯.
這兩個不同的翻譯.
都有不同的聖經學者支持.
所以很難選擇.
不過無論你是選A或者選B.
從聖經的教導裡面.
兩個都通.
兩個都對.
從約翰福音裡面.
兩個都對.
你沒有說一個是錯的.
另一個是對的.
如果你真的要逼我選.
我也是選不到.
我覺得難以取捨.
真的要我選.
我的傾向是both end.
兩個都對.
故意含糊.
有時候聖經就是這樣.
一個含糊的一句話.
有兩個不同的解釋.
兩個解釋都是對的.
於是如果兩個都對.
從救贖邏輯來說.
就是from B to A.
總要有一個先行的.
誰先行.
耶穌先行.
從耶穌的活水聖靈.

$^{1041}$從耶穌流出來.
先由耶穌的肚子那裡.
咕嚕咕嚕的出來.
接著流向信徒那裡.
來喝東西的人.
接著流進他們的體內.
於是耶穌的信徒.
都裡面咕嚕咕嚕.
又有活水的泉源.
活水的江河流出來.
聖學者說.
就如經上所說的.
你要去找舊約聖經.
但是這裡沒有一段.
舊約聖經是特別說這件事.
無論你翻譯是A還是B.
都不是限制在某一些聖經.
而是整個舊約聖經裡面的一些教導.
都會指向.
耶穌生命裡面有活水出來.
信徒出來.
有聖靈的活水出來.
無論如何.
有一段聖經是跑不掉的.
當你開始想.
從耶穌的身體裡面有活水流出來.
有一段聖經你是跑不掉的.
那段聖經是出埃及記第十七章.
有一個盆石.
那個盆石流出活水.
你怎麼都跑不掉的.
我想和你讀這一段的聖經.
出埃及記第十七章.
我們將第一節讀到第七節.
你聽我讀 有些是沒有打出來的.
以色列人傳會眾在利菲汀安營.
但百姓沒有水喝.
沒有水喝.
小事而已.
百姓就與摩西爭鬧說.

$^{1081}$給我們水喝吧.
摩西對他們說.
你們為什麼與我爭鬧呢.
你們為什麼私探耶和華呢.
越來越生氣.
第三節.
百姓在那裡口渴要喝水.
就向摩西發怨言說.
你為什麼把我們從埃及領出來.
使我們和我們的兒女.
以及畜生全喝死呢.
摩西就呼求耶和華說.
我要怎樣對待這百姓.
他們差一點要用石頭打死我了.
第五節.
耶和華對摩西說.
你帶著以色列的幾個長老.
走在百姓的前面.
手裡拿著先前擊打尼羅河的像.
去吧.
你想想.
當時如果你是百姓.
你怎樣呢.
閃吧.
那個像一打下來.
全地都裂開.
第六節.
他說看啊.
我上帝要在河裂的盤石那裡.
站在你的面前.
你要擊打盤石.
水就會從盤石裡流出.
吸百姓渴.
摩西就在以色列的長老眼前.
這樣行了.
最後一節.
他給那地方起名叫做馬薩.
又叫做米利巴.
因為以色列人在那裡爭鬧.
並試探耶和華說.

$^{1121}$耶和華是否在我們中間.
我鼓勵你.
下次讀這段聖經的時候.
不要想想你自己人生裡面有些什麼.
上帝是否和我同在.
然後埋怨上帝.
你一這樣說.
你就死了.
你就落地獄了.
究竟發生了什麼事.
這本來是一件很小的事情.
但最終變成了一宗血案.
我解釋給你聽.
如果你熟《釵及記》的時候.
你知道前面幾章發生了什麼事.
我數給你看.
第十二第十三章.
逾越節.
上帝最後咬手瓜.
咬贏那一下.
第十災.
頭生的.
法老頭生都死了.
接著第十三章尾.
耶和華帶他們出埃及.
有雲柱.
早上有雲柱.
晚上是火柱.
你大哥.
幾百萬人.
你想想多壯觀.
你心裡面應該.
沒有信心都變有信心.
第十四章第十五章.
過紅海.
二百萬人過紅海.
摩西之歌.
還會唱歌.
過了紅海之後發生什麼事.
有苦水.

$^{1161}$去到一個地方.
有水.
不過很苦.
於是怎麼樣.
不要緊.
百姓向摩西怨言.
上帝搞不搞定.
搞定.
馬上苦水變甜.
好不好.
第十五章.
第十六章.
百姓又向摩西和阿倫怨言.
沒東西吃.
OK.
沒問題.
嗎哪.
OK.
鵪鶉.
要吃肉也行.
要吃麵包也行.
我們經過了第十二到十六章.
所有的問題.
上帝No Problem.
即日解決.
OK.
你現在去到第十七章.
沒水而已.
和上帝說一聲WhatsApp.
他馬上秒回.
為什麼要搞到這麼大件事.
究竟發生什麼事.
發生的事情.
如果你真的去讀.
屬靈的意義.
就是他們的埃及偶像的毒引起.
未過紅海還算好.
一過了紅海.
回不了頭.
死了.

$^{1201}$一過紅海開始埋怨.
為什麼.
我們的偶像怎麼辦.
我們以前.
埃及我們也拜的.
我們現在帶不出來.
耶和華不准帶偶像.
我現在沒有偶像.
沒有粉在身.
我現在渾身不自在.
於是他們要想辦法.
怎樣可以回去埃及.
於是連續三章.
他向摩西.
向亞倫.
向上帝發怨言.
接著來到這裡.
第二節.
他說百姓就與摩西爭鬧.
原來爭鬧在希伯來文裡面.
是一個法律的用詞.
代表我們要開庭控告上帝.
控告上帝是什麼.
不負責任.
違背婚約.
上帝根本就沒有能力.
醫保我們的子女.
醫保我們.
又和我們結婚.
又將我們從我們以前的娘家拿走了.
壞上帝.
我現在要和他離婚.
你明白嗎.
所以他在玩怨言.
去到這個地方.
你開始明白.
其實當日上帝.
本就可以將他們就地正法.
現在是誰壞人.
你想離婚就離.

$^{1241}$看你怎麼死.
我現在就這樣打死你也可以.
但是上帝選擇了做一件事.
上帝站在你的面前.
我只是告訴你.
站在你的面前.
是一個法律上的用詞.
就是他要成為那個被告.
他要成為那個受刑者.
接著上帝站在他們的面前.
要替他們受刑.
上帝站在盤石那裡.
上帝成為了盤石.
接著他說你來擊打吧.
擊打這個盤石吧.
上帝替他們受刑.
上帝心裡面知道的東西.
不是我死就是你亡.
如果上帝不做這個動作.
上帝不是說.
我站在你面前.
我站在盤石之上.
你打盤石吧.
那天死的就是以色列人.
200萬人全死.
全該死.
在那天上帝取締了他們子民的地位.
領受了他們該受的刑罰.
以致他們可以活下來.
整個故事是說什麼呢.
是一個福音的寓言.
是一個十字架的影子.
因為耶穌成為那個盤石.
故事就這樣看.
盤石流出活水.
其實故事背後還有另一層意思.
盤石打出寶血.
是耶穌被他們擊打.
以致有血他們才可以保命.
他們不但得到血的拯救.

$^{1281}$還得到活水的拯救.
沒有流血就沒有赦罪.
從活水的盤石.
變成活的寶血的盤石.
那個盤石就是耶穌.
我想帶你去看新約的版本.
講這麼多前書第十章第一節.
保羅的說話.
保羅說弟兄們我不願意你們不知道.
我們的祖先從前也是喝一樣的靈水.
所喝的是出於跟隨著他們的靈盤石.
那個盤石照猶太人的傳統.
自從那天開始就一直跟著他們.
有輪子的.
自己走的.
不用搬的.
去到哪裡就有活水.
不用插電的.
喜歡出水就出水.
那個盤石就是基督.
好正啊.
但請你不要誤會.
不要以為有了耶穌做盤石代罪.
於是叛逆上帝的以色列人就可以得到廉價的恩典.
好像耶穌替他們贖罪那天.
沒事了 以後有什麼就敲那個石頭算了.
不是啊 第五節.
但他們中間多半是上帝不喜歡的人.
所以他們倒閉在曠野裡.
這些事都是我們的尷尬.
使我們不要貪戀惡事.
像他們貪戀的一樣.
你說多半.
其實是多多多多多多多多多多半都死的.
你想想你數一下那些不用死的人.
我數來數去都是約書 加納 摩西差不多.
你有沒有數到其他人是多半不用死的呢.
上帝說上帝不喜歡他們.
God was not pleased with them.
於是他們倒閉在曠野那裡.

$^{1321}$我再加多一句.
死後被送往永恆的地獄裡面.
所以我說為什麼你讀這些故事的時候.
你不要將自己讀進去.
讀進去你就和他一起死.
因為聖經告訴我們.
這班人是上帝所不喜歡的.
不喜悅的.
於是他們死在曠野.
於是他們不能進入安息.
於是他們永遠會上帝上帝的審判.
到了這裡可能有人會問 等一下 林牧師.
上帝不喜歡.
說真的 除了耶穌之外.
上帝是誰都不喜歡的.
我們系統神學是這麼教的.
於是我憑恩典得救.
我恩信耶穌 稱義對不對.
舊約的異色異人都是恩信稱義.
我們今天都是恩信稱義.
對不對 不對.
不對在哪裡.
不對是因為舊約的人根本沒有信.
起碼這班人.
99.9999\%的人都不信.
他們不是真的.
他們是以假亂真的.
舊約的異色異人.
不停地用嘴巴讚美耶和華.
心根本不是跟隨耶和華的.
是假的.
所以上帝不是不喜歡他們.
所以你不要以為全世界都不喜歡.
上帝哪裡會喜歡 除了耶穌.
於是我躲在耶穌後面.
我就可以上天堂.
你不要想得這麼簡單.
這班人就是想得這麼簡單.
以為他可以躲在上帝的聖殿裡面.
見下祭上帝 母他父.

$^{1361}$一定要接他上天堂.
沒有啊.
他們根本就沒有信.
從來沒有信過.
於是我們一直都知道.
我們需要講的真正的信.
在聖經裡面講真正的信.
有三個神學上.
有三個角度去明白.
第一Notitia.
第二Assensus.
第三Fiducia.
Notitia是知識明白上面.
我相信以色列人有沒有的.
有沒有Assensus.
有沒有認同.
有沒有承認自己的罪.
一點點.
我猜一半半.
但是有沒有Fiducia.
有沒有信靠毀身.
沒有.
死最後那件事.
於是講這麼多.
前書第十章第六節這麼說.
他們成了我們的尷尬.
不要貪戀惡事.
不要貪戀埃及.
不要貪戀世界.
如果我們一手離敬拜上帝.
一手心裡面貪愛世界.
你知道我們成為了什麼.
成為了以色列人.
你不要以為有耶穌替你死.
有活石.
有活水的盤石.
除非你有真正的迴轉悔改.
真正的信靠毀身.
有真正的信心.
你是不會得救的.

$^{1401}$這個是保羅給我們.
非常沉重的警告.
我想帶你回到.
《約翰福音》第七章的場景.
你知道當時有很多人.
圍著耶穌.
耶穌講那句話的時候.
很多人聽.
議論紛紛的人群當中.
其實橫看點看.
只有兩組人是比較有盼望.
會真信耶穌.
其他都是假的.
我帶你去看兩組人.
第一組人是警衛.
第七章第四十五節.
警衛們回到祭司長和.
伯利賽人那裡.
他們對警衛說.
你們為什麼沒有帶耶穌來呢.
警衛回答說.
我們從來沒有像他這樣說話的.
第一組是有望的人.
是聖殿的警衛.
他們違抗了.
祭司長和伯利賽人的命令.
原因是什麼呢.
因為我們去到聽到耶穌教導.
聽完五分鐘就聽十分鐘.
十分鐘十五分鐘.
聽完那個鐘.
從來沒有人像他這樣說話.
於是就沒辦法下手.
聖殿的警衛.
你要知道這群人是長住在聖殿裡面.
聖殿是什麼人出沒的地方.
是那些名人拉比.
很出名的老師.
拉比在哪裡說話的.
在聖殿.

$^{1441}$尤其是大節旗.
所以這群警衛其實長年累月.
什麼名人拉比沒聽過.
什麼法律專家沒聽過.
什麼神學教授沒聽過.
全部都聽過.
接著他們的總結是.
No one spoke like him.
有沒有留意他們不是說.
沒有人做他做的神蹟.
他不是說神蹟.
他說教導.
是不是有前途呢.
如果他說沒有人做他做的神蹟.
那我看不起他們.
但是這群警衛.
有希望有前途.
為什麼.
因為他說No one spoke like him.
沒有人像他這樣.
甚至去到一個地步.
他們違抗命令.
當然更有盼望的.
我們也知道最後.
最後信的是尼哥德姆.
第五十節.
七章五十節.
其中有尼哥德姆.
才從前去見過耶穌的.
對他們說.
不先聽本人的口供.
查明他所做的事.
難道我們的律法.
還定他的罪嗎.
儘管尼哥德姆.
第二組有盼望的人.
其實你要知道.
他當時是位高權重.
他是以色列工會的一份子.
身份越高.

$^{1481}$風力越大.
越高風力越大.
但是他竟然願意為耶穌說一句公道話.
他還沒相信.
他真的還沒相信.
但是他已經開始清醒了.
他開始清醒了.
他知道他們戴上有色眼鏡.
他說我們可不可以公道一點.
儘管他以前見過耶穌.
我們記得.
第三章裡面.
他被耶穌狠狠的勸戒過.
耶穌憤了他很厲害.
但是他竟然可以謙卑受教.
他反覆思想耶穌的說話.
當他聽到警衛說.
No one spoke like him.
他心裡說.
我知道.
我那一晚跟他說完之後.
我心裡也是這樣.
於是他繼續跟隨耶穌.
最後在耶穌的十字架面前.
他俯伏敬拜.
但是我們去到整個楊福音.
最高潮的高潮.
耶穌要成全了.
要成全所有的東西了.
子彬拿槍擲他的肋旁.
立刻有血和水.
從他的肋旁.
從他的身體上面流出來.
我想你一起跟我去思想.
這個場景.
駐彭哲的水禮和酒祭.
記不記得.
有水有酒.
倒在聖殿的祭壇上.
又從聖殿的祭壇上流出來.

$^{1521}$耶穌是誰?.
上帝的羔羊.
耶穌是誰?.
非人所做的永遠聖殿.
又有水又有酒.
又有他的寶血.
我要你想想以西結書第47章.
有末世的聖殿裡面.
有活水的江河.
流向東邊.
從哪裡?.
從祭壇的旁邊流出來.
帶著血.
復活.
醫治這個世界.
使這個世界更生.
我想你想想撒迦利亞書第14章8-9節.
活水的聖靈.
從耶路撒冷的聖殿流出.
以致活水流到所有的地方.
信主的人會變成怎樣?.
耶和華要做全地的主.
全地要充滿著敬拜耶和華的人.
不是該人一般隨意獻奢東西的敬拜者.
是阿伯一般獻賞最好.
獻賞所有的敬拜者.
你知道的.
尼哥德姆來了.
下一個出來就是尼哥德姆.
尼哥德姆是哪一種?.
是阿伯還是該人?.
當然是阿伯.
你明白嗎?.
你開始了.
耶穌說.
是子信他的人.
要領受聖靈說的.
那時還沒有賜下聖靈.
因為耶穌還沒有得到榮耀.
到了此時此刻.

$^{1561}$到了第19章的時候.
那句話可以倒過來寫.
因為耶穌已經開始得榮耀.
此時此刻聖靈可以開始降臨.
因為人子耶穌得著榮耀.
怎樣得著榮耀?.
他高掛在十字架上.
在十字架上耶穌大聲呼喊.
呼喊什麼?.
先不要成了.
我渴了.
耶穌說我渴了.
耶穌說凡渴的可以來.
為什麼?.
因為我為你承擔了.
你不能承擔的渴.
耶穌會替我們承受地獄的火.
所以你來你不需要再渴.
這個渴不只是說.
我心裡面真的很不滿足.
人生真的很空虛.
寂寞的.
我想要一些慰藉.
神經病.
太小了.
如果耶穌只是要給一些慰藉來.
他根本不需要說我渴了.
他說我渴了是因為.
他要在十字架上為我們承擔地獄的火.
於是說完我渴了之後.
他說成了.
耶穌得到榮耀.
於是活水來了.
可以開始了.
水喉開一點.
還沒開到大.
要去聖靈降臨日才開大.
但已經開始小小了.
活水來了.
聖靈來了.

$^{1601}$生命來了.
不再飢渴.
不再死亡.
醉得赦免.
復活重生.
痊癒起來.
走進安息.
所有東西在十字架裡面.
通了.
耶穌和尼哥德姆說.
摩西在曠野怎樣舉蛇.
人子也要照樣被舉起來.
使一切順他的.
到時你看著我被舉起來.
你被我舉起來.
你見到我的時候.
你信我.
你就可以得著永生.
接著耶穌和他說.
信我的人.
福中將流出活水的江河來.
這是尼哥德姆當時的經驗.
肚子咕嚕咕嚕.
不是肚子痛.
是聖靈到了.
所以你可以想像.
當時的尼哥德姆.
肚子開始咕嚕咕嚕.
沒辦法了.
急了一定要去.
第十九章第三十八節.
若失來.
求彼拉多.
要把耶穌的身體領去.
尼哥德姆也來了.
帶著約一百斤的活藥和沉香.
他們照猶太人的規矩.
用細麻布加上香料.
把耶穌的身體裹好.
我們說過很多次這個故事.

$^{1641}$大家都有印象.
當日.
聖靈來了.
重生了.
兩個天國的嬰兒.
一嬰兒剛出生.
厲害.
他們有什麼特徵.
他們離開黑暗.
追隨真光.
從前是在暗中.
現在告訴他.
叛國罪我都和你死過.
於是他們出來見光.
跟隨耶穌.
勇敢的跟隨耶穌.
第二.
他們奉上所有.
他們帶來的是一百斤的活藥.
七十五磅的活藥.
七十五磅等於.
七十五年的人工.
普通人的人工.
很搭水.
奉上所有.
敬拜君王.
最後他們拋棄世界.
單求天國.
因為如果照我們所知道.
他們用手親自埋葬耶穌.
他們第二天無法參加逾越節的敬拜.
他們的老闆.
他們的同事.
整個猶太人的公會.
會知道為什麼尼哥德姆今天不出來.
因為他骯髒.
因為他埋葬耶穌.
這個人前途盡毀.
按照教會的傳統.
不是聖經裡面.

$^{1681}$聖經裡面找不到他的下場.
但是按照初期教會的傳統.
尼哥德姆.
信了耶穌.
他是彼得和約翰為他受洗.
受洗之後就被猶太的領袖.
逐出猶太的公會和議會.
後來還被人逐出.
聖食耶路撒冷.
耶路撒冷都不可以.
在那裡住.
也有另外一個傳統說.
他因為信耶穌.
失去了所有的財富和地位.
但是他得著了耶穌.
得著了耶穌的國.
耶穌的義.
聖靈不單止臨到尼哥德姆身上.
聖靈也來到耶穌.
他自己的十二門徒.
和其他跟隨他的人.
我一定要帶你看完這一段.
我們才可以結尾.
第二十章第十九節.
七日的第一日.
耶穌來站在當中對他們說.
願你們平安.
門徒說了這話.
他就把手和立旁給他們看.
門徒一看見.
耶穌就喜樂.
耶穌說.
父怎樣差遣了我.
我也現在怎樣差遣你們.
說了這話.
就向他們吹了一口氣.
說領受聖靈吧.
耶穌將他的手和他立旁給你們看.
是十字架的記號.
因為唯有耶穌的十字架.

$^{1721}$才可以為我們帶來平安.
願你們平安.
這句話之前是不能播出的.
是沒有平安的.
有十字架有平安.
有十字架有生命.
有十字架有喜樂.
有十字架有使命.
他們不單止得到平安.
喜樂.
生命.
他們得著使命.
福音的使命.
人生的使命.
最大的使命.
父怎樣差遣我.
我們可以想像父怎樣差遣耶穌.
我現在也.
照樣同樣地差遣你們.
接著說了這話.
耶穌向他們吹了一口氣.
說領受聖靈.
聖靈是活水.
聖靈現在又是靈風.
上帝的吹氣落在門徒的身上.
為什麼要領受聖靈.
聖靈對我們有什麼好處.
聖靈對我們有什麼用處.
答案在前一節.
父怎樣差遣了我.
我也照樣差遣你.
聖靈是用來做什麼的.
聖靈來使我們可以領受上帝的差遣.
這個是約翰福音裡面的大使命.
The Great Commission in the Gospel of John.
於是我想你們這樣想.
聖父聖子差遣他的聖靈來到我們中間.
不是要幫助你和我.
成就我們在地上所要的意願和夢想.
Do you get that.

$^{1761}$很多時候我們說聖靈阿方.
我很害怕我明天要考試.
我明天要見工.
我找不到工作幫助我.
主阿幫助我.
聖靈來啊.
No.
聖父聖子差遣聖靈來到我們中間.
不是要幫助我們成就我們地上所要的意願和夢想.
而是要終極改造我們的生命.
使我們成為耶穌基督的真兄弟.
真姊妹.
真兄弟.
耶穌的兄弟是這樣的.
現在聖靈來.
使他成為真兄弟真姊妹.
以致我們可以和耶穌一起起飛.
奔跑天路逆風而行.
以上帝的真光在黑暗的世界裡面.
作天國的見證.
我想用一個電視節目.
兩套電影的名字.
去幫我們總結活水凝風在我們身上的工作.
全部都是美國的東西.
第一個電視節目叫做.
Extreme Makeover.
終極的改造.
聖靈來要終極改造我們.
Extreme Makeover.
And I mean extreme.
真的終極.
不是,改一點點吧.
其實你已經是很好的人了.
不好意思,你現在信了耶穌.
希望你每個禮拜抽兩個小時來教會.
這樣就可以了.
你已經是很好的人了.
No, I'm sorry.
終極的改造.
完全的改造.

$^{1801}$你想想耶穌的兄弟.
耶穌的兄弟.
猶大.
雅各.
終極的改造.
你想想耶穌的門徒.
約翰.
雅各.
你想想彼得.
全部被耶穌終極的改造.
第二.
Galaxy Quest.
這個是飛上天的電影.
這個是喜劇來的.
電影來的,我很喜歡看這套戲.
是講飛天的.
Never give up, never surrender.
永不放棄.
永不投降.
最後.
Into Darkness.
真光見證.
我們說耶穌被門徒.
在他黑夜裡面.
在海裡行走.
為什麼?因為他要猜他們進入黑暗.
為他做見證.
\newpage



\section{路加福音 16:19-31}
\label{sec:032156}
\textbf{錢財無效時}~~~~俞大鳴牧師
\newline
\newline
連結: \href{http://www.church.com.hk/acms/content.asp?site=cdc\&op=show\&type=product\&code=032156\&layout=sermon}{\texttt{032156}} ~~~~ 語音日期: 2026-05-17
\newline
\newline
\hyperref[sec:032155]{\small{< < < PREV SERMON < < <}}
~
\hyperref[sec:index_chronic]{\small{[返順時目]}}
~
\hyperref[sec:032157]{\small{> > > NEXT SERMON > > >}}
\newline
\newline
路加福音 16:19-31
\newline
\begin{longtable}{cl}
\hline
\hline
章節 & 經文 (和合本修訂版)\\
\hline
16:19 & \begin{tabularx}{0.7\textwidth}{X} 「有一個財主穿著紫色袍和細麻布衣服，天天奢華宴樂。 \end{tabularx} \\ \\ \relax
16:20 & \begin{tabularx}{0.7\textwidth}{X} 又有一個討飯的，名叫拉撒路，渾身長瘡，被人放在財主門口， \end{tabularx} \\ \\ \relax
16:21 & \begin{tabularx}{0.7\textwidth}{X} 想得財主桌子上掉下來的碎食充飢，甚至還有狗來舔他的瘡。 \end{tabularx} \\ \\ \relax
16:22 & \begin{tabularx}{0.7\textwidth}{X} 後來那討飯的死了，被天使帶去放在亞伯拉罕的懷裡。財主也死了，並且埋葬了。 \end{tabularx} \\ \\ \relax
16:23 & \begin{tabularx}{0.7\textwidth}{X} 他在陰間受苦，舉目遠遠地望見亞伯拉罕，又望見拉撒路在他懷裡， \end{tabularx} \\ \\ \relax
16:24 & \begin{tabularx}{0.7\textwidth}{X} 他就喊著說：『我祖亞伯拉罕哪，可憐我吧！請打發拉撒路來，用指頭尖蘸點水，涼涼我的舌頭，因為我在這火焰裡，極其痛苦。』 \end{tabularx} \\ \\ \relax
16:25 & \begin{tabularx}{0.7\textwidth}{X} 亞伯拉罕說：『孩子啊，你該回想你生前享過福，拉撒路也同樣受過苦，如今他在這裡得安慰，你卻受痛苦。 \end{tabularx} \\ \\ \relax
16:26 & \begin{tabularx}{0.7\textwidth}{X} 除此之外，在你們和我們之間，有深淵隔開，以致人要從這邊過到你們那邊是不可能的；要從那邊過到這邊也是不可能的。』 \end{tabularx} \\ \\ \relax
16:27 & \begin{tabularx}{0.7\textwidth}{X} 財主說：『我祖啊，既然這樣，求你打發拉撒路到我父家去， \end{tabularx} \\ \\ \relax
16:28 & \begin{tabularx}{0.7\textwidth}{X} 因為我還有五個兄弟，他可以警告他們，免得他們也來到這痛苦的地方。』 \end{tabularx} \\ \\ \relax
16:29 & \begin{tabularx}{0.7\textwidth}{X} 亞伯拉罕說：『他們有摩西和先知的話可以聽從。』 \end{tabularx} \\ \\ \relax
16:30 & \begin{tabularx}{0.7\textwidth}{X} 他說：『不！我祖亞伯拉罕哪，假如有一個人從死人中到他們那裡去，他們一定會悔改。』 \end{tabularx} \\ \\ \relax
16:31 & \begin{tabularx}{0.7\textwidth}{X} 亞伯拉罕對他說：『如果他們不聽從摩西和先知的話，就是有人從死人中復活，他們也不會信服的。』」 \end{tabularx} \\ \\
[1ex]
\hline
\hline
\end{longtable}
$^{1}$願至聖母的平安.
我們先一起祈禱.
賜給我們主耶穌基督的清晨.
讓我們可以一同聚集.
敬拜你.
也聆聽你的話.
求聖靈保守帶領我們.
讓我們清楚明白.
你指華語的訊息.
幫助我們能夠在生活中.
實踐出主你的教訓.
奉主耶穌基督的名祈禱.
前一個星期,不是上一個星期.
我們看了《路加福音》.
十六章一至十三節.
那個是一個機警管家的比喻.
我不知道大家有什麼感受.
也不知道大家會不會開始學習.
做錢的主人.
做上帝的僕人.
還是把錢看得比天大呢.
今天的經文.
是從十六章十九節開始的.
心水清的弟兄姊妹可能留意到了.
我們跳了幾節聖經.
十四節到十八節.
我們跳了去.
但是從十四節.
我們知道.
其實當主耶穌講給我們的時候.
那個場景是有法利賽人在場的.
並且聖經形容他們.
愛錢.
愛錢的法利賽人.
他們對耶穌講的.
機警管家比喻的反應是什麼呢?.
他們是「譏笑耶穌」.
為什麼他們會「譏笑耶穌」呢?.
很好笑嗎?.
機警管家比喻有什麼好笑呢?.

$^{41}$其中一個可能性是.
他們聽到一個不義的管家.
因為他行他人職貴.
用了老闆的錢.
為自己謀福利.
他們只是聽到不義.
聽到管家不義.
他們沒有聽到一個機警的管家.
他們可能認為.
我們應該遵守上帝.
我們應該遵守上帝法律的人.
學習做人的道理.
不可能向不義的人學習.
所以他們就「譏笑耶穌」.
另一方面.
他們也可能認為.
耶穌比喻的那位管家.
對那些債人可以說是過分慷慨.
他們熟悉摩西的律法.
他們知道律法上是有教導人.
我們應該顧念窮人.
要奉獻給上帝.
尤其是十分之一的奉獻.
他們法律上的人會嚴格遵守.
但是他們不能夠認同那位管家這麼慷慨.
再加上他們只慷慨人家的財物.
而不是自己的財物.
因此他們完全不能夠.
接受耶穌所說的比喻.
「譏笑耶穌」.
竟然會用這樣的比喻.
來教導人.
教導他們門徒學習.
主耶穌怎樣回應呢?.
主耶穌來到世界上.
是要拯救罪人.
因此主耶穌沒有放棄這些法理錯人.
在第十五章其實已經說了三個比喻.
失而復得的比喻.
其實是責備和教導法理錯人.

$^{81}$教導法理錯人.
你們應該反思自己的價值觀.
在十六章十五至十八節.
他講了一些教訓.
指出法理錯人的錯誤.
提醒他們應該有天國的價值觀.
然後再為他們講今日經文的比喻.
教導他們錢財的魅力.
始終有一天會過期無效.
我們一起思想主耶穌的教導和比喻.
在十九節.
今日經文十六章.
十九節至二十一節.
就介紹了他們的比喻中有兩個人物.
亦都介紹他們在世上的時候.
是怎樣過日子的.
其中一位主角.
講明是一位財主.
很有錢.
但是沒有說他叫什麼名字.
我們不知道財主叫什麼名字.
只知道他是一個有錢人.
其實如果故事當中所有人的名字都有記載.
唯獨不記載財主的名字.
就表明這個故事的人.
認為這個財主是不值得被紀念的.
故事中有另外一個人.
有名的.
叫拉薩路.
如果這個比喻是純屬虛構的話.
主耶穌給黑衣拉薩路的名字是很有意思的.
拉薩路的名字的意思是.
上帝是我的幫助.
即是說他是一個倚靠上帝的人.
他的名字是被記錄下來的.
財主是一位不倚靠上帝的人.
他的名字亦都不會被紀念.
這兩個人在世生活的方式.
亦都形成一個非常強烈的對比.
財主是怎樣?.

$^{121}$財主穿著紫色的細麻布衣服.
當時來說紫色是皇家的顏色.
有錢人才可以穿.
好像今天很多人一樣.
喜歡穿名牌的衣著.
名師設計的衣服.
從前我行路比較好的時候.
我太太都帶我去行公司.
去到那些LV手袋店.
真的令我大開眼界.
一個手袋要賣幾千元都很貴.
其實現代人和這位財主.
心理都是一樣.
既然有錢.
我應該穿得好一點.
讓別人知道我有錢才行.
是一種虛榮心的驅使.
希望得到別人羨慕的眼光看著我.
為自己爭面子.
我留意到有些衣著.
以前都很流行做一個蝸頭.
我記得有一個牌子的衣服.
蝸頭是一個人在打馬球.
不知道大家認不認識這個牌子.
這個牌子的衣服通常都很貴.
但以前蝸頭是釘在這裡.
不記得左邊還是右邊.
總之小小一個.
你留心看就看到.
但現在新的不是.
釘得大大個.
怕你看不到.
你老花都看到了.
釘得這麼大個.
讓人家容易看到他穿的那件.
是名牌衣服.
貴價衣服.
車胎一樣.
以前看到的奔馳汽車.
只是車頭有個圓的三個X.

$^{161}$其實都很型的.
凸了出來.
但現在不用了.
而在車頭的鬼面罩上.
弄了大大的三個X.
好像怕人家看不到.
此外耶穌還形容這位財主.
天天做什麼呢?.
奢華燕樂.
如果用今天的生活方式來說.
可能在海邊有間大屋.
每天可以看無敵海景.
天天魚翅撈飯.
飲鮑魚雞湯.
出海有豪華遊艇.
想遠一點的地方有私人飛機.
這些就是財主的窮奢極侈的生活方式.
拉撒路又如何呢?.
這個以教上的人好像比較淒慘.
他的生活很淒慘.
他說他渾身長瘡.
一長瘡就想起藥伯.
早上在查藥伯.
藥伯不但渾身長瘡.
還很癢很麻煩.
在醫藥不發達的時代.
渾身長瘡可能令他不能工作.
所以二十節就說.
要被人放在財主門口.
如果他不是盲的.
就不應該要常常被人帶到財主門口.
他不是說被人抬到財主門口.
只是帶到財主門口.
可能他的眼睛也受到瘡的影響.
以致他需要被人帶到財主門口.
他去到財主門口.
希望得到財主的憐憫.
拿一些碎渣.
我們叫冷飯菜汁.
希望財主憐憫他.

$^{201}$給他一些冷飯菜汁.
使他能夠繼續活下去.
但經文沒有記財主憐憫他.
反而記載了有隻狗去同情他.
去舔他.
如果說這兩個人在世上的生活.
財主就好像活在天堂.
而拉薩路就好像活在地獄.
我相信不會太誇張.
不是每個依靠上帝的人都像拉薩路一樣活得那麼慘.
我相信如果每個信耶穌的人都像拉薩路一樣.
可能就沒有人肯做基督徒了.
好了,比喻就很快進入人生的另一個階段.
說到兩位主角都死了.
二十二節就對比他們死後的情況.
這一次先說拉薩路.
拉薩路是怎樣的?.
他死了就被天使帶去放在亞伯拉罕的懷裡.
天使好像招待員一樣.
把拉薩路帶到亞伯拉罕的懷裡.
懷裡其實可以說是一個重要的位置.
是主人和抬舉客人的位置.
如果我們記得最後晚餐的畫面.
使徒約翰就在耶穌旁邊.
耶穌的懷裡.
所以當耶穌說有人要出賣我的時候.
使徒約翰就打手勢暗示.
不是,是使徒彼得打手勢暗示.
問師父和耶穌究竟是誰門徒要出賣他.
這是懷裡的重要性.
亞伯拉罕在哪裡呢?.
亞伯拉罕是主耶穌之前兩千年的人物.
是一位與上帝有約的人.
是猶太人的先祖.
他的肉身早已死了.
但他的靈魂在這個比喻裡依然活著.
應該是一個相當舒適的地方.
並沒有形容他們很淒慘.
二十二節繼續說.
財主也死了.

$^{241}$真是太好了,上帝真是公平的.
有錢人也一樣會死.
有錢人也不會保證他不會死.
財主死了之後被人埋葬了.
他被人埋葬了.
卻沒有被天使帶走.
沒有天使出現.
他的靈魂去了哪裡呢?.
二十三節說.
我們知道財主在陰間受苦.
陰間就是所有與上帝無關的人.
死後要去的地方.
那裡是不是一個好地方呢?.
從二十三至二十六節給我們看到.
信上帝與不信上帝的人.
死後在永恆世界的對比.
財主是不依靠上帝的人.
他死後就去到陰間.
那裡是一個很痛苦的地方.
我想多說一件事.
陰間不是地獄.
地獄是在耶穌再來審判之後.
才會把跟隨魔鬼的人一起扔進地獄.
或者地獄是用燒著的硫磺火壺來形容地獄.
陰間可以說是一個等候審判的地方.
不過陰間其實仍然是非常痛苦.
並且在陰間的人感覺到自己的痛苦.
感覺不到的就不是痛苦.
他是感覺到自己的痛苦.
在二十三至二十五節的三節經文裡.
每一次提到陰間都是非常痛苦的地方.
我們發覺財主去到陰間.
不能夠再過窮奢極侷的生活.
他的財富已經可以說是過期無效.
沒什麼東西可以吃了.
在那裡他受到極大的痛苦.
他偶然舉目一看.
聖經的用法是「舉目」.
習慣了猶太人的說法.
舉目一看,他看到什麼呢?.

$^{281}$他看到阿布拉罕和拉薩路在另一個地方.
那裡應該很舒適.
他們沒有痛苦,不像他,他有痛苦.
以致財主要求阿布拉罕可憐他.
「快點,拉薩路過來,點點水,讓我不用這麼口渴」.
他生前可能喝慣了明酒.
現在很痛苦,一滴水也好.
可能已經很滿足了.
他生前也可能習慣了吩咐僕人幫他辦事.
於是他又習慣了叫阿布拉罕.
他不敢叫阿布拉罕過來.
「阿布拉罕,你派拉薩路來幫幫我吧」.
但是阿布拉罕用了兩個理由來拒絕他的要求.
第一,他說他生前享過福.
拉薩路也同樣受過苦.
第二,財主生前沒有倚靠上帝,沒有遵行上帝說的行善.
沒有用財富.
好像十六章九節說.
「要藉著那不義的錢財結交朋友」.
他生前不做這些事.
拉薩路生前每天都在你門口,你都沒有幫過他.
現在你又想拉薩路來幫你.
拉薩路也不會幫助財主.
拒絕財主的要求的第二個理由就是.
陰間和阿布拉罕的地方原來有深淵隔絕.
死後的靈魂都無法跨越的鴻溝去到另外一邊.
有時深淵可能令他想著上下分界.
以為天堂在上面,地獄在下面.
其實平衡分界都可以.
去年我有機會去新疆.
拍了一幅大峽谷的照片.
不止一幅,有很多幅.
去新疆有很多大峽谷.
這是其中一幅.
其實深淵可以是中間的結構.
因為地裂的緣故將兩邊分開了.
中間真的有個深淵,很深的,其實在旁邊的大峽谷.
財主在無可奈何的情況下.
他自己已經沒有救了,已經去了陰間.
很苦,就這樣苦了.

$^{321}$但他總算有些親情.
既然拉薩路又不可以來幫他解汗.
不可以給他水.
可不可以派拉薩路去警告他還有五個兄弟.
不要去到那麼痛苦的陰間.
但是阿巴安說他們應該聽的.
是摩西和先知們的話.
如果他們不相信摩西和先知的話.
即使拉薩路復活去告訴他們.
他們也是不會相信的.
事實上阿巴拉罕的判斷是正確的.
如果我們記得約翰福音第十一章.
耶穌其實是駛到另外一位拉薩路復生過的.
他就是馬大和瑪利亞的弟弟.
當時祭司見到因為拉薩路復活了.
很多人信了耶穌.
那些祭司會不會又信了耶穌呢?.
有一個活生生的重寫律復生的人坐在那裡.
那些祭司為甚麼不信耶穌呢?.
他們不單止不信耶穌.
他們想殺了拉薩路.
讓他不能做見證.
有些人是這樣的.
有復生的人坐在那裡.
但是他們都不肯相信耶穌的身份和能力.
所以就想殺了他.
耶穌講這個比喻.
其實是要為那些愛財的發利財人講的.
他們很喜歡錢.
雖然有很多發利財人不喜歡耶穌.
但是耶穌沒有放棄那些發利財人.
十五章他們自以為是.
就好像父親的大兒子一樣.
耶穌就講了三個比喻去提點他們.
耶穌知道發利財人愛財.
十六章就講了兩個比喻去提點他們.
錢財是會過期無效的.
有一天變廢紙.
對死了的人來說錢財就是廢紙.
當我們有錢財的時候.

$^{361}$我們應該學習的是機警的管教.
早作準備結交朋友.
以便那些朋友在錢財失效的時候會照顧我們.
從這個比喻中其實我們可以學到幾件事.
第一.
錢財的功用只有在今生財主生前可以用.
可以過奢華的生活.
但是死後他的錢是帶不走的.
也不能用錢換一滴水.
因此當我們活在世上的時候.
有機會的時候.
使用錢去做善事.
幫助有需要的人.
第二.
人是會死的.
所以每個人都同意.
但是死後不是一了百了.
如果死後是一了百了就算了.
不用說了.
我們盡量享受今生就算了.
但是死後不是一了百了.
從這個比喻中讓我看到.
人的靈魂是會去到不同的地方.
信上帝的人.
是會被天使帶到一個有福的地方.
有親人的地方.
就好像拉薩路.
去到先祖亞伯拉罕那裡.
不信上帝的人.
就去到陰間.
那裡是一個只有痛苦的地方.
我們要相信耶穌的講論.
從來沒有人去過天堂或陰間.
只有耶穌是從上帝的世界.
來到我們人間.
在上帝的道成了肉身.
來到這個世界.
祂所說有關另一個世界的事都是真實的.
假如你沒去過新疆.
我去過.

$^{401}$我告訴你.
新疆的情況你都會相信我.
在陰間的人.
是獨自受苦的.
拉薩路就有一個先祖.
財主有沒有?.
沒有.
一個人在那裡.
去到陰間是沒有人可以幫你的.
間中我不知道大家有沒有這個經驗.
在一些未信主人傳福音的時候.
有些人會這樣說.
我真的不去天堂.
我很多朋友都不在天堂.
都去陰間.
我不如去陰間跟他們開采打麻將也好.
我想強調的是.
在陰間.
每個人都是獨自受苦的.
朋友之間是不能夠交流的.
有一次我聽到另外一個人這樣說.
我很愛我老婆.
她未信耶穌.
已經死了.
她當然在陰間.
我也要去陪她.
我不能夠去天堂.
我不想去天堂.
這個真是很偉大的愛情.
我也很感動.
值得人尊敬.
不過我跟她說.
你愛太太當然很好.
如果你太太真的很愛你的話.
我相信她會像財主那樣.
她其實不會希望你去陰間陪她.
你去陰間也沒有用.
你去陰間只是受苦.
你幫不到她.
安慰不到她.

$^{441}$當我們仍然在生的時候.
假如我們所愛.
有我們所愛的人.
我們是不會想他去陰間受苦的.
我們就要早日跟他傳福音.
讓他們認識上帝.
讓他們接受耶穌.
做他們的主和救主.
讓他們可以有把握.
死後由天使帶他們去一個好地方.
最後一件事.
當我們活在人世間的時候.
有時我們都會覺得世事很不公平.
有權有錢的人.
就可以隨便欺壓人.
拉撒路是代表倚靠上帝的人.
但他在人間卻是受苦.
財主不倚靠上帝.
但他有錢.
不過他又不照顧拉撒路.
有時我們看看我們現今的世界.
在巴勒斯坦.
黎巴嫩的人被以色列人打到殘廢.
以色列人還要賣了加沙的土地.
讓巴勒斯坦人保證他們沒有土地.
一定沒有機會建立巴勒斯坦國.
以色列人覺得自己國土的北面.
有黎巴嫩的珍珠島.
令他們覺得沒有安全感.
於是就動用軍隊入侵黎巴嫩的南面.
說明要那邊的人快點走.
我們也不擔心以色列人.
我們只要你們40公里的地.
40公里以北你們走.
在美洲美國又以高科技和軍隊.
擄劫了委內瑞拉的總統.
去了美國受審.
然後又入侵伊朗.
這些情況都讓我們想起兩千多年前.
傳道者所說的話.

$^{481}$在傳道書四章一節說.
看下受欺壓的流淚.
無人安慰.
欺壓他們的人.
有權勢.
也無人安慰.
有時有些人覺得我們所傳的上帝.
要麼就是無能.
要麼就不是公義的上帝.
你有能力為何世上這麼多不公義的事你不去阻止.
你有公義的話為何你不去處理.
有些人因為這樣而不信上帝.
他們不一定是普通信徒.
我看過一些受過神學訓練甚至神學院的講師.
都會由於看到世界上的災難而從此不信上帝.
但主耶穌的比喻讓我們看到.
人死後有上帝的安排.
大反轉是會出現的.
就好像瑪利亞的尊主頌裡面所說.
他叫有權並的失為,叫卑賤的升高.
他叫飢餓的飽餐美食,叫富足的空手回去.
上帝的報應不一定是今生.
也可能是來生.
實際上來生的報應是更加嚴厲的.
更加有可能是永恆的.
所以上帝如果不在今生報應我們認為是惡的人.
其實是要他們在來世受到更重的刑罰.
因此如果我們現在有人覺得自己正在受委屈.
我們仍然是要仰望上帝.
保羅說不要置於甚遠.
寧可給主的憤怒留地步.
因為經上帝說主說申冤償我,我必報應.
不一定是今生.
但上帝一定會報應.
如果有人覺得那些戰爭罪犯.
攬炸醫院學校的兇手.
似乎仍然在逍遙法外.
要記得今生不報.
來生上帝一定會報.
並且報應是更加嚴厲,更加長久的.

$^{521}$大家一起祈禱.
你會公義審判萬民.
我們為那些在黎巴嫩,巴勒斯坦的難民祈禱.
求主席聯合國和其他組織.
供應食物和水給他們.
使他們略得安慰.
求主教導我們善用錢財.
給我們機會向親人朋友傳福音.
使他們有來生的保障.
奉耶穌基督的名祈禱.
阿們.
\newpage



\section{使徒行傳 4:1-22}
\label{sec:032157}
\textbf{聖靈本色}~~~~司徒贊傳道
\newline
\newline
連結: \href{http://www.church.com.hk/acms/content.asp?site=cdc\&op=show\&type=product\&code=032157\&layout=sermon}{\texttt{032157}} ~~~~ 語音日期: 2026-05-17
\newline
\newline
\hyperref[sec:032156]{\small{< < < PREV SERMON < < <}}
~
\hyperref[sec:index_chronic]{\small{[返順時目]}}
~
\hyperref[sec:032158]{\small{> > > NEXT SERMON > > >}}
\newline
\newline
使徒行傳 4:1-22
\newline
\begin{longtable}{cl}
\hline
\hline
章節 & 經文 (和合本修訂版)\\
\hline
4:1 & \begin{tabularx}{0.7\textwidth}{X} 彼得和約翰正向百姓說話的時候，祭司們、守殿官和撒都該人來了， \end{tabularx} \\ \\ \relax
4:2 & \begin{tabularx}{0.7\textwidth}{X} 就很煩惱，因為使徒們教導百姓，傳揚在耶穌的事上證明有死人復活， \end{tabularx} \\ \\ \relax
4:3 & \begin{tabularx}{0.7\textwidth}{X} 於是下手拿住他們；因為天已經晚了，就把他們押在拘留所到第二天。 \end{tabularx} \\ \\ \relax
4:4 & \begin{tabularx}{0.7\textwidth}{X} 但聽道的人有許多信了，男人的數目約有五千。 \end{tabularx} \\ \\ \relax
4:5 & \begin{tabularx}{0.7\textwidth}{X} 第二天，官長、長老和文士在耶路撒冷聚集， \end{tabularx} \\ \\ \relax
4:6 & \begin{tabularx}{0.7\textwidth}{X} 又有亞那大祭司、該亞法、約翰、亞歷山大，和大祭司的親族都在那裡。 \end{tabularx} \\ \\ \relax
4:7 & \begin{tabularx}{0.7\textwidth}{X} 他們叫使徒站在中間，問他們：「你們憑甚麼能力，奉誰的名做這事呢？」 \end{tabularx} \\ \\ \relax
4:8 & \begin{tabularx}{0.7\textwidth}{X} 那時，彼得被聖靈充滿，對他們說：「民間的官長和長老啊， \end{tabularx} \\ \\ \relax
4:9 & \begin{tabularx}{0.7\textwidth}{X} 倘若今日我們被查問是因為在殘障的人身上所行的善事，就是這人怎麼得了痊癒， \end{tabularx} \\ \\ \relax
4:10 & \begin{tabularx}{0.7\textwidth}{X} 那麼，你們大家和以色列全民都當知道，站在你們面前的這人得痊癒，是因你們所釘在十字架、神使他從死人中復活的拿撒勒人耶穌基督的名。 \end{tabularx} \\ \\ \relax
4:11 & \begin{tabularx}{0.7\textwidth}{X} 這位耶穌是：『你們匠人所丟棄的石頭，已成了房角的頭塊石頭。』 \end{tabularx} \\ \\ \relax
4:12 & \begin{tabularx}{0.7\textwidth}{X} 除他以外，別無拯救，因為在天下人間，沒有賜下別的名，我們可以靠著得救。」 \end{tabularx} \\ \\ \relax
4:13 & \begin{tabularx}{0.7\textwidth}{X} 他們見彼得、約翰的膽量，又看出他們原是沒有學問的平民，就很驚訝，認出他們曾是跟耶穌一起的； \end{tabularx} \\ \\ \relax
4:14 & \begin{tabularx}{0.7\textwidth}{X} 又看見那治好了的人和他們一同站著，就無話可駁。 \end{tabularx} \\ \\ \relax
4:15 & \begin{tabularx}{0.7\textwidth}{X} 於是他們吩咐他們兩人從議會退出，就彼此商議， \end{tabularx} \\ \\ \relax
4:16 & \begin{tabularx}{0.7\textwidth}{X} 說：「我們當怎樣辦這兩個人呢？因為他們誠然行了一件明顯的神蹟，凡住在耶路撒冷的人都知道，我們也不能否認。 \end{tabularx} \\ \\ \relax
4:17 & \begin{tabularx}{0.7\textwidth}{X} 但為避免這事越發在民間傳揚，我們必須威嚇他們，叫他們不可再奉這名對任何人講論。」 \end{tabularx} \\ \\ \relax
4:18 & \begin{tabularx}{0.7\textwidth}{X} 於是他們叫了兩人來，禁止他們，再不可奉耶穌的名講論或教導人。 \end{tabularx} \\ \\ \relax
4:19 & \begin{tabularx}{0.7\textwidth}{X} 彼得和約翰回答他們說：「聽從你們，不聽從神，在神面前合理不合理，你們自己判斷吧！ \end{tabularx} \\ \\ \relax
4:20 & \begin{tabularx}{0.7\textwidth}{X} 我們所看見所聽見的，我們不能不說。」 \end{tabularx} \\ \\ \relax
4:21 & \begin{tabularx}{0.7\textwidth}{X} 官長為百姓的緣故，想不出任何法子懲罰他們，只好威嚇一番就把他們釋放了；這是因眾人為了所行的奇事都歸榮耀與神。 \end{tabularx} \\ \\ \relax
4:22 & \begin{tabularx}{0.7\textwidth}{X} 原來經歷這神蹟醫好的人有四十多歲了。 \end{tabularx} \\ \\
[1ex]
\hline
\hline
\end{longtable}
$^{1}$旗鼎姐妹早晨.
今日我們講到的主題是.
《聖靈本色》.
我們今天是看《仕途行傳》.
教會給我的題目是講勇氣.
講勇氣的時候.
我想到一個經典的電影.
叫《英雄本色》.
大家有沒有看過.
《英雄本色》.
這齣電影好看.
不單止是因為子彈飛得好漂亮.
更吸引人的是劇情.
講到裡面每一個角色.
他們都不怕死.
好像每一個人都有九條命.
他們都不怕死.
為什麼他們不怕死呢.
原來他們每個人都有心目中的正義.
他們要守住.
好像張國榮.
他是一個警察.
維護這個社會的秩序.
就是他要守護的正義.
而狄龍和Mark哥.
雖然他們是黑社會.
但由於周潤發太有型.
還有他死都死不去.
所以我們的觀眾.
很容易都會支持.
很愛這個角色.
但是我們再回到《仕途行傳》的時候.
我們看到《仕途行傳》裡面.
都有很多英雄.
他們個個都不怕死.
數數有肥力.
有史提反.
還有巴拿巴.
《仕途行傳》講得最多的.
就是保羅和彼得.

$^{41}$講得最多的就是這兩個角色.
但你看到這些人不怕死.
他們是在捍衛同一個義.
他們是在捍衛神的義.
所以他們是不怕死的.
而我們看到彼得和保羅.
最後真的殉道.
他們的勇氣不是一個口.
不是一個講字.
是真的願意甘於去維護神的義.
而不怕失去他們的生命.
今天我們的經文主要是講彼得.
我們回顧一下彼得的生平.
彼得的生平是有很多英雄本色的瞬間.
彼得是吃過野粥的.
很有義氣的.
而主耶穌也幫他起了一個花名.
叫做雷子.
真的是很雷的.
主耶穌被捕那一晚.
他大膽到一刀手起刀落.
把大祭司的僕人耳朵切下來.
真的不是這麼容易的.
你讓我也手震了.
把別人的耳朵切下來.
還要在羅馬冰冰的面前.
把別人的耳朵切下來.
真的沒有膽量也做不到.
這個是彼得.
再加上我們看看彼得的生平.
很大膽的.
他勇字走先的.
胸口掛勇字.
雷電交加.
夜晚凌晨三點至六點.
最黑暗的時候.
門徒在大風大浪.
見到主耶穌在水面行走.
最大膽的是彼得.
說我不怕.

$^{81}$你們個個都怕.
我不怕的.
然後就走出去水面.
說主耶穌我要和你一起走.
然後見到風一吹.
就沉下去了.
然後他還很多話說.
拍胸口.
主耶穌說個個都會跌倒.
我是不會跌倒的.
拍胸口說話很有勇氣.
你們有沒有說過這句話.
弟兄姊妹.
這句話可能我們今天也有說.
我們身邊的弟兄姊妹.
怨若主耶穌.
我不會的.
我們有沒有說這句話.
都要反省一下.
如果我們有說的話.
還有說什麼.
就算我為你死.
我和你一起死.
我都甘心.
說的時候很勇.
但我們見到.
很有名的故事.
三次不認主.
所以我們見到.
人靠血氣裝出來.
扮出來的英勇.
原來在一些關鍵的時刻.
那個本色.
我們裡面那個本色.
怯懦的本色.
依然會反映出來.
今日經文.
繼續幫助我們.
去看一看彼得.
他由英雄本色.

$^{121}$到最後.
他站在一些.
很出名有權有勢的人面前.
都站立得穩.
正靈本色.
那種勇氣在心裡湧出來.
怎樣他可以做到呢.
在五旬節前.
五旬節後有一個很巨大的反轉.
一個很巨大的反轉.
這個就是他經歷了.
耶穌復活的能力.
彼得經歷了.
耶穌復活的能力.
完全變了另外一個人.
願意今日我們看完這段經文的時候.
我們都得著這份能力.
好不好.
在經文之前我們來個禱告.
在天父親.
多謝你派你的獨生子.
為我們上十架.
給我們有這個寶貴的救恩.
父下我們不願意今日只是得著這個救恩.
父下求你幫助我們得著聖靈的能力.
以至於我們在生命當中遇到很多的挫折難關.
我們都是不怕的.
我們都是願意靠著聖靈的能力.
去為你做美好的見證.
求你給我們有今日這樣的能力.
在經文中得力.
父下求你的靈與我們同在.
保守我們.
共主明求.
Amen.
我們一起來進入經文.
今日的經文其實是很豐富.
這裡這樣說.
使徒在傳揚耶穌的時候.
那些祭司,撒都哥人,大粒佬.

$^{161}$全部突然離開.
他們很討厭.
這個討厭不是生理上的討厭.
不是說不喜歡這群人.
這個討厭的原文是說.
他們很害怕.
即是說搶了他們的地盤.
他們覺得很害怕.
害怕什麼呢?.
就是害怕這群人搶了他們的信徒.
因為他們傳揚耶穌.
還要說耶穌復活這件事.
因為他們盯著耶穌上十架.
而他們又傳揚耶穌復活.
然後就抓了他們.
放了在拘留所.
其實這群人很恐怖.
因為他們有權力.
他們有地方的權力.
說抓就抓.
很恐怖.
抓了他們.
等到第二天.
就召他們來跟他們說話.
但聽到的人很多信了.
添到五千人.
而這五千人正正是他們覺得.
非常妨礙的原因.
很害怕.
很討厭.
到第二天.
他們就召集所有的權貴.
議員 長老.
全部聚集在那裡.
很大件事.
所有有頭有臉的人都聚集在一起.
又有大祭司.
蓋亞法 約拿丹 亞歷山大.
和所有屬於大祭司家族的人.
都坐在那裡.

$^{201}$他們叫使徒站在他們前面.
就查問.
這個查問不是普通的審問.
這個查問是帶著一些專業的術語.
就好像在法庭裡面.
他們派一個最大的大狀.
最厲害的.
最會說話.
最有口才.
最會法律.
最會律法的這個人去查問他們.
而他問這個問題.
都是很大的陷阱在裡面.
陷阱在裡面.
他就問這個問題.
就好像現在在法庭裡面說.
我現在問你.
你只是回答我什麼.
是 還是不是.
他問這個問題.
也是有這樣的陷阱在裡面.
你究竟是用什麼能力.
奉什麼人的名去作這件事.
一個很簡單的問題.
他就是想對方回答誰.
回答奉主耶穌的名.
如果對方這樣回答的時候.
就踩中這個陷阱.
因為根據《申命記》13-18章說.
如果你奉一個已經判罪的人.
因為耶穌已經被釘十架.
你奉耶穌的名去做這件事.
我就可以說你是假先知.
因為你奉這個人的名去做這件事.
但這個人已經被釘十架.
所以他問這個問題.
很有專業水準.
很有氣勢.
但是我們看看彼得怎樣回答.
彼得怎樣回答.

$^{241}$那時候彼得充滿聖靈.
他第一句說話是說.
民間的議員和長老.
我們看到彼得已經不是動刀動槍的彼得.
其實這裡說的民間的議員.
很有氣質很有禮貌.
彼得說Your Honor.
一個本身動刀動槍的彼得.
走到法官面前說Your Honor.
第一句說話Your Honor.
他很有條理地回答這個問題.
今天如果我們在這個病人身上所行的善事.
善事是一個關鍵詞.
這個善事彼得回答得很有水準.
這個善事捉住了猶太人律法的關鍵詞.
因為猶太人相信只有神才可以作善事.
彼得回答問題的時候.
不是說他奉誰的名.
因為他問他奉誰的名.
他回答的時候不是先說他奉誰的名.
他是先將善事這個事擺在前面.
如果你是因為這個善事而審問這件事.
這個人是因著什麼而得救.
那麼然後他才回答.
那麼你們眾人和以色列民眾.
就應該知道是因著你們所釘十架的.
上帝從中所揀選起來的.
活起來 復活起來 拿撒勒 耶穌的名.
所以彼得不是害怕去回答奉誰的名.
他也不是迴避這個問題.
而是他將善事這個字先擺在前面.
然後他一點都不帶害怕.
還要說是你們釘上十架的主耶穌.
還要是復活過來的.
就是因為這個人的名.
這位本身不能走路的人.
現在可以站在你們面前.
康健無病.
就是因為這個人.
醫治了這個本來不能走路的人.

$^{281}$然後他繼續說.
這個耶穌是你們這群匠人所藐視的.
他是不害怕的.
他說是你們藐視的耶穌.
現在成為了防國石的一塊.
然後他繼續說.
因著別人並無拯救.
天下也沒有別的名次在人間.
這句話弟兄姊妹很重要.
這句話根本是犯了他們的褻瀆罪.
他說天下間沒有別的名次在人間.
我們必須靠著而得救的.
這個場景很像耶穌被審問的夜晚.
如果大家有《馬可直擊》.
你就知道那群人是奈何不了耶穌.
如果耶穌不親自說自己是神的兒子.
坐在神的權能右邊.
他們是不可以釘耶穌的.
是耶穌親自將自己送上十架.
而彼得當晚三次不認主.
他在這裡說這句話.
彼得說這句話是很大的風險.
這正正是褻瀆罪.
他說天下沒有別的名.
他根本是說.
你現在要我的命,我也不怕.
彼得到了這個時候.
他已經去到一個絕對不怕.
你要我的命,你就拿去吧.
他真的有耶穌當晚的風範.
這個彼得完全是聖靈充滿.
所以這群人見到彼得,約翰坦然無懼.
見到他們坦然無懼.
又理會他們是沒有學問的外行人.
彼得真的是外行人.
但這一刻他說的話帶有權力.
還要像律師一樣.
他捉住善事這個字.
他們就奈他不何.
因為只有神才是善.

$^{321}$英文是good deed.
只有神才是善.
所以他們就稀奇.
又說這個人不是石頭爆出來的.
因為他之前曾經和耶穌生活在一起.
他真的是當時的彼得.
但為何今天突然變成像律師一樣.
又有口才又懂得說這些.
他們很稀奇.
接著我們來看看.
他們又見到.
其實這個時候彼得和約翰得了民心.
他們很有民心.
民意是站在他們那邊的.
所以這些議長就不能反駁他們.
其實我們今天看新聞也是.
你得民意.
如果你口頭上被他們捉到漏洞.
他依然可以定你的罪.
但當時彼得和約翰.
不單止民意站在他們那邊.
而那個理也是站在他們那邊.
所以這一刻那班人就選擇了.
先不去捉拿他們.
因為全城的民意都站在彼得那邊.
所以這個時候他們避免.
繼續避免傳揚耶穌這件事.
發佈在人間.
所以就要恐嚇他們.
他們失去了民意.
又失去了道理.
口頭上又說不過他們.
所以就恐嚇他們.
叫他們先不要再說耶穌.
彼得在這個時候肯不肯聽.
彼得當然不聽.
彼得就說這句很出名的.
我們所見的所聽的.
我們不能不說.
聽你們還是聽上帝呢.

$^{361}$在這個時候.
在上帝面前對不對.
你們自己判斷吧.
我們見到真是一個大的翻轉.
在彼得三次不認主的那一個晚上.
他面對的只不過是一些.
普通的市民和婦孺.
他們問彼得認不認耶穌.
彼得已經嚇到不敢認了.
但當他經歷了.
耶穌復活的大能之後.
今天站在公會.
站在一些有頭有臉的人面前.
他竟然可以自指亨昌.
句句擇地有聲地反駁.
這個是耶穌復活的大能.
這個不是彼得扮出來的.
我們再看看.
其實彼得的生命裡面.
有很多軟弱的.
他在登山的變象裡.
其實他講的語無倫次.
他不知道自己在講什麼.
一時又說為主死都可以.
但你現在見到他被聖靈充滿的時候.
他在公會面前.
他那種口才是超自然的.
還有我們見到在空墳墓.
他不明白什麼意思.
見到那個空墳墓.
他都不知道自己在信什麼.
但是當聖靈降臨的時候.
經歷過耶穌復活的大能之後.
他很清楚自己在信什麼.
天下沒有別的名.
他很清楚他很知道.
經歷過耶穌復活的大能之後.
我們再做一些對比.
彼得不只是這個時候.
那麼聖靈充滿.

$^{401}$他還指證了阿拿尼亞.
指證了他們的不是.
指證了他們說謊.
你知道在教會要指證人.
是需要勇氣的.
就算你明知道那個人做得不對.
你要指證他.
你是需要勇氣的.
很多時候在教會.
多一事不如少一事.
不講就不錯.
講多就講錯.
但是靠著聖靈的能力.
原來我們要學會.
不是為了自己的緣故.
不是為了自己的preference.
而是為了真理的緣故.
有時候我們是需要.
夠膽去講真話.
接著我們還見到.
彼得被人打.
他不但不還手.
他反而喜樂.
歡歡喜喜地離開.
被人打.
接著還要怎樣.
到最後生命盡頭.
等死的時候.
他竟然還可以睡著.
《史獨行傳》第十二章說.
到最後臨死那一刻.
他都睡得安然入睡.
好像主耶穌在船上.
風浪當中安然入睡.
所以我們見到.
聖靈的本色.
給我們帶來的力量.
是很值得我們去追隨的.
我們今天用三點.
去看看我們怎樣.

$^{441}$可以活出聖靈的本色.
第一點我們講到破碎.
破碎得剛強.
對焦就可以得本色.
不只是這樣.
我們還要操練.
操練得韌性.
我們逐點來看看.
破碎得剛強.
不知道大家認不認識.
鄧紫棋.
認不認識.
他很出名.
不是因為他唱歌好聽.
不只是因為他唱歌好聽.
是因為他那種生命力.
最近他在台灣開了演唱會.
對於一些唱歌好聽的人來說.
其實在台上講福音.
是很大的挑戰.
其實他是不需要的.
因為票都已經賣出去.
別人進來聽你唱歌.
但是他在裡面.
將一個很完整的.
分享他的生命.
信仰給他的生命的力量.
真的散發了出來.
雖然每個人都可能會繼續跌倒.
我不知道有一天他會不會跌倒.
但這一刻.
我看他的分享的內容.
和他那一種.
聖靈在他身上的能力.
是看得到的.
如果大家有機會.
去聽一下他的分享.
他其中一句說話.
就是我們要承認自己的破碎.
才可以接受別人的破碎.

$^{481}$然後.
我們才可以.
學會去愛.
他也有這樣說.
原來年輕時候的勇敢.
不是勇敢.
那些叫無知.
就好像彼得一樣.
他年輕時候的勇敢.
原來不是勇敢.
那些叫無知.
我們看到主耶穌.
其實也很幫助彼得.
去擁抱他的破碎.
因為主耶穌.
復活的時候,他再見彼得.
他很專門地.
帶彼得回到.
兩個場景.
早幾個星期.
Pastor Anders也說過一篇道.
就是說.
耶穌怎樣帶彼得.
回到這兩個場景.
第一個場景就是雨.
耶穌在他以前的生命裡.
怎樣供應他.
怎樣幫助他,經歷他.
第二個場景就是炭火.
炭火就是.
怎樣出賣耶穌的場景.
一見到炭火.
彼得就會想起.
自己怎樣出賣耶穌.
怎樣出賣耶穌.
但是我們看到.
我們的主,他幫助彼得.
去.
回去面對彼得的破碎.
不是要.

$^{521}$Put him to shame.
不是要彼得覺得自己很.
慚愧,耶穌要帶他.
回去面對他這些軟弱.
不是為了要.
奚落他,也不是為了.
要彼得跟他說.
對不起,因為.
在這個對話裡,耶穌沒有要求.
彼得跟他說對不起.
而耶穌要的也不是.
這件事.
耶穌三次去問彼得.
你愛不愛我.
很明顯就是要提醒彼得.
他三次背叛他的時刻.
但是.
耶穌不是.
要彼得說對不起.
因為每一次當.
彼得回答他的時候.
耶穌問他你愛不愛我.
彼得說,耶穌你知道我喜歡你.
其實當時.
阿加皮和弗雷奧這個字.
愛和喜歡是沒有分別的.
是同一個意思.
但是彼得偏偏.
不想用耶穌用的字.
去回答.
因為他心有愧.
但是.
當彼得.
每一次去回答耶穌的時候.
耶穌.
要彼得什麼.
你去餵養.
我的羊.
你去牧養我的羊.
你去餵養我的羊.

$^{561}$意思就是說.
耶穌告訴彼得.
我依然相信你.
耶穌每一次.
彼得回答他的時候.
耶穌跟他說.
我要支持你.
我要鼓勵你.
我要相信你.
我愛你.
這個是我們的主.
弟兄姊妹.
當神幫助我們.
面對我們破碎的時候.
他不是要我們.
說我們自己有多不堪.
耶穌要在當中.
要我們去面對.
我們的破碎,而在當中去醫治我們.
而我們的神.
也是一個公義的神.
他不會說以前的事不要說了.
這個也不是我們的神.
他不是說以前的事.
不要說了,抹掉它.
你現在就向前走,他也不是這樣說.
我說我要你記得.
你做錯的是哪裡.
但不要緊.
我要建立你.
我要affirm你.
我要信任你,我要醫治你.
這個是我們的主.
其實心理學上.
有很多.
心理學的始祖.
很多有名的.
榮格這些人.
他們也有說這些道理.
要我們去面對我們的情緒.

$^{601}$心理學家說.
佛洛伊德說.
未表達出來的情緒.
他們不會離開.
你要和自己的情緒對話.
你不去acknowledge.
這些情緒.
他們就會被活埋.
活埋的通常.
都是不好的情緒.
我們開心的時候.
表露出來沒有問題.
但我們不開心的時候.
把情緒壓下去.
我們的weakness.
破碎我們喜歡埋了它.
佛洛伊德說.
不要埋了它.
你硬埋了它.
這些情緒會以更猙獰的面目.
回來影響你的生命.
而榮格.
榮格比佛洛伊德小.
19歲.
是他的後輩.
本來佛洛伊德說.
把衣缽.
傳承給榮格.
榮格也有說.
潛意識.
他是其中一個.
比較早的人說.
潛意識是怎樣影響你的人生.
我們不要被潛意識.
控制我們的人生.
簡單來說.
我們要去認識我們的潛意識.
認識我們的破碎.
認識我們的weakness.
面對它.

$^{641}$很可惜的一件事.
就是榮格.
在他自己的寫作中.
提及過.
他曾幾何時激怒.
佛洛伊德暈倒.
激怒他暈倒.
因為他們兩個.
有一些分歧.
他們學術的理念上.
有一些分歧.
導致他們的關係破碎.
破裂.
佛洛伊德已經不認.
這個徒弟.
致死他們的關係.
沒有被修補.
我們看到.
我們的主耶穌.
不單止明白.
心理學.
心理學的大師.
我們也看到他們不可以靠心理學.
去修補或面對.
他們的負面情緒.
他們說就說得很厲害.
但他們面對不到.
他們有很多的苦毒在裡面.
他們的關係.
但我們的耶穌不單止明白.
知道我們要面對.
面對我們的破碎.
他還修補.
修補.
我們這個神主動的.
回來和我們修補這個關係.
是勝過一切的心理學.
勝過一切的心理學.
所以.
鄧紫棋也是這樣說.

$^{681}$因為他認識了我們的神.
他說這個世界是沒有完美的愛.
如果有完美.
就不需要去愛.
如果有完美就不需要去愛.
我們的神.
在我們還作罪人的時候.
已經是愛我們.
我們不完美.
但我們的神愛我們.
願意幫我們修補這個關係.
所以PJ.
常常在他的講章裡也有說.
我們要Be real with God.
Be real with God.
我們如果和神說話.
很虛偽的時候.
神是不懂得怎樣去回答我們.
所以今天我們.
和神再一次去聊天的時候.
我們真的要挖我們的破碎出來.
真的要求神.
求神幫助我們.
未受浸禮的弟兄姊妹.
真的.
我們不要以為我們靠自己.
修補好自己之後.
才將自己的生命交給神.
我們是在抹殺神.
幫助我們的能力.
我們要將自己還在破碎的時候.
將自己交給神.
讓神去修補我們.
其實我們每一個得救的人.
都是一個.
我們的身份已經很尊貴.
我們是上等的.
名茶.
我們很有味.
但是功夫茶.

$^{721}$怎樣沖功夫茶出來.
很多功夫.
少一步都不行.
少一步味道就出不來.
所以有時候我們基督徒.
得救的生命.
不等於被聖靈充滿的生命.
你們每一個得救的弟兄姊妹.
你們已經是很寶貴神的兒女.
但是怎樣才可以完全.
讓聖靈充滿我們.
讓我們更加.
發掘我們茶的本色.
我們聖靈的本色.
我們一起去探討.
好不好.
在《使徒行傳》這樣說.
怎樣才算對焦的本色呢.
《使徒行傳》這樣說.
我們一起讀好不好.
《使徒行傳》第一章十四節.
我們一起讀好不好.
一二三.
這些同著.
並耶穌的弟兄.
都同心合意地.
行切禱告.
《使徒行傳》裡面說到.
什麼時候門徒.
被聖靈充滿呢.
是他們一群人的行切禱告.
聖靈充滿.
不是一個人.
特別走出來發光.
聖靈充滿是說.
神的靈在一群人.
弟兄姊妹的身上.
一起閃耀出來.
在《使徒行傳》.
他們說.

$^{761}$當門徒同心合意禱告的時候.
神的靈就在他們當中.
當同心合意傳福音.
得救的人數.
天天加給他們.
榮耀神的時候.
神的同在降臨.
所以原來對照.
是要一群人去做的.
我感恩.
我們教睦團隊.
很成熟.
對我來說.
是一個很好的影響.
我們有不同的恩賜.
不同的崗位.
性別.
男人和女人的想法很不同.
年齡.
我們可以說.
99..
老實說.
99\%.
我們的看法.
都是一致的.
但當然有些位置.
有不同的看見.
不同的看見.
我感恩我們每一個人.
都有這樣的心智.
有不同的地方.
我們要.
放下禱告.
放下自己的喜好.
我們禱告.
然後再回來.
看看神的心意在哪裡.
當我們願意.
放下自己的喜好.
的時候.

$^{801}$我們見到神真的會帶領.
我們怎樣才能看到.
是不是神的旨意.
不是永遠說我的看法是對的.
因為神的聖靈.
降臨在我們眾人身上.
我們要同感一靈.
同感一靈的時候.
我們就看到神的心意.
彰顯在我們的教會.
我們很感恩會恩教會.
在建堂之後.
我們每個禮拜四.
風雨不改.
我們有祈禱會.
除了感恩節和聖誕節.
我們休息了.
我們每個禮拜四.
風雨不改.
我們有祈禱會.
當然不是所有弟兄姊妹.
都有默契.
但這個禱告的香爐.
是沒有停止過.
我們在過去一年.
每個禮拜四.
我們數算神的恩典.
我們經歷神的帶領.
我們看到神的同在.
這是聖靈降臨的.
一個表現.
所以被聖靈充滿.
不是一個個人的旅程.
是一個集體.
同行的恩典.
而聖殿神學.
告訴我們.
會幕的神學.
教曉我們一件事.
進入神的同在.

$^{841}$是一群人一起對焦.
原來.
進入神的同在.
不是一個偶然的感動.
神的感動.
不是偶然降臨在你身上.
原來.
進入神的同在.
是要經歷一個屬靈的操練.
是一個屬靈的操練.
我們看到地上的會幕.
不會過時.
因為希伯來書說.
這是天上敬拜的影儀.
所以這套路徑.
是不會過時.
今天我們進入神的同在.
如何預備自己進入神的同在.
都是要經歷.
這段路程.
賣一點廣告.
下星期第二課.
我們會有一個.
講敬拜的workshop.
在當中我們會.
細講.
這個會幕的每一個station.
但今天我有一個蓋覽.
一個蓋覽.
這裡有五個station.
進入神同在之前.
要經歷五個station.
第一個就是.
這個獻祭的爐.
每一個station.
它起碼有兩個意思.
雙重意義.
第一個意義就是耶穌一次的.
一次的奉獻.
獻上自己的身體.

$^{881}$所以我們今天不用在這裡殺豬殺牛.
我們不需要再獻動物的祭.
第二個meaning.
第二個意思就是說.
原來我們今天還要獻什麼祭.
有什麼祭.
羅馬書十二章.
活祭.
我們今天要繼續去獻.
這個活祭.
所以我們敬拜之前.
我們要為我們的奉獻去禱告.
這個是我們的奉獻.
不單止是金錢的奉獻.
是我們整個生命的奉獻.
所以這個是我們.
定期要go through的東西.
不是一次就搞定.
這個洗盤.
很多學者.
都會解釋為.
這個在預表浸禮.
在我們加入這個.
會幕.
會幕是我們教會的大家庭.
我們加入這個大家庭之前.
我們要先願意獻祭.
認罪獻祭.
將自己交給神.
透過洗禮.
我們去加入教會.
然後我們才可以.
去serve.
會幕裡面的其他的station.
第三個.
而這個浸禮.
第二個意思.
每定期我們要邀請聖靈.
去更新.
去清洗我們.

$^{921}$我們要繼續邀請聖靈.
在我們生命當中去潔淨我們.
所以我們手煮餐的時候.
我們就要做這些事.
第三個station.
就是這個陳設餅.
陳設餅的意思.
就是團契.
12個餅.
12疊餅.
代表以色列12個支派.
今天是代表.
和弟兄姊妹的團契.
加入教會.
我們要有團契.
團契不是一個bonus.
原來這條路徑.
是神教我們.
怎樣親近祂的方法.
親近祂是不可以skip團契.
團契之後.
才有金燈台.
金燈台不只是代表聖經.
金燈台是神的話語.
怎樣照亮我們的生命.
我們不只是讀聖經.
不只是認識神的話語.
是要讓神的話語.
改變我們的生命.
但你不讀神的話語.
又怎樣讓神的話語.
改變你的生命呢.
所以讀經是必然的.
一定的.
但不只是讀聖經.
是要讓神的話語改變我們的生命.
然後我們才可以同心合意.
去到金香爐.
一起同心合意.
我們的禱告便會直達到神的寶座面前.

$^{961}$我們便進入神的同在.
所以.
未受浸的弟兄姊妹.
我都願意鼓勵你.
這是一個開始.
浸禮是真的一個開始.
不是你告訴別人.
我已經很成熟.
我已經懂得給神的話語光照.
原來你在給神的話語光照之前.
你是要先願意.
將自己的生命交給神.
如果不是神的話語是光照不到你的.
有很多人聽完.
可能聽完說到.
touch不到我的heart.
為什麼呢.
原來我們skip了前面.
這幾個station.
我們要先願意.
將自己的罪交給神.
叫聖靈去清洗我們.
有團契.
我們才可以被神的話語光照.
但神的旨意.
永遠是要祂的子民.
as a group.
我們整個教會一起去活出祂的旨意.
不是一個個別.
聖靈充滿的行為.
時間關係.
希望大家可以參加.
下一個禮拜第二堂.
關於敬拜的.
workshop.
所以當日.
神的同在如何充滿.
今天聖靈也可以.
充滿我們的心.
但我們要有屬靈的操練.

$^{1001}$來到第三點.
原來.
聖靈充滿.
要不停的.
我們要不停的.
有屬靈的操練去keep up.
我們見到.
彼得早期.
聖靈充滿.
已經不怕死.
但原來.
時間一轉眼.
過了大概.
十二年至十七年之後.
彼得默會.
默了很久.
一開始在士徒人轉第十章的時候.
彼得已經知道神已經啟示他.
我說潔淨的.
你不要說不潔淨.
已經告訴彼得.
你去和外邦人吃飯.
福音要傳到給外邦人.
所以這十年期間.
彼得已經很習慣性的.
和外邦人吃飯.
他已經習以為常.
和外邦人吃飯.
但過了十多年之後.
我們在《迦太史》看到.
保羅要公開的.
去指責彼得.
因為發生了一件事.
因為.
突然間來了一群人.
一群身份尊貴的人.
在雅各那裡.
來了一群人.
身份尊貴.
看來像猶太人.

$^{1041}$在這個時候彼得竟然.
突然間.
害怕了那群人.
那群人來了之後.
經文說.
彼得突然間.
害怕了那群人.
那群已經受了割禮的人.
彼得因為害怕這群人.
他就決定.
暫時.
不和外邦人一起吃飯.
他突然間害怕了這群人.
很奇怪.
一個本身動刀槍的人.
聖靈充滿.
什麼都不怕.
你拿了我的命我也不怕.
為神作見證有口才.
到最後.
十多年之後.
突然間變回害怕這群人.
不知道他們害怕什麼.
所以我們看到.
聖靈就算充滿你.
這刻你為主很大膽.
去做很多事.
這件事是可以持續的.
所以弟兄姊妹.
我們要保持互相承諾.
我們要有教會的生活.
就算有時候.
牧者說得出牧者做不到.
這是事實.
說是很容易.
做是很難.
所以現在Pastor Andrews.
他有一個項目.
就是在說.
怎樣去面對邪教.

$^{1081}$怎樣去.
宣揚自己的信仰.
所以他有一個課程.
活因教會之門徒當自強.
我們要.
越多門徒越好.
越多人能夠辨別是非.
就越好.
如果當初我可能都會明白.
為什麼彼得突然間害怕了那群人.
不知道他們怎樣想.
有時候指證他們會不會令他們跌倒.
他有很多考量.
但是.
很寶貴的一件事.
就是當保羅指責他.
保持他有權的時候.
彼得是沒有反駁.
他們的關係.
也是沒有破碎.
後來彼得寫信的時候.
他都稱保羅是.
我親愛的弟兄保羅.
你看到彼得是很受教的.
但是.
如果教會當中.
越多人知道是非.
越多人知道真理.
就不需要.
傳道人.
更多的更正.
因為有時候更多的更正.
未必有用.
所以最好是大家.
當自強.
邪教.
很多時候人們以為.
異端邪教是從外面進來的.
可能突然間一個人走進來.
其實不是這樣的.

$^{1121}$因為現在我們經常在網上看到很多東西.
一些異端邪教的教導.
進入了我們的心.
我們再將這些東西帶回教會.
變成一個新的常態.
異端不是就這樣走進來.
在提摩太前後書中.
都是說.
其實異端的教導.
很多時候是在我們教會當中.
有人吸收了外面的教導.
再帶進來.
很多時候是這樣.
所以我們真的要當自強.
我們一起在神的話語裡.
下功夫.
不單止.
是對抗異端.
我們自己的馬轄得穩的時候.
我們更多的門徒的時候.
我們更多生命的改變的時候.
我們更多人.
可以活出聖靈本色的時候.
得救的人數.
神會天天加給我們.
我們不願意.
我們得救的生命.
有時候我們得救我們回教會.
但是大風大浪一吹過來的時候.
我們第一個反應就是.
害怕.
平安.
這個不是.
耶穌給我們的身份.
和見證.
我們一起的.
剛強.
我們一起有屬靈的操練.
我們每一個.
station.

$^{1161}$在會幕的station.
我們都好好地去經歷好好地去操練.
我們釋放這個剛強.
釋放這個活力.
還有一樣東西我不記得說.
那個陳設餅.
是每個禮拜換一次.
每個禮拜換一次.
我們不要輕看.
我們每個禮拜回到神的教會.
去敬拜.
這件事.
這個不是一個option.
神告訴我們.
祂喜歡我們這樣做.
祂喜悅我們這樣做.
我們按照神給我們的方法.
去親近神.
我們不要按照我們自己的方法去親近神.
我們一起來禱告.
早天父親.
我們帶著感恩的心.
來到你的面前.
求你幫助我們去.
面對我們的破碎.
幫助我們不要.
強行將這些東西.
收起來.
我們今天和你去對話的時候.
給我們有這樣的勇氣.
給你去.
進入我們的破碎.
給你去醫治.
幫助我們去親近你的時候.
按照你教導的方法去親近你.
讓我們有這個屬靈的操練.
讓我們在生命裡面.
得著聖靈的力量.
以至於我們.
走出來的那份膽色.

$^{1201}$不是我們自己裝扮出來.
而是你.
在我們的心裡面.
真真實實的.
賜予我們的力量.
求你幫助我們去釋放.
這個.
聖靈的能力.
給我們有生命的改變.
奉主耶穌得勝名字祈求.
\newpage



\section{列王記上 4:1-5:18}
\label{sec:032158}
\textbf{繁榮中的伏筆,國度中的危機}~~~~袁惠鈞牧師博士
\newline
\newline
連結: \href{http://www.church.com.hk/acms/content.asp?site=cdc\&op=show\&type=product\&code=032158\&layout=sermon}{\texttt{032158}} ~~~~ 語音日期: 2026-05-17
\newline
\newline
\hyperref[sec:032157]{\small{< < < PREV SERMON < < <}}
~
\hyperref[sec:index_chronic]{\small{[返順時目]}}
~
\hyperref[sec:032159]{\small{> > > NEXT SERMON > > >}}
\newline
\newline
列王記上 4:1-5:18
\newline
\begin{longtable}{cl}
\hline
\hline
章節 & 經文 (和合本修訂版)\\
\hline
4:1 & \begin{tabularx}{0.7\textwidth}{X} 所羅門作全以色列的王。 \end{tabularx} \\ \\ \relax
4:2 & \begin{tabularx}{0.7\textwidth}{X} 這些是他的官員：撒督的兒子亞撒利雅作祭司， \end{tabularx} \\ \\ \relax
4:3 & \begin{tabularx}{0.7\textwidth}{X} 示沙的兩個兒子以利何烈、亞希亞作書記，亞希律的兒子約沙法作史官， \end{tabularx} \\ \\ \relax
4:4 & \begin{tabularx}{0.7\textwidth}{X} 耶何耶大的兒子比拿雅作元帥，撒督和亞比亞他作祭司， \end{tabularx} \\ \\ \relax
4:5 & \begin{tabularx}{0.7\textwidth}{X} 拿單的兒子亞撒利雅作宰相，拿單的兒子撒布得作祭司和王的顧問， \end{tabularx} \\ \\ \relax
4:6 & \begin{tabularx}{0.7\textwidth}{X} 亞希煞作管家，亞比大的兒子亞多尼蘭掌管服勞役的工人。 \end{tabularx} \\ \\ \relax
4:7 & \begin{tabularx}{0.7\textwidth}{X} 所羅門在全以色列有十二個官員，供給王和王室的食物，每年各人供給一個月。 \end{tabularx} \\ \\ \relax
4:8 & \begin{tabularx}{0.7\textwidth}{X} 這些是他們的名字：在以法蓮山區有便‧戶珥； \end{tabularx} \\ \\ \relax
4:9 & \begin{tabularx}{0.7\textwidth}{X} 在瑪迦斯、沙賓、伯‧示麥、以倫‧伯‧哈南有便‧底甲； \end{tabularx} \\ \\ \relax
4:10 & \begin{tabularx}{0.7\textwidth}{X} 在亞魯泊有便‧希悉，他管理梭哥和希弗全地； \end{tabularx} \\ \\ \relax
4:11 & \begin{tabularx}{0.7\textwidth}{X} 在多珥山岡有便‧亞比拿達，他娶了所羅門的女兒她法為妻； \end{tabularx} \\ \\ \relax
4:12 & \begin{tabularx}{0.7\textwidth}{X} 在他納和米吉多，以及靠近撒拉他拿、耶斯列下邊的伯‧善全地，從伯‧善到亞伯‧米何拉直到約緬的另一邊有亞希律的兒子巴拿； \end{tabularx} \\ \\ \relax
4:13 & \begin{tabularx}{0.7\textwidth}{X} 在基列的拉末有便‧基別，他管理在基列的瑪拿西子孫睚珥的城鎮，巴珊的亞珥歌伯地的六十座大城，各有城牆和銅閂； \end{tabularx} \\ \\ \relax
4:14 & \begin{tabularx}{0.7\textwidth}{X} 在瑪哈念有易多的兒子亞希拿達； \end{tabularx} \\ \\ \relax
4:15 & \begin{tabularx}{0.7\textwidth}{X} 在拿弗他利有亞希瑪斯，他也娶了所羅門的一個女兒巴實抹為妻； \end{tabularx} \\ \\ \relax
4:16 & \begin{tabularx}{0.7\textwidth}{X} 在亞設和亞祿有戶篩的兒子巴拿； \end{tabularx} \\ \\ \relax
4:17 & \begin{tabularx}{0.7\textwidth}{X} 在以薩迦有帕路亞的兒子約沙法； \end{tabularx} \\ \\ \relax
4:18 & \begin{tabularx}{0.7\textwidth}{X} 在便雅憫有以拉的兒子示每； \end{tabularx} \\ \\ \relax
4:19 & \begin{tabularx}{0.7\textwidth}{X} 在基列地，就是亞摩利王西宏和巴珊王噩之地，有烏利的兒子基別，他一個官員管理這地 。 \end{tabularx} \\ \\ \relax
4:20 & \begin{tabularx}{0.7\textwidth}{X} 猶大人和以色列人如同海邊的沙那樣多，都吃喝快樂。 \end{tabularx} \\ \\ \relax
4:21 & \begin{tabularx}{0.7\textwidth}{X} 所羅門統治諸國，從大河到非利士地，直到埃及的邊界。所羅門在世的日子，這些國都向他進貢，服事他。 \end{tabularx} \\ \\ \relax
4:22 & \begin{tabularx}{0.7\textwidth}{X} 所羅門每日所用的食物：三十歌珥細麵，六十歌珥粗麵， \end{tabularx} \\ \\ \relax
4:23 & \begin{tabularx}{0.7\textwidth}{X} 十頭肥牛，二十頭草場的牛，一百隻羊，還有鹿、羚羊、麃子，以及肥禽。 \end{tabularx} \\ \\ \relax
4:24 & \begin{tabularx}{0.7\textwidth}{X} 所羅門管理整個大河西邊，從提弗薩直到迦薩，以及大河西邊的諸王，屬他的四境盡都平安。 \end{tabularx} \\ \\ \relax
4:25 & \begin{tabularx}{0.7\textwidth}{X} 所羅門在世的日子，從但到別是巴，猶大和以色列各人都在自己的葡萄樹下和無花果樹下安然居住。 \end{tabularx} \\ \\ \relax
4:26 & \begin{tabularx}{0.7\textwidth}{X} 所羅門擁有給戰車用的四萬個馬棚，還有一萬二千名騎兵。 \end{tabularx} \\ \\ \relax
4:27 & \begin{tabularx}{0.7\textwidth}{X} 這些官員各按自己的月份供給所羅門王，以及一切與他同席之人的食物，一無所缺。 \end{tabularx} \\ \\ \relax
4:28 & \begin{tabularx}{0.7\textwidth}{X} 他們各按其分，把給馬與快馬吃的大麥和乾草送到指定的地方去。 \end{tabularx} \\ \\ \relax
4:29 & \begin{tabularx}{0.7\textwidth}{X} 神賜給所羅門極大的智慧和聰明，以及寬闊的心，如同海邊的沙。 \end{tabularx} \\ \\ \relax
4:30 & \begin{tabularx}{0.7\textwidth}{X} 所羅門的智慧超過所有東方人的智慧，和埃及人一切的智慧。 \end{tabularx} \\ \\ \relax
4:31 & \begin{tabularx}{0.7\textwidth}{X} 他的智慧勝過萬人，勝過以斯拉人以探，以及瑪曷的兒子希幔、甲各、達大。他的名聲傳遍四圍的列國。 \end{tabularx} \\ \\ \relax
4:32 & \begin{tabularx}{0.7\textwidth}{X} 他作箴言三千句，詩歌一千零五首。 \end{tabularx} \\ \\ \relax
4:33 & \begin{tabularx}{0.7\textwidth}{X} 他講論草木，從黎巴嫩的香柏樹直到牆上長的牛膝草，又講論飛禽、走獸、爬行動物和魚類。 \end{tabularx} \\ \\ \relax
4:34 & \begin{tabularx}{0.7\textwidth}{X} 地上凡曾聽過他智慧的君王，都派人來；萬民都有人來聽所羅門的智慧。 \end{tabularx} \\ \\ \relax
5:1 & \begin{tabularx}{0.7\textwidth}{X} 推羅王希蘭是大衛平生的好友。希蘭聽見以色列人膏所羅門接續他父親作王，就派臣僕到他那裡。 \end{tabularx} \\ \\ \relax
5:2 & \begin{tabularx}{0.7\textwidth}{X} 所羅門也派人到希蘭那裡，說： \end{tabularx} \\ \\ \relax
5:3 & \begin{tabularx}{0.7\textwidth}{X} 「你知道我父親大衛因四圍的戰爭，不能為耶和華－他神的名建殿，直等到耶和華使仇敵都服在他腳下。 \end{tabularx} \\ \\ \relax
5:4 & \begin{tabularx}{0.7\textwidth}{X} 現在耶和華－我的神使我四圍太平，沒有仇敵，沒有災禍。 \end{tabularx} \\ \\ \relax
5:5 & \begin{tabularx}{0.7\textwidth}{X} 看哪，我吩咐要為耶和華－我神的名建殿，是照耶和華向我父親大衛說的：『我必使你兒子接續你，坐你的王位，他必為我的名建殿。』 \end{tabularx} \\ \\ \relax
5:6 & \begin{tabularx}{0.7\textwidth}{X} 現在，請吩咐人在黎巴嫩為我砍伐香柏木，我的僕人必幫助你的僕人。至於你僕人的工錢，我必照你所定的給你。你知道，在我們中間沒有人像西頓人那樣擅長砍伐樹木。」 \end{tabularx} \\ \\ \relax
5:7 & \begin{tabularx}{0.7\textwidth}{X} 希蘭聽見所羅門的話，就很高興，說：「今日耶和華是應當稱頌的，因為他賜給大衛一個有智慧的兒子，治理這眾多的百姓。」 \end{tabularx} \\ \\ \relax
5:8 & \begin{tabularx}{0.7\textwidth}{X} 希蘭送信給所羅門，說：「你派人向我所提的那事，我已聽見了；論到香柏木和松木，我必照你一切的心願去做。 \end{tabularx} \\ \\ \relax
5:9 & \begin{tabularx}{0.7\textwidth}{X} 我的僕人必把這木料從黎巴嫩運到海裡，我會把它們紮成筏子浮在海上，運到你告訴我的地方，在那裡拆開，你就可以收取；你也要照我的心願做，把食物給我的家。」 \end{tabularx} \\ \\ \relax
5:10 & \begin{tabularx}{0.7\textwidth}{X} 於是希蘭照所羅門的心願，給他香柏木和松木； \end{tabularx} \\ \\ \relax
5:11 & \begin{tabularx}{0.7\textwidth}{X} 所羅門給希蘭二萬歌珥麥子，二十歌珥搗成的油，作他家的食物。所羅門每年都是這樣給希蘭。 \end{tabularx} \\ \\ \relax
5:12 & \begin{tabularx}{0.7\textwidth}{X} 耶和華照著所應許的賜智慧給所羅門。希蘭與所羅門和平相處，二人彼此立約。 \end{tabularx} \\ \\ \relax
5:13 & \begin{tabularx}{0.7\textwidth}{X} 所羅門王從全以色列挑取服勞役的人，徵來的人有三萬， \end{tabularx} \\ \\ \relax
5:14 & \begin{tabularx}{0.7\textwidth}{X} 派他們輪流每月一萬人上黎巴嫩去；一個月在黎巴嫩，兩個月在家裡。亞多尼蘭管理他們。 \end{tabularx} \\ \\ \relax
5:15 & \begin{tabularx}{0.7\textwidth}{X} 所羅門有七萬扛抬的，八萬在山上鑿石頭的。 \end{tabularx} \\ \\ \relax
5:16 & \begin{tabularx}{0.7\textwidth}{X} 此外，所羅門有三千三百個監督工作的官長，監管百姓做工。 \end{tabularx} \\ \\ \relax
5:17 & \begin{tabularx}{0.7\textwidth}{X} 王下令，他們就鑿出又大又貴重的石頭來，用以立殿的根基。 \end{tabularx} \\ \\ \relax
5:18 & \begin{tabularx}{0.7\textwidth}{X} 所羅門的工匠和希蘭的工匠，以及迦巴勒人，把石頭鑿好，預備了木料和石頭來建殿。 \end{tabularx} \\ \\
[1ex]
\hline
\hline
\end{longtable}
$^{1}$列王記的第五講.
其實我應該講的是.
合一王朝系列的第五講.
因為我將列王記分開有五個系列.
第一個系列就是合一王朝的系列.
第二個系列就是北國王朝的系列.
第三就是南國王朝的系列.
然後就是以利亞傳系列.
以及以利沙傳系列.
今天要跟大家講的.
就是列王記上的第四章和第五章.
講到這裡又要講了.
我其實講這麼多的系列.
都是鼓勵大家來讀經.
如果你們每個星期聽我五篇的道.
以及讀相關的經文的話.
你應該兩年半就能讀完整部聖經.
而且應該對部聖經有很深入的認識.
因為我每一個系列的第一講.
都是講到一些的介紹.
有很多的資料.
有很多的背景.
幫助大家能夠明白這一部書.
如果你是這樣來讀經的話.
因為我是用字面來解釋聖經.
我是用聖經來解釋聖經.
你會發覺當你讀聖經的時候.
是非常之一致.
完全是沒有矛盾.
今天我給我這篇道的主題.
叫做繁榮中的伏筆.
國度中的危機.
我曾經跟大家講過.
列王記的作者文筆非常之精彩.
他在他的文筆當中給我們看到.
雖然所羅門的王朝.
外表看起來太平盛世.
但是原來內裡是危機四伏.
因為作者是有很多很多的伏筆.
以致我們知道最後帶領所羅門.

$^{41}$就是跟隨他的妃嬪去拜偶像.
使到以色列分為南北兩國.
我今天的大綱是非常非常之簡單.
一共是兩點.
第一點就是所羅門組織他的政府.
在列王記的第四章.
第二點就是所羅門準備起聖殿.
在列王記上的第五章.
大衛很有管理的恩賜.
他身邊也有很多的謀士.
例如迦德先知和拿丹先知.
都是他的謀士.
所羅門在聖經沒有說他有什麼謀士.
但是神就很公道.
神給所羅門有智慧有權柄.
但是他也需要在他身邊有很好的官員來輔助他.
所以在列王記上的四章一至六節.
講出了他官員的名單.
第一個出現的就是大祭司亞撒尼亞.
亞撒尼亞是撒督的孫.
在很多聖經版本裡面.
會告訴我們亞撒尼亞是撒督的兒子.
你說為什麼這麼奇怪.
原來希伯來文的「辯」這個字.
記不記得辯雅文.
右手之子.
辯這個字就是解釋兒子.
也可以解釋孫.
所以很多的翻譯都說亞撒尼亞是撒督的兒子.
其實他是撒督的孫.
又有奇怪的事.
在列王記上的四章四節.
又講到撒督和亞比亞他都是祭司長.
那究竟亞撒尼亞是大祭司還是撒督和亞比亞他.
其實當所羅門在位的時候.
撒督已經去世.
亞比亞他是以利的後人.
其實這兩個人都不簡單.
被我們看到是大祭司的家族.
聖經曾經預言在撒母耳上的第二章.

$^{81}$告訴我們.
撒督的家族將會代替以利的家族.
成為以利的大祭司.
現在在列王記上的四章四節.
居然將撒督和亞比亞他放在一起.
亞比亞他就是以利的後人.
有沒有看到這兩個大祭司的家庭.
放在一起.
就被我們看到.
這兩個大祭司都曾經為大衛服務.
所以在經盟約他們列出來.
還有第二個重要的人物.
就是彼拿亞.
彼拿亞曾經幫助所羅門排除異己.
殺死了藥鴨和侍媒.
大家記不記得.
他現在成為了所羅門的元帥.
他出生在祭司的家庭.
但他後來成為了一個軍人.
還有第三個人.
是宰相亞撒尼亞.
哇 有人拗頭了.
剛才說大祭司是亞撒尼亞.
為什麼現在宰相又是亞撒尼亞.
這兩個人是同一樣的名字.
其實是兩個人.
在聖經裡面告訴我們.
他的父親是拿丹.
我們知道在以色列.
很多人都叫做拿丹.
拿丹先知就是拿丹.
大衛也有一個兒子叫做拿丹.
究竟他的父親是誰.
聖經沒有告訴我們.
但這個人是所羅門派他去管理十二個地區.
是所羅門所分出來的十二個地區.
那你說這十二個地區要來做什麼.
這十二個地區.
不是照著十二支派來分.
當然有些地區就是根據十二支派.

$^{121}$例如阿徹 以撒加.
還有那佛他里.
但也有其他的地區.
例如多爾山干 阿魯柏.
和馬哈廉.
這些不是十二支派.
那為什麼要分出十二個地區.
一年有多少個月.
十二個月.
每個月這些地區都要負責皇室的食物一個月.
這個就是主要的原因.
如果你留意的話.
這十二個地區有沒有加上猶大支派.
只有猶大支派沒有加在裡面.
這個就出現了問題.
有很多其他的支派就不會開心了.
為皇家準備食物.
你猶大支派不需要出一份力嗎.
你猶大支派就是皇的支派.
但不只是你有皇的支派.
是皇的支派.
變亞敏支派也是.
第一位皇是誰.
是素羅王.
這個就出現了日後有很大的問題.
我不會把這十二個地區一一讀出來.
我會留在PowerPoint那裡.
如果大家有興趣可以下載來看.
還有第四個人物.
又是拿丹的兒子.
一個兒子做宰相.
另外一個做什麼.
是以色列的管家.
他就是阿希薩.
他也是所羅門的朋友.
還有第五個重要的人物.
就是亞多尼蘭.
我們講大衛的功績的時候.
大衛王朝都講過亞多尼蘭.
或者是亞多蘭.

$^{161}$他是管理奴役的.
奴工.
尤其是起聖殿的人.
等一下我們會見到.
他是有很悲慘的收場.
有很多的人.
很多學者都說.
所羅門的王朝.
那些官員的結構.
和埃及一模一樣.
但其實如果你注意的話.
所羅門王朝的結構.
和大衛王朝是一模一樣.
不一定要抄埃及人.
經文說.
無論是猶太人,以色列人.
都如同海邊的沙那麼多.
列王記上的四章二十節.
這個就是神給亞伯拉罕的應許.
說他的子孫好像天上的星.
好像地上的沙一樣那麼多.
他的後裔的應許只不過是部分.
神給亞伯拉罕的應許.
除了後裔還有.
是版圖還有地的應許.
在述列王記上的四章二十一節.
說這個版圖是會從大河去到斐理士地.
直到埃及的邊界.
在述詩篇的七十二篇.
也就是所羅門所作的詩歌裡面.
所羅門也有說過.
他要執掌權柄.
這個他是誰?.
是彌賽亞.
所以詩篇的七十二篇.
是彌賽亞國度的寓言.
他是要從這海到那海.
這個地是要從這海到那海.
從大河直到地極.
住在曠野的必在他面前下拜.

$^{201}$他的仇敵必要添土.
他斯和海島的王要進告.
士巴和西巴的王要獻禮物.
諸王都要叩拜.
萬國都要事奉他.
有沒有看到?.
萬國都要事奉他.
這個絕對是彌賽亞的寓言.
就好像肥納比書的二章告訴我們.
萬室都要敬拜.
萬口都要承認.
這個版圖有沒有實現到?.
由這海到那海.
我跟大家講過.
西邊的領域.
去到斐理士人的地方.
所羅門的時候.
沒有去到大海.
這個要什麼時候才實現?.
在千禧年國.
就是彌賽亞的國度才會實現.
如果地上的應許都不實現的話.
我們怎麼能夠知道.
將來天上的應許會被實現.
所羅門一共有七百個妻子.
三百個妃嬪.
還有很多官員和常客.
你光是想餵飽這些人都不簡單.
你說是不是?.
有常客來你都要照顧的.
去到列王記上的第十章.
有個士巴女王.
一帶就帶了幾百人來.
來參觀所羅門的國度.
你突然間要給幾百人吃東西.
好像我們五十週年.
幾百人去到.
不是那麼容易找到的.
有個地方給大家吃東西.
絕對不簡單.

$^{241}$在列王記上的四章.
二十二至二十三節.
說所羅門每天所用的食物.
細麵是三十哥兒.
粗麵是六十哥兒.
肥牛十隻.
草腸的牛二十隻.
羊一百隻.
還有鹵,羚羊,鵰子與肥禽.
這個是一天的食物.
你不要以為是一個月的食物.
或者一年的食物.
你說一個哥兒有多少?.
差不多六十加侖.
三十哥兒的細麵就是一千八百加侖.
一天你都別說.
真的是非常大的一個資源.
當然所羅門也有很多的附庸國.
是會付出這些代價.
但是主要就是由那十二個地區來負責.
這個也是原因.
為什麼所羅門死了之後.
十個支派就起來造反.
反抗羅波安.
就是所羅門的兒子.
所羅門不單只需要龐大的資源來支持他的王朝.
他也需要龐大的軍隊.
這些軍隊都是由地區的官員負責來供給.
這些戰馬是從哪裡來的?.
是從埃及來的.
列王記上的十章二十八節非常昂貴.
這個也是作者有一次讓我們見到所羅門的軟弱.
他是依靠兵馬多過依靠耶和華.
就是新命紀的第十七章第十六節.
此事王不可為自己多添馬匹.
也不可使百姓回埃及去.
為要家添他的馬匹.
因耶和華曾吩咐你們說不可回那條路去.
神救贖以色列人出埃及.
圍爐之家.

$^{281}$是要救贖他們.
救贖了你之後.
你又回到圍爐之家.
還要去買馬匹來依靠馬匹而不依靠神?.
所以神在新命紀的第十七章.
是警告君王千萬不要這樣做.
而在聶王紀上的四章二十五節.
說所羅門在世的日子.
從彈道別士巴的猶太人和以色列人.
都在自己的葡萄樹下和無花果樹下安然居住.
在所羅門的時候.
神還是祝福以色列人.
當他們守約的時候.
神也會守約讓他們在應許之地安居樂業.
但當他們背弱.
神甚至將他們趕出應許之地.
這就是聶王紀讓我們見到以色列人的歷史.
神不單止給所羅門有智慧.
還給他財富尊榮.
神還應許所羅門.
在他在世的日子.
聶王中沒有一個能比他.
所羅門的智慧是超越了所有當時出名的人.
包括東方和埃及的人.
所有的智者.
其實在以色列都有很多智者.
有以探,有希曼.
記不記得以探和希曼.
是大衛王朝.
大衛差派他們管理音樂的那兩個官員.
他們是非常有智慧的.
這個以探.
他另外有一個名字叫做耶道頓.
他就是作詩篇89篇那個.
希曼是作詩篇第88篇的那個.
另外在歷代志上的二章六節.
還列了有兩個智者.
就是甲郭和大拉.
無論是以探,希曼,甲郭,大拉.
沒有一個人是比所羅門更有智慧.

$^{321}$所羅門作了《箴言》一共有三千句.
他寫詩歌是寫了有一千零五首那麼多.
但是在我們手頭上已經沒有那麼多.
很多大部分三千句的詩歌.
都已經失傳了.
在《箴言》裡面我們只有所羅門所作的《箴言》.
大概六百個.
而詩篇在著我們詩篇的當中.
72篇和127篇都是所羅門的作品.
我們也都遺失了另外的文獻.
是所羅門記在列王記上的十一章四十一節有提到.
另外拿丹,阿希雅和易多也有寫了一部叫做《寓言書》.
是有關所羅門記載在歷代志下的九章二十九節.
也是已經失傳了.
我們知道所羅門在聖經裡面作了一些什麼書呢?.
就是《箴言傳道書》和《所羅門之歌》或者《雅歌》.
我們看到在這些書本裡面.
所羅門對自然科學有他的研究.
有他獨特的見解.
他的父親大衛是欣賞神的創造.
在他的詩歌裡面大衛不斷讚美神的創造.
但是所羅門是欣賞神有關大自然科學.
他對這些大自然科學有洞察力,有觀察力.
神給所羅門有智慧能夠理解到這些花草樹木.
飛禽走獸和昆蟲水族.
在《傳道書》的第二章五節.
說所羅門建造了庭院和公園.
他在這些花園裡面做什麼呢?.
欣賞大自然生物的成長.
能夠了解自然的科學.
所羅門的名譽傳到很遠.
很多的君王都會派使者來看所羅門的寶藏.
要欣賞所羅門的智慧.
聆聽他的智慧.
去到《列王記》上的第十章.
所羅門就是一個例子.
所羅門在位的期間.
外表看起來非常太平盛世.
當他登上了王位.
到他奉獻聖殿這段時間.

$^{361}$似乎他是與神同行.
他是尋求神的喜悅.
但是作者不斷在《列王記》裡面給我們有伏筆.
例如在《列王記》上的三章一節.
說所羅門在興建聖殿和王宮的時候.
居然將王宮放在聖殿的前面.
這個就讓我們看到.
作者可能認為所羅門的優先秩序出現了問題.
遲些去到《列王記》上的第六章和第七章.
作者說所羅門用了七年的時候興建聖殿.
用了多少年來興建王宮?.
十三年.
有沒有看到這個比例是二比一?.
將相配的時間放在自己的王宮.
而聖殿只有七年.
就讓我們看到問題所在.
而在《列王記》上的三章二節.
說所羅門和他的子民.
仍然在休壇上來敬拜神.
我們知道以色列人應該在哪裡敬拜神?.
應該在這裡回望.
回望在機片那麼遠.
但是大衛已經將約瑰搬回耶路撒冷.
建了一個帳篷.
如果以色列人要去敬拜神.
也可以去約瑰.
為什麼仍然要去休壇?.
休壇是俗世的人敬拜偶像的地方.
這裡又讓我們看到另外一個問題.
所羅門的確沒有他父親對神有專一的信心.
他最後被他的妻妾帶領他去拜偶像.
以致他離棄了神.
我們已經講完了第一點.
就是第四章裡所羅門組織政府.
第二點是所羅門準備建聖殿.
在詩篇的一百三十二篇三至五節.
大衛說.
「我必不准我的帳幕也不上我的床榻.
我不容我的眼睛睡覺.
也不容我的眼睛打頓.

$^{401}$直等我為耶和華尋得所在.
為亞國的大能者尋得居所」.
大衛真的很想為神來建聖殿.
在《說詩篇》的二十七篇第四節.
大衛又說.
「我有一件事我曾求耶和華.
我仍要尋求.
就是一生一世住在耶和華的殿中.
瞻仰祂的榮美.
在祂的殿裡求聞」.
大衛真的很有耶和華,很有神的心.
專心希望為神建聖殿.
但是神就說.
「你不要為我建聖殿了.
你滿手鮮血.
你不如繼續去東征西討.
擴張以色列人的領土.
你的兒子不同.
所羅門就可以為我來建聖殿」.
我們講大衛傳的時候也說過.
其實大衛已經為所羅門.
準備好一切建聖殿需要的東西.
從摩西起.
其實我剛才也說過.
以色列人也在回望,敬拜.
但是當時以色列人還在曠野漂流.
到處走.
所以他們敬拜神的地方.
回望是一個帳篷.
可以跟著以色列人到處走.
但是現在以色列人已經不再是遊牧民族.
而他們四周圍的人.
四周圍的外邦國家.
都有聖殿來敬拜他們的假神.
所以以色列也很需要有一座很宏偉的聖殿.
來敬拜他們萬有的真神.
萬君的耶和華.
在所羅門的第四年.
在位第四年.
就是公元前的九八六十六年的二月.

$^{441}$也就是我們四月至到五月間.
所羅門就開始了他的工程來建聖殿.
所羅門強調他能夠為神建聖殿.
不是他自己.
完全是神和大衛所納的約.
神應許在大衛的皇位上.
將會有一個後人.
會為耶和華來建聖殿.
這就是所羅門.
如果你看到在《列王記》裡面.
這是一個非常重要的一點.
就是大衛之約.
有人會說.
神會住在人手所造的殿裡面嗎?.
我記得在《述使徒行傳》的七章和十七章.
經文也有說.
「神創造宇宙和其中萬物的神.
既是天地的主.
就不住在人手所造的殿」.
為何神要建一座聖殿呢?.
神是否真的住在裡面呢?.
歷代志下的七章十六節又怎麼說呢?.
「神說現在我以選擇者殿.
分別為聖使我的名永在其中.
我的眼我的心也必常在那裡」.
你說大件事.
聖經是否有矛盾?.
是否有衝突?.
一會又說神不住在人手所造的殿.
一會又說神的名在那裡.
其實完全沒有矛盾.
神說這個聖殿很重要.
我會特別關注這個聖殿.
因為我的名聲也會在那裡.
但是聖殿能否限制住神呢?.
絕對不能夠.
神是無處不在的.
所以絕對沒有矛盾.
現在神的靈住在哪裡呢?.
在新約.

$^{481}$是住在我們每一個信徒的身上.
所以我們每一個信徒.
我們都要很小心我們的行為.
所羅門也清楚指出.
他興建聖殿不是為了紀念他父親.
不是他父親的紀念碑.
這一座的聖殿用來做什麼呢?.
是榮耀耶和華名的聖殿.
我說在《列王記》裡面.
是充滿了作者的伏筆.
其實在《列王記》上的五章到第十章.
似乎作者又有另外的伏筆.
今天的經文第五章.
就是說所羅門準備興建聖殿.
和推羅王希蘭來立約.
推羅王希蘭是在庶北邊的一個國王.
他就是推羅王.
推羅在以色列的北邊.
推羅王是男性.
另外有一個王是士巴女王.
在《列王記》上的第十章.
他是女性.
他是在以色列的南邊.
有沒有看到這兩個王.
一首一尾是一個首尾結構.
就是讓我們看到兩位王.
一個是男一個是女.
一個在北邊一個在南邊.
這個首尾結構的第二層.
都是有關聖殿.
在《列王記》上的第六章.
是所羅門建造聖殿.
在《列王記》上的七章.
十三節到九十九節.
是所羅門奉獻聖殿.
兩個都是有關聖殿.
一首一尾就是這個結構的第二層.
而經文最中心的地方.
作者指出.
在《列王記》上的第七章.

$^{521}$一至十二節.
所羅門建造自己的王宮.
有沒有看到.
所羅門的優先次序.
似乎作者在暗示.
真的出現了問題.
所羅門的確為自己的王宮.
是很有私心.
建造聖殿需要物力,人力.
所以首先所羅門需要準備木材.
我們講大衛傳的時候.
我們也講過.
大衛其實已經為了.
建造聖殿的人力,財力,物力.
作好準備.
整個建造聖殿的藍圖.
都已經給了所羅門.
而推羅王聽見所羅門.
聽見所羅門承繼了大衛.
就派使者來到.
在《列王記》上的五章一節.
說希蘭愛大衛.
希蘭王其實和大衛很有交情.
這兩個國家我們見到.
都有文約是很親密.
但是你想想現在發生什麼事.
推羅在哪裡?.
尼巴嫩.
和以色列打仗打得乒鈴噹噹.
但當年原來兩國是這麼友好.
有這麼密切的關係.
所羅門說他的父親.
不能夠為耶和華建造聖殿.
要等到耶和華令以色列的仇敵.
都伏在他的腳下.
這就是我剛才說.
為什麼神母叫大衛來為他建造聖殿.
在《列王記》上的五章四節.
所羅門說.
現在耶和華我的神.

$^{561}$使我四周平安.
四周平安.
沒有仇敵.
沒有災禍.
所以所羅門就為神來建造聖殿.
但我們知道原來大衛在位的時候.
都曾經有太平盛世.
在撒母耳夏的第七章第一節.
經文說.
耶和華使他使大衛安靜.
不被四周的仇敵擾亂.
就在這個時候大衛說.
我要為神建造一座聖殿.
亦在這個時候.
神跟大衛說.
你不要為我建造聖殿.
我為你建造你的家室.
就是跟他納了大衛之約.
因為神知道大衛的安靜不會持續.
那你說.
所羅門的太平盛世會不會持續?.
其實本來應該持續的.
如果他繼續信靠神.
我相信是會持續.
但是很可惜.
在列王記上的第二章.
我們看到他做什麼?.
在二章的十二節.
經文說得很清楚.
神已經堅固鞏固了所羅門的王位.
但是在第二章的尾.
他仍然靠自己的手段.
自己的智慧.
來擊倒,來攻擊他的敵人.
排除異己.
他是靠自己的能力.
很有趣.
在希伯來文的攻擊,擊倒是哪個字?.
paga這個字.
但是災難呢.

$^{601}$神本來令到他沒有災難.
災難這個字是pega.
是同一個字根.
本來神令到他是沒有災難的.
但是因為他靠自己的智慧.
自己的能力.
去攻擊他的敵人.
去擊敗他的敵人.
paga這個字.
現在神令到他有災難.
在列王記上的十一章.
十四節,二十三節,二十五節.
都說神興起敵人.
來攻打所羅門.
來攻擊所羅門.
於是所羅門就和希蘭合作.
來興建聖殿.
尼巴嫩就攻擊以色列人.
鄉柏木.
所羅門答應給推羅有食物.
推羅王希蘭很開心.
他讚頌耶和華.
因他次給大衛一個有智慧的兒子.
列王記上的五章七節.
相信他日後會後悔他說了這句話.
我們稍後會見到.
所羅門也希望得到一個巧手的工匠.
來幫他做聖殿的工具.
那些家具.
於是希蘭王就派了烏蘭這個人.
烏蘭的父親是推羅人.
他的母親是彈支派的人.
也就是說是猶太人.
他是一個非常有天才的工匠.
他是做了聖殿裡面的兩條銅柱.
如果我們記得.
聖殿裡面有兩條很出名的銅柱.
一條叫波亞斯.
記不記得波亞斯這個人.
另外一條叫做瓦根.

$^{641}$而烏蘭也做了在聖殿裡面.
很多金屬的家具.
在建造會幕的時候.
神幫助摩西.
能夠召集所有需要的工人來建會幕.
現在同樣神是幫助所羅門.
召集所有工人來建聖殿.
根據這個合約.
所羅門每年要給推羅.
要有墨紙,橄欖油.
日後還要給很多的誠邑給推羅.
還要給推羅人工資.
給工人的工資.
就是大麥,小麥,橄欖油和酒.
都給很多的.
那推羅給以色列人有什麼呢?.
香柏木和松木.
就是這麼多.
造成木佛.
造圓海在地中海.
由尼巴嫩一直去到約柏這個地方.
就是現在的瓦佛.
然後再經過三十五里運輸去建聖殿的地方.
來建聖殿.
作者再次指出.
所羅門是擁有從神而來的智慧.
和推羅人來做交易.
在外表看上來.
推羅人佔盡上風.
又有這麼多的食物.
工資都給了.
將來還要給誠邑給推羅人.
以色列人只拿到什麼?.
木材而已.
好像推羅人佔盡上風.
但遲些我們就會見到.
去到列王記上的第九章.
推羅王對所羅門給他的誠.
非常之不滿.
似乎這一單的交易.

$^{681}$有利於以色列人.
這個可能就是從神而來的智慧.
所羅門是很有外交的手腕.
在尼巴嫩曾經出土了一副石棺材.
讓我們見到古代外交的記載.
這個石棺材在1855年在西頓出土.
上面有斐理基的文字.
記載了所羅門王和推羅王.
一起合作建聖殿的事.
讓我們見到當時的推羅.
真的非常富有.
而當時的以色列也是非常繁榮.
兩個國家一起合作.
來建造一個在世上.
可以說是最重要的建築物之一.
希蘭和所羅門立約.
你說是一件好事還是一件壞事?.
在列王記上的五章十二節.
說這是一個文約.
如果我們記得在《申命記》的七章一至四節.
神曾經叫以色列人.
不要和當地的迦南人立約來通婚.
你說這條條例是禁止和迦南人立約通婚.
現在推羅不是迦南人.
但是遲些我們會見到在列王記上的十一章.
當作者譴責所羅門娶外邦妻子的時候.
就是用《申命記》的七章一至四節來罵他.
可見這一節的經文不單止是給迦南人.
日後阿哈王也會娶了西頓的公主耶洗別.
而被聖經譴責.
阿哈王遲些又會和亞蘭王變哈達來立約.
亦是為神所不滿不起悅.
所以我們見到.
其實《申命記》的七章一至四節.
不單止是給迦南人.
而作者在此有如此細緻入微的描述.
亦可能是另外一個伏筆.
我們講完了所羅門準備材料來建聖殿.
他亦準備了工人.
我們知道建聖殿需要很多的木材.

$^{721}$需要斬很多的木.
運輸很多的木.
需要大量的人力.
我們如果記得的話.
大衛在述撒姆以下的二十四章.
曾經數點人民.
而不為神所喜悅.
所以他那次數點人民沒有完成.
但未完成的人民的數點已經知道.
當時有130萬的藍丁.
而現在所羅門只用了3萬的藍丁來建聖殿.
所以不到總人口的2.3\%.
而經文告訴我們.
1萬人在尼巴嫩工作一個月.
和希臘的工人一齊來建聖殿.
然後兩個月在自己家.
你說不錯.
哪裡有這樣的工我都去做.
一個月去工作.
兩個月在家.
你現在先不要這麼開心去做.
你等一下聽完才作出決定.
至於以色列人的工資.
是否已經包括在推羅人的工資裡面.
這已經是無從稽考.
我們不知道.
有很多解經家都說.
所羅門是強迫以色列人來建聖殿.
因為希伯來文Mers這個字.
其實是用在虎翼的上方.
但我們要知道.
這個希伯來文.
也可以解釋做僕人.
效勞甚至是進貢.
所以不一定是虎翼的.
而在列王記上的九章二十至二十三節.
亦都說所羅門沒有令到以色列人成為奴僕.
你說為何作者要這麼遲才說.
不早一點說.
這裡可能又有另外一個伏筆.

$^{761}$雖然Mers這個字.
不一定是解釋做虎翼.
但是永遠或長期的服役.
這就是原因讓我們看到.
所羅門可能長期勞民傷財.
所以令到以色列人非常不滿.
叫苦連天.
這裡亦都讓我們看到.
就是撒母耳.
記不記得以色列人曾經要求撒母耳給他們一位王.
撒母耳怎樣說.
你們要一位王.
這位王將來就會勞動你們.
苦待你們.
會徵從稅.
這些事在所羅門的身上都發生了.
剛才說一個月在利巴嫩工作.
兩個月在家.
你說很享受有這樣的工作可以做.
其實經文.
一個月在利巴嫩就很清楚了.
起聖殿.
兩個月在家.
在你自己的家還是在所羅門的家.
這裡作者的意見都不合一.
如果一個月所羅門叫他們起聖殿.
兩個月為所羅門起他自己的家.
這樣就很合理了.
這樣就是勞民傷財了.
為何人民會這麼辛苦.
兩個月在家享受生活不會勞民傷財.
他就是要人們一個月起聖殿.
兩個月起皇宮.
而這個比例剛好和.
所羅門用了七年起聖殿.
十三年起皇宮.
有沒有看到這個也是二比一的比例.
一個月起聖殿.
兩個月起皇宮.
又是一樣.

$^{801}$所以有很多學者都認為.
所羅門可能起皇宮和聖殿是同時來舉行.
可能聖殿是早一點.
所羅門也和外邦人的人口做了調查.
因為起聖殿都需要很多外邦人.
調查到一共是有十五萬的外邦人.
要他們來做甚麼呢?.
切割石頭,挖石頭,運輸石頭.
在十五萬人當中.
七萬人是運輸石頭.
八萬人是挖石頭,切割石頭.
而且還有三千三百個督工.
在署列王記上的五章十六節.
還有二百五十個高級督工.
在署列王記下的八章十節.
為何需要這麼多督工呢?.
你想想.
切割了石頭.
每塊都要差不多大小.
然後去到皇宮合併.
就能夠天衣無縫.
不簡單的.
需要很多人力,物力,督工.
這個在將來我們見到.
聖殿被毀的時候.
因為本來聖殿的牆是塗了金的.
是有金造的.
而當聖殿一燒著的時候.
這些金都跑到縫隙裡.
而這些縫隙是細到.
那些金屬是可以跑進去的.
但你想找金屬出來就找不到.
以至將來的羅馬兵要怎樣呢?.
將所有聖殿裡面的磚頭.
全部打散.
去拿這些金出來.
所以我們見到當時.
做聖殿.
他們的工程是做得不錯的.
那你說這些外邦人在哪裡來的?.

$^{841}$以色列人進入迦南地.
神叫他們殺光所有迦南人.
但他們沒有殺.
有很多這些後裔.
就成為了以色列人的奴僕.
有人說聖經是種族歧視.
在這裡你有沒有見到.
給他們甚麼工作?.
建聖殿.
有機會來認識神.
其實這裡還有一個舊俗的圖畫.
就是以色列人代表甚麼?.
就是從埃及圍爐之加避神.
是舊俗的人.
在他們當中還有甚麼人?.
是未信的.
好像我們身邊也有很多未信的人.
他們仍然是罪的奴僕.
就好像這些外邦人一樣.
所以不單止神給他們機會來接觸自己.
這也是一個舊俗的圖畫.
我剛才說.
所羅門因為他很專制勞動人民.
徵重稅.
所以令到人民非常之反感.
而當他的兒子羅波安.
繼承了所羅門成為以色列人的王.
他又不肯放鬆在人民身上的重壓.
那又怎樣?.
南北國就分裂了.
這個羅波安的頭頭重點.
他就是差遣管府役的主管.
記不記得我剛才跟你們說過.
這個叫做亞多蘭.
或者亞多尼蘭.
這個就是所羅門出名的官員之一.
他派他去管那些以色列人.
結果那些以色列人就用石頭.
將亞多尼蘭活活地打死.
而這個羅波安就害怕得馬上上了車道.

$^{881}$回到耶路撒冷.
這些記載在列王記上的十二章.
我將來會跟大家說.
神給所羅門有很大的智慧.
但是很可惜.
所羅門很多時候.
他仍然是靠自己的智慧來行事.
一個例子就是我剛才所說.
在列王記上的二章十二節.
已經是說所羅門的國位已經是鞏固.
但是在第二章的尾.
他仍然用自己的手段.
用自己的方法去排除異己.
這個就讓我們看到.
他其實很多時候是沒有用到神的智慧.
今天接觸的時候.
想跟大家來看一個三點.
都是跟神的聖殿有關.
第一,我們看到非信徒的不信者的參與.
無論是在尼巴嫩工作的工人.
就是推羅國的工人.
又或者在以色列本土的外邦人.
都是不信主的.
是兩群的外邦人.
但是神就准許他們為自己來建聖殿.
為什麼呢?.
因為神的殿是萬國禱告的殿.
在以賽亞書的五十六章.
瑪提福音的二十一章.
和盧卡福音的十九章.
神是願萬人得救.
不願一人沉淪.
我們看到在這個例子裡面.
神就是容許這些外邦人來建聖殿.
有機會來認識自己.
很多時候我們在教會也是一樣.
很多時候我們都是很謹慎.
有很多工作我們不應該讓未信的人來做.
我也同意.
例如港道,領司.

$^{921}$有很多事情不適合未信的人做.
但是也有很多其他的事情.
可能可以適合,可能可以讓未信的人來做.
以致我們能夠讓他們有機會來經歷神,來接觸神.
這就是第一點.
第二點也是和聖殿有關.
今天說所羅門建聖殿.
我們知道這個聖殿.
是會被尼布格尼撒.
被巴比倫的軍隊毀滅.
毀了這個聖殿.
但是神也是透過.
甚至尼布格尼撒被神稱為他的僕人.
但是同樣神是會利用波斯王,塞魯士.
來釋放以色列人.
給他們有機會重建聖殿.
我們看到無論是尼布格尼撒或者是塞魯士.
都不是信神的人.
但是神仍然可以用他們來成就祂的旨意.
有時神甚至可以用一些我們根本不欣賞的人來管教我們.
甚至神也可以用那些不屬於祂的人.
來成就祂的旨意.
其實這就給我們有很大的信心.
因為我們相信的這位神是萬有的主宰.
一切都是在祂的手裡面.
祂是有主權的那位神.
所以就能夠幫助我們.
我們只要靠著我們的禱告.
繼續有信心地來事奉祂.
一切都掌握在神的手裡面.
還有第三點.
就是我今天的主題.
「繁榮中的伏筆,國度中的危機」.
在外表看上來.
所羅門的王朝.
太平盛世非常繁盛.
但是內裡危機四伏.
作者很多的伏筆都讓我們看到.
這就是屢外不一致.
在耶穌基督的時候.

$^{961}$祂也罵法利賽人.
是甚麼呢?.
粉色的墳墓.
外表看上來冠冕堂皇.
很漂亮.
內裡是裝死屍的.
冤奔濫囚.
我們呢?.
內裡裝甚麼?.
不是死屍.
我們是裝了聖靈在裡面.
我們的身體是聖靈的殿.
但我們仍然都需要裡表合一.
我們要活出我們裡面的聖靈.
在我們的行為那裡.
以致別人見到我們的好行為.
就會將榮耀歸給我們的神.
所以最後一點.
我們都需要裡表合一.
來榮耀神.
願神祝福他自己的還願.
我們低頭禱告.
讓我們有列王記相的第四章和第五章.
讓我們見到所羅門組織政府.
也讓我們見到.
祂準備聖殿.
從中也讓我們學習到.
我們是需要將我們的信心.
放在你的身上.
因為你就是那位主管一切的神.
無論聖殿的備位.
被建造或者再被建造.
一切都掌握在你的手裡.
甚至你會用那些不屬於你的人.
來成就你的旨意.
這實在是給我們很大的鼓勵.
願你讓我們繼續禱告為你的事公.
繼續出力.
只需要做好我們的本份.
其他我們都不需要擔心.

$^{1001}$因為你就是那位主管萬有的真神.
但我們還要學到.
我們需要做一個屢表合一的人.
因為我們裡面是有聖靈住在我們裡面.
願意我們都能夠活出.
聖靈要我們做的事.
讓路給聖靈.
而不是我們自己的旨意.
願意我們都能夠.
是靠我們的行為表現出.
我們是一個真正.
靠我們的行為表現出.
我們是一個真正得救的人.
將那榮耀頌讚歸給你.
我們這樣禱告.
奉主耶穌基督的名義求.
阿們.
\newpage



\section{約書亞記 18:3}
\label{sec:032159}
\textbf{聖經人聖經事:你們耽延不去得}~~~~黃克勤牧師博士
\newline
\newline
連結: \href{http://www.church.com.hk/acms/content.asp?site=cdc\&op=show\&type=product\&code=032159\&layout=sermon}{\texttt{032159}} ~~~~ 語音日期: 2026-05-22
\newline
\newline
\hyperref[sec:032158]{\small{< < < PREV SERMON < < <}}
~
\hyperref[sec:index_chronic]{\small{[返順時目]}}
~
\hyperref[sec:032160]{\small{> > > NEXT SERMON > > >}}
\newline
\newline
約書亞記 18:3
\newline
\begin{longtable}{cl}
\hline
\hline
章節 & 經文 (和合本修訂版)\\
\hline
18:3 & \begin{tabularx}{0.7\textwidth}{X} 約書亞對以色列人說：「耶和華－你們列祖的神所賜給你們的地，你們耽延不去得，要到幾時呢？ \end{tabularx} \\ \\
[1ex]
\hline
\hline
\end{longtable}
$^{1}$各位弟兄姊妹,平安!.
希望你身體健康,靈性進步.
今天又一齊讀聖經.
今天讀的聖經是.
「你們擔言不許得」.
看看聖經人和聖經的事情.
這次的故事是說.
約書亞說那些以色列人.
你們擔言不許得.
這是《生命說書》的網絡.
跟大家一起分享的.
你們擔言不許得.
上帝給你的,你也擔言不許得.
就被罵了,指責他了.
「以色列的全會眾都聚集在士羅.
將會幕設立在那裡.
那地已經被他們制伏了.
以色列人中其餘的七個支派.
還沒有分級他們地業」.
第十八章的開始是這樣的.
「全會眾聚集在士羅.
士羅是一個很重要的地方.
等同今天.
或者是蘇羅大衛王國的時候.
等同是耶路撒冷首都.
當時整個以色列會眾.
已經進入了迦南.
很大陣容的.
那裡真是每個民族在裡面都怕他.
因為他一進來是百多萬人.
再聽他們以前的歷史.
怎樣出埃及過蘆葦海,紅海.
行曠野.
然後耶利哥一大叫.
叫成又被他們攻破.
然後南方四個王.
五個王不夠他打.
北方無數個王都被他戰勝.
很厲害的.
他們聚集在一個地方.

$^{41}$後來就一路發展.
就用了士羅做他們的一個.
一個重要的首都的地方.
他們就將會幕設在士羅.
很長的時間.
這個地方那地已經被貪污制服了.
這個地是什麼意思呢.
這個真的有些難明.
究竟這個地是指士羅被制服了.
所以他們用了.
還是指整個應許地迦南地呢.
如果照上下文看.
還有些地你未拿.
這個地應該不只是指士羅.
是整個應許之地迦南地.
已經被他們制服了.
贏了 拿到手了.
說到是了才叫制服.
以色列中其餘有七個支派.
還未分到地.
還未給他們地.
不是還未給他們.
是還未拿 還未分.
譬如這樣.
譬如我們這間房很大.
我們七個人在這裡住.
十二個人在這裡住.
誰住誰的位置.
大家要商討的.
要劃定的.
有些大隻些 有些細隻些.
有些睡台底都可以.
每個地方是不同的.
其餘七個支派還未分地給他們.
我們看看怎麼做.
約書亞對以色列人說.
留意.
「耶和華你們列祖的上帝.
所賜給你們的地.
你們擔言不許得.

$^{81}$要到幾時呢」.
所以這句說話是.
是責問來的.
約書亞說.
「以色列人你們留意.
耶和華神給了地.
賜給你們的地.
你們擔言不許得.
不去拿.
要等到幾時.
等甚麼.
你們擔言不許得」.
這就是我要跟大家說的第一點.
很多時候.
電影節目有沒有發覺到.
有很多事情.
神已經為我們預備了.
開了門.
或者等著我們開門進去.
預備好了.
我們遲遲都不能去.
或者我們上班的時候.
有一個人已經看來.
「你去教會的嗎?.
去教會好像挺好的」.
「我兒子也想去主日學就好了」.
人家都開了聲.
你都不擔言.
不去傳福音.
有些人就說.
「是呀,我留在教會.
不過自己沒去過.
不敢進去.
有人陪我就好了」.
「可以,我陪你」.
你不要擔言.
傳福音不要擔言.
又或者是.
「我都很想多些讀聖經的.
教會開禪經班,你又不去.

$^{121}$有些甚麼你又不做.
都沒有機會」.
沒有甚麼機會.
是你擔言不許得.
很多事奉都是.
很多傳福音都是.
很多時候我們和神的關係.
都是我們擔言了.
不去拿.
於是我們得不到.
擔言不許得.
你們這樣要到何時?.
等甚麼?.
所以第一個.
我想和大家分享的就是.
不要擔言.
神很多時候給我們一些東西.
不要懶理.
反正都已經到手了.
全地,整個迦南地在我們手上.
我們怎樣分,慢慢來.
大家都是兄弟一場.
不趕的.
不是的.
不是的,要去拿.
約書亞就.
神就和約書亞說.
叫這些人你們擔言不許.
要等到何時?.
是不是等我要收回?.
很多時候你擔言得久.
神收回的.
「哎呀,為甚麼你收回的?.
不等我」.
等你很久了.
好了,第四次.
「你們每支派當選舉三個人.
我要打發他們去.
他們就要起身走遍那地」.
留意一下.

$^{161}$紅色字,藍色字,黑色字.
都是聖經和學本的經文.
綠色字就是我的註解.
「你們就要起身走遍那地.
按著各支派應得的地業.
寫明就回到我這裡來」.
你們每個支派找三個人.
這個選舉就是指派.
找三個人.
這三個人是怎樣找的?.
選舉投票.
原來聖經也有投票選舉的.
原來聖經教我們選舉.
是與不是.
很多時候.
聖經不是很選舉的.
這個選舉的意思是.
是選一些人.
正如你今天好像我這樣選.
今天選這件衣服.
選這個西裝.
選這個東西.
其實就是我選出來的.
是我有主權的.
好,那些族長.
或者那些族的支派裡面的長老.
就選三個人.
起碼選這三個.
「行得路,捱得先得」.
當時要走路的.
要他走遍那地.
懂得思考.
有概念.
每個支派當選舉三個人.
就「做ye 」.
剛剛上一節就說.
你們不要擔言.
如果不是神收回.
或者那些人.
發展了對抗你.

$^{201}$你就麻煩了.
趁現在我們制服了他們.
快點去拿.
等到什麼時候.
如果不是我們失了這個機會.
神收回你沒有了.
擔言.
接著就說「做ye 」.
第二就是找三個人.
每個支派找三個人.
就將那個地方走遍.
然後按著各支派應得的地頁.
寫明就回到我謝女郎.
說的是什麼呢?.
走遍整個迦南地.
然後按著.
原來按著這兩個字挺有趣的.
要留意.
按.
喂.
怎樣分那些地呢?.
是否將整個迦南地畫上九宮格.
大家量度好.
然後大家公正了.
每一個佔多大.
12個支派.
不是的.
早幾天說嗎哪西支派說.
我們人多.
我們支派強盛.
人多.
生畜又多.
所以大家如果一個尖溝分地.
我們一份不夠的.
我大隻.
要吃兩碗飯.
你小隻一碗飯就夠.
或者你還小隻半碗飯就夠.
所以人頭飯是公正.
但並不公平.

$^{241}$我要吃多點才飽.
於是大家又分.
嗎哪西是大些地方.
這個按字.
就是按照或者朝著.
著就是pear.
pear是口.
按著大家的口.
來到按著各支派應得的地.
寫了出來.
寫了或者叫翻譯.
或者叫畫圖.
什麼意思呢.
就是說.
大家每一個支派.
派三個人.
一起周圍走遍的地方.
然後照著你自己的人口.
人口有多少.
於是將自己要的地業.
這個地業.
就是財產產業.
就是按著你的人口.
和你的財產.
來畫出來.
人口譬如.
有些人口是小些的.
舉例子.
如果一會兒看到地圖.
你會知道.
變雅敏支派人少.
所以人口少.
所以當時分的地是小些的.
另外就是.
他的產業少.
什麼產業.
牛羊少.
他就要少些地方.
人口很多的嗎哪西.
他就要大些地方.

$^{281}$這樣分的話.
是按著.
就是跟著他們的人口.
有多少口人.
就是各支派應得的地業.
這個地業.
就是應得的地業.
是指他有多少財產.
有多少牛羊.
有多少帳幕.
就是寫明.
或者叫翻譯.
找人去走遍他.
寫好回來.
做功課.
就是三個字.
做功課.
大家拿那些地.
做功課.
第五字.
他們就要將地分為七份.
在七份之外.
為什麼七份呢.
十二支派嘛.
分七份.
其實因為尤大已經在南方.
住在他的境內.
原來尤大已經拿了南方了.
約瑟加仍在北方.
那約瑟加在十二支派是沒有約瑟的.
因為約瑟分成兩個兒子.
以法蓮和馬勒西.
而馬勒西很強盛.
正如馬勒西.
這個字就是昌盛的意思.
哇.
真是約瑟沒有改錯.
他第二個兒子的名字.
神使我在這個地方昌盛.
以法蓮就是忘記.

$^{321}$我忘記神令.
我忘記以前的痛苦.
現在使我在這個地方昌盛.
就是昌盛.
馬勒西說.
我可不可以要打這地方.
未入迦南的時候.
馬勒西就說.
這裡很適合我的平原.
約旦河東岸平原.
又是很寬闊.
很適合我牧養羊群,牛群.
不如我不入迦南了.
我在這裡,我要這裡.
摩西就說.
行,不過你一定要和以色列一起進去.
打完仗你才出來.
不要馬上停在這裡,算了.
仗要打的.
好,我們一起去打完才回來.
約瑟家就在北方.
你要將地方分成七份.
寫明就拿來我這裡.
大家協議好來給我.
然後我就在神面前為你們占溝.
占溝之後說.
大家都整好了.
然後我就在神面前禱告.
看看我透過這個占溝.
拿出來是不是.
神的心意在這個占溝裡面表明出.
是不是和你們想的一樣.
如果是就好了.
利未人在你們中間沒有份的.
沒有分地的,沒有分.
為什麼?.
因為他們公爺和華祭司的職任.
利未人要做祭司.
聖殿靈修的工作.
敬拜的工作.

$^{361}$而做這個工作本身就是貪媒的產業.
所以他們是沒有地的.
分七份.
在十二支派裡面.
利美是沒有地分給他的.
好.
迦德支派裡面支派.
和嗎哪西半個支派.
已經在約旦河東得了地業了.
所以他們是沒有的.
迦德沒問題,大家聽過了.
十二支派的名字.
裡面其實就是以前的舊的中文聖經.
譯作「劉便」.
第一個兒子,劉便西勉利美猶大.
劉便是大兒子,雅各的大兒子.
劉便現在中文翻譯叫做「呂便」.
為什麼呢?.
因為「劉便」不好聽.
令人想了別的不好的東西.
「劉便」.
於是就改了「呂便」.
改的名字都很好笑.
因為「呂便」之後.
中文廣東話「呂便」.
是不是有「外便」.
「呂便」出「便」.
就「呂便」,不過不要緊.
於是你會見到.
有嗎哪西約瑟加.
猶大在南方.
然後嗎哪西迦德裡面.
就是嗎哪西的半個支派.
就在北方.
和約旦河的東岸有地業了.
因此剩下的就是七個支派.
這裡大家拿了.
然後約瑟加第十八章第十一節.
便雅敏支派最先抽.
按著宗族尖溝所得的地.

$^{401}$就在猶大約瑟子孫中間.
他們七個支派抽的.
便雅敏是第一尖溝.
就得了他們.
然後到第十九章.
第一節西面又抽了.
然後到西部輪又抽了.
以撒迦又抽了.
第四第五.
阿撒切是第五.
拿忽他利是第六個尖溝.
但就是第七個.
於是抽完了.
大家留意一下這幅地圖.
要先留意以色列十二支派分的地.
是這樣的.
利美會有土地.
散居在十二支派之中.
在十二支派裡面做祭祀的工作.
靈修,敬拜等等.
約瑟由兩個兒子.
以法蓮和馬來西代替了.
所以你沒有見到約瑟支派的地.
不過有見到以法蓮和馬來西的地.
這是以色列的地圖應許之地.
這裡是約旦河流下來.
加利利海繼續流下來.
去到死海.
當時的死海是未斷的.
水是充足的.
現在斷了流.
阿徹支派在這裡.
如果這個時候.
今天叫推羅西頓這個地方.
黎巴嫩那邊.
不是.
這樣就高一點.
阿徹在這裡高一點.
由北面的部分黎巴嫩到這裡.
阿徹支派,西部論.

$^{441}$旦有兩個地方.
北面拿了一些.
其實旦是分到這裡的.
旦支派.
拿忽他利是這個綠色的地方.
加利利海.
以薩加可能人不是很多.
所以在這裡小小的.
西部論又不是很多.
可能人口按著口.
和他的財產.
馬來西是厲害的.
這裡是一條約旦河.
在進來的時候.
部分馬來西就說.
我們想在這裡牧羊.
平地.
不如我們不進去.
我們住這裡.
可以.
不過你要進去打完仗.
一起進去打完仗之後.
再抽籤.
馬來西就拿了這個地方.
所以整個馬來西.
就在約旦河的東和西都有地.
馬來西是最大面積的.
一個以色列支派.
然後就加德.
又在約旦河的東.
然後裡面又在約旦河的東.
變亞敏就在這裡小小的.
以法連就是紫色這一塊.
旦也不大.
猶大就大了.
猶大支派很多人.
整個猶大.
西面支派就在猶大當中.
四面都是猶大.
好了.

$^{481}$變亞敏在裡面.
大家看到了.
一定要記得這條約旦河.
和地中海的海邊.
約旦河東.
約旦河西.
這樣.
好了.
如果看.
咦.
今天的耶路撒冷.
是在哪個支派的地方裡面呢.
是在變亞敏裡面的.
那麼.
耶路撒冷在變亞敏支派的地裡面.
那麼這個就是他們的地.
上面支派的分地了.
那麼就.
沒有了利美.
那麼就有一個.
以法連和馬來西去補回的.
那麼約旦河有兩個.
好.
上面支派分地.
不過這個圖我想給大家看到就是.
如果隱約看到非利士.
非利士就是這幾個城.
就是阿薩德,阿薩基倫,加薩,加德.
這些就是這一帶.
海邊的路.
沿海.
非利士.
是時時和以色列作對的.
在這個地圖也是拍出來的.
只是有點模糊.
所以要加一些字上去.
猶大支派是佔了非利士地的.
應該是他的.
不過有另一幅圖.
大家看到了.

$^{521}$一樣了.
那麼這張圖我特別要講一下了.
這些沒有的.
但,阿,切.
但,阿,切推落在這裡.
阿,切.
西布倫.
以薩加.
以薩加有個字在這裡.
但是紅色一點的.
西布倫很好.
最西是西布倫了.
那麼馬來西就.
哇.
河東.
河西.
黃色也是馬來西的.
白色這一塊就是以法連的.
但就是綠色.
深綠色這一塊.
然後呢.
猶大支派是這樣的.
黃色這個.
那麼裡面就在約旦河的東邊.
裡面就是柳邊.
迦德在這裡.
然後馬來西.
那麼現在看到了.
拿伏塔利.
好了在這裡我特別要看一下.
其實這張圖我想指出.
當時他們雖然叫做.
十二支派分了地.
有一個特別的地方.
本來如果上一張圖.
你會看到.
非利士地在猶大裡面的.
但是這張圖呢.
猶大是並沒有包括非利士這個地的.
阿什特 迦德 阿什 基倫 加薩等等.

$^{561}$加薩地帶就是今天非利士的地方.
那麼是沒有畫在猶大裡面的.
那就是說這兩張圖有分別.
分別就是猶大得著了土地之後.
這個非利士地應該是屬於他的.
不過他沒有得著.
他沒有去打.
打贏他沒有得著.
或者是贏了之後又退了.
所以非利士一直都還在這裡.
那綠色字呢.
我就想大家留意一下.
原來綠色字就是圍著.
以色列十二支派的那些外族.
如果你看一下路聖經下來.
就是這幾個外族.
是經常和以色列人打仗.
為害不淺的.
就是他們沒有去制服他們.
然後自己就被騷擾了.
就時打仗了.
非利士在這裡.
那麼你見到大衛也要和非利士打仗.
數萊也是.
那麼這裡呢.
加特也有一個人出來.
哇,九尺高.
對人哥利亞.
那麼時時和亞馬力人爭戰.
亞馬力就在非利士下面.
這個位置就是亞馬力了.
那麼以東就是約亞國的哥哥.
就是以素,以東人.
那麼如果你看《先知書》.
以東人又是時時後後兄弟幫得好.
都是和以色列經常爭鬥的.
摩亞就在這個位置.
就是羅德的後裔了.
那麼然後阿滿也是羅德的後裔了.
今天這個地方已經變成約旦了.

$^{601}$那麼約旦的首都叫做安曼.
安曼阿曼.
那麼你見到就是阿滿.
就是阿滿.
安曼就是這個位置了.
然後北邊還有亞蘭.
亞蘭王又經常在北邊搞以色列的.
那麼亞蘭的關係是好一點的.
但是經常和他打仗.
有時候會一起打.
是敵是友.
我想講出這幾個綠色的牌.
是告訴我們以色列沒有去收服他們.
終於變成自己很多禍根了.
以色列人按著他的警誡分完了地業.
就在他們中間將地給亂的兒子約書亞為業.
那麼十二個支派都有地了.
那麼亂的兒子約書亞.
都當機一功有沒有地給他呢.
是按照耶和華的吩咐.
將約書亞自己希望的.
是以法尼這個山地.
庭拿西拉這個城給約書亞.
約書亞就收建那個城住在當中了.
所以約書亞和迦勒這樣有個城.
記不記得迦勒要求的城是什麼.
希伯倫.
就是這個基列阿巴.
得了之後改名希伯倫.
聯合這樣的意思.
那麼這個約書亞就說.
我歡喜以法尼山地庭拿西拉.
就做我的住的地方了.
這個就是祭司以利亞薩和亂的兒子.
有個祭司現在出現了.
有個祭司就叫以利亞薩.
是全國的大祭司了.
和亂的兒子約書亞.
這個就是管信仰宗教了.
這個是管政治軍隊了.

$^{641}$並以色列各支派的自己的族長.
都一起在士羅回望門口.
又說是面前尖溝分的地.
剛才那個地圖大家看到了.
第三,他們把地分完了.
第一神就說.
你不要擔言了.
擔言我收回了.
擔言你拿不到了.
而且你丟久了.
東西就會壞了.
要等到什麼時候.
於是他們約書亞叫他們起來工作.
每一族派三個人.
每個支派派三個人.
到處去看好,量好,做量地官.
看好跟著人口等等.
自己的財物有多少.
就劃了地了.
然後在神面前尖溝.
就確實是神給的了.
於是這樣他們就把地分完了.
「他們擔言不許得」.
就是今天要跟大家看的題目.
記得,弟兄姊妹.
不要擔言.
神說你為什麼要擔言?.
那個人要信主.
是我安排了機會在你旁邊出現的.
你不要擔言.
要快點跟他傳福音.
教會有一個需要要有人做的.
你不要擔言.
感動來到了,不要擔言.
有一個宣教工作想叫你去.
不要擔言.
有一個事奉綱位給你.
不要擔言.
有些經文要讀.
想你查一下經文.

$^{681}$不要擔言.
想你進修校國.
不要擔言.
想你做一些事情.
不要擔言.
就是這樣的了.
第二就是馬上做.
自己做不到都叫人做.
派人去做.
於是畫好圖.
清清楚楚在神面前求印證.
好了,然後第三.
工作就是這樣做完了.
就是這樣神樂歡喜的.
不要擔言.
這個是今天跟大家分享的.
今天的故事就是以色列分地.
分了地之後.
就變成有他們的地圖了.
就是這樣了.
安排好.
其實四景平定.
十分感恩.
那些外族當時都沒有什麼搞作.
後來約書亞不在.
哎呀.
就軍心大亂.
出現到事思時代了.
到時我們再看吧.
我們一起祈禱.
你是按著你應許.
將這一塊應許之地.
真的給了以色列.
是按著他們十二支派分了的.
我們看到你的信實廣大.
你說過會做的就會做.
你應許的就會臨到.
多謝你.
你是給這一班本來都不會打仗.
做奴隸出身的人的下一代.

$^{721}$就得到這塊地.
完全不是因為他們好打.
是每個當地的民族都害怕.
都知道他背後的耶和華上帝厲害.
就在以色列當中顯出神自己的能力.
顯出你的高貴.
顯出你的大能.
讚美主.
幫助我們.
當你叫我們做事.
要得一些東西的時候.
讓我們不要擔心.
我們很多惰性.
怕辛苦.
好日烏魯.
幫助我們努力去.
得主給我們的地.
得主給我們的土地.
得主給我們的經歷.
得主給我們的事奉.
讚美主多謝你.
禱告奉靠耶穌基督明求.
阿們.
\newpage



\section{使徒行傳 15:36-41}
\label{sec:032160}
\textbf{起了爭端,彼此分開}~~~~李應新牧師
\newline
\newline
連結: \href{http://www.church.com.hk/acms/content.asp?site=cdc\&op=show\&type=product\&code=032160\&layout=sermon}{\texttt{032160}} ~~~~ 語音日期: 2026-05-17
\newline
\newline
\hyperref[sec:032159]{\small{< < < PREV SERMON < < <}}
~
\hyperref[sec:index_chronic]{\small{[返順時目]}}
~
\hyperref[sec:032161]{\small{> > > NEXT SERMON > > >}}
\newline
\newline
使徒行傳 15:36-41
\newline
\begin{longtable}{cl}
\hline
\hline
章節 & 經文 (和合本修訂版)\\
\hline
15:36 & \begin{tabularx}{0.7\textwidth}{X} 過了些日子，保羅對巴拿巴說：「讓我們回到從前宣揚主道的各城，看看弟兄們的情況如何。」 \end{tabularx} \\ \\ \relax
15:37 & \begin{tabularx}{0.7\textwidth}{X} 巴拿巴有意要帶稱為馬可的約翰同去； \end{tabularx} \\ \\ \relax
15:38 & \begin{tabularx}{0.7\textwidth}{X} 但保羅認為不宜帶他去，因為馬可從前在旁非利亞離開他們，不和他們一起工作。 \end{tabularx} \\ \\ \relax
15:39 & \begin{tabularx}{0.7\textwidth}{X} 於是二人起了爭執，甚至彼此分手。巴拿巴帶著馬可，坐船往塞浦路斯去； \end{tabularx} \\ \\ \relax
15:40 & \begin{tabularx}{0.7\textwidth}{X} 保羅則揀選了西拉，也出發了，蒙弟兄們把他交於主的恩典中。 \end{tabularx} \\ \\ \relax
15:41 & \begin{tabularx}{0.7\textwidth}{X} 他就走遍了敘利亞、基利家，堅固眾教會。 \end{tabularx} \\ \\
[1ex]
\hline
\hline
\end{longtable}
$^{1}$我們一起相聚 一起同行.
今天牧師為何不講《爾賽亞書》.
因為是延伸來的.
上一次我講的時候是約《爾賽亞書》第49章第6節.
我不知為何我很想講完這裡去做一個對應.
第49章講到什麼呢?.
講到今天沒有誣陷你《愛本聖經》.
因為我客行遲了回來.
沒有誣陷你《愛本聖經》.
第49章如果你讀的話.
亞國眾支派復興 使以色列中得保存的可以歸回.
是什麼呢?.
尚為小事.
什麼是大件事呢?.
大件事是.
我還要使你作外邦人的光.
叫你施行我的救恩.
直到地極.
施行救恩直到地極.
這些才是大件事.
其實是小事.
我們大小不分.
於是乎很多苦惱在生命當中.
不知你有沒有這個感覺.
多數是大事不著重.
小事經常針對著.
家裡經常爭吵.
如是.
或者我們自己生命都是.
樣樣都要.
你又想男朋友又帥又有錢又疼愛你.
韓劇才有這些東西.
平時在街上很少見.
但你樣樣都要.
你都要分哪樣是最重要.
哪些呢?.
有就好了.
有都是好的.
不是說不好.
沒有的話.

$^{41}$那就沒有了.
在這裡又講回頭.
牧師我接下來.
報告都有寫了.
今晚我走了.
今晚我就放假.
直到六月中才回來.
那麼大.
那麼激動.
喂.
我很少放假.
你覺不覺得我.
我幾乎鞠躬盡瘁.
死而後已.
很少放的.
不要妒忌.
很少放.
就這樣說.
我回去.
我女兒正在考A級.
我正在考A級.
其實她將會讀大學.
其實讀大學的話.
很貴的.
英國讀大學.
是傾家蕩產的.
你的樣子好像覺得很少錢.
你的樣子全部都覺得很少錢.
對於我來說.
傾家蕩產.
我心裡面覺得傾家蕩產.
妹又讀大學.
哥又讀大學.
傾家蕩產.
我死定了.
不過我再想多一步.
我馬上豁然開朗.
你是否想她不讀.
是否考得差一點.
就讓你不用傾家蕩產.

$^{81}$只剩下很多錢.
你是否想這樣.
我想完之後就說.
豁然開朗.
可以嗎.
感受到嗎.
馬上開朗了.
我寧願傾家蕩產.
她傾家蕩產比她更好.
爹地.
你既然傾家蕩產.
我幫你.
我考差一點.
入不到U.
你安樂了.
不用.
千萬不要想這些.
可以嗎.
哪樣比較重要.
比較重要.
先走.
其他沒有那麼重要.
有就好.
沒有就無所謂.
因為天父總包底.
我記住.
天父總包底.
可以嗎.
今天跟隔壁的說.
天父總包底.
天父是包底的.
知道嗎.
天父包你底.
所以我怕甚麼.
做甚麼事.
我們應該無關.
天父包底.
回看聖經.
聖經說甚麼.
「過了些日子」.

$^{121}$「保羅對巴拿巴說」.
「我們可以回到從前宣傳主導的國城」.
「看望弟兄警方如何」.
「過了些甚麼日子」.
「原來保羅和巴拿巴」.
「第一次福音旅程去了哪裡」.
「去了很多地方傳福音」.
「然後建立教會」.
「建立教會」.
「這個也是保羅按照主的心意」.
「這個最重要的事」.
「我們那個不是廉價福音」.
「所以我們要付出沉重代價」.
「其實他們去傳福音是很多代價的」.
「如果你看保羅書上說」.
「又有沉船的危險」.
「又有風浪的危險」.
「又有盜賊的危險」.
「各樣危險」.
「去完一次旅程之後」.
「帶了一些人信主」.
「不是就這樣就算了」.
「他很希望他們甚麼」.
「希望他們可以得到建立」.
「所以他們要看望弟兄警方如何」.
「又說回頭」.
「其實那些情況是挺困難的」.
「傳完福音就有很多外邦人信主」.
「外邦人信主又有問題」.
「一大堆會湧出來的」.
「有甚麼問題」.
「原本主要是猶太人信主」.
「他們就守律法」.
「平時都守開的就繼續守」.
「但那麼多外邦人信主守不守律法」.
「守又守多少」.
「哪些不用守」.
「於是又煩一輪」.
「就去耶路撒冷大公會議」.
「然後又派保羅回去」.

$^{161}$「告訴安提哥教會」.
「哪些要守哪些不用守」.
「哪些要守」.
「淫亂和那些偶像就要守」.
「其他的可以不守」.
「可以嗎」.
「於是又派他回去」.
「過了些日子就過了些日子」.
「過了一段時間」.
「又說回頭」.
「剛才所說的」.
「問題不是很多嗎」.
「很多問題」.
「難道你想沒有外邦人信主」.
「外邦人信主就有很多問題」.
「正常是吧」.
「有問題」.
「沒有問題」.
「沒有問題才有問題」.
「可以嗎」.
「因為神會包底」.
「於是就派他們回去」.
「把問題的答案解答完之後」.
「又準備第二次福音旅程」.
「要做栽培了」.
「這裡說到栽培」.
「還想多說一件事情」.
「他們看望弟兄警方如何」.
「其實他們已經做了一些事情」.
「要跟進的」.
「其實成立完教會之後」.
「在《少林寺傳》第十四章二十三節」.
「選立了長老禁食禱告」.
「即是說他們走之前」.
「已經解決了」.
「在各個傳福音據點之後」.
「不過都要看看弟兄警方如何」.
「報導完之後」.
「問題就有很多出現」.
「不過又說回」.

$^{201}$「我們近年需要面對的問題」.
「究竟如何傳福音呢」.
「因為二零二七年」.
「可能我們報告宣傳」.
「叫你去那些動員大會」.
「下個星期好像就是了」.
「我們又會搞福音盛會」.
「二零二七年在哪裡搞」.
「香港搞大型報道聚會」.
「最先的時候」.
「我們計劃是想去啟德的」.
「去啟德」.
「厲害嗎」.
「想起都厲害」.
「不過太厲害」.
「不是我們沒有錢給」.
「其實我估我們都籌到的」.
「如果要去啟德」.
「不過覺得未必一定會」.
「預訂到場」.
「因為很多人跟你爭」.
「你估不估得每個星期都有傳播會」.
「我開頭想」.
「我以為那裡是用來養草的」.
「原來不是」.
「那麼多人」.
「每個星期都差不多滿了」.
「還有些明星我都不認識」.
「那些國內又有來又有甚麼來」.
「我真是背」.
「其實一半認識的」.
「是那些林憶蓮認識的」.
「其他那些不認識的」.
「但原來都爆了幾場」.
「所以估計」.
「你爭到的機會不是很確定」.
「機會不高」.
「所以哪裡比較高機會」.
「大球場蘇廣寶那個」.
「沒有人要」.

$^{241}$「不不不都有人要」.
「我們要」.
「所以基本上都很穩陣」.
「但政府未批了說可以的」.
「但我們估都沒有人跟你爭」.
「所以都很穩陣」.
「11月出2027年」.
「不過又說回頭」.
「近年不知道你有沒有感覺」.
「爭論是否要大型報道會」.
「爭論得很厲害」.
「日夜爭論都有」.
「就說這個成效不高」.
「又跟進不到」.
「其實我覺得」.
「這裡聖經都有提到一樣東西」.
「我都覺得很想跟你看」.
「我繼續講下去」.
「我跳下去再看」.
「我讀一段經文」.
「經文是這樣說的」.
「過了些日子」.
「保羅對巴拿巴說」.
「我們可以回到從前宣傳主導的國城」.
「看望弟兄警方如何」.
「巴拿巴帶著馬可坐船往居比路去」.
「保羅揀選西拉也出去」.
「蒙弟兄們把他交在主的恩中」.
「他就走遍敘利亞,基尼加」.
「堅固宗教會」.
「我這樣讀你覺不得順」.
「是否都幾順」.
「聽不聽得到很順」.
「聽完我讀」.
「感覺都沒有違和感」.
「你閉上眼聽」.
「我當然沒有誣陷你我那篇聖經」.
「就很順」.
「大家有否留意到」.
「留意到甚麼」.

$^{281}$「沒了兩節」.
「沒了三十七,三十八」.
「我希望你退回PowerPoint」.
「退回前面讀經那裡」.
「三十八節」.
「剛才我沒有讀到哪裡呢」.
「就是巴拿巴有意」.
「帶稱呼馬可的約翰同去」.
「但保羅因為馬可從前」.
「在龐菲利亞離開他們」.
「不和他們作供」.
「就以為不可帶他去」.
「三十九」.
「於是二人起了爭論」.
「彼此分開」.
「我跟著才讀巴拿巴」.
「明白」.
「剛才那裡我跳過了」.
「那裡沒有讀」.
「但順的」.
「如果我作聖經的話」.
「我就不作在那裡了」.
「我就照我剛才的版本讀就可以了」.
「但我現在呢」.
「但仕途行轉著者」.
「就不是這樣寫」.
「他加了這裡」.
「不是加進去」.
「他放了這件事件進去」.
「可加可不加對於我來說」.
「可以嗎」.
「不加這件事是順的」.
「但加了是他的意圖」.
「可不可以」.
「明不明白我」.
「感不感受到我想說的問題」.
「這件事是可以說的」.
「不過這樣說」.
「重點是甚麼呢」.
「重點呢」.

$^{321}$「就是原來教會是不介意」.
「將我們的爭論」.
「我們的起了爭論」.
「和彼此分開告訴你」.
「我們不需要擺一些」.
「矛盾在地氈底的」.
「OK的擺出來」.
「但我一會兒再說」.
「我再引申我剛才說的」.
「最近很多人罵你」.
「報道會是否應該這樣做」.
「我覺得都OK的」.
「百花齊放是很OK的」.
「應該個人報道」.
「個人報道都幾好」.
「我們其實香港有個機構」.
「叫做街報」.
「有聽過嗎」.
「街頭報道」.
「那個是」.
「剛剛回到天父那裡那個」.
「柳振平牧師」.
「走遍那麼多地鐵站」.
「他是新鮮事出」.
「最熱心就是柳振平牧師」.
「我都覺得佩服」.
「他每個地鐵站」.
「早上他選擇上班時段」.
「在那裡喊咪高峰」.
「即是街頭報道」.
「走遍那麼多個地鐵站」.
「非常之厲害」.
「這隻都幾好」.
「如果你覺得大型報道」.
「這隻夠小型」.
「走遍各街頭」.
「但亦都有些人不是很喜歡的」.
「有些人批評他是甚麼呢」.
「就覺得這隻報道方法」.
「比較cheap cheap的」.

$^{361}$「即是在那裡走來走去」.
「真是大減價」.
「即是那些呢」.
「他說太像那些」.
「即是形象不太好」.
「這隻街頭報道」.
OK 既然是這樣.
「友誼報道來」.
「即是我們不如友誼報道」.
「但我們都幾流行」.
「叫你回來帶些朋友回來」.
「我們教會友誼報道」.
「即是我們有搞這樣搞那樣活動」.
「叫你回來」.
「沒有人說了吧」.
「導遊」.
「這隻會影響友誼的」.
「我們太hard sell」.
「迫人回來教會搞這些」.
「我們其實要professional些」.
「即是如果你說那些」.
「我們又要搞又分不上水準」.
「我們教會有時都未夠水準」.
「那些東西又要去那些」.
「如果要跑步又要去專門跑步」.
「去那些搞甚麼講座」.
「去那些專門的」.
「還有我們監人來的話」.
「就很hard sell」.
「總之就不是很好」.
「原本有友誼」.
「都搞到那些友誼沒有了」.
「那不是更壞嗎」.
「有聽過嗎」.
「我也有聽過有人這樣說」.
「真是混亂」.
「那即是怎樣」.
「o下」.
「甚麼都不要搞」.
「我真是暈了」.

$^{401}$「最不好就是這樣」.
「甚麼都不搞」.
「其實我的看法是甚麼呢」.
「希望你體諒」.
「我的看法是各有各搞」.
「你喜歡搞甚麼就搞甚麼」.
「但是務要全盜」.
「無論得是不得是」.
「這是主給我們的使命」.
「你說這隻不好」.
「但是我想逐個逐個講」.
「尤其是最後那隻」.
「你說我們hard sell」.
「不是很好」.
「你看教會」.
「我們太古城堂」.
「很長時間都未到」.
「你怕都怕別的樣子」.
「你擔心你不夠努力」.
「叫人回來教會好過」.
「你何時我們這隻教會會甚麼呢」.
「就是會hard sell」.
「所以我又講回頭」.
「各有各困難」.
「但不是彼此祝福」.
「其實這裡都是實踐這件事」.
「於是二人雖然起了爭論」.
「彼此分開」.
「不過到最後他們有沒有傳福音」.
「照傳」.
「照跟著神的心意」.
「保羅就選了阿西拉」.
「巴拿巴就選了馬可」.
「兩個都去了不同的地方」.
「一個去俱比路」.
「一個就去了敘利亞,基尼加」.
「到最後都能夠堅固宗教會」.
「行不行」.
「這亦都是最初他們的賞望」.
「沒有defeat到的purpose」.

$^{441}$「不過又講回頭」.
「很多時候我們叫完教之後」.
「就忘記了原本叫甚麼教」.
「不知道你有沒有這情況」.
「叫著叫著都不知道你叫甚麼」.
「我們最初是叫甚麼的」.
「最初想做的事沒有做到」.
「其實這招我覺得有些兒子都學了」.
「有些小朋友,嬰兒」.
「你走去說他沒有禮貌」.
「他又講了別的樣子」.
「岔開來講又走來哭」.
「defeat了甚麼」.
「安慰他不要哭」.
「然後你忘記了罵他」.
「他們都懂得用這些招數」.
「小心提防」.
「免得我們忘記了最初」.
「最初就是要傳福音」.
「最後還是要傳福音」.
「就分開傳」.
「你有你的方法」.
「我有我的方法」.
「但我們不要傳到」.
「無論得時與不得時」.
「都要傳到為時」.
「不過在此說到」.
「我們也說回頭」.
「有時傳福音都不一定很順利的」.
「有時都是有困難在當中的」.
「譬如我們剛剛」.
「我們短宣隊」.
「待會看片你會看到」.
「我們去了澳門短宣」.
「我們其中」.
「我們播道會」.
「其實以前都挺多宣教士去澳門的」.
「其中有一檔事我就沒去過」.
「不過我有支持過」.
「我跟他以前的服務生都挺友好的」.

$^{481}$「那個叫寧亞諮宣教士」.
「就搞成人教育中心」.
「我昨天說到」.
「我分享說而已」.
「馬太說我去過」.
「馬太你比我厲害」.
「我只是說而已」.
「我都沒去過」.
「澳門成人教育中心」.
「澳門有幾檔事」.
「其實都搞得不太順利成功」.
「所以我知道十年八年前」.
「不知道多少年前」.
「總之澳門成人教育中心倒閉了」.
「聽說過」.
「有沒有些難過的感覺」.
「是呀 你有份去」.
「以前開張的時候好像有份去」.
「我真的沒去過」.
「不過無論如何」.
「我全部都是聽說而已」.
「聽他說開頭挺好生意」.
「但後來怎樣搞就倒閉了」.
「聽說又有些不開心什麼的」.
「總之總有些事」.
「只是沒生意就不止倒閉了」.
「還有些事」.
「總之不知道甚麼事」.
「最後就倒閉了」.
「心裡都有點難過的這件事」.
「澳門果然是個福音硬土」.
「這麼難做的」.
「最後寧雅芝孫老師去了另一個工場」.
「我想有些弟兄姊妹都知道」.
「我都繼續有follow她」.
「無論如何我有跟進」.
「但總之沒有了這個工場」.
「上次我們真的有再去澳門」.
「我去澳門的時候探訪了不同的機構」.
「包括我們播導會」.

$^{521}$「慈愛家庭」.
「詳情的話應該是星期三晚來到我們祈禱會」.
「就會詳盡羅姑娘講給你聽」.
「慈愛家庭」.
「小小一題就是」.
「我們由我們的孫教士」.
「戴文牧師發起的」.
「智障人士孩子的家長家庭」.
「來參加敬拜崇拜的地方」.
「叫做慈愛家庭」.
「聽說他都有計劃過」.
「即是他開的時候都有點老了」.
「已經計劃了」.
「就要找succession」.
「聽說要找董姑娘去幫他」.
「誰知董姑娘」.
「戴文牧師後來就退休了」.
「回到美國了」.
「董姑娘原來都是多做一年而已」.
「她都退休了」.
「我們香港教會就是這樣」.
「找人succession」.
「原來都很老了」.
「差不多是時候了」.
「不過原來天父是有預備的」.
「到時羅姑娘會講給你聽天父的預備」.
「這是第一站」.
「慈愛家庭」.
「第二站我們就去澳門聖經學院」.
「其實我這次去」.
「其實我開頭不去都可以的」.
「我去是想探一探一個」.
「灣仔唐時代的團友」.
「那個叫林牧師」.
「那個林牧師我小時候」.
「他已經是過去澳門的」.
「他好像去的時候都未結婚」.
「剛剛結婚而已」.
「個個都挺」.
「他的子女都大了」.

$^{561}$「他的兒子有個女兒都結婚了」.
「在那裡做了幾十年」.
「我說幾十年都未去過」.
「我們去探他都要去一下」.
「我們太古城堂去」.
「所以我特意去而已」.
「聽他分享一下澳門的情況」.
「聽完之後」.
「我突然間覺得」.
「人生充滿希望」.
「因為我們不行」.
「剛才說了甚麼」.
「天父會」.
「爆底了」.
「未講完即是怎樣」.
「原來他也當是這樣的」.
「黃永信牧師」.
「你聽過他的名字吧」.
「你不知道他是誰」.
「聽他的名字都覺得他厲害」.
「創立澳門聖經學院」.
「當年原來他在北美籌款」.
「買了一棟舊樓」.
「即是澳門」.
「等於香港有些戰前唐樓那些」.
「沒有電梯的」.
「走幾層的那些樓」.
「用來做聖經學院的」.
「聽說是這樣的」.
「挺好的」.
「樓上用來做宿舍」.
「中間用來做圖書館」.
「樓下用來做課室」.
「是否都挺完美」.
「不過這些是不規範地區」.
「才搞這些東西而已」.
「原來後來澳門又規範化這些東西」.
「這些地方是住宅來的」.
「住宅又不可以用來做學校」.
「要拿牌照」.

$^{601}$「你拿不到牌照」.
「連申請都不能申請」.
「你這樣完全不對」.
「順便額外說一下」.
「剛才有些教友跟我說」.
「喂!牧師」.
「有地方賣」.
「隔壁間國際學校賣了」.
「知不知道」.
「隔壁間學校不知多少億賣了」.
「牧師!我們又買了一些」.
「好像街市一樣」.
「隔壁賣學校不等於這裡賣這裡」.
「是不同的」.
「隔壁很便宜」.
「你看人家的牌照」.
「人家的牌照應該是學校的牌照」.
「學校的牌照的價錢當然跟這些不同」.
「轉用途要補地價」.
「無論如何不要說太長了」.
「插回去」.
「剛才說到哪裡」.
「澳門聖經學院黃永信牧師買舊樓」.
「政府又說要規範化」.
「學校要找一些商業」.
「要取牌照等等」.
「總之又要合乎防火等規範」.
「才可以做到學校」.
「問題又來了」.
「問題真的頗大」.
「問題大在於甚麼呢」.
「你馬上去買了一層樓」.
「首先要買了一層樓才可以買樓」.
「因為沒有錢」.
「買完樓又要裝修」.
「又要入紙申請去處理」.
「你不是馬上可以變出來的」.
「同意嗎」.
「這些事要搞多久」.
「說不定」.

$^{641}$「最壞打算是要停一年」.
「停一年」.
「沒有辦法」.
「很難延續得到這麼多事」.
「是不是都很頭痛」.
「剛才說了甚麼」.
「天父會有安排的」.
「我們靠天父的」.
「原來」.
「原來甚麼」.
「我們的成人教育中心倒閉了」.
「最適合賣給他」.
「我們重點不是因為有樓在」.
「重點是甚麼」.
「我們的牌都齊了」.
「就是教育牌全部都搞定了」.
「即用」.
「天父是有預備的」.
「同意嗎」.
「看到林太我真的忍不住了」.
「我們播道會專門做這些事」.
「我們當年在廣州的時候」.
「我們就是Pioneer」.
「我們Pioneer是甚麼」.
「買福音船回來」.
「不過買完之後」.
「我就叫啟明號」.
「我上次我有講過」.
「你有沒有印象講過」.
「啟明號」.
「有些人好像未聽過」.
「我一講啟明號之後」.
「完了崇拜林太就走來跟我說」.
「啟明號不是你們的」.
「是我們的」.
「你們教會叫甚麼名字」.
「你又忘記了」.
「是東方基督教會」.
「李氏牧師」.
「李氏牧師的啟明號」.

$^{681}$「那隻啟明號是我們賣給你們的」.
「就是我們做不來 倒閉了」.
「所以賣給你了」.
「不過我做不來而已」.
「他們東方甚麼基督教會」.
「我忘記了名字」.
「水上基督教會」.
「差不多吧」.
「你們發揚光大」.
「後來買多兩隻」.
「有隻運了來油麻地」.
「香港仔還有一隻」.
「我都認識得少」.
「怎會不感受到」.
「我做不來而已」.
「但我開了個頭」.
「我開了個頭之後」.
「不是我」.
「那個是誰」.
「寬跨林牧師」.
「Anyway」.
「我們開了個頭」.
「因為上午有計劃安排帶領」.
「現在造就了事工」.
「於是他們去了香港仔和油麻地」.
「創立教會」.
「林太在這裡回教會信主」.
「所以我想說回頭」.
「天父是有預備」.
「我們只需要忠心做好工作就可以了」.
「吵架少少吵難以避免」.
「不過吵完之後」.
「繼續要做主的工作可以嗎」.
「巴拿巴和保羅」.
「都沒有忘記最初的使命」.
「繼續努力地傳揚福音」.
「求神幫助我們」.
「講到這裡的時候」.
「我還想再提一件事」.
「人生也是如此」.

$^{721}$「首先講」.
「我們的爭議是難以避免的」.
「因為其實兩個都講得通的」.
「怎樣講得通呢」.
「其實每個人都講得通你的道理」.
「你感受到嗎」.
「有時你又對他又對」.
「不過我想說人生也是如此」.
「人生如此是甚麼呢」.
「我們最初的目標」.
「究竟神給我們的目標是甚麼呢」.
「教會的目標是很清楚的」.
「教會的目標呢」.
「就是傳福音」.
「要叫大使命」.
「還有呢」.
「實踐彼此相愛」.
「我們彼此相愛愛人如己」.
「就是大誡命」.
「我們大使命加大誡命都要進行」.
「這樣就完結了」.
「量度這兩件事就夠了」.
「我們人生呢」.
「其實都差不多」.
「我們人生呢」.
「榮耀上帝的名」.
「將天父的慈愛帶給我們身邊的人」.
「同不同意」.
「我們完成了使命」.
「那麼中間有沒有些困難呢」.
「難免有些困難」.
「有些挑戰」.
「最近我又遇見一位弟兄」.
「剛剛初信主的」.
「他的情況呢」.
「我們都為他代禱叫Kelvin」.
「Kelvin弟兄呢」.
「好像我都同你講過」.
「就被咽癌在英國管理醫」.
「英國就話要切jor 隻眼」.

$^{761}$「跟住再醫」.
「香港話打jor 缺無線」.
「睇下得唔得」.
「不過最後咩都係唔得」.
「都係要切」.
「咁當然弟兄心裡面就好沮喪」.
「跟住前路會點呢」.
「我同佢分享ge 時候」.
「我自己有第一身個人經驗」.
「我第一身個人經驗係咩呢」.
「我爸爸都係Cancer」.
「我爸當年係中學」.
「其實同佢仔而家差唔多」.
「初中ge 時候我爸係搶癌架」.
「咁但係我覺得雖然有疾病」.
「但我地可以沒有欠缺」.
「頭先講gwo 兩樣ye 我地做到架喎」.
「邊兩樣ye 話」.
「榮耀上帝」.
「同埋將神ge 愛帶比我地身邊ge 人」.
「患病帶唔帶到呀」.
「患病都可以架」.
「我更加鼓勵弟兄」.
「可能真係會冇jor 隻眼睛」.
「但我地仍然靠天父ge 恩典同埋力量」.
「可以去到想去ge 地方」.
「做到想做ge 事情架喎」.
「我地可以生命面沒有欠缺架喎」.
「Mandy係咪呀」.
「到到都去到」.
「有人推咪得啦」.
「唔係 自己黎都得」.
「神俾我地力量」.
「其實我同佢講架」.
「弟兄 其實你都可以展現在你ge 兒女ge 面前」.
「其實你個仔」.
「將來都會面對唔同ge 艱難架啦」.
「而家係你demonstrate」.
「梗係點樣去應對而家面對ge 挑戰」.
「你ge 仔係睇住你架」.

$^{801}$「神俾力量你」.
「眼可能真係會冇jor 」.
「但係生命存在係度我格外講返多樣ye 」.
「其實偉大兄都有D 沮喪架」.
「係佢沮喪ge 時候」.
「我話D ye 係好容易諗ge 」.
「點樣諗呢」.
「調轉黎諗就得架啦」.
「調返轉 gwo 個唔係你」.
「gwo 個係你個仔」.
「咁你想佢點呀」.
「咁成件事就豁然開朗」.
「唔知你感唔感受到我點諗」.
「即係而家我仔話比我聽佢有病」.
「咁我想佢點」.
「咁我都咁」.
「咁大致就OK架啦」.
「即係我知道前面跟住應該點行落去」.
「OK」.
「我地可以冇有險決」.
「可以豐足」.
「就其實我想一個見證架」.
「我以前教會裡面有位姊妹架」.
「佢呢就又係cancer」.
「又係D 仔好細」.
「仔都好細ge 時候cancer」.
「咁呢就醫員又犯法」.
「醫員又犯法」.
「breast cancer」.
「咁但係佢開頭最開頭」.
「佢呢第一次患癌ge 時候」.
「就個仔女都係讀緊中學」.
「咁佢同我講呀」.
「係我ge 目標呢」.
「我ge 目標就希望呢」.
「就希望睇到佢地中學畢業啦」.
「因為佢係好差ge 」.
「成日都犯法ge je 」.
「好快又犯法」.
「但係我去gwo 間教會牧會ge 時候呢」.

$^{841}$「佢仔女都大學畢業jor 架喇喎真係」.
「超額完成架喇喎」.
「跟住呢」.
「咁我係問你而家想點呀真係有咩呀」.
「最好就見到佢地咩啦即係gwo D 啦真係」.
「不過佢就同我講」.
「呢D ye 無恥勁ge je 」.
「天父帶我行得幾遠就幾遠啦真係」.
「咁呢就後來呢就佢」.
「但係呢位姐妹厲害ge 地方係咩呢」.
「佢係好鍾意茶經ge 」.
「佢次次D 茶經都翻足架真係」.
「每一次佢翻ge 時候呢」.
「其實大部分都翻到」.
「同埋佢好特別」.
「雖然我地都有D 醫護人員係係教會」.
「佢話佢D file好差架喎真係」.
「但唔知點解個樣又唔差架喎佢」.
「即個樣即係面色又好紅潤架喎真係」.
「好靚個下架喎真係」.
「總之佢又翻到黎」.
「不過大家都知佢病得好緊要架je 」.
「又翻發牙咁樣D 架真係」.
「咁呢咁但係佢次次翻到黎茶經ge 時候」.
「我地身邊ge 弟兄姊妹會話」.
「見到佢咁準時翻」.
「我地唔敢遲到激勵埋我地去到想去ge 地方」.
「做到想做ge 事情」.
「咁但係最後尾總有一日會翻天家架咁」.
「所以我呢都為佢主持jor 安息禮拜」.
「我本來想拎個本比你養下我地姐妹」.
「咁佢最後尾臨走ge 時候」.
「佢仔已經考jor accountantge 牌架」.
「個女呢就去jor 澳洲training做pilot」.
「咁所以神的恩典係佢生命ge 當中毫無欠缺」.
「我地都鼓勵弟兄同樣」.
「而家面對咁大ge 挑戰」.
「我地生活如常可以毫無欠缺」.
「我地展現住靠著那加給我地力量ge 凡事都能做」.
「有時好多ge 苦難困難面對你ge 挑戰」.

$^{881}$「唔知你有冇覺得唔知點搞」.
「今日牧師兩個提醒你諗o下」.
「如果係你個仔咁你想點」.
「你想佢點」.
「第二 信靠信服」.
「因為主一定有安排預備」.
「親愛天父 求你保守每一個弟兄姊妹」.
「我地按照你ge 心意 信靠信服」.
「縱然世界紛擾」.
「但我地向著你 指引我地飆桿直跑」.
「得你ge 掌上 得你ge 認許」.
祈禱教徒奉耶穌名靠.
阿們.
\newpage



\section{約書亞記 1:1-9}
\label{sec:032161}
\textbf{給所有母親的寄語:「剛強壯膽」}~~~~馮永和博士
\newline
\newline
連結: \href{http://www.church.com.hk/acms/content.asp?site=cdc\&op=show\&type=product\&code=032161\&layout=sermon}{\texttt{032161}} ~~~~ 語音日期: 2026-05-10
\newline
\newline
\hyperref[sec:032160]{\small{< < < PREV SERMON < < <}}
~
\hyperref[sec:index_chronic]{\small{[返順時目]}}
~
\hyperref[sec:032162]{\small{> > > NEXT SERMON > > >}}
\newline
\newline
約書亞記 1:1-9
\newline
\begin{longtable}{cl}
\hline
\hline
章節 & 經文 (和合本修訂版)\\
\hline
1:1 & \begin{tabularx}{0.7\textwidth}{X} 耶和華的僕人摩西死了以後，耶和華對摩西的助手嫩的兒子約書亞說： \end{tabularx} \\ \\ \relax
1:2 & \begin{tabularx}{0.7\textwidth}{X} 「我的僕人摩西死了。現在你要起來，和眾百姓過這約旦河，往我所要賜給以色列人的地去。 \end{tabularx} \\ \\ \relax
1:3 & \begin{tabularx}{0.7\textwidth}{X} 凡你們腳掌所踏之地，我都照我所應許摩西的話賜給你們了。 \end{tabularx} \\ \\ \relax
1:4 & \begin{tabularx}{0.7\textwidth}{X} 從曠野和這黎巴嫩，直到大河，就是幼發拉底河，赫人的全地，又到大海日落的方向，都要作你們的疆土。 \end{tabularx} \\ \\ \relax
1:5 & \begin{tabularx}{0.7\textwidth}{X} 你一生的日子，必無人能在你面前站立得住。我怎樣與摩西同在，也必照樣與你同在；我必不撇下你，也不丟棄你。 \end{tabularx} \\ \\ \relax
1:6 & \begin{tabularx}{0.7\textwidth}{X} 你當剛強壯膽，因為你必使這百姓承受那地為業，就是我向他們列祖起誓要給他們的地。 \end{tabularx} \\ \\ \relax
1:7 & \begin{tabularx}{0.7\textwidth}{X} 只要剛強，大大壯膽，謹守遵行我僕人摩西所吩咐你的一切律法，不可偏離左右，使你無論往哪裡去，都可以順利。 \end{tabularx} \\ \\ \relax
1:8 & \begin{tabularx}{0.7\textwidth}{X} 這律法書不可離開你的口，總要晝夜思想，好使你謹守遵行這書上所寫的一切話。如此，你的道路就可以亨通，凡事順利。 \end{tabularx} \\ \\ \relax
1:9 & \begin{tabularx}{0.7\textwidth}{X} 我豈沒有吩咐你嗎？你當剛強壯膽，不要懼怕，也不要驚惶，因為你無論往哪裡去，耶和華你的神必與你同在。」 \end{tabularx} \\ \\
[1ex]
\hline
\hline
\end{longtable}
$^{1}$當我們想到母親節的時候.
還有母親節的大餐.
因為她真的在我們的生命裡負上太多.
很多的愛護和支持.
所以慶祝母親節是很應該的.
因為做母親其實是十分辛苦的一件事.
回看我太太的時候.
我發覺我們大家都關心我們的女兒.
但是她比我有多一個母親的心.
就是女兒無論有什麼頭雲生慶 考試.
就算不計分那一隻她都會緊張.
所以其實這就是神賜給母親的一個意念.
就是愛護她的子女.
其實不是容易做的.
因為通常我們說女人不容易做.
其實做母親更加困難.
為什麼呢 因為其實母親一定要同時兼顧.
很多時候都是工作和家庭.
如果有一些親子關係有緊張的.
或者要照顧一些有特殊需要的子女.
這樣的情況就會更加困難和複雜.
所以其實一個不容易做的人生.
但是又很漫長.
好像打了一場仗又打另一場仗.
我看到我太太一年級 二年級.
我女兒三年級.
好像每年都在打仗.
接著我就覺得人生真是漫長.
所以我想用一段聖經去鼓勵.
我太太也鼓勵你們.
就是約書客記第一章的一至九節.
在漫長的人生裡面的時候.
我們會經歷生命上每一個的高低.
每一個的起伏.
有得有失 有成有敗.
我們怎樣去過一個.
可以說是有意義和德性的人生呢.
其實在摩西死了之後.
帶領以色列人重返迦南地這個重任.
其實就落到去.

$^{41}$曾經是摩西的弟子或者助手.
現在被神揀選作為接班人的約書客身上.
但如果我們很細緻地看聖經.
特別是《申命記》的時候.
你就會發覺.
這是一個不可能的任務.
為什麼呢.
因為當你看《申命記》第七章第一節的時候.
當中就會說到.
耶和華你神領你進入.
要德為業之地.
從你面前講出很多的國民.
就是什麼呢.
是客人 甲家殺人.
阿摩尼人 迦南人.
比利洗人.
希美人 耶布斯人.
共七國的民都比你強大.
我們要注意聖經說什麼.
它不是說七國的人加在一起的時候.
才比你強大.
而是每一個民族他要面對的.
都比他更加強大.
我們可以想像是.
一個好像三年級的小朋友 我女兒那樣.
會有一群中四中五的大朋友.
在裡面跟他打架.
那你想到結果會是怎樣的.
所以其實這個差不多是不可能的任務.
接著再看看《申命記》第九章第一節二節.
它當中說到以色列你當聽.
你今天要過約旦河.
去趕出比利強大的國民.
要得著廣大堅固高得頂天的承邑.
當中原文說到.
Cities great and fortified up to heaven.
那文是阿納族人來的.
又高又大.
People great and tall.
是你所知道的.

$^{81}$也曾經有人指著他們說.
誰能夠在阿納族人面前站立得住呢.
敵人又多又強大.
城池又高又堅固.
易守難攻.
其中的居民以善戰出名.
誰能夠在阿納族人面前站立得住.
絕非當時只慣於.
目速的以色列人所能夠對付.
所以其實約書亞在當中面對什麼呢.
面對一個不可能的任務.
他其實都很知道.
如果我們要打贏一場戰爭.
他可能都會這樣想.
如果我們要攻陷一個城池.
即是打贏一場戰爭的時候.
在天時地利人和.
都配合的環境.
又加上一些運氣.
可能會做得到.
但現在不是.
你要面對無數的城池.
很多的戰爭.
你要長長都要得勝.
否則你的民族就會有生存的危機.
所以其實為什麼神在約書亞記第一章九節的時候.
重複了三次說.
你不要懼怕.
其實約書亞在做什麼呢.
他很害怕.
他是懼怕.
為什麼呢.
是沒有勝算的.
是一個不可能的任務.
親愛的弟兄姊妹.
親愛的母親.
很多時候我們在我們的生命裡面.
是面對一個沒有勝算的人生.
差不多是不可能的任務.
很多時候我們發覺.

$^{121}$一件事就是說.
我的子女很困難.
從幼稚園的時候已經好像打仗一樣.
同學慢慢逼他長大.
現在才三年級而已.
什麼時候到中學.
接著你看看.
現在才中學什麼時候到大學呢.
接著大學出來了.
做事找工作可能有問題.
怎樣可以看著他過.
已經年紀大了.
其實在江中裡面的時候.
我們都會想到這些事情.
究竟你的人生有沒有勝算呢.
如果你覺得你沒有把握.
你在母親節的歡樂下.
忍一忍你覺得恐懼的時候.
那這篇講道就適合你.
我們發覺一個很奇特的事.
就是當我們看這篇訊息的時候.
我們發覺一件事就是說.
約書亞是面對很困難的事.
但是神從來沒有說過.
那些誠邑有多高.
還有那些敵人有多強大.
在一章一節九節裡面.
他差不多沒有說任何話.
為什麼呢.
是因為這些困難和這些不行.
在神的眼裡面不是最重要的問題.
在約書亞來看.
在一個沒有勝算的人生裡面的時候.
神臨到他很簡單的說了.
你只要去遵行兩個條件.
你的人生就會得勝.
你就會不斷地得勝下去.
你就會不停地打成長.
兩個條件.
但是當我很細緻地看這兩個條件的時候.

$^{161}$我就發覺.
這兩個條件在我們今天同樣的適用.
我們現在看看這兩個條件究竟是什麼條件.
第一個條件是記載在第一段聖經.
這段聖經其實分為兩段.
第一段是二至五節.
接著是六至九節.
第一個得勝的條件在第一段聖經第二至五節裡面.
當中講到.
第一件事神要做的就是你要起來過河.
祂說我的僕人摩西死了.
所以現在你要起來過河.
在希臘文七十四譯本的原文聖經裡面.
過河是一個主要的動詞.
但是在文法上你會發覺一件事.
如果你誠具看的時候.
是一個條件你要先達到的時候.
你才能過河的.
那個條件就是你要起來.
就是說過河只能夠發生在起來之後.
如果你沒有起來的時候.
你基本上就不能夠過河.
接著聖經裡面就繼續講.
在當中裡面祂說.
起來是什麼時候神要我們起來呢.
是現在.
這樣在哪裡起來呢.
其實根據上文後義講的時候.
最明顯的.
其實有兩重的意思.
最明顯的意思是.
要約書亞在失去了神偉大僕人摩西的哀痛裡面起來.
約書亞可以說是摩西親密無間的弟子.
是他悉心栽培的接班人.
他們關係很密切.
甚至你說好像父子都不為過.
失去摩西對約書亞的打擊一定是很大.
就好像失去一個親人一樣.
在我們的生命裡面的時候.
我們會經歷損失.

$^{201}$我們可能會損失一個對我們很重要的人物.
一段很重要的感情.
一個很重要的金錢的數目.
我們通常經歷損失的時候.
我們覺得很多時候我們未必走得下去.
我們不能夠再起來.
我們那種隱隱然然的傷痛常常在我們裡面.
因為那個人對我太重要.
那件事對我太重要.
沒有了他我什麼都做不到.
這個失敗的經歷對我太過重要.
他使我留下陰影.
讓我能夠很難再向前走.
當我們經歷損失的時候.
我們會消沉.
我們會很不開心.
但會發覺一件事.
當你慢慢回過神來的時候.
你就會發覺生命是仍然繼續下去的.
世界仍然是運轉的.
人生仍然都要向前努力的向前走.
所以當約書亞哀悼的時候.
就好像《申命記》第34章第8節裡面說.
摩西死了之後.
以色列人為摩西哀哭了三十天.
但當他哀悼之後.
聖經裡面怎麼說呢.
約書亞第一章第二節裡面就說到.
神的說話是臨到約書亞.
你要起來.
你過河.
因為當你過河.
當你起來的時候.
你就會揭開你人生新的一頁.
這個新的一頁.
是你沒有想像過的.
所以你一定要起來.
無論在多困難的情況下.
你靠著我給你的能力去起來.
如果你這樣做的時候.

$^{241}$我就會與你同在.
在你平生的日子.
就沒有一個敵人會在你面前站立得住.
你要起來.
除了在摩西的損失.
損失了摩西之下起來之後.
其實還有一個最明顯的地方.
我們要在那裡起來.
就是懼怕.
在約書亞裡面的時候.
他一定很懼怕.
所以神在第六節到第九節裡面.
三次說你要剛強壯膽.
你不要懼怕.
為什麼他會懼怕呢.
其實我相信最重要的原因就是說.
他有一個認識.
就是除了他的環境.
他覺得是不可以得勝的環境.
在一個沒有勝算的人生之外.
其實他隱隱然然一定會想到一件事.
摩西是無可取代的.
摩西是我的師父.
有些困難只能夠摩西先去解決.
我不是摩西.
我從來沒有使紅海的水分開.
我從來沒有面對面見過神.
我從來沒有吩咐十災臨到埃及地.
我怎麼能夠完成.
連摩西都完成不了的事情呢.
我怎麼帶領以色列人進入迦南地呢.
我不行的.
很多時候我們都有類似的迷思.
在很多母親而言.
當她面對工作和事奉.
和子女家庭的時候.
在當中的時候.
很多時候第一個感覺就是.
阿神我已經很忙.
家庭和子女已經吃光我所有的時間.

$^{281}$如果我子女和我親子關係有問題的時候.
耗盡我一切的心力.
你還要我做這樣做那樣.
在城裡面有那麼多要求和條件.
我真的不行.
你找別人吧.
是不是我們覺得很多時候我們不行.
是不是在當中覺得.
這些事情只有別人才能做到.
我做不了.
但是神在這裡怎麼說呢.
很奇怪的.
神當祂的說話臨到約書亞的時候.
其實祂的意思是.
你起來.
你由你的傷痛和損失裡面起來.
你從你覺得你不行的情況下.
你起來.
當你這樣做的時候.
你不單止會揭開你人生新的一頁.
你會做一些連摩西都做不到的事.
當約書亞起來的時候.
他過河的時候.
他就見到神的榮耀.
當他來到約旦河的河水的時候.
他就會發覺一件事.
就是河西好像摩西一樣.
紅海怎麼在摩西兩邊分開.
今天約旦河的河水.
在他的面前兩邊分開.
接著他就看到了神的榮耀.
但是神的能力遠遠沒有這樣簡單.
約書亞的巔峰.
他的宿命巔峰是記載在約書亞記第十章裡面.
就是當基片投降了.
那時候耶路撒冷王和很多阿摩利王開始很害怕.
他們想聯合起來對抗以色列人的時候.
他們就組織了一個很強大的軍隊.
然後神就叫約書亞起來.
有一次的起來去面對那個軍隊.

$^{321}$約書亞擊打那個軍隊的時候.
他就做了一個禱告.
他就說太陽你要停在基片 月亮你要停在阿倫谷.
直到我向仇敵報仇的日子.
接著很奇妙地.
聖經的最後面就說.
太陽停在基片 月亮停在阿倫谷.
若有一日之久沒有急速落下.
直到以色列人和耶和華向仇敵報仇的時候.
滿足之後.
是摩西都沒有做過的事.
摩西都做不到的事.
所有的事情是豈止是一句說話.
起來.
你是不是覺得很多的事牽制著你.
很多時候我們的生命 實在在當中裡面.
我們太多的顧慮.
這裡有事 那裡有事.
家裡其實已經很困難.
工作已經很多事要做.
神拿了很多事未完成.
我真的做不到你要我做的事.
我不行.
我被太多的罪菌綁著.
在以前的日子裡面的時候.
有些傷痕.
我一直現在都覺得痛.
我一直都再起不到身.
我在depression裡面的時候.
我不能夠起來一生.
我真的不行.
你找個別人.
是不是你覺得不行呢.
當約書亞覺得自己不行的時候.
當約書亞覺得很懼怕的時候.
當約書亞面對一個沒有勝算的人生的時候.
當約書亞面對很大的困難的時候.
神沒有叫他去看著困難.
神沒有叫他去看著一個沒有勝算的人生.
神說在你沒有勝算的人生裡面.

$^{361}$在你很大的困難的日子裡面.
你要做一件事.
你要專心定精仰望我.
不要看這些困難.
不要看著這些沒有勝算的人生.
現在起來.
過河.
而當你肯憑著我的能力去這樣做的時候.
我就會與你同在.
因為我的同在.
你就會得勝.
你就會做到一些.
連摩西都不會做得到的事.
這個就是你真正的人生.
親愛的弟兄姊妹.
在這種十年的經歷.
不單單只是以前的約書亞經歷過.
在中國內地會的創辦人戴德生先生.
即是今日的海外基督使團OMF.
當他未蒙召出發到中國宣教的時候.
其實他已經有個心很想去那裡宣教.
所以有一次跟一個牧師談到他神呼召.
他去宣教的事.
這個牧師就問他.
你打算怎樣預備他呢.
接著戴德生就說.
我完全一點都不知道.
我想好像那時候的十二使徒和七十個門徒.
往猶太地傳訪的方法一樣.
我們的腰帶不要帶金錢.
行路不帶口袋.
只單單憑信心神會供應我一切所要.
那時候這個牧師就很慈愛地.
把他的手放在戴德生的肩頭上.
他就說年輕人.
當你年紀大一點的時候.
你就會變得聰明一點.
這種想法只有耶穌基督的時代才行得通.
在現在的世界是不行的.
即是牧師這樣說.

$^{401}$但是戴德生在以後的日子裡面.
他憑信心開展了內地會的事供.
他從來沒有跟別人募捐.
他只是單單仰賴神的供應.
在他去了中國傳福音大約二十年之後.
去到1900年代的時候.
內地會的同工.
在中國宣教的人數是超過了一千人.
在很多年之後.
他再說起這件事的時候.
他就說我現在的年紀比當時大得多.
但不見得比當時更加有智慧.
我是越來越深信一件事.
如果我們是照著神給門徒的指示.
和應許去做的時候.
我們今天在這個世界上.
都是行得通的.
親愛的朋友.
親愛的弟兄姊妹.
和親愛的母親.
你是不是覺得不行呢.
在約書亞的時期.
幾千年前.
神跟約書亞說是行得通的.
在一百年前當戴德生.
過去海外宣教的時候.
神再說是行得通的.
在今天.
神的眼目仍然擦遍全地.
要顯大能去幫助一些向他心存誠實的人.
你是不是覺得不行呢.
你是不是覺得不可能的任務呢.
你是不是覺得沒有勝算的人生呢.
你是不是覺得我在極大的困難裡面.
已經沒有可能得勝呢.
但在當中神說.
起來.
你試一下起來.
當你肯憑著信心起來的時候.
當你按著我吩咐起來的時候.

$^{441}$當你過河的時候.
你就會見到我的榮耀.
這是神給我們第一個得勝的條件.
第二個得勝的條件記載在哪裡呢.
是記載在第二段聖經.
在第六節至第九節.
當中裡面好像很簡單.
是謹守遵行神的命令.
第六節至第九節這段聖經好像很簡單.
但其實作者是運用了古代近東.
一個很常用的寫作手法去寫這段經文.
如果不了解這種手法的時候.
我們其實很難掌握這段經文的重點所在.
當時的歷史學家通常是用一個漸進式的寫作手法.
即是說他記載一件事的時候.
特別是重要的事情的時候.
他會重複寫幾次.
首先第一次的時候會簡略去描述整件事發生是甚麼事.
接著他會再描述同一件事件一次兩次.
而每一次重複的時候.
都會簡化之前所講的東西.
接著再加上額外的資料.
使每個讀者看的時候.
會有對事情更加深入的了解.
和明白更加多的真相.
而通常額外加上的部分就是重點所在.
所以這種手法其實在摩西五經裡面.
特別是在創世記裡面是常常引用到的.
而約瑟亞記的作者在寫作第一章第一至九節的時候.
亦都採取了一個類似的手法.
我們可不可以去看下一張片.
我想是之前那兩張是不是.
不是 是啦 再去多一張.
沒有錯 謝謝.
我們會發覺一件事就是.
其實在一至九節裡面的時候.
作者是重複了兩次神的命令和神的應許.
是同一個命令 同一個應許.
那個命令就是你要起來過河.
去得到迦南地做你的產業.

$^{481}$然後那個應許就說.
如果你敢做的時候 我就和你同在.
你會發覺這個作者是重複描述了這兩次.
第一次的時候.
所以你會看到第一次是一至五節.
第二次是六至九節.
所以你會發覺第二次的描述.
第六節是對比到第二至四節.
而第九節是對比到第五節的.
這樣的時候其實新增的部分在哪裡呢.
就是第七至第八節.
第一次的描述是沒有說的.
也是這次重點所在.
所以這也是神給我們的條件.
要求我們要有的條件.
是第二個條件.
你怎樣可以得勝的人生.
那個條件是甚麼呢.
是謹守進行神一切的律法.
不可以偏離左右.
是第七節所說的.
我們下去下一章.
而第八節就很清楚的去解釋到.
怎樣可以做到這件事呢.
第一 律法書不可以離開你的口.
第二 是要晝夜思想.
就是要不停地熟讀.
謹記和思想神的說話.
才可以進行.
我們一定要明白.
當時以色列人其實是怎樣看聖經的時候.
我們才會明白到這個聖經真正的意思.
因為近時的方法和現在大不相同.
近幾千年前在約桑那時代的時候.
所有的經文是寫在羊皮卷上面.
那時候蒲草寺應該還沒發明.
羊皮卷上面.
所以其實不是很多人可以有一本書.
因為有一本書是很貴的.
你要殺了一隻羊.

$^{521}$然後拿一層皮.
然後再用很小的字寫在羊皮卷上面.
如果一篇寫不完就要殺多一隻.
所以不是很多人做得到這件事.
所以很多人是沒有書的.
通常一個民族或者一個群體.
通常只有一本書.
而那個拉比在安息日或是先帝國的時候.
就會讀出來.
而當他讀出來的時候.
他會期望什麼呢.
他會期望我讀一句的時候.
會眾會跟著我去讀那句.
讀到大家熟悉為止.
然後你就會記在心裡面.
所以其實這裡說的是什麼呢.
其實這律法書不可離開你的口.
是你要上上朗讀.
在沒有書卷的情況下.
將它記在心裡面的意思.
而當你這樣做到之後.
你就能夠晝夜去思想.
因為你已經有聖經在裡面.
我記得當我1985年回香港的時候.
我那時候是first degree畢業.
在加拿大的溫莎回香港.
那時候其實我覺得很孤單.
因為我不是在香港信主.
也沒有教會.
回到香港的時候.
沒有人告訴我要回哪間教會.
也沒有什麼支持.
所以那時候很自然的時候.
就希望找回之前回到香港的弟兄姊妹支持.
所以我走去那時候找我的中學老師.
約他出來吃飯.
我中學老師那時候出來吃飯的時候.
他都化了濃妝出來.
穿了一套suit.
我心裡面想.

$^{561}$要不要啊 現在不去上班了.
現在我和他吃飯了.
好了 接著大家寒暄完之後.
接著隨口說.
問他 現在你要回哪間教會啊.
接著他就看著我一陣子.
他問 什麼叫教香港可以的嗎.
我沒有回教會很久了.
接著我看著他.
那時候我20歲出頭.
沒有什麼熟齡經歷.
看著他.
他說我小學老師來的.
接著我說.
哇 他回來香港早幾個月而已.
這麼快就消失了.
接著我心裡面就想.
我絕對不會像他那樣.
那時候我做了一件事.
就是我開始上班之前.
我起了一個小時.
在當中讀經圖稿.
那時候很多時候.
我是意猶未盡的.
一個小時不夠.
所以那時候我做什麼呢.
那時候是沒有手提電話的情況之下.
我是用一些卡片.
抄下我覺得意猶未盡的經文.
放在卡片上.
當我搭巴士去地鐵站的時候.
當我搭地鐵的時候.
我就在巴士和地鐵站的時候.
拿張卡片出來看.
看 不停地看 不停地看.
慢慢的時候.
我發覺一件事.
就是我沒有刻意去這樣做.
但是那些經文會記在我裡面.
其實就是跟他現在做的一樣.

$^{601}$我很老實告訴你.
我接著的時候.
經歷很多的困難人生.
很多時候是那些經文.
pop up去幫助我的.
我們再看看這段細緻的經文的時候.
你就會發覺第六至九節裡面.
分開三段的訊息.
就好像剛才那個表那樣.
但是每一段的時候.
他都用了一個剛強壯膽開頭.
但是第二段就不是.
第二段是加重了語氣.
用了你只要剛強.
要大大的壯膽.
接著的時候.
真的很明顯.
明眼人都知道.
就是說這是一個重點.
這是重點所在.
所以我再看的時候.
我又發覺.
先前的時候我不是很明白.
後來才知道.
就是中間那段.
不過是要你勤讀聖經記住它而已.
為什麼要剛強壯膽呢.
看聖經的時候會很害怕嗎.
為什麼要剛強壯膽呢.
之前看的時候不是很明白.
但是當人生的歷練多了的時候.
就開始明白.
是.
當你要謹守遵行神的說話的時候.
你要勇敢.
在我仍然做會計師的日子裡面的時候.
會計師這一行.
其實有一個不成文的規定.
其實這是一個.
很多人都會做.

$^{641}$但是不會報告給上面聽.
我那時候在一間大的.
國際的會計公司做.
那時候我的學生.
我們很忙.
通常的時候.
會計一年做連結的時候.
就會通宵達旦地做.
我曾經試過48小時沒睡覺.
就在那裡做.
在客戶辦公室做.
所以你知道我可以多忙.
那麼.
壓力又大又忙的時候.
通常很多人都會做一件事.
就是.
我師兄曾經叫我.
你不用做那些東西.
你說做了就可以了.
你只要剔除它就可以了.
就是說.
我們本來要做50個樣本.
用不著做50個.
你做了第一個.
證明它還在的系統.
接著你剔除其他40個就可以了.
那你變成你會做快很多.
大家都不用熬夜了.
其實基本上.
是很多時候.
我都被人要求做一件事.
我也很老實告訴你.
因為那時候.
我常常將神的說話.
放在我的口袋裡.
接著拿出來看.
使得我有勇氣去玩弄.
我告訴你一個很殘酷的事實.
就是在香港這個社會裡面.
或者在世界裡面.

$^{681}$你做所謂正義的行動.
是不會有人欣賞你的.
還有你的表現報告寫得會很差.
我是不止一次被人.
你想一下.
老闆叫你做.
你說NO.
他可以怎樣都寫好你.
所以其實我是不止一次.
拿過一些很差的表現報告.
但是親愛的弟兄姊妹.
當你起來.
當你過河的時候.
你就會看到神的榮耀.
從來沒有一份表現報告.
影響到我.
我在那方面是年年升職.
直到我離開的時候.
從來沒有影響過我.
反而是寫我壞的那些學生.
我發覺後來就升不了.
所以這不是有直接關係的.
這是一個事實.
我很老實告訴你.
你是不是覺得不可能.
你是不是覺得很難.
如果你不把神的說話放在心裡.
你是不會有這樣的勇氣去做.
如果你不禱告神去幫助你的時候.
你不感覺到神跟你同在同去的時候.
你是不會做得到.
是神加能賜力給你的時候.
你才能夠可以在每一次.
你要選擇神還是世界當中.
你去選擇神.
所以你是需要勇氣.
是需要剛強壯膽.
而且不是剛強壯膽.
是只要剛強大大壯膽.
所以其實這段聖經裡面的時候.

$^{721}$整體來說神在說些什麼呢.
其實神說了一個訊息.
祂面對一個很害怕的若書雅.
祂面對一個若書雅.
這個若書雅看到她的生命裡面沒有勝算.
她基本上失敗的頻率.
可以去到100\%.
如果照人家看的時候.
沒有可能得勝.
在這樣的情況之下怎樣走下去呢.
在不行的情況之下怎樣走下去呢.
在失去了重心人物之後.
這個民族怎樣走下去呢.
神就在當中裡面說.
你成為這個民族的重心人物.
你怎樣成為呢.
不是因為你很有能力.
不是因為你很有智慧.
不是因為你很有社會的地位.
不是因為你很有錢.
而是今天我要你專心仰望我.
不是看著你的困難.
不是看著你有沒有勝算.
不是看著你行不行.
如果你在這樣的情況之下.
你仍然肯起來的時候.
我就能夠使你過河.
如果你肯過河的時候.
你就會看到我的榮耀.
你就會揭開你生命新的一頁.
你就會做到.
一些連摩世都做不到的事.
親愛的弟兄姊妹.
在幾千年前.
當神和約瑟亞說這段說話的時候.
在一百年前.
神和戴特生重複說同一樣的說話.
親愛的弟兄姊妹.
親愛的母親.
今天神仍然說話.

$^{761}$祂的說話仍然在空氣裡面.
挑戰著每一個人.
當你靜下來的時候.
你就會聽到神的說話.
神的說話沒有改變過.
因為神是不改變的.
今天祂的要求仍然是.
起來.
和.
當你這樣做的時候.
你就會揭開你人生新的一頁.
不在乎你幾多歲.
不在乎你經歷著什麼困難.
不在乎你覺得自己行不行.
你就會看到神的榮耀.
你就會做到一些連摩西都做不到的事.
在大約一百年前.
有一個母親帶著她很小的兒子.
只有幾歲大.
去欣賞帕雷德夫斯基的演奏會.
一個被當時認為是蕭邦之後.
波蘭最偉大的鋼琴家.
那母親想帶著兒子去看.
因為她想鼓勵兒子在學琴的造詣上進步.
當他們到達會場入座的時候.
母親就看到人群裡面有一個相熟的朋友.
跟著她就走過去跟他寒暄幾句.
但母親所不知道的是.
她才期望她兒子坐在那個位置等她.
但她不知道的是.
這個年紀的男孩最喜歡探險.
所以她趁她媽媽看不到的時候就覺得撿到了.
所以她就開始到處去探索.
後來她就看到一道門寫著.
No admittance 不准進入.
然後她就打開進去了.
當劇院的燈光開始轉暗.
音樂會即將開始的時候.
母親回到座位的時候才發現兒子不見了.
然後她開始到處去問人.

$^{801}$見不見我兒子.
開始不知所措的時候.
突然她發覺台上的帳幕向左右分向.
接著燈光集中帶在台上要演奏的.
很出名很名貴的鋼琴上.
接著她媽媽就發現.
她兒子是坐在那個琴椅上.
接著她驚恐的時候.
她發覺她兒子在彈她最喜歡的那首.
Twinkle Twinkle Little Star.
接著她媽媽第一個想法就是.
糟糕了 這次要大大出醜.
沒有人知道那個是我兒子.
當她這樣想的時候.
帕雷德夫斯基就入場.
他很快走到鋼琴的旁邊然後彎下身.
他在小朋友的耳邊說了一句話.
Don't quit, keep playing.
不要放棄繼續彈下去.
接著這個偉大的鋼琴家就伸出他的左手.
在低音的部分開始彈奏低音的配樂.
為Twinkle Twinkle Little Star彈奏配樂.
接著再將他的右手繞過小朋友.
在高音的領域上開始彈奏高音的配樂.
很快這首本來很平庸.
一首很簡單平平無奇的.
Twinkle Twinkle Little Star.
就變成了那些人從來沒有聽過這麼好聽.
一個很完美的段落.
這是一個從來沒有想像過的意外的場景.
那時候就震撼了在場每一個觀眾.
所以這件事是記入了歷史裡面.
親愛的弟兄姊妹.
親愛的朋友.
親愛的母親.
你是不是覺得自己不行?.
你是不是覺得在當中你最多只能夠做到.
就是彈Twinkle Twinkle Little Star?.
你是不是覺得.
在你的人生裡面是一個沒有勝算的人生?.

$^{841}$你是不是覺得在很大的困難之下.
在母親節日的時候.
你竟然要咬緊牙根去吃完那頓飯然後走下去?.
不要這樣想.
在當中裡面.
如果你將你的生命再一次交給神的時候.
神就可以使到你的Twinkle Twinkle Little Star.
成為一首從來沒有人聽過最完美的樂章.
在你的生命裡面.
在你經歷的人生當中.
在你所經歷的高低起伏當中.
每一個的得和失.
每一個的勝和敗.
祂就在那裡跟你同在同行.
只要你肯的時候.
祂就在那裡跟你同在同行.
祈求我們能夠因為你的緣故.
看到你的榮耀.
主啊,幫助我們成為.
在你今這個世代裡面的見證人.
使用我們每一個的人.
將再願意獻上生命給你的.
祝福他們.
使用他們.
帶領他們經過人生的高低起伏.
聽我們的禱告.
奉耶穌的名.
阿們.
\newpage



\section{詩篇 62:1-12}
\label{sec:032162}
\textbf{實在靠得住}~~~~陳永就牧師
\newline
\newline
連結: \href{http://www.church.com.hk/acms/content.asp?site=cdc\&op=show\&type=product\&code=032162\&layout=sermon}{\texttt{032162}} ~~~~ 語音日期: 2026-05-17
\newline
\newline
\hyperref[sec:032161]{\small{< < < PREV SERMON < < <}}
~
\hyperref[sec:index_chronic]{\small{[返順時目]}}
~
\hyperref[sec:032163]{\small{> > > NEXT SERMON > > >}}
\newline
\newline
詩篇 62:1-12
\newline
\begin{longtable}{cl}
\hline
\hline
章節 & 經文 (和合本修訂版)\\
\hline
62:1 & \begin{tabularx}{0.7\textwidth}{X} 我的心默默無聲，專等候神，我的救恩從他而來。 \end{tabularx} \\ \\ \relax
62:2 & \begin{tabularx}{0.7\textwidth}{X} 惟獨他是我的磐石，我的拯救；他是我的庇護所，我必不大大動搖。 \end{tabularx} \\ \\ \relax
62:3 & \begin{tabularx}{0.7\textwidth}{X} 你們大家攻擊一人，使他被殺，如歪斜的牆、將倒的壁，要到幾時呢？ \end{tabularx} \\ \\ \relax
62:4 & \begin{tabularx}{0.7\textwidth}{X} 他們彼此商議，要把他從高位上拉下來；他們喜愛謊話，口雖祝福，心卻詛咒。（細拉） \end{tabularx} \\ \\ \relax
62:5 & \begin{tabularx}{0.7\textwidth}{X} 我的心哪，你當默默無聲，專等候神，因為我的盼望是從他而來。 \end{tabularx} \\ \\ \relax
62:6 & \begin{tabularx}{0.7\textwidth}{X} 惟獨他是我的磐石，我的拯救；他是我的庇護所，我必不動搖。 \end{tabularx} \\ \\ \relax
62:7 & \begin{tabularx}{0.7\textwidth}{X} 我的拯救、我的榮耀都在於神；我力量的磐石、我的避難所都在於神。 \end{tabularx} \\ \\ \relax
62:8 & \begin{tabularx}{0.7\textwidth}{X} 百姓啊，要時時倚靠他，在他面前傾心吐意；神是我們的避難所。（細拉） \end{tabularx} \\ \\ \relax
62:9 & \begin{tabularx}{0.7\textwidth}{X} 人真是虛空，人真是虛假；放在天平裡就必浮起，他們一共比空氣還輕。 \end{tabularx} \\ \\ \relax
62:10 & \begin{tabularx}{0.7\textwidth}{X} 不要仗勢欺人，也不要因搶奪而驕傲；若財寶加增，不要放在心上。 \end{tabularx} \\ \\ \relax
62:11 & \begin{tabularx}{0.7\textwidth}{X} 神說了一次、兩次，我都聽見了，就是能力屬乎神。 \end{tabularx} \\ \\ \relax
62:12 & \begin{tabularx}{0.7\textwidth}{X} 主啊，慈愛也是屬乎你，因為你照著各人所做的報應他。 \end{tabularx} \\ \\
[1ex]
\hline
\hline
\end{longtable}
$^{1}$文安!弟兄姊妹!.
好開心呀,你重複講道就第一次.
通常來親都按目禮.
好開心呀,上次何牧師提到.
來按親都三個,幾好呀.
三三不盡.
願神興起更多更多的工人在我們當中.
再一次自我介紹,播道總會總幹事陳永舟.
亦都代表播道總會和大家來到去文安.
願上帝的恩惠和慈愛伴隨著大家.
好開心呀,來到大家當中.
我早堂又見到一個.
之前認識的弟兄姊妹.
是不是叫中堂呀,剛才那個堂.
不知是甚麼堂.
又見到一個,這堂稍後完了會否又見到.
好開心呀,好親切.
主內的一家.
提到親切呢.
弟兄姊妹在總會的服侍.
都是叫我咒目的.
因為我們有很多德高望重的陳牧師.
有時開會叫陳牧師,我都不知道叫我,我還坐在這裡.
我們以知識別,有時開會開到很晚,大家都睏了.
叫下咒目,醒神就知道叫我咒目.
我真的很喜歡別人叫我咒目.
因為行章咒目就快死了.
提一下自己要傳福音,我真心話的.
我們行章咒目就快死了.
努力去傳福音.
都是重要的,傳道人有時.
以我們的職業就是.
傳福音救靈魂,讀經祈禱,這些固然是需要的,所以都要提一下.
正如剛才所說,工作就.
不是第一天.
出來就做牧師.
因為我有些.
會友有時都很天真地問,你是不是一出生就做牧師,我當然不是.
我做過很多工作,才神呼召我.
在一個社區教會服侍,在彩雲村做了一整年.

$^{41}$我算上做社工,做傳道,我差不多三十年了.
神又呼召我.
去總會服侍.
播道一家,在這個崗位裡.
去不同的播道會服侍,真是我自己很希朗的士逢德.
正如剛才所說,不是一出生就做牧師,做過很多工作.
我從前做事真的很深印象,不知道大家有沒有遇過一些同事.
比較卸膊,斤斤計較.
一叫他做事都未開口,這件事不是我跟開的.
跟同事說,我找誰做,大家怎樣配合.
他真的這樣跟我說.
找誰做事.
誰誰靠得住.
你知不知道下一句是什麼.
嘩,原來你們有這麼多豬奶在身邊.
我以為是我說的,原來你們也有豬奶在身邊,我真的不知道.
我又沒有考究,為什麼不是牛公在身邊,而是要豬奶呢.
他真的這樣說,誰靠得住,豬奶都會上樹.
這就證明那個人不可靠,對不對.
事實是真的,你待人接物.
其實你這個人可不可靠,負不負責任,真是一個很重要的生命質素.
今天我們宣講詩篇第62篇.
只能提醒我們神是可靠的那一位.
是值得我們放心投靠的那一位.
我們再做一個祈禱.
我們祈禱,願聖靈來到我們工作.
父愛願你這樣來幫助,帶領我們進入你的話語.
在道裡面我們能夠體會到上帝你自己的愛,是何等的長闊和高深.
並且讓我們能夠可以這樣實踐.
我們這樣同心的祈禱,奉耶穌基督的聖名而求.
阿們.
大衛的詩.
家級詩班長.
照耶穌頓的做法.
標題之後,經文開始.
唯獨等候神,我的心才有安息,我的拯救從祂而來,唯獨祂是我的盤石,我的拯救,祂是我的保障,我必不自動搖.
你們大家攻擊一個人,要把他殺害如同拆毀傾斜的牆.
張倒的壁.
要到幾時呢.
他們彼此相義,無要把他從尊位上拉下來.

$^{81}$他們喜愛方話.
口須祝福.
心卻就坐.
有一個小奶奶表示停一停.
我們讀聖經,我們知道詩篇是在聖經裡面最長篇幅的經卷.
150篇.
如果你揭實體的聖經,差不多你中間一揭,一定是詩篇的了,一定中的了.
亦都看到有編修的痕跡,有卷一,卷二,卷三,卷四,卷五.
遠到摩西,嘩!上古的時候.
近到可能是大衛.
我們看到是有個編修的痕跡.
而這個卷一,卷二,卷五,很明顯和五經.
摩西五經是正在對照的,我想這個不用怎麼爭辯,是猶太的傳統.
卷二,是最多大衛的詩篇.
很多時候會有個標題,就像這裡,大衛的詩.
當然大衛的詩是大衛寫的詩,大衛喜歡的詩.
還是為歌頌大衛而寫的詩呢,我們又不知道.
不過肯定是一個很密切的關係.
這裡還提到趙爺道頓,這個趙爺道頓是在歷代寺找到的,他真的是詩班長.
所以從這些蛛絲馬跡來看,這篇詩篇和大衛的關係一定非常之緊密.
甚至你代入在大衛的生平來看,又將他從尊位拉下來,人們又謀害他.
自然就令我們想起阿沙隆的叛變.
他的兒子都背叛他,眾叛親離的時候,他就逃難了.
如果你這樣代入這篇詩篇的時候.
你看的情況,你會明白到詩人的心境會很不同.
他正面對一個非常之巨大的艱難.
面對艱難的時候,他不是怨天尤人.
他是提醒自己,也提醒跟隨他的同伴,我們懷著感恩的心.
去依靠神.
我們如果讀和合本,我講到的就是用新譯本,和合本是我們很熟悉的經文.
我的心默默無聲,尊等候臣.
新譯本的譯法就將他導致了,他先放一個尊.
尊等候臣.
唯獨等候臣.
是同一個字來的.
所以有些譯本翻譯成真等候臣,真正的真.
我無要,我承言,我確實這樣的意思.
他告訴大家知道,我們可以強調的是什麼呢?.
是強調我們可以依靠上帝,而當我這樣依靠上帝的時候,我心裡就有平安了.
所以那個依靠.

$^{121}$和平安是貫穿了這一首詩篇裡面的.
待會你讀下去,你會讀到的,第1至4節.
讀到第5至8節,差不多是重複的,特別是第1,2和第5,6節.
唯獨安息貫穿了這首詩篇的主題.
接著之後,重複也出現了一些詞彙.
盤石,拯救,保障,不致動搖.
盤石就比喻恆久,堅固,可靠.
和下文可以對照的.
是張島的牆,要倒的壁,嘩,長島眾人推,推了一份.
這邊搖搖欲墜,嘩,這邊堅固穩妥,他做了一個對比,告訴你.
至於拯救,是作者親身經歷過的.
所以他才有這樣的信心,他安靜來等候,不僅如此,他也呼籲其他人學他.
相信得過上帝的拯救.
保障有避難所的意味,保護.
難免對敵人的攻擊,惡意的加害.
他接著問,要到幾時呢?.
這個要到幾時呢?.
的片語,在詩篇裡面很多時候我們也讀到的,其實也是我們祈禱的心境來的.
很多時候我們祈禱,神啊,你為我開路,某份工作,或者神啊,你為我開路,某個事情.
但等極也等不到的時候,你自己也會這樣祈禱.
要到幾時呢?等到幾時呢?.
不過你看看這裡也不是這樣,如果你看看這裡的上盲和下理.
他是被攻擊,被人加害.
詩人反而說,你們不怕上帝嗎?.
你們還不收手?.
你們還不停?要到幾時呢?.
他不是一種控訴,他反而好像是勸那些人不要再作惡了,你們做這樣的事情.
這些人其實也挺可怕的,他有野心,要加害他,要將他從尊位拉下來.
不過更加令人心寒的是,這班人說謊,口就祝福.
心卻就坐.
這令我想起文革的時候有幾句很出名的話.
那時候每個人都拿著一本毛語錄,現在叫小紅書的毛語錄.
語錄不離手,萬歲不離口,天天說好話,背後下毒手.
不是我講的,是周恩來講的,你上網找得到的.
天天說好話,背後下毒手.
如果現在用廣東話講,即是說甚麼?.
笑騎騎.
你身邊真是很多.
我以為我講這些說話,原來身邊有這麼多笑騎騎放毒蛇的人.
明槍是二檔,對嗎?.

$^{161}$陰險是難防.
他面對這樣的人.
但他也說你們要收手了,你們不怕上帝要到何時呢?.
這些人處處逼人.
不過面對這樣的困難逆境,內心的躁動很容易就駕馭了我們.
你這樣害我,人不犯我,我不犯人.
你犯我,我當然要跟你拼,這是正常的.
或者拼完就走,也是對的.
詩人說我的心默默無聲,專等候神,他選擇了等候神.
我相信.
這是他跟上帝有一份很親密的關係.
他意識神的同在.
他知道上帝這個時間.
要他做甚麼.
這是跟神建立回來的關係.
我們做牧者常常勉勵弟兄姊妹,你要靈修,你要親近神.
你要意識神的同在.
在生活裡面.
而不只是在書本當中.
我們有沒有錯過了在生活裡面上帝.
對你微小的帶領呢?.
有個男人很餓,家裡沒東西吃,沒錢,他想出去碰碰運氣,一出去真的被他撞到.
撞到一個斬柴的巢夫.
有沒有東西吃?.
沒有,只有柴給你.
我不是要柴.
他拒絕了人的好意,氣沖沖地繼續走.
又見到一個.
賣煮食器具的人,一個小販,又問他有沒有東西吃?.
沒有東西吃,不過我有些器具,煮食器具送給你,我都不是要這些.
由他拒絕了人的好意,再向前走,終於見到一個小販.
賣菜,賣肉,賣米那些,有沒有東西吃?有,送些給我可以嗎?可以.
不過.
你有沒有柴火呢?.
你有沒有器具可以煮食物呢?全部都是生的.
他才發覺原來他錯過了沿途.
很多很多的好意.
我們常常都說我們要熟能操練.
當然我不是很喜歡用操練這個字,我比較喜歡用習慣,你培養一個習慣.
現在經常說養生.

$^{201}$我們都要養心的.
每天早上起床的第一件事,我一定不看手機.
為什麼呢?因為一開手機已經過百個WhatsApp了.
其實很多的研究都知道,太多的資訊會癱瘓了你的.
你一個人最清醒最清明的時候.
就是去親近神.
我也是這樣的,不用很長,二十分鐘到半個小時,你每天都是這樣做.
你打一本聖經,可能一節兩節.
祈禱,默想神的話語,然後安安靜靜,好像撒慕爾那樣.
請說.
捕人敬聽.
用十分鐘去等候神.
去聆聽神.
你真的不要小看,有很多時候.
你以為很重要的東西,原來不是,在內心裡面原來是別的東西.
那你就能夠敏感到上帝的同在.
耶穌說我和你的關係是什麼?葡萄樹和枝子.
來了我,你不能作甚麼.
當然這個慕容的比喻,在我們今天又不是容易明白的,不知道大家有沒有去過葡萄園.
我試過回內地,我看過葡萄園,我明白為什麼耶穌會這樣說.
如果今天耶穌在這裡跟我們說,我和你的關係是什麼?.
是手機和充電器的關係啊.
你每晚睡覺,你最著緊的是什麼?.
充電,你上了床又不記得充電,你都回去充電.
我自己也是,我試過沒有帶錢包,出去地鐵.
差不多入閘才知道,沒帶錢包,沒關係,手機照過.
我沒帶手機,結果碎碎地都要走回家,都要拿回手機,為什麼?因為所有的日程,所有的聯絡.
都在那部手機裡面.
親愛的弟兄姊妹,你和耶穌的關係就是這樣.
來了我,你不能作甚麼.
求主幫助我們.
和祂有這份感慨的關係.
意識神的同在.
讓神介入你的生活裡面.
我想這個很重要.
經文繼續下去.
唯獨等候神,我才有那份盼望.
唯獨等候神,我心裡才有安息,因為我的盼望是從祂而來.
唯獨祂是我的盤石,我的拯救,祂是我的保障,我必不自動搖.
我的拯救,我的榮耀都在於神.

$^{241}$我堅固的盤石,我的避難所都在於神.
中文啊.
你們要倚靠祂,在祂面前傾心吐意.
因為.
神是我們的避難所.
正如天主所說,你看到第二個段落的第五和第六節,幾乎和第一和第二節.
是重複平衡的.
唯一不同的字,就是由拯救轉為盼望.
而這個盼望的基礎是從拯救而來的.
明顯你讀完這篇詩,你會知道上帝還沒救他,他還在等候.
但在等候的過程裡面,盼望就成為一個很重要的力量.
因為上帝以往這樣來拯救我,現在雖然還在等候,但我仍然有這份盼望.
而這個盼望是基於那個信心的.
因為他經歷過上帝,他知道上帝是真的,上帝不會離棄他,不會放棄他.
在神面前默默無聲.
這個等候不是懦弱的表現.
第七節接著說榮耀,表示有份量.
當然可以用在神的當中,但由太陽與月這個字都會指到一些人的名聲.
社會的地位,名譽等等.
人們在這裡突出了上帝的榮耀.
接著第九,第十節其實是說人的短暫.
做了一個對比.
接著又再重複出現了,堅固的盤石,比喻那份穩妥,安全.
避難所,表示狂風暴雨的安全.
詩人被敵人攻擊,密謀,造反,要將他從尊位拉下來.
我們做過的事都知道,你不在那個位置,你沒有了那個位置,你就沒有所有東西了.
但是詩人也有一個洞察.
他說我今天坐到這個位置,其實都不是你,是因為上帝給我.
這份是一份屬靈的洞察力.
他沒有被那些成就蓋過.
而是他知道我今天在這個位置,都是上帝給我的.
因此他呼籲那些和他在一起的中文的人跟他講.
有些是管轄百姓.
大衛帶著這班人走難的時候.
告訴他們知道神是可靠的,唯獨神是我們的保障.
這是肺腑之言來的,你可以傾心吐耳.
你依靠的對象不要放在那些金錢,物質,甚至你以為有權位的人身上.
你要傾心吐耳,你依靠的對象是上主.
你跟祂可以有密切的關係,不需要隱藏.
你可以將你的苦澀和眼淚跟祂說.

$^{281}$你也可以將你的喜樂和憧憬向祂禱告.
基於這樣的理由,詩人要勸大家不要依靠物質的東西,或者人的力量.
你只要單單依靠上帝.
依靠的意思是無論什麼時勢,環境,所有時間當中,我都願意尋求求敏神.
而詩人的秘訣就是在於向神傾心吐耳.
詩人面對一個很艱難的環境,頹垣敗瓦,甚至有人惡意加害.
但是在這樣的環境裡,他仍然選擇盼望.
專心等候神,這份不是無奈,也不是一個逃避.
而是因為他有了一個正確的方向.
他內心有一份踏實,他不急於回應一個這麼惡劣的環境.
他有一份安息和把握.
因為神告訴他,你要有盼望.
在盼望當中,在艱難的裡面.
詩人這樣的宣告對你有什麼提醒呢?.
你又是否正正在一個艱難裡面呢?.
你有沒有這份從盼望而來的力量呢?.
上兩個月我去了日本開會.
每四年,亞洲的播道會聚在一起.
每兩年,就像梅花間中,會有青年的播道會大會.
明年我們會在香港開會,所以賣定這個廣告.
去年我們輪流開會,有五個國家.
有地區,還有香港,菲律賓,馬來西亞,新加坡,還有日本.
今年加上澳洲,所以我們改了名,叫亞太播道聯盟.
這些全部都是美國播道會在亞洲的席堂.
已經成為了一個中派,當然也有一些觀察員.
很開心,我們香港也有三四十個弟姐妹去.
我們說民族之夜,所以穿著中國服.
有些可能都認識,東九,林偉俊也有去,跟我們帶隊一起.
富士山拍的,這張照片不是AI來的.
先打卡,告訴大家我真的有去.
在會議裡面,除了好像每天的研習.
工作坊,或一些陪靈聚會之外,最令我享受的.
是跟來自不同國家的弟兄一起有團契.
去分享上帝在不同地方,祂的臨在.
我不會說日文,他也不會聽廣東話.
不過拿起電話打開AI翻譯真的能做到,真的能聽九成.
所以我聽得懂祂說什麼.
其中我跟小組分享的時候,這位弟兄,這位哲義弟兄.
我們開小組,他不是傳道人,他是社工.
他跟我們分享他怎樣信主的見證.

$^{321}$話說他在日本做社工,大家知道.
日本也有外勞,雖然不是很多,但也有的.
主要是菲律賓外勞.
人離鄉賤,很慘,去到別的地方.
他做社工的時候,當時就想幫這些外勞.
希望能夠跟他們設計一些社會政策之類.
有心到一個地步,直接去到菲律賓馬尼拉的大學做研究.
希望了解這個民族,各方面來幫他.
馬尼拉出名其中一個地方,不知道大家知不知道,叫垃圾山.
人們真的去倒垃圾,倒到山那麼高.
當中很多人在垃圾堆裡,去實況.
掙扎求存,環境多惡劣.
在那裡鋪一張布,在那裡就吃飯了.
在這樣的地方做田野研究,做問卷調查.
很自然,問卷的設計,你生活得開不開心.
不開心有什麼政策你想要,有什麼可以幫你,大概是這樣.
做了幾天,去到垃圾山,有一天他看到一個小妹妹迎面而來.
他就想做訪問,問你開不開心.
大家都猜到,怎麼會開心,那麼小就在那裡撿垃圾.
誰知道那個小妹妹回答他,我開心啊.
他呆了,我說為什麼你開心?.
他說因為我有主耶穌.
哇,他說那個衝擊大到他無法形容.
那個小妹妹沒理由騙他,對不對?.
他說為什麼宗教的能力可以改變一個人在惡劣的環境中.
仍然有這麼大的力量呢?.
於是他改了研究的題目,結果就去研究宗教.
研究宗教,研究基督教.
研究基督教,研究完之後就信了耶穌.
結果娶了一個菲律賓老婆.
就回到日本的教會裡來事奉.
親愛的弟兄姊妹,你面對著什麼困難?.
你的困難比一個在垃圾山十方的小女孩更大嗎?.
但是有盼望而來的力量就可以去面對.
求主幫助你.
給你這一份從盼望而來的力量.
經文繼續下去.
「卑賤的人虛空,高貴的人虛假.
放在天平上升起,給空氣還輕.
不要倚靠那些強暴得利的財物.

$^{361}$也不要倚靠那些搶奪回來的錢財.
如果財寶加增,你的心也不要倚靠它.
神說了一次,我兩次聽過.
能力屬於神.
主啊,慈愛也屬於你.
你必照國人所作的報應國人」.
看到最後一句,你知道報應還沒到.
如果到了,他就Hallelujah了,對不對?.
所以他在等候,在等候報應,審判來到.
當事人呼籲大家信靠神之後.
他就突顯另一個角度,或者另一個視野.
什麼不可靠呢?.
而正正那些是我們以為最可靠的人,就是人和錢.
卑賤的人,高貴的人做了一個比較.
這個不是要貶低那些人.
如果用成語說就容易明白一點.
無論你是販夫走卒,還是達官貴人.
無論你什麼的社會階層都好.
原來都是好像空氣那麼輕,毫無重量.
這個虛空和《傳道書》的虛空的字一樣.
他想帶出生命的那種短促.
接著說到世人以為金錢是最可靠.
甚至不擇手段,貪戀不義之財.
私人提醒,不要啊,你不要落入這些網羅裡面.
不是啊,我自己賺到錢啊.
都提醒你不要啊.
你不要將那個倚靠的對象錯置了.
不是錢,是給你金錢的那個上帝.
虛空,短暫,一陣風.
放在天平上面,稱都稱不到.
這個天平的意象很明顯是從真言而來.
世上有惡人的.
但是怎麼兇惡都好,都會消逝.
你倚靠他,何等愚蠢.
生命真是短促的.
早幾個月,大家知道.
弘福苑大火,舉世震驚.
不是誇張,我有些在外面的大英姐妹都WhatsApp我.
你有沒有事啊?我不是,我不是住那裡.
看到電視機上面的畫面.

$^{401}$我們一班執委馬上開會,總會馬上開會.
我們說不行,這個城市很大的創傷,我們要做點事.
於是馬上出一個通告,呼籲我們五,六十間的報道會,堂會.
機構,學校等等,我們募捐.
我們籌款賑災.
我們沒做過的,在香港沒做過的.
我自己做慈惟部長,以前在菲律賓我試過,但在香港我沒試過.
不懂得做,不過不要緊,我們知道我們本於愛.
我們就去服侍有需要的人.
短短一個禮拜籌了二百多萬.
最後我們籌了差不多七,八萬的.
全部用在災民身上.
行政費我們出,一個線都派給他們.
開頭我們想得比較簡單.
因為我們報道會在不同的區份都有堂會.
我們就聯絡不同的區份的堂會.
總有一些大英姐妹可能認識到一些災民.
又或者我們有些學校的網絡,NGO的網絡等等,總之就幫.
你要填表,可能一萬,可能八千就盡量幫.
現金給他.
開頭那個禮拜我們做了二十多個個案.
我們都覺得OK啦,這樣做沒問題.
誰知道做到第二個禮拜,那個畫面完全不同了.
我們收到很多的電話打上來.
是不是有援助金,是不是有錢派.
我坐在接待處的同事,我看著電話一放下,接著就響一放下.
真的是這樣.
接著呢,很多很多的申請表進來.
我們收那些電郵是用影印機的.
我不知道是否有關係,影印機都壞了,太多了.
接著有人走上來按鐘敲門.
其實我們是辦事處,尖沙咀的辦事處.
沒有禮堂,沒有這麼多房間.
是不是有援助金,我們很有需要.
我說為什麼發生什麼事.
原來有人將我們的資料放在WhatsApp的群組.
很簡單,大家都明白,一傳十,十傳八,八傳千.
那些一會一社工又將我們那些申請表放在群組.
幫他們填了,一傳十,十傳八,八傳千.
哇,那些人湧過來.

$^{441}$我真的祈禱,我靈生我祈禱.
神啊,你到底要我們怎樣呢.
我們怎樣去回應這麼需要的人呢.
大家知道不可以我拿現金走去派,不行的.
肯定出事的,大家知道一定出事的.
你兼顧不了,每一天都去想我們明天要怎麼做.
要派支票,不然你跟進不了.
始終我們要很謹慎地用善款.
哇,支票簿都沒有了.
大家知不知道,要拿支票簿有時候要兩個星期.
於是我們打電話下去.
問那些堂會,有沒有支票簿,先借幾本來.
有有有,不是借幾本,你借幾十萬來.
你借錢才有用,你借支票簿有什麼用.
跟堂主任說,跟他執事會,你先借來.
我們真的明白,大家明白,不明白.
我們不是騙徒.
沒有支票,我有錢都沒有用.
我每天就是這樣,今天去明天,去決定明天的事情.
又不能叫所有的災民去所有的教會,大家想像得到.
不行的,你說湧過千人,哪間教會都承受不了.
於是又聯絡,他們又四散.
有些去到洪水橋,於是我們打電話聯絡恩福元朗堂.
有些去到南區,於是我們又打電話聯絡港福堂.
他們有個服務中心,有些去到啟德那邊.
我們又打電話聯絡啟福堂,我當天就是做這些事情.
你又不能夠每次太多人,可能真的三十到五十差不多.
太多人真的應付不了.
為什麼?因為我們不只是想派支票給他們.
我們也想接觸他們,告訴他們.
神沒有忘記你們,我們也沒有.
我們很想為你祈禱,於是又要聯絡傳道人.
或者教會的領袖,童工,義工在那裡待命.
別人來到的時候,你真的有人接待他.
哇,很多的眼淚,很多的故事.
因為我也走到最前線,我當然也做了一份.
拿完支票,簽了名跟他祈禱.
我試過有一個災民.
他真的這樣跟我說,不停地在哭.
他說那時候政府說那些遇難者有一些恩恤.

$^{481}$給一些錢那些家庭有人離世的人.
若干的金額都是一筆錢.
他說他真的有一個遠房親戚.
一年見一次的打電話來.
聽說有一個人離世了,有錢分.
可不可以算上我一份.
哇,你有這樣的親戚,你也要哭.
他真的在哭,我說不要哭.
因為是男生來的,我拍著他.
我說不要哭,不要哭,在那裡祈禱.
這樣也有人性的光輝.
不只是這麼陰暗.
有些你聽的新聞都是,拍門叫人走.
走了之後就失散了,不見了.
哇,然後在教會拿支票撞回來.
你想想,真的抱頭痛哭在那裡.
然後牧師就擁抱他們,不要哭.
我們一起為你們祈禱.
恍如隔世.
我自己有些個案,我都真的很觸動.
我試過有一次跟一個災民來祈禱.
祈完禱之後.
他說牧師可不可以再為我祈禱.
我說當然可以了,為什麼.
他說這十幾天我沒試過這麼平安.
我真的沒試過這麼平安.
然後他說,他的兒子下個禮拜就考試.
如果不是考試.
大家做家長就知道放早,放早兩天而已.
這樣就燒死他了,就沒了.
人生那種生死,你看,真是言於一善.
然後他小時候回過教會學校,我跟他傳福音.
他最後信了耶穌.
就是這樣.
你想想,如果你是那個災民.
你見到自己的房子一直在燒,你想怎樣.
你什麼都想拿,對不對.
但結果怎樣,什麼都拿不到.
這個豈不是都是你和我的寫照嗎.
我們出生帶了什麼來.

$^{521}$當我們離世的時候,你帶了什麼走.
生命是很短促的.
生命是很輕的.
但是上帝是猜測他讀《新愛子耶穌基督》.
成就了這個沉重的救恩.
新元十教告訴我們他們世人的愛.
求主幫助我們.
讓我們告訴身邊的人知道.
神愛你.
福音是好的.
是值得我們傳頌的.
面對生命的短促.
先人告訴我們不要倚靠那些財富.
你以為有勢力的人.
你要倚靠的是神.
並且這裡說,上上倚靠祂.
神說了一次,我聽了兩次.
於是又釋經問題了.
到底神說了一次還是兩次.
不過我們說你讀詩篇不可以這樣讀.
詩篇是一首歌,是一個韻律.
好像一個層次.
或者我用個比喻,好像走樓梯一樣.
一級,兩級,三級.
有個遞進.
告訴你重點是什麼.
能力在乎神.
的確面對人生裡的艱難.
甚至我們很想報應快點來.
但先人知道不是他決定.
能力在乎神.
當然我們知道耶穌基督回來.
我們每一個都要在基督的台前交仗.
信耶穌都要受審判.
你不要以為走得掉.
不過神已經稱我們為義.
我們能夠在這個救恩的當中.
事實上我們很想不用等到那一天.
最好快點有審判.
快點有報應.

$^{561}$有時上帝的管教真的來得很快.
神的報應可以來得很快.
這正正反而是我們需要學習的那份謙卑.
因為能力在乎神.
不是我們做到的.
先人提醒我們要專心倚靠.
放心投靠.
因為神的能力是無限的.
祂的慈愛也是無限的.
這個慈愛現在有些譯本.
想強調前面是堅定的愛.
不只是愛這麼簡單.
是堅定的愛.
是真正的愛.
告訴我們知道.
身處職場.
我們每天面對我們的生活.
其實很多時候都有很多試探.
特別現在你看新聞報道.
說原來有些事情可以這麼黑暗.
我自己有個弟兄.
他做那些測量的.
即是建樓.
有時跟他聊這些話題.
他說你現在才知道.
我見識少 不懂這麼多東西.
不過他真的有次跟我說 我也很深刻.
他說牧師 你讀聖經吧.
我說是的 讀聖經.
你看《但以利書》.
但以利是被人推下獅子坑的.
他也是一輩子被人推下獅子坑一次.
他說我上班每天都被人推下獅子坑.
你為我祈禱吧.
我說真的 祈禱吧.
原來你們的行業這麼黑暗.
我說祈禱吧.
真的.
我也真的有個肢體.
他做內地的私募基金.

$^{601}$有次跟我說.
他說他收到一張支票.
寫在他的抬頭.
下面簽了名.
不寫銀碼任你填.
不過你要做某件事.
當然是有些不義的事情.
我又好奇.
我問他最後有沒有要那張支票.
他說當然沒有.
還能回來.
有時真的是這樣.
在生活裡面我們面對很多的試探.
我們有沒有這樣依靠神呢.
提到依靠.
我想起我一位弟兄.
他本身是做一間公司部門的主管.
真的做得不錯.
我見過他很長時間的工作.
不是經常轉工.
做得很有貢獻.
誰知換了一個CEO.
就將他從一個部門調去另一個部門.
而那個部門.
全公司都知道是沒有發展和等執.
他真的很不服氣.
我說我又不是沒有建樹又不是沒有貢獻.
為什麼你要這樣對我.
當然很多的職場故事都是這樣.
原來安插了他的親信.
故事大概是這樣.
他真的電郵去總公司投訴.
都不得要領.
你又不是什麼事.
他祈禱.
其實正常人都要不跟他拼過.
要不就寧謀高就.
但他祈禱上就叫他等.
他又是敬虔人.
他說你祈禱上就叫他等.

$^{641}$他就等.
但不知道等什麼.
不過回到他自己工作的崗位.
他就做回自己應有的事.
他又不會搞其他的一些手段.
總之就忠心繼續做好工作.
這兩個月不夠.
這個CEO就被人解僱了.
然後很快又調回他原本.
做那個部門的主管.
他都覺得很奇怪.
發生什麼事呢.
你知道辦公室是沒有秘密的.
字字枕枕的可以聽到.
原來這個CEO跟他一個下屬有染.
搞婚外情.
公司說不行.
這個你自己做.
帶領你的操守問題.
於是就解僱了他.
他才發現原來他以為等得很辛苦的時候.
上帝其實已經在工作.
參加電影節目.
我不知道你面對著什麼艱難.
我們剛才唱的詩歌.
《投靠在主的翅膀任下》.
我也想敬拜的團隊幫我們彈這首詩歌.
為什麼呢.
如果你看詩篇62篇.
令人感動的地方就是.
很容易產生共鳴的.
很真實的.
當惡人當道的環境.
我要掙扎求全已經很艱難.
你可以放棄.
甚至消極厭世.
又或者.
你屈服.
同流合污.
不過詩人就選擇了第三條路.

$^{681}$第三條路就是.
倚靠耶和華.
我的心才有安息.
詩人的選擇.
你們有什麼提醒呢.
我為弟兄姊妹祈禱.
我們當中有些可能在工作上.
真的面對很大的艱難和抉擇.
求聖靈讓他們有清楚的辨別和帶領.
好讓他們能夠回到職場的.
事奉的場景裡面.
仍然地來倚靠神.
主要我們也為當中.
可能有些正在面對抉擇的.
求聖靈指引他,帶領他.
因為你是又真又活的神.
願你這樣來成為我們的力量.
主啊,求你這樣來施恩.
幫助我們.
又願你充滿在我們當中.
讓我們有那份投靠和確信.
主啊,多謝你.
因為你是我們的高臺.
你是我們的山寨.
你是我們的避難所.
我必不致動搖.
求你幫助我們站得穩.
穩站在這個真理當中.
願你垂念.
越南我們眾人在你面前的禱告.
奉耶穌基督的聖名而求.
阿們.
\newpage



\section{約翰福音 4:46-54}
\label{sec:032163}
\textbf{信心的建立}~~~~曹子光牧師
\newline
\newline
連結: \href{http://www.church.com.hk/acms/content.asp?site=cdc\&op=show\&type=product\&code=032163\&layout=sermon}{\texttt{032163}} ~~~~ 語音日期: 2026-05-10
\newline
\newline
\hyperref[sec:032162]{\small{< < < PREV SERMON < < <}}
~
\hyperref[sec:index_chronic]{\small{[返順時目]}}
~
\hyperref[sec:032164]{\small{> > > NEXT SERMON > > >}}
\newline
\newline
約翰福音 4:46-54
\newline
\begin{longtable}{cl}
\hline
\hline
章節 & 經文 (和合本修訂版)\\
\hline
4:46 & \begin{tabularx}{0.7\textwidth}{X} 耶穌又到了加利利的迦拿，就是他從前變水為酒的地方。有一個大臣，他的兒子在迦百農病了。 \end{tabularx} \\ \\ \relax
4:47 & \begin{tabularx}{0.7\textwidth}{X} 他聽見耶穌從猶太到了加利利，就來見他，求他下去醫治他的兒子，因為他兒子快要死了。 \end{tabularx} \\ \\ \relax
4:48 & \begin{tabularx}{0.7\textwidth}{X} 耶穌對他說：「若不看見神蹟奇事，你們總是不信。」 \end{tabularx} \\ \\ \relax
4:49 & \begin{tabularx}{0.7\textwidth}{X} 那大臣對他說：「先生，求你趁著我的孩子還沒有死就下去吧。」 \end{tabularx} \\ \\ \relax
4:50 & \begin{tabularx}{0.7\textwidth}{X} 耶穌對他說：「回去吧，你的兒子會活！」那人信耶穌所說的話，就回去了。 \end{tabularx} \\ \\ \relax
4:51 & \begin{tabularx}{0.7\textwidth}{X} 正下去的時候，他的僕人迎面而來，說他的兒子活了。 \end{tabularx} \\ \\ \relax
4:52 & \begin{tabularx}{0.7\textwidth}{X} 他就問甚麼時候見好的。他們對他說：「昨天下午一點鐘熱就退了。」 \end{tabularx} \\ \\ \relax
4:53 & \begin{tabularx}{0.7\textwidth}{X} 他就知道這正是耶穌對他說「你的兒子會活」的時候；他自己和全家就都信了。 \end{tabularx} \\ \\ \relax
4:54 & \begin{tabularx}{0.7\textwidth}{X} 這是耶穌從猶太回到加利利後所行的第二個神蹟。 \end{tabularx} \\ \\
[1ex]
\hline
\hline
\end{longtable}
$^{1}$我有一個朋友,很多年前他喜歡用一隻手機的牌子叫做Samsung.
他總是覺得,不知道什麼原因,三叔公的手機的電池特別快玩完.
不知道你有沒有同感,每一個人不同.
所以他對於三叔的牌子,這個牌子打了折扣.
於是沒多久他就轉用了另一隻水果牌子的手機.
就發現他的電池質素就明顯改善了很多.
所以從此之後就轉換了,不再用三叔這個牌子.
就轉用了水果這個牌子.
當你相信那個產品的時候.
往往是因為你曾經用過那個產品.
所以你投他信心的一票.
大家不知道有沒有試過肚痛.
大家都知道Benny是最容易肚痛的一個人.
有人經常介紹我不吃不同的東西.
譬如補交叉丸.
又或者有些人就介紹,不如你試試日本那隻大什麼牌子的紫蛇丸.
有沒有試過?.
你沒聽過?這隻聽聞是挺好的東西.
有些人就說,不好呀,你試試那隻臭到不分軒那隻什麼樂器牌那隻東西.
那你就會更加快些痊癒.
這種信心是建基於別人用過,推介給你的.
又或者我們很怕被人放飛機.
又或者很怕被人背後捅你兩刀.
我們對一個人的信心往往是建立在兩件事上.
第一,他有沒有曾經承諾過一些東西.
做不做到他的承諾.
第二,他有沒有實力去兌現.
而今天讀的這個經文.
這個大神,我們待會讀的這個經文裡面.
有個官,有個大神所面對的.
正是一個關於承諾和實力的東西.
我們來看看約翰福音第四章46到54節.
找到的時候我們一起讀出來,請讀.
耶穌又到了加利利的加拿.
就是他從前變水為酒的地方.
大神,他的兒子在加伯隆患病.
他聽見耶穌從猶太到了加利利.
就來見他,求他下去醫治他的兒子.
因為他兒子快要死了.
耶穌就對他說.

$^{41}$若不看見神即其事,你們總是不信.
那大神說.
先生,求你趁著我的孩子還沒有死就下去.
耶穌就對他說.
回去吧,你的兒子活了.
那人信耶穌所說的話就回去了.
正下去的時候.
他的僕人說.
他的兒子活了.
他們說.
未時而就退了.
他便是耶穌對他所說.
今天想和大家看三樣東西.
第一樣東西就是.
信心通常都是經過痛苦來萌芽的.
特別想你留意第46和第47節.
經文告訴我們.
有一個人去到耶穌面前.
來求耶穌在他兒子還沒有死的時候.
趕快下去救他一條命.
這個來到耶穌面前的男人.
他是一個父親.
不止,聖文告訴我們.
他更加是一個大神.
很可能就是當時這個地區裡.
希律王的內閣的重臣.
因為當時加伯隆是一個邊界的重鎮.
有很多高官在那個地方駐守.
他可能就是其中的一個.
根據約翰福音第二章第十二節告訴我們.
耶穌其實原來在六到九個月之前.
曾經去過加伯隆這個地方.
而四章四十六節也告訴我們.
這個加伯隆的大神的兒子患病.
患病這個字,希律文這個字.
是一個未完成式.
換句話說,他是持續進行著患病.
意思就是說,這個兒子很有可能不是突然間病.
而是一直都在病.
只是現在面臨近死亡,瀕臨死亡的邊緣.

$^{81}$換句話說,當耶穌在六到九個月之前.
去到加伯隆這個地方.
這個兒子,這個官員的兒子.
其實可能正在病.
或者未去到垂危.
這樣當耶穌去到他家門口附近的時候.
這位的爸爸從來都沒有去找過他.
是不是沒有聽過耶穌這個名字呢?.
還是他太忙,太busy沒辦法去找到耶穌呢?.
還是他從來沒有想過.
他從來都沒有想過,他要找這個草根不過的猶太人的老師.
初出道,名氣是不是還沒有呢?.
到底值不值得去光顧他呢?.
更有可能的就是.
他身為一個高官.
已經有充足的醫療資源給他用.
例如,希臘羅馬的醫術當時非常之高.
還有猶太傳統又不同的醫術.
再加上當時在近東一帶埃及的醫術非常之有名.
而且是傳到街知巷聞的.
甚至更加有不同宗教的儀式.
他什麼都可以嘗試.
他有大把資源,有錢,有勢力.
甚至有所有的門路.
不過直到有一天.
他發現所有的辦法都不work.
耶穌已經離開了加伯隆.
他的兒子也撐不下去.
去到一個這樣絕境的地步.
無路可行的時候.
他就謙卑,親自去找耶穌.
一個這麼有身份的官員.
為了自己的兒子.
從加伯隆去到迦拿這個地方找耶穌.
兩個地方相隔多遠呢?.
大家知不知道我們今天是馬拉松嗎?.
馬拉松今天是半馬加全馬.
全馬大家知不知道有多少公里?.
多少?.
42公里是嗎?.

$^{121}$半馬,好像不止吧?.
全馬有40多公里?.
Dickson是跑過全馬的.
45公里,我猜的.
半馬應該是一半吧?.
一半的話應該是20-23公里左右.
其實23公里遠嗎?.
我經常問人們走10公里路.
不要玩了,10公里路走很久的.
現在這個世界有車,有巴士.
樣樣都可以,但腳就不可以走.
其實有多遠呢?.
加伯隆去到迦拿是足足35公里左右.
是35公里.
走路的話,我查過Google Map.
走得快的話都要七八個小時.
七到八個小時才走得遠.
大概你都知道他是高官.
你幾時會見到一個高官走這麼遠呢?.
不要玩了.
我經常聽到網上的人說他走路.
簡直浪費力氣,沒可能走路的.
他一定是坐交通工具.
就是一些馬或者駱駝.
但因為地形的緣故,極有可能是坐雷.
雷都要四到五個小時.
所以應該一大清早就出發了.
大家有沒有試過為了一件事.
拼到一個地步,完全放低了自己的身份.
這位大神放低了自己的官威.
親自去見一位在他眼中頂多都是一個木匠.
和民間老師的耶穌.
還要求祂.
換作是你呢?.
你可能放不放得下你心裡面最帥的那一面.
放下你的自尊,有沒有辦法呢?.
信心不是在冷氣房裡面想出來的.
是趕著那三十五公里的路上的汗水,眼淚.
這樣營造出來的.
磨練出來的.

$^{161}$我們一起看看第四章第四十七和第四十九節.
這兩節經文非常特別.
四十七節就說.
他聽見耶穌從猶太到加利利就來見他.
求他去醫治他的兒子.
四十九節那位大神就說.
先生求你趁著我的孩子還沒死的時候就下去.
你有沒有發現兩節經文兩次講到這個爸爸去求耶穌.
而且兩節都強調死.
非常明顯.
這個爸爸所承受的就是死亡快要親生他的兒子.
是極其恐懼,極其害怕,極其絕望的處境.
我絕對相信我們做父母.
這裡有很多人做父母.
如果可以選擇.
我們都寧願和他交換身份.
有問題的是我們,病的是我們.
這個爸爸,所以我們可以想象.
他多趕急去到耶穌的面前去求耶穌.
我認識一位太太.
一位太太都幾年輕,不是很大的太太.
醫生告訴她,她患了癌症.
情況一點都不樂觀.
而且有擴散的跡象.
她曾經好多次埋怨.
為什麼?為什麼是我?Why me?.
就是在人生最絕景,最不開心的時間.
她認識到一些教會的朋友.
在她鼓勵之下.
再一次回到教會那裡.
再一次去找耶穌.
雖然病況一點都沒有改變.
但是她的心境改變了.
她和很多人說.
如果不是這個痛苦的病患.
她根本不會再遇到耶穌.
這個正正就是大臣的故事.
如果沒有痛苦.
如果兒子沒有病成這樣.
如果他手上還有很多籌碼.

$^{201}$有很多辦法的話.
難道他會去找耶穌?.
難道他會放下自己的身份.
去自己的身段.
去親自去求救?.
不知道大家知不知道.
戈爾夫球是何時發明的呢?.
二百多年之前.
原來當他製造戈爾夫球的時候.
原本戈爾夫球是滑溜溜的.
是平的.
是完全沒有凸凸的.
所以人們用極力打出去.
就發覺效果很奇怪.
是打得到.
但是速度又不是很快.
而且又不是很軟.
結果就很奇怪地.
這支球越打越舊.
但是打舊了的時候.
不知道為什麼.
窮人就撿到舊的球.
接著就發現越打越厲害.
而且越軟又有力.
有錢人就開始學東西了.
就發現原來窮人撿到這些球之後.
可以這樣.
於是就開始改良戈爾夫球.
後期人們就將戈爾夫球浮.
變成凹凸不平.
大家知不知道這個凹凸不平的面.
有多少凹凸位呢?.
有多少個凹凸位呢?.
有沒有人戴眼鏡呢?.
馬上不敢了.
是有三百三十幾個.
一個球原來有這麼多個凹凸位.
人生就有很多這些凹凸位.
很多不平滑的經歷.
正是因為這樣的磨練.

$^{241}$才能夠練出成熟的生命.
信心的種子.
是可以從我們所經歷的痛苦裡面.
成長出來.
孕育出來.
痛苦的來臨.
告訴我們.
到底我們應該怎樣看人生.
教我們去尋求生命的意義.
是一個給我們對神的信心萌芽的大好機會.
但是請我們不要忘記.
請我們不要忘記.
我們不需要邀請祂來.
苦難來到的時候.
我們更加不需要跟神說Come on more and more.
我們不需要這樣.
因為.
祂會自然找上門.
祂會自然敲你的門.
在你不知不覺的時候.
祂就不請自來.
硬闖入我們的生命的世界裡面.
但是我們當時.
如果是這樣的時候.
我們要知道.
我們可以從.
面對痛苦的裡面.
培植出對神的信心.
這個正是我們對神的信心萌芽的大好機會.
還有第二樣東西.
這段經文提醒我們.
信心是必須要自己去親自去尋找.
我們相信.
這位的爸爸.
這位的高官.
在他人生裡面.
應該從來都未正式主動去見過耶穌.
他沒有直接聽過耶穌.
任何一篇的講道都未定.
所有他對耶穌得來的東西.

$^{281}$都是人家告訴他.
他對耶穌的理解.
他對耶穌的聽的講法.
都是全部.
都不是他一手資料.
都是二手資料得來.
但是.
我們要知道.
信仰是必須要自己親自去經歷.
不是人告訴我們.
而是我們必須要自己親身去體會.
我記得我和大家分享過.
我好喜歡看北極光.
雖然我人生未看過北極光.
我曾經.
我曾經心裡面想.
北極光有多震撼.
真是我都好想知.
但我心裡面想.
不如google一下.
就了解自己的心事.
不知你敢不敢google一下.
出來的那幅圖畫.
和你真實見到的圖畫.
有沒有不同呢.
我以為沒什麼分別.
直至我有一年去西安.
我第一次帶隊去看兵馬俑.
那次人生我從來沒去看過兵馬俑.
我常常都心裡面一堆石頭.
so what.
當我入去看的時候.
我真是震撼了一下.
原來是這麼偉大.
可以是這麼震撼性.
大家都知.
我都去帶很多聖地團.
我一直帶聖地團的時候.
以前.
我帶聖地團.

$^{321}$我常常都好想去兩個地方.
第一個就是曠野.
沒什麼團帶你去曠野.
坦白講.
因為危險.
還有沒什麼團會帶你去看所羅門銅礦.
因為遠.
那個近紅海的地方.
我有一九年的時候.
一九年我最後一次帶聖地團.
我記得.
我現在都沒帶過.
一九年去最後一次的時候.
我特意設計了一個團.
是去這兩個地方.
我一去完這兩個地方之後.
嘩.
你真是靚到不得了.
比起之前google的圖畫.
完全是超越了想像力.
其實.
這個正正是我們有時候實體.
和你在google的看法.
是完全不同的.
有些像什麼.
有些人在玩什麼.
黑神話的悟空.
又或者youtube裡面的開箱片.
你會覺得自己很懂裡面的東西.
但是你根本是沒有手感.
你只是看沒用的.
你要有那個手感才行.
就好像你父母的whatsapp account那樣.
那個不是你的.
你需要有自己.
要登入你自己的account才行.
信仰都是一樣.
信仰不能夠靠父母.
父母只是開一條路給你進去.
你是要靠自己.

$^{361}$去尋求這個信仰.
是要和耶穌建立這個關係.
信仰是沒有自動轉帳.
也都不能夠靠課金.
是因為這個緣故.
我們就能夠明白.
耶穌在第四章第48節那裡.
祂說你們若不看見神蹟奇事.
總是不肯信.
來這裡就是約翰邦.
唯一一次用這兩個字走在一起.
神蹟再加上奇事.
兩個字放在一起.
就叫人想起.
猜及記第七章第三到第四節那裡.
神童摩西說.
雖然他在摩埃及這個地方.
不斷行神蹟奇事.
法老見到這麼多神蹟奇事.
但是法老是怎麼樣.
都不肯相信.
耶穌好像不太客氣.
但是其實.
祂是在挑戰.
身邊周圍的人和這個父親.
祂說你來找我.
是因為想追求我.
還是想看神蹟.
還是你真的想認識我這個人.
信心.
是個人和神的相遇.
不是神蹟的消費.
這位爸爸.
面對著耶穌的認識.
基本上大部分都是二手貨.
所以當他來到耶穌的面前的時候.
他心裡面已經預設了.
耶穌應該怎麼去醫.
甚至他相信.
只要耶穌去.

$^{401}$就能夠醫得好.
耶穌真的很諷刺.
祂就不去.
基文告訴我們.
祂就不去.
你以為耶穌這樣去醫.
耶穌就不是這樣去醫.
每一個人都假設了.
耶穌是這樣這樣這樣.
信仰不是二手.
這個看法.
他這個二手的看法.
限制了他對上帝的認識.
和上帝能力的認識.
但耶穌的回應.
他遠遠超越了他所想所求.
根本就不需要他自己親自出去.
而且他一離開.
馬上就康復.
還有我想你留意.
這段經文.
用了三個很不同的希臘文字眼.
去描述這個兒子.
爸爸在四章四十九節那裡.
他叫他兒子做Padion這個字.
這個字是一個非常親密的稱呼.
帶著一種爸爸的那種.
心裡面很疼愛兒子的感覺.
瀑人在第四章四十五節.
就用了另一個字.
Pais這個字.
其實就是一個很熟悉.
很親切的字眼.
好像一個長期相處了一段時間的關係.
但耶穌和約翰.
他用了另一個字.
Heroes這個字.
兒子這個字.
一個充滿著尊嚴的字眼.
耶穌看重兒子的生命.

$^{441}$比父親所求的更深更廣.
很多時候我們求的是病好得翻.
考得好一點.
過得關.
耶穌給的是什麼呢.
是生命.
是身份.
耶穌重視的是你的人是怎麼樣.
這裡我們見證到.
耶穌對生命的重視.
就算是一個從來沒見過的小孩.
他都一樣尊重他的生命.
甚至每一個人在座的每一個人.
每一個片段.
經歷的每一個關口.
都一樣緊張.
都一樣關心.
他都很想我們每一個生命.
和我們每一個生命相遇.
所以生命要自己親自和耶穌相遇才行.
就好像在楊方第二到第四章這個段落裡面.
約翰用了兩個迦拿的神跡.
一頭一尾.
好像是括弧一樣.
括住了中間兩個人物.
第一個人物就是尼哥底姆.
他半夜去找耶穌.
這個有錢的男士.
接著在第四章出現第二個人士.
就是撒瑪利亞.
一個邊緣的人.
是一個婦人.
甚至他是一個窮一點的人.
他在正午這個時候去找耶穌.
這位爸爸.
這個今天的主角這個爸爸.
當他走投無路的時候.
他都去找耶穌.
每一個人的處境是不同的.
但是他們都是親自去遇見耶穌.

$^{481}$沒得代價的.
沒得代勞的.
不知道大家有沒有聽過現代科技呢.
有一種眼鏡叫做Encoma.
是專門給色盲的人士設計的.
有一位年輕的盲人叫做.
Aaron William Mell.
他天生就有嚴重的紅綠色盲.
從小到大他都知道世界有的顏色.
他都知道的.
因為他聽過老師講的.
草是綠色的.
他女朋友經常告訴他.
最美的玫瑰是紅色的.
有無數的人說.
這個日落多麼壯觀多麼美麗.
甚至他都知道.
他自己都可以感覺到紅綠燈在哪個位置.
他擁有了對顏色所有的知識.
但在他眼裡面.
世界不是這些顏色.
世界是灰灰黃黃的顏色.
就好像退色的舊相片一樣.
他對顏色的認識.
完全是建立在別人說甚麼.
他就相信甚麼的層面.
直至到他生日那天.
有個朋友送了一副Encoma給他.
送了這副給色盲的人戴的眼鏡給他.
他站在海邊.
戴著這副眼鏡那一刻.
他看見原本朦朦朧朧的海天.
他就發現.
他的顏色怎麼會這麼燦爛奪目呢.
他就當堂跪在地上抱頭痛哭.
他後來就說.
我以前以為自己知道了紅色是怎樣.
原來我根本不知道.
直至這一刻.
我才真正和這個世界connect.

$^{521}$很多時候我們回教會.
我們都是聽別人說.
神的愛有多動聽多精彩.
但是你自己有沒有親身.
戴上這副信心的眼鏡.
親自去經歷一次.
神在我們生命裡的那種顏色呢.
信仰不是研究顏色.
而是你親自走進那個彩色的世界裡.
就好像聖經裡面的大神一樣.
他現在不再停留在純粹聽人說耶穌是怎樣怎樣.
而是他自己親自起行.
親自去找耶穌.
親自和祂相遇.
生命就是要這樣.
怎樣我們可以和上帝相遇呢.
最簡單最簡單的就是和祂聊天.
傾心吐耳去禱告.
說了你心裡所說的話.
這個是相遇的起始點.
還有第三樣東西.
信心就是必須直跟在神的話裡.
這裡包括了三樣東西.
第一樣東西就是.
如果我們的信心是直跟在神的話裡.
那麼我們自然就會信服我們的主.
這個神蹟是和水變走的神蹟有一個共通點.
就是原來耶穌開始的時候都不是那麼願意去理的.
不知道你有沒有記得.
在水變走的神蹟裡面.
耶穌和祂的媽媽說過什麼話.
希望你今天不要和你媽媽這樣說.
因為今天是母親節.
是不是母親節.
所以不要說了這句話.
聽聖經這樣說就算了.
婦人我與你有何干呢.
我的時候還沒到.
所以今天媽媽叫你洗會的時候.
你千萬不要當我說這句話.

$^{561}$今天的這段經文.
這個父親第一次求了之後.
耶穌48節就說了.
若不按見神蹟你們總是不信.
是不是很相似.
不過這個父親沒有放棄.
他繼續去求.
所以耶穌第四章五十節就說.
回去吧.
你的兒子活了.
我想問大家.
當他走的時候.
他的兒子活了沒有.
活了.
他知不知道.
他不知道的.
他還沒見到.
但他竟然就回去.
我們要記得.
這個父親去迦拿的目的.
就是要帶耶穌去.
要帶耶穌回去他的家.
去醫治他的兒子.
但結果呢.
結果是他自己一個人回去.
因為他信服了耶穌所說的話.
換作我們就很不同了.
可能會跟耶穌說.
不是吧耶穌.
你要我走.
他已經康復了.
遙遠發功.
我聽都沒聽過啊.
可不可以啊耶穌.
不要耍我了.
你隨便吧.
你隨便把藥給我.
你可不可以跟我一起去.
隨便寫些東西給我都好.
香港人最明白的.

$^{601}$你給了醫生五百元.
醫生說.
回去吧.
你的病康復了.
你跟著就說.
肯定是回水吧.
然後就在網絡平台罵爆他.
香港人很明白的.
英國醫生就一定不是.
英國醫生就不給藥給你.
你打電話去的時候.
你現在排第四十八個.
等等等到零的時候.
聽電話.
復了八點鐘.
不是八點鐘太晚了.
五點鐘去看一個訪問.
五點鐘去的時候.
當你去到見那個醫生.
醫生問你.
你還可以啊.
沒什麼啊.
你休息多點就可以了.
回去吧.
你有沒有試過這樣去看醫生.
我真的試過了.
你知道他的答案是什麼嗎.
我那時候腰背痛.
他給了我一個答案.
不是一個字.
是單字.
Unlucky.
我真的瘀青了.
那是什麼意思.
有沒有藥吃啊.
我等了很久.
看到這個appointment.
等了很久看醫生.
你給我一個這樣的答案.
然後他說.

$^{641}$沒多久你就康復了.
如果你真的要吃藥.
出街買.
那些什麼Panadol.
Ibuprofen.
幫到你了.
我用得著去看你嗎.
所以我明白的.
我明白的.
這裡看醫生和香港看不一樣.
香港看醫生.
你這麼多錢.
起碼都給我幾劑藥.
是不是.
耶穌這樣去.
這個爸爸竟然都聽.
信心植根在上帝的說話裡面.
其實是信服的根基.
這個爸爸信服耶穌的說話.
還有第二樣東西.
如果我們的信心植根在神的說話裡面.
我們的心靈就會輕散.
我們一起讀51和52.
請讀.
正下去的時候他的僕人迎見他.
說他的兒子活了.
他就問什麼時候見好了.
他們說昨日未時夜就退了.
這裡有一個很有意思的細節.
僕人就出來迎見這個爸爸.
告訴他.
你的兒子康復了.
昨日下午就康復了.
一點鐘左右就退燒了.
昨日.
什麼叫昨日.
就是已經過了一晚.
是不是.
我們要知道.
加拿大去到加巴隆.

$^{681}$大概三十公里.
剛才說了.
走得快一點.
八個鐘就搞定了.
但是我們剛才說了.
他應該是坐雷仔去的.
照計.
如果他一點鐘左右見完耶穌.
當遲了半個鐘.
一點半鐘就走人了.
當.
他什麼時候.
他什麼時候回家.
他騎雷回去的話.
六點鐘可以回去吃晚飯.
是不是.
你明白我的意思嗎.
每個人都會趕回家.
看看兒子康復了沒有.
是不是真的康復了.
但是他說昨日.
那就是代表他沒有回家.
如果今時今日.
老婆就打電話.
打爆你的電話.
就召喚你.
去哪裡了.
這麼久都沒回來.
但是他竟然休息了一晚.
為什麼.
他不是急著回家.
為什麼.
因為耶穌的說話在他心裡.
他手裡拿著什麼.
他不是一袋藥袋.
沒有藥拿回去.
沒有任何秘方拿回去.
甚至連耶穌都沒有和他一起回去.
但是他有一件事拿了回去.
就是耶穌的那一句說話.

$^{721}$你的兒子活了.
這句耶穌的說話.
已經足夠讓他安定下來.
這就是walking by faith.
not by sight.
就好像你在網上.
訂了一份禮物給你媽媽.
因為今天是母親節.
貨還沒到手.
你當然害怕.
但是原來order confirmed.
不過現在order confirmed.
都要小心一點.
要delivered才行.
就是送了貨.
那你才安心.
還要拍張照片給你看.
這樣才可以百分百安心.
沒錯.
照片都看了.
現在安心了.
信心的最高境界.
就是當事情還沒發生.
你已經安心地去睡覺.
去吃飯.
去慢慢走.
有了神的話.
安放安定在我們心裡.
心靈就會平安.
我們不需要再恐懼.
最後第三件事.
如果我們的信心.
是直跟著神的話.
我們就會影響其他人.
有時候我們會覺得.
我信不信有什麼關係.
但是不要忘記.
你的信心.
其實從來都未必是你一個人的事.
你看看四章第五十三節.

$^{761}$我們一起讀一讀四章第五十三節.
請讀.
他便知道這正是耶穌對他說.
你兒子活了的時候.
他自己和全家就都信了.
哇 全家人.
連同僕人全部都相信耶穌.
我們要留意.
這位父親的家人.
其實遠遠早於這位爸爸.
見證了兒子康復.
是不是.
你好像沒有反應.
是.
因為他爸爸離開耶穌的時候.
那一刻.
很遠的兒子已經康復了.
是不是.
所以他的家人.
在他兒子旁邊的家人.
沒有理由不知道.
他們第一時間就應該知道.
照道理他們應該猜到.
一定是這個醫治.
一定和耶穌有關.
因為他爸爸就出去了找耶穌.
而這麼巧.
他的兒子就在這時候就康復了.
不過.
但聖經在四章第五十三節那裡.
重複說多一次.
你兒子活了.
然後才強調.
這個爸爸是先相信.
然後全家才相信.
換句話說.
雖然全家比爸爸更早見證到神蹟.
但帶領全家人信主的.
是這個爸爸的信心.
一個人的信心.

$^{801}$可以打開了一道寬闊的門.
讓全家都一起進來.
你的信心.
是會影響你的朋友.
你的家人.
你的親戚.
甚至你的同事.
你的同學.
甚至你從來沒有預計到的人.
大堂的弟兄姊妹.
我都和大家分享過.
有一個病「白石」的經歷.
我不知道青少年有沒有聽過這個經歷.
我曾經有一個病叫「白石」.
這個病就在嘴唇角這個位置.
就會不斷擴散出去.
我剛讀完神學之後.
年青有為的時間就發現了這個病.
我非常不開心.
靠著讀聖經.
我就祈禱醫治了.
到現在都暫時沒有任何復發的跡象.
我經常都問.
到底我這個病.
可以康復.
有什麼特別的東西呢.
我一開始都不明白.
直到有一次.
我在一個地方講完道之後.
有一個男生他讀醫科的.
他走來找我.
他說Benny.
可不可以和你在一個房間裡聊聊天.
我就和他在一個房間裡聊聊天.
然後他脫了衣服給我看.
我背後全部都是白色.
他說剛才你講完道的時候.
我很受感動.
他說你可不可以和我分享一下你的心路歷程.
沒錯.

$^{841}$在那一刻開始.
我就知道.
神當年的醫治.
不只是為了我自己.
而是為了有一天.
用我的經歷.
去開一扇門.
去鼓勵和安慰另外.
和我有同樣問題的人.
當我們的信心.
植根在神的話的時候.
這一份信心的力量.
是遠遠超過我們所想像的範圍.
這位大神的故事.
其實是我們每一個人的故事.
他在痛苦裡面出發.
因為他無路可走.
結果他爭著去找耶穌.
但最後他竟然信服耶穌的說話.
空手回去.
因為他選擇相信.
結果他的兒子真的活著.
他的全家都信了.
信心的建立.
從來都不是一個輕鬆的過程.
他是可以在痛苦裡面萌芽生長.
在親自尋求裡面成長.
也更加在神話語裡面摘土根基.
今天倘若你也在經歷痛苦.
感覺很難受.
也許你的信仰一直都是二手.
一直只是靠著.
其他人的經歷來撐著.
又或者你對神的說話.
一直都不是很信的.
有疑慮的.
不知道可不可以放心.
來去交托的.
希望今天.
我們可以從這位爸爸身上.

$^{881}$學習得到怎樣.
在主耶穌的裡面.
建立到真正的信心.
我們一起禱告.
從父親的經歷.
我們去學習信心.
主我們知道.
信心這個功課.
不是容易學.
因為我們人生裡面.
有太多的關口.
可以令到我們沒有信心.
所以我們需要.
從聖經裡面.
來去學習.
在這個爸爸的身上去學習.
求你將你的話語.
放在我們的心裡面.
聽我們禱告.
奉耶穌基督恩主的命求.
Amen.
\newpage



\section{士師記 3:8-31}
\label{sec:032164}
\textbf{士師記析讀 (5) 三師匯}~~~~梁永善牧師
\newline
\newline
連結: \href{http://www.church.com.hk/acms/content.asp?site=cdc\&op=show\&type=product\&code=032164\&layout=sermon}{\texttt{032164}} ~~~~ 語音日期: 2026-05-13
\newline
\newline
\hyperref[sec:032163]{\small{< < < PREV SERMON < < <}}
~
\hyperref[sec:index_chronic]{\small{[返順時目]}}
~
\hyperref[sec:032165]{\small{> > > NEXT SERMON > > >}}
\newline
\newline
士師記 3:8-31
\newline
\begin{longtable}{cl}
\hline
\hline
章節 & 經文 (和合本修訂版)\\
\hline
3:8 & \begin{tabularx}{0.7\textwidth}{X} 所以耶和華的怒氣向以色列發作，把他們交給美索不達米亞王古珊‧利薩田的手中。以色列人服事古珊‧利薩田八年。 \end{tabularx} \\ \\ \relax
3:9 & \begin{tabularx}{0.7\textwidth}{X} 以色列人呼求耶和華，耶和華就為以色列人興起一位拯救者來救他們，就是迦勒的弟弟基納斯的兒子俄陀聶。 \end{tabularx} \\ \\ \relax
3:10 & \begin{tabularx}{0.7\textwidth}{X} 耶和華的靈降在他身上，他就作了以色列的士師。他出去爭戰，耶和華將亞蘭王古珊‧利薩田交在他手中，他的手戰勝了古珊‧利薩田。 \end{tabularx} \\ \\ \relax
3:11 & \begin{tabularx}{0.7\textwidth}{X} 於是這地太平四十年。基納斯的兒子俄陀聶死了。 \end{tabularx} \\ \\ \relax
3:12 & \begin{tabularx}{0.7\textwidth}{X} 以色列人又行耶和華眼中看為惡的事。耶和華使摩押王伊磯倫強大，攻擊以色列，因為他們行耶和華眼中看為惡的事。 \end{tabularx} \\ \\ \relax
3:13 & \begin{tabularx}{0.7\textwidth}{X} 伊磯倫召集亞捫人和亞瑪力人到他那裡，他就去攻打以色列，佔據了棕樹城。 \end{tabularx} \\ \\ \relax
3:14 & \begin{tabularx}{0.7\textwidth}{X} 於是以色列人服事摩押王伊磯倫十八年。 \end{tabularx} \\ \\ \relax
3:15 & \begin{tabularx}{0.7\textwidth}{X} 以色列人呼求耶和華，耶和華就為他們興起一位拯救者，便雅憫人基拉的兒子以笏，他是個慣用左手的人。以色列人託他送禮物給摩押王伊磯倫。 \end{tabularx} \\ \\ \relax
3:16 & \begin{tabularx}{0.7\textwidth}{X} 以笏打造了一把兩刃的劍，長一短肘，綁在右腿上衣服裡面。 \end{tabularx} \\ \\ \relax
3:17 & \begin{tabularx}{0.7\textwidth}{X} 他把禮物獻給摩押王伊磯倫。伊磯倫是個很肥胖的人。 \end{tabularx} \\ \\ \relax
3:18 & \begin{tabularx}{0.7\textwidth}{X} 以笏獻完禮物的時候，就把抬禮物的人送走。 \end{tabularx} \\ \\ \relax
3:19 & \begin{tabularx}{0.7\textwidth}{X} 但他自己卻從靠近吉甲的雕像那裡轉回來，說：「王啊，我有一件機密的事要奏告你。」王說：「迴避吧！」於是所有侍立在他左右的人都退去了。 \end{tabularx} \\ \\ \relax
3:20 & \begin{tabularx}{0.7\textwidth}{X} 以笏來到王那裡，那時他獨自一人坐在陰涼的頂樓。以笏說：「我有神的話向你報告。」王就從座位上站起來。 \end{tabularx} \\ \\ \relax
3:21 & \begin{tabularx}{0.7\textwidth}{X} 以笏伸出左手，從右腿上拔出劍來，刺入王的肚腹。 \end{tabularx} \\ \\ \relax
3:22 & \begin{tabularx}{0.7\textwidth}{X} 劍柄連同劍刃都刺進去了，肥肉夾住了劍刃。他沒有把劍從王的肚腹拔出來，糞便就流出來了。 \end{tabularx} \\ \\ \relax
3:23 & \begin{tabularx}{0.7\textwidth}{X} 以笏出到門廊，把王關在樓門裡面，就上了鎖。 \end{tabularx} \\ \\ \relax
3:24 & \begin{tabularx}{0.7\textwidth}{X} 以笏出來之後，王的僕人就來了。他們觀看，看哪，樓門鎖住，就說：「他必是在陰涼的房間裡大解。」 \end{tabularx} \\ \\ \relax
3:25 & \begin{tabularx}{0.7\textwidth}{X} 他們等得不耐煩，看哪，樓門仍然不開，就拿鑰匙打開樓門，看哪，他們的主人已經倒在地上死了。 \end{tabularx} \\ \\ \relax
3:26 & \begin{tabularx}{0.7\textwidth}{X} 他們耽延的時候，以笏就逃跑了。他經過雕像那裡，逃到西伊拉。 \end{tabularx} \\ \\ \relax
3:27 & \begin{tabularx}{0.7\textwidth}{X} 他到了那裡，就在以法蓮山區吹角。以色列人跟隨他從山區下來，他在他們前面引路， \end{tabularx} \\ \\ \relax
3:28 & \begin{tabularx}{0.7\textwidth}{X} 對他們說：「緊跟著我！因為耶和華已經把你們的仇敵摩押交在你們手中。」於是他們跟著他下去，佔據了摩押對面約旦河的渡口，不准一人過去。 \end{tabularx} \\ \\ \relax
3:29 & \begin{tabularx}{0.7\textwidth}{X} 那時，他們擊殺了約一萬摩押人，都是強壯的勇士，連一個也沒有逃脫。 \end{tabularx} \\ \\ \relax
3:30 & \begin{tabularx}{0.7\textwidth}{X} 那日，摩押在以色列手下制伏了。於是這地太平八十年。 \end{tabularx} \\ \\ \relax
3:31 & \begin{tabularx}{0.7\textwidth}{X} 以笏之後，有亞拿的兒子珊迦，他用趕牛的棍子打死六百非利士人。他也拯救了以色列。 \end{tabularx} \\ \\
[1ex]
\hline
\hline
\end{longtable}
$^{1}$今晚和大家去看三思匯.
我們先看回在《時事記》裡面.
給我們看到一個循環.
這個循環是很悲哀的.
就是他們國中太平.
但是接著他們就放蕩了.
接著又會失敗了 然後他們犯罪.
結果神就融入外邦來統治他們.
他們向神來有呼求.
神就興起時事來幫他們.
我們看看在《時事記》的第二章的第七節裡面.
就記載了這件事.
我們一起來看看.
我請Yvonne跟我們讀出來 好嗎?.
約書亞在世和約書亞死後.
那些見耶和華為以色列人所行大事的長老.
還在的時候 百姓都事奉耶和華.
你聽不聽得到?很寶貴.
約書亞在世和他剛剛死了之後.
但是那個世代那些未死的人.
一去經歷到神的作為.
他們就去事奉耶和華.
你不要忘記了.
進去迦南地的人除了約書亞和迦納之外.
基本都是大概四十歲以外的.
有些可能是五,六十歲.
為什麼呢?因為第一代出埃及的人.
他們不斷埋怨神缺乏信心.
所以那一代出埃及的成年人.
全部死在曠野.
所以進去的時候是第二代.
我們很相信出埃及第一代的小孩子.
即是他們所生的孩子.
可能是十多歲.
那些都不用被神帶入去剿滅.
即是說進迦南地的人.
可能是四十歲或者到六十歲多些.
那時候他們經歷到.
神怎樣在曠野祝福他們.
雲柱,火柱帶去護衛他們.

$^{41}$他們有嗎哪來去享壽.
亦有鵪鶉經過摩西擊打盤石出水.
所以他們很深體會到神的能力.
他們持守他們的信仰.
樂於事共神.
但是一個很大的可惜.
就是他們沒有將他們的信仰.
很好的來跟下一代分享.
以致這個信仰不能夠承傳下去.
這才是悲哀.
我們看第二章第十節.
那世代的人也都歸了自己的列祖.
後來有別的世代興起.
不知道耶和華也不知道.
耶和華為以色列人所行的事.
這句說話提醒我們.
你看到後世代的人.
已經不知道耶和華為以色列人所做的事.
因為可能他們沒有好好地去教育他們.
講解,分享讓這個信仰得以承傳.
我不斷在這幾年裡面都很強調.
我們重視的不要只是兒女的學業.
我們更加重要的是他們的信仰.
我經常說我們不能夠看到我們的兒女.
甚至我們的孫子只是一世的.
我們看不到他們最後的成長.
因為有一天我們可能會走了.
但如果我們能夠從小培育他們認識神的時候.
你就要放心了.
因為是由神看著他們.
不是我們.
你留意到他們受迦南宗教影響.
敬拜迦南不同的神.
這就離棄耶和華.
所以引致到神的憤怒.
我們看看第二章第十四節.
「耶和華的怒氣向以色列人發作.
就把他們交在搶奪他們的人手中.
又將他們賦予四位仇敵的手中.
甚至他們的仇敵別面前再不能站立得住」.

$^{81}$原來他們不斷犯錯.
所以神勸他們都不聽的時候.
神就容讓他們被外敵入侵.
甚至是神將他們交在他們的敵人手中.
搶奪他們的財物,農作物.
甚至管轄他們.
他們沒有能力去反抗.
甚至他們想走.
但坦白說.
走到甚麼地方呢?.
沒有神的同在.
怎會有安全呢?.
正如詩篇139那篇很清楚地說.
我們怎樣可以去逃避神呢?.
很寶貴的一首詩篇.
甚至我們之前曾經播放過這首詩篇改編的歌.
我們看看第七到第十二節.
「我往哪裡去逃 躲避你的靈.
我往哪裡逃 躲避你的面.
我若升到天上 你在哪裡.
我若在陰間下塔 你也在哪裡.
我若展開清晨的翅膀 飛到海極居住.
就是在哪裡 你的手必引導我.
你的右手也必扶持我.
我若說黑暗必定遮蔽我.
我周圍的亮光必成為黑夜.
黑暗也不能遮蔽 我是你不見.
黑夜卻如白晝發亮.
黑暗和光明 在你看到時一樣」.
這裡很清楚地指出.
我們其實沒有辦法去逃避神的.
你打算走 希望去一個和平或平安的地方.
如果沒有神同在的時候.
其實不會有平安的地方.
當他們面對痛苦的時候.
他們向神來有呼求.
神的憐憫就發動了.
就興起時詩拯救他們.
不是時詩有能力去幫助他們.
是神透過時詩去幫助他們.

$^{121}$所以時詩在世一切的日子.
耶和華就拯救他們脫離逃敵的手.
他們因為受欺壓受到困擾.
他們哀聲來呼求神的時候.
所以神後悔了.
神就來答應了.
時詩第2章第18節.
「耶和華為他們興起時詩.
就與那時詩同在.
時詩在世的一切日子.
耶和華拯救他們脫離逃敵的手.
他們因受欺壓擾害.
就哀聲歎氣.
所以耶和華後悔了」.
神應允了他們的哀求禱告.
這裡提到神的後悔.
我們經常這樣說.
神是無所不知的.
怎麼會有後悔呢?.
不是人嘛.
有時候我們的父母都說.
「哎呀 我早知道就不生你了」.
這麼激氣 是不是?.
鼎爺可能也有這樣的感受 是不是?.
好了 但這個神的後悔.
並不是指神不知道這樣.
而是在感受上感到傷痛.
當時詩拯救他們脫離了囚敵.
那些日子他們就歸回平靜了.
但是當災難一過去.
他們又再犯錯.
而且他們犯的錯是越來越嚴重.
而他們比列祖所做的更甚.
去事奉靠拜別神.
總是不斷絕夾來還感惡行.
好 我們看一看《釋述2:19》.
「自史詩死後 他們就轉去行惡.
比他們列祖更甚.
去事奉靠拜別神.
總不斷絕還感的惡行」.

$^{161}$記住 罪只會越犯越多.
越犯越嚴重.
你不會說因為我犯了一次罪.
有些好處.
好了 我就會收手了.
而且會越來越頻密.
如果不悔改的時候.
我們所面對的苦楚.
自然就會越來越大.
這也是今天會不會有一些信徒的經歷呢?.
有時候我們離開神.
好了 後來神拯救.
但我們繼續.
後來沒多久又再犯錯的時候.
終於有一天你不能夠回頭.
因為罪的那種影響力.
捆綁力是非常大的.
所以神就容讓那些外族.
成為他們審判的工具.
第二章第十九到二十三節.
於是耶和華的怒氣向以色列人發作.
他說:因者問違背我吩咐他們列祖所守的約.
不停從我的話.
所以若輸臘死的時候所剩下的各族.
我必不再從他們面前說出.
為要藉此試驗以色列人.
看他們肯照他們列祖謹守.
遵行我的道不幸.
這樣耶和華樓下各族.
不將他們速速說出.
也沒有加賦若輸臘的手.
有沒有留意到.
是神要留下一些試驗給他們.
今天神有時候都要我們來做一些測試.
希望我們會成功.
好了 樓下是哪幾族呢?.
我們看士師記第三章第一到第四節.
耶和華樓下這幾族為要試驗.
那不曾知道與加林爭戰之時的以色列人.
好叫以色列的後代.

$^{201}$又知道又學習未曾曉得的戰事.
所留下的就是非利士的五個首領.
和一切迦南人,西頓人.
並住尼巴嫩山的希末人.
從巴力克門山直到哈馬口.
樓下這幾族為要試驗.
若以色列人知道他們肯聽從.
耶和華藉摩西吩咐他們列祖的大命不幸.
這就是神特意容讓有幾族留下來.
讓他們有挑戰,讓他們有經歷.
因為這一代的人.
入了加林地的人可能沒有經歷過以往的經歷.
所以不少的以色列人.
沒有在約書指派他們的地方居住.
因為他們覺得打不贏.
相反地反而住在外邦人的中間.
我們看看士師記第三章第五節.
「以色列人竟住在迦南人,赫人,阿摩利人,.
比利時人,希末人,耶布斯人中間」.
耶布斯人是「竟」字.
以色列人竟住在迦南人,赫人這些地方裡面.
為甚麼呢?原來是很現實的.
因為迦南地的人.
他們的文化較為長遠.
他們在那裡住了很久.
即是說他們的地方已經發展得不錯.
然後他們有很好的設備.
或者城市較為繁榮.
於是這些人打不贏再去那裡住.
正如我們一過來英國的時候.
特別是在香港來的那些.
你不會走到一個很偏遠的地方.
杳無人煙的地方.
你一定會怎樣呢?.
可能先住在一些市中心.
即是接近我們那個叫伯奴.
有一個地方我經常說像甚麼呢?.
就是叫沙田瑞和園.
在火車站附近來的.
那些樓的形式就像瑞和園那樣.

$^{241}$我當年即是幾年前去探會友.
同一棟樓裡面有五個家庭.
我就不需要去很遠.
逐個逐個家庭.
後來有三個家庭.
直接集中在一起見面.
我們剛剛到的時候.
肯定住在最旺的地方.
好了 住定了.
然後再找找哪些地方可能遠一點的.
但已經有車的時候就不再有問題.
但剛剛來的時候沒有車.
即是說同樣這個道理.
以色列人他們進去迦南地的時候.
他們打五贏仗的時候.
這些地方這麼繁榮的時候.
那沒有問題.
那就不如住在這裡當中.
但一住在這裡當中的時候.
就很容易被他們的文化感染.
好了 不單單這樣.
他甚至讓他們的兒女和他們通婚.
我們看看《士師記》第三章的第六節.
取他們的女兒為妻 將自己的女兒.
嫁給他們的兒子 並事奉他們的神.
有沒有留意到他們竟然來跟外族人通婚.
我們要搞清楚.
其實外族人通婚是沒有什麼問題的.
《聖經》沒有阻止我們跟外族人通婚.
最重要就是這些外族人的信仰是怎樣.
很簡單舉一個例子.
拉合 拉合不是希伯來人.
但是他是外邦人.
還有路德 路德是摩訶人.
但這兩個人都成為什麼呢?.
成為耶穌的先祖.
提摩泰 他是混血兒.
他的父親是什麼?.
他的父親是歐洲的人.
甚至在建聖殿的時候.

$^{281}$也有一些混血兒.
他們是被神利用.
去做會幕和做聖殿的工作.
所以其實不是血緣或者種族問題.
是這些種族的信仰是怎樣.
我們有不少人現在身在外地.
甚至我現在認識的.
已經來了這裡三,四十年的老華僑.
他們的子女都是怎樣呢?.
都是跟當地人 即是外國人結婚的.
他們的孫子都是混血兒.
因為在這個文化底下.
很難去阻止你的兒女跟外國人交往.
譬如我的孫子讀了一間學校.
同班只有一個是香港人.
全部都是外籍人士.
所以他們一定會這樣.
當他長大一點的時候.
可能他所認識的男朋友或女朋友.
都是外籍人士.
那是不是不可以結婚呢?.
不是 最重要的是他們的信仰是怎樣.
這個是很重要的.
所以我記得我一講.
事事記是我們過去所忽略的.
是很著緊在家庭的問題.
我們要留意.
我們不要以為他們拜的那些假神.
就好像我們父母拜菩薩一樣.
很簡單 在香港來的電影節目.
可能你的父母會拜觀音.
初三會去七宮廟 去投籤.
其餘的時間對他生命沒有什麼影響.
拜菩薩不是說對嗎?.
那個影響性不是那麼大.
但是在迦南宗教就不同了.
除了建築物有神壇之外.
巴黎和亞瑟拉是生殖之神.
他們的名字是複數的.
即是有很多種類型的神.

$^{321}$他們的敬拜方式是非常淫亂.
所以他們敬拜神之後.
就會男男女女發生關係.
因為這樣才是土神的喜悅.
令到宗族或者家族很混亂.
可能你生了小孩都不是你經手.
同時他們要放縱.
所以神特別地阻止他們.
不可以跟他們這群戒嘔症的人結盟.
他們敬拜其他的神.
即是表示他們失去了對耶和華神的忠誠.
所以神也沒有義務對以色列人持守忠誠.
這就是守約和違約的意義.
一方面違約.
另一方面不需要承擔約裡的條款.
這是很重要的.
我們看了這個悲哀的循環.
然後我們看《時詩記》第三章七到三一節.
記載了三個時詩的事情.
所以我們叫三思會.
這三個時詩的記載較為簡單.
看起來也不錯.
原來基督徒講藉此就引到我們較熟悉的時詩.
亦是較多人討論的時詩.
但亦是越來越差的時詩.
下星期我們會看底波拉.
好 看第一個時詩.
第一個時詩叫《俄陀列》.
我們看第三章第七到第九節.
前面就講到為何他們會面對苦境.
因為他們犯罪 忘記神.
神就容讓他們面對米索波大米王.
古生尼撒田.
他們受苦百年然後向神呼求.
神就興起一個人.
這個人就叫做尼撒田.
尼撒田有兩個名字.
這個人就叫做俄陀列.
根據第一章十一到十三節記載.
俄陀列是迦勒的兄弟.

$^{361}$基勒斯的兒子.
我們曾經看過.
在上一次迦勒說.
誰可以來攻佔那些城.
我就會將我的女兒亞撒給他為妻.
結果俄陀列就成功了.
所以我們可以推想到.
俄陀列做士師的時候.
若舒死了不久.
因為迦勒仍然存在.
迦勒和若舒是同時代的人.
這個時候.
他似乎是士師裡最好最值得信賴的一位.
因為作者將他和信心的偉人.
專一跟隨神.
不敢按日再闖高峰的迦勒.
這個並排我上一次看迦勒的時候.
我自己也很深感受.
我也期望自己像迦勒一樣.
我們看到征服迦南地除了迦勒之外.
俄陀列是屬於第一位有領導能力的領袖.
第一章十三節很簡單.
只是尼達說到他奪取了那個城.
這個城是激烈西忽的城市.
好像很簡單.
但其實很肯定是一個驚心動魄的戰爭.
我曾經說過一篇信息.
叫《無悔當初》.
大衛幫助基伊拉城和菲律賓人打仗.
聖經也用一節經文說.
大衛打贏了菲律賓人.
拿了他們的物品.
還解救了基伊拉城.
只是一句.
你記住聖經的焦點從來不是放在戰爭裡面.
所以往往會輕輕帶過.
但不代表這場戰爭容易處理.
他在俄陀列列在士師的第一位.
當然絕對不是因為時間的問題.
但因為他是最好的士師.

$^{401}$在他之後的士師來說.
是一代不如一代.
最後記載的士師是參宣還是不知所謂.
但諷刺的是.
這位偉大的士師.
這麼優秀的領袖.
聖經很少用蝙蝠去描繪他.
因為士師的樣子我們看到.
當時社會是怎樣違背了神.
人民任意而行才有這樣的現象.
我們看看以色列人在埃及.
慰勞四百多年.
但可悲的再一次成為另一個王的奴隸.
我們再看剛才所看的第八節.
將會交在米索波大米王古山利薩田的手中.
以色列人服侍古山利薩田八年.
在希伯來語裡.
米索波大米這個名字具有.
古薩相配邪惡者的意思.
即利薩田也可能是一個.
附加給古薩的別名.
他是以邪惡來殘忍出名的.
這個王就是來自米索波大米.
西北面的亞蘭地區.
但我們不是很肯定這個王的身份.
但卻可以肯定.
他不是來自迦南附近的區域.
即神容許一位來自遠方的國王.
對他的子民施予審判.
你可以看到神是多憤怒.
但神是更加渴望子民可以回心轉意.
受苦之後出現這個循環.
第九節.
以色列人呼求耶和華的時候.
耶和華就為他們興起一位拯救者救他們.
就是迦納兄弟基納斯的兒子王托列.
以色列人向神發出這個呼求.
他們痛苦的時候.
這個呼聲達到神的耳中.
神就憐憫他們.

$^{441}$他們終於發覺到所謂的假神不能幫助他們.
他們就轉向神呼求.
當子民在繁榮盛世的時候.
事情順利的時候就沒有注視到神.
但神就使用苦難提醒他們.
不可以忘記自己的身份和地位.
更加不可以忘記仰望祂.
接著我們看看為何祂會贏呢?.
第十節.
耶和華的靈降在他身上.
他就作了以色列的史詩.
出去征戰.
耶和華將米索波大米王古撒列.
古山尼撒田交在他手中.
他變成了古山尼撒田.
這裡很清楚指出.
真正的領導者.
真正的戰勝者.
不是某一個史詩.
是神.
是神的靈降在他身上.
他才能夠成功.
這裡記住舊約裡面的靈.
未必等於三維一體的聖靈.
這個靈是一種能力賦予給他.
但是三維一體的聖靈是有違格的.
他給人有力量.
這個力量驅使人有能力去面對更大的事.
聖經裡面沒有詳細的去記述.
阿託列怎樣來打贏這場仗.
我們都不便推測.
總而言之就是神給他有力量.
就戰勝了這個邪惡的力量.
解救了這班子民.
就擺脫了古山尼撒田的核制.
不再成為他們的奴隸.
這是第一個史詩.
就是阿託列.
第二個史詩就是以弗勒.
這個士師記載得特別多一點.

$^{481}$如果你真的看下去.
這段經文是很好笑的.
我們看看以色列人他們冥頑不靈.
這個時候並沒有吸取到痛苦的教訓.
又再一次踏上了自毀之路.
《士師記》三章第十二節.
以色列人又行耶和華眼中學為惡的事.
耶和華就使摩格王以機倫強勢攻擊以色列人.
有沒有留意到「又」字.
以色列人又行耶和華眼中學為討厭的事.
即是表示他們又重蹈覆轍.
又再繼續犯錯.
這段的記述其實很戲劇性.
用很多諷刺的筆法來描繪.
他是一個刺客.
不知道你有沒有來看過.
當年麗的件事有十大刺客.
很厲害.
但我最記得是王羽所做的.
即是大刺客.
穿白色的衣服.
今次聖經就用以弗來刺一個很肥的王.
伊基倫.
伊基這個字是代表了小牛的意思.
但這個王是非常肥.
即是神經那隻小肥牛.
他懲罰他們.
伊基倫和古山利薩田相比之下.
其實他的威脅性是少很多.
但被殺的過程.
我們看不到他是一個戰士.
作者藉著伊基倫.
王公欠缺的安全措施.
進一步表示到這個王其實是很愚昧.
但神竟然用一個可以很愚昧的小肥牛.
成為管理以色列人.
甚至管教以色列人的工具.
可見權柄存在神的手中.
一個肥胖的伊基倫欺壓他們.
對以色列人是一個極大的侮辱.

$^{521}$即是一個很厲害的人.
來對付你.
你就沒辦法去抗拒.
但一個震騰騰的阿婆.
走來都打贏你的時候.
你真是很受羞辱.
教生去解釋就是他是一個有權力的人.
因為希伯來文「肥胖」這個字.
不僅是一個人很肥 賤辱橫生.
更加說他生活得很強壯.
但我仍然相信他是一個很肥的人.
爾忽憑著他周詳的計劃.
靈敏的頭腦和有爆炸性的力量.
如何去對付爾忽.
即是對付伊基倫.
一會兒我們會看到.
他就拿一把利劍壓一壓進肚子裡.
插到裡面.
甚至連脂肪都穿透了.
所以這個武器是從一塊很幼的金屬造出來的.
是沒有劍柄的.
所以是即時可以插進去.
我在日本旅行.
幾個月之前.
我都見到一些這樣的武器.
但我忘記了拍照.
我想買 但很貴.
對我來說200多元買一條這麼幼的.
我打算回來給大家看.
但無論如何他的智勇.
和之後登場的巴拉機電參宣.
有很大的對比.
我們看看神如何來引用這個人.
這個人有什麼特點.
是《悉記》三章第十五節.
「以色列人呼求耶和華的時候.
耶和華就為他們興起為拯救者.
就是變亞敏人基拉的兒子爾忽.
他是左手便利的.
以色列人託他送禮物給摩訶王伊基倫.

$^{561}$這個爾忽原來很特別.
是用左手的.
我們再看看一些背景.
以色列人被擊敗的地方是中樹城.
中樹城很清楚.
指的是耶利哥城.
作者用一個很諷刺的手法指出.
他進入耶利迦南第一場戰勝的地方.
就是耶利哥城.
但現在在這個地方失敗.
原因是他們不信服離開神.
而且他們受苦的時日是一次比一次長.
因為他們犯的罪是一次比一次嚴重.
記不記得去服侍古山利薩田服侍了多少年.
是八年.
但現在去服侍伊基倫是多少年.
是十八年.
多了十年.
而且伊基倫比古山利薩田弱很多.
神讓他們在一個沒有能力的統治者之下受更多的苦.
藉此使他們有吸取教訓.
這個時候亦顯出.
管理他們的權柄.
其實不是在壓迫者的手上.
而是在神的手上.
讓我們一次又一次地看到.
他們不斷地有犯罪不信服神.
於是神的管教一次比一次嚴重.
一次比一次厲害.
當以色列人向神呼求神的念文又再發動.
於是興起以弗成為他們的拯救者.
以弗的名字顯示出.
他所處的處境其實是非常惡劣的.
因為他的名字希伯來文的含義是.
到底君王在哪裡.
他父母取這個名字.
就是將以色列人的沮喪表露出來.
神本身就是以色列人的王.
但以色列人卻說我們沒有王.
可見他們離開了神是非常遠.

$^{601}$結果就陷在黑暗的深淵之中.
以弗的名字成為他們無望的表示.
以弗是來自辯雅文字派.
有很特別的.
原來辯雅文這個字的意思是甚麼呢.
就是右手之指.
右手是甚麼呢.
是代表了權力的象徵.
右手的力比較大.
我早上有做運動.
今早去做運動.
我去舉啞鈴.
每邊都是舉十下.
但我發現我右手舉的時候.
是舉得久一點.
而且沒那麼累.
但左手.
那兩個啞鈴的重量是一樣的.
但我左手是比較累的.
因為我習慣用右手.
所以力量就在右手那裡.
辯雅文是右手之指.
是代表了權力.
但這個字派竟然出了一個人.
是用左手的.
那時候一般人.
大部份人你和我都是用右手.
左手之擔是一些衛生工作.
有些地方缺乏紙巾.
他們的時候.
他們左手的功用.
就是用來清潔的.
所以確有一手.
但如果那些國家的人.
和你用左手來握手.
就是表示侮辱你.
很多時候都侮辱的.
很簡單舉個例子.
台灣的國民黨.
現在的領導是女人.

$^{641}$他們就去國內.
縱然和習近平見面.
但你有沒有留意到.
一行十八人.
只是坐很普通的車去見習近平.
其實這種就是侮辱你.
自取其辱 將你矮化.
你有沒有看BBC.
就是剛剛早個兩個小時拍了.
當特朗普下飛機的時候.
嘩 多麼多人歡迎.
晚上還有很多小朋友.
歡迎 歡迎.
他坐的車是多長 是多大.
想起國民黨的一個所謂主席.
坐的普通車.
都是一個羞辱.
接待規格不同.
現在特朗普會和習近平行天壇.
封了.
但英國的首相.
給街上帶著你去走.
只是僱工.
所以留意到我想說.
如果一個人在一個國家用左手和另一手.
就是侮辱你.
左手是什麼呢.
便利那些人做什麼呢.
拿盾牌來擋住.
然後用右手去揮劍.
所以右手的防禦力.
就帶來了.
很多城市.
他們的城牆設計都是怎樣呢.
都是根據右手的習慣來建造防禦的坡度.
爬上去這些斜坡的方向.
習慣用左手攻擊.
如果你沒有盾牌的地方.
你左手拿盾牌.
他就專門在你右邊.

$^{681}$你攻上去的時候.
在右邊帶人去打你.
你就要較為難去擋.
所以你留意到.
用左手的人.
在這時候就反而變為有利.
他可以在極短時間內攻擊敵人.
這個行動其實不單單是爾忽個人的行動.
是以色列人忍受不了這個壓迫.
於是他們就要做一個刺殺伊基倫的計劃.
他藉著送禮物給王.
然後進行暗殺.
這個重要的任務就交給爾忽.
為什麼.
他們相信爾忽一定是很靈敏.
很有智慧.
能夠擔此重任.
你想想如果他刺殺失敗的時候.
他們會遭到報復.
他們會受到的壓力更加大.
所以這次的行動真是孤注一擲.
此可勝不可敗.
為什麼要選擇爾忽.
剛才我說除了他是一個很有能力的人外.
還有一樣是怎樣.
他是用左手的.
當他進攻檢查的時候.
那些人.
我想大家都去過出境.
當你要搜身.
那些人都會拍你兩邊的手.
然後拍你兩邊的大腿掃到腳下面.
同樣都是.
那時候是怎樣呢.
一個慣用左手的人.
當他進攻的時候.
他的武器一定會怎樣.
是放在他右膝的內側.
即是小腿那邊.
當他拍他的時候.

$^{721}$檢查他一定是檢查他兩邊.
但不會摸他的右邊腳的內側.
是不是.
那時候.
所以他用左手.
就方便他將利器.
放在他右腳的內側.
好 我們看看他帶了什麼武器.
士師記第三章第十六節.
以忽打了一把兩刃的劍.
長一手.
戴在右腿上衣服裡面.
好 戴在右腿上.
衣服裡面.
所以那些手要摸他.
搜身的時候.
就沒那麼容易搜到這裡.
就較為容易將武器運入宮中.
好 我們看看那個刺殺情況是怎樣.
第十八節到二十三節.
以忽獻完禮物便將台禮物的人打發走了.
自己卻從靠近吉格鑿石之地回來.
說:王啊 我有一件機密事就告你.
王說會被罷.
於是左右侍立的人都退去了.
以忽來到王面前.
王獨自一人坐在樑樓上.
以忽說:我奉神的命報告你一件事.
王就從座位上站起來.
以忽便伸左手.
從右腿上拔出劍來.
刺入王的肚腹.
連劍把都刺進去了.
劍把肥肉夾住.
他沒有從王的肚腹拔出來.
且穿通了頭身.
以忽就出到油廊.
將樓門盡都關鎖.
你有沒有留意到以忽其實很有智慧.
他首先將禮物呈獻給伊基倫.

$^{761}$接著就走了.
走了 你有沒有留意到.
這些禮物一定是很多 很重.
而且不止他一個人帶去.
是有些人親自帶去.
將所有幫手的人.
抬了這些禮物給王之後.
他們就走了.
離開了.
但是他自己也暫時離開了.
可能在一段很短的時間.
王一定是看那些是甚麼禮物.
其後他個人又再轉頭說.
我有些機密的事要告訴王.
記住.
王和守衛已經鬆懈下來.
因為第一次他進攻的時候.
搜他身.
可能搜不到甚麼出來.
甚至很可能.
他第一次未必會帶著武器.
但是聖經說其實已經有帶了.
然後再一次回去的時候.
你再回頭.
自然就會鬆懈了.
而伊基倫也毫無智慧.
因為當以忽說.
王 我有些秘密想告訴你.
他沒有這個警覺性.
竟然拆走了那些侍衛.
即是那些奧地格.
於是將自己暴露在以忽的前面.
以忽說.
我有些很重要的情報.
但不想大聲告訴你.
王叫你來.
我細聲說不提防.
於是就被以忽拿出刀插進去.
插進去的時候.
有沒有留意到那些描繪很有趣.

$^{801}$是從前面進入.
後面出.
是刺穿了他的肚腹.
他想拔出來.
拔不到.
因為被那些肥肉夾住.
是不是很特別的描繪呢.
然後他離開之後.
更加將城門鎖了.
最好笑.
那些守衛看到城門關了很久.
以為王在裡面大解.
所以沒有打擾他.
為何這麼久都沒有出來呢.
打開門的時候才知道他死了.
但有沒有留意到.
那扇門是在外面鎖的.
不是內鎖.
是以忽將門鎖了.
然後離開.
這班守衛都很愚蠢.
所以有時候.
愚蠢的王就有愚蠢的跟從者.
那些神博都是這麼愚蠢.
如果王是大解.
他沒理由在外面鎖門.
明不明白這個道理.
所以你留意到.
很可悲.
但這一段的記載.
後來你會留意到.
如果曾經有上我的識讀大衛就會知道.
當約翰回到王宮.
他知道當時挾天子義令之後.
的蘇羅元帥有求成求靠大衛.
他就妒忌.
然後他就馬上叫人追上去.
就叫他回來.
他說大衛王有事要跟你說.
這個元帥亞歷爾就不妨有詐.

$^{841}$因為見完大衛.
大衛對他很禮遇有加.
所以他打算約爹去跟他說.
然後約爹就說來來.
有事要跟你說.
就叫世澈上來.
然後箍著他的頸.
一刀插進去插死他.
一個元帥就這樣死了.
其實就是歷史在重複.
我們看看經文裡面的記載.
剛才我們看了二十三節.
二十三節就記載到.
以弗出了由郎.
將門樓盡刀關鎖.
好 二十四節到二十五節.
「以弗出來之後.
王的僕人到了.
看見門樓關鎖.
就說他必是在樓上大解.
他們等煩了.
見仍不開門.
就拿鎖匙開了.
不料他們的主人已死.
倒在地上」.
剛才已經說了.
沒有理由真的大解.
沒有理由在外面下鎖.
所以這些人真的很蠢.
這些細節有時我們會忽略了.
這段經文有兩次提及.
他經過的地方.
《士師記》第三章第十九節.
和三章二十六節.
「自己他從靠近吉甲.
鑿石之地回來說.
我有一件機密事就告你.
我說會被罷.
於是左右士立的人都退去了.
他們擔言的時候.

$^{881}$以以弗就逃跑了.
經過鑿石之地逃到西伊拉」.
有沒有留意到.
特別提到吉甲鑿石之地.
剛才不是說了.
顯然是離開了.
然後再回來.
但原來是從鑿石之地.
這個吉甲尼拉去見王.
以不殺王之後當然要走.
經過鑿石之地.
好像很簡單.
但原來這些地方.
都是迦南人拜偶像的地方.
即是有一個隱喻是甚麼呢.
就是這些偶像不能夠救基伊倫的性命.
這些偶像是無能的.
唯有耶和華才是真實.
接著我們看看.
這個結果帶來甚麼呢.
舌斯記第3章27到30節.
「到了就在以法連山地吹過.
以色列人隨著他下了山地.
他在前頭引路.
對他們說.
你們隨我來.
因為耶和華已經把你們的仇敵摩訶人.
交在你們手中.
於是他們跟著他下去.
把手約旦河的渡口.
不容摩訶人一人過去.
那時擊殺了摩訶人約有一萬.
於是強壯的勇士沒有一人逃脫.
這樣摩訶就被以色列人制服了.
過中太平八十年」.
當殺了伊基倫王之後.
他就馬上呼召以色列人向摩訶進攻.
所以是整體配合.
才能夠得勝.
不只是以法.

$^{921}$如果我們留意以法沒有將暗殺的成功歸於自己.
是我殺了他.
你留意到以法很清楚地說.
是神藉著他的手令他們得勝.
所以說到.
「你們隨我來.
因為耶和華已經把你們的仇敵摩訶人.
交在你們的手中」.
他很強調是神的作為.
所以看到神興起義忽.
作為以色列的領袖.
是有他自己特別地揀選.
這個領袖要很清楚地知道自己要付出.
但是成功與否.
是由神來掌管.
和神來將力量加給他的.
這是我們每一個事奉者.
都要有的心態.
我們要努力.
但是要仰望神.
要彼此配合.
這是永恆不變的真理.
在義忽的帶領下.
他就有八十年的平安.
我們看了兩個時詩.
「摩訶列」.
義忽第三個時詩.
只有一節經文.
就叫做「三加」.
這個時詩某個程度.
我們都不知道它是一個怎樣的時詩.
我們看看時詩結第三章第三十一節.
義忽之後有雅蘭的兒子三加.
他用港牛的棍子打死六百非利士人.
他也救了以色列人.
有沒有留意到.
聖經只是用了這一節的經文.
去講述三加的事.
它說他在義忽之後.
記住.

$^{961}$不是代表義忽死了之後才輪到他.
可能是什麼呢?.
義忽成為時詩之後.
他再參與一些其他以色列的戰役.
在底波拉的歌頌裡面.
我們見到他在底波拉的時期仍然活著.
所以義忽和三加.
可能是同一時間來作時詩.
最特別的是什麼呢?.
三加並不是一個以色列人的名字.
他可能是來自互利人的民族.
但經文說他是阿那的兒子.
阿那可能只是一個閒頭.
是當時敬拜一位戰爭和生育的女神.
叫做阿那.
如果一個人特別尊重這個女神.
就會將她的名銜加在他身上.
所以三加應該是來自一個完全外邦人的家庭.
但是這樣.
讓我們看到原來神都仍然可以用他.
來拯救以色列人.
這顯出神的挑選.
神的大愛是遍及每一個人.
不論他的種族.
不論他的膚色.
最重要的是他沒有心.
三加用那些鋼牛棒.
來打死六百個非利士人.
這是第一個使用人的力量.
去對抗非利士人的事蹟.
記不記得我說.
《士師記》的作者很有技巧.
原來三加對照最後一個的事詩.
如是就是參宣.
因為參宣都是憑自己一個人的力量.
去擊殺非利士人.
但兩者是一個強烈的對比.
三加雖然是一個外邦人.
但他從非利士人手中拯救以色列人.
但參宣雖然是以色列人.

$^{1001}$但他的言行和外邦人一樣.
從參宣身上我們看到.
以色列人比外邦人更加差.
甚至以色列人的領袖都展現了.
被外國人的文化,習俗,色彩所影響.
是很可悲的.
所以最後提醒大家.
我們要好好地著緊自己的身份.
不要被這些不好的信息污染了.
我們也要更加好好地教導我們的下一代.
讓他們對神有認識.
讓他們親身經歷神.
正如他們第一代還是約書亞年代.
他們去事奉耶和華.
但他們沒有好好地教育下一代.
結果令他們迷失.
第二方面我們需要有好的領袖.
但我們要認定真正的領導是我們的神.
有時我們陷入兩種極端.
有時甚麼都是神.
是神的祝福.
而輕忽了領袖所付出的努力和代價.
另一方面很可能是個人崇拜.
而輕忽了神的恩典.
俄台烈的故事結尾讓我們看到.
他戰勝了敵人.
國中太平四十年.
所以好的領袖.
他們自然有帶給我們好的生活.
但我們要確信.
好的領袖背後是來自神的.
我仍然再說.
神會幫助我們.
但人有這個責任.
在剛剛我放上去的.
我忘記了是永善論壇.
他講到有個人不知是否善意.
我相信是討論.
他講到甚麼都是有神預備的.
所以好像人沒有甚麼責任要承擔.

$^{1041}$我跟他說.
他講藥室.
他說我用了人的道德的看法.
去看一件事.
去解釋聖經.
譬如我問他.
藥室的哥哥.
去妒忌藥室.
波特弗的太太.
去誣衊藥室.
然後藥室要坐監.
他說不是.
就是因為這樣存故.
所以藥室就不要生哥哥的氣.
藥室成為宰相.
一切有神的安排.
我說波特弗的太太有沒有責任?.
藥室的哥哥有沒有責任?.
甚至我問.
神是安排了耶穌基督.
如果這個世界為我們死.
代替我們受這個罪的刑罰.
但猶太出賣耶穌.
你又說是神的安排.
猶太有沒有罪?.
猶太要不要承擔責任?.
所以我們要記住.
原來我們不要什麼都說是神.
又或者什麼都是人.
我們兩者要配合.
人都有他自己的責任.
我們一起來祈禱.
天父我們多謝你.
讓我們透過你的話語.
成為我們的提醒.
以色列人他們在平安穩妥的時候.
他們忘記你.
他們在苦難當中.
然後才呼求你.
神求你幫助我們.

$^{1081}$不是這樣去利用你.
不過我們更加見到.
有時候我們沒有好好教育.
我們的下一代的時候.
引來的傷痛.
超乎了我們所想.
求你幫助我們.
雖然有時候我們教導我們的兒女.
特別當我們已經長大的時候.
是更加難以跟他們分享.
但是叫我們不可以放棄.
我們已經有那個責任.
懇求你幫助我們.
叫我們懂得守望我們的下一代.
或者再下一代.
讓我們用我們的生命.
用我們的經歷.
用我們的分享.
都成為他們的激勵.
也懇求神你幫助我們.
讓我們知道我們要努力.
但不是只靠我們的努力.
而是靠著神你給我們的恩典.
願神你賜福給我們每一個人.
多謝你給弟兄姊妹有學無的心.
也求你祝福.
在YouTube來大陽去收聽.
收看的人.
都同樣祝福他們.
讓他們生命都有成長.
多謝你.
聽我們在面前的祈禱交托.
祈禱交托 奈奉主德性的名.
阿們.
\newpage



\section{哥林多後書 5:18-20}
\label{sec:032165}
\textbf{作基督的大使 (華)}~~~~葛國光牧師
\newline
\newline
連結: \href{http://www.church.com.hk/acms/content.asp?site=cdc\&op=show\&type=product\&code=032165\&layout=sermon}{\texttt{032165}} ~~~~ 語音日期: 2026-05-17
\newline
\newline
\hyperref[sec:032164]{\small{< < < PREV SERMON < < <}}
~
\hyperref[sec:index_chronic]{\small{[返順時目]}}
~
\hyperref[sec:032166]{\small{> > > NEXT SERMON > > >}}
\newline
\newline
哥林多後書 5:18-20
\newline
\begin{longtable}{cl}
\hline
\hline
章節 & 經文 (和合本修訂版)\\
\hline
5:18 & \begin{tabularx}{0.7\textwidth}{X} 一切都是出於神；他藉著基督使我們與他和好，又將勸人與他和好的使命賜給我們。 \end{tabularx} \\ \\ \relax
5:19 & \begin{tabularx}{0.7\textwidth}{X} 這就是：神在基督裡使世人與自己和好，不將他們的過犯歸到他們身上，並且將這和好的信息託付了我們。 \end{tabularx} \\ \\ \relax
5:20 & \begin{tabularx}{0.7\textwidth}{X} 所以，我們作基督的特使，就好像神藉我們勸你們一般。我們替基督求你們，與神和好吧！ \end{tabularx} \\ \\
[1ex]
\hline
\hline
\end{longtable}
$^{1}$這個世界.
真的是需要愛.
我們看見打開報紙.
到處都是紛爭.
在基督的裡面我們可以建立一個彼此相愛.
同心合一的教會.
在基督裡面我們可以成為一個平安的.
充滿和諧的一個地方.
這是教會在這個世界當中.
一個非常重要的角色.
剛才我們讀的這段聖經裡面.
其實如果我們看哥林多後書.
前面一節大家會比較熟悉.
若有人在基督裡.
若有人在基督裡他就是新造的人.
就是已過一切都變成新的了.
在之前保羅在講什麼呢.
在十四節開始保羅講了說.
原來基督的愛激勵我們.
叫我們如今活著不再為自己活.
來為那位替我們死而復活的主而活.
所以保羅講到原來基督的愛激勵我們.
使我們活著為他而活.
接著他就說若有人在基督裡.
怎麼樣.
他就是新造的人.
就是已過一切都變成新的了.
然後接著他就講這一段話.
一切都是出於神.
他藉著基督使我們與他怎麼樣.
和好.
然後呢.
他說又將確認與他和好的職分怎麼樣.
賜給我們.
所以在這段聖經裡面一直在強調一件事情.
就是與神怎麼樣.
和好.
這是神的一個旨意.
這是神最重要的一個旨意.
就是神要使我們與他和好.

$^{41}$並且他也要藉著我們.
把這個確認與神和好的這一個職分.
他託付給我們.
這是今天我跟大家分享的.
首先我們來看一件事情.
什麼叫做和好.
和好的意思.
這個字在希臘文是兩個字合在一起的.
這個apō.
是front的意思.
從一個.
那麼接著下面那句.
這一句話呢.
它是改變.
或者是exchange.
交換.
他在講什麼呢.
他在講.
一個.
他強調的是一個關係的改變.
所以和好是一個關係的改變.
人跟神之間的關係.
人跟人之間的關係.
都需要和好.
所以這個和好如果直接把它翻譯就是.
離開本來敵對的一個情況.
回到原來和諧的關係.
親愛的獨子妹.
我們跟神之間.
原本是一個敵對的情況.
在我們不認識神的時候.
想想看.
我們跟神的關係.
是處在一個不認識祂.
甚至很多時候抵擋祂的一個情況.
所以神的心意是要.
使我們在基督的裡面與神怎麼樣.
和好.
與神和好.
第一件事情.

$^{81}$也就是說.
如果今天我們還不認識這位耶穌基督.
藉著耶穌基督在實價上面.
祂所做的工作.
祂要使世界上許多還不認識祂的人.
能夠跟祂和好.
亞當夏娃離開神.
使得人跟神之間的關係斷絕了.
我在這個門訓班跟大家講.
怎麼樣傳福音.
我不曉得門訓班畢業的你有沒有還記得.
是不是.
把.
人跟神之間的關係為什麼會斷絕.
為什麼.
因為什麼.
因為罪.
所以換句話說.
人跟神之間本來的關係是怎麼樣.
怎麼好法.
我們現在感覺不到人跟神的關係很好.
我們看到的到處都是不好的現象.
但是丁字輩.
上帝在創造人的時候.
亞當跟夏娃本來跟神之間是可以直接來往的.
亞當夏娃來到神的面前.
她不需要做任何的遮掩.
對不對.
她完全沒有任何的懼怕.
就好像一個孩子來到爸爸媽媽面前的時候.
她完全沒有任何的懼怕.
她可以很釋放的來到神的面前.
但是如果你碰到一個爸爸媽媽非常兇惡的.
常常會打孩子的.
請問孩子碰到爸爸媽媽會如何感覺.
有沒有這個.
家裡面爸爸一回來馬上全家都安靜了.
有沒有.
甚至還沒有回來的時候.
這個媽媽會說:爸爸快回來了.

$^{121}$有沒有.
所以這意思就是爸爸是什麼.
爸爸回來了會教訓你的是不是.
所以千萬你不要這樣講.
因為牧師來了.
就是我們跟神之間的關係本來是好的.
人神創造人的時候.
他看見一切甚至宇宙萬物都是好的.
他看著他所造的.
他說好 好 好.
他一直在說好.
他看著都是美好.
一切都是美好.
但是今天你看見.
世界不僅僅是人跟神之間的關係不好.
而且你看見整個宇宙萬物.
都處在一個很不好的情況之下.
但是神說他藉著基督.
要使我們與他和好.
所以這是唯一的道路.
只有藉著基督我們可以跟神和好.
沒有別的路.
沒有別的方法.
藉著基督神要使我們跟他之間重新和好.
就是唄.
所以這是第一個.
那第二個.
不僅僅是信耶穌的時候.
我們跟他要重新和好.
我們因著耶穌基督我們和好了.
對不對.
好 今天你信耶穌.
我們信耶穌很多年了.
請問你跟神的關係好嗎.
問一下你旁邊.
你跟你跟神的關係好嗎.
好嗎.
什麼叫關係好.
什麼關係好.
一年見一次面.

$^{161}$關係會好嗎.
一個禮拜見一次面.
關係會好嗎.
如果你一個禮拜才跟上帝見一次面.
關係會好嗎.
我很懷疑.
是不是.
所以當我們說與神和好的時候.
信耶穌是一個開始.
那不是一個結束.
當我們說要與神和好.
神的心意是要我們跟他之間的關係.
恢復到這樣起初創造的時候.
上帝跟亞當夏娃.
他真的是可以說.
他們中間沒有任何的組合.
他們可以面對面的.
他們可以彼此之間很自然的交通.
這是神要跟我們之間的關係.
他不是因為今天基督教規定說.
禮拜天要來主日聚會.
所以你信了耶穌之後.
大概基督徒都會告訴你說.
你禮拜天要來聚會.
這已經算是OK的了.
但是進一步你會發現.
如果有一些人會告訴你說.
你信耶穌之後.
不僅要禮拜天聚會.
還要來參加小組聚會.
有沒有.
好不好.
但是現在又怎麼樣.
我們基督徒跟神的關係.
只有一個禮拜一次兩次嗎.
我們跟他之間的關係.
應該是怎麼樣.
主耶穌說.
祂是葡萄樹.
我們是枝子.

$^{201}$對不對.
你們要藏在我的裡面.
我就藏在你們的裡面.
所以祂要我們建立跟祂是一個.
住在祂裡面的關係.
什麼叫住在祂裡面.
你今天早上住在祂裡面嗎.
當你一起床的時候.
你住在祂裡面嗎.
是嗎.
當你早上起來醒過來的時候.
你第一件事情是怎麼樣.
先拿手機.
先起來的時候.
第一件事情說.
主啊.
我愛你.
有嗎.
有.
丟著沒.
這不是天方夜譚.
你最愛的是誰.
是你兒子嗎.
是你女兒嗎.
是葛牧師啊.
這叫甜言蜜語.
我要相信你.
什麼時候再吃你的牛肉麵.
神的心意要我們與祂和好.
丟著沒.
我們一生我們所追求的.
就是希望我們能夠跟主.
跟主之間的關係.
越來越親近.
這是為什麼我們要主日聚會.
這是為什麼要禱告會.
這是為什麼我們要門徒訓練.
這是為什麼我們要小主聚會.
這一切的一切.
如果不是讓我們能夠跟神之間的關係更好.

$^{241}$你所做的這些浪費時間.
親愛的同志們.
想想看我們一天當中花多少時間.
在手機上面.
在電腦上面.
在與神無關的事情上面.
但是神祂的心意是要我們跟祂和好.
恢復到起初跟祂之間.
那樣子的一個美好的關係裡面.
不僅僅如此.
神祂也要我們與人和好.
我們來看看.
以弗所說第二章十四到十五節.
我們一起來讀.
來,預備,請.
因祂使我們和睦.
將兩下合而為一.
拆毀了中間隔斷的牆.
以自己的身體廢掉冤仇.
就是那既在律法上的規條.
為要將兩下藉著自己造成一個新人.
如此便成就了和睦.
以弗所說這一段聖經裡面.
講到說因祂使我們.
祂是誰啊.
祂是指的基督.
原文裡面就是說.
因祂使我們和睦是.
祂是我們的和睦.
He is our peace.
不是祂make us peace.
是祂就是和睦.
祂就是那個和平.
換句話說.
藉著祂我們跟神之間能夠和好.
藉著祂我們才能夠中間.
我們怎麼樣.
彼此和好.
我想我們今天上過恩愛婦孺期.
或者是你上過什麼一些.

$^{281}$家庭的一些輔導課程.
都聽過一個叫鐵三角吧.
聽過嗎.
沒聽過的記得一下.
什麼叫鐵三角.
三角懂嗎.
那麼這邊是夫這邊是妻.
那上面是什麼.
是神對不對.
所以這個鐵三角的意思是什麼.
當夫跟夫妻同時都跟神的關係接近的時候.
兩個人的關係是怎麼樣.
就越接近是不是.
就這麼簡單.
什麼時候當我們跟神的關係.
越來越接近的時候.
夫妻之間的關係也就會越來越接近.
所以這反映一件事情.
為什麼今天很多的夫妻吵吵鬧鬧.
不能夠彼此相愛呢.
是不是也反映一件事情.
因為很多的家庭裡面沒有神.
甚至可能基督徒的家庭裡面.
神在你們中間.
不是在做主.
基督是我家之什麼.
主不是.
因為丈夫說我是一家之主.
太太說什麼.
我是伯子.
所以你看這些歪理一大堆有沒有.
清楚什麼.
但是如果我們都有同一位主.
當夫妻都把主當作主的時候.
每一個都跟神之間的關係接近的時候.
當我們跟神的關係越接近的時候.
我們彼此之間的關係就越來越接近.
我這次到義大利的一個.
他們是全義大利的南部的.
他們有一個叫復興教會.

$^{321}$是他們義大利的一個.
算是你要稱它中派也可以.
他們都是溫州人.
所以他們是有總會制的.
這個復興教會他們大概有.
可能差不多三十間左右.
在全義大利.
所以羅馬也有.
拉布勒斯這些地方都有.
他們姐妹退休會就在拉布勒斯舉行.
拉布勒斯那邊的聚會.
有一天我聚會結束.
主日聚會剛結束.
馬上就有一個姐妹就告訴我說.
拉布勒斯你可不可以來為一個木道友禱告.
她就帶著我出去教會.
然後就找到那個木道友.
然後把她先生也找來了.
我一看那個太太就是愁眉苦臉.
當然這個姐妹很熱心希望我為她禱告.
但我就問她說你要為她禱告什麼.
她就說禱告她姓祖.
禱告她姓祖.
那麼我就代她姓耶穌.
可是那個姐妹又把她的先生帶過來.
說是先生剛姓祖.
也就是說要我為他們夫妻禱告.
這個兩個事情怎麼連在一起.
但是當時候聖靈在我心裡面就感動.
所以我就為他們禱告.
我說主啊.
求你讓他們夫妻兩個人同姓一位主.
以至於他們能夠在基督徒裡面彼此相愛.
很奇妙.
我在講的是什麼.
我本來不知道他們夫妻之間有問題.
後來是我這樣禱告之後.
這個太太好像就突然覺得就開始就很有意願.
要跟著我一起來禱告.
後來才知道.

$^{361}$我聽了一位帶我去為他們禱告的姐妹才知道.
夫妻兩個人已經是水火不容了.
兩個人幾乎可以說是叫同床異夢.
但是細節我就不講了.
幾乎可以說已經到達兩個人要分開的一個階段.
先生說剛姓的主.
但是他在跟我講話的時候.
講了很多一些東西.
都在講自己怎麼樣怎麼樣怎麼樣.
就在炫耀自己.
所以我想這個人大概姓耶穌.
可能還有很長的路要走.
所以那一天我就為他們夫妻兩個人禱告.
禱告完之後告訴他們一件事情.
我特別跟這個先生講.
我說你信了耶穌.
你一定要好好的愛你太太.
對不對.
聖經是不是這樣說.
意思就是說.
他一直在跟我講言下之意就是講他太太怎麼不好不好.
我不知道你們夫妻兩個人關係怎麼樣.
但是我先問一件事情.
那你愛她嗎.
你有沒有好好的愛她.
是不是.
就像基督愛教會.
為教會怎麼樣.
死.
為教會捨棄.
做弟兄的.
做丈夫的.
有沒有愛姐妹.
愛到一個地步為他捨棄.
有沒有.
當他們兩個人都同樣要信耶穌.
我要讓他們看見這是一個開始不是結束.
Amen.
同樣在座的弟兄姊妹.
你信主不知道多少年了.

$^{401}$請你永遠記得.
我們跟神的關係.
沒有一個結束的一天.
要一直地追求.
什麼時候我們跟神的關係好了.
我們夫妻之間.
我們家庭之間.
父母親跟兒女之間的關係.
才可能越來越好.
Amen.
這是一個根本的道路.
在彼此之間.
我們要和好.
保羅說他要使我們和睦.
藉著基督拆掉了我們中間隔斷的牆.
並且將兩下藉著自己造成一個新人.
這是耶穌基督已經做的事情.
接著呢.
我們要不要做什麼事情.
彼得前書三章十一節怎麼說.
在這裡說要尋求和睦怎麼樣.
要一心追趕.
所以我們的責任是什麼.
要一心追趕.
追趕的意思不是被動.
追趕的意思是主動.
但是你會發現我們.
我們華人大概通常都比較被動.
但是我們賺錢都很主動.
可是你會發現在我們人際關係裡面.
我們常常是被動的有沒有.
我們如果這些夫妻關係如果有了問題.
會主動要去尋求兩個人和好嗎.
通常我們都是.
要不然就熱戰.
要不然就冷戰.
除了熱戰跟冷戰沒有別的戰法.
是不是.
所以一輩子forever and ever.
從新竹的結婚的第一年.

$^{441}$一直到結婚的第五十年.
到最後也沒什麼好吵了對不對.
最後就變成是兩個人不吵了.
不吵的意思不是沒有問題.
是懶得跟你吵了.
這不是根本的路.
神要我們學習操練一心的追趕.
好那麼那我們需要操練.
從我們的.
從我們彼此的關係.
我們學習多說造就人的好話.
不要說批評論斷智者的話.
根據研究.
批評指責是人際關係.
特別是婚姻關係的殺手之一.
批評指責也會使人沮喪失去盼望.
一個孩子從小天天聽父母親批評他.
通常都沒有信心.
是不是.
想想看可能你自己的父母親.
常常跟你說哎呀這個孩子啊.
他這個沒有出息.
他是我們家讀書最不好的有沒有.
或者是說哎呀你看他長得這個樣.
將來怎麼能夠嫁出去呢.
這個很奇怪.
我們中國人好像很喜歡假謙卑.
我記得以前我這個考上.
以我們那個就是功課來講最好的學校.
在我們那個地方讀大學就已經不簡單了.
考上好學校那更少對不對.
他應該很高興啊對不對.
但是我媽媽在跟著我一起到菜市場.
我帶她去買菜的時候.
碰到鄰居就說哎呀.
聽說你兒子考上了好大學啊什麼什麼.
他怎麼說呢.
他說哎呀哎呀沒有了沒有了.
沒有什麼叫沒有啊.
有就是有嘛.

$^{481}$這沒有.
好像我們好像覺得說.
哎呀這叫假謙卑懂嗎.
現在就準備我們需要學習彼此鼓勵.
在家庭中如此.
夫妻之間常常彼此鼓勵.
所以每一次聚會一開始我要大家彼此說.
跟你在一起真好幹嘛.
為什麼為什麼要為什麼講這句話.
是不是.
是的我們在一起真好.
哎呀大聲一點嘛.
跟你在一起真好.
跟你旁邊說你真好.
你真棒.
你一輩子可能講的話.
聽到的話都不是這樣子.
所以你需要花更多的時間更多的力量.
操練學習說造就的話.
要把我們中國人那種說負面的話.
批評論斷的話否定的話.
這種的文化完全的剷除.
若有人在基督裡他就是怎麼樣.
心造的人.
就是一個你過去爸爸媽媽講話的口吻.
你不要再讓他受他影響.
你過去爸爸媽媽批評論斷你的否定你的.
你不要再把他帶到你的下一代去.
這個宗族我們不容許他繼續的傳承.
我們要傳承好的東西.
在基督的裡面一切都變成新的了.
我們說話行為說話要說恩言.
說造就人的好話.
造就人的好話原文是什麼意思呢.
造就什麼意思.
Build you up.
不是put you down.
給你旁邊再說一句Build her up的話好不好.
如果你不會說我舉一些例子給你聽.
你今天看起來好帥啊.

$^{521}$你今天看起來好美啊.
我真慶幸能夠嫁給你啊.
我真慶幸能夠娶到你啊.
下一輩子不是沒有下一輩子.
如果有下一輩子我還要娶你啊.
辛苦你了.
你回去可以把它多寫下來好不好.
然後對著鏡子練習.
講的時候要慢慢操練慢慢帶著微笑唱.
不要依附那個婉娘面孔.
不僅僅是讓我們的口說恩言.
讓我們有恩典的行動.
我昨天回來坐飛機坐了11個多小時.
因為我知道我到的時間是下午.
所以我就知道要避免時差的方法.
就是我在飛機上盡量不睡覺.
懂我意思嗎.
所以我在飛機上面就大概只瞇了一個小時當作午睡.
因為我從義大利出來的時候是早上.
我們是早上的飛機.
所以我就當作是午睡的時間.
睡了一小時.
其他時間我就警醒預備講章.
當我回到洛杉磯的時候.
我們喜樂的小強弟兄來接我.
義大利航空的東西.
其實西方的這些航空公司都一樣.
那個餐點實在不怎麼樣.
雞肉幾片跟土豆.
然後用microwave圍出來.
那就是我們的中餐還是晚餐搞不清楚.
早餐呢.
因為我們到的時候應該算是早餐.
早餐就是一個sandwich還是冰的.
連溫熱都不溫熱.
服了他們.
去義大利才知道旅館裡面早上的白煮蛋是冰的.
我第一次才發現.
竟然是冰的.
這個文化不同.

$^{561}$我要講的就是說.
實在是在上面吃那些東西.
幾乎等於是沒吃.
那麼到了洛杉磯是幾點呢.
洛杉磯出來出關已經一點半快兩點了.
我要公開表揚小強.
小強來接我.
我一坐進車子裡面.
我看到一個便當飯盒.
但我也不知道確定.
我也不確定是不是我的.
對不對.
可能是他嘛.
所以我就忍耐著不講話.
開了一會之後.
小強就說.
他說牧師啊.
不好意思.
我剛剛在店裡.
因為他的親戚在快餐店.
他說我剛剛在店裡匆匆忙忙的.
打了一個便當給你.
心裡好溫暖啊.
心想這麼好的人.
怎麼還沒有結婚呢.
你們家女兒還沒有結婚.
這個我推薦給你.
小強在不在.
在啊.
小強在後面.
不好意思.
丟著妹.
其實這就是家.
教會不管再大.
教會一定是一個家.
親愛丟著妹.
讓我們建立一個像家的教會.
一個彼此相愛.
彼此關心的教會.
不要把那個大公司.

$^{601}$大教會的那種.
那種不知不覺.
人際之間的關係.
變成是一種世公關係.
帶到我們中間來.
我們需要成為一個彼此相愛的教會.
才會有神的同在.
你的家要成為一個彼此關懷的家.
夫妻之間久了之後.
很容易就把對方當作是什麼.
當作是工人.
明天倒垃圾趕快.
你談戀愛的時候會這樣叫嗎.
不會.
那為什麼婚姻生活久了之後.
就變成是老公就應該是倒垃圾.
叫他倒垃圾的時候.
也不會講一聲.
可不可以麻煩你.
有沒有.
來跟我講一遍.
可不可以麻煩你倒垃圾.
可以請你幫我倒杯水嗎.
可以嗎.
老夫老妻還客套著幹嘛.
正是老夫老妻.
更不要失掉彼此之間的尊重.
你會發現婚姻當中.
很容易造成彼此之間關係的.
就是不尊重.
人際關係也是如此.
尊重是一個最基本的.
所以我們要使人和好.
我們必須成為一個和好的團體.
如果今天教會彼此紛爭.
教會吵吵鬧鬧.
你傳什麼和好的道理.
所以神要使我們與人和好.
這是神一再提醒我們的.
保祿在這裡接著說.

$^{641}$他說什麼呢.
他說在十架上面滅了冤仇.
便藉著十架.
使兩下歸於一體與神和好.
並且後面說.
傳和平的福音給你們遠處的人.
也給近處的人.
近處的人就在我們這邊.
我出去的這三個禮拜.
我交代門訓班.
我不在的時候.
請你們利用這一段時間.
上課時間去傳福音.
我那天看到我們姐妹門訓班.
他們很聽話.
就邀請了一些福音朋友.
就一起在哪裡.
在教會嗎.
那個朝上.
在哪裡.
在公園.
他們就邀了一些福音朋友一起在公園.
當作是一個傳福音的聚會.
弟子們.
這個不是應該因為功課要求.
傳福音應該變成是我們每一個人的責任.
你知道我們這次到義大利去了很多地方.
每到一個地方.
我一定看到誰.
耶和華競爭人.
他們就在那邊擺攤位.
然後放著書籍等等這些.
有興趣過去的人他們就跟他們談.
爭執無論在哪個地方.
無論得失不得失.
他們總在傳.
我幾乎沒有看到基督教的教會或者機構.
在那邊傳福音.
這是很令人感慨的一件事情.
義端反而在那邊很積極的傳福音.

$^{681}$那麼就在我們周遭的人.
其實你可能比牧師更有機會碰到很多不信主的人.
這些人是你的責任.
就在近處的人.
我們要把福音帶給他.
也要把福音傳給遠處的人.
這次前面的時間.
除了我們教會有大概十位左右.
還有別的教會的二十位左右.
一起我們到義大利.
你知道義大利很多一些的古代的歷史.
包括保羅在那個地方被斬頭.
包括保羅在那個地方.
還有一些中世紀的歷史都在那邊.
你可以看見.
鬥獸場.
基督徒在那邊.
在地下墳場.
這些為主受苦的基督徒.
當時候他們怎麼樣子在那邊.
我要講的不是這些所看到的東西.
我要講的是說.
當我們到這些地方去的時候.
我們跟一般的.
雖然我們是請當地的一個旅行團.
但是那個旅行團知道我們都是基督徒.
所以他特別找了一個也是基督徒.
是一個羅馬教會的弟兄.
那麼一路上面.
我就跟他講說.
我說每天早上我們要晨禱.
因為在旅館沒有時間晨禱.
所以我們在巴士上面晨禱.
每一次大概可能至少有半個小時時間.
我們每天早上晨禱.
還有呢.
我心裡面也想到說.
因為這也是我第一次.
在義大利也好.
不同的地方.

$^{721}$有些地方我也沒去過.
所以我也不知道那邊狀況如何.
不像以前我們到新約或者到聖地的時候.
我去過.
所以我知道在哪些地方.
那個地方可以唱歌.
哪個地方可以一起在那邊獻詩等等.
所以這一次算是比較新的一次.
但是呢.
結果我們就到了每一個地方.
我們到了每一個地方.
我們都在那邊做見證.
就唱詩歌.
我們每一次唱完.
旁觀的人都有人在拍手.
甚至還有人看了之後.
顯然他們也是基督徒.
他們還跟著一起唱.
所以我們到哪邊都唱祈恩典.
因為我們想這個大概是全世界.
基督徒都會唱的一首歌.
另外一首歌我們也常唱的就是10,000 Reasons.
因為這首英文歌.
可能一般基督徒都會比較熟.
而且在戶外唱也蠻好的.
所以我們大概主要就唱這兩首歌.
兩歌走天下.
所以以後我們師班也可以這樣子走出去.
真的.
當我們在那邊唱的時候.
在那許許多多世界各地來的人面前.
當我們唱完之後.
我們一起說.
耶穌愛你 Jesus loves you.
Amen.
請大家準備.
這個福音.
這個使人和好的福音.
也要帶到遠處去.
不管你今天到哪裡去.

$^{761}$神藉著我們.
要把這個和好的福音帶出去.
我們暑假有短暄隊.
即使你不是短暄.
你要出去旅遊.
你也帶著一個使命.
你就是和平的使者.
你就是那個把這個和好的福音帶出去的使者.
Amen.
你在那邊吃東西的時候.
你在那邊給小費.
你跟那個餐廳的那個waitress.
在那邊談的時候.
你知道嗎.
你有很多的機會可以跟他們來傳福音做見證.
Amen.
不要把我們的人.
把基督徒的生活分開了.
好像我們的基督徒變成是什麼.
我在教會裡面就是基督徒.
除了教會之外我就變成什麼.
我就變成世界上的人.
那你人格會分裂.
不是的.
我們無論到哪裡去.
我們就是基督的使者.
我們就是基督的大使.
跟你旁邊說你是基督的大使.
那麼.
在那邊.
這次去呢.
我也同時有很多次的機會去那邊講道分享.
我一到.
其實一到禮拜五晚上.
羅馬東區的華人教會.
他們青年團體就邀請我講道.
那麼我20年前去過那邊教課.
也在那邊的青年團體講道.
那個青年團體那時候大概有100多人.
後來這100多人的傳契呢.

$^{801}$他們的傳道人就出去開了一個教會.
也是在他們這個系統的.
所以大部分的年輕人都跟著他出去開教會了.
所以現在這些是新的一批.
大概也有五六十個人.
他們說這個主日聚會多的時候.
差不多也有100人.
我很難得看見.
我們這邊的第二代幾乎都到英文堂去.
是不是.
他們那邊.
就是所謂的第二代.
土生土長在義大利長大.
會講義大利文.
義大利文是他主要的語言.
可是他們也會講中文.
還會講溫州話.
所以很厲害.
所以他們在那邊聚會的時候.
就是講中文 講普通話.
所以很難得看到在華人教會有這麼多年輕人.
海外的華人教會.
所以跟他們分享.
那麼那一天分享完了之後.
他們覺得很好.
結果隔一個禮拜五又說要我再去一次.
所以我去了兩次.
去跟這些年輕人分享.
師母大概有PO一些在禱告會群組裡面.
有一些年輕人很可悟.
他們聽了之後特別講挑戰他們人生的意象.
他們就提到說他們很可悟.
他們剛從埃及短宣回來.
所以他們很火熱的.
想要真是在教會裡面.
能夠帶動年輕人一起宣教.
所以我們看見這些年輕人.
他們真是我們未來需要好好提拔培養的.
那麼接著.
後來我們也到.

$^{841}$剛才我提到.
在他們義大利南部的姐妹的退休會.
從羅馬往南邊.
拉布羅斯.
還有這些西西里島.
各個地方的都有.
那麼他們各地的姐妹們聚集在一起.
有三天兩夜的分心陪靈會.
所以在這個聚會當中跟他們分享信息.
那麼這是他們姐妹們.
每一年都會在那邊有這樣子的一種分享.
分心的聚會.
眺望他們的靈命.
那麼接著有一個禮拜.
就是在過去這個禮拜.
我就在他們的復興神學院.
那邊也是我們這邊正道神學院幫他們開的.
等於算是在那邊開了分校.
所以是叫做復興神學院.
有在實體的也有在線上的.
大概有25位26位左右.
所以一個禮拜的課一天講7個小時.
講到後來我實在是.
我說都是主的恩典.
我今天還能夠講到.
講完之後最後一天的晚上.
下午通常慣例.
我會帶他們一起禱告.
那個禱告會是足足兩個小時.
那麼帶著他們沒有停的.
帶他們一起禱告追求.
唱詩歌禱告這樣子.
看見這些弟兄姊妹們也很可慕.
你要知道.
溫州教會他們是很少傳道人的.
他們大部分都是弟兄姊妹們.
代持的出來服事.
所以加上他們都是做生意的.
大概有的初中畢業.
甚至有的小學畢業就出來工作了.

$^{881}$就出來打拼出來做生意了.
所以他們就沒有什麼太多的機會去受教育.
所以他們一邊服事就感覺到自己很不夠.
所以他們就希望說能夠有神學裝備.
所以這些學生.
他們不一定將來是全時間傳道人.
但是他們都是在教會的裡面.
他們叫做講道人.
或者他們叫傳道.
其實他們的傳道就是.
他們派單的到各個地方去講道.
你到羅馬你到拉布羅斯.
每一個禮拜不同的派單制.
這是溫州教會的一個特色.
神也使用溫州教會.
藉著他們就在義大利.
全義大利還有西班牙.
法國這是我知道這三個地方.
大概主要的都是溫州人的教會在那邊建立.
所以藉著他們能夠藉著他們的工作.
藉著他們去做生意.
就把福音帶到了各地.
當然教會真是看見說我們何等需要被打開.
因為溫州的教會.
它很自然是會向溫州人傳福音.
甚至有的還是用溫州話.
最起初大家早年的時候還有溫州語聚會.
現在比較沒有了.
現在因為東北的人也來了.
或者也有一些其他的人來了.
所以他們開始就用普通話為主了.
親愛的同學們.
同樣的我們很多時候我們也會有界限.
我們很多被我們自己限制了.
我們華人教會.
我們就很容易想到我們就跟華人傳.
他們那邊也是.
但是那天我就跟他們提到.
我後來問了一下他們才知道.
那時候我到羅馬教會去講道的時候.

$^{921}$我看到他們教會的旁邊的人幾乎都是黑人.
對不起.
就是說幾乎都是膚色是黑的.
我就問他們說.
我說這些人是什麼人啊.
他就告訴我說.
他說這些人大部分是孟加拉人.
還有呢.
還有就是巴基斯坦人.
後來才知道說.
原來孟加拉是除了華人以外.
最大的南亞的就是亞洲的移民.
再來就是巴基斯坦.
那巴基斯坦孟加拉都是所謂的伊斯蘭教的.
對不對.
他們都是伊斯蘭教.
都是信這個伊斯蘭教的.
所以這些的移民.
他們在那個地方.
通常都是幫這些華人餐廳.
或者是幫他們做工的.
或者是他們自己開個店在旁邊.
那麼我就跟他們提醒說.
這些人都已經到你們家門口了.
你們應該要把福音傳給他們.
Amen.
現在我們華人需要被神突破.
Amen.
我們覺得這些人非我族類.
同樣的,為什麼我們教會.
當初我們要做阿富汗難民的工作.
為什麼.
阿富汗難民都到了我們家門口了.
誰把福音帶給他們.
是不是.
我們可以把福音帶給他們.
所以這個福音不僅僅傳給我們附近周遭進出的人.
同文化的人.
這個福音也要透過我們.
把它傳給我們周遭這些異文化的不同族裔的人.

$^{961}$甚至是遠方的這些異文化的人.
求主使用我們的教會.
讓我們教會真是成為一個使萬民做主門徒的教會.
Amen.
教會存在一定要宣教.
教會存在一定要傳福音.
教會如果不傳福音.
小主如果不傳福音.
你不傳福音.
你就是一個死人.
我必須講得更嚴肅一點.
因為神給你的這個福音.
神給你的恩典.
在你的裡面是死海.
沒有流出去.
只有當你傳福音的時候.
你才會感受到那種神福音的那個火.
才會眺望起來.
所以保羅在剛才我們讀的經文裡面.
保羅他這樣說.
我真盼望我們能夠像他一樣.
他說什麼.
保羅說我們做基督的使者.
就好像神借我們勸你們一般.
我們替基督求你們與神和好.
保羅在這裡用勸這個字.
這個勸原文的意思是很強烈的懇求.
很強烈的一種懇求要求.
那麼這個求呢.
也是懇求的意思.
換句話說.
什麼叫懇求.
這個先生啊.
有空來教會啊.
這是不是懇求.
什麼叫懇求.
你要不要示範一下.
給你旁邊看看.
你一定要去啊.
太重要了.

$^{1001}$你一定要信啊.
現在為什麼.
這就是保羅的心態.
保羅迫切到一個地步.
他說我替基督求你們.
我懇求你們信耶穌吧.
所以保羅在這地方懇求.
懇求到一個地步.
而且在原文裡面.
這是一個持續的.
持續不斷的一個動作.
表明說保羅在這個地方.
他說我持續不斷的.
持續不斷的.
我一直一直的懇求你們.
信耶穌吧.
教會不傳福音.
就會變成是外功外行.
內功內行.
聽懂嗎.
聽懂不懂.
聽不懂再講一遍.
教會不傳福音.
就會變成外功外行.
內功內行.
懂嗎.
外戰.
外功.
外戰.
就是說.
我們如果不出去傳福音.
我們就外行了.
外戰.
內外功吧.
功夫的功.
OK.
懂這意思就好了.
然後內功怎麼樣.
功是什麼功.
功極的功.

$^{1041}$內功內行.
懂嗎.
今天很多教會不傳福音.
大家在裡面.
你看我不順眼.
我看你不順眼.
開起會來.
檢討這個檢討那個.
但是每個人只叫別人.
做自己一根指頭不肯動.
只是會開會.
只是會講話.
只是會說.
但是從來不傳福音的教會.
大有人在.
我這次去義大利.
我看到一些教會的問題.
我心裡面很感慨.
二十年前我去一個教會.
在那邊服事的時候.
就聽到他們之間有紛爭.
二十年後這個紛爭.
不僅僅沒有停止.
還越來越烈.
歐祖啊.
生命何其短暫.
浪費了二十年的光陰.
教會在那麼好的location.
但是教會沒有向外發展.
沒有對外傳福音.
只是花功夫在彼此之間的內鬥.
親愛的弟兄姊妹.
你的家是不是也如此呢.
你的小組千萬不要變成如此.
我們要成為一個彼此相愛.
同心合意傳揚福音的教會.
我們要成為一個彼此相愛.
彼此鼓勵.
一起來成為榮耀的一個見證.
讓別人看到我們這個家庭就羨慕.

$^{1081}$他們也要成為基督徒.
\newpage



\section{路加福音 10:1-12}
\label{sec:032166}
\textbf{莊稼已熟,工人在哪裡?}~~~~慎勇牧師
\newline
\newline
連結: \href{http://www.church.com.hk/acms/content.asp?site=cdc\&op=show\&type=product\&code=032166\&layout=sermon}{\texttt{032166}} ~~~~ 語音日期: 2026-05-17
\newline
\newline
\hyperref[sec:032165]{\small{< < < PREV SERMON < < <}}
~
\hyperref[sec:index_chronic]{\small{[返順時目]}}
~
\hyperref[sec:032167]{\small{> > > NEXT SERMON > > >}}
\newline
\newline
路加福音 10:1-12
\newline
\begin{longtable}{cl}
\hline
\hline
章節 & 經文 (和合本修訂版)\\
\hline
10:1 & \begin{tabularx}{0.7\textwidth}{X} 這些事以後，主另外指定七十二個人，差遣他們兩個兩個地在他前面，往自己所要到的各城各地去。 \end{tabularx} \\ \\ \relax
10:2 & \begin{tabularx}{0.7\textwidth}{X} 他對他們說：「要收的莊稼多，做工的人少。所以，你們要求莊稼的主差遣做工的人出去收他的莊稼。 \end{tabularx} \\ \\ \relax
10:3 & \begin{tabularx}{0.7\textwidth}{X} 你們去吧！看！我差你們出去，如同羔羊進入狼群。 \end{tabularx} \\ \\ \relax
10:4 & \begin{tabularx}{0.7\textwidth}{X} 不要帶錢囊，不要帶行囊，不要帶鞋子；在路上也不要向人問安。 \end{tabularx} \\ \\ \relax
10:5 & \begin{tabularx}{0.7\textwidth}{X} 無論進哪一家，先要說：『願這一家平安。』 \end{tabularx} \\ \\ \relax
10:6 & \begin{tabularx}{0.7\textwidth}{X} 那裡若有當得平安的人，你們所求的平安就必臨到那家，不然，將歸還你們。 \end{tabularx} \\ \\ \relax
10:7 & \begin{tabularx}{0.7\textwidth}{X} 你們要住在那家，吃喝他們所供給的，因為工人得工錢是應當的；不要從這家搬到那家。 \end{tabularx} \\ \\ \relax
10:8 & \begin{tabularx}{0.7\textwidth}{X} 無論進哪一城，人若接待你們，給你們擺上甚麼食物，你們就吃甚麼。 \end{tabularx} \\ \\ \relax
10:9 & \begin{tabularx}{0.7\textwidth}{X} 要醫治那城裡的病人，對他們說：『神的國臨近你們了。』 \end{tabularx} \\ \\ \relax
10:10 & \begin{tabularx}{0.7\textwidth}{X} 無論進哪一城，人若不接待你們，你們就到大街上去，說： \end{tabularx} \\ \\ \relax
10:11 & \begin{tabularx}{0.7\textwidth}{X} 『就是你們城裡的塵土黏在我們的腳上，我們也當著你們擦去。但是，你們該知道神的國臨近了。』 \end{tabularx} \\ \\ \relax
10:12 & \begin{tabularx}{0.7\textwidth}{X} 我告訴你們，在那日子，所多瑪所受的，比那城還容易受呢！」 \end{tabularx} \\ \\
[1ex]
\hline
\hline
\end{longtable}
$^{1}$我們同心祈禱.
在主日的時候帶領我們來到你的面前.
讓我們用詩歌.
頌讚你 也都讓我們.
心靈誠實去敬拜你.
更加求天父你帶領我們.
讓我們的心向你敞開.
願你的話語進入到.
我們的心裡面 成為我們生命.
成長和轉變的一個.
力量 也都求天父.
你看顧保守我們.
求主你的寶血.
潔淨求主的聖靈.
在我們當中自由的運行.
讓我們的心.
傳言向主敞開.
求主你親自.
對我們每一個人說話.
奉耶穌基督的聖名 阿們.
其實我有一個問題.
你知道有時候牧師.
又不是說有空.
牧師是需要常常動腦筋的.
我常常想一些問題.
比如說我想.
為什麼我們應該教會.
應該怎樣去看待.
宣教的工作 其實你知道.
這個題目 你一看這個題目.
牧師你又是說猜拳 又是說宣教.
為什麼你每次都說.
這件事 我們想一個問題.
為什麼教會.
要有猜拳事工.
教會應該如何.
來看待猜拳事工.
在教會當中.
宣教猜拳.
這個事工 是不是.

$^{41}$在教會當中眾多事工.
的一個環節.
抑或是我們基督徒信仰生活當中.
的一個選項.
也就是說 宣教猜拳.
是不是對於我們這些信主的人來說.
是一個選項.
你可以宣教 你可以不宣教.
你可以參與 你可以不參與.
是不是這樣.
抑或宣教是在教會.
必備的一個項目 一個配搭.
就好像我們廣東人說.
連晚戰隊.
人有我有 每個人都有宣教部.
我們沒有宣教部是不行的.
所以我們弄一個猜拳部.
宣教部.
我們如何來看待宣教這件事.
我想.
我們如何去定位猜拳.
宣教在教會眾多事工.
當中的一個位置.
這個很重要.
我想如果找到一個正確的答案.
我需要認真的思想.
一個問題.
就帶出另一個問題.
也就是說我們必須思想.
我們要成為耶穌基督的門徒.
我們今時今日成為了基督徒.
我們究竟.
是為了什麼.
可能以前的老師說.
你們信耶穌得永生.
對不對.
對.
當然對 你不要懷疑.
信耶穌是得永生.
但是信耶穌是否只是為了得永生.

$^{81}$不是吧.
你們都100分.
我不用說了 我下台了.
我們要明白.
我們夢召成為.
耶穌門徒的究竟.
是為了什麼.
我們回頭看看.
我們主耶穌的心意究竟是什麼.
我們從聖經裡面.
去找尋答案.
因為你說我說 他說沒用.
我們看看神是如何說.
聖經是如何告訴我們.
如果我們從福音書來看.
我們可以見到一個非常清楚的.
一個漸進式的.
一個善奴.
主耶穌來到這個世上.
去呼召人悔改.
跟隨他.
對不對.
神愛世人.
甚至將他的獨生子賜給他們.
叫一切信他的不致滅亡.
反得永生.
17字怎麼說.
因為神差得如此降世.
不是為要定世人的罪.
乃是要世人因他得救.
所以我們都是.
蒙恩蒙召.
得蒙救贖的人.
我們蒙召悔改跟隨他.
然後我們見到主耶穌讓.
這些蒙召的人去見證.
和經歷主的豐盛.
他呼召了門徒之後.
接下來做什麼.
就讓他們跟隨耶穌.

$^{121}$經歷他們的見證 經歷主的豐盛.
然後在過程當中.
不斷地幫助他們.
去明白神的旨意.
並且去端正.
他們的三觀.
人生觀 世界觀 價值觀.
登山寶訓.
是一個很典型的例子.
當眾人來到.
耶穌面前要求福的時候.
為什麼要求福.
如果你看馬太福音第四章.
25節打後.
成千上萬的人來找耶穌.
當中有很多不同類型的人.
有得病的.
有沒飯吃的.
有分家產的 有很多問題的.
有心裡憂鬱的.
各種各樣的人都來找耶穌.
求什麼 求一個福字.
然後才去到馬太福音第五章.
那個我們很熟悉的.
登山寶訓.
你不要理他八福也好.
九福也好.
那些神學家去爭論 我們無謂浪費時間.
八福也好 九福也好.
你有沒有留意到這八福.
這九福指向的是什麼.
指向的是天上的福氣.
其實這個是不合邏輯的.
當眾人來尋求耶穌的時候.
他們尋求的是地上的福氣.
但是耶穌給他們的答案.
是什麼 是天上的福氣.
是什麼.
就是要轉變他們的三觀.
如今活著的不再是我.

$^{161}$乃是基督在我裡面活著.
藉我的生命去榮耀主的名.
我們人生的價值觀去改變.
但是你看到耶穌帶領門徒.
這三年多的時間.
不斷發生的就是.
門徒與主耶穌之間的衝突.
這種衝突不是正面的衝突.
是在思想上的衝突.
在認知上的衝突.
甚至耶穌的.
我們稱之為大弟子.
是誰.
彼得很喜歡說話的那個.
連耶穌快要釘十字架.
他還說錯話.
拍硬胸口說.
主耶穌你不用死.
我為你死.
其實不要怪他.
因為聖經說.
其他人也都是這樣說的.
只不過他口快.
大頭說.
其他人門徒說.
我願意為你死.
結果有沒有為主死.
雞叫之前.
三次不應轉.
是一個很大的爭戰.
為什麼耶穌要來.
目的是什麼.
他不明白為什麼耶穌呼召他.
跟隨他.
在這個過程當中.
耶穌也激勵他們.
持守真理.
要過一個聖潔的生活.
並且要訓練他們成為與神同工的人.
差遣他們出去.

$^{201}$做訓練.
做短訓的訓練.
要讓他們與神同工.
要將屬天的恩賜.
賞賜給門徒.
讓他們能夠在神的榮耀當中.
可以共享.
但這還不是目的.
為什麼有這麼多的改變.
這麼多的訓練.
這麼多的機會.
最終的目的是什麼.
是要差遣他們出去.
完成大使命.
當他們去.
完成大使命的時候.
就是他們經歷得勝的時候.
這才是目的.
這才是主耶穌呼召門徒的心意.
其實你們也知道.
我剛才已經問了你們.
你們已經回答了我.
在理性上我們也很清楚.
我們信耶穌不只是為了得永生.
那是為了什麼.
是為了大使命.
為什麼耶穌升天的時候.
如果我們信耶穌就得永生的話.
那最好的結果是什麼.
我們信耶穌的那一刻.
我們與主同在不是很好嗎.
去到父神那裡不是很好嗎.
何必在地上.
勞苦受煩.
但這不是耶穌的心意.
《約翰福音》17章.
等一下你們會看到一點.
耶穌的心意.
就是將我們留在地上.
去完成一個使命.

$^{241}$在完成使命的過程當中.
經歷他的豐盛.
讓我們的生命更有價值.
剛才那首詩歌我很喜歡.
最後那首特別是.
藉我的生命.
去榮耀主名.
你的生命價值是不是提升了.
是提升了的.
這個是主的心意.
在《約翰福音》17章第18節.
其實主耶穌很清楚地.
告訴我們.
去呼召門徒.
為的是要差遣我們.
所以我們知道《約翰福音》17章.
可以稱之為大祭司的禱告.
他在那裡禱告.
當然他是和天父之間的禱告溝通.
你怎樣差遣我到世上.
我也照樣.
差他們到世上.
這個是大祭司的禱告.
當然我們稱之為耶穌是大祭司.
他的禱告當中.
一個點睛之筆.
也告訴我們.
他為什麼要呼召我們.
為的是要差遣我們.
差遣我們做什麼.
完成大使命.
差遣我們去什麼.
藉著我們的生命.
去榮耀主名.
弟兄姊妹們.
我們剛才不是在唱歌.
我們剛才是在禱告.
我們剛才是在敬拜.
雖然我們沒有拍硬胸口.
但是我們隱形的手已經拍硬胸口.

$^{281}$主祀我生命.
榮耀主名.
口唱心和.
讚美主.
雖然我們是在唱歌.
但其實不是唱歌那麼簡單.
我們是在向主禱告.
是一個回應.
耶穌呼召門徒.
為的是要差遣我們.
已不鎖書.
早兩年我寫了一本小冊子.
也說到已不鎖書.
已不鎖書告訴我們.
主耶穌拯救我們.
呼召我們來跟隨他成為門徒.
不只是讓我們得永生.
不只是讓我們在生活當中.
得著各樣屬靈的福氣.
也不只是讓我們參與各種.
屬靈的聚會.
活動 搞得我們很忙碌.
這不是主的心意.
我們相信耶穌.
不要經常想著我們得得得.
原來.
耶穌呼召人.
悔改來跟隨他.
成為他的門徒.
是要讓我們.
成為永耀神的見證.
並且被差遣.
去參與.
大使命.
使得萬民.
成為主的門徒.
這就是為什麼在馬太福音.
和馬學福音.
高峰的時候 都停留在哪裡.
都停留在門徒.

$^{321}$受主差遣.
弟兄姊妹們.
我們聖經應該很熟悉.
但我們有沒有看到這件事.
聖經福音書究竟.
告訴我們什麼.
不只是讓我們知道耶穌的生平.
我們的知識豐富了.
不只是這樣.
更加是要明白.
為什麼神讓我們看到.
他的心意是什麼.
他的計劃是什麼.
弟兄姊妹們.
為什麼我們要積極.
努力的參與.
和推動.
宣教和差遣的事工.
因為.
這是我們夢召.
成為主耶穌門徒的意義.
和我們領受的使命.
是我們的使命.
就算我們每一個夢召成為.
耶穌門徒的人.
基督徒.
其實基督徒是別人稱呼我們.
其實我們真正的身份.
我們是耶穌的門徒.
你有沒有見過有些基督徒說.
我不是門徒 我是基督徒.
理論上不應該有這件事.
所以大使命是給誰的.
給牧師的.
給長老的.
給執事的 給理事的.
給童工的.
不是.
各位耶穌的門徒.
就是你 就是我.

$^{361}$只不過我們.
扮演不同的角色.
但我們必須要清楚.
這是我們的使命.
所以.
宣教不只是.
教會生活和建造當中.
一個重要的項目.
也不是一個配件.
從根本上來說.
更加只是教會生活.
建造 甚至是教會.
存在的意義.
價值和目的.
我講得再直白一點.
若果一個教會.
被設立在這個世上.
卻不做福音的工作.
不做猜拳宣教的工作.
其實他就失去了.
他的意義.
環顧當今的世界.
其實莊家.
真的已經熟了.
可以說熟透了.
但另一個殘酷的現實.
就是的確.
很缺少做工的人.
讓我們從當年.
主耶穌的教導.
了解和明白主對當今.
門徒和教會的心意.
同樣我們.
願意去遵行.
主耶穌的吩咐.
我相信教會.
必將成為更有活力的.
更健康的教會.
首先我們來看.
福音和祥.

$^{401}$極為的廣闊.
今天我們讀的經文.
第一節.
第一節開頭.
說「這事」以後.
「這事」是什麼.
其實如果回想上面對上的經文.
就是講到之前發生的一件事.
就是耶穌登山變相.
我們都知道登山變相.
登山變相這個事情之後.
這個登山變相.
也是主耶穌在.
傳道旅程當中的一個轉折點.
因為從登山變相之後.
祂就下山.
祂下山.
就帶領門徒.
邁向耶路撒冷.
不是去過節.
不是去度假.
邁向耶路撒冷.
邁向各國的十字架.
因為這個是祂的使命.
主耶穌所有的行程.
都是.
目標.
定睛在祂的使命當中.
所以轉折點.
就是從登山變相之後.
祂知道時間就到了.
祂就奔向.
祂的使命.
祂的使命是什麼.
十字架要成就的救贖之恩.
在馬太福音第十章.
第五節和馬可福音第六章.
第七節.
都記載了主耶穌曾經.
差派過十二個門徒出去.

$^{441}$短宣.
但這次在路加福音這裡告訴我們.
主耶穌在這個時候差遣了多少個人.
七十個人.
當然有抄本.
是說七十二個人.
我們不需要太關注具體人數是多少.
重點是什麼.
我們見到什麼.
就是當耶穌知道自己成就.
救贖計劃的時間臨近.
祂知道自己就要被釘十字架.
所以我們見到.
主耶穌就加快去培訓.
和操練門徒的速度.
以前差派十二個.
十二個.
我這次差派七十個.
抓緊時間.
好好地訓練他們.
培養他們.
培訓他們.
我們看到的是什麼.
我們看到的是.
主耶穌差遣門徒出去傳福音的.
那種迫切性.
我們從祂的智慧行間.
我們可以體會得到.
要收的莊稼多.
但工人卻少.
之前.
每一次差派十二個.
但不夠我們這次要差派七十個.
加倍.
數學厲害的你可以數出多少倍.
加倍這個訓練.
為什麼.
因為傳福音是何等的迫切.
也讓我們見到.
福音的和長.

$^{481}$是何等的廣闊.
藉我此恩福.
我也很喜歡這首歌.
其中一首.
很喜歡.
開頭歌詞讓我們看到一件事.
在這個世界上有很多.
將要失去的靈魂.
他們需要聽到福音.
在很多迷失的羊.
需要被召聚.
歸回到神的羊圈.
但是如果我們沒有.
那種迫切感的話.
我們就唱歌.
在這個世界上.
多少靈魂要消失.
我已經得救了.
關我什麼事.
你不要這樣說.
你當然不敢這樣說.
但你心裡的態度就是這樣.
不關我的事.
我們連那種唱歌的迫切感都沒有.
我記得這首歌是.
比較舊的歌.
比較老的歌.
我們不要覺得.
舊的歌都是錯的.
有很多都是要保留的.
當年我想起.
我蒙召的時候.
也是這首歌的感動.
是那種迫切感.
當我們有一種迫切感的時候.
你才真正能夠找到.
你生命的目標和方向.
現在教會最大的問題.
是什麼.
我學了一個詞.

$^{521}$(笑聲).
你們笑得這麼大聲.
(笑聲).
心有同感.
誰躺平.
我們.
你不要說別人.
指別人最厲害.
別人他躺平.
三隻手指指著你.
我們習慣了.
所以求主幫助我們.
當我們在福音的事公上.
當我們在恩典當中.
當我們躺平的時候.
我們一定沒有這種迫切感.
我們要留意一件事.
這個迫切感.
在兩千年前.
就已經存在了.
耶穌說這番話的時候.
是什麼時候說的.
昨天說的?前天說的?.
將近兩千年前.
說的.
兩千年前.
耶穌已經說.
要收的莊稼多.
莊稼已經熟了.
什麼時候熟?不是現在才熟.
那時候就已經熟了.
兩千年前就已經熟了.
每一個時代.
最缺乏的是什麼?.
不是缺乏莊稼熟不熟.
缺乏的是收割的工人.
我們都中了撒旦的詭計.
最好有一個推卸責任.
推給撒旦.
當然是他的什麼鑊.

$^{561}$當然是他的責任.
但是我們都上了當.
他們讓我們已經習慣了.
在恩典當中躺平.
讓我們習慣了.
在我們生活當中.
他給我們很多的理由.
讓我們忽略了.
這個大使命.
是與我們息息相關的.
我們看到在任何的時代和地區.
福音的需要.
都是迫切的.
基督徒 教會.
我們應當隨時保持這種迫切感.
否則我們就會退步.
就會落後.
其實在屬靈的事情上.
在福音的事公上有一個原則.
就是不進則退.
你說我們教會.
其實都沒怎麼後退.
但又沒怎麼前進.
其實這樣就已經退後了.
至少讓我到今時今日為止.
我還有這樣的意識.
不要讓我自己停下來.
不要讓我.
停留在一種安寧當中.
很多時候.
師母現在已經習慣了.
師母有時.
過去這二十多年.
我有兩三次.
轉換工場.
其實我的方向是很清楚的.
師母說.
好像你.
不夠穩定.
好像你一安定下來.

$^{601}$你就想要突破.
我說是呀.
教會沒有四平八穩.
我已經看到將來的方向.
我覺得對我來說.
這樣的人生不夠精彩.
當然這不是我想的.
神有神的一個大靈.
我以後有機會再講.
神有很清楚的大靈.
經常都在做工.
我今次去三月份.
意大利.
你們可能有人去過.
參觀羅馬和保羅.
和彼得被關押的水路.
有沒有人去過.
沒去過.
你去過.
裡面沒有水了.
有個見證人.
竟然在水裡見到彩虹.
奇怪.
我看到的時候.
彩虹.
我就立即拍照.
被我拍到.
因為我也是在尋求.
神的一個大靈.
神就是用彩虹給我一個印記.
很多時候不是說.
我們自己想做什麼.
而是我們要學習.
怎樣去尋求聖靈的大靈.
但我們的心很重要.
心已經躺平.
就算聖靈在你面前.
講來講去.
你都聽不到.
所以最重要的是什麼.

$^{641}$我們的心.
讓我們保持一種活力.
和你的年紀無關.
我現在七老八十.
七十八十歲都算年輕.
七十歲叫壯年.
以前我們教會.
有張長者席.
不夠九十歲.
都不敢埋位.
所以求主幫助我們.
那種迫切感.
福音的和長有多大.
我們看看數字.
因為這個年代喜歡看數字.
在2026年5月.
最新提供的這個數據.
世界人口已經.
將近接近83億.
增長真的很快.
我們看到.
基督教目前.
你留意這個用詞.
目前還是.
全球最大的宗教.
大概24億.
為什麼說.
暫時還是.
因為有其他的信仰已經追上來.
追得速度很快.
迅速追上來.
是什麼.
伊斯蘭教.
我們目前來說.
算是包括了24億人.
包括天主教 東正教.
和我們這個新教.
我們算是新教.
大概佔全球.
人口30\%.

$^{681}$數字我就不說了.
你自己看 天主教人數.
東正教人數 新教人數.
大概是多少.
按照人口的分裂.
這個國家是多少.
美國是多少.
美國是包括天主教和新教.
這個數字.
看起來也不錯.
自己看自己也不錯.
然後穆斯林人口最大.
佔大多數的土耳其.
信徒只有.
8000.
只是佔萬分之一.
所以你很清楚.
看到這個比例.
其中有一個數字.
中國的基督徒是多少.
有人說.
1.5億.
我很負責任地說.
我不是說不想有1億基督徒.
但我可以說這些數字.
是吹水.
誇大.
毫無意義的誇大.
但保守一點來說.
比較準確一點來說.
大概6000多萬.
比較準確一點.
但我們看到.
在穆斯林的地區.
我們舉一個例子.
土耳其在土耳其比較熟.
佔地區.
萬分之一.
等於說很少很少.
所以我們相信.

$^{721}$主耶穌的日子越來越近.
但世界上.
做一個數字可以清楚地減.
加減法.
83億減24億是多少.
59億.
至少還有59億將要失去的靈魂.
需要主耶穌的福音.
和救恩.
其實受到經濟.
政治戰亂很多方面的影響.
人心越來越.
徬徨不安.
對未來越來越迷茫.
我們看到人們.
活在心靈的黑暗之中.
苦苦掙扎.
要尋找人生的出路.
要尋找光明的出路.
所以為什麼這幾年有這麼多人走線.
這個名詞也是這幾年才有的.
不過好聽一點.
以前都有的.
以前叫偷渡.
講話要含糊一點.
叫走線.
走線就是篤出.
篤出也是早期的.
以前說大陸篤出是香港.
篤出過河.
過那條羅湖橋.
為什麼.
因為他們要尋找人生的出路.
要尋找光明的出路.
在這幾年.
不要說這幾年.
就說今年.
我們這三次的出外.
都感受到這個桑架.
真是熟了.

$^{761}$我們不是去殺種.
聖靈已經殺了種.
我們要去收割.
我們去意大利.
我很快講一講.
我沒有願意講建制.
很快講一講.
我們去意大利的時候.
我們的導遊是基督徒.
剛剛信了主.
但是他的太太還沒有信.
剛剛我們去到佛羅倫薩的時候.
他住在那裡.
所以他太太來探班.
其實不試探我們,試探他.
探我們的導遊.
就在一個很簡短的吃飯的時候.
我們就把握一個機會.
向他傳福音.
其實不用傳.
直接就問他你要不要信耶穌.
誰知他就肯.
這樣就收割了桑架.
就帶他去做決戰祈禱.
這叫做冷手.
收割現戰隊.
團團都是這樣的.
今年四月份.
我們去土耳其.
當然土耳其和他開車的司機.
和導遊.
都是土耳其人.
司機又不會講中文.
導遊會講一點.
他們很羨慕我們.
後來那個導遊跟我說.
司機很喜歡我們這一團.
他說從來沒有見過我們這麼友善的團隊.
而且他很喜歡聽我們唱歌.
你知道我們唱什麼歌嗎.

$^{801}$在車裡.
唱聖詩.
唱讚美詩.
然後他說.
你們講了一句說話.
和剛才那首詩歌.
裡面有一句講得很清楚.
是我們生活當中的見證.
可以影響到人.
對我們來說很正常.
但是對司機來說很大的衝擊.
他說我沒試過.
接一個團.
幾天之後.
我的車就好像沒人坐過.
意思是什麼.
很乾淨.
就憑這件事.
他就覺得我們與眾不同.
當然他知道我們是基督徒團隊.
他用土耳其文寫了一個.
朝聖團.
這個不是我叫他寫的.
是他自己寫的.
因為我們參觀七教會.
有大灣區廣東省.
粵東.
之類.
當然我們習慣在車裡唱詩歌.
讚美禱告.
有一天.
大家都很累.
從廣州開車到惠州.
兩三個小時.
突然導遊就說.
你們不要睡.
你們可不可以唱歌.
司機很喜歡聽你們唱歌.
其實不是說.
司機他也喜歡聽.

$^{841}$唱什麼歌.
唱聖詩.
然後我們快要回來的時候.
我們在深圳結束.
從汕頭汕尾開車.
回到深圳的路上.
很長的路.
正在下雨.
導遊問我可不可以有一個不情之請.
我說什麼事.
他說你們可不可以再唱兩首歌.
我想錄下來.
我說哪兩首.
「耶和華的心」.
第二首他選擇了什麼.
「寶貴十價」.
我說當然是「女犯應」.
他說好.
他既然要錄下來.
我們要常常聽.
當然路途當中.
我們也有跟他分享福音.
也是給他一個福音的插子.
要明白真理.
他很認真在讀.
他說不好意思.
他知道我是牧師.
不好意思.
我們很累.
所以晚上就睡不著.
看不多.
但這個團結束之後.
我會繼續看.
當然我也有問他.
你信不信耶穌.
我現在還做了一些.
我認真尋求.
但你可以看到.
他主動尋求.
這個真理.

$^{881}$說明什麼.
我們去是收割.
很多時候我們只不過是一句說話.
你要不要信耶穌.
很多時候他就要.
去年我們去華夏之旅.
去西安的導遊.
一問他.
眼淚飆了.
很奇怪.
其實不是奇怪奇妙.
你要不要信耶穌.
流眼淚.
我要.
提醒我們.
我們要去.
無論得時不得時.
總要去傳福音.
正如主所說的.
要收的莊家多.
表明福音和祥的需要.
是極為廣闊.
至少有五十九億.
我們要去.
我們再看看.
福音和祥的挑戰也很大.
第三次說你們去吧.
我猜你們出去如同楊高進入狼群.
不過指出宣教猜拳.
一定會遇到很多挑戰和困難.
會感到無助.
甚至有時會遇到一些危險.
被人排斥你.
被人嘲笑你.
甚至可能.
身體上的逼迫.
都有可能.
當然我們禱告.
希望這些事不要發生.
就好像主耶穌所形容的.

$^{921}$如同楊高進入狼群.
所謂進入狼群.
進入狼群未必真是.
你身體上有什麼損害.
但更多可能是精神方面.
比如說.
有時你會遇到一些來自於政府的.
或者一些群體的.
或者個人的抗議.
逼迫,抵抗,傷害等等.
甚至包括.
來自一些異端的挑戰和攻擊.
也包括會遇到各樣的艱難.
誘惑和攻擊.
但有一樣東西.
可以幫助我們克服.
就是當我們.
很清楚地將我們的目光.
我們的人生.
定標在這個大使命.
當我們對這個福音的工作.
是一種迫切感的時候.
其實這些困難你不會在乎.
就好像我們.
人生當中很多經歷.
如果有一樣東西你真的很想做.
甚至是影響你的.
起居飲食.
這是你的目標.
你身邊所遇到的事情.
對你來說就不算什麼.
我見到我們當中年長的不少.
可能最喜歡的.
我有很多牧師朋友.
都差不多這個年紀.
他們經常在我面前曬.
曬他們孫的照片.
臭孫.
當你的目標是很清楚.
要去臭孫.

$^{961}$你開車的路途上.
有時有些阻滯.
颳風,下雨.
汽車沒有油,爆胎.
還是不是問題.
不是問題.
你會想盡一切辦法解決問題.
直達目標.
因為你的目標很清晰.
我舉一個很小的例子.
我們就應該明白.
倘若我們對福音,對使命.
是能夠真是把持得很穩的話.
我們真是好像聖經所說的.
真是主耶穌吩咐.
我們將大使命當作是我們的使命.
有這種迫切感的話.
其實你生活當中的其他困難.
都已經不能再阻擋你.
繼續前進.
其實為何有時.
我們在這個世代上.
還會經常.
有時有些混亂.
因為很多時候.
碰到這些事情.
颳風,下雨,那裡打仗,那裡射飛彈.
我就這麼緊張.
我發覺.
同樣觀察我們的同工.
都會緊張.
為甚麼?.
因為他們的使命很清楚.
人生的價值很清楚.
他們的迫切感很清楚.
這種迫切感.
這種清楚人生的目標.
使他們不會受到環境任何的影響.
所以這才是我們說.
不要看環境.

$^{1001}$是定睛在主身上的意思.
我們要明白.
這句話.
不是要讓門徒知難而退.
它是給我們一個秘訣.
它是告訴我們.
我們出去做福音的事工.
去參與的時候.
無論你是禱告,金錢奉獻.
親身參與.
長宣也好,短宣也好.
你一定會遇到一些困難.
我們又說撒旦.
撒旦一定會興起一些困難.
難阻你.
但你不要怕.
越是難阻你.
越是說明你做的事是對的.
我們都很任性.
在生活當中.
你不要回答我是不是.
但我們唯一缺乏的.
就是我們在信仰上.
在真理上.
不夠任性.
我們在真理上.
很容易妥協.
所以為什麼我們的生活.
好像都是這麼多波動.
很多人說.
牧師你去土耳其你不怕嗎.
我怕什麼.
你不怕打仗嗎.
伊拉克打過來很久都射不到.
是不是.
不用怕.
明年我三月份還要去以色列.
去不去.
正在打的.
不用怕.

$^{1041}$照樣去.
當然它現在還沒有開放.
你想去都不讓你去.
因為我們知道.
我們在做什麼.
主耶穌給我們一個秘訣.
不是讓我們知難而退.
要讓我們什麼.
在我們戰騰大使命.
去宣教的時候.
會面對的處境.
但讓我們明白.
我們有一個秘訣.
就是緊緊的依靠主耶穌.
去面對和跨越這些挑戰.
而最終去經歷.
那種特性.
這也是我們生活當中的動力.
面對宣教的挑戰.
進行主耶穌所教導的.
我們就會經歷.
主的同在.
和奇妙的作為.
這是事實.
但關鍵是.
我們是否相信主耶穌的話.
在第四節.
當然我們今天沒有讀經文.
主耶穌在今天.
在《路亞福音》後面.
第四節到第九節.
你們出去的時候.
不要帶錢囊.
不要帶口袋.
不要帶鞋.
無論進到那一家.
才要說願這一家平安.
那裡若有當得平安的人.
你們所求的平安.
就必臨到那一家.

$^{1081}$不然就歸於你們了.
你們要住在那一家.
吃喝他們所供給的.
因為工人得工價是應當的.
然後第八節.
就要問你們了.
其實當我們讀到.
《路亞福音》這段經文的時候.
文字上我們都明白.
但在理性上.
我們很難接受.
我們出去宣教.
又不要帶錢囊.
又不要帶口袋.
又不要帶鞋.
我們每次出去.
裝備都非常.
是甲級裝備.
我就是甲級裝備.
丙級技術.
裝備好一點.
要裝備裝備裝備.
需要裝備嗎?需要.
但不需要那麼過份.
著急.
為什麼?因為我們要明白.
當我們出去做神的工作的時候.
誰會負責?.
你們都知道答案.
神會負責的.
你信不信?.
只有兩三個信.
信就要去走.
按人的常理.
常識.
其實很難做到.
有很多顧慮.
但關鍵是.
你相信耶穌的話嗎?.
你願意去經歷嗎?.

$^{1121}$但如果你願意經歷的話.
你相信聖經的話語.
如果你真的相信耶穌的話.
你會有福,你會蒙福.
因為我們的短宣主義就是這樣做.
以前我們在中國.
現在我們在土耳其.
短宣隊.
暑假又來了.
分組,然後打開地圖.
我們感動之後.
我們派不同的組別去不同的鄉下.
當然現在很方便.
不用走路,可以搭火車.
搭巴士去到.
當然你要走一段路,不是一步都不走.
去到村落,去到小鎮.
要入鎮.
怎麼辦?第一件事要做什麼?.
禱告.
祈求讓我們遇到應當遇到的人.
睜開眼睛.
啞門之後,應該遇到誰不知道.
憑信心就走.
我們就是這樣去宣教.
但每一次都經歷.
很多奇妙的經歷.
為病人的禱告.
有些簡直.
可以說是.
熟悉了裝甲,自己說要信耶穌.
每一團.
都有,好像我們剛剛過去.
今年年初.
去年有一個隊.
帶了兩個人去信.
在穆斯林地區,你去別人的鄉下.
很難想像.
他們都很難想像.
短訓隊都很難相信.

$^{1161}$他們想著隨意去應付一下.
就回來了,誰知道.
他們真的經歷,當人無助的時候.
他們才懂得依靠神.
我想這個就是.
我們的操練.
所以當他們.
相信耶穌的話語的時候.
他們願意去經歷的時候.
他們才是真正蒙福的.
最有效的宣教士訓練在哪裡.
不是神學院,雖然我做了很多年神學教育.
但我自己很深的體會.
最有效的宣教士訓練.
是要進行主耶穌的教導.
主怎麼說.
我們就怎麼做.
這個是我們要學習的.
進行主耶穌的教導.
才能夠完成大使命.
才能夠活出挑戰多.
恩典更多的精彩人生.
最後.
我們講到的是.
福音和祥是缺少工人.
要收的莊稼多.
作工的人少.
所以你們當求莊稼的主.
打發工人出去收他的莊稼.
其實很坦率來講.
宣教,猜拳.
歷來都是教會最不容易推動的.
核心事工.
請留意我的用詞.
是最不容易推動的.
卻是核心的事工.
缺少宣教工人.
更是自古就存在的問題.
其中的原因有很多.
比如說.

$^{1201}$從人的角度.
有很多的擔心和疑慮.
時機不成熟.
我們教會還很弱小.
其實這些都不是理由.
有些是因為對教會的角色和使命.
模糊不清.
失去了大使命的焦點.
有些是出於主耶穌.
對主耶穌和自身信仰.
冷淡退後的原因.
當我們躺平的時候.
當我們失去了這種迫切感的時候.
大使命關我什麼事.
牧師你在台上說說.
唱歌就一起唱.
唱完之後.
關我什麼事.
事實是面對廣闊的.
福音和祥.
工人欺缺.
而且分配非常不均勻.
我們看數據.
數據可以幫助我們.
了解這個情況.
這個皮尤研究中心.
這個都很出名的.
美國專門做這個.
Research的中心.
他發佈了一個數據.
讓我們很驚訝.
在過去的.
10年,2010年到2020年.
這個數據說.
基督徒的增長都是5\%.
你看這個數字.
還比較安慰.
這個數字就很令人驚訝.
穆斯林的增長率.
在過去的10年是20.7\%.

$^{1241}$佛教徒增長.
12\%.
無神論者增長16.5\%.
基督教增長率才是5\%.
雖然我們的基數比較大.
但如果數學家.
一計就計到了.
最新的報告已經出來了.
去到2050年.
回教徒的.
數字會超過基督徒.
2050年.
不是2025.
2050再過24年.
可能還會快.
這個世界人口評論.
最新的統計.
這個數字.
已經超過20億.
早前.
應該是去年年尾.
他們在馬來西亞有一個宣教大會.
穆斯林的宣教大會.
超過1萬多位.
宣教士出席.
出席.
穆斯林的宣教士.
只是很少的一部分.
而我們呢.
宣教大會.
穆斯斯姆長老.
不斷派擔獎鼓勵大家來啊來啊.
會眾加起來都不夠1萬.
幾千人.
參加一個大會已經很厲害.
但你看穆斯林.
只是宣教士.
穆斯林宣教士出席1萬.
所以這個數字.
很驚人.

$^{1281}$我們可以看.
其實這些網上全部都有.
如果我們有心去看的話.
花點時間去看.
我們就知道其實我們.
教會的狀況是怎樣.
分配不均勻.
而且最重要就是.
我看到一份香港的報告.
就是資源的分配.
也是不平衡.
我看到這個數字不錯.
它是全球基督徒的奉獻.
大概是一千三百億美金.
每年.
不錯啊.
一千三百億美金.
整個港幣的數量.
人民幣的數量很大.
但是只有3.2\%的奉獻.
是用在猜拳.
用在跨文化宣教的.
更少.
我深有體會.
因為我們機構都籌款的.
我們知道跨文化的宣教.
籌款確實很難.
但是有一件事我們可以看到.
其實我們當中有很多.
我想你們有做公司.
有做企業的.
我們做一個比喻.
如果一個公司一個企業.
包括教會都是.
財政的規劃.
其實我們看資金預算的流向.
我們就可以預測到你這個公司.
這個企業你工作的重點.
和方向是什麼.
國家也是如此.

$^{1321}$國家突然間軍工企業預算特別大.
講明是什麼.
講明未來要發展軍工事業.
這個很正常的一個比例.
同樣我們看教會的財政規劃.
我們可以看到教會.
是怎樣看待宣教事工.
其實答案都在我們心中.
回過頭來.
之前我們所講的.
宣教對於教會來講.
是不是只是一個配搭的事工.
如果是的話.
當然太公分豬肉.
分一塊給他就可以了.
但是如果我們真的很清楚.
大使命宣教是教會存在的目的.
和主要的方向的話.
你的資源的投資投入.
就不同了.
面對港府福音的禍場.
工人和資源的奇缺.
就好像主耶穌.
在門徒所吩咐的.
要收的莊稼多.
做工的人少.
所以當求莊稼主.
打發工人出去收他的莊稼.
所以求主幫助我們.
我知道時間到了.
其實我講這麼多你都明白.
我想講什麼.
再講多就慘了.
求主幫助我們.
端正我們的目標.
不要忽略了.
大使命在我們.
信仰當中的位置.
你其他都不記得.
不要緊 記得一句.

$^{1361}$我們千萬不要在.
神的恩典當中躺平.
希望我們能夠.
聖靈喚起.
我們的迫切感.
把握時機.
去收莊稼.
我們機構一二年成立.
十年之後神才開我們的眼.
去進入跨文化的領域.
這個是我們的突破.
你們很多.
知道我們的網頁.
可以上我們的網頁去看.
所以求主幫助我們.
讓我們能夠在未來的日子.
能夠.
行就主的心意當中.
在這個世界上.
有很多事我們可以做.
但有一樣事情.
是我們每一個屬主的門徒.
必須要做的就是大使命.
我們去參與.
扮演好我們的角色.
很多種又是一線的.
又是背後做後勤的都是工人.
求主來使用我們.
我們同心祈禱.
賜下你的話語.
我們很熟悉的經文.
你都求聖靈來喚醒我們.
讓我們能夠積極的參與.
宣教猜拳的事工.
不只是因為進行.
主耶穌的旨意.
而得蒙主的喜悅和紀念.
你都可以開闊我們的眼界.
少許的付出帶給別人.
意想不到的祝福.

$^{1401}$主耶穌說.
莊稼已經熟了可以收割了.
我們願意.
回應主的差遣.
用不同的方式去參與.
猜拳的工作.
倘若我們願意的話.
就用心靈和誠實.
對主耶穌說.
主啊我在這裡請差遣我.
請使用我.
求主使用我們.
在座每一位聆聽你的話語.
並且進行一天.
奉耶穌基督的聖名.
\newpage



\section{羅馬書 1:16-23}
\label{sec:032167}
\textbf{懺悔錄}~~~~梁誠德牧師
\newline
\newline
連結: \href{http://www.church.com.hk/acms/content.asp?site=cdc\&op=show\&type=product\&code=032167\&layout=sermon}{\texttt{032167}} ~~~~ 語音日期: 2026-05-17
\newline
\newline
\hyperref[sec:032166]{\small{< < < PREV SERMON < < <}}
~
\hyperref[sec:index_chronic]{\small{[返順時目]}}
~
\hyperref[sec:032168]{\small{> > > NEXT SERMON > > >}}
\newline
\newline
羅馬書 1:16-23
\newline
\begin{longtable}{cl}
\hline
\hline
章節 & 經文 (和合本修訂版)\\
\hline
1:16 & \begin{tabularx}{0.7\textwidth}{X} 我不以福音為恥；這福音本是神的大能，要救一切相信的，先是猶太人，後是希臘人。 \end{tabularx} \\ \\ \relax
1:17 & \begin{tabularx}{0.7\textwidth}{X} 因為神的義正在這福音上顯明出來；這義是本於信，以至於信。如經上所記：「義人必因信得生。」 \end{tabularx} \\ \\ \relax
1:18 & \begin{tabularx}{0.7\textwidth}{X} 原來，神的憤怒從天上顯明在一切不虔不義的人身上，就是那些行不義壓制真理的人。 \end{tabularx} \\ \\ \relax
1:19 & \begin{tabularx}{0.7\textwidth}{X} 神的事情，人所能知道的，原顯明在人心裡，因為神已經向他們顯明。 \end{tabularx} \\ \\ \relax
1:20 & \begin{tabularx}{0.7\textwidth}{X} 自從造天地以來，神的永能和神性是明明可知的，雖然眼不能見，但藉著所造之物就可以了解看見，叫人無可推諉。 \end{tabularx} \\ \\ \relax
1:21 & \begin{tabularx}{0.7\textwidth}{X} 因為，他們雖然知道神，卻不把他當作神榮耀他，也不感謝他。他們的思想變為虛妄，無知的心昏暗了。 \end{tabularx} \\ \\ \relax
1:22 & \begin{tabularx}{0.7\textwidth}{X} 他們自以為聰明，反成了愚昧， \end{tabularx} \\ \\ \relax
1:23 & \begin{tabularx}{0.7\textwidth}{X} 將不能朽壞之神的榮耀變為偶像，仿照必朽壞的人、飛禽、走獸、爬蟲的形像。 \end{tabularx} \\ \\
[1ex]
\hline
\hline
\end{longtable}
$^{1}$張超雄議員.
主席.
\newpage



\section{使徒行傳 1:6-11}
\label{sec:032168}
\textbf{耶穌升天}~~~~黃健羒牧師
\newline
\newline
連結: \href{http://www.church.com.hk/acms/content.asp?site=cdc\&op=show\&type=product\&code=032168\&layout=sermon}{\texttt{032168}} ~~~~ 語音日期: 2026-05-17
\newline
\newline
\hyperref[sec:032167]{\small{< < < PREV SERMON < < <}}
~
\hyperref[sec:index_chronic]{\small{[返順時目]}}
~
\hyperref[sec:032169]{\small{> > > NEXT SERMON > > >}}
\newline
\newline
使徒行傳 1:6-11
\newline
\begin{longtable}{cl}
\hline
\hline
章節 & 經文 (和合本修訂版)\\
\hline
1:6 & \begin{tabularx}{0.7\textwidth}{X} 他們聚集的時候，問耶穌：「主啊，你就要在這時候復興以色列國嗎？」 \end{tabularx} \\ \\ \relax
1:7 & \begin{tabularx}{0.7\textwidth}{X} 耶穌對他們說：「父憑著自己的權柄所定的時候和日期，不是你們可以知道的。 \end{tabularx} \\ \\ \relax
1:8 & \begin{tabularx}{0.7\textwidth}{X} 但聖靈降臨在你們身上，你們就必得著能力，並要在耶路撒冷、猶太全地和撒瑪利亞，直到地極，作我的見證。」 \end{tabularx} \\ \\ \relax
1:9 & \begin{tabularx}{0.7\textwidth}{X} 說了這些話，他們正看的時候，他被接上升，有一朵雲彩從他們眼前把他接去。 \end{tabularx} \\ \\ \relax
1:10 & \begin{tabularx}{0.7\textwidth}{X} 他升上去的時候，他們定睛望天，看哪，有兩個人身穿白衣站在他們旁邊， \end{tabularx} \\ \\ \relax
1:11 & \begin{tabularx}{0.7\textwidth}{X} 說：「加利利人哪，你們為甚麼站著望天呢？這離開你們被接升天的耶穌，你們見他怎樣升上天去，他也要怎樣來臨。」 \end{tabularx} \\ \\
[1ex]
\hline
\hline
\end{longtable}
$^{1}$剛剛過去的禮拜四是復活節之後的四十一.
也是傳統教會紀念耶穌升天的日子.
因為升天日通常都是禮拜四.
所以大家平日上班的時候.
一般教會不會有什麼特別的聚會.
因為大家來不及.
不過這是提醒我們.
耶穌基督是結束了在地上的時宮.
升到天上.
耶穌升天是表明.
從此以後沒有人在地上看見祂在走動.
或者耶穌在哪裡.
聽祂說話.
也不能夠有人宣稱.
耶穌的化身.
再來世界.
第一次基督來的時候就是猶太人耶穌.
現在這個時代就是我了.
基督來到了.
我就是祂的化身.
我今天是在韓國的.
韓國真的有些邪教是這樣宣稱的.
第一次基督來的時候就是猶太人.
第二次基督來的時候就是韓國人.
我們不是這樣看.
我們認為耶穌基督是同一位耶穌.
祂來到世界上面.
也是同一位耶穌會升天.
將來來到世界.
再來世界的時候.
也是那位耶穌.
所以《約翰福音》第三章第十三節.
除了從天降下的人子.
沒有人升過天.
這裡是指人子耶穌基督.
本來是來自天上的.
好像剛才我們所唱的詩歌最後那一首.
原本就是由天上來的.
祂來到之後.
祂就是神的話語.

$^{41}$是神的道.
這個道.
藉著瑪利亞生來成為人.
在世界上度過了三十多年.
升天.
就是回到祂原來充滿榮耀的地方.
《哲堂傳》第七章第五十六節也說.
是基督教會第一個殉教士.
殉道士斯提法.
他快被人用石頭打死的時候.
他說.
我看見天開了.
人子站在神的右邊.
如果你問我.
神是怎樣的.
祂究竟坐的位置是大理石還是木頭.
耶穌升天之後.
祂穿什麼顏色的衣服.
什麼款式.
如果你問我這些問題.
我就沒有答案.
我不知道.
我也不會問這些問題.
但是我會問.
究竟耶穌基督升天.
對我們今天作為一個信徒.
有什麼意義.
對我今天的生活.
對我今天的生命.
對我的信仰.
有什麼影響.
這個是我的問題.
首先我覺得.
耶穌升天帶給我們永生的盼望.
永遠生命的盼望.
上兩個主日.
我們讀聖經關於耶穌被釘十字架.
之前跟門徒所說的話.
門徒那時候是怎樣.
是憂愁.

$^{81}$因為耶穌三番四次都說.
祂就快要死了.
就快被人捉了.
被人釘十字架了.
這班門徒很擔心.
因為過去那三年多.
跟著耶穌去哪裡.
他們就去哪裡.
要做什麼.
就問先生.
他就告訴你.
怎樣說.
去哪裡.
耶穌就告訴他們.
所以三年多.
他們都是以耶穌為頭.
他們是一些門徒跟著祂.
如果耶穌快要死了.
那怎麼辦.
我們做什麼.
去哪裡.
所以憂愁.
好像說沒有了頭.
真的散了.
不知道要做什麼.
但是.
門徒很歡.
跟著耶穌說.
我會再回來.
很快回來.
跟著過了幾天.
哇 耶穌復活了.
他們那時候很開心.
耶穌復活就是說.
他們以後不用擔心了.
耶穌仍然和他們同在.
仍然可以帶領他們.
有什麼話仍然可以問祂.
當然是通過聖靈對他們的感動等等.
所以門徒那時候.

$^{121}$是有一個這樣的轉變.
耶穌勝過死亡.
祂所教導.
所談論的.
信我的人永遠不死.
這就再清楚確定他們的信心.
所以很開心.
我們的生命也是這樣.
因為相信耶穌基督的話.
我們應該是一個喜樂的生命.
信耶穌的人.
看生命是一個直線.
我們開始的時候.
就是在母親的懷裡面.
承孕開始.
有身體.
有靈魂.
有一天我們會死亡.
我們的身體就會腐爛.
沒有了.
但是我們的靈魂.
繼續.
存到永遠.
當然這是很簡單的一句表示.
聖經也有說.
我們離開世界之後.
有一天要復活.
有一天要審判.
主宰我們等等.
這些我們就要.
但是.
始終我們覺得生命是一個直線.
所以我們不相信輪迴.
不相信我們變成另外一個生命.
變成另外一個人.
或者變成另外一個存在.
我們仍然是.
所以耶穌復活之後.
祂將手上的釘痕給門徒看.
也叫門徒.

$^{161}$特別是多瑪.
你不相信我.
你伸手指來探入我的肋旁的洞.
有支槍刺了它.
有個洞在那裡.
你摸一下.
我不知道.
三天會不會結了個疤.
或者怎麼樣.
你不知道.
但是那個意思就是說.
我仍然是那個同一個人.
釘痕還在.
肋旁的傷口還在.
我仍然是那個耶穌.
被釘十字架的那個耶穌.
而且祂在門徒面前吃東西.
這個更加難解釋.
我們也不問這些問題.
只不過重點就是說.
那個仍然是耶穌.
所以這段聖經也是說.
當時他們看著耶穌升天.
旁邊有兩個穿白色衣服的就是天使.
跟他們說.
耶穌怎樣被接上天.
將來也都會怎樣再回來.
所以重點不是說.
祂再回來的時候.
一定在以色列.
或者在巴大利那個地方.
要不然我們要坐飛機去那裡.
重點是說.
同一位耶穌.
祂從永遠到永遠.
都是存在的耶穌.
最後一天祂要回來.
如果是這樣的解釋.
這樣的了解的話.
我們今天仍然活著.

$^{201}$今天怎樣生活呢?.
如果我仍然是我.
不會變成另一個人.
我的生活態度是怎樣呢?.
就是說一個負責任的生活.
很嚴謹地面對每一天.
因為我們的品價.
我們的道德判斷.
永遠成為我的一部分.
不會變的.
成為你的一部分.
今天的說話,行為,思想.
成為我的一部分.
不會刪除的.
想一想都有點害怕.
今天做過的事,說過的說話.
成為我的一部分.
不可以被消滅.
特別是馬太福音第十二章說.
神審判的時候我們說過的說話.
句句都要供出.
憑我們的說話.
判斷我們是義人還是罪人.
害不害怕?.
真的有點害怕.
我們有沒有說過別人的閒話?.
有.
我既會彈又會唱.
有些人都是這樣的.
我們說彈人,批評人.
或者唱.
唱壞他們.
如果你像我一樣.
既會彈又會唱的時候.
真的有點害怕.
如果有一天.
耶穌來的時候.
審判說我說過的說話.
怎樣去批評人.
全部都記錄在裡面.

$^{241}$現在就說.
真的有些可以想像.
你說按個鍵就刪除了.
不是的.
在雲端裡面.
你不能夠刪除的.
在那裡的.
你不懂得怎樣去取回.
這個時候我們可以想像.
真的有可能.
我們所做過的一切.
所說的話.
全部都是記錄在裡面.
害不害怕?.
不害怕.
因為聖經說.
羅馬書第八章.
這一章的聖經.
通常我提外話.
是聖經所謂百答的章節.
你人生無論遇到什麼問題.
信仰搖動也好.
被人欺負也好.
什麼都是.
你回去看看羅馬書第八章.
看完一次的時候.
你就有很大的安慰.
今天我只提這個沙爾節.
上帝既不顧惜自己的兒子.
即是耶穌基督.
為我們眾人捨了他.
豈不也把萬物和他.
一同白白地賜給我們?.
誰能控告上帝所揀選的人?.
就是我們.
有上帝稱他們為異了.
誰能定他們有罪?.
有基督耶穌已經死了.
而且復活了.
現今在上帝的右邊.

$^{281}$也替我們祈求.
所以不用怕.
耶穌已經替我們死.
替我們贖罪.
現今在上帝的右邊.
替我們祈求.
只要我們知道自己犯罪的時候.
馬上承認.
認罪悔改.
耶穌會替我們向神求情.
他剛剛信主而已.
他不太懂.
原諒他吧.
或者他信了主很多年.
不過現在有點不小心.
或者人生遇到一些事情.
原諒他吧.
我們有一位主幫我們求情.
神一定會赦免我們.
聖靈也會幫助我們改過.
所以如果遇到同樣的試探.
來到的時候.
聖靈會提醒我們.
上次這個事情出現過了.
它來過了.
小心一點.
我們就可以勝過這個試探.
更加是充滿喜樂.
不僅不怕.
還充滿喜樂.
有盼望面對將來.
我媽媽在2018年在香港過世.
之前一兩年.
每逢說電話.
或者我們回到香港看她的時候.
她都是說這句話.
她說我很感恩.
現在什麼都有了.
兒子大,女兒大.
連孫子都大了.

$^{321}$我什麼都不用擔心.
你們也不用擔心.
我準備好了.
天父何時來接我走.
我就走了.
她離開世界的時候是96歲.
我們全家都相信.
她是回到天父那裡.
她很安心.
她相信耶穌所說的話.
耶穌說我去完事為你們預備地方.
我必再來接你們到我那裡去.
我在那裡.
叫你們也在那裡.
所以有一天我們會離開這個世界.
我們會回到父神那裡.
耶穌已經為了我們預備好地方.
所以什麼行李都不用帶.
也不能夠帶.
就算你能夠帶.
也未必合用.
聖經也說.
神為愛他的人所預備的.
是眼睛未曾看見.
耳朵未曾聽見.
人心也未曾想過.
想都未想過.
那你就不用想了.
只知道什麼都預備好了.
什麼都預備好的時候.
你真的什麼都不需要帶.
到時就搬去那裡就行了.
這是耶穌基督升天的時候.
帶給我們一個永恆生命的盼望.
我們這個世界上面.
短暫的時間是會離開.
但有一天我們會回到天父永恆的家中.
第二就是提醒我們有一個見證的責任.
第八字這裡說.
聖靈降臨在你們身上.

$^{361}$你們就必得著能力.
並要在耶路撒冷.
猶太傳地和撒瑪利亞.
直到地極作我們的見證.
如果看地圖的時候.
耶路撒冷再遠一點就是猶太傳地.
然後到撒瑪利亞北面的地方.
去到加利利.
到世界的地方作它的見證.
路加在這裡的記載.
比馬太稍為簡單一點.
沒有大使命那些.
但特著一個大使命的重點.
就是作耶穌的見證.
更加準確一點就是說.
是作耶穌復活的見證.
耶穌復活.
升天證明祂是基督.
是神的兒子.
祂就是神.
耶穌復活.
升天也表明祂來到世界上的任務完成了.
回到天上.
這個任務就是完成救贖.
凡是信祂的人都可以蒙恩.
聚得赦免.
並且成為神的兒女.
我們作為見證人.
就是將我們自己蒙恩得救的故事.
告訴別人.
正如在法庭上.
如果你是做證人的話.
將你所看到的.
你感受到的.
再按照事實回答.
多餘的東西不用說.
Tell the truth.
Nothing but the truth.
只是說你知道的.
你看見的.

$^{401}$說了出來.
通常那些.
如果你有做過證人.
就知道一些律師,法官提醒你.
你的答案只有三個.
一是.
二不是.
三不知道.
就行了.
當然有要求.
或者解釋一下另外一件事.
但基本上很簡單.
你體驗的就是.
有沒有看見?有.
有沒有做過?有.
沒有就沒有.
或者不知道.
其實意思就是說.
倘若.
如果我們沒有經歷過的時候.
沒有經歷在神面前知道自己有罪.
認罪悔改.
我們沒有辦法向人傳講這個.
如何解釋什麼叫做罪德赦免.
你不解釋得到.
你解釋得到只不過是.
諸葛學先生或者看書教你怎麼說.
但你說的東西.
你沒有經歷過.
你說不到的.
如果我們自己沒有感受到.
你不能夠回答這個問題.
但如果你有的話.
就簡單了.
別人問你.
是不是這樣?是.
有沒有經歷過?有.
如果沒有就沒有.
不知道就是不知道.
所以不用怕.

$^{441}$如果他問我聖經的問題.
神學問題.
我不懂得回答怎麼辦?.
有些人怕.
你就說不知道.
我也有很多事情不知道.
你千萬不要以為.
我信主月來.
這些問題解決了.
就沒有問題了,不是的.
你會越來越多問題.
因為好像越挖越深.
很多的問題.
你想到看聖經.
第一次看不到的問題.
第二次你會看到.
第三次你會看到.
所以問題會越來越.
你不知道就不知道.
所以我們在神經面前.
如果是為主作見證.
不用擔心.
你經歷過,你就說.
我經歷過神的恩典.
經歷過我的聚德赦免.
經歷過豐盛的生命.
經歷過神的同在.
甚至有些人未必明白.
如果不是很明白.
你就不用說了.
譬如有些人講方言.
聖靈充滿.
如果聽的人不明白.
你不用說這些.
如果他明白.
你可以說的.
可以分享.
在亞洲斯里蘭卡.
即是以前的錫蘭.
即是印度南面的大島.

$^{481}$斯里蘭卡有一位牧師.
DT Niles.
他說過一句名言.
Evangelism is just one beggar telling another beggar.
where to find bread.
傳福音就是一個乞丐.
告訴另一個乞丐.
哪裡可以找到麵包.
大家可以想像.
斯里蘭卡,印度,巴基斯坦.
即南亞很多地方.
到處都有乞丐.
如果你是做遊客的時候.
一下車.
一群乞丐圍著你.
大的小的問你拿錢.
生命最基本的需要就是麵包.
吃的東西.
也是很多人最起碼的期望.
今天如果有得吃就好了.
如果今天找到麵包就好了.
這個是什麼意思呢?.
我們傳福音的時候.
不是有個態度.
我比你聖潔.
Holier than you.
不是說我什麼都懂.
不是說我比較熟悉.
在道德上面.
我比較高一級.
現在我們告訴你.
你應該怎麼做.
應該怎麼信.
傳福音不是這樣的態度.
而是說.
我曾經也是這樣.
我也是貧窮.
我也是有缺乏.
但現在我找到食物.
找到耶穌.

$^{521}$你去祂那裡.
都可以得到生命所需要.
好像這裡.
約翰福音第六章.
耶穌說.
我就是生命的糧.
或者是生命的麵包.
到我這裡來的絕不飢餓.
信我的永不乾渴.
耶穌基督那裡.
才有我們生命所需要的.
去到祂那裡.
就可以得到豐盛的生命.
得到生命的需要.
那人家去不去祂那裡.
這些我們不能夠強迫.
你最多是為祂祈禱.
希望祂去到耶穌那裡.
得到生命的需要.
但你不能夠強迫祂.
也不能夠騙祂.
去耶穌那裡吧.
你的問題都可以解決的.
你的病又得到醫治了.
你找工作又可以有工作了.
不是這樣的福音.
而是說.
我自己以前的需要.
我生命上面的經歷.
現在在耶穌那裡.
得到解決.
我們以前是怎樣.
可能都是有缺乏的.
可能我們是有錢.
但買不到法樂.
錢買不到法樂.
或者我們事業上成功了.
但是不滿足.
你有一間店.
有一間店舖.

$^{561}$又有第二間店舖.
又有一疊樓.
又有第二間.
又有第二疊樓.
但是不滿足.
好像當然是越來越多更好.
或者別人稱讚你.
你真的有福了.
有兒有女又有孫.
你自己才知道.
兒女是不是時時跟你有互動.
會不會是.
真的到了母親節,父親節.
或者你生日.
如果搬到外地的話.
打個電話回來.
恭喜一下.
平時沒有來往.
你都不知道他做什麼.
他又不知道你做什麼.
別人說你有福.
但是你自己才知道.
是不是福.
其實有電話.
寄張卡來都不錯.
現在是歷史時代.
又或者你有工作.
有收入.
但是你每一天都是做同樣的事.
都是同樣的活動.
同樣的動作.
沒有什麼選擇.
沒有什麼出路.
到退休了.
多好.
又有社會安全了.
基本上生活上沒有問題.
但是剩下的日子怎麼樣.
一天過一天.
很多人都是退休的時候.

$^{601}$真的又孤單.
又不知道做什麼.
我希望你不要糊塗.
一天過一天.
一定要有一個目標.
充實的過每一天.
在耶穌基督裡面.
就有豐盛的生命.
這個生命是怎麼樣.
我們存著感恩的心來吃飯.
享受神所賜各樣美好的事物.
好像我剛才說.
我媽媽這樣.
每一天都是感恩.
真的.
如果你信了耶穌之後.
每一天起床的時候.
能夠呼吸.
在這個地方.
沒有戰亂.
你知不知道打仗的地方.
那些人怎麼樣.
整晚睡不著.
突然之間.
砰一聲.
好像我們有時候下大雨.
雷雨.
打雷雨.
整晚都是這樣.
而且不是一晚.
你想想烏克蘭那些居民.
已經幾年了.
你想想中東那些人.
幾年都是這樣.
時時都會有驚醒.
沒有睡好.
如果你有一覺好睡.
一覺睡到天亮.
感恩的心.
如果你是經濟獨立.

$^{641}$過一個負責任的生活.
不用靠救濟.
不用靠別人幫忙.
不用借錢.
也是神的恩典.
在教會服侍.
在社會上的貢獻.
在家庭裡面能夠照顧家人.
這些都是為了創造主而做的.
我們的努力.
永遠被神紀念.
這是豐富的信念.
並且我們三個教會.
在T4C裡面.
我們幾多弟兄姊妹在一起.
有什麼事打個電話.
他們可以幫忙.
他們為我們祈禱.
我們這些是.
在神裡面.
在耶穌基督裡面的恩典.
如果你找到.
如果你經歷過.
跟別人來傳講.
就是見證.
就是講你的故事.
你的故事.
你的經歷.
在座當中.
我相信所有相信耶穌的人.
都是有你的故事.
每一個人的故事都不一樣.
如果你能夠誠懇地跟別人分享.
這個就是做耶穌基督的見證.
如果你還未接受耶穌.
做你生命的救主.
但有興趣聽我們的故事.
你可以在散會問.
隨便一個教會的會友.
你怎樣信耶穌的?.

$^{681}$他將他的故事講給你聽.
每一個人都有一個得救的故事.
每一個人有一個轉變的故事.
英國的女王.
在2032年過世.
在英屬香港的時候.
叫她老闆娘.
她最受人尊敬的地方是什麼呢?.
就是她將所有的事情.
看為一種責任.
她做的事情是一種責任.
大家知道做皇帝,做女王.
不是那麼舒服的.
沒有了自由.
你一天下來見什麼人.
去什麼地方.
全部都幫你安排了.
你沒有選擇了.
所以做皇帝.
一舉一動.
別人又拍你照片.
而且那些鏡頭那麼長.
那麼大.
你都不知道他在那裡.
拍了你的照片.
你說什麼話.
錄下來.
你又不敢亂說話.
行動都要小心一點.
沒有了自由.
真的好像綁手綁腳.
特別是女王經歷二次大戰.
還有經濟危機,冷戰.
很多林森的事情在她任內那麼多年期內都發生.
但是別人給她評價.
她在每一個危機當中.
她都是為了國家的榮耀.
為了結合全國的公民.
代表國家履行這個責任.
她沒有投訴.

$^{721}$這是她的責任.
她沒有說這件事太難做了.
我就退位了.
讓兒子做吧.
她繼續做到年老的時候.
我們事奉神見證主.
也應該想到這是我們的責任.
應該做.
好像保羅教導提摩太也一樣.
這個責任已經交付給我.
你喜不喜歡順利不順利.
盡力去做就行了.
你不要問那麼多.
不要擔心別人的反應.
不要擔心你講了,你分享過.
別人信不信耶穌.
這是他和神之間的問題.
我們要不要得到這個救恩.
我們自己可以很肯定.
別人要不要是他們的決定.
這是我們傳福音.
我們見證主應該有一個心態.
但起碼我們盡了自己的責任.
有什麼機會,有什麼場合我們可以見證.
不一定每次都是由頭講起.
只要講你自己的故事就行了.
不需要依據聖經怎麼講.
別人有興趣的是你怎樣去信耶穌.
你的生命有什麼改變.
信了耶穌之後.
在你的生命裡面有什麼福氣.
這是我們做見證的責任.
第三方面就是耶穌升天.
提醒我們自己有一個尊貴的身份.
千萬不要以為神在他榮耀當中.
他只是隨隨便便.
找些人好像執醜一樣.
「保羅,你又做宣教了」.
「你是做什麼呢?」.
神不是這麼隨便的.

$^{761}$而是他想過了.
看中你,揀選了你.
在耶穌基督裡面揀選了你.
才有這樣的一個遺訓.
所以如果你有成為神的兒女.
千萬不要隨隨便便.
這是你尊貴的身份.
在這個過程當中.
我們叫做神的呼召.
不是隨便的.
是刻意的.
揀選了你.
《羅馬書》第八章第十六節.
又是第八章.
聖靈和我們的心同證.
我們是神的兒女.
聖靈自己與我們的靈一同見證.
我們是神的兒女.
既然是兒女.
就是後子.
就是神的後子.
和基督同作後子.
第三十節亦都是.
祂所預定的人.
祂又召他們來.
所召來的人.
祂又稱他們為兒.
所稱為的兒.
又會叫他們得回要.
想想我們和基督耶穌.
一同成為神的後子.
當然我們說.
耶穌是神的獨生子.
我們就是因為祂領養我們回來.
有少許分別.
但是基本上我們和耶穌基督.
是同樣成為神的兒女.
享受神家庭裡面的福氣.
享受祂的同在.
有這個遺訓在當中.

$^{801}$我們得到神的保護.
得到神的照顧.
好像這裡所說我們得稱為兒.
算為無罪.
在神的眼中算是一個好人.
並且會得到榮耀.
我想在座當中.
大部份人不是在美國出世.
所以沒有資格做總統.
美國出世才可以做總統.
但是可能很多人都像我一樣.
經過申請入籍的過程.
成為公民.
我們留在這裡.
我記得當時在芝加哥我們入籍的時候.
大概有一百人左右.
在法官宣誓.
法官在宣誓之前說了幾句話.
從今以後你不要以為自己是二等公民.
你和美國出生的美國公民是一樣的.
你有一樣的權力.
當然也有一樣的義務.
但是你是受法律的保障.
所以你不是二等的公民.
你就是公民.
和他們一樣都是公民.
那時候感覺上真的鬆了口氣.
以後不需要再擔心.
簽證過期.
不用擔心什麼事就遞解出境.
有法律的程序.
有法律的保護.
成為公民就是有這樣的身份.
我們在神面前是什麼身份?.
彼得前書第二章說.
你們是被揀選的一族.
是君尊.
即是所謂皇家.
皇家的祭司.
是神聖的國度.

$^{841}$是屬於上帝的子民.
是神國度裡面屬於神的公民.
所以我希望大家在神的面前.
能夠有這樣的自尊.
我們皇家的祭司.
是服侍萬王之王.
萬主之主.
我們一起的時候.
是神的國度.
我們是公民.
祭司和子民.
是有兩個很主要的身份的象徵.
祭司的標記是什麼?.
事奉.
以前十二支派當中.
利美的支派.
他們不用去耕田.
是其他的支派奉獻.
來支持他們的生活.
他們主要是什麼?.
就是事奉神.
有了會幕之後.
有了聖殿之後.
他們就專門在聖殿當中.
來事奉神.
所以以弗所書第四章第八節說.
所以有話說.
得升上高天的時候.
努力來俘虜.
將各樣的恩賜想吸引.
他所賜的有仕徒.
有先知.
有傳福音的.
有牧者和教師.
為要裝備聖徒.
做事奉的工作.
建立基督的身體.
後面那幾節.
大家應該會背了.
今年教會的主題經文.

$^{881}$今早我去洗手間的時候.
洗手間出來的時候.
就貼在這裡.
他所賜的有仕徒.
有先知.
有傳福音的.
有牧師.
有教師.
為要裝備聖徒.
來做事奉的工作.
建立基督的身體.
耶穌升天提醒我們.
要使用各樣的恩賜.
在各樣的事情上.
來事奉神.
建立基督的身體.
就是教會.
這個也是.
哥林多前書第十二章.
所提到的.
同一個身體上.
不同的部分.
不同的肢體.
要彼此配搭事奉.
建立這個身體.
事奉神不一定要動手.
要做很多事.
為教會祈禱.
為弟兄姊妹祈禱.
為你的家庭.
同事.
鄰居.
祈禱.
也是事奉.
事奉有很多的方式.
如果你還行得到.
動得到.
當然是做.
但是未必做到的時候.
有很多方式可以事奉神.

$^{921}$也不一定是在教會裡面.
在外面.
在你家裡面.
工作的地方.
都可以事奉神.
這是一個.
我們做祭司的標記.
事奉神.
我們的一生.
就是為了事奉神.
做子民的標記.
做公民的標記.
是甚麼呢?.
就是追求屬天的事.
二弗所書第二章說.
祂又使我們和耶穌基督.
一同復活.
一同在在天上.
這不是說將來的復活.
我們和耶穌一樣.
同一個時間復活了.
一同在在天上.
當然了.
時間上.
我們不是說回到以前.
就是在意義上.
或者地位上.
我們是一同在在天上.
有崇高的地位.
哥羅世書第三章說.
既然你們已經.
與基督一同復活.
就當求上面的事.
那裡有基督.
在在上帝的右邊.
你們思考上面的事.
不要思考地上的事.
思考上面.
因為我們的地位上.
已經是在天上.

$^{961}$和耶穌基督一起.
所以這兩段經文.
都是提醒我們.
在這個時代.
這個世界裡面.
我們已經是.
心態上面.
已經是和耶穌基督在天上.
想到天上的事.
地上的事是甚麼?.
有錢就買豪宅.
開個靚車.
穿個名牌.
要賺更多的錢.
擁有更多.
如果我有一百萬.
我就買這些東西.
擁有這些東西.
這是世上的事.
天上的事是.
你有基本生活的需要.
就知足了.
然後將錢財.
用在神的國度裡面.
用在福音的事供.
用在那些有需要的人身上.
將你的財富放在天上.
所以有少少分別.
地上的事是.
在親朋好友的面前.
你可以誇口.
別人就讚賞.
真是厲害.
成功.
有名氣.
被人看一張卡片.
印得密密麻麻.
又甚麼董事長.
天上的事就是.
不在乎別人的稱讚.

$^{1001}$我不是說拒絕.
有人稱讚你.
多謝一聲.
但不在乎他們稱讚.
有沒有人看到.
不要緊.
有沒有人讚.
不要緊.
只是在乎主說.
你一個忠心善良的僕人.
這是天上的事.
我們要多些思考.
地上的事是甚麼.
福祿壽而康.
有沒有看你的筷子.
印了福祿壽而康.
或者五福臨門.
以前我都說過.
你上網查一查.
知道甚麼是五福.
地上的事是.
如果這樣.
我們就有福了.
有錢.
有地位.
有兒女.
有孫.
天上的事就是.
要認識神.
遵行神的話.
愛神愛人.
與主同行.
多些思考.
神是怎樣.
多些認識神.
地上的事是.
強調你擁有多少身家.
有多少的靈.
有多少的本事.
做多少的事.

$^{1041}$天上的事是.
強調你是怎樣的人.
你的品格.
成名的過程.
有沒有在你的生命裡.
所以地上的事.
是看你擁有多少.
Having.
或者你有多少能力.
Doing.
你做了多少事.
幾個本事.
但是天上的事是.
你是怎樣的人.
Who you are.
你是怎樣的人.
你的品格是怎樣.
所以有分別的.
聖經教導我們.
不要想太多地上的事.
多些想想天上的事.
因為我們是屬於.
上帝的子民.
是公民.
是天國的公民.
所以你有沒有這樣的標記.
你有沒有這樣的身份.
被揀選的一族.
是皇家的祭司.
是神聖的國度.
是屬神的子民.
所以今天我們紀念.
耶穌基督復活之後.
四十天.
升天.
我們相信主有一天.
要再來.
第一次來.
祂是為了救贖萬民.
救贖全世界的人.

$^{1081}$祂再來的時候.
就是審判的機會.
審判的時候.
我們不一定.
會等到祂再來.
但是無論什麼時候.
我們都不用怕.
不用怕被定罪.
或者是刑罰.
耶穌已經救贖了我們.
只要我們持定我們的信仰.
靠著聖靈的力量.
來遵行神的話.
我們可以活在.
永恆的生命當中.
我們仍然活在.
這個世界上.
一定要去履行我們的責任.
見證神的愛.
見證主的復活.
我們要珍惜我們的身份.
更多的事奉神.
更多的思考上面的事.
不要思考地上的事.
以前我們在世界上生活的時候.
我們依然享受到.
豐盛的生命.
所以永遠的生命.
不是說我們死後才開始.
是現在就開始了.
現在這個生命.
就可以傳到永恆當中.
但願我們能夠在主的裡面.
有這樣的攀望.
有這樣的喜樂.
我們一起同心祈禱.
在這個世界上.
我們不是糊塗的.
一天過一天就算了.
但是為了你的緣故.

$^{1121}$為了天國的緣故.
我們能夠怎樣去計劃.
我們的一生.
計劃我們這一天要做的事.
求主幫助我們.
不一定要做很多事.
而是說我們的生命當中.
要成為一個怎樣.
你所喜悅的人.
求主幫助我們.
藉著耶穌基督升天的一個紀念.
我們有永遠的攀望.
亦都是提醒我們自己.
有一個屬天的生命.
有屬於天國裡面.
公民的一個責任.
求主去做憐憫.
垂聽我們的禱告.
付諸耶穌基督的聖願祈求.
多謝.
\newpage



\section{以弗所書 6:10-24}
\label{sec:032169}
\textbf{在基督裡作剛強的人}~~~~葉慧敏傳道
\newline
\newline
連結: \href{http://www.church.com.hk/acms/content.asp?site=cdc\&op=show\&type=product\&code=032169\&layout=sermon}{\texttt{032169}} ~~~~ 語音日期: 2026-05-17
\newline
\newline
\hyperref[sec:032168]{\small{< < < PREV SERMON < < <}}
~
\hyperref[sec:index_chronic]{\small{[返順時目]}}
~
\hyperref[sec:032170]{\small{> > > NEXT SERMON > > >}}
\newline
\newline
以弗所書 6:10-24
\newline
\begin{longtable}{cl}
\hline
\hline
章節 & 經文 (和合本修訂版)\\
\hline
6:10 & \begin{tabularx}{0.7\textwidth}{X} 最後，你們要靠著主，依賴他的大能大力作剛強的人。 \end{tabularx} \\ \\ \relax
6:11 & \begin{tabularx}{0.7\textwidth}{X} 要穿戴神所賜的全副軍裝，好抵擋魔鬼的詭計。 \end{tabularx} \\ \\ \relax
6:12 & \begin{tabularx}{0.7\textwidth}{X} 因為我們的爭戰並不是對抗有血有肉的人，而是對抗那些執政的、掌權的、管轄這幽暗世界的，以及天空靈界的惡魔。 \end{tabularx} \\ \\ \relax
6:13 & \begin{tabularx}{0.7\textwidth}{X} 所以，要拿起神所賜的全副軍裝，好在邪惡的日子能抵擋仇敵，並且完成了一切後還能站立得住。 \end{tabularx} \\ \\ \relax
6:14 & \begin{tabularx}{0.7\textwidth}{X} 所以，要站穩了，用真理當作帶子束腰，用公義當作護心鏡遮胸， \end{tabularx} \\ \\ \relax
6:15 & \begin{tabularx}{0.7\textwidth}{X} 又用和平的福音當作預備走路的鞋穿在腳上。 \end{tabularx} \\ \\ \relax
6:16 & \begin{tabularx}{0.7\textwidth}{X} 此外，要拿信德當作盾牌，用來撲滅那惡者一切燒著的箭。 \end{tabularx} \\ \\ \relax
6:17 & \begin{tabularx}{0.7\textwidth}{X} 要戴上救恩的頭盔，拿著聖靈的寶劍—就是神的道。 \end{tabularx} \\ \\ \relax
6:18 & \begin{tabularx}{0.7\textwidth}{X} 要靠著聖靈，隨時多方禱告祈求，並要為此警醒不倦，為眾聖徒祈求。 \end{tabularx} \\ \\ \relax
6:19 & \begin{tabularx}{0.7\textwidth}{X} 也要為我祈求，讓我有口才，能放膽開口講明福音的奧祕， \end{tabularx} \\ \\ \relax
6:20 & \begin{tabularx}{0.7\textwidth}{X} 我為這福音的奧祕作了帶鐵鏈的使者，讓我能照著當盡的本分放膽宣講。 \end{tabularx} \\ \\ \relax
6:21 & \begin{tabularx}{0.7\textwidth}{X} 今有親愛、忠心服事主的弟兄推基古，為了你們也明白我的事情和我的景況，他會讓你們知道一切的事。 \end{tabularx} \\ \\ \relax
6:22 & \begin{tabularx}{0.7\textwidth}{X} 我特意打發他到你們那裡去，好讓你們知道我們的情況，又讓他安慰你們的心。 \end{tabularx} \\ \\ \relax
6:23 & \begin{tabularx}{0.7\textwidth}{X} 願平安、慈愛、信心從父神和主耶穌基督歸給弟兄們。 \end{tabularx} \\ \\ \relax
6:24 & \begin{tabularx}{0.7\textwidth}{X} 願所有恆心愛我們主耶穌基督的人都蒙恩惠。 \end{tabularx} \\ \\
[1ex]
\hline
\hline
\end{longtable}
$^{1}$好 弟兄姊妹 由上年開始.
25年1月開始.
我用已忽所書和大家分享主的訊息.
我想大家都應該留意到.
每次我講道都是講已忽所書.
來到今天.
一共今天走過17次的講道都是講已忽所書.
在過去16次.
每一次不知道大家有沒有留意.
我的主題都是用「在基督裡」作為主題的開始.
「在基督裡」這本書帶給我們什麼訊息呢?.
保羅帶我們經歷了三個層次.
我們可以一起再回顧一下.
第一至第三章.
保羅讓我們看見我們在基督裡的身份.
他明白我們在基督裡是蒙愛,蒙揀選的地位.
而第四,第五章是講到我們有一個新的生命.
保羅很想我們知道怎樣去脫去舊人.
活在光明裡面.
而第三個層次就是當我們有了新的生命之後.
五章二十一至六章九節.
提醒我們在不同的關係裡面.
怎樣去讓福音實踐出來.
透過夫婦的相處.
透過親子的相處.
透過主僕的相處.
這些很靠近的經驗.
這些很靠近的生活.
我們一起去體會神的教導.
來到今天六章十節至二十四節.
保羅加了一個很重要的提醒.
他提醒我們.
這些美好的生命和關係.
並不是在溫室或者不是在無菌的地方去培養出來.
而是在戰場裡面活出來.
我們很多時以為工作面對的壓力.
家庭面對的張力.
人際關係這些只是普通的生活問題.
但是保羅提醒我們.
其實這些是屬靈真賤的核心.

$^{41}$所以他用這段經文來總結.
「已失所整卷」的書卷.
不是要嚇我們.
而是很想裝備我們.
教我們怎樣在真實的世界裡面.
靠主站立得穩.
所以今天我給他的講題是.
在基督裡我們要做一個剛強的人.
求主去幫助我們一起先祈禱.
主我求你在裡面去幫助我們.
我們多謝你.
多謝你在創世以前已經揀選了我們.
你讓我們去體會這一份被愛.
我們多謝你.
我們都是一群很多問題的人.
但是天父你愛我們.
你賜下了獨生的兒子耶穌基督.
為我們的罪寫在十字架上.
以致我們能夠有一個新的生命.
以致我們能夠去經驗在裡面的新生活.
主我求你今天去幫助.
來到醫院書卷的最後一個段落的時候.
主我求你幫助我們無論聽的人講的人.
我們一起去學習.
我們怎樣去經驗做一個被你剛強的人.
主求你在裡面去看顧人靈.
我們多謝你.
我們將會禱告交託.
奉主明求.
阿們.
我今天採取的版本是和合本修訂版.
我都想和大家先讀一下這幾節經文.
我們一起讀.
一二三.
最後你們要靠著主.
依賴他的大能大力.
作剛強的人.
要穿戴上帝所賜的全副軍裝.
好抵擋魔鬼的詭計.
因為我們的真善.

$^{81}$並不是對抗有血有肉的人.
而是對抗那些執政的掌權的.
管轄這幽暗世界的.
以及天空靈界的惡魔.
第十節告訴我們.
剛強不是來自我們自己.
而是在主裡面被剛強.
第十節一開始提到最後.
在新約裡面最後這個字有兩個用法.
可以說一個用法是引入一個新的題目.
又可以作為一個段落的總結.
在六章十節這裡.
可以說它作為四至六章的一個總結.
指明信徒他經歷了新的生命之後.
他經歷在關係裡面怎樣去活出信仰之後.
其實我們要面對的是一場屬靈真戰.
保羅首先說.
依賴他的大能大力作剛強的人.
愚民的意思更加正確的是.
要被剛強起來.
這是一個被動式.
也是一個持續的動作.
剛強不是靠努力去撐住.
不是靠性格硬朗.
也不是靠經驗的累積.
而是不斷從神領受能力.
讓神托住我們的生命.
值得留意的是.
保羅提到在主裡這個字.
他說在主裡被剛強起來.
在主裡這句說話已不鎖書出現超過二十次.
也是全卷書神學的核心.
所以你們會明白.
為什麼我每次講道.
主題都是在主裡的一個主題.
保羅不是叫我們要向主祈求力量.
然後再靠自己去行.
而是他講到一個位置.
我們本身已經是在主裡.
位置先於行動.

$^{121}$身份先於努力.
換句話說.
我們剛強的起點.
不是我要更加努力尋求神.
而是我已經在主裡.
讓主的能力在我身上運行.
問題不是我們可不可以更加堅強.
而是我有沒有讓神使我剛強起來.
很多弟兄姊妹都會知道.
這個星期我都在一個很瘋癲的狀態.
因為在三個星期之前.
我預備這幾個星期的聚會的時候.
我才發現原來所有事情是在同一個星期發生.
我自己那一刻才恍然大悟.
原來十一號,十四號和前一天,今天.
全部在同一個星期發生.
在人的層面是很忙碌.
壓力也不少.
但是感恩的是.
在每一天的禱告,靈修.
安靜在主裡面的時候.
上帝都提醒我.
不是將焦點放在我準備的聚會裡面.
也不是我自己很努力去撐著.
而是要在上帝裡面去支取力量.
我這幾個星期都體會.
被剛強起來不是一個我努力達到的狀態.
而是一個我每一天重新選擇的位置.
留在主裡面.
讓上帝托著我的生命.
我不知道每一個聚會我做得好不好.
這個也不是我最著重的東西.
我都希望頂展會能夠享受.
但是在這幾個星期裡面.
我的體會就是與主同工那份美好.
我們繼續看經文六章十一節.
他說要穿戴上帝所賜的全副軍裝.
好抵擋魔鬼的詭計.
保羅在原文用了三個發音相似的字.
去將三個動作連結起來.

$^{161}$這三個發音相似的字.
在原文裡面發音相似的字是.
被剛強穿戴能夠抵擋.
結構很清楚.
穿戴神的軍裝就是被剛強的方式.
被剛強才能夠抵擋仇敵.
這一套神的全副軍裝.
在希臘文裡面.
是指到羅馬重裝步兵的完整裝備.
由頭到腳一件都不能少.
保羅特別用全副這個字來提醒我們.
不可以只選擇自己喜歡的裝備.
有些人只是選擇有信心.
但是不紮根真理.
有些人識神的道.
但是缺乏了福音的平安.
這些都是不完整的裝備.
會在真戰的時候露出破綻.
而一套全副的軍裝有兩個意思.
第一個意思是神給的.
來源不是來自人.
第二個意思是屬於神的.
我們穿戴了這副軍裝之後.
是代表神出戰.
神的軍裝不是外在的裝備.
而是我們內在的屬靈生命.
所以每一位信徒都可以穿得上.
而十二節更加提醒我們.
為什麼要穿戴這副軍裝.
那個真戰的本質是什麼.
經文這樣說.
因為我們的真戰並不是對抗有血有肉的人.
而是對抗那些執政的掌權的.
管轄者幽暗世界的.
以及天空靈界的惡魔.
保羅提醒我們我們並不是跟屬血氣的真戰.
魔鬼最厲害的手段不是張牙舞爪嚇我們.
而是用詭計.
一些看不到的謊言和控告.
在面對家庭的張力的時候.

$^{201}$魔鬼會控告我們的家人.
會告訴我們他沒得救了.
你付出那麼多都沒用.
又或者控告我們自己.
你都是一個失敗者.
你不配得到這份愛.
無論是在工作裡面.
在家庭裡面.
還是在教會服事裡面.
魔鬼都用同一個方法詭計.
去動搖我們對神的信任.
扭曲我們的身份.
所以我們的對手並不是我們工作上的老闆.
並不是我們的配偶.
不是我們的子女.
也不是在教會裡面.
我們不喜歡的弟兄姊妹.
如果我們看錯了戰場的話.
一定會用錯方法.
只有認清這場真戰是關乎真理與謊言的角力.
我們才明白為何有需要穿上神的全副軍裝.
所以剛強的起點.
是我們要知道自己的位置.
不是努力.
當我們在主裡面.
我們就會被剛強.
當我們被剛強.
我們就會穿起這副軍裝.
當我們穿起這副軍裝.
我們就能夠站立得穩.
面對真正的敵人.
接著看看這副軍裝的裝備.
13至17節.
描述了這副軍裝是怎樣的.
我們一起讀.
123.
所以要拿起上帝所賜的全副軍裝.
好在邪惡的日子能抵擋仇敵.
並且完成了一切還能站立得住.
所以要站穩了.

$^{241}$用真理當作帶子束腰.
用公義當作護心鏡遮胸.
又用和平的福音當作預備走路的鞋穿在腳上.
此外要拿順德的盾牌.
用來抗壓惡者一切消著的箭.
要帶上救恩的頭盔.
拿著聖靈的寶劍.
就是上帝的道.
保羅在這段經文裡.
不斷重複的就是要抵擋.
要站立得住.
要站穩.
他沒有說你要打敗你的仇敵.
他是說你要站立得穩.
這個分別是很重要.
如果我們的目標是要打贏.
我們會活在一個壓力裡.
我夠不夠強.
我會不會輸.
但如果我們的目標是站穩的話.
我們就會明白.
基督徒的德性不是我們有多強.
而是我們在壓力裡.
在試探.
在混亂裡.
仍然是不會倒下.
剛強的人不一定贏每一場戰.
不一定贏每一場仗.
但是一定能夠站穩得住.
我們繼續看保羅形容的這副軍裝.
第一件的器具是真理的腰帶.
羅馬士兵束腰是為了行動自由.
如果長袍不束起.
一出戰場就會跌倒.
真理不是靜態的知識.
而是要讓我們在戰場裡能夠走得動.
當我們對神的認識是夠紮實.
面對試探和謊言.
我們才有方向有立場走得出去.
沒有真理的根基的人.

$^{281}$不是被敵人打倒.
而是自己跌倒.
第二件是公義護心鏡.
用來抵擋控告的保護.
護心鏡保護的是人的心臟.
魔鬼最喜歡攻擊的地方.
因為一旦攻擊到你的心臟.
你就會死.
他不斷的控告.
你都算是教徒嗎?.
你服事是這樣.
神怎會縱容你?.
很多教徒站不穩.
不是因為他們不愛主.
而是一直活在這些控告裡.
保羅說的公義.
不是我們自己的義.
是耶穌基督的義.
所以當魔鬼控告我們的時候.
可能說你不配.
我們可以回應.
是的,我們不配.
但耶穌基督配.
祂的義已經遮蓋了我們.
公義的護心鏡.
不是證明我們從來沒有失敗.
而是提醒我們.
即使我們失敗過.
我們仍然是屬於主耶穌基督.
聖經裡面有一個很真實的例子.
就是彼得.
彼得是耶穌最親密的門徒之一.
是第一個認耶穌為基督的人.
但在耶穌被埋,被捕那一晚.
他三次否認不認識耶穌.
之後彼得出去痛哭.
這不是單單是他後悔.
而是他生命裡面一個徹底的崩潰.
我們猜想彼得的內心在想什麼.
他可能在想.

$^{321}$我這樣的人.
我連跟耶穌理都做不到.
我還有沒有資格做一個門徒.
但復活的耶穌沒有放棄彼得.
約翰福音21章.
耶穌三次問彼得你愛我嗎.
對應彼得三次的否認.
這是一個完全的復活和恢復這段關係.
耶穌不是問彼得你上次做得好不好.
而是問他你愛不愛我.
彼得重新的站立.
不是因為他要證明自己夠好.
而是因為耶穌主動來接納他恢復他.
弟兄姊妹這個就是公義護心鏡的真正意思.
當我們面對一些的控告.
壓倒我們.
不是叫我們證明自己.
而是我們要很清楚知道.
我們要回到主耶穌面前.
因為耶穌的義已經為我們付上代價.
第三件器具是福音的平安鞋.
是一對任何地形都站得穩的鞋.
羅馬士兵的鞋是有釘的.
設計是在任何地形都不會滑倒.
不是一對漂亮的鞋.
而是一對穩陣的鞋.
福音的平安不是環境平靜的平安.
而是無論我們雙腳踏在任何地形.
都站得穩的平安.
當失去了工作.
關係破裂.
前途不明朗.
或者有很多恐懼.
福音不是叫我們感覺良好.
而是要我們知道我們的根基是在主耶穌基督裡面.
不是在環境裡面.
然後是信心信德的盾牌.
保羅提到盾牌.
其實是在表達一種長方形的大盾牌.
大到一個士兵蹲下.

$^{361}$這個盾牌是可以完全遮住士兵的.
羅馬的士兵作戰的時候.
每位士兵站在一起.
將盾牌彼此連接.
形成了堅固的防線.
用來抵擋敵人的攻擊.
這個意象提醒我們兩件事.
在個人的層面.
信心是我們面對懷疑和恐懼的時候的保護.
在群體的層面.
我們的信心是教會防守的一部分.
當我們放下盾牌.
不單止是我們自己受傷害.
還會令到群體都出現缺口.
接著是說救恩的頭盔.
十七折上.
頭盔保護的就是我們的頭.
思想和身份的認知.
魔鬼最想攻擊的是我們的身份.
你是不是真是神的兒女?.
你的救恩是不是真的?.
救恩的頭盔就是我們救恩的確據.
我是神的兒女.
這個身份沒有人可以奪去.
當這個確據穩固.
魔鬼針對我們的身份去攻擊.
就沒有辦法入侵我們的思想.
最後一件器具就是聖靈的寶劍.
留意一件事 弟兄姊妹.
剛才所說的五樣東西.
其實全部都是防禦性的裝備.
但是寶劍是整副的軍裝裡面.
唯一拿來進攻的兵器.
保羅不是叫我們靠我們的口才去對抗仇敵.
不是用我們的方法去對抗仇敵.
而是用神的道去主動拆穿這些謊言.
回應試探.
主耶穌在曠野被試探的時候.
祂做了一個最好的示範.
魔鬼每一次出招.

$^{401}$耶穌的回應都是經上記著說.
耶穌不是用感覺 不是用道理.
而是用神的話去回應魔鬼.
弟兄姊妹.
我不想犯罪.
我不想犯罪這句話是一個感覺.
經上記著說這句話是一個事實.
感覺會動搖.
但是神的話是不會改變.
一個不認識神話語的信徒.
就等於手上拿著一把鈍的劍.
甚至沒有拿著東西的手.
或者可以防守一時.
但是面對試探的時候.
是沒有辦法主動的還擊.
所以我也想我們一起去思考一個問題.
我們現在手上那把劍.
是一把鋒利的劍.
還是一把根本拿來切普通東西都切不到的劍.
我們對神話語的認識.
是否足夠讓我們在最黑暗的時刻站穩.
我這樣問不是想大家覺得內疚覺得羞愧.
而是很想發出一個邀請.
我們真的要磨利這把劍.
最近我和牧師溝通的時候.
我們都很想去有一些聖經的課程.
目的其實就是很想磨利大家的劍.
最初我們背負著一種壓力.
起點租多了兩個單位.
大家會問你們拿來做什麼.
我們又交代不到就很麻煩.
最初我們就想在那裡搞課程.
但日間搞課程來來去去都是那幾個.
所以我們才想開一個晚上的課程.
讓大家可以參與的機會和空間都大一點.
早兩天我和牧師說有很多人報名.
很多人報名?多少人?.
我說十三個.
十三個很少.
我們多少人?弟兄姊妹.

$^{441}$起碼都八三人.
我們一起去努力.
讓這把劍可以磨利一點.
以致我們在面對不同的艱難和挑戰裡面.
我們知道怎樣去回應.
知道怎樣去還擊.
有關詳情你們可以看看稍後的報告.
總結一個剛強的人要被全面的軍裝去裝備好.
六件軍裝缺一不可.
真理不會讓我們自己跌倒.
公義讓我們不會活在控告裡面.
福音的平安讓我們不會被環境去動搖.
信心讓我們不會被謊言去擊倒.
救恩的確據讓我們不要失去身份.
神的道讓我們能夠主動回應不同的挑戰.
不同的問題.
剛強不是感覺強.
而是在生命裡面被全面的裝備.
求主去幫我們.
第三部分是說到一個剛強的生活.
很需要的一個元素.
我們一起看十八至二十節.
我們一起讀.
123.
要靠著聖靈隨時多方禱告祈求.
並要為此警醒不倦.
為眾聖徒祈求.
也要為我祈求.
讓我有口才能放膽開口.
講明福音的奧秘.
我為這福音的奧秘.
作了戴鐵鏈的使者.
讓我照著當盡的本份放膽宣講.
十八節說.
要靠著聖靈隨時多方禱告祈求.
祈禱是貫穿整個真正的生命節奏.
十八節的祈禱這個字.
是在接駁前面軍裝的段落.
在文法上.
禱告這個字是指穿軍裝的方式.

$^{481}$或持續的狀態.
不是一個獨立的第七個步驟.
換句話說.
保羅的意思是.
我們穿這副軍裝的整個過程裡.
都是在祈禱當中進行.
沒有禱告.
軍裝是死的.
有禱告.
這副軍裝才是活的.
就好像我們的手機一樣.
我們會用一些好的手機.
有沒有人是用iPhone17的.
我們的手機再好的功能.
更好如果沒有訊息.
收不到.
Disconnect.
它只是一塊很貴的磚頭.
同樣我們認識很多的真理.
我們有信心.
我們有救恩的確據.
但如果我們沒有祈禱.
我們和能力的源頭斷了線.
面對的真戰.
就會覺得有力沒有用.
保羅用四個詞語去形容這種連線.
靠著聖靈.
不是靠自己去想.
而是博通和聖靈的頻道.
隨時不是開會才祈禱.
不是吃飯才祈禱.
不是翻崇拜才祈禱.
而是上上去禱告.
多方.
不只是用一種形式.
而是可以用任何的形式.
什麼事都可以和神去傾.
祈求將真實.
我們自己的軟弱.
講給上帝知.

$^{521}$今日我們的祈禱生活又是怎樣呢.
求主去幫我們.
我們的祈禱.
保羅更加提醒我們.
這個祈禱是站立得穩的關鍵.
是一個警醒不倦的禱告.
屬靈真戰最大的危險.
不是我們太弱.
而是我們不覺得自己在打仗.
當人失去了警醒.
就會鬆懈妥協.
慢慢就會跌下去.
禱告最大的挑戰.
不是我們不懂得祈禱.
而是站一會就會停.
保羅說要不倦.
不要放棄.
因為屬靈真戰不是一日完成.
禱告不可以只是一時的熱心.
在十九世紀英國.
有一個人叫做喬治·慕勒.
慕勒年輕的時候.
他是一個浪子.
他會偷竊,說謊.
是一個放蕩不羈的人.
但當他遇見耶穌之後.
他決定要用一生去證明一件事.
神是信實的.
祈禱是真實的.
他在英國Bristol開了一間孤兒院.
收容了成千上萬無家可歸的孤兒.
但他立了一個原則.
從不公開籌款.
從來不告訴別人自己有需要.
只是將一切的需要帶到神的面前.
他相信如果神呼召他做這件事.
一定會公認.
幾十年來沒有一個小朋友是餓著睡覺.
有一次早餐時間.
廚房沒有食物了.

$^{561}$慕勒叫小朋友坐下.
他帶小朋友做了一個感恩的禱告.
但廚房沒有食物.
為什麼還要感恩?.
他相信主的公認.
在祈禱完結後.
有人敲門.
送了足夠所有小朋友吃飽的麵包和牛奶.
他將每一個禱告都記錄起來.
超過了五萬個禱告的經驗.
他的生命令我最感動的.
不是這些即時的答案.
在禱告裡面得到這些即時答案.
而是他有五個朋友.
他為他們的得救去禱告.
第一位朋友在他祈禱之後幾個月就信了主.
第二第三位在他祈禱之後幾年就信了主.
第四位在他祈禱之後幾十年就信了主.
第五位朋友為他祈禱了六十年.
直到慕勒臨終.
這個人都沒有信耶穌.
但他不斷地為他禱告.
這位朋友在他喪禮之後.
決了自信主.
弟兄姊妹這個經歷提醒我們.
禱告不倦不是因為我們看見答案.
而是因為我們認識那位垂聽禱告的神.
他不是禱告到見到結果就停.
他是禱告到生命的終結.
這個朋友神在他生命完結之後.
成就了上帝的應許.
保羅提醒我們要警醒不倦.
為眾聖徒祈求.
慕勒就是這樣用他的一生.
示範了這句說話的意思.
真正剛強的人不是祈禱很多.
而是祈禱到底.
保羅更加是十八節加上為眾聖徒祈求.
這句說話不是一個附帶的補充.
而是屬靈真戰裡面一個很重要的提醒.

$^{601}$前面說到信德的盾牌.
我提到每一個士兵拿著盾牌排成一直線.
就可以保護著大家.
作戰的時候士兵彼此成為一個堅固的防線.
這種陣型厲害的地方.
不是因為某一個士兵特別強.
而是每一個人都沒有單獨暴露在敵人面前.
只要有一個缺口出現.
整個防線就受到威脅.
而敵人往往就是專門找這些破口.
教會的屬靈真戰都一樣.
神從來沒有叫我們信徒孤軍作戰.
而是很想我們彼此扶持.
彼此守望.
彼此代求.
當我們缺少了為弟兄姊妹祈禱.
就好像這個防線出現了一個缺口.
讓身邊的人獨自面對攻擊.
但當我們彼此代求的時候.
就好像將這些盾牌彼此連結起來.
讓整個群體都得到保護和堅固.
但很多時候我們的祈禱都停留在自己的需要裡面.
我們都可以嘗試花時間了解我們身邊的弟兄姊妹.
有什麼需要祈禱.
去為身邊的弟兄姊妹的生命作出守望.
去一起建立堅固的保障.
保羅亦都做了一個很實際的榜樣.
十九至二十節他提出.
也要為我祈求.
讓我有口才.
能夠放膽講明福音的奧秘.
保羅是使徒.
是宣教士.
是寫聖經的人.
但他都會說請你為我祈禱.
弟兄姊妹.
我們要為我們身邊的弟兄姊妹祈禱.
要為教會的領袖祈禱.
要為同宮去禱告.
沒有了大家的祈禱.

$^{641}$其實我們很容易有缺口.
求主去幫我們.
最後我們去到二十一節至二十四節.
剛強的根源.
我們是活在恩典的群體裡面.
保羅講完屬靈真戰.
講完全副軍裝.
講完禱告之後.
他沒有停在戰場的畫面.
而是將信徒帶到群體.
關係,愛和恩典.
他在二十一節提到.
有個人叫推基古.
一個很真實的人.
他特別去探訪信徒.
安慰信徒的心.
他不只是給我們有軍裝.
神也給我們有彼此.
屬靈真戰不是一個人去打.
軍裝不是一個人去穿.
禱告也不是一個人去求.
神用人去安慰我們身邊的人.
用群體去托住群體.
你身邊有沒有這樣的推基古呢?.
你又有沒有成為別人的推基古呢?.
保羅最後祝福信徒.
平安,愛心,信心.
正正呼應前面的軍裝.
平安對應福音的孩.
信心對應那個盾牌.
而愛心是貫穿整副軍裝的核心.
軍裝不是冰冷的盔甲.
而是活在愛,信心,平安裡面的生命.
而整卷書的收結是一個字.
恩典.
我們能夠站立得穩是恩典.
我們能夠祈禱是恩典.
我們能夠穿上這副軍裝也是恩典.
我們能夠被剛強,能夠愛主也是恩典.
我們從來都不是靠自己去打贏這場仗.

$^{681}$是上帝自己托著我們每一步.
惟有耶穌基督在十字架上為我們死.
惟有耶穌基督有復活的能力.
惟有耶穌基督陪伴我們走過我們每一步.
\newpage



\section{創世記 26:1-11}
\label{sec:032170}
\textbf{應許的傳承——神在危機中的保守}~~~~甘雪芬傳道
\newline
\newline
連結: \href{http://www.church.com.hk/acms/content.asp?site=cdc\&op=show\&type=product\&code=032170\&layout=sermon}{\texttt{032170}} ~~~~ 語音日期: 2026-05-16
\newline
\newline
\hyperref[sec:032169]{\small{< < < PREV SERMON < < <}}
~
\hyperref[sec:index_chronic]{\small{[返順時目]}}
~
\hyperref[sec:032171]{\small{> > > NEXT SERMON > > >}}
\newline
\newline
創世記 26:1-11
\newline
\begin{longtable}{cl}
\hline
\hline
章節 & 經文 (和合本修訂版)\\
\hline
26:1 & \begin{tabularx}{0.7\textwidth}{X} 那地有了饑荒，不是亞伯拉罕的時候曾有過的那次饑荒，以撒就到基拉耳，非利士人的王亞比米勒那裡去。 \end{tabularx} \\ \\ \relax
26:2 & \begin{tabularx}{0.7\textwidth}{X} 耶和華向以撒顯現，說：「你不要下埃及去，要住在我所指示你的地。 \end{tabularx} \\ \\ \relax
26:3 & \begin{tabularx}{0.7\textwidth}{X} 你要寄居在這地，我必與你同在，賜福給你，因為我要將這一切的地都賜給你和你的後裔。我必堅定我向你父親亞伯拉罕所起的誓。 \end{tabularx} \\ \\ \relax
26:4 & \begin{tabularx}{0.7\textwidth}{X} 我要使你的後裔增多，好像天上的星，又要將這一切的地賜給你的後裔，並且地上的萬國都必因你的後裔得福， \end{tabularx} \\ \\ \relax
26:5 & \begin{tabularx}{0.7\textwidth}{X} 因為亞伯拉罕聽從我的話，遵守我的吩咐、誡令、律例和教導。」 \end{tabularx} \\ \\ \relax
26:6 & \begin{tabularx}{0.7\textwidth}{X} 於是，以撒住在基拉耳。 \end{tabularx} \\ \\ \relax
26:7 & \begin{tabularx}{0.7\textwidth}{X} 那地方的人問起他的妻子，他就說：「她是我的妹妹。」原來他害怕說「我的妻子」。他想：「或許這地方的人會因利百加殺我，因為她容貌美麗。」 \end{tabularx} \\ \\ \relax
26:8 & \begin{tabularx}{0.7\textwidth}{X} 他在那裡住了一段很長的日子。有一天，非利士人的王亞比米勒從窗戶往外觀看，看哪，以撒在撫愛他的妻子利百加。 \end{tabularx} \\ \\ \relax
26:9 & \begin{tabularx}{0.7\textwidth}{X} 亞比米勒召以撒來，說：「看哪，她實在是你的妻子，你怎麼說『她是我的妹妹』呢？」以撒對他說：「因為我想，恐怕我會因她而死。」 \end{tabularx} \\ \\ \relax
26:10 & \begin{tabularx}{0.7\textwidth}{X} 亞比米勒說：「你向我們做的是甚麼事呢？百姓中有一個人幾乎要和你的妻子同寢，你就把我們陷在罪中了。」 \end{tabularx} \\ \\ \relax
26:11 & \begin{tabularx}{0.7\textwidth}{X} 於是亞比米勒命令眾百姓說：「凡侵犯這個人，或他妻子的，必要把他處死。」 \end{tabularx} \\ \\
[1ex]
\hline
\hline
\end{longtable}
$^{1}$各位晚安.
願上主的恩典常與眾人同在.
我們同心的禱告.
親愛的主.
能夠一同去敬拜你.
是我們何等不配的福氣.
願為帶著一份渴慕的心.
去聆聽上主你的教導.
求聖靈親自在我們當中的工作.
以至無論講的.
聽的一同去蒙福.
聽我們的禱告.
奉靠我救主耶穌基督.
寶貴的聖名去求.
阿們.
大家可能都會聽過.
有些叫百年老譜.
因為不想一些很寶貴的東西.
去失傳.
一代接一代傳承下去.
窩福踏入了五十週年.
大家有沒有想過.
我們都是在一場.
年年不斷的.
屬靈接力賽那裡.
我們當中無論身處什麼崗位.
什麼的角色.
我們有意無意.
都是在傳承著一些東西.
我們小子愛言.
待會吃什麼都好.
但是有些神學立場我們要去持守.
每間教會是完美.
總有強有弱.
在接棒交棒的過程裡.
大大小小的決策裡.
我們要盡上人的責任.
但我們不得不承認.
現實裡邊.
我們有軟弱.

$^{41}$我們有限制.
我們需要依靠上帝的恩典.
帶領我們去跨過.
願意藉著今天的經文.
《創世記》.
二十六章一至十一節.
我們一起去思想.
我們去學習.
在《創世記》裡.
假如我提到以薩這個人物.
大家想起什麼敘事呢?.
亞伯拉罕獻以薩.
主角是誰?.
亞伯拉罕多一點.
以薩被雅各欺騙長子的名分.
主角又是誰?.
雅各多一點.
確實在《創世記》裡.
這個族長的敘事裡.
以薩不像亞伯拉罕,雅各,約瑟.
有戲劇性的經歷.
但是以薩的角色.
仍然有很獨特的位置.
神安放給每一個生命有不同的份.
以薩是神給亞伯拉罕.
憑應許所生的兒子.
又別於黯血器所生的以薩瑪利.
當由亞伯拉罕領受萬國恩利的後裔德福.
這個應許.
到雅各成為以色列的始祖.
以薩繼承.
傳遞了這個盟約.
確保了信仰血脈的延續.
今天看的二十六章.
在《創世記》裡.
一章比較清晰.
以以薩為主角的敘事.
我們先來概念一下二十六章的主題.
剛才大家聽過二十六章的上半段.
有饑荒.

$^{81}$神重新許的應許.
以薩在基拉爾認妻作妹妹.
內容是否似曾相識?.
在哪裡聽過?.
不難發現.
以薩有著父親亞伯拉罕事跡的影子.
這樣的鋪排.
有什麼意義?.
與此同時.
如果我們看看二十六章的上一章是二十五章.
和下一章的二十七章是說什麼的呢?.
二十五章是寄述.
雅各.
如何藉著餅和豆湯和耳掃.
奪取長子的名分.
而下一章的二十七章是寄述.
雅各.
去扮耳掃.
欺騙老眼昏花的以薩的祝福.
兩章一起開成說就好了.
為什麼中間要插入二十六章.
寄述神將納入的福分.
從亞伯拉罕傳到以薩呢?.
這樣的鋪排有什麼意義呢?.
今天我們先來看二十六章的上半部一至十一節.
下次我們會看下半部.
一至十一節可以分為兩個部分.
寄述到以薩面對兩個危機.
和神在當中如何保守.
我們先來看一至六節.
這是一個外在的危機.
看當時的事件.
讀者提到有饑荒.
強調不要和亞伯拉罕日子.
那一次的饑荒.
混為一談.
讀者刻意.
一開始就將以薩和亞伯拉罕.
結連在一起.
他們有著很近似的境遇.

$^{121}$以薩會如何回應這個饑荒?.
會否像父親一樣.
我們重溫一下.
昔日亞伯拉罕遇上饑荒時.
他如何做?.
上文十二章.
亞伯拉罕去了埃及.
雖然充足.
饑荒時去到一些比較理想的地理環境.
很合情理.
那麼神又如何吩咐以薩呢?.
神吩咐以薩.
不要學校的父親去到埃及.
要住在他指示的地方.
就是基拉爾.
看回二十章的上文.
原來亞伯拉罕曾經在基拉爾暫居過.
如果看地圖.
基拉爾當時是屬於非利士領土的一部分.
離開饑荒的地方其實不是很遠.
相對埃及水源充足.
又遠離饑荒的地區.
按地理環境.
如果去埃及當然安全很多.
不可以去埃及.
這個吩咐對於以薩來說.
是一個信服神的考驗.
這個吩咐要帶著一些應許.
在三至四節.
神應許說.
如果你願意寄居在這個地.
我必與你同在,賜福給你.
因為我要將一切的地都賜給你和你的後裔.
我會堅定地向你父親亞伯拉罕所起的誓.
我要使你的後裔增多好些天上的星.
我要將一切的地賜給你的後裔.
並且地上的萬國因你的後裔得福.
如果我們看看《尚文》十二章.
大家還記得嗎?.
其實那時候.

$^{161}$神也曾經吩咐過亞伯拉罕.
要怎樣呢?.
離開本地本族父家.
到上帝指示的地方去.
當時上帝也有一樣.
給很多賜福的應許給亞伯拉罕.
以薩領受的呼召.
和當年神給亞伯拉罕的呼召.
其實很相似.
神給以薩的應許.
比起當初給亞伯拉罕的應許.
有更具體的描述.
以薩為何會蒙福呢?.
這是一個很有趣的原因.
我們值得花點時間去看看.
《尚文》二十五章說.
亞伯拉罕死了以後.
神賜福給他的兒子以薩.
亞伯拉罕的離世.
需要有後裔持續的傳承.
讓神的約可以一路應驗.
以薩就被揀選.
以薩要做些甚麼?.
好像沒甚麼要做.
他要履行的吩咐就是寄居在這個地方.
神賜給以薩的福份有些甚麼呢?.
有些是現在的.
我必與你同在賜福給你.
有些是將來的.
是講給他的後裔.
但為何上帝要賜福給以薩呢?.
他說:我必堅定.
我向你父親亞伯拉罕所起的誓.
神銳意要履行他過去的承諾.
也因為亞伯拉罕聽從我的話.
遵守我的吩咐,誡令,律例和教導.
是因為基於亞伯拉罕對父神的忠誠.
我們看見一個忠心的亞伯拉罕.
帶來蒙福的以薩.
神重申他的盟約.

$^{201}$將福份從亞伯拉罕傳承到他兒子以薩身上.
神要將亞伯拉罕祝福萬民.
如果我們看看第四節,第五節.
上帝給以薩賜福的內容.
我們發現由一個人發展到一個大國家的過程.
需要三個很重要的元素.
第一個要口裔增多.
口裔多才可以建立國民.
需要有很多的地圖.
這是要建立國土.
既然建立的是上帝國.
這群子民一定要去遵行上帝的說話.
這就是國法.
國民,國土,國法.
在舊約的五經裡構成了一個上帝國的雛形.
到了新約.
這個元素就在耶穌基督裡被重新定義.
舊約裡國民是以色列民.
國土是迦南地.
國法是很多的律例規條.
到了新約的時候.
國民是我們這一群信靠耶穌基督的人.
國土甚至國林在我們心裡終末新天新地完全的實現.
國法.
耶穌基督.
祂重新去承傳了律法真正的意義.
愛神愛人.
由聖靈內在我們心裡去引導.
我們看見.
神基於祂的信實和恩典.
承諾一定會成就祂的應許.
在救贖歷史裡.
神父召亞伯拉罕.
應許會成為大國.
萬族因此可以得福.
同樣.
神應許大位的後裔永遠作王.
最終致向耶穌基督成就的救贖.
一切都是白白的恩典.
神的應許是白白的恩典.

$^{241}$但是.
人要透過信心去信靠信服神.
才能領取神應許的福分.
人怎樣去回應.
決定他自己個人和他所屬的世代.
可否經歷其中的福分.
丙姐妹.
在一場連綿不斷的.
接棒交棒的過程裡.
盼望我們不只是單單.
我自己蒙福就很滿足.
盼望我們接棒接得好一點.
以致我們交得好一點.
剛才說過.
神吩咐以撒要住在他所指示的地方.
基拉爾.
以撒聽了他的說話.
就著這一點.
以撒仿效他的父親.
做一個信服上帝吩咐人.
最後.
饑荒的問題有沒有得到解決.
我們先睹為快.
夏文是怎樣.
以撒在那裡耕種.
那一年有百倍的收成.
以撒將會發現.
即使身處很危難的饑荒.
上帝總能夠保守屬於他的人.
做個小結.
神因為亞伯拉罕的信心.
忠誠賜福給下一代的以撒.
神的英雄一路在傳承.
一直在推展.
接著我們就會問.
以撒的生命又是否能夠成為.
上帝可以使用的合用器皿.
將神的祝福一路傳承下去.
我們繼續看下去.
最初的問題是什麼.

$^{281}$危機的問題.
饑荒的問題.
神解決了這個饑荒.
是否沒有其他的挑戰呢.
七至十一節.
以撒面對另一個危機.
這是來自他自己內心.
來自他自己內心的驚恐.
我們看回第七節.
當時他在基拉爾.
其實基拉爾人很平常.
問及他的妻子一句.
但以撒感到很大的威脅.
認妻子為妹妹.
他怕因為妻子太漂亮.
人們會殺他.
神吩咐以撒在基拉爾.
跟他說過什麼.
「我必與你同在」.
以撒踏實地經歷過.
在饑荒裡有保守.
但當他身處在.
非利士領土的時候.
他就害怕.
上帝所揀選的僕人.
有小信的時候.
第八節.
說他在那裡住了一段很長的時間.
沒有事情發生.
更加印證了以撒.
其實是過於擔心.
直至到以撒和利伯家的.
親暱行為被發現.
就東窗事發了.
第十節提到.
阿彼美勒和以撒這樣說.
百姓裡有一個人.
幾乎要和你的妻子同寢.
你就將我.
陷在罪裡面.

$^{321}$這其實是很諷刺的.
反而你看到.
阿彼美勒比以撒更加著緊.
要捍衛.
不要犯罪.
更加著緊要警告.
以撒我們不可以犯罪.
在十一節.
看到他那種.
預旨.
反映出他為人很公義.
他不想百姓犯罪.
凡是.
侵犯這個人.
或他的妻子都要將他處死.
神曾經應許以撒說.
萬國因他的後裔.
會得福.
但當以撒身處在.
地圖基拉爾的時候.
他的行為幾乎.
讓這個國家遭殃.
所以其實是一個很諷刺的.
做法反而讓一個.
以彼美勒的王.
教回以撒.
應妻為妹妹.
這個事件你們聽過多少次.
記不記得.
起碼有三次.
第三次.
為什麼有三次那麼多.
我們很快簡單概念一下.
第一次是亞伯蘭機方的時候.
去到埃及.
第二次是他寄居在基拉爾.
第三次是以撒在基拉爾.
發生了什麼事呢.
第一次.
撒拉已經被帶到.

$^{361}$法老的宮殿.
甚至法老已經給他聘禮.
其實很危險.
聘禮都給了.
第二次撒拉都被人帶走了.
而今次第三次.
他們住了一段很長的時間.
沒有事發生.
特別想提一提.
第20章26章那兩個.
阿彼米勒是不一樣的.
阿彼米勒好像當時的一些名號.
好像一些.
法老那些.
上帝如何介入.
在事件裡面呢.
第一次是降災.
去擊打法老和全家.
第二次在夢裡面.
去啟示.
第三次是阿彼米勒王.
他自己無意中見到以撒.
和李伯家一些親戚的行徑.
結果是怎樣.
第一次.
法老吩咐亞伯蘭和他的妻子.
將他所有的東西都被他送走了.
第二次.
所有的地都要在他前面.
你認為好的可以在那裡住.
第三次命令.
凡侵犯這個人這個妻子的.
要將他處死.
三次的事件.
我看到危機.
一次比一次少.
我看見神在那裡保守住.
他被揀選的後裔.
那份純潔.
三次事件的.

$^{401}$平衡並列.
主要不是講一些道德的教導.
我們看見.
人一次又一次的失敗.
大神一次又一次的保守.
我們做個小結.
七至十一節.
神越過人的失敗.
去保守他的應許.
可以得以成就.
危機來自不同的.
源頭.
有時是客觀的環境.
有時來自人的威脅.
但最重要的是.
我們自己的問題.
其實很多危機的事.
不到我們去想.
有時很有趣.
如果繼續看26章的下半段.
很開心.
上帝在饑荒裡.
給他100倍的修成.
做了大富翁.
接著非利士人.
妒忌他.
我們有時會覺得.
很有趣.
祝福我,為何不祝福到底.
看下文就發覺.
神容讓人.
在面對各樣的挑戰.
各樣的處境裡.
考驗,操練.
我們對神的信心.
和信服.
蒙神揀選的人.
我們有時很剛強.
但總有時會很軟弱.
今日經文讓你看見.

$^{441}$以撒他面對兩次的危機.
第一次.
是一個很外在的客觀環境.
因為饑荒.
所以他要求問神.
他願意聽從神的吩咐.
以致可以蒙恩.
上帝在當中的保守.
第二次.
是一個內在的危機.
他很害怕.
以致他欺騙.
妻子是妹妹.
一個很大的危機裡.
上帝在當中繼續保守.
以致難了惡事的發生.
以撒對特出的品質.
其實是很信服的.
留意他被父親.
獻上繁祭的時候.
他沒有反抗.
不過與此同時.
他是一個很膽小的人.
如果下次有機會我們看.
在非利士地.
當井被霸佔的時候.
以撒一再去引揚.
我們人有缺點.
神所揀選的器皿.
我們都是脆弱的.
但神願意與我們同在.
越過我們的失敗.
讓祂的應許得以成就.
今天學習的經文.
讓我們去看見.
亞伯拉罕的.
信心和忠誠.
以致上帝願意賜福.
給下一代的以撒.
而以撒也就在.

$^{481}$亞伯拉罕的腳中.
我們很遲會問下一個問題.
以撒下一代的生命.
以素和雅各又如何呢?.
是否可以繼續成為.
神合用的器皿.
將福分一路傳承下去呢?.
所以.
二十六章就夾在.
二十五章.
二十七章的中間.
以素的問題是.
他繼承了長子的名分.
他尚且興翰.
對聖約的繼承權.
他看來應該不太重視.
雅各的問題是.
他要不擇手段去奪取.
長子的名分.
他的品格究竟可不可以.
去重視他.
順從上帝的吩咐.
能夠符合聖約的期望呢?.
在各樣的危機裡面.
如果不是有.
上帝當中保守.
應許的話.
上帝的聖約如何可以傳承下去呢?.
今日的經文.
讓我們看見.
在與神的交往裡.
上帝很希望.
我們能與祂建立一個更深的關係.
神不用靠人.
才能成就祂的旨意.
但神很希望人.
可以與祂同工.
經歷神的恩典.
在當中建立一個更深厚的關係.
神的計劃.

$^{521}$當然不會使人難阻.
但人如何回應呢?.
就會視乎.
人能否蒙福呢?.
就視乎人如何回應.
這是一個.
感動的見證.
認識兩位.
小弟兄.
他們在一個.
公園.
有一個冬菇亭.
做一個小吃店.
生意不太理想.
他們要自己.
做一些兼職.
以幫補平時的生活.
那裡的品流.
很複雜.
除了平時.
購物之外.
會被人罵.
甚至被人打.
當我約滿的時候.
我以為.
我會儘快離場.
他們竟然繼續續約.
為什麼呢?.
原來那裡.
有些牧師.
教會.
有些崇拜.
有些福音工作.
他們的小吃庭.
成為了後盾.
例如水電.
雖然那裡.
有些品流複雜的人.
會被人罵.
但有很多南亞裔的小孩子.

$^{561}$很喜歡他們.
跟他們一起玩.
小弟兄這樣說.
學歷不高 錢也沒有.
還有什麼可以輸得起.
沒有什麼可以輸的.
為什麼我不繼續.
我使命未完.
為什麼我不繼續.
我覺得.
很感動.
一個開始.
微不足道 很平凡的生命.
可以成為上帝使用的器皿.
神明白我們的限制.
我們有軟弱.
我們所做的.
很多時候都微不足道.
但神體寵我們.
在當中忠誠地.
我們學習.
愛上帝 靠上帝.
為上帝而活.
烏福發入了.
五十週年.
在這個.
連綿不斷接棒交棒.
屬靈的接力賽裡.
大家找到你的角色嗎.
在面對什麼危機或挑戰.
我們很容易跟你.
來一大堆挑戰.
我們確實有很多外在問題要面對.
我們堂主任仍然空缺.
我們堂子的安排.
我們士工的策劃.
我們童工跨代的配搭.
我們有些教務童工.
身心靈的健康需要.
等等.

$^{601}$不過更加重要的.
既然教會的使命.
是要去擴展上帝的國度.
我們要去審視.
究竟我們體重的是什麼.
我們問我們回來教會.
真正感到雀躍的是什麼.
我們最大的危機.
會不會其實是.
我們回來教會.
做一些不是上帝要我們做的事.
我們在傳承什麼.
我們的接棒交棒.
接什麼交什麼.
這個反而是我們需要.
去深思說.
我們緊張的事.
剛才構成上帝國三個元素.
國民.
國土國法的框架去思想.
國民.
什麼是上帝國的國民.
這班去信靠耶穌.
跟從耶穌的門徒.
我苦心自問.
我忍受在很多事工裡.
有時自己很忙的時候.
希望快快把事情.
搞定,交功課我就很安樂.
在所難免.
但我們其實不是.
搞宗教活動.
我們是要建立一班.
信靠跟從基督的門徒.
我們的國土.
我們的社區,我們的職場.
我們的家庭.
都可以是我們的和場.
國法.
這是上帝國.

$^{641}$上帝的國竟然去聽上帝的說話.
就是最重要的國法.
教導,聆聽.
信服神的話語.
這是我們責無旁貸.
面對個人的挑戰.
打著人的限制.
我們確實需要依靠上帝的恩典.
帶領我們去走過,跨過.
我們期盼.
我們不會輕率抬舵.
但我們難免.
有失敗,軟弱的時候.
盼望在當中的時候.
我們不會過於震驚.
我們不會過於失望,而被拖垮.
我們信得過,只要我們捉緊上帝.
仍然有信心,有勇氣.
繼續走下去.
神的計劃,當然不會受到人的攔阻.
但人在當中.
怎樣去回應呢?.
就決定他能不能夠蒙福了.
我們在高峰五十週年.
這個年年不斷接棒.
交棒的屬靈接力賽中.
我們在接一個什麼棒?.
我們在交一個什麼棒?.
我們能否成為上帝一個合用的器皿.
以致讓上帝可以將這份福分.
一直傳承下去.
我們一同同心的受望.
一起的禱告.
我們謙卑來到您的庚前.
主讓我們何等的蒙恩.
我們能夠領受上帝您救贖的恩情.
他們擺上的生命.
將這個美好的福音.
向我們傳上.
當我們蒙福之後,我們盼望獨善其身.

$^{681}$我們更加去將這份福氣.
傳給我們身邊的人.
主,我求你讓我們.
有那份心智.
更加去映清楚.
我們作為天國子民的身份.
以致我們能夠成為.
上帝合用的器皿.
可以一路將這份福氣.
傳承,願主你祝福.
願主你大靈.
成我們的軟弱,我們有很多的虧欠.
讓上帝親切在心中的.
無助,成為我們的鼓勵.
聽我們的禱告.
奉靠我救主耶穌基督.
寶貴性命去求.
\newpage



\section{使徒行傳 1:1-26}
\label{sec:032171}
\textbf{使命接力:靠聖靈而行的群體}~~~~羅翠雲傳道
\newline
\newline
連結: \href{http://www.church.com.hk/acms/content.asp?site=cdc\&op=show\&type=product\&code=032171\&layout=sermon}{\texttt{032171}} ~~~~ 語音日期: 2026-05-17
\newline
\newline
\hyperref[sec:032170]{\small{< < < PREV SERMON < < <}}
~
\hyperref[sec:index_chronic]{\small{[返順時目]}}
~
\hyperref[sec:032172]{\small{> > > NEXT SERMON > > >}}
\newline
\newline
使徒行傳 1:1-26
\newline
\begin{longtable}{cl}
\hline
\hline
章節 & 經文 (和合本修訂版)\\
\hline
1:1 & \begin{tabularx}{0.7\textwidth}{X} 提阿非羅啊，我在第一本書中已論到耶穌從開頭所做和所教導的一切事， \end{tabularx} \\ \\ \relax
1:2 & \begin{tabularx}{0.7\textwidth}{X} 直到他藉著聖靈吩咐所揀選的使徒後，被接上升的日子為止。 \end{tabularx} \\ \\ \relax
1:3 & \begin{tabularx}{0.7\textwidth}{X} 他受害以後，用許多確據向使徒顯明自己是活著的，在四十天之中向他們顯現，並講說神國的事。 \end{tabularx} \\ \\ \relax
1:4 & \begin{tabularx}{0.7\textwidth}{X} 耶穌和他們聚集的時候，囑咐他們說：「不要離開耶路撒冷，但要等候父的應許，就是你們聽見我說過的。 \end{tabularx} \\ \\ \relax
1:5 & \begin{tabularx}{0.7\textwidth}{X} 約翰是用水施洗，但過了不多幾天，你們要在聖靈裡受洗。」 \end{tabularx} \\ \\ \relax
1:6 & \begin{tabularx}{0.7\textwidth}{X} 他們聚集的時候，問耶穌：「主啊，你就要在這時候復興以色列國嗎？」 \end{tabularx} \\ \\ \relax
1:7 & \begin{tabularx}{0.7\textwidth}{X} 耶穌對他們說：「父憑著自己的權柄所定的時候和日期，不是你們可以知道的。 \end{tabularx} \\ \\ \relax
1:8 & \begin{tabularx}{0.7\textwidth}{X} 但聖靈降臨在你們身上，你們就必得著能力，並要在耶路撒冷、猶太全地和撒瑪利亞，直到地極，作我的見證。」 \end{tabularx} \\ \\ \relax
1:9 & \begin{tabularx}{0.7\textwidth}{X} 說了這些話，他們正看的時候，他被接上升，有一朵雲彩從他們眼前把他接去。 \end{tabularx} \\ \\ \relax
1:10 & \begin{tabularx}{0.7\textwidth}{X} 他升上去的時候，他們定睛望天，看哪，有兩個人身穿白衣站在他們旁邊， \end{tabularx} \\ \\ \relax
1:11 & \begin{tabularx}{0.7\textwidth}{X} 說：「加利利人哪，你們為甚麼站著望天呢？這離開你們被接升天的耶穌，你們見他怎樣升上天去，他也要怎樣來臨。」 \end{tabularx} \\ \\ \relax
1:12 & \begin{tabularx}{0.7\textwidth}{X} 有一座山，名叫橄欖山，離耶路撒冷不遠，有安息日可行走的路程。那時，門徒從那裡回耶路撒冷去， \end{tabularx} \\ \\ \relax
1:13 & \begin{tabularx}{0.7\textwidth}{X} 他們一進城，就上了所住的樓房；在那裡有彼得、約翰、雅各、安得烈、腓力、多馬、巴多羅買、馬太、亞勒腓的兒子雅各、激進黨的西門，和雅各的兒子猶大。 \end{tabularx} \\ \\ \relax
1:14 & \begin{tabularx}{0.7\textwidth}{X} 這些人和幾個婦人，包括耶穌的母親馬利亞，和耶穌的兄弟，都同心合意地恆切禱告。 \end{tabularx} \\ \\ \relax
1:15 & \begin{tabularx}{0.7\textwidth}{X} 那時，有許多人聚會，約有一百二十名，彼得在弟兄中間站起來，說： \end{tabularx} \\ \\ \relax
1:16 & \begin{tabularx}{0.7\textwidth}{X} 「諸位弟兄，聖經的話必須應驗。聖經中，聖靈曾藉大衛的口預先說到那領人來拿耶穌的猶大； \end{tabularx} \\ \\ \relax
1:17 & \begin{tabularx}{0.7\textwidth}{X} 他本來算是我們中的一個，並且得了這一份使徒的職任。 \end{tabularx} \\ \\ \relax
1:18 & \begin{tabularx}{0.7\textwidth}{X} 這人用他不義的代價買了一塊田，以後身子仆倒，肚腹崩裂，腸子都流出來。 \end{tabularx} \\ \\ \relax
1:19 & \begin{tabularx}{0.7\textwidth}{X} 住在耶路撒冷的人都知道這事，所以按著他們當地的話把那塊田叫亞革大馬，就是「血田」的意思。 \end{tabularx} \\ \\ \relax
1:20 & \begin{tabularx}{0.7\textwidth}{X} 因為《詩篇》上寫著：「願他的住處變為廢墟，無人在內居住。」又說：「願別人得他的職分。」 \end{tabularx} \\ \\ \relax
1:21 & \begin{tabularx}{0.7\textwidth}{X} 所以，主耶穌在我們中間出入的整段時間， \end{tabularx} \\ \\ \relax
1:22 & \begin{tabularx}{0.7\textwidth}{X} 就是從約翰施洗起，直到主離開我們被接上升的日子為止，必須從那常與我們一起的人中，立一位與我們同作耶穌復活的見證。」 \end{tabularx} \\ \\ \relax
1:23 & \begin{tabularx}{0.7\textwidth}{X} 於是他們推舉兩個人，就是那叫巴撒巴，又稱為猶士都的約瑟，和馬提亞。 \end{tabularx} \\ \\ \relax
1:24 & \begin{tabularx}{0.7\textwidth}{X} 眾人禱告說：「主啊，你知道萬人的心，求你從這兩個人中指明你所揀選的是哪一位， \end{tabularx} \\ \\ \relax
1:25 & \begin{tabularx}{0.7\textwidth}{X} 去得這使徒的職任；這職位猶大已經丟棄，往自己的地方去了。」 \end{tabularx} \\ \\ \relax
1:26 & \begin{tabularx}{0.7\textwidth}{X} 於是眾人為他們搖籤，搖出馬提亞來；他就和十一個使徒同列。 \end{tabularx} \\ \\
[1ex]
\hline
\hline
\end{longtable}
$^{1}$弟兄姊妹平安.
今天很感恩能夠在此與大家同領聖餐.
特別是與殷描的肢體.
讓我們在聽到之先都一同低頭祈禱.
「聖恩主,我們多謝你.
是因為你將我們納入了救恩之內.
聖靈求你親自來幫助我們.
謙卑我們的心.
今天專心聆聽你的話.
讓你的話進入我們的生命.
更新我們.
我們禱告奉主耶穌基督的名字祈求.
阿們」.
有人說「創業難,首業更難」.
一個企業要將做得好的傳統價值.
承傳下去靠的往往是一個好的團隊.
大家齊心擁抱同一個使命意象.
並且一代一代延續下去發揚光大.
在商場上也有一個成功例子.
可能大家都很熟悉的李錦記.
他們在1888年起家.
至今已傳到第六代的傳人運作.
並且業務不斷發展,蒸蒸日上.
作為基督徒.
天國的起點是由耶穌開始.
但亦由一代一代的信徒延續下去.
時至今日.
福音的使命亦交給了我們.
今天我們開始看一卷新的書卷.
《仕途行傳》.
選讀這卷書的原因.
是因為想配合今年教會的主題.
讓我們都能夠從第一間教會.
第一代信徒的經歷來得著激勵.
好讓我們今天好好的把握時機.
將神的國延伸出去.
讓更多人認識上帝,得著福音.
《仕途行傳》的作者是路加.
他分別寫了《路加福音》和《仕途行傳》.
在《仕途行傳》的第一章一節.

$^{41}$我們知道收信人是提亞非羅.
和第一本書《路加福音》的收信人是一樣的.
提亞非羅應該是當時的一個高官.
是一個有權勢的人.
在《路加福音》裡路加稱呼他為大人.
但去到《仕途行傳》.
似乎兩人的關係變得親密了.
所以沒有加上大人的尊稱.
今天我們會看完《仕途行傳》的第一章全章.
這章發揮著一個承前啟後的作用.
將耶穌地上的工作和信徒日後的事工連接起來.
首先我們會看到耶穌將天國擴展這個使命.
親手交給他的門徒.
由我到你們.
一章一節裡路加說.
我在第一本書已經論到耶穌從開頭所做和所教導的一切事.
這裡的開頭是指在《路加福音》裡講到耶穌地上的工作.
同時這些工作也是教會的開頭.
使徒是接續著耶穌在世時所做的工作繼續下去.
因此耶穌的使命也是日後教會的使命.
第三節說到由耶穌死後復活直至到他升天的期間有四十天時間.
耶穌復活後就向信徒顯現.
除了向他們證明他真的復活了之外.
第三節說他也講述神的事.
復活的耶穌把握這四十天時間為門徒做了一次特訓.
預備他們接續他在地上的工作.
耶穌的焦點是放在神的國那裡.
不過我們看看第六節.
門徒在耶穌升天之前最後一次聚集的時候.
竟然這樣問耶穌.
「主啊,你就要在這時候復興以色列國嗎?」.
問題反映著他們的眼光仍然停留在地上的事.
而不是天上的事.
他們的關注是以色列何時復國.
我們這群以色列人何時才可以吐氣揚眉.
在耶穌的門徒裡仍然有人未能夠掌握耶穌的心意.
與上文第三節「耶穌以神的國為焦點」有一個很大的落差.
第七節耶穌回答他們.
「父憑著自己的權柄所定的時候和日期.
不是你們可以知道的」.

$^{81}$意思是這些事由天父決定.
你們沒有權知道.
意味著他們所緊張復國的這件事.
並非神這個時候要他們所關注的.
今天我們會不會也像這些門徒.
曾經有一些問題去問上帝.
但神沒有給你答案.
可能是關乎你的人生大事.
也可能是關乎社會,國家,世界的問題.
這些找不到答案的問題.
可以纏繞著我們揮之不去.
結果帶給我們沮喪甚至憤怒.
有沒有想過當神不把答案給你的時候.
是想告訴你其實你不需要知道這個答案.
因為這個不應該是這個時候你的生命焦點.
如果能夠轉個角度用上帝的視野.
可能可以幫助我們找到此時此刻的生命,使命和意義.
教會裡不少熱心服侍的弟兄姊妹.
都是曾經經歷過生命的難關.
都問過一些找不到答案的問題.
但當他們換個角度之後.
反而可以用他們的經歷來服侍上帝,祝福其他人.
不知道你知不知道教會有一個恩樂小組.
當中服侍的肢體都是自己或家人曾經患過癌病.
作為過來人,他們曾經都經歷過恐懼,迷惘,難過.
但同時因為他們是過來人.
他們很明白患者和他們家人的需要.
當他們帶著千國的眼光.
本來那些很難的經歷.
反而今天成為他們服侍他人的資源.
第八節那裡,耶穌交給他們的使命.
不單止是他們當時身處的耶路撒冷一帶.
更加是地域上由近到遠去擴展,直到地極.
稍後在《仙途行傳》裡面.
你會看到這間教會的發展.
真是由耶路撒冷開始.
去到撒瑪利亞和猶太的地方.
再到小亞細亞,歐洲的羅馬,再延伸出去.
第八節,這節的經文,其中有一個很重要的部分.
就是耶穌邀請門徒,你要作我的見證.

$^{121}$「作」這個字應該翻譯為「彰示」,是一個未來式.
耶穌基督要我們成為祂的目擊證人.
然後告訴其他人福音的大能.
而在《仙途行傳》裡面.
不少信徒正是在苦難,逼迫之中來見證上帝.
就是那些很困難的處境.
他們仍然堅守信仰.
忠心達成主交給他們的任務.
信服上帝的命令.
你試想想,如果你是當時的門徒.
現在耶穌在你面前,就跟你說.
今天你要成為我的見證人.
不單止在這裡,還要去到地極.
你會有什麼反應?.
可能你會覺得這個責任很大.
這個擔子太重了,怎麼做得到?.
連耶穌你也是留在耶路撒冷加里尼一帶傳道.
但我們要傳到地極,怎麼可能?.
究竟能夠擔起福音使命的能力.
是從哪裡而來的呢?.
我們進入經文的第二點.
原來要完成神的工作.
不是靠人的努力.
而是從聖靈裡得力.
第二節提到耶穌藉著聖靈.
吩咐所揀選的仕途.
事實上,福音書裡的耶穌.
與聖靈的關係是很密切的.
當祂受洗的時候,有聖靈降臨.
當祂開始事奉之仙,祂被聖靈充滿.
其實耶穌也是充滿著聖靈的能力.
來展開祂的事奉.
第四,第五節.
當耶穌把大使命的福音棒交給門徒之前.
祂就祝福他們要等候.
第四節說.
「不要離開耶路撒冷,但要等候父的應許.
就是你們聽見我說過的,約翰是用水來施洗.
但是過了不多幾天,你們要在聖靈裡受洗.
神應許聖靈會降臨.

$^{161}$時間就是在幾天之後.
就是在仕途行傳第二章所記述的五秦節的時間.
那時候就好像第八節那裡的寓言.
門徒就會滿有聖靈的能力.
也有帶著主的權柄.
因為聖靈降臨在你們身上.
你們就必得著能力」.
起樣福音的十四章十六至十七節說到.
聖靈是耶穌升天之前向父求的一位保衛師.
聖靈永遠與他的門徒同在.
聖靈的出現是代替了基督在世上的工作.
去引導門徒,幫助門徒.
教導他們去承傳耶穌所給予的使命.
今天如果我們忽視聖靈的工作.
很容易就會用自己的能力來承擔上帝的事工.
沒有看重聖靈的教會.
很容易高舉人的成功.
以為專業,經驗,能力就可以達成天國的事.
結果就會引來兩個極端.
第一,人就會變得自負,驕傲.
因為覺得是自己的能力做到的.
另一個極端,就是人變得自卑.
只看到自己不行的地方.
其實兩個情況都是拒絕聖靈與我們同工.
在窩福三十週年的時候.
出了一本感恩見證集.
叫做《窩福的故事》.
我想問問在座各位有多少人看過?.
都有些手.
我在作為同工的時候被派一本.
最近翻看的時候.
發現窩福初期的歷史裡.
其實神跡處處.
你試想一下.
一間剛剛成立的教會.
怎樣可以在這麼短時間內.
建立更多的分堂?.
窩福成立於1976年.
五年之後,即是1981年.
它開設第一間分堂.

$^{201}$再過兩年,1983年.
同一年開了兩間分堂.
分別在九月開了景福堂,在屯門.
十月開了康福堂,在沙田.
每個月開一間.
幾乎是沒有可能的事.
我相信這是聖靈的能力.
臨到當時的教會.
臨到一群愛主的弟兄姊妹當中.
而他們信服上帝,獻上自己.
不辭勞苦地去拓展福音的事工.
當時去屯門,去沙田.
沒有今天這麼方便.
但他們就是這樣.
讓聖靈來與窩福同工.
從溫窩福的歷史提醒我們.
今天我們繼續要謙卑的等候神.
倚靠神.
我們每一個信徒的生命.
都需要充滿聖靈的能力.
求神賜給我們聖靈.
求神幫助我們倚靠聖靈.
並且有那份勇氣,力量.
繼續延續祂賜給我們天國的使命.
經文第九至十一節.
記述耶穌將使命交付給門徒後.
接著祂就升天了.
耶穌升天的時候.
是被一朵雲將祂接去的.
然後祂就從門徒的眼前消失了.
看著耶穌升天消失.
門徒當時的反應是怎樣的呢?.
第十節說到他們當時定睜望天.
可能那一刻太震撼了.
他們呆了.
接著第十一節.
有兩個穿著白衣的天使.
出現在他們面前.
跟他們說.
「加利利人,你們為何站在那裡望天呢?」.

$^{241}$相信當時的門徒百感交集.
一方面看見耶穌消失了.
心裡難過不捨.
沒多久之前祂才從死裡復活.
跟他們一起四十天.
很好的.
現在突然間又沒有了.
而且消失之前.
祂還將這麼大的任務交給他們.
怎麼辦呢?.
很惆悵.
在第十一節的下半節.
天使就跟這群門徒說.
你們看見祂怎樣升上天去.
祂也會怎樣來臨.
天使在這裡給這群門徒一個定心丸.
耶穌不是離棄他們.
祂是會再回來的.
並且祂的回來跟祂的升天是一樣的.
是這些門徒可以親眼看見的.
耶穌的再來.
對門徒來說是一份安慰.
帶給他們盼望.
同時也鼓勵他們要努力去做好.
耶穌交給他們要做的事.
將來當耶穌再回來的時候.
他們要向耶穌交仗.
耶穌升天了.
沒有了耶穌的門徒.
要做些什麼呢?.
現在沒有了耶穌.
他們不可以再做旁觀者了.
他們又呆望耶穌的消失.
不知道該怎麼辦.
直到當他們定神的時候.
這群門徒就變得積極起來.
由旁觀被動者的角色.
變成行動者.
積極地去預備聖靈的來臨.
在第十二節至二十六節的段落.

$^{281}$記載了這群門徒在預備的時候所做過的事.
就是禱告和補損仕途.
第十二至第十四節.
記載門徒按照耶穌的吩咐.
他們沒有散開走路.
他們去到耶路撒冷一個樓房裡.
同心合意地行切祈禱.
這群祈禱者的名單.
在第十三節清楚列出.
除了出賣耶穌的加勒人猶大之外.
其餘十一個使徒的名字都記下了.
而且還有幾個婦人.
包括耶穌的母親瑪利亞和耶穌的兄弟.
這群人代表了當時和耶穌很親近.
目擊耶穌死,復活,升天.
一直很忠心地跟隨耶穌的一群人.
在這個等候的時間.
他們遵循耶穌的教導.
留在耶路撒冷.
同心合意聚集一起來禱告.
經文第十五至第十八節.
交代了猶大的下場.
第十五至第十八節說到.
一個當時曾經三次不認耶穌.
被耶穌挽回的門徒彼得.
他竟然勇敢地站起來.
擔起了領袖的角色.
他還在大約120人面前對他發言.
他說出賣耶穌猶大的下場是悲慘的.
他凋棄了使徒的職事.
為了三十塊錢作惡.
他的結果就是土腹崩裂.
腸子都流出來.
他因此而死.
亦因此而一身蒙羞.
耶穌立十二個使徒的職份.
是重要的事奉崗位.
為了幫復活主作見證.
現在猶大已經失去了這個職份.
要找合資格的人來補償.

$^{321}$這是神的心意.
並且有聖經為證在第20節提到.
這十二個門徒.
亦標誌著舊約裡以色列十二個支派的延續.
在舊約裡十二個支派代表以色列的整體.
而在新約裡.
十二個使徒亦帶著一定的代表性.
因此在第21至26節.
彼得帶領當時的門徒商討補選使徒的事.
彼得首先列出候選人的條件.
就是這個人必須常常與他們同在.
並且由施洗約翰起.
直至耶穌被接升天為止.
他都一直見證著這些事發生.
這樣的人才是一個無可置疑的見證人.
第23節提到.
符合這些條件的.
他們找到兩個合資格的人.
就是巴薩巴和瑪提亞.
提名之後眾人一齊來禱告.
求神來選擇合適的人.
最終透過姚簽.
姚出瑪提亞.
他與11個使徒並列.
補上這個位置.
姚簽在聖靈降臨之前.
是一個合宜的做法.
在舊約作決定的時候.
都會用這個方式.
不過稍後當聖靈降臨之後.
就不需要再這樣做了.
不過在此想強調.
選出瑪提亞.
最重要的不是姚簽.
而是之前的同心祈禱.
由第12節至26節.
這個大段落.
我們見證著這群本來膽怯軟弱.
望著耶穌升天法外的門徒.
怎樣轉變過來.

$^{361}$他們由旁觀者變成行動者.
但他們的行動沒有比上帝快.
他們是以禱告為先.
才去求問上主的心意.
一次又一次的同心祈禱.
祈禱是一切事公的基礎.
這群門徒首先向我們見證的.
就是怎樣用祈禱尋找上帝的心意.
怎樣將祈禱成為他們作決定.
成為他們日後事公的基礎.
科福每逢星期三晚上8點至9點.
都有祈禱會.
你有沒有參加呢?.
這次我不敢叫你舉手了.
幾年前當我們在這個禮堂重啟.
星期三晚祈禱會的時候.
有八多人在這裡參加.
但到今天每個星期三晚上.
實體和網上參與的人數.
都只有十多人.
而且當中還有幾個是同工.
求主去帶領我們.
激勵我們去效法這群門徒.
以禱告為先.
以禱告來復興我們的生命.
以禱告來復興我們的教會.
總結《使徒行傳》的第一章.
描述了耶穌升天前.
將拓展天國的使命.
交予給門徒.
教導他們要倚靠聖靈的能力.
而不是自己的能力去承擔.
經文給我們看見一群.
忠心跟隨耶穌的門徒.
如何由旁觀者變得勇敢.
走出來保衛.
更重要.
他們信服基督.
以禱告為先.
行切為神要成就更大的事.

$^{401}$來作出準備.
今天作為窩福的一份子.
上帝也將這個.
拓展天國的接力棒傳給了我們.
你願不願意效法門徒.
勇敢去接棒呢?.
今日窩福的處境.
唐主任離開了.
消失了.
我們沒有唐主任.
但神祭藉著這個處境.
呼召我們每一個起來.
我們要承擔今日教會的使命.
就好像當日的門徒一樣.
有時我們也很想.
唐主任快點來.
何時有?.
這也是我們今天沒有答案的問題.
但肯定的是.
耶穌說祂會再回來.
祂給我們這個拓展天國.
傳揚福音的職份.
我們有沒有看重呢?.
當耶穌再回來的時候.
我們會像猶大那樣.
羞辱不敢見主.
因為我們撇棄了我們的職份.
或我們像其他的使徒.
戰戰兢兢.
但我們仍然去努力開展上帝的工作.
當見到耶穌的時候.
我們可以很開心.
很興奮地去報告.
神的國是這樣這樣.
由我們的教會那裡開展.
我們和上帝一起來感恩.
在經文的大綱上.
我列了三個反思的問題.
盼望大家也花一點時間.
去思考一下.

$^{441}$今天你的事奉.
你的角色.
你的群體.
我給大家一點時間安靜.
沉澱一下今天的訊息.
思考一下這些問題.
然後我會帶大家有一個禱告.
(禱告).
讓我們一起來對禱告.
我們求聖靈光照我們.
幫助我們去反思.
我們今天的事奉.
我們的角色.
我們和教會的關係.
在這裡學習.
在這裡成長.
當今年我們一起慶祝五十週年的這個時刻.
她們辛勞地來對服侍主.
靠著聖靈的能力.
建立了今天的窩福堂.
直至現在.
或者過去我們有些時間.
只是做一個旁觀者.
但今天我們求聖靈充滿我們.
賜下那份熱情和能力.
去承擔這個使命.
去繼續擴展主你的國度.
求你加我們力量.
讓我們能夠活出一個與蒙恩相稱的生命.
每天帶著敬畏的心.
喜樂的心.
來對服侍主,服侍人.
我們這樣祈禱.
奉主耶穌基督的名字祈求.
阿們.
\newpage



\section{約書亞記 20:2}
\label{sec:032172}
\textbf{聖經人聖經事:設立逃城}~~~~黃克勤牧師博士
\newline
\newline
連結: \href{http://www.church.com.hk/acms/content.asp?site=cdc\&op=show\&type=product\&code=032172\&layout=sermon}{\texttt{032172}} ~~~~ 語音日期: 2026-05-22
\newline
\newline
\hyperref[sec:032171]{\small{< < < PREV SERMON < < <}}
~
\hyperref[sec:index_chronic]{\small{[返順時目]}}
~
\hyperref[sec:032173]{\small{> > > NEXT SERMON > > >}}
\newline
\newline
約書亞記 20:2
\newline
\begin{longtable}{cl}
\hline
\hline
章節 & 經文 (和合本修訂版)\\
\hline
20:2 & \begin{tabularx}{0.7\textwidth}{X} 「你吩咐以色列人說：『你們要照我藉摩西所吩咐你們的，為自己設立逃城， \end{tabularx} \\ \\
[1ex]
\hline
\hline
\end{longtable}
$^{1}$大家好,又一齊讀聖經.
《生命述書》的網絡跟大家看聖經人,聖經事.
這次看到在以色列人進入迦南地之後.
全部定了,十二支派都分了地.
大局已定,算是頗為太平的.
特別是開頭那段時間.
約書亞也還在,神很眷顧他們.
靈性相當好,後來就差了.
在這時間,神叫以色列人做了一件事.
就是設立逃城.
逃走的城,可以躲藏的城.
設立逃城並不是叫他們玩捉迷藏.
是城可以被人躲藏.
逃進去城裡面,避免一些事.
但不是玩的,是很認真逃避索命的人.
要拿他的命,逃城看看.
約書亞結第二十章.
是耶和華曉遇約書亞的.
叫他設立逃城.
做什麼的呢?.
「你吩咐以色列人說:你們要照著我,正摩西所曉遇你們的,為自己設立逃城」.
逃走用的城,其實是做什麼的呢?.
逃城很多解經家認為是特登神做的,一來是為了公平,為了對人命的愛惜.
但逃城也可以說是某方面象徵耶穌基督.
象徵救贖,象徵白白的救贖和清義.
先說說逃城.
「使那無心而誤殺人的,可以逃到哪裡?.
這城可以作你們逃避報血仇人的地方」.
是無心而誤殺了人.
當日是以牙還牙,以眼還眼,以命還命的時代.
大家都覺得這樣才公道,其實也有公道的地方.
例如以眼還眼,以牙還牙,什麼意思呢?.
如果不小心弄盲人家的眼,對方的家人可以要求弄盲你一隻眼.
以牙還牙,打架打甩人家的牙,要打甩人家的牙.
其實你可以說,這麼野蠻的,賠錢不行嗎?.
當然可以談,但最主要也算是一個保障.
打我兒子,打到他盲了一隻眼.
有時我可能帶兄弟強壯去打你家人.
打爛你家人,打死你兒子,如果遇到他,將你兒子殺死.
因為很生氣,但只能保障他以眼還眼,以命還命,以牙還牙.

$^{41}$你不能傷害得多,但有時無心而誤殺人.
他可以合法地逃命,逃去屠城.
在這些城可以做逃避,報血仇人的地方.
即是躲進屠城,他就不能追溯你的命.
只能在外面等你出來,做什麼呢?.
其實就是被無心而誤殺人.
例如在我這裡劈柴,拿著斧頭劈.
斧頭斷了,那塊鐵飛了出去.
剛好打到另一個人,他暈倒,死了.
又或者我拿刀斬樹枝,刀飛了出來,插死了人.
又或者我無心地斬東西,掉下來打死一個人之類.
總之就是無心而誤殺人,家人要追溯他的命.
這個人就可以逃到屠城,逃避報血仇人的地方.
那殺人的要逃到這些城中的一座城.
設立了城,去到的時候站在城門口.
站在城門口的原因就是.
因為當日以往那些管理城的人,他們是坐在城門口的.
一來看著誰進來,什麼人在哪裡來.
或者在一個固定的地方,人們要去那裡報告訊息.
有什麼要判斷,他打爛了我的東西,我要他賠償,他不肯賠償.
去找人做審判的時候,做公道的時候,都去城門口.
所以在城門口做官的人,做判斷的人就在那裡.
去到城門口,就將你的事給城內的長老講給他聽.
長老是否一定坐在城門口呢?可能是,或者不是.
可能有當值的坐,又或者有其他官員坐.
不過,你是來幹什麼的?我就是不小心誤殺了人.
我知道這是屠城,所以我來這裡想避一下.
你等等,我先叫長老來,或者先帶他去長老那裡.
他們就把他收進城裡,給他地方,使他住在他們中間.
他裡面的長老知道之後,如果他是誤殺人的.
就要給他地方住,供應他住,還有讓他住在這裡.
就收留了他,那不就很划算了,殺了人就跑去那裡.
然後說,哎呀,我誤殺,這樣就不是了.
若是報血仇的人追了他來,長老就不能夠說.
這樣嗎?他在這裡,你去捉他吧,不行.
長老不可將他交在報血仇的手裡.
因為他是素無仇恨,這是條件來的.
素無仇恨,是無心殺了人的.
什麼意思?都說了,斬個斧頭飛了出去就打死了人.
這些是無心殺了人的,因為他是素無仇恨的.

$^{81}$或者他安裝了陷阱,捉野獸,誰知道會掉下去.
素無仇恨,這樣就算了,留意.
他要住在那城裡,他不能離開.
第二,他要做的不是誤殺,我要找地方躲,你收留我吧.
是要先見了長老,他還要站在匯眾面前聽審判.
即是說他不是就這樣走去算的,不是這麼簡單的.
你不要利用這樣躲起來,我無心殺了人.
我跟那個人素無仇恨,不過那個斧頭飛了出去.
真是不好意思,我都很內疚,唯有跑來這裡.
他要聽審,要審他,真是假的.
你怎會這樣做,你的斧頭飛了多久,那個人走到哪裡.
當時怎樣發生,就要問他,在匯眾面前.
即是說有旁聽的人,一起聽.
第二,站在匯眾面前聽審判,這個步驟是一定要做的.
不是就這樣隨便說,我是誤殺,這樣行不行.
如果就算經過審判發覺是了,他就要在城裡不可以出來.
等到那時的大祭司死了,當年整個以色列只有一個大祭司.
你要在城裡住到多久呢,就要住到大祭司死了.
當時的大祭司很重要,他們覺得大祭司代表一個年代.
大祭司的年代過了,新的年代開始,你有機會可以出來.
大祭司死了,殺人的才可以回到本城本家.
就是他所逃出來的那城,他由本城本家在那裡.
說殺了人,不殺,所以跑出來.
等大祭司死了,他才可以回到他本身自己的家鄉.
問題就是,如果大祭司年紀很老了.
走進逃城那裡避,很短時間大祭司死了.
他就可以回到家裡,回到家裡家人還可以追殺他.
其他對方的家人是不可以了,因為他進了逃城.
大祭司死了,他就可以回到家裡,沒事了.
那就很值得了,因為他是無心殺人.
神要眷顧庇護這些無心殺人,不能說這件事是意外.
如果大祭司新上任,很年輕又很強壯,那怎麼辦呢?.
等,我死他都還沒死,那怎麼辦呢?.
那你就先死吧,出不來了,這裡講得很清楚.
於是以色列人就拿了城出來.
在拿忽他利支派的地方有基底斯.
以法連的地方有士劍.
猶大的山地有希伯倫.
以前叫基列亞巴,大家記得了.
呂便支派,即是劉便支派,在平原有比色.

$^{121}$在迦德支派有拉末.
嗎哪西支派有哥蘭,有六個地方.
這六個地方分佈在哪裡,以色列的屠城呢?.
地中海,約旦河,加利利海,死海.
約旦河的西有基底斯,士劍,希伯倫.
這個地方,這個地方,這個地方.
哥蘭在北邊,拉末在這裡,比色在這裡.
這個比較平原,約旦河這邊是平原,這邊是山地.
有六個地方,可以跑去那裡.
借刀是為以色列眾人和他們中間寄居的外人所分定的誠意.
以色列人可以跑去,住在那裡的外人.
有誤殺發生,無冤無仇的,不過意外.
他都可以跑去屠城,使誤殺的人可以屠到那裡.
不死在報血仇的人手中.
留意,等他站在匯眾面前聽審判.
留意,其實他是進了去,聽一次審判.
看來你都真的誤殺,好吧.
你就不小心意外發生,你住在這裡.
但是如果追殺他的血,追殺他的命的人來到的時候怎樣呢?.
知道他跑去那裡,哎呀,你是不是陳大文?.
他一斧頭飛出來,打死了我兒子.
我現在要他償命,血債血償.
這次不得了,他在我們這裡,現在我們先定下審判.
叫陳大文出來,叫匯眾出來,長老坐定在這裡.
好,你們雙方對質一下,說一下究竟發生了什麼事.
如果發覺原來他是騙了我們.
那個死者的家人來到說這件事.
原來不是誤殺,不是意外.
是謀殺有意的,哦,這樣不得了.
他就會將那個人交給他們,以命償命.
但是如果真的誤殺,人家可以挑戰你.
再審,好了,等他站在匯眾面前聽審判.
當時就是想保護一些意外不小心出了人命的人.
不至於以命償命,就看生命是否重要.
其實這個屠城,當時是有這樣的規矩神定的.
有些不值得,由於我們心裡覺得.
以牙還牙,以牙還牙,以命償命.
就讓他走進去,這樣就很不值得.
是呀,那誰拿主意,神拿主意的.
我想保護他,你們不要再追討了.

$^{161}$這個是神的心意,當時是這樣.
有沒有傻一點,不傻一點,當時是這樣.
大家都一致的,這樣就可以了.
因為這樣的條例,不只是保護那個人.
將來你也有機會可以無意.
無愁無怨,意外殺死了人,不知道.
譬如這樣,你也有這個優惠.
也可以有這個逃命的地方,就是屠城.
所以其實屠城是有某一個象徵的意義.
就是象徵你進去追命的人.
就不能夠再殺你了.
就好像進去也是一個逃避的地方.
有少少象徵著耶穌基督是保護了我們.
逃避了索命的人的血債.
有象徵的意義.
設立這個屠城,也當成他們的辦法是神設立的.
我當然不敢挑戰,當時可能覺得很合理也說不定.
沒得說了,意外也說不定.
說這件事給大家知道,當時是有這樣的事發生的.
就是設立了六個屠城.
設立屠城就是他們當時的一個法例.
或者一個practice.
是一個辦法給無意殺害了人的人.
有一個可以避開的機會.
只是意外.
你說是文明,也是很文明的.
也很對的,是對於生命的愛惜.
也是象徵著耶穌基督,好像屠城一樣.
任何的人去到裡面都是被赦免的.
當然屠城就是被那些無意殺人誤殺的.
而我們這些有罪的人去到耶穌基督那裡.
是同樣被保護,被尋索.
我們命的就不能夠追討到我們了.
是新的一個生命的開始.
我今天跟大家分享到這裡.
設立屠城是以色列歷史裡面曾經出現過.
都挺有趣的東西.
跟大家分享一下他們當時.
是神叫他們有這個習俗的.
有這個辦法與規條的.

$^{201}$好,我們一起來祈禱吧.
天父上帝,雖然我們今天用我們現今的看法.
未必完全贊同這個屠城的做法.
但是當日在以色列人當中.
他們是接受的,也樂意接受神的命令.
也保護了一些意外無意傷害到別人生命的人.
以免他因此而死亡.
當然其他的,未傷到生命的.
就以眼還眼,以牙還牙,或者以錢來代替了.
不過在今天我們看到這個屠城.
在那個時候,可能那些人是覺得很好,很接受的.
都是對生命的保護.
主也幫助我們,雖然未完全明白.
不過我們都欣賞上帝來當日做這件事.
也都原來當日的人是很接受的.
是覺得這樣其實是進步與文明的.
主我們幫助,讓我們心胸更廣闊.
有時候我們都會被人無意傷害到.
有些人是他的成長過去背景.
令到那個人三尖八角,說話都傷害到人的.
有些人是很容易去罵人的.
很多人是很容易去傷害其他人.
是有意甚至是著意去傷害人的.
我們都曾經經歷過這些被傷害.
主是幫助我們,今天你就是我們的屠城.
我們可以走進去.
也都是知道每一個人在基督裡面都是被赦免的.
就是我們來到你面前,都是公平的.
求主給我們有這個寬闊的心胸.
誰聽祈禱,同靠耶穌基督命求.
阿們.
\newpage



\section{以弗所書 3:1-21}
\label{sec:032173}
\textbf{在受限中活出無限}~~~~駱瑞麟傳道
\newline
\newline
連結: \href{http://www.church.com.hk/acms/content.asp?site=cdc\&op=show\&type=product\&code=032173\&layout=sermon}{\texttt{032173}} ~~~~ 語音日期: 2026-05-24
\newline
\newline
\hyperref[sec:032172]{\small{< < < PREV SERMON < < <}}
~
\hyperref[sec:index_chronic]{\small{[返順時目]}}
~
\hyperref[sec:032174]{\small{> > > NEXT SERMON > > >}}
\newline
\newline
以弗所書 3:1-21
\newline
\begin{longtable}{cl}
\hline
\hline
章節 & 經文 (和合本修訂版)\\
\hline
3:1 & \begin{tabularx}{0.7\textwidth}{X} 因此，我—保羅為你們外邦人作了基督耶穌囚徒的，替你們祈禱。 \end{tabularx} \\ \\ \relax
3:2 & \begin{tabularx}{0.7\textwidth}{X} 想必你們曾聽見神賜恩給我，把關切你們的職分託付我， \end{tabularx} \\ \\ \relax
3:3 & \begin{tabularx}{0.7\textwidth}{X} 用啟示讓我知道福音的奧祕，正如我以前略略寫過的。 \end{tabularx} \\ \\ \relax
3:4 & \begin{tabularx}{0.7\textwidth}{X} 你們讀了，就會知道我深深了解基督的奧祕； \end{tabularx} \\ \\ \relax
3:5 & \begin{tabularx}{0.7\textwidth}{X} 這奧祕在以前的世代沒有讓人知道，像如今藉著聖靈向他的聖使徒和先知啟示一樣， \end{tabularx} \\ \\ \relax
3:6 & \begin{tabularx}{0.7\textwidth}{X} 就是外邦人在基督耶穌裡，藉著福音，得以同為後嗣，同為一體，同為蒙應許的人。 \end{tabularx} \\ \\ \relax
3:7 & \begin{tabularx}{0.7\textwidth}{X} 我作了這福音的僕役，是照著神的恩賜，是照他運行的大能賜給我的。 \end{tabularx} \\ \\ \relax
3:8 & \begin{tabularx}{0.7\textwidth}{X} 雖然我比眾聖徒中最小的還小，他還賜我這恩典，讓我把基督那測不透的豐富傳給外邦人， \end{tabularx} \\ \\ \relax
3:9 & \begin{tabularx}{0.7\textwidth}{X} 又使眾人都明白甚麼是歷代以來隱藏在創造萬物之神裡的奧祕， \end{tabularx} \\ \\ \relax
3:10 & \begin{tabularx}{0.7\textwidth}{X} 為要在現今藉著教會使天上執政的、掌權的知道神百般的智慧。 \end{tabularx} \\ \\ \relax
3:11 & \begin{tabularx}{0.7\textwidth}{X} 這是照著神在我們主基督耶穌裡所完成的永恆的計劃。 \end{tabularx} \\ \\ \relax
3:12 & \begin{tabularx}{0.7\textwidth}{X} 我們因信耶穌，就在他裡面放膽無懼，滿有自信地進到神面前。 \end{tabularx} \\ \\ \relax
3:13 & \begin{tabularx}{0.7\textwidth}{X} 所以我求你們，不要因我為你們所受的患難喪膽；這原是你們的光榮。 \end{tabularx} \\ \\ \relax
3:14 & \begin{tabularx}{0.7\textwidth}{X} 因此，我在父面前屈膝— \end{tabularx} \\ \\ \relax
3:15 & \begin{tabularx}{0.7\textwidth}{X} 天上地上的各家都是從他得名的－ \end{tabularx} \\ \\ \relax
3:16 & \begin{tabularx}{0.7\textwidth}{X} 為要他按著他豐盛的榮耀，藉著他的靈，使你們內心的力量剛強起來； \end{tabularx} \\ \\ \relax
3:17 & \begin{tabularx}{0.7\textwidth}{X} 又要他使基督因著你們的信住在你們心裡，使你們既在愛中生根立基， \end{tabularx} \\ \\ \relax
3:18 & \begin{tabularx}{0.7\textwidth}{X} 能夠和眾聖徒一同明白基督的愛是何等的長、闊、高、深， \end{tabularx} \\ \\ \relax
3:19 & \begin{tabularx}{0.7\textwidth}{X} 並知道這愛是超過人的知識所能測度的，為要使你們充滿神一切的豐盛。 \end{tabularx} \\ \\ \relax
3:20 & \begin{tabularx}{0.7\textwidth}{X} 神能照著運行在我們心裡的大能充充足足地成就一切，超過我們所求所想的。 \end{tabularx} \\ \\ \relax
3:21 & \begin{tabularx}{0.7\textwidth}{X} 願他在教會中，並在基督耶穌裡，得著榮耀，直到世世代代，永永遠遠。阿們！ \end{tabularx} \\ \\
[1ex]
\hline
\hline
\end{longtable}
$^{1}$各位弟兄姊妹,平安!.
開始的時候,我們先做一個祈禱.
求你這時候,你的靈就與我們同在.
耳講的,和聽的,都能夠同感一靈.
同得造就.
禱告奉主耶穌基督的聖名.
阿們.
有人呢,形容作家.
就是那些什麼.
爬格子動物,大家有沒有聽過.
即是說,它在一個有限的紙張裡面.
一行一行地來書寫.
很多人類很偉大的.
歷史,偉大的作品,都是這樣誕生出來的.
現在不同了,電腦發明了.
作家成了.
鍵盤戰士.
就好像那些鋼琴家一樣.
不斷地敲擊那個鍵盤.
但同樣地,亦都產生了很多偉大的樂章出來.
但當保羅,他寫這個《以弗所書》的時候.
他不是在一個很舒適的環境裡面寫的.
而是在羅馬監獄裡面.
即是說,在一個有限的間隔裡面.
更加有士兵來看守.
加上當時的紙張亦都很貴.
那些所謂蒲草紙,羊皮卷那樣.
所以亦都不可以隨便地來更改.
即是說,他在當時裡面是很受限制的.
甚至根本是我們按法來想像.
但很奇怪的就是.
從他的字裡行間.
我們看不到他有一絲的絕望.
他的視野更加是非常之遼闊.
那為甚麼一個人在這個監獄裡面.
竟然能夠寫出這個這麼偉大的訊息呢.
由於保羅的處境,真是值得我們今天來學習的.
在人生的限制裡面.
未必能夠限制神的豐盛的恩典.
今天就讓我們從我們剛才讀的經文.

$^{41}$我們一同來思想一下.
我會分三個層面來跟大家一起讀下去.
首先第一點就是有限的困局.
開頭第一句說話.
保羅就用了「避囚」這個字來形容.
即是說他很明白福音本身就是一條受限的道路.
人是沒有辦法來迴避的.
換句話說,不是因為保羅他能力不足.
而是因為他是忠於神的托父.
大家有沒有想過.
這位足跡曾經是走遍亞非拉的宣教巨人.
但現在他的活動範圍.
只不過是在羅馬監倉裡面的一個間隔.
並且是受到極大的限制.
也許這也是我們生命裡面.
我們所謂的常受的一些限制.
有時候這些限制,這些鎖鏈.
對我們來說叫做時間.
例如我們每一天的工作,我們的家庭.
我們很多繁瑣的事不停地追趕著我們.
令我們可以說是透不過氣來.
而有時候我們這種限制,這種鎖鏈就叫做身體.
例如一場突然間的大病.
或者隨著我們每一個人年紀的越來越增長.
我們都變成了senior.
身體可能會越來越有些毛病,有些衰弱.
常常我們自己都會有些意識到.
我們心靈雖然是十分之願意.
但肉體卻是很限制.
而最可怕的就是.
當這種限制大到一種程度的時候.
我們可能會感覺到很沮喪,很枯乾.
所以保羅在第十三節那裡就說到.
「不要因我所受的患難而來喪膽」.
喪膽這個字在原文來說是很氣餒,很灰心的.
這是一個希伯來的思維.
是在說一個心靈萎縮的一個狀態.
這種狀態,這種被困的感覺.
其實也是對我們最殘酷的一個打擊.
所以很多時候我們作為基督徒.

$^{81}$我們會有一個很大的誤區.
以為有神的同在.
我們就應該沒有所有的限制.
那到底是不是呢?.
其實我有一位老友.
他曾經也是一個.
在年輕的時候跟我說他.
有很大的抱負.
要創一番事業.
但是很可惜.
他因為家庭的債務.
以及種種的責任和壓力.
他被迫要為一間.
他很不喜歡做的公司來打工.
並且做了十年.
他跟我說.
「我上班就像坐牢一樣」.
「我的鎖鏈就是我每月的賬單」.
「我根本上找不到出路」.
「我看見的就是我的生命」.
「好像漫漫一樣一天一天來枯乾」.
這種被困住的感覺.
其實不單只是打工.
其實對我們每一個移民的人來說.
其實都有很大的感受.
大家在這個新的環境裡面.
你可能處處受限制.
不像你移民以前那麼得心應手.
如果有人說.
「為什麼你不喜樂呢?」.
「為什麼你不對神有信心呢?」.
聽起來有時候會覺得有點冒犯.
但是如果生命就是這樣的時候.
如果保羅也是這樣想的時候.
他可能只會帶著一個很遺憾.
來死在這個獄中.
但他卻不是這樣想.
所以你看到在第一節.
你會看到他用了一句修昔的語言來講述.
他說.

$^{121}$「我是為你們外邦人」.
「作了基督耶穌的被囚」.
這句話.
換句話說.
就是說保羅.
他不覺得他是凱撒的囚犯.
而是他是基督的奴僕.
大家有沒有覺得很不同呢?.
所以他在第二到十二節.
他就不再講監獄了.
他轉而講他需要承受的積分.
他是受神所委託.
要將歷世以來.
蘊藏在創世萬物.
創造萬物之神裡面的奧秘.
來講述出來.
奧秘這個字.
從來都不是講到.
那些令人無法明白的密碼.
或者那些言學.
意思就是一個.
本來就是隱藏在神的心意裡面.
的一個偉大的救贖計劃.
現在就在基督耶穌裡面.
徹底地向我們打開.
保羅其實在前一,第二章裡面.
也是有提及的.
雖然現在保羅.
他是被監獄的牆捆鎖住.
但是神的計劃.
正正是透過保羅的文字.
祂的禱告和祂的受苦.
向天上的執政掌權.
來展示祂神的百般智慧.
大家應該記得.
以弗所其實.
裡面是有世界最大的亞底瑪耶神殿.
靈界的氣氛其實非常重.
他這樣講其實是非常有意思的.
他在告訴我們.

$^{161}$神的應許是不會受限制的.
他更加有一種翻轉的能力.
保羅的受限.
正是展示出神在榮耀的舞台裡面.
來工作.
保羅更加在第八節裡面.
用了「七筆頭,unsearchable」的封城.
來形容基督.
我們都知道.
人的資源是很有限的.
無論你多有錢也好.
你的銀行戶口.
那本簿仔有多少個零.
但是也有盡頭的.
是不是?.
我們說回人的體力.
無論我們能夠抬多少百斤.
都是有見頂的時候.
但是基督的封城.
他是用「不可測量」來形容的.
為什麼說.
福音的價值不是取決於.
全港福音人的外面的境遇.
而是基督的恩典.
和基督的智慧.
保羅儘管他失去了所有的社會地位.
但是他跟其他信徒一樣.
他有一個特別的身份.
這個身份是什麼呢?.
就是他在第六節說到.
「在基督耶穌裡.
藉著福音.
得以同為後製.
同為一體.
同蒙應許」.
這就是保羅所說的奧秘.
他用了三個「同」字來表達的.
其實比上一章說到.
外邦人和猶太人要合而為一.
其實是更加進一步.

$^{201}$更加來說出神的心意.
就是說外邦人不需要先成為猶太人.
而直接在耶穌基督裡成為神的子民.
沒有什麼困鎖可以限制他.
保羅在監獄裡雖然是這樣.
但是他仍然可以跟外面的信徒.
連為一體.
福音的奧秘.
在這裡你看到.
是衝破了人肉體的限制.
並且神對人的應許.
對信徒的應許.
是一直存留到今天.
這裡也讓我想起.
在十九世紀有一位非常著名的盲眼詩人.
他叫做Fanny Crosby.
她出生之後六個星期.
就因為一次醫療事故而傷目新明.
按照人的邏輯.
她一生注定要在黑暗裡度過.
但是她曾經說過.
她說.
如果我能夠再選擇.
我會再選擇生來核硬.
因為當我到了天上.
我首先會看到的就是.
我親愛的主耶穌基督.
這種視野.
你看到她是完全顛覆了生命的限制.
不僅如此.
她一生人創造了八千多首贊吏詩.
在她內在的世界裡.
其實你看到她比任何開眼正常的人.
更加要遼闊.
神的無限.
使她在黑暗裡.
能夠成為千萬人的亮光.
所以保羅在第二十二節.
他就說道.
信徒我們因此要放膽無懼.

$^{241}$我們要篤信不疑地來到神的面前.
有人又說.
難道保羅你不是在坐牢嗎?.
你說得這麼口響.
但是保羅在這裡其實是要信徒.
他不要因為他坐牢而後退.
而是要憑著信心.
要相信神的應許.
神的計劃是多麼深奧.
而其實保羅這段話.
也同時在挑戰我們.
我們是否也願意接受.
神對我們受限的現狀呢?.
你會不會覺得.
這種現狀.
是有祂的旨意呢?.
但有時候我們必須要承認.
要捉緊神的應許.
其實談何容易呢?.
如果我們只是停留在現實的困苦裡.
我們常常會感到絕望.
但是如果我們停留在屬靈的理想裡.
其實我們可能會覺得.
我們好像脫離了現實一樣.
但是這正正是保羅在十四至二十一節.
他屈膝向神祈禱.
向我們展示的.
那裡他說神是無限的.
神的無限不是要帶領我們.
逃離我們的限制.
而是要進入我們.
並且要充滿我們這個受限的氣命.
所以在第十六,十七節裡.
保羅懇切地向神祈求什麼呢?.
按著祂豐富,豐盛的榮耀.
藉著祂的靈.
叫你們心裡的力量.
這裡心裡的力量.
原文就是說.
是內在的人.

$^{281}$Inner man.
來當強起來.
使基督因你們的信.
駐在你們心裡.
叫你們的愛心有根有機.
這句話其實是整段經文的轉捩點.
不是我們怎樣去突破我們這個限制.
而是基督是要進入我們這個限制.
環境雖然限制了我們外面那個人.
但是神的大能.
祂是藉著聖靈.
駐在我們裡面那個人那裡.
並且在那裡運行.
所以當基督是透過人不斷來信教祂.
而活在我們每一個人的心裡面的時候.
我們內面這個人.
現在可能叫做小宇宙.
物理的空間可能是很小.
但是心靈的空間卻是無邊無際.
我們再不需要受外面的環境擺佈.
監獄的牆壁.
雖然是限制了保羅的肉體.
但是它是沒辦法限制.
神在他們裡面的那個作為.
因此保羅在第十八及九節那裡.
他寫出他祈求裡面對神的感受.
他感受到什麼?.
他感受到基督的愛是多麼的長闊高深.
弟兄姊妹.
在平常的日子裡.
其實我們最喜歡追求的是什麼?.
可能是長和闊.
就是說我們事業更加寬闊.
我們成就更加廣.
但是萬難來到的時候.
我們外面的長和闊可能會受到限制.
但是這個時候.
神要開啟你另外生命的兩個屬靈的深度.
就是高和深.
其實保羅在第十七節那裡已經告訴我們.

$^{321}$他要我們的愛心.
是要有根有基.
什麼是根呢?.
不用多說了.
是樹的樹根.
就是說當樹的樹枝受到限制之後.
樹要逼著向下扎根.
這樣才能夠站穩.
基當然是指這個房子的地基.
一間房子如果需要建得高.
一定要挖得深.
其實這個也是說出限制的真相.
當外面的環境.
將你的長和闊壓縮的時候.
聖靈就會藉著這個限制.
來將你的生命往下挖,往下扎根.
好讓你的生命能夠承載.
神那個漆不透的豐盛.
就是說是高和深.
找我這裡有一次.
我去一間基督教的老人院做探訪.
並且也是參加他們的崇拜.
在那裡我沒有看到有什麼光輝烈烈的見證.
也沒有什麼驚天動地的神蹟.
我只看到的是一群頭髮發白.
行動很緩慢.
有些甚至是坐著輪椅.
有些需要參附.
有些甚至記憶力已經衰退.
連自己的名字都忘記了的老人家.
在人的眼裡.
他們的生命好像越來越受限制一樣.
但是在他們的臉上.
根本上看不到有一絲的恐懼和怨恨.
反而看到他們很平靜,很安穩.
整個氛圍都是這樣.
他們可能知道.
他們自己是慢慢走進人生的終點.
但是同時.
他們也可能明白到.

$^{361}$他們的終點就是天家了.
那一刻我突然間明白到.
真正的豐盛是什麼呢?.
不是外面的自由.
而是裡面有耶穌基督的愛.
這個愛是超過我們所能夠理解.
而不能夠測度的.
就正如這個加拉太書.
保羅所說的.
「現在活著的不再是我.
乃是基督在我裡面活著.
並且我如今在肉身活著的.
是因信神的兒子而活」.
這句話.
可能也是講出.
一個信耶穌人的心底裡的說話.
最後.
保羅.
就懇求神的大能.
就是活在信徒的心裡面.
在20-21節那裡.
當我們不再聚焦我們的生命是怎樣走出這個限制.
而是專注於讓神的豐盛來充滿我們的時候.
神就會透過我們這個有限的生命.
來彰顯祂運行我們裡面的大能.
就好像一條.
平時很平靜的河流一樣.
平時你根本很難看到祂那種澎湃的力量.
但是當這條河流突然間被中斷了.
水流從高處落下.
波瀾壯闊的場景.
那條大瀑布.
就完全地將神的隱藏能力.
彰顯出來了.
各位弟兄姊妹.
我們的人生.
我們沒有人喜歡.
有一個這麼大的落差.
因為我們很怕.
跌得很慘.

$^{401}$很傷.
但是如果沒有這個這麼大的落差.
我們永遠看不到神隱藏的智慧.
和基督那個漆不透的豐富.
但是神的能力.
不是用來讓我們去逃避我們所有的限制.
而是好像保羅那樣.
雖然他坐在監獄裡面.
是有能力來寫出.
那個影響我們人類文明兩千多年的書信.
並且向靈界.
向這個執政掌權的.
來展示神的那個智慧.
各位弟兄姊妹.
也許你現在也是在這個病痛中身陰.
也許你現在正在做那份很枯燥的工作.
也許你在這個破碎的關係中徘徊.
但是如果你能夠.
將基督那份.
長闊高深的愛.
來活出來.
你就在見證神與你的同在.
而這份影響力.
同樣是會穿透現狀.
是直達到永恆的.
總結一下.
我們剛才所說的.
我們會發覺.
保羅的身體雖然是受到困鎖.
但是他的心靈卻是進入了神的永恆的奧秘裡面.
他帶出的福音的信息.
是超越了這個時空.
是直到今天的.
會說原來神不是先除去人的限制.
才去彰顯祂的榮耀.
很多時候.
神是藉著限制來讓我們經歷祂的無限.
所以真正限制我們生命的.
未必是我們的環境.
而是我們的裡面.

$^{441}$是不願意給基督留下空間.
但是當基督住在我們裡面的時候.
監獄可以成為聖所.
鎖鏈可以成為見證.
各位電影姐妹.
由於我們今天每一個人的生命裡面.
都是有著無法揮去的種種的限制.
但是保羅的經歷.
使我們看到.
限制不是無限的對立面.
限制正是無限需要降臨的地方.
正如在哥林多後書十二章九節那裡.
神對保羅說.
祂說我的恩典是夠你用的.
因為我的能力.
是在人的軟弱上.
是演得完全.
所以我們之所以能夠在這個受限裡面.
能夠活出這個盼望.
不是因為環境是改變了.
而是因為主活著.
並且是活在每一個人的心裡面.
讓我們一起讀出.
保羅對神的頌讚第二十和二十一節吧.
由我這個同理來禱告.
慈悲人我的天父.
當我們來到你們面前的時候.
我們承認我們.
正是身處於各種的患難裡面.
因為你的榮耀從未受阻.
求你就賜給我們.
智慧和啟示的身體.
讓我們能夠活得更加優雅.
讓我們能夠活得更加強壯.
讓我們能夠活得更加堅強.
讓我們能夠活得更加勇敢.
智慧和啟示的靈.
讓我們在這個受限裡面.
能夠看到你緝不透的豐富.
親愛的主耶穌.

$^{481}$求你更加照著豐盛的榮耀.
藉著聖靈.
使我們心裡的力量能夠剛強起來.
求你更加帶領我們進入你的同在.
讓你長闊高深的愛.
來充滿我們這個受限的氣鳴.
並且更加可以湧流出來.
主啊,我們每一天都會面對不同的挑戰.
我們將它交在你的手裡.
願你蘊幸在我們裡面的大能.
成就一切遠超我們所想所求.
願一切榮耀都歸給你.
直到永遠.
禱告奉主耶穌基督的聖名.
阿們.
\newpage



\section{啟示錄 2:18-29}
\label{sec:032174}
\textbf{推雅推喇:成長卻容忍錯誤的教會}~~~~林穎豪傳道
\newline
\newline
連結: \href{http://www.church.com.hk/acms/content.asp?site=cdc\&op=show\&type=product\&code=032174\&layout=sermon}{\texttt{032174}} ~~~~ 語音日期: 2026-05-24
\newline
\newline
\hyperref[sec:032173]{\small{< < < PREV SERMON < < <}}
~
\hyperref[sec:index_chronic]{\small{[返順時目]}}
~
\hyperref[sec:032175]{\small{> > > NEXT SERMON > > >}}
\newline
\newline
啟示錄 2:18-29
\newline
\begin{longtable}{cl}
\hline
\hline
章節 & 經文 (和合本修訂版)\\
\hline
2:18 & \begin{tabularx}{0.7\textwidth}{X} 「你要寫信給推雅推喇教會的使者，說：『神的兒子，那位眼睛如火焰、雙腳像發亮的銅的這樣說： \end{tabularx} \\ \\ \relax
2:19 & \begin{tabularx}{0.7\textwidth}{X} 我知道你的行為：愛心、信心、勤勞、忍耐；又知道你末後所行的善事比起初所行的更多。 \end{tabularx} \\ \\ \relax
2:20 & \begin{tabularx}{0.7\textwidth}{X} 然而，有一件事我要責備你，就是你容忍那自稱是先知的婦人耶洗別教唆我的僕人，引誘他們犯淫亂，吃祭過偶像之物。 \end{tabularx} \\ \\ \relax
2:21 & \begin{tabularx}{0.7\textwidth}{X} 我曾給她悔改的機會，她卻不肯悔改她的淫行。 \end{tabularx} \\ \\ \relax
2:22 & \begin{tabularx}{0.7\textwidth}{X} 看吧，我要使她病倒在床上。那些與她犯姦淫的人若不悔改他們的行為，我也要使他們同受大患難。 \end{tabularx} \\ \\ \relax
2:23 & \begin{tabularx}{0.7\textwidth}{X} 我又要殺死她的兒女，眾教會就知道，我是那察看人肺腑心腸的，我要照你們的行為報應各人。 \end{tabularx} \\ \\ \relax
2:24 & \begin{tabularx}{0.7\textwidth}{X} 至於你們其餘的推雅推喇人，就是一切不隨從這教訓，不明白他們所謂撒但深奧之理的人，我告訴你們，我不會再把別的擔子放在你們身上。 \end{tabularx} \\ \\ \relax
2:25 & \begin{tabularx}{0.7\textwidth}{X} 你們只要持守那已經有的，直到我來。 \end{tabularx} \\ \\ \relax
2:26 & \begin{tabularx}{0.7\textwidth}{X} 那得勝又遵守我命令到底的，我要賜給他權柄制伏列國； \end{tabularx} \\ \\ \relax
2:27 & \begin{tabularx}{0.7\textwidth}{X} 他必用鐵杖管轄他們，如同打碎陶器， \end{tabularx} \\ \\ \relax
2:28 & \begin{tabularx}{0.7\textwidth}{X} 像我也從我父領受了權柄一樣。我又要把晨星賜給他。 \end{tabularx} \\ \\ \relax
2:29 & \begin{tabularx}{0.7\textwidth}{X} 凡有耳朵的都應當聽聖靈向眾教會所說的話。』」 \end{tabularx} \\ \\
[1ex]
\hline
\hline
\end{longtable}
$^{1}$在今日的靜悼儀先給大家一個處境題想像一下.
如果今天教會來了一個新朋友.
這個新朋友就說自己回了很多年教會.
不過這一兩年神就呼召他.
要來到教會宣講悔改的信息.
神就特別在66卷聖經以外.
跟他有些新的啟示 特別只是跟他講.
要他跟教會講.
一些新的啟示.
他很想來到教會跟教會的弟兄姊妹講.
如果你今天作為站在門口的招待師時.
你會怎樣呢.
會怎樣呢.
找我?.
我放假.
又或者不是招待師時.
可能你是在上班途中.
你坐火車 坐巴士.
有個人走過來跟你講.
你回哪裡教會的.
很想跟你教會的弟兄姊妹講.
神給了我很多特別的信息.
很想跟你們講.
會怎樣做呢.
這是一個很難的處境.
究竟我們要怎樣做呢.
可能我們心裡都會想.
教會都是打開門.
歡迎所有人來聚會.
我沒理由拒絕他來教會 是不是呢.
又或者我們很多都是罪人.
起初我們本身都不是很懂事.
都是回來教會聽一聽.
我們才慢慢信主.
可能教會的愛 教會的接納可以融化它呢.
可能我們都會這樣想.
懷著愛心的去到邀請.
又或者堅定的拒絕.
的確是很兩難.
同樣這也是昔日推瓦推拉教會所面對的一個掙扎.

$^{41}$所以今日都讓我們透過耶穌基督寫給教會第四封信.
這封寫給推瓦推拉教會的信.
學習主耶穌基督教導我們怎樣面對這種兩難的處境.
推瓦推拉是一個什麼地方呢.
跟我們上星期說的別加摩是差天共地.
上星期我們說別加摩是一個羅馬帝國.
一個行政的中心.
又做一些錢幣 發行錢幣.
又有很大的圖書館.
所以今日趙牧師走去圖書館 是不是呢.
不過推瓦推拉就很不同.
跟大城相比.
推瓦推拉只是一個連接著大城別加摩和薩迪的中心點.
中間一個小鎮 人口只有三萬人左右.
如果我們想像 可能給大家一個觀念.
如果今日我們由悉尼開車去到肯培拉 或者去到墨爾本.
中間我們總會經過一個小鎮叫做哥本.
大家有沒有印象 那裡有一隻很大隻的羊.
一個小鎮 哥本也是三萬人左右.
一個這樣的小鎮.
這個小小的城鎮裡面的教會.
它在聖經裡面留名.
耶穌再次在信的開頭.
向這間教會頒下他的稱讚.
如果大家有留意.
耶穌寫給其餘六間教會的信.
對比之下我們發現.
主耶穌基督在這裡推瓦推拉教會的稱讚.
可以說是無懈可擊.
耶穌其實數得到的他都讚了.
推瓦推拉教會有什麼好呢.
有好的行為 有愛心 又有信心.
有勤勞 有忍耐 還有善行.
耶穌最後再加一句.
末後所行的善事比起初所行的更加多.
即是他們的善事是越做越好 越做越多.
就好像今天我們講題所說.
推瓦推拉教會不單止是一間對內對外.
持守著信仰 見證著傳道.
都值得嘉許的教會.

$^{81}$他們更加是一間持續成長的教會.
想像一下.
一間每個星期都有人信主的教會.
小組充滿關懷互助.
信徒又很熱心在社區參與探訪 扶貧.
又有忍耐 抵得住異教徒的逼迫.
一間這樣的教會.
豈不是我們眾人都很羨慕的教會.
然而主耶穌基督出現的形象是什麼形象.
祂如何形容自己在十八節.
祂說自己是「眼目如火焰 腳像光明瞳的神之子」.
這正正是但義理書十章曾經記載過的形象.
耶穌如火焰的雙眼.
代表祂可以洞察一切隱藏的是與非.
任何的虛假 任何的罪惡.
都在這個火焰目光的監察之下無所遁形.
而祂好像光明瞳的腳.
在聖經裡代表著審判和權柄.
在推也推拉.
這間眾人都所佩服或我們都很欣賞的教會.
主耶穌竟然所展現的.
是一個如此威嚴 如此嚴肅的姿態.
並且剛才讀經的時候.
不知道大家有沒有發現.
今天我們所讀的經文有十二節那麼多.
是耶穌基督寫給七間教會的信裡.
最長的一封信.
最嚴厲的一封信.
究竟推也推拉教會.
剛才說得那麼好.
有什麼缺點呢?.
原來推也推拉的教會.
和我們之前所認識的以弗所教會.
第一封信以弗所教會.
剛剛好相反.
以弗所教會就是.
他們的教義很嚴謹.
把異端都趕出教會.
但是他們卻是將起初的愛心離棄了.
大家如果記得的話.

$^{121}$但是推也推拉教會.
他們有行為.
剛才說有愛心 有信心.
勤勞忍耐.
越做越多善事.
但是.
但是耶穌斥責他們.
容讓那自稱先知的婦人耶洗別.
教導她的僕人.
引誘他們行奸人.
食祭偶像之民.
我們就想.
這個耶洗別是誰呢?.
學者一再翻查當時希臘羅馬時期的文獻.
當中一直找不到.
當時有出現過一個女人叫做耶洗別.
而且當時耶洗別.
這個也不是猶太人會用這個名字.
去為自己的女兒命名.
因為這個名字「耶洗別」.
在猶太人的歷史中.
是一塌糊塗的.
在舊約《列王記》上記載.
耶洗別是當時以色列王阿哈的老婆.
這位皇后卻不是一位好的皇后.
這位皇后不僅在以色列.
引入了對巴力的偶像崇拜.
她也殺了很多先知.
很多無辜的人.
最終敗壞了整個以色列.
帶來極大的屬靈災難.
所以我們可以想像.
這個推瓦推拉的耶洗別.
可能不是真的有一個女人叫做耶洗別.
而是有一個具體的假先知.
只是用耶洗別代入了她的名字.
可能是一股自稱為先知的勢力.
以至教會中慢慢出現了一些追隨她的信徒.
也漸漸成為了一股不斷蔓延的思潮.
聖經經文記載.

$^{161}$耶洗別為推瓦推拉教會帶來假的教導.
引誘他們行奸淫 製偶像之物.
彷彿上星期我們說到的.
別加摩教會一模一樣.
然而如果推瓦推拉教會的信徒.
真的有做這些行奸淫.
和廟祭行淫 吃製偶像的食物.
如果他們真的有這樣做.
我們又會想.
為什麼耶穌會在信的開頭.
稱讚他們有好行為呢?.
或者將耶穌基督.
在推瓦推拉教會的責備中提到.
行奸淫 吃製偶像之物.
和人行淫.
其實我們可以將這些.
不是理解為他們真的有做這些行為.
而是理解為他們對主的不忠.
理解為對信仰上的奸淫.
可能更加妥當.
就好像舊約何西雅書一樣.
同樣都用淫婦去類比.
那些不敬拜神.
當時轉而敬拜偶像的以色列一樣.
我們就想.
究竟推瓦推拉的問題是什麼呢?.
如果不是行淫.
他們的罪是什麼呢?.
推瓦推拉教會的罪.
不在於他們自己發明了異端.
最大的問題.
也不是他們隨從假教訓.
最大的問題是他們容忍耶誓別.
這個假先知.
是容忍他留在推瓦推拉.
去教導眾人.
耶穌在二十一節說.
他已經給了耶誓別.
這個假先知悔改的機會.
但是他不肯悔改.

$^{201}$這句說話讓我們知道什麼呢?.
其實這句說話.
可能跟我們暗示.
其實這個假先知.
已經在教會有一段長時間了.
但是這段時間.
沒有任何的信徒.
甚至沒有教會的領袖做任何的舉動.
「容忍」這個字.
在希臘文裡面.
有著「容許」的意思.
「撇下」的意思.
「任憑」的意思.
即是任由他.
放他在那裡.
甚至有著一個「寬恕」的意思.
可能他們也會覺得.
我們作為教會.
我們是有愛心的.
有人進來.
我們不要論斷他.
我們讓他說吧.
推瓦推拉教會這種沒有底線的包容.
這種缺乏了真理的愛心.
使他們的包容到了一個地步.
連信仰的根基也放棄了.
豈不知「容忍」錯誤.
其實是對整個群體的殘忍.
這種毅靈性的假教導.
將教會慢慢地摧毀.
摧毀他們真實的信仰.
使這些信徒不知不覺.
與神越走越遠.
所以教會如果沒有底線.
「容忍」錯誤.
其實是在破壞整個群體.
所以耶穌接下來在二十二節.
開始宣判假先知的結局.
耶穌要這個假先知病餓在床.
聖經原本的意思是.

$^{241}$將他扔到床上.
這個昔日他用來行淫.
用來誘惑教會信徒的地方.
很諷刺地成為了.
他要受審判受懲罰的地方.
而那些追隨他.
那些不悔改的人.
二十三節形容這些人甚至是黨內.
或者如果大家有些聖經譯本看到.
寫著「兒女」兩個字.
這群人可能甚至已經.
將這個假教師尊到一個位置.
好像父母一樣去敬重,尊崇,仰慕.
這些假先知的追隨者.
都要受到審判.
而這次審判的床.
也被具象化為大患難.
甚至是二十三節所說的「殺戮」.
推亞推拉教會.
原來仍然有一群信徒沒有跟隨假先知.
不明白他們所謂「撒旦森奧之理」的人.
當然這裡所寫的「撒旦森奧之理」.
不是說來自撒旦的假教道是很深奧的道理.
只是其實聖經用來諷刺他們.
看似很深奧,很高深,很合理.
但實質上是將人導向死亡.
主耶穌說.
祂不將別的擔子放在他們身上.
只要他們繼續持守現在已經在做的.
直等到主再來的日子.
倘若推亞推拉教會.
去到堅守先前所說的.
他們做得好的事.
他們的行為,他們的愛心,信心,勤勞,忍耐.
並且他們堅定不跟隨假先知這些假教道.
他們持守著信仰的真理.
然後怎樣呢?.
耶穌要給他們另一個很不同的結局.
耶穌要賜給他們三個應許.
第一個是制服列國的權柄.

$^{281}$第二個是跟耶穌基督一同的作王.
27節說到鐵匠打碎瑤瑚的瓦器.
這些其實正是舊約詩篇第二篇的內容.
詩篇第二篇說到一個被神揀選的君王.
甚至是指向彌賽亞的王權.
耶穌說正如他從父神領受這種權柄.
推亞推拉教會要有份於這種權柄.
和耶穌一同作王的權柄.
第三個應許.
耶穌說要將神星賜給他們.
這是古今東時期一個勝利權柄的象徵.
記不記得我們剛才說.
推亞推拉是一個怎樣的地方?.
是一個像哥本的鎮.
一個彈丸之地.
但是主耶穌基督應許.
即使是一個這麼小的地方的世界會.
他將要有制服列國的權柄.
和耶穌一同作王的尊榮.
好像神星一樣被眾人所仰望.
推亞推拉這個教會.
這間小教會,小群體.
對他們來說真是難以想像.
哪有可能.
一個小教會,一個小地方.
制服列國被眾人所仰望.
但是這就是主耶穌基督.
這個大君王對於這間教會.
如果他們做得好.
給他們一個很好的獎賞.
當我們看完這一封.
主耶穌基督給推亞推拉教會的信.
然後我們將目光從推亞推拉教會.
放回在我們身上的時候.
我們就會問.
關我什麼事呢?.
即是我們要做些什麼呢?.
大家記不記得開頭我說的那個處境.
即是那個教會有人進來.
跟你說有些事很想和弟兄姊妹分享.

$^{321}$這裡講台可以上來嗎?.
你會怎樣做呢?.
再問多一次.
比特尼遇到他.
走進來跟你說.
你剛剛站在門口他跟你說.
你會怎樣做呢?.
就找我.
很抱歉 弟兄姊妹.
主耶穌在信裡面說.
總要持守.
直等到他來的是誰呢?.
主耶穌基督在啟示錄裡面.
寫給七間教會的七封信裡面.
有一句說話是不斷不斷地重複.
同樣的字眼.
在二十九節說到.
聖靈向眾教會所說的話.
凡有耳的就應當聽.
要持守信仰.
守住教會真理的界線的.
原來不是教會領袖.
不只是牧者 長老 執事.
是凡有耳的.
主耶穌說的 不是我說的.
主耶穌說教會裡面.
每個信徒的職責和本分.
教會有了真理的界線.
才可以不犯.
好像推也推拉的教會一樣的罪.
是容讓假先知的入侵滲透帶來破壞.
「容讓」這個字所表達的.
就是即使我們已經明知對方是錯的.
但我們容忍.
我們任憑他繼續做他想做的事.
這彷彿是我們今天時常聽到很多的包容 寬容.
我們這個時代.
我們將包容視為最高的道德準則.
我們很怕得罪人.
我們很怕被人貼上一些口鼓板.

$^{361}$不寬容的標籤.
於是我們在真理上慢慢地退讓.
我們在聖潔上打折扣.
我們就好像推也推拉教會一樣.
我們外面很有愛心.
我們有事工.
但在人際關係裡.
我們為了迎合那些不信主的朋友或伴侶.
我們默默地將我們聖經所說的道德準則降低.
甚至我們附和這個世界對於婚姻.
對於性別.
對於道德價值那種扭曲的看法.
我們心裡想.
我們不要禁古板.
做出來的樣子不要這麼抗拒其他人.
我們用愛去包容他們.
我們用愛心作為包裝.
但實質上推動我們這樣做.
為什麼要推動我們要包容.
其實不是因為我們真的有愛心.
而是我們內心的懦弱和恐懼.
我們害怕衝突.
弟兄姊妹不要覺得我在說你們很軟弱.
很懦弱.
其實因為我自己也非常明白.
非常心痛感受.
因為其實我也是一個很內向的人.
我也是一個很重關係的人.
我也是一個本性上很怕衝突的人.
但耶穌今天呼召我.
安排我今天做講員.
耶穌先用這段經文.
這篇訊息提醒我.
祂叫我不要走不要逃避.
不要害怕得罪人.
不要害怕要帶來衝突.
因為你勇敢站在這裡是為著真理.
是為著主的道.
是為了穩固教會基督的身體.
沒錯有時我們本能會很害怕.

$^{401}$但讓我們學習.
為了教會.
為了真理的界線.
為了主耶穌基督.
是我們勇敢的站立.
拒絕錯誤的教導和價值觀.
不過說到這裡.
另一個問題又浮現了.
究竟什麼是對的.
什麼是不對的呢?.
有問我.
其實就像上個月我們說的.
以弗所教會的那封信一樣.
其實已經說過了.
好像以弗所教會一樣.
就是用聖經去檢驗不同的教導.
就好像今天你來到教會聽到.
其實今天站在台上的講員.
他所說的話.
跟你一直讀聖經所領受所明白的.
會不會不同呢?.
又或者會不會你曾經讀過的經文.
跟今天講員所說的完全相反呢?.
所以弟兄姊妹.
教會真理的界線在哪裡?.
如何分辨教導的對與錯.
究竟站在台上是可信的牧者.
還是假教師呢?.
其實這個責任在於你們自己.
究竟你有多認識聖經?.
究竟你有多認識神的話語?.
又或者究竟你有多花時間去了解.
去分辨這些教導呢?.
除了讀經.
另一個問題就是.
當我們每天接收很多不同的資訊.
很多不同的思想價值觀.
我們會不會接收的時候.
聽到的時候都多想一些.
這是否符合神的心意呢?.

$^{441}$有時我們來到教會.
見到教會上的講員.
我們回到教會.
講完講什麼.
我們就聽.
這是好的.
我也知道.
說宣的大家很尊重我們牧者.
亦感受到大家很愛我.
不過我想說的是.
教會的講員.
所說的話是否一定是正確的?.
是否站在講台.
他說的話就一定全然反映神的心意呢?.
我也是人.
我可能有時也會說錯話.
又或者.
當今天我們在網上.
社交媒體.
見到很多很多不同的貼文.
有些貼文可能寫明自己是基督教的內容.
又或者有些貼文可能出自基督教的機構.
但我們看的時候.
我們會否就像一塊海綿.
吸水一樣.
全部都吸進去.
是否一定是正確的?.
還是我們也要按照.
神賜給我們的智慧.
悟性.
神賜給我們的腦袋.
去過濾一下?.
今天.
今天我們的腦袋.
裝的是什麼呢?.
我們想得最多的是什麼呢?.
現在是金融週末.
現在我們買什麼呢?.
又或者我們下星期.
去哪裡玩呢?去哪裡吃呢?.

$^{481}$究竟我們今天的腦袋是裝什麼的呢?.
今天聖經提醒你和我.
我們要多些用.
神創造給我們的腦袋.
承載神自己的說話.
分辨我們每天所聽到.
所看到的資訊.
以致我們每一個人.
不單只是牧者.
不單只是領袖.
我們每一個人.
都持守我們信仰的界線.
耶穌基督給推也推拉教會這封信.
是對今天教會一個很嚴厲.
卻充滿恩典的提醒.
我們固然要追求愛心,善行.
要追求教會的成長.
但這一切絕對不能建立在.
一個妥協真理的基礎之上.
真正的愛不是懦弱.
和恐懼的縱容和欽鵬.
聖經說真正的愛是.
「不喜歡不義,只喜歡真理」.
真正的包容不是對罪惡.
對假教導視而不見.
而是勇敢的在愛裡說誠實話.
作為基督的身體.
我們每個信徒對守護真理.
持守信仰的界線都責無旁貸.
求主幫助我們.
不但在行為,愛心上不斷增長.
也在真道上用神的話語.
裝備我們的腦袋.
作一個有分辨能力的信徒.
可能有時要堅守真理.
站穩界線都會害怕,會痛.
甚至可能要付一些代價.
可能我們會被人孤立.
被人誤解為不近人情,思想落後.
但是讓我們謹記.

$^{521}$耶穌基督的應許.
即使多微小也好.
卻是有制服列國.
與基督一同作王得勝.
所以讓我們一起堅守.
主所託付真理的界線.
直等到他們的日子.
我們同心抬頭祈禱.
寫給推亞推拉教會的信.
再次與我們每一個人說話.
你說「凡有意的就應當聽」.
主也承認我們有時會因為軟弱.
因為害怕得罪人.
因為害怕衝突.
所以我們選擇妥協.
我們用廉價的包容.
去掩蓋對真理的堅持.
所以求主赦免我們.
也求主如火焰的眼目.
監察我們的內心.
光照我們裡面的隱惡.
求聖靈賜給我們.
屬靈的智慧和勇氣.
讓我們不但追求外在的事工和熱心.
我們更願意花時間在你的話語裡.
深深地紮根.
求主幫助我們每一個頂子妹.
不單不只是依賴牧者.
而是我們自己也擔當起.
守護教會真理界線的責任.
是用你賜給我們的悟性.
去明辨每一天我們所聽到.
所看到的不同的資訊.
不同的教導.
主耶穌賜給我們有力量.
去堅守你的真道.
勇敢地拒絕任何似是而非.
錯誤的價值觀.
但願我們這個屬你的群體.
能夠在真理裡站立得穩.

$^{561}$成為這個世代裡發光的星辰.
直到你再來的日子.
禱告是奉主耶穌基督得勝的名字祈求.
拜拜.
\newpage



\section{士師記 4:1-24}
\label{sec:032175}
\textbf{士師記析讀 (6) 謙招損,滿受益}~~~~梁永善牧師
\newline
\newline
連結: \href{http://www.church.com.hk/acms/content.asp?site=cdc\&op=show\&type=product\&code=032175\&layout=sermon}{\texttt{032175}} ~~~~ 語音日期: 2026-05-20
\newline
\newline
\hyperref[sec:032174]{\small{< < < PREV SERMON < < <}}
~
\hyperref[sec:index_chronic]{\small{[返順時目]}}
~
\hyperref[sec:032176]{\small{> > > NEXT SERMON > > >}}
\newline
\newline
士師記 4:1-24
\newline
\begin{longtable}{cl}
\hline
\hline
章節 & 經文 (和合本修訂版)\\
\hline
4:1 & \begin{tabularx}{0.7\textwidth}{X} 以笏死後，以色列人又行耶和華眼中看為惡的事。 \end{tabularx} \\ \\ \relax
4:2 & \begin{tabularx}{0.7\textwidth}{X} 耶和華把他們交給在夏瑣作王的迦南王耶賓手中；他的將軍是西西拉，住在夏羅設‧哈歌印。 \end{tabularx} \\ \\ \relax
4:3 & \begin{tabularx}{0.7\textwidth}{X} 以色列人呼求耶和華，因為耶賓王有鐵的戰車九百輛，並且殘酷欺壓以色列人二十年。 \end{tabularx} \\ \\ \relax
4:4 & \begin{tabularx}{0.7\textwidth}{X} 有一位女先知底波拉，是拉比多的妻子，當時作以色列的士師。 \end{tabularx} \\ \\ \relax
4:5 & \begin{tabularx}{0.7\textwidth}{X} 她住在以法蓮山區拉瑪和伯特利的中間，在底波拉的棕樹下。以色列人都上到她那裡去聽審判。 \end{tabularx} \\ \\ \relax
4:6 & \begin{tabularx}{0.7\textwidth}{X} 她派人從拿弗他利的基低斯把亞比挪菴的兒子巴拉召來，對他說：「耶和華－以色列的神吩咐你：『你要率領一萬拿弗他利人和西布倫人上他泊山去。 \end{tabularx} \\ \\ \relax
4:7 & \begin{tabularx}{0.7\textwidth}{X} 我必使耶賓的將軍西西拉率領他的戰車和全軍往基順河，到你那裡去，我必把他交在你手中。』」 \end{tabularx} \\ \\ \relax
4:8 & \begin{tabularx}{0.7\textwidth}{X} 巴拉對她說：「你若同我去，我就去；你若不同我去，我就不去。」 \end{tabularx} \\ \\ \relax
4:9 & \begin{tabularx}{0.7\textwidth}{X} 底波拉說：「我一定會與你同去，然而你在所行的路上必得不著榮耀，因為耶和華要把西西拉交給一個婦人的手裡。」於是底波拉起來，與巴拉一同往基低斯去了。 \end{tabularx} \\ \\ \relax
4:10 & \begin{tabularx}{0.7\textwidth}{X} 巴拉召集西布倫人和拿弗他利人到基低斯，跟他上去的有一萬人。底波拉也同他上去。 \end{tabularx} \\ \\ \relax
4:11 & \begin{tabularx}{0.7\textwidth}{X} 摩西岳父何巴的後裔，基尼人希百離開了基尼族，到靠近基低斯的撒拿音橡樹旁支搭帳棚。 \end{tabularx} \\ \\ \relax
4:12 & \begin{tabularx}{0.7\textwidth}{X} 有人告訴西西拉：「亞比挪菴的兒子巴拉已經上了他泊山。」 \end{tabularx} \\ \\ \relax
4:13 & \begin{tabularx}{0.7\textwidth}{X} 西西拉就召集所有的鐵戰車九百輛和隨從的全軍，從夏羅設‧哈歌印出來，到了基順河。 \end{tabularx} \\ \\ \relax
4:14 & \begin{tabularx}{0.7\textwidth}{X} 底波拉對巴拉說：「起來，今日就是耶和華把西西拉交在你手中的日子。耶和華豈不在你前面行嗎?」於是巴拉下了他泊山，跟隨他的有一萬人。 \end{tabularx} \\ \\ \relax
4:15 & \begin{tabularx}{0.7\textwidth}{X} 耶和華使西西拉和他一切的戰車，以及全軍潰亂，在巴拉面前倒在刀下。西西拉下了車，徒步逃跑。 \end{tabularx} \\ \\ \relax
4:16 & \begin{tabularx}{0.7\textwidth}{X} 巴拉追趕戰車、軍隊，直到夏羅設‧哈歌印。西西拉的全軍都倒在刀下，一個也沒有留下。 \end{tabularx} \\ \\ \relax
4:17 & \begin{tabularx}{0.7\textwidth}{X} 只有西西拉徒步逃跑到基尼人希百之妻雅億的帳棚，因為夏瑣王耶賓與基尼人的希百家和平共處。 \end{tabularx} \\ \\ \relax
4:18 & \begin{tabularx}{0.7\textwidth}{X} 雅億出來迎接西西拉，對他說：「請我主進來，進到我這裡來，不要怕。」西西拉就進了她的帳棚，雅億用被子將他蓋住。 \end{tabularx} \\ \\ \relax
4:19 & \begin{tabularx}{0.7\textwidth}{X} 西西拉對雅億說：「我渴了，求你給我一點水喝。」雅億就打開裝奶的皮袋，給他喝，再把他蓋住。 \end{tabularx} \\ \\ \relax
4:20 & \begin{tabularx}{0.7\textwidth}{X} 西西拉對雅億說：「請你站在帳棚門口，若有人來問你說：『有人在這裡嗎？』你就說：『沒有。』」 \end{tabularx} \\ \\ \relax
4:21 & \begin{tabularx}{0.7\textwidth}{X} 西西拉疲乏沉睡了。希百的妻雅億取了帳棚的橛子，手拿著錘子，靜悄悄地到他那裡，將橛子從他的太陽穴釘進去，直釘到地裡。西西拉就死了。 \end{tabularx} \\ \\ \relax
4:22 & \begin{tabularx}{0.7\textwidth}{X} 看哪，巴拉追趕西西拉，雅億出來迎接他，對他說：「來，我給你看你要找的人。」他就進入帳棚，看哪，西西拉已經倒在地上死了，橛子還在他的太陽穴中。 \end{tabularx} \\ \\ \relax
4:23 & \begin{tabularx}{0.7\textwidth}{X} 那日，神在以色列人面前制伏了迦南王耶賓。 \end{tabularx} \\ \\ \relax
4:24 & \begin{tabularx}{0.7\textwidth}{X} 從此，以色列人的手對迦南王耶賓越來越強硬，直到將迦南王耶賓剪除。 \end{tabularx} \\ \\
[1ex]
\hline
\hline
\end{longtable}
$^{1}$好 今天我們會看到《時事記》.
我們查到第六次了.
這次的題目是較為特別一點的.
就叫做「謙招選滿受益」.
中國人有一句話怎麼說呢?.
就是「滿招選謙受益」.
意思就是如果你自滿,自傲,太驕傲的時候.
你只會招致虧損.
其他人就不會滿意你.
甚至看輕你.
而一個謙虛的人.
他就會有益處了.
因為他樂於聽別人的指引.
他願意向別人學習.
更加不會逞強.
對他來說是一個祝福.
所以也可以套用在我們的生活和信仰上.
如果我們有看過.
若書當他面對耶利哥城很難打贏的時候.
他去倚靠神.
結果可以打到耶利哥城全部倒塌了.
但後來他自以為自己很厲害.
去打艾城.
他說那個人很少.
派二,三千人上去就行了.
結果就在艾城那裡失敗了.
而我自己也有這個經歷.
在我這麼多年超過五十年的事奉.
有幾次我自己也面對一些失敗或者打擊.
都是當我以為自己很厲害.
飄飄然的時候.
這樣就出問題了.
所以真的學到.
原來我們需要謙卑倚靠神.
神是敵黨驕傲的人.
但是次因比謙卑的人.
但今天我的題目剛剛相反.
就是謙招選.
但滿就受益.
當然.

$^{41}$我這個滿和謙是有不同的解讀的.
透過《時詩記》.
我們就可以看真理了.
我們這次先看《時詩記》.
第四章第一到第九節.
如果你記得.
之前有三詩會.
有三個時詩.
但是這三個時詩都算是好的.
特別是《訛陀列》.
但是講的很少.
但是我們最熟悉的就是《底波拉》.
我們一起來看第四章第一到第九節.
好 先請Yvonne替我們讀出來.
以色列人就呼求耶和華.
有一位女星之名叫底波拉.
是拉比多的妻.
當時在以色列的史詩.
她住在以法連山地拉瑪和伯特利中間.
在底波拉的中樹下.
以色列人都上她那裡去聽判斷.
她打發人從拿佛塔利的基底詩.
將亞比·羅恩的兒子巴拉·趙鳥蘭對她說.
耶和華以色列的神吩咐她說.
「你率領一萬拿佛塔利和西布倫人上塔珀山去.
我必使耶賓的將軍西西拉率領他的車輛和全軍往基遜河.
到你那裡去.
我必將他交在你手中」.
巴拉說:「你若同我去我就去.
你若不同我去我就不去」.
底波拉說:「我必與你同去.
只是你在所行的路上得不著榮耀.
因為耶和華要將西西拉交在一個婦人手裡.
要使底波拉,比萊和巴拉一同找基底詩去了」.
多謝二位默默地讀出這九次經文.
我們先看一看背景.
這次戰爭影響著以色列人的經濟和合一.
為什麼呢?我們看一看地圖就會知道.
因為以色列平原分隔了古以色列的中部.
和北部的各支派.

$^{81}$而當時迦南人和菲律賓人控制了這個地方.
他們如果要走這些交通大道的時候.
就會遭到搶劫和婦女被姦污.
我們下星期會看底波拉的詩歌時.
就會講到原來在這些大道上.
他們是不敢帶他們行走的.
為了避免這個危機.
他們就要繞道而行.
所以他們在經濟和政治上陷入了困境.
所以如果他們得勝的話就可以解決這個問題.
而這次的戰爭.
也顯示了以色列各支派在政治上.
比以前更加有合一的意識.
為什麼呢?我們看一看.
在《史詩紀五章》十四到第十八節裡面.
我們下次會繼續看下去.
有一班人他們是願意來參與的.
我們看一看《史詩紀五章》十四到第十八節.
你有沒有留意到特別紅色字的part1呢?.
就列出了有以法連,有辯雅敏,有西布倫,.
有以撒迦,還有是拿伐他利.
「拿」字沒有開在這裡.
這時候你就會看到他們是有一個合一.
但是我想挑戰大家.
以色列有十二支派.
另外一些支派是沒有出席的.
是哪些支派呢?.
為什麼他們沒有出席呢?.
下次我們會看.
這時候大家看到政治是很現實的.
如果是對自己經濟利益有關的時候.
什麼的道義或者藥都可以棄如敝履.
但無論如何我們看到原來政治就是這樣.
現實就是這樣.
你有利益就可以在一起.
沒有的話就算了.
所以你留意到.
以前可以是朋友.
現在可以是敵人.
以色列人你留意到.

$^{121}$這次有幾個支派聯合成為一個聯盟.
因為對大家的經濟都有幫助.
還有就是兄弟情誼的手足.
他們的對手就是夏所王.
就是在夏所作王的加勒蒙·耶賓.
和他的將軍西西賴.
這次敵人跟過去的敵人有些不同.
以前的敵人是居住在範圍以外的人.
但是現在夏所王卻是住在迦南境內.
所以造成更大的困擾.
外敵你容易去抵禦.
但是住在裡面的時候.
內患就難擋了.
他擁有戰車900輛.
剛才我們看到了.
而加勒蒙是一個聯盟的首領.
據考古學的尼達去發現.
夏所就是古代巴勒斯坦地小有的大城.
發掘出來的城包括了.
停放戰車的地方有多少公頃呢?.
是80公頃.
很厲害.
80公頃.
大家都知道.
所羅門是猶大國.
或者整個以色列國最輝煌的時間.
我們去旅行一定會去看米吉多這個地方.
就是所羅門去養馬的地方.
但是米吉多所羅門的防衛城.
養馬的地方有多少公頃呢?.
是8公頃.
大家看到嗎?.
夏所王有80公頃養馬.
但是所羅門這麼繁榮的時候只有8公頃.
可以想到大家的勢力很懸殊.
當時戰車是最厲害的武器.
而且有900輛那麼多.
所以它的勢力是非常龐大的.
相對以色列真是手無寸鐵.
他們沒有做鐵器的兵器.

$^{161}$我們看看13章第19到20節.
這是撒母耳記上的.
那時以色列全地沒有一個鐵匠.
因為菲律斯人說.
恐怕希伯來人製造刀槍.
以色列人要磨鋤,泥,鼓,鏟.
就下到菲律斯人那裡去磨.
所以你留意到他們原來沒有鐵器.
你知道中國曾經有一段時間.
譬如奧運或者一些特別元首來的時候.
是不准賣刀的.
賽刀也不准賣.
就怕有事發生.
原來當時菲律斯人就收了所有鐵器.
所以他們就沒有兵器來對抗.
但是對手有900輛戰車.
他們被欺壓20年.
他們向神呼求.
神就興起士師來拯救他們.
但是我們不要單單看一節經文.
其實這些是以色列人自招.
記不記得《士師記》有一個模式.
就是他進入迦南地之後.
和外邦人來有殘防.
好好的敬拜神.
甚至有屢勸不改.
於是神就容讓外邦人來攻擊他們.
壓制他們.
當他們受苦醒覺的時候.
向神呼求.
求神來憐憫他們.
於是神就會興起一些士師.
帶領他們作出反擊.
成功之後他們又享太平了.
很可惜在安寧日子裡面.
他們又固態服盟.
又再離開神.
又再重複這些方法.
神就用一樣的再重複.
就是向外敵帶進去攻擊他們.

$^{201}$所以是一個很無奈的循環.
我們看看《士師記》第四章第一到第二節.
已忽死後.
以色列人又行耶和華眼中看為惡的事.
耶和華就拜他們賦予.
在加索作王的迦南王耶賓手中.
因為他們士師死了之後已忽.
他們又行耶和華眼中看為惡的事.
這個循環會不會.
我已經多次提到.
當我們基督徒的寫照.
當我們有事的時候.
我們會求告神.
求神幫助.
並且許下一些承諾.
但是當神來幫助我們渡過難關之後.
我們沒有機會將這些承諾拋之於路後.
愛去犯錯.
結果神無奈又要出一個愛的懲罰.
所以當我們回轉.
很希望今天我們每一個參與者.
甚至在YouTube看《茶經》的人.
都不要做這些事.
不要重蹈子民所犯的錯.
好了.
我們今次會看看《女權當道》.
當子民呼求神的時候.
神就興起底波拉.
我們要留意下星期我們會特別講《女權當道》.
底波拉特別被標籤為什麼呢?.
是拉比多的妻子.
《聖經》稱祂為什麼呢?.
稱祂為先知.
是代表神彌大人去傳訊息教導子民的.
但都是作事司.
就是要去審判的.
我們看看《事司》第四章第五節.
「她住在以法連三地拉瑪和伯特利中間.
在底波拉的中樹下.
以色列人都上她那裡去聽判斷.

$^{241}$她身兼至少三職.
甚至四重身份.
她是人的妻子.
她是先知.
她是事司.
還可以是人的母親.
在那時男權當道的時代裡面.
顯然她是一個很突出的人物.
其實神不會因為那個人的性別.
而做出選擇.
一定是男的.
不會選女的.
今天最重要的是什麼呢?.
留意我強調這句.
是你的心志和願不願意為神擺上.
如果你有這份心志願意為神擺上.
神是不會理你的性別的.
也會有理由縱容你.
作為妻子.
她要照顧自己的家庭.
如果有孩子的時候.
她當然要照顧自己的孩子.
更要承擔雙重的職份.
就是事司和先知.
所以委實是不容易的.
但是深信神必定會加力.
但是不得不承認.
現在很多的教會.
他們的領導人都是什麼呢?.
都是女多男少.
都是女士們.
她又要照顧家庭.
又要來打理家務.
但都做很多教會的工作.
這就是我去非制的時候.
非制教會的執事.
有沒有留意到.
右邊那個是牧師.
1,2,3,4,5.
五女一男.

$^{281}$即是五個女士.
一個男士是執事.
他們都不容易的.
因為我們探訪他們.
他們自己又要打理店鋪.
又要照顧孩子和丈夫.
還要出席教會祈禱會事奉.
所以很佩服他們.
反而有時候我們男士是有些失色的.
看到底波拉去呼召領導.
底波拉接受到神的巧遇.
就吩咐巴拉來解救以色列人.
我們看看《士師記》第四章第六到第七節.
「他打發人從拿弗他利的基底斯.
將亞比羅恩的兒子巴拉召了來.
對他說:.
『也護話以色列的神吩咐你說.
你率領一萬拿弗他利和西布倫人上他撲山去.
我必使耶賓的將軍西西拉.
率領他的車輛和全軍往基遜河到你那裡去.
我必將他交在你手中』」.
究竟巴拉是甚麼人呢?.
聖經沒有詳細記載.
《士師記》也沒有直接稱巴拉為「事司」.
但以色列人的傳統.
就將他列為「事司」的其中一位.
但是巴拉如果我們這樣推測.
可能是一位出身很顯赫的知名人士.
《士師記》作者特別提到他父親的名字.
巴拉是來自拿弗他利的基底斯.
這個地方和古代的通商要道米吉多.
只是數英里的距離.
所以他很可能擁有了非常豐厚的財富和名望.
因為他能夠號召一萬人來和他同行.
不要忘記這一萬人不是和他去旅行.
而是去打仗.
可能要付上生命的危險.
但都要有人來跟從他的時候.
所以我們不可以輕忽巴拉.
他也一定有他的功勞在當中.

$^{321}$而事實上神都可以用人的資源來完成他的計劃.
正如我說的.
較為有錢的人他們可以有較多的奉獻.
來支持這樣的事工.
就像上次我說的.
當年我們名人堂要買地方的時候.
我們只有60萬.
但是要買600萬的物業.
但是我一個月就已經籌到100萬出來.
而其中一個家庭一次過奉獻50萬.
因為他是廠家來的.
他很愛主的.
那時候他就願意奉獻50萬.
因為他有錢.
有神的祝福.
他就去做更多的工作.
去支持更多的福音工作.
所以這是很重要的.
所以巴拉應該就是拿伐他尼的統治階級.
因此底波拉就來帶入他特別的蒙神指示.
去呼召他出來.
當然底波拉一定有神的指示.
才會選擇選擇巴拉.
巴拉原本的名字是什麼呢?.
意思就是閃電 很快的.
快速的意思.
但是有時候名字是剛剛相反的.
正如我說過.
我家裡有三個兒子.
我爸爸我叫永善.
我是最兇的那個.
我大弟弟叫永安.
但是他的健康較為弱一點.
永明不要說了.
某個程度讀書不是最厲害的那個.
所以你留意到有時候名字是剛剛相反的.
巴拉原本的名字叫閃電.
但是他很多東西是猶豫不決的.
在整個行動當中.
你會更加留意到.

$^{361}$他往往是遲人一步的.
我沒有怎麼看如來神掌.
如來神掌的高佬傳.
就叫做東島祥尼.
東島祥尼每次都是慢人一步的.
遲了一步被一個婦人殺了.
他就失去了他的榮耀.
但是在原文裡面.
你留意到是很特別的.
它說到.
豈不知耶和華以色列的神吩咐你說.
你留意到在中文聖經裡面.
沒有這個字翻譯出來.
但是其實.
底波拉跟他說你不是知道耶和華以色列神吩咐你要這樣做嗎?.
即是說巴拉絕對不是第一次聽到底波拉.
叫他去打仗.
他就去.
你想一想.
如果他第一次聽到底波拉叫他打仗.
他會不會即時有回應呢?.
要找命搏的.
還要呼召這麼多人跟從他.
絕對不容易.
所以巴拉應該一早已經有神的呼召.
不過他猶豫不決.
後來底波拉來的時候.
就引證了他這個呼召.
正如當年我很清楚神呼召我要做傳道人的時候.
其實我有很多很多的呼召.
不是一次.
但是我都拒絕.
後來我要求神給我一些特別的經歷.
然後真的有這些特別的經歷的時候.
特別的經歷是甚麼呢?.
好像很簡單.
你給我看傳福音不就很重要了嗎?.
但接著那一年我第一次去宣道園.
做福音營的導師.
接觸一些東頭村的小朋友.

$^{401}$我沒有去過東頭村的七層大廈.
也在那一年我第一次回中國大陸.
我是甚麼呢?.
我是在香港出生的.
在香港長大的.
我18,19歲才第一次回中國大陸.
就是文革的後期.
我看到中國人民需要福音的時候.
這些就引證了我的呼召.
最後我決定去獻身.
所以我絕對相信巴拉園在原文裡面.
豈不知這個字很重要.
即是說神已經向他作出這個呼召.
只是他沒有回應.
所以神息差派雷先知底波拉再一次提醒他.
即是神不是不給他機會.
縱使他沒有即時回應.
但神仍然不放棄.
而藉著底波拉給他明確的呼召.
並且怎樣呢?.
給他有勝利的預告.
因為很清楚地說.
我已經將他交在你手裡.
這個呼籲裡面我們看到兩方面.
神呼召巴拉帶領以色列人.
特別是拿伐塔利和西布倫人.
上塔珀山預備作戰.
塔珀山去以色列的地方我們都去過.
不是很高.
是在耶斯列平原一帶.
耶斯列平原是甚麼呢?.
是一塊大到一望無際的平原.
我第一次看到耶斯列平原的時候.
心裡很震動.
那個年代還未有手機.
我拍照要這樣拍.
然後將相片放出來.
然後黏在一起.
因為是很遼闊的.
這個地方就是啟示錄所說的.

$^{441}$哈米吉多頓大戰的地方.
當你真的踏上聖地的時候.
你就看到原來有二萬萬的紅軍.
去到這個地方作戰.
真的可以容納這麼多人.
拿破侖也在這裡打過仗.
所以這個平原有戰車是非常非常有利的.
因為沒有甚麼樹木.
在平原裡面.
敵人沒有地方可以躲藏.
在山上有茫茫的樹林.
他們就難以發揮所長.
一個能征慣戰的耶賓王和手下猛將西沙.
一定會知道那個平原是對地有利的.
是一個好開戰的地方.
但是神清楚的宣告.
祂會有使西西拉率領全軍.
去到基遜河.
這個河水大部分都是乾涸的.
是一條小溪.
只有冬天雨季的時候.
才會有水流下來.
但是有水的時候.
就會令到地方成為沼澤.
泥濘滿佈.
敵軍就會出問題了.
但是你有沒有留意到.
神很清楚的說.
「我必將他交在你手中」.
即是說.
縱使敵軍如何武器精良.
人數眾多.
但是這場仗還沒有打.
勝負已經定了.
我們看看.
謙招選.
這個歷史事件除了看到神奇的作為之外.
因為祂使雨水驟然降下.
使洪壯的軍隊.
潰敗在神的子民眼前.

$^{481}$但是得勝的.
當然是靠著神的帶領之外.
另外一個重點是什麼呢?.
就是很重要.
如果我們應該承擔的責任.
但是我們不去承擔.
我們退縮.
甚至我們棄權.
結果損失的是我們自己.
當底波拉呼召巴拉領軍攻擊敵人的時候.
很清楚地告訴你.
神已經逝去得勝了.
甚至指出這場仗是如何打的.
先教他們上塔珀山.
好了.
但是巴拉竟然提出一個要求.
第四章第八節.
巴拉說.
你若同我去 我就去 你若不同我去 我就不去.
當我讀到這段的時候.
我就想起一個小朋友.
叫他做事.
你和我一起去 我就去.
你不和我去 我就不去.
這樣很離譜.
作為一個被揀選的領袖.
應該有勇氣去面對的.
但是為什麼他要拒絕呢?.
之後領軍.
但是他有一個附帶條件.
就是要底波拉和他一起去.
為什麼會這樣?.
是不是因為他謙卑呢?.
不是.
其實這是假謙卑.
真正的原因.
他是膽怯.
他沒有充足的信心.
雖然可以召喚一萬人.
其實也為數不少.

$^{521}$聖經沒有記載西西拉有多少軍隊.
但是相信絕對不會少.
因為戰車也有900輛.
所以他們所持的武器.
相差太遠了.
底波拉的宣告.
應該是真實的.
但是他感到很矛盾.
所以他就說不行.
一定要和我去.
如果底波拉在那個時候.
有些保留.
我不去 他就害怕了.
但是如果底波拉肯去的話.
他就覺得有保證了.
這個呼召從他口出.
所以他去的時候一定會很穩妥.
不會只是你叫我去.
我又不去做.
所以一定要和我去.
信心不足.
尋求確實的證據.
或者憑據.
絕對不是一個問題.
正如稍後.
應該是下兩次.
我們會看機電.
當神呼召機電的時候.
他也要憑據的.
當神就使火.
從盆石裡面出來.
燒了一些肉和無膠餅.
所以原來第一次會有火燒了的製肉.
不是以利亞在加密山.
是機電.
我們稍後會看.
還有.
機電還要繼續.
陽無濕和乾.
在《時史記》第六章十一到四十節.

$^{561}$我們稍後會看這件事.
但是你留意到機電將焦點.
不是放在人的身上.
要放在神的身上.
神要將你的指示清晰讓我知道.
還有.
如果你不和我去.
我就不去.
這句說話大家熟不熟呢?.
是否似曾出現過呢?.
如果一直一直一直.
在過去這幾年裡.
有參加我的查經班.
在一兩年前.
我會查舊約的祈禱.
那時候.
你記不記得摩西也有這樣的禱告.
摩西再上山.
為子民尼達去代求.
求神不要滅他們.
為他們做金牛毒的罪.
求神赦免.
神就答應不即時刑罰子民.
但卻說.
他不會再親自帶領這班百姓.
即入迦南地.
改由師姐來帶領.
《出禮及記》第33章第二第三兩節.
我要看見師姐在你面前.
面出迦南人.
阿摩利人 赫人 比利西人.
希末人 耶布斯人.
領你到那流乃與民之地.
我自己不和你們上去.
因為你們是眼著頸看的百姓.
恐怕我在路上把你們滅絕.
神後終於說.
我不會和你們上去.
我會找師姐帶你們上去.
其實是一樣的.

$^{601}$免得我一怒起來你們又再犯錯.
我就滅你們.
但摩西竟然有一個堅持.
就是你不去我也不去.
我們看看摩西的呼求.
15,16兩節.
摩西說:你若不親自護我同去.
就不要把我們從這裡領上去.
人在何事上得而知道我和你的百姓.
在你眼前蒙恩.
若是因你願我們同去.
使我和你的百姓與地上的萬民有分別嗎?.
摩西和子民有很大的分別.
子民不想神直接和他們說話.
不如摩西.
神和你說甚麼你告訴我吧.
然後神說不和他們去.
由師姐帶他們去.
我想以色列人覺得沒問題的.
但摩西說不行.
如果你不去我就不去了.
他有這個堅持.
結果是甚麼呢?.
第17節.
耶和華對摩西說:你者所求的我也要行.
因為你在我眼前蒙了恩.
並且我按你的名認識你.
是不是很寶貴?.
結果摩西這個堅持.
神是答應了.
所以神就繼續和他們同去.
來到基甸.
摩西,巴拉三者都是求憑據.
但分別在哪裡呢?.
基甸和摩西是要神自己給他們保證.
但巴拉雖然受了神的呼召.
和底波拉再次確認這個呼召.
並且作出保證.
但他們卻要用人來保證.
要底波拉同往.

$^{641}$他才會覺得有信心.
這時候才願意去做.
其實我們稍後會看到.
在戰爭的過程裡.
我們看不到底波拉有甚麼參與.
我們看不到底波拉做了甚麼.
結果巴拉失去了他應有的榮耀.
神很清楚宣告敵人的將領.
不是敗在他的手下.
而是死在一個婦人的手上.
這個婦人就會得到榮耀和稱頌.
第四章第九節.
「底波拉說:我必與你同去.
此是你在所行的路上得不著榮耀.
因為耶和華要將西西拉交在一個婦人手裡」.
神很清楚透過底波拉說.
你這樣 你得不到榮耀.
西西拉這個將軍將會死在一個女人手上.
這個宣告可以說是對巴拉的一個羞辱.
就是因為他的自謙和膽怯.
而招致損失.
他損失了他應有的榮耀.
所以謙招損.
這個給我一個很大的感慨.
今天我們要得到的信心的保證是甚麼呢?.
是我們本身擁有的人力,財力,物力.
是不是?.
還是一些有名望的人.
來與我們在一起.
就做保證.
你看有很多人和那些政要.
拍完照放這麼大的一張.
放在自己的辦公室那裡.
然後就說我們的Facebook又會怎樣?.
我明天都會去倫敦見一個在香港很出名的人.
他剛剛來了英國.
明天他約了我吃飯.
但我就不會說我跟誰誰吃過飯.
他要拿這些照片拍照之後我都會和他拍照.
但我絕對不會放上任何那些大眾媒體那裡.

$^{681}$因為始終都是一個人.
但有些人就要拿這些人去抬高自己.
這個是很敗獄.
巴拉要的是底波拉.
但昔日摩西要的是神而不是其他人.
結果摩西的成就和巴拉的成就.
是不是有很大的差距呢?.
我想問大家心裡有答案嗎?.
今天未查經之前你對巴拉熟不熟悉?.
應該不熟悉.
但你一定會對摩西有這個熟悉 是不是?.
巴拉因為對神選照的懷疑.
結果他的榮耀就轉移到一個婦女身上.
在一個愛邦女子的身上.
是不是很出人意表呢?.
神的計劃其實往往都是出人意表的.
西沙拉的死應該是巴拉的榮耀.
但現在卻成為他的損失.
一切就是因為什麼呢?.
他太過自謙.
記住.
對我們現今有什麼提醒呢?.
今天神的事工是需要人跟祂配合的.
是要神人共憤的一起努力的.
神會向屬於祂的人發出呼召.
甚至可能是很多次的呼召.
但人因為懼怕畏縮而棄權的話.
神是不會勉強的.
但神的工作是絕對不會因為那個人的退縮而受影響的.
你不做而已.
神就可以叫另外一個人去做.
正如大家很熟悉的《耳私帖記》.
我們看看《耳私帖記》第四章第十四節.
此事李若閉口不言.
猶大人必從別處得解脫.
蒙拯救.
你和你父家必自滅亡.
焉知你得了皇后的位份.
不是為現今的機會嗎?.
第一句很重要.

$^{721}$就是末底改叫以斯帖向王尼達去求.
免他們猶太人成族滅亡.
但是以斯帖就很害怕.
不知道行不行.
因為王母召我很久了.
如果王母召我入宮的時候.
我進去隨時會被斬頭的.
我對她說.
如果你這個時候閉口不言.
我們猶太人都一樣會得救.
不過不是由你去幫.
這個時候神都會興起另外一些人.
這個很重要.
所以今天我們要清晰.
如果你不向你的家人.
不向你的朋友,同事傳福音.
難道他們永遠都沒有機會聽到福音嗎?.
絕對不會.
我很相信神會透過另外一些人.
傳福音甚至帶領他們去信主.
所以今天我們沒有承擔這個責任.
那些帶領去信主的人.
他們會得到神的獎賞.
但是我們沒有承擔這個責任.
我們得不到神的祝福.
所以求神幫助我們.
如果是神的工作.
但是我們拒絕的.
尼達去事奉.
記住那事奉不會因為你的拒絕而停止.
神是會興起另外一些人去完成的.
只是那些自謙而棄權的人.
就會失去了寶貴的經歷.
少了神的同在,恩典.
失去了修成的喜樂.
失去了見到果實出現的那種恩悅.
更加失去了他的榮耀.
為神工作當然不是為自己得榮耀.
如果是這樣的時候.
絕對是神不會祝福的.

$^{761}$但是如果人願意為神擺上.
神必定會將榮耀交給他.
今天但願我們不會因為膽怯.
小信,恐懼而棄權.
這個最終損失的只是我們自己.
但是很可悲.
今天我們見到不少人.
不單單在事奉上棄權.
更加在自己應有的崗位上棄權.
這個不單單是損失.
更加是其他人的損失.
譬如作為丈夫的.
應該是妻子的頭.
姐妹應該不是很服這句話.
要做屬靈的帶領.
一家之主要好地用真理.
教導自己的兒女培育他們.
但可惜他們棄權.
或者因為自己屬靈心靈的脆弱.
又或者因為他們忙,很累放棄.
結果引來很多問題.
你不去好些教他們.
有其他那些人教他們.
但教他們的未必是正確的.
可能就會帶來他們的失敗.
損失會很多的.
所以求神幫助我們.
我們看看《神人共憤》.
聖經很簡單.
只是用一節來描繪這一場戰爭.
我們看看《釋詩紀》第四章第十五節.
「耶和華使西西拉和他一起車輛全軍遣亂.
在巴拉面前被刀殺敗.
西西拉下車步行逃跑」.
這是很簡單的描繪.
但我們肯定這句話的成功.
是神和人一起的努力.
共同奮鬥的結果.
巴拉雖然有不足的地方.
但我們不可以完全否定它的.

$^{801}$它最終都有召喚子民和它並肩作戰.
並且由它作領導.
第四章第十節.
「巴拉就知在西部旅人和那伐他里人到基底斯.
跟他上去亦有一萬人.
底波拉也同他上去」.
他們依從神的吩咐就上了他朴山.
在這一刻有人通風報信.
告訴西西拉.
敵人上了他朴山.
於是他就領大軍.
去到基信河預備包圍他朴山.
這個山很矮 不是很高.
好了 當我們要預備攻擊的時候.
我們看看第十二到第三節.
「有人告訴西西拉說.
阿比拿恩的兒子巴拉已經上他朴山去了.
西西拉就聚集所有的鐵車900輛和跟隨他的全軍.
從外邦人的下落前出來到了基信河」.
好了 巴拉就聽命帶兵下去迎戰那900輛戰車西西拉.
故此西西拉一定有信心的行動.
但是第十五節那裡很清楚說.
就是怎樣呢?.
就會令到他全軍覆沒.
我們一起來看十四,十五兩節.
「底波拉對巴拉說.
你起來 今日就是耶護華將西西拉交在你手的日子.
耶護華豈不在你前頭行嗎?.
於是巴拉下了他朴山 跟隨他有一萬人.
耶護華使西西拉派一切車輛全軍遣亂.
在巴拉面前被刀殺敗.
西西拉下車步行逃跑」.
900輛戰車一定還有很強的龐大的敵人.
但是為甚麼可以打贏這場仗呢?.
第十五節說到耶護華使西西拉.
跟他一切的車輛全軍賄亂.
但我相信那時候你記住.
一定不是雨季.
否則的話.
一位曉用善戰很厲害的將軍.

$^{841}$絕對不會猜到.
如果下雨會遭逢大劫的.
但是我們透過下個星期我們會詳細看.
第五章第四到第五二十到二十一節.
我們看到神做了甚麼工作.
我們一起來看第四到第五二十到二十一節.
我們在第五章底波拉詩歌裡面.
我們就知道原因在哪裡.
原來神是突然之間.
讓旱季裡面變成了傾盆大雨.
而且是突然之間下得很厲害.
令到河水氾濫.
於是他們就亂了.
還有就是因為河水氾濫.
令到那些泥土鬆軟.
那些鐵車就會泥足深陷.
突然之間河水可以氾濫.
給大家看一張圖片.
這個就是飛濟的圖片.
是他們給我的.
試過有一次突然之間.
早上下了一場大雨.
他們的河水就不斷地上升.
上升到哪裡呢.
大家可以看到右邊的圖.
地上全部淹沒.
你可不可以想像到.
河水氾濫得多厲害.
原來是短短一個多兩個小時.
就已經淹沒地上.
甚至上到二樓.
所以你留意到.
使得自此的戰車.
竟然成為他們的負累.
因為車一定陷在泥濘之中.
不能夠動彈.
結果遭到大敗.
而西沙拉下車就慌忙逃跑.
其實以色列人只是做甚麼呢.
只是做追殺的工作.

$^{881}$這個時候.
但他們仍然要努力.
我們看看士氣第四章第十六節.
巴拉追趕車輛軍隊.
直到外邦人的下落前.
西沙拉的全軍都倒在刀下.
沒有留下一人.
以色列人的軍力雖然薄弱.
但卻戰勝了強大的迦南盟軍.
這是一個神蹟.
可見戰爭的勝敗並不是靠自己的力量.
而是靠著神的工作和人的配合.
所以神只是做人不能做的事.
人可以做的.
必定自己要承擔責任.
我最近講學手爭.
都特別有一句說話.
We do the best.
We don't do the rest.
最後我們看滿.
就是受益.
西沙拉逃跑去到他的友好聯盟.
就是腓力基人希伯的妻子.
瓦易的家裡.
第四章第十七節.
為什麼會走到希伯的帳篷呢?.
原來耶賓和希伯間和好.
除了沒有磨擦之外.
更加有聯盟的意思.
原來腓力人通常住在巴勒斯坦的南部.
在《士師記》第一章才有.
而這一刻.
希伯間在巴勒斯坦的北部.
但彼此之間.
為什麼一個南一個北都會有關係呢?.
就是做生意.
在菲濟.
就是Ottweiler這個地方.
有不少中國人.
因為賺到錢.

$^{921}$所以他們移居太平洋島嶼.
由我去探訪他們.
去那裡做短訊.
而更加重要的是.
因為他們要入籍澳洲是很難的.
但原來入籍菲濟的島嶼很容易.
不是很大的數目就可以入籍.
入籍之後.
他們拿著菲濟或者太平洋島嶼的護照.
要入籍澳洲就簡單很多.
只要多走一圈.
所以雖然兩地的距離很遠.
你想想.
澳洲的下面的太平洋島嶼.
和中國有一段距離.
但很多中國人.
為什麼會這樣呢?.
就是因為有經濟的關係.
所以他們是以銅鐵業為生.
夏所王的戰車.
就是用鐵去做.
所以一定有密切關係.
所以西沙拉逃亡到希伯加.
然後.
他在這裡想躲起來.
結果被一個女士.
瓦易所殺.
西沙拉求瓦易給他水喝.
但瓦易卻把奶給他.
其實是奶油.
這是可能稀釋了的母語.
有催眠的作用.
正如有些人失眠.
臨睡前喝杯暖暖的奶.
令你可以好睡一點.
西沙拉因為戰爭.
因為逃跑已經很累.
很累.
所以又相信聯盟.
再加上瓦易視他為上賓.

$^{961}$為什麼呢?.
我們下次會看.
他會拿一些很漂亮的盤.
來盛上奶油.
給他.
又用被子給他蓋住.
讓他受保護甚至被封起來.
所以他就放心了.
誰知.
我們看看第四章21節.
西沙拉疲乏沉睡.
希伯的妻瓦易取了帳篷的缺紙.
手裡拿著錐子.
竊輕悄悄地到他旁邊.
將缺紙從他旁邊釘進去.
釘入地裡.
西沙拉就死了.
真的不要輕視女士.
你想想有多可怕.
他睡了一覺.
累了.
好了.
他拿了一支什麼呢?.
那些缺紙.
如果你去露營.
就是釘在地上.
然後拉著繩子.
就是那些缺紙.
悄悄地走過去.
然後.
在他旁邊那裡.
然後拿著錐子.
一插進去.
一插進去的時候.
不但突穿了一邊.
竟然是穿過了他的頭腦.
將那個錐子釘在地上.
即是當你左邊進.
右邊出.
釘在地上.

$^{1001}$那種力度是多麼大呢?.
所以西沙拉就死了.
當巴拉在董德瓦在家裡的時候.
這個強悍的敵人.
已經被一個寂寂無聞.
毫不顯眼的弱質女子殺了.
解救了以色列人的困境.
下次會看多一點.
或者可以說奇怪.
為什麼瓦易會這樣做呢?.
據推論.
瓦易其實在兩難之間.
如果他按和約接納西沙拉.
而西沙拉那個時候是喪家之犬.
如果瓦易不接待他.
不要啊.
你不要來.
你走啊.
始終是一個弱質女流.
所以西沙拉可能會發怒.
不接待我.
就會遷怒於他.
他就陷入一個危險的境地裡.
第二方面就是.
一個古代的時候.
一個陌生男性進入異性的家裡.
使那個女性受到懷疑和蒙羞.
有個男人去了家裡.
不知道發生什麼事.
不知道是不是有什麼問題.
對他來說是一個很大的羞辱.
而更重要的是我們相信瓦易.
也是一個忠信於耶華的信仰者.
他同樣看到神如何眷顧以色列人.
他的丈夫又和以色列人有血緣關係.
我們看看《史詩紀》第四章第十一節.
波西厄府河伯的後裔基里人希伯.
曾離開基里族.
到靠近基底斯撒拿欽的橡樹旁支撻帳棚.
原來是摩西的厄府的後裔.

$^{1041}$就是希里人 即是希伯.
所以他們會知道以色列人的歷史.
對神都有一定的認識.
而那時西西拉正在戰敗.
如果巴拉來到的時候.
他又沒有辦法去解釋.
為什麼要收藏西西拉的話.
也會引來殺身之禍.
即是巴拉說你窩藏我們的敵人.
這樣就有事了.
所以他在這兩難之間.
就將西西拉殺了.
這個除了解決了難題之外.
又可以成就以色列人的勝利.
因為他滿有信心.
不懼怕.
因為如果他不能殺死.
這個驍勇善戰 恐怖有力的西西拉.
他自己會有殺身之禍的.
所以他的決斷和猶豫不決的巴拉.
恰恰是相反的.
因為他滿有信心.
他有力量和智慧從神而來.
結果他成功了.
他滿就受益 得到獎賞.
事實上.
及後他得到以色列人的歌頌和表揚.
這個是下星期我們會看的.
《時詩記》第五章二十四節.
比眾婦人多得福氣.
比著帳篷的婦人更蒙福祉.
他得到很大的祝福.
其實底波拉的預言.
一早已經讓我們知道.
瓦易是神手中的器皿.
他成就了計劃.
最後我們做了一個結論.
巴拉雖然膽怯.
但我仍然說他小心.
他要求底波拉和他同行.

$^{1081}$他才去打這場仗.
事實上我們看不到底波拉在這場仗裡.
除了叫他下山之外.
他做了甚麼.
但巴拉仍然有付出.
雖然他失去了榮耀.
但我仍然相信他有神的祝福,眷顧.
所以以色列人仍然將他列在.
「時詩」的行列裡.
正如保羅指出.
如果我們用心去做的時候.
會得到獎賞.
敷衍的時候我們不會有獎賞.
但仍然會得救.
哥林多前書第三章色列度第十五節.
若有人用金,銀,寶石,草木和街.
在這根基上建造.
各人的工程必然顯露.
因為那日子要將它表明出來.
有火發現.
這火要試驗各人的工程怎樣.
人在那根基上所建造的工程若存得住.
他就要得獎賞.
人的工程若被燒了.
他就要受虛損.
自己卻要得救.
雖然得救.
乃將從火裡經過的一樣.
所以你留意到今天.
如果我們沒有用心去做.
即是做表面的事情的時候.
得不到神的祝福.
因為神都要試驗我們.
馬亦對神有信心.
不懼怕敵人的強悍.
而成功殺敵.
結果他得到榮耀和神的祝福.
所以謙你就是愚蠍.
你會招損.
但你滿有信心的話.

$^{1121}$你會受益.
我們會問究竟這場戰爭.
誰是士師呢?.
底波拉?.
巴拉?.
大聖經從來沒有提到是士師.
還是瓦易?.
但其實真正的士師是神.
祂掌管一切.
祂容讓子民受二十年的苦.
因為他們遠離神.
祂選照底波拉,巴拉,瓦易.
祂引領西沙拉到基遜河.
使大禹突然之間落下.
河水泛濫使大軍潰敗.
祂使用一個弱質女子.
戰勝一個元帥.
祂才是掌管一切.
有屬靈的領袖.
長者帶領著我們.
給我們意見是好的.
始終他們有一定的成熟生命和經驗.
如果我們將焦點放在這些人身上.
就有危機.
因為人始終是人.
有軟弱有不足的地方.
曾幾何時我們見到一些牧者.
也會使我們失望.
甚至搖頭嘆息.
如果我們信心建基於某些人身上.
當那個人失敗失落的時候.
會影響我們的信心.
帶我回教會的基督徒同學.
我還未正式信主.
他因為拍拖.
女朋友不喜歡就離開了教會.
幸好我的信心不是建基於他.
是建基於神.
所以焦點放在神的身上.
這就是巴拉和摩西最大的分野.

$^{1161}$讓我們用信心仰望這位掌管歷史的主.
一句很寶貴的話.
是我經常提醒自己的.
我的靈應當努力前行.
這就是《時詩記》第五章二十一節.
然後靠主你給我們的信心.
我們可以面對這個挑戰.
救到我們能夠.
懂得怎樣從你而來.
這份力量去承擔我們應有的責任.
多謝你.
讓我們看到聖女華語的寶貴.
也幫助我們能夠好好的去研讀.
以致我們的生命更加茁壯.
祝福不單單今天我們直接來大學研讀的.
參與的弟兄姊妹.
更加如果有弟兄姊妹或朋友.
他們來大學在YouTube收看的時候.
也同樣向神祈福他們.
聽孩子在面前的祈禱交托.
奉主德性的名.
阿們.
\newpage



\section{馬太福音 7:13-29}
\label{sec:032176}
\textbf{兩道門,你揀jor 邊道門?}~~~~王正傑牧師
\newline
\newline
連結: \href{http://www.church.com.hk/acms/content.asp?site=cdc\&op=show\&type=product\&code=032176\&layout=sermon}{\texttt{032176}} ~~~~ 語音日期: 2026-05-24
\newline
\newline
\hyperref[sec:032175]{\small{< < < PREV SERMON < < <}}
~
\hyperref[sec:index_chronic]{\small{[返順時目]}}
~
\hyperref[sec:032177]{\small{> > > NEXT SERMON > > >}}
\newline
\newline
馬太福音 7:13-29
\newline
\begin{longtable}{cl}
\hline
\hline
章節 & 經文 (和合本修訂版)\\
\hline
7:13 & \begin{tabularx}{0.7\textwidth}{X} 「你們要進窄門。因為通往滅亡的門是寬的，路是大的，進去的人也多； \end{tabularx} \\ \\ \relax
7:14 & \begin{tabularx}{0.7\textwidth}{X} 通往生命的門是窄的，路是小的，找到的人也少。」 \end{tabularx} \\ \\ \relax
7:15 & \begin{tabularx}{0.7\textwidth}{X} 「你們要防備假先知。他們到你們這裡來，外面披著羊皮，裡面卻是殘暴的狼。 \end{tabularx} \\ \\ \relax
7:16 & \begin{tabularx}{0.7\textwidth}{X} 豈能在荊棘上摘葡萄呢？豈能在蒺藜裡摘無花果呢？憑著他們的果子，就可以認出他們來。 \end{tabularx} \\ \\ \relax
7:17 & \begin{tabularx}{0.7\textwidth}{X} 這樣，凡好樹都結好果子，而壞樹結壞果子。 \end{tabularx} \\ \\ \relax
7:18 & \begin{tabularx}{0.7\textwidth}{X} 好樹不能結壞果子，壞樹也不能結好果子。 \end{tabularx} \\ \\ \relax
7:19 & \begin{tabularx}{0.7\textwidth}{X} 凡不結好果子的樹就砍下來，丟在火裡。 \end{tabularx} \\ \\ \relax
7:20 & \begin{tabularx}{0.7\textwidth}{X} 所以，憑著他們的果子就可以認出他們來。」 \end{tabularx} \\ \\ \relax
7:21 & \begin{tabularx}{0.7\textwidth}{X} 「不是每一個稱呼我『主啊，主啊』的人都能進天國；惟有遵行我天父旨意的人才能進去。 \end{tabularx} \\ \\ \relax
7:22 & \begin{tabularx}{0.7\textwidth}{X} 在那日必有許多人對我說：『主啊，主啊，我們不是奉你的名傳道，奉你的名趕鬼，奉你的名行許多異能嗎？』 \end{tabularx} \\ \\ \relax
7:23 & \begin{tabularx}{0.7\textwidth}{X} 我要向他們宣告：『我從來不認識你們，你們這些作惡的人，給我走開！』」 \end{tabularx} \\ \\ \relax
7:24 & \begin{tabularx}{0.7\textwidth}{X} 「所以，凡聽了我這些話又去做的，好比一個聰明人把房子蓋在磐石上。 \end{tabularx} \\ \\ \relax
7:25 & \begin{tabularx}{0.7\textwidth}{X} 風吹，雨打，水沖，撞擊那房子，房子總不倒塌，因為根基立在磐石上。 \end{tabularx} \\ \\ \relax
7:26 & \begin{tabularx}{0.7\textwidth}{X} 凡聽了我這些話而不去做的，好比一個無知的人把房子蓋在沙土上。 \end{tabularx} \\ \\ \relax
7:27 & \begin{tabularx}{0.7\textwidth}{X} 風吹，雨打，水沖，撞擊那房子，房子就倒塌了，並且倒塌得很厲害。」 \end{tabularx} \\ \\ \relax
7:28 & \begin{tabularx}{0.7\textwidth}{X} 耶穌講完了這些話，眾人對他的教導都感到驚奇， \end{tabularx} \\ \\ \relax
7:29 & \begin{tabularx}{0.7\textwidth}{X} 因為他教導他們正像有權柄的人，不像他們的文士。 \end{tabularx} \\ \\
[1ex]
\hline
\hline
\end{longtable}
$^{1}$大兄姐妹 有閘門和寬門 你選擇了哪個門呢.
我想做個調查 你覺得你自己已經選擇了閘門的請舉手.
你覺得可以的了 我又不會讓你去檢查一下 是不是.
那麼認為自己已經選擇了寬門的請舉手.
OK 那麼其他沒有舉手的呢 走了後門了.
我希望透過這段經文給我們一個審查.
究竟我們今天的生命選擇了哪一道門.
兩道門你選擇了哪一道門 是寬門還是閘門.
當我們看今天這段經文 其實是登山補品最後一個呼召.
大家都知道馬太福音的五至七章就是主耶穌很著名的登山補品.
這裡講述了跟從主的人悔改之後應該有的生命究竟是怎樣.
所以我們一看登山補品的內容.
我們會看到主耶穌會吩咐跟隨他的人要作炎作光 要愛仇敵.
跟隨他的人 他的義是要勝過民事化理錯人的義 也不可以憂慮 也要愛神愛人.
所以我們看到整個登山補品.
我們看到耶穌一步一步去描繪作為一個天國子民.
究竟你應該有的究竟是怎樣的一個樣式.
而登山補品更加不單止是叫跟隨他的人要聽從耶穌基督的教訓.
更重要的是 主耶穌要呼召要進入天國子民的一種生活方式.
信仰不只是認知 而是在說我們與神立日後 我們要活出一個生命見證.
所以弟子們 今天我們回教會.
我們不是在說我們究竟知道多少聖經 我們認識多少聖經.
聖經的教導不是停留在我們的腦袋裡面.
不只是一個理性的認知 而是在說我們究竟聽了之後.
我們讀了之後 我們查了之後.
究竟我們有多少聖經的道理我們能夠活出來.
這個才是更加重要.
這個才是天國子民生命的方式.
所以當我們去看今天這段經文.
主耶穌一開始就很清楚地吩咐跟隨他的人要進閘門.
閘門是一道很容易走的路.
門是寬闊的 路是大的 進去的人是多的.
為甚麼呢 這道門很簡單.
其實入閘門就代表著甚麼呢.
你是不去進行天父旨意的人.
所以你進入這道閘門 你不需要悔改.
你跟隨這個世界的價值觀.
而這道門是在說甚麼.
你看上去好像很自由 你很舒服 你喜歡怎樣就怎樣.
但它是怎樣呢 漸漸的離神越來越遠.

$^{41}$而最重要那句是甚麼.
結局是我們會引導滅亡.
接著閘門 主耶穌告訴我們.
進閘門是一條通往生命的路.
閘門是通往生命的 為甚麼呢.
我們去進行天父的旨意.
在這道門裡面有主的同行 有恩典去扶持.
是有永生的盼望 但主耶穌告訴我們.
這條路是少的 找到的人都是少的 為甚麼呢.
因為要悔改 要捨己 要背十家 要信復神.
等於說到這裡 看到寬門和閘門.
對於一般人來說 我們看到的時候.
我們很容易選擇的.
大部份的人 你會選擇容易的事做 還是選擇困難的事做.
連埃連今晚 警方都會向我們呼籲.
你要鎖好門窗 因為會很容易有盜竊.
因為那些賊都是取而不取難.
所以我們要去選擇 我們會選擇甚麼.
選擇容易的事做 還是選擇困難的事做.
選擇容易的事做的話 請舉手.
多了些人.
選擇困難的事做的話 請舉手.
沒有舉手的又是甚麼.
不走路 坐在這裡最好.
這麼難走路 都不知道走不走得好.
所以主子說了給我們聽 有寬門和閘門.
閘門是通往新明 但是路是少的.
找到的人都少 因為要付代價.
最近有一幅警世圖 關於教會的.
這幅警世圖是說 教會正在邁向死亡的七個徵兆.
第一個徵兆是甚麼呢.
崇拜已經變成一個表演 而非參與.
即是人回到這裡已經不是一個敬拜者.
而是只是作為一個旁觀者去觀望.
所以回到這裡 我們不會透過唱歌去讚美神.
我們回來會留意些甚麼呢.
留意一下國賢剛才有沒有走音.
他剛剛是否拍子入 有沒有唱慢了.
然後伴唱有沒有走音 打鼓是否正確拍子.
Powerpoint有沒有打慢了 音控有沒有控制不好.

$^{81}$主席祈禱有沒有甩甩咳咳等等.
我們回來崇拜就變成一個看表演的人.
但我們每一次回來教會去敬拜其實很簡單.
我們都是用心靈承受去敬拜.
透過詩歌唱出的歌詞來讚美我們的主 我們的神.
這才是最重要的.
所以第一個徵兆就是當崇拜變成表演而非參與.
第二件事就是那些訪客來了一次之後就消失了.
我記得早前我們在那紅團形容教會很有邊界感.
我們對那些新朋友很有邊界感就是我們都不會怎樣理會他們.
所以很有邊界感.
所以很多新朋友來了一次之後就消失了.
第三個徵兆就是同一班人坐在同一個位置上.
有十年已經不變了.
你幫我看看你旁邊那個是否已經坐了十年不變的位置.
我覺得炳雄和兆娟應該是超過十年了.
我每次都好像看見你們坐在那裡.
Patrick都好像是.
你們坐在那個位置上很詳情.
這些人回來就是用一種安逸取代了穩定.
教會一選擇安逸我們就會停滯 不會向前.
第四就是逃避衝突就變成了一種新的文化.
因為我們都很怕衝突.
特別是回來教會都要衝突我們是很不喜歡的.
有兩方面 一方面就是因為怕衝突我們就客氣氣.
我們不會講真話.
但另一邊就是因為我們太喜歡講真話了.
所以就打死了很多人.
有時候我們都不知道應該怎樣去處理.
但當我們逃避衝突變成一個新文化的時候.
就缺乏了一種真實性的合一.
就並不是一種真正的聯合.
我們心裡就失去了一個真正的交流.
第五樣的徵兆就是下一代正在流失.
如果沒有了年輕的一代 教會就沒有未來.
弟兄姊妹 我針對的不只是年輕的一代.
譬如年輕人 教會沒有下一代.
不是說年輕人會流失 或者兒童會流失.
我更加在講的是我們其他的年齡層的人.
我們都沒有了下一代去承接.

$^{121}$我們有資深的基督徒 究竟有沒有新加入的基督徒呢.
這個都是下一代.
第六樣的徵兆就是什麼.
只是有士工卻失去了使命.
我們可能很喜歡搞很多士工.
士工一個接一個 一個接一個.
但其實我們不知道做來做什麼.
然後可能都是為做而做.
究竟我們做這個工具.
究竟我們最終的目標是什麼呢.
其實我們是不是要去祝福這些群體呢.
祝福年輕人 祝福長者等等.
我們要找回一個使命.
第七個徵兆就是追求舒適勝過馴服.
特別是當我們信奈了.
其實我們都希望安安樂樂.
回來教會開開心心.
有團契 有主一學 有崇拜 就很舒服.
當我們說要付代價.
我們的生命要去毀身的時候.
我們就會感到很困難.
回來崇拜的椅子坐得舒服都很重要.
昨天我們在轉浸舉行婚禮.
昨天在婚禮裡面 我坐在台上.
轉浸那張椅子 比土浸那張椅子更舒服.
這裡其實很重要 說的是什麼呢.
我們會不會回來教會慢慢追求安逸.
多於想要去馴服 要去毀身呢.
等下我和你分享這幅警世圖.
我不是要和大家分享.
我們一起去檢查土浸.
究竟我們是否都有齊這七個徵兆.
我們是否教會都要邁向死亡.
我不是在說這件事.
我很想告訴大家的是.
當我去看見土浸的時候.
我看見到三件事.
有一部份的人可能真是.
我們會享受安逸 享受穩定.
因為可能人生都去到某個階段.

$^{161}$我們都希望穩穩定定 開開心心的去回教會.
另一群人我看見的是很疲累 很疲乏.
有很多的事奉.
已經不斷地做 不斷地做.
很疲乏 很疲累.
第三類是一班可能已經很挫敗 很失望的.
可能過去曾經參與過很多的服侍.
在事奉裡都很賣力.
但可能遇到挫敗.
可能又受到一些批評.
可能對教會會有失望等等.
我最近聽了一段對話.
是一個以前回土浸的年青人.
現在已經不再年輕了.
已經結婚 有小朋友.
但他沒有再回來土浸.
他分享的是.
為什麼現在不再回教會.
因為過去的日子.
每次回來教會都是工作.
工作完又繼續工作.
工作都不是問題.
但最後他覺得很沒有被務養.
回來教會工作是重要的.
但會中你們有被務養 都是很重要的.
所以今天我們看見的是.
究竟我們是否正在進一個閘門.
我們是否正在順服天父的旨意.
我們是否好好事奉.
是否好好做好務養的工作.
這是我們今天需要認真去面對.
所以基督教的一位作家C.S.Lewis.
他說過一番話.
我們進閘門是真的要去付代價.
他說我選擇宗教不是為了讓自己開心.
他說我知道一瓶波多爾紅酒.
就能令到他開心.
如果你想找一個讓你感到舒服的宗教.
他絕對不會推薦基督教.
因為他要表達的就是.

$^{201}$其實我們的信仰就是一份的委身.
是要付代價.
當你要去進行天父的旨意.
你就是要去付代價.
主耶穌吩咐跟從他的人要進閘門.
是要付代價.
是要委身.
這是一個必然會出現的事情.
但是這一道門是通往生命的.
所以去到這裡弟兄姊妹.
我想再問大家一次.
弟兄姊妹.
你究竟選擇了哪一道門.
你選擇了閘門.
還是寬門.
你選擇了進閘門.
你去進行天父的旨意.
還是你選擇了一道寬門.
沒有去聽從天父的旨意呢.
然後接著.
再進一步.
主耶穌就提醒我們.
他說要防備假先知.
他們的外表是怎樣呢.
是很像一隻羊的.
裡面是怎樣呢.
殘暴的狼.
等一會你幫我看一下你過來的那個人.
他像一隻羊還是像一隻狼.
我經常說在教會裡面.
我們有兩類人.
有兩類假先知.
一類假先知是甚麼呢.
就是他處心積慮的假先知.
他真是錯誤的引導人走向寬門.
另一類是甚麼呢.
我們不知不覺會成為了假先知.
我每次說很多人都接受不了.
因為甚麼呢.
你有誰可以接受得到你旁邊那麼可愛的弟兄姊妹是一隻狼呢.

$^{241}$你再看一下她是不是很可愛.
弟兄姊妹.
當我們回到教會.
我們就很容易怎樣.
很容易不知不覺就變成了一個假先知.
假先知最大的問題是甚麼呢.
就是將人帶離去閘門.
帶離開了去閘門.
沒有去聽從天父的旨意.
走向寬門.
說一些令人聽了很舒服的說話.
告訴別人不需要悔改.
不需要改變生命.
信仰只停留在一個表面的上面.
很容易的弟兄姊妹.
我們是不知不覺可能會做了假先知.
是甚麼時候呢.
當你很愛心的時候.
當你有過份的包容的時候.
當你可能過份體諒的時候.
我們就會怎樣呢.
不小心的引領人走向寬門.
譬如當有人說.
他不想十一奉獻.
我們會跟他說甚麼.
我不用十一奉獻的.
教會也有錢的.
多你那幾千元不多.
少你那幾千元不少.
教會沒有錢也沒有問題的.
天父有很多錢的你放心.
又或者可能有時.
當我們說你要彼此相愛.
不要緊的.
你愛不下那個人.
你就不要跟他坐.
不要坐在他旁邊.
愛不到就愛不到.
不要勉強自己.
愛別人吧.

$^{281}$又或者可能叫人去傳福音.
你不用傳的.
你不傳福音而已.
你不傳.
最重要是天父要他相信.
天父要他相信他就會相信的.
你不傳還有很多人會傳.
放心不用傳的.
各位姐妹.
可能有時我們都不知不覺.
因為我們可能有過份的愛心.
有過份的體諒.
有過份的包容.
我們會不會都變成一個假先知.
引導人走向寬門.
沒有去聽從天父的旨意呢.
接著主耶穌就說.
荊棘上不能摘到葡萄.
摘泥裡豈能摘到無花果.
其實荊棘和摘泥.
這些都是一些沒用的植物.
如果當它生長在葡萄樹裡面也好.
無花果樹那裡也好.
就會阻礙了那些樹的成長.
影響它們的收成.
同樣的是.
主耶穌在這裡要引用荊棘和摘泥.
就是要比喻.
在我們生命裡的假先知.
引領我們離開那個閘門.
走向寬門.
不去聽從天父的旨意的時候.
就會影響我們成長.
影響我們的收成.
所以主耶穌接著用了四句話.
說完和沒說好像沒什麼分別.
好樹就結好果子.
壞樹就結壞果子.
好樹就不能結壞果子.
壞樹就不能結好果子.

$^{321}$是不是你們聽完也懂的.
對不對.
很簡單 其實要說什麼.
如果你是好樹的.
你就會結到好果子.
如果你是壞樹的.
你只能夠結出壞果子.
所以他告訴我們.
好樹結好果子就是真正跟從主的人.
他的生命慢慢流露出.
有愛 有溫柔 有聖潔 有馴服 有憐憫和忠心.
這是一個真實的生命的改變.
不是一個表面化的信仰.
而壞樹結壞果子就是.
生命沒有真正被主去改變.
即使他外表表現得很屬靈.
也無法結出神所喜悅的果子.
果子是怎樣.
是生命真實的流露.
不能夠長期地去假扮.
假扮不出來.
但是鄧小姐.
這裡最重要的問題是什麼.
這些壞樹會結壞果子的樹.
最終是要怎樣.
撞下來.
凋在火裡.
如果鄧小姐.
今天當你走向寬門.
當你的生命成為一棵壞樹.
你結壞果子.
最終樹是要怎樣.
就是撞下來.
凋在火裡.
所以耶穌進一步說.
憑他的果子就可以認出這些假先知.
耶穌說憑著他們的果子.
就可以認出他們來.
果子不是表面的敬填.
而是一種持續生命的見證.

$^{361}$所以當我們看見我們的生命有改變.
我們就知道它真正有悔改.
鄧小姐這裡要留心.
其實我們回教會越久.
我們會越容易出現一個問題.
就是我們生命越來越少改變.
生命越來越少成長.
通常說生命會成長會改變.
通常在什麼時候你見到最多.
初信.
朕禮寫見證的時候.
當你過了初信和朕禮的時候.
你就很少再說自己生命有什麼成長.
有什麼更新有什麼改變.
越來越少的.
你問問你旁邊那個最近生命有什麼成長.
他最多會告訴你老了一歲.
最近膝蓋有點痛.
腰骨有點不舒服.
睡得不太好.
我們會越來越少說生命有什麼成長.
我們很容易變成一群.
只是講信仰的人.
但真正的門徒是要怎樣.
不單止是懂得講信仰.
更加是要活出信仰.
你的生命是真的經歷了悔改.
你生命真的有改變.
真的有成長.
我認識一位姐妹.
在教會認識她的時候.
我覺得她的人很友善.
和她相處很舒服.
亦都很客氣.
大家一起合作事奉.
合作得很好.
很配合得到.
所以我都很欣賞這位姐妹.
去到她受浸加入教會的時候.
在浸禮上她就要講見證.

$^{401}$她見證的題目是.
她說這個見證是一個超級港女的見證.
她說她從小到大都是一個超級的港女.
這個港女是怎樣呢.
她去到哪裡都是女王.
由她成長的時候.
對著父母都是一個超級港女.
到認識了她老公.
都是一個超級港女.
到生了女兒.
她都是一個超級港女.
去到任何時候都是罵人的.
在家裡經常都罵她老公.
這件事看不過眼.
那件事就罵她.
然後女兒出生之後.
都是不斷在罵她.
但她說當她相信耶穌之後.
她生命真的經歷了很大的改變.
她沒有再罵她老公.
沒有再罵她女兒.
她沒有再做一個超級的港女.
我聽到這裡我很感動.
也很為她感恩.
我是為她老公和女兒感恩.
因為神真的改變了她的生命.
從此她老公和女兒就不用被人罵了.
但最重要那件事.
最感恩那件事是什麼.
在當天這個浸禮上.
她的老公和女兒.
都舉手決志相信耶穌.
是因為他們真的看見她的老婆.
這個媽媽.
她相信了之後她的生命的改變.
弟兄姊妹.
從她見到的好果子.
我們認出它是什麼.
是好處來的.
然後接著主耶穌再進一步.

$^{441}$祂說了一段很震撼的話.
祂說不是每一個稱呼我主呀主呀的人.
都能進天國.
原來弟兄姊妹.
你的外表屬靈.
不代表你真正是屬於神.
你有因此去服侍.
但也不代表你有真正的去認識主.
對我們來說很簡單.
我們回來教會.
我們見到什麼人是覺得他特別屬靈.
很認識聖經.
很熟悉聖經的我們覺得他很屬靈.
走到回來的樣子很端莊的那些很屬靈.
你幫我看看你旁邊那個端不端莊.
以前基督徒有個樣子.
穿著白衣 長裙那些一定是姊妹.
摔衣服西褲那些一定是弟兄梳了個彈托頭.
基督徒好像是有個樣子的.
我們都會從這些很屬靈的外表.
他講話可能溫柔.
祈禱裡面有十幾段經文在裡面.
有些金句.
這些就很屬靈.
我們都會看這些.
又或者我們見到一些很有恩賜的.
出來做個主席 領師.
我們見到他真的很多恩賜.
做很多事都是很好的.
我們就覺得這些是很屬靈.
這個是我們很正常的一般的看法.
但天主教這裡告訴我們.
不是每個稱呼我主也主也的人都能進天國.
外表屬靈的不等於他真正屬神.
有些人可能是很懂聖經的.
很熱心去聚會.
很懂得講那些屬靈的說話.
但是他的生命沒有真正去馴服神.
他外表可能只是一些宗教的行為.
不能夠代替了他內心.

$^{481}$究竟有沒有真實的去毀身.
另一種有恩賜的服事.
不等於他真的代表他真正的去認識主.
經文告訴我們耶穌說什麼.
他說有人奉主名去傳道講.
鬼行以來很多的恩賜.
但是耶穌說我從來不認識你.
問題不是在於恩賜有多少.
而是究竟這些人有沒有真正去進行天父旨意.
講到這裡我們很容易理解.
因為其實我們都知道.
單純你講你有一個屬靈的外表.
或者你有很多恩賜.
真的不代表你真正屬神的人.
我們最近基督教都有一些人.
我們見到他是很有恩賜.
他們很屬靈的.
甚至他們成為領袖的.
但是他們的生命有很多軟弱跌倒.
你們見到你們的左邊.
是一位基督教著名的作家.
他承認他曾經有八年的婚外情.
中間那位是新加坡最大教會的主任牧師.
他曾經因為挪用教會的資金.
去支持他太太的演藝事業.
最終他落得的是牢獄之苦.
坐了牢三年.
最後那位是韓國最大教會的主任牧師.
同樣都是教會的錢銀的問題.
最終被判有罪.
所以我們都很明白.
不是單單講你有一個屬靈的外表.
或者你有很多恩賜.
就代表你是真正屬神的人.
但是主耶穌說對於這些人.
我從不認識你.
主耶穌在這裡引用了.
你們這些作業的人離開我吧.
因為耶和華聽了我哀哭的聲音.
這個是來自詩篇,律篇,八字.

$^{521}$他在這裡要表達什麼.
當時詩人他要表達.
就是因為那些不認識神的人.
攻擊他,傷害他.
然後最後他發出這句話.
就是說你們這些作業的人離開我吧.
因為他們都是耶和華的仇敵.
耶和華已經聽了他哀哭的聲音.
同樣在這裡要告訴我們.
這些假先知是耶和華的仇敵.
他引領人離開去進行天父的旨意.
所以縱使他們有很屬靈的外表.
有很多恩賜都好.
主耶穌說我從不認識你.
所以在這裡.
等姐妹.
主耶穌告訴我們.
真正的門徒是什麼人.
是一些聽道.
而且有馴服,願意去行的人.
登山寶訓的核心不是停留在.
我們只是一個頭腦上的知道.
不是我們聽了一篇道我們有感動.
不是我喜不喜歡聽誰講完他講什麼道.
或者查哪一段聖經.
我們很喜歡哪一段聖經.
而是需要的是.
我們願意去行愛神,愛人.
饒恕,誠實,憐憫.
這些不是一些口號.
而是我們真是每天活出來的生命.
弟兄姊妹.
過去一段時間.
我去了不同的團契.
帶了很多的茶經.
有很多的分享.
我都聽到有兩件事是.
很多時弟兄姊妹提出是最困難最困難的事.
是天父的旨意是要去遵行.
但是很難做到.

$^{561}$所以沒有做到.
有兩件事.
第一件事知道是什麼嗎.
第一件事是饒恕.
當講到要饒恕一些可能傷害過你.
出賣過你.
攻擊過你的人.
我們要饒恕他們呢.
很困難.
第二件事知道是什麼嗎.
愛仇敵.
我見到那個仇敵.
我看見他的前面.
又看見他的後面.
我怎樣去愛他呢.
我不攻擊他已經算很好.
所以很難.
但是弟兄姊妹.
主宋告訴我們.
真正的門徒.
是當我們聽到之後.
我們願意去順服和忠恆.
這個才是真正的門徒.
然後最後一個部分.
耶穌就用了兩個蓋房子的人.
來作為整篇登山寶訓的一個總結.
這兩個人他們有一個共通點.
他們都聽過耶穌的說話.
但是最大的分別是什麼呢.
就是一個聽了就去做.
一個聽了就不做.
聽了就去做的人是什麼呢.
是一種將根基建在盤石上的人.
耶穌稱他們是什麼呢.
是聰明人.
因為他們將他們生命的根基.
建造在盤石上.
盤石就是主耶穌.
真正穩固的信仰.
不是我們查過多少次的聖經.

$^{601}$我們聽過多少篇的道.
我們背熟了多少段的經文金句.
而是我們究竟有沒有按照神的話.
來去活.
來去行道.
而當我們能夠將我們的生命根基.
建在耶穌的上面的時候.
告訴我們.
風浪是仍然會有的.
困難是仍然會存在的.
但是生命是怎樣呢.
是不會倒塌.
因為我們的根基都在耶穌基督的裡面.
另一個人就是聽了就不去行的人.
他們的生命根基就建在沙土上.
聽了主的話就不去行的人.
就好比那些無知的人.
平時看來沒什麼問題.
等於聽了神的話.
然後沒有去做.
是沒什麼問題的.
你看看你旁邊那個.
他真的沒什麼問題.
他沒有聽神的話.
真的沒什麼問題.
但是當風吹雨打.
水沖來的時候.
這些聽了神的話.
沒有去做的時候.
就會怎樣呢.
整間屋就倒塌了.
我們人生總會有風浪.
有壓力.
有試探.
有痛苦.
有失望的時候.
如果我們的信仰.
只是一些知識.
感覺或者習慣的時候.
當遇到這些風浪的時候.

$^{641}$我們就很容易倒下.
所以這裡很簡單.
聰明人和愚昧人.
聰明人就是主耶穌告訴我們.
聽見祂的話就去行的.
你就是聰明人.
你就好比就是.
將房子建在盆石上.
風吹雨淋.
房子總是不倒.
但是如果你聽了神的話.
不去行的人.
你就是一個愚昧人.
你的房子就是建在沙土上.
風吹雨淋.
房子倒塌.
並且倒塌得很大.
所以登山寶訓是一個呼召.
是一個完整的生命的旅程.
我們不是停留在.
純粹就是聽到.
領受.
有感動等等.
然後我們讓每一次聽見主的話.
就成為我們一個生命改變的起點.
從聽見到馴服.
從馴服到我們的根基穩固.
這個才是耶穌在這裡.
要描述一個聰明人的路.
是應該怎麼走.
所以今天我們要聽見的時候.
我們要學習去馴服.
當我們馴服.
我們去行道的時候.
以致我們生命的根基能夠穩固.
這個才是要成為一個聰明人的路.
是要怎麼走.
最後我和你分享兩個圖.
你見到黃色這幅圖.
是一些長輩圖.

$^{681}$為什麼我形容他長輩圖呢.
因為長輩看到是很同意他這些道理的.
但是頭一段其實不是我最想講的.
我最想分享的是最後那幾句.
不過我都照講.
讓大家看完長輩圖都有點同意.
這幅長輩圖說什麼呢.
不管在什麼地方上班.
請記住.
工作不養閒人.
團隊不養懶人.
沒有哪個行業的錢是好賺的.
賺不到錢賺知識.
賺不到知識賺經歷.
賺不到經歷賺閱歷.
只有先改變自己的態度.
才能改變人生的高度.
是不是長輩特別覺得很合適.
最後那幾句我反而覺得很重要.
他說.
讓人迷茫的原因只有一個.
那就是本該拼搏的年紀.
卻想得太多.
做得太少.
弟兄姊妹.
當我們聽到天父的旨意.
當我們聽到主耶穌的吩咐.
我們會不會都是.
想得太多.
做得太少.
當主耶穌說我們要彼此相愛的時候.
你看見你旁邊的弟兄姊妹.
你就開始在想.
究竟我有沒有跟他們打招呼呢.
然後想著想著.
我不是很記得他們叫什麼名字.
然後又想著想著.
我又不記得問他們做哪一行.
又再問他們會不會不開心呢.
然後.

$^{721}$我又不知道有什麼話題要跟他們說.
問不問他們好呢.
最後我們想了這麼多之後我們決定什麼.
就不再問了.
我們會不會很多時候在教會裡面.
我們聽到天父的吩咐.
聽到他的旨意.
我們都選擇是.
想得太多.
做得太少呢.
所以我很喜歡這套戲.
大家知不知道這套戲是什麼戲.
是的.
壯志凌雲.
不過這個是第二集.
第二集是獨行俠.
其中有一句對白我很喜歡.
就是.
Don't think.
Just do.
弟兄姊妹.
其實今天我們要去回應天父的吩咐.
按照天父的旨意去行.
我們很多時候都是.
要.
Don't think.
Just do.
可不可以幫我和你旁邊的人說.
Don't think.
Just do.
說普通話大家不是很喜歡.
說英文應該可以了吧.
所以弟兄姊妹當我們能夠聽到.
天父的吩咐.
他的旨意的時候.
讓我們都Don't think.
Just do.
當我們能夠聽了就去行.
我們就做了一個聰明人.
所以今天弟兄姊妹.

$^{761}$主耶穌仍然是呼召我們.
呼召我們.
究竟你要去選擇哪一道門呢.
你要去做這個選擇.
閘門不是一條容易的路.
是一條要付代價的路.
是一條毀身的路.
但它是通往生命.
有平安和天國的路.
以致我們真的能夠.
當我們走這條路的時候.
我們成為一棵好樹.
我們能夠結出好果子.
能夠成為一個聰明人.
當我們聽到神的吩咐的時候.
我們就照著去做.
以致我們的根基是建立在盤石上.
縱使風吹雨打.
我們仍然能夠站立不穩.
站立不住.
弟兄姊妹.
今天你要選擇哪一道門呢.
你會選擇閘門還是呼門呢.
盼望讓我們一起.
都去回應那位愛我們的主.
選擇這道閘門.
走這條屬天的道路.
讓我們一同低頭.
我們有一段禱告的時間.
能夠遇見你.
認識你.
得到你的救恩.
主今晚我們來到你面前.
我們懇求你的愛.
你的話語再一次去激勵我們.
幫助我們的生命.
我們真的能夠可以選擇.
去走在閘門的裡面.
幫助我們能夠去聽.
能夠去馴服.

$^{801}$能夠可以去活出天父你的旨意.
最多我們可能會面對很多困難.
很多挑戰.
有時候我們也會有很多的軟弱.
主我們懇求你幫助我們.
讓我們仍然能夠堅持.
走在閘門的裡面.
讓我們真的成為一個聰明人.
我們聽了你的吩咐就照做.
以致我們生命的根基真的建立在盤石上.
遇到風吹雨打.
我們都能夠可以堅定堅韻.
主我們就懇求你親自去幫助我們.
帶領我們.
讓我們生命時刻刻都是警醒.
我們能夠可以有悔改有更新.
以致我們真的見出好果子來.
讓我們的生命不單止我們自己得著祝福.
更一樣我們能夠可以成為我們身邊人的祝福.
求主你去幫助使用我們每一個.
主我們多謝讚美你.
我們在這裡祈禱交托.
是奉教主耶穌基督的聖名致祈求.
阿們.
\newpage



\section{約翰福音 20:19-23}
\label{sec:032177}
\textbf{從囚禁到釋放}~~~~林玉英牧師
\newline
\newline
連結: \href{http://www.church.com.hk/acms/content.asp?site=cdc\&op=show\&type=product\&code=032177\&layout=sermon}{\texttt{032177}} ~~~~ 語音日期: 2026-05-24
\newline
\newline
\hyperref[sec:032176]{\small{< < < PREV SERMON < < <}}
~
\hyperref[sec:index_chronic]{\small{[返順時目]}}
~
\hyperref[sec:032178]{\small{> > > NEXT SERMON > > >}}
\newline
\newline
約翰福音 20:19-23
\newline
\begin{longtable}{cl}
\hline
\hline
章節 & 經文 (和合本修訂版)\\
\hline
20:19 & \begin{tabularx}{0.7\textwidth}{X} 那日（就是七日的第一日）晚上，門徒因怕猶太人，所在的地方門都關了。耶穌來，站在當中，對他們說：「願你們平安！」 \end{tabularx} \\ \\ \relax
20:20 & \begin{tabularx}{0.7\textwidth}{X} 說了這話，他把手和肋旁給他們看。門徒一看見主就喜樂了。 \end{tabularx} \\ \\ \relax
20:21 & \begin{tabularx}{0.7\textwidth}{X} 於是耶穌又對他們說：「願你們平安！父怎樣差遣了我，我也照樣差遣你們。」 \end{tabularx} \\ \\ \relax
20:22 & \begin{tabularx}{0.7\textwidth}{X} 說了這話，他向他們吹一口氣，說：「領受聖靈吧！ \end{tabularx} \\ \\ \relax
20:23 & \begin{tabularx}{0.7\textwidth}{X} 你們赦免誰的罪，誰的罪就得赦免；你們不赦免誰的罪，誰的罪就不得赦免。」 \end{tabularx} \\ \\
[1ex]
\hline
\hline
\end{longtable}
$^{1}$各位弟兄姊妹平安.
今日教會舉行的嘉賓是教會成禮.
當中有很多是第一次參與教會聚會的來賓.
還是你的朋友邀請你在我們當中.
我們代表教會來歡迎大家.
願你們平安.
我們一起聆聽上帝的話之前.
請同心低頭祈禱.
願聖靈與我們同在.
在我們的心裡來運行.
主你今天要我們聽的是什麼話.
願你打開我們的心來領受.
讓講的和聽的都同蒙你的月立.
祈禱奉主耶穌 祈禱明球 阿們.
不知道大家有沒有試過.
將自己藏起來的經歷.
不想見人 不想被人找到.
不知道你有沒有這樣的經驗.
回想那個歷程.
你想想你做了什麼自己要躲起來.
又或者你經歷了什麼恐怖的事.
你要將自己收起來.
原來每一個人心裡都有一道關得很緊的門.
有時我們以為關上門.
是為了防備外面給我們的傷害.
不過當我們這樣做的時候.
實際上我們是將自己困在一個籠子裡面.
這個籠子可能是一個.
覺得自己是一個失敗的人的籠子.
又可能是一個我不願意寬恕那些傷害我的人的籠子裡面.
總之就不知道為什麼.
總是覺得自己有時很辛苦.
今天經文符的耶穌告訴我們.
祂要穿透我們這道關得很緊的門.
祂要帶來的不是一個審判或者定罪.
而是能夠醫治我們的傷痕.
展開綁著我們在籠子裡面的一份寬恕.
我們回到聖經.
大家試想像一下.
當時門徒將自己帶到房間裡面.

$^{41}$門關得很緊.
他們最怕的可能不是來搜捕他們的猶太長官.
他們最怕見到的可能是耶穌.
因為三天之前耶穌被敲問,被鞭打,被侮辱.
被釘在十字架上最痛苦的時刻.
這班門徒雞飛狗走.
只顧著自己的性命.
四處逃奔.
為了不要受牽連.
甚至死後都說「我不認識耶穌」.
或者選擇不出聲.
所以現在他們心裡充滿了內疚,羞愧和不安.
當然我們知道這是七日的第一天.
剛剛那些婦女告訴他們.
埋葬耶穌的墳墓空了.
他們說耶穌復活了.
「什麼事啊?」.
於是那班門徒在這方面很害怕.
很羞愧,內疚.
也很困惑,不安.
突然間耶穌出現在他們當中.
我們試想像一下.
如果你當時在場.
你會有什麼反應?.
試想像一下.
如果當時我們跟那班門徒在一起的時候.
你有什麼反應呢?.
你的心會想些什麼呢?.
其實我們可能跟門徒一樣.
情緒很複雜.
很害怕.
很慚愧.
想找地方躲起來.
不過耶穌竟然沒有一句質問.
沒有一句責備.
「彼得,你為什麼說不認識我?」.
「你們一個個扔下我」.
沒有啊,耶穌沒有說.
祂只是連續對門徒說了兩次.
「願你們平安」.

$^{81}$這句說話表達了耶穌對門徒的寬恕.
寬恕不是因為對方做對了.
而是因為耶穌選擇去愛.
耶穌這句說話所表達的接納.
徹底粉碎了門徒的羞愧感.
和內心的不安.
各位靈姐妹,各位來賓.
你是不是覺得自己不夠好.
或者過去實在有太多的污點,太多的失敗.
我們不敢來接近上帝.
今日經文告訴我們.
是耶穌祂主動.
是耶穌祂做主動.
祂想穿透我們的心.
祂要給我們的是醫治和平安.
而不是定罪.
耶穌向門徒顯示祂復活的身體.
在另一處經文記載祂鼓勵門徒.
走來摸祂手上,腳上和肋旁側邊的疤痕.
耶穌這樣做不是要記得這些痛苦和祂所受過的傷害.
「你看,你害得我變成這樣」.
不是.
我們知道我們每一個人都無法抹除記憶.
因為發生了的事就是發生了.
耶穌這樣做是要告訴門徒.
祂不是鬼魂來的.
祂是真正復活.
祂的身體真正復活了.
不過最重要的是祂要告訴門徒.
祂要藉著這些疤痕告訴門徒.
痛苦為祂帶來的傷害.
不單止不會將祂困在苦毒的牢籠當中.
不會將祂困在苦毒的牢籠裡.
祂要轉化這些傷害成為祝福的記號.
在座的很多頂智妹信耶穌很久了.
又或者如果你未信的來賓.
你可能在一些學校.
小學,幼稚園,中學.
都聽過約瑟的故事.
約瑟原本是天之驕子.

$^{121}$他受著他爸爸很大的寵愛.
所以他十個哥哥氣得他要死.
不過約瑟很天真.
他不知道.
他沉浸在這個溫室裡.
當他出去曠野.
爸爸叫他出去看看他的哥哥會否很安全.
因為他們在曠野來牧養他們的羊.
遠遠看到我們的弟弟.
嘩!機會來了.
爸爸不在這裡.
於是打算將他殺了.
不過最終不要殺了他.
我們有什麼利益?.
將他賣了去埃及做奴隸.
你想像一下原本你是天之驕子.
突然間你的兄長.
你的哥哥這樣對你.
你都驚魂未定.
就賣了去做最低賤的奴隸.
聖經沒有形容約瑟當時是怎樣.
不過只是形容他做好他眼前的事.
他很努力的來到主人家裡做好他自己.
去得到主人家的賞識.
當然聖經說是上帝幫助他.
讓他再一次成為天之驕子.
一人之下 萬人之上.
他的主人什麼都不知道.
因為將一切交給他打理.
不過不知道什麼原因.
他的主母陷害他.
於是他再一次來到他被他的主人.
再一次來到他跌入谷底.
在監獄裡做監等.
聖經沒有說約瑟這一刻.
他心裡是怎樣.
他腳踏實地 硬著頭皮.
去接納他現在這個狀況.
他努力做好自己.
一樣 上帝幫助他.

$^{161}$他努力之餘 得到監獄長的賞識.
又好像之前一樣.
監獄長將所有的東西交給他.
所有的囚犯都交給他來打理.
又是一人之下 萬人之上.
經過一段日子之後.
他得到一個機會去和法老解夢.
順利將這個夢的解釋告訴約瑟.
正如法老說.
如果不是順利的靈在你裡面.
你沒有可能去解得通這個夢.
於是他又再次成為天之驕子.
好像過山車一樣.
一次 兩次 三次.
他再次成為一人之下 萬人之上.
被法老去差派.
去管理整個埃及的事務.
成為埃及的宰相.
跟著好像夢的說法一樣.
七年的瘋年 跟著就是七年的荒年.
饑荒不單止臨到埃及.
也是臨到巴勒斯坦整個地區.
他的哥哥 他們的糧食都沒有了.
於是在爸爸的祝福下.
他們要來到埃及來買糧食.
嘩 一眼當然認得出他們的哥哥.
仇人不見面 份外都是清明.
不過若失 我們要留意經文.
他不是馬上跟哥哥相認.
他也不是心裡的苦毒.
充分了他的頭腦.
所以我認為他早就沒有被這些苦毒囚禁.
他用各種方法去試探他的十個哥哥.
究竟是不是從以前的錯誤去改正過來.
當他知道自己的哥哥.
很懊悔之前所做過一切的事情.
並且很悉心地去照料.
虐失同父異母的弟弟和很孝順的爸爸.
於是就跟他們相認.
跟哥哥相認.

$^{201}$嘩 哥哥當然很害怕.
以往自己做在自己弟弟身上.
是一件多麼不對的事.
很害怕他會懲罰自己.
不過這句說話.
以下的說話是虐失自己說的.
「你們不要因為以前做在我身上.
去自怨自恨.
你們原本的心是不好的.
但卻是上帝藉著這件事.
要拯救我們全個家族」.
虐失去了一件不好的事.
來體會一個正面的祝福.
一定是上帝將他這個痛苦已經轉化了.
並且鼓勵哥哥回到他們原本的地方.
回到他們的家.
去接爸爸來.
接所有的家人來到埃及去享福.
因為還有幾年的饑荒.
寬恕不代表我們能夠抹除記憶.
不過卻是能夠將痛苦去轉化.
第21節告訴我們.
當門徒看到耶穌的傷痕的時候.
他們就喜樂了.
這個傷痕證明了.
耶穌真真實實從死亡中復活.
死亡不能困住耶穌.
痛苦不能毀滅祂.
反而成為祂走上榮耀的道徒.
這個也告訴我們.
在基督裡面.
我們過去被傷害過的經歷.
是可以從這個致命傷.
轉化為醫治自己和其他人的良藥.
我的老師霍元年姐妹寫過一本書.
這本書說.
「饒恕果真如此容易」.
饒恕真的這麼容易?.
是不容易.
是很不容易的事.

$^{241}$不過絕對有可能.
饒恕不是一句對不起.
對不起就了事.
不是.
當中包含了真誠的道歉.
願意努力改變自己錯誤的行為.
對曾經傷害過的人作出賠償.
和自己的生命要有更生變化.
才可以走上關係的復活.
走入藥室.
知道哥哥他們不單有悔意.
並且已經改正了自己的錯誤行為.
於是他相認.
他們的關係就開始走上復活的道路.
約瑟一早就已經在他心裡啟動了寬恕.
所以你曾經的傷害是不會繼續去苦害他.
耶穌能夠醫治是因為他曾經受過.
最深之痛的苦楚.
當我們願意將傷痕交給他的時候.
又願意踏出寬恕的步伐.
第一個被釋放的竟然是我們自己.
聖靈會幫助我們走過這段寬恕的旅程.
也會幫助我們讓那段痛苦的過去.
轉化成為能夠安慰.
能夠鼓勵其他人生命的見證.
耶穌在第十三節這樣說.
「你們赦免誰的罪 誰的罪就得到赦免」.
「你們不赦免誰的罪 誰的罪就不得到赦免」.
《新標的無後本》的翻譯是.
「你們赦免誰的罪 誰的罪就赦免了」.
「你們留下誰的罪 誰的罪就留下了」.
這是一個警示.
當我們抓住別人的錯不放.
當我們拒絕寬恕對方的時候.
其實我們是將那份仇恨和罪惡.
留在我們自己的生命當中.
成為了苦毒殘害我們自己.
不願意寬恕傷害我們的人.
到頭來竟然是囚禁了我們自己.
我們的心靈會被那個傷害我們的人.

$^{281}$扣押在牢籠當中.
不想想 又無端端會飄出來.
不想碰見他 又碰到正.
令到我們苦不堪言.
怎麼辦呢?.
寬恕那個這麼壞的人.
豈不是便宜了他?.
請聽我說.
這真是生命的實況.
當我們願意踏上寬恕的旅程的時候.
第一個得益的.
得到釋放的是我們自己.
然後當對方願意真誠的道歉.
願意努力去改變他那個錯誤的行為.
願意對曾經傷害過我們的傷害.
作出合理的賠償.
我們見到他生命有更新的轉化.
整個的饒恕旅程才算是完成.
然後我們可以走上復和的關係.
走入藥室和他這麼多位的哥哥.
我的另一位老師陳全華牧師.
寫了一本書叫「真饒恕」.
有個「真」字在前面也有「假」.
「真饒恕」.
當中的見解非常到位.
所以鼓勵大家可以買來看.
在第21至22節.
耶穌吹了一口氣.
他賜下聖靈給門徒.
猜怕他們赦免其他人.
特別是那些傷害過我們的人.
我們的仇敵.
基督徒在世界上不是爭論誰對誰錯.
而是要成為和平的使者.
要在我們的家庭.
我們的職場.
我們的社區.
我們的世界.
特別是在教會當中.
去締造和平.

$^{321}$今天我們有21位加入教會聖禮的弟兄姊妹.
所以除了梁法太太的弟兄姊妹.
你是被上帝差派成為和平使者.
大家都一樣是.
我們是上主和平的使者.
你願意接受這份任務嗎?.
不過你不願意.
也要接.
因為你是基督的門徒.
你是信耶穌的.
和平不是姑息罪惡.
維持表面的和諧.
你好我好.
私底下暗箭亂飛.
多麼虛偽.
這不是和平.
聖經說愛是喜歡真理.
讓對方知道自己錯了.
傷害了自己.
其實也是一份愛.
當然.
在我們告訴對方.
他做錯了甚麼.
以及傷害了自己甚麼事之前.
我們要讓對方感受和體會到.
你是愛他.
你才好說誠實話.
這個聖經是教導.
用愛心說誠實話.
的確要感受到你的愛.
你才好說誠實話.
否則的話.
只會變成向對方發洩我們的義憤和怒氣.
只會傷上加傷.
所以我們要謹慎行事.
求上帝幫助.
寬恕不是軟弱.
而是從基督領受了赦免的權柄.
主動選擇去放手.
讓自己和對方都得到真正的自由.

$^{361}$更讓破裂了的關係.
有機會得到修復.
邁向真正的復和.
容易嗎?.
當然不容易.
而且要看他對你的傷害有多深.
所以千萬不要隨便跟旁人說.
你寬恕他.
真的不容易.
不過你要為他祈禱.
所以耶穌向我們吹一口氣.
要賜聖靈給我們.
成為我們的教練.
我們的導師.
我們的鼓勵者和安慰者.
在不容易的寬恕旅程之中.
有勇氣向我們傷害過的人尋求寬恕.
也願意冒險踏出一步.
去寬恕那些曾經傷害過我們的人.
各位靈姐妹 各位來賓.
如果你心裡有掙扎.
請你來找我.
我一定會與你同行.
漢啊.
耶穌正站在你的羞愧的門之外.
對你說「願你平安」.
你摸一摸耶穌的傷痕.
祂正向你發光.
耶穌要醫治你最深的痛.
我們一起走.
去領受聖靈的氣息.
向上你說.
我願意寬恕那個.
我們以為永遠都無法寬恕的人.
因為聖靈一定會幫助我們.
各位靈姐妹 各位來賓.
請你在心裡去想.
我們當中.
有誰願意放下心裡的那把枷鎖呢?.
有誰願意去領受.

$^{401}$使人和平的使命呢?.
記住.
我們決定去相信耶穌的時候.
我們的立志就是要進行祂對我們的吩咐.
願上帝感動我們.
幫助我們.
以下我要帶領大家祈禱.
請大家在心裡記下那位傷害過你的人的名字.
又或者你曾經傷害過的人的名字.
將這個人帶到上帝的面前.
讓我們一起祈禱.
主耶穌 托成你帶著你的傷痕而來.
我們求你醫治我的過去.
使我不再被仇恨,內疚所囚禁.
而是帶著你的平安走出這道門.
去祝福這個撕裂了的世界.
上帝我們知道每個痛苦都有它的意義.
所以求你向我們啟示.
我們曾經所受過的傷害對我們有什麼意思.
我們更求你轉化這些深深的傷害和痛苦.
成為祝福我們自己和祝福其他人的傷疤.
天父請你幫助我.
我們能夠向傷害過的人去真誠道歉.
努力去改正自己的錯誤.
願意作出賠償和在聖靈的導引之下.
去更新自己的生命.
上主啊.
求你賜給我一顆願意寬恕的心.
按著聖靈的導引.
走上寬恕這位曾經深深傷害過我的人.
以致我們同得真正的釋放.
我這樣禱告 藉主耶穌基督的名而求.
阿們.
\newpage



\section{使徒行傳 2:1-21}
\label{sec:032178}
\textbf{我們今天一同得這能力吧!}~~~~高國雄牧師
\newline
\newline
連結: \href{http://www.church.com.hk/acms/content.asp?site=cdc\&op=show\&type=product\&code=032178\&layout=sermon}{\texttt{032178}} ~~~~ 語音日期: 2026-05-24
\newline
\newline
\hyperref[sec:032177]{\small{< < < PREV SERMON < < <}}
~
\hyperref[sec:index_chronic]{\small{[返順時目]}}
~
\hyperref[sec:032179]{\small{> > > NEXT SERMON > > >}}
\newline
\newline
使徒行傳 2:1-21
\newline
\begin{longtable}{cl}
\hline
\hline
章節 & 經文 (和合本修訂版)\\
\hline
2:1 & \begin{tabularx}{0.7\textwidth}{X} 五旬節那日到了，他們全都聚集在一起。 \end{tabularx} \\ \\ \relax
2:2 & \begin{tabularx}{0.7\textwidth}{X} 忽然，有響聲從天上下來，好像一陣大風吹過，充滿了他們所坐的整座屋子； \end{tabularx} \\ \\ \relax
2:3 & \begin{tabularx}{0.7\textwidth}{X} 又有舌頭如火焰向他們顯現，分開落在他們每個人身上。 \end{tabularx} \\ \\ \relax
2:4 & \begin{tabularx}{0.7\textwidth}{X} 他們都被聖靈充滿，就按著聖靈所賜的口才說起別國的話來。 \end{tabularx} \\ \\ \relax
2:5 & \begin{tabularx}{0.7\textwidth}{X} 那時，有從天下各國來的虔誠的猶太人，住在耶路撒冷。 \end{tabularx} \\ \\ \relax
2:6 & \begin{tabularx}{0.7\textwidth}{X} 這聲音一響，許多人都來聚集，各人因為聽見門徒用他們各自的鄉談說話，就甚納悶， \end{tabularx} \\ \\ \relax
2:7 & \begin{tabularx}{0.7\textwidth}{X} 都詫異驚奇說：「看哪，這些說話的不都是加利利人嗎？ \end{tabularx} \\ \\ \relax
2:8 & \begin{tabularx}{0.7\textwidth}{X} 我們每個人怎麼聽見他們說我們生來所用的鄉談呢？ \end{tabularx} \\ \\ \relax
2:9 & \begin{tabularx}{0.7\textwidth}{X} 我們帕提亞人、瑪代人、以攔人，和住在美索不達米亞、猶太、加帕多家、本都、亞細亞、 \end{tabularx} \\ \\ \relax
2:10 & \begin{tabularx}{0.7\textwidth}{X} 弗呂家、旁非利亞、埃及的人，並靠近古利奈的利比亞一帶地方的人，僑居的羅馬人， \end{tabularx} \\ \\ \relax
2:11 & \begin{tabularx}{0.7\textwidth}{X} 包括猶太人和皈依猶太教的人， 克里特人和阿拉伯人，都聽見他們用我們的鄉談講論神的大作為。」 \end{tabularx} \\ \\ \relax
2:12 & \begin{tabularx}{0.7\textwidth}{X} 眾人就都驚奇困惑，彼此說：「這是甚麼意思呢？」 \end{tabularx} \\ \\ \relax
2:13 & \begin{tabularx}{0.7\textwidth}{X} 還有人譏誚，說：「他們是灌滿了新酒吧！」 \end{tabularx} \\ \\ \relax
2:14 & \begin{tabularx}{0.7\textwidth}{X} 彼得和十一個使徒站起來，他就高聲向眾人說：「猶太人和所有住在耶路撒冷的人哪，這件事你們要知道，要側耳聽我的話。 \end{tabularx} \\ \\ \relax
2:15 & \begin{tabularx}{0.7\textwidth}{X} 這些人並不像你們所想的喝醉了，因為現在才早晨九點鐘。 \end{tabularx} \\ \\ \relax
2:16 & \begin{tabularx}{0.7\textwidth}{X} 這正是藉著先知約珥所說的： \end{tabularx} \\ \\ \relax
2:17 & \begin{tabularx}{0.7\textwidth}{X} 『神說：在末後的日子，我要將我的靈澆灌凡血肉之軀的。你們的兒女要說預言；你們的少年要見異象；你們的老人要做異夢。 \end{tabularx} \\ \\ \relax
2:18 & \begin{tabularx}{0.7\textwidth}{X} 在那些日子，我要把我的靈澆灌，甚至給我的僕人和婢女，他們要說預言。 \end{tabularx} \\ \\ \relax
2:19 & \begin{tabularx}{0.7\textwidth}{X} 在天上，我要顯出奇事，在地下，我要顯出神蹟，有血，有火，有煙霧。 \end{tabularx} \\ \\ \relax
2:20 & \begin{tabularx}{0.7\textwidth}{X} 太陽要變為黑暗，月亮要變為血，這都在主大而光榮的日子未到以前。 \end{tabularx} \\ \\ \relax
2:21 & \begin{tabularx}{0.7\textwidth}{X} 那時，凡求告主名的都必得救。』 \end{tabularx} \\ \\
[1ex]
\hline
\hline
\end{longtable}
$^{1}$頂尖的平安.
請祈禱.
讓我們今天能夠回到電當中.
有這麼舒適的環境.
能夠去細心的聆聽.
開我們的耳朵.
讓我們專心.
開我們的眼睛能夠看見.
你的榮耀.
你的力量.
你的盼望.
更以我們的手腳行動.
回應聖靈你的提醒.
和你的愛.
就如部人口中的言語.
頂尖姊妹請深思以念.
蒙聖靈的啟迪.
蒙主的悅立.
禱告奉耶穌基督的名求.
阿們.
各位頂尖姊妹.
今天是聖靈降臨節.
你見到聖壇和木色的帶都換上紅色.
象徵著上帝的火在我們當中.
今天和大家分享的題目是環繞著能力的.
我想你們先想一想在這一刻.
你最想要的是什麼能力呢.
如果任你選擇.
什麼能力你最想擁有呢.
這一刻.
在現今世界裡面.
不同人在追求不同的能力.
學生可能追求的是.
學習的能力.
寫文的能力.
接著不要太睏.
通宵都可以做到功課的能力.
又或者甚至是那種.
不上課不做功課.
拿A的能力.

$^{41}$這都不是能力.
學生就追求學習的能力.
上班的職場就追求競爭力.
我們升職.
加薪找工作.
找到自己合適的崗位.
國家就追求經濟.
軍事的能力等等.
而我們環顧我們的世界.
我們發現很多的戰爭.
氣候的變遷.
或者在我們個人生命裡面的憂慮.
那些破碎無力感的時候.
我們就會再問.
那些能力是不是真的有用呢.
真的可以不讀書.
拿A是不是真的有用呢.
真的可以不斷工作.
效率越來越提高.
是不是真的很有意思呢.
國家有最大的經濟和軍事的力量.
那又怎樣呢.
聖靈降臨節是很特別的日子.
聖靈降臨節不只是一個.
歷史的宗教節日.
這麼簡單.
聖靈降臨節是告訴我們.
聖靈降臨了.
是一個得著能力的慶典.
一個充全能力的慶典.
所以今天我們有一個很特別的題目.
就是我們今天一同.
得著能力嗎.
希望你今天完了崇拜之後.
你得到一些不同的能力.
更重要這個能力真的幫助到.
你屬靈生命的成長.
因為聖靈降臨就是教會的開始.
我經常說的.
其實就是教會的生日.

$^{81}$今天就是我們慶祝我們這個群體.
要努力了.
首先我們看看今天的經課.
我們去到.
整個能力的來源.
是在福音書裡面.
當福音書裡面開始.
在福音剛才大家聽到的.
二十章裡面.
講到復活之後.
耶穌就在門徒當中顯現.
我經常都問這個問題.
耶穌為什麼要顯現呢.
復活然後升天不就行了.
耶穌的顯現是要讓門徒.
和那些見到顯現的人.
看見一些歷史歷代.
都沒有機會看見的東西.
他們是獨獨的單獨的.
去看見復活的主的力量.
所以當耶穌.
在那一天復活的日子.
七天的頭一天晚上.
門徒所在的地方.
門徒當時很害怕.
因為很害怕猶太人.
因為耶穌被捉拿.
很多門徒都四散.
最親近耶穌的門徒的首領.
很多人躲起來.
門關上了.
門關上了之後.
耶穌就站在他們當中.
跟他們說願你們平安.
整個的狀況.
就是當時門徒.
很怕猶太人.
他們的人生失去了焦點.
失去了重心.
前路茫茫不知怎麼辦.

$^{121}$因為過去幾年他們追隨的.
他們跟從的這位耶穌.
被捉拿被釘甚至已經死了.
他們追求的價值.
可能過去幾年很努力.
很想跟別人分享的價值.
這一刻好像沒有價值一樣.
因為耶穌已經不在他們當中.
這個可能也是我們.
有一段時間生命中的寫照.
希望不是這一刻的寫照.
因為有時我們都會關上我們的心門.
不想被外界再傷害我們.
關在自己的恐懼當中.
躲起來算了.
耶穌見到他們第一句.
就是說願你們平安.
接著很重要.
第22節說.
說完這話之後就向他們吹噓.
說你們受聖靈.
我們每次想起聖靈降臨.
就一定想起.
少林寺很轟烈的.
舌頭的火.
臨到當中辣著.
很長的經文.
但是原來在聖經裡面.
記載門徒受聖靈.
就不是要等到.
聖靈降臨那天之前.
已經和他們在一起.
當然只是在房間裡的那群門徒.
而這種受聖靈就是.
耶穌向他們吹噓你們受聖靈.
沒有很轟烈.
沒有突然有火有風.
很轟烈的畫面.
沒有.
不過耶穌將這份勇氣.

$^{161}$信心灌注了.
門徒當中.
在聖經裡面說這個零.
愚民的字眼.
和氣息呼吸的字眼一樣.
風的字眼一樣.
當舊約一開始創世的時候.
上帝創造人.
創造亞當的時候.
他就將靈氣吹在亞當的鼻孔裡面.
聖經裡面說.
亞當就成為了有靈的活人.
原來人的開始.
是由受上帝的靈開始的.
你有沒有記得這件事.
人為什麼會認為人和其他動物不同.
就是因為有靈.
這個靈就是上帝的氣.
上帝將他的氣息.
吹入亞當的鼻孔.
耶穌再一次.
做這個行動.
將這句說話.
化成一個氣息.
向門徒吹氣.
門徒就受了聖靈.
耶穌在這一刻.
就是一個新的創造.
將這個力量帶給這一群.
很害怕躲在.
很懼怕的門徒當中.
記住留意.
這個畫面沒有很轟烈.
風吹雨打很厲害.
門徒仍然很害怕.
他們都繼續躲起來.
多瑪在下星期發生的.
看到耶穌復活的事.
就是他們繼續躲起來.
去到聖靈降臨很光亮才有更大的勇氣.

$^{201}$不過這一刻.
聖靈已經透過耶穌.
進入門徒當中.
他們已經得到一份平安.
以致他們夠有勇氣繼續生活.
可能有時候也會害怕.
頂至我們領受了聖靈.
有些人可能還是.
好像沒有什麼變的外形.
沒有什麼特別的情況.
但其實內心已經得到主的平安.
好了.
過了五十天.
去到今天.
去到聖靈降臨的日子.
這份平安慢慢轉化成為很大的力量.
剛才大家聽到的.
《使徒行傳》第二章.
二章一到四節就是講那件事.
再一次大家想像一下當日的畫面.
在五旬節的日子.
門徒聚集在一起.
忽然天上有響聲.
好像大風的吹過.
充滿了他們坐在那間屋.
接著.
蝎頭好像火一樣顯現.
落在他們身上.
他們就被聖靈充滿.
按照聖靈所賜的口才.
講起別國的話.
這裡在講門徒受聖靈.
那個畫面.
有大風.
有蝎頭.
聖靈說有如蝎頭的火.
即是聖靈是蝎頭.
聖靈是蝎頭.
好像火一樣燒著.
接著門徒就講方言.

$^{241}$今時今日.
不同教會有不同的表達.
講方言都有不同的理解.
但在第一代.
一開始門徒領受方言.
其實只有一個很重要的目的.
就是宣講上帝的偉大.
因為聖經裡面講到.
有一群不同的人.
很多地方.
其實當時很多人聚集在耶路撒冷.
在那個地方.
因為他們來自世界各地.
有不同的朋友.
有些是海外回來的.
可能是一些猶太人.
但去了海外生活.
因為他們有不同的語言.
不同的鄉談.
好像一班人走在一起.
但有不同的語言.
我們現在認識的世界各地的語言.
或者我們在國內有不同的鄉口話.
這是一個聚集.
而當時聖靈的第一個工作.
一個很重要的工作.
當門徒被聖靈充滿之後.
他們就能夠得到那份能力.
就是一份溝通的能力.
就是一份復和的能力.
讓那些語言不通的人.
來自世界各地不同背景.
不同膚色.
不同方言的人.
一起去見證上帝的作為.
其實正正就是對照.
我們熟悉的創世記巴別塔故事.
上帝就變了口音.
讓人分散.
但這一刻不同的口音.

$^{281}$雖然繼續有不同的口音.
但透過聖靈不同的鄉談.
不同的口音.
他們都能夠去明白體會上帝的話語.
這就是聖靈.
其中一個很重要的能力.
丙姐妹當日.
聖靈就透過這班門徒.
給了他們和不同人溝通的力量.
包括和那些.
他們看不起的人.
和他們完全不認識的人.
因為聖經裡面跟著.
《少林主義》二章十七到二十一節就跟著說.
裡面用舊約耳書的語言說.
我要將我的靈囂冠.
煩有血氣.
裡面說了三組不同的人.
一組叫少年人.
要見到異象.
老年人要做異夢.
那些僕人和使女.
都要被上帝的靈囂冠.
要說語言.
這是什麼意思呢.
這裡說的是.
那些少年人能夠有異夢.
他們找到他們的將來.
能夠望見將來.
聖靈其中一個很重要的力量就是讓門徒望見將來.
剛才不是說他們很害怕.
很害怕到要躲起來.
沒辦法怎麼辦呢.
給到力量他們有勇氣走出來.
我不怕猶太人的迫害.
那怎麼辦呢.
原來這裡說的是上帝讓他們看到將來.
看見上帝的國度要展現.
看見門徒要建立教會.
看見有更多人.

$^{321}$要看見盼望.
這是其中一個很重要的力量.
老年人要做異夢.
年紀大很多時候.
都習慣了.
各位電影姐妹你上了教會多久.
可能很長的日子當中.
我相信很多電影姐妹上了教會比我的日子更長.
可能每個禮拜崇拜的時候.
睡醒都要上教會.
坐下沒什麼期望.
聽完就去喝茶.
我們可能很習慣.
年紀越大越不講夢想.
夢想就是年輕人講的.
小學生中學生講我的志願.
你在長者中心.
做長者工作的電影姐妹你會不會跟他們講我的志願呢.
比較難.
其實是好的.
不過有人說講我的志願不如講我的遺願.
其實可能是同一件事.
志願可能都是遺願.
我們每一個人都走向上帝的方向.
不過我們經常覺得老人家就是這樣.
習慣了.
正常.
年紀大了不記得事情是正常.
年紀大了開始老花.
看不清楚.
正常.
年紀大了.
我們沒有期望過年紀大了的人.
老了的人仍然可以重拾對生命的熱誠.
這裡講老人家都要有義盟.
要有意象.
一定繼續有方向找到熱誠.
找到生命要追求的事情.
另外在講侍女和僕人都要被上帝的靈囂冠.
要講寓言.

$^{361}$因為打破了當時一些階級的枷鎖.
僕人和侍女當時來講.
他們看不起的人.
或者在身邊不礙眼甚至覺得是財產.
多過覺得是人的人.
但上帝都看重這些人.
各位弟兄姊妹.
聖靈的能力不只是給少數人.
聖靈的能力不只是給牧師.
傳道.
給一些領袖.
讀書厲害的人.
不是.
上帝的靈是給每一個渴慕上帝的人.
每一個願意追求上帝.
渴慕上帝.
希望得到希望熱誠.
並且能夠打破隔膜枷鎖的人.
聖靈的能力不是讓我們.
自我膨脹覺得多好多厲害.
是要讓我們有勇氣.
跨出我們的舒適圈.
去宣告.
耶穌基督的救恩.
這個才是聖靈給我們的能力.
就是靠著祂的能力.
能夠成為更多人的祝福.
去到經課.
哥林多前書今天提到的.
除了在講聖靈裡面.
好像小彭傳那種爆發力.
很多人信耶穌很偉大之外.
也都看見聖靈給到.
這班基督徒群體教會的.
持續的能力.
哥林多前書十二章四到七節.
聖經裡面剛才講過聽過.
因此有分別聖靈祂是一位.
即是有分別主祂是一位.
功能有分別上帝祂是一位.

$^{401}$聖靈在人裡面運行的一切的事.
聖靈顯在個人身上.
要叫人得益處.
少爺保羅透過.
聖經裡面提醒我們.
因此有很多的.
但聖靈是一位.
因此有不同有即是不同功能不同.
但是目標只有一個.
就是叫人得益處.
弟兄姊妹.
我們很多時候想起.
聖靈的能力信仰的能力.
中國的能力真的只是想起那些領袖.
今天我站在港台宣告.
宣講和分享.
你可能只是坐在那裡聽.
但其實我們每一個人在每一個位置裡面.
都能夠發揮.
少爺保羅很清楚地提醒.
哥林多教會因為哥林多教會覺得自己挺厲害.
所以保羅想提醒他.
是你們厲害不過都要各國去配合.
如果整個身體都是眼.
那誰去做事.
整個身體都是手那如何合作.
其實就是在提醒我們.
每人有各人有不同的角色.
你在上帝的國裡面是甚麼角色.
你在教會當中是甚麼角色.
你如何在你的角色裡面.
發揮上帝給你的恩賜能力.
而能夠讓更多人.
有益處呢.
這個才是我們追求的.
不是在想我們做得多好.
我今天要和大家分享訊息.
所以講到要預備多精彩的講章.
而我在想上帝擺我這個位置.
以訊息作為服侍的時候.

$^{441}$我如何發揮我的恩賜.
可以讓大家都體會明白.
我們都要一起有力量.
去建立上帝的國.
你在你的位置裡面.
又如何發揮上帝給你的恩賜.
當每一個人.
每一個人在我們的教會當中群體.
都能夠按著上帝給我們的恩賜.
去發揮的時候.
就能夠有持續力地繼續去見證.
聖靈一臨在.
聖靈降臨在.
當日是一日的爆發.
之後沒有每日都不斷的火啊.
蝕頭啊風啊.
不過繼續爆發的力量.
就是讓門徒各自將他們的恩賜去發揮.
所以教會就一直這樣發展.
一直這樣向世界各地去拓展.
今天我們呢.
我們在當中又在經歷什麼呢.
我們有時.
可能會誤會.
以為有些能力就是.
我做到別人做不到的事.
這些就是能力.
因為我做一些事別人做不到.
但是聖靈的能力在教會運作的時候.
是讓我們看見.
自己不足.
和看見他人的價值.
聖經裡面提到.
我們不區是猶太人希臘人.
自主的都是受一個聖靈成為一個身體.
讀起來很簡單.
但是對當時初期教會來說.
是一個震撼的描述.
這個震撼一點都不低過.
聖靈降臨那種畫面的震撼.

$^{481}$因為當時對他們來說.
猶太人就是猶太人.
猶太人以外的就是不是猶太人.
世界上只有兩種人.
但是.
竟然斯圖普羅透過聖靈的教導說.
就算是猶太人.
希臘人.
甚至是那些為奴的.
作為奴隸的.
和那些主人都是受同一個聖靈.
這個很震撼.
因為對很多人來說聖靈不是我們.
甚至教會頂尖的都會想.
上帝只是向教會說話.
上帝只是感動教會屬上帝的人.
我是基督徒.
上帝不應該向我多說一些話.
不是信徒的人.
那些未認識上帝的人.
上帝不會跟他們說話.
他們都是受同一個聖靈.
同一個身體.
所以很重要.
因此我們要互相配搭.
在教會裡面要配搭.
在教會外面都同樣要配搭.
以致我們發現能力目的不是為了競爭.
這個世代有很多能力是為了競爭.
我要跑贏旁邊那個去展示我的能力.
並且得到應有的好處.
能力是為了互相補足.
上帝的能力.
聖靈的能力在教會當中.
是要讓我們互相補足.
上帝記住這句說話.
有些事情是我的強項.
有些事情是你的強項.
所以你的強項可以補足我的弱點.
我的強項也可以讓你的弱點被補足.

$^{521}$因為這是上帝所賜的.
記住.
我們互相補足這件事.
在教會裡面是要發生的.
我常常跟上帝說.
我年紀大了很多事都做不到.
有些人做到.
你可以服侍很好我就沒有機會了.
但我常常跟頂尖說.
你可以在你的位置去盡忠.
上帝給你有的可能是退休有時間.
請你每天為牧者祈禱.
每天為教會禱告.
你說我沒有很多能力.
做很多事不會說不要緊.
你就安定地去接受並且去跟身邊的人見證分享.
每一個人都有他的角色.
互補才能夠讓教會去成長.
在詩篇第104篇.
大家聽到啟營的經文裡面.
是這樣說的.
當上帝對著.
這些惡如怪獸.
上帝的管理時.
聖經裡面說.
你張手他們就得未食.
你掩面他們就驚惶.
你收回他們的氣他們就死亡.
歸於塵土.
這個描述很有趣.
這個描述是在說上帝的主權.
也在說上帝的靈上帝的氣.
除了向人給我們有靈.
成為有靈的活人之外.
在萬物當中都有上帝的氣息.
聖靈不只是在教會當中.
聖靈是要運行在.
整個的傳遞.
包括我們.
可能知道.

$^{561}$他們未認識上主的朋友.
又或者包括世界的萬物.
宇宙的萬物.
上帝的靈都在裡面.
祂是讓平衡出現.
祂是讓生態的平衡出現的上帝.
祂是讓我們在當中體會.
祂是我們看顧的上帝.
所以詩篇裡面提醒我們.
上帝不只是向人.
去將祂的氣息.
去擺上.
是動物甚至我們不知道什麼名的動物.
上帝仍然看顧.
亦都在祂的掌權裡面.
所以上帝的氣息同樣地告訴我們上帝的掌權.
頂尖的今天我們有.
這個機會在這裡繼續有氣息去敬拜.
事奉其實是上帝的恩典.
亦都是上帝的靈在當中.
能力的彰顯.
今天我們和大家思考一個題目.
叫做我們一起去.
領受這個能力.
我不知道你聽到剛才說的.
不同的經課有什麼體會.
可能我們仍然是.
在軟弱當中第一步.
我們要去承認我們的軟弱難處.
好像門徒一樣他很害怕.
他們就去到主人面前承認.
上帝就臨在他們當中.
我們可以向主人呼求.
求上帝吹氣.
給他的靈在我們當中.
是我們的生命的力量.
我們更要求.
上帝幫助教會.
我們能夠放下我們的成見.
互相有支援服務.

$^{601}$能夠去互相合一.
今天在這個.
聖靈降臨節當中.
我想起.
很想和大家去分享.
聖靈降臨.
每一年我都會在這裡.
有機會去舉辦音樂會.
為什麼呢.
因為以前沒有.
為什麼就要解釋.
第二就是因為我覺得.
在這個禮拜堂裡面.
音樂是一個很重要的元素.
更重要的就是聖靈降臨節.
是最好機會和人分享的日子.
每一年教會有幾個日子.
很值得分享的.
我們一般只分享聖誕節.
所以平安夜我們都會有很多人回來.
每一年聖誕節我們都很喜歡分享.
復活節都可以的因為都有長假期.
我們可以有很多的分享.
聖靈降臨節就好像沒什麼特別的.
又沒有假期.
明天公假不是因為聖靈降臨節.
稍後報告禱告都有分享的事.
聖靈降臨節就好像很靜悄悄.
在教會發生了.
因為我們從來沒有想過.
這是什麼特別的節日.
有些教會甚至不說.
復活聖誕就夠了.
每一年聖靈降臨節.
我都是再一次被反省.
因為我經常都想聖靈降臨.
就是在說上帝的能力.
聖靈降臨節就是教會的開始.
我在教會裡面有什麼.
的遺憾令到教會繼續去發揮.

$^{641}$以致能夠吸引更多的朋友.
在我們當中一起去.
經歷上帝的豐富呢.
在今天可能當中有參與教會很久的.
弟兄姊妹.
你怎樣在當中去體會經歷.
上帝的能力呢.
可能有第一次參與教會.
甚至可能你還在信仰裡面去考慮的朋友.
你有沒有想過.
今天就是你經歷聖靈.
進入你生命當中.
改變你能夠成為你生命的主.
成為你的救主的開始呢.
聖靈的能力是什麼能力呢.
是溝通 是關懷.
是互相支援.
是建立教會.
聖靈的能力也告訴我們.
我們可以放膽去擁抱.
去信任.
這位張開手的基督.
在十字架上張開手的基督.
他是我們的十字架.
同樣他好像張開手等待我們的擁抱一樣.
請祈禱.
讓我們堅持.
讓我們能夠好好地服侍.
有力量地面對我們懼怕的事情.
有勇氣地向前.
向身邊的人去分享.
更有溝通的能力.
以致我們能夠和不同的人.
甚至那些我們從來沒有想像過.
我們能夠和他分享信仰的人.
去分享信仰.
主一求你擁抱我們.
讓我們生命.
那份力量.
能夠讓我們堅持.

$^{681}$繼續向前.
勇往直前.
無論路有幾窄門有幾窄.
我們就努力.
更讓我們並肩成為教會.
一起去見證你的榮美.
願你的能力在我們當中.
願你恩代賜福.
讓我們一起去承受.
你所賜的恩典,力量,勇氣.
並且能夠去見證.
榮耀你的名.
多謝主聽我們禱告.
\newpage



\section{約書亞記 1:8}
\label{sec:032179}
\textbf{一切從神的話開始}~~~~王偉德傳道
\newline
\newline
連結: \href{http://www.church.com.hk/acms/content.asp?site=cdc\&op=show\&type=product\&code=032179\&layout=sermon}{\texttt{032179}} ~~~~ 語音日期: 2026-05-24
\newline
\newline
\hyperref[sec:032178]{\small{< < < PREV SERMON < < <}}
~
\hyperref[sec:index_chronic]{\small{[返順時目]}}
~
\hyperref[sec:032180]{\small{> > > NEXT SERMON > > >}}
\newline
\newline
約書亞記 1:8
\newline
\begin{longtable}{cl}
\hline
\hline
章節 & 經文 (和合本修訂版)\\
\hline
1:8 & \begin{tabularx}{0.7\textwidth}{X} 這律法書不可離開你的口，總要晝夜思想，好使你謹守遵行這書上所寫的一切話。如此，你的道路就可以亨通，凡事順利。 \end{tabularx} \\ \\
[1ex]
\hline
\hline
\end{longtable}
$^{1}$靖姐妹早晨 早晨.
上一次見大家是半年前多一些.
今天再一次和大家分享神的話語 很感恩.
多謝主席和敬拜隊來帶領我們敬拜.
一起用詩歌,用禱告向神獻出讚美.
也多謝孫震環基督教會和陳全道的邀請.
可以今天和大家分享神的話語.
可以在這裡一起說.
聽神的話語是很美好的一件事情.
是很重要的事情.
所以今天和大家分享的宣講主題.
叫做「一切從神的話語開始」.
首先我們一起祈禱.
來對公敬將我們的心交託.
讓我們有聆聽主導的心.
對準上帝 一起祈禱.
多謝天父讓我們一起在主日裡面.
同時是聖年多倫節的主日.
也同時是教會洗禮的一個主日.
也是我們一群頂智媒每個禮拜.
奉你的名聚會的一個主日.
我們來到神你的家 神你的殿.
來敬拜你.
求主你閱立我們的敬拜.
也求你預備我們的心.
求你預備我們的腦去聆聽,聽得明.
預備我們的心來聽得進去.
來放在心中 因為一生過後由心發出.
也求你讓我們能夠活出你的道.
讓我們一生行在神你的旨意當中.
一切從神的話語開始.
弟兄姊妹.
無論我們在當中已經信了主幾十年了.
我自己也信了主四十年了.
當中有些可能信了主五十年了.
又或者那些十年豐富的資深基督徒.
或者在這五年之間.
這五年其實也挺特別的.
在英國見到一些華人或者香港人.
這五年也是一個挺大的分水嶺.

$^{41}$特別是經歷很多的生活.
適應,挑戰,恩典都有的.
又或者是今天來到水禮禮堂的中央.
見證他們的生命 交託給主.
還有加入教會.
我想問一下 也想關心有沒有今天會受浸,受洗的.
請在節目當中歡迎你.
有沒有?今天也有回來.
然後1點45分就有水禮 是不是?.
我們也聽清楚了.
無論在哪裡 全部人都是對的.
因為剛才說到五十年,這五年.
或者你是剛剛水禮.
但是也可能信了很久 有一個新的開始.
我想宣告一件事.
無論我們的生命是怎樣不同的階段.
無論是我們加起來.
可能我們面對很多日新月異的科技.
又或者在政治環境.
或者身邊有一些不安或者戰亂.
正如剛才敬拜帶領我們去思考.
有很多戰亂,很多不安.
但是也都無論我們信仰的年資有多深,有多長.
或者短 或者可能我們在人生下半場.
又或者可能純粹你是信仰有一個新的一頁.
到底我們會怎樣走下去呢?.
我想今天藉著一個真理跟大家分享.
就是無論你信主多久 長或短.
或者你的環境 我們的環境是怎樣變遷.
人生在不同的階段.
我想說一切是從神的話語開始.
一切從神的話語開始.
首先我們會看神的話語.
要先從我們能夠思想神的話語開始.
我們要得著這個屬靈的真的智慧.
得著這個智慧 是真的智慧.
所以我們一起看今天第一段經文.
一起讀一次好嗎?.
預備起.
「這律法書不可離開你的口.

$^{81}$總要晝夜思想.
好使你謹守遵行這書上所寫的一切話.
如此道路就可以亨通.
凡事順利」.
「約書亞記一章八節」.
這是在舊約.
約書亞進入迦南美地 神的應許.
上帝先吩咐約書亞.
大家很熟悉吧?.
就是不是吩咐很多細節的東西.
我們譬如來到一個新地方.
有一個人生的階段轉變.
可能我們要想的東西.
未必是這段經文提醒我們的內容.
可能我們要關心的是一些經濟上的安排.
住屋的安排.
或者學校的安排.
很多東西需要安排.
但不是的.
當約書亞進入迦南的時候.
上帝吩咐第一件事.
祂叫他「這律法書不可離開你的口.
總要晝夜思想.
好使你謹守遵行書上所寫的一切話.
如此道路就可以亨通.
凡事順利」.
我們本末倒置有時.
就是我們先想亨通.
然後凡事順利了.
我們就有時間晝夜思想.
以及謹守遵行神的話.
也說不定的.
當然這是互相提醒.
我也有時是這樣.
今天聖經提醒我們.
就是先晝夜思想.
不過我想退後一步.
會不會大家都很想那種亨通和順利?.
大家想不想?.
早前也有很多反應也挺多的.

$^{121}$亨通,順利.
我們也很喜歡亨通和順利.
我記得甚麼叫不亨通呢?.
有一次我想去Pottyshed.
由Bristol去探望一個家庭.
已經談好了.
飯他們也煮好了.
我的車去到那條橋.
那條橋要過一個短的河.
沒多久我發覺交通停了.
然後那個公路牌有個交叉紅色的.
然後見到警察車在旁邊飛過.
你估計應該會發生甚麼事呢?.
然後我在車上就這樣停留了.
你估計一個小時還是四個小時還是八個小時.
這裡的安排和我們過往的經歷不是很不同的.
它不會開一條路給你疏導的.
它會封鎖整條路.
因為有些well-being的情況.
即是有些情況發生在橋上.
所以他們不亨通.
我就很經歷到.
不亨通就是這個意思.
其中一樣東西很不亨通.
即是不順利.
於是過了一個多兩個小時.
我就打電話給或者聯絡那個家庭.
「我今天探不到你」.
「不好意思,我要再約時間」.
再過了一段時間.
我見到那些車開始掉轉頭.
我就跟著他們掉轉頭逆線.
然後接駁回去那條Motorway.
即是另外一個Motorway.
然後就離開了.
我們很經歷了不亨通.
不過聖經說的亨通順利.
到底是不是我們平時想的那種.
我照那個翻譯這樣說.
就是是不是發達,有錢,有權.

$^{161}$或者是贏過所有人呢?.
因為我看那個原文的時候.
他給我的那個字就叫做seco.
seco這個字其實那個亨通,成功.
就不是指剛才所說的那個解釋.
它更深的層次的意思就是.
留意我的電視節目.
就是有洞察力.
有智慧.
能夠分辨.
做事謹慎.
當一個人能夠明白神的心意之後.
能夠行得正和行得穩.
這個就是原文原本的意思.
不一定是我們原先想的那種亨通和順利.
那個屬靈的層次就是.
首先思想神的話.
然後明白神的心意.
然後行事有智慧.
然後人生就行得通達,行得穩.
真正的亨通就是.
生命和神是對齊的.
你的生命和神是對齊的.
然後就自然產生的結果.
就會亨通和順利.
那什麼叫做對齊呢?.
我們會繼續看下去.
是的.
在我們一直學習做基督徒的時候.
我們會發覺在思想神的話.
有時候我們會想.
會不會好像世界那種方法.
去講順利,講亨通呢?.
雖然我以前做科技的公司.
但其實現在在座的每一個人都很先進科技的.
很多都有很先進的手機在手.
我們都很講求資訊,數據,data.
algorithm 即是演算.
或者是要很優化那些東西.
要很厲害,要做得很強.

$^{201}$有AI那些東西.
我們以為要掌握這些技術或者是數據.
就能夠掌控那種成功.
但上帝和約書亞說.
他說重點不是你有多強就成功.
有多強就會有多順利.
不是有那種經濟上,科技上或者財力.
那幾樣東西就叫你人生很成功.
而是他說要每天都思想成的話.
這樣才會有真正的智慧和方向.
上帝感動我講這一段經文的時候.
其實是在提醒我.
「不可離開我的口,總要晝夜思想」.
特別是頂尖會留意「晝夜」這個字.
大家睡覺睡得好不好?.
The well-being.
我有時因為擔心.
或者有時世界戰亂.
或者有一個階段政治很多轉變.
都會令我們有些憂心.
睡得不太好的時候.
在每一個晚上.
那個「晝夜」.
這樣才是主的.
每一個晚上.
都有很多不同的經文提醒我.
那些經文從何而來呢?.
最常見的就是我們一起讀的主禱文.
我鼓勵頂尖智慧記得.
我鼓勵有多少個記得背得到主禱文?.
看到了 可以放下.
我在那些晚上經常都會背.
「我們在天上的父,爺,女,尊」.
總有一節.
有一句是會認到那個晚上的憂心.
那個憂慮或者是思慮.
或者那種肯定,那種安穩.
如果不記得主禱文.
我有時用耶穌禱文.
「主也憐憫我」.

$^{241}$可能就是這麼簡單.
在這一刻你都可以坐在位置上.
都可以禱告.
「主也憐憫我」.
其實就是那些時候.
所以那些「晝夜」思想成的華語.
就是「幫助到我」.
另外剛才主播有說.
有沒有心裡面思想成的華語.
你都要有經文思想才可以.
有沒有一些經文是記得的?.
可能是你從小到大.
我今年信主四十年 剛剛好.
我記得小時候.
那時候是teenage.
那時候有些三福,四律.
栽培那些百貨.
我們小時候就很多要上這些.
後來雖然做了三福的領隊.
但是我們都要背那些經文.
小時候就不喜歡.
現在長大了就很treasure.
很開心有這些經文要背過.
這些成為了放在這裡.
那些屬靈的寶庫.
放在思想的層面.
變成在那些晝夜的時候.
那些經文有時就會出來.
或者有時人生遇到一些處境.
在家庭,在工作,在社會,在政治.
有很多環境.
在教會,在超市.
遇到的弟兄姊妹.
遇到不信的,未信的家人,鄰居.
有時這些經文就在腦裡面浮出來.
弟兄姊妹這個就是當時.
上帝提醒若書.
你進入迦南地,你進入一個新的地方.
又或者你已經在這個地方幾十年.
仍然都要面對新的挑戰.

$^{281}$或者你剛剛這個時候.
今天受洗,加入教會.
認真的,我記得我受浸的時候.
我的分享見證其中說了一句.
是認真的.
小時候十幾歲的時候.
我說了一句這樣的話.
因為這個決定僅次於你決自信主.
也僅次於有一個禮儀.
你是沒有辦法知道你正在參加的.
那個是甚麼禮儀.
是很重要的禮儀,就是你的安息禮.
但浸禮就是你可以仍然很生活地.
去見證給別人.
所以思想神的話包括有很多經文.
你要記得那一句.
背經,小時候就背經.
很認真地背經.
弟兄姊妹,我們思想神的話之餘.
上帝提醒我們.
我們要遵行神的話.
有神的話之後,你要遵行.
這裡我寫了一個副題.
叫「沒有Google Map,只有腳前的燈」.
沒有Google Map.
腳前的燈是甚麼意思呢?.
神是不會給我們一個Postcode.
不會給我們一個地址.
你就由起點去到終點.
你就去到那個SN二十幾.
不會的.
神只會給你下一步的光.
在古代這段經文.
或者講之前我們一起讀了.
預備起.
「你的話是我腳前的燈.
是我路上的光」.
很清楚地在講腳前.
是我們路上.
在古時這段經文寫的時候.

$^{321}$背景其實我們走路的人.
只是靠一盞小油燈.
很小的油燈.
不像現在.
現在我們要甚麼光都可以有甚麼光.
車頭燈可以發明到雷射.
可以射到很遠.
智能的.
也有街燈.
但是古時只是一盞小油燈.
你可以想像黑夜中.
那盞燈只能夠照到腳前一個很小的範圍.
很小的.
小油燈.
我想暗過你手機那盞小燈.
它不會.
那意味著.
腳前幾步你只能夠看到很少.
你不能夠看到整條路.
也控制不了前面的環境.
也不知道轉角有甚麼危險.
你的安全也不夠.
你只能夠一步一步跟著光.
那個小油燈給你的光.
上帝的話語給我們的時候.
它不是一張人生的全景地圖.
上帝當然知道我們的心靈.
但是我們能夠知道的.
就不是一張全景的地圖.
不是一個Google Map.
那個比喻就是.
它不會直接給你一個postcode.
就算它給了你一個postcode.
就算你知道你將會去到那個地方.
然後再下一步怎麼辦?.
再下一步怎麼辦?.
向東還是南,西還是北走?.
很多東西都不知道.
信心不是看得見的全部.
而是相信此時此刻.

$^{361}$神所給你的下一點的光已經足夠了.
所以跟我們過往所想的那種.
剛才所說的經濟,科技,財力.
那幾樣最重要.
我們人生思慮繁榮了的地方.
是有些不同的.
或者應該說是整個根源性的不同.
是完全不同的.
我們英文叫做paradigm shift.
即是完全不同的一個思維.
我們在人生.
既然光是照亮我們前面.
我們天天靈修,天天讀經.
那應該都很清楚.
走下一步.
所以當我預備的時候.
有另外一個經文.
突然間正靈感動我.
提醒我要防備,弟子們.
亦都是跟光有關.
那就是一個叫做恨光的心態.
恨不要,不喜歡光.
《約翰福音》三章二十節.
作惡的人恨光.
不肯接近光.
恐怕他們的罪恨暴露出來.
聖經除了很正面地鼓勵我們.
亦都有負面的地方.
聖經不會負面.
不過人是會負面.
就是面對經文.
面對神的光的時候.
我們會不肯接受.
所以弟兄姊妹.
我們無論是信主多久的老信徒.
或者是剛剛起步的新信徒.
我們彼此鼓勵,彼此提醒.
就是有時我們寧願.
好像會想停留在黑暗的焦慮當中.
黑暗一定是帶給你焦慮.

$^{401}$因為看不到出路.
亦都不願意去讀神的話語.
因為神的話語有時.
會照光我們裡面的私慾.
我們不信任神的話語.
有時我們會小信.
未必不想信.
不過我們可能很小信.
我們身邊有時會有些朋友有病.
大小的病.
有時甚至會有癌症.
在NHS排隊.
去照鏡或者做手術.
在那個時候.
這盞腳前的燈路上的光又出現了.
因為在那個時候.
很多時我們會發覺.
人生的計劃瞬間歸零.
沒有很遠的計劃.
三年後怎樣呢?.
兩年後怎樣呢?.
我們望著那個錶光去跑的時候是怎樣呢?.
沒有了.
我看不到明天.
所以我們看不到那幾年後會發生甚麼.
完全沒有了人生地圖.
每天早上醒來.
其實我們只能夠求神.
給我們每一天的恩典.
去活好每一天.
那個就是腳前的燈.
在路上的光.
光就會一步一步地展現出來.
和大家分享一個見證.
關於上帝都讓我們學習遵行神的道.
我們向一位老人家.
其實那個是我的岳母.
傳福音了很多年.
在上次我來Swindon華人基督教會.
9月到今年5月.

$^{441}$中間12月尾我們就回了香港.
其實目的只有一個.
就是陪他去做手術.
陪岳母做手術.
所以我們整個旅程在香港.
每天都是在醫院.
我意思試探病時間.
沒有時間去其他地方見很多朋友.
因為他很需要我們禱告.
因為他還沒有順序.
我們上機的時候已經懷著一個.
很希望他順序.
在病床.
即是我們下機就去到醫院.
接著太太問他.
你信不信耶穌?.
其實這個問題.
心大略問了很多年.
祈禱了很多年.
你信不信耶穌?.
耶穌未必醫好你的.
不過耶穌會給你平安.
給你安穩.
給你盼望.
給你其他東西.
未必一定給你醫好.
因為他要做一個脊椎.
一個大手術要開.
其實他坐了輪椅一段日子.
這個手術對他都很重要.
他竟然回答.
信.
我們真的沒有想過這麼多年.
當然我們預備好.
我們在英國特意帶了這個十字架給他.
這個十字架在那一刻送了給他.
在病床當中.
這個十字架在英國.
我記得好像.
有沒有人知道哪裡可以買?.

$^{481}$我不記得是不是某些教堂裡的書室.
書店有時候是大中小的.
這個是我以前在神學院讀乳木的一科.
其中一科是要讀醫院牧羊.
他教我們這個很滑的十字架.
是令到病人可以握緊它.
令到他能夠掌握.
亦都握緊神的應許.
如同握著上帝給他的應許.
亦都很安全地放在病床邊.
沒有利器不沾的.
可以放在身邊.
這個很適合病人.
老人家或者有病的任何人.
有一位我們很資深的導師提醒我.
其實你可以給這個十字架.
我就把我唯一一個在英國的十字架.
搭飛機拿回香港.
在病床給了他.
他從此就握緊舉起.
這是三隻手.
三隻手指我們三個在床邊祈禱.
與神的話語有什麼關係呢?.
在我們聖母交投的時候.
誰帶領祈禱呢?.
我太太說「當然是你啦」.
我們祈禱的時候說哪一段經文呢?.
我很覺得我們唯獨經文.
才是我們最穩妥的依據.
有時都會是人的感動,情緒.
那些都很重要.
因為是神給我們的.
唯獨經文歷久常新.
幾千年都無變.
這一刻記得.
「你若口裡承認耶穌為主於心裡信神.
叫他從死裡復活就必得救」.
「口裡承認耶穌為主於心裡信神」.
所以口信心也都信.
「叫他從死裡復活就必得救」.

$^{521}$那一刻神的話語就成為病房.
很真實的安慰和生命得救的確具.
一個生命的翻轉.
一切都是從神的話語開始.
請姐妹我們在思想神的話語.
和遵行神的話語那裡.
我想再提多一次.
我們需要記得神的話語.
所以平時我們要記,要讀經.
得來我們還要背經文.
趁我們還記得.
趁你還記得經文.
我現在記得的經文.
很多都是我很年輕的時候記得的.
到後來再記就開始越來越難記.
記東西越來越不容易記.
順帶提一提.
其實將經文變成詩歌.
正如我們剛才唱了一首.
由經文轉為詩歌.
其實很多希伯來傳統.
他們承傳的信仰.
經文是很重要的一件事.
它令我們能夠將一些難背的經文.
變成詩歌.
變成詩歌就可以一代一代傳下去.
這是原先聖經為什麼有詩歌題材的原因.
他們就是想傳下去.
未有印刷術的時候.
其實都好像我們中國人說的.
月光光照地堂連山.
這些歌你不需要看譜.
但你小時候有人唱.
你就一代一代傳下去.
所以這個順帶一提.
整個話語成為了一些詩歌.
亦都令到特別.
我以前在神學院音樂科.
老師教我們.
我們有很多現代音樂.

$^{561}$但亦都有一些很古.
不要說古.
不要用古這個字.
比較歷久常新都唱了很多代.
那時候老師提醒我們.
教會裡面有時候有些老人家.
他唱了一些詩歌聽了幾十年.
他來到教會可能只記得.
《奇因》電影舉例.
或者大家都有唱過.
《三不重》舉例.
或者是還有一些詩歌.
亦都學一些新的.
《傳遞至高時我處》.
OK 很好.
他又會學新的.
但亦都有舊的詩歌.
令到教會無論信主幾十年.
信主住幾年.
或者剛剛認識耶穌的人.
都可以一起將這個信仰.
傳統承傳下去.
一路加下去.
一路累積信仰下去.
越來越踏實.
越來越根深蒂固.
越來越建基在整個話語裡面.
很好.
我們都繼續努力.
第三點想和大家分享.
「留思想」之後就去到「遵行」.
去到第三就是「活出神的話語」.
我們一起讀.
《瑪利諾福音》五章十六節.
預備起.
「如果我們也多謝你的少數人士.
你耽誤了我們天地的多難題.
將你的賢慧和光榮所尊賞」.
請姐妹們在這裡.
由腳前的燈.

$^{601}$去到世上的光.
我這裡有個副題叫「互相保光」.
由詩篇.
剛才的詩篇.
去到《瑪太福音》.
神的吩咐升級了.
光不再單單是讓我們自己.
去照自己的路,個人的引導.
光變成了透過你去照亮別人.
是一個群體的見證.
不再是自己.
而是照亮別人.
不再是照你的一生.
而是群體的見證.
我想用一個在英國很生活化.
我希望就是一個例子.
和鄰居之間互相照明的例子去分享.
大家知道冬天四點鐘.
三,四點就會黑了.
這個不單是我們要學習.
其實那些本地的英國朋友.
有時我的同事都會說.
其實都很不適應的.
這麼早就黑了.
所以冬天他們會去西班牙.
去其他地方玩.
其實都不單是我們不適應.
四點後全黑.
在冬天.
想像一個畫面.
你的鄰居.
因為我們上年九月到現在五月.
我們搬了一個比較永居.
永遠的一個居所.
都叫做一個場.
我們認識一些新的鄰居.
你試想像一個畫面.
當他要推一個黑桶要倒垃圾.
四點後都未煮晚飯.
第二天他就推出一個黑漆漆的全黑.

$^{641}$推一條短的路出去街邊.
完全黑的.
其實他們都是新搬來的.
我們剛好這間屋的兩邊.
都是剛剛新搬前後幾個月.
是重新認識的.
而另外我們.
我和太太開車放工.
特別是太太放工.
四點後六點回來.
舉例.
停車場有一點距離.
那個停車區.
黑漆漆的全黑.
又鎖了車門.
走過來的時候全黑.
你有很多想像到底安不安全.
會不會踩到一些骯髒的東西.
或者是一些垃圾.
卡到安全.
天都黑了.
什麼都看不清.
後來我們就安裝了一盞智能的燈.
在我們門口.
剛剛他們倒垃圾的位置.
他一出來就亮了.
他就可以走出去.
放下黑桶.
然後回來.
他很開心.
他問我.
那盞燈在哪裡買的.
我也很開心地分享了Amazon的連結給他.
這個是我們的那盞燈.
不用一個星期.
三天後就到貨了.
他不用一個星期.
他就安裝了他那邊的風扇.
螺絲上實了.
所以接著到我們泊車的時候.

$^{681}$當時還是一月.
所以還是很黑的.
泊完車照亮了整個泊車的區.
他出去停車倒垃圾.
我們泊車再回家.
大家沿路都是兩盞燈都照亮了.
有一次還沒到晚上.
有隻貓經過.
整個這邊都亮了.
然後再走到這邊.
全部都亮了.
整條街都亮了.
我們互相互照.
他用英文.
他很欣賞我們那盞燈.
我就跟他說英文.
We cover each other.
我們互相保光.
互保.
後來不單止我們兩家人.
兩個neighbor.
剛才說的愛倫如己.
愛倫社.
舊約的所有的律例.
都歸納到一句.
就是愛倫如己.
愛人是翻譯成人.
其實再準確一點是倫.
倫社的倫.
neighbor.
他還挑戰你.
不是愛遠一點的那個.
他挑戰你愛近一點的那個.
倫社.
你的鄰居.
你坐在你旁邊的弟兄姊妹.
這個很近.
或者在家裡.
諸如此類.
自己可以應用.

$^{721}$一妖通白.
在英國這個生活化的例子.
互相照耀.
就cover each other.
於是我們的路就很安全.
這個就是我們做世上的光.
一個很真實的寫照的一個比喻.
從此不單止我們兩家人.
現在聽到有時.
再遠一點的那間屋.
因為它的sensor照顧到.
他們有人探望他.
在門口聊天.
又光了他們.
照光了他們.
於是我們聽到突然有些人在聊天.
看出去.
他們在閒談.
這個以前是不會發生的.
現在sensor有時間關掉了.
我看到他.
他特意這樣.
再走遠一點.
令到他又光了.
又光多一段時間.
那個畫面很有趣.
有時.
當然貓經過他都照耀他.
當然我知道.
說笑的.
貓是不需要那個燈的.
他的夜光能力應該很強.
其實.
我記得.
如果大家還想聽.
我還有很多故事.
我記得我上一份工作.
上一份工作是做charity.
幫那些老人家學科技.
其中一家人.

$^{761}$一個老人家我被委派到去.
因為他們是housebound.
他們不可以出街.
他們是有disability.
有輪椅那些.
不可以出街.
但是他還可以跟他的家人聯絡.
我被委派去到那個家.
那位伯伯是九十多歲.
看上去這裡應該未必有九十多歲的朋友.
我今.
九十多歲了沒有?.
好了.
和早堂不同.
多謝你.
九十多歲.
都來崇拜聽到.
好了.
說回這裡.
今早我就說笑.
四十多加五十多就很多.
現在三十多加六十多就很多.
九十多真的很少.
那個伯伯.
他的故事長話短說.
他其實都繼續和他的孫子 叔叔.
有聯絡的.
他有個ipad.
用那個視像.
facetime.
有聯絡的.
都保持著.
這個伯伯很厲害.
可能家人安排給他都不一定.
但是.
那個九十多歲的伯伯.
到我探望他那一年.
其實心裡有數.
他覺得到底我能不能過下一個.
過到這個冬天.

$^{801}$下一個.
春天他還在不在這裡呢?.
還在不在生呢?.
不知道.
但是他的孫子和他的叔叔.
經常都很想他.
你可不可以留下你的故事給我們聽.
你一生人.
做基督徒.
或者你一生在英國.
你的人生是怎樣過.
有什麼可以留給我們.
當故事來看.
但是他不懂寫字.
年紀這麼大.
於是我就在那六次裡面.
四至六次.
我就教了他一個.
其中一個Voice to Text.
即是語音變成文字的功能.
他不是錄音.
他是語音.
當然現在AI也做到了.
當時只有這個科技.
他就很厲害.
他學到了.
他用回他的iPad.
就將他的文字.
讀他的英文.
變成準確的英文文字.
因為他的英文是很純正的英文.
我們學會了北方的英文.
就令到他的英文是.
但是Google.
為什麼我說純正呢.
因為Google能夠完全.
百分之一百.
將他講的東西變成百分之一百的文字.
他開始講.
我年輕的時候.

$^{841}$曾經被委派.
差到一個地方.
士兵戰爭的地方.
他要被派到那裡.
士兵去上陣.
有一次他也說多一點.
他和一個士兵坐在一起的那輛車.
就去到一個前線一點點的地方.
想去打探探聽一些東西.
突然被人發現.
然後他們不知如何逃亡.
他講到這裡就沒有講太多了.
但是他這些說話全部都記錄了.
文字記錄了.
還有他也是一個.
信主很久的基督徒.
因為在我們言談裡面.
雖然我們不會主動講基督教信仰.
因為我們是Charity.
但是他先講起.
那我就去接.
他講到原來他住的那個區.
在Bristol的那個區.
是以前約翰·韋斯利.
最初開始在街上講道的那個區.
所以他說他那個.
不知祖父還是哪一代在上面.
曾經親耳在街上聽到.
約翰·韋斯利講道.
即是你知道.
秦道就是在那裡開始.
所以我相信他的人生故事.
都包括很多信仰的經歷.
人生的一些信仰的挑戰.
這些成為了很多寶貴的.
承傳給他們的兒子,孫子,神.
主給我有這個機會.
去將這個光分享給他們.
當然我們都有講信仰,有講經文.
這個行動教了他這麼好的一件事.

$^{881}$他亦都很能夠掌握.
將它變成一個更大的光.
寫了一些故事.
然後就祝福到他的家族.
他的兒女,他的孫,他的神.
甚至可以再祝福其他朋友.
我們的光很微小.
但是當我們傳出去.
分享了就會更加大.
我們的鄰舍亦都是跟他分享了那個光.
他亦都再分享.
那個光再光照整個社區.
整個community.
整條街.
一路光下去.
在教會亦都是.
所以無論你是信了主很多年.
信了主這幾年.
或者你只是剛剛初信.
或者剛剛五個禮拜.
五十年都好.
我們在主裡面是沒有分彼此.
我們每一天都繼續去學.
活出神的話語.
在英國這個地方.
有時是光明 有時是很黑暗.
有時是戰亂.
有時是inflation 即是通脹.
我們面對很多事情.
都願我們都是一班互相曝光的人.
因為我們心願.
可以頂著門在教會亦都可以互相曝光.
我們聚在一起.
就成為一盞明亮的燈台.
明亮的燈台 聚在一起.
成為一個明亮的燈台.
如果再講故事就快超時了.
如果大家願意再聽.
我想講多一個.
大家不介意吧.

$^{921}$講多一個.
就是一個曾被我們每個人的帶領和預備都不同.
這個經文亦都令我想起.
以前記得的經文.
是可以應用在不同的場景和情況.
這個經文就是.
源於有一個弟兄.
他曾經都有在外國生活的經驗.
對他來說在英國的適應不算很困難.
但他沒有將順利當作理所當然.
反而常常主動熱心去幫助其他新來到英國的朋友.
以及去適應新生活 解決各種困難.
我們看著他最適合的經文就是.
朋友乃時常親愛.
弟兄乃患難而生 是真言.
我記得這個經文.
我為甚麼會記得呢.
其實都是因為我小時候十幾二十歲.
在不知道是基督還是突破.
有些那麼小張卡片的書籤.
有兩隻羊仔的那些.
那些是一疊的.
在一個膠盒裡面.
其中一些是我們互相送的.
團契很喜歡送給互相.
那個經文在那裡.
我小時候記得.
竟然這個時候用得上.
朋友乃時常親愛.
弟兄為患難而生.
我們也成為了別人患難而生的弟兄.
這也是一個很美好的互相曝光.
另外還有很多經文可以去應用的.
譬如在.
如果弟兄姊妹可以記得.
其實有很多經文.
包括曾經有一個姊妹跟我說.
她也試過曾經很崩潰.
有些日子很崩潰.
不知怎樣的時候.

$^{961}$那段經文救了她.
她記得一部份.
不過她想起那段經文是來自詩篇.
143.
她說:耶和華求你快快應允我.
我已絕亡.
求你不要掩面不顧我.
不然我必和下墳墓的人無異.
求你讓我再清神.
就聽見你的慈兒愛兒.
因為我信靠你.
求你指示我當走的路.
因為我的心仰望你.
她很崩潰,很絕望.
有段經文在那個時候出來.
聖靈提醒她.
有些經文你是小時候.
或者現在開始背,開始讀.
你記下了它.
在適當時間聖靈提醒你.
最後有一個結語.
我很想套用我其中一位老師.
我未讀全時間神學之前.
讀兼職在中神.
石昌牧師之楊依.
她在我們放工之後.
經常在桌子上等我們晚飯.
一起吃飯的時候去交流.
然後就上她的課.
一九九幾年,二千年那些日子.
我懷念她.
因為她也去世了幾年.
我想用這個去做一個總結.
就是回歸聖賢之召喚裡面其中一段.
請聽我為大家讀出來.
你的生命是否建立在神的話語上.
若然不是.
你也不會得到從神而來的使命和肯定.
你只是自己活著.
唯有道在你的生命中.

$^{1001}$我們必須回到神的話語去.
重新建立這個基礎.
唯有神的話語才可以成為我們生命的基礎.
丙姐妹.
神早已預備你,預立你.
神預備了你.
祂喜悅你這個人.
祂要栽培你成為一編度.
你要開放你自己.
要多聆聽神的話.
可能你已經聽過不少神的話.
但神的話是否能夠進入你的心呢?.
你首先要校正方向.
使你一生都以神的話為你生命的基礎與依據.
我們同心討共.
你的話語成為了我們一生很穩妥的物.
讓我們定在神的道當中.
讓我們思想神的話語.
讀經讓我們可以多些背經.
讓這些話語能夠放在我們的腦裡.
記得.
神亦都求你讓我們去遵行神的話語.
讓你的一生都進行.
我們不偏離.
因為晝夜思想.
這樣我們的生命就會亨通.
凡事順利.
與神的心意,神的旨意對齊.
與神對齊.
我們的生命與神對齊.
一生都不偏離.
亦都願我們的光都照在人前.
不單止是我們自己的生命.
亦都互相補光.
照在人前.
叫他們看見我們好運運.
便將榮耀歸給我們天上的父.
我們都將你的話語.
恭敬地交在你面前.
如果有神.

$^{1041}$如果有很多人的說話.
求主你都剔除.
承諾在我們每一個弟兄姊妹的心裡.
都是你寶貴的話語.
我們祈禱交託.
外交專利 多謝.
\newpage



\section{哥林多前書 13:4-8}
\label{sec:032180}
\textbf{愛的四重奏}~~~~羅杰才牧師
\newline
\newline
連結: \href{http://www.church.com.hk/acms/content.asp?site=cdc\&op=show\&type=product\&code=032180\&layout=sermon}{\texttt{032180}} ~~~~ 語音日期: 2026-05-10
\newline
\newline
\hyperref[sec:032179]{\small{< < < PREV SERMON < < <}}
~
\hyperref[sec:index_chronic]{\small{[返順時目]}}
~
\hyperref[sec:032181]{\small{> > > NEXT SERMON > > >}}
\newline
\newline
哥林多前書 13:4-8
\newline
\begin{longtable}{cl}
\hline
\hline
章節 & 經文 (和合本修訂版)\\
\hline
13:4 & \begin{tabularx}{0.7\textwidth}{X} 愛是恆久忍耐；又有恩慈；愛是不嫉妒；愛是不自誇，不張狂， \end{tabularx} \\ \\ \relax
13:5 & \begin{tabularx}{0.7\textwidth}{X} 不做害羞的事，不求自己的益處，不輕易發怒，不計算人的惡， \end{tabularx} \\ \\ \relax
13:6 & \begin{tabularx}{0.7\textwidth}{X} 不喜歡不義，只喜歡真理； \end{tabularx} \\ \\ \relax
13:7 & \begin{tabularx}{0.7\textwidth}{X} 凡事包容，凡事相信，凡事盼望，凡事忍耐。 \end{tabularx} \\ \\ \relax
13:8 & \begin{tabularx}{0.7\textwidth}{X} 愛是永不止息。先知講道之能終必歸於無有；說方言之能終必停止；知識也終必歸於無有。 \end{tabularx} \\ \\
[1ex]
\hline
\hline
\end{longtable}
$^{1}$不知不覺因應今天我的講題.
「愛的四重奏」.
所以有四個人獻奏,獻唱.
這個「永遠的奇妙的愛」.
我後面聽就唱得非常好.
是不是.
真是值得讚賞一下的.
鼓勵一下他們.
常州堂的歌詠小奏是很特別的.
我每次回來聽.
我有些其他教會也有聚會.
也有聽.
這是上帝給我們的一個恩賜.
我選了這段經文.
已經熟悉得不需要再講解了.
因為幾乎懂得背了.
太多場合裡面主要是婚禮.
或者有些地方也是.
都會用它來講的.
可以說是一個相當經典的.
但我今天主要不會講這段經文.
但這段經文是我整個講道裡面的一個基礎.
什麼是愛呢?.
有一首歌其中一句就說.
什麼是愛呢?.
合則是愛.
離則有恨.
好像很直白的.
什麼是愛呢?.
合則是愛.
我和你合就是愛.
不合就分開了.
那就是有恨.
當然不贊成.
合則是愛.
愛其實不是一種偶然.
或者幸運的遇上.
那種的驚喜.
或者是一種很浪漫的激情.
很多時候我們都用一種很浪漫的激情.

$^{41}$來去形容或者代表愛.
愛是什麼呢?.
愛是一種生命質素的表現.
甚至你可以這麼說.
它是一種屬於上帝生命質素的表現.
《聖經》說.
神就是愛.
神是用愛來去呈現他自己的.
愛就好像藝術一樣.
好像智慧一樣.
好像修養一樣.
好像品德一樣.
它是.
這些的匯集.
是一種很複雜的能力.
質素的一種匯集的表現.
你要掌握藝術.
掌握智慧.
你要達到有修養.
你才能夠達到一個程度.
其實真的很不容易.
是相當難的.
就像唱歌一樣.
你只要去到一個水平.
你是相當需要下功夫的.
而且它就在人的關係裡面表現出來.
當我們跟別人相處.
別人感覺到有愛的時候.
他不只是說我喜歡他.
而是他跟你的相處當中.
有了一種有素養的.
有內涵的.
有溫度的.
有質素的.
一種這樣的情感上.
或者關係上的流動.
所以愛和真理.
和公平和良善.
和和平.
很多元素都是組合在一起的.

$^{81}$就好像寶萊在哥林多前書.
他用了很多元素來組合.
愛就有這麼多種元素在裡面.
所以它不是很簡單地發生.
它也不能夠很簡單地持久.
哥林多前書.
這種論愛的表達.
可以這麼說.
他把一個很完美的人格.
表寫出來.
你又要有忍耐.
又要有恩慈.
又要有那種不嫉妒的心.
你又要不自誇.
又要不張狂.
一個人做到這樣.
你說他的水平.
他的品格是多高呢.
他又不求自己的益處.
他又不輕易煩惱.
他又不計算別人的惡.
他又不會做那些不好的事.
只喜歡真理.
一個絕對絕對是上等的人.
相當之成熟的人.
或者完整的人.
今天我的題目主要是.
用呂奕詩她的一本書.
她說的四種愛.
我就是說這四種愛.
親愛有四種愛.
人間有四種愛.
第一就是親愛.
第二就是友愛.
第三就是情愛.
第四就是對世界和上帝的愛.
這四種愛可以說是概括了.
我們在愛的範疇裡面.
四種類型的關係.
它是重要的.

$^{121}$我們不應該偏頗.
我只是有親愛.
我只是愛我家人.
或者我只是愛我的朋友.
情愛還有對世界和上帝的愛.
今天是母親節.
所以我也選擇了.
跟母親節有少許相關的.
最重要的.
愛.
其中的就是父母的愛.
骨肉的情.
這種是基於血緣關係.
血緣是一種先決的條件.
它甚至近乎於是一種無條件的愛.
做父母的對兒女的那種愛.
是真的沒有什麼理性的.
沒有什麼底線的.
我們做兒女的.
可能未必有這麼強烈的感受得到.
但是如果你反過來的時候.
我也見過不少的父母.
似乎特別是女性比較多.
當然有些男性也是很強烈的.
但是那種付出.
那種耐性.
那種程度是我們一般人很難想像得到.
接受得到.
明白得到.
這種親愛的特點是不平等的.
它不是交換式的.
不是你給我我就給你的.
它是很自然地發生的.
因為你是我的兒女.
我愛你就是當然的.
也不是因為你好.
也不是因為你乖.
也不是因為你帶給我一種光榮.
完全沒有這些的.
我們是超越這些的.

$^{161}$他就愛她.
想盡方法去照顧她.
她甚至連那個相處的機會.
通常我們的感情就是相處出來的.
我們生她出來這樣對著她.
從小時候一塊粉這樣.
養大她.
有那種情感在那裡.
甚至她是不用的.
沒有的.
最近看到一個影片.
就說.
有一個人在街上.
就看到另外一個人跟她很像樣.
她很奇怪.
於是消息慢慢地傳回她父母那裡.
她有一對雙胞胎.
帶雙胞胎到她家裡.
哥哥跟弟弟都在家裡.
那沒問題的.
但是.
她怎麼那麼像樣.
於是就找回那個.
這個做媽媽的很奇怪.
怎麼會那麼像樣.
為什麼會那麼像樣.
於是就去做親子鑑證.
驗到99\%.
她們是同一個血緣的.
那就是她了.
再驗另外一個.
那個原來是調戀了.
她一出生的時候就調戀了.
抱了第二個回來.
就養大了.
那她應該對這個沒什麼感情.
一天都沒對過.
不是.
一知道的時候.
就完全是另外一個世界了.

$^{201}$這個就是我的兒子.
是我的骨肉.
親情是這樣.
她不是因為她乖.
因為我跟她相處得久.
是這樣的形式.
但是親愛.
就不是依附形的.
不是依附形的.
孩子由依附父母.
到她慢慢就會離開父母了.
會離開父母.
兄弟姐妹.
也會自己獨立組織家庭.
是不是.
所以親愛會慢慢地分開.
不等於她不好.
親愛就好像神的愛.
把自己賜給我們.
我們自己會成長.
自己會獨立.
有種關係的黏連性.
是很重要的.
也很必然的.
這是第一類的愛.
就是親愛.
第二類的愛就是友愛.
友愛就是朋友之間的愛.
是基於條件形成的.
有條件才可以做朋友.
大家是同意的.
是有條件才可以做朋友.
如果沒有條件.
我們可以彼此之間.
有某一種.
但是做不到朋友的.
做朋友真的要講條件.
條件如何.
你決定了.
那個友愛的質素如何.

$^{241}$你也決定了.
你跟哪些人有友愛.
這種關係.
所以.
我想認識更多好朋友.
你要做好一點.
朋友在某一個範疇裡面.
好像唱歌.
或者玩音樂.
同一個水平才可以玩.
做運動也是.
你在那個水平裡面.
你不一定是最好.
起碼在那個水平裡面.
才可以玩.
友愛是一個這樣的性質.
所以.
條件如何.
也決定了.
你友愛的質素.
友愛的深度.
友愛的長度.
所以有些人.
他有很多朋友.
有些朋友是很短暫的.
一會就過去了.
要調整.
友愛又要好像栽種一樣.
我今天太太進來了.
她喜歡種植.
在外面那個小小的.
我不敢叫花園.
那麼多地方不是花園.
小小的地方.
還要搞來搞去.
值得其樂的.
我們看她.
也是很多東西有趣的.
你要不斷的修理.
要施肥.

$^{281}$要管理她.
她才會長大.
友愛也是.
一段關係.
一段友情.
你是要經營的.
一個人如果他那個.
心靈的世界是豐富.
也夠深度的話.
他結識的朋友.
那個友愛的關係.
同時也可以更加有深度.
也很持久.
他們不需要.
經常都要喝茶.
經常都要見面.
經常一起.
那就很開心.
一分開就沒事了.
一起就開心.
不在的時候.
那就沒了.
我們也知道.
有不少人是這樣的.
簡簡單單.
誰知道.
晚上去喝完.
吹完水就回去.
那就沒了.
因為他們在很淺層那裡.
是有那個關係.
沒有去深層的關係.
但是如果一個人有深度.
他們溝通的東西不同.
那種的共鳴不同.
彼此珍惜對方的程度.
或者是尊敬對方的程度.
又不同.
叫做一見如故.
是不是.

$^{321}$他們可以.
他們知道在這個人間裡面.
找到好像他一樣的想法.
能夠這樣交流的人.
其實不多.
很珍惜.
所以這個人也是很特別.
所以如果我們想有好的友愛的關係.
都一定要將自己作出一種提升.
當然也要對人有一種接納.
和願意跟別人交流.
有愛的人.
是一種幸福的人.
第三種愛是情愛.
情愛就是人到了青年期的時候.
那些愛與夢轉變.
就開始有發情期.
於是就會有一些竇牙夢.
就可以開始對異性有好感.
開始追求.
就是情愛.
以前的社會.
如果發展了感情.
拍了拖.
就會結婚.
結了婚就是終身的了.
我們基本認定這個設定是好的.
是合理的.
不過現代的社會似乎有一點變形.
就是.
大家都知道現在的男女關係.
婚姻狀態的情況.
但是情愛呢.
這也是一個很重要的愛.
人建立家庭.
和一種彼此之間.
去建立生活的一種愛.
情愛不是必然發生的.
不是必然發生的.
有些人不是很多人.

$^{361}$有些人沒有.
有幾種情況.
我也想有.
不過沒有.
我也不能夠強硬自己.
我找不到.
我也很想找到另一半.
我找不到.
找不到又怎樣?.
我站起來閉上眼睛.
嫁了她.
閉上眼睛.
只不過回來了.
不行的.
有些人甚至是.
他真的沒有這種衝動也好也不好.
你不能夠說.
這個人真的有問題.
要見醫生.
是有.
可以這麼說.
普遍.
多數人一個階段去到的時候.
他就會有這種情愛的嚮往和追求.
人類才會結合.
才會發現下一代.
這是一個過程.
它是一種很健康的.
很正常的事情.
並不是必然發生.
不是說我沒有了情愛.
那我就.
人生就沒辦法完成了.
肯定是出了問題了.
主耶穌也沒有結婚.
祂不一定的.
但是祂是在我們的世界裡.
祂是應該容許祂發生.
也讓祂自然發生.
有些人很低貶.

$^{401}$認為這是不需要的.
當然不是.
情愛也是很重要的.
情愛可以一分為二去解讀.
如果我們用中文來說.
有情和有愛.
情和愛是不同的.
情就是感情.
我喜歡你.
感情.
那你怎會和一個人彼此相愛呢.
因為你喜歡他.
但你喜歡他之後.
你和他拍拖之後.
你要結婚的時候.
你站在教堂那裡.
你要立約的時候.
就不跟你說喜歡了.
你不喜歡也不行.
是說愛.
愛和喜歡是不同的.
愛是說責任的.
是不是.
我現在不喜歡你了.
你走開吧.
你發了誓.
你承諾了.
愛是在責任裡面來呈現的.
愛是不求自己的益處.
所以當我們真的愛一個人的時候.
他和情是不同的.
他是超越了情.
很慘,我和他相處的時候只有責任.
一點感情都沒有.
很困難.
在人性上來說真的很困難.
但是你要看到.
有些情況就是.
你會看到.
有些人去照顧一個人.

$^{441}$或者照顧伴侶.
伴侶已經不能表達任何.
對另一半的照顧.
或者好等等.
因為他出現了一些狀況.
但另一半是可以持續地.
我應該說過這個著名的故事.
不是著名,是我喜歡的故事.
是真實的故事.
是.
你們應該不認識.
葉特生是70年代的一個人.
他因為病了,寫了一本書叫《浮過生命海》.
一本小書.
小書裡面記載了一個故事.
他自己有一個堂叔.
堂叔二十幾歲就和堂嬸結婚.
一對逼人.
誰知婚後沒幾年.
堂嬸就腦中風.
癱瘓了.
堂叔就要照顧這個人.
睡在那裡.
一邊工作一邊照顧.
後來過了一段時間.
再收養一個女兒.
誰知整個文章裡面說到.
最後堂嬸都得享一個.
壽壽有六十幾歲.
即是三十歲就照顧了六十幾歲.
這個堂叔就是這樣照顧了這個.
嫁了過來,娶了她.
幾年而已.
他最後有幾個字.
什麼是愛情.
恩情.
道義.
存在.
哪裡.
我很記得.

$^{481}$什麼是愛情.
愛情就是我對你的承諾.
我會堅守到底.
執子之手遇住佳路.
終身不遇.
所以他是超越了情.
情是要有互動才有的.
但愛好像是這樣灌注下去.
愛情是毀身的.
是終生的.
愛情的重要.
當然是關乎兩個人.
亦都關乎下一代人的幸福.
正常的情況來講.
愛和情是要一齊存在的.
但是我們不可以只是重情而輕愛.
第四項就是對上帝.
對世界和上帝的愛.
這個不是一種出於關係的.
或者出於條件的.
當你出世的時候世界就在你身邊.
不是出於關係.
不是出於條件.
而是你領悟不出來.
我們的心智開通.
我們的心智成熟.
發現到原來這個世界關我事.
原來這個世界是養育我的.
從這裡我再看到.
這個世界是有一個創造者的.
我們要對這個世界做甚麼呢.
就是有一種不同的態度.
我們在這個世界生活.
我們未必懂得愛這個世界.
我們現代人對這個世界的傷害其實是很大的.
但這個已經是一種人力很難去回轉.
個人的力量很難去回轉.
我們自己個人在態度上知道.
他明白到原來這個世界和上帝是先於自己.
這個是一種覺悟.

$^{521}$世界是客觀存在的.
它不會改變的.
因此我們對這個世界和對上帝的愛.
是可以奉獻一些或貢獻一些.
去回報他.
因為他怎樣養育我們.
怎樣幫助我們.
所以我們應該要關顧或愛惜這個世界.
聖經有一句話不要愛世界和世界上的事.
有時被誤解說不要愛這個世界.
不要理會這個世界.
我說這是被誤解的神學.
那個愛是說你不要貪戀.
或者在當中來虛耗你的光陰.
所以說你不要體貼肉體.
不要愛惜身體.
心理學家伊弗洛姆提出了愛有四個元素.
我們說有四個範疇.
有親愛,有友愛,有情愛,有對世界的愛.
我們說有四個元素.
其實對每一種愛裡面都會有四種元素.
我覺得他這個說法值得推薦.
值得介紹給大家.
這四個元素就是.
第一個是關懷.
第二個是責任.
第三是尊重.
第四是理解.
所有的愛包括親愛,友愛,情愛,對世界的愛.
都應該有這四種元素.
就是你關懷,你擔當責任,你尊重和你理解.
關懷就是你善待它.
很簡單.
你適時的去.
譬如他需要問候的時候.
你適時的去問候他.
幫助他.
表達你對他的重視.
提出你能夠付出的幫助.
這就是關懷.

$^{561}$不難理解的.
關不關懷.
我們關不關懷我們的朋友.
我們關不關懷.
有沒有關懷別人的心呢.
我們都不理的.
我們只顧自己.
你怎麼可能有愛呢.
一定要有關懷的一種情懷.
第二就是責任.
在這種關係裡面我有什麼責任呢.
我是不是一個可以託付的人呢.
成為一個可以託付,可靠的人.
這也是一種很重要的責任.
就是別人可以將責任交給你.
你去擔當那個責任.
需要照顧的.
你去照顧他.
需要去解決問題的.
你去幫他解決問題.
這個責任是很重要的.
所以做妻子的,做丈夫的,做父親的,做母親的.
即是做兒女的.
我們有責任心.
責任感是很重要的.
不過現在的責任感好像是低落了一點.
人們都是斜薄的多.
不用做就裝作不知道之類.
這是一種歪風.
責任很重要.
怎麼去呈現愛呢.
責任.
準時到也是一種愛.
約了人準時去到.
表示一種尊重.
是一種愛.
尊重.
尊重是你給對方有他自己的空間,位置.
我愛你.
我幫你規劃了.

$^{601}$你跟我去走就行了.
你不要走另一條路.
這條路是我幫你安排的.
不可以的.
這條路是最好也不行的.
你可以提出你的意見.
你可以幫助他.
你可以輔助他.
上帝待人也是這樣.
如果上帝說不行.
你怎樣怎樣.
根本人沒有空間可以走.
為什麼今天人有這麼多的問題呢.
上帝對人真是有尊重的.
他有愛.
但是他真的容許人去做出他的決定.
尊重在愛裡面來說是很重要的.
我們要留意一下自己.
有沒有因為愛一個人的緣故.
我們去限制一個人.
強迫一個人.
強迫對方.
限制對方.
因為我為你好.
強迫這個限制.
又不可以說是完全不行的.
去到什麼程度應該這麼說.
我們去到什麼程度.
你給意見給到什麼程度.
然後你說沒辦法了.
既然你也是這樣堅持.
我已經一而再再而三.
我是勸你.
或者阻攔你.
你都不行.
那我沒辦法了.
然後真的沒辦法了.
尊重也不容易學的.
不容易做得到的.
但是尊重是很重要的.

$^{641}$特別有些時候.
對方是做一些對方喜歡的事.
或者對方想做的事情.
而那些事情.
甚至構不成好壞對錯.
有時候會這樣.
你不要做了.
關你什麼事呢.
你給一些空間.
去做一些他想做的事.
這是他的世界.
你認為不好的事.
但是對他來說他有趣味.
那他就可以做.
或者你可以容許他做.
尊重是一個很重要的元素.
在愛裡面.
理解.
需要被明白.
或者明白對方.
原來你明白人是很重要的.
人是需要被明白的.
如果你發現一個人很明白你.
你會很開心的.
你同意嗎.
他明白我.
你理解就好像被人發現被人接納.
所以如果一個人.
他不想理解你.
不想理解你.
你的感覺是什麼呢.
他無視我.
但是一個人耐心聽你講.
給你分析.
給你理解.
你覺得這個人是很接納我.
理解很重要.
所以去聆聽.
去理解.
去嘗試了解他.

$^{681}$你未必幫到他.
他未必想你幫他.
但是你理解他.
你知道一個人理解你.
你的心裡面就有一種.
很實在的感覺.
被理解.
愛有四個元素.
今天我們講的就是八點.
有四種愛.
有親愛.
有友愛.
有情愛.
有對世界和上帝的愛.
我們不要太窄.
我們要闊一點.
你看到原來在愛的世界裡面是這麼闊的.
這四種愛裡面我們都會有四個元素.
有關懷.
有責任.
有尊重.
有理解.
記住我們慢慢去學一些.
每個人都在學習的過程當中.
如果沒有這些的質素.
我們是不會發生.
或者很難發生.
或者發生了都不能夠持久.
不能夠深化.
哥林多前書的十三章.
四至八節的內容.
其實就是包含在這裡.
這些的內容就是包含在哥林多前書.
四至八節這裡.
我們不是在情緒上的.
感覺上的.
是我們人生.
我們生命中的一種火和光.
我們能夠令人家活得好.
活得更安全.

$^{721}$更溫暖.
我們也能夠見證上帝.
成就自己.
和別人的生命.
所以我們不可以缺乏他.
也不可以忽視他.
我們應該是要花多些心機和時間.
去發現他.
去培養他.
好不好.
我們一起祈禱.
真愛主我們求你賜恩幫助.
讓我們都能夠在生命當中.
有更多的愛.
\newpage



\section{歌羅西書 3:1-4}
\label{sec:032181}
\textbf{神榮耀的倒影——神的子民}~~~~黃紹權牧師
\newline
\newline
連結: \href{http://www.church.com.hk/acms/content.asp?site=cdc\&op=show\&type=product\&code=032181\&layout=sermon}{\texttt{032181}} ~~~~ 語音日期: 2026-05-03
\newline
\newline
\hyperref[sec:032180]{\small{< < < PREV SERMON < < <}}
~
\hyperref[sec:index_chronic]{\small{[返順時目]}}
~
\hyperref[sec:032182]{\small{> > > NEXT SERMON > > >}}
\newline
\newline
歌羅西書 3:1-4
\newline
\begin{longtable}{cl}
\hline
\hline
章節 & 經文 (和合本修訂版)\\
\hline
3:1 & \begin{tabularx}{0.7\textwidth}{X} 所以，既然你們已經與基督一同復活，就當求上面的事；那裡有基督，坐在神的右邊。 \end{tabularx} \\ \\ \relax
3:2 & \begin{tabularx}{0.7\textwidth}{X} 你們要思考上面的事，不要思考地上的事。 \end{tabularx} \\ \\ \relax
3:3 & \begin{tabularx}{0.7\textwidth}{X} 因為你們已經死了，你們的生命與基督一同藏在神裡面。 \end{tabularx} \\ \\ \relax
3:4 & \begin{tabularx}{0.7\textwidth}{X} 基督是你們的生命，他顯現的時候，你們也要與他一同在榮耀裡顯現。 \end{tabularx} \\ \\
[1ex]
\hline
\hline
\end{longtable}
$^{1}$在聆聽主的說話以先.
我們一起再次同心低頭禱告.
將我們的生命交給我們的主.
我們的神.
父神多謝你給我們再次有生命氣息.
踏足在地上.
行走在你的創造裡面.
能夠讓我們不單享受上帝的創造的榮美.
並且讓我們在這個地方.
能夠帶著一個使命.
能夠去真正活出.
反映到上帝你的榮美.
求你給我們這份委託.
可以妥善地去實踐.
也祈求上主你給我們有充足的智慧.
和恩賜.
好使我們能夠在這個地上.
能夠活出合神心意的生命.
求你要帶領.
無論是聆聽的,說的.
線上的,線下實體的.
我們要祈求上主你.
每一個都透過聖靈去祝福我們.
有充足的聰明智慧.
來去了解,明白.
並且領受上主你的說話.
成為我們屬靈生命的一部分.
好使我們都能夠活出,行出.
並且說出準確上帝你的救恩.
你的道理.
你的生命的特性.
讓每一個能夠接觸到主你的說話.
都同蒙救恩的祝福.
禱告奉靠主耶穌基督的名字而行.
阿們.
我以前小時候因為家裡不是很有錢.
所以每得到一樣東西我們都會很珍惜.
尤其是小孩子沒什麼機會經常去買玩具.
我們家裡有個特色.
不是隨便買一件玩具給你的.

$^{41}$我家其實街口附近有幾間玩具店.
但是有份看又沒份買.
見到同學有份買玩具.
真的很羨慕.
我父母通常都這樣說.
你考試考到好成績.
我就買一件玩具給你.
的而且確有時偶有僥倖.
可以考得好成績.
買到一件半件的玩具.
但是那件玩具你買了之後.
你預了有很長時間玩.
你不會經常有第二件新的玩具出來.
所以我們在很小的時候.
學會珍惜每一樣東西.
就算那件玩具玩爛了之後.
我們不會扔掉.
我們會做一樣功課.
就是拆開這件爛了的玩具.
有什麼都還可以拿出來繼續玩.
所以這個習性其實我都習慣了.
在現在家裡的東西壞了.
我不會馬上扔掉.
因為我會去拆一下.
看看裡面還有什麼可以用.
譬如我們小時候買了一輛小車.
我記得以前買過一個遙控車.
不是無線遙控.
你們玩過那些遙控車.
是不是有條電線扭著的那些.
那些很高級的.
有條電線扭著可以遙控.
向左向右轉彎.
就這麼多了.
還有可以走.
壞了之後怎樣呢.
我們就拆了.
拆的時候拆了個摩打出來.
就多東西玩了.
因為拆了個摩打.

$^{81}$只要你有了電芯.
你就可以想出很多玩法.
除此之外.
那輛車有個車頭燈.
又有一塊擋風玻璃.
雖然那些東西全部都是膠的.
但是我們就很珍惜的.
有一塊鏡子.
那輛車什麼倒後鏡.
有一塊小鏡子.
樣樣拆出來.
那我就想一下這塊鏡子有什麼用法呢.
這個車頭燈有什麼用法呢.
其實有很多玩法的.
譬如車頭燈來算.
我們就會怎樣玩呢.
就會發現原來車頭燈會有.
一個彎形的設計.
那我們叫convex lens.
突出的鏡.
我沒有記錯.
英文叫convex.
concave一粒的.
convex就是突出來的.
有什麼玩法呢.
原來可以用它來聚焦.
趁好天的日子.
我有個弟弟.
跟我一起伏在家裡的床窗口邊.
有出猛陽光的日子.
就會找車頭燈的鏡頭那塊膠.
來玩聚焦.
玩著玩著那個聚焦不太準確.
家裡就有個老人家用來看東西的放大鏡.
就拿個放大鏡出來玩.
原來放大鏡有很多東西玩的.
放大鏡不只是放大看東西.
放大鏡可以用來寫字的.
懂不懂.
是嗎 牧師可以寫.

$^{121}$是呀 你在太陽底下用個放大鏡.
對準位置.
它就會聚焦就有熱力了.
那你有沒有辦法不燒穿張紙.
因為燒穿張紙當然是很大件事.
在家裡有火災.
媽媽一定知道會出事.
不行的.
要控制好.
就在那裡燒字眼出來.
這樣就玩你一兩個小時了.
好天有得玩.
不好天沒得玩.
好天還有一樣東西可以玩的.
記不記得我剛才說.
拆車出來的時候.
有些小鏡子.
反光鏡其實那些哪裡是鏡.
那些是膠.
倒了一層反光的東西.
那我就有很多東西玩了.
又趁好天的日子.
兩兄弟就佈在家裡的窗口邊.
對著太陽.
然後原來可以反射.
反射有什麼玩法.
很悶是嗎.
不是呀.
家裡附近有些野鳥.
那你怎樣用鏡子照那隻鳥.
照到那隻鳥.
那隻鳥就被你嚇跑了.
原來還可以.
用鏡子反照那些光.
可以在對面樓那麼長的地方.
反照可以寫一下字.
照一下別人的房子.
其實也不是什麼好東西.
看別人的房子有什麼.
雖然看不到.

$^{161}$但是.
每一樣東西我都很珍惜的.
但是.
在這個過程裡面.
不知不覺就學了很多東西.
很多物理上的原理.
原來.
陽光可以有這麼多東西玩法的.
原來.
有適當的媒體.
就能夠將它去反照或者投射.
去實踐到.
這個陽光你原本想不到的作用.
所以在我小時候.
我曾經想過.
我以為我將來一定會做工程師的.
不過當然大家都知道結果了.
當你到了五十幾歲.
我現在五十六歲還是五十七歲了.
差不多吧.
工程師就沒有你的份了.
但是.
扮一下工程師也好.
有時候在家裡.
弄一下東西.
拆一下東西.
修理一下東西.
都可以滿足一下自己的.
所謂工程師癮.
這個人生.
不過.
從這些儀式遊戲.
或者儀式的經歷.
就讓我開始看見.
其實在屬靈生命上.
其實我們都在做一樣.
就是正在反照上帝是怎樣的一位上帝.
而這個正正也是創世記.
在一開始神做人的時候.
上帝做人的目的.

$^{201}$上帝做這個創造.
然後在最後做人的時候.
其實也是要讓人代表上帝.
在祂的創造裡面治理者地.
而這個治理就是要反映出上帝的性情.
而今天當我們去看這段聖經的時候.
就正正有一個很類似的陳述.
由歌羅西書的作者保羅提了出來.
原來我們作為神的子民.
我們的人生其實是一個上帝榮耀的倒影.
或者將祂反照.
將祂在這個世界的地上能夠活出來.
彰顯出來.
那個榮美仍然是上帝的.
但是透過我們的存在.
歸回有上帝的生命的時候.
我們就能夠在這個地上將祂反映出來.
所以今天我們一起去再次透過.
歌羅西書的第三章.
第一至第四節.
也是我們教會今年的主題.
談論到我們怎樣與基督一同去榮耀.
我們這位上帝.
或者一同在基督裡面.
彰顯神的榮耀.
讓我們能夠一同在這段經文裡面.
掌握,成為我們每一個基督徒.
走往後的日子的一個很重要的原則和動機.
讓我們一起再次打開.
歌羅西書的第三章第一至第四節.
讓我們一起去讀一下這段經文.
雖然不是很長.
但是裡面的意思和講述.
是非常深入的,有意義的事情.
希望我們一起去讀這段經文.
去更掌握內容想說些什麼.
請大家一起打開.
歌羅西書第三章一至到第四節.
找到之後我們讀.
所以你們若真與基督一同復活.

$^{241}$就當求在上面的事.
哪裡有基督坐在神的右邊.
你們要思念上面的事.
不要思念地上的事.
因為你們已經死了.
你們的生命與基督一同藏在神裡面.
基督是我們的生命.
祂顯現的時候.
你們也要與祂一同顯現在榮耀裡.
這段聖經有兩個字眼.
是講到方向的問題.
上面和地上面.
即是下面.
如果你在中間的話.
究竟講些什麼呢?.
似乎裡面講的東西.
好像有些宗教宣言.
但是往往這些宗教宣言.
你們可能看完之後.
未必真的知道它想講些什麼.
我們今天仔細去看看.
提提前文講些什麼.
如果大家還記得上一堂.
我們講到前文在第二章尾的時候.
我們談論到一個很重要的題目.
就是教會如果要去尋求社會認同.
其實是一件危險的事.
並且將教會拖入一個危機的當中.
當教會要嘗試去與世界接軌的時候.
甚至於遷就這個世界的時候.
其實無非是反映出.
今天的信徒在教會裡面.
其實他們缺乏了一份屬靈的骨氣.
信不過上帝的說話.
信不過上帝的啟示.
於是乎他們就要用巴結的手法.
或者巴結這個字眼.
你未必喜歡.
很負面.
很多時候教會就用了另一個字眼.

$^{281}$我們要服侍世界.
服侍社群.
因為基督來服侍這個世界裡面.
請大家留心.
基督來到不是服侍這個世界.
這個是很多教會講出來.
去哄信徒的一種.
我可以說是一種叫做掩眼法.
在你不覺意的時候.
以為真是教會是要來服侍社會.
這種其實是一種對世界的妥協.
希望教會在這個世界當中.
能夠賺回一些人的認同.
而使到教會可以有影響力.
這種其實是屬地上的一種手段.
所以我們上次講到有關於.
那些所謂purpose driven的教會官.
Mega church model.
Community church movement.
還有角色形式很多的.
包括靈恩派的教會.
他們都是以討好人的方式.
來嘗試讓教會繼續生存下去.
他信不過教會的生存.
是直根於上帝的說話.
他信的只是教會能夠用方法吸引人來.
教會才會成功地存在.
所以教會很多時候.
你會看見今天我們去看的.
不是教會要傳的福音的內容是甚麼.
而是專門研究福音傳播的方式.
方式不是錯.
但很多時候這些方式成為主體.
而缺乏了我們對於福音的實質內容掌握的時候.
我們變成本末倒置.
花了很多心機力量去想方法.
甚麼方法才能夠吸引人來教會.
人家喜歡聽.
而缺乏了內容.
變成了教會越來越空洞.

$^{321}$這個就是傳道書裡面講的虛空.
虛空這個字眼其實知不知道甚麼意思.
是無氣捉不到的東西.
空虛的東西.
好像存在,不過捉不到.
教會很多時候不是不知道要做甚麼.
而是不願意去做一些應該要做的事.
保羅對於哥羅西教會有他們的信.
是值得感恩的.
但是保羅也很擔心他們當時面對的情況.
因為在教會裡面有很多人.
傳很多小學.
那些小學回到一些傳統裡面.
回到一些似是而非,很基本的事情.
講到卦象,節旗.
這一類型的問題.
好像讓人回到一個.
合神心意的敬拜的空間和環境.
但是卻是不知不覺地陷入了.
被那些人去利用這些小學.
去操控了他們的信仰生活.
所以保羅不得不明明說明給哥羅西教會知道.
這些人在你們當中.
用這些方法去控制你們的信心.
將你們的信心投靠在他們的身上.
而不是投靠在上帝的身上.
投靠在這些環節,節旗的活動當中.
慶典裡面.
而使得你墮落了都不知道.
因為人在快樂當中是更容易墮落.
所以保羅提醒他們.
千萬不要漏了有些人是用這些方法.
在教會裡面去引誘你們.
去導致你們墮落.
而保羅這種做法.
講得這麼明明白白.
的而且確有人是不喜歡的.
不過既然全道的工作加在了保羅身上.
加在了你和我的身上.
我們就必須要有這份骨氣.

$^{361}$敢承擔挑戰.
當真理陳述出來的時候.
人縱使不喜歡也好.
但只要是上帝的說話.
是對人永恆生命有益處的事.
我們不應該害怕去做.
這個就是上文.
我們所看到保羅的陳述.
來到第三章第一至第四節裡面.
保羅不單止自己親身做這個見證.
並且指出他不要走這個.
熟悉的世界那一套.
不跟這一套去走.
就是以討人喜悅或者搞關係.
和人去聯絡.
然後嘗試用吸引人回來這個群體.
他不走這一套.
因為為什麼.
這一套不是帶人去上帝裡面.
只是帶人去到教會或者這個群體裡面.
保羅帶他們走一條什麼道路呢.
應該往什麼方向去走呢.
怎樣走呢.
在什麼基礎上走呢.
這個正正是今天這一段.
一至到第四節的主題.
我會分為兩個部份和大家去看.
第一個部份從第一和第二節裡面.
保羅先指出我們要有一個正確的方向.
去尋索和思想.
我們先看第一節和第二節說什麼.
他在這裡這樣說.
所以你們若真與基督一同復活.
就當求在上面的事.
哪裡有基督坐在神的右邊.
你們要思念上面的事.
不要思念地上的事.
你會在這裡讀完之後.
你全部字眼都知道在說什麼.
但是其實內容你可以一點都不明白.

$^{401}$因為這裡說了很多東西.
什麼上面啊.
右邊啊.
什麼地上的事啊.
你完全不知道他們想說什麼.
是不是想說一些.
我們要離開這個世俗.
脫離這個世界.
然後我們才是真正信主呢.
有人可以這樣詮釋的.
是不是這個意思呢.
保羅在這裡指出.
信耶穌不是一件紙上談兵的事.
也不只是一個思想的運動.
這一點我要很小心提出.
很多信徒在這個信主的過程當中.
常常都以為信主只是一個心而已.
心好就行了.
心態正確就行了.
其實信耶穌.
不只是一個心態上的改變這麼簡單.
信主是一個整個人的人生取向.
和他根基的徹底改變.
改變在基督耶穌的身上.
所以在前文講過.
所有這些小學.
其實應該要講的是講基督.
不是講小學.
你要認識的是基督.
基督的屬性是怎樣.
基督是怎樣的一位.
因為你的基礎在他身上.
你整個人的人生取向.
也是建基於這個基礎.
因此如果當一個人聲稱自己信了主.
甚至受了「浸」.
所以我們教會受浸是很嚴格的.
不是你說浸我就讓你浸.
你說是嗎?聖經不是這樣的.
說浸就浸.

$^{441}$如果我們是這麼輕忽去看受浸的事.
就造成很多教會裡面.
很多浸了而仍然是死亡.
屬靈上死亡的人.
教會的角色受浸到要幫助你清楚.
知道你是在信甚麼.
因甚麼而去浸.
但是如果一個人只是聲稱自己信了主.
或者已經受了浸.
但是仍然他的人生方向.
他做事的原則以及他的基礎.
用著同一套在世界同樣的思想的話.
或者他所作的一切行為,想法,宣告.
都仍然是高度與.
這個世界不信主的人的想法是一樣的話.
我可以很肯定地跟大家說.
這個人仍然未清楚他自己的認罪悔改.
他是未相信主的.
而他只是在做一個掛名的基督徒.
那為甚麼這麼多人喜歡做掛名的基督徒.
因為做掛名的基督徒.
他們想得到的就是成為基督徒有的好處.
仍然是一個利益的交易與上帝.
所以他們死都拿著.
我要做基督徒.
但是當他要做基督徒.
又做不足真是基督徒應有的生命的時候.
可惡的就是很多教會.
就將他們用了另外的方法幫他們包裝.
幫他們去假信主.
幫他們做假的掛名基督徒.
講得清楚一點.
不是幫他們做假的掛名基督徒.
幫他們做掛名基督徒.
做假基督徒.
教會不應該做這些事.
保羅其實想提出給哥羅西教會.
就是這一點.
在你教會裡面就有一批這樣的人.
存在在教會當中.

$^{481}$當一個真正信主的人.
在整個人生的方向和基礎上.
有一個基本的根本上改變之後.
保羅在這裡就很清楚的指出.
一個真正信主的人.
他必然是會怎樣.
他的思想行為都會往天上來去尋索.
所以在第一節裡面講.
所以你們若真與基督一同復活.
這裡所強調的真與基督的一同復活.
是前文所談論到.
每個信主的人.
他的生命是隨著基督而更生過來的.
他的復活是一個更生的生命.
並且有一個更生全新的身體.
可以賦予給他們的.
在這樣的生命之下的人.
保羅講到他們的思想行為.
他們的基礎就是會怎樣.
當求在上面的時.
但在保羅的神學上.
這個陳述可以說是革命性的陳述.
因為保羅很多時候所講的只是.
在屬世上我們怎樣與基督一起去活出.
是一個橫向的發展.
沒錯,我們可以有一個橫向的發展.
但保羅在講到比哥羅西教的情況就是.
你們必然會去.
為到上面基督的事去思想.
為甚麼?.
因為基督復活了.
基督今天不是活在地上面.
基督活在天上面.
所以我們就會去經常想.
現在天上面的主.
祂是怎樣想的呢?.
祂是怎樣的做法呢?.
當我做一件事情的時候.
天上的主.
你在上面也會怎樣告訴我祂怎樣做呢?.

$^{521}$就是這個緣故.
好像.
有沒有試過暗戀一個人?.
你又笑,當然有了.
Eric都有少少陰陰嘴笑.
有點暗戀.
暗戀一個人是怎樣的?.
你會經常想著他.
他現在會在哪裡呢?.
有沒有試過.
我用一個男生的角色去講.
如果我暗戀了一個.
我讀男校,假設.
完全是假設,不是事實.
假設我讀男校的.
我暗戀了隔壁校的女校有個女學生.
那你會怎樣?.
經常想著.
她上學了沒有呢?.
她是怎樣的呢?.
她是不是跟我搭同一架巴士呢?.
她哪一架巴士呢?.
她上學了沒有呢?.
到放學的時候.
你又想,放學了.
我要放學搭巴士回家了.
她會不會在巴士站等巴士呢?.
你很早就會去到巴士站第一時間部署.
她會不會在那裡呢?.
有沒有機會見她一眼呢?.
就是這樣的心態.
一個基督徒.
屬神的子民.
他心繫於基督的事上.
經常想著基督會怎樣做呢?.
這裡保羅提出這一點.
今天如果,這個若真.
所以保羅寫個if.
若真.
逢親假設陳述的說話.

$^{561}$我們就要問.
假設的陳述有哪一類?.
我認識經學,我講過很多次.
假設的陳述在聖經裡面.
總共有四類可能的解釋.
其中一類就是假設一定不出現的.
如果假設,譬如假設你是一隻小鳥.
即是你不是小鳥.
這種假設是讓你去代入一個情況的思想.
有些假設是甚麼呢?.
一定出現的.
如果你是一個男人.
我要向著一個男人講.
如果你是一個男人.
那你一定他是一個男人.
你當然當他是男人.
接著講以下的事情.
這是一定出現的情況.
有些假設是會出現的也不會出現.
這裡是哪一類?.
就是第三類.
第四類我不講了.
免得擾亂給大家.
由每一個哥羅西教會的信徒自己去考慮.
如果你是基督的同一同復活的.
你的生命就會和基督連結起來.
你想想每一樣東西都是基督會做些甚麼.
當然不是的,不用想了.
你一定不會想.
你想的不是基督想的.
你想的東西是甚麼?.
接著在第二節裡面講.
你想的東西全部都是地上的事.
今天教會有多少人.
是真的和基督有一個復活的關係.
而經常都是想基督的事呢?.
還是想地上的事呢?.
承繼著前文2章第20節裡面.
談論到哥羅西教會已經和基督一同.
同死同埋葬同復活在這樣的情況之下.

$^{601}$所以保羅在第三章第一節裡面.
進一步去強調.
哥羅西的教會會因此和基督一同.
有一種盼望升到在這個永恆當中.
在天上的生活.
當然信徒沒有馬上去在天上生活.
上帝讓我們活在地上仍然有一段日子.
因為上帝創造你在地上.
甚至上帝容許你和我今天活在地上.
是有個意義的.
意義是怎樣呢?.
是教導我們將福音傳開.
帶領呼召那些屬神的子民回來.
這是上帝放我們在地上現在的意義.
將來在新天新地裡面的時候.
就不同了.
不需要這一部分回到創世記.
原本上帝創造地上.
讓我們在祂的創造裡面去榮耀.
顯出管理上帝的創造這個身份.
那個身份又不同了.
但是現在上帝放我們.
沒有馬上升到天上.
但是我們的心卻有了這個條件.
來不斷思想在天上.
基督所作的事情是應該怎樣去做.
心態上,行為上,取態上,決定上.
我們怎樣去取捨呢?.
這個就是保羅在這裡所強調.
「當求在上面的事」.
那你說牧師我都當求在上面的事.
很多的事我都祈求了.
我往天上,往上帝,往基督的想法想了.
不過還有一點.
請留意.
這裡談論到的是.
哪裡有基督坐在神的右邊.
為甚麼有句這樣的說話.
甚麼來的?.
無端端坐在上帝的右邊關我甚麼事?.

$^{641}$坐在上帝的右邊.
其實這個陳述是早在.
列王記上第二章第十九節當中.
指到能夠活在上帝的右邊.
是一份極度光明的榮譽.
當一個人.
以地上的君王系統來計算.
如果一個人能夠在君王的右邊.
這個人就是君王的宰相大臣.
他是能夠有一個位份.
來和皇帝一起參與政事.
作出意見.
並且去幫助這個王.
去成就一個合王心思意念的一個行動或決定.
而這裡所指到哪裡有基督坐在神的右邊.
其實是告訴我們知道.
主耶穌基督現在活在上帝的右邊.
原來是有份和上帝一同參與.
怎樣在這個創造繼續下去.
包括現在這個拯救的計劃.
那個進展是怎樣去決定.
而這裡所講到的.
每一個屬神的子民.
特別是這裡講到的.
哥羅西教會的信徒.
他們能夠思想上面的事.
換句話來說.
他們和基督一樣.
與上帝一同參與.
上帝在這一刻.
我們在這個地上要參與的事情.
上帝在這一刻.
希望在這個地上要做到的.
就是拯救屬神的子民.
而這個陳述.
正正在保羅其他的書信裡面.
亦都有類似的介紹.
譬如在《羅馬書》第八章第34節.
同樣講到基督坐在神的右邊.
其實是講到.

$^{681}$基督怎樣去一同與父神參與.
在今天這個救恩的活動當中.
所以上帝在這一刻.
其實不是停止了那個事奉.
基督仍然在這個救恩上面是在參與的.
因為我們能夠與祂同埋葬,同復活.
所以我們應該在參與這份慰憤的上面.
我們有這樣的機會.
我們珍不珍惜呢?.
在舊約的神話當中.
曾經在烏加列神話當中.
亦都談論到有兩個祂的左右臣子.
一個叫Copha.
一個叫Casus.
兩個在巴黎風雨的神的手上.
他們同樣在右邊去參與.
今天我們作為屬神的子民.
能夠藉著基督的復活.
我們可以參與在上帝這個救恩的計劃裡面.
我們會不會因此而感到榮耀呢?.
或者感到光榮的一分子.
最近大家亦都參加我女兒的婚禮的婚宴.
大家有沒有留意到主家席上坐了甚麼人?.
我不怕告訴你,除了我們兩個家人的父母.
雙方父母之外.
我們的長輩在上面.
能夠邀請你去坐在同台的主家席上.
是一種榮譽.
亦都反映出你對能夠被坐在主家席上的人的尊重.
《馬太福音》第八章十一節.
二十二章第一節裡面.
同樣談論到主耶穌怎樣邀請人一同坐席.
大姐,姐妹.
能夠和基督耶穌一同坐席是一份榮譽.
是一種光榮.
我們做基督徒有沒有在這份光榮裡面.
去感受到那份被尊重.
和那份更緊密和上帝的關係.
上帝欣賞你和我呢?.
我有我信,祂有祂繼續做祂的神.

$^{721}$這種情況就是我們對上帝的信是一種交易.
我和祂沒有建立關係.
但如果你欣賞上帝能夠讓你和祂一同坐席的話.
你不要錯過.
不過很可惜的是在聖經裡面.
同時間告訴我們有些人仍然思念地上的事.
這是保羅說的.
保羅說這句話有什麼背景?.
其實保羅在提議給我們看.
什麼人只是思念地上的事.
在耶穌的年代也有一個人.
在上帝在耶穌邀請坐席的時候.
有份參與其中.
但他最後仍然是思念地上的事.
那個人叫做加勒人猶大.
在聖經裡面.
如果大家還記起耶穌基督在最後的日子.
去在耶路撒冷要預備吃逾越節的筵席的時候.
祂邀請門徒去預備這個筵席.
你記不記得?.
門徒也肯預備這個筵席.
是一份榮譽.
所以事奉不是苦差.
雖然事奉的過程可能有些事是要麻煩.
要去遷就.
要去用心思,時間,用體力.
譬如我們教會.
我們教會是什麼教會?.
招降猛賊的教會.
是麻煩.
但是你有沒有為了在這個機會裡面.
我很記得有一個姐妹.
她在她年終的日子裡面.
她跟我說一句.
牧師我很想來教會事奉.
我說你來我們教會事奉做什麼?.
我說我什麼都做.
我什麼都做.
然後她告訴我她身體不舒服.
我更加不敢叫她什麼都做.

$^{761}$我們教會做的都是粗重功夫.
又要搬桌子.
她說我什麼都做.
我搬東西都可以.
我派司鑿都可以.
各位姐妹.
真正信主的人不計這些有多辛苦.
只求上主的心思能夠得到滿足.
所以我怎麼看一個真正信主的人.
不只是純粹他講出來有多漂亮.
而是他的行動反映給我們看.
能夠有這個為分參與.
但是卻是在這個過程中.
仍然沒有離開地上一切.
加勒人猶大,你記不記得?.
甚至在最後晚餐.
跟主踢逾越節的日子的時候.
他有份吃的.
甚至耶穌醒他吃一塊麵包.
你記不記得?.
點了給他吃.
但是怎麼樣?.
吃完之後.
聖經記載.
那個邪惡的心態在他裡面發作.
接著耶穌說.
你去做你的事.
他馬上就走出去.
走出去的時候.
還包裝到被其他門徒以為他走去.
節制窮人.
各位姐妹.
教會裡面有多少這些.
在思想地上的人.
活在教會當中.
保羅在這裡就提醒這一點.
你如果要去榮耀上帝的.
麻煩你真的有一個.
跟上帝和耶穌的關係.
有一個復活的關係.

$^{801}$使得你能夠思想上面的事.
能夠有份參詳.
基督和上帝一起.
商議的拯救計劃.
不是你喜歡怎麼做就怎麼做.
教會不是你喜歡搞什麼娛樂活動.
使得別人回來教會.
就等於成功的教會.
NO!.
完全不是這件事.
是上帝和基督.
一起與你可以有份參詳的.
監察你所做的事.
你要思想上面的事,保羅.
不要只是仍然想地上的事.
這次教會失敗.
就是因為教會裡面有很多掛名的徒.
仍然在思念地上的事.
我有時覺得是很可惜的.
這類的信徒可以說是.
捉到鹿都不會脫角.
臨進入對救恩的大門口的時候.
雖然腳有半隻踩了進去.
他都要把這隻腳拔出來.
踩在地上.
在外面的世界當中.
跑自己喜歡的路.
可不可惜?.
保羅在這裡就提出這個可惜的現象.
「不要這樣,哥羅西教會」.
「弟兄姊妹,你和我都不要這樣」.
「不要這樣」.
那怎樣好呢?.
這就是第三和第四節的經文.
第三,第四節的經文裡面.
講了一樣東西很特別的.
原來我們在一項掩護之下.
來到顯揚神.
我讀一次,好嗎?.
第三,第四節.

$^{841}$請大家拿著聖經跟著我去讀.
跟著我去看,好嗎?.
好,第三節.
「正因為你們已經死了」.
「你們的生命與基督一同藏在神的裡面」.
「基督是我們的生命」.
「祂顯現的時候」.
「你們也要與祂一同顯現在榮耀裡」.
我讀多一次英文的新英文翻譯給大家聽.
「正因為你們已經死了」.
「你們的生命與基督一同藏在神的裡面」.
「基督是我們的生命」.
「祂顯現的時候」.
「你們也要與祂一同顯現現在榮耀裡」.
牧師一隻字我都懂得解釋.
我也不知道他想講什麼.
我講給大家聽這段聖經想講什麼.
「怎樣與基督在裡面一同榮耀神」.
是不是教會的主題?.
「與基督同顯他的榮耀」.
怎麼做?.
是不是「風和日麗」.
很美的環境.
然後我們搞幾堂詩歌敬拜.
或者搞幾堂報道會.
請幾個很出名的福音歌手出來唱歌.
然後大家很高興.
然後就順了主.
是不是這樣?.
似乎不是.
保羅在這裡當向完哥羅西教會.
提供了一個清楚的人生方向.
你要往上面去思想.
你要以上面主耶穌基督的心思意念.
作為基礎來作你往後人生的所有決定.
甚至他提出了這個警誡.
不要思念地上的事.
不要好像加勒人猶大那樣.
在救恩的門裡面.
甚至在教會裡面.

$^{881}$你仍然是不得救的.
不要這樣.
提出了這個警告之後.
保羅在這裡再一次強調.
我們的改變是由於與基督同死.
和先前所說的同活復活.
而塑造出來的.
當我們在這個改變之下.
我們的思想就會一直去思察上主的做法.
當我們在這樣的情況之下.
原來我們和基督一同藏在神的裡面.
這句說話精彩了.
你再不是一個人活.
你是和基督一同同活.
祂死你和祂一起死.
相聚而死.
Good.
這個正正是上帝和耶穌基督想達到的目的.
但不是完在那裡.
原來繼續在復活的裡面讓我們一同在祂的裡面.
藏身在上帝的當中.
即基督徒和耶穌一同活在上帝的裡面.
原來是有耶穌基督的藏在祂裡面.
這個字眼我要多作解釋.
這個字眼藏在神的裡面.
這個字眼在希臘文的版本.
就用了「Karupatai」.
這個字眼是甚麼意思?.
畫本的翻譯就是藏在神的裡面.
但這個字其實有另一重的意思.
除了說隱藏之外.
還可以翻譯為掩護.
知不知道甚麼是掩護?.
這個字眼首先我講出它的含義在哪些地方有說.
在舊約詩篇第27篇第5節.
和詩篇第31篇第20節裡面.
都同時有提出.
原來在舊約裡面.
原來神對祂的子民,對祂的選民.
相信祂的人是願意.

$^{921}$當他們有困厄的時候.
或者當他們要進行敬拜,事奉.
或者甚至打仗.
上帝會伸出掩護的行動.
我讀給大家聽.
詩篇第27篇第5節.
如果你有聖經都可以打開.
不是很難的.
你們手機按一按就搞定了.
詩篇第27篇第5節這樣說.
因為我遭遇患難.
祂必暗暗地保守我.
在祂情紙裡把我藏在祂帳篷的隱密處.
將我高舉在盤石上.
當私人遇見患難的時候.
原來上帝暗暗地在保守他.
將他收在上帝的房屋裡面.
保護著他.
另外一節在哪裡?.
詩篇第31篇第20節.
我們打開不如又一起讀.
不要那麼單純.
讓我讀吧.
我們一起讀好不好?.
詩篇第31篇第20節.
我找到一起讀了.
你必把他們藏在你面前的隱密處.
免得遇見人的計謀.
你必暗暗地保守他們.
在情紙裡免受口舌的爭鬥.
原來上帝當擺我們在爭戰的時候.
在困厄的時候.
上帝動手來將我們施行掩護.
擺我們藏在一個地方.
免得我們遭難.
弟兄姊妹.
原來這個隱藏在上帝面前的陳述.
有一種掩護的作用.
原來上帝讓我們今天活在地上的時候.
沒錯.

$^{961}$你會在信仰上受到很多的挑戰.
你講真理的時候.
人家會否定.
有人會討厭你講的話.
甚至好像我上個星期那樣說.
那些purpose driven.
community church movement.
包括這個所謂mega church model.
你看看普世的教會.
有多少在嚮往的時候.
這些其實是異端思維的開始.
講了之後有多少人討厭你.
極多人討厭你.
或者有人說.
以馬來尼尼不足道.
採他都傻了.
不需要理他了.
不要緊的.
你不理我不要緊.
但將來上帝理你的時候.
你出事就不要說.
我作為一個傳道人不要說.
靚姐妹.
有沒有見過打仗.
有一套電影很有趣.
叫做雷霆救兵.
Saving Private Ryan.
Ryan我很熟悉.
是我女婿的Ryan.
這套電影裡面.
講到男主角帶著一隊小隊.
在諾曼第搶灘的時候.
收到一個訊息.
說要去救一個軍人.
他叫做Ryan.
因為他的兄弟全部當兵.
只剩下他一個.
原來打仗有一個這樣的.
我不知道那個國家.
是不是每個國家都這樣.

$^{1001}$說這個美國.
美國原來有這樣的.
不想一個家庭自子絕孫.
至少仍然有一個生還.
差一個小隊.
如果我沒記錯應該是八個人.
去救一個士兵.
而他被困在敵方的深入之處.
知道他還沒死.
但是你們去找他.
最後犧牲多少個人.
去救這個人回來.
各位兄弟姐妹.
不要想著上帝.
把我們在這個世界傳福音.
叫我們單打獨鬥.
No.
上帝不是叫我們單打獨鬥.
而是基督和你一起去砌的.
基督會作為掩護.
將我們放在上帝的隱藏之處.
當他要動我們的時候.
可能那些人會很不滿意.
要襲擊我們.
事奉主的時候有人反對你的.
尤其是你按照聖經的說話去講述的時候.
連魔鬼的爪牙都來的.
哪些魔鬼的爪牙來的.
就是哥羅西教會裡面傳那些小學的人.
他們傳那些什麼節旗啊.
慶典啊.
那類東西.
他們會搞你.
你是錯的.
不要講那麼多啊保羅.
接著保羅就不講了.
算了.
不要得罪人多.
稱呼人少.
留得青山在.

$^{1041}$哪怕沒有柴燒.
這些很多基督徒都經常講的.
傳福音.
不要得罪那個人.
以後都沒有機會跟他說.
我告訴你.
你不用怕得罪那個人.
如果那個人有聖靈在他裡面的時候.
你今天講的東西就算他反對.
將來他會回心思想.
我不想講第二個.
我只是用我自己做見證.
我信主其實是一個神蹟來的.
你講得這麼厲害啊牧師.
是啊.
是一個神蹟來的.
你知道我第一次接觸基督徒.
我多討厭基督徒啊.
又說我有罪.
我沒有罪的嘛.
我這麼勤力讀書.
我很記得我第一次傳.
有人跟我傳福音.
就是我在香港大學的學生會的一個飯堂裡面.
溫書的時候.
有一個大學生走來跟我說.
你有罪啊.
我多討厭他啊.
接著他還給我那本《屬靈4定律》.
我拿完後扔了垃圾桶.
後來我想了想.
不如這樣吧.
先撿起來.
從垃圾桶拆出來.
然後拿回去給一間天主教的神父.
我不怕告訴你哪個神父.
香港前教區總主教.
我給他看.
你猜他跟我說什麼.
不用理他們.

$^{1081}$如果我要不信主.
我大可以不信主.
加上我有香港天主教.
總區總主教.
肯定我.
不用理他們.
後來我來了加拿大.
我信了主.
回去見到這位主教.
你猜他跟我說了什麼話.
厲害了.
信了耶穌了.
你作為天主教的總主教.
我信了耶穌.
不是應該開心嗎.
他覺得我好像背叛了他.
其實不是.
我很尊重這位長者.
我的老師.
為什麼會有這樣的改變.
不是一個人.
而是上帝的說話.
當你打開聖經的說話的時候.
你會見到自己的罪是怎樣一回事.
今天你說他未必喜歡聽.
他憎恨你.
不需要擔心的.
如果有聖靈在他裡面內住.
他會回心思想的.
如果你說他憎恨你.
然後他以後沒機會再聽福音.
Good.
這就是上帝命定他不得救.
為什麼你不可以讓上帝去決定呢.
是誰人得救就誰人得救嗎.
所以教會今天有很多.
是人塑造出來的假得救人士.
當我們去事奉的時候.
原來耶穌和我們一起.
他就好像那套戲裡面的男主角一樣.

$^{1121}$他在任何的時間裡面.
為了去救那個唯一家庭死剩的士兵的兒子.
他叫其他的隊友一起掩護他去打.
向前去救.
雖然他受了很多的傷害.
在很大的驚慌裡面.
你看回那套戲.
但是在種種掩護之下.
幫助他的團隊.
終於救回那個家庭唯一死剩的士兵.
看完節目.
這段聖經說了什麼.
當我們能夠與基督一同得到上帝的掩護之下.
基督就是我們的生命.
他顯現的時候.
他和我們同在的時候.
我們就一同顯出上帝的榮耀.
上帝的目的.
基督為什麼要和我們一起.
因為我們一起去履行上帝的目的.
今天很多的信徒很擔心.
我事奉我受傷害了.
那些人不喜歡我.
那些人不喜歡我就不喜歡聽福音.
那些人不喜歡你.
不代表他一定不喜歡那個福音.
只不過上帝要那個人怎樣去喜歡那個福音.
怎樣去接受那個福音.
這是上帝的靈的工作.
而保羅對於哥羅西教會.
作出這一連串的行動.
是要告訴他們知道.
我們作為教會的一份子.
是和基督一同去榮耀神.
而且基督才是那個主力.
因為他就是使我們可以從死.
以至復活一同同在的那位主.
這就叫真正信耶穌.
任何時間都與主一起.
而主耶穌成為那個事奉的主力.

$^{1161}$來讓我們可以參與其中.
我們只不過是一個輔助形式.
我們的作用只是在反映上帝這個工作.
有機會這樣參與.
應該是一份光榮.
為什麼引而我講.
我小時候就在玩那些鏡子在反射陽光.
那塊鏡子不會發光.
真正的能力在那個太陽身上.
我們的屬靈的能力是在上帝父神身上.
透過耶穌基督讓我們可以參與.
而我們可以反映出來.
就好像夜晚的月亮一樣.
前晚我記得師母跟我說.
快點快點來.
我說什麼事.
我正在寫講章.
很忙.
打擾我思想.
不是很重要.
快點來看天的月亮多美.
當你看到那麼漂亮的月亮的時候.
你會不會對那個充滿力量的太陽.
有更大的讚歎呢?.
作為信徒的.
就是應該這樣.
與主一同同功.
當主跟我們的生命同在.
當主顯現的時候.
就是我們跟祂一起去反映出.
上帝的榮美.
和上帝拯救那份心願.
而且上帝的光榮比這個還要多.
怎樣可以知道更多.
將來逐漸逐漸顯現給我們知道.
我發覺今天很多的教會其實是打錯仗參錯戰.
他以為在這個地上只要有本聖經.
然後他以為在這個地上進行的征戰.
就是要在這個地上爭取教會的權益.
爭取教會在這個世界上的影響力.

$^{1201}$爭取在甚至這個世界裡面搞所謂的公義的抗爭.
或者在這裡搞很多對於教會不合法的法律的一些抗爭.
搞聯署.
對不起.
這些全部是off.
大家足球會知.
這叫offside.
你根本錯了位.
越位.
真正我們在打一場戰爭.
是針對罪的戰爭.
這個就是上帝要我們和基督一起打的仗.
沒錯.
搞聯署.
你可以寫一封信給你的國會議員.
去對一些法律的意見的不同意.
但是你和我絕對不能阻止到這個世界走向罪惡這個傾向.
我不是在這裡做一個政治宣言.
不過我很討厭Liberal Government.
為什麼.
因為他們是最喜歡擦鞋的政黨.
第二個我最討厭的就是NDP.
是一個由基督徒一個牧師所創立的一個政黨.
完全是以討好人心思成為他們政黨的核心.
現在另外剩下的那些黨.
有多少還有腰骨的.
沒有多少.
隨時過黨.
所以我罵到另外一個政黨的領袖狗血淋頭就是這個原因.
為什麼要這樣說.
因為作為基督徒要講的是真理.
要講的是事實.
而不是和別人去喊殺.
你也改變不了那些人的邪惡.
但是上帝卻把我們放在一個邪惡的世界裡.
在這個邪惡的世界裡面上帝給我們掩護.
讓我們在這個邪惡的世界裡面學會怎樣參與.
當有任何攻擊的時候.
他來掩護.
掩護我們去履行上帝的目的.

$^{1241}$而今天的教會說不需要掩護.
上帝只需要將所有的惡行摧毀就行了.
我告訴你上帝不會這樣做.
上帝要使這個世界越來越可怕.
容許它可怕.
讓人的罪擴展到一個可怕的地步.
然後末期就來到.
所以那些祈禱.
祈求上帝世界和平.
我告訴你對不起這不是上帝的心思.
我不是要去打仗.
但是我知道這個世界很自然就會去打仗.
沒有和平.
而當這個世界盡顯沒有和平.
盡顯沒有公義的時候.
你和我在這個世界.
就是要作這個見證.
哥羅西教會.
你有沒有這樣做.
在教會裡面有這樣的人.
你有沒有抵抗他們的能力.
至少你不參與就其中.
有不少的人在過去我這麼多年的事奉裡面.
經常都認為以馬來理華人浸信會是一群不跟著潮流走的教會.
很不合群啊以馬來理華人浸信會.
不過在我自己對聖經的理解.
和作為一個信徒.
我明白到.
上帝不是要一間只是一間教會.
是期望得到這個世界認同讚許的教會.
才是一間合神心意的教會.
上帝喜歡的教會是要有他的樣式.
在他的啟示裡面所塑造的教會.
而保羅就在這裡將這一點提醒了.
哥羅西一間已經有很好的信仰的教會.
仍然需要參與在這場屬靈的爭戰裡面.
這場的爭戰可能是非常激烈.
甚至可能會有損傷.
但是神卻是透過基督施予掩護.
而這個掩護是讓我們更加準確去實踐上帝的心事.

$^{1281}$所以掩護不是讓我們逃離困厄.
而是在困厄裡面行上帝的心事.
弟兄姊妹.
這樣才能夠榮耀上帝.
你們明白嗎?.
一起同心禱告.
父神多謝你.
多謝你給我們在這段聖經裡面保羅的啟示.
多謝你讓我們知道你透過基督.
和我們一起去打這場屬靈的爭戰.
多謝你透過主耶穌基督來締造各式各樣的掩護.
好使我們在重重挑戰的當中.
仍然沒有失去事奉的焦點.
就是上帝你的說話.
拯救的心願的展現.
求你的帶領.
讓我們在這個過程當中有毅力.
能夠緊緊捉著基督的掩護.
來完成上主你的心思.
今天的教會我們同樣為祂禱告.
有多少信徒在這些事情上扭曲.
以為只要脫離困難困厄.
就等同於上帝你的榮耀彰顯.
但我們卻看見的是.
真真正正讓我們得到上帝的心思意念的彰顯.
正正就在這些困厄的過程當中.
求你幫助我們.
納準我們的心思靈性.
在你的真理啟示之下.
生根奠造.
帶領每一個願意走這條屬靈道路的信徒.
都能夠在基督的復活.
一同邁向一條榮神的道路當中.
禱告奉靠主耶穌基督得勝名字而求.
\newpage



\section{約書亞記 21:2}
\label{sec:032182}
\textbf{聖經人聖經事:利未的城邑}~~~~黃克勤牧師博士
\newline
\newline
連結: \href{http://www.church.com.hk/acms/content.asp?site=cdc\&op=show\&type=product\&code=032182\&layout=sermon}{\texttt{032182}} ~~~~ 語音日期: 2026-05-23
\newline
\newline
\hyperref[sec:032181]{\small{< < < PREV SERMON < < <}}
~
\hyperref[sec:index_chronic]{\small{[返順時目]}}
~
\hyperref[sec:032183]{\small{> > > NEXT SERMON > > >}}
\newline
\newline
約書亞記 21:2
\newline
\begin{longtable}{cl}
\hline
\hline
章節 & 經文 (和合本修訂版)\\
\hline
21:2 & \begin{tabularx}{0.7\textwidth}{X} 在迦南地的示羅對他們說：「從前耶和華曾藉著摩西吩咐給我們城鎮居住，以及城鎮的郊外供我們牧養牲畜。」 \end{tabularx} \\ \\
[1ex]
\hline
\hline
\end{longtable}
$^{1}$大家好,今天跟大家講利味的誠邑.
不是說利未人沒有地嗎?.
在以色列十二支派分了迦南地,應許之地的時候.
利未人說沒有地,因為他們沒有地的產業.
因為他們服侍神,代表著百姓和上帝連結.
幫他們做祭司獻祭,或者處理所有靈修敬拜的事務.
這是一個很寶貴屬靈的產業,好過有塊地.
所以利未人分佈在所有十二支派裡面,沒有自己的地.
這樣就很不值得,如果你是利未人.
會不會覺得自己流離失所,到處求人?.
不是的,在這裡是處理了的.
生命樹書的網絡跟大家一起看聖經.
聖經人,聖經事,這次聖經的事有利味的誠邑.
就看看神如何處理這件事.
誠邑就說屠城,神如何處理那些意外無愁無怨殺死人.
給他們一個屠城,有機會獲得生命,也好.
約書記也有一章《那時》,這次就看看利未人.
利未人的眾族長,利未人有很多人,十二支派之一.
利未人的後代,利未人的眾族長來到祭司以利亞薩.
和亂的兒子約書亞,並以色列各支派的族長面前.
這裡出了一些以前沒什麼留意到,但他們開始上位了.
是什麼呢?以往你看到摩西帶領以色列人.
跟亞倫一起,後來亞倫先死了.
然後就到摩西,摩西也在尼泊爾去世了.
交給約書亞帶以色列人入迦南.
打了很多場仗,平定了,十二支派也分地了.
一開始大家還不分,約書亞說,神給你們的地.
你們在這裡擔言不去拿,等到什麼時候呢?.
於是剩下未有地的七支派就去分了地.
十二支派都齊了,說明利味支派是沒有地的.
為什麼呢?因為他們要處理敬拜,靈修等等的事.
因此沒有地的,敬拜的處理帶領所有以色列人敬拜.
這些事情已經是一個很寶貴的產業,沒有土地的產業了.
也要住的,這裡,利未人的眾族長就來到了.
不只是找約書亞,亂的兒子約書亞.
為什麼要特別說亂的兒子呢?因為當時的人,再說一次.
因為當時的人是沒有姓的,只有一個名字.
要書亞,要書亞,要書亞,要書亞,約書亞.
約書亞其實跟耶穌的名字是一樣的,要書亞.
跟他的希伯來文名一樣,要書亞.

$^{41}$要書亞救主,我們很少改自己的名字.
什麼名字,我叫王救主,這樣說是不行的.
當時是很多人歡喜跟從神去幫助別人的.
要書亞就是這個名字了,所以他一定要說亂的兒子.
或者捕魚的約書亞那些,亂的兒子約書亞.
他是領袖,進去平定地方之後,各樣敬拜都開始了.
士羅做他們的首輔,於是回望約季又在士羅.
這個時間出現了一個大祭司,祭司的頭叫以利亞薩.
約書亞處理領袖,處理軍事,政務.
而處理屬靈的事,就由大祭司,或者祭司長.
或者祭司的頭,以利亞薩去負責了.
另外在這麼多地方,這麼多人口.
現在就要以色列各支派的族長一起處理了.
因此他們很多時候都會在一起的.
可能不知道會不會有定期的會議都不一定.
多久見一次了,要記得事情要處理了.
這次要處理的是利未人眾族長提交的一件事.
一個申請,申請是什麼呢?.
於是他們一起,祭司長.
為什麼要祭司長呢?因為說到利未人嘛,敬拜的嘛.
一定要找他們的頭出來,就拿來講了.
祭司長,聯的兒子,約書亞他們的領袖.
然後各支派的族長一起決定了.
再加上南地的士羅,就是他們的首都,當時的在這裡.
對他們說,第一點就是對他們說,誰對他們說呢?.
第一點對他們說是利未人的族長.
利未人的眾族長,不是一個.
就去到找他們,就在士羅這個地方對他們說.
這個他們指誰呢?.
這個他們就是指祭司以利亞薩,約書亞和各支派的族長.
對他們說,為什麼特別要這樣都成為第一點呢?.
其實我們有些訴求,有些東西是需要的.
哎呀,埋在心裡,哎呀,這樣怎麼行呢?.
可能有些利未人就會,家人會說.
爸爸,或者老公,為什麼我們沒有地的?.
哎呀,我認識的朋友,親戚,全部都有自己的地的.
我們利美族沒有,真是慘了,可能這樣會有壓力的嘛.
希望交在神手中,希望他們知道,這樣也好,很靈性.
不過,不如你把心裡的說話,直接對他們說出來.
那你也要表達一下,有自己個人很真誠的需要.

$^{81}$不妨表達出來,說出來.
怎麼處理,能不能申請,行不行,也要說出來.
於是他們也挺好的,就去跟領袖說,跟祭司罪代表說.
跟領袖約書說,跟各族的族長說.
因為各族裡面都有利未人,所以利未人的代表就來說.
對他們說,不要在心裡,不要以為別人會知道.
不要以為別人會懂得做,是要說了才知道.
哎呀,原來這樣,是的,沒有想到,行行行,我們先處理一下.
從前耶和華藉著摩西吩咐.
摩西的時候,神透的摩西吩咐是說了什麼呢.
是給我們誠邑居住的,都說你沒有地了.
我不是要地,利未人是說,以前神透的摩西是告訴我們.
我們是有誠邑可以居住的,給我們一些誠,不是地.
給我們住的,而且誠邑就包括了誠邑的郊野.
就因為我們都要牧養的,我們都要牧羊的.
都要養牛,放羊的,是來牧養我們的牲畜.
他們很真誠,說出一個以前上帝透過摩西的吩咐.
就是說,會有誠邑,不是地,有誠邑給我們利未人的.
而且包括了誠邑外面的郊野,讓我們可以放牧.
哦,這樣是喔,有些人可能就說,是啊,我都記得真是這樣說的.
哦,這樣是吧,他們很好啊,我說不用了.
我們給你屋子住就算了,他們很好啊.
其實是很難做的,譬如如果我已經分地了.
有些人在城裡面住著,現在城要給你.
就是說,我要搬遷他們出來去別的地方.
一定會被人罵得糊塗,一定會被人說,怎麼搞的,你這樣怎麼做事的.
一早點又不說,可能是有很多困難的.
不過,那班的領袖很好啊,以色列人的族長啊.
祭司,以利亞薩啊,代表啊,約書啊這樣.
於是以色列人呢,大家同意,第二.
照耶和華所吩咐的,哇,這一點真是令到我很感恩,很激動.
就是在說土地誠一喔.
竟然間,以色列人聽到之後,好啊,既然是耶和華的吩咐.
好,我們就照耶和華所吩咐的,留意這幾個字,紅色的.
從自己的地業中,將以下所寄的誠一和誠一的郊野給了利未人.
以下所寄的就是下面跟著來,那些地方.
不如我們就將我們,哇,這幾個字真是要命啊.
從自己的地業中,將一些我們的城分了給利未人.
那我們的人怎樣啊,那我們就可能逼一些,要去別的城,搬屋了.
哇,不是吧,你搬得了,就是很生氣的嘛.

$^{121}$是給了我住的嘛,剛剛才佈置好,你又說要收,又說要搬遷,不是吧.
哇,很大件事,如果今天很大件事,當時都大件事,都不舒服的啦.
不過,既然是耶和華所吩咐的,好啦,我們就照耶和華所吩咐的.
以色列人照耶和華所吩咐的,哎呀,我心裡見到真的很感恩.
這個第二點就是聽話,於是以色列人聽話.
從自己地業裡面撕那些城出來,哇,真是痛啊.
那就將下面記錄了的誠一和誠一的郊野給了利未人.
好,參考,藍色字參考,參考初一及記第六章.
為什麼呢?因為我才想跟大家讀些人名而已.
第六章記載了初一及記第六章十六節,利美種子的名字.
是按著他們的後代記載下面.
大家知道利美的爸爸是雅各,雅各有十二個兒子嘛.
劉變西面,利美尤大,尤大第四,利美第三.
劉變即是呂變第一,西面第二.
記載下面,利美有什麼兒子的,這些是雅各的孫子.
利美的兒子有甲順,有哥核,然後有米拉尼,三個兒子.
而利美一生的壽數是一百三十七歲,這一百三十七歲裡.
利美生了三個兒子,就是甲順,哥核和米拉尼,三個人.
接著我們見到這三個人名,甲順的兒子按著家室.
有納的兒子有納尼和侍妹,哥核的兒子有暗蘭.
紅色字的,因為暗蘭和一陣子有關係.
以斯哈,希伯倫,希伯倫不是地方來的嗎?是名字來的,也可以做名字.
有烏洩,哥核一生的壽數是一百三十三歲.
利美的兒子哥核生了一個孫子,叫暗蘭和這些.
暗蘭特別紅色字有關係.
米拉尼的兒子就是茉莉和武氏,也是利美家的.
按著利美的後代數出來的.
第十節,暗蘭娶了她父親的妹妹約基比為妻.
這樣可以嗎?兒子娶了爸爸的妹妹.
我們突然間覺得少,不過合法的,在當時也是合法的.
也不覺得太突然,因為當時的人口不是很多.
男女結婚或者是合理的,如果沒有婚姻的關係.
她跟別人結婚也可以.
暗蘭娶了她父親的妹妹約基比為妻.
暗蘭生了阿倫和摩西,阿倫和摩西的爸爸就是暗蘭.
暗蘭的爸爸就是哥核,哥核的爸爸就是利美.
利美的爸爸就是雅各,雅各,利美,哥核,暗蘭.
然後就到阿倫和摩西,阿倫是摩西的哥哥.
暗蘭的歲數是一百三十七歲.
這個紀錄是這樣的,是家譜,族譜,然後那代人的首領就是暗蘭.

$^{161}$這一族就是他生存了多久.
於是生了阿倫和摩西,所以阿倫和摩西是利美族.
是做祭司的,阿倫後來成為以色列的大祭司.
為了哥核族尖溝,為了哥核這一族利未人,跟他們抽籤.
抽籤什麼城可以給他們,留意留意.
他們真的很尊重利未人,也很聽城的吩咐.
其實如果你要將你的城給別人,就是心疼的.
是親戚來的,利美族都是我們雅各的後代一家親.
沒有的時候就是一家親,有的時候就是我給你一個城,好痛啊.
給你一塊餅也痛,不要說城了.
為哥核族尖溝,就是抽籤,看看我們給他們什麼城.
這一族的利未人,利未人的祭司,阿倫的子孫.
從猶太大支派,西面支派,變雅敏的支派.
哥核族抽到比亞倫族的城.
從哪裡拿出來的?猶太大支派,西面支派,變雅敏的支派.
在他們的地獄裡,按尖溝得了多少座城?.
十三座城,一座算了.
教區也是他的,城可以種田,畜牧的地方也給他.
利美族專心唱歌,專心獻祭,搞好敬拜就可以了.
用不著地啊,來吧,你和我買,我們奉獻錢給你的.
十分一給你的,你和我買牛羊就可以了,不用自己養.
要什麼城,我給你地方住,我開幾間屋給你住.
是神吩咐給城和教區牧羊的,牧放牧的.
第二點,他們照耶華的吩咐,我最感恩,最欣賞的地方就是這裡.
哥核族分了十三座城,哥核其餘的子孫.
哥核的亞倫族人在三個支派裡拿了十三座城.
哥核族的其餘子孫在以法連旦和馬來西亞半個支派地又分了十座城.
還是哥核的後代已經分了二十三座城?.
二十三座城,很慷慨,也很聽神的吩咐,很捨得,很愛這班利未人.
另一個兒子除了哥核之外,甲舜的子孫也是利未人.
在以撒加支派,亞瑟支派,拿布泰利支派.
拿到住在巴山的馬來西亞半個支派的地業裡又抽籤了十三座城.
甲舜的子孫又拿了十三座城.
米拉尼的子孫按著他們的宗室來拿.
里面支派,加德支派,西普倫支派裡面又抽籤給了十二座城.
對於這三個利美族一共給了很多個城.
以色列人照耶和華籍摩西所吩咐的.
將城邑的郊區尖岸,尖溝分級利未人.
特別用螢光的字很欣賞.
竟然將自己的地按耶和華的吩咐給他們.

$^{201}$給城和城的郊外.
我住了的人要請他出來他就要住了.
城已經沒有了肉都赤了.
但是他們這麼照耶和華吩咐.
所以神很欣賞很祝福他們.
約書記二十一章四十一節.
終於很長,聖經簡一點,到這裡一跳.
利美總共在以色列人的地業中所得的城.
第三點,共四十八座並屬城的郊野.
利美族人說我們照著摩西當日的吩咐.
我們可以拿到城邑來住,我們也很需要.
而且還有郊外的地可以放牧.
以色列人說好啊,既然是這樣,不要緊.
我們一起眾族長來抽籤.
我們這麼多個城全部寫在瓦片上.
城一個城一個名,扔下去然後抽.
這個給哥葛的亞倫族,這個給哥葛的其他人.
這個給米利拉,這個給甲遜.
四十八座城,很慷慨.
他們這麼肯奉獻給利未人住.
其實是奉獻給神.
他們是照顧這一群在他們當中擔任牧職.
獻祭的祭司家族.
哇,很好啊.
這些城四處都有屬城的東西交.
城城都市如此,四十八座城全部都有交東西給.
用紫色字表示重要.
在這裡我特別想講第三點,共四十八族.
所以真的很照顧利未人.
我想神的心是很開心的.
他們為了上帝喜悅.
他們照顧以色列人照顧城的吩咐.
於是利美的城邑有四十八座.
連旁邊的郊野.
哇,好啊.
為了這群以色列人心裡很感恩很激動.
他們真的很聽神的吩咐.
聽話,神是祝福的,一起祈禱吧.
多謝我們的主.
我們讀聖經心裡是很激動的.

$^{241}$因為要將自己的城交給利未人.
免費的,城還得了旁邊的郊區.
可以放牧.
竟然以色列十二支派給了利未人四十八座城.
看到他們的慷慨,看到他們對神的尊重.
看到這個既然是耶和華神的吩咐.
他們就做了,好捨得.
今天有時我們奉獻多五毫子.
多五元心裡都很赤.
他們是將自己有的城要送給利未人.
免費獲得,而且還要將自己裡面的人處理安排他們調遷.
主說這些都是很難的,不過他們不怕難.
他們就這樣將自己有的城一片給主.
是給利未人用.
哇,好感恩,好激動,是好欣賞.
主願意我們都學習到當日的以色列百姓.
對主的吩咐是這樣遵從.
聖經講了幾次,他們是照耶和華吩咐的.
幫助我們都學習,是照耶和華的吩咐去行事.
將你的命令給我們,將你的心意給我們.
我們都這樣照著你的吩咐去實行.
給我們力量,垂聽祈禱,奉耶穌基督名求,阿們.
\newpage



\section{創世記 4:3-9}
\label{sec:032183}
\textbf{活出召命——生命的覺醒與承擔 (華)}~~~~梁明財牧師
\newline
\newline
連結: \href{http://www.church.com.hk/acms/content.asp?site=cdc\&op=show\&type=product\&code=032183\&layout=sermon}{\texttt{032183}} ~~~~ 語音日期: 2026-05-24
\newline
\newline
\hyperref[sec:032182]{\small{< < < PREV SERMON < < <}}
~
\hyperref[sec:index_chronic]{\small{[返順時目]}}
~
\hyperref[sec:032184]{\small{> > > NEXT SERMON > > >}}
\newline
\newline
創世記 4:3-9
\newline
\begin{longtable}{cl}
\hline
\hline
章節 & 經文 (和合本修訂版)\\
\hline
4:3 & \begin{tabularx}{0.7\textwidth}{X} 過了一些日子，該隱拿地裡的出產為供物獻給耶和華； \end{tabularx} \\ \\ \relax
4:4 & \begin{tabularx}{0.7\textwidth}{X} 亞伯也把他羊群中頭生的和羊的脂肪獻上。耶和華看中了亞伯和他的供物， \end{tabularx} \\ \\ \relax
4:5 & \begin{tabularx}{0.7\textwidth}{X} 卻看不中該隱和他的供物。該隱就非常生氣，沉下臉來。 \end{tabularx} \\ \\ \relax
4:6 & \begin{tabularx}{0.7\textwidth}{X} 耶和華對該隱說：「你為甚麼生氣呢？你為甚麼沉下臉來呢？ \end{tabularx} \\ \\ \relax
4:7 & \begin{tabularx}{0.7\textwidth}{X} 你若做得對，豈不仰起頭來嗎？你若做得不對，罪就伏在門前。它想要控制你，你卻要制伏它。」 \end{tabularx} \\ \\ \relax
4:8 & \begin{tabularx}{0.7\textwidth}{X} 該隱與他弟弟亞伯說話 。 二人正在田間時，該隱起來攻擊他弟弟亞伯，把他殺了。 \end{tabularx} \\ \\ \relax
4:9 & \begin{tabularx}{0.7\textwidth}{X} 耶和華對該隱說：「你弟弟亞伯在哪裡？」他說：「我不知道！我豈是看守我弟弟的嗎？」 \end{tabularx} \\ \\
[1ex]
\hline
\hline
\end{longtable}
$^{1}$好,弟孫姐妹平安!.
我跟各位來看一個聖經的人物,從他的生命當中來提醒我們,給我也給各位弟兄姊妹這個提醒..
這個人物可能是一個很不出名的,當然可能也很出名的吧,他真的是算是在聖經裡面的一個小人物..
可能有人我們喜歡改那個英文名字,我猜沒有多少人會用他的名字作為你的英文名字,我們大概都想改那個什麼David,Paul,Peter那些,大概沒有人會用這個人物的英文作為他的英文名字..
他的兄弟的那個名字就有些人也會用,我不曉得你會不會記得有誰是用他的兄弟的那個英文名字,有些人是認識的..
剛才這個題目是"活出照命",是我們從這個人生命裡面的覺醒跟承擔..
我不曉得你還記得剛才隱姨帶領我們在開始的禱告裡面,她曾經求神幫助我們,讓我們都有一個怎麼樣的心..
她求神給我們有一個清醒,可能我們有點都還想睡覺,還沒睡醒,是讓我們有一個清醒跟喜樂的心..
如果你還記得剛才隱姨主席帶領我們禱告的時候,就是求神給我們..
其實這個今天講到覺醒,其實就是讓我們怎麼樣去保持我們有一個清醒的心,這個是對我們每一個人的生命都很重要..
有時候我們真的像還在睡覺的那個狀態裡面,有時候可能性真的是有點像在喝醉酒或者是那種狀態裡面..
所以我想今天從聖經人物裡面幫助我們去怎麼樣多做一個生命的覺醒跟生命裡面的承擔..
好,我們來看這個人物是誰了.我會分兩次在崇拜講到裡面講,所以今天不是會講完這個經文的..
我們在下下禮拜以後,六月第一個禮拜天,我還會在崇拜當中繼續講.所以今天沒講完,大家放心.從這個小人物身上,我想是幫助我們有這樣子的提醒,.
也幫助我們去做一個,我怎麼樣說,這是一個我叫做反面的教材.它是一個不太好的生命,所以就讓我們從這個人不好的生命裡面怎麼樣成為我們一個很重要的提醒..
是一個反面的教材,不是叫各位來學這個人的生命,是從這個人的生命我們能夠有更多的提醒..
好,這個人叫什麼名字?蓋尹.蓋尹是誰?是亞當跟夏娃的第一個兒子,所以有一個叫他的兄弟.待會我們會認識他多一點..
我們看到,我邀請各位來讀一次前半段,我們今天要講的大概的經文,《創世記》第四章三到九節,我們來認識一下蓋尹,也認識一下他的兄弟..
看到了,我邀請各位,我們一起來讀這段經文.蓋尹拿地理的出產為貢物獻給耶和華,亞伯也將他羊群中投生的和羊的脂油獻上..
耶和華看中了亞伯和他的貢物,只是看不中蓋尹和他的貢物.蓋尹則大大發怒,變了臉色..
耶和華對蓋尹說:你為什麼發怒了?你為什麼變了臉色了?你若行得好,豈不蒙怨那?你若行得不好,罪咒伏在門前.他便怜慕你,你卻要制服他..
蓋尹與他兄弟亞伯說話.在庭間,蓋尹起來打他兄弟亞伯,把他殺了.耶和華對蓋尹說:你兄弟亞伯在那裡.他說:我不知道,我只是看守我兄弟的..
蓋尹有一個兄弟,蓋尹其哥哥.他的兄弟是誰?亞伯.你認識誰的英文名字叫亞伯的?你不認識?在教會久一點的應該認識的.亞伯的英文是什麼?Abel,A-B-E-L.誰是Abel?最有能力的牧師,林培喬牧師..
他大概十一月的時候會回來,所以你就多認識他一點.他是我們在普羅馬俗拜跟我們一起服侍一段很長時間的一位牧師,他就叫Abel,是亞伯..
今天我們主要是講蓋尹,我們大概只會講一點點的亞伯.我們都知道蓋尹是哥哥,他有一個弟弟叫亞伯..
潘王基在我們當中能夠從蓋尹身上做得不好,他是一個生命不好的人,我們就需要從他的生命,從他的身上來學會,我們應該有那種生命覺醒的操練..
蓋尹,我們認識一下他的心態跟行為.他是一種怎麼樣的心態跟行為呢?只是看不中蓋尹和他的供物,蓋尹就大大的發怒,變了臉色..
這個變了臉色的意思,大概對我們中國人來說我覺得是很有意思的.變了臉色的意思跟什麼有關?跟你的臉有關..
他不是說他的臉色當然也會,可能當一個人發怒的時候,他的臉會發紅,不是發光,可能是小心一點老血管充血,那就糟糕了,他的臉會發紅..
其實他也不是這個意思.他不是說講到發怒的時候,蓋尹的臉可能都會發紅,但最終他想表達的意思是什麼呢?他其實有另外一個更準確的意思是說,他的面子,他的臉掉下來..
就好像女生可能比較多一點,男生也有,有時候我們都會去保護我們的臉要放一副面膜在臉上,我不是要讓你看到你太太和你身邊的女生戴了面膜的突然之間掉下來,會有這種感覺,還是你自己經歷過..
這個臉就像他的人的面子,不是那個面膜.所以,他的變了臉色的意思是什麼?他的面子掉下來是什麼意思?沒面子,他是大大的發怒,其實最主要是因為他沒面子,他的面子掉下來..
這個是為什麼呢?剛才我們說蓋尹是哥哥,亞伯是他的弟弟.但當耶和華在接受他們的限制的時候,我們回去再看那個..
耶和華神悅納他看中他喜歡誰所獻的職務,他喜歡看中亞伯的,他喜歡悅納亞伯所獻的.他沒有看中,也沒有悅納蓋尹..
所以這個對蓋尹來說就是,他是那個長子,他是哥哥,他是本身在家裡面,或許在當時候他應該是一個蒙神悅納,他是家裡面的代表,但他所獻給耶和華的職務沒有被悅納..
所以他第一個問題出來了,就是他覺得他自己沒有面子.這個有點,我們簡單說就是丟臉吧,他臉真的是掉下來,他沒有面子..
所以這個沒有面子就讓他生氣了.所以在他臉色有沒有變紅我們不要管,他就覺得他自己沒有面子的時候,他就發怒,他就生氣..
所以第一個是,他覺得他自己是長子,他是哥哥,但他沒有被神悅納他所獻的職務,所以他就覺得自己沒有面子..
所以他另外一個反應就變成,他的怒氣來了,而且這個怒氣是在聖經形容,他沒有受到控制.這個沒有受到控制就是說,他讓這個怒氣一直發作下去,一直泉流在他的心裡面,他沒有覺醒,他沒有保持清醒的狀態..
他反而是在讓怒氣放在他的生命裡面,他沒有察覺這個是一個成為罪他的問題,而且他那個怒氣是沒有受到控制,而且帶來的是他後來犯了殺人的罪..

$^{41}$在舊約講到該隱的時候,其實沒有多少的經文有講到該隱.我們反而是在新約的時候,在約翰一書裡面,把它引用了該隱這個例子,幫助我們學一個很重要的功課..
約翰一書的第三章,我想邀請弟兄姊妹們來讀,我們來讀這一段的經文,預備來."不可上該隱,他是屬拿惡者,殺了他的兄弟.為什麼殺了他呢?因自己的行為是惡的,兄弟的行為是善的.凡恨他弟兄的,這是殺人的..
你們曉得,凡殺人的,沒有永生存在他裡面.".
他一直走在一個錯誤的路上..
約翰一書的主題是什麼呢?如果大家有機會,回去看新約的約翰一書.約翰一書我們很熟悉的一些經文,特別是在剛結婚的..
約翰一書有引用那個約翰一書的經文,我不曉得你還記不記得,應該記得吧,那個新郎.約翰一書引用的是神是愛.約翰一書其實有一個很重要的講道,神是愛..
我們作為基督徒就要在神的愛裡面彼此相愛,不單只是丈夫要愛妻子.他們引用那個經文,其實他更要強調在教會裡面,我們都是弟兄姊妹,在基督裡面我們都要彼此相愛..
所以,約翰一書裡面講到的主題是神是愛,他引用該引作為一個例子,是要警告我們說,該引他走了一條錯的路..
我們待會會看到,該引用了一個很特別的例子來跟神對話.當該引控制不了祂的怒氣的時候,祂就犯罪了..
所以,每一個人我們都是真的有七情六慾,我們都會有生氣的時候.但我不曉得各位,你生氣的密度多少?可能你對不同的人有不同的一種對待,有些人對某些人可能是特別容易生氣,有些人他就覺得比較容忍..
但這裡有點提醒我們就是說,該引祂的生氣發怒的時候,讓祂去犯罪的時候,是不是我們就很危險呢?.
我們在當中會不會看到,有時候有些人要貼一句話,那句話蠻多在不同的場合會發現的.最近我在教會哪一個門口也看到,有一句話叫"莫生氣"..
有些是寫了一段很長的話,就是在餐館,餐廳,我有時候也看到有些人在"莫生氣",然後是怎麼怎麼的.是不是我們做人的時候,有時候跟人相處的時候,不要生氣了,生氣有時候會很危險,就會變成怒氣,發怒的時候控制不了,就有很危險的事情可能會發生,有可能上該引..
所以,我們做基督徒,尤其我們做基督徒,我們有時候都有點覺得別人知道我們是基督徒的時候,就說"哎,你不要那麼生氣,你是基督徒,所以我不想你給人家說過你是這樣子.".
所以,在我們的生活當中,我們是不是不能夠生氣呢?我們做為基督徒,尤其是應該要和和氣氣,要溫柔一點,不要很發怒,是不是我們就是要這樣子"莫生氣"?.
當然,我們有一些情緒,這個人最常情,每一個人都有,只是我們要怎麼樣去好好的跟這個情緒相處,或者說怎麼樣去好好的去面對我自己的情緒..
所以聖經怎麼樣說呢?我相信在我們有機會讀聖經的時候,我們也可以看到聖經在莫生氣的教導是怎麼樣的.我跟各位來看一段經文,一份所述的第四章二十五到二十七節..
如果有人告訴你,基督徒不能生氣的話,你大概可以跟他說這段經文.我們還是人,我們還是有時候會受到一些情緒需要表達的時候,我們還是可以生氣..
其實生氣是要怎麼樣去好好的去控制也好,或者說怎麼樣去表達,也是一個我們要學習的.我們來讀一次怎麼樣聖經教導我們在生氣,在發怒的事情上,我們應該怎麼樣做?.
所以,你們要棄絕謊言,個人與靈色說實話,因為我們是互相為肢體,生氣卻不要犯罪,不可含怒到日落,也不可給魔鬼留地步..
生氣就應該停了,不要再那麼常放在心裡面那麼久.如果該隱,他真的是把他的怒氣放在他的生命裡面,他沒有放下他的怒氣,沒有放下對弟弟的那種生氣..
所以,這個在我們聖經的教導也是很清楚告訴我們說,我們還是可以有發怒,我們還是可以有生氣的時候,當然我們也要小心一點,不要亂發脾氣,不要亂發怒..
但這個還給我們一個很好的教導,就是不要把怒氣放在你的生命裡面太久,或許有時候真的是你永遠放在你的生命當中.這告訴我們,我們的怒氣,我們的生氣應該好好去面對,希望能夠處理得比較好,不要讓這個怒氣成為我們生命的控制..
所以,在該隱的生命裡面,他的怒氣就是成為對他生命的一種控制..
除了面子,怒火以外,該隱的生命要覺醒的其中一個地方是,他是一個推涉責任的人..
如果你在香港生活一段時間,你就知道一個我們常用的名詞叫"蟋蟀",這個蟋蟀就是說,有些人他們是很喜歡把所有的責任或者說所有的不好的,他背的黑鍋都擺在別人的身上,跟他完全沒有關係..
這個就是蟋蟀,他很會不用背黑鍋的那種人.他是推涉他應該負的責任..
所以該隱也是這樣子,當神問他:你的兄弟亞伯在哪?該隱怎麼樣回答?他說:我豈是看守我兄弟的人呢?.
其實就是該隱把那個責任,把那個黑鍋推掉.其實在當時候沒有多少人,當時候如果是亞當,夏娃,聖經所記載的,我不曉得他們應該結婚還是沒結婚呢,至少那時候只有亞當,夏娃跟該隱,亞伯,就這四個人在這個地上了,我們假設是這樣子..
那如果不是該隱的問題,是誰的問題呀?是不是亞當跟夏娃把那個亞伯不曉得怎麼樣弄?.
所以這個就是該隱他那種生命沒有覺醒的地方,他是一個推涉責任的人.他的回應是說,我豈是看守我兄弟的那個人?.
他連他兄弟的名字也不願意提,他不願意提到我豈是看守亞伯的人呢?.
當然,我想告訴各位,在該隱的生命裡面,他不單只是應該看守他的弟兄,他的兄弟,他更應該是一個怎麼樣的人?他應該是像約翰一書所說的,他應該是愛他兄弟的那個人..
其實神對他的要求不單只是要看守他的兄弟,其實神要求我們是彼此相愛,不單是教會裡面的彼此相愛,更是一個家庭裡面..
可能在你家裡面,不一定所有人都是基督徒,但我相信神給我們的一個很清楚的,很重要的就是我們要彼此相愛,因為我們的神是愛.我們不單只是愛我們教會裡面的弟兄姊妹,更要愛我們身邊的家人,甚至說你身邊其他的朋友,我們都需要愛他們..
所以這個不單只是說,我只是看守我弟兄弟的.其實神對我們的要求是比看守更高,是愛我們的兄弟,愛你身邊你的家人,愛你教會裡面的弟兄姊妹.這才是從聖經裡面更重要的一個吩咐,比看守更大.這是愛..
所以約翰遺書,我們來看另外兩節的經文.這裡就是告訴我們說,除了神是愛的源頭,我們更需要在神的愛裡面去愛我們的弟兄姊妹..
我們一起來讀這段約翰遺書的第三章的十跟十四節."一旦彼此相愛,就是你們從起初所聽見的命令,愛弟兄,就曉得是已經出死入生了,沒有愛心的人住在此中..
"如果該隱他有對雅伯的愛,我相信就像一個出死入生的人,他可以離開在被怒火控制的人生當中,他離開這個怒火可以帶來死亡那種災禍的時候,他能夠出死入生..
"如果他有愛的話,他就能夠真的像經文所講的,出死入生,不會犯下殺人的罪.".

$^{81}$所以今天,我們在該隱這個生命的覺醒,反省當中,我們更多的去問我們自己,我們是不是一個不願承擔責任的人?我們會不會很容易推涉我們應該負的責任?我們有沒有在神的愛裡面去彼此相愛?.
我們不單止是要看守,我們更需要去愛我們的弟兄,愛我們身邊弟兄姊妹,我們的家人..
除了面子,怒火,推涉責任,邂逅,還有另外一個在該隱生命當中我們要被提醒的是什麼呢?.
該隱是有一種損人不利己的心態.我不曉得在你的人生經歷當中,或許在工作裡面,在職場當中,很有機會會遇到這樣子的人..
當然我們盼望的是說,第一,我不應該是這樣子損人不利己的人.但我們也會有時候真的遇到在身邊,特別是跟工作有關,很多時候有些人真的只是會損人,他不單沒有對自己有益處,反而是有點兩敗俱傷..
我跟各位來看這個尤大叔,他這邊也引用了該隱.他是一種怎麼樣的心態,在這個經文有更明顯的..
這是從現代中文譯本的聖經裡面我們拿出來的,跟我們平常用的合文不一樣,但我覺得它的意思更能夠明顯..
我們一起來讀這段經文,尤大叔的第十到十一節."可是,這些人僅毀謗他們所不了解的事,像野獸一樣,隨著本能去做那些毀滅自己的事,這些人要遭殃了..
他們跟該隱走同一條路.這種人,他們就像該隱走同一樣的路,就是他們隨著自己的本能,就是有不同的心態,他們做那些損人不利己..
當看到別人有的,他們也不自己沒有,他就不想別人有."這個很明顯就是該隱.他覺得雅伯他的獻祭是神所約納的,他覺得他自己沒有被約納,他寧願把雅伯殺掉,也不想他能夠得到神的約納..
這種損人不利己,一派兩散.當然更嚴重的就是剛才我說到的,有可能是兩敗俱傷.不單只是損人不利己,他自己也受虧損,害人的最終也是把自己損害了,所以是兩敗俱傷..
所以,在我們身邊有可能有時候會碰到這樣子的人,我們求主來保守我們.當然我們有時候真的很難知道誰是這樣子的心態,但求主幫助我們也保守我們,也從我們自己心做起..
我們不要做那種像該隱的損人不利己,別人你自己沒有,也不想別人有,以致你跑去傷害了其他的人.求主來憐憫我們..
從該隱這個反面教材裡面,我們看到他為了自己的面子,他在生命裡面那種全流不斷的怒火和怒氣,他那種推涉責任和損人不利己的心態,都是我們要好好的被提醒的,我們不要重蹈該隱所走的路..
就像剛才尤大叔說,他們跟該隱走同一條路.我們作為基督徒,我們不要走該隱的路,我們應該以該隱作為我們的見解,我們要被提醒,我們要多在生命去覺醒,我們是否走該隱的路..
在你的生活裡面,我不曉得你會不會真的有人很坦白,坦誠地跟你說,你應該多一點清醒,多一點警覺,提醒你會不會為了面子,會不會有怒火在你的生命當中,會不會推涉你應該負的責任,會不會你做了一些損人不利己的事情..
有時候,我們不會自動自覺去省察我們自己的.這種省察生命的心和力,是需要我們求神來幫助我們,多有積分的心和力..
在《耶利米蘇》告訴我們說,我們人心比萬物都詭異,壞到極處,誰能悉透呢?所以人很多時候沒有這種自覺能力的話,就很危險..
在《約翰一書》,剛才我們都提到不少的經文.《約翰一書》有一段的經文,我覺得很重要.它告訴我們說,在我們的生命裡面,因為我們是屬於神的,不是屬於那惡者..
如果我們是靠近那個惡者,靠近那個魔鬼撒旦的話,我們就沒有多少的力量.但如果我們記得我們是屬於神的,而且是因為神已經勝過了世界.在我們每一個基督徒的裡面,我們有神的生命,有聖靈在我們生命裡面,那是在比世界的怎麼樣?更大..
神在我們裡面比在世界的惡者更大.這就是我們的力量,我們的能力的來源.不是靠我自己,是靠著神,我們多有一份自我省覺,我自己常常去省察我的行為是怎麼樣..
所以這就是我們應該要做的,求神幫助我們.在該演其中一個很明顯的,我想跟各位提醒,也是對我自己常常提醒的,該演是在一個家裡面的長子,長子代表一個身份很尊貴,地位很高的一個人..
常常我們在我們的生活裡面,如果我們去到某一個位置,我們覺得我自己已經很重要了,我自己已經是單位的領導了,我這個身份,有時候我沉浸在這個身份裡面的,我就很難有那種覺醒的能力,覺醒的心..
有時候我們怎麼樣可以保持清醒,保持警覺呢?真的是別人是可以提醒的,但很多時候別人的提醒可能成為你不喜歡聽的話..
所以,我們對自己的那種省察,對自己生命的覺醒是很重要的.求主幫助我們,多有這份自我省察,覺醒的心態.尤其是當我們不當什麼的領導,你真的是長子,你真是很有尊貴身份,地位的,.
你更要小心,更要醒覺你自己會不會像該隱走在該隱的道路上..
這個是我想在第一講的時候跟各位來領受神給我們的信息.我在兩個禮拜以後會繼續續講,下面其他跟該隱有點相似的作為,我們繼續在當中領受的信息..
我們要常常更多的自我醒覺,察覺我們有沒有在像該隱這樣子的不好的心態.無論是為了面子,無論是我們心裡面還有不同的怒火,或許我們是推涉責任,我們甚至有一些是選人不利己的心態的時候,求主寬恕我們,求主憐憫我們,求聖靈在我們的心來幫助我們有這樣子省察覺醒的能力..
也讓我們繼續的依靠你,信靠你.因為我們能夠有這樣子的覺醒的能力,能夠繼續走在一條正路上,不是靠我們自己的所做多少,真的是要完全依靠在我們生命裡面,那一位比外面世界更大的..
因為有你在我們裡面,有聖靈在我們生命當中,我們就能夠去走在正路上.也讓我們真的被提醒神是愛,我們應當彼此相愛,我們應該愛我們身邊的人,願主將一份愛繼續放在我們的生命裡面,謝謝主確定我們的禱告,靠耶穌基督領求,阿們..
\newpage



\section{腓立比書 4:7}
\label{sec:032184}
\textbf{平安,平安:生命的守護}~~~~陳淳牧師
\newline
\newline
連結: \href{http://www.church.com.hk/acms/content.asp?site=cdc\&op=show\&type=product\&code=032184\&layout=sermon}{\texttt{032184}} ~~~~ 語音日期: 2026-05-24
\newline
\newline
\hyperref[sec:032183]{\small{< < < PREV SERMON < < <}}
~
\hyperref[sec:index_chronic]{\small{[返順時目]}}
~
\hyperref[sec:032185]{\small{> > > NEXT SERMON > > >}}
\newline
\newline
腓立比書 4:7
\newline
\begin{longtable}{cl}
\hline
\hline
章節 & 經文 (和合本修訂版)\\
\hline
4:7 & \begin{tabularx}{0.7\textwidth}{X} 神所賜那超越人所能了解的平安，必在基督耶穌裡，保守你們的心懷意念。 \end{tabularx} \\ \\
[1ex]
\hline
\hline
\end{longtable}
$^{1}$何俊仁議員.
主持人:當弟兄姊妹在主日崇拜之中,被上帝的懷與激勵,滿滿帶著盼望和喜樂踏出禮拜堂,回到自己的生活當中,.
回到那個可能充滿著很多投訴,壓力,衝突,甚至有很多滿天的壞消息的世界,.
你這份喜樂究竟可以維持多久呢?我們可能很多時候都是星期一滿懷盼望,.
星期一回到現實裡面,開始基於現實,而總是滿心憂慮,然後持守著你的衝勁,你的毅力,.
去到星期二,你可能已經在重壓之下,滿身疲乏,疲憊,然後你的喜樂,你的盼望,這一切慢慢消失無蹤,.
就好像你很辛苦地賺了一大筆的財富,但是因為沒有保險箱,因為沒有保安系統,.
結果一出門就被生活的小偷偷得一乾二淨,洗劫一空..
保羅在菲勒比書四章十七節那裡說:.
「神所賜出人意外的平安,必在基督耶穌裡保守你們的心懷意念.」.
其實在這裡神有一個應許,透過保羅宣講,他說:「平安,平安,上帝要設立一個最強大又最安全的系統保守我們.」.
「平安,平安,上帝要設立一個最強大又最安全的系統保守我們.」.
我們這火非常容易擔驚受怕,甚或冰吹草動,風吹草動之下,我們就會動搖的心..
在訊息之前,我們一起同心吐槽..
神所賜出人意外的平安,必在基督耶穌裡保守你們的心懷意念..
平安,平安,上帝要設立一個最強大又最安全的系統保守我們..
神所賜出人意外的平安,必在基督耶穌裡保守你們的心懷意念..
平安,平安,生命的守護是從追逐外在的安穩,去到屬天的連結的跨越..
很多時候我們提到平安,一般人的直覺就是沒有壞事發生,就只有平安..
如果今天我開車,沒有出任何意外,順暢,平安..
如果今天做了身體檢查,報告沒有任何問題,平安..
今個月我的收支平衡,不會入不敷支,平安..
在我們的邏輯裡面,我們所有的平安,都基本上是因為外在的環境平靜,或者我們所有事情順暢,所以帶來的結果..
不過你有沒有留意,這種平安有一個致命點,致命傷,就是它是被動..
意思就是,它需要依賴著一些條件性的滿足,無論是錢,無論是健康,無論是人的認同,或者辦事的順暢度..
如果我們這樣的平安,能夠只是建立在這一連串的假設,或者這一連串的外在條件,.
難怪我們的心,好像一座蓋在沙土上的房子,水來土掩,風吹草動,波浪大地蓋過,我們的平安,就好像玻璃碎片那樣碎了..
保羅在這裡經文一開始,他用了一個非常關鍵的片語,這個片語就是:神所賜的平安..
原文是這樣寫的,神所賜的平安..
當然這個原文用的是一個所格,意思就是,可能你對這個名詞不太理解,.
但它的用意是什麼意思呢?就是它想解釋平安的出處..
平安的出處不是來自環境,平安的出處不是來自資源,平安的出處不是來自你的IQ,你的人力,物力..
平安的出處是來自神自己,是來自屬神的平安..
所以,平安這個字,Aveni,希臘文,本身不是講希羅文化的平靜和安靜..
保羅用的時候,他用的是引用白萊文的salon,平安的用詞..
這個用詞包含了完整,整全,完備,更加有一種重新連合的狀態..
就好像一個打碎了的瓷碗,它支離破碎,但是神用它一種神聖的膠水將它重新接合起來..
這個就是福音,這個就是我們在過去的罪惡當中,因為自我中心,因為離開上帝,.
所以我們的生命都支離破碎,陷在焦慮和與上帝隔絕當中..
但是我們透過耶穌基督重新與神連結,透過耶穌基督的救贖,透過耶穌基督的福音,.

$^{41}$我們可以重新接通這個狀態,這個就是希伯來文裡面說的平安..
這個跨越非常重要,因為平安不是指你外在的環境沒問題,.
不是你外在資源沒問題,不是你裡面的底子夠厚,.
或者甚至有人力物力支持你,不是這些..
平安是指你與神的關係沒問題,你與神的關係連接上..
屬神的平安無需依靠世界性,世俗性的條件,.
意味著你不再需要看環境做人,甚至你不需要看自己主觀或者被人影響的意欲或者能力..
三個星期前我們說,我們喜樂是在主裡面找到那個位分..
今天我們很清楚知道,因為這種與上帝的接通有這個位分,.
所以我們擁有無限的供應,包括喜樂,包括平安..
這份平安是神所賜給人出人意外的平安..
世界給你的平安是有價值的,有標價的,有標籤的..
你要透過努力,你要透過做事成功,你要透過保持健康,.
你要透過你能夠做一些事情對人有益,有幫助,你要討人喜歡,你要討人歡喜,.
然後你才能夠透過你一切的付出,換取到一點點的安慰,一點點的安恕..
但是苦心自問,我們付出一切,總是好像換不到那份安樂..
那份安樂總是好像表面上隨手可得,好像我們付出一切就能換到,.
但是總是吊著我們,總是令到你和我身心俱疲..
但是神的平安是恩典,神的平安是神主動賜下,.
這種平安就好像高山的清泉從山上流下,.
無論山下城市的人有多少的喧嚷,多少的紛爭,多少的擾亂,.
甚至城裡面的人有多少的乾燥,多少的乾旱,多少的爭鬥,.
但是這種水的源頭清泉,仍然是永不枯竭,仍然是流轉在我們每個人心中,.
將我們心中那份的煩躁能夠洗滌一空,這個就是第一重的跨越..
我們的平安不再是環境的產品,我們的平安是因為身為神而屢,而獨有的產業..
所以,我邀請大家,當你下次面對混亂的人和事,.
甚或力有不退而帶來的心中沮喪,你問一問自己,.
你的心現在是和現實環境的信號連上了,.
抑或是你和神連上了那個源頭..
神的信號透過聖經的話語成為我們的安慰,成為我們的幫助,成為我們的堅定..
只要源頭接對了,我們的平安就守住了..
神所賜出人意外的平安,必在基督耶穌保守你們的心懷意念,平安平安..
生命的守護,第二個層次的跨越,是從依靠理性的分析,去到屬靈超越邏輯,悟性的跨越..
神所賜的這份平安有一個標籤,這個標籤保羅稱之為出人意外的平安..
不是你猜不到,我們看中文我們很容易就猜不到,對不起..
我在解釋之前先問一下,你上一次你的不安,或者你的失眠,.
或者你心口裡面覺得硬一硬,好像悶悶不樂,甚至心有鬱結,.
不知道怎麼辦的時候,你會做些什麼?你會做些什麼?你的不安,你會做些什麼?.
我們很多時候,我們大腦馬上開啟一個全足運轉的模式,.
我們馬上分析風險,我們馬上做最壞打算,我們馬上試圖想通整件事來龍去脈,.

$^{81}$我們馬上找一下究竟我有什麼支援,究竟我有什麼資源可以應對,.
然後如果我想通整件事來龍去脈,原來是這樣我得罪他,原來是因為這樣,.
我們心裡面就好像有一份平安,我們這份平安好像有應對的方法就有平安,.
好像找到方法,知道原因就有平安,但是我們這樣的平安其實是將它建基於我們的理解能力,.
我們資源多寡,甚或是我們的能力多少,甚或我有多少人力物力,.
識人比識字好,所以我識多些人,我就有多些把握,是不是這樣呢?.
事實證明人生遇到很多事情,不如意事十常八九,總是不如其來的噩耗,.
總是無法解釋的苦難,總是不明不白的病痛,總是不知道為什麼關係失去,.
這一切都遠遠超過我們大腦能夠分析的負荷能力,.
所以當你和我想不通,想不到,借不到,找不到的時候,你和我心中平安就崩潰..
保羅在這裡用了一個出人意外的字,這個字非常特別,不是你猜不到的字,.
這個字在希臘文是有三個字組成,是一個片語,是一個形容的片語,保羅用了一個強調的排列方式,.
他的寫作方法非常特別,他是先寫平安,所以他說神所賜的平安,第一句,.
或者不是第一句,第一句的一段落,然後後面出人意外的,出人意外的其實剛才我們說了三個字,.
所以直接翻譯是什麼呢?就是那超越所有心智的,還沒完的,語法告訴你,.
他一定指向前面那個平安,神所賜的平安,就是那超越所有心智的平安,.
所以在原文語法裡面是有兩個平安的,保羅故意將這個形容放在後面,.
這樣寫法,語法上面,文法上面,就是代表著將聚光燈聚焦在中間平安這個字,.
他的絕對能力上面,他不只是一個抽象的名詞,他是一個凌駕一切,超越人類理性邏輯的屬靈能力,.
所以這裡的超越這個字,希臘文是在說地理用語,他其實是在說一種高過甚或是凌駕其上的意思,.
想像一下,如果你開著車,在迂迴的盤旋的山路上面,山路邊沒有欄杆,.
你不斷轉一個彎又一個彎,但在這個時候你的路充滿大霧,你的視野最多不夠兩米三米,.
你的車高速地滑行,你會不會害怕?這份迷茫就好像我們人生一樣,我們身處濃霧當中,.
但是如果我們能夠站高過迷茫,我們能夠高過霧,我們在高峰上面看下去,.
你就會發現雖然雲霧依舊,但是我們的心就安定下來,這個就是保羅的用字,.
所有當然就是包括所有,一切,一切是什麼?心智,這個心智說的是悟性,理智,思維,邏輯,分析能力,.
就是我們剛才說賴之為平安的工具,我們就是透過這些擁有,甚至因此得著世界平安的方法,.
所以保羅很清楚告訴我們,平安,神的平安是一種超然的覆蓋,是一種壓倒性的碾壓,.
一切所有你擁有的能力,它是凌駕於你大腦的邏輯之上,所以這個第二重的跨越,.
世界的平安是基於你對人情世故的掌握,基於你對前途的把控,基於你對資源運用的能力,.
但是屬天的平安是來自於悟性的凌駕和跨越,這樣意味著無論你和我人生順或逆,.
甚或我們處身人生的低谷當中,甚或我們看不透想不明白為什麼有這個病會發生,.
甚至我們完全想不到為什麼那些關係會去到這麼糟糕的情況,甚或我們不明白為什麼自己的心情會去到這麼差,.
甚至不知道怎麼從泥沼裡面走出來,在這個時候你對前路毫無頭緒,你對前面毫無把握,毫無保障,.
從世界的經驗來說,你就是一個失敗者,從世界的經驗來說,你就是沒辦法幫助自己的人,.
從世界的經驗來說,你無藥可救,從世界的經驗來說,你根本會落在一個沮喪和失望的深淵當中,.
不過,今天的經文很清楚告訴我們,即或不然,即或我們是這樣,你仍然可以擁有那份平安,.
因為這份平安是神所賜下出人意外的平安,這份平安不是靠你腦袋分析出來,不是靠你擁有的資源堆積出來,.
不是透過其他人擁抱你,而是令你飄搖而出來,保羅在這裡告訴我們,.
我們不需要想通一切,不需要擁有很高的IQ或者很高的EQ而才擁有平安,只需要你和我在神面前領受這份高過理性思考帶來的平安,.

$^{121}$因為在你和我未出生之前,在你和我未懂得怎樣思考之前,在你和我未掌握到足夠資訊之前,.
在你和我未堆積到有足夠資源幫助我們渡過難關的能力之前,神已經賜下平安給每一個屬祂的人,.
如果我們三個星期前說喜樂是一個持續住在基督裡面,憑意志作出選擇,跟隨主的行動,.
這節經文是很清楚地告訴我們,在我們未擁有選擇意志,未懂得選擇,未懂得相信之前,.
神已經從天上降下這份高過你和我頭腦的平安,作為你和我的後盾,作為你和我的盾牌,幫助我們面對人生的波浪,人生的危險,.
這不是代表你要放棄你的理智,不需要用你的智力,祂只需要告訴你,神的主權比你和我的智力更寬廣,.
我們用盡我們的大腦,用盡我們的所有資源,用盡我們的所有方法,當你去到極限,沒辦法解決困難的時候,.
神所賜的平安,就在這個時候,在我們內心接管,弟兄姊妹,這個就是我們對上帝的信靠,.
這個就是我們行事為人,對得起主,作為人生管家,應有的本分,你是不是這樣信靠上帝呢?.
平安在最混亂的時候,反而顯得最真實,它不是在我們冷靜的時候分析出來的結果,.
它反而反過來,它在我們混亂之間,紛亂的過程裡面,驚濤駭浪的過程當中,你發現你依然能夠靠主站立得住,.
這個是我們每一個弟兄姊妹,每一個基督徒,我們經歷的日常神蹟,.
這個神蹟讓我們再一次經歷神系凌駕萬有的平安,能夠幫助我們身在困難裡面,.
想無可想的時候,可以安然放下重擔,對著主說,主啊,我真的想不通,我真的想不明白,.
我真的用盡所有我的方法,你給我的所有資源,我盡了責任,但是我仍然是舉手投降,.
在這個時候,我懂得也能夠因為耶穌基督而信靠你,因為你是超過一切,你是超越一切的智慧,.
你的愛完全涵蓋過我的不安,聖經說,神所賜出人意外的平安,.
必在基督耶穌裡保守你們的心懷意念,平安,平安,我們是要擺脫被動情感的干擾,.
去明白神已在主動守護的跨越,弟兄姊妹,如果我們找到平安的源頭,.
剛才我們說了平安來自上帝,我們又看到平安的邏輯是我們盡了所有一切努力,.
我們懂得附上我們應盡的本分之後,我們懂得放下跟隨主,現在我們即將要進入一個最關鍵的實踐,.
這份平安在我們生活裡究竟如何運作呢?我們會否濫用它呢?或者我們會否不用它呢?.
我們很多時候對平安有一個很大的誤解,我們覺得平安是軟綿綿,安安靜靜,平平淡淡,.
就好像現在這一刻,我們總是好像寂靜無聲,我就平安了,我們以為平安很像我吃完飯,.
閒時無事,無憂慮的情況之下,我走去公園散步,心無波瀾,.
甚或你覺得好像在湖水裡面,看著毫無漣漪,完全沒有波浪,覺得波平如鏡,平平安安,.
你就覺得這樣叫平安,如果你是這樣想,難怪生活的攻擊來到,難怪那些不公平的批評來到,.
難怪你對未來不知的時候的恐懼,難怪你對前景毫無把握的時候,.
甚或有一些引誘,情慾的衝擊來到你面前,你的利益的引導,.
以至你明知要違背神的原則的時候,拍打你內心,在你門口不斷敲門,不斷搖動你,.
那個時候我們的血壓會上升,我們的呼吸會急速,我們的心會加快,然後你會產生喘氣,.
在這個時候你會怎樣做?你會大膽一點喘氣,幫助自己平復,一般情況都有效的,.
但是你知不知道心理焦慮是有另一種叫搶呼吸的,就是你越吸越覺得好像呼吸不到,.
然後越呼吸不到就越吸多一點,吸多一點就越害怕,越害怕就越吸不到,然後好像窒息一樣,.
你有沒有試過?你有沒有試過你見過別人這樣做?你知道怎樣做嗎?.
我們給他一個紙袋,吸他自己的二氧化碳,發現原來大腦是正在呼吸的,那他就平靜了,.
我這樣說不是叫你下次試和不試,我只是告訴你我們習慣是這樣,.
其實我們在很慌亂的時候,我們很容易就覺得我們平安無事,.
因為我們剛剛說我們平安是感覺,我們平安是平靜,所以一呼吸急促一點,.
一思想混亂一點,我們就覺得我們平安消失了,因為我們心裡面有一個潛意識,.

$^{161}$我們以為平安是我們要努力維持的情緒和心情,.
所以你努力維持它,它變化,你就會覺得失控,.
但是從這段經文很清楚地保羅告訴我們,平安是一個很震撼的應許,.
平安是用來保守我們的心懷意念,保守這個字,希臘文不是一個溫柔的字,.
不是一個平靜的字,不是一個平淡的字,它是一個充滿力量的軍事術語,.
它的原意就是說駐紮守衛,是軍人擔當哨兵守護的用意,.
在當時菲律賓城是一個重要戰略經濟價值的城市,.
所以羅馬駐城他們要猜探最精銳最強悍的軍隊,在門口站著駐防站崗,.
這些精銳部隊唯一的目標就是守護,他們守住門口,擋住一切外來的侵略者,.
確保敵軍無法滲進來破壞,所以它不是觀光遊客的影像對象,.
它不是裝飾,它不是拿來立威,它拿來最重要的是守護,.
好讓城中的百姓可以安然居住,在裡面能夠享受真正的平安和自由,.
這個就是第三個跨越的重要,平安不是你需要努力守護的心情,.
平安不是你付出一切去交換的物件,平安不是你擁有多少資源可以購買的代價,.
平安是神親自差派一支精銳部隊守護你的人生,.
平安不是你維持的情緒維持的代價,平安是神哨兵讓守衛守護你的人生,.
這支屬天的精銳部隊守護的範圍非常具體,.
它是守護心懷意念,心懷意念中文四個字,原文是兩個字,.
涵蓋兩個很重要的部份,心懷指的是你的情感的中心,.
是說你會痛會怕會焦慮會徬徨那個地方,而意念是說我們的想法思想邏輯,.
甚至我的判斷甚至我的決定,這份平安就是一個專業的軍人哨兵,.
當你恐懼衝入你內心的時候,當焦慮試圖鑽進你大腦的時候,.
這個哨兵站在門口攔截他們,他告訴這些壞消息,告訴這些攻擊你的聲音,.
對不起這顆心已經被神的主權掌管了,接管了,神不准你進入去破壞,.
如果我們說喜樂是活在主裡面得著那份財寶財產,.
這樣今天神就是親自差派來的保安系統衛士哨兵幫助我們守護這份財產,.
不被生活的小偷偷走,不被你和我的私慾所遺失,.
說到這裡我需要可能和大家有一個重要的提醒,.
如果我們說平安是衛士是哨兵是軍人,哨兵站崗是有準則的,.
這個準則決定了幫助我們,究竟哨兵站不站在這裡,.
這個決定在於我們行事為人是不是符合聖經真理的標準,.
我們行事為人是不是對得起主,我們或者在生活裡面,.
在弟兄姊妹生活當中,不知道你有沒有聽過,或者不知道你有沒有說過,.
你說我祈禱之後心裡很有平安,所以我決定做這個決定,.
或者你說我做了這個決定,人家覺得很詫異,問你你真的這樣做決定?.
是的,我祈禱之後心裡滿有平安,有沒有聽過這些對話,.
或者你有沒有說過,當你聽到這樣的說話的時候,.
你可以說什麼,你除了說Amen,沒有什麼好說的了,.
他出了王牌了,高戰了面戰牌,你沒有挑戰過,.
我與上帝已經同政,我的決定就是上帝的決定,.

$^{201}$弟兄姊妹,我需要今天嘗試解釋,或者嘗試演繹一下,.
其實有三個錯誤的平安錯覺,是與那份來自上帝屬天的平安有不同的,.
我試著說一下,看看你會不會回心微笑,.
第一個,叫做慾望滿足的舒暢感,慾望滿足的舒暢感,.
是因為你有些東西想得到,所以就快得到了,都明明將會得到了,.
只要我伸手就拿到了,所以你心滿意足,你很自然就漠視一切,.
總之心中舒暢這份叫做平安,第二,那些叫做逃避現實的解脫感,.
因為你可以選擇一條比較容易的路,逃脫了,不用面對一些麻煩,.
不用面對一些衝擊,不用守護著那些真理,.
所以你終於鬆一口氣,不關我事了,和我無關了,.
所以我心裡面那份平靜,你誤以為平安,.
如果前面那兩者你都容易分辨的,第三種我相信很多基督徒都有點挑戰的,.
包括我和你,我都是會遇到麻煩的,那個叫做追求不凡的成就感,.
我們都會選一些高難度的挑戰,向所有人,不是所有人,先向自己,.
向自己證明我是不凡的,我不是平庸的,我是被上帝所用的,.
上帝你塑造我,你當然要用我,所以我做一些高難度動作,做一些有挑戰性的動作,.
平庸的事不輪到我做的,讓平庸的人做吧,然後我熱血沸騰,.
心裡面充滿著那份傲氣,充滿著上帝等我來吧,.
這一份滿足令你以為是平安,剛才這三種都是錯覺來的,.
都不是真正的平安來的,但很多時候我們卻拿來做擋箭牌,迴避自己的責任,.
真正的平安是攻克己身,叫身伏我,總叫基督在我身上照常顯大,.
唯有這樣持守真理,行事為人,對得起主的穩妥感,.
這個才是真正的平安,個人慾望會因為滿足而感到舒暢,.
但它會冷卻的,逃避責任,你會鬆一口氣,但責任總會追上來的,.
麻煩總會繼續來的,迴避真理,逃避問題,以為自己力能勝天,.
驕傲無視一切,你總會撞到頭崩額裂的,.
但是神所賜的平安,守護的卻是你和主之間的連結,.
和你對主的忠貞,甚至對你身為管家應有的本份,.
不是過分的浪費,也不是過分的收買,是恰如其分的應用上帝所賜給你,.
所以我們或者會因為做了一些合乎真理的決定,.
依然環境沒有改變過,仍然是波濤洶湧,前面滿有艱難,.
你為此甚至要付出一些代價,大到旁人側目,.
但你心裡面有一份安妥,有一份平安,因為你守住真理,.
你很清楚這是上帝叫你做的,你很清楚這份坦然,.
是因為你持守真理,對得起主,你站立就在神的心意當中,.
你清楚知道你不是驕傲,你不是狂傲,你不是無知,.
你是因為相信上帝,所以你夠膽繼續做下去,.
這份對得起主的穩妥才是真正的平安,.
因為平安這個哨兵,這個守衛,守護的是真理的道路,.
是當我們意念,當我們想法,情慾偏離幾路的時候,.

$^{241}$他就會提醒我們,他就不會站在我們的情緒當中,.
所以我需要鼓勵大家,有不安無須恐懼,.
因為不安正正是我們思想,我們行事為人,.
是不是對得起主,不安這個信差是告訴你,.
你可能偏行了幾路,可能你力不能及,.
可能你需要求援,向人,向神求援,悔改,.
以致你和我不要偏行幾路,不要走捷徑,.
這個哨兵最重要的重點是讓我們清楚知道,.
持守真理我就得著平安,違背真理我就得著不安,.
這個真正的平衡不是在於你的感覺,.
真正的平衡在於你行事為人,是不是對得起主,.
你有沒有守住真理的底線,有沒有平衡自己過於樂觀,.
而不顧一切,只有我們恰如其分選擇行在主的真理當中,.
讓我們的心懷意念保持在對得起主的頻譜上,.
這份屬天的守護就永不設缸,.
它會幫我們擋住我們因此而帶來的恐懼,私慾,甚至無知,.
甚至權衡利弊的怯懦,這些雜音會從我們腦中消失而去,.
因為我們都站在上帝要我們走的道路當中,保守我們,.
弟兄姊妹,如果我們說平安,平安,.
生命的守護是神為我們生命設立最大最安全的系統,.
這樣容許我今天用四句說話來總結今天的講道,.
喜樂盼望平安好,屬天連結安領保,.
超越邏輯平添普,真理守護安穩度,.
可能平安好不用解釋,然後安領保也應該不用解釋,.
你安心領就可以了,你只要信奉耶穌就可以領到,.
平添普又是什麼呢?你真是不靠譜了,.
平添普我想表達的是,當我們大腦想不通,.
我們自己胡亂荒腔走調的時候,.
當我們自己以為自己可以擁有一切去走我人生路的時候,.
神的平安,超然的覆蓋,增上那份屬天的信號的頻譜,.
讓我們走在恩典的樂章當中,按著恩典的曲調而走下去,.
按著上帝的譜曲,你和我或高或低,仍然心中沒有平安,.
透過屬靈的悟性跨過我們所有理性分析利弊的勇氣或者怯懦,.
讓我們看見上帝就是我們人生的出路,.
什麼是安穩道呢?我們生活當中有很多情慾,.
很多自己的情緒,甚至有很多的怯懦,甚至有很多我們的計算,.
生活的小偷會偷去我們一切的安樂,平安,.
真理守護安穩道,幫我們堵塞所有的破口,.
當我們回到家庭裡面面對所有的瑣事,.
當我們回到職場面對那些讓你頭痛,讓你麻煩,讓你焦慮,.

$^{281}$甚至不知道怎麼處理那些人和事,.
甚或你面對你的引誘誘惑,想你走人生的捷徑,.
這個時候,讓你記得今天這段經文,.
神不是讓你一個人面對,神是讓你和我已經得到一個很巨大的守護,.
他派駐了這個守護者,這個守兵,就是出人意外的平安,.
只要我們立志行事為人,按著真理對得起主,.
我們這份平安就能夠成為我們生命最大的護盾,.
這份平安就能夠成為我們生命當中最堅固的堡壘,.
讓我們同心吐槽..
親愛的阿爸父,我們來到你面前需要承認,.
我們的平安總是隨著環境起起伏伏,.
我們太容易就讓生活的不如意牽著我們鼻子走,.
我們太容易就為了旁邊的人而迷失了我們應該奔跑的方向,.
我們對未來總是焦慮不安,.
主啊,求你讓我們有勇氣,讓我們學會積惑不言,.
即使想不通,即使不知道怎麼做,.
但是我們仍能夠憑順領受你這份高過一切的平安..
多謝你,差派平安好像哨兵一樣駐守在我們內心,.
幫助我們心懷意念都較正對著聖經的真理,.
求你幫助我們,透過永不斷線的Wi-Fi,.
與你連繫,在你的保守底下,.
我們能夠安然居住,不再擔驚受怕..
我們能夠如此祈禱,乃奉我主耶穌基督的聖名,.
阿部.
\newpage



\section{使徒行傳 2:1-4}
\label{sec:032185}
\textbf{歡迎聖靈的風吹來 (華)}~~~~劉彤牧師
\newline
\newline
連結: \href{http://www.church.com.hk/acms/content.asp?site=cdc\&op=show\&type=product\&code=032185\&layout=sermon}{\texttt{032185}} ~~~~ 語音日期: 2026-05-24
\newline
\newline
\hyperref[sec:032184]{\small{< < < PREV SERMON < < <}}
~
\hyperref[sec:index_chronic]{\small{[返順時目]}}
~
\hyperref[sec:032186]{\small{> > > NEXT SERMON > > >}}
\newline
\newline
使徒行傳 2:1-4
\newline
\begin{longtable}{cl}
\hline
\hline
章節 & 經文 (和合本修訂版)\\
\hline
2:1 & \begin{tabularx}{0.7\textwidth}{X} 五旬節那日到了，他們全都聚集在一起。 \end{tabularx} \\ \\ \relax
2:2 & \begin{tabularx}{0.7\textwidth}{X} 忽然，有響聲從天上下來，好像一陣大風吹過，充滿了他們所坐的整座屋子； \end{tabularx} \\ \\ \relax
2:3 & \begin{tabularx}{0.7\textwidth}{X} 又有舌頭如火焰向他們顯現，分開落在他們每個人身上。 \end{tabularx} \\ \\ \relax
2:4 & \begin{tabularx}{0.7\textwidth}{X} 他們都被聖靈充滿，就按著聖靈所賜的口才說起別國的話來。 \end{tabularx} \\ \\
[1ex]
\hline
\hline
\end{longtable}
$^{1}$讀一段聖經,在史徒雄文第二章第一到第四節..
這時候五旬節到了,門徒都聚集在一處,.
忽然從天上有響聲下來,好像一陣大風吹過,.
充滿了他們所坐的屋子,.
又有舌頭如火焰顯現出來,分開落在他們個人頭上,.
他們就都被聖靈充滿,按著聖靈所賜的口才說起別國的話來..
今天是五旬節,也是教會傳統當中的聖靈降臨日..
在五旬節那一天,門徒們聚集在一起的時候,.
忽然間從天上有大的響聲,好像大風吹過一樣,.
又有舌頭如火焰顯現,分開落在他們個人頭上,.
他們就都被聖靈充滿..
聖靈是誰呢?聖靈就是我們所信三位一體真神中的一位..
我們所信的神是三位一體的神,.
聖父,聖子,聖靈,是一位神,但有三個不同的位格..
而聖靈卻是這三位一體真神當中與我們關係最密切的那一位..
因為今天真正住在我們裡面的是誰呢?是聖靈..
聖父,聖子是藉著聖靈住在我們裡面,.
而且聖經說聖靈住在我們裡面,永遠不離開我們,Amen?.
所以我常常說你千萬不要害怕聖靈,.
跟你朋友說不要害怕聖靈..
如果你害怕聖靈,聖靈永遠與你同在,你怎麼辦呢?.
No, no, no, no,不要害怕..
聖經講到聖靈的名字就是寶慧師,.
寶慧師翻譯出來就是安慰師跟輔導者..
哇,想想看,這難道不是我們生命當中最需要的兩件事情嗎?.
在我們每天生活當中,我們多麼需要得到安慰,.
多麼需要得到生活的指引..
聖靈就是安慰師,就是輔導者,指引者,.
而且祂住在我們的裡面,直到永永遠遠,Hallelujah!.
聖靈就是與我們關係最密切的一位神,.
而且聖靈在我們裡面是賜給我們能力的神,.
生活的能力,事奉的能力,.
祂在我們生命裡面能夠帶來翻轉生命的大能..
所以聖經常常把聖靈稱為是神的氣息,.
當神的氣息一吹來,就能夠恢復我們的心靈..
當神的聖靈一吹來,就能夠更新我們的生命,.
並且扭轉我們的環境..
那在以西節書第37章裡面,神讓以西節看見一個異象,.
那就是在一個平原裡面,充滿了骸骨..
而聖經說這種骸骨,這些骸骨是極其的枯乾,.

$^{41}$一個非常絕望的景象..
神就在以西節37章第三節,神就問以西節一個問題,.
說人子啊,這些骸骨能復活嗎?.
以西節怎麼看?這不可能的,不可能的事..
但他的回答卻很有意思..
神問他說,這些骸骨可以復活嗎?.
他說,主耶和華,你是知道的..
換句話說,我不知道,我怎麼看都不可能..
但是主啊,我相信你是知道的,.
因為在你,永遠沒有難成的事..
我們看見第九節,神就吩咐以西節說,.
人子啊,你要像風發預言,.
主耶和華入主的氣息要從四方吹來,.
吹在這些被殺的人身上,使他們活過來..
所以第十節,我們以西節就按照神所吩咐的,.
像風發預言,叫風從四面吹來,.
結果神,就讓神的氣息吹進這些骸骨當中,.
這些骸骨就活了,而且成為一支極大的軍隊..
Hallelujah!.
我就想到,今天我們不是在同樣的光景當中嗎?.
每當我們面對這樣一個充滿動盪的世代,.
多少時候我們面對環境,心中充滿憂慮跟恐懼?.
多少時候面對我們的前途,心中充滿徬徨?.
因為不知道何去何從?.
很多時候我們就像這個平原當中的骸骨一樣,.
我們多麼需要神的氣息吹來,.
讓我們的靈魂能夠甦醒..
我們需要聖靈的風吹來,.
讓我們生命能夠得以更新..
我們歡迎聖靈的風吹到我們中間,.
能夠重新帶來翻轉的大能..
那你說我們要怎麼歡迎聖靈的風呢?.
今天早上從三方面跟大家來分享..
首先第一個,我們一定要渴慕聖靈來充滿..
跟你們說要渴慕聖靈來充滿..
今天我們如果想經歷聖靈的風,.
我們一定要心中渴慕聖靈..
我很喜歡四篇四十二篇第一節的話,.
所以說神啊,我的心竊慕你,如同鹿竊目欺水..
這一首說的渴慕,不是一種普通的渴慕,.

$^{81}$是一種很迫切的渴慕,像鹿竊目欺水一樣..
你要知道當鹿竊目欺水的時候,.
就算那是在遙遠的鹿,.
就算在危險的環境,.
都不能夠攔阻他要跑到水源去得到水..
因為那不是可有可無的,那是個必須的..
只要他知道自己裡面有缺乏,.
只要他知道那個水源在哪裡,.
他一定會勇往直前..
同樣的,今天我們若是想渴慕神,.
我們也要真的看見這兩件事情..
一方面我們要看見我們自己的缺乏,.
如果一個人都不覺得自己有缺乏,.
他不可能來渴慕神..
其次,他需要看見神擁有答案,.
他需要看見神的豐富,.
他需要看見神足夠滿足他心裡一切的需要,.
否則他也不會來渴慕神..
在聖經當中我們看見耶穌的母親,.
在迦納的婚宴上,她看見一個缺乏..
什麼缺乏呢?他們沒有酒了嘛..
所以她就來到主耶穌的面前.為什麼?.
因為她相信耶穌有答案嘛..
那個管會堂的雅魯,在他家中,.
他看見一個缺乏.什麼缺乏?.
他的女兒生病,快要死了..
他來到耶穌基督面前,為什麼他來?.
因為他相信耶穌能夠醫治嘛..
那個患了十二年血肉病的婦人,.
在她自己身上,她看見一個缺乏..
什麼缺乏?就是她花盡所有一切,.
都沒辦法治好她的病..
所以她要來到耶穌面前,.
而且不顧一切要抓住耶穌,觸摸耶穌..
為什麼?因為她相信耶穌是她唯一的盼望嘛..
弟兄姊妹,今天你的心可目神嗎?.
你知道很多時候我們之所以不可目,.
正是因為第一個,我們看不見自己的缺乏嘛..
我們自以為很富足,自以為很聰明,.
自以為擁有很多,所以我們不大需要嘛..

$^{121}$在啟示錄當中,你看見主怎麼責備老底家的教會..
主耶穌斥責老底家教會說,你們自以為富足,.
一向都不缺以為..
所以主說,去買買眼藥膏,擦擦自己眼睛去看看吧..
啟示錄第三章十六節,.
你看見主怎麼對老底家的教會說,.
你既如溫水,也不冷,也不熱..
換句話說,一點可目的心都沒有..
不冷不熱,一點可目的心都沒有..
主說,我就必從我口中怎麼樣,把你們吐出去啊..
很多時候我們就是這樣,不冷不熱..
看不見我們自己的缺乏..
不但看不見自己的缺乏,我們更看不見神的豐富耶..
我們不確定,神真的擁有我們的答案嗎?.
所以他們覺得我們必須來依靠自己嘛..
依靠自己的聰明跟才智來解決自己的問題..
第二點,多少時候我們必須承認,.
我們之所以無法經歷神的恩典跟作為,.
正是因為我們容許我們的心剛硬啊..
我們依靠自己,不來依靠神..
這為什麼在以西節述第三十六章二十六節裡面,.
主這樣說,我要賜給你們一個新星啊,.
將新的靈放在你們裡面..
又從你們肉體中除掉什麼石心啊,.
賜給你們肉心..
這也講到神必須要除去我們那個石頭剛硬的心啊..
要重新給我們一個柔軟的肉心啊..
我們才可能來渴慕神嘛..
否則我們的心就像石頭一樣剛硬..
我們無法經歷神的恩典跟作為來的..
我還記得很多年前,我第一次飛到亞特蘭塔去的時候,.
我還記得當飛機要準備降落在亞特蘭塔的時候,.
我從窗外往下看,哎喲,我覺得這亞特蘭塔何等的美麗啊..
因為就在綠油油一片,看不到任何房子,.
所有房子都是藏在樹叢的裡面..
不像你飛到加州的時候,往窗外一看,全是黃土一片..
可是到亞特蘭塔就,哇,怎麼這麼美麗的城市啊..
我就充滿了期待,哎,這美麗的城市..
一降落,弟兄姊妹們來迎接我..
他們就告訴我,哎,劉慕斯啊,我們要帶你去參觀一個名勝..

$^{161}$我說,哦,好啊,參觀名勝..
他們說,要去哪裡呢?.
要帶你參觀叫石頭山..
蛤?我說,這石頭山,我看綠油油一片,什麼叫石頭山?.
所以這是一個大石頭嗎?我是很難想像啊..
他們就帶我去了,真的是一個石頭山啊..
很高的一個大石頭,就是個石頭山啊..
他們就帶我爬到山頂上..
當我爬到山頂上,往下一看,.
啊,真的,四圍充滿了茂密的樹林,.
唯獨石頭山上寸草不生..
啊,在那裡,我的心就在想,.
哎,難道這裡沒有任何雨水降臨嗎?.
No,事實上,同樣的雨水也是降臨在這個地方..
可是因為這個地方全是石頭嘛,一點水分都不能吸收..
所以別的地方都是綠色的,綠油油的,.
只有這裡什麼都長不出來..
你知道那天,當我站在石頭山上山頂的時候,.
我心裡感觸很深..
佐德林就提醒我,.
會不會我們的心也像這個石頭山一樣,.
是實心,不是肉心啊..
以至於我們無法吸收啊,.
神所賜給我們一切的恩典跟祝福..
哎,不是神不施恩啊,不是神沒有祝福啊,.
而是我們的心這樣的剛硬啊,無法領受啊..
第二點,今天我們真的需要重新來渴慕神,.
渴慕聖靈來充滿我們,我們需要神..
那就像舊約時代的以利沙一樣嘛,.
當她知道她老師以利沙被接升天的時候,.
她知道光靠她自己的力量,.
是無法抵擋當這個時代一個黑暗詮釋..
她知道她必須要依靠在她老師身上那個聖靈,.
所以她才一路跟隨以利沙嘛,不是嗎?.
從吉甲到伯特利,從伯特利到耶利哥,.
從耶利哥到約旦河邊..
為什麼她緊緊跟隨,一路跟隨?.
因為她渴慕嘛,她知道她必須要得到聖靈的能力,.
否則她無法生活,否則她無法事奉..
隨列王經下經二章第九節,.

$^{201}$在約旦河邊之後,伊麗亞就對以利沙說,.
我未曾被接去離開你,你要我為你做什麼,只管求我..
以利沙毫不猶豫,立刻回答,願感動你的靈,怎麼樣?.
加倍感動我..
你可以看見她的渴慕,不只是在伊麗亞身上的聖靈,.
而是加倍的恩高要降臨在她身上..
你知道我們讀這段聖經的時候,.
我們常常只看見以利沙求加倍的恩高,.
卻沒有看見是因為以利沙心中加倍的渴慕..
加倍的恩高,是因為加倍的渴慕..
跟你們說,加倍的恩高是因為加倍的渴慕..
今天當我們面對這樣一個充滿危機的時代,.
我們不是同樣需要向神來呼求嗎?.
我們需要聖靈來充滿我們,.
我們需要聖靈的恩高跟能力來面對每一天的生活..
因為只有聖靈的恩高能夠斷開沉重惡,.
只有聖靈的恩高能夠在我們生命的困境當中為我們親自來開路,.
唯有聖靈的恩高能夠斷開一切黑暗詮釋的困鎖..
所以聖經裡面告訴我們說,.
我們要大大向神張口,神就大大充滿我們..
所以我常說,你要大大張口,神就大大充滿..
你小小張口呢?小小充滿..
不張口呢?不充滿..
你看神的道理多簡單,真的..
大大張口,神大大充滿..
你小小張口,就小小充滿..
不張口,當然無法充滿..
我們求神來幫助,我們要渴慕聖靈..
更多渴慕聖靈的作為在我們生命當中..
第二個,我們要依靠聖靈的帶領..
各位要依靠聖靈的帶領..
撒迦利亞書第四章第六節這裡講得很清楚,.
他對我說,這耶華只是所羅八博的萬軍之耶華,.
說不是依靠勢力,不是依靠才能來,.
是什麼?依靠我的靈方能成事..
在這裡,聖經很清楚地告訴我們,.
在我們每天的生活當中,.
我們要學習來依靠聖靈,.
而不是依靠人的勢力跟才能..
什麼是勢力跟才能?.

$^{241}$勢力跟才能就是人所以為要成功的要素..
所以今天當你在世界上,.
無論你做什麼事情,.
你問的問題是什麼呢?.
你想要成功的話,.
你所問的問題一定是我能不能做?.
我的交情夠不夠?.
我的關係好不好?.
我有多少的資源?.
我有多少的本錢?.
這是我們所問的問題..
這就是勢力跟才能..
但聖經很清楚地告訴我們說,.
神的道路跟人的道路是不一樣的..
事實上,神的道路是高過人的道路..
人以為一定要依靠勢力跟才能才能成功,.
但神說不..
不是依靠勢力,.
也不是依靠才能,.
是依靠神的靈方能成事..
那你說什麼叫依靠神的靈呢?.
那就是在你生活,思想,行為,.
無論想做任何一件事情的時候,.
你要開始學習問心的問題了..
不再是問說我能夠做什麼,.
而是主啊,你的心意是什麼?.
主,你要我做什麼?.
問題不再是我要做什麼,我能做什麼,.
而是主啊,你要我做什麼,你要我怎麼做?.
Amen?.
我要告訴你,.
在生活當中,.
一個依靠勢力跟才能的生活,.
是跟一個依靠神的靈的生活,.
是截然不同的生活方式..
而所帶來的結果,也是截然不同的..
一個生活在勢力跟才能裡面的人,.
他永遠受到的是自我的限制嘛..
因為你所依靠的,就是勢力跟才能..
有一天你的勢力跟才能,.

$^{281}$會成為你自中的一個綑索跟攔阻..
但一個依靠聖靈的人不一樣..
永遠是超越人的有限,.
能夠成就神偉大奇妙的作為來..
這為什麼地獄裡面在每天的生活當中,.
我們要學習依靠聖靈..
在凡事上學習問新的問題,.
來尋求聖靈的引導..
因為每當我們在人生的一個困境,.
在事奉的關口的時候,.
只有聖靈的引導,能夠帶來我們生命的突破..
Amen?.
這就好像我們很熟悉的五餅二魚的神蹟嘛..
當時的門徒面對著五千人的會眾,.
而也聽見耶穌很清楚給他們的吩咐,.
你們給他們吃吧..
若是當時的門徒,五千人怎麼餵啊?.
對他們來講真的一籌莫展..
我們看見門徒們真的用盡自己的勢力跟才能..
飛利不是說嗎?.
這是野地,買不到東西啊..
就算能夠買,我們也沒錢啊..
沒錢沒電,不可能啊..
他們用盡自己的勢力跟才能,一籌莫展嘛..
可是約翰福音第六章第六節怎麼說呢?.
這裡說耶穌說這話是要試驗飛利..
下面一句話怎麼說?.
他自己,他自己怎麼樣?.
緣知道當如何行得..
阿利路亞,主知道..
主知道要怎麼做,主知道我們的答案..
第六節,今天你可知道在人生的過程當中,.
你看見無論你遇到怎麼樣的難處跟困境,.
主早已經知道你的困境,.
而且早已經知道你的答案在哪裡嗎?.
你不需要再絞盡腦汁,.
用盡自己的辦法..
你所需要做的,是尋求祂的指引,.
依靠祂的帶領..
神的靈就要親自引導你穿越一切的困境,.

$^{321}$經歷神的突破,Amen?.
第三,要順從聖靈的感動..
給你們還有順從聖靈的感動..
我告訴你,當聖靈在運行的時候,.
最重要一件事情,.
我們一定要順從聖靈的帶領,.
我們才可能經歷聖靈的工作..
這就好像在五旬節的時候,.
當聖靈降臨在門徒身上的時候,.
他們必須要順從聖靈的感動,.
他們要開口..
開口之後,使律師範第二章第四節,.
他們就被聖靈充滿,.
按照聖靈所賜的口才,說起別國的話來..
他們必須開口,才能夠說出方言..
他們如果閉口,方言就無法出來了..
在使律師範第二章十四節裡面,.
彼得和十一個使徒,他們受了聖靈感動之後,.
你看見在這裡,彼得和十一個使徒怎麼樣?.
站起來,高聲說話..
他們必須順從聖靈的感動,.
站起來高聲說話,.
然後使律師範第二章四十一節,.
我們看見這裡,於是領受他話的人就受了喜..
那一天,門徒約天內有多少人?.
三千人..
他們站起來說話的時候,.
三千人歸向耶穌基督..
我們看見經歷聖靈的能力,.
跟順服是有絕對密切的關係..
今天我看見很多的基督徒,.
只喜歡聽有關聖靈的事情,.
喜歡討論關於聖靈的事情..
但事實上,.
今天只有當你真正開始順從聖靈的感動的時候,.
聖靈的能力才能夠被釋放出來的..
你們,這就好像有一次,.
當主耶穌來到耶利哥的時候,.
在旁邊就有個瞎子叫巴迪曼,.
他在路旁就喊叫,.

$^{361}$大衛的子孫可憐我吧,.
大衛的子孫可憐我吧..
其實,當時耶穌正在趕路,.
所以旁邊的人就對這個瞎子說,.
哎呀,安靜點,.
你們一看夫子很忙,很忙碌..
可是沒想到就在這個時候,.
耶穌就為他停了下來..
我相信當時是有個聖靈的感動在他裡面,.
耶穌就是順從了聖靈的感動,.
他停留下來,.
結果神蹟就發生了,.
那個瞎子就得了醫治..
第二點,今天我們很容易,.
在我們生活跟事奉當中,.
我們忙忙碌碌,忙忙碌碌,.
可當神的聖靈感動我們的時候,.
要我們停留下來為神做一件事的時候,.
請問我們願意順從嗎?.
也許你說,啊,牧師,那是耶穌啦,.
當然聖靈會感動他..
No, no, no, no,.
今天聖靈也照樣常常會感動你,.
感動我的,Amen?.
你說真的嗎?當然真的..
你還記得每一個主日崇拜,.
結束的時候牧師為你的祝福嗎?.
你知道每個禮拜領受日,.
你知道這祝福是在祝福你什麼東西嗎?.
恐怕你每個禮拜聽一聽都不知道我在說什麼..
我每個禮拜講同樣一句話,.
你都不知道我在說什麼..
我怎麼為你們祝福的?.
代表主耶穌基督的恩惠,.
上帝的慈愛,.
聖靈的感動,.
常與牧師同在?.
沒有,常與眾人,.
你們同在..
每個禮拜都為你祝福,.

$^{401}$讓聖靈感動你耶..
聖靈會感動你,.
而且常常感動,Amen?.
跟你們說,聖靈常常感動你的..
你說怎麼感動法呢?.
其實聖靈常常透過我們的心思意念在感動我們嘛..
你知道我不是在說一些很玄的事情,不是的..
常常聖靈會感動,.
有時候你想起來,應該打個電話給一個弟兄姊妹,.
可是時間很晚,明天再說吧..
結果明天呢,你就忘記了,不是嗎?常常有的..
有時候在小組裡面有個新來的朋友,.
你覺得應該去探訪一下,.
正準備出門的時候,天下降下大雨,.
算明天再說,後天天晴再說,.
沒想到天晴就忘了嘛..
多少時候聖靈不是沒感動我們,.
聖靈常常感動,.
只不過我們常常消滅聖靈的感動,不是嗎?.
你知道我以前常常覺得我自己是一個很順服神的人..
直到有一次在聚會當中,.
我聽到一位講員問一個問題,.
他說請問你,.
從你開始知道神要你做一件事情,.
到你真正去做,中間要花多少時間啊?.
哇,我從來沒這樣想過..
從我知道神要我做一件事情,.
到我真正去做,中間要花多少時間?.
我還記得那天晚上我回到家裡面,.
我想了很久,.
我就發現原來其實我是一個非常緩慢的跟隨者..
多少次聖靈說往前走,.
我總是很猶豫,考慮再三..
多少時候聖靈說交託給我吧,.
我總是很不放心,不敢交託..
我以為我是個很順服神的人,.
我那天晚上才發現其實我離神的心很遙遠..
弟兄姊妹,你呢?.
想過沒有?.
從你知道神要你做一件事情,.

$^{441}$一直到你真正去做,.
請問中間要花多少時間啊?.
我所說的事情不見得是轟轟烈烈的大事,.
什麼要你到非洲去宣教,沒有,不見得..
可能是生活的很小的事情而已..
也許神在提醒你,.
每天早點起來親近我吧..
也許神要你做這很小的事情,.
向你身旁周圍的人說聲對不起吧..
也許神要你做的只不過重新起來,.
不要再灰心了,來事奉神吧,.
重新參與吧..
當神提醒你要做一件事情,.
請問你,到你真正去做,要花多少時間?.
我告訴你,當神的靈開始運行的時候,.
你必須學習順從聖靈的感動..
什麼時候當你順從聖靈的感動,.
你就能夠經歷聖靈的能力,Amen?.
弟兄姊妹,聖靈是住在我們裡面的神..
這位聖靈祂明白神對我們的心意..
這位聖靈也不斷地按照神的心意會來感動我們..
只有當我們順從聖靈的感動的時候,.
我們就能在神的手中成為神跡的器皿了,Amen?.
約翰福音第七章三十七到三十八節,.
這時候節期的末了,就是最大之日,.
耶穌就站起來高聲說,.
人若渴了,可以到我這裡來喝..
信我人,就如經上所說,.
從他腹中要流出活水的江河來..
這講得多麼清楚..
耶穌說,人若渴了,可以到我這裡來喝..
信我人,從他腹中要流出活水的江河來..
很有意思,耶穌沒有說,人若聖潔了..
耶穌沒有說,人若完全了..
如果耶穌說,人若聖潔了,人若完全了,.
那你我可能一點資格都沒有..
但耶穌說,人若渴了..
非常簡單的條件,渴了..
第二句,今天你的心渴慕嗎?.
你的心甘渴嗎?.

$^{481}$我們需要重新歡迎聖靈來,.
吹盡我們的生命..
就像第一個五旬節一樣,.
那門徒們在那裡等候禱告的時候,.
聖靈就降臨在他們中間..
同樣,今天我們弘心回到神的面前,.
求神賜給我們新的渴慕,.
渴慕聖靈來充滿我們..
在每天的生活當中學起來依靠聖靈,.
跟隨聖靈,順從聖靈的感動,.
我們就能夠經歷神的得勝,.
經歷神的突破恩告..
我們一起禱告..
謝謝主,你降下聖靈在我們中間..
神的話說,.
這個聖靈是神所應許給神的兒女的..
我們謝謝主..
今天這位聖靈正是住在我們裡面,.
每天生活當中與我們同在..
原諒我們主,.
多少時候我們忽略了這個住在我們裡面的聖靈..
原諒我們多少時候我們按照自己的心意在做事..
但今天我們願意回轉到你面前來,.
我們再次歡迎聖靈的風重新吹來,.
就像第一個五旬節一樣,.
吹進我們的生命當中,.
再來更新我們,恢復我們,.
來翻轉我們的生命..
我們歡迎聖靈再次來掌權..
歡迎聖靈在我們生命當中再次來做更新的工作..
我們謝謝主..
今天早上求您的聖靈大大的降臨..
聖靈的火重新點燃我們生命的火焰,.
讓我們能夠重新為你而活..
我們謝謝你,我們歡迎你..
Amen.
\newpage



\section{創世記 2:7}
\label{sec:032186}
\textbf{人是被創造,而非被製造}~~~~招世超牧師
\newline
\newline
連結: \href{http://www.church.com.hk/acms/content.asp?site=cdc\&op=show\&type=product\&code=032186\&layout=sermon}{\texttt{032186}} ~~~~ 語音日期: 2026-05-24
\newline
\newline
\hyperref[sec:032185]{\small{< < < PREV SERMON < < <}}
~
\hyperref[sec:index_chronic]{\small{[返順時目]}}
~
\hyperref[sec:032187]{\small{> > > NEXT SERMON > > >}}
\newline
\newline
創世記 2:7
\newline
\begin{longtable}{cl}
\hline
\hline
章節 & 經文 (和合本修訂版)\\
\hline
2:7 & \begin{tabularx}{0.7\textwidth}{X} 耶和華神用地上的塵土造人，將生命之氣吹進他的鼻孔，這人就成了有靈的活人。 \end{tabularx} \\ \\
[1ex]
\hline
\hline
\end{longtable}
$^{1}$今日第一講,我就想,不如我都試一試.
試什麼呢?.
我將我今次講到的題目.
第一講是,人是被神創造的,而不是被人製造出來的.
還有,我將這個主題,我想講什麼的重點.
交了給Fred.
我說,Fred你的手腳比較靈活,我不懂這些東西.
你幫我可不可以製造一個30秒的視頻.
就根據我這個題目,以及我的主題.
Fred真的很厲害,很快就幫我搞定了.
先看看怎麼樣.
30秒而已,是不是很厲害呢?.
不知道在座裡面,有沒有誰沒有用過AI的?.
看看有多厲害.
亦因為這個緣故,你可以看到現時AI人工智慧.
對我們社會造成的影響.
有時真的讓我們有感覺,現在AI發展的速度不單是快.
而且好像真的很厲害.
厲害到讓人的腦轉得更快,這是我們無可爭議的.
有很多事情,似乎人是比不上AI的.
其實看是什麼來的.
今天想跟大家看一節聖經.
創世記第二章第七節,回到原本.
創世記第二章七節,那裡說什麼呢?.
那裡說:耶和華上帝用地上的塵土做人.
將生命的氣吹進他的鼻孔.
這個人就成為有靈的活人.
當然,當我們看聖經的時候,我們都知道.
今天你見我,我見你,你忘我,我忘你.
大家都知道,神創造了一個人.
用泥土造成的.
當然,不是做了一個泥公仔.
這個泥公仔會說話,會走路.
好像我小時候看木偶奇遇記,Pinocchio.
那個木偶會說話的,會走路的.
不,神不單是將一個血肉之軀放在人的身上.
而最重要的是什麼呢?.
聖經說,神向人的鼻孔吹口氣.
他就成為一個有靈的活人.
人和人工智慧其中一個主要分別是什麼呢?.

$^{41}$人是有靈魂的,聖經說得很清楚.
AI有沒有?.
AI當然沒有.
所以,當我們說到人有靈魂的時候.
到底這個是什麼意思呢?.
簡單來說,英文叫做SOUL.
這個靈魂裡面包含的就是.
原來這個靈魂包括人有心智,有情感.
有像神一樣的創新和創造能力.
人還懂得愛人,還有被愛的需要.
這些都是靈魂裡面所包含的外行之一.
我剛才說到這大堆,AI裡面有沒有呢?.
AI是沒有這些東西的.
所以今天我們希望在這個講道系列開始的時候.
首先和大家去重溫一下.
聖經告訴我們人到底是什麼呢?.
你怎樣去看人呢?.
人對你來說到底是什麼一回事呢?.
人生存在這個世界上面.
我們在教會裡面經常說.
有什麼價值?.
做人有什麼意義?.
首先,我們看聖經的時候.
聖經告訴我們.
原來人是按照神的形象而被創造的.
所以因為這個緣故.
人是十分尊貴的.
我們看《創世記》第一章.
這個比較詳細一點.
說到人為什麼這麼尊貴呢?.
神說,神自己說.
我們要照著我們的形象.
按照我們的樣式去做人.
27節,於是就採取行動了.
神就照著自己的形象做人.
乃是照著自己的形象做男做女的.
整個形象呢.
都要解釋一下.
以防萬一信仰數目不夠多.
剛剛來教會.

$^{81}$不是說神的外表.
不是說神的肉體.
因為其實聖經裡面.
聖經是說得很清楚的.
在新約聖經說.
神是什麼?.
神是個靈.
神是個靈體.
所以人的肉眼是不可以見到的.
神沒有肉體.
神是沒有外形的.
所以當聖經說.
人是按照神的形象而被神創造的時候.
這裡說的形象是說什麼呢?.
是說到神的屬性.
神的屬性就像我剛才所提過的.
人的理智.
人的創造力.
人有道德感.
有正義感.
有愛和被愛的那種需要.
這些就是聖經裡面所說到的神形象的一部分.
是不是全部呢?.
我不知道.
但起碼我們能夠在聖經裡面找到.
知道的.
起碼有這幾方面.
人是按照神的形象而被造.
人是擁有神這種屬性.
所以為什麼人變得這麼尊貴呢?.
就是這樣的原因.
其他的生物.
都是被神創造的.
但是他們不是按照神的形象造.
他們沒有神這樣的屬性.
所以人和其他被造的物相比.
就高了一截.
是很不一樣的.
比其他的東西尊貴.
聖經告訴我們.

$^{121}$我們是一個有靈,有生命,擁有神屬性的活人.
所以.
如果我們有這樣的認知.
和你真的這樣相信.
聽就聽很多了.
但是我們是不是真的這樣相信呢?.
如果我們有這樣的認知和真的相信的時候.
這樣的話.
無論現在我們的肉體.
我們的心智.
是健康.
或者殘缺.
我們都應該尊重.
第一.
神的形象不是說形體.
不是說有血有肢體.
神的形象是說神的屬性.
而每個被造的人.
都擁有神的屬性.
所以無論我們見到一個人.
哎呀,很醜很醜的樣子.
或者生出來沒手沒腳.
或者長大之後.
哇,這個人變得很壞.
或者心理不健康.
有殘缺.
但是他是一個人.
如果是這樣的話.
我們應該怎麼做?.
應該去尊重.
就不應該因為我們見到這些不足和缺陷.
覺得這些人的生命不值得存在這個世界上.
生命如同螻蟻.
或者我們中國人說生命如同草芥.
No!.
那些極權者會這樣看人.
人的生命是不值錢的.
不要說那些有殘缺的.
就算一般的平民百姓都不值錢.
我們沒有價值的.

$^{161}$當我們的生命不是一回事.
可以隨意去踐踏的.
No!.
如果我們相信聖經.
相信聖經經常說.
我們對人有一個很不一樣的看法.
人的生命是尊貴.
大家記不記得.
其實在美國生活了幾十年.
其中有一樣東西.
我很欣賞美國人.
是什麼呢?.
就是對生命對人權的那種尊重.
大家同意嗎?.
記不記得年初美伊開戰的時候.
有兩個飛機師.
有一個飛機.
美國的戰機被人打下來了.
兩個飛行員.
兩個飛機師.
要跳傘逃生.
一個很快就救了.
另外一個就被困在伊朗的叢林裡面.
大家還記不記得這件事.
美國用了上百號的人去搜索.
還要在其中的過程當中.
有兩架戰機因為受損毀.
自己去炸毀.
不想留下高科技給伊朗人.
到第二個受傷的飛行員.
被拯救回來之後.
全個國家鋪天鋪地去報道.
特朗普總統說.
We got him!.
全美國都很開心快樂.
有沒有人去投訴?.
有沒有搞錯?.
為了救一個受傷的人.
用這麼多的人力.
浪費一兩個飛機.

$^{201}$浪費這麼多錢.
這些是納稅人的錢.
就是為了救一個人.
有沒有搞錯?.
有沒有人這樣去投訴?.
沒有人這樣去投訴.
也不敢這樣去投訴.
不怕被人扔死嗎?.
除此之外.
大家知不知道美國人.
每一年用了億萬美元.
去幫助去資助傷殘.
和有智障的人.
不知道我們的親戚朋友當中.
如果有這樣的情況的小朋友.
有沒有受惠過呢?.
我認識很多人.
真的很好.
美國照顧得他們很好.
還有一件事我曾經跟大家說過.
我很欣賞美國人.
領養小朋友.
不是自己沒有子女.
自己有自己的小朋友.
他去領養孤兒.
不單是領養孤兒.
領養什麼孤兒?.
身體有殘疾的孤兒.
我做牧師也不會這樣做.
我對這些人.
真的十分十分敬佩.
因為他們看重.
每一個人.
每一個人的生命.
所以弟兄姊妹.
我們信耶穌.
我們明白聖經.
我們怎樣看我們身邊的人.
不知道我們看不看到.
這個人生命裡面.

$^{241}$從神而來的那種價值.
不知道我們會不會看到.
這些人在我們自己覺得.
這個人不值得我們去愛和關心.
但是原來神也愛他.
神也關心他.
我們看不看到他生命當中的價值.
和有沒有想想有什麼方法.
可以幫助他去認識神.
「神愛世人」這四個字.
在教會裡面.
耳熟能詳的講得很離譜.
我們沒什麼感覺.
這四個字.
「神愛世人」.
聽起來很簡單,對不對?.
其實裡面所包含的意義是十分的豐富.
「神愛世人」.
「神愛世界」裡面.
每一個人.
無論你覺得他是好還是壞.
健康還是有殘缺.
聖經說「神愛世人」.
「神愛世上每一個人」.
為什麼?.
剛才聖經告訴我們.
因為每一個人都是按照神的形象而被創造.
只不過我們看到神的創造.
後來因為人的背影不聽話.
神的創造受虧損.
神學上叫做「完美的虧損」.
神創造的人是完美的.
不過完美受到虧損.
因為人的自由意志濫用了.
故意不聽神的話.
所以.
因為完美受到虧損.
為什麼今天我們見到有些人出生.
沒有手沒有腳.
有殘疾.

$^{281}$或者長大之後.
人心變得比萬物都詭詐呢?.
就是因為人不聽神的話.
犯罪.
離開神的緣故.
但是弟兄姊妹.
如果我們相信上帝.
如果我們相信聖經.
我們相信人是按照神的形象而被造的.
人雖然是受神的刑罰.
聖經說「被改出伊甸園」.
但是這個神性受到虧損.
但是在人性裡面.
仍然存在有神的性情在當中.
所以為什麼我們才能夠跟神溝通呢?.
所以為什麼我們見到人.
雖然是那麼軟弱那麼有問題.
甚至是那些罪大惡極的人.
他都需要愛.
和需要許得他愛人呢?.
既然我們明白了這個道理的話.
今天我們要彼此鼓勵.
無論是誰都好.
如果神把人帶到我們面前.
無論你喜歡他不喜歡他也好.
最低限度我們學會去尊重他.
有能力的時候去照顧他.
幫助他.
特別是那些遠離上帝.
還未認識神的人.
因為他們每一個都是神所創造.
而且是在創造當中比萬物都尊貴.
為什麼呢?.
再說一次.
因為這些人是按照神的形象造的.
我們現在說人工智慧.
說AI.
人工智慧是什麼呢?.
剛才說了人是什麼.
人工智慧是什麼呢?.

$^{321}$我不是科學家.
如果我說錯你們幫我矯正吧.
對我的認識來說.
人工智慧只不過是人所創造.
是神所創造的人.
製造出來一些運算的系統.
不知道說得對不對.
人工智慧不懂得去愛人.
不過我們看到有些人傻得.
愛上了那些有人工智慧的機器人.
真是沒辦法.
我們看看.
這個也叫做AI造.
人工智慧和人類智慧之間.
有什麼分別呢?.
有什麼不同呢?.
我們看到.
人工智慧的驅動力是什麼呢?.
是生物的本能.
是情感.
是人的意識和創意.
這些人工智慧有沒有呢?.
沒有.
人工智慧的驅動力是什麼呢?.
是簡訊法.
人工智慧的驅動力是什麼呢?.
是數據和晶片的算力有多高.
人的智慧.
是擅長做什麼呢?.
是擅長做那個.
我的眼睛真的不行.
複雜的決策.
情感的共鳴.
empathy.
和有因果的推理.
而人工智慧擅長的是什麼呢?.
是海量.
很多數據的分析.
模式的識別.
和快,高速的計算.

$^{361}$人的智慧的識別能力是很高的.
為什麼呢?.
因為可以在一個全新的陌生環境.
未知的情況下.
繼續生存,人是可以隨機應變的.
但是人工智慧在適應力上來說.
其實是有很多限制的.
當他遇到一些沒有受過訓練.
或者資料庫裡面沒有的資料的時候.
就很容易有幻覺.
他們的幻覺.
就是沒有效果.
做不到你要求他做的事情.
而人類的智慧.
出錯的機制在哪裡呢?.
為什麼會出錯呢?.
因為人會疲倦.
Burn out, Exhaustion.
情緒不好.
於是就很容易.
急急忙忙做事.
做錯決定.
但是人工智慧出錯的地方是什麼呢?.
如果他的資料庫裡面沒有那種資料.
他沒有受過那方面的模式的輸入.
很多時候他就會出錯.
他的答案就會有偏差.
所以我們簡單做一個比較的時候.
你會發覺其實.
人工智慧在很多的地方.
還是比不上人的.
但是今天有越來越多人.
特別是那些尊崇科學的人.
是有意無意地高舉了AI那種能力.
而有意無意地貶低了.
人類那種的智慧和價值.
現實環境裡面來說.
甚至有些大學生大家都知道他找不到工作.
為什麼呢?.
因為AI出現了.

$^{401}$所以有時候他們都很depressed的.
很抑鬱的.
因為覺得自己連AI都不如.
AI是人造出來的.
但是AI搶光我們的工作.
當然黃仁勳不是這樣說的.
黃仁勳說不會搶光你的工作.
但是現在真的很多人沒有工作.
Facebook要layoff八千人了.
layoff了沒有呢?.
是一個很大的衝擊.
諸多兄弟姐妹.
回到我們今天的主題.
聖經告訴我們.
人為什麼這麼尊貴呢?.
第一.
人的尊貴是因為.
是誰做出來的?.
創造.
宇宙萬物的主宰創造出來的.
神做出來的東西尊不尊貴?.
我做出來的東西.
和總統做出來的東西.
大家都是一起做一樣的東西出來的.
不過一個是愛了總統做,誰尊貴些?.
除此之外.
人是按照神的形象去做.
尊不尊貴?.
剛才說過沒有任何東西是按照神的形象做的.
為了人之愛.
所以.
人的生命應該被尊重.
但是除此之外.
人的生命.
它的價值在哪裡呢?.
人的生命.
神做出來的不是自己當玩具玩的.
神做人有原因有目的.
神做人的價值.
人的價值到底在什麼事情上面呢?.

$^{441}$第一段聖經比較長一點.
在希伯來書第二章.
第五到第九節.
這裡說.
我們所說將來的世界.
神原來是沒有交給天使管理.
但有人在經上某處.
這是詩篇第八篇.
怎麼說呢?.
人算什麼理景顧念他?.
世人算什麼理景眷顧他?.
第七節.
你叫他比天使微小一點.
賜他榮耀尊貴為冠冕.
並將你所做的都派他管理.
叫萬物都伏在他的腳下.
叫萬物都伏他.
就沒有剩下一樣不伏他的.
只是如今我們還不見萬物都伏他.
唯獨見那成為比天使小一點的耶穌.
或亦作唯獨見耶穌暫時比天使小.
因為受死的苦.
就得了尊貴榮耀為冠冕.
叫他因著神的恩.
為人人償了死命.
我讀完一次.
你不知道他說什麼.
因為又說到他又說到耶穌.
有時候有點混淆.
我慢慢跟大家去解釋.
聖經說.
神是按照自己的形象去創造人.
當他創造人之後.
你看回《創世記》告訴我們.
神是賦予人有權並有能力.
去做什麼?.
去管理其他受造之物.
這個工作也很尊貴.
變成萬物疊口.
可以管理其他所有一切被造之物.

$^{481}$但是剛才這段聖經希伯來書.
忽然之間話鋒一轉.
就指出一個現實.
這個現實是什麼?.
只說如今.
我們還不見萬物都伏他.
聖經開頭是這樣說.
《創世記》說讓你去管理萬物.
為什麼忽然之間說到現在.
到寫希伯來書的時候.
我們還不見萬物都伏在人的管治底下.
為什麼會變成這樣?.
其實我剛才也解釋過.
因為人犯罪.
因為人離開上帝.
偷吃禁果.
不聽神話.
得罪了神.
所以.
因為人犯罪的結果帶來的是什麼?.
就是失去了亞當.
一樣管理天地萬物的能力.
洗一切被造物.
包括人自己.
每天是步向滅絕和死亡.
為什麼?.
因為人犯罪之後.
不單單失去了管理萬物的能力.
人更加是為了生存.
有意無意的破壞了自己的生態環境.
污染了神所創造的世界.
所以我們也知道.
我們沒有留心.
很多生物已經是動物.
是瀕臨絕種.
有些已經絕種了.
現在熱帶森林越來越少.
Ozone越來越大.
可能我們為了每天忙碌的生活.
都沒有去理會這些世界上發生什麼事情.

$^{521}$這些都是我們生活的環境.
但是變成這麼有問題.
我們沒有時間去理會.
以人自己.
也走上了自我毀滅的道路.
為什麼美國要打伊朗?.
我不知道是真是假.
起碼他公開說.
那個原因是什麼?.
我們不可以容許這個邪惡的國家.
擁有核武器.
因為我只是聽說的.
真是我也不知道.
因為伊朗的觀念.
世界要被毀滅.
才可以有重生有再生.
神才會再來.
所以這麼邪惡的.
專門支持恐怖分子的國家.
一擁有核彈的時候怎麼辦?.
這樣搞不定的.
其實重點是什麼?.
因為大家都知道.
這個核武器真是很厲害.
真是可以毀滅全世界.
你不去好好去管理去管治的話.
世界每一個人都很危險.
所以諸多頂智妹.
我不知道我們代學裡面.
有沒有無道的朋友.
當我們了解到這樣的情況的時候.
難怪有些人心裡會擔心.
會不會有一天.
被人研製出來的AI.
人工智慧會操控人類.
會主宰人類的命運呢?.
這部電影的橋段已經出過很多了.
不過如果我沒有誤解聖經的話.
正如我那時候所說.
這樣的情況.

$^{561}$我相信在電影的橋段裡面.
是會出現的.
為什麼我這樣說呢?.
這個不單單是我剛才說過.
AI是被人設計和創造出來的計算模式.
是人做出來的.
而最重要的是什麼?.
神的應許.
神的設計.
是叫人去管理世上的萬物.
包括人自己做出來的AI.
你看看牧師說人軟弱.
人壞了.
人控制得罪神.
人破壞了這個世界.
沒錯.
不過不要忘記.
世界的主宰不是人.
是神.
神才是世界的主宰.
神只不過是把管理世界的權柄.
賜給了人.
世界的主宰仍然是上帝.
而上帝其中一個特性是什麼?.
神是愛的源頭.
所以神看到人變成這樣.
有沒有袖手旁觀?.
會不會袖手旁觀?.
翹起手看你怎麼死.
這個不聽話.
神會不會這樣做?.
神很生氣,很不開心.
因為人不聽話,得罪了他.
不過因為神愛人的緣故.
神沒有袖手旁觀.
所以我們看回聖經.
又是話鋒一轉.
轉到第九節說什麼呢?.
希伯來書第二章第九節說.
唯獨見那比天使少一點的耶穌.

$^{601}$因為受死的苦.
就得了尊貴榮耀為冠冕.
叫他恩著神的恩.
為人人賞他死命.
結果承擔了人犯罪所帶來的一切刑罰.
因為耶穌的救贖.
人就可以和神和好.
並且超越創造主和被造物的關係.
人本來和神的關係就是創造主和被造物.
就是這麼簡單.
但當神透過耶穌基督的救贖.
讓我們和神和好之後.
我們有一個超越創造主和被造物的關係.
人成為了什麼?.
人成為了被上帝領養的子女.
凡接到他的,就是信他命的人.
神就似權並做神的兒女.
聖經是這樣說的.
以前我們沒有兒女這個身份.
是園丁來的.
代神管理世界.
現在成為神的兒女.
可以領受神一切的尊貴和榮耀.
弟兄姊妹,這個身份比以前尊貴很多倍.
所以,追求丁子們.
我們信耶穌之後.
我們所擁有的生命和和神的關係.
不是AI人工智能可以取代.
或者是可以擁有的.
大家同意嗎?.
AI人工智能只能夠在人的指示底下.
為人搜集資料,數據去做推算,去分析.
給一些結論給你.
提供一些可叉可的答案.
使用他所擁有的資料.
都是人給他的.
和裡面的程式給我們一些建議.
但是人工智慧沒有自己的思想.
人工智慧是不能夠管理人的.
這個起碼是現時的信念.

$^{641}$所以主要是弟兄姊妹.
人存在是有價值的.
你同不同意嗎?.
到現在我希望你同意.
人的存在是有價值的.
不過我們要搞清楚.
人的價值到底是什麼?.
我們應該怎樣去衡量一個人的價值呢?.
看他賺多少錢?.
做官做到多高級?.
還是其他什麼呢?.
對我來說,當我們看聖經的時候.
人的價值再不是因為.
人是神創造的高峰.
這個我剛才說過.
人的價值也不是因為.
神將一個管理大地世界的責任.
交給了人.
人的責任到底是什麼呢?.
有一次聖經我跟大家說過很多次.
我們基督徒是蒙簡選的族類.
我們成為一個特別的族群在世界上.
你們是被簡選的族類.
是有君尊的祭司.
是聖潔的國度.
是上帝的子民.
我們價值來自哪裡呢?.
上帝要召我們離開黑暗.
進入輝煌的光明裡面.
來宣揚上帝奇妙的工作.
最後頂著門我們成為基督徒之後.
原來我們在世界上有一個神聖的任務.
就是要去宣揚天父爸爸.
拯救我們脫離死亡.
進入他的榮耀當中的人典.
這個是我們成為基督徒之後.
在地上的使命.
所以當我們存在在世界上的時候.
每個基督徒都應該知道.
我們擁有這樣的使命.

$^{681}$我們應該做這樣的事.
很多時候我們覺得這是牧師做的事.
這是那些熱心傳道的基督徒做的事.
No.
如果你是這個族類的一部份.
你是其中的一員.
我們就有責任去宣揚.
呼召我們出黑暗.
進入奇妙光明者的奇妙的工作.
今天我們跟大家看聖經是在這方面的.
告訴我們我們可以怎樣去完成上帝給我們的使命.
彼得前書第三章十五節.
一齊讀一讀第三章十五節.
彼得前書一來.
要有敬畏基督的心.
以他為主.
有人要求你們解釋心裡的盼望.
要隨時準備答辯.
再來一次.
要有敬畏基督的心.
以他為主.
有人要求你們解釋心裡的盼望.
要隨時準備答辯.
當然怎樣去完成上帝給我們的使命.
去傳揚他為主做見證.
有很多經文可以幫我們.
不過今天時間關係.
我只看這一節聖經.
兩點.
第一件事.
要有敬畏基督的心.
以他為主.
如果你想為主做見證.
我們要用我們的生命做見證.
但是如果我們沒有敬畏耶穌基督的心.
我們不能以耶穌基督為主.
我們就不要想我們能夠成為一個美好的見證.
為什麼我特別選擇這一節聖經.
因為我發覺今天很多人信耶穌.
在教會裡面只是想得到神的安慰和幫助.

$^{721}$我發覺有很多人其實心目中不敬畏神.
或者不夠敬畏神.
很多基督徒的心裡面.
怎樣看神的呢.
我好像用過一個這樣的例子.
就好像你們放在車尾箱的儲物檯.
或者我們現在流行的後備電池.
教會辦公室有一個.
平時放在角落.
沒電了.
生活要倒胎了.
要拿個儲物檯.
沒電了.
exhaustion.
什麼都沒了.
拿個後備電池來充電.
我們不知不覺在不自知的情況下.
我們會把神當成儲物檯.
或者是我們後備的電池.
很多基督徒求神幫助我們.
在祂的面前自我反省.
今天我們是敬拜上帝.
很多基督徒在心態和行為上面.
真的不是很敬畏神.
就像我剛才所說.
很多基督徒只需要神幫助的時候.
才會去祈禱 才會去教會多一點.
平時沒事呢.
讀經祈禱都沒有.
沒有怎樣做靈修.
神放在一邊.
我們對祈禱 對查經 對敬拜.
沒錯 我們坐在那裡.
很多時候我們心不在那裡.
我們只是來盡一個基督徒的責任.
星期天又去教會聚會 敬拜上帝.
我們對信仰很馬虎.
就算我們參加崇拜坐在那裡.
我們沒有一個莊重真誠的心.
我不是要指責大家.

$^{761}$我只是希望我們能夠在神的面前.
大家彼此提醒.
神是眼不能見的.
剛才牧師也說過.
神是視覺力.
你見不到的.
我們是來敬拜神的.
還是來聽牧師講道.
聽師班唱歌.
不是 聽師班祈禱.
你有沒有預備好來見神.
你抱什麼態度來敬拜.
你怎樣參與敬拜.
你平時的生活是怎樣的.
求神憐憫我們 幫助我們.
心態上 有沒有敬畏上帝呢.
自己做個檢討.
如果你說心態是很難講的.
行為上.
行為上我們有敬畏上帝嗎.
我們有沒有輕視聖經的教導.
今天我發覺很多人.
說謊當吃生菜.
隱瞞真相欺詐.
為了保護自己.
那是無所謂的.
有時我們不單不怕神的審判和刑罰.
我們還有很好的解釋.
我們生活在恩典時代.
新約不是在恩典時代嗎.
新約不是在講神的愛嗎.
神必定會寬恕我們的軟弱和罪過.
如果我們這樣想的時候.
原來我們把神的恩典.
變成了我們放縱自己的藉口.
一個不敬畏神的基督徒.
沒有把神當作自己的主.
於是他的信仰.
只停留在一個初信.
或者表面做了很多基督徒的習慣的生命.

$^{801}$生命的本質有沒有改變.
我們問問自己.
信耶穌到現在.
我們的生命有什麼改變沒有呢.
除了多了基督徒的習慣之外.
這個我們需要在神的面前.
很坦誠地做一個自我反省.
我們說信耶穌.
我們擁有耶穌的生命有多少.
我們的心有敬敬畏神.
我們的行為能不能榮耀上帝.
一個不敬畏神的基督徒.
有時變成我們像拜神拜佛的人一樣.
把神偶像化和工具化.
弟兄姊妹.
求神幫我們.
求聖靈指引我們.
不要忘記聖經告訴我們.
神是一個絕對公義.
絕對聖潔的神.
神的審判不會不見了.
神必定會審判.
當我們缺乏這種認知和警覺性的時候.
我們就會變成一個孖孖苦苦的基督徒.
你很少會見到上帝在你身上的恩典.
你會聽到一些很刺激的見證.
但這些見證不發生在你身上.
你的生命也不能彰顯出.
你是按神的形象而被造的榮耀.
但相反.
如果我們有一個敬畏上帝的心.
以神作為我們生命的主宰.
我們的生命就不同.
人可以看得出來.
不是自己說我現在不同了.
我們受浸講見證都要講.
信耶穌之前怎樣.
信耶穌之後怎樣.
很多時候我們就停留在那裡.
有沒有進步?.

$^{841}$有沒有繼續去改變?.
求神幫助我們.
如果你真的以基督耶穌為主.
你敬愛上帝的話.
你很自然地.
你的生活會不斷被聖靈去無助改變.
問題就是.
當聖靈要無助改變你的時候.
你願不願意去信服.
願不願意去聽.
星期三.
我和一班教務同工.
有一個訓練班.
六個星期的.
剛好最後一個星期.
我和他們講了幾件事.
第一 怎樣改善你的講壇.
第二 怎樣做領導.
第三 怎樣建立你的團隊.
在講道的時候.
我跟他說.
其中一樣東西很重要.
你要知道你的會眾是什麼人.
你怎樣去看你的會眾.
你怎樣去看人.
我和他說了五件事.
這是我和他分享我的信念.
我預備講道的時候.
都會去體驗這件事.
第一 我相信.
上帝渴望與每個人.
建立一個親密的關係.
你相不相信?.
如果你不相信就糟糕了.
為什麼神會叫耶穌來呢?.
因為祂想和我們建立一個和好.
親密的關係.
第二件事.
上帝對每個人都有一個計劃.
你相不相信?.

$^{881}$我相信.
無論你是誰.
神在你的心靈當中.
為你設計了一個.
特別為你設計的計劃.
第三 每個人的成長經歷.
和需求都不同.
這個不會反對.
大家都會同意.
第四 每個人都應該尋找.
發現並活出上帝給他的使命.
這是我們會點的意象.
人與神的交會點.
第五 尋找並發現.
神創造你的目的.
你相不相信這件事?.
如果你相信的話.
你知道是什麼味道嗎?.
神如何模擬你?.
神如何造你?.
神如何帶領你?.
神如何放你在這裡?.
你的使命是什麼?.
最後 當然不是最後.
這是我個人的看法.
每個前來敬拜上帝的人.
你們來敬拜上帝.
都應該信服上帝的旨意.
不是聽牧師的話.
信服聖經的啟示.
聖靈給我們的感動和帶領.
主佬弟兄姊妹.
AI人工智能.
不會跟神有親密的關係.
AI不會有自我的意識.
和有自己的計劃.
不會有自己的經驗.
AI只能搜集資料.
累積,歸納,推算人的經驗和需要.
AI沒有使命感.

$^{921}$AI也不會聽從神的指示.
AI只會按照自私有限人的指示.
去給我們提供一些答案.
AI更加不會主動去為神做見證.
但是我們相信神.
經歷過神的恩典.
擁有過基督生命的人.
我相信我們很自然就會.
你的生命流露.
你有耶穌的生命在當中.
別人會留意到,看到的.
你的生命就自然可以成為.
分享信仰的一個渠道.
AI可以幫我們什麼呢?.
AI可以幫我們準備.
講些什麼?.
有什麼內容?.
怎麼整理次序?.
AI可以提供一些意見.
你怎麼分享法律好一點呢?.
有一個這樣的指引給我們.
但是AI不能代我們做這些事情.
給我們的鼓勵是什麼呢?.
要隨時做好準備去答辯.
解釋我們心裡面盼望的緣由.
先一陣子,如果你要看牧師的牧人心語.
有沒有看牧人心語?.
有沒有收到?.
如果還沒有收到就sign up吧.
每個禮拜我們有一個牧人心語.
我記得和大家分享過.
有一晚我去成年Angel開完會.
天黑了,我要走了.
八,九點鐘的時間.
在parking lot遇到一個姐妹.
剛剛探完她爸爸回來.
講個故事給你們聽.
在parking lot遇到一個姐妹.
剛剛探完她爸爸回來.
她爸爸說早上不太好.

$^{961}$可能要走了.
我們師母受感動.
我們馬上去探望她.
去到安老院的時間.
一個九十九歲的傭工.
忽然之間很清醒.
我再和她講福音.
她就在那一晚決志送署.
第一天的姐妹你有沒有留意到?.
我有沒有準備?.
沒有的.
完全沒有準備.
你可以argue的.
你是牧師已經準備好了.
不是的.
這個提醒是什麼呢?.
每個基督徒都要常作準備.
當神要給我們機會的時候.
順從才能夠引導.
你就能夠幫助帶領一個人.
決志送署.
所以主要提醒姐妹.
求神幫助我們.
如果我們敬愛神.
我們自然會把握機會.
去為主做見證.
用我們的生命和生活.
去告訴別人.
我們在世上唯一的盼望.
就只有上帝.
沒有上帝.
我們的生活會變得很慘.
人工智慧可以做這些事嗎?.
No,不可以.
有人這樣說.
這個也是AI出來的.
A mirror can reflect a face.
But it never becomes a person.
它是一塊鏡子可以反照出一個人的面貌.
或者一個影像.

$^{1001}$但是它永遠不能夠變成一個人.
它只是一個影子.
它不是一個真人.
而我們呢.
我們是一個有血有肉有靈魂.
神所創造的人.
而且是按照神的形象做出來的.
我們從神的裡面領受了責任.
現在不單單是管理這個世界.
被人弄得冤七八糟.
現在我們的使命是什麼呢?.
要告訴別人神的愛和恩典.
得救的盼望是來自耶穌基督.
這個就成為今天你和我存在.
在這個世界裡面一個主要的使命.
我們怎樣回應神的呼召呢?.
我們一起祈禱.
因為今早你再次提醒我.
原來我們人是這麼尊貴的.
如果我們在我們的生活圈裡面.
我們認識一些人.
我們覺得他不值得我們去愛.
不值得我們去尊重.
這個人太差了.
讓我們能夠退一步.
用神的眼鏡再看多一次.
這個是神所創造的人.
我們不同意他的行為.
不同意他的表現.
不同意他說的話.
但是這個神所創造.
而且用他繼續留在這個世界當中.
教我們怎樣去尊重.
去愛護.
引領這些人回到這個面前.
天父,我們知道原來每個人都有價值.
因為你創造每個人都有各個目的.
主要願意你使用我們.
去幫助人發現你創造他們的目的.
讓他們能夠活得有意思.

$^{1041}$讓他們開心快樂.
不只是每人勞勞碌碌.
為生活而去忙碌.
天父上帝更急其求你幫助我們.
用我們的生命成為一個美好的見證.
事事做好準備.
將人邀請人帶到神的使人的寶座面前.
求天父你聽我們的禱告.
奉耶穌基督的聖名.
\newpage



\section{歌羅西書 3:18-4:1}
\label{sec:032187}
\textbf{群策群力所建造的基督群體}~~~~黃紹權牧師
\newline
\newline
連結: \href{http://www.church.com.hk/acms/content.asp?site=cdc\&op=show\&type=product\&code=032187\&layout=sermon}{\texttt{032187}} ~~~~ 語音日期: 2026-05-24
\newline
\newline
\hyperref[sec:032186]{\small{< < < PREV SERMON < < <}}
~
\hyperref[sec:index_chronic]{\small{[返順時目]}}
~
\hyperref[sec:032188]{\small{> > > NEXT SERMON > > >}}
\newline
\newline
歌羅西書 3:18-4:1
\newline
\begin{longtable}{cl}
\hline
\hline
章節 & 經文 (和合本修訂版)\\
\hline
3:18 & \begin{tabularx}{0.7\textwidth}{X} 你們作妻子的，要順服自己的丈夫，這在主裡面是合宜的。 \end{tabularx} \\ \\ \relax
3:19 & \begin{tabularx}{0.7\textwidth}{X} 你們作丈夫的，要愛你們的妻子，不可虐待她們。 \end{tabularx} \\ \\ \relax
3:20 & \begin{tabularx}{0.7\textwidth}{X} 你們作兒女的，要凡事聽從父母，因為這是主所喜悅的。 \end{tabularx} \\ \\ \relax
3:21 & \begin{tabularx}{0.7\textwidth}{X} 你們作父親的，不要惹兒女生氣，恐怕他們會灰心。 \end{tabularx} \\ \\ \relax
3:22 & \begin{tabularx}{0.7\textwidth}{X} 你們作僕人的，要凡事聽從你們肉身的主人，不要只在眼前服事，像是討人喜歡的，總要心存誠實，因為你們敬畏主。 \end{tabularx} \\ \\ \relax
3:23 & \begin{tabularx}{0.7\textwidth}{X} 你們無論做甚麼，都要從心裡做，像是為主做的，不是為人做的； \end{tabularx} \\ \\ \relax
3:24 & \begin{tabularx}{0.7\textwidth}{X} 因為你們知道，從主那裡必得著基業作為賞賜。你們要服侍的是主基督。 \end{tabularx} \\ \\ \relax
3:25 & \begin{tabularx}{0.7\textwidth}{X} 行不義的人必受不義的報應；主並不偏待人。 \end{tabularx} \\ \\ \relax
4:1 & \begin{tabularx}{0.7\textwidth}{X} 你們作主人的，待僕人要公正，因為知道，你們也有一位主在天上。 \end{tabularx} \\ \\
[1ex]
\hline
\hline
\end{longtable}
$^{1}$各位弟兄姊妹,網上的弟兄姊妹,大家早安.
讓我們再次回來,在主的說話面前.
我們一起同心去預備.
我們心思意念,來明白認識我們的主愛之神.
我們一起同心低頭禱告.
父神多謝你,讓我們可以再次來到崇拜的地方.
一起將祂心緊貼上主你的胸膛.
很想我們能夠從當中去細聽上主你的心思意念.
求你的帶領,讓我們藉著你的啟示.
更明白上主你的心.
也讓我們學曉怎樣能夠在這個地上.
活出一個合神心意的生命.
讓人見到我們的時候.
就讓他們可以同時讚歎.
這位認識上主的人.
他們怎樣在地上顯出上主的榮美.
所以祈求主你給我們有智慧.
求你讓我們在這份事奉裡面有所擔當.
並且能夠竭力做好.
我們知道很多時候我們或許會有一些挑戰.
或者會重新受到我們本身罪的誘惑.
從而我們在上主你的見證上面有所虧欠.
祈求主你的赦免,也讓悔改的心.
與追求主道的心思能夠並行.
讓我們帶著守潔心清的生命.
來到上主你這條永恆的道路上面一直奔跑.
多謝你積著聖靈的引導.
讓我們能夠在今天再次明白主你的道理.
讓每一個追求主道的人.
都能夠在今天同夢主你的啟示的祝福.
禱告奉靠主耶穌基督得勝名字而求.
阿們.
一件事物或者一個人如果只是單獨存在的話.
往往只不過是一種閉門造車.
做出來的成果有限.
或許做出來的東西甚至是蕩然無存.
同樣地神國的子民.
也都不是單獨存在在這個世界的裡面.
我們是透過和不同人不同的關係.
來去積著溝通來往.

$^{41}$來去彼此的建立.
最重要的是在這個過程當中.
讓我們一起去摸準神的心思意念.
叫我們可以憑著上帝的聖情.
來去做成一間真是合神心意的屬靈群體.
簡稱「教會」.
這個才是我們步向屬靈的永恆國度的一個方向.
所以今天我想透過歌羅西書的第三章.
從第十八節到第四章的第一節.
一起思想一個題目叫做.
「群策群力所建立的基督群體」.
這是一個透過多方多向性的一個群體的互相的交往.
來去建立起來的.
亦都在這個交往的過程當中.
讓我們能夠去折射出.
上帝在我們身上被作為選民的一個生命力.
怎樣能夠彰顯出來.
讓我們再次翻開這段聖經.
歌羅西書第三章第十八節.
至到第四章第一節.
讓我們打開之後我們再自己讀多一次.
好使我們在聆聽主的說話的過程當中.
能夠對經文的理解更加的全面和準確.
請大家一起打開歌羅西書的第三章十八.
我們讀到第四章第一節.
翻到最後我們一起讀.
「你們作妻子的,當信服自己的丈夫,.
這在主裡面是相宜的.你們作丈夫的,要愛你們的妻子,.
不可輔待他們.你們作兒女的,要凡事聽從父母,.
因為這是主所喜悅的.你們作父親的,不要惹兒女的氣,.
恐怕他們失了志氣.你們作僕人的,.
要凡事聽從你們肉身的主人,不要只在眼前事奉,.
像是土人喜歡的,總要全心誠實敬畏主,.
無論做什麼都要從心裡做,像是給主做的,.
不是給人做的.因你們知道從主那裡必得著基業為賞賜,.
你們所事奉的乃是主基督,那行不義的,.
必受不義的報應,主並不偏待人.你們作主人的,.
要公公平平地待僕人,因為知道你們也有一位在天上.」.
在這段聖經的上面,其實歷世歷代的教會.
都會將它定為所謂的家女,即屬靈家的家女的家規,.

$^{81}$我們簡單叫一個Household Code..
這類型的經文不少的,除了我們剛剛所讀過的.
《歌羅西書》第三章十八節至第四章一節之外,.
我們在聖經其他保羅的書信,或者其他的仕途的,.
或者門徒他們的記載都有的,.
包括以弗所書的第五章二十二節至第六章九節,.
提摩太前書第二章八節至十五節,.
第六章的一至第二節,提多書的二章一至十節,.
以及彼得前書二章十三節至三章七節,.
都同時在聖經裡面被指為.
Household Code,或者Household Tables,.
屬靈的家規..
當讀到這個陳述的時候,.
尤其是去到華人教會就會覺得.
好像我們傳統華人家庭裡面所講的家規一樣,.
我不知道大家家裡有沒有家規這件事,.
我很記得我小時候,.
因為我們每個星期都會去探望阿公一次,.
其試探望阿公就是借題,.
發揮就是跟表弟妹一起玩,玩足一日,.
那個就是最主要的目的,.
但是在玩的過程當中,.
有時候阿公就坐在那裡看著我們,.
跟小孩子一樣玩,我就覺得阿公很奇怪,.
因為他坐在那裡只是看著我們不做事,.
而且他正正座位的後面牆就掛著一幅,.
因為我媽媽姓潘,我阿公姓潘,.
他就寫著一幅叫做潘氏家規,.
好像一個大老爺坐在一幅大的警告說話的標題之前,.
很有一個共產思維的感覺,.
總之看見阿公就一定看見家規,.
自從我開始懂事,學中文,.
那些家規是很複雜的,很壓韻的,.
現在究竟那個家規去了哪裡,我想沒有人知道,.
不過如果我記得他其中一件事,.
他就說要以姓潘家為榮耀,.
我心想厲害,原來家規第一件事就是不要收家,.
講得很簡單很清楚的一句講完,.
之後下面一大連串的東西我其實不記得了,.
不過這句話我就很清楚,.

$^{121}$其實家規挺好的,學中文字上面的字都挺深的,.
但是學著學著就懂得有這樣的字,.
原來可以這樣講的,.
於是乎這種類似的思想,.
其實在華人教會對於這一類型的所謂淑寧家規,.
都有一個類似的詮釋,.
他們的詮釋很多時候就是說,.
這些就是聖經裡面講給教會知道,.
教會應該要怎樣處理教會裡面的事情,.
教會的運作也要按照這些家規來執行,.
於是乎這些家規就衍生了很多不同類型的內容,.
譬如在我們之前剛剛上一課,.
或者這段聖經之前講過的有關於在教會裡面有些什麼活動,.
我們看回第三章的第十六節,.
他在這裡這樣講,.
當用各樣的智慧把基督的道理,.
豐豐富富地藏在心裡,.
用詩章頌詞靈歌,.
彼此教導,互相勸勉,.
心被恩感,歌頌神,.
於是乎就成為了教會崇拜學的家規,.
如果你聽過我上一課的講道,.
也都我講得很清楚給大家聽,.
其實這一部份只不過是衍生在前文所談論到的,.
有關於信徒他們之間的關係,.
怎樣是透過一連串的歌頌,.
來反映出他們上帝的感恩,.
以及他們將會怎樣奉主的名,.
來去活他們的生命生活,.
並不是以作為崇拜學為基礎的經文,.
這個前文所講的,.
而這個主題也都引申到今天我們去看的另一個大部份,.
也都同樣是談論到有關於在基督裡面的群體,.
他們的生活是怎樣,.
但是在這段經文裡面講的東西,.
就引發在今時今日這個充滿苦解運動,.
以及所謂的家庭倫理觀的deconstruction,.
摧毀的狀態之下,.
很多人就很反感這段聖經,.
其中為了自圓其說,.

$^{161}$很多的牧者當面對這段聖經的時候,.
最直接處理的方法就是不要講,.
不講就沒有問題,.
不會惹起教會的紛爭,.
因為教會裡面最重要是合一,.
合一在於大家不要吵架,.
最緊要準時奉獻,.
回來聚會,.
就是這麼多,.
失卻了這段聖經原本想講的核心主題,.
所以他們不講,.
所以就當沒事發生,.
有些牧者就忠心一點,.
他就說,.
我都講的,.
這些聖經很重要,.
段段都很重要,.
但是他講出來的時候,.
就將這段聖經演繹成為,.
這段聖經是為了維持人際關係的核心內容,.
但是當講到維持人際之間的核心內容的時候,.
就惹起更多的紛爭,.
尤其是劈頭第一句十八節講,.
作為妻子當馴服丈夫,.
馴服這個字眼我們一會再詳細解釋,.
不過簡而言之馴服這個字眼,.
就是完全跟從,.
就是你丈夫叫你做什麼,.
你就要做什麼,.
妻子就頂不住了,.
我們這裡很多妻子,.
如果你丈夫叫你做什麼,.
你就做什麼,.
你就頂不住了,.
你以為我是奴隸嗎?.
我嫁了你,.
我整個人終生都給了你了,.
你現在這樣來叫我還要聽你,.
什麼都聽你講,.
姐妹很難接受了,.

$^{201}$不用講了,.
接著下一句,.
你們作丈夫的,.
愛你的妻子,.
不可苦耐她們,.
這句比較中聽,.
弟兄,.
我愛你,.
所以跟著聖經,.
你要聽我說,.
聽我說我就愛你,.
這些全部將聖經內容扭曲了,.
沒有前文後理,.
沒有背景,.
然後騰空拿出來講,.
於是有些牧者知道這樣講不行,.
在教會會引起公憤,.
姐妹又不妥我,.
弟兄好像很拒我,.
結局怎樣?.
牧師是男性,.
所以我要安納女牧師,.
教會就開始摧毀了,.
大家看不看到,.
你不懂解聖經的時候,.
不單止你在經文上誤會上,.
最重要的是將淑玲的家摧毀,.
既然是這樣,.
有些精明的牧者這樣講,.
這段聖經其實是講昔日的舊世界,.
當時的世界裡面,.
男作為家女的主,.
女的做輔助,.
於是有這樣的介紹,.
用這段聖經嘗試說順服,.
不要太緊張,.
我們不是這樣用的,.
弟兄姊妹,.
我剛剛以上講了這一大堆東西,.
都是教會對這段聖經的誤解,.

$^{241}$尤其是對家規,.
Household Codes,.
或者Household Tables,.
一個不了解的詮釋,.
所造成的種種問題,.
如果你說這段聖經只是為了以前的信徒寫的,.
這段聖經也不需要講了,.
今天不合用了,.
講一句大家順服,.
彼此順服就算了,.
不要講內容了,.
不用講了,.
散會,.
是不是這樣呢?.
因此,.
今天我想透過這段聖經,.
一方面幫助大家去準確了解,.
這一類型的經文,.
如何去詮釋聖經裡面的所謂屬靈家規,.
那個主旨是什麼?.
第一樣要搞清楚,.
第二,.
它詳細的內容的介紹,.
是想指什麼來講,.
免得大家在這裡對聖經誤解的同時,.
對上帝席著保羅所寫下來的這一連串的介紹,.
讓我們失去了對上帝說話的主權的尊重..
我們將這段聖經分為兩個大部分去看,.
第一個的部分,.
就是從第三章十八節至到第二十一節,.
在這段聖經裡面,.
談論到的第一樣東西,.
是從屬基督的家庭關係運作,.
我們再看看這段聖經說什麼,.
留心聽,.
我讀一次,.
十八節至二十一節,.
你們作妻子的,.
當信服自己的丈夫,.
這在主裡面是相宜的,.

$^{281}$你們作丈夫的,.
要愛你們的妻子,.
不可輔代他們,.
你們作兒女的,.
要凡事聽從父母,.
因為這是主所喜悅的,.
你們作父親的,.
不要惹兒女的氣,.
恐怕他們失了志氣..
當我們讀完一次之後,.
我們會清楚看見整段的聖經,.
其實是談論到一個家庭,.
是有基督在他們當中的,.
請大家留意,.
多次談論到,.
妻子,.
丈夫,.
兒女,.
父親等等這些關係,.
都談論到在主裡面,.
有些因素在裡面,.
簡單來說,.
從這段聖經當中,.
主的角色是非常重要,.
作為一個屬靈的家庭,.
一切作主,.
不是誰做家裡的主人,.
而是我們的主,.
耶穌基督,.
或者用當時保祿的時代,.
就是上帝,.
祂的啟示就是一切的衡量標準,.
以及祂的最終的目的,.
祂將它展示在這段聖經裡面,.
讓我們去明白到,.
原來你家裡面和妻子關係,.
和丈夫關係,.
和兒女關係,.
作為父親的關係,.
你怎樣在裡面活出有主的心思,.

$^{321}$這個就是家規裡面最重要的主旨,.
不是維繫和諧,.
而是以達成主的啟示的標準,.
和祂的目標,.
我們看看第一個部分,.
第十八節,.
所談論到有關於信服丈夫的問題,.
「你們作妻子的,.
當信服自己的丈夫,.
這在主裡面是相宜的.」.
什麼叫做信服呢?.
在希臘文的版本裡面,.
信服這個字眼,.
是利用了一種叫做.
imperative immutable voice的文化寫出來,.
為什麼我要將它寫在大家的程序表裡面,.
特別強調是希臘文的文法,.
是以imperative immutable voice來說,.
第一,imperative命令式的說話詞體,.
命令式的字眼,.
我們很多時候就以為是一種.
不可以逆轉的理在上,.
壓下來給我們有必須進行的做法,.
這個解釋對於imperative命令式的語句,.
是一種很大的誤解,.
在希臘文裡面,.
這個命令式imperative,.
除了是指導一個人必須進行這個含義之外,.
他也有另一重的意思,.
就是說他給一個很鐵定的指導方向,.
你可以選擇的,.
有一句廣東話叫做意見接受,.
跟著是什麼,.
態度照舊,.
imperative是可以意見就給你了,.
你可以態度照舊的,.
不過,希臘文精彩之處,.
就是他加了一個叫做middle voice,.
在希臘文裡面,.
和很多的語文裡面都未必有,.

$^{361}$尤其是中文沒有的,.
中文講voice這個東西,.
叫做語調,.
語調要用長篇大論去討論的,.
有時沒錯,包含了一個字眼的裡面,.
譬如默想,.
默想這個字就是想,.
靜靜地自己一個想,.
你聽清楚,.
為什麼多了自己,.
中文這個自己,.
默想這個自己的行為,.
是含在字眼裡面,.
但很多時候很多字眼是沒有的,.
中文很難去表達voice,.
一個叫做語調的問題,.
希臘文就不同了,.
希臘文一定要有個voice,.
寫到明在字眼裡面,.
而精彩之處這裡所用到的就是,.
除了是命令式的文法詞彙之外,.
他再加上用middle voice,.
什麼叫middle voice?.
middle,.
就是中間,.
除了middle voice有什麼?.
有主動的,有active的,.
有被動的,被人催動的,.
middle就是很簡單的,.
自己,.
自己裡面,.
原來保羅在這裡說,.
不是給一個命令的妻子去馴服丈夫,.
而是談論到,.
妻子是要甘心自願地去馴服,.
她的丈夫給她的指示,.
整段聖經的解釋就很不同了,.
當我們去看回和合本,.
和合本最引人爭議的就是,.
你們作妻子職,.

$^{401}$當馴服自己的丈夫,.
這個當字好像是沒得選擇,.
但是其實當這個陳述,.
我相信翻譯的時候那個已經很盡力了,.
其實當這個當字,.
如果我們將它取締不要的,.
就是說你們作妻子的,.
甘心樂意馴服自己的丈夫,.
是陳述到作為妻子的,.
她是心甘情願地去馴服,.
聽自己的丈夫的指示,.
但是丈夫的指示不是她個人的意思,.
請留意這句話,.
這在主裡面是相宜的,.
所以和合本的翻譯引發很多人討厭讀這段聖經,.
我不喜歡當這個字,.
好像強迫我要做,.
不是,.
這裡強調的是你自己甘心願意去聽丈夫的話,.
為什麼你要甘心情願去聽你丈夫的話,.
因為你的丈夫是在主裡面發出這個指示給你們,.
如果你的丈夫給你的指示不是主裡面合宜的事情的話,.
你應該和他說,.
這些不是上主的心思的話,.
那你就不是不馴服他,.
你是在愛護你的丈夫,.
幫助他能夠使你和他在關係上,.
你是甘心願意聽你丈夫發出合主心意的話,.
而這個做法是上帝喜悅的,.
是合宜的,.
所以和本所謂合宜這個陳述,.
又是相宜,又是令人有一種無伶兩可的思想,.
準確一些的翻譯,.
其實這句話,.
這在主裡面是相宜的,.
好一些的翻譯就是,.
做到一個程度是有所銜接,.
又或者在這件事上做的時候,.
有恰當的性質,.
妻子去甘心聽丈夫的話,.

$^{441}$是因為丈夫所發出的東西,.
是連上帝都覺得恰當,.
當你肯聽他的話,.
這個就是恰當的square,.
即是恰當又恰當,.
在恰當的命令之下,.
或者指示之下,.
或者意願之下,.
你肯恰當甘心去委身在裡面,.
所以這個不是一個人可以做到的事,.
一定有個丈夫,有個妻子,.
所以為什麼在希臘文裡面,.
丈夫和妻子這兩個字眼都是用眾數,.
你不信自己查,.
找blue letter bible,.
我很歡迎的,.
其實我們教的弟兄姊妹好,.
包括網上會有,.
都是很好,.
真的會去核證,.
你不是純粹聽牧師講的東西,.
都是屬於二手研究,.
雖然我這些全部都是一手研究,.
不過我講給你聽已經是二手,.
你再去印證一下,.
都是用眾數,.
什麼意思?.
這個是普世應有的,.
屬靈家庭裡面,.
一個家庭,丈夫和妻子的關係的描述,.
而耶穌成為一個很重要的因素,.
所以如果大家還記得,.
在第一章第十五至二十節裡面,.
我們去強調談論到,.
基督是一切關係的永恆,.
超級屬靈的接著劑,.
有基督在家庭裡面的時候,.
夫妻關係就會建立得越來越緊密,.
但前提是,.
大家都是以主基督,.

$^{481}$作為合宜的事情來決定,.
聖經就不是單方面的,.
我剛才講,.
所以十九節保羅,.
flip this coin to another side,.
將銀反過來再看,.
剛剛看完字那邊,.
我現在反過來看宮那邊是怎樣說的,.
宮是甚麼?宮就是丈夫,.
十九節說,你們作丈夫的,.
要愛你們的妻子,.
不可輔待他們,.
首先這裡直接是指到丈夫,.
丈夫你去愛你的妻子,.
這個字眼也是用了命令式的middle voice,.
imperative with middle voice,.
來寫出這個動詞,.
丈夫去愛你的妻子,.
這個是一個強烈的呼籲,.
但這個呼籲是你甘心樂意,.
自己願意的,.
不是被迫的,.
丈夫怎樣愛你的妻子,.
是不是他要包包,.
你多買兩個包包給他,.
問問在座那些弟兄,.
做丈夫那些,.
你怎樣愛你的妻子,.
他要甚麼你給他甚麼,.
我要一輛SUV,.
我要馬上買房,.
你就愛我了,.
是不是這樣看的,.
妻子說,.
牧師你今天說話對聽,.
我喜歡吃壽司,.
丈夫馬上買壽司,.
好像昨天那樣,.
橫風橫雨都要冒著去買,.
那你就愛我了,.

$^{521}$是不是這樣看,.
看看聖經怎樣說,.
詞文已經講了在主裡面的東西,.
這些不用講了,.
但接著說,.
不可輔待他們,.
甚麼叫輔待他們,.
原來你愛你的妻子,.
不是他喜歡包包,.
你買個包包,.
Gucci就Gucci,.
LV就LV,.
不是這樣,.
現在是否流行這隻牌子,.
好像不是,.
不知現在流行甚麼,.
不是你喜歡飲飲食食,.
你滿足他就等同於愛他,.
聖經都講,.
你不要輔待他,.
輔待這個字眼,.
在希臘文原文的意思,.
即是去毒害他,.
嘩,.
如果聽到這裡,.
即是說丈夫不整死你,.
就已經是害你了,.
你們在笑,.
這裡想講的是甚麼,.
不要在他的肉體和精神上,.
去輔待他的意思,.
這是在輔待這個字眼,.
所涵蓋的含義,.
聽清楚,.
不要在肉體和精神上,.
去輔害你,.
在這裡讓我插句口說,.
我記得我在結婚之前,.
我們會去做婚前輔導,.
其中有一個婚前輔導的導師,.

$^{561}$他講了一句說話,.
我很深刻,.
他說夫妻之間,.
是包括肉體和思想上的交流,.
夫婦之間,.
切勿以肉體和思想上,.
去要脅你的配偶,.
跟著在婚前輔導裡面,.
有一個這樣的說法,.
如果你用性的關係,.
來要脅你的丈夫,.
或者要脅你的太太的話,.
你的婚姻一定不好過,.
因為這個傷害是不單止肉體的問題,.
而且是包括思想和靈性的問題,.
鄧小平,.
我相信這個導師,.
婚前輔導導師,.
因為他是基督教機構,.
那時候我們去了一個機構,.
那時候機構還好的,.
現在機構就不好了,.
他講的,.
我相信就是在這段聖經裡面出的,.
當然了,.
因為他的輔導的session,.
不只是基督徒才去,.
很多人都去,.
不同的信仰的人都會去,.
不信的都去,.
因為那些是給錢的,.
那些叫婚前輔導,.
但的而且確,.
作為丈夫的,.
你切勿在思想和肉體的上面去毒害你的太太,.
這個就為之愛你的妻子,.
不是提供物質叫做愛,.
而是你關心到她整全一個人,.
從這段聖經裡面,.
就延伸到作兒女之間的關係,.

$^{601}$二十節,.
他說你們作兒女的,.
要凡事聽從父母,.
因為這是主所喜悅的,.
首先這裡指兒女,.
就不只是指你自己的兒女,.
要聽你的說話,.
因為大家不要忘記,.
當你有兒女的時候,.
其實你也在做人兒女,.
如果你想你的兒女,.
懂得怎樣做兒女的話,.
就要看你怎樣對你父母,.
怎樣去照顧或者去信服你的父母,.
聽清楚,.
凡事聽從父母,.
這個字眼之後會再用,.
在第二十二節,.
不過用在第二個人的關係上,.
就是僕人和主人的關係,.
我一會兒再說,.
凡事聽從,.
所有的事都要聽從,.
今天在這個充滿網上的社交媒體,.
所傳遞的訊息就是,.
你是你個人的意見,.
你是你個人的主宰,.
你根本不需要你們的父母,.
所以現在說家庭解體就是這個原因,.
你越用的社交媒體多,.
你就長期和人的關係斷絕,.
你越和網上社交媒體結交,.
你就越和真正和人的關係破裂,.
你看看你身邊的人,.
他天天看Facebook,.
你看那些社交媒體,.
越看得多那些人,.
你看看他和人之間的關係是怎樣,.
聽人說話都有問題,.
你和他說話,.

$^{641}$他一定誤會你的,.
誤會你的機會九成九,.
他根本以自己為主宰,.
我不喜歡這個人,.
我就反過來了,.
每一條片,.
每一幅圖,.
他只看幾秒鐘,.
他沒有深思考慮過,.
這個聖經說什麼,.
作兒女凡事聽從父母,.
是因為這件事主所喜悅的,.
聽清楚,.
如果父母給你的指示,.
如果不是主所喜悅的,.
主怎會喜悅你去信服他的父母,.
因此你的父母作為父母的,.
他給你的指示是主所喜悅的,.
而你肯去聽的話,.
主就會喜悅這個行為,.
喜悅的而且是你的心思,.
所以同樣這個字說到凡事聽從,.
都強調要說一件事,.
是主動地去做這件事,.
因為你看到你的父母,.
是從主所喜悅的意思而發出指示,.
所以這裡的字眼,.
同樣用了命令式的語句卻轉了語氣,.
不是middle voice,.
而是active voice,.
是你主動地去欣賞到父母,.
給你那個是主所喜悅的心思,.
而你主動地去跟從,.
因此來到第21節,.
你們作父親的,不要惹兒女的氣,.
恐怕他們失了志氣,.
在這裡很強調,.
說到作為父母,.
你要主動地不要惹兒女的氣,.
這句說話有些人就這樣傳譯,.

$^{681}$不要惹即是不要惹怒子女,.
子女要什麼,.
不要騙他,.
他喜歡什麼就給他,.
如果是這樣就變成縱容了,.
這段聖經想說的是什麼?.
不要惹兒女的氣這個陳述,.
是指導我們不要使子女去反感,.
什麼反感?.
是你給他們的訊息,.
引起他們反感,.
當他們反感的時候會怎樣?.
會有可能,.
所以恐怕不一定出現,.
恐怕他們失了志氣,.
他們會失去信心,.
或者會失去什麼?.
使自己覺得沒用,.
簡單來說,.
父母的角色是作為善導兒女,.
使他們成熟地起來,.
曉得作成熟的思考和決定,.
既然是這樣的,.
作為父母的,.
尤其是作為父親的,.
因為以一個父系的世界來介紹,.
包括父母親的,.
作為父親的,.
你的角色是鼓勵他們做他們能夠做到的事,.
幫助他們去善導他們,.
去成就他們的世界,.
他們所作的事情,.
當然你作為父親的,.
你給他們什麼意見,.
就是前面所說的,.
主所喜悅的事,.
整段聖經說的包括一個很重要的核心元素,.
就是雙方彼此,.
英文叫Reciprocal Relationship,.
作為一個屬靈的家庭,.

$^{721}$有基督在其間的,.
他最重要的角色就是維繫著所有,.
無論作妻子,丈夫,兒女,.
以及父親的,.
你都要以基督的標準,.
以及他的喜悅來管理,.
維繫,並且連結起來,.
不是人的標準,.
所以為什麼有些人對這段聖經很反感呢?.
就是因為他們覺得,.
丈夫給我的命令是他的個人意思,.
妻子不喜歡聽丈夫的話,.
因為她自己喜歡怎樣就怎樣,.
兒女是怎樣,.
因為他個人喜歡做些什麼,.
於是父母就不要管了,.
由他,縱容他,.
這種情況就搞到關係在家庭裡面,.
特別一個基督徒的家庭裡面,.
關係破裂,.
保羅在這裡說這一番連串的說話,.
是要讓哥羅西教會的信徒明白,.
一個大都會的哥羅西城,.
你們作為基督徒,.
你們怎樣活在一個這樣的世界當中,.
成為見證上主,.
怎樣建立一個家庭的真正的實體,.
給這個城裡面的人看到,.
這個做法,.
就是我們已經講了兩堂,.
怎樣榮耀神,.
你要榮耀神,.
不是只是說的,.
我已經講過很多次了,.
不只是在嘴裡說,.
而是實實在在你的生活裡面,.
家庭裡面建立出來的,.
你因為主喜悅的緣故,.
所以你在家庭裡面,.
發出一些要求或者指示,.

$^{761}$而你的太太,.
你的兒女,.
是願意在這裡跟從的,.
包括你自己,.
也是要鼓勵他們成長,.
更加在主裡面成熟,.
這個才是一個屬靈家庭,.
基督徒家庭用今天說,.
能夠見證主的原因,.
不是你私人的意願,.
不是你個人的標準,.
來量度,.
所以當大家都能夠彼此在主裡面降伏的時候,.
就會出現一個很特別的現象,.
而在聖經就記載在其他經文裡面,.
第一段經文,.
馬太福音第十六章第25節,.
請大家一起打開好不好,.
我們看看,.
讀一讀好不好,.
馬太福音第十六章第25節,.
很多人不明白這句聖經怎麼解,.
以為只是用在耶穌基督身上,.
其實不是,.
我們一起讀,.
馬太福音十六章25節,.
因為凡要救自己生命的必喪掉生命,.
凡為我喪掉生命的必得著生命,.
讀多段,.
哥林多前書第十章24節,.
我們由23節開始讀,.
找到位一起讀,.
凡事都可行,但不都有益處,.
凡事都可行,但不都造就人,.
無論何人不要求自己的益處,.
乃要求別人的益處,.
這段聖經說什麼?.
讀多句,讀第33節,.
33節說什麼?.
我們一起讀,.

$^{801}$就好像我凡事都叫眾人喜歡,.
不求自己的益處,.
只求眾人的益處,.
叫他們得救,.
弟兄姊妹,當我們讀完這三節經文後,.
你應該得出一個結論,.
當如果妻子,丈夫,兒女,.
所有的人都能夠在基督合他心意,.
認為相宜的狀態下去做的時候,.
沒有個人意見,.
只是以達到上主的心思的時候,.
你做出的表現就是為了你的家人的靈性好處而著想,.
當你願意放下自己的意見,.
去成就你的家人能夠在主裡更成熟的時候,.
你這個願意的付出,.
好像沒有了自我,.
這個和世界的觀念完全相反,.
這個世界什麼都是我自我,.
我自己成功了就好了,.
我用我成功去帶協其他人,.
這個是世界的看法,.
聖經的看法是我不重要,.
我最重要幫助到你能夠成就你的靈性好處,.
所以我為你的好處,.
我願意放下自我,.
簡單地說,.
好像犧牲了自己的喜好,.
滿足了你在對方的屬靈上的需要,.
譬如我作為丈夫,.
我為了我太太的靈性的好處,.
我願意放下我自己的好處來為她付出,.
我為了我子女的好處,.
所以我願意放下我自己的好處,.
來滿足他們在屬靈上的滿足,.
當我這樣做的時候,.
他們也看到這一點,.
願意甘心回應的話,.
太太甘心願意,.
兒子兒女願意主動來聽你的說話,.
他們也可以用放下他們自己的好處,.

$^{841}$來滿足作為丈夫的,作為父親的,.
他的最大好處,.
在靈性上的好處,.
這幅圖畫,.
就是藉著犧牲來成就整個團體的滿足的豐富,.
這個是這個世界的人母的思維,.
是逆這個世界的思維,.
但是作為信徒的,.
這個就是聖經的教導,.
你家庭的成長,.
在這個世界裡面,.
就是每個人出去挖到多少東西回來,.
賺到多少東西回來,.
我家庭就會好了,.
其實實質上不會的,.
你看很多家庭,.
子女賺到錢之後,.
怎麼會回來請爸爸吃一頓飯,.
有沒有,.
你請他吃飯就有,.
是不是,.
他什麼時候請你吃一頓飯呢,.
如果他肯呢,.
這樣就有得救了,.
我不敢說什麼,.
我上個星期,.
和我女兒,.
去弄一些東西,.
弄完之後,.
我就說要去,.
我家裡發現有老鼠,.
我就手遍整個多倫多的馬琴區,.
嘗試找老鼠藥這東西,.
聽聞用老鼠藥是違法的,.
其實我也知道的,.
不過儘管一試看看有沒有老鼠藥,.
吃了不會違法,.
去到哪裡都說沒有,.
買了那些老鼠藥,.
吃了會口渴,.

$^{881}$老鼠藥喝水會死,.
沒有即食即死那種,.
那我沒辦法了,.
買了老鼠膠回來黏,.
黏完之後,.
心裡面不是很安樂,.
為什麼呢?.
因為我家裡有老鼠,.
我就無所謂,.
我與鼠同活了幾十年了,.
我家裡以前的睡房,.
晚上睡的老鼠在那裡走,.
那些老鼠不是這裡的老鼠,.
這裡的老鼠叮那麼大隻,.
很大隻的,貓那麼大隻,.
心裡不是很開心,.
周到我去巡一巡店鋪,.
我女兒看到我,.
爸爸你真的很囂張,.
她去巡一巡,.
她去巡護膚品的部門,.
她看到護膚品的抹面濕紙巾,.
她說爸爸我買包給你,.
我買包濕紙巾幹什麼,.
護膚嘛,.
我女兒很體貼我的心想,.
買不到老鼠藥回去,.
解決不了老鼠的問題,.
我就會很頭痛,.
因為師母頭痛就是我頭痛,.
有兩隻老鼠黏在家門口,.
現在等停雨去處理,.
當你願意為對方付出的時候,.
對方看到你也願意付出,.
這樣的淑齡家庭就是保羅所期望,.
是透過付出,.
不是為了自己立多少東西而建立的一個家庭,.
你作為基督徒,.
是否同樣用這種,.
好似捍衛損失,.

$^{921}$而去成就你家庭裡合主心意建立的基礎呢?.
看看耶穌基督,.
祂何嘗不是這樣,.
看看父神,.
何嘗不是這樣,.
祂願意犧牲兒子來成就父神的喜悅,.
父神的喜悅就是拯救那些願意悔改的人,.
父神願不願意付出損失他的兒子,.
祂願意,.
耶穌肯不肯付出?.
又肯,.
兩父子彼此的付出就成就了舊恩,.
你和我就有機會歸回上帝,.
你看看,.
犧牲不是損失,.
犧牲是贊取了合神心意的事情得以成就,.
我再強調多一次,.
犧牲如果為個人某些人的目的就不對,.
但如果為上帝的心思,.
俠義的,.
上帝喜悅的,.
你的犧牲就有意義,.
不然的話,.
你的犧牲可能會變成一種縱容,.
保羅在這裡給了一個這樣的門路,.
給哥羅西教會知道怎樣榮耀神,.
第二個部分,.
第三章二十二節至第四章第一節,.
請聽我讀出,.
你們作僕人的,.
要凡事聽從你們肉身的主人,.
不要只在眼前事奉,.
像是討人喜悅的,.
總要全心誠實敬畏主,.
無論做什麼都要從心裡做,.
像是給主做的,.
不是給人做的,.
因你們知道從主那裡,.
必得著基業為賞賜,.
你們所事奉的乃是主基督,.

$^{961}$那行不義的,.
必受不義的報應,.
主並不偏待人,.
你們作主人的,.
要公公平平地待僕人,.
因為知道你們也有一位主在天上,.
這段聖經很多時候被用作,.
為如何在今天的社會裡面打工,.
和你的上司和你老闆之間的關係,.
是不是想說這些呢?.
其實這段聖經想說的是,.
一個從屬基督的賓主關係的運作,.
先前說到一個家庭裡面,.
妻子,.
丈夫,.
兒女,.
父親之間的運作,.
現在轉了進去,.
說一個僕人和主人之間,.
在屬靈的關係上,.
如何運作,.
前文談論到,.
原來每個屬神的子民,.
都是神的選民,.
既然上帝的選民,.
不只是規限在猶太人身上,.
在哥羅西這個社群當中,.
包括很多非猶太人的僕人,.
或者主人,.
會在這個環境裡面出現,.
但是在這裡,.
保羅籠統地去將一個這樣的,.
勸喻說給,.
作為僕人的去聽,.
他說第一樣,.
凡是聽從你肉身的主人,.
同樣在這裡用同一個的陳述,.
命令式的主動語,.
作為僕人,.
是主動地去服從,.

$^{1001}$你那個主人所說的一切的命令,.
這樣打工不就是做奴隸,.
不過請留意,.
留意,.
要凡是聽從你肉身的主人,.
不要只在眼前事奉,.
22節用了事奉這個描述來介紹,.
我要強調,.
這個所謂事奉這個字眼,.
是一個詮釋者,.
或者翻譯聖經的時候,.
這個翻譯聖經的人去到理解,.
這個僕人不是一般的僕人,.
他視乎他的主人好像事奉神那樣看,.
我可以強調一點,.
你可以將事奉這個字眼剔除,.
他強調的是你在他眼前,.
夠了,.
你在主人不在場的時候,.
主人在場的時候,.
你才去聽他的說話,.
你是在主人不在的時候,.
都同樣去聽他的說話,.
不要說是否事奉,.
是否信徒先,.
這種態度不是因為某一個人,.
你才有什麼回應,.
你聽那個人,.
聽那個指示,.
不是因為那個人管著你,.
你才去聽,.
而是因為你對那份工作的忠誠,.
而去做,.
不要說什麼,.
說我自己,.
我做傳道人做了多少年,.
我忘記了,.
我1996年開始做傳道,.
現在多少年了,.
30年了,.

$^{1041}$很驚人,.
這30年裡面,.
教會幾年有辦公室,.
之後沒有辦公室,.
我長期在家裡工作,.
在家裡工作,.
今天你是否work from home,.
你work from home的時候怎麼穿,.
穿睡衣?.
還有work from home的時候,.
沒有人在家裡,.
師母上班,.
我的子女上學,.
我只有一個人,.
有沒有大兄姐每天打電話來問,.
牧師你在做什麼,.
我聽說有間教會更差,.
有辦公室的,.
還要在辦公室的牧者辦公室上面,.
裝一個閉路電視,.
長期監察著牧師傳道人在做什麼,.
有沒有?沒有,.
你們有沒有試過我不講道,.
有沒有試過我的講道打折扣,.
隨便信口說的,.
我寫一堂道要接近20個小時,.
大兄姐們,.
是什麼驅使我這樣做,.
因為我要忠誠於我這份事奉的工作,.
上帝委託我,.
忠誠於上帝的委託,.
所以我去做,.
上帝監察著你,.
上帝監察著我,.
怎麼監察,.
怎麼看,.
你不懂的了牧師,.
我不懂,.
當我不懂,.
那你懂,.

$^{1081}$那你告訴我,.
上帝怎麼看著我,.
不是看著你才去做,.
你做的那些工作,.
是不是有人看著你才做得好的,.
如果不看著你,.
你不做得好的,.
保羅在這裡講道,.
作為僕人去聽主人講的話,.
不是因為主人看著你才去做,.
討人喜悅這個意思,.
你想領功在主人上面,.
你看我多好,.
或者其他的公司同事看著我,.
你看我多聽話,.
我的老闆要加薪給我,.
你們沒得吵了,.
不是這樣,.
是因為你作為僕人專注做你這份工作,.
大家請要明白古晉東文化裡面,.
僕人奴隸,.
應該原文這個字眼,.
奴隸可以有兩種,.
一種是被賣的,.
沒得說自己選擇離開的,.
另一種就是自願甘心毀身在主人裡面,.
即使他自由可以選擇不同的人可以做他主人,.
甚至可以選擇不做勞部,.
這兩種無論是哪一種也好,.
被迫的也好,.
沒得選擇也好,.
或者有得選擇的也好,.
因為你做得那份工作你就專心去做,.
有些人說牧師你又做核數師,.
不要講得這麼高級,.
核數員,你做核數的很大權,.
我可以核數我要你罰你錢也行,.
我不罰你錢也行,.
全部我自己決定,.
有些人說算了做少些不用那麼煩,.

$^{1121}$不是的,.
你做得那份工作就盡忠做好那份工作,.
基督徒就是這樣的性情,.
你做得音響在教會就搞定音響,.
你當日搞不定就拿回家搞定,.
就算數隨便,.
就算我寫得唐道,.
我可以寫多些寫短些,.
但是你寫少了沒了主題,.
變了對上帝不住,.
你有沒有這樣想過?.
所以保羅這裡說,.
總要全心誠實敬畏主,.
你是誠實地去面對主,.
你做這一切工作,.
你願意遵行做好你那份職責,.
是因為你無時無刻都是願意誠實對著我們的上主,.
因此我們作為神的選民,.
不是說有時見到有事做就做好些,.
牧師不在就隨便些,.
或者在你的工作裡老闆不在,.
今天放假我就放任些,.
做少些,.
隨便些,.
不是的,.
仍然做足你要做的事,.
是隨時隨地每一刻,.
你的心思都是盡忠於你所得到的委託,.
在這種情況下,.
保羅不只是正說,.
他有行動告訴23節,.
無論做什麼都要從心裡做,.
心裡做的意思不是問心無愧那麼簡單,.
而是盡力盡心地去做,.
好像你為主所做的一樣,.
這個主有人可以將他翻譯成為你的主人,.
但我覺得更準確的翻譯就是向著你的上主,.
你的神,.
你去介紹,.
努力去做好,.

$^{1161}$那份工作,.
你不是只是純粹做了過到人就算數,.
你是給上主去做,.
所以為什麼?.
像是給主做的不是給人做的,.
經文就規限了我們這裡的理解,.
是向著上帝去做,.
向基督而去做,.
所以他解釋24節,.
因你們知道從主那裡必得著基業為賞賜,.
你們所事奉的乃是主基督..
很清楚,.
不過在這裡我想多講一點,.
他說因為你們知道從主那裡必得著基業為賞賜,.
基業這個字眼講到inheritance,.
產業,.
或者講得更直接為產,.
講到這個字眼往往有兩重意思一定包含在裡面,.
第一一定包括有死亡,.
什麼時候會領遺產,.
有人死了就有人領遺產給你,.
第二就是他用了未來式去講這番,.
領受基業,.
從主那裡必得著基業為賞賜,.
這裡所講的是一個未來式,.
講到人原來會有終結的一日,.
或者有死亡會參與介入其中,.
但是在裡面有基業作為賞賜,.
如果當我們明白到這一點的時候,.
似乎好像遙不可及的事,.
將來的事將來的事算了,.
所以努力事奉好我那份工作,.
稱職做好成果將來的事,.
不是的,.
所以第一件事講到身份問題,.
誰才可以領受基業,.
死了為什麼要給你基業,.
為什麼要分家產給你,.
因為你是神的兒女,.
你是選民一份子,.

$^{1201}$你有份分享這個國家裡面的好處,.
第二不是將來的事,.
是這一刻你已經可以得到,.
在彼得後書第三章八至九節裡面有句經文講,.
親愛的弟兄,.
有一件事你們不可忘記,.
就是主看一日如千年,.
千年如一日,.
主所應許的尚未成就,.
有人以為他是擔言,.
其實不是擔言,.
乃是寬容你們,.
不願有一人沉淪,.
乃願人人都悔改,.
領子們,.
上主叫我們做好僕人的工作,.
目的其實給我們可以分享這個基業,.
而這個基業不是將來的事,.
現在我們已經有份參與在其中,.
其中一樣東西,.
就是我們可以參與在這個悔改裡面..
從這段聖經裡面給我們看見,.
保羅作為僕人的要求,.
似乎是為著僕人的最大益處著想,.
僕人要付出,.
主人不用付出嗎?.
主人大哂嗎?.
不是,.
我們再看第25節,.
25節我們一起讀,.
那行不義的必受不義的報應,.
主並不偏待人,.
你們作主人的,.
要公公平平地待僕人,.
因為知道你們也有一位主在天..
為什麼我到這裡會特別讀到這一節呢?.
因為第25節的這一節,.
是將僕人種種的行為,.
連帶到主人之間的關係,.
什麼將僕人和主人連起來呢?.

$^{1241}$是二,.
僕人去對主人做事,.
因為願意委身,.
願意投入,.
縱使主人沒有看著你,.
沒有人看著你,.
你也願意忠心的,.
因為你已經是直著悔改,.
進入到二的身份裡面,.
所以在這裡第25節說,.
那行不義的必受不義的報應,.
這句話是對僕人和主人說的話,.
所以馬上跳進去,.
主不偏待人,.
這個作為主人你就要醒目了,.
怎樣醒目?.
不要偏待,.
怎樣不偏待?.
你們作主人的,.
要公公平平地待僕人,.
因為二你們要知道,.
天上有一位主,.
而將僕人和主人混在一起,.
而這個關係原來是反映那些被上帝招載他其中的人,.
和這位上主的關係,.
換言之我們今天在這個地上,.
作為打工一族,.
你和你的主人或者你的老闆,.
上司的關係,.
你應該做到的,.
是好像我們和上主之間,.
那個二的關係的建立,.
不是利益的建立,.
講到這裡有些弟兄姊妹就,.
牧師你真是很清高,.
講大道理,.
打工當然求財,.
講心,.
講真,.
真是要講心,.

$^{1281}$你去做一份工作,.
如果你和你的上司,.
他日日都是以不義來對待你的,.
你等他炒你,.
很慘的,.
我有次和湯全道聊天,.
湯全道講了一句說話,.
他一生人被人炒過一次魷魚,.
這件事我前日才和他談,.
昨日星期六我見到他,.
昨日我在他家探他時,.
他跟我說這番說話,.
他說他一生人被人炒過一次,.
你知不知道被人炒那種痛苦是一生的,.
因為你一生被人覺得你是無能,.
你辭工反而好些,.
你辭工就是我不適合做,.
老闆沒有hard feeling,.
你不用擔心,.
如果他純粹是用錢去交代,.
不要說不講心,.
弟兄姊妹,.
我們作為信徒會在這個地上,.
當我們見到這些種種關係,.
我們有無要透過這些關係去反映出.
我們和神之間的關係呢?.
如果你的上司,.
你的老闆,.
你的同事全部都當你像奴隸般踐踏,.
奚落你,.
你其實給一個機會他們在犯罪,.
最慘的傷害是對於你永恆,.
有些智慧去處理,.
從這段聖經給我們看見,.
整段經文,.
從第十八節至第四章第一節,.
不是講一個歷史文化傳統,.
怎樣在裡面得以維繫,.
不是的,.
這段經文是一段超越時代歷史文化空間的一段經文,.

$^{1321}$是談論到人與人之間的心思,.
怎樣在其中去反映出我們和神之間的關係,.
為什麼要講這個題目的保羅?.
因為這個正正就是榮耀神的重要而可行的實踐方案,.
所以每一個人,.
他的人與人關係出問題,.
永遠都是那個人,.
因為他們與上帝的關係有問題,.
那個人與上帝的關係越疏離,.
他與人的關係越疏離,.
登山寶訓馬太福音第七章第十二節這樣說,.
無論何時你們願意怎樣待你們,.
你們也要怎樣待人,.
因為這是律法和先知的道理,.
在這段聖經裡面講到,.
你想人家怎樣待你,.
最起碼你要首先明白到,.
你怎樣去待人,.
而這個道理怎樣去待人,.
就是在律法和先知的道理,.
即是上帝的啟示,.
當人與神之間的關係調整的時候,.
就會與人與人之間的關係開始調和,.
而這個過程不是靠獲取多少,.
而是看你怎樣按照使人成熟的需要去付出,.
我作為父親,.
我不介意對子女有付出,.
但我肯定一件事,.
我所付出的東西是對於她在靈性上有栽培,.
如果我付出了是縱容她的,.
我不會做,求主幫助我們,.
我們是一間教會,.
是一個群體,.
不是天主教那些要靠修煉修道修出來,.
而是透過與弟兄姊妹彼此聯絡交通而塑造出來,.
所以去退修會,.
我們有幾日走在一起生活下,.
你就會看到原來弟兄姊妹是這樣的,.
然後在過程當中,.
我們如何彼此透過付出建立大家的屬靈生命,.

$^{1361}$求主幫助我們,.
讓我們一起同心低頭祈禱,.
天父多謝你讓我們在這段聖經中學習,.
也求主你的大靈,.
我們知道我們生命上有一個地方,.
不是單靠我們自己一人磨練出來,.
雖然以馬萊利我們的回眾有實體有虛擬,.
在網絡上與我們一起聚會,.
讓我們能夠繼續有彼此交通彼此聯絡,.
更加我們值得珍惜的是我們每一個弟兄姊妹,.
都能夠將我們心裡面真正所看見所面對的困難,.
都能夠願意彼此分享,.
好使我們能夠彼此去鼓勵,.
是一個合主心意的群體,.
我為了網上的每一個弟兄姊妹,.
為了我們實體崇拜的每一個弟兄姊妹代禱,.
沒錯,雖然我們在崇拜的裡面,.
時間可能比其他教會的牧者講得長,.
但也都比我們有各方的操練,.
讓我們能夠可以在主義的真道上面,.
能夠有更透徹明晰的理解,.
求主你的帶領,.
無論聽的講的,同蒙主你恩福,.
多謝你,我們同心禱告,.
奉靠主耶穌基督,得勝名字而求,.
阿們..
\newpage



\section{約書亞記 22:34}
\label{sec:032188}
\textbf{聖經人聖經事:證壇——證明耶和華是上帝}~~~~黃克勤牧師博士
\newline
\newline
連結: \href{http://www.church.com.hk/acms/content.asp?site=cdc\&op=show\&type=product\&code=032188\&layout=sermon}{\texttt{032188}} ~~~~ 語音日期: 2026-05-25
\newline
\newline
\hyperref[sec:032187]{\small{< < < PREV SERMON < < <}}
~
\hyperref[sec:index_chronic]{\small{[返順時目]}}
~
\hyperref[sec:032189]{\small{> > > NEXT SERMON > > >}}
\newline
\newline
約書亞記 22:34
\newline
\begin{longtable}{cl}
\hline
\hline
章節 & 經文 (和合本修訂版)\\
\hline
22:34 & \begin{tabularx}{0.7\textwidth}{X} 呂便人和迦得人給這壇起了名，因為這壇在我們之間見證耶和華是神。 \end{tabularx} \\ \\
[1ex]
\hline
\hline
\end{longtable}
$^{1}$大家好,又來到讀聖經的時間.
新衛書鼓勵讀聖經.
和大家一起研究聖經人物,聖經故事.
看看背後上帝做了些什麼.
今次講一個壇.
壇是用來獻祭的.
今次有一個壇是以色列人建造的.
但他們不是用來獻祭的.
而是用來證明,證據,證實一些事情.
這個壇是用來做證的.
所以叫做「正壇」.
大家第一次聽是不是呢?.
可能我們很少聽到.
也不是很多人.
很少傳道人或牧師傳道.
講這件事的.
今次和大家看看「正壇」.
「正壇」是證明耶和華是上帝.
記得座壇.
證明耶和華是上帝.
要用一座壇來證明耶和華是上帝的嗎?.
看看原聖經是怎麼說的.
以色列十二師派分了地.
你會知道其實.
我們很多時候都會想.
約旦河的西岸.
約旦河的東岸我們很少想的.
原因就是因為今天的以色列.
全部的土地都在約旦河西.
和地中海沿岸之間.
北就去到黎巴嫩.
和亞述.
以前的亞述.
今天叫做敘利亞.
而約旦河的東岸這邊.
今天就是約旦.
約旦這個國家.
因此我們經常說以色列.
的概念都是以約旦河的西邊.
但其實當以色列人.

$^{41}$由埃及經過曠野.
然後進入迦南.
在這裡窺探一下.
然後再出去.
然後繞到這裡進入耶利哥.
正式進入約旦河的時候.
經過約旦河.
就拿了這些地方.
有幾個支派.
馬來西亞,迦德和里面.
即是里面.
他們整個以色列.
走到這裡要進去的時候.
特別馬來西亞就說.
可不可以我們就用這個地方.
我們不進去迦南地.
不過河了.
為什麼?.
因為我們這裡見到是平地.
而且是很適合畜牧.
我們的牛羊又多.
人口又多.
不如我們留在這邊.
迦德和里面都可以一起.
所以以色列的十二支派分的地.
有三個支派.
是在約旦河的東邊.
都說了.
現在東邊我們不覺得是以色列.
因為今天東邊是約旦.
敘利亞和黎巴嫩.
以色列很小的.
大家留意到.
這裡去到死海.
北邊比較多城市.
南邊叫南地.
就是什麼的曠野.
南地是曠野.
大家知道支派名在這裡.
綠色這些.

$^{81}$非利士,阿瑪利,以東,摩訶,阿滿.
和上面的亞蘭.
全部都是外族人.
雖然摩訶和以東都是以色列的後裔.
以東是以素.
亞國的兄弟.
摩訶和阿滿是羅德的女兒.
兩個女兒的後人.
我們看看這裡.
當時約書亞召了裡邊人.
迦德人嗎哪西半支派的人來.
對他們說.
耶和華僕人摩西所吩咐你們的.
你們都遵守了.
我吩咐你們的你們都聽從了.
看看地圖.
因為約書是叫裡邊人.
迦德人和嗎哪西半支派.
即是兩個半支派.
為什麼呢?.
原因是進去已經平定了.
分地了.
分好了.
在這個時候.
其實嗎哪西黃色這一塊.
是這麼大的.
嗎哪西是最地最大的地方.
不過有約旦河在中間.
所以有約旦河一半的嗎哪西在這邊.
河西一半的嗎哪西在河東.
現在約書亞叫嗎哪西.
迦德裡邊支派的人來見他.
於是叫他們來.
跟他們說.
摩西吩咐你們的.
你們都遵守了.
我吩咐你們的你們都聽從了.
那麼久的日子.
你們總沒有撇離你們的弟兄.
用紫色的紙表示這些說話是重要的.

$^{121}$那麼多日子你們都沒有離棄你們的弟兄.
直到今天都守了耶和華你們神.
所吩咐你們當守的事守的誡命.
要你們做的事打的仗都做了.
你們那麼久都沒有撇離你們的弟兄.
說什麼成十二支派一起打仗.
哪有誰撇下誰.
其實是有背景的.
主要原因就是.
當嗎哪西迦德裡邊.
整個以色列走到這裡預備進入耶利哥的時候.
我剛才講過了.
他們說這裡的地很大很漂亮很平.
適合我們在這裡牧養生畜.
好不好我們的家就定在這邊.
我們不進去.
應許之地就在這裡埋便.
沒有計這裡出便.
不過我們想在這裡.
我們不尋求裡面那些.
我們想在外面.
摩西就說好.
既然你們這樣你們就在這邊.
分了這邊的地給你.
佔領了你就拿.
不過.
不要因為你們在這邊.
就不一起進去打仗.
你們都要放下家人和牛羊.
然後你們的士兵壯丁.
一起和其他支派進去迦南地打仗.
獲得這個地方.
然後才出來.
他們說好合理.
應該這樣.
於是他們就放下家人和所有東西.
三個支派的壯漢一起去打仗.
到現在就叫他們來.
以前叫你做的你都做了.
而且沒有撇離你的弟兄一起去打仗.

$^{161}$直到今天守了神叫你們守的誡命.
太好了.
所以如今你們上帝照著.
祂所應許的.
使你們弟兄得享平安.
現在你們可以轉回你們的帳篷.
回到你們的家.
你們的帳篷就是回到你們的家.
現在我回家了.
回家了.
我不會說我回我的帳篷.
因為他們住在帳篷.
你可以回到你們的家.
就是到去耶和華的那一天.
摩西約旦河東.
賜你們為業的地方.
你們可以回去了.
打完仗平定了.
多謝你們沒有離棄弟兄一起打.
雖然打的地方不是你們的.
你們都這麼賣力打.
非常之好.
可以回去了.
不過你們要切切謹慎地.
記住.
第二要說的是遵行耶和華的誡命.
第一.
你們都遵守了.
仆人摩西所吩咐你的事.
我吩咐你的事.
都遵守了.
是這麼聽話.
不要自己躲在家裡.
我們已經是約旦河東了.
我們不進去河西了.
不跟你一起打了.
你努力一點吧.
我在這裡祝福你.
不是.
你要遵守摩西的吩咐.

$^{201}$遵守上帝的吩咐.
是一起進去打仗的.
好.
既然是這樣.
你們沒有離棄弟兄.
好啊.
現在你可以回家了.
打完仗了.
不過你回去了.
仍然要謹慎地.
遵守摩西吩咐你們的誡命律法.
要遵守誡命律法.
最重要的一個字.
要愛耶和華.
你們的神.
要遵守.
要愛.
要行他的道理.
要守他的誡命.
要專靠他.
盡心盡性.
事奉他.
約書也挺厲害的.
這個字說得很清楚.
你要謹慎.
遵行耶和華的誡命.
愛耶和華.
行他的道.
守他的誡命.
專心靠耳靠他.
和盡心盡性.
事奉他.
心就是裡面的生命.
性就是意志力.
盡心盡意志.
盡心盡意.
順心事奉耶和華.
你們要遵守.
第二點是.
雖然現在回家了.

$^{241}$不過無論去到哪裡.
你都要專心.
謹慎.
遵行神的話.
愛耶和華.
行他的道.
守他的誡命.
專心依靠他.
不要依靠自己的聰明才智能力.
武功.
要用你所有的東西.
能力.
所有的能力.
去事奉耶和華.
要記住.
為什麼呢.
為什麼要特別提醒他們呢.
其他的支派都在.
因為他們要回到約旦河的東.
約書亞和其他百姓就在西部.
很遠的地方.
不是今天有互聯網.
又有社交媒體.
打個電話.
沒有的.
真的不知道還會不會再見面.
因為很遠我都不會去一次.
而且行曠野路去.
郊野的路.
很麻煩的.
不是隨便去.
所以今朝一別.
真的不知道會不會重逢.
在世上.
因此.
你記住.
記住.
長篇地講.
遵行耶和華的律法誡命.
愛耶和華.

$^{281}$行他的道.
守他的誡命.
靠他事奉他.
記住.
於是約書亞為他們祝福.
打發他們回到自己的帳篷.
因為臨別贈言.
所以要講得很清楚.
對他們說.
你們帶很多財物.
生畜和金銀銅鐵.
很多衣服回到你的帳篷.
這都是你們和我們一起.
整個十二支派一起打仗的時候.
從仇敵奪利.
你拿回去.
不要只是自己用.
回去和那邊看守著的弟兄.
一起分享.
他們放下了.
奸爺仔那在約旦河的東邊.
牲口都放在那裡.
是不是全部男人過來呢.
絕絕絕大部份來了.
但都要留少少人幫忙.
粗重事都要男人做.
因此.
要把你們所有奪取到仇敵的東西.
拿回去和弟兄分享.
都幾好.
約書亞這裡都講了.
大家要公平點.
分享公道點.
其實他和你看守財物.
看守你的家人.
都是要的.
你們要獲得的.
要和我們分.
交給他們.
裡面人家德人和馬來西半個支派.

$^{321}$半個支派靠近約旦河.
一帶迦林地就在那裡過河.
他們回去之後.
在約旦河那裡.
靠近約旦河那一帶迦林地帶.
在約旦河那邊.
捉了一座壇.
那座壇看上去是高大的.
不介意我又按上去看看圖.
他們回去打完仗.
過河了.
馬來西加德三個地方.
黃色,藍色,綠色,黃色.
三個地方淺黃.
他們累積到這裡.
大概在哪個位置.
就在靠近約旦河.
我想應該是中間一點吧.
總之就是在這個附近.
不知是哪個位置.
不過是這邊.
河的東邊.
就在這個位置.
建了一個壇.
遠遠地看到.
這個壇很高大.
建了一個壇.
捉了一座壇.
那壇看著高大.
以色列人聽見裡面.
加德馬來西支派的人.
靠近約旦河邊.
在迦林地屬於以色列人的那邊.
捉了一座壇.
留意第十二節.
很生氣.
傳媒中一聽見.
就聚集在視羅.
要上去攻打他們.
過河攻打馬來西支派.

$^{361}$要攻打他們.
即是說.
河東邊的弟兄聽見.
就要攻打西邊的弟兄.
為什麼?.
因為他們聽見.
他們在那邊建了一座壇.
而且很大.
不是吧?.
剛剛離開.
跟我們一起打完仗.
跟隨爺爺說在這裡得勝.
你回去馬上建一座壇.
拜第二個神.
很生氣.
因為他們以為他們拜第二個神.
很生氣.
豈有此理.
於是要打他們.
大家跳了二十四節.
由十二節跳二十四節.
二十四節.
那邊三個支派的人說.
我們行者是並非無故.
不是隨便的.
是故意做的.
為什麼?.
因為我們過了河西邊這邊住.
離開了東邊最多人.
我們同宗同族的兄弟姊妹.
我是怕我們的子孫日後.
你們的子孫日後跟我們的子孫說.
你們與耶和華以色列的上帝有何干涉呢?.
留意一下,要有很清晰的頭腦.
我們這座壇是故意建的.
建來做什麼呢?.
因為怕日後你們的子孫對我們的子孫說.
你們與耶和華以色列的上帝有何干涉呢?.
為什麼無緣無故說你們有信耶和華?.
我們住在河西邊才是以色列地方.

$^{401}$你們住在東邊.
又說耶和華有什麼干涉呢?.
現在我們這一代很熟悉.
但可能一代一代過去之後.
你們的子孫跟我們的子孫開始久不見.
不認識.
以為我們是外族人.
以為我們是河東.
以為我們是外族人.
以為你們都拜我們的耶和華.
是不是耍花樣?.
是真的假的?.
關你什麼事?.
我們只是敬拜神.
我們怕你們的子孫不認識我們的子孫.
大家都是敬拜耶和華的.
所以我們故意建一座壇.
那座壇不如為自己築一座壇.
是為獻梵祭,不是為獻別的祭.
所以我們建這座壇不是來獻祭的.
是要在你我中間和你我的後人中間作證據.
所以我們建這座壇做一個紀念.
好叫我們有一載在耶和華面前獻梵祭.
平安祭和別的祭事奉祂.
免得你們的子孫日後對我們子孫說.
你們與耶和華無份了.
其實我們建這座壇.
都是想證明給以後的歷代聽.
其實我們都是敬拜耶和華.
是和你們一樣.
我們都是同源同宗同族的.
信仰都是一樣的.
我們就是這樣想.
於是當時過去大興問罪之師的祭司菲利克.
和他的領袖和他的軍長.
聽見裡面的加德人馬來西亞人這樣說.
都以為美啊,多好啊.
好啊,其實這樣好啊.
哎呀,真是誤會了.
你們又沒有書信講清楚.

$^{441}$我們聽到你建這座壇就很生氣.
打算過來打你.
這麼快就離開耶和華,這樣不行.
原來是這樣,是怕著將來我們兩族人.
我們河兩邊的同宗同族.
不認識對方.
這樣就誤會了.
好啊,建這座壇表示大家都是相信神的.
以色列人,河東河西的人都以為是美事.
於是他們回去回報.
第三十三節回去回報.
原來河東那三個支派是這樣想.
明啊明啊,好啊,就稱頌神以事為美.
不再說去攻打他們.
裡面的加德人毀壞他們的地方.
就稱頌神.
啊,真是神真是好東西.
感動了他們懂得這樣做.
於是裡面的加德人就蓋了這座壇的名字.
就叫做正壇.
正名的正.
意思就是說這個壇在我們中見證明耶和華是上帝.
這個壇的象證明.
其實意思就是我們都是相信.
耶和華是我們共同敬拜的上帝.
信仰是一樣的.
我們沒有偏離的.
不過太遠.
其實太遠我們都不能夠時時回來視來的會望.
因此我們在那邊同樣是敬拜主的.
都同樣是獻祭的.
大家都明白.
保持著我們的信仰.
大家是一樣的.
免得過了幾代大家就不明白了.
哦,原來是這樣,收到收到.
好啊,這件事好.
所以出現了一個壇叫正壇.
證明我們中間都是相信耶和華.
獨一的神.

$^{481}$是共同的俗裔,共同的信仰.
這樣也不錯.
我在這裡看到這件事的時候.
我心裡也覺得很開心.
就是他們一直依靠神去打仗打仗打仗.
經過這麼多場戰爭之後.
勝利,贏了.
拿到那塊地了.
按照神的吩咐拿到了.
神的應許給他們拿到了.
他們是很注重.
我信靠神,很開心.
其實真的相當不錯.
而且還為將來想.
我們的子孫可能會忘記.
所以我們就做這個壇,叫正壇.
我們有很多時候都不想將來的.
有很多時候我們今代敬拜神.
我看到很多弟兄姊妹今天很惆悵.
因為他們作為父母.
小孩子小小一個跟他們上主教學.
很乖啊,聖經朗朗上口.
不過很多到中學,很多到大學.
不肯去教會了.
不肯去.
甚至部分不肯承認自己是基督徒.
這些父母痛心到不得了.
怎麼會這樣呢?.
一代就變了,還沒過已經變了.
你說我孫子還會不會去教會?.
真的很憂心,怎麼辦?.
他們厲害,以前的人看得很遠.
將來我們要建一座壇.
令到我們可以記得將來我們相同信仰.
無論你們我都是信靠獨一的神爺華.
因此弟兄姊妹我鼓勵大家.
如果你有疑慮的.
你都要在你家建一座正壇.
即是你在家都要很敬田地生活,敬拜.
讓小朋友看到你的禱告.

$^{521}$看到你自己對神的尊敬.
而且要帶他一起有家庭的祭壇.
即是有家庭的聚會,家庭的祭壇.
我們一家人是敬拜神的.
將這個很重要的思想放進去.
可能後來他都會被同學朋友.
返去大學之後有什麼信仰的衝擊.
不過他仍然是跟隨神的.
而且可能他脫離了,衝了去.
但他會醒覺回頭.
所以小時候就要踏好.
他們一回家就馬上想我們要建一座壇.
否則我們很久都過不了河.
去不到會幕那裡敬拜.
我們會脫離的.
無論是跟那邊弟兄的關係脫離.
甚至我們跟神的關係都脫離.
所以我們要建一座壇來提醒對面的人.
提醒我們的人.
我們都是在中間能夠證明耶和華是我們的神.
這點很重要的.
記得要為自己的子女做一座家庭的祭壇.
這就是我們敬拜的,這個很重要.
當我在台灣做宣教工作的時候.
有件事讓我知道原來這麼難突破他們的民間信仰.
初一十五.
我們南方人叫做「造衙」的時候.
初一十五會做什麼呢?.
一家人就會在一間房子.
當然房子不是大廈.
以前都是一間間透天三層樓一家人住.
房子都是平房比較矮.
他們就會在自己房子的門外露天.
就會開張桌子.
放些香爐插香.
放隻雞.
放些祭物.
然後排好隊.
然後輩分由爺爺到爸爸到孫到兒子到孫這樣排.
就排好隊.

$^{561}$男尊女卑排好輩分.
所以孫子和兒子就排在後面.
忠賢平時見到爸爸嘩嘩嘩嘩嘩嘩嘩嘩嘩嘩嘩嘩嘩嘩的粗口.
或者罵人很絕,媽媽的嘴巴很厲害.
見到爺爺不好,不過呢,經常說話沒有一句真.
不過去到這一天這個時候,見到他們全部都非常嚴肅.
一柱清香,三支香zar 在頭頂,向著祭品三個宮.
然後就過去插好支香,然後到爺爺之後到兒子.
兒子之後到孫子,孫子之後到室,一個個排過去.
全部都拿著那柱香,神情肅穆十分嚴重.
那些小朋友就見到爸爸平時多麼不好.
但這個神情肅穆的一刻,他們知道這件事是認真的.
這件事是認真的,這個就是我們拜的神.
你知道爺爺離開世界,爸爸離開世界,到他自己排頭位.
他就會叫兒子孫子要跟著我,一定要這樣做.
這個就是我們拜的神,所以他們世世代代都拜.
覺得是我們家族,這個祭祀的過程.
就是一個正壇,不過是正台,台在門口.
不過就是這樣,敬拜天地,敬拜各樣假神.
哇,這麼嚴肅,起了壇,就會帶到他一家人世世代敬拜.
當你叫他信耶穌,我不信的,我們家不拜這些的.
我們家是有自己的神,你那些不是我們拜的.
多謝了,不用了,不要再拜我了.
哇,刀錯不入,所以你會知道在家有正壇是多麼重要的.
真的,希望你家都有家庭的聚會,有家庭的正壇.
家裡不要缺少十字架,有些擺著擺著.
對你整個家人的信仰都很有幫助.
好了,我們一起祈禱.
天父上帝多謝你讓我們看到這段聖經.
故事很長,不過很有意思.
主我們都立志要在家做一個壇.
高舉禮叫所有人都知道.
耶和華是我的神,是我們家的神,是我們的神.
讓我們站立得穩,一代一代傳下去.
如果我們是教徒家庭,讓我們築起家庭的祭壇.
帶領小朋友一代一代忠心不移的跟隨我們的主.
榮耀你,以你是全能獨一,是我們敬拜的神.
願你得到榮耀,願你成為我們家的主.
禱告奉耶穌基督名求,阿們.
\newpage



\section{歌羅西書 4:7-14}
\label{sec:032189}
\textbf{使徒保羅的監獄合照}~~~~羅錫堅牧師
\newline
\newline
連結: \href{http://www.church.com.hk/acms/content.asp?site=cdc\&op=show\&type=product\&code=032189\&layout=sermon}{\texttt{032189}} ~~~~ 語音日期: 2026-05-24
\newline
\newline
\hyperref[sec:032188]{\small{< < < PREV SERMON < < <}}
~
\hyperref[sec:index_chronic]{\small{[返順時目]}}
~
\hyperref[sec:032190]{\small{> > > NEXT SERMON > > >}}
\newline
\newline
歌羅西書 4:7-14
\newline
\begin{longtable}{cl}
\hline
\hline
章節 & 經文 (和合本修訂版)\\
\hline
4:7 & \begin{tabularx}{0.7\textwidth}{X} 推基古是我親愛的弟兄，忠心的僕役，和我一同作主的僕人；他要把我一切的事都告訴你們。 \end{tabularx} \\ \\ \relax
4:8 & \begin{tabularx}{0.7\textwidth}{X} 我特意打發他到你們那裡去，好讓你們知道我們的情況，又讓他安慰你們的心。 \end{tabularx} \\ \\ \relax
4:9 & \begin{tabularx}{0.7\textwidth}{X} 我又打發一位親愛忠心的弟兄阿尼西謀同去；他也是你們那裡的人。他們會把這裡一切的事都告訴你們。 \end{tabularx} \\ \\ \relax
4:10 & \begin{tabularx}{0.7\textwidth}{X} 與我一同坐牢的亞里達古問候你們。巴拿巴的表弟馬可也問候你們。關於他，你們已經得到指示；他若到你們那裡，你們要接待他。 \end{tabularx} \\ \\ \relax
4:11 & \begin{tabularx}{0.7\textwidth}{X} 稱為猶士都的耶數也問候你們。奉割禮的人中，只有這三個人是為神的國與我作同工的，也是使我心裡得安慰的。 \end{tabularx} \\ \\ \relax
4:12 & \begin{tabularx}{0.7\textwidth}{X} 有一位你們那裡的人，作基督耶穌僕人的以巴弗問候你們。他禱告的時候常為你們竭力祈求，願你們能站穩而成熟，充分確信神一切的旨意。 \end{tabularx} \\ \\ \relax
4:13 & \begin{tabularx}{0.7\textwidth}{X} 他為你們、老底嘉和希拉坡里的弟兄多多勞苦，這是我可以為他作見證的。 \end{tabularx} \\ \\ \relax
4:14 & \begin{tabularx}{0.7\textwidth}{X} 親愛的醫生路加和底馬問候你們。 \end{tabularx} \\ \\
[1ex]
\hline
\hline
\end{longtable}
$^{1}$我很喜歡看書.
中外的小說我都看了很多,那些名著.
其中一本.
就是.
三劍客.
這本書是.
法國的.
作家大仲馬寫的.
Three Musketeers.
這三位劍客.
很.
合一,他們堅貞忠誠.
志願共享.
性格配搭.
共同進退.
我不懂法文,我看英文,其實中文也有.
如果大家有空,我鼓勵你去看,你會很開心.
我又喜歡看《三國演義》,有沒有人看過《三國演義》的,羅冠中有沒有?.
有,書不是看戲.
《三國演義》有三位.
劉備,關公,張飛.
桃園傑義,生駛世文.
融合了君臣之義,因為劉備是皇帝,君臣之義和兄弟之情.
經得起極端環境考驗.
剛膽相照,彼此支持.
這兩本書都有一個共同點,就是很珍貴的友情.
我想到我從小開始在.
中學,大學.
在石油公司工作,後來在教會裡面事奉.
直到現在,上帝都為我預備一些和我很合得來.
很支持我彼此援助的友情的朋友,很寶貴.
我相信你們也有.
今天要跟大家講的就是保羅.
剛才葦紅跟我們讀《亂經文》.
如果你說,怎麼講呢?.
其實我想的主題就是.
保羅在監獄裡面的一個合照.
這幅就是保羅和八個人的合照.
你說羅木斯,那時候沒有相機,你怎麼拍出這樣的東西呢?.
我是用AI做的.

$^{41}$在這裡我要講一下AI有沒有用.
有信徒問羅木斯,AI好不好用?應不應該用?.
會不會誤導我們?.
AI其實很簡單,就是一個很大的資料庫.
全世界所有的人,所有的意見都放在裡面.
就是這麼多.
如果你.
簡單來說,現在你想問AI問題,大致上是一個資料庫.
如果你問一個問題,AI就在資料庫裡面找一個答案給你.
這個答案就是大部分人認為對的答案.
這個答案會不會正確呢?.
未必的,因為大部分人認為對的,可能上帝不認為對的.
我做過一個實驗.
我預備這篇導讀的時候.
我就問AI.
耶穌復活.
是否歷史事實.
清不清楚?.
有些人說是真的,有些人說是假的.
其實就幫不了你.
所以AI是有用的.
如果不是的話就做不到這幅畫出來.
但是.
AI用來給你去找一些數據.
或者我記得一首詩,不知道是誰作的.
朝池白帝採雲間,你打上去.
原來是李白.
李白因為被皇帝赦免,所以很開心,李白這麼多詩最開心是他.
你知道這些就OK.
但是如果你要找聖經,研究聖經.
研究聖經,你要帶一個禱告的心.
讀聖經,然後求神的靈,聖靈.
幫助你去了解.
上帝今天.
要你在這段聖經裡面.
得到什麼.
這個.
AI就對不起,幫不了你,因為AI不是上帝.
所以我們.
AI有用的.

$^{81}$但是就不要亂用.
不要每樣都找AI,以為它有用.
我今天骨痛.
肩膀痛一打,AI說我癌症,現在還沒到這個地步.
所以尤其是看聖經就要小心一點.
今天我希望大家用一個禱告的心來看看.
保羅這一段經文,這裡就是哥羅西斯最後一段.
六月我都會講讀,我又講第二樣東西,又不是講哥羅西斯.
因為哥羅西斯沒有第五章,第四章已經講完.
我們就要看看.
究竟.
保羅身邊這八個人.
網上這八個人.
他們是怎樣的呢.
你會發現其實都很多好東西.
很多好東西你學習到.
第一個他提到的就是推基古.
我一切的警方,推基古會告訴你們.
他是我所愛的弟兄,是忠心的僕役,也是在主裡作僕人的.
我派他到你們那裡去,使你們知道我們的警方,並且安慰你們的心.
這個推基古就和保羅並肩事奉很久.
而且他特別是做.
差不多保羅的.
信差,把一些供應品,比如支持一些有需要的教會,或者保羅的信.
傳給教會,他是做那樣東西.
他是一個忠心.
一個忠心能夠安慰人的人.
那時候保羅就在坐牢,第一次坐牢.
他第一次坐牢.
比他第二次好一點.
政府.
給他租了一間屋.
然後就.
准人探他.
不過.
保羅被控的罪,就是顛覆叛國.
嚴不嚴重呀朱先生.
朱先生舊時在入境處,其實不是很關事.
不過他都知道這些例子.
顛覆叛國.

$^{121}$他坐牢.
你經常去探他.
可能你都有問題.
保安局看了一下,推基古經常去探他,捉他審他.
所以其實.
這八個人,他群懷保羅.
就.
其實有風險.
但推基古就不怕.
他.
和他做.
譬如那些信.
監獄書信,那些就是他帶去的.
還有他需要.
保羅需要的東西,他又帶來.
這個推基古就是其中一個他講的.
第二呢.
第二個叫做歐尼西姆.
這個是,如果你看腓利門書就認得他.
他原本是腓利門.
的奴隸.
那時的奴隸不是每個黑人都是奴隸,那時的奴隸.
羅馬帝國有三分一人以上是奴隸.
有些人說有半人.
這個奴隸.
他可能犯了些事,可能有什麼原因我不知道.
就逃走了.
去了羅馬.
去了羅馬就認識了保羅,信了主,現在在監獄.
保羅在監獄裡面.
他都照顧他.
但如果你看腓利門書,後來保羅.
寫了封信.
要歐尼西姆回到腓利門.
向腓利門認錯.
那封信裡面,他又叫腓利門.
饒恕弟兄.
結果怎樣呢?結果.
腓利門饒恕了他.
最後這個歐尼西姆.

$^{161}$竟然做了以弗所的.
監督.
監督,bishop,主教.
你看到信耶穌.
人生一個新的開始.
所以如果你未信耶穌.
你說我現在.
這樣的年紀,一把年紀.
又好像沒什麼作為.
那你信耶穌的時候,舊事已過.
都變成新.
這樣一個新的開始,還不遲.
所以歐尼西姆就有這樣的情況.
第三個就阿里達古,這個就很厲害.
阿里達古是一個希臘信徒.
來自帖撒羅尼迦加.
他竟然和保羅一起坐牢.
原來當時羅馬,他第一次坐牢有這樣的例子.
如果你坐牢需要人照顧他,你可以找你的奴隸.
或者你的僕人.
跟你一起坐.
證明他是你的近身侍從就可以了.
這個阿里達古.
volunteer成為保羅的奴隸,可以和保羅一起坐牢.
然後就可以照顧他.
這不就很偉大.
還有,在紹行傳十九章記載.
有一次.
保羅去傳福音,全城騷動.
他們抓住了保羅.
和阿里達古.
打他們,所以其實他是和保羅.
共患難的一個人.
一個很堅定的人.
我們其實是為他感恩.
還有馬可.
保羅提過兩次馬可,關於馬可你們受了吩咐.
他若到你們那裡,你們要接待他.
這個馬可很特別.
馬可曾經.

$^{201}$跌倒離開了保羅.
那件事是這樣的.
在紹行傳十五章.
馬可.
和他們傳福音.
做著做著頂不住,就走了.
回到耶路撒冷.
巴拿巴.
過了一段時間,加拿巴就想帶馬可.
和約翰.
帶馬可和保羅一起再出發.
但是保羅就認為不應該帶他們去,因為他從前在.
龐菲利亞那裡離開了他們.
不和他們一起作供.
於是巴拿巴和保羅之間引起一個紛爭.
各持己見.
以至分手.
巴拿巴就帶了馬可.
保羅.
就帶了另一個去了其他地方.
帶了西拉.
我這裡要講講巴拿巴.
巴拿巴是一個很厲害的人.
聖經很少提及他.
但是他的人際關係和憐憫的心很好.
他是第一個.
保羅決志信耶穌之後.
阿拉尼亞就開眼禱告.
但是.
耶路撒冷的.
使徒就不敢.
接待他,害怕他,因為他迫害教會.
巴拿巴是第一個推薦保羅出去宣教.
所以其實他是保羅的恩人.
今次就和保羅意見不合,很可惜.
當馬可.
離開他們,跌倒的時候,巴拿巴就挽回他.
所以巴拿巴其實是一個很有心挽回的人.
我覺得巴拿巴.
很值得學習,聖經沒有很直接的提及,但其實他是一個很偉大.

$^{241}$很謙卑.
很有心腸.
的一個.
領袖.
但是在這裡保羅坐牢的時候.
馬可已經改變了.
馬可又陪伴他.
在監獄裡面.
服侍他.
他.
在這裡稱讚馬可和叫他們.
我曾經講過.
他如果到你們那裡,你們要接待他.
我看到.
就算跌倒了都可以迴旋.
下一個就是尤士多,尤士多的耶穌.
也問候他們.
尤士多和馬可在八個人當中,是這兩個猶太人.
因為他們兩個都受過國禮.
只有.
其他的全部都是外邦人.
只有這幾位,這些人是神.
國和同工,成為他的安慰.
這八個人.
在保羅患難的時候.
他是安慰他.
第二,當你患難的時候.
有沒有人可以和你聊天,有沒有人可以鼓勵你.
這個就很重要.
我們教會.
存在.
有這樣的.
很重要的功能,就是彼此相顧.
我看到很多新聞.
尤其是在日本最多.
老人家死了.
子女很久都不知道.
死了很久,半年都沒有人發現他.
就是這樣的情況.
其實這個很可惜.

$^{281}$教會就不應該有這樣的事情發生,所以我們希望大家.
能夠有更多的,更緊密的彼此聯繫,彼此相顧.
接著就是以巴弗.
保羅說:以巴弗問候你們,他是你們那裡的人.
是基督耶穌的僕人.
他禱告的時候,常常竭力為你們祈求.
喂,你試過禱告,竭力為人祈求沒有?.
有些信徒就睡覺之前,睡覺的時候蓋著一張被子禱告.
第二天早上起床才說阿們.
這個就不會是竭力,睡著了.
但是這個以巴弗就竭力為信徒祈求.
好使你們.
在神的一切旨意中完全站穩,滿有堅定的信念.
我可以為他作見證.
他為你們.
和在腦底加希拉波的人.
多愁無苦.
這個以巴弗.
以巴弗的名字.
其實意思就是禱告者.
一個禱告的人,我希望大家都學習他做一個禱告的人.
由路加醫生,路加醫生就是那位作《史道行傳》和.
路加福音.
馬可福音就是這個馬可寫的,所以其實在白雷裡面有兩個人都寫了.
聖經的書卷.
只有路加在我這裡.
你找到馬可就把他帶來,這次提摩太后書就是第二次的坐監.
是最後一卷的書來的.
沒多久就被處決.
斬頭.
因為他犯了叛國罪.
對羅馬帝國.
顛覆羅馬帝國,真是很冤枉.
所以看到其實這八個人.
是成為保羅的.
死黨.
還有一個叫底馬.
因為底馬.
貪愛世界.
他原本好端端,在保羅第一次坐監的時候都陪他.

$^{321}$等到.
第二次的時候他可能已經有點害怕.
都是走了.
結果.
或者未到第二次.
他已經離開了貪愛世界.
去了帖撒羅尼迦.
所以看到裡面有一個馬可就是迴轉的.
有一個就沒有.
最後就散了.
其實.
都給我們看到人的軟弱.
但是神的恩典.
和我們弟兄姊妹之間.
沒有一個人是完美的.
你要找一個很完美的人.
要找一個很完美的教會找不到,有沒有人覺得我們全印度是很完美的.
應該有八成九成但都未到.
所以我們一定要.
要竭力彼此.
來禱告.
我們.
做一些反思.
在保羅的生命.
裡面這個團隊在他最後.
那幾年這個團隊對他很重要.
他有什麼特點呢?這個團隊第一就是支持牧者.
和他並肩作戰成為他的慰藉.
安慰.
很重要.
這個團隊亦都是彼此相愛.
互相配搭彼此顧念.
保羅在以弗所書第四章都教導我們.
他說.
我借為主被囚禁的.
勸你們行事為人.
要配得上你們所蒙的膚照.
凡是謙虛.
溫柔.
忍耐用愛心彼此寬容.

$^{361}$以和睦聯繫.
竭力持守聖靈所賜的合一.
這個教會.
是耶穌基督的身體.
我們在當中.
既然我們每個人都知道自己不是完美.
我們彼此之間.
就要溫柔忍耐.
寬容.
謙卑.
當你謙卑的時候你就容易得寬容別人.
因為如果你覺得自己我是最厲害的.
我認了第二就沒有人敢認第一.
很難謙卑.
要竭力保守聖靈所賜的合一.
還要全身的投入.
每個信徒.
在教會裡面都應該完全投入這個團隊裡面.
不要一則薄.
我們的投入有一個很重要的就是要透明.
什麼叫透明呢.
就是願意透露你自己軟弱那一面.
保羅.
第十八節都說你要紀念我的困所.
他的軟弱就是被人.
鎖在監獄裡面.
被人告他顛覆.
所以就.
做不到.
還有要和解.
有個和解的心.
好像.
歐尼西姆.
和馬可.
他們都經歷過憂患.
經歷過跌倒.
而他們.
是得著寬恕.
Reconciliation.
和解,彼此之間和解.

$^{401}$保羅都學到東西他都和解,原本他不要馬可.
他都知道其實阿巴拿巴是對的.
他就接受馬可.
而且兩次.
這幾節聖經裡面兩次提到馬可.
所以我們.
弟兄姊妹之間要有一個.
和解的心.
千萬不要憤怒.
憤怒的人太多.
因為耶穌基督.
說得很清楚.
如果你.
你不饒恕你的弟兄,上帝都不饒恕你.
上帝.
上帝饒恕了你的,上帝將你的罪從過去到將來都饒恕了,但是耶穌基督.
嚴重到如果你不饒恕你的弟兄.
還說我實實在在的鬧鬧鬧鬧鬧鬧鬧鬧鬧鬧鬧鬧鬧鬧鬧鬧鬧鬧鬧鬧鬧鬧鬧鬧鬧鬧鬧鬧鬧鬧鬧鬧鬧鬧鬧鬧鬧鬧鬧鬧鬧鬧鬧鬧鬧鬧鬧鬧鬧鬧鬧鬧鬧鬧鬧鬧鬧鬧鬧鬧鬧鬧鬧鬧鬧鬧 bou.
心裡面就不要恨弟兄.
他是壞的.
他是麻煩的.
他借了你的錢沒有還到的.
你都要還他.
還要毀身基督.
我們要.
我們在教會裡面.
不是消費者.
不是回來聽一個好精彩的講道.
而是我們是彼此毀身毀身基督.
推基輔這個名的意思就是可靠.
Trustworthy.
信得過的.
以巴弗我提過就是禱告者.
要可靠的.
要一個禱告的心.
有些人羨慕保羅有八個這麼好的朋友.
甚至你羨慕三劍客的三位.
或者是.
三國演義的劉國鑫的三位.
其實主人所謂我們預備這個教會.

$^{441}$是叫我們享受到.
真正的友誼.
我們這裡.
我們教會如果班齊有四人左右.
其實.
你應該找到不止八個.
當然你坐監的時候有多少人探你.
你入醫院應該多些人探你.
我們做一個.
做一個的應用就是支持牧者.
我們怎樣在全因堂裡面實踐.
這幾節聖經教導支持牧者.
我們全因堂現在只有一個牧者就是康惠明牧師.
他做了.
十年差不多.
接近十年.
身心疲累.
很欣賞執事給他有三個月的.
按釋.
你不要以為自己很偉大.
我們總會.
做五年.
總會的同工.
我曾經在總會做總按釋的時候.
我們有.
半年的按釋.
如果你不拿的話足夠十年.
所以三個月都需要.
我們要支持他好像這八位支持保羅一樣.
我曾經有一次我和我太太開了一間新的教會.
很多人都是在外地來的.
我安排一些領袖一些組長都是在外地來的.
新的.
幾個月.
其實有點危險.
但是不安排就不行因為疫情發展得很快.
有一次我病了.
發燒.
我不是新人我大人發高燒.
有幾天都沒有看電郵.

$^{481}$那時候沒有WhatsApp.
只有電郵.
幾天都沒有看電郵.
四天.
有個信徒打電話來.
牧師你快點看電郵散了散了散了.
原來這幾個組長就不太行.
我一看到.
當堂燒都沒有賣.
馬上打電話要處理.
所以我其實很.
我真的求主給我們教會.
有成熟的信徒領袖成熟的執事.
等到江牧師要病的時候都可以安心的病.
安心的病.
保羅在柯林多後書.
十一章二十八節他說.
除了外面這些麻煩事.
宗教會令到我掛心的事天天壓在我身上.
天天的意思.
就是不停止的一直壓在我身上.
我相信江牧師.
是經歷過這些在過去的十年.
畢業.
2013年.
不止十年.
這間教會就是九年多但他總共.
事奉了十三年.
我自己其實都經歷過很多這些.
我.
幾個月前臨危受命.
本來就在台灣做宣教full time.
現在.
董事會說.
麻煩你接任多做一份full time.
兩份都是義工來的不過就是full time的.
那麼.
為什麼義工呢因為沒有錢.
那麼.
就要幫我們處理.

$^{521}$這些.
警察會的一些問題.
那麼其實.
很辛苦.
不只是時間上的問題而是一些事.
一些關係上的事.
很難處理.
真的我.
天天壓在身上.
其實一個傳道人.
一個傳道人就不是打.
不是打工.
你放工就不用理.
因為他.
24小時.
好像保羅說宗教會令我掛心的事天天壓在我身上.
所以我希望大家體諒江牧師.
他休息了一個半月.
求主叫他的身體健壯.
你主動的愛他.
和他一起事奉.
如果有些事你不贊成你體諒他.
他不是完美.
你都不是完美.
有時候不知道哪個比較好.
你不知道哪個比較好.
那麼你.
應該信服的就信服.
除非和真理有抵觸.
江牧師的真理呢.
肯定沒有問題.
所以我們要照顧.
要體諒.
這個牧人.
你想想保羅這八個老友.
在他被控顛覆叛國的時候都仍然黏著他.
仍然願意冒險.
所以.
我們都應該有這樣的心智.
彼此相愛,彼此相愛就講了很多.

$^{561}$但是.
講就容易做就難.
謙虛溫柔忍耐愛心和睦.
這是雙方面的.
我們期待.
別人愛我們,為什麼這個教會這麼冷酷的.
你自己又沒有愛人.
你沒有愛人呢.
別人怎麼愛你呢.
記得有一次我公司派我去英國.
去做一年的事.
打官司,我做一個specialist.
打官司的時候.
在英國參加那裡的華人教會.
第一次去.
到處有人握手.
有位姊妹跟我說.
我們英國就不流行握手,你來到好像你擁了這間教會那樣跟我們握手.
我們很開心,你這麼快才投入.
第一個meeting.
其實.
我想所有教會都是耶穌基督.
好像7-11.
街頭街尾兩間都是.
都是大老闆.
我去教會裡面.
當然彼此握手,我們怎知道你英國不流行,我們加拿大流行.
但是很開心.
所以我們要那個.
彼此相愛那裡很重要.
全身投入.
其實.
聖恩堂有很多.
其實.
很多東西做出來的.
你看我們那個碧埔版.
我知道Dennis他來這裡的時候,他很高興見到他自己的照片在這裡.
其實海鮮教士.
很需要我們為他祈禱.
我們.

$^{601}$我們很多的事工很需要我們做.
我們的兒童工作,我們搬到北角,我們怎樣重新起步呢.
我很欣賞寶珠她在這裡.
上帝猜拜她來.
和我們一齊.
每一個來到我們當中的都是神恩典.
沒有coincidence,沒有巧合的,你今天在這裡.
是上帝放你在這裡.
他要你聽到一些東西,他要你講到一些東西,要幫到其他人.
我們要全身投入.
香港教會面臨老化的現象.
播道會.
在五六年前有2萬7千人.
現在就變成1萬7千.
有些人就出了國.
去了外國.
有些人.
就上了網上,就未下來.
有些人.
就散了,沒有回教會.
有些人就回了天家.
因為人會老的.
我和我太太住.
在香港住酒店,在台灣住酒店.
我在香港住大酒店旁邊的大酒店.
那間大酒店.
星期六就很忙.
所以其實我看到人生.
如果我現在做牧師很方便,走兩步就去到那裡,去主持安息禮拜.
不過現在很多人不流行在大酒店做.
在醫院直出.
所以其實.
我們.
面對老化的.
教會老化現象.
我們要去找青年人.
去跟他們傳福音.
想用珠本的方法.
你不要說.
牧師我來世怎樣跟16歲的傳福音.

$^{641}$不要氣餒.
魔鬼就會令到我們氣餒.
你不要理會播導會好不好,總之你好.
我們全心投入.
毀身忌妒.
不怕勞苦.
因為這場屬靈的戰爭.
不是我們打的.
是上帝打的.
我今日毀身忌妒.
其實我73歲.
其實是辛苦的兩份工.
但都做.
因為上帝叫我做.
我都跟大家提過.
我們.
中國歷史裡面有很多人為國家.
捐助很多,連老都出去打仗.
所以我想到.
我們做基督徒.
上帝給你有生命氣息.
你就去做.
上帝會允許你.
你的日子如何.
力量如何.
我們要學.
保羅那七位那八位朋友.
照顧他.
不怕艱難.
彼此建立.
我們一齊低頭祈禱.
多謝你的話語奇妙.
保羅的生命裡面.
在周圍那八位.
他提到的.
朋友身上.
得著榮耀.
我們看到教會就因為你這樣.
這班人.
而得著建立.

$^{681}$我們感恩.
我們為全恩堂眾弟兄姊妹禱告.
特別為康牧師禱告.
在他休息期間求你自己.
叫他身心靈都建壯.
能夠得靈.
求你賜給我們弟兄姊妹有智慧有愛心有謙卑.
彼此配合.
願你的靈引導我們潔淨我們.
讓我們有平安起來.
身心靈都得著建壯.
我們這樣祈禱奉耶穌基督的名.
阿們.
\newpage



\section{羅馬書 3:21-31}
\label{sec:032190}
\textbf{最好的福氣}~~~~潘建堂牧師
\newline
\newline
連結: \href{http://www.church.com.hk/acms/content.asp?site=cdc\&op=show\&type=product\&code=032190\&layout=sermon}{\texttt{032190}} ~~~~ 語音日期: 2025-10-26
\newline
\newline
\hyperref[sec:032189]{\small{< < < PREV SERMON < < <}}
~
\hyperref[sec:index_chronic]{\small{[返順時目]}}
~
\hyperref[sec:032191]{\small{> > > NEXT SERMON > > >}}
\newline
\newline
羅馬書 3:21-31
\newline
\begin{longtable}{cl}
\hline
\hline
章節 & 經文 (和合本修訂版)\\
\hline
3:21 & \begin{tabularx}{0.7\textwidth}{X} 但如今，神的義在律法之外已經顯明出來，有律法和先知為證： \end{tabularx} \\ \\ \relax
3:22 & \begin{tabularx}{0.7\textwidth}{X} 就是神的義，因信耶穌基督加給一切信的人。這並沒有分別， \end{tabularx} \\ \\ \relax
3:23 & \begin{tabularx}{0.7\textwidth}{X} 因為世人都犯了罪，虧缺了神的榮耀， \end{tabularx} \\ \\ \relax
3:24 & \begin{tabularx}{0.7\textwidth}{X} 如今卻蒙神的恩典，藉著在基督耶穌裡的救贖，就白白地得稱為義。 \end{tabularx} \\ \\ \relax
3:25 & \begin{tabularx}{0.7\textwidth}{X} 神設立耶穌作贖罪祭，是憑耶穌的血，藉著信，要顯明神的義；因為他用忍耐的心寬容人先前所犯的罪， \end{tabularx} \\ \\ \relax
3:26 & \begin{tabularx}{0.7\textwidth}{X} 好使今時顯明他的義，讓人知道他自己為義，也稱信耶穌的人為義。 \end{tabularx} \\ \\ \relax
3:27 & \begin{tabularx}{0.7\textwidth}{X} 既是這樣，哪裡可誇口呢？沒有可誇的。是藉甚麼法呢？功德嗎？不是！是藉信主之法。 \end{tabularx} \\ \\ \relax
3:28 & \begin{tabularx}{0.7\textwidth}{X} 所以我們認定，人稱義是因著信，不在於律法的行為。 \end{tabularx} \\ \\ \relax
3:29 & \begin{tabularx}{0.7\textwidth}{X} 難道神只是猶太人的嗎？不也是外邦人的嗎？是的，他也是外邦人的神。 \end{tabularx} \\ \\ \relax
3:30 & \begin{tabularx}{0.7\textwidth}{X} 既然神是一位，他就要本於信稱那受割禮的為義，也要藉著信稱那未受割禮的為義。 \end{tabularx} \\ \\ \relax
3:31 & \begin{tabularx}{0.7\textwidth}{X} 這樣，我們藉著信廢了律法嗎？絕對不是！更是鞏固律法。 \end{tabularx} \\ \\
[1ex]
\hline
\hline
\end{longtable}
$^{1}$《錄音》或《錄影》的版權是屬於基督教宣導會不覺堂所有.
我們的屬靈傳統之一是以真理教導為本.
歡迎您將此《錄音》或《錄影》的連結傳給親友.
懇請注意未經授權請勿轉錄或剪輯本《錄音》或《錄影》.
甚願好像《使徒行傳》二十章三十二節下所說.
「這道能建立你們,叫你們和一切誠誠的人同得基業」.
讓我們再次禱告.
因為您不單止賜給我們寶貴的救恩.
您也賜下寶貴的話語.
讓您的話語成為我們生命的光.
讓您的話語在我們聚會的時候.
就是您與我們同在.
您要向我們說話.
求主開啟我們的心.
讓我們能夠明白您寶貴的福音.
就能夠因著您所賜的福音.
感受到那種福分和喜樂.
我們同心願望您禱告奉耶穌基督的名字.
阿們.
剛才所讀的一段經文最後一句就說.
大衛這樣說:「過犯被赦免.
罪孽被遮蓋的人是有福的.
主不算為有罪的人是有福的」.
這句話是提醒我們.
我們說信耶穌不是錦上添花.
我很好,我可以再好一點.
信耶穌這裡說的是.
關乎罪德赦免.
是關乎到我們生死存亡.
我們今天所讀的羅馬書.
是保羅寫給羅馬教會的一封書信.
當時猶太人都在羅馬帝國之下.
帝國的版圖非常大.
當時來說,羅馬帝國就是以羅馬這個地方為中心.
建造了四通八達的道路網.
是方便軍隊的調動,貿易,文化交流.
讓人感覺到當時在帝國裡.
從任何一個地方都可以很容易直達中央羅馬的首都.
所謂「條條大路通羅馬」.
是一個事實.

$^{41}$這句話如果站在地理上和交通上.
甚至經濟文化上都是正確的.
但是如果說「條條大路通羅馬」.
作為一種殊途同歸的一種說法.
對福音來說就不正確.
因為要得到永生,只可以信靠耶穌.
羅馬書向我們展示出一個真理.
就是基督所成就的救恩,是恩順清儀.
不是靠律法,不是靠人的行為.
或者任何人的功德.
這個理念是基督教信仰的核心.
甚至是判別它是不是異端的其中一個很重要的教義.
沒錯,很多人信耶穌.
可能都是不同的人.
他們都是透過不同的經驗來做開始.
或者入門的方法有不同.
切入點有不同.
但是很重要的.
信靠耶穌,聚德赦免是不可以去繞過的.
所以今天我們所重讀的《仕途信經》.
我們都會說有說到聚德赦免.
我們也說「我信永生」.
在近代的教會有這樣的現象.
或者也有不少的信徒反省多一點.
都會看到一件事.
如果我們說「信耶穌」只是說「得永生」.
「上天堂」是不足夠的.
好像整件事都只是將來的事.
今生的生活都很重要.
所以這個說法基本上是正確的.
但循著這種說法,有些情況又會出現極端.
就認為全福音不一定.
或者不用說「聚德赦免得享永生」這些.
因為一說這些,就好像馬上將福音簡化.
但其實一條橋是不可能沒有起點和終點的.
正如我們早前在唐慶主日的講題裡.
都說到「以信為始,以信為終」很重要.
基督徒的人生的確不只是這個起點.
他決志舉手,然後就進入天堂.
就是這麼簡單.

$^{81}$中間那段路就是我們今天的生活.
是很重要的.
我們的使命是非常重要.
這個是有立足點的.
我們的使命不是為了培育一種使命感.
我們能夠實踐使命的感覺良好.
而製造一種使命.
完全不是.
我們的使命是基於清清楚楚的真理的事實.
是基於福音本身所反映的.
人人都需要的.
人生真實而終極的需要.
我們從羅馬書裡面學習福音的真理的時候.
我們常常聽.
甚至經常重複「因信正義,信耶穌,信耶穌」.
就覺得是否羅馬書其實是寫給微信的人呢?.
如果你有讀清楚.
你就知道它寫的不是街上那些微信的人.
羅馬書作為一封新約教會的書信.
是寫給教會的.
通用的不是像福音單章那樣派.
而是給當時已經信了主耶穌的基督徒.
不單止是給當時在羅馬外邦的地方.
那些從來不認識真神的人.
他們已經歸向了基督那些外邦的信徒.
也是給從前已經認識耶和華.
懂得舊約聖經的.
本身已經被稱為神子民的那些猶太人.
歸信了耶穌.
所以羅馬書是有兩方面的對象.
而羅馬書由第一章開始.
裡面是第十六節.
我們讀過的就指出.
猶太人和外邦人在基督的救恩面前都是一樣的.
什麼是一樣呢?.
就是在罪的問題上是一樣的.
猶太人不能夠自誇.
不能夠因為他們是神的子民就可以自誇.
而猶太人他們所誇的很多次是說.
他們由亞伯拉罕的年代就蒙神揀選.

$^{121}$他們是神的子民.
他們有上帝的律法.
引導他們進入神的恩典和神的國度.
當他們這樣自誇的時候.
其實這封書信的作者保羅.
他本身也是一個猶太人的領袖.
受過很嚴格的律法訓練.
他指出一個很重要的真相.
神是一位賞善罰惡的神.
無論是不是猶太人.
其實事實上是沒有真正的疑人.
是說連一個都沒有.
因為律法的功用原本就是叫人知罪.
伏在神的審判之下.
沒有人能夠真正滿足律法真正的要求.
所以律法不可以叫人真正的可以去稱義.
反而只能夠叫人知道自己做不到.
自己是一個罪人.
但在這個時刻.
就是神為所有的人提供永生的出路.
我們剛才所讀的21節開始說.
但如今神的義已經在律法之外顯明出來.
有律法和眾先知作見證.
那就是神的義直著相信耶穌.
臨到一切相信的人身上.
事實上這並沒有差別.
神的義在羅馬書裡面的意思就是.
是從神而來的賜給我們那個義.
又叫我們可以稱義.
是神賜給我們的.
而這個途徑也是公平的.
是符合神的公義.
因為無論是猶太人或者外邦人.
這個途徑都是一樣是沒有分別.
不單止在罪的問題上沒有分別.
在出路方面猶太人和外邦人兩者都沒有分別.
一條路和猶太人說的.
說到要遵行律法這種途徑是完全不同的.
不過雖然完全不同.
但是又不是完全新的發明.

$^{161}$因為經文說到.
舊約的律法和先知已經是為這件事作見證.
為它作見證的意思只是等它出現.
那現在就出現了.
基督就是彌賽亞.
這條路就是順靠耶穌基督的路.
去向一切相信的人開放.
猶太人和外邦人都是一樣.
是沒有分別.
分別是什麼呢.
他說你先是猶太人然後到外邦人.
審判的先是猶太人.
成為神的子民的又是先是猶太人.
所以這個福音雖然好像是有先有後.
但是在本質上在途徑上兩者是沒有分別的.
保羅基督說因為人人都犯了罪.
虧缺了神的榮耀得靠著神的恩典.
藉著在耶穌基督裡的救贖無常地稱議.
因為猶太人和外邦人都犯了罪.
沒有達到神創造人的時候原本的期望.
是不忠神的期望目標.
但是藉著這條新的路.
猶太人和外邦人都可以得到上帝的赦免.
恢復和神之間和好的關係.
能夠有這條這樣的出路.
是因為神對世人有百般的容忍.
在基督來臨之前是寬容人的罪.
也就是說沒有因為人的罪而毀滅這個世界.
以致有機會得到救恩.
是讓基督犧牲的死可以贖清我們的罪.
可以得到赦免.
而我們就可以憑著信心去領受這份恩典.
剛才那句話是無常地稱議.
稱議的意思就是藉著基督的救贖.
可以在神的面前站立得住.
沒有被判為有罪.
甚至是判為無罪.
是因為基督承擔了我們的罪.
而無常地去稱議.
無常的意思就是不用付出任何代價.

$^{201}$就好像我們今天處理合約的時候.
如果合約簽了之後.
如果毀約.
通常都是要一些賠償.
但有些情境就是說這件事很難搞.
如果雙方都不需要任何賠償.
這件事就容易解決.
就會採取一種無常解約.
這個大家都可以理解得到.
人與人之間在合約上無常的一種解約.
是一種協議.
雙方同意.
但無常地稱議是神伯伯賜給我們.
是因為我們根本付不起這個代價.
因為羅馬書告訴我們.
罪的公價乃是死.
一付代價就是完蛋了.
那就不行了.
況且就算付代價.
也不等於真正處理到罪的問題.
因為這個本來就是我們應該受的審判.
是應有的結果.
結果就是死亡.
只有耶穌基督作為神的兒子.
祂的死才能夠為我們贖罪.
接著說有什麼可誇的呢?.
一點也沒有.
根據什麼準則說沒有呢?.
是根據行為的準則嗎?.
不.
而是根據信心的準則.
因為我們認定人稱義是靠著信心.
不是靠著遵行律法.
難道神只是猶太人的神嗎?.
祂不也是外族人的神嗎?.
主的,祂也是外族人的神.
因為神只有一位.
既然人是不需要付出這個代價.
那就是沒有任何功勞.
所以就不可以誇了.

$^{241}$所以接著在第四章裡面也有說到.
如果人是做完工作得到報酬.
那是應得的.
那就不算是恩典了.
所以實際上.
律法是要叫人知罪.
知道自己有罪.
唯有信靠耶穌.
才可以去到稱義.
猶太人和外邦人都一樣沒有分別.
保羅再一步去到說.
如果只是好像猶太人那樣說.
要遵行律法才可以蒙神的約立的話.
這樣的話.
那這位神就只是猶太人的神.
因為祂給的方法只是猶太人才有.
但其實祂同時也是外邦人的神.
這不是說外邦人的偶像.
而是外邦人也是神創造的.
其實原本都是屬於神的.
只不過他們不知道是誰創造祂.
所以神是約立世上所有的人.
包括猶太人和外邦人.
神是所有人的神.
不只是猶太人的神.
而接納他們的是因為他們的信心.
所以既然人稱義是藉著信心.
而不是靠行為.
行為無可誇.
所以有很多人問過這個問題.
信心可不可以誇.
因信稱義.
這件事可不可以誇呢.
我就想起一件事.
大家可以這樣去想像.
這樣想像.
可能很多人有這種經驗.
就好像一個六十多歲的人.
穿高跟鞋很自然.
去複診.

$^{281}$複診之後.
醫生就說我有兩個消息告訴你.
兩個消息很多人害怕.
通常都是說一個好的一個壞的.
說誰先.
不過醫生就說.
這兩個都是好消息.
安樂了.
你隨便說吧.
兩個都是好消息.
第一個好消息就是.
今天跟你檢查.
你所有的指數都控制得很好.
可以繼續服藥.
回來複診就可以了.
那病人就問.
第二個好消息呢.
他說你下次再來複診的時間就免費了.
病人就說.
醫生你真是好人.
不得我的錢.
醫生就說不是我好人.
因為你下次來的時間.
已經可以有醫療券.
所以你不是很嚴重.
所以你給的錢不是很多.
夠給的了.
應該台下的哥哥姐姐有經驗.
病人就說.
我下次來其實可以介紹未足月.
未足月大家明白嗎.
未到生日的日期.
是不是未足月也有呢.
醫生就說.
政府只是計年份.
看醫生那天的年份.
減你出生的年份.
所以明年的時候.
一月一號就有了.
所以總之你六十五歲.

$^{321}$生日那年.
一月一號就可以有了.
那真是好.
那病人就問.
那我需不需要先申請.
到時候才來可以有醫療券呢.
醫生就說.
不用.
下次你來複診的時候.
我幫你辦好了.
你反正都要拿身份證出來.
都不用.
你來複診就可以了.
什麼都不用做.
這個病人可以免費看醫生拿藥.
未足月都可以享用.
連手續都不用自己做.
只是去醫生那裡複診.
給他治療就可以了.
不需要付任何代價.
只需要去接受就可以了.
彭遜.
接受耶穌基督.
就真的只是說.
接受這件事情.
你清楚它是真實.
清楚這不是禮物.
接受就可以了.
既然只是接受.
又有什麼可誇呢.
所以猶太人和岸邦人.
都是得到上帝的恩典.
不是靠自己的功德.
世人在神的面前.
乃是恩順稱義.
和律法無關.
如果和律法無關的時候.
如果從猶太人的角度裡面.
這讓我們借著信心.
使律法生效嗎.

$^{361}$絕對不是.
乃至反而是要鞏固律法.
那是不是律法完全沒有任何意義呢.
保羅說不是.
絕對不是.
因為耶穌基督.
要在十字架上為我們受死待贖.
正正是因為.
彰顯出律法背後上帝的公義.
律法是很有意義的.
保羅是要指出在舊約裡面.
亞伯拉罕作為信心之父.
是一個恩順稱義很明顯的例子.
如果一個人靠行為能夠稱義的話.
哪個行為好呢.
在猶太人心目中.
亞伯拉罕就非常好.
這是一個好的例子.
所以就要看舊約裡面.
亞伯拉罕究竟是因為什麼而稱義呢.
是不是因為他的行為呢.
保羅說不是.
雖然他真的行為很好.
比世上同期的人起碼是很好.
但是神約束他不是基於這一點.
保羅繼續說.
這樣我們說我們欲身上的祖宗亞伯拉罕得到了什麼呢.
若亞伯拉罕自靠行為稱義.
他有可誇之處.
可自在神面前.
他並沒有可誇之處.
基本上怎麼說呢.
亞伯拉罕相信神.
就算他為意.
上帝約束亞伯拉罕.
是在亞伯拉罕有任何行動之先.
而是單單因為他相信的.
上帝所說的應他.
就在那個時間裡面.
稱他為義.

$^{401}$所以稱義不是在於人做了什麼.
而純粹只是上帝的憐憫.
另一個例子就是貞耀.
貞耀大衛論道.
那不靠行為而被神算為義的人的福分.
他說.
過犯被赦免.
罪孽被遮蓋的人自有福敵.
處不算為有罪的人自有福敵.
八八稱義.
這個就是大衛他所經驗到.
得到蒙赦免.
得到釋放.
所以人生最大的福分.
這裡說到就是.
就是夢神的憐憫.
被他接納.
最近他看到一位姊妹的見證.
這位姊妹出生於一個書香世代.
從小就讀很多書.
讀的都是中外的名著.
在這些名著裡面.
都會有一些西方的人物.
似乎他看到一些內容.
有些人在祈禱.
願上帝保佑我.
求主赦免我的罪.
等等這些說話.
這些說話都挺有趣的.
這位姊妹當時當然未相信.
她說有本聖經看看也好.
不過她身處的地方.
不容易找到聖經.
所以又漸漸地去到.
好像亡卻了一樣.
所以長大之後.
她已經獲取了很多的尊嚴資格.
都考了.
但是當她去一個地方.
要去遠一點的地方工作.

$^{441}$她發揮這麼高的學歷的時候.
她發覺那個地方很陌生.
面孔很陌生.
突然間她覺得整個人很孤單.
很害怕.
在迷茫的時間.
很痛苦的時間.
有一個基督的朋友.
帶她回教會.
他說不如你回教會吧.
認識一下我們基督徒的朋友.
她記得她小時候.
她說看書也想有本聖經.
那就去教堂吧.
當她第一次步入教堂的時候.
她就很感受到那個氣氛.
當她回想的時候.
這個叫什麼呢.
就是感覺到上帝的聖潔.
上帝的尊貴.
感覺到自己很微小.
但是很得安慰.
因為在教堂裡面.
看到有幾個字.
有些金句.
信忘愛.
這幾個字.
覺得很親切.
覺得這個就是一直以來.
她很想有的生活方向.
她回想的時間.
進了主之後.
她進主之前都很積極.
很樂觀.
很開朗.
看事情很多時候都比較豁達.
但是當她認識了耶穌之後.
她發覺.
她所得的喜樂是完全不同的.
和以前是不一樣的.

$^{481}$她覺得自己完全是一個新做的人.
當她有聖經的時候.
她就很勤力去看.
因為從小就很隨便去看.
聖經是什麼.
雖然她小時候讀過很多書.
包括醫學.
經濟學.
文學.
心理學各類的書.
但是她感受到.
當她讀的時候.
聖經是特別的不同.
特別對她的生命.
有一種影響力.
當她很想發怒的時候.
聖靈就藉著神的話語提醒她.
愛是寬容的.
大家想到吧.
可能是前書.
愛是不計較的.
如果你不寬容別人.
想一想.
你也被神寬容過範了.
神都原諒你了.
你也原諒人吧.
所以聖經說話就是這樣去安慰她.
支持她.
給她心靈裡面那種力量.
跟她以往看那麼多書的經歷完全不同.
她回憶她很多年的生命.
無論是順境還是逆境.
她都經歷了上帝的帶領.
有一段時間.
因為她有那麼高的學歷.
賺了很多錢.
人人都覺得有房有車等等.
但其實內心有一份虛空.
覺得人生的方向都不知道是什麼.
甚至曾經經歷過生意失敗之後.

$^{521}$跌到谷底.
情緒很低落.
但當她每天讀聖經.
安靜的祈禱思想.
大閉日.
回到教堂裡面崇拜.
工作之餘.
有時間在教會的餐廳幫忙服侍.
她感受到生命裡面.
雖然外在的環境有很多困難.
但生命裡面充滿了喜樂.
有一次當她遇到一位舊同學.
跟她聊天就發覺.
她有很嚴重的心理毛病.
情緒的毛病.
甚至有自殺的傾向.
她跟家人的關係又非常緊張.
似乎姐妹跟她聊天的時間.
就跟她分享上帝的話.
當這位同學離開要回去的時候.
姐妹就送一本聖經給她.
後來這位同學又繼續回到教會.
不再想自殺的事情.
這位姐妹回想這個經歷.
是神的帶領.
是讓這個舊同學找到上帝.
是上帝去拯救這個人.
而這個舊同學也分享了.
她進入到教堂.
心靈裡面又輕鬆.
看聖經的時候.
心靈就平靜.
這些平安的經歷.
都完全不是靠律法可以有的功效.
是神的憐憫.
是神的恩典.
臨到我們.
是因著我們單純的原意去接受.
所以今天羅馬書的訊息裡面.
讓我們明白.

$^{561}$救恩是神的義.
是用公平公正的途徑賜給所有人.
沒有分別.
無論猶太人和外邦人都沒有分別.
救恩是神的恩典.
信仰是恩典.
律法,行為甚至連本身都沒有可掛之處.
單單是上帝的憐憫.
最後,什麼是福氣呢?.
是我們因著神的憐憫.
我們生命有改變.
這也是我們每一位基督徒已經經驗了.
這個福氣我們每一個信主的人都有.
而且是會一生的伴隨我們.
願意我們都能夠好好珍惜這份寶貴的福氣.
\newpage



\section{以賽亞書 61:1-4}
\label{sec:032191}
\textbf{破碎世代的城市宣教}~~~~羅泰然牧師
\newline
\newline
連結: \href{http://www.church.com.hk/acms/content.asp?site=cdc\&op=show\&type=product\&code=032191\&layout=sermon}{\texttt{032191}} ~~~~ 語音日期: 2026-05-24
\newline
\newline
\hyperref[sec:032190]{\small{< < < PREV SERMON < < <}}
~
\hyperref[sec:index_chronic]{\small{[返順時目]}}
~
\hyperref[sec:032192]{\small{> > > NEXT SERMON > > >}}
\newline
\newline
以賽亞書 61:1-4
\newline
\begin{longtable}{cl}
\hline
\hline
章節 & 經文 (和合本修訂版)\\
\hline
61:1 & \begin{tabularx}{0.7\textwidth}{X} 主耶和華的靈在我身上，因為耶和華用膏膏我，叫我報好信息給貧窮的人，差遣我醫好傷心的人，報告被擄的得釋放，被捆綁的得自由； \end{tabularx} \\ \\ \relax
61:2 & \begin{tabularx}{0.7\textwidth}{X} 宣告耶和華的恩年和我們的神報仇的日子；安慰所有悲哀的人， \end{tabularx} \\ \\ \relax
61:3 & \begin{tabularx}{0.7\textwidth}{X} 為錫安悲哀的人，賜華冠代替灰燼，喜樂的油代替悲哀，讚美為衣代替憂傷的靈；稱他們為「公義樹」，是耶和華所栽植的，為要彰顯他的榮耀。 \end{tabularx} \\ \\ \relax
61:4 & \begin{tabularx}{0.7\textwidth}{X} 他們必修造久已荒涼的廢墟，建立先前淒涼之處，重修歷代荒涼之城。 \end{tabularx} \\ \\
[1ex]
\hline
\hline
\end{longtable}
$^{1}$我想問問各位,大家知不知道今天是甚麼主日?.
猜拳主日,講到猜拳大家會想起甚麼工作呢?.
我們教會有甚麼工作?.
有些訪孫隊是嗎?.
有些內地事工.
我們過往參與海外猜拳工作,甚麼地方呢?.
印尼?柬埔寨?.
印尼就過氣了,當然過是對的,很多年了.
我們現在開始完成了,在工場工成新退.
開始尋求一個新的方向.
其中柬埔寨或泰北,遲些也有訪孫隊,大家可以參加.
這是一些猜拳的工作.
當然猜拳的定義,特別是幾十年前.
我們說的美德之民或美德之地.
是一個很重要的猜拳方向.
向著從來未聽過福音或沒有教會的群體.
一個群族,未必是一個國家.
國家裡面也有很多不同的群體.
所以我們也會覺得,柬埔寨的方向是首選.
我們可以去這些地方,做一些全方位的工作.
今天我當然是說猜拳.
但今天是說城市宣教.
城市宣教,不像我們平常所說.
是想到很遠方的地方.
或者跨文化的宣教.
或者去到偏遠的地方,我們做福音的需要.
今天我們想大家一起去放我們的眼光.
放回來我們現在身處的香港這個地方.
香港是一個很繁華的國際都市.
擁有很多美好的條件.
繁忙的金融中心.
有很熱鬧的街道,多元的文化.
這個城市看來是充滿了活力和機會.
但在繁華的背後.
卻是隱藏了一代年青人很深層的破碎.
在2019年社會事件之後.
這班年青人的心靈是留下了一些落印.
他們會有很多的憂慮,抑鬱.
甚至睡不著覺,對未來的迷惘和不安.
在20到40歲這班成青人士裡面.

$^{41}$當然一個地方來講應該是一個年青中堅的力量.
但在香港這班年青人面對很多的職場上的壓力.
在住房上,在居住上有很沉重的負擔.
家庭和人際的關係都充滿了掙扎.
還有很多的孤單感在當中.
剛才主席幫我們讀了一個以賽亞書61章1到4節.
是以色列民經歷過被擄荒涼重建的處境裡面.
才宣告神的安慰和復興.
這段經文的回應是當時有大人去被擄回歸.
回去的時候的困境.
但同時間耶穌基督出來傳福音事奉的時候.
他引用了一段經文.
宣告祂是一位受高者.
所以神是對以賽亞的呼召.
成為我們今天對教會的呼召.
而耶穌基督的使命.
也是給我們看到祂身體力行裡面的實踐和具體的榜樣.
所以今天我們要將傳人的福音帶給破碎的人.
要將城市變成神醫治的禍場.
在這樣的處境裡面.
所以城市宣教的呼召和期待是格外的清晰.
在繁華的都市裡面.
我們要看到破碎的生命.
在撕裂的世代裡面.
我們要去醫治帶來盼望.
還有在文化社會的衝擊洪流裡面.
我們要活出基督的愛和真理.
所以今天題目是破碎世代的一個城市宣教.
分開三方面.
第一方面就是被差遣到破碎的生命.
第二就是重建生命的廢墟.
第三就是興起成為城市之光.
第一方面講到被差遣到破碎的生命.
在第一節裡面講到.
耶和華的靈在我身上.
因為耶和華用膏告我.
叫我去報好訊息給貧窮的人.
差遣我去醫治傷心的人.
報告被告的得釋放被困所的得自由.
大家有沒有留意到開始是主耶和華的靈.

$^{81}$主的靈.
這個是聖靈的差遣.
宣教從來不是人自己去用自己的熱心.
用自己的方法去實踐.
你記得在上行傳裡面.
在教會建立宣教.
是聖靈在各人的中間做感動的工作.
差遣合適的人出去做猜拳的工作.
做宣教的工作.
所以宣教的源頭是聖靈的差遣.
請我們去留意.
當然我講的.
你們很快就會想到.
我講的差遣.
得到神的靈去感動.
通常是什麼人呢.
就是宣教士.
做了宣教士.
他們對一個地方有所感動.
他們接受了神學和宣教的訓練.
學習如何去適應不同的文化.
甚至學當地的語言.
然後去差遣.
去到那個地方去傳福音.
但是這個似乎和我們很遠了.
因為神沒有這樣呼召我.
我只是間中會在差遣的祈禱會裡.
聽到一些宣教士回來的分享.
就是這麼多.
在這裡講神的靈要去感動我們.
其實剛才阿加皮斯班唱的那首詩歌.
成為了這個introduction.
大家記不記得他唱什麼歌.
「呼召之地」.
裡面有一句經常說的.
「傳來召我的呼聲」.
你聽完有沒有感動.
這首詩歌是我們一起獻給神的.
你有沒有留意到.
這個也是聖靈的呼召.

$^{121}$原來聖靈是呼召我們眾人的身上.
不是說單單我們有一個海外的宣教士.
而是我們身處一個地方.
我們都有一個宣教的行動.
待會我會講一講什麼叫城市宣教.
會清晰一點.
所以這裡說的對象是什麼呢.
華本說的是謙卑的人.
修訂版就叫貧窮的人.
因為貧窮更加清楚指出.
他們的情景是缺乏的.
他們是物質上缺乏.
當然心靈上也是缺乏.
他們是受壓迫.
他們是心靈破碎的人.
我們要做些什麼呢.
目的是要猜險我們去醫好傷心的人.
報告被擄的得釋放.
在罪裡面得自由.
這是神要在我們身上.
感動我們去做的事情.
所以這個福音不是留給某些人去做.
而是我們每個人都可以.
在我們所處的環境裡面去實踐.
今天我會集中說一點.
20到40一個群體.
因為我們在近年裡.
經歷了19年的事件之後.
發現這個群體在教會裡面.
全個香港的教會都是.
是一個斷層.
因為很多的離開了.
有些是移民.
有些是信心軟弱了疏離了.
在教會裡面.
這件事情發生跟教會有些不和.
離開了教會.
但是這群群體是怎麼樣呢.
今天的光景是怎麼樣呢.
它們的外表是堅強.

$^{161}$但是內心是疲乏的.
大家有沒有看過標籤.
是這樣的.
如果有一個貨物.
很厲害.
這個我沒有叫你弄的.
標籤是什麼.
Best Year.
Handle with Care.
很好.
我今天早上就說了.
我今天早上就做了出來.
我還在想.
這個杯子有一個列行是什麼呢.
讓你猜一下.
那就不用猜了.
大家看到那個標籤了.
這個貨物裡面.
就算你去旅行.
是托運,飛機的行李.
那些東西要小心輕放.
你就會貼這個給你.
或者你收到一些貨物.
都會有一個這樣的標籤在裡面.
原來他們就是這樣.
容易破碎的人.
容易破碎,小心輕放.
小心一點處理他們.
為什麼呢?.
他們情緒困擾很普遍.
他們職場壓力很巨大.
對未來是無力感.
對信仰是疏離.
這個是有數據顯示的.
香港有20到40的年輕人.
在19年社會事件之後.
普遍他們面臨心理健康惡化.
職場的壓力加劇了.
住房的困難.
高企的租金.

$^{201}$對未來的無力感仍然充滿.
政府的青年發展藍圖.
還有2025年的統計年刊裡面.
提供了一些具體的數據.
心理健康和情緒的困擾裡面.
是按照這個青年數據平台.
有40\%的青年人表示.
在過去一年裡有焦慮和抑鬱的狀況.
最明顯的是失眠的情況上升了.
還有幸福感下降了.
他們青年人的幸福感指數.
在23到25年期間.
是持續低於香港平均的指數.
如果你不是20到40的班.
你的幸福感可能高過一點.
雖然我也不是很幸福.
但是他們還低一點.
所以跟一些20多歲的青年人聊天.
他們覺得香港很迷茫.
未來不知道會怎樣.
不知道人生的方向會怎樣.
整個環境和氛圍.
好像讓他們進入了一個.
一個充滿大霧霧景的山谷的感覺.
另外就是職場的壓力.
在這25年來講.
統計顯示青年失業率是7\%左右.
高於整體人口平均的失業率.
他們群體的工資增長緩慢.
所以生活和通脹的差距會慢慢增加.
所以他們又發覺.
職場的晉升機會不是很容易.
所以有很大的壓力在當中.
住房的樓價高企.
所以數據又這樣顯示.
家庭收入和樓價比率是不合理的.
如果你買樓置業.
你估計差不多負擔是超過你收入的幾成.
七成.
七成是很厲害的.

$^{241}$遠高於國際的標準.
租金的壓力又如何?.
租金當然是低一點.
但租金和收入的比例又佔幾成呢?.
四成.
都很厲害.
影響了他們生活的質素.
還有他們結婚的計劃.
所以對年青的夫婦.
他們面對育兒,經濟的壓力.
總是覺得生活很艱難,很困苦.
說到一些心靈的情況.
信仰和社會的狀況.
他們的信仰是疏離的.
據調查顯示.
超過一半的年青人對宗教的信仰.
是保持了距離.
同時間心理上.
他們對心理健康和靈性支持的需求.
是增加的.
你覺得好像很矛盾.
如果你覺得教會.
不要說其他宗教.
教會是一個幫助他們的心靈.
可以得到舒緩.
或者幫助他們心靈重建.
但為什麼他們對教會有疏離感呢?.
我們從教會的角度去看.
是不是我們很多的事情.
是不能夠讓他們得到幫助呢?.
這一點我們會反省.
但是一代的年青人.
仍然是很關心公共的事務.
關心某些弱勢的群體.
他們參與環保,參與公益活動都會參加.
因為這些都是能夠他們寄託.
而不再是直接的政治鬥爭.
我就想說回城市宣教.
什麼叫城市宣教呢?.
剛才我說的猜拳宣教.

$^{281}$意思就是福音去到未得之民.
或者未得的地方.
不論是一個小的群體.
或者某一個小的地方.
從這個概念開始.
你們想起.
就算是一個城市裡面的結構都很複雜.
城市裡面有很多不同的階層的人.
貧富懸殊的城市.
有更加的看得到不同的階層.
不同的需要.
或者不同的城市裡面.
有不同的少數民族.
香港也有很多不同的少數民族.
在香港的生活.
還有一些.
按照職業來說.
那些我們不容易接觸到的人.
他們未必在教會開放的時間裡面.
可以參與到教會.
很多人星期天要上班.
很多人可能平時要上班.
可能才是.
就是晚上.
甚至有些是通宵的.
是沒辦法接觸到的人.
這些都是我們可以說.
城市宣教對一個大的宣教的縮影.
同樣有一些基層的人士.
或者某些階層的人士.
某些少數民族的人士.
某些年齡層的人士.
是我們接觸不到的.
不容易接觸到的.
教會不容易接觸到的.
這個就是我們說城市宣教的需要.
城市宣教的核心就是說.
在這個都市.
文化,社會,經濟,政治等等的.
不同的情況裡面.

$^{321}$在教會的角度.
我們要活出基督的見證.
帶來一個屬靈的更新.
還有一個社會轉化.
所以這個是一個.
全人生命的一個使命.
所以.
先不要說怎樣接觸他們.
而在說什麼途徑去做呢.
剛才所說的.
以賽亞在這個時代裡面的人.
他們的心靈是破碎的.
所以在靈性的層面裡面.
是在破碎的世代裡面.
是要求神去憐憫他們.
是推動敬拜,禱告.
在社群的層面裡面.
我們要進入社區裡面.
回應人群的需要.
所以我們有外展的工作.
福音御公的工作.
去展現基督的愛.
這個都是我們要做的事情.
在文化的層面裡面.
今天我們更加需要.
用創意和科技的方法去傳揚福音.
最近清善都跟我們說.
我們會不會有一個Instagram.
我們可以分享我們的見證.
因為現在年輕人.
是不會用Facebook的了.
大家知道嗎.
還有有些人都未必用WhatsApp.
去傳遞的.
直接用Instagram裡面的訊息.
去傳遞那些訊息.
因為經常看那些東西.
就說不如那裡傳遞一些訊息給你.
所以我們都要知道.
他們的情況.

$^{361}$但是這個文化層面.
這個接觸.
都是要日新月異.
我們要看得明白.
因為這裡裡面說.
我們要走進他們生命的當中.
去到這一群人裡面.
去進入他們裡面.
甚至正是要我們去接觸.
這些貧窮的人.
所謂貧窮的人.
是心靈脆弱的人.
這個心靈需要的人.
在二十到四十個群體裡面.
就是這個情況.
所以在成雪宣教的第一步.
是要看見破碎的人.
神要差遣教會成為一個依治者.
要去幫助人.
不是在某一些的事工.
某一些可能是日常的.
每週的一些事工去做到.
我聽到一個故事很感人.
教會裡面有一個初信者.
回到教會短短時間.
但後來發覺患了癌症.
癌症身體越來越軟弱.
一開始都覺得.
是一個比較後期的癌症.
是挺困難的.
開始沒有回教會.
但有一位弟兄知道.
他差不多每天去探訪他.
因為這個患了癌症的初信者.
他一個人生活.
不是很大年紀.
四五十歲.
這個弟兄很關心.
很願意去幫助他.
知道他沒有什麼人可以幫忙.

$^{401}$去煲湯.
去跟他聊天.
如是者一段時間.
但為什麼後來才發現.
這種愛心默默的可以體驗出來.
後來這個患癌症的弟兄竟然好起來.
竟然可以再回教會.
才知道原來一直有個這樣的弟兄.
默默的去幫助他.
去關心他.
我相信這個弟兄是傳遞了.
一個很重要的良藥.
喜樂激心乃是良藥.
讓這位患癌症的弟兄不是孤單.
不是無患.
而是孤寂中等死.
帶給他有信心.
帶給他有盼望.
帶給他有喜樂.
他就好起來.
當然這個初診患癌的弟兄.
都有看醫生.
有接受治療.
但我相信這位弟兄所付出的.
是神的恩典.
所以我們需要更多的人可以起來.
好像二賽瓦華爺爺說的.
「靈在我身上」.
好像詩班唱的詩歌.
我聽到那些呼召的聲音.
我們就去實踐.
城市宣教就是走進他們生命的裡面.
第二方面就是重建生命的廢墟.
這段經文裡面.
你會看到那些人.
真的是很廢的.
很慘的.
宣告耶華因憐.
我們向神報仇的日子.
安慰所有悲哀的人.

$^{441}$為神安悲哀的人.
賜華冠代替灰燼.
喜樂的油代替悲哀.
讚美唯依代替憂傷的靈.
稱他們為公義樹.
是神所栽植.
要彰顯神的榮耀.
你看到那些人.
很慘的.
又悲哀又灰燼.
即是頭.
毫無一種生氣.
又憂傷的靈.
這些.
是好像打敗仗的感覺.
這些是在這個.
世代裡面.
好像不被認同.
不被接受.
或者很多的失敗挫折.
在當中.
這些人的破碎的心.
當然剛才我說.
這些二十歲四十歲的年青人.
當然不會被你看到.
外表仍然很堅強.
但是仍然標貼了一個標籤.
Vegeta.
很容易碎的.
我們要看到這些人.
聖經怎樣說.
我們在這裡要做些甚麼呢.
是要宣告耶華的因年要臨到.
甚麼是因年呢.
因年是神自己恩典的日期.
日子.
當然這是代表一個釋放和更新.
罪裡面得到釋放.
生命裡面得到更新.
這些人可能是在罪的當中.

$^{481}$這些人可能是在壓迫的當中.
這些人可能是在困苦的當中.
但是他們不能夠抬起頭來.
他們的心靈是破碎.
他們很多時夜間.
夜闌人靜的時候.
他們的心靈是充滿悲哀.
但我們怎樣去做呢.
華官代替灰燼.
起落的油讚美的衣.
這裡似乎好像給一些很實質的.
好像衣著給他們看到.
這裡當然是這樣形容.
其實意思是華官.
當然可能是一個很好的頭飾.
是代表身份上的恢復.
讓他們從羞辱裡面變得尊榮.
我跟他們分享過.
之前我看一些叫做.
甚麼全民造星的節目.
你們有沒有看的.
你們不看ViuTV的.
全民造星有很多年輕人.
有時是男的有時是女的.
你會發覺他們背後.
最好看不是看他們唱歌表演好不好.
其實那班人背後是有很多破碎的.
他們裡面是有很多艱難的.
在社會裡面工作是不如意的.
或者讀書又不是理想的.
但是他們借了一個機會.
能夠經歷一下實踐自我.
或者是一個突破自己.
所以他們看到成功的時候.
他們淚流滿面.
我會覺得很多年輕人是這樣的.
我們教會都需要這樣的情況.
讓這些淚流滿面的年輕人來到我們當中.
是可以感受到我們對他們的愛.
突破和支持.

$^{521}$即使不是給了他們什麼.
但是他們能夠體會到.
得到尊重.
他們在舞台裡面是得到認同.
雖然是有競爭有贏輸.
但是他們能夠完成舞台一次.
已經是覺得是一個成就.
今天我們是很需要.
給我們的一代新一代年輕人.
是這樣感受在裡面.
後期教會是充滿愛和恩典.
我覺得更加容易讓他們感受到.
但是我們沒有這扇門打開.
所以華官要代替灰燼.
羞辱變為尊榮.
喜樂的柔.
喜樂的柔是代表一個.
心靈情緒裡面的醫治.
我們去國內去關顧.
有時去外展.
我們都會帶一瓶柔給人的.
活絡油.
給老人家很適合.
我現在都用的.
活絡油塗一下筋骨.
都精神爽利.
這裡說的喜樂柔.
是說心靈裡面的柔.
表面上是需要.
身體上是需要.
心靈更加需要.
這裡你看到.
我們會不會成為喜樂的柔呢?.
是在他們當中.
好像那個弟兄那樣.
常常陪伴他.
他就好像每天跟他.
高抹他一樣.
用那些柔.
讓他的心靈得到復甦.

$^{561}$以致有喜樂的心靈.
身體都能夠恢復起來.
讚美為醫.
在這裡說.
神是去更新人內在生命.
怎樣去更新內在?.
只有神才能夠更新我們生命.
我們怎樣去讓神更新我們生命.
一定放在敬拜的裡面.
敬拜是更新我們生命.
你來每一次崇拜.
是否更新你的生命呢?.
你會說不知道.
看看哪裡說得怎樣.
是不是這樣呢?.
當然這個敬拜讚美.
你自己有沒有去開口敬拜讚美呢?.
但是我沒有這樣的氛圍.
我們唱詩詩歌熟不熟呢?.
但當然這個也是提醒我們.
能夠在這個方面幫助人.
有美好的敬拜的生活.
重點這個重建生命的廢墟.
就是要將羞辱變為尊榮.
悲哀變為喜樂.
要讓一班二十到四十流失很高的人.
在這個教會裡面找回歸屬感.
得回這個同行.
剛才我們開始還在唱一首歌.
那首歌是唱什麼歌呢?.
《使我作你和平之子》.
我們唱完就忘記了.
《使我作你和平之子》.
那個詞是誰作的呢?.
有詞才有歌.
是法蘭西斯.
或者天主教裡面是聖方濟各.
聖法蘭西斯的土文.
很出名的.
他是個什麼人呢?.

$^{601}$他是個有錢人.
但他竟然一次經歷到神的帶領.
他覺得他要關心身邊貧窮的人.
特別有一次他騎著馬.
在很窄的路裡面.
見到一個麻風的病人.
在那裡好像有一種期望.
他就很掙扎.
因為麻風病是很容易傳染的.
他掙扎.
那時他已經信了主.
究竟馬就快些打聲.
過了他先.
還是落低.
看有什麼需要.
他掙扎底下.
最後就是落了馬.
靠近他.
麻風病人.
可能他需要一些式文不奇.
但更加需要什麼呢?.
就是人的接納.
因為那時的人是扔出去的.
讓他自生自滅.
好像已經被社會拋棄了.
法蘭西斯當時很感動.
受著聖靈的差遣.
走上去.
勇敢地抱著他.
他就說.
我抱著的時候.
我看到他的面容.
那個面容不是誰?.
是耶穌基督.
所以他就從此就開始建立.
在一個服侍裡面.
專做貧窮人.
聖方濟各.
就是今天在大公教會裡面.
做這些關懷的工作.

$^{641}$所以我們在這一刻.
這套經文裡面.
就是要我們成為和平之子.
我們要去關心我們周邊的人.
在困苦絕望裡面.
給他們有希望.
我們需要這樣去做.
二十到四十歲的年輕群體.
是容易破碎.
無論是未信主的非信徒.
或者是已經信了主的.
在我們教會裡面.
我們都會看到他們裡面的情景.
當然我們在一個牧羊裡面.
近年我們推行小組的牧羊.
是用更少的人數去彼此關懷.
職稱裡面我們已經開了幾個小組.
人數一個組不會很多.
不會像以前團契那樣.
二十多人三十人四十人.
但是團契人多了.
後來都會歸回人少.
因為關顧不了.
那些人失落了.
每個團契都變成小組.
但是仍然用團契的模式.
但是現在的小組是不同的.
是一個完全彼此的牧羊.
還有關心.
還有在神的說話裡面成長.
我們用一個小組的牧羊.
更多的同工的接觸.
關懷.
彼此的接待.
門徒的訓練.
推進他們能夠建立.
這個屬靈的知識.
所以城市宣教的靈性牧羊.
是包括了一個敬拜的運動.
剛才議員經文說.

$^{681}$安慰所有的人.
悲哀的人.
這些有需要的人.
以讚美為意.
代替了這些受傷的靈.
所以我們有更新的崇拜.
有美好的敬拜.
今天我們也準備開一個.
二十到四十的成青人士的崇拜.
最快會一月末開始.
我們是會去讓一班人.
能夠有歸屬感.
有認同.
有彼此同行.
更加需要更.
我們要很多破碎的人.
心靈的年青人可以來到的地方.
是得到他們神的關心.
得到從神而來.
那個力量.
就好像那些傳聞造星.
這些人破碎.
淚流滿面.
來到這裡得到一個認可.
認同.
所以這個敬拜我們是需要.
今天我們也在發問卷.
問一下他們.
究竟什麼時間最好呢.
禮拜日下午三點兩點.
禮拜六黃昏晚上八點.
或者禮拜五晚上.
讓他們自己看看.
覺得什麼時候最適合.
我們一路去看看怎樣去.
更加的合適的時間.
有美好的敬拜.
我們更加的能夠去敬拜裡面.
我們大家都知道.
早前月會已經通過了.

$^{721}$我們將樓下小禮堂重新裝修.
昨天我們有機會和建設顧問.
在裝修上談.
這個叫做黑盒劇場的做法.
讓這個敬拜裡面.
我們用台燈影的.
更加美好的方式.
能夠去配合詩歌.
配合廣道.
能夠一體化地去.
集中讓人去敬拜.
我們過去兩年.
青城部去到台灣的交流.
我們都體會到這種敬拜.
我自己都曾經去過一兩次.
三次.
我都覺得我很能夠投入當中.
雖然我不是年輕人.
但我自己都很投入當中.
因為你會很專注在.
那個訊息的部分.
詩歌的部分.
讓我們更加可以.
現在很流行說的.
沉浸式的一種敬拜模式.
我們當然是在學習.
我們未必做得很好.
但我們要嘗試.
我們做好這件事.
目的是要養這些破碎的心靈.
更加觸摸到神自己的心靈.
另一件事.
這個城市宣教裡面.
那個屬靈的策略.
是在悔改認罪的禱告裡面.
所以我們更多要推行悔改認罪的禱告.
因為如果我們不是去潔淨心靈.
不是我們去嘗試認罪悔改.
我們在坦然無拘來到神的面前.
我們不能夠抓緊我們跟神的關係.

$^{761}$因為很多時候.
剛才我說過他們的信仰是疏離的.
一方面他們好像覺得不受認可.
不受認同.
是不是有些事做得不好.
他們心裡面都有很多的.
是很複雜的想法.
但是在神的面前.
可能也不能夠坦然無拘.
但是人仍然很容易.
在一個城市裡面.
我們知道城市有很多誘惑.
城市有很多很多的事情.
令到我們分心.
所以我們很多時候.
我們一是不專注在神那裡.
一是我們受到誘惑.
會跌倒.
因為我們好像有很多的荊棘.
纏繞著我們.
不能夠去親近我們的神.
很多的問題會令到我們這樣.
但是這個認罪悔改的禱告.
來到神的面前.
我們每次去清潔自己.
每次去反省我們自己.
讓我們能夠可以.
控制我們跟神的關係在裡面.
所以這個我們要推動.
我仍然是呼籲大家.
來到禮拜五早上.
認罪悔改的祈禱會.
我要感恩.
執事和公一起去帶領.
我們所謂教會領袖.
已經是帶頭的了.
不是我們你們做吧.
我們在後面支持你.
我希望你們都能夠參與.
你參與除了為你自己.

$^{801}$是為整個教會.
我們教會裡面需要.
常常有人為我們整個群體.
去認罪悔改.
你有沒有記得約伯.
他每天都為他的子女.
去禱告獻祭.
他都希望.
能夠可以將最好的獻給神.
我們都需要這樣的行動.
我需要你們做這件事.
你跟我沒有什麼關係.
主的靈在我身上.
猜全工作.
我不會去到海外.
你可以回來和祂禱告.
我需要這些人.
當然有靈修的資源.
有這些聖經的學習.
我們現在有更多的課程出現.
大家知道有很好的神學院的老師.
教授我們.
除了主學.
我們還有其他培訓的機會.
都可以參與.
待會我再報告分享.
第三方面.
興起為城市.
成為城市之光.
他們必修造久已荒涼的廢墟.
建立先前荒涼之處.
重修歷代荒涼之城.
這個城市可能是說耶路撒冷.
歷代這麼長的時間.
已經是荒廢了.
你留意到一開始.
不是用神做主體.
剛才是注定靈做神做主體.
現在用他們做主體.
這裡說什麼呢?.

$^{841}$不是神自己重建.
是你們的子民重建.
興起他們的子民重建.
當然我們要重建這一代的人.
他們要興起自己.
裝備自己.
能夠在城市裡成為城市之光.
在教會裡成為一個明亮的燈台.
所以看到永久久矣荒涼的廢墟.
很長時間代替了.
一代一代經歷破碎.
特別是他們被擄七十年.
經歷了幾代的時間.
荒涼的地方就是城市.
但不只是地理上.
是整個的生命,家庭,社會,文化,宗教等等.
都荒涼了.
但是教會在城市裡是怎樣呢?.
這個群體.
是要活在城市中間的教會.
因為我們有活潑的道.
我們有活潑的生命.
但這個不是留在四面牆裡.
而是我們能夠在城市裡.
發出成為光和炎.
教會不是去逃避城市.
也不是給姐妹們回來.
教會是在逃避.
在城市裡受到的艱難.
反而我們是被差遣進入城市裡.
我們是重新得力.
我們被興起,我們被造就.
所以我們盼望我們能夠.
有一個崇拜小組的模樣.
有一個禱告的開始.
我們是要建立一個群體.
他們這班人成為小組的領袖.
成為職場的宣教士.
他們要成為他們家庭.
年輕家庭的牧者.

$^{881}$他們有年輕的,細小的兒女.
他們要能夠有很好的建立的家庭.
至少用神的話去培育他們的小朋友.
他們成為社區裡的服事者.
他們很喜歡參與社區的服事.
他們能夠參與.
他們成為城市文化的見證人.
用Instagram.
用其他的平台.
去分享他們的見證.
所以城市宣教是社群影響力很大的.
所以我們要看到.
我們要能夠去到這件工作裡.
嗯.
作為一個教會.
是一個開放的教會.
跟這個城市.
跟這個社區是連繫.
我們應該.
這道門.
應該有很大的間隔在裡面.
嗯.
今天.
很多的教會.
一九年之後.
都做了很多的.
共享的空間.
開放的教會.
或者你在家附近.
都遇到一些教會裝修了.
見到他們有些不同的情況.
是容易些進去的.
是容易些.
好像歡迎你進入當中的.
嗯.
這個是.
城市宣教是要進入社區.
或者教會跟這個社區裡面.
不應該有道.
就是.

$^{921}$嗯.
什麼.
一種很大的阻隔裡面.
嗯.
好像我們尖浸.
我不知道大家怎麼看.
是不是平時過來.
一看好像暗暗的.
好像是不是.
是不是還開門的呢.
或者是方不方便進來的呢.
嗯.
很多人都想找個地方去坐下.
休息下這樣.
最近我去參加一個教會.
他開了一個咖啡室.
嗯.
他真的是免費招待的.
他一進去裡面.
有些什麼.
可能是在做.
那些叫什麼.
嗯.
送貨那些呢.
就是物流那些呢.
那些快遞那些呢.
休息一下也好.
或者有些人呢.
是會用一些時間.
好像去到咖啡室裡面.
上下網啊.
做下事啊這樣.
那他們都是一個.
好的喘息的地方給他們.
我們教會會不會有這樣的機會呢.
我們小禮堂.
我們輔堂都開放了.
但是慢慢要前面.
要光芒的.
慢慢前面呢.

$^{961}$是要去看到一個需要.
大家知不知道.
不是很多人知道.
我們就是一至五.
其中禮拜四就停了.
十二點到兩點.
我們開放輔堂給人家用.
小禮堂.
但是你們也未必知道.
我們就很.
當然我們設了一個.
熱拉架給人家看到.
但是人家看到.
未必是很敢進來的.
因為還要按個鐘啊.
是不是真的可以進去的.
所以就不是很多人會進來.
但是如果我們很開放.
很清晰,很光明就很容易進去.
我們曾經去過一些教會.
是很容易.
你自然進去都不知道的感覺.
這種情況.
因為有很大的空間.
走著走著去休息.
可以上網.
這些就是一個共享的空間.
但是這些需要我們的心靈廣闊.
需要神的靈去激勵我們做一些事情.
是我們需要付出的.
我們不是說人進來我們就算了.
我們還有很多的配套.
很多的部門.
我們還要想很多的創意.
怎樣去讓人來到中間的時候.
更加有讓他們認識的聲音,訊息.
我們不只是打開門就完了.
我們需要很多義工.
所以今天都很開心.
聽到說如果將來.

$^{1001}$我們都有個咖啡局的時候.
有位姊妹姐跟我說.
我準備學一下.
怎樣做好那部機器.
用咖啡出來.
我認識一間教會.
他們招請一個童工.
但是他有幾十個義工.
因為都要很多幫手.
不是說義工來一整天.
可能一個小時,兩個小時.
半個小時.
都是成為了很不好的幫助.
在當中.
這些是需要我們更加的參與.
或者我們有機會再舉辦一些音樂聚會.
歡音的餐會.
更加可以接觸那些人.
我剛才說.
希望一月在這個地方有一個崇拜.
我們童工已經在想.
一月之前如果裝修好.
我們可能有Alpha Cross開始.
讓你們帶領一些微信的朋友.
藉著有晚餐,有詩歌.
Alpha Cross是一個整體的.
一個福音的完整的學習.
有一套片去看的.
很現代.
我們有很美好的一個電視幕牆.
讓你們看到一個很好的情況.
但需要我們有這樣的心.
需要我們能夠參與當中.
城市宣教成為我們生活的方式.
所以需要大家奉獻.
奉獻你的身心靈.
奉獻你的時間.
奉獻你的才幹.
奉獻你的金錢.
因為我們遲些都會說.

$^{1041}$這個工程是需要費用的.
都需要籌集.
所以我們每一個人都可以.
在這個破碎的世代裡面.
成為一個城市宣教的事奉者.
希望我們教會.
更加在我們這個世代裡面.
在香港這個地方裡面.
能夠幫助更多有需要世代的人.
我今天是集中說24/20的群體.
因為我們現在集中在工作那裡.
但不代表其他人不需要.
同樣當這個教會開放的時候.
任何的人士,任何年齡層.
任何階層的人.
我們都會歡迎去進度當中.
感受到基督的愛.
還有那個力量在當中.
去同心祈禱.
因為你自己的靈.
是要在我們的身上.
你擁有高我們.
要我們能夠有口才去傳福音.
要我們能夠有力量去為主作供.
要我們能夠同樣.
首先我們致潔.
被你所使用.
直至我們真的看到.
這個你悅納人的恩憐.
讓我們配合這個時代的需要.
趕上這個腳步.
我們教會成為你合用的器皿.
禱告.
感恩奉耶穌命給.
阿們.
\newpage



\section{以弗所書 2:10}
\label{sec:032192}
\textbf{建立健康的使命教會 (15) 得發揮的恩賜}~~~~吳磐德牧師
\newline
\newline
連結: \href{http://www.church.com.hk/acms/content.asp?site=cdc\&op=show\&type=product\&code=032192\&layout=sermon}{\texttt{032192}} ~~~~ 語音日期: 2026-05-24
\newline
\newline
\hyperref[sec:032191]{\small{< < < PREV SERMON < < <}}
~
\hyperref[sec:index_chronic]{\small{[返順時目]}}
~
\hyperref[sec:032193]{\small{> > > NEXT SERMON > > >}}
\newline
\newline
以弗所書 2:10
\newline
\begin{longtable}{cl}
\hline
\hline
章節 & 經文 (和合本修訂版)\\
\hline
2:10 & \begin{tabularx}{0.7\textwidth}{X} 我們是他所造之物，在基督耶穌裡創造的，為要使我們行善，就是神早已預備好要我們做的。 \end{tabularx} \\ \\
[1ex]
\hline
\hline
\end{longtable}
$^{1}$我們在講道之先.
我們再一次同心低頭祈禱.
創造天地萬物的神.
創造我們生命的主.
我們求天父你的靈.
現在向我們說話.
求你感動我們.
幫助我們.
看見主你對我們的創造.
又明白神你對我們的呼召.
讓我們知道怎樣去回應神你的大靈.
我們將以下.
沿讀聖經話語的時間.
恭敬的有望給天父.
聖靈求你自由的運行.
幫助我們.
去明白真理.
又活出真理.
祈禱感恩.
奉主耶穌名求.
阿們.
所以當時就說2022年世界盃阿根廷表現出色.
球員發揮得好好.
什麼叫做發揮得好.
你想想當時原本比賽時間.
他只能夠三對三打平手.
一個悶局.
但加時竟然讓阿根廷四比二.
就在加時的時間贏了.
所以其實真的需要臨場的發揮.
包括技術的發揮.
就是推進形點.
就是個別球員的技術.
身體的發揮就是個別球員跑得快不快.
接球準不準.
心理的發揮就是輸了球.
和這麼緊張.
當你四比二加時的時候.
那種大的壓力底下.
他怎麼能夠保持一個心理狀態.

$^{41}$可以贏得球.
頭三樣都是個人的.
但是最後那兩樣是團隊的.
如果你有多好的技術.
多快也好.
如果你不能夠在應該傳球給人的時候沒有傳到.
又或者你應該在帶球的時候.
你又停下來慢的時候.
團隊就發揮不了.
教練也有戰術的.
你也知道一打到加時就經常帶球員回來.
說一下戰術是怎麼樣的.
可能個別球員也會覺得教練說得不對.
但是他還是要信服.
然後大家發揮到團隊和戰術的發揮.
所以當我想起世界盃的時候.
雖然我不是一個球迷.
我哥哥也是.
他從小到大都是球癡.
所以我有一點黏到.
不是球癡而是黏到一點.
一些足球的東西.
但是我想想和教會真的有很大的關係.
你想一想.
神秘教會每一個人就是你和我在座之中.
都有一份禮物.
上兩個星期江峰才和大家說.
這份禮物已經給了你.
你需要去發掘.
需要去訓練.
中文是訓練.
去培育他.
然後去發揮擺得對的位置.
其實球員也是.
如果擺得不對的位置就很笨.
所以當你看這幅圖畫的時候.
我不知道你自己覺得在這個球隊裡面.
恩泉大家庭就好像一個恩泉大球隊一樣.
這個大家庭的身體.
你就是其中一份子.

$^{81}$你是否能夠盡情發揮所賦予你的潛質呢?.
你怎麼回答這個問題呢?.
你覺得自己懷才不遇呢?.
你覺得自己是什麼狀態呢?.
或者你正在做的是不是已經發揮到盡情的發揮呢?.
其實要回答這個問題.
首先你要有一個更深的認知.
就是究竟你是不是神獨特的創造.
大家怎麼回答這個問題呢?.
你覺得是的話就舉手吧.
我們大部分的人都會覺得是.
但是有時候我們會覺得不是.
為什麼呢?.
我們覺得神是獨特的創造.
但是我好像在這裡沒什麼用呢.
或者和別人不合.
或者我的才幹也不一定適合在這裡.
各種很多的原因.
究竟神的心意是什麼呢?.
我的生命有什麼意義呢?.
神究竟把我在這裡做什麼呢?.
可能有時候也會有很多這些疑問.
甚至我可以怎樣給神用呢?.
你也會有這些疑問.
但是你回答這個問題的重要性就是.
其實沒有人可以代替你.
這是關鍵.
因為你是獨特的.
所以就算你覺得自己很普通.
有人比你厲害.
就算你覺得你只是羨慕別人的恩賜和能力.
覺得自己沒什麼恩賜沒什麼能力.
又或者你不知道我們怎樣可以在教會中給神使用.
你覺得自己很普通.
但是我想邀請你的就是.
透過今天的經文.
神已經向我們說話了.
提醒我們的.
不要埋沒我們自己的恩賜.
因為我們的存在不是偶然.

$^{121}$而是神特別的設計.
所以如果要發揮.
其實一隊球隊.
輸球就因為沒有發揮得好.
臨場發揮得不好就輸球.
其實教會也是.
為什麼歷世歷代的教會.
我們要去成就神.
其實給我們的大使命也有缺陷.
就因為當我們不是能夠.
全部的教會所有的人.
都能夠臨場發揮到我們團隊.
發揮到你在當中最適合的角色.
發揮是將能力精神意志.
充分的表現出來.
也將內在生命的性質素和能力表現出來.
換句話來說.
發揮這個字的意思.
就是原本已經有的東西去發揮.
你沒有的東西就不能發揮.
你有的東西.
譬如你這首歌.
你看他學琴.
彈到八級.
他表演的時候.
他要那八級才行.
你看我現在彈不出來.
我來來去去都是一級都未考.
雖然以前是一首歌的鋼琴歌.
演奏皇帝.
八歲的時候.
但現在連那首歌都不記得怎麼彈.
所以你原本有的東西.
你不發揮不去發展.
不去栽培.
就不會發揮出來.
而且你發揮是要釋放你有的東西.
因為已經存在你生命裡面.
而且是產生一個合適的果效和影響力.
你願不願意回應神.

$^{161}$因為其實保羅在今天一節的經文裡面.
是附著一個很重要的神學觀念.
是發表一個邀請.
也是一個命令.
我們是他所造之物.
在基督耶穌女造成的.
就是要我們行善.
就是神早已預備好要我們做的.
如果你覺得這段經文.
你看過很多次了.
很熟悉了.
但如果你有一個意念.
就是我們是受造之物.
那當然了.
我們每個人都是受造之物.
常常覺得我們很特別.
很重視我們.
但為什麼我們仍然不去發揮最佳狀態.
在一個團隊裡面.
因為我們覺得.
神都做他.
又做隔壁那個.
又做另外的人.
不一定要我.
所以要回答的問題是.
你是獨特的.
你的重要性在這裡.
就是他所造之物究竟在說什麼.
在基督耶穌女造成又是什麼意思呢.
其實原文裡面我不會花太多時間去解釋.
如果你現在上網Google.
都可以找到很深入的希臘文的釋經.
但最重要的意思就是.
我們神所創造的創造這個字.
英文的字是poem.
那個字根是poema.
在希臘文裡面.
其實就是神的創造.
是一個傑作.
更加精準的翻譯就是.

$^{201}$你是神的傑作.
這樣就很不同了.
傑作的意思就是說.
好像詩歌.
好像詩詞.
poem是怎樣的.
詩詞是需要.
可能是對聯.
可能是要有些各種的.
寫詩當中的詩情畫意.
裡面一些的安排.
所以你的創造.
你的創造你.
你所造之物.
你是神的傑作.
是獨特的傑作.
這個傑作是沒有人相同的.
而在基督裡面造成的意思.
在原文裡面其實是說.
是為使命去被創造.
你這個傑作.
為什麼神要這麼多心思.
這麼多時間.
這麼多的精準的傑作去創造.
去塑造你的生命.
就是為了基督成就了救贖的恩典.
基督的使命.
所以你是為了基督的使命.
而是神操作你成為他的傑作.
既然他已經這麼花時間精神去做了這個傑作.
他就已經.
這個關鍵在這裡.
那個字是prepare beforehand.
就是說神已經預備了你要走的.
不是說到時候再想一想.
做了出來再想一想.
怎麼做呢?.
放在哪裡呢?.
不是.
他做的時候已經有一個目的.

$^{241}$在想.
他之前未創造你的時候.
他已經預先為你預備了道路.
你的回應就要行在其中.
行這個字很重要.
行善這個行字.
因為在原文裡面說的是一個生活方式.
不是一個聚會.
不是一次的服事.
不是一件的事情.
而是一個事奉的方式.
生活的方式.
換句話來說.
你是神的傑作.
你是為了祂而活.
而你放在哪個位置.
你怎麼去服侍祂.
祂已經預備好給你的.
你的責任就可以找出來.
要找到.
上個星期.
星期一.
5月17號.
我又開心.
但又帶著一個不開心的時間.
為什麼呢?.
因為剛剛我兒子.
開心的事情就是.
5月17號他在UC Berkeley畢業.
讀Ethnomusicology.
他拿了博士學位.
研究人類學和音樂.
當中對文化的影響.
他那天畢業.
我不開心的地方是什麼呢?.
就是剛剛錯過了上個星期的退休營.
因為我也知道.
當中我也期待神要在當中.
做很多奇妙的工作.
但因為我兒子在2,3月才告訴我.

$^{281}$就是那天.
學校出了一個計劃.
所以撞到正.
但我想引用他的故事.
跟大家訪訪.
這個恩賜的發揮.
他今年.
不告訴你他的年紀.
不過你數一下.
他的生肖.
是第三次.
完成了第三次.
所以他現在畢業了.
也有一個兒子.
但是.
在26年前.
2000年的時候.
我是第一年回香港信眾教會.
做傳道.
當時是第3個月.
因為我是9月加入的.
12月.
牧師就叫我說.
你說聖誕節到了.
其實很大壓力.
在大教會裡要講道.
腳都軟了.
那天講的道.
就是今天這節經文.
不過就肯定.
不是和那天講的道一樣.
但都是那節經文.
真理不變.
但內容演繹有些不同.
這張圖在我心裡.
常常看到也想送給大家.
就如上兩星期.
剛才講的禮物.
當時我講神送給你恩賜的禮物.
聖誕節.

$^{321}$神送給我們最大的禮物就是耶穌基督.
而耶穌基督就送了.
十零的恩賜給你.
這就是你生命的禮物.
當時真的很多弟兄姊妹配合.
因為我祈願土.
想用這個比喻.
於是我就邀請弟兄姊妹.
他們很好,幫我想想怎樣去做POP子.
特意有些做木工的.
就搭了一個架.
就是這個禮物架,是木的.
為什麼呢?因為崇拜開始的時候.
我就要兒子10歲.
進去裡面,捐進去裡面.
想想就把他捉住.
你知道我講得很長.
把他捉住就是這麼長的時間.
10歲了當時.
所以一定要夠高夠穩.
不會推著推著,不穩的時候推倒.
還有就夠空間讓他呼吸.
所以這些都是恩賜.
彼此的配搭.
才發生得到的那個禮物.
然後在講到結尾的時候.
他就出來了.
就彈手敲鍊.
因為在他10歲.
能夠彈的手敲鍊已經是最好的.
他將最好的貢獻給神.
不過就是.
不會現在彈的.
大人的時候就不會這樣彈.
但是最精彩的地方.
就是恩賜發掘.
就是你要夠膽用.
神就夠膽用你.
因為我女兒.
大女兒當時.

$^{361}$大女兒.
比我兒子大五年.
是五年的.
所以就大姐姐.
就想想怎樣可以當她走出來的時候.
她帶著另外兩個妹妹.
一個就好像六歲.
和四歲我忘記了.
四歲.
她彈什麼呢.
但是姐姐很好.
為了他們能夠四重奏.
就是最小的那個.
咦咦咦咦咦咦咦咦咦咦咦咦咦.
這樣就搞定了.
於是有四重奏.
就變成了這個音樂.
我想說的就是.
這其實是一幅最美的圖畫在我的生命裡面.
也想送給大家.
其實教會裡面就是這樣.
不是要最好的.
而是神要你當時的最好.
不是別人看為最好.
而是神在你生命裡面.
你當時的最好.
並且不要真的在禮物盒裡面.
把自己焗死.
你出來透透氣.
用一用.
十歲當時.
我們剛剛回香港.
她就因此在這個盒子裡.
發揮了出來.
但我不知道當時好像一個寓言一樣.
七年後.
她就參加一個音樂比賽.
這個音樂比賽.
就.
是一個Tom Lee的音樂比賽.

$^{401}$是成年人的.
怎知道我做爸爸當然告訴他.
兒子明天就算輸了.
也不要哭.
人家那些人個個都比你大.
個個都二十多三十歲.
你才十七歲.
已經預備了他的心情.
做一個好的父親通常都是幫助子女.
知道如果失敗也不用哭.
但竟然神真的很奇妙.
在那天他就拿了獎.
所以South China Morning Post.
南華早報就訪問他.
當時他說了一句話.
對我現在也很深刻.
其實他想藉著音樂去改變世界.
我心想.
傳統華人父母.
音樂是賺不到錢的.
可能要老婆養你多過你養他.
但是.
這個就是他的心智.
所以他中間.
他也經歷過.
又轉讀新聞學.
又轉讀社工.
最後就是人民學家音樂.
其實走了很長的一段歷程.
他太太也很支持他.
一直走這段路.
所以他的生命.
也被神使用.
他在教會服侍.
我記得他在UT讀碩士的時候.
教會在Downtown.
去一些酒吧傳福音.
去講一些音樂文化.
透過他講完.
牧師又可以.

$^{441}$傳福音.
幫他傳福音.
原來神可以這樣使用每一個人.
當他願意的時候.
我很希望.
這個故事也幫助大家.
去思想究竟神在你生命裡面.
你是神的傑作.
兜兜轉轉可能你沒有機會發揮.
或者甚至有人告訴你你是不行的.
預了失敗的.
但是神已經預備了一些道路給你.
你堅持跟隨神.
也相信神的時候.
神在那個過程當中要改造你.
建造你.
因為如果我們將.
正話體的希臘原文.
這個字的深意.
你就會見到新普及.
譯本聖經將原文.
將希臘原文.
譯得比較精準一些.
容易明白一些.
就去講到直接說.
這個不是我譯的而是新普及.
譯本,日本聖經是新普及譯本.
譯的就是我們是神的傑作.
直接講.
他已經.
在基督耶穌裡.
將我們創造成為新人.
所以我們信了耶穌.
我們又成為新人.
被神創造.
目的是什麼呢?.
這個傑作的目的就是去.
行就是成為一種生活的方式.
無論在我們的.
學業,我們的事業.

$^{481}$在我們的事奉,好像我剛才所分享.
我兒子那個音樂的歷程一樣.
他當時.
被人訪問的時候,他就說.
I am at the beginning of my discovery of music.
所以我也不知道他那時17歲.
說這些話.
他剛剛開始去發掘.
原來發掘完之後.
神會怎樣用他.
我希望發掘也是你的起點.
因為神早已為我們計劃.
要做的善事.
這些善事就是.
神心意當中的事.
所以神塑造了你.
我們是神的傑作.
已經在基督耶穌裡面造成.
包括了什麼呢?.
這個傑作有時候我們只是很窄的看.
在教會服侍,不是的.
我們的恩賜,我們整個人.
我的負擔.
每個人有不同的負擔,有些人對兒童.
有些人對接待.
有些人對教導,有些人去帶領小組.
有些人對代展.
每個人有不同的負擔.
在社會上你們也有不同的負擔.
有環保.
有社會的議題,經濟.
各種,你會有你自己的負擔.
財幹.
你的學歷,你的過去的經歷.
性格.
有時候你說,我的性格好像不適合做這件事.
但是神創造了你.
某一種的性格.
在某一些的崗位裡面.
正正是這個性格所需要的.

$^{521}$一些價值觀.
都是因為你成長而帶來的一些價值觀.
價值觀是我們的成長很多時候帶來給我們的.
這些價值觀.
如果需要更新的,就被神更新.
成為基督的價值觀.
如果神就已經創造你這些價值觀.
而沒有和聖經原則違反.
這些價值觀.
就是你獨特的.
你的背景也沒有人能和你一樣.
你在哪裡家庭長大.
你在哪裡出生.
你做過什麼.
你的經歷,你的失敗.
你的挫敗,你的成就.
所有的經歷.
所以這個就是獨特的你.
神是透過我們整個生命.
去做了一個傑作出來.
就是你.
所以我們在教會.
有一個叫做生命之約的課程.
叫做事奉德力的傳言.
去年我們有一個課程上了.
如果大家上去那個APP.
我也鼓勵大家.
去年9月至12月.
上過課的學生.
你可以鼓勵其他的朋友.
你就得益了.
你可以去鼓勵他們.
就是我們要去發掘自己.
神這個傑作.
究竟我應該是透過我.
各方面生命的不同的方面.
去開了我生命的藥匙.
去更加去.
找到神要我做的工作.
也都在當中.

$^{561}$這個課程裡面.
用SHAPE我們的特色.
包括我們恩賜,熱忱,才幹.
性格,經歷.
怎樣去發掘自己.
以致能夠有焦點的.
去走你的人生.
所以我鼓勵你.
我們可能在.
明年我們會再開這個課程.
你現在可以報名.
因為我們知道哪些弟兄姊妹會參加.
你就可以.
到時我們開課程的時候.
就會通知你.
神獨特的創造.
我很想你記住這個圖.
你不是.
大量生產的.
你不是神流水作業.
升的生產.
你是神個別.
獨特巧妙.
奇妙的創造.
這個認知.
和把握是重要的.
因為你會知道.
為什麼保羅後來.
會繼續去用.
恩賜和基督的身體.
去鼓勵我們.
就是當我們知道.
我們是神獨特的創造.
我們已經擁有了某些東西.
我們就需要繼續給神去.
裝備我們.
給神去塑造我們.
給神去建立我們,改變我們.
在.
《羅馬書》第十二章四至八節.

$^{601}$我們一起讀這段經文.
這段經文是很熟悉的,但我想大家一起讀這段經文.
正如我們.
一個身子上.
有好些肢體.
肢體也不都有一樣的用處.
這樣.
我們許多人.
在基督裡是一個身體.
互相聯絡作肢體.
按著所得的.
恩典.
或有不同的恩賜.
或說預言.
要按著信心的程度說預言.
或服事的.
要專一服事.
或教導的要專一教導.
或勸勉的要專一勸勉.
施捨的要誠實.
賜禮的要殷勤.
憐憫的人的要樂意.
當你讀這段經文的時候.
我很想你看到幾樣.
很重要的東西.
就是為何保羅在.
二弗所書二章十節裡.
先說了你是神的傑作.
因為你知道你是神的傑作.
而你這個傑作.
就擺在基督的身體裡.
那你的.
獨特性在這個身體裡.
究竟應該怎樣才能夠.
發揮出最大最佳的功能呢.
肢體也不都有一樣的用處.
所以需要彼此的.
配搭.
而最關鍵我們需要明白的.
不是我們要做什麼.

$^{641}$而是神要你做什麼.
所以你的起點不是.
我有什麼厲害.
我有什麼要做.
而是你問神或問弟兄姊妹.
你有什麼認同我.
你看到我在神面前.
你怎樣在整個教會團隊當中.
有一個球隊的配搭.
還有教練.
去幫助球隊.
能夠有好的配搭.
就按著所得的恩典.
我們有不同的恩賜.
恩典和恩賜其實是.
同一個希臘原文的字.
Charisma.
是恩典的字.
Charisma的.
希臘文是Charis.
所以.
這兩個字是overlap的.
換句話說你的恩賜是神給你的恩典.
是神賜給你的.
所以神做你 神裝備你.
所有都是由神做起點.
而既然是恩典.
我們怎樣去回應.
我們有沒有浪費神的恩典.
我們有沒有浪費了.
祂已經給你的東西.
有時候我們.
看過自己太高 又不行.
看過自己太低 又會浪費了.
神給你的恩典.
所以我們求主憐憫.
我們在團隊當中.
我們能夠彼此去印證 彼此去服侍.
知道原來.
神已經給了這些恩典給我們.

$^{681}$而我們在這一刻.
我們怎樣去回應祂.
我也想大家留意.
就是.
他在下面說.
或說預言.
就按信心的程度說預言.
或服侍 或教導.
或勸勉.
要專一.
詩寫的要怎樣 致禮的要怎樣 憐憫的要怎樣.
你有沒有留意 很特別的說.
這裡其實是說給我們聽.
什麼呢.
我們要以為已經是完全人.
或者已經是專一.
或者已經是誠實.
已經是恩勤 已經是樂意的人.
才去服侍.
而是神已經給了你這個恩賜.
你已經可以去參與.
你就開始.
神其實要在服侍的過程當中.
你運用這些恩賜的過程當中.
是在造就你.
去建立你.
去改變你.
去讓你熟齡成熟和成長.
讓你熟齡成熟和成長.
大家不要忽略這段經文.
很多時候我們用這段經文.
就是我們要有什麼恩賜.
就放在哪裡去事奉.
那個當然是對的.
但不單止是這樣.
如果我們以為我們要很成熟.
或者我們要怎樣才能夠做.
其實.
只要我們在團隊當中.
弟兄姊妹的印證是重要的.

$^{721}$我記得.
很多年前.
1987年的時候.
我跟牧師說.
我一直在做生意.
忽然間我跟他說.
我想全時間讀神學.
他第一件事問我.
你有沒有參與教會的.
外展服侍.
或者對於某些.
因為當時做生意很忙的.
所以經常到處去旅行.
但我只知道神一直呼召.
這是另外一個故事.
但是弟兄姊妹.
不是挑戰我而是提醒.
你說你從做生意之前.
你已經被神呼召.
現在你做生意這麼忙.
你在教會沒什麼服侍.
忽然間你說要做到那裡.
傳道人已經是全時間.
好像有一個很大的跨越.
雖然我當時的主任牧師.
是我以前青少年的導師.
但他都會覺得.
好像有很大的跨越.
別人的印證是重要的.
所以我當時.
先去外展部服侍.
在當中去學習.
其實我不熟悉那些服侍.
但就在當中被神造就.
所以很多時候.
神是透過我們的第一步.
入門之後去塑造我們.
去造就我們 去建立我們.
在這裡我想特別提及.
因此是重要的.

$^{761}$我們要發掘因此.
但是我們的因此.
未明顯 甚至是需要繼續去發掘.
就是被訓練去鍛鍊.
去培育.
但這裡也有提到專一.
可能字面上你會覺得.
是不是我們專心做一件事.
不要又做A又做B.
當然不是要大家做到耗盡.
只不過.
專一這個字原文.
不是解釋只是做一件事的意思.
專一這個字的意思.
就是說你做的事.
譬如說寓言 服侍.
教導 勸勉.
你要專一.
意思是你要單純.
不計算.
不求回報的去服侍.
和分享.
這就是專一這個字原文的意思.
當你這樣做的時候你會看到.
我剛才所說的.
就是你要去講寓言.
你要有信心 你的信心大小.
按著那裡去做.
服侍有時候我們可能會比較.
可能我們服侍的.
初初我們做服侍的時候.
可能有不同的心態.
是不是.
但是專一的意思就是.
你要在服侍的過程裡面.
要被練正.
你單純為神為人.
不計較 不再看其他東西.
不會少少.
一些的挫敗或者被人得罪.

$^{801}$就已經.
有不同的回應.
而是要專心的.
單純的 不計算的.
不求回報的去服侍和分享.
是不是這樣的人.
才可以去服侍呢 不是.
保羅奇在講.
你見不見到他在這裡所講.
就是他有很多不同的恩賜 你有這些恩賜之後.
你去參與服侍.
你參與服侍的時候要怎樣.
就是你要留心自己.
是不是不夠單純呢.
是不是不夠不計算呢.
是不是可以寬恕別人.
或者能夠不計回報呢.
是不是不會比較呢.
這些都是在詛咒我們.
內在生命的過程.
所以神裝備我們.
第一節我想大家很清楚明白.
神裝備我們.
不單止是要我們.
去做.
去按恩賜.
服侍一個崗位.
更重要的是.
在服侍的過程裡面 我們不斷被詛咒.
面對挑戰的時候.
不斷被神去.
無助我們克服這些挑戰.
改變我們的內心.
這個裝備.
也是在事奉當中的過程.
以致我們.
能夠成為更合神心意.
被他使用的人.
在經文裡面也講到.
就是.

$^{841}$這讓我們許多人.
在基督裡就是一個身體.
互相聯絡.
作肢體.
這裡很重要的一個觀念.
讓我們能夠清楚明白.
愚民的意思是.
我們不是獨立的個體.
而是在基督裡面.
彼此連結的.
彼此需要的.
彼此建立的.
所以獨特的你.
如果不在服侍當中.
彼此去琢磨的時候.
其實沒辦法去成為.
一個基督身體.
所以教會.
是基督身體的活現.
而不是一個聚會.
所以大家千萬不要想.
我們來到教會聚會.
其實來到教會就是基督的身體.
如何被活現出來.
發揮得最好.
所以每個人要發揮自己的恩賜.
教會就不單止是聚會的地方.
而是活著的基督的身體.
基督徒不是單獨的存在.
而是生命彼此的結連.
你可以想像.
為何保羅要用身體.
做比喻.
因為身體.
缺一不可.
你也不想.
你的身體沒有了一部分.
哪裡都不要.
只要身體裡面任何一部分.
停止運作.

$^{881}$整個身體都受傷害.
這是保羅在這裡想我們明白的.
既然你是獨特的你.
你不會想傷害這個身體.
也想活出身體的意義.
所以我再一次問你.
當你看到這幅圖畫的時候.
在這幅圖畫裡.
你也看到自己.
究竟你的獨特性.
其實在整個團隊.
恩傳的大球隊.
大家庭裡.
究竟你是中衛.
前鋒.
後裔 龍門.
還是教練.
還是一個戰略部.
其實.
這些都是基督身體裡要發揮的.
所以你要問自己.
你的恩賜.
不是為了自己.
而是為了基督的身體.
所以需要別人的印證.
需要和別人一起合作.
你也在那方面.
你也問自己.
你在哪方面最自然的.
去祝福別人.
和和別人配搭.
哪些是你比較容易.
作為一個起點.
願你在那裡開始.
因為神要透過這個起點.
去裝備你.
最後就是神裝備你.
我們在一個基督身體.
的時間.
神要使用你.

$^{921}$就叫我們能夠去行.
祂早已為我們.
計劃好的善事.
這裡的善事.
是包括.
活出神的旨意 信服神.
活出基督的樣式 聖潔的生活.
這是善事.
這個善事不是說好人.
很多時候我們想.
善功就是好事.
其實是說活出神的旨意.
運用恩賜去服侍 去建立教會.
建立別人 建立基督的身體.
成全弟兄姊妹.
裝備弟兄姊妹.
第四就是世界之中去做見證.
這些都是善功.
不是靠行為去換取我們的救恩.
而是活出一種神的旨意.
使得恩賜.
能夠祝福別人的生命.
成為了.
我們生活的方式.
我們真的.
做一次認定在神的裡面.
其實神在你未知之前.
祂已經在建立著你.
究竟有多少是你未知.
而是你留意到神在你生命當中的工作.
所以最後的提醒是.
神已經預備道路.
我們要走在.
那個道路當中.
你有沒有這個尋求的心.
剛才Abbie的分享.
和我們最近在執事會.
我們看了一本書叫領袖的養成.
在結束的時候.
用一個一分鐘的.

$^{961}$總結這本書的學習.
和剛才所說的很類似.
無獨有偶 我們沒有合作過.
其實首先有入門.
其實每一個人的服事.
有一個入門.
不過你的入門是入對了還是入不對.
其實神總之你願意.
神是會使用你.
入門之後.
神會訓練你.
在技巧上.
在你工作上的需要.
也要在人際關係上.
去學習.
現在有很多不同的人.
各方不同意見.
究竟我們怎樣去.
聯繫神.
而神的心意是怎樣呢.
其實.
我們作為一個跟隨神的人.
去事奉祂.
我們首先的起點是要參與.
參與之後神會裝備我們.
去訓練我們.
透過實際的技巧.
也透過生活的一些智慧.
聖經的原則.
透過人際關係當中的複雜性.
去塑造我們 建立我們.
以致有一天我們成熟.
成長 我們能夠分辨到.
神的心意.
神的帶領在每一刻.
所以很鼓勵大家.
其實如果你沒有.
投入參與在事奉當中.
你沒有貢獻.
獨特的你在這個.

$^{1001}$大群體裡面的時候.
其實你是錯過了.
一個屬靈成長的機會.
錯過了經歷.
神更深的機會.
經歷一些你過去.
未曾經歷過的事.
所以在服事當中不單止是要完成.
那個工作 那個事工.
而最重要是被神去塑造.
成為他可以使用的.
領袖.
有時你說 領袖這個名字.
很令人害怕.
但這本書告訴我們.
其實領袖就是一個.
影響別人的人.
其實每一個基督徒.
大使命裡面說的我們聽甚麼.
所以你們要去 使萬民作我的門徒.
就是一班影響別人的人.
其實你和我.
不論影響一個人.
在小組裡面一個關心 一個代討.
那個是一個影響.
那就是一個角色.
你也可以帶一群人.
在整個教會裡面去帶領.
所以神是要使用你.
造就你 然後放你在合適的崗位.
假如在教會裡面.
最需要的是領袖.
而是你最不願意.
做的就是領袖.
那你可能會想 我能不能做到這件事呢?.
其實神關心的.
正如Abby所說.
是你是否願意.
神關心的是.
當你願意的時候.

$^{1041}$祂可以建立你 可以造就你.
不要浪費你的生命.
我記得我很多年前.
在差不多20年前開始.
去參加一些教會研討會講領袖.
到現在都有的.
叫做Global Leadership Summit.
很多時候就說.
如果你是領袖.
或者你覺得你自己可以為神去做一些影響力.
請你踏前一步.
不是馬上去到那個位置.
但是你真的踏前一步.
為什麼呢?.
因為那是一個承擔.
那是一個經歷.
你不要的不是權力.
而是承擔了神託付你的責任.
分擔了你的責任.
然後神就在那個責任當中.
去建立你 造就你.
所以教會其實未來的發展.
最需要的是一群.
同心合意的領袖.
也是一群願意被神塑造.
成為領袖的人.
每一個都是.
其實每一個都可以.
不要浪費你的生命.
要好發揮神給你的潛質.
最後就是.
神塑造你 神裝備你.
神使用你 你留意就是.
一切從神而起的起點.
只要你願意.
神就可以塑造你 裝備你 使用你.
所以你怎樣去回應呢?.
而我們整個的.
健康教會.
建立健康的知名教會的系列.

$^{1081}$其實頭那三樣就是.
敬拜是要以神為中心.
我們的人生是信服祂.
就是我們的小組.
是需要一個.
互相守望 彼此扶持.
建立屬靈生命成長的小組.
我們的屬靈生命.
靈命的塑造 行成長的靈命.
我們的內在生命.
怎樣去被神塑造 建立成長.
這三方面是健康教會最基本的基礎.
就好像今天.
二弗所書第二章十字一樣.
就是神塑造你 神裝備你.
透過這三方面.
但你不要停留在那裡.
你要去到神使用你.
就是神給了你恩賜.
並且祂要有值得你去傳揚祂的恩典.
我們在下個星期開始.
會講到關於你怎樣去.
一起去建立一間.
健康的教會去成就神的使命.
我們都有一個回應.
就是主啊.
怎樣?.
但我們願意求你使用我們.
因為我們不願意埋沒.
我們的恩賜.
因為你拯救我們.
不是因為只是要脫離罪惡.
神呼召.
你和我.
活出祂已經為你預備的人生.
我願意.
成為你未來幾個星期的禱告.
就是究竟.
你是怎樣去回應神.
\newpage



\section{使徒行傳 11:1-30}
\label{sec:032193}
\textbf{同感一靈,突破規限的合一}~~~~余大峰傳道
\newline
\newline
連結: \href{http://www.church.com.hk/acms/content.asp?site=cdc\&op=show\&type=product\&code=032193\&layout=sermon}{\texttt{032193}} ~~~~ 語音日期: 2026-05-24
\newline
\newline
\hyperref[sec:032192]{\small{< < < PREV SERMON < < <}}
~
\hyperref[sec:index_chronic]{\small{[返順時目]}}
~
\hyperref[sec:032194]{\small{> > > NEXT SERMON > > >}}
\newline
\newline
使徒行傳 11:1-30
\newline
\begin{longtable}{cl}
\hline
\hline
章節 & 經文 (和合本修訂版)\\
\hline
11:1 & \begin{tabularx}{0.7\textwidth}{X} 使徒和在猶太的眾弟兄聽到外邦人也領受了神的道。 \end{tabularx} \\ \\ \relax
11:2 & \begin{tabularx}{0.7\textwidth}{X} 等到彼得上了耶路撒冷，那些奉割禮的信徒和他爭辯， \end{tabularx} \\ \\ \relax
11:3 & \begin{tabularx}{0.7\textwidth}{X} 說：「你竟進入未受割禮之人當中，和他們一同吃飯！」 \end{tabularx} \\ \\ \relax
11:4 & \begin{tabularx}{0.7\textwidth}{X} 彼得就開始把這事逐一向他們解釋，說： \end{tabularx} \\ \\ \relax
11:5 & \begin{tabularx}{0.7\textwidth}{X} 「我在約帕城裡禱告的時候，魂遊象外，看見異象，有一塊好像大布的東西降下，四角吊著從天縋下，直來到我跟前。 \end{tabularx} \\ \\ \relax
11:6 & \begin{tabularx}{0.7\textwidth}{X} 我定睛觀看，見內中有地上四腳的牲畜、野獸、爬蟲和天上的飛鳥。 \end{tabularx} \\ \\ \relax
11:7 & \begin{tabularx}{0.7\textwidth}{X} 我還聽見有聲音對我說：『彼得，起來！宰了吃。』 \end{tabularx} \\ \\ \relax
11:8 & \begin{tabularx}{0.7\textwidth}{X} 我說：『主啊，絕對不可！凡污俗或不潔淨的東西從來沒有進過我的口。』 \end{tabularx} \\ \\ \relax
11:9 & \begin{tabularx}{0.7\textwidth}{X} 第二次，有聲音從天上回答：『神所潔淨的，你不可當作污俗的。』 \end{tabularx} \\ \\ \relax
11:10 & \begin{tabularx}{0.7\textwidth}{X} 這樣一連三次，然後一切就都收回天上去了。 \end{tabularx} \\ \\ \relax
11:11 & \begin{tabularx}{0.7\textwidth}{X} 正當那時，有三個從凱撒利亞差來見我的人，站在我們所住的屋子門前。 \end{tabularx} \\ \\ \relax
11:12 & \begin{tabularx}{0.7\textwidth}{X} 聖靈吩咐我和他們同去，不要疑惑，還有這六位弟兄也跟我一起去，我們進了那人的家。 \end{tabularx} \\ \\ \relax
11:13 & \begin{tabularx}{0.7\textwidth}{X} 那人就告訴我們，他如何看見一位天使站在他家裡，說：『你派人往約帕去，請那稱為彼得的西門來， \end{tabularx} \\ \\ \relax
11:14 & \begin{tabularx}{0.7\textwidth}{X} 他有話要告訴你，因這些話你和你的全家都可以得救。』 \end{tabularx} \\ \\ \relax
11:15 & \begin{tabularx}{0.7\textwidth}{X} 我一開始講話，聖靈就降在他們身上，正像當初降在我們身上一樣。 \end{tabularx} \\ \\ \relax
11:16 & \begin{tabularx}{0.7\textwidth}{X} 我就想起主的話如何說：『約翰用水施洗，但你們要在聖靈裡受洗。』 \end{tabularx} \\ \\ \relax
11:17 & \begin{tabularx}{0.7\textwidth}{X} 既然神給他們恩賜，像在我們信主耶穌基督的時候給了我們一樣，我是誰，能攔阻神嗎？」 \end{tabularx} \\ \\ \relax
11:18 & \begin{tabularx}{0.7\textwidth}{X} 眾人聽見這些話，就不說話了，只歸榮耀給神，說：「這樣看來，神也賜恩給外邦人，使他們悔改得生命了。」 \end{tabularx} \\ \\ \relax
11:19 & \begin{tabularx}{0.7\textwidth}{X} 那些因司提反的事遭患難而四處分散的門徒，直走到腓尼基、塞浦路斯和安提阿。他們不向別人講道，只向猶太人講。 \end{tabularx} \\ \\ \relax
11:20 & \begin{tabularx}{0.7\textwidth}{X} 但內中有塞浦路斯和古利奈人，他們到了安提阿也向希臘人傳講主耶穌的福音。 \end{tabularx} \\ \\ \relax
11:21 & \begin{tabularx}{0.7\textwidth}{X} 主的手與他們同在，信而歸主的人數很多。 \end{tabularx} \\ \\ \relax
11:22 & \begin{tabularx}{0.7\textwidth}{X} 這風聲傳到耶路撒冷教會的人耳中，他們就打發巴拿巴到安提阿去。 \end{tabularx} \\ \\ \relax
11:23 & \begin{tabularx}{0.7\textwidth}{X} 他到了那裡，看見神所賜的恩就歡喜，勸勉眾人要立定心志，恆久靠主。 \end{tabularx} \\ \\ \relax
11:24 & \begin{tabularx}{0.7\textwidth}{X} 這巴拿巴原是個好人，滿有聖靈和信心，於是有許多人歸服了主。 \end{tabularx} \\ \\ \relax
11:25 & \begin{tabularx}{0.7\textwidth}{X} 他又往大數去找掃羅， \end{tabularx} \\ \\ \relax
11:26 & \begin{tabularx}{0.7\textwidth}{X} 找著了，就帶他到安提阿去。他們足有一年和教會一同聚集，教導了許多人。門徒稱為「基督徒」是從安提阿開始的。 \end{tabularx} \\ \\ \relax
11:27 & \begin{tabularx}{0.7\textwidth}{X} 當那些日子，有幾位先知從耶路撒冷下到安提阿。 \end{tabularx} \\ \\ \relax
11:28 & \begin{tabularx}{0.7\textwidth}{X} 內中有一位，名叫亞迦布，站起來，藉著聖靈指示普天下將有大饑荒；這事在克勞第年間果然實現了。 \end{tabularx} \\ \\ \relax
11:29 & \begin{tabularx}{0.7\textwidth}{X} 於是門徒決定，照各人的力量捐錢，送去供給住在猶太的弟兄。 \end{tabularx} \\ \\ \relax
11:30 & \begin{tabularx}{0.7\textwidth}{X} 他們就這樣做了，託巴拿巴和掃羅的手送到眾長老那裡。 \end{tabularx} \\ \\
[1ex]
\hline
\hline
\end{longtable}
$^{1}$弟兄姊妹主女平安.
很高興可以和大家一起分享神的話語.
開始之前有個祈禱.
真的,八天父親是主耶穌.
我們一起去你面前.
主耶穌求你幫助我們.
讓我們在你的話語當中.
在聖靈的帶領下.
讓我們去學習.
特別是在今天這個世代.
弟兄姊妹之間.
往往是口裡說合一.
心裡卻沒有利我.
弟兄之情.
幫助我們能夠在不同的事情.
不同的看法當中.
能夠同感同心.
一起在基督裡面前行.
將二十四頁傳言擺上.
求主你託大.
祈禱加特奉主名而求.
阿們.
今年的主題.
教會的主題是什麼?.
厲害了.
沒有看.
在基督裡同心前行.
我不知道你們大哥哥看這個主題.
這個主題不用說了.
我們一向都是在基督裡同心前行.
是不是呢?.
一說.
是呀.
今年是再被提醒.
再被提醒.
真的要在基督裡同心前行.
我們不要說得遠了.
兩夫妻都是基督徒.
是不是兩夫妻一定是同心呢?.
是不是一定前行呢?.

$^{41}$會不會是前行但不是同心.
但是同心但是分開行.
或者同心是有前有後呢?.
兩夫妻都是.
子女都會不會呢?.
大家可能都是回教會的.
是不是?.
但是可能不同教會.
但是不同教會不代表不同心.
究竟同心前行.
是不是必然.
但是若然必然.
就不需要整個主題了.
所以究竟我們怎樣去理解同心前行呢?.
今天我想特別是透過.
大家也知道有一個講道系列.
這次是我們講到《四學行傳》.
我們被吩咐講第十一章.
實際上大家去看的時候.
第十一章和第十一章是相連的.
因為第十一章是講第十章所發生的事.
但是是跟不同的人講.
是非常重要的.
但是十一章是全章一至三十節.
你可以說.
我讀了十八節.
其實之後十九節,三十節.
是同樣精彩的.
但是因為時間關係.
我就調一調.
有一點點講道的模式不同.
不是順序的.
先講了十九節,三十節.
講了十九,三十節的時候.
就要提出一個很重要的問題.
究竟在基督裡同心前行是和誰人講.
可能是和我們是基督徒.
但是有沒有問過.
其實基督徒這個名是從何而來.
什麼叫基督徒呢?.

$^{81}$大家都可能會問.
基督徒就是跟從基督的人.
就叫做基督徒.
是不是?.
我就想告訴你.
其實在《四學行傳》一至到第十章.
是未出現這個字眼的.
即是說在絕星那天之後.
起碼十多年.
是沒有出現這個字眼完全.
第一個問題.
第二個問題.
如果你的時光隧道去到那個時間的時候.
你說我是基督徒.
基督徒?什麼基督徒?.
他就完全沒有這個字眼.
第二個問題.
哪些人開始給這個名字給基督徒呢?.
兩個選擇.
是基督徒自己給這個名字給自己.
還是不是基督徒的人去給這個名字給基督徒呢?.
前者舉手.
後者舉手.
不知者舉手.
基督徒這個原文.
即是希臘文.
希臘文.
它是叫做Christianos.
不是Christian.
是Christianos.
當時當時是希臘世界.
羅馬地方.
希臘世界.
Christianos其實有兩個字組成的.
你要明白兩個字組成.
你就知道為什麼來了.
第一個字是希臘文的Christos.
Christos大家知道了.
Christos就是受高者.
當然我們現在就說基督.

$^{121}$受高者.
特別是基督教就指向基督.
但是他接著的尾.
是其中一個拉丁文的尾.
因為當時羅馬地方是說拉丁文的.
普遍說拉丁文.
所以拉丁文的尾.
就是Ianus.
I-A-N-U-S.
但是在希臘文的發音.
U是沒有這個發音的.
所以它就變成Ianos.
所以你發覺一個頭一個尾.
用希臘文的尾.
去黏在一起.
就變成Christos加上Ianos.
就變成什麼呢.
就叫做Christianos.
最初最初叫Christianos.
這個字呢.
是由外邦人去說給他們聽.
是由安提阿的人說出來的.
這個問題來了.
為什麼要突然間說一個這樣的字眼.
去形容這一班的人.
這一班人有什麼特別.
他不是說他是猶太人.
或者說他信耶穌的人.
為什麼叫Christianos.
為什麼特別這個字眼.
因為在安提阿.
發生了一件很歷史性的事.
我現在就要說一說.
這個歷史性的事.
從這個歷史性的事開始說.
經文呢.
我一開始就由十九節開始說.
在十九節他這樣說.
那些因史提反事件.
遭受苦難而四散的門徒.

$^{161}$四提反事件大家記得了.
四提反是誰.
四提反就是第六章.
發現基督徒在猶太人當中.
是說希伯來文或者阿南文的猶太人.
和說希臘文的猶太人.
在他們大家吃飯的時候.
那些菜是分配得不是那麼平均.
特別是孤兒寡婦那些.
所以有些說希臘.
對於哪些人不平均呢.
是對說希臘話的那些.
所以說希臘話的那些基督徒.
就開始有些微言.
和那些說猶太阿南文.
或者說希伯來文的那些猶太人.
有些微言.
所以結果就選了七個執事.
去管理那些飯食.
其中一個執事叫做四提反.
後來四提反也是傳福音.
所以執事也傳福音.
大家記住執事也傳福音.
不只是牧者才傳福音.
傳福音之後.
大家就是被逼白.
被猶太人逼白.
逼白之後.
所以特別是說這些希臘話的信徒.
就四散.
四散之後.
他們分別去到哪個地方呢.
你看看那個箭咀.
去到這個腓里基.
即是我們叫做腓里市.
或者舊約叫做腓里市.
就是那個沿海的地方叫腓里基.
那個海島就叫做塞浦路斯.
再上面就圈著安提柯.
去了這三個地方.

$^{201}$但是他們去到那裡.
特別是要註明什麼呢.
是他們不對別人傳講.
只對猶太人傳講.
到這一剎那來說.
他們怎樣理解馬太福音的大使命呢.
我曾經也講過.
他們的理解就是.
只不過是去到萬國萬民.
當中的猶太人傳福音.
所以就只是向猶太人去傳講.
但是很奇怪.
其中有一些塞浦路斯人和古里內人.
塞浦路斯人在哪裡呢.
這個小道.
古里內人在哪裡呢.
是現在的利比亞.
利比亞就是北非.
但是他們在安提柯那裡.
這兩班人去到安提柯.
在安提柯那裡.
他竟然開始對希臘人傳講主耶穌.
和希臘人傳講主耶穌.
這個就未是歧視.
為什麼呢.
因為耶穌來之前.
當時的猶太人.
都會有和愛邦人傳揚神.
如果當時的愛邦人信了神.
耶穌還沒有來.
就是耶穌之前.
愛邦的人信了神.
覺得猶太教的神都不錯.
他們相信.
那麼他們叫什麼.
信了神.
他們在聖經叫什麼.
God-fearer.
敬畏神的人.
所以在新約見到一個敬畏神的人.

$^{241}$他是一個敬畏神的人.
高利樓就是一個敬畏神的人.
就是因為那個高利樓.
但是如果敬畏神的人.
為什麼叫敬畏神的人.
因為他信了猶太教的神.
但是如果他再進一步.
他想歸依入猶太教.
他起碼要做兩個動作.
第一個動作.
如果是男士要做國禮.
第二樣東西.
他以後吃東西.
就要很檢點.
所有不潔之物不能吃.
而不潔之物不能吃.
也和不潔之人不能吃.
據說有兩樣東西去到手.
和不潔的人有些距離.
因為不潔的人吃不潔之物.
更加不可以去他家吃東西.
現在又到時光倒流.
想問大家一個問題.
今天你時光機坐到那裡.
十一章之前.
我問你.
你見到耶穌.
直接說耶穌.
不要說彼得保羅.
直接說耶穌.
耶穌說.
我今天在家裡吃火鍋.
過來吃.
耶穌去不去你家吃飯.
十一章之前.
去還是不去.
去不去.
去.
他有沒有犯一個潔淨的律例.
讓你問一下.

$^{281}$耶穌的門徒去不去.
所有這些人去不去.
十一章之前.
他如果進去.
去到你家.
和你吃.
他就犯了所有潔淨的律例.
他不可以再說守律法.
他只是這樣東西.
都已經被守律法的人大肆宣揚.
你連律法你都不守.
所以答案是不會的.
但是怎麼全副音呢.
全副音不同.
我和你聊天.
在街上.
中開地方和你吃.
我不是那麼親近.
還有吃也好.
到達可能.
我自己的袋子.
可能會帶著一些.
那些叫什麼.
潔淨了的食物.
餅也好.
各方面也好.
我自己都帶著.
有一點點事.
現在你去closing woollies.
有個特別尊貴的.
我們北岸更加知道.
北岸有些聖埃特.
特別多猶太人區.
所以特別會去到聖埃特的woollies.
有個櫃.
Cultured food.
Cultured food就是潔淨的食物.
當時就是他們去傳副音給這些.
希臘人.
是他們開始的.

$^{321}$但是竟然沒有提名字.
但是就是一班這樣的人.
去傳副音給他們.
而主一手和他們同在.
信而歸主的人.
就多起來了.
這件事就傳到耶路撒冷教會的耳中.
他們就派巴拿巴和安提奧去.
派巴拿巴去.
巴拿巴是什麼人呢.
下面可以看到.
巴拿巴是一個.
本來叫約瑟.
使徒叫他巴拿巴.
因為巴拿巴的名字是安慰者.
所以你可以說.
巴拿巴是一個橋樑.
兩個不同的群體.
幾個橋樑接駁他.
他本身也是在塞浦路斯出生.
剛剛好.
塞浦路斯出生的.
他就和塞浦路斯人說.
就問過究竟.
為什麼你傳副音給他們.
所以他去到那裡.
看見神所施的恩.
他就歡喜了.
原來他見到神職.
神施的恩.
他的心就歡喜.
全面眾人堅心靠主.
巴拿巴是個好人.
滿有聖靈和信心.
於是很多人就歸了主.
總部的事.
耶路撒冷總部.
聽聞安提奧竟然有些信徒.
猶太信徒.
傳副音給外邦人.

$^{361}$巴拿巴去看看.
巴拿巴去看過究竟.
而巴拿巴是充滿聖靈的.
授權去看.
聽了就覺得.
真的有神的靈同在.
所以他接受了他們.
非常歡喜.
信了他們很多.
後來.
不單止這樣.
巴拿巴還專門去大掃.
去找誰呢.
去找掃羅.
講回第九章.
第九章就講掃羅.
怎樣在達瑪士格的路上.
去遇見主.
所以遇見主之後.
他就被人逼罷.
當時被猶太人逼罷之後.
他就去了大掃.
所以巴拿巴知道.
掃羅應該找他來幫忙.
所以就去找他.
就帶他來到安提奧.
足足有一年的時間.
他們就一同在教會聚會.
教導了很多人.
門徒就開始在這裡被稱為.
應該說被稱為.
基督徒.
就由安提奧開始.
我問大家一個問題.
他們一同在教會聚會.
有沒有吃東西.
你怎樣知道有吃東西.
怎樣知道有吃東西.
我們聚會不吃東西的.
你今天崇拜.

$^{401}$崇拜就這樣吃東西.
你崇拜不一定吃東西的.
但當時的聖餐.
吃不吃東西.
吃著東西是聖餐的.
當時聚會有沒有聖餐.
有多少聚會就有多少聖餐.
因為聚會就是聖餐.
所以是有吃東西的.
問題來了.
竟然這班的外邦人.
竟然.
行了國禮沒有.
經文沒有.
竟然和他們吃東西.
這班人就變了新種類.
他不是以往猶太人.
不是以往的外邦人.
因為以往的猶太人和外邦人.
合不合在一起.
可以合在一起.
不能合在一起聚.
更加不能合在一起吃.
全都破了.
因為這樣所以外邦人看到.
不知如何歸類.
之前怎樣稱基督教的信徒.
稱他為Nazarene.
即是說你是跟隨.
或者你是Nazarene.
那撒勒人.
或者你是加利利人.
來稱呼那些信徒.
因為你所信的耶穌就是從那裡來的.
但還沒有出基督徒這個名字.
在那裡開始.
因為他是第三類的人.
因為之前沒有人這樣混在一起.
他們兩班人是不可能一起聚會的.
人們看到是完全不明白的.

$^{441}$素羅的背景我不說了.
素羅本身是非常爭議性的.
在經文裡面說到.
他上了耶穌之後.
是被猶太人攻擊.
不單止猶太人攻擊.
甚至乎當素羅上了耶穌之後.
回到耶路撒冷.
去接近門徒.
門徒是很害怕他的.
也不相信他是門徒.
而是靠誰去接待他.
都是巴拿巴.
所以巴拿巴一個人.
他同時間是將這兩班人.
和又是將那兩班人.
都是透過巴拿巴這個橋樑.
去給他們大家.
幫他們做一個同心合意的一個同行者.
今日教會要同心同行.
很需要在我們當中.
找多些巴拿巴.
你覺得你是不是巴拿巴?.
你做不做這個角色?.
在小組的衝突.
在教會的事奉的衝突.
在大家不同心大家爭拗的時候.
你覺得你做不做巴拿巴的角色?.
如果教會一個巴拿巴的角色都沒有.
一個巴拿巴的牧者都沒有.
一個巴拿巴的執事都沒有.
一個教會怎樣同心怎樣同行.
但若然教會有很多巴拿巴.
就很不同.
但巴拿巴是一個什麼人?.
充滿聖靈和信心的人.
所以今年的主題.
我們說要在基督律同心同行.
讓我們每人都是學習做巴拿巴.
所以在這個地方.

$^{481}$首次基督教開始真真正正跨民族.
跨語言,跨地域.
一個很具體的表現.
不單止這樣.
講完前面那一段之後.
他接著那一段.
他說那時有幾位先知從耶路撒冷.
來到安提河.
其中有一個名叫阿加布站起來.
藉著聖靈指出天下將有大饑荒.
這件事在格魯丁時期發生.
於是門徒決定按照各人的力量去捐款.
好送給住在猶太的弟兄.
他們就這樣行了.
由巴拿巴到索羅.
經手送到長老門那裡.
現在再講一講.
是哪些人去幫哪些人呢?.
是安提河的信徒.
安提河是什麼信徒?.
主要是愛邦信徒.
而愛邦信徒其實之前是被猶太人分類的.
不和他在一起的那些信徒.
但現在反而援手去幫助.
同心同行.
不單止你和他同心.
也要他和你同心.
哪怕曾幾何時你和他表示敵意.
但當大家同心同行的時候.
他都會反過來幫助他.
換句話說.
在安提河的愛邦弟兄.
他當不當猶太人是自己人呢?.
當不當在耶路撒冷的肢體是自己人呢?.
他當是自己人來幫助.
出錢又出力.
換句話說.
是真真正正這兩個地方的教會.
不同種族的信徒.
能夠透過這件事件顯示出.

$^{521}$雖然他們兩個地方不能吃飯.
但有事相救.
一定會捐款去相救.
這是一個非常好的圖畫.
也透過巴拿巴和索羅一起拿去.
好了.
這個可以看到.
以弗所書所講的兩下合而為一.
當然在以弗所書來講.
是講愛邦人和猶太人.
但其實再深入看.
其實猶太人之間都是合一的.
為什麼這樣說呢?.
在耶穌之前.
在耶穌來到的時候.
當時的猶太人已經有不同的幫派.
他們合不合一呢?.
是發覺不合一的.
有很多爭拗的.
有犯下錯誤的人.
有殺刀該人.
有些是藏在山洞裡的.
我們現在叫做.
山洞的那些人.
所以當時是很不合一的.
猶太人之間有很多爭拗.
有些是較近希臘人.
有些是沒那麼近希臘人.
但其實這是兩下合而為一.
在聖靈底下.
他們都會合在一起.
好了.
說到這裡.
就要問一個問題了.
為什麼之前那些塞浦路斯人.
他們為什麼會轉變成希臘人呢?.
因為之前在十一章發生了一些事情.
事出必有因.
就要去回帶.
看看之前發生了什麼事情.

$^{561}$之前發生了什麼事情呢?.
在十一章頭說.
就是今天的經文.
就是使徒和猶太的弟兄們.
聽說外族人受了神的毒.
彼得去到耶路撒冷的時候.
那些守國禮的人.
就跟他們爭論.
你竟然去跟未受國禮的人.
跟他們一起去吃飯.
我想問大家一個問題.
這班守國禮的人是什麼人呢?.
彼得有沒有守國禮呢?.
彼得是不是屬於守國禮的人呢?.
這裡所指的守國禮.
不單止是守國禮.
他好像快財一樣.
他自己守國禮.
他都堅持所有信神的人.
即是那些敬畏者.
他都要守國禮.
他們都要皈依猶太教.
在這班基督徒來說.
他們的救我.
也是帶到他們身上.
他們都要皈依基督教.
而其中一個皈依.
就是他們要做守國禮.
但是要記住.
這班人已經接受了神的道.
接受了神的道.
他唯一未有的.
是未做國禮.
所以他就走過去跟彼得說.
彼得去到耶路撒冷的時候.
彼得彼得.
他跟他們爭論.
當他跟彼得爭論的時候.
他跟彼得都是守國禮.
兩者現在有沒有合一?.

$^{601}$大家開始爭論.
有沒有合一他們兩個?.
彼得跟他們有沒有合一?.
他們兩個現在沒有什麼合一.
因為已經爭拗了.
而爭拗什麼呢?.
就是接受了神的道的人.
他們竟然未受國禮.
而你跟他吃飯.
用今天的表達是什麼意思呢?.
就是說.
他們信耶穌了.
但是信了耶穌.
不代表你可以跟他吃飯.
今天跟你說這句話難聽嗎?.
就是說信了耶穌.
我跟他也有距離.
我久不久也會聽到.
不要把我放在他那一組.
他信了耶穌.
信了耶穌又怎樣?.
不要把我放在他那一組.
有他沒有我.
他也是信主.
他也許是初信.
或者他也有些掙扎.
那又怎樣?.
有他沒有我.
大家合不合一?.
同不同心?.
前不前行?.
是做不到的.
但這群人正正就是在做這件事.
你竟然跟未受國禮的人去吃飯.
眾人聽見這些話.
為什麼會跳到十一章十八節呢?.
中間插著彼得回述他發生什麼事.
說完這番說話之後.
突然間這群人聽完之後.
眾人聽見這番話就默默無聲.

$^{641}$飽榮耀歸給神.
說這樣看你了.
神也是將悔改的心賜給他們外族人.
使他們得生命.
是不是馬上緩和了氣氛?.
馬上不同了.
問題是為什麼態度突然間會一百八十度轉變呢?.
就是中間夾著餡了.
而夾著餡.
就是彼得去講述發生了什麼事.
之前發生了什麼事.
而當他聽到發生了什麼事之後.
連這群幫他爭拗的人.
都默默無聲.
歸榮耀給神.
現在這十一.
大家合不合一?.
合一了.
所以我問大家一個問題.
剛才說的問題是.
有他沒有我.
他如果在這組我就不在這組.
在這樣的情況之後.
突然間過幾天之後.
他就說好吧好吧大家一起吧.
那你問他發生了什麼事呢?.
就是中間的餡了.
大家有沒有經歷過餡呢?.
有沒有經歷過有他沒有我呢?.
或者被人說有你就沒有他.
有沒有經歷過呢?.
不要舉手了.
太尷尬了.
兩夫妻都可能會的.
今天你出去外面睡.
那就是有你沒有我了.
床就是我睡而已.
你出去外面睡.
經歷了什麼呢?.
就是經歷我們所講的.

$^{681}$愛邦人的五重節.
原來五重節我們經常說是猶太人經歷.
在第二章經歷.
但是原來在第十章.
神同樣是讓一群愛邦的信徒.
哥利留.
他們在彼得面前.
是經歷了聖靈降臨.
而就是因為這件事件.
當彼得再複述的時候.
他們就百分之百改變.
今天我的題目是講什麼呢?.
同感.
我的題目是什麼呢?.
同感一靈.
同心前行.
怎麼能夠同心呢?.
原來同心要包括同感.
是感什麼呢?.
就是感.
原來聖靈感動他們.
正如他這樣感動我們.
聖靈降臨在他身上.
正如聖靈降臨在我身上.
既然聖靈都這樣去恩.
神都賜聖靈給他們.
都賜給我們.
我們是誰呢?.
但這就是你和我要問的問題.
當你對著一個.
今天基督徒因為什麼事會吵架.
你想一下.
你有多少時候是因為神學.
和你小組組員吵架.
有多少?.
你講了十次.
和你吵架有多少次是因為神學不同呢?.
吵架的.
很少的.
但會因為什麼而吵架.

$^{721}$如果以神學來說.
就是他只講他的神學.
或者將他的神學蓋過其他所有人的神學.
有他講沒人講.
這樣小組就不耐了.
會不會有這樣的情況?.
但這個我覺得不是多數.
更加多數是什麼呢?.
我懷疑.
是大家性格的差異.
大家的互相不諒解.
大家的以我為標準去看你.
以你為標準去看我.
而當你以你為標準去看我的時候.
我覺得是被審判.
被藐視.
被忽略.
我沒有自我.
我太過遷就你.
但我遷就下遷就下忍不住.
另一個就是太過自我.
亦都覺得他自己的我是很對的.
其實是救我.
新的我他沒有考慮別人.
他仍然那麼大聲.
仍然那麼獨裁.
仍然是一言堂.
仍然是很自誇.
仍然是很自大.
我怎樣和他同行.
會不會有這樣的情況呢?.
可能是.
但是今天的經文告訴他.
我們是不是每個人都要.
同樣好像那個經文那樣說.
要聖靈降臨.
我們大家都講謊言.
就可以了.
如果是這樣的時候.
我就要問這個問題了.

$^{761}$若然聖經我們這樣理解.
那很自然問這個問題.
他們因為這樣都能夠跨語言.
跨種族.
跨飲食的文化.
都能夠合一.
我們不就都要等聖靈降臨.
問題就來了 聖靈降臨了沒有?.
如果用他這個邏輯.
那你懂得講謊言嗎?.
懂得講謊言舉手.
我懷疑在當中有些人懂得講謊言.
你笑的那個人是你.
但是你懂得講謊言都不敢舉手.
因為可能懂得講謊言和不懂得講謊言.
在某些群體來講.
都成為一個不能同心前行的.
一個的阻礙.
會不會啊?.
所以就算是你都不敢舉手.
同不同意啊?.
這個問題你都不會舉手.
因為同意和不同意.
也都成為障礙.
就是原來很多很多障礙.
很多很多不同的規限.
那你為什麼要同心前行呢?.
那問題就反過來了.
若然用這個邏輯的話.
要大家都要懂得講謊言.
啊 原來你也懂得講謊言啊.
哎 我也懂得講謊言啊.
那所以大家就同感一靈了這樣.
但是你不懂得講謊言啊.
我又不懂得講謊言.
那大家就沒辦法了.
又分離我.
是不是這樣應用呢?.
究竟今天我們信徒之間.
怎麼能夠同心同感呢?.

$^{801}$我猜啊.
當我們大家去相交的時候.
而不只是說飲飲食食.
去分享其實.
聖靈怎樣活在你當中.
怎樣跟你說話.
怎樣感動你.
又真又實地說出來.
你的掙扎.
聖靈提醒你.
怎樣帶來你喜樂.
帶來你嘆望.
當你跌倒的時候怎樣扶起你.
怎樣去拆敗使者在你身邊.
怎樣叫你學習天悲.
學習要去聖潔.
學習要去公義.
你將這些聖靈的提醒.
或者掙扎.
你跟別人分享的時候.
而對面那個人就說.
啊 原來你也有這些掙扎啊.
我也有啊.
哎呀 原來你這樣子去跟從神的.
我也是這樣子去跟從神的.
大家就突然間.
不是因為大家的性格而分離我.
而是反而哪怕性格.
仍然大家好像connect了.
會不會呢.
夫婦之間相處之道.
會不會.
不是說你家 我家.
你又怎樣煮飯.
我又怎樣吃飯.
不只是說這些細節的問題.
因為未信的人都會說這些細節的問題.
是不是啊.
都會為這些細節的問題離婚.
但是信的人就是說.

$^{841}$是.
我是這樣子出生的.
我家裡說話很大聲的.
你家裡是很溫文的.
不說話的.
我一吵架就在那裡爆.
你一吵架就死忍的.
一個就爆 一個就忍.
越爆就越忍.
越忍就越爆.
結果忍那個才爆.
爆那個就爆到發瘋了.
但是如果你在主內的時候跟他說.
是啊.
原來我是一個很爆的人.
但是主剛剛提醒我.
叫我為著愛你的緣故去忍.
而忍那個人就說.
是啊 聖靈也提醒我.
為著你的緣故.
我將我的苦.
我的看法.
我夠膽去說出來.
但是我忍了一輩子.
我沒有了自我.
但是我為了主的緣故.
我去說出來.
你聽不聽到.
會不會這樣.
就是夫婦之道.
和感一靈.
用聖靈的感動.
去開始建立那個夫婦關係.
而是這樣的時候.
弟兄姊妹.
可不可以當我們討論.
無論是518.
那些時間怎樣安排.
而不是講細節.
而是講.

$^{881}$我在當中.
聖靈怎樣感動我去安排呢.
聖靈怎樣感動我要去考慮別人.
或者放下自己.
當在小組的時候.
是啊.
他進來這個小組.
其實很難相處.
和他性格很不同.
但是.
我可不可以和他嘗試去.
分享.
原來.
我是很怕一些這麼大聲的人.
因為我爸媽這麼大聲.
所以當你在小組這麼大聲的時候.
我就看著好像你爸爸媽媽一樣.
當我在家裡經常被老婆罵的時候.
而你在小組又好像很權威一樣.
我就會.
覺得你好像.
很有權威的太太一樣.
會不會有這樣的情況.
將那個掙扎.
去講出來.
但是很難做的.
難不難做.
彼得都很難做.
但是神要透過三次.
去和彼得講.
但是三次和彼得講的時候.
原來之前.
在彼得未知之前.
他已經和高利樓講.
這些是用AI GEN出來的.
試試玩一下.
COPYRIGHT GEN出來的.
神他怎樣做呢.
他第一件事他不是找彼得.
他找一個外邦的敬畏神的.

$^{921}$這個白夫長.
高利樓.
先和他講.
在他的角度來講是第一天.
但是當他聽了之後.
天使就說你這樣做.
你吩咐那些人.
那些人要去.
由哥薩尼亞去到約伯.
這個城市.
去到找一個做皮匠的西門.
裡面有一個叫彼得的西門.
逐層逐層去.
去到那裡.
你就要和他聊天.
他會有人和你講.
而他呢.
這個高利樓.
他又聽.
所以又派了.
兩個僕人.
再加上一個丁丁.
三個人.
在這段經文很多三字.
在舊約.
很多三字什麼意思呢.
就是神的手.
神的工作.
三次三次三次.
所以呢.
就是這樣由.
這個哥薩尼亞就走到約伯那裡.
所以這裡就在那個.
他要走大約一天的時間.
二十個小時的時間.
到第二天.
就去到他門口了.
但是問題是這幾個人.
這三個人去到門口.
同時間.

$^{961}$同時間的故事講什麼你知道嗎.
好笑.
同時間就講彼得.
肚餓.
他就在天台那裡.
可能不知道在天台那裡.
很肚餓.
肚餓在想什麼.
肚餓就想起吃東西.
想起吃東西的時候.
突然間天上面.
有一塊大布下來.
四條四邊掛著.
那塊大布下來.
裡面看到什麼呢.
看到有一些四腳的.
四腳不是四腳的食物.
四腳的畜生.
然後有一些鳥.
然後有一些爬蟲類.
四腳的畜生.
包不包括豬.
有沒有.
這裡有圖畫嗎.
有豬.
你們喜不喜歡.
我打算明天去.
一個弟兄姊妹家裡吃飯.
他問我弄什麼送他.
我打算弄紅燒豬手.
但是彼得他不吃豬手的.
果實.
豬紅等一下不用說了.
所以他miss out了.
中國人弄紅燒豬手很棒的.
但是他miss out了.
為什麼呢.
他不吃豬的.
是不潔的.
所以看到這樣的情景.

$^{1001}$我無論是那些族物.
族物或者不潔之物.
我都沒有碰.
我從來都沒有碰.
什麼是族物呢.
不潔之物大家知道了.
豬就是不潔之物.
但是什麼是族物呢.
族物本來是潔的.
但是潔的物.
和不潔的混在一起.
或者碰過.
或者和不潔的人碰過.
或者處理過.
他都成為由不潔就變為族了.
我用這個舉例.
你在家裡.
你切豬肉有那塊板.
切生牛肉有那塊板.
切完之後.
牛肉本來是潔的.
但是你拿去切豬肉.
然後再切牛肉.
那你豬肉牛肉.
豬肉就不潔了.
牛肉本來是潔的就變成了什麼.
變成了族了.
我們現在都有這樣的情況.
你家裡切東西.
有沒有熟的找一塊板切.
生的找一塊板切.
有沒有混在一起.
混在一起有沒有事發生.
什麼事發生.
肚子痛嗎.
肚子痛未必發生的.
被老婆看到.
在我的狀況.
我老婆不在這裡.
我老婆不在這裡可以進言.

$^{1041}$如果我被我老婆看到.
就馬上有事發生了.
潔和不潔不能放在一起的.
生和不生的很多次都放在一起的.
所以不可以懶的.
不可以懶的.
好了,講回這裡.
就是看到這樣的情況之下.
而彼得就說我重來.
我重來.
重來是什麼意思.
就是彼得去救我的規限.
原來今天我們每一個人.
有一個教會.
都很多很多時候會講重來.
我未信主的人的時候.
我從來不會和這些人相交往.
我只喜歡.
我喜歡打球.
我喜歡和打球的人.
我喜歡人士就喜歡和不喜歡人士.
我有品味的就喜歡和有品味.
我喜歡煮菜就喜歡和煮菜.
我從來都是這樣的.
但是信了耶穌之後.
可能要將你的重來.
要因為見到這個異象去放下.
在主裡面.
以前的重來.
就變為今天的突破.
以前的重來.
變為今天的突破.
希望大家去問一問自己.
有沒有什麼人.
你從來都和他隔不離.
見他後面.
你已經爭他前面.
而今天講到合一的課題.
特別講到今天這個故事.
和我們中國人是很接近的.

$^{1081}$就是什麼呢.
如果你和一個弟兄姊妹.
你能夠願意和他同桌吃飯.
而有商量.
你們都很同心.
都很同行.
都很同感.
剛才聽到勁敗小組之前說.
他說什麼.
嘆茶.
我就後來問他.
你待會中堂會不會去嘆一嘆咖啡.
他就說會.
我就這樣聽到這句話.
今天就講到.
我今天就講這個例子.
你覺得他們一起嘆咖啡.
同不同心嗎.
這個動作都不容易的.
但如果突然間有一個.
我希望沒有.
但我都是猜的.
我希望沒有.
我為什麼要講出一個Trump(特朗普).
我希望沒有.
就是.
不要剩下一個.
他如果沒空就來一件別論.
因為剩下一個.
你就怎樣同心.
怎樣能夠同行.
是.
有時可能相處是困難.
同感.
若然神感動你.
是.
很難相處.
但不要這樣.
叫上祂一起.
你有這個感動.

$^{1121}$而你有遵從的時候.
你就真真正正.
成為他們所講的基督徒.
你問我.
宇傳道.
你有沒有這樣的挑戰.
我就告訴你.
我以前未做傳道人之前.
我認識這麼多基督徒.
這麼多當中.
我還有選擇.
可以.
有些從來我都沒怎樣跟他們傾談的.
我可以慢慢慢慢.
只是跟這些人傾談.
我以前都有些選擇權.
當然其實不對.
我現在發覺不對.
不過我都有選擇權.
但如果做了傳道人之後.
你有沒有選擇權.
其實你都有的.
不過你假裝沒有.
其實你知道那些弟兄姊妹.
是很難搞的.
你又表示愛心.
但心裡面.
最好不要隨便吃飯.
咖啡也不要喝.
是不是?.
但神的感動就是.
祂是不是你的羊.
你是不是要和祂同行.
是的時候.
你的從來大.
還是神的感動大.
因為.
神都是如此去拯救祂.
所以我想聽到的就是.
主啊.

$^{1161}$求你幫我.
都聽到這個弟兄.
這個姐妹.
她內心怎樣經歷神.
只要我聽到.
她是真心相信.
真心相信.
我但願可以和她一起分享.
我又如何真心相信.
亦都真實的掙扎.
而她說.
是的如願道.
我和你一起.
都很難受的.
又喝了東西說.
又不聽人說.
但其實可能這句說話.
這句說話.
令我難受的說話.
可能打破了.
兩者之間.
的人的感悟.
用東西一起祈禱.
重現了.
一個五旬節的事件.
當大家見證到的時候.
她的心就被融化.
而當.
守國禮的.
猶太信徒.
聽到之後.
她的態度都是轉變.
由爭論.
去變為.
她滿心的喜樂.
去讚美你.
主啊我們很想經歷的.
不是在教會分離我.
而是卸下自己的重來的時候.
靠著聖靈.

$^{1201}$突破這個規限的時候.
靠著我們大家去分享.
在主裡面的經歷.
真實的經歷.
真實的掙扎的時候.
我們能夠無分離我.
再次的是同心.
同感.
前行.
求主幫助.
求主帶領.
求主赦免.
祈禱交給我.
主耶穌基督得勝靈以求.
阿們.
\newpage



\section{傳道書 4:4-5:7}
\label{sec:032194}
\textbf{關係的失衡}~~~~譚文倩莊傳道
\newline
\newline
連結: \href{http://www.church.com.hk/acms/content.asp?site=cdc\&op=show\&type=product\&code=032194\&layout=sermon}{\texttt{032194}} ~~~~ 語音日期: 2026-05-24
\newline
\newline
\hyperref[sec:032193]{\small{< < < PREV SERMON < < <}}
~
\hyperref[sec:index_chronic]{\small{[返順時目]}}
~
\hyperref[sec:032195]{\small{> > > NEXT SERMON > > >}}
\newline
\newline
傳道書 4:4-5:7
\newline
\begin{longtable}{cl}
\hline
\hline
章節 & 經文 (和合本修訂版)\\
\hline
4:4 & \begin{tabularx}{0.7\textwidth}{X} 我見人因彼此嫉妒而有一切的勞碌和各樣工作的成就，這也是虛空，也是捕風。 \end{tabularx} \\ \\ \relax
4:5 & \begin{tabularx}{0.7\textwidth}{X} 愚昧人抱著雙臂，自食其肉。 \end{tabularx} \\ \\ \relax
4:6 & \begin{tabularx}{0.7\textwidth}{X} 一掌滿滿而得享安靜，勝過兩掌滿滿而勞碌捕風。 \end{tabularx} \\ \\ \relax
4:7 & \begin{tabularx}{0.7\textwidth}{X} 我轉而觀看日光之下有一件虛空的事： \end{tabularx} \\ \\ \relax
4:8 & \begin{tabularx}{0.7\textwidth}{X} 有人孤單無雙，無子無兄弟，竟勞碌不息，眼目也不以財富為滿足。他說：「我勞碌，自己卻不享福，到底是為了誰呢？」這也是虛空，是極沉重的擔子。 \end{tabularx} \\ \\ \relax
4:9 & \begin{tabularx}{0.7\textwidth}{X} 兩個人總比一個人好，他們勞碌同得美好的報償。 \end{tabularx} \\ \\ \relax
4:10 & \begin{tabularx}{0.7\textwidth}{X} 若是跌倒，這人可以扶起他的同伴；倘若孤身跌倒，沒有別人扶起他來，這人就有禍了。 \end{tabularx} \\ \\ \relax
4:11 & \begin{tabularx}{0.7\textwidth}{X} 再者，二人同睡就都暖和，一人獨睡怎能暖和呢？ \end{tabularx} \\ \\ \relax
4:12 & \begin{tabularx}{0.7\textwidth}{X} 若遇敵攻擊，孤身難擋，二人就能抵擋他；三股合成的繩子不易折斷。 \end{tabularx} \\ \\ \relax
4:13 & \begin{tabularx}{0.7\textwidth}{X} 貧窮而有智慧的年輕人，勝過年老不再納諫的愚昧王， \end{tabularx} \\ \\ \relax
4:14 & \begin{tabularx}{0.7\textwidth}{X} 那人從監牢裡出來作王，在國中原是出身貧寒。 \end{tabularx} \\ \\ \relax
4:15 & \begin{tabularx}{0.7\textwidth}{X} 我見日光之下所有行走的活人，都跟隨那年輕人，就是接續作王的那位。 \end{tabularx} \\ \\ \relax
4:16 & \begin{tabularx}{0.7\textwidth}{X} 他的百姓，就是他所治理的眾人，多得無數；但後來的人還是不喜歡他。這也是虛空，也是捕風。 \end{tabularx} \\ \\ \relax
5:1 & \begin{tabularx}{0.7\textwidth}{X} 你到神的殿要謹慎你的腳步；近前聽，勝過愚昧人獻祭，他們不知道自己在作惡。 \end{tabularx} \\ \\ \relax
5:2 & \begin{tabularx}{0.7\textwidth}{X} 在神面前你不可冒失開口，也不可心急發言；因為神在天上，你在地上，所以你的話語要少。 \end{tabularx} \\ \\ \relax
5:3 & \begin{tabularx}{0.7\textwidth}{X} 事務多，令人做夢；話語多，顯出愚昧。 \end{tabularx} \\ \\ \relax
5:4 & \begin{tabularx}{0.7\textwidth}{X} 你向神許願，還願不可遲延，因他不喜歡愚昧人，你許的願應當償還。 \end{tabularx} \\ \\ \relax
5:5 & \begin{tabularx}{0.7\textwidth}{X} 你許願不還，不如不許。 \end{tabularx} \\ \\ \relax
5:6 & \begin{tabularx}{0.7\textwidth}{X} 不可放任你的口使肉體犯罪，也不可在使者面前說是錯許了。為何使神因你的聲音發怒，敗壞你手所做的呢？ \end{tabularx} \\ \\ \relax
5:7 & \begin{tabularx}{0.7\textwidth}{X} 多夢多言，其中多有虛空，你只要敬畏神。 \end{tabularx} \\ \\
[1ex]
\hline
\hline
\end{longtable}
$^{1}$我想問問大家,人與人之間的關係,你覺得甚麼是最重要?.
是否請飲茶?.
不是?.
沒理由喔.
若果我告訴你,是良好的溝通,你同不同意嗎?.
但甚麼才是良好的溝通呢?.
若讓你定義,甚麼才是良好的溝通呢?.
聆聽,OK,這是很重要的一部分.
我想至少要包括清楚的表達.
剛才亦提到我們要用心聆聽.
好奇對方所說的話.
真誠嘗試理解,或不明白時會問問題.
去釐清對方所說的話的意思.
以及要有合適的回應.
但很多時候是否誰又有空,有那麼多時間.
所以說完,你聽到就聽到,明白就明白,不明白就算了.
接著才會出現那些說你又不聽,聽你又不明白.
明白之後怎樣?.
又不做,做又會做錯,錯了又會怎樣?.
錯了又會賴我,會出現很勞氣的情況.
今天我們繼續講傳道書.
今天的經文是四章四節至五章七節.
關係的失衡,即是失平衡.
我在家裡預備這段經文時.
看到我一頭霧水,我就和我老公說.
傳道者是否有ADHD,是否有過度活躍症.
為何這段經文好像跳下這裡跳下那裡.
還是他不是很夠份,思想有點凌亂.
所以我一路看的時候,究竟他想說甚麼.
當然可能是我的問題,是我自己功力未夠深厚.
若果將傳道書三至六章抽出來.
前面的三章一至二十二節,如果你記得的時候.
是講定時,講時間怎樣限制人同時顯出神的智慧和主權.
後面那段五章十節到六章九節,是講財富和享樂.
我們人怎樣掌控,擁有和滿足.
夾在中間大約四章一節到五章九節裡面.
是講六種不同的關係.
這六種不同的關係就好像一層膠水.
將前面那段講定時和後面那段講財富這兩個大主題.
好像將它們拉在一起.

$^{41}$提醒我們無論是時間的限制還是財富的追求.
都離不開關係這個元素.
就是想講關係是多麼重要.
其實一點都不凌亂.
上次四章一至三節,傳道者講了第一個關係是甚麼關係.
是壓迫者和被壓迫者之間的關係.
今天的四章四節到五章七節.
就轉向人同勞碌,人同人之間,人同權力,人同神之間的關係.
我們先一起低頭祈禱.
主啊,我們心存感恩,我們人算甚麼.
你竟然顧念我們.
主啊,求你謙卑我們的心.
去聆聽你要跟我們講的說話.
奉主耶穌基督的聖名祈求.
我先跟大家讀一讀經文四章四至十二節.
我見人因彼此嫉妒而有一切的勞碌.
和各樣工作的成就.
這也是虛空,也是暴風.
如每人抱著雙臂自食其肉.
一掌滿滿而得享安寧.
勝過兩掌滿滿而勞碌暴風.
有人孤單無雙,無子無兄弟.
竟勞碌不息,眼目也不以財富為滿足.
他說我勞碌,我自己卻不享福.
到底是為了誰呢?.
這也是虛空,是極沉重的擔子.
兩個人總比一個人好.
他們勞碌同得美好的報償.
若是跌倒,這人可以扶起他的同伴.
倘若孤身跌倒,沒有別人扶起他來.
這人就有禍了.
再者,二人同睡都暖和.
一人獨睡怎能暖和呢?.
若遇仇敵攻擊,孤身難擋.
二人就能抵擋他.
三股合成的繩子,不易截斷.
第四節這裡一開始說一種人.
他勞碌得來的成就.
都成為了其他人嫉妒他的原因.
而這個嫉妒Jealousy.

$^{81}$亦都成為了其他人追求成功的動力.
傳道者說,這些都是虛空,都是暴風.
傳道者指出,如果我們人的勞碌.
是由比較和疾動去推動回來的時候.
其實是很沒意思的.
正因為這個原因,所以有些人說.
我不玩這些遊戲.
我不喜歡做一個欺壓人的人.
我亦都不喜歡競爭.
所以我索性慢活.
這些是什麼人?.
有沒有聽過「躺平」這個詞彙?.
「躺平」是中國大陸近年出現的一個詞彙.
意思是在一個高競爭的社會裡面.
人選擇不做一些要求高的工作.
不追求物質.
亦都不去嘗試回應社會對他的期望.
他選擇用一個最低消費的方式去過日子.
其實這些是什麼人?.
其實就是我們常說的寄生蟲.
第五節,他說「如每人抱住雙臂,自食其肉」.
這種人懶到一個點.
他連將食物放入自己嘴裡都懶.
他最終只會由內向外去毀滅自己.
亦都拖累身邊的人和社會.
第六節,「一掌滿滿而得享安寧」.
意思是辛勞工作,但亦要懂得休息.
恰當的比例.
勝過兩掌滿滿而勞碌暴風.
是在說第八節所提到的.
兩隻手不停,只有個「做」字.
又孤零零一個人的工作狂.
傳道者在這裡說這個人沒有伴侶.
沒有子女,只是不停地「做」.
不斷地去囤積賺回來的錢.
從來都沒有停下來想.
或者他都沒有時間去想.
我其實這麼辛苦是為了誰呢?.
傳道者的質疑是要帶我們去思想一下.
在躺平和工作狂之間.

$^{121}$生命裡會否有些東西是更值得我們追求的呢?.
第九節說「兩個人總比一個人好」.
老老實實,有時一個人都幾好.
哇,ERICA第一個up 頭.
我自己都喜歡有時一個人自己出去看場戲.
喝杯咖啡,找個餐館自己吃東西.
去享受一下我自己的me time.
但享受一個人的時間和孤單是兩回事.
舊約甄賢十八章二十四節說.
「朋友太多的人必受損害」.
朋友多,我們經常說這個人是人緣好,聽得多.
為什麼這裡說「朋友多必受損害」呢?.
在這字經文裡的原文.
「朋友」這個字其實有兩個不同的字根在裡面.
第一個朋友的名字叫做Rare.
意思是好鄰舍.
即是這個朋友可以幫你送外賣.
但有事起來的時候,他可以告你上法庭.
同一個字代表一大堆的關係.
這種朋友的關係可能只是停留在.
我們大家都知道對方很多事情的層面.
每次見面都是我們說一下我最近在做什麼.
你最近又在做什麼,誰最近又在做什麼.
我們只是不停在增加資訊.
很膚淺的,沒有深度的.
對我們大家之間的關係也沒有什麼好處.
甄賢在這裡說.
如果你身邊只有這種朋友的話.
最終這些關係會毀滅了你.
之後聖經用了另一個動詞叫Ahav.
這種朋友不單只是陪伴你的人.
而是一個真正關心你生命的人.
當你遇到困難的時候.
他可能比你最親的人更加快出現在你家門口.
當你做錯,走錯路的時候.
他又不會怕得罪你而選擇不出聲.
因為是真正不想你死的朋友.
他給你的提醒雖然有時真的很難聽.
坦白說我其實有誰喜歡被提醒,被糾正.
好像我每一次寫完講章.

$^{161}$我就會講一次給我老公聽.
心裡我經常都希望他讚我.
但結果每一次他都說媽媽在這裡這樣.
媽媽在這裡這樣.
我就覺得他一邊講我一邊不斷在自己縮小.
但在這個過程裡我又在想.
我知道我老公很認識我.
他也很喜歡聖經.
所以他覺得他那個出聲很重要.
我就會慢慢振習起來.
好好地繼續預備講章.
忠言逆耳.
但這些是最愛疼你的人.
不過能夠做到這一步的確不是很簡單.
最近我們家庭團契的媽媽組.
我們的一次聚會大家才在這裡分享.
當我們很坦白地去講自己的感受.
有時其實真是一個很大的冒險.
因為可能你講完出來之後.
你會被誤解.
你會被批評.
甚至受傷.
特別是當我們坦誠之後.
可能換回來的是一些冷淡的回應.
被論斷.
甚至失去一些關係.
結果會怎樣.
我們人就會越來越不出聲.
我們就會學會保護自己.
我們不敢再分享.
但我在這裡反思.
如果在教會這個大家庭裡面.
我們都不能夠衷心去分享.
我們不能夠傾心吐耳.
另一方面我們不去學習聆聽的功課.
我們學習不要去批評.
我們學習接納.
甚至當對方的選擇.
不是我想最好的選擇也好.
我都去尊重對方的選擇.

$^{201}$甚至支持對方.
我們弟兄姊妹的關係.
只會停留在一個很表面的情況下.
我們教會是很難活出.
弟兄姊妹關係的真正意義.
我想問問你.
當我一直這樣說的時候.
你腦海裡面有沒有出現.
一個或者一些這樣的朋友.
或者當你有事的時候.
你想不想起有哪一個.
會立刻來到你家門口.
倒過來.
你又是不是你朋友的這種朋友.
傳道者在這裡說.
可能會有的.
maybe 一個.
兩個.
三個可能已經是極限.
為甚麼.
因為大部份時間.
我們人與人之間是很多競爭.
很多competition.
很多比較在裡面的.
你知不知道男人特別明顯.
因為男人通常是會透過比較.
來定義自己.
為甚麼.
因為裡面有一個潛在的意識.
就是當我比你好的時候.
我就找到我自己的價值.
這就是傳道書裡面所說的競爭.
競爭本身不是一個問題.
而是當競爭的目的.
是為了我要壓住你.
這個欺壓的動作的出現.
甚麼是欺壓的動作.
而競爭為甚麼會變成不公義.
就是當你要贏.
你想人羨慕你.

$^{241}$你就要壓住人.
為達到目的的時候.
我們就會被試探.
口不對心.
會走捷徑.
會偷工減料.
會詭詐.
會說謊話.
我們會踩人.
我們會妥協自己的誠信.
誠信是甚麼.
誠信是我講出來的東西.
和我心裡想的東西是一樣.
我答應了你的事.
就算我吃虧我都會做.
再加上一樣.
就是當你做錯的時候.
你會認錯.
你會求原諒.
這個就是誠信.
但當你只是活在一個.
只是競爭的社會的心態裡.
你整天只顧著要贏.
你還會不會.
選擇講真心的話.
會守承諾.
會認錯.
答案是有時都會的.
但當有利益衝突的時候.
你就可能選擇不會.
其實講到底.
欺壓不單止是想說.
我們要好過人.
而是有意無意.
我們要壓低對方.
為了要保護.
去鞏固我自己.
因為當別人處於一個比我低的位置.
我就自己愈來愈感覺良好.
同時我亦不斷餵養我自己的心態.

$^{281}$另一個智慧真言.
一條繩 兩條繩.
最堅固的是三條繩繞在一起的繩.
如果兩個人比一個人好.
三個人就會更加好.
人生很多事情的確很虛空.
但有些事情相對來說沒有那麼虛空.
例如建立關係.
我們彼此同行.
追求群體的生活.
比起單純追求財富,成就,享受.
是有更長遠的價值.
因為當我們遇到攻擊的時候.
你很需要一,兩個人在你身邊.
跟你一起走.
去明白你,去撐你.
在你軟弱的時候.
透過安慰,提醒和保護.
成為你生命裡面的一個穩定的力量.
再者在權力,制度和群體的社會裡面.
我們更深的焦慮.
就是今天我們是被需要的.
但可能明天已經沒有人記得你.
今天你是重要的.
下一刻你已經被取代.
現實是很殘酷的.
接著十三至十六節.
他說貧窮而有智慧的年輕人.
勝過年老不再立年的原煤王.
那人從監牢裡出來作王.
在國中原是出身貧寒.
我見日光之下所有行走的活人.
都跟隨那年輕人.
就是接續作王的那位.
他的百姓.
就是他所治理的眾人.
多得無數.
但後來的人還是不喜歡他.
這也是虛空.
也是暴風.

$^{321}$這裡說的智慧是甚麼智慧.
這裡說的智慧是謙卑受教.
然後將選擇,行為和結果或後果放在一起去想.
我再說一次 就是謙卑受教.
然後將選擇,行為和結果或後果放在一起去想.
當中的重點是謙卑受教.
不要自以為聰明.
這裡說的年老的王.
他有甚麼優勢.
他年紀大.
所以怎樣.
年紀大會怎樣.
食鹽多過你食米.
很有經驗.
他是王.
有權.
為何他是愚昧.
不是因為他有老人癡呆症.
而正正是因為他恃著自己食鹽多過你食米.
所以他不再聽人說.
他不再願意接受教導.
不願意再接受勸告.
這個王出身卑微.
坐過監 有些似虐室的故事.
但有一天他都會怎樣.
有一天他都會老.
到時他所治理的百姓.
就會轉會去支持另一個年輕人.
成為新一代的明星.
跟著這個明星慢慢都會過氣.
再有新的明星出現.
沒多久這個明星又會再老.
又再被取代.
即是說你以為你自己很厲害嗎.
有一天你都會被取代.
這個世界沒說沒了誰不行.
你以為自己年紀越大.
你就越有智慧.
有地位就代表成功.
傳道者在這裡提醒.

$^{361}$當你以為自己已經懂得.
不需要再學.
其實就是開始衰落的時候.
碩士虛空 也士暴風.
我們基督徒的生命成長.
是一生之久的事.
不是說去到65歲.
你退休的時候就停下來.
甚至我們牧師傳道人又如何.
牧師傳道人就更加應該是一個願意受教的僕人.
誰說去到某個年紀.
讀了很多書.
或者去到某個位置.
我就已經懂得.
不用再學.
不用再去成長.
五章一至七節.
你到神的殿要謹慎你的腳步.
要向前聽聖餘媒人獻祭.
他們不知道自己在作惡.
在神面前你不可無失開口.
也不可心急發言.
因為神在天上你在地上.
所以你的話語要少.
事無多令人做夢.
話語多顯出愚昧.
你向神許願.
還願不可遲疑.
因他不喜歡愚昧人.
你許的願應當償還.
你許願不還.
不如不許.
不可放任你的口使肉體犯罪.
也不可在使者面前說是錯許了.
為何使神因你的聲音發怒.
敗壞你所做的呢.
多夢多現.
其中多有虛空.
你只要敬畏神.
在這裡開始.

$^{401}$我很想大家要留心.
因為傳道者在這裡.
要給一個非常重要的智慧訊息.
交到給我們.
其實之前那段.
傳道者的語氣.
都是比較傾於勸勉.
但由這裡開始.
他的語氣是嚴厲了很多.
傳道者可以很嚴厲地去批判.
人來到神的聖殿的時候要謹慎.
要想清想楚.
才能開口.
我們不要用我們的說話.
我們的口.
去說一些話出來去嘗試操控神.
嘗試控制神的作為.
尤其輕率許願和不守承諾.
在神面前不要霧失開口.
不要心急發言.
話語要少是甚麼意思.
你不要來到神的面前.
一卵嘴在說.
我們通常對著神的時候.
說些甚麼.
就是告訴神.
我需要你應養我甚麼.
我有甚麼需要.
我求你這樣這樣.
我們說的東西.
全部都是要神為我們做一些事.
去達到一些事.
你祈禱的時候.
整個祈禱是否只是不停地說.
自己想神怎樣怎樣的祈禱.
真正的敬拜.
揭露了我們人的自我中心.
和神尊榮之間的差距.
一個巨大的差距.
因為真正的敬拜.

$^{441}$不是我們人用說話.
用一些宗教活動.
或者儀式去跟神做交易.
或者我們去嘗試用祈禱.
去展示我們自己有多麼的敬虔.
來到神的面前.
第一件事是要做甚麼.
去聽禱.
不是.
去聽敬拜的音樂.
也不是.
最重要的是.
聽神跟你說話.
神啊 我來到你的面前.
因為我需要聆聽.
你要跟我說的話.
無論是透過講道的訊息.
還是透過詩歌裡面的歌詞.
祈禱.
還是跟弟兄姊妹相交.
神啊 我好想聽到你的聲音.
因為在我生命裡面.
現在發生著一些事情.
我很需要你跟我說話.
經文教我們.
「近前聽」.
打開耳朵去聽.
甚麼叫聽.
在這裡聽的這個字.
compare to我們英文的listen.
或者中文聽這個字.
是有一個更深層的意思.
在聖經裡面.
不單只是在說聆聽.
而是在說我們是主動地.
謙卑下來去聆聽.
接著是帶著一個信靠.
信服的心去回應.
說得直接一點.
是說我們要用心去聽.

$^{481}$所以敬拜的優先次序是.
我們謙卑去聆聽神要跟我們說的話.
而聽了之後.
是會有回應的.
多夢多言其中多有虛空.
你只要敬畏神.
講就講了很多次.
但是如果我們沒有一個敬畏神的心.
其實都是虛空的.
都是沒有意思的.
之前四章四至十六節.
經文提到錯誤的地上敬拜.
我們人將成就,財富,權力.
當作是神.
五章一至七節這裡.
錯誤的殿中敬拜.
是說我們人用嘴巴去敬拜上帝.
但其實可能只是一場宗教的表現.
這段經文傳道者提到.
其實歸根究底.
說了這麼多種不同的關係.
其實都是在反映我們跟神之間的關係.
因為敬畏神不單只是一個態度.
而是我們基督徒的生命核心的方向.
因為當我們人的生命是被神所塑造的時候.
我們就不會再被我們的自我中心.
不斷地去操控.
神的標準影響我們的言語.
影響我們所做的選擇.
影響我們對每一種,每一段關係的態度.
影響我們怎樣用我們的時間.
和我們生命的優先次序.
影響我們的價值觀.
換句話來說.
敬畏神的生命是用讓神去塑造我這個人.
怎樣去做人.
而一個人怎樣去做人.
就決定了他怎樣對待他身邊的關係.
關係的失衡.
由探討日光之下.

$^{521}$我們人以自我中心所經營的關係.
轉移去思考我們人和神之間的關係.
人和人之間的關係的失衡.
不單只是在說社交的問題.
根本上就是人和神之間的關係失衡的結果.
因為歪了,彎了的心.
我們容不二會嘗試用言語,儀式,承諾去操控人.
操控上帝.
結果我們人就由欺壓人延伸到欺壓上帝.
傳道者智慧的導向.
我們人要選擇由敬畏神主導的方向去做起點.
因為敬畏神是所有的關係準則和決定最有智慧的方向.
我們一起低頭祈禱.
七海天父我們多謝你.
多謝你今天藉著你的話語提醒我們.
再一次提醒我們不要停留在一個只有一張嘴的狀況裡面.
你再一次提醒我們.
要我們思考我們來到你的面前的時候.
我們是否謙卑,是否專心.
我們是否在聆聽你要跟我們說的話.
主啊求你幫助我們明白.
我們人與人之間或許我們跟很多我們生命裡面.
不同的關係裡面的失衡.
其實都是因為我們跟你之間的關係.
有些事要搞好.
因為唯有我們追求去敬畏你.
我們願意在你面前學習.
好像耶穌那樣去愛人,去尊重人.
主啊我們需要你的幫助.
求你幫助我們來到你面前的時候.
不要說那麼多無謂的話.
教我們活出一個真實的生命.
叫我們心存謙卑誠實.
行事為人都是土裡喜悅的.
願你幫助我們每天都是帶著敬畏和信靠的心去跟隨你.
奉主耶穌基督的聖名祈求.
Amen.
\newpage



\section{約書亞記 23:11}
\label{sec:032195}
\textbf{聖經人聖經事:約書亞贈別 (1) 愛耶和華}~~~~黃克勤牧師博士
\newline
\newline
連結: \href{http://www.church.com.hk/acms/content.asp?site=cdc\&op=show\&type=product\&code=032195\&layout=sermon}{\texttt{032195}} ~~~~ 語音日期: 2026-05-26
\newline
\newline
\hyperref[sec:032194]{\small{< < < PREV SERMON < < <}}
~
\hyperref[sec:index_chronic]{\small{[返順時目]}}
~
\hyperref[sec:032196]{\small{> > > NEXT SERMON > > >}}
\newline
\newline
約書亞記 23:11
\newline
\begin{longtable}{cl}
\hline
\hline
章節 & 經文 (和合本修訂版)\\
\hline
23:11 & \begin{tabularx}{0.7\textwidth}{X} 你們要分外謹慎，愛耶和華－你們的神。 \end{tabularx} \\ \\
[1ex]
\hline
\hline
\end{longtable}
$^{1}$大家好,又一齊讀聖經啦.
新衛書宣教網絡為你預備的聖經人,聖經事.
今次是看約書亞贈別.
贈別是為了我只有這麼短一行字.
我要壓縮一點.
這是臨別贈言.
約書亞年紀大了.
臨別贈言跟他的百姓.
他所愛的百姓講說話.
有1,即是有2,有3,真的有的.
約書亞臨別贈言的第一.
就是叫他們愛耶和華.
愛耶和華.
好了,看一下.
約西記23章.
耶和華使以色列人安靜.
不與四圍的一切仇敵爭戰已經多日.
約書亞年紀老邁了.
就擺以色列眾人的長老,族長,審判官.
並官長都召了來.
對他們說:我年紀已經老邁了.
好了,留意一下.
耶和華使以色列人安靜.
這個安靜呢.
是不與四圍的一切仇敵爭戰已經多日.
這個安靜的意思.
其實就譯作安靜是不對的.
應該譯作休息.
Luwak.
Luwak這個字呢.
希伯來文Luwak呢.
是休息的意思.
神讓以色列人打了一輪仗.
獲得了全地之後.
迦南地英雄之地得到了.
辛苦到不得了.
打了一場仗.
現在神呢.
是讓他們.
不與一切仇敵爭戰.

$^{41}$已經是叫他們休息下.
讓他們休息.
中文就譯作安靜.
Luwak呢.
是休息.
已經多日.
即是用多日的日.
是很多日.
一日一日.
很久的意思.
而約書客呢.
也都年紀老了.
就叫了所有人來.
就講臨別贈言.
臨別贈言第一句是什麼呢.
就是我年紀已經老了.
要接受自己老的嘛.
我都接受自己老的嘛.
年紀老了.
耶和說.
你們的神因你們的緣故.
向那些國所行的一切事.
你們親眼看到了.
以往神帶我們從埃及出來.
一切的過程你都看到了.
還記不記得耶利哥一叫.
那城倒塌了.
記得.
很厲害.
其實為你們征戰的是.
耶和說.
你們的神.
大家要記得.
其實呢.
現在我們很平安.
很久了.
佔領了這個地方.
在這裡住了.
想一想.
多厲害.

$^{81}$我們多厲害.
真的全地都被我們打爛了.
被我們進來住了.
嘩.
很強.
不要自以為是.
不要自以為很厲害.
其實呢.
一切你都看到了.
是為你們征戰的是耶和說.
你們的神.
完全是神的工作來的.
所以.
記得了.
對於神.
你要大大壯膽相信祂.
第二.
謹守遵行寫在摩西律法書上的齊華.
不過偏離左右.
也沒有一來到二.
因為牌的經文的時候.
最重要.
愛耶和華那裡.
還沒到.
所以這是第二.
謹守遵行寫在摩西律法書上的齊華.
不要偏離左右.
第三.
第七節.
不可與你中間留意.
所剩下的國民才集.
他們的神.
你們不可以提他的名字.
不可以向他起誓.
不可以事奉.
這個事奉是敬拜的意思.
不可以敬拜.
叩拜他們.
叩拜.
不可以.

$^{121}$不可以和他們才集.
只要你們做什麼.
照你們.
到今天所行的.
是專靠耶和華.
你們的神.
倚靠神.
你們要分外謹慎.
愛耶和華.
你們的神.
這個才是我要講的第一點.
約書亞臨別贈言的第一.
就是叫他們.
愛耶和華.
什麼因素呢.
有因素的.
就是怎樣愛耶和華.
愛耶和華.
你們記得嗎.
是神用祂達能的手帶我們出來.
進這個迦林地.
得勝所有的王.
不是我們的工作.
是神的大能.
你要愛耶和華.
所以.
要怎樣愛.
耶和華怎樣愛.
愛神多一點.
怎樣愛神.
謹守遵行寫在摩西律法書上的誓話.
如果你說.
愛耶和華要怎樣做.
第一.
謹守遵行律法.
這個是約書亞教他們的.
就是神透過摩西給他們的律法書.
我愛耶和華.
祂好愛.
怎樣呀.

$^{161}$買多點水果獻給祂.
不是.
我要多獻幾隻牛.
都好.
不過.
最重要的.
神說什麼呢.
我不是喜愛祭祀.
我不是喜愛那些牛羊.
我喜愛你們的順服的心.
所以呢.
如果你說我要愛耶和華.
要做的就是.
謹守遵行寫在摩西律法書上.
是用遵行律法表達你的愛的.
因此.
你說主耶穌時時罵那些法利賽人.
其實法利賽人是.
因為他虛偽.
自以為義.
我遵行了.
我叻個你.
自以為義.
主耶穌是說他這件事.
其實法利賽人都不錯的.
他真的謹守遵行.
摩西律法書上所寫的東西.
不差的.
他們都愛耶和華的.
不過.
他自以語話其他人不好.
這個是主耶穌罵得最多的.
愛耶和華怎樣愛.
愛耶和華.
你跟著律法書去做.
然後第二就是.
不可以與剩下的人.
國民殘雜.
財雜.
這個是什麼意思.

$^{201}$就是因為.
你記得他們進去的時候.
神本來叫他們將所有的人都殺死.
生畜都不要.
後來.
那些財寶可以拿.
人就不要了.
為什麼.
因為怕他們會.
弄來弄去.
還在傳流.
會帶你們去拜他們的神.
這樣就糟糕了.
所以不要漏.
但是有些人.
他是沒有殺死的.
如果你有留意的話.
起碼未得之地很多.
前好幾次說了.
未得之地那些人.
他們是沒有將他完全滅絕的.
於是.
留下很多非利士人.
他們拜另一種神.
男性人生如美的大隻佬.
海上上來的民族.
他們拜非利士的人叫大菌.
還有很多拜巴力.
很多拜亞瑟拉.
很多拜摩洛.
拜不同的神.
後來就弄壞了他們.
所以約書亞特別叫.
你們要臨別贈言.
要愛的話.
怎樣愛呢.
第一.
守摩西的神.
透過摩西給我們的律法.
第二.

$^{241}$不要和靜下來沒有殺盡的國民.
殘襲他們的神.
你不可提他們的名字.
提名字都不可以.
所以第十六篇就說.
我不提別神的名號.
我連祂的名字都不提.
不可以指著那些神起誓.
不可以事奉敬拜那些神.
不可以.
記得.
要照耶和華所吩咐.
專靠耶和華.
如果你要尋求幫助.
專靠耶和華.
只找一個就不好.
巴力找找.
摩洛找找.
不要和那些人一起.
這樣就行了.
分愛謹慎.
愛耶和華你們的神.
第十二節都很重要.
你們若稍微轉去.
綠色字不要看.
你們若稍微轉去.
與你們中間所剩下的這些國民.
聯絡彼此結親.
互相往來.
這樣就不行了.
甚至和剩下的國民.
不可以聯絡.
不可以結親.
所以以色列人一直不和外族通婚.
不可以.
只可以自己支派之間結婚.
不要和他們聯絡.
不要抱頭抱頸.
彼此結親.
互相往來.

$^{281}$攬親熱.
攬親密好朋友.
你們稍微轉去.
我解釋一下「稍微轉去」.
為何「稍微轉去」呢?.
「稍微」這個字是「soup」.
「soup」這個字其實是「返回」.
兩個字「稍微」已經是「返回」.
「轉去」這個字希伯來文又是「soup」.
即是甚麼意思?.
又是返回.
所以如果翻譯就是.
你們若返回返回.
與他們聯絡.
我就不歡喜了.
約書林別正賢就說.
不要和本地未殺的人.
剩下的人與他們聯絡.
如果你是.
轉回轉回返回返回.
這樣就不行了.
為何會翻譯成「稍微」呢?.
翻譯得不夠好.
不過這個「soup返回」.
第一個時式是不定式.
即是如果你們返回與他們聯絡.
本來就是要將他們殺了.
不過你返回與他們聯絡.
是講得很普遍的.
普遍意思就是說.
沒有講時間.
沒有講地點.
沒有講內容.
就是如果你返回與他們做朋友.
這個就是不定的時式.
而「轉去」呢.
就是如果你們返回.
是將來返回.
你何時返回也好.
或者將來返回.

$^{321}$所以第二個「soup」是將來式.
就是如果你們返回.
或者將來返回.
與他們剩下的國民聯絡.
彼此結親互相往來.
這樣就糟糕了.
中文就是想說.
你想都不要想.
一點點返回都不行.
你就要確實知道.
你的神就不再將他們.
從你當中說出來了.
神就會.
如果你與他們來往.
結盟,結親.
一起敬拜那些神.
神就會不給面子了.
不會再與你想辦法.
趕他們走的了.
那些人反而成為你們的網絡.
一個機關.
機欖意思就是一個機關.
讓你掉下去.
陷阱.
它成為一個網.
綁著你成為陷阱.
「蠟上的鞭」.
即是鞭在你的身上.
一條痕.
「蠟上的鞭」.
「眼中的刺」.
很辛苦.
經常眼中有條刺.
眼睛都睜不大.
直到你們在耶和華.
你們神所賜給你.
這個美地上滅亡.
所以約書亞最臨別贈言的.
第一個就是.
「愛耶和華」.

$^{361}$怎樣愛呢?.
「愛耶和華」.
怎樣愛呢?.
第一.
遵行摩西.
神透過摩西給的律法.
第二.
你要謹慎.
不要與這些.
靜下來的人來往.
保持自己的分別與清潔.
否則.
神會讓他們.
侵蝕你.
信仰上侵蝕你.
而最終.
他們會成為你們的害處.
那時候.
這樣的話.
你轉來轉去.
神會讓你們在.
祂賜給你的地上.
滅亡.
於是叫他們.
臨別贈言.
要記得「愛耶和華」是正面的.
不要走回頭.
怎樣愛呢?.
遵行律法.
第二.
分別為聖.
不要與當地人一樣.
「愛」的意思就是.
遵行律法.
分別為聖.
講得很清楚.
否則.
神會讓你滅亡.
大家一起祈禱吧.
多謝主.

$^{401}$透過約書亞都警戒我們.
他臨別贈言.
叫所有人來告訴他.
要愛耶和華.
想起你曾經帶領我們.
經過了人生那麼多高低.
好像出埃及行曠野.
經過死亡.
洗禮.
過躍誕河.
然後來到這個應許之地.
幫助我們愛你.
進行你的說話.
不與這個世界.
聚集一起.
讓我們有所分別.
主幫助我們.
讓我們專心跟隨你.
專心愛你.
專心倚靠你.
加我們的力量.
祝福我們.
祈禱奉耶穌基督名求.
阿們.
\newpage



\section{以賽亞書 56:1-8}
\label{sec:032196}
\textbf{安息有時}~~~~文麗芳傳道
\newline
\newline
連結: \href{http://www.church.com.hk/acms/content.asp?site=cdc\&op=show\&type=product\&code=032196\&layout=sermon}{\texttt{032196}} ~~~~ 語音日期: 2026-05-24
\newline
\newline
\hyperref[sec:032195]{\small{< < < PREV SERMON < < <}}
~
\hyperref[sec:index_chronic]{\small{[返順時目]}}
~
\hyperref[sec:032197]{\small{> > > NEXT SERMON > > >}}
\newline
\newline
以賽亞書 56:1-8
\newline
\begin{longtable}{cl}
\hline
\hline
章節 & 經文 (和合本修訂版)\\
\hline
56:1 & \begin{tabularx}{0.7\textwidth}{X} 耶和華如此說：「你們當守公平，行公義；因我的救恩臨近，我的公義將要顯現。 \end{tabularx} \\ \\ \relax
56:2 & \begin{tabularx}{0.7\textwidth}{X} 謹守安息日不予干犯，禁止己手不作惡，如此行、如此持守的人有福了！」 \end{tabularx} \\ \\ \relax
56:3 & \begin{tabularx}{0.7\textwidth}{X} 與耶和華聯合的外邦人不要說：「耶和華將我和他的子民分別出來。」太監也不要說：「看哪，我是枯樹。」 \end{tabularx} \\ \\ \relax
56:4 & \begin{tabularx}{0.7\textwidth}{X} 因為耶和華如此說：「那些謹守我的安息日，選擇我旨意，持守我約的太監， \end{tabularx} \\ \\ \relax
56:5 & \begin{tabularx}{0.7\textwidth}{X} 我必使他們在我殿中，在我牆內，有紀念碑，有名號，勝過有兒有女；我必賜他們永遠的名，不能剪除。 \end{tabularx} \\ \\ \relax
56:6 & \begin{tabularx}{0.7\textwidth}{X} 「那些與耶和華聯合，事奉他，愛他名，作他僕人的外邦人，凡謹守安息日不予干犯，又持守我約的人， \end{tabularx} \\ \\ \relax
56:7 & \begin{tabularx}{0.7\textwidth}{X} 我必領他們到我的聖山，使他們在我的禱告的殿中喜樂。他們的燔祭和祭物，在我壇上必蒙悅納，因我的殿必稱為萬民禱告的殿。 \end{tabularx} \\ \\ \relax
56:8 & \begin{tabularx}{0.7\textwidth}{X} 我還要召集更多的人歸併到這些被召集的人中。這是召集被趕散的以色列人的主耶和華說的。」 \end{tabularx} \\ \\
[1ex]
\hline
\hline
\end{longtable}
$^{1}$各位弟兄姊妹早晨.
我們一齊有個祈禱開始.
我們一齊看一看祈禱.
慈悲人愛天父我們多謝你.
讓我們今天在星期日的早上.
我們可以回到教會裡面.
跟弟兄姊妹一齊敬拜你.
是何等的恩典.
主啊求你讓我們心裡面.
常存著感恩的心.
當我們去聆聽主的教導的時候.
主啊就求你開我們的心.
開我們的耳朵.
求主使用你視內為你所預備的講章.
讓無論聽的講的都得著教導.
我們今天禱告奉主耶穌基督得勝的名祈求.
阿們.
崇拜講道我們一直都是講以賽亞書.
主要都講以賽亞書.
今天已經去到五十六章了.
五十六章是一個新的開始.
為什麼這樣說呢.
因為我們知道.
以賽亞書的作者分為第一以賽亞.
第二以賽亞和第三以賽亞.
今天就是第三以賽亞的篇章開始了.
就是第五十六到第六十六章.
第三以賽亞這一段的經文.
主要的讀者是什麼時代的人呢.
跟之前兩個已經不同了.
是被擄回歸時期的人.
在進入我們講到經文之前.
我想交代一些經文的背景.
讓大家可以了解多一些.
我們假設第三以賽亞的對象.
是被擄回歸之後的人.
住在耶路撒冷.
按學者的研究.
究竟是什麼時候.
有沒有一個精準的年份.

$^{41}$有不同的說法.
我們補充多一些資料.
以色列人被擄回歸.
其實主要分為三次.
第一次回歸是主前538年.
索羅巴伯帶著第一批百姓回耶路撒冷.
第二次回歸是主前458年.
是以色列人帶第二批回耶路撒冷.
第三次回歸我們比較熟悉.
就是主前445年或444年.
由尼希米帶第三批百姓回耶路撒冷.
並且重建城牆.
按經文的描述.
當時以色列人的情況是怎樣呢?.
按經文的描述.
以色列人的時間大約是什麼呢?.
是公元前500年左右.
或是公元前六世紀末到公元前五世紀初.
我們簡單假設是公元前500年左右.
當時的社會情況是怎樣呢?.
公元前500年左右的耶路撒冷.
基本上是波斯帝國統治下的一個地方.
它沒有獨立的王國.
沒有自己的王室.
只是屬於波斯帝國轄下的一個行政區.
當時猶太人已經從巴比倫被擄回去了.
是剛才說的第一批人已經回去了.
他們的人數是多少呢?.
人數大概是五萬五千人左右.
但是有四萬九千八百九十七人.
當中有一些核心的就是四萬二千三百六十人.
其餘的就是一些奴僕.
一些歌唱者.
而是整個群體一起回來的.
想想有五萬人的時候.
由巴比倫到耶路撒冷.
回到他們自己的地方.
這個時期耶路撒冷正在重建之中.
第二聖殿是何時開始興建的呢?.
就是537年.

$^{81}$即是這班人回去之後.
有一段時間就開始興建聖殿了.
去到516年才完成.
如果你看看那些先知書.
為什麼會用這麼久來興建聖殿呢?.
就是因為他們有很多事情要做.
所以他們就把興建聖殿放在一邊.
當時的城市和聖殿雖然開始慢慢的恢復.
或者其實五百年的時候.
其實聖殿已經興建好了.
但是規模,人口和經濟實力.
遠遠不及他們大胃王的年代.
或者所羅門王的年代.
由於歸回也面對了很多外來的壓力.
內部又很貧窮.
社會秩序有很多的問題.
所以耶路撒冷並不繁榮.
或者當時的生活並不算是好的.
在信仰上這個時期的核心就是.
去回歸,重建,重整他們的生活.
百姓主要都是環繞聖殿,律法.
和敬拜去建立他們自己的群體.
說了這麼多.
究竟56章到66章.
究竟讀者是誰呢?.
這個部分的議題是主要的對象.
當然剛才我們說的那五萬人.
就是那幫猶太人.
那些人去被擄回到自己的故土.
正是從新建立自己宗教和社會群體.
社會生活的一幫人.
他們正在經歷他們這一段階段的時候.
其實他們面對很多的壓力.
面對很多的困難和衝突.
他們被擄回到自己的情況.
發覺那個地方不是他們想像中那麼好的.
是一片很荒漠,很貧窮的地方.
經濟當時很混亂.
各種的文化和宗教都夾雜在一起.
失去了他們可以想像的.

$^{121}$在巴比倫的生活雖然不好.
但是都很平穩的生活.
可能可以說是好的生活.
但是他們回到自己的故土的時候.
他們好像開荒牛一樣.
經文的讀者不單單是這一幫人.
也有其他生活在耶路撒冷這個地方的人.
包括那些沒有被擄的以色列人.
那些人為什麼不被擄呢?.
因為他們不夠重要.
因為他們不夠富有.
或者他們沒有什麼才能.
所以他們被留下在耶路撒冷這個地方.
他們苦苦掙扎了幾十年.
現在看著他們以前民族的人回到這裡.
好像打鑼打鼓般回來.
好像期望著帶來一些新的改變.
但是卻是混亂一片.
他們還帶著巴比倫的文化.
還有他們可能帶著巴比倫的物資.
來到炫耀一番.
可以想像那些留下來的人的心情是怎樣呢?.
還有這個地方其實已經被很多外族人居住.
在這裡寄居.
裡面耶路撒冷這個地方.
更加有些無權無勢的外邦人住在這個地方.
他們面對這些猶太人來到這個地方.
要重新建立他們的宗教和社會.
令他們感到很迷茫.
因為他們是被排擠的一群.
這段經文.
我們今天讀的經文.
這一群被邊緣化的人.
就是一群外邦人.
或者中間有一些被排擠的.
好像一些太監.
或者一些窮人.
或者一些受壓迫的人.
今天我們讀的經文.
就是五十六章第一到第八節.

$^{161}$經文告訴我們.
要守安息日.
是作為上帝子民身份的一個認同.
還有一個蒙福的管道.
也是上帝重新塑造他們的子民.
五十六章一到二節是這樣說.
耶穌說:如此說你們當守公平行公義.
因我的救恩臨近.
我的公義將要顯現.
謹守安息日而不干犯.
禁止己守而不作惡.
如此行如此持守的人便為有福.
經文一開始就焦聚他們的子民.
就說你們要當守公平行公義.
大家留意.
在經文裡面.
沒有很熟悉的說.
以色列你要停.
以前常常說.
以色列啊.
很明顯有改變.
五十六章有改變.
就是這群受眾.
不單單是血統上的以色列人.
上帝講明了救恩.
施行審判的日子就快到了.
謹守安息日而不干犯.
禁止己守而不作惡.
如此行如此持守的人.
就會得到祂的賜福.
守公平行公義是先知一直的教導.
我們今天暫時不看.
我們集中去看守安息日的吩咐.
經文告訴我們.
謹守安息日是很重要的.
是神賜福的原因.
你信不信得到上帝的福.
如果你想得到.
你就要守安息日.
究竟怎樣才是守安息日呢.

$^{201}$守安息日這個吩咐.
相信大家都不陌生.
在十誡開始.
神就有這個命令.
但我們有沒有遵守過呢.
對於主前500年耶路撒冷的猶太人.
他們為何不遵守呢.
上帝叫他們遵守.
即是說他們沒有遵守.
守安息日是不是單純地.
像我們這樣每天.
每星期放下.
有一天不工作.
單純地去教會.
就是守安息日呢.
為何上帝一定要我們做這件事呢.
我們現在這個世代.
我們有時候說.
我們是24/7工作的.
那守安息日是不是實際呢.
如果要滿足上帝的要求.
我們是不是有時.
我們都會感到力不從心呢.
我們可不可以.
我們怎樣可以看守安息日.
作為上帝恩典的渠道.
怎樣去看安息日.
不單單是一條規條.
而是上帝給我們的祝福呢.
如果我們看尼希米記十三章的話.
我們認識到回歸的初期.
以色列人是面對很多的事.
他們百廢待興.
他們當時耶路撒冷很荒涼.
他們要重建城牆.
重建聖殿.
他們還要重建自己的家園.
他們經濟還有很大的壓力.
很多人都忙於耕種.
經商去謀生.

$^{241}$他們會在安息日裡面去做買賣.
又要耕田又要運貨.
加上土地很貧瘠.
收成亦都很有限.
還要交稅給波斯帝國.
他們覺得過安息日.
好像我們經常說.
好像浪費了賺錢的時間.
他們更加會面對外族的威脅.
不單止攻擊.
還要阻撓他們的重建.
持續騷擾他們的生活.
加上當時很流行和外族通婚.
因為當時那個地方已經很久沒有.
以色列人居住.
很多外族的人在那裡住.
很多人和外族通婚.
導致信仰很混淆.
安息日作為神立約的記號.
逐漸對於他們來說很模糊.
初期回到去的祭司和尼美人.
其實他們都有很多的適應.
所以他們都未完全到位.
律法的教導中斷了幾十年.
他們沒有正式崇拜的生活.
守安息日.
守節旗.
守潔淨的規矩.
其實對於他們來說很生疏.
要到以斯拉和尼希米回來的時候.
他們才會重新拿律法書出來讀.
改革混婚和安息日.
買賣這些陋習.
可見他們這個時候的以色列人.
其實他們沒有遵守到律法的教導.
另一方面他們面對的是什麼呢.
面對的是靈性和身份的危機.
經過了被擄之後.
有些人開始懷疑.
守約是否有福呢.

$^{281}$上帝不是應許他們回歸的時候.
祝福他們嗎.
但好像還沒有實現.
他們有些人就轉為去務實主義.
他覺得解決現時的困難和問題是優先.
他們追求宗教或者信仰上的一些行為.
他們慢慢放下了.
跟隨了附近的人.
在這裡生活的人的節奏.
跟隨了他們的風俗和習慣.
尼希米記五章更加指出.
社會裡有不公平和內部的剝削.
有錢的人壓制窮人.
欠債的人要賣兒女還債.
窮人要日夜工作.
安息日也沒辦法安息.
領袖由貴族開始他們就是這樣去做.
下層的百姓更加沒辦法.
要被逼中日的去工作.
在這些資料看來.
安息日似乎那個時候被遺忘了.
所以賽亞書五十六章一至八節.
正正是針對這個處境.
強調守公平,行公義,守安息日.
是活出聖約子民身份一個重要的一環.
經文三至六節提到兩種人.
不知道你們有沒有留意到.
第一種人是和耶和華聯合的愛邦人.
第二種是太教.
顯明上帝的說話都給兩班人去聽.
五十六章第三節這樣說.
耶和華聯合的愛邦人不要說.
耶和華必定將我從他民中分別出來.
以往愛邦人被以色列人認為.
是神恩典之外的人.
他們的血統上不是選民.
所以他們不能夠承受上帝的福.
根據申命記二十三章三至八節.
確實是寫著阿們人,摩訶人.
他們永遠都不能夠進入耶和華的會.

$^{321}$但是經文這裡有不同.
經文裡面跟他們說.
不要說耶和華.
必定將他們從神子民裡面分別出來.
亦都要告訴他們.
耶和華現在就要將他們歸入神子民裡面.
五十六章六節更加說.
這些愛邦人可以事奉神.
作神的僕人.
就好像以色列人一樣.
說到太監.
其實太監基於申命記二十三章.
亦都有這樣的記載.
「凡愛臣受傷或被閹割的.
不可入耶和華的會」.
亦都這樣說.
太監不能夠參與神敬拜的活動.
是宗教上被排擠的人.
其實以色列是沒有太監這個職位.
為什麼會有提到太監呢?.
可能是當時以色列.
有一段時間很拜偶像.
好像夜洗別的時候.
將愛邦的信仰都帶了進來.
當時就連這些職位都帶了進來.
這裡所排擠的就是這類型的人.
或者在被擄的時間.
被擄歸回的時間.
可能他們因著外族的影響.
要成為.
其實太監很多時候.
都是一些高官.
可能是某些生活的原因.
要令他們成為太監.
神在這段時間.
針對他們說.
經文說.
太監不要說我是菇鼠.
本來太監給人的感覺是沒有希望的.
因為他們沒有後代.

$^{361}$經文就用菇鼠來形容.
但是神這樣說.
那些謹守我安息日.
揀選我所喜悅的事.
持守我約的太監.
我必使他們在我殿中.
我牆內有紀念有名號.
比有兒女的更美.
我必使他們永遠的名不被剪除.
他們可以在神的殿裡面.
永遠被紀念.
比有兒女的更加美.
當人覺得希望有兒女.
藉著有兒女有後代.
可以令自己的名.
留在歷史裡面的時候.
神卻不是用後代的方式.
而是用另一種方式.
令太監的名被紀念.
永遠都不會廢去.
神不是想廢掉律法.
而是表明.
在神這個新造的群體當中.
裡面的忠誠和敬拜.
取代了血緣和生理的因素.
我們可以假設.
當時很多人因為戰亂.
而成為被排擠的人.
但在這個時候.
他們重新建立以色列群體的時候.
這群人重新被接納.
並且被恢復他們的尊嚴.
留意的是.
不是所有太監在神的殿.
都會永遠被紀念.
而是那些謹守神的安息日.
揀選神可喜悅的事.
和持守神的約的太監.
才會有永遠的名被紀念.
不是所有外邦人都能成為子民.

$^{401}$而是那些與耶和華聯合的外邦人.
聯合是什麼意思呢?.
就是連結的意思.
就是與耶和華連結的外邦人.
他們有事奉神.
他們可以作神的僕人.
我們可以看到.
在這裡說了些什麼呢?.
這裡似乎重新定義了.
什麼人才是神的子民.
這裡說的是.
無論是以色列人或外邦人.
當他們謹守安息日.
揀選神所喜悅的事的時候.
他們都能成為神立約的子民.
可能這群人有些以前是被排擠的.
但今天神都撥他們過來.
收納他們成為自己的子民.
經文第七第八節說到.
神會怎樣對待這些子民呢?.
神說會必領他們去他的聖山.
在神的裡面滿有喜樂.
並且會閱納他們為他.
在壇上獻上的飯祭和平安祭.
神的殿必稱為萬民禱告的殿.
祂不但招聚被趕散的以色列人.
祂更加招聚其他人.
一併成為祂的子民.
愛上帝守安息日的太監.
外邦人可以被神重新接納成為神的子民.
是不是實在很大的應許和祝福呢?.
我總結一下剛才說的話.
今天的經文告訴我們.
上帝的救恩臨近.
神要以公平和公義去招聚一切.
願意持守安息日親近祂的人.
不分猶太人或外邦人.
上帝不是按血統或人的地位去揀選人.
而是按著他怎樣去回應上帝的救恩.
用公義和忠誠去接納那些人成為祂的子民.

$^{441}$就是那些曾經被排擠過的人.
只要他愛耶和華.
持守神的約.
事奉祂.
那些人都可以納入神的家.
因此上帝子民應有的生活方式是怎樣呢?.
就是行公義.
行有過聖潔的生活.
並且謹守安息日,不要作惡.
還要打破排擠和隔閡.
宣告那些與耶和華聯合的人.
都可以在神的殿裡有名分被紀念.
上帝的國度是普世性的.
並且祂持續地要將人帶入祂的救恩計劃裡.
成為屬於祂的群體.
今日的經文告訴我們.
上帝救恩臨近.
神是要以公平和公義.
朝聚一切願意持守安息日.
親近祂的人.
不分猶太人和外邦人.
上帝這樣跟我們說這番話.
究竟給我們有什麼意義呢?.
上帝我們應該怎樣去成為.
我們本身是外邦人.
我們怎樣去成為上帝的子民呢?.
我們有沒有恆公義?.
有沒有過聖潔的生活?.
有沒有謹守安息日?.
有沒有愛別人如同愛自己呢?.
但是說了這麼久.
究竟怎樣才算做守安息日呢?.
是不是單純七天裡.
有一天我們休息什麼都不做.
這樣就叫做守安息日呢?.
其實安息日這個意義.
在聖經裡面是一個很重要的課題.
今日只是稍稍的.
我們大家去重溫一下.
或者去記憶一下.

$^{481}$究竟安息日是一個什麼日子.
不知道大家想起安息日.
第一樣東西選擇.
想起什麼呢?.
一定是創世的時候.
第七日神定為.
停工就安息,將那天定為聖日.
開始了有安息日的概念.
讓我們知道.
其實安息日不單單是休息的日子.
去提醒我們世界是上帝創造的.
這個上帝是會停下來.
去見到這個世界去創造的圓滿.
和這個秩序是有規律的.
世人不應該.
我們現在去看回安息日.
我們不應該覺得.
我們靠人才可以維持這個世界.
要運作得好.
最後的其實是上帝去賜福和保守.
我們這個世界才可以運作.
我們才可以得以生存.
去到後期一點.
到他們出埃及的時間.
出埃及和申命記裡面的記載.
安息日加了一些意義.
就是要紀念他們蒙救贖.
以色列人曾經在埃及做奴隸.
當這群人被釋放之後.
神要他們守安息日.
因為神要他們讓自己.
讓家人.
讓奴僕和寄居的人.
都可以一起休息.
安息日不單單只是宗教上.
一些敬拜和敬拜的操練.
更加是一個解放的記號.
表明上帝是反對奴役.
並且他的子民.
要活在自由和憐憫之中.

$^{521}$到摩西律法的時候.
安息日就成為十誡其中的一條誡命.
也是上帝和子民之間.
文約的一個記號.
祂要求人停止日常的生產.
和自我中心的生活.
好讓他們把時間重新歸給上帝.
重新學習去信靠神.
而不是靠自己.
所以安息日也叫我們知道.
我們不是主人.
神才是我們的主人.
到了先知書的時候.
安息日又有進一步的意義.
它的意義又擴展了.
不單只是一些外在的規條.
我們要去跟從.
而是裡面更加包含著.
要公義,要敬虔,要馴服.
真正的安息日.
不單單是我們放下手上的工作.
就是完成了安息日.
而是包含了我們的心態.
我們怎樣去敬拜.
或者我們和人怎樣去相處.
違反了上帝的公義和秩序.
本質上也是破壞了安息日的精神.
舊約的安息日.
不單只是不工作一天.
它更加是叫我們知道.
我們要紀念上帝的創造.
我們經歷了救贖.
活出文約,追求公義.
並且預告最終在上帝裡面.
我們得到安息.
到了新約.
有什麼不同呢?.
新約的時候.
讓你記得什麼呢?.
關於安息日.

$^{561}$其實在新約的時候.
也延續了安息日.
作為上帝賜福,休息,敬拜.
信靠主的一些意義.
但是有些焦點又改變了.
它就將一條規條.
要守的規條轉向了.
我們要在基督裡面得到真正的安息.
在福音書裡面.
耶穌並沒有否定安息日的本身.
反而宣告.
人子是安息日的主.
記不記得?.
表示安息日仍然是屬於主的.
它有它的屬靈的意義.
新約保留了安息日.
作為敬拜,停止勞動.
轉向上帝這個概念.
不過不再單單說因為要守律法.
而去守安息日.
在福音書裡面.
耶穌多次在安息裡面.
醫病,趕鬼,行善.
祂指出安息日是為人的益處而設.
而不是用來壓制人的.
這裡令到安息日.
從一條規條化的律例典章.
回到它原本的目的.
就是要人得生命.
並且得到恢復和自由.
新約特別將安息日提升為.
對著基督的一些面向.
聖經裡面.
馬太方寫一章28節.
我們很熟悉.
耶穌說.
凡勞苦擔重擔的人.
可以到我這裡來.
我就使你們得安息.
真正的安息.

$^{601}$只可以在基督裡面才找得到.
所以安息日不單單是每星期休息一天.
更加預表了.
藉著耶穌基督.
我們可以進入和神和好的關係.
享受那份上帝賜給我們的安息.
到書信的部分.
在歌羅西書2章16到17節是這樣說.
所以不拘在飲食上或截其月屬安息日.
都不可讓人論斷你們.
這些完事後事的影兒.
那形體卻是基督.
初期的教會對應該守哪一天.
未有很完全的.
大家有不同的說法.
這裡保羅要提醒他們.
重點不是哪一天.
而是你們不要讓人彼此被論斷.
你們要持守愛心和信心.
因為重要的不是哪一天.
而是安息日的精神.
是追求我們在基督裡面得到的那份安息.
所以基督徒應該在我們生活當中.
去實踐這份從耶穌基督給我們的安息.
新約不是取消安息日的核心精神.
而是將它由影子帶到實體.
從外面的規條帶到基督裡面的安息.
究竟今天我們怎樣去遵守安息日才是最合宜呢?.
其實以冊書對安息日的教導.
如果應用在今天生活的時候.
重點並不是我們一些的宗教行為.
而是我們怎樣將我們敬拜的生活.
我們和人中間的關係.
或者我們怎樣在社會裡面的生活都連結在一起.
守安息日就是要基於神的愛.
將神喜悅的公平和公義.
實踐在我們生活的之中.
神呼召我們在日常生活裡面.
活出誠實,憐憫,謙卑.
我們還要關顧弱勢的群體.

$^{641}$因為上帝看重的是公義.
這個公義不單單是個人的品格.
更加包括整個的群體.
我們是不是按著神的心意去生活呢?.
信徒不是單單在教會裡面.
在團契裡面.
在崇拜的裡面去實踐公義.
更加要在家庭裡面.
在職場裡面.
在社會裡面.
去讓人看到公義.
安息日到今天.
仍然都是我們是不是屬於上帝指紋的一個印記.
叫我們知道.
讓人看到我們的生命.
就知道我們是基督徒.
以下書58章13節.
聽我讀.
今日我們怎樣去守安息日呢?.
我們應該都好像以往的人一樣.
有固定的時間.
分別來歸給上帝.
這裡經文告訴我們.
安息日應該做什麼呢?.
調轉你的腳步.
為什麼呢?.
你應該停止你的工作.
向目標為本的節奏.
那個生活的方向.
停下來.
專心的去敬拜神.
去讀經,祈禱,去休息.
去在家庭裡面.
跟家人相聚.
去過憐憫人的生活.
重點的不是守這些外在的規條.
而更加重要的.
是我們要跟上帝的節奏合拍.
可能你的生活很快很快地走.
但是安息日裡面.

$^{681}$上帝就是要你停下來的腳步.
跟祂合拍子.
讓上帝調節你的生活的節奏和方向.
讓你去反省.
你自己在生活上面.
關係上面.
在社會上面.
有沒有行公義.
有沒有做神喜悅的事.
我們有沒有關懷弱勢的社群.
有沒有拒絕社會上剝削的事件.
有沒有在不公義的事的面前.
我們有勇氣的去指出呢?.
以賽斯的作者.
沒有告訴我們.
我們可以靠著我們的德行.
去取決上帝.
而是如果我們行公義.
我們可以把它看成神和我們同行的標記.
當我們行公義的時候.
神就是和我們同行.
當我們在安息日的時候.
我們禱告我們悔改.
我們聆聽上帝的話語的時候.
我們會被上帝的靈去改變.
更新改變我們.
我們之後就將這個改變.
化成具體的行動.
將上帝的愛 真實和具體感.
在我們生活裡面彰顯出來.
剛才說.
以賽亞書第五十六章的對象.
是主前五百年的人.
當時斯里人生活在很多困難的裡面.
他們一方面要重建聖殿.
重建家園 重建群體.
去面對經濟的壓力.
人生的安全.
對未來的充滿困惑.
耶和華明白他們的心.

$^{721}$也知道他們的困難.
因此耶和華呼召他們.
調轉他們的腳步.
呼召他們要守安息日.
要和上帝同行.
今天.
你可能在某些方面.
和主前五百年的人.
都是一樣.
面對很多的困難.
面對很多的挑戰.
你的腳步可能很快很快地可以走.
要達成這些 要達成那些.
今天.
上帝呼召你.
你要調轉你的腳步.
去跟上帝走.
你願不願意呢?.
安息日叫我們知道.
我們是屬於上帝.
守公平 行公義.
叫我們知道.
神是和我們同行的.
在尼加書.
六章八節這樣說.
西人啊.
耶和華以指示你何為善.
他向你所要的是什麼呢?.
只要你行公義.
好憐憫.
全謙卑的心.
與你的神同行.
求主親自引領我們走面前的路.
\newpage



\section{腓立比書 1:27}
\label{sec:032197}
\textbf{腓書回顧 (4) 相稱}~~~~盧偉心弟兄
\newline
\newline
連結: \href{http://www.church.com.hk/acms/content.asp?site=cdc\&op=show\&type=product\&code=032197\&layout=sermon}{\texttt{032197}} ~~~~ 語音日期: 2026-05-23
\newline
\newline
\hyperref[sec:032196]{\small{< < < PREV SERMON < < <}}
~
\hyperref[sec:index_chronic]{\small{[返順時目]}}
~
\hyperref[sec:032198]{\small{> > > NEXT SERMON > > >}}
\newline
\newline
腓立比書 1:27
\newline
\begin{longtable}{cl}
\hline
\hline
章節 & 經文 (和合本修訂版)\\
\hline
1:27 & \begin{tabularx}{0.7\textwidth}{X} 最重要的是：你們行事為人要與基督的福音相稱，這樣，無論我來見你們，或不在你們那裡，都可以聽到你們的景況，知道你們同有一個心志，站立得穩，為福音的信仰齊心努力， \end{tabularx} \\ \\
[1ex]
\hline
\hline
\end{longtable}
$^{1}$各位聽眾朋友,歡迎大家來到十格基督教會.
很感恩我們能夠繼續看肥辣比書.
之前我們用了十四篇道講完整本肥辣比書.
然後去年年底我們開始回顧.
我們的大題目是肥書回顧.
一方面我們是回顧肥辣比書.
另一方面我希望透過肥書回顧.
讓我們每一個都成為一隻肥豬.
就是FAT.
我們Faithful 忠誠委身.
Available 隨時候命.
Teachable 我們樂於受教.
我們有一個肥辣比書的超濃縮版.
我希望透過這卷書當中重要的句子和詞語.
去幫助大家記得整卷書大概是講甚麼.
A部份保羅的文案.
「瀑人,瀑人致函聖徒,願主恩歸你們」.
「瀑人」和「主恩」兩個字眼是重要的.
「瀑人」希利文是Doulos.
是指向整封信的中心基督之歌.
因為這個字眼會在基督之歌再出現.
另外「主恩」是指基督的恩典.
B「那日必成者功」三合一榮耀果.
「那日必成者功」是講一章六戰.
保羅的宣告.
保羅說上帝在信徒的心裡開始一個工作.
到了基督再來的日子.
他必定完成這個工作.
而三合一是指保羅為信徒祈求的三種屬靈能力.
是愛心,真認識和審查力.
這三合一的屬靈能力不斷增長.
並且結滿果子.
就可以讓榮耀歸給父上帝.
C是保羅自己的見證.
活著就是基督.
死活基督顯大.
保羅分享到他的近況.
他說他雖然被軟禁.
但是為著福音能夠更加興旺.
為著弟兄們能夠更加放膽傳道.

$^{41}$他非常高興.
然後保羅分享他自己的內心世界.
他說無論是生是死.
總要讓基督在他身上照常顯大.
這部分是保羅的勉勵.
天文雙稱福音.
無畏暫穩同敬.
合一謙卑他顧.
以基督心為心.
天文雙稱福音就是我們剛才讀過的一章二十七節.
保羅說這個是最重要的.
你們行事為人.
要與基督的福音去相稱.
接著保羅勸勉信徒.
要無畏暫穩同心競鬥.
所謂競鬥是為福音齊心努力的縮寫.
因為希臘文是Synophilia.
意思是一班運動員參與對際比賽.
為著勝利齊心努力的意思.
合一謙卑他顧.
是二章一至四節.
大家可能都會記得.
各人看別人比自己重要.
各人不要單顧自己的事.
也要顧別人的事等等.
而以基督心為心就是二章五節.
我們要思想耶穌心中所想的.
二部份就是肥勒比書的核心.
基督之歌.
八幹信服至死.
受苦進入榮耀.
所謂八幹是指耶穌基督.
從本與上帝同等.
一路虛己.
一路下降.
直到去全心信服.
以至於死且死在十字架上.
但基督受苦是為要進入榮耀.
他復活了.
並且得次超乎萬明之名.

$^{81}$萬失都要跪拜.
萬口都要宣認.
耶穌基督是主.
歸榮耀給付上帝.
去到D'保羅再一次面禮信徒.
戰經作不怨言.
顯主道明光照.
因為這部份我們稍後會講.
所以我們先跳過.
然後在X部份.
保羅舉了兩個例子.
提摩泰求主事.
以巴賭命近死.
第一 他提到提摩泰.
如何不求自己的事.
去真心關懷肥勒比信徒.
第二 保羅提到以巴弗提.
他如何為基督的工作.
就是他如何支援保羅宣教.
賭上自己的性命近乎於死.
去到C'.
保羅再一次講自己的見證.
氣萬事 贏基督.
應主召 竭力追.
他講到他以往是如何靠肉體.
到現在因為基督.
他明白地上的萬事都是糞土.
所以他定義丟棄萬事.
單單要贏得基督.
不過保羅說他自己並不完全.
他其實只是應主召 竭力追.
應主召的意思是指保羅.
他如何回應基督叫他悔改的呼召.
保羅提到他以前是迫害教會.
他故意用到迫害.
希臘文是Dioko這個字.
現在他說他悔改了.
所以他竭力追求基督.
同樣是Dioko同一個字.
去到D'2.

$^{121}$保羅再一次提醒信徒.
效榜樣防受敵.
天文 後主 榮耀.
保羅勉勵肥勒比信徒要效法他.
也要留意那些跟他們一起.
效法保羅的人.
因為那些人是口不對心的十價仇敵.
那些人的特徵是專以地上的事為念.
相反 如果我們是真正天上國民.
我們會以天上的事為念.
專心等候救主從天降臨.
等候擁有和基督一樣的榮耀身體.
主近了同心證 無慮祈求永恆.
是四章的內容.
在主已經近了這個前提之下.
保羅勸勉信徒要靠主站立得穩.
他特別提醒到尤雅蝶和秦都基.
要在主裡面同心.
要回想他們以前怎樣跟保羅一起.
Sunafilio 一起為福音齊心努力.
也一無掛慮地為了永恆的事情.
去到B'.
知足不求鬼贈 果多榮耀歸附.
但保羅不想他們誤會.
或者錯估了他們奉獻的作用.
所以他去強調他其實自己是知足的.
也表明他不求鬼贈.
但另一方面 菲勒比信徒的奉獻.
其實是他們的天國存款.
是歸在他們帳上的果子.
這些果子能夠讓榮耀歸附上帝.
最後 保羅問候每位聖徒主因與你同在.
他用「願主耶穌基督的恩惠與你們的靈同在」.
作最後的問安和祝福.
呼應著一章二節的文案.
「願恩惠平安 從我們的父上帝.
和主耶穌基督歸給你們」.
菲勒比書最重要的部分就是中間.
「基督之歌」.
所以我們這個回顧系列.

$^{161}$我們會按這個結構由外至內進行.
之前我們已經看了ABC C' B'和A'.
兩個月前開始.
我們進入了D D'和D'2的部分.
原本上次我說打算用兩篇道去講這個部分.
但後來在預備這篇道的時間.
我發覺我需要三篇道才足夠.
所以如果大家還記得上次那篇道叫做天文.
今次這篇道叫做商程.
而下一篇道我們會講爭戰.
所以今天這篇道我決定不要太貪心.
我只是分兩個部分.
第一「與主福音相稱 誕生登基凱旋」.
第二「薦經作不怨言 顯主道明光照」.
第一個部分「與主福音相稱 誕生登基凱旋」.
一章二十七字「最重要的是」或者是「只要」.
「你們行事為人 要與基督的福音相稱」.
保羅說這一章二十七字的吩咐.
是全卷書最重要的.
甚至這個最重要可以翻譯為「只要」.
他說你們只要行這件事.
就是你們行事為人 與基督的福音相稱.
上一篇道我們特別提到行事為人這個字.
原文是Politurme.
意思是做公民的意思.
所以上一篇我們是集中思想.
天國公民這個身份的重量.
我們說到在《出埃及記》十九章的時間.
上帝已經在西奈山.
他已經說要從萬民當中分別以色列人出來.
作屬上帝的子民.
上帝說他們要成為祭司的國度 聖潔的國民.
而這個在舊約西奈山的應許.
是在伽勒山上面實現的.
是耶穌基督用自己的寶血埋屬一群罪人.
將他們從罪裡面釋放出來.
在啟示錄裡說.
耶穌基督的血埋屬了他們之後.
使他們成為國度 作父上帝的祭司.
這就是天國公民這個身份的重量.

$^{201}$今天我們會集中思想另外兩個詞語.
一個是福音 另外是相稱.
大家可能記得這幅圖畫.
相稱這個字在原文是Axios.
意思是一樣的重量 或者一樣的價值.
所以當時看到Axios的時候.
他們會想起天平.
上次的圖畫是想像一下.
天平一邊放的是基督的福音.
另外一邊是我們的行事為人.
然後保羅說天平應該是平衡的.
信徒的生命應該是能夠與基督的福音同等重量.
同樣的價值.
我們必須要問一個問題.
究竟福音有多重.
與基督福音的相稱.
其實是不是一個不可能的任務.
我們的生命究竟怎樣才可以與基督福音同等重量.
或者我們一步一步去想這個問題.
我首先問大家究竟什麼是福音.
如果你要跟一個未信的朋友介紹福音.
你會怎樣說呢.
然後我們再思想究竟福音有多重.
如果你在我們當中聽林牧師講到已經有一段時間.
你可能會記得福音這個詞是一個專有名詞.
字面的意思是好消息.
不過不是一個普通的好消息.
在古希臘羅馬文化裡面.
福音這個字是指一個君王.
一定是關於君王的.
關於君王的什麼呢.
一 誕生.
二 就位登基.
三 戰勝敵人.
海船如歸的公開宣告.
所以基督的福音是講耶穌基督這位天國的君王.
祂的誕生.
祂到城獄身.
降世為人的好消息.
祂的登基.

$^{241}$祂在十字架上登基作王的好消息.
祂的海船.
祂怎樣戰勝罪.
戰勝死亡.
海船而歸的好消息.
誕生 登基 海船.
其實也是菲勒比書的核心.
基督之歌的脈絡.
這裡說.
他本有上帝的形象卻不堅持自己與上帝同等.
反倒虛己.
取了奴僕的形象.
成為人的樣式.
說的是耶穌基督的誕生.
既有人的樣子.
就謙卑自己.
全心信服.
以至於死.
且死在十字架上.
這是耶穌基督的登基.
所以上帝把他升為至高.
由此給他超乎萬名之上的名.
是一切天上的.
地上的.
和地底下的.
因為耶穌的名.
萬室都要跪下.
眾口都要宣認.
耶穌基督是主.
歸榮耀.
給付上帝.
這就是耶穌基督的海船而歸.
究竟福音有多重.
今天我想和大家離開菲勒比書.
我們會看很多不同的經文.
我們會和大家思想.
福音究竟有多重.
我們會思想主耶穌基督.
這位天國君王的誕生.
登基和海船.

$^{281}$我們先看誕生.
希伯來書十章.
因為公牛和山羊的血不能除罪.
所以基督到世上來的時候就說.
「積物和禮物不是你所要的.
但你曾給我預備了身體」.
基督到世上來的時候.
原來這裡是在說耶穌基督的誕生.
這裡說到耶穌基督為何要誕生.
是因為世人都犯了罪.
是因為罪的公教乃是死.
因為死後有審判.
但是公牛和山羊的血不能除罪.
希伯來書這句話.
在我們聽來可能沒有什麼感覺.
但對於一個猶太人來說.
這是一個非常非常震撼的宣告.
為什麼呢.
因為這裡的公牛和山羊的血.
是代表著憲制.
而憲制是上帝親自設立的.
由《出埃及記》的回望.
到《利未記》的種種憲制.
祭司要怎樣宰殺生畜.
怎樣灑血.
怎樣焚燒.
都有很仔細的要求.
這些所有一切都是上帝吩咐的.
特別是在每年的贖罪日.
大祭司都要帶著贖罪祭的血.
進入至聖所為百姓贖罪.
但震撼的是這裡說.
原來這一切都是沒有用的.
這一切的憲制都不能除罪.
因為這些祭全部都是影子.
預表著耶穌基督這個實體.
所以這裡正正是因為.
公牛和山羊的血不能除罪.
所以基督才要來到世上.
基督到世上來的時候.

$^{321}$說了一句話.
祭物和禮物不是你所要的.
但你曾給我預備身體.
其實這句話出自詩篇四十篇.
祭物和禮物你不起愛.
犯祭和贖罪祭非你所要.
那時我說.
看啊 我來了.
我的事在經卷上已經記載了.
我的上帝啊.
我樂意照你的旨意行.
雖然詩篇四十篇是大衛所寫的詩.
但當中的「我」其實是指向耶穌基督.
所以當我們綜合詩篇四十篇.
和希伯來書十章這兩段經文的時候.
似乎意思是耶穌基督說.
父啊.
你所要的不是牛或者羊的生命.
你所要的是一個信服的生命.
你所要的是一個人完全地.
徹底地.
毫無保留地去相信你.
愛慕你.
樂意照你的旨意行.
而這份信服.
必須透過一個身體去活出來.
一個會餓.
會渴.
會累.
會痛.
會流血.
會死的身體.
而這就是「道成為肉身」的意義.
上帝為耶穌基督預備一個身體.
是要讓祂能夠按經卷上的記載.
並且樂意照上帝的旨意而行.
行什麼呢?.
我們很快會有答案.
不過我們希望大家先記住.
「我」.

$^{361}$這裡耶穌基督的「是」.
在經卷上已經記載.
稍後我們會看到另外一個例子.
現在我們先回到希伯來書.
剛才我們看了十章四節.
五節那裡.
「但你曾給我預備了身體」.
去到第十節.
「我們憑著這旨意.
藉著耶穌基督僅此一次獻上祂的身體.
就得以成聖」.
「我們憑著這旨意」.
「旨意」是指上帝的旨意.
上帝的旨意是什麼呢?.
是耶穌基督獻上的祂的身體.
以致我們.
即是我和你.
得以成聖.
這是上帝的旨意.
耶穌基督降生在世上.
道成為肉身.
是為了有一天.
使得你和我這些罪人.
可以得以成聖.
而獻上你的身體.
這是上帝的旨意.
我在馬槽裡面的小小嬰孩.
他小小的手.
小小的腳.
預備好了.
是為了有一天.
為了我們這些罪人.
被釘在木頭上.
這就是上帝的旨意.
這個小小的嬰孩.
他小小的頭.
小小的身體.
預備好了.
是為了有一天.
為我們這些罪人.

$^{401}$戴上荊棘的冠冕.
他的肋旁被茅去刺透.
流出血和水.
來洗淨和重生我們這些罪人.
這就是上帝的旨意.
所以親愛的弟兄姊妹.
究竟福音有多重?.
這就是福音的重量.
而這只是說到誕生.
還只是誕生.
如果誕生已經讓我們感到震撼的時候.
登基就更加會顛覆我們對作王的一切想像.
耶穌基督的登基.
不像地上的君王.
他不在皇宮裡登基.
不在寶座上登基.
不在萬民歡呼中登基.
他的登基發生在十字架上.
我們看詩篇22篇.
是一首大衛的詩篇.
但聖靈感動大衛寫下了一些超越他時代的說話.
一些唯獨在耶穌基督身上才應驗的說話.
就像剛才詩篇40篇所說的.
耶穌基督的事在經卷上已經記載.
這就是其中一個例子.
我的上帝我的上帝為什麼離棄我.
但我被眾人羞辱被百姓藐視.
凡看見我的都癡笑我.
他把自己交託給耶和華.
讓耶和華救他吧.
我如水被倒出來.
我的骨頭都脫了賤.
他們摘了我的手我的腳.
他們分我的外衣為我的內衣抽籤.
弟兄姊妹當我們讀到這些說話.
你會發現這裡不是在說大衛的經歷.
大衛沒有被人摘手摘腳.
大衛的衣服也沒有被人抽籤去分.
大衛這裡所寫的是十字架的寓言.
耶穌在十字架上也親口引用了這篇詩篇的第一節.

$^{441}$我的上帝我的上帝為什麼離棄我.
而在那一刻.
他承受的不單單是肉體上的痛苦.
他承受的是更大的痛苦.
是與天父上帝分離的痛苦.
那位從耿古到永恆.
與聖父緊密相連的聖子.
在那一刻.
就是那個那無罪的.
替我們成為罪的那一刻.
被父神轉念不顧.
地上的君王登基時候.
穿的是華麗的王袍.
戴的是黃金的冠冕.
接受的是萬民的歡呼.
但是我們的君王耶穌基督.
他登基的時候沒有華麗王袍.
他是赤身露體.
他接受的是懸該我們這些罪人.
應該承受的羞辱.
他沒有黃金的冠冕.
他頭上戴的是荊棘編成的冠冕.
是羅馬冰冰的嘲笑和戲弄.
他沒有得到萬民的歡呼.
他得到的是眾人的譏諷.
詩篇22篇說.
「凡看見我的都痴笑我」.
就應驗在路加福音的記載.
無論是官長,士兵.
還是銅釘的犯人.
都譏諷他.
他們說你若是猶太人的王.
救救你自己吧.
不過我們要知道.
大衛所寫的十字架語言.
其實還有下文.
敬畏耶和華的人.
要讚美他,榮耀他,懼怕他.
因為他(耶和華).
沒有藐視,憎惡受苦的人.

$^{481}$也沒有轉念不古.
地的四極都要歸順他.
列國的萬族都要在你面前敬拜.
因為國度屬於耶和華.
他是管理列國的.
主所做的事必傳給後代.
他們(敬畏耶和華的人).
必來傳他的公義給尚未出生的子民.
所以在十字架上.
雖然主耶穌基督被眾人羞辱.
被百姓藐視.
但耶和華沒有藐視.
沒有憎惡這位受苦的僕人.
在十字架上主耶穌基督呼喊.
我的上帝!我的上帝!.
為什麼離棄我?.
但最終耶和華沒有永遠的.
轉念不顧這位受苦的義僕.
在十字架上當主耶穌基督登基作王.
他用自己的血.
從地的四極.
列國的萬族當中.
買熟了一群子民.
歸屬於耶和華的國度.
降服在這位十字架君王的皇權之下.
這裡也呼應著.
菲勒比書的Polytheum.
而這群被買熟的子民.
要成為天國的公民.
他們要做些什麼呢?.
要把主所做的事傳給後代.
要把上帝在十字架上所彰顯的公義.
傳給尚未出生的子民.
所以請問我們弟子們.
究竟福音有多重?.
這是福音的重量.
這是登基的重量.
但我們還未完成福音的全部內容.
因為這個登基的君王.
他還要戰勝敵人.

$^{521}$海船而歸.
耶穌基督的海船包含了.
復活,升天,再來這三個層次.
而在復活,升天,再來這三個層次的海船當中.
重要的是他不是獨自得勝.
之前他是獨自降生.
他是獨自登基.
但是在海船當中.
他要帶著他的子民一同進入海船.
海船是一個集體的維度.
海船是和我們有關.
首先海船是開始於復活.
耶穌基督從死中復活.
證明他自己勝過了罪和死亡的權勢.
死亡不能捆綁他.
憤怒不能囚禁他.
撒旦以為當耶穌基督寫在十字架上的時候.
他贏了.
但復活是證明撒旦沒有得勝.
一切仍然在上帝的計劃和掌握當中.
更重要的是耶穌基督的復活.
不是單單關於他自己.
他的復活是和我們有關的.
我們看羅馬書第六章.
所以我們藉著洗禮歸入他的死.
和他一同埋葬.
是要我們行事為人都有新生的樣子.
所以我們經歷與他一樣的死.
也將經歷與他一樣的復活.
我們的救人和他同釘十字架.
不再作罪的奴隸.
我們若與基督同死.
我們信也必與他同活.
所以不要讓罪在你們必死的身上掌權.
使你們順從身體的私慾.
也不要把你們的肢體獻給罪作不義的工具.
倒要讓從死人中活著的人.
把自己獻給上帝.
並把你們的肢體獻給上帝作義的工具.
剛剛在希伯來書十章.

$^{561}$我們看到上帝的旨意.
是耶穌基督一同獻上他的身體.
以致我們得而成聖.
而在羅馬書六章.
保羅有一個相應的吩咐.
我們要與基督同死同活.
我們同樣也要獻上身體.
保羅說我們要將肢體獻給上帝作義的工具.
但保羅為什麼可以這樣吩咐呢?.
因為有一個事實已經改變了.
保羅說當我們與基督同死同復活的時候.
我們的舊人已經釘死了.
罪的權勢被打破了.
我們已經不再是罪的奴隸.
我們已經被釋放.
正正是因為這個事實改變了.
所以保羅說你們現在行事為人.
要有新生的樣式.
你們不可以再像以前一樣.
將身體獻給罪.
你們已經復活得新生.
現在你們要將身體獻給上帝作義的工具.
什麼是作義的工具呢?.
是完全被上帝所用.
做上帝所喜悅的事.
然後去到海船的第二個層次.
升天.
升天也與我們有關.
因為使徒行傳告訴我們.
復活後的耶穌被升天.
然後賜下聖靈.
在使徒行傳二章彼得說.
祂既被高舉在上帝的右邊.
又重復受了所應許的聖靈.
就把你們所聽見所看見的囂冠下來.
耶穌基督升天之後.
賜下聖靈.
使每一個屬祂的人能夠重生.
聖靈更新了我們的心思意念.
使我們有能力活出.

$^{601}$與福音相稱的生命.
而且耶穌基督升天之後.
又時刻為我們去代求.
在羅馬書八章保羅說.
誰能定他們的罪呢?.
有基督耶穌已經死了.
而且復活了.
現今在上帝的右邊.
也替我們祈求.
祈求什麼呢?.
希伯來書七至九章告訴我們.
耶穌基督作為更美滋藥的忠報.
祂是永遠流傳不可更換的大祭司.
祂為我們除去致死的行為.
好讓我們事奉那位永生的上帝而去祈禱.
所以在希伯來書七章二十五節說.
所以凡靠著他進到上帝面前的人.
他都能拯救到底.
因為他長遠活著為他們祈求.
最後海船將會在基督再來的日子中完成.
耶穌基督這位海船的君王.
不是獨自一個人海船.
他要買熟一群子民.
從各支派,各語言,各民族,各邦國當中.
將他們買來歸給上帝.
這群人要在海船當中與基督有份.
他們要成為國民和祭司歸屬於上帝.
在第五章說到.
他們唱新歌說:你佩拿書卷佩揭開它的人.
因為你曾被殺.
用自己的血從各支派,各語言,各民族,各邦國中.
買了人來歸給上帝.
又使他們成為國民和祭司歸於我們的上帝.
他們將在地上執掌王權.
大聲說:曾被殺的羔羊.
配得權能,豐富,智慧,力量,尊貴,榮耀,頌讚.
大家看得見嗎?.
這位海船的君王.
他不是孤獨地站在榮耀當中.
他身邊其實是有無數.

$^{641}$從萬國萬民當中買熟回來的子民.
一同進入最終的海船.
這些人成為國民和祭司歸給上帝.
在基督再來的時間.
他們要與基督一同地上執掌王權.
去到那時海船才真正完成.
因為不是正正君王耶穌得勝.
而是整個上帝的國.
從創世以來到末世.
所有的選民都與耶穌基督這個君王一同得勝.
因為剛才我們講的內容非常豐富.
我想我們要重溫總結一下.
如果我們要回答什麼是福音.
究竟福音有多重.
我們要記得福音是講君王的誕生.
希伯留書十章告訴我們.
因為世人的罪無法用牛羊的血來解決.
所以耶穌基督來世上的時間.
上帝為他預備身體.
上帝的旨意是耶穌基督獻上他的身體.
以致我們成為我們的贖罪者.
第二 君王的登基.
詩篇22篇告訴我們.
這位君王登基不在皇宮裡.
在十字架上.
他被羞辱 被氣絕 手腳被砸.
承受了與上帝分離的痛苦.
他在那裡登基作王.
用自己的血買贖一群子民.
歸給上帝成為天國的公民.
最後是君王的凱旋.
這位君王從死裡復活.
勝過罪的權勢.
然後他升上高天 賜下聖靈.
使我們有能力活出與福音相稱的生命.
最後他要再來.
去到那一天會率領所有的子民.
從各族 各國 各民當中買贖回來的人.
一同進入最終的凱旋.
而這就是福音.

$^{681}$這是福音的重量.
我們生命如何與福音相稱呢?.
中國人有句說話.
「死有輕於鴻毛 有重於泰山」.
如果一個普通人的死都可以重於泰山.
那麼耶穌基督的福音.
只會比喜瑪拉雅山更加重.
因為當中包含了上帝的兒子.
為我們而死.
還有他之後的整套計劃.
實在不能夠去形容.
所以當古來吩咐我們要行事為人.
而與基督福音相稱的時候.
我腦中最初浮現的是這幅圖畫.
天平的一邊.
好像喜瑪拉雅山那麼重的福音.
另一邊我們一群信徒.
很竭力地跳來跳去.
大家一二三跳.
希望用我們的重量.
讓天平可以平衡一點.
但發現天平民風不動.
一點也不動.
親愛的兄弟姐妹.
你是否還在這些跳上跳下的信徒當中.
還是你已經放棄了.
心裡想根本是mission impossible.
不如躺平吧.
如果你有這個想法.
我可以告訴你.
躺平也是不符合聖經的.
對的 如果我們靠自己.
無論我們如何努力跳上跳下.
我們的生命都永遠無法.
與那位為我們誕生.
在十字架上登基.
又海船復活的君王相稱.
我們可以同等重量是不可能的.
但我們也不可以躺平.
因為當保羅在一章二十七節.

$^{721}$跟我們說這門行事為人與福音相稱的時候.
他連接著第二章.
我們今天首讀的經文.
他馬上告訴我們.
如何可以做到.
就是第二部分.
「戰經作不怨言,顯主道明光照」.
其實這裡是三個應用一個的後果.
「戰經作不怨言,顯主道」.
是三個的應用.
而「明光照」是後果.
讓我們先看戰經作.
菲勒比書二章十二節.
「所以我親愛的,你們向來是信服的.
就當恐懼戰經完成你們自己得救的事.
因為是上帝在你們心裡運行.
使你們又立志又實行.
為要成就他的微意」.
在十二節和合本譯作.
「這樣看來」這個字.
其實更加直接翻譯為「所以」.
有「所以」就有「因為」.
「因為」是什麼呢?.
大家記得十二節之前是什麼嗎?.
是基督之歌.
六至十一節基督之歌.
剛才說過了.
基督之歌其實是說福音.
誕生,登基,海船.
因為耶穌基督的誕生,登基,海船.
所以我們應當恐懼戰經.
完成自己得救的事.
這裡完成自己得救的事.
不是我們要通過努力得著自己的救恩.
我們知道保祿很清楚教導的.
我們得救是本福因.
因著信.
不是出於自己.
也不是出於行為.
人在得救的事上沒有任何功勞.

$^{761}$人不可以藉著任何方式去賺取救恩.
什麼時候人以為自己的善事.
自己的功德.
去維持自誇的時間.
或者以為他們靠守律法.
可以自稱為儀的時間.
他就與基督的義隔絕.
所以完成自己得救的事的意思.
是持續不斷努力去實現.
並活出自己的救恩.
所以雖然我們不能夠通過行為去賺取救恩.
但是另一方面.
一個已經重生得救的人.
他必然會有重生得救之後的行為表現.
所以保羅的意思是.
我們的生活應該能夠彰顯出.
上帝在我們身上的拯救大能.
我們應該不斷努力地活出.
已經重生的新生命.
就是單單只求上帝的國和上帝的義.
但是為什麼我們努力地活出新生命的時間.
要恐懼戰經呢?.
第十三節給了我們答案.
因為是上帝在你們心裡運行.
使我們又立志又實行.
為要成就祂的美意.
這裡運行和實行原文其實是同一個字.
原來我們之所以努力地活出重生得救的生命.
不是因為我們突然睡醒了.
突然想通了.
我們真的做了基督徒.
不要這麼壞了,做好一點吧.
不是的.
我們的立志是因為上帝在我們心裡面動工.
我們的實行是因為上帝在我們裡面運行.
而這個就是我們恐懼戰經的原因.
因為是創造天地的主宰.
海船的君王.
永生的上帝.
祂竟然住在我們裡面.

$^{801}$是全能的上帝在我們裡面.
持續不斷地努力實現我們的救恩.
是全能的上帝在我們裡面.
以致我們能夠行事為人.
與基督的福音相承.
所以恐懼戰經其實是包含了.
心感不配.
受寵若驚.
認真謹慎.
全力以赴.
全能的上帝怎樣在我們裡面呢?.
其實剛剛第一部分已經講過了.
當耶穌基督升天之後.
上帝將聖靈敲冠下來.
住在我們裡面.
而聖靈在我們裡面.
又立志又實行.
所以剛剛那幅天秤的圖畫應該是這樣的.
天秤一邊是重如喜瑪拉雅山的福音.
另一邊是一群聖靈內住的信徒.
他們穿上全副的軍裝.
為要進行征戰.
這個天秤之所以平衡.
一邊是基督的福音.
另一邊是聖靈的工作.
是祂讓信徒又立志又實行.
為要成就上帝的美意.
上帝的美意是甚麼呢?.
就是祂的子民,祂的軍隊.
要在這個世代當中好像明光照耀.
不過我們跳得太快了.
我們先看不怨言.
如果你記得「不發怨言,不要起爭論」這句話.
不是在說字面上的意思.
保羅不是吩咐你們不要經常發牢騷.
不要埋怨的意思.
保羅的意思是.
不要像他們.
誰是他們呢?.
不要像出埃及時期的以色列人那樣.

$^{841}$這樣即是怎樣呢?.
不要像他們一樣貪戀世界.
不要像他們一樣自以為宜.
不要像他們一樣興旱救贖.
為甚麼我們知道發怨言是說那些以色列人呢?.
因為其實在2章15節裡.
保羅是在引用申命記32章的摩西治國.
申命記32章.
「這乖僻彎曲的世代.
向他行了敗壞的事.
因著他們的弊病或者亦作瑕疵.
他們不再是他的兒女」.
去到20節.
「我要轉念離開他們.
看他們的結局如何.
他們是極乖謀的世代.
是不忠實的兒女」.
如果你記得申命記32章摩西治國的背景.
其實是摩西臨死之前.
第二代以色列人將要進入迦南.
耶和華就吩咐摩西去寫這首歌.
教他們去唱.
目的是甚麼呢?.
目的是將來以色列人進入迦南地之後.
當他們偏向別臣.
當他們離棄耶和華的時候.
他們會說「使這首歌成為耶和華指責以色列人的見證」.
所以摩西治國其實是上帝對以色列人的判詞.
verdict.
現在大家也看到.
其實同顏色的字就是保羅的引用.
綠色的字.
以色列人是不忠實的.
保羅勸我們要誠實無畏.
藍色的字.
以色列人是乖僻彎曲的世代.
保羅說我們要身處在彎曲的世代當中.
紫色的字.
以色列人因為他們的弊病.
因為他們的瑕疵.

$^{881}$他們不再是上帝的兒女.
保羅就勸我們要作上帝無瑕疵的兒女.
所以在這個引用和對比之下.
我們很有信心相信.
法院言是在說川普時期的以色列人.
所以如果你用聖經軟件去搜尋法院言這三個字.
你就會發現七個以色列人法院言的事件.
第一件事.
川普記十五章馬拉的苦水.
第二件.
川普記十六章.
他們沒有食物.
他們上帝賜下馬拉.
第三件事.
川普記十七章.
他們沒有水喝.
上帝吩咐摩西擊打盤石流出水.
第四件事.
在文蘇記十一章他們覺得馬拉食厭了.
他們求肉吃.
上帝賜下很多鵪鶉給他們.
第五件事.
十探子.
其實十二個探子.
不過有十個回報壞消息.
並且說他們要回埃及.
第六件事.
文蘇記十六章.
有一群可拉.
有一黨他們二百五十幾人.
他們跟摩西說.
你為什麼要高舉自己.
我們所有人都是上帝的聖民.
最後.
文蘇記二十一章.
大集會舉同社.
所謂的大集會.
是說他們將之前的怨言綜合起來.
炒在一起.
然後發出一個終極嫌棄的怨言.

$^{921}$我讀一讀.
文蘇記二十一章五節.
百姓向上帝和摩西法院言說.
你們為什麼把我從埃及領上來.
使我們死在曠野呢.
這裡沒有糧食.
沒有水.
我們也厭惡這淡而無味的食物.
所以是大集會.
炒在一起.
而在這七個法院事件當中.
我有三個觀察.
第一個觀察.
這七宗事件.
其實橫跨了以色列人在曠野的三十八年.
頭三宗.
發生在剛剛過了紅海.
出埃及的初期.
中間的閹嗎哪球肉吃.
十貪子回埃及.
其實是發生在出埃及的第二年.
大家可能記得.
他們在西奈山逗留了一年.
而這兩件事.
就發生離開西奈山之後不久.
而在十貪子回埃及之後.
以色列人就開始在曠野漂流.
然後去到可拉黨.
台自基這個事件.
聖經沒有明確地說明時間.
但是很多學者都估計.
它是發生在漂流曠野的初期或中期.
去到最後的大集會.
舉銅蛇.
發生在漂流曠野的末期.
那些已經是第二代的以色列人.
在進入迦南之前.
所以換而言之.
其實法院言言不是一時間的本應.
而是從頭到尾.

$^{961}$持續不斷的本應.
第二個觀察.
大家都可能知道.
七這個數字.
在聖經裡面是代表完整.
充足的意思.
所以七宗的法院言言事件.
是表示他們是毫無保留.
全力以赴.
完全充足地去本應.
因為如果你心水清的話.
可能你會記得.
其實以色列人法院言就有其他的情況.
例如離開了埃及.
過完紅海之後.
有士兵追他們.
記不記得他們怎麼說的.
我讀一讀給你聽.
他們對摩西說.
難道埃及沒有墳地.
你要把我們帶來死在曠野嗎.
你為什麼這樣待我們.
將我們從埃及領出來呢.
很明顯是法院言.
但是摩西在這件事上.
沒有寫下法院言這幾個字.
我不肯定摩西是否故意.
儲起這七件事.
然後就只在這七件事上.
寫以色列人法院言.
來表達其實他們是毫無保留.
全力以赴這個意思.
但是我們會看到.
以色列人他對上帝的本應.
是由頭到尾持續不斷.
並且是毫無保留.
全力以赴完全這樣的本應.
第三個觀察.
上帝是怎樣回應這樣的本應.
一方面他會審判.

$^{1001}$因為我們的上帝是吞滅的火.
他不以有罪為無罪.
所以我們會看到.
譬如在第六事件的時候.
是他會使地裂開.
吞滅荷拉一黨.
在第七個事件.
他會使火蛇進入百姓當中.
咬死許多人.
但另一方面.
他是有憐憫,有恩惠的上帝.
他有豐盛的慈愛和信仇.
他竟然負上重家.
為罪人迴轉去開路.
我相信大家都知道.
沒食物次嗎哪當中的嗎哪.
是指向耶穌基督的十字架.
如果你上個星期在這裡.
聽林牧師說.
榮耀君王活水靈風.
當中去解釋.
出埃及的十七章 沒水鶴擊盤石事件的時候.
當中的盤石.
也是指向耶穌基督被擊打.
和刺下聖靈.
更加重要的是.
在這七宗發怨言事件的結尾.
在以色列人從頭到尾.
持續不斷毫無保留.
全力以赴的叛逆的最後.
上帝的回應是什麼呢?.
上帝的回應是.
高舉銅蛇.
大家知道這是耶穌基督的十字架.
是耶穌親自對尼哥德姆說的.
摩西在曠野怎樣舉蛇.
人子也必須照樣被舉起來.
要使一切信他的人.
都得永生.
親愛的朋友 親愛的弟兄姊妹.

$^{1041}$我想告訴你們.
其實我們和這些以色列人沒什麼分別.
其實我們每一個都曾經.
甚至或者現在.
仍然是從頭到尾.
持續不斷.
毫無保留 全力以赴.
去叛逆上帝的人.
我們每一個都同樣需要.
依靠耶穌基督的十字架.
來獲得赦罪和救贖.
所以當保羅他吩咐我們.
吩咐我們不要仗貪門的時候.
同樣適用於我們.
我們不可以像以色列人一樣貪亂世界.
我們不可以像他們一樣自以為宜.
我們不可以像他們一樣興旱救贖.
來到這裡我想和大家讀讀.
約翰一書三節的經文.
無獨有偶都是二章十五節前後.
所以如果我們以後看到.
菲勒比書二章十四節不要發怨言.
不要仗貪門的時間.
我們可以聯想約翰一書二章十五節.
讓我們一起讀.
不要愛世界和世界上的東西.
若有人愛世界.
愛父的心就不在他裡面了.
因為凡世界上的東西.
好比肉體的情慾.
眼目的情慾和今生的驕傲.
都不是從父來的.
而是從世界來的.
這世界和世上的情慾都要消逝.
唯獨那遵行上帝旨意的人永遠長存.
最後讓我們回到菲勒比書.
顯主道和明光照.
二章十五節.
作上帝無瑕疵的兒女.
你們在這世代中要仗明光照耀.

$^{1081}$將生命的道顯明出來.
兄弟姐妹,這就是第三個應用.
顯主道.
保羅說我們要仗明光照耀.
並且將生命的道顯明出來.
保羅用顯明這個字眼.
本身是包含了兩個意思.
是雙重的意義.
第一個意思是堅守.
hold fast.
我們要緊緊捉祝福音這個生命的道.
我們不可以被這個彎曲背後的世界去同化.
這個世界不斷告訴我們.
錢是最重要的.
成功是最重要的.
自己的感受是最重要的.
但我們應該要堅守生命的道.
我們要單單只求上帝的國和上帝的義.
第二個意義是高舉.
hold forth.
我們不單止自己收起來去堅守.
我們更加要向這個世代顯明出來.
我們應該要用我們的生命.
用我們人生的抉擇.
向這個世界去宣告.
宣告什麼?.
宣告聖經是真的.
宣告罪是真的.
宣告審判是真的.
地獄是存在的.
最重要的是.
福音是真確的.
耶穌真是答案.
迴轉必定得救.
我們必須去單單只求上帝的國和上帝的義.
因為一個人不能事奉兩個主.
而當我們這樣做的時候.
保羅說後果是什麼?.
明光照耀.
不是我們自己發光.

$^{1121}$是我們反射那個真光的光輝.
就好像尼哥德姆這個從前的黑夜人.
領受聖靈之後他要見光.
我們也要同樣的明光照耀.
不過大家要留意.
光來到世間.
世人因自己的行為善惡的.
不愛光倒愛黑暗.
令他們的罪就是在此.
凡作惡的便恨光.
並不來就光.
恐怕他的行為受責備.
但行真理的必來就光.
要顯明他所行的是靠上帝而行.
保羅預先警告我們.
加上約翰.
明光照耀的後果.
不是我們會讓這個世界.
歡迎,欣賞,做得好.
大家記得上星期林牧師提到.
尼哥德姆的下場嗎?.
在教會的傳統當中.
他是被趕出猶太議會.
甚至要離開耶路撒冷.
最後他因為放膽傳揚福音而殉道.
但就算是死.
保羅說你們要將明光照耀.
要將生命的道顯明.
堅守,高舉出來.
所以下一篇的題目是僧戰.
第一節我們回到最初這幅天秤的圖畫.
我們小小回顧一下.
天秤的一邊是寵魚.
喜馬拉雅山的福音.
君王的誕生.
君王十字架上的登基.
君王的復活.
升天再來.
天秤的另一邊是我們的生命.
當我們靠自己.

$^{1161}$無論我們怎麼跳來跳去.
都無法移動天秤.
但保羅告訴我們一個好消息.
因為不是我們在跳.
是上帝在我們裡面運行.
因為耶穌基督升天.
祂賜下聖靈.
上帝住在我們裡面.
使我們又立志又實行.
使我們能夠恐懼戰驚地活出救恩.
使我們不發怨言.
不會像當年以色列一樣貪戀世界.
自以為異.
輕看救贖.
上帝是聖靈.
使我們能夠堅守高舉生命的道.
讓我們將明光照耀.
這就是與福音相稱的奧秘.
不是靠自己努力跳來跳去.
靠聖靈的帶立.
但不是全靠聖靈自己躺平.
而是恐懼戰驚地立志實行.
是因為我們心感不配.
受寵若驚.
所以我們更加要認真謹慎.
全力以赴.
更加重要的不是我們孤軍奮戰.
是各位海船的君王.
藉著聖靈.
與我們同工同在同行.
下一次回來我們會看增展.
保羅告訴我們.
我們不是單獨作戰.
我們是一整支軍隊.
是一隊運動員.
一同為福音齊心努力.
我們一同祈禱.
因為你的福音不是一個輕飄飄的好消息.
你的福音是重於喜瑪拉雅山的重量.
預備了那個將要被獻上的身體.

$^{1201}$承受了我們應該承受的羞辱.
升上高天.
癡下聖靈.
並且要再來率領我們進入最終的海選.
主愛我們承認.
我們知道我們靠自己.
沒有辦法從這樣的福音去上程.
我們曾經好像那些跳上跳落的人.
又或者曾經是躺平放棄的人.
其實是你在我們裡面運行.
是聖靈住在我們裡面.
使我們立志,使我們實行.
所以求你幫助我們.
不要像他們.
不要像當年的以色列人一樣.
貪戀世界.
自以為是.
輕看這個救贖.
幫助我們恐懼戰經般.
活出救恩.
幫助我們堅守高舉這個生命的道.
在這個彎曲不謀的世代當中.
反射你愛子的光芒.
好像明光一樣照耀.
願榮耀歸給你.
直到永遠,仆主耶穌基督,明知彼得,阿們..
\newpage



\section{提多書 1:1-3:15}
\label{sec:032198}
\textbf{提多書:行為手冊}~~~~劉少康牧師
\newline
\newline
連結: \href{http://www.church.com.hk/acms/content.asp?site=cdc\&op=show\&type=product\&code=032198\&layout=sermon}{\texttt{032198}} ~~~~ 語音日期: 2026-05-20
\newline
\newline
\hyperref[sec:032197]{\small{< < < PREV SERMON < < <}}
~
\hyperref[sec:index_chronic]{\small{[返順時目]}}
~
\hyperref[sec:032199]{\small{> > > NEXT SERMON > > >}}
\newline
\newline
提多書 1:1-3:15
\newline
\begin{longtable}{cl}
\hline
\hline
章節 & 經文 (和合本修訂版)\\
\hline
1:1 & \begin{tabularx}{0.7\textwidth}{X} 神的僕人、耶穌基督的使徒保羅，為了使神的選民信從與認識合乎敬虔的真理— \end{tabularx} \\ \\ \relax
1:2 & \begin{tabularx}{0.7\textwidth}{X} 這真理是在盼望那無謊言的神在萬古之先所應許的永生， \end{tabularx} \\ \\ \relax
1:3 & \begin{tabularx}{0.7\textwidth}{X} 到了適當的時機，藉著傳揚福音，把他的道顯明了；這傳揚的責任是按著我們的救主神的命令交託給我的— \end{tabularx} \\ \\ \relax
1:4 & \begin{tabularx}{0.7\textwidth}{X} 我寫信給在共同的信仰上作我真兒子的提多。願恩惠、平安從父神和我們的救主基督耶穌歸給你！ \end{tabularx} \\ \\ \relax
1:5 & \begin{tabularx}{0.7\textwidth}{X} 我從前把你留在克里特，是要你將那沒有辦完的事都辦妥，又照我所吩咐你的，在各城設立長老。 \end{tabularx} \\ \\ \relax
1:6 & \begin{tabularx}{0.7\textwidth}{X} 若有無可指責的人，只作一個婦人的丈夫，兒女也是信主的，沒有人告他們放蕩，不受約束，就可以設立。 \end{tabularx} \\ \\ \relax
1:7 & \begin{tabularx}{0.7\textwidth}{X} 監督既然是神的管家，必須無可指責、不自負、不暴躁、不酗酒、不好鬥、不貪財； \end{tabularx} \\ \\ \relax
1:8 & \begin{tabularx}{0.7\textwidth}{X} 卻要樂意接待外人、好善、克己、正直、聖潔、節制， \end{tabularx} \\ \\ \relax
1:9 & \begin{tabularx}{0.7\textwidth}{X} 堅守合乎教義的可靠之道，就能將健全的教導勸勉人，又能駁倒爭辯的人。 \end{tabularx} \\ \\ \relax
1:10 & \begin{tabularx}{0.7\textwidth}{X} 因為也有許多人不受約束，說空話欺哄人，尤其是那些奉割禮的人。 \end{tabularx} \\ \\ \relax
1:11 & \begin{tabularx}{0.7\textwidth}{X} 這些人的口必須堵住，因為他們貪不義之財，將不該教導的事教導人，敗壞人的全家。 \end{tabularx} \\ \\ \relax
1:12 & \begin{tabularx}{0.7\textwidth}{X} 克里特人中有一個本地的先知說：「克里特人常說謊話，是惡獸，貪吃懶做。」 \end{tabularx} \\ \\ \relax
1:13 & \begin{tabularx}{0.7\textwidth}{X} 這個見證是真的。為這緣故，你要嚴厲地責備他們，使他們在信仰上健全。 \end{tabularx} \\ \\ \relax
1:14 & \begin{tabularx}{0.7\textwidth}{X} 不要聽猶太人無稽的傳說和背棄真理之人的命令。 \end{tabularx} \\ \\ \relax
1:15 & \begin{tabularx}{0.7\textwidth}{X} 在潔淨的人，凡物都潔淨；在污穢不信的人，甚麼都不潔淨，連心地和天良也都污穢了。 \end{tabularx} \\ \\ \relax
1:16 & \begin{tabularx}{0.7\textwidth}{X} 他們宣稱認識神，卻在行為上否認他；他們是可憎惡的，是悖逆的，不配做任何好事。 \end{tabularx} \\ \\ \relax
2:1 & \begin{tabularx}{0.7\textwidth}{X} 至於你，你所講的總要合乎那健全的教導。 \end{tabularx} \\ \\ \relax
2:2 & \begin{tabularx}{0.7\textwidth}{X} 勸老年人要有節制、端正、克己，在信心、愛心、耐心上都要健全。 \end{tabularx} \\ \\ \relax
2:3 & \begin{tabularx}{0.7\textwidth}{X} 又要勸年長的婦女在操守上恭正，不說讒言，不作酒的奴隸，用善道教導人， \end{tabularx} \\ \\ \relax
2:4 & \begin{tabularx}{0.7\textwidth}{X} 好指教年輕的婦女愛丈夫，愛兒女， \end{tabularx} \\ \\ \relax
2:5 & \begin{tabularx}{0.7\textwidth}{X} 克己，貞潔，理家，善良，順服自己的丈夫，免得神的道被毀謗。 \end{tabularx} \\ \\ \relax
2:6 & \begin{tabularx}{0.7\textwidth}{X} 同樣，要勸年輕人凡事克己。 \end{tabularx} \\ \\ \relax
2:7 & \begin{tabularx}{0.7\textwidth}{X} 你要顯出自己是好行為的榜樣，在教導上要正直、莊重， \end{tabularx} \\ \\ \relax
2:8 & \begin{tabularx}{0.7\textwidth}{X} 言語健全，無可指責，使那反對的人，因說不出我們有甚麼不好而自覺羞愧。 \end{tabularx} \\ \\ \relax
2:9 & \begin{tabularx}{0.7\textwidth}{X} 要勸僕人順服自己的主人，凡事討他的喜悅，不可頂撞他， \end{tabularx} \\ \\ \relax
2:10 & \begin{tabularx}{0.7\textwidth}{X} 不可私竊財物；要凡事顯出完美的忠誠，好事事都能榮耀我們救主神的教導。 \end{tabularx} \\ \\ \relax
2:11 & \begin{tabularx}{0.7\textwidth}{X} 因為，神救眾人的恩典已經顯明出來， \end{tabularx} \\ \\ \relax
2:12 & \begin{tabularx}{0.7\textwidth}{X} 訓練我們除去不敬虔的心和世俗的情慾，在今世過克己、正直、敬虔的生活， \end{tabularx} \\ \\ \relax
2:13 & \begin{tabularx}{0.7\textwidth}{X} 等候福樂的盼望，並等候至大的神和我們的救主耶穌基督的榮耀顯現。 \end{tabularx} \\ \\ \relax
2:14 & \begin{tabularx}{0.7\textwidth}{X} 他為我們的緣故捨己，為了要贖我們脫離一切罪惡，又潔淨我們作他自己的子民，熱心為善。 \end{tabularx} \\ \\ \relax
2:15 & \begin{tabularx}{0.7\textwidth}{X} 這些事你要講明，要充分運用你的職權勸勉人，責備人。不要讓任何人輕看你。 \end{tabularx} \\ \\ \relax
3:1 & \begin{tabularx}{0.7\textwidth}{X} 你要提醒眾人，叫他們順服執政的、掌權的，要服從，預備行各樣善事。 \end{tabularx} \\ \\ \relax
3:2 & \begin{tabularx}{0.7\textwidth}{X} 不要毀謗，不要爭吵，要和氣，對眾人總要顯出溫柔。 \end{tabularx} \\ \\ \relax
3:3 & \begin{tabularx}{0.7\textwidth}{X} 我們從前也是無知、悖逆、受迷惑，作各樣私慾和宴樂的奴隸，在惡毒、嫉妒中度日，是可恨的，而且彼此相恨。 \end{tabularx} \\ \\ \relax
3:4 & \begin{tabularx}{0.7\textwidth}{X} 但到了我們救主神的恩慈和慈愛顯明的時候， \end{tabularx} \\ \\ \relax
3:5 & \begin{tabularx}{0.7\textwidth}{X} 他救了我們，並不是因我們自己所行的義，而是照他的憐憫，藉著重生的洗和聖靈的更新。 \end{tabularx} \\ \\ \relax
3:6 & \begin{tabularx}{0.7\textwidth}{X} 聖靈就是神藉著我們的救主耶穌基督厚厚地澆灌在我們身上的， \end{tabularx} \\ \\ \relax
3:7 & \begin{tabularx}{0.7\textwidth}{X} 好讓我們因他的恩得稱為義，可以憑著永生的盼望成為後嗣。 \end{tabularx} \\ \\ \relax
3:8 & \begin{tabularx}{0.7\textwidth}{X} 這話是可信的。我願你堅持這些事，使那些已信神的人留心行善。這都是美好且對人有益的。 \end{tabularx} \\ \\ \relax
3:9 & \begin{tabularx}{0.7\textwidth}{X} 要遠避愚拙的辯論、家譜、紛爭和因律法而起的爭辯，因為這都是虛妄無益的。 \end{tabularx} \\ \\ \relax
3:10 & \begin{tabularx}{0.7\textwidth}{X} 分門結黨的人，警戒過一兩次後就要拒絕跟他來往； \end{tabularx} \\ \\ \relax
3:11 & \begin{tabularx}{0.7\textwidth}{X} 因為你知道這樣的人已經背道，常常犯罪，自己定自己的罪了。 \end{tabularx} \\ \\ \relax
3:12 & \begin{tabularx}{0.7\textwidth}{X} 我打發亞提馬或推基古到你那裡去的時候，你要趕緊往尼哥坡里來見我，因為我已經決定在那裡過冬。 \end{tabularx} \\ \\ \relax
3:13 & \begin{tabularx}{0.7\textwidth}{X} 你要趕緊給西納律師和亞波羅送行，讓他們沒有缺乏。 \end{tabularx} \\ \\ \relax
3:14 & \begin{tabularx}{0.7\textwidth}{X} 我們的人也該學習行善，幫助有迫切需要的人，這樣才不會不結果子。 \end{tabularx} \\ \\ \relax
3:15 & \begin{tabularx}{0.7\textwidth}{X} 跟我同在一起的人都向你問安。請代向在信仰上愛我們的人問安。願恩惠與你們眾人同在！ \end{tabularx} \\ \\
[1ex]
\hline
\hline
\end{longtable}
$^{1}$各位聽眾節目平安好.
我會用一堂的時間.
跟大家來看提多書.
提摩泰前書,提摩泰後書.
和提多書被稱為教目書信.
提摩泰前書和後書.
可以說是講到教目的搜索.
而提多書是一本行為的搜索.
保羅在第二次和第三次行程裡.
帶了提多一起去報道.
所以可以說提多是他的助手.
是他的傳道人.
而早年的時候.
保羅在第一次行程的時候.
就遇見提多.
就帶了他去信耶穌.
所以在提多書一章第四節裡說.
「因著信,做了我真兒子的提多」.
好像提摩泰一樣.
都叫他做「真兒子」.
所以提多和保羅有很緊密的關係.
然後保羅第二,第三次.
帶著提多去探訪不同的教會.
例如哥林多教會.
然後他就把一間教會的責任交給提多.
這間教會就在一個島上.
這個島我們和合本的翻譯.
就叫做「格里底」.
但是很多現在的新的譯本.
就叫它做「克里特」.
大家都知道在地中海那裡有一個島.
就叫做「克里特島」.
所以提多就做了那裡教會的牧師.
他都碰到很多的難處.
因為那間教會不成熟.
剛剛成立沒多久.
所以提多責任很重.
他要幫這間教會去設立監督,設立長老.
所以保羅在三章聖經裡面.
第一章就講到做監督有些什麼的職份.

$^{41}$什麼人才可以有資格去成為監督.
所以就給提多很多的教導,很多的提示.
還有特別提醒他.
你要很小心.
不要讓教會裡面有些假教師進來了.
那些假的教師在傳異教的道理.
在傳一些異端的道理.
所以保羅就教提多.
要怎樣去分辨.
怎樣要守住真道.
然後到了第二章.
就講到在一個家庭裡面.
你會見到長者.
你會見到年幼的.
有男有女.
你怎樣去面對.
怎樣去處理很多家庭之間的關係.
今天我們做牧師的.
很多時候都要處理的.
都是家庭一些的問題.
因為家庭裡面有時都會有些糾紛的.
那些長者覺得年輕的不尊敬他.
年輕的又覺得那些長者不明白他們.
所以很多問題需要去處理.
特別婚姻.
一男一女結了婚之後.
很多時候都會有不同的衝突.
大家有很多不同的看法.
結婚的時候.
可能想著我怎樣改對方.
但是沒有什麼可能的.
我們只能夠遷就.
只能夠彼此的和諧.
彼此一切同心同合.
我結了婚之後.
都和師母在很多方面是需要磨合的.
特別我性急.
有時有些臭脾氣.
所以師母都要很多的忍讓.
很多幫助我去成長.

$^{81}$看到我自己的真面貌.
然後怎樣我去改正自己一些.
很性急.
很容易發脾氣.
很喜歡說話.
很硬.
這些都是我們很需要在婚姻裡面.
大家彼此學習的.
然後在第二章開頭是講到家庭方面.
後面開始講到社會方面.
特別在第三章.
保羅就提到.
你要信服掌權者.
因為在一個社會的制度裡面.
有執掌權並的人.
我們應該信服他們.
又為他們祈禱.
讓他們成為合神心意的人.
不要做一些不合神心意的事.
還有保羅又提醒.
他說信徒是應該做正經的事.
什麼叫做正經的事呢?.
就是你應該有正當的行業.
你有你自己的工作.
你要做正當的事.
來謀生.
你的生活就會得到神的帶領.
神的祝福.
神的恩代.
所以保羅就是這樣勸勉帶領提多.
這個年輕的童工.
讓他可以不斷地來學習.
特別在第二,第三次行程裡面.
保羅就帶了他去哥林多.
後來又去了以弗所.
然後又叫他快點去尼哥波納去見保羅.
因為保羅定義是要在那裡過冬天.
所以他說提多你快點來尼哥波納那裡見我.
因為我很想跟你在一起.
所以這個是保羅在晚年的時候.

$^{121}$寫信給提多的.
他提醒提多.
你要小心假教師.
你要避免一些猶太人.
那些無稽之談.
他們偏離真道.
你要.
他們偏離真道就喜歡那些愚昧的爭論.
哇 討論來討論去.
都是一些沒有什麼益處的爭論.
你要小心.
然後你特別挑選教會領袖的時候.
一定要很謹慎.
提一些可以造就的人.
可以擔當責任的人.
成為教會的領袖.
因為如果你立錯了教會的領袖.
就對教會有很大的影響.
所以我們透過提多書.
看見保羅那個牧者的心象.
他不單止關心這個年青的童工提多.
他更加關心的是.
教會的前景是怎樣.
他很希望教會是神的家.
是永生神的家.
教會要在地上強壯.
教會的根基要建立得好.
教會裡面的信徒要不斷地成長.
然後他們一起追求長進.
那麼神就會喜歡.
於是教會可以在地上永耀神的名.
今天我們是地上的教會.
無論在哪一個城市.
無論在世界哪一個地方.
我們只有一個目的.
教會是屬於永生神的家.
是真理的主石和根基.
教會的根基必須要牢固.
教會必須要穩固起真理上.
然後教會站立得穩.

$^{161}$就可以造就信徒.
叫信徒同樣強大.
然後信徒得著造就.
就可以繼續將福音傳開.
我記得以前在我帶領教會的時候.
就用Win, Build, Send.
這三個英文字.
得人得更多的人.
領人歸主.
然後去建立他們.
建立信徒.
叫他們生命穩固.
遭咒滿途.
叫他們得到訓練.
然後就Send.
差他們出去做宣教的工作.
去到各地建立教會.
廣傳福音.
將福音傳遍天下.
所以我們必須要得人建立.
又差人.
這就是教會一個循環.
讓我們不斷在教會一個很重要的課題上.
我們去學習.
求神幫助我們.
讓我們在教會裡得到好的造就.
我們成長.
我們一起努力為神工作.
願天父祝福給各位.
\newpage



\section{士師記 5:7}
\label{sec:032199}
\textbf{士師記析讀 (7) 女權當道}~~~~梁永善牧師
\newline
\newline
連結: \href{http://www.church.com.hk/acms/content.asp?site=cdc\&op=show\&type=product\&code=032199\&layout=sermon}{\texttt{032199}} ~~~~ 語音日期: 2026-05-27
\newline
\newline
\hyperref[sec:032198]{\small{< < < PREV SERMON < < <}}
~
\hyperref[sec:index_chronic]{\small{[返順時目]}}
~
\hyperref[sec:032200]{\small{> > > NEXT SERMON > > >}}
\newline
\newline
士師記 5:7
\newline
\begin{longtable}{cl}
\hline
\hline
章節 & 經文 (和合本修訂版)\\
\hline
5:7 & \begin{tabularx}{0.7\textwidth}{X} 以色列農村荒蕪，空無一人，直到我底波拉興起，興起作以色列之母！ \end{tabularx} \\ \\
[1ex]
\hline
\hline
\end{longtable}
$^{1}$我們上個星期看到滿朝選謙受益那句說話.
就是剛剛我調轉了.
就是謙朝選滿受益.
我們看到巴拉被神呼召.
叫他去為以色列人作戰.
但他一定要帝波拉和他在一起.
原來就是膽怯,懼怕.
好像很謙卑那樣.
但是因此就引來了損失.
就是最終榮耀不能歸給他.
如果滿有信心的時候.
反而會得到益處.
上次我們看到.
他們要面對迦南王耶賓.
他有900輛戰車.
一個很龐大的軍隊.
但終於打贏了仗.
而主要一個人物就是.
帝波拉.
但上次我們較少著墨在帝波拉身上.
焦點就放在那場戰爭裡面.
但今天我們會探討戰爭.
得到勝利兩個很重要的人物.
就是兩個女性.
還有我們會看.
接著是帝波拉之歌.
《士兵記》第四章記載了戰爭.
但第五章是一個完全不同的題材.
也是整卷的《士兵記》裡面.
唯一一次用詩歌的題材出現.
當中提及兩個女性.
第一個就是帝波拉.
她是以色列人的領軍者.
另一個女士就是瓦易.
她是殺死元帥西西拉的.
還有一個人.
但沒有記名的.
就是西西拉的母親.
我們稍後會看.
這兩個人.

$^{41}$即是前兩者.
獻出了一己之力.
打退了以色列人的敵人.
我們看到似乎是女權當道.
但原來選擇一些人.
最重要是她有沒有心事奉神.
性別不是最重要的.
我們千萬不要用這段經文說.
女性很重要.
稍後我們會探討.
因為他們忘恩.
所以帶來了惡果.
我們先看《士兵記》第四章第二節.
耶護說:就把他們賦予在下所作王的甘納王耶賓手中.
他的將軍是西西拉.
住在外邦人的夏羅茨.
以色列人因為享受了一段很長時間的太平日子.
如果我們看回第三章第三十節.
他們享受了八十年的太平日子.
於是又再顧太復萌.
特別是八十年之後.
又有另一代的新人起來.
那些又離棄神.
那些又行爺說眼中捍衛惡的事.
不僅是違反了神的誡命.
而且執迷不悟.
所以神再一次讓他們受外敵入侵.
他們再一次成為奴隸.
因為他們忘記了.
其實他們選擇去敬拜神所敵對的偶像的時候.
他自動就會成為神的敵人.
這時候就會面對神的懲罰.
所以我們要看到.
原來問題的關鍵並不在於神是否站在子民那邊.
最重要的就是子民是否站在神那邊.
這才是貞潔的所在.
當我們遇到困難,壓力,挑戰的時候.
我們要好好地檢視.
究竟是神容許.
還是刻意讓我們兩手磨練.

$^{81}$藉著這樣使我們得到更加成熟呢?.
就好像弱室一樣.
他沒有犯錯的.
但是被出賣.
後來還要坐牢.
但原來神就要去詛咒他.
還是因為我們犯了錯.
違背了神的指令.
所以才會面對這種光景呢?.
下星期我們會看到.
基澱如何來質問神.
為何你不去保護我們呢?.
他沒有檢討這個因素.
所以我們要分辨究竟是神的懲罰.
還是神的磨練.
當我們面對這些困難的時候.
神是永恆不變的.
只要我們順從祂的時候.
我們一定是永不動搖.
夏所王耶賓極可能是一個至高君王的身份.
管轄著迦南地區的重迦南國王.
而當時夏所是一個非常繁榮的大都市.
但我們又要看若書雅曾經毀滅過這個城市.
我們看看《若書雅記》第十一章第十節.
「當時若書雅轉回奪了夏所.
用刀擊殺夏所王.
素來夏所在者諸葛鐘是為首的」.
看不看到?.
素來夏所在諸葛鐘是為首的.
即是夏所是一個非常非常重要的地方.
但若書雅曾經毀滅過夏所.
亦殺了當時的王.
但因為夏所是一個商業貿易很重要的地位.
所以也是列國的人要去搶奪這個地方.
若書雅征服了夏所之後.
後來迦南人又再一次搶回這個夏所.
因為這個地方掌握了整個迦南地區的經濟命脈.
控制了它的時候.
就可以控制到當時所有商業的運作.
所以我們看出這一次出戰的是.

$^{121}$夏所王的元帥西西拉.
他是一個很勇猛的元帥.
我們再看回上次我們看過的經文.
《士師記》第4章第3節.
「耶賓王有鐵車九百兩.
他大大欺壓以色列人二十年.
以色列人就呼求耶和華」.
這次我們的焦點不是放在那九百兩的戰車裡面.
而是放在那個年份裡面.
大大欺壓了他們二十年.
他們才懂得求問神,呼求神.
二十年絕對不是一個小的日子.
但是他們竟然二十年之後才懂得去呼求神.
可見如果人要還耿叛逆的時候.
是可以冥頑不靈的.
甚至有時候真的不去到谷底的時候.
他是不會有反省的.
所以某個程度是自作孽的.
我們現在就看兩個女性.
第一就是底波拉.
綜觀過去的歷史.
以色列人從來沒有像今次那樣打一場大仗.
這是他們整個民族歷史以來.
第一場很巨大的戰爭.
以前不是在利菲町打過的.
但是利菲町的戰役絕對不及這一場的戰役.
然後約書亞進去之後.
就滅了耶利哥城.
但是不是打仗的.
後來也打艾城.
所以這一場是很特別的戰爭.
但是這次真正的領導者是一個女性.
是一個毫無畏懼的女性.
去領導一群人.
來與迦南的大元帥作戰.
《士師記》的作者.
是以女先知的身份帶出底波拉的影響力.
雖然摩西的姐姐米利欽.
《聖經》也說她是一個女先知.
但是《聖經》對米利欽的描繪是很少的.

$^{161}$甚至指出米利欽因為有妒忌的心.
後來在背後批評摩西.
結果神就要懲罰米利欽.
米利欽就患了大麻瘋.
但是《士師記》的記載底波拉的篇幅是很多的.
我想問問大家.
你停一停想一想.
究竟出現了多少個女先知?.
但是出現的女先知.
沒有認識她有多少呢?.
但無論如何底波拉是最出名的.
我給大家看看教會裡面還有哪幾個女先知.
麻煩Yvonne幫我讀《聶王記下》二十二章第十四節.
「於是祭司希臘加和阿希金.
阿加波,沙番,阿薩雅.
都去見女先知古勒戴.
古勒戴是掌管禮服沙龍的妻.
沙龍是哈爾哈斯的孫子.
特瓦的兒子.
古勒戴住在耶路撒冷第二區.
他們請問於他」.
總之就是有一個女先知叫做古勒戴.
好了 接著黎希米記第六章第十四節.
又記載了另一個女先知.
「我的神娜多比亞,坦巴拉.
女先知羅亞底和其餘的先知.
要叫我去懼怕.
求你紀念他們所行的這些事」.
這個女先知羅亞底真是差的.
她站在敵人那邊勾結外國勢力.
來對付黎希米.
但她都被稱為女先知.
好了 以賽亞書第八章第三節.
「我以賽亞語妻子.
願問著女先知同實她懷孕生子.
也會話就對我說.
給她起名叫馬克爾沙拉勒克斯伯斯」.
當年我們考聖經的時候.
就要記得很清楚.
就是以賽亞的兒子.

$^{201}$所以如果由我讀就會順很多了.
就是馬克爾沙拉勒克斯伯斯.
OK 但是你留意到以賽亞的妻子.
都被稱為女先知.
但是沒有記她的名字.
這就是舊約裡面所說的女先知.
就是不是很多.
但是這麼多人來講.
底波拉的身份是最獨特的.
因為除了先知之外.
她更加是以色列人的領導者.
好 我們看看《士斯記》第四章第五節.
「家住在以法連三地.
拉瑪和伐特利中間.
在底波拉的中樹下.
以色列人都上他那裡去聽半段」.
以色列人有問題的時候.
就去那裡請底波拉作出判斷.
他就坐在中樹下有樹蔭.
就可以遮擋猛烈的太陽.
來聽訴訟作出判斷.
連那棵樹都被稱為底波拉的中樹.
可見他的影響力是多大呢?.
下星期我們會看到基甸.
他的父親也有棵樹以他命名的.
但是無論如何.
底波拉的中樹.
就可以反映底波拉的影響力.
他住在以法連三地.
拉瑪和伐特利的中間.
根據第一章第二十二到二十六節裡面記載了.
這是當年約瑟加征服了黑人.
獲得重大勝利的地方.
而黑人就是以色列所有敵人之中.
最具經驗的戰士.
以法連三地是位於以色列的中央.
因此所有的以色列人.
都能夠到那裡尋求判斷.
有這麼多人尋求判斷的時候.
可見他一定是一個有崇高品格,正直.

$^{241}$有公平判斷.
而且有好明星.
有好的明星.
這些人是要尊重他的.
底波拉一定是一個好明星的人.
但是無論如何.
大家要接受.
特別是姐妹們.
在那個年代裡面.
女性始終有她的限制.
女性仍然是丈夫的財產.
所以被隱蔽.
所以在第四章裡面.
會講到底波拉是拉比多的妻子.
同樣瓦易.
她被重複稱為是希伯的妻子.
丈夫的名字多次出現.
就證明了當時的女人.
她的名譽與丈夫的榮辱.
是有密切的關係的.
底波拉被稱為以色列之母.
第五章第七節.
以色列人中的官長停職.
直到我底波拉興起.
讓我興起作以色列的母.
她的興起並不是當時的社會正常現象.
底波拉不單單是一位先知,事師.
她更加是一位母親.
所以她身兼政治,事奉,家庭.
多重身份的婦女.
她應該沒有受過任何軍事訓練.
當時沒有女性參戰.
神興起的領導者大部分是男性.
好像我們看過的俄陀列,以弗.
但神也會使用一些女性.
但不得不承認.
在底波拉時代的領導者.
顯然不是女性.
所以由女性去領軍,對抗敵人.
是很罕有的事.

$^{281}$當然神有主權.
去呼召任何一個祂所定義使用的人.
當然這個人被呼召的時候.
與祂本身的質素和信心.
是有密切的關聯.
神絕對不會任意妄為.
隨便起用一個人.
可以看到祂的本質,潛質,這是很重要的.
當我們要誠實去面對這張聖經的時候.
就會看到底波拉是被當時的社會環境.
文化來限制著.
雖然祂扮演了一個很顯著的領袖角色.
但因為祂是一個女人.
所以並不具備了當時成為事司的必要條件.
祂沒有軍事作戰的能力.
更加沒有親自帶領軍隊上前線去廝殺.
底波拉是以色列人的領袖.
祂的突出正如我們上次所看.
是巴拉的餘孽和失敗.
用現今的角度去看.
就好像男女不得平等.
我承認男女是不得平等的.
時移世易.
現在女性的地位一定是高過男性的.
這是我親身感受到的.
我在外面可能是風光.
回到家裡是滄桑的.
因為我家裡有兩個女性都比我高.
但是記住.
在當時的文化來說.
祂是一個現實.
所以我們絕對不要用現在的標準.
去量度數千年前的情況.
很多女權主義者.
特別是強調底波拉的角色.
但我絕對不是貶低底波拉的身份.
性別的對照並不是時移世易的主要主題.
作者在故事中對男女的對照.
主要偏重於男女在原始社會裡的真實領導能力.
而不是強調兩者之間的性別差異.

$^{321}$在當時的社會.
女人根本沒有機會接受軍事訓練.
所以史記裡沒有一個女人是直接參與戰爭的.
神使用這些女性.
就是為了要對照那些男士的餘孽和退縮.
所以這段經文強調的是領導能力.
而不是性別的差異.
神沒有呼召底波拉上戰場.
祂徵召的是什麼?.
上次我們看到是巴拉.
因為巴拉具有作戰的能力.
我們不要忽略巴拉.
祂很輕易就招聚了一萬名的跟隨者.
所以祂本身就是一個很卓越的軍事領袖.
但是神所看重的.
不單單是誰做軍事領袖.
而是那個人的內心是怎樣的.
在史記中女性所受到的限制.
已經不是今天我們無論英國或者香港.
對女士所面對的問題.
現在很多女性是強過男性的.
而且很多時候在社會,教會的領導地位上.
我們已經看到女性的領導能力或者人數.
比男性佔多.
我不知道你怎樣看.
以往冥恩堂很多屆以來.
都是女的執事多於男的執事.
所以今天女性可以自由接受.
各種不同的領袖訓練.
與數千年之前是絕對不同的.
神會興起一些最不合資格的人.
來執行祂的工作.
這個人的條件最重要的是什麼?.
他要尊重神.
甘心樂意為神擺上.
以及有信心.
在經文中我們看不到.
神對底巴拉有什麼特殊的呼召.
但是祂竭盡所能的成果.
肯定是神給祂的呼召和祝福.

$^{361}$所以你看到一些很偉大的宣教士.
未必有博士學位或者很高的學歷.
但是神卻是使用他.
所以有時候我反而會看.
甚至作為牧師.
不稱自己為牧師.
就要用博士來稱自己.
我覺得牧師的身份才是最重要的.
還有記住.
神用一個人不只是看他的學歷.
我們看看瓦易.
這次我們看到女人的表現.
是代替了男人應有的角色.
雖然瓦易的丈夫希伯這個名字.
出現在經文中.
但是真真正正的主角就是瓦易.
那時候兒童底巴拉.
他比巴拉強.
瓦易一直使他的聲名蓋過他丈夫.
瓦易殺死西西拉.
其實應驗了底巴拉的預言.
《士師記》第四章第九節.
「底巴拉說:我必與你同去.
只是你在所行的路上急不著榮耀.
因為耶和華要將西西拉交在一個婦人手裡.
於是底巴拉起來.
與巴拉一同往基底斯去了」.
這段經文上次有參加.
就會知道神父召巴拉去做事司.
但他就說要底巴拉.
你一定要與我要去.
你不去我就不想去.
但是底巴拉肯去.
但他就說你的榮耀.
將會交在一個婦人的手上.
這個就一早已經證實了.
神已經有特別的安排.
瓦易殺死西西拉的表現.
就超越了巴拉.
我們看看.

$^{401}$基里人希伯和摩西有親屬的關係.
因為瓦易是基里人.
我們看看《士師記》第一章第十六節.
「摩西的內卿是基里人.
他的子孫與猶大人一同來了中樹城.
往西拉德爾南的猶大放野去.
就住在民中」.
摩西的內卿在原文中是外府.
是基里人.
他居住的地方很近戰場.
是一個很重要的戰略地位.
但是為何與摩西這麼有密切關係的人.
竟然與迦南王耶賓有這麼好的交往呢?.
這很可能是在摩西的世代不久之後才發生.
即是說摩西的親屬.
亦肆無忌憚地與迦南人有交往.
他們已經喪失了辨識敵人和真實信仰的能力.
這是很可悲的.
但是現實很簡單的舉個例子.
我信耶穌.
但不代表我的兄弟姐妹會信耶穌.
因為我很少與他們交往.
因為我中二已經離家出走.
去了我外婆的家住.
然後沒多久我就進了神學院.
所以我與我的兄弟姐妹的關係較為疏離.
例如我的妹妹會怎樣呢?.
她是拜菩薩的.
不等於我信耶穌.
她們一定會信的.
所以就是這樣.
接著我們看看《士師記》第四章17至18節.
看看西西拉如何死在亞裔的手上.
只有西西拉步行逃跑 到了基里人希伯之妻雅益的帳篷.
因為下嫂王耶賓與基里人希伯家和好.
雅益出來迎接西西拉.
對她說:請愛主進來 不要懼怕.
西西拉就進了她的帳篷.
雅益用被子將她遮蓋.
這裡看到來到基里人希伯之妻雅益的帳篷.

$^{441}$一定要冠以她丈夫的名字在前面.
然後說到原來下嫂王耶賓與基里人希伯家和好.
希利人應該與以色列人關係好.
但為何會與下嫂王和好呢?.
上次我們已經看過.
可能是做生意.
因為經濟利益.
所以可以與一些人建立關係.
這是我們上次所看到的.
雅益可能因為西西拉的明星.
曾經與她交往.
所以認識這個將軍.
於是她就用被子蓋住西西拉.
可能因為當時是冬天.
或者因為要假裝.
要將她收藏.
不要讓別人知道.
解除了她心裡的那種防禦.
西西拉想喝水.
我們上次有看過.
但雅益給她奶喝.
及後她睡了之後雅益就用帳篷的缺紙.
將她帶去.
釘穿了她的雲箏.
戳到地上.
原來顯赫一時的將軍.
竟然在毫無還手的情況下死去.
這位南征北討.
征服了不少地方.
樹立了無數帳戶.
釘了無數缺紙的將軍.
竟然被一枝缺紙.
釘穿了她的頭而死.
所以你會看到這是一個很大的諷刺.
我想說.
《士師記》是用很多.
很特別的諷刺的說話.
去出現的.
以色列人的領軍者巴拉.
似乎不知道西西拉的去向.

$^{481}$及後他需要雅益告訴他.
指引他西西拉在哪裡.
他才會知道.
《士師記》第四章二十二節.
巴拉追趕西西拉的時候.
雅益出來迎接他說.
老伯 我將你所尋找的人給你看.
他就進入帳棚.
看見西西拉已經死了.
倒在地上.
缺紙還在他品中.
你看看.
原來巴拉又遲了一步.
他去找西西拉在哪裡呢?.
雅益就告訴他.
這個除了應驗了底波拉的預言之後.
榮耀歸給一個婦人之外.
我們都見到巴拉錯失了良機.
所以我們看到一個很重要的教導.
如何能夠成為合神心意的人呢?.
就是那個人願意信服神.
但如果我們帶有條件的心態來服侍神.
神就很可能限制了那個人.
對那個事供的參與度.
特別是神如果給我們機會去服侍的時候.
但如果我們沒有珍惜 缺乏了信心.
那麼我們服侍的那一刻.
應該有的榮耀都會被打折扣.
但其實最重要的是我們要知道.
每一場戰爭的勝利.
那個真正的得勝者.
真正的英雄就是神.
巴拉沒有使西西拉的軍兵去賄亂.
沒有使他的戰車陷在泥沼裡面.
一切都是神的作為.
賴以設下陷阱來去殺西西拉的行動.
更加顯出神能夠藉著人認為不可能的方法而成功.
所以記住 人的資源並不是最可靠.
有人靠車 有人靠馬.
但我們卻要靠如何?.

$^{521}$九百輛的戰車 記不記得?.
是一個很龐大很強勁的軍事力量.
但都不能夠成為他們永恆的倚靠.
即是暫時西西拉的倚靠.
但這些戰車不能產生作用的時候.
西西拉馬上變成一個膽小如鼠的逃亡者.
神要讓人民知道.
雖然他們在兵器和科技上比敵人落後很多.
但無論軍備或戰鬥能力如何變化.
最後的勝利都是屬於願意信服神的人.
神並沒有為以色列人提供勢均力敵的戰車.
來對抗西西拉.
反之是用一條很小的 不是很長的鋳.
來將西西拉的生命結束.
所以我們看到.
不要小看自己.
如伍秉儀如那個小孩子.
獻上這些看耶穌的神.
同樣我們不要只是將目光.
注視在我們有多少的資源上面.
反而是要將它轉眼仰望我們的神.
教會也是的.
有時教會很著緊我們現在有多少人.
做多少事 我們有多少經濟.
我們有多少儲備就可以做多少事.
甚至他們是很好的.
教會一定要有四個月的儲備.
不可以少於這四個月的儲備.
否則怎麼發工資呢?.
但有時我說我不會只看這些.
有需要就先用它.
弟兄姊妹就會懂得如何奉獻.
我們看第一部分.
就是看了底波拉和瓦易.
接著我們會看我們平時很少看的一章聖經.
第五章就是底波拉之歌.
因為是查經.
所以我很希望整卷的時詩記.
逐章逐章跟大家看.
也很少人講道的時候.

$^{561}$就用這篇歌來作為講道之用.
所以因為是查經.
所以我們會研讀下去.
這章的聖經似乎是重述前一章的戰爭.
但如果我們細心研究.
就會發現這章有些細節.
是前一章沒有提及到的.
這首詩歌的作者應該就是底波拉.
因為在第五章第七節裡.
他使用了單數的「我」字來成為代平詞.
這首歌是以婦女作為中心的.
所以風格較為獨特.
其中主要的角色就包含了底波拉,瓦易.
這幾位同學我跟大家講過.
還有就是西西拉.
不知道什麼名字的母親.
這首歌就是描述戰爭的過程.
主要由預備戰爭到最後勝利.
所以這首詩歌是從事件的發展來做分段.
這首歌也帶出了迦南諸神的無能.
他不能夠拯救那些異教徒.
而這首歌集中刻畫每一個人物.
這首歌與以色列人過紅海之後.
米利欽去歌頌神有很多相似的地方.
這首歌有很多平衡之處.
所以底波拉可能也藉著.
被米利欽作的歌傳下來.
是受了這方面的影響.
於是寫下了這首詩歌.
他依循出埃及的傳統.
唱出神戰勝異教神明和敵人的偉大作品.
底波拉將以色列人得勝.
視為另一次全新的出埃及的經歷.
每一個人都參與其中.
在神的計劃當中要各盡其職.
領導者要忠心的領導.
跟隨者願意信服.
整體投入才有成功的果效.
所以底波拉不只是歸功於自己.
我們看看《士師記》第五章第三節.

$^{601}$君王啊要聽.
王子啊要擇耳而聽.
我要向耶和華歌唱.
我要歌頌耶和華以色列的神.
所以你留意到的時候.
這裡很清楚的說到甚麼呢?.
君王你要聽.
王子你要聽.
要聽甚麼?.
就是要去聽我向神的歌頌.
底波拉將他們的成功視為一個群體.
所以你留意到第五章第二節.
因為以色列中有軍長率領.
百姓甘心犧牲自己.
所以你應該要稱頌耶和華.
就是說這個是有軍長帶領.
有百姓甘心的犧牲.
彼此配合一個整體的效果.
所以領導是有非常重要的責任和影響力.
如果他們不好.
他們的失敗.
跟隨者就會同樣陷入這個失敗當中.
所以領導必定要自己.
是先完全的肯去跟隨神.
他們的責任就是要尋求和理解神的心意.
要信服而行.
否則只會錯誤地帶領群眾.
陷入一個失敗的境地.
教會也是一個整體.
如果只是一,兩個會友.
愛主,跟隨神,信服神.
但其餘的人就斷了手.
不理會,不參與的時候.
教會是絕對不會成長的.
如果大家能夠同心合力的話.
仰望主.
教會就能夠有好的發展.
當然教會的領導牧者.
我深深體會到.
是非常非常重要的.

$^{641}$然後第三節就說到.
君王啊.
你要聽王旨.
你則以耳聽.
大家就這樣讀的時候.
會覺得.
好吧 很簡單.
但是你有沒有細心想一想.
《時詩記》.
經常有一個主題就是.
當時以色列中會有王.
為什麼會提到君王和王旨呢?.
其實很簡單的解釋.
記不記得.
在《申命記》裡面.
神都頒下了一個指令.
如果君王來守神的律法.
遵照神的吩咐去行.
這個國家就會強盛.
如果君王偏離神.
拜偶像的時候.
而又屢勸不改.
國家就會滅亡.
但是那個時候.
他還在漂流.
哪有國家.
怎會有君王呢?.
但是這裡就是.
在神的啟迪之下.
就好像在《申命記》裡面.
預兆了將來的君王應該要怎樣.
神同樣用他的靈.
啟迪了給提柏拉.
所以會寫下這樣東西.
當然另一個解釋.
就是因為他們有十二支派.
他們已經分散在不同的地方居住.
各族各個支派都有自己的政治形態.
所以支派裡面的領袖.
管理那個地方.

$^{681}$就好像有君王和王旨那樣.
這次我們看第四節和第五節.
耶和華你從西域出來.
就以東地行走那時地震天漏雲也落雨.
山見耶和華的面就震動.
西那山見耶和華的以色列神的面也是如此.
這裡好像是一個詩歌的描繪.
講一些地方.
提柏拉向他的子民呼籲後.
就以整首詩納為根基.
根基是甚麼?.
就是神的榮耀和能力.
西異就在南邊.
而以東人的領土上面.
西異的旁邊就是西奈半島.
我們很熟悉.
而西奈半島就是耶和華.
初次向摩西頒下律法的地方.
基里人居住就在摩押營的地區.
就是以東人的鄰居.
即是右邊.
這時候你留意到.
所以他提及到雲也落雨.
就預兆到甚麼呢?.
西西拉的軍隊將會受到毀滅.
敵人會被暴雨所造成的急流所沖走.
因為上次出不到拍攝點.
其中一張拍攝點就拍到.
當時在斐濟的島嶼裡面.
短短一個上午不到幾個小時下雨.
雨水突然之間漲得很厲害.
去到屋的二樓.
地上全部被淹沒.
這就是上次拍攝點裡面.
有這張圖片讓我們看到.
所以本來基遜河是很乾涸的.
很少一條小溪.
但當下雨下來的時候.
突然之間洪水就暴漲.
按照傳統的說法.

$^{721}$迦南的眾神是來自北方的高山.
而耶和華是從南方而來.
頒布律法的西拉山就在南部.
所以被稱為以色列的神.
就表示祂特別眷顧祂的子民.
這也表示神是戰士.
神童來自西以北面迦南區域.
和西西拉直接為敵.
因為神是由南方西拉而來.
而且耶和華神就在迦南諸賈神之上.
祂的能力和那些賈神是無法比擬的.
底玻拉就用迦南宗教的言語去描繪神.
表達神就是獨一的真神.
巴力的賈神.
根本沒有辦法去掌管大自然.
生育和人類的命運.
所以這場戰爭.
不單能見到肉體廝殺.
而是一場神的明星和信仰的屬靈爭戰.
當時子民已經和外族融和了.
第四章第一節.
以色列人遊行耶和華眼中看為惡的事.
所以作為一個屬靈的領袖.
一定要將子民帶回到正途上.
這時候他的詩歌就充滿了抗衡文化的力量.
我們看第六和第七節.
在亞蘭之子生家的時候.
又在瓦易的日子帶道無人行走.
都市擾渡而行.
以色列中的官長停職.
直到我底玻拉興起.
等我興起作以色列的母.
透過這兩節經文我們看到一件事.
一個人一個社會的墮落並不是朝夕突發的事.
是在漫長的歲月之下失落.
正如黃福遠這次的大火好像很突然.
但現在你看聽證會.
各個部門都不去承擔自己的責任.
以及沒有去理會這些事.
是制度上的壞.

$^{761}$但最後就會出了這些惡果.
所以記住原來墮落不是一朝一夕的.
是慢慢慢慢的去腐敗.
這段經文特別提到生家和瓦易.
兩者絕對不會是同一個時期.
因為以弗死後.
以色列人過了八十年的平安日子.
所以他們才再度陷入困境中.
因為拜假神.
所以你留意到.
生家是他們所屬的地區幫助以色列人.
及至他們犯罪行惡就陷入困境中.
神就興起底波拉帶領他們出困境.
所以生家和瓦易很可能是同屬的世代.
但以弗就屬於前一代的事先.
這裡特別提到甚麼呢?.
大盜無人行走.
是描寫了當時淒慘的境況.
因為通商的大盜的命脈.
已經被迦南人控制住.
切斷了以色列人的經濟命脈.
所以他們不能行走大盜.
大盜無人行走.
我們再看第五章第八節.
以色列人選擇登陣.
真正的士就臨到城門.
那時以色列人四萬人中.
豈能見盾牌槍矛呢?.
這裡講出以色列人當時的困境.
因為他們沒有武器.
盾牌就是保護軍人的武器.
槍矛就是攻擊敵人的利器.
但他們兩樣都沒有.
這時他們縱使有一定的人數.
四萬人.
但就沒有足夠的武器去攻打.
難道赤手空拳可以來戰勝.
滿有很多利器兵器的敵人嗎?.
所以你留意到他們怎樣都不能和敵人爭戰.
但真正令他們戰敗的原因.

$^{801}$不是因為他們缺乏武器.
是因為他們選擇了新的神去敬拜.
即是他們離棄了耶和華.
在約瑟克時代.
他們都不需要用矛槍和盾牌的兵器.
就能夠戰勝敵人.
耶利哥的倒塌不是因為他們有精良的武器.
是因為他們信服沉默無聲.
走了十三次.
於是就發出這個呼喊.
城就倒塌了.
所以離棄神就是他們失敗的主要原因.
當他們敬拜迦南地生育的眾假神.
但他們的生育沒有因為這樣而增加.
因為在戰爭和災難打擊之下.
人口出生率就低了.
正如我們之前說過.
香港的生育率是不斷降低的.
差不多壓到底了.
比其他國家更加少.
這時候為甚麼這麼多人都不肯生孩子.
我現在都認識在這幾年結婚的人.
到現在都是不肯生孩子的.
無謂害人 是不是.
我們都曾經說過.
在索罕頓兒子讀大學.
但才懷孕 是不是很厲害呢.
因為他們生活穩定.
可能就會面對.
曾經我們看過大衛.
大衛當他真是作王的時候.
他才生兒育女.
所以你留意到原來當時.
他們拜巴黎這班以色列人.
但他們生育竟然減少.
底窩拉是一個被神重用的領袖.
但一個人不能發揮很大的果效.
就需要人來配合跟隨.
所以這首詩歌裡面.
他沒有提及巴拉的威脅.

$^{841}$一定要他來帶去陪伴.
更加沒有提及亞裔擊殺西西拉.
而使巴拉被人輕描.
他有很寬大的胸襟.
他樂意將德性的功勞.
跟人來分享.
我們看看第十二十三節.
底窩拉啊輕起輕起.
你當輕起輕起.
唱歌 阿彼羅恩的兒子巴拉啊啊.
你當奮興 努力你的敵人.
那時有餘靜的貴族.
和百姓一同下來.
耶護華當林為我攻擊勇士.
你留意到他特別有提到.
有貴族和百姓一齊下來.
但你有沒有留意到.
是誰叫底窩拉輕起呢?.
底窩拉啊輕起輕起.
你當輕起.
即是叫底窩拉幫手.
這個描述在第四章沒有出現過.
似乎在以色列人當中.
有一班人很想主動去解決問題.
於是他們覺得底窩拉是一個合適的人.
於是請他挺身而出.
跟隨者和領導者有合一的心智.
於是就能夠有成功的機會.
所謂「時窮節乃賢」.
往往在大時代.
在惡劣的環境中.
我們見到人性的光輝和醜惡.
有些人會在凶險中.
仍然堅守立場.
甚至挺身而出.
有些人就會變成沉默.
效忍.
但有些更差.
放棄了自己的原則.
轉態.

$^{881}$去幫助邪惡的勢力.
底窩拉之歌.
就帶出了這兩類人的差別.
這就是我們平時所看不到的.
為何會說這些呢?.
因為底窩拉曾經呼召一些人.
有些人參與.
有些人退縮.
甚至有些人跟外敵入侵.
我們看看這兩段經文.
在第十四到第十八節裡面.
有哪些人來參與?.
有哪些人來拒絕?.
你留意到我用紅色字的.
就是參與中的.
藍色字的.
就是拒絕,參與.
我們請電話同學讀出來.
有根本在亞馬力人的地.
從二百年下來的.
亞敏在民眾跟隨你.
有掌權的從瓦吉下來.
有持杖檢點民訴的從西布倫下來.
以撒加的首領與底窩拉同來.
以撒加怎樣巴拉也怎樣.
眾人都跟隨巴拉沖下平原.
在樓邊的溪水婆.
有心中定大志的.
你為何坐在羊圈內.
聽群中吹笛的聲音呢?.
在樓邊的溪水婆.
有心中設大樓的.
激烈人安居在約旦河外.
旦人為何等在船上.
阿切人在海口靜坐.
在港口安居.
西布倫人是拼命敢死的.
那佛塔利人在田野的高處也是如此.
這裡讓我們看到.
原來有底窩拉家鄉的二百年.

$^{921}$還有辯雅敏,嗎哪西,瓦吉.
西布倫,以撒加,那佛塔利人.
就作了忠心勇敢的回應.
他們參與一起的戰爭.
但是樓邊激烈加德,旦.
你記住旦是一個支派.
是最不便的.
以及阿切這些人停留不動.
他們不理會同胞的需要.
不會作出輔助.
這些不脫鍵的情況.
及後是不斷地出現.
有些支派對於自我的利益的堅持.
好像樓邊只是坐在舒適的陽圈裡.
聽著吹笛的聲音.
激烈因為地理位置較遠.
而且安居在約旦河外.
而不加以援手.
而旦,阿切就在海口和船上.
他們顯然是為了商業的利益.
而不願意挺身.
為自己的同胞,手足而戰.
當時比利基人是商業上很厲害的人.
所以和他一齊合作的話.
或者和他一齊做生意是有好處的.
而且他更加是船手.
所以他們安定下來後.
為了自己的利益盤算.
而不會像進入迦南前後那樣團結.
從第五章的第八節裡.
我們見到當時有四萬人可以來打仗.
二十章的第十七節.
講到他們可以除了有四萬人拿刀之外.
還有四十萬人.
多厲害.
但是真真正正跟隨巴拉的只有一萬人.
所以看到底波拉大聲疾呼.
自願參戰的真的不是很多.
所以他如耶穌所說.
被召的人多.

$^{961}$但是被選上的人少.
但是肯回應呼召的人委實不多.
但是當然有時候多是否真的好呢?.
有時候未必的.
下次我們再看參選.
就會知道多少勝多了.
所以作為領袖.
我們需要作出一個很正確的判斷.
有多少人幫忙才是好呢?.
我們看看戰爭的過程是怎樣.
這段經文就是對戰爭的描繪.
在上一堂裡我們已經查過了.
所以在這裡只是簡略的作為解釋.
我們看看士詩第五章第十九節.
君王都來爭戰.
那時迦南諸王在.
米吉多水旁的他立爭戰.
卻未得擄掠銀錢.
其實這段經文我們知道.
其實不是單單是迦南王耶賓人來攻擊以色列人.
他應該邀請了很多同盟一齊來作戰.
他們以為一定很容易就能摧毀以色列人.
就好像俄羅斯攻擊烏克蘭.
他們攻擊之前已經預了兩三個星期就搞定了.
預備開慶祝會.
誰知道現在離足深陷.
他們竟然沒有去征服以色列人.
更加沒有辦法從以色列人當中.
掠奪任何財富.
因為神介入.
天上地下都一齊向他們宣戰.
到二十一節.
西蜀西蜀從同天上爭戰.
從其軌道攻擊西西拉.
寄信古河把敵人沖沒.
我的靈牙應當努力前行.
這給我們一個很大的諷刺.
迦南人向巴黎敬拜.
求巴黎賜下雨水.
使他們的農作物可以有好的收成.

$^{1001}$當然適量的雨水.
使土地肥沃和農產品會有豐收.
但太多的雨水.
反而令這些農作物會淹死.
況且這一次的雨水.
就成為那九百兩戰車的大敵.
上次我們已經提及.
寄信河一直都是很少水.
好像一條小溪流.
但西西拉是一個很有經驗的猛將.
如果他看見河水高漲的時候.
他絕對不會坐上這些戰車.
指揮軍隊前進.
所以這次是突然之間傾盆大雨.
使寄信河的水位急速地上漲.
使西西拉的軍隊陷入困境.
因為河水猛漲更加令他們混亂.
正如他們的先祖過紅海一樣.
神都會徹底消滅埃及的戰車.
《聖經》14章23到28節.
不過我們不讀了.
講到神怎樣使得埃及軍兵在紅海被淹沒.
藉著這件事讓他們重新認識.
原來耶和華神就是掌管自然.
創天造地的神.
也證明了迦南人所拜的神是假的.
雖然是神幫助他們.
但他們仍然要努力.
所以底波拉特別的尼達強調這句話.
我的靈要努力前行.
就像昔日神血滅約書下的《剛強壯膽》一樣.
第一章第六節約書下記.
你當剛強壯膽.
因為你必使這百姓承受那地為業.
就是我向他們列祖起誓.
應去賜給他們的地.
以色列人的奮戰並不是在於戰場.
那是在淤泥之中去追殺敵軍.
因為神為以色列人展開了勝利的前奏.
祂的指紋要一起來完成這個勝利的樂章.

$^{1041}$接著有句話很特別.
第五章二十三節.
耶和華的使者說.
應當就助米羅斯大大就助其中的居民.
因為他們不耐幫助耶和華.
不耐幫助耶和華攻擊勇士.
這裡出現了一個我們不知道的名稱.
就是米羅斯.
格羅斯可能我們比較熟悉.
這節究竟是講人名還是地名呢?.
沒有一個肯定的答案.
但這一節經文就指出.
不同神同一陣線的就是與神為敵.
已經直接引出下一節.
來讚賞雅逸的經文.
但我們會很奇怪他的描述.
因為他們不願意幫助耶和華.
因為他們不來幫助耶和華.
神是否需要我們來幫助呢?.
當然絕對不需要.
但神卻渴望他的子民會來幫助他.
我想問.
做媽媽的.
你需要你的子女來幫忙拿碗拿碟去開飯嗎?.
你連飯四sung 一湯都煮出來了.
難道你自己不可以拿碗拿碟嗎?.
但如果你的兒女過來幫你放碗放碟.
你的心情會怎樣呢?.
你一定會很開心.
所以神也很渴望他的子民來參與他的工作.
站在他那一邊.
我們神是沒有中立的立場的.
不是與神為友.
就是與神為敵.
這就是主耶穌所說的.
《馬達福音》十二章第三十節.
「不與我相合的,就是敵我的;.
不與我受罪的,就是分散的」.
所以主耶穌也是這麼說的.
如果不是與我相合就是與我為敵.

$^{1081}$接著他們就會來讚賞雅逸.
《士師記》五章二十四到二十七節.
我們就不讀出來.
大家看著這個部分.
怎樣去描繪他怎樣殺死西西拉耶.
雅逸應該站在加拿大那邊.
但他卻是幫助以色列人.
這就很特別了.
所以我們看到原來生命的定論.
並不在乎種族或國家的起源.
而是在乎那個人的行動.
或許你在某個國家的人民.
你的血緣是這樣.
但如果國家做錯事.
去攻擊別人的時候.
你未必要站在這個國家的立場.
俄羅斯人就是這樣.
有些俄羅斯軍隊不想打仗.
後來他們就去投奔對方.
所以你留意到.
這個行動會影響你一生.
蒙福或者受咒詛.
雅逸和拉合有很多相似的地方.
雖然拉合對以色列的神認識不深.
但他憑著信心的行動得到拯救.
也因為成為信心偉人系列之一.
更成為耶穌的祖先.
雅逸也是.
我們看到雅逸全面掌管了局面.
而西西拉一步一步的慢向死亡.
正因為他信心的行動.
使他比眾婦人多得福氣.
多蒙恩澤.
這段就看到另外一個母親了.
我們一起來看.
雖然時間過了一點.
但我們照樣看28到31節好嗎?.
從聖靈中呼叫說.
他的戰車為何耽然不來呢?.
他的車輪為何走得慢呢?.

$^{1121}$聰明的宮女安慰他.
他也自言自語的說.
他們莫非得財而婚.
每人得了一兩個女子.
西西拉得了采依為勞物.
得奏花的采依為略物.
這采依兩面受花.
乃是披在被擄之人頸上的.
耶和華願你的仇敵都這樣滅亡.
願愛你的人與日頭.
出現光輝烈烈.
這樣過中太平四十年.
《底波拉之歌》最後一段.
就帶出西西拉的母親的心態.
這首詩歌是以女性角度出發的.
在這首詩歌裡面.
女性只是西西拉的母親.
是一個無名的婦女.
她為人怎樣我們不知曉.
但這段經文就帶出了.
母親是怎樣對自己的兒子的愛.
西西拉的母親似乎並不知道.
兒子死亡的消息.
我們只是看到一位慈母.
在窗邊盼望自己的兒子得勝回歸.
好像浪子的父親.
一直都盼望浪子回歸一樣.
透過她母親對兒子的期望.
亦都看出.
西西拉是一個很殘酷無比的人.
她以掠奪大量的婦女財物聞名.
而她都可以分享兒子的戰利品.
當然她的士兵更加可以.
分享這些戰利品作為獎賞.
所以她們繼續忠心地去作戰.
這首歌特別以兩個女人作為對照.
瓦易和西西拉的母親作為對照.
如果母親只是專顧自己的兒子.
為了養成熱心的戰士.
而得到財富.

$^{1161}$而忘記了將他帶到神的面前.
是不得昌盛的.
反之一個竭盡所能服侍神的女人.
是一生蒙福.
最後以讚美作為詩歌的結束.
這個結束的訊息再一次提醒以色列人.
他們是屬神的子民.
跟隨神的一定會得到福奇.
而不是受了禍患.
最後我們知道.
當人願意以神所賜予的力量.
回歸到呼召的時候.
那個人就能夠擔當神所交給他的責任.
而且神也會定這個人.
必定和這個人同在.
成為他的幫助.
而那些以各種理由去推搪.
而不願意回應神呼召的人.
永遠沒有辦法來與神的事工有份.
我們必定要成為一個神的召命參與者.
而不是成為一個旁觀者.
或者遵行祂的旨意.
是一條艱辛而漫長的路.
但我們仍然靠著神給我們的力量.
要努力前行.
我們來到低頭祈禱吧.
天王多謝你讓我們能夠去研讀你的話語.
求神你親自來眷顧.
幫助我們.
讓我們知道原來神你所揀選的.
不是在乎性別.
是在乎我們有沒有心.
願意我們能夠懷著一個謙虛.
謙卑.
但是有信心依靠你.
以致我們能夠和神你配合.
願神你透過你的話語.
神我們的幫助提醒.
聽我們在面前的祈禱交托.
拉奉主德性的名.

$^{1201}$阿們.
\newpage



\section{使徒行傳 2:1-47}
\label{sec:032200}
\textbf{五旬節:教會的誕生}~~~~鄭兆奇牧師
\newline
\newline
連結: \href{http://www.church.com.hk/acms/content.asp?site=cdc\&op=show\&type=product\&code=032200\&layout=sermon}{\texttt{032200}} ~~~~ 語音日期: 2026-05-24
\newline
\newline
\hyperref[sec:032199]{\small{< < < PREV SERMON < < <}}
~
\hyperref[sec:index_chronic]{\small{[返順時目]}}
~
\hyperref[sec:032201]{\small{> > > NEXT SERMON > > >}}
\newline
\newline
使徒行傳 2:1-47
\newline
\begin{longtable}{cl}
\hline
\hline
章節 & 經文 (和合本修訂版)\\
\hline
2:1 & \begin{tabularx}{0.7\textwidth}{X} 五旬節那日到了，他們全都聚集在一起。 \end{tabularx} \\ \\ \relax
2:2 & \begin{tabularx}{0.7\textwidth}{X} 忽然，有響聲從天上下來，好像一陣大風吹過，充滿了他們所坐的整座屋子； \end{tabularx} \\ \\ \relax
2:3 & \begin{tabularx}{0.7\textwidth}{X} 又有舌頭如火焰向他們顯現，分開落在他們每個人身上。 \end{tabularx} \\ \\ \relax
2:4 & \begin{tabularx}{0.7\textwidth}{X} 他們都被聖靈充滿，就按著聖靈所賜的口才說起別國的話來。 \end{tabularx} \\ \\ \relax
2:5 & \begin{tabularx}{0.7\textwidth}{X} 那時，有從天下各國來的虔誠的猶太人，住在耶路撒冷。 \end{tabularx} \\ \\ \relax
2:6 & \begin{tabularx}{0.7\textwidth}{X} 這聲音一響，許多人都來聚集，各人因為聽見門徒用他們各自的鄉談說話，就甚納悶， \end{tabularx} \\ \\ \relax
2:7 & \begin{tabularx}{0.7\textwidth}{X} 都詫異驚奇說：「看哪，這些說話的不都是加利利人嗎？ \end{tabularx} \\ \\ \relax
2:8 & \begin{tabularx}{0.7\textwidth}{X} 我們每個人怎麼聽見他們說我們生來所用的鄉談呢？ \end{tabularx} \\ \\ \relax
2:9 & \begin{tabularx}{0.7\textwidth}{X} 我們帕提亞人、瑪代人、以攔人，和住在美索不達米亞、猶太、加帕多家、本都、亞細亞、 \end{tabularx} \\ \\ \relax
2:10 & \begin{tabularx}{0.7\textwidth}{X} 弗呂家、旁非利亞、埃及的人，並靠近古利奈的利比亞一帶地方的人，僑居的羅馬人， \end{tabularx} \\ \\ \relax
2:11 & \begin{tabularx}{0.7\textwidth}{X} 包括猶太人和皈依猶太教的人， 克里特人和阿拉伯人，都聽見他們用我們的鄉談講論神的大作為。」 \end{tabularx} \\ \\ \relax
2:12 & \begin{tabularx}{0.7\textwidth}{X} 眾人就都驚奇困惑，彼此說：「這是甚麼意思呢？」 \end{tabularx} \\ \\ \relax
2:13 & \begin{tabularx}{0.7\textwidth}{X} 還有人譏誚，說：「他們是灌滿了新酒吧！」 \end{tabularx} \\ \\ \relax
2:14 & \begin{tabularx}{0.7\textwidth}{X} 彼得和十一個使徒站起來，他就高聲向眾人說：「猶太人和所有住在耶路撒冷的人哪，這件事你們要知道，要側耳聽我的話。 \end{tabularx} \\ \\ \relax
2:15 & \begin{tabularx}{0.7\textwidth}{X} 這些人並不像你們所想的喝醉了，因為現在才早晨九點鐘。 \end{tabularx} \\ \\ \relax
2:16 & \begin{tabularx}{0.7\textwidth}{X} 這正是藉著先知約珥所說的： \end{tabularx} \\ \\ \relax
2:17 & \begin{tabularx}{0.7\textwidth}{X} 『神說：在末後的日子，我要將我的靈澆灌凡血肉之軀的。你們的兒女要說預言；你們的少年要見異象；你們的老人要做異夢。 \end{tabularx} \\ \\ \relax
2:18 & \begin{tabularx}{0.7\textwidth}{X} 在那些日子，我要把我的靈澆灌，甚至給我的僕人和婢女，他們要說預言。 \end{tabularx} \\ \\ \relax
2:19 & \begin{tabularx}{0.7\textwidth}{X} 在天上，我要顯出奇事，在地下，我要顯出神蹟，有血，有火，有煙霧。 \end{tabularx} \\ \\ \relax
2:20 & \begin{tabularx}{0.7\textwidth}{X} 太陽要變為黑暗，月亮要變為血，這都在主大而光榮的日子未到以前。 \end{tabularx} \\ \\ \relax
2:21 & \begin{tabularx}{0.7\textwidth}{X} 那時，凡求告主名的都必得救。』 \end{tabularx} \\ \\ \relax
2:22 & \begin{tabularx}{0.7\textwidth}{X} 「以色列人哪，你們要聽我這些話：拿撒勒人耶穌就是神以異能、奇事、神蹟向你們證明出來的人，這些事是神藉著他在你們中間施行，正如你們自己知道的。 \end{tabularx} \\ \\ \relax
2:23 & \begin{tabularx}{0.7\textwidth}{X} 他既按著神確定的旨意和預知被交與人，你們就藉著不法之人的手把他釘在十字架上，殺了。 \end{tabularx} \\ \\ \relax
2:24 & \begin{tabularx}{0.7\textwidth}{X} 神卻將死的痛苦解除，使他復活了，因為他原不能被死拘禁。 \end{tabularx} \\ \\ \relax
2:25 & \begin{tabularx}{0.7\textwidth}{X} 大衛指著他說：『我看見主常在我眼前，他在我右邊，使我不至於動搖。 \end{tabularx} \\ \\ \relax
2:26 & \begin{tabularx}{0.7\textwidth}{X} 所以我心裡歡喜，我的舌頭快樂，而且我的肉身要安居在指望中。 \end{tabularx} \\ \\ \relax
2:27 & \begin{tabularx}{0.7\textwidth}{X} 因你必不將我的靈魂撇在陰間，也不讓你的聖者見朽壞。 \end{tabularx} \\ \\ \relax
2:28 & \begin{tabularx}{0.7\textwidth}{X} 你已將生命的道路指示我，必使我在你面前充滿快樂。』 \end{tabularx} \\ \\ \relax
2:29 & \begin{tabularx}{0.7\textwidth}{X} 「諸位弟兄，先祖大衛的事，我可以坦然地對你們說：他死了，也埋葬了，而且他的墳墓直到今日還在我們這裡。 \end{tabularx} \\ \\ \relax
2:30 & \begin{tabularx}{0.7\textwidth}{X} 既然大衛是先知，他知道神曾向他起誓，要從他的後裔中立一位坐在他的寶座上。 \end{tabularx} \\ \\ \relax
2:31 & \begin{tabularx}{0.7\textwidth}{X} 他預先看見了，就講論基督的復活，說：『他不被撇在陰間；他的肉身也不見朽壞。』 \end{tabularx} \\ \\ \relax
2:32 & \begin{tabularx}{0.7\textwidth}{X} 這耶穌，神已經使他復活了，我們都是這事的見證人。 \end{tabularx} \\ \\ \relax
2:33 & \begin{tabularx}{0.7\textwidth}{X} 他既被高舉在神的右邊，又從父受了所應許的聖靈，就把你們所看見所聽見的，澆灌下來。 \end{tabularx} \\ \\ \relax
2:34 & \begin{tabularx}{0.7\textwidth}{X} 大衛並沒有升到天上，但他自己說：『主對我主說：你坐在我的右邊， \end{tabularx} \\ \\ \relax
2:35 & \begin{tabularx}{0.7\textwidth}{X} 等我使你的仇敵作你的腳凳。』 \end{tabularx} \\ \\ \relax
2:36 & \begin{tabularx}{0.7\textwidth}{X} 故此，以色列全家當確實知道，你們釘在十字架上的這位耶穌，神已經立他為主，為基督了。」 \end{tabularx} \\ \\ \relax
2:37 & \begin{tabularx}{0.7\textwidth}{X} 眾人聽見這話，覺得扎心，就對彼得和其餘的使徒說：「諸位弟兄，我們該怎樣做呢？」 \end{tabularx} \\ \\ \relax
2:38 & \begin{tabularx}{0.7\textwidth}{X} 彼得對他們說：「你們各人要悔改，奉耶穌基督的名受洗，使你們的罪得赦免，就會領受所賜的聖靈。 \end{tabularx} \\ \\ \relax
2:39 & \begin{tabularx}{0.7\textwidth}{X} 因為這應許是給你們和你們的兒女，並一切在遠方的人，就是給所有主—我們的神所召來的人。」 \end{tabularx} \\ \\ \relax
2:40 & \begin{tabularx}{0.7\textwidth}{X} 彼得還用更多別的話作見證，勸勉他們說：「你們當救自己脫離這彎曲的世代。」 \end{tabularx} \\ \\ \relax
2:41 & \begin{tabularx}{0.7\textwidth}{X} 於是領受他話的人，都受了洗；那一天，門徒約添了三千人。 \end{tabularx} \\ \\ \relax
2:42 & \begin{tabularx}{0.7\textwidth}{X} 他們都專注於使徒的教導和彼此的團契，擘餅和祈禱。 \end{tabularx} \\ \\ \relax
2:43 & \begin{tabularx}{0.7\textwidth}{X} 眾人都心存敬畏；使徒們又行了許多奇事神蹟。 \end{tabularx} \\ \\ \relax
2:44 & \begin{tabularx}{0.7\textwidth}{X} 信的人都聚在一處，凡物公用， \end{tabularx} \\ \\ \relax
2:45 & \begin{tabularx}{0.7\textwidth}{X} 又賣了田產和家業，照每一個人所需要的分給他們。 \end{tabularx} \\ \\ \relax
2:46 & \begin{tabularx}{0.7\textwidth}{X} 他們天天同心合意恆切地在聖殿裡敬拜，且在家中擘餅，存著歡喜坦誠的心用飯， \end{tabularx} \\ \\ \relax
2:47 & \begin{tabularx}{0.7\textwidth}{X} 讚美神，得全體百姓的喜愛。主將得救的人天天加給他們。 \end{tabularx} \\ \\
[1ex]
\hline
\hline
\end{longtable}
$^{1}$OK,今天就都很快了,五月了,是不是?.
這個月已經五月了,有沒有覺得很快?.
過去那五個月你做了什麼呢?.
數一數,1,2,3,4,5,都不知道,接著就五月了,是不是?.
五月的時候,今天想特別跟大家去講一個題目,.
這個題目叫做五旬節,有沒有人聽過五旬節?.
有聽過,是不是?有沒有人在五旬節教會的以前?.
有,OK,那就對了,這些是五月份,.
五月份就是五旬節,當然不是五月份才回教會,是不是?.
五旬節在聖經裡面其實有一個特別的意思,這個節期,.
當然它和Mother's Day有時候在現實上競爭不了,是不是?.
但是五旬節在聖經裡面的地位是不可以忽視的,.
所以都跟大家去講一講,除了復活節,我們是有五旬節,.
五旬節可以說是教會的一個誕生的日子,.
英文叫做Pentecost,Pentecost,你們應該聽過這個字,是教會誕生的一天..
如果你記得耶穌基督復活之後,祂不是馬上升天的,是不是?.
祂復活之後,祂怎樣?.
耶穌基督復活之後,祂和門徒顯現,.
在這個《使徒行傳》裡面告訴我們,祂和門徒大概顯現了多少天?.
有沒有人記得?.
是的,是40天的,在那裡寫,.
那就是說,祂一個多月復活之後,一個多月在那裡,.
在這個40天裡面,祂有好幾句話是特別在聖經裡面記載的,.
當然了,祂有講很多東西,因為在那40天裡面,.
祂就和祂的門徒去談論天國的事,在聖經裡面說..
如果你記得耶穌基督在這40天裡面,.
首先祂就說過,祂就說大使命,.
祂說所有人,你們要去,使萬民作我的門徒..
因為上帝很願意每一個人,.
包括我們今天在座的每一位,都成為祂的門徒..
而作為祂的門徒,祂這裡說到要奉聖父,聖子,聖靈的名,.
給他們施洗,施浸,當然了,這個也是很應該的..
作為耶穌的門徒,我們順而受洗,.
也都將耶穌基督所教訓的一切,我們一起都去教訓其他的人..
所以我們彼此有查經,有祈禱,OK?.
耶穌教導我們很多東西..
這個是耶穌基督臨升天之前,其中一個大使命,.
願意我們每一個人都能夠得著的使命,.
而這個使命是一個不平凡的使命,雖然是幾句說話..
如果你作為門徒,你聽到耶穌基督講這個使命的時候,.

$^{41}$你會馬上怎樣呢?.
可能有幾種想法,第一種想法就是,.
都不知道怎樣做的,即是怎樣呢?.
第二種想法就是可能很積極,.
我要馬上巴不得將所有人都成為他的門徒,.
抽去全福音,我不知道你是哪一種..
但是其實在這四十天裡面,耶穌基督除了講,.
說你們要遵行大使命之餘,.
他也都講了幾句說話,.
他說大家不要急,我們一起讀,好嗎?.
他說他們聚集的時候,有一個門徒問耶穌,.
他說主啊,你要在這個時候使以色列復國嗎?.
即是你留在這裡四十多天,那你是不是現在要復國呢?.
然後耶穌又說一句,.
他說:父憑著自己的權柄所定的時期,時間或日期,.
你們不必知道,.
但聖靈降臨在你們身上,.
你們就必得著能力,.
並要在耶路撒冷,猶太傳地和撒瑪利亞直到地極,.
作祂的見證..
耶穌基督除了頒佈大使命在這四十天裡面,.
臨升天之前,祂還說,.
大家雖然有著這個使命,.
但是你們要怎樣?.
你們要等一等,等什麼呢?.
祂說要等聖靈的降臨,等聖靈的降臨..
在這個等候的期間,.
當然我們要知道他們不是就這樣百無聊賴,.
或者不知道他們做了什麼,.
上網購物等等..
他們那段時間是怎樣的呢?.
耶穌叫他們等,他們就等,.
但是等的時候,他們是行切的祈禱,.
因為他們都想得著這一份所應許的能力,.
聖靈的降臨..
所以,聖靈的降臨和五旬節有什麼關係呢?.
原來,聖靈的降臨,教會的誕生,.
就是在五旬節的那一天..
其實五旬節是一個什麼節日呢?.
這個節日不是在新約突然間彈出來的..

$^{81}$這個節日在舊約裡面,.
這個節日在舊約聖經其實叫做七七節..
在舊約的背景,大家都知道了,.
舊約有逾越節的,是吧?.
你知道有逾越節,.
逾越節是什麼事?.
大家知道什麼事嗎?逾越節..
逾越節就是耶穌,不是耶穌,.
逾越節就是神拯救以色列人,.
叫他們在門上吐血,.
然後他們就可以避過最後的一災,去拯救..
所以代表逾越節pass over,.
是代表神拯救他們的子民,pass over,.
就是奪命的天使pass over了這扇門,.
這樣我們就pass over,就像我們打啤牌一樣,pass..
逾越節在舊約裡面,大家都很清楚,是吧?.
但其實過了五十天之後,.
他們真的出了埃及,過了紅海,.
五十天之後,就是上帝叫他們在五旬節裡面去慶祝..
就好像剛才那幅圖畫一樣,.
其實五旬節是一個豐收的節日,大豐收..
在那天新約的時候,很多猶太人會去耶路撒冷朝聖,.
就是這一天,第五十天..
所以五旬節其實也可以說是第五十天,.
即是復活節之後的第五十天..
這兩個節日,我知道了,pass over和pentecost,.
逾越節和五旬節,那怎麼辦?.
這兩個節日在舊約裡面有什麼神學意義,.
以致在新約裡面都能夠實現出來呢?.
我們就說在舊約,逾越節就代表神的拯救他們的子民,.
五旬節就代表神要建立他們的子民,供應給他們..
而其實到了新約的時候,這兩個節日就不得了了..
逾越節,記不記得發生什麼事?.
在新約的時候,耶穌基督的受死和復活就在逾越節那段期間,.
亦都代表pass over,是神的拯救,耶穌是神的告約..
而在五旬節,舊約就說是一個豐收神供應的日子,.
而在五旬節的新約裡面,就是聖靈降臨在教會裡面,.
以致教會能夠在馬可樓上所有門徒領受到,.
然後他們就建立教會,.
以致他們可以去遵行那個大使命,行不行?.

$^{121}$由耶穌基督復活,到耶穌基督升天,.
到祂升天之後,聖靈的降臨..
這些一切的事,我們看到上帝的作為是遠超我們的想像,.
祂的計劃也是非常震撼當時的每一個人..
我們在士堂傳的第二章一至四節,我們先讀那段經文,好嗎?.
以致我們知道五旬節的時候在新約發生了什麼事,OK?.
我們一起讀吧.五旬節到了,門徒都聚集在一處,忽然充滿了他們坐在裡面的,.
又有火焰般的勢頭顯現出來,他們就被聖靈充滿,.
就按著聖靈所賜的口才,說起別國的願..
OK,好,這個就是士堂傳的第二章一至四節,短短四節的經文,.
就說出了在五旬節那天發生了什麼事..
五旬節那天發生了一件很奇特,很不尋常的事..
什麼事?在毫無預警的情況下,空中突然間傳來一股很大的聲音,.
而這個聲音是狂風的聲音,沒有人知道它從何而來,.
但這個聲音充滿了當時的馬可樓..
而士堂傳的第三章,剛才你也有讀到,.
那天不單止有一個大聲音,好像風一樣,也有火,.
這個火好像什麼?好像勢頭一樣,落到每一個人身上當時,.
在那一刻,當時每一個信徒就被聖靈充滿,.
而當他們每一個人被聖靈充滿之後,他們的特異功能是什麼呢?.
不是特異功能,就是突然間被聖靈充滿了他們做了一件事,.
那件事是什麼呢?是說別國的語言,.
這個不是方言,這個是別國的語言..
這三件事,在五旬節來說是一件奇特的事發生,.
我記得有一次我和弟兄姊妹去祈禱的時候,.
在一個房間裡祈禱,然後有東西在震動,.
我有震動,大家知道是什麼事嗎?.
就是地震,為什麼早不地震,晚不地震,.
為什麼當我們祈禱的時候地震呢?.
我們一不祈禱,它就不震了,.
我們說不震了,我們繼續祈禱,.
我們一祈禱,又震了..
有時候在五旬節他們在祈禱,聚集在一起,.
就見到這些異象,三個事情,.
三個事情其實可以用三個向度,三個問題,.
第一個問題是他們聽到什麼?.
第二個問題是他們見到什麼?.
第三個問題是他們說了什麼?.
在那天,他們聽到什麼?.
他們看到什麼?.

$^{161}$他們說了什麼?.
首先第一個,他們聽到什麼?.
他們聽到一個很大聲的風一樣的聲音,.
我們就問為什麼這個聲音,.
聲音有很多種,.
為什麼要像風一樣呢?.
像風一樣的聲音是怎樣的呢?.
呼呼聲,就是那些風聲..
其實在舊約聖經當中,.
如果你是記得,.
上帝是怎樣跟那些舊約的人說話的呢?.
譬如約伯,他是在旋風中跟他說話的,.
但如果是旋風,我真的不會扮,.
都是呼呼聲..
而在新約裡面,你又會記得,.
當耶穌和尼哥底姆講論到耶穌基督,.
即是上帝的愛,.
他講到聖靈,.
他就說聖靈就像風一樣,.
你們聽見這個風聲,.
但你不知道它從哪裡來,.
聖靈就是這樣,.
它不是不在這裡,.
但你總是不知道,.
所以其實風這個代表,.
無論在舊約還是新約,.
都會看到這個大聲像風一樣,.
其實就是代表聖靈的那種來臨..
而無獨有偶,.
在舊約的希伯來文和新約的希臘文,.
風這個字就是靈的意思,.
這是文字上的解說..
所以在這裡看到,.
當我們看到五神節他們聽到像風一樣的聲音的時候,.
我們明白是為什麼..
剛才說他們見到什麼?.
除了聽到有風的聲音,.
他們見到的是有火一樣,.
不是,是舌頭一樣的,.
即是像火一樣的,.

$^{201}$即是舌頭般的火..
我們又問,為何會是用火呢?.
如果你熟悉聖經,.
你會知道火在聖經當中,.
經常都是象徵神的同在..
摩西在《創埃及記》第三章被呼召的時候,.
他見到什麼?.
他見到燒著的荊棘,.
即是有火,.
當然這個燒都燒不著,.
即是不會燒爛的荊棘,.
只是看到有火而已..
如果你有記得,.
當以色列人離開埃及,.
去到曠野的時候,.
這四十年,.
夜晚是什麼帶領著他們的呢?.
就是火柱,.
雲柱就是早上,.
夜晚就是火柱..
所以對於猶太人來說,.
當時他們看到這個火,.
好像舌頭般的,.
去到他們身上的時候,.
其實他們很熟悉的,.
他們很知道的,.
風又來了,火又來了,.
即是代表什麼呢?.
即是耶穌基督所說的聖靈,.
就是來了,真是來了..
這件事不得了,.
當它來到他們身上的時候,.
他們就可以去說別的語言,.
剛才聖經說,.
你覺得很奇怪,.
為什麼要說別的語言呢?.
當然當時的國家不是那麼多,.
我真的想問一下你,.
你懂得說多少語言呢,大家?.
有沒有人懂得說多過五種的?.

$^{241}$很多,是吧?.
有沒有人?.
有沒有人懂得說多過三種呢?.
有沒有?.
一種,好嗎?.
不說話的,大家?.
原來大家不說話的,OK..
其實不要很流利地說,.
譬如說你好,hello這樣,.
其實我想在座當中,.
你們能夠說十個國家語言都不奇怪,.
hello而已,.
可以嗎?.
不可以?.
試一試,好嗎?.
廣東話,不用想了,你好..
國語可以嗎?難不倒你,你好..
上海話,.
隴河,隴河,有些人懂得了..
很簡單,.
這個韓文,.
第四種了,你們已經懂得四種了..
日文,こんにちは,真是可愛,是吧?.
五種了,.
泰文,.
sawadee,.
sawadee krab是跟女性說的,.
sawadee krab是跟男性說的,OK?.
我去過泰國,.
但是我沒有說過這些,OK..
sawadee,OK..
你已經記住了,.
還有1,2,3,4,5,6種了,OK?.
Hindi,即是印度,.
你們到處都Namaste,Namaste,是吧?.
即是你見到你到處那些,是吧?.
隔壁鄰舍,你接兒子,那些同學,.
Namaste,他應該會回應你的..
那Spanish呢?.
是吧,小孩一出來,.

$^{281}$你就問他,Hola,.
這樣就行了,Hola..
你又懂了,你現在多少種了?.
1,2,3,4,5,6,7,8種了,已經..
然後French,.
Bonjour,真是,.
沒錢了 請你去磅水吧.
這樣就行了 見到人就去磅水.
這個可能難一點.
希臘 有沒有人去過希臘.
有 應該是他教過你的.
忘記了 不要緊 很容易記的.
希臘那個叫Yasu.
Yasu 是吧.
其實就像耶穌一樣.
見到希臘的人就像耶穌一樣.
耶穌 耶穌 耶穌 耶穌.
這樣就行了.
希伯來文 Hebrew就很容易.
大家都知道.
Shalom.
Shalom 平安.
你撿起來有多少種.
十幾種語言 大家都懂.
只是說Hello 是吧.
當然了 聖靈降臨.
是不是要聖靈感動大家才會說呢.
當時為何門徒要懂得說這麼多語言呢.
其實在當中我們發覺.
有幾個真理.
上帝希望我們掌握到.
祂不是只想我們去旅遊的時候.
懂得說你好.
首先.
上帝 聖靈是將來自不同背景.
和不同的人.
是文化.
即是Live stage.
怎樣也好 不同的.
聚焦在一起.

$^{321}$即是成為一個群體.
這個是上帝的心意裡面.
怎樣去將這群人結為一呢.
就是用福音的大能.
其實我們Living Grace.
即是我們教會的形成.
都是這樣的.
雖然我們現在講廣東話.
但其實我們整個教會.
都有不同文化 不同背景.
我們有Youth講英文.
有Mandarin講國語.
即是我們有不同文化.
我和你可能是不同背景的.
年齡層我們都有不同的.
我們有Children的小孩.
我們有Senior.
我們有長者.
我們的思想很多時候很不一樣.
但上帝透過聖靈.
繼續建立教會的時候.
祂就將不同的人要放在一起.
這件事我覺得是不容易的事情.
我不知道你最近有沒有看.
我們那個Living Room的短片.
最近拍了一個UFO.
有沒有人看過.
新聞都可能看過.
現在早一輪新聞.
美國的房友有暴露了一些UFO.
說他用F16戰機打低了那些UFO.
我不知道是不是真的.
UFO.
如果你看過那個Living Room.
他第一個問題問什麼.
有沒有外星人的呢.
第二個問題立刻問.
如果有的時候我們歡迎不歡迎他呢.
歡迎不歡迎他呢.
如果有外星人.

$^{361}$他來到這裡Living Bay崇拜.
歡迎不歡迎他呢.
那個Living Room真的搞笑.
當然歡迎他.
但當然他來到的時候.
就嚇到那些人都走了.
當說到UFO的時候.
我們不是在這裡.
這一刻不是想大家討論.
有沒有UFO的存在.
而其實都不需要說到UFO這麼遠.
有時候我們當中.
都真的有些外星人在裡面.
就是和大家不是很合得來.
或者就是.
就是不是同類.
你總是覺得他有些UFO的感覺.
你真的想問他是誰請求你的.
但其實.
這些都是上帝的心意.
他給我們看到聖經的時候.
他是用不同的人.
講不同的語言.
不同的文化.
不同的背景要聚在一起.
目的就是要這樣去榮耀他自己.
因為上帝的愛是普世的.
是每一個人.
不是UFO.
就是無論他是什麼形狀的人.
上帝都很愛.
第二點我們看到就是聖經.
告訴我們教會也是.
始於信徒同心合意.
和耐心等待馴服神的時刻.
去建立教會.
因為你記得在五重節.
剛才聖靈降臨之前.
我們就說其實他們是同心合意去禱告.
這件事情我想是.

$^{401}$每一個真正是一個教會的教會.
他應該是這樣去做.
我們每一個教會.
都是從禱告開始的.
當然我們不是在馬可樓.
我們那次在Louis和Candace家裡.
那個不知道什麼樓.
每次去到都很多東西吃.
但是除了很多東西吃.
我們會一起敬拜.
一起祈禱.
尋求神的旨意之餘.
我很記得有好幾次的時候.
我們一直在祈禱的時候.
弟兄姊妹會哭泣和流淚.
而那些哭泣和流淚不是很哀傷那種.
是一種被聖靈去感動.
被聖靈去醫治那些淚水.
一起祈禱的時候所說的.
都不是出於我們人的說話很多次.
我們會安教會是在禱告當中.
我們繼續被神去醫治.
繼續在禱告當中我們被神合一我們.
所以今天如果你沒有機會來我們的祈禱會.
我歡迎你來.
我也希望你會來.
感受一下我們禱告當中那份力量和甜味.
和聖靈那份同在.
而第三點我看到.
當聖靈給力量那些人能夠說不同的意見的時候.
我們發覺他們不是一些特別厲害的人.
而其實他們是平凡的信徒.
聖靈所賜能力是給平凡人用的.
但是這班平凡人就如我們一樣.
但其實也承載著一個不平凡的使命.
要使萬民作他的門徒.
要奉聖父,聖子,聖靈的名.
給他們施浸絕對不是一個人能做到的事情.
其實我們每一個人都很平凡的.
你認不認同?.

$^{441}$大家覺不覺得你都是一個平凡人?.
當我們見到每一位平凡弟兄姊妹無力當中.
他們又有很多很不平凡的事情.
因為上帝給他們恩賜.
給他們的能力.
給他們的所有都獻上出來.
就變得很不平凡.
我們有個小組,一個非聖職小組.
這個叫Beer.
就是Beer.
就是喝酒的.
Beer.
那個組當然是一班頂.
我那次沒有機會去.
下次記住告訴我,我又怎樣可以去到.
後來我聽他們說,當然很開心.
不是說喝醉了,燙肚子,不是那些.
說到阿星.
我們都知道阿星是星奶茶.
有時有些奶茶,就是阿星的奶茶.
你覺得那些奶茶一點都不平凡.
但那個Beer組,Anders跟我說的時候.
他說原來阿星.
他懂得調酒.
那個Beer組其實不是喝啤酒的,是喝wine的.
那個組應該叫Wine.
但我不知道,我只是聽他這樣說,是不是真的.
可能我不知道,可能我的資訊錯.
有人說他早上是做奶茶店的,賣奶茶.
晚上是做酒店的,是不是.
阿星你晚上是做酒店的,我不知道.
怎樣也好,很多不平凡的東西.
在我們平凡人身上是會見到的.
原因就是聖靈賜能力給一班平凡的人.
今天如果你覺得你是平凡,不要緊.
我們就是平凡,但當我們肯擺上.
整件事就不平凡.
去承接這個不平凡的使命.
我們教會不是靠人的血氣.
不是靠人的能力去建立.

$^{481}$我們知道我們是靠神親自帶領.
一起服侍,一起相愛,一起去建造.
當我們很開心的時候.
其實我們外面是有很多需要.
我上兩個星期四在祈禱會當中.
我都叫弟兄姊妹為我禱告.
為我隔壁有一個鄰舍禱告.
因為那個鄰舍我認識了很久.
我現在住那裡住了七年.
他差不多跟我一起在那裡開始.
他是一個年輕人.
那個星期四晚上我們在祈禱.
那個星期四早上我聽到.
八點多有一個年輕人在叫.
很大聲地叫,很痛苦地叫.
在那裡叫了幾分鐘,很厲害.
當然我們隔壁的鄰舍就會衝過去看看發生什麼事.
下飛機的時候看到他躺在那裡.
就是我們隔壁鄰舍的年輕人.
他躺下了,他說他從上面跳下來.
但是他還是不死.
他說你現在就殺死我.
You kill me, just kill me directly.
他只說了這句話.
我們問他什麼事,他只說了這句話.
當然我們報警諸如此類.
其實他是一個大好青年.
他有幾年去了哈佛讀書.
那間屋租了給別人.
但是都很麻煩.
有時候租屋你都不知道租了什麼給別人.
然後搞到那間屋不知道怎麼辦.
接著他回來.
我說你讀完書來哈佛?.
他說是的.
他在那裡工作,在Tesla工作.
我說不錯啊.
畢業了,拿了MBA.
在Tesla工作,不錯啊.
我想都想不到他那天會在那裡叫得這麼厲害.

$^{521}$其實他問過我去哪個教會.
因為他不相信.
每次見到我都很喜歡跟我談信仰.
可能他知道我是Pastor.
所以我叫弟兄姊妹幫我禱告.
我希望如果他回來.
如果他真的回來,他真的清醒的時候.
我可以minister他.
我可以用神的愛.
用弟兄姊妹禱告的力量.
讓他知道.
神不要放棄他.
我想說的就是我們每間教會.
怎樣才有聖靈的能力.
我們一班平凡的人怎樣被神使用.
其實就是因為.
在這個世界有很多事情.
不平凡的事情在發生.
不尋常的事情在發生.
當聖靈感動了弟兄姊妹當時的門徒.
他們可以說到別的語言之外.
其中有一個使徒叫彼得.
他在那個時候.
他很勇敢地站出來.
去講道.
上星期Pastor生也說過.
當時他在公會面前很勇敢.
但在公會之前他又很勇敢地被聖靈充滿.
去講道.
他怎樣說呢.
他說彼得和十一個使徒站起來.
他高聲對眾人說.
猶太人和所有住在耶路撒冷的人.
你們應當明白這事.
也應該留心我所說的.
這些人並不是照你們所想的喝醉了.
為什麼他們可以說到別國的語言呢.
是不是喝醉了呢.
他說現在只不過是上午九點鐘.
怎會喝醉呢.

$^{561}$他說以色列人.
我們一起讀吧.
以色列人.
請聽聽這幾句話.
正如你們所知道的.
神已經藉著那撒冷人耶穌.
在你們中間施行大事.
大能,奇事,神職.
和你們證明祂是神所立的.
祂叫做神的定子和御資.
被交了出來.
你們就藉不法之徒的手把祂釘死了.
神卻把死的痛苦解除.
使祂復活了.
因為祂不能被死亡.
祂被驅禁了.
彼得這篇講道的核心是什麼呢.
很明顯.
就是耶穌基督.
耶穌基督就是祂的主.
耶穌基督已經在十字架被釘死.
也被你們釘死了.
但是祂已經復活了.
整個訊息完整的中心.
就是耶穌基督的中心.
當時來說.
這是不容易的一篇講章.
我現在講好像很輕鬆.
但當時不是口震震也說不出來.
但我相信彼得沒有口震.
耶穌基督是主的中心訊息.
其實是我們會恩教會的一個中心思想.
當大家要宣講上帝恩典的時候.
你不能離開耶穌基督的恩典和恩惠.
當然我們有個人的見證.
但是否只有男人才可以這樣講這篇道呢.
是否只有彼得才可以呢.
姐妹們 你可不可以呢.
女人可不可以講道呢.
為甚麼呢.

$^{601}$聖父和聖母就來安納了.
其實他們一向都有講道.
也有人問牧師你來安納他們.
他又講道.
但好像聖經說.
不可以的.
在提摩太前書第二章第十二節.
姐妹們讀吧 拜託.
我看得不太清楚.
保羅說我不許女人講道.
也不許她嚇管男人.
只要沉靜.
我看到弟兄們在笑.
不要笑得那麼開心.
不許女人講道.
牧師 你讓聖父和聖母講道.
不適合聖經.
聖經是需要查考的.
有另一處經文.
在哥林多前書.
他這樣說.
女人祈禱或講道的時候.
如果不蒙頭就收肉自己的頭.
我們先不要說他用甚麼蒙頭.
他講道 他們兩個用甚麼蒙頭.
這明明可以說女人祈禱和講道.
即是可以的.
即是有這樣的事.
在哥林多教會就可以.
剛才提摩太前書.
提摩太在哪個教會.
在以弗所教會.
即是以弗所教會不可以.
哥林多教會就可以.
是不是這個意思.
還是怎樣.
是可以還是不行.
如果你看保羅的書信.
你會看到保羅頗尊重女性.
無論是伯羈拉 阿居拉.

$^{641}$即是這對夫婦.
他們都一起教導.
當中有一個執事叫Phoebe.
Phoebe在不在.
即是他讚賞Phoebe等等.
到底是否可以講.
其實我們發覺.
保羅剛才在提摩太前書.
他說我不許女人講道.
他怎樣說.
我不許女人講道.
也不許他們嚇管男人.
只會尋證.
還加了一句.
這句大兄讀算了.
節目真的很傷心.
因為先做的是阿當.
後做的是阿娃.
不過阿當受了引誘.
不是阿當受了引誘.
而是女人受了引誘.
陷入了這個過份.
大家都知道發生什麼事.
講來講去都是那個結果.
怎樣也好.
特別要這樣寫.
然後又說不許講道.
其實我們怎樣查考.
盛景怡的意思是.
如果看回提摩太前書.
很多蛛絲馬跡.
講到當時爾弗所教會.
實在有一群姊妹.
很離譜.
和哥林多教會都很離譜.
如果你看哥林多前後書.
但當時.
他竟然在爾弗所說不許女人講道.
當時的姊妹.
是有些事情.

$^{681}$都看回前後文.
看得很快.
剛才是第二章.
第二章的前面.
因此我願男人沒有憤怒.
沒有爭論.
舉起手到處禱告.
很明顯.
當時當然有男人憤怒.
當然有男人在爭端.
在爭論處處理.
他就不用講.
我願那些男人不要那麼帥.
不要那麼有禮貌.
他不會.
男人就有問題.
接著他說.
照樣我也願女人.
在那時以端正.
閒淑.
自律為裝飾.
當時有些女性.
不少.
用什麼裝飾自己.
捲髮.
金飾珠寶.
周大福.
或名貴衣裳.
手袋.
為裝飾.
當時有一群.
頗有錢.
當時練法.
不容易.
我不知道那些珠寶.
怎樣也好.
當時有些.
應該是以.
善行為裝飾.
但他們這樣裝飾.

$^{721}$自稱為敬畏神.
他們覺得.
不要弄外表.
你看外表.
好像有點不對勁.
但你再看第五章.
同一本書.
又是講異乎所.
他怎樣說.
至於年輕寡婦.
他們的問題.
他們懶惰.
慣了.
挨家走右.
不但懶惰.
還要.
割舌.
好說閒話.
好管閒事.
說不該說的話.
很嚴重.
因此我願年輕寡婦.
快些結婚.
管理家務.
這個很重要.
不要給敵人.
有辱罵的機會.
因為有些人已經跟隨.
殺彈了.
整個提摩太前書.
本來要提醒提摩太.
你們教會有很多異端邪說.
你們要小心.
而誰散播異端邪說.
其實是一群幕後黑手.
是男人.
他們使用了當時.
一群年輕寡婦.
做了很多.
弄死弄死的事.

$^{761}$甚至當時.
有人講道.
搞得很亂.
以致其實.
就算更亂的高林多教會.
他也沒有說不讓女人講道.
但在異忽所當時.
其實是什麼意思.
他說不讓女人講道.
不是一個超越時空的教導.
不是應用在每一個教會.
每一個群體的教導.
所以我第一個問題.
問給不給.
大家都說給.
為什麼呢.
聖經既然說不行.
我們要明白.
如果有人問你的時候.
希望你能夠將聖經.
很清楚地告訴人.
到底發生什麼事.
以致聖靈的能力降臨到每一個人的身上.
當聖靈降臨到我們每一個人的身上.
希望我們能夠結出.
聖靈的什麼.
果子.
你有沒有聖靈的果子.
人愛喜樂和平.
你有沒有人愛喜樂和平.
如果你沒有的時候.
你很難.
你沒有這個果子.
一粒麥子跌在地裡.
要生出很多果子.
人要喜樂和平.
然後怎樣.
會不會有恩慈良善信實.
對神有沒有信實.
對人有沒有恩慈.

$^{801}$有沒有良善.
惡世能得.
在逆境當中.
有沒有溫柔忍耐節制.
人家反對你的時候.
你可不可以溫柔一點.
在很艱難的時候.
你可不可以繼續忍耐.
在你想放縱的時候.
你能夠有節制.
昨天.
我和一班弟兄.
LG Cycling.
我有機會和他們一起去踩單車.
炒車那些.
不要說了.
最後有一段.
是要上斜路.
那段斜路.
我覺得都.
挺勤的.
當然.
我有輛好一點的車.
好一點.
我已經轉了最low的gear.
我由無氣踩到斷氣.
但不要緊.
我和自己說.
我要忍耐.
我要有聖靈果子.
我要忍耐.
但環境很困難.
為什麼呢.
不單止弟兄已經不見了.
在前面.
還在旁邊.
有一個黑人.
他半跑步和走路之間.
他比我快.
比我快.

$^{841}$你怎樣去忍耐和堅持.
我想問弟兄姊妹.
在那時候你只是懷疑.
你懷疑什麼呢.
你懷疑.
當然先懷疑車.
當然是賴的.
懷疑車是不是不行.
個個都拿著一些凹頭.
然後就懷疑.
不關車的事.
是自己的主意不行.
沒有Daniel那麼厲害.
懷疑後.
不要緊繼續去.
我不理那人.
他走路比我快.
應該下車推.
我不想推.
終於見到那班弟兄.
等我多開心.
我們這樣上斜路.
現在是不是下山.
苦盡甘來.
我們這條路叫苦盡甘來.
弟兄說.
是呀.
更加甘.
我還要上.
有時是這樣的環境.
你以為在一個.
很艱難的環境裡面.
你忍耐.
誰知你會見到旁邊的人.
會影響你.
你去到終點.
誰知原來苦盡甘來.
是更加甘.
當然他開玩笑.
當然直衝下去.

$^{881}$我們在信徒生命當中.
也是.
聖靈給予你的能力.
就是在很困難的環境.
你仍然保持友誼.
喜樂和平.
繼續說.
大家敬你.
恩慈良善信實.
溫柔忍耐節制.
一起祈禱.
我們願你今天去更新我們.
充滿我們改變我們.
讓我們每一天.
能夠有聖靈你的充滿.
因為你以聖靈降臨.
以致我們會靈教會.
能夠誕生.
我們深信聖靈你必定繼續帶領.
以致我們能夠活出聖靈的果子.
願以神.
你在我們的生命和教會中.
得罪一切榮耀.
祈禱奉主名求.
\newpage



\section{加拉太書 5:2-6}
\label{sec:032201}
\textbf{靠律法 VS 靠信心 (華)}~~~~葛國光牧師
\newline
\newline
連結: \href{http://www.church.com.hk/acms/content.asp?site=cdc\&op=show\&type=product\&code=032201\&layout=sermon}{\texttt{032201}} ~~~~ 語音日期: 2026-05-24
\newline
\newline
\hyperref[sec:032200]{\small{< < < PREV SERMON < < <}}
~
\hyperref[sec:index_chronic]{\small{[返順時目]}}
~
\hyperref[sec:032202]{\small{> > > NEXT SERMON > > >}}
\newline
\newline
加拉太書 5:2-6
\newline
\begin{longtable}{cl}
\hline
\hline
章節 & 經文 (和合本修訂版)\\
\hline
5:2 & \begin{tabularx}{0.7\textwidth}{X} 我—保羅告訴你們，你們若受割禮，基督就對你們無益了。 \end{tabularx} \\ \\ \relax
5:3 & \begin{tabularx}{0.7\textwidth}{X} 我再指著凡受割禮的人確實地說，他有義務遵行全部的律法。 \end{tabularx} \\ \\ \relax
5:4 & \begin{tabularx}{0.7\textwidth}{X} 你們這要靠律法稱義的是與基督隔絕，從恩典中墜落了。 \end{tabularx} \\ \\ \relax
5:5 & \begin{tabularx}{0.7\textwidth}{X} 至於我們，我們是靠著聖靈，憑著信心，等候所盼望的義。 \end{tabularx} \\ \\ \relax
5:6 & \begin{tabularx}{0.7\textwidth}{X} 因為在基督耶穌裡，受割禮不受割禮都沒有功效，惟獨使人發出仁愛的信心才有功效。 \end{tabularx} \\ \\
[1ex]
\hline
\hline
\end{longtable}
$^{1}$我們加勒泰書中間隔了一段時間.
可能大家都已經忘了.
還記得嗎.
記得什麼.
羅馬書,加勒泰書是保羅很重要的兩卷書信.
馬丁·路德稱為羅馬書跟加勒泰書.
可以說是他因性稱義改教的非常重要的兩本書.
那麼加勒泰書比較深奧一點.
雖然只有六章.
但是加勒泰書保羅的書信有個特色.
就是他前面講教義.
後面就一定講生活.
所以一定把教義講清楚了.
然後一定把這個教義要應用在生活當中.
這也是我們講到或者我們聽到我們讀經一個很重要原則.
我們把神的話弄清楚.
但是不是停留在那個地方.
我們一定要讓神的話能夠實行出來應用出來.
跟你旁邊說好好聽到.
保羅在羅馬書第五章.
今天我們要看到羅馬書第五章第一.
第一節到十六節.
可以說是整個他教義的一個完結篇.
也可以說是他講到因性稱義的這個重要主題的一個結論.
從十七節開始他就要講生活方面的.
保羅還記得我們花了很多時間在講五章第一節是不是.
今天基督既然已經釋放了我們.
我們已經從盧浦的惡就不再被盧浦的惡所挾制了.
所以我們講到說我們今天在基督的裡面我們已經得到了自由.
從哪裡得自由.
從律法的當中得自由.
從這個罪惡的捆綁當中得自由.
所以怎麼樣子我們能夠得自由呢.
在基督的裡面我們可以得自由.
藉著真理可以使我們得自由.
聖靈在哪裡,哪裡就得自由.
所以親愛的幼稚妹,你自由了嗎?.
問你旁邊的人,你自由了嗎?.
你剛纔在敬拜的時候你自由嗎?.
自由嗎?.

$^{41}$如果你真自由你應該盡情的敬拜.
你盡情的讚美.
你如果還在受到過去傳統.
或者一些不正確的這種敬拜觀影響你.
你就很難釋放,是不是?.
當我們真正釋放的時候你就不會懼怕人.
當我們真是在靈裡面敬拜的時候.
我不知道你沒經歷過,我經歷過我知道.
那真是一個喜樂得不得了的一個敬拜.
那真是一個盡情的一個敬拜,Amen?.
那是我們所要的,Amen?.
所以保羅在接著第五章第二節開始.
第五章第二節開始.
他說我告訴你們.
若受割離基督就與你們無異了.
保羅這邊又在提割離.
因為保羅之所以寫加拉太書.
就是因為當時候有從耶路撒冷來的母會.
母會很大.
有母會來的這些猶太的基督徒.
他們雖然信了耶穌.
但是他們還活在猶太教的那種令法的觀念裡面.
所以他們來跟加拉太的信徒說.
他說你們光信耶穌不夠啊.
你們還需要怎麼樣?.
受割離.
他們還告訴他們要受割離.
好了,割離是什麼?.
割離是什麼?.
以色列的男子第八天就要受割離,是不是?.
謝謝,不然我不知道我在跟誰講話.
以色列的男子要受割離.
割離這件事情對以色列人來講.
那是非常重要的一件事情.
割離是律法當中的一個規定.
割離代表著舊約時代神跟人的一個利約.
好了,所以簡單的說割離就代表著律法.
割離就代表著要靠律法來稱義.
所以保羅為這件事情.
從一開始加拉太書的時候他大力的講.

$^{81}$他說今天有人來跟你們傳一個別的福音.
看到沒有?.
這是很嚴肅的事情.
今天如果有人來跟你傳別的福音.
不是我傳給你們的.
不是聖經所告訴你們的這個福音.
那保羅說這是不得了的事情.
所以保羅畫這麼多的篇幅.
五章的聖經在這邊講這件事情.
那保羅在這裡接著講什麼?.
他說,他講到這個.
你如果受割離會帶來什麼結果?.
他講第一件事情.
你如果受割離,基督就與你們無異了.
基督就與你們無異了.
這個改教的.
我們都知道改教的是馬丁路德.
但是還有一位其實大家比較不清楚.
加爾文.
加爾文是長老會.
就是這種改革中的長老會的.
可以說都是尊加爾文為他們的可以說算師父,師祖.
加爾文說,他說加拉太人中間的律法主義者.
想要教他們絕德.
他們可以同時既有耶穌.
還可以保持與上帝的律法關係.
保羅對他們說.
這並不是一個向他們開放的選項.
恩典體系和律法體系不可兼容.
凡是想要得到半個基督的人.
必定全然失去他.
瞭解我意思嗎?.
保羅,所以加爾文在這裡講到的意思就是說.
你不可能說你今天靠著律法.
又要靠著基督.
你一定只能夠靠著基督.
這是唯一的道路.
沒有別的選項.
只有因著相信耶穌,耶穌基督.
我們纔可以得救.

$^{121}$纔可以成癮.
確定嗎?.
非常確定.
要很重要的確定.
上次我給大家看過美國做的調查.
很多的基督徒不清楚這點.
那麼我們的總統川普.
在2025年的時候.
8月19號.
他在跟烏克蘭總統.
扎文斯基還有歐洲領導人會面之後的第二天.
他就跟福斯新聞的一個節目聯繫.
這個節目當中他被訪問.
他在節目當中就表示說.
他說他自己聽說.
在天堂的功德榜上.
他處於圖騰柱的最底層.
什麼叫圖騰?.
英文叫做At the bottom of the totem pole.
這什麼意思?.
其實英文這個是一個用語.
它的意思就是說.
就是在地位最低,最不重要的意思.
所以他說他在天堂的那個功德柱裡面.
是當At the bottom就是最底部的.
因此他的言下之意是什麼呢?.
他的言下之意說.
他說他如果能夠促成烏克蘭跟俄羅斯的和平.
每個禮拜拯救7000人.
他說這將成為他進入天堂的理由.
請問在座的可以嗎?.
可以嗎?.
可以 好.
那麼現在他覺得他自己應該可以進天堂了.
我們來看看他的訪問.
幫我開一下.
這是在2026年的早餐會.
國家早餐會.
我們有個軍隊.
他們都看起來像湯姆克魯伯格.

$^{161}$請問一下美國.
請問阿聯酋.
他們正在談判.
他們不想讓我們進去.
我們有個大艦隊要去阿聯酋.
那是很多假新聞.
因為我上次很開心.
我們有個大艦隊.
有60萬人.
他們說我永遠不會進入天堂.
我永遠不會.
我開始開始開玩笑.
你不能和他們開玩笑.
因為他們寫了一句話.
紐約時報的前線報道說.
特朗普在質疑他人的生命和生命意義.
不 我只是開玩笑.
我認為我應該做.
我不是個完美的選舉人.
但是我做了很多好事.
為完美的人做.
川普宣佈了兩項宗教相關計劃.
一個是邀請美國民眾.
在5月17號.
在華府國家大考評禱告.
另外川普也表示教育部.
將針對公立學校的禱告權.
發佈指導方針.
但他預料這個禱告方針.
將會面臨訴訟.
TVB的新聞.
江輝 美國華府報導.
Ok 好.
好 那麼.
所以.
川普總統可以上天堂嗎.
他進步了.
他從去年.
不 2025年到今年.
他已經覺得他應該可以上天堂了.

$^{201}$請問怎麼樣可以上天堂.
啊.
Ok.
所以做了很多好事.
就可以上天堂嗎.
其實.
不少的基督徒也有這種觀念.
我們華人的基督徒.
受到佛教的背景的影響.
其實我們很多時候也會有這種觀念.
不一定是只是在所謂的上不上天堂的問題.
很多時候我們覺得我們應該要做好.
所以有時候會聽到一些弟兄姊妹說.
他說我覺得我如果再做得更好的話.
我應該可以上天堂.
那什麼叫做做得更好.
那好的標準在哪裡.
你懂我意思嗎.
所以保羅在這裡告訴我們一件事情.
如果你今天還要靠律法來稱義.
你就與基督無異了.
接著保羅講的第二件事情.
他說如果你要行割離.
你就要行整個律法.
他的原話是這樣說.
我在指著凡受割離的人確實的說.
他是怎麼樣.
欠著行全律法的債.
那我們接著來練這個雅各蘇二章十一節.
第四 我們一起來練 好不好.
來 預備 請.
雅各是在講行為 對不對.
羅馬蘇講陰性稱義.
那雅各蘇同時他也告訴我們說.
你不要光是講說你陰性.
但是你的信心要有行為.
所以雅各蘇在這裡.
雅各蘇講的這個律法的時候.
他提了一點.
他說凡遵守全律法的.

$^{241}$你在一條上犯.
你一條上面你沒有遵守.
你一條上面你跌倒.
你就是犯了重條.
這什麼意思啊.
我們有沒有拿罰單的經歷 有沒有.
好比說你今天超數了.
哎呀 今天趕著要來主日聚會.
所以呢 我就Speedy 哎呀 這非常正當的理由.
是為主啊.
好了 結果呢.
警察伯伯在後面 噠噠噠 就抓到你了.
你跟他說 警察伯伯啊.
我現在是要去教會.
而且我是牧師.
我要去講道.
時間來不及了.
我是一個good driver.
我過去三十年駕駛的經歷.
你去查一下DMV.
去看一下 我是good driver.
我大部分的時間.
我都是遵守交通規則的.
我就這麼一次啊.
警察伯伯 饒了我吧.
警察伯伯怎麼說.
我不care 你過去三十年是good driver.
你今天Speedy.
你就是Speedy 懂我意思嗎.
所以保羅意思就是說.
你犯了這一條.
你過去再遵守了多少條.
你都叫做犯罪.
你都叫做犯了律法.
所以保羅是一個非常懂得這個律法的.
所以他用一個辯證的一個方式.
在跟猶太人在講.
那麼我們中國人背景比較多熟悉佛教.
你知道佛教有很多跟這個行為有關係的.
對不對.

$^{281}$去西方極樂世界.
怎麼樣能夠去呢.
其中有一個大部分的佛教的都強調了一件事情.
就是要守戒.
持戒或守戒.
好了.
那麼這個戒呢.
有五戒.
有八齋戒.
有比丘裡戒.
有比丘戒.
含不當當加起來大概兩百多條戒.
但是現在我不是要教你怎麼守戒.
我現在告訴你是說最基本的佛教.
他要守五個戒.
來看看這哪五個戒.
來我們起來練好不好.
來.
不殺身.
還有呢.
不偷盜.
不邪淫.
不妄語.
還有呢.
不飲酒.
好這五戒你通通都遵守的.
舉手.
都沒有.
按照佛教的道理來講.
這五戒是最基本的戒.
也是最重要的戒.
特別前四個戒.
不殺身.
不偷盜.
不邪淫.
不妄語.
這四個戒當中.
你只要有一個戒.
你沒有遵守.
Out.

$^{321}$你沒有辦法上天堂.
你沒有辦法進西方極樂世界.
重不重要.
好了.
那我請問大家.
所以你看到沒有.
這只是兩百多條戒的當中的一個.
這個五個戒.
你一個戒你沒有犯.
你一個戒你不能夠遵守.
你就等於是犯了重戒.
這個意思懂嗎.
不是要等到你所有的戒都遵守了.
你才能夠去西方極樂世界.
而是怎麼樣.
你只要有一個戒你沒有遵守.
你就不及格.
你就沒有辦法進去.
瞭解我意思嗎.
所以保羅在這裡用這樣子一個觀念來說.
如果我們今天還在強調割離.
我們就必須要遵循權律法.
好了.
這是第二點.
第三個.
我們起來練這段經.
第四點來請.
你們這要靠律法稱義的.
是與基督隔絕.
從恩典中墜落了.
你們這要靠律法稱義的.
是與基督隔絕.
更厲害了.
前面是與基督怎麼樣.
五義現在是什麼.
與基督怎麼樣.
隔絕.
從恩典中墜落了.
你有沒有去過.
有多少人去過意大利.

$^{361}$去過羅馬.
有沒有去看過.
你去羅馬一定要去一個地方.
叫做聖彼得大教堂.
有沒有.
去過吧.
好.
那麼你知道聖彼得大教堂.
跟馬丁路德改教.
是有一個很密切的關係.
什麼關係呢.
首先我們來看看.
聖彼得大教堂.
可以說是世界上面.
人類歷史上面.
這個造價最昂貴.
花費的錢財最多的其中的建築物之一.
他一共花費了120年來建造.
120年了.
前前後後.
花費的金額呢.
大概是以今天的幣值來算.
因為有一些專家幫他算過.
大概是50億到330億美金.
那弟兄姊妹.
請問這麼多的錢從哪裡來的.
1515年的時候.
教宗利奧十世.
他為了籌措聖彼得大教堂的這個鉅額的資金.
因為他們當時候這個資金.
當然是第一個就是由教廷的國庫.
還有這些王公貴族的.
他們的這個交納的這些的貢獻.
另外就是這些朝聖者.
他們的奉獻.
除此之外.
還有很多的一筆是從哪裡來.
從當時候他們販賣贖罪券.
當時候呢.
那麼教廷就派了這麼一個人.

$^{401}$有沒有.
這個人.
就是右邊這個人.
那麼這個人呢.
教廷派他來.
他是德國名字.
但是把它翻成英文叫做Tessel.
Johann Tessel.
那這個特赦勒.
有的把它翻成特赦勒.
他被這個教廷指派.
就到了馬丁路德威丁堡.
他在那個地方的時候.
他就被指派到那個地方去販賣贖罪券.
那麼這個人極具煽動性的.
他用了一個極具煽動性的商業口號來宣傳.
他用的這一句話.
這個英文是說.
當然是德文翻成英文.
.
中文的意思是說.
當銀幣投入錢箱.
叮噹一響.
靈魂便瞬間從煉獄飛昇天堂.
在原文的德文裡面.
它是押韻的.
我覺得中文翻得很好.
也有押韻的味道.
所以馬丁路德看見他自己的會眾弟兄姊妹.
很多人就跑去跟他買贖罪券.
他心裡當時非常的憤怒.
這個就是導致馬丁路德.
他要改教的原因.
他立刻接著他就寫了95條大綱提綱.
然後貼在威丁堡上面.
其中有好多條都跟贖罪券有關係.
那麼沒有想到當他把這個貼出來之後.
整個的歐洲反應很響亮.
因此就產生了改教的運動.
親愛的同志們.

$^{441}$馬丁路德在這裡強烈的指責.
恩典是不可能不可以用買的.
耶穌基督的恩典是無法用金錢買贖的.
而今天你竟然用金錢交易來買這個上帝的恩典.
這是對耶穌基督的定死死之下是何等大的一個羞辱.
親愛的同志們.
如果我們今天還要靠律法來稱義.
那麼我們就與基督隔絕了.
從恩典中墜落了.
對不對.
求主幫助我們.
這是為什麼保羅在這裡這麼樣子的嚴肅.
甚至可以說嚴厲的在這裡責備.
提醒加拿大的信徒.
好了.
在這裡我們一直在講說這個律法.
是不是.
那跟我們有什麼關係呢.
我們今天也不會去買贖稅券了.
但是親愛的同志們.
我們今天會不會也還活在律法主義裡面呢.
我認識有一個弟兄.
他結婚的時候.
結婚初期的時候.
大概有很多年的時間.
跟太太有矛盾有衝突.
其實其中一個很大的衝突就是.
他來一次的原生家庭.
父親對他的要求是非常高.
近乎完美主義.
所以有一次他太太跟我抱怨說.
她說我老公他要求.
我們家的廚房必須要乾乾淨淨的.
而且要有光.
我們家的廚房也有光.
是油的光.
所以她的意思就是說.
她老公要求非常非常的乾淨.
所以你一般的媽不叫媽媽.
你知道廚房 廚房懂我意思吧.

$^{481}$你叫媽媽媽的乾淨嗎.
媽不乾淨.
特別一些污積連著加上有一點油什麼的.
那個東西要怎麼辦.
所以她老公就告訴她.
我示範給你看怎麼樣擦.
她是這樣子拿一塊馬布.
然後那個肥皂粉.
在地上這樣一塊一塊的這樣子刷的這樣子.
要求的非常的非常的高.
太太非常的痛苦.
親愛的同志們.
你有沒有一個老公是這樣子的.
你有沒有一個老婆是這樣子的.
親愛的同志們.
所以你會發現.
很多時候我們是活在什麼.
我們活在律法的裡面.
而不是活在恩典裡面.
在家裡面我們需要多一些恩典.
上個禮拜我們講說要多說恩言.
有沒有.
這個禮拜過去這禮拜有操練嗎.
來現在再操練一下.
跟你旁邊說你真棒.
我好欣賞你啊.
所以保羅在這裡提醒我們一件事情.
他說如果我們活在律法的中間.
我們甚至我們也要求別人活在這個律法的裡面.
你會帶來一個結果.
你會發現你就活在那種永遠你達不到律法的要求.
你就活在那種永遠你是緊緊張張.
你很擔心做父母親的也是如此.
你如果一直是用這一種的.
這種律法的要求.
小孩子考試考成績95分.
你說怎麼才考95分.
為什麼你不考100分.
小孩子給你看他的成績單.
我這一學期是A-.

$^{521}$那你怎麼說.
很棒.
Michael從此開始要這樣做.
你說很棒.
對不對.
但是今天我們中國父母親看到A-怎麼樣.
你為什麼不是A.
你為什麼不A+.
是不是.
我們好像我們期待的是一個完美的.
我們把我們自己的完美主義放在孩子的身上.
放在自己的配偶身上.
放在我們自己弟兄姊妹的身上.
所以呢.
別人很痛苦.
其實你快樂嗎.
你快樂嗎.
別人痛苦的時候別人會讓你快樂嗎.
不會吧.
所以夫妻之間也好.
婆媳之間也好.
父母親跟兒女之間也好.
你會發現我們兄弟姐妹之間也好.
我們很多時候.
如果我們用恩典相待.
就像聖經上說以恩慈相待.
跟你旁邊說我們要以恩慈相待.
好.
接下來第五節第六節.
保羅就開始講到.
他說我們怎麼樣子要憑著信心來稱義.
首先他第一個.
他講說受割離不受割離全無功效.
這句話的意思呢.
除了表明受割離.
你受割離並不能夠使你稱義.
他同時也表明了不受割離呢.
不受割離呢.
也沒有辦法使你稱義.
懂這意思嗎.

$^{561}$意思就是說稱義不稱義跟割離無關.
跟遵守立法沒有關係.
但是我們很多時候我們華人社會里面.
其實這一段聖經裡面.
那麼有一些的學者在解釋這個的時候.
他把這一段的經文裡面更擴大.
或者說更豐富了.
就講到說保羅在這個地方其實強調的意思.
不單單只是是不是靠立法.
還是靠信心的稱義的問題.
他在這地方要強調的就是說.
受割離是我們猶太人的一個identity.
懂我意思嗎.
所以他在這邊要讓我們看見的是.
保羅講這句話是要讓我們看見.
當我們信了耶穌之後.
我們不應該用我們的民族,種族,膚色,語言.
或者我們的這些的關係.
或者是我們的社會背景,經濟背景,地位.
所有這一切都不能夠成為我們中間的割離.
懂我意思嗎.
割離是使得猶太人跟外邊人分開的一個印記.
那什麼是使我們跟那些非我族類分開的呢.
我曾經提過.
在洛杉磯有個教會.
它就叫做.
什麼什麼名字我就不講.
叫做臺灣人教會.
OK 好 那請問.
臺灣人的教會.
這個教會的名字叫臺灣人教會.
那意思就是說不是臺灣人怎麼樣.
你就不要進來.
是吧 對不對.
那難不成我們今天再來一個上海人的教會.
上海人.
我們再開一間山東人的教會.
再開一間四川人的教會.
我都可以去.
就這麼一種意思.

$^{601}$所以同樣的.
在我們今天的北美的教會里面.
我們需要學習.
不管你從哪裡來的.
我們在基督裡面都是一家人.
我前面一個福祉的教會.
那個是二十年前.
那個時候剛好洛杉磯來了一批這種叫做.
就是都住在那種客棧.
叫家庭旅館.
家庭旅館 懂嗎.
就是他們把一個house租下來.
把它改成像旅館一樣.
其實就是.
哇 那裡面住十幾個人啊.
本來甚至有時候租個二十個人.
這些人在那邊可能一天只要五塊美金.
這樣子的費用.
所以很多人去住那邊.
這些人來了.
通通都是非法移民.
都是想要留下來這個身份.
好了.
那麼當時候我們教會有一對夫婦.
他們就很火熱.
他們看到這個需要.
他們就自己開始.
就這夫妻兩個人就去他們當中接他們啊.
到教會來啊.
幫他們去弄什麼什麼家居.
弄這些東西.
就是藉著幫助他們.
然後就帶領他們的信耶穌.
他在的那個團序呢.
是一個臺灣人的團序.
意思就是說.
其實我們早期的華人教會都這樣子.
對不對.
這個可能團序都是臺灣來的.
所以都是臺灣人.

$^{641}$那麼這弟兄把剛開始帶了幾個人來的時候.
大家還OK.
為什麼.
反正這幾個人不影響我.
後來越來越多.
越來越多.
哇.
本來是25人團序.
到後來變成50人的團序.
一半都是中國大陸來的.
這時候這些臺灣人開始感覺有壓力了.
開始有壓力了.
那他們就開始.
本來以前就說.
我們反正團序分小組嘛.
對不對.
那你們就弄你們自己小組.
我們是臺灣小組.
我們臺灣小組.
到後來變成這個弟兄就說.
哎呀.
我們大家一起團體一起來做這個志工吧.
他說哦.
一起來做什麼意思啊.
那就是說我也跟你一樣去接送他們啊.
到最後.
整個這個團體就被解放了.
聽懂了嗎.
整個團體就被解放了.
聽懂嗎.
哈雷路亞.
你們都不哈雷路亞嗎.
今天我們根本沒這個問題了.
是不是.
如果你還有這種種族的觀念.
只有這種所謂我是哪裡人的觀念.
你就在基督的恩典上面.
你沒有得到完全.
你要看見.
耶穌基督的恩典是白白的.

$^{681}$是給所有的人的.
Amen.
我們華人彼此之間會分你是哪裡人.
我是哪裡人.
那你想想看.
如果今天來了一個我們非洲朋友呢.
我們會說那個外國人.
有沒有.
那請問誰是外國人.
他是老外.
那請問你是老幾.
你有沒有發現.
我在意大利的時候.
我的學生們每次都說.
哎呀那些老外.
我說那你們是什麼.
你們纔是老外啊.
對不對.
可是你會發現我們很不知不覺的.
在我們的背景的裡面文化裡面.
我們就這樣區分了.
這種區分在我們中國的社會非常的強.
是不是.
所以你會發現這個公司裡面有一幫是叫上海幫的.
那你不是上海人.
最好考慮你要不要進入這個公司.
是不是.
我們很多時候會一問.
你是哪裡人.
老鄉.
這一老鄉起來.
我們大家就在一起.
就好像覺得我們這一群是比較特別的.
有沒有.
親愛的同志們.
但是親愛的同志們.
當你認識的耶穌之後.
請你既得一切主.
在基督裡面.
不分是猶太人.

$^{721}$犀利人都在基督的裡面成為一了.
成為一了.
那神造了一個新的人.
那個與佛所說講的那個新人就是指的教會.
這個教會里面不分是所謂的希臘人.
或者是非希臘人.
不管是猶太人.
或者是非猶太人都在基督的裡面成為一了.
所以我們受割離不受割離.
全無功效.
而且呢.
此人.
我們起來讀這段聖經.
我們靠著聖靈.
憑著信心.
這裡接連一再地講.
我們靠著聖靈憑著信心.
這邊講了信心.
然後後面又說使人深發仁愛的信心.
所以在這裡.
我們稍微解釋一下.
他說我們靠著聖靈憑著信心.
等候所盼望的意.
什麼意思.
等候所盼望的意.
這個意是你直接看這個中文.
這個表示說這個意.
你到了沒有.
有了沒有.
沒有.
對不對.
所以你盼望.
對不對.
那我請問你.
那保羅說我們因信稱義.
那你信了沒有.
那所以你稱義了沒有.
那為什麼你還要盼望.
為什麼你還盼望.
盼望表示你還沒有得著.

$^{761}$是不是.
所以這個就是在新約神學裡面的一個很重要的觀念.
叫做already but not yet.
已然和未然.
我們因著相信耶穌基督.
anytime.
你相信耶穌基督.
你就稱義了.
對不對.
但是這個意.
要一直等到主耶穌再來新天新地來臨的時候.
這個意纔會是完全的實現.
也就是說我們在神面前的那個地位.
稱義的地位.
纔是完完全全的落實了.
那我比較喜歡用得救跟得勝來區別.
這就是為什麼在我們的教會里面.
我非常強調這一點.
很多基督徒只知道得救.
不知道接下來要做什麼.
沒有錯.
得救很重要.
美國63\%的人說自己是基督徒.
但是隻有36\%的人.
他們知道他們自己會去天堂.
這是上次的調查.
那表示說有一半的人他不清楚他去哪裡.
你懂意思.
或者他沒有把握他去哪裡.
這是很可怕的一件事情.
是不是.
聚會聚了一輩子.
竟然不知道自己可不可以進天國.
好.
那今天用什麼.
我要問大家一件事情.
再問大家一件事情.
你有把握可以進天國嗎.
我不要你舉手了.
因為可能你也覺得不好意思.

$^{801}$你就不舉手.
但是你一定要問你自己這個問題.
那我再問.
這件事情太重要了.
你知道.
在這個山東大復興的時候.
這個神使用一位挪威來的一個旅宣教士.
單身旅宣教士叫做孟教士.
他為中國的復興禱告了23年.
一個人.
孤孤單單的一個人禱告了23年.
後來神有一天告訴他說.
神說我要來複興中國.
所以他知道神已經應許了他的禱告.
好.
make a long story short.
那麼孟教士呢.
他常常在帶領聚會的時候.
他其中一個常常講的一個信息.
就是重生得救.
他常常問別人的一句話也是問.
你重生了嗎.
那麼有一次他到一個教會去講道.
那麼這個教會的牧師呢.
聽說孟教士要來.
他知道孟教士最喜歡問人家.
就是你重生了沒有.
那他實在覺得自己也沒有把握.
所以呢.
那一天當孟教士講道完了之後呢.
通常都會到後面.
像牧師一樣跟大家握手問安.
對不對.
好了.
他一看在禱告當中眼睛閉起來.
眾人眼睛閉起來.
他就一個箭步.
他就跳窗戶.
那個年代的窗戶是.
有點像是落地窗的.

$^{841}$有一點點的高度.
懂意思嗎.
那麼所以一打開.
你要跨個步就過去了.
他那一天就這樣子.
從一個窗戶裡面就這樣跨.
一打開然後就跨出去.
他想要趕在孟教士出來之前.
沒有想到他正跨的時候.
他有一點碰到了怎麼樣.
就稍微就這樣突下去了.
這時候.
當他一抬身起來的時候.
孟教士就在他面前說.
你重生了嗎.
哇.
那一天孟教士的你重生了嗎.
這個牧師重生了.
親愛的兄弟妹妹.
你重生了嗎.
問問你旁邊你重生了嗎.
重生得救因性稱義.
講的是同一件事情.
所以怎麼樣子才能夠重生呢.
怎麼樣子才能夠稱義呢.
怎麼樣子.
因著信.
那這個信.
羅馬書第一章十細節.
我們一起來讀好不好.
因為神的義.
寫明出來.
這義是本於信.
也如經上所記.
義人必因信得生.
好.
這一串真經裡面你注意看.
他說這個義是本於信.
以至於信.
原文的意思就是說.

$^{881}$from faith to faith.
是不是.
我們今天在買機票的時候.
他還問你說from.
你說哪裡.
L A X.
to 哪裡.
Rome OK.
或者是這個Paris.
好.
你從哪裡到哪裡是不是.
所以這裡說.
他說本於信.
以至於信的意思就是說.
信心是你的起點.
信心是你的終點.
懂嗎.
也就是說我們的信.
我們的信仰從信心開始.
因著相信我們稱義.
也因著相信信心.
我們得勝.
我們成勝.
Amen.
信心扮演一個非常重要的角色.
在我們整個基督徒的一生當中.
我們是從信以至於信.
from faith to faith.
所以從這個角度來講.
弟兄姊妹.
你不能夠停留在.
當我信了耶穌之後.
我受洗之後.
我就停擺了.
你的信心沒有增加.
你的信心不但不增加.
還可能越來越減退.
剛信的時候覺得禱告很有信心.
信主久了之後.
越禱告越沒信心.

$^{921}$甚至不禱告.
有沒有.
因為你跟神的關係是不可能停止的.
一個baby生出來之後.
如果他真正的重生了.
真正的有生命的話.
這個生命一定要成長的.
怎麼成長.
就好像財神的嬰孩.
怎麼樣.
愛慕臨來.
你愛慕臨來嗎.
還是你愛慕手機.
弟兄姊妹.
當我們愛慕主的話.
那是一個正常基督徒應該有的現象.
當我們喫了主的話.
主的話就使得我們成長.
我們的信心因著主的話就增加了.
Amen.
聖靈在我們讀聖經的時候幫助我們.
使神的話不單單是成為一個字句.
也使它成為經意.
所以神的話 神的靈.
兩個彼此配合.
使得我們的生命越來越成長.
Amen.
好 最後.
信心必定帶來相愛.
這段經文裡面.
它說唯獨使人生髮仁愛的信心纔有功效.
NASB 這個是中文的新義本的翻譯.
NASB 的翻譯就是.
把Faith Walking Through Love.
Walking Through Love.
這個信心是透過.
這個什麼意思叫 Walking Through Love.
這個意思就是說.
信心是透過愛心的行為.
來表現出來.

$^{961}$Amen.
所以一個真實的信心一定有愛的行為.
所以保羅在問安的時候.
常常把信 望 愛 怎麼樣.
放在一起談.
你沒辦法把這三者分開的.
信 望 愛 這三個完全是.
它們是相在一體的.
它們是三合一的.
三合一咖啡.
你不能夠把它分開來.
它們三個人彼此相連的.
所以真實的信心一定會帶出愛的行動.
愛的行為.
有一句話.
我們應該都聽過.
說家就是講愛的地方.
不是講理的地方.
聽過嗎.
那你們家是講愛還是講理.
你會發現剛結婚的夫妻常常都是什麼.
講理.
是不是啊.
講理.
今天輪到你洗碗了.
為什麼你每次都要拖到.
拖這麼久.
拖到隔天.
你這樣很不負責任.
我們明明好好.
葛牧師在幫我們做婚前輔導的時候.
我們講好的.
做家是每個人一天的.
你不守規矩.
對不對.
當對方這樣說的時候.
你會怎麼反應.
是是是.
小的不對.
求你赦免.

$^{1001}$對嗎.
你會怎麼反應.
請告訴我.
你們結婚這麼多年.
都不知道怎麼反應.
怎麼反應.
Michael.
找葛牧師.
當他這樣反應的時候.
他跟你講理.
對不對.
那你會怎麼樣.
你也會跟他講理.
是不是.
你也是啊.
你也不是隻有我.
你也是.
你看看那一天.
他記得很清楚.
那一天你也沒有.
輪到你洗碗.
你也沒有洗碗啊.
好了扯平.
扯不平.
講理永遠講不通.
懂這意思嗎.
所以.
這個英國的劇作家.
瓦爾德.
他曾經講過一句話很經典.
特別弟兄要聽.
Women are meant to be loved.
not to be understood.
女人是用來被愛的.
不是用來被理解的.
Amen.
所有的弟兄說 Amen.
我結婚了這麼多年.
我是完全瞭解.
我是一個比較理性的人.

$^{1041}$我只能講比較.
我不算是很理性.
我是一個比較理性的人.
你沒有一個道理.
你要我這樣做.
很難.
但是婚姻生活當中.
當我講理的時候.
我就發現不work.
是不是.
對別人肯work.
對自己的配偶往往不work.
是吧.
好.
所以怎麼辦呢.
所以家裡不是講理的地方.
家裡是什麼.
講愛的地方.
親愛的同樣的.
在教會里面也是如此.
教會.
如果都是規定.
Policy.
By law.
你不可以這樣子.
門要隨時關.
我不是說教會不要有規矩.
不要.
我也是很看重規矩的.
不要誤會我意思.
但是我的意思就是說.
親愛的兄弟妹.
當你教會的裡面.
充滿了這些規規條條的時候.
請問.
弟兄姊妹來到這個教會.
他會如何感覺.
緊緊張張.
是不是.
但是如果我們見了面的時候.

$^{1081}$我們大家在一起的時候.
我們以恩慈相待.
我們常常彼此說恩言.
我們常常彼此包容.
彼此瞭解.
今天你來晚了.
是不是發生什麼事情了.
是不是身體不好了.
有沒有.
但是你今天一個弟兄姊妹來晚了.
你心裡想.
你怎麼今天來這麼晚了.
你看到沒有.
兩種反應帶來什麼不同的結果.
懂我說的意思嗎.
一個是他知道他被瞭解.
另外一個他覺得.
你都別問我為什麼會來晚.
我正被炒魷魚了.
我心情這麼不好.
我今天掙紮了半天.
我纔想到我還是來教會吧.
但是你一見到我就說.
你為什麼來晚了.
弟兄姊妹.
家讓我們的教會也成為一個家.
讓你的教會.
讓你的小主成為一個家.
讓我們的家.
讓我們的教會成為一個講愛的地方.
講恩典的地方.
Amen.
大聲一點.
Amen.
親愛的弟兄姊妹.
我們何等需要.
因為主就是那位恩典的神.
Amen.
祂是道成的肉身.
充充滿滿有恩典.

$^{1121}$有真理的神.
我們不單單要好好地遵循神的話.
神的話是我們的真理.
By the way.
我要區別一下.
神的話是我們的是真理.
是我們要遵守的.
而不是所謂的龜橋.
因為龜橋.
它不一定是符合神的話.
懂這意思嗎.
龜橋會受到時代的影響.
它可能應該要改變.
所以我們需要有真理.
但是我們也需要恩典.
Amen.
求主讓我們建造一個.
充充滿滿有恩典有真理的教會.
Amen.
求主幫助我們.
讓我們學習更多地去體恤別人.
去幫助別人.
去原諒別人.
去接納別人.
從你的家中.
從你的配偶.
從你最親密的同工開始.
讓我們大家一起來操練.
Amen.
\newpage



\section{羅馬書 15:1-6}
\label{sec:032202}
\textbf{彼此擔當}~~~~張健庭牧師博士
\newline
\newline
連結: \href{http://www.church.com.hk/acms/content.asp?site=cdc\&op=show\&type=product\&code=032202\&layout=sermon}{\texttt{032202}} ~~~~ 語音日期: 2026-05-24
\newline
\newline
\hyperref[sec:032201]{\small{< < < PREV SERMON < < <}}
~
\hyperref[sec:index_chronic]{\small{[返順時目]}}
~
\hyperref[sec:032203]{\small{> > > NEXT SERMON > > >}}
\newline
\newline
羅馬書 15:1-6
\newline
\begin{longtable}{cl}
\hline
\hline
章節 & 經文 (和合本修訂版)\\
\hline
15:1 & \begin{tabularx}{0.7\textwidth}{X} 我們堅強的人應該分擔不堅強的人的軟弱，不求自己的喜悅。 \end{tabularx} \\ \\ \relax
15:2 & \begin{tabularx}{0.7\textwidth}{X} 我們各人務必要讓鄰人喜悅，使他得益處，得造就。 \end{tabularx} \\ \\ \relax
15:3 & \begin{tabularx}{0.7\textwidth}{X} 因為基督也不求自己的喜悅，如經上所記：「辱罵你的人的辱罵都落在我身上。」 \end{tabularx} \\ \\ \relax
15:4 & \begin{tabularx}{0.7\textwidth}{X} 從前所寫的聖經都是為教導我們寫的，要使我們藉著忍耐和因聖經所生的安慰，得著盼望。 \end{tabularx} \\ \\ \relax
15:5 & \begin{tabularx}{0.7\textwidth}{X} 但願賜忍耐和安慰的神使你們彼此同心，效法基督耶穌， \end{tabularx} \\ \\ \relax
15:6 & \begin{tabularx}{0.7\textwidth}{X} 為使你們同心同聲榮耀我們主耶穌基督的父神！ \end{tabularx} \\ \\
[1ex]
\hline
\hline
\end{longtable}
$^{1}$主席.
陳偉業議員.
主席.
湯家驊議員.
陳偉業議員.
湯家驊議員.
陳偉業議員.
湯家驊議員.
陳偉業議員.
湯家驊議員.
陳偉業議員.
湯家驊議員.
\newpage



\section{民數記 11:1-35}
\label{sec:032203}
\textbf{百姓之怨言和摩西之不滿}~~~~林誠信牧師
\newline
\newline
連結: \href{http://www.church.com.hk/acms/content.asp?site=cdc\&op=show\&type=product\&code=032203\&layout=sermon}{\texttt{032203}} ~~~~ 語音日期: 2026-05-24
\newline
\newline
\hyperref[sec:032202]{\small{< < < PREV SERMON < < <}}
~
\hyperref[sec:index_chronic]{\small{[返順時目]}}
~
\hyperref[sec:032204]{\small{> > > NEXT SERMON > > >}}
\newline
\newline
民數記 11:1-35
\newline
\begin{longtable}{cl}
\hline
\hline
章節 & 經文 (和合本修訂版)\\
\hline
11:1 & \begin{tabularx}{0.7\textwidth}{X} 百姓發怨言，惡言傳達到耶和華的耳中。耶和華聽見了就怒氣發作，耶和華的火在他們中間焚燒，燒燬營的外圍。 \end{tabularx} \\ \\ \relax
11:2 & \begin{tabularx}{0.7\textwidth}{X} 百姓向摩西哀求，摩西祈求耶和華，火就熄了。 \end{tabularx} \\ \\ \relax
11:3 & \begin{tabularx}{0.7\textwidth}{X} 那地方就叫做他備拉，因為耶和華的火曾在他們中間焚燒。 \end{tabularx} \\ \\ \relax
11:4 & \begin{tabularx}{0.7\textwidth}{X} 他們中間的閒雜人動了貪慾的心；以色列人又再哭著說：「誰給我們肉吃呢？ \end{tabularx} \\ \\ \relax
11:5 & \begin{tabularx}{0.7\textwidth}{X} 我們記得在埃及的時候，不花錢就可以吃魚，還有黃瓜、西瓜、韭菜、蔥、蒜。 \end{tabularx} \\ \\ \relax
11:6 & \begin{tabularx}{0.7\textwidth}{X} 現在我們的精力枯乾了。除了這嗎哪以外，在我們眼前甚麼都沒有。」 \end{tabularx} \\ \\ \relax
11:7 & \begin{tabularx}{0.7\textwidth}{X} 嗎哪好像芫荽子，看上去如同樹脂的樣子。 \end{tabularx} \\ \\ \relax
11:8 & \begin{tabularx}{0.7\textwidth}{X} 百姓四處走動撿取嗎哪，把它用磨磨碎或用臼搗成粉，在鍋中煮了做成餅，滋味好像油烤餅的滋味。 \end{tabularx} \\ \\ \relax
11:9 & \begin{tabularx}{0.7\textwidth}{X} 夜間露水降在營中，嗎哪也隨著降下。 \end{tabularx} \\ \\ \relax
11:10 & \begin{tabularx}{0.7\textwidth}{X} 摩西聽見百姓家家戶戶在帳棚門口哀哭。因此， 耶和華的怒氣大大發作，摩西看了也不高興。 \end{tabularx} \\ \\ \relax
11:11 & \begin{tabularx}{0.7\textwidth}{X} 摩西對耶和華說：「你為何苦待僕人？我為何不在你眼前蒙恩，竟把這眾百姓的擔子加在我身上呢？ \end{tabularx} \\ \\ \relax
11:12 & \begin{tabularx}{0.7\textwidth}{X} 這眾百姓豈是我懷的胎，豈是我生下來的呢？你竟對我說：『把他們抱在懷裡，如養育之父抱著吃奶的嬰孩，一直抱到你起誓應許給他們祖宗的土地去。』 \end{tabularx} \\ \\ \relax
11:13 & \begin{tabularx}{0.7\textwidth}{X} 我從哪裡拿肉給這眾百姓吃呢？他們都向我哭著說：『給我們肉吃！』 \end{tabularx} \\ \\ \relax
11:14 & \begin{tabularx}{0.7\textwidth}{X} 我不能獨自帶領這眾百姓，這對我太沉重了。 \end{tabularx} \\ \\ \relax
11:15 & \begin{tabularx}{0.7\textwidth}{X} 如果你這樣待我，倒不如立刻把我殺了吧！我若在你眼前蒙恩，求你不要讓我再受這樣的苦。」 \end{tabularx} \\ \\ \relax
11:16 & \begin{tabularx}{0.7\textwidth}{X} 耶和華對摩西說：「你要從以色列的長老中為我召集七十個人，就是你所認識，作百姓的長老和官長的，領他們到會幕，使他們和你一同站在那裡。 \end{tabularx} \\ \\ \relax
11:17 & \begin{tabularx}{0.7\textwidth}{X} 我要在那裡降臨，與你說話，把降給你的靈分給他們。他們就和你分擔帶領百姓的擔子，免得你獨自承擔。 \end{tabularx} \\ \\ \relax
11:18 & \begin{tabularx}{0.7\textwidth}{X} 你要對百姓說：『你們要為了明天使自己分別為聖，你們將有肉吃。因你們哭著說：誰給我們肉吃呢？我們在埃及多麼好！這聲音傳到了耶和華的耳中，所以他必給你們肉吃。 \end{tabularx} \\ \\ \relax
11:19 & \begin{tabularx}{0.7\textwidth}{X} 你們不只吃一天、兩天、五天、十天、二十天， \end{tabularx} \\ \\ \relax
11:20 & \begin{tabularx}{0.7\textwidth}{X} 而是整整一個月，直到肉從你們的鼻孔噴出來，使你們厭惡。因為你們厭棄那住在你們中間的耶和華，在他面前哭著說：我們為何出埃及呢？』」 \end{tabularx} \\ \\ \relax
11:21 & \begin{tabularx}{0.7\textwidth}{X} 摩西說：「跟我在一起的百姓，步行的男人就有六十萬，你還說：『我要把肉賜給他們，使他們可以整整吃一個月。』 \end{tabularx} \\ \\ \relax
11:22 & \begin{tabularx}{0.7\textwidth}{X} 難道宰了羊群牛群，就夠給他們嗎？或者把海中所有的魚都捕來，就夠給他們嗎？」 \end{tabularx} \\ \\ \relax
11:23 & \begin{tabularx}{0.7\textwidth}{X} 耶和華對摩西說：「耶和華的膀臂豈是縮短了嗎？現在你要看我的話向你應驗不應驗。」 \end{tabularx} \\ \\ \relax
11:24 & \begin{tabularx}{0.7\textwidth}{X} 摩西出去，把耶和華的話告訴百姓，並從百姓的長老中召集七十個人來，叫他們站在會幕的四圍。 \end{tabularx} \\ \\ \relax
11:25 & \begin{tabularx}{0.7\textwidth}{X} 耶和華在雲中降臨，對摩西說話，把降給他的靈分給那七十個長老。靈停在他們身上的時候，他們就說預言，以後卻沒有再說了。 \end{tabularx} \\ \\ \relax
11:26 & \begin{tabularx}{0.7\textwidth}{X} 但有兩個人仍在營裡，一個名叫伊利達，一個名叫米達。他們本是在那些登記的人中，卻沒有到會幕那裡去。靈停在他們身上，他們就在營裡說預言。 \end{tabularx} \\ \\ \relax
11:27 & \begin{tabularx}{0.7\textwidth}{X} 有一個年輕人跑來告訴摩西說：「伊利達和米達在營裡說預言。」 \end{tabularx} \\ \\ \relax
11:28 & \begin{tabularx}{0.7\textwidth}{X} 嫩的兒子約書亞，年輕時就作摩西的助手，說：「請我主摩西禁止他們。」 \end{tabularx} \\ \\ \relax
11:29 & \begin{tabularx}{0.7\textwidth}{X} 摩西對他說：「你為我的緣故嫉妒嗎？惟願耶和華的百姓都是先知，願耶和華把他的靈降在他們身上！」 \end{tabularx} \\ \\ \relax
11:30 & \begin{tabularx}{0.7\textwidth}{X} 於是，摩西回到營裡去，以色列的長老也回去了。 \end{tabularx} \\ \\ \relax
11:31 & \begin{tabularx}{0.7\textwidth}{X} 有一陣風從耶和華那裡颳起，把鵪鶉從海上颳來，散落在營地和周圍；一邊約有一天的路程，另一邊也約有一天的路程，離地面約有二肘。 \end{tabularx} \\ \\ \relax
11:32 & \begin{tabularx}{0.7\textwidth}{X} 百姓起來，整天整夜，甚至次日一整天，都在捕捉鵪鶉。每人至少捉到十賀梅珥，各自擺在營的四圍。 \end{tabularx} \\ \\ \relax
11:33 & \begin{tabularx}{0.7\textwidth}{X} 但肉在他們牙間還未咀嚼時，耶和華的怒氣向百姓發作，用極重的災禍擊殺百姓。 \end{tabularx} \\ \\ \relax
11:34 & \begin{tabularx}{0.7\textwidth}{X} 那地方就叫基博羅‧哈他瓦 ，因為他們在那裡埋葬了貪慾的百姓。 \end{tabularx} \\ \\ \relax
11:35 & \begin{tabularx}{0.7\textwidth}{X} 百姓從基博羅‧哈他瓦起程，到哈洗錄，就住在哈洗錄。 \end{tabularx} \\ \\
[1ex]
\hline
\hline
\end{longtable}
$^{1}$將這首美好的詩歌呈獻在神的面前.
但願我們每一個人都能夠將我們的生命.
如此地獻給上帝.
如果昔日的以色列人.
是有一個心口如一的這種獻生命給主的話.
我就相信《民數記》十一章就不會出現.
《民數記》十一章我們看見.
他們其實剛剛出埃及大概一年多的時間.
如果你看看那個時間表.
第十章是說他們是開始.
由西乃山那裡起行的.
那是他們第二年的第二月.
而在當時其實他們是經歷過.
頭一年他們出埃及見證過過紅海.
而摩西也領受了十屆.
也照著神的吩咐建造了會幕.
其實各樣的事情我們都可以說.
是你看到神的恩典多麼地浩大.
但竟然在這張聖經我們看見三件事出現了.
而這三件事.
你看表層的原因.
都只是一些所謂怨言不滿等等.
但矛頭到最後其實是直指上帝自己.
我們如果留意來看百姓的怨言.
他們在抱怨些什麼呢?.
聖經沒有說清楚.
不過有可能是否缺乏水呢?.
但他們在第一年的時候不是沒有缺乏過水.
他們一出埃及第三天.
他們就在曠野開始埋怨.
沒有水喝了.
誰知道水泉是馬拉的地方.
誰知道水是苦的.
神吩咐摩西如何將種樹枝放在裡面.
水就由苦變甜.
神行很多神蹟.
後來他們去到利山定的地方的時候.
也是一樣.
他們沒有水.
神就叫摩西擊打盤石.

$^{41}$結果他們有水喝.
他們在這個曠野的時候.
他們又抱怨他們缺乏食物.
我們留在埃及就好了.
他們這種抱怨.
神是為他們設下什麼呢?.
是每天都有嗎哪.
去到一星期.
去到安息日之前.
他們總之就拿夠兩天.
如是者一年多了.
已經發生過很多這些事.
你發覺其實百姓的怨言.
其實本身表面的因素是一些環境上的問題.
但實際因素是他們內心.
那種對神的不信服.
我們一起來讀第一三節.
「百姓發怨言.
惡言傳達到耶和華的耳中.
耶和華聽見就怒氣發作.
耶和華的火在他們中間焚燒.
燒毀形的外圍.
百姓向摩西哀求.
摩西祈求.
耶和華火就熄了.
那地就叫塔比拉.
因為耶和華的火曾在他們中間焚燒」.
我們看到這裡沒有說怨言是什麼.
但我們可以推斷有可能.
是不是因為不夠水喝.
這個惡言傳到神的耳中.
神其實我們說什麼在地上祂都知道.
神是監察人心的主.
而耶和華聽見之後祂怒氣發作.
祂是針對百姓.
這個怨言的背後.
其實就是針對神的引導.
而當百姓看到神發怒.
神燒毀了形的外圍.
神沒有燒死他們.

$^{81}$只是形的外圍.
百姓就向摩西哀求.
結果神就止息了這些火.
但那個地方就叫塔比拉.
塔比拉的意思就是火焚之地.
火燒起來.
火燃燒的其實給我們是一個警告.
是一個很危險的警告.
沒有燒掉他們的十二營.
我們在《民數記》之前.
這十章裡面都試過有好幾次.
都給了地圖給大家看.
東南西北.
他們每個方位都有三個支派去紮營.
如果我們看見神提醒他們.
怨言不是一件小事.
怨言可以破壞一個群體.
好像保羅在《菲勒比書》說.
不要發怨言.
接著又怎樣?.
就起爭論了.
原來怨言可以蔓延的.
怨言如果你不是針對事情.
你變成是針對一個人.
或者針對可能是一個領導者.
或者是一個管理者.
你將你的怨氣,憤怒投射在他身上.
其實到最後是解決不了事情的.
我們要針對事.
不要針對人.
我們這樣說.
是不是不可以作出任何的投訴呢?.
其實你在家庭.
你在明天上班辦公室的地方.
或者你要見客人.
你可能淋淋種種的事.
包括在教會.
你可能都見到一些事.
可能辦得不夠最好.
你要去改善那些事情.

$^{121}$我們是要透過合情合理的表達方式.
千萬不要用一些激進.
甚至是違法的手段等等.
我們是需要將這些所謂投訴,建議.
我們要用一些進言.
將一些想你自己改善的建議.
能夠好好的透過一個渠道去表達.
大家更加不要忘記.
在提出這些意見之後.
一定要祈禱.
為了那件事我們要扮好.
要糾正那些事情.
所以是要對準神.
我們求恩.
求神幫助我們去解決那些事情.
同樣叫當事人都明白.
我們一起去糾正那些需要糾正的事情.
第二方面我們看到百姓的貪欲.
表層的原因好像是缺乏肉食.
但這是很表面的.
實際的原因是他們厭棄神供應的猛男.
我們在第四節看到.
他們中間的閒雜人動了貪欲的心.
閒雜人是指甚麼人?.
他們出埃及的時候.
不只是以色列這十二個支派的神的百姓.
在當中他們有一些寄居者.
一些外族人.
是和以色列聚在一起的.
神吩咐他們.
你們都要紀念孤兒寡婦.
也要紀念你們當中的外族人.
這些外族人在當中.
不是住在那十二個營的支派裡面.
他們是住在營的外圍的地方.
但就是這些閒雜人.
他們開始說.
我們天天只吃這些猛男.
他們覺得厭倦.
所以他們很想怎樣?.

$^{161}$他們很想希望能夠吃肉.
他們是聳動那些以色列人.
哭哭啼啼.
我們記得在埃及的時候.
不用花錢就可以吃魚.
有黃瓜,西瓜,韭菜,蔥,蒜.
我們中國人都很熟悉的.
現在我們的精力枯乾.
除了這猛男以外.
在我們眼前甚多都沒有.
我不知道大家會不會覺得.
你試過這種味道沒有?.
真是很眷戀.
我和師母在月初的時候.
因為去了以色列.
不是以色列.
是意大利.
意大利參加我乾女兒的婚禮.
他們現在年輕人的婚禮.
是喜歡去其他生活以外的地方.
他說那裡辦婚禮比在三藩市便宜.
我們的婚禮是四天的.
第一天去歡迎晚餐的時候.
是吃薄餅的.
第二天就要去吃意大利面.
不過你是自己找麵粉.
自己搓.
搓完自己吃.
很噁心也好.
人家做的薄餅應該是一條條的.
薄餅.
你又做了那些不知道變成什麼.
不過我們混在一起吃的.
我們一群人一起很開心的.
不過如是者.
去到第三天又吃意大利面.
第四天又薄餅.
跟著我跟師妹已經說.
哎呀 離開了.
離開這個結婚的地方.

$^{201}$我們要找一些中國菜去吃.
我們已經心裡有些不好了吧.
還有幾天的旅程.
我們希望不要每天都是吃這些吧.
而且我們發覺越吃越糟.
搞不定.
會很肥.
我們中國人是很懂得吃的一個民族.
我們也會抱怨.
但這些抱怨.
有時候我們說.
是因為我們餵食而已.
或者你有你就可以比較.
以色列人也是.
不過這種比較是針對.
上帝明明不是因為他們缺乏食物.
是因為他們厭棄神的供應.
究竟這些嗎哪是什麼呢?.
是什麼東西來的?.
我們看看第七到第九節.
他是這樣形容的.
嗎哪好像芫茜籽.
其實是一種植物的種子.
它是很小粒的.
是白色的.
就好像夜間.
他們要下雪.
那些雪粒一樣的.
他們四處去撿取這些嗎哪.
然後磨成粉.
就在鍋裡面煮成餅.
可以用橄欖油.
烤這些餅.
其實很香口的.
好像印度人吃的那些Naan.
那些薄餅.
當夜間露水降在營中.
嗎哪也會隨著降下.
他們早上會去撿拾的.
這是聖經的一些插圖畫家.

$^{241}$他們嘗試表達給我們看.
但我們看到今天.
我們有沒有埋怨神的供應?.
我們比較別人的豐富.
我們來到有貪欲.
這種貪欲就是你貪圖神未有應許的東西.
神已經為我們擺設好筵席.
主耶穌教我們的祈禱.
是為你每天的飲食.
不是為你未來一年.
或者你未來五年.
你不是.
需要的我們緊緊跟隨上帝.
是一個信心的歷程.
人生的道路都是一樣.
我們在任何階段都需要學習這個東西.
今天我嘗試用這個反思.
來看我們的婚姻.
其實婚姻對我們每一個弟兄姊妹.
我們如果有經歷過.
這是上帝對我們一個美麗的恩典和祝福.
在美國有一個長壽節目叫做《初見即婚》.
這個節目是吸引一些單身人士.
他們因為現在這個時代.
有時可能工作員未必可以遇到.
你的如意郎君.
或者你所謂一直在想的白雪公主的出現.
但是這個節目其實是很危險的一件事.
他要那些人素未謀面.
但是他們經過了.
有大概五位心理輔導專家.
他們看看他們兩個人.
男和女之間的背景.
然後將他們撮合.
他們一來到.
他們就簽一個結婚的協議書.
是根據當地的法律.
真的是成婚的.
由第一晚就開始住在一起.
之後的八個月時間.

$^{281}$他們就去看.
如果之後你決定了之後.
你要離婚都可以的.
結果有三分二的人就離婚了.
成功率到最後只是14\%.
我告訴大家.
我們今時今日.
我們的婚姻.
不是需要這樣交在.
一個好像盲婚硬嫁的狀態下.
我們需要大家有認識.
需要大家有觀察.
大家彼此用一個恆久忍耐的心.
我們兩夫妻一起共諧年齡之後.
要白頭到老走完這個婚姻歷程的話.
你們必須就好像保羅所說的.
要恆久忍耐.
又有恩慈.
你不要說我死忍爛忍.
我忍夠你了.
我試一下測試一下.
我們中間有沒有人已經結婚超過30年.
可不可以舉手給我看.
哇 都很多.
為你們感恩.
30年了.
有沒有結婚是10年以下的呢?.
10年以下的 有沒有?.
有啊.
我們給些年輕人鼓勵他們.
其實我們婚姻真是有一個歷程的.
我今年年尾會退休.
不過有一位很有心的教會的肢體.
他說要贊助我.
將我十多年前的那本《兼諜情深》.
這個《兼諜情深》是第二版的.
他說要贊助我去出版.
到時我會鼓勵.
有誰還沒有可以去買.
買完就奉獻給教會.

$^{321}$就是這樣.
我希望這是一個小小的心意.
但裡面講的東西.
都是十多年前提的東西.
和今時今日沒有分別.
我們需要的東西其實都是一樣的.
不過不知道為什麼現在的設計師.
就由《兼諜情深》變成這樣.
不知道是不是說我年紀大了.
會比較成熟好一些.
所以其實內容是一樣的.
不過是第三版.
我盼望大家都能夠好好去學習.
我們中國.
我們大家都公認.
《萬世詩表》孔子在《禮記》.
《譚公下》這個古卷裡面.
有一個很出名的說話.
是應用在婚姻裡面都可以的.
叫「同仇之禮所以人向黨也」.
他說我們能夠結為那些鄉親.
那些黨派等等.
是因為大家有一種結盟的背後.
就好像在一起去共患難.
這個「同仇」是一個這樣的意思.
是大家去有同理心.
去面對貧病,喪親,冤屈,孤獨等等.
是給予尊重和關懷.
但是你要約束你的行為.
在面對苦難的認知和體恤.
這個同理心.
都需要我們學習「以禮相待」.
當我們如果應用在兩夫妻.
你說是不是坐了牢呢?.
做了老襯呢?.
「同仇」嘛.
其實不是這樣的意思.
「同仇」對夫妻來說.
即是說你是失去了單身的時候.
部分的自由.

$^{361}$你是進入面對好像一些束縛和限制.
特別面對經濟的壓力.
養兒育女,家務等等.
其實需要兩方面大家要有協商的.
要大家共同一起去努力的.
不要令到對方不安.
你需要維護對方的尊嚴和體面.
所以你不要說.
你現在結了婚三十年,四十年.
你才說我要和她離婚.
因為我再不能忍受她的壞習慣.
你明知道的.
頭七年你都知道她的壞習慣是改不了的.
改不了是用愛,包容,忍耐,善意的提醒等等.
如果都不能改變她.
你就為她祈禱.
讓這些壞習慣不要惡化.
你需要將這些壞習慣慢慢昇華.
其實兩個人一粒一特不需要相同.
但需要大家彼此有這種愛.
而且這個同仇是表示大家有一種情在裡面的.
不是冷漠無情的.
不是困在一個監獄用怨恨冷漠去控制對方.
這個變成家暴.
這個變成兩夫妻的關係一定會破裂.
我們需要用一種體諒,關心,尊重.
大家一起共行.
婚姻的路途裡.
風霜雨雪.
有時就算不是常態也必然在我們人生歷程裡會經歷的.
但正正在那些困難,極大的患難來臨的時候.
我們要搭建這種避難所.
就是夫妻不只是一起共享所謂的榮譽,富貴的時候.
當有患難的時候不是各自飛.
患難的時候更加要手執之手.
與子皆老.
你不要放棄.
你要和他一起去面對.
我們看到摩西的不滿.
這是第三方面要處理的事情.

$^{401}$當時他的抱怨就是.
每個人在他們的帳幕門外哀哭.
因此瑤華的怒氣大大發作.
因為上帝知道他們抱怨的矛頭.
不只是指向摩西.
而是指向摩西背後的上帝.
摩西自己也受不了.
被那些百姓.
他不是現在才發覺他們的怨言.
第一年已經是這樣.
不同的事件都出現了.
摩西現在他向神就開始.
好像有一點發脾氣.
他一年有四個問題在這裡跟神傾訴.
他說:你為什麼苦待僕人呢?.
如果我在你面前.
你給我蒙恩.
如果是蒙恩.
你為什麼要將這些百姓的擔子加在我身上呢?.
第三個問題.
他說:這些百姓喜事我懷的胎.
不是我生下來的.
是上帝來的.
但是你竟然對我說.
你要將他們抱入懷.
好像養育他們的爸爸.
好像餵他們吃奶的嬰孩.
一直抱到你喜誓應許.
他們祖宗的土地去.
如果在這個場景.
百姓他們有一個生命的真正的悔改和回轉.
他們其實由那個地方西乃山.
進入去應許之地.
你最多走三個月就可以了.
是有的.
但是結果之後發生什麼事呢?.
就是後文我們十三十四章.
十二個探子探完之後.
個個都說不可以去.
只有約書亞和加納.

$^{441}$就是這兩個他們說.
上帝應許我們.
為什麼我們不去呢?.
我們要憑信心.
不用怕那些亞納族人.
他們高大衛猛.
但是上帝應許給我們.
不用怕.
我們是靠上帝的說話.
來領受上帝的恩典和祝福.
但是我們看到百姓的情況.
他們整個回眾的屬靈的狀態.
你就明白了.
整個屬靈的心境.
他們有很多這些埋怨,不信任.
和厭棄神住在他們中間.
那你怎樣去打敗那些亞納族人呢?.
根本沒有可能.
所以摩西也猜不到.
原來之後還要撐多差不多39年的時間.
摩西接著說下去.
第十三節.
他說我從哪裡可以拿這麼多肉給那些百姓吃.
但是他們非要纏著我.
給我們肉吃.
我不能獨自帶領這些百姓.
這個對我太沉重了.
接著摩西最後的殺手簡說.
如果你這樣對我.
我求上帝你不如現在馬上殺了我.
求死啊.
摩西的壓力爆煲了.
到現在這個時候.
好像以利亞經歷過加密山的大大的得勝之後.
被耶誓別一句說話.
明日日落之前我不取你命.
我就不叫耶誓別了.
是這個意思.
結果竟然經歷屬靈這麼大的復興的以利亞.
他也經歷了精神沮喪,懼怕.

$^{481}$我們看到上帝很有憐憫.
上帝很有恩情.
有些聖經學者將這一段經文.
將「我」和「你」變成一個對話.
這個對話其實是一個.
頭尾呼應的一個文學結構.
就是管理百姓.
第十一節和十四十五節.
就是摩西那四個問題擺在這裡.
向神抱怨.
我一個人搞不定.
這些管理太重了.
但上帝竟然要求摩西的就是餡.
就是中間十二十三節.
帶領百姓和要慰養百姓.
我們做教會,做管理,做帶領.
是一個屬靈的領袖.
就不只是那些行政管理.
最重要是要將神的子民,神的百姓.
他們屬靈的心態.
是要牧養,慰養他們.
讓他們懂得怎樣去敬畏神.
這個心態教得正的話.
其實其他任何事務性.
各樣的挑戰.
那些麻煩事是迎刃而解的.
但那些沒有解決的時候.
每個人裡頭的內心.
都是自以為是,自以為王.
各人偏行幾路的話.
這個很可怕.
這樣就會出大問題.
我們看見神怎樣回應摩西的不滿.
十六,十七節.
神對摩西說很簡單的.
就是這個解決方案.
祂說你在長老當中召集七十個人.
這班長老在《出埃及記》二十四章.
神已經吩咐摩西揀選了七十個長老.
他們曾經上過山.

$^{521}$他們見證過神與摩西所立的約.
但是這裡說是另外七十個人.
當中可能有重複.
但不是全部都是之前的七十個長老.
有部分是.
祂說你要找一些你認識.
作百姓的長老和官長的.
帶他們來到會幕.
叫他們和你一起站在那裡.
我就要在那裡降臨和你說話.
我要將降給摩西你的靈分給他們.
他們就和你分擔帶領百姓的擔子.
免得你獨自承擔.
一語中的直接去解決摩西的問題.
我一個人搞不定.
我怎麼可能去承擔.
這百姓有多少?.
只是男子也有六十萬.
他們說整個會眾有二百多萬人.
這麼龐大的一個數目.
但是我們看到神願意賜下同功.
叫摩西你得著這個重擔.
可以卸下分擔.
當摩西執行神的吩咐的時候.
你發現神就在這個會幕降臨.
然後他的靈就確是分在這七十個長老.
停在他們身上的時候.
他們就說預言.
之後他們沒有再說.
這個說預言表明.
這一班長老是神所引證的.
是那個領袖.
是去分擔摩西管理的責任.
不過這裡有一個小插曲.
就是在這個片段裡面.
七十個長老有兩個人就留在營中.
沒有去到會幕.
就是伊利達和米達.
他們說預言與其他會幕的六十八位長老是相同的.
約書亞在這個時候知道這件事.

$^{561}$他就很緊張.
他為摩西的緣故就嫉妒他們.
為什麼?.
他說怎麼可以這樣呢?.
如果在會幕以外.
也有人可以得到在會幕裡這種與神相遇.
得到的祝福的話.
這樣會不會這兩個人在搞小圈子.
他們要僭奪摩西的領導權呢?.
所以約書亞想為摩西出頭.
約書亞是摩西的副手.
他是很明白摩西這個領袖為人極其謙卑.
但摩西你看到夏文被我們看見.
他的心胸頗大.
摩西說我寧願百姓都是神的先知.
是神的邊先知.
摩西期望百姓的身份.
是建立在神的靈的上面.
要真正作先知的.
不僅是懂得自理.
他們要明白上帝的心意.
他們本身的生命是經過神的靈的充滿.
對付,拆毀,重建.
不是照我們的心意去作每一件事.
是願意我們去信服在神的面前.
我們和我們的領袖.
和我們的同工一起去面對.
帶領百姓這個重任.
絕對不簡單.
所以我們看見如果用今天現代的說法.
教會不僅是執事,傳道,牧師.
我們是需要全民皆兵.
而且是全民皆將.
有兵有將.
因為我們是要打一場屬靈的爭戰.
我們要領人歸主.
要建立弟兄姊妹作主的真門徒的話.
整件事是魔鬼絕對不喜歡.
魔鬼要攻擊我們的家庭.
要在我們的事業,生意.

$^{601}$要我們接受很多試煉,考驗.
在我們的兒女或配偶家庭身上.
可能都會出現亂子.
但在這些情況下.
讓我們不要懼怕.
我們要知道我們要作真門徒的時候.
你就要被聖靈去充滿.
每一天你真的去讀聖經.
而且你去祈禱.
我們過去一年多了 今年十一月畢業了.
教牧和一些執事上門了.
我們大家都非常興奮和開心.
開心的原因是什麼?.
因為我們知道每一天神的話.
都對我們有說話.
我們將這些領袖寫聖靈修的札記.
去跟我們的小組弟兄姊妹一起分享.
原來神就是這樣.
今早我自己看的經文就是.
神跟亞伯拉罕說.
「明日這個時候.
你的妻子就懷孕了」.
我看完之後.
「明日這個時候.
我的妻子才沒有懷孕?」.
你讀聖經有時候不要靈異化.
或者騙語地讀.
但這句話精彩的地方就是.
這時候是神命定的時候.
是上帝命定的旨意.
祂所喜悅的計劃和旨意.
一定會成就的.
不用你擔心.
所以當你知道.
上帝的心意是要你全家歸主的.
問題就在這裡.
你有沒有為你未信的家人緊張過.
半點時候.
你有沒有為你還未信主的同事.
日夜相對的.

$^{641}$一個同事你緊張過.
你為他祈禱.
每個月教會基本上都有福音崇拜.
你試過作出邀請沒有?.
這對我們來說.
就是我們必須生命要被破碎.
我們不要停留只是.
很喜歡讀經.
很喜歡查經.
很喜歡聽聖經.
很喜歡這個講道.
那個講道.
很喜歡在這裡下午茶.
我們在這裡茶餘飯後.
在這裡評論.
人家正評.
我們就在這裡評論.
這個牧師講的那句很好.
那個牧師講的那句很精彩.
只會講.
但你不會行的話.
弟兄姊妹.
我們仍然停留在這個西乃山下面.
摩西和張祿返回營中.
他們的不滿得到完全的解決.
我們感恩.
神怎樣去回應百姓的問題呢?.
百姓表層的原因是他們缺乏肉食.
但心情的原因是他們貪肉.
他們不信靠神.
當神供應給這些肉他們吃.
吃到他們一個地步.
說是鼻孔噴出來.
你們那些覺得麻辣好像不夠和味.
現在給你們吃足一個月.
他們又吃到厭了.
在嫌棄神的供應.
其實在二十節就講出那個精髓的地方.
他們其實是嫌棄住在他們中間的耶和華.
他們的神.

$^{681}$二十節百姓在哭訴.
我們為何出埃及.
這個就是問題的所在.
他們還有人在埋怨.
埋怨什麼?.
沒有水喝,沒有東西吃.
其實背後就是.
為何我們要出來抗疫?.
為何我們不要留在埃及?.
是這件事.
摩西在這個情況.
他壓力爆煲之餘.
他也懷疑神是否可以供應這麼多人.
吃一個月的肉.
但神說.
我的膀胱喜事縮短.
現在你就要看我向你應驗不應驗.
二十三節所講.
我們看解決了摩西的問題之後.
神還解決了百姓的問題.
就是他們吃東西的問題.
他透過一陣風.
將海邊很多的鵪鶉.
刮來散落在營地的周圍.
很誇張.
請你留意.
他講的這一陣風.
在《原文希伯來文》裡面.
其實和神的靈在摩西身上的「靈」.
「Rua」這個字.
是一模一樣的.
神解決我們人的問題.
永遠是用神的靈去解決.
今天我們人很多時候.
只是看見那些問題很表層.
環境上,經濟上,心靈上.
但真正叫一個人迷失.
是他的心靈離開了上帝.
他以為用自己的聰明智慧.
以為用自己的謀略.

$^{721}$他是叫神走開一邊.
我來做主人.
這個才是問題的癥結.
當神用他的靈去解決的時候.
你要吃.
我用我的靈.
一呼.
真是千萬應.
你看看這裡.
很誇張.
究竟有多少鵪鶉來了.
這個形容很誇張.
在31至33節.
他講到一邊在這個營的一邊.
有一天的路程.
另一邊也有一天的路程.
所有這些鵪鶉被大風刮過來.
就停在那裡.
這些鵪鶉我查究過.
原來它們不是擅長飛翔的鳥.
它是擅長在地上走的.
走在那裡蹬地上的穀物.
大家有沒有人吃過鵪鶉?.
鵪鶉其實營養很豐富.
蛋白質很高.
其實這些肉很香口.
但是在這些情況下.
你說的每邊有多長?.
是說13至14公里這麼長.
嘩.
然後有多高?.
嘩 不是一層.
是兩層.
有三尺.
這裡說是二爪.
地面又有二爪.
為什麼說它們重疊降臨?.
一層也不夠.
兩層也不夠.
三層也不夠.

$^{761}$三尺高.
三英尺.
一爪就是一尺半.
由我手肘到中指的位置.
一尺半.
兩爪就是三尺.
所以整件事很誇張.
如果你要拍電影.
我不知道怎麼拍.
但是難怪.
這些百姓一看到怎麼辦?.
平時撿那些碎碎的嗎哪嗎哪.
現在一隻隻這些鵪鶉.
嘩 太好吃了.
太好吃了.
拿了拿了.
他說每人至少捉到10荷梅爾.
荷梅爾是一個容量.
1荷梅爾是說多少呢?.
1荷梅爾是說220公升.
所以10荷梅爾是說二千多公升.
二千多公升是多少呢?.
不知道怎麼形容給你聽.
不過大家試過入那些汽油.
如果找那些罐.
那些起碼都5公升.
或者再大一點可能10公升.
如果10公升的話.
你說400多個這些罐.
是再那麼多的鵪鶉.
擺在你整個車房都滿了都不止.
問題在這裡.
你一個人能吃多少.
你吃不完會怎麼樣.
臭的了.
這個情況就是這樣.
因為他們貪.
只是貪.
嘩 這麼多東西.
不花錢的.

$^{801}$拿啊拿啊拿啊.
拿了 囤積了這些.
結果其實是病毒在當中發生.
他們還在這裡吃完.
還沒消化的時候.
還沒咀嚼的時候.
爺爺說他老氣向他們發作.
用極重的災禍擊殺這些百姓.
我最後完結的時候.
要大家留意.
經文給了我們一些很重要的線索.
神處理百姓和摩西的問題.
全部都是有根有據.
當怨言和貪欲.
神怎麼樣呢?.
神是用災禍去擊打.
災禍在原文的聖經.
其實就是一種.
好像重拳的一個追擊.
當年神用十災.
就是不容許那些以色列人離開埃及人.
這種災禍擊打他們.
現在神同樣用災禍去擊打那些.
自認不要離開埃及的以色列人.
神是要管教他們.
神要審判懲罰他們.
而摩西他透過招集那七十個人.
來分擔他們管理的問題.
百姓透過同一個動作.
去捕捉那些暗巡.
捕捉這個原文.
和那個招集是一個字.
是同樣的字.
神是要解決他們食慾的問題.
但是要承受神的審判.
所以這個地方叫做「基博羅哈他雅」.
是兩個希伯來文.
「哈他雅」就是貪欲者的意思.
貪心的人和他第四節講到.
他們被那些色人所煽動的這些貪欲.

$^{841}$是有相同的字根.
其實是首尾呼應.
去總結神對百姓怨言的審判.
所以神審判他們.
是因為他們沒有尊重神的帶領.
和他們貪念神引導摩西.
這七十個人.
對於我們今天要學的三個功課.
我們不要抱怨神的引導.
我們要數算神的恩典.
要禱告.
我們不要放縱貪欲.
我們要為神今天每天給予我們的晚餐.
我們要感恩.
我們不要孤單事奉.
要讓神透過團隊去分擔我們的重擔.
我們很感恩.
我上週末已經在這裡和大家分享了.
2027年1月1日.
陳全華牧師會成為我們繼任的一個新的堂主任.
我和他很相熟.
我們這幾年在總會都有很多機會一起事奉.
我也很欣賞他帶領這個門女.
其實實在是對我們教會一個很大的祝福.
我很盼望我們每一個弟兄姊妹.
和我們的執事以及傳道同工.
一起和祂同心.
去面對新的下一里路.
神要帶領祝福,講福.
要建立我們不單止在這個地區.
我很深信神會透過我們很多弟兄姊妹.
你們生命的栽建和培育.
是要成為更多在不同職業界別裡面.
其他甚至微信的朋友的祝福.
你做生意的.
你有很多的你的客人.
可以透過你的生命的見證去祝福他們.
所以在結束的時候.
希望大家看到.
我們看到百姓不是因為缺水缺糧.

$^{881}$他們不斷發怨言.
其實和貪這些玉石.
是因為厭棄神的帶領和恭迎.
但是神很體貼摩西和百姓的軟弱.
祂主動和他們興起同工,分擔他們.
牧養我們知道.
教牧和長執是必須要同心同工.
我們要分擔自理和牧養群養這個重責.
弟兄姊妹請你為我們禱告.
我已經和陳牧師有約了時間.
我們在這半年裡面.
我們會通過很多的一起協商.
我會將很多事情和他交代.
我們盼望能夠讓我們整個過程都非常順暢.
最後.
當人只追求玉石.
我們的眼目離開這個回望的上帝.
我們的貪欲叫我們從敬拜之地.
走進墳墓之谷.
盼望弟兄姊妹你們不是成為這一種的回眾.
你們要成為一個真門徒.
去跟隨主.
得神喜悅的回眾.
我們一起祈禱.
天父我們感恩.
我們獻上禱告.
求主賜恩.
大大賜福給廣福堂.
給我們眾執事.
眾教牧同工.
我們帶領弟兄姊妹.
一起照著神的美意.
我們建立你所喜悅的門徒.
一起叫我們前行的日子.
能夠榮神益人.
領人歸主.
也建立更多能夠高舉榮耀上帝之名的.
你自己的兒女.
讓我們在這個社會上.
在我們的家庭上.

$^{921}$在各樣的公職上.
我們都能夠貢獻國家.
貢獻社會.
在教會裡外見證.
都能夠榮神益人.
同心合意祈禱.
奉主名求.
阿們.
\newpage



\section{約書亞記 23:16}
\label{sec:032204}
\textbf{聖經人聖經事:約書亞贈別 (2) 勿事奉別神}~~~~黃克勤牧師博士
\newline
\newline
連結: \href{http://www.church.com.hk/acms/content.asp?site=cdc\&op=show\&type=product\&code=032204\&layout=sermon}{\texttt{032204}} ~~~~ 語音日期: 2026-05-27
\newline
\newline
\hyperref[sec:032203]{\small{< < < PREV SERMON < < <}}
~
\hyperref[sec:index_chronic]{\small{[返順時目]}}
~
\hyperref[sec:032205]{\small{> > > NEXT SERMON > > >}}
\newline
\newline
約書亞記 23:16
\newline
\begin{longtable}{cl}
\hline
\hline
章節 & 經文 (和合本修訂版)\\
\hline
23:16 & \begin{tabularx}{0.7\textwidth}{X} 你們若違背耶和華－你們神吩咐你們所守的約，去事奉別神，敬拜它們，耶和華的怒氣必向你們發作，使你們在他所賜給你們的美地上迅速滅亡。」 \end{tabularx} \\ \\
[1ex]
\hline
\hline
\end{longtable}
$^{1}$各位電影師妹.
又一齊看看約書亞怎麼說了.
《聖經人》仍然是約書亞.
《聖經的事》就是約書亞的臨別贈言.
贈別的第二.
今天《生命說書家》網絡跟大家說一下這件事.
就是臨別贈言第二點.
就是物是奉別神.
這個物是奉別神.
不是直接經文來的.
是我改進去.
特別是這麼短的位置擺五個字.
物是奉別神.
是綜合了今天要說的話.
今天說的話比較短的.
那麼大家留意到.
就是一說到約書亞記.
次字都很長.
有半個小時.
哇!這麼長.
真的有些東西聽到很悶.
這次短一點.
是因為件事很長.
經文很長.
用的時間比較多.
好了 約書亞記臨別贈言.
就是第二.
物是奉別神.
約書亞記來到二十三章了.
是不是差不多了.
成卷約書亞記看完了.
可能我們平時少看一點的.
這次看完了.
厲害.
十四節.
「我現在要走世人必走的路」.
「你們是一心一意地知道」.
「耶和華你們神所應許賜福與你們的話」.
「沒有一句落空」.
「都應驗在你們身上了」.

$^{41}$約書亞說話雖然他年紀大.
都很厲害的.
其實做一個屬靈領袖.
能夠能言善道.
都是一個很重要的.
做領袖的條件.
好了.
約書亞就說了.
「我現在要走世人必走的路」.
那是什麼意思.
就是我快要死了.
我很老了.
就是個個都會走這條路了.
我的以前的老師.
又或者我以前的上司.
約摩西就走了.
摩西的哥哥阿倫又走了.
現在到我帶大家幾十年.
到現在分好地.
我都差不多要走了.
「走世人必走的路」.
看節目你都要記得.
你和我都會走世人必走的路的.
我都覺得我很快要走這條路了.
我都要走的了.
沒得抗拒的了.
於是問題就是.
預備好未呢.
回望的時候.
這一生人做了多少事.
這條路走了.
有沒有走了冤枉路呢.
有沒有浪費了時間呢.
有沒有一些事情還沒做呢.
現在我要走世人必走的路.
而你們還沒走完的.
還沒走完就要留意了.
還沒走完就要做什麼.
你要回想一心一意地知道.
其實是一個回想來的.

$^{81}$就是耶和華你們的神.
所應許賜福你們的說話.
是沒有一句落空.
所有應許都是確實是成就了的.
跟爸爸在香港.
我要將一個兒子給你.
我要將很多子孫給你.
我要將一塊地給你.
全部應驗.
跟著誰知道去了埃及做奴隸.
四百年是神預言的.
又應驗.
四百年是要出來又應驗.
於是就帶他們出來了.
就差遣僕人摩西帶他們出來.
跟著經過千辛萬苦.
進到應許之地.
又全部所講的都應驗.
然後分地給他們.
應驗了每人都有地方安居樂業.
昨天說的.
就四景平靜神讓他們休息了很多的日子.
現在約書亞就說.
我的時間都差不多了.
不能夠再陪大家一起走了.
那這條路我會先走一步.
那你們繼續.
記得 你知道嗎.
神對你們的說話.
全部沒有一句落空.
都應驗在你們身上.
所以這位神是確實的.
耶和華以色列的神是確實的.
好了 不過耶和華你們的神.
應許的一切福氣怎樣臨到你身上.
耶和華也可以照樣使個人的禍患臨到你們身上.
直到將你們從耶和華你們神所賜的那美地上.
除滅.
這句說話很重的.
第一樣要留意的.

$^{121}$好東西 我們含著甘索時出世.
又一來到就不用做什麼.
只是這樣哀哀 城倒了.
我們得了耶利哥城.
對付這個皇帝.
去打他不是皇帝.
這個皇 這個土皇的時候.
打他怎樣都不用打.
神差遣黃蜂將他們全部走掉.
全部死掉.
死的人比我們用刀去殺更多.
哇 落薄.
哎呀 打死所有人.
打死了 又贏了一個城.
做了很多事其實是神做的.
我們確實知道.
若說話說你記得了.
所有說話都應驗的.
神的祝福不要落空.
句句應驗.
你們知道了.
好了 不過.
第一點我臨走的時候臨滅鎮賢.
就是要知道福氣怎樣臨到你.
也可以照樣離開你的話.
就是禍患可以臨到你.
福氣臨到照樣禍患都可以臨到.
如果你們離開耶和華神的話.
祂甚至可以將你們在所次的這美地.
除滅.
不要以為自己.
我是上帝的百姓.
厲害了 大了.
金鑰匙含著幾條.
全得了 還吃別人自己母栽種的葡萄.
吃自己母栽種的果子.
嘩 多好啊.
神怎樣祝福你.
是可以照樣叫禍患臨到你.
是約書下的第一個.

$^{161}$是第二個臨滅鎮賢的第一點.
好了 第二個就是.
若違背耶和華你們神吩咐你的約.
去事奉別神.
神為何要和我們以色列人這麼好.
我們是祂的子民.
為何要你是子民.
因為亞伯拉罕的緣故.
所以你含著金鑰匙.
但其實為何要和亞伯拉罕特別好.
因為和亞伯拉罕有約.
這個約就可以了.
有約嘛 純手約的.
是 手約斯詞外.
不過去到西乃山.
大家不要忘記了.
神又和我們再納了一個約.
就是給了戒命我們.
那個就是戒命的約.
就是耶和華神和我們的關係.
就因為這個西乃山之約.
如果你違背耶和華這個西乃山的約.
去事奉別神.
除外意外不可以有別的神.
不用有別的神.
你又去跟別的神.
你就會違背耶和華的約.
第二點 記住不要違背這個律法的約.
第三點.
如果這樣怒氣會向你們發作.
你們就在他所賜的美地上.
速速滅亡.
上一段說在這個美地上除滅.
沒有 抹了它.
這裡說在美地上.
速速滅亡.
都會這樣的嗎.
他不是給了我們嗎.
是呀 是給了.
但是你要守約.

$^{201}$千萬不要去事奉別神.
叩拜那些神.
如果你去事奉別神.
叩拜那些神.
怒氣就會向你發作.
這個怒氣這個字.
通常希伯來文就是鼻子.
鼻 nose.
那為什麼要鼻呢.
怒氣就是嗯.
即是生氣.
那些氣會噴出來.
耶和華的鼻就會向你發作.
那你就會在他所賜的美地.
速速滅亡.
不是慢慢的.
是速速很快就沒有.
上面那一點.
是說在美地上將你除滅.
這個是說有與沒有.
將你沒有了.
這裡是說速度.
速速地沒有了.
是很快.
不是慢慢借.
不是陰乾.
是可以很快沒有了.
因此記得.
就是不要違背耶和華的約.
以免他的怒氣.
他的鼻向你發作.
最重要就是記得.
福氣怎樣來到.
福氣可以離開.
禍患可以臨到.
是在乎你有沒有手藥.
以至整個怒氣會發作.
這時你會看到這裡.
就是約書亞和他們臨別贈言.
就是約書亞說.

$^{241}$我要走世人必走的路.
現在老了.
就不再陪到你們了.
即是說很多事情我不再跟你們說了.
以前我就會.
慢著慢著.
大家要留意.
經常提醒你.
現在我走了.
可能沒有人提醒你了.
或者有其他人提醒.
不過無論如何.
我作為領袖與父親.
我要提醒大家.
千萬不要.
記得.
這個就是約書亞的臨別贈言.
第二.
記不記得第一是什麼呢?.
第一就是告訴他們.
你們要愛耶和華.
只是昨天說的.
愛耶和華.
怎樣愛呢?.
愛耶和華就是.
不要跟本地剩下的人殘雜.
不要混在一起.
不要結婚.
不要嫁娶.
不要有太多來往.
以免他影響了你.
你愛耶和華就要守他的誡命.
歡喜讀他的命令.
你經常說.
誡命這麼難守.
校規這麼好.
如果你用愛來做.
就是我很愛我這個家庭.
媽媽爸爸說的話.
我要跟著.

$^{281}$這是愛的表達.
要做這些.
要做那些.
又要我做功課.
其實他賺錢很辛苦.
就是想你做功課.
將來可以有更好的前途.
你要用愛來欣賞.
就很喜歡做功課.
要整個限制.
我要這樣.
要這樣.
不要吃這個.
不要吃那樣.
這樣.
其實他不想你吃那麼多糖.
將來牙痛而已.
到牙痛的時候就明白了.
我小時候就這樣.
兩隻大牙.
一個洞在旁邊.
明明是說不要吃那麼多糖.
吃了就含著睡.
發個美夢.
下去之後牙齒都麻了.
哎呀.
好慘.
痛了.
牙痛慘不過大病.
經歷過很多.
小心.
不要亂玩.
在門口小心.
不就夾住手指.
一夾住手指.
嘩.
痛到眼淚直接飆出來.
說了給你聽的.
愛爺爺.
守他的誡命.

$^{321}$愛你的家.
愛你的父母.
是守他的說話.
這個臨別贈言第二.
第二就是.
不要事奉別臣.
就是今天你和大家說的.
剛剛說了.
你要知道.
就是.
臣是可以怎樣祝福.
也可以怎樣.
比禍還.
而且在所給你的恩典當中.
速速滅亡.
這個是給他們的臨別贈言.
也警惕了我們的生活.
我們同心祈禱.
天父上帝多謝你.
你就是那麼守施慈愛.
我們在約書亞的身上.
他臨別贈言.
他提到.
耶和華臣你給以色列人的.
每一句應許.
是沒有一句落空.
哦.
這個真是一個很大的恩惠.
很大的恩典.
叫我們很感恩.
當約書亞提醒我們.
人生都是有一天會走完的.
要在人生裡面.
除了要愛耶和華.
我們要進行你的革命.
免得你可以祝福.
你又可以比禍還我們.
而且在我們的福氣當中.
速速滅亡.
幫助我們更多跟隨你.

$^{361}$遵從你.
愛你.
祈禱奉靠耶穌基督的名.
阿們.
\newpage



\section{歌羅西書 1:1-14}
\label{sec:032205}
\textbf{在基督裡}~~~~鄭卓雄牧師
\newline
\newline
連結: \href{http://www.church.com.hk/acms/content.asp?site=cdc\&op=show\&type=product\&code=032205\&layout=sermon}{\texttt{032205}} ~~~~ 語音日期: 2026-05-24
\newline
\newline
\hyperref[sec:032204]{\small{< < < PREV SERMON < < <}}
~
\hyperref[sec:index_chronic]{\small{[返順時目]}}
~
\hyperref[sec:032206]{\small{> > > NEXT SERMON > > >}}
\newline
\newline
歌羅西書 1:1-14
\newline
\begin{longtable}{cl}
\hline
\hline
章節 & 經文 (和合本修訂版)\\
\hline
1:1 & \begin{tabularx}{0.7\textwidth}{X} 奉神旨意，作基督耶穌使徒的保羅，和我們的弟兄提摩太， \end{tabularx} \\ \\ \relax
1:2 & \begin{tabularx}{0.7\textwidth}{X} 寫信給歌羅西的聖徒，在基督裡忠心的弟兄。願恩惠、平安從我們的父神歸給你們！ \end{tabularx} \\ \\ \relax
1:3 & \begin{tabularx}{0.7\textwidth}{X} 我們為你們禱告的時候，常常感謝我們主耶穌基督的父神， \end{tabularx} \\ \\ \relax
1:4 & \begin{tabularx}{0.7\textwidth}{X} 因為聽見你們對基督耶穌的信心，並對眾聖徒有的愛心。 \end{tabularx} \\ \\ \relax
1:5 & \begin{tabularx}{0.7\textwidth}{X} 這都是因著那給你們存在天上的盼望，它就是你們從前所聽見真理的道，就是福音； \end{tabularx} \\ \\ \relax
1:6 & \begin{tabularx}{0.7\textwidth}{X} 這福音傳到你們那裡，也傳到普天下，並且繼續增長，不斷結果，正如自從你們聽見福音，真正知道神恩惠的日子起，在你們中間也是這樣。 \end{tabularx} \\ \\ \relax
1:7 & \begin{tabularx}{0.7\textwidth}{X} 這福音是你們從我們所親愛、一同作僕人的以巴弗學到的。他為我們作了基督的忠心僕役， \end{tabularx} \\ \\ \relax
1:8 & \begin{tabularx}{0.7\textwidth}{X} 也把聖靈賜給你們的愛告訴我們。 \end{tabularx} \\ \\ \relax
1:9 & \begin{tabularx}{0.7\textwidth}{X} 因此，我們自從聽見的日子就不住地為你們禱告和祈求，願你們滿有一切屬靈的智慧和悟性，真正知道神的旨意， \end{tabularx} \\ \\ \relax
1:10 & \begin{tabularx}{0.7\textwidth}{X} 好使你們行事為人對得起主，凡事蒙他喜悅，在一切善事上結果子，對神的認識更有長進。 \end{tabularx} \\ \\ \relax
1:11 & \begin{tabularx}{0.7\textwidth}{X} 願你們從他榮耀的權能中，得以在一切事上力上加力，好使你們凡事歡歡喜喜地忍耐寬容， \end{tabularx} \\ \\ \relax
1:12 & \begin{tabularx}{0.7\textwidth}{X} 又感謝父，使你們配與眾聖徒在光明中分享基業。 \end{tabularx} \\ \\ \relax
1:13 & \begin{tabularx}{0.7\textwidth}{X} 他救了我們脫離黑暗的權勢，遷移到他愛子的國度裡。 \end{tabularx} \\ \\ \relax
1:14 & \begin{tabularx}{0.7\textwidth}{X} 藉著他的愛子，我們得蒙救贖，罪得赦免。 \end{tabularx} \\ \\
[1ex]
\hline
\hline
\end{longtable}
$^{1}$特別是禹堂的弟兄姊妹們一起敬拜,事奉神.
T4C的中文名就是「雙城華人基督教會」.
我將它理解為「雙城華人基督的教會」.
這是一個主義的教會 一個屬於基督的教會.
是由一群在基督裡的人所造成.
在明州雙城這個地方事奉神.
進行基督給我們的使命 服侍華人社區.
在這個社群當中傳揚基督 領人歸住.
為了要達到神對我們的心意.
我希望在接下來的一段時間裡.
和大家一起從哥歌羅西書.
我們學習怎樣去宣揚基督的超越性.
和祂的充足性 祂的supremacy 祂的sufficiency.
講到基督的超越就是說 耶穌基督在萬有之上.
超越一切物質界和靈界的事物.
又加力量給在基督裡的人.
包括在座的每一位弟兄姊妹 每一個基督徒.
使你和我能夠把各人在基督裡.
完完全全的引導神面前.
求神使用歌羅西書的信息.
教導你和我怎樣體會神賜給信耶穌的人.
在基督裡的身份和豐盛的生命.
保佑你和我能夠活出.
我已經與基督同釘十字架.
現在活著的不再是我.
乃是基督在我裡面活著.
我在基督裡和基督在我裡面成為共存的屬靈實際.
今天我們要開始思想.
基督的福音怎樣帶給你和我.
在基督裡的信,望,愛.
求神眺望我們愛主的心.
激勵你和我為主而活.
活出我們是基督的僕人.
因為我們都是主耶穌用重價買來的.
神要使用祂的僕人建立基督的國.
就把我們帶到這個基督的國度裡面.
作神國的子民 做基督的僕人.
求聖靈幫助我們.
明白這些是真理.
懂得怎樣配合神的作為.

$^{41}$遵行神的旨意.
歌羅西書是一封信.
按照當時代的格式.
寫信人首先是自我介紹.
奉神旨意作基督耶穌使徒的.
保羅和兄弟提摩泰.
保羅講到他自己是有耶穌基督使徒的身份.
而這個身份是神按照自己的旨意.
神揀選了保羅做基督的使徒.
保羅和提摩泰寫信給在哥羅西的聖徒.
在基督裡聖潔中心的弟兄們.
講明這些收信人是住在哥羅西的聖徒.
聖徒就是被神分別為聖 歸給神的.
預備要去做神所要用他們做的事情的僕人.
哥羅西是位於小亞細亞.
就是現在的土耳其西邊的一個城市.
在寫信的時候 保羅是在羅馬的監獄裡面.
他從來沒有去過哥羅西.
他不認識這些收信人.
但是他很親切的問候他們 祝福他們.
願因為平安 從實我們的父 歸於你們.
這一群已經在基督裡的人 他們是對主忠心.
保羅為他們向神求的恩惠 不是得救的恩典.
因為這幫都是已經蒙恩得救 在基督裡的人.
但是叫他們能夠承受神給他們的每日此刻.
新的恩典 可以持續忠於所信的福音.
逐步的來到經歷傳遞的救恩.
享受從神來的平安.
是超越平安無事 無憂驚那種的平靜.
無事無幹那種的平安.
更加是我們活在神的同在裡面.
得著從神而來的平靜,安穩.
神透過善知義錯下提醒我們.
我們基督徒生活能夠有得力的秘訣在於平靜,安穩.
因為我們得救在乎歸回安息.
我們得力在乎平靜,安穩.
看你和我都能夠在基督裡經歷到.
祂能夠使我們完全的滿足 有真正的穩妥.
讓我們一起低頭去仰望神.
天父我們感恩你自己在基督裡面.

$^{81}$揀選了我們 拯救了我們.
讓我們能夠來到你自己的根前.
獻上我們的敬拜,我們的讚美.
因為你實在是那位超越一切的神.
也是那位可以讓我們得到真正滿足的神.
求神幫助孩子,幫助眾弟,卿姐妹.
在敬拜當中 將我們的眼目定精在你的身上.
我們單單的來去仰望你.
單單的求主你自己向我們施恩.
接受聖靈 將神的話語解釋清楚明白.
補足人的不足.
以致你自己的話語是沒有爛掉的.
成為我們生命中的糧 成為我們心中的滿足.
求主施恩 聽孩子的禱告 奉耶穌基督的名.
保羅告訴哥羅西的信徒們.
我們為你們禱告的時候.
他的同工為了這些在基督裡的弟兄姊妹祈禱的時候.
保羅感恩的對象是父神.
他感恩的內容是神子在信徒生命中的作為.
因為他聽見你們在基督耶穌裡的信心.
並向眾信徒的愛心.
是為了給你們存在天上的盼望.
保羅從以巴弗聽見了在哥羅西信徒的生活.
他們可以活出在基督耶穌裡的信,望,愛.
而且指出他們有這樣的生命的表現.
根源就是他們在基督耶穌裡的信心.
他們憑信心接受神差到這個世上.
做救主的彌賽亞就是耶穌.
就是那位為我們信心創始成終的耶穌.
在祂裡面我們找到神所賜的信心.
基督耶穌不單止是我們信靠的對象.
祂也是你和我 我們信心得到成全.
我們所信的得以實現的根據.
在基督耶穌裡面的信心.
其中一個的表現就是對眾信徒持有的愛心.
這是一個超越情感的愛.
或者親情的愛的Fileo那一種的愛情.
乃是對所有被神分別出來,歸於神的人.
生發出好像基督那樣無條件神聖的愛.
只會在這種的愛加倍 Love.

$^{121}$只會在同有神的生命.
在基督裡的群體當中會表明出來.
是跟隨主的人必須追求的一個生活表現.
因為主耶穌曾經對他的門徒說.
我賜給你們一條新命令.
乃是叫你們彼此相愛.
我怎樣愛你們 你們也要怎樣相愛.
主耶穌不單止吩咐我們在基督裡的人.
我們要彼此相愛.
很清楚的告訴我們我們要怎樣彼此相愛.
就好像祂愛我們一樣.
基督是怎樣愛我們?基督為我們寫命.
我們從此就知道何謂愛.
求神幫助我們也能夠實踐.
能夠活出為弟兄寫命的愛.
基督徒可以信,能夠愛.
都是由於那給你們存在天上的盼望.
這盼望是你們從前在福音真理的道上聽過的.
就是說信和愛兩者都是基於.
那個給我們保存在天上的盼望.
基督徒的信仰 你和我的信仰.
是本於信而治於信.
我們不單止是靠信心入門.
我們既因信清義.
就藉著我們主耶穌基督得與神相和.
重建與神和好的關係.
就要持續的對神有信心.
因為二人必因信得心.
所以神透過使徒保羅指示信徒.
只要你們在所信的道上恆心.
基督徒能夠相信到底.
是基於神給我們保存在天上的盼望.
這個盼望神已經給了我們.
現在是暫時保留在天上.
等到最後的實現.
聖經說的盼望不是有一個樂觀的心態.
總是覺得明天會更好.
也都不是人自己想像出來的想望.
我們很多時候有一些wishful thinking.
但我們知道不一定會實現.

$^{161}$但聖經所說的盼望是植根於.
神的話語和祂的應許.
加上我們知道神是信實無變.
從以色列人的歷史裡面我們看到.
也有說神所應許給他們的.
似福給以色列人的話.
沒有一句落空.
都應驗在他們的身上.
我們只能把握我們基於聖經的盼望.
是必然會實現的 是不會落空的.
這盼望是你們從前在福音真理的道上聽過的.
是信徒曾經聽過的.
包含在基督的福音裡面的真理之道.
不是外交 不是附帶.
是關乎耶穌基督的福音的基要核心的真理.
基督的福音這個好消息.
就是神賜給人一個全被的救恩.
耶穌不單止是洗淨了我們過去的罪.
撫清你和我的罪債.
也給了信徒對未來的盼望.
是有永恆意義的盼望.
這是我們持守相信耶穌的動力.
會影響我們現今生活的盼望.
是一個活潑的 是一個Living Hope.
是一個活潑的盼望.
這福音到了你們那裡.
正如到了全世界 結果只增長著.
正如在你們中間一樣 結果只增長.
自從你們聽見了.
並實實在在認識神恩惠的日子以來.
都是如此.
這福音到了你們那裡.
是一個現在時態的分支.
Present Participle.
表示這福音仍然在他們中間.
仍然繼續產生著作用.
正如在全世界各地一樣.
無論他們從誰的口裡聽到福音.
都會產生同樣的效應.
因為這福音本是神的大能.

$^{201}$要救一切相信的.
福音本身就有能力能夠去救人.
能夠使所有相信基督的福音的人.
都可以得救.
無論在什麼地方.
基督的福音同樣都是在繼續結果子.
聖靈內在相信福音接受耶穌的人裡面.
使信徒的生命不斷地被更新.
結出聖聖的果子.
並且福音被增長.
是一個被動式的動詞.
在世界各地都被傳開.
令更多的人歸順耶穌.
主張得救的人天天加給他們.
是神自己的作為.
我們是承受神的恩典.
成為祂所使用的器皿.
將得救的人天天加到神的家中.
造成信徒群體在質量和數量上都有增進.
福音的能力和果效.
是自從你們聽見了並實實在在認識神恩惠的日子以來.
都是如此.
這福音就是在你們中間不斷地結果子和增長.
關鍵在於他們不單聽見了基督的福音.
而且他又用信心與所聽見的打成一片.
聽見和聽從是連結在一起.
信心和行為是並行是並重.
就會因著真理確實地認識了神的恩典.
真知道神傳遞的恩惠.
神子耶穌是基督福音的中心.
因為恩典和真理都是由耶穌基督來的.
在座還未信耶穌的朋友.
假如你想聚德赦免與神和好.
你渴慕成為神的兒女.
就要聽就要信福音.
因為信道是從聽道來的.
聽道是從基督的話來的.
弟兄姊妹 你和我如果渴望基督的福音.
在我們中間不斷地結果和增長.
就要更深入更全面的認識耶穌基督和祂的福音.

$^{241}$並且要有順從的行動.
與我們所聽見的福音真理連接起來.
才會產生這些美好的果效.
人要聽信基督的福音.
但是未曾聽見祂怎能信祂呢?.
沒有傳道的怎能聽見呢?.
本來指出在哥羅西的基督徒.
是從以巴弗聽見福音.
這福音也就是你們從我們所親愛的.
一同作僕人的以巴弗那裡學到的.
祂為你們作了基督中心的僕役.
也把你們在聖靈裡的愛心告訴了我們.
保羅介紹以巴弗是我們所親愛的一同作僕人的.
是和保羅和提摩太堅並肩一起做奴隸的.
有一個譯文更貼切的翻譯是一起做奴隸.
是奉派做僕役.
僕役就是講到那個詞的意思是.
祂有一個特別的職份 有特別的任務.
去講傳基督的福音.
以巴弗是一位來自哥羅西的基督徒.
曾經告訴我們 當保羅在以忽所.
我們知道他在那裡曾經有差不多三年的時間.
或者三年以上的時間.
他的同工在鄰近的城市在那裡傳福音.
兩年之久叫他一起住在亞西亞.
無論是猶太人,希臘人都聽見主的道.
包括在亞西亞的哥羅西城.
以巴弗很可能就是在那個時候.
他得以聽到福音 相信福音 順從基督.
以巴弗也都將你們在聖靈裡的愛心.
告訴了保羅和提摩太太 告訴了他的同工.
向保羅匯報哥羅西教會的狀況.
除了表揚信徒活出基督裡的信亡外.
也都講明了教會群體當時正面對.
一些錯誤教導的困擾.
促使保羅寫這封信給那幫.
蘇美蒙面的主內弟兄姊妹.
保羅對所聽見你們在基督耶穌裡的信心.
和在聖靈裡的愛心都確信不疑.
因為他知道以巴弗是一位忠信的僕役.

$^{281}$忠於主耶穌所託付的使命.
信實可靠地服侍他所事奉的對象.
以巴弗是保羅的同工 一起做主的奴隸.
保羅勉勵哥羅西的信徒 也勉勵你和我.
要在我們所信的道上恆心.
認定以巴弗所傳的福音是從神那裡來的真道.
持守以巴弗所學的真理.
在假教師的錯誤教導和搞擾當中.
要屹立在穩固的根基上 堅定不移.
今天我們會在一個知識爆炸的資訊時代.
各式各樣的潮流思想.
充斥了大眾傳媒和網絡.
加上有很熱門的人工智能.
做出來的一些似是而非的事象.
你看了都不知道究竟是真人還是人工智能.
製作出來的一個影像.
信徒需要神思明辨.
明辨我們所聽見的信息是來自什麼人.
是否基於神的話語 是否合乎聖經的真理.
被神分別為聖屬基督的信徒群體.
是被建造在使徒和先知的根基上.
就是以先知和使徒所寫下來.
被正典的聖經為我們信仰的基礎.
並且由基督耶穌自己為防角石.
防角石就是有時候翻譯為頭塊的石頭.
當一個人建屋的時候 首先放下來.
成為那間屋的一個角落.
那個角落就定準了那間屋會建在哪裡.
那間屋的方向在哪裡.
你和我無要與基督看齊.
因為接下來放下來的那些基礎.
放下來的那些石頭 都要和這個防角石看齊.
那間屋才會建得整齊整齊.
我們要和基督看齊.
要和其他的基督的僕人一起做他的奴僕.
做基督的忠信的僕人.
為者原父 我們自從聽見的日子起.
就為你們不停地禱告祈求.
願你們對他的旨意得到充滿的認識.
有全被屬靈的智慧和領悟.

$^{321}$保羅聽見哥羅西教會的情況之後.
他看到神的作為已經顯明在這幫信徒群體當中.
他沒有停頓為他們祈求.
在他日常的禱告裡面 他不斷地紀念他們.
為他們的需要帶到神的面前.
保羅祈求的內容.
是信徒被關乎神旨意的整全知識所充滿.
就是讓神的旨意充滿他們的思想.
不留下任何空間給那些錯謬的言論.
面對很多的錯謬思想橫流的時候.
我們更需要的是花時間在那些正確的.
在神的話語 在神的真理上面去紮根.
也不再花很多無謂的時間在網上去聽.
去看很多不知道從哪裡來的一些信息.
我們要對神的旨意有全面的認識.
知道耶穌基督的福音是全被的.
是足以滿足我們一切的需要.
就不需要在基督以外別有所求.
這樣的了解是藉著一切屬靈的智慧和悟性而達到.
靠著聖靈在信徒心中的工作.
引導我們進入一切真理.
能夠領悟神對你 對我的旨意.
而且由智慧選擇 按著神的心意去行.
明白驚威耶和華是智慧的開端.
凡遵行祂的命令的 是聰明人.
因為除了神的令 沒有人知道神的事.
我們所領受的不是世上的令.
那是一種神來的令.
叫我們能夠知道神開恩賜給我們的事.
就持定我們所信的福音真理.
存著盼望來面對明天.
聖靈幫助你和我可以充分明白神的旨意.
最終的目的是使你們行事為人.
對得起主 凡事蒙他喜悅.
對得起主不是說我們的言行.
我們的好行為足以報答或者償還.
主耶穌為你 為我所付的代價.
那是說我們行事為人 與蒙召的恩相承.
配合神恩召我們的旨意.
跟從耶穌的榜樣 追求滿足神的心意.

$^{361}$我們就能夠蒙神喜悅.
接著保羅說明.
希督的僕人蒙主喜悅的行事為人有甚麼表現?.
第一 在一切的善事上結果子.
在一切良好美善的事情上有所表現.
活現出聖靈所結的果子.
就是人愛 喜樂 和平 忍耐 恩切 良善.
信實 溫柔 節制.
神將賦予屬祂的人的生命質素能夠活現出來.
第二 對神的認識越發增長.
不停的透過遵行神的旨意體驗到神的真實.
就是神話語的信實.
讓我們看到我們對神的認識在被增長.
保羅為信徒不停的禱告祈求.
願他們對神的旨意得到充滿的認識.
同時那些圈信徒不單是被增長.
我們自己要去努力追求.
配合聖靈在我們生命中的工作而得以增長.
第三 依照祂榮耀的大能得著一切能力.
是依照神的大能.
不在乎我們多強或者多弱.
乃是在於我們的神必照祂榮耀的豐富.
在基督,耶穌,女士們一切所需用的都充足.
神加添我們的力量.
是要你和我凡事能夠堅韌,恆久忍耐.
這個用翻譯做堅韌的字.
是指在惡劣的環境,在很多的壓力底下.
能夠忍受得著,能夠堅定不動搖.
而能夠恆久忍耐.
是用另一個字有時候翻譯做忍耐.
但是是指當我們受別人的傷害.
我們能夠一而再的忍受痛苦.
持續以恩慈對待對方.
向傷害冒犯你的人表現出無偽的愛心.
是聖靈的感化所產生的愛心.
不是有口無心的 只是lip service.
只是有口無心的在那裡講.
那是真心誠意 神要的不是我們嘴裡面的.
有口無心的一些話語.
是真心誠意的為神透過祂的僕人.

$^{401}$是祂使我們有資格分享基督.
分享聖徒在光明中的基業.
天父已經賜給所有在基督裡的人.
包括外邦人在基督耶穌裡.
藉著福音德義同為後製同為一體同蒙英雄.
同為後製就是有資格和聖徒一起承受基業.
在舊約裡面我們看到神賜給以色列人應許之地.
就是他們用來立國的迦勒美地.
是他們世代相傳的基業.
如今對新約的基督徒來說.
在基督裡面 因為祂做了新約的宗保.
神給我們有一個新的約.
屬神的人也都會承受地圖.
叫蒙召的人得著所應許永遠的產業.
為的是要成為基督國度的領土.
當耶穌再來在地上建立神國的時候.
承認說世上的國成了我主 我主基督的國.
祂要作王直到永永遠遠.
耶穌基督要以王的身份向神的子民宣告.
你們這蒙我父賜福的.
可耐承受創世以來為你們所預備的國.
與眾聖徒在光明中同得基業.
當耶穌展開在地上的使弓的時候.
聖經說到祂來到迦勒尼宣傳神的福音.
神的福音是甚麼?.
就是日期滿了 神的國近了 你們當悔改信福音.
神旨意的核心是在人間建立神的國.
神在祂的子民當中要掌權 要作王.
保羅為信徒祈禱.
期盼他們被關乎神旨意的知識所充滿.
滿心知道神的國即將降臨在地上.
所以你們要謹慎.
因為不知道你們的主是哪一天來到.
就要時時審查我們自己的行事為人.
是否配得起主 有沒有凡事蒙他喜悅.
在善事上多結果子.
在對神的認識上越來越被增長.
從神禮的能力也都不斷的被加強.
更加是越來越帶著喜樂的心.
我們知道這些不是只能用嘴巴說.

$^{441}$這些不是唱高調.
因為祂救了我們脫離黑暗的權勢.
父神已經將你和我.
從那些執政的掌權的管轄者.
幽暗世界 空中屬靈氣的惡魔的權勢中.
拯救出來.
我們已經親身經歷到.
我們不再在他們這些罪惡的權勢的控制底下.
我們不再做罪的奴僕.
既然已經脫離撒旦的轄制.
神就給我們有自由.
去選擇順從神做誰的奴僕.
如今 現今你們既從罪裡得了釋放.
做了神的奴僕 就有誠誠的果子.
那結局就是永生.
神給我們有自由 去選擇做神的奴僕.
去事奉神 去信服神.
因為神不但救我們脫離黑暗.
祂同時帶領了你和我進入光明.
就是把我們遷到祂愛子的國裡.
是光明的國度.
是和這個黑暗的權勢完全相反.
而且絕對是更好,更美的結局.
就是你和我因為信耶穌.
就不自滅亡,反得永生.
你和我被神遷到祂愛子的國度裡.
使得神的旨意得以成全.
是透過神指耶穌而達成的.
我們能夠脫離黑暗的權勢.
進到基督的國度裡.
完全是靠著在這兒子裡.
我們得著贖放,得了罪惡的赦免.
我們得夢救贖和釋放,重獲自由.
是需要負極高昂的贖價.
是沒有人能夠負得起.
詩人說到有人曾經想贖他自己的弟兄.
替他來將贖價給神.
但結果因為贖生命的價值極高.
只好永遠罷休.
但今天靠著神的恩典.

$^{481}$恩著祂愛我們的緣故.
我們可以靠著神的愛治.
基督獻了一次永遠的贖罪祭.
為全人類捨命流血.
使所有信耶穌的人罪過得以赦免.
除去了人與神之間罪的阻隔.
使我們既恩順清義.
就藉著我們的主耶穌基督得與神相和.
可以成為在基督裡的人.
是基督國度的國民.
我們都犯了罪.
規缺了神的榮耀.
如今及蒙神的恩典.
因基督耶穌的救贖就白白的清義.
在基督裡的救恩.
對人來說 對你和我來說是白白的.
是免費的 但卻不是廉價的.
是極高昂的價值.
而是我們付不起的代價.
唯有神子為你為我付上這個極大的代價.
才使我們直者愛子的血德蒙救贖.
過犯得以赦免.
乃是照著祂豐盛的恩典.
將你和我想不配得到的好處賜給我們.
這就是恩典.
恩典是將你和我不配得到的好處.
不配得到與神和好的基礎.
就是白白的.
不需要我們付任何代價的給了我們.
請問你是怎樣回應主的恩.
你是怎樣回應神的愛?.
神從創立世界以前在基督裡揀選了我們.
祂又透過基督的福音指示世人.
藉著相信接受神子耶穌.
承傳了的救贖恩典.
就可以享受做神兒女的權柄.
與神重建和好的關係.
神又藉著聖靈在基督裡的人生命中的作為.
聖靈在你和我生命中的作為.
使我們可以得到在基督裡的順望愛.

$^{521}$作為我們在福音真理道上站立得穩的根基和動力.
使我們的生命不斷地被更新.
我們不斷地成長.
經歷基督的福音在我們當中的結果子和增長.
使基督在我們身上被顯大.
基督的僕人.
無論是保羅,是伊巴菲,是提伯太,是你,是我.
我們都是基督的僕人.
我們都要做基督忠信的僕役.
承認自己生命的主權是屬於神的.
是祂將我們買回來的.
是祂將我們買熟回來的.
屬於神 一起去做神的奴隸.
忠於神託付的使命.
到處廣傳福音.
使萬民作主的門徒.
由於關心信徒靈明的成長.
基督的僕人為眾信徒不主的禱告,祈求.
使他們可以充分明白神的旨意.
將我們遷到愛子的國度裡.
也要效法基督的行事為人.
追求凡事 蒙他喜悅.
被神裝備好.
去承受基督的基業.
去成為神國度裡有貢獻的成員.
我不知道你有沒有聽過一個叫做.
約翰·牛頓這個人的名字John Newton.
但是很可能 無論你有沒有聽過這個人的名字.
你很可能都聽過他寫的一首詩歌《奇異恩典》.
《奇異恩典》是一首不單止很受歡迎.
但是也觸動很多人的一首詩歌.
因為約翰·牛頓將他自己的生命的經歷寫下來.
在他相信耶穌之前.
約翰·牛頓是一個粗魯的講話.
講很多的荒謬.
講粗口 生活很放蕩 很墮落的一個船員.
他在船員當中可以說是相當有成就.
他終於成為了一艘陸地船的船長.
但是他的生活一直在這些罪惡.
這些不好的事情裡面.

$^{561}$但是感恩他有一個慈愛的母親.
在他十年幼的時候.
就向他講述耶穌基督的故事.
他終於有機會從一本書看到.
一本叫《孝法基督》的書 他看到.
認識到耶穌基督 耶穌基督的救恩.
他相信接受耶穌基督做他個人的救主和生命的主.
從此之後 約翰·牛頓在九片英國各地.
傳揚基督的福音.
到他晚年的時候 他已經身體軟弱到一個地步.
他每次站港台他都需要有個人站在他身邊.
有個助手在那裡輔助他.
他的眼睛幾乎看不到東西.
他講話已經變得聲音柔弱 音聲細氣.
不是大聲大氣的在那裡講.
但是他只要一息尚存 沒有任何的東西.
去阻止他去宣講基督的福音.
有一個星期天 他在一間教會講道的時候.
他重複的講了一句話.
耶穌基督是寶貴的 Jesus Christ is precious.
他的助手就細聲跟他說 你已經講了兩次了.
你不要再重複講了 老人家好像囉囉嗦嗦.
可能講過一樣東西.
約翰·牛頓突然間轉過身來.
對著他的助手很大聲的講.
是 我真的已經講過兩次了 我還要再講.
這位當時德高望重的老傳道人.
他再次的開口 再次的大聲疾呼.
耶穌基督是寶貴的 Jesus Christ is precious.
弟兄姊妹 你有沒有真的認為耶穌基督是寶貴的?.
你和我的生活有沒有具體反映出.
祂在我們生命中是寶貴的?.
祂在你和我的心中是居首位,是掌王權.
如果不是 是否因為你和我還沒有認識到祂的超越性.
還沒有體驗到基督能夠滿足我們一切的需求.
我們內心最深的渴望?.
還沒有信耶穌的朋友.
你知道基督是多麼的寶貴嗎?.
正經地告訴我們 反而接待他的接待耶穌.
就是信他的人 信耶穌的人.

$^{601}$他就賜給他們權柄 造神的兒女.
今天你從神的僕人保羅所寫的書信裡面.
聽到基督的福音.
請問你有沒有願意信靠耶穌以成為神的兒女?.
你有沒有渴望能夠進入基督的國度裡面?.
求神幫助你 求神幫助我.
我們一起低頭 祈禱.
希望你感恩 因為你所給過我們的是傳遞的恩典.
在基督裡我們可以成為一個新造的人.
我們是一個新造的人 我們是一個無礙的人.
求神幫助我們 真正能夠明白基督的福音.
有何等大的能力 要救一切相信的人.
讓我們看到做基督的僕人.
是怎樣的盡心盡性 忠於你的託付.
承擔你給我們的使命 靠著你給我們的力量.
去做你所喜悅的事.
以致我們真是能夠為基督的國度.
裝備更多的人才 裝備更多的神的僕人.
在當中一起事奉 一起將你的心意帶到地上.
求主使人 特別紀念我們當中還沒有信主的.
但願基督的朋友們 求神你自己開他們的眼.
真是讓他們看到基督的超越性 基督的充足性.
讓他們都好像我們一樣 來到你的根前.
你的憐憫 蒙你自己的恩典.
成為神國度裡面的人.
求主施恩 聽我們祈禱 奉主耶穌基督的名.
Amen.
\newpage



\section{約翰福音 4:19-30}
\label{sec:032206}
\textbf{真正的敬拜}~~~~駱柱成牧師
\newline
\newline
連結: \href{http://www.church.com.hk/acms/content.asp?site=cdc\&op=show\&type=product\&code=032206\&layout=sermon}{\texttt{032206}} ~~~~ 語音日期: 2026-05-24
\newline
\newline
\hyperref[sec:032205]{\small{< < < PREV SERMON < < <}}
~
\hyperref[sec:index_chronic]{\small{[返順時目]}}
~
\hyperref[sec:032207]{\small{> > > NEXT SERMON > > >}}
\newline
\newline
約翰福音 4:19-30
\newline
\begin{longtable}{cl}
\hline
\hline
章節 & 經文 (和合本修訂版)\\
\hline
4:19 & \begin{tabularx}{0.7\textwidth}{X} 婦人對他說：「先生，我看你是一位先知。 \end{tabularx} \\ \\ \relax
4:20 & \begin{tabularx}{0.7\textwidth}{X} 我們的祖宗在這山上敬拜神，你們倒說，應當敬拜的地方是在耶路撒冷。」 \end{tabularx} \\ \\ \relax
4:21 & \begin{tabularx}{0.7\textwidth}{X} 耶穌對她說：「婦人，你要信我。時候將到，你們敬拜父，既不在這山上，也不在耶路撒冷。 \end{tabularx} \\ \\ \relax
4:22 & \begin{tabularx}{0.7\textwidth}{X} 你們所敬拜的，你們不知道；我們所敬拜的，我們知道，因為救恩是從猶太人出來的。 \end{tabularx} \\ \\ \relax
4:23 & \begin{tabularx}{0.7\textwidth}{X} 時候將到，現在就是了，那真正敬拜父的，要用心靈和誠實敬拜他，因為父要這樣的人敬拜他。 \end{tabularx} \\ \\ \relax
4:24 & \begin{tabularx}{0.7\textwidth}{X} 神是靈，所以敬拜他的必須用心靈和誠實敬拜他。」 \end{tabularx} \\ \\ \relax
4:25 & \begin{tabularx}{0.7\textwidth}{X} 婦人對他說：「我知道彌賽亞—就是那稱為基督的—要來；他來了，會把一切的事都告訴我們。」 \end{tabularx} \\ \\ \relax
4:26 & \begin{tabularx}{0.7\textwidth}{X} 耶穌對她說：「我就是，正在跟你說話呢！」 \end{tabularx} \\ \\ \relax
4:27 & \begin{tabularx}{0.7\textwidth}{X} 正在這時，門徒回來了。他們對耶穌正在和一個婦人說話感到驚訝，可是沒有人說：「你要甚麼？」或說：「你為甚麼和她說話？」 \end{tabularx} \\ \\ \relax
4:28 & \begin{tabularx}{0.7\textwidth}{X} 那婦人留下水罐，往城裡去，對眾人說： \end{tabularx} \\ \\ \relax
4:29 & \begin{tabularx}{0.7\textwidth}{X} 「你們來看！有一個人把我素來所做的一切事都說了出來，難道這個人就是基督嗎？」 \end{tabularx} \\ \\ \relax
4:30 & \begin{tabularx}{0.7\textwidth}{X} 他們就出城，來到耶穌那裡。 \end{tabularx} \\ \\
[1ex]
\hline
\hline
\end{longtable}
$^{1}$大家住內平安.
親愛的朋友.
如果你是第一次來到.
或者是在網上的.
我們也很歡迎你跟我們一起去敬拜.
我今天說的是真正的敬拜.
上次我就說到.
撒瑪利亞婦人.
是說到一到十八節的.
今天我們繼續上一次的講題.
就是耶穌和撒瑪利亞婦人的一個談話.
當時猶太人對撒瑪利亞人.
當時有很深的偏見.
撒瑪利亞人不僅是文化不同.
他們在信仰上.
以及他們的血統,身份上.
都是被猶太人去稀罕的.
人很多時候都是劃界線的.
所以他們之間就有一種民族界線.
宗教的界線.
以及人與人之間的界線.
而耶穌基督.
祂就是跨越了這一切的界線.
專門去撒瑪利亞找這個婦人.
不是祂必須要經過撒瑪利亞.
才能去到祂要去的目的地.
而是祂看到有一個靈魂在那裡.
這個靈魂很需要祂.
所以耶穌基督可以說是為了祂而去的.
那時候是正午.
是很熱的.
那時候很少人出來打水.
但是這個婦人偏偏就是中午的時候才出來打水.
因為她的背景.
有些可以說不是很能見光的.
但是她就是要避開這些人的目光.
所以她是一個被人排斥的.
被人稀罕的.
也是活在那種有些自卑.
有些羞恥的那種警況當中.

$^{41}$她可以說她的背景是有.
在經文看到是有五個丈夫的.
現在有這個.
耶穌基督說不是祂的丈夫.
還要是五個半丈夫就變了.
但是耶穌基督在這個時候是沒有稀罕祂的.
因為耶穌基督來.
祂說不是召義人悔改.
是要召罪人悔改.
所以主是看重祂的.
不是因為祂的背景如何.
而是因為祂有一個寶貴的靈魂.
耶穌基督沒有一開始就責備祂.
反而說.
請你給我喝點水.
耶穌基督很尊重祂.
耶穌基督可以說不是真正需要水.
因為後面會說到.
耶穌基督沒有提到耶穌基督喝水.
所以這個婦人受了尊重.
但是她有防備心.
充滿防備.
所以當她有自卑去了之後.
受人尊重之後.
她的防備又出來了.
能夠推翻民族和身份之間的對立.
猶太人和撒瑪利亞人.
你既然是猶太人.
為什麼香港的撒瑪利亞婦人要喝水.
所以你會看到.
很多時候人不是看對方說了些什麼.
反而是一直說的事情是要保護自己的.
但是耶穌基督沒有跟她爭拗.
反而祂要將這個婦人帶回到真理的核心.
祂沒有跟她爭拗.
耶穌就說.
你若知道神的恩賜.
和跟你說給我水喝的.
這個是誰的時候.
你就必早求祂.

$^{81}$祂也必一早給你活水.
而且接著就說.
凡是喝這個井水的還要再喝.
但是人喝我所賜的水就永遠不喝了.
這句話正是指著婦人的生命狀況的一個實際問題.
人都是這樣的.
明星不可以真正滿足別人.
金錢也不可以.
情慾也不可以.
但是當你人不斷去追求這一切的時候.
你的心是很難得到滿足的.
你越追就越渴.
你越喝這些東西你就越渴.
因為這些東西不能夠真正使人滿足.
反而是使人越喝越渴.
所以真正缺的.
很多時候真正缺的不是地上的水.
而是生命的水.
當婦人開始渴望活水的時候.
耶穌就突然跟她說.
你去叫你丈夫來.
到我這裡來.
因為耶穌基督知道她真正的問題是什麼.
不是那些外在的需要.
而是她生命的真正需要.
婦人說我沒有丈夫.
耶穌沒有羞辱她.
反而肯定她說你沒有丈夫是不錯.
但是你已經有五個丈夫了.
現在有這個不是你的丈夫.
耶穌基督沒有說是是是.
就順從祂說.
反而是將事實說出來.
將事實說出來的時候.
祂沒有傷害別人.
所以耶穌基督祂講真理.
但是祂不是為了羞辱別人.
貶低別人.
而是要拯救別人.
這是我們上一次所講的經文.

$^{121}$今天我們進入另外一個主題.
上一次講到耶穌基督.
祂是怎樣有智慧.
講到活水的問題.
今天是講到敬拜.
當耶穌基督指出祂真實這個光景的時候.
這個婦人的反應.
她不是悔改.
你有五個丈夫.
現在有這個也不是你的.
她反而馬上就要轉話題了.
她說先生.
我看出你是先知.
你是先知來的.
因為只有先知.
才有這種熟齡的洞察力.
你可以看透別人的心.
也看出我的忍情是什麼.
你也知道我的羞恥的時候.
但是她沒有馬上悔改.
她明知道耶穌基督是先知.
是神所差來的.
但是她沒有悔改.
她也沒有說.
主啊,我錯了.
求你赦免我.
或者我願意悔改.
求你救我.
她沒有講這些.
可能她會思想.
如果我信悔改的話.
離開這些情慾的事的話.
那我以後的生計怎樣呢.
我的生活怎樣呢.
如果我離開這個男人.
我就沒有依靠.
如果我離開了他.
他會不會傷害我呢.
我是一個婦人.
所以一個罪人.

$^{161}$他的思維不是一下子就悔改的.
反而是他的天然不會主動讓他悔改.
所以你會看到.
她沒有馬上悔改.
我們傳福音.
很多時候也是.
我們向人傳了很多次.
福音講清楚了.
也知道自己是不對的.
但這些人很多時候.
未信耶穌基督的人.
也不是馬上悔改的.
所以人聽見真理之後.
不是馬上去降伏.
很多時候是逃避.
拖延和轉移.
今天我們也是這樣.
聽見真理之後.
不一定願意去面對自己.
很多時候可能說.
這個道理不錯.
但我也想過得開心一點.
做人而已.
開心一點就算了.
我們相信也好.
不要太著迷.
相信也好.
何必搞得自己這麼辛苦呢?.
很多未信的人是這樣.
如果信徒呢.
很多時候我們不願意離開罪惡的時候.
也是這樣.
當撒瑪利亞婦人知道.
耶穌是先知的時候.
她沒有馬上悔改.
她很聰明.
她馬上就轉移話題.
她的口才很好.
我記得在讀大學的時候.
有個同學她的口才也很好.

$^{201}$每次輪到她搞.
我們是同一個宿舍的.
每次輪到她一起搞衛生的時候.
她都是不做的.
基本上就是逃避責任.
所以宿友都會提醒她.
輪到你了.
她很聰明.
如果那個是喜歡讀書的.
先不要說這件事.
我們今天要做甚麼作業呢.
說著說著她就說了其他.
如果是跟她一起喜歡看電視的人.
她就說大家一起看電視就好了.
先不要說這些.
所以她會說回電視.
到底男主角追到女主角了沒有.
所以聊著聊著就說了其他東西.
後來這個同學做了保險經理.
很年輕就做了一個區的經理.
她的口才很了得.
撒瑪利亞婦人其實也是很聰明.
她很快就看出耶穌是先知.
而且她比黎哥帝母更早一步.
黎哥帝母只是說拉比.
你只是一個先生.
但是撒瑪利亞婦人說你是先知.
而且她的口才很好.
她馬上就轉移話題.
土高先知呀先知.
土高到底在哪裡呢.
到底是在這個基利森山.
還是在耶路撒冷呢.
在我們撒瑪利亞這個地方.
還是在你們猶太的南國呢.
經文是這樣說的.
我們的祖宗在這座山裡拜.
你們卻是在說.
應當在耶路撒冷裡拜.
她真正的意思是說.

$^{241}$我們是不一樣的.
你是這樣的.
我是這樣的.
你們有你們的傳統.
我呢.
我有我的傳統.
你們有你們敬拜的地方.
我也有我敬拜的方式.
今天我們每間教會都很不同.
可能敬拜方式不同.
詩歌的風格不同.
我們是很古典的.
是很傳統的.
年輕一代有現代的音樂.
我們有傳統的音樂.
禮儀習慣也不一樣.
所以宗派可不可以有宣導會.
浸信會,播道會,長老會.
不同的宗派可不可以有?.
可以.
但是宗派主義可不可以有?.
什麼是宗派主義?.
就是將我這個宗派.
我這些教義絕對化的.
我這裡只是純正的真理可以說.
我只是真正屬神的.
如果你看聖經你就知道了.
我是屬保羅的.
你是屬亞伯羅的.
我是屬基法屬彼得的.
後來又有人說.
我是屬基督的.
那個分門別類就出來了.
我是屬基督的.
是不是真的我是屬基督的呢?.
是不是我就是最屬靈的呢?.
但其實背後很多時候.
不是真正是屬靈的.
而是只有我只是屬基督的.
你們呢?.

$^{281}$你們好像不是那麼純正.
我只是真正的屬靈.
但是當你有這樣的驕傲的時候.
你的屬靈就變成屬靈了.
可以說是靈感的靈.
驕傲是使人的靈性倒退的.
今天也是一樣.
有人說我是屬長老會的.
我是真信會的.
我們是宣道的.
我們這裡才有道.
我們是宣道.
我們這裡才有道.
是不是這樣呢?.
所以你們都是宗派.
但是屬基督的人就說.
你們都是宗派.
我就不一樣了.
如果你有中國大陸背景你就知道.
有些聚會所.
他很難跟別人合作.
因為他就是強調.
不是屬於任何一個宗派的.
不是屬於你宣道會的.
不是屬於你真信會的.
也都不是長老會的.
反而形成一種很驕傲的心態.
就是說.
我.
我就是純正的.
你們都不是.
所以我們很難合一.
我和你沒什麼商量的可以說.
所以宗派可以有.
宗派主義不可以有.
你一有.
你就絕對化了.
你就驕傲了.
真正的合一.
不是說.

$^{321}$真正的合一也都不是說.
隨便混在一起.
隨便吧.
大家都是教會.
不是這樣.
應該是在主的道那裡合一.
當如果你在真理的道.
站立不穩的時候.
一味追求那種合一是不可能的.
也是危險的.
我們可以在不同的宗派那裡敬拜神.
不是說.
我一定要來這個教會.
一定要屬於那個宗派.
而是我們真正相信福音的人.
相信耶穌基督的福音的人.
純正的福音的人.
都是在基督裡面的.
凡是信基督福音的人.
都是主內的弟兄.
所以我們真正盼望的.
不是那種宗派婢子那種抬舉.
高舉自己.
而是每個教會都回到聖經.
在真道那裡合一.
在真理那裡.
撒瑪利亞夫人其實很聰明.
她的口才真的很好.
她一開始就說.
你是猶太人.
我是撒瑪利亞夫人.
你是男人.
我是女人.
你有你們的傳統.
我有我們的傳統.
後來也提到耶穌基督.
她真是很聰明.
所以她就馬上轉換話題.
你說你們在耶路撒冷敬拜.
我們在撒瑪利亞的基尼心山上敬拜.

$^{361}$原先是說丈夫的問題.
一下子就說到先知的問題.
敬拜的問題.
在哪個山頭去敬拜的問題.
宗派的問題.
但是耶穌基督很有智慧.
祂沒有跟著他走.
也沒有用傷害的話再去詆毀他.
你有五個丈夫.
你現在為什麼又轉換話題.
剛剛才說到丈夫的問題.
你馬上就轉換話題了.
耶穌基督沒有這樣做.
沒有去詆毀人.
相反耶穌基督很有智慧.
撒瑪利亞夫人說什麼.
她就從真理的角度去回應什麼.
耶穌基督是真理的本體.
祂也是智慧的本體.
祂本來就是真理.
祂本身就是智慧.
耶穌基督對撒瑪利亞夫人說.
要到耶路撒冷敬拜.
其實這個不是猶太人的傳統.
也不是先知自己的意思.
不是他們自己的意思.
是上帝在律法當中的吩咐.
要到耶路撒冷敬拜.
因為出埃及記說得很清楚.
神曾經命令以色列人的南丁.
每年三次要到神面前去朝拜祂.
當時敬拜的中心就是耶路撒冷.
而且神允許他們.
當他們去敬拜神的時候.
他們的家庭得到保護.
他們的牛羊不會被搶去.
他們的妻兒也不會被欺壓.
也不會被別人擄去.
這是神給予那些敬畏祂的人的保護和平安.
後來南北國分開了.

$^{401}$南國,北國因為政治的原因.
北國人不再去耶路撒冷.
所以問題不是在上帝的命令改變了.
而是北國人沒有照著律法去行.
但是與此同時.
北國的先知.
例如以利亞,以利沙.
他們都沒有對百姓說.
你們不要去耶路撒冷.
神有新命令了.
這些先知沒有這樣的教導.
雖然以利亞,以利沙.
聖經沒有提到他們去耶路撒冷.
可能他們真的沒有做那種榜樣.
但是上帝依然重用他們.
為什麼?.
因為他們抓住一個很重要的原則.
是聖經所啟示的一貫原則.
他們不是死守在那個字句那裡.
反而是沒有廢掉那個真理.
所以當撒瑪利亞夫人開始討論宗教的時候.
到底應該在哪裡敬拜?.
在基列森山還是在耶路撒冷?.
耶穌基督說了一句很重要的話.
你們拜父不在這個山.
也都不在耶路撒冷.
這是很大的一種釋放.
因為耶穌基督將人對神的那種敬拜.
從地域性,不是基列森山,撒瑪利亞.
也都不是耶路撒冷.
在地域性當中要釋放出來.
有些信徒以為.
特別是有些極端的靈因.
他說:在越高的山.
越神聖的地方.
我就越靠近神.
如果你是廣州來的.
你知道有白雲山.
在香港有獅子山.
獅子山比白雲山還要高.

$^{441}$因為它有495米.
白雲山382米.
獅子山比白雲山高.
是不是離神近一點.
所以神就特別聽我禱告一點.
是不是我家比你家格雷倫社高一個台階.
所以神更聽我禱告.
當然不是.
因為耶穌基督早已經在兩千年前.
已經說了一句很重要的話.
敬拜不是在這座山.
也都不是在耶路撒冷.
為什麼?.
因為耶穌基督已經把真正敬拜.
從那種地理,空間,時間.
這些限制釋放出來了.
耶穌基督沒有將神的真理.
限制在一個民族的裡面,民族主義的裡面.
不是在基里森山.
也都不是在耶路撒冷.
只可以敬拜.
真正對神敬拜的.
是在屬靈的層面去進行的.
真正的敬拜是超越了時間和空間的.
所以耶穌基督其實是告訴撒瑪利亞婦人.
你真正的問題不是在這座山.
也都不是在這個地方.
你真正的問題是在你根本上不夠明白.
不夠認識神祂的話的精義是什麼.
聖經不是只停留在字句裡面.
是要捉住祂的屬靈的精義.
為什麼?.
為什麼律法曾經要求人上耶路撒冷?.
因為舊約敬拜的制度一早就提到了.
為什麼以利沙祂沒有去.
上帝反而一樣重用祂呢?.
因為這些先知一早就看到那個實質了.
因為神真正看重的不是地理的本身.
也都不是你在耶路撒冷還是在撒瑪利亞.
而是你和神的關係是什麼.

$^{481}$看重的是神與人的關係.
所以耶穌基督沒有陷入猶太人與撒瑪利亞人的民族的爭鬥.
那個山頭爭鬥,那個宗教的爭鬥.
耶穌基督沒有用那種民族主義去維護真理.
也都沒有用那種宗派的精神去處理這些屬靈的問題.
既然全地都可以.
敬拜上帝是不是就是今天我們就不需要回來敬拜呢?.
在家中YouTube多方便.
那麼遠,那麼早起床.
我不想了,在家中算了.
上帝,聖經沒有這樣的教導.
因為聖經說你們不可以停止聚會.
聚會是很重要的.
不是因為那個地方特別神聖.
而是因為聖徒的一個聚集.
教會不是靠建築物而成為聖結的.
信徒是奉主的名聚會的時候.
主就在我們中間.
我有一個同學.
彭長老認識的,叫做夏份子牧師.
他在Montrey Park那邊.
蒙市那邊公園租了一個地方.
是政府的場地.
每星期日晚他就在那裡聚會.
但是其他時間,其他機構用.
包括星期日早上,下午都還是其他機構用.
所以原本是一個普通的政府建築.
但是當聖徒在那裡敬拜的時候.
那個地方就分別出來了.
為什麼?.
不是因為政府的建築本身很聖結.
它只不過是一個建築.
而是因為聖徒在那裡敬拜.
有聖徒敬拜,有主的同在.
那個地方就成為聖所了.
所以教會不是一個建築物.
教會也不是禮派堂的本身.
教會是蒙教屬的.
被主分別為聖的人.
是奉主的名聚集的一個群體.

$^{521}$這個叫做教會.
所以我們一定要明白.
聖山,聖地不是那個地方天然好聖.
聖字是什麼意思?.
就是分別出來的.
分別出來歸給神的.
這個就是聖了.
因著神的同在.
因著被分別出來.
因著寶血的潔淨.
因著聖靈的充滿.
歸屬於上帝.
希伯來文「聖」這個觀念.
第一次出現的時候.
上帝對摩西說.
你要把你腳上的鞋脫下來.
因為你站的這個地是聖地.
那塊地為什麼成為聖地?.
不是因為那塊沙漠地特別神聖.
而是因為聖潔的主在那裡.
主在哪裡?.
聖就在哪裡?.
因此上帝將我們分別出來為聖.
我們無論在哪裡敬拜.
只要用心靈和誠實去敬拜.
那裡就成為一個敬拜的地方.
不是在這座山.
也不是在那個地.
因為上帝是靈.
當神的兒女用心靈.
藉著聖靈,誠實去敬拜的時候.
神就和他們一起.
閱納他們的禱告.
閱納他們的敬拜.
那個地方因此成為歸主為聖的地.
所以「聖」是這樣.
真正拜父的.
就是耶穌基督所說的「我的父」.
拜父的就是耶穌基督所說的「我的父」.
這位父就是上帝.

$^{561}$認識耶穌基督.
認識「子」的.
你才可以認識「父」.
不認識「子」的.
你不可以認識「父」.
因為耶穌基督來到這個世界.
就是要使我們認識這位「父」.
在天上的「父」.
成為我們的「父」.
也都成為我們的「父」.
耶穌基督來.
要讓這位「父」也都成為我們的「父」.
只有一位「父」.
就是獨一的上帝.
這位「父」是萬有的源頭.
祂創造一切.
祂也都掌管一切.
不是局限在耶路撒冷.
也都不是局限在薩瑪利亞的基尼森山.
祂不是人想像的那種「山神」.
也都不是那塊土地公公爾爾.
祂不單是猶太人的上帝.
祂也是中國人,美國人,全世界人的上帝.
因為全世界都是這位上帝所造的.
祂是真正的上帝.
耶穌對薩瑪利亞婦人說.
「你所敬拜的,你不知道」.
「我們所拜的,我們知道」.
這裡讓我們看到.
真正的信仰和敬拜.
是要和真理的知識連在一起的.
今天很多的廟宇.
你問他們「為甚麼你這樣信?」.
可能有些人會說.
「我祖宗都是這樣信的」.
「我媽媽我奶奶都是這樣信的」.
「那我就這樣信」.
「你拜甚麼神?」.
「我拜黃大仙」.
「香港最有名的」.

$^{601}$「誰教你的?」.
「我啊,媽媽教我的」.
「你媽媽又是誰教的?」.
「我奶奶教的」.
「所以為甚麼要信黃大仙?」.
「黃大仙」.
「香港人都知道有求必應」.
「有求必應所以我又信了」.
「他南方一帶很出名」.
「香港特別出名」.
「廣東人也都知道」.
「北方人開始有些蒙查查了」.
「美國來到美國,少之又少」.
所以他可能答不出為甚麼要信.
為甚麼要信神是全地的主.
很多宗派缺乏一種普世性的理性.
即是神應該管理一切.
但他們就有他們的局限.
所以再問多些的時候.
他們就答不出.
但基督信仰.
基督教好的神學.
是建立在上帝聖經的基礎之上.
是合乎理性的方式.
系統地將它講明白的.
耶穌基督講出最偉大的敬拜定義.
真正拜父的要用心靈和誠實去拜祂.
中文翻譯為心靈和誠實.
英語是In Spirit, In Truth.
在真理和靈的裡面.
在聖靈和真理的裡面.
意思是在聖靈的引導之下去敬拜神.
在真理的原則的裡面去敬拜神.
今天很多人以為我信甚麼神都可以.
我要拜觀音.
觀音就成為我的神.
我要拜門神保護我家和我的屋.
這麼多生物我要拜門神.
看著門口.
如果有些人去坐郵輪.

$^{641}$我要拜海神.
有甚麼風浪怎麼辦.
坐飛機我要拜雷公.
只要心性就好了.
很多人甚至覺得.
拜得神多自有神庇佑.
但聖經所講的這個信.
不是你信甚麼就是甚麼.
真正的信是對真理的忠誠.
你要回歸到真理的裡面.
回歸到那個一直永恆不變的真理裡面.
這個才是信.
我這台車是白色的.
有時很容易就開花了.
泊車一下就開花了.
很多時候就掉了漆.
所以我就貪便宜.
我又不會噴過它.
我這台車是二十年了.
所以要補漆.
在網上買很便宜的.
就要揀色.
白色的你不會揀黑色.
因為不對.
黃色很近了開始.
但都不對.
都對不上.
白色是不是一定是對的.
不是的.
原來白色還有分米白.
象牙白.
珍珠白.
奶白.
很多的白.
所以我這台車是珍珠白的.
差一點就不一樣.
你塗了也不像樣.
塗漆一定要按著它原廠的顏色是怎樣出的.
你的製造商也要按標準去調色.
所以是按著原來的顏色去對色才行.

$^{681}$同樣信不是憑感覺.
也不是差不多就行.
真正的信是要回到真理裡面.
你要按真理.
所以.
人的理性應該要回到真理裡面.
我們的感情要回到神的性愛裡面.
意志要回到上帝永恆的旨意裡面.
這個字眼叫信.
信不是隨便.
我信什麼就信什麼.
不是憑感覺是什麼就什麼.
耶穌說真正拜父的.
必須要藉著真理和聖靈.
聖靈是什麼?.
聖靈是三維一體當中的大三維格.
在這裡耶穌就將聖靈和真理連在一起.
他們是不分開的.
今天很多教會的失敗.
在事奉上,在行為上.
很多信徒在事奉上,在行為上失敗.
很多時候是因為缺乏真理和聖靈.
兩樣東西是分開的.
因為聖靈是真理的靈.
他們不是分開的.
聖靈啟示真理.
就是神的真理.
所以真理的靈亦都是聖靈.
將真理啟示下來.
聖靈亦都用真理去光照我們每一個信徒.
每一個人.
聖靈亦都用真理去教導我們信徒.
讓我們回到真理的裡面.
亦都引導他的兒女進入真理的裡面.
所以聖靈和真理一直都不是分開的.
因此聖靈藉著真理去啟發人.
藉著真理去改變人的生命.
又將人帶回到真理的裡面.
真理和聖靈是絕對不可以分開的.
很多時候我們都講聖靈,聖靈.

$^{721}$如果你聽到有些信徒說.
「聖靈叫我這樣做,聖靈叫我這樣做」.
但是你聽他們的話的時候.
沒有那個真理的根基.
很多時候他們是憑著那種感覺去做事的.
但是另外一些人.
天天讀經那個原則,原則,原則.
那個規條,規條.
但是他們停留在頭腦那種規條裡面.
沒有依靠聖靈.
所以一邊是聖靈.
一邊是真理.
但是另一邊是沒有聖靈的.
聖靈和真理不是分開的.
如果你將它分開.
很多時候我們就落在很多的迷惑當中.
當我剛剛來美國的時候.
我一次去到很新的那個時候.
去到一間教會.
和一個負責教導的聊天的.
我就講到一個.
我曾經在山上.
在大陸山上那些山讀過書.
他就講了一句.
山上禱告會感覺很親近神的.
就算不用聖經.
也都覺得神特別同在的.
我一聽我就.
這句話印象很深刻.
因為他所講的.
和我所思想的.
好像對不上.
在山上禱告,敬拜.
環境比較安靜.
是舒服一點.
是容易安靜.
很容易進入那個敬拜.
不受騷擾.
這一點我覺得是對的.
但是真正的敬拜.

$^{761}$不應該離開聖經的真理.
聖經的真理和聖靈.
它們不是分開的.
兩樣東西應該是結合的.
真正的信.
真正的愛.
生命的根基.
如果我們不是在聖靈的啟示真理來信.
我們的信就變成.
很多時候只不過是我們一個人的自信.
那個不是聖經所講的信.
真正的信是對真理的忠誠.
你要回到真理的裡面.
如果我們不按真理而行.
我們的行為.
很多時候就按著我們自己的自由行.
聖經說我們要用真理去促要.
用真理去促要.
意思就是說.
你一切所行的.
應該在真理的裡面.
那個規範的裡面.
靠著聖靈.
靠著聖靈去行出來.
上帝講過.
不可以為自己雕刻偶像.
也不可以做什麼形象.
不可跪拜那些像.
這裡講到「做」.
也講到「拜」.
為什麼重點不是那個「做」呢?.
而是那個「拜」呢?.
因為聖經上帝也吩咐人做那個基魯壩.
基魯壩也是一個像.
放在藥櫃的上面.
所以形象本身不一定就是犯罪.
不一定.
真正的問題在哪裡呢?.
就是你做這些受造的物.
你要來「拜」.

$^{801}$這個就出問題了.
十誡強調的.
不要為自己去做像.
去敬拜.
因為這個是對神的一個褻瀆.
為什麼?.
因為上帝是那個「靈」.
不是物質.
創造物質的.
本身祂就不是物質.
創造世界的.
祂本身不是世界.
創造萬有的.
祂比一切萬有都重要大.
祂可以超乎萬有.
高過萬有.
所以萬有就算加在一起.
都不可以比上帝大.
上帝是「靈」.
祂不是那個物質.
祂不是物質.
祂就不是我們這個肉眼.
這個物質的肉眼所看得到的.
而是祂所造這個物質的世界.
我們可以看得到.
看得到的.
這個不是上帝.
真正的上帝.
上帝是「靈」.
祂是看不到的.
我們今天敬拜一位看不到的神.
因為這位看不到的上帝.
創造了這個看得到的世界.
這個在世界裡面.
你很難用肉眼去看得到上帝.
但不等於沒有上帝.
我們都很喜歡去迪士尼公園玩.
但你看不到那個工程師.
但不等於沒有這個工程師.
因為做迪士尼的是那個工程師.

$^{841}$但迪士尼本身.
祂不是自己將自己造出來的.
上帝不是這個世界.
上帝創造萬有.
但不是萬有的本身.
祂是超越了萬有.
超乎萬有之上.
上帝不單止創造物質.
上帝也創造了「靈」.
為何這樣說?.
祂說「靈」.
上帝是「靈」.
我們創造了.
我們也是有「靈」的.
上帝的「靈」是無限的.
我們呢.
我們只不過是人.
是有限的.
我們的「靈」是有限的.
上帝是自有永有的.
但我們被造.
是因著被造而存在的.
我們是很有限的.
神與人之間.
所以.
一個是無限的.
一個是有限的.
所以本質上來說.
是不是完全相同的.
可以說是不同的.
因為上帝造我們成為有「靈」的.
活人的時候.
神與人就有溝通的可能.
晚堂有位詩人叫李商人.
他形容情侶的默契是.
「身無彩鳳相非也」.
「心有靈犀一點通」.
什麼意思?.
就是男女之間的情愛.
心靈的融合.

$^{881}$是人的「靈」裡面的.
那種很高層次的.
那種高層次是動物沒有的.
因為動物.
你家裡的寵物,狗是沒有的.
因為牠沒有「靈」.
所以牠是沒有辦法.
人的「靈」那些最深層次的東西.
但人與人之間是可以的.
上帝造人的時候.
將「靈」放在人的裡面.
所以人就可以與上帝溝通.
這就成為與上帝溝通的一個可能.
所以人可以向上帝禱告.
藉著耶穌基督.
所以我們藉著聖靈的光照.
明白神對我們心中的指引是什麼.
「靈」與「靈」是可以交通的.
理性可以對真理去了解.
就成為人和敬拜上帝的一個渠道.
因為人有理性可以明白真理.
因為「靈」.
人有「靈」.
上帝也有「靈」.
所以人可以敬拜這位上帝.
所以耶穌基督說.
上帝是一個「靈」.
所以拜祂的.
必須要用心靈和誠實去拜祂.
又說因為父要這樣的人去拜祂.
上帝用聖靈去引導.
祂所做的「靈」.
引導人的「靈」.
用真理去引導人的理性.
所以我們可以用.
當藉著聖靈和真理去敬拜上帝.
聖靈所啟示的真理.
是讓人明白上帝的心意.
聖靈啟示真理.
讓人去明白上帝的心意.

$^{921}$所以人和上帝又可以交通.
因為你可以明白神的心意.
聖靈也藉著祂的真理.
讓人去明白.
所以人和神.
可以有一個敬拜的關係.
動物不可以.
今天我們來到神的面前.
不是憑著民族的血統.
也不是因為我們站在哪裡.
不是因為我們在教堂裡面.
是憑著真理和聖靈.
拜祂的人.
必須要用心靈和誠實去拜祂.
要按著聖靈.
按著真理去拜祂.
所以問題是.
很多時候我們知道真理.
但很多時候我們不懂得去回應真理.
所以耶穌基督在這段經文.
講了另外一番很重要的題目.
將人帶回到真正的真理裡面.
耶穌基督說.
「事後將到,如今就是了」.
拜父的要用心靈和誠實.
耶穌的回答不是說.
一下子被婦人滲入相對觀念裡面.
相對主義的裡面.
祂不是說每個人都有自己的真理.
怎樣拜都可以.
祂沒有這樣說.
耶穌基督是將人帶回到.
絕對的真理裡面.
是永恆不變的真理裡面.
真正拜父的.
真正拜父就是.
有真正的也有.
不是那麼真正的.
耶穌約翰福音多次提到.
「事後將到,如今就是了」.

$^{961}$約翰福音出現了很多次.
人活在歷史的進程當中.
各有各的表達.
各有各的主張.
這是很常見的.
但是有一天.
一定會回到一個獨一的原則裡面.
這個獨一的原則要顯明出來.
一切都回到真正的原則裡面.
真正的時刻一定會來到.
「事後將到,如今就是了」.
耶穌基督說「事後將到」.
其實跟撒瑪利亞婦人心中所期待的.
「未塞亞要來」.
這個「事後將到」.
「未塞亞要來」.
是很相同的.
是很相通的.
只不過一個知道.
現在就是了.
耶穌基督一早就知道了.
婦人還不知道.
當「未塞亞要來」的時候.
婦人都知道.
一切的問題都會清楚了.
只是她不知道.
她不知道站在她面前的這一位.
就是「事後將到」.
就是要來的那一位.
「如今就是了」.
如今站在她面前的.
就是她一直所盼望的.
她就是「未塞亞」.
她不知道.
很多時候人知道有一天會來.
有一些事會發生.
知道人生有終點.
知道有永恆.
我們信徒知道有一天會見主.
但問題是很多時候.

$^{1001}$我們沒有將很長遠的一天.
你覺得很長遠的一天.
和今天拉在一起.
婦人現在就是了.
沒有將很長遠的一種觀念.
和今天見到的耶穌基督.
結合在一起.
我記得前幾個月.
我的老師和我住在一起的老師.
因為他不做支架的時候.
醫生說「你只有半年了」.
他知道他的日子有限.
我也知道他不是那麼長.
他的使命已經完成了.
教到神學.
該教的教了.
現在沒有辦法再教了.
太老了.
很多東西都壞了.
很難動了.
說話又說不清了.
已經做的做了.
既然知道那一天.
該安排的也安排好了.
各地的教區安排好了.
各樣的事交代好了.
臨終的事也交代好了.
所以他離世前的兩天.
我看到他這樣.
做了一個禱告.
如果神真的合你心意的話.
求主你釋放了你的婦人.
讓她安息.
結束她地上的奴婦.
兩天之後.
那一天來到了.
但是很多時候.
一下子那一天就來了.
很多時候我們的知覺.
不是將概念當中的那一天.

$^{1041}$和今天要做的事.
拉在一起的.
有一天我們也會這樣有一天.
但是我們沒有將將來的一個.
我們一直在我們頭腦的將來的概念.
和今天連在一起.
耶穌基督說.
婦人如今就是了.
意思就是你所等待的那一天.
不是很遙遠的.
不是遙遠的觀念.
彌賽亞已經實實在在地站在你面前了.
問題不是今天.
但是真正的問題是.
你今天願不願意回應耶穌基督.
你要回應祂.
而這位耶穌基督不是突然出現在你面前的.
因為聖經一早就啟示了.
祂是上帝在歷史當中早已安排好了.
一早就預備了耶穌基督這個救恩.
救恩是從猶太人出來的.
不是表示耶穌基督將救恩.
將自己限制在民族裡面.
不是.
因為神雖然是透過一個民族.
將救恩帶到全世界.
但是神沒有將自己限制在民族裡面.
救恩不是限制在一個民族裡面.
耶穌基督也不是將祂自己限制在一個民族裡面.
神是藉著這個民族.
成就一件超民族的工作.
可以說是一個超民族的旨意.
所以有地區性的源頭.
沒錯.
耶穌基督是從猶太人出來的.
有時間的開始.
沒錯.
因為耶穌基督是按照神所定的時間來到地上的.
但神所成就的是超過了這些時間和空間.
超越了猶太人的局限.

$^{1081}$超越了歷史的長河.
今天這個時代很多事情都很真假難分.
我們看一部電視的時候.
很多時候你會在集數那裡看到.
本故事純屬虛構.
耶穌基督不是純屬虛構.
人工智能做的事情也是純屬虛構.
人工智能你看到的.
很真實地看下去.
在虛擬的空間裡.
你動來動去都可以.
但那個好像是真的.
好像而已.
但耶穌基督的信仰就不是這樣.
耶穌基督的信仰不是虛構的神話劇集.
他是按照真實來到歷史的裡面.
所以聖經一直都告訴我們.
在什麼時候君王要執政.
在什麼地方彌賽亞就要出來了.
耶穌基督不是活在虛擬的世界裡面.
不是去到虛擬的世界.
不是在虛擬的世界裡面.
是真實地來到人類的歷史當中.
他是在奧古斯都的年間降生的.
在伯利恆那裡的.
這個是歷史可以考據的.
可以證明的.
在永恆當中所定的永恆的救贖計劃.
是在時間,時空的歷史中真正地實現出來的.
永恆的旨意已經定了.
在地上要實現出來.
因此耶穌基督說.
「事後到了,如今就是」.
意思是這件事不是遙遠的.
現在我已經來到了.
現在已經進入到地上.
進入到歷史.
彌賽亞已經站在你面前了.
真正與神的關係.
不是建立在耶路撒冷,撒瑪利亞這種空間的限制.

$^{1121}$是建立在耶穌基督的裡面.
你對耶穌基督有多少的認識在裡面.
是超越了時間和空間那種真實的裡面.
但是這個救恩又實實在在地藉著猶太人而來的.
所以耶穌基督說.
因為救恩是從猶太人出來的.
不是神指揮猶太人.
而是透過猶太的民族將救恩帶給世界.
如果你一直追著這句話去研究的話.
你會看到上帝對亞伯拉罕說過.
萬國要因你的後裔得福.
所以這個後裔是單數的.
就是指耶穌基督.
亞伯拉罕有千千萬萬的後裔.
就是指這位特別的,單數的耶穌基督.
所以耶穌基督是在猶太民族降生的.
但是不是單單為猶太民族而來的.
是為萬民而來的.
神所說的一定要成就的.
神所成就的一定是在時間和空間那裡去實現的.
是我們在地上去實現的.
所以耶穌基督說.
你所拜的你不知道.
我們所拜的我們知道.
因為救恩就是從猶太人出來的.
上帝呼召亞伯拉罕.
亞伯拉罕生耳薩.
但是耳薩瑪利他不在這個簡善的當中.
耳薩生雅各.
耳素被棄絕.
後來雅各改名為以色列.
有十二個支派.
之後就不是將整個支派去除掉了.
沒有將一批一個種族去除掉.
但是如果你再看到後面的時候.
啟示錄的時候.
有一個支派被除名.
就是彈支派.
而頂上就是約瑟的兩個兒子.
為什麼?.

$^{1161}$因為彈支派是最早帶領以色列人拜偶像的.
所以拜偶像是很大的罪.
這個聖經是一直在強調的.
當耶穌基督將這些真理講得清清楚楚明白的時候.
這個婦人又開始走了.
她不願意被那種真理去綑綁.
於是她就說.
我知道耶律賽亞會來.
就是那個稱為基督的那個要來.
他來了之後.
要將一切事都講清楚.
她還是停留在有一天.
但是有一天她不相信是今天.
有一天微賽亞會來.
她不相信現在站在她身邊的就是.
所以她的意思是說.
不要說了.
說到這裡.
總有一天我們明白的.
你是先知.
沒錯.
你說得很好.
但是有一天那位來到了.
一切大家都清楚.
不用再爭辯了.
於是耶穌基督馬上就將她拉回頭.
祂說.
這位跟你說話的.
就是你所盼望的.
你所等候的.
就是等候我.
我現在就是站在你面前.
很多時候就是這樣.
我們總以為有一天神會將恩賜給我們.
其實我們一直都沒有察覺到.
神的恩賜一直都在我們的身邊.
祂用祂的恩典環繞著我們.
我們很多時候是沒有這種屬靈的察覺的.
今天所領受的.
可能是神已經一早給了我們的恩典.

$^{1201}$我們常常以為要等到有一天.
那一天才得到那份很大很大的恩典.
你有沒有想過.
每一天都是恩典.
今天所領受的.
可能是神一早就為你預備了一切.
很好的恩典.
那一天.
其實就是今天一天.
這是很大的恩典.
所以耶穌基督說.
事後將到.
如今就是.
神的工作.
不只是在將來.
神的工作現在就在我們生命當中進行.
神的工作.
一直都是在我們身上的動工.
神的恩典一直都環繞著我們.
我們以為要很大的場面.
神才會跟我們說話.
不一定.
現在可能就是神已經跟你說話了.
我們以為要遇到貴人的時候.
才是好機會.
但是你有沒有想過.
每一次跟人交談.
都是每一次神給我們的恩典.
每一天都是神的恩典.
可能你的兒子.
你的孫子很調皮.
很可能他就是將來的牧師.
這不是一個天方夜譚.
在我們教會當中就有人是這樣.
小時候很調皮.
他跟兒子打架而已.
從小玩到大.
現在呢.
現在很敬虔.
做了牧師.

$^{1241}$當時你就要察覺到.
這一切都是神的恩典.
要把握好機會.
做好經典的榜樣和教導.
當我們有這樣的察覺的時候.
察覺神是真實的.
神現在就是我們生命當中的動工的時候.
很多時候我們就是在這個過程當中去成長.
當撒瑪尼亞夫人.
要跟耶穌結束.
他們談話要結束的時候.
門徒突然就回來了.
他們看到耶穌基督和撒瑪尼亞夫人坐在一起.
按照當時來說.
耶穌基督不應該跟祂一起.
而且說這麼多話.
因為拉比不是在公開的場合.
跟人去隨便聊天的.
但是因為天氣這麼熱.
因為耶穌基督在井旁游泳.
這個又是在情理之中的.
所以門徒雖然看到.
覺得有點驚訝.
但是沒有說你為什麼要跟祂說話.
沒有這樣說.
因為這是正常的關係.
耶穌基督做事都是規規矩矩.
很有智慧的.
就在這個時候.
突然撒瑪尼亞夫人逃跑.
走了.
這個走和之前一直逃跑不一樣.
為什麼?.
因為祂知道祂已經不可以再逃跑了.
祂幾次轉換話題.
幾次想將話題帶開.
但是耶穌基督指出祂的罪.
祂講宗教.
講敬拜.
講這座山.

$^{1281}$講耶路撒冷.
耶穌基督一次一次都將祂帶回到真理裡面.
耶穌基督已將祂這個人看透了.
祂知道祂過去的婚姻是怎樣的.
祂現在的生活是怎樣的.
這麼混亂的.
祂的內心是怎樣的空虛.
多麼的饑渴.
耶穌基督.
剛剛所說的.
真正的敬拜不是在宗教的地點.
是在心靈和誠實的裡面.
這句話直接就碰觸祂的生命.
為什麼?.
因為祂的一生太多的不誠實.
祂的一生當中太多的事不是在真理裡面.
太多的事反而是在放縱情慾那裡.
任意而行的裡面.
今天祂必須要回到真理的裡面.
耶穌基督說.
「事後將到,如今就是」.
就是說你不要再等了.
有一天神工作.
現在神工作已經開始了.
耶穌基督已經在你面前了.
雖然撒瑪利亞人和猶太人.
有很多不同的不同.
有敬拜地方的不同.
但是他們都是盼望這一位耶穌基督的到來.
撒瑪利亞人長久被人們輕看,被排擠.
他們盼望以色列恢復.
也都恢復他們撒瑪利亞人成為以色列人的身份.
他原先只不過是想出來打水的.
但是經文告訴我們.
他的水貫子他都不要了.
他就走了.
水貫子的貫對他的生活是很重要的.
他就丟下了.
這個細節告訴我們.
他生命的優先次序已經改變了.

$^{1321}$他天天都去奔波.
天天都去躲避那些人.
天天都去避開那些人.
人知道他的過去.
只是沒有當面跟他說.
他中午才出來打水.
想保留一點面子.
但是當他真正遇到耶穌基督之後.
他不再去躲藏了.
他直接衝回城裡.
大聲說.
你們來看.
有一個人將我掃來所向的.
我掃來所向的吳光彩.
一切的事都給我說出來了.
莫非這個就是基督.
他很有智慧說話.
這句話很震撼.
因為在當時的文化當中.
沒有人願意揭醜.
特別是一個婦人.
她怎會願意將自己的名聲.
作為一個代價去做這件事.
但是這個婦人竟然將自己的羞辱.
公開在人前.
一個真正悔改的人.
她願意放下自己.
為主去做見證.
真正的敬拜.
不是停留在內心那種感動.
而是有生命的那種真正的改變.
一個人在基督裡面得到釋放之後.
敬拜不僅是停留在教堂裡面.
每天都回來.
我們今天回來敬拜.
領受了神的赦罪.
真理和力量.
一到六.
我們可能大部分時間都沒有回來.
在世界那裡傳福音.

$^{1361}$為主做見證.
為主而活.
全地都可以敬拜.
不僅是在教堂裡面.
不是圍在一個四幅牆裡面.
我們前段時間有一個.
是中國大陸來的.
在河北那邊.
不知是河南還是河北.
他來到就是要照顧他的兒子.
因為他的兒子撞車.
已經不能動了.
只剩下呼吸,插上管.
沒有反應.
那個父親很痛苦.
有時來到教會.
我們星期四值班.
他又來到那裡.
我們抱怨說.
為什麼教堂不是二十四小時開放呢?.
讓我們回來親近神呢?.
為什麼不是每天都開著門給我?.
我這麼痛苦.
我要回來這裡.
我要找一個地方躲起來.
去敬拜神.
我家裡又不方便.
所以我們星期三.
我們都一直為他禱告.
一年多之後.
他的兒子很奇蹟地.
神很有恩典.
讓他的兒子恢復得很好.
現在回到中國了.
這家人在這個過程當中.
有很多軟弱.
但是他們也想過逃避.
但是藉著神.
他們依然在地上為主做見證.
隨時禱告.

$^{1401}$隨時讚美.
因為真正的敬拜不是局限在一個地方.
而是用心靈誠實去敬拜.
所以基督徒是一群領受恩典以後.
繼續去傳揚救恩的人.
想不到這位撒瑪利亞婦人.
一開口全城就都來到了.
她原來是一個這麼有影響力的人.
她這麼會傳道的人.
她真的好像一個宣道家,一個報道家一樣.
在某種程度來說.
若拿一開口.
使尼尼微全城的人悔改.
她不在若拿之下.
在她的人生當中.
山重水復.
但是神讓她柳暗花明.
當她遇到耶穌基督的時候.
她的生命就變了.
東正教有一個傳統.
撒瑪利亞婦人被尊稱為光明的人.
因為她遇到耶穌基督.
黑暗進入到光明.
逃避進入到.
變成一個見證.
從打水變成到傳揚活水.
真正的敬拜不是在哪間教堂.
哪個宗派.
哪種形式.
真正的敬拜是在真理的裡面.
在聖靈的裡面.
去敬拜我們這位父.
我們現在就要回應耶穌基督.
讓我們的生命因著耶穌基督而改變.
我分享到這裡.
\newpage



\section{腓利門書 1:1-25}
\label{sec:032207}
\textbf{腓利門書:接納你的弟兄}~~~~劉少康牧師
\newline
\newline
連結: \href{http://www.church.com.hk/acms/content.asp?site=cdc\&op=show\&type=product\&code=032207\&layout=sermon}{\texttt{032207}} ~~~~ 語音日期: 2026-05-27
\newline
\newline
\hyperref[sec:032206]{\small{< < < PREV SERMON < < <}}
~
\hyperref[sec:index_chronic]{\small{[返順時目]}}
~
\hyperref[sec:032208]{\small{> > > NEXT SERMON > > >}}
\newline
\newline
腓利門書 1:1-25
\newline
\begin{longtable}{cl}
\hline
\hline
章節 & 經文 (和合本修訂版)\\
\hline
1:1 & \begin{tabularx}{0.7\textwidth}{X} 為基督耶穌被囚的保羅，同弟兄提摩太，寫信給我們所親愛的同工腓利門、 \end{tabularx} \\ \\ \relax
1:2 & \begin{tabularx}{0.7\textwidth}{X} 亞腓亞姊妹，和我們的戰友亞基布，以及在你家裡的教會。 \end{tabularx} \\ \\ \relax
1:3 & \begin{tabularx}{0.7\textwidth}{X} 願恩惠、平安從我們的父神和主耶穌基督歸給你們！ \end{tabularx} \\ \\ \relax
1:4 & \begin{tabularx}{0.7\textwidth}{X} 我在禱告中記念你的時候，常為你感謝我的神， \end{tabularx} \\ \\ \relax
1:5 & \begin{tabularx}{0.7\textwidth}{X} 因聽說你對眾聖徒的愛心，和你對主耶穌的信心。 \end{tabularx} \\ \\ \relax
1:6 & \begin{tabularx}{0.7\textwidth}{X} 願你與人分享信心的時候，能產生功效，讓人知道我們所行的各樣善事都是為基督做的。 \end{tabularx} \\ \\ \relax
1:7 & \begin{tabularx}{0.7\textwidth}{X} 弟兄啊，由於你的愛心，我得到極大的快樂和安慰，因為眾聖徒的心從你得到舒暢。 \end{tabularx} \\ \\ \relax
1:8 & \begin{tabularx}{0.7\textwidth}{X} 雖然我靠著基督能放膽吩咐你做該做的事， \end{tabularx} \\ \\ \relax
1:9 & \begin{tabularx}{0.7\textwidth}{X} 可是像我這上了年紀的保羅，現在又是為基督耶穌被囚的，寧可憑著愛心求你， \end{tabularx} \\ \\ \relax
1:10 & \begin{tabularx}{0.7\textwidth}{X} 就是為我在捆鎖中所生的兒子阿尼西謀求你。 \end{tabularx} \\ \\ \relax
1:11 & \begin{tabularx}{0.7\textwidth}{X} 從前他與你沒有益處，但如今與你我都有益處。 \end{tabularx} \\ \\ \relax
1:12 & \begin{tabularx}{0.7\textwidth}{X} 我現在打發他回到你那裡去，他是我心肝。 \end{tabularx} \\ \\ \relax
1:13 & \begin{tabularx}{0.7\textwidth}{X} 我本來有意將他留下，在我為福音所受的捆鎖中替你伺候我。 \end{tabularx} \\ \\ \relax
1:14 & \begin{tabularx}{0.7\textwidth}{X} 但不知道你的意見，我不願意這樣做，好使你的善行不是出於勉強，而是出於自願。 \end{tabularx} \\ \\ \relax
1:15 & \begin{tabularx}{0.7\textwidth}{X} 他暫時離開你，也許是要讓你永遠得著他， \end{tabularx} \\ \\ \relax
1:16 & \begin{tabularx}{0.7\textwidth}{X} 不再是奴隸，而是高過奴隸，是親愛的弟兄；對我確實如此，何況對你呢！無論在肉身或在主裡更是如此。 \end{tabularx} \\ \\ \relax
1:17 & \begin{tabularx}{0.7\textwidth}{X} 所以，你若以我為同伴，就接納他，如同接納我一樣。 \end{tabularx} \\ \\ \relax
1:18 & \begin{tabularx}{0.7\textwidth}{X} 他若虧負你，或欠你甚麼，都算在我的賬上吧， \end{tabularx} \\ \\ \relax
1:19 & \begin{tabularx}{0.7\textwidth}{X} 我必償還。這是我—保羅親筆寫的。我並不用對你說，甚至你自己也虧欠我呢！ \end{tabularx} \\ \\ \relax
1:20 & \begin{tabularx}{0.7\textwidth}{X} 弟兄啊，希望你使我在主裡因你得益處，讓我的心在基督裡得到舒暢。 \end{tabularx} \\ \\ \relax
1:21 & \begin{tabularx}{0.7\textwidth}{X} 我寫信給你，深信你必順服，知道你所要做的，必過於我所說的。 \end{tabularx} \\ \\ \relax
1:22 & \begin{tabularx}{0.7\textwidth}{X} 此外，還請給我預備住處，因為我盼望藉著你們的禱告，必蒙恩回到你們那裡去。 \end{tabularx} \\ \\ \relax
1:23 & \begin{tabularx}{0.7\textwidth}{X} 為基督耶穌與我一同坐監的以巴弗問候你。 \end{tabularx} \\ \\ \relax
1:24 & \begin{tabularx}{0.7\textwidth}{X} 我的同工馬可、亞里達古、底馬、路加也都問候你。 \end{tabularx} \\ \\ \relax
1:25 & \begin{tabularx}{0.7\textwidth}{X} 願主耶穌基督的恩與你們的靈同在。 \end{tabularx} \\ \\
[1ex]
\hline
\hline
\end{longtable}
$^{1}$各位弟兄姊妹,平安!.
很高興今天和大家來看《肥利門醫卷書》.
《肥利門醫卷書》在聖經裡很少有一章經卷.
《肥利門醫卷書》是比較簡單.
但裡面的內容卻是很寶貴,很豐富.
講到有一個主人,他有一個奴隸.
誰知後來那個奴隸信了耶穌.
保羅也是帶有錢的信徒信耶穌.
保羅就勸他們.
既然大家都是兄弟,大家都是信主的弟兄.
那你可不可以接納這個奴隸呢?.
雖然他曾經可能得罪過你.
做了一些錯事,他就走了.
那你可不可以把他挽回過來.
接納他,接納他,當弟兄那樣.
所以保羅就提出一個頗高的要求.
希望這個阿彌西某.
或者現在新的譯本就譯他做阿彌西某.
這個主人希望他原諒.
接納他以前走了的奴隸.
就是這個阿彌西某.
就是他的奴隸阿彌西某.
而這個主人就是「肥利門」.
就是「肥利門書」這個主角.
讓我首先來介紹一下「肥利門」.
「肥利門」就住在哥羅西這個城裡面.
所以我們都知道有一本書叫做「歌羅西書」.
那「肥利門」他是一個比較富有的信徒.
即是家裡面比較有錢.
所以他可能就養了一些奴隸.
其中一個奴隸就是這個阿彌西某.
誰不知這個阿彌西某.
可能就是拿了主人一些財物.
就走了.
這個主人當然很不開心.
你拿了我的錢又走了.
於是就看一下怎樣去處理.
後來這個阿彌西某就遇見保羅.
結果保羅就帶了他順主.
在現在研究聖經比較新的一些派別.

$^{41}$就不是說這個阿彌西某是拿了主人的錢.
而是說他和主人在相處上面.
出現了一些問題.
所以他就走了.
保羅他覺得既然大家都信了耶穌.
他就很想做中間人去調停.
大家都知道現在有些叫做調停人.
我都認識有幾個牧師.
都去讀了調停人的課程.
就拿到調停人的資格.
這個調停人有時可以做中間人.
處理有些人有些財務糾紛.
或者有某方面的爭執.
他就去處理聽聽雙方怎樣說.
就提出一個和好的方法.
讓大家彼此和好.
所以保羅他也都向這個肥利門.
提出一個很高的要求.
就是你可不可以接納阿彌西某.
做你的弟兄.
特別是在肥利門書第十六節.
保羅就這樣要求的.
讓我來少少介紹一下這本書.
這本書雖然很短.
開頭第一到第三節是文安.
保羅就希望肥利門可以在福音的事上幫助他.
由第四節一直到第七節.
他說你有愛啦.
有信心啦.
你要做一些美善的事.
保羅很懂得讚人.
他說你有愛.
有信心.
要做美善的事.
然後由第八節一直到二十節.
就為了阿彌西某去求情.
他說你以前這個奴隸.
曾經做了一些可能令你很不開心的事.
但現在他信了耶穌.
他是我們的弟兄.

$^{81}$你可不可以以愛的緣故去接納他呢?.
因為愛的緣故去接納他.
今天我們在教會裡面.
雖然有些不認識的弟兄姊妹.
甚至有些可能是過去得罪過你的弟兄姊妹.
你能不能夠因為愛的緣故去饒恕他.
去包容他.
在主裡面去接納他呢?.
這個都是我們今天需要去學的一個功課.
然後就由第二十一節到第二十五節.
就一個結束的說話.
保羅他雖然知道他沒有權柄.
去吩咐腓利門做這件事.
所以保羅只是站在一個請求腓利門.
你可以考慮一下.
會不會你原諒這個以前的奴隸.
你接納他成為主來的弟兄姊妹.
你用愛心去這樣做.
沒錯,可能他帶來的是一些經濟的損失.
特別如果他接納了他成為弟兄.
那以後這個亞尼西茂就不再是奴隸.
但是你作為基督徒.
你願不願意付出一些代價呢?.
今天同樣我也想問各位.
你有沒有為基督的緣故.
你願意付上一些代價呢?.
可能在你工作的環境裡面.
有些人曾經得罪過你的.
或者你在工作裡面你又是上司.
有些下屬可能是不合作的.
那你今天會不會因為愛的緣故.
去原諒他們,去接納他們.
去彼此一起同心共作呢?.
所以《肥利門書》都傳遞一個很重要的訊息.
都是很值得我們今天去思想的.
我們有時也碰到一些人.
曾經令我們不開心的.
你能不能夠出於對基督的愛.
你去原諒他,去接納他.
重新跟他和好.

$^{121}$以致你能夠成為他一個主內的兄弟.
大家一起同心走這條天路.
普羅對肥利門的一個請求就是這樣.
他說你是主內的弟兄.
現在這個奴隸雖然做了一些錯事.
但現在因著主的緣故.
因著愛的緣故.
你接納他,你和他一起同行.
你成為他的一個榜樣.
你成為一個愛的源頭.
你傾出愛,接納你的弟兄.
今天都是一個充滿紛爭,充滿怨恨.
充滿很多不和的一個社會.
我們周圍都有很多人曾經得罪過你.
你沒有每一位主的緣故.
你原諒他,接納他.
求主都加我們力量.
讓我們成為一個充滿愛心的弟兄姊妹.
去寬恕,去接納曾經得罪我們的人.
讓我們在愛上學習更多的功課.
願天父親自祈福彼國.
\newpage



\section{哥林多前書 12:12-13:13}
\label{sec:032208}
\textbf{「將勤補拙」的肢體相交}~~~~王啟恩牧師
\newline
\newline
連結: \href{http://www.church.com.hk/acms/content.asp?site=cdc\&op=show\&type=product\&code=032208\&layout=sermon}{\texttt{032208}} ~~~~ 語音日期: 2026-05-24
\newline
\newline
\hyperref[sec:032207]{\small{< < < PREV SERMON < < <}}
~
\hyperref[sec:index_chronic]{\small{[返順時目]}}
~
\hyperref[sec:032209]{\small{> > > NEXT SERMON > > >}}
\newline
\newline
哥林多前書 12:12-13:13
\newline
\begin{longtable}{cl}
\hline
\hline
章節 & 經文 (和合本修訂版)\\
\hline
12:12 & \begin{tabularx}{0.7\textwidth}{X} 就如身體是一個，卻有許多肢體，身體的肢體雖多，仍是一個身體；基督也是這樣。 \end{tabularx} \\ \\ \relax
12:13 & \begin{tabularx}{0.7\textwidth}{X} 我們無論是猶太人是希臘人，是為奴的是自主的，都從一位聖靈受洗成了一個身體，並且共享這位聖靈。 \end{tabularx} \\ \\ \relax
12:14 & \begin{tabularx}{0.7\textwidth}{X} 身體原不只是一個肢體，而是許多肢體。 \end{tabularx} \\ \\ \relax
12:15 & \begin{tabularx}{0.7\textwidth}{X} 假如腳說：「我不是手，所以不屬於身體」，它不能因此就不屬於身體。 \end{tabularx} \\ \\ \relax
12:16 & \begin{tabularx}{0.7\textwidth}{X} 假如耳朵說：「我不是眼睛，所以不屬於身體」，它也不能因此就不屬於身體。 \end{tabularx} \\ \\ \relax
12:17 & \begin{tabularx}{0.7\textwidth}{X} 假如全身是眼睛，聽覺在哪裡呢？假如全身是耳朵，嗅覺在哪裡呢？ \end{tabularx} \\ \\ \relax
12:18 & \begin{tabularx}{0.7\textwidth}{X} 但現在神隨自己的意思把肢體一一安置在身體上了。 \end{tabularx} \\ \\ \relax
12:19 & \begin{tabularx}{0.7\textwidth}{X} 假如全都是一個肢體，身體在哪裡呢？ \end{tabularx} \\ \\ \relax
12:20 & \begin{tabularx}{0.7\textwidth}{X} 但現在肢體雖多，身體還是一個。 \end{tabularx} \\ \\ \relax
12:21 & \begin{tabularx}{0.7\textwidth}{X} 眼睛不能對手說：「我用不著你。」頭也不能對腳說：「我用不著你。」 \end{tabularx} \\ \\ \relax
12:22 & \begin{tabularx}{0.7\textwidth}{X} 不但如此，身上的肢體，人以為軟弱的，更是不可缺少的； \end{tabularx} \\ \\ \relax
12:23 & \begin{tabularx}{0.7\textwidth}{X} 身上的肢體，我們認為不體面的，越發給它加上體面；我們不雅觀的，越發裝飾得雅觀。 \end{tabularx} \\ \\ \relax
12:24 & \begin{tabularx}{0.7\textwidth}{X} 我們雅觀的肢體自然用不著裝飾；但神配搭這身子，把加倍的體面給那有缺欠的肢體， \end{tabularx} \\ \\ \relax
12:25 & \begin{tabularx}{0.7\textwidth}{X} 免得身體不協調，總要肢體彼此照顧。 \end{tabularx} \\ \\ \relax
12:26 & \begin{tabularx}{0.7\textwidth}{X} 假如一個肢體受苦，所有的肢體就一同受苦；假如一個肢體得光榮，所有的肢體就一同快樂。 \end{tabularx} \\ \\ \relax
12:27 & \begin{tabularx}{0.7\textwidth}{X} 你們是基督的身體，並且各自都是肢體。 \end{tabularx} \\ \\ \relax
12:28 & \begin{tabularx}{0.7\textwidth}{X} 神在教會所設立的：第一是使徒；第二是先知；第三是教師；其次是行異能的；再次是醫病的恩賜，幫助人的，治理事的，說方言的。 \end{tabularx} \\ \\ \relax
12:29 & \begin{tabularx}{0.7\textwidth}{X} 難道個個都是使徒嗎？難道個個都是先知嗎？難道個個都是教師嗎？難道個個都是行異能的嗎？ \end{tabularx} \\ \\ \relax
12:30 & \begin{tabularx}{0.7\textwidth}{X} 難道個個都是有醫病的恩賜嗎？難道個個都是說方言的嗎？難道個個都是翻方言的嗎？ \end{tabularx} \\ \\ \relax
12:31 & \begin{tabularx}{0.7\textwidth}{X} 你們要追求那更大的恩賜。我現今把最妙的道指示你們。 \end{tabularx} \\ \\ \relax
13:1 & \begin{tabularx}{0.7\textwidth}{X} 我若能說人間的方言，甚至天使的語言，卻沒有愛，我就成為鳴的鑼、響的鈸一般。 \end{tabularx} \\ \\ \relax
13:2 & \begin{tabularx}{0.7\textwidth}{X} 我若有先知講道的能力，也明白各樣的奧祕，各樣的知識，而且有齊備的信心，使我能夠移山，卻沒有愛，我就算不了甚麼。 \end{tabularx} \\ \\ \relax
13:3 & \begin{tabularx}{0.7\textwidth}{X} 我若將所有的財產救濟窮人，又犧牲自己的身體讓人誇讚，卻沒有愛，仍然對我無益。 \end{tabularx} \\ \\ \relax
13:4 & \begin{tabularx}{0.7\textwidth}{X} 愛是恆久忍耐；又有恩慈；愛是不嫉妒；愛是不自誇，不張狂， \end{tabularx} \\ \\ \relax
13:5 & \begin{tabularx}{0.7\textwidth}{X} 不做害羞的事，不求自己的益處，不輕易發怒，不計算人的惡， \end{tabularx} \\ \\ \relax
13:6 & \begin{tabularx}{0.7\textwidth}{X} 不喜歡不義，只喜歡真理； \end{tabularx} \\ \\ \relax
13:7 & \begin{tabularx}{0.7\textwidth}{X} 凡事包容，凡事相信，凡事盼望，凡事忍耐。 \end{tabularx} \\ \\ \relax
13:8 & \begin{tabularx}{0.7\textwidth}{X} 愛是永不止息。先知講道之能終必歸於無有；說方言之能終必停止；知識也終必歸於無有。 \end{tabularx} \\ \\ \relax
13:9 & \begin{tabularx}{0.7\textwidth}{X} 我們現在所知道的有限，先知所講的也有限， \end{tabularx} \\ \\ \relax
13:10 & \begin{tabularx}{0.7\textwidth}{X} 等那完全的來到，這有限的必消逝。 \end{tabularx} \\ \\ \relax
13:11 & \begin{tabularx}{0.7\textwidth}{X} 我作孩子的時候，說話像孩子，心思像孩子，意念像孩子；既長大成人，就把孩子的事丟棄了。 \end{tabularx} \\ \\ \relax
13:12 & \begin{tabularx}{0.7\textwidth}{X} 我們現在是對著鏡子觀看，模糊不清；到那時，就要面對面了。我如今所認識的有限，到那時就全認識，如同主認識我一樣。 \end{tabularx} \\ \\ \relax
13:13 & \begin{tabularx}{0.7\textwidth}{X} 如今常存的有信，有望，有愛這三樣，其中最大的是愛。 \end{tabularx} \\ \\
[1ex]
\hline
\hline
\end{longtable}
$^{1}$在我小學三年級 年中的成績表.
老師對我有一個評語.
這個評語是什麼呢.
今天的講題 將勤補拙的肢體相交.
主要經文在哥林多前書十二章.
這個評語是什麼呢 這個評語是四個字的.
和今天的講題很有關係 大家一起讀出來.
將勤補拙 不知道大家如何理解.
我印象很深刻 我有兩個理解.
第一個是正面的 正面的意思是老師讚賞我的勤力.
就是說我可以彌補我不足之處.
第二個是負面的 我表現得不好.
有不足的地方 需要更加努力去彌補這些不足.
當時確實是有些不開心 負面的情緒比較大.
但是我始終到今時今日 都沒有忘記這一句話.
我自知自己是資質平庸的.
但是我提醒自己要努力不懈.
因為努力必定會有收成.
我最近問了AI 什麼意思.
AI就這樣說.
他說這一句是一則強調後天努力可以克服先天限制的成語.
事實上都是用來鼓勵 透過加倍付出時間和心力去達成目標.
AI就這樣說.
這是一句鮑爾詞 即是讚賞的 不是貶義的.
所以當年的老師其實是對我讚賞的.
在肢體相交的事情當中.
弟兄姊妹 我們其實都是需要將勤補拙.
我相信將勤補拙的將勤是在說什麼呢.
就是我們要事實上是依靠我們的神.
而補拙的意思就是幫助我們克服天生自我中心那種傾向.
亦都是說到犯罪的傾向.
以至到教會可以有美好的肢體生活.
讓神得到祂所得的榮耀.
一個美好肢體相交的基礎是什麼呢.
我們一起讀這兩節經文.
就如身體是一個 卻有許多肢體.
身體的肢體最多 仍是一個身體.
基督也是這樣.
我們無論是猶太人 是希臘人 是維路的 是自主的.
都從一位聖靈受洗 成了一個身體.

$^{41}$並且共享這位聖靈.
經文講到人的身體和肢體的關係.
肢體雖然是多 但是不可以分割的.
因為是同屬於一個身體.
保羅用這個比喻講到很明白.
這裡有六個字就是基督也是這樣.
都是一樣的.
意思就是說基督的身體 即教會都是一樣.
而當中的肢體是哪些人呢 大家認為.
就是我們這群信徒.
所以你跟隔壁的朋友說 你也是肢體.
是不可以分割的.
所以無論你覺得自己的作用很大.
又或者你覺得自己是沒什麼作用的.
但是都是有份於這個肢體.
是不可以分割的.
我們是在一起的.
這是第一個重點.
同屬於一個身體是不可以分割的.
而第二個同屬就是同屬於一位聖靈.
同屬於一位聖靈的洗.
意思就是說我們信主那一剎那.
我們已經有聖靈內著.
正如《耳法所書》一章十三節這樣講.
既然信他 就受了所應許的聖靈為印記.
所以我們信了耶穌 我們每一個人已經有聖靈內著.
使我們成為基督的一個部份.
並且在這裡講到我們可以共享這個聖靈.
意思就是說我們可以在經歷上.
去分享聖靈的同在和大能.
同時之間我們可以領受.
從聖靈而來的大大小小的恩賜.
所以我們不可以說我沒有聖靈.
大英姐妹你和隔壁左右說.
你不可以說你沒有聖靈.
我們每一個人都有聖靈.
問題就是我們用聖靈和我們裡面的靈去溝通.
還是你將它隔絕.
但是無論如何我們都有聖靈.
就好像我們一定有肉身父母一樣.

$^{81}$這是一個事實 是無可推諉的事實.
在現今可能一些很不幸的孤兒.
他們都有父母一樣.
恩賜的目的是什麼呢.
主要有兩個.
第一個就是我們一起讀.
聖靈彰顯在各人身上.
是要使人得益處.
使人得益處.
意思就是說使整個救贖群體.
就是我們這一班人.
不是為了自我標榜自己的滿足.
而出發點是對神和人那種的服侍.
最終的目標就是為了去榮耀我們的神.
另外第二個目標.
看到這個身體是裝滿人的.
就是你你我我.
他說這裡所設的有使徒.
有先知 有傳福音的.
有牧者 有牧師.
為要去裝備聖徒做事奉的工作.
建立基督的身體.
紅色的字.
所以很明顯這裡說到我們要建立基督的身體.
但我們不可以忽略之前那兩句.
之前那兩句寫什麼.
為了要裝備聖徒做事奉的工作.
就是說這些所有所有的恩賜.
有這些恩賜的人.
是需要裝備聖徒做事奉的工作.
那聖徒是誰呢.
大家認為.
就是你你我我.
那跟旁邊的左右說 是你 是我.
我們是聖徒.
蒙恩基督教的信徒都是聖徒.
經文很清楚的說.
教會領袖.
就是我其中一個.
不是代替聖徒或者信徒去事奉.

$^{121}$這裡說得很清楚.
是要裝備信徒.
使他們可以做事奉的工作.
重點的真理就是這五個字.
就是說全民應該是一起去事奉的.
所以事奉不是一些個別的人.
或者少數核心同工.
屬靈精英去做的事情.
或者是責任.
而是全體信徒的本分和義務.
為什麼.
剛才說得很清楚.
因為我們每一個信徒都是肢體.
都是屬於這個身體.
都有從聖靈而來.
大大小小的那種恩慈.
所以在這個真理下.
我們是需要有美好的相交.
以致我們可以互相配合.
合一的心去事奉我們的神.
事奉我們其他的人.
問大家會不會有挑戰.
在這個過程當中.
一定有.
好多好多的挑戰 好多好多的難是阻礙著的.
人生有挑戰是必然的事.
但問題就是我們怎樣將這些挑戰.
去化為機遇.
同時之間可以去促進我們的靈命成長.
就是這個問題.
主要有兩個情況.
在經文那裡說了.
第一個 慶賀自己.
事不關己.
這裡這樣說.
假如腳這樣說.
我不是手 所以不屬於身體.
他不能夠因此就說.
我不屬於這個身體.
如果耳朵這樣說.

$^{161}$我不是眼睛.
所以不屬於這個身體.
他都不能夠因此就說.
我不屬於這個身體.
經文好 擲地有聲這樣說.
縱然好像看起來.
手是比較靈活一點.
同時之間好像做一些很繁複的.
操作或工作.
或者可以令到腳會自慚形愧.
覺得自己的功能比較低.
但腳不可以說.
喂 我不屬於這個身體.
因為他本身已經是屬於這個身體.
我們本身已經是屬於這個身體.
因為我們每一個人信了耶穌.
已經有聖靈的印證.
已經chop印了.
而事實上.
腳是有手不能夠做的功用.
腳要用來行路.
但大家都可以說.
反過來用手行路就可以了.
可以嗎.
可以的 不過很艱辛.
一般都不會這樣做.
大家認為盲腸有沒有用.
大家可能會說.
吃完飯不要跳 跳窄的時候.
就會有盲腸炎.
有些醫學的證明.
就是這樣說.
盲腸是有很重要的功能.
能夠幫助人體免疫系統的運作.
一個防禦的角色.
可以儲存和修復腸道益生菌.
防止特定的感染.
所以不應該隨便割掉.
同樣的 弟兄姊妹.
千萬不要覺得自己只是一個小畜生.

$^{201}$一個小組員.
一個崇拜的小會眾.
覺得自己沒什麼恩賜.
又或者覺得沒什麼可以做到.
又或者覺得自己不被需要.
或者不被去重視.
而去輕看自己.
很多人事奉.
少了自己也不要緊.
大家千萬不要這樣去想.
因為剛才也說過.
神蹤每個人都是獨特的.
你跟鄰居說 你是很獨特的.
你是很特別的.
都有它的作用.
是不可以被替代的.
而我們需要的就是這四個字.
將勤補拙.
怎樣將勤補拙呢.
第一個就是.
無論你的恩賜有.
你覺得自己有多小.
你總是有可以用的地方.
你可以做的勤於去盡自己的本份.
忠心為主 拜上你那份的恩賜.
這個應該是謙卑先.
這是第二點 待會再說.
謙卑先 接受自己的不完美.
不完美 不需要強求成為別人.
還有一個很重要的.
很想和大家分享.
你自己不要拆毀自己.
總有一天你不完美.
你不要跟自己說 我還是不行了.
我還是沒有什麼才能 才剛.
別人都看不起我 等等等等.
而自暴自棄.
妄自誹謗.
甚至有一個極端.
就是很想做一些事情.

$^{241}$去證明自己是可以的.
但很多時事實告訴我們是適得其反.
有一次我有一個犯罪.
就是吃飯的聚會.
不是那種犯罪.
吃飯的聚會.
就是我們一班同學.
上完課就有一個晚餐.
教授就和我說.
Larry不如你祈禱謝飯吧.
當時同台又有一兩個牧師在.
我就和他說.
我不是牧師 是不是不應該我去謝飯.
接著那位教授就說.
Larry你不要妄自誹謗.
你不要妄自誹謗.
等等.
沒有人可以拆毀你的.
甚至如果有些人想去拆毀你.
如果你不拆毀自己.
是沒有人可以拆毀你的.
所以我們不要輕看自己.
第二個就是剛才說了.
無論恩賜有多小.
總是有個地方你可以去用的.
你要做的就是勤力.
盡自己的本份.
忠心為主.
擺上你那份恩賜.
你可以做的是什麼.
勤勞點去服侍.
多點去禱告.
多點去傾聽別人.
多點去連結等等.
還有你可以去做的.
就是捕捉別人的需要.
大家不要小看.
一句的讚賞.
一句的鼓勵.
一句的關心.

$^{281}$可以產生很大.
即是你意想不到的.
很大很大.
那種效果.
我曾經做牧者.
時常都會關心別人.
都會鼓勵別人.
甚至有些電視節目做得好.
我會讚賞他.
特別是有些在困難當中.
我都會鼓勵.
很多時候他們會回來跟我說.
啊 Larry.
真是多謝你.
那天鼓勵我.
那天為我禱告.
我已經過了那些困難.
所以弟兄姊妹.
這些可以去做.
是一句話.
一句讚賞.
一句鼓勵.
一句關心.
其實都可以產生很大.
很大很大的作用.
同時之間.
另外一個可以做的就是陪伴.
說回二十年前.
二十年前我曾經參加過一個研討會.
是成為分享的嘉賓.
這個研討會是說到.
怎樣去幫助一些在職的基督徒.
那天是去介紹一個項目.
我在這個項目當中.
都是有少許的研究.
和有參與在當中.
於是那個機構就找了我去做分享.
完成了之後.
那個負責人就打電話給我.
Larry 有些事情想和你說.

$^{321}$什麼事情.
出來談吧.
接著他就說.
殘忍的事情來了.
我認為.
什麼事情.
原來是有問卷調查的.
問回那天的研討會.
大家有什麼感覺.
覺得好還是不好.
同時之間.
對於分享嘉賓.
你有什麼意見呢.
這次那個負責人就和我說.
Larry 你可以不聽.
如果你不聽.
我不會說的.
接著我說.
古惟勇說.
當然要聽.
我心情很緊張.
大家緊不緊張想聽.
想聽是吧.
OK 好.
既然大家這麼想聽.
我就和大家分享.
評語只有兩個字.
膚淺.
對了.
當時好像當頭棒喝.
我當時真的笑不出來.
跟著和自己說.
是不是自己不行呢.
是不是自己是.
所謂道行未深.
在這個項目裡面.
我認知不夠多.
甚至我曾經想過.
我都不再在眾人面前.
去分享或者說話.

$^{361}$但是姐妹.
人很容易被人一句話.
去打沉的.
千萬不要這樣.
大家覺得我有沒有被這句話打沉呢.
應該沒有吧.
很明顯我今天可以站在這裡.
和大家分享.
忠賢耕那位弟兄.
給評語我的弟兄.
是出於好意.
我相信也是.
神沒有去撇下我.
沒有說.
Larry 你這麼差.
這麼差.
你說話這麼膚淺.
祂繼續去使用我.
之後也有很多機會.
今天我可以在這裡和大家分享.
全是恩典.
但是我不覺得自己仍然很厲害.
因為我仍然覺得自己.
口齒不靈.
我只是盡上我自己的本份.
去好好地預備.
去和大家分享.
第二個.
就是自視過高.
目中無人.
這裡說到什麼呢.
就是說到眼睛不能夠和手這樣說.
我用不著你.
頭也不能和腳這樣說.
我用不著你.
接著就說.
經文就這樣說.
不但如此.
身體上的肢體.
人以為軟弱的.

$^{401}$更加是不可以缺少.
弟兄姊妹.
眼和頭都有一個共通點.
共通點是什麼呢.
就是引導手和腳.
去做動作.
那種的作用.
看起來就好像比較重要.
因此好像比較高.
但是如果沒有了手和腳去配合.
眼和頭都不可以發揮它的功效.
所以弟兄姊妹.
我們是不可以擺出.
縱然之間.
你有多厲害.
有多高超.
多優秀.
等等等等.
你的事奉有多有果效.
等等.
你傳福音有多厲害.
我們都不可以擺出.
一副我用不著你.
那種的態度.
而當時哥林多教會.
確實是有這種的情況.
弟兄姊妹.
人心沒有改變過.
這種的情況.
時刻都會發生在教會裡面.
我見過不少.
這種的優越感.
就像這幅圖.
上面看著下面.
那種優越感.
會產生什麼不好的作用呢.
就是這樣.
這個人.
站起來都很困難.
意思就是說.

$^{441}$他會將整個身體.
帶來一些機能上的衰退.
甚至是癱瘓.
因為這種的態度.
不讓一些肢體去貢獻.
特別不讓一些軟弱的肢體去貢獻.
所以沒有了他的貢獻.
沒有了這些肢體的貢獻.
身體自然就會衰退.
自然會退化.
有一條珍珠項鍊.
很漂亮的項鍊.
掛在一個女士的頸上.
這個女士就在街上行走.
在陽光的底下.
閃閃發光.
路過的人都讚口不絕.
很漂亮 很漂亮 很漂亮.
珍珠聽到.
然後就很開心.
引以為豪.
然後跟那條線說.
線啊 線啊.
如果不是我使用你.
你什麼用處都沒有.
於是那條線就開始.
嘰哩咕嚕 怒氣沖天.
那條線就斷裂了.
斷裂了 那些珍珠就掉在地上.
他失去了他引以為傲的光彩.
OK 弟兄姊妹.
光林多前書十二章二十二節.
就是講到這件事.
不但如此.
身體上的肢體.
引以為軟弱的.
是不可以缺少那個的意思.
所以自視過高的肢體.
也是需要做這件事.
就是將勤補.

$^{481}$但是這個絕不是那個絕.
為什麼用這個絕呢.
有意思的 弟兄姊妹.
就是我想突顯.
這種引以為絕的優越感.
其實都是需要調教的.
怎樣調教呢 第一.
當然是放下自己的身段.
聖經也有這樣說.
我們要看別人給自己.
什麼啊 強.
我們要多點欣賞別人.
我們不要自命我們是不凡.
為什麼呢.
因為所有的恩賜.
都是從同一位聖靈而來的.
我們沒有什麼比別人厲害.
同時之間.
我們要認清楚.
無論你的恩賜有多大.
你總是有不能夠勝任的地方.
第二個.
我們要勤於彼此去照顧.
這裡這樣說.
身體上的肢體.
我們不雅觀的.
愈發裝飾得雅觀.
我們雅觀的肢體.
自然就不需要裝飾了.
彼此不需要裝飾了.
神配搭這個身子.
把加倍的體面.
給那些有缺乏的肢體.
免得身體不協調.
總要肢體彼此相顧.
首先問問大家.
人的天性.
是不是對於身體不好的部分.
會加倍注意和照顧.
務求裝飾得好.

$^{521}$變得漂亮一點.
雅觀一點.
給大家一個例子.
相信大家都會認同.
不少人去做面部護理.
有沒有.
為什麼我們要去做面部護理.
就是想漂亮一點.
雅觀一點.
昨天不是說面部護理.
說錯了.
我是說什麼.
我是說醫美.
有一位姐妹給我意見.
我聽意見.
她說不是醫美.
醫美說的是整容.
我們沒有去整容.
我們只是做面部護理.
等等等等.
有時候我們都會去遮掩身體上.
一些不雅觀的部位.
我最近生了一顆暗瘡.
不知道大家有沒有看到.
暗暗的.
現在已經暗了.
我跟自己說.
生暗瘡其實已經是好事.
對於我來說.
因為代表青春.
沒錯.
大家說得好.
但同時我也很想這個印.
快點消失.
有沒有什麼方法.
可以介紹給我.
怎樣去快點消失這個印.
各位姐妹.
我們對待一些比較軟弱的肢體.
也是一樣.

$^{561}$我們必須要很樂意特別去照顧他們.
所以肢體總要彼此相顧.
之前所說.
因為這是神所配搭的身體.
是有神的主權在他的手裡.
我們要明白清楚這一點.
所以我們唯有這樣.
教會才會真正的合一.
另外一個第三點.
我們需要將勤補拙的是什麼.
就是勤於接納包容.
我們要建立不是拆毀別人.
在耶利米蘇二十四章六節說到.
耶和華對那班被擄去到巴比倫的子民.
以色列人這樣說.
我要建立你們.
不是拆毀你們.
這是神的心意.
對我們每一個人也是的心意.
所以我們沒有這樣的權利去拆毀別人.
因為神的心意是建立.
不是拆毀.
人很容易去看低別人.
不知道大家同不同意.
說兩個例子.
如果你坐巴士或者坐地鐵.
旁邊坐了一個不斷咳嗽的人.
他不戴口罩.
你會怎樣.
你會這樣嗎.
會這樣嗎.
這邊 會這樣嗎.
大家會這樣嗎.
未必一定很興旺.
但你會覺得有沒有搞錯.
我通常會這樣.
我出街會多帶一個後備口罩.
後備口罩就是為了預備這些人.
如果他咳嗽的話.
又沒有戴口罩.

$^{601}$如果他坐在旁邊.
我就會給他一個口罩.
第二個.
地鐵裡面有關外座位.
有一次我見到一個.
應該是小學生.
他跳蚱蜢進來地鐵裡面.
坐在關外座位.
同時我見到前面有一個比較老的老人家.
他沒有讓座位給老人家.
大家又會怎樣呢.
可能你都會看著他.
我就鼓起勇氣.
跟這個小朋友這樣說.
小朋友 你剛才跑進來.
這麼有精力 這麼有活力.
前面有個老人家.
你會不會讓個位給他坐.
他很需要這個座位.
你猜這個小朋友怎樣做.
馬上站起來 彈起來.
讓了這個位.
我們不去說.
這個小朋友可能是不知道的.
我們很容易會輕看別人.
上次這個惡目典禮.
在這個惡目典禮的分面.
謝牧師這樣去介紹我.
Larry 和任何人都可以合拍.
一起工作.
無論他的性格是怎樣.
無論其他人的年齡層是怎樣.
我都可以和他們一起去事奉.
接著就說 Larry 真是一個高人.
大家有沒有聽過.
在這個惡目典禮應該有聽過.
兄弟姐妹 多謝謝牧師對我的褒獎.
我不可以說他是高人.
如果說他是高人的時候.
真是高高在上.

$^{641}$覺得自己有很高很高.
很高尚的恩賜.
其實我不是.
我不覺得自己有很高尚的恩賜.
我覺得這是神給我的恩賜.
其實我很想做的就是.
努力去活出和所有的人彼此和好.
和神有和好的生命.
就是這樣.
因為聖經也有說.
我們要盡力使人和睦.
兄弟姐妹.
我不是想要做一個很好的先生.
大家所說的話.
我全部都接納了.
我要去認同別人所說的話.
我沒有了自己的主見等等.
我不是要做一個很好的先生.
而是我很想去竭力.
做一個使人和睦的人.
最後一個總結給大家.
美好的肢體相交的精義是什麼呢.
大家也很清楚.
彼此的相愛.
這裡說到你們要追求更大的恩賜.
我現今把最妙的道指示給你們.
這節經文是保羅.
在參考教會恩賜和肢體運作之後.
所作出的總結.
我們看看這些字.
追求是什麼意思呢.
這是一個命令式.
不是說你做就行不做就行.
是一定要去做的.
是要竭力去做的.
那更大的恩賜是說什麼呢.
一般學者認為能夠造就教會.
造就更多人的恩賜.
就是那種恩賜.
你就去求多一點吧.

$^{681}$最後說到最妙的道.
什麼是最妙的道呢.
十二章三十一節.
是第十二章最後一節的經文.
接著是第十三章.
高林多前書十三章.
相信大家都很熟悉.
最妙的道是說什麼呢.
一個字 愛.
大家可能都會背.
愛是恆久忍耐又有恩賜.
愛是不述道不自誇等等.
所以最妙的道是說愛.
而全段的經文是說這件事情.
這十六個字.
這一節的經文是說這十六個字.
保羅去鼓勵信徒.
我們不應該滿足在.
追求炫耀性的恩賜.
或者把自己的恩賜.
炫耀地說自己是可以.
或者說自己是好.
或者是厲害等等.
而是要追求.
更能夠造就教會.
那種更大的恩賜.
還有.
這個高超的方向.
還有去正確運用.
一切屬靈恩賜的途徑.
就是非愛莫屬.
所以弟兄姊妹.
無論你有多大的恩賜.
你覺得自己的恩賜是多小.
如果沒有了愛作為基礎的話.
在神的眼中就是一文不值.
就正如十三章一節這樣說.
我們一起讀一二三.
命的羅 享的拔一樣.
所以去到最後將勤補拙.

$^{721}$最高超的境界.
就是非愛莫屬.
而愛是凌駕一切的恩賜.
我最近看了一篇文章.
很想和大家分享.
這篇文章是一個國內.
高中三年級的學生寫的.
他準備升大學.
一開始他這樣說.
我不捨得叫我媽媽 爸爸.
買手機給我.
因為我終於知道.
錢是怎樣來的.
接著他又說.
準備升大學的時候.
我很想我爸媽買手機給我.
因為我所有同學都有.
我拿了出來之後.
我覺得自己完全不同了.
很有面子.
爸媽最初不肯買給他.
於是他就絕食抗議.
在絕食抗議當中.
他覺得自己很委屈.
為什麼我生在.
這麼普通的家庭當中.
為什麼爸媽不愛我呢.
最後媽媽心軟.
跟著跟他說.
這樣也行 開出一個條件.
我出七成 你出三成.
三成怎樣拿回來呢.
你跟我和我爸一起去幹活.
一起去工作.
做什麼呢.
搬運買賣的工作.
這樣你就賺回一些錢.
兒子聽到當然說.
好啊 三成而已.
工作有什麼難呢.

$^{761}$第一天四點起床.
不是下午 早上.
凌晨四點起床.
去110公里外的市集.
一個批發市場.
去買西瓜.
爸爸開車就去到那裡.
去到那裡.
原來要搶西瓜.
搶一些比較漂亮的.
新鮮的 紮實的西瓜.
很不容易.
買了一車的西瓜.
但買了一車的西瓜.
是要搬上車的.
跟著他就說.
我搬了十幾個.
已經躺在那裡.
不懂得動了.
但爸爸媽媽仍然繼續搬西瓜.
很不容易搬完西瓜.
去到零售的地方.
或者市場裡面.
賣西瓜.
賣西瓜的途中.
都很多人來.
一個一個的買.
有些人要買半個西瓜.
跟著他就自告奮勇.
我幫忙切西瓜.
大家都猜到了.
他都不切.
切了幾刀.
就切身自己.
雙手.
於是爸爸媽媽看到.
就很心急.
很緊張.
去照顧他.
為他包紮.

$^{801}$那個時候.
他才第一次見到.
媽媽的手是這樣的.
粗糙.
老人家的手.
滿佈傷痕.
裂痕的手.
但媽媽仍然在臉上掛著笑容.
跟著他就問媽媽.
為什麼你經常都笑嘻嘻.
媽媽就說.
降活嘛.
哪裡有不受傷的呢.
就算是多累.
我都不能將累放在臉上.
免得你.
五個字.
惹你不開心.
我不想惹你不開心.
這個惹你不開心.
這五個字.
那個學生就說.
我呆了.
整個人呆了.
我不懂得去表達.
最後.
九點半.
買完西瓜.
就說.
媽媽很開心.
今天生意特別好.
賺了三百七十一元.
同時之間.
他看到媽媽的臉.
是汗流滿面.
那塊臉.
所以他就說.
他意識到.
這三百七十一元.
人民全部都是辛勞.

$^{841}$汗水和愛.
於是他就忍不住.
哭了出來.
媽媽就將錢轉給他.
他就將錢退回給媽媽.
因為他發現.
這些錢他花不了.
他花不了這些錢.
然後他就說.
這場手機爭奪戰.
我輸了.
但我都贏了.
贏了什麼呢.
我最深刻.
就是這三件事.
第一個.
輸了輕飄飄的虛榮.
即是一些炫耀.
他贏得沉甸甸的懂得.
首先我想說.
他那種絕食抗議的行動.
是代表什麼呢.
我覺得是代表輕看自己.
他覺得爸爸媽媽不重視他.
甚至他可能會覺得.
這個家不屬於我.
好像我剛才所分享.
會不會覺得這個肢體不屬於我呢.
大家不妨.
在這個故事裡面去聯想一下.
接著來到這裡.
輸了輕飄飄的虛榮.
即是炫耀.
卻贏得沉甸甸的懂得.
即是覺得.
如果我們覺得自己有很高超的恩賜.
很叻很掂.
而看不起別人.
這些都是炫耀的東西.
輕飄飄的.

$^{881}$毫無作用.
但是.
如果我們懂得這件事的話.
我們就是一生受用.
一生受用.
接著他又說到第二點.
他說愛不是理直氣壯索取的底氣.
而是我們能夠看到.
藏在父母裡面不容易的愛之後.
發自內心的那種心痛和體諒.
他開始體諒到.
觸摸到.
爸爸媽媽那種愛.
弟兄姊妹.
我們有沒有進入耶穌的心腸裡面呢.
去觸摸他的愛呢.
有沒有發自內心.
又或者我們有沒有體會到.
十字架上面.
耶穌基督所受的痛楚呢.
當我們讀一些捷經文.
和看回耶穌的釘痕手.
大家有什麼感受呢.
因他受的懲罰.
我們得平安.
因他受的鞭傷.
我們得醫治.
同時之間.
這一隻的傷痕手.
這一對的傷痕手.
大家有什麼感受呢.
各位學生的反省.
對於我們今天.
我們有什麼提醒自己呢.
第三點.
我發現以前的我.
是多自私和虛榮.
我可以寫出很漂亮的文章.
他懂得感恩.
他觸摸到愛.

$^{921}$他願意付出愛.
他已經不要這部手機了.
看到那個轉變嗎.
那個轉化.
所以我覺得最深最深的.
我自己的領受就是說.
這個少年人.
少年人那種覺醒.
是比任何的大道理都來得深刻.
我們有經歷.
才會有成長.
我們有經歷才會有成長.
有一位牧師和學者牧師.
叫Timothy Keller.
他分享過.
什麼是真正的愛.
在申命記七章七至八節.
他有講到.
他講到什麼呢.
耶和華愛這一班以色列人.
不是因為他們是列國之中.
最偉大的民族.
其實他們是最細小的一個國.
然後耶和華說.
我愛你們.
因為我定義去愛你們.
這個就是我愛你們的原因.
就是這麼簡單.
真正的愛是我愛你.
只因為我愛你.
不是因為你有多厲害.
或者你有多差.
你有多優秀或者你多差.
甚至我也愛你.
而這個愛是說恩典的愛.
所以如果有你的親人.
特別是老公.
如果太太問你.
你愛我媽嗎.
你千萬不要這樣回答.

$^{961}$我愛你因為你漂亮.
為什麼不要這樣說呢.
因為她會反問你.
如果老了不漂亮怎麼辦呢.
你愛她就因為你愛她.
因為她是你的老婆.
就是這麼簡單.
所以弟兄姊妹.
如果沒有愛.
任何事物都會失去意義.
教會不是一班完全的人的聚集.
我們在座之中.
這一個是完全的人.
但願我們每一個人.
都能夠對主耶穌基督的愛.
有深入的覺醒.
好像這個少年人一樣.
並且體會到.
我們每一個肢體.
都是耶穌基督所配搭的.
都是有它的用處的.
以致我們這班不完全的人.
可以因著愛.
去努力的去連結在.
耶穌基督的身體裡.
彼此相顧.
將勤補拙.
去建立耶穌基督的身體.
使教會和更多人.
得到造就.
耶穌基督得到榮耀.
因為畢竟.
這一個家.
都是大家的.
這一個家都是大家的.
\newpage



\section{約書亞記 24:15}
\label{sec:032209}
\textbf{聖經人聖經事:約書亞贈別 (3) 我家事奉神}~~~~黃克勤牧師博士
\newline
\newline
連結: \href{http://www.church.com.hk/acms/content.asp?site=cdc\&op=show\&type=product\&code=032209\&layout=sermon}{\texttt{032209}} ~~~~ 語音日期: 2026-05-28
\newline
\newline
\hyperref[sec:032208]{\small{< < < PREV SERMON < < <}}
~
\hyperref[sec:index_chronic]{\small{[返順時目]}}
~
\hyperref[sec:032210]{\small{> > > NEXT SERMON > > >}}
\newline
\newline
約書亞記 24:15
\newline
\begin{longtable}{cl}
\hline
\hline
章節 & 經文 (和合本修訂版)\\
\hline
24:15 & \begin{tabularx}{0.7\textwidth}{X} 若你們認為事奉耶和華不好，今日就可以選擇所要事奉的：是你們列祖在大河那邊所事奉的神明，或是你們所住這地亞摩利人的神明呢？至於我和我家，我們必定事奉耶和華。」 \end{tabularx} \\ \\
[1ex]
\hline
\hline
\end{longtable}
$^{1}$各位平安.
希望我們身體健康 靈性進步.
今天我們來到約書亞臨別贈言的第三.
我家事奉神.
你知道約書亞最出名的一句話是什麼.
至於我和我家.
我們必定事奉耶和華.
你們可以選擇你們事奉的神家.
不過我和我家.
我們必定事奉耶和華.
這個就說了我家事奉神.
因為不夠位所以縮短了字.
是約書亞臨別贈言的第三.
臨別贈言第一要愛耶和華.
昨天說的就是臨別贈言第二.
就是不要事奉別神.
今天就是我家事奉神.
今天很長的十五節聖經.
大家真的要打起精神.
不過我也是讀的.
甚至聖經二十三章第一節.
留意.
約書亞將以色列的眾支派聚集在士堅.
士堅就是他們當日民族的首都.
也是繪幕放在那裡.
很久耶路撒冷還沒有出來.
耶路撒冷出現成為首都.
是大衛王之後的.
以前就是士堅.
照了以色列的長老,族長,審判官.
並官長來.
他們就站在神面前.
不是站在約書亞面前.
約書亞叫他們來.
大家站在神面前.
不是站在我面前立正.
站在神面前立正.
約書亞就對眾民說.
其實不只是叫他們來.
是叫百姓都來.

$^{41}$耶和華以色列的神遇此說.
他將整個說話向他們說.
這是我今天要和大家分享的第一點.
就是以色列的神遇此說.
大家要留意.
神跟你說的話.
大家要留心.
來到一起站在神面前聽.
清清神神.
有些教會今天的崇拜.
我覺得很欣賞.
就是那些教會有一個文化語習慣.
就是一讀到一讀聖經.
人們會立即站起來.
為什麼呢?.
神說話當然站在那裡.
就不像大爺那樣.
躺在那裡讀這麼長時間.
不是吧?.
就是這樣幾點鐘.
不是的.
一讀到神的說話.
每個人站起來.
肅然起敬.
神的說話.
他們就站在神面前.
現在耶和華以色列的神.
遇此說.
大家要聽清楚了.
「古時你們的列祖」.
說「古」.
就是說「古」.
這十五節聖經是說「古」的.
由以前說到現在.
「古時你們的列祖」.
是誰呢?.
亞伯拉罕.
和他的兄弟拿學.
和亞伯拉罕和拿學的父親.
他拉.

$^{81}$是住在大河那邊事奉別神.
解釋一下.
不是事奉耶和華嗎?.
以前他們在大河.
大河是什麼呢?.
幼發拉底河.
幼發拉底河和底亞里斯河.
兩條河在新月沃圖.
就是美索不達米亞.
他們住在吾爾烏爾.
吾爾是這兩條大河的出口地方.
去到波斯灣.
是以前那個地方.
今天就是伊拉克.
那裡後面就是巴比倫.
那裡還沒有出國.
他們是未認識耶和華的.
事奉別神的.
後來認識耶和華.
然後神叫他離開吾爾城.
離開.
我帶你去一個我應許給你的地方.
其實大家要知道.
亞伯拉罕離開吾爾.
其實紅色的字.
就是亞伯拉罕的父親他拉.
他的父親他拉才是最先離開吾爾.
帶亞伯拉罕和拿學一起的.
就是他拉.
拿學早死.
於是就帶了他的兒子羅德一起.
就是亞伯拉罕的兒子.
父親他拉.
他拉這個人物不見經傳.
但是他帶領亞伯拉罕和他們家人.
是最先一家人離開吾爾城的.
後來他拉死了.
亞伯拉罕帶了.
在大河那邊以前他們事奉別神的.
我將你們的祖宗亞伯拉罕從大河那邊帶來.

$^{121}$領他走遍迦南全地.
使他的子孫眾多.
將以撒賜給他.
講故事.
他拉,亞伯拉罕,以撒.
由幽伐拉底河那邊開始出發吾爾城.
一直下來.
故事是這樣.
然後又將以撒生了雅各和以掃.
把雅各和以掃賜給以撒.
不是以掃先的嗎?.
是.
但是只是選擇雅各.
所以是將雅各和以掃排先.
雅各和以掃給以撒.
又將西爾山.
即是以東地.
給以掃為業.
而雅各和他的子孫就下到埃及去.
做了奴隸四百年.
四百年,四百三十年.
好了.
我就猜險摩西亞倫.
就照我在埃及中所行降災.
與埃及多少個災?.
十個災.
然後將你們帶出來.
大家還記不記得?.
這個「鼓」大家記得嗎?.
是呀是呀是呀.
約書說是代表神在向我們所有官長長老百姓講「鼓」.
聽到.
然後我領你們列祖出埃及.
他們就到了紅海.
記得紅海的意思.
又叫蘆葦海.
很有可能不是今天的紅海.
可能是在北比較地中海的南的一個小的湖泊.
也不一定.
那就叫蘆葦海.

$^{161}$可能是誤用了紅海.
埃及人帶領車輛,馬兵追趕你們.
列祖直到紅海.
你們列祖哀求耶和華.
他就使你們和埃及人中間黑暗了.
又使海水淹沒埃及人.
大家記得這些事嗎?.
聽爸爸講過.
聽爺爺講過.
因為他們上一代經歷這些事的時候.
那一代已經在曠野死了.
我上埃及所行的事.
你們親眼都見過.
其實他們沒有見過.
我是聽過而已.
不過不要緊.
當是知道你們在曠野也住了許多年.
是呀是呀這個知道了.
因為他們在曠野是第二代.
之前的事可能他們未必見過.
可能見過.
有人見過.
我領你們到約旦河東摩利人所住之地.
這個點特別紅了.
約旦河東.
這個字多了.
約旦河東阿摩利人那裡.
於是你們和他們爭戰.
我就將他們夾在你的手.
你就得了他們的地為業.
我將他們在你面前.
將他們滅絕了.
所以阿摩利人.
這個東字沒有的.
阿摩利人就沒有了.
第一個剷除了的人.
然後亦都去到有摩押王西貝的兒子巴勒.
巴勒是另一個王.
他就起來要攻擊以色列人.
就打發人召了比爾一個先知巴蘭.

$^{201}$大家聽過了.
就叫巴蘭來救助以色列人.
救助你們.
我不肯聽巴蘭這個先知的禱告.
所以他倒為你們連續仆.
好笑了.
巴勒就請巴蘭先知去救助以色列人.
預備攻打他.
救了他 你給我救了他.
你好靈.
誰知巴蘭一開口.
阿神就令他不斷講祝福的說話.
氣得巴蘭不得了.
巴蘭說對不起.
我真是講不到.
我一開口自然就祝福了.
我都不知為何.
上帝真是厲害.
耶和華.
於是我多給你些錢.
又好啦.
怎知又是祝福.
年年祝福.
這樣我便救你們脫離巴勒的手.
你見到第一個是阿摩利人.
然後是摩訶人.
巴勒有兩個.
阿摩利人 摩訶人.
然後你們過了約旦河.
去到耶利哥.
遭遇到耶利哥人.
阿摩利人.
比利西人.
迦南人.
黑人.
甲加薩人.
希美人.
耶布斯人.
數一數.
1 2 3 4 5 6 7 8.

$^{241}$計算剛才那兩個.
十個什麼什麼人.
都是與你們爭戰.
然後我將他們交你們手裡.
都勝過他們十個民族.
然後.
我打發黃蜂飛在你們前面.
將阿摩利人的兩個王從你們面前趕走.
這個字.
如果我查過字典.
就是「鏈」.
這麼有趣.
不是讀「貝」嗎.
有點像.
就是「鏈」.
什麼叫「鏈」呢.
就是我們廣東話的音.
就是你做什麼.
「鏈車」.
即是飆車很快.
你「鏈車」.
「鏈」即是.
「鏈」他出去.
其實就是出力將他推出.
「鏈出」.
而不是用你們的刀.
不是用你們的弓.
是我用黃蜂趕走阿摩利人.
你記得嗎.
記得記得.
很害怕.
我們都很害怕他執我.
其實現在站在耶和華面前的百姓.
他們是聽過的.
見過見過.
經歷過的.
我賜給你們土地.
地圖.
非你們所收治的.
那些地圖這麼漂亮.

$^{281}$都不是你們平整的.
我賜給你們的誠意.
都不是你們建造的.
你們就住在其中.
無端端就來拿到別人的城.
拿到別人的地.
而你們吃的就不是你們栽種的.
葡萄園和橄欖園的果子.
是你去到已經在這裡.
是我預備好.
前面那十種人.
都拍好了.
我給你的.
記不記得.
記得記得.
記得.
知道知道.
記得.
現在你們要敬畏耶和華.
誠心實意地事奉他.
這是第二點要講.
敬畏耶和華.
將你們列祖在大河那邊.
和埃及所事奉的神.
是除掉.
去事奉耶和華.
其實約書亞講到這裡.
臨別證言第三.
要事奉耶和華.
是什麼呢?.
就是你記不記得以前.
我們講古的時候.
最初連他拉亞伯拉罕.
都是在吾爾事奉當地的神.
是否認識耶和華.
神就帶他們來這裡.
帶你們有亞伯拉罕以撒雅各.
一直到現在.
你們十二支派.
贏了十個王.

$^{321}$十種民族.
無數的王.
抗日的不說了.
都已經這麼多國.
然後.
你記得.
還是拿進來的地方.
不是你經營開的.
那些果子都不是你種的.
我還派王豐幫你打仗.
記不記得.
印象深刻.
所以你現在要敬畏耶和華.
是第二點要講的.
敬畏耶和華.
誠心實意事奉祂.
不要擇假意.
隨便早做便可以.
熬湯隻雞便可以.
去教會便可以.
時間到,走,喝茶.
不是的,是誠心實意.
做做給上帝看.
祂知道的.
因為你不誠心實意祂知道.
是將你們列祖.
大河那邊埃及.
以前幼發拉底河那邊.
和埃及那些神.
你不要再拜.
除掉.
在埃及拜什麼?拜金牛.
那你除掉.
去事奉耶和華.
現在.
你們以事奉耶和華.
不好.
如果.
今天大家來公公道道.
面對面講清楚.

$^{361}$在神在上面.
如果你們以事奉耶和華為不好.
OK.
那你就選擇你要事奉的.
是以前在幼發拉底河.
最早先哈拉阿巴拉罕區.
去事奉那些神呢?.
還是我們這個地方.
阿摩利人那些神呢?.
即是迦南的神.
有很多個的.
還是以前在埃及的金牛呢?.
金牛毒呢?.
不過你們選吧.
我想告訴大家.
我先做一個榜樣.
第三條.
至於我和我家.
我們必定事奉耶和華.
我和我家.
其實我和我家就是我家.
我家就是我.
所以.
不要只是我.
我家必定事奉耶和華.
我將它縮短.
我家事奉神.
這個是今天要和大家講的.
亦都是約書亞臨別贈言的時候.
第三.
我家事奉神.
請大家回想.
整個經歷是這樣的.
神帶領我一直到今天.
到今天約書亞就說.
神在人面前.
大家自己想吧.
要事奉誰.
你可以事奉那些神.
你自己做個決定.

$^{401}$不過你要知道.
上一次臨別贈言第二講就是說.
你事奉別神.
你就速速滅亡了.
神的鼻子就向你發作了.
你要記得神給你很多恩典.
要愛耶和華.
就是臨別贈言第一.
愛耶和華.
臨別贈言第二.
勿事奉其他神.
第三.
我家事奉神.
好像約書亞一樣.
希望大家都記得.
下定決心永不改變.
我就是跟從主事奉祂的了.
你可以選擇.
不過.
至於我和我家.
我們必定事奉耶和華.
而且誠心實意地事奉祂.
不要點點檔.
不要求求其神知.
這樣就可以了.
祂知道我忙.
知道我小孩多.
知道我忙.
知道我要搶第一桶金.
祂明白的.
不是的.
神不明白的.
未必一定說.
這樣就算了.
不是的.
神很嚴厲的.
神有很親和的一面.
好像父親一樣.
但這個父親同時是嚴父.
又是慈父.

$^{441}$都在神身上一起.
所以我們甘心樂意.
我家事奉耶和華.
一起祈禱.
天父上帝我們在你面前.
一起立志.
我們的家.
我自己和我的家.
都要事奉耶和華.
因為你是配我們事奉的神.
你曾經為我們生命裡面.
做過很多事.
在人類歷史裡面.
做過很多事.
整個救贖裡面.
由創世以先.
已經做了很多事.
將來你還是審判的主.
有一天.
你的兒子耶穌基督.
要回來審判這個世界.
將我們帶進永恆裡面.
至於我們今天.
再次重新立志要事奉你.
是我和我家一起事奉你.
認認真真的事奉主.
要討你的喜悅.
我們愛耶和華.
愛你.
垂灘祈禱 奉耶穌基督名求 阿們.
\newpage



\section{路加福音 19:1-10}
\label{sec:032210}
\textbf{撒該:尋回真我}~~~~陳傳華牧師博士
\newline
\newline
連結: \href{http://www.church.com.hk/acms/content.asp?site=cdc\&op=show\&type=product\&code=032210\&layout=sermon}{\texttt{032210}} ~~~~ 語音日期: 2026-05-24
\newline
\newline
\hyperref[sec:032209]{\small{< < < PREV SERMON < < <}}
~
\hyperref[sec:index_chronic]{\small{[返順時目]}}
~
\hyperref[sec:032211]{\small{> > > NEXT SERMON > > >}}
\newline
\newline
路加福音 19:1-10
\newline
\begin{longtable}{cl}
\hline
\hline
章節 & 經文 (和合本修訂版)\\
\hline
19:1 & \begin{tabularx}{0.7\textwidth}{X} 耶穌進了耶利哥，要從那裡經過。 \end{tabularx} \\ \\ \relax
19:2 & \begin{tabularx}{0.7\textwidth}{X} 有一個人名叫撒該，作稅吏長，是個財主。 \end{tabularx} \\ \\ \relax
19:3 & \begin{tabularx}{0.7\textwidth}{X} 他要看看耶穌是怎樣的人，只因人多，他的身材又矮，所以看不見。 \end{tabularx} \\ \\ \relax
19:4 & \begin{tabularx}{0.7\textwidth}{X} 於是他跑到前頭，爬上桑樹，要看耶穌，因為耶穌要從那裡經過。 \end{tabularx} \\ \\ \relax
19:5 & \begin{tabularx}{0.7\textwidth}{X} 耶穌到了那裡，抬頭一看，對他說：「撒該，快下來！今天我必須住在你家裡。」 \end{tabularx} \\ \\ \relax
19:6 & \begin{tabularx}{0.7\textwidth}{X} 他就急忙下來，歡歡喜喜地接待耶穌。 \end{tabularx} \\ \\ \relax
19:7 & \begin{tabularx}{0.7\textwidth}{X} 眾人看見，都私下議論說：「他竟然到罪人家裡去住宿。」 \end{tabularx} \\ \\ \relax
19:8 & \begin{tabularx}{0.7\textwidth}{X} 撒該站著對主說：「主啊，我把所有的一半給窮人；我若勒索了誰，就還他四倍。」 \end{tabularx} \\ \\ \relax
19:9 & \begin{tabularx}{0.7\textwidth}{X} 耶穌對他說：「今天救恩到了這家，因為他也是亞伯拉罕的子孫。 \end{tabularx} \\ \\ \relax
19:10 & \begin{tabularx}{0.7\textwidth}{X} 人子來是要尋找和拯救失喪的人。」 \end{tabularx} \\ \\
[1ex]
\hline
\hline
\end{longtable}
$^{1}$今天講題是「殺該,尋回真我」.
我們怎樣看自己.
以及怎樣認知自己的人生.
其實就決定我們的生命是怎樣走.
我再講一次.
我們的生命是怎樣走.
是因為我們怎樣看自己的人生.
怎樣認知自己的人生.
就會決定我們人生怎樣走.
今天我們講「殺該」這個事件.
正正是講這個東西.
不過開始前當然有些破冰的問題問一下大家.
我不知道大家去旅行的時候.
有沒有試過在酒店留下一些.
忘記拿走的重要東西.
有沒有試過?.
真的留下了.
很多人點頭,有些是零頭.
如果你真的試過,check out.
哎呀,然後走了,忘記拿走.
有沒有人試過這件事?.
舉手先,都有20-30\%.
根據hotel.com在上年做過報告.
就研究訪問了400間酒店.
搜集資料究竟客人常常留在房間.
是包括些什麼,看是不是你留的東西.
他就發現很多人留的是尿道,可以俗稱.
手機充電器,有沒有人試過?.
留下了,真的有.
也有變壓器.
他說有些人留下了髒衣服.
我這樣看其實他不是留下的.
是送給你的,他不要搞了,拜拜.
有化妝品留下過,也有.
也有些是灌洗用品,可能很多人真的留下了.
不過這裡也很普遍.
待會說一些很特別的東西.
我不知道你們有沒有試過留下這麼特別的東西.
他說曾經有人留下什麼東西呢?.
就是一隻勞力士手錶.

$^{41}$嘩,看看旁邊那隻手有沒有手錶.
然後越看越有趣,是一輛豪華汽車的車鑰匙.
和成交的文件.
成交了,忘記拿走.
待會你看看,如果你開車來,你的車還在停車場.
我昨天說的時候,徐徐說.
曾經有人留過車鑰匙在教會.
你看看你的椅子有沒有.
接著就有趣了,這個你一定沒聽過.
我也是剛剛搜集資料聽過.
就是有人留了汽車的車輪.
這件事真是大整蠱.
顯然所思.
接著就更奇怪,你看看你兩隻腳有沒有留下.
有人留下兩隻腳的石膏.
打完石膏留下了,真的很奇怪.
不知道有沒有醫生在這裡,我在想為什麼會這樣.
接著留下的我一定不會.
因為我沒有,就是一堆現金.
嘩,什麼事呢.
接著很慘,留下的東西.
就是寵物很慘.
就是留下了一個寵物蜥蜴.
有沒有人養在家裡.
最後一個,今天想說的是最慘的.
在這個破冰的時段.
很可憐,留下了一隻小雞.
我不知道你有沒有留下一些有趣的東西.
或者貴重的東西.
當我做這個PowerPoint的時候.
Melissa幫我做的時候.
她說我也有留下東西.
她說當她結完婚.
動物月的第一晚走了.
你猜她留下什麼呢.
結婚戒指.
嘩,真是大件事,這件事真是不行.
當然我們留下的東西也可以找到花.
我們可能也買得到.
不過,今天我想說的主角.

$^{81}$她不見了什麼,她留下了什麼呢.
她留下了整個人生.
她發現她找不到自己.
在經文裡面第十九章第十節描述.
她是一個迷失了失狀的人.
一個人為什麼會迷失呢.
大概可能有五種狀態.
第一種狀態是她方向上的迷失.
她可能有學歷的,有工作的.
也是在讀好的書.
但是她想說我這樣活下去為了什麼.
找不到人生的方向.
第二個可能是迷失.
我們今天也會的.
是情緒和內在迷失.
她不是失敗的.
不過她感到疲憊.
她外表很正常.
都會跟你笑.
會跟你吃飯.
但是她裡面是空蕩蕩的.
她一個人安靜的時候.
她沒有喜樂的.
她會麻木了.
第三個可能性的迷失.
她是靈性和價值的迷失.
她不是不知道對錯.
不過總是走不動.
她很努力想突破.
誰知道又回頭.
她慢慢就開始妥協了.
那種失敗.
妥協了我走不動的那種艱難.
第四種是關係群體中的迷失.
她不是沒有人坐在她旁邊.
你看看你旁邊有沒有人.
有的.
不過她可能覺得很孤單.
很孤立.
不被理解.

$^{121}$家裡裡面.
群體裡面.
她不是沒有位置.
不過好像突然間沒有了身份.
第五種.
不是自己.
是被人.
被人見到.
被裁定.
沒得救.
有沒有聽過人說.
你沒得救.
有的.
沒有了.
你是這樣的了.
她突然間這種人覺得.
哇.
每個人都跟她說.
你不行了.
不行的了.
你不神用.
你沒用.
做事你沒用的了.
你完全是放棄了的.
我不知道.
撒該是哪一種.
或者哪很多種.
不過我一會兒介紹她的背景.
可能你就會更加明白.
究竟為甚麼.
耶穌在這段時候.
會跟撒該有這種互動交往.
亦都想想我們人生.
有沒有都跌入在某種迷失中.
藉著今天的經文.
可能我們都找到.
一些我們不見了的東西.
我們先看看背景.
究竟撒該被人掙扎.
他迷失了甚麼呢.

$^{161}$我們看看《路加福音》第十九章一到二節.
我們一起讀.
一二三.
耶穌進了耶利哥.
正經我的時候.
有一個人名叫撒該.
作祟彌彰.
始終都在此.
在這裡描述的地方.
耶穌進了哪裡.
耶利哥.
祂是甚麼經過.
我不知道你知不知道.
耶利哥.
耶穌去完之後.
下一站去哪裡.
下一站.
當然是天后.
下一站是甚麼.
不是天后.
下一站是耶路撒冷.
其實耶穌已經去到.
祂在地上最後一個星期.
你想想.
祂進入耶路撒冷.
就是最後幾天.
祂就會上十甲.
一個人.
如果剩下幾天.
你一定把握時間.
要做你很想讓人知道的事.
耶穌把握那個時間.
揀選了撒該這個事件.
跟我們說一些很重要.
我們不可以沒有了.
亦都不可以在最後階段.
不認識的真理.
究竟是甚麼.
好了.
祂真的介紹了一個人.

$^{201}$這個人的名字叫甚麼.
撒該.
他做甚麼的.
稅利.
還是掌.
他有沒有錢.
有.
好了.
短短幾句.
你開始會知道.
這個背景又是甚麼.
我們把握最後時機.
耶穌教甚麼.
這個稅利其實是一個包稅箱.
甚麼叫包稅箱.
包銷箱.
你聽過了.
即是包銷的.
總之你賣了.
賣完之後給我錢.
其他東西歸你.
包稅箱當日的背景是甚麼.
當日羅馬政府叫它包了.
一個地方的稅款.
5億.
例如亂說.
即是今天的尺度.
5億以上你收到.
全歸你所有.
可以嗎.
可能收8億,9億.
所以你會明白.
開始有些事會發生.
第一.
這些稅利.
即是收稅的人.
是貪心的.
是濫權的.
鍾寶思狼牙.
而這樣的角度你看.

$^{241}$你就知道.
如果錢他越多.
就代表你的錢會越少.
因為他在哪裡拿.
你那裡拿.
他的屋越大.
你的屋就怎樣.
越小.
他吃的東西越好.
你吃的東西又怎樣.
越差.
很明顯的.
就好像當日六七十年代香港.
我知道是警察一樣.
周圍收片.
收多少是沒有限制的.
跟著他上繳之後就全部他.
這種人.
你喜不喜歡.
你會不會跟他做朋友.
你真是好了.
不會的.
但是你會很害怕.
你很害怕得罪他.
你又不會很明顯表達.
但是你知道.
最好有距離.
第二.
他一定是被藐視和邊緣的.
這種人沒有人想要的.
為甚麼.
因為他是一個叛徒.
他幫羅馬政府做事.
是賣國賊.
可以嗎.
他差到一個地步.
連猶太人會堂不可以進入.
即是教會不可以進入.
不會讓他進入.
你會不會不讓人進入教會.

$^{281}$不給的.
為甚麼.
當日甚麼人不可以進入會堂.
例如外邦人.
土匪.
妓女.
可以嗎.
看不起的人.
他不可以進入.
他是自己人.
不過不可以進入.
所以他另一個層面就是.
不單止別人討厭你.
跟你吃飯交易打招呼都不想.
因為我不跟你同一類.
第三.
他迷失根源.
不是他沒有物質世界.
他會住豪宅.
他有很多工人.
有很多美食.
不過.
他是一個有錢的孤兒.
沒有關係.
或者我們以前的年代很喜歡說.
窮得只剩下甚麼.
錢.
物質世界給不到我們平安喜樂.
我們需要的是關係.
所以他被他的身份和金錢囚禁了.
當然我們今天不是稅吏長.
他不單是稅吏長.
他是財主.
我們都不是.
不過我們某程度都被我們有些身份.
我們對物質世界的追求囚禁了.
我們有些事是做不到的.
好了.
今天.
我們當然好像剛才說.

$^{321}$我們都知道我們不是國際.
但是我們可能都知道這個行業不被認同.
或者我們的身份是被否定的.
我們都有這種狀況在今天.
如果再這樣看我們剛才那五個問題.
可能你都知道薩哥他在面對甚麼.
又或者我們可能今天都面對某程度的渴求.
他的渴求當然是突破他人生的身份限制.
他的渴求當然是可以超越到.
他那種和人與人之間和神的隔膜.
今天說十節聖經很少.
十節裡面你會找到很多重要的答案.
也是耶穌把握這個機會想和我們說的話.
我們一起讀這十句十節的經文.
這些字細一點有原因的放在這裡.
我們一起讀.
一二三 耶穌盡了耶利哥.
當天我接受了耶利哥的教訓.
他就急忙下台.
悶悶氣氣地接待耶穌.
從人看見到司祭議論說.
他頂到罪人家裡去住宿.
實在站不住對著說.
主啊 我不想有這一般的恐懼.
我若有犯了信.
就還他四倍.
耶穌說 今天有恩到料者家.
因為他接受了哈拉三世之說.
人死亡 唯一的尋求.
只有受到罪人.
這十節經文裡面說了很多東西.
我先帶大家看一下.
你看到有多少個角色.
第一個角色.
耶穌第二個角色.
撒戈是第三個角色.
夢多.
中人.
中人是吧.
他提到的角色.

$^{361}$所以夢多在附近.
不過沒有明顯提到.
好了 三個角色.
三個角色都做同樣.
同一樣的動作.
是什麼動作.
這個字經常出現.
重複又重複的.
看見.
你跟隔壁的說你很厲害.
找到了.
三個角色.
都是不斷看 看 看.
看到什麼呢.
我今天很想說.
那件事不同的角色看見.
結果可以不同.
我們看看看到什麼.
結果又是怎樣.
他們看見了什麼呢.
之後怎樣影響他們的人生態度呢.
而這件事.
從我一開始說.
我怎樣尋回真我.
我們怎樣看人生.
就會帶來我們怎樣活人生.
有什麼關係呢.
怎樣可以活出我們豐盛人生.
三個要訣.
今天立刻進入第一個.
第一個就正是撒戈這個人的角色.
他的漢代表什麼.
他的漢是代表對耶穌的渴求.
我們看看這兩節你就更加明白.
第三第四節我們一起看.
一二三.
他要看漢 耶穌是怎樣的人.
這個有錢人.
這個人很不喜歡的人.
他很想見到耶穌.

$^{401}$但是因為他怎樣見不到.
矮.
矮你很見到的.
今天在台前也有一個.
看不到的.
除了Jimmy那麼高.
梁洪釋那麼高.
於是怎樣.
爬上什麼.
桑樹.
你想想他這樣做.
是他不怕丟臉.
他不怕沒有體面.
他一個人不喜歡的人.
人拒絕的人.
但是因為一件事.
他竟然做第一步.
就是因為他想見到耶穌.
他不怕被人笑.
哪有人爬樹見人.
還是你這樣.
很奇怪的.
第一件事我們要學習.
撒戈寧願失去面子.
都不願意錯過見耶穌的機會.
我記得我初信的時候.
我是男校.
我最喜歡去一個會.
就是不去教會.
我就去一個會叫做馬會.
我當時跟我的同學說.
要去教會是很困難的.
很瘀青的.
我們年代流行說.
所以以前跑馬是星期六的.
如果你知道.
有沒有倒過馬.
前陣子跑馬是星期幾的.
怎麼都說不知道.
是星期日.

$^{441}$現在是星期日.
以前是星期六的.
我們青年團是星期六的.
有人帶我去中午.
我就怎樣呢.
去了馬會跟一班舊朋友.
當時是同學.
跟著再去教會.
教會還沒完的時候.
我再去馬會可以嗎.
會合他們.
他們住在一起.
我又不是很想告訴別人我的身份.
但是慢慢我發現.
回到教會太急人了.
就是尋找出.
但是那的確是一個心理關口.
就是我們原本.
我們的身份要有突破.
不容易的.
今天很多人都仍然在公司,在家.
不想表明他追求主的身份.
為什麼呢.
為什麼會這樣呢.
我也明白的.
實在不容易.
我記得沈祖堯醫生.
當天說SARS的時候.
他提到.
當時醫生又開始感染.
護士又感染.
病人又會死的時候.
他就跟神說.
我的科技已經不能做任何事了.
但是他跟神祈禱.
就是主不要讓我的無知.
他已經是全球權威.
他是很厲害的.
他的分數是很高的.
他說不要讓我們的無知.

$^{481}$令到人失去生命.
於是他就.
召喚了一個祈禱會在醫院.
接著醫生來.
他才發現.
原來你是基督徒.
原來人們都不肯表明身份.
為什麼呢.
我們一定要問一個問題.
我們要表明身份.
我們要告訴別人.
我們真是渴求追求主.
這種撒該他都願意的.
是很重要的.
撒該為什麼會呢.
我大概可以猜到.
他為什麼這麼想見耶穌.
第十九章.
耶穌已經去到最後上十家了.
所以我們知道耶穌.
祂的傳聞是人們都聽到的.
五餅如魚餵飽多少人.
五千.
男人而已.
加上女人一萬以上.
祂可以醫治人.
死人復活.
瘸腿可以行走.
污鬼被趕出.
盲可以見到.
是人都知道的.
不過.
我相信還有一樣東西可以觸及到祂.
就是祂知道耶穌十二個門徒.
其中一個叫做.
馬太.
馬太是做什麼的.
瑞麗來的.
行家來的.
這個行家竟然.

$^{521}$不收稅.
去信耶穌.
是不是瘋了.
錢都不要.
有沒有不要錢的.
舉手.
都要嘛.
錢都不要.
瘋了嗎.
當然有些東西的.
但是他見到這個馬太.
又跟了耶穌很好.
錢都可以放下.
他一定有很多東西在他的腦海.
所以他很想追求主.
很真實.
我記得在我其中有一次的.
宣教短宣.
我就有一個星期去了一個地方.
整個星期.
就是教聖經.
由早上九點教到六點.
九點到六點.
到了星期二.
星期二左右.
我就在廚房.
就跟他的廚師聊天.
我說為什麼你在這裡服侍.
那裡不是講廣東話的.
不過不告訴你說什麼文.
免得你知道在哪裡.
他就說.
其實我以前內戰的時候.
是信耶穌.
內戰.
我說什麼事.
他說內戰的時候.
一路走走走走.
就見到教會進去了.
然後見到一本聖經.

$^{561}$他就在那裡問那位牧者.
我可不可以拿一本聖經來看.
他拿了來看.
看了看了.
他說看到第四十三節.
他就順道注意了.
我說.
哪一卷聖經.
哪一章.
還有哪一章的第四十三節.
他說你聽錯了.
不只是第四十三節.
是我看了整本聖經的第四十三節.
他在那裡遇見主.
他解釋不了那種渴慕.
並且他在裡面.
在那個很艱難的環境.
但是他知道上主的同在.
還有那種平安.
他沒有辦法推翻.
他的渴求讓他知道上主的實在.
我當刻見到他.
簡直高仁在廚房.
我想問他一個問題.
其實教神學是不是應該由你來教呢?.
原來高仁到處都是.
重點是我們渴慕主.
你摸一摸你旁邊那個是不是高仁.
高仁不是我們讀不讀得到神學.
是我們對主多麼渴慕.
對上帝在我們生命的塑造有多麼著緊.
撒該他本來是一個.
他自己都未必喜歡的人.
不過他囚禁在那種物質和地位中.
他知道.
他聽到二手的奇蹟不滿足.
耶穌是迷失者的出路.
是有能力改變人生命.
讓人盼望.
在那一刻他又知道.

$^{601}$聽到人的二手經歷不滿足.
我們要成長經歷必須要是一手.
並且不是很多年前在上上都體驗.
有人問我怎樣可以經歷主.
當日未有聖經.
今日我們透過聖經裡面.
不是那些文字.
在裡面我們可以體驗背後上主.
並且與他相遇.
當然今日耶穌都會直接出現.
我不知道你有多渴求主.
有什麼阻止你可以與上帝相遇.
可能是我們忙碌的工作.
可能是我們懶惰.
可能是iPhone.
可能是社交媒體.
或者愛理不理.
我一定有機會遇見你.
當日撒該知道他只有一個機會.
耶穌就在那裡經過只有一次.
他不見耶穌就沒有了.
今日我們仍然都需要這種渴慕.
每天早上我們都可以遇見主.
在靈修裡面.
在安靜裡面.
我們在崇拜.
敬拜裡面我們是遇見主的.
這個很重要.
昨晚我都是.
在這裡看著看著.
看到最近很觸動我的是UAP.
即是UFO.
我看的時候我馬上想一件事.
我們都要作回應.
在神學上.
聖經上是什麼來的呢.
我看到都快十二時要睡覺.
今早要講多話.
因為我要挖出時間遇見主.
撒該的看見是代表.

$^{641}$他的渴求定睛在耶穌身上.
各位弟兄姊妹朋友.
你渴不渴求遇見耶穌.
你渴不渴求經歷耶穌.
我們不是一個形式.
我們是一個親身緊密的接觸.
第一個.
撒該的看見.
渴求與主相遇.
第二.
是耶穌的看見.
迎接耶穌接納人的一刊.
你不要覺得很簡單.
我解釋一下為什麼.
我們先讀第五第六節.
預備起.
耶穌到了那裡抬頭一看.
到了那裡撒該拜看來.
敢於看見主在耶穌面前.
撒該望向耶穌.
恭祝耶穌接待耶穌.
很奇怪的這個場面.
很多人的.
我們知道耶穌出去.
幾千,一萬人走一走.
走著走著.
耶穌你只是想望一下祂.
突然間祂在你下面停下來.
望著你.
你都很害怕.
你不是和我算賬嗎?.
我會堂都不能進去.
我來遠望一下你.
你不是都想趕我走嗎?.
是不是我臭名遠播到一個地步.
你都不放過我.
我已經在樹上看著你.
你都要.
你通常走都這樣走.
你不是突然望上面.

$^{681}$這個動作本身已經很奇怪.
不僅如此.
祂還大聲說.
撒該.
搞錯.
你為什麼叫我的名字.
讓所有人都知道我在這裡.
不是這麼好的.
這件事真的很奇怪.
我也常常有時候.
在街上有些人突然叫我.
我都走一走.
我記得十多年前有一次.
這次是傾斜的.
香港是比較好.
有一次和我家裡.
兒女,父母,外父去歐洲玩.
去到意大利一個小鎮.
我們在玩耍,聊天.
突然間.
有一個人在我後面說.
Pastor John.
我走一走.
我望一望.
接著說Pastor John.
我說嗨.
接著.
其實我不記得他是誰.
我就說.
說了一句廢話.
來玩嗎?.
接著他說完之後他就繼續.
和他的朋友走了.
我女兒一看就知道.
Daddy.
看你就知道.
不記得他是誰了.
這也不夠奇怪.
上年我們和敬拜隊去越南玩.
上機那一刻.

$^{721}$還在越南上機那一刻.
即是走了.
突然有人在我旁邊說.
Pastor John.
我走一走.
是我.
我說哦.
是你.
十幾年前我碰過你的.
在意大利.
真是他.
我老婆在旁邊.
接著我說.
哇.
我很期待他告訴我他的名字.
哇.
他又不說.
接著我說.
哇.
又說一句廢話.
來玩嗎?.
我不知道他是不是在聽.
我很開心.
大家和我打招呼.
不過有時候我真的不記得.
你去提一下你的名字.
和你在哪個教會.
我真的嚇了一跳.
為什麼會這樣.
很害怕.
很奇怪.
不僅如此.
他叫你的名字.
真是奇怪.
所以我相信.
那一刻.
薩戈是很驚嚇的.
不是我們今天那麼期待的那種.
因為他本身所做.
不過更奇怪的是.

$^{761}$那一刻.
耶穌叫他.
不是帶來審判.
那一刻.
他說快下來.
今天我必住在你家裡.
他就歡歡喜喜接待耶穌.
這句更加奇怪.
剛才我說過.
他連會堂都不進去.
人都不會和他聊天.
竟然這麼多人圍著耶穌.
在他下面停下來.
薩戈下來.
我今天在你家裡.
那種震撼.
是奇怪的.
我不知道大家喜歡不喜歡看演唱會.
我都蠻喜歡的.
最近又臨一年.
又有MC.
Michael Jackson.
他不會復活的.
你幻想一下.
如果你出去看演唱會.
突然間萬多人.
或者有些五萬人.
他停了音樂.
拿個咪指著你.
遮你.
大家看到你.
你會嚇到.
那一刻真的很震撼.
我想說真的很震撼.
但是他震撼得來.
他很想告訴你.
耶穌知道誰見到他.
誰在做夢.
耶穌知道.
誰來到他面前.

$^{801}$只是擠擠湧湧.
而誰真的有心見到他.
你想想.
一萬人.
為什麼只有他.
今天我們也是.
你聽到.
我相信你真的很想見到耶穌.
應該是吧.
你問你旁邊那個.
應該是吧.
我的意思是.
神知道誰真的想見到耶穌.
誰在滾手機.
誰的樣子真的在聽.
但是他心裡想.
什麼時候結束.
很痛啊.
耶穌知道的.
耶穌會去到你那裡.
所以為什麼耶穌要放這件事.
在臨走前.
很想告訴你.
這個很重要.
虛心的人有福.
因為天國是他.
虛心在他面前.
飢渴慕爾人有福.
因為他們必得飽足.
我聽過很多次.
例如有些過來人.
吸毒了二十年三十年.
什麼戒毒都失敗.
在他最失落的時候.
耶穌不是罵他.
他跟他說.
我要愛你到底.
那一刻他有力量.
起來.
又有一個過來人跟我分享.

$^{841}$又是坐牢搞了很久.
他在街上很慘的時候.
耶穌跟他說你是巴拿巴.
巴拿巴是誰.
是幫助別人的.
在聖經裡面.
他說那天我已經爛了.
躺在那裡都聽到.
結果他起來.
他真的幫助別人戒毒.
很奇妙的.
不僅如此.
你說有些很壞的人.
都會這樣經歷.
但是有些.
不是在那方面需要.
都遇見主的.
他都被上主幫助.
引導他內心.
並且經歷他.
經歷一個徹底翻新的生命.
前陣子我跟我其中一個博士學生.
在教育院教的時候聊天.
他在寧德時代工作.
寧德時代如果大家知道.
他的鋰電池.
是全球50\%是他的.
我就問你們的科技是哪裡來的.
當然他們有一班科學家.
但是他說最主要的科技是同一個人來.
這個名字我叫John Goodenough.
我最喜歡John is Goodenough.
不是我加上去.
他真的叫John Goodenough.
我當然再仔細找他的過去.
他的科技.
他在2019年拿了美國.
拿了諾貝爾獎.
他也是一個很創新的鋰電池.
一個科學家.

$^{881}$他怎樣徹底改變人類.
首先他徹底被神改變.
他在牛津大學工程系鋰電池.
在裡面開發的時候拿了諾貝爾獎.
有很多不同的報道.
明報報道很多報道.
有一個報道很有趣.
他怎樣描述呢.
他鋰電池之父.
是手機筆記電腦等等.
他是一名基督徒.
他加入了更多諾貝爾獎的基督徒行列之一.
紐約芝加哥一個非常有名的大學這樣記載他.
他說他的辦公室.
我不知道你辦公室會放些什麼.
我一定放一些我覺得最重要的東西.
或者有一張桌子.
他說看到有最後晚餐的大掛煎.
有很多宗教的文物.
還有他和他太太多年不同旅行收集紀念品.
你知道信仰對他很重要.
很多人覺得科技和信仰是抵觸的錯.
科技是上帝創造的.
怎麼會有抵觸呢.
當我們在主裡面.
我們可以更加有這種創造力.
他就這樣分享.
他小時候原來爸爸也是教授.
而且是大班教會.
但是他當時還沒有和神連繫.
他說有一天晚上.
他經歷了上主的真實.
他怎樣經歷呢.
當他一直體驗不到神的真實.
他突然間明白了神的愛.
明白的時候.
他說這個就是聖經一直在說.
我第一次經歷了愛.
於是他興奮到一個地步.
把床跳起來.

$^{921}$十幾歲而已.
接著就要不停分享神就是愛.
我想說什麼呢.
原來我們的原動力.
就是始於神就是愛.
無論薩哥的改變.
他有能力改.
不是我靠自己的能力.
我想可以超越我的微小.
我的黑暗面.
不可以的.
如果可以的話.
世界不會搞得越來越差.
耶穌也不用來.
或者我們不是很有問題的.
很大件事.
但是我們很想找到人生的路途路徑.
我們要靠著主在我們裡面開啟.
接著他還沒說完.
接著他繼續說.
這種感覺.
神是愛感的伴隨他.
八十幾年.
十幾歲.
一百歲死.
於是乎給他力量開始探索這種超級電池.
並且他演繹愛人如己的愛倫社.
我們愛倫社的探訪.
好的.
我們愛倫社的家人.
愛倫社好的.
他領受了愛倫社.
發明一些環保的東西.
保護世界.
我們知道了.
現在汽油車很污染.
電池車是零污染.
他就這樣領受到.
接著補充就是.
科學家所做的.

$^{961}$是了解大自然.
以為大自然服務.
接著他繼續提到.
聖靈監察一顆開放的心.
有能力將我從世界的局限可以釋放出來.
今天是聖靈降臨節.
聖靈降臨是顛覆生命的重要起始點.
婷小妹.
我不知道你去到第二點.
你聽到些甚麼.
你有甚麼尋求的渴求.
那個不是別人的故事.
那個是要你和神建立的那種關係.
當我們可能已經建立心的關係.
又或者我們好像很認識.
但都是很抽離的話.
我鼓勵你.
你一定要和神建立一個密切的關係.
又或者你很想帶別人.
經歷主不是就這樣講耶穌.
是與他和神建立關係.
怎樣建立.
很簡單.
例如.
我很想和你分享一下.
我怎樣經歷主的見證.
我和你分享幾分鐘.
如果他想聽.
你和他祈禱最重要.
叫開放心.
讓耶穌直接和他說話.
用經文用神的同在非常之重要.
今天我們經常留在一個形式化的空殼裡.
那裡改變不了生命.
撒該告訴我們.
他因為這樣.
他一個翻天覆地的改變.
John is good enough.
John good enough.
一樣.

$^{1001}$這個原動力改變了.
不單止是他.
改變了今天.
最後一個漢奸.
就是眾人.
眾人的漢奸代表失敗的.
代表失望.
代表絕大部份人的回應.
絕大部份人的回應.
見到耶穌.
都是拒絕的.
我們看看第七第八節.
一起讀一二三.
眾人看見都私下議論說.
他見到罪人打亂了祖宗.
誓我斬上罪處之人.
我把所有的一半給同人.
我一方沾了水就能克死你們.
剛才我說到很多人在耶穌那裡都是擠擠湧湧.
這班人見到不是見到耶穌對人接納.
讓人改變都見到.
反應就是投訴.
搞錯耶穌.
都說你無了道.
你知不知道他是誰.
他家住.
我們都不去.
你搞什麼.
其實我這樣說都已經很有禮貌.
因為這班人最後釘死了耶穌.
最大問題耶穌就是.
有這種力量改變人心.
但那些不想被改變的.
就會不斷指控他.
指控到一個地步.
最後他死在十字架.
這一幕之後就禮了.
但是撒該的改變是奇怪的.
他的改變就是.
他把他一半的財產分給窮人.

$^{1041}$騙了誰給他四倍.
即是想說什麼呢.
即是說如果我們人生想改變的.
第一我們必須要見到自己的限制和不足.
我們看不到是一定不會信耶穌的.
耶穌是奢侈品來的.
多餘的.
如果我們見到我們的不足.
見到我們原來還有一些潛能是發揮不到的.
需要我們可以去到突破的話.
我們一定要第二個要與耶穌相遇.
與耶穌相遇的時刻.
我們經歷到第三就是神的愛.
而不是審判.
我們才有力量改變第四.
和我們以往有一個很重新的顛覆改變.
我們才能活出我們真正的人生.
當然如果我們種種聽人之後.
我們沒有改變的審判就來了.
即是說如果我經歷這四個步驟.
我跟著跟大家說我本來有五個女朋友.
現在見到耶穌這麼好.
我減一個我繼續有四個.
我又結了婚的.
這個是假的.
我正在訛詐的.
我正在說謊的.
我正在做一些事情.
神不喜悅.
我接受主人.
我繼續做同一樣東西.
是假的.
這班眾人就是這樣.
他們說我跟著主人.
不過耶穌的教導我不要.
我們要分清楚.
我們要知道不同的步驟.
所以一開始我說.
如果我們很想.
我們的人生.

$^{1081}$我們的迷失中有力前行.
不是用我們自己的力量.
做不到的.
佛教也上了三法印四聖體.
五溫六根八淨度十二因.
超越再去獵盤.
做不到的.
靠自己的力量做不到.
當我們謙卑在主前主下求你.
進到我內心開放.
並且讓我體驗你.
我肯定你是一個不一樣的人生.
甚至我這樣說.
很多人以為自己信了主.
都未信主.
因為我見不到他人生後面.
有這種顛覆性改變.
還是依然固惡.
早前我記得有一次在小組.
分享一個題目叫做.
「把序連根拔」.
當時每個組員在那裡笑笑.
聊聊玩玩在那裡說.
接著有一個組員開始說.
我說說我懶.
所以我想勤力一點.
這樣都可以的.
我都要做多點運動.
要健康都可以的.
說到第三個組員突然說.
我想處理我的罪.
就是我婚前性行為.
哇 大家呆了.
本來在笑 突然靜默了.
接著那個組員說.
我不跟你玩那麼多.
說到這個題目.
就認認真真回應.
大家還呆了.
接著第四個說.

$^{1121}$我也說一些真話.
我還在吃大麻.
大家還在靜.
接著第五個組員說.
哇 你設定的標準那麼高.
我不知道分享些什麼好.
當然不是每個人都會這樣.
但我們都有些東西想回應的.
接著有一個說.
我做銷售的.
我都知道有時候我買的東西.
其實不應該買.
老闆要我去貨.
我還沒跟客人說 這個好啊.
貴一點更好啊.
他用多幾年的.
功能又多.
其實他知道假的 賣不出.
那些分享掙扎.
有些人說.
其實我整天在公司工作.
都在樓下打邊爐.
圍爐.
我們要回應就要給些真話出來.
我覺得很好.
薩蓋對我來說是一個偉人.
他大可以躲起來.
他大可以.
尼哥底冇夜媽媽找耶穌.
雖然耶穌未必見他.
他大可以自己覺得.
依然固我在那個層面裡面.
享受自己擁有的東西.
自己暗暗地跟天父祈禱.
不是的.
信仰是清清楚楚在主前.
這個認真回應我們一個很重要的內容.
就是我們的生命.
今天究竟是不是跟神有關係.
跟著那個祖賢.

$^{1161}$有最後一個祖賢分享.
我家爸爸剛剛電騙.
被人騙了幾十萬.
那些祖賢就說 嘩 那麼蠢 爸爸.
現在個個都說不要被人電騙.
還被人騙.
跟著那個祖賢說.
不是我爸爸差.
是我做得不好.
我沒有關心我爸爸.
直至到被人騙了才知道.
什麼是尋回真我.
就是尋回主開始.
隨之而來我們就變化.
隨之而來我們就回應主.
原來我們裡面本身都有神的形象.
我們裡面都有上帝給我們美好的生命.
只是我們很多的物質.
很多的外在.
很多的自我.
很多的驕傲.
遮蓋了.
我們要改變.
我們的生命要有著這種意義.
必須要從與主連繫開始.
這樣最終讓我們跟神的關係.
可以真正的富足起來.
曼科多倫神學家說.
我們要實現一個人身為人的潛能.
其實我們就正正要反映神的形象.
今天最困難就是.
我們的形象.
我們的自我.
比神的形象.
來得更大更深.
讓我們見不到上帝.
在我們生命美好的計劃.
總結.
我們稱祂為女.
最後這兩節經文是這樣說的.

$^{1201}$我們一齊讀.
一二三.
耶穌說.
「今天救一道了者家」.
「因為他人是罷了他人之子」.
「人子來為要尋找拯救失喪的人」.
尋回真我.
被神接納的身份.
悔改的行動.
有使命的人生.
是必然接踵而來出現的.
當我們關心我們的生命.
我們的靈命.
我們的情緒.
肉體.
使命.
都會隨之而改變.
尋回真我是爭取人生的開始.
有神有使命有富足有承擔.
有平安有情有關係.
願意活出的秘訣.
渴求迎接改變.
這才是精彩的.
最後有一個小的比喻.
我就完成今天的分享.
從前有一個國王有四個太太.
國王快要死了.
病重.
臨死的時候.
他就跟他最寵愛的第四太太說.
「你一直都是我最好的」.
「我很珍惜你」.
「我愛你」.
「我給你最多」.
「你願不願意陪我一起死」.
「陪我一起進入來生」.
第四太太說.
「對不起」.
「我恐怕不能」.
說完之後他走了.

$^{1241}$他很傷心.
於是他跟第三太太說.
他也很愛她.
經常帶著她到處列國炫耀.
「你願不願意陪我死」.
三太太說.
「不行」.
「人生太好了」.
「你死之後」.
「我告訴你我會改嫁」.
三太太還沒死.
馬上哭了.
心如刀.
跟著他跟第二太太說.
「二太太你一直都很支持我的」.
「我想問你願不願意陪我一起死」.
「進入下一世」.
二太太說.
「對不起」.
「這次我真的幫不了你」.
「我已經經常幫你」.
「你甚麼需要我都幫你」.
「不過你放心」.
「我會幫你好好安排葬禮」.
「好好看看」.
「然後葬禮送你一程」.
國王很絕望.
在那一刻.
大太太大老婆就說.
他沒有問她.
大老婆突然說.
「我願意跟你一起去」.
「無論你去哪裡」.
「我都跟你一起去」.
原來他一直最忽視的一位太太.
也是最少照顧的一位太太.
原來是最深愛她.
他很懊悔.
原來他一直都忘記.
好好關心大太太.

$^{1281}$這個故事其實是說人生的四樣東西.
四太是我們的身體.
我們無論怎樣喜歡她打扮她.
她都沒有辦法陪我們入永恆.
物質世界會完結.
三太是我們的財富和地位.
你怎樣炫耀她.
有什麼車有什麼好.
兒女在哪裡讀書都好.
最後她都會改嫁的.
貴客一位.
我很抱歉.
二太是我們最好的家人.
最好的朋友.
最支持我們.
不過極其量.
在我們死後都為我們準備安息禮拜.
大太是誰.
是我們的靈魂.
我們常常都忘記了她.
無所謂啦.
我有空理一下你.
是我們和上帝的關係.
是我們的靈性.
原來這個靈性.
才是影響我們人生.
最深最廣的地方.
撒該告訴我們.
他找到他的靈性.
而這件事.
就好像我一開始說.
我們怎樣看人生.
就影響我們怎樣看人生.
剛才我說.
這個經文有三個角色.
其實有四個.
最後一個角色.
就是你.
你望一望你旁邊.
你說你是第四個角色.

$^{1321}$你怎樣回應今天的經文.
讓上主親自感動你.
我們一起祈禱.
我不知道你學不學求主.
你學不學求與主相遇.
可能你來過主教會.
或者你聽過很多次耶穌.
或者甚至你已經是基督徒.
但是你說主.
我真的很想這一種.
很入心的連繫.
叫你的名字.
也真的會.
你也願意耶穌住在你家裡.
如果你其中一個.
很想為你禱告.
你很想找到你人生的使命.
有承擔有力量.
如果你其中一個.
將你的手放在你的心.
我為你禱告.
天王我們多謝你.
讓你這樣向我們說話.
讓我們真正知道我們是誰.
不是別人說.
不是我們自己說.
是從你而開始.
讓我們經歷一個徹底改變的人生.
在一個聖靈來的時刻.
讓我們真正一個生命的革命.
生命的巔峰.
讓我們活出神的形象.
你放在我們哪裡.
我們就是哪裡.
你要我們怎樣去行.
我們怎樣去行.
一切都是源自我們對你的渴求.
對你的開放.
願你幫助.
多謝主.

$^{1361}$預立我們的禱告.
奉耶穌名靈教.
\newpage



\section{以弗所書 4:1-16}
\label{sec:032211}
\textbf{沒有輪心的車輪走不動}~~~~禤凱榮牧師
\newline
\newline
連結: \href{http://www.church.com.hk/acms/content.asp?site=cdc\&op=show\&type=product\&code=032211\&layout=sermon}{\texttt{032211}} ~~~~ 語音日期: 2026-05-31
\newline
\newline
\hyperref[sec:032210]{\small{< < < PREV SERMON < < <}}
~
\hyperref[sec:index_chronic]{\small{[返順時目]}}
~
\hyperref[sec:032212]{\small{> > > NEXT SERMON > > >}}
\newline
\newline
以弗所書 4:1-16
\newline
\begin{longtable}{cl}
\hline
\hline
章節 & 經文 (和合本修訂版)\\
\hline
4:1 & \begin{tabularx}{0.7\textwidth}{X} 我為主作囚徒的勸你們，既然蒙召，行事為人就要與你們所蒙的呼召相稱。 \end{tabularx} \\ \\ \relax
4:2 & \begin{tabularx}{0.7\textwidth}{X} 凡事要謙虛、溫柔、忍耐，用愛心互相寬容， \end{tabularx} \\ \\ \relax
4:3 & \begin{tabularx}{0.7\textwidth}{X} 以和平彼此聯繫，竭力保持聖靈所賜的合一。 \end{tabularx} \\ \\ \relax
4:4 & \begin{tabularx}{0.7\textwidth}{X} 身體只有一個，聖靈只有一位，正如你們蒙召，是為同有一個指望而蒙召， \end{tabularx} \\ \\ \relax
4:5 & \begin{tabularx}{0.7\textwidth}{X} 一主，一信，一洗， \end{tabularx} \\ \\ \relax
4:6 & \begin{tabularx}{0.7\textwidth}{X} 一神－就是萬人之父，超越萬有之上，貫通萬有，在萬有之中。 \end{tabularx} \\ \\ \relax
4:7 & \begin{tabularx}{0.7\textwidth}{X} 我們每個人蒙恩都是照基督所量給每個人的恩賜。 \end{tabularx} \\ \\ \relax
4:8 & \begin{tabularx}{0.7\textwidth}{X} 所以有話說：「他升上高天的時候，擄掠了俘虜，將各樣的恩賜賞給人。」 \end{tabularx} \\ \\ \relax
4:9 & \begin{tabularx}{0.7\textwidth}{X} 既說「他升上」，豈不是指他曾降到地底下嗎？ \end{tabularx} \\ \\ \relax
4:10 & \begin{tabularx}{0.7\textwidth}{X} 那降下的，就是高升遠超越諸天之上的，為要充滿萬有。 \end{tabularx} \\ \\ \relax
4:11 & \begin{tabularx}{0.7\textwidth}{X} 他所賜的有使徒，有先知，有傳福音的，有牧者和教師， \end{tabularx} \\ \\ \relax
4:12 & \begin{tabularx}{0.7\textwidth}{X} 為要裝備聖徒，做事奉的工作，建立基督的身體， \end{tabularx} \\ \\ \relax
4:13 & \begin{tabularx}{0.7\textwidth}{X} 直等到我們眾人在信仰上同歸於一，認識神的兒子，得以長大成人，達到基督完全長成的身量。 \end{tabularx} \\ \\ \relax
4:14 & \begin{tabularx}{0.7\textwidth}{X} 這樣，我們不再作小孩子，中了人的詭計和欺騙的法術，被一切邪說之風搖動，飄來飄去。 \end{tabularx} \\ \\ \relax
4:15 & \begin{tabularx}{0.7\textwidth}{X} 我們反而要用愛心說誠實話，各方面向著基督長進，連於元首基督， \end{tabularx} \\ \\ \relax
4:16 & \begin{tabularx}{0.7\textwidth}{X} 靠著他全身都連接得緊湊，百節各按各職，照著各體的功用彼此相助，使身體漸漸增長，在愛中建立自己。 \end{tabularx} \\ \\
[1ex]
\hline
\hline
\end{longtable}
$^{1}$很多年前我在香港認識了一個朋友.
他是開眼鏡店的.
他就要求他的員工全部都要穿一件白色的醫生袍.
我跟他開個幻想.
嘩 你這間店真的不錯.
連醫生都來你的店做兼職.
他就說 你不知道的了.
我的員工如果穿了一件醫生袍的話.
客人對他的信心就會大增.
生意額也會高一些.
這樣想也挺有道理的.
人總是很有趣的.
人總是覺得穿了某些制服之後.
那些制服無論是醫生也好.
或者是警察,消防員等等.
你穿了那套制服之後.
我們對他的信心就會大一些.
我們就會覺得他有某些行為表現.
都會合乎他的制服.
我想問大家.
基督徒有沒有制服的呢?.
我們有沒有的?.
不知道大家有沒有覺得.
人其實對基督徒都有某些行為上的期待.
在今天的經文裡面.
保羅就告訴我們.
我們作為基督徒的行為表現.
需要與我們蒙召的恩相稱.
大家覺不覺得做基督徒.
在我們的言談舉止.
我們的一言一行.
都有某些別人對我們的期望.
或者我們都應該有一些準則的.
大家有沒有試過與某些你所認識的人說.
我是基督徒的時候.
他們是怎樣反應的?.
你不用說我也知道你是基督徒.
又或者會不會是.
你是基督徒來的?.
大家是屬於哪一種?.

$^{41}$今天我們所讀的經文.
不斷告訴我們在教會裡面的合一.
其實在聖經裡面不斷的說.
不斷的強調.
究竟甚麼是合一呢?.
合一通常來說.
應該是指我們一班人.
為了某些共同的目標.
一同的努力.
這個就叫做合一.
在很多的機構裡面.
一個共同的目標是很清晰的.
很明顯的.
舉個例子.
商業的機構.
他們的共同目標是甚麼?.
賺錢.
軍隊呢?.
軍隊的共同目標是甚麼?.
打勝仗.
現在世界盃快要來臨.
一個足球隊的共同目標是甚麼?.
贏球.
這些共同目標是相當的明顯和清晰的.
而所有你會發覺這些共同目標.
其實某個程度上.
都和機構裡面的成員的利益掛鉤.
對於商業的機構.
公司賺錢.
員工很有大機會加薪拿花紅.
在軍隊裡面打贏仗的時候.
不單止自己的性命可以保住.
國家都可以得到平安與安全.
足球隊呢?.
足球隊贏了球之後.
發達了.
名利兼收.
那些甚麼東西都找你來做代言.
那教會呢?.
教會有甚麼共同目標?.

$^{81}$有沒有的?.
《聖經經文》今天在第十二節.
保羅告訴我們.
這裡有一個我們可以知道的共同目標.
就是為甚麼?.
要成全聖徒.
各盡其職.
建立基督的身體.
基督的身體這裡所指的.
就指教會或者神的國度.
合因不單止是說.
教會裡面每一個人都是一模一樣的.
說的東西想的東西做的東西都是一模一樣.
不是這樣.
而是強調.
我們通過在神的恩典裡面.
我們有共同合作配合的精神.
我們能夠互相的毀身.
縱然我們每一個人都是不同的性格.
不同的經驗.
不同的背景等等.
故此一個合一的群體.
就不是好像一盤散沙那樣.
而是各自互相的配搭.
以致我們可以向著共同的目標來進發.
這就是合一的教會裡面.
不過教會和外面的一些商業群體是很不同的.
因為當我們基督徒為了我們教會共同的目標努力.
做好我們本份的時候.
你會發覺.
不會有很多明顯可以看得到的利益.
有沒有?.
我們在教會裡面做.
自從有沒有什麼實質的利益.
雖然有人曾經這樣說.
你們基督徒這麼熱心傳福音.
你帶了人上來.
所以是否有擁婚?.
你們為何這麼熱心?.
有沒有?.

$^{121}$你們跟家裡的朋友.
傳福音是否有擁婚?.
給你一個擁抱.
有沒有?.
大家都知道是沒有的.
既然是沒有擁婚.
又沒有很明顯物質上的利益的時候.
為什麼我們還要這麼努力傳福音?.
為什麼我們還要這麼努力邀請我們的親朋戚友來教會?.
為什麼?.
在這裡保羅講了一個很本質上的原因給我們聽.
保羅在這裡說.
當我們聚在一起的時候.
神就用和平讓我們彼此聯絡.
以致我們能夠竭力保守.
聖靈所賜合而為一的心.
我們第一樣要搞清楚的就是.
保羅在這裡說.
我們又要竭力保守聖靈所賜合而為一的心.
即是告訴我們什麼?.
教會的合一.
不是靠我們來重新製造這個合一.
合一是怎樣來的?.
是神賜給我們的.
我們的工作不是要製造.
從沒有變有.
來做這個合一.
而是什麼?.
神已經賜給我們這個信仰群體一個合一.
而我們的工作,我們的責任就是.
持守.
保持著這個合一.
而我們要保持這個合一.
我們要清楚明白.
我們不是參加一個俱樂部.
而是我們是被植入神的家裡.
在這裡.
保羅告訴我們.
七個關乎神所賜給我們的生命.
本質上就有這個合一的事實.

$^{161}$而這個七種合一的意義.
讓基督徒能夠被綁在一起.
被組合在一起.
來去發揮這個信徒在教會裡面的責任.
好了.
大家要明白.
我們的合一並不是來自大家有同一個想法.
同一個意見.
有同一個喜好.
吃的東西都一樣.
不是這樣.
而是在說我們在耶穌基督裡面.
有七種不同的合一的事實.
讓我們在這裡看一看.
保羅告訴我們第一個原因就是.
因為我們有共同的身體.
這個身體是指教會.
即是我們同屬一個屬靈的家庭.
我們都是屬靈家庭的一部份.
同一個性靈.
我們都是受同一個性靈所感動.
來接受領受這個新的生命.
同一個指望.
指望甚麼.
我們每一個基督徒都指望將來我們.
承受神在天上為我們預備的產業.
我們都是同一個主.
我們說主耶穌.
我們都是同一個主.
事奉同一個主.
同一個信.
我們都是信耶穌基督所成就的救恩.
同一個洗禮.
我們都是通過奉耶穌的名得做這個洗禮.
以致我們可以分享進入基督的死亡與復活裡面.
而同一個神.
清清楚楚.
我們所信的都是唯一天上獨一的三一真神.
沒有其他的生世.
我想大家想像一下.

$^{201}$好像一個車輪.
大家都見過.
踩單車的時候有一個車輪.
車輪是一個很簡單的組合.
中間有一個輪芯.
然後一些橫枝連繫著輪芯和外圍.
大家想一想.
如果輪的中間最重要的部份爛了或很弱的時候.
又或者那些橫枝不是連接著輪芯的時候.
你猜這個輪還能不能走得動.
就算真的能走得動也不能走得遠.
聖經告訴我們.
耶穌基督就好像在這個輪的中心的位置.
將所有的東西連在一起.
而當時我們所看到七個關於這個合一的真理的事實.
就是定義了我們連於元首基督耶穌的事實.
就好像這些橫枝.
可能有不同的長度不同的角度.
但是都是連於這個中心.
所以我們基督教是以耶穌基督為中心的意思.
而不是建基於我們自己個別的利益上.
所以雖然我們不知.
雖然我們全福音都是大人歸信耶穌.
沒有湧昏也沒有什麼很明顯很清楚的利益.
但是我們仍然願意投入來維持這個合一.
為什麼呢?.
因為這個就是我們所領受這個新生命本身的本質.
有這個合一的本質.
也是神的命令.
也是我們的責任.
作為基督身體的一部分責任.
還有當我們合一的時候.
我們個別的屬靈生命也都可以得到獎盡.
以致我們每一天更像我們的主耶穌基督.
好了.
接著保羅也跟我們再繼續講下去.
當我們有了明白了這七種合一的條件之後.
保羅更進一步告訴我們.
我們行事為人需要與我們蒙召的恩相稱.
他再具體地說了四個生命.

$^{241}$我們應該有的質素.
來實際上活出這個合一.
哪四種質素呢?.
謙虛,溫柔,忍耐,寬容.
用我們每一種質素來看看.
第一個,謙虛.
什麼叫謙虛?.
謙虛的意思.
不是看低自己.
而是少想一點自己.
這個才叫做謙虛.
所以無論我們有千百個理由覺得我們比其他人厲害.
又或者我們可以有實際上很多原因覺得我們比不起其他人.
但是這裡保羅告訴我們一個很重要的事實.
我們每一個人來到耶穌基督面前都是什麼?.
因著耶穌基督的恩典來得到這個救恩.
所以每一個人來到神的面前的時候.
沒有一個比另外一個更加優越.
或者一個比另外一個更加差勁.
沒有.
在神的面前我們每一個人都是蒙恩的罪人.
在保羅的日子或者甚至今天都是一樣.
講謙虛這件事的話.
通常都會讓人有感覺為什麼你這麼軟弱.
講什麼謙虛.
通常都是比較負面一點的.
不過很有趣.
在聖經裡面.
聖經裡面謙虛其實是一個很強大的力量.
因為它能夠容讓我們.
不是時時刻刻都在想著自己的好處.
相反來說.
如果一個自高自大的人.
是怎樣的?.
經常覺得自己很了不起.
他會期望別人來服侍他.
所以自大你可以說是合一的剋星.
有自大就沒有合一.
第二樣看的就是溫柔.
溫柔是什麼呢?.

$^{281}$搞清楚.
溫柔不是指內弱.
不是指軟弱.
而是一種受到駕馭的力量.
這是很重要的分別.
所以溫柔的人不是軟弱無能的人.
大家搞清楚.
反而溫柔的人是堅強的人.
為什麼?.
因為他選擇不使用凌駕於別人之上的權利.
也許有這樣的權利.
但是他沒有用這種權利.
不會一意孤行.
不會當有不開心或不滿意的事情的時候.
就會冒犯.
反擊,攻擊別人.
不會有這種情況.
我記得有一次.
我去一個食肆訂了午餐.
我叫了午餐之後.
大家都已經吃過了.
我給你一張票.
我看著號碼.
自己找位子坐.
坐下來之後看著牌子顯示.
什麼時候到我了.
我看一看.
過了我的號數.
原來我錯過了.
我立刻出去拿餐.
出到的時候.
原來我的餐還沒出.
原來他們不是按數碼.
不是按號數的排列來出餐的.
總之你看到號碼出來了.
反正已經出了.
舔都舔你.
免得回去.
在這裡等.
等餐來.

$^{321}$幾分鐘都還沒到.
沒理由等了幾分鐘回去.
餐應該快要來了.
繼續等.
五分鐘,十分鐘,十五分鐘.
等的時候.
我看見那個女生.
那個在前台服侍.
在廚房出餐的時候.
出了餐之後.
給餐的顧客.
一個人可以很有效率地工作.
很快,很有效.
很細心,很仔細.
那個女生看見我在這裡.
很耐心地等.
終於我的餐來了.
我都不知道等了多久.
來到的時候.
那個女生第一句跟我說.
不好意思.
真是不好意思.
要你等這麼久.
如果你是我.
你會說什麼.
我就跟她說.
你真是很有效率.
很快,很細心,很準確地工作.
你知道她怎樣回答我.
她反應如何嗎.
她看一看我.
然後就低著頭.
繼續做她的事.
一句話都沒說過.
說什麼這麼沒禮貌.
我一點都不覺得她沒禮貌.
我反而覺得在那個時候.
她不知道該說什麼.
我相信她是沒想過有人會這樣跟她說話.
她不知道她的反應如何.

$^{361}$繼續做她的事.
這個是實際上.
很多時候會發生的.
你要我等這麼久.
我當然會投訴你.
你為什麼要我等這麼久.
叫你的經理來見我.
在今天的消費者權益裡面.
也希望我們可以這樣做.
但是溫柔的意思是什麼.
溫柔的意思就是我們願意放棄我們可以有的權利.
而將這個權利轉化成為一種安慰別人的力量.
這個就是溫柔.
好,忍耐.
什麼叫忍耐.
忍耐不是單單忍受環境.
我忍著等我老婆化妝這麼久都搞不定.
或者忍耐等我老公去廁所這麼久.
或者我忍著等我的菜這麼久都沒到.
聖經裡面所講的忍耐不是這種忍耐.
而是在講我們在以弗所書第四章所講的.
忍耐的意思其實是講我們信徒在教會裡面.
如何互相的相待.
它是指到一種能力.
能夠去忍受別人的弱點.
他的失敗.
他跟我的不同.
而不放棄他們.
仍然視他們為我們的弟兄姊妹.
好像賽跑一樣.
你走得很快.
後面的人走得很慢.
你會怎樣.
你會不會回頭看.
你走得這麼慢,走快點.
還是.
你願意等待他.
甚至幫他.
這個就是那個分別.
這個就是忍耐.

$^{401}$在這個世界裡面.
今天的文化裡面.
很多時候都是什麼.
當我們受到別人的批評的時候.
我們怎樣.
就覺得很生氣.
當我們遇到一些不方便的時候.
投訴.
當別人與我們不同意見的時候.
我們怎樣.
跟他疏遠.
不同意見就朋友.
但是一個忍耐的人.
不會因為某一些不方便.
或者不開心.
不喜悅.
就即時爆發.
最近大家有沒有聽過新聞.
香港有一個中學的校長在新加坡.
因為有些不開心,爆發.
爆發完之後.
他要承擔很大的後果.
雅各在這裡給了我們一個很好的忠告.
雅各他怎樣說.
他說你們國人要快快的聽.
慢慢的說.
慢慢的動腦.
在教會裡面.
如果我們要維持教會的合一.
忍耐極其重要.
為什麼.
因為教會裡面所有的人都不是完全的人.
我們每一個人都不同.
我們每一個人都跟其他人有不一致的地方.
如果我們個個都同意我們的時候.
合一很容易.
但問題當教會不是這樣的時候.
你有你的看法.
我有我的意見.
教會就需要有忍耐.

$^{441}$以至我們能夠持守教會的合一.
第四個寬容.
寬容我們要搞清楚.
很多時候我們都覺得寬容.
不要看得這麼緊張.
將你的要求降低一點.
不是.
寬容的意思不是降低我們對真理的持守.
寬容的意思是要讓我們活出韌點.
保羅在這裡說.
教會大家要記住.
教會從來不是完美的人活在一起.
教會是一群蒙恩的罪人.
倚靠神的恩典嘗試活在一起.
所以我們經常說教會不要搞錯.
教會不是聖人的博物館.
什麼意思?.
就是不存在.
教會從來都是一群罪人的醫院.
大家來到這裡都是一個被醫治.
走在神聖路上的一群人.
當然保羅千萬不要誤會.
保羅這裡不是說我們對罪.
就是直眼開直眼閉.
不要這麼緊張.
不要要求高.
有時候對這些事情忍氣吞聲.
不要理他.
聖經所說的寬容不是說這些.
所以當我們看到別人明明犯罪的時候.
你千萬不要跟他說.
我看著主耶穌奉上.
我忍你.
聖經所說的寬容不是這樣.
因為這種寬容是消極甚至是錯的.
保羅這裡所說的就是.
當我們有不同的地方.
當我們可能相處的時候出現困難的時候.
保羅這裡的教導是什麼?.
用愛心互相寬容.

$^{481}$沒有愛的時候.
告訴大家就沒有寬容.
這麼簡單.
但是我們也要留意.
在彼得前書裡面所說.
所以他說我們需要彼此切實相愛.
因為愛能遮掩許多的罪.
所以剛才我所說.
沒有愛就沒有寬容.
但是我們也要留意一樣東西.
就是有愛的地方不一定沒有摩擦.
不一定沒有衝突.
真正的愛是能夠讓我們在衝突裡面.
仍然願意彼此接納彼此包容.
為什麼?.
因為保羅在哥羅西茲說.
我們需要彼此饒恕彼此包容.
為什麼?.
因為耶穌基督已經饒恕了我們.
既然耶穌基督饒恕了我們的時候.
我們怎能夠不饒恕得罪我們的人.
這就是寬容背後的比較.
我們能夠寬容別人的力量.
這就是我們互相寬容的意思.
當我們明白了我們領受新的生命.
有著這種合一的本質.
也有著能夠活出這個合一的聖情之後.
保羅再進一步告訴我們.
在這樣的情況.
我們每一個信徒就會連於元首基督.
也靠著他來互相割愛割戚.
來發揮彼此相助的功能.
以致屬靈的身體可以漸漸增長.
保羅這裡說的是什麼?.
百折割愛割戚.
很明顯告訴我們什麼?.
告訴我們教會裡每一個人都有一個職份.
大家知不知道你在教會裡的職份是什麼?.
你不知道,有嗎?.
沒有吧?.

$^{521}$有些人可能會謙虛上身.
牧師兄我什麼都不懂.
什麼都不會做.
你不要期望我對教會有什麼貢獻.
首先你要搞清楚.
剛才我們說過什麼?.
謙虛.
不是貶低自己.
貶低自己不是叫做謙虛.
或者這樣問大家.
你覺得教會裡的會友是否可有可無?.
是不是?.
或者用另一個角度問大家.
你覺得身體裡有沒有部份是沒用的?.
有.
什麼?.
盲腸.
你不知道盲腸是沒用的嗎?.
盲腸是真的沒用的.
我上網做一個調查.
做一個研究.
他說沒錯.
很多人都以為盲腸是沒用的.
可以隨時切去割去.
沒問題.
不過原來有調查.
再進一步研究.
發現原來盲腸沒有很重大的功用.
但是原來它是幫助免疫力和腸部健康的.
換言之.
或多或少.
身體每一個部份都有它的功能.
同樣.
教會裡每一個成員都應該有它的功能.
不是的牧師.
我什麼都不懂.
真的嗎?.
保羅在第七節告訴我們.
我們各人蒙恩.
都是照基督所量給各人的恩賜.

$^{561}$這裡講得很清楚.
各人.
就是說我們每一個人都在神的面前.
領受祂所給過我們.
為我們度身訂做的恩賜.
所以神是期望什麼?.
神是期望我們每一個信徒運用神給過我們的恩賜.
服侍它.
貢獻教會.
如果神給了我們恩賜.
我們拒絕.
不用它.
或者甚至不想知道它是什麼.
我給大家聽聽.
將來有一天我們在神面前交仗.
不過當然.
我們都有很多客觀的環境因素.
可能是健康問題.
可能是年紀問題.
或者我們在某一個人生階段要負的責任.
也許我們不能夠在教會承擔一些很重要很巨大的責任.
但是你會留意.
教會總是有一些地方.
無論是什麼人都可以參與.
什麼?.
為教會祈禱.
為你周圍的人祈禱.
你每一個星期來到教會.
坐在禮堂的敬拜.
都會成為別人的祝福.
當你坐在旁邊的那個你不是太熟悉的時候.
跟他打個招呼.
將你在神面前所領受的恩典.
告訴他.
這一種都是一種互相配搭的事奉.
敬拜本來就是一種事奉.
所以當我們依靠神的能力.
運用神所賜給我們的恩典.
讓我們連於我們的主耶穌基督.
這樣.

$^{601}$聖經告訴我們.
我們自己的屬靈生命就會得著長進.
耶穌基督的身體就會得著成長.
就好像我們剛才所說的那個車輪.
當我們每一個人好像那個橫枝.
連於那個車輪的中心部分的時候.
這個就是教會的合一.
這個車輪就能夠行走.
這個車輪就能夠同一個方向.
向著教會的目標來進發.
所以弟兄姊妹.
希望今天的信息能夠提醒大家.
我們必須要連於我們所相信的耶穌基督.
共同努力.
持守聖靈已經賜給我們的教會合一.
以致我們可以跟其他的信徒.
一同互相學習.
互相包容.
去成長.
去建立耶穌基督的身體.
為要讓神的榮耀得著高舉.
活出這個真理.
以致你在教會當中能夠得著你當得的榮耀.
多謝你給我們一切這樣的祈禱,交托.
奉耶穌基督得勝的名字而求.
阿們.
\newpage



\section{哥林多前書 8:1-13}
\label{sec:032212}
\textbf{因基督的愛而造就人}~~~~馮志成牧師
\newline
\newline
連結: \href{http://www.church.com.hk/acms/content.asp?site=cdc\&op=show\&type=product\&code=032212\&layout=sermon}{\texttt{032212}} ~~~~ 語音日期: 2026-05-31
\newline
\newline
\hyperref[sec:032211]{\small{< < < PREV SERMON < < <}}
~
\hyperref[sec:index_chronic]{\small{[返順時目]}}
~
\hyperref[sec:032213]{\small{> > > NEXT SERMON > > >}}
\newline
\newline
哥林多前書 8:1-13
\newline
\begin{longtable}{cl}
\hline
\hline
章節 & 經文 (和合本修訂版)\\
\hline
8:1 & \begin{tabularx}{0.7\textwidth}{X} 關於祭過偶像的食物，我們曉得「我們都有知識」，但知識使人自高自大，惟有愛心能造就人。 \end{tabularx} \\ \\ \relax
8:2 & \begin{tabularx}{0.7\textwidth}{X} 若有人自以為知道甚麼，他其實仍不知道他所應當知道的。 \end{tabularx} \\ \\ \relax
8:3 & \begin{tabularx}{0.7\textwidth}{X} 若有人愛神，他就是神所認識的人了。 \end{tabularx} \\ \\ \relax
8:4 & \begin{tabularx}{0.7\textwidth}{X} 關於吃祭過偶像的食物，我們知道「偶像在世上算不得甚麼」；也知道「神只有一位，沒有別的」。 \end{tabularx} \\ \\ \relax
8:5 & \begin{tabularx}{0.7\textwidth}{X} 雖然在天上或地上有許多所謂的神明，就如他們中間有許多的神明，許多的主， \end{tabularx} \\ \\ \relax
8:6 & \begin{tabularx}{0.7\textwidth}{X} 但是我們只有一位神，就是父，萬物都出於他，我們也歸於他；並只有一位主，就是耶穌基督，萬物都是藉著他而有，我們也是藉著他而有。 \end{tabularx} \\ \\ \relax
8:7 & \begin{tabularx}{0.7\textwidth}{X} 可是，不是人人都有這知識。有人到現在因拜慣了偶像，仍以為所吃的是祭過偶像的食物；既然他們的良心軟弱，也就污穢了。 \end{tabularx} \\ \\ \relax
8:8 & \begin{tabularx}{0.7\textwidth}{X} 其實，食物不能使我們更接近神，因為我們不吃也無損，吃也無益。 \end{tabularx} \\ \\ \relax
8:9 & \begin{tabularx}{0.7\textwidth}{X} 可是，你們要謹慎，免得你們這自由竟成了軟弱人的絆腳石。 \end{tabularx} \\ \\ \relax
8:10 & \begin{tabularx}{0.7\textwidth}{X} 若有人見你這有知識的在偶像的廟裡坐席，而這人的良心是軟弱的，他豈不放膽去吃那祭過偶像的食物嗎？ \end{tabularx} \\ \\ \relax
8:11 & \begin{tabularx}{0.7\textwidth}{X} 因此，基督為他死的那軟弱弟兄，也就因你的知識沉淪了。 \end{tabularx} \\ \\ \relax
8:12 & \begin{tabularx}{0.7\textwidth}{X} 你們這樣得罪弟兄，傷了他們軟弱的良心，就是得罪基督。 \end{tabularx} \\ \\ \relax
8:13 & \begin{tabularx}{0.7\textwidth}{X} 所以，食物若使我的弟兄跌倒，我就永遠不吃肉，免得使我的弟兄跌倒了。 \end{tabularx} \\ \\
[1ex]
\hline
\hline
\end{longtable}
$^{1}$各位朋友 各位弟兄姊妹 早安.
今天可以在正道用神的話.
和大家一起分享祂是怎樣改變我的生命.
正如剛剛傳道所說 我是個牧師.
我亦在家庭更生協會服侍了很多年.
有教會裡面或者我們身邊的朋友.
都很羨慕我太太.
他們說你又好 老公這麼好.
又做牧師 心地善良.
又去帶加薪 肯定哄得很開心.
你知道我太太怎樣承認嗎?.
你猜猜.
我太太說是的 她的確是個好的牧師.
她在加薪事奉得很好.
但這是她的角色.
這是她的工作.
這是她的事奉.
但不代表她是個好丈夫.
神給她很多東西不代表她會用.
她這樣說並不是沒有原因的.
因為我自己在教會服侍的時候.
我慢慢去讀聖經的時候.
我講道的時候.
我越來越發現.
我可以給到我太太的愛真的很有限.
有一天我很坦誠地和我太太說.
其實我知道我是很愛你的.
但我發現我愛你的能力.
是去不到一天二十四個小時.
一個星期七天.
不間斷地愛.
我發現我沒有這個能力.
我有這樣的心 但我做不到.
我發現天上地下.
唯有耶穌可以這樣做到.
我太太回答我.
我嫁的是你 不是耶穌.
所以各位朋友.
如果你是我 你會怎樣回答?.
你給我一點辦法.

$^{41}$我們還未揭聖經.
所以你不要指望有答案.
不過你回答吧.
我覺得我沒有問題.
我已經很好.
我忠心服侍教會.
愛弟兄姊妹 愛上帝.
也愛太太.
但我發現我有局限.
我是個罪人.
我沒有可能全天候地愛著她.
我很坦誠.
她那天多讚我兩句.
我會愛她多一點.
她那天不是給我面色看.
她那天只是沒有對著我笑.
我總覺得我愛她好像有點虛偽.
但從我太太的角度來說.
不是應該這樣嗎?.
我嫁給你.
聖經也說了.
你們做丈夫愛太太.
你要為太太捨命.
我都還沒有要你的命.
你就說NO.
你就卸膊給耶穌.
你會怎樣回答?.
我覺得我們兩邊都對的.
不過如果我太太這樣問我.
弟兄 幫幫忙.
或者你說.
哎呀 Jason 你今天只是騙我.
今晚我太太問我怎樣回答.
有個弟兄給我一個辦法.
我覺得是挺好的.
他說 Jason 有沒有想過.
現在的AI這麼好.
你要坦誠一點.
將你所有寫過的講章.
講過的道.

$^{81}$你所有的靈修筆記.
全部調給AI.
將這個APP.
設在你太太的手機上.
有時候你太太找不到你.
或者你們剛剛吵完架.
她又不方便跟你說話.
你又不方便跟她說話.
那怎樣呢?.
打開那個AI.
她有什麼想跟你說的.
通過AI回答她.
AI發個copy給你.
你又知道那個AI的你回答了什麼.
可以補鑊的.
行不行?.
有人搖頭不行.
有沒有人行的?.
有人搖頭.
其實現在我覺得AI是很影響我們每天的生活的.
因為好用.
我坦白說好用.
甚至乎這篇講章.
你都可以懷疑一下我有沒有用AI的.
因為它可以幫我糾集很多的資料.
用不同的神學家的角度去看經文.
當我正是很開心地.
使用AI這麼好的功能的時候.
我都面臨另一種挑戰.
我的弟兄姊妹跟我講.
她說牧師啊.
我覺得有一天我打電話給你.
你沒有接我的電話.
接著我打開那個AI.
我叫她扮演一個基督教的牧師跟我聊天.
我跟她聊了五分鐘而已.
就很滿足了.
所以我那天就沒有再call你了.
如果你教會有這麼好的弟兄姊妹.
好不好啊Nelson.

$^{121}$不會再煩你了.
他們有什麼問題自己用AI解決了.
行不行呢.
值得想一想.
AI的確是給我們很多知識.
現在AI已經去到博士級的知識水平.
那為什麼還要牧師呢.
牧師有時候又不是24小時7天on call.
有時候你半夜兩點起床.
睡不著很煩.
你試試call牧師.
看牧師接不接你的電話.
不會的嘛.
那AI會的.
我約牧師還要請牧師給錢給牧師請喝咖啡.
問AI又不用錢.
究竟知識豐富了.
是不是真的可以解決我們所有問題.
或者是不是我知識豐富了.
我就可以滿足到我太太.
今天我們回去聖經看一看.
其實人對知識的那種追求.
或者人對知識的某一種.
我稱之為迷信.
不是我們這個時候才有.
不是因為有AI我們才有.
其實在兩千年前.
哥林多教會裡面已經是這樣了.
所以我們在哥林多前書.
我們如果要了解保羅說什麼.
其實是先要了解.
究竟他們住在一個什麼樣的城市裡面.
首先哥林多這個城市.
它被稱為一個富庶的城市.
因為它是一個港口.
它是連結歐洲和亞洲的一個城市.
兩個港口.
它很有錢.
因為做貿易.
它屬於一個羅馬帝國下面的城市.

$^{161}$每個城市都有一個他們認為的神明.
去祝福或者保護這個城市.
而哥林多城的保護神.
就叫做阿弗洛蒂尼.
其實是愛的意思.
愛神.
這個愛的神廟.
或者愛神所居住的神廟.
裡面就有住著一千名妙妓.
或者我們叫她妙妓.
其實他們對外面宣稱.
她們是聖女.
神聖的聖.
這樣.
那又怎樣呢?.
這個時候我們沒有這個文化.
現在你回教會也好.
或者去到佛廟也好.
你不會見到這些聖女在那裡.
沒有的.
但當時是怎樣呢?.
當時候的文化.
其實你可以想一想.
當日以色列人出埃及的時候.
他們在曠野.
不是弄了一個金牛毒出來的.
弄金牛毒的時候.
他們指著金牛毒說.
這個就是帶我們出來.
從埃及出來的耶和華.
接著他們做什麼呢?.
接著他們就敬拜金牛毒.
但為什麼金牛毒不是我們裝香拜神般敬拜.
聖經說的是.
那些以色列人就在曠野那裡.
在金牛毒的面前.
坐下吃喝.
起來喚灑.
和合本將它翻譯得很文雅.
其實是坐下吃喝.

$^{201}$真的是坐下吃肉.
但起來喚灑的意思不是小朋友玩遊戲.
是在說男女交合.
你說好像很醜很噁心.
為什麼會這樣呢?.
是因為當時他們對神明的敬拜文化是什麼呢?.
如果這個神是管生殖的.
生人多的時候.
你怎樣令這個神明喜悅呢?.
就是在神明面前去吃肉.
和去做這些男女交合的事情.
令這個神明喜悅.
令這個神明喜悅之後.
他就會祝福這個民族.
令這個民族的生殖的人口會增加.
這個當時他們的文化是這樣.
所以去到哥林多的時候.
或者在那個歷史背景下.
哥林多的人也做類似的事情.
他們富庶.
他們有錢.
他們要增加人口來做貿易.
提供城市的服務.
那又怎樣呢?.
於是他們男男女女.
都去到這個神廟裡面.
去吃肉.
去行男女之事.
但他們這樣做不是滿足自己的慾望.
他們這樣做是深信他們要取悅那個神明.
令這個神明祝福這個城市.
以至於每一年.
這個神廟會做一次巡行.
一千個這些聖女.
就會在哥林多的大街小巷巡行.
令整個城市知道.
這個城市是由這個神明.
這個愛神保護的.
那個場景如果用我們中國人明白的文化.
你可以說那一天整個哥林多走馳玉林.

$^{241}$甚至當時有一個很獨特的.
希臘的文字.
動詞.
大家不用懂讀了,我寫給大家看.
那個動詞的意思就是.
描述那些人生活很奢侈和墮落.
好像哥林多人那樣放蕩生活.
你可以知道那時哥林多這個城市是一個怎樣的城市.
某個層面上我剛剛看到你們有宣教士在巴黎.
在法國.
其實法國的文化我覺得有點相似.
很富有.
但是這些生活很放蕩.
浪漫之都.
我們現在說浪漫之都.
但是你又知道裡面很多藏污立舊的東西.
在這裡.
這些人在哥林多信主了.
但是你要知道這些人是從小到大.
包括可能他小時候見過爸爸.
去過神廟吃喝玩樂.
甚至他自己年輕的時候.
也和他的朋友去過神廟吃喝玩樂.
現在他信了主.
他成為基督徒.
那你覺得他會怎樣呢?.
他會停下.
不敢做呢?.
還是繼續會做呢?.
你猜猜.
如果他放在今天我們的華人教會.
我相信他會停下.
因為我們的牧者.
弟兄姊妹說.
不可以這樣做的.
你這樣做是道德敗壞的.
但是當中也有些人.
你要知道.
他是從小到大都思空見慣這些東西.
他覺得如果我不這樣做.

$^{281}$我們的城市會不會.
就是不幸運的.
如果用一些我們現在明白的位置.
我也接觸過這些朋友.
我不知道你身邊有沒有.
有沒有一些人小時候.
或者你的朋友.
小時候被迫.
因為你家人幫你做的.
將你過契給了觀音.
有沒有這些人.
如果有的話.
你會明白那種張力.
因為他們很害怕.
耶穌好.
不過我小時候過契給了觀音.
耶穌多好.
到時候我不再認觀音.
那有沒有什麼不幸運的事情呢?.
你跟他說.
不會的.
你牧師跟他說.
不會的.
不要緊耶穌是萬有的主.
一切天上地下.
只有祂是真神.
你那天說得很爽.
但是如果那天不幸運.
真的有時候有些機會.
那天如果剛好.
過契給觀音的弟兄姊妹.
初信主的.
他回到家裡.
剛好那天他決志.
回家之後突然間爆水管.
你覺得他怎麼想呢?.
你會覺得他說.
幸好今天順了主.
爆了水管.
馬上發現不會浸濕整個房子.

$^{321}$或者說.
哎呀.
是不是我今天順了主之後.
觀音懲罰我呢?.
保羅現在在處理這樣的張力.
在教會裡面.
我希望大家讀下去的時候.
讀得到保羅那種掙扎和糾結.
那保羅怎麼看哥林多教會呢?.
我們要花點時間去看看.
保羅和哥林多教會的互動.
在史萊恆的十八章提到.
是保羅在宣教旅程建立了哥林多教會.
所以他是哥林多教會的創始人.
後來他在哥林多住了十八年之後離開.
一直保持著和哥林多教會的書信來往.
但是他大部分是寫信的.
所以你在臨前的第五章.
保羅說我以前寫過信給你們說什麼什麼.
你也知道那時候沒有電話和電郵寫信.
經過一段時間.
之後就寫了我們今天讀的哥林多前書.
後來可能是這個書信給了哥林多之後.
可能哥林多人好像更加反感.
或者更加不喜歡保羅.
於是保羅不得不再次去哥林多.
這是第四點.
但是那次探訪就更加差.
保羅就說那次是很哀傷的探訪.
那次探訪完之後保羅離開哥林多.
寫了一封書信給他.
那封書信保羅說是一個流淚的書信.
哥林多後書第二章提過這封書信.
之後就說我再也不去了.
後來去哥林多後書就說.
我或者也會來的.
因為你有所改變.
我們相信在流淚書信之後.
哥林多教會是有改變的.
所以保羅在哥林多後書提過.

$^{361}$我或者會去.
你看到保羅自己親自建立這個教會.
但是後來這個教會好像不服保羅的教導.
保羅要教他們.
但是又愛他們.
所以在哥林多前書第四章.
保羅是這麼稱哥林多教會裡面的弟兄姊妹.
他說我把你們當作我親愛的兒女一樣.
現在他說的不是牧師和會眾的關係.
是父親和子女的關係.
因為保羅說我在耶穌基督裡.
借著福音生了你們.
保羅說我帶領信主的.
我肯定很愛你.
所以哥林多前書後書.
不是一個神學家寫給神學生的一份教學大綱.
而是一個父母對伴役子女的一份愛的書信.
希望子女糾正過來.
希望子女成熟.
但是這種不是出於父母憤怒的宣洩.
而是充滿愛的一種表達.
如果你們看見保羅和哥林多教會.
這份愛恨交纏的關係的時候.
你會明白舊約裡面.
耶和華神和以色列人的關係也很相似.
以色列人一次又一次的伴役.
但是耶和華上帝一次一次去管教他們.
赦免他們 令他們迴轉.
令他們繼續回到神的懷抱.
但是好景不長.
以色列人又一次伴役.
神再次召喚他們回來.
在這種糾纏裡面.
其實不斷被我們醞釀著某一樣東西.
就是你覺得以色列人真的不懂得神的律法嗎?.
你覺得哥林多人真的不知道保羅所宣稱的福音嗎?.
但是為什麼以色列人知道神的律法.
領受了神的律法.
哥林多人也聽了福音成為基督徒.
但是為什麼他們的樣子.

$^{401}$或者他們表現出來的東西.
不是按照他們知道的東西來的?.
好像他們不懂得福音.
以色列人好像不懂得神.
但是他們覺得他們懂得.
因為我們有十誡.
我們有摩西的律法.
我們有先知講道.
哥林多人說「我回教會的」.
「我有牧師的 保羅你帶我信主的」.
保羅說「不是」.
因為在保羅收到哥林多人寫給他的信.
即是教會寫給保羅的信裡面.
提到一樣東西.
就是關於祭過偶像的食物.
這些祭過偶像的食物.
就是剛剛我和你說過的.
在神廟裡面那些.
獻祭過給愛神的肉.
那個焦點其實是在說.
那些肉我們怎麼處理?.
可不可以吃?.
正在問這件事.
保羅先給一個很肯定的回應.
保羅說「我們曉得」.
「我們都有知識」.
可能是當時候他們的人的口頭禪.
即是大家都懂得一些東西.
我們都很懂得東西.
保羅說「是啊 你也懂 我也懂 大家都懂」.
不過保羅馬上就有一個轉折.
他說「知識使人自高自大」.
「愛人 愛卻能造就人」.
那你一看這句話.
你不理會上下文.
你很容易誤解.
誤解就說.
「知識叫人自高自大」.
「愛就可以造就人」.
如果我們真的這樣按字面去信的話.

$^{441}$你千萬不要送你的兒子去補習.
你送他去補習叫他「自高自大」.
「愛他就可以了」.
「他愛我就足夠了」.
這樣就很糟糕了.
因為我們不理會經文上下文.
我們不理會是什麼.
就按著字面去解釋.
接下來保羅又說了兩類人在教會裡面.
第一類人就是.
有些人自以為什麼都知道.
即是自以為知道.
「那麼他應該知道的 他還是不知道」.
那我就問他究竟知不知道.
接著我再問他究竟保羅要他知道些什麼.
其實用我們的文法來說.
保羅就說教會有些人懂一點點.
扮代表.
在說這件事.
但是我要追問的.
他懂些什麼呢?.
他又在代表些什麼呢?.
所以保羅暫時沒有回答.
接著第二件事.
教會裡面有另外一類人.
這類人比較好.
這類人是愛神的.
而且神是認識他的.
但他又不說這些人懂些什麼.
很有趣吧?.
保羅你在說什麼呢?.
接下來我和我老婆要商量.
要不要送我兒子去補習.
說些什麼呢?.
我繼續看下去.
首先保羅會先去回答.
什麼叫做我們都有知識.
我們究竟知道些什麼.
在原文裡面.
知識,知道是同一個紙巾.

$^{481}$保羅繼續在四到六節說.
我們會知道一件事.
是大家都知道的.
教會裡面人人都知道.
甚至教會外的人.
他們知道.
不過他們不相信.
不同意而已.
他們知道的.
知道什麼呢?.
關於吃西方偶像的食物.
就是飲食.
我們知道三件事.
西方的偶像算不得什麼.
就是沒有偶像.
也知道神只有一位.
沒有別的神.
只是說是一個真神.
其他都是假的.
第三件事.
一樣意思.
雖然天上地下有所謂的神明.
猶如很多神明,很多主.
當時他們叫的神除了什麼神.
都是以神為主人的意思.
保護祂.
愛神也是神.
保護祂.
保羅說這三句.
四至五節就說.
只有一位神.
第六節.
保羅在神上面有少許具體化.
他在說.
我們真的只有一位神.
就是父萬有同祂而來.
我們為祂而存在.
人和神的關係.
在這裡有第一個角度.
我們活著是為了這位天父耶和華.

$^{521}$接著保羅說.
我們也有一個主.
就是耶穌基督.
萬物都是借祂而存在.
我們也是借祂而存在.
就是說我們為神而活的時候.
還要靠著耶穌才有這個能力.
有這個生命.
有這個可能為神而活.
否則我們沒有能力為神而活.
所以保羅首先在說.
那麼多人.
大家都知道一些事情.
你也知道.
我也知道.
就是說.
只有一位神.
耶和華上帝.
只有一位主.
耶穌基督.
大家沒有異議.
所以剛剛保羅提到的兩類人.
無論是知道的人.
或者愛神的人.
都知道這些知識.
接著保羅說.
那麼我們知道知識又如何呢?.
我們知道這麼多事情.
會成為一個怎樣的人呢?.
我們今天可以用人工智能來解經了.
那麼會不會令我們的愛心多一些呢?.
保羅首先指出.
在第二類人裡面.
他們不是說不懂這些東西.
他們不是很相信這些東西.
所以他所說的.
並非所有人都有這知識.
他的知識是如何運用這些知識.
有些人雖然知道.
耶和華是一位獨一的神.

$^{561}$耶穌是獨一的主.
但是這些人從小.
用我們的說法.
過氣給觀音.
他從小到大.
要吃祖父祖母祭祖先神台給他吃的東西.
我就吃過了.
我小的時候.
回到我祖父祖母的家.
牆角有一個格子.
有一個逢佳列代的牌位.
逢年過節.
那裡有三杯茶三杯酒.
一個蘋果.
有隻雞腿在上面.
每一次我回去.
因為剛剛好我是長子的孫.
我爸爸是大兒子.
我又是大兒子.
我家裡只有一個.
我又是大兒子.
我奶奶會很鄭重其事的.
將那三杯茶.
因為小孩子不能喝酒.
三杯茶.
蘋果和雞腿放在我面前.
吃吧.
小孩子說很肚餓.
當然吃吧.
為什麼只有我吃.
其他的堂妹都不能吃.
你明白嗎?.
我想分給他吃.
他說不行不行.
只有你能吃.
習慣了嘔脹.
我就說為什麼只有我能吃.
因為你是長子的孫.
而且這些是我們祖先的口水.
舔過的.

$^{601}$你喝了之後.
吃了之後讀書會好一點.
習慣了嘔脹.
我從小我奶奶不會騙我的.
我奶奶不會害我的.
但是我就在這種環境長大.
直到現在.
我做牧師.
但是如果剛剛信主的時候.
也掙扎過一下.
他看到這些食物的時候.
他信心就軟弱了.
他說.
他也能吃這些東西嗎?.
那就是說我們要神的幫手.
又要主的幫手.
也要祖先的幫手.
很合理的.
好了 第十節.
這是文化結構的問題.
他講過第二種人的另一種.
如果有人看見你有知識的人.
這些是有知識的.
神只有一位.
坐在偶像廟裡坐直.
什麼叫偶像廟裡.
就是那個愛神的神廟.
就是他們從小到大去的地方.
他們的生活習慣沒有改變過.
繼續去那個廟裡吃東西.
那他就會害怕.
喂 你們不是只有一位神嗎?.
為什麼你要去愛神那裡呢?.
有知識的人就這樣回答.
沒有神的.
愛神的廟這些是騙你的.
哪裡有愛神呢?.
這些肉低於什麼呢?.
什麼都沒有.
只是一塊肉而已.

$^{641}$那你就吃吧.
沒事的.
有知識的人.
這樣向那些信心軟弱的人解釋.
使自己的行為合法化.
保羅是如何回應這些有知識的人呢?.
第八,第九節.
他就說.
既然你有知識.
知道這些在凡人眼中獻給偶像的食物.
是沒有意義的.
沒有什麼神明的力量.
那你吃這塊肉.
對你又沒有什麼益處.
尤其現在在澳洲.
吃多一塊肉.
你就多跑兩圈.
不然就是脂肪肝,心臟病.
你不吃又沒有損失的.
第九節他說.
然而你要小心.
千萬不要讓你的權力.
成為叛島弱者的絆腳石.
就是說你吃不吃是你的自由.
大家都知道是一塊肉而已.
不過你要知道.
有一些人.
那些從小給觀音的那些人.
他們知道這個知識.
不過不是很相信這個知識的人.
看到你這樣.
他們心裡不安.
他們就以為.
哦!即使成為了基督徒.
我繼續還要靠著吃那些肉來祝福我,保護我.
就是說我同時需要腳踏兩船.
左腳就是耶穌的教會十字架.
右腳就是觀音祖先舔過的口水.
保祿就說.
喂!你的知識沒有問題.

$^{681}$你信你的知識沒有問題.
不過那些人還沒有準備好.
那些人仍然軟弱.
你明白那種張力嗎?.
我曾經是這樣的.
我和太太說.
喂!我可以很愛你.
不過我做不足.
只有耶穌可以做得這麼好.
你覺得我太太心裡不明白.
耶穌是那位愛她,救她的主嗎?.
你覺得我太太心裡不知道.
耶穌可以二十四小時,七天,三百六十五天.
不間斷地愛祂,保護祂,帶領祂,感動祂嗎?.
我太太有這個知識.
不過我太太說.
但我嫁給的是你.
我還是不可以單靠我的理性.
我不可以單靠我自己的乾巴巴的頭腦知識.
我去經歷到神的愛.
不是說沒有神的愛不足夠.
所以她會挑戰我的.
她會說.
那你為什麼不想想.
神使用你作為一個渠道愛我,祝福我呢?.
這些就是將軍了.
我們這些很熟悉的.
那你不能反駁的.
是的,神可以使用我來愛我的太太.
關心我的太太.
令她經歷到耶穌基督是多麼愛她.
我不可以拒絕耶穌這麼做.
也是這個意思.
我不可以那時候跟太太說一篇道理.
你相信我吧,我真的愛你.
聖經是這麼說的.
聖經已經讓我娶到你了.
那還不是愛你?.
我不愛你,娶你幹什麼?.
但是我太太有她的需要.

$^{721}$她也知道我所說的全部是對的.
所以有一次我太太跟我開玩笑.
她說:其實你已經說了這麼多次.
你又懂得希伯來文,希臘文這麼多愛的寫法.
你懂得這些之後.
你覺得你愛我多了一點呢?.
那一刻我問了問我自己.
好像真的沒有.
因為我說出口心裡一寒.
我懂這麼多東西,研究這麼多東西.
但是我沒有辦法.
也幫不了我所說的我愛的對象.
感受得到我對他的人與人之間的愛會多一點點.
我很掙扎.
所以保羅在這裡也在提醒哥林多教會裡面的人.
那些很有知識的人.
他們可能認識很多神學.
他們可能講道講得很好.
他們對知識的理解很全面,很深入.
很廣闊,融會貫通.
不過保羅就說:你們這一批人.
好像對另外一批知道福音的人.
好像沒有什麼幫助.
接下來保羅沒有停留在人與人的關係.
保羅將它上升到我們與基督的關係.
11至13節.
保羅就說:如果我們不在意那些所謂軟弱的人.
就說:你多讀幾次聖經.
你禱告多一點.
你成熟一點.
你信主多一點.
我們就可以一起去神廟吃肉.
你不軟弱.
我就得到自由.
你為什麼要用你的軟弱.
限制我的自由?.
保羅對這些人說.
如果你這樣跟你弟兄姊妹說話.
如果你這樣對初信說話的時候.
保羅就說:那些軟弱的人.

$^{761}$你要記住耶穌基督是為他們而死的.
耶穌連命都給了他們.
使他們可以生命改變更生.
但是當你說你的自由,你的權力,你的知識.
很重要的時候.
你是弄死了他們.
你是叛倒了他們.
12節他就說.
當你這樣犯罪對不起弟兄.
當你這樣打擊他的軟弱良心.
就是犯罪,干犯基督.
這個就很要命了.
現在說的不是你和弟兄的關係.
是你和耶穌基督的關係.
你以為你堅持的是對的.
的確你堅持的知識本身沒有問題是對的.
這是一位神.
沒有偶像的.
所以那些肉放在金錠的像面前.
那塊肉形成那塊肉.
吃與不吃沒有問題.
吃都是你的自由.
但是那些軟弱的不是這樣看的.
如果你不再理會那些軟弱的人.
繼續按著你的知識去行使你的自由的時候.
那些軟弱的人就開始覺得.
我也要吃那塊肉.
我不吃那塊肉是不行的.
保羅說當你這樣做.
其實你就是不在意弟兄.
你就是干犯基督.
因為基督為他而死.
如果讀到這裡.
你就開始明白.
保羅在第二節說.
他說「某些人自以為知道什麼.
然而他應該知道的.
還是不知道」.
你開始明白保羅在說什麼了.
保羅在說.

$^{801}$某一些人.
自己神學很厲害.
或者背聖經很厲害.
很容易完全明白福音.
都會如流.
但是當你面對那些軟弱的弟兄的時候.
他們也知道你的知識.
不過他們不是100\%相信的時候.
你會行使你的自由.
做你喜歡做的事情.
不顧他們的時候.
你以為你已經背下了.
說清楚福音.
不過你做的恰恰.
剛剛就是和福音的主旨相反.
因為你不在意那些軟弱的人.
你不在意那些神為他而死的罪人.
所以你應該知道的.
還是不知道.
你應該知道福音是要令一個罪人得以悔改的.
你是知道福音是令一個人.
不再被偶像所箝制.
可以活出新的生命的.
你應該知道這些.
但是當你行使著你自以為是的自由的時候.
你剛剛好就令那個人遠離福音.
所以你應該知道的.
你還是沒有知道.
所以保羅在第二節說的是.
這一類人這種情況.
於是到了第三類人.
你就會說.
他們也要成長的.
難道我一輩子遷就著他們嗎?.
十三節.
「所以如果食物使我的弟兄跌倒」.
在說第二類的那些.
知道的但不太相信的那些.
「我就永遠不吃肉.
免得使我的弟兄跌倒」.

$^{841}$保羅說這句話.
現在不是在說他立了願望.
我不再吃肉.
他在說.
如果食肉使我的弟兄跌倒.
我就不吃.
但是換一個情景.
我不想改聖經.
不過我希望你明白.
我們不可以單單將經文變成律法去應用.
我和你換另一個情景.
保羅立刻會轉tie .
轉什麼呢?.
如果有人又是知道福音.
不太相信的.
他們的想法可能是這樣.
這些人沒有出現那麼多教會.
不過如果他們這樣想.
那些不是拜國偶像的肉而已.
你不去吃那些肉.
其實你也是覺得那些肉是有力量的.
不過那些力量是邪惡的.
所以你不能碰.
只有耶穌基督的聖餐.
主的身體,主的血.
才是聖潔的.
所以你不想污染你的身體.
你就只吃聖餐.
不吃那些肉.
如果有人這樣的疑問.
這個是假如的.
如果假如那麼多教會有人這樣的疑問.
保羅你會怎麼說呢?.
為了這一類的弟兄.
我就和他一起吃肉.
讓他知道那些肉是沒有問題的.
沒事的.
吃了不會有癌症的.
你明白我的意思嗎?.
所以這個時候保羅不是說.

$^{881}$將吃肉不吃肉變成律法.
保羅說的是.
你吃不吃肉的原因.
不是你自己來定.
而是由那個信心軟弱的人來定.
你會說是不是這樣的?.
Jason你將這個拉得很高來說.
我們再看看哥林多前書有另一段經文.
也說飲食的.
在十章的二十三到三十三節.
那裡也說的是吃的肉是否可以吃.
那些肉有少許不同.
這次的肉是擺在市場的.
不是在神廟裡面.
就是你在烏爾斯湖買回來的那些.
或者是擺在你面前的.
保羅說的.
如果有些未信主的朋友想請你去他家吃飯.
你肯去的話.
讓你吃什麼就吃什麼.
接下來他說.
如果那時候你買個肉回家.
很可憐的.
你買肉回家之後.
這個肉祭祖先的.
保羅說你做什麼?.
為了他說那句話.
他又未信主.
你怎麼處理這塊肉?.
扔掉它.
是和牛來的.
這樣就扔掉它?.
人家說不行的.
你知道嗎?.
我現在.
待會我們可能有些茶點時間.
大家談一談.
我們基督徒可不可以吃那些哈拉的食物?.
清真食物.
可以吃嗎?.

$^{921}$交給傳道去處理.
我不管了.
我們一樣在環境裡面.
可以吃嗎?.
你吃又怎樣?.
不吃又怎樣?.
我希望你明白.
保羅說的又不是停留在食肉上.
保羅接著在第十章33節說.
我在一切事情都討所有人喜悅.
包括未信主的人.
包括初信不穩固的人.
為了什麼?.
為了他們得救.
所以保羅已經不是停留在吃不吃肉.
不是停留在你改不改生活習慣.
而是說.
我做的這些事情.
是堅固詛咒著人.
還是將人從福音那裡推開去?.
所以我就將這次的題目定為.
是因為基督的愛去詛咒人.
因為基督的愛不代表沒有知識.
保羅已經說了.
我們有的知識是什麼?.
神只有一位.
基督只有一位.
基督的救贖令我們不再一樣.
這是我們的知識.
但是單有這個知識.
不足以幫助軟弱的人生命成長.
單有這個知識.
不足以令我太太有足夠的安全感.
相信我是愛她的.
我還可以寫希伯來文.
希涅文的愛的不同的寫法.
So what?.
如果這個時候你看完這個之後.
你再看哥林多前書十三章那裡.
愛的篇章那裡.

$^{961}$我們很多人結婚用的那個.
愛是恆久恩來.
愛是恩慈.
愛什麼什麼的時候.
但是我鼓勵你.
那段經文.
那個其實嚴格來說.
在上下文不是在說婚姻的.
不過你可以用在婚姻.
你看前面看一點.
十三章的一到四節說.
倘若我有先知廣道的能力.
沒有愛.
倘若我可以說多少多少謊言.
沒有愛.
倘若我可以供客企室幾身.
叫身服我.
但是沒有愛.
倘若我將一切都獻給.
送給窮人.
但是沒有愛.
保羅已經在說.
如果沒有愛.
即使我做這麼多事.
我懂這麼多事.
都是沒有意義的.
當然保羅說的這個愛.
不是男女之愛.
保羅說的這個愛.
是耶穌基督對我們的拯救.
是這個愛.
如果是這樣.
我就開始在想.
我們現在教會的生態是怎樣.
或應該是怎樣.
保羅說過.
如果知識沒有愛.
是叫人吹脹的.
即是你們應該知道的.
香港人應該明白什麼叫做吹脹.

$^{1001}$即是一個氣球.
體積很大.
不過重量很輕.
懷疑不實.
我們是有知道一些知識的.
這些知識是沒有錯的.
一位神一位主.
但是如果這個知識沒有愛.
就會叛倒弟兄.
成為響的撥.
那些東西.
乾巴巴沒有用的.
如果有知識沒有愛的話.
還要是干犯基督.
自己因為福音的緣故遠離福音.
還搞到其他人遠離福音.
保羅在說這件事.
於是保羅開出一個藥方.
即是在哥林多裡面.
保羅說你當效法我.
正如我效法基督.
什麼意思?.
什麼叫效法基督?.
你效法基督.
做五餅二魚?.
你效法基督.
在風浪裡行走?.
你效法基督.
你去到Darling Harbour.
舉起你手中的聖經.
海水在你面前分開?.
你效法基督是什麼東西?.
我完全不反對大家效法基督.
你效法基督是什麼東西?.
有些東西我們效法到.
有些效法不到.
我們總是要效法一些我們做不到的東西.
然後說我們效法不到.
當然啦.
你做不到.

$^{1041}$你效法來做什麼?.
有些你做到的.
你只是不去效法.
效法什麼東西?.
為那個軟弱的放棄自己自由.
為那個未信初信的.
你不要用盡你的自由.
什麼意思?.
在教會裡面.
我都會將不同的人分在不同的位置.
左下最底.
右上最高.
這個甚至乎我不是說要停留在基督徒那裡.
我覺得我們現在這個世代.
無論基督徒也好.
不是基督徒也好.
都是一些無心無力的人.
缺乏意願的人.
什麼人呢?.
躺平.
大家聽過了.
就是那些.
不要緊的.
我幹嘛這麼拼.
我這麼拼都買不起樓的.
是不是?.
我覺得拿救濟挺好的.
有的.
有這樣的想法.
不是他不上進.
而是說.
他讀完書.
去大學畢業.
出來找工作.
找不到工作.
因為AI取代了他的位置.
那他就去.
或者現在他還想去麥當勞.
不過可能過多兩三年.
麥當勞全部機器人送餐了.

$^{1081}$他覺得他很難.
他很辛苦.
千萬不要說他們沒用.
千萬不要說他們不上進.
只不過處境完全不同了.
那他們活著的意義在哪裡呢?.
有的.
不過要幫他一起找出來.
第二類就是有心無力.
他有意願的.
他想成長.
他想過得好一點.
不過他覺得他能力不夠.
這個就是保羅提到的那種.
初信那種.
有一張提示.
初信不過不夠疏力那種.
第三種是有心有力.
他們那些已經搞定了.
你可以說這種可能是保羅說的第一種人.
他們有知識的.
也在用這些知識.
過得不錯的.
但是保羅繼續挑戰他們說.
等一下.
如果你是有心有力.
但是你不可以將你的有心有力.
轉成一個有心出力.
什麼叫做有心出力?.
就是說不僅僅你是願意跟從耶穌.
不僅僅你是活出生命的使命.
你還可以幫那些有心無力的人.
一起找到他的潛力.
一起使他們的生命開始改變.
用我們教會的教護.
或者很傳統的角度.
Discipleship.
門徒訓練.
就是說.
保羅說的.

$^{1121}$你千萬不要以為你可以自給自足.
有了知識就行了.
保羅說如果你不去關心.
造就主裡面的弟兄姊妹的話.
你不幫助其他人在信仰中成長的時候.
你還沒有準備好.
甚至你還沒有完全知道福音是什麼.
你自以為知道而已.
其實還不知道.
我就稱之為有心給力的人.
再向下一個階段就是.
給心給力.
你不僅僅是幫助那些有心無力的人.
你還可以影響那些沒有心無力的人.
第十章保羅提到那些人.
他們去市場買肉.
他們在飯桌上和你一起吃飯.
那一類人.
那一類就是沒有心無力的人.
因為他們完全不在意.
你是否基督徒.
他們只希望每天可以吃喝玩樂.
我們仍然要向那些人去接觸他們.
去和他們傳福音.
保羅其實在說.
整個教會形態.
整個健康的教會.
或者合神心的教會.
應該是我們不斷地向下一個階段移動.
我不是說整個教會都要盡心盡力.
有心有力給心給力的人.
不是在說這個.
是在說動態的移動.
是在說可能一開始你沒有心無力.
不過你給自己一個月時間.
一年的時間變成有心無力.
你是本來有心無力的.
你給自己一個月時間.
一年的時間變得有心有力.
是一級一級地上去.

$^{1161}$這個叫生命.
這個叫做福音改變的生命.
看看怎麼做就交給傳道了.
下個節目你們的教會你們的策略.
不過方向就有一個方向了.
因為我們是出於這樣的緣故去造就人的.
我們不是出於自己很厲害很厲害.
我們完全是出於耶穌基督的愛.
因為耶穌基督是為罪人而死的.
耶穌基督不是為那些讀神學好的人而死.
耶穌基督不是為那些背金句很厲害的人而死.
耶穌基督是為那些軟弱的人而死.
這個就是福音.
接著你會再追問我.
那Jason啊.
那你怎麼回答司母呢?.
你們有沒有問過.
你們有沒有辦法給我.
司母在說.
我嫁給的是你.
我不是嫁給耶穌.
你縱然說很多偏倒.
你縱然說很多愛的研究.
你感覺不到啊.
我做了兩件事.
一件事是表個態度.
另一件事呢.
就要成為日常生活了.
首先表個態度.
你真的做不到的嘛.
但剛剛不是說了.
生命的重點不是你的狀態在哪裡.
生命的重點是你怎樣移動.
你怎樣進入下一個階段.
不停地進入下一個階段.
我和我太太說.
第一點.
我真的做不到.
第一點是我坦誠愛做人.
第二點.

$^{1201}$我不放棄 我不躺平.
我努力向那個24小時7天.
不斷地去追求去逼近.
你有什麼特別喜歡我做的.
你告訴我.
我去做.
有什麼你不想我去做的.
即使我很喜歡.
我也不去做.
你立個態 你表個態.
第二件事.
我就.
每一天.
當然你們家裡有自己的特點.
不要照樣生搬硬套.
不過你可以試一個星期.
你已經很厲害了.
每一天.
我用5分鐘時間.
放下我所有的東西.
即使今天要來宣導會講道.
昨晚我都要放下所有的東西.
花5分鐘.
抱著我太太.
什麼都不做.
是什麼都不做.
什麼都不想.
看著她.
只看著她.
5分鐘.
好不好 姐妹.
你給點回應.
如果你老公今晚不會做.
你給點回應.
總之我太太.
是很瘦的.
我太太很瘦的.
什麼都不做.
真的什麼都不做.
即是我將她視之為一種.

$^{1241}$愛的表達.
加上一種屬靈操練.
什麼意思.
如果一個生勾勾.
實的的人在這裡.
你都沒有辦法專注在他身上.
5分鐘.
你跟我說.
哎呀 精神牧師.
我很愛主的.
主時時刻刻與我同在.
我感覺到的.
我不相信.
一個人在這裡.
你都不能集中5分鐘.
就想著去揭手機.
我怎麼會相信你.
你可以專注在神的話語裡面.
我想這個就是一個造就人.
我想這個剛剛好.
可以讓我的太太.
不單單在知識上知道.
她的老公很愛她.
耶穌很愛她.
而是在實際上.
在一個實實在在的經歷裡面.
她經歷到.
耶穌很愛她.
以至於她的老公可以很愛她.
當然.
這個只是我自己和我太太探索出來.
屬於我們的方式.
你們有你們的方式.
我也很希望你們可以探索得到.
你說:哎呀 不行啊.
牧師 我好像很愚鈍.
我不是很擅長這些.
容許我用植入式廣告.
你可以問傳道.
最近一次.

$^{1281}$加辛夫婦營什麼時候舉辦.
你在裡面肯定令你滿載而歸.
絕不失望.
你們身邊可能有些人去過.
你可以問一下他們的意見.
如果他們有差評的時候.
可以找我退款.
這個可以這麼說的.
我希望你們可以意識到.
我們當下的教會.
其實不再缺乏一些.
很龐大知識.
很豐富的AI.
我們會用AI 但是不夠.
我們教會可能會變得越來越專業.
好的 但是不夠.
我們教會應該有那種愛.
那種捨己的愛.
所謂捨己的愛.
是顧及那些軟弱的人的愛.
是為了他們不準備.
所以我們要來放下我們某些的自由.
我很欣賞你們教會.
那種為將來的群體做的預備.
因為教會在未來十年.
會發展Mockdale.
Rockdale.
會令你這裡更加多的居民.
所以當下你們要搞好這裡.
裝修.
令這個環境令人更加適應.
令那些新進來的居民.
更容易融入這個環境.
我很欣賞你們.
但是在這個時候.
我不是你們的牧者.
所以我說話可以硬一點.
牧者也很愛你們.
很疼愛你們.
我就直接繞道來了.

$^{1321}$在你們裝修的時候.
肯定有各式各樣不同的意見.
有人喜歡計劃A.
有人喜歡計劃C.
有些人覺得你們的計劃都不好.
我這個最好.
在出現這個時候.
我鼓勵大家.
你們的計劃都是好的.
但是記住你們的計劃都是知識.
請你用愛將你的知識表達出去.
請你用愛稍稍放下你的知識.
如果大家都可以知道什麼是知識.
同時又放下自己的知識.
外面的人就知道.
這裡是耶穌基督的教會.
他們真的不是因為這個建築物有多美麗而來.
當然美麗一點.
是比較喜歡來的.
不過如果他們知道.
在這裡大家是有商有量的.
在這裡原來大家是互相體諒遷就的.
如果在這裡可以互相聆聽的.
在這裡我們可以放下自己.
令其他人得到祝福的話.
我想這個才是教會真正的吸引力.
這個才是福音的本質.
願神祝福你們.
在你們教會發展的時候.
成長的時候.
你們看到的不是自己的靈明的成長.
而是看到在外面.
很多人未聽過福音.
未經歷過福音.
我們身邊有很多弟兄姊妹.
仍然有他們的掙扎有他們的軟弱.
可能他們的立場和看法跟你不一樣.
但是為他們禱告.
為他們的緣故.
放下你所堅持的某一些.

$^{1361}$不是必要真理的事.
這個就是福音.
耶穌是為他們而死的.
當然也是為你而死的.
耶穌是為我們而死的.
我們一起禱告.
天父我們來到您面前的時候.
我們感恩.
在我們最不配的時候.
主你揀選了我們.
在我們因為不安全.
我們要堅持自己的某些原則的時候.
主你首先是為我們寫明.
是你先在我們面前作出榜樣.
主你成就你的國的方式.
不是成為一個沒有霸權.
沒有能力.
而是你用最卑微的方式.
或者用人間覺得最恥辱的方式.
為我們寫明為我們死.
而你的復活也給我們盼望.
知道你是萬有的那一位.
知道重點不是我們贏了什麼.
重點是你為我們寫了什麼.
主要也是在感動裡面.
你也感動到我.
在我們去處理很多人和事的時候.
去謹記一句.
當我正在堅持我是否正確的時候.
主你提醒我.
幫助我.
再問一句.
如果我輸了這個弟兄.
如果我輸了這個姐妹.
贏了這個世界.
那又如何呢?.
奉耶穌的名.
阿們.
\newpage



\section{使徒行傳 20:13-38}
\label{sec:032213}
\textbf{基督裡合而為一}~~~~謝友德牧師
\newline
\newline
連結: \href{http://www.church.com.hk/acms/content.asp?site=cdc\&op=show\&type=product\&code=032213\&layout=sermon}{\texttt{032213}} ~~~~ 語音日期: 2026-05-31
\newline
\newline
\hyperref[sec:032212]{\small{< < < PREV SERMON < < <}}
~
\hyperref[sec:index_chronic]{\small{[返順時目]}}
~
\hyperref[sec:032214]{\small{> > > NEXT SERMON > > >}}
\newline
\newline
使徒行傳 20:13-38
\newline
\begin{longtable}{cl}
\hline
\hline
章節 & 經文 (和合本修訂版)\\
\hline
20:13 & \begin{tabularx}{0.7\textwidth}{X} 我們先上船，起航往亞朔去，想要在那裡接保羅；因為他是這樣安排的，他自己本來打算要走陸路。 \end{tabularx} \\ \\ \relax
20:14 & \begin{tabularx}{0.7\textwidth}{X} 他既在亞朔與我們相會，我們就接他上船，來到米推利尼。 \end{tabularx} \\ \\ \relax
20:15 & \begin{tabularx}{0.7\textwidth}{X} 我們從那裡開船，第二天到了基阿的對岸；再下一天，在撒摩靠岸，又過了一天，到了米利都。 \end{tabularx} \\ \\ \relax
20:16 & \begin{tabularx}{0.7\textwidth}{X} 因為保羅早已決定要越過以弗所，免得在亞細亞耽延，他急忙前行，假如可能的話，在五旬節前能趕到耶路撒冷。 \end{tabularx} \\ \\ \relax
20:17 & \begin{tabularx}{0.7\textwidth}{X} 保羅從米利都打發人往以弗所去，請教會的長老來。 \end{tabularx} \\ \\ \relax
20:18 & \begin{tabularx}{0.7\textwidth}{X} 他們來了，保羅對他們說：「你們自己知道，自從我到亞細亞的第一天，我怎樣跟你們相處， \end{tabularx} \\ \\ \relax
20:19 & \begin{tabularx}{0.7\textwidth}{X} 怎樣凡事謙卑，以眼淚服侍主，又因猶太人的謀害經歷試煉。 \end{tabularx} \\ \\ \relax
20:20 & \begin{tabularx}{0.7\textwidth}{X} 你們也知道，凡對你們有益的，我沒有一樣隱瞞不說的，或在公眾面前，或在每一個人的家裡，我都教導你們， \end{tabularx} \\ \\ \relax
20:21 & \begin{tabularx}{0.7\textwidth}{X} 不論猶太人和希臘人，我都已證明他們當在神面前悔改，信靠我們的主耶穌。 \end{tabularx} \\ \\ \relax
20:22 & \begin{tabularx}{0.7\textwidth}{X} 現在我被聖靈催迫要往耶路撒冷去，雖然不知道在那裡會遭遇甚麼事， \end{tabularx} \\ \\ \relax
20:23 & \begin{tabularx}{0.7\textwidth}{X} 但知道聖靈在各城裡向我指證，說有捆鎖與患難等著我。 \end{tabularx} \\ \\ \relax
20:24 & \begin{tabularx}{0.7\textwidth}{X} 我卻不以性命為念，只要走完我的路程，完成我從主耶穌所領受的職分，為神恩典的福音作見證。 \end{tabularx} \\ \\ \relax
20:25 & \begin{tabularx}{0.7\textwidth}{X} 「我素常在你們中間到處傳講神的國；現在我知道，你們眾人以後不會再見到我的面了。 \end{tabularx} \\ \\ \relax
20:26 & \begin{tabularx}{0.7\textwidth}{X} 所以我今日向你們作證，你們中間無論何人死亡，罪不在我。 \end{tabularx} \\ \\ \relax
20:27 & \begin{tabularx}{0.7\textwidth}{X} 因為神一切的旨意，我並沒有退縮不傳給你們的。 \end{tabularx} \\ \\ \relax
20:28 & \begin{tabularx}{0.7\textwidth}{X} 聖靈立你們作全群的監督，你們就當為自己謹慎，也為全群謹慎，牧養神的教會，就是他用自己血所買來的。 \end{tabularx} \\ \\ \relax
20:29 & \begin{tabularx}{0.7\textwidth}{X} 我知道，在我離開以後必有兇暴的豺狼進入你們中間，不顧惜羊群。 \end{tabularx} \\ \\ \relax
20:30 & \begin{tabularx}{0.7\textwidth}{X} 就是你們中間也必有人起來，說悖謬的話，要引誘門徒跟從他們。 \end{tabularx} \\ \\ \relax
20:31 & \begin{tabularx}{0.7\textwidth}{X} 所以你們要警醒，記念我三年之久，晝夜不斷地流淚勸戒你們各人。 \end{tabularx} \\ \\ \relax
20:32 & \begin{tabularx}{0.7\textwidth}{X} 現在我把你們交託給神和他恩惠的道；這道能建立你們，使你們和一切成聖的人同得基業。 \end{tabularx} \\ \\ \relax
20:33 & \begin{tabularx}{0.7\textwidth}{X} 我未曾貪圖一個人的金、銀或衣服。 \end{tabularx} \\ \\ \relax
20:34 & \begin{tabularx}{0.7\textwidth}{X} 你們自己知道，我靠兩隻手工作來供給我和同工的需用。 \end{tabularx} \\ \\ \relax
20:35 & \begin{tabularx}{0.7\textwidth}{X} 我凡事給你們作榜樣，叫你們知道應當這樣勞苦，扶助軟弱的人，又當記念主耶穌的話，說：『施比受更為有福。』」 \end{tabularx} \\ \\ \relax
20:36 & \begin{tabularx}{0.7\textwidth}{X} 保羅說完了這些話，就和大家跪下來禱告。 \end{tabularx} \\ \\ \relax
20:37 & \begin{tabularx}{0.7\textwidth}{X} 眾人痛哭，抱著保羅的頸項跟他親吻。 \end{tabularx} \\ \\ \relax
20:38 & \begin{tabularx}{0.7\textwidth}{X} 叫他們最傷心的，就是他說「以後不會再見到我的面」那句話。於是他們送他上船去了。 \end{tabularx} \\ \\
[1ex]
\hline
\hline
\end{longtable}
$^{1}$鄭英天媒早晨.
很多謝我們的主席.
非常之流暢地讀了這麼多的經文.
亦都多謝大家這麼有耐性地聽著他讀.
亦都看著你讀.
大家都會記得我們早幾個禮拜.
一直都是講舌頭行船的經文.
這個是我們在今年教會的主題裡面.
其中一個很主要的部份.
今天是我們這個系列講舌頭行船的最後一次.
下個禮拜六月的時候.
我們就會看已忽所書.
然後七月我們就會看肥立比書.
都是循著同樣教會的主題來一起學習.
我們的主題是在基督裡同心前行.
所以今天這個經文的一部份的主題.
都是講到怎樣同心前行.
我們再次一起祈禱.
多謝我們天上的父.
有你的恩典臨在我們每一個人的身上.
亦都是臨在我們的教會.
因為教會是屬於你的.
是主耶穌所設立的.
在我們所做的一切.
都是依靠你.
都是跟從你的旨意.
都是為了成就你在這個世上.
你的奇妙的救恩的大計.
願你繼續教導我們.
願聖靈在我們每一個人的心中.
在我們的教會裡面.
繼續帶領我們同心前行.
祈禱奉耶穌基督的名.
我本來也有想過.
這段經文這麼長.
不如不要讀完.
在崇拜裡面要讀這麼長.
也是一件困難的事.
但是看回這個經文.
又好像很難把它切開.

$^{41}$尤其是十三節到十六節.
這麼多的地名我們都不認識.
現在的地圖也找不到這些地名.
為什麼要讀這些地名呢?.
以前我也經常說.
舊約很多人名地名.
我們都很容易會想跳過它.
讀得快一點.
讀聖經要讀這麼長的時間.
好像有點困難.
但是這些地名.
就是我們學習的其中一個部分.
而且是一個重要的部分.
所以十三節到十六節.
講這些地名我們都要稍微看一下.
然後後面那一段就講到保羅.
他怎樣去事奉神.
用他自己一生的事奉.
來鼓勵當時以弗所的長老.
所以下個星期開始讀以弗所書.
也有一個很好的連繫在那裡.
我們看十三至十六節的地名之前.
我們就回頭重溫一點.
以致我們可以容易一點去總結主題.
士旁傳一至十二章.
是一個重要的段落.
因為這裡講到福音在猶太人的世界.
由耶路撒冷開始.
然後去到薩瑪利亞.
然後去到外邦人的世界.
這個可以說在士旁傳頭那段的時候.
已經表明了.
尤其是一章第八節那裡.
就是用這樣的進程.
由耶路撒冷去到薩瑪利亞.
去到地極全世界.
這個就是福音的開始.
然後十三至二十八章.
就講到福音去到外邦人的世界.
大家都會記得福音去到亞細亞.

$^{81}$然後去到希臘馬其頓.
然後十六章就講到馬其頓的異象呼召.
之後十九至二十八章就講到保羅抵達羅馬的過程.
最終結束在二十八章三十一節.
我們可以將《仕途行傳》分兩個大的段落來看.
是更加容易知道究竟那裡是記錄了些什麼呢.
看回我們十三至十六節今天讀的經文.
我就特別將那些地名就弄了黃色.
給大家都是容易些去看得到.
我就不再讀了.
但是提出那些黃色的地名.
首先就講亞索.
然後就米推利尼.
然後就基亞.
然後在薩摩.
然後在米利都.
然後去到越過以弗所.
沒有進去.
反而是叫以弗所的長老.
差不多兩天的路程.
去到保羅的地方.
跟他見面.
然後就講他會希望趕在五旬前回到耶路撒冷.
大概看一看這個地圖.
希望大家都看得黃色那些.
特別大家比較熟悉的那些地名.
我們又看回耶路撒冷開始.
耶路撒冷開始福音一路傳播.
尤其是保羅一路經過很多地方.
這裡有些地方現在已經不再存在.
在我們現時的地圖裡.
但是給大家容易些看.
福音由耶路撒冷開始.
一直保羅三次宣教的旅程.
然後我們剛才讀過的那些地名.
有一些我找出來寫了中文.
譬如提到米利都,以弗所,米推尼尼.
特羅亞今天這一段經文沒有講.
特羅亞是一個很特別的地方.
因為有人聽著聽著就睡著了.

$^{121}$跌下樓不單止跌下了椅子,跌下樓下.
要保羅救醒他,死過翻生.
然後就肥勒比大家熟悉的地方.
馬其頓去到羅馬.
看一看這樣的地圖.
尤其是我們回想保羅的三個宣教旅程.
我們可以想像得到究竟是.
我想用這個有多少的困難.
有多少的障礙,有多少的艱辛.
那些艱辛不單止是因為當時沒有飛機.
就算坐船都要坐很久.
我還認識一些年長的弟兄姊妹.
他們60年前從香港來澳洲的時候要坐船.
那些不是大家去遊覽的郵輪.
上到郵輪那裡有很好的設備.
聽說,我從來沒有機會去.
聽說24小時都有娛樂節目.
24小時你隨時想吃東西.
都有些餐廳是開的.
又有各種的國際的美食.
當時就算有船坐都是一些很危險的船.
大家看電影都看過了.
但保羅在那13年的裡面三次的宣教.
由主後46至58.
這13年的裡面.
他要在這麼長途的旅程裡面.
而且不是順利的.
不是享受世界的.
而是屢次被恐嚇被毆打被監禁.
困難重重.
但是他仍然堅毅沒有退縮.
他怎樣堅毅沒有退縮呢.
他和這些已乏所有的長老.
去講他一生的事奉.
去勸勉他們將自己事奉的人生.
再給他們看到.
他怎樣可以成為他們的榜樣.
那裡形容得很詳細.
我們平時很快地跳過這些事蹟.
就算保羅怎樣在監獄裡面.

$^{161}$被人抓了然後又唱歌.
然後又帶領很多人信主.
這些很多時候我們是逐一段逐一段看.
我們很多時候讀聖經都是這樣.
都要逐段逐段來.
不會說一次過.
全卷的士徒行傳一次過讀完.
所以我們平時想保羅怎樣.
由猶太的世界去到外邦人的世界.
而且他是帶著一個很特別的使命.
來做這件事.
這件事提醒給我們知道的就是.
他事奉的信心和能力是怎樣來的.
有時候我們都會想.
要保羅經歷這麼多的困難.
要三次這樣的旅程.
在地中海這一帶的世界裡面.
他要這麼多的艱辛.
才可以將福音帶到外邦人的世界.
他的能力我相信是.
因為他的呼召他的儀仗.
是要成為外邦人的使徒.
而且很重要的就是.
他不只是單單想神要我向外邦人去事奉.
而且在他的心裡面.
一定很強烈的呼召和儀仗就是.
神要猶太人和外邦人在基督裡面合一.
這個在這份所書第二章形容得很詳細.
今天就不會看了 下個月才看.
我們平時就是想.
保羅總之是要將福音.
由耶路撒冷一路傳開.
然後去到地極.
這個是一個很重要的呼召.
很重要對神給他的儀仗的回應.
但是我們比較少去看一下.
在神的救恩計劃裡面.
在神對於整個世界人類的歷史進程裡面.
他要看到的不單只是猶太人要歸主.
不單只是外邦人要歸主.

$^{201}$不只是這兩個好像是不同的群體.
要兩個群體都要信主.
不只是這樣.
我們比較少看的就是.
這兩個群體要在基督裡面合一.
如果我們細心去想的時候.
原來這個神的旨意.
是對於全世界.
對於他建立教會.
對於他要拯救全人類.
是一個很重要的元素.
當然每一個靈魂得救.
都是神非常體重的.
也是他給我們的使命.
就是要萬民都成為他的門徒.
我們很多時候的焦點就是.
總之是萬民.
就是所有的人都要有份.
但是原來在這個裡面.
在基督裡面的合一.
是一個我們時時會忽略了的元素.
在這個合一的元素裡面.
我就更加體會到.
合一對於耶穌基督的教會.
是多麼重要.
每一個靈魂都這麼重要.
每一個靈魂我們都要帶領他們.
明白這個救恩是多麼重要.
但是更加整體來說的時候.
我們有時候一跳就跳到整體.
總之全人類都要信主.
然後在啟示錄的形容裡面.
萬族萬民都聚集在神的面前.
一起去唱歌敬拜他.
我們看到這個大圖畫.
但是這個大圖畫裡面.
原來合一是很重要的.
既然全人類.
這兩個猶太人和外邦人.
這兩個群體.

$^{241}$都要在基督裡面合一的時候.
我們就更加體會到.
就算在我們自己個別的教會裡面.
整個靈福堂這個教會.
甚至是我們教會的三種語言的會眾.
每個會眾裡面.
都要有那種合一的元素.
是多麼重要.
我們時時口是會說要合一的.
看約翰福音第17章.
這個就是耶穌為他的門徒的合一的祈禱.
但是從猶太人和外邦人.
在基督裡面的合一這個角度來看.
甚至連猶太人和外邦人.
都要在基督裡面合一的時候.
我們個別的教會的合一.
是更加顯得看到耶穌他的心意.
然後在教會裡面.
不同的語言的會眾裡面的合一.
也都讓我們看見.
耶穌是多麼體重這件事.
所以普羅他是成為外邦人的仕途.
他要將福音從耶路撒冷經過這麼多的旅途.
要將福音帶到外邦人的世界.
他要經歷那些困難.
被人打被人扔在野外的地方.
各種各樣的困難.
他都要堅毅不屈地將福音傳出去.
不只是我們所想的將福音傳遍地極這麼簡單.
當然這個傳遍地極.
在這句話裡面不可以包含所有福音的內容.
但是在普羅的事奉裡面.
在普羅的宣教裡面.
就更加看到合一的重要.
所以也都不只是看到.
本來福音在猶太的世界經歷重重的困難.
就將福音帶到外邦人的世界.
那些都是很重要.
但是在所有這些裡面.
原來合一的元素是我們有時忽略了.

$^{281}$而且是我們不應該再忽略的一件事.
所以這樣再看的時候.
就去到後面那段.
普羅怎樣去鼓勵以弗所的長老.
當然這一段的說話.
不只是給以弗所的長老去思想去學習.
我們都相信聖經裡面所記載的說話.
都是給我們所有的耶穌基督的門徒.
我們都不可以說因為我又不是長老.
所以就算是全個教會的人都要聽.
但是長老要聽多一點.
執事要聽多一點.
教務要聽多一點.
我們可以有時會這樣想.
但是寫在聖經裡面的說話.
都是對我們很重要的糞便.
所以這一個大段就是.
普羅和以弗所的長老所說的話.
他總結自己的事奉的一生.
來勸勉教會要怎樣事奉他.
在第十八節.
我簡單這樣說.
就是那些以弗所的長老.
來到見普羅的時候.
他總結來說.
我在你們中間事奉的時候.
我的為人是怎樣.
我凡事都謙卑而且流淚.
經歷猶太人的謀害.
經歷了很多的試煉.
但是我仍然將要講的福音.
全部都講給你們聽.
他說或在眾人面前.
或在各人家裡.
我都教導你們.
他已經是將他的全部.
教了所有的門徒.
表明自己怎樣是全然獻上.
全心事奉.
也為了福音的緣故.

$^{321}$對所有事奉的人.
就是所有他所事奉的人.
都全然付出毫無保留.
將他從神那裡領受的所有的東西.
教導了他所事奉的人.
然後到了二十一節.
他再強調一樣東西.
又對猶太人和希臘人.
證明當向神悔改.
順靠我主耶穌基督.
他這樣的寫法特別再提.
我是對猶太人和希臘人.
都證明我們需要向神悔改.
要順靠主耶穌基督.
他不只是用那些.
很包含所有的向你們.
向全世界的人.
他要特別再講猶太人和希臘人.
猶太人和非猶太人.
就是我們所謂的外邦人.
然後二十四節他最後說.
我卻不以性命為念.
也不看為寶貴.
只要行完我的路程.
他講到這種道別.
似乎已經是講到永別.
因為他知道他之後不會再去以弗所.
不會再見到這一班的長老.
所以他可以叫做表明心跡.
完全這樣.
他的事奉的心智的態度.
這樣來表達出來.
這也是對我們.
再次是一個非常好的教導提醒.
就算我們不是保羅.
就算我們不是特別領受了.
神給我們怎樣向外邦人去事奉.
但是在現在這個時代.
甚至我們要將福音帶回去給猶太人.
當時福音從猶太的地方開始.

$^{361}$由耶路撒冷開始.
去到地極.
是經過了這麼多保羅的艱辛.
然後向外邦人的世界傳福音.
但是現在這個時代.
我們除了對這一班非猶太人傳福音之外.
甚至要將福音帶回給猶太人.
雖然有很多猶太人不承認這件事.
但是我們從聖經裡面看到.
我們是要將福音帶回給猶太人.
以至耶穌基督所看到的.
所渴望見到的猶太人和外邦人.
在耶穌基督裡面都成為合一.
一部份我們教會的儀仗.
當我們預備要搬進Lincoln 518的地址.
這個新堂裡面.
之前我們都提過.
神比靈福堂的儀仗裡面.
其中一部份就是.
我們要將福音帶回去給澳洲的白人.
我們所謂的西人.
在之前的禹堂的教會.
甚至在去年我們講道的時候.
都已經是重複地說.
在澳洲的白人.
他們在十九世紀的時候.
已經是向中國人傳福音.
當時在中國南邊的地方.
很多的當時所謂的賣豬仔.
當然有些是正正式式去做工.
不是每個人都被人賣豬仔的.
但有很多人去東南亞的地方.
也都原來有很大部份的.
當時在澳洲的人.
是來掘金的人.
當時的澳洲白人.
向這些來掘金的中國人.
傳福音.
到了今天.
有些像猶太人那樣.

$^{401}$福音是從那裡開始.
但是現在他們很多的猶太人.
失去了這個福音.
澳洲的白人.
福音可以說是.
向中國人來掘金的時候.
他們向中國人傳福音.
但是現在很普遍來講.
不是全部澳洲人都不信主.
也不是全部澳洲人的教會.
都不再信福音.
但是我們看到的需要就是.
我們要將福音帶回去給澳洲的白人.
以致我們更加看到.
神給我們的猶太人和外邦人.
都要在基督裡面合一.
那一個耶穌基督的異象.
我們不詳細講了.
提到29至38節那裡.
保羅提醒以弗所的長老.
也都提醒我們所有的信徒.
將要面對的挑戰和攻擊.
當日在保羅的世界裡面.
有羅馬政府的逼迫.
有猶太人的逼迫.
現在我們所身處的世界.
我們所面對的挑戰和攻擊.
那個種類是很不同.
但是大體來講原則都是一樣.
我們要知道我們面對什麼挑戰和攻擊.
我們要預備好自己.
在最後35節保羅這樣講.
他說我凡事給你們作榜樣.
所以這段經文給我們的反思就是.
究竟我們行事為人.
是不是與蒙召的恩相稱.
我們在以弗所書再次會出現.
與這裡的詞句是互相的關聯.
究竟我們的行事為人.
是不是與蒙召的恩相稱.

$^{441}$以至我們能夠榮耀天父和激勵別人.
這個不單止是因為保羅.
他有個特別的使命.
要做外邦人的使徒.
他凡事都給你們作榜樣.
而是我們所有的信徒.
都有領受同樣的一個使命.
這個可以說是我們人生的使命.
就算我們不是保羅.
就算大家坐在這裡.
今天不是一定要去做教務.
不是一定要做宣教士.
但是那個行事為人.
要與蒙召的恩相稱.
是所有的信徒都一樣的.
在任何的事奉之先.
我們自己要有這種心智.
就算你說神有沒有呼召我去做傳道人.
有沒有呼召我去做宣教士.
所以差不多就可以了.
差一點也無所謂了.
但是這個神給我們每一個人的使命.
都是一樣的.
然後我的事奉是不是充滿.
或者出自愛心,謙卑和捨己.
這個也是保羅給我們的勸勉.
給我們所有的人.
你的事奉.
不只是說你會參加師班.
不只是說你做團契的職員.
或者小組的組長.
不只是說你有沒有參與崇拜裡面.
做主席,做師事.
事奉是什麼呢?.
很多時候我們都強調.
事奉所謂有廣義的意思.
也都有所謂狹義的意思.
比較叫做狹義的意思.
就是剛才我說的那些事奉.
或者做一些事情的崗位.

$^{481}$那個工作的內容.
但是廣義的事奉.
我們都提過.
就算你崇拜的時候.
不只是參與會眾一起唱歌.
那個唱歌為什麼要唱歌呢?.
那個唱歌不是說我預備聽道.
所以要唱歌.
絕對不是這件事.
我們唱歌的時候.
是一個事奉.
是一個參與去讚美神.
去重揚我們的神的偉大.
祂的權能.
祂在歷世歷代的救恩的工作.
我們去感恩.
我們去將這些唱出來.
去讚美神.
所以每一位.
每一位在座的.
你去參加一個崇拜的時候.
你的唱歌.
你的祈禱.
這些就是你的事奉.
所以我們完全不可以輕看.
我們參與這個敬拜的事奉.
所以絕對不是只是在台上.
帶領我們唱歌的.
這些敬拜小組.
或者獻詩的我們的師班.
他們就是在做這個敬拜的事奉.
如果我們要用在做這個事奉的時候.
我們應該是包括我們所有的人.
就算你自己在家裡讀聖經祈禱.
你會唱詩歌.
這個就是你的事奉.
所以不是可以說.
我都沒有什麼事奉.
你是不可能沒有事奉的.
所以這些的提醒.

$^{521}$是保羅給我們所有的人.
他怎樣表明他的心跡.
講出他事奉的時候.
是充滿愛心.
出自愛心.
所以他到結尾的時候.
保羅就在那裡記錄.
他跟長老一起祈禱.
他跟長老擁抱.
他跟長老痛哭道別.
這個就是保羅的事奉.
這個就是保羅事奉的心態.
這個就是保羅事奉的果子.
當然他當日在依弗所那裡.
是揀選這些長老一路去教導他們.
培育他們怎樣做一個教會的.
牧養的工作.
但是同時.
看到他跟這些依弗所來的長老.
是有一種多麼深厚的感情.
他叫這些依弗所的長老.
你來找我吧.
我沒空去找你們.
你們就馬上來找他.
找他的時候.
在這個道別裡面.
很深刻的那種感情流露.
所以在教會裡面.
我們不只是要做這些工作.
在教會裡面.
不是完成一個事奉的崗位的要求.
在這個教會裡面.
我們不只是看著前面的目標.
五一八我們要好好計劃.
我們要好好使用.
我們要做完裝修.
我們要籌到足夠的款項來結帳.
收到這些單的時候.
我們的財政要在銀行裡面.
看一下我們是不是夠錢給.

$^{561}$這些是教會的工作.
但是在這些所有的工作裡面.
不可以忽略的就是.
保羅所說的這種.
他的事奉.
他成為所有人的榜樣.
他的事奉是出於愛心.
充滿愛心.
他的謙卑和寫記.
我們今天不可以再詳細說了.
大家一定聽過很多.
讀了不少.
最後的一個反思.
猶太人和外邦人在基督裡面的合一.
對於我們現今的教會.
究竟有什麼啟迪.
有什麼重要.
剛才我所說.
耶穌基督對於外邦人和猶太人的合一.
對神的心來說.
是多麼重要.
既然外邦人和猶太人的合一是重要.
我們在教會裡面.
更加要看到.
是非常非常之重要.
願意保羅他的陳述.
他自己事奉的人生.
陳述他自己事奉的態度.
和他怎樣經歷了這麼多的逼迫.
這麼多的困難.
他仍然沒有放棄.
仍然是很堅忍地要完成神所給他的使命.
我們也要看回.
耶穌基督同樣給我們每一個人.
每一個人.
都有這樣的人生的使命.
我們一起祈禱.
我們再次多謝天上的父.
你愛我們.
創造我們的生命.

$^{601}$給我們的生命可以得著豐盛的救恩.
各樣的恩典.
你都賜給我們.
從來沒有吝嗇.
從來沒有藏起來.
隱藏不給我們.
所以主我們多謝你.
在我們這樣行走這條天路旅途的時候.
求你感動我們.
我們領受了你的恩典.
不只是留給自己.
不只是留給自己的家人.
而是將你的恩典.
向全世界的人去宣告.
我們要不單止以口宣告.
我們要以我們的生命去宣告.
去見證.
主我們又多謝你.
讓我們透過保羅的事奉.
見到合一.
對於我們的主耶穌來說.
是多麼的重要.
在聖經很多不同的地方.
讓我們重覆重覆的見到.
合一的重要.
但很多時候這個合一.
好像有一種抽象的感覺.
主求聖靈幫助我們.
我們既然看見主耶穌.
你甚至對外邦人和猶太人的合一.
是這麼重視的時候.
主我們作為一個教會.
應該怎樣回應你的呼召.
怎樣去實踐你的異象.
求神你在我們教會裡.
特別給我們有恩典.
特別是我們見到前面.
我們還有很多的挑戰.
我們還有很多未完成的使命.
願主你在我們心裡繼續感動.

$^{641}$祈求禱告奉靠救主耶穌基督的命.
阿們.
\newpage



\section{約書亞記 24:24}
\label{sec:032214}
\textbf{聖經人聖經事:我們必事奉耶和華}~~~~黃克勤牧師博士
\newline
\newline
連結: \href{http://www.church.com.hk/acms/content.asp?site=cdc\&op=show\&type=product\&code=032214\&layout=sermon}{\texttt{032214}} ~~~~ 語音日期: 2026-06-29
\newline
\newline
\hyperref[sec:032213]{\small{< < < PREV SERMON < < <}}
~
\hyperref[sec:index_chronic]{\small{[返順時目]}}
~
\hyperref[sec:032215]{\small{> > > NEXT SERMON > > >}}
\newline
\newline
約書亞記 24:24
\newline
\begin{longtable}{cl}
\hline
\hline
章節 & 經文 (和合本修訂版)\\
\hline
24:24 & \begin{tabularx}{0.7\textwidth}{X} 百姓對約書亞說：「我們必事奉耶和華---我們的神，聽從他的話。」 \end{tabularx} \\ \\
[1ex]
\hline
\hline
\end{longtable}
$^{1}$各位電影姐妹大家好.
一起讀聖經.
今次可以說是約書亞記的最後一章.
經過一段時間,講聖經人物,約書亞,摩西死了然後到約書亞.
約書亞帶領以色列人入迦納,一直打,然後得到整個應許之地.
迦納地.
分了地.
屬靈上又設立了.
祭司.
在士劍安放好藥櫃.
等等.
來到一個地方,約書亞記的結束.
結束的時候,約書亞已經很老,年紀老了,就把他們叫來.
語重心長,講了很多東西.
這裡就是約書亞記裡面的人.
約書亞,士就是.
來到這裡,他們最後立志.
事奉耶和華.
約書亞記二十四章.
二十四節.
二十四二十四,如果這次聖經.
你想要你.
自己立志都不錯的,約書亞記二十四二十四.
百姓回答約書亞說.
我們必事奉耶和華.
我們的上帝.
聽從他的話.
當日,約書亞就與百姓立約.
在士劍為他們訂立律例典章.
留意了.
是我要講的第一點.
就是我們必事奉耶和華.
我們的上帝.
聽從他的話.
在約書亞記尾,約書亞已經很老了,差不多的時候.
最後一次跟我們講.
語重心長的說話.
我和我家必定事奉耶和華.
你們可以選擇別的,不過你要想一下.
耶和華神就是帶你們離開埃及那個,不止.

$^{41}$他從亞伯拉罕一直數,數到他們在埃及為奴.
然後摩西帶他們出來,經過曠野.
然後入迦南,戰勝那麼多個王,到現在.
其實是神透過約書亞.
容易識別百姓,他自己是百姓講的.
今天你們就選擇.
敬拜什麼神.
大河那邊,右法拉底河那邊的神.
還是埃及的神金牛.
還是其他.
百姓就說,我們必事奉耶和華.
是我們的上帝,耶和華是我們的上帝.
這裡就不同了,耶和華是你的神,不是,是我們的上帝.
而且還聽從他的話.
因此留意.
這裡解釋了,什麼叫事奉.
事奉就是.
聽從他的話.
事奉也包括了敬拜的.
事奉這個字.
是有敬拜的含義,服侍他.
有敬拜,聽話的含義.
我們會事奉他,我們會服侍.
和聽從他的話.
很清楚了.
當日約書亞跟他們立約,說既然如此.
我就在我們現在的首都士堅.
跟你們訂立律例典章.
律例典章是什麼呢?其實摩西都講過了,不用再重新抄寫一次.
其實律例典章.
就是跟他們再重覆講一次.
還有做一些.
文字,或者皮卷.
約書亞將這些話都寫在上帝的律法書上.
寫下了,我們都會.
聽神的話的.
寫在律法書上.
然後又將一大塊石頭納在像樹下.
耶和華的聖所旁邊.
耶和華的聖所就是回眸.

$^{81}$我們未有機會詳細講回眸,或者講回眸的介紹.
聖所指回眸,回眸裡面有.
外院,聖所,至聖所,至聖所裡面有藥櫃.
應該在回眸旁邊有很多樹.
其中大部分都是象樹.
於是就在回眸旁邊一塊大石.
放了一塊大石.
納了一塊大石,是什麼來的.
可能是寫著字,我們都遵守上帝.
我們神耶和華的律法.
眾百姓納石,約書亞語百姓納石.
記下了,我們要事奉耶和華.
約書亞對百姓說,看呀,看呀.
這石頭可以向我們作見證.
這個石頭納在這裡.
是一個見證.
是因為聽見了耶和華所吩咐我們的一切話.
享讓你們,背棄你們的上帝.
這石頭就可以向你們作見證.
納這塊石.
納這塊石有什麼好處呢?這塊石應該不只是一塊石.
要挑一些字上去,如果你是約書亞.
你會決定這塊石.
會挑什麼字呢?.
今晚可以在神面前禱告.
主啊,我是約書亞,現在要和百姓.
挑一塊石,紀念說.
我們都一起納字,眾百姓納字.
遵行你的話,這裡寫什麼字呢?.
可能就是.
以色列眾百姓同心納石作為見證.
我們.
會跟從耶和華神.
事奉他.
聽從他的話.
然後約書亞與眾百姓納石.
幾時,可能是這樣寫的.
約書亞說,這塊石頭是一個見證.
如果將來你背棄上帝的話.
這塊石頭就會在控告你.

$^{121}$背棄上帝,你看這塊石頭.
你就會知道,我們失守了,我們沒有跟隨這塊石.
我們當日所納的願,我們沒有跟隨耶和華.
而且如果我跟隨的時候,看到就說不對了,我們一路歷世歷代都是跟隨的.
這塊石頭當我們離棄了神的時候.
不跟隨他,他就會成為一個.
控告.
與警惕.
都好的.
於是約書亞打發百姓,好,到我們今天這個儀式到此為止.
大家各歸自己的地業.
又去到那裡聚集完.
大家回去了.
第二點跟大家說的就是死了.
這是以後耶和華的僕人.
就是亂的兒子.
約書亞.
正一百一十歲.
就死了.
這就是我第二點要說的,第一點就是我們施封耶和華是一個立志來的.
第二點就是死了.
約書亞.
一百一十歲.
正一百一十歲.
死了.
這個正一百一十歲死了,就是說.
是告訴我們約書亞的生命是在上帝的.
手裡面掌控著.
正一百一十歲.
死了.
其實日子不是你自己說了算的.
就是這天死的,其實都是神的一個.
給他很完整的意思.
以色列人就將他撞在他地業的境內,就是以法連山上的.
庭拿西拉.
在.
加實山的北邊.
撞了他,這個地點庭拿西拉你記得.
不記得不要緊.
就將他撞在那裡.

$^{161}$看著.
約書亞死了.
是第二點.
接著看下去,第三點先不說,先讀了.
約書亞在逝和約書亞死後.
那些知道耶和華為以色列所行.
諸事的長老還在的時候,以色列人.
是事奉耶和華的.
第三點跳一跳去看一下,以色列人從埃及所帶出來.
約瑟的骸骨.
哇.
這個約瑟的骸骨.
多久了.
第二點就是死了,約書亞死了.
然後約瑟有沒有死了?當然死了,現在只剩下骸骨.
約瑟的骸骨應該在棺木裡面都沒什麼剩的,因為.
他已經死了四百多年.
這個棺木裡面.
應該都是約瑟的骸骨.
可能如果以當日埃及的辦法,他是埃及的.
宰相.
可能會祭過他的身體.
變成今天叫做一個木乃伊,這樣帶出來.
不過都四百多年了.
應該都.
很舊了.
死了.
約書亞死了.
約瑟其實都死了四百多年了.
約瑟其實在這裡看到骸骨.
還是我想說的第二點,就是死了.
第二點我想說死亡.
約書亞死了.
約瑟當然都死了四百多年了,在埃及出來可能.
是用埃及的木乃伊方法,把他祭了的古代方法.
約瑟呢.
都葬在事件,約瑟是重要人物來的.
沒有他,沒有以色列人保存著生命直到今天.
葬在事件哪裡呢?就是雅各從前,即是他爸爸.
用一百塊銀紙.

$^{201}$他爸爸時代用一百塊銀紙.
在事件的父親.
哈梅的子孫所買的那塊地.
這塊地應該很值錢,雅各時代都一百塊銀紙.
一百片銀.
所以都是很貴的了.
就做了約瑟子孫的產業.
約瑟葬在那裡.
重要人物都死了.
阿倫的兒子以利亞薩也死了,第二點就是死得人多了.
死了.
以利亞薩是阿倫的兒子,阿倫是摩西的哥哥.
年紀比摩西大,所以他先在河爾山死了.
然後摩西就在尼泊爾山死了.
現在約書亞又死了.
再記起以前一個很古老影響他們的人物約瑟.
他的鞋骨也死了很久.
現在阿倫的兒子以利亞薩也死了.
葬在他兒子菲利克.
得到二百年山地的小山上.
請大家留意其實人的生命無論你當時多厲害.
約書亞的豐功偉績,真是約書亞要說豐功偉績.
很厲害.
但是他都死了.
約瑟當時顯赫一時埃及宰相.
除了皇帝之外沒有最大的,是他最大的.
又離開了世界.
當時的大祭司阿倫又死了,大祭司連兒子以利亞薩也死了.
摩西也死了.
一代一代過去的領袖原來都會死的.
我也在想自己什麼時候死.
真的會死的,死了又怎樣.
死了之後你知道中國人叫人走茶涼.
人走了那杯茶慢慢就凍了,凍了就沒人喝了.
一朝天子又一朝臣了.
下一代怎麼辦呢.
是長江後浪推前浪.
還是說一代不如一代呢?不知道.
死了.
所以如果你是依靠著一個人.

$^{241}$你就發覺不依靠多久.
很快他就死了.
很厲害的一時都死了.
又輸了長命一百一十歲死了.
那個八十歲還很精壯的迦勒.
死了.
你看到他們都是死了.
所以你不要依靠人,要依靠神.
因為領袖都是會死的.
我想和大家看的.
反而是.
第三點.
第二點是說孩冠,第三點這句說話真的很警惕性.
約書亞在世仇和約書亞死後.
約書亞在世.
即是和約書亞在一起的那些.
和約書亞死了他們都沒死的那些人.
那些就是知道耶和華為以色列人所行諸事的.
第三點就是.
那些長老.
那些長老呢.
都叫長老,也不是小的.
那些長老.
和約書亞曾經同一個世代.
見過了,約書亞死後他們又沒死.
他們知道了耶和華為以色列人所行諸事的人.
長老.
留意.
還在的時候,這句真的聽到有點心up .
還在的時候.
以色列人是事奉耶和華的.
那這班長老.
第二點是說那些老人家,全部的領袖都死了.
那這班和他們同世後生的人,還在的時候.
以色列人都是事奉耶和華的,因為他們會叫那些以色列人.
我們要事奉耶和華,看回那塊石,記得嗎.
我們一起納石的.
是啊,我們要事奉耶和華.
但他還在的時候就是事奉耶和華,後面這句.
沒有.

$^{281}$那他們之後,那些長老死了之後,還有沒有事奉耶和華呢.
人走有沒有茶涼呢.
一朝天子,還是不是一朝臣呢,以後是怎樣呢.
接不接觸到呢.
是不是這個傳承到呢,傳下去.
他又承接到呢.
有講傳承,有講承傳的,傳下去他承接到呢.
不知道.
一個問號沒有說到.
其實呢.
如果照看下去.
創出利民生,約書亞記.
之後到士思記.
沒有了這些這麼大的領袖,在這裡全局,整個以色列民眾,十二支派.
都叫得住的人沒有了.
以色列人很多就開始不事奉耶和華了.
這個就是士思記的記錄來的.
就不事奉耶和華.
所以是那班長老還在.
第二點那些老人家,那些英雄,那些首領.
敬虔的人都死了.
剩下的長老在這裡,都還開得住.
不過還在的時候,都事奉緊實的之後.
可能那塊石都不見了.
那你忘記了.
當人不跟隨耶和華的時候,他們自然慢慢就會將那塊石丟了,不如我們拆了它.
拆了它,因為.
礙眼.
不好看.
他們之後.
看下去,接著的士思記.
很可惜,他們就沒有那麼.
那麼著意事奉耶和華,很多時候是拜其他.
當地的神.
好,今天和大家看到這裡.
主要要看的.
就是記得.
我們必事奉耶和華.
各位弟兄姊妹,大家要記得,如果有下一代,你要.
很努力的.

$^{321}$不停的為他們祈禱.
不停的自己作榜樣,帶領他們.
堅持著信仰的原則.
他們慢慢會跟隨你的.
有時覺得你很古老,很僵化,不跟隨.
他們離開上帝,但你一直禱告神,是會帶他們回來的.
一定會的.
這是上帝一定會做的事,而我也見過不少.
在他們少年青年的時候,可能放蕩不羈.
奧古斯丁就是.
很多都是,但後來.
三十多四十歲迴轉.
很多都是這樣的.
多數有父親母親不停為他們祈禱.
你就算不在.
你一定要預先預備好.
可能他們就預備得不夠了,都預備的,要立塊石.
不過呢,可惜以色列人都是.
一個.
很活躍的百姓.
都好像我們這樣的.
記得.
我們必事奉耶和華,記得要事奉耶和華,跟隨他,同心祈禱.
只要我們立志.
重新立志,要跟隨你.
有時候我們立了志忘記了,立了石都忘記了.
不過求聖靈不停地提醒我們要跟隨你.
要事奉你.
因為你是我們的神.
我們跟隨你,聽從你的話.
讓我們有一個樂意跟隨你的心.
有一個喜愛跟隨你的心.
求你的靈吸引我們.
垂聽禱告.
預備好我們的下一代.
祈禱奉耶穌基督名叩.
阿們.
\newpage



\section{馬太福音 28:16-20}
\label{sec:032215}
\textbf{大使命的謎思}~~~~林家棋宣教師
\newline
\newline
連結: \href{http://www.church.com.hk/acms/content.asp?site=cdc\&op=show\&type=product\&code=032215\&layout=sermon}{\texttt{032215}} ~~~~ 語音日期: 2026-05-31
\newline
\newline
\hyperref[sec:032214]{\small{< < < PREV SERMON < < <}}
~
\hyperref[sec:index_chronic]{\small{[返順時目]}}
~
\hyperref[sec:032216]{\small{> > > NEXT SERMON > > >}}
\newline
\newline
馬太福音 28:16-20
\newline
\begin{longtable}{cl}
\hline
\hline
章節 & 經文 (和合本修訂版)\\
\hline
28:16 & \begin{tabularx}{0.7\textwidth}{X} 十一個門徒往加利利去，到了耶穌指定他們去的山上。 \end{tabularx} \\ \\ \relax
28:17 & \begin{tabularx}{0.7\textwidth}{X} 他們見了耶穌就拜他，然而還有人疑惑。 \end{tabularx} \\ \\ \relax
28:18 & \begin{tabularx}{0.7\textwidth}{X} 耶穌進前來，對他們說：「天上地下所有的權柄都賜給我了。 \end{tabularx} \\ \\ \relax
28:19 & \begin{tabularx}{0.7\textwidth}{X} 所以，你們要去，使萬民作我的門徒，奉父、子、聖靈的名給他們施洗， \end{tabularx} \\ \\ \relax
28:20 & \begin{tabularx}{0.7\textwidth}{X} 凡我所吩咐你們的，都教導他們遵守。看哪，我天天與你們同在，直到世代的終結。」 \end{tabularx} \\ \\
[1ex]
\hline
\hline
\end{longtable}
$^{1}$弟兄姊妹平安.
我忘了有這個講題.
為什麼我這麼清楚.
因為那天講到那個是我.
所以當日的經文.
其實都是用今天訊息的經文.
當日的訊息分享.
提到大使命和我們基督徒之間的關係是什麼.
我們是沒辦法去逃避.
今天我再一次和大家去思想大使命這一段的經文.
其實會不會我們有一些.
我們需要再一次被提醒的地方.
關於大使命這段經文.
我們的印象.
第一個印象就是你們要去.
同時我們也當你們要去作為一個主要的命令.
你們要去的這個說話.
我們很自然就會聯繫到.
我就是和人去傳福音.
所以我們一提到大使命.
我們就必然會聯繫到.
一定要和傳福音有關.
我否定.
大使命是和傳福音有關的.
我不會否定的.
不過弟兄姊妹.
如果我們只是集中在你們要去這個行動的說話之中的話.
可能我們理解耶穌基督所吩咐的命令.
可能還有一段距離.
甚至我們可能對於耶穌基督所吩咐的大使命.
我們理解了三分之一.
讓我們一起去思想.
19至20節裡面的經文.
其實我們有沒有什麼忽略了.
邀請大家先去打開今天的經文.
馬太福音28章16至20節這段經文.
給一點時間大家去打開.
馬太福音28章16至20節.
使萬民作我的門徒.
奉父子聖靈的名給他們施洗.

$^{41}$凡我所吩咐你們的.
都教導他們遵守.
看啊.
我天天與你們同在.
直到世代的終結.
在耶穌基督所吩咐門徒的命令之中.
又或者我們去看一下.
第19至20節的上半節裡面的大使命內容.
其實它有幾個的動詞.
第一個是去.
第二個是使人作門徒.
第三個就是吸人施洗.
然後就吩咐教導遵守.
但是在這六個的動詞裡面.
只有一個.
只有一個是命令語的動詞.
就是使人作門徒.
那其他動詞的意思是甚麼呢.
容後大家再去和大家去思想.
正如剛才提到大使命的核心.
其實原來就是要使人作門徒.
亦都是大使命的最終目標.
使人作門徒.
他用了命令語去表達.
其實發揮了一種很迫切的作用.
就是在這個終末的時間.
要實踐這個行動.
不能再等了.
弟兄姊妹.
首先我們要搞清楚一件事.
就是耶穌基督這個命令.
是要使人作門徒.
而不是使人作基督徒.
大家可能聽完這個說話.
一頭問號.
是這個意思嗎.
我們由信耶穌那一刻.
其實我們已經被稱為基督徒.
《使徒行傳》都是這樣記載的.
當然這裡要帶出的就是.

$^{81}$我不是說基督徒和門徒.
是彼此對立.
又或者哪一個的地位高一點.
哪一個的地位低一點.
沒有這樣的意思.
不過耶穌在這個命令裡面.
他呼召人切切實實.
是要呼召人成為他的門徒.
也是基督徒要邁向的目標.
耶穌吩咐11個門徒.
要使人作他的門徒.
而不是就這樣吩咐門徒.
你就這樣出去傳福音吧.
你就這樣出去傳福音.
叫人信耶穌.
信我就可以了.
這個不是耶穌他所吩咐的核心.
但我們要理解什麼叫做使人作門徒.
我們要明白定義是什麼.
很簡單.
我們有模特兒.
我們有參考的模特兒.
我們可以參考回.
耶穌呼召門徒的片段.
無論是彼得,雅各,約翰.
他們都是撇下所有.
跟從耶穌.
做瑞利的馬太.
同樣都是.
當耶穌向他發出.
來跟從我這個邀請.
馬太同樣都是撇下所有.
包括一個可以讓他發財的職業.
他跟了耶穌.
他放下了這個職業.
去跟耶穌.
所以使人作門徒這個命令.
就要讓微信者.
去走上十一個使徒.
所要走的路.

$^{121}$就是要放下所有.
忠心跟隨.
弟兄姊妹.
我們同樣都要留意.
耶穌這個命令的對象.
是要使萬民作他的門徒.
使萬民作門徒是一項挑戰.
同時也是信服.
以及向耶穌忠心的途徑.
為什麼這樣說呢.
首先我們要去看一下.
萬民這個字.
這一組字.
又或者正確一點來說.
這一組字.
在馬太福音裡面.
有其他的地方都有出現過.
就是二十四章九字.
二十四章十三字.
二十五章三十二節.
和我們今天的經文裡面.
除了二十五章三十二節裡面提到.
萬民都要來到耶穌寶座面前.
要被分別出來之外.
我們去大家去揭一揭.
二十四章九到十四節的經文.
我們一起去揭開.
馬太福音二十四章九到十四節的經文.
在前一點的.
我們看一下馬太福音二十四節.
二十四章九到十四節的經文.
心漸漸冷淡.
但堅忍到底的忠必得救.
這個天國福音.
要傳遍天下.
對萬民作見證.
然後終結才來到.
我們看見在末後的日子.
作門徒.
是要面對陷入患難之中.

$^{161}$又要面對被殺害.
更加亦都會因為耶穌的名.
而被萬民憎恨.
然而.
門徒處於一個艱難受苦的局面之中.
在第十三節.
耶穌卻是說要堅忍到底.
甚至乎要對這一班憎恨門徒的萬民.
要作見證直到世界終結.
做門徒不是只是我信耶穌就完.
不是就這樣停在這裡就算了.
做門徒是要面對要被萬民所憎恨.
同時亦都要在這些憎恨門徒的萬民當中去作見證.
而這個不單單只是一次性的見證.
而是在這些艱難局面裡.
仍然要堅忍一路走下去.
一路在這班萬民的面前去作見證.
所以大使命的核心.
除了叫人撇下一切跟隨耶穌.
並且都需要去為著基督的緣故.
要被萬民憎恨.
卻又要在萬民之中.
使他們作基督的門徒.
要實踐大使命的核心.
其實包括了三個的元素.
耶穌用了三個的分詞來去表達.
就是去吸引思洗和教導.
第一個的元素去.
我們不難理解.
就是要使人作耶穌的門徒.
我們要知道.
我們必須要有去這個行動才可以做到.
而去這個分詞.
它裡面的作用其實是.
修飾著使人作門徒這個命令語.
所以弟兄姊妹.
你們要去這個分詞.
不單單只是叫我們去傳福音.
更貼切的意思.
是要使人作耶穌的門徒.

$^{201}$簡單來說.
我不是只是叫人信耶穌.
我們一向誤解.
以為大使命我只是去叫人信耶穌.
叫人信耶穌只不過是第一步.
也是入門級.
我們要做的是怎樣去使人作門徒.
即使我們可能因為基督的緣故.
我們被人憎恨.
我們仍然放膽宣講上主的道.
讓這個使命一代傳一代.
這個才是我們要去的目的和目標.
弟兄姊妹可能你們會覺得.
你們要去不只是叫人信耶穌嗎.
嘩 說得這麼嚴重.
信耶穌而已.
不用去到這麼盡吧.
每次一來到就跟人說.
你扔下所有.
我昨天開玩笑.
跟某某位學生說.
用一個例子.
他們下個禮拜考試.
他們接下來這個禮拜考試.
我說你立即跟耶穌.
你扔下書本.
放下你沒有溫書.
你跟耶穌.
嘩 不用這麼盡吧.
玩到這麼盡.
信耶穌而已.
真的要去到這麼盡嗎.
還要說到.
我要被人憎恨這麼嚴重.
嘩 不用這麼誇張吧.
信耶穌而已.
但是弟兄姊妹.
這一刻.
我可以告訴大家.
要進行大使命真的要這樣.

$^{241}$布弗華裡面提到.
當基督去呼召一個人.
他是呼召他去為他死.
這個死.
可能好像第一代的門徒.
要撇下一切跟隨主.
也可能好像馬丁路德那樣.
離開修道院進入世界.
如果我們從布弗華去理解.
作門徒的意思.
大使命所傳的.
不再只是恨福音.
恨福的恨.
恨福音.
也不只是單單的論述.
信耶穌得永生.
而是去傳一個.
為耶穌死的福音.
這個死不是只靠我們自己去成就.
而是靠著耶穌基督.
和耶穌基督和死.
和復活的恩典.
就好像加拉太書.
二章二十字裡面說到.
現在活著的不再是我.
乃是基督在我裡面活著.
並且如今在肉身活著.
是因信上帝的兒子而活.
他是愛我 為我捨己.
對於弟兄姊妹來說.
我就這樣叫人信耶穌也很難.
更何況你設一個這麼高的門檻.
這麼高的難度.
怎麼跟呢.
當然我們不會馬上向一個.
微信者一開口說.
你要信耶穌.
你要為他死.
我斷固你說完這句話.
肯定會說.

$^{281}$我連耶穌是誰都不知是誰.
你跟我說我要為他死.
但是.
當我們跟對方.
傳福音的時候.
當然我不是說.
我們過去就這樣叫人信耶穌.
是不對的.
我們都需要去.
有少許簡單的入門的方法.
但是當我們去跟對方傳福音的時候.
而當他願意信主了.
他願意信主了.
我們感到歡天喜地是.
很自然不過的事.
但是.
我們不要只是停留在這個地步.
我們不要只是停留對方信主的階段.
大部分教會可能出現了一個普遍現象.
就是當有新朋友又或者微信者進入教會之中.
這位微信者又或者新朋友成為了教會中眾人的焦點.
會受到很熱烈的歡迎和會被接待.
而這位微信者又或者新朋友.
在一段日子之後.
他被聖靈感動歸主.
教會裡面的弟兄姊妹.
很自然地說很感恩.
你終於信主了.
很開心你終於可以信主了.
直到這個果子.
這個果子真的可以收割了.
不過高興過後.
高興過後.
這個歸主的果子.
隨著過去的日子.
連日過去.
好像漸漸被忽略.
當初被重視.
被關心的程度.
好像和他剛剛進來教會.

$^{321}$微信之前.
大不如前.
為什麼我歸主.
我信耶穌之後.
沒有人關心我.
為什麼我剛剛來到微信主的時候.
他是很多人去關心我.
但是現在呢.
為什麼會這樣呢.
這個狀況.
這個果子的狀況.
他是成為了教會裡面.
我們叫做.
屬靈棄嬰.
他信了主之後.
好像被人遺忘了.
好像沒有人關心他.
好像沒有再很熱烈地歡迎他.
弟兄姊妹.
大使命的經文.
核心是要使人作門徒.
而我們第一個元素.
就吩咐我們要去.
要去不是只是.
領人歸主這麼簡單.
而是要使.
那些歸主者.
去作基督的門徒.
我們需要幫助他們.
進入下一個階段.
以及讓他成為.
基督的門徒.
讓他延續耶穌基督.
所吩咐的大使命.
一代傳一代.
我們去看看大使命第二個元素.
就是.
吸人施洗.
這個分詞的作用.
是補充.

$^{361}$使人作門徒這個命令.
每一間教會.
最少都有兩個很重要的禮儀.
其中一個就是洗禮.
縱然每一個的宗派.
每一間的教會.
都有不同的洗禮官.
即使又拿我和趙孫.
做一個例子.
我們的洗禮官是不同的看法.
不同的觀點.
但是我們各自都會.
共同認同一樣東西.
就是洗禮不是.
只是一個禮儀.
更加不是.
信仰的終點.
而是每一個基督徒.
信仰的開始.
一步一步.
這樣走近更似耶穌基督.
的樣式的道路.
為什麼大使命第二個元素.
是吸引人思洗.
而不是就這樣洗人信耶穌.
因為洗禮本身就象徵.
一個跟隨基督的人.
他願意展開.
他作為.
跟隨者的生命.
開始.
並且透過這個禮儀.
向世界公開宣認.
他願意歸入.
基督的死.
以及和他一同.
埋葬.
意思就是.
我們要去幫助這一班.
歸主者去治死老我.

$^{401}$就好像.
即是潘博華所講.
就是基督呼召人.
去為他死.
要透過洗禮.
去重新進入.
三一神.
和三一神聯合的.
新生命的開始.
並且要和這個.
老我去斷絕.
這個就是要吸引.
思洗的原因.
作門徒第三個.
元素就是教導.
這個分詞同樣.
都是補充.
洗人作門徒.
這個命令.
耶穌吩咐.
向萬民作出教導.
那應該教什麼.
我們看回.
二十八章二十節的上半節.
裡面這樣說.
凡我所吩咐你們的.
都教導他們遵守.
如果我們原文.
去直譯的時候.
就是教導他們一切.
我所命令你們的.
耶穌.
很清楚地吩咐.
門徒要向萬民.
教導的內容.
就是要教導他們去遵守.
遵守本身這個字.
就是用了一個不定式.
就是infinitive.
去帶出教導.

$^{441}$教導的目的.
目的就是要他們遵守.
遵守什麼呢.
耶穌吩咐.
門徒要教導萬民.
去遵守一切.
我所命令過門徒的內容.
當中包括.
登山寶訓的教導.
天國的比喻.
群體的倫理生活.
門徒的代價等等.
弟兄姊妹.
我們更加要留意一組的詞語.
就是一切.
又或者和修版本裡面的.
凡所這個字.
一切.
這一組的詞語.
出現了五次.
列在.
這個powerpoint上給大家看到.
幾節經文其實都是.
表達所有的意思.
所有的意思.
就是說耶穌向門徒.
所講過所做過.
所教過的一切.
都要100\%如實.
這樣教導萬民去遵守.
沒有選擇的空間.
沒得選擇.
沒得我喜歡選擇.
我不教這個我不教那個.
我沒得.
我只是教天國的倫理和應許.
但我卻是不教.
不教大家.
不教大家揹起.
自己的十字架.

$^{481}$耶穌的吩咐.
耶穌的命令.
不是這樣去理解.
而是我們.
既要宣講悔改的訊息.
也都要講.
得救的應許.
更加不可以忽略.
作門徒.
要為主捨己的代價.
不單單.
只是去教.
我們更加要讓.
這班歸主者去遵守.
去實踐.
去遵行.
耶穌所有的教導.
兄弟姐妹.
我們明白其實.
大使命對於我們每一個信徒.
都是責無旁貸.
不過我們真的要明白.
真正明白大使命.
真正的意義.
耶穌這個命令.
不是只是叫人信耶穌.
聽清楚.
不是只是叫人信耶穌.
而是叫人去成為.
他的門徒.
所以我們要清楚.
明白大使命不再只是.
停在叫人.
傳福音.
信耶穌.
而是我們還要幫助.
這一班願意歸主的人.
去受洗.
歸入基督並且.
教導他們遵守.

$^{521}$耶穌所有的吩咐.
讓這班歸主者.
一步一步成為.
基督的門徒.
對於弟兄姊妹來說.
不要說做耶穌的門徒.
我連去傳福音.
我都害怕.
不要一下子設那麼高門檻.
不要耍我.
當然.
這個不是.
一步登天的事.
弟兄姊妹這個不是一步登天的事.
亦都是.
剛才所提到.
要洗人成為.
耶穌基督的門徒.
不是一朝一夕的事.
很漫長漫長地去做.
是一個.
信仰的進程.
我們先從.
第一步即是第一個元素.
開始就是要去.
我們明白到去的.
困難可能有部分的弟兄姊妹.
會覺得我很害怕.
我不敢出去.
我很怕被人拒絕.
甚至連我自己身邊.
未信主的親人.
都覺得很害怕.
他們會聽我說.
大孫說.
上主的福音.
從來都不是一個人.
可以做到.
是需要其他人去互相.
效力配搭.

$^{561}$我們看回《士徒行傳》的時候.
保羅起初宣教.
他是和巴拿巴一起出去.
及至去到中段.
保羅和巴拿巴分道揚鑣.
但他們都不是.
各自一個人去.
繼續去履行他們的福音使命.
巴拿巴會帶上馬可.
保羅會帶上.
西拉繼續他們的.
福音工作.
而且耶穌.
差遣門徒的時候.
都是兩個兩個.
出去.
所以大使命第一個元素.
你們要去.
不是只是一個人.
就這樣出去.
需要有同路人.
一同出去.
就算是我們自己.
未信的親人.
我們都可以找我們的隊友幫手.
以我家裡為例.
家裡的人.
未信主.
要我和家裡人.
直接傳福音.
避免不了會吵架.
這是很真實的.
但是.
家裡需要.
一個神隊友.
而這個神隊友就是我太太.
我太太就成為了.
我家裡的福音.
的一個切入點.
她是一個好幫手.

$^{601}$所以其實.
弟兄姊妹.
大使命第一個元素.
你們要去.
其實需要有同伴.
有隊友幫手.
就好像.
一場足球比賽.
中間那個.
是今年.
英超入球.
最多的那個.
他叫克蘭達.
這張照片是幾年前的.
只不過是.
他左邊的隊友離開了.
但是其實克蘭達.
多好球都好.
就算那幾年他還是神射手.
入了很多球都好.
其實他左手邊那個.
沒有他的助攻.
餵球給他.
他幾好球他怎麼入.
幾好球都入不到.
他左手邊那個是迪布尼.
迪布尼是出名.
派牌的球員.
我們俗稱派牌的球員.
沒有他派牌.
很難一個人去闖關.
去入那麼多球.
在.
李恆大使命裡面.
其實我們都需要有.
共同的戰友.
去奮鬥.
有多一個.
多一個.
去一起和我們去宣講.

$^{641}$上主的說話.
有多幾個.
一起和我們去工作.
去共同努力福音的工作.
我們相對.
沒那麼膽怯.
李恆弟兄姊妹.
我們需要有被拒絕的.
勇氣又或者好像.
馬太福音24章.
提到我們要有.
被萬民所憎恨的.
勇氣.
有時我們很害怕被拒絕.
又或者害怕得罪別人.
我們不敢去接觸.
那些未信者又或者我們的家人.
但當.
雖然我們向未信主的家人.
或者其他朋友去講福音.
未必下下.
說我憎死你.
我憎到你入骨.
未必下下去到這個地步.
但我們.
不多不少可能都會.
面對一些的情況.
我們可能會面對被人冷嘲熱諷.
最常聽到.
一句說話就是.
信耶穌.
食到飯嗎?.
你給到飯我食嗎?我信耶穌.
當然.
亦都.
不單止.
有一件事就是.
很多朋友可能都真是.
會表達自己很憎基督教.
很憎基督教.

$^{681}$面對這些的情況我們可能.
真是會很害怕.
但是.
我們要想一下.
其實我們向這一班.
未信者又傳福音.
其實.
其實就是我們.
為著基督的緣故.
去要承受.
並且需要跨過的難關.
但使命不是.
停留在去的階段.
而是還有下一步.
特別是一些剛剛歸主者.
他們仍然需要被關心.
以及去陪伴他們.
讓他們不會感到我信主之後.
我好像被人遺忘了.
成為了一個.
好像屬靈棄嬰.
他們仍然需要有人幫助他們.
去繼續走那個.
漫長的信仰的路.
以致他們能夠.
受洗歸入基督.
當然.
受洗歸入基督不是終結.
而是開始.
要讓到這一班洗禮的.
信徒成為基督的門徒.
需要教導耶穌所.
吩咐的一切.
有時候我們可能.
為了迎合對方的口味.
避免嚇怕對方.
我們可能只是說.
信耶穌得永生.
信耶穌好一好.
這一類的福音.

$^{721}$但我們好像對於.
犯罪要悔改.
要背十字架.
你要為耶穌受苦.
這些的教導.
我們可能淡化了.
又或者我們避重就輕.
如果是這樣的話.
我們沒有辦法讓這班歸主者.
踏上成為基督門徒的路.
我們不能夠.
只教導.
我們信耶穌成功.
那一面.
我們不去說.
信耶穌其實要.
為耶穌受苦.
這個不是一個整全的.
信仰.
就好像之前引用潘博華的說話.
基督呼召人來.
是叫他為他死.
要好像第一代的門徒.
一樣離開家庭.
和工作來跟隨耶穌.
每一天都必須.
和罪和惡去爭戰.
每一天都必須.
遭受試探.
每一天都要重新為基督受苦.
整全的信仰.
和耶穌教導.
就是這樣.
而且是必須.
以至歸主者能夠明白.
他面對的困難.
原來.
就是耶穌所吩咐.
所教導過.
所經歷過的一樣.

$^{761}$這種信仰經歷.
所經歷的真實.
能夠幫助他.
去轉化他們的生命.
我們的角色.
就是要和他們一起去學習.
一起去同行.
一起去分享.
如何去經歷耶穌基督的教導.
弟兄姊妹.
大使命的真諦.
萬民作門徒.
當中需要我們.
去吸引思洗.
以及教導他們.
遵守耶穌一切的命令.
但是在履行.
大使命的任務之中.
不是只有我們自己.
一個人去工作.
重要的就是那位死而復活.
勝過天上.
地下地底下.
一切權勢的耶穌基督.
他每一天.
都和我們同在.
直到世界的終結.
我們不是只是一個人.
單打獨鬥.
所以弟兄姊妹.
我們不要害怕.
我們覺得我們自己沒有能力.
去履行這個使命.
但是.
耶穌和我們一起.
祂差遣我們.
祂讓我們有能力.
去履行這個使命.
今天.
耶穌要問你.

$^{801}$不是問你願不願意.
去傳福音.
而是問你.
你願不願意.
一生作祂的門徒.
並且.
讓其他人.
一樣作祂的門徒.
弟兄姊妹.
你願意嗎.
給一點時間.
安靜大家.
去想想.
耶穌問你.
你願意.
一生作祂的門徒嗎.
你願意.
讓其他人.
一樣作祂的門徒嗎.
讓我們有一點.
安靜的時間.
耶穌基督我們都.
你把你的救恩.
去臨到我們.
你叫我們.
真是得到這份救恩.
你在這裡.
在這裡.
真是得到這份救恩.
令到我們有復活的盼望.
當你.
復活後.
升天前.
向門徒頒布.
大使命的時候.
你不只是叫我們.
就這樣出去.
傳福音那麼簡單.
而是要叫萬民.
作你的門徒.

$^{841}$對我們來說.
我們每一個人都可能很害怕.
有膽怯甚至乎.
可能真是未必.
願意.
去踏出這一步.
但是你說.
天上地下.
一切的權柄都已經賜了給你.
你給了這份力量.
給了這個同在.
和我們一起.
你叫我們不是單靠自己的能力.
去做.
萬民作門徒.
而是你每一天都和我們同在.
每一天都賜給力量.
賜力量給我們.
去使我們有這份勇氣.
去走出去.
去向萬民.
宣告上帝你的福音.
耶穌基督你的福音.
求你給我們有這份勇氣.
去讓一班.
可能真是未信的人.
又或者初信的人.
可以一步一步讓他們成為.
跟隨基督.
作主基督的門徒.
我們想要禱告奉主耶穌基督.
得勝明子祈求.
\newpage



\section{約拿書 1:1-17}
\label{sec:032216}
\textbf{神要找的人}~~~~張淑嫻牧師
\newline
\newline
連結: \href{http://www.church.com.hk/acms/content.asp?site=cdc\&op=show\&type=product\&code=032216\&layout=sermon}{\texttt{032216}} ~~~~ 語音日期: 2026-05-31
\newline
\newline
\hyperref[sec:032215]{\small{< < < PREV SERMON < < <}}
~
\hyperref[sec:index_chronic]{\small{[返順時目]}}
~
\hyperref[sec:032221]{\small{> > > NEXT SERMON > > >}}
\newline
\newline
約拿書 1:1-17
\newline
\begin{longtable}{cl}
\hline
\hline
章節 & 經文 (和合本修訂版)\\
\hline
1:1 & \begin{tabularx}{0.7\textwidth}{X} 耶和華的話臨到亞米太的兒子約拿，說： \end{tabularx} \\ \\ \relax
1:2 & \begin{tabularx}{0.7\textwidth}{X} 「起來，到尼尼微大城去，向其中的居民宣告，因為他們的惡已達到我面前。」 \end{tabularx} \\ \\ \relax
1:3 & \begin{tabularx}{0.7\textwidth}{X} 約拿卻起身，逃往他施去躲避耶和華。他下到約帕，遇見一條船要往他施去。約拿付了船費，就上船，與船上的人同往他施，為要躲避耶和華。 \end{tabularx} \\ \\ \relax
1:4 & \begin{tabularx}{0.7\textwidth}{X} 耶和華在海上颳起大風，海就狂風大作，船幾乎破裂。 \end{tabularx} \\ \\ \relax
1:5 & \begin{tabularx}{0.7\textwidth}{X} 水手都懼怕，各人哀求自己的神明。他們把船上的貨物拋進海裡，為要減輕載重。約拿卻下到艙底，躺臥沉睡。 \end{tabularx} \\ \\ \relax
1:6 & \begin{tabularx}{0.7\textwidth}{X} 船長到他那裡，對他說：「你怎麼還在沉睡呢？起來，求告你的神明，或者神明顧念我們，使我們不致滅亡。」 \end{tabularx} \\ \\ \relax
1:7 & \begin{tabularx}{0.7\textwidth}{X} 船上的人彼此說：「來吧，我們來抽籤，看看這災難臨到我們是因誰的緣故。」於是他們就抽籤，抽出約拿來。 \end{tabularx} \\ \\ \relax
1:8 & \begin{tabularx}{0.7\textwidth}{X} 他們對約拿說：「請你告訴我們，這災難臨到我們是因誰的緣故呢？你做甚麼行業？你從哪裡來？你是哪一國的人？屬哪一族？」 \end{tabularx} \\ \\ \relax
1:9 & \begin{tabularx}{0.7\textwidth}{X} 他說：「我是希伯來人，我敬畏耶和華，天上的神，他創造了滄海和陸地。」 \end{tabularx} \\ \\ \relax
1:10 & \begin{tabularx}{0.7\textwidth}{X} 那些人就大大懼怕，對他說：「你做的是甚麼事呢？」原來他們已經知道他在躲避耶和華，因為他告訴了他們。 \end{tabularx} \\ \\ \relax
1:11 & \begin{tabularx}{0.7\textwidth}{X} 海浪越來越洶湧，他們就問他說：「我們當向你做甚麼，才能使海浪平靜呢？」 \end{tabularx} \\ \\ \relax
1:12 & \begin{tabularx}{0.7\textwidth}{X} 他對他們說：「你們把我抬起來，拋進海裡，海就會平靜了；我知道你們遭遇這大風浪是因我的緣故。」 \end{tabularx} \\ \\ \relax
1:13 & \begin{tabularx}{0.7\textwidth}{X} 然而那些人竭力划槳，想要把船靠回陸地，卻是不能；因風浪越來越大，撲向他們。 \end{tabularx} \\ \\ \relax
1:14 & \begin{tabularx}{0.7\textwidth}{X} 於是他們求告耶和華說：「耶和華啊，求求你不要因這人的性命使我們滅亡，不要使流無辜人血的罪歸給我們；因為你－耶和華隨自己的旨意行事。」 \end{tabularx} \\ \\ \relax
1:15 & \begin{tabularx}{0.7\textwidth}{X} 他們把約拿抬起來，拋進海裡，海的狂浪就平息了。 \end{tabularx} \\ \\ \relax
1:16 & \begin{tabularx}{0.7\textwidth}{X} 那些人就大大懼怕耶和華，向耶和華獻祭許願。 \end{tabularx} \\ \\ \relax
1:17 & \begin{tabularx}{0.7\textwidth}{X} 耶和華安排一條大魚吞下約拿，約拿在魚腹中三日三夜。 \end{tabularx} \\ \\
[1ex]
\hline
\hline
\end{longtable}
$^{1}$弟兄姊妹平安.
很高興今天又來到大家的中間.
我想問一下今天我們的題目是「神要找的人」.
和尋找有關.
不知道大家有沒有找過東西.
有沒有試過尋找東西.
經常都忘記東西放在哪裡.
通常大家會找些什麼東西呢.
有找過什麼東西嗎?.
有沒有試過尋找人呢?.
有沒有試過尋找東西呢?.
通常現在不見的東西是什麼呢?.
眼鏡?手機?.
手機應該很少不見.
因為手機很重要.
但有時很有趣.
手機不見其實只是在家裡.
即是突然在客廳.
「咦?去哪了?」.
原來可能在電視機上.
或者在櫃底之類.
有時夫婦也很有趣.
很多時候不知道為什麼.
大家有沒有經歷過.
或者可能家人.
譬如我和我先生說.
可能他問我.
「那些東西在哪裡?」.
我說「在那裡」.
打開櫃就看見了.
然後打開櫃「沒有啊」.
「看不見啊」.
「沒有嗎?明明在那裡」.
「就在這裡,為什麼看不見?」.
很多啊.
不知道為什麼大家看到的東西不一樣.
我們就叫「開眼」.
有時突然「開眼」就看見了.
或者有時我們有些人.
譬如找工作.

$^{41}$有些人不知道找的過程當中.
有沒有一些特別深刻的.
很困難的.
或者其實很多波折的.
或者有沒有一些.
最後未必找得到的.
有沒有一些.
當中可能有很多的.
感受和情緒.
每個尋找的故事.
都有不同的過程和不同的結局.
有時長有時短.
有時結局是開心.
有時不一定是.
但原來不單止人會尋找.
神也會尋找.
雖然神是掌管萬有的神.
但他也會.
經歷尋找的過程.
神尋找的對象.
通常是什麼?.
是人.
今天的經文.
《約拿書》一章裡面.
神也在尋找一些人.
是什麼人?在過程中有什麼發生?.
我們一起來看.
剛才我們有讀過經文.
我們可以再詳細去看一下.
有沒有人在我們當中.
是未看過《約拿書》.
四章的?.
一二三四章的.
有沒有人是未看過《約拿書》.
這一卷的聖經?.
大約都看過.
雖然未必記得很細緻.
《約拿書》是十二小先知書.
其中一卷.
只有四章.

$^{81}$講述北國先知若拿.
被神呼召到亞述國的.
尼尼微城裡面.
傳講神話語的經過.
而今天的經文是第一章.
《約拿書》的一節.
記載了耶和華對若拿的吩咐.
剛才我們有讀.
耶和華記說.
那米泰的兒子若拿.
說:起來到尼尼微城去.
警告其中的居民.
因為他們的惡行已經達到我面前.
其實在《約拿書》.
雖然是小書卷.
但它有很豐富的內容.
還有很多的修辭手法.
其實是很有趣的.
其中一個就是.
一個大字.
在《約拿書》裡面有很多個大字.
為什麼是這樣?.
讓我們注意.
這個城不是一個普通的城.
是一個很大的城.
在當時來說真的很大.
在神面前也很大.
當然在第四章的時候.
可能會有分曉.
我打算每章會講一次.
講一篇.
有人說《約拿書》也是一個諷刺文學.
諷刺.
因為它有很多對比.
有很多令人驚訝.
甚至和大家所認為的常理.
有相反的東西.
例如在一至二字裡面.
已經有一個驚訝.
出現了.

$^{121}$意外出現了.
通常在聖經裡面.
如果神吩咐他服人說話的時候.
通常下一句是什麼?.
神吩咐以徐.
或者是以利亞.
或者是吩咐摩西.
然後就會說「於是他就起來」.
然後就去做.
按著神吩咐去做.
這是怎樣呢?.
但是《若拿》就完全不同.
它可視.
所以用了一個轉折詞.
可視《若拿》就起來.
它不是去做神要它做的事情.
它是逃走.
在神的面前去逃走.
逃到他斯那裡.
躲避耶和華.
它去到約柏找到一艘船.
要開往他斯.
《新譯本》特別註明.
上船原文是下船.
因為大家知道中文我們說上船.
但其實它是下船.
特別用到這個詞跟後面.
它的狀態有關的用詞.
下船.
向下走.
要和他們.
和他斯離開耶和華.
它是想離開耶和華.
這個也挺有趣.
大家想想.
大家讀的時候.
《若拿》也有形容過耶和華是什麼.
創造天地海洋.
既然神是創造天地海洋.
你可以逃去哪個地方.

$^{161}$可以逃避到神呢.
為什麼它會覺得.
它去到他斯就可以離開耶和華呢.
我們可以看看這個地圖.
有點小.
其實它有個圈圈.
在這裡.
有個尼尼尾城在這裡.
尼尼尾城在這裡.
看到嗎.
尼尼尾在這裡.
約柏在這裡.
這個距離.
遠嗎.
去到這裡.
去到比較遠的地方.
上船.
如果在這個地圖看.
好像是走下去.
已經是一個比較遠的距離.
它本身在猶大地.
在以色列地.
所以.
尼尼尾在那邊.
它不想去那邊.
所以就下來這邊.
但這個不夠誇張.
看下一幅圖.
約柏上船.
目的地是他斯.
他斯在哪裡呢.
他斯在這裡.
亞述在這裡.
尼尼尾就是亞述其中一個城.
看到嗎.
其實很遠.
剛才說的那個.
他們那個地方應該在.
約柏應該在.
約達這邊.

$^{201}$剛才那個距離是.
尼尼尾和上船的地方.
在這裡.
但他想去到這裡.
去到西班牙.
很遠.
東離西又很遠.
叫你去東還是去西.
反其道而行.
所以他的逃避.
他是離開不到神.
所創造的世界.
但他說神要我去這裡.
我逃不去那裡.
他逃避神給他的任務.
那麼.
約拿收到神清晰的吩咐.
他起來但不是進行.
是反其道而行.
目的地是他斯.
很遠.
神會不會讓他去到呢.
神是沒有讓他去到.
但其實這個遙遠的距離.
顯明了約拿.
他唯命的決心.
我們有時會.
這樣說是一個唯命的先知.
通常先知是聽命.
先知是聽命.
因為先知是一個.
神的百姓的一個.
代表.
也是一個榜樣.
所以先知應該是聽命.
而約拿唯命的原因.
是一直到約拿書第四章.
才去揭露.
第一章主要是描述.
約拿對神的一個吩咐的反應.

$^{241}$和他行動帶來的即時影響.
曾經有人問我.
約拿是否真先知.
為何真先知會這樣唯命.
如果是用內地他們很喜歡說的.
用詞就叫做潑逆.
那麼潑逆.
約拿是否真先知.
我們怎樣知道他是真先知.
有認證的.
列王記下十四章二十五節.
這裡特別提到.
神做的一些事.
神曾經預言過一些.
是藉著他的僕人.
加特希忽人.
亞米太的兒子約拿先知所說.
所以你說約拿書是約拿寫.
是不是他自己說自己是.
其實這裡有說.
新約也有引用.
也稱呼約拿是先知.
所以約拿是一個真的先知.
他真的會聽神的吩咐去講預言.
大家也會去聽.
有記錄.
除了尼尼美宣告.
神也透過約拿做其他預言.
所以他的身份是真的.
一個唯命的先知.
是不是真的很難理解.
或者.
一個逃避神的人.
或者一個逃避神的基督徒.
又難不難理解.
弟兄姊妹我們有沒有見過.
逃避神的基督徒.
我們自己.
信主之後有沒有逃避過神.
或者你信主之前也逃避過神.

$^{281}$很久了.
拖很久.
不用信主的.
遲些再想想.
再想想.
可能有些人很小的時候已經聽過福音.
但一直抗拒.
因為那種原因.
讓我先賺夠錢.
到某些事情發生.
才終於願意接受.
耶穌基督為救主.
得著神的救恩.
信主的人也會逃避神.
剛才說了.
明知道神的心意要我們做一些事.
或者不做一些事.
但我們總是有些原因拖延.
是不是.
有沒有試過.
簡單來說.
有些人知道要向家人傳福音.
但很難.
見到家人就開不了口.
見到朋友.
聖靈說不要.
很難很尷尬.
可能是這些.
有些可能是更加重要的事.
我聽過.
就算是教會的.
夢召見證也不少.
有時神父召他們上來.
第一時間說不是馬神不好.
經過很多的印證.
或者一次又一次.
有時很有趣的.
常常在.
廣道裡.
聚會裡.

$^{321}$靈修裡不斷出現同一段經文.
提醒他一些事.
有時是神正在追趕.
追尋你.
就像是士記裡.
基甸也是.
祂一再求聞耶和華的印證.
陽無變濕 濕變乾 乾變濕.
祂才願意去做.
通常當我們逃避神的時候.
我們也有這樣的經歷.
我們的狀態通常是怎樣.
是不是覺得自己靈敏很高.
是不是意志很高昂.
應該不是.
是有些低沉.
有些乏力.
好像若拿的情況一樣.
四至五節.
這裡說.
耶和華在海上忽然掛起大風.
海中狂風大作 船幾乎要破裂.
大家有沒有遇過.
風浪在船上.
或者你坐飛機的時候有沒有遇過顛簸的氣流.
我試過坐飛機的時候.
那個氣流.
好像機動遊戲.
很多人都害怕.
有些人會暈 會嘔.
船也是 如果剛好遇著大風浪.
我們通常都是在安全的環境上.
但萬一.
長途船遇到風浪的時候.
有沒有人會暈船.
會的 是吧.
有些人其實沒有什麼動作也會暈.
但是祂當時.
應該是很誇張.
如果你想像一下.

$^{361}$運用一下我的想像力.
一艘船在海上.
而它幾乎要破裂.
它是一艘長途船.
剛才看到它可以去很遠.
應該是一艘大船 有很多貨在上面.
就好像大家想像.
當年的鐵達尼號.
當然沒有那麼大.
不過類似.
如果你在船上.
而風浪大到.
船幾乎要破裂.
你可以想像激烈的程度.
連水手都很害怕.
水手在船上生活.
所以他們各人向自己的神.
大聲呼求 為了減輕船的重量.
就把船的東西拋在海裡.
希望可以減輕一點.
讓它渡過難關.
但是約拿做了什麼.
又有另一個給我們的驚喜.
它又是不安常地出牌.
正常來說.
一個人如果在這種情況.
這麼危急的時候.
如果飛機上遇到氣流.
可以先醒一下.
隨時有什麼事應變.
或者可以幫一下身邊的人.
我們基督徒.
如果有足夠的能力.
去保護自己之餘.
可以照顧同行和身邊的人.
但是約拿.
他下了船.
然後他就向下走.
下了船艙.
下是很重要的.

$^{401}$顯示給我們看.
約拿整個人的心靈狀態.
也是一個屬靈的狀態.
他的身體不斷向下走.
象徵他的靈明.
他的心靈.
都是向下走.
下到船艙的底.
下到最底的艙.
大家看到這麼大風浪的時候.
他居然可以下到船艙底.
還要做什麼.
睡覺.
睡覺就好像再向下.
沉下去.
在夢裡面.
所以我們看到.
大家想一下.
除了外面狂歡大作之外.
剛才他說水手都很害怕.
大家都很害怕.
大家大聲呼求神.
我們常用的話.
就好像求神拜佛.
大聲喊.
不斷地.
忙忙碌碌地.
將貨物拋下船.
其實很嘈吵.
一片忙亂.
而且當時腎上腺素急速上升.
因為現在要救命.
但若嵐的狀態竟然可以沉睡.
好像他和周遭環境.
我們叫斷線.
完全不配合.
這個畫面.
好像他不是身處奇境一樣.
這裡告訴你.
當人逃避神的時候.

$^{441}$身心會不斷下沉.
使自己和外界有一個隔絕的狀態.
你有沒有試過.
逃避神到一個地步.
合上耳朵不想聽.
當然耳朵沒事.
但其實你的心靈耳朵已經關閉.
人家說什麼都聽不進耳.
有沒有見過.
自己陷入過這種情境.
突然有個執著.
很想要的時候.
人家說什麼都聽不到.
神的話都聽不到.
其他人的需要都看不到.
有時我們因為自己的困境.
因為自己和神之間.
有些基督徒會生神氣.
會怪神.
有時遇到事的時候會.
因為這些狀態.
他已經不能夠.
關懷到身邊的人.
可能在這個時候.
自己好像已經不關心自己.
其實約拿不單不關心身邊的人.
他們也不關心自己的性命.
好像沒所謂.
怎樣就怎樣.
今天有一個詞.
我們對這個叫什麼.
躺平.
我覺得像擺爛.
已經被躺平再進一步擺爛.
所以.
我們覺得很奇怪.
因為一個先知神的僕人.
我們認為神的僕人.
應該對人有憐憫.
有愛心有關心.

$^{481}$關心身邊的人.
但現在約拿連自己的性命也不顧.
對身邊的人的需要也自若罔聞.
但神透過他身邊的其他人.
叫醒他.
這是身體上的叫醒.
但神會一步一步.
使他再次慢慢地醒過來.
船長.
現在派出一個船長.
走到他面前和他說.
你怎麼了.
為什麼你會這樣.
為什麼你會睡得這麼熟.
或者神紀念我們.
使我們不至於喪命.
人就是這樣.
他希望你求你的神.
我求我的神.
總有一個聽的.
我們就能夠得救.
其實也是一個很大的諷刺.
一個外邦人的船長.
提醒神的先知去求告神.
神的先知應該是.
眾人都未懂得去求告神.
之前已經去求告神.
就好像傳道牧者.
帶領大家去尋求神.
但現在又相反.
若嵐好像是一個沉睡的狀態.
有一個外邦人去.
提醒他.
在這個時候.
其他水手都用他們的方式.
去求問神.
他們就抽籤.
想知道是因為哪個緣故.
就抽中了若嵐.
然後就和若嵐說.

$^{521}$請告訴我們.
其實很客氣的.
我們是哪個緣故.
你是做什麼的.
從哪裡來的.
因為他們知道這件事跟他們有關.
就抽籤告訴他們.
他們就說我是希伯來人.
我敬畏耳和天上的神.
創造海洋和陸地.
又是另一個諷刺.
在那個場景他自己想睡覺.
又不求告神.
他一開口就說.
我是希伯來人.
很有驕傲在裡面.
還有他說.
我敬畏的是.
天上的神.
創造海洋和陸地.
但他自己又逃避一個這麼偉大的神.
他說自己敬畏神.
但其實他又一直逃避神.
有些人就害怕.
因為知道他躲避耳和.
若嵐告訴他們.
海浪又越來越洶湧.
他們又和他們說.
怎樣處置你才可以叫海浪平靜.
其實若嵐又告訴他們.
不知道大家有沒有看.
荷里活的大片.
譬如你想像.
讓大家能夠理解那種情況.
譬如荷里活大片裡面.
有一個角色.
和主角.
剛好不幸地和他一起上了飛機.
飛機被人劫持.
原來飛機的人.

$^{561}$是想找這個男主角.
譬如前FBI.
CIA通常都是.
這些情節.
我問他們.
我得罪了一個全世界最大的.
組織.
他們知道.
得罪的是這麼偉大的神.
他們很害怕.
然後他們就問他.
如何處置.
若嵐竟然回答.
這個解決方法.
其實很不合理.
他們說讓我進海.
在一個.
這麼大的海浪.
浪濤洶湧的海.
一個人如果進了海.
等於甚麼.
死了一條.
正常是這樣想.
你懂游泳也沒用.
在這麼大的海浪裡.
若嵐就等於叫他們.
殺了我.
好像是這樣的意思.
但是那些人不想.
那些人都不想.
因為他們覺得很難接受.
你看到那些水手.
他們不想.
他們仍然想靠自己的方式.
去靠岸.
為甚麼那些水手.
他們都想活命.
他們都想.
找出解決的辦法.
可以令風浪平靜.

$^{601}$或者引導自己得救.
為甚麼他們都不願意.
照若嵐所說去做.
當然他們不想殺人.
他們有良心.
水手也會愛惜人的性命.
雖然他們是外邦人.
這是另一個諷刺.
若嵐好像不愛惜自己的性命.
說你們拋我下海.
但其實水手說.
我們做不出這樣的事.
怎可以這樣做.
但他們不得已.
因為海浪就像.
把訊息發給他們.
越來越催逼的時候.
他們沒有其他辦法.
海浪不斷洶湧.
我的態度是.
大家可以嘗試去想想.
他很認識神.
他有關於神的知識.
其實捲了牙說之後.
再告訴我.
他和神的關係很親密.
他也很認識神.
可以和神交流.
他一邊說.
敬愛神 知道神的大能.
但他也沒有相認一致的行動.
而是他為了自己的原因.
他繼續去做自己想做的事.
14至17節.
那些水手不得已這樣做的時候.
他們就祈禱.
呼求耶和華.
我們懇求你不要因為這個人叫我們索命.
也不要將無辜人的血歸我們身上.
因為你耶和華是照自己所起義的行事.

$^{641}$他們要表明.
是你要我們這樣做.
不是我們想的.
我們已經表明我們絕對不想這樣做.
不過現在沒有辦法.
知道這是神的心意.
海浪也作證了.
所以一定要做.
所以他們就將約拿抬起來.
投在海裡 路途就真的平息了.
於是其實這是一個很大的見證.
讓那些水手知道.
原來神真的這麼厲害.
耶和華神真的這麼厲害.
約拿所信的神真的這麼厲害.
原來海浪真的是他的緣故.
和他掌管風浪.
所以那些人就大大敬畏耶和華.
向他獻祭和許願.
另一方面.
其實約拿.
在海裡.
差點淹死.
神就安排一條大魚將約拿吞了.
與福沖三日三夜.
在這裡神施行了兩個極大的神蹟.
一個是.
約拿掉進海裡.
但是不致於死的神蹟.
約拿掉進海裡.
但是他沒有死.
有一條魚吞了他.
當然我們很難想像那條魚怎樣吞了他.
裡面可以三日三夜存活.
而另一個.
是外邦水手.
這班水手或船上的人.
可能不止水手還有其他人.
他們透過這個不信服的.
真的先知.

$^{681}$經歷到神的大能榮耀和拯救.
以致他們歸順神認識神.
這班外邦人本來是求告自己的神.
呼求自己的神.
但是在經歷這件事之後.
他們呼求耶和華.
這可能是他們人生第一次.
向耶和華神去呼求.
因為他們覺得是希伯來人的神去呼求.
神要找的人.
在這裡是誰呢?.
是唯命的先知.
約拿,他要把他找出來.
因為他逃避神.
神好像追著他找他出來一樣.
同時其實神也要找.
這些原本不認識他的外邦人.
約拿他因為自己的原因.
他不願意遵行神的吩咐.
但是神仍然透過狂風暴雨.
透過不信的船長和水手.
去提醒約拿去呼求神.
回到神的面前.
而在約拿面臨被淹死的危機的時候.
神透過一條大魚.
將他吞下肚子.
保護他,拯救他的性命.
也讓他有一個可以反思.
向神禱告的時間.
所以有第二章的禱告.
而神透過約拿的不信服.
也讓原本信服外邦神明的外邦人.
去認識他.
經歷他,敬畏他.
這不是意外.
不是順便.
而是神的智慧和大愛.
他可以使用人沒有想到的.
讓人去認識他.
歸向他.

$^{721}$我們以為先知要聽命.
才能夠被神使用.
其實神仍然能夠使用先知的不聽命.
這就是神的智慧.
神要找的人.
神不會放棄.
而且必定會找到,會尋回.
神要找回那些正在逃避他.
身心靈下沉的人.
神要召回那些雖然認識他.
但不願意聽命的人.
神要尋回那些未曾聽聞他.
行在迷惘中的人.
你覺得你是哪一種.
或者你身邊有哪些人.
今天你的心靈狀態是怎樣.
你最近一次.
感覺和神親近.
是甚麼時候.
有很多弟兄姊妹.
有時跟我說.
我覺得我靈明經常都不好.
最近都不好.
但我一直向下沉.
有時甚至會進入一種.
屬靈的沉睡狀態,好像不願意起身.
一直逃避神.
不願意回到他面前.
有時逃避不一定要睡覺.
但現在通常.
我發覺很多人逃避.
都是用一種方式.
最方便.
用最多時間.
看手機.
明明很多人跟我說.
起身或睡前.
沒有甚麼時間.
有時說納定心智.
今天要親近神.

$^{761}$要靈修禱告.
過了一個小時.
兩個小時.
問他看了甚麼.
其實講不出.
沒甚麼養份.
看些好笑的.
短片.
有時有些人逃避.
好像心裡有個催逼.
要看聖經.
一拿起手機.
開了個抖音.
或其他東西.
沒看.
其實心裡都在逃避神.
有些人.
心裡很怕神會提醒他.
他有明知.
但他又怕神提醒他.
又不可以不做.
有些人.
越陷入這個循環.
越看手機越內疚.
一方面有責備.
為何我又看手機.
越想逃避.
又越看.
進入一個惡性的循環.
如果我們都在這個情況.
我們靠神.
放下手機.
一起身不要馬上伸手.
拿手機.
伸手禱告.
這是我們的習慣.
大家是不是一起身摸手機.
好像有很重要的東西.
我經常覺得好笑.
個個看手機好像有幾百萬上落.

$^{801}$好像有些很重要的東西.
有時有的.
但不是那麼多時間有.
其實我們可能.
都是在逃避一些東西.
我們不要再逃避.
可能今天.
這個訊息都是提醒你.
我們要.
在生活當中.
分別為性.
把時間分出來.
和神獨處.
安靜在神的面前.
去靈修.
這才是我們最需要的.
當你和神的.
關係恢復的時候.
身心靈才有真正的安穩和.
血色.
你認為的難題其實在神面前.
可能我們心態改變的時候.
我們就能夠面對.
就算環境不是即時改變.
但神會幫助我們.
又或者我們會不會慢慢.
不知不覺間我們發覺我們對身邊的環境.
和人.
失去了一種關心.
變得冷漠.
可能我們被今生的思慮.
擠住了.
以致我們信仰好像變成一個空殼.
又或者我們被.
不快的事增多.
讓我們心冷淡.
可能我們看到新聞.
我們很多的激氣.
然後我們慢慢很多東西不想理.
雖然我們會承認我們是.

$^{841}$基督徒,但我們沒有.
活出在主裡面的關心和愛心.
我們有多久沒有將神.
介紹給人.
沒有將人帶到神的面前.
沒有向人傳福音,沒有為人禱告.
我們有沒有將信仰而行一致的活出來.
讓人可以從我們身上.
看到,經歷到神的作為.
如果有其他在這裡聽的弟兄姊妹.
或者你的朋友.
如果還沒有認識耶穌基督的話.
如果還沒有得著.
祂的救恩的時候.
今天神正在尋找你.
願你能夠認識祂.
歸向祂.
讓你的罪得赦免,並且得著永生.
與神復和.
去體會神的愛和恩憐.
在萬太福音十二章三十九.
至四十一節裡面,耶穌基督都提到.
約拿這個事蹟.
耶穌回答說.
邪惡淫亂的世代要求神蹟.
除了約拿才知這個神蹟.
不再有別的神蹟給這世代了.
約拿在大禹福中.
三日三夜,同樣人子也要在地裡.
三日三夜.
審判的時候.
你們未人要與這世代的人.
一同站起來.
他們要定這世代的罪.
因為他們聽見約拿的宣告就悔改了.
看那比約拿更大的.
就在這裡.
即使約拿有不聽命的時間.
逃避神的時間.
但神仍然使用他.

$^{881}$甚至他的.
三日三夜在禹福裡面.
預表了主耶穌.
三日三夜在地裡的神蹟.
我們的耶穌基督.
比約拿更大.
我們都是屬於他的.
耶穌基督.
也是繼續去尋找人.
正是逃避神.
身心靈下沉的人.
只要我們呼求神.
就能夠重新在主裡面得力.
雖然認識他,但不願聽命的人.
只要倚靠神.
就能夠再次為神作見證.
未曾聽聞他.
行在迷茫中的人.
只要歸向神.
就能夠得著無價的救恩.
讓我們一起在神面前.
有一時間安靜的時間.
我們低頭.
在神面前.
安靜我們的心.
在今天的經文.
今天的講道裡面.
神有沒有.
在你說話.
有沒有給你一些提醒.
有沒有看到.
神一直的.
要尋找你.
神很希望.
你能夠每天.
跟他有這份.
交流.
能夠去.
經歷他的愛.
恩典.

$^{921}$也能夠成為.
愛的管道.
去將這份愛.
分享給身邊的人.
接下來有.
兩分鐘的時間.
讓我們在神面前.
去禱告.
我們在禱告.
你就是那位尋找人的神.
透過約拿的經歷.
讓我看到.
你是怎樣不離不棄的.
去挽回我們.
招回我們.
尋回我們.
並且你也愛過一些.
未相信你的人.
你也很願意.
他們也能夠去經歷你.
歸向你.
因為我們也是.
被你沉著的人.
我們.
也是迴轉歸向你的人.
就求你幫助我們.
在生活當中.
我們也常常去親近你.
我們也去.
聽你的話.
讓我們聽到你的話的時候.
你有呼召我們的時候.
我們就去回應.
我們就去遵行.
讓我們能夠將你的愛.
將你的福音帶給.
我們身邊的人.
也都可以帶到.
你要我們去的更遠的地方.
我們祈禱交托.

$^{961}$是奉主耶穌命求.
\newpage



\section{腓立比書 1:12-30}
\label{sec:032221}
\textbf{一切為福音!}~~~~王永揚牧師
\newline
\newline
連結: \href{http://www.church.com.hk/acms/content.asp?site=cdc\&op=show\&type=product\&code=032221\&layout=sermon}{\texttt{032221}} ~~~~ 語音日期: 2026-05-31
\newline
\newline
\hyperref[sec:032216]{\small{< < < PREV SERMON < < <}}
~
\hyperref[sec:index_chronic]{\small{[返順時目]}}
~
\hyperref[sec:032222]{\small{> > > NEXT SERMON > > >}}
\newline
\newline
腓立比書 1:12-30
\newline
\begin{longtable}{cl}
\hline
\hline
章節 & 經文 (和合本修訂版)\\
\hline
1:12 & \begin{tabularx}{0.7\textwidth}{X} 弟兄們，我要你們知道，我所遭遇的事反而使福音更興旺， \end{tabularx} \\ \\ \relax
1:13 & \begin{tabularx}{0.7\textwidth}{X} 以致御營全軍和其餘的人都知道我是為基督的緣故受捆鎖的； \end{tabularx} \\ \\ \relax
1:14 & \begin{tabularx}{0.7\textwidth}{X} 而且那在主裡的弟兄，多半都因我受的捆鎖而篤信不疑，越發放膽無所懼怕地傳道。 \end{tabularx} \\ \\ \relax
1:15 & \begin{tabularx}{0.7\textwidth}{X} 有些人傳基督是出於嫉妒紛爭；有些人是出於好意。 \end{tabularx} \\ \\ \relax
1:16 & \begin{tabularx}{0.7\textwidth}{X} 後者是出於愛心，知道我奉差遣是為福音辯護的。 \end{tabularx} \\ \\ \relax
1:17 & \begin{tabularx}{0.7\textwidth}{X} 前者傳基督是出於自私，動機不純，企圖要加增我捆鎖的苦楚。 \end{tabularx} \\ \\ \relax
1:18 & \begin{tabularx}{0.7\textwidth}{X} 這又何妨呢？或是假意或是真心，無論如何，只要基督被傳開了，為此我就歡喜。我還要歡喜， \end{tabularx} \\ \\ \relax
1:19 & \begin{tabularx}{0.7\textwidth}{X} 因為我知道，這事藉著你們的祈禱和耶穌基督的靈的幫助，終必使我得到釋放。 \end{tabularx} \\ \\ \relax
1:20 & \begin{tabularx}{0.7\textwidth}{X} 這就是我所切慕、所盼望的：沒有一事能使我羞愧；反倒凡事坦然無懼，無論是生是死，總要讓基督在我身上照常顯大。 \end{tabularx} \\ \\ \relax
1:21 & \begin{tabularx}{0.7\textwidth}{X} 因為我活著就是基督，死了就有益處。 \end{tabularx} \\ \\ \relax
1:22 & \begin{tabularx}{0.7\textwidth}{X} 但是，我在肉身活著，若能有工作的成果，我就不知道該挑選甚麼。 \end{tabularx} \\ \\ \relax
1:23 & \begin{tabularx}{0.7\textwidth}{X} 我處在兩難之間：我情願離世與基督同在，因為這是好得無比的； \end{tabularx} \\ \\ \relax
1:24 & \begin{tabularx}{0.7\textwidth}{X} 然而，我為你們肉身活著更加要緊。 \end{tabularx} \\ \\ \relax
1:25 & \begin{tabularx}{0.7\textwidth}{X} 既然我這樣深信，就知道仍要留在世間，且與你們眾人一起存留，使你們在所信的道上又長進又喜樂， \end{tabularx} \\ \\ \relax
1:26 & \begin{tabularx}{0.7\textwidth}{X} 為了我再到你們那裡時，你們在基督耶穌裡的誇耀越發加增。 \end{tabularx} \\ \\ \relax
1:27 & \begin{tabularx}{0.7\textwidth}{X} 最重要的是：你們行事為人要與基督的福音相稱，這樣，無論我來見你們，或不在你們那裡，都可以聽到你們的景況，知道你們同有一個心志，站立得穩，為福音的信仰齊心努力， \end{tabularx} \\ \\ \relax
1:28 & \begin{tabularx}{0.7\textwidth}{X} 絲毫不怕敵人的威脅；以此證明他們會沉淪，你們會得救，這是出於神。 \end{tabularx} \\ \\ \relax
1:29 & \begin{tabularx}{0.7\textwidth}{X} 因為你們蒙恩，不但得以信服基督，而且要為他受苦。 \end{tabularx} \\ \\ \relax
1:30 & \begin{tabularx}{0.7\textwidth}{X} 你們的爭戰，就與你們曾在我身上見過、現在所聽到的是一樣的。 \end{tabularx} \\ \\
[1ex]
\hline
\hline
\end{longtable}
$^{1}$亦是我在尖沙咀浸信會的最後一次講道.
我們有機會在尖沙咀浸信會受按 成為牧師.
多謝主任牧師羅太 牧師的提攜和指導.
多謝辦事處的幹事 工友 一同協助.
讓我來到這裡 學習 成長.
更加得蒙弟兄姊妹的照顧.
在主裡面 雖然在不同的地方.
但我們仍然是主代一家 我們江湖再見.
在正道之前 我們同心祈禱.
帶領我們進入真理.
恭應將以下環節交託仰望.
求主保守 引領祈禱.
賴奉耶穌基督 聖命祈求 阿們.
被監禁或囚禁這段時間.
來向這四間教會或個人 所寫的書信.
當時保羅身處在凌雨.
在監獄裡 不能失去自由.
但聽聞教會發生問題 內部不和.
在信仰裡面 面對挑戰 受不下異端的攻擊 誘惑.
他未能親身來到當中.
於是他來收書 用書銜.
來按當時有專職的遊帝人員.
去送上用蒲草紙所寫成的書銜.
親身來信道 去收信人.
大多數是教會領袖.
教會領袖收到書信後.
就會在公開的場合.
就像現在這樣.
公開來宣讀.
讓姐妹們一起來聆聽.
保羅所寫的書信當中的教導.
而這些書信 被教會收集 傳遞.
經歷了一段時間之後.
在大公教會的會議.
在書信 書卷的致正.
和聖靈的引導下.
來視為正典.
然後被收納 成為我們今天的新約聖經.
昔日 保羅被囚在羅馬.
等候羅馬皇帝薩戈親自來審判.

$^{41}$就像今天 我們去不服這個地方法院.
我就上訴到高等法院.
高等法院再上訴到終審法院.
終審法院一錘定音.
這個就不能再犯案.
保羅就運用他自己是羅馬公民的身份.
上訴到當時的皇帝凱撒.
來審訊他的案件.
但由冰冰看守 放在其他地方.
容許其他人來探訪.
就像之前疫情一樣.
我們要不在家居隔離.
要不送去方艙社區隔離設施.
聽到名字 或者看到電視直播.
或者去現場直播.
看到那些貨櫃場.
看到那些還在流水.
風會吹進來的那些情景.
當時保羅也是不遑多讓.
沒有此景的等待.
羅馬皇帝不是說你見就見.
去國事這麼多.
一個人一個人走來.
馬上見你 馬上審你案.
他就不得了.
排期都不知道排到多久.
在無此景的等待.
是需要人照顧和接濟.
在這樣的情景下.
去收書給這個菲律賓的教會.
一方面要多謝他.
送來的物資.
供應他身體上的需要.
更加多謝他.
差派伊巴忽提到現場照顧他.
伊巴忽提病好了.
就帶了這封書信.
回到菲律賓的教會.
道謝.
如果你是保羅.

$^{81}$有機會寫這封書信.
給你的米飯班主.
真是靠他.
吃粥吃飯真是靠他.
那你會寫些什麼呢?.
你想想.
如果你代入保羅的情景.
你會寫些什麼呢?.
我記得我小時候.
我們死坑死底.
依宿蝕.
將物資寄去.
帶給國內的親友.
然後他們就收書.
一封信函下來.
最少是兩頁紙.
第一頁千字文.
千多謝.
萬多謝.
滔滔不絕江水.
寫到水蛇充滿腸.
多謝.
但後來那一頁才是真章.
是說什麼呢?.
我還想我們還有什麼缺陷.
有什麼需要.
有什麼不足和缺乏.
但宗觀.
腓立比書和其他的三封監獄書信.
保羅都沒有列出.
沒有窮舉.
他在物質上的缺陷和需要.
你會覺得很出奇.
保羅是不是天生異稟.
不用吃不用住.
我今天講到的兩個重點.
因為保羅是以福音為先的關注.
以及他是以福音為先的實踐.
第十二節.
弟兄們.

$^{121}$我願意你們知道.
我所遭遇的事更叫福音興旺.
是回應九到十一節我所禱告的.
就是你們看愛心和各樣知識上多而又多.
使你們能夠分辨是非.
作誠實無過的人.
直到基督的日子.
並靠著耶穌基督.
結滿人意的果子.
叫榮耀歸於神.
保羅就是作為一個現身說法的榜樣.
是怎樣在知識和見識上多而又多.
能夠分辨是非.
能夠結滿人意的果子.
叫榮耀歸於神.
保羅願意肥立彼教信徒知道的.
是有關於我的事.
或者是發生在身上的事.
要由第十三節和十四節補充.
這個事就是指保羅被囚.
被捆綁.
而福音興旺.
就是說進展.
就是說在這個情況.
竟然是相反.
是因減不加.
保羅超脫了現實的處境.
從被囚的環境.
看到事實的真相.
從自身的限制.
看到福音不被爛了.
這個就是迴響聖傳.
二十一章開始.
他由耶路撒冷來還願.
然後在耶路撒冷聖殿裡面被捉拿.
一幕一幕.
怎樣千辛萬苦排除萬難.
來到羅馬這個地方.
先是被人毆打.
然後被下在監牢裡面.

$^{161}$然後去面對很多不同的官員的質詢.
他曾經在巡撫面前.
在阿基伯王分封王面前.
來自辯.
去漏夜被帶到凱撒利亞.
逃避猶太人的追殺.
到二十一章真的上路了.
去到羅馬這個地方之前.
經歷海難.
船隻被拋來拋去.
在大風大浪幾乎沒頂.
最終能夠去到瑪爾他島來登陸.
但又被蛇咬.
但又無餓無愛.
最終來到羅馬等候審判.
作為讀者.
當然很想知道保羅上訴的結果.
我以前初信的時候.
因為哲行傳就好像章回小說一樣.
一幕又一幕.
我一口氣看完又再看.
情節很吸引很緊湊.
就好像看金庸的武俠小說.
《神鵰俠侶》.
楊過來到絕情谷.
見到小龍女克子在懸崖上說.
十八年後再相見.
我就想十八年後會不會真的出現.
整個故事發生就是情節說這些.
十八年間楊過怎樣去生活.
到最後他終於能夠去重逢.
同樣作為讀者的我們.
都很想知道保羅上訴的結果.
但偏偏路加醫生不風一轉.
在最後的時間.
在哲行傳二十八章來到港島.
保羅在自己所租的房子住了足足兩年.
凡來見他的人.
他全都接待放膽講論神國的道.
將主耶穌基督的事教導人.

$^{201}$並沒有人禁止.
可以看到路加很明白.
保羅的關注不在自身.
他更加關心福音的發展.
按一般人的看法.
失去自由必是令到福音的工作大受限制.
就好像先前疫情.
教會不能公開聚會.
我們都只是網上來看崇拜.
網上來參與團契.
出來行山都不可以.
都不可以一齊拍照.
不可以公開在會.
但這是否難阻到福音的工作呢?.
當真.
果然.
透過網絡不同的時區的人.
可以一齊相見.
我們看到.
由肢體帶領在遠在澳洲的家人.
信耶穌.
透過網絡.
然後將資訊交給當地的教會.
由澳洲的牧者跟進.
最後他的親戚.
能夠在澳洲受浸加入教會.
我也知道有某些教會.
在疫情的時間.
向他們的會員派發物資.
在會員換証的時間.
180個會員.
來到教會換証.
要填表,拍照等等.
就在填表,拍照等待的時間.
是怎樣呢?.
我們講三幅.
早,中,晚180個人用了一個月.
作過講.
大家也知道牧者不講名字.
忙到抽筋.

$^{241}$但是很開心.
這180人裡面有30個人.
決志信了耶穌.
這個疫情.
人受限制.
是否成為難阻呢?.
重點是有沒有決心和胸襟.
13節.
「以致我受的困所.
在禦刑全軍和其餘的人中.
已經顯明是為福音的緣故.
保羅要上告當時的皇帝蓋薩.
他是重義的囚犯.
由禦林軍來看守.
當時駐軍在羅馬的禦林軍.
大概有9千到1萬人.
我們去計算.
保羅怎能夠向這麼多人傳福音呢?.
如果用早,中,晚的說法.
一天早,中,晚三間.
每一間向三個人傳福音.
一個星期做六天.
即是一個年52.14個星期.
他能夠向2815個人傳道.
9千人要3.2年才能傳完.
有沒有可能呢?.
真的勸棄了.
向每個士兵談道根本是不可能的.
最有可能的解釋.
是保羅在千辛萬苦排除萬難.
終於來到羅馬.
他向皇帝上訴.
是因為信仰的緣故.
定估都不會為子虛烏有的事情.
冒死來送人頭.
保羅堅持信仰誓不低頭.
不委曲求全.
更加放棄自由.
被囚於獄中.
這份剛毅不屈.

$^{281}$贏取了士兵的認同.
他的事跡廣泛在禦林軍中流傳.
我以前借去木湖一位童工講述.
他在希腹讀神學.
作為神學生的見證.
他和一班弟兄到深水Po 龍屋服侍.
即是龍屋住在龍裡面.
你猜他會做什麼服侍?.
到深水Po 龍屋服侍.
義工.
你猜他會做什麼服侍?.
派飯?.
遊鐵閘?.
不是,是洗廁所.
他說場景污穢化.
污糟化.
好像人間地獄有多污糟.
因為坐廁塞.
有些人垃圾照扔下去.
塞滿了.
漏出來滿死.
那些男人就地屙.
大小二面在廁所裡.
就地屙.
他說他去清潔.
戴著兩個口罩.
都聞到一股味道.
真是受不了.
他說以前都不知道為什麼有這樣的能耐.
做到這件事.
整個禮拜都吃不下飯.
他做完之後.
他說那種情景.
他一路清潔.
在洗坐廁的時間.
有些人在旁邊屙尿.
這樣去清潔.
龍屋的人.
那些男士.
在做的時候.

$^{321}$那些人在說髒話.
在罵.
為什麼這麼污糟.
自己就射到到處都是.
類似這樣.
然後.
那些人.
離開的時候.
你猜龍屋的人.
怎樣跟我們說.
他說.
我們欣賞.
他說你們這些耶穌佬.
跟我們非親非故.
山長水遠.
無離拉更.
走來幫我們做義工.
你們不是瘋子.
就是愛心爆棚.
我們都做不出.
每天看到我們都受不了.
你走來可以這樣.
我們是極之欣賞.
你們的愛心付出.
這些東西應該怎樣.
業主做的.
你這麼傻幫我們做.
我們真金白銀交租.
業主應該提供一個廁所.
可以去到的廁所.
不是應該你們來做的.
所以我們是欣賞.
我們童工.
說了這麼多年之後.
依然是記憶猶新.
那個情景.
他說不能夠磨滅.
我相信.
這班龍屋的居民.
他這一天.

$^{361}$唯有有心之念.
他還記得這班傻佬.
這班耶穌佬.
竟然這樣走來.
用愛心來服侍.
因為一切都是為了福音.
每天有不同的士兵.
來到囚室近距離的.
看看保羅何許人也.
保羅的言行舉止.
成為儒林軍茶餘飯後的話題.
保羅就好像一場戲.
保羅就成為了示範單位.
贏取士兵的信任.
一切都是為了福音.
對於保羅以福音為先的態度.
保羅的羅馬信徒.
有什麼反應呢?.
十四節提到.
「並且那在主裡的弟兄.
多半因我受的菌所就篤信不疑.
越發放膽傳神的道母所懼怕」.
保羅對信仰的堅貞.
傳遞感染了信徒.
比較之下更加積極.
更加勇敢傳揚福音.
基督信仰因為保羅這個罪犯.
成為大眾的議題.
正常來說羅馬教會怕被連累.
會盡量減低他的存在感.
試圖與保羅劃清界線.
相反.
羅馬教會的熱心.
是保羅所感染的.
保羅以福音為先的態度.
反映他在得悉羅馬信徒.
信仰群體傳福音後的反應.
剛才主席讀經文來講.
當有兩群人用不同的動機.
來傳揚福音.

$^{401}$有人是結黨並不誠實.
要加珍保羅菌所的苦楚.
但有些人是真心知道.
保羅是為福音的緣故.
但保羅的說法是.
無論怎樣.
基督究竟被傳開了.
為此我就歡喜.
第一群人是受保羅感召.
積極投入傳福音.
是可以理解的.
第二群人到底是什麼人呢?.
更加有不同的說法.
但從保羅書信裡見到.
保羅對別的福音.
對異端.
他都是當面職責.
強烈反對.
如果這些人說福音的內容.
有問題的話.
保羅不會不斥責他們.
保羅的斥責.
是講到他們傳福音的動機.
保羅作為愛來者.
名聲竟然如日中天.
他的名聲在猶太人當中.
一直被惡意重賞.
說他教導人不遵守律法.
保羅每到不同的地方.
都有一班宗教分子.
來跟著他.
來破壞他的福音工作.
挑起仇恨.
阻礙福音的工作.
保羅最終用自己是羅馬公民的皇牌.
來到羅馬這個地方.
第三天.
羅馬的猶太人就找上門.
經文裡說.
要保羅解釋被攻擊的原因.

$^{441}$詩行短二十八章二十二節.
「但我們願意聽你的意見.
如何因為這教門.
我們曉得是到處被毀謗」.
可能這群人是嫉妒保羅的成功.
更加不願意見到保羅的影響力.
繼續擴大.
為了爭奪市場的佔有率.
他們比先前更加積極傳揚福音.
保羅有沒有為自己自辯呢?.
有沒有說他們去搶奪自己的功勞呢?.
沒有.
保羅一詞定音說到.
基督究竟是被傳開.
為此我就歡喜.
回響路加福音.
九章四十九到五十節.
約翰走來.
向耶穌領功.
我阻止了那個人奉你的名來趕鬼.
因為他不是和我們一起的.
耶穌怎麼說呢?.
耶穌說.
「不要禁止他們.
因為不抵擋你們的.
就是幫助你的」.
保羅的情操是.
「成功不需因我.
但成功只需由我」.
一切都是為了福音.
保羅不單以福音為先作為考量.
更加實踐以福音為先.
保羅等待審判的處境.
引發了他深刻的反思.
他主要擔心兩件事.
第一是監禁和審判的結果.
如果我們從事行傳來看到.
他坐了兩年多.
我們有讀者的視覺.
知道他兩年多.

$^{481}$但是保羅不知道的.
保羅不知道自己要等多久.
大家都試過排街症.
有時候真的不知道要等多久.
每個人都走了.
明明我是最早來的.
為什麼每個人都走了.
每次都是走成我呢?.
不知道要等多久.
第二是死亡的可能性.
因為上告該殺.
是有機會死刑的.
因為來到最高的法院.
去宣判死刑的話.
就不能被推翻.
是一定執行.
隨著等待的日子.
面對的壓力不斷增加.
保羅能夠擴大地談論生死.
不是他心臟強大.
而是願意選擇以福音為先.
生死這個課題.
正正是表述他以福音為先.
20-21節說.
「照著我所切勿所盼望的.
沒有一事叫我羞愧.
只要凡事放膽.
無論是生是死.
總叫基督在我身上照常掩戴」.
保羅的喜樂和自信.
是同一種態度的兩種表現.
他相信靠著神的恩典.
他的願望能夠達成.
為什麼他說到自己的生死.
好像與自己無關.
好像是時下的年青人.
去回應其他人.
回應老師.
回應家人的血面的時候.
他會說.

$^{521}$「你忍我吧.
我管你死」.
或者是否他收到內幕的消息.
收到風.
可以大步跨過.
等待結果是令人勞累.
等待去報告.
醫生的檢查報告.
身體的報告.
在手術室外.
去等候家人.
去接受手術出來.
這個等待是令人勞累.
患得患失.
與保羅面對人生生死之間的處事態然.
進入了平衡時空.
因為他確信自己已經得救.
因為我知道.
「這是藉著你們的祈禱.
和耶穌基督之靈的幫助.
終必要我得救」.
問題的重點.
是「得救」是指什麼?.
保羅談到他得救的語氣堅定.
這是帶有自信的意味.
「得救」可以指身體救贖.
宣判無罪.
因為羅馬政府仍然不能區別.
基督教和猶太教之間兩者的差別.
所以包容猶太教的前提下.
對保羅的案件可以寬大處理.
也可以解釋保羅面對審判結果.
無愧無懼.
生死皆可.
繼續拘禁或者被判處極刑.
一切都是處許可.
有神的心意和計劃.
「得救」也可以指終末的救恩.
保羅可以引用七十字譯本.
「勃不自羞愧」.

$^{561}$明顯是借用約伯記十三章十三到十八節.
約伯深信在主面前去陳述他的案件時.
終必得救.
這三個解釋都有聖經學者支持.
就經文的文氣來說.
「照著我所切勿我所盼望的.
沒有一時叫我羞愧.
凡事放膽.
無論是生是死.
總叫基督在我身上照常顯大」.
保羅不在乎上訴羅馬皇帝終審結果.
他更在乎透過他生死.
面對生死的表現.
成為耶穌基督活培的見證人.
來到當時的凱撒皇帝面前.
這就是他見證基督的機會.
他能夠將他的信仰陳明.
他善忘神彼的恩典.
在這個時刻.
可以宣認信仰的時候.
「不自蒙羞」.
羅馬的審判對他來說是一個機會.
經過生命的歷練.
信仰的沉澱.
去作出信心的宣告.
保羅是表達一種向死而生的決絕.
在這段落結束的時間.
保羅以假設可以選擇生或死的難題.
他如何在生死之間作出選擇.
為主殉道與處同在.
是好得無比的.
或者存留性命繼續與肥立比教會並肩作戰.
兩個都是好的選擇.
最終他選擇生.
若果可以的話.
他捨棄自己的想法.
寧願為信徒的遠故.
繼續為主作歌.
活在主內.
不論動機如何.

$^{601}$保羅就在任何選擇.
任何情況.
他的生命去到盡頭.
他所關心.
他所實踐的.
基督在他身上照常顯大.
見證神恩惠的福音.
近來在網絡見到一則新聞.
是拍攝一位88歲的奶奶和孫子.
陪同一口氣.
走過14層樓梯.
回到已經被燒到昏黑的家居.
兩個人帶著頭燈.
因為沒有電.
房子是黑的.
來到現場.
拿回覺得有紀念價值的物品.
限時限刻.
每個人都只能有限的時間.
在家裡做什麼呢?.
這個婆婆.
帶著頭燈.
還在家裡.
借去倖存的電子琴.
做什麼呢?.
在彈《耶和華是愛》.
她沒有激情地控訴.
沒有狂呼大叫.
只是堅定地咒惡.
面對無情的大火.
天災人禍.
老姐妹堅定地宣告.
《耶和華是愛》.
比什麼大道理更加令人信服.
她見證了.
活出了.
福音.
\newpage



\section{約翰福音 7:37-39}
\label{sec:032222}
\textbf{榮耀君王:世界的光}~~~~林鍵牧師
\newline
\newline
連結: \href{http://www.church.com.hk/acms/content.asp?site=cdc\&op=show\&type=product\&code=032222\&layout=sermon}{\texttt{032222}} ~~~~ 語音日期: 2026-05-30
\newline
\newline
\hyperref[sec:032221]{\small{< < < PREV SERMON < < <}}
~
\hyperref[sec:index_chronic]{\small{[返順時目]}}
~
\hyperref[sec:032223]{\small{> > > NEXT SERMON > > >}}
\newline
\newline
約翰福音 7:37-39
\newline
\begin{longtable}{cl}
\hline
\hline
章節 & 經文 (和合本修訂版)\\
\hline
7:37 & \begin{tabularx}{0.7\textwidth}{X} 節期的最後一天，就是最隆重的一天，耶穌站著，喊著說：「人若渴了，到我這裡來喝！ \end{tabularx} \\ \\ \relax
7:38 & \begin{tabularx}{0.7\textwidth}{X} 信我的人，就如經上所說：『從他腹中將流出活水的江河來。』」 \end{tabularx} \\ \\ \relax
7:39 & \begin{tabularx}{0.7\textwidth}{X} 耶穌這話是指信他的人要受聖靈說的；那時還沒有賜下聖靈，因為耶穌還沒有得到榮耀。 \end{tabularx} \\ \\
[1ex]
\hline
\hline
\end{longtable}
$^{1}$歡迎大家來到我們的「周獨崇大」.
上年7月我們開始研讀.
《新約聖經》當中的《約翰福音》.
我們講到系列的題目叫做「榮耀君王」.
我們花了半年的時間.
去完成《約翰福音》一至四章的宣講.
接著今年的一月份.
我們開始研讀《約翰福音》第五章.
四個星期前我們來到《約翰福音》的第七章.
這個是《約翰福音》第七章的一個大圖畫.
從故事的時序來看可以分為三大段.
第一段是猶太人的駐彭捷近了.
但還沒到.
背景是駐彭捷.
第二是節期.
七天八天的節期已經過了一半.
耶穌到了耶路撒冷.
這是第二段.
第三段是節期的最後一天.
在節期的最後一天.
耶穌說了兩句話.
一個是「我是生命的活水」.
另一個是「我就是世界的光」.
兩個星期前我們讀了3A.
3A就是「我是生命的活水」.
嚴格來說3A並不是《約翰福音》裡面.
《約翰福音》裡面有七個.
耶穌說「我是」.
「我是XXX」.
「我是生命的糧」是第一個.
「我是生命的活水」並不是第二個.
祂不是七個「我是」的一個.
是因為在故事當中.
耶穌其實沒有明確聲明.
大聲吶喊說「我就是生命的活水」.
是我們讀完祂那段說話之後.
做出來的總結.
但是下面那個3B.
就是耶穌真的自己說的.
祂說「我就是世界的光」.

$^{41}$我想你知道一件事.
就是3A和3B.
其實是分不開的.
如果你問我的話.
我會告訴你.
3A和3B不是兩個啟示.
而是兩個連在一起的啟示.
我們可以重新整理說.
3A加3B就是.
「我是生命活水」再加「世界之光」.
是同一句說話來的.
於是「世界之光」和「生命活水」.
是同一樣東西.
而另外一件事你要知道的.
就是「世界之光」和「生命活水」.
都需要等待「聖靈降臨」才可以實現的.
如果你兩個星期之前在這裡.
你聽我說「生命活水」的時候.
我們說了很明顯「生命活水」是在說「聖靈」的.
於是「聖靈」不來.
「生命活水」是出不來的.
我想你用同一樣的眼光去想.
除非「聖靈」來.
如果不是.
耶穌是沒有辦法變成「生命之光」的.
「聖靈」不來.
耶穌是沒有辦法變成「世界之光」.
「生命之光」.
說到這裡可能有人都讀過《約翰福音》.
會問我林牧師.
你好像漏了一段說話.
漏了一個故事.
其實這個週中有.
我想起碼有三,四個人問過我這個問題的.
為什麼我們不說《約翰福音》第八章一至十一節.
那個是我最喜歡的配音故事.
有一個年輕女孩犯了奸淫被人抓了.
然後耶穌說.
「我也不審判你了」.
有人要試探耶穌.

$^{81}$叫耶穌拿石頭扔死她.
然後耶穌說.
在地上寫了一些字.
然後耶穌說叫她.
Go and sin no more.
不要再犯罪了.
林牧師為什麼你不說那個故事.
那個故事就夾在3A和3B之間.
不說這個故事的原因.
原來是在最早,最可靠的聖經古卷裡面來說.
這個故事是沒有出現在《約翰福音》裡面的.
於是大部分的聖經學者其實都相信這個故事.
就算是真實的.
真的有發生的事情.
真的有遇見過這個犯奸淫的女子的話.
都是後來的人加入去的這個故事.
在《約翰福音》裡面.
是沒有這個故事存在的.
於是我們要skip over不讀這個故事.
對於我們而言.
那個很動人的故事.
其實打斷了我們的思路.
令我們覺得生命活水和世界之光是兩張聖經.
你的聖經是第七章,第八章.
兩個故事其實不是的.
其實兩個故事是同一個故事來的.
同一天發生的.
前後腳的.
兩個故事應該是同一張聖經裡面的.
正如我所說.
之前我剛才有提示過.
這兩個故事是雙胞胎來的.
甚至是那種黏在一起的雙胞胎.
不是前後腳出來的那種雙胞胎.
是兩個人黏在一起的雙胞胎.
他們加起來.
才是耶穌的自我啟示.
於是兩個星期之前.
我們讀的只是前面的一部分.
現在我們要看後面的一部分.

$^{121}$接著我們加起來.
你就會知道耶穌究竟在說什麼.
我想帶你快快去溫習一下生命的活水.
接著再連結今天我們新要學的世界之光.
我們回到生命活水開始.
節期的最後一天.
就是最隆重的一天.
耶穌站起來喊著說.
這個背景是猶太人的駐彭寺.
他們是紀念上帝怎樣照顧了他們在曠野裡面40年.
於是他們一個星期裡面.
他們住在彭寺裡面.
好像紀念當年上帝怎樣帶他們出埃及.
接著進入迦南.
在節期的最後一天.
聖經告訴我們.
七天的節期裡面.
每一天大祭司都要用一個金水壺.
去帶領一個隊伍去西羅亞祠那裡拿一些水.
接著經過水門來到耶路撒冷聖殿的獻祭的地方.
就是他們殺動物獻祭的地方.
去除罪的地方.
接著將那個水倒進去祭壇裡面.
有聖經學者告訴我們.
當他們倒水的時候.
他們都會獻上酒祭.
於是有水有酒.
酒是紅色的.
酒是葡萄酒.
於是一起有水和酒倒進聖殿的祭壇裡面.
你可以想像這是一個何等榮耀的福音圖畫.
因為當耶穌掛在十字架上裡面的時候.
那些兵兵一戳他的時候.
他也有水和酒.
或者他的血都流出來了.
於是我們已經在這個故事裡面.
看到一個福音的圖畫.
指向最後耶穌在十字架上面.
為了罪人除罪而作的犧牲.
又有學者告訴我們.

$^{161}$來到第七天的時候.
同一樣的水禮會重複七次.
接著當他做完七次之後.
耶穌大聲的宣告.
他說以下的說話.
「我就是生命的活水」.
但是也有學者說.
這個最後的一天未必是第七天.
而可能是第八天.
第八天叫做擦鵬日.
就是已經鑄鵬了七天.
接著可以擦鵬.
第八天仍然是一個很開心的時候.
但是第八天已經再沒有水禮.
無論是第七天也好.
第八天也好.
基本上就是在水禮要結束的時候.
耶穌突然間站起來.
他說「我就是活水」.
「我要取締了鑄鵬節裡面的水禮」.
「你以後再不需要做水禮」.
「因為我來了」.
「我要成為那個水禮」.
耶穌說「我再不是一時一刻的水禮」.
「而是一個源源不絕的活水江河」.
耶穌這麼說.
「人若渴了可以到我這裡來」.
「生我的人就如經上所說」.
「從他的福眾將流出活水的江河來」.
這是一個很驚人的宣告.
究竟是什麼意思呢?.
作者約翰嘗試給我們解謎第39節.
約翰說.
耶穌的意思是.
「耶穌這話是指著信他的人」.
「要領受聖靈所說的」.
「那時還沒有賜下聖靈」.
「因為耶穌還沒有得到榮耀」.
原來所謂的活水江河.
是聖靈的降臨.

$^{201}$將會有一個新的時代.
聖靈的時代.
聖靈將要全面降臨在.
如果你們記得兩個星期前我所說的.
每一個信上帝的人的身上.
都不是一個一小撮領導級別的人.
好像在舊約聖經裡面.
我們有摩西,約書,大衛這些人.
獨有的經驗.
而是在新約裡面.
從最小的到最年長的.
從最體弱的到最強壯的人.
每一個信耶穌的人都會領受聖靈.
聖靈的同在,聖靈的經驗.
聖靈的加靈.
但這裡約翰說.
「然而那時還沒有賜下聖靈」.
「因為耶穌還沒有得到榮耀」.
那時是在主鵬子的時候.
主鵬子的時候聖靈還沒有來.
「因為耶穌還沒有得到榮耀」.
所謂得到榮耀就是上十字架.
還有六個月.
再過六個月耶穌就要得到榮耀.
耶穌會掛在十字架上.
再過六個月就有血和水從祂的身體流出來.
再過六個月.
聖靈就會降臨在門徒的身上.
耶穌要上十字架.
耶穌要為罪人替死贖罪.
耶穌第三天要從死裡復活.
耶穌要得勝.
耶穌要升天作王.
然後耶穌要差遣祂的聖靈.
將祂為我們得著的贖罪.
擺進每一個相信祂的人身上.
榮耀十字架.
聖靈降臨.
大家可不可以記住這一點.
有兩件事我們一直在等著.

$^{241}$約翰告訴我們.
你一定要等到那天.
事情才會成就.
你需要榮耀的十字架.
你需要聖靈的降臨.
這兩者不單止是耶穌和信祂的人.
福中流出生命活水的關鍵.
你要福中.
我們要領受聖靈的話.
我們要福中流出光水.
必須要有耶穌上十字架.
必須要有聖靈降臨.
同時這兩件事也是耶穌和信祂的人.
要變成為世界之光的前提.
就是說如果耶穌要成為世界的光的時候.
祂必須要先得到榮耀.
然後聖靈需要降臨.
如果我們要參與.
耶穌裡面成為世界的光的一部分.
你去到最後最後你就知道了.
原來世界的光不單止是耶穌.
是關於我們的.
如果我們要參與世界的光的時候.
耶穌必須要上十字架.
然後聖靈必須要臨在.
我們才有可能參與成為世界的光.
我想帶你去到《讓步音》第八章第十二節.
我們現在從第七章跳到第八章.
但是你要記得.
那個《讓步音》第八章的宣講.
我是世界的光.
跟從我的必不在黑暗裡面走.
卻要得著生命的光.
我要你記住這句話.
和第七章三十九節是連在一起的.
我想你要知道.
如果八章十二節要承傳的話.
上面的影子要拿開變成實體.
現在還不可以.
因為上面好像你去一家燒臘店還沒開幕.

$^{281}$上面那個罩還是罩著的.
還沒正式開幕的.
因為直至耶穌得榮耀上十字架.
直至聖靈降臨的時候.
耶穌還沒可以成為世界的光.
生命的光.
耶穌的宣講是這樣.
第八章第十二節.
我就是世界的光.
跟從我的必不在黑暗裡面走.
卻要得著生命的光.
我想你留意一件事.
世界的光其實是等於生命的光.
世界的光等於生命的光.
成世界.
你得著生命的光.
然後把它成到全世界都是.
那就是世界的生命光.
或者生命的世界光.
在這裡如果你熟悉聖經的話.
你知道作者約翰.
彷彿帶我們回去創世記第一章.
創世記第一章上帝起初創造世界的時候.
先做些什麼的.
Let there be light.
起初上帝創造世界的時候.
第一樣創造的是光.
有光之後才開始有生命.
有花草樹木.
有動物.
最後創造了人.
所有的東西從光而來.
這個是創世記第一章上帝創造的時候.
Let there be light.
有光之後才有生命.
然而現在我們在說的是.
耶穌的真光來到世界的時候.
當耶穌到城肉身降生成人.
整個世界和創世記第一章的世界.
是完全不一樣的.

$^{321}$起初創世記第一章裡面有黑暗.
那個黑暗是一個物理上的黑暗.
是一個物質上的黑暗.
現在這個在說的黑暗是.
我們犯罪之後.
我們叛逆我們罪孽裡面的黑暗.
如果要將墮落的罪人.
就是我們這些已經是亞當犯罪之後.
墮落了生下來的人.
從罪惡黑暗裡面拯救出來.
並要帶領我們進入一個明光與生命.
不是一時一刻的明光與生命.
而是永遠長存的明光與生命.
如果上帝要將我們從黑暗裡面拿出來.
然後帶我們進入明光與生命.
我們需要的再不是起初創造的大能.
起初上帝說一句話就搞定了.
我們現在需要的是榮耀十字架的救贖.
因為罪人都犯了罪.
虧缺了上帝的榮耀.
罪的功近乃是死死後且有審判.
上帝必須要為我們支付我們的罪債.
這個黑暗和起初的黑暗不一樣.
我們現在需要耶穌上十字架為我們支付罪債.
然後我們需要聖靈降臨重生我們.
因為我們死在罪中.
我們需要重生的大能.
令我們再活在上帝的面前.
唯有榮耀的十字架才可以支付我們的罪債.
唯有聖靈的重生才可以使我們從罪中復活起來.
我想帶你去看看我們今天讀的聖經大圖畫.
我們讀了八章第十二節到第二十節.
但這次我想跟你說的時候.
用一個倒數的形式去看這段聖經.
由最尾開始.
因為最尾才告訴你耶穌說話的時候究竟在哪裡.
那個場景.
我們五四三二一這樣讀下來.
聖殿銀庫房.
時候還未到.

$^{361}$你父在哪裡.
你門不認識.
差我來的父與我同見證.
我自己作見證.
見證還是真的.
接著去到第十二節.
我們倒過來讀.
我是世界光.
跟從我的人.
可以得生命.
我想你留意第四點.
他們不信耶穌.
他們不信的是猶太的領袖.
於是他們質問.
你的父親在哪裡.
他們不認識耶穌.
他們不愛慕耶穌.
他們也不認識.
耶穌的天父.
耶和華上帝.
舊約裡面的上帝.
如果他們認識愛慕耶和華上帝.
他們很自然的會認識愛慕上帝獨一的兒子.
耶穌基督.
在這裡我們問究竟.
發生了些什麼事.
答案是聖經告訴我們.
這群罪人.
以色列的領袖.
他們黑了眼.
他們心又硬了.
於是他們拒絕耶穌.
他們的眼睛黑暗.
他們看不見耶穌是誰.
他們的心很硬.
他們不想馴服耶穌的權力之下.
於是他們拒絕耶穌.
但是耶穌卻深愛他們.
耶穌告訴他們.
我是世界的光.

$^{401}$耶穌最終會嘗試用他榮耀的十甲.
來溶化他們的心.
耶穌要用他榮耀的十甲.
去打開他們黑暗的眼睛.
我今天想和你們讀的其實只有兩節.
一個是第20節.
一個是第12節.
第20節.
整個故事的大大大大場景.
這些話是耶穌在聖殿的銀庫房裡.
教導人的時候說的.
聖殿的銀庫房.
即是猶太人耶路撒冷聖殿的銀庫房.
是耶穌宣告的背景和場景.
我們問為什麼耶穌要去宣告在銀庫房.
是不是和錢有關係.
不是.
其實很簡單.
就是因為那裡特別人多.
請學者告訴我們.
在那裡有十三個奉獻箱.
每個都有它特殊的用途.
而這些奉獻箱.
可以說是聖殿裡面婦人庭院.
Court of Women裡面的一個很重要的部分.
這個就是聖殿裡面的婦人庭院.
你看到銀庫在哪裡.
銀庫在這裡.
你剛才拍的照片你就知道了.
銀庫是有遮陰的地方.
他們去奉獻的.
這個大環境是Court of Women.
這個地方是以色列人男女老幼相聚的地方.
是整個聖殿的中心.
我可以說是整個聖殿最大的舞台.
這個就是Court of Women.
外面那裡只有外邦人的地方.
猶太人是不會停留在外面的.
猶太人全部進去.
這個是男女老幼可以去的地方.

$^{441}$再進去只有男人可以去.
只有祭司可以去.
於是你說哪裡人最多.
就是這裡.
於是耶穌選擇在那裡成為祂的一個大舞台.
然後在那裡.
祂還說我是世界的光.
於是在駐彭捷的最後一天.
耶穌進前來.
大聲的吶喊說.
我就是世界的光.
跟從我的人.
必不在黑暗裡面走.
卻要得著生命的光.
然後約翰又加了一句.
他說當時沒有人捉拿他.
因為他的時候還沒有到.
那個所謂的時候是十甲.
上十字架的時候.
還有六個月還沒有到.
榮耀十甲的時候還沒有到.
反過來說.
因為榮耀十甲還沒有到.
所以生命之光也還沒有到.
除非他上了十字架.
他沒有辦法成為生命之光.
世界之光的時候就更加不用說了.
還有一段時間.
為什麼?.
世界之光要等到耶穌上了十字架.
才有生命之光.
接著再等到耶穌出凶寒霧.
復活升天.
聖靈降臨.
這樣才有世界之光.
生命之光.
可以祂上了十字架就有.
可以救了一兩個.
但是要成為世界的光.
需要聖靈的降臨.

$^{481}$然後聖靈在他的門徒身上.
帶領門徒.
將生命之光傳到地極.
這樣就成為世界的光.
所以你開始明白.
整個故事的流程是什麼.
我們正在等待耶穌的十字架.
我們正在等待聖靈的降臨.
我想帶你去仔細看看.
耶穌的所在地.
約翰沒有跟我們介紹.
因為他假設我們都認識的.
他假設我們都是猶太人.
我想你仔細看看.
耶穌所在的地方.
升天的銀庫在哪裡.
我們再看看銀庫.
銀庫就在這裡.
這些有遮掩的地方就是銀庫.
但是這裡.
他的所在地.
也是婦人的庭院.
The Court of Women.
我想你留意.
在婦人的庭院裡面.
有什麼特別的地方.
你有沒有留意到.
婦人的庭院裡面.
有四條很高很高的支柱.
這個是約翰假設你會知道的東西.
聖經學者告訴我們.
四條很高的支柱.
大概是70~75呎高.
你想一下有多少個幽冥.
乘十左右吧.
十個NBA的玩家.
很高很高.
每一條柱子上面.
如果你留意一下.
有四個金所做的碗.

$^{521}$金碗.
每個金碗是用來裝橄欖油的.
每個金碗可以裝.
大概五加侖的橄欖油.
於是一共是十六個金碗.
加起來是八十個加侖的橄欖油.
這些所有的一切.
是要預備一個萬眾期待.
舉世矚目.
每年在駐彭捷連續六天晚上所舉行的火禮.
這個火禮就叫做.
聖殿點燈的照明禮.
所以你到了晚上就準備火禮了.
現在天氣變黑了.
就看著這四個地方去進行火禮.
連續六天的傍晚.
他們會先派四個年輕人.
你要爬上七十呎高的樓梯.
你不能找一些年紀大一點的.
要找四個年輕人拼命的.
上去聖殿那裡點燈.
如果大家知道.
點燈的時候.
那盞燈應該是需要有一個燈芯.
我不知道你知不知道.
如果有一個蠟燭應該有一條線.
原來他們告訴你.
當時拿來做那些金碗上面做芯.
有一個東西去燒著它.
那個燈芯.
他們是用祭司的舊衣服.
OK.
很有趣的.
祭司的舊衣服拿來做燈芯.
然後找四個年輕人爬上去.
點了那些核心.
於是有油在那裡.
就開始燃燒.
我想你記住這個圖畫.
甚至我想你記住那些仔細的點.

$^{561}$究竟用什麼做那個燈芯.
做那個蠟芯.
再說的是當時聖殿的外牆.
其實裝了很多黃金的金片.
於是你可以想像在晚上紅紅的烈火.
黃金反射著.
所帶出來的特別效果.
比海洋公園還是迪士尼樂園的晚宴好看.
有人說整個耶路撒冷.
因為耶路撒冷聖殿在山上.
整個耶路撒冷每一個角落在晚上.
就是那六天.
一年一度六天.
可以看到山上聖殿發光.
榮光.
這就是耶穌說我是世界的光.
生命的光宣告的場景和背景.
我們看不到的.
假設你是知道的.
不過.
醒學者告訴我們.
他相信耶穌不是在紅紅的烈火燒的時候.
大家在看火的時候.
去宣告我是世界的光.
因為這樣誰會聽他的.
他是等到第七天.
第七天火就熄滅了.
第七天是沒有火禮的.
在第七天大家說昨天好看.
或者第八天.
前兩天的火禮好看.
不過很慘.
今天就不能再看了.
在他的火禮已經完結的時候.
耶穌站起來說.
我是世界的光.
耶穌說我將要取締這個火禮.
如果你記得.
兩個星期之前.
我都說耶穌將水禮做完的時候.

$^{601}$再沒有水禮的時候.
你要等下一年的時候.
耶穌出來說.
我是生命的活水.
什麼意思.
你不用再追求這個水禮.
因為我會取締這個水禮.
同時之間耶穌說.
我將要取締這個火禮.
我要燃點耶路撒冷.
我要燃點耶路撒冷的聖殿.
如果你是一個基督徒.
你聽慣了我的講道.
你已經知道最後的謎底是什麼.
所有的事情其實是指向六個月之後.
耶穌將要在耶路撒冷聖殿以外的地方.
成為一個很巨大的火焰.
要照明整個世界.
耶穌說.
我就是.
我就是上帝在聖約中.
舊約聖經裡面獨有的名字.
我就是.
如果大家聽慣了我的講道.
你就知道.
我們有時候說舊約聖經的上帝叫耶和華.
這是一個pronunciation的問題.
這是一個你讀它的時候.
你讀它做耶和華.
但是它的意思.
耶和華的名字的意思就是「我就是」.
於是當耶穌說「我就是」的時候.
其實祂是在借用舊約聖經的上帝的名字.
祂說你知道我是誰.
我父差我來.
我與我父以合為一.
同為一神.
我就是舊約聖經裡面的耶和華上帝.
這是耶穌的第一句.
我就是.

$^{641}$接著祂說我就是世界的光.
接著祂又說我就是生命的光.
這裡我們有三個極其重要的字句.
我就是.
我們有世界光.
我們有生命光.
這三個詞句其實帶我們.
回到舊約聖經裡面.
很多很多的經典經文.
如果你想不到的話.
你想不起的時候.
是因為你不熟悉舊約聖經.
但是如果你讀過舊約聖經.
你記得舊約聖經裡面說「我是」.
「我上帝的名字」.
你記得舊約聖經裡面說「世界的光」.
你記得舊約聖經裡面說「生命的光」的時候.
你突然間腦海裡面爆了很多很多.
舊約聖經裡面的語言.
我今天.
時間關係.
想帶你回去舊約聖經的五本書裡面.
五卷的舊約聖經裡面有八段舊約的聖經.
我想帶你去看一下.
耶穌是怎樣承傳「我是世界生命光」的這個預言和應許.
而這個「我是世界生命光」的這個應許.
和他十字架和聖靈當臨有什麼關係.
不用怕.
因為時間關係.
我只可以很粗略地跟你閱讀那些經文.
然後很快幫你做一個總結.
我先帶你去看第一本書.
第一本書很簡單.
第一本.
《創世記》.
《創世記》一開播就說「光」.
「起初上帝創造天地」第一章第一節.
「地是空虛混沌.
深淵上面一片黑暗.
上帝的靈運行在水面上.

$^{681}$上帝說要有光就有光.
上帝看光是好的.
於是上帝就把光與暗分開.
上帝稱光為晝 稱暗為夜.
光是上帝首要的創造.
光是上帝本性的彰顯.
你想到上帝就是光.
上帝不會是黑暗.
黑暗是撒旦.
撒旦是魔鬼.
上帝是光.
光就是上帝的本質.
光也是生命的傳源.
如果你讀下去.
你知道繼續下去.
我們沒有時間去看.
開始做植物.
接著有動物.
接著有人.
所有的東西.
今天科學會告訴你.
與光是有關係的.
耶穌說要承傳光照黑暗.
生命誕生的應許和預言.
你可不可以這樣想.
第一本書.
現在帶你去看第二本.
第二本書是《出埃及記》.
《出埃及記》裡面一定有三段聖經.
與光有關係.
我必須要帶你去看.
《出埃及記》第二本書的第一段聖經.
第三章第二節.
耶和華的使者在荊棘的火焰中.
向他顯現他是摩西.
摩西觀寒寒.
荊棘在火中焚燒.
卻沒有燒盡.
Burning bush大家都熟悉.
摩西說.

$^{721}$他們若說他叫什麼名字.
上帝說我差遣你去拯救以色列人出埃及.
我要對他們說什麼.
上帝說這是我的名字.
我是自有永有.
那自有永有的差派我到你們那裡去.
這裡自有永有是上帝的名字.
他代表了我是 我就是.
I am.
這個上帝的名字.
大家可不可以明白.
自我永有其實是一個意思的翻譯.
它不是原本的意思.
如果你直接看希伯來文.
你說那句算什麼.
我是.
上帝的名字叫我是.
但是因為我們這樣聽我是好像聽不下去.
於是我們將它異譯.
變成自我 自有 永有的上帝.
當然這個是真實的.
但是其實那個原文是我是.
上帝也說上帝是我是.
在火焰當中啟示自己.
我們一想光的時候.
你很自然想著按個電掣開燈.
但是你要想一下在耶穌的時候.
人們一想光的時候是和火有關係的.
當你回到舊有聖經你不是想光.
你是想火.
火和光是同一件事來的.
耶穌來是要承傳火焰光中.
上帝臨在的應許和預言.
第二段《猜疑及記》的聖經.
又說走在他們前面.
日間用雲柱引領他們的路.
夜間用火柱照亮他們.
使他們日夜都可以行走.
日間的雲柱 夜間的火柱.
從不離開百姓的面前.

$^{761}$第十四章第二十四節.
出埃及的時候.
破曉的時分.
耶和華從雲柱火柱中瞭望.
使埃及的軍隊混亂.
埃及人的軍隊在追著他.
他們剛剛過紅海.
還在追著來.
上帝怎樣去營救他的子民.
通過他的火柱.
去使他們混亂.
又說上帝在黑夜中.
以火柱臨在.
照亮引導他的子民.
並且以火柱保護他們.
替他們抗敵.
耶穌來為要承傳.
耶和華上帝.
以火柱臨在照亮引導.
保護抗敵的應許和預言.
我知道有點複雜.
但你明白你在讀舊約的東西.
其實所有的東西.
耶穌來都要承傳的.
於是我們在讀的是舊約的聖經.
但我們在想的是新約的應許.
第三段.
埃及記第四十章第三十四節.
那時雲彩遮蓋回眸.
耶和華的榮光充滿著回眸.
摩西不能進入回眸.
因為雲彩停在旗上.
耶和華的榮光充滿著帳幕.
第三十八節.
在他們所行的路上.
在以色列全家的眼前.
白天耶和華的雲彩在帳幕上.
黑夜有火在雲彩的中間.
耶和華上帝在帳幕完成.
起初只是雲柱.

$^{801}$只是火柱.
現在有帳幕了.
當帳幕完成的時候.
雲彩的榮光來到充滿著帳幕.
從前的雲柱火柱.
和現在耶和華新造的帳幕合成一體.
我想你開始明白.
上帝一直在等待這個帳幕的完成.
因為這個帳幕是代表了.
祂和祂子民一同在的一個地方.
是耶和華上帝的家.
於是現在雲柱和帳幕成為一體.
臨在照亮引導保護抗敵.
耶穌來要成全什麼.
又說上帝榮耀充滿帳幕的預言和應許.
你聽不明白不要緊.
到時候我最後幫你總結的時候.
你會更加聽得明白.
我們看了兩本書.
《創世記》和《猜一級》.
現在我們看第三本書.
第三本書是《歷代志下》.
不用怕.
這段聖經大家一定熟悉的.
雖然這本書大家是不熟悉的.
《歷代志下》第七章第一節.
所羅門祈禱完畢.
就有火從天降下來.
燒盡繁祭與祭物.
耶和華的榮光充滿著聖殿.
現在再不是那個帳幕了.
是聖殿了.
是所羅門所造的聖殿了.
那火降下第七章第三節.
耶和華的榮光在殿上的時候.
以色列眾人看見.
就在石板上俯伏敬拜.
稱謝耶和華.
耶和華本為善.
他的慈愛永遠長存.

$^{841}$如果大家熟悉舊的聖經就知道.
起初上帝帶他們出埃及.
是先火柱.
接著他們造了一個帳幕.
於是火柱和帳幕合而為一.
但是火柱和帳幕只是在曠野裡面造的.
現在上帝的子民已經到達了應許之地.
當時臨時的帳幕.
現在可以變成一個永久的聖殿.
於是他們說.
看啊.
榮光充滿著聖殿.
接著眾人看見這個榮光.
就俯伏敬拜這個聖殿.
他們稱讚耶和華.
稱讚耶和華是什麼呢?.
耶和華本為善.
是什麼意思呢?.
什麼耶和華本為善呢?.
不是 如果你看過希伯來文.
耶和華是good.
good是代表什麼?.
當上帝創造天地的時候.
連續七次說.
每一天都是good.
所以當我們說耶和華本為善.
當我們說耶和華是很好的時候.
我們在說耶和華做了一個新的創造.
你記得起初我們每個人都住在一個完美的創造裡面.
起初我們沒有一個人會死會病.
但是我們因為亞當的墮落.
世人都犯了罪.
虧損了上帝的榮耀.
最低功肯乃至死.
我們現在回不去之前那個很好的創造.
現在上帝的聖殿完畢了.
上帝的聖殿是什麼?.
上帝的聖殿是上帝很好的創造.
我現在沒有時間去說究竟聖殿是什麼.
但是我想告訴你.

$^{881}$聖殿其實是代表了新創造.
聖殿是新創造的入口.
當我們去敬拜上帝.
來到上帝為我們預備的聖殿.
是通過這個聖殿我們進入了一個新的世界.
保羅會說.
若有人在基督裡新的創造.
基督是那個新的聖殿.
我們可以進入這個聖殿.
我們可以進入那個完美的世界.
再沒有死亡.
再沒有罪惡的世界.
於是在這個榮耀的聖殿裡面.
耶和華的慈愛永遠長存.
不是一刻長存.
不是一刻短暫的存在.
是永遠長存.
這裡說的慈愛是祂的救贖.
Hazard.
又說的Hazard.
又說的Emmet.
祂的慈愛,祂的信實.
在聖殿裡面永遠長存.
耶穌來要承傳耶和華上帝.
榮耀聖殿的應許和預言.
那個聖殿 耶穌帶給我們將來的聖殿.
是一個很好的新創造.
耶穌要給我們的聖殿.
是耶和華上帝慈愛,信實,救贖的所在.
當我們看見耶穌的新聖殿的時候.
我們應該怎樣?.
俯伏敬拜.
永遠的敬拜.
OK.
好 我們看了三本書.
我們現在看第四本書.
和光有關係的.
和世界的光有關係.
第四本書是《以賽亞書》.
我們要看兩段《以賽亞書》.

$^{921}$第一段《以賽亞書》第九章.
我想如果你在教會內.
你會發現每個聖誕節.
我們基本上都會讀這段聖經的.
「外邦人居住得加利利地.
在黑暗中行走的百姓看見了一個大光.
住在死任之地的人有光照耀他們.
For unto us a child is born.
因有一英孩為我們而生.
A son was given.
有一子賜給我們.
Hand of Messiah」.
如果你去聽的時候.
「政權必擔在他的肩頭上.
他的名稱為奇妙賊士.
全能的上帝.
永在的父和平的君.
他的政權和平安必加增無窮」.
我想你看看這裡.
這裡說死任之地.
將會有一個光,大光去照耀他們.
這個死任之地是什麼地方?.
上面告訴我們是加利利.
誰娶加利利出來的?.
耶穌.
加利利是耶穌的家鄉.
加利耶穌開始福音事供的地方.
在加利利死任之地.
有一個很大的光明可以照耀著.
接著《以賽亞》有一個預言.
他說「For unto us a child is born.
there is a child that is born for us」.
我想你留意一件事.
當這個嬰孩生下來.
當這個嬰孩長大以後.
要反過來.
是我們要為他而活.
他為你而生.
他長大了之後.
是你為他而活.

$^{961}$你要俯伏敬拜他.
你要為他而活.
為什麼?.
因為他必在大位的寶座上.
治理他的國.
從今直到永遠.
這個我沒有打出來.
是第九章第七節的延續.
因為這個孩子長大了之後.
他會成為上帝應許的君王.
永遠的君王.
他要審判這個世界.
他要拯救這個世界.
耶穌來為要承傳應許君王.
永遠君王的應許和預言.
耶穌不單止是以色列人的君王.
耶穌也是死任之地.
外邦人的明光救主.
耶穌是全世界的救主.
這是第一段的《以賽亞書》.
我想帶你去看第二段的《以賽亞書》.
你更加可以看到.
耶穌是全人類的救主.
第四十二章第六節.
「我」耶和華「憑公義召你.
要藉著你和百姓立約.
使你成為萬邦之光」.
The light to the nations 萬邦之光.
「開盲人的眼.
領囚犯出監獄」.
不是你想的那些監獄.
「領坐在黑暗中的出地牢」.
這個才是一個屬靈的監獄.
活在罪惡和黑暗當中的人.
可以通過萬邦的光.
開眼睛走出監獄.
第四十九章第六節.
「你作我的僕人.
使雅各眾之派復興.
使以色列歸回.

$^{1001}$然而此事尚小」.
還不夠大.
「我還要使你作萬邦的光.
使你施行我的救恩直到地極」.
這是耶穌說「我是世界的光」的背景.
萬邦的光出現過兩次.
萬邦的光要使雅各以色列人復興.
使以色列人歸回.
但是還不夠.
他要施行救恩直到地極.
什麼是救恩.
救恩就是「開盲人的眼睛.
讓囚犯出監獄」.
是屬靈的監獄.
屬靈的眼睛.
「令坐在黑暗中的人出地牢」.
有沒有想過拒絕耶穌的猶太人.
基本上就是坐在這個監獄裡面.
因為他們核眼.
因為他們心死.
因為他們坐在監獄裡面.
他們拒絕耶穌.
他們要繼續活在黑暗裡面.
他要活在他們最終.
他不願意悔改.
但是耶穌卻不斷地向他們宣告.
希望用祂的光去照明他們黑暗的地方.
耶穌來.
是要承傳萬邦之光的應許和預言.
耶穌要使雅各復興.
要救贖萬邦.
耶穌要開黑暗人的眼睛.
耶穌要使罪人認識祂.
服侍祂.
愛慕祂.
獨一的真神 獨一的救主.
最後一本書.
第五本.
我們讀過《創世記》.
我們讀過《出埃及記》.

$^{1041}$我們讀過《歷代治下》.
我們讀過《二才亞書》.
最後一本.
是《撒迦利亞書》.
小書.
但是很難讀.
但是不要緊.
兩個星期之前我們讀過的.
這段聖經.
信不信由你.
《撒迦利亞書》第十四章第五節.
「我的上帝必降臨」.
這是一個末後的意象.
將來的意象.
「所有的聖者與你同來.
所知的那一天.
沒有白天沒有黑夜.
到了晚上仍有亮光」.
接著那兩節我們讀過的.
兩個星期之前.
「在那一天必有活水從耶路撒冷出來.
往東海流往西海流.
東下都是如此.
也或說要作全地的王.
哪一天也或說必為獨一無二」.
在這裡你看到.
生命活水和世界的光.
兩個末世意象.
二合為一.
是不是.
所以我沒有騙你.
3A 3B還記不記得.
3A 3B都是三來的.
它們不是兩個故事.
是一個故事.
不是兩個意象.
是一個意象.
你去看《撒迦利亞書》.
撒迦利亞一直跟你說.
世界的光和生命活水.

$^{1081}$是同一樣東西來的.
前面的部分.
末世的意象.
《撒迦利亞書》第十四章.
第五節 第七節.
「再沒有白天黑夜」.
意思是什麼.
有一個持續的亮光.
照明整個世界.
後面 Part B.
說將來會有一個活水長流.
在這個活水長流的時候.
活水去到哪裡.
光去到哪裡.
耶和華要作全地的王.
全地.
以前是黑的.
現在被照明了.
全地以前是死的.
現在有活水長流.
於是活水到的地方.
光到的地方.
會有人俯伏敬拜耶和華.
讓耶和華作全地的王.
怎樣俯伏敬拜.
耶和華必為獨一無二.
你記不記得我為什麼.
在兩個禮拜之前.
耶和華是獨一的 Shema o Israel.
以色列最大的解密.
以色列的上帝和華上帝.
是獨一的上帝.
你要盡心盡性盡意愛祂.
當我們說耶和華是獨一的時候.
我們說耶和華將有一班.
盡心竭力敬拜祂的人.
這個就是光到的地方.
水到的地方.
必有真正敬拜耶和華.
盡心盡性盡意愛祂.

$^{1121}$好似阿伯般敬拜者的人.
耶穌來.
耶穌是世界的光.
耶穌來就要承傳這個.
亮光照黑暗的應許和預言.
耶和華要做全地的王.
通過耶穌.
耶和華要做全地的王.
要領受獨一盡心的敬拜.
我用一張圖表.
可以總結我們前面那五本書.
我們所看過的舊約預言.
我這樣說.
起初創天地.
有光有生命.
創世記.
我師耶和華.
榮光象脈行.
採求其.
榮光滿聖殿.
現在再不是象脈.
是聖殿.
慈愛永長存.
歷代志下.
光照死蔭地.
瀑人萬邦光.
耳採阿書.
最後.
明光照黑夜.
盡心愛上帝.
剎加利亞說.
所有舊約的預言和應許.
大家記得.
都需要榮耀的十價.
才可以成全.
你不要這樣讀完之後.
上帝動手.
需要榮耀的十價.
所有的預言.
所有的應許.

$^{1161}$也都需要聖靈的降臨.
才可以完成這兩樣東西.
是必須要的.
然後耶穌才可以成為.
世界的光.
好了.
差不多要結束了.
在我結束的時候.
之前.
我想帶你再回去.
約翰福音.
這次我們看約翰福音.
我們不看舊約.
我們看新約.
我想你看一下.
耶穌的榮耀十字架.
和聖靈的降臨.
是如何會成全.
這五個舊約書卷裡面的預言和應許.
OK.
首先.
我們看.
創世記裡面的預言和應許.
起初創天地.
有光有生命.
你去到約翰福音的時候.
都是跟你說光和生命.
但是有些不一樣的了.
約翰福音第一章第一節.
大家都很熟的了.
「太初有道.
道與上帝同在.
道就是上帝.
萬物都是藉著他做的」.
提醒你之前發生了什麼事.
但是現在他要再做另一個新的創造.
第一章第四節.
「在他裡面有生命.
這生命本是人的光.
現在光照在黑暗裡面.

$^{1201}$再次照.
第二次照.
和之前不一樣的.
那光是真光.
來到世上照亮所有的人」.
正如我之前有說過.
這裡的黑暗和光照黑暗.
和創世記第一章的黑暗是不一樣的.
這裡的黑暗是叛逆罪孽裡面的黑暗.
你要照亮這個黑暗.
你需要榮耀的十字架.
才可以拯救最終的黑暗.
你也需要有聖靈的降臨.
才可以重生這些死在最終的罪人.
這是耶穌怎樣承傳創世記.
然後是出埃及記的預言和應許.
「我師耶和華.
榮光將莫行」.
在出埃及記裡面.
上帝在帳幕當中與上帝的子民同行.
我想看看耶穌怎樣承傳.
因為福音第一章第十四節這樣說.
「道成了肉身」.
現在耶穌親自降臨在這個世界.
帳幕在我們中間.
我知道你的反應是注重不當.
但是你看到的原文是.
帳幕在我們當中.
祂充充滿滿地有慈愛和信仇.
然後下面還說.
我們有見過祂的榮光.
這是一個榮耀的帳篷.
耶穌來到這個世界.
耶穌成為了一個榮耀的帳篷.
「道成了肉身」.
帳幕在我們當中.
這個榮耀的帳幕要帶領上帝的子民.
前往應許之地.
我想問問你的應許之地在哪裡.
再不是耶路撒冷.

$^{1241}$再不是中東某一個聖地.
是庫那裡去.
於是第十四章第六節這樣說.
耶穌說「我就是道路真理.
生命若不藉著我.
沒有人能到庫那裡去」.
耶穌要帶我們通過這個帳篷.
通過祂自己.
帶我們去庫那裡.
庫那裡是怎麼去的.
是通過十字架的道路.
那條道路是十字架的道路.
榮耀十架才有道路.
唯有聖靈降臨才有生命.
榮耀十架是那條道路.
聖靈降臨是賜我們生命.
帶我們去庫那裡去.
OK 時間關係.
我想帶你去看看第三.
耶穌是怎樣承傳歷代志下的預言和應許.
歷代志下是什麼.
聖殿.
榮耀充滿於聖殿.
然後慈愛永長存.
我們看一節而已.
《約翰福音》第二章第十九節.
那個帳篷現在變成了一個永遠的聖殿.
你們拆毀這殿.
耶穌說我三日內要把它重建.
在歷代志下裡面.
我們見到一個充滿榮耀.
慈愛長存的聖殿.
不過舊約的聖殿只是一個影子.
唯有替我們受死復活的耶穌基督裡面.
在耶穌基督裡面.
在耶穌基督裡面有上帝的羔羊.
上帝的羔羊給我們有血和聖靈的洗.
我們才可以永享上帝的慈愛和信實.
上帝的慈愛信實盡在耶穌基督榮耀的十字架裡面.
那個才是永遠的聖殿.

$^{1281}$於是耶穌說拆毀這聖殿.
Physical Temple.
我三日內要把它重建.
耶穌在說的是指著祂的身體.
祂的受死與祂的復活.
於是在祂的受死與復活裡面.
我們找到一個永遠有上帝慈愛信實的聖殿.
接著以賽亞書的預言.
以賽亞書說光照在死刃地.
愛邦人的地方會見到一個大光.
服人成為萬邦的光.
我想帶你去看我們很熟悉的.
撒瑪利亞人婦人的故事.
因為這個光現在不單是在以色列人的地方.
是去到愛邦人甚至雜種人的地方.
耶穌的光照明撒瑪利亞的婦人.
耶穌叫她找她的老公來.
她說我沒有老公.
耶穌這樣說.
你已經有過五個丈夫.
現在跟你一起的不是你的丈夫.
這是耶穌對她的罪孽的照明.
然而他們沒有拒絕耶穌.
耶穌在那裡再留了兩天.
接著通過耶穌的宣講.
現在撒瑪利亞人的承認.
他們口裡的承認.
我們信因為我們親自聽見了.
知道這人真是世界的救主.
死刃之地再不單是加利利愛邦人的地方.
連撒瑪利亞雜種的地方.
都得到耶穌的明光照耀.
於是在撒瑪利亞那裡.
如果你記得不單止有明光照耀.
還有耶穌活水聖靈的認許.
因為耶穌和那個女人經常在說活水.
耶穌給他們的何止是明光照耀.
耶穌還給了他們活水傳言.
通過榮耀的十誡.
通過聖靈的降臨.

$^{1321}$耶穌將要成為一個名副其實的世界救主.
最後撒迦利亞的預言.
撒迦利亞說有明光會照黑夜.
再沒有晚上早上.
這個明光會一日二十四小時照著你.
而照出來的結果是甚麼呢.
就是你會盡心愛上帝.
我想帶你去看尼哥德姆的故事.
有一個法利賽人大家都很熟悉的.
名叫尼哥德姆.
他是猶太人的官.
他是猶太人的神學家.
他是一個道德人.
他是一個律法者.
願意他在黑夜裡面來找耶穌.
耶穌告訴他.
你不要以為你可以靠你自己的好行為得救.
耶穌說你們必須要重生.
光照在這個黑暗者的身上.
尼哥德姆最後如何得到耶穌的重生.
答案有兩邊.
一個他仰望被舉起來的人子.
另外他領受了血和聖靈的洗禮.
於是那一天.
明光照黑夜.
尼哥德姆從此盡心愛上帝.
如果大家記得他離開了黑暗.
他追隨了真光.
他奉上他所有.
他戴的乳香抹藥.
是75年人工所可以換取價值高昂的東西.
全部敬拜.
擺在耶穌的腳前.
然後他自己拋棄他的前途.
他追求耶穌的天國.
這個是活水和真光帶給我們的生命.
我想和你一起安靜在上帝面前.
去俯伏敬拜我們獨一的救主.
世界的真光.
為了拯救罪人的緣故.

$^{1361}$真光耶穌被黑暗所包圍.
被黑暗所勝過.
你記得馬可福音.
怎樣說耶穌被釘十字架那一刻.
第十五章第三十三節.
到了正午.
全地都黑暗了.
真光被黑暗所包圍.
被黑暗所勝過.
直到下午三點.
三點鐘耶穌大聲呼喊著.
我的上帝.
為什麼要離棄我.
馬可福音第十九章.
講到耶穌被釘在十字架上那個時刻.
他這樣說 他有一件內衣是沒有縫的.
誰穿無縫的內衣.
是祭司.
不要撕開.
他有一件內衣是祭司的內衣.
不要將它撕開.
可以拿來做燈芯.
可以拿來做蠟芯.
可以燃點.
燒盡上帝的憤怒.
讓它成為拯救世界的真光.
Do you understand what is happening?.
那件衣服使它燃點.
被上帝的憤怒所燃點.
它要燒盡上帝的憤怒.
以致它可以成為拯救世界的真光.
於是他說我喝了.
耶穌又說.
成了.
他就低下頭來.
斷了氣.
於是五十天之後.
紹行傳第二章第一節至第四節.
五巡節.
五巡節就是五十天之後的意思.

$^{1401}$到了.
一陣大風吹過.
聖靈來了.
充滿了他們所作的整個屋子.
又有蝕頭魚火焰火焰來了.
分開落在他們每個人的身上.
他們都被聖靈充滿.
我們說的是耶穌的門徒.
所賜的口才.
他們說起別國的話來了.
有上帝的火焰榮耀降臨.
臨到新約的教會裡面.
臨到基督的身體末世的聖殿.
你明白嗎.
那個聖殿已經再不重要.
耶路撒冷的聖殿已經再不重要.
在耶穌的受洗活裡面.
信耶穌的人成為了那個新的聖殿.
你看到上帝的榮耀.
現在降臨在他們的身上.
有火有光降臨在.
新的末世教會的身上.
他們都被聖靈充滿.
說起別國的話來.
什麼意思.
他們將要成為萬邦的光.
世界的光.
不是有口才.
可以做多幾單生意.
懂多幾個英文廣東話普通話.
做多些生意.
NO.
為什麼他們可以說起別國的話.
因為他們不單止要成為生命的光.
他們要成為萬國生命的光.
聖靈降臨在你們的身上.
你們必要得著能力.
要在耶路撒冷猶太全地.
撒瑪利亞直到地極.
作我們的見證.

$^{1441}$最後一節一定要和你分享.
馬太福音第五章第十四節.
耶穌這樣說.
你們是世界的光.
承造在山上.
a city on a hill.
令你想起什麼耶路撒冷.
a city on a hill.
耶路撒冷就是一個城在山上.
你們是世上的光.
或者你們是世界的光.
你們是新耶路撒冷.
你們將取替現在的耶路撒冷.
你們是世上的光.
承造在山上.
a city on a hill.
是不能隱藏的.
到第十六節給你們聽.
你們的光也要這樣照在人前.
叫他們看見你們的好行為.
把榮耀歸給你們天上的父.
我們現在聽到這件事.
好行為是我們應該要做好人.
是我們應該工作卓越.
但你忘記了好.
是上帝獨有的形容詞.
好行為是上帝的行為.
好行為是上帝的作為.
叫你去好行為是叫你做上帝的旨意.
而上帝的旨意就是我們不為這個世界賣命.
而是我們為上帝盡忠.
Two cities have been formed by two loves.
the earthly by the love of self.
and the heavenly by the love of God.
奧古斯丁的名言.
有兩種愛造就了兩個都市.
一個是因愛自己而形成的世俗之城.
另一個是因愛上帝而形成的天國之都.
You are the light of this world.
你是世上的光.

$^{1481}$City on the hill.
你們要做的是上帝的旨意.
不要為世界賣命.
卻要盡忠天國.
我們天父多謝你.
這是一個不容易宣講的訊息.
然而我希望當大家想通想透.
所有的圖畫與connection的時候.
整個福音.
並且福音的形容.
都展露在我們的眼前.
多謝你給我們1189張聖經.
或者我們真的覺得嫌多.
真的覺得多了幾本書.
真的覺得為什麼要這麼厚.
然而我們回去看每一本書裡面.
都充滿著福音的影子.
耶穌是生命的光.
耶穌是世界的光.
原來這些話語.
全部都已經在舊約聖經裡面.
映射預言著.
當耶穌宣講這句話的時候.
祂是帶我們回到舊約裡面.
很多關於祂的預言.
在十字架上.
在聖靈的降臨裡面.
祂將要承傳世界的光.
十字架成為生命的光.
聖靈的來臨.
將生命的光變成世界的光.
求主你幫助我們.
現在輪到我們.
與耶穌一起.
因為你們是世界的光.
承造在山上.
A city on a hill.
求主讓我們用我們的生命.
我們對天國的盡忠.
成為了世界的光.

$^{1521}$祝福我們的教會使用我們.
我們敢祈禱 奉主耶穌基督的名禱.
\newpage



\section{詩篇 29:1-11}
\label{sec:032223}
\textbf{風暴之上的平安:化解焦慮的三個行動 (華)}~~~~卓爾君牧師
\newline
\newline
連結: \href{http://www.church.com.hk/acms/content.asp?site=cdc\&op=show\&type=product\&code=032223\&layout=sermon}{\texttt{032223}} ~~~~ 語音日期: 2026-05-31
\newline
\newline
\hyperref[sec:032222]{\small{< < < PREV SERMON < < <}}
~
\hyperref[sec:index_chronic]{\small{[返順時目]}}
~
\hyperref[sec:032224]{\small{> > > NEXT SERMON > > >}}
\newline
\newline
詩篇 29:1-11
\newline
\begin{longtable}{cl}
\hline
\hline
章節 & 經文 (和合本修訂版)\\
\hline
29:1 & \begin{tabularx}{0.7\textwidth}{X} 神的子民哪，你們要將榮耀、能力歸給耶和華，都歸給耶和華！ \end{tabularx} \\ \\ \relax
29:2 & \begin{tabularx}{0.7\textwidth}{X} 要將耶和華的名的榮耀歸給他，要敬拜神聖榮耀的耶和華。 \end{tabularx} \\ \\ \relax
29:3 & \begin{tabularx}{0.7\textwidth}{X} 耶和華的聲音在眾水上，榮耀的神打雷；耶和華打雷在大水之上。 \end{tabularx} \\ \\ \relax
29:4 & \begin{tabularx}{0.7\textwidth}{X} 耶和華的聲音大有能力，耶和華的聲音滿有威嚴。 \end{tabularx} \\ \\ \relax
29:5 & \begin{tabularx}{0.7\textwidth}{X} 耶和華的聲音震碎香柏樹，耶和華震碎黎巴嫩的香柏樹。 \end{tabularx} \\ \\ \relax
29:6 & \begin{tabularx}{0.7\textwidth}{X} 他使黎巴嫩跳躍如牛犢，使西連跳躍如野牛犢。 \end{tabularx} \\ \\ \relax
29:7 & \begin{tabularx}{0.7\textwidth}{X} 耶和華的聲音使火焰分岔。 \end{tabularx} \\ \\ \relax
29:8 & \begin{tabularx}{0.7\textwidth}{X} 耶和華的聲音震動曠野，耶和華震動加低斯的曠野。 \end{tabularx} \\ \\ \relax
29:9 & \begin{tabularx}{0.7\textwidth}{X} 耶和華的聲音驚動母鹿落胎，樹林也脫落淨光。凡在他殿中的，都述說他的榮耀。 \end{tabularx} \\ \\ \relax
29:10 & \begin{tabularx}{0.7\textwidth}{X} 耶和華坐在洪水之上為王；耶和華坐著為王，直到永遠。 \end{tabularx} \\ \\ \relax
29:11 & \begin{tabularx}{0.7\textwidth}{X} 耶和華必賜力量給他的百姓，耶和華必賜平安的福給他的百姓。 \end{tabularx} \\ \\
[1ex]
\hline
\hline
\end{longtable}
$^{1}$最終準備平安.
標桿人生的作者 華里克牧師.
他對於Life 生活.
他有一個非常生動的猜解.
他說Life L-I-F-E.
開頭是Learning 學習或教育.
教育可以提升我們的生活.
我想這是我們華人.
最能夠深刻體會的.
我不知道你有沒有遇過那種同學.
他在考試之前說.
哎 我都沒看 怎麼辦.
結果他考出來全班第一.
跟你旁邊的人說.
你是不是那種人啊.
很討厭啊.
但你又拿他沒辦法.
我們有一個專有名詞叫學霸.
所以我們當中應該有很多學霸的.
但即使不是學霸.
Life的結尾還有一個E.
就是Experience 經驗.
職場的經驗 服事的經驗 人生的經驗.
這些經驗同樣可以將我們的生活.
推向更高處.
可是當我們既有了L也有E.
人生是否就高枕無憂呢.
別忘了一個殘酷的事實.
在Life的正中間.
還有I F If 佔據了一半.
這個If字面上是如果.
但是在現實生活中它叫做萬一.
各位可以回想一下.
那些熬夜加班的夜晚.
那個家裡櫃子裡面一整排的維他命.
還有那個為孩子申請大學.
傷透腦筋的策略規劃.
我們做足了這一切準備.
都是在防備人生當中的那個If.
萬一.

$^{41}$但現實往往是就在我們覺得一切順理成章.
覺得那個If根本不可能發生的時候.
它就發生了.
可能是半夜突然響起了手機鈴聲.
可能是醫生面色凌重.
遞來了一份體檢報告.
甚至有可能是孩子突然開口說.
我不要再去教會了.
弟兄姊妹朋友們.
不管我們的教育有多高.
經驗多豐富.
不管我們的人生走了多遠.
仍然有If.
那個萬一的風暴會來臨.
這個就是生活 Life.
今天我們要來讀一首在風暴當中.
誕生的詩篇.
當作者大衛寫下這段文字.
他的生命正遭遇一場猛烈的風暴.
但是大衛的反應出人意料.
他在風暴中所感受到不是毀滅的恐懼.
而是聽見神充滿榮耀的歌聲.
所以接下來我們要透過大衛的親身經歷.
來學習如何在突然的風暴當中.
找回那份屬天的平安.
有三個方面.
第一個方面.
用敬拜的聲音吞沒焦慮.
跟你旁邊的講.
用敬拜的聲音吞沒焦慮.
弟兄姊妹當你眼前突然閃了一下.
然後聽見巨大的雷聲.
你有什麼感覺.
在這種可怕的自然力量展現出來的時候.
內心會不會有一種敬畏.
大雷的聲音轟隆一聲巨響.
那個聲音大到什麼程度.
大到可以把客廳窗戶震得嘎嘎響.
然後街上的汽車警報器就開始亂叫.
但是就在那一瞬間.

$^{81}$你有沒有發現一件非常奇妙的事情.
就在那轟隆一聲炸開的那一秒.
世界上其他的聲音就全部消失了.
你原本開得很大聲的電視.
聽不見了.
原本你的配偶正在睡睡念.
聽不見了.
甚至你內心那個煩躁的聲音.
也全部被無情的蓋過去.
所以在當下你能聽到的.
就只有讓人從心裡深深敬畏的那個雷聲.
而這個就是大衛在詩篇29篇的開始.
他要為每一位陷入焦慮的人所描繪的畫面.
他要向我們展示一個絕對壓制的屬靈力量.
讓我們一起來大聲念這一段經文.
大家一起來念.
要大聲.
如果太小聲我要叫你們再念一次.
大聲一起請.
神的眾子啊.
你們要將榮耀能力歸給耶華.
歸給耶華.
要將耶華的民所當得的榮耀歸給他.
以聖潔的裝飾敬拜耶華.
耶華的聲音發在水上.
榮耀的神打雷.
耶華打雷在大水之上.
耶華的聲音大有能力.
耶華的聲音滿有威嚴.
這首詩篇是一篇極具戰鬥性的屬靈宣告.
因為在大衛的時代.
當時的人普遍敬拜一種很強勢的偶像.
名叫巴利.
在迦南的神話裡面.
巴利被尊為風暴之神.
人生性巴利.
掌管天上的雷電,狂風.
然後帶來豐沛的雨水.
但大衛在這裡宣告說.
不!你們都搞錯了.

$^{121}$在天上打來的不是什麼巴利.
那個能夠掌控風暴不是那些虛假的偶像.
而是耶和華.
是耶和華大能的聲音.
運行在那個象徵死亡跟混亂的大水之上.
在詩篇29篇當中.
有幾個字出現了7次.
就是耶和華的聲音.
意思是我們的神不需要像民間宗教的偶像.
還需要拿出什麼法寶道具.
還要這個差遣天兵天將來收腰.
更別說我們常常在新聞裡面看見.
當在大火大水當中.
這些偶像還要人去搶救的.
我們的神只是高高坐在.
那個榮耀的寶座上.
發出他的聲音.
那個帶有絕對權柄震動天地的雷聲.
就足以瞬間鎮壓宇宙間所有的混亂跟恐慌.
神的聲音就像是在演唱會裡面.
當那個最強的重低音一吹下去的時候.
你已經不是單純在用耳朵聽了.
你甚至能夠感覺那個全身的骨頭.
在那個強大的頻率當中就微微震動.
當那個最具統治力的聲音一出現.
全場那些尖銳的高頻的噪音就完全聽不見了.
弟兄姊妹,這個就是神的聲音所帶來的絕對威力.
但是現實中的我們卻常常讓另外一種聲音.
佔據了我們的腦海.
就是焦慮的聲音.
焦慮就是那個心裡那些不受控的雜訊.
我們腦海裡面那些焦慮的聲音.
不像打雷那樣震耳欲聾的.
但就像夏天夜裡面關燈之後.
一直在你耳邊嗡嗡作響的那個蒼蠅蚊子.
你揮手打不到它,開燈也看不到它.
但是那種尖銳突然而來的高頻噪音.
有一種可怕的魔力.
它能夠讓你徹底失眠到天亮.
其實仇敵精心設計的焦慮就像那樣的.

$^{161}$事實上你知道嗎?.
聖經所記載的鬼王叫別西普.
他就是蒼蠅造型的.
他無孔不入的.
所以整天就在我們的腦海尖叫.
萬一這一波財源名單有我.
房貸怎麼辦?.
萬一股市澎湃.
我辛苦存的退休金怎麼辦?.
萬一體檢報告的紅字降不下來怎麼辦?.
萬一我找不到理想的對象怎麼辦?.
其實這些的聲音日復一日就在我們的心裡散播著.
你不夠好.
你的處境非常危險.
你的人生快要完蛋的那種焦慮.
而這種靈魂的噪音污染會一點一滴的侵蝕我們的平安.
更糟糕的是現代人處理這種噪音的方法.
往往是選擇用另外一種噪音來蓋過去.
我們就瘋狂地滑著短視頻.
我們毫無節制地熬夜追劇.
或是大把時間投在遊戲當中打怪.
或是在網路上面不斷地下單買東西.
我們試圖用強烈的視覺跟聽覺的刺激.
來暫時麻痺那顆快要超載崩潰的心.
當地兄姊妹有用嗎?.
當你拔下耳機關掉屏幕的那一瞬間.
那隻焦慮的蚊子是不是立刻又飛到你的耳邊了?.
另外一方面.
也有人試圖透過放空讓自己心無掛慮.
教會有一首經典詩歌叫《安靜》.
有沒有唱過《安靜》.
但其實很多人誤解了聖經當中.
安靜跟休息的真正含義.
結果《安靜》唱著唱著在現實生活就變成踏平.
其實這一首詩歌出於詩篇46篇.
它根本不是教導我們雙手一攤什麼都不做的消極躺平.
詩篇46篇徹頭徹尾是一首充滿張力的戰鬥詩篇.
而第十節的下半句.
更把我們為什麼能夠安靜的那個底氣講得非常清楚.
詩篇46篇第十節說你們要休息.

$^{201}$要知道我是神.
後面說什麼.
我必在外邦中被尊崇.
在遍地上也被尊崇.
什麼叫做被尊崇?.
被尊崇就是被人所敬拜.
弟兄姊妹我們絕對不是用放空.
來對付仇敵帶來的恐懼跟焦慮.
恰恰相反.
我們是用神那個極具威嚴的震懾力.
直接震碎惡者的謊言.
大衛在詩篇29篇裡面就為我們示範這個強而有力的鼠鱗反擊策略.
對抗噪音的唯一有效方法.
就是要釋放那個更有權柄的聲音.
我們必須刻意讓那個神大有能力滿有威嚴的聲音.
充滿在我們的思想裡面.
當真理的話語像重低音的雷聲炸開的時候.
那些成天在你腦海裡面說謊.
定罪的蚊子就會全部被震落到地面上.
弟兄姊妹你還記得以利亞的故事嗎.
上帝的先知以利亞一人單挑450位巴黎先知.
再加上400位亞瑟拉先知.
比賽的規則簡單粗暴.
大家各自竹壇禱告.
誰的神能夠從天上降下烈火來燒掉積物.
誰就是掌管宇宙的神.
所以巴黎的先知先開始.
從清晨到下午整整850人在加密山頂狂呼亂叫.
瘋狂跳舞甚至拔出刀槍瘋狂割自己的身體搞得血濫噴唉唉叫.
當時的情景可以用三個字形容叫吵死人.
但結果是什麼.
聖經告訴我們說.
從午後直到現晚季的時候他們狂呼亂叫.
卻沒有聲音沒有音韻的也沒有理會的.
這個就是焦慮的噪音的真相.
很吵.
但沒有結果.
接下來輪到以利亞上場.
他沒有拿出擴音器來蓋過對方的聲音.
他不過做了一個尊崇神的禱告.

$^{241}$他說耶華.
求你今日使人知道你是誰.
就在以利亞話音剛落的那一瞬間.
轟的一聲開天闢地的巨響.
帶著上帝那個絕對威嚴的烈火從天而降.
在零點一秒瞬間燒盡祭壇的流毒木材.
連祭壇的石頭地上的塵土全部燒成灰燼.
在那一瞬間.
全世界都安靜了.
就像現在.
這就是耶和華的聲音.
那個從天上而來的重低音.
能夠瞬間清空我們靈魂裡面所有的噪音.
弟兄姊妹神不要我們做被動的受害者.
神呼召我們起來.
要奪回屬靈生命的那個音控台.
當你覺得心煩意亂.
當你覺得快被四面八方的壓力壓得快無法呼吸.
我強烈地邀請你做兩件事情.
第一個你大聲朗讀神的話.
讓你的耳朵真實聽見神的雷聲.
流牧師在舟山禱告會給出更好的建議.
你趕快背經吧.
而且不要在心裡默默地讀.
用導讀的方式讓真理的聲音從外至內.
震碎我們心中的焦慮.
第二件事情你能做的敬拜.
請你放一首宣告神能力的詩歌.
這個時候你不要放那種哀怨的詩歌.
你需要放的是那種宣告神主權的屬靈戰歌.
把音量開大聲一點.
跟著詩歌大聲地唱出來.
讓神的雷聲奪回那個心中的主控權.
你要相信.
當那個坐在寶座之上的神怒吼的時候.
你生命當中那些所有的焦慮跟風浪.
除了乖乖平息之外.
別無其他選擇.
第二個.
拆毀虛假的安全感.

$^{281}$跟你旁邊的講拆毀虛假的安全感.
在詩篇29篇5到6節說.
耶華的聲音震破箱柏樹.
耶華震碎黎巴嫩的箱柏樹.
它還使之跳躍如流瀆使黎巴嫩和西連跳躍如野牛瀆.
這一場屬靈風暴.
現在帶著神震動天地的權柄從地中海面登陸.
一路向北橫掃黎巴嫩的山脈.
而這場風暴的首要目標.
叫做黎巴嫩的箱柏樹.
首先我們要理解黎巴嫩的箱柏樹是什麼.
它絕對不是你在路邊看到的行道樹.
在古代黎巴嫩的箱柏樹.
可以說是木材中的勞斯萊斯.
這種樹能夠長到12層樓高.
這個樹幹粗壯到需要幾個人合抱的.
它的根會深深扎進土地的最深處.
有無與倫比的穩固.
而且它的紋理很細緻.
香氣很濃郁的.
所以當所羅門建造聖殿.
當埃及的法老打造豪華的巨船.
用的全是箱柏木.
所以對當時的人來講.
箱柏樹就是尊貴跟安全的代名詞.
但是當耶和華的聲音一響.
發生什麼事.
經文兩次的強調.
耶和華的聲音震破箱柏樹.
在希伯來文當中.
這不是指樹被折斷而已.
這是指樹被炸裂成為碎片.
上帝連一根手指頭都沒動.
只是發出聲音.
這個讓人類視為最堅固的象徵.
就輕而易舉被炸成碎片.
那我們就要問.
箱柏樹是招誰惹誰.
神為什麼要把它炸成碎片.
其實在這裡.

$^{321}$箱柏樹是代表人所依賴的安全感.
是人最習慣去倚靠的資源.
所以箱柏樹可能是你所引以為傲的工作.
箱柏樹可能是你精心規劃的投資理財策略.
箱柏樹也可能是你多年累積的人脈跟關係.
我們人常常看著我們這些親手栽種的箱柏樹.
就想說這下我穩了.
我的人生非常有保障.
結果人就不知不覺開始倚靠這些箱柏樹.
去遠離那個創造箱柏樹的神.
弟兄姊妹這一陣子油價高漲.
任何財源的風吹草動.
都會讓人心中焦慮的不得了.
但會不會當你一打開手機銀行的APP.
你發現你的帳戶餘額.
在第一後面還有好幾個零.
你的肩膀馬上就鬆下來了.
會不會在那一個瞬間.
這個手機銀行的APP所帶給你的平安.
比起你在房間裡面禱告一小時.
所得到的平安還快呢.
還大呢.
那麼這個時候銀行帳戶.
是不是就變成我們的箱柏樹了.
所以有時候神會降下風暴.
親手把我們的箱柏樹震碎.
神要親手打破我們虛假的安全感.
好讓我們懂得去抓住那個真正的磐石.
就是神他自己.
接著風暴他沒有停在北方的山脈.
他繼續南下橫掃迦底斯的曠野.
經文描述了一個更驚心動魄的畫面.
在詩篇29篇7到9節這邊說.
耶華的聲音使火焰分岔.
耶華的聲音震動曠野.
耶華震動迦底斯的曠野.
耶華的聲音驚動母鹿落胎.
樹木也脫落盡光.
凡在他殿中的都稱說他的榮耀.
弟兄姊妹啊.

$^{361}$當你看到這個經文這裡所描述的.
像不像我們人生當中被風暴掃過去的現場.
當我們工作挫敗.
家庭的危機.
或者意外來襲的時候.
我們可能會覺得自己就像經文講.
那隻早產的母鹿一樣.
驚慌失措.
或者像那片被剝光的樹林一樣.
感到很羞恥.
我們會質問神說.
為什麼要震碎我的生活啊.
為什麼要讓我看起來如此狼狽啊.
但是答案就藏在第9節說.
凡在他殿中的都稱說他的榮耀.
請思考一下這個對比.
殿外是破碎的香柏樹.
是嚇壞的母鹿.
是赤裸的曠野.
但在殿中.
在樹林的真實裡面.
所有人都發出讚嘆榮耀.
為什麼.
因為當那些巨大的香柏樹倒下之後.
你的視野才終於被清空了.
以前的那些成就,財富,人際關係.
擋住了我們的視線.
讓我們誤以為他們就是一切.
現在樹倒了.
我們終於可以毫無阻礙的看見.
祂高高坐在大水之上的神.
從那一刻開始.
我們才真正踏上信靠神.
與神同行.
經歷根心變化的旅程.
就像在新約當中有一個人名叫掃羅.
他早年的履歷表是頂級香柏木.
他的血統很純正.
他是名師門下的高材生.
他覺得自己超級穩的.

$^{401}$他還是一棵自帶殺氣的香柏樹.
他為什麼大發熱心.
要除掉那些他看不順眼的基督徒.
直到他走在大馬士革的路上.
那一道讓火焰分岔的聲音.
直接把掃羅這一棵驕傲的樹震碎了.
他瞬間撲倒在自己變成一個.
需要人拉著手才能走路的瞎子.
但這個也是神對他做過最溫柔的事情.
因為神震碎了舊的掃羅.
是為了在廢墟中建造出偉大的保羅.
經歷了這一次的破碎.
保羅寫下震撼千古的告白.
在斐利比書三章八節說.
不但如此我也將萬事當作有損的.
因我以認識我主基督耶穌為至寶.
我以為他丟棄萬事.
看作糞土為要得著基督.
親愛的弟兄姊妹.
如果你此刻感覺你緊抓不放的那棵香柏樹.
被無情的破碎了.
當原本的安全感破碎了.
請不要絕望.
更不要急著找另外一棵香柏樹來依靠.
聽聽那個雷聲背後的聲音.
神不是要毀滅.
而是要改變.
請記得如果不折斷舊的掃羅.
永遠不會誕生新的保羅.
如果你覺得生活正被劇烈震動.
不要害怕.
神從未停止祂的工作.
祂在預備你的人生.
可以承載更大的榮耀.
這個時刻我們就要學習第三點.
轉眼仰望神的寶座.
跟你旁邊講.
轉眼仰望神的寶座.
大衛看著眼前這個被風暴吹藏光禿禿的世界.
他的心裡沒有絕望.

$^{441}$他反而被聖靈充滿.
寫下這首詩篇最偉大的神學宣告.
在詩篇29篇第10節說.
洪水氾濫之時耶和華坐著為王.
耶和華坐著為王直到永遠.
洪水這個希伯來文在整本聖經當中.
舊約聖經.
除了這個地方之外.
只出現在另外一個特定的地方.
那就是創世記挪亞時代.
那一場滅世的大洪水.
那是一場無情吞沒整個世界.
終結人類文明.
讓一切受造物都陷入死亡的毀滅性洪水.
但即便如此.
大衛人看到一個極其榮耀的事情.
就是耶和華坐著為王.
請記住這個神的姿勢.
當洪水跟風暴氾濫的時候.
我們的神並沒有在天堂的走廊上急著滿頭大汗.
他沒有焦慮地對天使長米加勒說.
快快快去拿沙包去拿水桶.
沒有的.
聖經清楚地告訴我們說.
他坐著就表示絕對的掌控.
坐著就表示完全的安息.
坐著就表示沒有任何力量可以動搖神的統治.
這個對我們今天充滿焦慮的人來說.
是何等大的安慰.
不管你人生所經歷的洪水有多麼恐怖.
不管你有多麼崩潰.
多想尖叫.
你必須知道那位創造並掌管整個宇宙的神.
此刻依然安穩地坐在他的寶座上面.
你還記得耶穌跟門徒在加利利海上所遭遇的風暴嗎.
彼得,雅各,約翰.
這些從小在加利利海上討生活的漁夫.
嚇得魂飛魄散.
但在這個危急時刻.
耶穌在做什麼.

$^{481}$聖經說他在睡覺.
弟兄姊妹.
睡覺是比坐著還更高的境界.
焦慮最直接的副作用就是睡不著.
但是耶穌卻安穩睡覺.
直到門徒很慌張地搖醒他.
而面對風暴耶穌只淡淡地說了幾個字.
住了吧.
靜了吧.
瞬間風就止住了.
海就大大的平靜.
為什麼.
因為當造物主開口.
所有受造的風浪.
哪怕再狂.
都必須乖乖地趴下聽話.
弟兄姊妹.
風暴再大.
神的寶座也從來沒有搖晃過一釐米的.
神沒有失去過半秒鐘的控制權.
英國的偉大講道王子斯木真.
他說.
神在狂風暴雨中掌舵.
就像在平靜無波的湖面上一樣.
輕鬆自如.
你要知道.
斯木真.
其實他長期被嚴重的憂鬱症所困的.
但即便如此.
每當他憂鬱的時候.
他遵從神他禱告.
他相信神掌管一切.
當他全心遵從神的時候.
神就讓他每一次的事奉都充滿能力.
弟兄姊妹.
我們要學習像斯木真一樣.
重新認識到神的主權.
神與我們不同的.
神是完全超越的.
祂是那個永恆無限全能的主宰.

$^{521}$只有當我們深刻認識到.
祂是這位宇宙的創造主.
我們才能夠在狂風暴雨當中.
真正的安息.
當然我知道有些人會抗議說.
我知道神有至高主權.
但是祂懂什麼.
祂高高在上.
祂享受天使的敬拜.
祂怎麼知道我現在的苦難有多痛啊.
弟兄姊妹.
基督教跟世界上其他宗教.
哲學最大的不同.
就在於我們的神.
不是遠在天邊冷眼看戲的神.
當我們質疑上帝.
你懂什麼.
請你看看十字架.
主耶穌為了拯救我們.
祂親自來到這個破碎的世界.
祂經歷極度的貧窮.
祂被門徒猶大出賣.
彼得三次不認祂.
祂被不公平的政治跟宗教體制羅織罪名.
被鞭打到血肉模糊.
最後赤身肉體被釘在那個木頭上.
在十字架上耶穌甚至經歷與天父分離的終極痛苦.
祂大喊我的神.
我的神為什麼離棄我.
弟兄姊妹.
十字架是人類歷史上最大的惡.
受造物竟然企圖把創造主釘死.
但是十字架同時也是神展現愛的最高時刻.
弟兄姊妹我要說.
我們所信靠的.
是一位為我們背負傷痕的神.
他與我們同受苦難.
他不但懂我們在苦難當中的感受.
而且他復活的大能要帶領我們進入一個真實的應許.
就是詩篇29篇的最後一節.

$^{561}$那邊說耶華必賜力量給他的百姓.
耶華必賜平安的福給他的百姓.
詩篇29篇以大有能力的雷聲作為驚心動魄的開場.
但結尾卻是令人嚮往的Shalom平安.
因為真正的平安從來不是沒有風浪.
真正的平安是確切知道是誰在掌管風浪.
歷史上有一個人為我們示範了這個風暴之上的平安.
在第二次世界大戰時期的德國著名的神學家潘霍華.
他因為反對納粹的暴政被關進集中營裡面.
你可以想像那個壓力嗎.
每天聽著外面的槍聲.
不知道明天輪到的會不會是自己.
所以在這種徹底失控隨時都可能喪命的環境.
如果平安是建立在環境上面.
那潘霍華早就精神崩潰了.
但是他深知到神依然掌權.
他有沒有要焦慮 有啊.
他有沒有曾經對神質疑 也有.
但是他沒有任憑讓那個焦慮擴散.
相反的他成為那一座集中營裡面黑暗中的光.
他依靠神去安慰那些恐懼的囚犯.
為他們禱告.
他遵從神在牢房裡面寫下那首著名的詩歌.
所有美善力量.
那首詩歌說.
所有美善力量都奇妙遮蓋.
無論如何都期盼那安慰.
後面說.
在晚上 早上 每個新的一天.
神都將與我們同在.
而目擊者回憶就在潘霍華.
被帶往絞刑架的那一天早晨.
他的神色非常的平安.
面對死亡的狂風暴雨.
他對身旁的夥伴微笑的說.
這就是終點了.
但對我而言.
這是生命的開始.
這就是看見神坐在寶座上.
所產生那個無人可以奪走的平安.

$^{601}$願神賜福你們也有相同的平安.
我們一起來禱告.
天父上帝.
因為我們深知道.
你是那位大有能力.
充滿憐憫的神.
也許在我們的生命當中.
常常經歷到這些焦慮的時刻.
也許在生命當中.
我們常常向神問說.
神啊 為什麼.
為什麼會發生這個事情.
為什麼要震碎我的生活.
為什麼要拿走我手裡的一切.
但主說 孩子啊.
我都在.
不要害怕.
雖然眼前是一片的曠野.
但是你將看見那坐在寶座上的神.
祂的寶座從未動搖.
祂的旨意要留到萬代.
祂是那位掌管你生命.
掌管你生活的神.
祂是擁有你一切應許的神.
祂說 孩子啊 不要害怕.
當我們專心的跟隨祂.
我們必然看見.
那榮耀與威嚴的同在.
主我們謝謝你.
我們在任何的環境當中.
我們要繼續來敬拜你.
我們在任何的環境當中.
我們要繼續來讚美你.
我們願你的聲音.
來震動我們的心.
震碎那焦慮.
震碎那仇敵的謊言.
讓我們重新站立起來.
能夠為主來宣告.
為主來征讚.

$^{641}$主我們謝謝你.
我們這樣禱告.
奉主耶穌基督的聖名.
Amen.
\newpage



\section{瑪拉基書 2:1-9}
\label{sec:032224}
\textbf{與神有約}~~~~石偉漢牧師
\newline
\newline
連結: \href{http://www.church.com.hk/acms/content.asp?site=cdc\&op=show\&type=product\&code=032224\&layout=sermon}{\texttt{032224}} ~~~~ 語音日期: 2026-05-31
\newline
\newline
\hyperref[sec:032223]{\small{< < < PREV SERMON < < <}}
~
\hyperref[sec:index_chronic]{\small{[返順時目]}}
~
\hyperref[sec:032225]{\small{> > > NEXT SERMON > > >}}
\newline
\newline
瑪拉基書 2:1-9
\newline
\begin{longtable}{cl}
\hline
\hline
章節 & 經文 (和合本修訂版)\\
\hline
2:1 & \begin{tabularx}{0.7\textwidth}{X} 現在，眾祭司啊，這誡命是給你們的。 \end{tabularx} \\ \\ \relax
2:2 & \begin{tabularx}{0.7\textwidth}{X} 萬軍之耶和華說：「你們若不聽，若不放在心上，將榮耀歸給我的名，我就使詛咒臨到你們，使你們的福分變為詛咒；其實我已經詛咒了你們的福分，因你們不把誡命放在心上。 \end{tabularx} \\ \\ \relax
2:3 & \begin{tabularx}{0.7\textwidth}{X} 看哪，我要斥責你們的後裔，把糞抹在你們臉上，就是你們祭牲的糞。人要把你們和糞一起抬出去， \end{tabularx} \\ \\ \relax
2:4 & \begin{tabularx}{0.7\textwidth}{X} 你們就知道我頒這誡命給你們，使我與利未所立的約可以常存。這是萬軍之耶和華說的。 \end{tabularx} \\ \\ \relax
2:5 & \begin{tabularx}{0.7\textwidth}{X} 我曾與他立生命和平安的約。我將這兩樣賜給他，使他存敬畏的心；他就敬畏我，懼怕我的名。 \end{tabularx} \\ \\ \relax
2:6 & \begin{tabularx}{0.7\textwidth}{X} 真實的訓誨在他口中，他的嘴唇中沒有不義。他以平安和正直與我同行，使許多人回轉離開罪孽。 \end{tabularx} \\ \\ \relax
2:7 & \begin{tabularx}{0.7\textwidth}{X} 祭司的嘴唇當守護知識，人也當從他口中尋求訓誨，因為他是萬軍之耶和華的使者。 \end{tabularx} \\ \\ \relax
2:8 & \begin{tabularx}{0.7\textwidth}{X} 你們卻偏離正道，使許多人在這訓誨上絆跌。你們破壞了我與利未人所立的約。這是萬軍之耶和華說的。 \end{tabularx} \\ \\ \relax
2:9 & \begin{tabularx}{0.7\textwidth}{X} 所以我使你們被眾百姓藐視，看為卑賤；因你們不遵守我的道，在律法上看人的情面。」 \end{tabularx} \\ \\
[1ex]
\hline
\hline
\end{longtable}
$^{1}$好 弟兄姊妹主內平安.
家中的弟兄姊妹都平安.
好 我們一齊先有祈禱.
親愛的主 我們多謝你.
你留下寶貴的話語.
特別你藉著舊約.
最後一卷《函瑪拉基書》.
都仍然對我們生命有很多提醒.
幫助我們明白舊約.
幫我們明白我們能夠活出新約.
讓我們能夠被你幫助.
我們能夠成為一個有力的生命.
我們求你聽我們祈禱.
奉耶穌的名字而救 阿們.
在演奏會中最關鍵的是指揮.
有一次有一個指揮.
在管弈樂團裡做指揮.
指著不同的樂器.
彈奏不同的聲量,音樂,大小聲等等.
那次的指揮很興奮.
很想帶動整個表演的高潮.
然後樂團就越來越大聲.
就抬高了.
見到別人看不到.
就再抬高一點.
結果抬得太高.
聲音一下子就掉到地上.
將整個演奏的高潮變成了冰點.
原來指揮帶領整個樂團到達音樂的高峰.
但因為指揮的跌倒.
令整個音樂成為了一個低點.
整個音樂就停頓了.
其實教會都是這樣的.
我們讀馬拉基斯.
一直說的是上樑不正下樑壞.
教會的領導出現問題.
就會帶動到下樑的會眾,百姓,神的子民.
同樣出現問題.
所以上次我們說到A4的敬拜.
當中說到回歸之後的神的子民.

$^{41}$他們的生命的狀況.
今天我們會繼續看多一點.
特別是上次說到上樑.
我們今天會看多一點上樑.
一些以色列領袖的問題.
我們下次會再看下樑.
今天我們用到的經文是.
馬拉基斯二章一至九節.
其實這個經文都有一個結構.
上次Alex講到三文治結構.
大家有印象嗎?.
不過上次Alex講的是日式三文治.
我今天是用港式三文治.
無論日式還是港式都是那一種.
重要的三文治本身就是.
有一塊麵包在上面.
一個餡在中間.
下面又有一塊麵包.
那個麵包是什麼呢?.
我們稍後會讀到的A段.
就是二章一至四節.
就是神警告祭司背若帶來的奏作.
而另一個麵包就是到了第八第九節.
神指控祭司背若必有後果.
而當中那個餡在二章五至七節.
回顧著若和神的應許和要求.
所以你看到我為什麼用港式三文治.
其中特色是什麼呢?.
我們是吃腿蛋治的.
腿蛋治不是只有蛋.
不是蛋治.
腿蛋治有雞蛋和火腿.
所以今天講到的內容.
先由內而外再由外而內.
聽懂了嗎?.
我們先講三文治餡裡面的蛋.
再講兩塊三文治的麵包.
再看三文治裡面的火腿.
聽懂了嗎?先由內而外再由外而內.
這樣看下去.

$^{81}$所以我們用這個方式去看.
整個經文裡面一至九節的段落.
簡單來說我們會分開三個重點.
三個重點包括第一個重點就是立若.
第二個重點就是背若.
第三個重點就是踐若.
踐即是實踐的意思.
不是踐踏.
是實踐若.
現在這個速度應該可以了吧?.
OK,我控制住.
我們先看中間的部分.
中間的餡裡面有蛋和火腿.
我們特別抽了蛋的部分去看.
這裡特別看第五至七節.
讓我再一次讀一讀.
「我曾與他立生命和平安的若.
我將這兩樣賜給他.
使他全敬畏的心.
他就敬畏我.
懼怕我的命.
真實的憤悔在他口中.
他的嘴唇中沒有不疑.
他以平安和正直與我同行.
使許多人迴轉離開罪孽.
祭司的嘴唇當守護知識.
人也當從他口中尋求憤悔.
因為他是萬君之爺和華的使者」.
這幾節經文追溯以色列人和神之間.
最開始的美好的狀況.
是什麼狀況呢?.
就是神與他們納若.
事實上我們讀聖經.
我們會知道.
在聖經裡面說了很多不同的若.
特別是在舊若裡面有很多不同的若.
包括我們知道.
羅亞方舟裡面有彩虹之若.
我們經常想起彩虹.
你們記得.

$^{121}$我們聖經裡面說的彩虹是七個顏色的.
但是有些同性戀群體.
借用了這個概念.
用了六個顏色去表達彩虹.
跟我們說的不一樣.
但是你們知道.
神與羅亞納若的彩虹遺跡.
表示上帝不再以洪水滅大地.
接著我們會讀到有亞伯拉罕的若.
說他的後裔和應許之地的福氣.
接著我們會讀到有摩西的若.
說律法和聖潔的要求.
到今天我們有新若的部分.
事實上很多時候神為何要與人納若呢?.
就好像今天如果你要租一間房子.
你為了安全會跟他簽一份契約.
這是對雙方的承諾和保證.
所以你會簽一個租屋契約.
結婚你會簽婚書.
有一個婚約.
代表兩夫婦之間.
會在婚姻裡面有一個彼此的承諾和承擔.
所以上帝也很想與人納下一個盟約.
以致讓人與神有一個關係.
有一個責任和承擔.
以致能夠互相實踐若的精神.
這就是若的意念.
而我們知道神與人納若的最大分別是什麼呢?.
神是遠遠高於一切.
人是何等的卑微.
不過高於一切的上帝甘願與普通人納一份契約.
本來祂有一個優勢可以納這個契約.
不過祂願意與我們納一個平等的契約.
讓我們能夠參與在祂的契約裡面.
而今天我們讀到的部分是說到這個契約.
又不是說羅亞之約.
又不是說亞伯拉罕之約.
是說一個什麼契約呢?.
與誰納的呢?.
我們在第一,第四和第八節.

$^{161}$我們看到有幾個提及的人物.
包括第一節是說到眾祭司.
一群的祭司.
第四節是說到與利未納的約.
而第八節更加清楚地解釋.
是與利未人納的約.
所以今天我們讀到的一至九節的部分.
正好是說神願意與利未人和祭司立約.
當然有些人很熟悉聖經.
也有些人不太熟悉聖經.
利未人就是一群亞國的子孫.
即十二個兒子裡面其中一個就是利未.
利未所生出來的就是利未的支派.
所衍生出來的利未支派.
而當中的祭司是什麼呢?.
祭司就是當中的利未人裡面.
其中屬於亞倫和亞倫後代的那群人.
才特別被列為祭司.
才有機會成為一個祭司.
所以你簡單想想.
利未人就不一定是祭司.
但祭司就應該一定是利未人.
不過無論你是利未人還是祭司.
很明顯就是上帝有這個約的時候.
特別針對這群利未人和祭司這個支派裡面.
一群人所納的一個約.
所以你看到這幅圖.
你知不知道這兩個人有什麼分別?.
一個祭司.
一個大祭司.
有什麼分別?.
祭司協助大祭司亞倫.
第一代就是亞倫.
祭司就會穿左邊的白色服飾.
而大祭司亞倫就會穿右邊的服飾.
這些大家可能都看過很多的服飾.
不過你知不知道亞倫.
他都會穿左邊的衣服.
知不知道?.
什麼時候要穿這個衣服?.

$^{201}$原來他要為百姓贖罪之前.
他都有一段時間要穿這件衣服去面對上帝.
所以原來都有這樣的情況.
但無論如何我們知道.
祭司這群人物,這群利未人,這群特別神.
讓他們能夠進入和其他十一個支派不一樣的約裡面.
特別如果你再看我們現在讀經計劃.
剛剛讀到什麼?.
文授記.
你讀到文授記第三,第四章.
其實是不是開始多一點聽到行軍了?.
因為講了這麼久.
終於要走了.
以色列人到文授記才開始繼續推展.
他們講完會幕的東西.
繼續推展文授記.
開始講到第二年.
出亞歷史第二年開始的行動.
而當中特別講到行軍.
軍隊有六十萬大軍.
六十萬二十歲以上的男丁.
我都慢慢講.
六十歲以上的男丁.
行軍的陣營.
記不記得第三章.
三個三個支派這樣走.
第一個是什麼?.
就是猶大支派帶領著三個.
第一隊.
第二隊就是劉變.
有三個支派.
但當未講到第三隊之前.
原來講到有一群人去搬那些會幕的器具.
工具,物品,木材等等.
原來就是在講利美.
所以利美的支派就排在中間.
接著再排隊的時候.
第三隊就是以法連支派帶領著的三個支派.
最後就是彈支派帶領著的四個支派.
所以如果走一條直線的話.

$^{241}$一路數下來.
第一隊,第二隊.
利美加,聖殿的所有東西.
接著第三,第四隊.
原來是一條這樣的線這樣走.
接著如果去到停下來.
即是走會幕走到差不多.
停下來的時候怎麼辦呢?.
紮好營就紮了一個會幕的位置.
接著其他的支派就圍著會幕.
分了東南西北.
而其中在中間又是那些利未人.
所以你見到利未人是很不一樣的.
跟其他的支派領袖是不一樣的.
再簡單一點問.
其實你見到那十二個支派.
因為當中利未人已經歸了神.
即是去服侍神.
所以當中就分了以法連和麻拿西.
再分了兩個.
就算是十二個了.
但是你留意這十二個支派.
他們主要是有什麼用的呢?.
為什麼列陣走的時候走頭和走尾.
當擺營的時候要擺圍著會幕呢?.
這幫人有什麼用呢?.
這六十萬.
第一章民數記你們讀經運動.
考考你們有沒有溫書.
幫助利未人操練一下你們的聲色.
第一章裡面說什麼?.
他們為什麼要數點人數?.
行軍打仗嘛.
所以這六十萬大軍.
都是其中一個數點.
他們選擇他們的原因.
他們要行軍打仗的藍丁.
但是你想一下.
利美是不用打仗的.
你想一下打仗和不打仗有什麼分別.

$^{281}$打仗的話有機會面對什麼?.
死亡.
受傷死亡.
但中間的利美其實是什麼?.
如果用一份功去說.
是筍功.
他在裡面是安全的.
因為有其他人保護.
如果連他都保護不了.
全部人都死定了.
所以他是保護.
但他還有一個重要的特點.
他分別為聖.
是專心服侍上帝的.
你想一下是不是很棒呢?.
不用出去上班.
不用去會計師樓工作.
只是專心服侍上帝.
所以Jenny很有福氣.
雖然他是三分之一的利未人.
原來是一個這樣的人.
原來上帝選了利未人.
能夠參與這個約.
其實他們是蒙福的一群.
這個約是什麼呢?.
聖經裡面提到神和人立約.
其實沒有特別說什麼生命和平安的約.
我們通常說羅亞的約.
亞伯拉罕的約.
摩西的約.
甚至有時候說大衛的約.
但沒有特別說什麼生命和平安的約的概念.
其實這是一個什麼概念呢?.
其實你要看民數記.
又是民數記.
大家讀經.
讀到第25章就要留意一下.
第25章說什麼呢?.
因此你要說.
漢雅我將我平安的約賜給他.

$^{321}$這是他和他的後裔.
你知道說哪些人了.
他和他的後裔永遠當祭司職任的約.
因他為了上帝而嫉妒.
他為以色列人贖罪.
這裡正好在民數記的時候.
說當中在以色列的陣營裡有一些混亂,淫亂.
以致你看看第25章的部分.
有一些拜偶像巴勒的問題.
以致後來神讓菲尼哈代表神.
去對付這些背叛上帝的人.
所以神就說菲尼哈.
現在將平安的約給了他.
這個約就讓他永遠當祭司的職任.
這裡也衍生到後來.
當然是說一群屬於利未人的祭司.
一群屬靈的敬拜上的領袖.
生命和平安的約.
簡單來說可以想到包括什麼呢?.
生命就是你的性命.
生命的約.
指的是本身以這種態度和與上帝建立關係的那群人.
他們會有一個生命上帝賜給他們.
同時平安就是指人與神的關係.
包括你與神之間有和平.
你與神之間的關係好.
也讓人與神的關係好.
生命和平安的約可能包括三個重點.
包括自己領受這份約的人.
他的生命能夠與神的關係有很好的生命.
第二 他也讓人與神有很好的關係.
第三 他也讓所有其他人都能與上帝有很好的關係.
這三個方向.
如果生命和平安的約是指這三個方面.
我們就很明白究竟這個約有什麼目的.
就是我們剛才讀到的三文治的內涵.
包括第一 一定要敬畏神.
第五 生命和平安的約是什麼呢?.
將這兩樣生命和平安賜給他們.
要使他們全敬畏的心去敬畏神.

$^{361}$去懼怕神的命.
大家記得我們之前在講第一章的時候.
我們對上帝的關係是什麼呢?.
父親對子女 主人對僕人.
子女對父親就是孝敬.
僕人對主人就是敬畏.
就是說 當這些領袖能夠得到生命和平安的約的時候.
他們也要幫助別人去敬畏上帝.
幫助別人處理我們如何與神的關係.
第二 第六節說到.
真實的訓誨在他口中而行.
他的嘴唇沒有不義.
以平安和正直與我同行.
這個人他自己也能夠按照律法的要求去生活.
他領受這份約的時候.
他能夠有一份真實的生命和上帝建立關係.
以致他活出來的時候.
他的口 他的嘴唇.
他走出來的都是按照律法去生活.
就是說對他自己是有幫助的.
第三 他也要負責將知識去教導別人.
使別人能夠從他口中尋求上帝的心意.
即是說他能夠幫助他能夠與上帝和平和好 平安.
原來就是這樣.
所以原來第一個部分我們要看到的就是一個怎樣的狀況呢.
就是原來今天上帝要與人立約.
並不是要你很辛苦.
弟兄姊妹.
今天上帝與你納了一份生命和平安的約.
其實是要祝福我們.
剛才我們說到外面那十二個支派要打仗.
要冒生命危險.
但原來他能夠成為中間的那些利未人.
那些祭司能夠坐在營地中間.
近會幕 近神的一個位置的時候.
其實是一個祝福 是一件美事.
今天原來上帝讓我們與他立約.
是一件祝福 是一件美事.
但很多時候今天我們與神立約.
簡單來說就是我們今天能夠成為基督徒.

$^{401}$我們只是有這樣的身份.
我們有時候也會覺得很麻煩.
今天有時候我們覺得做基督徒很辛苦很麻煩.
好像很多事情都不行.
很多事情都不可以.
我們聽到很多人說.
你又不可以這樣 又不可以那樣.
以致我們覺得做基督徒有時候是很辛苦的事.
弟兄姊妹 我們有時候搞錯了.
我們覺得做基督徒很辛苦.
是因為我們不明白.
原來上帝給我們一個基督徒的身份.
我們有一個約與上帝立約的關係.
其實是一件美好 是一件祝福的事.
我們不要搞錯.
如果我們以為做得很辛苦.
其實是我們誤會了.
所以有時候有些基督徒覺得.
信耶穌要做得好是重擔.
但我想告訴你.
如果你珍惜到這份這麼特殊與你的關係.
你就不會覺得重擔.
昨天有沒有看球賽?.
不關你的事 看那些球賽.
只是看網球而已 世界盃還沒到.
昨晚看什麼球賽呢?.
昨晚看NBA.
昨晚NBA很重要.
昨晚NBA是說西岸的決賽.
西岸的決賽是打七場的.
賽制是七場.
最先拿到四場勝利.
就是能夠出線成為西岸的總冠軍.
昨天是打第七場.
第七場就是前面的六場是三比三.
就是雷霆對馬刺.
三比三.
你想一下.
很重要的決勝關鍵.
而當中主場就是雷霆.

$^{441}$雷霆做主場.
希望能夠在自己的地方拿到比賽的冠軍.
這個西岸的冠軍.
當中在電視機看到全場都穿藍色的衣服.
雷霆的主場是整個藍色的海.
因為他們是藍色 馬刺是黑色的.
全部都是藍色.
原來每個人都穿著一模一樣的衣服.
嘩 你知不知道多有氣勢.
但是你想一下.
這麼傻的.
特意買件衣服去捧自己的球隊.
會不會是重擔.
我花錢買這麼貴的衣服進去.
還要花這麼貴的門票買進去.
會不會覺得很辛苦.
為什麼不覺得.
要花這麼多錢 不覺得.
是的 我們真的不覺得.
我們甘心樂意花這麼多錢買這麼貴的門票.
因為這些門票.
我估計要過千元以上才能進去看.
美金來的 因為看這些門票是很貴的.
還要穿著一模一樣的款式.
沒有自己的衣服.
你看其他比賽 每個人都穿自己的衣服去.
但是這些比賽是全部一模一樣.
為什麼這麼奇怪.
為什麼不會覺得很辛苦.
不能穿自己喜歡的衣服.
因為他這一刻穿著這件衣服.
是尊貴的身份.
代表了我支持的那隊.
我是支持的那隊 我是那個立場.
你明不明白.
原來整個形勢很不同.
我穿這件衣服是代表我尊貴的身份.
我是支持這隊.
我和你們是同一隊的.
我站在雷霆那邊.

$^{481}$今天如果我們都想想.
神和我們立約都是這樣.
挺好的.
我們掛著基督徒的名字.
不是很辛苦.
我甘心樂意掛著基督徒的名字.
你也是 我也是.
全部基督徒站在一起的時候.
我們人海挺有氣勢的.
輸人不輸陣.
氣勢都贏了.
我們要認真想一想.
我們有沒有想一想.
神給我們這種身份.
我們一起掛著這種身份.
其實不是重擔.
這是福氣.
神給我們的福氣這麼近.
能夠代表祂.
掛著祂的波衫.
代表神的波衫.
代表神的制服.
代表神的軍隊.
代表神的憲制者.
哇 何等美好.
這是第一個我們要認真想一想.
立約是一件好事.
但是.
是否每個人都珍惜呢?.
接著看到三文治的兩塊麵包.
麵包有多少塊?.
啊 又記得我剛剛說的兩塊麵包.
所以我們先看.
兩塊麵包包括了一至四節.
一塊麵包.
八至九節另外一塊麵包.
先看一至四節的麵包的部分.
就是我們說的「背若」.
「背若」.
這裡說什麼呢?.

$^{521}$「現在眾祭司這個誡命是給你們的.
萬軍之謀.
話說你們若不聽.
不放在心上.
不把張榮耀歸給我的名.
我就使咒詛臨到你們.
使你們的福分變為咒詛」.
其實我已經咒詛了你們的福分.
因為你們不把誡命放在心上.
「看啊 我要斥責你的後裔.
把份抹在你們的臉上.
就是你們自生的份.
人要把你們和份一起抬出去.
你們就知道我頒這誡命給你們.
使我與你們所納的約可以相傳.
這是萬軍之謀 話說的」.
這裡說了一個很有趣的問題.
我們既然知道我們有一個身份.
我們有一個尊貴的身份.
理論上我們應該很尊貴地去珍惜這個身份.
不過有時候又不是.
特別這裡說了什麼呢?.
他說到這群祭司有三個「不」.
「不聽」.
「不放在心上」.
「不把榮耀歸我的名」.
說到一開始的時候.
其實可能他們也不是想這樣.
不過就是由開始不聽那一刻.
開始不聽就開始不放在心上.
不放在心上就開始不歸榮耀給神.
原來他們那種狀況.
是不知不覺隨著那種狀態的轉變.
由不聽 不放在心上.
以至到最後不把榮耀歸給神.
一路一路惡化.
有時候我們做基督徒的生命.
不是即時變差的.
我們很少會很熱心地回教會.
突然有一天完全不回教會.

$^{561}$我完全抗拒神.
很少會這樣.
很多時候是一步一步走下去.
原來是這樣的.
上班也是這樣.
這裡沒有人喜歡上班.
只有你一個.
只有Rose一個.
你也不喜歡上班.
老闆也不喜歡上班.
什麼世代.
原來我們很喜歡上班.
我自己也很喜歡上班.
因為我上班就是回教會.
上班是很開心的.
但有時候我們一開始上班是有熱誠的.
我拼勁地挖掘一番事業.
做一些驚天動地的事.
希望有好的將來.
但上班到一個程度.
日以繼夜 夜以繼日之後.
就開始沒什麼反應了.
這裡說.
慢慢地將你的上班.
本來好好的將來變成例行公事.
辦公變成辦公.
有時候上班做事也不上心.
時間到就可以走了.
四點半就可以走了.
可以這樣的.
當然有些公司容許.
但有些四點半的心已經走了.
肉體在這裡五點鐘.
心已經在這裡.
九點鐘一上班已經埋罩.
打完卡就可以離開了.
原來是這樣的.
其實老闆的學習不是這樣的.
就是說.
在想很多事情的時候.

$^{601}$原來我們不知不覺的心態轉變了.
以致到.
本來有好好的.
但最後又不是那麼好.
這裡說上有政策.
下有對策.
所以本來是一個好好的狀況.
上帝用你和你納藥.
不過慢慢不知不覺.
你的對策就是什麼呢.
容忍自己一路一路一路轉變.
最簡單.
你想一下.
讀經運動.
其中一個很重要的元素是什麼.
聽YouTube.
聽YouTube的聲音.
聽聽你聽的音樂.
剛才崇拜播放的音樂.
播放音樂的時候.
開始慢慢很集中地聽.
聽著聽著.
哎呀這個音樂.
行了.
安靜了兩分鐘.
不用工作了.
我做其他事情了.
慢慢由不聽的那一刻開始.
當你聽那些聖經.
盧家福音.
不知不覺聽著聽著.
沒有聽到那些內容了.
不放在心上.
所以讀經計劃.
記住這個計劃不是要令你辛苦.
而是要讓你得到祝福.
所以這裡正正好就是.
那種不去三個不.
令到一路一路越來越變差的過程裡面.
最後有什麼結果呢.

$^{641}$有什麼後遺症呢.
這裡說到.
包括第一個就是什麼呢.
福分變奏座.
福分會變奏座.
這裡你要留意聖經說.
一開始的時候.
是這麼寫的.
他說.
我找回那句.
他們說你們弱.
你們弱不聽.
你們弱不放在心上.
聽起來是不是好像還沒發生.
你弱過這樣就會這樣.
你如果才會這樣.
但你留心後面說什麼.
留心後面說.
其實我已經詛咒了你們的福分.
即是說什麼呢.
如果你這樣.
神就會這樣對你.
但其實你已經被詛咒.
即是代表什麼呢.
其實你已經在做這些事了.
你已經做了.
你不知道.
你不知道你不聽.
不放在心上.
不把榮耀給神.
你不知道.
所以神給你在聽祂吩咐的過程裡.
那些福氣.
祂給你立約的福氣.
原來你不知不覺就已經讓祂流失了.
這個不是假設性的問題.
而是發生了的問題.
這班祭司就是將一個這麼好的福氣.
做祭司 親近神 服侍神.
不知不覺成為例行公事.

$^{681}$以致他們為了要符合百姓的期望要求.
就在做百姓想要的事.
而不是上帝想要的事.
所以他們的福氣沒有了.
成為了州座.
越是坐得高位.
州座越大.
因為你越是虧欠上帝.
接著說你的後裔要被斥責.
為什麼呢.
因為你的後代.
被你這樣影響之下.
他們都不會做到好的.
你想想你今天如果上教會.
輕輕地上學.
你都不要期望你的子女都很認真.
原來你做祭司就這樣.
你都不要期望之後的祭司會更加好.
你讀士師記就知道了.
上了史詩篇常見班就知道.
去到後來那些利未人又取妾.
將人砍成十二粒 當祭物般砍.
變質了.
那些後裔不可以有好的狀況.
所以說要斥責他的後裔.
甚至後面說將糞末抹在臉上.
你想像一下你作為祭司.
其中一個很重要的功用是什麼呢.
你將那些祭身砍開之後.
特別是有一些祭身要燒在壇上.
而一些不要的東西要放在外面.
其中包括了什麼呢.
這裡說述序祭.
包括了公牛的皮和所有的肉.
頭,腿,內臟.
就是全公牛.
內臟後面還有一個糞.
糞在禮儀上是不潔淨的.
這些獻祭動物的糞在禮儀上是不潔淨的.
所以他們要搬到營外清潔的地方倒掉.

$^{721}$倒灰之處放在柴上用火焚燒.
要全部不要.
糞是這樣處理的.
在這個述序祭裡面.
《利未記》第四章你應該都讀過了.
現在說什麼呢.
將那些糞抹在誰的臉上呢.
抹在祭司的臉上.
就是說他們是不是聖潔呢.
已經成為不潔了.
這些要獻祭的一幫人.
這些要處理獻祭的一幫人.
他們已經被糞便沾到臉上了.
你想想 這麼骯髒的.
如果我們今天說骯髒的.
甚至乎在他們的禮儀上是不潔淨的.
他們不可以再留在這個會幕裡面做獻祭了.
他們要離開 要出去.
甚至乎要去營外了.
所以這裡說的是什麼呢.
原來背若是帶來這樣的一個後果.
接著下面那塊麵包是怎樣呢.
第八 第九節.
「你們卻偏離正道.
使許多人在這份悔上半跌.
你們破壞了我與利未人所立的約」.
這是萬軍耶和華所說的.
「所以我使你們被眾百姓藐視.
看為卑賤.
因你們不遵守我的道.
在律法上看人的情面」.
看別人的頭做事了.
結果怎樣呢.
這裡說第二個問題.
就是下面那塊麵包是說什麼呢.
「偏離正道」.
他們的背若.
上面是說三個不.
下面是說什麼呢.
就是歪了.

$^{761}$「偏離正道」.
離開了應該要走的路上面.
以致到怎樣呢.
這裡說.
第一 你會半跌人.
因為很多人跟著你走.
你走去哪裡.
別人跟著你走.
他們不會的.
所以當你偏離正道的時候.
你都會令到很多人.
都在這個正道上面偏離.
這裡說甚至是半跌.
即是卡到人.
卡到人即是令到他.
整個跌倒.
這個跌倒是什麼呢.
就是破壞了別人和上帝的和平.
剛才我們說到.
生命和平安的約.
其中一個平安的部分是說什麼呢.
你讓人和神和好.
讓人和神和平.
現在原來你令到人和神不得和平.
使人半跌.
這是一個大的問題.
這是一個上帝不容許的問題.
你不單止令人跌.
這裡還說什麼.
你令自己跌.
為什麼呢.
原來是使己被藐視.
因為原來當你不認真去看待你的身份.
不認真看待你的事奉.
不認真看待你的領袖的角色的時候.
百姓一見到你.
看你為卑賤.
看你不配.
以致他看你的份量大大降低.
你記得我們說什麼叫榮耀.

$^{801}$榮耀就是重量.
那樣東西是值一公斤重的.
你給他一公斤重的重量.
就是榮耀.
藐視正正好就是什麼.
本來你應該有一公斤的.
但是現在只能夠給你的重量.
可能只有半公斤.
零點一公斤.
你完全拿不回你應有的重量.
你是被藐視.
即是你累己累人.
原來一個背弱的狀況就是累己累人.
作為一個領袖.
我們不謹慎我們的生命.
不謹慎我們的行為.
不謹慎我們怎樣做一個示範的作用呢.
我們確實會帶領別人離開上帝的正道.
所以很多人會說.
哇 基督徒是這樣的.
別人這樣說.
其實他看東西是有問題的.
不過.
他有時這種說法是有道理的.
因為有時我們真是這樣.
我們真是這樣.
最近在香港有一宗新聞.
不知道你有沒有聽過.
就是說一個中學的校長.
嘩 為甚麼每個人都顫抖.
大家都很update一下.
一個中學的校長.
他帶一群學生.
去到新加坡.
即是郊遊團.
應該是開心的事.
這個校長.
其中能夠在youtube裡面.
排這麼高名次.
被這麼多人搶著看的原因是甚麼呢.

$^{841}$是因為這個校長.
在新加坡.
帶著校巴要停在路邊的時候.
就跟那些security的警衛保安發生衝突.
發生衝突.
一般人都會有.
我們都會有時跟別人有衝突.
不過.
這個校長在衝突的時候.
就對著那些保安說髒話.
不斷罵那些保安.
還叫那些學生繼續拿 不要理他.
作為一個校長.
他有個很重要的責任是甚麼.
要以身作則.
他在學生面前說髒話.
第一個問題.
在學生面前說髒話.
罵這些保安.
這些保安正在盡他們的職責.
是第二個問題.
可以談.
可以爭取.
但不是這樣的方式.
結果甚麼呢.
結果現在最新的.
當然被人放上網.
放了影片.
剛才問大家.
大家個個都說看過了.
最新就說.
他都要為這件事承擔的時候.
他都拍了個短片.
去公眾道歉.
亦都說.
他要跟校董會請辭.
他辭了這個職位.
因為他知道他其身不正.
他現在這個行動.
已經被網民藐視.

$^{881}$而且他叛奪了這間學校的學生.
以後大家一想.
這間屯門嘟嘟嘟嘟嘟的學校.
就真的看不起這間學校.
我們起碼是這樣嗎.
我們說做校長要這樣.
我們做基督徒不用這樣嗎.
我們一些行為表現.
都會帶來一些衝擊.
會否叛奪別人.
亦都讓自己被人藐視.
所以這個問題.
我們要認真去想.
這裡說.
背若不一定是公開的否認.
因為很多時候我們背若離棄神.
或者對神的關係不好.
往往是由心的冷淡開始.
我們要問你的心在哪裡呢.
剛才說得好.
九時就已經回家了.
十時的心又回家了.
十時崇拜.
十時的心已經回家了.
你的心在哪裡呢.
你的心是否火熱呢.
是否冰冷呢.
是否不冷不熱呢.
存在於我們沒有留心.
我們的心有退化有冷淡的時候.
我們馬上去燃點我們的心的火熱.
我們馬上處理我們心裡面的問題.
以致不要以為我們.
由一步去到下一步.
以致我們一路一路向神說拔.
這是第二點.
我想大家去思考一下.
我們會否不知不覺背棄了這個若呢.
第三點就是去實踐這個若.
去實踐這個若的時候.

$^{921}$我們回到中間的部分.
究竟我們怎樣去思考這個若的狀況呢.
就要思考一下.
其實上帝既然這個若是好的.
既然我們背棄這個若是不好的.
那我們要思考一下.
應該怎樣去行好這個若呢.
在第七節說.
「祭司的嘴唇是當守護之色.
人當從他口中尋求分配.
因為他是萬君之爺和華的使者」.
這裡說的是.
原來是有出路的.
當你似乎越來越遠離神.
背若的時候.
你是有出路.
你只需要回去恢復那個過程.
恢復那個若的關係就可以了.
而這個呢.
是怎樣做得到呢?.
其中一個關鍵是什麼呢?.
重視你的身份.
這裡說.
原來這個是萬君之爺和華的使者.
原來這裡說.
我們有一個使者的身份.
一個大使的身份.
作為神的大使.
你去到哪裡呢?.
你都是代表了神的.
就好像你今天是美國大使館的職員.
你去到世界各地的國家.
你為什麼這麼有權柄呢?.
因為你代表美國.
來到別人的國家裡面.
所以今天我們也要想想.
我們有沒有重拾我們的身份呢?.
這裡特別說的是.
這個聖若裡面.
我們對我們屬靈的領袖.

$^{961}$或者對我們每一個的信徒.
都是有要求的.
而這個裡面.
我們一定要搞清楚.
我們在當中的使者的身份.
以致我們能夠活出來的時候.
能夠得到別人的尊重.
這裡有一本書.
是由陳奕謙牧師所寫的.
叫《為何教會不像樣》.
裡面說的是什麼呢?.
他說教會患上一種叫教會官的貧乏症.
教會官的貧乏症.
他說很危險.
因為根據他幾十年來的經驗.
他去到很多不同的華人教會.
他發現真的有一個結論是什麼呢?.
就是原來不少毀身建立教會的人.
原來都不知不覺缺乏了一個合乎聖經.
簡單,具體,可行.
而又被測試的健康教會官.
就是說我們今天的教會觀念.
可能我們不知不覺都被扭曲了.
我們對於怎樣去活出一間教會的樣式.
那個標準原來我們不知不覺都沒有了.
以致到今天很多教會.
有時候他們裡面的實質都不是很健康.
所以失去了主動焦點和準繩.
受到不同的潮流和轟毒的一些個案.
甚至乎是成功神學所影響.
事實上今天很多教會都用成功神學的標準.
成功神學的標準其中一個標準是什麼呢?.
就是人數.
一間教會是否成功就以人數來決定.
所以你是巨型教會.
你只要知道你的教會有過萬人.
甚至十多萬人.
你就是一間成功的教會.
所以你能夠去到一間成功的教會.
你就覺得自己是一個成功的基督徒.

$^{1001}$那就錯了.
去到所謂成功的教會.
首先成功是什麼?.
成功都不知道.
但你以為這樣就是成功的基督徒.
你又是錯了.
基督徒的本身是在於你自己的生命.
能否與神和好 與人和好.
能夠與自己和好.
要搞清楚.
不要讓教會不像樣.
教會不是一個建築物.
不是說這個建築物不像樣.
而是我們這群人像不像樣.
像不像樣 能否交出來.
交出來的意思就是.
第二個重點 我們再看.
怎樣去做呢?.
原來他說.
第一次說.
他說原來這個誡命是給你們的.
要實踐這個約.
一定要清楚這個誡命是上帝給我們的.
給我們的.
這是他和利美所立的約.
是可以常傳的.
其實沒有到期的.
其實上帝和人立了約.
沒有延期的.
而很多時候是因為人毀約.
所以這個約才不能夠有效實踐.
上帝和我們立了一個永恆的約.
所以我們可以回到過去.
怎樣回到過去呢?.
就是我們要重建我們的身體.
剛才說重拾你的身份.
就是使者.
現在說什麼呢?.
就是要重建你的身體.
就是一個好像祭司一樣的身體.

$^{1041}$這個身體是怎樣呢?.
這裡說.
包括我們從右邊看一下.
包括你的耳朵.
要聽從所傳的誡命.
你的耳朵要張開.
因為我們人很多時候有選擇性聆聽.
叫做selective listening.
不知道你覺不覺得.
有某些人跟你說話.
你是聽不到的.
為什麼大家笑得出來.
有共鳴.
有些人說某些話.
是聽不到的.
接著我們就說.
剛剛不是說了嗎?.
接著他說你沒有說.
是不是這樣.
我發現很多時候都是這樣.
這種狀況不是因為你是年老.
有些人說是不是你老了才這樣.
不是.
是你的心.
你的耳朵沒有張開.
就算年輕的也會有.
因為這個不喜歡跟這個人聊天.
他說的話我都聽不到.
明明是同一句話.
他說是不行的.
接著外面的人說是可以的.
你明白嗎?.
說同一件事而已.
外面的人說總是比較好聽的.
這裡牧師說.
總是聽完就算了.
如果對神的說話都是這樣.
就糟糕了.
所以第一個經文告訴我們.
就是要聽從所傳的誡命.

$^{1081}$要打開你的耳.
另外一個是什麼呢?.
你的口要怎麼處理呢?.
就是要真實的訓誨在他的口中.
即是你的口要活出來.
要說一些真實的訓誨.
真是神的說話.
那你想想你平時說的神說話.
準不準繩呢?.
有時候我們背金句都背不到.
我不知道這裡是怎樣的.
我小時候去教會.
我小時候也不是很小.
去教會是怎樣的?.
要背聖經的.
每次去團體要背金句的.
接著做栽培又要背金句的.
今天我們教會有信仰栽培.
我都要求栽培員.
和被栽培的人一起背金句的.
有沒有跟你背金句的?.
Johnny.
還沒開始.
你們已經上了幾堂了.
要大家背金句是互相抵禮的.
因為我們有時候都會善忘.
年紀大了.
有些金句都背得「甩甩甩甩」.
要背金句的.
你的口中要說得出來的.
但現在我們沒有這些習慣了.
因為手機很方便.
但如果有一天沒有手機在身.
手機都可以沒電的.
你背不背得出多少金句出來呢?.
你的口有沒有神的說話.
你的嘴唇有沒有不義呢?.
不單是有沒有聖經在你的口裡.
而且你講的說話.
會不會不知不覺.

$^{1121}$含有一些歹毒的說話.
一些咒詛的說話.
一些不公平的評語在別人身上呢?.
你的心要怎樣呢?.
要傳敬畏的心.
要傳懼怕神的心.
要懼怕神的命.
你害不害怕神?.
如果一會兒打風下雨 走雷閃電.
你敢不敢向天發誓.
向神發一個誓.
確保你發誓之後.
你都不怕他有閃電劈中你.
我們很害怕這些事情.
但如果我們害怕神.
就不是怕他劈不劈中你.
神有很多方法對付你.
但我們有沒有敬畏他呢?.
他的手使許多人迴轉離開罪孽.
我們要帶領人.
帶領人遠離罪.
因為你要走正路.
你將人帶回來.
走這條路.
是這條路才對.
不要走錯.
我們今天都很喜歡.
指責別人走這條路.
我發覺.
我去不知哪裡.
有人說你這樣走.
很多提議的.
但如果我們能夠在信仰上.
讓人走正路.
讓人指向一條明燈.
這樣很好.
最後是你的腳怎樣走呢?.
「以平安和正直與我同行」.
和神走路要有平安和正直.
平安 記住這就是Shalom.

$^{1161}$和神的關係.
怎樣和神走路呢?.
和神一路繞著膠走.
還是和神一路踏著膠走.
你和神怎樣走路呢?.
手牽手?.
還是說.
你有你行.
他有他行.
好像逛街一樣.
這邊是女的看女裝.
男的看男裝.
分開走.
終點先見.
會不會這樣呢?.
這裡說.
真的要我們提醒我們.
最後一個問題就是問.
這段聖經.
和誰說的呢?.
一開始我們問.
這段聖經一至九節和誰說的呢?.
祭司.
利美的子民 利未人.
特別是當中的祭司.
特別是回歸後的這班祭司.
可能你會問.
這段聖經其實與你無關.
你也不是祭司.
如果從這樣說.
就算你今天當牧師是祭司.
應該是牧師跟你說的.
牧師你自己跟自己說道德.
但是留意.
我們現在是一個什麼樣的教會?.
知道這是一個什麼樣的宗派教會嗎?.
浸信會.
浸信會有好幾個重要的.
他們的.
不是條文叫什麼.

$^{1201}$重要的觀念.
浸信會有什麼重要的觀念?.
不如你告訴我.
開浸信會.
你們正在開浸信會.
不是開宣導會.
信徒皆祭司.
浸信會很強調的是什麼?.
信徒皆祭司.
有很多不同的中文.
有不同的很重要的說話.
但是浸信會.
其實很強調這件事.
就是信徒皆祭司.
所以為什麼浸信會的體制.
是要有投票呢?.
最後是全部人投票.
去決定教會的方向.
因為每個人都有資格.
不是有些人高一些.
有些人低一些.
是大家都有同樣的均等份.
去投票去為教會.
去走那一步做決定.
這個不知道是不是.
這個沒有說是對還是錯的一種做法.
但起碼這裡說的一個觀念就是.
每一個信徒.
其實就有祭司的職份.
所以今天這段聖經.
雖然是在說跟祭司說.
所以如果你是信徒的話.
今天就跟你說了.
那我想問今天這段聖經是不是跟你說的?.
因為你認定你是信徒.
信徒有一個祭司的職份.
讓我們都能夠在我們生命裡面.
去持守這個職份.
所以今天的反思.
第三個反思就是.

$^{1241}$信徒皆祭司.
我們被稱為什麼呢?.
君尊的祭司.
我們沒有意識到.
在自己身上帶著神使者的身份.
這個托付給我們今天在我們的職場.
在我們的家庭.
在我們的生活.
在我們居住的環境裡面.
做這個見證.
最後是時候要背金句了.
今天的結束要背金句.
背不背得到?.
背金句背這個.
就是彼得前書二章九節十節.
不過時間關係.
背一下第九節.
最常背的應該都背得到.
我現在隨便抓一個人出來背得到.
你為什麼搖頭?.
通常搖頭的就是被人抓出來的.
背不背得到?.
不過我們背開和學本的版本.
不要緊 你試一下和修本的版本.
一起讀一下.
開始.
不過你們是被揀選的一族.
是君尊的祭司.
是神聖的國度.
是上帝的子民.
要使你們宣揚那一召.
你們出黑暗入奇妙光明者的美德.
你們從前不是子民.
現在卻成了上帝的子民.
從前未蒙憐憫.
現在卻蒙了憐憫.
彼得前書這句.
這裡面就是在說.
原來今天我們是一個君尊的祭司.
所以浸信會都引用這段聖經.

$^{1281}$去說我們是一個君尊的祭司.
所以這段話.
今天的經文都對我們有幫助.
讓我們能夠活出祭司的積分.
當然我們更加明白.
我們都是神聖的國度.
屬神的使者.
屬上帝的子民.
是祂的百姓.
圍著會幕的那一群.
我們要怎樣?.
要宣揚那一召我們出黑暗入奇妙光明者的美德.
我們就要宣揚那一位.
召我們出黑暗的那一位.
這是我們的身份.
這是我們的責任.
求主幫助我們.
我們與神有一個約.
是神主動邀請給我們有這個約的福氣.
願意讓我們珍惜有約的身份.
我們不要覺得辛苦.
我們樂於在這個約裡面.
以致我們拒絕一些.
令我們背棄上帝的行動.
而我們真的行出上帝的心意.
我們一起祈禱.
親愛的天父.
因為我們知道.
我們本來真的算不進你的約裡.
我們又不是以色列人.
更加沒有什麼尊貴的身份.
又沒有什麼背景.
更加不要說我有多少錢.
有多少學歷.
但主人你偏偏揀選了我們.
讓我們不配成為你的兒女.
我們有蒙基督的救贖.
讓祂的寶血能夠洗淨我們一切的罪.
以致我們今天成為君尊的祭司.
願主你幫助我們每一個.

$^{1321}$讓我們守著我們的身份.
尊重我們的身份.
看重我們的身份.
以致我們今天能夠成為主你的大使.
走出去不同的地方.
成為主你使用的.
願主你幫助我們.
聽我們祈禱.
奉耶穌名字而來.
阿們.
\newpage



\section{哈巴谷書 3:16-19}
\label{sec:032225}
\textbf{義人必因信得生}~~~~陳雅莉傳道
\newline
\newline
連結: \href{http://www.church.com.hk/acms/content.asp?site=cdc\&op=show\&type=product\&code=032225\&layout=sermon}{\texttt{032225}} ~~~~ 語音日期: 2026-05-31
\newline
\newline
\hyperref[sec:032224]{\small{< < < PREV SERMON < < <}}
~
\hyperref[sec:index_chronic]{\small{[返順時目]}}
~
\hyperref[sec:032226]{\small{> > > NEXT SERMON > > >}}
\newline
\newline
哈巴谷書 3:16-19
\newline
\begin{longtable}{cl}
\hline
\hline
章節 & 經文 (和合本修訂版)\\
\hline
3:16 & \begin{tabularx}{0.7\textwidth}{X} 我聽見這聲音，身體戰兢，嘴唇發顫，骨中朽爛，在所立之處戰兢；但我安靜等候災難之日臨到那上來侵犯我們的民。 \end{tabularx} \\ \\ \relax
3:17 & \begin{tabularx}{0.7\textwidth}{X} 雖然無花果樹不發旺，葡萄樹不結果，橄欖樹也不收成，田地不出糧食，圈中絕了羊，棚內也沒有牛； \end{tabularx} \\ \\ \relax
3:18 & \begin{tabularx}{0.7\textwidth}{X} 然而，我要因耶和華歡欣，因救我的神喜樂。 \end{tabularx} \\ \\ \relax
3:19 & \begin{tabularx}{0.7\textwidth}{X} 主耶和華是我的力量，他使我的腳快如母鹿，又使我穩行在高處。這歌交給聖詠團長，用絲弦的樂器。 \end{tabularx} \\ \\
[1ex]
\hline
\hline
\end{longtable}
$^{1}$各位姐妹平安.
又可以一齊回來神的殿裡面.
一齊聆聽祂的話語.
一齊去透過神的話.
我們來更認識我們的主.
今天我們一齊來看哈巴谷書.
不知道大家是否陌生這本書.
哈巴谷書其實是在舊約裡面的.
在舊約的聖經裡面.
舊約聖經裡面我們知道有仙姿書.
仙姿書裡面又有分大仙姿和小仙姿.
我們又知道大仙姿和小仙姿.
不是因為仙姿小一點.
那個仙姿大一點.
所以叫大小仙姿.
是因為那本書的長短.
它的篇幅比較長就是大仙姿書.
篇幅比較短就是小仙姿書.
在舊約裡面有十二卷小仙姿書.
其中一卷就是我們今天要看的哈巴谷書.
哈巴谷書說它是小仙姿書.
都合理的 因為它只有三章 很短.
那哈巴谷書是誰寫的呢?.
在哈巴谷書的第一章第一節就有說.
哈巴谷仙姿所看見的默示.
它就記錄下來.
我們知道當時的仙姿.
就是耶和華神的代言人.
他是舊約的仙姿.
它所記載的就是當時以色列國.
南國 猶大的信息.
當時在以色列地已經分為了南北兩國.
原本是一個國.
不過後來分裂 成為北國和南國.
北國繼續叫以色列 南國就是猶大.
其實在哈巴谷做仙姿的時候.
北國已經是亡國了.
剩下來的就只有猶大南國.
哈巴谷仙姿就在當時發佈信息.
去說神要她說的話.

$^{41}$或者說我們看哈巴谷書的時候.
我們會發現不只是記錄一些神要她說的話.
更加是她和神之間的對話.
所以哈巴谷書的三章.
頭兩章就記載了哈巴谷仙姿的疑問.
她有很多不明白的地方 她要去問上帝.
上帝也回答她 和她說話.
第一二章都是記載她和上帝的一些對話.
去到第三章就記載了哈巴谷仙姿的一個禱告.
嚴格來說 它是一個禱告.
也是一首詩歌來的.
是可以唱的 當時應該是可以配上音樂.
可以唱 是一首詩歌來的.
剛才主席帶領我們讀的那一段.
最後就是整個第三章.
這本書的最後那幾節.
也是我們說是哈巴谷仙姿一個信心的宣告.
當我們今天進入來看這本書的時候.
我們會發現 它不單只是一個.
記錄了當時仙姿和上帝的對話.
或者她的一個禱告.
我們更加發現 我們從第一章走到第三章.
是一個旅程來的.
是一個由黑暗走去光明.
由混亂走去最終安息的旅程來的.
我們一邊看 你們就會發現.
進入哈巴谷書 一開始的時候.
我們發現很特別.
它沒有交代任何的背景 年份什麼的.
當時是誰在當皇帝 沒有.
它在書卷的一開始沒有提及任何背景.
一開始就說.
記載了仙姿對上帝一連串的質問.
仙姿很多事情不明白.
她不明白 她要問神.
頭三節就已經看到.
她跟神說 我呼求.
為什麼你不應運.
要到什麼時候.
我呼喊 有暴力 為什麼你不拯救.

$^{81}$你看不見有罪孽嗎.
你有毀滅 有兇惡 有爭執 有紛爭.
不斷地發生.
你為什麼做事不理.
她真問神.
仙姿去問神.
她說 律法無效.
公理無彰顯.
惡人為困異人 你看不到嗎.
不知道今天你這樣看的時候.
有沒有覺得 這麼熟口熟面的.
都不是當時南郭猶帶的警方這麼簡單.
好像到今時今日.
我們仍然看見這些事情.
可見聖經不老套.
有些師妹跟我說很怕看舊約.
舊約很悶.
怎麼看都看不明白.
其實不是的.
舊約有很多歷史在裡面.
有很多神要讓我們看見.
祂是怎樣帶領祂的子民.
怎樣引領祂的子民.
怎樣教導祂的子民.
給我們透過仙姿書今日來看看.
仙姿很不明白.
她很不明白.
她看到猶大國中有這些事情發生.
不是外面的人打他們.
不是外面的人打猶大的國.
而是有這些罪孽 爭執.
而是自己人打自己人.
是在南郭猶大國中.
他們本是同胞.
他們同是神的子民.
但竟然在自己的國裡有紛爭.
有欺壓的事 有暴力.
有惡人去圍困異人.
不是被別人欺負.
乃是神的子民自己敗壞.

$^{121}$他們沒有守神的律法.
沒有行神的公義.
當時大概是猶大王約西亞死後的時間.
其實哈巴谷仙姿和耶利米仙姿.
是同時期的仙姿.
如果我們翻查一下耶利米書.
就會發現當時的情況真的很可怕.
有形容到當時猶大充斥著罪惡.
充斥著很多不公義的情況.
哈巴谷仙姿在當時看到這些狀況.
她就問上帝為何你可以坐視不理.
仙姿在一個很絕望的境況下.
她帶著一個很焦急的心情.
她要向上帝呼求.
她關心公義.
因為她看到為何上帝竟然容許.
在她的子民中間長期有人處於不義的裡面.
她不明白.
她覺得沒理由還不懲罰他們.
為何還不出手擊打他們.
仙姿在第四節提到.
她說律法無效.
公理沒有彰顯.
她指出以色列民已經不再遵守神的律法.
他們行事為人的標準也不再跟從神.
公理被惡人扭曲.
是很可怕的事.
今天.
這可能也是我們今天的膚淺.
當我們看到很多難以承受.
難以接受的現實的時候.
我們可能也會問.
神啊 你去了哪裡.
公理何在.
公義在哪裡.
今天可能也是我們的膚淺.
或者我們再仔細看.
我們發現她是在說同胞.
她是在說內部有互相殘殺.
可能也讓我們想到.

$^{161}$可能我們的國家與國家之間.
但再近一點不是.
是家裡面.
是自己人打自己人.
可能我們家裡面也有.
不知道會不會也有.
大哥欺負弟弟.
姐姐又反妹妹.
我們也不希望看到這些情況.
如果你的家裡面有這些情況發生.
你會怎麼求呢.
你當然希望這個現象不要再發生.
如果我們說一個代討事項.
希望我們的家不要再有這些事發生.
你可能就會求.
如果可以很神職地.
上帝明天就打一支針進入我們家裡面.
全部一睡醒都變得很有愛心.
很包容 很接納 很欣賞.
然後大家就很有愛 很和諧.
粗俗一點說.
最好快點轉死性.
快點轉死性.
最好明天早上就發生.
又或者我們可能也有另外一個祈禱的方向.
就是巴不得那些惡人快點死掉.
如果他們全死掉.
懲罰他們罪有應得.
那就世界和平了.
是的.
我們可能希望最好突然間天下太平.
又或者作惡的人罪有應得.
如果我們這樣祈禱.
我們又會希望神用什麼方法去回應我們呢.
你會希望神怎麼回應你呢.
可能我們會希望.
神最好打暈他們.
明天早上睡醒就變成好人.
又或者.
神可不可以快點懲罰他們.

$^{201}$我們期望上帝怎麼回應我們這些禱告呢.
又或者.
在當時的哈巴菊仙姿.
她去呼求的時候.
她會希望神怎麼回應她呢.
她也可能會像我們這樣想.
希望神可不可以快點令那些惡人突然間變善良.
可不可以全部壞事都突然間消失.
可不可以那些罪的友人和罪的人.
最好有很多自然的後果.
大病一場.
下天災 塌樓 砸鬼死他.
毀滅惡人.
可能是這樣.
仙姿一再問.
為什麼神不聽我們的呼吸.
為什麼你不理我們.
你看不到嗎.
為什麼你不應運.
為什麼你不拯救.
為什麼你坐視不理.
仙姿希望神怎麼做呢.
神的回答.
又怎麼超出了仙姿的預期呢.
結果神回答仙姿.
神跟她說.
意思是什麼呢.
意思是.
你只看到猶大.
你只看到你國中發生的事.
我現在要你去看看這個世界.
神說你要驚奇.
要再驚奇.
因為祂要做一件你難以置信的事.
神說儘管有人說.
你都不會相信這件事.
祂將要做一件什麼事呢.
就是神說我一定會興起加勒底人.
加勒底人是什麼人呢.
聖經形容他們是殘忍暴躁之民.

$^{241}$他們通行遍地.
霸佔不屬於自己的住處.
加勒底人其實是他們的敵人.
聖經形容他們是很可怕.
很殘暴的.
加勒底人其實是當時位於巴比倫的南部.
剛才我們說過北角二七已經亡國了.
當時是被亞述所亡的.
亞述的勢力當時已經開始衰落了.
而加勒底的王朝當時開始鞏固.
而最終加勒底人控制了整個巴比倫.
成為了巴比倫王朝的領袖.
我們知道.
如果我們一直看歷史.
我們知道最後猶大.
南角猶大是亡在巴比倫的手上.
神說我要興起一些殘暴不仁的人.
去打你們.
神給了一個很滿意的答案.
明天就天下太平.
人人變成善良.
神竟然說.
我要興起一些比你們更邪惡的人.
來懲罰你們.
這是什麼道理呢?.
神為什麼要這樣做?.
哈伯高先知不明白.
甚至神的回答.
我相信令當時的哈伯高先知更加困惑.
不問好過問.
問完更困惑.
你的答案完全超乎我的想像.
哈伯高先知提到.
他和神投訴的時候.
他說沒有公理.
公理沒有彰顯.
但神卻說審判和威權都要出自迦勒底人.
原來的文字公理和審判是同一個字.
在原來的文字裡面顯示了.
當先知說沒有公理.

$^{281}$神就說我要用邪惡的人來.
顯彰顯祂的公理.
你聽到會不會覺得很困惑呢?.
這是什麼道理?.
很難令人相信.
神怎可以用比我們更邪惡的人來審判我們?.
先知認為他們自己是比較公義的.
先知對神說.
你的眼目清潔.
你不看邪惡都不害奸惡.
為什麼你會看著人行詭詐呢?.
祂說惡人.
祂在說迦勒底人.
祂滅給自己公義的人.
祂在說自己.
祂在說猶太國.
祂覺得自己是比較公義的.
雖然已經有很多邪惡在我們裡面.
但祂仍然覺得自己至少比他們公義一點.
你怎可以用一些比我們更不公義的人來懲罰我們?.
先知不明白.
先知不明白這是什麼道理.
祂知道上帝是聖潔的.
祂知道祂是清潔的.
祂知道祂不看邪惡不看奸惡.
但祂沒有辦法.
難以接受.
上帝為什麼會這樣做?.
可能今天.
當我們帶著一些疑問去求問上帝的時候.
甚至乎我們也會帶著一些質問去問上帝.
究竟發生什麼事?.
為什麼會這樣?.
你有沒有想像過.
上帝會給你一些超出你所能夠預期的答案?.
是你不能夠理解.
不能夠明白的答案給你.
先知得到上帝的答案之後.
他發現原來警方不是會變得更好.
而是會變得更差.

$^{321}$現在都沒有忘記.
原來將來我們會忘記.
會被人踐踏.
先知之前就覺得.
上帝有沒有做事?.
為什麼你做這麼少的事?.
國裡面有邪惡.
為什麼你不理?.
現在他就覺得上帝做得太多了.
超出了.
這樣很過分.
當先知知道神會這樣出手的時候.
他發現原來他們將會被敵人完全的擄掠.
將要經歷亡國.
他明白不明白.
神的心意是什麼呢?.
或者今天.
當我們看到我們不能夠明白上帝的作為的時候.
我們又能不能夠體會.
神原來最終的心意.
是為了要他的民回轉歸向他呢?.
神是要藉著外邦的人.
去擊打他們.
以致他們能夠懂得重新歸回.
他們深深悔疚自己遠離神的律法.
不跟從神的律法.
以致神一再預言你們不回轉.
一定會亡國.
最終是為了要這班以色列民.
神的子民重新回轉認定他.
先知所關心的是.
上帝的公義有沒有在以色列人當中得到彰顯.
而神的回答超出了先知的預期.
或者有時候來到這些時候.
我們真的不明白上帝為什麼可以這樣.
我們不明白為什麼可以超出我們的預期.
或者這些時候我們只有必須謙卑去承認.
我們實在太有限了.
我們怎麼能夠跟創天造地的神.
跟無所不知,無處不在,無所不能的上帝去相比呢?.

$^{361}$我相信在我們的腦袋裡面.
最不能夠理解的是.
這位崇高過於諸天的上帝.
竟然愛我們至心.
竟然為著要我們回轉.
要我們歸向祂.
去到新約.
祂竟然用祂的愛子來為我們贖我們.
是我們無法理解,無法明白的.
不過我們一直再看下去.
我們就知道上帝不是一個不公不義的神.
祂興起迦勒底人.
同樣祂也將會審判迦勒底人.
當先知得到這個答案的時候.
祂就說:我要站在瞭望台上面.
我要看耶和華要對我說什麼.
祂可以用什麼話向祂申訴.
祂不明白,祂就願意跟神對話.
我要跟神申訴.
結果神回答他:好啊.
神回答他:你要將這些默示清楚地寫下來.
我要跟你說的話,你寫下來.
讓人容易朗讀.
結果在整個二章我們會發現.
寫下了所有神將要懲罰.
將要怎樣懲罰迦勒底人.
而在這個懲罰之前.
上帝應許先知和安娜.
安人自高自大,身不正直.
不過,為義人必因他的信得生.
這是上帝給先知的應許.
給當時的民的應許.
是,雖然你們將要經歷亡國.
很慘痛,將會是你們不能承受的.
但是,為義人必因他的信得生.
上帝應許他們只要保持忠心.
持守他們的信心.
行在義中.
當審判來臨的時候.
他們仍然必然得救.

$^{401}$在《新約聖經》裡面.
三次引用哈巴谷書裡的這一句.
義人必因他的信得生.
我今天不打算逐節跟大家看.
不過,在《新約聖經》裡面三次的引用.
我們都發現,其實來到新約.
當時的信徒都是面對非常困難的境況.
當時有逼迫,有患難.
面對很多不公不義的情況.
《新約聖經》的信息仍然是要為當時的子民.
帶來莫大的安慰.
義人必定會因他的信得生.
今天,同樣來到今天.
我們知道信心不是一刻.
好像信了主,就是那一刻.
信心是一生的旅程.
是你一直持守,不放棄這個信心.
持續下去信靠我們這位上帝.
你必然得救.
當我們落在苦難,困境,絕望當中.
上帝仍然要鼓勵我們.
仍然要安慰我們.
用內心等候祂.
堅持持守你的信心.
上帝應許,祂必不遲然.
到了時候,祂的公義必要彰顯.
兄弟姐妹,我今天不知道.
李嘉靖面對什麼樣的考驗.
或許你面前有很大的困境.
很大的艱難.
讓你感覺到絕望.
讓你感覺到沒有出路.
看不到明天.
又或者可能你經歷過傷害.
經歷過苦難.
你感覺到難過,憂傷,愧疚,無奈.
又或者可能你面前有非常大的挑戰.
你將要克服.
讓你感覺到難以承受.
讓你戰驚恐懼.

$^{441}$你要相信上帝今天.
仍然要用祂的話語安慰你.
祂應許你.
二人一定會因祂的信得生.
面對巨大的困難.
神不會不知道.
我們每一天所面對的是什麼.
但祂仍然想我們用耐心等候.
持守我們的信心.
我們要知道.
神在掌權.
不只是在我們看得見祂在做事的時候.
神的掌權是.
當祂完全沉默.
在我們完全看不見祂有什麼作為的時候.
祂仍然在掌權.
你能不能相信.
不只是上帝做些事情給你看的時候.
你覺得祂在掌權.
乃是祂完全好像沒有做事的時候.
祂仍然在掌權.
神在我們無論知道不知道.
明白不明白.
了解不了解的時候.
祂仍然是痛觀萬有的神.
我們只要憑信心相信.
二人一定會因祂的信得生.
來到這裡 哈巴菊先知很有趣.
當他聽到上帝說將會怎樣懲罰加勒底人.
審判加勒底人的時候.
不代表他們不需要面對.
他們南國猶大.
仍然需要面對他們需要的懲罰.
但當他聽到.
加勒底人神都不是不公不義的.
神都會審判的時候.
他就說.
為耶和華在他的殿中.
全地都當在他面前肅靜.
先知聽到.

$^{481}$原來神都會審判加勒底人.
先知那一刻就對上帝產生出一個敬畏.
他覺得全地都要在神的殿面前.
在神的面前肅靜.
當哈巴菊先知來到聖殿的時候.
有些人說他是不是沒眼睛看了.
逃避現實去了聖殿躲起來呢?.
不是的.
我相信他知道在神的殿裡面.
無論外面有多大的崩壞.
他都能尋著真正的安穩.
甚至有些人去研究這段聖經的時候說.
當時耶路撒冷可能已經陷入困境.
當時聖殿可能已經被毀了.
哈巴菊先知未必真的進入了一個物質的聖殿.
有人形容可能先知是進入了一個.
屬靈的眼光所看見的建築裡面.
他是用信心相信.
外面無論多破碎 這個世界無論多破碎.
人仍然可以在耶和華裡面得安息.
先知降伏在神的掌權裡面.
他明白惡人興起 惡人都會衰敗.
一切都在神的掌權裡面.
所以去到最後的部分.
今天剛才主席帶我們讀的部分.
我們看到他願意傳言的信服.
當進入第三章 先知的禱告的時候.
先知其中有一句說.
求神在發怒的時候 以憐憫為懷.
神要發怒.
不過他仍然求神可不可以以憐憫為懷.
並且最後先知說.
我聽見這些聲音.
就是將要攻打他們的聲音來到.
他說他會身體戰驚 嘴唇發顫 骨中扭爛.
你能夠想像 當我們面對困境.
面對將要來臨的困境的時候.
確實內心會很戰驚.
先知說他在所立之處戰驚.
面對逆境將要來臨 確實是很害怕.

$^{521}$但先知竟然能夠說.
帶我安靜等候.
他等候什麼呢?.
他等候災難的日子來到.
那些侵犯我們的民來到.
先知願意安靜等候神的作為.
並且發出最深的頌讚.
接著是我們很熟悉的經文.
我們當作是金句來的.
我們會背的.
我們一起讀好不好?.
雖然無花果樹不發旺.
葡萄樹不結果.
橄欖樹也不收成.
田地不出糧食.
圈中絕了羊.
棚內也沒有牛.
然而我要因耶和華寬恩.
因救我的神喜樂.
主耶和華是我的力量.
祂使我的腳快如舞鹿.
又使我穩行在高處.
是的 確實.
叫人怎能夠發出頌讚呢?.
在打仗的日子 被攻打的日子.
今天也是.
所有戰爭的國家.
無花果樹都不發旺.
葡萄樹都不結果.
橄欖樹都不收成.
田地沒有糧食.
很荒涼.
羊牛都沒有了.
你能夠想像.
被攻打的日子.
將要面對亡國的日子.
是多麼可怕的事.
但先知竟然仍然說.
然而我要因耶和華寬恩.
因救我的神喜樂.

$^{561}$今天.
盼望我們透過神的話語.
我們更深明白.
更深認識祂.
我們能夠仍然確信.
我們的上帝在掌權嗎?.
即使什麼都沒有了.
一無所有.
今天我不知道你面對什麼困境.
什麼艱難.
前面的路會不會變得很荒涼.
但我願你仍然仰望神.
無論將來的環境如何.
我們仍然堅信.
上帝在掌權.
有時對於上帝的作為.
我們總是欠缺一點耐性.
但我們仍然相信.
我們的上帝.
是不受時間空間所限制.
我們的上帝看千年如一日.
看一日如千年.
所以我們.
雖然這一刻.
眼前見到的是什麼都沒有了.
但我們仍然.
靠著我們的上帝.
在那個終極的盼望裡.
可以欣賞.
可以今天就欣賞.
我們的神怡歡呼.
先知在最後說主耶和華是我的力量.
祂使我的腳化如舞鹿.
能夠使我穩行在高處.
這一句其實是出自大衛的詩篇.
在詩篇18,33節裡說到.
在這篇詩的標題裡.
當耶和華拯救大衛脫離他的手的.
和掃來的手的日子的時候.
他寫這首詩.

$^{601}$去向耶和華說話.
是大衛的勝利之歌.
雖然是困難.
有和患有艱難.
剛才的詩歌很好.
有北風有巨浪.
有和患有艱難.
我藏身主任下.
昂然報國.
聽聞舞鹿有很獨特的本領.
鹿有超廣角的視覺.
和很寬大的體.
我們正常人類.
我們看到的大概是180度.
如果我們要看我們後面的東西.
我們就要轉頭.
但鹿的視覺是300度.
他只看不到最後面的60度.
很厲害吧.
所以鹿能夠在極度危險.
危機四伏的地方.
即使他低頭找東西吃.
他都能夠看到超廣角.
以致他能夠迅速回應.
逃避危險.
而且他的腳,他的體.
能夠走在超級陡峭.
非常斜的斜坡裡.
他都能夠穩行.
他在大自然裡覓食.
是一種奇景.
我想給大家看一段短片.
(英文翻譯).
是不是很厲害.
我第一次看的時候.
覺得站在這裡都很害怕.
站都站不住了.
看下去都心驚膽顫.
但他竟然能夠跑.
聖經說.

$^{641}$主耶和華是我的力量.
祂能夠使我們的腳快如舞鹿.
穩行在高處.
我們人生的逆境.
都如那些陡峭的斜坡.
如果我們依靠自己.
我們一下子就滾下去.
唯有依靠我們的神.
相信祂是我的力量.
認定祂是我的力量.
祂能夠使我們在人生的逆境裡站穩.
這是逆境裡的讚歌.
無論你將要面對什麼困難.
什麼艱難.
無論是身體的軟弱.
家人的需要.
甚至乎你面對世界的崩壞.
你覺得很可怕.
轉一個視角.
轉一個角度.
你能夠仍然在逆境中站穩.
其實哈巴菊書短短三章.
如果你有機會再回去看.
它講述了哈巴菊不只是在說苦難如何面對.
它是在說一個屬靈的旅程.
去看見上帝的作為.
從第一章開始.
還記得剛才一開始第一章.
哈巴菊才知道是怎樣開始這本書的嗎?.
她用質問.
用疑問去開始這本書.
到結束.
她用頌讚.
我們即使是很混亂的世代.
甚至乎我們的人生都很混亂.
但當我們在黑暗裡.
看見耶和華的掌權的時候.
即使動盪不安.
即使世界失序.
都不能夠讓我們跌倒.

$^{681}$從舊約到新約.
神仍然應許我們.
二人必因祂的信得生.
你仍然能夠靠著耶和華.
靠著耶和華寬恩.
有耶和華神而喜樂.
我們一起禱告.
賜給我們愛的上帝.
多謝主讓我們看見.
哪怕世界已經崩壞.
哪怕動盪的世代.
讓我們感到不安.
但我們仍然能夠.
不因著環境改變.
不因著風暴停止.
才能夠令我們感到安穩.
乃是單單因著你.
因著有神你在我們中間.
有人有神你在我們的中心.
我們不是依靠勢力.
不是依靠任何外在的力量.
也不是依靠人的能力.
乃是依靠坐穩在聖殿中的耶和華.
求主你幫助我們.
讓我們有超越的眼光.
現實雖然讓我們感到可怕.
但事實是你坐著為王.
你仍然掌權.
求主你讓我們看見.
不單看見現實.
更看見超越的事實.
又懇求主你.
讓我們深知.
終有一天主要來審判.
主要來審判這個世界.
願你發怒的時候.
以憐憫為懷.
恩待所有因信而行義的人.
多謝主你的應許.
凡所行義的人.

$^{721}$二人必因祂的信得生.
求主你讓我們今天仍然依靠你.
歡欣喜樂.
奉耶穌基督名字祈求.
阿們.
\newpage



\section{約書亞記 23:1-16}
\label{sec:032226}
\textbf{傳承敬虔與恩典}~~~~方兆雄牧師
\newline
\newline
連結: \href{http://www.church.com.hk/acms/content.asp?site=cdc\&op=show\&type=product\&code=032226\&layout=sermon}{\texttt{032226}} ~~~~ 語音日期: 2026-05-31
\newline
\newline
\hyperref[sec:032225]{\small{< < < PREV SERMON < < <}}
~
\hyperref[sec:index_chronic]{\small{[返順時目]}}
~
\hyperref[sec:032227]{\small{> > > NEXT SERMON > > >}}
\newline
\newline
約書亞記 23:1-16
\newline
\begin{longtable}{cl}
\hline
\hline
章節 & 經文 (和合本修訂版)\\
\hline
23:1 & \begin{tabularx}{0.7\textwidth}{X} 耶和華使以色列人從四圍所有的仇敵中得享安寧，已經有很多日子了。約書亞年紀老邁， \end{tabularx} \\ \\ \relax
23:2 & \begin{tabularx}{0.7\textwidth}{X} 就召了全以色列的眾長老、領袖、審判官和官長來，對他們說：「我年紀已經老邁。 \end{tabularx} \\ \\ \relax
23:3 & \begin{tabularx}{0.7\textwidth}{X} 耶和華－你們的神因你們的緣故向這些國家所做的一切，你們都親眼看見了，那為你們作戰的是耶和華－你們的神。 \end{tabularx} \\ \\ \relax
23:4 & \begin{tabularx}{0.7\textwidth}{X} 看，我已經把所剩下的列國，連同從約旦河起到大海日落的方向，我所剪除的列國，都抽籤分給你們各支派為業了。 \end{tabularx} \\ \\ \relax
23:5 & \begin{tabularx}{0.7\textwidth}{X} 耶和華－你們的神必將他們從你們面前趕出去，使他們離開你們，你們就必得他們的地為業，正如耶和華－你們的神向你們所應許的。 \end{tabularx} \\ \\ \relax
23:6 & \begin{tabularx}{0.7\textwidth}{X} 你們要大大壯膽，謹守遵行寫在摩西律法書上的一切話，不可偏離左右。 \end{tabularx} \\ \\ \relax
23:7 & \begin{tabularx}{0.7\textwidth}{X} 不可與你們中間所剩下的這些國家往來。你們不可提他們神明的名，不可指著它們起誓，不可事奉它們，也不可敬拜它們。 \end{tabularx} \\ \\ \relax
23:8 & \begin{tabularx}{0.7\textwidth}{X} 只要緊緊跟隨耶和華－你們的神，就像你們直到今日所做的。 \end{tabularx} \\ \\ \relax
23:9 & \begin{tabularx}{0.7\textwidth}{X} 因為耶和華已經把又大又強的列國從你們面前趕出；直到今日，沒有一人能在你們面前站立得住。 \end{tabularx} \\ \\ \relax
23:10 & \begin{tabularx}{0.7\textwidth}{X} 你們一人必追趕千人，因為耶和華－你們的神照他向你們所應許的，為你們作戰。 \end{tabularx} \\ \\ \relax
23:11 & \begin{tabularx}{0.7\textwidth}{X} 你們要分外謹慎，愛耶和華－你們的神。 \end{tabularx} \\ \\ \relax
23:12 & \begin{tabularx}{0.7\textwidth}{X} 你們若斷然轉離，緊緊跟隨你們中間所剩下的這些國家，彼此結親，互相往來， \end{tabularx} \\ \\ \relax
23:13 & \begin{tabularx}{0.7\textwidth}{X} 就要確實知道，耶和華－你們的神必不再將他們從你們面前趕出；他們卻要成為你們的羅網、圈套、肋上的鞭、眼中的刺，直到你們在耶和華－你們神所賜的這美地上滅亡。 \end{tabularx} \\ \\ \relax
23:14 & \begin{tabularx}{0.7\textwidth}{X} 「看哪，我今日要走世人必走的路了。你們要一心一意知道，耶和華－你們神所應許要賜給你們的一切福氣，沒有一件落空，都應驗在你們身上了。 \end{tabularx} \\ \\ \relax
23:15 & \begin{tabularx}{0.7\textwidth}{X} 耶和華－你們的神所應許的一切福氣怎樣臨到你們身上，耶和華也必照樣使各樣災禍臨到你們身上，直到他把你們從耶和華－你們神所賜給你們的這美地上除滅。 \end{tabularx} \\ \\ \relax
23:16 & \begin{tabularx}{0.7\textwidth}{X} 你們若違背耶和華－你們神吩咐你們所守的約，去事奉別神，敬拜它們，耶和華的怒氣必向你們發作，使你們在他所賜給你們的美地上迅速滅亡。」 \end{tabularx} \\ \\
[1ex]
\hline
\hline
\end{longtable}
$^{1}$各位弟兄姊妹平安.
在今天和大家繼續看約書亞記來到尾聲的傾文的時候.
不知道大家會不會留意到近來這幾個月.
尤其是在大街或者重要的商場,外場,巴士,港鐵.
都好像有某一類的廣告多了一點.
看看有沒有見過那些句子.
一起寫下傳承故事.
又為家的未來掌握財政先機.
世代無限的承傳.
亦都說到跨越時光轉動未來.
這些廣告過往好像都不是這樣賣法的.
但是就好像它轉了一個風向.
不知道是不是想打進一班累積了若干的財富.
很關注家人,後人能不能在當中受惠的潛在客戶.
加上近幾個星期有些娛樂圈相關的新聞.
亦都有一些後輩不務正業攤大手辦.
問年事已高的爸爸拿錢還理直氣壯.
很讓人關注到自己所勞苦積攢在地上的財寶.
是不是被用得合宜呢.
《聖經》關於財富的教導.
不論在舊約的《智慧文學》.
例如《傳道書》.
新約的《方庵書》.
主耶穌基督都很豐富的教導.
很警世的.
但是今天我們所讀的《約書日記》第23章.
所看重的就是比財富更加重要的傳承.
就是信仰的傳承.
是敬畏神生命的傳承.
這個已經超越了個人,家庭.
更加是關乎整個民族的.
整個信仰群體的.
而《約書日記》23章.
也是來到書卷的尾聲.
就是講到約書日的年老.
離世之前對以色列人民的叮囑.
究竟這位屬靈長者.
對下一代要用心腸的勸勉.
是包括些什麼呢.
也都對今時今日的我們.

$^{41}$有些什麼提醒呢.
不如我們就先讀《約書日記》.
第23章一到三節.
看看當中有沒有引入發生什麼事呢.
預備起.
耶穌說使以色列人安靜.
不如四圍的一切仇敵爭戰.
已經多日.
約書日年紀老邁.
就把以色列眾人的長老,族長.
審判官,評官長都召了來.
對他們說.
那年紀已經老邁.
耶穌說你們的神.
因你們的怨辜.
向那些國所行的一切事.
你們親眼看見了.
因為為你們爭戰的.
之而禍禍你們的神.
在我繼續講解之前.
讓我們先有祈禱吧.
愛我們的天父上帝.
你的話語.
讓不單止當時的人.
如何經歷神的恩典.
後來讀到的人.
都能夠了解到神的大能.
以及我們當中要注意的事情.
天父願以你的聖靈.
開啟我們的心竅.
讓他們能夠明白真道.
並且有這份心智.
樂意謹守遵行.
我們祈禱.
奉我主耶穌基督的聖名祈求.
阿們.
在這版經文裡面.
我們看到.
耶和華使以色列人安靜.
不予四圍的仇敵爭戰已經多日.

$^{81}$在當中一句說了他們的場景.
回顧他們曾經.
不論是在過約旦河的時候.
的各種爭戰.
那些已經大部分完成了.
終於有安靜的日子了.
情況有些改變.
約書亞再不是帶他們出出入入.
到處打仗的那位.
而是說到年紀老邁.
嘉儀弟兄姊妹有讀過約書亞的.
有沒有印象約書亞活到多少歲.
即是安息他的命有多長呢.
有沒有印象.
斯班有沒有印象.
是110歲.
約書亞結婚兩部分.
頭十二章就說到他們打仗.
到完結打仗分地的時候.
有記錄到迦勒.
另一位領袖.
就是那位和約書亞.
初初年輕的時候.
就是窺探迦蘭地.
回來報喜訊.
十二個探子裡面的兩個之一.
那時迦勒八十五歲.
約書亞應該和迦勒差不多年紀.
所以由分地到約書亞.
他接近臨終的時候.
計計數數.
110減85.
大概二十二十五年的時間.
耶和華史以色列人安靜.
沒有用四圍的仇敵爭戰.
當然他們當時的地還沒有拿完.
一來還沒有那麼多人.
沒有那麼需要.
但是約書亞看出.
他自己世上的日子不多了.

$^{121}$所以他要將這個訊息.
重要的訊息交大給.
當時的長老 族長.
審判官 領袖.
約書亞提醒他們.
那時是耶和華為他們爭戰的.
約書亞首先重提的是神怎樣帶領.
怎樣拯救那些以色列人.
回顧神的作為.
因為假如人以為.
我們今天有那麼和平.
是我們那天打仗回來的.
是自己功勞的話.
那就會覺得自己有這個權利.
去用自己的方式去生活.
但是假如他們要記得.
是神為他們爭戰的時候.
知道這是恩典的時候.
他們就需要留意和珍惜.
約書亞亦提醒他們.
一個很重要的屬靈次序.
並不是神的子民.
你們做得好.
神才賜這塊地給你.
不是 倒過來的.
是神為他們打仗.
神賜了給他們這個應許地.
然後就要求他們當中.
按神和他們所立的約來信服.
甚至他們得到的地.
都是抽籤抽出來的.
神按著他向列祖.
即是亞伯拉罕.
以撒雅各所立的約.
憑著他自己的信實.
經過幾百年之後.
帶領以色列人進入這個應許之地.
在這個恩典當中.
神的子民需要記得守約.
今時今日我們信了主耶穌的.

$^{161}$已經蒙恩的都是一樣的.
主耶穌祂先拯救了我們.
而且還要是唯有.
我們還在做罪人的時候.
唯我們死.
神的愛向我們顯明了先.
接著我們會應神的恩典.
來信服的.
所以在提醒的過程當中.
約書亞特別強調.
你們親眼看見.
神怎樣為他們征戰.
是他們親身經歷的.
不論見到耶利哥的城牆倒塌.
不論是他們曾經有位阿干犯罪.
為了他們帶來失敗.
神怎樣審判.
神怎樣赦免再重新帶領.
接著甚至在征戰的時候.
看見日頭已經停住.
以致他們的勝利能夠完整.
對於今時今日來說.
我們也需要回望神怎樣帶領你.
也要回望神怎樣帶領城鎮.
很多時候回想起.
在自己迷茫.
道路不明白的時候.
怎樣來依靠神的保守.
來供應.
來開路.
而這些經歷.
這些故事.
就是我們需要承傳.
亦是有得要來承傳給人的根本.
就是主耶穌基督的故事.
神在我們教會身上恩典的故事.
所以我們要提醒的.
我們所傳承的.
不是我們自己的功勞.
或是我們的習慣.

$^{201}$習慣有時好有時壞.
時代都不同了.
這樣做事怎可以同一個方法呢.
是要承傳神怎樣透過我們.
去造成他的工作.
以致將來的人.
新信的人.
他們生命遇到困難的時候.
知道神原來曾經這樣做過.
有能力的.
我們都要依靠他.
有這個效果在.
而另一方面.
又想回頭.
不知道大家會不會有同感.
人都有份傾向.
希望自己一生所建立的.
所付出的.
一來被認同.
被認可.
二來都希望有人能夠承繼.
發揚光大.
不論是你所做的生意.
你的教學.
你行業裡面的貢獻等等.
當然.
在這些傳承裡面.
有些是很有意義的.
有些可能是自私的.
有些譬如中國歷史裡面.
都有這些例子的.
秦始皇統一六國.
這個功業夠大了吧.
當然他們都希望二世.
三世.
萬世相傳.
但歷史同樣給我們看到.
這些建立在.
可能很多人的痛苦.
身上的權柄.

$^{241}$可以一時之間衰落.
而人之所以重視傳承.
相信都是知道自己的有限.
也都希望一些.
值得你珍惜的價值.
信念能夠傳承下去.
甚至這個我回想自己小時候.
都有一次這樣的學習和經歷.
在中學生的時間.
我不知道那時候的中學生和現在.
我們看到的中學生有什麼不同.
但那時候好像我們要承擔多點責任.
我在學校裡面有參與.
和負責音樂學會.
有一點像教會音樂部一樣.
學校大事大事.
什麼都在當中有份參與的.
當年雖然有音樂老師.
但基本上校風是鼓勵學生.
去承擔和主導的.
所以我中四的時候.
已經登入全校的歌詠團.
就是choir orchestra.
樂團的總指揮.
百多二百人和現在差不多.
到了中六去澳洲升學的時候.
我在外地.
我頭一個月心想.
現在音樂學會會怎樣.
因為學期不同.
英國的制度.
或者北半球是九月開學.
澳洲那邊是兩三月開學.
所以就很擔心.
究竟音樂學會會怎樣.
在異地頭一個月.
念念不忘.
於是我也不知道當時有什麼力量.
就是一口氣.
寫了一封長信給接任的師弟.

$^{281}$詳細解釋音樂學會怎樣運作.
當中的理念.
怎樣練歌詠團.
怎樣練樂團.
怎樣指揮.
圖文並茂.
怎樣處理人事.
哪些人信得過.
哪些人要小心之類.
二十歲還不夠中學生.
什麼什麼寫了九頁還是十頁紙.
還要是那些學生會出的五六毫米間距.
很密密麻麻的那些紙.
但是一年多之後.
學校.
因為之前那位老師退休.
來一位新的音樂老師.
就將所有東西都改了.
在當中的學習.
我體會到.
原來不同的人的時候.
所有東西都可以一時間改變.
但是我也很希望那種緊張.
那種精神.
那種認真.
那種為學校服務.
或者是當中的傳承.
都可以讓他們師弟能夠明白.
這個不知道是不是約書亞.
他當時對以色列民的一份心智.
所以在第六至第八節裡面.
看得到他祝福以色列人.
有什麼需要特別注意緊守的呢.
約書亞不是教他們種田.
怎樣種地.
怎樣牧羊.
就可以多一點羊群.
用來獻祭也好 不是自己著數的話.
他要祝福他們什麼呢.
我們一起來讀.

$^{321}$預備起.
所以你們要大大壯膽.
謹守遵行寫在摩西律法書上的一切話.
不可偏離左右.
不可與你們中間所剩下的這些國民參集.
那們的臣.
你們不可提他的名.
不可指著他起誓.
也不可事奉叩拜.
只要照著你們到今天所行的.
專靠耶和華你們的臣.
你們要大大壯膽.
不知道會不會讓我們想起約書亞的開頭.
第一章.
當他要承繼摩西擔任去帶領以色列人的時候.
臣親自對他說.
你要剛強壯膽.
謹守遵行我僕人摩西所吩咐的一切律法.
不可偏離左右.
到幾十年之後約書亞年老了.
將要離開世界了.
他將同樣的說話交給繼續生活的人.
約書亞他一生得力的秘訣.
不是在於有什麼軍事的秘方.
或者怎樣牧養 怎樣種地.
可以多點出產 讓生活更加豐富.
而是好處可以.
要叮囑他們要謹守臣的話.
約書亞忠心之處就是他不單止在臣那裡領受這份命令.
自己進行.
也著意地將這個命令傳遞給以色列人.
謹守.
Human Sharma 這個字.
是有留心.
進行.
保守 看守的意思.
即使臣的話語並不是掛在牆上的裝飾品.
基督是我家之主.
不是一個橫額屬靈的口號.
而是生命的道路和方向.

$^{361}$要留意 進行.
好像你保存這個寶物.
放在保險箱裡面珍而重之一樣.
所以既然人生的方向到老了.
所以約書亞要提醒他們要不可偏離左右.
究竟有什麼東西容易讓以色列民.
會偏離左右的呢?.
假如那時候已經過了二十多年安靜的時間.
大家都很敬畏神 大家都經歷過神的恩典.
都著意怎樣走在神的道路當中.
有什麼東西會讓他們偏離呢?.
第七節那裡講到.
他們不可以跟其他國民.
拜其他神的人去參集.
我相信這裡並不是種族歧視.
或者民族優越感等等.
二十世紀或者是之前的世界裡面那種的.
不好處.
而是特別提到在信仰上面.
以色列民需要留意分別為勝.
問題不是在於外邦人的血統.
正如他們剛剛過河打耶利哥城的時候.
怎樣知道情況呢?.
就是外邦的女子拉俠為他們帶路.
問題不是血統.
而是外邦人那種敬拜他們的偶像.
和那種腐敗的生活方式.
假如以色列民跟了他們的時候.
當然神是會審判的.
回想我們今時今日.
我希望在座沒有.
當然常識都告訴我們.
不可以從星期天回來教會.
然後又出去拜神之類的.
是不可以的.
但是現在二十一世紀.
科技日新月異的時候.
偶像已經不一定在廟裡面了.
任何那些取代神支配我們的心靈.
甚至決定你的價值觀的東西.

$^{401}$都可能成為了你的偶像.
一個簡單的例子.
大家回想一下.
平時大家上班或者退休的時候.
或者有很多的世界去做.
有空閒的時間你會做些什麼呢?.
究竟你會默想.
看屬靈書籍.
裝備一下自己.
支援一下自己的心靈.
或者鍛鍊一下自己的身體去爬山.
或者做運動.
還是可能你沉迷著.
投資升升落落的.
升升跌跌的.
或者沉迷著網上的買東西.
這個有什麼新款的東西呢?.
哪裡便宜呢?.
價格都隔了.
不是輕鬆的娛樂的問題.
分別之處你可能會分到的就是.
那個東西有沒有佔據了你的心靈.
充塞了你的生命.
而影響到你和其他人.
或者和神的關係.
這是此其一.
第二.
當約書亞刻意提醒他們要剛強壯膽的時候.
也都需要留意我們在這個世界上面.
價值觀和信仰裡面都有很多的不同的.
可能世界的人.
都是追求名利,著數.
各樣的享受,享樂.
或者就是在講誠信,吃虧的時候.
我們怎樣堅守真理呢?.
這個同樣都需要很大的勇氣.
約書亞特別提到他們要留意.
不可以參雜這些外邦的價值觀.
換了今時今日.
我們就要留意了.

$^{441}$有這份敏銳了.
究竟我們怎樣生活.
是否被社會的潮流.
或者網絡裡面各樣的聲音或者噪音.
個人的慾望.
取代了神在我們生活上的地位呢?.
是甚麼掌管著我們的生命呢?.
所以大大壯膽.
不一定要打甚麼大仗.
或者做甚麼轟烈的事情.
甚至簡單到上班的時候.
不講謊話.
不埋沒良心.
或者在家庭的關係裡面.
有些可能是不論恩親或者是怨親.
比較難相處的.
當中能不能夠仍然有一份寬恕的心呢?.
這些都需要勇氣的.
這些都需要神給我們的力量.
而經文第八節裡面所提到的.
相信是中心點了.
那裡就講到我們要.
專靠義和華我們的神.
專靠這個字我們可以稍為思想.
因為原文達伯的意思.
就是緊緊的結合,結連.
這個字同樣用在《創世記》第二章.
形容夫妻二人的關係.
人要離開父母.
與妻子連合.
二人成為一體的那種結連.
就是達伯這種.
所以專靠義和華.
就不是口講.
或者表面的忠誠.
而是一種深刻.
遵守盟約的關係.
緊緊依附,不離不棄的.
約書下這裡不單止叫以色列人.
不要拜偶像.

$^{481}$積極的就是要緊緊連於義和華.
而這份的聯繫.
在十一節後來就提到.
你們需要憤愛,謹慎去愛,和你的神.
希伯來文的「愛」所指的.
是忠誠,委心,愛慕,信服.
正如《申命記》第六章.
那裡安西跟我們說.
以色列你要聽.
義和華我們神是獨一的主.
你要盡心,盡性,盡力愛義和華你的神.
約書下同樣將《申命記》裡面的核心.
交給以色列人.
因為愛有時可以變成責任,習慣,程序.
不論是每個星期我們的崇拜.
或者你回城長邦或者個人的聚會.
你每個星期都是這樣的.
回來吧.
但是神所要的.
看來不單止是我們有沒有忠心.
來去分辨為聖的時間去敬拜祂.
更加看我們心有沒有分辨為聖去歸神.
主耶穌在約翰福音那裡這樣說.
祂對他們門徒說.
「你們若愛我,就必遵守我的命令」.
對我們今天的提醒.
我們今天還愛不愛主呢?.
仍然會不會以神的話語捍衛寶貴呢?.
我很欣賞一些弟兄姊妹很喜歡查經的.
很愛慕神的話語.
也很欣賞弟兄姊妹願意,樂意.
將神的愛傳遞給人.
在社區裡面有很多的參與.
來見證基督.
這些都是愛的表現.
在這裡也分享一位浸信會的先賢.
很值得我們學習.
雖然他很年輕.
何顯利.
假如不知道的話可能會以為他是男生的名字.

$^{521}$不是.
是女宣教士.
顯利中學也是因為這位宣教士.
來紀念他的精神來建立的.
何顯利 Henrietta.
他1835年的時候.
十七歲的時候.
也做了很多人生裡面很重要的決定.
那一年他和宿牧師.
Reverend Yehud Louis Struck 結婚.
兩個人同年領受呼召去海外宣教.
接著他們被差往遠東.
先後在馬來西亞檳城.
新加坡和澳門服侍.
而他們兩夫婦在1842年.
來到香港.
在五月份建立香港第一間浸信會.
翻查記錄.
原來也是香港第一間的基督教會.
有得比較的.
譬如 St. John's Cathedral 海外宣教是1849年開始的.
在1842年.
他們兩夫婦已經來到香港.
來建立教會.
那時空坐的是在1842年.
那時還在打鴉片戰爭.
《南京條約》在1842年.
八月的時候割讓了香港.
那時還在打仗.
還沒打完仗.
在1841年英國侵佔了香港.
他們兩夫婦在澳門看到這方向的需要.
於是就過來.
還沒打完仗.
這位年輕的宣教士.
也是為女孩提供教育的先驅.
早在澳門辦學的時候.
他們的理念就是.
生一個男孩就要生一個女孩.
對於百多年前中國人的社會當中.

$^{561}$不會給女子接受教育的情況當中.
那個算是創先學.
但是他.
可以看到1844年.
就安息主懷.
那年27歲.
黃韻多少歲了現在.
不是不是.
你應該多過27歲.
可以想到.
因為很年輕的生命.
他死了之後兩年.
1846年.
也在美國出版了.
何韻莉宣教士的回憶錄.
可以找回她的日記.
那些日記.
她對宣教存有浪漫崇高理想的人的描述.
打破他們的夢幻.
描述真實宣教生活是充滿難題的.
充滿焦慮,危險,不穩定.
也會面對試探.
其中一篇1838年.
她所寫的日記.
1838年她多少歲呢.
也是20歲左右.
她形容那天幾個晚上之前.
在祈禱會裡面.
有一位剛出來宣教的宣教士提到.
當他離開美國的時候.
以為自己已經撇下世界了.
跟隨主耶穌了.
世界在背後,十格左前頭.
他發現原來事實並非如此.
他在這裡的分享同樣有很多試探.
他的罪的心理研釋可以引導他走迷路.
何宣教士記述.
確實宣教士只不過是普通的男女.
並不會比其他基督徒更加好.
在這裡.

$^{601}$那時候應該是在澳門.
都會被時尚的社交圈子所包圍.
可以想像到在殖民的地方裡面.
那些商人,官之類.
都是來來回回.
他們又會收到舞會的邀請.
舞會不是猜吹的舞會.
是那些去桌球的舞會.
這些我們當然不會接受.
他寫了.
其實我認為在這裡.
使得我們注意世俗事物的事情.
比我們在家鄉的時候還要多.
宣教士同樣會面對很多的挑戰困難.
甚至罪的誘惑.
是一個很困難的社會場景當中.
有很多人向你招手的時候.
你怎樣去應對呢.
所以約書人對當時的斯里文有一個嚴重的警告.
我們繼續讀這個十二十三節.
預備起.
你們若稍微轉去.
與你們中間所剩下這些國民聯絡.
彼此結親.
互相往來.
你們要確實知道.
耶和華你們的神.
必不再將他們從你們眼前趕出.
他們卻要成為你們的網絡機構.
勒上你邊眼中的刺.
直到你們在耶和華.
你們神所賜的這美地上滅亡.
一個很嚴厲的警告.
首先當中提到.
假如稍微轉去與國民聯絡.
聯絡也很客氣.
因為十二節的聯絡這個字.
與剛才第八節我們看到的專靠.
耶和華的專靠.
創世的第二章.

$^{641}$夫妻的聯合.
這個字是同一個字來的.
假如你不再與專靠或聯繫.
你與那些外邦.
信義教的人聯絡的話.
最後那些人.
就會成為你的網絡被困住.
成為你機構的意思.
即是陷阱 失去自由.
邊帶來痛苦.
刺帶來長久的折磨等等.
即使將會被罪核制.
在當中承受不單止神而來的審判.
以及所帶來的各樣的惡果.
其實罪惡時常都是這樣欺騙人的.
現在我們看見社會裡面.
很多的騙案或電騙.
槍城都是這樣的.
起初吸引人.
很方便 很合理 很划算.
但是最後成為你的網絡.
人以為自己可以走一些精面.
貪心.
但是竟然被這個陷阱所控制.
或者甚至輕一點.
或者退一點.
人以為可能只是買東西.
或者是一些投資.
但是被人騙了一副身家.
或者人以為在世上所做的事情.
都應該被人尊重.
但最後被自我中心驕傲所蒙騙.
那個心不知不覺已經偏離了神.
所以若須阿廣警告他們.
假如你們稍微轉去.
特意加這個陽節稍微.
好像飛機航行的時候.
可能有些儀器故障.
或者沒有現在那麼先進.
上一世紀空難都是這樣.

$^{681}$起初偏離一點點.
後來迷失了方向.
不知去哪了.
然後就墜機了.
人離開神.
很少都是一兩天之間.
忽然做一個決定.
更多的都是逐漸慢慢的偏離.
今天先不要祈禱.
今天趕上班.
先不靈修.
晚上先吧.
太累了 明天先吧.
於是慢慢的妥協.
成為了我們遠離神的習慣.
若須在這裡警告他們.
因為這位信實的神.
不單止賜福給他們.
但假如他們杯葛犯罪的話.
你同樣都會降禍給他們.
這個對於我們就算今時今日.
在主耶穌基督裡面.
有著神異類的身份.
承受著這份救恩的人.
都要小心.
雖然我們知道.
神會赦免人的罪.
但都要理肯認罪.
知道自己其實是犯了罪.
當中需要迴轉.
需要回到神那裡才可以.
神的恩典是賜予.
但也要接受.
神的審判也是真實的.
對罪絕不苟且.
若須在第十五節提醒他們.
若說你們神所應許的一切福氣.
如何臨到你們.
假如你們轉向其他神明的時候.
也會照樣寫各樣的禍患.

$^{721}$臨到你們身上.
這個呼應著摩西在申命記三十章.
講到的.
我今日呼天喚地向你們作見證.
將生死禍福陳明在你們面前.
所以你們要揀生命.
使你和你的後裔都得存活.
對於新玉來說.
福音也有這份兩面的.
大家可能都記得.
又是三章十六節.
剛看到有個錯字.
看看大家找不找到哪個字錯了.
神愛世人.
甚至將他的獨生子賜給他們.
叫一切信他的.
不致滅亡.
字字應該沒有那邊的字.
不致滅亡.
不致於滅亡.
反得永生.
但是同一章聖經之後兩句就講到.
信他的人不被定罪.
不信的人罪已經定了.
因為他不信神獨生子的名.
在當中對神的福音.
怎樣的回應.
都帶來兩方面的結果.
不單是身份上的不同.
也都是結果上的不同.
就算我們覺得自己信了耶穌.
信了幾十年.
回到教會很久了.
我們需要檢視自己.
有沒有像額外書所說的.
信心沒有傷害的行為.
這是一個死的.
沒有果效的信心.
既然神已經給了我們恩典.
已經應許了.

$^{761}$在主耶穌基督裡面.
所有的都會成為是的.
那時候我們就需要留意.
黏著主耶穌基督.
連繫於主耶穌基督.
以致我們不單止在律約中有份.
也都在神的恩典當中有份.
行事為人.
要與蒙召的恩典相稱.
好,最後了.
在結語裡面.
約書亞記23章裡面.
年老約書亞對後來的人.
說出這份屬靈的遺囑.
提醒他們是義和和為你們征戰的.
這個就是你們屬靈裡面的恩典.
所以你們要謹守.
遵行他的說話.
要專心靠他.
不可偏離左右.
這個希望都是我們今天.
需要領受的訊息.
我們需要小心.
我們會稍微轉去.
與世界來參集.
要留意究竟我們有什麼可以承傳給別人.
神在我們生命當中.
個人的恩典,個人的故事.
就是我們可以承傳的,給別人的.
當後來的人那些信仰上.
經歷困難.
生命當中經歷很大的掙扎.
或者迷失的時候.
回想起來.
不是喔.
那時候哪個哪個導師.
哪個哪個.
他這樣經過的.
不如我試試祈禱吧.
這樣就可能改變了很多人的生命.

$^{801}$我們要專心靠主.
也希望能夠成為承傳.
可以承傳的財產.
最後也和大家分享一下.
昨天禮拜六早上.
我們所差出去.
去柬埔寨服侍的林于謙宣教士.
他回來述職分享.
當中讓我感動的是.
不單止他和太太.
在工場上經歷神的恩典帶領.
更加讓我感動的.
就是那天坐滿六樓禮堂的弟兄姊妹.
他們關心差傳.
支持宣教.
很希望他們能夠.
不單止將來有機會參加短宣.
更加有可能神差遣他們成為宣教士.
這個前方後方的協作.
就是我們需要傳承的福音使命的精神.
也希望最後以.
23章第8節的經文.
成為我們今天的祝福.
希望大家記得.
我們一起來讀.
只要照著你們到今天所行的.
專靠耶和華你們的神.
讓我們一起祈禱.
持我信實的天賦.
因為你向子民所說的.
「應許」沒有一句落空.
也祝我們如何帶領以色列人進入應許之地.
也在主耶穌基督裡面.
為我們帶來豐盛的生命.
完備的救恩.
神下求你赦免我們忘記你恩典的罪.
有時我們生活安穩的時候會鬆懈.
事奉贖責的時候.
很容易靠經驗.
或者世界的聲音吸引我們的時候.

$^{841}$我們會稍微轉離你.
求神你藉著聖靈.
鞭策我們.
提醒我們要謹守.
遵行你的說話.
不可偏離左右.
讓我們專心靠你.
緊緊連繫於你.
讓我們也有這份愛.
謹慎的心.
愛你的心.
讓我們在你裡面所承受的恩典.
承受你所賜的福氣.
能夠繼續傳承給後來信主耶穌基督的人.
讓大家都能夠可以.
忠心活在你的約裡面.
我們祈禱奉主耶穌基督的聖名祈求.
阿們.
\newpage



\section{士師記 2:12}
\label{sec:032227}
\textbf{聖經人聖經事:惹耶和華發怒}~~~~黃克勤牧師博士
\newline
\newline
連結: \href{http://www.church.com.hk/acms/content.asp?site=cdc\&op=show\&type=product\&code=032227\&layout=sermon}{\texttt{032227}} ~~~~ 語音日期: 2026-06-01
\newline
\newline
\hyperref[sec:032226]{\small{< < < PREV SERMON < < <}}
~
\hyperref[sec:index_chronic]{\small{[返順時目]}}
~
\hyperref[sec:032228]{\small{> > > NEXT SERMON > > >}}
\newline
\newline
士師記 2:12
\newline
\begin{longtable}{cl}
\hline
\hline
章節 & 經文 (和合本修訂版)\\
\hline
2:12 & \begin{tabularx}{0.7\textwidth}{X} 他們離棄領他們出埃及地的耶和華－他們列祖的神，去隨從別神，就是四圍列國的神明，向它們叩拜，惹耶和華發怒。 \end{tabularx} \\ \\
[1ex]
\hline
\hline
\end{longtable}
$^{1}$弟子們,稍後一日.
生命樹又跟大家看聖經人,聖經事.
爺爺和華發怒.
嘩!怎麼搞的.
仍然是《士斯記》歷史書.
是講述聖經的事的.
《士斯記》第二章.
從前約書亞打發以色列百姓去的時候.
令到他們進了迦南地.
然後各歸自己的地業.
佔據土地.
第六節就說.
他們進去.
拿了十二支牌.
拿了他們的地.
各歸自己的地業.
佔據了.
拿了.
第七節約書亞在世.
還有很重要的.
紫色字,重要的.
約書亞死後.
他在世與死後.
還有.
那些見爺和華.
為以色列人所行大事的長老.
還在的時候.
先是摩西.
然後到約書亞.
約書亞在那裡帶領著他們.
很順服,很好.
很強,信心很好.
領袖很好.
然後約書亞一百一十歲死了.
當時那些.
在曠野出世的那一代人.
是經歷到曠野的後半.
然後入迦南.
看到了,過約旦河.
他們沒有經過蘆葦海.

$^{41}$紅海,因為他們.
上一代已經死了在曠野.
他們經歷約旦河的.
大神跡.
嗎哪的神跡.
他們全部都經過.
然後石頭出水.
他們喝過.
然後約旦河的大神跡.
過河乾地.
再一個個城南方的黃.
南方的黃.
那些見耶和華為.
以色列人所行大事的長老.
還在的時候.
他們還在說這些事.
第一點要說的是.
百姓都事奉耶和華.
這件事很重要.
第一點要說的是.
因為有約書亞.
有當時的長老在那裡.
所以百姓.
都事奉耶和華.
很開心.
在當時.
戴得住.
耶和華的僕人聯的兒子.
約書亞110歲就死了.
第九節.
就葬了.
在哪裡呢?.
以法連他自己的成女.
在嘉實山的北邊.
約書亞死了.
好了.
那世代的人也都歸了自己的列祖.
即是說.
上面所說的.
那些見耶和華.

$^{81}$為以色列人所行大事的長老.
還在.
不過現在.
他們約書亞死了.
他們都歸了自己列祖.
即是說.
都死了.
出事了.
以色列開始離棄耶和華.
為什麼會這樣?.
因為.
一看就知道.
其實上一代.
沒有將下一代教好.
沒有將他們.
傳和承接得好.
你說沒有傳.
可能有傳.
但是下面承接不到.
又或者沒有傳.
下面沒得承接.
那世代的人都歸了自己的列祖.
後來有別的世代興起.
什麼別的世代?.
就是下一代.
即是包括了下一代再下一代.
那一代的人過去.
這一代新興的世代.
他們竟然.
這個士師記.
是描述他們.
不知道耶和華.
也不知道耶和華曾經做的事.
他不知道耶和華是不知道什麼?.
是不知道耶和華.
為以色列人所行的事.
都知道耶和華的.
不過以前做過這麼厲害的事嗎?.
不知道.
譬如我們中國.

$^{121}$很多近代.
很多偉人.
都為國捐軀犧牲.
很艱苦.
自己怎樣出力.
保護國家.
不知道.
我爸爸時代說.
日本人時代怎樣怎樣.
保家衛國.
不知道.
沒有看過.
更加不懂.
沒有人講歷史.
於是他們不知道.
於是以色列.
開始離棄耶和華.
離棄的意思是離開.
好了.
以色列人行耶和華.
眼中捍衛惡的事.
什麼惡事?.
是不是殺人放火?.
這個惡事是耶和華.
眼中最不喜歡的一件事.
就是去事奉.
事奉子敬拜.
朱巴力.
很多個巴力.
巴力不是只有一個嗎?.
巴力是一個.
迦南地的人.
像我昨天說的.
很多人.
很多族人他們沒有殺.
繼續在那裡.
神很憤怒.
他們將會成為你們.
臘上的荊棘.
和你們的網.

$^{161}$朱巴力.
他們的神全部叫巴力.
巴力的意思是Bael.
Bael就是主.
是當地迦南的土語.
Bael就是他們的主.
不同的.
當地的民族.
迦南的民族.
他們有不同的神.
他們都統稱為Bael.
離棄了.
令他們出埃及地的耶和華.
離棄了阿撒.
阿撒是離開.
不再跟隨.
離開.
曾經帶領他們.
出埃及的耶和華.
就是他們獵了的神.
去叩拜.
叩拜即是說.
祈福敬拜別神.
周四為獵國的神.
結果他們惹禍.
惹禍說法老.
法老.
Kaas.
讀錯了.
Kaas.
去惹禍法老.
其實法老在舊約裡面.
聖經希伯來文.
很多時候都是講鼻子.
很生氣.
噴氣.
這裡用一個字就是表示苦惱動氣.
生氣.
苦惱.
我花這麼多錢.

$^{201}$去栽培他兒子.
為什麼他現在這樣.
吸毒.
賭錢.
我花這麼多錢栽培他.
好苦惱兼生氣.
Kaas.
這個字是苦惱兼生氣.
苦惱兼生氣.
就是法老.
苦惱就生氣.
豈有此理.
這麼辛苦.
你竟然這樣.
惹禍說法老.
惹這個字是可以不翻譯的.
是特別加上去的.
就是他們四位的神.
就是四位的朱巴力.
於是不用看那個惹字.
惹禍說法老.
於是神覺得很苦惱.
你知道如果惹禍說.
苦惱.
生氣.
我們就慘了.
今天跟大家講到這裡.
結果是以色列人.
惹惹禍說法老.
苦惱兼動氣.
神就不會就這樣算了.
他就會.
做事.
要小心.
在這裡特別要講.
我們今代要向你的子女.
你的下一代.
傳語承.
讓他們知道.
惹禍曾經在我們身上做過事.

$^{241}$你不怕多講自己的經歷.
多講自己的經歷.
我怎樣經歷神.
讓小朋友聽聽.
他們都會覺得我們世世代代.
都是跟隨我們的神耶和華.
我都是信主耶穌.
我跟隨耶穌.
去傳語承.
以色列人在出埃及之後.
一直走到這裡.
傳承失了.
惹惹禍說法老.
最糟糕的是.
當地人可以做苦工.
他們拜神都很有趣.
有些美女.
有些牛羊留用.
這樣就糟糕了.
開始成為他們的.
往來與荊棘.
終於令他們拜了.
諸巴力.
惹惹禍說法老.
是一個警惕.
求主幫助我們.
看到這個警惕.
可能我們沒有傳承.
沒有再看神歷世歷代.
做了什麼.
於是我們的信仰失了.
愛對神的愛失了.
對神的敬拜與尊重失了.
終於我們被荊棘.
勒住綁住.
痛苦不堪.
終於我們被網羅捉到.
跌入魔鬼陷阱之中.
主求你警惕我們.
讓我們更加敏銳.

$^{281}$跟隨你.
誰聽祈禱奉耶穌基督名字.
阿姆.
\newpage



\section{路加福音 19:11-27}
\label{sec:032228}
\textbf{在王離去與再來之間:作忠心的僕人}~~~~梁成裕牧師
\newline
\newline
連結: \href{http://www.church.com.hk/acms/content.asp?site=cdc\&op=show\&type=product\&code=032228\&layout=sermon}{\texttt{032228}} ~~~~ 語音日期: 2026-05-31
\newline
\newline
\hyperref[sec:032227]{\small{< < < PREV SERMON < < <}}
~
\hyperref[sec:index_chronic]{\small{[返順時目]}}
~
\hyperref[sec:032229]{\small{> > > NEXT SERMON > > >}}
\newline
\newline
路加福音 19:11-27
\newline
\begin{longtable}{cl}
\hline
\hline
章節 & 經文 (和合本修訂版)\\
\hline
19:11 & \begin{tabularx}{0.7\textwidth}{X} 眾人正聽見這些話的時候，耶穌因為將近耶路撒冷，又因他們以為神的國快要顯現，就接著講了一個比喻， \end{tabularx} \\ \\ \relax
19:12 & \begin{tabularx}{0.7\textwidth}{X} 說：「有一個貴族往遠方去，為要取得王位，然後回來。 \end{tabularx} \\ \\ \relax
19:13 & \begin{tabularx}{0.7\textwidth}{X} 他叫了自己的十個僕人來，交給他們十錠銀子，說：『你們去做生意，直到我回來。』 \end{tabularx} \\ \\ \relax
19:14 & \begin{tabularx}{0.7\textwidth}{X} 他本國的百姓卻恨他，打發使者隨後去，說：『我們不願意這個人作我們的王。』 \end{tabularx} \\ \\ \relax
19:15 & \begin{tabularx}{0.7\textwidth}{X} 他得了王位回來，就吩咐叫那領了銀子的僕人來，要知道他們做生意賺了多少。 \end{tabularx} \\ \\ \relax
19:16 & \begin{tabularx}{0.7\textwidth}{X} 頭一個上來，說：『主啊，你的一錠銀子已經賺了十錠。』 \end{tabularx} \\ \\ \relax
19:17 & \begin{tabularx}{0.7\textwidth}{X} 主人對他說：『好，我善良的僕人，你既在最小的事上忠心，你有權柄管十座城。』 \end{tabularx} \\ \\ \relax
19:18 & \begin{tabularx}{0.7\textwidth}{X} 第二個來，說：『主啊，你的一錠銀子已經賺了五錠。』 \end{tabularx} \\ \\ \relax
19:19 & \begin{tabularx}{0.7\textwidth}{X} 主人也對這個說：『你管五座城。』 \end{tabularx} \\ \\ \relax
19:20 & \begin{tabularx}{0.7\textwidth}{X} 又有一個來說：『主啊！看哪，你的一錠銀子在這裡，我把它包在手巾裡存著。 \end{tabularx} \\ \\ \relax
19:21 & \begin{tabularx}{0.7\textwidth}{X} 我向來怕你，因為你是嚴厲的人：沒有放的，也要去拿；沒有種的，也要去收。』 \end{tabularx} \\ \\ \relax
19:22 & \begin{tabularx}{0.7\textwidth}{X} 主人對他說：『你這惡僕，我要憑你的話定你的罪。你既知道我是嚴厲的人，沒有放的也去拿，沒有種的也去收， \end{tabularx} \\ \\ \relax
19:23 & \begin{tabularx}{0.7\textwidth}{X} 為甚麼不把我的銀子存在銀行，等我來的時候，連本帶利都取回來呢？』 \end{tabularx} \\ \\ \relax
19:24 & \begin{tabularx}{0.7\textwidth}{X} 於是他對那些站在旁邊的人說：『把他這一錠奪過來，給那有十錠的。』 \end{tabularx} \\ \\ \relax
19:25 & \begin{tabularx}{0.7\textwidth}{X} 他們對他說：『主啊，他已經有十錠了。』 \end{tabularx} \\ \\ \relax
19:26 & \begin{tabularx}{0.7\textwidth}{X} 主人說：『我告訴你們，凡有的，還要給他；沒有的，連他所有的也要奪過來。 \end{tabularx} \\ \\ \relax
19:27 & \begin{tabularx}{0.7\textwidth}{X} 至於我那些仇敵，不要我作他們王的，把他們拉來，在我面前殺了！』」 \end{tabularx} \\ \\
[1ex]
\hline
\hline
\end{longtable}
$^{1}$今天我們分享的題目是.
在往未去和在來之間.
我們做忠心的僕人.
如果大家認識我.
知道我年輕的時候曾經在英國留學.
讀書的時候也去過不同的宿舍住.
我很感恩.
我也很疼我的爸爸媽媽.
所以無論我住不同的宿舍.
他們都對我愛惜和愛護.
每一次我要入住不同的宿舍.
他們都會為我預備不同宿舍的物品.
日常的用品.
衣服,被子,清潔的用品.
每一次他們離開的時候.
他們都會叮囑我的.
你要好好照顧自己.
房間要收拾得整齊.
要好好清理.
每一次都是這樣叮囑我的.
但是作為年輕人的我.
是十分冷惰和拐彎抹角的.
所以每一次同學來到我的房間.
這個不是真的我的房間.
這個是設計圖片而已.
房間拍不回來.
但是很多時候同學來到我的房間.
就說你的房間是這樣的.
就是房間很亂.
很骯髒.
東西到處放.
而且有時連用過的廚具.
都還沒拿去洗.
你可能會問.
你什麼時候清潔的呢?.
就是爸爸媽媽每一次來探望我的時間.
當我知道爸爸媽媽來探望我的時候.
我就會預留時間.
好好收拾一下房間.
清潔房間.

$^{41}$預備爸爸媽媽的到訪.
有一次我和同學在喝東西.
喝得很開心.
突然間我收到一個訊息.
就是我爸爸給我的訊息.
他說兒子.
我們現在到了倫敦.
明天來找你.
他就說現在到了倫敦.
明天來找我.
於是我就立刻扔下那杯東西.
跟同學說我要走了.
同學就說你這麼急.
你去哪裡?.
我就這樣回答.
我跟他說.
I need to go and wash.
同學就問.
就是說我要去清潔.
同學就會這樣問我.
Wash what?.
你清潔什麼?.
我就這樣回應他.
I need to wash the clothes.
I need to wash the room.
And I need to wash myself.
就是說我要去洗衣服.
我要清理好房間.
以及我要清潔好自己.
然後就扔下.
衣服到處放.
衣服又有味道.
所以一天之內要做完所有事情.
清潔衣服,房間.
包括自己.
為什麼?.
因為要預備家裡的到訪.
這是設計圖片.
也是設計圖片.
我的房間也應該不太乾淨了.

$^{81}$但是今天我想說什麼呢?.
其實我想說主耶穌基督已經復活升天.
並且我們的主耶穌將會榮耀地回來.
但是今天當我們去等候主耶穌基督再回來的時候.
我們要怎麼好好等候.
以及預備好我們的生命.
今天我們會一起去看《路加福音》.
十九章十一至二十七節一段經文.
我們一起學三個很重要的功課.
當我們知道主耶穌基督要榮耀地再回來的時候.
我們又怎麼去等候.
並且預備好自己呢?.
今天我們學的第一個功課.
是我們要有正確的國度期待.
我們首先一起讀十九章的十一至十四節.
我們一起讀.
一二三.
「眾人正在聽見這些話的時候.
耶穌因為將近耶路撒冷.
又因他們以為神的國快要顯出來.
就另設一個比喻說.
有一個貴就往遠方去.
要得國回來.
便叫了他的十個僕人來.
交給他們十錠銀子.
說:你們去做生意.
直等我回來 他本國的人卻恨他.
打發使者隨後去說.
我們不願意這個人作我的王」.
這裡一開始告訴我們.
眾人正在聽見這些話的時候.
聽見什麼話呢?.
大家記不記得上個星期.
陳牧師說哪一個人的記載?.
撒該!.
上個星期我們聽過陳牧師說撒該的記載.
正正就是耶穌講完有關撒該的言語之後的時間.
在這裡聖經告訴我們.
耶穌因為將近耶路撒冷.
將近耶路撒冷是什麼意思呢?.

$^{121}$這句話很重要.
因為原來將近耶路撒冷.
是帶著很濃烈的彌賽亞期待.
原來耶路撒冷是王作王國度顯現的地方.
特別那時候是逾越節的前後.
所以原來當時有很多人對耶穌有一個誤解.
這個誤解是什麼呢?.
他們以為神的國快要顯出來.
什麼為之誤解了神的國快要顯出來.
就是他們以為耶穌會來做地上的王.
所以他們以為耶穌快要到耶路撒冷了.
到了耶路撒冷他們就會趕走了羅馬軍隊.
推翻了羅馬帝國.
然後就在地上做猶太人的王.
我們就在耶路撒冷會看到這個地上猶太人的王.
但是這是當時人對耶穌的誤解.
因為我們知道耶穌在耶路撒冷做什麼的.
受死的.
所以耶穌為了矯正他們的誤解.
就講了一個比喻.
這個比喻一開始是說有一個貴就往遠方去.
貴就是出身顯貴的意思.
這裡要講一講當時羅馬政府的背景.
原來在當時的政治架構裡面.
如果你是一個地區性的封王.
就是你做一個地區性的王.
是需要去到羅馬.
接受羅馬軍王的授權.
就是說今天你要做這個地方的王.
因為整個地方是羅馬帝國統治的.
所以你要做這個地方的王.
你要首先去到羅馬.
接受羅馬軍王的授權.
然後你回來就做這個地方的王.
所以在耶穌時期的大希律.
和他的兒子阿基努.
都曾經要這樣往遠方去的.
往遠方去做什麼呢?.
就是得著這個王權回來統治這個地方.
而耶穌就是用這個比喻.

$^{161}$去矯正他們的誤解.
耶穌想要講的就是神的國已經臨到了.
但是還沒有完全顯現.
為什麼呢?.
因為耶穌要先受苦.
先離去.
然後再榮耀地回來.
耶穌講這個比喻.
猶太人會很明白的.
因為猶太人知道耶穌是說我會暫時離開.
然後會回來作王.
耶穌將會在十字架上犧牲.
被埋葬.
三日之後復活.
然後升天.
然後帶著國度和權柄.
榮耀地回來.
而更加重要的就是.
當人等待的時候.
人要做什麼呢?.
在這裡告訴我們.
主人交了十錠銀紙給僕人.
十錠銀紙是多少呢?.
會不會是十塊呢?.
十塊現在大概可以做什麼呢?.
可以想到的就是過去對面金端燒賣.
是的.
十塊可以買一串燒賣.
但是不要外賣.
口袋外賣是要監運的.
所以十塊就讓你買到一串燒賣.
那麼十錠銀紙是不是十塊呢?.
不是.
原來十錠銀紙是一千錢銀紙.
一千錢銀紙是多少呢?.
一千錢銀紙是大概一天的工資.
所以是差不多三年的工資.
就是說主人很信任僕人.
將很重要的事情交給了僕人.
就是說今天我們在等候的時間.

$^{201}$我們是需要去做神的工作的.
我們要運用我們的資源.
我們的恩賜.
我們的時間.
我們的關係.
去做神的工作.
我們怎麼做神的工作呢?.
我們第二點會說的.
但是今天我想說的就是.
我們首先要有一個正確的國度期待.
因為如果我們沒有一個正確的國度期待.
我們很容易放棄.
當時的猶太人也有一個錯的國度期待.
他們以為耶穌來了.
做政治的皇.
趕走羅馬軍隊.
推翻羅馬帝國.
然後在地上作王.
但是他們這個是錯誤的期待.
而今天我們也可能有一個錯誤的期待的.
我們今天可能怎麼會有一個錯誤的期待呢?.
我不知道你有沒有聽過人這樣跟你說的.
我也不知道你信的那個神是不是真的.
如果你信那個神是真的.
為什麼世界還有這麼多不公義啊?.
為什麼世界還有這麼多禍患和災難啊?.
為什麼會這樣啊?.
你的神是不是真的?.
就正正在過去社會運動.
以及疫情期間.
有一個年輕人來找我.
他說肥姓啊.
為什麼會這樣呢?.
如果耶穌是真的.
為什麼這個世界還有這麼多的不公義呢?.
如果這個神是真的.
為什麼我們會有疫情.
為什麼會有禍患.
為什麼會有災難.
但是我們要有一個正確的國度觀.

$^{241}$今天我們信了耶穌.
這個世界仍然會有不公義.
仍然會有災難.
仍然會有災害.
因為我們仍然等著耶穌回來.
不代表我們信了耶穌.
這個世界會馬上變好的.
同樣有時候我們信了主之後.
你可能會覺得為什麼我的生命沒有改變呢?.
曾經有一個肢體來問我.
他說梁牧師.
為什麼我信了主之後.
我的生命沒有改變呢?.
我就以為他覺得自己的生命沒有改變.
我就說你作為一個門徒.
你的生命就會改變.
我說不是說這些.
他說你不是說這些.
你是說哪些呢?.
為什麼我信了主之後.
我的工作仍然是這麼難打的呢?.
我說你想說這些.
然後他再問.
另外一個人問我.
是啊.
我信了耶穌之後.
生意都沒有變好的.
我曾經聽過一個肢體這樣分享的.
他信耶穌的時候他是做生意的.
他一信了耶穌.
他的生意就更加差了.
然後他就去找一個帶他信耶穌的肢體.
一個很資深的肢體.
就跟他說為什麼我信了主之後.
我的生意繼續會差呢?.
這個資深的肢體就跟他說.
他說你要搞清楚今天我們信的是一個什麼神.
如果今天你覺得一信了耶穌.
工作馬上變好.
然後你信了耶穌.

$^{281}$生意就馬上上揚.
我們有一個錯誤的角度觀.
有時候我們可能都會經歷我們的身體不好.
但是信了主之後.
都仍然沒有改變.
今天我們知道耶穌在我們生命裡面.
祂跟我們同行.
更加比我們永生的盼望.
但是有時候確實是未必會改變.
但是我們要知道耶穌是會再回來的.
有時候我們錯誤的角度觀可能是在教會的.
你有沒有試過可能想.
為什麼教會那些人.
為什麼教會那些人會這樣呢?.
我很喜歡Pastor Pam講一句話.
他說其實在教會裡面.
都是蒙恩的罪人.
我們都是罪人.
大家聽過我分享我的見證.
我很感恩.
我信了主之後.
我馬上戒賭了.
接下來6月4日我都會講的.
但是你知不知道有一樣東西.
我信了主我用了四年才戒賭的.
你知道是什麼嗎?.
髒話.
我戒了四年了.
你想一下我剛剛信主的時候.
我忍不住在教會講髒話.
接著有些人就會說.
哇!你的神是不是真的?.
你看看那個胖子.
他說髒話.
我們每一個人都還在誠信的路上.
我們都不完美.
我們要繼續去跟隨耶穌.
讓我們的生命誠信.
但是有時候我們有個錯誤的價值觀.
我們來到教會.

$^{321}$我們期待每一個人都是完美的.
你摸一摸你旁邊一個.
你問他他完不完美的?.
如果旁邊一個說完美的.
你舉手.
不是開玩笑的.
就是說.
是啊.
我們有個正確的角度觀.
我們仍然會面對世界的不公義.
我們仍然會看到世界會有戰爭.
禍患和災害.
但是我們要有信心.
因為我們知道耶穌將要榮耀地回來.
而當耶穌回來之前.
我們要作主的工作.
在進入第二點以先.
我想讀完第十四節.
我們一起讀.
一,二,三.
他本國的人卻恨他.
打發使者隨後去說.
我們不願意這個人作我們的王.
這裡在比喻裡面.
當貴州去了見王的時候.
這裡說他本國的人就恨他.
然後打發使者隨後去.
其實是做什麼呢?.
其實就是要跟君王說.
我們不要這個人做我們的王.
這件事情確實在耶穌的時期是發生了的.
剛才我們說阿基努.
阿基努就是希律的兒子.
他去了羅馬.
去接受羅馬的授權.
但是因為猶太人不喜歡他.
所以去投訴他.
說我們不想這個人做我們的王.
而耶穌也說有些人不想我做他的王.
而這些人的結局是怎麼樣呢?.

$^{361}$我們會看第三點.
但是我們到了今天的第二點.
就是原來當我們等候主回來的時候.
我們要按著神的託付去忠心做好神的工作.
我們一起從第十五節讀到第二十三節.
我們一起讀.
一二三.
他既得國回來.
就吩咐叫那領銀子的僕人來.
要知道他們做生意賺了多少.
頭一個上來說.
主啊!你的一定銀子已經賺了十定.
主人說:好!良善的僕人.
你既在最小的事上有忠心.
可以有權並管十座城.
第二個來說.
主啊!你的一定銀子已經賺了五定.
主人說:你也可以管五座城.
又有一個來說.
主啊!看啊!你的一定銀子在這裡.
我把它包在手巾裡存著.
我原是怕你.
因為你是嚴厲的人.
沒有放下的還要去拿.
沒有種下的還要去收.
主人對他說:你這惡僕.
我要憑你的口定你的罪.
你既知道我是嚴厲的人.
沒有放下的還要去拿.
沒有種下的還要去收.
為什麼不把我的銀子交給銀行.
等我來的時候.
連本帶利都可以要回來呢?.
在這裡.
一開始告訴我們這個貴咒.
德國回來了.
就是那些反對他的人.
未能夠成功攔阻他作為他們的王.
這個王回來的時候.
他就要問那些得著銀子的僕人.

$^{401}$究竟你們去做生意賺了多少呢?.
這裡告訴我們那個王.
他會回來問那些僕人.
你有沒有好好運用我給你的東西?.
我交託給你的東西.
你有沒有好好運用?.
因為原來這些財富是不屬於僕人.
僕人是要照著主人的心意去走的.
接著我們就看到有三個僕人.
我們看第一個.
第一個在第十六節.
第一個來說:主啊.
你的一定銀子已經賺了十錠.
你的一定銀子.
代表這個僕人很清楚知道銀子不是他的.
是屬於主人的.
而他是要按照主人的心意去運用這些銀子.
也強調這個僕人不是靠自己的本事.
但是他是按照主人的吩咐.
遵從主人的命令.
也相信主人是會回來的.
所以他很忠心去做.
他知道自己是一個管家的身份.
是暫時去保管和運用所託管的金錢.
而他也知道當主人回來的時候.
他是要向主人交帳的.
而當他這樣的時候.
主人的回應是什麼呢?.
主人說好.
良善的僕人.
他說你在最小的事上有忠心.
可以管十座城.
這裡主人說你是良善的僕人.
因為你在最小的事上有忠心.
你可能會問什麼是最小的事呢?.
最小的事就是財富的運用.
可能這一刻你會說.
原來財富運用也是小事.
那你知道什麼是大事嗎?.
大事是什麼?.

$^{441}$是神國的事.
所以在這裡.
主人說你在最小的事上.
就是你在運用我交給你的事上有忠心.
將來我會給你大事.
你可以管十座城.
就是你在神國裡面的責任.
但是更加重要.
這裡主人說.
當你忠心的時候.
你是在神國裡面有份.
而今天這個人也可以去分享主人的王權.
所以原來今天我們忠心.
就能夠在神國裡面有份.
而第二個又是怎樣呢?.
我們看第十八節.
第二個就說.
主啊.
你的一定銀子已經賺了五定.
我們要看看主人的回應.
主人的回應是什麼?.
是不是責怪他?.
有沒有搞錯?.
只有五定.
剛才那個十定.
不是.
主人不是這樣責怪他.
因為這個比喻的重點.
不是說我們人有多少的才能.
而是說我們有沒有對神忠心.
所以這個主人完全沒有責怪他.
他也沒有將這個人和之前那個僕人去比較.
因為這裡強調的是人的忠心.
而這個僕人也同樣去認定了主人的主權.
他說你的一定銀子.
我已經忠心這樣運用了.
所以今天我們有一個很重要的反思.
今天我們不是要跟別人比較的.
所以你跟隔壁那個說.
我不是跟你比較的.

$^{481}$你也不需要跟隔壁那個比較.
但是你要問自己.
今天神給我們的那份.
我們有沒有好好這樣運用.
而到了第十九節.
這個主人說.
你也可以管五座城.
這個忠心的僕人同樣得到獎賞.
同樣得到神國的福分.
同樣分享這個王的王權.
好了,我們看了兩個僕人.
接著我們看第三個僕人.
第二十節告訴我們有第三個僕人.
他來到這裡說.
主人,你的一定銀子在這裡.
因為我將它包在毛巾裡面存著.
他以為我只要原封不動地包著它.
表面上應該都沒有虧損,是不是?.
但是問題是什麼?.
問題是他完全沒有遵從過主人的吩咐.
主人的吩咐給一定銀子是做什麼的?.
做生意的.
他做什麼?.
他什麼都沒有做.
就這樣將它包在毛巾裡面.
所以好像沒有虧損.
但是這個僕人的問題就是.
他完全沒有做過主人要他做的事情.
他甚至沒有拿去投資.
他沒有放在銀行.
他完全沒有讓這一定銀子發揮任何功效.
而他的藉口是什麼呢?.
我們看二十一節.
他說我原是怕你.
因為你是嚴厲的人.
他將主人描繪成為一個苛刻的人.
他根本就不認識主人的心.
不僅如此.
他還這樣跟主人說.
他說我知道你是沒有放下.

$^{521}$還要去拿.
沒有種下.
還要去收.
這裡他在暗示主人.
你是不是不勞而獲?.
主人其實完全不是不勞而獲.
因為銀子是屬於誰的?.
是屬於主人的.
所以當他這麼說的時候.
主人的回應是什麼呢?.
主人說你這惡僕.
原來今天當我們什麼都不做的時候.
就變為你這惡僕了.
我們今天是不是惡僕?.
你看看你旁邊那個是不是惡僕?.
他有沒有什麼都不做?.
不用舉手的.
但是今天我們要搞清楚.
原來當我們什麼都不做的時候.
主人就會說你這惡僕.
但不僅如此.
主人說你既然這麼說我.
他不是承認他自己是這樣.
主人說你既然這麼說我.
我就用你所說我的去定你的罪.
為什麼呢?.
他說我要憑你的口定你的罪.
主人絕對不是承認.
但是主人說如果你覺得我是一個嚴厲的人.
你就快點主動去做生意吧.
你說我嚴厲.
如果我真的嚴厲.
你還不快快去做生意?.
接著主人同樣說.
如果你覺得我是那些沒有放下.
還要去拿.
沒有種下.
還要去收.
主人不是承認他自己是這樣.
主人不是這樣子.

$^{561}$但是主人說如果你真的以為我是這樣子.
你會更加主動去做事.
你起碼也將錢放在銀行等我回來.
確實當時將錢放在銀行也是有風險的.
但是做生意有沒有風險?.
有風險.
但是最重要.
這個主人就是說.
你的問題是.
你連走一步你都不願意.
所以今天我們千萬要想清楚我們的生命.
我們有沒有忠於做好神交給我們做的事情.
如果我們什麼都不做的時候.
我們就會被人稱呼利者惡僕.
所以今天我們要問我們自己.
我們有沒有按著主的託付.
去忠心做好神的工作呢?.
我們要問自己.
我們有沒有好好運用神給我們的恩賜.
但是你問我.
在雅庚堂裡面.
我看到很多很美麗的生命.
因為我看到弟兄姊妹.
你們在不同的崗位上面.
都為神擺上.
單單走進來崇拜.
我們有我們的迎新隊.
我們多謝他們好不好?.
我們有我們的影音隊.
多謝他們好不好?.
我們有我們的敬拜隊.
也就是少年敬拜隊.
我們多謝他們.
我們有我們的代禱隊.
我們多謝我們的代禱隊好不好?.
接著呢.
我們有 Comic Con 啊.
接著如果你走進去兒童崇拜.
你會看到兒童導師.
多謝我們兒童導師好不好?.

$^{601}$你在少年崇拜.
你會看到少年導師.
你在長者崇拜.
你會看到我們長者的義工.
接著你會看到愛倫社.
我們愛倫社的義工.
還有當我們每個禮拜返小組的時候.
我們的組長,副組長.
我們的 helper.
多謝他們好不好?.
如果我漏了.
等一下你來找我.
我請你吃東西好不好?.
但是我想說.
弟兄姊妹每個星期的預備.
其實他們都要提早到場.
他們也會預先祈禱.
他們很多更加是留後去收拾.
可能沒有人跟他們說過一聲多謝.
所以如果下次你看到福士的弟兄姊妹.
你跟他們說一聲多謝好不好?.
但是弟兄姊妹就是一年.
五年.
十年.
十二年了我們十二歲.
弟兄姊妹就是這樣衷心的去服侍.
同樣我們有沒有好好運用我們的時間呢?.
最近我認識一位門徒.
他的生命令我很感恩.
他是一個工作很忙的人.
但是他知道做了門徒.
我們會跟你一起去做時間管理的.
接著我們會一起去看看我們的時間運用.
接著有一天就說了.
他說我們不可以只有教會的朋友.
因為如果我們只有教會的朋友.
我們就沒有機會傳福音了.
所以當我們做時間管理的時候.
我們說我們要挖時間.
去繼續跟微信的朋友去建立關係.

$^{641}$因為如果你聽過Pastor Pam教.
我以前就試過這樣.
以前我試過打給我的朋友.
誰知道打給他的時候.
朋友就說.
Eugene.
是不是教會有報道會啊?.
為什麼?.
因為平時不見你打來的.
教會有報道會你就打來的.
我們平時要去跟朋友建立關係.
這個門徒他平時工作真的很忙.
但是他就每一個月.
挖一個時間.
去跟微信的朋友吃飯.
是每一個月一定挖定一個時間.
去跟一個微信的朋友吃飯.
你問我.
我看到他的生命.
這樣的時間運用很感恩.
另外我們又可不可以運用我們的關係呢?.
我認識另外一個門徒.
他工作也是很忙的.
但是他很主動在公司關心壓力大的同事.
你說我這麼忙.
放工我就馬上歸心似箭回家了.
但是他會主動留下去關心壓力大的同事.
為他祈禱.
然後邀請他們回來教會.
我們有沒有好好做好我們的工作呢?.
今天我們在我們的工作裡面.
我們有沒有用心.
並且誠實.
不做詭詐的事.
我們有沒有好好運用我們的金錢呢?.
我認識有一位姿態.
他的收入不算高的.
但是他就做好每一個月的十一奉獻.
不僅如此.
他會留下說這些是給宣教.

$^{681}$這些是給扶貧.
這些是給人有緊急的時候使用.
你問我很感動.
因為他不是說.
總之我洗剩多少有多就拿來奉獻.
洗得了就不奉獻.
不是的.
他是做好他的十一奉獻.
並且留下每一個月是給宣教.
每一個月是給扶貧.
每一個月還留下這些是給誰有緊急的時候去奉獻.
你問我很感恩.
我們有沒有好好地去管理我們的資源.
今天我們要問.
我們有沒有按著主的託付.
去忠心做好神給我們的工作呢?.
而今天.
去到今天分享的第三點.
我們要為最後的判決做決定.
我們從第24節讀到第27節.
我們一起讀.
一二三.
就對旁邊站著的人說.
奪過他者一定來.
給乃有十定的.
他們說.
主啊.
他已經有十定了.
主人說.
我告訴你們.
凡有的還要加給他.
沒有的連他所有的也要奪過來.
至於我那些仇敵.
不要我作他們王的.
把他們拉來.
在我面前殺了罷.
這裡也要講講當時一些羅馬政府的背景.
原來當時在羅馬的政制裡面.
如果有一個地方官員他做得不好的話.
他的處分就是被革職了.

$^{721}$而他管理的地方.
就會交給其他人管理.
而這個人的結局就是被放逐.
以後不可以再回到自己的家園和國土.
在這裡講了一件事情.
原來這個處罰.
是會讓他跟他的家永遠隔絕的.
所以在這裡.
耶穌也用這個比喻說.
今天我們也要為我們的永生做決定.
今天我們也要為我們面臨最後的審判做決定.
而今天的結局是.
究竟我們在神國裡面有沒有份.
如果我們看看二十五節.
當主人把惡瀑的一丁給了十丁那個的時候.
那些人就說.
主啊.
祂已經有十丁了.
但是這裡是在講.
原來我們要為我們的永生做一個決定.
我們要為我們永生究竟在神國裡面有沒有份去做一個決定.
因為如果我們不這樣去跟隨耶穌的時候.
就好像這個瀑人那樣.
祂有的一丁也給了那個有十丁.
代表祂在神國裡面沒有份.
而今天我們要做一個決定.
就是究竟我們在神國裡面有沒有份.
所以第二十六節這裡講.
主人說.
我告訴你們.
凡有的還要加給他.
沒有的連他所有的也要奪過來.
這裡給我們兩個生命很重要的提醒.
第一個提醒.
是原來我們現在要為我們的生命負責任的.
我們要為我們將來的審判.
在今生做一個負責任的決定.
而第二件事.
今天如果我們不忠心地去做.
不做好主託付給我們的時候.

$^{761}$將來審判的時候.
我們要受到神的責罰.
而第三就是趁我們還有時間的時候.
就為我們的永生去做決定.
過去的一個星期.
因為是一個長週末.
我不知道大家有沒有筆上.
有的舉手.
我也筆上了我先舉手.
你也筆上你舉手好不好?.
我也知道有不少弟兄姊妹也筆上的.
我不知道你有沒有經歷回來的時候.
你的經歷是怎麼樣的.
因為我也看到很多圖片.
知道如果你長週末筆上的時候.
你回來也不容易的.
我自己又怎麼樣呢?.
我自己就以為自己很厲害.
因為我也計劃了筆上.
所以我早就在幾個星期前.
已經買了高鐵的票了.
我是一早買了星期一晚上七點多.
從深圳回到香港的票.
我以為自己怎麼樣呢?.
很穩固了.
到處跟別人說我不怕的.
我買了高鐵票了.
我會回到香港的.
在六點多的時候.
我就在一個商場裡面.
於是我就在計劃我的回程.
我就按那個程式.
那個程式就說.
從這裡去高鐵站23分鐘.
我就覺得.
真的十分穩固了.
現在六點.
去了23分鐘.
於是我就計劃過了.
趁還有45分鐘的時候.

$^{801}$我就跟老婆說.
這個時候我們就坐計程車去高鐵站.
就很穩固了.
我老婆就很相信我了.
於是我們就上了計程車.
上了計程車之後我就按了程式.
突然間那條線從綠色變成了紅色.
你知道變成紅色是什麼嗎?.
極級塞車.
然後從23分鐘轉到45分鐘.
哇.
要坐45分鐘.
我說.
這件事有點危機了.
可能趕不上那班車.
趕不上那班車是不會有第二班的了.
那一晚就留在深圳了.
你可以排關口的.
但是我太太已經跟他說了.
我不會跟你排關口的.
很緊張地看著計程車.
我就跟計程車司機說.
黃敏鋼車.
他說知道.
然後他就不停地繞路.
但是他繞到哪裡呢?.
他也只有一句話.
啊.
這麼多.
極級塞車.
無論怎麼繞都塞.
然後我看著時間一分一秒地過去.
我真的很擔心究竟我能不能上車呢?.
於是我就跟司機說.
還差多遠呢?.
他說沒辦法.
這麼多.
這次真是怎麼辦呢?.
但是他很好.
到了那個位置.

$^{841}$他就跟我說.
你現在塞車走?.
我說好啊.
謝謝你.
然後我跟我太太拿了行李.
就走去高鐵站了.
我沒試過這樣走去高鐵站.
我試過這樣走去關口.
我也沒試過這樣塞車塞去高鐵站.
拿了行李就走了.
到了那裡.
還有時間入閘的.
但是如果你坐過高鐵你就知道了.
其實他只是給你十五分鐘入閘.
就是他開之前十五分鐘才開的.
但是我進去還有幾分鐘.
我記得還有很多人還沒入閘.
然後我就跟太太說.
其實也不能入閘的.
為什麼還有這麼多人還沒入閘呢?.
然後就在這個時候.
開了另外一個閘.
然後那些人.
我就跟老婆說.
跟著那裡.
最後感恩了.
我們上到車了.
安全地七點多就回到香港了.
然後就看到一些朋友還沒下來的.
就看到十一點多關口的照片了.
我想說什麼呢?.
今天我們也不要以為自己真的這麼穩固.
我們有沒有反思我們的生命.
我們不要以為我們已經拿定了天國的入場券了.
我們要問清楚自己.
我們要反思我們今天我們的生命.
我們有沒有為最後的判決作好決定.
而最後當我們看第二十七節.
這裡說.
至於我那些仇敵.

$^{881}$就是剛才我們十四節那裡讀過.
那些不想這個人作王的人.
最後他們的結局是把他們拉來.
在我面前殺了吧.
在這裡我們要做兩個很重要的反思.
第一件事我們要問自己.
究竟我們有沒有忠心去承擔主給我們做的事.
今天當我們說我們等候耶穌回來的時候.
當耶穌交託了給我們不同的恩賜.
我們有時間.
我們有資源.
我們有關係.
究竟我們有沒有去做神的工作.
而第二件事給我們的反思.
是今天我們是不是只是名義上的基督徒.
但是原來裡面實際是在拒絕神的主權.
我們可能都洗了禮.
成為了會友.
但是我們在我們的婚姻關係.
在我們的金錢運用.
在我們的事業選擇上面.
我們從來沒有問過神的心意.
不單止是這樣.
有時候我們明知道聖經教導是這樣.
但是我們反方向去走.
我們都有基督徒的表面.
你看一看你旁邊那個.
他一定有基督徒的表面.
因為今天他坐在教會.
他一定有基督徒的表面的.
但是今天我們是不是只是表面那些.
主啊!主啊!.
你問我.
我每一個星期都反思的.
為什麼?.
因為我每一個星期都站在那裡.
敬拜的時候.
主啊!.
你剛才有沒有看到我舉手.
主啊!主啊!.

$^{921}$我們每一個都口說主啊!主啊!.
但是我每一個星期都問自己.
究竟今天我的心是不是真的願意.
耶穌做我的王.
因為今天如果我們只是口說主啊!主啊!.
但是我們生命不是真正被耶穌作王.
將來我們的結局.
都是跟神國沒有分的.
所以在總結的時候.
我們請敬拜罪.
在今天這條聖經裡面.
我們看到有三種人.
第一種人.
就是將那筆錢包在手巾裡面.
他們以為包在手巾裡面就沒有吃虧.
所以我們可能都是坐在那裡.
什麼都不做.
以為這樣是OK的.
但原來今天聖經告訴我們.
這樣是不OK的.
另外第二.
有一群人.
他們可能像我們這樣.
嘴巴都承認耶穌是主.
但原來心裡不想耶穌作王.
我們今天都要問.
我們是不是都是口說主啊!主啊!.
但原來我們的心沒有讓耶穌作王.
第三種人.
就是那種在小事上忠心的人.
我們不要跟別人比較.
我們不要看旁邊的人在做什麼.
反過來我們看看今天神給了我們什麼.
我們有沒有忠心去為神擺上.
而今天我們要問自己.
我們有沒有好好將神給我們的東西.
好好地運用.
而今天我們決定做哪一種人.
\newpage



\section{士師記 6:1-40}
\label{sec:032229}
\textbf{士師記析讀 (8) 要有憑據}~~~~梁永善牧師
\newline
\newline
連結: \href{http://www.church.com.hk/acms/content.asp?site=cdc\&op=show\&type=product\&code=032229\&layout=sermon}{\texttt{032229}} ~~~~ 語音日期: 2026-06-03
\newline
\newline
\hyperref[sec:032228]{\small{< < < PREV SERMON < < <}}
~
\hyperref[sec:index_chronic]{\small{[返順時目]}}
~
\hyperref[sec:032230]{\small{> > > NEXT SERMON > > >}}
\newline
\newline
士師記 6:1-40
\newline
\begin{longtable}{cl}
\hline
\hline
章節 & 經文 (和合本修訂版)\\
\hline
6:1 & \begin{tabularx}{0.7\textwidth}{X} 以色列人又行耶和華眼中看為惡的事，耶和華就把他們交在米甸手裡七年。 \end{tabularx} \\ \\ \relax
6:2 & \begin{tabularx}{0.7\textwidth}{X} 米甸的手戰勝以色列；以色列人躲避米甸人，就在山中挖洞穴，挖洞建營寨。 \end{tabularx} \\ \\ \relax
6:3 & \begin{tabularx}{0.7\textwidth}{X} 每當以色列人撒種之後，米甸、亞瑪力和東邊的人都上來攻打他們， \end{tabularx} \\ \\ \relax
6:4 & \begin{tabularx}{0.7\textwidth}{X} 對著他們安營，毀壞那地的農作物，直到迦薩，沒有給以色列留下食物，牛、羊、驢也沒有留下。 \end{tabularx} \\ \\ \relax
6:5 & \begin{tabularx}{0.7\textwidth}{X} 因為那些人帶著他們的牲畜和帳棚上來，像蝗蟲那樣多；人和駱駝無數，都進入境內，毀壞全地。 \end{tabularx} \\ \\ \relax
6:6 & \begin{tabularx}{0.7\textwidth}{X} 以色列因米甸的緣故極其窮乏，以色列人就呼求耶和華。 \end{tabularx} \\ \\ \relax
6:7 & \begin{tabularx}{0.7\textwidth}{X} 以色列人因米甸的緣故呼求耶和華的時候， \end{tabularx} \\ \\ \relax
6:8 & \begin{tabularx}{0.7\textwidth}{X} 耶和華就差遣先知到以色列人那裡，對他們說：「耶和華－以色列的神如此說：『我曾領你們從埃及上來，從為奴之家出來， \end{tabularx} \\ \\ \relax
6:9 & \begin{tabularx}{0.7\textwidth}{X} 救你們脫離埃及人的手，脫離一切欺壓你們之人的手。我從你們面前趕出他們，把他們的地賜給你們。 \end{tabularx} \\ \\ \relax
6:10 & \begin{tabularx}{0.7\textwidth}{X} 我對你們說，我是耶和華－你們的神。你們住在亞摩利人的地，不可敬畏他們的神明，但你們卻不聽從我的話。』」 \end{tabularx} \\ \\ \relax
6:11 & \begin{tabularx}{0.7\textwidth}{X} 耶和華的使者到了俄弗拉，坐在亞比以謝族約阿施的橡樹下。約阿施的兒子基甸正在醡酒池那裡打麥子，為了躲避米甸人。 \end{tabularx} \\ \\ \relax
6:12 & \begin{tabularx}{0.7\textwidth}{X} 耶和華的使者向基甸顯現，對他說：「大能的勇士啊，耶和華與你同在！」 \end{tabularx} \\ \\ \relax
6:13 & \begin{tabularx}{0.7\textwidth}{X} 基甸對他說：「主啊，請容許我說，耶和華若與我們同在，我們怎麼會遭遇這一切事呢？我們的列祖告訴我們：『耶和華領我們從埃及上來』，他那奇妙的作為在哪裡呢？現在耶和華卻丟棄了我們，把我們交在米甸人的手掌中。」 \end{tabularx} \\ \\ \relax
6:14 & \begin{tabularx}{0.7\textwidth}{X} 耶和華轉向基甸，說：「去，靠著你這能力拯救以色列脫離米甸人的手掌。我豈不是已經差遣了你嗎？」 \end{tabularx} \\ \\ \relax
6:15 & \begin{tabularx}{0.7\textwidth}{X} 基甸對他說：「主啊，請容許我說，我怎能拯救以色列呢？看哪，我這一支在瑪拿西支派中是最貧寒的，我在我父家又是最微小的。」 \end{tabularx} \\ \\ \relax
6:16 & \begin{tabularx}{0.7\textwidth}{X} 耶和華對他說：「我與你同在，你就必擊敗米甸，如擊打一個人。」 \end{tabularx} \\ \\ \relax
6:17 & \begin{tabularx}{0.7\textwidth}{X} 基甸對他說：「我若在你眼前蒙恩，求你給我一個證據，證明是你在跟我說話。 \end{tabularx} \\ \\ \relax
6:18 & \begin{tabularx}{0.7\textwidth}{X} 求你不要離開這裡，等我回來，將供物帶來，供在你面前。」他說：「我必等你回來。」 \end{tabularx} \\ \\ \relax
6:19 & \begin{tabularx}{0.7\textwidth}{X} 基甸去預備一隻小山羊，用一伊法細麵做了無酵餅，將肉放在籃子裡，將湯盛在壺中，帶到他那裡，在橡樹下獻上。 \end{tabularx} \\ \\ \relax
6:20 & \begin{tabularx}{0.7\textwidth}{X} 神的使者對基甸說：「將肉和無酵餅放在這磐石上，把湯倒出來。」他就照樣做了。 \end{tabularx} \\ \\ \relax
6:21 & \begin{tabularx}{0.7\textwidth}{X} 耶和華的使者伸出手裡的杖，杖頭一碰到肉和無酵餅，就有火從磐石中出來，吞滅了肉和無酵餅。耶和華的使者就從他眼前消失了。 \end{tabularx} \\ \\ \relax
6:22 & \begin{tabularx}{0.7\textwidth}{X} 基甸見他是耶和華的使者，就說：「哎呀！主耶和華啊！因為我真的面對面看見了耶和華的使者。」 \end{tabularx} \\ \\ \relax
6:23 & \begin{tabularx}{0.7\textwidth}{X} 耶和華對他說：「安心吧，不要怕，你不會死。」 \end{tabularx} \\ \\ \relax
6:24 & \begin{tabularx}{0.7\textwidth}{X} 於是基甸在那裡為耶和華築了一座壇，起名叫「耶和華沙龍」。這壇至今還在亞比以謝族的俄弗拉。 \end{tabularx} \\ \\ \relax
6:25 & \begin{tabularx}{0.7\textwidth}{X} 那夜，耶和華對基甸說：「你要把你父親的公牛，就是那七歲的第二頭公牛取來，並拆毀你父親為巴力築的壇，砍下壇旁的亞舍拉， \end{tabularx} \\ \\ \relax
6:26 & \begin{tabularx}{0.7\textwidth}{X} 在這堡壘頂上整整齊齊地為耶和華－你的神築一座壇，將第二頭公牛獻為燔祭，用你所砍下的亞舍拉當柴。」 \end{tabularx} \\ \\ \relax
6:27 & \begin{tabularx}{0.7\textwidth}{X} 基甸就從他僕人中選了十個人，照耶和華吩咐他的做了。他因怕父家和本城的人，不敢在白天做這事，就在夜間做。 \end{tabularx} \\ \\ \relax
6:28 & \begin{tabularx}{0.7\textwidth}{X} 城裡的人清早起來，看哪，巴力的壇被拆毀，壇旁的亞舍拉被砍下，第二頭公牛獻在築好的壇上， \end{tabularx} \\ \\ \relax
6:29 & \begin{tabularx}{0.7\textwidth}{X} 就彼此問：「這是誰做的事呢？」他們尋找查訪之後，就說：「這是約阿施的兒子基甸做的事。」 \end{tabularx} \\ \\ \relax
6:30 & \begin{tabularx}{0.7\textwidth}{X} 城裡的人對約阿施說：「把你的兒子交出來，我們要處死他，因為他拆毀了巴力的壇，砍下了壇旁的亞舍拉。」 \end{tabularx} \\ \\ \relax
6:31 & \begin{tabularx}{0.7\textwidth}{X} 約阿施對站著敵對他的眾人說：「你們是為巴力辯護嗎？你們要救它嗎？誰為它辯護，就在早晨把誰處死吧！巴力如果是神，有人拆毀了它的壇，就讓它為自己辯護吧！」 \end{tabularx} \\ \\ \relax
6:32 & \begin{tabularx}{0.7\textwidth}{X} 所以那日人稱基甸為耶路巴力，意思是：「他拆毀了巴力的壇，讓巴力與他爭辯吧。」 \end{tabularx} \\ \\ \relax
6:33 & \begin{tabularx}{0.7\textwidth}{X} 那時，所有的米甸人、亞瑪力人和東邊的人都聚集在一起，過了河，在耶斯列平原安營。 \end{tabularx} \\ \\ \relax
6:34 & \begin{tabularx}{0.7\textwidth}{X} 耶和華的靈降在基甸身上；他吹角，亞比以謝族都聚集跟隨他。 \end{tabularx} \\ \\ \relax
6:35 & \begin{tabularx}{0.7\textwidth}{X} 他派使者走遍瑪拿西，瑪拿西人也聚集跟隨他。他又派使者到亞設、西布倫、拿弗他利，他們也都上來會合。 \end{tabularx} \\ \\ \relax
6:36 & \begin{tabularx}{0.7\textwidth}{X} 基甸對神說：「你如果真的照你所說的，藉我的手拯救以色列， \end{tabularx} \\ \\ \relax
6:37 & \begin{tabularx}{0.7\textwidth}{X} 看哪，我把一團羊毛放在禾場上，若單是羊毛上有露水，遍地都是乾的，我就知道你必照你所說的，藉我的手拯救以色列。」 \end{tabularx} \\ \\ \relax
6:38 & \begin{tabularx}{0.7\textwidth}{X} 一切果然發生了。次日早晨基甸起來，把羊毛擰一擰，從羊毛中擠出露水來，裝滿一碗的水。 \end{tabularx} \\ \\ \relax
6:39 & \begin{tabularx}{0.7\textwidth}{X} 基甸又對神說：「求你不要向我發怒，我再說一次，讓我用羊毛再試一次，但願羊毛是乾的，遍地都有露水。」 \end{tabularx} \\ \\ \relax
6:40 & \begin{tabularx}{0.7\textwidth}{X} 這夜，神也照樣做，遍地都有露水，只有羊毛是乾的。 \end{tabularx} \\ \\
[1ex]
\hline
\hline
\end{longtable}
$^{1}$今晚我們會來看時詩記.
我相信很多基督徒對時詩記裡面.
特別有兩個時詩是很熟悉的.
一個就是《機殿》.
一個就是《參宣》.
因為《機殿》我們知道.
它以少勝多.
大家回去想一想.
究竟《機殿》帶著三百人對抗多少敵軍.
但是事後《機殿》的發展.
很少人會繼續去看.
特別在崇拜的講道裡面.
是沒有怎麼提及的.
但是《機殿》從成功慢慢慢慢就滑落了.
就失敗了.
但是我們因為賢經.
所以每一章都會去研究的.
所以我們會由《時詩記》第六章.
看到第九章.
看《時詩記》《機殿》的興.
然後看它的衰.
特別是家庭的慘劇.
作者用了很大的篇幅.
來記述去蒙召的過程.
如果你留意到.
用一章的聖經.
反而講打仗是較為少的.
也是今天我們所探討的.
我不知道大家有多少人會知道.
我如何蒙召做傳道人.
如果今天夠時間的話.
我也想跟大家分享一下.
我如何被主呼召成為一個傳道人.
現在已經過了五十年了.
希望有一天我離開世界的時候.
都可以立旨傳召太神奇妙的大靈.
我們很快的看回今天第六章裡面.
主題就是要有憑據.
你留意到《機殿》.
它不斷要拿憑據.

$^{41}$每一次都是和二字有關.
二和夜晚黑.
我們遲些會來看.
我們看回《持史記》裡面很特別的.
大家都很熟悉的.
就是那個悲哀的循環.
因為以色列人有行謠.
眼中害迷惡的事.
所以神就將它交在米殿人的手中.
有七年這麼久.
第六章第一節又講到.
不單單作者這樣說.
神親自來聲討他們.
第六章第十節.
麻煩阿爾灣.
又對你們說:我是耶和華你們的神.
你們住在阿摩利人的地.
不可敬畏貪民的神.
你們竟不聽從我的話.
六章一節.
「竟不聽從我的話」.
「聽」.
你留意到.
這個「聽」不單單是你聽了頭腦的知識.
當日那班子民都有對誡命的知識.
但是有沒有實踐出來.
這個才是很重要的.
我們就看看究竟他們所面對的是甚麼挑戰.
舌詩第六章.
我們由第二節讀到第五節.
麻煩阿爾灣.
以色列人有行耶和華.
眼中害迷惡的事.
耶和華就把他們交在米殿人手裡七年.
米殿人壓制以色列人.
以色列人因為米殿人就在山中.
刮穴,刮洞,建造城寨.
以色列人每逢殺種之後.
米殿人,阿瑪力人和東方人都上來攻打他們.
對著他們安然毀壞土產直到家殺.

$^{81}$沒有給以色列人留下食物.
牛,羊,鹿也沒有留下.
因為那些人帶著生畜,象,蟠來.
象王蟲那樣多.
任我諾托無數都進入國內毀壞全地.
如果你留意到米殿人其實跟他們有一點血緣關係.
因為摩西的岳父也是米殿人.
他們不斷地來搶掠他們的農作物.
所以他們很害怕.
就會來刮一些山洞.
然後把他們收藏在當中.
他們不僅搶掠他們的農作物.
甚至會怎樣呢?.
就是去破壞他們的焦土政策.
搶掠完之後你留意到.
還會來破壞他們的土產.
而拿了食物之後.
不要忘記.
牛,羊,鹿是他們最重要的牲口.
牛,羊是用來獻祭的.
以及用來食用的.
不是所有的牛,羊都是用來獻祭的.
鹿是用來運輸的.
但是將這些動物,這些生畜全部掠走.
食物就沒有了.
大家可以想像到.
當時的情況是多麼惡劣.
他們在經濟上是受著嚴重的損失的.
而以色列人很諷刺.
他們那時候是敬拜巴力的.
巴力就是希望他們可以有豐收.
但是現在他們不僅沒有好的收成.
就算有也被人搶了.
更加缺乏了平安穩妥的生活.
他們渴望得到的東西是完全得不到的.
神就審判他們.
米顛人欺壓了他們七年.
終於他們就向神有呼求了.
神是滿有憐憫的心.
於是就帶他們去呼召基丁.

$^{121}$在第六章的第七節裡面.
就有記載這件事.
我們先來看《釋述記》第六章第六到第十節.
我們要留意.
待會大家看的時候.
基丁的蒙召其實和摩西是很相似的.
好 我們先看這段經文.
以色列人因米顛人的緣故極其困乏.
就呼求耶和華.
以色列人因米顛人的緣故呼求耶和華.
耶和華就差遣先知到以色列人那裡對他們說.
耶和華以色列的神如此說.
「我曾領你們出從埃及上來.
出了遺牢之間.
救你們脫離埃及人的手.
並脫離一切欺壓你們之人的手.
把他們從你們面前趕出.
將他們的地賜給你們.
又對你們說.
我是耶和華你們的神.
你們住在那麼利人的地.
不可敬畏他們的神.
你們竟不聽從我的話」.
你有沒有留意到神是差派先知責備他們.
不是單單是事先的作者.
再一次提到那個歷史事實.
就是他們在埃及的時候.
他們在遺牢之地.
被埃及人虐待.
那時某個程度他們都不是拜偶像.
但是在政治上成為廉價勞工被虐待.
但現在他們是自討苦吃.
因為他們敬拜巴黎.
所以神就來責備他們.
你們竟然不聽從我的話.
當他們向神父呼求的時候.
神就憐憫的心發動.
興起基丁成為事先幫助他們.
但基丁在一個極不尋常的地方裡.
蒙神的呼召.

$^{161}$第十一節開始.
讓我們看到基丁如何給神父呼召.
十一,十二兩節.
(故此我要答).
Sorry 這是出埃及的.
剛才我說的就是因為他們受苦.
就向神來哀求.
記不記得我剛才開始說.
與以色列人出埃及.
就是基丁和摩西的避兆很相似.
以色列人也是在苦難當中向神呼求.
神就興起摩西.
我們先看出埃及的第三章第九到第十節.
(現在以色列人的哀聲達到我耳中).
(我也看見埃及人怎樣欺壓他們).
(故此我要打發你去見法老).
(使你可以將我的百姓以色列人從埃及領出來).
是不是一樣呢?.
就是他們的哀聲呼求達到神的耳中.
所以神就呼召了摩西.
去帶領以色列人出埃及.
而上面也看到.
這幫以色列人因為被米顛人欺壓.
第六節尾的時候就呼求耶和華.
所以耶和華就差遣他的使者.
待會我們會看到.
就是去呼召基丁.
好 我們來看看基丁的避兆.
基丁的避兆你留意的地方.
是很特別.
是絕不尋常的地方.
我們看看第十一到第十二節.
(耶和華的使者到了俄弗拉).
(坐在阿比爾謝族人約阿斯的像樹下).
(約阿斯的兒子基丁正在走在那裡打麥子).
(為要防備米顛人).
(耶和華的使者向基丁顯現對他說).
(大能的勇士啊 耶和華與你同在).
這裡就說到.
他說在約阿斯的像樹下.

$^{201}$記得上次我們看到.
底波拉都在一棵樹的下面.
來去做審判的.
那棵樹被稱為底波拉的樹.
這次是阿比爾謝族人約阿斯的像樹下.
留意.
約阿斯是基丁的父親.
約阿斯竟然是一個拜巴力的祭司.
竟然自己的父親是一個拜偶像.
他說在走在那裡打麥子.
我們看看俄弗拉是甚麼地方.
俄弗拉可能在什瓦利亞以南大概十公里的地方.
是屬於嗎哪西之派.
約阿斯的像樹.
因為像樹的樹蔭很大.
樹大很遮蔭.
所以成為加林人敬拜神.
以及求問神的地方.
也被視為這棵是神聖的樹.
而基丁的父親.
應該是巴力祭壇裡面的主持人.
去守護這棵樹.
所以被稱為約阿斯的像樹.
可見神揀選人未必一定以他的家庭背景作為條件.
我自己也很大一個證明.
我當年家裡是一個非常迷信的家庭.
拜很多的神.
我記得跟大家分享過.
我讀幼稚園開始的時候.
是要拜孔子的.
家裡要帶我去文武廟拜文帝.
但是我竟然可以成為傳道人.
是否很想像不到呢?.
原來基丁的父親是拜巴力的祭師.
接著我們看看.
當時的人打麥在什麼地方呢?.
是在空曠的地方.
麥就鋪在這些空曠的地方.
有牛或者雷就拖著牌子滾.
然後把穀和殼弄出來.

$^{241}$因為是空曠的地方.
大風吹就可以吹走那些穀.
就是穀殼.
然後他們就會收成.
然後撥那些麥子.
但是因為基丁怕被菲律賓人來帶入去搶.
所以他們就在酒鑿裡面.
酒鑿是什麼呢?.
就是一個山洞裡面.
通常是由盤石造成的洞.
裡面藏著葡萄.
就用腳來帶入去踩那些葡萄.
讓那些汁液流出來.
然後藏在瓶子裡.
在陰涼的地方裡面.
山洞裡面就是陰涼的.
讓它可以慢慢變成酒.
所以酒鑿就是在山洞極深的地方.
但是你留意到基丁現在不是在.
平時的空曠地方裡面打麥.
它是在酒鑿裡面.
那時候就容易躲藏.
因為它很害怕.
你留意到它是怕的.
怕米電人來帶入去搶它的東西.
雖然基丁躲藏不想被人發現.
但是神的使者在這個時候都找到它.
向它顯現.
並且向它問安.
但是這個問安是什麼呢?.
我不知道是真的讚賞它還是諷刺它.
大能的勇士啊基丁.
神跟你在一起.
但是這個所謂的勇士其實是搞底了.
躲在酒鑿裡面.
當然我們也可以說神看到基丁裡面的特質.
所以這樣稱呼它.
但是我們一直看以後的經文.
我們並不覺得基丁是一個非常勇敢的人.
當然最終它後來有這個勇氣.

$^{281}$帶著三百人去面對多如海沙的敵人.
但是整個過程上.
其實基丁不是一個很勇敢的人.
作者暗示基丁很可能像摩西那樣.
可以帶領以色列人出困境.
問題在於基丁有沒有像摩西那樣.
對以色列人的後代有長遠的影響呢?.
某個程度都有.
不過是負面的影響.
記不記得第十二節下面說到.
耶和華與你同在.
這個英雄同樣呼召摩西的時候.
出現過的.
但是當基丁看到耶和華的使者跟他說.
耶和華與你同在的時候.
他馬上來反駁.
我們看看第六章第十三節.
基丁說:主啊 耶和華若與我們同在.
我們何至遭遇這一切敵事呢?.
我們的列祖不是向我們說.
耶和華領我們出埃及上來嗎?.
他那樣奇妙的作為在哪裡呢?.
現在 他的丟棄我們.
將我們交在米顛人手裡.
丟棄了我們 將我們交在米顛人手裡.
你有沒有留意到.
基丁並不像我們之前所看的先知那樣.
譴責以色列人拜偶像 犯錯.
今天也是類似我們的情況.
很多時候我們有些問題的時候.
有困難.
我們不是去檢討.
為什麼會有這些事情發生在我們身上.
而是將責任推卸在神那裡.
或者推卸在人那裡.
甚至埋怨神為什麼不去保護我們.
不去眷顧我們呢?.
所以其實當我們遇到不如意的事情的時候.
當我們遇到衝擊的時候.
首先要做的不是去埋怨神.

$^{321}$而是去檢視為什麼會有這件事發生呢?.
這樣才對我們會有幫助.
否則我們永遠都是在埋怨當中.
不去檢討.
對我們的成長是不會有任何的幫助的.
好了 我們看看神怎麼去回應.
第六章第十四節.
耶和華官向基丁說.
你靠著你這能力去從米顛人手裡拯救以色列人.
不是我差遣你去的嗎?.
你留意到「官漢」.
這個「官漢」就是神彌達人一直留意著他.
然後特別的彌達人去指出.
不是你啊.
雖然我挑選你 但不是你.
不是你的能力可以來到成功的.
而是要靠著我.
他就問基丁要求他作出一些肯定的回覆.
但是那時候基丁知道是什麼呢?.
覺得自己很渺小.
不能夠承擔責任.
你看看他拒絕的藉口是什麼呢?.
基丁說:主啊 我有何能拯救以色列人呢?.
我家在嗎哪西之派中是至貧窮的.
我在我父家是至微小的.
你留意到他馬上的彌達人去拒絕.
正如當天神呼召摩西的時候.
摩西馬上拒絕了.
我們看看摩西拒絕的是什麼?.
再看第三章第十一節.
摩西對神說:我是什麼人竟能去見法老.
將以色列人從埃及領出來呢?.
是不是呢?相對都是一樣的.
當天神呼召摩西的時候.
摩西說:我是什麼人啊?.
我怎麼可以來見法老呢?.
當然摩西不是不知道自己是什麼人.
他的意思就是我什麼身份.
不要忘記當時摩西是什麼?.
是一個流亡海外的被通緝的兇手.

$^{361}$埃及王是一個.
那時候是國勢最盛.
以及是管理著很多地方的王.
一個流亡海外的兇手去見最強的王.
哪裡有這份信心啊?.
所以我是什麼人呢?.
基澱也說我有什麼能力拯救以色列人呢?.
但這是一個藉口.
我在馬來西亞支派中是最貧窮的.
我在富家裡面是最微小的.
如果你留意到「大能勇士」這個詞.
也可以譯作「大財主」.
記不記得我剛才說.
基澱的父親是巴黎的祭司.
他家裡其實也是一個有錢的家庭.
我們稍後會看到.
基澱晚上叫他做一些事情.
是帶著十個僕人去的.
不是所有.
只是其中一部分.
他家裡至少有十個以上的僕人.
家裡會不會很窮呢?.
是絕對不會的.
只不過他在推搪而已.
稍後有時間講我自己被神呼召的情況.
但我很清楚先講一講.
我記得在大會堂音樂廳裡面.
見到神樂園的聚會裡面.
我就很清楚感受到.
之前已經有了.
不過我更加肯定.
神呼召我 要我做傳道人.
我就好像基甸 好像摩西那樣.
所以我有屬靈的特質.
我馬上拒絕了.
我拒絕是甚麼呢 就說.
哎呀 不行的.
我又喜歡玩 又喜歡開玩笑.
都不像傳道人.
那你叫隔壁那個吧.

$^{401}$我突然間很謙卑.
我不行的.
但你在這個刻子裡面.
當時也也問也沒有.
覺得自己也很行的.
反應又夠快.
但是要我做傳道人 當然拒絕了.
一定有些藉口.
所以你留意到.
基甸這個是一個藉口來的.
在這個時候.
而且米殿人經常攻擊子民受苦之後.
他還有自己的生口.
你不要忘記.
一會兒我們會看到.
他會帶去獻給神.
可見基甸的家族是一個顯赫的家族.
而且是一個有錢人家.
絕對不是一個貧窮的人.
所以他說我在馬來西亞是自貧窮的.
這是藉口來的.
往往他們都是以自己的能力.
和背景去看神的事工.
但是神怎麼看呢?.
第六章第十六節.
耶和華對他說:.
「我與你同在.
你就必擊打米殿人.
如擊打人 一人一樣」.
「我與你同在 你就不用擔心」.
你擊打他們.
縱使他們有很多軍兵.
但好像擊打一個人一樣.
不要忘記神同樣將這個應許.
來給摩西的.
記不記得當摩西說我是甚麼人.
後來他繼續拒絕.
然後他又說那些人不知道你是怎樣.
神又將自己自幼並幼向他啟示.
然後他說那些人不信我的.

$^{441}$神就給摩西行那些神蹟.
接著就是這件事了.
第三章第十二節.
「神說:我必與你同在.
你將百姓從埃及領出來之後.
你們必在這山上伺服我.
這就是我打發你去的證據」.
記不記得?.
「我必與你同在」.
再一次出現了.
不僅是給摩西.
給摩西的接班人約書也一樣有這個應許.
約書是第一章第九節.
「我豈沒有吩咐你嗎?.
你當剛強壯膽.
不要懼怕也不要驚慌.
因為你無論往哪裡去.
也會說你的神必與你同在」.
原來神與事奉者同在這個應許.
不僅是給基丁.
不僅是給摩西.
更加不僅是給約書.
其實是給我們每一個事奉神的人.
只要我們願意倚靠神的時候.
神都會與我們同在.
在同一個時代裡面.
神都沒有放棄祂自己所屬的人.
所以倚靠神來幫助我們.
所以很在乎那些被召者.
他願不願意用信心.
踏上這條應許,滿有福氣的路.
這才是最重要的.
每當我們面對一些我們不願意做的事情.
我們必定有很多堂而皇之的藉口.
好像很正確.
甚至很謙卑來推卻.
譬如有些人上去煮了很久.
叫他洗禮.
「還未可以的,我還未成熟.
我還有很多漏閘.

$^{481}$所以不行,如果我洗禮的時候.
我怕影響了別人對神的看法」.
這就好像很正確.
但是你有沒有去改?.
未叫你洗禮的那一刻.
你又不會去想這些事情.
所以這只是一個藉口.
其實裡面就是不想承擔.
好了.
基澱就要要求有憑據.
膽怯,小信的基澱.
面對神的使者的宣告.
他沒有信心.
他就要一些很實質的憑據.
「我要有,我才可以考慮」.
所以就向這個使者發出這個要求.
然後神也答應.
留意在第七節裡面.
「求你給我一個證據」.
好.
我們看看這段較長的經文.
麻煩了Yvonne.
17到21節.
「基澱說:我若在你眼前無恩.
求你給我一個證據.
使我知道與我說話的就是主.
求你不要離開這裡.
等我歸回.
將禮物帶來供在你面前.
主說:我必等你回來.
基澱去預備了一隻山羊羔.
用一衣髮,細面作了毛餃餅.
將肉放在湯內.
把湯盛在壺中.
帶到橡樹下.
獻在使者面前.
神的使者吩咐基澱說.
將肉和毛餃餅放在這盤石上.
把湯瀝出來.
下午這樣行了.

$^{521}$耶和華的使者伸出手內的匠.
匠頭挖了肉和毛餃餅.
就有火從盤石中出來.
燒盡了肉和毛餃餅.
耶和華的使者也就不見了」.
你留意到基澱就說我要拿憑據.
你給我一個憑據.
我就知道原來你是神所差派的.
你不要走.
等我回來.
這個使者說好 我等你回來.
基澱就預備了一隻新羊羔.
一衣髮,細面.
做毛餃餅.
也有些湯.
不要忘記.
原來這樣是需要一段時間的.
他不是已經有個雪櫃.
等肉在那裡.
到時開個雪櫃拿出來.
然後在微波爐加熱.
即是說基澱要去宰一隻羊.
所以你見不見到.
在米顛人裡面不斷去攻擊他們的時候.
他仍然是有一些生畜.
所以他家裡絕對不是一個窮的家庭.
他要宰羊.
然後要去潔淨他.
還要竹壇用火去焚燒.
我們會見到.
又要去做毛餃餅.
所以他用了很長時間.
還有一衣髮,細面等於多少呢?.
一衣髮至少十公升.
甚至超過二十公升.
他就說我打算拿這些來給你吃.
來帶去送給你.
但這是超過一個人可以用的份量.
都頗禮遇神的使者.
神的使者就將基澱的食物.

$^{561}$叫他放在盆石上.
然後用匠頭打下去.
結果就焚燒了這些肉和麵.
還有不要忘記.
是之前淋了一些湯在上面.
這時候其實這些禮物.
給這個使者的禮物.
就變成了向神呈獻的濟生.
有沒有留意到這個證據呢?.
就好像摩西被蒙召一樣.
摩西怎樣被神孽呼召呢?.
就是他見到荊棘著火.
但荊棘卻沒有燒爛.
於是就引起他的好奇.
走過去看.
然後神就說這是聖地.
你要脫下你的鞋.
同樣這次的呼召都是很特別的.
不要忘記.
在《出拉急記》第三章第二節裡面.
就記載了摩西的蒙召.
「耶和華的使者從荊棘裡火焰中.
向摩西顯現.
卻沒有燒毀」.
所以你留意到.
兩者都有相似的地方.
昨天在這本記載第十八章裡面.
都重複出現.
神給先知以利亞去捉了個壇.
巴力先知就叫到聲嘶力竭.
但是他們的祭肉都沒有被火燒.
但是以利亞就把祭肉放在壇上.
然後還要用水淋它.
水就蓋過那些肉.
然後就呼喚.
結果有火從天而降.
不僅燒了那些祭肉.
連祭壇都燒爛了.
但是有沒有留意到.
這就顯示出.

$^{601}$當日他們拜的巴力是完全廢的.
是沒有用的.
但是這次很不同.
這次如果我們看回的時候.
原來是有火從盤石出來.
燒盡了肉和無價餅.
不是由天降火.
是在盤石裡面出火.
我想大家有沒有去過日本吃一隻牛肉.
我很喜歡在日本吃的牛肉.
有一種牛肉是怎樣的呢?.
它燒到很燙的一塊類似磚.
即是一塊石.
然後它就切開了那些牛肉.
你自己在這塊很燙的石上.
就燙那些牛肉.
牛肉不可以太熟.
你帶去吃的.
即是那塊石已經有很多熱力在當中.
但這次在盤石裡面.
你帶去出火.
燒了祭肉.
所以你看到它是很特別的.
所以機電實在有潛力.
同摩西甚至以利亞.
成為一個很偉大的拯救者.
機電知道神親自來帶去向他顯現.
他就很害怕.
很害怕神來帶去會消滅他.
因為人不可能來帶去見到神的使者.
但神給機電一個應許.
你放心.
不用懼怕.
你必不至死.
他及後還為耶和華捉了一座壇.
就起名為「耶和華沙龍」.
神賜平安的意思.
我們看看這段經文.
《士師記》第6章第22到第24節.
「機電見他是耶和華的使者.

$^{641}$就說:愛在主耶和華啊.
我不好了.
又為我的面看見耶和華的使者.
耶和華對他說:你放心.
不要懼怕.
你必不至死.
於是機電在那裡為耶和華捉了一座壇.
起名叫「耶和華沙龍」.
就是「耶和華賜平安」的意思.
這壇在亞比爾濟族的黃佛拉直到如今.
你留意到那個使者已經不見了.
機電很害怕.
他死了.
耶和華神說:我不好了.
我現在見到神你的使者.
神母向他直接顯現.
但神用一種方法跟他說.
你不用怕.
你放心.
你不必死.
於是就捉了這座壇.
透過這個神跡.
讓機電感受到神的能力.
和祂的恩典.
祂捉了一座壇.
讓祂的家族和後人知道神同在的恩典.
然後神又要他做一件事.
就是要他回歸起點.
這座壇就是「耶和華沙龍」.
就是平安.
所以以色列人的平安.
是由機電的個人經歷開始.
因此.
很深體會到領袖的生命是很重要的.
這個起點是甚麼呢?.
這是一個千古不變的定律.
機電願意稍稍信服.
但神沒有立刻叫他召集以色列人去打這場仗.
神是要機電從他的家庭開始.
來回歸信仰的起點.

$^{681}$作為一個領袖.
他必定要先對付好自己.
潔淨好自己.
甚至他的家庭.
記不記得《新約聖經》講到教會裡面的牧者.
監督.
甚至是執事.
他都要好好處理好自己的家庭.
如果連自己的家都管不好的時候.
他怎樣來帶人去管理教會呢?.
所以神就要他來對付家裡的偶像.
正如記不記得雅各在比奴伊勒那裡.
來帶人去見神之前.
他都要他的家人將那些偶像帶去除去.
今次神要機電拆毀他父親為巴黎所做的祭壇.
還要砍下祭壇旁邊的木偶.
還要在盆石上.
這可能就是他們的祭壇.
來獻上第二隻的牛給耶和華.
好 所以當我們認識到對神了解多了之後.
我們就有更大的使命和信心.
來面對更大的挑戰.
這是非常重要的.
我們看看《事實》第六章第25到26節.
「當那夜耶和華吩咐機電說.
你取你父親的牛來.
就是那七歲的第二隻牛.
並拆毀你父親為巴黎所做的祭壇.
砍下祭壇旁邊的木偶.
再借盆石上整整齊齊地.
為耶和華你的神做一座壇.
將第二隻牛獻為飯祭.
用你所砍下的木偶作柴」.
神要機電獻第二隻牛成為飯祭給耶和華.
我們有沒有留意到這個基載教是特別的.
為甚麼要獻第二隻牛呢?.
神很清楚指明是要選第二隻.
是七歲的牛.
我們都知道過去我們獻祭給神的時候.
在律法上是要求生的.

$^{721}$是最好的.
但為甚麼會七歲這麼大呢?.
大家想一想.
為甚麼會這樣?.
這似乎刻意要用拜巴黎的宗教儀式.
來諷刺他們.
好 我們看看《列王記》上.
當時以利亞和那些巴力先知鬥法的時候.
是怎樣的?.
都是要將牛切好的.
好 我們看看《列王記》上18章23到24節.
「當給我們兩隻牛獨.
巴力的先知可以挑選一隻.
切成塊紙.
放在柴上不要點火.
我也預備一隻牛獨.
放在柴上也不點火.
你們求告你們神的名.
我也求告耶和華的名.
那降火獻影的神就是神.
眾民回答說這話甚好」.
大家記得在加密山上的對峙.
巴力獻第一隻牛.
但不成功.
因為沒有火降下來.
但以利亞用第二隻牛.
見到神來帶他去降火.
所以第一隻牛是給巴力.
第二隻牛是給神的.
同樣這件事.
雖然是後期的事.
都有這個意義.
因為第一隻牛可能獻給了神.
所以第二隻牛.
還有.
剛才叫大家想一想.
為何是獻七歲的牛呢?.
記不記得.
他們被米顛人欺壓了多少年.
然後向上呼求.

$^{761}$是七年.
所以這隻牛是他們開始被欺壓的時候.
來帶牠出生的.
有這個特別的意思.
所以基澱必定要為自己的家族.
連續七年的來帶牠去背道.
向神來認罪.
以色列人的罪惡.
藉著基澱敬拜的行動.
得到赦免.
基澱要拆毀巴力的祭壇.
還要砍斷亞瑟拉的木偶.
顯示迦南諸神的虛假.
沒有能力來拯救自己.
因為神沒有辦法去保護自己.
那些用木製造的木偶.
就成為獻祭給神.
火燒的材料.
所以藉著這個勝利.
讓我們知道未來的仗.
勝負已經定了.
但是不要忘記我講.
卑殿其實是個有些懼怕.
有些害怕的人.
他怎樣呢?.
我們看看他怎樣做事.
《慈悲記》第6章第27節.
「基澱就從他僕人中挑了十個人.
照著耶和華吩咐他的行了.
他因怕富家和本城的人.
不敢再罷就行者事.
就在夜間行了」.
看到我一開始講的時候.
夜間是不斷在經文裡面出現的.
現在基澱他白天不敢.
去破壞巴西祭司巴力的祭壇.
於是晚上他去做.
顯出他的膽怯和小信.
但是神要基澱在盤石上做祭壇.
這個盤石是要塞在高處的地方.

$^{801}$不能夠隱藏的.
目的就是要彌大祂彰顯神的榮耀.
也要迫使基澱坦然地去面對群眾.
和面對反對勢力.
所以他起初是躲在酒窖裡面.
在山洞裡面.
但是神要在顯現的高處裡面去設壇.
除了顯示神的榮耀之外.
也是要塗造基澱.
正如患了三十八年血流病的那個女人.
是摸耶穌衣裳的睡紙.
得到康復的.
那麼耶穌說是誰摸我?.
難道耶穌不知道嗎?.
但是耶穌要逼這個人出來面對他自己.
所以祂現在要迫使基澱表明祂自己的立場.
祂終於成功地去除了那些偶像.
神沒有因為祂的膽怯而怪責祂.
因為我們很深信.
生命和信心的成長是一個歷程.
今天神是不斷忍耐等候我們的成長.
這是很感恩的.
今天的事奉也是由最基本做起.
在事奉的過程當中.
當然會有不順利.
會有失敗.
但是如果我們心仍然向著神的時候.
神繼續會祝福我們.
這是一個很重要的屬靈教導.
因為當時神的指紋.
並不是否定耶和華的真實.
記不記得當日以利亞也這麼說.
你們心懷異議要到何時呢?.
你們一邊拜巴力.
一邊又來拜耶和華神.
他們想左右逢源.
他們相信巴力也會賜給他們祝福.
但是神是要我們絕對的.
專一的來跟隨祂.
所以神就要機電拆除這個偶像.

$^{841}$再在這裡之上重建這個祭壇.
所以以色列人得的拯救.
不是脫離米電人的核制度開始.
原來的起點是什麼?.
是回到神的身邊.
這兩個起點都很重要.
就是機電要面對自己家庭.
拜偶像這件事.
要破除這個偶像.
以色列人要得到的自由.
不是靠打贏仗.
而是回到神的裡面.
所以這個很重要.
今天我們同樣有沒有正視.
我們屬靈生命的那個忍疾呢?.
我們不會否認神的真實.
但是有時候我們卻是愛世界.
跟世界糾纏著.
我們會不會喜歡看劇集.
多於看神的話語呢?.
原來今天有時候神只是在我們生命當中.
一個很次等的地位.
但如果我們不除去我們內裡的那種隱惡.
我們很難給神使用的.
好了.
機電拆除了這個偶像之後.
就陷他爸爸在一個困難當中.
但這次這個爸爸真的愛自己的兒子.
他就要維護自己的兒子.
我們看看《時詩記》第6章28到30節.
「城裡的人清早起來見巴力的壇拆毀.
壇旁的木偶坎下.
第二隻牛獻在新竹的壇上.
就彼此說:這是是誰作的?.
他們訪查之後就說:.
這是約阿斯的兒子機電作的.
城裡的人對約阿斯說:.
將你兒子交出來好似死他.
因為他拆毀了巴力的壇.
坎下壇旁的木偶」.

$^{881}$當伊斯利文看見他們所拜的偶像被拆毀的時候.
就要追究是誰做的?.
我們看不見.
機電挺身而出.
說是我做的.
然後跟著為神來做見證.
沒有.
或者他已經躲藏在某些地方了.
這顯示機電怎會是一個大能的勇士?.
伊斯利文宣告要殺死機電.
可見他們拜偶像已經將他們的心靈眼睛弄盲了.
他們寧願殺死自己的同胞.
來維護他們對巴力的敬拜.
但他們又因為受欺壓.
於是又要向神來呼求 這是不是矛盾?.
他們同時事奉兩個主.
作為一個拜巴力的祭司.
他自己所拜的偶像竟然被拆毀.
又要將他兒子交出來 他怎麼辦呢?.
這次約阿斯就有他的智慧.
我們看時史記第六章第31節.
約阿斯回答詹卓 攻擊他的眾人說.
「你們是為巴力爭論嗎?你們要救他嗎?.
誰為他爭論?趁早將誰致死?.
巴力若果是神 有人拆毀他的壇.
讓他為自己爭論吧」.
在這裡我們看到機電的父親終於改邪歸正了.
他為了維護自己的兒子甚至要改變自己的信仰.
他再次站在耶和神那邊.
約阿斯很可能身為巴力的祭司.
但面對自己的兒子被殺的時候.
愛子情切 他也不能不站在兒子那一方.
你留意到約阿斯就為了兒子.
他為兒子辯護了.
愛子情切 他不得不站在兒子那邊.
為兒子講說話 跟他們爭辯.
他很聰明的 你有沒有留意到.
他絕對沒有提到機電所做的事.
也不提這班子民 特別是俄弗拉人的要求.
約阿斯沒有將自己的兒子交給別人.

$^{921}$他甚至竟然宣佈拜巴力的人就應該要面對死亡.
讓機電和巴力的敬拜者安危.
放在你們各自敬拜的神手裡面.
為甚麼呢?.
他將問題提到一個很高層次.
如果巴力是一個真神.
巴力就親自去對付機電.
誰拆毀他的像.
這個神那麼有能力就去動他.
這時候是一個很實際的挑戰.
用甚麼我來動手呢?.
結果那些以色列人在無言反駁之下.
就將機電稱為耶路巴力.
表示約阿斯的論述回應.
耶路巴力就是讓巴力為自己爭辯的意思.
所以機電就多了一個新名.
《士師記》第6章第32節.
所以當日人稱機電為耶路巴力.
以施說他拆毀巴力的壇.
讓巴力與他爭論.
沒錯.
結果你留意到.
不是機電自己出來.
是他爸爸維護著他.
但帶出一個真理.
如果誰是神是真的.
由神親自去聲討他.
結果後來人們就稱他為耶路巴力.
最後就是大家很熟悉的.
就是再次的尋求一些憑據.
經過這一次激烈的潔淨行動之後.
機電終於接受了神的呼召.
成為以色列人的領袖.
《六章》第34,35兩節.
耶和華的靈降在機電身上.
他就摧割阿彼以者族都聚集跟隨他.
他打發人走遍馬來西地.
馬來西人也聚集跟隨他.
又打發人去見阿切人,西布倫人,拿伐他利人.
他們也都出來與他們會合.

$^{961}$神的靈降在機電身上.
所以阿彼以者群族就跟隨者.
就摧割了機電.
原來跟隨機電的還有馬來西.
阿切,西布倫和拿伐他利.
這些人也跟隨這個領袖.
所以機電的成功絕對不是來自他本身的魅力.
或者他父親的社會地位.
但很清楚地說到.
是因為神的靈在他身上.
機電看見這麼多人來跟隨.
應該有信心吧?.
但他仍然缺乏信心.
他竟然再次要求神給他一個憑證.
就是神要來帶他做甚麼呢?.
就是要神給一些證據給他.
大家很熟悉.
就是把羊毛放在地上.
就說「神,明知道羊毛要濕,地下要乾」.
神真的答應了.
大家又想一想會不會有問題呢?.
於是有一個更加深的難度是甚麼呢?.
就是羊毛要乾.
但地下要濕透.
因為地下是平地.
日間會有霧水.
所以羊毛吸收了霧水.
所以會覺得羊毛吸收了霧水.
那就濕了.
但如果地下乾的話.
就比較難了.
即是容易了.
但如果羊毛是乾的.
全部地下是濕的.
因為中東地方是很曬的.
霧水一吐出來就會沒有了.
但他現在第二次要求是甚麼呢?.
就是周圍都濕的.
這個難度就高很多了.
所以讓我們看到一件事是甚麼呢?.

$^{1001}$就是他沒有這份信心.
再要憑據.
神會要求人超越他在世上所需要的力量.
來服侍他.
他要求蒙召的人.
用自己現有的力量來依靠神.
所以口講無憑.
他就指著神來起誓.
神本身就是一個保證.
只有軟弱的人.
他一再忍耐.
看到神忍耐這些軟弱的人.
來建立他們.
可見基甸經常懷著一個質疑的心.
去面對他.
但這些行動很多時候成為一些信徒尋求神心意的一個範例.
但記住聖經的作者.
並不是用這件事.
讓我們在這方面作出指示.
他只是引用.
在神的律法當中唯有一種證據.
去求神的指引在舊約裡面.
就是污淩和討命.
但現在我們都不需要.
因為聖靈在我們當中來工作.
求神彰顯祂的旨意的時候.
我們需要懷著一個謙卑.
求神憐憫的心來向神要求.
神呼召每一個人成為祂的兒女.
但在生命的階段當中.
可能發出第二種特別的呼召.
即是今天我們每一個參與查經的.
或者在YouTube聽的弟兄姊妹.
你成為神的兒女.
這是第一個呼召.
但可能有些特別的心意在某些人身上.
好像在我身上.
神要我來的.
祂奉獻為全套人.
這個計劃我們怎樣去回應呢?.

$^{1041}$我再說.
就是不會要求以超越事公所需要的能力來服侍他.
今天神要求蒙召的人.
是以自己現有的力量.
但再依靠神來服侍他.
如果我願意全心全意.
獻給主,被主使用的心智.
就是神所體寵.
我們一定要確認與神同在.
神與我們同在.
這就是成功的因素.
所以求神來幫助我們.
神呼召我們事奉祂.
我們有時候往往都會感到懼怕.
怕自己力有不逮.
又怕自己困難.
這就成為我們事奉的絆腳石.
但很多時候我們要求很多很多的引證才肯踏上.
在機電,蒙召,戰勝的過程當中.
我們見到神的忍耐.
慢慢賠了一個祂合用的器皿.
所以今天我們所信的神.
都要求我們全然獻上.
毫不保留.
甚至讓我們在家人和朋友當中.
陷入兩難的困境裡面.
神絕對不容讓敵人存在.
令我們雙手將敵人彌帶入去.
命令我們將敵人彌帶入去擊倒.
今天神都要將我們帶出隱密的走廊.
然後將我們帶上戰場.
這位要利我冒險的神.
稱呼我們做大能的勇士.
是希望讓我們見證神的奇妙.
讓我們能夠有能力擊敗長期以來的壓迫者.
就是我們的仇敵撒旦.
所有人得到自由.
很願意在我們信心不足之下.
求神憐憫我們.
給我們一些憑據.

$^{1081}$記住人和神的關係當中.
懷疑神比懼怕敵人更具破壞性.
很希望我們能夠得到憑據之後.
使我們有更深的經歷.
接受神的呼召.
更願意你在神的磨練下.
成為一個真正大能的勇士.
更願意我們用信心接受神給我們痛在的應許.
有信心面對前面的挑戰.
我們一起來祈禱.
天父多謝你讓我們有機會學習.
讓我們要的不只是一些憑據.
要的是神你自己和我們在一起.
讓我們面對你給我們的呼召.
不是懷著膽怯的心.
而是我們知道主要你同在.
我們就一定有這個力量.
祈求主你自己幫助我們.
透過你的話語給我們更多的反省.
求主你自己親自祈禱祝福.
多謝你.
聽我們在你面前的禱告.
祝福日後收看YouTube的弟兄姊妹或朋友.
都讓我們同樣得到祂的幫助.
祈禱拉奉主德性的名.
阿們.
\newpage



\section{詩篇 61:1-8}
\label{sec:032230}
\textbf{領我到更高的磐石 (華)}~~~~高志翔牧師
\newline
\newline
連結: \href{http://www.church.com.hk/acms/content.asp?site=cdc\&op=show\&type=product\&code=032230\&layout=sermon}{\texttt{032230}} ~~~~ 語音日期: 2026-05-31
\newline
\newline
\hyperref[sec:032229]{\small{< < < PREV SERMON < < <}}
~
\hyperref[sec:index_chronic]{\small{[返順時目]}}
~
\hyperref[sec:032231]{\small{> > > NEXT SERMON > > >}}
\newline
\newline
詩篇 61:1-8
\newline
\begin{longtable}{cl}
\hline
\hline
章節 & 經文 (和合本修訂版)\\
\hline
61:1 & \begin{tabularx}{0.7\textwidth}{X} 神啊，求你聽我的呼求，留心聽我的禱告！ \end{tabularx} \\ \\ \relax
61:2 & \begin{tabularx}{0.7\textwidth}{X} 我心裡發昏的時候，要從地極求告你。求你領我到那比我更高的磐石， \end{tabularx} \\ \\ \relax
61:3 & \begin{tabularx}{0.7\textwidth}{X} 因為你是我的避難所，是我的堅固臺，使我脫離仇敵。 \end{tabularx} \\ \\ \relax
61:4 & \begin{tabularx}{0.7\textwidth}{X} 我要永遠住在你的帳幕裡！我要投靠在你翅膀下的隱密處！（細拉） \end{tabularx} \\ \\ \relax
61:5 & \begin{tabularx}{0.7\textwidth}{X} 神啊，你聽了我所許的願；你將產業賜給敬畏你名的人。 \end{tabularx} \\ \\ \relax
61:6 & \begin{tabularx}{0.7\textwidth}{X} 求你加添王的壽數，使他的年歲存到世世代代。 \end{tabularx} \\ \\ \relax
61:7 & \begin{tabularx}{0.7\textwidth}{X} 願他在神面前永遠坐在王位上，求你預備慈愛和信實保佑他！ \end{tabularx} \\ \\ \relax
61:8 & \begin{tabularx}{0.7\textwidth}{X} 這樣，我要歌頌你的名，直到永遠，天天還我所許的願。 \end{tabularx} \\ \\
[1ex]
\hline
\hline
\end{longtable}
$^{1}$弟兄姊妹大家早安.
有沒有覺得這個題目太偉大了.
當我們遇見困難患難的時候.
竟然可以讓我們從目前的絕望看到永恆的盼望.
到底是什麼樣的心境.
什麼樣的能力讓我們可以有這樣的遠光呢.
我們一同來禱告.
父神我們讚美你.
你是與我們同在的神.
願主的聖靈今天早上再次地接著你的話語.
讓我們看見信仰的真實樣貌.
讓我們每一位都能夠來依靠你.
主啊能夠來親近你.
把我們真實的重擔卸在你的面前.
主啊讓我們在你的話語裡面得到安息.
得到力量.
願聖靈來感動我們.
保守孩子的口.
能夠說的和你的心意.
禱告奉靠耶穌基督的名.
你是與我們同在的神.
願主的聖靈再次地接著你的話語.
讓我們每一位都能夠來依靠你.
你是與我們同在的神.
願主的聖靈來感動我們.
保守孩子的口.
能夠說的和你的心意.
願主的聖靈來感動我們.
保守孩子的口.
願主的聖靈來感動我們.
保守孩子的口.
願主的聖靈來感動我們.
保守孩子的口.
願主的聖靈來感動我們.
保守孩子的口.
願主的聖靈來感動我們.
保守孩子的口.
願主的聖靈來感動我們.
保守孩子的口.
願主的聖靈來感動我們.

$^{41}$保守孩子的口.
願主的聖靈來感動我們.
保守孩子的口.
願主的聖靈來感動我們.
保守孩子的口.
願主的聖靈來感動我們.
保守孩子的口.
願主的聖靈來感動我們.
保守孩子的口.
願主的聖靈來感動我們.
保守孩子的口.
願主的聖靈來感動我們.
保守孩子的口.
願主的聖靈來感動我們.
保守孩子的口.
願主的聖靈來感動我們.
保守孩子的口.
願主的聖靈來感動我們.
保守孩子的口.
願主的聖靈來感動我們.
保守孩子的口.
願主的聖靈來感動我們.
保守孩子的口.
願主的聖靈來感動我們.
保守孩子的口.
願主的聖靈來感動我們.
保守孩子的口.
願主的聖靈來感動我們.
保守孩子的口.
願主的聖靈來感動我們.
保守孩子的口.
願主的聖靈來感動我們.
保守孩子的口.
願主的聖靈來感動我們.
保守孩子的口.
願主的聖靈來感動我們.
保守孩子的口.
願主的聖靈來感動我們.
保守孩子的口.
願主的聖靈來感動我們.

$^{81}$保守孩子的口.
願主的聖靈來感動我們.
保守孩子的口.
願主的聖靈來感動我們.
保守孩子的口.
願主的聖靈來感動我們.
保守孩子的口.
願主的聖靈來感動我們.
保守孩子的口.
願主的聖靈來感動我們.
保守孩子的口.
願主的聖靈來感動我們.
保守孩子的口.
願主的聖靈來感動我們.
保守孩子的口.
願主的聖靈來感動我們.
保守孩子的口.
願主的聖靈來感動我們.
保守孩子的口.
願主的聖靈來感動我們.
保守孩子的口.
願主的聖靈來感動我們.
保守孩子的口.
願主的聖靈來感動我們.
保守孩子的口.
願主的聖靈來感動我們.
保守孩子的口.
願主的聖靈來感動我們.
保守孩子的口.
願主的聖靈來感動我們.
保守孩子的口.
願主的聖靈來感動我們.
保守孩子的口.
願主的聖靈來感動我們.
保守孩子的口.
願主的聖靈來感動我們.
保守孩子的口.
願主的聖靈來感動我們.
保守孩子的口.
願主的聖靈來感動我們.

$^{121}$保守孩子的口.
願主的聖靈來感動我們.
保守孩子的口.
願主的聖靈來感動我們.
保守孩子的口.
願主的聖靈來感動我們.
保守孩子的口.
願主的聖靈來感動我們.
保守孩子的口.
願主的聖靈來感動我們.
保守孩子的口.
願主的聖靈來感動我們.
保守孩子的口.
願主的聖靈來感動我們.
保守孩子的口.
願主的聖靈來感動我們.
保守孩子的口.
願主的聖靈來感動我們.
保守孩子的口.
願主的聖靈來感動我們.
保守孩子的口.
願主的聖靈來感動我們.
保守孩子的口.
願主的聖靈來感動我們.
保守孩子的口.
願主的聖靈來感動我們.
保守孩子的口.
願主的聖靈來感動我們.
保守孩子的口.
願主的聖靈來感動我們.
保守孩子的口.
願主的聖靈來感動我們.
保守孩子的口.
願主的聖靈來感動我們.
保守孩子的口.
願主的聖靈來感動我們.
保守孩子的口.
願主的聖靈來感動我們.
保守孩子的口.
願主的聖靈來感動我們.

$^{161}$保守孩子的口.
願主的聖靈來感動我們.
保守孩子的口.
願主的聖靈來感動我們.
保守孩子的口.
願主的聖靈來感動我們.
保守孩子的口.
願主的聖靈來感動我們.
保守孩子的口.
願主的聖靈來感動我們.
保守孩子的口.
願主的聖靈來感動我們.
保守孩子的口.
願主的聖靈來感動我們.
保守孩子的口.
願主的聖靈來感動我們.
保守孩子的口.
願主的聖靈來感動我們.
保守孩子的口.
願主的聖靈來感動我們.
保守孩子的口.
願主的聖靈來感動我們.
保守孩子的口.
願主的聖靈來感動我們.
保守孩子的口.
願主的聖靈來感動我們.
保守孩子的口.
願主的聖靈來感動我們.
保守孩子的口.
願主的聖靈來感動我們.
保守孩子的口.
願主的聖靈來感動我們.
保守孩子的口.
願主的聖靈來感動我們.
保守孩子的口.
願主的聖靈來感動我們.
保守孩子的口.
願主的聖靈來感動我們.
保守孩子的口.
願主的聖靈來感動我們.

$^{201}$保守孩子的口.
願主的聖靈來感動我們.
保守孩子的口.
願主的聖靈來感動我們.
保守孩子的口.
願主的聖靈來感動我們.
保守孩子的口.
願主的聖靈來感動我們.
保守孩子的口.
願主的聖靈來感動我們.
保守孩子的口.
願主的聖靈來感動我們.
保守孩子的口.
願主的聖靈來感動我們.
保守孩子的口.
願主的聖靈來感動我們.
保守孩子的口.
願主的聖靈來感動我們.
保守孩子的口.
願主的聖靈來感動我們.
保守孩子的口.
願主的聖靈來感動我們.
保守孩子的口.
願主的聖靈來感動我們.
保守孩子的口.
願主的聖靈來感動我們.
保守孩子的口.
願主的聖靈來感動我們.
保守孩子的口.
願主的聖靈來感動我們.
保守孩子的口.
願主的聖靈來感動我們.
保守孩子的口.
願主的聖靈來感動我們.
保守孩子的口.
願主的聖靈來感動我們.
保守孩子的口.
願主的聖靈來感動我們.
保守孩子的口.
願主的聖靈來感動我們.

$^{241}$保守孩子的口.
願主的聖靈來感動我們.
保守孩子的口.
願主的聖靈來感動我們.
保守孩子的口.
願主的聖靈來感動我們.
保守孩子的口.
願主的聖靈來感動我們.
保守孩子的口.
願主的聖靈來感動我們.
保守孩子的口.
願主的聖靈來感動我們.
保守孩子的口.
願主的聖靈來感動我們.
保守孩子的口.
願主的聖靈來感動我們.
保守孩子的口.
願主的聖靈來感動我們.
保守孩子的口.
願主的聖靈來感動我們.
保守孩子的口.
願主的聖靈來感動我們.
保守孩子的口.
願主的聖靈來感動我們.
保守孩子的口.
願主的聖靈來感動我們.
保守孩子的口.
願主的聖靈來感動我們.
保守孩子的口.
願主的聖靈來感動我們.
保守孩子的口.
願主的聖靈來感動我們.
保守孩子的口.
願主的聖靈來感動我們.
保守孩子的口.
願主的聖靈來感動我們.
保守孩子的口.
願主的聖靈來感動我們.
保守孩子的口.
願主的聖靈來感動我們.

$^{281}$保守孩子的口.
願主的聖靈來感動我們.
保守孩子的口.
願主的聖靈來感動我們.
保守孩子的口.
願主的聖靈來感動我們.
保守孩子的口.
願主的聖靈來感動我們.
保守孩子的口.
願主的聖靈來感動我們.
保守孩子的口.
願主的聖靈來感動我們.
保守孩子的口.
願主的聖靈來感動我們.
保守孩子的口.
願主的聖靈來感動我們.
保守孩子的口.
願主的聖靈來感動我們.
保守孩子的口.
願主的聖靈來感動我們.
保守孩子的口.
願主的聖靈來感動我們.
保守孩子的口.
願主的聖靈來感動我們.
保守孩子的口.
願主的聖靈來感動我們.
保守孩子的口.
願主的聖靈來感動我們.
保守孩子的口.
願主的聖靈來感動我們.
保守孩子的口.
願主的聖靈來感動我們.
保守孩子的口.
願主的聖靈來感動我們.
保守孩子的口.
願主的聖靈來感動我們.
保守孩子的口.
願主的聖靈來感動我們.
保守孩子的口.
願主的聖靈來感動我們.

$^{321}$保守孩子的口.
願主的聖靈來感動我們.
保守孩子的口.
願主的聖靈來感動我們.
保守孩子的口.
願主的聖靈來感動我們.
保守孩子的口.
願主的聖靈來感動我們.
保守孩子的口.
願主的聖靈來感動我們.
保守孩子的口.
願主的聖靈來感動我們.
保守孩子的口.
願主的聖靈來感動我們.
保守孩子的口.
願主的聖靈來感動我們.
保守孩子的口.
願主的聖靈來感動我們.
保守孩子的口.
願主的聖靈來感動我們.
保守孩子的口.
願主的聖靈來感動我們.
保守孩子的口.
願主的聖靈來感動我們.
保守孩子的口.
願主的聖靈來感動我們.
保守孩子的口.
願主的聖靈來感動我們.
保守孩子的口.
願主的聖靈來感動我們.
保守孩子的口.
願主的聖靈來感動我們.
保守孩子的口.
願主的聖靈來感動我們.
保守孩子的口.
願主的聖靈來感動我們.
保守孩子的口.
願主的聖靈來感動我們.
保守孩子的口.
願主的聖靈來感動我們.

$^{361}$保守孩子的口.
願主的聖靈來感動我們.
保守孩子的口.
願主的聖靈來感動我們.
保守孩子的口.
願主的聖靈來感動我們.
保守孩子的口.
願主的聖靈來感動我們.
保守孩子的口.
願主的聖靈來感動我們.
保守孩子的口.
願主的聖靈來感動我們.
保守孩子的口.
願主的聖靈來感動我們.
保守孩子的口.
願主的聖靈來感動我們.
保守孩子的口.
願主的聖靈來感動我們.
保守孩子的口.
願主的聖靈來感動我們.
保守孩子的口.
願主的聖靈來感動我們.
保守孩子的口.
願主的聖靈來感動我們.
保守孩子的口.
願主的聖靈來感動我們.
保守孩子的口.
願主的聖靈來感動我們.
保守孩子的口.
願主的聖靈來感動我們.
保守孩子的口.
願主的聖靈來感動我們.
保守孩子的口.
願主的聖靈來感動我們.
保守孩子的口.
願主的聖靈來感動我們.
保守孩子的口.
願主的聖靈來感動我們.
保守孩子的口.
願主的聖靈來感動我們.

$^{401}$保守孩子的口.
願主的聖靈來感動我們.
保守孩子的口.
願主的聖靈來感動我們.
保守孩子的口.
願主的聖靈來感動我們.
保守孩子的口.
願主的聖靈來感動我們.
保守孩子的口.
願主的聖靈來感動我們.
保守孩子的口.
願主的聖靈來感動我們.
保守孩子的口.
願主的聖靈來感動我們.
保守孩子的口.
願主的聖靈來感動我們.
保守孩子的口.
願主的聖靈來感動我們.
保守孩子的口.
願主的聖靈來感動我們.
保守孩子的口.
願主的聖靈來感動我們.
保守孩子的口.
願主的聖靈來感動我們.
保守孩子的口.
願主的聖靈來感動我們.
保守孩子的口.
願主的聖靈來感動我們.
保守孩子的口.
願主的聖靈來感動我們.
保守孩子的口.
願主的聖靈來感動我們.
保守孩子的口.
願主的聖靈來感動我們.
保守孩子的口.
願主的聖靈來感動我們.
保守孩子的口.
願主的聖靈來感動我們.
保守孩子的口.
願主的聖靈來感動我們.

$^{441}$保守孩子的口.
願主的聖靈來感動我們.
保守孩子的口.
願主的聖靈來感動我們.
保守孩子的口.
願主的聖靈來感動我們.
保守孩子的口.
願主的聖靈來感動我們.
保守孩子的口.
願主的聖靈來感動我們.
保守孩子的口.
願主的聖靈來感動我們.
保守孩子的口.
願主的聖靈來感動我們.
保守孩子的口.
願主的聖靈來感動我們.
保守孩子的口.
願主的聖靈來感動我們.
保守孩子的口.
願主的聖靈來感動我們.
保守孩子的口.
願主的聖靈來感動我們.
保守孩子的口.
願主的聖靈來感動我們.
保守孩子的口.
願主的聖靈來感動我們.
保守孩子的口.
願主的聖靈來感動我們.
保守孩子的口.
願主的聖靈來感動我們.
保守孩子的口.
願主的聖靈來感動我們.
保守孩子的口.
願主的聖靈來感動我們.
保守孩子的口.
願主的聖靈來感動我們.
保守孩子的口.
願主的聖靈來感動我們.
保守孩子的口.
願主的聖靈來感動我們.

$^{481}$保守孩子的口.
願主的聖靈來感動我們.
保守孩子的口.
願主的聖靈來感動我們.
保守孩子的口.
願主的聖靈來感動我們.
保守孩子的口.
願主的聖靈來感動我們.
保守孩子的口.
願主的聖靈來感動我們.
保守孩子的口.
願主的聖靈來感動我們.
保守孩子的口.
願主的聖靈來感動我們.
保守孩子的口.
願主的聖靈來感動我們.
保守孩子的口.
願主的聖靈來感動我們.
保守孩子的口.
願主的聖靈來感動我們.
保守孩子的口.
願主的聖靈來感動我們.
保守孩子的口.
願主的聖靈來感動我們.
保守孩子的口.
願主的聖靈來感動我們.
保守孩子的口.
願主的聖靈來感動我們.
保守孩子的口.
願主的聖靈來感動我們.
保守孩子的口.
願主的聖靈來感動我們.
保守孩子的口.
願主的聖靈來感動我們.
保守孩子的口.
願主的聖靈來感動我們.
保守孩子的口.
願主的聖靈來感動我們.
保守孩子的口.
願主的聖靈來感動我們.

$^{521}$保守孩子的口.
願主的聖靈來感動我們.
保守孩子的口.
願主的聖靈來感動我們.
保守孩子的口.
願主的聖靈來感動我們.
保守孩子的口.
願主的聖靈來感動我們.
保守孩子的口.
願主的聖靈來感動我們.
保守孩子的口.
願主的聖靈來感動我們.
保守孩子的口.
願主的聖靈來感動我們.
保守孩子的口.
願主的聖靈來感動我們.
保守孩子的口.
願主的聖靈來感動我們.
保守孩子的口.
願主的聖靈來感動我們.
保守孩子的口.
願主的聖靈來感動我們.
保守孩子的口.
願主的聖靈來感動我們.
保守孩子的口.
願主的聖靈來感動我們.
保守孩子的口.
願主的聖靈來感動我們.
保守孩子的口.
願主的聖靈來感動我們.
保守孩子的口.
願主的聖靈來感動我們.
保守孩子的口.
願主的聖靈來感動我們.
保守孩子的口.
願主的聖靈來感動我們.
保守孩子的口.
願主的聖靈來感動我們.
保守孩子的口.
願主的聖靈來感動我們.

$^{561}$保守孩子的口.
願主的聖靈來感動我們.
保守孩子的口.
願主的聖靈來感動我們.
保守孩子的口.
願主的聖靈來感動我們.
保守孩子的口.
願主的聖靈來感動我們.
保守孩子的口.
願主的聖靈來感動我們.
保守孩子的口.
願主的聖靈來感動我們.
保守孩子的口.
願主的聖靈來感動我們.
保守孩子的口.
願主的聖靈來感動我們.
保守孩子的口.
願主的聖靈來感動我們.
保守孩子的口.
願主的聖靈來感動我們.
保守孩子的口.
願主的聖靈來感動我們.
保守孩子的口.
願主的聖靈來感動我們.
保守孩子的口.
願主的聖靈來感動我們.
保守孩子的口.
願主的聖靈來感動我們.
保守孩子的口.
願主的聖靈來感動我們.
保守孩子的口.
願主的聖靈來感動我們.
保守孩子的口.
願主的聖靈來感動我們.
保守孩子的口.
願主的聖靈來感動我們.
保守孩子的口.
願主的聖靈來感動我們.
保守孩子的口.
願主的聖靈來感動我們.

$^{601}$保守孩子的口.
願主的聖靈來感動我們.
保守孩子的口.
願主的聖靈來感動我們.
保守孩子的口.
願主的聖靈來感動我們.
保守孩子的口.
願主的聖靈來感動我們.
保守孩子的口.
願主的聖靈來感動我們.
保守孩子的口.
願主的聖靈來感動我們.
保守孩子的口.
願主的聖靈來感動我們.
保守孩子的口.
願主的聖靈來感動我們.
保守孩子的口.
願主的聖靈來感動我們.
保守孩子的口.
願主的聖靈來感動我們.
保守孩子的口.
願主的聖靈來感動我們.
保守孩子的口.
願主的聖靈來感動我們.
保守孩子的口.
願主的聖靈來感動我們.
保守孩子的口.
願主的聖靈來感動我們.
保守孩子的口.
願主的聖靈來感動我們.
保守孩子的口.
願主的聖靈來感動我們.
保守孩子的口.
願主的聖靈來感動我們.
保守孩子的口.
願主的聖靈來感動我們.
保守孩子的口.
願主的聖靈來感動我們.
保守孩子的口.
願主的聖靈來感動我們.

$^{641}$保守孩子的口.
願主的聖靈來感動我們.
保守孩子的口.
願主的聖靈來感動我們.
保守孩子的口.
願主的聖靈來感動我們.
保守孩子的口.
願主的聖靈來感動我們.
保守孩子的口.
願主的聖靈來感動我們.
保守孩子的口.
願主的聖靈來感動我們.
保守孩子的口.
願主的聖靈來感動我們.
保守孩子的口.
願主的聖靈來感動我們.
保守孩子的口.
願主的聖靈來感動我們.
保守孩子的口.
願主的聖靈來感動我們.
保守孩子的口.
願主的聖靈來感動我們.
保守孩子的口.
願主的聖靈來感動我們.
保守孩子的口.
願主的聖靈來感動我們.
保守孩子的口.
願主的聖靈來感動我們.
保守孩子的口.
願主的聖靈來感動我們.
保守孩子的口.
願主的聖靈來感動我們.
保守孩子的口.
願主的聖靈來感動我們.
保守孩子的口.
願主的聖靈來感動我們.
保守孩子的口.
願主的聖靈來感動我們.
保守孩子的口.
願主的聖靈來感動我們.

$^{681}$保守孩子的口.
願主的聖靈來感動我們.
保守孩子的口.
願主的聖靈來感動我們.
保守孩子的口.
願主的聖靈來感動我們.
保守孩子的口.
願主的聖靈來感動我們.
保守孩子的口.
願主的聖靈來感動我們.
保守孩子的口.
願主的聖靈來感動我們.
保守孩子的口.
願主的聖靈來感動我們.
保守孩子的口.
願主的聖靈來感動我們.
保守孩子的口.
願主的聖靈來感動我們.
保守孩子的口.
願主的聖靈來感動我們.
保守孩子的口.
願主的聖靈來感動我們.
保守孩子的口.
願主的聖靈來感動我們.
保守孩子的口.
願主的聖靈來感動我們.
保守孩子的口.
願主的聖靈來感動我們.
保守孩子的口.
願主的聖靈來感動我們.
保守孩子的口.
願主的聖靈來感動我們.
保守孩子的口.
願主的聖靈來感動我們.
保守孩子的口.
願主的聖靈來感動我們.
保守孩子的口.
願主的聖靈來感動我們.
保守孩子的口.
願主的聖靈來感動我們.

$^{721}$保守孩子的口.
願主的聖靈來感動我們.
保守孩子的口.
願主的聖靈來感動我們.
保守孩子的口.
願主的聖靈來感動我們.
保守孩子的口.
願主的聖靈來感動我們.
保守孩子的口.
願主的聖靈來感動我們.
保守孩子的口.
願主的聖靈來感動我們.
保守孩子的口.
願主的聖靈來感動我們.
保守孩子的口.
願主的聖靈來感動我們.
保守孩子的口.
願主的聖靈來感動我們.
保守孩子的口.
願主的聖靈來感動我們.
保守孩子的口.
願主的聖靈來感動我們.
保守孩子的口.
願主的聖靈來感動我們.
還是我們常常會像一個.
很像小孩子一樣.
爸爸抱我.
爸爸幫助我.
爸爸我在危難的時候.
我該怎麼辦.
我有很多次在我人生中.
遇到一些困難的時候.
我就是跪下來禱告.
主啊 怎麼辦.
主啊 幫助我.
甚至我說有一次危及到.
如果這個災難不解決的話.
會不會我們公司所有的data.
全部 特別是研發的.
所有的圖紙.

$^{761}$所有的機房裡面的重要資料.
都沒有辦法回復.
這麼嚴重.
甚至股價可能就會往下跌.
在這種光景之下.
我只能求神.
主啊 主啊 你來拯救.
禱告 禱告 禱告著.
硬碟可以恢復正常.
雖然有備份.
但是當下還是非常的緊急.
這件事情可能只有我知道.
還有我的神知道.
祂是我患難中隨時的幫助.
我們都認識的戴德森.
他曾經在一個光景之下.
1969 1869年.
他寫信給他的親愛的妹妹.
他跟她說.
我想跟你分享一件事情.
是我最近的一個喜樂.
他的一個得著.
但是在此之前.
戴德森他揭露了.
他靈魂中的黑暗.
他說 親愛的妹妹 親愛的.
過去六到八個月裡.
我一直在思索萬千.
深感自己和我現在的使命.
都需要在我的靈魂中.
擁有更多的聖潔 生命和力量.
戴德森他說.
我的工作越來越忙了.
真的是太多事工了 太多了.
他感受到他需要更多的聖潔.
更多從神來的生命和力量.
於是他感受到.
有另外一股自私的力量.
在他生命中拉扯.
他說 我感受到.

$^{801}$我個人的需要始終是首要的.
也是最大的.
我覺得自己是一個忘恩負義.
我感受到這是一種危險.
我覺得我自己沒有更親近神.
我反而犯了更多的罪.
他說 綠色的字是.
綠色的字是他很努力的禱告.
他說 我禱告.
我進食 我苦思.
我努力 我下定決心.
我更勤奮的研讀聖經.
尋求更多的時間 獨處和夢想.
但這一切都怎樣呢.
無濟於事.
弟兄姊妹.
請問你要怎麼幫助他.
在座的小組長.
在座的牧者同工們.
在座的長老.
我們怎麼樣幫助.
這一個已經很努力禱告.
很努力進食.
很努力下定決心.
很努力讀聖經.
很努力花時間與神獨處.
他發現這一切都無濟於事.
他有沒有顛覆我們的信仰神學呢.
有沒有顛覆我們的.
屬靈知識和經驗呢.
他說 幾乎每天每時刻.
罪的意識都壓得我喘不過氣來.
長輩的神他已經心裡非常清楚.
知道神的話.
什麼是對的 什麼是錯的.
然後他看見自己仍然是一個敗壞的人.
他想盡用他自己的神學.
用他自己過去的裝備訓練.
他也沒有辦法解決.
他裡面為惡的動力.

$^{841}$還是自私的 還是會犯罪.
你看他犯了什麼罪.
他說 我知道.
只要我能藏在記錄裡.
一切都會好起來.
這是他的神學 沒有錯.
但是我做不到.
兄姐妹.
你會不會因此而感覺到安慰呢.
你會不會覺得.
這樣的一個宣教士.
我都沒有辦法像他一樣.
有這樣的熱誠去宣教.
甚至我都沒有辦法像他一樣的.
有這樣的屬靈的操練.
他都有這樣的信心.
甚至他都覺得.
只要我能藏在記錄裡.
一切就會好起來.
可是他承認.
我做不到.
兄姐妹.
我們今天是不是也要在主的面前承認.
主啊.
我做不到.
主啊 有太多我知道.
我做不到的.
我們有沒有很多時候.
我們指責人說.
你應該要悔改.
你應該要這樣.
其實有很多時候.
如果我們好好自省的話.
有很多我們離神的道甚遠.
是嗎.
我們都還在背負著.
很多的罪.
很多的重擔.
然而當這些都沒解決的時候.
我們就會偽裝成.

$^{881}$我們都是聖潔的.
我們都是得勝的.
我們都是喜樂的神的兒女.
可是好好檢視我們內在的時候.
我們裡面是破碎的.
戴德森說.
我每天都以禱告開始.
決心片刻不離他.
他下定決心了.
但繁重的職責.
他非常多的事工.
有時令人精疲力盡.
不斷的干擾.
也常常使我忘記神.
在這種環境下.
人的神經會變得怎樣呢.
異常的緊張.
更容易產生意怒.
負面情緒.
甚至說出刻薄的話語.
這些都更難控制.
戴德森他發現.
他這個人是精疲力盡的.
他被許多的事物干擾.
讓他忘記神.
甚至他裡面有一個負面的情緒.
常常說不好聽的話.
他說我每一天都充滿.
罪惡和失敗以及無力感.
你看一個神的僕人.
他還可以知道.
我每一天都充滿罪惡和失敗.
有時候我們軟弱到一種程度.
是我們根本沒有意識到.
我們每一天都在犯罪.
不是嗎.
他說我的確有決心.
卻不知如何付諸行動.
他說我恨自己.
我恨我的罪.

$^{921}$我沒有能力去戰勝他.
兄姐妹,怎麼辦呢.
如果人若單靠自己的能力.
我們就會錯過上帝的能力.
我們就會錯過.
聖靈要賜給我們能力.
大衛他不僅要神做他的磐石.
大衛他回想起來.
神曾經做過他的避難所.
神曾經做過他的兼顧台脫離仇敵.
所以他決定.
他的心不是只是要有一個磐石在那裡.
好像有一個信仰的力量在那裡.
好像戴德森知道上帝在那裡.
上帝一定能夠幫助我.
但是裡面沒有感受到上帝的幫助.
弟兄姊妹.
我們的信仰生活是不是這樣呢.
你會不會覺得.
有時候我們要傳講的.
我們要傳講的耶穌基督.
講著講著我們發現.
我好像不是過這種得勝的生活.
是吧.
所以當我們每次去傳福音的時候.
你會不會覺得.
我跟人家說信耶穌有平安.
但是我裡面有平安嗎.
心虛吧有時候.
我想戴德森他就是這種處境.
所以他在這張他給他妹妹的書信裡面就寫到.
他說我真的是需要救贖.
但是我怎麼跟世人講.
有一位救贖主在我的裡面.
他幫助我得勝.
戴德森他都不覺得他得勝了.
他覺得他傳講這個福音很虛偽.
弟兄姊妹.
我們人生我們的信仰.
會不會是不是真實的情況就是如此呢.

$^{961}$大衛讓我們看見他不是只是要神做他的盤石.
他說我要投靠在那翅膀的隱秘處.
翅膀的隱秘處在舊約我們知道.
在這個藏墓裡面的約櫃.
有那個基督伯的天使.
他的有翅膀.
而翅膀下面就是神同在的.
施恩座的就是代表神同在的地方.
大衛他想要投靠的是.
那個藏墓裡面那個隱秘處.
那個隱秘處是神的同在.
但是你想想看.
在瑪哈涅那個地方.
有神的約櫃嗎.
沒有.
當時在逃難的時候.
本來這個撒度.
祭司他們抬著這個約櫃.
要跟大衛一起逃難的.
大衛說你是先知.
你留在那裡.
沒你的事.
所以約櫃就留在耶路撒冷.
沒跟著大衛走.
大衛他想要起來投靠神.
在詩篇61篇.
我們看見兩個上帝的寶藏.
一個是磐石.
一個是翅膀.
磐石好像是給我們客觀的寶藏.
好像是外面的拯救.
好像代表神穩固的根基.
那一種能力.
但是翅膀.
永遠沒你想到翅膀.
翅膀想到什麼呢.
不是烤雞翅膀.
是耶穌說.
我曾經想要.
耶路撒冷啊.

$^{1001}$耶路撒冷.
我曾經像母雞.
想要聚集小雞一樣.
把你們聚集起來.
但是你們怎麼樣呢.
你們不願意啊.
神想要用祂的翅膀來遮蓋我們.
所以約櫃的那個基督伯.
他是有翅膀的.
以賽亞他也看見.
神坐在他的寶座上的時候.
那個翅膀是垂下來的.
很大的翅膀.
它可以從天上垂到地上.
弟兄姊妹.
那是多大的翅膀.
可以把整個地球給遮蓋.
那個翅膀是大的.
那翅膀.
大衛在這個翅膀裡面.
他找到了安慰.
他找到了醫治.
他找到了神的同在.
我們看詩篇裡面.
大衛說我要投靠在你翅膀的印象.
等到災害過去.
他在那裡找到了神的平安.
神成為他的遮蓋.
大衛也找到了親密.
大衛說我要在你的翅膀印象歡呼.
為什麼呢.
因為上文是這樣寫的.
因為你曾幫助我.
大衛知道誰能夠幫助他.
是那一位有翅膀的神.
能夠遮蓋他.
大衛可以在神的翅膀下怎麼樣呢.
喜樂.
所以你看大衛在發昏的時候.
他想到的是.

$^{1041}$我可以在神的翅膀的印象.
我可以在那裡歡呼喜樂.
你在詩篇63篇.
再翻兩頁就看見了.
大衛可以在那裡歡呼.
可以在那裡喜樂.
另外一個就是.
大衛可以在那裡.
像我現在一樣.
可以得到樂和的水.
大衛口渴的時候.
他可想的是神.
不是水.
但是大衛他說.
我在神的裡面.
可以得到樂和的水.
兄弟姐妹.
戴德森在他人生的光景.
最低落的時候.
他那時候其實也陷入到.
一種長期的憂鬱裡面.
他的世公.
然後他同工背叛他.
當時有一些揚州的教案.
發生這件事情.
他妻子病重.
甚至他的兒子.
他的孩子相繼的死亡.
他就陷入到這種長期的憂鬱.
還有靈性的掙扎裡面.
在這種光景裡面.
他的信仰沒有辦法幫助他.
就在這個危急的時候.
有一封信改變了他.
這也是來自於一個.
同樣做宣教士的麥卡錫.
寫信給他.
他告訴他說.
真正的讓你的信心增長.
不是靠你努力自己去追求的.

$^{1081}$而是安息在那位信使主的身上.
也就是說.
如果我們要成聖.
不是靠著自己的意志去奮鬥.
而是純粹是安息在基督獻經的能力裡面.
所以戴德森因為他的這樣的一個信息.
讓他體會到.
我是需要安息在基督的裡面.
我需要.
我真的是要把我的重擔交給他.
讓他來帶領我成聖.
而不是我用我的意志.
用我的心裡面想盡的辦法.
要成為更好一點的人.
弟兄姊妹.
你用你的意志是做不到的.
我曾經在信仰的生命中裡面.
當我認識我的主.
是我生命中唯一的主的時候.
我也做了一個勵志.
當時我勵志什麼呢.
我每天都要讀聖經.
弟兄姊妹.
你有沒有做過這樣的勵志.
你有沒有讀聖經.
有嗎.
我當時我真的是讀聖經.
我真的是天天讀聖經.
晚上如果我發現我今天沒有讀聖經.
我還要爬起來讀完再睡覺.
後來我發現.
不行.
我一定要找出我生命中最容易讀聖經的時候.
那就是吃午飯之後.
吃午飯之後.
因為我那時候在重考.
在讀書.
非常壓力大.
沒有辦法很routine的有一個很規律.
於是我就是.

$^{1121}$我每天一定要吃飯吧.
我就是吃完中飯的時候.
我就開始讀聖經.
就開始拿出注釋書.
把重點就抄在這個聖經上面.
我靠我自己的意志.
我每天要讀聖經.
果然還是可以做到.
但是我就發現裡面的罪.
裡面的慾望.
怎麼樣呢.
沒辦法消除.
我還記得當時上癮的一個八點檔.
我每天回到家的時候.
就心裡面在想那個八點檔.
我要不要看.
因為家人在看.
因為太好看了.
看了兩集之後發現很好看.
所以我每次都在那裡掙扎.
我要不要看呢.
主啊 幫助我.
掙扎 弟兄姊妹.
掙扎 掙扎 掙扎.
然後有一個.
坐在我重考教室的.
有一個隔壁排的女生.
她長得很嬌小.
非常的可愛.
每次都過不去的時候.
都想多看她兩眼.
但是我的神就告訴我說.
你不可以戀慕偶像.
你知道那種掙扎嗎.
弟兄姊妹.
就是當我們知道.
神的話語是這樣.
準神是這樣.
哎呀 我裡面的慾望.
我裡面的老我是敗壞的.

$^{1161}$怎麼樣辦法勝過呢.
我靠我自己的意志.
我可以.
但是我會發現我好苦啊.
我好像過著一種修行的生活.
我好像發現.
一旦有這個.
罪惡的念頭來的時候.
我馬上要認罪悔改.
主啊 赦免我.
我會在這種緊張壓力當中.
就像戴德森所寫的.
這很真實.
戴德森他發現.
他需要的是安息.
他需要的是那位主.
真的是帶領他的主.
他說我已經與基督一同死.
一同埋葬了.
我要相信.
我現在活著就不再是我.
乃是基督在我裡面活著.
我現在相信.
我已經向罪死了.
他說我仍然像以前一樣會犯罪.
但基督的靈在比以前.
以往任何時候都更真實.
弟兄姊妹.
他的磐石活化在他的面前.
真理的聖靈.
藉著他的話.
真的活在他的面前.
這位磐石.
真的是在他生命中軟弱的時候.
成為他穩固的根基了.
他發現他不需要再說謊.
他不需要再犯罪.
他也不覺得他裡面有犯罪的動力.
他說我的良心變得更加的什麼呢.
溫柔.

$^{1201}$罪惡立刻被看見.
被承認.
被赦免.
平安和喜樂.
和謙卑.
立刻被恢復.
弟兄姊妹.
我們都知道在基督裡的重要.
我想葛牧師常常的講這個.
我們要在基督裡.
我們在城島.
我們常常也說.
我們在基督裡.
弟兄姊妹.
但是我們很少停留在基督裡.
我們就像戴德森一樣.
很多的事工.
很多的忙碌.
我們沒有停留在主的同在裡.
或者是.
我們沒有讓那個磐石.
常常在我們的面前.
所以大衛說.
我常把你擺在我面前.
你在我右邊.
我就不知怎麼呢.
不知動搖了.
他有時候還會說.
我不知很動搖.
顯然他還是會有動搖的時候.
弟兄姊妹.
我們生命中會有軟弱的時候.
但是神到底在你生命中.
關鍵的時刻.
祂是不是你的救主呢.
祂在你跟你的另外一半.
關係好像到一個冰點的時候.
你甚至要想要離婚的時候.
有那個念頭出來的時候.
是你自己的意願.

$^{1241}$在帶領你走人生.
還是上帝的法則.
在帶領你走人生的道路呢.
恐怕很多時候.
我們還是按照我們自己.
任憑我們的腳步.
任意而為.
所以我們就經歷不到.
在困難的時候.
神能夠使這個苦情翻轉.
戴德森說.
不管壓力有多大.
真正的壓力在哪裡呢.
是擋在你與上帝的中間.
還是祂將你推向.
更親近神的中心呢.
我們都會遇到苦難.
都會遇到在地急昏頭的時候.
但是到底這樣的一個處境.
這樣的一個心昏.
發昏的時候.
祂是擋在我跟上帝中間呢.
還是祂把我推向更靠近神呢.
最後我們看見.
大衛他從一個逃亡者.
他轉變成為一個永恆的敬拜者.
兄弟姐妹.
大衛他可以用他的眼光看見.
他失去的是王位.
他失去的是王的尊嚴.
他失去的是.
以色列百姓對他的尊崇.
他失去的是敬拜神的權利.
他失去的是他幸福的家園.
失去了他的財富.
甚至失去了他的夢想.
大衛失去了這一切.
大衛失去了.
大衛有沒有可能會失去.
他為神間諜的願望呢.

$^{1281}$很有可能對不對.
因為他在已經這種光景之下.
還有可能回耶路撒冷嗎.
以色列人的心都歸向他兒子亞撒隆了.
大衛還有第二次機會嗎.
如果照世人的標準.
大衛的眼光看到的應該是失望.
應該是要埋怨神.
為什麼.
你不是跟我說好的嗎.
當我歲數滿足的時候.
我的兒子要接續我做王.
而且他的王位是到永遠.
兄弟姐妹.
大衛沒有失望.
我們看見大衛他是起來.
他說.
神啊.
你越是聽了我許的願.
你將產業賜給敬畏你的人.
大衛在那裡.
他回想到.
他所禱告的內容不就是大衛之約嗎.
就是神跟他所立的約啊.
你要加添王的壽數.
他的年歲必存到世世.
他必永遠坐在神面前.
願你預備慈愛和誠實保佑他.
大衛心裡面應該很清楚知道.
他兒子的王位沒有所謂的永遠的.
是吧.
誰知道.
誰會認為你的兒子的性命是活到永遠的.
沒有吧.
但是神有給他屬靈的眼光和看見.
雖然這個啟示在當時並不是很完全.
但是大衛相信.
上帝要興起的那一位王.
就是我們的主.
要稱為主.

$^{1321}$然後他的國位呢.
是永永遠遠直到永遠的.
他要坐王到永遠.
這一位王是誰.
耶穌基督.
大衛在很久之前.
他就有這個信心的眼光.
他就來讚美神.
神你所說的是對的.
雖然有永遠這個字.
人世間哪有永遠呢.
神說永遠.
神你說永遠就是永遠.
神跟他說的是.
你的家和你的國必在我面前永遠堅定.
你的國位也必堅定直到永遠.
這裡有家永遠.
國永遠.
國位是永遠的.
耶穌基督才是真正永遠的君王.
阿們.
大衛的信心今天傳承到哪裡.
大衛的信心今天一直傳承.
一直跟我們有很深的連結.
如果大衛在那時候他就不信了.
他就自我了結了.
今天也沒有彌賽亞.
今天也沒有耶穌基督.
大衛的信心今天仍然持續.
好像傳承到我們今天一樣.
他的信心是歷史的證明.
神保守.
神自己保守他要成就的.
兄姐妹.
我們的信心不是建立在我們所眼見的.
是不憑眼見的.
而應該建立在神的應許當中.
阿們.
我們是不是神的兒女.
我們是不是神所寶貝的.

$^{1361}$我們是不是神會在我們的裡面供應我們.
是不是有火水江河在我們裡面呢.
我們是不是要遠離一切的罪惡呢.
神能不能夠聖靈可不可以幫助我們.
制止我們身體的惡行呢.
可以.
兄姐妹.
我們要把神的應許當應許.
你就會活得有信心.
阿們.
耶穌基督是我們的王.
他才是那個磐石.
他才是那個翅膀.
在歷史當中我們就看見.
上帝聽了大衛的禱告.
他讓這個互塞來迎接他的時候.
大衛就說你去.
你去到那個亞沙隆那裡去破壞他.
所以亞沙隆就接受了互塞的這樣的一個建議.
本來亞西多佛當時他的策略是對的.
聖經中說.
這是因為葉赫華定義破壞亞西多佛的良謀.
看到沒有.
亞西多佛的良謀是最有效的.
當時亞西多佛就建議亞沙隆說.
給我一萬兩千的部隊.
我當天晚上我就可以拿大衛的命.
其他的人就可以平安的回.
都歸向你.
只要殺大衛一個人就好了.
這是其實是對的.
可是上帝卻使用互塞.
破壞亞西多佛的計謀.
然後亞西多佛認為.
這個互塞你的建議更好.
我照著你所說的.
我可以殲滅所有大衛.
還有大衛的勇士.
所有都可以殲滅.
所以神定義要破壞這個惡魔.

$^{1401}$這個計謀.
所以神他出手了.
神真的是出手了.
大衛當時有看見嗎.
大衛還沒看見.
可是他就在信心裡面敬拜神.
他就在他的信心裡面宣告.
這樣我要歌頌你的名.
直到永遠好天天還我所許的願.
兄姐妹,聖經告訴我們.
我們既領受了永不振動的國.
我們就應當感恩.
照著神所喜悅的.
用虔誠敬畏的心來事奉神.
敬拜也就是事奉.
兄姐妹,你得著這永恆的國度了嗎.
在永恆的國度裡面.
你會不會遇到風浪呢.
誰才是你所事奉的神呢.
當你遇到困難的時候.
你是向天國的國王懇求.
還是向他抱怨呢.
你願意讓他成為你的遮蓋嗎.
我們請敬拜團隊預備.
在戴德森最難的時候.
他在最後的信裡面說到的結論是什麼.
他跟他妹妹說.
唯有基督才能使我們脫離罪惡.
才能使我們真正的服侍他.
我們一起來再讀一次.
唯有基督才能使我們脫離罪惡.
才能使我們真正的服侍他.
阿們.
\newpage



\section{詩篇 73:1-28}
\label{sec:032231}
\textbf{在聖所中明白人生}~~~~張健庭牧師博士
\newline
\newline
連結: \href{http://www.church.com.hk/acms/content.asp?site=cdc\&op=show\&type=product\&code=032231\&layout=sermon}{\texttt{032231}} ~~~~ 語音日期: 2026-05-31
\newline
\newline
\hyperref[sec:032230]{\small{< < < PREV SERMON < < <}}
~
\hyperref[sec:index_chronic]{\small{[返順時目]}}
~
\hyperref[sec:032232]{\small{> > > NEXT SERMON > > >}}
\newline
\newline
詩篇 73:1-28
\newline
\begin{longtable}{cl}
\hline
\hline
章節 & 經文 (和合本修訂版)\\
\hline
73:1 & \begin{tabularx}{0.7\textwidth}{X} 神實在恩待以色列那些清心的人！ \end{tabularx} \\ \\ \relax
73:2 & \begin{tabularx}{0.7\textwidth}{X} 至於我，我的腳幾乎失閃，我的步伐險些走偏； \end{tabularx} \\ \\ \relax
73:3 & \begin{tabularx}{0.7\textwidth}{X} 因為我嫉妒狂傲的人，我看見惡人享平安。 \end{tabularx} \\ \\ \relax
73:4 & \begin{tabularx}{0.7\textwidth}{X} 他們的力氣強壯，他們死的時候也沒有疼痛。 \end{tabularx} \\ \\ \relax
73:5 & \begin{tabularx}{0.7\textwidth}{X} 他們不像別人受苦，也不像別人遭災。 \end{tabularx} \\ \\ \relax
73:6 & \begin{tabularx}{0.7\textwidth}{X} 所以，驕傲如鏈子戴在他們項上，殘暴像衣裳覆蓋在他們身上。 \end{tabularx} \\ \\ \relax
73:7 & \begin{tabularx}{0.7\textwidth}{X} 他們的眼睛因體胖而凸出，他們的內心放任不羈。 \end{tabularx} \\ \\ \relax
73:8 & \begin{tabularx}{0.7\textwidth}{X} 他們譏笑人，憑惡意說欺壓人的話。他們說話自高； \end{tabularx} \\ \\ \relax
73:9 & \begin{tabularx}{0.7\textwidth}{X} 他們的口褻瀆上天，他們的舌毀謗全地。 \end{tabularx} \\ \\ \relax
73:10 & \begin{tabularx}{0.7\textwidth}{X} 所以他的百姓歸到這裡，享受滿杯的水。 \end{tabularx} \\ \\ \relax
73:11 & \begin{tabularx}{0.7\textwidth}{X} 他們說：「神怎能曉得？至高者哪會知道呢？」 \end{tabularx} \\ \\ \relax
73:12 & \begin{tabularx}{0.7\textwidth}{X} 看哪，這就是惡人，他們常享安逸，財寶增多。 \end{tabularx} \\ \\ \relax
73:13 & \begin{tabularx}{0.7\textwidth}{X} 我實在徒然潔淨了我的心，徒然洗手表明我的無辜， \end{tabularx} \\ \\ \relax
73:14 & \begin{tabularx}{0.7\textwidth}{X} 因為我終日遭災難，每日早晨受懲治。 \end{tabularx} \\ \\ \relax
73:15 & \begin{tabularx}{0.7\textwidth}{X} 我若說「我要這樣講」，就是愧對這世代的眾兒女了。 \end{tabularx} \\ \\ \relax
73:16 & \begin{tabularx}{0.7\textwidth}{X} 我思索要明白這事，眼看實係為難， \end{tabularx} \\ \\ \relax
73:17 & \begin{tabularx}{0.7\textwidth}{X} 直到我進了神的聖所，思想他們的結局。 \end{tabularx} \\ \\ \relax
73:18 & \begin{tabularx}{0.7\textwidth}{X} 你實在把他們安放在滑地，使他們跌倒滅亡； \end{tabularx} \\ \\ \relax
73:19 & \begin{tabularx}{0.7\textwidth}{X} 他們轉眼之間成了何等荒涼！他們被驚恐滅盡了。 \end{tabularx} \\ \\ \relax
73:20 & \begin{tabularx}{0.7\textwidth}{X} 人睡醒了，怎樣看夢，主啊，你醒了也必照樣輕看他們的影像。 \end{tabularx} \\ \\ \relax
73:21 & \begin{tabularx}{0.7\textwidth}{X} 因此，我心裡苦惱，肺腑被刺。 \end{tabularx} \\ \\ \relax
73:22 & \begin{tabularx}{0.7\textwidth}{X} 我這樣愚昧無知，在你面前如同畜牲。 \end{tabularx} \\ \\ \relax
73:23 & \begin{tabularx}{0.7\textwidth}{X} 然而，我常與你同在；你攙扶我的右手。 \end{tabularx} \\ \\ \relax
73:24 & \begin{tabularx}{0.7\textwidth}{X} 你要以你的訓言引導我，以後你必接我到榮耀裡。 \end{tabularx} \\ \\ \relax
73:25 & \begin{tabularx}{0.7\textwidth}{X} 除你以外，在天上我有誰呢？除你以外，在地上我也沒有所愛慕的。 \end{tabularx} \\ \\ \relax
73:26 & \begin{tabularx}{0.7\textwidth}{X} 我的肉體和我的心腸衰殘；但神是我心裡的力量，又是我的福分，直到永遠。 \end{tabularx} \\ \\ \relax
73:27 & \begin{tabularx}{0.7\textwidth}{X} 看哪，遠離你的，必要死亡；凡離棄你行淫的，你都滅絕了。 \end{tabularx} \\ \\ \relax
73:28 & \begin{tabularx}{0.7\textwidth}{X} 但我親近神是於我有益；我以主耶和華為我的避難所，好叫我述說你一切的作為。 \end{tabularx} \\ \\
[1ex]
\hline
\hline
\end{longtable}
$^{1}$弟兄姊妹平安.
我們再一次低頭祈禱.
天父願你的話語再一次的教導我們.
讓我們聽到你的聲音.
願你自己讓我們好像私人的經歷一樣.
我們進入致勝所當中來思想我們的人生.
奉耶穌基督的名祈求 阿們.
今天是我們三十天一起讀靈修書的時候最後的一講.
所以就選擇了詩篇第七十三篇來思想.
詩篇第七十三篇的時候.
一開始的時候作者就說.
神實在恩待以色列那些清心的人.
至於我 我的腳幾乎失險.
我的腳險些滑跌.
一開始的時候作者就說到.
他幾乎在信仰上跌倒的一個經歷.
我想這個也是挺值得有意思的來反省.
當然這個寫詩篇的人如果你看那個標題.
稱為阿薩.
他應該是當時帶領以色列人敬拜的.
如果在教會來說可能是一個師班長.
他應該是一個祭司.
帶領一個祭司去唱詩.
在神的殿當中敬拜的一個領袖.
你可以想像一個師班長的靈命.
都應該有相當程度的造詣.
如果不是他怎麼能夠帶領人一起去敬拜神呢.
所以他應該有相當程度的認識神.
跟神有一個相當親密的關係.
這個也是一個很合理的假設.
事實上他應該是一個這樣的人.
但是即使一個這樣的.
如果以我們今天來說.
即使是一個這樣的屬靈人.
他都面對著一個.
他都有遭遇過就是說.
他的信仰都好像落在一個很動搖的一個境地.
我們就看看究竟發生什麼事呢.
為什麼一個師班長或者一個屬靈敬拜的領袖.
他的信仰.

$^{41}$他反省他的信仰都會有一個這麼大的危機呢.
他就這麼說.
我見惡人和狂傲人在平安就心懷不平.
他們死的時候沒有疼痛.
他們的力氣卻也壯實.
他們不像別人受苦.
也不像別人遭災.
所以驕傲如鏈子.
打在他們的汗上.
強暴像衣裳遮住他們的身體.
他們的眼睛因體叛而突出.
他們所得的過於心裡所想的.
這裡就解釋了究竟什麼原因.
令到他的信仰有很大的動搖.
他的信仰很大的動搖.
其中他就發覺到一個.
令到他心裡很不平安的事情.
就是說為什麼這個世界會這樣的呢.
他每天在聖殿那裡帶領這些師班員.
一起來去敬拜神.
講道神的幾樣的慈愛.
幾樣的公義.
但是好像一離開了聖殿之後.
如果今天用我們來說.
一離開了教會之後.
為什麼這個世界好像.
完全跟他唱詩的內容不同呢.
那些惡人又說很好.
生活又順利.
又不見惡人.
說個個都是周身疾病.
而且他還講道.
他們的驕傲如鏈子.
打在他們的汗上.
他好像帶著一條頸鏈.
不是黑社會那種.
帶著金鏈那種.
但是理論上就好像那種.
不斷的show off表現.
很驕傲 很沙塵.

$^{81}$看不起那些敬虔的人.
信仰的人.
而且他說.
他們強暴如衣裳.
遮蓋他們的身體.
而且你真的完全可以想像.
又驕傲又帶著一條鏈.
又強暴.
真的像那些黑幫大佬.
帶著一條金鏈走出來.
小朋友 幹什麼.
在吵嗎.
打你兩拳 信不信.
就是這樣.
他說現在的世界是這樣的.
所以他說.
眼睛因體半而突出.
所得的 果於心裡所想的.
你說肥就不好.
但是那時候不是這樣的.
那時候怎樣.
吃都沒得吃.
有多少肥人.
沒有肥人的.
你看那些古代的片.
基本上肥的.
就是那些貪官污吏.
是不是.
真的.
因為那時候吃都沒得吃.
只有肥的.
那些肥的人一定是.
很豐富 豐富到.
其他的人是羨慕到.
羨慕到心裡很想的.
所以他說.
所得的 果於心裡所想的.
就是說.
他多到.
很多人想都得不到.

$^{121}$他說他多到.
不用想都更多.
所以你就明白.
他的作者阿薩.
他是非常之不安.
不單是不明白.
其實是不安.
他心裡感覺到很糾結.
為什麼信仰講那套東西.
他一邊在聖殿裡敬拜.
大人唱詩.
但是一走出外面.
整個世界根本就不是這樣.
他心裡就很糾結.
接著他再說.
他們譏笑人.
憑惡意說欺壓人的話.
他們說話自高.
他們的口舌逐上天.
他們的舌毀破全地.
所以神的文歸到那裡.
可盡你滿杯的苦水.
他們說神怎能曉得.
自高者豈有知識呢.
這裡其實是很諷刺.
他說那些人.
不單是做人就惡霸.
又可以得到很多他們想要的東西.
因為你做那些欺壓人的人.
反而得到很多他們想要的東西.
而且不單是他們的道德生活.
他信仰那裡是肆無忌憚的.
他自己不信神就好了.
他還說不如有神吧.
神怎會知道呢.
如果有神他也不知道.
神怎會看這麼多東西呢.
所以他們不斷地去毀謗神.
但是毀謗神沒有後果.
他看到的不是毀謗神.

$^{161}$馬上神懲罰他.
下次不敢這樣說.
拜偶像的人就會.
拜偶像的人他們真的很害怕.
譬如拜完偶像之後.
如果說了一些得罪神明的東西.
他們就很害怕會受罰.
但是這些不信神的人反而不會.
而且他不單是自己去毀謗神沒有後果.
那些信耶穌的人在那裡就很慘.
所以神的民歸到那裡.
歸到這裡喝盡滿杯的苦水.
那怎麼辦呢.
他們又好像各樣東西都很順利.
很興盛.
萬事如意.
還財就手.
還財不是好的財.
不是正當手段得的財.
他真的很就手.
而且他周圍詆毀神.
去踩基督徒.
很得意.
所以作者心裡真的很難過.
不單是難過.
其實他對於整個信仰的基礎.
都好像面對一個很大的挑戰.
究竟怎麼樣.
究竟發生什麼事.
究竟是不是真的這樣的神.
你是不是真的在看著這些.
所以他感覺到很受委屈.
他說.
看那也就是惡人.
他們既是常享安逸財寶便加珍.
我實在徒然潔淨了我的心.
徒然洗手表明無辜.
因為我終日遭災.
每早晨受懲治.
所以他感覺到很難受.

$^{201}$他說.
這些就是惡人.
你看.
我們經常.
聖經有句話.
惡人必不得平安.
但這裡怎麼說.
我們常享安逸.
你說平不平安.
我又不知道.
但我表面來看.
他真的很順風順水.
接著他說財寶便加珍.
可能你還在優柴優米的時候.
他可能還財.
就不知道怎樣.
用什麼手段.
可能很順利地得到.
事實上.
真的.
那些可能做正當生意的人.
永遠都沒有.
做偏門生意的人.
那麼發達.
永遠都是.
是不是.
那麼.
他說.
他就覺得.
如果.
我對比自己的時候.
他就覺得.
好像很吃虧.
跟從神.
好像很吃虧.
我實在徒然.
潔淨了我的心.
徒然洗手.
表明無辜.
是不是白做.

$^{241}$是不是.
就是.
你完全可以代入到他的心情.
他也很真誠地去表白.
究竟.
為什麼他心裡這麼糾結.
所以他一開始說.
我的腳幾乎滑倒.
我好像去了滑地.
信仰的基礎.
現在.
好像.
沒什麼.
找不到力.
信仰的基礎.
好像現在被人挖空了.
因為.
我終日遭災.
我每天還要遇到.
可能受到不好的人的欺凌.
如果你.
舉例來說.
你做一個.
好的人.
或者做一個.
善良的人.
你每天做生意.
被黑社會來.
又搞亂你.
又收你錢.
又什麼.
我真的終日遭災.
每天都好像受罪.
每天都好像受懲罰.
那.
所以跟從神.
真的在.
從這樣的角度來說.
是真的好像.
很徒然.

$^{281}$徒然的意思就是說.
做了些事得局吉.
究竟有什麼.
收穫呢.
跟從神.
好像是很.
根本.
不切實際.
徒然洗手表明無辜.
徒然潔淨了我的心.
我潔淨了我的心.
但是那個邪惡的人.
每天都.
張牙舞爪的.
都很好.
我何須要這樣呢.
是不是.
於是他就說.
我若說.
我要這樣說.
這就是以奸詐代理的眾子.
我思索怎樣明白這事.
眼看實際為難.
直等到進入.
神的聖所.
思想他的結局.
當.
詩人這樣說的時候.
其實.
後面就開始有一個轉接點了.
後面就開始.
在詩篇的下半段.
是很不同的.
但是前半段這裡.
你看到他講了十四節.
有.
就是頭那兩節.
不是隱然的.
基本上都講了.
十三節.

$^{321}$是講那些.
很繪形繪聲的.
那些心裡面的掙扎.
所以在這裡.
首先.
當還沒講到他有什麼出路之前.
我就想大家停下來思想.
你想想.
這是一個什麼人.
一開始已經講過了.
他是一個師班長.
帶領人來到.
每天在聖殿裡敬拜的.
首先這樣說.
如果你單單看到第十四節這裡.
他講這些話.
合不合意呢?.
有沒有想過?.
合不合意?.
喂.
你帶人敬拜.
你心裡面有這些.
你都不要講出來.
人家是.
人家聽你這樣講.
還要相信耶穌.
是不是?.
你說.
哇.
這樣搞.
你在這裡反宣傳.
倒霉.
你做師班長.
你起碼也講一些屬靈的東西.
是不是?.
你說.
哇.
我看到惡人怎樣怎樣.
但是.
他都不會說沒有後果的.

$^{361}$他咒罵神.
神遲早.
被他一家人都受害.
起碼你這樣講.
心理都沒有那麼不平衡.
是不是?.
人家這樣聽你講.
下次不來了.
還用來敬拜嗎?.
事實上就是要面對這樣的處境.
首先.
你如果這樣停下來.
你有個看法.
你會發覺其實很不合適的.
可以這樣講這些東西.
是不是?.
在這樣的角度來說.
其實很不合適.
但是停下來.
又問問大家.
剛才講的那些東西.
你有時候有沒有都這樣覺得呢?.
都有的.
是不是?.
那就是.
究竟什麼回事呢?.
首先.
我覺得這裡.
這裡首先就是.
很真誠的講出了一個事實.
就是.
你要.
去真實的.
經歷到信仰的力量本身.
你是先要面對.
自己裡面真實的掙扎.
你不承認自己有什麼掙扎.
可能你剛才他講的那些東西.
你都有想的.
我旁邊的同事.

$^{401}$天天在這裡沉迷上班的時間.
又看股票.
接著就炒到盤滿缽滿.
我上班的時候.
做事當然沒有看股票.
接著說這次股災虧到眼吞吞.
接著你又覺得不公平.
是不是?.
就是.
這些天天都看到的.
只不過我們都不.
回來教會不要講這些.
為什麼?.
因為呢.
很不屬靈的.
你回來教會講.
哇.
我看到旁邊那個炒股票.
炒到盤滿缽滿.
我的心都不知道有多妒忌.
就是.
我都想發財.
我妒忌是犯罪.
不過我不妒忌了.
不過我還是很羨慕.
就是.
你有這樣的想法.
不過你回來教會就不要講.
就是說.
世界是非我家.
就是.
盡量.
這些東西.
就是.
好像.
大家都.
收起來.
大家都自己面對就好了.
是不是?.
但是事實上.

$^{441}$久而久之.
信仰變成了.
你只會用一些.
屬靈的術語.
基督教的術語.
來去.
在信.
大家一起信的人面前.
在信徒面前.
來去.
就是.
大家都講一些.
大家都對的合理答案.
或者一個標準答案出來.
但是.
你真正的貪心.
有沒有被神對付呢?.
根本你沒有的.
因為你根本也不想去面對.
在信仰的層次.
你沒有面對過.
你每天繼續上班.
看到別人炒股票.
或者看到別人做一些不好的事情.
然後你又心裡就.
誰叫我信了耶穌.
我不可以和他這樣做.
如果我不是信耶穌.
我和他這樣做.
就怎樣怎樣.
就是.
就.
處於一個.
很.
可以這麼說.
好像很精神分裂的狀況.
就是.
是.
一方面.
每天在世界當中.

$^{481}$很多事情.
你見到世界就是這樣.
那些惡人真的樣樣順利.
說他肥就當然不好.
現在的人就說他瘦.
就是怎樣.
那些不信神的人.
他不用去教會.
不用什麼.
每天禮拜天去健身.
做到很fit.
他剛才說他肥到硬突.
他說不是.
現在說不信耶穌就很fit.
這樣.
可能是這樣.
但是.
你心靈裡面.
真實信仰.
是不是有出路呢.
或者不是單純信仰有出路.
而是真實.
可不可以經歷到那個屬靈的轉變呢.
那個就是一個問題.
很多時候我們.
需要明白一點.
這個作者本身.
而且記載在聖經那裡.
他就做了一個示範.
即使是一個屬靈的人.
一個跟從神的人.
他每天活在世上的時候.
他也有他的糾結.
有他的掙扎.
他也見到這個世界那些東西.
他真的.
這個不單是不明白.
他不是不明白.
他明白.
他就是看得太明白.

$^{521}$所以心裡面覺得很難受.
信仰.
第一步.
你要經歷到神對你的生命的改變.
第一步.
我覺得這裡是很重要的.
就是你go through你自己內心.
究竟掙扎著些什麼呢.
你有沒有想過.
你內心裡究竟掙扎著些什麼呢.
如果你內心裡自己掙扎著些什麼.
你都不願意去面對的時候.
神怎樣可以幫到你呢.
信仰的力量.
在乎神.
在你生命那種改變的時候.
你究竟自己掙扎著些什麼呢.
你很嫉妒人.
你又不承認.
你有時候對人.
很在乎人.
你的自尊心很低.
很在乎人家.
那個時候不斷地.
好像心裡邊上上下下被人扯著.
信仰那裡.
你不去神那裡.
神也幫不到你.
所以各方面的問題.
我想這個就是第一步.
阿薩這個敬拜的作者.
他給我們看到.
即使在敬拜的詩歌當中.
即使是他自己的當中.
他是活在一個真真實實.
罪人當道的世界.
那個世界不是完全美好的.
而事實上.
他活在這樣的時候.
他喝盡了滿杯的苦水.

$^{561}$但是.
問題就是.
究竟有沒有出路呢.
他說:我若說我要這樣講.
這就以奸詐待你的眾子.
我思索怎能明白這事.
眼看實是為難.
等我進了神的聖所.
思想他們的結局.
這裡呢.
就講出了.
一個挺有趣的轉變.
他剛才講的那些東西.
完全是他真實的想法.
他沒有掩飾的.
但是呢.
他講了出來的時候.
講了出來就發覺.
不對.
不對喔.
如果是這樣的時候.
好像我在騙人.
以詭詐待你的眾子的意思是什麼.
意思就是我在騙人.
其實這個.
如果用現在的說話來講.
就是說那個牧師在台上講道.
就像我這樣.
講道說.
多好多好.
但是他自己就很不值得.
隔壁的鄰居怎樣.
怎樣發達.
這樣.
其實.
他自己反省.
我不是真的在台上賣假藥.
其實是這樣的意思.
所以我不是以詭詐待你的眾子.
意思就是.

$^{601}$這樣的意思.
我不是好像在騙那些.
來跟我一起敬拜的人.
於是他就覺得不對.
於是他說.
我思索怎能明白.
這件事眼看實在為難.
他就陷入一個很.
兩難的狀況之後.
於是他就.
有一條出路.
他說.
等我進入神的聖所.
思想他們的結局.
他需要去到神的面前.
去想.
究竟.
究竟這個世界.
究竟一切.
是不是真的這樣呢.
我覺得這個就是非常重要.
我們.
剛剛過了這三十天.
和我們一起去.
去靈修.
去尋求神的.
聖所的時候.
其實就是在做這件事.
事實上.
我們每天活在一個世界當中.
很多事情是不斷的衝擊著你的信仰.
好像.
那個真實的.
和你信仰的理想.
是好像.
表面上是不合的.
這時候.
你心裡面很糾結的時候.
你又沒有出路.
只有兩條路.

$^{641}$可以走.
第一條路就不要跟從神.
有人是這樣做的.
他不是說不信耶穌,他仍然信耶穌.
不過只不過.
平時的生活就不會想神的事.
不跟從神,不會做神的事.
自己走自己.
禮拜日.
敬拜都可以的.
這是第一條路.
這樣就.
沒那麼為難.
因為平時.
人找游泳,你也找游泳.
沒那麼為難.
這是第一條路.
第二條路是怎樣呢?.
第二條路就是.
我要有一段時間.
好像要去.
進入聖所和神見面.
在神的面前.
去反省.
去面對我生活上面.
所遭遇的.
內心的掙扎.
這就是第二條路.
我們30天讀那本書.
不是說30天讀完就完的.
其實就是希望我們去建立.
我們真的需要.
有時候要停一停.
每天有一些時間.
進入到去神的聖所當中.
去整理.
因為要面對神.
才可以整理得到.
如果不是.
要不就走一條.

$^{681}$好像.
掛名.
做基督徒,星期日基督徒那條路.
如果不是,你也沒得出路.
所以它就.
第二條路就是.
要進入去聖所.
每天有一些時間去思想.
所以那本書叫流動的聖所.
就是要去思想.
在神的面前去反省.
這個是很重要的.
兄弟姐妹.
有一次我和一個朋友.
去聊天.
他告訴我.
他這陣子.
禱告什麼,怎麼禱告.
我一聽的時候.
我都覺得.
我都不懂得怎麼回應,因為都很特別.
我跟他說.
你知不知道我這陣子.
在掙扎什麼.
我都略知.
我都和他很熟.
你這陣子.
整天都在.
他不是炒股票,是炒外匯.
在香港的時候.
美國就很少人.
但香港人.
你知道有時候買不同的貨幣.
其實轉價.
不知道他炒那些.
我知道你炒外匯.
那些.
因為.
細節那些不說了.
所以我這陣子.

$^{721}$這樣禱告.
我講來聽聽.
神啊.
你知道我很想發達.
我也想賺多點錢.
你知道.
多點錢有個傍身.
和家裡都是好的.
所以我真的很想.
你幫我賺多點錢.
投資.
每件事都很順利.
這樣禱告.
我聽到他這樣講.
接著.
我都不懂怎樣回應.
即時.
好像退後一步.
不懂回應.
接著他再講.
他說.
神啊,你都知道我的心.
我貪錢.
那些事.
我不說你都知道.
不過我現在正正式式.
在你面前講.
因為我發覺.
我自己就改不了.
你覺得這樣不對.
不如你幫我改.
我自己這樣祈禱.
其實我都覺得.
很不對勁.
但我的心常常想這些.
我祈禱求上那些.
唯一.
你幫我.
當聽完下半段.
的時候.

$^{761}$我就很欣賞.
事實上.
很多時候我們的禱告.
本身.
我們不是這樣講.
我們有時候即使知道.
都是簡單輕描淡寫.
神啊,你都知道.
你教導,聖經教導,你不要貪心.
幫我,我不貪心就算了.
但事實上.
你完了禱告.
一張開眼睛之後.
然後眼睛仍然是.
一個S同兩棟.
看很多東西.
一樣的.
因為你.
其實你沒有高輸過.
你真實內心的掙扎.
你都沒有將.
掙扎真實的帶到.
神的面前.
以致到神很多東西是碰不到.
碰不到.
那些房門基本上是關上了.
即是我.
看到那個弟兄我很欣賞.
事實上.
慢慢慢慢發覺他有改變.
他自己都.
後來分享.
我發覺以前有些.
看得很.
心裡面抓得很緊的東西.
慢慢慢慢神給我.
有不同的看法.
而事實上真是這樣.
你的.
真實的掙扎.

$^{801}$你不講神都知道.
但是你.
願意真實的去面對.
譬如你一個人.
表面就好像和他.
以前.
有過節.
表面就好像.
覺得沒什麼.
但其實你心裡面那條刺根本都沒拔.
但在人面前.
當然要表現一個好的形象.
好像覺得很大方.
諸如此類.
你根本那條刺根本沒拔.
根本你自己.
一碰一碰那裡就痛.
即是背後.
所以有些小動作又出來了.
即是你自己可能都.
沒留意但身邊的人其實看到.
那.
但如果你高呼.
那些掙扎.
因為你自己那條刺根本插.
有時你都知道.
刺插進去很久不拔.
那怎麼辦.
那些肉.
生了黏住了.
自己拔不到.
一碰一碰就痛得飛起來.
怎麼拔.
而事實上是需要神.
才可以幫你拔.
但你不想碰那裡.
沒事沒事.
這樣根本上.
很多.
我們其實不是.

$^{841}$根本我們每個人.
都是一樣這樣.
我們就有很多這些傷口.
很多這些刺.
但我們在信仰那裡.
如果我們要經歷.
神的生命的改變.
但是沒變的.
沒變.
如果你回來教會的時候.
當然你學很多聖經.
但你懂得說是一回事.
你真正讓你的生命.
給了聖經來對付.
你真正要去.
面對它的時候.
又是另一回事.
所以.
這個作者.
他就說.
不行了.
我要帶到聖所那裡.
每天我們都需要將我們的生命.
帶到聖所那裡.
唯有在神的面前去思想.
人生又有不同的結局.
你實在把他們安在滑地.
使他們掉在沉淪之中.
在表面來看.
那個黑社會大佬.
這樣走出來.
又帶著一條驕傲的金鏈.
然後又說.
一有不對就要.
就說.
信不信我動你.
那種強暴的時候.
他很不順眼.
但在神面前來看的時候.
他就發覺到有些東西.

$^{881}$他自己是看不到的.
就說.
你把他們安在滑地.
使他們掉在沉淪之中.
多少時候.
你只看到人在這個世界.
那些順利的時候.
特別不跟從神.
就特別容易順利的時候.
你只看到那方面.
但你看不到.
那些不跟從人.
不跟從神的時候.
那種人心裡面.
那種不安.
那種沉淪.
事實上.
很多時候.
不是代表.
那個.
越順利.
物質越豐富.
那個人就越不沉淪.
你知不知道.
在.
其他國家我就不知道.
在美國或西方國家.
那些.
吃毒品的人.
當然.
後來吃到出事.
工都沒了.
就變了.
無家可歸.
當然是走一條路.
但是很多時候.
那些吃毒品的人.
不是窮到窿的人吃毒品.
是那些.
初初的時候.

$^{921}$因為很無聊.
有錢就.
就.
但是心裡面很空虛.
去派對.
或者吃毒品.
聽過一個.
真實的例子.
是真實的.
是Vivian告訴我的.
Vivian.
我忘記那次她去什麼訓練.
那個訓練.
她就說.
有一個警察.
跟我們分享.
不知道說什麼.
但是說了一個很真實的.
案例.
她說你不要以為吃毒品的人.
是沒事的.
其實警察還想帶出這一點.
但是她的故事就是這樣.
很真實.
她說她認識.
有兩個.
專業人士.
是那些.
都很高收入的.
專業人士.
但是他們就很無聊.
因為.
又沒有信仰.
最多去派對.
很無聊.
兩夫妻表面正常.
又很好.
很無聊.
她說不如這樣.
我們試一下吸毒.

$^{961}$兩夫妻約在一起.
買些毒品在家.
試一下吸毒.
但是那個警察就說.
他平時沒有.
他只是試一次.
那些毒品.
有些不太順.
那個.
老婆.
就破壞了腦.
接著.
就很有暴力傾向.
接著後來.
不說了.
整個家庭都毀了.
兩個人的職業.
都毀了.
這個.
只不過是一個.
想說出一件事.
你只看到表面風光的東西.
我們只看到這些.
你看不到人心靈裡面的空虛.
但是.
事實上好像有一個.
沉淪的狀況.
你是看不到.
人心靈裡面的空虛.
所以作者就說.
當他去到神的面前.
好像他對於人生.
不是只看到表面那麼簡單.
多了一層的.
靈悟.
看到那些人即使是.
他順利又怎樣.
他很多東西.
吃到眼睛凸了出來.
又怎樣.

$^{1001}$可能他在沉淪當中.
更加他就說.
他們轉眼之間成了何等的荒涼.
他們被驚恐滅盡了.
就說到.
你表面看他很順利.
但是其實他內心那種恐懼.
可能你不知道.
事實上.
越是.
有錢的人的家.
越多監控.
那些camera.
這個是.
不真的事實.
你沒錢你都沒錢裝.
那麼多鏡頭.
這個是事實.
是事實.
心靈裡面那種恐懼.
你擁有的時候.
怕失去.
不說那個打劫.
就算你所擁有的東西.
即使你擁有很多.
你都怕失去.
那種又是另一種.
恐懼.
人生就好像.
你不斷要用自己.
去積蓄一些東西.
用自己的力量不斷去維持.
很累的一個過程.
那種你怕失去那種恐懼.
你又看不到.
你只看到別人順利.
是不是.
所以當先人.
他就再說.
人睡醒了怎樣看夢主.

$^{1041}$你醒了也必照樣.
看他們的影像.
更加他就看到.
原來在神的面前.
去真實的.
去看人生的時候.
他發覺原來有一個觀點.
他以前是沒有的.
就是人生命.
是不是有重量.
影像就沒有重量.
特別是他用輕看他們.
發夢就像影像.
影像是沒有重量的.
因為假的.
那些東西.
一關燈就沒有了.
關了報應機就沒有了.
人生那些追逐的東西.
其實擁有那些東西.
是有多少重量呢.
有多少重量呢.
事實上.
現在就很少人.
去想的.
但以前很多人.
因為人很短命.
其實以前的人.
直到.
未發明盤尼西林.
即是未發明那些antibiotics.
你知道antibiotics.
很後才發明的.
我忘記了是第一次還是第二次.
世界大戰才發明的.
因為那時候.
打仗死的人多的時候.
才mask production.
那些antibiotics.
出來的.

$^{1081}$之前那些人一有一點病.
就沒有抗生素就死了.
所以以前的人平均壽命.
你知不知道多少歲.
三十多歲.
我可能.
你聽的早一點.
我看的是四十多歲.
四十多歲.
我看到大部份.
我自己都過了.
所以以前的人.
因為他們很短命.
所以他們.
經常都想.
他們的生命的價值.
在他們的墓碑.
會.
怎樣記得他們的人生.
人人在墓碑.
通常有些外國人.
唐人就沒有.
唐人是會很少copy給你.
但外國人會自己想一句.
然後寫下去.
那時候.
他們很年輕的時候.
就已經想好將來那句.
要寫什麼.
要做一個怎樣的人.
是被人記得到.
我覺得某個程度.
其實就講人生那種重量.
你留得低的是什麼.
因為你重量的意思就是有些東西.
是紮在那裡的,留得低的.
即是.
你表面看.
很豐盛,那些惡人很厲害.
賺到賠滿半滿.

$^{1121}$但是他留得低的是什麼呢.
人生當中的重量.
是有沒有呢.
然後他就再說.
因而我的心裡.
發酸,肺腑被刺.
當他想到這幾方面的時候.
他進入到聖所.
不單只是給世界.
不斷告訴那些灌輸.
給他們知道,這樣就是順利.
這樣就是人生,這樣就是豐盛.
他去到神面前.
再一次來看人生的時候.
他看到剛才那些.
心靈的重量.
看回剛才那些.
沒有神的人的心靈那種驚恐.
各方面的時候.
他就說.
他心裡好像突然間.
有一種酸溜溜的感覺.
那種感覺就是.
肺腑被刺.
突然間好像整個人都.
很大的震撼.
就是.
究竟我的人生又要怎樣呢.
我要選一個什麼樣的人生呢.
他說.
我只是愚昧無知在你面前.
蓄慮一樣的意思.
不只是罵人.
禽獸不如.
蓄慮一樣的意思,蓄慮人和動物的分別.
就是說.
沒有跟神聯繫.
因為人.
和動物最大的分別就是.
人的靈魂是可以跟神聯繫的.

$^{1161}$如果我沒有跟神聯繫.
那我就跟.
禽獸一樣.
其實意思是這樣.
很大的震撼.
我只是看到人生.
只是這個世界表面各樣的東西.
或者只是今生的時候.
我就跟禽獸一樣.
事實上.
禽獸不是說.
只是降低的.
其實是很可愛的.
你家裡有養過寵物.
那些寵物要你.
要你拿你的心.
有時候還多過.
一個真實的人.
黏著你.
每天撐著你.
不是說不可愛.
而是說沒有了跟神聯繫的部分.
沒有了跟神聯繫的部分.
人生沒有了跟神聯繫的部分.
那靜下來的.
又是什麼呢.
他說.
然而我常與你同在.
你殘符著我的右手.
所以詩人.
當他進入至聖所.
去反省人生的時候.
他就發覺.
跟神聯繫的部分是很寶貴的.
我可以每天和神同在.
讓神帶領著我.
這個他不是說.
你常與我同在.
這裡是說什麼.
是我常與你同在.

$^{1201}$誰是主誰是.
什麼.
這個反過來說.
很多時候我們說主與你同在.
但其實這裡他說.
我與主同在.
這裡是說出了一個選擇.
是說出了一個.
主動的一個揀選.
主與你同在是一個願望.
我們有時候見到人.
主神與你同在.
這是一個.
問候的一個願望.
但我與神同在.
就是我的選擇.
你可以選擇.
當你與神同在的時候.
神就會帶著你的右手.
你不會怕.
你選擇神同在.
神不會帶你.
然後就纏著你的右手.
什麼意思呢.
你要以你的訓詞引導我.
以後必接我進.
榮耀裡.
神的說話.
神自己的教訓.
經常都會帶著我.
所以每天去到靈修.
去到那個聖所.
你的人生慢慢慢慢.
你發覺有神帶你走的路會不同.
你整個人生會不同.
你是在預備自己.
將來去進榮耀.
這裡所說.
你表面來看那些惡人很榮耀.
但再在神面前的時候.

$^{1241}$他就知道.
將來有一個榮耀.
我現在在預備.
神帶我進去.
人生在世.
其實我們是很短的.
有時候.
我都去想.
在耶穌在世的時候.
就叫過.
三個人復活.
其中一個.
就是拉薩路.
很出名的.
《約翰福音》十一章.
拉薩路是耶穌的朋友.
應該是一個.
正直成年壯年的男士.
他就患了急病.
進了墳墓死了.
耶穌就叫他出來.
拉薩路就是歷史上.
有名有姓的.
一個人.
是從死那邊.
真實復活.
這是無可置疑的.
因為耶穌其他那兩個復活.
都沒有仔細提他的名字.
因為那兩個小朋友.
拉薩路是有的.
所以他們是可以追溯.
那些猶太人.
在《約翰福音》的時候.
他們知道.
拉薩路就是那個拉薩路.
就是死的那個.
但是.
有時候我就問了我的問題.
拉薩路被耶穌復活了之後.

$^{1281}$他有沒有死呢?.
你有沒有想過?.
會不會死呢?.
應該都死吧.
如果不死.
會不會看電影.
那些Highlander.
不是吧.
我沒有聽過.
應該會死的.
如果拉薩路真的死了.
最多都是幾十年.
耶穌復活.
當他三十歲.
跟耶穌差不多年紀.
那時候的人.
長命都是八十歲.
那他就多活五十年.
又死了.
那我為什麼會說這個事情呢.
就是說如果你寄望今生.
即使你死了.
耶穌把你復活.
都會再死的.
因為都是很短的.
但我們的寄望不是今生.
他會接我們去到榮耀裡面.
求神幫助我們.
在結束的時候.
所以詩人就說.
除你意外在天上我有誰呢.
除你意外在地上.
我又沒有所愛無的.
我的肉體和我的心腸雖殘.
但神是我心裡的力量.
又是我的福分直到永遠.
詩人當他看清楚人生.
看清楚自己.
看清楚這個世界的時候.
他就知道他愛什麼.

$^{1321}$有時候你看不清楚什麼.
你不知道愛什麼.
有時候.
有些人.
看過什麼食物製作那些.
他就知道不愛吃什麼.
有時候我們就看不清.
我愛的那些東西.
看清你就不愛.
現在讓我們去知道.
每天在神面前.
看清楚這個世界.
看清楚我們自己的內心.
看清楚我們自己的掙扎的時候.
我們知道.
真正.
神是真的很可愛.
而當日.
最後每個人都要死亡.
這個世界一切的東西都會完結.
但是他說親近神.
是與我有益的.
因為神就是我們的避難所.
神是我們的避難所.
一開始的時候.
這個世界好像不斷的侵入.
信仰的基礎.
好像被掏空了.
他自己都好像.
很容易跌倒了.
但是有什麼方法呢.
就是進入聖所.
這個聖所就是一個避難所.
好像我們每個人的心靈.
好像有一個房間.
那裡是除了給神的.
沒有東西可以侵入.
那個就是我們每個人的避難所.
有神幫助我們.
\newpage



\section{士師記 2:19}
\label{sec:032232}
\textbf{聖經人聖經事:士師記的惡性循環}~~~~黃克勤牧師博士
\newline
\newline
連結: \href{http://www.church.com.hk/acms/content.asp?site=cdc\&op=show\&type=product\&code=032232\&layout=sermon}{\texttt{032232}} ~~~~ 語音日期: 2026-06-02
\newline
\newline
\hyperref[sec:032231]{\small{< < < PREV SERMON < < <}}
~
\hyperref[sec:index_chronic]{\small{[返順時目]}}
~
\hyperref[sec:032233]{\small{> > > NEXT SERMON > > >}}
\newline
\newline
士師記 2:19
\newline
\begin{longtable}{cl}
\hline
\hline
章節 & 經文 (和合本修訂版)\\
\hline
2:19 & \begin{tabularx}{0.7\textwidth}{X} 但士師一死，他們又轉去行惡，比他們祖宗更壞，去隨從別神，事奉叩拜它們，總不放棄他們的惡習和頑梗的行為。 \end{tabularx} \\ \\
[1ex]
\hline
\hline
\end{longtable}
$^{1}$大家好,一起讀聖經,生命樹陪著你.
每一天都是聖經人,聖經時.
看看上帝在後面做甚麼工作.
認識祂多些.
亦看到人自我認識頗軟弱.
今天繼續看《士師記》.
由於我們已入到《士師記》.
在歷史書裡約書記之後.
下一卷就到《士師》.
《士師記》有一個很有趣的惡性循環.
不停繞圈,而且是負面的.
所以是惡性循環.
看看.
《士師記》第二章第十四節.
藍色是第一點.
看看字先.
「耶和華的怒氣向以色列人發作」.
螢光了的就是要講的主題重點.
第一點.
「耶和華的怒氣向以色列人發作」.
「就把他們交在搶奪他們的人手中」.
「又將他們賦予四位仇敵的手中」.
「甚至他們在仇敵面前再不能站立得住」.
看看這一點.
「耶和華的怒氣向以色列人發作」.
發作的時候是怎樣發作的呢?.
就是將他們交在搶奪他們的人手中.
讓被搶奪的人成功搶到他們的東西.
怎樣搶到他們的東西呢?.
由於以色列人附近沒有聽到神的話.
沒有跟著神當日跟他們納入迦南地納的約.
「我送你們進去,幫你們趕走當地的人」.
「你們要麼殺了他們,要麼全趕走」.
「或者牛羊生畜都不要」.
「最重要是毀壞他們拜神的壇」.
「不可以跟他們來往,不可以通婚」.
「不可以事奉他們的神,記得」.
他們會這樣進去,誰知沒有做.
沒有做的時候,其實是.
「耶和華的怒氣向以色列人發作」.

$^{41}$就是將他們交在搶奪他們的人手中.
剩下的當地的民族.
他們站穩之後.
他們或者在附近的民族走進來搶他們的東西.
怎樣搶呢?何時搶呢?.
如果要搶你的農作物.
會不會在你播種的時候搶呢?.
當然不會.
等你收割好的時候去搶.
因為這樣才一年半載的.
老虎的耕種一下子搶完.
他根本不用做,馬上跑上山.
哭得不得了,又沒有東西吃,肚子餓了.
怎算呢?還殺死了一些人,或者被他們擄走.
本來是神保護他們的.
不過神的怒氣向他們發作.
就交在搶奪他們的人手中.
或者交在四圍的仇敵手中.
而且令到他們站立不住.
即是說輸了,打仗輸了,沒得打,很慘.
他們無論往何處去.
耶和華都以災禍攻擊他們.
正如耶和華所說的話.
猶如耶和華向他們所起的誓.
令到他們極其困苦.
你看到第一點說的是.
神向他們發怒的時候.
是將他們交給敵人,搶奪他們的人.
而且令到他們去到哪裡都受災攻擊.
或者神直接攻擊,或者其他敵人攻擊他們.
而令到他們死的時.
就是極其困苦,不單止困苦,是極其困苦.
很慘,很痛.
這就是神讓他們遭遇到的.
因為神的怒氣向他們發作.
所以弟兄姊妹要留意.
當我們沒有聽神的話.
或者不照著做.
神要我們學得好.
他就會讓某些災淋到我們當中.

$^{81}$於是我們就學到東西了.
那時候就極其痛苦.
極其痛苦就會反省.
反省了,但如果你早些跟著做.
就不需要經歷這些極其痛苦.
我們很多時候都是這樣.
極其痛苦.
於是,神為什麼要你這樣.
於是,當他們求神了.
於是第二點來了.
「耶和華興起事司」.
「事司就拯救他們脫離搶奪他們的手」.
「他們卻不聽從事司,竟隨從叩拜別神行了邪淫」.
「速速地偏離他們滅祖所行的道」.
「不如他們滅祖順從耶和華的命令」.
這裡講得簡單了.
當他們極其痛苦之後,他們會做一件事.
「神啊,好慘啊,為什麼你不理我們啊」.
神就說「好,既然你學到東西,懂得回頭了」.
「耶和華興起事司」.
「事司」的英文字是Judges.
是指很多審判官的意思.
事司這個字的希伯來文.
就是法官,審判官,拯救者,官員,長官.
事司並不是全國的.
是12個支派裡的某個支派的領袖.
當時約書人不在了,那些大人物不在了.
於是每一個支派就在做自己的事.
不過跟隔壁支派是兄弟,有很多來往.
「哎呀,好慘啊,其他的兄弟未必幫到」.
很痛苦,極其困苦的時候.
他們哀求上帝,神就為他們興起.
自己支派裡的事司.
即是一個小規模,小型的民族英雄.
這些事司通常都很有信心,信仰很好的.
事司就拯救他們,脫離搶奪他們的手.
不過無奈.
不過一對抗完之後.
他又不肯聽事司說了.
竟然又叩拜別神.

$^{121}$轉得很快,極其痛苦的時候.
神就說「為什麼你這樣對我們,我們願意跟隨你」.
然後神就給他一個事司興起.
幫了他們,救了他們.
然後立刻轉頭跟從別神叩拜,即是敬拜.
而且敬拜神行了邪淫,什麼行了邪淫呢?.
因為其實敬拜其他的很多當地的神.
都會有廟祭,男祭,女祭在做什麼呢?.
原來當時他們會覺得.
怎樣可以跟神更加親近呢?.
可以跟敬拜的時候有的男祭,女祭交合.
於是就會跟神更加親近.
有這樣的感覺.
拜別神又行了邪淫.
速速地偏離他們列祖所行的道.
不像他們的列祖這樣順從耶和華的命令.
又來?剛剛有事司,你又這樣了.
於是耶和華為他們又興起事司.
而且神是「與那事司同在」.
事司在世上的一切日子.
耶和華會做什麼呢?.
是拯救他們脫離受敵的手.
他們因受欺壓妖害.
哀聲嘆氣.
所以耶和華就後悔了.
「神啊,我們真是對不起你啊,對不起啦」.
「哎呀,對不起對不起,我們又忘記了你啦」.
「又去拜別的神啦,真是不好啦」.
「求你救我們啦」.
耶和華就後悔.
好像對兒子一樣.
「哎呀,是不是這樣,是不是要打呀?」.
「是呀,不過哎呀,細力一點啦」.
「其實我都不應該的,對不起呀,爸爸」.
這個後悔的意思.
神的後悔.
後悔即是撤回祂對他們的打擊.
後悔的意思是.
我轉念幫助他們.
這個就是事司不停出現的事情.

$^{161}$我們看看吧.
第三點.
「及至事司死後,他們就轉去行惡」.
其實這班人很有趣.
耶和華興起事司.
事司第一點是耶和華怒氣發作.
因為他們拜其他神.
於是敵人就來了.
敵人來了就慘了,於是神就興起事司.
事司一幫他戰勝敵人,他就不聽從了.
快到這樣.
拜別神又行了邪任.
於是神又打擊他.
耶和華安排的事司是跟事司同在的.
於是事司在世上的日子就拯救了他們.
不過他們事時都未拯救完.
又壞了,又受欺負迫害.
事司又起來打.
最後那些哀聲「哎呀,好慘啊」.
於是神就說「好吧,我又找事司幫你」.
不過通常最後都去到第三點.
「事司死後,他們就轉去行惡」.
一死之後就馬上轉身.
「行惡,比他們列祖更甚」.
「又去事奉偷拜別神」.
「總不斷絕橫梗的惡行」.
於是這裡不停地發生的事是什麼?.
第一點,他們就轉去跟從別神.
於是耶和華怒氣發作.
其他敵人就來搶奪攻打他們.
然後他們就很慘了.
極其困苦,於是就哀聲求告耶和華.
耶和華就說「好吧,我就興起個事司來幫你守」.
事司就幫他們打仗.
「哎呀,太好了,我們跟隨主」.
不過一平定之後,又轉身去跟隨別神.
為什麼呢?因為他沒有把別神,廟,人,祭司.
那些假的,不好的,邪的東西全部毀滅.
所以一得勝之後,轉頭又變壞了.
於是神又叫敵人打他.

$^{201}$一打他就說「哎呀,好慘呀,好痛呀」.
於是極其困苦,於是神聽到他們的聲音.
「好吧,仔即係仔,好,我又叫個事司幫手」.
又與事司同在,又打退敵人.
然後他們一轉身又來.
最壞的時候是什麼呢?.
就是連事司都死了.
全部轉身去跟隨別神.
絕不斷絕橫梗的惡行.
所以這個形容是橫梗.
簡直是一邊頸.
嘩,真是不聽話.
這個就是事司裡面不停地出現的一個惡性循環.
令到他們很慘.
這個惡性循環是什麼呢?.
就是拜其他的假神.
神很不開心.
於是耶和華怒氣發作.
怒氣發作就被其他的外族進來入侵.
於是神就用外面那些其他的民族.
好像一條皮鞭那樣打他.
鞭打他,懲罰他,打醒他.
一打,「嘩,好痛,醒了,神呀,對不起」.
於是「好,我找個事司幫你」.
一幫了,馬上轉身.
又拜過那些假神.
於是事司又幫手.
何時事司死了,馬上轉身.
「啊,好慘啊」.
於是神又用外族鞭打他.
其實我早一兩次就講了.
神為什麼那麼「咩」呢?.
當他們入迦南的時候.
要他們將所有人趕走滅絕.
然後甚至牛羊都不要.
用你自己帶進來的就可以了.
就不要顧惜.
因為那些將會成為你.
「立旁上的荊棘」.
「嘩,如果這裡被荊棘刺著會好痛」.

$^{241}$或者成為你們的網羅.
捉飛鳥的網羅.
你一跌進去陷阱就會網住你.
走不掉.
真的會這樣的,很清楚.
誰知道他們硬是不聽.
於是其實為什麼要滅絕那些迦南人呢?.
因為其實迦南人的惡.
神是要有一天懲罰他的.
就用以色列人時期到了.
時機到了就用以色列人進去.
就鞭打他們.
懲罰他們的惡.
但是當以色列頑皮的時候.
神就同樣用外族的人.
成為鞭打的皮鞭.
懲罰的皮鞭.
又令以色列人好一點.
但是以色列人終於出現的.
就是一個惡性循環.
是非常還緊.
還緊就是說.
不肯令.
很倔強地去拜那些假神.
其實是自己做了算.
當日沒有聽神說.
於是這些外族人.
和那些假神成為.
勒住他們.
勒旁的荊棘.
成為他們的網.
網住他們.
那麼士司祭的惡性循環.
跟大家看到這裡.
是不停地循環的.
頑皮了,神的怒氣發作.
然後他們哀求.
神就給個士司.
令他們好一點.
但是過了之後又作惡.

$^{281}$作惡神的怒氣發作.
於是又有外族來侵擾.
然後又哀求.
哀求神又回轉.
又給個士司.
就是這樣不停地惡性循環.
那麼我們的人生都是這樣的.
教徒的生命,屬靈生命.
很多時候都是這樣.
希望我們打破惡性循環.
一路不要這樣圈圈地走.
走很久以為很遠.
因為我們要離開.
打開惡性循環,一條直線.
一條直線.
就跟隨主了.
這樣就最好了,我們祈禱吧.
天父求你幫助我們.
士司祭是教導我們一件事.
原來.
祂是指出.
原來我們人.
是這麼環境活躍的.
只要平順一點.
我們就會.
轉了頭去看那些.
虛假的神.
看那些人生的享受.
快感,就不記得.
望向我們的神,祂想我們怎樣.
那些若忘記了.
於是神的怒氣會向我們.
發作,發作之後.
我們又想,怎麼會這樣.
這麼慘,求你憐憫我們吧.
於是神又給我們.
一個恩典.
給我們一位士司,帶領我們可以.
轉頭,然後我們又環境過.
於是你的怒氣又發作.

$^{321}$我們人生很多時候都是.
不停地繞圈,繞圈.
其實神是想我們一條直線.
向著你,一直奔跑.
這樣就最好了.
求主給我們.
給我們定睛在你身上.
求主耶穌幫助.
求天父幫助.
求聖靈幫助.
誰聽祈禱,奉靠耶穌基督的名.
阿們.
\newpage



\section{民數記 12:1-16}
\label{sec:032233}
\textbf{從比較到復和的兄姊弟情}~~~~何嘉敏傳道
\newline
\newline
連結: \href{http://www.church.com.hk/acms/content.asp?site=cdc\&op=show\&type=product\&code=032233\&layout=sermon}{\texttt{032233}} ~~~~ 語音日期: 2026-05-30
\newline
\newline
\hyperref[sec:032232]{\small{< < < PREV SERMON < < <}}
~
\hyperref[sec:index_chronic]{\small{[返順時目]}}
~
\hyperref[sec:032234]{\small{> > > NEXT SERMON > > >}}
\newline
\newline
民數記 12:1-16
\newline
\begin{longtable}{cl}
\hline
\hline
章節 & 經文 (和合本修訂版)\\
\hline
12:1 & \begin{tabularx}{0.7\textwidth}{X} 摩西娶了古實女子為妻。米利暗和亞倫因他娶了古實女子就批評他， \end{tabularx} \\ \\ \relax
12:2 & \begin{tabularx}{0.7\textwidth}{X} 他們說：「難道耶和華只與摩西說話嗎？他不也與我們說話嗎？」耶和華聽見了。 \end{tabularx} \\ \\ \relax
12:3 & \begin{tabularx}{0.7\textwidth}{X} 摩西為人極其謙和，勝過地面上的任何人。 \end{tabularx} \\ \\ \relax
12:4 & \begin{tabularx}{0.7\textwidth}{X} 忽然，耶和華對摩西、亞倫和米利暗說：「你們三個人都出來，到會幕這裡。」他們三個人就出來了。 \end{tabularx} \\ \\ \relax
12:5 & \begin{tabularx}{0.7\textwidth}{X} 耶和華在雲柱中降臨，停在會幕門口，叫亞倫和米利暗。二人就出來， \end{tabularx} \\ \\ \relax
12:6 & \begin{tabularx}{0.7\textwidth}{X} 耶和華說：「你們要聽我的話：你們中間若有先知，我－耶和華必在異象中向他顯現，在夢中與他說話； \end{tabularx} \\ \\ \relax
12:7 & \begin{tabularx}{0.7\textwidth}{X} 但我的僕人摩西不是這樣，他在我全家是盡忠的。 \end{tabularx} \\ \\ \relax
12:8 & \begin{tabularx}{0.7\textwidth}{X} 我與他面對面說話，清清楚楚，不用謎語，他甚至看見我的形像。你們為何批評我的僕人摩西而不懼怕呢？」 \end{tabularx} \\ \\ \relax
12:9 & \begin{tabularx}{0.7\textwidth}{X} 耶和華向他們怒氣發作，就離開了。 \end{tabularx} \\ \\ \relax
12:10 & \begin{tabularx}{0.7\textwidth}{X} 當雲彩從帳幕上離開時，看哪，米利暗長了痲瘋，像雪那麼白。亞倫轉向米利暗，看哪，她長了痲瘋。 \end{tabularx} \\ \\ \relax
12:11 & \begin{tabularx}{0.7\textwidth}{X} 亞倫對摩西說：「我主啊，求你不要因我們愚昧，因我們犯罪，就將這罪加在我們身上。 \end{tabularx} \\ \\ \relax
12:12 & \begin{tabularx}{0.7\textwidth}{X} 求你不要使她像那一出母腹、肉已侵蝕了一半的死胎。」 \end{tabularx} \\ \\ \relax
12:13 & \begin{tabularx}{0.7\textwidth}{X} 於是摩西哀求耶和華說：「神啊，求你醫治她！」 \end{tabularx} \\ \\ \relax
12:14 & \begin{tabularx}{0.7\textwidth}{X} 耶和華對摩西說：「她父親若吐唾沫在她臉上，她豈不蒙羞七天嗎？現在要把她隔離在營外七天，然後才領她回來。」 \end{tabularx} \\ \\ \relax
12:15 & \begin{tabularx}{0.7\textwidth}{X} 於是米利暗被隔離在營外七天；百姓沒有起程，直等到米利暗回來。 \end{tabularx} \\ \\ \relax
12:16 & \begin{tabularx}{0.7\textwidth}{X} 以後百姓從哈洗錄起行，來到巴蘭的曠野安營。 \end{tabularx} \\ \\
[1ex]
\hline
\hline
\end{longtable}
$^{1}$各位弟兄姊妹平安.
很感恩可以再次回到大堂分享訊息.
我其實有一個弟弟.
我也是一個頗疼弟弟的姐姐.
我記起以前生日買禮物給他的時候.
我會拿很大部份我的零用錢.
去買一輛很漂亮的車給我弟弟.
不是說真的那些車.
即是玩具車.
價值大概是二百五十元左右.
而大家又想不想猜一下.
我弟弟小時候又買些什麼給我呢.
其實我弟弟現在也很好.
我不是說他不好.
他就會在他的錢包裡.
拿一個五元的硬幣給我.
沒有多沒有少.
是他小時候給我的心意.
其實今天我們在看聖經的故事裡面.
我們也會看到一個姐姐和弟弟的關係.
還有哥哥.
其實姐姐和哥哥也很疼弟弟.
但是今天相信大家看到.
提到的題目的時候.
也會看到當中有一些比較和疾度.
就讓我們透過今天的經文.
文述記十二章.
這個故事裡面.
我們有一些的亮光和學習.
讓我們一起同心吐槁.
親愛的天父.
你也讓我們能夠成為兄弟姐妹.
成為弟兄姊妹.
我們可以在這份的.
無論是有血緣或者沒有血緣的關係.
的家庭關係裡面有很多的學習.
並且可以在這個關係裡面.
彰顯上帝你的愛.
求主打開我們的心.
打開我們的眼.

$^{41}$打開我們的耳朵.
讓我們能夠聽到.
上帝你今天要跟我們說的話.
也求主高沒你孩子的口.
讓我能夠講得清楚又明白.
我們這樣吐槁奉主名求.
阿們.
今天我們會繼續看文述記的十二章.
我們上次看到其實以色列民他們.
開始有些埋怨和投訴.
我們可以看到摩西.
他在領導上面其實是面對著.
種種的張力和挑戰.
上次我們看到百姓他們發怨言.
到今天我們會看到摩西的哥哥和姐姐.
他們產生了一些妒忌對弟弟摩西.
並且最後我們會看到.
上帝他親自的保護.
臨到摩西他的僕人身上.
讓我們一起首先去讀以下的經文.
我們123就起了123.
摩西娶了古實女子為妻.
米尼暗和雅倫因她娶了古實女子.
就批評她.
她們說難道耶和華只與摩西說話嗎.
她也不與我們說話嗎.
耶和華聽見了摩西對日記的祈禱承諾並不再任何人.
華對摩西雅倫和米尼暗說.
你們三個人都說話.
我會看著你.
你們三個人都說話.
你可不可以再民主中國.
你再回望問候.
耶和華對摩西雅倫跟米尼暗說.
你們三個人都說話.
我會看著你.
你們三個人都說話.
我會看著你.
你們三個人都說話.
我會看著你.

$^{81}$你們三個人都說話.
我會看著你.
你們三個人都說話.
我會看著你.
你們三個人都說話.
我會看著你.
你們三個人都說話.
我會看著你.
你們三個人都說話.
我會看著你.
你們三個人都說話.
我會看著你.
你們三個人都說話.
我會看著你.
你們三個人都說話.
我會看著你.
你們三個人都說話.
我會看著你.
你們三個人都說話.
我會看著你.
你們三個人都說話.
我會看著你.
你們三個人都說話.
我會看著你.
你們三個人都說話.
我會看著你.
你們三個人都說話.
我會看著你.
你們三個人都說話.
我會看著你.
你們三個人都說話.
我會看著你.
你們三個人都說話.
我會看著你.
你們三個人都說話.
我會看著你.
你們三個人都說話.
我會看著你.
你們三個人都說話.
我會看著你.

$^{121}$你們三個人都說話.
我會看著你.
你們三個人都說話.
我會看著你.
你們三個人都說話.
我會看著你.
你們三個人都說話.
我會看著你.
摩西是一個極其謙和的人.
我們剛才讀經文的時候.
我們甚至見到摩西那種謙卑.
在英文的NIV譯本.
謙和也譯做Humble.
他那種謙卑是勝過地面上的任何人.
不知道大家身邊有沒有認識一些.
你認為都是很謙卑的人.
為什麼一個人可以那麼謙卑呢.
謙卑這個特質不是我們每一個人最想追求的事.
因為有時候我們會誤解.
謙卑可能是作為弱者的意思.
但我們在摩西這個故事裡面.
我們見到上帝正正就是很喜歡.
這些那麼謙卑的人.
其實越謙卑的人上帝越重用.
首先我們見到摩西他和上帝的那份關係.
是他會馴服上帝的指示.
就算他覺得自己可能不夠好去做.
但是他都馴服去做.
我們見到摩西的領導.
其實是將忠心放在上帝的話語當中.
而且他的領導不是自我的.
他是清楚的出於上帝的呼召.
並且摩西是一個為百姓代求的領導.
當百姓闖禍的時候.
他不會先想自己.
而是他謙卑的在神面前為百姓代求.
摩西那份謙卑.
就是讓他可以馴服跟隨神.
一個很重要的特質.
作為上帝的僕人.

$^{161}$我們見到摩西在面對批評的時候.
他並沒有想要去還口.
要去Defend自己.
而是他的回應就是不回應.
單單安靜在上帝的面前.
我們也嘗試把這個.
有米利暗和弟弟的鏡頭.
好像倒帶回到摩西還是BB的時候.
我們見到米利暗在出埃及的2章已經有描述到他.
在當時法老的迫迫下.
所有希伯來人的兒子都要扔進尼羅河裡浸死.
相信這個故事大家都很熟悉.
摩西的姐姐為了確保摩西BB的安全.
她當時在一個籃子裡放了一條河.
確保了她的安全.
當法老的女兒見到籃子裡的BB.
她就動了慈心要去照顧BB.
米利暗 摩西的姐姐.
她也很聰明地想了一個策略或方法.
她就向法老的女兒提出.
她要先照顧BB.
最後她找了希伯來人的女兒去照顧.
其實就是找摩西的親生母親去照顧摩西.
從而摩西真的避過一劫.
所以我們可以見到米利暗小時候已經是一個有勇氣的姐姐.
是一個有策略的人.
她的那種看顧其實是很實在的.
雖然我們在經文讀到.
米利暗的姐姐和哥哥在批評她.
但那份在《出埃及記》裡的姐弟情其實是很實在的.
我相信在座有兄弟姐妹的會眾.
我們會有愛在這種兄弟姐妹的關係裡.
雖然隨著時代的變化.
社會給我們的壓力.
可能這份愛的關係會變得越來越模糊.
我回想起自己有沒有一些嫉妒我弟弟的一幕.
我認知自己這樣做的批評是不對的.
其實每一次我要去說教完批評完弟弟之後.
其實我都有一種不太平安和憤怒在自己的裡面.
當我有這個審查的時候.

$^{201}$我現在和弟弟FaceTime的時候.
我就會勒住自己的舌頭避免去批評.
反而是要說一些造就他的說話.
把所有的事情帶回上帝的面前.
因為深信神一定會帶領他的前路.
雖然不知道他覺不覺.
但是我都盼望弟弟回來香港的時候.
我都可以和他講清楚.
當我自己有被聖靈提醒的時候.
我真的感受到和弟弟的溝通和關係輕散了.
並且沒有了以前那種的苦毒.
不用好像很怒氣沖沖地和他講話.
是上帝改變了我的心.
說回經文.
我們在民數記十二章見到米尼暗和亞倫.
他們在批評弟弟.
我們剛才都有讀到.
哥哥和姐姐和弟弟在比較.
他和弟弟在比較什麼呢.
其實他和弟弟摩西在比較他們的能力.
和弟弟比較和神的關係誰比較近呢.
究竟誰和神的關係比較近呢.
可能有時候在教會.
我們都無意中會有這些比較.
我們可能不知不覺.
潛意識在我們當中.
會不會有時候和身邊的弟兄姊妹去比較呢.
會不會不知不覺都會有一些屬靈的驕傲.
其實教會.
我們未必每個人都有血緣關係的弟兄姊妹.
但是我們在教會.
回來教會其實坐在你旁邊的.
都是你的弟兄姊妹.
我們的內心會不會有時候都有以下的想法呢.
我在教會有很多的事奉.
做得比你多.
我和神的關係一定比較近.
或者我們每天都有靈修.
還會發給別人看.
分享.

$^{241}$而我見到旁邊的弟兄姊妹.
好像沒有經常靈修.
有時候我們會不會有這些比較呢.
因為當我們去比較的時候.
如果我們不去在神的面前去梳理去認罪的時候.
其實會產生那種的嫉妒.
而當嫉妒去發芽的時候.
其實是會轉化成為一個行為.
而在這裡我們見到米尼暗和雅倫.
那個行為就是他們去批評.
批評弟弟摩西.
在另一邊廂的摩西.
被人批評的摩西.
我們見到他其實是柔和謙卑.
摩西其實都沒有馬上去回應.
哥哥和姐姐的批評.
如果我們是摩西的話.
你聽到哥哥和姐姐去批評你的時候.
不知道大家會怎樣反應呢.
如果我是摩西.
我是會馬上去維護自己.
當然不是.
我娶另外一半.
都是好的.
可能會有很多東西去維護自己.
但是摩西在這裡並沒有這樣去做.
我們在第三節.
我們見到摩西他謙卑的特質.
其實和米尼暗和雅倫成為一個強烈的對比.
其實米尼暗和雅倫的背後.
其實都有一些驕傲.
因為其實當人有嫉妒有妒忌的時候.
其實都是來自自己的驕傲.
我要比這個人好.
這個人不准比我好.
有種令我要超越他的感.
要超越他.
或者你很想得到這個名銜.
這個人不准得到.
是我的.

$^{281}$其實這個社會.
是不是都不停地鼓吹這件事呢.
我們要去爭贏一些東西.
我們要爭贏一個partnership.
我們要爭贏一個deal.
所以如果我們被這個世界的風氣帶著走的話.
其實都會影響到那份兄弟姐妹的感情.
而當有這些超越感.
在我們的群體當中的時候.
其實很容易會造成分裂和分化.
而這個絕對不是上帝喜悅的事情.
我們繼續再看第四節第六節的時候.
我們在第四節我們見到摩西.
他其實都沒有回應的時間.
或者其實也可以解讀為.
他根本都不需要回應.
因為今次耶和華神親自出手.
神就對他們三個弟弟,哥哥,姐姐去講.
叫他們去到會幕那裡.
我們見到耶和華在雲柱的降臨.
雲柱象徵著神的同在.
上帝特別指著要召阿倫和米尼暗.
即是哥哥和姐姐出來.
神並且在那裡要他們聽他講.
我們再看, 上帝在這裡表明.
他一定會在異象向先知顯現.
上帝在說他和先知的溝通方式.
和他和他僕人摩西溝通的方式.
是有些不同的.
我們首先看到, 他說他和先知.
普通的先知就會透過異象和先知顯現.
在夢中向他們說話.
神首先在這裡指出.
他和一般的先知他們交流的方式.
那種方式就是神會在異象當中.
讓先知去認識他.
第二就是在夢中.
先知會聽到神向他們說話.
但我們在第七節可以清楚見到.
強調摩西的獨特性.

$^{321}$摩西和神的連結方法很不同.
第七節描述摩西.
他是在全以色列人當中最忠誠.
正因為摩西有這個特別的身份.
所以神不是用他和一般先知溝通的方式.
上帝是面對面和他說話.
這樣的說話方式當然比在夢中說話更加清晰.
面對面和你做一個夢.
去得到一個訊息.
我們可能也要猜一下那個夢是否真的.
是否在說甚麼.
但這裡代表著摩西和上帝的關係的親密度.
親密度就像你和家人面對面的交流.
並且神說他清清楚楚.
他說他不用謎語和摩西溝通.
代表摩西不需要猜測.
可以直接聽到神和他說的話.
上帝也說摩西會見到神的形象.
所以總括而言.
上帝和摩西的交流和溝通方式.
其實遠勝一般的先知和一般神的代言人.
所以神就問米利安和亞倫.
你們為甚麼要批評他的僕人摩西而不感到懼怕呢.
我們再看下去的時候.
我們見到上帝就向亞倫和米利安發怒氣.
上帝因為米利安和亞倫對他僕人摩西的批評.
而向他們發怒.
不知道大家認為上帝發怒的時候是怎樣呢.
我們在這裡見到上帝發怒的時候.
是充滿著威嚴和充滿著權柄.
上帝他刑罰了米利安.
當亞倫轉身去望米利安的時候.
他看到米利安怎樣呢.
是長了麻風.
經文描述到米利安的麻風是像雪一樣的白.
其實是在形容那個形狀.
那些麻風好像雪的形狀.
好像碎片一樣.
所以我們所說所做的事.
其實是會帶著後果的.

$^{361}$我們在十一章見到.
上一次我們看到.
以色列人他們埋怨的後果.
就是上帝用火燒毀營的外圍.
因為他們用惡言在發怨言.
所以我們真的要謹慎我們心裡面所想的事情.
如果我們沒有好好梳理心裡面的感受.
或者所想的意念.
其實可能會變成我們口裡面.
不知不覺口裡面說出來的事情.
所以我們真的很需要每天.
要有一些安靜時間.
我們回到神的面前.
去安靜去梳理.
究竟心裡面是有一些不舒服.
或者有一些地方我們要跟神去說.
我們不要讓忙碌.
不要讓我們生活的忙碌.
輕輕地帶走這些問題.
我們真的要回到上帝的面前去審測.
去反思和心思意念上.
被神的話語去更新.
願上帝的話語不斷地塗鑿我們的生命.
免得我們說了一些不合神心意的說話的事.
因為我們所說所做.
上帝是看著知道.
祂亦都聽到.
於是我們見到亞倫和摩西的對話.
亞倫在對話當中其實他都很清楚.
知道他和米利安對摩西的批評其實是不對的.
因為他說他犯了罪.
亞倫稱呼摩西為我主.
表示他對摩西的尊重.
並且為他和米利安的投訴去道歉.
亞倫叫摩西不要將這些罪加在他們身上.
因為他見到米利安身上的麻風.
在侵蝕他的身體.
相信我們有看福音書裡面的時候.
我們都有見到耶穌醫治麻風病人的情況.
其實麻風病是很辛苦的.

$^{401}$麻風病人其實是那些邊緣化的人.
沒有人想走過去因為那時候會傳染.
而且那些麻風其實會一直去侵蝕你的身體.
甚至如果嚴重的話其實是會死的.
於是我們見到摩西的回應是怎樣呢.
他有沒有說哥哥姐姐你這樣說我是罪有應得的.
摩西沒有.
摩西一句疑問都沒有.
他就立刻哀求上帝求神去醫治姐姐米利安.
其實摩西在整個段落裡面.
在第十三節其實他是第一次說話.
就是求神去醫治他姐姐.
所以我們見到剛才都有提及過.
在《出拉急》的描述裡面.
當年米利安姐姐就拯救了摩西.
他用聰明的策略讓摩西可以在一個安全的環境裡成長.
到現在摩西總算回報了姐姐求神去醫治他.
其實難道兄弟姐妹的關係不是這麼微妙嗎.
當日那個姐姐去幫弟弟妹妹照顧.
同樣地今日可能弟弟妹妹也會去關懷哥哥和姐姐.
於是神要管教米利安.
要將他在營外關鎖七日.
我們在民數記的五章二至三節看到.
其實耶和華去吩咐摩西.
你要吩咐以色列人擺一切患麻風的,患柳症的,和因屍體而不潔症的.
都送到營外去.
無論男女你都要把他們送到營外.
免得他們玷污了他們的營.
這是我住在他們中間的地方.
所以我們看到米利安他要潔症之後.
才能夠帶他回到營裡.
在這七日裡面.
百姓都因為米利安他犯了的罪不能夠起程.
所以他們要等到米利安回來.
他們才可以起程.
所以看到領袖的行為.
都會對屬靈群體造成影響.
並且他們本應要繼續進入應許地的行程.
亦都延誤了七日.
去到總結.

$^{441}$我們看到米利安他因為比較.
在他的心裡面發芽.
米利安和雅倫.
他投訴自己.
不是都能夠像摩西一樣.
可以和上帝有那種親密的關係.
可以和他去說話嗎.
其實比較是一種不安全感.
是覺得自己不夠好.
所以要透過和別人去比較.
去justify自己一些東西.
但是米利安他在這裡.
他比較的心態.
就導致了他對摩西的批評.
而這個是神不喜悅的.
並且當我們比較.
作出投訴批評之後.
我們會見到神一定會管教他的子民.
而第三我們見到.
其實今日究竟神.
是怎樣去看雅倫和米利安.
即哥哥和姐姐呢.
其實他們在神的眼裡面.
又是怎樣呢.
其實上帝都很喜悅這兩位的兄子.
因為他們都在做上帝的聖功.
他們都在服侍神.
所以當我們有比較的意念的時候.
我們真的要專注在上帝的身上.
問上帝其實是怎樣去看我們呢.
而我們會發現.
其實神都是一樣的那麼愛我們.
因為我們見到.
其實上帝他愛的屬性是不會改變.
如果我們真的專注在上帝的特質.
他的愛裡面.
其實是能夠消除了比較的心.
我們真的要有感恩的心.
因為所有的恩賜財幹都是從神那裡來的.
如果我們知道.

$^{481}$所有的能力都是從神那裡來的話.
那我們為甚麼要比較呢.
因為其實所有東西都是從神那裡來.
其實我們是不需要和隔壁那個去比較.
因為一切都是出於上帝的手.
主人蘇利已經為我們去犧牲.
他的愛讓我們能夠有能力去愛身邊的人.
讓我們可以不需要透過比較去愛身邊的人.
今天透過這一段的兄子弟情.
我們又怎樣能夠回應.
我們怎樣去活出.
上帝想我們活出那份的兄弟姐妹的關係呢.
我們每一個都是得蒙上帝恩典的僕人.
今天我們應該怎樣去活出.
好像摩西那種謙卑的態度呢.
最後就送一個Love給大家.
讓我們都是.
首先Love L.
我們要去Love God.
我們要渴望和神有更深的關係.
其實甚麼是謙卑呢.
謙卑其實是渴望和神有更深的關係.
第二我們就是要去O.
我們要去遵守上帝的教導.
當我們的心時常去閱讀上帝的話語的時候.
我們會自然地很想去做上帝喜悅的事情.
然後V我們就看到我們要Value.
上帝擺在我們生命裡面的關係.
我相信每一個的關係.
在我們生命裡面都不是偶然的.
上帝擺一些的關係.
無論是有血緣還是沒有血緣的弟兄姊妹.
其實都是想要我們去珍惜.
透過這些關係我們彼此的相愛.
彼此的互補不足.
彼此的鼓勵.
彼此的同行扶持.
最後一樣就是E.
我們要活出愛.
我們要Exercise這個Love.

$^{521}$並且當我們真的活出了這份愛.
我們不只是看口說.
我們活出這份愛的時候.
是因為上帝給予我們愛的動力.
所以最後我們要將榮耀歸給神.
我們要高舉上帝的名.
或許我們今天覺得兄弟姐妹的關係.
可能仍然有些地方很難去復和.
可能仍然有張力在裡面.
但盼望我們都一同聚焦在上帝的身上.
我們謙卑地給上帝去教導,去塗鑿我們的生命.
讓我們能夠活出愛.
化解彼此的仇恨和嫉妒和比較.
用上帝看待我們的眼界.
去看待我們自己.
看待得合二.
讓我們同心吐告.
親愛的天父.
教導我們比較發牙到投訴紛爭的結果.
求主你幫助我們.
縱然我們在不同的關係裡面.
可能仍然有張力.
可能仍然有我們要學習的功課.
但我們深信當我們信靠主你的時候.
你一定會賜我們力量去面對.
並且求主賜我們更多的愛心.
因為每一個關係都是上帝賜給我們.
讓我們好好珍惜這些關係.
讓我們在地上都能夠活出上帝你的愛.
在我們身邊的人.
我們這樣同心吐告.
奉主名求.
\newpage



\section{民數記 12:1-16}
\label{sec:032234}
\textbf{嫉妒與毀謗的可惡}~~~~吳安強牧師
\newline
\newline
連結: \href{http://www.church.com.hk/acms/content.asp?site=cdc\&op=show\&type=product\&code=032234\&layout=sermon}{\texttt{032234}} ~~~~ 語音日期: 2026-05-31
\newline
\newline
\hyperref[sec:032233]{\small{< < < PREV SERMON < < <}}
~
\hyperref[sec:index_chronic]{\small{[返順時目]}}
~
\hyperref[sec:032235]{\small{> > > NEXT SERMON > > >}}
\newline
\newline
民數記 12:1-16
\newline
\begin{longtable}{cl}
\hline
\hline
章節 & 經文 (和合本修訂版)\\
\hline
12:1 & \begin{tabularx}{0.7\textwidth}{X} 摩西娶了古實女子為妻。米利暗和亞倫因他娶了古實女子就批評他， \end{tabularx} \\ \\ \relax
12:2 & \begin{tabularx}{0.7\textwidth}{X} 他們說：「難道耶和華只與摩西說話嗎？他不也與我們說話嗎？」耶和華聽見了。 \end{tabularx} \\ \\ \relax
12:3 & \begin{tabularx}{0.7\textwidth}{X} 摩西為人極其謙和，勝過地面上的任何人。 \end{tabularx} \\ \\ \relax
12:4 & \begin{tabularx}{0.7\textwidth}{X} 忽然，耶和華對摩西、亞倫和米利暗說：「你們三個人都出來，到會幕這裡。」他們三個人就出來了。 \end{tabularx} \\ \\ \relax
12:5 & \begin{tabularx}{0.7\textwidth}{X} 耶和華在雲柱中降臨，停在會幕門口，叫亞倫和米利暗。二人就出來， \end{tabularx} \\ \\ \relax
12:6 & \begin{tabularx}{0.7\textwidth}{X} 耶和華說：「你們要聽我的話：你們中間若有先知，我－耶和華必在異象中向他顯現，在夢中與他說話； \end{tabularx} \\ \\ \relax
12:7 & \begin{tabularx}{0.7\textwidth}{X} 但我的僕人摩西不是這樣，他在我全家是盡忠的。 \end{tabularx} \\ \\ \relax
12:8 & \begin{tabularx}{0.7\textwidth}{X} 我與他面對面說話，清清楚楚，不用謎語，他甚至看見我的形像。你們為何批評我的僕人摩西而不懼怕呢？」 \end{tabularx} \\ \\ \relax
12:9 & \begin{tabularx}{0.7\textwidth}{X} 耶和華向他們怒氣發作，就離開了。 \end{tabularx} \\ \\ \relax
12:10 & \begin{tabularx}{0.7\textwidth}{X} 當雲彩從帳幕上離開時，看哪，米利暗長了痲瘋，像雪那麼白。亞倫轉向米利暗，看哪，她長了痲瘋。 \end{tabularx} \\ \\ \relax
12:11 & \begin{tabularx}{0.7\textwidth}{X} 亞倫對摩西說：「我主啊，求你不要因我們愚昧，因我們犯罪，就將這罪加在我們身上。 \end{tabularx} \\ \\ \relax
12:12 & \begin{tabularx}{0.7\textwidth}{X} 求你不要使她像那一出母腹、肉已侵蝕了一半的死胎。」 \end{tabularx} \\ \\ \relax
12:13 & \begin{tabularx}{0.7\textwidth}{X} 於是摩西哀求耶和華說：「神啊，求你醫治她！」 \end{tabularx} \\ \\ \relax
12:14 & \begin{tabularx}{0.7\textwidth}{X} 耶和華對摩西說：「她父親若吐唾沫在她臉上，她豈不蒙羞七天嗎？現在要把她隔離在營外七天，然後才領她回來。」 \end{tabularx} \\ \\ \relax
12:15 & \begin{tabularx}{0.7\textwidth}{X} 於是米利暗被隔離在營外七天；百姓沒有起程，直等到米利暗回來。 \end{tabularx} \\ \\ \relax
12:16 & \begin{tabularx}{0.7\textwidth}{X} 以後百姓從哈洗錄起行，來到巴蘭的曠野安營。 \end{tabularx} \\ \\
[1ex]
\hline
\hline
\end{longtable}
$^{1}$聽姐妹早晨.
有人見到這個題目的時候問我.
牧師你是不是打錯字.
他說是不是誹謗.
應該是窒道要誹謗的可怕.
為什麼你寫可惡呢.
我沒有打錯字的.
一會兒我會解釋為什麼.
我覺得是可惡多過可怕.
窒道呢.
心理學家告訴我.
是一種很複雜的情緒.
通常的來源是在我們心裡面.
和其他人的比較.
和競爭的心態.
但是在這種情緒裡面.
又夾雜著一些羨慕.
焦慮.
憤怒.
和自卑的情緒.
我覺得窒道往往發生在我們和周遭親近.
或者條件很相似的人身上.
如果差很遠我們是沒有窒道的.
因為不同的.
很多時候很相似很相近.
於是窒道就會發生.
窒道在希伯來文的意思是「忌邪」.
窒道可以講翻譯.
或者也可以譯作「熱心」和「激動」.
《聖經》當中其實也有提到神也有窒道.
所以有時候這個字不一定完全是百分百的負面.
當是和合本譯作「忌邪」多一點.
當譯到神的那種窒道的心的時候.
不過神的窒道和人不同.
神的忌邪和忌惡是不同的.
不過神的窒道和人不同 神的忌邪的神.
祂的窒道不是出於祂的自私和惡意.
就是說神的窒道不是只為了祂自己的好處.
而是出於祂聖潔的屬性.
目的是引導人走向正確的道路.

$^{41}$遠離邪惡和偏離真理的事物.
聖經說神是一個忌邪的神.
神是一個窒道的神.
就是當祂見到人犯罪離開祂的時候.
做一些事情跟祂背道而馳.
祂就會覺得很憤怒,很不開心.
然後祂就會想辦法將你帶回頭.
但是人的窒道不同.
祂讓你看到因為別人的成功.
而心裡面產生一種焦慮和不安的情緒.
而這種情緒佔滿了我們整個的心.
讓我們越看到那個人的成功.
我們就越不開心.
越看到那個人順利.
而我們心裡面就越覺得不開心.
所以窒道就好像一個人.
不斷在你耳邊說這句話.
「為什麼是他?為什麼不是我?」.
不知道大家有沒有說過這句話.
升職的時候可能會問.
「為什麼是他?不是我?」.
拿到一百分,我九十九分.
我真的有一次在大學說過.
我有一個同學很厲害的.
有一次我跟他差不多.
我拿回來之後我很高分.
我很開心.
我就問他,他比我多一分.
我打算跟他說.
「這次我終於贏你了」.
而我心裡面就說.
「為什麼是你?不是我?」.
「又是我?」.
「為什麼是這樣?」.
這個說話其實很多時候.
都在我們生命當中出現.
今天我們看的那段經文.
米利暗和亞倫.
一個做姐姐,一個做哥哥.
但是他們兩個人的說話.

$^{81}$正正是將以上窒道的意思.
完全的表露了出來.
文俗記十二章第二節.
「難道耶和華只與摩西說話嗎?」.
「他不也與我們說話嗎?」.
就是米利暗和亞倫說.
「神只是跟摩西說話嗎?」.
「都不跟我們說話嗎?」.
其實他們知道.
神都跟他們說話.
因為米利暗是一個女先知.
亞倫是大祭司.
其實神都透過夷蒙.
透過很多方式來跟他們說話.
所以他們的意思是.
既然我們都聽到神的說話.
為什麼摩西是領袖.
為什麼祂跟我們不同.
為什麼祂這麼特別.
所以米利暗和亞倫所說的.
正正是為什麼是祂.
為什麼不是我們.
當然了.
三兄弟姐妹就算了.
三個人自己在家裡吵就算了.
不要理他們.
但是我們明白.
這不是一個簡單的家庭不和事件.
因為是兩位當時以色列當中.
很重要的宗教領袖.
就是大祭司和先知.
對於神給摩西特別有的權柄.
他們心裡面感到不滿.
他們覺得為什麼神.
只是跟摩西這麼親近.
親近我們.
為什麼我們都不可以跟神這麼親近.
這件事件.
這個妒忌的起因.
可能.

$^{121}$我們不百分百肯定.
可能是追溯到上一章的經文所說的事件.
就是摩西根據耶和華的吩咐.
將權力分配給七十個長老.
我想大家記得了.
他將權力分配給七十個長老.
可能在這樣的安排下.
使亞倫和米利欽的權力似乎降低了.
因為以往可能摩西處理不到的事.
就是米利欽和亞倫.
但是現在有七十個長老.
他們不需要再參與當中很多治理的工作.
所以對於亞倫和米利欽來說.
他們覺得.
當摩西將這七十個長老.
將權力分配給他們之後.
好像在某個程度裡面.
使得米利欽和亞倫的權力減退了.
就因為權力慾的追求.
引發了這兩位同時都是宗教裡面.
很重要的領袖妒忌的心.
於是米利欽和亞倫.
就起了奪取領袖位置.
或者起碼他們覺得.
我們應該和摩西有同一的位置.
同一個權柄.
同一個話事權.
米利欽和亞倫因為妒忌.
摩西被主使用而重傷他.
讓我們想到聖經真言十四章三十節.
裡面這句話是這麼說的.
心中安靜的是肉體的生命.
疾到是骨中的瘤癩.
那個意思是什麼?.
那個意思就是說.
這個經文告訴我.
提出一個警告.
他說.
疾到這種情緒.
帶來非常嚴峻的負面影響.

$^{161}$這種藏在我們心裡面深處的罪惡.
不知覺間慢慢會從.
隱藏內心深處的動機.
發展成為外在擴散的行動.
就是說原來這種動機.
不會長久留在我們心裡面.
讓它在我們心裡面慢慢累積.
然後我們就會出了行動.
去傷害我們所疾到的人.
內心的瘤癩慢慢從內體的表現.
從人的行動當中呈現出來.
正如《箴言二十七章》第四節所說的.
憤怒為殘忍.
怒氣為狂瀾.
兩者都很可怕.
但他說唯有疾到誰能抵得住呢?.
這個經文的意思就是說.
憤怒和怒氣.
好像被巨大的風所吹動的海浪一樣.
極具破壞性.
我想過去我們見到了.
打風那些浪.
你的餐廳很好的.
玻璃都吹爛了.
所有的東西都吹爛了.
憤怒和怒氣.
就好像被巨大的風浪所吹動的海浪一樣.
很具破壞性.
不過.
疾到比憤怒和怒氣更可怕.
因為兩者可以抵擋.
但是.
疾到卻是無可抵擋的破壞力量.
米利暗當年作為姐姐.
當弟弟到了三個月就收不起來了.
媽媽把他放在一個箱子裡放在爐的當中.
想看看弟弟究竟怎麼樣.
於是這個姐姐很關心這個小寶寶.
於是她就走去看.
躲起來看看.

$^{201}$究竟我這個弟弟怎麼樣呢?.
當年米利暗很疼愛這個弟弟.
很關心他.
想方設法來保護他.
但是.
現在她無法容忍這個弟弟成為神重用的領袖.
所以聖經說我們要保守我們的心.
勝過保守一切.
很疼愛這個弟弟.
很關心他.
來到這一刻.
她是不能夠接受弟弟成為神重用的領袖.
於是.
她就梳擺亞倫一起的眾雙摩西.
一起的偉榜.
去破壞了團隊的和諧.
這就是為什麼我在主題裡面.
不說疾到和偉榜可怕.
可怕的事情我們會避開的.
我們不會走下去的.
你見到一條飯鏟頭.
你不會過去.
嘩!挺好玩的.
不會的.
因為很可怕.
但是你見到那條蛇沒有毒.
你又喜歡養蛇.
香港有些人養蛇的.
那你會跟他玩的.
因為你知道他不可怕.
可惡的地方就是米利暗不知道.
他不知道他心裡面的疾到已經出來了.
他只覺得為什麼他不是我.
為什麼摩西只有我不行.
這個就是我覺得.
可惡的地方就是我們不在意.
我們不為意.
但是這個疾到的心就進了我們的裡面.
所以他沒辦法容忍這個弟弟.
成為神重的領袖.

$^{241}$於是梳擺亞倫.
結果.
摩西,亞倫,米利暗.
三個非常和諧的團隊.
為神做著非常美好事情的時候.
這個和諧的團隊被破壞.
而從歷史發展來說.
我們見到之後發生一件很嚴重的事情.
不知道是不是真的因為這件事而發生的.
就是跟著米利暗和亞倫對摩西的疾到毀謗之後.
就出現了可拉一黨的人的叛變.
可拉因為疾到摩西亞倫的屬靈權柄.
營私結黨.
煽動群眾反叛.
這個是後面文述的第十六章.
到時候我相信同工會詳細跟大家看.
所以我們見到疾到是有傳染性的.
剛才牧師祈禱的時候祈禱剛果.
那個伊波拉病毒的傳染性很強.
疾到一樣.
他說當這種情況一旦在屬靈群體當中出現.
就會好像猩猩之火一樣傳播得極廣.
就好像傳染病一樣.
會帶來非常大的破壞和傷害.
但是問題就是我們不懷疑.
神已經告訴我們米利暗和亞倫.
先從摩西的個人生活開始批評他.
這個說話好像很簡單.
十二章一節.
他說因他娶了古實女子就批評他.
古實如果用現在的字眼說非洲.
不過這裡應該不代表摩西娶了一個非洲的太太.
因為我們知道摩西年輕的時候為了逃避法老的追殺.
曾經走到米甸的地方.
神已經告訴我們.
他就娶了米甸祭司的女兒.
叫西波拉為妻.
我們看到春夏日記第二章.
但是當摩西回到埃及的時候.
當他在神的命令下.

$^{281}$回去要帶領以色列人離開埃及的時刻.
西波拉沒有和他一起回去.
但是我們看到初初本來是一起去的.
但是在中途西波拉就想起.
長子未行割禮.
於是就幫長子在路上行了割禮.
很多人就相信.
因為長子甲訊行了割禮.
不能夠繼續行程.
於是西波拉就留下來.
繼續去幫長子養傷.
以及預養他的小兒子以利爾.
直到以色列在西乃山安營.
摩西的岳父知道摩西已經離開埃及.
在西乃山.
於是就帶著西波拉和兩個兒子.
去到西乃山和摩西團聚.
在這個時候.
有可能米利和亞倫才知道.
摩西原來娶了一個米甸女子為妻.
於是他們就趁這個機會來圍捕摩西.
當然聖經這裡說摩西娶了古實的女子為妻.
簡家告訴我們有兩個可能.
一是西波拉已經死了.
所以摩西再娶了一個古實女子為妻.
古實也有機會指揮塞俄比亞.
即是他另娶了一個女士.
第二是西波拉仍然是摩西的妻子.
古實女子在貶義的稱呼.
說他的太太皮膚很幽黑.
黑得好像一個非洲人.
兩者都是批評摩西娶外邦女子為妻.
即是他覺得你為何不是娶猶太人?為甚麼?.
但我們要留意的是當米利亞這樣去.
去毀謗摩西的時候.
我們看這聖經要留意的是.
這個時候神的律法未頒布.
而根據《創世記》二十五章一到二節.
我們知道其實最重要的就是米甸族.
其實是亞伯拉罕的後裔.

$^{321}$他不是外邦人.
其實他是亞伯拉罕的後裔.
都是希伯來人.
所以摩西娶的是亞伯拉罕的後裔.
娶的是希伯來人.
這個民族.
所以基本上是沒有問題.
只要摩西娶的不是迦南族的女子就可以.
作為神的大祭司和阿倫.
和作為神的女仙子米利亞.
借著摩西他們覺得他娶了外邦女子為妻.
這件事來毀謗摩西.
但他們毀謗不是為了糾正摩西的錯誤.
而是指出摩西為何娶這個女子為妻.
不是想摩西知道他做錯事.
而是出於嫉妒摩西的屬靈權柄.
他們想通過摩西的責備.
表達他們自己作為以色列的領袖.
擁有和摩西同等的權柄.
所以為什麼第二節說難道耶和華至於摩西說話嗎?.
他不也與我們說話嗎?.
米利亞作為女仙子.
他真的擔任過耶和華的代言人.
在《春夏六記》十五章二十到二十一節.
神感動他.
他就唱詩歌讚美神.
而阿倫作為大祭司.
也的確在以色列人的生活當中.
他有判別神的心意.
在《文述》第二章我們見到.
他們也與摩西一同親自聽過神的聲音.
就是神告訴他們應該如何安寧處理事情.
是的.
這兩個人的確.
耶和華不僅與摩西說話.
也與他們說話.
就是因為這個緣故.
米利亞和阿倫認為既然神不僅與摩西說話.
也與他們兩人說話的時候.
為什麼摩西可以擁有與他們不一樣.

$^{361}$特別的屬神權柄.
為什麼摩西的領導權在他們之上?.
所以《文述》十二章第二節的說話.
才是米利亞和阿倫.
想透過批評摩西娶古實女為妻的真正目的.
但是他們很想抬高他們與摩西同等地位的領導權.
當然同時超越那七十個長老.
做回最高,最頂尖的三個領袖.
這不僅是人類負面的一些情感,一些感受.
更加是一種世俗的價值觀.
世俗的人就覺得.
我都聰明.
你也聰明.
為什麼是你才做那個位置?為什麼不是我?.
但是在屬靈的角度來說.
他不僅是一種世俗的價值觀.
其實也反映出我們靈性低落的狀態.
這裡顯示了米利亞.
當他有這樣的心態的時候.
他與神的關係其實可能已經差了一些的情況.
因為當他與神的關係好.
他很明白.
不是高低.
而是在帶領整個意識形態的群體當中.
神是將對的人放在對的位置裡面.
一起配搭好.
你想想一個領袖,一個先知,一個大祭司.
其實各司其職做好那樣的事情.
就能夠將整個意識形態的群體帶進神的角度裡面.
但是他們不是這樣看.
他們看的是人的價值觀.
就是為什麼我低於你.
所以習道通常與自私與貪婪相關.
米利亞不甘心只是做女先知.
又慫恿亞倫不要甘於只是做大祭司.
因為這些都是宗教上的遺訓.
他們還在政治上有領導權.
這是摩西猶的.
當時帶領以色列人的時候.
這種貪婪的心態.

$^{401}$使得他們不能夠學習到為別人的幸福而感恩.
也都難以對其他人的成功有喜悅和感受.
如果我們心裡面存著這種感恩的感覺的時候.
我們見到人成功,見到人順暢的時候.
我們不會覺得開心.
因為我心裡面覺得我都應該有這些東西.
習道的情緒也使人沒辦法真正專注於神的關係.
我們的眼會看著我們習道的人.
我們會羨慕他得到外在的祝福.
但是他有沒有注意到原來我自己.
其實都有很多祝福在我生命裡面.
但是我不記得去數算.
我只是不滿意其他人得到.
這樣的人或者可以跟哀哭的人同哭.
即是那個人很懶的時候.
但是他沒辦法跟喜怒的人一同的喜怒.
米利亞和亞倫已經忘記了.
神是因為摩西的緣故才興起他們的.
因為神興起摩西的時候.
摩西就說我不行,很多事情都不行.
神說不要緊,你有亞倫.
你有家裡的哥哥姐姐可以幫助你的.
就是說在神的計劃當中.
神一直將他們擺在一個位置.
就是成為摩西的協助者.
大家千萬不要看協助者不重要.
沒有協助者其實那個主要角色都做不完他的事情.
都沒辦法成功.
協助者和帶領外人其實是一起去成就那件事情.
他們忘記了神將他們擺的是什麼角色.
不是平起平坐的領導者.
但是米利亞和亞倫不願意做幫助者的角色.
然後聖經記載了兩句很奇怪的說話.
在文宿記第十二章第二節和第三節.
就說耶和華聽見了.
然後就說摩西為人極其謙和.
勝過地面上的任何人.
所以我們說一個人.
如果我跟大家說吾安強牧師是全世界最謙卑的牧師.
我說的吾安強牧師.

$^{441}$嘩 你搞錯了 你自己這樣說.
有人說摩西為什麼突然間寫一句這樣的東西.
還是最謙和的.
我們看看兩句很奇怪的說話.
摩西為人極其謙和.
勝過地面上任何的人.
表達的是摩西面對這樣的指責.
他沒有為自己辯護.
也沒有好像先前一樣.
當他被人埋怨誤會這樣向神申訴.
他沒有.
謙和的希伯來文在這裡的反應是謙卑.
但是在其他經文.
他是指那些貧窮困乏.
社會地位低微.
受苦或者羞辱的人.
所以這裡其實除了說摩西是一個謙和的人.
其實也指摩西在這個時候是一個舉足輕重的領袖.
但是他因為百姓而受苦.
摩西沒有得到過他應該受的尊重.
反而不斷地面對百姓的那種埋怨.
帶來他很多的苦澀.
剛剛上一章.
在塔比拉先因為百姓發怨言.
使耶和華聽見就怒氣發作.
使火燒到他們中間.
不過神很大恩典.
燒到營的邊界就沒有燒下來.
那是一個警告.
然後.
在基博羅哈他雅.
百姓有埋怨沒有東西吃.
上次林牧師講到的時候提過了.
他們有埋怨沒有東西吃.
最後神用最重的災殃擊殺他們.
這兩件事使得摩西作為領袖.
受到非常沉重的壓力和打擊.
他愛這班百姓.
這班百姓他一而再再而.
其實不是叛逆他.

$^{481}$是叛逆神.
這就是帶來神的懲罰.
摩西覺得很難過.
從基博羅哈他雅去到哈洗祿.
有人計算過大約是四十八公里.
就是一天裡面可以走完的路程.
就是說在這麼短的路程裡面.
只是一天而已.
摩西曾經受到的壓力和打擊還沒有復原.
一天之內他可能還在沉澱整個過程.
如何去處理.
很快地他又要面對自己的姐姐和哥哥的惡待和羞辱.
所以摩西在這個時候.
在哈洗祿.
他受到前所未有的打擊和焦慮.
他沒有能力為自己辯護.
所以為什麼聖經用了一句話.
「耶和華聽見了」.
當然世人所說的「耶和華聽見了」.
為什麼這裡特別記載「耶和華聽見了」呢?.
這句話的意思其實不是說.
耶和華只是聽到米利安和亞倫的誹謗.
這句話學者告訴我們有一個很重要的意思.
就是說神決定他要介入去處理這件事情.
神聽到很多人說的話.
但是神不一定介入.
但是這裡很清楚.
神聽到他要介入來處理這件事情.
所有的說話神都聽見.
但是不一定介入.
神這次介入.
作摩西的辯護.
在一方面摩西無能為力去做辯護.
所以而這個誹謗是一個非常之嚴重的誹謗.
因為摩西是神自己設立.
要帶領以色列人進迦南地的領袖.
如果在這個情況裡面.
摩西被打倒了.
又或者這個領袖祭司先知的團隊.
因為這件事而拆夥.

$^{521}$而很多的嚴規.
很多的問題就會產生.
成就不了神的工作.
摩西和艾利暗暗都肩負著帶領以色列人進來迦南地.
所以他們三個是不是同心很重要.
於是神親自在會幕當中審斷這件事.
民俗的十字掛神忽然在他們當中出現.
忽然這個詞是說.
然後說怒氣發作.
很快就介入這件事情當中.
因為神自己都知道.
這件事不能夠蔓延出去.
若說蔓延出去.
就會成為整個以色列群體帶來更大嚴重危機.
即是說當三個領袖各有不同的看法.
他們就可能有不同的人跟著他們.
我們見到後來可拉一黨也出了這個問題.
所以神很快就介入這件事.
神親說將米利,阿倫和摩西叫出來.
要他們去到會幕那裡.
因為會幕是神與人相見的地方.
亦都是神彰顯祂自己旨意的地方.
亦都是代表神臨在和統治以色列民的地方.
會幕是當時代表神的臨在.
一個很重要的地方.
神就好像在法庭當中.
為米利,阿倫和摩西的批評作出一個公義的裁決.
神沒有否定米利,阿倫.
有關祂作為先知.
可以從異象當中在夢裡得到神的啟示.
所以民數記十二章第六節是這樣說.
「你們要聽我的話.
你們中間若有先知.
我耶和必在異象中向祂顯現.
在夢中與祂說話」.
意思就是說.
是的.
米利,阿倫你說得沒錯.
我不是只跟摩西說話.
我也跟你說話.

$^{561}$不過跟著神就說.
「但摩西與你們不是同一個層次」.
神說摩西是祂親自揀選的僕人.
如果用一個字眼.
就是神說摩西在我的家當中.
祂是管家.
祂也是僕人.
和你們一樣.
不過我給祂一個位分.
祂是一個管家.
祂要帶領著整個的情況.
祂作為帶領整個以色列家.
離開埃及進入迦納的代表人物.
祂代表的不是祂的榮耀.
祂代表的是神自己.
要透過祂的帶領.
讓整個以色列民看到神的榮耀.
所以神說摩西與祂見面的時候.
不需要透過任何媒介.
不需要義目.
不需要異象.
而是好像朋友一樣.
文述第十二章七到八節這樣說.
「但我的僕人摩西不是這樣.
他在我全家是盡忠的.
我與他面對面說話清清楚楚.
不用謎語.
他甚至看見我的形象」.
當然這裡不是說面對面.
你給他看看背部.
他可以看見我的形象.
這正是神在西乃山的時候.
指著摩西所說的話.
在《春艾求其》三十三章十一節.
要說與摩西面對面說話.
好像人與朋友說話一樣.
神說摩西是直接領受祂的說話.
不需要透過異象.
不需要透過夢境.
如果用一個更直接的方式表達.

$^{601}$神說摩西是我的入室弟子.
我把全套功夫轉給他.
你們只不過是我眾多弟子中的一個.
都是弟子.
所以不需要嫉妒.
因此對於米利任和亞倫.
對摩西的嫉妒和毀謗.
神的判詞是這樣說的.
「你們為何批評我的僕人摩西而不懼怕?」.
這句話的意思是.
你們不是攻擊毀謗摩西.
你們是出言攻擊我.
神說摩西在這個位置.
是我把他放在這裡.
你們攻擊他.
其實即是說.
神說你是不滿意我.
你是不滿意我所做的事.
所以我們說.
嫉妒和毀謗的可惡.
就是我們不留意.
原來當我們嫉妒一個人.
說他一些不好的事.
其實.
原來按照聖經所說.
我們是在控訴神.
覺得他不公.
覺得他做得不好.
所以神這個判語很嚴重.
隨著判語而來的是一個警告性的懲罰.
聖經說在文述第十二章第十字.
「米利任長了大麻風」.
像雪那樣白.
當亞倫見到米利任長了大麻風.
他可能想起.
神為了向埃及法老展示自己的大能.
而在摩西手上施行了相同的神蹟.
大家還記得嗎?.
摩西把手伸出來.
任長了大麻風.

$^{641}$所以作為大祭司.
這個時候他就知道.
神是要將他的大能.
向他連這三個神所擺的領袖展示出來.
作為大祭司亞倫.
這個時候可能應該被提醒.
他是應該為以色列百姓贖罪的.
但是他竟然認為.
他應該和摩西有同等的地位.
所以這一刻.
神透過米利任的大麻風提醒了亞倫.
他犯了罪.
他需要去倚靠摩西向神求憐憫和恩典.
於是這位曾經一度支持米利任的亞倫.
就來到摩西面前向他認罪求寬恕.
文述第十二章十一到十二字是這樣說的.
「我主啊,求你不要因我們愚昧.
因我們犯罪.
就將這罪加在我們身上.
求你不要使他像那一出無福.
體以侵蝕了一半的死胎」.
即是他來到摩西面前向摩西說.
你寬恕他.
因為米利任犯了罪.
神用麻風作為警告性的懲罰.
使米利任和亞倫重新認識.
和確定摩西在整個以色列群體當中.
獨特的領導地位.
獨特的領袖.
我就想一下.
如果我是摩西.
純粹從人性來說.
前天才被百姓激撐死.
今天你們不支持我算了.
還來攻擊我.
我就會這樣想.
我暫時不為你祈禱.
等你捱幾天.
使善道氣.
然後再跟別人說.

$^{681}$神啊,你原諒他吧.
面對哥哥亞倫的認罪.
懇求面對姐姐的病患.
摩西沒有絲毫的拖延.
不容易的,如果我們從人性去想.
我們就看到摩西其實很明白神的心.
他沒有拖延.
他沒有覺得我做一些事.
讓姐姐受苦.
受苦幾天才印象深刻.
再疾度四星期.
哎呀,不要來了.
上次因為疾度出了這些事.
摩西不計前嫌.
他馬上就為那個曾經因為妒忌.
而毀謗他的哥哥和姐姐向神哀求.
他說神啊,求你醫治他.
我相信民俗的十二章十三節.
這句很簡短的說話.
正正是道出了摩西為人極其謙和的名稱.
他不介意哥哥姐姐所說的話.
摩西不計前嫌.
不過神的懲罰定有界限.
即是說神的懲罰是有它的界限.
雖然摩西已經去求了.
但神說不行.
因為米利暗的懲罰太快結束.
就會淡化了毀謗事件的嚴重性.
神將這個懲罰保留一個星期.
米利暗就要離開營地.
住在其他的地方.
在那裡他有機會反思.
同樣這個禮拜的時間.
也都讓所有的以色列人.
作為整個一個屬靈的群體.
都要反思一件事情.
就是人如果反叛背叛神的計劃.
就要招致懲罰.
我相信不是在說誰是領袖的問題.
而是在神的計劃當中他們要信服.

$^{721}$否則就會招致懲罰.
以致他們學會在未來的日子.
在他們整個行程當中.
好好地按著神的心意向前行.
既然嫉妒和毀謗是神所不喜歡的.
但又很容易發生在屬神的群體裡面.
我們可以怎樣來應對呢?.
在上一章經文我們見到的.
就說摩西設立了七十個長老.
他們在會幕當中.
但有兩個沒有出現.
一個叫伊利達 一個叫米達.
六十八個去到會幕裡面就受感講預言.
但這兩個人沒有順從摩西的指示.
去到會幕集合.
但聖經就說.
縱然他們不在會幕裡面.
他們只在自己的帳幕當中.
他們都是受感說話.
跟其他長老一樣.
當時有一個年輕人.
是摩西的幫手 叫約書亞.
他見到這件事就忿忿不平.
他覺得這七十個長老不順從摩西的說話.
不去到會幕裡面.
不應該可以這樣做的.
沒有理由他們都可以受感說預言.
應該是被禁止的.
摩西是這樣回應約書亞.
他說:你為我的緣故疾到人嗎?.
惟願以和我的百姓都受感說話.
願以和我把丹麗靈降在他們身上.
這是摩西在神面前一個極其謙和的表徵.
他說:我有什麼特別?.
為什麼是我才可以講預言?.
甚至為什麼是這七十個長老?.
他說:所有的人.
他說:他但願整個以色列民族.
都能夠跟神有一個非常親密的關係.
以致神都可以使用他們.

$^{761}$能夠受感說預言.
所以摩西跟約書亞說.
不要疾到.
不要這樣誹謗人.
你要順從.
因為這是神自己給的恩告.
是神自己給這七十個長老.
一個很特別.
是印證他們是神自己所選召的領袖.
無論是面對輕者的誹謗.
或者其他人.
都能夠得到神的祠福上面.
我們見到摩西都有一個謙和的態度來面對.
他分明大量面對別人.
擁有好像他一樣的才幹和恩賜.
他也不計較那些背後暗箭傷他的人.
當我們看新約聖經的時候.
也從主耶穌面對法利賽人和民事.
因為疾度而產生的敵視.
而對耶穌惡言相向的時候.
我們見到.
耶穌也不把時間和精力.
花在跟他們的鬥爭當中.
耶穌核心的行動.
就是他要繼續執行天父的旨意.
成為他生活的核心.
成為他學習的指引.
今天在我們生命當中.
有時的確我們都面對別人的疾度誹謗.
但是聖經裡面.
無論從摩西從耶穌的生命當中.
我們見到他們不反駁.
不為這件事爭執.
他重要的就是.
我是否真的走在神的心意當中.
繼續好好地成就神要我所成就的事情.
疾度和誹謗的可惡.
在於我們一不小心.
就會令我們花上無比的心思.
在崇靈群體當中建立好的關係.

$^{801}$一刻間就被破壞盡盡.
真的,弟兄姊妹.
我覺得在崇靈群體當中.
我們在團契,在小組,在相交裡面.
我們真的花很多時間,心機.
才建立親密的關係.
但是一句的誹謗.
一個疾度的心.
就可以完全拆毀那個關係.
所以讓我們將多些的心思.
放在耶穌基督的身上.
讓我們多些為祂所成就的救贖而感恩.
因為我們自己.
你和我.
都是因為耶穌基督的救贖.
然後我們可以在這個地方.
一切的來敬拜.
讓我們多些為這件事感恩.
讓我們也多些為祂賜給我們.
有服侍的恩典而感恩.
我們在不同的崗位事奉.
正如我剛才所說.
神將我們在不同的崗位.
沒有了你的崗位.
神的國度就真的受到虧損.
所以我們一起來配搭.
也讓我們多些為身邊那些.
擁有比我更多,更大.
和更顯眼的恩賜的弟兄來感恩.
因為神興起他們.
又讓神的國度能夠有更好的發展.
當我們這樣做的時候.
才能夠讓基督的謙卑.
牢牢的住在我們的心裡面.
成為港族從撒旦而來的雙生子.
我真的覺得嫉妒和誹謗.
是雙生兄弟.
哪裡有嫉妒.
哪裡就會有誹謗.
所以《雅各書3章16》說.

$^{841}$有嫉妒何處就有困擾和紛爭.
所以求主憐憫我們.
用祂的柔和和謙卑.
去指釋在我們心靈當中.
躁動的嫉妒的心.
又求主幫助我們.
用祂的溫柔,恩賜的手.
去醫治如果我們曾經被嫉妒所傷的時候.
求神都憐憫醫治我們.
於是我們在神的國度裡面.
我們繼續前行.
繼續成就神的工作.
最後和大家分享網上見到的幾句說話.
我覺得是很好的.
作為我今天講道的結束.
在網上我看到.
如果你看到這四種徵兆.
可以讓你檢視自己是否落入一疾渡的網絡裡面.
第一.
我們刻意減少對他人的讚美.
明明想讚他.
但抓住不出聲.
平時見他做完事.
「做得不錯」.
「OK」.
第二.
我們沒辦法真心為朋友的成功獻上祝福和感恩.
你祈禱的時候.
你心裡面沒辦法為他感恩.
為什麼?.
我們要問.
如果你看到他真的成功.
為什麼沒辦法為他感恩?.
第三.
看到有些人出事.
我們會在內心開心地吶喊.
「他活該」.
都說他了.
看著他.
「只能掉下去了」.

$^{881}$第四.
我們更擅長批評而不是讚美.
求神幫助我們.
我們一起禱告.
傳領我們上帝.
我們真的帶著謙卑的心來到你的面前.
有時候當我們說我們帶著謙卑的心.
我們也要問問自己.
其實我是否驕傲?.
我是否看別人不順眼?.
或者好像聖經裡面.
米利暗和亞倫所說.
「為什麼是他而不是我?」.
神幫助我們.
因為耶穌基督.
你整個生命都是一個柔和謙卑的生命.
你就是用柔和謙卑的生命.
贏得人來跟從你.
祂幫助我們.
在我們罪性當中.
嫉妒很容易出現.
撒旦也找了很多機會.
在我們不知覺之間.
將種子放在我們心中.
讓它萌芽.
讓它成為毀謗別人.
甚至我們自己都以為.
我用愛心說誠實話.
甚至我都以為.
我只是在說一些真實的事情.
再求你的聖靈.
真的在我們心中常常提醒我們.
因為建立關係真的不容易.
但是嫉妒帶來的毀謗.
很快就很容易破壞了我們的關係.
主是憐憫我們.
讓我們這一群在你面前的群體.
我們真的以愛彼此的對待.
讓我們以愛心彼此的關懷.
主是讓我們真的能夠.

$^{921}$一方面接受你給我們在屬靈上的慰訓.
又讓我們帶著欣賞的心.
欣賞每一個在我們身邊的弟妹.
因為主是你自己將她擺在我們身邊.
亦都一同為你的身角來努力.
我們仰望主.
我們禱告奉基督耶穌的名義而求.
\newpage



\section{以弗所書 2:1-10}
\label{sec:032235}
\textbf{神極大的救恩 (華/粵)}~~~~姚臣泰長老
\newline
\newline
連結: \href{http://www.church.com.hk/acms/content.asp?site=cdc\&op=show\&type=product\&code=032235\&layout=sermon}{\texttt{032235}} ~~~~ 語音日期: 2026-05-31
\newline
\newline
\hyperref[sec:032234]{\small{< < < PREV SERMON < < <}}
~
\hyperref[sec:index_chronic]{\small{[返順時目]}}
~
\hyperref[sec:032236]{\small{> > > NEXT SERMON > > >}}
\newline
\newline
以弗所書 2:1-10
\newline
\begin{longtable}{cl}
\hline
\hline
章節 & 經文 (和合本修訂版)\\
\hline
2:1 & \begin{tabularx}{0.7\textwidth}{X} 從前，你們因著自己的過犯罪惡而死了。 \end{tabularx} \\ \\ \relax
2:2 & \begin{tabularx}{0.7\textwidth}{X} 那時，你們在過犯罪惡中生活，隨從今世的風俗，順服空中掌權者的領袖，就是現今在悖逆的人心中運行的邪靈。 \end{tabularx} \\ \\ \relax
2:3 & \begin{tabularx}{0.7\textwidth}{X} 我們從前也都生活在他們當中，放縱肉體的私慾，隨著肉體和心中的意念去做，和別人一樣，生來就是該受懲罰的人。 \end{tabularx} \\ \\ \relax
2:4 & \begin{tabularx}{0.7\textwidth}{X} 然而，神有豐富的憐憫，因著他愛我們的大愛， \end{tabularx} \\ \\ \relax
2:5 & \begin{tabularx}{0.7\textwidth}{X} 竟在我們因過犯而死了的時候，使我們與基督一同活過來—可見你們得救是本乎恩— \end{tabularx} \\ \\ \relax
2:6 & \begin{tabularx}{0.7\textwidth}{X} 他又使我們在基督耶穌裡與他一同復活，一同坐在天上， \end{tabularx} \\ \\ \relax
2:7 & \begin{tabularx}{0.7\textwidth}{X} 為要把他極豐富的恩典，就是他在基督耶穌裡向我們所施的恩慈，顯明給後來的世代。 \end{tabularx} \\ \\ \relax
2:8 & \begin{tabularx}{0.7\textwidth}{X} 你們得救是本乎恩，也因著信；這並不是出於自己，而是神所賜的； \end{tabularx} \\ \\ \relax
2:9 & \begin{tabularx}{0.7\textwidth}{X} 也不是出於行為，免得有人自誇。 \end{tabularx} \\ \\ \relax
2:10 & \begin{tabularx}{0.7\textwidth}{X} 我們是他所造之物，在基督耶穌裡創造的，為要使我們行善，就是神早已預備好要我們做的。 \end{tabularx} \\ \\
[1ex]
\hline
\hline
\end{longtable}
$^{1}$好,我們今天會繼續看二福所書.
我們繼續會看二福所書.
正如剛才我們所唸的這段經文.
我們剛才已經讀了這段經文.
那我把它命名為神極大的救恩.
我將這個經文標題為神極大的救恩.
因為我希望我們看見說神為我們所成就的是.
如此大,如此長闊高深的一件事情.
我希望大家能夠看見神為我們成就的何等大,何等高深的救恩.
對不起,可能需要後台的弟兄姊妹幫忙一下.
這個有時候不操作.
所以很簡單,除非你從神領受了生命.
否則你無法為神而活.
這段說話,除非你從神領受了生命.
否則你無法為神而活.
這個應該很明顯,對吧?.
這個很明顯,是不是呢?.
我們希望說我們在座的每一位弟兄姊妹.
都是蒙恩得救的弟兄姊妹.
我希望我們在座當中每一個都是蒙恩得救的人.
所以不是你做什麼,以致你成為教會的一份子.
不是,因為你做了什麼,成為教會的一份子.
也不是你在教會多久,以致你成為神的兒女.
也不是你在教會多久,以致你成為教會的一份子.
所以我們常常問弟兄姊妹的時候.
起碼我會學習說,我不是問你信主了沒有.
所以很多時候我都會問弟兄姊妹.
不是問你信了主多久.
我會問你重生了嗎?.
我會問你是不是重生了?.
因為如果你問他信主的話.
他說我信了,我曾經.
他就會追溯到他以前的一段時間.
因為如果你問他是不是信了主的時候.
他就會返回某個時刻告訴你.
但是我們關心的是.
你到底有沒有神的生命在你裡面.
但是我們其實關心的就是.
有沒有神的生命在一個人的生命裡面.
所以我們今天看經文的時候.

$^{41}$我們可以回顧一下我第一章講到的一些內容.
所以我們現在要重溫我們上一堂的時候.
我們講過什麼內容呢?.
所以保羅為教會的禱告.
他是希望說我們看見.
開我們的眼睛看見神的恩兆有何等的指望.
保羅為他們祈禱讓我們體驗什麼呢?.
體驗神的恩兆在當中有何等的指望.
並且他在聖徒中的積業有何等豐盛的榮耀.
以及在聖徒中得的積業有何等豐盛的榮耀.
也讓我們看見他向我們這信的人.
所顯的能力是何等的浩蕩.
以及讓那群人知道.
他向我們這信的人所顯的能力是何等的浩蕩.
所以你發現說.
神給我們的.
如果簡單來歸納這三個禱告的內容.
所以如果我們簡單的將保羅的禱告內容是這樣的.
你發現說關鍵是.
同一個有沒有關鍵詞.
在三句話同時出現.
有沒有看見在三段話裡面有一個同樣的詞出來呢?.
當然不要講神啦,他之外.
你發現有一個詞是「何等」.
就是「何等」.
所以你發現說.
現在關心的不是你知不知道的問題.
所以這個不是你知不知道的問題.
現在是知道神的有何等.
How big, how deep and how wide.
何等長闊高深.
所以我們今天回到這裡的時候.
你發現說一開場第二章.
他就說從前你們因性.
因著自己的過犯罪惡而死了.
而在二章一開始的時候.
第一句說話就是從前.
在這句話當中.
我們都是在這個從前裡面.
同意嗎?.

$^{81}$而我們在這句話當中提醒我們.
以前同樣都是在從前這個光景.
所以對於義父所教會的弟兄姊妹.
對於你我.
我們從前都是這樣子.
對於義父所的弟兄姊妹.
對於你我一樣.
從前我們同樣是這個情況.
所以如果你好好讀一到三節的話.
你會發現說.
這幅圖畫真的是很悲慘.
所以你仔細去看一至三節的經文來講的時候.
其實這個光景也是一個很悲慘的光景.
所以保羅在這裡就想啟示我們說.
我們在過犯罪中死了.
是因為我們在過犯罪惡中死了.
是已經死了.
死沒有程度上的.
死其實沒有一個程度上的死.
所以死就是死.
同意嗎?.
死就是死.
大家同意嗎?.
在廣東有沒有死佛,死宗.
不懂有沒有程度上的不同.
不過死就是死.
我們應該明白到死就是一個死的光景.
所以如果你看這裡一開場就告訴我們說.
死它成為我們自己一個人的一個故事的話.
我們就要留意說.
到底這段聖經保羅給我們看見.
人在神之外到底是怎麼樣的情況.
如果這個故事一開始就告訴我們.
已經是死了的情況.
究竟我們的生命的故事又是怎麼樣的講述呢?.
所以死是我們每一個人在基督之外的一個結局.
這個是每一個在基督以外的人生命的結局.
所以我不曉得你怎麼跟人家傳福音.
我不知道你們怎麼跟人家傳福音呢?.
我們等一下回來看.

$^{121}$不過我先解釋一下死的意思.
死的意思就是分開或者分離.
死的本身的意思就是分離分開的意思.
所以我們聖經當中有講到三個層面.
第一個層面就是靈魂的死.
所以我們也曾經講過.
聖經裡面講到死是分開三個層面.
第一個層面就是靈魂的死.
所以靈魂的死的意思就是說.
如果你不與神有關係的話.
你是永遠與神分開的.
所以如果你與神沒有任何關係的時候.
你就與神永遠的分開分離.
這個就是靈魂的死.
所以我再問你們死是什麼意思.
你們要回答我是嗎?.
我再問你們死的定義是什麼呢?.
就是分離分開的意思.
那肉體的死我們就很清楚了.
肉體的死我們就很清楚明白.
就是靈魂和身體分開.
靈魂和身體的分開.
這叫做肉體的死.
這個是肉體的死.
那又永遠的死.
永遠的死呢?.
我們發現永遠的死這個意義不是時間上.
而是放在我們整個宇宙裡面.
是沒有辦法改變的一個事實.
你不能從這裡到那裡.
這個不單止是一個時間線上的分離.
也是在宇宙性上我們看到的.
永遠不會改變的一個事實.
所以我們發現說死帶來一個結局.
就是我們其實是身不由己.
而死帶來一個真實的事情就是身不由己.
所以你的慾望到底是為了成就誰的旨意?.
你的慾望是為了成就誰的旨意?.
那你發現在《伊夫所書》這段告訴我們說.
我們是順服空中掌權者的首領.

$^{161}$而在《伊夫所書》這段經文提醒我們就是.
我們是順服空中的掌權者.
隨從今世的世遇.
隨從今世的風俗.
就是現今在悖逆之子心中運行的邪靈.
很可怕.
的確是很可怕.
所以如果你不相信主.
如果你不接受祂的話.
你跟祂沒有關係.
你就是在誰的手中.
所以如果你跟神沒有關係的時候.
你究竟在誰的手中呢?.
就是在撒旦.
給他做奴隸.
其實就是成為撒旦的奴隸.
那你發現當你放縱私慾.
這裡告訴我們說.
你就說放縱私慾.
隨從肉體和心中所喜好的去行.
所以如果我們是放縱私慾.
隨從心中所思想的去行.
所以變成說我們整個人是被什麼控制啊?.
被我們的私慾.
我們在沒有想到什麼.
什麼是美的事情.
所以其實當中我們就被那些私慾去控制.
不能夠幸福.
所以你發現說最終我們的去向.
我們是招惹神的憤怒.
而最後我們的去向就是招惹神的憤怒.
所以不要看錯我們是無法自拔的.
大家不要看錯.
其實我們是真的不能自拔的.
好 如果這是人在為信主之前的光景.
如果這個是人未信主之前的光景.
那我請問.
當你們向一個人.
這樣子的人傳福音的時候.
你該怎麼傳?.

$^{201}$所以如果我們向一個這樣的人.
去傳福音的時候.
我們應該怎麼去傳呢?.
你們要怎麼傳?.
你會怎麼去傳呢?.
死到不死的.
你要怎麼傳?.
死不死.
你會怎麼去傳呢?.
在聖經當中.
如果有一個很奇妙的詞.
我可以說是然而.
或者是但是.
如果聖經中有一個很奇妙的詞語.
我會說是然而 神.
或者是但是.
那你發現英文有一個書.
甚至寫成But God.
英文就是But God.
你知道嗎?.
在聖經當中.
如果你追溯到任何一個詞.
任何一段出現這個詞的時候.
你發現意思就是神介入.
祂做了一個非人能做的事情.
而我在聖經當中.
每次看到這個詞出現的時候.
我們看到神的介入在當中.
是做了一些人不能夠做的事情.
所以你注意到But God.
後面就帶出神的作為.
而每次看到But God.
然而神之後.
就是神的一個行動出現.
所以憐憫還有神愛我們的大愛.
是神的憐憫和祂的大愛.
庸子妹.
我再次說.
保羅為教會的禱告.
在第一章的後面.

$^{241}$它告訴我們關鍵是何等.
而第一章我們說到保羅的那段禱告.
關鍵詞就是何等.
所以有時候當你知道.
你們長大的過程.
不曉得你們多大才體會到父母的愛.
我不知道大家在成長過程當中.
大概幾多歲你開始體會到神對你的愛.
你父母對你的愛呢?.
你們大概幾歲?.
大概幾多歲的時候體會到?.
十歲?.
二十歲?.
三十歲?.
每個人都不一樣.
其實每一個人都不一樣.
我們都能夠體驗的.
就是當我們體驗到的時候.
你發現說你會很感動.
同意嗎?.
而我相信當我們真的體驗到父母的愛的時候.
我們都很感動.
因為你發現說父母的愛.
到底對你有多大.
因為你會體驗到父母對你的愛是何等.
是幾大.
庸子妹.
神的愛.
我們用言語.
確實是很難形容.
同樣神對我們的愛.
我們都沒有辦法用言語去形容.
但是當你細看神給我們的話的時候.
當我們細讀神給我們的說話的時候.
我們就能夠體驗到說.
神自己為我們所做的是何等大的事情.
我們就能夠體驗到.
神為我們所做的是何等的大.
庸子妹.
我剛才講了.

$^{281}$死人你要怎麼辦?.
剛才我說了.
死的人會怎麼辦呢?.
怎麼做呢?.
只有一個辦法.
就是叫他活過來.
那個方法就是叫他活過來.
所以你千萬不要跟死人打.
起來呀!起來呀!.
踢他呀!打他呀!.
有沒有用?.
因為死的人你不可以搖他.
打他就叫他復活.
是不是?.
你應該怎麼做呢?.
神做的事情是超自然的事情.
神所做的是超自然的事情.
祂叫人從死裡復活.
所以祂告訴我們.
使我們與基督一同活過來.
一同復活.
一同坐在天上.
庸子妹.
保羅用的詞非常非常的關鍵.
而保羅這段經文用的詞也是很關鍵的.
不是復活之後.
你可以隨心所欲做我要做的事情.
不是復活之後你可以隨心所欲做你想要做的事情.
而且復活是給我們一個新的生命.
不再是在我們舊的生命做一個補貼.
也不是一個復活的生命.
不是將我舊的生命補貼進去.
所以你發現基督教不是一個心靈雞湯.
所以基督教不是一個心靈雞湯.
明白什麼是心靈雞湯嗎?.
大家明白什麼是心靈雞湯嗎?.
就是你該做這個讓你的心感受比較好一點.
基督教和任何宗教都不一樣.
先告訴你的信息是你做不到.
基督教首先要告訴我們你做不到.

$^{321}$跟任何宗教不一樣就在這裡.
跟任何宗教不一樣的地方就在這裡.
其他宗教都是你要做做做拼命做.
為的是期待希望以後有更好的結果.
因為很多其他宗教就說你做做做.
做到有一天你可能會得到你想要的改變.
關鍵是可能.
這個他們提到的只是可能.
所以你說回教徒把自己炸死的話.
他是希望什麼呢?難道他只是為了犧牲嗎?.
而我們有時看到死士把自己的生命炸死了.
他們只是為了什麼呢?.
他也是想著一個極樂的世界.
以後有極多的老婆多大的很有錢的境界.
因為他們期待將來可能會有一個很好的生活.
有其他的一個豐盛在當中.
當我們認識到神為我們所做的事情的時候.
我覺得我們要好好地看說.
特別是在第七節.
我們一起讀第七節好嗎?.
當我們看到神給我們的豐盛的時候.
讓我們能夠明白.
所以我們一起讀第七節的經文.
一起讀.
為要把他極豐富的恩典.
就是他在基督耶穌裡.
向我們所施的恩賜.
顯明給後來的世代.
你知道什麼是show off?.
你知道什麼是顯擺出來嗎?.
你廣東話怎麼說?.
演ye .
所以有一天.
神就是要把歷世歷代信靠他人.
全部給天上的地上和地底下.
全部人去看.
看什麼東西呢?.
而神其實就是將歷世歷代的人.
向天上地上一起表現出來.
顯明這些人的厲害.

$^{361}$是要顯明那幫人的厲害嗎?.
這些人有多有能力?.
絕對不是.
要世世代代.
天上地下都認定一點.
是他做了一個奇蹟.
是向世世代代天上地下顯明神所做的事蹟.
讓人看到他的恩典是豐富.
是讓人看到神的恩典極其豐富.
他向我們的恩賜.
英文叫什麼?.
Kindness.
就是神向人的那種恩賜.
Kindness.
雍子妹.
你們有沒有領受過恩賜的人的一個.
Have you received kindness before?.
你們有沒有接受過人的一些好行為.
幫助感動到你?.
應該有吧?.
大家有沒有經歷過.
想到那個向你給你很好的恩典的人呢?.
當神說我要顯明.
把這些人全部給歷世列來看的時候.
他就是要讓人看見說.
這些都是我從萬人中買熟過來的.
而神要向世人所顯示的就是.
這幫人是他去買熟回來的.
雍子妹.
當你想到這件事情的時候.
絕對不是只有你信主之後.
只有你得到這個永生而已.
大家要知道的一樣就是.
不是單單只是你.
在得到這個救恩之後.
得到的這個永生這個事情.
我在開始已乏所恕的時候.
我一直提醒你們.
你們是一個重要的.
負數的一個詞.

$^{401}$對吧?.
而我開始和大家分享已乏所恕的時候.
已經提醒大家.
我們要看到的就是.
你們是一個負數詞.
我自己從聖經來看.
我們千萬不要沉迷在說.
我信主就是我自己得到而已.
而我們要大家不要沉迷一樣就是.
當我認為我信主就是我的事情.
我得到的事情.
所以我們在主日學的時候.
我們最近講到這個次要神學的介紹.
最近我們主要學的時候.
我們開始看次要的神學觀的時候.
所以什麼是Baptism?.
什麼是浸禮?.
所以在聖經當中浸的意思是什麼意思?.
聖經裡面講到浸的意思是什麼意思?.
就是聯合的意思.
當你把他Baptize.
他一定後面有一個詞就是.
Into.
進入什麼東西呢?.
Into Christ 在基督裡面.
所以在聖經當中每次講到浸禮的時候.
就是Into.
就是進入的情況.
就是你和基督聯合.
也同時把你進入什麼?.
洗入教會.
同時是將你受洗之後歸入教會.
不是T4C是地方教會.
但是歸入什麼?.
他自己的普世教會.
就是他的身體.
而我們教會只是一個地方性的教會.
而是歸入到一個普世性的教會.
所以嚴格來講.
當一個人信了主之後.

$^{441}$他不去教會的話.
有沒有問題?.
所以每一個人信主之後.
他們回教會有沒有問題呢?.
有沒有問題?.
有沒有問題呢?.
沒有問題啊.
我再問一次啊.
如果一個人信主之後.
不去教會有沒有問題?.
對不起.
如果一個人信主之後.
不去教會有沒有問題?.
很大的問題啊.
是一個很大的問題.
你發現說.
反覆的重複.
在基督裡.
所以你告訴我.
你告訴其他人的時候.
你千萬要告訴他.
你可以問他說.
你是不是在基督裡.
因為在基督之外沒有指望.
是完全的無可希望.
我們教第八節到十節的時候.
我們發現說.
很重要的.
我們得救.
是本乎恩也因之信.
第八節的時候.
也開始提醒我們.
很重要的事情.
是本乎恩.
因著信.
所以什麼不是.
救恩的.
What is not salvation.
什麼不是救恩.
而救恩裡面.

$^{481}$什麼不是救恩呢?.
所以他說.
不是出於自己.
也不是行為.
免得人自誇.
不是出於自己.
不是出於行為.
免得有人自誇.
那我問你們.
什麼是一個精彩的見證.
所以我問大家.
什麼是一個精彩的見證呢?.
我發現有些從小.
特別是你們當中.
有些是從小在香港教會信主.
或者在中國教會信主.
我也知道.
在我們當中.
有些是從小回教會.
無論是香港還是中國.
我們對於從不信到信.
好像很模糊.
會不會?.
我們很多時候.
從不信到信的過程.
有點模糊.
或者你聽別人的見證.
多精彩.
他從一個吸毒的.
變成說.
全部放棄.
神戒毒.
讓他完全的尾聲.
獻給神.
我們有時聽到.
很厲害的見證.
我覺得.
他是一個從吸毒的.
到放棄一切.
後來去服事主.

$^{521}$離開這一切.
你會羨慕這樣的一個改變嗎?.
你相信.
你會羨慕這樣的改變嗎?.
那我告訴你.
好消息啊.
我們以前都是.
死到不死的人啊.
我要告訴大家一個好消息.
一件事情就是.
我們從以前.
都是一個死的人.
如果死就是死的話.
只有一個辦法.
就是神叫你活過來.
如果死是真正的死的時候.
唯一的方法.
就是將那個人復活.
所以不要浪費時間.
跟你旁邊的人.
跟位於你身邊的人.
分享你的故事.
你的故事.
所以不要浪費很多時間.
跟其他人分享你的故事.
所以不要擔心.
你的故事沒有那麼精彩.
也不要擔心.
你的故事不是很精彩.
你仔細想的時候.
神做的是不得了的事情.
你要思想的就是.
神在你生命當中.
做了一個不得了的事情.
那你告訴我.
什麼是「是」.
神做的「是」.
究竟神在我們生命當中.
完成了什麼事情呢?.
當你留意的話.

$^{561}$你會發現說.
完全是祂的作為.
所以如果你留心去看的時候.
完全是神自己所做的事.
那請問.
在這段生命中.
有沒有記住到.
你做了什麼?.
而我們細心去看.
經文裡面.
有沒有提到.
我們做了什麼呢?.
Zero.
是完全沒有.
如果神不開恩.
如果神不主動.
做了祂的救恩.
祂開恩.
把你召回來的話.
不可能得救.
所以如果不是神先去工作.
把你救回來的時候.
你根本不可能得救.
所以本福恩恩澤信.
是一個很重要的.
很簡單.
但是是非常非常寶貴的一個信息.
所以本福恩恩澤信.
雖然很簡單.
但是是一個很寶貴的事情.
所以你千萬不要告訴別人.
你信主.
你要怎麼樣信主.
你要做什麼改變.
你要像樣.
才信主.
所以我們不要告訴別人.
怎麼去信主的時候.
要做什麼.
要改變什麼.

$^{601}$是先信主後改變.
同意嗎?.
先信主後改變.
大家同意嗎?.
所以關鍵是要清楚得救.
所以關鍵是.
我們要清楚自己得救.
所以我們關乎的是.
一個新生命.
到底有沒有的問題.
所以我們要清楚.
就是新生命的事情.
那有人來了教會幾年.
有沒有人問他一個關鍵的問題.
你到底是不是重生了.
所以有人來了教會幾年.
我們有沒有問一個很基要的問題.
就是他是不是已經重生了.
就是你看醫生.
你喜歡醫生告訴你真話還是假話.
如果你看醫生的時候.
你希望醫生跟你講真話還是假話.
你希望醫生告訴你什麼.
如果我有癌症.
如果.
我還不知道.
我最近沒有做身體檢查.
我看起來很健康.
不過是看而已.
所以如果我說我自己有癌症.
最近我都沒有做身體檢查.
我什麼都不知道.
但如果是這樣的事情.
我希望醫生怎樣呢?.
那你會希望醫生跟你講說.
你都好好的.
好健康.
什麼都沒事.
你想像剛才那樣.
醫生跟他說話嗎?.

$^{641}$還是你希望醫生跟你.
馬上一針見血地跟你講說.
你們要馬上對付.
還是你想醫生很坦白.
一針見血就告訴你.
你應該怎樣開始治療.
聽牧師所經歷的.
你想如果醫生稍微晚一點點的話.
就很危險.
我們要知道.
如果我們聽醫生所說的話.
遲了一點點的時候.
結果就可能很不一樣.
所以意思是什麼呢?.
診斷一定是在先的.
對嗎?.
所以我們要知道診斷在前.
然後才有醫治的方案.
之後才有治療方案出來.
那請問你告訴一個人說.
他是罪人的時候.
你有沒有跟他講.
他真正的光景是什麼?.
所以你告訴一個人.
他是一個罪人的時候.
你其實沒有告訴他.
他真實的光景是如何.
如果一個人對罪不清楚的話.
他對得救肯定有問題的.
所以如果一個人對罪.
是不清晰了解的時候.
他可能對得救的事情.
就不是很清晰.
所以他會以為來教會就好.
所以他就以為.
回來教會就好了.
我就跟T中姐一起嘻嘻哈哈.
我就跟大英正妹一起嘻嘻哈哈.
T中姐我們關心的一個是.
神在耶穌基督裡面的新創造.

$^{681}$我們知道每一個新生命.
就在神裡面是一個新的創造.
是祂的傑作.
是神的傑作.
英文叫Masterpiece.
英文就是Masterpiece的傑作.
所以改革中的覺醒就是在.
馬丁路德他們回歸聖經之後.
他們從神的話.
發現幾個重要的事情.
而改革中的覺醒.
我們知道當馬丁路德.
他們這一班人重回聖經真理.
他們所看到的事情.
我簡單的跟你們分享五個說啦.
我會跟大家簡單的講五個唯讀.
第一個唯讀聖經.
第一個就是唯讀聖經.
為什麼基督徒要讀聖經.
所以這就是為什麼基督徒要看聖經.
因為它是神要人知道的.
一個最清楚的途徑.
因為這是神讓人看到最清楚的途徑.
T中姐不要忘記.
神不再用.
不是說神不能.
神要人知道的啟示.
Revelation.
都是什麼.
是特殊的留給你我的字句.
他自己親自啟示的.
是神親自啟示給我們的字句.
我們知道聖經其實就是.
神要提示給人一個.
很清楚的一些字句字語.
所以你不用苛求說去做夢.
所以我們不需要乞求的.
就好像做夢.
或者是說有特別的經歷.
或者是一些特殊的經歷.

$^{721}$當一個人打開神的話的時候.
而當我們每次打開神的話語的時候.
他就是在看神對他本人要說的話.
而是神對這個人要說的話.
所以唯讀恩典.
唯讀恩典.
恩典就是根本不配的.
就是我們原本是不配的.
第三唯讀信心.
第三是唯讀信心.
你要為一個人怎樣信主.
一個人怎樣去信主呢.
難道只要信.
就能夠承受這一切的福音嗎.
難道單單一個信就得到這一切的福音嗎.
這就是信心.
最後我們信的對象.
唯讀基督.
而我們信的對象就是唯讀基督.
那你信主之後.
你覺得你的生命生活是為誰而活.
唯讀榮耀歸神.
沒有別的.
這個沒有其他了.
所以我們的錯覺.
我們有些錯覺.
你們有沒有聽過.
你跟人家分享福音.
我們就像在一個急流的河.
大家有沒有聽過.
我們分享救恩的時候.
我們說一個人好像掉進了河裡.
那將要一個萬丈的瀑布在前面.
而前面是一個急流瀑布.
所以我們傳福音就像我們.
丟一個救生圈或丟一個繩子.
所以我們傳福音就像將救生圈拋給那個人.
所以你要抓住.
你就可以得救.
而那個人就要拿著這個救生圈.

$^{761}$就會得救.
我請問你們.
我現在考試.
看你們有沒有聽得見.
請問.
你覺得這樣的福音.
這樣的故事.
對不對.
所以等於剛才所講的福音故事.
講得對不對.
你們只要做兩件事情.
是這樣子還是這樣子.
你們要這樣這樣.
有些是這樣子.
你究竟要怎樣回答.
我不明白.
你們要清楚一點.
到底是這樣子還是.
你們要這樣或者這樣.
OK.
所以如果我們嚴格來講.
是不對的.
嚴格來講.
這個畫面這個故事是不對的.
剛才我一開始想說.
一個人在基督之外是什麼人.
所以一個人在基督之外是怎樣.
誰呀.
死人.
死就是死.
沒有別的定義.
死就死.
所以你如果丟一個救生圈.
他怎麼辦.
所以丟一個救生圈給那個死人.
有沒有用.
沒有用.
因為他也不會抓.
因為沒辦法拿得到.
你把他勾到的話.

$^{801}$拉回來的話.
還是死.
就算你把他拉回來.
仍然是一個死人.
丟你明白我的意思嗎.
但是你跟著人家用這些比喻.
這些形容.
來跟人家講救恩的時候.
要小心.
所以每次我們用這些比喻.
這些情況.
跟人講救恩的時候.
我們要小心.
耶穌告訴我們.
這個完全是神自己的作用.
祂自己從萬人中.
祂就說.
不單是救.
祂叫我們復活.
這個東西.
人能夠做到嗎.
不可能.
所以我們要清楚知道.
救恩就是.
神在萬人當中.
減少我們.
不單是救了我們.
也將我們復活了出來.
這個是人做得到的嗎.
所以當然你.
明尼蘇達因為有時候太冷.
有時候會發生一些污染.
這個不是只有明尼蘇達.
有時候因為冬天太久.
水裡面的氧氣沒有了.
所以很多有魚死掉.
就是魚都死了.
我們有時候會知道.
在明尼蘇達走的時候.
如果長期冬天的狀況.

$^{841}$或者其他情況的時候.
水溫太低的時候.
水中間的氧氣都沒有了.
我們就會看到很多魚死了浮上去.
請問.
我告訴你.
這幅圖畫比較像樣一點.
這個圖畫其實比那個圖畫更像.
你告訴我.
如果你傳福音給這條魚的話.
你要怎麼傳福音.
所以如果你要將福音傳給死魚的時候.
你應該怎麼傳呢.
你人再怎麼說都沒有用.
你無論人怎麼去講都沒有辦法.
只可能你給牠一個新的生命.
讓牠活過來.
唯有就是你要給牠一個新的生命.
讓牠復活出來才行.
那第三個我想給你們看的是.
有沒有看過這幅圖畫.
大家有沒有看過這幅圖畫.
這幅是世界著名的著作.
是一個著作.
這幅世界著名的名畫.
我們不曉得這個女的是誰.
沒有人知道這個女士是誰.
所以人家說.
這是一個女孩戴著珍珠耳環的女孩.
但是我們都知道.
但是你們注意的話.
我們很多人都知道.
這位畫家是誰.
我們其實都知道這個畫家是誰.
是17世紀的Vermeer.
17世紀的一個畫家.
所以如果你真的看到.
是很美麗非常生動.
你看的時候這幅畫非常生動.
我的意思是說.

$^{881}$當你被基督得救之後.
我是說其實當我們被基督拯救之後.
人家看到你是看到誰.
所以人看到你的時候是看到什麼呢.
看到誰.
看到誰呢.
唯獨神得到榮耀.
同意嗎.
唯獨神得到榮耀.
同意嗎.
所以我們都是祂的Masterpiece.
因為我們都是神的傑作.
當人看到你的時候.
應該看到說.
哇你這個人以前是什麼爛樣子.
所以每次人家看到的時候.
就說你以前是多差呢.
什麼死樣.
怎樣死就怎樣死.
就現在怎麼活過來.
所以他們看到的是後面的這位畫家.
同意嗎.
他們應該看到後面的那位畫家.
所以不可能也沒有的自誇.
所以就不可以自誇.
也不可能我們自誇.
所以最後給你們一些再思.
我給大家做一些思想.
我們有些很清楚.
但是我需要我們一起再同一步.
有些東西好像很清楚.
但是希望我們能夠同一個Page.
同一個層面.
在基督之外沒有死亡.
我們在告訴人.
你要來教會.
你要來團契.
這些東西是幫助一個人.
能夠聽到好消息.
我們叫人回教會.

$^{921}$回團去的時候.
我們其實給他一個機會.
能夠去聽到這個福音.
但是直到一個人認清自己的光景.
是死的時候.
唯有當那個人認清自己是什麼光景.
是一個死光景的時候.
他只有一個辦法.
就是向神呼求.
而他自己唯一的方法就是向神呼求.
你們有沒有求過人.
真的苦求跪下來求.
你們有沒有求過.
你是不是會求.
是你有辦法的時候.
還是你沒有辦法的時候.
當你去跪求的時候.
是因為你有方法.
還是你沒有方法的時候.
肯定是你沒有路走的時候.
當然是你沒有路再去走的時候.
你才會跪下來.
所以你要跪下來.
庸子妹.
你有沒有在神面前跪下來說.
主啊.
我沒有除你之外.
別無拯救.
求你救我.
我們有沒有想到.
其實我們就在同樣的情況.
就在神面前跪下來.
告訴神我們再沒有任何的把握.
我們除你之外.
不要再說這個人好像很不錯.
如果基督救他的話.
他以後一定是為神所重用.
這句話千萬不要用.
所以我們要小心.
我們不可以再說這個人很好.

$^{961}$所以當他信主之後.
一定可以神重用他.
這個我們不可以這樣說.
當神要用一個人的時候.
一定是用一個能夠彰顯.
manifest他的榮耀的人.
是一個能夠彰顯神榮耀的人.
所以如果一個人得救之後.
他彰顯是我多厲害的話.
有沒有問題啊.
所以當一個信了神的人.
他看見自己的就是自己多厲害.
有沒有問題呢.
很大的問題.
這是一個很嚴重的問題.
所以聽我準備.
當一個人信主之後.
人家看到他就是什麼.
他manifest.
他就彰顯.
他本身就是神的見證.
所以我們其實一個信主的人.
讓人看見的就是神自己的榮耀.
看見背後的那位神.
這是如此的大救恩.
你稀奇嗎.
另一個就是如此的救恩.
你稀奇嗎.
我們很多時候信主之後.
我們對我們的救恩.
好像是麻木了.
我們有時候信了主多年.
或者一段時間之後.
我們對於救恩麻木了.
就像我以前跟我的那些朋友玩.
又愛說耶穌啊.
所以好像以前.
跟一班我自己熟悉的朋友.
又說耶穌.
丟著妹.

$^{1001}$老實講.
如果你從死裡被一個人救過的話.
你對那個人很大的恩情啊.
其實事實上就是.
如果你被一個人從死裡救回來的時候.
你會很感激他.
是一個很大的恩情在當中.
我本身幾乎死過幾次.
我自己本人其實已經死過幾次.
以後有機會跟那個man講.
有時是自找的.
以後會跟大家分享.
很危險的事情.
有時是自己找回來.
有時是很大的奇蹟.
你覺得我要為誰活.
你認為我是為誰而活呢.
那救恩你稀奇嗎.
救恩你稀奇嗎.
我很怕.
基督徒信主很久之後.
對救恩麻木了.
我其實很怕.
當一個人信主之後.
他對救恩麻木了.
羅馬書告訴我們一件事情.
我們的基督教是本乎神.
我們基督教是本乎神.
是從信以至於信.
意思是說憑信開始.
之後的路還是要憑信.
之後的路仍然是要靠著信心.
第三 活著是要為神而活.
第三是活著是為神而活.
我不知道你現在的中心是誰.
我不知道你現在人生的中心是誰.
你們問一下.
你們告訴我.
你們活著的誰是你們的中心.
你的中心是什麼.

$^{1041}$如果你的追求需要調整的話.
優先次序的優先次序需要調整的話.
優先次序調整的意思.
那還是有時間還是有機會.
我們仍然是有時間有機會這樣做.
第四 救恩很好行為有什麼差別.
第十節告訴我們我們原是祂的工作.
在基督耶穌裡造成的.
為要叫我們行善.
就是神所預備叫我們行的.
我們每一個蒙召得救的人都是要行善.
我們每一個蒙召的人都是要行善.
什麼是善事.
就是神要我們做的事情.
在得救之前好行為能不能救你.
能不能救.
就是不能救了 大家講得很清楚.
絕對不能救.
但是信主之後呢.
好行為是信主的見證.
因為新生命 祂有新的生命.
因為我們內在的新生命.
祂表現出來的 所想的 所愛慕的.
應該有一個調整.
祂所表現的是一個新生命的改變.
最後.
我們人生的舞臺.
最後我們人生的舞臺.
有時候我們在看說 哎呀 為什麼要受洗啊.
為什麼要受禁啊.
有時候我們會問 為什麼我們要受浸受洗呢.
你會留意到說 受禁的牧師會跟你講說.
這個是給在私人面前.
你在天使面前做見證.
牧師會說 你是在人面前.
在天使面前做這個見證.
如果你發現說 第七節告訴我們什麼呢.
這段經文第七節告訴我們什麼呢.
要給歷世歷代 天上地下所有的.
讓他們看見神的恩賜.

$^{1081}$是要給歷世歷代看到神的恩賜.
所以我在座的弟兄姊妹.
救恩是不得了的事情.
所以弟兄姊妹我們要看到救恩是一個很好的事情.
是非常極好的事情.
如果你的對基督徒的一個生活圈.
就是禮拜五或禮拜天來到教會.
然後完了就走人.
所以你的基督徒的生活圈就是說.
星期五到星期日你回來教會就走人.
弟兄姊妹 基督徒 基督教的一個生命.
絕不是只有這個層面.
我們要清楚知道基督教不只是單單這個層面.
當一個人認識到神在他身上所施的救恩是何等大的時候.
所以當一個人認識到神在他生命裡面.
所成就一個何等大的救恩的時候.
他會做什麼呢.
我們會看到他的感恩和讚美是很自然的.
壓都壓不住.
壓都壓不住他就會出來.
他一來就告訴你神多好.
他一來到就告訴你神是多好.
神是多偉大.
神是極偉大.
他是應當稱頌.
他是應當被稱頌.
我們一起禱告.
我們一起祈禱.
天父我們在你面前.
我們求你開我們的眼睛.
神我們在你面前的時候.
求你開我們的眼睛.
讓我們看見你為我們所成就的是何等大的救恩.
只因著靠著這個信.
也就是憑著你自己的恩典.
我們就進入如此蒙福的一個地步.
是憑著信藉著你的救恩.
我們能夠進入一個極大的福氣裡面.
讓我們心能夠享福.
讓我們心能夠降福.

$^{1121}$能夠明白我們現在一生的主人是誰.
讓我們能夠明白我們一生的主人是誰.
願意更多地為你而活.
願意我們能夠更為你而活.
讓我們天天心中充滿著感恩和讚美.
就是因著你做了奇妙的事情.
就是因著你所做的奇妙的事情.
我們如此禱告奉耶穌基督的聖名.
我們如此禱告 奉基督耶穌的聖名.
阿們.
\newpage



\section{約伯記 22:21}
\label{sec:032236}
\textbf{你要認識神}~~~~曹子光牧師
\newline
\newline
連結: \href{http://www.church.com.hk/acms/content.asp?site=cdc\&op=show\&type=product\&code=032236\&layout=sermon}{\texttt{032236}} ~~~~ 語音日期: 2026-05-31
\newline
\newline
\hyperref[sec:032235]{\small{< < < PREV SERMON < < <}}
~
\hyperref[sec:index_chronic]{\small{[返順時目]}}
~
\hyperref[sec:032237]{\small{> > > NEXT SERMON > > >}}
\newline
\newline
約伯記 22:21
\newline
\begin{longtable}{cl}
\hline
\hline
章節 & 經文 (和合本修訂版)\\
\hline
22:21 & \begin{tabularx}{0.7\textwidth}{X} 「你要與神和好，要和平，這樣，福氣必臨到你。 \end{tabularx} \\ \\
[1ex]
\hline
\hline
\end{longtable}
$^{1}$好 今日想同大家分享一節經文.
這節經文是來自舊約的約伯記二十二章的第二十一節.
我們一齊讀一讀這節經文好嗎.
約伯記二十二章的第二十一節.
我們一齊讀出來 請讀.
「你要認識神 就得平安 福氣也必臨到你」.
和你旁邊那個讀這節經文好嗎.
「你要認識神 就得平安 福氣也必臨到你」.
如果你是第一次聽這節經文的話.
很有可能你會想 嘩 這麼正.
我一直學武藝的那種平安 福氣.
原來秘訣這麼簡單.
只要我們認識神.
這樣這兩樣就會自動地衝上門.
2026年其實都過了一半.
不知道你覺不覺得平安.
何烏戰士 嘩 已經多少年了 四年有嗎 有了 不止.
四年都應該有了.
美國以色列對伊朗的軍事行動.
到現在都未畫上一個完整的句號.
現在香港人坐飛機都還要繞個圈.
怕了 免得搭中東的航線.
怕無端端中了導彈.
油價飆升到一個地步.
我不知道你知不知道一月份.
今年一月份到今年五月份.
即是相差大概四五個月.
你知不知道一月份平均的布蘭特原油多少錢.
是72元一桶.
五月份平均布蘭特原油的油價是103元一桶.
即是升幅42\%或以上.
嘩 如果人工有這樣升幅 真是不得了.
短短的四個月.
你試試油由一月和五月都相差很遠.
最便宜那間Costco.
我相信已經升了不少.
油價又飆升.
機票的捱貴價.
錢包的大縮水.
很多人都想找平安.

$^{41}$很多人都說可以給我們平安.
但結果我們得不到平安.
前幾天我聽到一個消息都幾震撼.
因為我以前都住過Sherfield.
Sherfield有位太太年輕 三十多歲.
沒有三十多歲 三十歲.
她大概凌晨兩點.
不知為何這麼晚 可能玩晚了.
在大學區裡面一個很旺的酒吧地方.
被人無端端槍殺了.
是那些人 可能以為target錯了.
其實他是完全無辜的事.
走過那裡 經過那裡就吃了一粒子彈.
那就沒了命.
平安可以是短短一兩秒之間就沒了.
但是如果我們讀聖經只讀這一句說話.
我們大概就會錯過了約伯記最震撼的地方.
我們說大家都看到這個真相.
讀約伯記 其實他要讀他的上下文.
當我們讀上下文的時候.
我們就會知道原來約伯記這一句說話.
是一句風涼話.
說這句話的人不是上帝 也不是約伯.
是遠道而來探望約伯的三位朋友的其中一個.
他的名字叫做以利法.
最諷刺的就是以利法說這句說話的對象是誰.
他對著約伯說.
大家知不知道約伯是哪位仁兄.
如果你常常去教會 你都聽過約伯這個響噹噹的人.
約伯你以為他是二打六.
約伯是根據新約的《雅各書》說.
他是苦難中最好的忍耐者.
厲害吧.
你不要說最好的忍耐 忍痛打兩支針就叫救命了.
但是約伯是苦難的忍耐者.
還有舊約的《以西結書》說到.
在古代最出色的最偉大的三個屬靈人之一.
就是約伯.
你說厲不厲害.
更厲害的地方就是.

$^{81}$約伯的第一章第八節.
上帝親自在撒旦面前點名稱讚約伯.
看一下第一章第八節.
他說地上再無人像他完全正直.
敬畏上帝遠離惡事.
你說他厲害不厲害.
是不是很出色的一位屬靈的偉人.
你可以想像一下一個圖畫頂尖的朋友.
以利沙竟然叫這個人去認識神.
就好像一個完全沒有讀過醫科的人.
走去一個外行人.
走去一個心臟科頂級的權威教授面前跟他說.
教授 你要認識一下心臟科.
簡直是開玩笑.
但是以利法為什麼會這樣說呢.
因為他心目中有一套根深蒂固的意識形態.
今天就很想藉著以利法和約伯說的這句話.
帶大家拆解一下這個認識神背後的.
三個不同的深層次的認識神的誤解.
無論是我們信得主的或者未信主的朋友.
我們心目中對上帝的誤解.
可能我們可以解得開.
第一層次就是要正確地認識神.
打破了一種所謂的交易式的信仰.
以利法口中的認識這個字.
"Sakan"這個字在白萊文的意思.
其實你可以指他認識.
但是原來二十二章第二節有另外一個翻譯.
這個翻譯就是指有益或者有利.
所以以利法他的前台詞是什麼呢.
他說約伯 你快點去親近神.
因為信神是對你有好處.
不知道你覺不覺得信仰對人是有好處的呢.
同不同意 贊不贊成.
信仰 如果你不贊成的話.
你今天沒理由坐在這裡.
你坐在這裡 你起碼認同.
信仰對人是有價值的.
除非你老公或者你老婆強行逼你來.
否則的話 你坐在這裡.

$^{121}$我估計你應該對信仰有益處.
這個命題不會反對.
不只是以利法這樣想的.
信仰對人是有益處.
現代的科學或者心理學.
都大量的研究贊成這件事.
例如 原來信仰可以幫人解除一些不好的習慣.
例如 污酒 例如 吸煙 賭博.
這些信仰可以幫到我們.
還有 宗教團體能夠提供一些社交的支援.
所以我們就明白為什麼.
宗教信仰是對人有益處.
另外有一項調查顯示.
原來有信仰的患精神病患的人士.
他們接受治療之後.
他們的效果顯著是比沒有信仰的精神病患者.
是快速痊癒多一些.
還有一個研究是非常重要.
特別對女士們.
如果你是50歲或以上.
不用舉手了.
50歲或以上的姐妹們或者婦女們.
原來你有宗教信仰.
你的死亡率是低於沒有信仰的人士20\%.
20\% 你想一下有多厲害.
原來有信仰是有這麼多益處.
最近我看到一個報告.
是2020年美國史丹福人類學專家.
就有一個專業的報告.
講到原來宗教信仰有什麼好處.
原來 第一就是心理和情緒會帶來積極的改變.
例如你祈禱的話.
你可以將一些不好的事情放下.
還有第二.
你的專注力和比較正面思維.
透過宗教的表達.
你會發現人的專注能力不單止是集中.
更加可以轉移到一些順利或者好的事情上.
那就不用想一些不開心的事情太多.
還有第三樣東西就是帶來平靜.

$^{161}$帶來你感覺被愛的一個感受.
還有第四 原來有信仰是會增加你免疫系統的.
不知道大家有沒有發現.
來英國之後你的免疫系統是高了還是低了呢.
我是一個親身示範來的.
我來了英國之後.
我的免疫系統是越來越差.
如果你發現你有HEY FEVER的話.
就證明了其中一樣非常明顯的原因.
你的HEY FEVER就會越來越差.
你的皮膚差了也是和免疫系統可能有關係的.
但是原來信仰是可以令你的免疫系統強化的.
還有第五樣東西就是.
原來經過MRI掃描之後.
發現有信仰的人的腦部的反應.
是和你有很多朋友的.
深度朋友聊天的時候那個反應是很相似的.
換句話說.
是會令到我們沒有那麼多寂寞的感覺.
還有第六.
如果信徒和神明顯的有好的關係的話.
他的滿足度.
人生是會好一點的.
是對於沒有信仰的人是更加好.
更加豐富的.
所以你記不記得信神或者信仰是對我們有益的.
但是不知道你有沒有留意.
有很多人所說的信仰有益.
不是這些.
是些什麼呢.
我做了教會這麼久真的見過無數.
第一.
信耶穌好啊 幫兒子找基督教學校.
我簽過無數這麼多東西.
簽完之後.
那些人沒多久就不見了.
因為有時候去基督教學校.
你回來教會會加分的.
所以那段時間你會見到那些人出現.
有些人說去教會好啊.

$^{201}$有益處啊.
方便找結婚對象.
年輕人不知道會不會是這樣.
因為教會裡面的人都叫純樸一點.
知不知道.
咦.
像是吧.
你看看你旁邊那些多純樸的樣子.
還有呢.
還有就是進教堂拍拍照多漂亮.
現在更方便.
因為你可以發現.
原來回到教堂之後.
可以找牧師幫你簽那些什麼.
搞居留啊.
搞入籍啊.
那些東西.
更加特別的地方就是.
有時候有些人會認為.
回到教會的時候.
找基督教的墳場更加方便.
英國是沒有的.
香港應該可能有些機會.
一旦搞定這些東西之後.
教會不一定會回來了.
寧願多點時間去旅行.
去一下這裡去一下那裡.
有空就回來教會.
如果我們對神的認識.
只是停留這樣的層次的話.
我們就將創造宇宙萬有的這位真神.
降格成為一個人生福利交易所的總經理.
這個正正就是以理法的一個問題.
他以為和神之間是一場交易.
我做得好.
你就賜福給我.
約伯出了事.
一定是他沒有交足貨.
所以問題就在他那裡.
但是約伯的第一章.

$^{241}$已經告訴我們了.
撒旦去到上帝面前.
講了一句很重要的話.
在第一章第九節.
約伯敬畏神.
豈是無辜呢.
撒旦的意思是.
人愛神.
只是因為一場交易.
神幫我.
我就相信了.
神幫我.
我為什麼要相信他.
神聽我祈禱.
我就相信他.
神聽我祈禱.
我即時就不回教會了.
不去教堂了.
不再相信神.
但是約伯用自己的行動.
粉碎了這個邏輯.
當你知道子女都死了的時候.
大部分的財產都沒了的時候.
他跪在神面前.
在一章第二十一節.
我們讀一讀這一節經文.
第一章二十一節請讀.
賞賜的是耶和華.
收取的也是耶和華.
耶和華的名是應當稱重的.
約伯說我敬畏上帝.
和上帝.
是否賜福給我.
是否保護我.
是沒有關係的.
是完全沒有關係的.
真正的有益.
不是想發達.
不是想消災解難.
不是想幫我任何的事.

$^{281}$真正的有益在哪裡.
使徒行傳第四章第十二節.
說了給我們聽.
我們一起讀一下這篇經文.
請讀.
除他以外別無拯救.
因為在天下人間.
沒有次下別的名.
我們可以靠著得救.
這一個就是利益.
這一個就是好處.
好處就是.
我們每個人都怕死.
但是我們可以征服死亡.
我們每個人都有罪.
解決不了罪惡這個問題.
但是耶穌.
祂可以幫我.
這是一個利益.
是一個好處.
信仰的好處.
就是令我們起碼是這樣.
不單只是這樣.
更加重要的就是.
我們認識到這位真神.
和祂建立一個好的關係.
這是第一樣東西.
第二樣東西.
認識神是打破了.
律法主義的恐懼.
認識這個字.
Saccan這個字.
其實有第二個不錯的翻譯.
它叫做「復和」.
或者是「修補關係」.
經文就跟著說.
它說「就得平安」.
原文更加準確就是說.
「跟他和好」.
這樣的話.

$^{321}$主動權就不同了.
跟他和好.
就是人做的主動.
神給的和平平安.
就是神的主動.
所以以理法真正的意思就是.
約伯.
你跟神之間出現了裂口.
你要主動走第一步.
就是向神修補這個關係.
跟他和好.
這樣神才會將福氣給你.
以理法的邏輯非常清楚.
受苦.
因為是你犯罪.
為什麼你受苦是約伯.
為什麼你家人受苦.
是因為你或者你家人犯了罪.
犯罪離開神.
所以就搞成這樣.
就算你沒有罪.
你表面沒有罪.
你一定有隱藏的罪.
甚至你家人其他人的罪.
所以你必須承認.
和要跟神建立這個關係.
不過.
以理法這個思維.
其實今天有很多人都是這樣.
有這樣的想法.
例如.
我不知道你有沒有這樣的經歷.
家人發生意外.
或者你自己本人發生意外.
你第一個反應就是說.
是不是我做錯了什麼.
有沒有試過呢.
有沒有試過.
又或者是.
是不是我上個禮拜沒有回教會.

$^{361}$所以上帝作弄我.
是不是上帝刑罰我.
當世界天災人禍的時候.
我們很多人就會想.
是不是人類得罪了上帝.
上帝要降罰你的世界.
舉一個例子.
大家還記不記得南馬海嘯.
南馬海嘯應該出世了.
看到大家還沒出世的應該在那邊.
在起點那邊.
應該還沒出世.
我估計應該還沒出世.
南馬海嘯的時候.
大家知不知道死了多少人.
我估計是有20萬人或以上.
當時有一個流行的說法.
是什麼說法呢.
都幾傳遍基督教的.
死的人大部分都是某一個種族的人.
或者某一個信仰群體的人.
你聽完這句話之後.
不知道你有什麼反應.
有什麼感受.
我聽完第一個反應就是.
不是吧.
不是沒有基督徒死嗎.
沒有基督徒死嗎.
很多人死了在南馬海嘯裡面.
你怎麼能夠因為他們拜偶像.
所以上帝全部要滅絕他們.
我們憑什麼.
可以這麼大膽去這樣說.
死了的基督徒又怎麼解釋呢.
上帝從來都沒有交代過.
為什麼南馬海嘯死了這麼多人.
為什麼我們有這樣的想法呢.
我年輕的時候.
我都會去祈禱會.
所以我先賣一賣廣告.

$^{401}$逢星期四.
大概至少三個禮拜.
晚上幾點鐘.
八點鐘.
我們有網絡的祈禱會.
這個禮拜就有了.
請記住.
我小時候都會去祈禱會.
晚上大概六七點左右.
我們自己的教會播道會.
有一次我就去祈禱會.
準備禁門鐘.
你知道香港大廈下面有個鐵閘.
你要按門鐘.
我準備按門鐘的時候.
就有一個男士.
年輕人來的.
就截停我.
他說.
我問.
他說.
你很面熟.
我說我不認識你.
他說你叫了我的女兒.
我說我什麼時候叫了你的女兒.
是啊.
我根本沒有女兒.
怎麼叫你的女兒呢.
他說.
他說我的女兒說是你.
我說你打給她.
我跟她談.
然後他真的在旁邊.
我們有個電話.
就是隔壁的店鋪.
他真的用那些搜撥的電話去打.
打完之後.
我說你給我跟她談.
他說我的女兒說就是你了.
然後我就說.

$^{441}$我不認識他.
發生什麼事.
他說你跟我來蘇屋村.
我們當時十多歲.
去了蘇屋村之後.
他就問我拿錢.
其實只是打劫.
何必搞這麼大件事.
最後就拿了五十元.
港幣.
是港幣五十元.
可想而知我當時不是有錢的人.
那我回到教會.
我都繼續回到祈禱會.
你又要欣賞我了.
一個年輕人經歷這樣的情況之下.
都繼續按門鈴.
上去回到祈禱會.
上去祈禱會又想著有些慰藉.
想著有些人說.
哎呀年輕人你真乖.
那都回到祈禱會.
誰知道第一句說話.
我當說回這件事之後.
那個姐妹跟我說.
她第一句說話說.
你今天有沒有靈修.
我真的惡了然後她答應我.
而當時我這麼傻的.
我就真的以為.
真的因為我沒有靈修.
所以上帝給了我這樣的教訓.
接著我回到家.
晚上想了想之後.
不是的.
那天是有靈修的.
為什麼都會這樣呢.
是不是.
我們很多時候有這種.
律法主義的精神.

$^{481}$這種意識形態.
很多人會覺得上帝是一些很嚴厲.
很冷漠.
很無情的執法者.
時時刻刻都要抓你的錯處.
看看你今天做錯了什麼.
看他明天又做錯了什麼.
做得不對嗎.
罰鬼你.
這個是我們的假設.
我們的觀念.
我們的上帝真的這樣.
抄你牌那個或者會.
我不知道大家有沒有被人抄牌呢.
有一次我去探人.
在聖詹姆醫院.
那天我就在停車場附近.
離遠看著一個抄牌人.
你知道抄牌人.
通常都會被人不太喜歡他.
我離遠看他有什麼動靜.
我當時是早上十一點五十五分.
在他站在一輛車的前面.
我就猜他應該是想抓這輛車.
抄這輛車.
應該是那個人買了張票是十二點鐘.
他就站在那裡五分鐘.
我看著他.
我真的研究這個人他做了什麼.
他到十二點的時候.
真的發出一張告票就給這輛車.
他真的要抓這個人這輛車.
一點點的德性都沒有.
一點點的恩典都沒有.
一步就抄你.
這個世界是這樣的.
你以為上帝都是這樣嗎.
我告訴你錯了.
上帝不是這樣.
有些人以為上帝是這樣.

$^{521}$如果你認為上帝是這樣.
今天透過這段經文.
我想告訴你.
不是這樣的.
約霸面對這三個朋友的指控的時候.
他覺得極其委屈.
但他選擇在上帝面前做一個真誠的人.
向上帝發出三條很真摰的問題.
第一條問題記載在第三章二十節.
他說受患難的人為何有光賜給他.
心中受苦的人為何有生命賜給他.
約霸問這條問題是人這麼脆弱.
上帝你為何給我們有生命.
你為何給我們有光明.
第二條問題在第七章第十七節.
他說人算什麼.
你竟然看他為大.
將他放在心上.
人這麼脆弱.
約霸說生命這麼短暫.
這麼渺小.
為什麼你盡這麼認真對待每一個人.
還有第三條問題在第十章第三節.
你所做的你又欺壓.
又藐視.
卻光照惡人的計謀.
這事你以為美呀.
約霸說你創造了我.
為什麼有弱惡事發生在我身上.
你看我的女兒我的子女.
他們都是多年輕.
這樣就沒有了.
如果你真的將他們放在心上.
我為什麼會這樣.
約霸不是挑戰為什麼有苦難.
他是在問一位這麼偉大的上帝.
為什麼這麼緊張.
這麼在乎每一個渺小脆弱的人.
而且又讓他們經歷這麼多惡事呢.
然後出乎意料之外.

$^{561}$大家如果讀完約霸記.
你就知道上帝親自出現.
和約霸來說話.
但是上帝不是派一個代表.
或者不是派一個先知.
或者不是透過夢境和約霸來顯現.
而是親自和約霸來說話.
這樣和古希臘裡面.
一部著名的悲劇叫做《普羅米修斯》.
這個悲劇是非常相似.
《普羅米修斯》說到他受苦的時候.
有神人走到他那裡來勸他.
和約霸記這三位老朋友.
一起圍攻約霸好像一模一樣.
但是最終希臘這個悲劇的結局.
就是至高無上的宇宙宙司.
這位神子沒有出現過.
只是派了一個使者來傳話.
那你就會發現有點像約霸記.
因為有第四個人物在約霸記最後.
是出現和約霸說了很多話.
不過約霸記最後.
原來上帝是不需要任何人.
幫他傳遞訊息.
因為他親自和約霸.
來到他面前來說話.
約霸的好朋友伊利法和約霸說.
你要主動去和上帝和好.
這個其實和福音剛好相反.
我們不是需要去神和上帝和好.
而是上帝主動降卑.
來到我們這個人世間.
去建造和好的關係.
這個正正就是耶穌基督的故事.
當人類還活在罪惡和絕望當中的時候.
無法靠著自己和神修補好關係.
上帝不是坐在高天.
讓我們每一個人搞定了.
然後才讓我們進入天國.
不是的.

$^{601}$上帝不是這樣的.
上帝是親自鍍成肉身.
成為了人.
進入了苦難的世界裡面.
被人出賣他.
你想過.
你有沒有想過我當時有什麼感受.
當這個小孩在我面前哭泣的時候.
他認錯的時候.
其實當他說sorry那一刻.
我心裡面已經沒事了.
是完全沒事了.
我作為他爸爸心裡面一早已經原諒了他.
他其實不需要做任何事情.
去彌補自己做錯的任何事情.
也不需要做任何宗教的行為.
去平復他所謂的內疚.
因為在那一刻.
我作為他爸爸仍然是他爸爸.
還是那麼愛他.
這就是一個這樣的關係.
頂尖的各位朋友.
這正正是天父的心腸.
以理法同約伯說.
約伯你的禍你要做好.
你要和神和好.
福氣先至來到.
其實約伯就好像被誤會的我另一個孩子.
他沒有偷來喝的.
而真正偷喝的那個.
我不需要他做任何事情去彌補.
他作為我兒女的身份是不會改變.
我仍然都是那麼愛他.
我們經常有個錯覺.
以為我做好一點.
再做好一點.
多回教會.
聽多點聖經.
我才可以達標.
其實不是的.

$^{641}$在上帝的角度.
不是等你去高攀我.
是當你還是滿身污穢的時候.
就好像寫滿了.
I am sorry.
的這個孩子.
當你還是成王成孔的時候.
耶穌基督在十字架上面.
一早就已經為你寫下原諒的福音.
還有第三樣東西.
就是認識神生命的那種深度的熟悉.
我們剛才說打破了交易式.
打破了這個律法主義的誤解之後.
認識神到底是什麼.
認識神其實就是一個深度的建立一個認識.
希伯來文這個實勘這個字.
其實根據很多字典告訴我們.
它是說一種習慣.
一種熟絡.
又或者一個緊密關係形成的.
它不是頭腦上的知識.
而是一種經過長時期的相處.
形成了一種習慣.
進而你可以彼此之間.
知道對方療癒指掌.
聖經舊約裡面出現了十幾次這個實勘這個字.
其中有兩次翻譯成為.
和這個習慣而來的熟絡是相關.
我舉兩個很明顯的例子你一定知道.
第一個地方是記載在《民數記》第22章第30節.
那裡就說到有一個先知的名叫做巴蘭.
他的奴就有一天跟巴蘭說.
說些什麼呢.
他說我不是你從小時直到今天所騎的奴嗎.
我常向你這樣實勘.
這樣這樣做嗎.
就是說如果你是認識我的話.
我從來都沒有試過這樣做.
這是第一個例子.
第二個例子大概你記得了.

$^{681}$大衛穿上掃羅的那件戰衣.
你記不記得大衛那個時候.
聖經裡面描述他.
大衛把刀誇在戰衣上面.
事事能走不能走.
因為掃來沒有穿慣實勘這個字.
就是說他從來都不是這樣的.
所以是從一個習慣而來的深度的熟悉.
你可想而知舉個例子如果你是一個彈鋼琴的人.
你跟一部琴相處了十年.
你每天都彈他的時候.
你有沒有感覺到那個鍵盤的熟悉度.
那種感情在這裡.
是經過時間的磨合.
經過風雨建立的這一份的關係.
約伯的這個朋友以利法.
他認為約伯應該要跟神建立一個深度的關係.
這樣就不會出事了.
這樣就不會犯罪了.
這樣你跟他家人就不會搞到這樣的田地了.
當然正如我剛才所說.
約伯不是因為跟神的關係不好.
也不是因為犯了什麼罪搞成這樣.
以利法他提到認識這個字實勘這個字.
其實是反映人真的很需要建立跟神深度的關係.
這件事確實也是對的.
在約伯記的故事的結束.
約伯再沒有聽到這三個朋友的聲音了.
這個以利法的聲音也沒有聽到.
出現的是在旋風裡面.
他聽到上帝來說話.
沉默的上帝.
他終於說話了.
我們見到這位上帝.
他沒有在法庭裡大家噴口水對質.
他完全沒有.
他帶約伯進入一個深度遊.
讓他在裡面說一些約伯從來沒有聽過的說話.
這個正正因為約伯跟他建立一個好的關係.
透過一些深度的關係.

$^{721}$透過摸索年月的經歷.
以至到他能夠跟上帝有這樣的親切的體會.
這個其實也是我們一個很重要的東西.
我們對上帝的關係.
不是別人怎麼跟你說.
也不是像銀行轉帳那樣.
發個Send字就按了出去.
或者這樣就過了帳.
完全不是這樣的.
而是我們需要時間來建立的.
所以你看到有時候我們要邀請一些我們的家人.
我們的朋友來教會崇拜.
有時候去一下閻惟軒.
有時候去一下魯友記.
有時候去一下男人侍宮.
日積月累這樣.
你就對神會更加多認識.
慢慢就建立了一種跟上帝的關係.
你看約伯.
他跟神有一個很深的經歷.
神跟他說完.
上帝帶約伯去看大自然界.
看動物界.
他問約伯.
我建立大地根基的時候你是在哪裡.
你知道那些野山羊什麼時候生產.
什麼時候出生.
你見到老鷹在天空的巢穴.
是誰幫他來到有這麼銳利的眼睛.
可以看到這麼多東西的.
上帝傳遞的訊息是.
約伯啊你連這些野獸.
這些星獸.
我都用這麼奧妙的智慧去設計和照管.
你又懂多少.
難道你認為我作為上帝這位創造者.
我最愛的人.
你所經歷的苦難.
會在我計劃之外嗎.
我們經歷的一切的事.

$^{761}$難道在上帝計劃之外嗎.
就好像一個世界級的鋼琴大師.
彈奏一首極其深奧的樂章.
如果一個不太懂音樂的小朋友就問這個大師.
大師啊.
為什麼你在這裡彈一個鋼琴調.
是不是彈錯了.
大師不會馬上跟他說樂理的.
可能就會跟這個小朋友說.
小朋友.
你不需要明白每一種樂理的.
你只需要坐下.
閉上眼睛.
相信我所彈奏出來的東西.
這個美麗的樂章就可以了.
約伯其實最後在第20章.
第42章第5節.
有一句非常重要的說話.
想和大家分享這節的《聖經》.
42章第5節我們一起讀出來.
請讀.
我從前風聞有你.
現在親眼看見你.
以你發口中.
那個神.
聽別人講的那位的上帝.
因果報應的那位的上帝.
大腦理論的上帝.
這個是風聞.
親眼看見你的就是.
當我們在苦難在旋風的裡面.
就好像約伯那樣.
他發現到上帝竟然親自臨到在我們身邊.
在我們不明的人生裡面.
我們遇見這位鋼琴大師坐在我們隔壁.
一個真正深度認識上帝的人.
他最真實的表現.
是在日常生活裡面.
約伯之後活多了140年.
在140年裡面.

$^{801}$他每天都繼續上班.
每天都繼續打理他的財務.
繼續每天對著同一個老婆.
繼續養育兒子女.
面對生活的種種柴米油鹽.
一個真正認識上帝的人.
不是經常飛上天的.
絕對不是的.
不是不吃人間煙火的.
而是他去到街市那裡.
都是誠誠實實.
在面對衝突的時候.
面對不開心的時候.
面對病痛的時候.
他會流淚的.
心底裡面.
但仍然對世界有一份不能奪去的平安.
因為上帝給了這個平安.
放在他心裡面.
因為他知道賞賜的是耶和華.
收取的都是耶和華.
弟兄姊妹朋友.
在結束之前.
我想和大家分享一個真實的故事.
這個故事是幾年之前.
有一位大學教音樂系的教授.
他本身是一個很敬虔的基督徒.
他先生是一位很傑出的外科醫生.
他們一家有一個十四歲的兒子.
非常之有音樂天份.
一家三口.
原本就過了一個令人羨慕的人生.
很有福氣的生活.
不過突然之間有一個暑假.
先生帶著兒子去衝浪.
突然之間被一股無情的浪撲過來.
先生為了救他的兒子.
就被海浪衝走了.
再沒有回來.
而兒子雖然救得回來.

$^{841}$但是因為腦部缺氧太久.
醒來之後.
智力只能夠有五歲左右的智力.
而且雙目都失明.
連媽媽都認不出來.
一天之內幾乎沒有了兩個人.
那幾個月裡面.
那些很熱心的基督徒.
走到他身邊和他說.
你要堅強.
上帝有他的美意在裡面.
有些就說.
你會不會有些隱藏的罪惡你還不知道.
所以上帝要這樣來教育我們.
有些就說.
約伯也是這樣.
堅強點 姐妹.
你要多點祈禱.
和神和好.
上帝就會醫治你.
這個媽媽聽完這一堆的東西之後.
她寫了一些日記.
她說.
這些宗教的口號.
我聽到每一句都像一巴掌打在我身上.
我說我問上帝.
你親手創造的這個家.
現在你又搞散了.
你以為美嗎.
我不明白背後的惡理.
我只是聽到我人生裡面的惡章裡面.
充滿了刺耳的聲音.
有接近一年的時間.
她沒有再回教會.
她覺得上帝很殘忍.
很冷漠.
直至有一天的深夜.
這個媽媽因為抑鬱.
沒辦法睡覺.
在客廳裡面一直哭.

$^{881}$突然之間她聽到兒子的房門打開了.
這個雙目生明的.
智力退化的一個十四歲的男子.
摸黑一步一步的走到客廳的鋼琴裡面.
這個兒子坐在鋼琴的椅子裡面.
用他已經看不到的眼睛.
但是又一般般靈活的手.
開始在鍵盤上彈.
當然了.
全部都是很難聽的音樂.
完全沒有神蹟.
絕對沒有.
聽到媽媽很感動.
但是她可以發現到.
當刻她有一個心境不同了.
她就坐在鋼琴裡面.
坐在鋼琴裡面.
然後就在高音區裡面.
坐在右邊.
當她彈每一音.
她就在另一邊彈另一音.
是她自己作上去的.
兒子彈得多難聽的樂章也好.
但是他右手邊.
就突然充滿了一個很美麗的樂章.
因為他是一個音樂的教授.
所以他用技巧去修飾.
去伴奏.
去調和.
幾分鐘之後.
客廳裡面響起的.
是一首很奇特.
又淒美.
但是又很好聽的樂章.
原本很雜音的音樂.
在媽媽右手的高音伴奏之下.
竟然變成一首.
充滿了深度的靈魂的樂章.
在這一刻.
她媽媽坐在琴椅裡面.

$^{921}$一邊彈.
眼裡邊一邊流眼淚.
那一晚.
她彷彿.
她看見上帝.
她彷彿看見上帝.
上帝藉著這個琴聲.
跟她說.
她說女兒.
你現在的人生.
確實充滿了命運無情的雜音.
你摸不透.
你也不明白.
但你不需要明白的.
因為我就坐在你旁邊.
我就是那個鋼琴大師.
我沒有離開過你.
我的右手就在高音區裡面.
為你所悲哀的人生.
去彈奏更好聽的音樂.
你只需要繼續信任我.
將你殘破你的雜音交給我.
我就能夠伴奏一首最美麗的樂章.
這個媽媽在那天晚上.
重新跟神和好.
她沒有得到一個雙倍的祝福.
坦白說.
但她卻是可以親眼經歷到上帝.
可以.
以前是風聞有上帝.
現在是親眼見證到上帝.
各位朋友 弟兄姊妹.
今天以理法的這句福音金句.
我們走了一個神學的旅程.
人生很多時候是很艱難.
2026年世界有戰爭.
我們的生命很多時候有疾病.
有意外.
有很多事情是我們沒有想過.
有很多的雜音.

$^{961}$有很多的痛苦出現.
但今天.
耶穌基督想跟我們每一個在座的人說.
認識我.
認識我不是一場交易.
認識我不是一個冷漠無情的因果公式.
而是我願意坐在這裡.
在你和你的生命的人生路程上.
我和你一起走.
親手修補你和上帝之間的關係.
你心目中的神到底是風聞有你的神.
還是親眼見你的神.
如果你今天願意走出一步.
放下你的框架.
邀請這位鋼琴大師坐在你人生生命裡面的時候.
將你的雜音調教的話.
我邀請你和我一起去祈禱.
我們一起閉上眼睛.
我們一起去禱告.
如果你願意耶穌基督.
進入你的人生裡面.
將你雜亂的陰域.
來到去改變的時候.
上帝和我們一起禱告.
我們一起呼喚我們的救主耶穌基督.
親愛的主耶穌.
因為你是那位創造宇宙萬有的主宰.
你也是來到這個世界上救贖我的主宰.
我承認我是一個罪人.
甚至我承認在這個世界上.
我不明白很多東西.
所以我今天將我的生命.
完全全交在恩主的手裡.
你帶領我走過人生的每一步.
沒錯.
可能我人生有很多雜音.
但是上帝.
你是那位偉大的生命岳章者.
你為我普奏美好的人生.
我將我的生命完全交在主的手裡.

$^{1001}$求你隨聽我的禱告.
奉耶穌基督 恩主的命求.
阿們.
\newpage



\section{士師記 3:1}
\label{sec:032237}
\textbf{聖經人聖經事:留下這幾族 試驗以色列}~~~~黃克勤牧師博士
\newline
\newline
連結: \href{http://www.church.com.hk/acms/content.asp?site=cdc\&op=show\&type=product\&code=032237\&layout=sermon}{\texttt{032237}} ~~~~ 語音日期: 2026-06-03
\newline
\newline
\hyperref[sec:032236]{\small{< < < PREV SERMON < < <}}
~
\hyperref[sec:index_chronic]{\small{[返順時目]}}
~
\hyperref[sec:032238]{\small{> > > NEXT SERMON > > >}}
\newline
\newline
士師記 3:1
\newline
\begin{longtable}{cl}
\hline
\hline
章節 & 經文 (和合本修訂版)\\
\hline
3:1 & \begin{tabularx}{0.7\textwidth}{X} 耶和華留下這些國家，為要考驗所有未曾經歷迦南任何戰役的以色列人， \end{tabularx} \\ \\
[1ex]
\hline
\hline
\end{longtable}
$^{1}$又一日了,我們又一齊讀聖經.
《生命樹葉》和大家一齊看聖經人士.
看看,認識一下神怎樣幫助他的百姓.
最緊要看完之後,除了更多認識神.
又更多認識我們自己.
今次看了一件事就是.
留下幾族,誰留下呢?.
神在以色列當中讓他們留下幾族.
不是神特意要留下的.
是要清剿他們.
不過他們留下,要有些不動,不趕.
神就說好,讓你留下幾族.
這樣就對了,既然你要留下幾族.
不趕他們走,不殺死他們,不清剿他們.
好,就讓你吧.
這樣也可以試驗以色列.
讓你留下幾族,試驗你自己.
看看你是不是這麼厲害.
其實神是否要試驗我們才認識我們呢?.
不用的,神知道,全知道.
只不過很多時候神給我們的試驗.
是讓我們知道自己怎樣.
很多時候我們以為自己.
我這麼高,跳得很高的.
這個欄在這個區區一個多尺的欄算什麼.
誰知道走過去一試的時候.
整個連人帶欄掉在那裡.
要試驗的.
以為自己很厲害,我拿得很重.
誰知道是拿起的.
不過一會兒就已經手軟了.
原來是沒有力的.
以為自己很厲害,意志力很堅強.
不過一眨眼就軟腳蟹,輸光了.
其實那些神讓我們以為自己很厲害.
留下一些東西,自恃.
結果那些東西就是神給我們.
以為做試驗我們自己.
讓我們自己知道.
神不需要試驗就知道我們了.

$^{41}$不過是讓我們自己知道.
我們先看史詩第三章.
耶和華說留下這幾族.
為要試驗那不曾知道.
與迦南人爭戰之事的新出世那群.
在迦南裡面的以色列人.
好叫以色列的後代.
又知道又學習未曾曉得的戰事.
說什麼呢?.
神留下幾族.
有一樣東西要做的.
第一,是要試驗他們.
剛才我都提了.
這個試驗就不是.
我先做試驗看看結果是怎樣.
我不知道,所以我要做試驗.
我不知道自己的血是怎樣.
所以我要驗血,做試驗.
但對神來說,祂不是要知道結果的.
祂早就知道了.
也不是要試驗才知道.
神一再就知道.
不過我剛才提了.
這個試驗是讓那些人知道自己的實力與結果.
試驗是讓那些不曾知道與迦南人爭戰之事的新人.
他們進迦南地經過很多場仗.
第一,為了耶利哥城.
第二,攻打艾城兩次才可以.
因為有一次阿干這個惹麻煩者令他們失敗了.
然後下面南部的五個王一起聯合攻擊他.
然後北部的王無數那麼多一起攻擊以色列.
打了很多場仗才獲得這些地方.
那些不曾知道與迦南人爭戰之事的以色列人.
一定是進了以色列之後才出世.
而且他們的上一代慢慢都好像早幾天讀的聖經.
已經死了.
他們都死了,但全城不夠好.
所以是另一代以色列人上來.
是不知道耶和華或認識不夠.
不會不知道耶和華,當然知道.

$^{81}$不過是認識不夠,心經力不夠.
所以要給他們一些試驗.
讓他們知道原來神是真神.
原來自己那麼差.
原來我的弱點在這裡.
原來我們是怎樣.
這是要留下記憶給他們試驗.
否則他們從來都沒有經歷過任何災難,患難,戰爭的日子.
於是他們是不會成長得好.
不夠精練,沒有磨練.
只不過一直都是含著金鎖匙出世.
他們不會了解這個那麼困難的世界,未夠成熟.
第二就是很叫以色列的後代.
就是這班新生代的以色列人.
知道和學習,知道即是認識和學習.
未曾曉得的戰事.
以往的以色列人全部能徵慣戰.
都是這樣,由埃及出來一直打到迦南.
但是這班在迦南出世的新生代.
是不懂得打仗的.
是沒有經歷的,是含著金鎖匙出世就已經很舒服.
讓他們知道,第一個就是試驗.
第二就是讓他們認識學習戰事是怎樣打的.
所以怎樣學習戰事怎樣打.
最好讓你打一下就知道.
如果你說學習戰事,我們在書本上學習是不行的.
打拳,學習這樣,如果你這樣來我就這樣擋.
最重要是兩個手遞起,保護著臉部.
這裡保護著胸口,如果這樣我就起一腳.
一定要打過才行,練到肌肉自然反應.
這樣會這樣會這樣會這樣.
他們未曾學習,所以他們要試驗又要學習.
這個就是留下幾軸的目的.
第一,試驗.
第二,知道又學習.
樓下有什麼人呢?神讓他們.
當然你留下的人,最好就是用那些人來試驗和讓你學習.
留下了非利士人有五個首領.
其實有五個城,五個首領,很厲害.
一班迦南人,一齊迦南人.

$^{121}$怎樣一齊呢?.
其實迦南地那個地方所有的人就叫迦南人.
裡面有很多族的.
不過就叫迦南人,統一來說就是迦南人.
譬如香港有多少區?18區.
你住在哪裡?我住在香港,你全都住在香港,我住在香港油麻地.
18區,油麻地在哪裡?很多人.
問題就是,油尖旺,你全都住在那裡,18區.
香港人,統稱迦南人.
西頓人是比較北的.
黎巴嫩山的希美人,又是北的.
巴力克門和巴力克門山直到哈馬口,這些全部都是北部的.
巴力,你知道的,是迦南人拜的神,統稱叫巴力主人.
巴力克門山意思就是Mount Hermon,是北部的.
由於那裡可能有很多外族在那裡拜巴力.
哈馬口也是北部的一個地方名.
那裡的人,留下這幾族就是要試驗以色列人.
知道他們肯不肯聽耶和華正摩西吩咐的誡命.
留下它,都說神都知道的,是讓他們知道.
原來我們沒有聽到耶和華透過摩西給我們的誡命.
我們就自己知道,以為自己好好的,靈性好勁,原來什麼都不懂.
原來是讓他們自己知道.
我們略為看看,當日以色列十二支派當中的迦南人有什麼人呢?.
很多,剛才也讀了很多個.
這裡有部份是講不了這麼多的,在這裡也未必找得出來.
還有什麼呢?迦南人有黑人,阿摩利人,比利西人,希美人,耶布斯人.
他們在哪裡呢?.
西頓人在這個位置,北面一點.
希美人,北一點,就是克門山.
另外外面少提到,其實後尾是有提到的,就是亞蘭人.
當然這些就是以色列有關的民族羅德的兩個女兒的後裔.
阿摩人,摩亞人,以東就是以素的後裔,也就是亞國那邊的.
另外,耶布斯人是在這一邊,就是屬於變亞敏茲派的耶布斯人.
耶布斯人就是大衛在跟耶布斯人買了一塊地.
跟耶布斯人買了一個和牆,在那裡在後尾起聖殿.
所以耶布斯應該是在大概這一塊.
在這個沿著地中海這一塊地,由丹之都,猶大.
有阿實特人,以實基倫,加薩,加特,很多.
這些地方就是非利士地,非利士人有五個王在這裡.
阿瑪利人在下面,很多,剛才提到的黑人都是在這個基拉爾.

$^{161}$就是以素,亞國的兒子以撒,以撒就有以素,亞國.
以素就娶了兩個當地人的女兒,一個就是黑人.
當時他們在基拉爾附近,就是這個位置.
所以黑人可能在這裡,其實黑人很多都在不同的地方.
黑人這個黑其實就是驚懼,fear的意思,懼怕.
可能那些人很害怕,那麼多種人,他們在做什麼呢?.
他們是娶他們的女兒為妻,以色列人娶他們的女兒為妻.
都說了不要和跟你才集在一起的那些人來往.
先是說不要來往,然後不要和他們成親.
不要拜他們的神,說得很清楚.
誰知他們還娶他們的女兒為妻,男的娶他們為妻.
而女的就嫁給他們的兒子,即是通婚,一通了婚即是怎樣?.
兩家人在一起,本來兩口子結婚是挺好的.
未結婚的時候兩個人自己,一結了婚,大家都多了一倍親戚.
哇,複雜異常,斷鏈了,起碼跑兩場,是不是這樣?.
叫多了一個人做爸爸,叫多了一個人做媽媽,很複雜的.
多了很多的,通婚之後他們就糟糕了.
因為通婚之後他們會將很多文化背景,宗教在一起.
拜了他們的神巴力為主.
第三點就是他們並且事奉他們的神.
這裡是最糟糕的,跟外族在一起的時候就發生了這些事.
回顧一次,神留下幾族,不是神特別留的.
是他們沒有趕走,沒有殺死,於是神既然這樣,好吧,你留下就留下.
不過這些留下的就會試驗給你自己知道,你是不是跟隨神的.
又給你學習一下要跟他們留下的幾族打仗.
是你們今所時出世,未打過仗,要學一下打仗的.
很多人,很多種族,然後他們不單止這樣.
第三點,他們跟他們一起結婚,然後事奉他們的神.
這裡出事了,神就非常不高興了.
你也知道,現在耶和華你都忘記了.
走去敬拜事奉那些神了,於是神的怒氣就開始發作.
他們就極其困苦,於是你又哀求神.
你不是說應許要幫助我們到永遠的嗎?.
神就說好吧,於是我給一個,我就回轉,我就給一個史詩給你.
轉一下頭又拜過.
所以耶和華說留下讓,照著他們這樣留下幾族.
是要試驗以色列的.
所以弟兄姊妹你也要留意到,有很多事情,你自己讓它留下.
那些讓它留下,你也知道不妥的事,神說了不妥的事,你讓它留下.
其實那幾樣事情就在試驗你.

$^{201}$當然也讓你知道你怎樣去應對這些挑戰.
也是學習到成長的,不過最重要不要氣絕.
不要回頭離開了神,你要跟隨耶和華.
不要叩拜,不要敬拜別的神,不要跟隨了它.
這是最為重要的,留下試驗一下,神讓我們知道自己多一點.
一起禱告.
主我們看到以色列他們沒有照著你的吩咐.
將留在迦南地裡面的人除滅.
終於那些人是令到他們歷世歷代都跟隨了不同的神.
很多時離棄你,神你向他們發怒很重要,最後還讓他們亡國.
漂流四方,那就是一個很悲慘的民族.
主要是你自己的選民,但是你也容讓他們這樣.
是要讓他們知道,試驗到它,讓他們學習到更多的成長.
主我們其實如果一早聽你的話,就不需要經歷那些困苦.
是極其的困苦,讓我們更多認識你.
就跟隨你,走吊直路,吸引我們.
順天祈禱奉耶穌名求,阿們.
\newpage



\section{以賽亞書 65:17-25}
\label{sec:032238}
\textbf{新天新地}~~~~羅婉明傳道
\newline
\newline
連結: \href{http://www.church.com.hk/acms/content.asp?site=cdc\&op=show\&type=product\&code=032238\&layout=sermon}{\texttt{032238}} ~~~~ 語音日期: 2026-05-31
\newline
\newline
\hyperref[sec:032237]{\small{< < < PREV SERMON < < <}}
~
\hyperref[sec:index_chronic]{\small{[返順時目]}}
~
\hyperref[sec:032239]{\small{> > > NEXT SERMON > > >}}
\newline
\newline
以賽亞書 65:17-25
\newline
\begin{longtable}{cl}
\hline
\hline
章節 & 經文 (和合本修訂版)\\
\hline
65:17 & \begin{tabularx}{0.7\textwidth}{X} 「看哪，我造新天新地！從前的事不再被記念，也不被人放在心上； \end{tabularx} \\ \\ \relax
65:18 & \begin{tabularx}{0.7\textwidth}{X} 當因我所造的歡喜快樂，直到永遠；看哪，因為我造耶路撒冷為人所喜，造其中的居民為人所樂。 \end{tabularx} \\ \\ \relax
65:19 & \begin{tabularx}{0.7\textwidth}{X} 我必因耶路撒冷歡喜，因我的百姓快樂，那裡不再聽見哭泣和哀號的聲音。 \end{tabularx} \\ \\ \relax
65:20 & \begin{tabularx}{0.7\textwidth}{X} 那裡沒有數日夭折的嬰孩，也沒有壽數不滿的老人；因為百歲死的仍算孩童，未達百歲而亡的算是被詛咒的。 \end{tabularx} \\ \\ \relax
65:21 & \begin{tabularx}{0.7\textwidth}{X} 他們建造房屋，居住其中，栽葡萄園，吃園中的果子； \end{tabularx} \\ \\ \relax
65:22 & \begin{tabularx}{0.7\textwidth}{X} 並非造了給別人居住，也非栽種給別人享用；因為我百姓的日子必長久如樹木，我的選民必享受親手勞碌得來的。 \end{tabularx} \\ \\ \relax
65:23 & \begin{tabularx}{0.7\textwidth}{X} 他們必不徒然勞碌，所生產的，也不遭災害，因為他們和他們的子孫都是蒙耶和華賜福的後裔。 \end{tabularx} \\ \\ \relax
65:24 & \begin{tabularx}{0.7\textwidth}{X} 他們尚未求告，我就應允；正說話的時候，我就垂聽。 \end{tabularx} \\ \\ \relax
65:25 & \begin{tabularx}{0.7\textwidth}{X} 野狼必與羔羊同食，獅子必吃草，與牛一樣，蛇必以塵土為食物；在我聖山的遍處，牠們都不傷人，也不害物；這是耶和華說的。」 \end{tabularx} \\ \\
[1ex]
\hline
\hline
\end{longtable}
$^{1}$鼎芝妹平安.
今天其實是講耳塞書的最後一講.
我不知道你們還記不記得第一講是講甚麼.
但我相信不緊要的.
今天你能夠記住你都能夠捉到.
因為過往應該都是經過頗長的時間.
都很長.
所以耳塞書其實都分了一個頗長的階段.
很多的識經書就分了三部份.
第一耳塞牙.
第二耳塞牙.
和第三耳塞牙.
分了三部份.
上次文姑娘都有講了少少.
分了三部份.
第一部份其實是講他被老的初期.
去到第二部份.
就是講被老後期和回歸的早期.
現在就講第三耳塞牙.
由56章到66章.
上星期文姑娘就講56章.
為甚麼會有第三耳塞牙.
第三耳塞牙的重點是甚麼呢.
上次講聖殿已經建好了.
聖殿建好了.
去到第三耳塞牙的時候.
其實聖殿已經建好了.
他們是身處波斯帝國的年代.
當聖殿建好了之後.
隨之而來就有很多的祭禮.
很多的祭禮.
也都形成了.
在過程中.
他們也都同樣面對一些挑戰.
雖然聖殿建好了.
有祭禮.
但其實他們另一個面對一個身份的挑戰.
那個身份的挑戰是甚麼呢.
就是面對波斯帝國.
這個大的帝國.

$^{41}$他們以前是一個國家.
但今時今日他們沒有了.
他們只是屬於波斯帝國.
裡面其中一個省.
耶路撒冷只是一個省.
所以他們在他們的身份.
或者各方面裡面.
其實有他們的張力.
也有他們的挑戰.
也有他們摸索不到.
其實到底是怎樣呢.
譬如說是波斯人.
屬於波斯還是不屬於波斯.
屬於猶太人.
但又沒有了一個國家.
所以其實是在他們重尋自己的身份和角色.
在各方面裡面.
有他們自己一些很難經歷的事.
所以去到第三二賽亞的時候.
其實是講聖殿的建好了.
其實在過程裡面.
就圍繞聖殿.
也圍繞他們去找回他們的角色.
找回他們的身份的時候.
很自然就回到信仰裡面.
回到信仰裡面.
還有回到摩西律法裡面的實踐.
我們再看下去.
他們就問.
我們是屬於耶和華的.
我們是屬於耶和華的子民.
那麼耶和華的子民有什麼特徵呢.
你看回在56至58年裡面.
聖殿獻祭.
敬拜.
另外就是守安息日.
上星期文工有講.
另外就是禁食.
即是神喜悅的禁食.
另外就是公平和公義.

$^{81}$都是56章.
所以其實是在56,57,58.
或者在那幾章裡面.
都是圍繞他們問.
如果我們是屬於耶和華的.
如果我們是屬於耶和華的子民.
其實我們有什麼特質.
我們是否在做這樣的特質.
而另外他們也在期盼.
期盼什麼呢.
就是綜合的新天新地.
在這裡其實有一個背景.
我們發覺他回去.
然後他興建聖殿.
好像很好.
但其實去到內裡.
整個當時在耶路撒冷是荒涼的.
荒涼的意思是.
其實回去的人是不多的.
回去的人不多.
同時你也想想.
他們當日被擄了去.
無論是巴比倫還是波斯.
其實他們已經是怎樣.
已經是settle下來.
可能我們容易理解是買屋.
或者買地買屋.
或者生兒育女.
其實已經是在那裡生根.
但他們有些人要再回去.
其實你說多不多人回去呢.
是不多人回去的.
回去的時候其實是百廢待興.
回去的時候也有荒涼在裡面.
同時你即管想想.
很簡單.
如果你的房子裝修完的時候.
裝修完那一刻.
即未計你清潔那一刻.
你會覺得真的很亂.

$^{121}$你見到也會頭痛.
你想想他回去耶路撒冷的時候.
很多東西他們要重建.
其實在經濟上有他們的困難.
在政治上他們有他們的困惑.
也有被壓迫.
另外一件事.
就是在道德上.
其實你看56,57,58章.
在道德裡面有很多混亂.
有很多欺壓.
有很多謊話.
在58章讀一兩節經文.
你就會知道.
在聖經裡面說.
看你們禁食的日子迎頭利益.
勒逼人為你們做苦工.
你們禁食卻互相憎恨.
以凶惡的拳頭打人.
聽到沒有.
你想想當日那種屬靈的狀況.
是不理想的.
所以就出現了.
為何會有65章.
就是在期盼新天新地的來臨.
亦都是在裡面.
以塞一直在說的時候.
或者以塞書在說的時候.
其實是一個在屬靈裡面的體見.
在祈禱裡面說出.
新天新地到底是怎樣.
我相信這個期盼.
可能都是今天我們.
你和我在面對今天的生活.
今天的政治.
或者今天所有不同的環境裡面.
上帝給我們一個新的出路.
和新的眼光.
我們一起禱告.
主當我們來到你面前.

$^{161}$去看新天新地的時候.
主求你幫助我們.
將我們固有的眼光放下.
或者將我們裡面的思維都放下.
讓我們單單在你裡面去看.
看到你的盼望.
看到你所跟我們說的話.
求真理的聖靈.
你親自開啟我們的心.
打開我們的眼.
以致我們能夠聽到.
能夠感受到.
能夠我們的心.
真的被你親自去觸摸.
以致主這個新天新地.
對我們今天來說是有意義的.
而不是一個好像遙不可及.
主你親自幫助我們.
聖靈你在當中與我們每一個同在.
我們感到高奮.
耶穌基督明求.
阿們.
我們下去看經文.
剛才有讀到.
你們要當因我所做的.
永遠歡喜快樂.
我不知道你們看十八十九節的時候.
有沒有留意到.
有些事常常重覆.
有沒有.
有是不是.
不要這樣說.
我用顏色給你們看.
重覆的.
歡喜快樂.
所喜所樂.
歡喜快樂.
必不再聽見哭泣的聲音.
和哀號的聲音.
你會見到.

$^{201}$他常常重覆.
重覆說那份喜樂.
在新天新地裡.
其中一樣.
第一樣的福氣.
他說的就是那份喜樂.
那個喜樂.
就是可喜可樂之成.
不是你喝的那種.
不過希望你下次喝的時候.
記起這件事.
在上帝裡面.
上帝給我們那種喜樂.
為什麼會是喜樂.
我相信.
其實也是對應當日.
他們面對著社會.
面對著宗教.
面對著那種壓迫.
或者面對著很多很多不同的事的時候.
他們心裡面有他們的愁苦.
或者不知道出路是怎樣.
但是今天.
純在新天新地裡.
透過先知告訴他們.
那個新天新地帶給他們的.
是一份喜樂.
這個喜樂.
也特別.
特別有紅色字.
那個「造」字.
那個「造」其實就是說創造.
這個「造」是.
這個創造和創世記的創造.
還有.
和去到啟示錄.
天上的耶路撒冷降臨的時候.
那個創造.
那個是一樣的.
都是那個字.

$^{241}$而這個「造」.
在《釋經書》裡看.
這個「造」其實是一個participle.
即是說.
是在進行中.
進行中.
即是說.
神所造的.
無論是這個城.
這個民.
這個「造」是一路一路.
繼續的.
不是停在那裡.
其實這個「造」字.
這個繼續.
帶給我們一樣.
我們心裡面覺得.
原來神造的東西.
不是做在那裡就放在那裡.
而是神造的時候.
讓我們的歡喜.
讓我們生命裡面有很多不同的經歷.
神都在當中.
去一路創造.
一路做.
以致我們的歡喜.
我們裡面的喜樂.
不是停留在昨天的經歷.
而是我今天.
今天我也在經歷.
神創造.
創造的時候帶給我的那份喜樂.
神的同在帶給我的喜樂.
而這個就是.
在《新天新地》裡面所說的.
其實這個《新天新地》.
和《啟示錄》的《新天新地》.
其實是有一點點不同的.
我們一路看下去.
你就會更加的明白.

$^{281}$好,這裡就是.
第一件事.
那個福氣是什麼呢.
就是那個喜樂.
這個城.
是一個可喜可樂的城.
也就是說.
也是對應當天跟他們說的.
你們回來覺得.
唉.
我撿了一支.
我這叫什麼.
撿了一支不好的籌.
回來.
有《聖經》書說.
當天他們回耶路撒冷.
是撿籌的.
這樣抽籤回來的.
撿到一支.
唉,嚇嚇籤了.
要回來了.
但是呢.
今天在這裡.
上帝跟他的子民說.
今天回來耶路撒冷.
其實這個是一個.
有喜樂的城.
在裡面縱使艱難.
是有喜樂.
縱使艱難.
但是呢.
不會再有那個哭泣.
哀號的聲音.
因為神在當中.
神跟祂的子民同在.
好,去到20.
20就是.
這裡說.
沒有幾天就夭折的嬰孩.
也沒有受數不滿的老童.

$^{321}$老者,sorry.
老者.
然後呢.
他也說.
百歲死的話.
仍然算是孩童.
也就是說很年輕.
也就是說.
是一個長壽.
一個健康在裡面.
而他說.
如果百歲死的罪人.
就是被咒了.
在猶太人裡面.
他們覺得.
你很早死.
是一個咒詛.
而在這裡.
神也說了一件事.
如果他真的這樣的話.
是神的咒詛.
因為在這裡.
是上帝給人.
有受數的滿足.
有健康.
我們看下去的時候.
在這個新天新地裡面.
神給人有長壽.
有健康.
而且長壽指的是那種永遠.
就是說.
讓我們看回一件事.
就是上帝給我們.
在當中所活的日子.
是神所喜悅.
並且是神祝福的.
也有神的公義在裡面.
這是第二個福氣.
第二個福氣.
我們再看下一個.

$^{361}$剛才我們有讀到.
是關於一個工作.
在新天新地裡面.
原來還有工作.
但這個工作是不同的.
我又把一些重複的字眼.
給你看到.
我們再讀的時候.
慢慢回答一下.
這裡到底第三樣東西.
第三樣在新天新地裡面.
神給我們的福氣.
是什麼福氣.
我們一起讀.
他們要建造房屋.
自己居住.
栽種葡萄園.
吃其中的果子.
他們建造的.
別人不得住.
他們栽種的.
別人不得吃.
因為我們的日子.
必像樹木的日子.
我選民親手.
勞碌得來的.
必長久享用.
他們必不徒然勞碌.
所生產的.
也不遭災害.
因為都是.
蒙耶和華賜福的後裔.
他們的子孫.
也是如此.
給了一些顏色給你看.
特別是因為.
我想把它的重複.
讓你更加容易去看到.
當你看完這件事的時候.
你會覺得.

$^{401}$這個工作是怎樣.
值得的.
就是說.
你做的事是有成果的.
別人能否掠奪你的東西.
是不可以的.
我想在職場裡面.
我們就更加明白.
有時你很努力做的事.
一下子被人拿走了.
當是他自己.
又或者.
有很多事.
不聲不響.
但不知為何.
那些東西落在別人身上.
你真的會覺得.
很勞氣.
你會覺得很不公平.
但是在這裡.
是說回.
我相信不只是工作.
而是神在說.
你付出的成果.
你付出的努力.
是神會有回報.
是神紀念的.
我們看下去.
工作之後是怎樣.
有一字講是.
你是可以享受.
你自己辛苦所得的成果.
並且.
其他人.
是不能欺壓.
別人的成果.
就是說.
屬於你的東西.
是屬於你的.
你所付出的.

$^{441}$神紀念.
沒有了貪婪.
沒有了掠奪.
沒有了欺壓.
沒有了不公.
沒有了一些.
你心裡面覺得.
是很多人洗滑手.
而得來的東西.
然後.
這個福氣是.
勞碌工作中.
你可以享福.
因為神親自賜予給你.
所以.
這個新天新地.
你會想到.
不是沒事做.
是在裡面.
神給我們的福氣.
是什麼.
你勞碌.
但你能夠去享福.
去享受你的成果.
我不知道你有沒有試過.
做完事之後.
你享受你的成果.
你會覺得很舒服.
很開心.
我曾經.
很久很久以前.
很久以前.
有時候很討厭做家務.
但是.
有幾次試.
做完之後.
坐下來.
真的很乾淨.
很舒服.
你會覺得.

$^{481}$真的好.
我覺得那次.
上帝跟我說.
你享受你勞碌得來的.
這是你的份.
自此之後.
我就覺得.
我也不是很埋怨.
當我能夠去勞碌.
但也能夠享受.
這是恩典.
所以在裡面.
新天新地裡面.
那個福氣.
第三個就是.
你可以享受.
別人不能夠拿走.
因為神為你.
量定這一份.
是給你的.
好了.
去到第四樣.
一共有四樣.
第四樣.
這個就是.
不如我們一起讀.
豺狼必與羊羔同食.
獅子必吃草與牛一樣.
塵土必作蛇的食物.
在我聖山的片處.
這一切都不傷人不害物.
這是耶和華說的.
有時候.
我看完這個.
你也會覺得.
真的嗎?.
不要說動物.
人和人之間.
只要有時候做了什麼.
都會受傷.

$^{521}$有看過這個鎖匙牌嗎?.
暴躁.
不要惹我.
今天你發覺.
你出街.
只要伸手一搖.
一群都是暴躁的人.
很害怕.
動不動就在街上吵架.
不要再說.
今時今日很流行的校長罵人.
你真的發覺.
上帝.
不要說動物.
人.
人都這樣.
但是.
在上帝裡.
神親自在他裡.
上帝讓我們一起去經歷.
在他生天生地裡.
這是他要扭轉的.
說到這個.
給你看一幅圖畫.
有一位畫家畫了出來.
我覺得挺有趣的.
小朋友摸了獅子的鬍子.
都是在表達這幅圖畫.
但我相信.
這幅圖畫也好.
這段經文也好.
這段經文背後要表達一件事.
在神的生天生地裡.
祂要賜福的是.
很想讓人不再爭鬥.
不再有殺戮.
暴力.
人與人之間的敵對能夠拿走.
那種傷害能夠拿走.
這是神很想在這裡.

$^{561}$告訴他的子民.
神要施行新的事.
這個生天生地的福氣是.
這件新的事.
你和我覺得不可能.
但在上帝裡.
上帝要施行他自己.
也在幫助我們.
在多元或很糟糕的人際關係裡.
但上帝要我們.
他的子民去學.
不比較.
不驕傲.
不爭奪.
以致我們能夠.
在復和的關係裡.
在恩典裡.
在神的愛裡.
在耶穌基督裡.
我們能夠和諧共處.
和平安.
所以當我們看完.
這個生天生地的四個福氣的時候.
喜樂.
長壽.
工作得回報.
和平或平安.
這四個在生天生地裡.
其實也是當日.
當日在以賽亞的年代裡.
因為我相信.
這個生天生地.
不只是說他將來很虛無的事.
而是當以賽亞書.
他當日的受眾.
就是那班被擄歸回.
和他們很努力.
去建整個民族.
整個信仰的時候.
是上帝釋放在當中.

$^{601}$給他們那份安慰和盼望.
也在裡面.
真的好像是.
Here I am.
我在這裡.
上帝你用你的話語.
用你的盼望.
安慰他的子民.
所以在以賽亞書裡.
你會發現.
有很多神罵他的子民.
但其實同樣也有很多.
神在呼籲他的子民回轉.
回轉歸向他.
神的救贖.
在以賽亞書裡面.
你一邊看一邊看的時候.
其實一直都有.
一直都有.
不只是來到生天生地才是這樣.
而是那個救贖.
那個提醒.
和那種.
跟他們說.
你回來.
回來我要復興你們.
我要更新你們.
更新我的子民.
這個就是生天生地.
其實同樣在復興和更新他的子民.
這個生天生地.
是在和當日的.
整個以色列民族說.
當日他們在面對那個環境.
面對其實我是誰.
我在耶路撒冷的時候.
亡了國.
然後縱使好像回來.
我重建聖殿.
但其實我是誰.

$^{641}$上帝你是不是還在.
這個就是帶來一份盼望.
就是上帝和他的子民說.
你們要回來.
我的救贖在這裡.
你們要回來.
我要復興你們.
我要重新去更新你們裡面.
這個更新就是.
無論在心裡賜給他們喜樂.
告訴他們.
神看顧他們.
告訴他們.
你今天很勞苦.
或者你問.
我做事是不是只有一個吉的時候.
神告訴他們.
不是只有一個吉.
是回報你們.
是有回報.
你是可以享受.
你勞碌所得來的東西.
所以.
在十三節.
這裡今天沒有讀.
但其實在八節至十六節.
是同樣在回應十七二十五節.
十三節.
他在這裡說.
主耶和華如此說.
我的僕人必得吃.
你們卻飢餓.
我的僕人必得喝.
你們卻乾渴.
我的僕人必歡喜.
你們卻蒙羞.
這裡是在說.
有兩類人.
我的僕人.
另外還有一個就是你們.

$^{681}$那到底是甚麼.
我的僕人是甚麼.
倚靠耶和華的人.
回歸到耶和華裡的人.
是屬於耶和華神的子民.
而你們呢.
就是他罵那班離棄了上帝的人.
因為在當日裡.
他們.
有人.
有很多的欺壓.
有很多奸詐在裡面.
剛才我讀了一段經文.
縱使是在禁食.
都可以打鬥.
禁食.
都可以有很多勾心鬥角在裡面.
即是說.
是徒有一個宗教的外表.
但是其實.
就好像.
聖經說.
我的子民.
他的嘴唇好像在親近我.
但他的心卻遠離我.
他在說的就是這件事.
所以.
主耶和華如此說.
我的僕人必歡喜.
你們必蒙羞.
這兩類人.
其實在說的是甚麼.
你選擇.
其實不是當日.
不單止是當日的以色列民要選擇.
今天.
我們坐在這裡聽的人都要選擇.
其實.
我們都有很多我們的憂暗.
又或者.

$^{721}$今天我們和神的關係是怎樣呢.
今天我們和神的關係是.
我坐在這裡.
但不代表我真的在親近祂.
今天我和神是怎樣呢.
又或者.
我只是在做很多很多事.
但我內裡是如何呢.
我內裡面屬靈的生命又是如何呢.
這隻經文提醒我們.
我的僕人必得吃.
你們卻飢餓.
今天你是飢餓的那個.
還是在上帝裡飽足的那個.
今天你的心是蒙羞的那個.
還是你在神裡你可以歡喜.
你今天.
做的所有事都是.
表面套路無功.
還是神說.
你做的事是可以有回報.
我紀念.
我紀念你所做的事.
這就是一個選擇.
在聖經裡面經常說選擇.
你怎樣選擇.
你要選擇生命還是選擇就坐.
所以去到24節的時候.
24節.
我們一起讀.
很美好的經文.
他們尚未求告.
我就應允.
正說話的時候.
我就垂聽.
嘩.
如果是這樣.
是不是很好.
那個好是因為.
當我們回轉歸向.

$^{761}$當上帝更新我們的生命的時候.
神說.
你未求我就已經應允.
因為你和我的心意相通.
你未求的時候已經應允.
你都未講其實我就已經垂聽.
因為你知道我的心意是怎樣.
這就是神一個最大最大的祝福.
神的同在.
所以.
回到最後.
才讀剛剛開頭那一節.
看呀.
我做新天新地.
看呀.
我做新天新地.
即是說看呀.
你們要看.
你和我都要看.
那個看是在屬靈裡面的看.
你看不看到.
上帝今天就是做一個新天新地.
以前是怎樣.
神不紀念.
也不再追了.
但很重要的一點就是.
今天.
我做這個新天新地.
同樣這個造節.
這個造節出現了幾次.
神一路做一路做.
一路在做你的生命.
無做你的生命.
一路在改變更新你的生命.
今天.
你願不願意給神去做.
給神去更新.
以致你能夠享受.
並且你擁有這個新天新地的價值觀.
還有.

$^{801}$真正最大的福氣就是神同在.
神和你一起.
你也知道.
這個新天新地裡面.
你得著神給你的喜樂.
你得著神給你的平安.
整個群體能夠和諧共處.
縱使有很多艱難.
但在神的裡面.
我們一起同行.
我想在《異塞俠書》.
其實當我們一路看的時候.
從來都不是平坦的道路.
是很艱難的.
真是很艱難.
那種艱難可能比我們今天更加艱難.
但也不可以這樣比.
有時我所面對的事.
或者你今天所面對的事.
你也覺得.
實在太艱難了.
但神知道.
神明白.
以致上帝釋放他的恩典.
釋放他的愛在當中.
神同在.
神給新天新地.
新天新地就是.
你能夠擁抱上帝給你的喜樂.
你能夠擁抱到.
神紀念你.
你所作的神紀念.
並且你能夠享福.
你能夠經歷到.
在上帝裡面那份新心靈的健康.
你能夠享受.
在上帝裡面.
神親自給你那份平安.
頂梓梅來到二賽亞書最後一講.
原因是新天新地.

$^{841}$是你能夠擁抱到.
我都能夠擁抱到.
以致今天我們踏入去的時候.
上帝也更新我們.
好,我們一起禱告.
主你說.
上帝不再紀念和追究我們過往所作.
今天就讓我們的眼目.
只是單單抬頭仰望你.
仰望在你裡面.
一切都是新的了.
並且仰望在你裡面.
上帝打開你的膀胱.
你就是很想我們經歷上帝的同在.
經歷你的愛.
經歷你的赦免.
經歷你的更新.
主,你幫助我們吧.
主,以致我們能夠真正在當中.
我們親自被你去無助.
這個造,這個創造是.
一路一路.
今天你無助我們.
明日,將來你也繼續無助我們.
以致我們活在你所喜悅的裡面.
多謝主.
幫助我們每一個人.
讓我們記住.
上帝,你就是我們的神.
我們除你以外.
我們沒有拯救,也沒有倚靠.
多謝主,多謝主耶穌.
我們感禱告,奉耶穌基督名靠.
阿們.
\newpage



\section{申命記 8:1-10}
\label{sec:032239}
\textbf{感恩 MODE (1) 記念思想}~~~~符氣松牧師
\newline
\newline
連結: \href{http://www.church.com.hk/acms/content.asp?site=cdc\&op=show\&type=product\&code=032239\&layout=sermon}{\texttt{032239}} ~~~~ 語音日期: 2026-05-30
\newline
\newline
\hyperref[sec:032238]{\small{< < < PREV SERMON < < <}}
~
\hyperref[sec:index_chronic]{\small{[返順時目]}}
~
\hyperref[sec:032240]{\small{> > > NEXT SERMON > > >}}
\newline
\newline
申命記 8:1-10
\newline
\begin{longtable}{cl}
\hline
\hline
章節 & 經文 (和合本修訂版)\\
\hline
8:1 & \begin{tabularx}{0.7\textwidth}{X} 「我今日所吩咐你的一切誡命，你們要謹守遵行，好使你們存活，人數增多，可以進去得耶和華向你們列祖起誓應許的那地。 \end{tabularx} \\ \\ \relax
8:2 & \begin{tabularx}{0.7\textwidth}{X} 你要記得，這四十年耶和華---你的神在曠野一路引導你，是要磨煉你，考驗你，為要知道你的心如何，是否願意遵守他的誡命。 \end{tabularx} \\ \\ \relax
8:3 & \begin{tabularx}{0.7\textwidth}{X} 他磨煉你，任你飢餓，將你和你列祖所不認識的嗎哪賜給你吃，使你知道，人活著，不是單靠食物，乃是靠耶和華口裡所出的一切話。 \end{tabularx} \\ \\ \relax
8:4 & \begin{tabularx}{0.7\textwidth}{X} 這四十年，你身上的衣服沒有穿破，你的腳也沒有腫。 \end{tabularx} \\ \\ \relax
8:5 & \begin{tabularx}{0.7\textwidth}{X} 你心裡要知道，耶和華---你的神管教你，像人管教兒女一樣。 \end{tabularx} \\ \\ \relax
8:6 & \begin{tabularx}{0.7\textwidth}{X} 你要謹守耶和華---你神的誡命，遵行他的道，敬畏他。 \end{tabularx} \\ \\ \relax
8:7 & \begin{tabularx}{0.7\textwidth}{X} 「耶和華---你的神必領你進入美地，那地有河流，有泉源和深淵的水從谷中和山上流出。 \end{tabularx} \\ \\ \relax
8:8 & \begin{tabularx}{0.7\textwidth}{X} 那地有小麥、大麥、葡萄樹、無花果樹、石榴樹，那地也有橄欖油和蜂蜜。 \end{tabularx} \\ \\ \relax
8:9 & \begin{tabularx}{0.7\textwidth}{X} 那地沒有缺乏，你在那裡有食物吃，一無所缺；那地的石頭是鐵，山中可以挖銅。 \end{tabularx} \\ \\ \relax
8:10 & \begin{tabularx}{0.7\textwidth}{X} 你吃得飽足，要稱頌耶和華---你的神，因為他將那美地賜給你。」 \end{tabularx} \\ \\
[1ex]
\hline
\hline
\end{longtable}
$^{1}$好,我們在更加進入思想的時候.
我們一同再次有個祈禱先.
我們一齊祈禱.
我們一同有個聚會,一同敬拜你.
我們向你獻上感恩.
我們感恩因為主你一直的看顧,保守我們.
而在今天,我們也在這個時刻.
我們也很渴望去一齊思想你的懷疑.
就求主能夠幫助我們去明白,去理解.
也讓我們能夠應用在我們的生命當中.
以致我們的生命成為一個充滿感恩的生命.
願主能夠帶領,奉主耶穌成名祈求,Amen.
好,我們今天思想經文同樣是一直以來.
我們思想的一段書卷就是《申命記》的一段經文.
《申命記》的背景我想不用再多說.
因為我們一直都是在說《申命記》.
不過我想再一次去說.
《申命記》對我來說就是一個摩西最後的提醒.
或者是一個老人家的一個提醒.
就是不斷地重複.
如果像我奶奶那樣.
我奶奶很喜歡每一次見到我.
就不斷地提醒我不同的事情.
其實她是重複重複地提醒.
譬如說我還沒結婚的就問我何時結婚.
現在結了婚了.
如果結了婚的就會問我何時生小孩.
生完小孩還問我何時逼你弟弟結婚.
其實很多很多很多東西去說的.
但是她會不斷地重複.
每一次打電話都是說同樣的東西.
對我來說《申命記》同樣的情況發生.
為什麼呢?因為不斷地提醒的同一樣東西就是.
我們要去遵守上帝的誡命.
不斷地提醒的是我們要去敬畏上帝.
包括我們今天所讀的這篇經文當中.
如果我們見第一節一開始都是在說這件事.
就是在說我今天所吩咐的一切誡命.
你們要謹守遵行.
當我們讀完之後.

$^{41}$我們覺得整個主題就是要謹守遵行這件事.
當然我的主題是在說思想.
就是在說後面的東西.
但是當我們讀這篇經文的時候.
很快的第一個印象或者第一個主題.
是在說要謹守遵行上帝的誡命.
而在整本《申命記》當中.
其實基本上就是不斷地提醒.
如果大家來到這個聚會.
可能有一段時間這一年的話.
其實之前那幾章都是在提同樣的東西.
除了一至三章可能在說出埃及之後.
發生了很多不同的事情之後.
第四章開始就會見到摩西就是這樣提同樣的東西.
如果說到這裡的話.
不如我們聽回以前的講到最完.
反正都是同一個提醒.
但是不是的.
因為同一個提醒.
不過在不同的原因或者不同的狀況之下.
去作出這個提醒的時候.
那個效果是有不同的.
那個理解也會有不同.
所以我們要看的就是.
究竟同樣的提醒背後.
為什麼摩西有這樣的提醒.
為什麼他有什麼不同的狀況.
或者他在這裡又再出同樣的提醒的時候.
究竟想給我們什麼訊息.
或者在上帝話語當中.
我們又可以得到什麼呢.
所以我們就要繼續去看.
當然這裡第一節就像一個主題.
今天所吩咐的一切誡命.
很明顯不是在說摩西的誡命.
而是上帝所吩咐的誡命.
透過摩西去給一班以色列民.
當然我們這邊有兩個理解.
可能是之前一些所說的誡命.
又或者可能接著下來的誡命.

$^{81}$因為他可以是一個開頭.
無論是前面.
或者總結前面.
又或者說後面的時候.
都是說緊守遵行.
然後就說原因.
做什麼.
或者有什麼需要去留意的.
我們就會看之後的經文.
在之後的經文當中呢.
我們就會看到一個很重要的東西.
就是你要紀念耶和華.
你的神在曠野引導你這四十年.
當我們去思想.
或者去思考關於.
以色列民在摩西五經當中的歷史.
我們會知道其實.
以色列民在出埃及之後.
有四十年在曠野當中.
而這個是在說摩西.
等他們完了這四十年.
要進入美地之前.
所作出的提醒.
現在在他進去之前.
快要完這四十年之前.
就是在提醒以色列民.
你要紀念你的神.
在曠野引導你這四十年.
這四十年.
很長.
這四十年.
說的是很多不同的狀況.
如果我們從開頭到尾.
可以看到的就是.
沒有食物 沒有水.
又冷又熱.
有很多的戰爭.
有很多的辛苦的地方.
又要搬搬抬抬.
有空就要走走.

$^{121}$有空就要休息.
很多不同的改變.
一直在改變.
當他看到這些狀況的時候.
其實很多人都有不同的看法.
有些人可能看為好.
但也有不少人.
因為這麼辛苦.
我們會覺得這不是一件好事.
紀念不一定看為一件好事.
你可以記得很多東西.
我們記得很多以前的東西.
但很多時候.
我們記得以前很多的東西.
是一些辛苦的東西.
我們就覺得.
這個不是一件好事.
但是在這裡.
摩西專登要提醒這件事.
他說你要紀念四十年.
發生過什麼事.
不是.
但到最後是在說.
我想用一句話去總結.
就是在說第四節說的那句話.
這四十年你的衣服沒有穿破.
你的腳沒有腫.
簡單來說就是說.
你什麼事都沒有.
你過了這四十年.
是呀 很辛苦.
很多事情發生.
但你正正常常.
你甚至沒有一件衣服是穿爛的.
當然我沒有試過.
穿一件衣服穿四十年.
我暫時都沒有到四十歲.
但是我們可以想到.
當我一件衣服可以穿很久的時候.
其實是一件很奇妙的事.

$^{161}$很恩典的事.
而這樣去說的時候.
就是在告訴摩西.
你們是在經歷很多的恩典.
當思念這四十年.
是呀 很辛苦 很累.
但是記得全部是有恩典的.
接著下來.
除了紀念這件事之外.
也要思想.
你當心裡思想.
耶和華你神管教你好像人管教.
二字一樣.
看這些事.
不只是看著上帝給我恩典.
但要思想什麼呢.
我們要有一個觀念.
上帝對我們好像父子一樣.
好像父親和兒子和女兒一樣.
說的是有一個關係.
說的是這一切都是說父親的管教.
父親的管教是說.
希望孩子能夠健康成長.
在前兩天.
昨天是四十五.
今天是四十六.
昨天我和全家.
加上兒子去了健康院.
要打針.
因為四個月大了.
這樣他才四個月大.
要打針.
其實在打針的時候.
我們要見姑娘.
見姑娘的時候.
我當時不知道.
原來是有東西要去檢查.
有些東西給他看看.
會不會跟著他看.
左右.

$^{201}$之後原來有些東西要給他看.
會不會自己去捉住他.
跟著就說.
哎呀 你兒子很好 合家.
原來有一種很開心的感覺.
我沒試過.
但我想說的是.
當見到兒子.
原來這樣成長.
然後別人說很厲害.
很棒的.
但我想這就是父親的管教.
希望我們自己的孩子可以成長.
不是說到時他什麼都做不到.
我們當然會擔心.
如果他說不合格.
我也有想過.
他只要說我兒子不合格.
我會怎樣呢.
我會罵他嗎.
我還沒說.
沒發生.
全部都沒發生.
但感恩的是.
他說合格.
見到他成長.
我很開心.
同一時間.
我想兒子也很開心.
因為他玩得很開心.
看著.
又有時間跟那個姑娘說.
說得好像很多東西.
但這些反應.
是成長.
很有趣的.
不只是父親.
不只是上帝開心.
但當孩子成長.
他有自己的生命的時候.

$^{241}$希望他能夠好好地活著.
希望他有一個好的生命.
紀念思想是為了.
能夠讓這些以色列人明白.
當上帝讓你經過所有的一切.
是希望你有一個好的生命成長.
當然之後說的是.
紀念思想之後.
是做什麼.
不是說只記得.
不是只想.
但後面下一步說的是第六節.
你要謹守耶和華利神的誡命.
遵行祂的道.
敬畏祂.
我特意在這裡.
記下了謹守這兩個字.
什麼是謹守.
就是小心地去守.
去遵行上帝的誡命.
要看清楚.
所謂的謹慎地去遵守.
對我來說就是.
你要明白.
你要知道規矩是什麼.
我想跟我玩遊戲是不同的.
我玩遊戲很多時候.
我不理會規則是什麼.
他講了一大堆.
我一邊玩一邊問.
我想不是一個好的參加.
我經常都很喜歡break the rule.
稍後有games night.
大家不要這樣.
要學得謹守.
要懂得去謹守.
要明白規則.
對我來說就是.
當你玩那個遊戲的時候.
如果你說謹守的話.

$^{281}$就是你要明白所有規則.
而你玩的時候.
你是跟著那個規則的.
你不是去break the rule.
你不是不知道那個規則是什麼.
你要知道.
然後去遵守.
這個是我們要知道.
要去遵守.
要去行.
而紀念思想其實就是為了這個.
為什麼紀念思想為了我們能夠.
推動我們去遵行上帝的誡命.
上帝的吩咐.
上帝的話語.
是有事要去做出來.
當然除了紀念思想之外.
也有另一個原因是給了我們.
因為在之後那節也有說.
因為耶和華利神領你進了美地.
當我們領你進美地的時候.
其實就是.
當我們看.
可能我們以為享受了恩典.
但我們心水清的都知道.
在這段經文當中.
以色列文是沒有進入美地的.
所以這裡才說在那個地方.
在那個地方會有河.
那個地方有泉.
為什麼?因為他們正在走過去.
而在中文當中是沒有"test"的.
這裡的"領"字是在說.
領著你出去.
領著你過去.
帶著你過去.
是這樣的一個過程.
簡單來說.
上帝是帶隊的那個.
那你就要怎樣呢?.

$^{321}$黏著祂.
因為祂是領隊.
說起這件事.
我想起我祖母.
我祖母在.
大概十七歲左右.
話說我祖母就叫我和她一起去旅行.
去泰國.
她來了香港之後.
就和我一起去報一個泰國團.
就一起去了.
為什麼會找了我呢?.
一,因為我祖母是一個很憎氣的人.
剛才我也提到.
所以我爺爺就說.
你不要搞我了.
不要叫我陪祖母去.
你去吧.
他就和我這樣說.
然後我就去了.
和她去呢.
有一件事.
我很記得的.
他經常都在說.
領隊去哪裡.
我們就去哪裡吧.
所以有時候去到商場.
譬如說.
領隊就說.
你現在有兩個小時的自由時間.
你可以走一走.
我打算走.
十七歲.
剛剛出來旅行.
當然想看看.
他說不要啊.
跟著領隊啊.
一定要跟著領隊.
要不就在這裡等.
不要走開.

$^{361}$我祖母就是這樣跟團的.
所以她是沒有什麼地方可以去.
她只是真的.
當然.
跟團是可以離開領隊.
因為我們可以享受其他地方.
但是對以色列人來說.
雖然知道上帝是他們的領隊.
但是他們不可以離開上帝.
不可以遠離上帝.
因為一.
上帝正在帶他們去那個地方.
如果不跟著的話.
會迷路.
但同時.
也有生命的危險.
為什麼呢.
我們回頭去第一節看看.
如果在第一節的時候.
其實也在說一樣東西.
就是要謹守遵行.
謹守遵行之後.
那裡的原因是什麼呢.
好叫你們.
存活.
人數增多.
且進去.
德耶和華向你們列祖起誓.
應許的那地.
有沒有想過.
為什麼會說好叫你們存活.
是不是如果我不謹守.
我就沒有那個存活.
就是不會有命.
我們要想的context就是.
原來當時以色列民.
當他們要進入之前.
其實是要面對很多不同的國家.
而這些國家是什麼呢.
當探子去那些國家去看的時候.

$^{401}$這些人是很厲害的.
而在第一至第三章也曾經提及過.
當摩西去走的時候.
打過不同的國家.
是因為上帝.
才能夠擊敗這一班人.
所以如果不跟著上帝的話.
連命都沒有.
是提醒以色列民.
你為了能夠存活.
你能夠繼續享受那個恩典的話.
你一定要跟著上帝.
為什麼呢.
這樣你才能去到那個美地.
所以謹守遵行是確保我們可以緊貼上帝.
貼著上帝去走那條道路.
當然在後面的那節經文.
說的是帶我們去那個美地.
有河,有泉,有源,有很多不同的東西.
貼著上帝不只是說.
我們可以去到那個美地.
但後面是說我們可以經歷上帝的恩典.
這些恩典只要我們緊貼上帝.
我們就能夠經歷.
我們就可以去見到,去看到.
甚至去享受.
因為第十節也說的是吃得飽足.
吃得飽足是什麼呢.
就是享受了.
你已經入了肚.
剛才我們吃了很多東西.
我現在很飽了.
嘴上還有蒜味.
因為有蒜味.
很享受.
是一個這樣的享受.
因為你享受了那些恩典.
而當我們經歷了恩典之後.
就要稱頌耶和華.
因為祂將美地賜給我們.

$^{441}$所以我們要去感恩.
當然就這樣看整段經文.
我想很快就可以看到.
原來有紀念思想.
為什麼有紀念思想.
就是因為我們可以緊守遵行.
然後緊貼上帝,經歷恩典.
然後感恩.
如果以以色列民這群人的經歷來說.
他開頭記得四十年的經歷.
記得那個關係.
然後因為這樣應該推動我們.
要去遵行上帝的誡命.
然後緊貼上帝.
因為上帝是帶領上帝.
是那個保守者.
然後我們就可以經歷恩典.
上帝賜給我們美地.
然後吃得飽足.
感恩.
但是是不是就這樣停在那裡呢?.
是不是我們享受完.
去到那個美地之後.
吃得飽足就完呢?.
值得留意的地方.
不是只是停在那裡.
因為如果我們只是停在.
吃得飽足就完的話.
很危險的地方就是.
很危險的.
很容易會忘記耶和華.
所以我們說的是要紀念.
為什麼要紀念呢?.
我們就要回頭去看.
二至四節的這段經文.
為什麼我會選擇紀念思想.
作為今天的主題呢?.
是因為在這段經文當中.
其實對我來說.
我覺得有時候很難理解.

$^{481}$它叫你這樣去想.
但其實不是一件容易的事.
不過為什麼要這樣去想呢?.
為什麼說苦練你.
試驗你.
要知道你心內如何.
行守你的誡命.
不幸之後要你知道.
就是人我著不得.
不試探究聖物.
凌是靠耶和華口裡.
所出的一切話.
為什麼要用四十年.
去教這件事呢?.
不是因為為了我們四十年.
經歷完知道了就完.
是因為在之後的生命當中.
仍然要記得這件事.
等同於什麼呢?.
譬如說我們上一年班的時候.
可能學了一些知識.
不代表只是停留在那裡.
因為這些知識會帶去二年班.
繼續去用的.
而當進入那個尾地.
其實不是停了.
是應該繼續去遵守而去用的.
所以應該是一個循環.
如果不是的話.
我們說吃得飽足的話.
我們就很容易停了.
好像在之後的經文.
十一至十四節就說.
你要謹慎免得忘記.
耶和華你的神.
不守他的誡命.
典章律例.
就是我今天所吩咐你的.
恐怕你吃得飽足.
建造美好的房屋居住.

$^{521}$你的牛羊家多.
你的金銀增添.
並你所有的都傳到家增.
你就心高氣傲.
忘記耶和華你的神.
就是將你從埃及的圍爐之家.
凌出來的.
提醒的是.
當我們得著那個恩典.
感恩了一次.
就不是停在那裡.
那問題.
回到剛才那個表.
我們應該怎樣去理解這件事呢?.
接下來是甚麼呢?.
當我思想不斷去想的時候.
我發覺就是一個循環.
接下來是應該.
紀念思想.
我吃得飽足.
不是說以前那些不記得.
但當我得著一個新的恩典的時候.
我繼續思想這些恩典.
而去繼續.
為甚麼要這樣做呢?.
試想想.
如果四十年之後.
只是紀念那四十年.
那麼往後的日子.
是否只是不斷地去想以前呢?.
不可以的.
如果我們只是不斷地想當年.
很多時候.
其實我們不滿足現在有的東西.
我們不懂得感恩.
對我來說.
要不停地感恩.
是要不停地有新的理由.
但問題是怎樣有新的理由呢?.
當你紀念思想完以前的事.

$^{561}$謹守進行完上帝的吩咐.
然後緊貼神的時候.
你沒有經歷恩典.
你就感恩.
你沒有新的恩典.
你就去感恩.
感恩mode on.
不是只是一次感恩就完結.
而是不斷地不停地感恩.
所以對我來說.
當我見到這件事的時候.
我發覺原來是很重要的.
不是只是說我吃飽了就完結.
但當我經歷完新的恩典的時候.
我可以放在感恩的清單上.
讓我紀念上帝在我生命當中的作為.
思想到原來上帝就是我的天父.
然後不斷地推動我去謹守進行上帝的吩咐.
從而我可以貼著上帝去經歷上帝的恩典.
很多時候我們失去那個感恩.
不是因為我們沒有那個恩典.
反而是因為我們沒有跟著上帝去行.
沒有跟著上帝去走那條路.
因而我們就失去了那個經歷.
而我們不知道要對什麼去感恩.
所以今天我很想邀請大家做的一件事.
就是我們不如嘗試跟隨上帝.
去經歷上帝給我們的恩典.
或者可能今天未必有的.
我們由最初的那個恩典開始看.
主耶穌怎樣去救助你和我呢?.
主耶穌怎樣去幫助我們呢?.
那個恩典的救恩應該是令我們明白.
上帝有一個恩典.
我們就紀念思想.
主就是帶我去那個美地的主.
我就謹守進行祂的吩咐.
然後慢慢經歷祂的恩典.
不斷地經歷新的恩典.
才會令到我們不停地去感恩.

$^{601}$讓我們祈禱.
你是我們的天父.
你好好的帶領我們.
去保守我們.
去引領我們.
你的旨意是美好的.
你的旨意是奇妙的.
你都有豐盛的恩典.
讓我們可以不停地去經歷.
讓我們一直以來不斷地遠離你.
主耶穌宣我們.
幫助我們能夠再一次回到你的身邊.
今天我們想回想上帝天父.
在我們生命當中所做的作為.
我們見到很多的恩典.
我們再次思想到.
主是我們的主.
天父是我們的父親.
不斷地帶領我們.
這一切願能夠成為我們的動力.
去謹守進行主你的教導.
主你的吩咐.
讓我們能夠緊貼你.
讓我們在你的道路當中.
能夠經歷更多你所預備的恩典.
以致我們可以不停地向主稱頌感恩.
主啊,我們很希望我們有一個不停感恩的生命.
就求主能夠幫助我們每一位.
賜福我們每一個.
帶領我們每一個.
再次向你感恩.
我們能夠有思想懷疑的時間.
就願主能夠賜福我們,帶領我們.
繼續帶領我們的生命.
任我們不斷地向你感恩.
Amen.
\newpage



\section{創世記 4:1-16}
\label{sec:032240}
\textbf{勝過嫉妒}~~~~何志滌牧師
\newline
\newline
連結: \href{http://www.church.com.hk/acms/content.asp?site=cdc\&op=show\&type=product\&code=032240\&layout=sermon}{\texttt{032240}} ~~~~ 語音日期: 2026-05-31
\newline
\newline
\hyperref[sec:032239]{\small{< < < PREV SERMON < < <}}
~
\hyperref[sec:index_chronic]{\small{[返順時目]}}
~
\hyperref[sec:032241]{\small{> > > NEXT SERMON > > >}}
\newline
\newline
創世記 4:1-16
\newline
\begin{longtable}{cl}
\hline
\hline
章節 & 經文 (和合本修訂版)\\
\hline
4:1 & \begin{tabularx}{0.7\textwidth}{X} 那人和他妻子夏娃同房，夏娃就懷孕，生了該隱 ，她說：「我靠耶和華得了一個男的。」 \end{tabularx} \\ \\ \relax
4:2 & \begin{tabularx}{0.7\textwidth}{X} 她又生了該隱的弟弟亞伯。亞伯是牧羊的；該隱是耕地的。 \end{tabularx} \\ \\ \relax
4:3 & \begin{tabularx}{0.7\textwidth}{X} 過了一些日子，該隱拿地裡的出產為供物獻給耶和華； \end{tabularx} \\ \\ \relax
4:4 & \begin{tabularx}{0.7\textwidth}{X} 亞伯也把他羊群中頭生的和羊的脂肪獻上。耶和華看中了亞伯和他的供物， \end{tabularx} \\ \\ \relax
4:5 & \begin{tabularx}{0.7\textwidth}{X} 卻看不中該隱和他的供物。該隱就非常生氣，沉下臉來。 \end{tabularx} \\ \\ \relax
4:6 & \begin{tabularx}{0.7\textwidth}{X} 耶和華對該隱說：「你為甚麼生氣呢？你為甚麼沉下臉來呢？ \end{tabularx} \\ \\ \relax
4:7 & \begin{tabularx}{0.7\textwidth}{X} 你若做得對，豈不仰起頭來嗎？你若做得不對，罪就伏在門前。它想要控制你，你卻要制伏它。」 \end{tabularx} \\ \\ \relax
4:8 & \begin{tabularx}{0.7\textwidth}{X} 該隱與他弟弟亞伯說話 。 二人正在田間時，該隱起來攻擊他弟弟亞伯，把他殺了。 \end{tabularx} \\ \\ \relax
4:9 & \begin{tabularx}{0.7\textwidth}{X} 耶和華對該隱說：「你弟弟亞伯在哪裡？」他說：「我不知道！我豈是看守我弟弟的嗎？」 \end{tabularx} \\ \\ \relax
4:10 & \begin{tabularx}{0.7\textwidth}{X} 耶和華說：「你做了甚麼事呢？你弟弟血的聲音從地裡向我哀號。 \end{tabularx} \\ \\ \relax
4:11 & \begin{tabularx}{0.7\textwidth}{X} 現在你必從這地受詛咒，這地開了口，從你手裡接受你弟弟的血。 \end{tabularx} \\ \\ \relax
4:12 & \begin{tabularx}{0.7\textwidth}{X} 你耕種土地，它不再給你效力；你必流離飄蕩在地上。」 \end{tabularx} \\ \\ \relax
4:13 & \begin{tabularx}{0.7\textwidth}{X} 該隱對耶和華說：「我的懲罰太重，過於我所能承當的。 \end{tabularx} \\ \\ \relax
4:14 & \begin{tabularx}{0.7\textwidth}{X} 看哪，今日你趕我離開這塊土地，不能見你的面；我必流離飄蕩在地上，凡遇見我的必殺我。」 \end{tabularx} \\ \\ \relax
4:15 & \begin{tabularx}{0.7\textwidth}{X} 耶和華對他說：「既然如此 ，凡殺該隱的，必遭報七倍。」耶和華就給該隱立一個記號，免得人遇見他就殺他。 \end{tabularx} \\ \\ \relax
4:16 & \begin{tabularx}{0.7\textwidth}{X} 於是該隱離開了耶和華的面，去住在伊甸東邊挪得之地。 \end{tabularx} \\ \\
[1ex]
\hline
\hline
\end{longtable}
$^{1}$今天我講到的也是教會的安排.
到今天要講這個題目.
雖然教會也給了我們一個自由度.
他說如果你不想講也可以講第二個題目.
但是當我自己去思考這次這個安排的時候.
講到一樣.
我們很多時候都會發生在我們心靈裡面.
是關於一個題目叫做「疾倒」.
即是「倒忌」.
但當然這個題目是怎樣去勝過我們的疾倒.
我們看回聖經.
很多經文是講到家庭裡面.
兄弟之間的疾倒.
無論今天我們所講的.
該隱阿伯.
我們比較熟悉的在創世記裡面.
而同樣也都看到摩西和阿倫和米利安.
甚至聖經裡面的約瑟.
教會裡面的約瑟.
他的哥哥倒忌他自己的弟弟.
今天我要特別看該隱和阿伯.
我們可能會比較熟悉的一個內容.
就是當該隱和阿伯獻祭的時候.
神越納阿伯卻不越納該隱.
為什麼他不越納該隱呢?.
甚至到了後期.
該隱就殺了自己的弟弟阿伯.
但是上帝仍然給該隱有個保護.
給他有個恩基.
他說沒有人可以殺你.
當然當我自己去思考這個題目.
和這段聖經的時候.
我就回想回到我自己.
當我回來香港.
我1985年回來香港事奉.
因為去到大概1992,93年的時候.
大概七年的事奉之後.
我就發覺有一天.
我自己忽然之間不太想回教會.
我不太想事奉.

$^{41}$我覺得好像很不開心.
當時我又不太懂得.
什麼叫做沮喪,抑鬱.
我都不太懂這些東西.
我整天覺得信耶穌怎麼會抑鬱.
信耶穌怎麼會沮喪.
但是真的不想回教會.
去到一個地步.
我真的想離開這個牧養的位置.
但是當時我就去找一些求助.
去找一個心理輔導的一位姐妹.
這位姐妹很細心.
每次去見她.
她都叫我從後門進出.
不要讓人看到我去見心理輔導.
因為如果被會友看到我去看心理輔導.
就會覺得我的牧師瘋了.
很害怕去看心理輔導.
所以每次都是從後門進出.
沒有人知道我去看心理輔導.
但是我第一次去.
我一去到.
她沒有什麼話要說.
她就幫我祈禱.
一祈禱我就哭了.
哭了大概三十多分鐘.
哭完之後.
她說可以了你走吧.
你下星期再來見我.
我說怎麼搞的.
你還沒說話就祈禱了.
我哭完之後她叫我走.
她說可以了你下星期再來見我.
如是者十二個星期.
幸好是免費的.
不然我真的覺得不太划算.
我哭了十二次.
最後一次十二次的時候.
在我哭的時候.
忽然間神跟我說了一句話.

$^{81}$我覺得當時我還沒有這樣的經歷.
神什麼聲音都沒有.
但是我感受到一個聲音.
祂說你和你大哥比較.
你嫉妒你的大哥.
怎麼會啊.
我大哥比我大十多年.
他是一個對我非常好的大哥.
我第一份玩具都是我大哥買給我的.
所以我覺得沒有理由.
於是我就跟這位心理輔導.
就是這位醫生跟他說.
我好像覺得神跟我說.
我和大哥比較.
我嫉妒我大哥.
他說可以了.
你下星期不用回來.
你好了.
我說好了.
他說是的.
我說你不要說話.
神不要說話.
神說話就可以了.
你回去想想這句話究竟是什麼意思.
我真的回去想.
究竟這句話是什麼意思.
我總覺得大哥對我很好.
我從來沒有想過我會妒忌我大哥.
但是當我真的想的時候.
我就發覺原來我大哥很厲害.
我大哥讀書一直都是考第一名.
甚至在上個世紀50年代.
他能夠考進台灣國防醫學院做醫生.
我覺得他簡直是很厲害的一個人.
在我自己心目中.
到他畢業那年在台灣.
他是全台灣所有醫學生的考試.
他是考第二名.
真的很厲害.
所以我一直覺得他真的很厲害.

$^{121}$其實我內心就跟他有些比較.
後來我才發覺.
原來我不止跟我大哥比較.
因為這個影響.
其實我經常跟其他教會比較.
大家明白我的意思嗎.
你知道我們教會旁邊有間.
在柏屯山的大學教會有很多人.
在那個年代我們教會有百多人.
講到教會就要排隊進去.
排隊進去我們只有百多人.
我們請都沒有人來.
他們就排隊進去.
差很遠的.
我覺得我又跟他們比較.
有時候見到外面的講員去講道.
為什麼請他們不請我.
原來我真的有些比較.
但是上帝就因為那一次.
我就開始真的在神的面前求主.
讓我去將這個比較的心.
求神慢慢慢慢將它沒有了.
其實我大哥對我真的很好.
過去十幾二十年他有時候會回來香港.
當然現在他年紀很大.
他都不會坐飛機.
他每次回來香港都來同福聚會.
每次同福聚會完之後.
他都會跟我說.
弟弟你都很厲害.
我說有什麼厲害.
你是醫師我是牧師.
雖然都有一個師一樣.
但是你是醫師我是牧師.
差很遠的.
這個世界是我的感覺.
他說不是.
你每次星期日都有幾千人聽你講話.
我在醫院只有五個人聽我講話.
其實他給我很大的安慰.

$^{161}$但是我覺得這個安慰.
就是那一次神已經給了我安慰.
他說你不要去比較.
你做回上帝要你做的事情.
原來家庭.
原來兄弟之間.
不只是戈隱爸爸.
不只是摩西和亞倫米利亞.
也不是約瑟和他們哥哥.
其實原來家庭那邊很多時候.
都會有這一種.
好像在心那邊一種的疾度.
當然我們還沒有去到.
好像戈隱這麼嚴重的地步.
於是我今天就透過這段聖經.
希望和大家想一想.
人為什麼會有些疾度.
但是怎樣去勝過這個疾度.
第一.
這段聖經提醒我們.
人的疾度的原因.
就是覺得神不公平.
但是我們怎樣去回應神不公平.
讓我們去勝過.
我就提醒先自我反省.
為什麼這樣說.
很多解經家就說.
神為什麼不約立戈隱的事情.
只是約立爸爸的事情.
很明顯.
戈隱是一個所謂農夫.
他去獻農作物.
而爸爸就是畜牧業.
去獻羊羔.
很多解經家就說.
當然他有血.
他一定沒有血.
他應該把農作物換成羊去獻.
就有血.
或者你看看 爸爸是獻最好的.

$^{201}$是獻頭身的羊.
但是我就想.
農作物怎樣去分頭身和不頭身.
一起出你怎知道哪一棵最頭身.
我就覺得.
他們這樣解經不一定是錯.
我不敢說他們錯.
但是我都會想.
為什麼神不約立戈隱.
就因為神和戈隱說了一句話.
他說第一.
戈隱怎會為你煩惱.
第二.
他沒有等戈隱去回應.
神就和他說.
如果你做得對.
我怎會不怨濫你.
大家明白我的意思嗎.
如果你做得對.
我絕對不會不怨濫你.
於是這句話讓我去反思.
為什麼神會和戈隱說.
如果你做得對.
我一定會怨濫你.
相對之下.
因為上帝就要戈隱去想一想他自己.
獻祭的問題.
獻祭的問題就是說.
你獻祭什麼.
當然.
我們希望是將最好的獻祭給神.
什麼叫做最好的.
在我們自己的定義那邊.
我深信.
戈隱都可能會覺得.
他自己獻給神是最好的.
都不奇怪.
但是神反而要戈隱去想.
他說我體重.
不只是你奉獻什麼.

$^{241}$我體重是你的心.
你奉獻的心.
加上你日常生活裡邊的行為.
你想一想.
各位頂智媒.
我經常都說.
在教會做基督徒很容易.
因為教會百分之九十幾是基督徒.
但是你去到世界就不同.
世界百分之九十幾不是基督徒.
換言之你在教會做一個基督徒.
跟在世界那邊做一個基督徒.
是一個很不同的回應.
在外面真的難很多.
有時候我開玩笑.
就是說你在外面借飯.
一旦吃完飯你碟飯不見了.
我們會這樣想.
所以我今天又要賣一個廣告.
這個廣告就是我生意上的一個QR Code.
是真的不是假的.
不要以為你拍了它有東西看.
我這個廣告就是說.
怎樣在日常生活裡邊.
成為一個活祭.
我這個課程就是講這件事.
怎樣在日常生活裡邊做一個活祭.
換言之我們真的要想.
不只是看在教會.
甚至我今天去想.
我做牧師做了幾十年.
我都會想究竟我做牧師.
我的心態是什麼.
是不是我覺得做牧師威一點.
是不是我自己覺得我挺好的.
我能夠令到教會有增長.
是不是這樣.
我都要想究竟我自己做一個牧者.
是回應上帝的召命.
還是我自己的私慾.

$^{281}$希望有些身份有些地位.
弟兄姊妹神要該人去想一想.
祂說如果你做得對.
沒有理由我不原諒你.
當然在同福內的弟兄姊妹知道.
我曾經二十幾年前公開說過.
我說我決定紅燈一定不會過馬路.
二十幾年前已經說了.
其實原來紅燈不過馬路是很難的.
你不要以為容易.
有一次我在中環.
星期一在中環.
你知道星期天中環沒什麼人.
我在地鐵出來.
在阿里山那邊太子巷那邊.
站在那裡.
想著紅燈不要過馬路.
等了好像很久.
又沒有人又沒有車.
其實應該過的.
沒有人沒有車.
我很想過的時候.
突然聖靈提醒我.
祂說你說過的.
站在那裡.
然後看到對面有位女士.
那位女士在對面.
她就站在那裡.
我就開始想.
是不是今天是星期天.
紅燈不會轉呢.
你知道有時候會這樣.
星期天沒有車.
可能不會轉.
於是又想過.
但星期天提醒我祂站在那裡.
我又不要過.
但去到最後就覺得.
是不是紅燈壞了.
人就會這樣想.

$^{321}$當我再想過的時候.
我又站在那裡.
其實很傻.
但去到一個地步.
終於紅燈變黑燈.
我只講三個字.
對面那位女士又過馬路.
中間她看著我.
何牧師.
嘩 多害怕.
嘩.
嘩 真的多害怕.
如果這位姐妹是回唐福堂的.
她就會這樣.
你講了天下無敵.
做了什麼.
有心無力.
嘩 真是這樣.
自己的心很軟.
有一次在時代廣場後面.
你知道時代廣場後面.
有一條馬路很窄.
有紅燈.
你不過馬路會被撞死.
有一次我覺得.
沒理由不過.
這麼多人都過.
但當我將過的時候.
後面有人拍拍我肩膀.
何牧師.
我就站在那裡.
原來紅燈不過馬路是很難的.
但有一次我在上海講道.
上海應該沒有人認識我.
我想著要上很多人.
過馬路.
我又想過馬路.
後面有人拍拍我.
何牧師.
我又站在那裡.

$^{361}$原來真的很難.
經過一年.
我終於覺得可以.
現在你怎樣推我.
我都不會紅燈過馬路.
經過一年的掙扎.
但如果一年你不去堅持這個掙扎.
你今天一樣.
我會紅燈過馬路.
當然現在大家比較害怕.
現在隨時會被警察抓.
好像要罰二千元.
你跟鄰居說小心點.
二千元給教會的那個給他.
即是說.
原來是不容易的.
我們日常生活裡面.
一件這麼細微的事.
其實都不容易.
神跟個人說.
你想一想.
如果你做得對.
我沒有理由.
不會越納你的欠債.
第二.
我們聽不到神的聲音.
相對我們要學.
去聆聽神的聲音.
第一次會.
神跟個人說.
如果你要做得對.
你就會被越納.
但如果拒絕做對的事.
你就要小心.
你做得對我會越納.
但你要想.
你肯不肯去做對的事.
我們很多時候.
神啊.
你不要跟我說話.

$^{401}$我聽不到你的聲音.
各位聽眾.
不需要這些神秘的.
因為整部聖經都是神的話.
當你肯讀聖經的時候.
就已經有神的話.
所以你千萬不要說.
沒有神的話.
千萬不要說沒有神的話.
因為整部聖經就是神的話.
就算我們今天開一個課程.
聆聽神的聲音.
我們都要聽完之後.
在神的話語那邊.
去看一看究竟我聽的東西.
有沒有違背聖經的原則.
換言之我們大前提.
我們要好好地去讀.
神的話語.
所以報了名的百多七百的人.
好好地從現在開始.
上課之前去讀聖經.
讀神的話.
我相信神.
對你說話.
一定不會違背聖經的原則.
以致到.
很多時候我們聽不到.
其實不是我們聽不到.
是我們不想聽.
所以神跟各人說.
如果你做得對.
我怎會不原諒.
但是你想一想.
為什麼你不肯做對的事.
大家明白我的意思嗎.
為什麼不去做對的事.
所以當我們去思考上帝.
給我們的說話的時候.
我們就要知道.

$^{441}$我有沒有好好地去讀神的話.
有沒有每一天.
早上去讀神的話.
哪怕我們每一天.
只讀一次聖經.
讀一次聖經.
你一年可以讀多少次聖經.
365次聖經.
換言之.
365次聖經一年.
就能夠在你的心靈裡.
提醒你.
所以聖經也說.
神阻擋驕傲的人.
你信不信.
當然我不會說自己謙卑.
我真的不敢說自己謙卑.
有時很多人.
突然會稱讚你的時候.
你也會覺得我挺厲害.
我在5月16日.
去澳門港獨.
晚上有個報道會.
我真的沒想過.
那晚的報道會在澳門.
在澳門出席人數.
有900人.
在澳門.
全澳門只有4000個基督徒.
只有4000個基督徒.
有900人來這個報道會.
就算當中有700個基督徒.
當中有700人.
當時我自己也飄飄然.
是不是我的名氣.
也會這樣想.
但不是.
因為有兩個個性.
一個是中年很聲音的個性.
又開始.

$^{481}$讓我謙卑.
不關我的事.
是他們兩個.
但真的.
當有人稱讚你的時候.
你的心.
也會有種.
很有趣的.
一種飄飄然.
我講完之後.
有個姊妹走過來.
我不認識她.
她不是說我講道好.
她說何牧師你唱歌也挺好.
那兩個個性在這裡.
我唱歌怎麼會好聽.
她說你挺好聽.
我又飄飄然.
原來我也等同那兩個個性.
我是不是應該去參加中年好聲音.
第一姊妹.
人就是這樣.
但聖經說什麼.
聖經說神阻擋驕傲的人.
所以我講五堂教導的篇.
我有一堂.
是怎樣去接受.
人對你的稱讚.
很多時候我們都會面對.
被人罵.
當然不開心.
被人稱讚怎樣去接受.
但當然去到最後.
受恨不去處理.
會帶來殺戮.
所以我們一定不可以.
讓殺戮.
留有任何餘地.
那個人聽完神說.
他沒有聽.

$^{521}$他聽不到.
他沒有去處理.
自己的那種煩惱.
因為神跟他說.
為什麼你要煩惱.
他沒有去處理.
於是他就想了一個辦法.
將爸爸殺了.
第一姊妹.
聖經曾經說過.
耶穌說.
你有沒有殺人.
當然沒有.
殺人還坐在那裡.
還這麼安然坐在那裡.
完全沒有.
但聖經說.
你有沒有試過發怒.
你有沒有試過.
對弟兄姊妹.
有一種仇恨.
約翰一書.
第三章第五章說.
如愛仇恨.
恨你的弟兄.
等同殺了弟兄.
怎麼搞的.
我們一起讀一讀.
馬太福音第五章.
這是新普及.
日本的法譯.
我覺得很有意思.
我們一起讀一讀第五章.
你們都聽過.
告誡我們祖先的話.
不可殺人.
殺人者要受審判.
可是我要說.
你們就是對人發怒.
也要受審判.

$^{561}$如果罵別人是白痴.
就要被帶上審判庭.
如果咒詛別人.
別懷疑他是傻.
他是白痴.
千萬不要說這種話.
說人白痴都會.
都會受審判.
當然.
難不倒他.
無所謂.
聖經也說過.
不可含怒.
不是說我們不可以生氣.
但我們要處理他.
處理我們的生氣.
一樣.
該人沒有處理.
他的憤怒.
去到最後.
就成為了一個.
殺人的動物.
其實不只是聖經說的.
很多坊間的心理學.
都會這樣說.
在網上有一個資訊說.
仇恨為什麼可以演變成.
殺機.
他說第一.
情緒長期的累積.
不會處理.
然後就殺了他.
第二.
失去理性.
其實理性地聽.
沒理由殺人.
但當你失去理性.
他就會殺人.
第三.
意識形態的極化.

$^{601}$為什麼這個社會.
這麼多行動.
有時行動去到兩極.
可能積聚了很多仇恨.
帶來一些.
行為上的表現.
弟兄姊妹.
我們千萬不要.
看輕.
看輕這個質度.
千萬不要看輕這個質度.
覺得我只是生他的氣.
有什麼所謂.
生他的氣就知道.
但你生他的氣.
不去處理這個.
生他的氣的時候.
真言.
真言很厲害.
有一句說話.
也在《新普及譯本》裡面說.
猜疑.
質度.
如同.
骨中的毒瘤.
嘩.
我們不要以為.
質度一下.
有什麼所謂.
如果我不和.
我哥哥那次.
那種內心的處理.
可能我還會和.
很多人比較.
弟兄姊妹.
我要說的不是今天的我.
我要說的是二十幾三十年前的我.
如果我不去處理.
當時我的內心.
這種質度.

$^{641}$我不知道今天會變成怎樣.
我也不敢說.
我今天會變成怎樣.
但上帝有祂的醫治.
有祂的提醒.
有祂的保護.
以致到.
我沒有讓.
這件事在我心裡.
繼續去.
積聚.
成為.
骨癌.
嚴格來說.
這就是骨癌.
骨中的癌症.
弟兄姊妹.
臣.
吾越立該人.
是要他去想一想.
換言之.
神體寵.
奉獻的人.
是多過.
你奉獻什麼.
神體寵.
人先.
祭物.
是後.
換言之今天我們在神教那邊.
上帝要我們去.
反省.
但我們很深信我們一定勝過.
為什麼?.
因為我們有神在.
我們有神在.
今天我們是信耶穌的人.
我們有神在.
問題是我們是否願意去依靠祂.
是否願意讓祂去主宰.

$^{681}$我們的人生.
所以我跟各位說.
不要被這個怒氣.
去控制你.
你要去.
將這個怒氣釋放.
以致你能夠.
去經歷到上帝.
更大更大更大的.
恩典.
所以今天.
我欣賞弟兄姊妹回到神的家.
欣賞弟兄姊妹.
願意去奉獻.
但是我覺得我們的神.
是更加欣賞我們.
這個人.
所以我們.
看著我們的家庭裡面.
每一個成員.
我們要學習.
去欣賞祂.
欣賞你的配偶.
欣賞你的上一代.
欣賞你的下一代.
我都說現在.
我有三個孫的時候.
我就覺得.
真是一個很大的挑戰.
三個男孫.
我正在上班.
大家不要以為我現在沒事做.
我正在上班.
但是我每天早上要照顧孫子.
兩個小時.
兩個小時真是一個很大的挑戰.
一個兩歲多.
一個四歲.
一個五歲.
抽兩個小時不是容易的.

$^{721}$跟他一起坐下.
他一秒鐘可以起床.
我就要一分鐘才起床.
你怎樣追他呢.
其實是一個很大的挑戰.
很多的學習.
怎樣去愛他們.
怎樣去欣賞.
怎樣去欣賞我們每一個的.
家中.
任何一個成員.
他一定有.
他的長處.
也有可能有他的弱點.
我們會不會.
為他的長處去感恩.
而為他的弱點去祈禱.
這個就是我們一生要學.
所以弟兄姊妹.
我們要記得神.
是一種奉獻的人.
所以我們是信徒.
加祭司.
我們每一個信徒.
在神的眼中.
都是祭司.
所以如果你在旁邊說是信徒.
你跟他說你是祭司.
你有祭司的身份.
但是祭司在舊約是利未人.
利未人他們不只是住在.
利未人自己中間.
他們要在不同的支派裡面.
分開住.
他們在教會.
在聖殿.
在會幕.
但是他們平時住在.
不同的宗族的中間.
因為他們平時的行為.

$^{761}$是更加重要.
所以我昨天.
去了一個晚宴.
我從來沒有經歷過.
這種晚宴.
能夠見到特首.
我覺得自己很了不起.
我從來都在電視.
見到她說話.
從來沒有見過真人說話.
我初初去到.
沒有人認識.
一群都是商界的人.
我都不太認識那些人.
但後來坐了一下.
突然有人拍我.
都有人認識我.
都開心的.
原來都有很多基督徒.
原來有很多基督徒.
在商界裡面.
我進不了商界.
但這群人進到商界.
其中一個WhatsApp.
他WhatsApp給我.
我們有很多基督徒在這個界別當中.
我們可以在這個界別裡面.
去為神做更多的事.
Amen.
所以不要看不起商界.
沒有基督徒.
很多基督徒.
所以如果.
我們能夠好好地.
在我們的生命裡面.
在我們的行為裡面.
哪怕我們只是一個家庭的主婦.
你所接觸的人.
可能多過我都不一定.
以致到.

$^{801}$不要看不起自己.
我們每一個人.
都是祭司.
我們每一個人都是事奉的人.
Amen.
願意去回應.
上帝給我的.
照明.
和給你的照明.
\newpage



\section{路加福音 17:1-10}
\label{sec:032241}
\textbf{門徒要順服}~~~~俞大鳴牧師
\newline
\newline
連結: \href{http://www.church.com.hk/acms/content.asp?site=cdc\&op=show\&type=product\&code=032241\&layout=sermon}{\texttt{032241}} ~~~~ 語音日期: 2026-05-24
\newline
\newline
\hyperref[sec:032240]{\small{< < < PREV SERMON < < <}}
~
\hyperref[sec:index_chronic]{\small{[返順時目]}}
~
\hyperref[sec:032242]{\small{> > > NEXT SERMON > > >}}
\newline
\newline
路加福音 17:1-10
\newline
\begin{longtable}{cl}
\hline
\hline
章節 & 經文 (和合本修訂版)\\
\hline
17:1 & \begin{tabularx}{0.7\textwidth}{X} 耶穌又對門徒說：「絆倒人的事是免不了的，但那絆倒人的有禍了！ \end{tabularx} \\ \\ \relax
17:2 & \begin{tabularx}{0.7\textwidth}{X} 人若把這些小子中的一個絆倒的，還不如把磨石拴在他的頸項上，丟在海裡。 \end{tabularx} \\ \\ \relax
17:3 & \begin{tabularx}{0.7\textwidth}{X} 你們要謹慎！若是你的弟兄犯罪，就勸戒他；他若懊悔，就饒恕他。 \end{tabularx} \\ \\ \relax
17:4 & \begin{tabularx}{0.7\textwidth}{X} 如果他一天七次得罪你，又七次回頭，說：『我懊悔了』，你總要饒恕他。」 \end{tabularx} \\ \\ \relax
17:5 & \begin{tabularx}{0.7\textwidth}{X} 使徒對主說：「請加增我們的信心。」 \end{tabularx} \\ \\ \relax
17:6 & \begin{tabularx}{0.7\textwidth}{X} 主說：「你們若有信心像一粒芥菜種，就是對這棵桑樹說：『你要連根拔起，栽在海裡』，它也會聽從你們。」 \end{tabularx} \\ \\ \relax
17:7 & \begin{tabularx}{0.7\textwidth}{X} 「你們當中誰有僕人耕地或是放羊，從田裡回來，就對他說『你快來坐下吃飯』呢？ \end{tabularx} \\ \\ \relax
17:8 & \begin{tabularx}{0.7\textwidth}{X} 他豈不對僕人說『你給我預備晚飯，束上帶子伺候我，等我吃喝完了，你才可以吃喝』嗎？ \end{tabularx} \\ \\ \relax
17:9 & \begin{tabularx}{0.7\textwidth}{X} 僕人照所吩咐的去做，主人還謝謝他嗎？ \end{tabularx} \\ \\ \relax
17:10 & \begin{tabularx}{0.7\textwidth}{X} 這樣，你們做完了一切所吩咐的，要說：『我們是無用的僕人，所做的本是我們該做的。』」 \end{tabularx} \\ \\
[1ex]
\hline
\hline
\end{longtable}
$^{1}$我們一起祈禱.
讓我們有你的清晨.
我們一同來敬拜你.
一同從你的話語中學習.
怎樣來跟隨主耶穌.
求聖靈幫助我們.
讓我們明白主耶穌教訓的意義.
共與耶穌基督一同祈禱.
阿們.
最近我們一連書上.
主耶穌所說的幾個比喻.
一方面我們知道那些比喻是.
斥責和提醒發利財人.
另一方面都是讓門徒有所學習.
因為我們要記得.
路加福音一開始.
就告訴我們.
這本書是寫給一個名叫.
提亞非羅大人.
提亞非羅是他的名字.
大人是路加對他的稱呼.
這個大人大概是一位政治家.
並且他是相當富裕.
以至他可以出錢請路加醫生.
去做詳細的調查.
研究過有關耶穌基督的事之後.
寫成書來報告他調查的結果.
這個提亞非羅.
是那麼願意出錢做這件事.
我相信他是一位信徒.
因為一章四節說.
目的是讓你知道.
你所學的道是正確的.
路加其實很喜歡用道這個字.
形容主耶穌所傳講的道理.
和他所做過的事.
我們看小覽轉.
特別多提耶穌的信仰是道.
由於提亞非羅是一位有錢的信徒.
路加似乎有意去教導他對金錢的看法.

$^{41}$故事從十二章到十六章.
有四個比喻.
都是和錢有關的.
包括無知財主的比喻.
基隱管家的比喻.
和上星期講的財主和拉達魯的比喻.
都是提醒有錢的信徒.
要把握機會.
善用金錢.
因為金錢是會過期無效的.
今天的經文.
剛才聽了一次.
主耶穌直接對門徒的教導.
作為主耶穌的信徒.
我們不只要接受耶穌做救主.
其實更加要接受耶穌做我們生命的主.
我們要去接受和進行耶穌的教導.
因為在《大使命》裡面講得很清楚.
「你們要去.
洗萬民作為門徒.
奉父子聖靈的名給他們施洗.
凡我所吩咐你們的.
都教導他們進守」.
所以我們不只要傳福音.
傳福音只不過是第一步.
我們要帶領人受洗加入教會.
之後還要花很長的時間.
耐心地去教導門徒.
去明白耶穌的教訓.
使信徒可以遵守.
這樣我們才算是進行了大使命.
今天的十節經文裡面.
其實提到幾個教訓.
如果我們單單讀一次或聽一次.
就好像覺得這幾個教訓互不相關.
不過如果詳細些看的話.
我可以總結它是.
「愛信望」或我們習慣說的.
「信望愛」都是有關的.
或者我自己對信望愛比較敏感.

$^{81}$有一天早上開車送師母去火車站的時候.
看到前面的車牌.
有三個英文字母.
「Heavy Gel」.
我就告訴師父「很特別」.
師父說「有什麼特別呢?」.
「Heavy Gel」代表信望愛.
Faith, Hope, Love.
我們現在來思考耶穌的教導.
一到三節主耶穌就面對世界.
警告門徒說.
「患倒人的事是免不了的」.
這裡好像是強調一些事情.
是一些事情使人犯罪.
但是人和事是不能夠分開的.
如果沒有人就不會有事情發生.
中國人說「人事人事,有人才有事」.
並且在下半節主耶穌說.
「但那叛倒人的有禍了」.
可見這一個教導是要提醒門徒.
行事為人要小心.
不要做一些會使人犯罪的事.
什麼事情會使人犯罪呢?.
要回答這個問題我們要想一想.
犯罪的事是指什麼呢?.
由於主耶穌在第二節說.
「凡日把這小子中的一個叛倒的.
還不如把磨石.
山在他的頸岸上,釣在海裡」.
我相信這個罪是非常嚴重的罪.
應該是使人不信上帝的罪.
在聖經記載的歷史故事當中.
其實都有很多使人犯罪.
不信上帝的例子.
舉一個例子.
在以色列人出了埃及之後.
其實幾個月後他們已經去到.
迦南地的對面.
可以進入迦南地.
當時摩西派了十二位密探.

$^{121}$去探路和了解地形的情況.
然後才進入迦南地.
但是那十二位探子當中.
探完地回來報告.
其中有十個探子.
煽動以色列人不信上帝.
還說不如我們回埃及.
打死摩西算了.
另外選一個領袖出來.
帶他們回埃及.
結果是怎樣呢?.
上帝很憤怒.
要整代以色列人都死在曠野.
這些十位探子不信上帝.
他們用言語去鼓勵其他人.
都不要信上帝.
這些就是使人叛島的罪.
他們死有餘辜.
在第一世紀時代是怎樣呢?.
保羅曾經討論過一個問題.
就是可不可以吃祭過偶像的肉.
我不知道大家的成長背景如何.
我的成長背景.
我媽媽經常都拜祖宗.
要煮些菜出來.
那些算不算拜偶像的肉呢?.
我作為一個信徒.
可不可以吃呢?.
保羅在這裡有一個很詳細的教導.
我讀一段經文給大家聽.
經文很清楚.
不用多解釋.
食物不能使我們更接近上帝.
因為我們不吃無損.
吃也無益.
可是你們要謹慎.
免得你們的自由.
竟成了軟弱人的絆腳石.
若有人見你這有知識的.
在偶像的廟裡坐直.

$^{161}$而這人的良心是軟弱的.
他豈不放膽去吃那祭過偶像的食物嗎?.
因此基督為他死的那軟弱弟兄.
也就因你的知識沉淪了.
你們這樣得罪弟兄.
傷了他們軟弱的良心.
就是得罪基督.
所以食物若使我的弟兄跌倒.
我就永遠不吃肉.
免得我的弟兄跌倒了.
在我們這個時代.
在這個城市.
某教會也曾經發生過一件事.
有一位教會領袖的太太.
帶領了一群來教會參加活動的朋友.
去參觀藍天寺.
她不是帶人去拜佛.
是帶人去一個旅遊觀光點.
但是有其他教會的領袖很不滿意.
要求那位丈夫.
不是丈夫帶去,是丈夫的妻子帶去.
但是他們要求丈夫要辭去教會領袖的地方.
其實他們是在應用保羅這個教訓.
雖然我們有自由去觀光.
但是作為教會領袖的太太.
帶人去參觀藍天寺可能會令人跌倒.
在另一間教會的青年領袖.
就曾經有這樣發生一件事.
他是青年領袖,青年人,還在大學裡面讀書.
他預備在大學裡面搞一個打麻將大賽.
發揚中國的國粹.
結果被牧師叫他來.
跟他談了很久,勸他不要搞這些活動.
為什麼呢?又不是賭錢,我們只是比賽.
西人也是這樣比賽,打橋牌也是這樣比賽.
為什麼打麻將不行呢?只是比賽.
結果牧師勸他不要搞.
因為這樣可能會令人跌倒.
打麻將贏了之後可能日後會沉迷於打麻將.
這樣就不太好了.

$^{201}$如果我們相信耶穌的話.
西人犯罪是會對別人造成永恆的後果.
西人永遠跟上帝不和.
我們就要小心自己的行為.
其實這是一種愛心的表現.
使我們小心行事為人.
不要令人犯罪.
講完要小心不要令人跌倒之後.
耶穌又教他們門徒.
彼此之間一定要彼此饒恕.
人與人之間相處一定會碰碰撞撞.
有時候有些誤會,矛盾.
甚至發怒也是很平常的.
在教會裡面也會發生.
因為教會也是由人組成的.
教會就不是天堂.
所以不要期望來到教會見到信徒.
每個人頭上好像有個光圈.
不是的.
我時時都說教會只是一個罪人的療養院.
我們每個人都有些病在療養.
故此耶穌也教門徒要學習,要饒恕人.
我曾經看過一本書.
裡面也提到怎樣處理憤怒.
內容也提到饒恕.
還引述了一位英國精神病院院長的話.
他說如果在他那間精神病院裡面的病人.
能夠有被饒恕的把握.
不是饒恕.
是被人饒恕的把握.
大概有一半人將來可以出院.
饒恕人和被饒恕都很重要.
如果我們確信上帝.
已經赦免了我們的罪.
我們也要確信我們被上帝饒恕.
這樣就會有一個很奇妙的心理治療的功效.
在第三節,第四節提到兩件事.
首先第三節是說有弟兄姊妹犯罪.
那個罪不一定是得罪你自己個人的.
如果有人知道弟兄姊妹犯罪.

$^{241}$譬如我們知道某個弟兄姊妹搞婚外情.
知道犯罪不一定是得罪我們自己的.
我們不要隻眼開隻眼閉.
假裝不知道.
也不要到處去跟別人說.
你知道那個弟兄姊妹犯了罪.
不要到處去宣揚他的罪.
耶穌教導要知情的人.
應該直接去勸戒那位犯罪的弟兄.
如果他悔改.
我們就饒恕他,接納他.
如果他不悔改怎麼辦?.
不悔改怎麼辦?.
在這裡沒有說清楚.
不悔改怎麼辦?.
但在《馬太福音》十八章十五到十七節的教導.
有說得清楚一點.
我建議大家可以對他說.
可不可以一起去約牧師輔導.
假設你已經跟他說了他不肯悔改.
我們可不可以一起去找牧師輔導呢?.
如果他同意就約牧師一起談談.
如果他不肯你可以再問他.
你介不介意我把這件事告訴牧師.
等牧師跟你談談.
如果你覺得弟兄所犯的罪是嚴重的話.
無論他介不介意你都要告訴牧師.
等牧師決定如何處理.
牧師一般都會替人保守隱私.
但如果犯了社會上的法律的話.
牧師都要向警方報告.
最近有新聞提到一位總督.
曾經是基督教的牧師.
做到很高級.
但他處理那些對兒童犯侵犯罪的人.
處理得不好結果都要辭職.
這些是嚴重的罪.
你沒有犯罪而不是直接得罪你.
都要去處理.
第四節就是一位弟兄犯罪直接得罪你.

$^{281}$怎麼辦呢?.
那位弟兄後來又來向你道歉.
主人都說你要饒恕他.
饒恕是很容易說的.
但又很難實行.
為什麼呢?.
因為我們是人.
人是有情緒的.
當我們被人得罪的時候.
我相信最快出現的情緒就是生氣.
這樣說我都可以.
很生氣.
故此我們要先處理憤怒.
然後平心靜氣之後.
才有可能做得到饒恕.
今天我沒有時間和大家詳細講饒恕的步驟.
其實要慢慢做.
大約有七個步驟.
如果你想知道詳細的方法.
弟兄姊妹你散會後來和我說.
我曾經講過一篇講道的題目.
叫做教我如何不恨他.
裡面就提到我們怎樣去饒恕人.
有七個步驟.
如果有人想聽的話.
可以告訴我.
我會發一個連結給你.
隨時聽都可以.
你喜歡睡覺前聽都可以.
什麼時候聽都可以.
小心自己的行為.
不要讓別人知道.
當別人得罪我們.
又向我們道歉的時候.
我們要饒恕.
其實這些可以總結為愛倫如己.
這是愛倫如己的實踐.
那些跟隨耶穌的使徒.
其實知道很難做.
很難做主耶穌的教導.

$^{321}$於是就求主加增他們的信心.
要加多些信心給我們.
否則怎樣做呢?.
你說的道理很難做.
其實這是一個很好的模範.
如果我們缺乏信心.
去照著主的教導去做的話.
我們真的要祈禱.
求主加添給我們信心.
當我們對主耶穌有信心的時候.
我們就有機會做到主耶穌的教訓.
主耶穌的回答是一個很明顯的誇張法.
「你們若有像一粒芥菜種子大小的信心.
就是對這棵桑樹說.
連根拔起來它栽在海裡 它也會聽從你們」.
芥菜種子在耶穌的時代被認為是最小的種子.
耶穌這句話的意思就是強調信修果孝.
是會出人意表的.
在澳洲很多人的家都有前後院.
有時都要斬樹.
其實都很麻煩的斬樹.
幾年前我都要出錢.
不是單單靠祈禱.
自己斬自己.
你自己去搬走了.
不是的.
我們都要出錢找人幫我們斬那棵樹.
要用了幾百元的.
我們不可以真的按著字面這樣說.
算了算了.
你快點連根拔起.
自己丟到邊邊去收草的箱子.
不可能做到這樣.
這個說法是一個誇張的說法.
這個說法有什麼意思呢.
其實信心的大小是很難量度的.
不過信心的對象就會決定信心的果孝.
對象很重要.
我曾經聽過一個故事.
有一隻蠍子想過河.

$^{361}$那條河都很寬.
他跳不過河.
又想過河怎麼辦呢.
他見到一隻青蛙.
跟青蛙說青蛙青蛙.
拜託你了我想過河.
你載我一程行不行.
青蛙說你是蠍子.
你是有毒的.
我載你不容易被你刺傷.
媽媽都告訴我.
叫我們要遠離你們.
蠍子對他說.
是正常的情形.
不過現在我想你載我過河.
我刺你幹什麼.
你載我過河.
我一定不會刺你的.
青蛙想了一想.
又覺得那隻蠍子說得有點道理.
況且媽媽又教訓過他.
助人為快樂之本.
於是小青蛙答應了.
好 媽媽 載你吧.
你上我背上吧.
蠍子就上去了.
到了河的對岸的時候.
蠍子都可以上岸了.
青蛙突然間覺得.
背上被蠍子刺了一下.
全身就劇痛.
他就問蠍子.
你為什麼刺我.
蠍子說我已經過河了.
況且我是蠍子.
要刺你是我的本性.
這個故事教導我們.
信心的對象是很重要的.
你搞錯對象.
可能有生命危險.

$^{401}$有了正確的對象.
然後有完全的信心.
我們就可以造成大事.
歷史上有很多信徒.
都是憑信心.
對上帝的信心.
來為上帝造成大事.
記得搞孤兒院那個.
慕勒先生.
大家都可能聽過他的故事.
有一天幾百個小朋友.
沒有食物了.
但是慕勒先生說.
你們照樣坐在餐桌上.
祈禱.
慕勒先生祈完禱之後.
就有一輛貨車.
送了一大車的麵包來.
於是那些孤兒就有得吃了.
信心的對象.
當我們有正確信心的對象.
憑信心去做上帝要我們做的事的時候.
神蹟就會發生.
故事的第六節的意思大概告訴門徒.
如果你們對我有信心.
願意遵從我的說話去跟人相處.
你們學習開口.
對那位向你道歉的弟兄姊妹說.
我原諒你 我饒恕你.
這樣就會帶來好的後果.
不要用信心不足做藉口.
也不要愛在心裡口難開.
要用信心實踐耶穌的教導.
就好像一位僕人聽從主人的吩咐辦事一樣.
我們就會有好的人際關係.
講了愛心信心.
第七節到第十節.
主耶穌就講了一個比喻提醒門徒.
假如有人做了耶穌的教導.
也不要驕傲.

$^{441}$以為自己真的厲害.
反而仍然要像一個僕人.
有謙卑信服信仇的態度.
對自己說我們是無用的僕人.
所做的只不過是我們應該做的.
其實主耶穌還講了很多有關主人回來.
會審判僕人的比喻.
其實我們都要盼望主會再回來.
所以主耶穌的教導是會帶來祝福的.
即使我們為了盼望得福而遵行主的教導都是好的.
但當我們有了祝福之後.
我們要小心不要自高.
以為靠自己的才幹而做得到.
以致我們能夠順利辦事.
沒有讓人跌倒也能夠饒恕那些得罪我們的弟兄姊妹和朋友.
使我們有和諧的人際關係.
時時都活得開心不激氣.
我總結一下今天經文所學到的功課.
就是在愛心,信心和盼望方面的實踐.
我們要對主耶穌有信靠順服的心.
進行他的教導.
要實踐愛倫如己的誡明.
與非基督徒或初信的弟兄姊妹相處.
我們要小心自己的行為.
不要使別人跌倒或不信耶穌.
只要有可能令人跌倒的行為.
我們自己都要小心.
有很多人不肯信耶穌.
其中一個最慣常聽到的理由是什麼呢?.
那個人說是基督徒.
行為比我還差.
怎麼信呢?.
所以我們真的要小心自己的行為.
別人看得到的.
與弟兄姊妹相處.
難免有時候有矛盾發生.
我們要學習向別人道歉.
要承認自己過犯.
被得罪的人就要學習饒恕.
讓我們的友誼建立在真理上.

$^{481}$以致我們有真正和諧的關係.
最後我們要記得主耶穌回來的時候.
或者我們去見主面的時候.
是會有審判.
我們要盼望得到主耶穌的稱讚.
認我們是又良善又忠心的僕人.
是遵從祂命令的僕人.
弟兄姊妹.
大家有沒有對其他弟兄姊妹有偏見.
或者有時候要蓄意避開某個弟兄姊妹呢?.
如果有的話.
我鼓勵你要馴服主耶穌的教導.
去跟得罪你的弟兄姊妹真誠對話.
尋求和解.
準備好饒恕別人.
和請求別人饒恕自己.
我們一起祈禱.
\newpage



\section{路加福音 17:11-19}
\label{sec:032242}
\textbf{感恩}~~~~俞大鳴牧師
\newline
\newline
連結: \href{http://www.church.com.hk/acms/content.asp?site=cdc\&op=show\&type=product\&code=032242\&layout=sermon}{\texttt{032242}} ~~~~ 語音日期: 2026-05-31
\newline
\newline
\hyperref[sec:032241]{\small{< < < PREV SERMON < < <}}
~
\hyperref[sec:index_chronic]{\small{[返順時目]}}
~
\hyperref[sec:032243]{\small{> > > NEXT SERMON > > >}}
\newline
\newline
路加福音 17:11-19
\newline
\begin{longtable}{cl}
\hline
\hline
章節 & 經文 (和合本修訂版)\\
\hline
17:11 & \begin{tabularx}{0.7\textwidth}{X} 耶穌往耶路撒冷去，經過撒瑪利亞和加利利中間的地區。 \end{tabularx} \\ \\ \relax
17:12 & \begin{tabularx}{0.7\textwidth}{X} 他進入一個村子，有十個痲瘋病人迎面而來，遠遠地站著， \end{tabularx} \\ \\ \relax
17:13 & \begin{tabularx}{0.7\textwidth}{X} 高聲說：「耶穌，老師啊，可憐我們吧！」 \end{tabularx} \\ \\ \relax
17:14 & \begin{tabularx}{0.7\textwidth}{X} 耶穌看見，就對他們說：「你們去，把身體給祭司檢查。」他們正去的時候就潔淨了。 \end{tabularx} \\ \\ \relax
17:15 & \begin{tabularx}{0.7\textwidth}{X} 其中有一個見自己已經好了，就回來大聲歸榮耀給神， \end{tabularx} \\ \\ \relax
17:16 & \begin{tabularx}{0.7\textwidth}{X} 又俯伏在耶穌腳前感謝他。這人是撒瑪利亞人。 \end{tabularx} \\ \\ \relax
17:17 & \begin{tabularx}{0.7\textwidth}{X} 耶穌回答說：「潔淨了的不是十個人嗎？那九個在哪裡呢？ \end{tabularx} \\ \\ \relax
17:18 & \begin{tabularx}{0.7\textwidth}{X} 除了這外族人，再沒有別人回來歸榮耀給神嗎？」 \end{tabularx} \\ \\ \relax
17:19 & \begin{tabularx}{0.7\textwidth}{X} 於是他對那人說：「起來，走吧，你的信救了你！」 \end{tabularx} \\ \\
[1ex]
\hline
\hline
\end{longtable}
$^{1}$我們先祈禱一下.
讓我們在主教的精神.
可以一齊來敬拜你.
去思想你的話語.
祈福你的聲明在我們心中動工.
幫助我們明白.
今天讀的經文.
其中的訊息.
幫助我們也懂得事神感恩.
奉耶穌基督的名祈禱.
Amen.
今天的主題是感恩.
我就想起.
美國有一個感恩節.
在每年的十一月.
二十八號.
都是美國感恩節的日子.
我就想.
為什麼美國人會搞一個.
感恩節呢.
原來感恩節.
最初是源自歐洲的.
農業傳統.
農民在秋收之後.
就舉行宗教儀式.
這種感恩的習俗.
隨著那些移民.
傳至美洲.
在1621年.
有一些清教徒.
在英國的清教徒.
其實他們都是基督徒.
但是當時英國的國家.
是聖公會.
他們就不太喜歡聖公會.
所以他們比較篤信聖經.
他們認為我們一切.
應該以神的名義.
來祈禱.
他們認為我們一切.

$^{41}$應該以聖經為本.
按聖經辦事.
這群清教徒就坐船.
五月花號.
很著名的船.
橫到大西洋.
去到他們認為的新世界.
他們當年.
就去到美國的.
馬薩諸塞州.
他們去到的時候.
其實是當地的.
冬季.
當時其實是很慘的.
第一代的移民是很慘的.
很多人是飢寒交迫.
隨著水土不服.
染病身亡的都有.
後來.
他們在當地.
原住民.
有一個部落.
叫做.
沃姆伯洛亞格.
的部落的幫助下.
他們就學習了.
當地人種玉米.
和捕魚.
一年之後.
他們的農作物.
有一個大豐收.
和原住民.
對他們的幫助.
和原住民一齊.
舉辦了三天的.
感恩宴會.
這個被認為.
美國感恩節的開端.
因此感恩節.
和人恩典的節日.

$^{81}$一般人.
都會有家庭團聚.
一起享用.
感恩的晚餐.
一定吃火雞.
表達對生命.
家庭朋友幫助.
與豐收的感恩.
所以.
我們教會其實每年.
選擇十一月都有感恩主日.
感恩主日.
都可能受美國的文化影響.
最近.
我看了一個.
發生在中國的小鎮.
一個小鎮的故事.
小鎮上有一個.
老張的修鞋像.
他.
舉腳行動不便.
在鎮上的一條巷.
擺攤.
擺了很多年.
一個下雨的晚上.
他收攤預備回家的時候.
突然看到一個.
五六歲的男孩.
捲縮在一個屋簷下.
被雨淋濕了身.
老張.
把這個小孩帶回家.
給他穿一件乾衣服.
煮一碗麵給他吃.
那個小孩.
叫小婉.
那時他爸爸媽媽.
去了外地.
打工.
跟著奶奶生活.

$^{121}$他下午和奶奶.
逛街的時候.
走散了迷路.
不知回家.
老張在鎮上周圍打聽.
終於找到他奶奶.
住在那裡.
將這個小孩送回家.
二十年後.
老張的修鞋譜還在.
但是歲月不饒人.
老張.
不知道是否患了.
白金遜症.
手開始顫抖.
當然是大不如前.
有一日傍晚.
有個年輕人.
拿著一雙鞋子.
放在他攤子上.
請他幫忙.
修理.
老張拿起這雙鞋子.
已經手顫抖.
拿針拿線.
穿針都穿不到.
沒辦法.
就跟年輕人說.
對不起.
年輕人.
我現在眼睛又矇.
修不好.
你去找其他人幫你修.
年輕人笑一笑.
說不要緊.
你慢慢修.
他蹲下.
跟老張穿線.
這個老張.
見到這個情景.

$^{161}$突然不知為何.
想起二十年前.
他也曾蹲下.
幫一個小孩.
問他.
你是否小院.
年輕人.
當時.
都很感動.
對呀.
張爺爺找了你幾年.
後來聽人說.
你還是在這擺攤.
從今日開始.
我來幫你.
華人很著重感恩.
有句說話說.
得人恩典千年記.
真正的感恩.
從來不需要華麗的承諾.
只是在某個.
普通的日子.
當你替我.
擋過雨.
如今我為你撐一把傘.
今日我們繼續看.
《路加福音》.
一個主耶穌向人.
施恩.
的故事.
在上星期.
我們看過主耶穌.
教導門徒.
有關信望愛.
的實踐.
我們看到.
這件事跡.
與信望愛有關.
而最後主耶穌.
亦提醒門徒.

$^{201}$要學習感恩.
我們先來.
看看這個故事.
然後再思想.
這個故事.
要教導我們的事.
在第十一節.
告訴我們這件事.
發生在哪裡.
在撒瑪尼亞和嘉利利.
中間的地區.
這個形容.
我看起來.
覺得有點奇怪.
我特意找一幅地圖.
出來看.
讓大家了解.
主耶穌.
出來傳道.
大部份時間.
其實在北面.
嘉利利的地方.
附近傳福音.
祂的總部.
應該在嘉伯龍.
如果我們看.
嘉利利湖.
最北面.
就是嘉伯龍.
那裡可以說是祂的總部.
所以如果我們看福音書.
看到祂回到嘉伯龍.
就知道了.
如果祂要到耶路撒冷.
耶路撒冷在南邊.
如果大家看地圖.
看到了.
往左手那邊看.
就找到耶路撒冷了.
祂這個旅程.

$^{241}$是由北向南走.
如果由北向南走.
為什麼不說.
經過嘉利利和薩瑪利亞.
中間的地區.
而要先說.
經過薩瑪利亞和嘉利利的地區呢?.
所以.
這就是奇怪的地方.
專家.
研究了之後.
很多專家都覺得奇怪.
為什麼要先提薩瑪利亞.
他們就.
得出一個結論.
他說經過的地方.
是薩瑪利亞和嘉利利交界的地方.
交界的地方.
也不奇怪.
但也應該先說嘉利利.
交界.
也是先經過嘉利利才進入.
薩瑪利亞.
但可能由於.
交界有點斜.
去了薩瑪利亞.
然後向東走.
所以就說.
去了薩瑪利亞和嘉利利的地區.
大概.
在地圖上.
藍色的地方.
去到那裡.
進入一個村子.
這個村子有沒有名字呢?我不知道.
聖經沒有提.
就說有十個麻瘋病人.
迎面而來.
遠遠地站著.
麻瘋病.

$^{281}$在耶穌的時代.
是屬於一種.
不治之症.
現在.
已經可以.
治療了.
已經可以靠藥物治療.
但當時就沒有.
醫藥沒有那麼發達.
病人不單止.
身體患病.
其實他們可能最大的.
痛苦是心靈方面.
是非常痛苦的.
一方面.
他們會被.
其他人歧視.
認為他們患了麻瘋病.
一定是犯罪的.
結果.
你得罪了上帝.
被上帝懲罰.
以致你患了麻瘋病.
最近.
我們《寓語查經》都在查.
藥百記.
藥百的三位朋友.
一致認為.
藥百是犯了罪.
你不犯罪.
怎麼會被上帝罰得那麼慘.
怎麼會有那麼悽慘的遭遇.
不可能,你一定是犯了罪.
所以他們就.
認為藥百犯了罪.
他們是朋友.
不是去罵他.
勸藥百.
快點認罪悔改.
你認罪悔改了就好了.

$^{321}$上帝的福音藥立即回來給你.
其實這些.
都是一般猶太人的神學觀念.
因此.
患了麻瘋病的人.
都會被人這樣看.
是罪人.
不是罪人怎麼會患麻瘋病.
此外.
根據摩西的.
律法.
患了麻瘋病的人.
要怎樣.
要被隔離.
他們知道.
可能上帝知道.
麻瘋病會傳染.
所以患了麻瘋病的人.
是不可以在社區裡面.
生活.
所以我們看到.
雖然耶穌進入了一個村子.
那些麻瘋病人.
應該不是住在那條村裡面.
他們平常.
可能是住在離開村.
比較偏遠的地方.
是沒有其他人住.
麻瘋病人就住在一起.
大家都患了麻瘋.
不怕傳染.
所以那些麻瘋病人.
是沒有正常的社交生活.
到安息日想去會堂.
聽聖經.
是不可以的.
沒有人歡迎你.
所以麻瘋病人要另外.
找地方生活.
假如麻瘋病人.

$^{361}$真的有些特別事.
不進村是辦不到.
一定要進村.
進村的話.
要怎樣做呢.
進村要一路走.
一路大聲叫.
麻瘋病人來了.
讓人們避開他.
不要接近他.
所以在那個時代.
患了麻瘋病.
是令人失去.
人性尊嚴.
他們.
這十個患了麻瘋病的人.
來找耶穌.
醫治他們.
他們都不敢太接近耶穌.
他們要遠遠地.
站著.
不要惹到耶穌有病.
這是他們的規矩.
不能夠跟人接近.
不過這十位麻瘋病人.
其實.
為何會走來.
見耶穌呢.
很可能他們已經聽過.
有關耶穌.
能夠醫病.
趕鬼,行神蹟的消息.
甚至.
他們可能會有家人.
知道耶穌.
是要去耶路撒冷.
應該會經過.
那條村.
於是那十個病人.
求醫深切.

$^{401}$一早就要在附近.
等.
他們的做法.
我想起一首詩歌的歌詞.
詩歌說.
生病的人.
會不會拒絕健康.
憂傷的人.
會不會拒絕安慰.
孤單的人.
會不會拒絕.
同伴.
從下面的經文記載.
我們知道這十個人當中.
其中有一位.
是撒瑪利亞人.
另外九位大概就是.
猶太人.
猶太人其實一般很歧視.
撒瑪利亞人.
但是當他們.
患上麻瘋病的時候.
與其.
孤單的來生活.
不如多一位同伴.
和同伴一起生活.
大家都是麻瘋病人.
大家都不會彼此歧視.
所以.
他們不會說.
你是撒瑪利亞人.
我們不侵略你.
他們和撒瑪利亞人.
一起生活.
他們.
很可能.
等了一段時間.
盼望耶穌的來臨.
終於看到耶穌來了.
一起.

$^{441}$向耶穌走過去.
他們遠遠.
站著.
因為他們是麻瘋病人.
不敢接近那些沒有麻瘋病的人.
他們.
也很可能.
早就約定.
我們見到的時候.
一個人叫.
可能聽不到.
要遠遠站著.
我們一起大聲叫.
耶穌老師可憐我們吧.
大聲叫耶穌.
就沒有什麼問題.
雖然他們.
以前可能沒見過耶穌.
這次只是第一次見到耶穌.
但是.
也可能有人告訴他們.
耶穌一般會穿的服飾.
以前.
服飾不是.
像我們今天有很多套.
隨時可以替換.
每天可以穿.
每天都可以穿不同的服飾.
不是的.
以前的人很少穿服飾.
不是很多服飾.
告訴他們耶穌穿什麼服飾.
而且耶穌穿的服飾.
可能是穿當時.
老師級別的人.
穿的服飾.
所以有點特別.
雖然沒見過耶穌.
也猜到他是耶穌.
可能有人告訴他們.

$^{481}$耶穌一般穿什麼服飾.
還有門徒.
耶穌不會一個人走.
一定是帶著門徒.
所以他們一見到.
就大聲叫耶穌.
因為他們相信.
這是他們唯一的機會.
他們真的很盼望.
自己能夠脫苦海.
重拾健康.
但是他們大聲喊.
大聲叫.
其實是表現出.
他們有信心.
這是他們對耶穌.
信心的表現.
所以稱呼耶穌.
沒問題.
但為什麼稱呼耶穌為老師呢?.
他們不是耶穌學生.
為什麼要稱呼他為老師呢?.
我相信.
或者我們當中.
有人知道.
猶太人.
稱呼那些懂得勝經的人.
我們聖經也叫老師.
在希伯來文.
或者亞蘭文.
也叫拉比.
這個字.
是可以譯作老師.
但在這裡.
這節經文的老師.
就不用拉比這個字.
用了另外一個希臘文的字.
這個字.
本來是用在軍隊裡面.
有長官的意思.

$^{521}$不是老師.
是官長.
有權發命令的人.
大部分英文譯本.
都不是譯作老師.
譯作Master.
有師父的意思.
大師.
師父.
路加醫生.
特別喜歡用這個字.
可能由於.
他的對象是希臘.
對象是希臘.
人的背景.
希臘文化背景.
記不記得路加醫生寫給誰.
提亞菲羅大人.
一聽名字就知道.
不是猶太人的名字.
他是一個有希臘文化背景的人.
所以他就用了.
這個字.
但這個字其實是比較.
尊敬人的.
稱呼.
稱呼人做Master.
師父.
以前我在香港的時候.
我們搭巴士.
如果想找司機.
跟他說幾句話.
就稱呼他司機.
或者客氣一點.
加多司機先生.
但後來我回香港的時候.
有人告訴我.
你不要稱呼他做巴士司機.
你稱呼他做巴士司機.
他會覺得不開心.

$^{561}$不高興.
我很尊重他.
他是巴士司機.
又不稱呼他做巴士司機.
要怎樣稱呼他才對呢.
你要稱呼他做車長.
奇怪.
怎麼叫車長.
可能我們稱呼那些.
駕駛飛機的人.
不是叫他做機師.
是機長.
所以駕駛巴士.
要稱呼他做車長.
就對了.
這些長大麻瘋的人.
稱呼耶穌做師父.
也反映出.
耶穌其實有權柄.
可以醫治他們.
這些麻瘋病人.
真的想.
耶穌醫治他們.
他們大聲喊.
耶穌老師可憐我們吧.
然後耶穌又沒有問他們.
你們想我幫你做什麼.
其實耶穌是很清楚明白.
這些人.
要求他們憐憫的意思.
就是想他們.
醫好他們的麻瘋病.
耶穌二話不說.
就吩咐他們.
你們去將身體.
給祭司檢查.
可能我們有些人.
以為.
祭司只是負責獻祭.
但如果我們讀.

$^{601}$利未記十四章.
就會知道.
原來祭司有很多工作.
除了獻祭.
還要去檢查病人.
第一.
你是否真的染了麻瘋病.
不單是人.
屋內如果發霉.
都要找祭司檢查.
屋內是否有麻瘋病.
類似的病.
第二.
麻瘋病人如果說.
我好了.
不用去隔離生活.
我要回到城裡住.
都要去見祭司.
祭司說你好.
你自己說好.
是沒有用的.
所以當耶穌吩咐.
他們要去給祭司檢查.
就是說.
可以醫好你們.
所以耶穌的吩咐.
其實包含了.
答應他們的要求.
憐憫他們.
醫治他們.
然後經文繼續說.
他們正去的時候.
就潔淨了.
他們聽到耶穌的吩咐.
不是大家先看看自己.
好了沒有.
我幫你看看你好了沒有.
不是的.
他們是馬上遵循耶穌的吩咐.
去見祭司.

$^{641}$而是一路走的時候.
一路發覺.
好了.
我的皮膚病好了.
我就想起.
在舊約時代.
當以色列人.
去到迦南地要過約丹河.
進入迦南地的時候.
有一段經文.
是這樣記載的.
在《約書下記》.
三章十五十六節.
那時正是收割的日子.
約丹河的水.
漲滿兩岸.
台約櫃的人.
到了約丹河.
台約櫃的祭司.
腳一入水邊.
那從上往下流的水.
就在很遠的地方.
在撒拉旦旁邊的.
亞當城那裡停住.
樹立城裡.
那往下拉巴海.
就是嚴海下流的水.
傳言中斷.
於是百姓在.
耶利哥的對面.
過了河.
神蹟是在我們.
憑信心.
有行動的時候.
就會出現.
發生果效.
這十位麻風病人.
開始的時候.
其實他們的信心.
可能只是.

$^{681}$聽人說.
耶穌是厲害的.
有能力趕鬼醫病.
他們聽到這些.
他們也非常盼望.
自己的病.
能夠得到醫治.
於是就準備好.
求耶穌憐憫他們.
然後他們對耶穌.
吩咐他們的話.
是有信心.
對耶穌有信心.
對耶穌說的話有信心.
願意進行的吩咐.
而當他們.
準備去見祭司的時候.
皮膚病好了.
他不是皮膚病.
好了之後就去找祭司.
而是他一路走的時候.
幾憤意.
發現皮膚病好了.
在.
十七章.
路加福音同一章.
耶穌教導門徒.
說.
你們若有.
信心像一粒芥菜種.
就是對這棵桑樹說.
你要連根拔起.
栽在海裡.
他也會聽從你們.
這裡所行的神蹟.
其實就是給門徒看到.
信心.
是可以使到麻瘋病人.
從病患中得到醫治.
從絕望中.

$^{721}$得到拯救.
這些麻瘋病人.
可以脫離.
他們那個.
非人化的生活.
重新投入社會.
過正常人的生活.
讓他們得醫治的.
盼望實現了.
最後.
有一個撒瑪利亞人.
還大聲.
歸榮耀給上帝.
其餘九個.
急著去見祭司.
其餘九個.
為什麼這麼急呢?.
我想.
可能他們要走比較遠的路.
他們要去耶路撒冷.
去找值班的祭司.
而撒瑪利亞人.
就不會去猶太那裡.
找祭司.
可能去撒瑪利亞.
他們的祭司.
但是.
他們是充滿了感恩.
這個猶太人.
這個猶太人的師父耶穌.
竟然都肯醫他們.
由於他們這樣做.
耶穌就加強了他們的信心.
「你的信救了你」.
當然.
這位撒瑪利亞人.
其實心裡面很清楚明白.
單單他個人的信心.
是不能夠使他得救的.
而是他信心的對象.

$^{761}$是正確的.
他對耶穌有信心.
再加上耶穌的憐憫.
和能力.
才能夠使他得醫治.
不過同時.
耶穌也很感慨地說了一句話.
「除了自己的外族人」.
耶穌沒有當他是自己人.
這個撒瑪利亞人.
外族人.
「再沒有別人回來歸榮耀給上帝嗎?」.
耶穌說這句話.
是要提醒門徒.
要學習感恩.
他用這位外族人.
來和那九位猶太人.
做一個對比.
這也是故意高抬這位外族人.
對上帝的心.
表揚這位外族人.
我們記得.
路加醫生寫給.
這位外族人.
外族人信神沒有問題.
感恩就行了.
從這段經文裡面.
讓我們學習三件事.
第一.
從那些患了麻風病的人.
我們可以學習.
盼望和信心.
盼望就不是.
坐在那裡.
甚麼都不用做.
而是我們會繼續留意.
世界上發生了甚麼事.
信心.
最重要是有正確的對象.
假如我們.

$^{801}$信錯了對象.
我們的盼望.
都會變成失望.
第二.
我們從耶穌學習.
怎樣愛人.
耶穌有時.
專門走某一條路.
其實都可能是有心.
去幫助那裡有需要的人.
就好像約翰福音.
第四章.
記載耶穌.
要返加利利.
那次的路程和這次相反.
他是從大地返回.
北面的加利利.
他是故意.
經過撒瑪尼亞.
去救一個女人.
那個女人的見證.
救了很多撒瑪尼亞人.
這次耶穌是要去耶路撒冷.
要經過.
撒瑪尼亞和加利利的邊界.
但他進入.
一條村莊.
就可能是故意的.
他也要了十位.
長大麻瘋的人.
他沒有問他們.
你們是甚麼人.
你們想我怎樣憐憫你們.
然後就揀出來.
你九個猶太人.
好吧,我依你,你是撒瑪尼亞人.
你回家吧,我不依你.
沒有.
耶穌愛所有.
需要得醫治的人.

$^{841}$那些來求救的人.
他不會先問他們的背景.
然後才決定要依誰.
第三.
最後一點.
我們要學那位撒瑪尼亞人的.
感恩的心.
我們的生命.
是上帝所賜的.
我們擁有的一切.
都是從上帝而來的.
最近主音穆斯.
告訴我.
教會今年三月.
收到有兩筆比較大的.
奉獻.
都是華人堂的會友奉獻.
藉著奉獻.
向上帝表達感恩是好的.
是值得鼓勵的.
但是我們.
還有其他的方式.
可以去表達.
感恩之情.
今天我想和大家分享.
一段經文.
記載在歌羅西書三章.
十二到十七節.
保羅是這樣教導我們的.
所以你們記住.
上帝的選民.
聖見蒙愛的人.
我中間跳了很多節.
當用.
各樣的智慧.
把基督的道.
豐豐富富的存在心裡.
用詩篇.
讚美詩,靈歌.
彼此教導.

$^{881}$互相勸誡.
以感恩的心.
歌頌上帝.
你們無論做什麼.
或說話.
或行事.
都要奉主耶穌的名.
藉著祂.
我鼓勵大家學習.
我們閒談的時候.
要分享.
大家最近在主律裡面.
有沒有學習到聖經.
有沒有對某段經文.
有新的了解.
很感恩.
有沒有領受過什麼恩典.
經歷到什麼真理.
保羅在這裡說.
要彼此鼓勵.
彼此勸勉.
也要一起歌頌上帝.
我們一起祈禱.
讓我們認識主耶穌.
藉著信祂而得到.
兒女的名分.
和天上各樣屬靈的福氣.
求聖靈帶領我們.
事實上.
也要樂意和別人分享.
我們從主所領受的福氣.
奉耶穌基督的名.
祈禱 阿們.
\newpage



\section{創世記 26:12-33}
\label{sec:032243}
\textbf{行在「恩約」中}~~~~秦民德傳道
\newline
\newline
連結: \href{http://www.church.com.hk/acms/content.asp?site=cdc\&op=show\&type=product\&code=032243\&layout=sermon}{\texttt{032243}} ~~~~ 語音日期: 2026-05-30
\newline
\newline
\hyperref[sec:032242]{\small{< < < PREV SERMON < < <}}
~
\hyperref[sec:index_chronic]{\small{[返順時目]}}
~
\hyperref[sec:032244]{\small{> > > NEXT SERMON > > >}}
\newline
\newline
創世記 26:12-33
\newline
\begin{longtable}{cl}
\hline
\hline
章節 & 經文 (和合本修訂版)\\
\hline
26:12 & \begin{tabularx}{0.7\textwidth}{X} 以撒在那地耕種，那一年有百倍的收成。耶和華賜福給他， \end{tabularx} \\ \\ \relax
26:13 & \begin{tabularx}{0.7\textwidth}{X} 他就發達，日漸昌盛，成了大富翁。 \end{tabularx} \\ \\ \relax
26:14 & \begin{tabularx}{0.7\textwidth}{X} 他有羊群牛群，又有許多僕人，非利士人就嫉妒他。 \end{tabularx} \\ \\ \relax
26:15 & \begin{tabularx}{0.7\textwidth}{X} 他父親亞伯拉罕在世的時候，他父親的僕人所挖的井，非利士人全都塞住，填滿了土。 \end{tabularx} \\ \\ \relax
26:16 & \begin{tabularx}{0.7\textwidth}{X} 亞比米勒對以撒說：「你離開我們去吧，因為你比我們強盛得多。」 \end{tabularx} \\ \\ \relax
26:17 & \begin{tabularx}{0.7\textwidth}{X} 以撒就離開那裡，在基拉耳谷支搭帳棚，住在那裡。 \end{tabularx} \\ \\ \relax
26:18 & \begin{tabularx}{0.7\textwidth}{X} 他父親亞伯拉罕在世的時候所挖的水井，在亞伯拉罕死後，都被非利士人塞住了，以撒就重新把井挖出來，仍照他父親所取的名為它們命名。 \end{tabularx} \\ \\ \relax
26:19 & \begin{tabularx}{0.7\textwidth}{X} 以撒的僕人在谷中挖井，就在那裡得了一口活水井。 \end{tabularx} \\ \\ \relax
26:20 & \begin{tabularx}{0.7\textwidth}{X} 基拉耳的牧人與以撒的牧人相爭，說：「這水是我們的。」以撒就給那井起名叫埃色，因為他們和他相爭。 \end{tabularx} \\ \\ \relax
26:21 & \begin{tabularx}{0.7\textwidth}{X} 以撒的僕人又挖了一口井，他們又為這井相爭，以撒就給這井起名叫西提拿。 \end{tabularx} \\ \\ \relax
26:22 & \begin{tabularx}{0.7\textwidth}{X} 以撒離開那裡，又挖了一口井，他們不再為這井相爭了，他就給那井起名叫利河伯。他說：「耶和華現在給我們寬闊之地，我們必在這地興旺。」 \end{tabularx} \\ \\ \relax
26:23 & \begin{tabularx}{0.7\textwidth}{X} 以撒從那裡上別是巴去。 \end{tabularx} \\ \\ \relax
26:24 & \begin{tabularx}{0.7\textwidth}{X} 當夜耶和華向他顯現，說：「我是你父親亞伯拉罕的神。不要懼怕，因為我與你同在，要賜福給你，也要為我僕人亞伯拉罕的緣故，使你的後裔增多。」 \end{tabularx} \\ \\ \relax
26:25 & \begin{tabularx}{0.7\textwidth}{X} 以撒就在那裡築了一座壇，求告耶和華的名，並且在那裡支搭帳棚；他的僕人就在那裡挖了一口井。 \end{tabularx} \\ \\ \relax
26:26 & \begin{tabularx}{0.7\textwidth}{X} 亞比米勒同他的顧問亞戶撒和他軍隊的元帥非各，從基拉耳來到以撒那裡。 \end{tabularx} \\ \\ \relax
26:27 & \begin{tabularx}{0.7\textwidth}{X} 以撒對他們說：「你們既然恨我，趕我離開你們，為甚麼又到我這裡來呢？」 \end{tabularx} \\ \\ \relax
26:28 & \begin{tabularx}{0.7\textwidth}{X} 他們說：「我們明明看見耶和華與你同在；因此就說，讓我們雙方彼此起誓，我們跟你立約， \end{tabularx} \\ \\ \relax
26:29 & \begin{tabularx}{0.7\textwidth}{X} 使你不加害我們，正如我們未曾侵犯你，素來善待你，並且送你平平安安地走。你是蒙耶和華賜福的！」 \end{tabularx} \\ \\ \relax
26:30 & \begin{tabularx}{0.7\textwidth}{X} 以撒為他們擺設宴席，他們就一起吃喝。 \end{tabularx} \\ \\ \relax
26:31 & \begin{tabularx}{0.7\textwidth}{X} 他們清早起來，彼此起誓。以撒送他們走，他們就平平安安地離開他去了。 \end{tabularx} \\ \\ \relax
26:32 & \begin{tabularx}{0.7\textwidth}{X} 那一天，以撒的僕人來，把挖井的消息告訴他，說：「我們得到水了。」 \end{tabularx} \\ \\ \relax
26:33 & \begin{tabularx}{0.7\textwidth}{X} 他就給那井起名叫示巴，因此那城名叫別是巴，直到今日。 \end{tabularx} \\ \\
[1ex]
\hline
\hline
\end{longtable}
$^{1}$各位靈智妹 大家好.
多謝 Pam 幫我們讀了一段這麼長的經文.
我們來到創世記第41講.
上一次《說分傳道》講了些什麼 大家記不記得.
因為講到以撒逃避饑荒.
逃難到非利士的基拉爾.
然後聲稱太太是他的妹妹.
這個情節在亞伯拉罕的事蹟中也出現過.
如果我們仔細看看聖經記載以撒的歷史.
其實你會發覺很多地方和亞伯拉罕很相似.
他們都是因為饑荒而逃難.
亦都去過基拉爾定居.
亦都聲稱自己太太是妹妹撒謊.
他們太太長期不育.
他們挖的井都是被人搶的.
他們都是和阿比米勒在別士巴訂立文約.
很多很多很相似的地方.
所以有些學者就以為摩西會不會搞錯.
是否他們重複記載了 擾亂了歷史呢?.
其實聖經這樣記載是非常有意思的.
是想告訴我們.
以撒和亞伯拉罕的人生軌跡是重疊的.
其實凸顯了一件事.
以撒延續著上帝對亞伯拉罕的約的線索.
亦都被人看到就算人多麼的軟弱.
重複的犯錯.
但是神對亞伯拉罕的應許.
仍然是跨世代堅定不移的實現中.
在這段經文我們看到.
神對亞伯拉罕的應許.
如何在很多不可能的當中仍然能夠逐步實現.
在看完經文之前我們一起禱告.
天父上帝 多謝你.
你似眼你的說話.
求你開我們的眼睛.
用你的角度去看你的說話.
求你讓我們看到你要我們看的東西.
讓我們明白神留這段聖經在今日.
對我們來說是甚麼心意.
多謝你 祈禱奉耶穌的聖名.

$^{41}$阿們.
創世記第26章開頭時.
以撒在那個地方耕種.
那個地方是基拉爾.
基拉爾大家看到那個位置是甚麼地方.
不用理會箭咀.
理會那個位置便可以了.
他說以撒在基拉爾這個地方耕種.
那一年有百倍的收成.
然後說切腹給他.
他有牛群羊群.
有很多僕人.
菲律賓人就襲到他.
成了大富翁 槍勝.
看到這裡的時候.
這些是很成功神學的感覺.
看到這麼好的.
跟耶和華說是無事抵.
我們試試看.
第一個不可能的事發生是甚麼呢?.
我們知道亞伯拉罕以撒.
他們不是一個農耕的民族.
是遊牧民族.
耕種對他們來說不在行.
但他們說甚麼呢?.
聖經說以撒在那地耕種.
一年竟然有百倍的收成.
是甚麼事發生?.
就算是農業專家.
你有豐收也好.
不會有百倍的收成吧?.
究竟聖經說百倍收成是否一個文學手法.
誇大呢?.
又不像.
因為羊群牛群有很多的僕人.
發達昌盛成為大富翁.
一年之間發生這麼多事.
以撒是甚麼人呢?.
他是一個難民.
逃避饑荒到基拉爾.

$^{81}$難民會發達在一個寄居的地方.
是否不可能呢?.
還有一件事是.
他們這個發達.
是令到菲律賓人憎恨他.
為甚麼他會這麼發達呢?.
是否以撒做了甚麼好事呢?.
又不是.
上文記載以撒做了甚麼.
好像一向都是這樣.
說謊.
他用自己的方法去保護自己.
從屬靈角度來說.
他是失敗的.
但是關鍵是什麼呢?.
耶和華賜福給他.
為甚麼耶和華賜福給他呢?.
他是失敗者.
因為甚麼呢?.
這個很重要.
有時候我看聖經的時候.
看細節的位置很有意思.
亞伯拉罕在世的時候.
他父親的僕人 華夷正.
被菲律賓人責罵得全都塞住了.
填滿了肚子.
聖經不厭其煩地記載.
他父親.
有意思吧?.
其實他是在延續.
他父親亞伯拉罕.
上帝對亞伯拉罕的約.
是因為上帝對亞伯拉罕的約.
上帝手約.
令他的後裔成為大禍.
所以上帝令他昌盛.
你說.
他那麼發達.
菲律賓人塞住了他的井.
這是不是惡作劇呢?.

$^{121}$這絕對不是幼稚的發洩.
不要以為他們真的看不過眼.
塞了井來作弄你.
不是的.
井是用來做什麼的?.
儲水.
打水.
在連降雨量只有三百毫米的中東地區.
井是重要的水源.
沒有了這些井和水源.
那是什麼呢?.
叫你離開.
不說發達.
你連生存都成問題.
所以你只能夠離開這裡.
不單止菲律賓人要離開這裡.
菲律賓王都出了聲.
菲律賓是什麼地方呢?.
可能我們很奇怪.
如果大家熟悉歷史的時候.
有些學者就會反駁.
《實經記載》這件事其實是不準確的.
為什麼呢?.
亞伯拉罕年代沒有菲律賓人.
菲律賓人是海上的民族.
到了十二世紀的時候.
他們才在地中海沿岸的地方定居.
以撒戈年代早了幾百年.
現在記載有菲律賓人.
《聖經》是假的.
不可信的.
是不是呢?.
其實不是的.
譬如我說.
宋朝的時候.
宋朝的皇帝逃難來到香港.
去到九龍城這個地方.
大家聽過這些歷史吧.
我這樣說對不對呢?.
對的.

$^{161}$但我可以抽清你的.
宋朝的時候都沒有香港這個稱呼.
更加沒有九龍城這個稱呼.
所以我正確來說.
宋朝皇帝逃難到官府牆那裡避難.
你不知道我在說什麼.
我說官府牆你不明白.
我說九龍城你就知道了.
同樣道理.
在摩西的年代.
大家都知道那個地方叫做非利士地.
你明白我的意思吧.
所以《聖經》說.
非利士人是沒有錯的.
用他們明白的語言.
去表達他們不明白的地方.
他們才能和這個歷史連繫得到.
所以非利士地就是那個區.
所以基拉爾就是當日的非利士地的範圍.
不單止非利士人和他們作對.
非利士王亞比米勒也有發聲.
因為亞比米勒已經出現過.
什麼時候出現過呢?.
他的父親亞伯拉罕也和亞比米勒立約.
亞比米勒怎麼這麼長壽呢?.
亞比米勒在閃族語言中.
是說「我父是王」.
其實這個就像埃及的法老.
羅馬的凱撒一樣.
是王的稱號.
所以不一定是一個名字.
亞比米勒就是王的稱號.
就像埃及的法老一樣.
所以當時的亞比米勒也有發聲.
他說「你離開我們吧.
因為你比我們強盛得多」.
這個譯法我覺得非常非常客氣.
看起來像一個勸告.
既然我們非利士人不歡迎你.
不如你們走吧.

$^{201}$安全起見.
不是的.
如果我們看回新譯本譯得比較貼切一點.
他說「你們走吧 離開我們吧.
因為你們比我們強大得多」.
如果用今天的語言來說.
其實亞比米勒是用行政命令驅逐以撒出境.
下令他走.
那走不走好呢?.
這是一個弔詭的位置.
因為你比我們強大得多.
那現在誰是強者呢?.
以撒是強者.
那意味著什麼呢?.
以撒的勢力比亞比米勒還要大.
如果以撒對著一個弱勢亞比米勒說.
「你老幾呀 應該誰走呀」.
他可以不走的.
我比你強.
反而我可以用我的強勢佔領你的地方.
成為我的領地.
歷史上有很多例子.
但是以撒離開那裡.
去了基拉爾谷.
以撒是一個強勢的以撒.
他放下他發跡的地方.
離開那裡.
其實離開是表示什麼呢?.
你在那裡建立了一切的事業.
一切的基業.
你離開意味著要放棄這些東西.
還要另外找地方落腳.
然後去到一個新的地方.
你要由頭來過.
之前我要挖井.
不是那麼容易的.
這個決定其實絕對絕對不容易.
比你不容易.
我想很少人會這樣做.
因為是一個強勢.

$^{241}$去到基拉爾谷.
基拉爾谷是什麼地方呢?.
大家看看地圖.
在基拉爾谷下面沒有戰墜.
我找了一個地方給你看.
是一條河流經.
其實這條是一條旱溪.
什麼意思呢?.
就是說到雨季才會下雨.
才會有水雨.
基拉爾谷有一個特質的地方.
就是雖然是一個乾旱的河谷.
但仍然可以有地下水源.
不過要挖得很深.
要懂得找才能找到.
不是那麼容易的一件事.
所以他不是不可以落腳.
但要很艱辛.
那一區也有牧人.
也是兵家必爭之地.
根據歷史記載.
這裡是軍事重要的地方.
去到基拉爾谷.
發生什麼事呢?.
為什麼以撒要離開這個地方呢?.
很多人都在想.
是不是以撒太軟弱呢?.
是不是以撒怕了比米勒呢?.
又不是啊 比米勒比他強.
其實如果大家記得.
神跟亞伯蘭說了什麼話.
創世記第十二章.
不如我們嘗試讀一讀.
我們一起嘗試大聲讀一讀.
神跟亞伯蘭說.
你要離開本地.
本族 本家.
我已嘗試你們的地位.
我已嘗試你們的貴大貨.
我已嘗試告發你們.

$^{281}$你們的命運貴大.
你要試人而得報.
我已到達目的.
我已嘗試回答.
我已經有多個人.
你再不回答我.
我便要人你而得報.
我已嘗試回答.
亞伯蘭信不信耶和華說.
從何而知?.
照做.
到了第二代.
你覺得上帝的應許還有沒有效?.
有沒有效? 有.
以撒信不信?.
怎知道以撒信?.
神說為你祝福的.
我必定賜福給他.
助紂你的.
我必助紂他.
上帝應許成為大國.
既然如此.
非利士人.
不滿意.
以撒趕走以撒.
以撒要不要自己出手?.
不用.
以撒的退讓表示什麼?.
是表示他相信上帝的應許.
是不是?.
他相信上帝出手.
既然上帝出手.
我需要出手嗎?.
我需要用我的強勢.
去對付這些非利士人嗎?.
不用.
我只要走在上帝的約裡便足夠.
所以.
他離開.
接著.

$^{321}$18-23節.
他父親亞伯拉罕在世時所畫的水井.
他將它重新畫出.
仍然叫他父親所取的名字.
為他們命名.
其實看聖經.
有時留意細節挺有意思.
他提到以撒的時候.
提到亞伯拉罕.
都滿意他父親的說話.
是不是? 剛才也是一樣.
表示什麼?.
表示這種約的承傳.
這個關係的承傳.
還有一個細節.
以前的井塞了.
現在再畫出來.
不僅再畫出來.
還做了一個很重要的動作.
照著他父親所取的名字命名.
用他父親命名這個井的叫法.
有什麼意義?.
這也是在延續亞伯拉罕所做的事.
雖然被人掩蓋了.
但仍然再一次復活.
這個井仍然存在.
仍然活著.
上帝的恩典.
你是殺不了的.
是不是?.
他是相信神的應許.
這也是繼承亞伯拉罕的一個線索.
亞伯拉罕畫了很多井.
為什麼要畫很多井?.
因為要很多井儲水.
要有足夠的水源.
但是十九節.
以撒吉牧人在谷中.
在基拉爾谷裡畫井.
就得到一口什麼井?.

$^{361}$活水井.
什麼是活水井?.
活水井就是畫得很深.
而你找到水墨.
什麼叫水墨?.
墨就是墨包的墨.
畫得適合的水墨.
然後畫得夠深.
才能拿到活水的井.
活水的井和其他井有什麼不同?.
其他井是用來儲水的.
下雨的時候你裝水.
給你用的.
但活水的井是自己湧出來的.
源源不絕地下水供應.
所以對於中東這麼乾旱的地方來說.
活水的井是最珍貴的資源.
是很多人夢寐以求要爭奪的東西.
所以就明白為什麼第二十節.
基拉爾人不會用土封這個井.
反而是搶奪這個井.
是不是很不一樣?.
這是很寶貴的.
原來上帝在以撒身上所做的.
是比亞巴拉克年代更進一步的賜福.
他們看到這個井有水的時候.
就說這個水是我們的.
他們可能是原住民.
覺得這個資源是我的.
為什麼要拿我的水?.
那麼以撒有什麼反應?.
有沒有跟他們揮霍過?.
他只是將這個名字起名叫埃塞.
埃塞是什麼意思呢?.
是象徵的意思.
其實以撒有沒有情緒呢?.
以撒不是盲目沒有情緒.
你搶就搶吧.
他有情緒.
是象徵 被人爭奪的東西.

$^{401}$但是他沒有揮霍過.
這是剛才的原因.
以撒的僕人又說這個井很棒.
他們又說這個井是象徵.
那表示這個是什麼井呢?.
又是一座水井.
如果不是的話就不用爭了.
這次以撒說的是西提拿.
西提拿這個字是有敵視有憎恨的意思.
這個字的字根和撒旦的字根一樣.
是邪惡的.
表明一件事就是對方的行為是不公義.
是邪惡的 為什麼呢?.
滑井的是我們.
辛苦工作的是我們.
但是我們辛苦工作的成果是怎樣的呢?.
他們略奪.
你想想.
以撒有沒有力量和他們爭奪呢?.
有沒有?.
為什麼你們額頭有呢?.
這群基拉爾的僕人.
連基拉爾王阿比米勒.
都說以撒比他們強盛得多.
何況這些平民百姓.
如果以撒要和他們爭奪的話.
一定贏我們.
但是以撒沒有這樣做.
為什麼沒有這樣做呢?.
大家要知道 因為上面有一個神.
以撒離開那裡.
真慘 又離開了.
有條件的 卻被人趕走.
慘不慘呢?.
趕走 流浪不是一件浪漫的事.
我記得我讀大學的年代.
幻想自己背著背囊.
走絲綢之路.
去中亞地方流浪.
那時二十多歲.

$^{441}$有一種很美麗的憧憬.
其實不切實際.
你以為很爽 很流浪.
如果今天三十多度.
我在火車站走回來的時候.
就知道為什麼.
其實流浪是一個什麼樣的生活呢?.
這個世界還有沒有流浪的民族呢?.
有啊.
有沒有聽過羅姆人呢?.
在歐洲叫吉卜賽人.
其實吉卜賽人是一個很負面的稱呼.
他們是羅姆人.
因為歐洲人都不知道他們的來歷.
以為他們是埃及來的.
所以叫Gypsies 吉卜賽.
給你什麼印象呢?.
去歐洲旅行的時候.
吉卜賽人的印象就是偷東西.
他們要你保護自己的財物.
其實羅姆人在歐洲這麼多世紀以來.
其實都是一個流浪的民族.
他們沒有自己的國家.
沒有自己的土地.
以一些篷車到處流浪去做寄居者.
寄居的生活是怎樣的呢?.
飽受歧視.
他們很慘.
因為他們到處流浪.
當年歐洲黑死病 鼠疫.
依照這麼多次的疫症.
歐洲人都當他們是過街老鼠.
因為他們到處流浪.
覺得他們是傳播病毒的.
污糟邋遢的.
他們靠做小偷.
行乞 街頭賣藝.
去過生活.
流浪者.
還有 薩瑪人有沒有聽過?.

$^{481}$就是巴瑤人.
他們在菲律賓 印尼 馬來西亞的海域.
巴里曼丹 索拉威西 綿藍老島那一帶.
去生活.
他們其實更慘.
稱為海上吉卜才人.
他們被迫離開陸地.
在海上住棚屋.
他們是靠岸上的接濟.
給錢 剋蝕.
在海上打魚.
他們沒有國籍.
因為他們沒有國籍.
所以沒有醫療 沒有教育.
沒有就業的機會.
原來 作為一個流浪的民族.
到處流浪 是一點都不爽.
流浪民族的下場就是這樣.
但是我們看看以撒.
以撒離開那裡 又滑又靖.
他們是否走這條路線呢?.
第二個不可能.
就是甚麼呢?.
他們離開那裡.
人們已經不再跟他們爭.
因為他們離開了.
向南走都很遠.
去到一個地方 滑靖.
他叫利河北 甚麼意思呢?.
這個地方就是耶和華給我們的寬闊之地.
是寬闊的意思.
我們能夠在這裡興旺.
一個流浪的民族.
不斷被人趕走.
他們竟然看到前景是興旺.
是否神的作為呢?.
他們接著就去別士巴.
別士巴是很有意思的地方.
最後我才說別士巴.
別士巴是甚麼地方呢?.

$^{521}$大家可以看到.
由基拉爾谷去別士巴.
其實利河北就是在南邊的地方.
好像是四和三那邊的地方.
想回別士巴.
他們回到別士巴的當晚.
耶和華親自向以撒說話.
又是說「你父親」.
來到這個特別的地方.
在這些細節的位置.
你會看到上帝顧念他的約.
「我是你父親亞伯拉雄的神.
你不用怕我與你同在.
要賜福給你.
因為我僕人亞伯拉雄的緣故.
使你的後裔增多」.
上帝一再向他保證.
很有趣的.
上面就是父親.
接著就是後裔.
而以撒的父親.
下面就是後裔的中間.
以撒就是傳遞上帝恩約的重要過程.
所以這一段聖經讓我們看到.
上帝的約是如何由亞伯拉雄.
一直延續到雅各.
一直延續到以色列的民族.
以撒就在那裡建了一座壇.
其實如果我們知道.
神與以撒和亞伯拉雄的約.
不只是十二章.
在十七章,十八章.
都有再次重申.
記不記得?.
大家應該讀過.
這次是直接向以撒說話.
傳媒都在說甚麼呢?.
「你的後裔增多」.
「我要成為大國」.
要留意「大國」這件事.

$^{561}$「大國」應該是一個國家的國民.
今天他們是流浪者.
你說跟「大國」這個應許是否很遠呢?.
上帝所說的話.
似乎好像沒有落實到.
是不是呢?.
我們一會兒看看是不是.
以撒在那裡築了一座壇.
求告耶和華的名字.
支搭帳篷.
嘩!很正!.
在說甚麼呢?.
神和人之間的關係.
上帝向他說話.
他有沒有回應呢?.
有!.
他的回應就是.
求告耶和華的名字.
這樣才能溝通.
不是上帝單方面跟他說話.
而是他也求告上帝的名字.
他跟上帝是有連繫的.
以致他能夠承受.
神給亞伯拉罕的福.
他是在恩約裡面.
他沒有離開上帝.
下一次.
你再看下去以後.
二十七章以後.
雖然是以撒的後裔.
以素他是離開了這條線.
走在神的恩約裡面.
你才跟上帝的約有份.
你不要離開上帝.
好了.
有一件事發生.
這件事是很受厚面的.
甚麼事呢?.
就是阿彼米勒跟他立約.
以前發生過阿彼米勒跟亞伯拉罕立約.

$^{601}$這次看看這件事是甚麼事.
阿彼米勒跟他的顧問.
阿胡撒.
和軍隊的元帥飛過.
從基拉爾來到以撒那裡.
做甚麼呢?.
看下去的時候知道.
要跟他們立約.
這件事是甚麼事呢?.
我發誓.
聖經有時記載一些仔細的位置.
很有意思.
它專門說顧問.
元帥.
這是甚麼陣容呢?.
用今天現代的說法來說.
這是一個國事訪問.
顧問的原文是朋友.
其實可以說是親信.
他帶領他國家的顧問親信.
阿胡撒.
和軍隊的元帥.
即是有政治上的高官.
還有軍方的最高領袖.
這是代表整個國家的實力.
飛過出現過.
在上一次和阿巴拉罕立約的時候.
飛過這麼長的時間.
所以聖經學者認為.
飛過不是一個人名.
是一個家族名.
去到立約.
所以其實他記載這個排場很有意思.
我們看到了甚麼呢?.
以撒是一個甚麼國家呢?.
不是的 它還不是一個國家.
不過上帝已經允許他們成為大國.
但今天.
竟然.
一個阿比米勒.

$^{641}$一個主權國家的領袖.
帶領他政府的內閣.
用國事訪問外交的排場.
去見一個以撒的家族.
似乎以撒已經有大國的影子了.
你有沒有留意到呢?.
好了 他說了甚麼呢?.
以撒就說.
你為甚麼來找我呢?.
你趕我走的.
他就說了甚麼呢?.
他們說.
他們不是阿比米勒.
是他們.
不是自己代表.
是代表他個人.
是代表他的國家.
他說我們明明見到.
耶和華和你同在.
因此就是說.
讓我們雙方彼此起誓.
我們跟你立約.
使你不加害我們.
正如我們未曾侵犯你.
來善待你.
並且送你平安離開.
你是甚麼呢?.
蒙耶和華賜福的.
你有沒有留意阿比米勒的說話?.
他開頭和拉尾.
都有一個共通的位置.
是在說甚麼呢?.
看到耶和華賜福給誰呢?.
以撒.
他看到這件事.
有甚麼意義呢?.
我們問.
以撒究竟有甚麼強勢.
使阿比米勒提出立約呢?.
我們要知道歷史上甚麼是立約.

$^{681}$立約通常是一個戰敗.
處於弱勢的一方.
要求強勢的一方.
保障和平.
才是立約.
現在是一個主權國家的領袖.
帶著軍方的元帥.
他們有軍隊 以撒沒有軍隊.
他們處於一個主動提出立約的一方.
即是弱勢的一方.
然後他們說甚麼呢?.
我明明看到耶和華給你痛哉.
這是很重要的.
我們說.
以撒的強勢強在哪裡呢?.
他有沒有做過甚麼風功偉績.
那種烈烈的事呢?.
沒有.
他所做的就是不斷的退讓.
為甚麼不斷的退讓的時候.
他還那麼強呢?.
正是因為這樣.
他不用自己的實力和作為.
他退出了這個爭執的可能性.
然後讓上帝的作為在人身上彰顯.
人就看到上帝的作為非同小可.
如果以撒用他的強勢為自己戰鬥.
人家看到甚麼呢?.
以撒你真是厲害.
你真是厲害.
我服了你.
不是的.
他一頭一尾的說話是說甚麼呢?.
是說耶和華厲害.
所以以撒最強的地方是甚麼呢?.
以撒最強的地方.
就是放下自己.
讓神工作.
所以我發覺.
詩篇有句說話是很對的.

$^{721}$「你們來看耶和華的作為.
耶和華是紫色戰爭.
你要休息,要知道我是神」.
要休息.
休息就是你不要那麼不人為的作為.
你靜止下來.
看上帝的作為.
要知道我是神.
以撒就是知道耶和華是神.
所以他不需要自己的實力.
正是因為他要紫色戰爭.
他要上帝的心而紫色戰爭.
他不是用武力.
去征服.
然後掠奪地方.
攻城掠地.
然後說這個地方是上帝賜給我的.
不是的.
他唯有這樣做的時候.
神說甚麼呢?.
「我的名字是愛邦被尊崇」.
你想想.
這句說話正正就是以撒的註腳.
阿彼米勒就是尊崇上帝的名字.
就是因為以撒.
他放下他自己的實力.
願意在神面前不作為.
讓上帝成就他的美意.
人就看到.
這是上帝的作為.
上帝的名字就能夠被愛邦人尊崇.
全地都被尊崇.
所以我覺得.
以撒的生命很精彩.
精彩的就在於.
他有實力.
但他不用的實力.
讓神彰顯上帝的實力.
這是謙卑的生命.
所以.

$^{761}$再看下去的時候.
還有一個妙的地方.
是什麼呢?.
以前的阿彼米勒和阿伯拉立約.
如果大家記得的時候.
立約的內容是說什麼呢?.
阿彼米勒要求.
他父親.
以撒的父親.
效忠於他.
對我忠誠.
善待我的後裔.
對我的國家忠誠.
忠誠是什麼意思呢?.
是一個強者要弱者.
臣服於他的意思.
但今天.
他的第二代.
阿彼米勒和立約有什麼要求呢?.
你不加害於我們.
你看到這個轉變嗎?.
現在.
他是要求.
得到.
以撒的保證.
保證我的安全.
現在是誰強呢?.
是以撒強.
不是再要求你效忠於我.
而是要求.
你不要加害於我.
在這些細節的位置.
我們看到上帝的恩若.
是否逐步地實現呢?.
大過影子.
越來越清楚.
所以.
阿彼米以撒擺設筵席.
和這些人一起吃喝.
擺設筵席一起吃喝.

$^{801}$真是他客氣了.
這是一個禮貌式.
擺設筵席是一個.
很有外交含義的行為.
是表示什麼呢?.
表示大家和平相處.
最近特集會.
北京都設宴.
彼此的氣勢.
是否一個.
和平.
平等的.
外交.
接著看到什麼呢?.
以撒送我們平平安安的時候.
你試試看.
擺設筵席是什麼呢?.
是東道主.
設宴款待你.
以撒再不是一個流浪飄蕩的.
處境.
以前你趕.
以撒走.
現在以撒在地上送你走.
是不是很妙呢?.
還有.
平平安安回去是什麼意思呢?.
就是因為.
他得到以撒的安全保證.
以撒這個強權.
願意.
答應不加害於你.
他就安心.
平平安安的走.
他不是怕.
以撒.
他是怕賜福給.
以撒的上帝.
那一天.
以撒的僕人又來懷井了.

$^{841}$得到水.
他們就叫別士巴.
什麼意思呢?.
他們還可以定居.
別士巴是一個很有意思的地方.
有什麼意思呢?.
他是亞伯拉罕.
以撒.
兩代好妙的連結.
別士巴是紀念了.
神和人的關係.
如果大家記得.
亞伯拉罕當年在別士巴.
種柳樹求告耶華的名字.
是不是?.
今天.
亦是什麼呢?.
他在別士巴那裡.
求告上帝.
然後.
紀念了他和外族的關係.
亞伯拉罕和.
菲利斯人在那裡用.
七隻羔羊納藥得到和平.
這裡.
亦都說了.
他在這個地方.
和菲利斯建立和平.
所以.
這個地方.
是以色列和外族.
恩怨情仇得到和平的見證.
所以日後.
以色列的國境.
由丹到別士巴.
這個.
說法非常有意思.
來到最後.
究竟這段歷史.
對我們有什麼意義呢?.

$^{881}$是否我們知道一個故事就算了?.
我們很多時候都聽到.
人的盡頭是神的開始.
有沒有聽過這句說話?.
其實為何如此?.
為何人.
到盡頭才肯被上帝.
動手?.
我們看到以撒比.
今天看到一件事.
在他強盛的時候.
他不需要走到盡頭的時候.
他甘願.
自己退下來.
收斂他自己的.
強勢.
在神面前謙卑.
讓上帝.
在他身上的作為.
讓他的軟弱.
成為榮耀.
成為上帝的工具.
這就是.
以撒比給我們看到精彩的地方.
還有.
以撒.
不是很魯皮.
因為他很魯皮.
所以要退讓.
他說這是撒旦的作為.
這是相爭.
他起的名字.
這是罪惡.
我們今天看到很多.
罪惡的東西的時候.
我們要立高拳頭.
我們要替天行道.
我們彷彿這個判官.
就是.
羅賓漢.

$^{921}$就是正義的執行者.
但是會不會這樣.
造就到.
你那種罪惡呢?.
我們看到不義的行為.
我們就怒火攻心.
因為我們是正義的.
但你試想想.
如果我們是正義的時候.
一個絕對公義的上帝.
會不會比我們更加公義呢?.
如果我們這麼憤怒的時候.
上帝會不會更加憤怒呢?.
上帝就出手了.
所以.
今天.
我們回望.
詩篇的說話.
我覺得是很對的.
我們不需要.
那麼憤怒.
止住怒氣.
不需要心懷不平.
以撒就是什麼呢?.
他止住怒氣.
等候耶和華必.
承受土地.
結果.
上帝最終.
就賜他們土地.
去到今天新約的時候.
保羅提醒我們.
我們信不信.
上帝.
是公義的呢?.
我們信不信.
上帝.
是會.
施行他的審判呢?.
我們信不信.

$^{961}$我們只不過是一個.
軟弱的人,我們不是判官.
我們盼望今天.
以撒的歷史.
都能夠給我們一些提醒.
好嗎?.
我們一個禱告.
求上帝我們多謝你.
我們人生.
會面對很多衝擊.
也會見到.
很多很多不公義的事情.
我們會.
為這些事情而憤怒.
但我們知道.
上帝是公義的神.
你會比我們更加憤怒.
求上帝教我們.
怎樣去面對.
這樣的情景.
尤其是.
上帝教導我們.
即使我們.
是有權有勢的時候.
我們所依靠的.
是神.
而不是我們自己的能力.
願我們願意學習謙卑.
讓上帝的.
名在我們的軟弱上.
顯為大.
多謝你,祈禱奉耶穌的聖名.
\newpage



\section{使徒行傳 2:14-36}
\label{sec:032244}
\textbf{仰賴基督與聖靈作見證}~~~~黃浩君傳道
\newline
\newline
連結: \href{http://www.church.com.hk/acms/content.asp?site=cdc\&op=show\&type=product\&code=032244\&layout=sermon}{\texttt{032244}} ~~~~ 語音日期: 2026-05-31
\newline
\newline
\hyperref[sec:032243]{\small{< < < PREV SERMON < < <}}
~
\hyperref[sec:index_chronic]{\small{[返順時目]}}
~
\hyperref[sec:032245]{\small{> > > NEXT SERMON > > >}}
\newline
\newline
使徒行傳 2:14-36
\newline
\begin{longtable}{cl}
\hline
\hline
章節 & 經文 (和合本修訂版)\\
\hline
2:14 & \begin{tabularx}{0.7\textwidth}{X} 彼得和十一個使徒站起來，他就高聲向眾人說：「猶太人和所有住在耶路撒冷的人哪，這件事你們要知道，要側耳聽我的話。 \end{tabularx} \\ \\ \relax
2:15 & \begin{tabularx}{0.7\textwidth}{X} 這些人並不像你們所想的喝醉了，因為現在才早晨九點鐘。 \end{tabularx} \\ \\ \relax
2:16 & \begin{tabularx}{0.7\textwidth}{X} 這正是藉著先知約珥所說的： \end{tabularx} \\ \\ \relax
2:17 & \begin{tabularx}{0.7\textwidth}{X} 『神說：在末後的日子，我要將我的靈澆灌凡血肉之軀的。你們的兒女要說預言；你們的少年要見異象；你們的老人要做異夢。 \end{tabularx} \\ \\ \relax
2:18 & \begin{tabularx}{0.7\textwidth}{X} 在那些日子，我要把我的靈澆灌，甚至給我的僕人和婢女，他們要說預言。 \end{tabularx} \\ \\ \relax
2:19 & \begin{tabularx}{0.7\textwidth}{X} 在天上，我要顯出奇事，在地下，我要顯出神蹟，有血，有火，有煙霧。 \end{tabularx} \\ \\ \relax
2:20 & \begin{tabularx}{0.7\textwidth}{X} 太陽要變為黑暗，月亮要變為血，這都在主大而光榮的日子未到以前。 \end{tabularx} \\ \\ \relax
2:21 & \begin{tabularx}{0.7\textwidth}{X} 那時，凡求告主名的都必得救。』 \end{tabularx} \\ \\ \relax
2:22 & \begin{tabularx}{0.7\textwidth}{X} 「以色列人哪，你們要聽我這些話：拿撒勒人耶穌就是神以異能、奇事、神蹟向你們證明出來的人，這些事是神藉著他在你們中間施行，正如你們自己知道的。 \end{tabularx} \\ \\ \relax
2:23 & \begin{tabularx}{0.7\textwidth}{X} 他既按著神確定的旨意和預知被交與人，你們就藉著不法之人的手把他釘在十字架上，殺了。 \end{tabularx} \\ \\ \relax
2:24 & \begin{tabularx}{0.7\textwidth}{X} 神卻將死的痛苦解除，使他復活了，因為他原不能被死拘禁。 \end{tabularx} \\ \\ \relax
2:25 & \begin{tabularx}{0.7\textwidth}{X} 大衛指著他說：『我看見主常在我眼前，他在我右邊，使我不至於動搖。 \end{tabularx} \\ \\ \relax
2:26 & \begin{tabularx}{0.7\textwidth}{X} 所以我心裡歡喜，我的舌頭快樂，而且我的肉身要安居在指望中。 \end{tabularx} \\ \\ \relax
2:27 & \begin{tabularx}{0.7\textwidth}{X} 因你必不將我的靈魂撇在陰間，也不讓你的聖者見朽壞。 \end{tabularx} \\ \\ \relax
2:28 & \begin{tabularx}{0.7\textwidth}{X} 你已將生命的道路指示我，必使我在你面前充滿快樂。』 \end{tabularx} \\ \\ \relax
2:29 & \begin{tabularx}{0.7\textwidth}{X} 「諸位弟兄，先祖大衛的事，我可以坦然地對你們說：他死了，也埋葬了，而且他的墳墓直到今日還在我們這裡。 \end{tabularx} \\ \\ \relax
2:30 & \begin{tabularx}{0.7\textwidth}{X} 既然大衛是先知，他知道神曾向他起誓，要從他的後裔中立一位坐在他的寶座上。 \end{tabularx} \\ \\ \relax
2:31 & \begin{tabularx}{0.7\textwidth}{X} 他預先看見了，就講論基督的復活，說：『他不被撇在陰間；他的肉身也不見朽壞。』 \end{tabularx} \\ \\ \relax
2:32 & \begin{tabularx}{0.7\textwidth}{X} 這耶穌，神已經使他復活了，我們都是這事的見證人。 \end{tabularx} \\ \\ \relax
2:33 & \begin{tabularx}{0.7\textwidth}{X} 他既被高舉在神的右邊，又從父受了所應許的聖靈，就把你們所看見所聽見的，澆灌下來。 \end{tabularx} \\ \\ \relax
2:34 & \begin{tabularx}{0.7\textwidth}{X} 大衛並沒有升到天上，但他自己說：『主對我主說：你坐在我的右邊， \end{tabularx} \\ \\ \relax
2:35 & \begin{tabularx}{0.7\textwidth}{X} 等我使你的仇敵作你的腳凳。』 \end{tabularx} \\ \\ \relax
2:36 & \begin{tabularx}{0.7\textwidth}{X} 故此，以色列全家當確實知道，你們釘在十字架上的這位耶穌，神已經立他為主，為基督了。」 \end{tabularx} \\ \\
[1ex]
\hline
\hline
\end{longtable}
$^{1}$弟兄姊妹平安.
先開紙紙 我們一起禱告.
自教交托 奉主明求 阿們.
話說我月頭去了兩個星期歐洲.
我還記得那天是五月十四日 星期四.
我和太太就身處在維也納的第三天.
當時很奇怪 明明是星期四平日.
不知為何街上很多人.
白電車的班次疏落了 五分鐘變成十分鐘一班.
到傍晚的時候.
我們去到超級市場想買支汽水飲.
誰知超級市場沒有營業.
我太太就問 其實今天是不是公眾假期.
我上網一找 原來那天是耶穌升天日.
當一刻 我第一個念頭是.
其實我做錯了甚麼 耶穌升天要懲罰我沒有汽水飲.
當然我回到酒店的時候就開始想.
我們香港的教會通常會慶祝聖誕節和復活節.
但可能因為公眾假期的關係.
耶穌升天節和上星期的聖靈降臨節.
好像往往都會被忽略.
彷彿耶穌降生 受死 復活 完成了救恩.
整件事就完結了.
無關乎 有時我們的信仰 屬靈 生命.
都只是停留在耶穌很愛我 耶穌救了我這個階段.
星期日 崇拜完 聽完道 在教會服事完.
就覺得耶穌很愛我 感覺很良好.
然後回到職場 回到我們生活的處境裡.
要為神作見證.
就處處感到無力和乏力 甚至是恐懼.
可能在我們意識裡.
我們都會覺得耶穌升了天 離開了.
丟下了我們在地上捱累.
面對現實的殘酷.
然後我們可能就要問自己一條問題.
耶穌升天到底和我這一刻在活在地上有什麼關係呢?.
我們是不是錯過了一些東西.
讓我們在生活上為神作見證 覺得很乏力呢?.
而今天的經文是《仕途行傳》二章十四至三十六節.
說到二千年前 耶穌升天之前.

$^{41}$就應許過門徒說.
會受聖靈的洗 聖靈降臨的時候.
他們就會得著能力成為祂的見證人.
但當在五場節聖靈降臨的時候.
門徒就用謊言宣講神大作為的時候.
他們面對的是外在環境充滿著對他們的嘲笑聲.
說他們灌滿了新酒.
同時其實在旁人的眼裡.
他們都是一般的平民百姓 講白一點就是一個小薯仔.
甚至他們是離開過 離棄過他們的老師耶穌.
但面對著外在環境的不友善.
同時自己只是一個小薯仔的情況下.
彼得竟然站起來挺身而出.
向群眾講了一篇道.
指向那一位復活升天的耶穌.
讓當時的人知道究竟在發生什麼事.
今天我們同樣可以透過彼得這篇道.
明白三個關於為神作見證的事實.
第一 我們可以透過聖靈的囂觀.
打破可以知道神的心意.
打破我們對自身的限制.
第二 靠著基督復活的真實.
幫助我們面對不同的客觀惡劣環境的恐懼.
第三 我們既然稱聖天作王的耶穌為主的時候.
我們可以得著生命的改變和屬靈的恩賜.
讓我們一起學習怎樣靠著聖靈.
怎樣靠著耶穌成為神在地上的見證人.
來到今天的經文 在十四十五節裡面.
面對議論紛紛抱有譏笑態度的群眾.
彼得和十一個門徒站了起來.
彼得大聲說.
猶太人和所有住在耶路撒冷的人.
這件事你們要知道 你要質疑聽我的話.
這些人不是你們所想像中的喝醉酒.
如果是喝醉酒的話.
現在才九點鐘 是早上九點鐘.
就太早了.
原來在猶太人的規矩裡.
朝聖者在早上祈禱和獻祭完之前.
是要禁食和禁酒的.

$^{81}$所以九點鐘喝醉真是太早.
彼得就用當時的生活常識打破了群眾的誤解.
否定了試圖喝醉酒的可能性之後.
彼得就立即將群眾的目光.
放在他們耳熟能詳的舊約裡.
解釋他們見到的一切事.
在十六節彼得說道.
這正是先知約爾所說的.
是要表達到群眾所見到的一切現象.
是應驗著約珥書二章二十八至三十二節的寓言.
然後十七至十八節.
在幕後的日子我要將我的靈囂冠繁血欲之軀的.
幕後的日子就是耶穌受死復活升天之後.
聖靈降臨工作那一刻開始.
同時對當時的猶太人來說.
更加意味著是彌賽亞時代的開始.
而在這個時代神的靈即是聖靈會囂冠繁有血氣的.
囂冠簡單來說就是一個農耕的比喻.
目的是要帶出神次下聖靈.
完全是神從上而下的恩典.
同時讓人得著新的生命.
那哪些人可以得著聖靈的囂冠.
在經文裡我們看到四組人.
你們的兒女.
你們的少年.
你們的老人.
我的僕人和婢女.
他們可以說喻言.
見異象.
做異夢.
意思就是不分男女老幼.
只要是神的僕人和婢女.
都可以得著聖靈的囂冠.
說喻言未必是預測未來.
我們可以理解成為得到神的啟示.
再向其他人說出神的心意.
鼓勵到和安慰到其他人.
另外在舊約裡.
神也透過異象和異夢啟示他的心意.
所以我們可以知道.

$^{121}$得到聖靈囂冠.
有著新生命的人.
就可以透過聖靈的啟示.
知道神的心意.
然後十九至二十節說到.
神要在天上地下.
行大的神蹟.
有血有火有煙霧.
太陽要變黑.
月亮要變血.
之後就是主大而光榮的日子.
這些現象都是在啟示錄裡.
世界終結之前.
主再來審判世界之前的先兆.
所以我們可以理解到.
由十七至聖靈囂冠屬神的指紋.
直到神在天上地下行大神蹟.
主在來為止.
這段時間所有神的指紋.
都會得著聖靈的囂冠.
而人也可以像二十一節所說.
凡求告主明的就必得救.
大英姐妹.
聖靈降臨在你們身上.
你們就必得著能力作我的見證.
你相不相信聖靈已經囂冠在你的身上.
今日聖靈降臨.
同時囂冠在我們每一個信徒的身上.
意思是今天就是約珥書所說的幕後日子.
在舊約的時代.
神的靈臨在只是選擇性和短暫性給少數的人.
但今日神卻是差派他的靈.
即是聖靈到一切信靠祂的人的身上.
包括男女老幼.
所有年齡層裡面.
去啟示神自己的心意.
這幅圖畫讓我們可以知道.
無論我們在工作上.
我們生活上.
還是任何處境上.

$^{161}$我們都可以透過聖靈.
去得到神的啟示.
讓我們知道神的心意.
所以可能在我們工作上.
當我們在工作上覺得職場很現實.
生活壓力很大.
時間永遠都不夠用的時候.
神其實仍然差派聖靈與我們同在.
啟示祂的心意給我們知道.
讓我們知道神的心意.
不是要我們在崗位上捱累.
成為KPI的奴隸.
成為辦公室複雜人際關係的犧牲品.
而是要在我們專業的崗位上.
察驗神的心意.
然後得著能力成為神在我們公司.
在我們生活上的見證人.
讓其他人可以求告主的名.
甚至年輕人都不需要說自己.
輩分未夠成熟沒有那麼多生活見證.
可以跟別人分享.
或者長輩也不需要以為自己年紀大了.
作見證傳道的工作留給年輕一輩.
因為聖靈已經囂貫在凡有血氣的.
所有神的僕人和婢女身上.
神沒有分男女老幼.
去揀選人啟示祂的心意.
讓我們所有人都可以成為祂在地上的見證人.
讓人可以求告主名.
得著神的拯救.
所以願我們都可以敏銳和信服聖靈的聲音.
遵行神的旨意.
在引用完《若以書》之後.
彼得在二十二至二十四節.
就挑戰當時的群眾去思考.
《若以書》的寓言應驗.
跟他們有什麼關係.
在二十二節.
彼得說到那撒勒人耶穌.
就是神以異能其事神跡.

$^{201}$向你們證明出來的人.
這些事是神藉著他在你們中間施行.
正如你們知道的.
無論當時的人也好.
現在的我們也好.
那撒勒人耶穌.
就是福音書裡面神藉著他.
行很多神跡異能.
這些事都證明了.
耶穌其實是有特殊的身份.
值得留意的是.
向你們證明出來的證明.
是被動式來的.
可以譯作為被設立.
《新按語》的譯本是這樣翻譯的.
神藉著耶穌行神跡其事.
是要向你們證明祂是誰.
所以其實耶穌是神所設立的彌賽亞.
當時的群眾或猶太人.
無論如何都心裡有數.
但是心裡有數的猶太人.
如何對待彌賽亞?.
在二十三節說到.
他們藉著不法之人的手.
將耶穌釘死.
不法之人就是在說.
摩西律法以外的人.
在這裡是說他們藉著.
彼拉多和羅馬的士兵釘死耶穌.
殺死神所設立的基督.
但是看似一切都是猶太人.
和外邦人共同行動去殺死耶穌.
彼得在這裡強調.
其實一切一切都是按著.
神確定的旨意和預知去進行.
其實神一早已經預先知道一切.
而且計劃一切.
目的是要成就祂的救贖計劃.
所以二十四節.
神為耶穌平反.

$^{241}$神將死亡的痛苦解除.
讓耶穌復活.
因為祂(耶穌).
原不能被死囚禁.
死亡的力量從來都不能和神對抗.
弟兄姊妹 耶穌行出神的心意.
在地上宣講天國的福音.
讓被擄的德釋放.
黑暗的德看見.
神也藉著祂行很多神蹟騎士.
去證明祂就是彌賽亞.
但是祂仍然會被猶太人.
藉著不法的人的手釘上十字架.
當耶穌上十字架的時候.
身邊的人還嘲笑祂說.
你救到人 救不到自己.
但是神說.
這些事都是按照祂的旨意和預知進行.
所有事都是祂的掌管之中.
那我們呢?.
如果我們在工作上.
我們在生活上行出神的心意和計劃的時候.
會不會都有機會.
面對很多外在環境的敵視.
會不會被人排斥.
被人嘲笑.
甚至被人誣衊呢?.
可能當我們遇到這些情況的時候.
都會誤以為神根本不在.
神在這個處境沒有祂的心意.
甚至是覺得神根本做不到事.
就好像前幾年DEI的風潮很厲害.
很多跨國的國際企業.
都在自己的公司增設了一個部門.
推廣多元性別同性戀.
視為一個社會責任.
甚至和招聘晉升和培訓掛鉤.
有大兄在北美的公司工作.
都分享到有同事要求其他同事簽一份東西.
去支持同性戀.

$^{281}$他面對的就是.
如果他不簽那份東西.
就會被人貼上落後.
歧視的標籤.
甚至有移民的家長都分享到.
自己面對的就是.
子女的學校推行多元性別.
不參加就會被人說.
被人邊緣化.
面對種種社會的壓力.
再成熟的基督徒要持守真理.
走出真理.
都可能會有信心動搖的時間.
覺得好像神失蹤了.
就好像當年門徒跟著耶穌.
看盡耶穌的神跡.
但見到耶穌上十字架的時候.
仍然四散.
被得仍然三次不認主.
但神在處境上仍然說.
所有事都在我掌握上.
同時我們在之後25至32節的經文裡.
我們見到神給了一個最終極的保證我們.
讓我們可以確信.
無論外在環境多麼惡劣.
神仍然掌管一切.
我們仍然可以靠著神去得勝.
在25至32節.
彼得引用詩篇16篇大衛的詩.
解釋為什麼死亡不可以將耶穌拘禁.
我們先看25至28節.
了解一下詩篇說什麼.
然後再看看這篇詩跟基督復活有什麼關係.
在25至28節.
大衛一開始已經認定他可以戰勝死亡.
因為他說道.
我看見主上在我的眼前.
祂在我的右邊.
使我不自動搖.
所以大衛心靈歡喜.

$^{321}$也不會懼怕未來.
而最重要的是27節.
原來神應許過.
必不將我的靈魂撇在陰間.
也不將你的聖者見撓壞.
陰間和撓壞.
都是指死亡的狀態.
而聖者也是指同神立約屬神的人.
所以兩句平行句.
是表達神不會掉低他在陰間.
他的身體也不會撓壞.
所以他一定會死而復活.
所以28節.
神一早已經啟示了生命的道路.
即是永恆的生命給他.
讓他(大衛)可以在神面前.
風起快樂.
講完詩篇16篇.
彼得在29-32節馬上說出事實.
坦白說.
大衛死了.
埋葬了.
而且他的墳墓直到現在.
都可以清晰可見.
既然大衛已經死了.
身體又撓壞了的時候.
大衛這篇詩一定不是在說他.
落到30-32節.
既然大衛是先知的時候.
他一定可以預先知道.
神對他後裔的安排.
就是詩篇132篇.
神發誓將他一位後裔.
安置在寶座上.
成為彌賽亞.
而耶穌的而且確就是大衛的後裔.
這就證明了兩件事.
第一 大衛預言那位不被撇在陰間.
肉身也不見扭壞的人.
一定就是經歷了死亡.

$^{361}$然後三日復活.
升天的耶穌.
因為神的而且確沒有將耶穌撇在陰間.
耶穌的身體也沒有撓壞.
第二 詩篇132篇的預言.
發生在大衛的子孫耶穌身上.
證明了耶穌就是那位聖者.
就是神向大衛發誓.
將他那一位後裔安置在寶座上的彌賽亞.
所以32節說到.
「因著神應許的實現.
神已經將耶穌復活.
所有曾經見過耶穌復活的人.
包括十二個使徒.
都是預言實現的見證人」.
弟兄姊妹.
彼得用了兩篇詩篇來解釋耶穌的復活.
證明了耶穌是可以戰勝死亡.
也證明了他是坐在為王的那位基督.
雖然我們知道.
耶穌掌管萬有.
連死亡都可以戰勝.
但同時我們也知道.
耶穌是應許我們.
以馬來尼尚與我們同在的神.
所以我們其實都可以靠著.
連死亡都可以戰勝的耶穌.
面對一切一切惡劣的環境.
去成為神的見證人.
耶穌在《馬太福音》十章說過.
「那殺人身體.
但不能滅人靈魂的.
不要怕他們」.
弟兄姊妹.
當我們在生活上.
因為持守真理而面對惡劣的環境.
可能被排斥 可能被嘲笑的時候.
我邀請你們.
抬頭仰望.
回想起這位連死亡都可以戰勝.

$^{401}$掌管萬有的耶穌基督.
然後放心去將處境面對的困難.
祈禱交託.
就像大衛的詩一樣.
「我看見主常在我的眼前.
祂在我的右邊.
使我不自動搖」.
讓耶穌復活的真實.
成為我們面對這個世界潮流的最後盾牌.
可以在惡劣的環境裡.
仍然持守真理 站立得穩.
成為神在這個世界的見證人.
最後在33至36節.
彼得將耶穌的復活.
和耶穌升天聯繫在一起.
彼得在33節說道.
「祂被高舉在神的右邊.
又重復受了所應許的聖靈.
就把你們所看見的.
所聽見的挑冠下來」.
耶穌被高舉到神的右邊.
意思是說耶穌已經升天.
在神的右邊.
擁有尊貴的身份.
而且祂領受了天父所賜下的聖靈.
並且挑冠在每一個信徒身上的果孝.
當時的群眾是聽到的 是見到的.
究竟他們聽到什麼呢?.
就是試圖用別國的話.
宣講出神的作為.
他們看到耶穌釘十字架時.
四散的門徒.
在群眾嘲笑聲之中.
可以勇敢地站起來.
他們聽到三次不認主的彼得.
在眾人面前宣講神的話語.
門徒的生命得到改變.
也可以用他們的屬靈.
因此宣講神的作為.
這些事當時的群眾是看在眼裡.

$^{441}$而最重要的是.
彼得在34至36節再次引用詩篇.
確立了耶穌是基督.
要讓群眾知道.
耶穌才是那一位主.
34至36節是這樣說的.
因為大衛沒有升上天.
但是大衛又說.
主對我主說.
你坐在我的右邊.
讓我使你的仇敵作你的腳凳.
所以以色列全家都應該知道.
你們釘十字架的那一位耶穌.
神已經立他做王為基督.
如果大家心水清.
都可能會記得.
在《馬可福音》裡.
耶穌都運用過這篇詩篇110篇.
暗示自己是理賽亞的身份.
而彼得引用詩篇110篇的邏輯.
其實很簡單.
第一 既然大衛沒有升上天.
耶穌又升上天.
詩篇110篇說到.
坐在神的右邊的一定不是大衛.
而是耶穌.
第二 大衛說到主對我主說.
大衛承認對方比自己大.
而以色列人也稱耶和華為主.
所以我們可以看到這個我主.
的身份一定比大衛大.
而且甚至可能是神.
第三 神留了我主在右邊.
賜一個寶座給他.
使他的仇敵作為他的腳凳.
意思是神給了他審判的權柄.
所以在36節彼得就下了一個結論.
所有以色列人.
作為神的選民.
都應該確實知道.

$^{481}$大衛說的我主.
升天坐在神的右邊的那一位.
就是他們釘在十字架的耶穌.
神已經立了他為主.
他是基督 彌賽亞.
彼得說到這裡就告一段落.
而當時的群眾聽到之後.
聽到自己有份釘死期盼已久的彌賽亞.
他們的反應就是擇心和悔改.
而這就是下篇道的事.
弟兄姊妹.
耶穌升天領受聖靈.
並且囂貫聖靈.
就像《約珥書》應許說到.
在幕後的日子.
神的靈會囂貫在每一個信徒身上.
受囂貫的人會得到神的啟示.
生命得著改變.
另一方面.
就像《以弗所書》四章所說.
耶穌升上高天的時候.
將各樣的恩賜賞給人.
耶穌基督透過聖靈.
分派各樣的屬靈恩賜給他的信徒.
當日試圖恩著受到聖靈的囂貫.
生命得著改變.
可以勇敢地用恩賜.
說未學過的愛國語言.
說神的作為.
今日我們同樣受到聖靈的囂貫.
生命得著改變.
可以運用恩賜為神作見證.
但問題是.
為何有時我們覺得自己的生命.
好像沒什麼改變.
又沒什麼恩賜可以運用為神作見證呢.
可能彼得一再強調.
耶穌升天坐著為王.
成為我們的主.
給了我們一些啟示.

$^{521}$王強調了.
耶穌坐在大衛寶座上.
統治萬國的權柄.
主強調了他絕對的主權.
有權柄去引領我們的生活.
而王和主的共通點.
都是我們要信服的對象.
所以為神作見證的關鍵.
可能就是不再將焦點放在.
我們的限制和軟弱.
和外在環境因素上.
反而是仰賴耶穌差派的聖靈的引導.
讓我們可以信服基督.
在今個星期的《牧者之聲》裡.
我都分享了一個在歐洲旅行的奇遇.
有些人都會知道.
我是一個很內向的人.
雖然沒有人相信.
我很不喜歡去認識陌生人.
基本上要我主動去聊不懂的人說話.
我很多時候都會輕忌.
相反我太太也有人知道.
很外向的.
如果不讓她出聲.
不讓她認識新朋友.
她也會輕忌.
而她也.
淑芬你知道.
很舒服.
而我太太也說得一口流利的國語.
在傳福音上很有恩賜.
當日我和太太.
坐著八人車.
去一個叫Hallstoop的地方.
我坐在司機的旁邊.
太太就在後面和其他人坐.
當去到第一個中途站廁所位置的時候.
太太一下車就說.
旁邊那三個是台灣人.
然後我就問她.

$^{561}$為什麼你不跟他們聊天認識一下.
我太太就說她不想.
但是在去廁所的途中.
聖靈就不斷感動我.
要我主動去認識他們.
但問題就來了.
我又不懂得跟陌生人聊天.
我的國語是差到一個地步.
我回大陸買一杯咖啡.
是會被人叫我講廣東話.
那個程度.
感恩的是.
去到廁所裡.
我看到那個台灣男人.
在找開關的水喉洗手.
原來他不知道那個是自動感應水龍頭.
所以我就把握著機會.
用英文打開了話題.
出去就慢慢認識.
然後邀請他們一起吃飯.
由於太太講得一口流利的國語.
而三個台灣人裡面.
有兩個是不太會講英文.
所以其實在飯桌上面.
主要都是太太跟他們聊天.
不好意思走音.
其實在聊天的過程裡.
我心裡很想有機會跟他們講信仰.
但在這個環境裡.
無端端講耶穌.
其實我自己也有點怪.
同樣感恩的是.
當表明了我們其實都是香港的傳道人之後.
傳福音很強.
講得一口流利國語的太太就找到位置.
起碼有十五分鐘可以分享信仰.
然後討論信仰.
詳情和反思.
大家可以看看這個星期的「牧者之聲」.
不過我想說的是.

$^{601}$為神作見證.
從來不需要我們是完美的人.
我很多限制.
人有內向.
國語又不好.
聖靈一樣感動我.
去認識這三位台灣的朋友.
那時候的環境.
也不是一個可以無端端講信仰的地方.
人也是剛認識.
但主要是神仍然製造機會給太太.
運用她的恩賜.
講述福音.
所以弟兄姊妹.
我們也可以時刻敏銳聖靈的帶領.
相信神對所有的環境都有掌控.
不再看自己的限制.
反而去運用我們的恩賜.
成為神在地上的見證人.
我們一起禱告.
提醒我們要為你作見證.
從來不需要我們是完美的人.
我們求你賜我們一個信服的心.
在面對環境的挑戰和自身的限制的時候.
懂得敏銳聖靈的感動.
深信你這位坐著為王的主.
讓我們今天開始勇敢地為你作見證.
聽我們禱告.
奉主明求 阿們.
\newpage



\section{士師記 3:9}
\label{sec:032245}
\textbf{聖經人聖經事:第一位士師 俄陀聶}~~~~黃克勤牧師博士
\newline
\newline
連結: \href{http://www.church.com.hk/acms/content.asp?site=cdc\&op=show\&type=product\&code=032245\&layout=sermon}{\texttt{032245}} ~~~~ 語音日期: 2026-06-04
\newline
\newline
\hyperref[sec:032244]{\small{< < < PREV SERMON < < <}}
~
\hyperref[sec:index_chronic]{\small{[返順時目]}}
~
\hyperref[sec:032246]{\small{> > > NEXT SERMON > > >}}
\newline
\newline
士師記 3:9
\newline
\begin{longtable}{cl}
\hline
\hline
章節 & 經文 (和合本修訂版)\\
\hline
3:9 & \begin{tabularx}{0.7\textwidth}{X} 以色列人呼求耶和華，耶和華就為以色列人興起一位拯救者來救他們，就是迦勒的弟弟基納斯的兒子俄陀聶。 \end{tabularx} \\ \\
[1ex]
\hline
\hline
\end{longtable}
$^{1}$《聖經人》《聖經事》今天來到一個很特別的人物.
又是一個人,聖經人物是第一位以色列的士司.
在《士司記》裡面的第一位士司就是俄陀涅.
新穎雪很想大家一起讀《聖經》.
今天向大家介紹第一位士司俄陀涅.
這個人物很重要的.
不過在這個經文裡面沒有提到俄陀涅.
做過什麼事的細節.
他怎樣帶兵打仗,他沒有說.
怎樣殺死對方的元首,他沒有說.
怎樣殺對方的王,他沒有說.
只說他勝過了.
看看《士司記》第三章.
循環開始.
以色列人行謠說.
「眼中看為惡的事」.
我看在眼裡覺得你這樣行是很差很錯.
惡的.
「看為惡的事」「眼中看為惡的事」.
「行謠說眼中為惡的事」.
在舊約聖經出現過很多次.
王,人行謠說眼中為惡的事,很多.
以色列人當時入了迦南不久後,約書死了.
長老都死了,連入了迦南的第一代都死了.
於是有新的一代人興起.
他們行耶和華眼中看為惡的事.
究竟怎樣惡法,行什麼不對.
忘記耶和華,他們不對之處.
神很不高興的是什麼呢.
是忘記耶和華他們的神.
第一是忘記了神.
神有那麼多恩典帶領他們仙祖出來.
忘記了,給了他們律法,他們不理.
不記得看,不重要.
而走去除了忘記.
這個惡第一樣就是忘記耶和華.
第二樣是去事奉巴力亞瑟拉.
這個朱巴力亞瑟拉是什麼呢.
都說了他們在迦南地這麼多民族.
他們都有自己當日不同敬拜的神.

$^{41}$通稱巴力亞瑟拉.
這個族有他的巴力,這個族有他的巴力.
巴力亞瑟拉的意思是主或主人.
所以叫做朱巴力亞瑟拉.
即是說各個不同的當地迦南族人.
他們迦南人是有不同的巴力亞瑟拉.
巴力亞瑟拉通常是管農業與生殖.
農業當然是他們很重要,賺錢的.
生殖也是他們很體重.
所以他們就拜巴力亞瑟拉.
怎樣可以得到巴力亞瑟拉的力量呢.
當然是在敬拜的時候和廟祭進行交合等等.
男祭女祭.
亞瑟拉是草樹叢,是掌管財富與幸福.
那麼有趣.
亞瑟拉就是樹叢,為何叫樹叢呢.
很多樹林,樹林是很隱密的地方,看不到裡面.
其實由於都說他們敬拜的時候.
他們會想和神更加親近.
怎樣能夠更加親近,令到巴力歡喜呢.
因為亞瑟拉是巴力的太太.
這個是主,這個是天后.
另外一些經文就是說.
阿瑟塔魯,就是阿瑟塔魯.
我們今天先說這些用的字.
亞瑟拉是樹林,樹叢.
這個阿瑟塔魯是星辰,是閃的星.
是有少少不同的,可能他們兩個太太也不一定.
不過都是亞瑟開始的.
亞瑟拉是樹叢.
由於他們會和那些敬拜的時候的男祭女祭發生交合.
去和神更加親近.
因此交合和女的,就叫做樹叢.
可能因為這樣比較隱蔽一點.
而且樹叢也是一種財富的象徵.
所以財富與幸福,很開心啊在樹叢裡面.
財富與幸福.
耶和華眼中為惡的時候是什麼呢.
是忘記耶和華,第一.
第二,竟然去事奉巴力與亞瑟拉.

$^{81}$這是耶和華眼中為惡的另一件事.
這個就是巴力的神像.
都是拿著一支棒,一支短的棍.
考古的時候覺得這是巴力當時的神像.
不過斷了一部分,只是瘦瘦的.
巴力.
這是巴力,也是拿著一支棒,像我這支棍是粗一點的,是短的.
是不是在這個過程中斷了呢?不知道.
兩個巴力神像是不是很相似呢?.
戴著尖頭的帽子,眼睛大大,一隻手在這裡,一隻手在那裡.
兩個神像都很相似.
這個也是考古出來的一個巴力的神像.
這個巴力的神像有點不同.
管理農業,牛角,拿著植物.
又一把劍.
這是另一個巴力的神像.
可能是不同民族中的巴力是有少少不同的.
這個就是亞瑟拉.
亞瑟拉是女性.
胸部比較大的亞瑟拉.
托著,表示生育,財富等等.
象徵.
這個也是亞瑟拉.
不過很有趣.
她是女性,身材很好,也有個胸.
另外,沒有穿衣服.
還有她的腳,青蛙腳.
很有趣,有幾隻貓頭鷹在旁邊.
拿著一些比較邪惡的物件.
這個就是亞瑟拉的樣子.
耶和華的怒氣於是向以色列人發作.
都是循環的.
他們行爺玩將,捍衛惡的事.
是利器耶和華.
第二,拜巴力亞瑟拉.
於是耶和華的怒氣向以色列人發作.
第二點,把他們交在美索不達米亞王.
美索不達米亞很遠.
是他們最早的漁居地.
有法拉底河,底加里斯河那邊的王.

$^{121}$很有錢,那邊有兩河流域很富裕.
古薩利薩田是一個很遠方的王.
不過勢力很大.
在美索不達米亞那邊有很多錢和軍兵.
一直打到迦南地.
如果你有印象中美索不達米亞河.
有兩條河.
有法拉底河,底加里斯河兩河流域.
中間美索不達米亞是兩條河去到波斯灣.
他們通常很富裕,河米去到北面.
在敘利亞黎巴嫩進攻以色列.
可能這樣就來到了.
這個古薩利薩田真是厲害.
他令到以色列傷亡慘重.
極其困苦,打不過他.
落在古薩利薩田手中.
神將他們交給遙遠的強大軍王.
古薩利薩田的手中.
以色列人服侍古薩利薩田多久?.
八年,這八年做奴隸就這樣做.
經常來哭哭啪啪搶搶殺殺.
他們服侍古薩利薩田八年.
我們說二次大戰,中國抗戰都是八年.
很長,八年,很慘.
古薩利薩田是兩個字.
古山,古山是什麼意思?.
不是很清楚,是地的名字.
利薩田是利薩田.
利薩田是雙倍的邪惡.
雙倍邪惡的古山就是他的名字.
他連名字都惡.
這位就是我們天下無敵.
無比邪惡,雙倍邪惡的王.
古山利薩田.
聽到名字都嚇一跳.
他將名字改為雙倍邪惡.
雙倍邪惡的古山.
古山是地名,可能是出處的地方.
古山利薩田服侍了他八年,很慘.
行爺說一件很惡的事.

$^{161}$第一,忘記耶和華.
第二,他們走去拜巴尼亞斯塔羅神.
於是將他交給一個外面的王.
用古山利薩田成為鞭,鞭打他.
他知道自己壞了,試驗他.
然後服侍了他八年.
第三,以色列循環.
以色列人呼求耶和華時.
「神啊,我們好慘,這個雙倍邪惡的古山利薩田.
對於我們真是好慘,服侍了他八年.
餓到我們死,打到我們這樣.
搶了我們的東西,搶了物件,搶了人,搶了食物.
我們實在好慘」.
哀求耶和華.
耶和華為他們興起一位拯救者.
於是他興起了一個拯救者,一個人.
救他們的是迦納的兄弟.
基納斯的兒子,阿託涅.
是甚麼來的呢?留意一下.
其實以色列第一個先知叫做阿託涅.
以色列的第一個士師叫做阿託涅.
阿託涅.
阿託涅,有親這個L尾的就是葛勒上帝.
上帝的奧夫尼,奧夫尼是獅子.
原來阿託涅這個名字很厲害.
如果你有子女也改這個名.
奧夫尼,奧夫尼L.
奧夫尼L就是上帝的獅子.
阿託涅打敗了古山利薩田.
雙倍邪惡的古山.
阿託涅的父親叫做基納斯.
Kednas.
Kednas是基納斯.
Kednas是獵人.
這個是獅子,這個父親是獵人.
而這個基納斯是很出名的迦納.
迦納就是當日進入迦南地的時候.
約書亞派了12個探子去打聽.
然後有10個回來說不行不行.
那些人很大隻,叫人不能打.

$^{201}$不夠他們來.
只有約書亞和迦納.
這個迦納是誰呢?.
當摩西帶著以色列人預備進入迦南的時候.
就派了12個探子去窺探.
窺探迦南地.
哇,10個探子回來說不行不行.
那些阿納族人很厲害.
很大隻,我們不能夠贏他們.
不要打了.
只有兩個,就是約書亞和迦納.
就說我們去打吧.
有耶和華幫助我們,我們足能得勝.
是這兩個這麼有信心的人.
迦納.
迦納這個人,這個字上次解釋過了.
講到他的時候.
聖經人聖經事,講聖經人有一次是講他.
迦納的意思呢.
沒有人喜歡講這個英文名的.
因為迦納的意思是狗.
狗很忠誠,真是厲害了.
可能他家裡都是獵人家庭.
所以迦納的兄弟.
不知道是哥哥還是弟弟就叫獵人.
迦納的叫狗,兒子叫獅子.
這樣的,上帝的獅子.
俄陀獵.
這裡只是記載俄陀獵,怎麼打法沒有講.
只是說他贏了這麼強的.
美索不達米亞王古山利薩賢.
然後說靈鋼在他身上.
通常都是這樣的.
他就是做以色列的士師.
士師在這裡解釋了.
士師的意思是savart.
士師的字,希伯來讀音savart.
就是法官,長官的意思.
士師就是這個意思.
所以其實他是judge,就是翻譯了法官的意思.

$^{241}$就是和他們判斷事情.
和帶領他們的長官出去爭戰.
耶和華就將美索不達米亞王古山利薩賢.
交在這個俄陀獵手中.
他贏了古山利薩賢.
於是釋放了統治了他們八年的外族人.
於是國中太平四十年.
國這個字其實不是國,是譯錯了.
是地方太平四十年.
於是地方太平四十年.
這裡沒有特別提到什麼支派.
所以很有可能是這麼厲害的黃金與方打來.
是整個迦南地都被他控制了.
八年,所以整個十二支派都很慘.
那地方自從俄陀獵出現之後.
就太平了四十年.
貢獻很大,第一個才知道.
迦勒兄弟的兒子,上帝的獅子.
不過之後俄陀獵死了.
太平之後,那個士師就死了.
四十年都會死的.
一死之後會怎樣?.
正如我們說的惡性循環.
他們又去拜第二個神.
這個俄陀獵這麼厲害.
四十年很太平,因為他在這裡.
大家跟著他的家族敬拜神.
但他死了,又惡性循環出現.
記得,如果問你以色列第一個士師是誰.
俄陀獵是什麼意思?.
上帝的獅子.
他爸爸叫什麼名?獵人.
另外一個兄弟叫迦勒狗.
整間都是獵戶,很有可能.
不過對神很有信心與忠心.
要幫助我們更多認識,一起祈禱吧.
主我們為了俄陀獵感恩.
你興起他,是以色列第一個士師.
可能他帶領的地方包括十二支派.
或者很廣闊的領域.

$^{281}$也沒有說哪一個支派出的士師.
當然他有自己的支派.
但似乎他帶領著整個以色列地.
因為王古山利薩田實在厲害.
主他就是這樣,得勝了.
是因為你的靈在他身上.
主要幫助我們這班服侍你的人.
為主工作的人.
更尋求你的靈在我們身上做更偉大的工作.
以至你的靈管理,你的能力能夠流出.
幫助我們都可以有更偉大的心願.
要成為上帝的獅子.
可以成為別人的祝福,幫助與保護.
主我們見到以色列人會退去,會過失.
會走著惡性循環.
有士師在,他們會好一點.
沒有在,馬上離開你,敬拜別神.
主要幫助我們對你有堅定的信心與跟從.
求主祝福,誰聽祈禱.
奉靠耶穌基督的名,阿們.
\newpage



\section{羅馬書 4:20-25}
\label{sec:032246}
\textbf{因信而剛強}~~~~莊明偉牧師
\newline
\newline
連結: \href{http://www.church.com.hk/acms/content.asp?site=cdc\&op=show\&type=product\&code=032246\&layout=sermon}{\texttt{032246}} ~~~~ 語音日期: 2025-11-02
\newline
\newline
\hyperref[sec:032245]{\small{< < < PREV SERMON < < <}}
~
\hyperref[sec:index_chronic]{\small{[返順時目]}}
~
\hyperref[sec:032247]{\small{> > > NEXT SERMON > > >}}
\newline
\newline
羅馬書 4:20-25
\newline
\begin{longtable}{cl}
\hline
\hline
章節 & 經文 (和合本修訂版)\\
\hline
4:20 & \begin{tabularx}{0.7\textwidth}{X} 仍仰望神的應許，總沒有因不信而起疑惑，反倒因信而剛強，將榮耀歸給神， \end{tabularx} \\ \\ \relax
4:21 & \begin{tabularx}{0.7\textwidth}{X} 且滿心相信神所應許的必能成就。 \end{tabularx} \\ \\ \relax
4:22 & \begin{tabularx}{0.7\textwidth}{X} 所以這也就算他為義。 \end{tabularx} \\ \\ \relax
4:23 & \begin{tabularx}{0.7\textwidth}{X} 「算他為義」這句話不是單為他寫的， \end{tabularx} \\ \\ \relax
4:24 & \begin{tabularx}{0.7\textwidth}{X} 也是為我們將來得算為義的人寫的，就是為我們這些信神使我們的主耶穌從死人中復活的人寫的。 \end{tabularx} \\ \\ \relax
4:25 & \begin{tabularx}{0.7\textwidth}{X} 耶穌被出賣，是為我們的過犯；他復活，是為使我們稱義。 \end{tabularx} \\ \\
[1ex]
\hline
\hline
\end{longtable}
$^{1}$《錄音或錄影》.
此錄音或錄影的版權.
是屬於基督教宣導會不各堂所有.
我們的屬靈傳統之一.
是以真理教導為本.
歡迎您將此錄音或錄影的連結.
傳給親友.
懇請注意.
未經授權.
請勿轉錄或剪輯本錄音或錄影.
甚願好像《仕途行傳》20章32節下所說.
這道能建立你們.
叫你們和一切誠誠的人.
同得基業.
弟兄姊妹早晨.
我們繼續思想羅馬書給我們的訊息.
上星期經文提到.
因信得稱為義的重要基督信仰道理.
因為信耶穌.
我們明明本身是罪人.
但神就看我們為無罪.
賦予我們寶貴的救恩.
我們繼續看著.
這個因信稱義的道理.
其實對我們的生命.
有什麼可以見得到的影響.
羅馬書四章的經文.
繼續以亞伯拉罕的生命作為例子.
講述人去信神對他的影響.
首先我們想看看.
亞伯拉罕這份對神的信.
其實在《創世記》15章.
當時亞伯蘭即是亞伯拉罕.
他當時未有自己的孩子.
但耶和華神就說.
他的後裔要像天上的星星那麼多.
而且他說是亞伯拉罕所親生的.
才成為他的繼承人.
《聖經》說.
即是《創世記》所提到的.

$^{41}$亞伯蘭就是信耶和華.
以此算為他的義.
其實今天羅馬書的經文也是這樣說.
亞伯拉罕相信神.
這就算為他的義.
神應許亞伯拉罕.
亞伯拉罕對神有信心.
其實這份信心.
是很大很大彰顯.
在亞伯拉罕獻兒子以殺的事情上.
神就是要亞伯拉罕在摩利亞山.
其實我們都很清楚.
去獻真正算為繼承人.
是以殺.
《聖經》說明.
神是要去試驗他.
這個時候亞伯拉罕已經經歷了很多事情.
其實他已經很相信.
神是要賜福他後裔梵多.
這個應許.
神說明了.
叫亞伯拉罕去獻以殺.
他就真的遵從神.
他一方面知道神要賜福他的後裔梵多.
這個應許.
但現在又要去獻唯一的兒子以殺.
那他怎麼會後裔梵多呢.
其實亞伯拉罕當時.
知不知道.
不知道.
他不知道神怎樣可以成就他的應許.
亞伯拉罕當時知不知道.
他不知道.
神是會將以殺復活.
是真的會用羊羔去代替以殺.
還是神再令亞伯拉罕.
即超過一百歲的亞伯拉罕.
又再生一子呢.
其實亞伯拉罕是完全不知道.
完全不知道.

$^{81}$亞伯拉罕只是相信.
然後跟從.
他相信神.
在《創世記》.
這裡的結果就是耶和華神再次肯定他.
再應許他.
甚至在聖經三處的經文.
包括這裡.
以賽亞書的經文.
亞伯拉罕是被稱為神的朋友.
被稱為神的朋友.
今天羅馬書講到這一份信.
不單止是要讓亞伯拉罕得到好處.
亞伯拉罕是一個榜樣.
經文講到.
「算他為義」這句話.
不是單單為他寫的.
也是為我們寫的.
我們就是那些要被算為義的人.
我們的信.
但是我真的覺得.
其實比起亞伯拉罕要去信.
好像簡單一點.
我們信的其實已經列明在聖經裡.
是歷史發生過的事.
這裡講到.
核心就是去信神使我們的主耶穌從死人中復活.
然後他也講到.
信的內容是.
耶穌被交出去受死.
為了我們的過半.
耶穌復活是為了我們清兒.
耶穌為我們受死.
也為我們復活.
其實這裡聖經清楚講到.
是耶穌受死復活.
其實是他的功勞.
去讓我們清兒.
所以不要太強調是我們的信.
要有內容.

$^{121}$我們才能夠信.
再看回第20節的經文.
他說人可以不信.
或者信會帶來不同的後果.
亞伯拉罕沒有疑惑.
他真的一步一步去跟從.
大家想想.
其實他走上山的路程去獻耳撒.
是很不容易的.
內心一定有很多的煎熬.
很多的掙扎.
但是他很剛強.
最終將榮耀歸還給神.
其實我們所相信的內容更加豐富.
比起亞伯拉罕所相信所知道.
我們所相信的是耶穌死而復活的內容.
我們是否都可以因著這份信.
生命一樣的剛強呢?.
是否有真正信耶穌的這份傾向呢?.
而且這一段亞伯拉罕的事跡.
也說明了一件事.
其實信.
很多時候我們是不會知道將來會發生什麼事情.
是會發展成怎樣.
我們不知道.
可以是一個我們不期望的結果.
超出我們想法的結果.
我們是沒有把握這件事會怎樣完成.
怎樣完結.
是完全沒有把握的.
但是信.
只是在信對象.
是一個我們信靠的對象.
我們認識這位信靠的對象.
知道祂是愛我們.
並且祂大有能力.
即使事情發展成怎樣都好.
信靠的對象.
祂是不會改變.
祂都是愛我們.

$^{161}$有能力的那位神.
弟兄姊妹,其實人生有很多起起跌跌,高高低低.
我們不要被事情.
被環境打倒了自己.
我們是不能夠對事情有把握.
有任何把握.
但是我們對信靠的對象是這位神.
就很有把握.
我們是對著這位神很有把握.
有祂.
我們的生命就可以顯得剛強.
其實真的有很多生命很剛強的信徒.
他們一生為神打美好的仗.
就算遇到一些難堪的事.
他們仍然努力信靠神.
其中一位令人敬佩的是澳洲人Nick.
尼克·胡捷.
我相信大家應該聽過他的事跡.
他天生沒有四肢.
但是他是一位基督教的報道家.
帶領很多人去信主.
他創立了勵志演講的公司.
他演過短片,出過單曲,音樂視頻等等.
其中一次他和神之間很重要的經歷.
是他看到聖經約翰福音第九章.
講到耶穌說那位身懷黑眼的人.
是要顯明神的工作.
他看到神在他的生命有他的計劃.
他生命很有毅力.
原來他十九歲的時候.
向學校推銷自己的演講.
一直一直被拒絕.
終於到了第幾次呢?.
第53次才成功開始了他的演講生涯.
他也成立了Life Without Limbs的事工.
積極去傳揚福音.
包括監獄,學生,工作等等.
現在仍然一直去為神去努力.
去為神去發光.
其實我們每一個人的生命很不一樣.

$^{201}$但也可以因著信靠神.
一生在各自的崗位上崗強.
在家庭裡面,職場,教會的服事,社區.
在各個不同的角色,崗位上去為主去發光發熱.
我們作為信神的人,就是對神有這份信.
生命就能持續地去崗強.
而且有這份信,聖經在告訴我們.
我們是永永遠遠在這份恩典之地.
經文說,我們憑著這份信.
就得以進入現在所站立的恩典當中.
其實得以進入這個字.
原本的文字是被引進的意思.
我們是被引進去到達這位尊貴的上帝面前.
是愛我們的神.
而另外被引進,其實就是要有一位引見者.
這位引見者就是耶穌基督自己.
這份與神和好的恩典.
直接使我們與神的榮耀.
特別是將來神的榮耀就有關.
神以後的榮耀,耶穌基督再來.
完全的管治.
我們是直接受惠者.
因為是愛我們的父.
去榮耀的管治.
再沒有任何邪惡的勢力.
沒有任何患難,沒有任何悲傷.
這份是你和我已經有的盼望.
是很實在的盼望.
其實為什麼這個時刻說這份盼望這麼重要呢?.
是因為經文說到人生當中.
我們必定會有的患難.
一定會有的患難.
其實說到患難,我們一般都不會喜歡的.
馬上就想避開.
但是經文竟然說.
信神的人能夠以患難誇口.
經文令我們感到很詫異.
它竟然帶出一份很積極的信息.
不是很不情願.
有患難的時候.

$^{241}$神給我們力量去面對.
不是單單這樣.
經文說到竟然可以因為患難而歡喜.
以患難誇口.
其實呂鎮忠的譯本都是說.
苦難而歡喜誇盛.
英文譯本都是一樣.
為什麼可以這麼積極呢?.
其實正正是和我們剛才所說的.
我們有的盼望有很大的關係.
首先說一下這裡的患難是什麼.
它一定不是說我們人因犯罪之後.
得罪了神,然後可能跌倒等等.
反而它有機會指到人因為跟從神.
而在世上有機會遇到的患難.
經文這裡說患難是生出人來.
同樣的忍耐這個字.
不是想說一種迫於無奈的忍.
不是被動或者負面.
這個字原文是有一份積極的信息.
是預備好的去堅持去堅忍的力量.
然後它說生出老年.
其實老年這個字是說一種性質.
本質或者品格.
英文是character.
在聖經裡面有很多經文是說.
信徒在患難裡面經過試驗.
生命就生出很寶貴的特質.
是內在的特質.
於我們是有益的,並且得到神稱讚.
就好像新約的這兩句經文.
信心經過試驗.
讓生命更加的寶貴.
但是這還不是這段經文最終想要表達的.
經文還有一句是說老年生出盼望.
這份基督徒真正跟從神的人.
生出來的性格特質是生出盼望.
其實最終是想指向盼望.
是一份不會落空.
不會使人羞愧的盼望.

$^{281}$當然這裡不是說基督徒熬煉出來的品格.
讓我們才有這份盼望.
不是的,這份盼望是一直存在.
剛才說過是因為基督的拯救.
讓我們有這份將來同德榮耀的盼望.
這裡經文是說老年生出盼望.
其實就是我們在神裡面.
那個生命不斷的生命成長.
但是知道根本不是我們的作為.
其實是在見證神在我們生命的工作的時候.
我們在自己的生命裡面.
見到神的工作.
我們就知道神的工作仍然要繼續.
他仍然要成就將來他必定得著榮耀.
而我們是有份的,要同德榮耀.
這份盼望我們是更加的確實.
更加肯定,其實一切仍然是相信神.
相信神現在在我們的生命當中工作.
也相信神將來要有榮耀的工作.
得著榮耀.
其實這個段落的經文.
給我們全新的一個眼光.
去看什麼是患難.
看各種的患難.
我們從來遇到患難當然是想避開.
盡早脫離患難.
這個是人之常情.
我們知道神也不是最終想我們有受苦.
患難不是終極.
將來也一定不會有這些苦.
但是地上有時有些患難是避不了的.
可以突然間有患病.
可以突然間家人有情況.
更加可以是我們因為堅持這份信仰.
就受到逼迫患難.
其實這個世界還有很多地方.
是因為有人要跟從耶穌基督.
而要受到患難.
就算你在職場上你去為主作見證.
有時因為信仰的緣故不去同流合污.

$^{321}$可能被人杯葛.
在這樣的環境工作你每天都覺得難.
其實都是患難.
但是經文想告訴我們.
患難其實不可怕.
我們是已經有永恆的盼望.
患難也拿不走我們這一份的恩典.
拿不走這一份榮耀的盼望.
相反在患難裡面.
就這段經文來說更加可以去跨性.
但是我想說這段經文.
只可以是基督徒當事人去領受去體會.
千萬不要運用這些經文在其他信徒身上.
或者我們放輕了其他人的患難.
我們不要.
真的作為信徒去體會這段經文.
這段經文是說當事人.
基督徒能夠積極去面對.
因為患難是真正驗證了.
跟隨神的信徒的這一份生命.
這竟然是一種互動.
信徒透過自己的生命.
就更加相信到神的工作.
也更加確信這份以後的盼望.
有一位經常帶著很多疑問的弟兄.
他經常去問神.
世界為何會有這麼多的苦難和患難.
他一直不太接受.
有一次他跟一位宣教士兩個人聊天.
這位宣教士回應神的呼召.
已經在一個很落後的地區.
服侍了神二十多三十年.
他們聊天的時候.
宣教士當然不知道.
弟兄有這麼多疑惑.
對神的疑惑.
宣教士正在分享自己的故事.
他分享到三次.
他自己三次很艱難的經歷.
三次都差點死的經歷.

$^{361}$都是在宣教工場發生.
有一次是弱疾在落後地區.
一次是新冠病毒肺炎.
是幾年前發生的.
一次是嚴重的癌症.
其實這位對苦難帶著很多疑問的弟兄.
很驚訝他的經歷.
有這麼多的患難.
但他最驚訝的是.
為何這位宣教士仍然可以這麼積極.
在宣教工場裡事奉.
無怨無悔.
其實這位信徒正在想.
如果經過這麼多的患難.
即使你繼續信神的時候.
都可能繼續勉勉強強.
不是對神很積極.
為神的國度很賣力的一位信徒.
他問這位宣教士.
宣教士回答.
他相信神不是一位要人受苦的神.
他是在信神的本性.
他在信神的對象.
不是因著環境.
不是因為變化.
又或者因為疾病,痛苦.
而影響他的信心.
是的,這個世界其實有很多的苦難,患難.
也有不少可以發生在基督徒的身上.
但我們相信的確可以讓我們很堅強.
和世人不一樣的.
是因為我們對神的信.
是這一份榮耀的盼望.
我們是絕對有把握的.
這一切都可以使我們更加正面,更加積極.
去面對生命很多各樣不同的患難.
其實是不容易的.
但是憑信我們可以一起去努力.
經文又再進而提到更多神的愛.
提到更多神的愛.

$^{401}$神的愛當然和剛才說的信和盼望絕對有關.
我們相信神的本質是愛.
他深切愛我們的時候.
我們就更加能夠信靠他.
而這段經文也提到神藉著聖靈.
讓我們更加深去體會神的愛.
以至於我們更加確信這份盼望.
肯定榮耀的盼望不會落空.
經文描述的是一份怎樣的愛呢?.
其實是我們很熟悉的經文.
但是這個是新漢語譯本.
不如邀請大家一起讀.
新漢語譯本我們少一點一起讀這段經文.
原來我們還軟弱的時候.
基督就按鎖定的時候.
為不敬虔的人死.
為義人死是少有的.
為好人死或許有人敢做.
但基督卻在我們還是罪人的時候.
為我們死.
神的愛就在此向我們顯明了.
羅馬書五章六至八節.
這裡為義人死的義人.
是守規則守律法的人.
而為好人死的好人.
是一個有愛心關心人的人.
也就是說他也有這些個體.
願意為義人或好人去死.
但是經文是在說.
我們原本是一個什麼狀態.
是不是在說的好人義人呢?.
不是,這裡是說我們原本的狀態.
是還軟弱的時候.
是不敬虔的人.
還是罪人的時候.
其實軟弱就是不能做得好.
其實三個詞都有接近的意思.
加強了想表達的內容.
我們都是得罪了神.
而且是一直在得罪神.

$^{441}$但竟然在這個時候.
基督就是神自己為我們死.
弟兄姊妹,這就是神對我們的愛.
救贖很核心的意義就彰顯了出來.
不是我們自己配得.
我們是一點都不配得.
是罪人的時候神愛我們而犧牲.
成就了很重價的救贖.
我們是靠著基督的血.
輕易被定為無罪.
而且本來聖潔的神.
其實他對罪人犯罪是會憤怒的.
但我們也因著基督的救贖.
免除了神的憤怒.
是很大的恩典.
經文也多次重申.
我們這一刻其實已經和神和好了.
弟兄姊妹,這個不可能是我們單方面.
可以做到的事情.
是神定義去拯救我們才可以.
是神的意願,是神的動機.
這裡更加說我們是以神為誇口.
的確是的,是祂的愛,祂就去拯救.
是祂的能力,祂必定能夠去承傳.
我們現在的堅信無法誇口我們自己.
其實我們由一開始的信.
到現在一直繼續去相信.
其實都是神持續不斷的工作.
是祂的工作,不是我們.
經文已經說了.
是神透過聖靈讓我們不斷體會到神的愛.
而同樣的,我們也知道是耶穌基督.
也不斷一直為我們代求.
我們就能夠一直去堅信.
其實既然持續都不是我們自己的工作.
我們就是真的一點都沒有可以誇.
況且其實是神的工作.
對我們是只有好處,沒有壞處.
為什麼?是因為是神的工作.
我們就不怕會失去這一份的信.

$^{481}$因為我們倚靠的是神,是神的工作.
我們試想想,如果我們是要靠著自己.
才能夠繼續去相信.
這個信肯定是不穩陣的.
我們自己怎會靠得住呢?.
我們想一想就知道.
直到我們將來去到天家,將來去到見主面.
都仍然是神的工作.
因此我們可以常常以神去誇口.
是神的功勞,我們以神為榮.
以神祂奇妙的救贖工作為榮.
基督徒就是可以有這一份被神愛充滿.
以神為榮的生命.
其實被愛充滿,以神為榮的生命會是怎樣呢?.
其實我想起一個小朋友.
他如果感受到父親是最好的.
夢愛的孩子,他感受到父親的愛.
他知道父親是最好的,最頂呱呱的.
這位孩子感受到爸爸的愛和保護的時候.
他會有很大的安全感.
更何況我們這位父親有無盡的能力和權能.
父親亦會看緊自己,看顧自己.
這位孩子亦會從心底裡感到絕對的滿足.
知道父親能夠滿足他,給他最合適的東西.
他有什麼需要他就可以找爸爸.
而且這位孩子以他父親為榮.
覺得父親所做的事都是好的,沒得彈的.
你想想這位小朋友會怎樣.
他真的有時會忍不住跟其他人說.
他爸爸有多厲害,做事做得多好.
他煮食又厲害,木工又好,水電又厲害.
又懂得做運動給我玩,陪我玩,很多很多.
他真的以他爸爸為榮.
其實一個蒙愛,以神為榮的生命就會很不一樣.
這一定是一個很堅壯的生命.
靈姐妹,我們自己的生命又會怎樣呢?.
你會想到像這位蒙愛的孩子一樣.
感受到父愛,以父親為榮.
還是有時我們仍然有很多的擔心在裡面.
可能在各樣的不滿足,不滿足的裡面.

$^{521}$仍然在找很多愛做的事,想去填滿自己的生命.
經文在提醒我們,其實我們的父.
我們真的有一位愛我們的父.
我們有什麼的時候,我們可以去這個父裡面去找就夠了.
若然我們未深刻去領受這份神的愛.
未深刻去體會我們所相信的神這位全能的.
做奇妙工作的上帝的時候.
求神,求聖靈不斷去幫助我們.
讓我們更相信,更深體會,經歷聖經的真理.
神的愛已經徹底顯明了,在這份的救恩當中.
當然不單止我們能夠在救恩裡面去領受這份的愛.
我們其實在日常生活裡面都能夠體會得到.
你又如何去體會這份的愛呢?.
是覺得很甘甜,還是覺得很軟?.
其實我真的覺得,這位神是我們無論大大小小的事.
都很願意去愛護我們,去看顧我們的神.
在最近一次旅行,其實我一家人,五個人.
每年暑假都會去一次旅行.
就在今年旅途的第二天,我們回到酒店.
太太就幫小孩洗眼鏡,我們三個小孩都帶著眼鏡.
洗完哥哥的眼鏡,哥哥去拿回眼鏡的時候.
他發現,托著鼻樑那塊膠沒有了.
那塊膠不是有個螺絲鎖住的那種.
是好像入榫的那種膠.
很可能是洗眼鏡的時候弄掉了.
然後太太就看到酒店的洗手盤,原來那個坑很大.
她就想,可能已經沖走了,那塊膠已經沒有了.
我們都盡力嘗試找找,但是都找不到.
我心想,如果哥哥整個旅行都看東西都是糊糊的.
太可惜了,我就真的去祈求神幫助.
在心裡面不停地祈求.
神就給了我一個意念,就不去洗手盤那邊找.
就去方言鏡的桌子附近找.
終於很快就在這張桌子附近的地上找到了.
我不知道你們是不是在不同的生活裡面.
我們都有去體會著神的這份愛.
我們可以浸淫在神的愛當中.
不單止在最大的愛,在這份救贖裡面.
我們已經看到了.
這份愛是想讓我們生活當中都去體會.

$^{561}$其實這份恩典,我們真的是一份恩典.
我們沒有任何東西可以誇.
只可以在一生當中,我們盡情去誇耀我們的神.
以我們的父神為榮.
是因為愛我們,是因為有能力的父神.
弟兄姊妹今天羅馬書的經文.
他再一次提醒我們.
其實這份救贖已經在這裡.
我們能夠一生去相信祂.
並且生活,我們平時的日常生活.
可以顯得剛強,在我們不同的崗位上面.
然後我們真的有一份永恆,不改變的盼望.
有著這份盼望,幫助我們可以經歷人生當中的高高低低.
我們憑信經過,甚至有患難的時候.
我們都不怕去面對.
亦都很願意我們每一個人都能夠上上去感受到神的這份愛.
願神去幫助我們去經歷這份愛.
去體會這份救贖的愛.
更加日常的生活當中,繼續去經歷這份愛.
以我們這位愛我們的父神,是有能力的父神.
上上以祂為誇口,以祂為榮.
我們一起低頭,有個祈禱.
是你賜下寶貴的救恩.
是當我們仍然不配的時候.
仍然有得罪你的時候,你已經賜下這份救贖.
讓我們能夠全然的去相信你.
求主你在我們每個人的一生當中.
繼續上上去加我們力量.
亦都讓我們即使在生命裡面有不同的高高低低.
我們仍然可以靠著對你的信心去跨城.
我們仍然能夠靠著對你的信心.
我們能夠去面對這些困境,可能是患難.
甚至是因著為了見證你的緣故.
主我亦都求你,上上以你自己的愛.
真的去浸淫我們.
讓我們去體會你這份救贖的愛.
體會你無時無刻看顧我們,愛我們.
你就是愛我們這位的父神.
我們以你爸爸為榮.
多謝你祈禱奉主耶穌基督的名,阿們.

\newpage



\section{創世記 12:1-9}
\label{sec:032247}
\textbf{放心……你的信治好你了}~~~~李均熊牧師博士
\newline
\newline
連結: \href{http://www.church.com.hk/acms/content.asp?site=cdc\&op=show\&type=product\&code=032247\&layout=sermon}{\texttt{032247}} ~~~~ 語音日期: 2026-06-07
\newline
\newline
\hyperref[sec:032246]{\small{< < < PREV SERMON < < <}}
~
\hyperref[sec:index_chronic]{\small{[返順時目]}}
~
\hyperref[sec:032248]{\small{> > > NEXT SERMON > > >}}
\newline
\newline
創世記 12:1-9
\newline
\begin{longtable}{cl}
\hline
\hline
章節 & 經文 (和合本修訂版)\\
\hline
12:1 & \begin{tabularx}{0.7\textwidth}{X} 耶和華對亞伯蘭說：「你要離開本地、本族、父家，往我所要指示你的地去。 \end{tabularx} \\ \\ \relax
12:2 & \begin{tabularx}{0.7\textwidth}{X} 我必使你成為大國，我必賜福給你，使你的名為大；你要使別人得福。 \end{tabularx} \\ \\ \relax
12:3 & \begin{tabularx}{0.7\textwidth}{X} 為你祝福的，我必賜福給他；詛咒你的，我必詛咒他。地上的萬族都必因你得福。」 \end{tabularx} \\ \\ \relax
12:4 & \begin{tabularx}{0.7\textwidth}{X} 亞伯蘭就遵照耶和華的吩咐去了；羅得也和他同去。亞伯蘭離開哈蘭的時候年七十五歲。 \end{tabularx} \\ \\ \relax
12:5 & \begin{tabularx}{0.7\textwidth}{X} 亞伯蘭帶著他妻子撒萊和姪兒羅得，以及他們在哈蘭積蓄的財物、獲得的人口，往迦南地去。他們就來到了迦南地。 \end{tabularx} \\ \\ \relax
12:6 & \begin{tabularx}{0.7\textwidth}{X} 亞伯蘭經過那地，直到示劍地方，摩利橡樹那裡；當時迦南人住在那地。 \end{tabularx} \\ \\ \relax
12:7 & \begin{tabularx}{0.7\textwidth}{X} 耶和華向亞伯蘭顯現，說：「我要把這地賜給你的後裔。」亞伯蘭就在那裡為向他顯現的耶和華築了一座壇。 \end{tabularx} \\ \\ \relax
12:8 & \begin{tabularx}{0.7\textwidth}{X} 從那裡他又遷到伯特利東邊的山，支搭帳棚；西邊是伯特利，東邊是艾。他在那裡又為耶和華築了一座壇，求告耶和華的名。 \end{tabularx} \\ \\ \relax
12:9 & \begin{tabularx}{0.7\textwidth}{X} 後來亞伯蘭漸漸遷往尼革夫去。 \end{tabularx} \\ \\
[1ex]
\hline
\hline
\end{longtable}
$^{1}$世英前平安.
這幾天這麼熱,有沒有熱到病?.
站在這裡也有點頭暈.
希望不會待會睡在這裡.
我們一起禱告預備心靈.
去聆聽上主的說話.
主在這幾天這麼酷熱的天氣下.
求主保守我們的身體.
也保守我們的心靈.
叫我們能夠常常的預備好.
去領受你的幫助.
領受你的恩典.
求主你在這一陣子時間與我們同在.
幫助我們.
讓我們心靈都聽到你的說話.
又叫我口中的言語心裡的思念.
當說出來的時候就成為你的道.
去幫助弟兄姊妹.
我們禱告,交託,仰望.
是奉主耶穌基督的名字祈求.
阿們.
今天我想和大家一起看看.
《經》課裡面的.
患十二年血流的婦女那一段的記載.
我想大家都很熟悉這段經文.
一個婦女患了十二年的血流.
或者在《新易和修本》裡面就說.
經血不止十二年.
她就偷偷在耶穌的背後.
想摸一摸耶穌的衣裳,衣裳,鎚子或三角.
想著摸一摸就會康復.
就會痊癒.
她真的痊癒.
而耶穌就轉身來問婦人.
或者我們今天的講題.
「放心,你的信治好你了」.
或者在《和本》裡面就譯作.
「放心,你的信救了你了」.
我稍後也會解釋一下這些用語的問題.
在《馬太福音》裡面.

$^{41}$其實有很多內容是重揀的.
如果大家想再了解一下.
譬如婦女在十二年的經驗等等.
你可以看《馬可福音》或《路加福音》同樣的記載.
不過其實這一段經文.
大家雖然很熟悉.
我不知道大家一邊看的時候.
會不會也有一個問題.
就是為什麼這個婦人要偷偷摸摸地接近耶穌呢?.
為什麼要這麼偷偷摸摸呢?.
我們很容易就會知道.
或者相信這段經文是講信心的.
摸摸衣服就會康復.
不過她的行動其實是十分費解的.
她沒有走到耶穌的面前.
也沒有開口去求.
她是等待沒有人注意她的時候.
在人群之中.
在後面去靠近.
然後想著輕輕摸摸衣角.
事情就會解決了.
她就可以離開了.
這個動作本身.
其實會不會也令我們要問一下.
為什麼她要這麼做呢?.
剛才我說大家可能也很熟悉這個故事.
傳統的理解認為.
在摩西律法裡面.
或者猶太人的習俗裡面.
血漏不止或者所謂經血不止或血漏病的身體狀況.
在律法上是視為不潔的.
不潔的人最主要的問題就是.
他不能在會堂或任何宗教場合出現.
因為這個不潔會傳染給其他人.
它不是一種疾病.
但是一種禮儀上的不潔.
他們當時認為不潔的人不可以參與祭祀.
參與這些宗教聚會.
而這種不潔是會傳染給別人的.
所以在這樣的情況下.

$^{81}$他就不能夠公開地接觸別人.
這個背景其實非常重要.
因為它幫我們明白到.
他在這個群體當中的位置是怎樣.
不過如果我們只是停留在因為他不潔所以他偷偷摸摸地接近耶穌.
去看這個故事的話.
我們其實就忽略了這個敘事或者這個故事.
這個醫治故事的一個很重要的張力.
因為它不是純粹因為律法規條所限制.
而可能其實他已經習慣了這樣的情況.
但不要期望別人見到他.
隱藏在人群當中.
他不是太相信自己會被別人正面地接納.
他不是完全拒絕幫助.
他只是容許自己.
用一種不驚動別人的方法下去得到幫助.
我不知道大家在平日或者日常的日子裡會不會遇到這樣的情況.
譬如我在宿舍的社監.
我經常會接觸到一些學生.
在宿舍裡面的學生.
現在的大學生,年輕人很多時候都會失眠.
或者我們有一個比較好聽的學名.
就是睡眠混亂症.
就是應該睡的時候他不睡.
他想睡的時候就睡不著.
要起床的時候他又不肯起床.
所以他不是失眠那麼簡單.
他是有一些睡眠混亂的情況.
但這些學生通常都不會很直接去問或者去尋求協助.
可能他會去診所.
就是中大診所.
他純粹說你給我一粒窩棉藥就可以了.
但當然醫生當然會問為什麼你會失眠.
什麼因素,壓力大,打機打得太多還是怎樣.
但他們不是很想和別人談他們的情況.
你給我一粒安眠藥我可以了.
我睡得著覺就可以了.
我不知道大家有沒有遇到這樣的情況.
又或者其實可能你家裡有個老人家.
有什麼事的時候.

$^{121}$就去買些中成藥或者西成藥.
總之買些去藥房的.
我怎樣怎樣去不到廁所.
自己吃了藥就算了.
就不會去看醫生.
美其名說我不想打擾別人.
又或者不想浪費錢.
但其實這是什麼心態呢.
不過除此之外我想提醒一下大家.
如果我們用今天的處境去理解這段經文的時候.
剛才是一個可能性.
但很容易我們會不知不覺地將故事讀成.
其實它是給一些特殊的人.
好像在故事裡面.
那個患血流12年經血不止.
他的病難以啟齒.
他覺得會影響別人.
所以他不去公開或者隱藏自己.
如果我們將這種說法.
變成指一些很特殊的人才對的話.
這段經文.
那我們就會誤解了這段經文.
或者將這段經文推到和我們離得很遠.
因為我們不是這樣的人.
我沒有什麼難以啟齒的疾病.
或者我不是這樣的狀態的時候.
這段經文就不適合我.
我們不去找一個什麼人是最像.
所謂患12年血流.
或者12年經血不止的女士.
我不是這樣.
所以這段經文好像與我無關.
我們就將這段經文慢慢變成.
這是他們的問題.
與我沒有什麼關係.
但其實你這樣看的時候.
就忽略了這段經文.
這個婦人的心態.
其實與今天我們很多人的心態很相似.
你可以說其實這就是人性經驗的一個狀態.

$^{161}$可能是現代人特別會更加似這種狀態.
或者甚至乎一會兒講下去.
可能大家會想到.
可能是一些後COVID的狀況.
因為在COVID的時候.
我們習慣了關上門自己處理好.
有些什麼東西.
不需要有什麼社交.
原來我們都能夠生存.
我們發覺我們可以關上門.
不用理會別人.
別人不用理會我.
不要讓別人理會.
有什麼問題.
上網問你AI就能解決.
我們有點肚子痛.
或者便秘地問你AI怎樣.
譬如剛才我說的失眠的問題.
那些同學真的上網問一下AI.
或者上網看看.
他們就有一些秘方.
話說你將那些退熱貼.
捆住在手指尾那裡.
你是會睡得好一點的.
你就會沒有失眠.
那些同學就說是很好的.
我沒有失眠的狀況.
所以我沒有試過.
我通常都倒頭大睡.
所以神學樓火警響的時候.
我都是不會起床的.
但是他們連去醫院.
拿一顆安眠藥都不需要.
自己搞就行了.
我自己上網.
不需要和別人任何相處.
或者討論交代任何東西.
我上網都可以解決我的問題.
於是耶穌就變成AI了.
摸一摸按一按.

$^{201}$就打算解決我的問題了.
這個會不會是我們普遍人的問題.
或者普遍人的狀態呢.
這個經文不是要我們找到.
誰最像這個婦人.
其實我們要意識到.
我們人在很多情況之下.
都有一種傾向.
就是願意去接近幫助.
同時要避免被看見.
是不是.
大家覺得是不是.
我們回到這個婦人的情況.
大家可以看看經文.
我們看到的.
她不只是有沒有信心的問題.
也包括她怎樣接近耶穌.
她想改變自己的處境.
但她不想被人看見.
被人帶出來.
她願意伸手.
但未必準備好被人認出.
她這個動作其實很微小.
很隱藏.
但反而她這個這麼小的動作.
就將她的心態放大了出來.
令我們更加去問為什麼她會這樣做呢.
如果你將這段經文.
放回整個敘事.
《馬太方九章》的脈絡裡面.
你會看到其實它夾在另一個故事的中間.
就是那個管會堂的主管.
在其他福音書裡面就稱呼他做「艾魯」.
他女兒急病.
可能甚至他認為她已經死了.
或者將近死亡.
於是他就找耶穌.
叫他快點來我家裡救她.
使她再活過來.
這個人公開地去耶穌面前.

$^{241}$甚至經文說他跪在拜耶穌.
在另一些福音書裡面說他跪在耶穌面前.
跟耶穌說我女兒剛剛死了.
或者將近死亡.
不同經文有不同的重點.
然後他公開地要求耶穌.
邀請他來到他自己的家裡.
來接受幫助.
如果這樣看的話.
他跟那個婦人的情況.
就剛好有一個很大的對比.
兩個畫面擺在我們面前.
一個偷偷摸摸.
不想被人見到 被人知道.
但是這個管會堂的.
他是很公開地 很大膽地.
去到耶穌面前邀請他來到他家裡.
甚至他家裡已經有一大群人在等待.
完全是很公開的.
是很社會性的.
跟這個婦人的動作.
心態正正是產生一個很強烈的對比.
一個是公開呼求.
一個是隱藏地接近.
一個是邀請耶穌進入他的家庭裡面.
一個想一想 摸一摸就打算離開.
我不知道大家是否體會到.
這兩個故事放在一起.
是否真的想帶出這種對比的張力呢.
回到那個沒有名字的婦人.
經文說她心裡想.
「我只要摸他的衣裳就必存如」.
其實這樣看 她是很有信心的.
摸摸衣裳也可以.
耶穌是有能力的.
她很相信耶穌的能力可以幫到她.
甚至她相信只要摸一摸耶穌的衣裳.
她就能夠得到這個能力.
這樣說當然沒有錯.
但是她的信心.

$^{281}$她好像停留在某一個狀態.
或者某一個地方.
她相信耶穌可以幫助她.
但是她又選擇一個.
不需要進入任何關係的方式.
去接近耶穌.
她可以觸摸耶穌.
但是她不願意耶穌觸摸她.
不願意耶穌看見她.
這個敘事本身的結構.
很多解經家都被我指出.
它是一個甲敘的效果.
之前之後就是艾魯或者管會堂女兒得救的故事.
然後突然間就插入了這個婦人的敘事.
當中產生很多的敘事張力.
或者很多敘事的技巧.
我們今天不詳細說了.
但是婦人的故事被放在會堂主觀的故事中間.
一邊是女兒急性即時的死亡.
一邊是一個患了流血經血不止十二年的故事的身體.
一個是急性的死亡.
一個是慢慢的流失.
還有其實很多東西可以對比的.
他們的社會地位.
他們本身公開和隱藏的狀態.
所以其實這兩個對比.
讓我看到不是有沒有信心的問題.
不是比較管會堂的人信心大一點.
還是患血流婦人信心大一點.
他們不是在比賽信心大小.
而是說究竟你有信心的時候.
你是怎樣去呈現這種信心的呢.
你是怎樣去接近耶穌.
怎樣去尋求耶穌的幫助的呢.
這個患血流婦人.
我們到現在都不知道她叫甚麼名字.
她雖然期待改變.
但是她是站在一個很安全的距離.
走得遠遠.
她不用開口.

$^{321}$不用對話.
不需要任何人知道.
也不需要耶穌那一刻跟她說甚麼.
她想著只要伸手她要得到摸到就會離開.
這裡就令我們想究竟她所追求的.
所謂「存余」是甚麼意思呢.
她說摸一摸我就會「存余」.
其實這個「存余」是甚麼意思呢.
甚麼叫真正的「存余」呢.
我講到這裡.
我不知道大家能不能夠對號入座.
在今天的信徒的處境裡面.
其實我們對這個狀態完全不陌生.
我們很熟悉信仰.
我們很熟悉信仰的語言.
我們懂得甚麼時候去禱告.
我們知道神可以幫我們.
我們知道有困難的時候.
你應該去到神的面前.
但是與此同時.
我們又未必習慣讓信仰真正進入那些最深.
最難講,最怕被人知道的地方.
我們不想與耶穌有甚麼正面的關係.
我們知道耶穌可以幫我們.
我們知道上教會是好.
但是我們不會開放自己.
進入任何關係.
認識我們身邊的人.
或者讓身邊的人去認識我們.
這個會不會是我們今天的狀態呢.
特別我剛才講到在後COVID的香港處境裡面.
會不會我們越來越能夠離群獨居.
有幫助就會按一按那個按鈕.
不需要任何關係.
不需要進入任何與人糾結的關係裡面呢?.
信仰慢慢變成一種可以被需要的資源.
而不是再是一段要逃入.
要被觸碰,要被改變的關係.
在這裡停一停.
我知道網上有很多朋友在看我們的直播.

$^{361}$這個可能都是因為COVID的時候.
大家都習慣了.
回家開了電視就可以崇拜了.
但是這樣的信仰.
是否就好像一個婦人那樣.
按一按那個按鈕我就得到幫助了.
但是她不願意讓人見到你.
不願意進入關係.
不願意讓牧師認識你.
或者不願意讓在座的弟兄姊妹認識你.
不願意進入那個關係.
她是要求信仰可以幫到你.
這個婦人的情況會不會和我們今天大部分人是相似的呢?.
所以你今天在電視機面前.
或者在YouTube旁邊看直播的.
你就考慮一下回來我們實體的崇拜裡面.
讓我們認識一下你.
你又認識一下我們.
如果這樣看的話.
這個不是一個古代.
因為甚麼不義而得到醫治的故事.
他們讓我們看到一個很熟悉的信仰姿態.
就是人可以一邊相信耶穌很有能力.
但是一邊又可以避免讓自己真正暴露在他面前.
這種姿態並不等於他沒有信心.
但是卻是顯示他信心是可以帶著一些距離的.
但是耶穌沒有放過這個女人.
祂轉了身.
祂看見這個本來想離開的婦人.
然後跟她說.
「女兒,放心,你的信救了你」.
這個說話表面上看來.
好像只是為一個醫治的故事畫上一個句號.
但是如果你留意當中的用字的話.
你就會發覺耶穌所說的.
其實比一個醫療結果更加深層.
這裡在《和學本》就譯作.
「你的信救了你」.
但是其實這個字「救」,「拯救」.
沒有錯,希臘文就是Social的字根.

$^{401}$其實這個字在同樣地.
就是那個婦人心裡面想.
我只要摸一摸耶穌就必痊癒.
是同一個字根.
「救了你」可以是從危險當中得救.
但是另一方面就是解「痊癒」.
身體,心靈的癒合.
所以英文就可能更加令我們.
能夠掌握這個字的一些意思.
在英文裡面整句說話.
耶穌的整句說話是譯作.
「Take heart, your faith has made you well」.
所以今天我們在醫療的服務裡面.
我們不是在說「醫治」.
「Cure」或者「Healing」.
「醫病」不是純粹說「醫治」.
我們在說什麼呢.
「Well-being」.
在世衛的組織裡面.
就有句這樣寫.
我們其實是追求人的「Well-being」.
「Make well」.
「痊癒」.
「得救」.
在這裡的字.
耶穌說「你痊癒了」.
「You make well」.
「Make you well」.
「Your faith makes you well」.
你可以回到「Well-being」裡面.
甚至我太太是一個臨終治療的社工.
就是「Positive Care」.
她的病人基本上不能醫好.
最大可能令他在世有一個好的「Well-being」.
這個其實都是一個痊癒的目的.
所以這個字既可以作為「醫治」.
也可以解作「拯救」.
究竟是這個婦人她想痊癒.
純粹只是想身體恢復正常.
還是可以進入一個更深更完整的狀態呢?.

$^{441}$可能這個婦人偷偷摸摸耶穌.
祂的痊癒就是「我身體無事」.
「我可以回復正常」.
但耶穌轉過頭來說.
「你的信心令你痊癒了」.
「所有的問題都解決了」.
所謂「身心捨靈」各方面.
都可以回復一個良好的關係.
祂對自己 對別人.
身體 心靈 社會的狀態.
都可以痊癒.
所以祂不是純粹偷偷被醫治.
而是因為祂的信心可以痊癒.
所以這個字「你的信救了你」.
或者在婦人心裡面.
「我摸一摸祂的身 失戀三宗的痊癒」.
就不是純粹痊癒那麼簡單.
這裡發生的不是身體的改變.
也包括一個人在孤立 隱藏.
不想被看見的位置.
被帶回到一段關係裡面.
他原本認為的痊癒.
可能只是祈求一個身體改變的結果.
但耶穌帶給他的是甚麼.
是一種更加深的痊癒.
祂不是讓他身體恢復正常.
而是讓他重新以一個可以被稱呼.
可以被接納的身份.
站在耶穌的面前.
他不是再一個沒有名字的婦人.
他在耶穌的眼中是女兒.
祂稱呼他.
在他隱藏的位置裡面抽他出來.
和他聊天對話.
讓他看見他不需要再躲藏.
就算他身體康復.
他的心靈 他的社會關係.
都要被醫治 被改變.
所以經文留下來的.
好像不是一個關於信心大小的問題.

$^{481}$而是一個關於信心位置的問題.
我們不是沒有信心.
更加重要的是.
我們信心是停在甚麼位置呢.
只是想接觸到耶穌的能力.
卻不想進入與耶穌的關係的一個位置.
還是願意被看見.
願意被帶到一個關係的地方呢.
很多時候我們都習慣.
在需要的時候才去到神的面前.
你都可能真的經歷過祂的幫助.
不過同時我們永遠都是保持一些距離.
我們不是很想太被看見.
不是很想我們某些部份被耶穌碰觸.
我們可以很宗教.
我們有屬靈的語言.
也可以在某些情況之下與神接觸.
例如聖餐你上到台上.
是吃耶穌的身體.
喝耶穌的血.
但是否純粹是一種信仰宗教的行為.
卻是透過禮儀與神.
與耶穌進入一個更加深層的關係呢.
到最後我們要問的.
不只是我們有沒有來到耶穌的面前.
而是當耶穌轉身來看見我們的時候.
我們是否仍然停留在一個想躲藏的位置.
又或者我們是否願意讓祂與我們說話.
重新定義我們.
讓祂帶我們進入那種被醫治.
真正痊癒的關係裡面呢.
我們一起祈禱吧.
主求你去幫助我們.
我們很多時候都是以你作為伴侶.
忠誠約有什麼事來找你.
但其實你很渴望與我們聊天,對話.
你想我們了解你,你也很想了解我們.
你很想我們與你進入一個不單止是醫療的關係.
它是進入一個生命的關係.
求主你幫助我們.

$^{521}$叫我們有勇氣向你創開.
叫我們有勇氣去面對自己的種種想隱藏的狀態.
幫助我們能夠進入與人,與神的關係裡面.
而且在這個孤立,孤獨,孤單的社會裡面.
一個人與人之間分離割裂的社會裡面.
我們生命真是能夠痊癒.
求主你幫助我們.
\newpage



\section{以弗所書 4:17-32}
\label{sec:032248}
\textbf{皇帝賜的新衣}~~~~禤凱榮牧師
\newline
\newline
連結: \href{http://www.church.com.hk/acms/content.asp?site=cdc\&op=show\&type=product\&code=032248\&layout=sermon}{\texttt{032248}} ~~~~ 語音日期: 2026-06-07
\newline
\newline
\hyperref[sec:032247]{\small{< < < PREV SERMON < < <}}
~
\hyperref[sec:index_chronic]{\small{[返順時目]}}
~
\hyperref[sec:032249]{\small{> > > NEXT SERMON > > >}}
\newline
\newline
以弗所書 4:17-32
\newline
\begin{longtable}{cl}
\hline
\hline
章節 & 經文 (和合本修訂版)\\
\hline
4:17 & \begin{tabularx}{0.7\textwidth}{X} 所以我這樣說，且在主裡鄭重地說，你們行事為人，不要再像外邦人存虛妄的心而活。 \end{tabularx} \\ \\ \relax
4:18 & \begin{tabularx}{0.7\textwidth}{X} 他們心地昏昧，因自己無知，心裡剛硬而與神所賜的生命隔絕了。 \end{tabularx} \\ \\ \relax
4:19 & \begin{tabularx}{0.7\textwidth}{X} 既然他們已經麻木，就放縱情慾，貪婪地行種種污穢的事。 \end{tabularx} \\ \\ \relax
4:20 & \begin{tabularx}{0.7\textwidth}{X} 但你們從基督學的不是這樣。 \end{tabularx} \\ \\ \relax
4:21 & \begin{tabularx}{0.7\textwidth}{X} 如果你們聽過他的道，領了他的教，因為真理就在耶穌裡， \end{tabularx} \\ \\ \relax
4:22 & \begin{tabularx}{0.7\textwidth}{X} 你們要脫去從前的行為，脫去舊我；這舊我是因私慾的迷惑而漸漸敗壞的。 \end{tabularx} \\ \\ \relax
4:23 & \begin{tabularx}{0.7\textwidth}{X} 你們要把自己的心志更新， \end{tabularx} \\ \\ \relax
4:24 & \begin{tabularx}{0.7\textwidth}{X} 並且穿上新我；這新我是照著神的形像造的，有從真理來的公義和聖潔。 \end{tabularx} \\ \\ \relax
4:25 & \begin{tabularx}{0.7\textwidth}{X} 所以，你們要棄絕謊言，每個人要與鄰舍說誠實話，因為我們是互為肢體。 \end{tabularx} \\ \\ \relax
4:26 & \begin{tabularx}{0.7\textwidth}{X} 即使生氣也不要犯罪；不可含怒到日落， \end{tabularx} \\ \\ \relax
4:27 & \begin{tabularx}{0.7\textwidth}{X} 不可給魔鬼留地步。 \end{tabularx} \\ \\ \relax
4:28 & \begin{tabularx}{0.7\textwidth}{X} 偷竊的，不要再偷；總要勤勞，親手做正當的事，這樣才可以把自己有的，分給有缺乏的人。 \end{tabularx} \\ \\ \relax
4:29 & \begin{tabularx}{0.7\textwidth}{X} 一句壞話也不可出口，只要隨著需要說造就人的好話，讓聽見的人得益處。 \end{tabularx} \\ \\ \relax
4:30 & \begin{tabularx}{0.7\textwidth}{X} 不要使神的聖靈擔憂，你們原是受了他的印記，等候得救贖的日子來到。 \end{tabularx} \\ \\ \relax
4:31 & \begin{tabularx}{0.7\textwidth}{X} 一切苦毒、憤怒、惱恨、嚷鬧、毀謗，和一切的惡毒都要從你們中間除掉。 \end{tabularx} \\ \\ \relax
4:32 & \begin{tabularx}{0.7\textwidth}{X} 要仁慈相待，存憐憫的心，彼此饒恕，正如神在基督裡饒恕了你們一樣。 \end{tabularx} \\ \\
[1ex]
\hline
\hline
\end{longtable}
$^{1}$大家都聽過皇帝的新衣的故事.
一個應該沒有人沒聽過的故事.
故事大概說一個皇帝被人騙了.
以為穿著一件只有聰明人才看到的外衣.
很厲害地走在街上.
殊不知原來他是沒有穿衣服的.
本來他應該是充滿羞辱的.
但卻以為自己很厲害地走在地上.
這個故事大家聽過.
不過大家有沒有聽過一個.
剛剛和這個故事相反的故事呢.
就是說原來有一個人.
穿著很華麗很漂亮的衣服.
卻是原來以為自己沒有穿衣服.
大家有沒有聽過這樣的故事.
不會這樣吧.
不過如果大家仔細想一層.
其實福音的故事豈不是有這樣的故事.
聖經告訴我們我們所信的神.
是萬主之主萬王之王.
而這個皇帝賜給每一個信他的人.
一件很高貴華麗的外衣.
不過很多時候我們走在人生路上的時候.
都忘記了這件外衣.
以為我們沒有穿衣服.
所以我們行事為人很多時候都不能夠活出.
這個新的生命應該給到我們的尊貴.
我們沒有.
我們今天一同來通過剛才所讀的經文.
來看一看究竟我們這個萬王之王.
賜給我們這個新的衣服.
我們一同來看.
一開始的時候經文告訴我們.
保羅很嚴厲地提醒我們信徒.
他說行事為人.
我們千萬不要像外邦人一樣.
存虛妄的心.
我想大家留意一下這裡.
保羅所說的外邦人.
並不是指種族上的非猶太人.

$^{41}$為什麼我這樣說呢.
因為如果我們看第十八節的時候.
這裡說到他們心地昏迴.
與神所賜的生命隔絕了.
就是說保羅這裡所說的外邦人.
是說到一些和神分開了隔離了的人.
一些不認識神和神隔離了的人.
大家知道書卷是寫給以弗所地方的信徒.
這一群信徒大部分如果不是全部的話.
都是外邦人.
都不是猶太人.
由於他們不是猶太人.
他們以前的生活當然是按照外邦人的生活方式.
而且他們住在一個地方叫以弗所.
以弗所是什麼地方呢.
以弗所千萬不要以為是一個無名小鎮.
看回歷史的考古.
以弗所在當時來說是一個經濟上.
政治上都很蓬勃.
很發達的一個海港.
但是這個地方雖然是這麼繁榮.
但是卻是一個遠離神的城市.
為什麼這樣說呢.
因為原來這個地方.
那裡有一個廟宇.
拜祭那個神叫Atemis.
這個神是一個生育的女神.
根據考古這個神廟.
很大很大的.
可以容納二萬四千的人在裡面聚集.
二萬四千人.
大家可不可以概念是一個多大的建築物呢.
大家如果來自香港的話.
沒有人沒有聽過紅館.
如果你是歌星的話.
最終必定要在紅館開演唱會.
你知道紅館可以容納多少人嗎.
紅館我上次查過.
可以容納一萬二千人.
就是說這個神廟.

$^{81}$是有兩倍紅館的大小.
很大很大.
也由於這麼大的原因.
所以這個神廟.
是古代七大奇蹟的其中之一.
由於這個神廟所拜祭的是一個生育的女神.
所以裡面拜祭的情況.
必定有的環節當然是燒香.
有很多音樂的.
很多娛樂性的節目.
還有一樣東西.
由於是生育的神.
當時這個神廟裡面有很多的廟祭.
去跟拜祭的人進行性交的活動.
大家可以想像一下.
所有的設計的活動.
都是五官的.
看得到.
聞得到.
感覺得到.
整個的程序.
其實就是將敬拜者.
推入一個情緒很高昂.
很狂熱的一個境況.
就是當時神廟的敬拜.
以弗所除了宗教上的不道德之外.
也是當時很多犯罪的人聚集的地方.
為什麼呢.
原來當時有一個傳統.
如果犯了法的人.
只要在你抓到之前.
走到神廟的時候.
就可以逃避刑罰.
就會成為犯罪者的一個庇護的地方.
他們不需要接受這個懲罰.
大家可以想像.
每個犯法的人都馬上跑去那裡.
所以久而久之.
那個神廟就成為一個罪惡溫床的地方.
以弗所雖然是一個很繁榮.

$^{121}$很有錢的地方.
不過同一個時間.
也是一個相當敗壞的地方.
講完這個歷史背景之後.
大家可以想像一個這樣的地方.
如果你在那個城市裡長大的市民.
今天信了耶穌成為基督徒.
現在叫你不要回到以前那種.
不道德的生活的時候.
你覺得容易嗎.
不容易.
城市每一個角落的人都是這樣犯罪的.
你叫我如何來持守基督徒.
有道德的生活表現呢.
你說如何的外邦的信徒.
如何的不被引誘.
回到以前那種生活呢.
怎樣呢.
其實這個問題都是今天我們要面對的問題.
保羅就是說.
由於這一幫遠離神的人.
他們從小到大都是在一個犯罪的城市裡.
來長大成長生活的.
所以他們心地昏迴.
心裡剛硬.
我想大家想一想.
什麼叫做心地昏迴.
什麼意思.
我們經常聽到一個形容詞.
就是認知障礙.
很多人經常說認知障礙.
什麼是認知障礙.
很簡單 顧名思義.
認知障礙就是指我們精神的狀態出現了問題.
以致我們不能夠健康地.
正常地來做一個認知的功能.
分辨不到事物.
這個就是認知障礙.
來心地昏迴.
你可以視之為一種屬靈的認知障礙.

$^{161}$什麼意思.
就是我們不能夠分辨到屬靈上的真理.
在這裡我們要留意.
保羅不是在嘲笑或者侮辱一些認知有問題的人.
因為我們都知道.
高學歷的人.
他仍然都可以是屬靈上的一個核子.
看不到東西的.
問題不是在於那個人的智商有多高.
而是問題在於那個人的屬靈上的辨別能力有多多.
根本的問題不是在腦子那裡.
而是在心那裡.
當人的心變得剛硬的時候.
人的認知能力就會變得昏暗.
大家都知道什麼叫白內障吧.
可能當中也有人做過白內障手術.
什麼叫白內障.
白內障就是我們眼睛裡面有一個水晶體.
漸漸混濁.
混濁到一個地步.
光進不去.
後果怎樣.
看不到東西.
這就是白內障的病.
當人被和神分離的時候.
你可以說神的光不能夠進入人的心裡面.
以致人的理性.
人的認知能力.
屬靈的判斷力.
變得昏暗.
就好像這個鏡片一樣.
光不能夠透過.
以致人不能夠看到東西.
這就是罪的結果.
罪不單止是令人的行為污染了.
也蒙蔽了人的思想.
罪人最大的問題不是在於這個人不夠聰明.
而是在於神的光不能夠進入這群人的心裡面.
以致看不清楚屬靈的真相.
以致分辨不到在神眼中的善惡對錯.

$^{201}$不是每個人都有善惡對錯.
是的,每個人都有善惡對錯.
我現在說的不能分辨是在神眼中的善惡對錯.
所以今天這麼多的戰亂.
你有你說,我有我說.
我覺得我對,我覺得你錯.
對方又是我覺得我對,我覺得你錯.
因為每一個人今天站在這個世界來看事物.
都是按著自己的準則來定對錯.
但是心地昏昧的人.
不能夠在神眼中能夠分辨到善惡對錯.
保羅就說良心既然喪盡.
就放縱私慾,貪行種種的污穢.
既然是這樣的情況,很惡劣.
我們怎樣才能將人從困局裡面抽出來呢.
怎樣呢.
保羅的建議就是說.
我們要脫去從前行為上的舊人.
脫去行為上的舊人.
大家腦海中立刻出現了一個很鮮明的影像.
在希利文裡面脫去這個舊人.
其實真的是脫了外衣的動作.
大家想像一下,你的衣服髒了,爛了.
你當然要脫掉它.
基督徒就將這個動作.
就是說我們以前那一種的行為模式.
在神面前是不妥當的,是污穢的.
我們就脫掉它.
保羅形容這個舊我.
就是被私慾所迷惑而變壞的.
我想大家留意這句說話.
被迷惑.
你覺得你有沒有被迷惑.
你有沒有今天活在世上被騙了.
有沒有這樣的感覺.
很多時候人總是覺得.
罪惡會讓我們感到開心快樂.
卻不知道原來罪惡是慢慢地將我們捆綁起來.
我們很多時候都覺得.
只要我們滿足了物質上的需要.

$^{241}$我們情慾上的需要.
在你管理其他人的死活.
最重要是自己有自己的好處.
人生在世就可以得到最大的好處.
人不為己,天知地明.
很多時候我們都會受這種思想影響.
卻不知道.
當我們有這樣的心態的時候.
其實這個心態就破壞了我們和其他人和神的關係.
保羅告訴我們.
今天你們不同了.
你們已經不同了.
你們已經聽過神的道.
知道祂的教.
學了祂的真理.
所以你需要脫去舊人的行為模式.
就是說現在你們已經知道什麼是對什麼是錯.
第一步要做的.
保羅說悔改.
悔改.
悔改是一個決定.
讓我們可以離開舊有的生活模式.
拿一個決定來反過來.
這個就是悔改.
這個就是除去舊人的意思.
除去舊人只是舊因的第一步.
保羅在這裡告訴我們我們需要走第二步.
第二步就是穿上新人.
又是很抽象.
什麼意思?.
什麼叫穿上新人?.
很多時候我們去評估一個人.
或者評估自己是不是一個好人的時候.
通常你會聽到什麼?.
我當然是好人.
又不嫖又不賭又不喝又不吹.
還不是好人?.
很多時候我們都是去衡量.
我們有沒有做一些不好的事情.
來決定我們是不是一個好人.

$^{281}$不過基督教不單止是這樣.
因為如果基督教單是去.
基於我們有沒有做一些不好的事情.
來決定我們的屬靈質素的時候.
我們錯過了很多的東西.
停止不做一些不好的事情.
並不足夠.
想像一下.
你現在在花園裡.
做得很辛苦.
滿身都大汗.
泥土都亂了.
一身大汗.
今晚我要去附一個晚宴.
你會怎麼做的?.
渾身大汗 渾身骯髒.
當然是脫光那些骯髒的衣服.
是不是?.
脫光那些骯髒的衣服之後怎麼辦?.
再去洗個澡是不是?.
你會怎麼做的?.
把那些骯髒的衣服脫光.
洗完澡之後怎麼辦?.
就去吃飯.
希望大家不是這樣吧.
你還沒穿衣服.
聽說最近的澳洲法例.
一絲不掛走在街上.
還沒合法的.
那你怎麼辦?.
當然是穿回乾淨的衣服.
同樣.
如果我們單是停止.
做一些不應該做的事情.
這個並不足夠.
我們要停止做一些不好的事情.
也要開始做一些正確公義的事情.
而兩者之間.
不做一些事情.
做一些事情.

$^{321}$中間.
還有一樣東西.
用那個例子.
脫光那些骯髒的衣服.
穿回乾淨的衣服.
中間還要洗澡.
洗澡是什麼?.
洗去我們的罪.
怎麼樣人才能夠洗淨自己的罪呢?.
大家記不記得.
保羅是由迫伯教會.
去到成為基督徒之間.
他經歷的事情是怎麼樣的呢?.
大家記得吧.
保羅這個人.
以前是迫伯教會的.
走在大馬色路上.
耶穌基督向他顯現.
顯現的時候.
這個保羅瞎了.
在他漆黑一片.
看不到東西的時候.
不斷反省他的一切.
於是就極為後悔.
聖經更要說.
有一個仙子亞拿尼亞.
被神派去見保羅.
接著亞拿尼亞見到保羅的時候.
就跟他說.
還等什麼呢?.
未來.
求告他的名.
就是耶穌基督的名受洗.
洗去你的罪.
不過不要搞錯.
不要因為這一句話.
就想著一定要受洗.
不受洗的話.
罪不得赦免.
洗禮這個禮儀.

$^{361}$本身並不能夠除去人的罪.
不過洗禮所代表的屬靈意義.
告訴我們.
我們的懊悔.
後悔我們之前的舊惡.
所做的一切.
而將我們的信心.
投靠在我們耶穌基督.
這個救主上面.
分享他的死亡.
分享他的復活.
這個就是背後的意思.
如果我們再看的話.
使徒彼得也提醒我們.
他說這個洗禮.
本不在乎除掉肉體的污穢.
是求在神面前有無愧的良心.
說得很清楚.
無愧的良心.
就是說我們通過這個洗禮.
背後所代表的屬靈意義.
我們被耶穌基督的寶血洗得乾淨.
我們的罪得都潔淨.
就像我們剛才洗澡的例子.
洗完澡之後乾淨了.
好了.
保羅就告訴我們.
這個洗澡.
在二十三節.
就好像洗澡的果效.
就是將我們的心智改換一生.
你說為什麼會這麼厲害?.
為什麼可以心智改換一生呢?.
聖經告訴我們.
洗禮很形象化地.
通過一個人浸在水裡面.
上來就好像我們與主同死同埋葬同復活一樣.
代表著那個屬靈的意思.
也都由於這樣.
我們已經與主耶穌基督同死去.

$^{401}$所以我們就能夠成為一個新造的人.
一個新造的人.
既然是新造的人.
我們的心智就是改換一生.
什麼意思?.
剛才用我們那個白內障的例子.
白內障的手術是做什麼的?.
將你那個看不到.
透不到光的水晶體.
拿出來製造一對新的水晶體.
清晰到不得了.
光就可以再一次進入你的視網膜.
讓你可以再一次看到東西.
這個洗禮所代表的重生.
就是讓我們重生一次.
與神有這個關係.
重生神的光可以進入我們的生命裡面.
以致我們能夠辨別什麼是對的.
什麼是不對的.
也都讓我們有能力可以做一些對的事情.
這個就是穿上新人的意思.
那你說.
牧師.
我們怎樣才知道我們穿上這個新人?.
剛才你所說的脫去舊人穿上新人.
很抽象.
我怎樣知道我有沒有真的穿上新人?.
大家覺得你有沒有現在穿上新人?.
怎樣知道?.
沒錯.
也許.
也許在觀念上真的挺抽象.
什麼叫脫去舊人?.
什麼叫穿上新人?.
不過.
保羅接著再給了一些例子.
在觀念上神學上是抽象.
但是在我們實際生活的層面一點都不抽象.
為什麼這樣說?.
因為保羅在這裡再跟著說了五個.

$^{441}$我們應該有穿上新人之後的行為表現.
我們逐一來看看.
第一.
不要說謊.
要說誠實話.
我們大家知道說謊不對.
雖然我們很多時候都說.
但是我們總是知道.
當子女說謊的時候.
我們叫他們不要說.
為什麼?.
因為我們知道說謊不好.
那我問你.
既然不好.
為什麼我們還要說謊?.
為什麼人仍然會說謊?.
當然有些好處.
說謊.
很多時候讓我們得到某些好處.
讓我們可以達到某些目的.
讓我們某些慾望可以被滿足.
不過留意保羅這裡所說.
叫我們不要說謊.
他並不是說剛才我說的原因.
他說的原因是什麼?.
他說的原因就是說.
因為我們是互相為肢體.
什麼肢體?.
我們基督徒是互為肢體.
是一個身體.
這個身體.
我們之前也說過.
就是耶穌基督.
耶穌在這個《約翰福音》的十四章說.
耶穌說.
我就是道路,真理,生命.
我就是真理.
就是說.
如果耶穌是真理的時候.
就是說他自己.

$^{481}$或者歸屬他的人.
都會說真理.
不會說謊.
除非我們不屬於耶穌基督.
不屬於耶穌基督屬於誰?.
不屬於耶穌基督就是屬於撒旦.
因為我們在《約翰福音》的八章四十四字說.
撒旦魔鬼本來就是說謊.
他本來就是說謊之人的父親.
當信徒互相說謊.
互相瞞騙的時候.
我們肢體的關係.
教會的合一就會受到破壞.
相反.
如果我們互相誠實的時候.
我們就能夠互相建立互信.
讓耶穌基督的身體能夠得以成長.
教會就能夠合一.
這就是不要屬謊.
第二.
不要懷恨.
要在日落前平息懷恨.
大家都知道.
憤怒是人性的一部份.
沒有一個人沒試過發脾氣.
沒有一個人都試過.
所以這裡保羅要留意.
保羅沒有說你千萬不可以犯罪.
完全不可以犯罪.
不可以發怒.
保羅沒有說過.
為什麼.
憤怒是人性的一部份.
耶穌基督都曾經憤怒.
問題不是憤怒本身.
問題是我們如何處理憤怒.
為了不正義的事情.
為了罪惡的事情而憤怒.
這是合理的.
但當然如果你亂發怒.

$^{521}$什麼都發怒的時候.
這是一個炸彈.
在地雷隨時會爆炸.
所以這裡重要的原則.
就是不要讓憤怒滋長和蔓延.
保羅說如果到日落之後.
你仍然的憤怒還沒處理好的時候.
其實就好像在家裡.
你有沒有試過晚上在家裡.
開著門睡覺.
開著門睡覺有什麼後果.
好像這裡.
蛇蟲鼠蟻.
狂風暴雨.
賊全部都可以自由出入.
沒有掩雞籠.
自由出入.
沒有一個會思考的人.
都會想有這樣的事情發生.
但這裡.
保羅這裡所說.
如果我們將憤怒.
日落之後.
仍然沒有處理好的時候.
其實就好像一間房子.
你開著門.
撒旦可以自由出入.
所以要說得很清楚.
弟兄姊妹.
不可給魔鬼留地步.
我們要在日落之前.
或者是更加的準確.
在短期的裡面.
我們需要處理好我們的憤恨.
第三.
不要偷東西.
要勞力做工.
把有餘分給有需要的人.
大家留意這是一個很漂亮的逆轉.
一個賊本來靠偷東西來維生.

$^{561}$現在不單止不再偷.
他還要勞力工作.
以致他可以有足夠的物資.
供給有需要的人.
一個以前偷東西.
靠拿的人.
現在成為一個給東西的人.
恩典.
不單能夠抑壓到罪惡.
也能夠將正面的能量.
引向成為愛人的動力.
我想大家想像一下.
一個水桶裝水.
如果水桶底有洞的時候.
你無論倒多少水.
水桶都不會裝滿.
無論你用什麼方式.
合法的方式也好.
不合法的方式也好.
嘗試裝滿水桶的時候.
你發覺都裝不滿.
人的心就是這樣.
永遠都不能夠被滿足.
但聖經告訴我們.
當我們成為基督徒之後.
我們不再是一個有洞的水桶.
我們會成為一個水管.
一個管道.
神的水仍然會進入我們的生命.
好像進入一個水管.
水進入水管之後會怎樣.
就看流到哪裡.
基督徒的人生就是這樣.
神讓我們的生命得著他的豐盛的時候.
不是儲起他.
我們在水管的時候.
我們就要將神領受的恩典.
傳給我們周圍的人.
成為一個流通的管子.
這個就是穿上新人.

$^{601}$其中一個很明顯的表現.
再看.
不要不說莊嚴.
只說造就人的好話.
中國人也說.
一言興邦.
一言喪邦.
人的說話.
是充滿巨大的能力.
可以建立.
也可以拆毀.
寶龍不是在說我們不要說髒話.
不要說一些污言穢語.
不是在說這件事.
寶龍這裡所說的.
不要說一些說話.
是帶來別人的傷害.
有些什麼說話帶來別人的傷害.
閒言閒語.
毀謗的說話.
涼薄的評語.
羞辱別人的說話.
這些統統都是.
寶龍這裡所說的污穢的言語.
因為這些言語不單止傷害了人.
也傷害了神的心.
為什麼.
聖靈如果在我們每一個.
信徒的生裡面的時候.
當我們做一些神不喜悅的事情的時候.
聖靈在我們生命裡面.
就會為我們擔憂.
所以不要說一些傷害別人的說話.
也傷透神的心的說話.
而作為基督徒.
不單止保羅勸勉我們.
不要說一些破壞的說話.
他更進一步說.
你們要說一些造就別人的好話.
叫聽見的人得益處.

$^{641}$我想大家特別留意.
保羅這裡說的是隨時說造就別人的說話.
什麼叫隨時.
隨時的意思就是說.
我們在不同的場合.
不同的境況.
向著不同的人.
我們需要用不同的說話.
為的是一個目的.
什麼目的.
叫那個聽見的人得著益處.
舉一些例子.
對於一些灰心失望的人.
我們要鼓勵他.
對於一些陷在很大的憂傷裡面的人.
我們要安慰他.
對於一些很徬徨很困惑.
沒有方向的人.
我們要引導他.
對於一些不可一世.
驕傲到不得了的人.
我們要糾正他們的狂妄.
對於一些掙扎著很絕望的人.
我們要重新給他希望.
這裡就是保羅說造就別人的說話.
因為我們所考慮的是那個人的需要.
而不是單是我們想說什麼.
所以有人說得很好.
基督徒的嘴不是垃圾桶.
我有什麼垃圾我就倒給別人.
基督徒的嘴不是垃圾桶.
是一個工具箱.
什麼意思.
按著人的需要.
用合適的語言去造就對方的生命.
這就是基督徒應該有的一個行為.
我們再看最後一個.
不存惡毒.
要彼此饒恕.
當人與人的關係很緊張的時候.

$^{681}$受到破壞的時候.
又或者當我們對人的憤怒.
沒有被處理的時候.
這裡保羅所說的系列的後果就會出現.
有什麼後果.
苦毒,惱恨,憤怒,讓鬧,毀謗.
和一切的惡毒.
這一大堆的東西其實都是關係的毒藥.
不過這些所有的一切.
都是以前屬於我們舊有生命的一些質素.
當我們信了耶穌之後.
我們有一個新的生命之後.
這些東西就不應該再存在我們新的生命裡面.
我想大家想像一下.
你揹著一個大背包.
背包上全部都是石頭.
你問為什麼這麼多石頭.
每一次我們與人的關係出現問題的時候.
想像一下.
你的背包多放一塊石頭進去.
很刺激吧.
有什麼問題.
可能是一個受傷害的評語.
一個出賣.
對人的失望.
或者一個惡意的抨擊.
惡意的羞辱.
你說人一生走來.
特別我們這裡都過了半百的人生.
人生一路走來.
你猜你背了多少個石頭在你的背包上.
越背得的時候.
人越走的時候就越疲倦.
因為背包太多石頭.
太重.
是不是.
這些石頭可能令到你.
晚上驚醒.
睡得不好.
我們需要讓神來處理我們背包上的石頭.

$^{721}$我們的寶來在第32節告訴我們.
怎樣才可以處理這些石頭.
他說我們要因此相待.
全憐憫的心彼此饒恕.
正如神在基督裡饒恕了我們一樣.
這裡所說的是饒恕.
原來就是這麼簡單.
我們要放下這些石頭的時候.
我們必須要饒恕.
我怎樣能夠饒恕.
他都沒有跟我道歉.
我還是憐憫的心.
我怎樣能夠饒恕.
聖經告訴我們.
寶來告訴我們什麼.
我們能夠饒恕.
因為神在基督裡已經饒恕了我們.
我們如果要學耶穌基督的榜樣.
要負他的額的時候.
我們必須要學習.
既然耶穌基督已經寬恕了我們的時候.
我們如何不能夠寬恕別人呢.
所以彼此饒恕.
就是將每一塊放在我們背包上的石頭.
拿出來放在耶穌基督的十字架下面.
讓我們可以得著一個釋放.
我們的神萬王之王萬主之主.
對於我們每一個真心相信的人.
他其實已經給了一件華麗高貴尊貴的外衣.
耶穌基督說.
我們每一個信徒都是披戴基督.
就是說耶穌基督就是我們這件外衣.
所以讓我們定睛在我們的天國.
神的國度.
讓我們在地上行走的時候.
記得耶穌基督就是我們這個尊貴的外衣.
以至我們走在人生路上的時候.
我們有應得的尊貴.
因為這個尊貴來自我們神如履的身份.
讓這個身份成為我們一個榮耀神的一個動力.

$^{761}$一個原因.
請我們低頭祈禱.
因為你這段經文很清楚地告訴我們.
主啊,你就是我們的主.
我們求你來幫助,求你恩代.
讓我們整個人生都能夠活出一個見證.
見證你的尊貴.
見證你的榮耀,你的大能.
以至我們走在人生路上的時候.
不會好像以前我們犯罪,得罪你的時候所做的一切.
求你恩代,求你大力.
讓我們有足夠的能量來活出你要我們活的人生.
多謝你讓我們一切祈禱奉主名求,阿們.
\newpage



\section{歌羅西書 4:1-18}
\label{sec:032249}
\textbf{建構在屬靈支援中的基督社群}~~~~黃紹權牧師
\newline
\newline
連結: \href{http://www.church.com.hk/acms/content.asp?site=cdc\&op=show\&type=product\&code=032249\&layout=sermon}{\texttt{032249}} ~~~~ 語音日期: 2026-06-07
\newline
\newline
\hyperref[sec:032248]{\small{< < < PREV SERMON < < <}}
~
\hyperref[sec:index_chronic]{\small{[返順時目]}}
~
\hyperref[sec:032250]{\small{> > > NEXT SERMON > > >}}
\newline
\newline
歌羅西書 4:1-18
\newline
\begin{longtable}{cl}
\hline
\hline
章節 & 經文 (和合本修訂版)\\
\hline
4:1 & \begin{tabularx}{0.7\textwidth}{X} 你們作主人的，待僕人要公正，因為知道，你們也有一位主在天上。 \end{tabularx} \\ \\ \relax
4:2 & \begin{tabularx}{0.7\textwidth}{X} 你們要恆切禱告，在禱告中警醒感恩。 \end{tabularx} \\ \\ \relax
4:3 & \begin{tabularx}{0.7\textwidth}{X} 同時，也要為我們禱告，求神給我們開傳道的門，能宣講基督的奧祕， \end{tabularx} \\ \\ \relax
4:4 & \begin{tabularx}{0.7\textwidth}{X} 使我能按著所該說的話將這奧祕顯明出來，我為此而被捆鎖。 \end{tabularx} \\ \\ \relax
4:5 & \begin{tabularx}{0.7\textwidth}{X} 你們要把握時機，用智慧與外人來往。 \end{tabularx} \\ \\ \relax
4:6 & \begin{tabularx}{0.7\textwidth}{X} 你們的言談要時常帶著溫和，好像用鹽調味，讓你們知道該怎樣應對每一個人。 \end{tabularx} \\ \\ \relax
4:7 & \begin{tabularx}{0.7\textwidth}{X} 推基古是我親愛的弟兄，忠心的僕役，和我一同作主的僕人；他要把我一切的事都告訴你們。 \end{tabularx} \\ \\ \relax
4:8 & \begin{tabularx}{0.7\textwidth}{X} 我特意打發他到你們那裡去，好讓你們知道我們的情況，又讓他安慰你們的心。 \end{tabularx} \\ \\ \relax
4:9 & \begin{tabularx}{0.7\textwidth}{X} 我又打發一位親愛忠心的弟兄阿尼西謀同去；他也是你們那裡的人。他們會把這裡一切的事都告訴你們。 \end{tabularx} \\ \\ \relax
4:10 & \begin{tabularx}{0.7\textwidth}{X} 與我一同坐牢的亞里達古問候你們。巴拿巴的表弟馬可也問候你們。關於他，你們已經得到指示；他若到你們那裡，你們要接待他。 \end{tabularx} \\ \\ \relax
4:11 & \begin{tabularx}{0.7\textwidth}{X} 稱為猶士都的耶數也問候你們。奉割禮的人中，只有這三個人是為神的國與我作同工的，也是使我心裡得安慰的。 \end{tabularx} \\ \\ \relax
4:12 & \begin{tabularx}{0.7\textwidth}{X} 有一位你們那裡的人，作基督耶穌僕人的以巴弗問候你們。他禱告的時候常為你們竭力祈求，願你們能站穩而成熟，充分確信神一切的旨意。 \end{tabularx} \\ \\ \relax
4:13 & \begin{tabularx}{0.7\textwidth}{X} 他為你們、老底嘉和希拉坡里的弟兄多多勞苦，這是我可以為他作見證的。 \end{tabularx} \\ \\ \relax
4:14 & \begin{tabularx}{0.7\textwidth}{X} 親愛的醫生路加和底馬問候你們。 \end{tabularx} \\ \\ \relax
4:15 & \begin{tabularx}{0.7\textwidth}{X} 請問候老底嘉的弟兄以及寧法，和她家裡的教會。 \end{tabularx} \\ \\ \relax
4:16 & \begin{tabularx}{0.7\textwidth}{X} 你們宣讀了這書信，也要交給老底嘉的教會宣讀；你們也要宣讀從老底嘉轉來的書信。 \end{tabularx} \\ \\ \relax
4:17 & \begin{tabularx}{0.7\textwidth}{X} 你們要對亞基布說：「務要完成你從主所領受的職分。」 \end{tabularx} \\ \\ \relax
4:18 & \begin{tabularx}{0.7\textwidth}{X} 我—保羅親筆問候你們。要記念我在捆鎖中。願恩惠與你們同在！ \end{tabularx} \\ \\
[1ex]
\hline
\hline
\end{longtable}
$^{1}$陳志全議員.
\newpage



\section{啟示錄 3:1-6}
\label{sec:032250}
\textbf{撒狄:好像活著卻是死的教會}~~~~趙崇健牧師
\newline
\newline
連結: \href{http://www.church.com.hk/acms/content.asp?site=cdc\&op=show\&type=product\&code=032250\&layout=sermon}{\texttt{032250}} ~~~~ 語音日期: 2026-06-07
\newline
\newline
\hyperref[sec:032249]{\small{< < < PREV SERMON < < <}}
~
\hyperref[sec:index_chronic]{\small{[返順時目]}}
~
\hyperref[sec:032251]{\small{> > > NEXT SERMON > > >}}
\newline
\newline
啟示錄 3:1-6
\newline
\begin{longtable}{cl}
\hline
\hline
章節 & 經文 (和合本修訂版)\\
\hline
3:1 & \begin{tabularx}{0.7\textwidth}{X} 「你要寫信給撒狄教會的使者，說：『那有神的七靈和七顆星的這樣說：我知道你的行為，就是名義上你是活的，實際上你是死的。 \end{tabularx} \\ \\ \relax
3:2 & \begin{tabularx}{0.7\textwidth}{X} 你要警醒，堅固那些剩下、快要死的，因為我發現你的行為，在我神面前沒有一樣是完全的。 \end{tabularx} \\ \\ \relax
3:3 & \begin{tabularx}{0.7\textwidth}{X} 所以，要記得你所領受和聽見的；要遵守，並要悔改。你若不警醒，我必如賊一樣來到；我幾時來到你那裡，你絕不會知道。 \end{tabularx} \\ \\ \relax
3:4 & \begin{tabularx}{0.7\textwidth}{X} 然而，在撒狄你還有幾位是未曾污穢自己衣服的，他們會穿白衣與我同行，因為他們是配穿的。 \end{tabularx} \\ \\ \relax
3:5 & \begin{tabularx}{0.7\textwidth}{X} 得勝的必這樣穿白衣，我也不從生命冊上塗去他的名；我要在我父面前，和我父的眾使者面前，宣認他的名。 \end{tabularx} \\ \\ \relax
3:6 & \begin{tabularx}{0.7\textwidth}{X} 凡有耳朵的都應當聽聖靈向眾教會所說的話。』」 \end{tabularx} \\ \\
[1ex]
\hline
\hline
\end{longtable}
$^{1}$今天訊息的題目是關於撒狄的教會.
而撒狄的教會特別之處是.
他好像是活著 卻卻是死的.
所以各位弟兄姊妹 各位朋友.
首先跟大家說一聲主內平安.
在剛纔的經文裡.
就是主耶穌基督親自跟撒狄教會所說的話.
這段話在啟示錄裡是一封書信.
先講述了教會的情況.
然後講述了耶穌祂希望教會怎樣去回應.
這封簡短的書信可以是平衡.
又可以是關懷.
也可以是圈面.
甚至乎是責備.
有一次在講道之後.
有一個肢體來問我.
如何可以像你們這樣去看聖經呢?.
這個問題我也思考了一段時間.
所以今天我想在這個訊息裡.
跟大家試一下新的活動.
試一下新的活動.
我會一次過交代這篇經文的歷史背景.
請留意是一次過.
再配合剛纔所讀的經文.
你們就有兩點資料.
可以去自己想象一下.
這些歷史背景和經文之間有什麼關係呢?.
讀聖經其實都是可以運用我們的創造力和幻想的.
其次第二件事.
查考聖經有時候都需要消化不少的資料.
如果有人認為一時之間因為資料太多.
實在消化不略.
不幸地發出低沉的呼吸聲.
麻煩你如果聽到的話.
輕輕拍一拍你身邊的肢體可以嗎?.
可以嗎?.
給個回應我可以嗎?.
輕輕拍的.
不要那麼大力.
加拉太書他這麼說.

$^{41}$若有人偶然被過犯所成.
你們屬靈的人.
就當用溫柔的心將他挽回過來.
一時之間我們不是挽回.
我們輕輕拍一拍他.
這樣就挽回一個小生命.
做個練習.
不如大家現在拍一拍旁邊各位.
醒一醒先.
勢勢力拍一拍.
運用一下我們的腦袋.
大家準備好了沒有?.
新活動.
互動一點.
薩迪是一個什麼地方?.
薩迪其實是一個很堅固的城堡.
建立在一個很陡峭的山坡上面.
因為它的地勢是很險要的.
在歷史裡面是很少被攻陷.
其中一次它失守.
就是因為城堡裡面的士兵.
實在太過大安主義.
以為一定不會出事.
結果有一晚.
敵軍一個士兵.
就如同賊一樣.
潛入城堡裡面.
打開了城門.
整個城市就被攻陷了.
薩迪附近還要有個河流.
這個河流並非一條普通的河.
而是一條有金沙流過的河.
不是我們現在喫的那些金沙.
是黃金的那些金沙.
這條河就令薩迪十分富裕.
整個城市又安全.
又富裕.
完全很理想.
除了黃金之外.
薩迪的衣服也很出名.

$^{81}$因為河流帶來有條件.
薩迪就盛產漂染和紡織.
當時達官貴人穿的衣服.
甚至很純潔白色的衣服.
叫做湯加長袍.
其實都可以由薩迪這個地方製造.
再出口到其他地方.
又安全又富裕又有商業.
因為種種的原因.
薩迪就被選擇成為.
建立雅底米斯大神殿的地方.
當時很多達官貴人.
都會來到這個神殿參拜.
參拜的時候也會消費一下.
第一世紀大部分的基督教會.
其實都仍然離不開猶太教的會堂.
英文是Synagogue.
在薩迪這個地方.
可能裡面的猶太人都幾乎有.
所以當時的會堂特別大.
特別有影響力.
是其中一個最有作為的會堂.
歷史就講到這麼多.
加上剛纔我們所讀的經文.
有沒有想象得到.
兩者之間可能有什麼關係呢?.
不知道大家有沒有感受得到.
有時候主耶穌基督所講的說話.
其實好像衣服一樣度身訂造.
耶穌今天也活著.
祂今天也會和我們說話.
祂的說話也是度身訂造的.
裝作一些就叫做很扎心.
但貼地的有時候都可以聽起來有點腰心腰肺.
求神幫助我們.
讓我們首先低頭祈禱.
親愛的恩主 求主禰幫助.
因為主耶穌基督禰活著.
因為禰的話語是活的.
亦因為這樣纔有能力.

$^{121}$父神求主禰的能力.
今天在我們的生命當中彰顯出來.
要是我們需要悔改的話.
求主禰幫助.
要是我們需要有新的思想.
求主禰指引我們.
要是我們活在黑暗之中.
求主禰光照.
因為父神我們不過是人.
主禰是蒼天造地的神.
我們今天在主禰的面前承認.
哪怕只是聽主禰的信息.
我們都需要聖靈的幫助.
軟化我們的心.
將我們造成賽弗的土壤.
好叫我們如同孩子在主禰的面前.
聆聽父神禰的說話.
求主禰幫助我們.
祈禱奉耶穌基督的名叫.
阿們.
一開始問大家.
天氣現在這麼冷.
大家喜不喜歡睡覺.
喜歡的舉手.
喜歡睡覺的舉手.
很多 我都喜歡.
哇 這麼多.
大家都真的很喜歡睡覺.
在很疲累的時候.
如果能夠一覺睡天光.
充滿電.
真的很舒服.
大家知不知道.
有什麼比起一覺睡天光更舒服呢.
就是即使你.
其實真的不是很累.
但是你都一覺睡天光.
充滿電.
第二天好像明天一樣.
公眾假期.

$^{161}$睡到自然醒.
不趕的.
那就很舒服了.
很幸福.
但是有時候我們要記住.
睡得太久.
都未必一定是好事.
長期臥床.
並非一件好事.
撒底教會.
就是一間.
遂在最終的教會.
耶穌基督祂很清楚整個情況.
經文就說.
乃有神的七靈和七星的說.
七靈是指舊約撒迦利亞書.
四章七節的七盞燈的含義.
而啟示錄一章.
也提到七星是代表守護教會的七位使者.
七靈和七星在希臘文.
其實是同樣的意思.
象徵的是聖靈的完全.
和聖靈的全知.
耶穌一開始告訴教會.
我什麼都知道.
我們的事.
耶穌所有都知道.
這位全知什麼都知道的神.
就和撒底教會說.
我知道你所做的事.
按明.
不是明字.
是明星.
按明星.
在別人的眼中.
你是很有生命有活力的.
但是實際上.
你卻是死的.
你要留意一下動詞.
你擁有生命的明星.

$^{201}$但你是死的.
大家覺得你擁有什麼比較重要?.
還是你是什麼比較重要?.
這句話是一句責備.
按明星撒底教會有生命.
因為他們富有.
可以有作為.
可以做善事.
如果我想象這個教會.
如果要搞教會裝修應該是怎樣呢?.
不需要實而不華.
華而不實都可以.
反正我們給得起.
和我鋪張一點.
誇張一點都沒有問題.
不過撒底教會.
卻是名不符實.
又沒有被逼迫.
又活在安逸之中.
其實很平常的.
給我們也是.
它既有黃金.
又有衣服.
又有神殿.
我們又找到.
我享受一下.
安逸一下.
犯法嗎?.
我又和別人在一起.
沒有被人逼迫.
不是我醒目嗎?.
有什麼問題?.
但耶穌說:我見你的行為.
在我神面前.
沒有一樣是完全的.
這句話你試一下.
在我神面前拿走了.
整句話.
短句.
都是一樣的通順.

$^{241}$在我神面前.
在這裡有個作用.
好像審判官一樣的景象.
撒底教會不是毫無作為.
它是做得不夠.
可能是做了表面.
為了名聲而做.
動機或者變了質.
在《耶穌所書》.
五章十四節.
就將遂教給予為死.
它裡面這樣說.
當理者遂在的人當醒過來從死裡復活.
基督就要光照你了.
就連主耶穌基督.
都曾經對一個死了的女孩說.
不用哭.
她不是死了.
她只是睡了.
耶穌說撒底教會.
她死了.
當然一方面指她沒有了.
她應有真正的生命.
但另一方面是指她睡著了.
她發夢一樣.
她真正的問題.
是她都不知道自己有什麼問題.
甚至都不知道自己有問題.
所以耶穌就說.
你要警醒.
你們仍然有幾個人在支持.
就快受不了了.
醒一醒.
兼顧那幾個人.
還可以有排忙.
最初我叫大家搭搭大家膊頭.
通常廣東話搭膊頭就是借錢.
不是什麼意思.
叫大家搭膊頭的意思.
正正就是整段經文.

$^{281}$就是叫我們要警醒.
要醒一醒.
要看清楚自己的光景是怎麼樣.
一套電影我很喜歡看.
我看過應該多過兩次.
已經看了三次.
這套叫做Life of Pi.
大家有沒有看過.
看過了.
《少年Pi的奇幻漂流》.
我很喜歡看.
主角Pi故事很有趣.
她家是經營動物園的.
很少人家裡做動物園.
她父親因為家境情況.
整家都要搬去加拿大居住.
動物是值錢的.
是家當.
全部一起要上船.
帶著動物出發去加拿大.
誰不知在中途遇上暴風雨.
故事是說Pi一個人.
她家裡人都離難了.
她一個人在大海里漂流的種種經歷.
其中一段是說Pi漂流了幾個星期.
去到一個滿布叢林的島裡面.
她看到島的時候很開心.
島上有食物.
有淡水.
很平靜.
她可以恢復體力.
浪費了這麼多天.
對比大海的風浪.
這個島很安逸.
十分適合休息和居住.
但有一天Pi在其中一朵花裡面.
逐一打開.
在花裡面她發現一些東西.
她再放近看一看.
大家記不記得是什麼?.

$^{321}$是一個人類的牙齒.
人類的大牙在花裡面.
為什麼有個牙在花裡面呢?.
由此她推論.
其實這個島以前有人來過.
並且沒有離開.
直至這個人在島裡面.
住到自己孤獨離開世界.
原來這個島表面上是一個很安逸的島.
很適合她在這裡生存.
但其實島上的植物是靠吞食.
在島上的生物來維持生命的.
所以表面上的安逸.
其實是極度危險.
這個島在電影裡面.
甚至對人生都可以有很重要的象徵意義.
撒迪教會就是住在一個像這樣的島上.
很富有,很安逸,沒有威脅.
但靈魂被慢慢吞噬.
自己更加的懵然.
不知原來自己正在喪失自己.
在這樣的歷史背景.
加上我們剛纔所讀的經文.
套用在今天的信徒身上.
到底在責備著什麼情況?.
在上一個星期.
一位老牧師跟我說.
他說澳洲某某教會做過統計.
這個教會很大膽.
做了這樣的統計.
原來大部分的會友都沒有祈禱和靈脩.
統計就是那個教會.
我們說沒做過這個統計.
不過我相信我們就應該不是這樣的.
是不是?.
祈禱靈脩你可以說是表面行動.
有人可以說我有心.
但缺乏祈禱靈脩.
會不會多少.
不是全部多少都反映著.

$^{361}$人對神的缺乏親近和荷爾蒙的心呢?.
基督徒的意思很簡單.
就是跟隨基督的人.
這個甚少親近神.
甚少看護神的人.
會不會是一個名不符實的基督徒?.
這個問題我留給大家自己去回答.
活在安逸.
套用在今天.
不是說追求安逸的基督徒有罪.
千萬不要誤會.
真正要說的是.
當我們安逸生活成為一個很高的標準.
和一個越來越終極的目標.
無論是在追求安逸的過程.
還是得到安逸之後.
我們忘記了.
神給我們安穩的生活.
不是給我們特別享受的.
更重要的是.
給我們有更多更大的空間去服侍神.
為神付出更多.
撒地教會的問題.
就是名不符實.
正正就是活在安逸.
而更大的問題.
是他們不知道自己有問題.
經文裡沒有說他有什麼問題.
他們睡著了 在做夢.
以為自己很OK.
不要說悔改.
他們連反省都可能.
要嗎?.
我為什麼要反省?.
人很大的問題.
往往不是自己有一個很大的問題.
人很大的問題就是.
其實我們從來都不知道自己有小小問題.
所以耶穌跟他們說.
你要醒一醒.

$^{401}$看看自己的現實.
自省一下.
不要在靈裡不斷地沉睡.
因為你以為自己睡著了.
其實撒地教會已經慢慢死了.
不知道大家身邊有沒有聽到.
一些平靜安穩深沉的呼吸聲.
我在臺上是聽不到的.
不如大家幫我拍一下身邊的人.
先叫醒一下他.
因為今天資料特別多.
特別難消化.
我們小時候.
大家都試過假裝睡覺.
父母叫我們早點睡.
但我們又不硬睡.
於是就上床假裝睡.
如果要叫醒一個真的睡著的人.
小時候有很多方法.
可以撩一下他鼻子的洞.
或者在他身邊一開火門.
就大力地拍手掌.
大力地叫火燭.
但如果一個人是假裝睡著的.
你怎樣叫醒他呢?.
拍多久都不醒.
他假裝睡著.
2012年.
華文世界出版了一本書.
很可愛.
很難買.
叫做《你永遠都無法叫醒一個裝睡的人》.
封面表面上是一隻駱駝.
實際上是一個人.
細心看一看.
表面上是一個駱駝.
實際上是一個人.
很閃玻璃.
作者叫周廉.
北京大學學士碩士.

$^{441}$香港中文大學博士.
曾於牛津大學訪學.
現任中國人民大學哲學院教授.
專門研究政治道德.
其實單單看這個行頭.
都已經值得我們瞭解一下這本書.
這本書講什麼呢?.
這本書就講社會上的一些事情.
其實大家都可能知道.
未必是真.
甚至是假.
知道是有問題.
但其實不少的人都會配合.
繼續演出.
甚至將一些虛假有問題的時候.
當成是真和沒有問題.
所以裝睡的人在這本書裡的意思.
就是不願意接受真相的人.
你勸他都沒用.
怎麼勸都沒用.
然而就算在這樣的處境下.
這本書是一個正面的訊息.
無論社會是怎樣.
其實人都可以繼續做一個好人.
是想帶出這個訊息.
兄弟姐妹.
各位朋友.
來到這個階段.
你覺得這本書.
和我們剛纔所講的教會和歷史.
會不會可以聯想到一點點關聯呢?.
耶穌叫撒達教會.
祂警醒.
祂去解釋警醒是什麼.
要回想你是怎樣領受.
怎樣聽見的.
又要遵守並要悔改.
回想遵守悔改三個動詞之中.
我們今天只會講悔改.
並且我暫時現在不會解釋.

$^{481}$留信息最後和大家深入探討.
接著看經文.
經文就好像一個故事.
第二節.
耶穌就叫醒教會.
第三節.
耶穌就說.
如果教會仍然不醒.
耶穌就說.
我必臨到你那裡.
如同賊一樣.
我何時臨到.
你也缺不知道.
這句話正好挑起撒達教會.
這個城堡被攻陷的歷史.
我會潛入來.
何時來.
你不知道.
同時兄弟姐妹.
你有沒有想起聖經哪個地方.
有相似是馬太福音二十四章二十五章.
當耶穌說.
如果知道賊何時來的時候.
在比喻裡他說.
主人就必定會警醒.
十童裡的比喻大家聽過很多次.
裡面提到你們要警醒.
因為那時辰那日子.
你們不知道.
這類型的比喻是帶出主在來審判的氣氛.
耶穌承認或者不承認.
那些自稱認識耶穌的人.
就在這個場合裡面.
啟示錄帶著撒達教會和我們.
進入這個背景和氣氛.
耶穌要撒達教會知道.
不要以為這樣下去是沒有結果的.
我現在來叫醒你.
如果你不醒.
我會如賊一樣.

$^{521}$在你沒有預備的時候.
臨門審判你.
不過耶穌不僅是嚇人的.
如果撒達教會肯悔改的話.
耶穌告訴他.
悔得的要穿白衣與我同行.
明知必不從生命賊上抹去.
白衣亦是回應這個城市的警方.
白色象徵著聖潔和高貴.
當時的達官貴人穿著湯加長袍.
就是這件白衣.
在慶祝典禮的時候.
穿著這件衣服.
就是身份的象徵.
與皇帝同行巡遊.
耶穌所說的其實是警醒的應許.
和不警醒的後果.
剛纔我提到的那本中文書.
裡面記載著一個很有趣的故事.
故事不是真的.
他說在幼兒園裡.
每次到五歲的時候.
其實很多小朋友都沒有睡著.
因為不想睡 想玩 心散.
他們只是假裝睡覺.
並且他們知道.
熬過了五歲的時間.
喫完午茶 喫餅乾和水果就很開心了.
但是在這個五歲的時間裡.
這群假裝睡覺的小朋友.
因為假裝要合上眼睛睡著.
久而久之.
有很多小朋友都真的睡著了.
甚至五歲時間過了也不醒.
下午茶和水果都開始了.
還在賴床 還不醒.
很確實.
這個故事想帶出什麼呢.
這個故事是想說.
社會上很多人都知道很多事是很不公的.

$^{561}$很多假的事情.
但是我們選擇裝睡.
視而不見.
我們可能會以為.
我正在控制著.
我知道自己什麼時候會醒.
I am in control.
以為夠鍾搞定.
醒了 全身而退 改過自新.
我不需要再這樣.
但是.
就好像那群假裝睡覺的小朋友一樣.
我真的睡著了.
不醒 時間到了 你要起床.
你自己都不知道.
忘記了要起床.
這個概念由社會帶到教會.
帶到我們的生命裡面.
有時候我們知道自己有什麼虛假.
我們知道自己有什麼問題.
我們視而不見.
和自己說.
時間到我們會悔改.
誰不知久而久之.
我們慢慢忽略了.
原來自己是有虛假的.
久而久之以為是正常.
在常態習以為常.
成為我們自己的一部分.
我們連最基本需要悔改的事情.
在我們視而不見之下成為常態.
我們由一個看見問題的人.
變成一個看不見問題的人.
變成一個完全睡著了.
看不到什麼時候醒的人.
我覺得耶穌很多時候很直接.
耶穌不需要我們寫一份論文.
去簡述我們對祂的意思.
祂沒有要求我們寫很長.
耶穌很多時候都是給我們一條是非題.

$^{601}$要麼是.
要麼是非.
要麼是A.
要麼是B.
撒迪教會就是面對同樣的是非題.
要麼是睡醒 警醒.
不要再假裝睡覺.
要麼是耶穌叫醒他們.
他們就繼續假裝聽不到.
繼續睡.
人人都是罪人.
我是罪人.
你是嗎?.
大家是罪人嗎?.
大家都是罪人.
重點不是我們是否罪人.
這不是重點.
重點是我們是一個假裝睡覺的罪人.
還是一個警醒的罪人.
所以經文叫我們要警醒.
我們是罪人 耶穌知道.
我們有什麼過犯 耶穌知道.
我們不過是塵土所造.
我們有軟弱.
不要說耶穌不知道.
經文說耶穌做過人.
祂體恤我們的軟弱.
祂知道.
很清楚.
耶穌同樣把一條是非題放在我們面前.
但耶穌不是考我們.
你知不知道該怎麼回答.
應該是是還是非.
不是.
耶穌是要我們誠實.
面對我們自己的答案.
耶穌叫撒迪教會悔改.
到最後.
我就想和大家分享.
悔改是什麼呢.

$^{641}$什麼一回事.
我們常常講悔改.
在《啟示錄》七封書信裡.
悔改這個字出現了七次.
我們明白什麼是悔改嗎.
現在我們用一個思辨,問答的方式.
我們不是直接講出生活的應用.
而是嘗試想清楚悔改是什麼.
盼望可以幫助我們走出來.
我們一同去思辨,問答什麼是悔改.
思想一定會有起點.
思想悔改的起點.
你不要讀完這些字.
不要這麼心急.
起點就是.
確信神必赦免認罪的人.
很合理.
約翰一書一章九節.
他這樣說.
我們若認自己的罪.
神是信實的,是公義的.
必要赦免我們的罪.
所以我們在神的面前悔改認罪.
問心其實我們不是來神的面前.
「神啊 求你賜我凝罰吧」.
我們不是這樣.
我們在神的面前是想要赦免.
如果我們懷疑神會不會赦免.
我們或許就會害怕了凝罰.
你懷疑神是否赦免.
自然就擔心會有凝罰.
這種擔心和懼怕.
或許會驅使我們逃避認罪.
甚至想象沒有凝罰的.
不用認罪的.
沒事的.
視而不見.
由恐懼催生而成.
不過這個確信會帶出一個問題.
你們都問過.

$^{681}$我都問過.
如果我再次犯罪.
又再次認罪.
神會不會赦免我呢?.
會不會呢?.
答案是會的.
為什麼?.
因為我們要承認另外一個確信.
其實就是確信我們一定會再犯罪.
無論你是犯同一個罪.
還是犯另外一種罪.
我們都一定會犯罪.
這個不是藉口.
這是坦然承認我們的實況.
你試想一想.
我儘量說慢一點.
如果神不赦免一個再犯罪的人.
而偏偏我們又一定會再犯罪.
結果就是沒有人能夠得到赦免.
而這個結論本身和約翰一書的經文相違背.
所以神是很有智慧的.
神就用最寶貴的代價救贖我們一切的罪.
包括了過去,現在和將來.
所以我們要回到剛纔第一個確信.
就是神必赦免認罪的人.
這是神信實的應許.
不過這個確信又再產生一個問題.
既然認罪就得赦免.
是否可以刻意縱容自己犯罪?.
這個問題不少人都問過了.
而這個問題本身其實是有問題的.
是一個有矛盾的問題.
因為認罪的意思本身就是承認自己所做的事有錯.
但刻意縱容是自己願意犯罪.
以罪為樂.
根本就不覺得那件事有錯.
既然不覺得有錯繼續去做.
我們就要問:認罪從何而來?.
認罪就變得不再成立.
留意我這說的是刻意縱容而不是肉體怨惡.

$^{721}$刻意縱容是我們願意的.
犯罪是帶來快樂的.
身體怨惡是自己不願意的.
每次犯罪之後自己會有愧疚的感受.
不想繼續.
換句話說.
一個長期刻意縱容罪惡.
以罪為樂的人.
就應該反問自己.
我是否真正認罪.
重點不是神不赦免.
而是第一個確信當中的認罪.
彷彿從來沒有發生過.
如果我們不應該以罪為樂.
我們的思想很容易又跳到另一個極端.
悔改是否等於我們從此不再犯罪呢?.
很多人以為悔改.
就是靠我們的決心和努力.
逼迫自己不要再犯罪.
然而這其實是撒旦的詭計.
叫我們將目光放在自己身上.
努力點 加油點就可以了.
而不是將我們的目光放在神身上.
當我們越依靠自己.
我們跟神的距離就越來越遠.
我們越覺得神好像一個考官.
兩隻眼睛盯著我們.
叫我們醒目點.
我們就越來越容易跌倒.
所以我們要回到剛纔第一個的確信.
第二個確信就是我們一定犯罪.
人天生有罪性.
犯罪的傾向其實是根深蒂固的.
沒有人能夠靠意志力克服.
正如保羅所說.
保羅知不知道他姓什麼呢?.
他是姓保羅 姓人來的.
他也說我所願意的善反而不做.
我所不願意的惡我倒去做.
誰能救我脫離者取死的身體呢?.

$^{761}$我想我們也是差不多的.
真正的悔改是求神幫助我們.
減少犯罪的傾向次數和嚴重程度.
拉近我們跟神的距離.
如果我們不依靠神的話.
我們是不可能不犯罪的.
在今生我們都是罪人.
我們仍然會因為軟弱而跌倒.
人有意外.
兒童照顧中心叫accident.
真正的悔改其實有幾個確信.
確信是神必定赦免認罪的人.
確信我們必定仍然再犯罪.
確信落在罪中不是真正的悔改.
確信我們自己不可能依靠自己.
唯有依靠神.
悔改其實就是承認.
我們一生都是在赦免和犯罪之間的.
徘徊一生我們都是在剛纔第一個確信.
和第二個確信之間徘徊遊走.
一個承認自己每日需要神的人.
這就是悔改的生命.
悔改不悔改.
又是神給我們的另一個是非題.
現在我給大家安靜片刻.
大家可以合上雙眼.
安靜片刻.
時間是屬於你自己和神的.
我邀請你在神的面前悔改.
不是單單為一件事悔改.
有時候是為了我們不斷跌倒而悔改.
亦在神面前承認我們需要神幫助.
親愛的恩主.
我們今天以悔改的心.
想看清楚我們是何許人.
如此來到主禰的面前.
護神哪怕是我們的心思.
我們的目光.
我們的說話.
我們的行為.

$^{801}$我們的價值.
我們對將來的計劃.
我們過去所做的事.
我們現在的軟弱.
各樣事情上都虧欠神禰的榮耀.
然而父神.
禰愛世人.
禰派禰的獨生子來到世界.
來拯救我們.
禰不單止來拯救.
禰更加來世上.
去體會我們的軟弱.
去明白我們的難處.
其實神禰真的沒有必要一定要這樣做.
然而神禰因為愛的緣故.
禰愛我們.
禰如此去拯救.
父神今天我們在主禰的面前.
求聖靈禰幫助.
給我們一個新的心.
一個軟弱的心.
這個心我們承認.
其實父神我們永遠都虧欠禰.
禰永遠都是施恩的主.
我們永遠都在主禰的面前.
是虧欠的一方.
我們永遠在主禰的面前.
都有需要悔改的地方.
父神我們有時候.
即使想悔改我們都軟弱.
真的由此可見.
我們真的不能自救.
但是神禰說禰是神.
禰說禰有能力.
禰說禰可以將山從這裡搬到那裡.
禰說過只要我們有信心.
如同大菜種一樣.
就什麼都可以.
今天求神禰給我們如此的信心.
叫我們不要依靠自己.

$^{841}$父神我們的心是想過聖潔親近禰的生活.
我們的心是想安靜平安.
討神的喜悅.
只是我們活在世上.
父神我們有些事情還沒有做到.
但是求神禰幫助我們.
今天在主禰的面前我們承認.
我們有很多犯罪的地方.
教導我們每一次在主禰的面前認罪.
在我們的心中更加感受到.
主禰的愛是何等的和睦高深.
教導我們立志在主禰的面前.
成為主禰想我們所做的一個人.
上位恰用的希命.
不是單單為了被主禰去用.
而是如此.
就是豐盛的人生.
我們祈禱奉耶穌基督來的名求.
阿們.
\newpage



\section{何西阿書 5:15-6:6}
\label{sec:032251}
\textbf{真誠的侮罪}~~~~趙承恩宣教師
\newline
\newline
連結: \href{http://www.church.com.hk/acms/content.asp?site=cdc\&op=show\&type=product\&code=032251\&layout=sermon}{\texttt{032251}} ~~~~ 語音日期: 2026-06-07
\newline
\newline
\hyperref[sec:032250]{\small{< < < PREV SERMON < < <}}
~
\hyperref[sec:index_chronic]{\small{[返順時目]}}
~
\hyperref[sec:032252]{\small{> > > NEXT SERMON > > >}}
\newline
\newline
何西阿書 5:15-6:6
\newline
\begin{longtable}{cl}
\hline
\hline
章節 & 經文 (和合本修訂版)\\
\hline
5:15 & \begin{tabularx}{0.7\textwidth}{X} 我要去，我要回到原處，等他們自覺有罪，尋求我的面；急難時他們必切切尋求我。 \end{tabularx} \\ \\ \relax
6:1 & \begin{tabularx}{0.7\textwidth}{X} 來，我們歸向耶和華吧！他撕裂我們，也必醫治；打傷我們，也必包紮。 \end{tabularx} \\ \\ \relax
6:2 & \begin{tabularx}{0.7\textwidth}{X} 過兩天他必使我們甦醒，第三天他必使我們興起，我們就在他面前得以存活。 \end{tabularx} \\ \\ \relax
6:3 & \begin{tabularx}{0.7\textwidth}{X} 我們要認識，要追求認識耶和華。他如黎明必然出現，他必臨到我們像甘霖，像滋潤土地的春雨。 \end{tabularx} \\ \\ \relax
6:4 & \begin{tabularx}{0.7\textwidth}{X} 以法蓮哪，我可以向你怎樣行呢？猶大啊，我可以向你怎樣做呢？因為你們的慈愛如同早晨的雲霧，又如速散的露水。 \end{tabularx} \\ \\ \relax
6:5 & \begin{tabularx}{0.7\textwidth}{X} 因此，我藉先知砍伐他們，以我口中的話殺戮他們；對你的審判如光發出。 \end{tabularx} \\ \\ \relax
6:6 & \begin{tabularx}{0.7\textwidth}{X} 我喜愛慈愛，不喜愛祭物；喜愛人認識神，勝於燔祭。 \end{tabularx} \\ \\
[1ex]
\hline
\hline
\end{longtable}
$^{1}$相信我們都會有一些經驗.
就是我們激怒人或傷害人的經驗.
有沒有曾荊棘怒過人或傷害過人的?.
可不可以舉起雙手?.
激怒人或傷害人?.
兩隻手.
我聽到羅忠說兩隻手.
有啦,多數都有.
我們除了激怒人之外.
也曾經試過被人激怒.
我們都會有一些這樣的經驗.
每當我們被人激怒的時候.
其實我們都會很想得到別人的道歉.
或者一個相應的賠償.
但是很多時候.
我們是得不到道歉或賠償.
當我們得不到這些道歉或賠償的時候.
其實情況就會令我們更加生氣.
更加覺得你傷害了我.
你為什麼還不道歉呢?.
好像有一種二次傷害的感覺.
而另一種情況.
其實在我們的關係上也會常常發生的.
就是一些虛假的道歉.
或者是得到一個不合理的賠償.
即是明明我受了傷害.
但是我明明應該得到的賠償.
卻是跟我們預期的少得多.
我們也會不高興.
所以有時候我們受傷的時候.
我們也會想得到一個真誠的道歉.
或者應有的賠償.
什麼是虛假的道歉呢?.
結了婚的弟兄姊妹.
相信也有試過跟配偶吵架或者爭執.
也應該有試過.
一般來說也會有一些這樣的爭執.
有時候吵到一段時間的時候.
其中一方想停.
停了吵架.

$^{41}$就可能會想用這招先道歉.
但是有時候會不經意的.
在道歉上加多了一個字.
「囉」.
看到大家也挺快的.
對不起「囉」.
一個「囉」字.
只會令到吵架更加強烈.
又或者可能你忍得住不加一個「囉」字.
那就可能「對不起」.
但是對方可能仍然繼續在吵.
因為覺得「對不起」不行.
那就繼續吵.
可能忍不住就會說一句.
「我都道歉了」.
「你還想我怎樣」.
是不是有點像笑.
好像笑的多數都是弟兄.
好像是.
接著說了這句之後就繼續吵下去.
原因是什麼呢?就是因為對方發現.
你剛才那個道歉.
其實不是一個真誠的道歉.
只是想息事寧人.
為什麼我這麼清楚呢?.
「事得多」.
所以我也試過.
我忍得住不說「囉」字.
但是也試過不少次說.
「我都道歉了」.
「你還想我怎樣」這樣的情況.
當然就繼續會有一個激烈的辯論.
或者爭執.
繼續爭論下去.
直到什麼呢?.
直到一個真誠的道歉出現為止.
悔罪在我們的基督信仰裡面.
是一個很重要的元素.
但是一個真誠的道歉.
應該是怎樣的呢?.

$^{81}$除了說我們不加一個「囉」字之外.
到底上帝其實在要求我們.
要有一個怎樣的真誠的悔罪呢?.
它裡面應該有什麼元素.
有什麼內涵呢?.
今天的經課.
荷西亞書五章十五至六章六節.
是上帝對以色列人.
一段很長很長的譴責.
由四章開始至十章.
其實都是上帝對以色列人的譴責.
我們又會問一個問題.
以色列又需要什麼?.
在荷西亞書最重要的譴責.
就是譴責以色列人拜偶像.
我們知道上帝在十誡裡面.
其實是說自己是一個忌邪的上帝.
即是說是一個妒忌的上帝.
因為妒忌的上帝就是不容許敬拜他的人.
在心裡面有第二個的神明.
在舊約裡面是不容許有第二個神明.
所以在舊約裡面.
是一個一神敬拜的信仰.
所以特別在荷西亞書裡面.
荷西亞書特別會說到.
以色列人拜偶像.
就是對上帝的不忠和不正.
什麼意思呢?.
說得俗一點就是.
以色列人令到上帝大怒舞.
所以荷西亞書.
就是記錄上帝呼召荷西亞去做先知.
然後就要荷西亞這個先知.
去迎娶一個淫蕩女子歌妙為妻.
去表達以色列國對上帝的那種不正.
當時以色列國其實是流行著一種.
混合主義的信仰.
即是什麼呢?.
就是將對耶和華的信仰.
和對巴力這個神明的信仰.

$^{121}$混合在一起.
所以當時的人就用敬拜巴力的方式.
也會去敬拜耶和華上帝.
也會將耶和華當成巴力去看待.
巴力是一個什麼神?.
巴力是一個風雨之神.
在迦南信仰裡面是一個風雨之神.
而巴力是有伴侶的.
即是說巴力是一個男神.
有性別的.
而他的伴侶就是大地女神.
可能是阿斯塔祿或者是阿瑟拉.
在當時一個農耕的社會裡面.
要農作物豐收.
就會想怎樣?.
就是想風雨和大地結合.
所以在當時巴力的敬拜是怎樣呢?.
就是對巴力的敬拜是人和這些妙技.
即是妙技.
以一個性行為去引誘巴力和大地女神結合.
然後當引誘到的時候就會下雨.
下雨的時候就會令大地得著滋潤.
然後有豐收.
所以以當時以色列人來說.
他們其實是分不清到底巴力和耶和華之間.
有什麼分別.
分辨不到巴力和耶和華的分別.
是誰的責任呢?.
舉一個例子.
如果有一個弟兄姊妹.
假設而已,不是說大家.
假設有一個梁法堂的弟兄姊妹.
她回來了很多年.
也是洗禮的,每個禮拜都回來.
但是年初三.
走去車公廟.
裝著頭柱香.
就說裝著頭柱香也沒什麼.
你覺得這個是誰的責任呢?.
是誰的責任呢?.

$^{161}$有些人就覺得這個弟兄姊妹也應該有責任.
是的,這個弟兄姊妹有責任.
但是同一時間.
其實我們教睦也應該有責任.
教好這個弟兄姊妹.
所以其實當時在學生的書裡面.
四章至十章的譴責當中.
也有一段不少的篇幅.
是指責當時的祭司.
當時的祭司只是行禮如意.
然後獻祭就獻祭.
沒有好好地教導.
當時的以色列人.
怎樣好好的敬拜耶和華.
而不隨從巴黎.
所以在裡面其實也是在譴責祭司.
沒有讓以色列人專注去敬拜上帝.
說了這些背景.
我們回到今天的經課.
如果我們去看今天的經課五章十五節.
五章十五節裡面.
其實我要去的這一句話.
是上帝自己說的.
上帝就說自己要回到原處.
即是以色列地.
為的是要讓以色列人自覺有罪.
尋求上帝.
在危難的時候可以切切尋求上帝.
這個情況就像什麼呢?.
就好像一個小朋友犯錯.
大人或者家長不會離開小朋友.
而是讓小朋友慢慢知道自己錯在哪裡.
再重新接納他.
一方面要告訴他.
你錯在這裡.
你錯在哪裡.
錯了什麼.
如果你肯改過.
我們就會重新接納他.
這樣做一方面.

$^{201}$就好像給一個小朋友一個安全網.
他知道不會完全被遺棄.
另一方面也可以不斷告誡這個小朋友.
他的錯處是什麼.
讓他可以改過.
所以上帝說回到原處.
其實是一個這樣的意思.
舉多一個例子.
可能是具體一點.
或者生活化一點.
可能讓大家明白.
前幾天.
又是說我和太太和兒子.
去了壽司郎吃晚餐.
如果有去過壽司郎吃晚餐的.
弟兄姊妹都知道.
壽司郎等位都要等一段時間.
但是我們那次很早.
大概等了十五分鐘左右.
就已經可以進去了.
怎知道在那十五分鐘裡面.
日齡搶了我太太的手提電話.
要看影片.
她說她忍不住.
她要搶.
太太就拉著手提電話.
就說要先問過她.
才讓她用.
但是她就死搶.
一路要搶 一路要搶.
最後到最後是怎樣.
我這個爸爸就出手.
拿走那部電話.
兒子就糾結.
哭了.
大叫.
就是一般的小朋友都會有的反應.
大叫到一個地步就是.
有壽司郎的職員走過來我們那裡.
就問你們OK嗎.

$^{241}$有沒有事.
很感恩.
就是上帝賜了一個.
冷靜的心給我和太太.
還有一個不太在乎別人眼光的狀態.
給了我們兩個.
讓我們可以回想.
沒事 她鬧彆扭而已.
讓我們慢慢做.
那時候就慢慢指導日齡.
你可以抱著我們哭.
你可以不開心.
但是如果你要拿這個電話.
你就一定要先問盡媽媽.
這樣其實一方面是讓她知道她自己的問題.
另一方面也知道.
她可以有一些.
在面對這些事情上面.
她可以有她的情緒反應.
但是她是不會被遺棄.
而上帝回到原處.
其實是同樣起了一個作用.
就是以色列人被罰.
上帝是沒有遺棄到他們.
同樣也是為了讓他們知道.
自己的錯誤.
但是問題是.
上帝這麼做.
上帝先走了一步.
以色列人有沒有因為這樣而悔改呢?.
我們看看經文.
經文提到以色列人回覆.
其實是六章一節開始.
其實是以色列人自己的回覆.
是這麼說的.
來我們歸向耶和華吧.
他撕裂我們也必醫治我們.
打傷我們也必包紮.
過兩天他必使我們蘇醒.
第三天他必使我們興起.

$^{281}$我們就在他面前得以存活.
我們要追求認識耶和華.
他如黎明必然出現.
他必臨到我們像金林.
像滋潤土地的春月.
如果我們看這段經文的時候.
我們就會覺得.
以色列人好像願意悔改.
感覺好像願意悔改.
好像願意悔罪回到上帝身上.
甚至說他要追求認識.
再說一次要認識耶和華.
不過如果我們細心思考這段經文.
我們就會發現.
其實以色列不是真心悔罪.
我們可以說是細思極恐.
有兩點我們可以留意.
首先其實以色列人.
他是為了利益而悔罪的.
為什麼這麼說呢?.
如果我們看回經文第一節的時候.
以色列人其實他們說.
是為了自己能夠得到醫治.
國家能夠復興.
他們可以得到醫治被包紮的時候.
他們才會悔罪.
其實他們是為了什麼?.
他們是為了自己生存所以悔罪.
說得好一點.
即是說得差一點是什麼呢?.
其實這個不是悔罪.
這個是一個投機.
你要給好處我.
你會給好處我.
我就悔罪了.
這個只不過是一個利益的輸送.
而不是一個真誠的悔罪.
真誠的悔罪並不會為了利益.
而是因為你單單知道你得罪了他.
然後就去認罪.

$^{321}$但是在這個狀態裡面.
以色列人卻是找上帝著數.
所以我們斷不能夠說.
這是一個悔罪.
去到我們不妨問問自己.
當我們信主的時候.
我們都會邀請大家.
去向上帝承認自己的罪.
但是突然之間.
我們是為了什麼而認罪呢?.
如果我們說.
我們悔罪之後.
是得不到永生的.
上不到天國的.
我們悔罪之後是得不到永生.
上不到天國的.
你還會不會悔罪呢?.
可能我們會覺得.
會得不到永生.
上不到天國的嗎?.
其實這些可以說是後一步去想.
真正的悔罪.
是因為我們意識得到.
認知得到.
我們是單純地得罪了.
我們愛我們的上帝.
單純地是因為我們知道.
我們得罪了愛我們的上帝.
所以而悔罪.
上天國,得永生.
是後一步的事情.
如果我們是看著.
因為我們想永生.
我們想得到天國.
我們想進入天國.
而去悔罪的時候.
其實我們和這群意識別人.
是沒有分別的.
所以我們單純的悔罪.
或者是一個真誠的悔罪.

$^{361}$或者是我們對上帝真誠的悔罪.
是應該只是因為.
我們的罪傷害了.
一個愛我們的上帝.
單單是這樣.
所以我們願意承認.
我們得罪了祂.
而不是因為.
我因為可以上天國.
因為可以得永生.
所以去悔罪.
所以我們不妨去思索一下.
其實我們的悔罪是怎樣的呢?.
我們是不是為了.
真的一些找上帝著數.
所以去悔罪呢?.
還是是因為我們知道.
我們單單是得罪了.
愛我們的上帝而悔罪呢?.
第二點是怎樣呢?.
第二點就是.
其實意識別人是按著.
自己心目中的神來悔罪.
而不是按著上帝真實的性情來悔罪.
為什麼這樣說呢?.
在六章三節裡面.
我們請Panel幫我按一下第三節.
在六章三節裡面是這樣說的.
就是意識別人指著上帝說.
祂好像黎明的出現.
好像金林,滋潤土地的春雨.
其實這些形容.
都是和天氣,風雨和龍奸有關.
意識別人口中說.
我們要認識上帝.
但其實他們心中所想的屬性.
是在想誰呢?.
其實他們是在想巴黎.
他們將巴黎這個風雨之神.
這個自然神的屬性.

$^{401}$放在上帝的身上.
情形就好像這樣.
如果我們得罪了一個人.
我們要知道我們自己錯了什麼.
到底我們錯的時候.
我們怎樣可以有一個真誠的道歉.
你要真誠的道歉.
我們會怎樣呢?.
當然是按著那個人的性情.
去計劃怎樣道歉.
就好像我們夫婦.
我們夫婦的時候.
如果你得罪了你的配偶.
你都會按著你的配偶的性情.
和按著他平時的要求.
他的計劃怎樣道歉.
但是如果你在這裡.
你得罪了你的配偶的時候.
但你心裡是在想另一個異性.
他的性情.
然後去計劃道歉.
這個我們應該大家都覺得.
這個一定是一個錯誤的計劃.
錯誤的道歉.
甚至乎是錯常加錯.
所以其實我們去看以色列人.
他得罪上帝.
但他心裡其實是在想巴力.
然後就硬把巴力的屬性.
放在上帝的身上.
其實我們就看到.
根本上以色列人.
是沒有一個真誠的悔改.
所以上帝是到了第四節開始.
就繼續責備以色列.
至於他們的慈愛.
在現代中文譯本.
「慈愛」這個字翻譯為「對上帝的愛」.
好像祖神的雲霧.
猶如速散的流水.

$^{441}$一瞬即逝.
上帝用一些天文的現象.
或者是一些天氣的現象.
去反擊以色列.
並且上帝繼續以先知的話語.
以及自己的話語.
去懲罰他們.
去講明自己的要求.
是「喜愛慈愛,不喜愛濟民」.
「喜愛人認識上帝,勝於凡濟」.
我們的上帝.
祂不是一個自然界的神.
祂是一個超越自然的神.
祂不會用一些巫術.
或者性行為.
去讓祂祈求祂的行動.
我們今時今日不會說.
祂裝著香.
然後祂就會為我們行動.
我們這個上帝是按著自己的心意.
去與人立約.
也指引人去遵守祂誡命的上帝.
所以向上帝真誠悔罪的第二個元素.
其實是認識上帝的聖情之後.
離開祂不喜悅的事情.
兄弟姐妹.
我們所信的這個上帝.
是一個厭惡罪的上帝.
因為祂是聖潔的.
祂是神聖的.
但是如果我們看聖經裡面.
我們也會發現.
祂也是一個愛的上帝.
祂愛著我們.
祂召叫我們與祂和好.
因為祂愛我們.
祂愛著我們.
祂召叫我們與祂和好.
當我們願意真誠悔罪的時候.
我們能夠藉著基督的福音.

$^{481}$得著赦免.
一瞬間我們會有聖餐禮.
在聖餐禮裡面.
我們會有悔罪.
也會有赦罪的時候.
我很想邀請弟兄姊妹.
在認罪的時候.
我們真的思索一下.
其實我們對上帝.
是不是一個真誠的悔罪呢?.
所以我剛才也邀請了劉執幫忙.
就會在認罪文之前.
會有一個比較長一點的靜默時間.
大概三十秒左右.
在這個時間裡面.
我很想邀請弟兄姊妹.
真的發掘自己的內心.
審查自己.
我們在上帝面前.
是怎樣得罪祂.
然後當我們帶著這份真誠的心.
去到悔罪的時候.
我們才去同頌這個認罪文.
好不好?.
讓我們可以在認罪的過程裡面.
帶著真誠的心.
讓上帝幫助我們去悔改.
願上帝賜福給我們.
讓我們可以成為一個真誠悔罪.
離開上帝不喜悅的行為的一個群體.
對.
\newpage



\section{士師記 3:15}
\label{sec:032252}
\textbf{聖經人聖經事:第二位士師 以芴}~~~~黃克勤牧師博士
\newline
\newline
連結: \href{http://www.church.com.hk/acms/content.asp?site=cdc\&op=show\&type=product\&code=032252\&layout=sermon}{\texttt{032252}} ~~~~ 語音日期: 2026-06-05
\newline
\newline
\hyperref[sec:032251]{\small{< < < PREV SERMON < < <}}
~
\hyperref[sec:index_chronic]{\small{[返順時目]}}
~
\hyperref[sec:032253]{\small{> > > NEXT SERMON > > >}}
\newline
\newline
士師記 3:15
\newline
\begin{longtable}{cl}
\hline
\hline
章節 & 經文 (和合本修訂版)\\
\hline
3:15 & \begin{tabularx}{0.7\textwidth}{X} 以色列人呼求耶和華，耶和華就為他們興起一位拯救者，便雅憫人基拉的兒子以笏，他是個慣用左手的人。以色列人託他送禮物給摩押王伊磯倫。 \end{tabularx} \\ \\
[1ex]
\hline
\hline
\end{longtable}
$^{1}$各位靖兄姊妹平安.
今天又和大家一起讀聖經.
在生命樹書教網絡為大家提供聖經人和聖經時.
看看聖經的人物 看看故事.
看看背後上帝做些什麼工作.
今天的題目是第二位事司已忽.
如果你有機會看聖經事司記.
就會看到很多位事司.
第一位就是昨天說的.
阿託涅 今天已忽第二位.
在這裡是事司記第三章.
我的藍色標題字是說段數 重點.
你會看到一二三是一個惡性的循環.
因為都是一個循環 未必惡性.
不過以色列經常都是這樣犯錯.
所以是一個不好的循環.
事司記第三章第十二節.
黑色字 紅色字都是聖經裡面和合本的字.
綠色字就是我的注釋 留意.
事司記三章十二節.
以色列人又行耶和華眼中看為惡的事.
耶和華就使摩訝王伊基倫.
強盛攻擊以色列人 解釋一下.
第一點 紅色字 留意藍色螢光.
以色列人又行耶和華眼中看為惡的事.
大概有什麼惡事呢?.
惡事主要是離開敬拜其他加拿大當地人的神.
這是神認為最惡的事.
最惡的意思不是對人惡死.
而是最壞 神看為可憎惡的事.
耶和華看為可惡的事就是這件事.
所以其實一早神就叫他們進去.
要殺死那些人 要趕走那些人.
牲畜都不要留 一定要拆掉他們的壇.
不要拜 否則我就很不喜歡.
但是沒多久又是了 沒殺光.
剩下很多 有些人就想.
他可以跟我們做奴隸 他可以做什麼.
結果就為自己留下禍根.
神說了很多次.

$^{41}$那些人終於會成為你立堂上的荊棘.
會成為你的網羅 綁著你.
他會帶你離開 會成為你的痛苦.
又行耶和華眼中看為惡的事.
耶和華說 就使摩訶王 摩訶這個詞.
摩訶就是他父親的意思.
其實摩訶是什麼呢? 就是路德的女兒.
有兩個女兒跟路德一起逃跑了.
所多瑪俄摩拉的火之後 逃跑了上山.
那兩個女兒就說 都沒有人會來這裡.
爸爸有老了 我們姐妹倆在這裡 怎麼可以有後裔呢?.
沒有人會上來的 於是他們就做了一件事.
看來這件事上帝並不歡喜的.
如果沒有就沒有了 但是他們覺得要有後裔.
要有後裔在當時的人是覺得很重要的.
他們做了一件事 就是灌醉爸爸.
然後跟他同房.
於是姐妹倆人都各自生了一個兒子.
一個叫阿滿 一個叫摩訶.
摩訶的意思就是這個小孩是來自他父親的.
當時摩訶王叫伊基倫·艾倫.
艾倫是很像小牛的意思.
名字很有趣 一個大王.
不過他的名字是小牛 就是牛仔.
大家知道我講過很多次了.
其實他出生的時候 名字不是自己改的.
是爸爸媽媽改的.
可能一出生就看到這個小孩.
咦 這麼像一隻牛仔的.
所以就說好像小牛.
伊基倫這個外邦的字就是這個意思.
另外很強盛 強盛他就攻擊以色列人.
那攻擊以色列人是為什麼呢?.
其實當然在人來看.
他攻擊另一個民族就是為了想搶掠的.
最好就是他有很多糧食 很豐富.
很多衣服 很多黃金 搶他.
這次不用自己做.
現在世界上很多大國.
有力量的國家民族都會搶別人的東西.

$^{81}$不用自己做 自己越來越強盛 不勞而獲.
他攻擊以色列人.
大家看到他攻擊以色列人是一個整體的.
看為整體的.
看看這個圖 剛才這個圖出過很多次.
這次我就盡量將他放大了幾個地方.
大家看到這個圖的.
就是以色列十月之下裡面.
混雜在這裡沒有驅趕離開的加拿大.
有好幾個的 其實不止這麼多個.
這次的主角就是右下方.
就是紅色底的摩亞.
摩亞他這次去進攻以色列.
他就不單止只是自己去.
那你知道了 如果怕自己不夠力.
一個打不贏 以色列都說了.
十月之派百多二百萬人.
不說笑的 又經歷過這麼多戰事.
越來越厲害的了.
所以他不敢自己去.
那怎麼樣呢?摩亞就連累了兩族人.
一個是北方的阿滿.
阿滿就是 阿滿其實今天是安曼.
約旦首都安曼這個字來的.
那就是連累他的兄弟民族.
摩亞阿滿都是羅德的後裔.
就是兩個女兒所生的兒子.
另外連累左邊近海一點的亞馬利人.
亞馬利人大家很多時候都聽過了.
三個民族一起連累都挺厲害的.
為什麼連累亞馬利人呢?.
因為兩個東一個西.
左右夾擊都挺麻煩的.
摩亞王都有點戰略的.
伊基倫就招聚阿滿人和亞馬利人.
去攻打以色列人.
佔據棕櫚城.
於是以色列人服侍摩亞王伊基倫十八年.
這個棕樹城大概就是以色列一個很重要的心臟位置.
可能是靠近吉甲那邊的.

$^{121}$於是就被他打贏了是攻佔了.
佔據了一個重要的城市十八年.
十八年是要以色列人服侍他.
服侍是不是擔擔台水呢?.
都是的 奴役.
另外就是要進貢糧食.
服侍他,聽話.
以色列人竟然是不夠他打或者沒有人出來.
有時候未必是不夠他打.
不過沒有人領導.
沒有人領導你就搞不定.
有時一個惡人拿著一把刀.
大家都服了不要動不要動.
他都能夠鎮壓你一段時間.
十八年都挺長的.
到了第二步.
當恆爺王說眼中下為惡的事第一點.
恆爺王之後很慘.
十八年在別人手下.
以色列人就呼求耶和華.
以色列人呼求耶和華的時候.
耶和華就為他們興起一位拯救者.
這個都是第二步.
通常都有的.
你看見幾個時詩都是這樣.
第一步第二步第三步.
連記錄的字眼都一樣.
呼求耶和華的時候.
耶和華就為他們興起一位拯救者.
我這裡沒有特別寫出來.
不過看看每次拯救者一位拯救者.
於是拯救者就變成名詞.
對一位拯救者.
但其實希伯來文是沒有「者」字的.
而拯救是動詞.
這樣就有趣了.
興起一位.
一位動詞就很有趣了.
所以中文就要譯作拯救者.
或者是一位拯救的.

$^{161}$但動詞是一個活的.
正在進行的一個拯救工作或者行為.
就是這個意思.
耶和華就為他們興起一個.
拯救的動作或者行為.
這次是「辯雅敏人」.
「基拉」的意思.
「以忽」.
「Et hud」.
「Et hud」的意思就是.
「我將獻上感恩」.
後面看到「感恩」後面有點點點點.
其實「Et hud」這個字有很多解釋.
大概最主要的解釋就是.
「我將獻上感恩」.
「我將被讚美」.
大概是這個意思.
是我將會有一些好的東西.
獻上感恩以忽.
「他是左手便利的」.
即是說他是左撇子.
用左手拿兵器的.
用左手的.
左手便利.
「以忽」是這樣記載關於他.
提一下「以忽」就是.
「我將獻上感恩」.
他是用左手的.
「辯雅敏」.
上面說一下「辯雅敏人」.
「基拉」的意思.
「辯雅敏」.
「Benjamin Lee」.
「Benjamin Lee」的意思很有趣.
「辯雅敏」是用右手的意思.
「辯雅敏」的意思是用右手.
出了一個民族裡的英雄.
「以忽」是用左手的.
他的漫畫應該很有名.
而且可以出來做一些事情.

$^{201}$做領袖.
做一些外交.
以色列人就托他送禮物.
給摩押王伊基倫.
送給小牛的王.
送禮物.
由下面看到.
「以忽」應該是一個代表以色列人的.
外交的人.
做到領袖的.
也是去見過伊基倫.
而伊基倫也有相當程度信任他.
托他送禮物.
其實「以忽」被上帝感動.
他打了一把兩刃的劍.
其實當時很有趣.
多數的人是用刀的.
用刀的.
也有劍.
是用刀為主的.
「刀」就是單刃的.
「刀」是單刃的.
他就做了一把兩刃的劍.
比較方便.
而且應該這把劍.
即是說刀和劍不同.
劍是幼一點.
兩邊都是鋒利的.
這樣叫劍.
大概長一爪.
一爪即是一尺半.
由指尖到手肘.
大概這裡是一尺半左右.
即是說是不長的劍.
短短的劍.
戴在右腿的衣服裡.
找一個被套.
綁在腿上.
右邊的腳.
很有趣的.

$^{241}$我不是很明白.
就是他左手.
左手的劍就綁在右腿.
在衣服裡.
也很有趣.
為什麼不綁在左腿.
垂手一拉就可以了.
不過也對的.
因為在右腿.
右腿應該綁在右腿內側.
左手伸進去.
就比較容易拔出來.
反而如果在左腿的左側.
有時要把手拔得很高.
才可以拔出來.
右腿的內側衣服裡.
一看到衣服就拔出來了.
他把禮物獻給魔顏王伊基倫.
帶了很多以色列的禮物.
可能什麼時候都要去.
生日又去一次.
什麼時候哪個兒子生日又去一次.
就要買他拔.
可能什麼時候又要進貢.
就去見伊基倫.
伊基倫見到他來.
是不覺得奇怪的.
沒有防備的.
就覺得這個也熟悉的.
原來伊基倫極其肥胖.
可能贏得多.
吃得多.
極其.
不是普通肥胖.
極其肥胖.
二忽用辦法.
獻完禮物變成抬禮物的人.
打發走了.
禮物是要抬的.
那什麼意思.

$^{281}$就是說很多.
很開心啊.
抬完禮物就.
好了皇上.
我們送完禮物了.
希望你月立了.
笑立了.
我們走了.
他走了.
變成抬禮物的人打發走了.
其實他也走了.
因為人家下面第十九節就知道了.
自己卻從靠近吉甲作石之地回來.
在一個礦石礦場那裡.
吉甲就是他們最先.
以色列人進入迦南地的第一個佔據的地方.
在那裡安營的.
一百多萬人在那裡安營的.
就是他經常回到那裡.
就是屯兵的地方了.
自己靠近吉甲作石之地回來.
離開了.
走了.
聰明啊.
現在又回來了.
就跟阿基律說.
皇上我有一件機密事要就告你.
皇上就說.
好了.
迴避吧.
於是左右事劃的人都退出去了.
也挺好的.
如果他在身邊磨來磨去不走.
又說這裡風景多漂亮.
哎呀哎呀.
然後呢.
我有件事跟你說.
可能引起懷疑.
他就走了.
然後才再回來.

$^{321}$皇上.
我有件重要的事跟你說.
於是皇帝相信他.
好了.
迴避.
於是所有人都離開了.
伊芬就來到皇上面前.
王獨自一人.
坐在樑樓上.
這個時間.
好啊.
有事要談.
當時其實應該收完禮物之後.
天氣很熱.
其實星期四天氣很熱.
而且是南面一點.
所以是更熱的.
曠野附近.
坐在樑樓上.
即是說.
起高了.
於是地面的熱.
沒那麼熱.
這麼高.
又坐在樑樓上.
於是.
涼一點.
伊芬就上去了.
就說.
我奉上帝的命.
是要來報告你一件事的.
這個人真厲害.
留意一下.
皇上.
我是奉上帝的命.
來跟你說一件事.
咦.
奉上帝的命?.
這個王也挺厲害的.
王就從座位上站起來.

$^{361}$站起來.
即是由上帝的代言人來說話.
於是他就站起來要聽.
這個王也算是不錯.
你知道.
這麼胖要站起來也很辛苦.
伊芬便伸了手.
從右腿上拔出劍來.
刺入王的肚腹.
這個就是伊芬的計謀.
就是讓他站起來.
站起來就說.
這樣就可以.
一站起來就演出他的大肚腩.
他其實是想刺入王的肚腹.
想了刺哪裡.
他沒有想刺心臟.
為什麼不想刺心臟呢.
我也不知道.
不過是想刺入王的肚腹.
好了.
誰知道這件事.
都被人在後來說的故事津津樂道.
就是這件事很有趣.
很好笑.
連劍把都刺了進去.
就是說.
他一刺進去.
一進去.
就連那把劍握手的地方.
都刺了進去.
而且劍被肥肉夾住.
劍被肥肉夾住.
嘩.
拔不出來.
很好笑.
都出力刺進去了.
因此大家可以留意津津樂道.
這把劍.
其實很多時候.

$^{401}$譬如我們看中國的劍.
在握手前面.
和劍的舌頭的縫隙中間.
是有一塊鐵.
是阻礙自己的手不會滑下去.
而戒傷自己的手.
這把劍顯然是沒有的.
就是說.
那個握手的地方.
那個握劍的地方.
可能是一塊木.
然後有些東西綁住.
前面就是一個印.
沒有一塊東西卡在中間.
因此他一刺進去之後.
他就拿著這個柄.
然後手應該是托著劍的尾.
用手把整把劍扯進去.
這樣才做到這個動作.
這個現象.
這把劍.
厲害.
參背有從黃的土伏伐出來.
且穿通了後身.
這把劍一尺半長.
一抓.
這樣出力.
還推著尾巴.
扯進去了.
拔不出來.
因為劍被肥肉夾住.
那把握手的手柄又進去了.
他就輸了.
於是就拔不出來.
伊忽就出去走廊.
將所有門都關上.
這個關鎖一定是在外面鎖的.
可能有條鏈把鎖塞住之類的.
伊忽出來之後.
黃的僕人剛剛就來了.

$^{441}$他們就回頭.
看見留門關鎖.
我這些警衛真是好笑.
看見留門都鎖住了.
他必是在樓上大解.
以為黃帝在上面做大事.
於是就說不要搞他了.
把門都鎖上了.
於是在這裡推測.
應該是黃上每一次辦事.
大解的時候都會關上門.
還被人鎖上了.
鎖在外面鎖的.
於是鎖的人在哪裡都挺有趣.
又沒有人在.
他們就等了很久.
仍見不開樓門.
就拿藥匙開了.
不料他的主人已經死了.
倒在地上.
哇 死了 怎麼辦.
追.
他們擔言怎麼辦.
怎麼辦.
現在怎麼做好.
當然是會擔言的.
追他.
就不知道他跑到哪裡去了.
於是這個時候.
都應該可能要請示皇后.
可能要問一下大將軍.
可能要問一下他們的宰相.
就像白子善那樣.
想辦法怎麼做.
於是這個時候就遲延了.
以發已經逃走了.
經過了回到吉甲.
去了那個石礦場之後.
逃到西伊拉.
到了就上了以發.

$^{481}$沿山地吹過以色列人.
隨著他下了山地他在前頭就引路了.
於是他以發連山地.
你見到其實如果大家不介意.
看回上面的圖畫.
看回這個圖畫.
大家就會見到.
其實.
卞雅敏在這裡.
摩押王在這裡.
他走上以發連山地.
到了以色列的中心點.
在那裡大喊.
大家一起來.
跟我一起去攻打摩押.
以色列人就出來了.
對他們說.
你們隨我來.
因為耶和華已經把你們的.
仇敵摩押人交在你們手裡.
於是他們就跟著他下去.
就把守約旦河的渡口.
不容摩押一人過去.
大家留意.
其實摩押是在約旦河的東邊.
是離開了以色列地的.
離開了迦南地的.
他們把守約旦河的渡口.
即是說.
其實不是把守著約旦河渡口的西邊.
是把守著約旦河渡口的東邊.
即是說守著約旦河.
讓摩押人如果被打.
要逃走走過來迦南地這邊.
是不能過來.
不能過河.
困住在摩押地本身那裡.
怕他們走回頭走過來.
不容摩押一人過去.
擊殺了摩押人.

$^{521}$大概一萬人.
都是強壯的勇士.
沒有一個強壯的勇士能夠逃脫.
即是說.
把他的軍力全軍盡滅.
不能再打了.
以後不會再來搞了.
摩押就被以色列人制伏了.
制伏了.
即是說他沒有把他們全部除滅.
制伏了.
把他所有能夠當兵的人.
一萬人殺了.
沒有軍力了.
其他.
阿滿,阿瑪力那些都害怕了.
這個時候.
摩押就被以色列人制伏了.
第三點.
國中太平八十年.
嘩.
很長了.
這個伊弗很厲害.
能夠做到這件事之後.
由這個形勢去判斷.
應該摩押人在當時是很強的.
再加上兩個盟友.
阿滿和阿瑪力.
使得以色列被他制伏了十八年.
都很厲害.
不能少說.
但現在平安了八十年.
就是左手士司伊弗的貢獻.
都很厲害.
你見到這個士司的特點之一.
就是他沒有說他如何帶領以色列人敬拜耶和華.
沒有說他怎樣的信心.
他很有謀略.
本身應該對神有足夠的信心.
耶和華已經將摩押人敵人.

$^{561}$交在我們手裡了.
所以他和神有一個好的關係.
他有這個信心.
有一個這樣的屬靈性的感覺.
他做這件事.
這個不見經傳.
也很少人知道.
不過今天跟大家介紹一下.
第二個.
那個循環是甚麼呢?.
以色列人行耶和華.
眼中黑外王的事.
於是神就用一個外族人做皮鞭.
鞭打他.
於是摩押人就來了.
將他制伏了十八年.
很慘.
他就哀求耶和華.
這一步走了.
耶和華就給他一個拯救的動作.
就是給他一個士司.
民族英雄起來.
這次是以弗.
就贏了敵人摩押.
然後他們又跟著他回神.
於是國中太平四十年.
接著你會見到下一個士司.
又是有這條路出現的.
就是那些人又離開上帝.
一平安之後就離開上帝.
所以其實以色列人也有些討人厭.
但是神對他們是很有憐憫與忍耐.
在這裡大家也要記得.
其實神對我們這個人也很有憐憫與忍耐.
其實如果看起來.
我們不比以色列人好很多.
就會在自己一平安的時間就亂來了.
很多時候都會這樣.
不過就慘了.
於是神用其他人鞭打我們.

$^{601}$於是我又乖了一會.
有拯救者很感恩.
然後我們又平安了.
接著平安內又行沙踏錯.
就是這樣的.
求主幫助我們.
是更多信靠,倚靠他.
好 我們到這裡一起禱告.
天父上下帝幫助我們.
我們見到士司裡面一個這樣的循環.
就是他們偏離你的道.
行耶和華.
眼中學為惡的事.
最主的惡就去敬拜別神.
然後他們被愛邦成為一條皮鞭這樣鞭打.
然後他們又回來哀求你.
神又憐憫他們.
為他們拯救.
然後又平安一段時間.
接著又來過.
幫助我們不會在自己人生與靈性的路程上.
跌入這樣的惡性的循環與陷阱之中.
求主警醒我們.
幫助我們.
誰聽祈禱.
放他去耶穌基督的名下.
\newpage



\section{馬可福音 3:1-6}
\label{sec:032253}
\textbf{忘我的愛}~~~~羅泰然牧師
\newline
\newline
連結: \href{http://www.church.com.hk/acms/content.asp?site=cdc\&op=show\&type=product\&code=032253\&layout=sermon}{\texttt{032253}} ~~~~ 語音日期: 2026-06-07
\newline
\newline
\hyperref[sec:032252]{\small{< < < PREV SERMON < < <}}
~
\hyperref[sec:index_chronic]{\small{[返順時目]}}
~
\hyperref[sec:032254]{\small{> > > NEXT SERMON > > >}}
\newline
\newline
馬可福音 3:1-6
\newline
\begin{longtable}{cl}
\hline
\hline
章節 & 經文 (和合本修訂版)\\
\hline
3:1 & \begin{tabularx}{0.7\textwidth}{X} 耶穌又進了會堂，在那裡有一個人，他的一隻手萎縮了。 \end{tabularx} \\ \\ \relax
3:2 & \begin{tabularx}{0.7\textwidth}{X} 眾人為了要控告耶穌，就窺探他會不會在安息日醫治那人。 \end{tabularx} \\ \\ \relax
3:3 & \begin{tabularx}{0.7\textwidth}{X} 耶穌對那手萎縮了的人說：「起來站在當中！」 \end{tabularx} \\ \\ \relax
3:4 & \begin{tabularx}{0.7\textwidth}{X} 他又問眾人：「在安息日行善行惡，救命害命，哪樣是合法的呢？」他們都不作聲。 \end{tabularx} \\ \\ \relax
3:5 & \begin{tabularx}{0.7\textwidth}{X} 耶穌怒目環視他們，因他們的心剛硬而憂傷，就對那人說：「伸出手來！」他把手一伸，手就復原了。 \end{tabularx} \\ \\ \relax
3:6 & \begin{tabularx}{0.7\textwidth}{X} 法利賽人出去，立刻同希律一黨的人商議怎樣除掉耶穌。 \end{tabularx} \\ \\
[1ex]
\hline
\hline
\end{longtable}
$^{1}$張國柱議員.
主席.
張國柱議員.
主席.
張國柱議員.
\newpage



\section{哥林多後書 5:14-19}
\label{sec:032254}
\textbf{建立健康的使命教會 (17) 立定心志作基督的使者}~~~~林耀楠牧師
\newline
\newline
連結: \href{http://www.church.com.hk/acms/content.asp?site=cdc\&op=show\&type=product\&code=032254\&layout=sermon}{\texttt{032254}} ~~~~ 語音日期: 2026-06-07
\newline
\newline
\hyperref[sec:032253]{\small{< < < PREV SERMON < < <}}
~
\hyperref[sec:index_chronic]{\small{[返順時目]}}
~
\hyperref[sec:032255]{\small{> > > NEXT SERMON > > >}}
\newline
\newline
哥林多後書 5:14-19
\newline
\begin{longtable}{cl}
\hline
\hline
章節 & 經文 (和合本修訂版)\\
\hline
5:14 & \begin{tabularx}{0.7\textwidth}{X} 原來基督的愛激勵我們；因我們這樣斷定，一人既替眾人死，眾人就都死了； \end{tabularx} \\ \\ \relax
5:15 & \begin{tabularx}{0.7\textwidth}{X} 並且他替眾人死，是叫那些活著的人不再為自己活，乃為替他們死而復活的主活。 \end{tabularx} \\ \\ \relax
5:16 & \begin{tabularx}{0.7\textwidth}{X} 所以，從今以後，我們不再按照人的看法來認識人，縱使我們曾經按照人的看法認識基督，如今卻不再這樣認識他了。 \end{tabularx} \\ \\ \relax
5:17 & \begin{tabularx}{0.7\textwidth}{X} 所以，若有人在基督裡，他就是新造的人，舊事已過，都變成新的了。 \end{tabularx} \\ \\ \relax
5:18 & \begin{tabularx}{0.7\textwidth}{X} 一切都是出於神；他藉著基督使我們與他和好，又將勸人與他和好的使命賜給我們。 \end{tabularx} \\ \\ \relax
5:19 & \begin{tabularx}{0.7\textwidth}{X} 這就是：神在基督裡使世人與自己和好，不將他們的過犯歸到他們身上，並且將這和好的信息託付了我們。 \end{tabularx} \\ \\
[1ex]
\hline
\hline
\end{longtable}
$^{1}$節目的朋友們,早安!.
問大家一個問題,如果生命可以再次選擇的話,你會考慮選擇些什麼呢?.
大家十幾秒鐘的時間,你想想,生命裡面有些什麼的情況,處境,或者是過去的選擇.
喜歡不喜歡,有沒有後悔,或者是成功或失敗的情況,你想想.
如果生命可以再次選擇的話.
哪些東西你會來考慮選擇呢?.
有幾多個想到自己有選擇的舉手?.
有幾多個覺得自己是不需要再選擇過的?.
其實生命不能回頭再重新選擇過,想過,或者只能夠向前.
如果你選擇了向左的,向左邊你可以再左一點,但你走一走左邊可以轉右邊一點,或者繼續向前走.
或者你選擇了向右邊走的時候,你走一走左邊可以轉左邊一點,或者繼續越來越右的,或者向前走.
總是向前不能回頭.
如果你生命裡面覺得自己是不需要選擇過的話,有兩個可能性,或者最少有兩個可能性.
一個就是你很滿意你現在的狀況,或者很happy,不需要再選擇過,安於自己現在的情況.
一個可能性就是,雖然現在的處境一般般,甚至很失望,很不理想,但你仍然堅持留在這個處境裡面.
因為你有信心,你知道自己沒有選錯.
即使在苦的裡面,在被埋怨的裡面,在不明朗的裡面,你仍然堅持著立定心智來做你現在手上所做的事情.
或者處在你現在的處境裡面.
立定心智做什麼事情呢?我們每一個人都會不同,但今天我們一起去立定心智做一件事.
今日的經文裡面只是讀了幾節聖經,其實在讀那幾節聖經之前,我們需要回頭多講一點.
因為今天讀的聖經裡面有四節可以,所以,所以什麼,所以什麼,所以什麼.
所以我們需要回頭多一點,看看那個因為,因為什麼的緣故,因為什麼的理由.
我們所要做的,接下來想要決定的那幾件事呢?.
我們將今日那幾節的經文,加上之前那幾節,從第五章第一節開始看.
看到今日第二十節的聖經,分為六段的段落來看這幾節的經文.
第一段是永恆的歸宿,有沒有人喜歡去露營的?.
夏天即將到來,去露營,去宿營,很喜歡,這間屋子也不錯,跟真正的屋子的樣子很相似.
露營的時候我們也會想一下,弄得舒舒服服,弄得好像一間真正的屋子一樣.
但你知道露營最終都要回家,露營只是短短幾天的時間,最後都要回到自己的家.
保羅跟我們說,因為我們知道在地上的帳篷如果拆毀了,神所造的居所將來會為我們預備.
不是人手所造的,而是在天上永存的.
這一句的說法,用帳篷來講這個比喻,對當代的以色列人來說就容易明白.
你可以想像以色列人在阿巴拉罕時代,他們已經離開自己的地方,在一個流離,拆帳篷,到處去居住不同地方的生活一段時間.
直到最後進入迦南地,約書亞帶著以色列人進入迦南地之後,仍然都沒有在一個永恆乾息的地方,只是在地上暫住的地方.
保羅用這個比喻,其實不是在講他們住的地方,而是在講他們的肉體,地上的肉體,這個帳篷,有一天都會拆毀的,有一天都會毀壞的.
但是我們天上有一個居所,有一間屋,不是人手所造的,而是天上永存的屋.
保羅說我們知道,因為我們知道我們這件事,我們期待著將來永恆裡面我們有一個新的身體,有一個永恆的居所在這裡,等候著我們.
這個知道對我們有什麼影響呢?就會影響我們今生的選擇,今生的取捨,今生的價值觀,今生的先後次序的排位不同.
永恆裡面的歸宿,影響我們今生很多的決定和選擇.
保羅也知道我們在這個帳篷裡面,其實過著一些勞苦嘆息的日子,每個人都有不同的時間都會生病,我們都會老.

$^{41}$如果我們知道天上有一個永恆的居所等我們,我們知道將來有一個復活了的身體等著我們的時候,我們就會對生命不再一樣.
我們不害怕生老病死的問題,我們就用不同的心情去面對今天的處境,保羅說我們並不是願意很快的去除了這件事.
我們不是很想早點死,我們只是想明白,希望可以快快的穿上將來在天上的居所那件外衣.
接著保羅也用外衣來做同樣的比喻,我們不是想早點離開的地方,我們只是期待著有一天可以披上一件新的外衣,新的身體,新的居所.
裡面有一個圖畫,圖畫裡面令我想起幾十年前認識的一個姐妹,她中年的時候有一個兒子和一個女兒,她的女兒只是不夠十歲,當時這個姐妹患上末期的骨癌.
醫生宣判了她只有幾個月的生命,雖然她的肉體正在受苦,可能她甚至想早點回到天家,不用再受地上的苦.
早點回到天上,在神的懷抱裡面,但是她女兒只有不到十歲的時候,她就和神祈求,求神給她多點時間,仍然存留在地上,可以看著她女兒的長大.
結果神有恩典給了她,她真的看到她女兒十四五歲的時候,她才離開這個家庭,回到神的家裡面.
神是那位為我們安排一切的神,我們幾時回去,幾時留下,幾時神會給多點時間,幾時神接我們回去,是神為我們安排,神不單止為我們安排時間.
神為我們安排在永恆裡面有個居所,還給我們聖靈作為一個憑據,神在耶穌基督裡面給我們新的生命,新的盼望,因為主耶穌基督的犧牲,好像我們聖餐幫助我們再次紀念一樣.
耶穌基督的犧牲,祂的死,祂的復活,祂給我們永恆裡面有份的盼望,基督,基督是聖靈在我們心裡面的一個憑據,讓我們知道,確實的知道將來的事是怎樣來發生.
知道一切是怎樣呢?得到主的喜悅,第二個命題,所以無論是住在室內或室外都好,我們都要立定志向,要得到主的喜悅.
原來我們知道剛才那種情況,知道我們永恆裡面有些什麼在等我們的時候,我們的自然反應,就應該過一些日子去得到主的喜悅,可能我們都有一些心裡面不一定十分的願意.
我們已經有了上天堂的門票,有了上天堂的確據,有了得到永生的承諾,有了一份盼望可以有新的生命,一切都得到了,為什麼我們不可以在今生裡面仍然會著我們自己想取捨所選擇的日子呢?.
為什麼我們要過一些得主的喜悅的生活呢?保羅都怕我們不明白,加多一些解釋在裡面,因為我們每一個人都必須要在基督的審判台前受審.
我們都知道人不能逃避將會受到神的審判這件事,聖經裡面說得很清楚,白色的大寶座是人要面對的審判.
原來基督徒神位我們為了預備了另一個審判,就是基督審判台前的審判,專為一些信了耶穌,得到救恩的人來面對的審判.
基督台前的審判和大寶座的審判是不同的,基督台前的審判其實就好像一個運動員在運動場上,向著最後的終點來努力.
最後在終點那裡,那裡就叫審判台,就是基督的審判台,原來基督在審判台那裡等著看我們每一個人完成我們在地上走的路.
不是和旁邊的人要比較,不是要比旁邊的人快,是要完成我們自己要完成的那條跑道,要走到最後審判台前,讓我們基督,我們主耶穌來作最後的審判.
難怪布羅華當走的路他已經跑過了,當守的路他已經守過了,布羅華就知道他需要跑這條路,知道將來永恆裡面有居所,有新的肉體.
直到我們將來真的要去到天國的寺院那一刻,中間原來他要經過一個審判的階段.
審判台你不用聽,好像很害怕,好像法官在那裡有什麼審判,是不是主耶穌很想罰我們,其實不是,我看審判台其實是神,耶穌基督在台那裡等著看我們過終點,是因為祂要給我們有獎賞.
是因為祂想見到我們努力完成我們今生所要完成的路,祂不希望見到我們拿到特天堂的門票之後站在跑道上悠悠閒悠悠,孖孖呼呼,最後都不知道為什麼沒有完成要完成的跑道.
沒有努力的去到基督審判台前,或善或惡,基督會來讓我們有各個審判和報應.
我們可能奇怪,為什麼保羅可以守得住呢,為什麼他可以跑完他的道呢,是因為他的生命自從預算了耶穌基督之後,他可以真正的為主而活.
他說原來基督的愛激勵了我們,原來基督的愛激勵了保羅的自己,基督的愛,基督的犧牲,基督的死,基督的復活,激勵了保羅願意為主而活.
他很確定的跟我們說,基督一人為眾人死,眾人就都死了,是向罪死,我們不再需要在罪的就坐,不再需要在罪的捆綁那裡,對眾人,我們因為耶穌基督的死的緣故,他犧牲的緣故,我們可以對罪說,不再需要受你的捆綁.
一人為眾人死,眾人就為基督而活,一人為基督,為眾人而死,眾人就為基督而活.
在小說裡面,武功高強的英雄人物拯救了一個美少女,通常美少女會怎樣回報呢?英雄救美在小說裡面是怎樣的呢?.
無以為報,以身相許,這是小說裡面所說的故事,在現實生活裡面我們是怎樣的呢?現實生活裡面我們得到主耶穌的拯救,主耶穌一人犧牲了自己,拯救了我們每一個,我們無以為報,那又怎樣呢?.
沒有人說的,以生命相交托,我們經歷了神的恩典,主的拯救,拯救了我們從罪裡面拯救了出來,我們無以為報,只能夠以我們生命來交托給我們的主,我們願意為主而活.
有沒有人還記得這件事講什麼呢?去年十一月的時候,杜洛斯四十週年交唱會永恆音樂在香港特地過來跟我們做一個分享報道的晚會,唱出心中那團火.
杜洛斯他們四十週年,Pastor Karen就唱足差不多四十週年,他的生命就是為主而活,唱出心中那團火.
他為主而活,為主而唱,那些隊友為主而彈,為主而講,永恆音樂裡面的梓盈和Gary,同樣的三十週年他們為主而活了一段很長的時間,放下自己的事業,為主而作功,這個就是為主而活的生命.
可能說我們有沒有這樣的恩賜,我不知道你有沒有,我就沒有了,我沒有唱的恩賜,也沒有彈的恩賜,那可以有什麼事奉呢?我們不可以像他們在台上那樣事奉,可能我們為了逃避事奉,逃避回應主的恩典的時候,給自己一個藉口,為了沒有那些恩賜,我們就不用了.
然後保羅告訴我們,其實我們和那些在台上事奉的,和其他的報道家,和其他的宣教士,和所有在主裡面事奉的人,都是一樣,都是一個新造的人.
保羅說,所以從今以後我們看人,我們看我們耶穌基督,和以前看的不再是一樣,和以前所看的不再是一樣,是怎麼不一樣呢?我們以前覺得怎樣呢?厲害的人就在台上事奉,厲害的人就在人面前影響世界.
而現在世界裡面影響人的,就是那些我們以為厲害的人,神裡面我們每一個人都是一樣,都是新造的人.
其實Karen的報道團的名字叫做Doulos,Doulos的意思是奴僕,不單只是僕人,是奴隸的僕人,是奴僕的意思,他們選了這個名字是代表他們的心聲,他們願意事奉他們的主,為主而活,以奴僕的身份來事奉.

$^{81}$我們和他們其實都是一樣,我們已經是新造人類了,我們可以為主而活,成為主的僕人,我們是主的跟從者,我們是神的兒女,是天國的公民,我們和那些其他努力事奉,為主而活的人,都是一樣,都是一樣.
又有人在基督裡,他就是新造的人,舊是已過,都變成新的了,不要再覺得自己是一個沒用的人,主可以大大讓我們,不要覺得我們在世上是一個很成功的人,主可以讓我們成為祂的奴僕.
第五點,新的使命,我們是新造的人,神就給我們一個新的使命,不單止神令到他自己和我們可以和好,要把這個和好的道理,這個信息教育給我們.
創世記裡面,神做了亞當,夏娃,給了他們一個大使命,是什麼?生養眾多遍滿全地.
今天我們這班新人類,神同樣給我們一個新的命令,我們新造的人類,神同樣給我們一個新造人類的命令,就是生養眾多遍布全地.
生養眾多,像這些嬰兒一樣,是屬靈的嬰兒,昔日神讓亞當,夏娃肉身上生養眾多遍布全地,今天我們,今天給神特別呼籲我們一班蒙恩的人,基督耶穌裡面一班新造的人,一個新的人類,就是生養眾多,屬靈裡面多一點的兒子,多一點的女兒.
將主耶穌和好的道理,將神和好的福音傳給更多的人,在靈裡面建立多一點屬靈的兒女.
雖然我們可能覺得我們沒有其他人事奉得那麼好,我們不可以為主而唱,為主而彈,為主而講,但是有一樣東西你可以做的,就是為主去探訪.
知不知道外教會有一個小城宣教隊,已經有兩年的時間了,兩年的時間裡面,他們是每年的夏天,每一個月都出隊去安省附近的小鎮,探訪那些餐廳,店鋪和老闆聊天.
現在我想聽一下他們幾位小城宣教隊友的心聲.
兩年裡面,我會覺得一方面覺得有機會將福音帶給人,但是首先始終都是要認識那個人.
到目前為止,兩年我覺得都做到的,到目前為止真的有些餐館的管主都當我們是朋友的,有時有時那些心事都跟我們講,所以是要很久去表白,可能每個人有些會快些有些會慢些,所以另外在同宮一起去的時候,也有很多的分享在路上.
我自己覺得是一個福音機會,如果突然之間去探訪,有時我們在車上都會聊很多東西的,所以都是另一個學習或者一個成長機會.
感恩由不認識到大家成為朋友,是小城宣教的生命見證,不覺得花了時間沒果效,因為宋土宰眾算不得什麼,唯有神叫他生長開花結果,榮耀頌讚歸於神.
很感恩在小城宣教參加了兩年,得著是跟別人傾談當中,學了很多東西,想到別人怎樣想我們去探訪他們的情況.
第二方面,其實不講福音也沒問題,最重要是跟他打好關係,讓他知道我們基督教來探訪他,主要是愛關心他,是最重要.
剛才那個姐妹說,就算不可以馬上傳福音給他,也不要緊,起碼可以保持到,建立到一份友情,一份關係.
上星期馮傳道跟大家講了3D,去跟那些人聯絡,去結交,去認識,有機會的將我們自己的生命故事跟他們分享.
我記得我們跟那些短宣隊去到中南美洲短宣的時候,很多時候有短宣隊出隊的隊員都很謙卑的.
他說我出去只是想幫忙煮飯,有些人說我出去只是想幫忙做司機,有些人說我出去只是想幫忙按PowerPoint.
我還記得我第一次出隊的時候,我就說我出去只是想幫忙按PowerPoint,覺得自己只是做一些手作.
但是去到工場有機會的時候,眾人都有機會跟弟兄姊妹一起,跟當地的那些人,那些不認識的陌生人.
講一圈幾句,進入認識一下他們,了解他們的生活,問一下他們的情況,都是短宣隊可以做的.
回來問一下我們自己,我們沒有那個恩賜可以站台,為主發聲,沒有那個機會去到外面做宣教,我們自己身邊有沒有機會做主的使者,為主做見證呢.
最後一個所以,所以我們做基督的特使,不是那些其他有恩賜的人,弟兄姊妹才做的,是我們這些有恩賜的人,同樣都可以承擔的一個使命.
還記不記得這一支旗幟,上一次沒有機會跟大家講我旗幟的意思.
紅色十字架代表什麼,大家都很心裡有數,代表耶穌基督為我們所釘的犧牲的十字架,祂留的寶血見證我們的罪,祂的犧牲,祂的死,祂的復活.
藍色代表天國的福音,祂帶來彰顯了神的信實,然後白色代表基督帶給這個世界和平,基督帶給這個世界潔淨,聖潔,這個就是基督徒的使命,就是主耶穌交給我們的使命.
其實保羅今天跟我們所講的,立定心智,無論道路有多艱難,無論結果是不是我們所期盼的那麼成功,我們仍然堅持在這個方法,在選擇裡面向前走.
因為主耶穌吩咐我們,祂吩咐的使徒保羅,就像保羅所講的,吩咐我們要成為主的使者,神交託給我們這個福音的使命,是要我們完成.
我們知道剛才所講的,我們永恆裡面,我們已經知道一定有歸宿等候我們,我們有新的身體,天上面有一個居所,我們就過著一些為主而活的生命,蒙主,得主的喜悅.
我們知道我們最後,最終每一個都要面對基督抬前的審判,我們希望這個審判是一個獎賞的審判,而不是一個懲罰的審判.
各位蒙族的人,我們都接受這個新的使命,神昔給亞當夏娃一個使命,今天給我們一個地上能夠傳講福音,建立更加多的屬靈兒女的一個生命,就是這個最後的基督特使的使者的身份.
要有一個使命,要有一個選擇,要有一個方向,要堅持,立定心智,努力向前,不論結果如何,我們知道只有一個,就是蒙主預立的生命.
一切向上祈禱,主人所有當你方言看我們生活的環境,我們生活周圍的人,我們見到的是甚麼呢?.
我們見不見到一些流離失所的人,見不見到一些還沒有認識你的人,見不見到一些屬靈生命裡面在死亡的人,見不見到你所給他們的使命.
我們都願不願意拿起這個旗幟,基督徒的旗幟,為主發聲,為主在我們走出去,為主我們願意開口,我們願意去探訪,我們願意去關心身邊的人.
這是主你所教的,就好像昔日你對在世上各個有錢的稅吏一樣,今日給我們這個愛心,同樣地接觸這個世界,接觸身邊你所帶來給我們的人.
求神的令來加力給我們,求你加力給我們每一個,奉主耶穌基督得勝名求,Amen.
\newpage



\section{以弗所書 1:3-14}
\label{sec:032255}
\textbf{無限靈福}~~~~謝友德牧師
\newline
\newline
連結: \href{http://www.church.com.hk/acms/content.asp?site=cdc\&op=show\&type=product\&code=032255\&layout=sermon}{\texttt{032255}} ~~~~ 語音日期: 2026-06-07
\newline
\newline
\hyperref[sec:032254]{\small{< < < PREV SERMON < < <}}
~
\hyperref[sec:index_chronic]{\small{[返順時目]}}
~
\hyperref[sec:032256]{\small{> > > NEXT SERMON > > >}}
\newline
\newline
以弗所書 1:3-14
\newline
\begin{longtable}{cl}
\hline
\hline
章節 & 經文 (和合本修訂版)\\
\hline
1:3 & \begin{tabularx}{0.7\textwidth}{X} 願頌讚歸給我們主耶穌基督的父神。他在基督裡曾把天上各樣屬靈的福氣賜給我們。 \end{tabularx} \\ \\ \relax
1:4 & \begin{tabularx}{0.7\textwidth}{X} 因為他從創世以前，在基督裡揀選了我們，使我們在他面前成為聖潔，沒有瑕疵，滿有愛心。 \end{tabularx} \\ \\ \relax
1:5 & \begin{tabularx}{0.7\textwidth}{X} 他按著自己旨意所喜悅的，預定我們藉著耶穌基督得兒子的名分， \end{tabularx} \\ \\ \relax
1:6 & \begin{tabularx}{0.7\textwidth}{X} 使他榮耀的恩典得到稱讚；這恩典是他在愛子裡白白賜給我們的。 \end{tabularx} \\ \\ \relax
1:7 & \begin{tabularx}{0.7\textwidth}{X} 我們藉著這愛子的血得蒙救贖，過犯得以赦免，這是照他豐富的恩典， \end{tabularx} \\ \\ \relax
1:8 & \begin{tabularx}{0.7\textwidth}{X} 充充足足地賞給我們的。他以諸般的智慧聰明， \end{tabularx} \\ \\ \relax
1:9 & \begin{tabularx}{0.7\textwidth}{X} 照自己在基督裡所立定的美意，使我們知道他旨意的奧祕， \end{tabularx} \\ \\ \relax
1:10 & \begin{tabularx}{0.7\textwidth}{X} 要照著所安排的，在時機成熟的時候，使天上、地上、一切所有的，都在基督裡面同歸於一。 \end{tabularx} \\ \\ \relax
1:11 & \begin{tabularx}{0.7\textwidth}{X} 我們也在他裡面得了基業；這原是那位隨己意行萬事的神照著自己的旨意所預定的， \end{tabularx} \\ \\ \relax
1:12 & \begin{tabularx}{0.7\textwidth}{X} 為要使我們，這些首先把希望寄託在基督裡的人，頌讚他的榮耀。 \end{tabularx} \\ \\ \relax
1:13 & \begin{tabularx}{0.7\textwidth}{X} 在基督裡你們聽見真理的道，就是那使你們得救的福音，你們也信了他，就受了所應許的聖靈為印記。 \end{tabularx} \\ \\ \relax
1:14 & \begin{tabularx}{0.7\textwidth}{X} 這聖靈是我們得基業的憑據，直等到神的子民得救贖，使他的榮耀得到稱讚。 \end{tabularx} \\ \\
[1ex]
\hline
\hline
\end{longtable}
$^{1}$提卿姐妹早晨.
我們剛才的經文主席都提過.
我們六月開始我們是在二佛所書裡面.
大家一齊學習.
按著我們教會的主題.
再次提醒大家.
在基督裡同心前行.
我們今天就開始二佛所書四次的學習.
我們再次一起祈禱.
多謝我們全能的天父.
我們在這裡學習你的心意.
你的話語.
學習你給我們的使命.
我們在世上怎樣去履行.
怎樣去遵從.
亦使我們的生命更加豐盛.
按著你的旨意.
我們參與你的救恩計劃.
亦都能夠享有你的永恆的天國的各種屬靈的福氣.
求聖靈繼續今日教導我們.
禱告奉耶穌的名 阿們.
剛才開始崇拜之前.
各團導跟我說.
牧土的內容就是今日講道經文的內容.
所以祈禱的時候已經講完了.
可以繼續講也可以不用講.
大家可以早些結束講道.
有時間唱回應詩歌.
多唱幾次聖餐的時候.
可以有時間去默想自己禱告.
事實上因為這段經文.
真是講道神給我們的屬靈福氣.
所以就用靈福來講.
教會叫做靈福.
我們不太清楚當年誰用靈福來做教會的名字.
真是有很多屬靈福氣.
尤其是無限多的福氣.
當然我們又祈求教會繼續有無限多的屬靈福氣.
在這些屬靈的福慶儀裡.
我們要逐一數出來是很重要的.

$^{41}$因為我們唱剛才那些詩歌.
其實是數了很多屬靈的福氣出來.
但因為是要唱歌.
那首歌就要一路唱一路唱.
沒理由中間突然停了.
所以我們唱這些詩歌的時候.
就沒有機會可以慢慢的仔細的.
想想或者領會一下那些屬靈的福氣究竟是怎樣的.
如果大家都想去做這件事.
就要自己回家繼續唱.
唱多幾次.
有時又有些歌不是那麼熟.
回家又不記得怎樣唱.
所以在我以前的教會.
有一位婆婆她是六十多歲才信主.
她就經常問我們要拿錄音.
等她回家就可以慢慢拿著歌詞自己去唱.
這個都是幾好的方法.
還有另外一種方法是.
怎樣領會神的恩典的偉大呢.
很多很多年前了.
我都不記得準確是多少年前.
但是很多年前了.
當時我有一個兒生.
他應該大約四五歲左右.
我姐姐他媽媽那次就說.
弟弟你是牧師來的.
你不如可不可以和我兒子講一下.
基督教的信仰.
他四五歲應該怎樣講好呢.
我大概就講了神是很大的.
你說他偉大他又不明白.
你說他很大他就容易明白.
那次剛好經過前面有一棟高樓.
三四十層的高樓.
他突然間問我.
舅父 神有沒有這棟高樓那麼大.
很難形容神究竟是否像這棟高樓那麼大.
還是更大呢.
我都不太記得當時怎樣回答.

$^{81}$總之意思是神大很多.
比這棟高樓大很多.
我經常想起這件事.
我們越能夠體會到.
領會到神的大能.
神是有多大.
很多時候就會直接影響到.
究竟我們信這位神.
信得多多 信得多實.
信得多真 信得多深.
因為我們越知道 越領會到.
這位神的大的時候.
我們就越有一種感恩的心.
當我們有感恩的心.
我們就更加有能力去跟從他.
更加有信心去服從他所說的話.
真的比我們有那種能力 有那種信心.
也有原因可以告訴自己.
我為什麼會信這位神.
我為什麼會這麼信他.
這麼聽他所說的話.
所以我們認識神越深.
跟從他的時候.
是更加有力.
更加不會覺得很乏力 很重擔.
要聽他所說的話的時候.
要很難.
這兩樣東西是很有關係的.
我們越深入的領會過.
或者經歷過神的時候.
我們就不會覺得.
聖經說了這麼多教訓 命令.
要做真的很難.
讀聖經要讀一下就沒問題.
聽講道的時候聽一下就沒問題.
查經的時候講一下都可以.
我們越能夠體會神的大能.
我們就不會覺得.
我們按著聖經所說的話.
或者我們去事奉這位神的時候.

$^{121}$就不會變成.
我要做到這些事情滿足神.
我就可以上天堂了.
或者我做了這些事情.
我跟他所說的 我聽他的話.
我就可以賺到多一些他給我們屬靈的福氣.
這個是我們很多時候都見到.
很多弟兄姊妹會掙扎的一些事.
所以今天這段經文裡面.
我就將所有我們讀到的屬靈福氣.
就把它們放在一起來看.
我就數五種的屬靈福氣出來.
我很多年前讀這段經文的時候.
就數七樣出來.
現在只剩下五樣不是少了兩樣.
不是說神收回 現在沒有了.
而是我覺得要記七樣好像比較難.
雖然是真的可以數七樣.
我們就將其中一些加在一起.
現在就用五個的屬靈福氣來看.
希望容易一點點會記得.
今天其中一個主題就是說.
我越知道我是神的兒女.
領受祂無限給我這麼多的屬靈福氣的時候.
我真的會體會在我未出生之前.
祂已經為我預備好這些屬靈福氣.
這個認識了解這種深信.
是令我們做基督徒的時候.
不會覺得很乏力.
不會覺得很難很辛苦.
因為我們帶著一個感恩的心的時候.
就更加能夠體會而有能力去跟從祂.
所以普羅今天告訴我們的這些.
可以說這麼多的屬靈福氣我們一向都知道.
但是有些時候因為我們信主不是很久.
沒有這麼深的體會或者經歷到.
有時候我們信主很久也好.
讀聖經的時候讀完又不是一定很了解得深入.
又不是很記得的時候.
都好像會差一點.

$^{161}$所以我們時時都要數出來.
我們唱詩歌,祈禱的時候都要數出來.
都是非常重要.
如果我們領會這些福氣不深的時候.
我們甚至去查經班,參加團契小組聚會.
甚至主日崇拜在這裡一起坐的時候.
我們都有機會當這些全部都是一種習慣.
所以做慣了就做.
有可能當是一種責任.
所以需要做的.
來崇拜的時候來團契的時候.
如果覺得好像參加一個社團的活動.
參加一個中親會.
我們可以怎樣去提升我們對神的了解呢?.
重溫這五個屬靈福氣是很重要的.
如果你時時覺得聽講道,查經是很悶的事.
有時候我們講道是有一點悶的.
除了這樣之外.
我們也要想想帶查經的那個帶得悶.
講道的那個講得悶.
要想想自己有沒有令到查經自己覺得悶.
或者那次講道自己覺得悶.
有一次很奇怪.
以前教會在唐人街大街很多人經過.
有一次很有趣的.
早上回到教會的時候.
見到門口放了一本書.
那時候我們很早還沒有開門的時候.
原來已經有人放在那裡.
到我們開門的時候.
我出去拿一本書.
這本書的題目很奇怪.
題目叫當你坐在那裡聽一篇悶的講道的時候.
有101樣東西你可以做.
你有沒有想過?.
當然第一樣可以做的就是睡覺.
或者第二樣可以做的就是拿個電話出來.
看看有什麼好看的.
一定有很多東西比聽道好看.
這本書裡面有101種.

$^{201}$今天不會講完.
我只講兩樣大家來參考一下.
看看有沒有用.
第45樣就是說.
你坐在你旁邊那個人的臉.
打他一巴掌.
不只是說他睡了所以打醒他.
他的作用就是除了打醒他之外.
你看看他會不會給你另一邊的臉打.
這是聖經的教導.
除了打醒他之外.
你還可以幫助他去反省一下.
馬太福音所教導的.
你有沒有願意做呢.
人家打你這邊的臉給他打那邊的臉.
這件事你可以參考一下.
然後第50樣.
第50樣就是.
你看著那個錶.
這個講完講到大約20分鐘.
你就預早set定.
一開始講到你就set.
到20分鐘那個鐘就會響.
就是那個錶的鬧鐘會響.
提醒那個講完已經講了20分鐘.
然後之後每5分鐘就響一次.
大家想不想試一下呢.
其實這件事就算你未試.
其實我們都已經預備了.
大家應該不會看到的.
因為到差不多的時候.
其實我們已經有定時的.
我們的司事長或者我們的幹事.
就會看著如果這次講到應該是.
譬如9點40分就應該完的.
就有人站在你後面.
就會舉一個牌就是說.
還有10分鐘.
然後5分鐘之後他又舉另外一個牌.
還有5分鐘.

$^{241}$當然有時候我們都不是很聽話.
或者我們講到非常專注的時候.
我們就不看他了.
所以有時候會過時的.
在這些裡面提醒我們.
時時我們數算神的恩典.
我們讀聖經有這樣讀.
詩歌有這樣唱.
越懂得數算這五種福氣的話.
就會越願意擺多些精神.
擺多些時間.
好好地讀聖經.
好好地崇拜.
好好地事奉.
所以在這裡講.
尤其是第三節開始的時候.
這裡有一個很特別的特色.
就是剛才我都幾乎想問Pastor Kenny.
如果你用這段經文來祈禱.
你的祈禱有多長?.
不知道大家剛才有沒有留意到.
讀這段經文.
由三節到十四節的時候.
用了多長的時間?.
我都嘗試去看一下時間.
是大約讀了三分鐘.
三分鐘的祈禱平時在牧土來說.
都不算很長.
但是保羅寫這段三至十四節裡面.
其實是一句話.
我們現在讀的聖經有很多標點符號.
但是其實保羅寫這段經文的時候.
三至十四節是一句話.
中間是沒有標點符號的.
幾乎是要一口氣讀完這十二節的經文.
你可以回家試一下.
保羅要這樣說.
因為他真的很感受到神的恩典.
他的屬靈的福氣是有多麼多.
多麼大.

$^{281}$多麼奇妙.
對我們的生命有多麼大的影響.
所以當我們能夠領會這些福氣的時候.
我們更加會明白他的恩典.
他的恩情是多麼豐厚.
以及知道我們怎樣去過一個感恩的生活.
第三節這裡已經講了主旨.
就是我們要去稱頌讚美我們的神.
他給我們的恩典福氣是這麼多.
我特意給大家注意的是應當稱頌.
其實這裡十二節的經文裡.
重複了幾次說要應當稱頌.
這裡就是說除了保羅要數出來.
神給我們多少屬靈的福氣之外.
就已經包含了.
既然我們知道神給我們這麼多屬靈的福氣的話.
所有這些都是應當稱頌的.
所以希望大家將來唱詩歌讚美神的時候.
真的要用心去唱.
因為詩歌裡面就唱出了神應當稱頌的所有恩典.
講來講去講了這麼久都還沒有講.
那些五個的恩典福氣究竟是什麼呢.
要這麼說.
因為每一個屬靈的福氣都是這麼豐富.
所以有人將其中一個屬靈的福氣可以寫一本書.
我們因為時間關係.
所以就要很快地略過每一個的福氣.
只是講一點點來提醒大家.
然後希望大家將來自己再深入地去研究和體會.
第一個的福氣就是.
神揀選我們成為聖潔沒有瑕疵.
不過揀選是神主動去做的事.
揀選了我們是因為祂的愛.
而揀選了我們之後.
使我們成為聖潔和沒有瑕疵.
這是要記住的第一個神給我們的福氣.
祂是揀選我們成為聖潔沒有瑕疵.
祂揀選我們是因為愛.
這裡有幾樣東西.
每一樣都可以慢慢去深入體會去研讀.

$^{321}$大家記住這個粗略的大綱就行了.
大家可能有很多的疑問.
如果你以前沒有認真去思考過.
你會有很多的疑問.
揀選是怎樣揀選的呢?.
用什麼方法揀選呢?.
然後揀選了我們成為聖潔沒有瑕疵.
什麼叫聖潔沒有瑕疵?.
是不是我們在世上生活的時候.
要成為一個很完美的十全十美的人.
這個是神揀選我.
於是我就要成為這樣十全十美的人.
是還是不是呢?.
中間有沒有位置可以走棧呢?.
這些就留給大家繼續去思考.
繼續去研讀.
第二個福氣我們成為了神的兒女.
這個是非常非常之重要.
因為我們明白了神揀選了我們.
其實是揀選我們除了剛才說成為聖潔沒有瑕疵之外.
是揀選了我們做祂的兒女.
這個是很重要的.
因為我們本來就不是祂的兒女.
但是祂創造我們的生命.
然後還揀選我們做祂的兒女.
有一點點好像是收養我們.
收養這件事尤其是在羅馬當時的文化裡面.
就有一樣我們很難明白也很難接受的事.
就是在當時的羅馬文化裡面.
媽媽生了一個BB出來.
然後爸爸就來看一下她.
因為以前不像現在這樣.
可以跟著進產房自己去剪磁帶.
生了然後爸爸才進來看一下.
他們奇怪的地方就是.
爸爸來看一下是兒子還是女兒.
看完之後他就可以做一個決定.
這個決定就是他想要還是不想要這個BB.
他可以當時去決定.
可以有很多因素的.

$^{361}$他想要女兒怎知道生了個兒子.
想要兒子怎知道生了個女兒.
生出來的BB是有什麼殘疾的地方.
生得不好還是生得太好.
他可以不要.
什麼原因都好.
他自己可以決定.
當時有這樣的奇怪的習俗.
如果他不要的時候.
就可以叫人拿他去市集那裡.
放在那裡等人來拿走他.
等人來收養他.
有少少這樣的文化背景.
所以當時保羅寫這些經文的時候.
大家都明白收養有多重要.
如果沒有人收養他.
BB放在那裡很快就會離開這個世界.
如果有人收養他.
收養了他就不是真的養他.
他雖然也要養他.
但他可以將他收了成為他家裡的奴隸.
這些奴隸可能要做苦工那種.
甚至他收了他之後不養他.
再將他賣出來.
賣給其他人.
他在市集撿回來的.
不用錢的.
然後他可以賣了他.
賣了他之後他就不理了.
賣了他之後BB將來可以成為一個奴隸.
甚至成為一個娼妓.
但在這個背景裡面.
保羅告訴我們.
神是收養我們.
收養了我們還稱為兒女.
當我們是他的子女.
真是收養了然後栽培.
然後很珍惜.
所以看到那個很大的對比.
更加我們體會神收養我們.

$^{401}$叫我們成為他的兒女.
是多麼重要.
多麼珍貴.
多麼奇妙.
是完全出於神的愛.
所以他收養我們成為他的兒女.
所以這裡再說.
神這個恩典的榮耀.
是要得著稱讚.
第三個福氣.
我們被救贖被赦免.
這個查經也好.
崇拜也好.
參加了幾年.
時時都會再提這件事.
但有時又是.
好像我今天這樣.
救贖很快就說完了.
大家未必有好好坐下來思想這件事.
使這件事能夠深深地藏在我們心裡.
因此而知道怎樣去跟從神.
怎樣去遵從神的說話.
我們被救贖.
以前大家都聽過.
那些奴隸.
當時很盛行有奴隸這種制度.
或者奴隸的市場.
所以在市集裡.
不只是買菜買肉那些.
隨時走進去.
可以買賣人口.
買賣奴隸.
所以他們非常熟悉這件事.
如果有一個人來到市集.
見到這個奴隸.
他喜歡無論什麼原因都好.
總之他喜歡他就可以去買了他.
買了他之後.
有時我們看那些.
中國古代的電影劇集都有.

$^{441}$本來有一個是被賣身.
去清流做娼妓.
但有一個人覺得.
他出於污泥而不染.
非常之好.
於是很想幫他贖身.
就是這種背景.
耶穌基督來到這個世界.
替我們贖身.
將我們賣熟.
而且他將我們賣熟的時候.
不是拿現金.
不是拿信用卡.
而是用他自己的生命.
是用他的血.
自己流出的血.
捨棄的生命.
來買熟我們.
買熟了我們之後.
他還知道我們被罪捆綁.
所以他赦免我們的罪.
這個也是神主動來做.
充充足足地賜給我們的.
這些我們要數算.
要唱出來.
要和周圍的人說.
無論周圍的人是信主也好.
未信主也好.
我們都要去說.
都要去讚頌神.
第四個福氣是使我們知道他的旨意.
你想像一下這個世界的創造主.
這個世界的救贖主.
他那種偉大奇妙.
是我們不能夠想像.
但是他竟然和我們說他的旨意.
就是說他有什麼秘密.
他有什麼心事.
他會和我們說.
這裡是不容易明白的.

$^{481}$因為是照著神的美意.
基督預先安排的美意.
使我們知道他旨意的奧秘.
不是只知道他的一些心事.
還要知道他的心事裡的奧秘.
這個奧秘是什麼呢.
大家留心聽第二第三章.
你就會知道這個奧秘究竟在說什麼.
這個奧秘我發覺很多人不明白.
信了耶穌很多年都不明白這個奧秘是什麼.
甚至老實說我聽過一些傳道人.
凡是說一些很難明白的東西.
他就說聖經說這些是奧秘.
所以我們不明白.
當然不容易明白.
你一定要深入去看去想.
總之我們不明白神的東西就是奧秘.
這裡其實不是這麼說.
所以你更加要留心聽第二第三章.
記得要來崇拜.
還有不要睡覺.
給大家留意一點.
十節說到都在基督裡同歸於一.
跟這個奧秘是很有關係的.
上次我們再有說到.
在基督裡合一究竟是什麼.
最後一個.
我們成為神的產業.
很多聖經的翻譯.
都是說我們會得到耶穌的產業.
但其實有一個更加真實.
更加全神或者更加正確的翻譯.
是應該說我們成為了神的產業.
凡是你自己的產業.
你都會很珍惜它很保護它.
它不會隨便亂花你所有的身家.
凡是屬於你的東西.
就是你的產業.
屬於你的東西你會很小心去保管.
小心去使用.

$^{521}$小心去收藏各樣東西都好.
我們是神的產業.
你想想這件事.
可以說是多麼難明.
多麼奇妙.
我們原來是神的產業.
我們是屬於祂的.
祂會保護我們.
祂會好好地使用我們.
不是當我們是奴隸般使用.
而是好好地使用.
然後還有聖靈作為因機和憑據.
我以前把這兩樣東西分開.
不過後來我覺得.
把它們放在一起.
更加全神.
我們是神的產業.
而且聖靈成為了一個因機和憑據.
那個因機的意思就是說.
如果你買了一幅地.
買了一棟樓.
買了一間屋.
就會扔印.
有些人用圖章.
或者是任何正式的一個印記.
扔印.
以前那些農業社會裡面.
如果這隻牛是屬於你的.
你會扔印.
如果這隻馬屬於你的.
你會扔印.
就是那種印記.
證明我們是屬於神的.
憑據的意思是.
有些是好像我們買樓的時候.
你就給了十分之一的訂金.
就證明了你是買了.
甚至好像我們買屋的時候.
你買了一間屋.
我問你那間屋是不是屬於你的.

$^{561}$你說是屬於你的.
不過原來不是完全屬於你的.
因為你還在供.
但是有了那個憑據.
扔了那個印.
就證明那個產業已經是你的了.
但將來有更多.
將來你要供完那間屋.
才完全屬於你的.
就好像我們現在在世上.
我們的基督徒.
我們是等候將來.
我們完完全全進入到.
所謂天堂的裡面.
我們也說.
那個天堂不只是將來才存在.
你在世上已經是在地若天.
在神的國裡面.
這五個屬靈的福氣.
其實還帶來無限的福氣.
這個憑據已經說了.
不多說了.
大家再看一看.
在這四節經文裡面.
重複地說.
神是應當稱重的.
祂的恩典是應當稱重的.
祂的榮耀是應當稱重的.
所以希望大家以後的崇拜.
不要只想著聽道.
最重要是因為神的說話.
所以講道之前利刀就行了.
我們經常強調.
我們唱詩歌是事奉的很重要的元素.
因為我們讚美神就是我們的事奉.
所以沒有人可以說.
我在教會沒有任何事奉.
這種想法是完全要改的.
因為你來唱詩是一個讚美的事奉.
你來祈禱是一個讚美.

$^{601}$祈求感恩的事奉.
神是主動使我們得到這些救恩.
所以我們更加明白我們的身份.
不再是罪的奴僕.
而是神的兒女.
我們被買熟.
我們是天國的子民.
我們是神的兒女.
我們是耶穌的門徒.
願大家繼續在這些各樣的屬靈恩典裡繼續成長.
我們禱告.
多謝聖靈住在我們生命裡.
讓我們更加肯定我們是屬於神的.
雖然我們沒有時間去深入思想.
聖靈住在我們生命裡的重要性.
怎樣帶領我們向前走.
怎樣讓我們有能力去遵從神的話.
怎樣讓我們有能力去事奉我們的神.
主當我們數出這些恩典的時候.
我們的生命就得著更大的豐盛.
我們過一個感恩的人生的時候.
亦都過一個事奉的人生.
我們祈求禱告奉靠救主耶穌基督的名.
阿們.
\newpage



\section{出埃及記 25:1-40:38}
\label{sec:032256}
\textbf{會幕}~~~~陳崇基牧師
\newline
\newline
連結: \href{http://www.church.com.hk/acms/content.asp?site=cdc\&op=show\&type=product\&code=032256\&layout=sermon}{\texttt{032256}} ~~~~ 語音日期: 2026-06-07
\newline
\newline
\hyperref[sec:032255]{\small{< < < PREV SERMON < < <}}
~
\hyperref[sec:index_chronic]{\small{[返順時目]}}
~
\hyperref[sec:032257]{\small{> > > NEXT SERMON > > >}}
\newline
\newline
出埃及記 25:1-40:38
\newline
\begin{longtable}{cl}
\hline
\hline
章節 & 經文 (和合本修訂版)\\
\hline
25:1 & \begin{tabularx}{0.7\textwidth}{X} 耶和華吩咐摩西說： \end{tabularx} \\ \\ \relax
25:2 & \begin{tabularx}{0.7\textwidth}{X} 「你要吩咐以色列人獻禮物給我。凡甘心樂意獻給我的禮物，你們都可以收下。 \end{tabularx} \\ \\ \relax
25:3 & \begin{tabularx}{0.7\textwidth}{X} 要從他們收的禮物是：金、銀、銅， \end{tabularx} \\ \\ \relax
25:4 & \begin{tabularx}{0.7\textwidth}{X} 藍色、紫色、朱紅色紗，細麻，山羊毛， \end{tabularx} \\ \\ \relax
25:5 & \begin{tabularx}{0.7\textwidth}{X} 染紅的公羊皮、精美的皮料，金合歡木， \end{tabularx} \\ \\ \relax
25:6 & \begin{tabularx}{0.7\textwidth}{X} 點燈的油，做膏油的香料、做香的香料， \end{tabularx} \\ \\ \relax
25:7 & \begin{tabularx}{0.7\textwidth}{X} 紅瑪瑙與寶石，可以鑲嵌在以弗得和胸袋上。 \end{tabularx} \\ \\ \relax
25:8 & \begin{tabularx}{0.7\textwidth}{X} 他們要為我造聖所，使我住在他們中間。 \end{tabularx} \\ \\ \relax
25:9 & \begin{tabularx}{0.7\textwidth}{X} 你們要按照我指示你的，帳幕和其中一切器具的樣式，照樣去做。」 \end{tabularx} \\ \\ \relax
25:10 & \begin{tabularx}{0.7\textwidth}{X} 「他們要用金合歡木做一個櫃子，長二肘半，寬一肘半，高一肘半。 \end{tabularx} \\ \\ \relax
25:11 & \begin{tabularx}{0.7\textwidth}{X} 你要把它裡裡外外包上純金，四圍要鑲上金邊。 \end{tabularx} \\ \\ \relax
25:12 & \begin{tabularx}{0.7\textwidth}{X} 要鑄造四個金環，安在櫃子的四腳上；這邊兩個環，那邊兩個環。 \end{tabularx} \\ \\ \relax
25:13 & \begin{tabularx}{0.7\textwidth}{X} 要用金合歡木做兩根槓，包上金子。 \end{tabularx} \\ \\ \relax
25:14 & \begin{tabularx}{0.7\textwidth}{X} 要把槓穿過櫃旁的環，以便抬櫃。 \end{tabularx} \\ \\ \relax
25:15 & \begin{tabularx}{0.7\textwidth}{X} 這槓要留在櫃的環內，不可抽出來。 \end{tabularx} \\ \\ \relax
25:16 & \begin{tabularx}{0.7\textwidth}{X} 要把我所要賜給你的法版放在櫃裡。 \end{tabularx} \\ \\ \relax
25:17 & \begin{tabularx}{0.7\textwidth}{X} 要用純金做一個櫃蓋，長二肘半，寬一肘半。 \end{tabularx} \\ \\ \relax
25:18 & \begin{tabularx}{0.7\textwidth}{X} 要造兩個用金子錘出的基路伯，從櫃蓋的兩端錘出它們。 \end{tabularx} \\ \\ \relax
25:19 & \begin{tabularx}{0.7\textwidth}{X} 這端錘出一個基路伯，那端錘出一個基路伯；從櫃蓋的兩端錘出兩個基路伯。 \end{tabularx} \\ \\ \relax
25:20 & \begin{tabularx}{0.7\textwidth}{X} 二基路伯的翅膀要向上張開，用翅膀遮住櫃蓋，臉要彼此相對；基路伯的臉要朝向櫃蓋。 \end{tabularx} \\ \\ \relax
25:21 & \begin{tabularx}{0.7\textwidth}{X} 要把櫃蓋安在櫃的上邊，又要把我所要賜給你的法版放在櫃裡。 \end{tabularx} \\ \\ \relax
25:22 & \begin{tabularx}{0.7\textwidth}{X} 我要在那裡與你相會，並要從法版之櫃的櫃蓋上，兩個基路伯的中間，將我要吩咐以色列人的一切事告訴你。」 \end{tabularx} \\ \\ \relax
25:23 & \begin{tabularx}{0.7\textwidth}{X} 「你要用金合歡木做一張供桌，長二肘，寬一肘，高一肘半， \end{tabularx} \\ \\ \relax
25:24 & \begin{tabularx}{0.7\textwidth}{X} 把它包上純金，四圍鑲上金邊。 \end{tabularx} \\ \\ \relax
25:25 & \begin{tabularx}{0.7\textwidth}{X} 供桌的四圍各做一掌寬的邊緣，邊緣周圍要鑲上金邊。 \end{tabularx} \\ \\ \relax
25:26 & \begin{tabularx}{0.7\textwidth}{X} 要為供桌做四個金環，把環安在四個桌腳的四角上。 \end{tabularx} \\ \\ \relax
25:27 & \begin{tabularx}{0.7\textwidth}{X} 環要靠近邊緣，以便穿槓抬供桌。 \end{tabularx} \\ \\ \relax
25:28 & \begin{tabularx}{0.7\textwidth}{X} 要用金合歡木做兩根槓，包上金子，用來抬供桌。 \end{tabularx} \\ \\ \relax
25:29 & \begin{tabularx}{0.7\textwidth}{X} 要用純金做桌上的盤、碟，以及澆酒祭的壺和杯。 \end{tabularx} \\ \\ \relax
25:30 & \begin{tabularx}{0.7\textwidth}{X} 要把供餅擺在桌上，常在我面前。」 \end{tabularx} \\ \\ \relax
25:31 & \begin{tabularx}{0.7\textwidth}{X} 「你要造一座用純金錘出的燈臺。燈臺的座、幹、杯、花萼和花瓣，都要和燈臺接連一塊。 \end{tabularx} \\ \\ \relax
25:32 & \begin{tabularx}{0.7\textwidth}{X} 燈臺兩旁要伸出六根枝子：這邊三根，那邊三根。 \end{tabularx} \\ \\ \relax
25:33 & \begin{tabularx}{0.7\textwidth}{X} 這邊枝子上有三個杯，形狀像杏花，有花萼有花瓣；那邊枝子上也有三個杯，形狀像杏花，有花萼有花瓣。從燈臺伸出來的六根枝子都是如此。 \end{tabularx} \\ \\ \relax
25:34 & \begin{tabularx}{0.7\textwidth}{X} 燈臺本身要有四個杯，形狀像杏花，有花萼有花瓣。 \end{tabularx} \\ \\ \relax
25:35 & \begin{tabularx}{0.7\textwidth}{X} 燈臺的第一對枝子下面有花萼，燈臺的第二對枝子下面有花萼，燈臺的第三對枝子下面也有花萼；燈臺伸出的六根枝子都是如此。 \end{tabularx} \\ \\ \relax
25:36 & \begin{tabularx}{0.7\textwidth}{X} 花萼和枝子都要和燈臺接連一塊，全是從一塊純金錘出來的。 \end{tabularx} \\ \\ \relax
25:37 & \begin{tabularx}{0.7\textwidth}{X} 要做燈臺的七盞燈，燈要點燃，照亮前面。 \end{tabularx} \\ \\ \relax
25:38 & \begin{tabularx}{0.7\textwidth}{X} 要用純金做燈剪和燈盤。 \end{tabularx} \\ \\ \relax
25:39 & \begin{tabularx}{0.7\textwidth}{X} 做燈臺和這一切的器具要用一他連得純金。 \end{tabularx} \\ \\ \relax
25:40 & \begin{tabularx}{0.7\textwidth}{X} 要謹慎，照著在山上指示你的樣式去做。」 \end{tabularx} \\ \\ \relax
26:1 & \begin{tabularx}{0.7\textwidth}{X} 「你要用十幅幔子做帳幕。這些幔子要用搓的細麻和藍色、紫色、朱紅色紗織成，並且以刺繡的手藝繡上基路伯。 \end{tabularx} \\ \\ \relax
26:2 & \begin{tabularx}{0.7\textwidth}{X} 每幅幔子要長二十八肘，每幅幔子寬四肘，全部的幔子尺寸都要一樣。 \end{tabularx} \\ \\ \relax
26:3 & \begin{tabularx}{0.7\textwidth}{X} 這五幅幔子要彼此相連；那五幅幔子也彼此相連。 \end{tabularx} \\ \\ \relax
26:4 & \begin{tabularx}{0.7\textwidth}{X} 在這一組相連幔子的末幅邊上要縫藍色的鈕環；在另一組相連幔子的末幅邊上也要照樣做。 \end{tabularx} \\ \\ \relax
26:5 & \begin{tabularx}{0.7\textwidth}{X} 這幅幔子上要縫五十個鈕環，另一組相連幔子的末幅上也縫五十個鈕環，環環相對。 \end{tabularx} \\ \\ \relax
26:6 & \begin{tabularx}{0.7\textwidth}{X} 要做五十個金鉤，用鉤子使幔子彼此相連，成為一個帳幕。 \end{tabularx} \\ \\ \relax
26:7 & \begin{tabularx}{0.7\textwidth}{X} 「你要用山羊毛織十一幅幔子來作帳幕的罩棚。 \end{tabularx} \\ \\ \relax
26:8 & \begin{tabularx}{0.7\textwidth}{X} 每幅幔子要長三十肘，每幅幔子寬四肘；十一幅幔子的尺寸都要一樣。 \end{tabularx} \\ \\ \relax
26:9 & \begin{tabularx}{0.7\textwidth}{X} 要把五幅幔子連成一幅，又把六幅幔子連成一幅，這第六幅幔子要在罩棚的前面摺上去。 \end{tabularx} \\ \\ \relax
26:10 & \begin{tabularx}{0.7\textwidth}{X} 在這一組相連幔子的末幅邊上要縫五十個鈕環；在另一組相連幔子的末幅邊上也縫五十個鈕環。 \end{tabularx} \\ \\ \relax
26:11 & \begin{tabularx}{0.7\textwidth}{X} 要做五十個銅鉤，鉤在鈕環中，使罩棚相連成為一個。 \end{tabularx} \\ \\ \relax
26:12 & \begin{tabularx}{0.7\textwidth}{X} 罩棚幔子餘下垂著的，那餘下的半幅要垂在帳幕的背面。 \end{tabularx} \\ \\ \relax
26:13 & \begin{tabularx}{0.7\textwidth}{X} 罩棚的幔子兩旁所餘下的，這邊一肘，那邊一肘，要垂在帳幕的兩邊，蓋住帳幕。 \end{tabularx} \\ \\ \relax
26:14 & \begin{tabularx}{0.7\textwidth}{X} 要用染紅的公羊皮做罩棚的蓋，再用精美皮料做外層的蓋。 \end{tabularx} \\ \\ \relax
26:15 & \begin{tabularx}{0.7\textwidth}{X} 「你要用金合歡木做豎立帳幕的木板， \end{tabularx} \\ \\ \relax
26:16 & \begin{tabularx}{0.7\textwidth}{X} 木板要長十肘，每塊板寬一肘半， \end{tabularx} \\ \\ \relax
26:17 & \begin{tabularx}{0.7\textwidth}{X} 每塊板有兩個榫頭可以彼此銜接。帳幕一切的板都要這樣做。 \end{tabularx} \\ \\ \relax
26:18 & \begin{tabularx}{0.7\textwidth}{X} 你要做帳幕的木板：南面，就是面向南方的那一邊，要做二十塊板。 \end{tabularx} \\ \\ \relax
26:19 & \begin{tabularx}{0.7\textwidth}{X} 在這二十塊板底下要做四十個帶卯眼的銀座；兩個卯眼接連這塊板上的兩個榫頭，另外兩個卯眼接連那塊板上的兩個榫頭。 \end{tabularx} \\ \\ \relax
26:20 & \begin{tabularx}{0.7\textwidth}{X} 帳幕的第二邊，就是北面，也要做二十塊板， \end{tabularx} \\ \\ \relax
26:21 & \begin{tabularx}{0.7\textwidth}{X} 和四十個帶卯眼的銀座；這塊板底下有兩個卯眼，那塊板底下也有兩個卯眼。 \end{tabularx} \\ \\ \relax
26:22 & \begin{tabularx}{0.7\textwidth}{X} 帳幕的後面，就是西面，要做六塊板。 \end{tabularx} \\ \\ \relax
26:23 & \begin{tabularx}{0.7\textwidth}{X} 帳幕後面的角落要做兩塊板。 \end{tabularx} \\ \\ \relax
26:24 & \begin{tabularx}{0.7\textwidth}{X} 下端的板是成雙的，上端要連在一起，直到頂端的第一個環子；兩塊板都要這樣，做成兩個角落。 \end{tabularx} \\ \\ \relax
26:25 & \begin{tabularx}{0.7\textwidth}{X} 一共有八塊板和十六個帶卯眼的銀座；這塊板底下有兩個卯眼，那塊板底下也有兩個卯眼。 \end{tabularx} \\ \\ \relax
26:26 & \begin{tabularx}{0.7\textwidth}{X} 「你要用金合歡木做橫木：為帳幕這面的板做五根橫木， \end{tabularx} \\ \\ \relax
26:27 & \begin{tabularx}{0.7\textwidth}{X} 為帳幕那面的板做五根橫木，又為帳幕後面，就是朝西的板做五根橫木。 \end{tabularx} \\ \\ \relax
26:28 & \begin{tabularx}{0.7\textwidth}{X} 板腰間的橫木，要從一頭通到另一頭。 \end{tabularx} \\ \\ \relax
26:29 & \begin{tabularx}{0.7\textwidth}{X} 板要包上金子，又要做板上的金環來套橫木；橫木也要包上金子。 \end{tabularx} \\ \\ \relax
26:30 & \begin{tabularx}{0.7\textwidth}{X} 要照著在山上所指示你的樣式，把帳幕豎立起來。 \end{tabularx} \\ \\ \relax
26:31 & \begin{tabularx}{0.7\textwidth}{X} 「你要用藍色、紫色、朱紅色紗，和搓的細麻織幔子，以刺繡的手藝繡上基路伯。 \end{tabularx} \\ \\ \relax
26:32 & \begin{tabularx}{0.7\textwidth}{X} 要把幔子掛在四根包金的金合歡木柱子上，柱子有金鉤，並且安在四個帶卯眼的銀座上。 \end{tabularx} \\ \\ \relax
26:33 & \begin{tabularx}{0.7\textwidth}{X} 要把幔子垂掛在鉤子上，把法櫃抬進幔子內；這幔子要將聖所和至聖所隔開。 \end{tabularx} \\ \\ \relax
26:34 & \begin{tabularx}{0.7\textwidth}{X} 又要把櫃蓋安在至聖所內的法櫃上， \end{tabularx} \\ \\ \relax
26:35 & \begin{tabularx}{0.7\textwidth}{X} 把供桌安在幔子的外面，供桌在北面，燈臺在帳幕的南面，和供桌相對。 \end{tabularx} \\ \\ \relax
26:36 & \begin{tabularx}{0.7\textwidth}{X} 「你要用藍色、紫色、朱紅色紗，和搓的細麻，以刺繡的手藝為帳幕織門簾。 \end{tabularx} \\ \\ \relax
26:37 & \begin{tabularx}{0.7\textwidth}{X} 要用金合歡木為簾子做五根柱子，包上金子。柱子有金鉤，又為柱子鑄造五個帶卯眼的銅座。」 \end{tabularx} \\ \\ \relax
27:1 & \begin{tabularx}{0.7\textwidth}{X} 「你要用金合歡木做祭壇，長五肘，寬五肘，這壇是正方形的，高三肘。 \end{tabularx} \\ \\ \relax
27:2 & \begin{tabularx}{0.7\textwidth}{X} 要在壇的四角做四個翹角，與壇接連一塊；要把壇包上銅。 \end{tabularx} \\ \\ \relax
27:3 & \begin{tabularx}{0.7\textwidth}{X} 要做桶子來盛壇上的灰，又要做鏟子、盤子、肉叉和火盆；壇上一切的器具都要用銅做。 \end{tabularx} \\ \\ \relax
27:4 & \begin{tabularx}{0.7\textwidth}{X} 要為壇做一個銅網，在網的四角做四個銅環， \end{tabularx} \\ \\ \relax
27:5 & \begin{tabularx}{0.7\textwidth}{X} 把網安在壇四圍的邊的下面，使網垂到壇的半腰。 \end{tabularx} \\ \\ \relax
27:6 & \begin{tabularx}{0.7\textwidth}{X} 又要用金合歡木為壇做槓，包上銅。 \end{tabularx} \\ \\ \relax
27:7 & \begin{tabularx}{0.7\textwidth}{X} 這槓要穿過壇兩旁的環子，用來抬壇。 \end{tabularx} \\ \\ \relax
27:8 & \begin{tabularx}{0.7\textwidth}{X} 要用板做壇，壇的中心是空的，都照著在山上所指示你的樣式做。」 \end{tabularx} \\ \\ \relax
27:9 & \begin{tabularx}{0.7\textwidth}{X} 「你要做帳幕的院子。南面，就是面向南方的那一邊，要用搓的細麻做院子的帷幔，長一百肘， \end{tabularx} \\ \\ \relax
27:10 & \begin{tabularx}{0.7\textwidth}{X} 院子要有二十根柱子，二十個帶卯眼的銅座。要用銀做柱子的鉤和箍。 \end{tabularx} \\ \\ \relax
27:11 & \begin{tabularx}{0.7\textwidth}{X} 北面的長度也一樣，帷幔長一百肘，要有二十根柱子，二十個帶卯眼的銅座。要用銀做柱子的鉤和箍。 \end{tabularx} \\ \\ \relax
27:12 & \begin{tabularx}{0.7\textwidth}{X} 院子的西面有帷幔，寬五十肘，帷幔要有十根柱子，十個帶卯眼的座。 \end{tabularx} \\ \\ \relax
27:13 & \begin{tabularx}{0.7\textwidth}{X} 院子的東面，就是面向東方的那一邊，寬五十肘。 \end{tabularx} \\ \\ \relax
27:14 & \begin{tabularx}{0.7\textwidth}{X} 一邊的帷幔有十五肘，要有三根柱子，三個帶卯眼的座。 \end{tabularx} \\ \\ \relax
27:15 & \begin{tabularx}{0.7\textwidth}{X} 另一邊的帷幔也有十五肘，要有三根柱子，三個帶卯眼的座。 \end{tabularx} \\ \\ \relax
27:16 & \begin{tabularx}{0.7\textwidth}{X} 院子的門要有二十肘長的簾子，用藍色、紫色、朱紅色紗，和搓的細麻，以刺繡的手藝織成；要有四根柱子，四個帶卯眼的座。 \end{tabularx} \\ \\ \relax
27:17 & \begin{tabularx}{0.7\textwidth}{X} 院子四圍一切的柱子都要用銀子箍著，要用銀做柱子的鉤子，用銅做帶卯眼的座。 \end{tabularx} \\ \\ \relax
27:18 & \begin{tabularx}{0.7\textwidth}{X} 院子要長一百肘，寬五十肘，高五肘。要用搓的細麻做帷幔，用銅做帶卯眼的座。 \end{tabularx} \\ \\ \relax
27:19 & \begin{tabularx}{0.7\textwidth}{X} 帳幕中各樣用途的器具，以及帳幕一切的橛子和院子裡一切的橛子，都要用銅做。」 \end{tabularx} \\ \\ \relax
27:20 & \begin{tabularx}{0.7\textwidth}{X} 「你要吩咐以色列人，把搗成的純橄欖油拿來給你，用以點燈，使燈經常點著； \end{tabularx} \\ \\ \relax
27:21 & \begin{tabularx}{0.7\textwidth}{X} 在會幕中法櫃前的幔子外，亞倫和他的兒子要從晚上到早晨，在耶和華面前照管這燈。這要成為以色列人世世代代永遠的定例。」 \end{tabularx} \\ \\ \relax
28:1 & \begin{tabularx}{0.7\textwidth}{X} 「你要從以色列人中，叫你的哥哥亞倫和他的兒子拿答、亞比戶、以利亞撒、以他瑪一同親近你，作事奉我的祭司。 \end{tabularx} \\ \\ \relax
28:2 & \begin{tabularx}{0.7\textwidth}{X} 你要為你哥哥亞倫做聖衣，以示尊嚴和華美。 \end{tabularx} \\ \\ \relax
28:3 & \begin{tabularx}{0.7\textwidth}{X} 要吩咐一切心中有智慧的，就是我用智慧的靈所充滿的人，為亞倫做衣服，使他分別為聖，作事奉我的祭司。 \end{tabularx} \\ \\ \relax
28:4 & \begin{tabularx}{0.7\textwidth}{X} 所要做的是胸袋、以弗得、外袍、織成的內袍、禮冠和腰帶。他們要為你哥哥亞倫和他的兒子做聖衣，使他們作祭司事奉我。 \end{tabularx} \\ \\ \relax
28:5 & \begin{tabularx}{0.7\textwidth}{X} 要用金色、藍色、紫色、朱紅色紗，和細麻去縫製。」 \end{tabularx} \\ \\ \relax
28:6 & \begin{tabularx}{0.7\textwidth}{X} 「他們要用金色、藍色、紫色、朱紅色紗，和搓的細麻，以刺繡的手藝做以弗得。 \end{tabularx} \\ \\ \relax
28:7 & \begin{tabularx}{0.7\textwidth}{X} 以弗得當有兩條肩帶，接上兩端，使它相連。 \end{tabularx} \\ \\ \relax
28:8 & \begin{tabularx}{0.7\textwidth}{X} 以弗得的精緻帶子，要以一樣的手藝，用金色、藍色、紫色、朱紅色紗，和搓的細麻縫製，與以弗得接連在一起。 \end{tabularx} \\ \\ \relax
28:9 & \begin{tabularx}{0.7\textwidth}{X} 要取兩塊紅瑪瑙，在上面刻以色列兒子的名字： \end{tabularx} \\ \\ \relax
28:10 & \begin{tabularx}{0.7\textwidth}{X} 六個名字在一塊寶石上，六個名字在另一塊寶石上，都按照他們出生的次序。 \end{tabularx} \\ \\ \relax
28:11 & \begin{tabularx}{0.7\textwidth}{X} 要以雕刻寶石的手藝，如同刻印章，把以色列兒子的名字刻在這兩塊寶石上，並把寶石鑲在金槽裡。 \end{tabularx} \\ \\ \relax
28:12 & \begin{tabularx}{0.7\textwidth}{X} 要把這兩塊寶石安在以弗得的兩條肩帶上，為以色列人作紀念石。亞倫要在耶和華面前把他們的名字帶在兩肩上，作為紀念。 \end{tabularx} \\ \\ \relax
28:13 & \begin{tabularx}{0.7\textwidth}{X} 要用金子做兩個槽， \end{tabularx} \\ \\ \relax
28:14 & \begin{tabularx}{0.7\textwidth}{X} 再用純金打兩條鏈子，像編成的繩子一樣，把這編成的金鏈扣在槽上。」 \end{tabularx} \\ \\ \relax
28:15 & \begin{tabularx}{0.7\textwidth}{X} 「你要以刺繡的手藝做一個決斷的胸袋，和做以弗得的方法一樣，用金色、藍色、紫色、朱紅色紗，和搓的細麻縫製。 \end{tabularx} \\ \\ \relax
28:16 & \begin{tabularx}{0.7\textwidth}{X} 胸袋是正方形的，疊成兩層，長一虎口，寬一虎口。 \end{tabularx} \\ \\ \relax
28:17 & \begin{tabularx}{0.7\textwidth}{X} 要在上面鑲四行寶石：第一行是紅寶石、紅璧璽、紅玉； \end{tabularx} \\ \\ \relax
28:18 & \begin{tabularx}{0.7\textwidth}{X} 第二行是綠寶石、藍寶石、金剛石； \end{tabularx} \\ \\ \relax
28:19 & \begin{tabularx}{0.7\textwidth}{X} 第三行是紫瑪瑙、白瑪瑙、紫晶； \end{tabularx} \\ \\ \relax
28:20 & \begin{tabularx}{0.7\textwidth}{X} 第四行是水蒼玉、紅瑪瑙、碧玉。這些都要鑲在金槽中。 \end{tabularx} \\ \\ \relax
28:21 & \begin{tabularx}{0.7\textwidth}{X} 這些寶石要有以色列十二個兒子的名字，如同刻印章，每一顆有自己的名字，代表十二個支派。 \end{tabularx} \\ \\ \relax
28:22 & \begin{tabularx}{0.7\textwidth}{X} 要在胸袋上用純金打鏈子，像編成的繩子一樣。 \end{tabularx} \\ \\ \relax
28:23 & \begin{tabularx}{0.7\textwidth}{X} 要為胸袋做兩個金環，把這兩個環安在胸袋的兩端。 \end{tabularx} \\ \\ \relax
28:24 & \begin{tabularx}{0.7\textwidth}{X} 要把那兩條編成的金鏈繫在胸袋兩端的兩個環上。 \end{tabularx} \\ \\ \relax
28:25 & \begin{tabularx}{0.7\textwidth}{X} 又要把鏈子的另外兩端扣在兩個槽上，安在以弗得前面的肩帶上。 \end{tabularx} \\ \\ \relax
28:26 & \begin{tabularx}{0.7\textwidth}{X} 要做兩個金環，安在胸袋的兩端，在以弗得裡面的邊上。 \end{tabularx} \\ \\ \relax
28:27 & \begin{tabularx}{0.7\textwidth}{X} 再做兩個金環，安在以弗得前面兩條肩帶的下邊，靠近接縫處，在以弗得精緻帶子的上面。 \end{tabularx} \\ \\ \relax
28:28 & \begin{tabularx}{0.7\textwidth}{X} 要用藍色的帶子把胸袋的環與以弗得的環繫住，使胸袋綁在以弗得的精緻帶子上，不致鬆脫。 \end{tabularx} \\ \\ \relax
28:29 & \begin{tabularx}{0.7\textwidth}{X} 亞倫進聖所的時候，要把刻著以色列兒子名字的決斷胸袋帶著，放在心上，在耶和華面前常作紀念。 \end{tabularx} \\ \\ \relax
28:30 & \begin{tabularx}{0.7\textwidth}{X} 又要將烏陵和土明放在決斷胸袋裡；亞倫進到耶和華面前的時候，要放在心上。這樣，亞倫在耶和華面前要把以色列人的決斷胸袋常常帶著，放在心上。」 \end{tabularx} \\ \\ \relax
28:31 & \begin{tabularx}{0.7\textwidth}{X} 「你要做以弗得的外袍，顏色全是藍的。 \end{tabularx} \\ \\ \relax
28:32 & \begin{tabularx}{0.7\textwidth}{X} 袍上方的中間要留一個領口，領口周圍的領邊要以手藝編織而成，好像鎧甲的領口，免得破裂。 \end{tabularx} \\ \\ \relax
28:33 & \begin{tabularx}{0.7\textwidth}{X} 袍子下襬，就是下襬的周圍要用藍色、紫色、朱紅色紗做石榴，周圍的石榴中間要有金鈴鐺： \end{tabularx} \\ \\ \relax
28:34 & \begin{tabularx}{0.7\textwidth}{X} 一個金鈴鐺一個石榴，一個金鈴鐺一個石榴，在袍子下襬的周圍。 \end{tabularx} \\ \\ \relax
28:35 & \begin{tabularx}{0.7\textwidth}{X} 亞倫供職的時候要穿這袍。他進入聖所到耶和華面前，以及出來的時候，袍上的鈴聲必被聽見，使他不至於死。 \end{tabularx} \\ \\ \relax
28:36 & \begin{tabularx}{0.7\textwidth}{X} 「你要用純金做一面牌，如同刻印章，在上面刻『歸耶和華為聖』。 \end{tabularx} \\ \\ \relax
28:37 & \begin{tabularx}{0.7\textwidth}{X} 要用藍色的帶子把牌繫在禮冠上，在禮冠的正前面。 \end{tabularx} \\ \\ \relax
28:38 & \begin{tabularx}{0.7\textwidth}{X} 這牌必在亞倫的額上，亞倫要擔當干犯聖物的罪孽；這聖物是以色列人在一切聖禮物上所分別為聖的。這牌要常在他的額上，使他們可以在耶和華面前蒙悅納。 \end{tabularx} \\ \\ \relax
28:39 & \begin{tabularx}{0.7\textwidth}{X} 要用細麻編織內袍，用細麻做禮冠，又以刺繡的手藝做腰帶。 \end{tabularx} \\ \\ \relax
28:40 & \begin{tabularx}{0.7\textwidth}{X} 「你要為亞倫的兒子做內袍、腰帶、頭巾，以示尊嚴和華美。 \end{tabularx} \\ \\ \relax
28:41 & \begin{tabularx}{0.7\textwidth}{X} 要把這些給你哥哥亞倫和他的兒子穿戴，又要膏他們，授予聖職，使他們分別為聖，作事奉我的祭司。 \end{tabularx} \\ \\ \relax
28:42 & \begin{tabularx}{0.7\textwidth}{X} 要用細麻布給他們做褲子來遮掩下體，從腰間直到大腿。 \end{tabularx} \\ \\ \relax
28:43 & \begin{tabularx}{0.7\textwidth}{X} 亞倫和他兒子進入會幕，或接近祭壇，在聖所供職的時候要穿上褲子，免得擔當罪孽而死。這要成為亞倫和他後裔永遠的定例。」 \end{tabularx} \\ \\ \relax
29:1 & \begin{tabularx}{0.7\textwidth}{X} 「這是你使他們分別為聖，作事奉我的祭司時要做的事：取一頭公牛犢，兩隻無殘疾的公綿羊， \end{tabularx} \\ \\ \relax
29:2 & \begin{tabularx}{0.7\textwidth}{X} 無酵餅、用油調和的無酵餅，和抹油的無酵薄餅；這些餅都要用細麥麵做成。 \end{tabularx} \\ \\ \relax
29:3 & \begin{tabularx}{0.7\textwidth}{X} 這些餅要裝在一個籃子裡，用籃子帶來，又把公牛和兩隻公綿羊牽來。 \end{tabularx} \\ \\ \relax
29:4 & \begin{tabularx}{0.7\textwidth}{X} 要帶亞倫和他兒子到會幕的門口，用水洗他們。 \end{tabularx} \\ \\ \relax
29:5 & \begin{tabularx}{0.7\textwidth}{X} 要拿服裝，給亞倫穿上內袍和以弗得的外袍，以及以弗得，又帶上胸袋，束上以弗得精緻的帶子。 \end{tabularx} \\ \\ \relax
29:6 & \begin{tabularx}{0.7\textwidth}{X} 要把禮冠戴在他頭上，將聖冕加在禮冠上， \end{tabularx} \\ \\ \relax
29:7 & \begin{tabularx}{0.7\textwidth}{X} 把膏油倒在他頭上膏他。 \end{tabularx} \\ \\ \relax
29:8 & \begin{tabularx}{0.7\textwidth}{X} 要帶他的兒子來，給他們穿上內袍。 \end{tabularx} \\ \\ \relax
29:9 & \begin{tabularx}{0.7\textwidth}{X} 要給亞倫和他的兒子束上腰帶，裹上頭巾，他們就憑永遠的定例得祭司的職分。又要授聖職給亞倫和他的兒子。 \end{tabularx} \\ \\ \relax
29:10 & \begin{tabularx}{0.7\textwidth}{X} 「你要把公牛牽到會幕前，亞倫和他的兒子要按手在公牛的頭上。 \end{tabularx} \\ \\ \relax
29:11 & \begin{tabularx}{0.7\textwidth}{X} 你要在耶和華面前，在會幕的門口宰這公牛。 \end{tabularx} \\ \\ \relax
29:12 & \begin{tabularx}{0.7\textwidth}{X} 要取些公牛的血，用指頭抹在祭壇的四個翹角上，把其餘的血全倒在壇的底座上。 \end{tabularx} \\ \\ \relax
29:13 & \begin{tabularx}{0.7\textwidth}{X} 要把所有包著內臟的脂肪、肝上的網油、兩個腎和腎上的脂肪，都燒在壇上。 \end{tabularx} \\ \\ \relax
29:14 & \begin{tabularx}{0.7\textwidth}{X} 只是公牛的肉、皮、糞都要在營外用火焚燒；這牛是贖罪祭。 \end{tabularx} \\ \\ \relax
29:15 & \begin{tabularx}{0.7\textwidth}{X} 「你要牽一隻公綿羊來，亞倫和他兒子要按手在這羊的頭上。 \end{tabularx} \\ \\ \relax
29:16 & \begin{tabularx}{0.7\textwidth}{X} 你要宰這羊，把血灑在祭壇的周圍。 \end{tabularx} \\ \\ \relax
29:17 & \begin{tabularx}{0.7\textwidth}{X} 再把羊切成肉塊，洗淨內臟和腿，連肉塊和頭放在一處。 \end{tabularx} \\ \\ \relax
29:18 & \begin{tabularx}{0.7\textwidth}{X} 要把全羊燒在壇上。這是獻給耶和華的燔祭，是獻給耶和華馨香的火祭。」 \end{tabularx} \\ \\ \relax
29:19 & \begin{tabularx}{0.7\textwidth}{X} 「你要把第二隻公綿羊牽來，亞倫和他兒子要按手在這羊的頭上。 \end{tabularx} \\ \\ \relax
29:20 & \begin{tabularx}{0.7\textwidth}{X} 你要宰這羊，取些血抹在亞倫的右耳垂和他兒子的右耳垂上，又抹在他們右手的大拇指和右腳的大腳趾上，然後把其餘的血灑在壇的周圍。 \end{tabularx} \\ \\ \relax
29:21 & \begin{tabularx}{0.7\textwidth}{X} 你要取些膏油和壇上的血，彈在亞倫和他的衣服上，以及他兒子和他們的衣服上；亞倫和他的衣服，他兒子和他們的衣服都成為聖了。 \end{tabularx} \\ \\ \relax
29:22 & \begin{tabularx}{0.7\textwidth}{X} 「你要取這羊的脂肪，肥尾巴、包著內臟的脂肪、肝上的網油、兩個腎、腎上的脂肪和右腿，這是聖職禮所獻的公綿羊； \end{tabularx} \\ \\ \relax
29:23 & \begin{tabularx}{0.7\textwidth}{X} 再從耶和華面前那裝無酵餅的籃子中取一個餅、一個油餅和一個薄餅， \end{tabularx} \\ \\ \relax
29:24 & \begin{tabularx}{0.7\textwidth}{X} 把它們都放在亞倫的手和他兒子的手上，在耶和華面前搖一搖，作為搖祭。 \end{tabularx} \\ \\ \relax
29:25 & \begin{tabularx}{0.7\textwidth}{X} 然後，你要從他們手中接過來，放在燔祭上，一起燒在壇上，作為耶和華面前馨香之氣；這是獻給耶和華的火祭。 \end{tabularx} \\ \\ \relax
29:26 & \begin{tabularx}{0.7\textwidth}{X} 「你要取亞倫聖職禮所獻公綿羊的胸，在耶和華面前搖一搖，作為搖祭；這份就是你的。 \end{tabularx} \\ \\ \relax
29:27 & \begin{tabularx}{0.7\textwidth}{X} 那搖祭的胸和舉祭的腿，就是聖職禮獻公綿羊時所搖的、所舉的，你要使它們分別為聖，是歸給亞倫和他兒子的。 \end{tabularx} \\ \\ \relax
29:28 & \begin{tabularx}{0.7\textwidth}{X} 這是亞倫和他子孫憑永遠的定例從以色列人中所應得的；因為這是舉祭，是從以色列人的平安祭中取出，作為獻給耶和華的舉祭。 \end{tabularx} \\ \\ \relax
29:29 & \begin{tabularx}{0.7\textwidth}{X} 「亞倫的聖衣要傳給他的子孫，使他們在受膏和承接聖職的時候穿上。 \end{tabularx} \\ \\ \relax
29:30 & \begin{tabularx}{0.7\textwidth}{X} 他的子孫接續他當祭司的，每逢進入會幕在聖所供職的時候，要穿這聖衣七天。 \end{tabularx} \\ \\ \relax
29:31 & \begin{tabularx}{0.7\textwidth}{X} 「你要拿聖職禮所獻的公綿羊，在聖處煮牠的肉。 \end{tabularx} \\ \\ \relax
29:32 & \begin{tabularx}{0.7\textwidth}{X} 亞倫和他兒子要在會幕的門口吃這羊的肉和籃子裡的餅。 \end{tabularx} \\ \\ \relax
29:33 & \begin{tabularx}{0.7\textwidth}{X} 他們要吃那些用來贖罪之物，好承接聖職，使他們分別為聖。外人不可吃，因為這是聖物。 \end{tabularx} \\ \\ \relax
29:34 & \begin{tabularx}{0.7\textwidth}{X} 那聖職禮所獻的肉或餅，若有剩餘留到早晨，就要把剩下的用火燒了，不可再吃，因為這是聖物。 \end{tabularx} \\ \\ \relax
29:35 & \begin{tabularx}{0.7\textwidth}{X} 「你要這樣照我一切所吩咐的，向亞倫和他兒子行授聖職禮七天。 \end{tabularx} \\ \\ \relax
29:36 & \begin{tabularx}{0.7\textwidth}{X} 為了贖罪，每天要獻一頭公牛為贖罪祭。你要為祭壇贖罪，使壇潔淨，並要用膏抹壇，使壇成為聖。 \end{tabularx} \\ \\ \relax
29:37 & \begin{tabularx}{0.7\textwidth}{X} 要為壇贖罪七天，使壇成為聖，壇就成為至聖。凡觸摸壇的都成為聖。」 \end{tabularx} \\ \\ \relax
29:38 & \begin{tabularx}{0.7\textwidth}{X} 「這是你要獻在壇上的：每天不可間斷地獻兩隻一歲的羔羊； \end{tabularx} \\ \\ \relax
29:39 & \begin{tabularx}{0.7\textwidth}{X} 早晨獻第一隻羔羊，黃昏獻第二隻羔羊。 \end{tabularx} \\ \\ \relax
29:40 & \begin{tabularx}{0.7\textwidth}{X} 獻第一隻羔羊時，要同時獻上十分之一伊法細麵，調和四分之一欣搗成的油，再獻四分之一欣酒作澆酒祭。 \end{tabularx} \\ \\ \relax
29:41 & \begin{tabularx}{0.7\textwidth}{X} 黃昏你獻第二隻羔羊，要照早晨的素祭和同獻的澆酒祭獻上，作為獻給耶和華馨香的火祭。 \end{tabularx} \\ \\ \relax
29:42 & \begin{tabularx}{0.7\textwidth}{X} 這要在耶和華面前，在會幕的門口，作為你們世世代代經常獻的燔祭。我要在那裡與你們相會，和你說話。 \end{tabularx} \\ \\ \relax
29:43 & \begin{tabularx}{0.7\textwidth}{X} 我要在那裡與以色列人相會，會幕就要因我的榮耀成為聖。 \end{tabularx} \\ \\ \relax
29:44 & \begin{tabularx}{0.7\textwidth}{X} 我要使會幕和祭壇分別為聖，也要使亞倫和他的兒子分別為聖，作事奉我的祭司。 \end{tabularx} \\ \\ \relax
29:45 & \begin{tabularx}{0.7\textwidth}{X} 我要住在以色列人中，作他們的神。 \end{tabularx} \\ \\ \relax
29:46 & \begin{tabularx}{0.7\textwidth}{X} 他們必知道我是耶和華－他們的神，是將他們從埃及地領出來的，為要住在他們中間。我是耶和華－他們的神 。」 \end{tabularx} \\ \\ \relax
30:1 & \begin{tabularx}{0.7\textwidth}{X} 「你要用金合歡木做一座燒香的壇， \end{tabularx} \\ \\ \relax
30:2 & \begin{tabularx}{0.7\textwidth}{X} 長一肘，寬一肘，這壇是正方形的，高二肘。壇的四個翹角與壇接連一塊。 \end{tabularx} \\ \\ \relax
30:3 & \begin{tabularx}{0.7\textwidth}{X} 要把壇的上面與壇的四圍，以及壇的四個翹角包上純金；又要在壇的四圍鑲上金邊。 \end{tabularx} \\ \\ \relax
30:4 & \begin{tabularx}{0.7\textwidth}{X} 要在壇的兩個對側，金邊下面做兩個金環，用來穿槓抬壇。 \end{tabularx} \\ \\ \relax
30:5 & \begin{tabularx}{0.7\textwidth}{X} 要用金合歡木做槓，包上金子。 \end{tabularx} \\ \\ \relax
30:6 & \begin{tabularx}{0.7\textwidth}{X} 要把壇放在法櫃前的幔子外，對著法櫃上的櫃蓋，就是我與你相會的地方。 \end{tabularx} \\ \\ \relax
30:7 & \begin{tabularx}{0.7\textwidth}{X} 亞倫要在壇上燒芬芳的香；每早晨整理燈的時候，他都要燒這香。 \end{tabularx} \\ \\ \relax
30:8 & \begin{tabularx}{0.7\textwidth}{X} 黃昏點燈的時候，亞倫也要燒這香。這是你們世世代代在耶和華面前常燒的香。 \end{tabularx} \\ \\ \relax
30:9 & \begin{tabularx}{0.7\textwidth}{X} 在這壇上不可燒別樣的香，不可獻燔祭、素祭，也不可獻澆酒祭。 \end{tabularx} \\ \\ \relax
30:10 & \begin{tabularx}{0.7\textwidth}{X} 亞倫每年一次要為壇的四個翹角贖罪。他每年一次要用贖罪祭的血為壇贖罪，作為世世代代的定例。這壇在耶和華面前是至聖的。」 \end{tabularx} \\ \\ \relax
30:11 & \begin{tabularx}{0.7\textwidth}{X} 耶和華吩咐摩西說： \end{tabularx} \\ \\ \relax
30:12 & \begin{tabularx}{0.7\textwidth}{X} 「你數點以色列人，計算人頭時，被數的每一個人要把他生命的贖價獻給耶和華，免得災殃在數點中臨到他們。 \end{tabularx} \\ \\ \relax
30:13 & \begin{tabularx}{0.7\textwidth}{X} 每一個被數的人要按照聖所的舍客勒，付半舍客勒，一舍客勒是二十季拉；這半舍客勒是獻給耶和華的禮物。 \end{tabularx} \\ \\ \relax
30:14 & \begin{tabularx}{0.7\textwidth}{X} 每一個被數的人，就是二十歲以上的，要將這禮物獻給耶和華。 \end{tabularx} \\ \\ \relax
30:15 & \begin{tabularx}{0.7\textwidth}{X} 富有的不必多付，貧窮的也不可少出，各人都要獻半舍客勒給耶和華，作你們生命的贖價。 \end{tabularx} \\ \\ \relax
30:16 & \begin{tabularx}{0.7\textwidth}{X} 你要向以色列人收這贖罪的銀子，用在會幕的事工。這要在耶和華面前為以色列人作紀念，作你們生命的贖價。」 \end{tabularx} \\ \\ \relax
30:17 & \begin{tabularx}{0.7\textwidth}{X} 耶和華吩咐摩西說： \end{tabularx} \\ \\ \relax
30:18 & \begin{tabularx}{0.7\textwidth}{X} 「你要用銅做洗濯盆和盆座，用來洗濯。要將盆放在會幕和祭壇的中間，盆裡盛水。 \end{tabularx} \\ \\ \relax
30:19 & \begin{tabularx}{0.7\textwidth}{X} 亞倫和他的兒子要用這盆洗手洗腳。 \end{tabularx} \\ \\ \relax
30:20 & \begin{tabularx}{0.7\textwidth}{X} 他們進會幕，或是走近壇前供職，獻火祭給耶和華的時候，必須用水洗濯，免得死亡； \end{tabularx} \\ \\ \relax
30:21 & \begin{tabularx}{0.7\textwidth}{X} 他們要洗手洗腳，免得死亡。這是亞倫和他的後裔世世代代永遠的定例。」 \end{tabularx} \\ \\ \relax
30:22 & \begin{tabularx}{0.7\textwidth}{X} 耶和華吩咐摩西說： \end{tabularx} \\ \\ \relax
30:23 & \begin{tabularx}{0.7\textwidth}{X} 「你要取上等的香料，就是五百舍客勒流質的沒藥、二百五十香肉桂、二百五十香菖蒲， \end{tabularx} \\ \\ \relax
30:24 & \begin{tabularx}{0.7\textwidth}{X} 和五百桂皮，都按照聖所的舍客勒；再取一欣橄欖油， \end{tabularx} \\ \\ \relax
30:25 & \begin{tabularx}{0.7\textwidth}{X} 以做香的方法調和製成聖膏油，它就成為聖膏油。 \end{tabularx} \\ \\ \relax
30:26 & \begin{tabularx}{0.7\textwidth}{X} 要用這膏油抹會幕和法櫃， \end{tabularx} \\ \\ \relax
30:27 & \begin{tabularx}{0.7\textwidth}{X} 供桌和供桌的一切器具，燈臺和燈臺的器具，以及香壇、 \end{tabularx} \\ \\ \relax
30:28 & \begin{tabularx}{0.7\textwidth}{X} 燔祭壇和壇的一切器具，洗濯盆和盆座。 \end{tabularx} \\ \\ \relax
30:29 & \begin{tabularx}{0.7\textwidth}{X} 你要使這些分別為聖，成為至聖；凡觸摸它們的都成為聖。 \end{tabularx} \\ \\ \relax
30:30 & \begin{tabularx}{0.7\textwidth}{X} 要膏亞倫和他的兒子，使他們分別為聖，作事奉我的祭司。 \end{tabularx} \\ \\ \relax
30:31 & \begin{tabularx}{0.7\textwidth}{X} 你要吩咐以色列人說：『你們要世世代代以這油為我的聖膏油。 \end{tabularx} \\ \\ \relax
30:32 & \begin{tabularx}{0.7\textwidth}{X} 不可把這油倒在別人身上，也不可用配製這膏油的方法製成同樣的膏油。這膏油是聖的，你們要以它為聖。 \end{tabularx} \\ \\ \relax
30:33 & \begin{tabularx}{0.7\textwidth}{X} 凡調和與此類似的膏油，或將它膏在別人身上的，這人要從百姓中剪除。』」 \end{tabularx} \\ \\ \relax
30:34 & \begin{tabularx}{0.7\textwidth}{X} 耶和華吩咐摩西說：「你要取香料，就是拿他弗、施喜列、喜利比拿，這些香料再加純乳香，每樣都要相同的分量。 \end{tabularx} \\ \\ \relax
30:35 & \begin{tabularx}{0.7\textwidth}{X} 你要用這些加上鹽，以配製香料的方法，製成純淨又神聖的香。 \end{tabularx} \\ \\ \relax
30:36 & \begin{tabularx}{0.7\textwidth}{X} 要取一點這香，搗成細的粉，放在會幕中的法櫃前，就是我和你相會的地方。你們要以這香為至聖。 \end{tabularx} \\ \\ \relax
30:37 & \begin{tabularx}{0.7\textwidth}{X} 你們不可用這配製的方法為自己做香；要以這香為聖，歸於耶和華。 \end{tabularx} \\ \\ \relax
30:38 & \begin{tabularx}{0.7\textwidth}{X} 為要聞香味而配製同樣的香的，這人要從百姓中剪除。」 \end{tabularx} \\ \\ \relax
31:1 & \begin{tabularx}{0.7\textwidth}{X} 耶和華吩咐摩西說： \end{tabularx} \\ \\ \relax
31:2 & \begin{tabularx}{0.7\textwidth}{X} 「你看，我已經題名召猶大支派中戶珥的孫子，烏利的兒子比撒列。 \end{tabularx} \\ \\ \relax
31:3 & \begin{tabularx}{0.7\textwidth}{X} 我以神的靈充滿他，使他有智慧，有聰明，有知識，能做各樣的工， \end{tabularx} \\ \\ \relax
31:4 & \begin{tabularx}{0.7\textwidth}{X} 能設計圖案，用金、銀、銅製造各物， \end{tabularx} \\ \\ \relax
31:5 & \begin{tabularx}{0.7\textwidth}{X} 又能雕刻鑲嵌用的寶石，雕刻木頭，做各樣的工。 \end{tabularx} \\ \\ \relax
31:6 & \begin{tabularx}{0.7\textwidth}{X} 看哪，我委派但支派中亞希撒抹的兒子亞何利亞伯與他同工。凡心裡有智慧的，我更要賜給他們智慧的心，能做我所吩咐的一切， \end{tabularx} \\ \\ \relax
31:7 & \begin{tabularx}{0.7\textwidth}{X} 就是會幕、法櫃和其上的櫃蓋、會幕中一切的器具、 \end{tabularx} \\ \\ \relax
31:8 & \begin{tabularx}{0.7\textwidth}{X} 供桌和供桌的器具、純金的燈臺和燈臺的一切器具、香壇、 \end{tabularx} \\ \\ \relax
31:9 & \begin{tabularx}{0.7\textwidth}{X} 燔祭壇和壇的一切器具、洗濯盆與盆座、 \end{tabularx} \\ \\ \relax
31:10 & \begin{tabularx}{0.7\textwidth}{X} 供祭司職分用的精緻禮服，亞倫祭司的聖衣和他兒子的衣服， \end{tabularx} \\ \\ \relax
31:11 & \begin{tabularx}{0.7\textwidth}{X} 以及膏油和聖所用的芬芳的香。他們都要照我所吩咐的一切去做。」 \end{tabularx} \\ \\ \relax
31:12 & \begin{tabularx}{0.7\textwidth}{X} 耶和華對摩西說： \end{tabularx} \\ \\ \relax
31:13 & \begin{tabularx}{0.7\textwidth}{X} 「你要吩咐以色列人說：『你們務要守我的安息日，因為這是你我之間世世代代的記號，叫你們知道我是耶和華，是使你們分別為聖的。 \end{tabularx} \\ \\ \relax
31:14 & \begin{tabularx}{0.7\textwidth}{X} 你們要守安息日，以它為聖日。凡干犯這日的，必被處死；凡在這日做工的，那人必從百姓中剪除。 \end{tabularx} \\ \\ \relax
31:15 & \begin{tabularx}{0.7\textwidth}{X} 六日要做工，但第七日是向耶和華守完全安息的安息聖日。凡在安息日做工的，必被處死。』 \end{tabularx} \\ \\ \relax
31:16 & \begin{tabularx}{0.7\textwidth}{X} 以色列人要守安息日，世世代代守安息日為永遠的約。 \end{tabularx} \\ \\ \relax
31:17 & \begin{tabularx}{0.7\textwidth}{X} 這是我和以色列人之間永遠的記號，因為六日之內耶和華造天地，第七日就安息舒暢。」 \end{tabularx} \\ \\ \relax
31:18 & \begin{tabularx}{0.7\textwidth}{X} 耶和華在西奈山和摩西說完了話，就把兩塊法版交給他，是神用指頭寫的石版。 \end{tabularx} \\ \\ \relax
32:1 & \begin{tabularx}{0.7\textwidth}{X} 百姓見摩西遲遲不下山，就聚集到亞倫那裡，對他說：「起來！為我們造神明，在我們前面引路，因為領我們出埃及地的那個摩西，我們不知道他遭遇了甚麼事。」 \end{tabularx} \\ \\ \relax
32:2 & \begin{tabularx}{0.7\textwidth}{X} 亞倫對他們說：「你們去摘下你們妻子、兒女耳上的金環，拿來給我。」 \end{tabularx} \\ \\ \relax
32:3 & \begin{tabularx}{0.7\textwidth}{X} 眾百姓就摘下他們耳上的金環，拿來給亞倫。 \end{tabularx} \\ \\ \relax
32:4 & \begin{tabularx}{0.7\textwidth}{X} 亞倫從他們手裡接過來，用模子塑造它，把它鑄成一頭牛犢。他們就說：「以色列啊，這是領你出埃及地的神明！」 \end{tabularx} \\ \\ \relax
32:5 & \begin{tabularx}{0.7\textwidth}{X} 亞倫看見，就在牛犢面前築壇。亞倫宣告說：「明日要向耶和華守節。」 \end{tabularx} \\ \\ \relax
32:6 & \begin{tabularx}{0.7\textwidth}{X} 次日清早，百姓起來獻燔祭和平安祭，就坐下吃喝，起來玩樂。 \end{tabularx} \\ \\ \relax
32:7 & \begin{tabularx}{0.7\textwidth}{X} 耶和華吩咐摩西：「下去吧，因為你從埃及領上來的百姓已經敗壞了。 \end{tabularx} \\ \\ \relax
32:8 & \begin{tabularx}{0.7\textwidth}{X} 他們這麼快偏離了我所吩咐的道，為自己鑄了一頭牛犢，向它跪拜，向它獻祭，說：『以色列啊，這就是領你出埃及地的神明。』」 \end{tabularx} \\ \\ \relax
32:9 & \begin{tabularx}{0.7\textwidth}{X} 耶和華對摩西說：「我看這百姓，看哪，他們真是硬著頸項的百姓。 \end{tabularx} \\ \\ \relax
32:10 & \begin{tabularx}{0.7\textwidth}{X} 現在，你且由著我，我要向他們發烈怒，滅絕他們，但我要使你成為大國。」 \end{tabularx} \\ \\ \relax
32:11 & \begin{tabularx}{0.7\textwidth}{X} 摩西就懇求耶和華－他的神，說：「耶和華啊，你為甚麼向你的百姓發烈怒呢？這百姓是你用大能大力的手從埃及地領出來的！ \end{tabularx} \\ \\ \relax
32:12 & \begin{tabularx}{0.7\textwidth}{X} 為甚麼讓埃及人說：『他領他們出去，是要降災禍給他們，在山中把他們殺了，將他們從地上除滅』呢？求你回心轉意，不發你的烈怒，不降災禍給你的百姓。 \end{tabularx} \\ \\ \relax
32:13 & \begin{tabularx}{0.7\textwidth}{X} 求你記念你的僕人亞伯拉罕、以撒、以色列。你曾向他們指著自己起誓說：『我必使你們的後裔像天上的星那樣多，並且我要將所應許的這全地賜給你們的後裔，讓他們永遠承受為業。』」 \end{tabularx} \\ \\ \relax
32:14 & \begin{tabularx}{0.7\textwidth}{X} 於是耶和華改變心意，不把所說的災禍降給他的百姓。 \end{tabularx} \\ \\ \relax
32:15 & \begin{tabularx}{0.7\textwidth}{X} 摩西轉身下山，手裡拿著兩塊法版。這版的兩面都寫著字，正面背面都有字。 \end{tabularx} \\ \\ \relax
32:16 & \begin{tabularx}{0.7\textwidth}{X} 版是神的工作，字是神寫的字，刻在版上。 \end{tabularx} \\ \\ \relax
32:17 & \begin{tabularx}{0.7\textwidth}{X} 約書亞一聽見百姓呼喊的聲音，就對摩西說：「在營裡有戰爭的聲音。」 \end{tabularx} \\ \\ \relax
32:18 & \begin{tabularx}{0.7\textwidth}{X} 摩西說：「這不是打勝仗的聲音，也不是打敗仗的聲音，我聽見的是歌唱的聲音。」 \end{tabularx} \\ \\ \relax
32:19 & \begin{tabularx}{0.7\textwidth}{X} 摩西走近營前，看見牛犢，又看見人在跳舞，就發烈怒，把兩塊版從手中扔到山下摔碎了。 \end{tabularx} \\ \\ \relax
32:20 & \begin{tabularx}{0.7\textwidth}{X} 他將他們所鑄的牛犢用火焚燒，磨得粉碎，撒在水面上，叫以色列人喝。 \end{tabularx} \\ \\ \relax
32:21 & \begin{tabularx}{0.7\textwidth}{X} 摩西對亞倫說：「這百姓向你做了甚麼呢？你竟使他們陷入大罪中！」 \end{tabularx} \\ \\ \relax
32:22 & \begin{tabularx}{0.7\textwidth}{X} 亞倫說：「求我主不要發烈怒。你知道這百姓，他們是向惡的。 \end{tabularx} \\ \\ \relax
32:23 & \begin{tabularx}{0.7\textwidth}{X} 他們對我說：『你為我們造神明，在我們前面引路，因為領我們出埃及地的那個摩西，我們不知道他遭遇了甚麼事。』 \end{tabularx} \\ \\ \relax
32:24 & \begin{tabularx}{0.7\textwidth}{X} 我對他們說：『凡有金環的可以摘下來』，他們就給了我。我把金環扔在火中，這牛犢就出來了。」 \end{tabularx} \\ \\ \relax
32:25 & \begin{tabularx}{0.7\textwidth}{X} 摩西見百姓放肆，因亞倫縱容他們，使這事成了敵人的笑柄， \end{tabularx} \\ \\ \relax
32:26 & \begin{tabularx}{0.7\textwidth}{X} 就站在營門前，說：「凡屬耶和華的人，都到我這裡來！」於是利未人都聚集到他那裡。 \end{tabularx} \\ \\ \relax
32:27 & \begin{tabularx}{0.7\textwidth}{X} 他對他們說：「耶和華－以色列的神這樣說：『你們各人把刀佩在腰間，從這門到那門，來回走遍全營，各人要殺自己的弟兄、鄰舍和親人。』」 \end{tabularx} \\ \\ \relax
32:28 & \begin{tabularx}{0.7\textwidth}{X} 利未人遵照摩西的話做了。那一天百姓中倒下的約有三千人。 \end{tabularx} \\ \\ \relax
32:29 & \begin{tabularx}{0.7\textwidth}{X} 摩西說：「今天你們要奉獻自己來事奉耶和華，因為各人犧牲自己的兒子和弟兄，使耶和華今天賜福給你們。」 \end{tabularx} \\ \\ \relax
32:30 & \begin{tabularx}{0.7\textwidth}{X} 第二天，摩西對百姓說：「你們犯了大罪。我如今要上耶和華那裡去，或許可以為你們贖罪。」 \end{tabularx} \\ \\ \relax
32:31 & \begin{tabularx}{0.7\textwidth}{X} 摩西回到耶和華那裡，說：「唉！這百姓犯了大罪，為自己造了金的神明。 \end{tabularx} \\ \\ \relax
32:32 & \begin{tabularx}{0.7\textwidth}{X} 現在，求你赦免他們的罪；不然，就把我從你所寫的冊上除名。」 \end{tabularx} \\ \\ \relax
32:33 & \begin{tabularx}{0.7\textwidth}{X} 耶和華對摩西說：「誰得罪我，我就把他從我的冊上除去。 \end{tabularx} \\ \\ \relax
32:34 & \begin{tabularx}{0.7\textwidth}{X} 現在你去，領這百姓往我所告訴你的地方去，看哪，我的使者必在你的前面引路。到了該懲罰的時候，我必懲罰他們的罪。」 \end{tabularx} \\ \\ \relax
32:35 & \begin{tabularx}{0.7\textwidth}{X} 耶和華降災與百姓，因為他們和亞倫一起造了牛犢。 \end{tabularx} \\ \\ \relax
33:1 & \begin{tabularx}{0.7\textwidth}{X} 耶和華吩咐摩西說：「去，離開這裡，你和你從埃及地領出來的百姓要上到我起誓應許給亞伯拉罕、以撒和雅各之地去；我曾對他們說：『我要將這地賜給你的後裔』。 \end{tabularx} \\ \\ \relax
33:2 & \begin{tabularx}{0.7\textwidth}{X} 我要差遣使者在你前面，把迦南人、亞摩利人、赫人、比利洗人、希未人、耶布斯人趕出 \end{tabularx} \\ \\ \relax
33:3 & \begin{tabularx}{0.7\textwidth}{X} 那流奶與蜜之地。但我不與你們上去，因為你們是硬著頸項的百姓，免得我在路上把你們滅絕。」 \end{tabularx} \\ \\ \relax
33:4 & \begin{tabularx}{0.7\textwidth}{X} 百姓一聽見這壞的信息，他們就悲哀，沒有人佩戴首飾。 \end{tabularx} \\ \\ \relax
33:5 & \begin{tabularx}{0.7\textwidth}{X} 耶和華對摩西說：「你對以色列人說：『你們是硬著頸項的百姓，我若在你們中間一起上去，只一瞬間，就必把你們滅絕。現在把你們身上的首飾摘下來，我好知道該怎樣處置你們。』」 \end{tabularx} \\ \\ \relax
33:6 & \begin{tabularx}{0.7\textwidth}{X} 以色列人離開何烈山以後，就把身上的首飾摘下來。 \end{tabularx} \\ \\ \relax
33:7 & \begin{tabularx}{0.7\textwidth}{X} 摩西拿一個帳棚支搭在營外，離營有一段距離，他稱這帳棚為會幕。凡求問耶和華的，就到營外的會幕那裡去。 \end{tabularx} \\ \\ \relax
33:8 & \begin{tabularx}{0.7\textwidth}{X} 當摩西出營到會幕去的時候，百姓就都起來，各人站在自己帳棚的門口，望著摩西，直到他進了會幕。 \end{tabularx} \\ \\ \relax
33:9 & \begin{tabularx}{0.7\textwidth}{X} 摩西進會幕的時候，雲柱就降下來，停在會幕的門前，耶和華就與摩西說話。 \end{tabularx} \\ \\ \relax
33:10 & \begin{tabularx}{0.7\textwidth}{X} 眾百姓看見雲柱停在會幕的門前，就都起來，各人在自己帳棚的門口下拜。 \end{tabularx} \\ \\ \relax
33:11 & \begin{tabularx}{0.7\textwidth}{X} 耶和華與摩西面對面說話，好像人與朋友說話。摩西回到營裡去，他的年輕助手嫩的兒子約書亞卻沒有離開會幕。 \end{tabularx} \\ \\ \relax
33:12 & \begin{tabularx}{0.7\textwidth}{X} 摩西對耶和華說：「看，你曾對我說：『將這百姓領上去』；卻沒有讓我知道你要差派誰與我同去。你還說：『我按你的名認識你，你也在我眼前蒙了恩。』 \end{tabularx} \\ \\ \relax
33:13 & \begin{tabularx}{0.7\textwidth}{X} 我如今若在你眼前蒙恩，求你將你的道指示我，使我可以認識你，並在你眼前蒙恩。求你顧念這國是你的子民。」 \end{tabularx} \\ \\ \relax
33:14 & \begin{tabularx}{0.7\textwidth}{X} 耶和華說：「我必親自去，讓你安心。」 \end{tabularx} \\ \\ \relax
33:15 & \begin{tabularx}{0.7\textwidth}{X} 摩西說：「你若不親自去，就不要把我們從這裡領上去。 \end{tabularx} \\ \\ \relax
33:16 & \begin{tabularx}{0.7\textwidth}{X} 現在，人如何得知我和你的百姓在你眼前蒙恩呢？豈不是因為你與我們同去，使我和你的百姓與地面上的萬民有分別嗎？」 \end{tabularx} \\ \\ \relax
33:17 & \begin{tabularx}{0.7\textwidth}{X} 耶和華對摩西說：「你所說的這件事，我也會去做，因為你在我眼前蒙了恩，並且我按你的名認識你。」 \end{tabularx} \\ \\ \relax
33:18 & \begin{tabularx}{0.7\textwidth}{X} 摩西說：「求你顯出你的榮耀給我看。」 \end{tabularx} \\ \\ \relax
33:19 & \begin{tabularx}{0.7\textwidth}{X} 耶和華說：「我要顯示我一切的美善，在你面前經過，並要在你面前宣告耶和華的名。我要恩待誰就恩待誰，要憐憫誰就憐憫誰。」 \end{tabularx} \\ \\ \relax
33:20 & \begin{tabularx}{0.7\textwidth}{X} 他又說：「只是你不能看見我的面，因為沒有人看見我還可以存活。」 \end{tabularx} \\ \\ \relax
33:21 & \begin{tabularx}{0.7\textwidth}{X} 耶和華說：「看哪，靠近我這裡有個地方，你可以站在這磐石上。 \end{tabularx} \\ \\ \relax
33:22 & \begin{tabularx}{0.7\textwidth}{X} 當我的榮耀經過的時候，我必將你放在磐石縫裡，用我的手掌遮掩你，等我過去， \end{tabularx} \\ \\ \relax
33:23 & \begin{tabularx}{0.7\textwidth}{X} 然後我要將我的手掌收回，你就可以看見我的背，卻看不到我的面。」 \end{tabularx} \\ \\ \relax
34:1 & \begin{tabularx}{0.7\textwidth}{X} 耶和華對摩西說：「你要鑿出兩塊石版，和先前的一樣；我要把你摔碎的那版上先前所寫的字，寫在這版上。 \end{tabularx} \\ \\ \relax
34:2 & \begin{tabularx}{0.7\textwidth}{X} 明日早晨，你要預備好了，上西奈山，在山頂那裡站在我面前。 \end{tabularx} \\ \\ \relax
34:3 & \begin{tabularx}{0.7\textwidth}{X} 誰也不可和你上去，整座山都不可見到人，也不可有羊群牛群在山下吃草。」 \end{tabularx} \\ \\ \relax
34:4 & \begin{tabularx}{0.7\textwidth}{X} 摩西就鑿出兩塊石版，和先前的一樣。他清晨起來，遵照耶和華吩咐他的，上西奈山去，手裡拿著兩塊石版。 \end{tabularx} \\ \\ \relax
34:5 & \begin{tabularx}{0.7\textwidth}{X} 耶和華在雲中降臨，與摩西一同站在那裡，宣告耶和華的名。 \end{tabularx} \\ \\ \relax
34:6 & \begin{tabularx}{0.7\textwidth}{X} 耶和華在他面前經過，宣告：「耶和華，耶和華，有憐憫，有恩惠的神，不輕易發怒，且有豐盛的慈愛和信實， \end{tabularx} \\ \\ \relax
34:7 & \begin{tabularx}{0.7\textwidth}{X} 為千代的人存留慈愛，赦免罪孽、過犯和罪惡，萬不以有罪的為無罪，必懲罰人的罪，自父及子，直到三、四代。」 \end{tabularx} \\ \\ \relax
34:8 & \begin{tabularx}{0.7\textwidth}{X} 摩西急忙俯伏在地敬拜， \end{tabularx} \\ \\ \relax
34:9 & \begin{tabularx}{0.7\textwidth}{X} 說：「主啊，我若在你眼前蒙恩，求主在我們中間同行。雖然這是硬著頸項的百姓，求你赦免我們的罪孽和罪惡，接納我們為你的產業。」 \end{tabularx} \\ \\ \relax
34:10 & \begin{tabularx}{0.7\textwidth}{X} 耶和華說：「看哪，我要立約，要在你眾百姓面前行奇妙的事，是在全地萬國中未曾做過的。你周圍的萬民要看見我藉著你所行，耶和華可畏懼的作為。 \end{tabularx} \\ \\ \relax
34:11 & \begin{tabularx}{0.7\textwidth}{X} 「我今天所吩咐你的，你要謹守。看哪，我要從你面前趕出亞摩利人、迦南人、赫人、比利洗人、希未人、耶布斯人。 \end{tabularx} \\ \\ \relax
34:12 & \begin{tabularx}{0.7\textwidth}{X} 你要謹慎，不可與你所要去那地的居民立約，免得他們成為你中間的圈套。 \end{tabularx} \\ \\ \relax
34:13 & \begin{tabularx}{0.7\textwidth}{X} 你要拆毀他們的祭壇，打碎他們的柱像，砍斷他們的亞舍拉。 \end{tabularx} \\ \\ \relax
34:14 & \begin{tabularx}{0.7\textwidth}{X} 不可敬拜別神，因為耶和華是忌邪的神，他的名是忌邪者。 \end{tabularx} \\ \\ \relax
34:15 & \begin{tabularx}{0.7\textwidth}{X} 你不可與那地的居民立約，因為他們隨從自己的神明行淫；祭他們神明的時候，有人邀請你參加，你就會吃他的祭物。 \end{tabularx} \\ \\ \relax
34:16 & \begin{tabularx}{0.7\textwidth}{X} 你為你兒子娶他們的女兒為妻，他們的女兒因著隨從她們的神明行淫，就引誘你的兒子也隨從她們的神明行淫。 \end{tabularx} \\ \\ \relax
34:17 & \begin{tabularx}{0.7\textwidth}{X} 「不可為自己鑄造神像。 \end{tabularx} \\ \\ \relax
34:18 & \begin{tabularx}{0.7\textwidth}{X} 「你要守除酵節，照我所吩咐你的，在亞筆月內所定的日期吃無酵餅七天，因為你是在亞筆月內出了埃及。 \end{tabularx} \\ \\ \relax
34:19 & \begin{tabularx}{0.7\textwidth}{X} 「凡頭生的都是我的；無論是牛是羊，一切頭生的公的牲畜都要分別出來。 \end{tabularx} \\ \\ \relax
34:20 & \begin{tabularx}{0.7\textwidth}{X} 頭生的驢可以用羔羊代贖。若不贖牠，就要打斷牠的頸項。凡頭生的兒子都要贖出來。沒有人可以空手來朝見我。 \end{tabularx} \\ \\ \relax
34:21 & \begin{tabularx}{0.7\textwidth}{X} 「六日你要做工，第七日要安息，即使在耕種或收割的時候也要安息。 \end{tabularx} \\ \\ \relax
34:22 & \begin{tabularx}{0.7\textwidth}{X} 在收割初熟麥子的時候要守七七節，又要在年底守收藏節。 \end{tabularx} \\ \\ \relax
34:23 & \begin{tabularx}{0.7\textwidth}{X} 你所有的男丁要一年三次朝見主耶和華－以色列的神。 \end{tabularx} \\ \\ \relax
34:24 & \begin{tabularx}{0.7\textwidth}{X} 我要從你面前趕走列國，擴張你的疆界。你一年三次上去朝見耶和華－你神的時候，必沒有人貪圖你的地。 \end{tabularx} \\ \\ \relax
34:25 & \begin{tabularx}{0.7\textwidth}{X} 「不可將我祭牲的血和有酵之物一同獻上。逾越節的祭牲也不可留到早晨。 \end{tabularx} \\ \\ \relax
34:26 & \begin{tabularx}{0.7\textwidth}{X} 土地裡上好的初熟之物要奉到耶和華－你神的殿。不可用母山羊的奶來煮牠的小山羊。」 \end{tabularx} \\ \\ \relax
34:27 & \begin{tabularx}{0.7\textwidth}{X} 耶和華對摩西說：「你要將這些話寫上，因為我按這話與你和以色列人立約。」 \end{tabularx} \\ \\ \relax
34:28 & \begin{tabularx}{0.7\textwidth}{X} 摩西在耶和華那裡四十晝夜，不吃飯不喝水。他把這約的話，那十條誡命，寫在版上。 \end{tabularx} \\ \\ \relax
34:29 & \begin{tabularx}{0.7\textwidth}{X} 摩西下西奈山。摩西從山上下來的時候，手裡拿著兩塊法版。摩西不知道自己臉上的皮膚因耶和華和他說話而發光。 \end{tabularx} \\ \\ \relax
34:30 & \begin{tabularx}{0.7\textwidth}{X} 亞倫和以色列眾人看見摩西，看哪，他臉上的皮膚發光，他們就怕靠近他。 \end{tabularx} \\ \\ \relax
34:31 & \begin{tabularx}{0.7\textwidth}{X} 摩西叫他們來，亞倫和會眾的官長回到他那裡，摩西就跟他們說話。 \end{tabularx} \\ \\ \relax
34:32 & \begin{tabularx}{0.7\textwidth}{X} 隨後以色列眾人都近前來，他就把耶和華在西奈山與他所說的一切話都吩咐他們。 \end{tabularx} \\ \\ \relax
34:33 & \begin{tabularx}{0.7\textwidth}{X} 摩西跟他們說完了話，就用面紗蒙上臉。 \end{tabularx} \\ \\ \relax
34:34 & \begin{tabularx}{0.7\textwidth}{X} 但摩西進到耶和華面前與他說話的時候，就把面紗揭下，直到出來。摩西出來，將所吩咐他的話告訴以色列人。 \end{tabularx} \\ \\ \relax
34:35 & \begin{tabularx}{0.7\textwidth}{X} 以色列人看見摩西的臉，他臉上的皮膚發光。摩西就用面紗蒙上臉，直到他進去與耶和華說話才揭下。 \end{tabularx} \\ \\ \relax
35:1 & \begin{tabularx}{0.7\textwidth}{X} 摩西召集以色列全會眾，對他們說：「這是耶和華吩咐你們遵行的事： \end{tabularx} \\ \\ \relax
35:2 & \begin{tabularx}{0.7\textwidth}{X} 六日要做工，第七日你們要奉為向耶和華守完全安息的安息聖日。凡在這日做工的，要被處死。 \end{tabularx} \\ \\ \relax
35:3 & \begin{tabularx}{0.7\textwidth}{X} 在安息日這一天，不可在你們一切的住處生火。」 \end{tabularx} \\ \\ \relax
35:4 & \begin{tabularx}{0.7\textwidth}{X} 摩西對以色列全會眾說：「這是耶和華所吩咐的話，說： \end{tabularx} \\ \\ \relax
35:5 & \begin{tabularx}{0.7\textwidth}{X} 要從你們當中拿禮物獻給耶和華；凡甘心樂意的，可以把耶和華的禮物拿來，就是金、銀、銅， \end{tabularx} \\ \\ \relax
35:6 & \begin{tabularx}{0.7\textwidth}{X} 藍色、紫色、朱紅色紗，細麻，山羊毛， \end{tabularx} \\ \\ \relax
35:7 & \begin{tabularx}{0.7\textwidth}{X} 染紅的公羊皮，精美的皮料，金合歡木， \end{tabularx} \\ \\ \relax
35:8 & \begin{tabularx}{0.7\textwidth}{X} 點燈的油，做膏油的香料、做香的香料， \end{tabularx} \\ \\ \relax
35:9 & \begin{tabularx}{0.7\textwidth}{X} 紅瑪瑙與寶石，可以鑲嵌在以弗得和胸袋上。」 \end{tabularx} \\ \\ \relax
35:10 & \begin{tabularx}{0.7\textwidth}{X} 「你們當中凡心裡有智慧的都要來，製造一切耶和華所吩咐的， \end{tabularx} \\ \\ \relax
35:11 & \begin{tabularx}{0.7\textwidth}{X} 就是帳幕、帳幕的罩棚、帳幕的蓋、鉤子、豎板、橫木、柱子和帶卯眼的座， \end{tabularx} \\ \\ \relax
35:12 & \begin{tabularx}{0.7\textwidth}{X} 櫃子、櫃子的槓、櫃蓋和遮掩的幔子， \end{tabularx} \\ \\ \relax
35:13 & \begin{tabularx}{0.7\textwidth}{X} 供桌、供桌的槓、供桌一切的器具和供餅， \end{tabularx} \\ \\ \relax
35:14 & \begin{tabularx}{0.7\textwidth}{X} 燈臺、燈臺的器具、燈和點燈的油， \end{tabularx} \\ \\ \relax
35:15 & \begin{tabularx}{0.7\textwidth}{X} 香壇、壇的槓、膏油和芬芳的香，帳幕門口的門簾， \end{tabularx} \\ \\ \relax
35:16 & \begin{tabularx}{0.7\textwidth}{X} 燔祭壇、壇的銅網、壇的槓和壇的一切器具，洗濯盆和盆座， \end{tabularx} \\ \\ \relax
35:17 & \begin{tabularx}{0.7\textwidth}{X} 院子的帷幔、柱子、帶卯眼的座和院子的門簾， \end{tabularx} \\ \\ \relax
35:18 & \begin{tabularx}{0.7\textwidth}{X} 帳幕的橛子、院子的橛子和繩子， \end{tabularx} \\ \\ \relax
35:19 & \begin{tabularx}{0.7\textwidth}{X} 以及聖所事奉用的精緻禮服，亞倫祭司的聖衣和他兒子的衣服，供祭司職分用。」 \end{tabularx} \\ \\ \relax
35:20 & \begin{tabularx}{0.7\textwidth}{X} 以色列全會眾從摩西的面前出去。 \end{tabularx} \\ \\ \relax
35:21 & \begin{tabularx}{0.7\textwidth}{X} 凡心受感動，靈被驅策的，都帶耶和華的禮物來，為要造會幕和其中一切的器具，以及縫製聖衣。 \end{tabularx} \\ \\ \relax
35:22 & \begin{tabularx}{0.7\textwidth}{X} 凡甘心樂意的，連男帶女都來了，各將金器，就是胸針、耳環、打印的戒指，和項鏈帶來，搖著金器的搖祭獻給耶和華。 \end{tabularx} \\ \\ \relax
35:23 & \begin{tabularx}{0.7\textwidth}{X} 凡有藍色、紫色、朱紅色紗、細麻、山羊毛、染紅的公羊皮、精美皮料的，都拿了來； \end{tabularx} \\ \\ \relax
35:24 & \begin{tabularx}{0.7\textwidth}{X} 凡願意獻銀和銅作禮物的，都拿禮物來獻給耶和華；凡有金合歡木可做各種用途的也都拿了來。 \end{tabularx} \\ \\ \relax
35:25 & \begin{tabularx}{0.7\textwidth}{X} 凡心中有智慧，可以親手紡織的婦女，也把所紡的藍色、紫色、朱紅色紗，和細麻都拿了來。 \end{tabularx} \\ \\ \relax
35:26 & \begin{tabularx}{0.7\textwidth}{X} 凡有智慧，心裡受感動的婦女都來紡山羊毛。 \end{tabularx} \\ \\ \relax
35:27 & \begin{tabularx}{0.7\textwidth}{X} 眾官長把紅瑪瑙和寶石，可以鑲嵌在以弗得與胸袋上的，都拿了來， \end{tabularx} \\ \\ \relax
35:28 & \begin{tabularx}{0.7\textwidth}{X} 又拿做香，做膏油，和點燈所需的香料和油來。 \end{tabularx} \\ \\ \relax
35:29 & \begin{tabularx}{0.7\textwidth}{X} 以色列人，無論男女，凡心裡受感動的，都帶甘心祭來獻給耶和華，為要做耶和華藉摩西所吩咐的一切工。 \end{tabularx} \\ \\ \relax
35:30 & \begin{tabularx}{0.7\textwidth}{X} 摩西對以色列人說：「看，猶大支派中戶珥的孫子，烏利的兒子比撒列，耶和華已經題名召他， \end{tabularx} \\ \\ \relax
35:31 & \begin{tabularx}{0.7\textwidth}{X} 又以神的靈充滿他，使他有智慧、聰明、知識，能做各樣的工， \end{tabularx} \\ \\ \relax
35:32 & \begin{tabularx}{0.7\textwidth}{X} 能設計圖案，用金、銀、銅製造各物， \end{tabularx} \\ \\ \relax
35:33 & \begin{tabularx}{0.7\textwidth}{X} 又能雕刻鑲嵌用的寶石，雕刻木頭，做各樣精巧的工。 \end{tabularx} \\ \\ \relax
35:34 & \begin{tabularx}{0.7\textwidth}{X} 耶和華又賜給他和但支派中，亞希撒抹的兒子亞何利亞伯能教導人的心。 \end{tabularx} \\ \\ \relax
35:35 & \begin{tabularx}{0.7\textwidth}{X} 耶和華使他們的心滿有智慧，能做各樣的工，無論是雕刻的工，圖案設計的工，用藍色、紫色、朱紅色紗，和細麻作刺繡的工，以及編織的工，他們都能勝任，也能設計圖案。」 \end{tabularx} \\ \\ \relax
36:1 & \begin{tabularx}{0.7\textwidth}{X} 比撒列和亞何利亞伯，以及一切心裡有智慧，蒙耶和華賜智慧和聰明，懂得做聖所各樣用途之工的人，都照耶和華所吩咐的去做。 \end{tabularx} \\ \\ \relax
36:2 & \begin{tabularx}{0.7\textwidth}{X} 摩西把比撒列和亞何利亞伯，以及那些蒙耶和華賜他心裡有智慧，心受感動願意前來做工的人都召來。 \end{tabularx} \\ \\ \relax
36:3 & \begin{tabularx}{0.7\textwidth}{X} 這些人就從摩西收了以色列人為建造聖所，以及聖所各用途之工而奉獻的禮物。每天早晨，百姓繼續把甘心祭拿來。 \end{tabularx} \\ \\ \relax
36:4 & \begin{tabularx}{0.7\textwidth}{X} 凡有智慧能做聖所一切工的人，都各自離開他們原本的工作前來， \end{tabularx} \\ \\ \relax
36:5 & \begin{tabularx}{0.7\textwidth}{X} 對摩西說：「百姓送來的禮物很多，已經超過耶和華吩咐建造之工所需要的了。」 \end{tabularx} \\ \\ \relax
36:6 & \begin{tabularx}{0.7\textwidth}{X} 摩西吩咐，他們就在營中傳令說：「無論男女，不必再為聖所的禮物做任何的工。」這樣才使百姓停止，不再拿禮物來， \end{tabularx} \\ \\ \relax
36:7 & \begin{tabularx}{0.7\textwidth}{X} 他們所有的材料已經足夠整個工程之用，而且有餘。 \end{tabularx} \\ \\ \relax
36:8 & \begin{tabularx}{0.7\textwidth}{X} 做工的人當中，凡心裡有智慧的，用十幅幔子做帳幕，幔子是用搓的細麻和藍色、紫色、朱紅色紗織成的，並且以刺繡的手藝繡上基路伯。 \end{tabularx} \\ \\ \relax
36:9 & \begin{tabularx}{0.7\textwidth}{X} 每幅幔子長二十八肘，每幅幔子寬四肘，全部的幔子都是一樣的尺寸。 \end{tabularx} \\ \\ \relax
36:10 & \begin{tabularx}{0.7\textwidth}{X} 他使這五幅幔子彼此相連，又使那五幅幔子彼此相連。 \end{tabularx} \\ \\ \relax
36:11 & \begin{tabularx}{0.7\textwidth}{X} 他在這一組相連幔子的末幅邊上縫了藍色的鈕環；在另一組相連幔子的末幅邊上也照樣做。 \end{tabularx} \\ \\ \relax
36:12 & \begin{tabularx}{0.7\textwidth}{X} 他在這幅幔子上縫五十個鈕環，在另一組相連幔子的末幅上也縫五十個鈕環，環環相對。 \end{tabularx} \\ \\ \relax
36:13 & \begin{tabularx}{0.7\textwidth}{X} 他又做了五十個金鉤，用鉤子使幔子彼此相連，成為一個帳幕。 \end{tabularx} \\ \\ \relax
36:14 & \begin{tabularx}{0.7\textwidth}{X} 他用山羊毛織十一幅幔子，作為帳幕上的罩棚。 \end{tabularx} \\ \\ \relax
36:15 & \begin{tabularx}{0.7\textwidth}{X} 每幅幔子長三十肘，每幅幔子寬四肘；十一幅幔子都是一樣的尺寸。 \end{tabularx} \\ \\ \relax
36:16 & \begin{tabularx}{0.7\textwidth}{X} 他把五幅幔子連成一幅，又把六幅幔子連成一幅。 \end{tabularx} \\ \\ \relax
36:17 & \begin{tabularx}{0.7\textwidth}{X} 他在這一組相連幔子的末幅邊上縫了五十個鈕環；在另一組相連幔子的末幅邊上也縫了五十個鈕環。 \end{tabularx} \\ \\ \relax
36:18 & \begin{tabularx}{0.7\textwidth}{X} 他又做五十個銅鉤，使罩棚相連成為一個。 \end{tabularx} \\ \\ \relax
36:19 & \begin{tabularx}{0.7\textwidth}{X} 他用染紅的公羊皮做罩棚的蓋，再用精美皮料做外層的蓋。 \end{tabularx} \\ \\ \relax
36:20 & \begin{tabularx}{0.7\textwidth}{X} 他用金合歡木做豎立帳幕的木板， \end{tabularx} \\ \\ \relax
36:21 & \begin{tabularx}{0.7\textwidth}{X} 木板長十肘，每塊板寬一肘半， \end{tabularx} \\ \\ \relax
36:22 & \begin{tabularx}{0.7\textwidth}{X} 每塊板有兩個榫頭可以彼此銜接。帳幕一切的板都是這樣做。 \end{tabularx} \\ \\ \relax
36:23 & \begin{tabularx}{0.7\textwidth}{X} 他做帳幕的木板：南面，就是面向南方的那一邊，做二十塊板， \end{tabularx} \\ \\ \relax
36:24 & \begin{tabularx}{0.7\textwidth}{X} 在這二十塊板底下做了四十個帶卯眼的銀座：兩個卯眼接連這塊板上的兩個榫頭，另外兩個卯眼接連那塊板上的兩個榫頭。 \end{tabularx} \\ \\ \relax
36:25 & \begin{tabularx}{0.7\textwidth}{X} 他在帳幕的第二邊，就是北面，也做二十塊板， \end{tabularx} \\ \\ \relax
36:26 & \begin{tabularx}{0.7\textwidth}{X} 和四十個帶卯眼的銀座；這塊板底下有兩個卯眼，那塊板底下也有兩個卯眼。 \end{tabularx} \\ \\ \relax
36:27 & \begin{tabularx}{0.7\textwidth}{X} 他在帳幕的後面，就是西面，做六塊板， \end{tabularx} \\ \\ \relax
36:28 & \begin{tabularx}{0.7\textwidth}{X} 在帳幕後面的角落做兩塊板。 \end{tabularx} \\ \\ \relax
36:29 & \begin{tabularx}{0.7\textwidth}{X} 下端的板是成雙的，上端連在一起，直到頂端的第一個環子；兩塊板都是這樣，做成兩個角落。 \end{tabularx} \\ \\ \relax
36:30 & \begin{tabularx}{0.7\textwidth}{X} 一共有八塊板和十六個帶卯眼的銀座，每塊板底下有兩個卯眼。 \end{tabularx} \\ \\ \relax
36:31 & \begin{tabularx}{0.7\textwidth}{X} 他用金合歡木做橫木：為帳幕這面的板做五根橫木， \end{tabularx} \\ \\ \relax
36:32 & \begin{tabularx}{0.7\textwidth}{X} 為帳幕那面的板做五根橫木，又為帳幕後面，就是朝西的板做五根橫木， \end{tabularx} \\ \\ \relax
36:33 & \begin{tabularx}{0.7\textwidth}{X} 他做了板腰間的橫木，從一頭通到另一頭。 \end{tabularx} \\ \\ \relax
36:34 & \begin{tabularx}{0.7\textwidth}{X} 他將板包上金子，又做板上的金環來套橫木；橫木也包上金子。 \end{tabularx} \\ \\ \relax
36:35 & \begin{tabularx}{0.7\textwidth}{X} 他用藍色、紫色、朱紅色紗，和搓的細麻織幔子，以刺繡的手藝繡上基路伯。 \end{tabularx} \\ \\ \relax
36:36 & \begin{tabularx}{0.7\textwidth}{X} 他又用金合歡木為幔子做四根柱子，包上金子，柱子有金鉤，又為柱子鑄了四個帶卯眼的銀座。 \end{tabularx} \\ \\ \relax
36:37 & \begin{tabularx}{0.7\textwidth}{X} 他用藍色、紫色、朱紅色紗，和搓的細麻，以刺繡的手藝為帳幕織門簾， \end{tabularx} \\ \\ \relax
36:38 & \begin{tabularx}{0.7\textwidth}{X} 又為簾子做五根柱子和柱子的鉤子，把柱頂和柱子的箍包上金子。柱子有五個帶卯眼的銅座。 \end{tabularx} \\ \\ \relax
37:1 & \begin{tabularx}{0.7\textwidth}{X} 比撒列用金合歡木做一個櫃子，長二肘半，寬一肘半，高一肘半。 \end{tabularx} \\ \\ \relax
37:2 & \begin{tabularx}{0.7\textwidth}{X} 裡裡外外包上純金，四圍鑲上金邊。 \end{tabularx} \\ \\ \relax
37:3 & \begin{tabularx}{0.7\textwidth}{X} 他又鑄了四個金環，安在櫃子的四腳上；這邊兩個環，那邊兩個環。 \end{tabularx} \\ \\ \relax
37:4 & \begin{tabularx}{0.7\textwidth}{X} 他用金合歡木做了兩根槓，包上金子， \end{tabularx} \\ \\ \relax
37:5 & \begin{tabularx}{0.7\textwidth}{X} 又把槓穿過櫃旁的環，以便抬櫃。 \end{tabularx} \\ \\ \relax
37:6 & \begin{tabularx}{0.7\textwidth}{X} 他用純金做了一個櫃蓋，長二肘半，寬一肘半， \end{tabularx} \\ \\ \relax
37:7 & \begin{tabularx}{0.7\textwidth}{X} 他造兩個用金子錘出的基路伯，從櫃蓋的兩端錘出它們。 \end{tabularx} \\ \\ \relax
37:8 & \begin{tabularx}{0.7\textwidth}{X} 這端一個基路伯，那端一個基路伯；從櫃蓋的兩端錘出兩個基路伯。 \end{tabularx} \\ \\ \relax
37:9 & \begin{tabularx}{0.7\textwidth}{X} 二基路伯的翅膀向上張開，用翅膀遮住櫃蓋，臉彼此相對；基路伯的臉朝向櫃蓋。 \end{tabularx} \\ \\ \relax
37:10 & \begin{tabularx}{0.7\textwidth}{X} 他用金合歡木做了一張供桌，長二肘，寬一肘，高一肘半， \end{tabularx} \\ \\ \relax
37:11 & \begin{tabularx}{0.7\textwidth}{X} 把它包上純金，四圍鑲上金邊。 \end{tabularx} \\ \\ \relax
37:12 & \begin{tabularx}{0.7\textwidth}{X} 供桌的四圍各做了一掌寬的邊緣，邊緣鑲上金邊。 \end{tabularx} \\ \\ \relax
37:13 & \begin{tabularx}{0.7\textwidth}{X} 他又鑄了四個金環，把環安在四個桌腳的四角上。 \end{tabularx} \\ \\ \relax
37:14 & \begin{tabularx}{0.7\textwidth}{X} 環靠近邊緣，以便穿槓抬供桌。 \end{tabularx} \\ \\ \relax
37:15 & \begin{tabularx}{0.7\textwidth}{X} 他用金合歡木做了兩根槓，包上金子，用來抬供桌。 \end{tabularx} \\ \\ \relax
37:16 & \begin{tabularx}{0.7\textwidth}{X} 他又用純金做了桌上的器具，就是盤、碟，以及澆酒祭的杯和壺。 \end{tabularx} \\ \\ \relax
37:17 & \begin{tabularx}{0.7\textwidth}{X} 他造一座用純金錘出的燈臺；燈臺的座、幹、杯、花萼和花瓣，都和燈臺接連一塊。 \end{tabularx} \\ \\ \relax
37:18 & \begin{tabularx}{0.7\textwidth}{X} 燈臺兩旁伸出六根枝子：這邊三根，那邊三根。 \end{tabularx} \\ \\ \relax
37:19 & \begin{tabularx}{0.7\textwidth}{X} 這邊的枝子上有三個杯，形狀像杏花，有花萼有花瓣；那邊的枝子上也有三個杯，形狀像杏花，有花萼有花瓣。從燈臺伸出來的六根枝子都是如此。 \end{tabularx} \\ \\ \relax
37:20 & \begin{tabularx}{0.7\textwidth}{X} 燈臺本身有四個杯，形狀像杏花，有花萼有花瓣。 \end{tabularx} \\ \\ \relax
37:21 & \begin{tabularx}{0.7\textwidth}{X} 燈臺的第一對枝子下面有花萼，燈臺的第二對枝子下面有花萼，燈臺的第三對枝子下面也有花萼；燈臺伸出的六根枝子都是如此。 \end{tabularx} \\ \\ \relax
37:22 & \begin{tabularx}{0.7\textwidth}{X} 花萼和枝子都和燈臺接連一塊，全是從一塊純金錘出來的。 \end{tabularx} \\ \\ \relax
37:23 & \begin{tabularx}{0.7\textwidth}{X} 他用純金做燈臺的七盞燈，以及燈剪和燈盤。 \end{tabularx} \\ \\ \relax
37:24 & \begin{tabularx}{0.7\textwidth}{X} 他用一他連得的純金做燈臺和燈臺的一切器具。 \end{tabularx} \\ \\ \relax
37:25 & \begin{tabularx}{0.7\textwidth}{X} 他用金合歡木做香壇，長一肘，寬一肘，這壇是正方形的，高二肘。壇的四個翹角與壇接連一塊。 \end{tabularx} \\ \\ \relax
37:26 & \begin{tabularx}{0.7\textwidth}{X} 他把壇的上面與壇的四圍，以及壇的四個翹角包上純金，又在壇的四圍鑲上金邊。 \end{tabularx} \\ \\ \relax
37:27 & \begin{tabularx}{0.7\textwidth}{X} 他在壇的兩個對側，金邊下面做了兩個金環，用來穿槓抬壇。 \end{tabularx} \\ \\ \relax
37:28 & \begin{tabularx}{0.7\textwidth}{X} 他又用金合歡木做槓，包上金子。 \end{tabularx} \\ \\ \relax
37:29 & \begin{tabularx}{0.7\textwidth}{X} 他按配製香料的方法製成聖膏油和芬芳的純香。 \end{tabularx} \\ \\ \relax
38:1 & \begin{tabularx}{0.7\textwidth}{X} 他用金合歡木做燔祭壇，長五肘，寬五肘，是正方形的，高三肘。 \end{tabularx} \\ \\ \relax
38:2 & \begin{tabularx}{0.7\textwidth}{X} 在壇的四角做四個翹角，與壇接連一塊，把壇包上銅。 \end{tabularx} \\ \\ \relax
38:3 & \begin{tabularx}{0.7\textwidth}{X} 他做壇的一切器具，就是桶子、鏟子、盤子、肉叉和火盆；這一切器具都是用銅做的。 \end{tabularx} \\ \\ \relax
38:4 & \begin{tabularx}{0.7\textwidth}{X} 他又為壇做一個銅網，安在壇四圍的邊的下面，垂到壇的半腰。 \end{tabularx} \\ \\ \relax
38:5 & \begin{tabularx}{0.7\textwidth}{X} 他在銅網的四角上鑄了四個環，用來穿槓。 \end{tabularx} \\ \\ \relax
38:6 & \begin{tabularx}{0.7\textwidth}{X} 他用金合歡木做槓，包上銅， \end{tabularx} \\ \\ \relax
38:7 & \begin{tabularx}{0.7\textwidth}{X} 把槓穿過壇兩旁的環子，用來抬壇。他用板做壇，壇的中心是空的。 \end{tabularx} \\ \\ \relax
38:8 & \begin{tabularx}{0.7\textwidth}{X} 他用銅做洗濯盆和盆座，是用會幕門前事奉之婦人的銅鏡做的。 \end{tabularx} \\ \\ \relax
38:9 & \begin{tabularx}{0.7\textwidth}{X} 他又做院子，在南面，就是面向南方的那一邊，用搓的細麻做院子的帷幔，一百肘。 \end{tabularx} \\ \\ \relax
38:10 & \begin{tabularx}{0.7\textwidth}{X} 帷幔有二十根柱子，二十個帶卯眼的銅座；柱子的鉤和箍都是銀的。 \end{tabularx} \\ \\ \relax
38:11 & \begin{tabularx}{0.7\textwidth}{X} 北面的帷幔一百肘。帷幔有二十根柱子，二十個帶卯眼的銅座；柱子的鉤和箍都是銀的。 \end{tabularx} \\ \\ \relax
38:12 & \begin{tabularx}{0.7\textwidth}{X} 西面的帷幔五十肘。帷幔有十根柱子，十個帶卯眼的座；柱子的鉤和箍都是銀的。 \end{tabularx} \\ \\ \relax
38:13 & \begin{tabularx}{0.7\textwidth}{X} 院子的東面，就是面向東方的那一邊，五十肘。 \end{tabularx} \\ \\ \relax
38:14 & \begin{tabularx}{0.7\textwidth}{X} 一邊的帷幔有十五肘，有三根柱子，三個帶卯眼的座。 \end{tabularx} \\ \\ \relax
38:15 & \begin{tabularx}{0.7\textwidth}{X} 另一邊也一樣，院子門口左右的帷幔也有十五肘，有三根柱子，三個帶卯眼的座。 \end{tabularx} \\ \\ \relax
38:16 & \begin{tabularx}{0.7\textwidth}{X} 院子四面的帷幔都是用搓的細麻做的。 \end{tabularx} \\ \\ \relax
38:17 & \begin{tabularx}{0.7\textwidth}{X} 柱子帶卯眼的座是銅的，柱子的鉤和箍是銀的，柱頂是用銀包的。院子一切的柱子都是用銀子箍著的。 \end{tabularx} \\ \\ \relax
38:18 & \begin{tabularx}{0.7\textwidth}{X} 院子的門簾是以刺繡的手藝，用藍色、紫色、朱紅色紗，和搓的細麻織的，長二十肘，寬也就是高五肘，與院子帷幔的高度相同。 \end{tabularx} \\ \\ \relax
38:19 & \begin{tabularx}{0.7\textwidth}{X} 門簾有四根柱子，四個帶卯眼的銅座；柱子上的鉤和箍是銀的，柱頂是用銀包的。 \end{tabularx} \\ \\ \relax
38:20 & \begin{tabularx}{0.7\textwidth}{X} 帳幕一切的橛子和院子四圍的橛子都是銅的。 \end{tabularx} \\ \\ \relax
38:21 & \begin{tabularx}{0.7\textwidth}{X} 這是帳幕，就是法櫃帳幕中物件的總數，是照摩西的吩咐，亞倫祭司的兒子以他瑪經手，利未人數點的。 \end{tabularx} \\ \\ \relax
38:22 & \begin{tabularx}{0.7\textwidth}{X} 凡耶和華吩咐摩西的，都是由猶大支派中戶珥的孫子，烏利的兒子比撒列去做的； \end{tabularx} \\ \\ \relax
38:23 & \begin{tabularx}{0.7\textwidth}{X} 與他同工的有但支派中亞希撒抹的兒子亞何利亞伯；他是雕刻師，也是設計師，又是用藍色、紫色、朱紅色紗，和細麻的刺繡師。 \end{tabularx} \\ \\ \relax
38:24 & \begin{tabularx}{0.7\textwidth}{X} 為聖所一切工作用的金子，就是所奉獻的金子，按聖所的舍客勒，一共是二十九他連得，七百三十舍客勒。 \end{tabularx} \\ \\ \relax
38:25 & \begin{tabularx}{0.7\textwidth}{X} 會中被數的人所獻的銀子，按聖所的舍客勒，一共是一百他連得，一千七百七十五舍客勒。 \end{tabularx} \\ \\ \relax
38:26 & \begin{tabularx}{0.7\textwidth}{X} 凡曾被數的，就是二十歲以上的人，共有六十萬零三千五百五十人。按聖所的舍客勒，每人半舍客勒，就是一比加。 \end{tabularx} \\ \\ \relax
38:27 & \begin{tabularx}{0.7\textwidth}{X} 一百他連得銀子是用來鑄造聖所帶卯眼的座和幔子下帶卯眼的座；用一百他連得鑄造一百個帶卯眼的座，每個帶卯眼的座一他連得。 \end{tabularx} \\ \\ \relax
38:28 & \begin{tabularx}{0.7\textwidth}{X} 一千七百七十五舍客勒是用來鑄造柱子的鉤，包柱頂，以及箍著柱子。 \end{tabularx} \\ \\ \relax
38:29 & \begin{tabularx}{0.7\textwidth}{X} 所奉獻的銅共有七十他連得，二千四百舍客勒。 \end{tabularx} \\ \\ \relax
38:30 & \begin{tabularx}{0.7\textwidth}{X} 這些銅是用來做會幕門口帶卯眼的座，銅壇、壇的銅網和壇的一切器具， \end{tabularx} \\ \\ \relax
38:31 & \begin{tabularx}{0.7\textwidth}{X} 院子四圍帶卯眼的座和院子門口帶卯眼的座，以及帳幕一切的橛子和院子四圍所有的橛子。 \end{tabularx} \\ \\ \relax
39:1 & \begin{tabularx}{0.7\textwidth}{X} 他們用藍色、紫色、朱紅色紗縫製精緻的禮服，在聖所用以供職；他們為亞倫做聖衣，是照耶和華所吩咐摩西的。 \end{tabularx} \\ \\ \relax
39:2 & \begin{tabularx}{0.7\textwidth}{X} 以弗得是用金色、藍色、紫色、朱紅色紗，和搓的細麻做的。 \end{tabularx} \\ \\ \relax
39:3 & \begin{tabularx}{0.7\textwidth}{X} 他們把金子錘成薄片，剪成細線，與藍色、紫色、朱紅色紗，以刺繡的手藝織在一起。 \end{tabularx} \\ \\ \relax
39:4 & \begin{tabularx}{0.7\textwidth}{X} 他們又為以弗得做兩條相連的肩帶，接連在以弗得的兩端。 \end{tabularx} \\ \\ \relax
39:5 & \begin{tabularx}{0.7\textwidth}{X} 以弗得的精緻帶子以一樣的手藝，用金色、藍色、紫色、朱紅色紗，和搓的細麻縫製，與以弗得接連在一起，是照耶和華所吩咐摩西的。 \end{tabularx} \\ \\ \relax
39:6 & \begin{tabularx}{0.7\textwidth}{X} 他們琢出兩塊紅瑪瑙，鑲在金槽裡，如同刻印章，刻上以色列眾子的名字。 \end{tabularx} \\ \\ \relax
39:7 & \begin{tabularx}{0.7\textwidth}{X} 他把這兩塊寶石安在以弗得的兩條肩帶上，為以色列人作紀念石，是照耶和華所吩咐摩西的。 \end{tabularx} \\ \\ \relax
39:8 & \begin{tabularx}{0.7\textwidth}{X} 胸袋是以刺繡的手藝，如同以弗得的做法，用金色、藍色、紫色、朱紅色紗，和搓的細麻縫製。 \end{tabularx} \\ \\ \relax
39:9 & \begin{tabularx}{0.7\textwidth}{X} 胸袋是正方形的，他們把它做成兩層，這兩層各長一虎口，寬一虎口。 \end{tabularx} \\ \\ \relax
39:10 & \begin{tabularx}{0.7\textwidth}{X} 他們在上面鑲四行寶石：第一行是紅寶石、紅璧璽、紅玉； \end{tabularx} \\ \\ \relax
39:11 & \begin{tabularx}{0.7\textwidth}{X} 第二行是綠寶石、藍寶石、金剛石； \end{tabularx} \\ \\ \relax
39:12 & \begin{tabularx}{0.7\textwidth}{X} 第三行是紫瑪瑙、白瑪瑙、紫晶； \end{tabularx} \\ \\ \relax
39:13 & \begin{tabularx}{0.7\textwidth}{X} 第四行是水蒼玉、紅瑪瑙、碧玉。這些都鑲在金槽中。 \end{tabularx} \\ \\ \relax
39:14 & \begin{tabularx}{0.7\textwidth}{X} 這些寶石有以色列十二個兒子的名字，如同刻印章，每一顆有自己的名字，代表十二個支派。 \end{tabularx} \\ \\ \relax
39:15 & \begin{tabularx}{0.7\textwidth}{X} 他們在胸袋上用純金打鏈子，像編成的繩子一樣。 \end{tabularx} \\ \\ \relax
39:16 & \begin{tabularx}{0.7\textwidth}{X} 他們又做了兩個金槽和兩個金環，把這兩個環安在胸袋的兩端。 \end{tabularx} \\ \\ \relax
39:17 & \begin{tabularx}{0.7\textwidth}{X} 他們把那兩條編成的金鏈繫在胸袋兩端的兩個環上， \end{tabularx} \\ \\ \relax
39:18 & \begin{tabularx}{0.7\textwidth}{X} 又把鏈子的另外兩端扣在兩個槽上，安在以弗得前面的肩帶上。 \end{tabularx} \\ \\ \relax
39:19 & \begin{tabularx}{0.7\textwidth}{X} 他們做了兩個金環，安在胸袋的兩端，在以弗得裡面的邊上， \end{tabularx} \\ \\ \relax
39:20 & \begin{tabularx}{0.7\textwidth}{X} 又做兩個金環，安在以弗得前面兩條肩帶的下邊，靠近接縫處，在精緻帶子的上面。 \end{tabularx} \\ \\ \relax
39:21 & \begin{tabularx}{0.7\textwidth}{X} 他們用藍色的帶子把胸袋的環與以弗得的環繫住，使胸袋綁在以弗得精緻的帶子上，不致鬆脫，是照耶和華所吩咐摩西的。 \end{tabularx} \\ \\ \relax
39:22 & \begin{tabularx}{0.7\textwidth}{X} 以弗得的外袍是以編織的手藝做的，顏色全是藍的。 \end{tabularx} \\ \\ \relax
39:23 & \begin{tabularx}{0.7\textwidth}{X} 袍上方的中間留了一個領口，領口的周圍織出領邊，好像鎧甲的領口，免得破裂。 \end{tabularx} \\ \\ \relax
39:24 & \begin{tabularx}{0.7\textwidth}{X} 他們在袍子下襬用藍色、紫色、朱紅色紗，和搓的細麻做石榴， \end{tabularx} \\ \\ \relax
39:25 & \begin{tabularx}{0.7\textwidth}{X} 又用純金鑄了鈴鐺，把鈴鐺釘在石榴中間，袍子下襬周圍的石榴中間： \end{tabularx} \\ \\ \relax
39:26 & \begin{tabularx}{0.7\textwidth}{X} 一個鈴鐺一個石榴，一個鈴鐺一個石榴，在袍子下襬的周圍，用以供職，是照耶和華所吩咐摩西的。 \end{tabularx} \\ \\ \relax
39:27 & \begin{tabularx}{0.7\textwidth}{X} 他們用編織的工為亞倫和他的兒子做細麻布內袍、 \end{tabularx} \\ \\ \relax
39:28 & \begin{tabularx}{0.7\textwidth}{X} 細麻布禮冠、細麻布精緻頭巾，和搓的細麻布褲子， \end{tabularx} \\ \\ \relax
39:29 & \begin{tabularx}{0.7\textwidth}{X} 又用藍色、紫色、朱紅色紗，和搓的細麻，以刺繡的手藝做腰帶，是照耶和華所吩咐摩西的。 \end{tabularx} \\ \\ \relax
39:30 & \begin{tabularx}{0.7\textwidth}{X} 他們用純金做一面聖冠上的牌，如同刻印章，在上面寫著「歸耶和華為聖」， \end{tabularx} \\ \\ \relax
39:31 & \begin{tabularx}{0.7\textwidth}{X} 又用藍色的帶子把牌繫在禮冠上，是照耶和華所吩咐摩西的。 \end{tabularx} \\ \\ \relax
39:32 & \begin{tabularx}{0.7\textwidth}{X} 會幕的帳幕一切的工程就這樣做完了。凡耶和華所吩咐摩西的，以色列人都照樣做了。 \end{tabularx} \\ \\ \relax
39:33 & \begin{tabularx}{0.7\textwidth}{X} 他們把帳幕運到摩西那裡，帳幕和帳幕的一切器具，就是鉤、板、橫木、柱子、帶卯眼的座， \end{tabularx} \\ \\ \relax
39:34 & \begin{tabularx}{0.7\textwidth}{X} 染紅公羊皮的蓋、精美皮料的蓋、遮掩的幔子， \end{tabularx} \\ \\ \relax
39:35 & \begin{tabularx}{0.7\textwidth}{X} 法櫃、櫃的槓、櫃蓋， \end{tabularx} \\ \\ \relax
39:36 & \begin{tabularx}{0.7\textwidth}{X} 供桌、供桌的一切器具、供餅， \end{tabularx} \\ \\ \relax
39:37 & \begin{tabularx}{0.7\textwidth}{X} 純金的燈臺、擺列的燈、燈臺的一切器具、點燈的油， \end{tabularx} \\ \\ \relax
39:38 & \begin{tabularx}{0.7\textwidth}{X} 金壇、膏油、芬芳的香、帳幕的門簾， \end{tabularx} \\ \\ \relax
39:39 & \begin{tabularx}{0.7\textwidth}{X} 銅壇、壇的銅網、壇的槓、壇的一切器具，洗濯盆和盆座， \end{tabularx} \\ \\ \relax
39:40 & \begin{tabularx}{0.7\textwidth}{X} 院子的帷幔、柱子、帶卯眼的座、院子的門簾、繩子、橛子，帳幕，就是會幕使用的一切器具， \end{tabularx} \\ \\ \relax
39:41 & \begin{tabularx}{0.7\textwidth}{X} 以及聖所事奉用的精緻禮服，亞倫祭司的聖衣和他兒子的衣服，供祭司職分用。 \end{tabularx} \\ \\ \relax
39:42 & \begin{tabularx}{0.7\textwidth}{X} 這一切工作都是以色列人照耶和華所吩咐摩西做的。 \end{tabularx} \\ \\ \relax
39:43 & \begin{tabularx}{0.7\textwidth}{X} 摩西看見這一切的工，看哪，耶和華怎樣吩咐，他們就照樣做了，摩西就為他們祝福。 \end{tabularx} \\ \\ \relax
40:1 & \begin{tabularx}{0.7\textwidth}{X} 耶和華吩咐摩西說： \end{tabularx} \\ \\ \relax
40:2 & \begin{tabularx}{0.7\textwidth}{X} 「正月初一，你要立起會幕的帳幕， \end{tabularx} \\ \\ \relax
40:3 & \begin{tabularx}{0.7\textwidth}{X} 把法櫃安放在裡面，用幔子將櫃遮掩。 \end{tabularx} \\ \\ \relax
40:4 & \begin{tabularx}{0.7\textwidth}{X} 把供桌搬進去，擺設桌上的器具。又把燈臺搬進去，點上燈。 \end{tabularx} \\ \\ \relax
40:5 & \begin{tabularx}{0.7\textwidth}{X} 把金香壇安在法櫃前，掛上帳幕的門簾。 \end{tabularx} \\ \\ \relax
40:6 & \begin{tabularx}{0.7\textwidth}{X} 把燔祭壇安在會幕的帳幕門前。 \end{tabularx} \\ \\ \relax
40:7 & \begin{tabularx}{0.7\textwidth}{X} 把洗濯盆安在會幕和壇的中間，在盆裡盛水。 \end{tabularx} \\ \\ \relax
40:8 & \begin{tabularx}{0.7\textwidth}{X} 又要在院子周圍支起帷幔，把院子的門簾掛上。 \end{tabularx} \\ \\ \relax
40:9 & \begin{tabularx}{0.7\textwidth}{X} 你要用膏油抹帳幕和其中所有的，使帳幕和一切器具分別為聖，就都成為聖。 \end{tabularx} \\ \\ \relax
40:10 & \begin{tabularx}{0.7\textwidth}{X} 又要抹燔祭壇和壇的一切器具，使壇分別為聖，壇就成為至聖。 \end{tabularx} \\ \\ \relax
40:11 & \begin{tabularx}{0.7\textwidth}{X} 要抹洗濯盆和盆座，使盆分別為聖。 \end{tabularx} \\ \\ \relax
40:12 & \begin{tabularx}{0.7\textwidth}{X} 你要帶亞倫和他兒子到會幕門口，用水洗身。 \end{tabularx} \\ \\ \relax
40:13 & \begin{tabularx}{0.7\textwidth}{X} 要給亞倫穿上聖衣，又膏他，使他分別為聖，作事奉我的祭司。 \end{tabularx} \\ \\ \relax
40:14 & \begin{tabularx}{0.7\textwidth}{X} 又要帶他的兒子來，給他們穿上內袍。 \end{tabularx} \\ \\ \relax
40:15 & \begin{tabularx}{0.7\textwidth}{X} 你怎樣膏他們的父親，也要照樣膏他們，使他們成為事奉我的祭司。他們受了膏，就必世世代代永遠得祭司的職分。」 \end{tabularx} \\ \\ \relax
40:16 & \begin{tabularx}{0.7\textwidth}{X} 摩西這樣做了；耶和華怎樣吩咐摩西，他就照樣做了。 \end{tabularx} \\ \\ \relax
40:17 & \begin{tabularx}{0.7\textwidth}{X} 第二年正月初一，帳幕就立起來。 \end{tabularx} \\ \\ \relax
40:18 & \begin{tabularx}{0.7\textwidth}{X} 摩西支起帳幕，安上帶卯眼的座，安上板，穿上橫木，立起柱子。 \end{tabularx} \\ \\ \relax
40:19 & \begin{tabularx}{0.7\textwidth}{X} 他在帳幕的上面搭上罩棚，把罩棚外層的蓋子蓋在其上，是照著耶和華所吩咐他的。 \end{tabularx} \\ \\ \relax
40:20 & \begin{tabularx}{0.7\textwidth}{X} 他把法版放在櫃裡，把槓穿在櫃的兩旁，把櫃蓋安在櫃上。 \end{tabularx} \\ \\ \relax
40:21 & \begin{tabularx}{0.7\textwidth}{X} 把櫃抬進帳幕，掛上遮掩櫃的幔子，把法櫃遮蓋了，是照耶和華所吩咐摩西的。 \end{tabularx} \\ \\ \relax
40:22 & \begin{tabularx}{0.7\textwidth}{X} 他把供桌安在會幕內，在帳幕的北邊，幔子的外面。 \end{tabularx} \\ \\ \relax
40:23 & \begin{tabularx}{0.7\textwidth}{X} 把餅擺設在供桌上，在耶和華面前，是照耶和華所吩咐摩西的。 \end{tabularx} \\ \\ \relax
40:24 & \begin{tabularx}{0.7\textwidth}{X} 他把燈臺安在會幕內，在帳幕的南邊，供桌的對面， \end{tabularx} \\ \\ \relax
40:25 & \begin{tabularx}{0.7\textwidth}{X} 並在耶和華面前點燈，是照耶和華所吩咐摩西的。 \end{tabularx} \\ \\ \relax
40:26 & \begin{tabularx}{0.7\textwidth}{X} 他把金壇安在會幕內，幔子的前面， \end{tabularx} \\ \\ \relax
40:27 & \begin{tabularx}{0.7\textwidth}{X} 又在壇上燒芬芳的香，是照耶和華所吩咐摩西的。 \end{tabularx} \\ \\ \relax
40:28 & \begin{tabularx}{0.7\textwidth}{X} 他又掛上帳幕的門簾。 \end{tabularx} \\ \\ \relax
40:29 & \begin{tabularx}{0.7\textwidth}{X} 在會幕的帳幕門口安設燔祭壇，把燔祭和素祭獻在壇上，是照耶和華所吩咐摩西的。 \end{tabularx} \\ \\ \relax
40:30 & \begin{tabularx}{0.7\textwidth}{X} 他又把洗濯盆安在會幕和祭壇的中間，盆裡盛水，以便洗濯。 \end{tabularx} \\ \\ \relax
40:31 & \begin{tabularx}{0.7\textwidth}{X} 摩西和亞倫，以及亞倫的兒子用這盆洗手洗腳。 \end{tabularx} \\ \\ \relax
40:32 & \begin{tabularx}{0.7\textwidth}{X} 他們進會幕或走近壇的時候，就都洗濯，是照耶和華所吩咐摩西的。 \end{tabularx} \\ \\ \relax
40:33 & \begin{tabularx}{0.7\textwidth}{X} 他在帳幕和祭壇的四圍支起院子的帷幔，把院子的門簾掛上。這樣，摩西就做完了工。 \end{tabularx} \\ \\ \relax
40:34 & \begin{tabularx}{0.7\textwidth}{X} 那時，雲彩遮蓋會幕，耶和華的榮光充滿了帳幕。 \end{tabularx} \\ \\ \relax
40:35 & \begin{tabularx}{0.7\textwidth}{X} 摩西不能進會幕，因為雲彩停在其上，耶和華的榮光充滿了帳幕。 \end{tabularx} \\ \\ \relax
40:36 & \begin{tabularx}{0.7\textwidth}{X} 每逢雲彩從帳幕升上去，以色列人就起程前行； \end{tabularx} \\ \\ \relax
40:37 & \begin{tabularx}{0.7\textwidth}{X} 雲彩若不升上去，他們就不起程，直等到雲彩升上去。 \end{tabularx} \\ \\ \relax
40:38 & \begin{tabularx}{0.7\textwidth}{X} 在他們所行的路上，在以色列全家的眼前，白天，耶和華的雲彩在帳幕上，黑夜，有火在雲彩中。 \end{tabularx} \\ \\
[1ex]
\hline
\hline
\end{longtable}
$^{1}$各位電影姐妹,大家好.
今次我們繼續我們的系列.
今次會有少許東西獎給大家.
所以如果你還沒有瓶子的話.
大家可以拿.
但這個瓶子裡面已經有好幾樣的物件在裡面.
我們希望今年裡面這個瓶子去提醒大家.
我們在上帝裡面有甚麼恩典.
今次是我們這個NEXT的第四講.
你見到那個Title是回眸.
我們就看看舊約聖經裡面講到的回眸.
你記得勝經是一個大的故事.
即是一個長篇劇有66集.
因為聖經裡面有66卷書.
由《創世記》到最後一卷書叫做《啟示錄》.
中間就是基督耶穌的來臨.
所以舊約裡面就算是一部分上半的長篇劇.
有39卷書.
所以你計它39集.
然後新約就有27卷書.
你就計它後面的27集.
所以我們由第一集看到最後一集的時候.
我們應該大概get到.
雖然每一集每一卷書.
講的東西是有些不同.
它們有一個主題是獨特的主題.
但是你串在一起的時候.
你就發現原來背後有一個編劇.
整個故事就呈現出來了.
整個大的故事.
其實就是介紹這位上帝是一個怎樣的上帝.
一個怎樣對待你的上帝.
一個跟你怎樣.
應該建立一個怎樣關係的一個上帝.
所以由頭到尾第一集到最後一集.
都是一個故事.
一個宏大的故事.
在故事裡面就去告訴你.
上帝告訴你我是怎樣的.
等於我們彼此認識的時候.

$^{41}$我們大家見到新朋友來到.
我們彼此介紹的時候.
你都會說甚麼.
由你出生開始說.
說到你現在.
我在哪裡出生 經歷過甚麼.
等等等等 有時間的話.
因為聽了你的故事的時候.
我就明白你是一個怎樣的人.
明白你的故事就是明白你.
明白上帝的故事也是明白和認識上帝.
所以我們第一次開始.
我們就由創世記開始.
直到現在我們看到出埃及記的後半段.
出埃及記的後半段.
我知道Alice傳道是在說出埃及記.
不過我比她快很多.
我們先用一個宏大的.
不用仔細的方法去看.
但我們今天看的是.
二十五章第四十章.
出埃及記.
我們看的時候.
我們首先問一個問題.
很多時候我們想像一個神.
我們想像一個.
很高超的.
很高的.
不能夠觸及的.
一個存在者.
或者那些古代的人.
他們就會想像.
神是一個怎樣的神.
他們想像神明.
任何一種神明.
他們就會想像一定在高山上.
因為你要攀上去.
你要爬上去.
你才能夠和他接觸.
他就不會下來.

$^{81}$這個不單是古代的人對.
宗教的信念.
其實現代人也是這樣的.
如果你是有一點.
宗教的背景.
你未相信基督之前.
未認識耶穌之前.
背景或者任何.
少許的概念也好.
神明總是遙遙不及的.
你是在上天的.
然後可能古代的人也好.
現代的人也好.
就弄一些大的寺廟.
或者在山上弄一些建築物.
然後你希望和那些人.
和他接觸到能夠攀山.
你上去和他.
有些接近.
無論民間信仰或者以前的都是一樣.
這樣的概念.
是很功利的.
因為是一個交換.
神是遠遠.
超越我們的.
然後我們需要用我們的方法.
走去祂那裡.
無論是藉著禮儀.
或者藉著一些不同的.
祭祀.
或者不同的方法.
我們希望去到神那裡.
但是在埃及記那裡介紹的神.
往往是相反的.
因為你記得以色列人.
在埃及出來之後.
上帝用一手拯救了他們.
他們就去了一個地方叫西乃山.
西乃山.
然後他們接納了.

$^{121}$上帝給他們的律法.
成為立約的子民.
然後到第25章開始.
一路講.
其實是講一些我們平時沒什麼興趣聽的東西.
就是弄一個東西出來.
弄一個所謂的會幕.
中文譯做會幕其實就是帳幕.
就是Tent.
25章第8節講什麼呢.
祂們要為.
就是上帝的子民.
要為我做一個聖所.
使我住在他們中間.
我將我的家.
住在你們裡頭.
建造我的家.
在你們裡頭.
住在他們中間.
你們要按照我的指示.
帳幕和一切的器具.
樣式照樣去做.
上帝就在.
第25章開始.
就一路很detail的.
很細節的去講會幕是怎樣建造的.
所以就.
對我們來說.
沒什麼興趣.
我記得有一次.
在神樂園的時候.
有一次我不記得去那裡.
我們有一個retreat.
就是和其他神樂生一起.
神樂生就很熟齡的.
我們早上當然有組會.
我們就在Airbnb.
大屋那裡.
我們就有小組的靈修.
不知道哪個人就選了這一段來讀.

$^{161}$就是25章.
就一讀就讀了4,5章.
如果你想知道.
細節的話你回去看一下.
我不會讀.
他一路讀一路讀.
就是讀那些會幕.
那些萬紙的curtain.
那些尺寸.
幾多爪幾多爪.
然後就做金燈台.
要多大多小.
全部都是尺寸的東西.
全部都是home depot的東西.
home depot就要量度.
幾多尺幾多寸.
讀得我們很悶.
那我心想.
什麼時候讀完.
這些就是我們沒興趣的東西.
讀完之後有一個弟兄就說得很對.
他說不如我們來祈禱.
我說祈禱?.
去完home depot怎麼會祈禱.
去完home depot.
全部都要看間隔.
量度一下.
kitchen counter有多大有多長有多闊.
接著我們不如祈禱.
他說有什麼好祈禱.
他真的祈禱.
那個弟兄就祈禱.
他說多謝主.
因為我們見到你的居所.
有幾個細節的安排.
他這樣祈禱的時候.
我覺得很有道理.
很有道理.
對我們現代人的耳朵來說.
我們沒什麼興趣去看這些東西.

$^{201}$但是呢.
上帝是建造他的居所.
在他的子民中間.
住在他們中間的時候.
他是很細微的很細節的.
去講出他是一個怎樣的神.
所以那些家具和間隔.
會幕的設計.
其實是一些.
建築的神學.
architectural.
theology.
用一些architect.
或者建築的.
藍圖的方法.
來表達他是一個怎樣的上帝.
我記得.
最近有一個朋友.
聊天.
他來了芝加哥很久.
我問他你有沒有去過architecture tour.
你知道我們.
navy pier那裡.
可以上船然後在chicago river上去.
然後去介紹不同的建築物.
有誰去過.
有誰沒去過.
然後我想沒理由沒去過.
不是吧.
我說這個是.
芝加哥的one of the best.
one of the best.
是不是.
在那個tour.
他記下所有的建築物.
然後告訴你.
芝加哥是一個怎樣的城市.
就是用這些architect.
to tell the story.
it's amazing.

$^{241}$去看一下上帝.
怎樣去設計他的居所.
在我們中間.
以色列南部有一個地方.
你開車去開幾個小時.
就是為了看這個東西.
就是一個1比1的.
匯膜的模型.
你見到一個匯膜.
其實就是純粹好像你平時露營.
出去露營的時候.
你開的營是一樣的.
就開一個tent.
在他的神的指紋中間.
對於我們.
對聖經不是很熟的.
我們就會很奇怪.
為何上帝會這樣做.
原來在這些.
architect.
architecture裡面.
他想表達一件事.
就是神學.
上帝是一個怎樣的上帝.
你見到.
這個匯膜的平面圖.
由右至左我們進去的時候.
右邊就是入口.
然後有一個外園.
外園就是長方形.
然後你慢慢進展.
進入最中心的地方的時候.
就變成一個perfect square.
一個完美的.
方形.
在古代近東的.
畫藍圖.
你找一個近東的結節計師去畫藍圖的時候.
他畫.
畫一個perfect square的時候.

$^{281}$就代表這個地方是最聖潔.
和最神聖的地方.
perfect square.
長和闊都是一模一樣.
但你走出來的時候就是一個長方形.
或者兩個正方形.
你在外園裡面.
見到甚麼設備呢.
第一個設備你見到的就是.
祭壇.
和銅盤.
這裡是在說祭祀的東西.
就是說你approach上帝.
進入去見上帝的時候.
你就要知道.
原來是有代價在.
那個代價就是上帝所提供的.
祭祀.
所以以色列神的子民.
要進去上帝這條路.
他就需要去獻祭.
你不可以空手隨隨便便來.
你要帶著梵祭.
其他的祭物等等.
在摩西五經裡面有提及.
就是上帝給一條路.
給人進入去.
他自己的童宅裡面.
你經過這個外園.
洗過手.
銅盤就是水盤.
祭壇.
其實就是燒烤的燒烤.
煮了之後.
燒了之後.
你就知道.
it cost something to approach.
然後你就去到另一個空間.
就是一個小的長方形空間.
那個就叫holy place.

$^{321}$或者聖所.
裡面有多少件傢俬.
有三件.
第一件就叫陳設餅桌.
table of showbread.
其實那裡就有.
十二塊的麵包在那裡.
十二塊的餅在那裡.
每天都要去.
弄出來的.
等於你今天去bakery買東西.
每天早上廚師就會.
弄些新鮮的麵包出來.
每一天.
在曠野裡面行走的時候.
以色列人.
會幕在他們中間.
神的會幕在他們中間.
祭司.
他們就是預備這些東西.
然後十二塊的麵包.
為什麼是十二塊.
因為十二支牌是整個神的子民.
每天就會放在那裡.
就是那些.
陳設餅.
就是去remind.
去提醒我們.
原來上帝是住在.
整個神的子民的中間.
去供應我們一切的需要.
然後第二件傢俬就是.
燈台.
燈台是每一天.
會添加橄欖油.
然後就要.
長年去燃亮它.
就代表上帝是光.
是不是.
然後我們就見到一個香壇.

$^{361}$Incense altar.
那裡就是點一些香.
或者那些.
以前有的spice等等.
他們會點了它.
然後那些香上升的時候.
就代表祈禱.
是上到上帝的面前.
上帝是聆聽我們禱告的.
如果你這樣看這些傢俬的設計的時候.
其實都是很practical的東西.
其實那些香是隨抽的.
因為你有barbecue.
你有祭祀.
所以那個地方要衛生的話.
有很多東西就需要有銅盤有水.
然後有香.
但是那個神學的意味.
就是說你進到這個聖所的時候.
你面對這個上帝.
祂是光.
但是祂願意去聽我們的禱告.
然後過了這個聖所.
就有一塊的.
簾.
一塊的布.
你揭開這塊布然後走到裡面.
你就會見到一件東西.
一個箱.
就叫做約櫃.
Ark of the Covenant.
這個箱裡面裝了些什麼.
其中一樣就是那個法板.
上一次我們講到律法的時候.
有兩塊法板.
一塊代表上帝這個國王.
還有另外一塊代表以色列國.
他們大家是.
建立了這個.
聖約的關係.

$^{401}$在這個盒上面.
就有兩隻天使叫做基魯柏.
其實就是上帝的座駕.
上帝的座駕.
這些全部的.
建築的.
方法,藍圖.
這些都是古埃及.
法老王去出征.
周圍出征的時候帶著軍隊.
帶著他的子民.
他的帳幕都在他們中間.
其實用這個背景來跟他們講.
現在上帝.
不再是法老.
耶和華是你的上帝.
現在你不再屬於.
以前埃及那個.
欺壓你的君王.
現在是屬於一個跟你.
恩典裡面立約的上帝.
是一個愛的statement來的.
是一個愛的statement.
然後每一次,每年.
只有一日叫做Yom Kippur.
Yom Kippur大概九月,十月的時候.
贖罪日的時候.
只有大祭司才能進去.
這個致勝所.
他就由外院然後結證了之後.
來到去.
致勝所裡面.
然後代表上帝的子民.
去獻祭.
整個的每一個器具.
其實都是講到上帝去赦免我們的罪.
去供應我們一切需要.
去聆聽我們.
和在各樣東西裡面去引導我們.
是跟我們同在的.

$^{441}$這個上帝不是一個高攀.
不能的上帝.
並且他是.
下來我們中間,住在我們中間.
將他的家放在我們中間.
的一個上帝.
所以二十九章這裡講什麼呢.
他說我.
耶和華要住在以色列人中間.
作他們的上帝.
他們便知道我是耶和華.
我他們的上帝.
是將他們從.
埃及領出來的.
我要住在他們中間.
你有沒有留意這裡講什麼.
經常講到住.
我的家.
在你們中間.
所以出埃及的時候.
以色列人所經歷的不是純粹說.
我們得到自由了.
離開苦難.
是進入關係的過程.
在西太山的山腳.
他們見到這個回眸.
去到建立的時候.
一個藍圖的時候.
你可以想像摩西在西太山.
拿到這個藍圖.
好像現在那些architect.
拿到一張很大張的印出來的紙.
一看這個平面圖.
一看裡面所有的設計.
所有的尺寸,所有的細節的時候.
他們知道原來這位上帝.
在恩典裡面來愛我們.
那是第一部分.
第二部分.
在出埃及記載的回眸的故事.

$^{481}$就出奇地.
中間那裡就有一個.
32至34張.
講到另一個故事.
就好像作者沒有什麼.
心思那樣.
突然間講到很多藍圖的東西.
25至31張.
講到回眸的設計,尺寸,home depot.
那些藍圖的東西.
然後突然32張講另一個故事.
就是以色列人在西乃山腳那裡.
拜金牛的事.
這個簡直是反高潮.
因為開頭講到.
一個statement of love.
你想像一下你的愛人.
很愛你.
他買了一個豪宅給你.
然後就你進去住.
他先將那些相.
和那些藍圖給你看.
全部都是合你心意的.
正不正.
很正.
下一步應該是什麼.
下一步應該入住.
或者開工整他.
誰知這個作者.
出埃及記就插了.
三張下去.
講到他們拜金牛.
這個故事我們不詳談.
講到以色列人等摩西下山.
拿藍圖的時候.
他們等不及了.
他們就叫阿倫不如我們現在整一隻金牛出來.
因為我們不想要耶和華上帝.
我們等不及了.
我們想拜這隻金牛.

$^{521}$這隻金牛當然在當時.
是一個物件的.
偶像.
所以.
耶和華就很生氣.
但其實.
今天的故事來說.
這隻金牛都代表我們很多不同的.
所謂金牛.
我們會崇尚物質.
我們一切都以物質.
作為我們決定的.
法則.
或者我們一些的情慾.
一些引誘我們的東西.
我們每個人心裡都有金牛.
這個作者.
特登在一個會幕的.
藍圖之後插這個故事.
其實想講些甚麼呢.
想講些甚麼.
因為在32至34章講到.
拜金牛的時候.
耶和華應該以公義去殲滅他的子民.
因為想想.
他用他的恩手拯救了你們出來.
拯救了你們出來.
拯救你們脫離法路的苦難.
然後你來到這裡的時候.
我給你一個表達愛的聲明.
是不是我說我愛你.
我愛你們.
但居然你.
去怎樣.
你背弱.
離經背道.
所以在這個時候.
其實是整個《出埃及記》裡面一個很大的危機.
是一個危機.
因為如果是這樣的話.

$^{561}$整個上帝所做的救贖計劃.
就會game over.
會不會game over.
你看看32至34章裡面.
摩西就做了一個很重要的角色.
所謂一個中保的角色.
他就在上帝面前求上帝去赦免他們.
因為上帝說.
不如我將這些子民全部殲滅.
然後在摩西你身上.
我興起一個新的子民.
好不好.
那摩西就威風了.
摩西就說不要.
因為如果是這樣的話萬國會怎樣看我們.
就會覺得我們很廢.
特別是耶和華很廢.
怎麼可能會這樣發生.
所以他就求上帝.
去用憐憫.
去恩典.
去赦免他的子民.
結果就出了一句.
舊約裡面最重要的一句statement.
就是認識這位耶和華神.
是怎樣的神.
這個statement有34章6至7節.
上帝說耶和華是大有憐憫.
大有恩典的神.
不輕易發怒.
且有豐盛的慈愛和聖神.
為千萬人存留.
不是審判.
是慈愛.
赦免罪孽.
過犯和罪惡.
在舊約裡面.
這個是最大的statement.
因為結識這位上帝.
不單止是遙遙不可及.

$^{601}$他是住在我們中間.
他不是一個遠遠不能及的上帝.
他住在我們中間.
並且這個上帝願意.
以恩典.
以慈愛不輕易發怒.
來對待我們.
所以他寧願決意的.
去赦免他的子民.
住在他們中間.
然後故事繼續發展.
上帝就說.
金牛的事件之後.
你們得到赦免.
你們.
為我建造我的家.
然後才開工.
做這個會幕.
宏觀的時候就發現.
其實作者.
想用一些很細節的東西.
用當時人明白的語言.
去介紹.
他自己是一個怎樣的神.
三十五至四十章.
就是摩西下山後.
負責這些百姓.
然後給這個藍圖給他們看.
你看看上帝就是這樣愛我們.
上帝願意.
下來我們中間.
住在我們這裡.
百姓居然有這樣的邀請.
雖然是有罪.
剛才那首詩歌.
有一句歌詞.
很觸動我.
千瘡百孔.
我們都是一班千瘡百孔的人.
但是居然.

$^{641}$我們得到上帝的赦免之後.
我們得罪上帝.
但是上帝赦免我們.
他居然還說什麼.
你照我的藍圖.
建造我的居所.
在你們中間.
所以百姓就動手來建造這個會幕.
藍圖是很長的.
三十五至四十章.
但是那裡和之前.
二十五至三十章完全不同.
二十五至三十章就是你要做這個.
你要根據藍圖做這個.
三十五至四十章.
就是他們做了什麼.
就是照樣.
下命令.
然後他們執行那個命令.
他們就照樣.
照著上帝的藍圖和樣式.
來做每一件事.
舉例二十五章三十一節.
這裡是說藍圖的.
第一部分左邊.
耶和華吩咐.
摩西和他們的子民.
要做一個.
純金.
做出來的.
燈台.
上面有酒杯等等.
然後放那些油燈.
這個是七燈台的建造.
那個藍圖.
你要照樣這樣做.
然後金牛的事件三十二至三十四章.
到了會幕工程開播.
開展整.
幾乎是同一句說話.

$^{681}$只不過之前是命令式.
現在是敘述式.
三十七章第十七節.
他就做了.
即是那裡有一個工匠.
他就用純金.
做出燈台.
然後有什麼.
exactly所有的東西.
燈台連接成一塊.
即是金燈台.
不是薄銷行出來的.
整塊真金這樣.
你見到開頭是命令.
告訴你那個藍圖是這樣.
現在他們就實際的.
來執行出來.
所以二十五章.
三十一.
二十五章至三十一章.
還有三十五至四十章.
其實好像照鏡那樣.
後面的exactly是之前blueprint.
那個藍圖所說的.
幾乎是一模一樣.
上帝說了.
上帝這樣吩咐他們.
就是這樣去建造.
你說你這個整個的宏觀.
你就見到.
上帝是一個滿有恩典的上帝.
並且他願意去.
邀請我們.
為他建立居所.
這個statement是一個.
最大的statement.
教會.
初時以色列人.
都算是初期的教會.
我們教會的前身.

$^{721}$和新約的子民.
同一個的子民.
我們居然有.
這樣的救贖.
還有.
可以參與建立.
上帝的居所.
教會都是千瘡百孔.
教會.
都是很多.
錯誤的地方.
很多時候我們有一種.
錯誤的觀念.
你聽見很多人都這樣說.
或者我們有時候都會這樣想.
教會都是.
我整理得好一點我才來.
我不來教會.
為什麼?因為我覺得我太多.
太多污穢的東西.
讓我先搞定.
我才配來.
教會.
配來上帝的面前.
很多未信主的人都是這樣想.
甚至很多信了主很久的人都是這樣.
所以我們就有個殼.
我們想搞定那個殼.
起碼我們表面上.
我們覺得我們都是很光鮮.
我們將那些東西隱藏.
但其實我心想.
上帝如果光照我們.
任何一個角落的時候.
我們個個都是千瘡百孔.
是污穢的.
我們個個都是不配.
有誰可以.
配來到上帝面前.
沒有一個人.

$^{761}$為到時上帝在恩典裡面.
邀請我們我們才能夠參與.
正因為這個緣故.
我們的參與就更加有價值.
因為我們經歷過上帝的恩典之後.
我們就明白.
這個上帝那麼愛我.
我豈能不參與.
我豈能仍然坐在這裡.
所以.
這裡會幕的建造裡面.
就有一些東西講到.
那些以色列民他們甘心奉獻.
奉獻到一個地步多到一個地步.
摩西說夠了我們.
太多了.
一個金燈台要多少金.
去建造.
大概十字架那麼高.
你當它.
多少斤也好.
他們捐獻的.
是多到一個地步.
可以建造幾十個金燈台.
摩西說夠了.
因為已經足夠.
就因為他們.
被上帝的恩典.
觸動過.
他被上帝的生命觸動過.
當你唱那些詩歌.
想起那些經文的時候.
我是一個千瘡百孔的人.
我是一個污穢不得的人.
但上帝居然在高天來到我中間.
住在我的生命裡面.
住在祂的子民中間.
這個上帝值得我們讚美.
是不是.
難怪.

$^{801}$去到最後的時候.
上帝的榮耀就充滿.
這個回望.
這就是《出雅及記》最後一章.
最後那幾個字.
他說那時雲彩就遮蓋著.
回眸.
耶和華的榮光就充滿著.
祂的障目.
使得摩西都不能夠進去.
因為雲彩停在旗上.
耶和華的榮光.
充滿這個障目.
沒有人有資格.
讓上帝住在.
祂的中間.
千瘡百孔.
但是祂在恩典裡面.
邀請我們.
甚至去建造祂自己的回眸.
我知道我們教會.
我們都來到一個.
時間.
十年的故事.
我們都去到一個擴展的時候.
是不是.
但請你.
請你知道.
當我們呼籲我們一起.
建立教會的時候.
請你知道.
我們是一班怎樣的人.
我們是一班不配的人.
既因為我們不配.
上帝的恩典.
卻來到我心靈裡面.
很真實的.
來到我心靈裡面.
所以我們豈不能.
豈能夠只是坐在這裡.

$^{841}$不去參與.
我們領受過上帝的恩典.
領受過上帝愛的人.
我們被祂觸動過的人.
我們更加需要.
建立這個群體.
讓上帝住在我們中間.
這是多大的privilege.
多大的權柄.
多大的福分.
是我們不配.
Yet.
It comes to you.
來到我們心靈裡面.
教會從來不是.
為完美的人而設的.
但卻是為一班千瘡百孔的人.
污穢不堪的人.
而經歷過上帝基督的恩典.
我們便來到祂.
組成一個.
上帝能夠住到的地方.
整個脈絡裡面.
你會發現上帝是.
越來越靠近人.
靠近我們這一班.
不配的.
難怪到了新約的時候.
約翰福音第一章.
約翰去介紹.
耶穌的時候.
耶穌來到我們中間.
約翰福音第一章.
由頭講到尾.
到第十四節.
道就是耶穌自己.
成了肉身住在我們中間.
充充滿滿有恩典有真理.
我們也見過他的榮光.
正是父獨生子的榮光.

$^{881}$那個住在我們中間.
Dwelling.
May His dwelling.
Among us.
住在我們中間.
這個字其實就是會幕.
Tabernacle.
所以你想直譯的話.
He made His tabernacle among us.
Or He tabernacled.
among us.
他已住在我們這些.
千瘡百孔的人.
的生命當中.
住在我們裡面.
這個真是.
上帝的榮光.
真是上帝的榮光.
不是神跡歧視.
很spectacular 放煙花的人.
很實在的是上帝說.
雖然你不配.
雖然我不配.
但我願意.
去救贖你.
住在你的心裡面.
住在你的生命當中.
你會發現世界.
一個誤信的世界.
需要的是什麼.
不一定是需要一個.
大家很光鮮.
很完美的教會.
世界需要見到的教會.
就是我們這班.
我們可能不配的.
但上帝的恩典.
在我生命裡面.
成為真實.
我親自經歷過.

$^{921}$被觸動過.
上帝在耶穌基督裡面.
給我們愛.
我們說在這個journey裡面.
是整個的.
大故事.
去認識這位主.
今天我們邀請你.
出去的時候.
你拿一塊布條.
為什麼布條.
因為今天說到匯木.
匯木有很多不同顏色的布條.
你都可以造成.
上帝在恩典裡面.
邀請祂的百姓,祂的子民.
一起去建造祂的居所.
你拿一條回去.
放在我們的.
所謂十週年的瓶子裡.
我們有幾行的.
物件已經有了.
海沙.
這麼多的.
亞伯拉罕的.
祝福.
然後有什麼.
有一塊碎片.
我們說到族長故事的時候.
個個都是破碎的人.
但是我們邀請上帝.
去使用他們.
叫萬國能夠得福.
然後我們有一些石子.
說到上帝的律法.
因為有一個新的身份.
我們作為上帝的子民.
我們就按照他的律法.
他的原則去生活.
今天就是布條.

$^{961}$因為上帝住在我們中間.
我們都有份.
有這樣的恩典.
可以參與.
建立上帝同在的地方.
教會要有上帝的同在.
就需要我們以心以力.
擺上.
最後的時候我想邀請你.
這個星期.
去為教會祈禱.
不要只是聽說.
教會需要.
擴展.
需要走多一步.
需要go next.
要go next.
我們英文有句說話.
out of sight, out of mind.
看不到.
我們就忘記了.
但我們基督徒呢?.
out of prayer, out of mind.
我們不去禱告,不去紀念的時候.
我們就當我們是旁觀者.
但今天的故事告訴我們.
原來上帝邀請你每一個.
去建立他自己的居所.
一個星期.
做一個功課.
拿一條布條.
為教會去禱告.
為教會我們不只是需要地方.
為教會的生命.
我們去禱告.
我們在信,忘,愛當中.
我們都要自己成長.
You have to go next.
And I have to go next.
上帝的榮耀.

$^{1001}$在我們中間.
\newpage



\section{尼希米記 13:1-31}
\label{sec:032257}
\textbf{更新}~~~~王正傑牧師
\newline
\newline
連結: \href{http://www.church.com.hk/acms/content.asp?site=cdc\&op=show\&type=product\&code=032257\&layout=sermon}{\texttt{032257}} ~~~~ 語音日期: 2026-06-07
\newline
\newline
\hyperref[sec:032256]{\small{< < < PREV SERMON < < <}}
~
\hyperref[sec:index_chronic]{\small{[返順時目]}}
~
\hyperref[sec:032258]{\small{> > > NEXT SERMON > > >}}
\newline
\newline
尼希米記 13:1-31
\newline
\begin{longtable}{cl}
\hline
\hline
章節 & 經文 (和合本修訂版)\\
\hline
13:1 & \begin{tabularx}{0.7\textwidth}{X} 在那日，百姓聽到人朗讀摩西的律法書，發現書上寫著，亞捫人和摩押人永不可入神的會； \end{tabularx} \\ \\ \relax
13:2 & \begin{tabularx}{0.7\textwidth}{X} 因為他們沒有拿食物和水來迎接以色列人，卻雇了巴蘭詛咒他們，但我們的神使那詛咒變為祝福。 \end{tabularx} \\ \\ \relax
13:3 & \begin{tabularx}{0.7\textwidth}{X} 以色列人聽見這律法，就與所有不同族群的人分別出來。 \end{tabularx} \\ \\ \relax
13:4 & \begin{tabularx}{0.7\textwidth}{X} 在這之前，與多比雅結親的以利亞實祭司，受派管理我們神殿中的庫房， \end{tabularx} \\ \\ \relax
13:5 & \begin{tabularx}{0.7\textwidth}{X} 為多比雅預備了一間大屋子，就是從前收存素祭、乳香、器皿，和照例供給利未人、歌唱者、門口守衛的五穀、新酒和新油的十分之一，以及歸祭司之舉祭的屋子。 \end{tabularx} \\ \\ \relax
13:6 & \begin{tabularx}{0.7\textwidth}{X} 當這一切事發生的時候，我不在耶路撒冷，因為巴比倫王亞達薛西三十二年，我回到王那裡。過了多日，我又向王告假。 \end{tabularx} \\ \\ \relax
13:7 & \begin{tabularx}{0.7\textwidth}{X} 我來到耶路撒冷，才知道以利亞實為多比雅所做、在神殿的院內為他預備屋子的那件惡事。 \end{tabularx} \\ \\ \relax
13:8 & \begin{tabularx}{0.7\textwidth}{X} 我非常憤怒，就把多比雅的一切家具都從屋子裡拋出去。 \end{tabularx} \\ \\ \relax
13:9 & \begin{tabularx}{0.7\textwidth}{X} 我又吩咐人潔淨這屋子，然後將神殿的器皿、素祭和乳香搬回那裡。 \end{tabularx} \\ \\ \relax
13:10 & \begin{tabularx}{0.7\textwidth}{X} 我發現利未人當得的份無人供給他們，甚至供職的利未人與歌唱的都各奔回自己的田地去了。 \end{tabularx} \\ \\ \relax
13:11 & \begin{tabularx}{0.7\textwidth}{X} 我就斥責官長說：「你們為何不顧神的殿呢？」於是我召集利未人，使他們在自己的崗位上供職。 \end{tabularx} \\ \\ \relax
13:12 & \begin{tabularx}{0.7\textwidth}{X} 猶大眾人就把五穀、新酒和新油的十分之一送入庫房。 \end{tabularx} \\ \\ \relax
13:13 & \begin{tabularx}{0.7\textwidth}{X} 我派示利米雅祭司、撒督文士和利未人毗大雅作司庫管理庫房，副手是哈難；哈難是撒刻的兒子，撒刻是瑪他尼的兒子；這些人都是忠實的，他們的職務是分派他們弟兄所當得的份。 \end{tabularx} \\ \\ \relax
13:14 & \begin{tabularx}{0.7\textwidth}{X} 我的神啊，求你因這事記念我，不要塗去我為神的殿與其中的禮儀所獻的忠心。 \end{tabularx} \\ \\ \relax
13:15 & \begin{tabularx}{0.7\textwidth}{X} 那些日子，我在猶大見有人在安息日踹醡酒池，搬運禾捆馱在驢上，又把酒、葡萄、無花果和各樣的擔子在安息日扛入耶路撒冷，我就在他們賣食物的那日警戒他們。 \end{tabularx} \\ \\ \relax
13:16 & \begin{tabularx}{0.7\textwidth}{X} 有一些住在城裡的推羅人也把魚和各樣貨物運進來，甚至在耶路撒冷，在安息日賣給猶大人。 \end{tabularx} \\ \\ \relax
13:17 & \begin{tabularx}{0.7\textwidth}{X} 我就斥責猶大的貴族，對他們說：「你們怎麼會做這惡事干犯安息日呢！ \end{tabularx} \\ \\ \relax
13:18 & \begin{tabularx}{0.7\textwidth}{X} 你們祖先豈不是這樣做，以致我們的神使一切災禍臨到我們和這城嗎？你們竟干犯安息日，使憤怒越發臨到以色列！」 \end{tabularx} \\ \\ \relax
13:19 & \begin{tabularx}{0.7\textwidth}{X} 安息日前一日黃昏的時候，我吩咐人把耶路撒冷城門鎖上；我又吩咐，不過安息日不准開門。我也派幾個僕人在城門口站崗，免得有人在安息日挑擔子進城。 \end{tabularx} \\ \\ \relax
13:20 & \begin{tabularx}{0.7\textwidth}{X} 於是商人和販賣各樣貨物的人，有一兩次在耶路撒冷城外過夜。 \end{tabularx} \\ \\ \relax
13:21 & \begin{tabularx}{0.7\textwidth}{X} 我警告他們說：「你們為何在城牆前過夜呢？若再這樣，我必下手辦你們。」從此以後，他們在安息日就不再來了。 \end{tabularx} \\ \\ \relax
13:22 & \begin{tabularx}{0.7\textwidth}{X} 我吩咐利未人潔淨自己來守城門，使安息日分別為聖。我的神啊，求你因這事記念我，照你豐盛的慈愛憐憫我。 \end{tabularx} \\ \\ \relax
13:23 & \begin{tabularx}{0.7\textwidth}{X} 那些日子，我又看見猶太人娶了亞實突、亞捫和摩押的女子為妻。 \end{tabularx} \\ \\ \relax
13:24 & \begin{tabularx}{0.7\textwidth}{X} 他們的兒女，一半說亞實突話，或其他種族的方言，不會說猶大話。 \end{tabularx} \\ \\ \relax
13:25 & \begin{tabularx}{0.7\textwidth}{X} 我就斥責他們，詛咒他們，打了他們幾個人，拔下他們的鬍鬚，叫他們指著神起誓：「你們不可把自己的女兒嫁給外邦人的兒子，也不可為自己和兒子娶他們的女兒。 \end{tabularx} \\ \\ \relax
13:26 & \begin{tabularx}{0.7\textwidth}{X} 以色列王所羅門不也在這樣的事上犯罪嗎？在許多國家中並沒有一位王像他，蒙他神喜愛，神立他作王治理全以色列。然而，連他也被外邦女子引誘犯罪。 \end{tabularx} \\ \\ \relax
13:27 & \begin{tabularx}{0.7\textwidth}{X} 我們豈能聽憑你們行這一切大惡，娶外邦女子干犯我們的神呢？」 \end{tabularx} \\ \\ \relax
13:28 & \begin{tabularx}{0.7\textwidth}{X} 以利亞實大祭司的孫子，耶何耶大的一個兒子是和倫人參巴拉的女婿，我就把他從我這裡趕出去。 \end{tabularx} \\ \\ \relax
13:29 & \begin{tabularx}{0.7\textwidth}{X} 我的神啊，求你記得他們的罪，因為他們玷污了祭司的職分，違背祭司和利未人的約。 \end{tabularx} \\ \\ \relax
13:30 & \begin{tabularx}{0.7\textwidth}{X} 這樣，我潔淨他們，使他們脫離屬外邦人的一切；我又分派祭司和利未人的班次，使他們各盡其職， \end{tabularx} \\ \\ \relax
13:31 & \begin{tabularx}{0.7\textwidth}{X} 按定期奉獻木柴和初熟的土產。我的神啊，求你記念我，施恩於我。 \end{tabularx} \\ \\
[1ex]
\hline
\hline
\end{longtable}
$^{1}$弟兄姊妹平安 昨晚睡得好嗎.
為甚麼這樣問呢 因為今天這篇道比較尖竭.
你剛才讀經的時候 你讀了十幾節.
我看見十幾節都已經睡著了.
又在看很多名字 又這樣又那樣.
看《尼希米基》的確我們平時很少去看.
所以會比較辛苦.
我前兩星期有一晚睡不著.
我很少睡不著這麼嚴重.
試了聽歌都睡不著.
聽完歌見睡不著就出去聽.
打機等自己睡得著.
打完機回去睡都是睡不著.
當時我想起一個好朋友 但當時已經三點多.
我想起這個好朋友 於是我就決定去聽他的講道.
結果我聽他的講道之後 在第一個點都差不多講完的時候.
我就已經睡著了.
所以我非常明白聽道除了可以帶給我們很好的教訓之外.
也有教好訓的.
所以弟兄姊妹今天如果你真的氣數已盡的話.
可以竭力進入安息.
但如果你也想跟他學習的話.
我們也努力一起去看看《尼希米基》十三章對我們有什麼提醒.
其實今天我們要講的經文是十三章的全章.
但剛才我已經自錄了只是十四節給大家讀.
因為我知道讀完真的會睡得著.
所以只是給大家讀十四節.
大家都見證到威力.
所以我們一起透過我講.
希望能夠讓我們一起思想經文對我們有什麼提醒.
今天我們要講的是更新.
其實當我們要去看第十三章的時候.
我們要看前面第十二章發生什麼事.
《尼希米基》十二章描繪著一幅令人很振奮的圖畫.
當時是很歡喜快樂的畫面.
因為當時城牆建成了.
如果大家對《尼希米基》有印象的話.
你也知道《尼希米基》是在講什麼.
《尼希米》是在帶領以色列人去重建城牆.
當時城牆建成了.

$^{41}$所以百姓就有一個建城牆的獻城禮.
所以百姓一起聚集.
你可以想像畫面是師伯在歌唱讚美神.
也有祭司在吹號.
男女老幼大家一起歡喜快樂.
《聖經》甚至這樣記載.
當時的情況是怎樣.
耶路撒冷裡面的歡樂聲音是怎樣.
聽到遠處.
是他們歡樂到一個情況.
聲音傳到遠處都知道他們很歡喜快樂.
這樣的一個狀況.
你可以想像當時建好這個城牆.
當時是一個很歡樂的氣氛.
但當我們今天看到第十三章.
你看到整個畫面突然之間完全不同.
當我們看見的時候.
我們很驚訝地發現.
《尼希米》開始在講一些有關罪的問題.
是在處理百姓他們犯罪.
違背了神的律法的問題.
所以我們不禁要問.
城牆不是已經建好嗎.
獻城禮不是已經完成了嗎.
復興不是已經來到以色列當中了嗎.
其實對於當時的百姓來說.
他們所認為的復興是甚麼.
就是城牆重建了.
他們的耶路撒冷可以得到保護.
他們的生命有保障.
他們的生命有供給.
這個對他們來說是他們所認為的復興.
他們以為能夠完聖城牆的興建.
就是完成了一個很偉大的工程.
但今天我們看《尼希米記》.
神卻要讓我們明白.
真正的復興不單只是在講城牆的建立.
而是更加在講神在我們人生命當中.
重新去掌權.
這是神更加關心,更加重視的復興.

$^{81}$所以《尼希米記》的焦點.
從來不只是在講我們要重建城牆.
重建耶路撒冷.
而是在講重建神的百姓他們的生命.
外在城牆的重建.
只需要五十二天就能夠完成.
但內在的更新是一生之久的功課.
我們的生命是需要不斷地被神去更新.
以至我們的生命真的能夠得著復興.
可能我們和以色列人都很相似.
我們以為的復興都是一些很外在的.
就是以為復興.
譬如我們以為的就是.
我們在農華國搬到去榮輝新城.
我們有一個更新,一個更大的地方.
我們教會就會復興.
我們教會以為的就是教會多了人.
多些人回來,教會有增長.
那就是教會復興了.
我們也以為的就是.
教會不斷有事工.
一個事工又接一個事工.
很多聚會.
我們就覺得教會很復興了.
或者可能我們有聚會當中.
我們回來大家都很熱鬧.
我們又覺得很復興了.
又或者可能我們有些問題.
我們需要去處理的.
我們去解決到.
我們又覺得教會是復興了.
現在近來有很多問題我們要面對.
我們又要.
可能有些維修.
有些工程.
這一兩個多月.
我們就要面對冷氣的問題.
所以如果你九點回來.
你就會差點被迫死.
所以我們都在面對這些問題.

$^{121}$我們都會以為可能解決了這些問題.
可能教會就得到復興了.
但是神告訴我們.
原來神更加看重的復興.
不單止外在的事情去處理了.
去解決了.
見到外在的這些東西.
有一些增長的徵狀.
神告訴我們.
祂更加看重的復興是說.
究竟我們這一班神的子民.
我們屬神的人.
我們的生命究竟有沒有改變.
神究竟是不是有在我們生命裡面.
重新去掌權.
我們是不是一班都去順服神的話的人.
我們有沒有持續去.
順服神去更新我們的生命呢.
這個才是神更加看重的復興.
所以利未提出了三個更新.
第一個更新是分別為聖的生命.
第二個更新是.
敬拜和安息的生活.
更加更新的是.
順服神話語的心.
我們先去看經文.
利未去提出一個分別為聖的生命的更新.
當時經文第十三章.
一開始我們會見到.
百姓再次去聽律法.
我們就這樣看經文.
我們見不到有些什麼.
但經文其實表達的意思是.
百姓當他們再次去聽律法的時候.
他們才重新發現.
原來神對他們這班百姓.
這班神的子民是有要求的.
要求他們成為一個聖義的群體.
就是他們不要讓阿滿人和摩洛人.
去進入神的會.

$^{161}$弟兄姊妹.
這裡是在說什麼呢.
原來是在說神的吩咐.
他們是知道的.
神的吩咐他們是聽過的.
但他們似乎好像忘記了.
好像對神的吩咐輕忽了.
面對著這個律法的吩咐.
他們都有一個很好的回應.
他們就馬上與一切的閒雜人.
所以說其實不單止有外邦人在他們中間.
還有一些閒雜的人.
可能來自不同的國家也說不定.
在這裡要強調的.
他們不可以讓阿滿人和摩洛人進入神的會.
和要和這些閒雜人絕交.
不是在說一種種族隔離的問題.
而是在說一個信仰的分別.
然後經文交代了一個原因.
他說阿滿人和摩洛人當初沒有去迎接摩西和以色列.
還竟然僱用巴蘭去就坐以色列人.
神的介入使得就坐變為祝福.
如果你想了解這段歷史.
你可以回去看民數記22-24章.
阿滿人和摩洛人最大的問題.
不單止是他們不接受神.
更加他們是與神為敵.
但竟然今天以色列人容許這些與神為敵的人.
去進入神的會.
而當以色列人聽到這個律法.
他們就和那些住在耶路撒冷的外邦人分別出來.
他們不再被外邦人的影響同化.
但為什麼以色列人不可以和外邦人有來往.
最大的原因不是因為種族問題.
而是因為信仰的問題.
所以之後你會看到他們更加會說到通婚的問題.
和外邦人的問題.
所以我們看到大祭司以利亞實.
他竟然容許阿滿人多比亞.
這位當初極力反對他建造城牆的人.

$^{201}$去住在聖殿的庫房.
剛才你有沒有看到經文.
我不知道你讀的時候是否已經很努力地讀.
所以沒有特別看到經文的內容.
剛才經文說給我們聽.
讓外邦人住進去聖殿的地方.
那個房間那個大屋.
本身是用來裝什麼的.
大家有沒有留意.
是用來獻授祭祀的東西.
是用來獻祭工作的一切物品.
第一點我們可以想像得到.
以色列人特別是祭司.
他們竟然墮落到一個地步.
是容讓他們本來用來獻祭的地方.
現在給一個外邦人去住在裡面.
所以這裡不是單單說一個管理的問題.
更加是在說一個宋人界線失守的問題.
他已經沒有理會神的吩咐.
對敬拜更加輕率.
然後尼希米的回應是.
他立即去結證聖殿.
將他們的家具全部丟出來.
尼希米在這裡不是出於個人的喜好.
而是因為他知道.
神的殿必須屬於神.
所以屬神的人要和外邦人分別為聖.
以色列人他們違背文約.
尼希米就指出.
發生這事的時候.
尼希米是身處在皇宮.
過了多天.
當他回到耶路撒冷.
才知道多比亞這件事.
經文形容尼希米是怎樣.
他去到一個路不可結的程度.
你可以想像到當時發生了什麼事.
尼希米走轉背回到皇宮.
你不要以為回到皇宮很長時間.
只過了幾天.

$^{241}$這群人又怎樣.
又重蹈覆轍.
又繼續離棄神.
沒有去聽神的吩咐.
走開了幾天.
好像我幫女兒溫書一樣.
盯著她的時候.
她就很專心溫書.
走一走到洗手間.
她就在看手機.
好像是這樣的情況.
因為尼希米也看到.
更重要的一個問題.
為什麼尼希米會去到一個路不可結的情況.
如果你看尼希米的前面.
九至十章提到.
原來以色列人不單止聽到神的吩咐.
知道神的標準.
知道神的要求.
以色列人更加在認罪之後.
簽了公約.
但竟然當尼希米離開耶路撒冷回到皇宮.
他們很快就怎樣.
就故態覆文.
就重蹈覆轍.
所以弟兄姊妹.
聖戰當時的光景是怎樣.
說了不代表會做.
就算簽了約.
也不代表他們有去做.
但弟兄姊妹這種狀況像不像我們.
神的吩咐我們知不知道.
可能很多我們都知道.
說了給我們聽.
但不代表我們會做.
不代表我們會做.
就算簽了約.
我們可能都經常和神立約.
但我們最後都沒有做.
弟兄姊妹.

$^{281}$年初的時候我們很習慣.
很多人都很喜歡和神立約.
弟兄姊妹你今年和神納了什麼約.
是不是每一天都有靈修.
是不是每一天都有讀經.
最經典很多人都會出現.
就是有人納至每一年都去讀經.
結果讀到出埃及記的時候.
永遠都出不到埃及.
還是甚至弟兄姊妹去到今時今日.
做了基督徒這麼多年.
都已經不再和神立約.
免得最後做不到.
我認識一個弟兄.
這個弟兄本身上班經常遲到.
每天都遲到.
有一年他在年初的時候.
就說今年納至每天上班都要準時.
不再遲到.
結果那一年他真的全年都上班準時.
沒有遲到過一天.
但當到第二年的時候.
他又開始每天都遲到.
遲到到一個地步就被人解僱了.
我就問他.
你上年不是納至每天都不遲到嗎.
他就說是呀.
上年而已.
今年沒有想過要納至.
弟兄姊妹.
我們是不是都好像以色列人一樣.
神可能都告訴我們他的吩咐.
但不代表我們會做.
就算我們和神納了約.
但我們最後還是沒有去跟.
沒有去做.
所以弟兄姊妹今天當我們看見.
當以色列人容讓.
外邦人住在聖殿裡面.
更將憲制的物品拿走.

$^{321}$他們容讓了一些不屬於神的東西.
擺在了神的聖殿當中.
同樣的弟兄姊妹.
我們都要思想的是.
在我們的生命裡頭.
其實我們有沒有一些原本都屬於神的位置.
我們擺了一些不屬於神的東西在裡面呢.
那些東西是甚麼來的呢.
是不是錢.
是不是成就.
是不是享樂.
權力.
你的面子.
甚至自我中心.
可能還有其他.
弟兄姊妹今天我們在說更新.
我們真的要希望.
在我們的生命裡頭能夠得到更新的時候.
是我們需要讓神在我們生命裡居首位.
當我們真的讓神居首位的時候.
我們生命才可以得到真正的更新.
我們的生命的更新才會開始.
然後再進一步.
尼希米他回來了.
他發現一件事就是.
那些人離開了聖殿.
聖殿的服侍幾乎停頓.
原因很簡單.
百姓沒有奉獻.
每個人只顧著甚麼呢.
只顧著太平.
城牆建好.
我們安全.
我們顧著自己的生活.
我們忽略了神的家.
尼希米責備那些班長.
他們為何離棄神的殿呢.
因為對於當時的人.
神已經不再是優先.
奉獻停止.

$^{361}$事奉到停止.
敬拜變成很次要的東西.
在這裡不是單純因為錢的問題.
而是在說敬拜的問題.
說屬靈的問題.
說他們生命的問題.
在這裡我們要留心.
我們很想說的.
我們不要單單留意.
為何百姓沒有做十一奉獻這麼簡單.
我們更加要看見.
為何利未人離開了聖殿.
為何沒有供給.
他們就不去服侍.
為何這班利未人會容讓事奉停止.
容讓聖殿裡面所有的服侍都會停頓.
這個我們更加要問.
同樣的姐妹.
今天我們回教會.
不只是單單說有沒有十一奉獻.
究竟教會.
我們有沒有容讓它.
事奉到停止.
敬拜變得次要.
這個不是單單說錢的問題.
是更加說我們生命的問題.
接著尼希米又再去處理安息日的問題.
百姓他們在安息日去做買賣.
他們搬運貨物.
他們追逐利益.
表面上看來只是工作和休息的安排.
這個對於我們來說我們很容易理解.
很簡單.
為何沒有回來崇拜.
我上班.
我有事要做.
這個我們可以理解.
但對於當時尼希米.
為何他要特別去處理.
其實是反映這班以色列人的信心問題.

$^{401}$因為安息日的核心.
不只是不工作這麼簡單.
不工作不是最重要的一件事.
最重要的是他們的信仰.
他們相信神是他們生命真正的供應者.
他們不要讓工作成為了他們生命的主.
生命的神.
所以面對著以色列人他們去違背律法.
尼希米就很嚴厲地糾正他們.
首先他恢復正常性的正常運作.
他去指責那些官長,貴族,祭司.
問他們為何要離棄神的殿.
然後眾百姓恢復十一奉獻.
去招聚那些利未人使他們照舊去供職.
重新派利米亞等人去恢復聖殿的運作.
也要幫助那些以色列人能夠重新守安息日.
他指責那些貴族.
你們怎麼行者惡事犯了安息日呢?.
他更要去吩咐那些守城門的人.
要去關閉城門.
免得有人繼續在安息日裡搬運貨物.
甚至去做買賣.
弟姐們,尼希米記在此要告訴我們.
耶路撒冷的復興不是單單修補城牆就足夠.
更重要的是以色列人的心要歸回.
要恢復和神的關係.
弟姐們,今天我們也一樣.
教會不是單單修理好教會.
處理好一些事務就足夠.
更重要的是我們每一個弟姐們.
我們的心歸回.
我們和神恢復關係.
弟姐們,你們的心在哪裡?.
我們有沒有歸回神?.
有沒有和神恢復關係?.
我認識一位弟兄,一位年輕的弟兄.
為什麼特別要和大家分享他的見證呢?.
因為近年很難得見到一個弟兄.
對信仰很認真,也很有追求.
這位年輕的弟兄從小到大都很喜歡飛機.

$^{441}$所以他讀大學時也去讀飛機工程.
畢業後很順利成長,去了做飛機工程.
突然之間有一天他辭職,轉行做其他工作.
很多人一知道的時候都覺得很錯愕.
特別是他身邊的好兄弟.
為什麼他會辭職呢?.
因為他竟然放棄了他從小到大的夢想工作.
當問他之後發現原來.
香港做飛機工程的公司只有幾間.
每一間星期日都要開工.
其中一間本來星期日都不用上班.
但最後這間公司都轉為輪班工作.
令他星期日無法回教會崇拜和服侍.
所以這位弟兄決定辭職.
他希望星期日能夠回教會崇拜和服侍.
弟兄姊妹,和大家分享見證.
我很想鼓勵大家.
不要離棄神的電.
不要離棄神的事.
我們仍然要回到教會.
我們要去敬拜神.
我們要去服侍神.
所以我們需要更新我們的敬拜和安息日的生活.
然後,弟兄姊妹看到百姓還有一個問題.
當時他們與外族通婚的問題.
其實祭司都出了這個問題.
他的孫子做了外邦人的女婿.
所以他們的兒女出現了一個問題.
他們已經開始不懂得說猶太人的話.
這不是一種語言的問題.
而是一種傳承信仰的問題.
是一個很大的危機.
所以弟兄姊妹很緊張.
因為他很擔心一代一代逐漸失去了屬臣子民的身份認同.
所以他斥責在他們當中的異族通婚.
他斥責的做法是很激烈的.
當時他斥責那些娶異族女子為妻的人.
他就坐和鞭打他們.
又拔掉他們的鬍鬚.
免得他們再去異族通婚.

$^{481}$也趕走了大祭司爾利亞實的孫子.
耶和耶大兒子和倫人.
和倫人參巴拉的女婿.
趕了他出去.
參巴拉是誰呢?.
你看回尼希米的前面就知道.
參巴拉是率領外邦勢力去嘲笑和威脅和阻撓猶太人重建城牆的工作.
所以尼希米他不單止去責罰那些娶了外邦女子為妻的人.
更加連祭司的孫子也趕走了.
尼希米他用這麼嚴厲的方法去糾正以色列人.
可能對我們來說覺得何必這麼激烈.
但等姐妹們想深一層.
其實你可以明白為什麼尼希米要用這麼嚴厲的方法去糾正以色列人.
因為雖然今天縱使城牆修建得好.
但如果以色列人他再次繼續犯罪.
得罪神.
不去遵守神的吩咐.
最後他們的結果是怎樣?.
都是一樣.
他們要面對審判.
他們都是要被神去懲罰.
外邦仍然攻打侵佔他們.
所以結局是會一樣.
所以我們明白為什麼尼希米要用這麼嚴厲的方法去糾正以色列人.
因為他用這麼嚴厲的方法去糾正以色列人是一個冒著生命危險的.
因為他會得罪了祭司.
得罪了官長和貴州.
但在這個時代正正是需要一個對神忠心的人去撥亂反正.
為的就只是讓以色列人去恢復和神的關係.
但姐妹們今天我們去看這一部分.
我不是要和大家去討論是不是和未信主的朋友結婚的問題.
但在這裡我們可以看到更重要的是.
這一群以色列人他們會娶外邦女子為妻.
最重要的問題是什麼?.
是他們沒有再去聽神的吩咐.
沒有再去謹守遵行神的律法.
這個才是他們最大的問題.
而尼希米這麼激烈的反應.
在當時和以色列人對神的態度這麼散漫.
成了一個很強烈的對比.

$^{521}$她更加做了一個親身的示範.
就是她身處在一個外邦人當中的情況下.
她仍然是對神忠心.
仍然去謹守遵行神的話.
所以姐妹們.
很多時候不是我們不知道神的旨意或神的吩咐.
而是似乎我們好像漸漸淡忘了.
好像在《雅各書》那裡這樣說.
一章二十二至二十四節那裡說.
此時你們要行道,不要單單聽道.
自己欺控自己,因為聽道而不行道的.
就像人對著鏡子看見自己本來的面目.
看見走後,隨即忘了他的相貌如何.
等姐妹神的話就像鏡子一樣.
使我們重新去看見生命需要調整的地方.
但是經文提醒著我們.
我們不要看見走後,隨即忘了我們的相貌如何.
我們要聽道,更加要行道.
所以我們明白李希米很關心.
他關心的從來都不單是說究竟這個民族有什麼利益.
究竟他們生活是怎樣.
我們更加關心的是這群神的子民他們屬靈的生命.
他們屬靈的未來.
如果這個屬靈的信仰不能夠傳承的話.
城牆再堅固都是沒有用的.
所以今天各位弟兄姊妹.
我們教會都面對著很大的挑戰.
我們要留心的是.
我們不單止我們自己要相信神.
我們更加需要去幫助下一代認識神.
不只是我們的兒女要去認識神.
更加說的是一代一代的信徒都要認識神.
我們不單止自己聽見神的話.
我們更加要去活出神的話.
真正的更新不是知識的增加.
而是說的是馴服的增加.
弟兄姊妹,在這裡我想多說一點.
我們經常都發現我們要聽到然後行到很困難.
我們覺得勝經的要求很難做得到.
昨晚我在南洪團大團契.

$^{561}$我們在說要保守我們的口.
其中我們分享到要講誠實的話.
當我們講到這裡的時候.
我們發現很大的困難.
因為我們都知道要講誠實的話.
是否很容易?.
有些難.
所以我們就去想.
我們要不講一些就不講一些.
有些東西不講就OK.
因為我們覺得講誠實的話很難.
很難去遵守到神的吩咐.
照著神的吩咐去做.
所以我們就想很多我們認為可以可行的.
講一些不講一些.
不知道是不是這樣都可以.
弟兄姊妹,從來神的吩咐.
神的教導.
神的律法.
不是只是我們單單聽了然後去做.
然後我們發現做不到就放下.
重點是講究竟我們是否願意去馴服.
大家記不記得主耶穌在喀西瑪利園的禱告.
當時主耶穌帶了三個門徒上去.
三個門徒上去後馬上怎樣?.
和你旁邊那個差不多都是睡著了.
對不對?.
他們睡著了之後,耶穌就在那裡禱告.
耶穌的禱告是怎樣?大家有沒有印象?.
耶穌是不是在那裡禱告說.
我很開心,我預備好了.
我要去完成神的吩咐,我要去承受救恩.
是不是這樣禱告?.
耶穌的禱告是怎樣的?.
求天父怎樣?.
拿走祂的富貝.
然而不要照祂的心意,要照天父的心意.
主耶穌的禱告是在講甚麼?.
在講馴服.
所以今天我們見到神的吩咐很難做.

$^{601}$我們見到神要我們做的事好像做不到.
重點不在於這件事有多難做,多做不到.
而是在於我們有沒有去馴服.
當我們馴服,我們就鬆心去完成神的吩咐.
教會還有一個挑戰就是.
不是我們要去知道多少聖經.
而是我們願意去進行多少聖經.
這是我們一直都明白的地方.
而最後,尼希米他用了他的生命.
來給我們看見.
當一班已逝的人,他們這麼墮落.
當他們的信仰變得這麼腐敗的時候.
他怎樣用他的生命,用他的見證來告訴我們.
他身處在外邦人當中的人.
他怎樣仍然對神忠心,謹守遵行神的吩咐.
尼希米他在十三章出現了四次的禱告.
這是尼希米記在傳說裡面最後的聲音.
我們看尼希米他沒有去誇耀自己有什麼成就.
也沒有去提及城牆有多偉大.
也沒有去講自己有什麼功勞.
他只是將自己去交託給神.
尼希米他自己成為一個榜樣.
他給我們看見他對神忠心.
對神忠心不是在乎處境或者際遇.
而是在於他在困難的時候仍然能夠信靠.
仍然去依靠神的幫助.
同樣弟兄姊妹,今天我們需要的是忠心.
忠心不是在說處境和際遇.
而是在於當我們面對困難的時候.
我們仍然信得過神,去依靠神的幫助.
當你覺得神的吩咐是很困難的時候.
我們不是選擇去逃避.
不是選擇去疏忽或者淡忘.
而是我們要求神去幫助.
我們仍然信靠得到神,然後依靠神的幫助.
尼希米更加告訴我們對神忠心的人.
而是求神的悅納.
所以弟兄姊妹,我們對神忠心.
從來不是因為我們獲得弟兄姊妹之間.
有很多的欣賞或者讚賞.

$^{641}$而我們更需要的是追求神的悅納.
十三章表面看來不是在說復興的問題.
但它是真的在提一個更深層的復興.
這個深層的復興是在說.
因為神不只是修補破損的城牆.
他是要更新百姓的生命.
所以弟兄姊妹,今天我們的生命.
需要被神去更新.
我們要被分別為性.
我們不要做和外邦人一樣的事.
和外邦人一樣的價值觀.
我們需要回到神的殿.
做好神的事.
也需要重新學習去馴服神的話語.
我們要有一顆馴服話語的心.
和大家分享一個見證.
大家都很認識的德蘭修女.
德蘭修女是十七歲的時候.
她回應神的呼召.
她去到印度的加爾各答去做老師.
原本做老師是可以過著很安穩的生活.
但當她走進去這個貧民窟的時候.
她就見到很多病倒在街頭.
被社會遺棄的人的時候.
她就決定離開這個安穩的生活.
離開她的工作.
她要去全心去服侍那些最貧窮的人.
所以在1950年的時候.
她就創立了人愛傳教修女會.
設立一個垂死病人的收容所.
也去照顧一個麻瘋病人的麻瘋病人中心.
也親身去照顧那些沒有人願意接近的人.
所以她曾經這樣說.
最痛苦的貧窮是當你感受不到別人的關愛.
和覺得被遺棄.
這個就是最痛苦的貧窮.
因此她不單去照顧別人的身體.
她更加用愛去醫治別人的心靈.
去到1979年她獲得諾貝爾和平獎.
但她將她所得到的獎金全部都捐出來.

$^{681}$繼續去幫助那些窮人.
她謙卑地說.
我只是神的僕人.
用我微小的成就她的偉大.
而到了德蘭修女晚年的時候.
近年我們會看到一些她的日記.
她心裡面的一些說話.
德蘭修女在晚年的時候.
曾經長期經歷靈性一個很黑暗和孤單的時間.
甚至她說她感受不到神的同在.
但即使如此.
她從未停止過服侍.
她依然每天都去中心照顧病人.
關懷那些弱勢.
將愛帶給那些最有需要的人.
她最後這樣說.
我們不能做偉大的事.
但我們可以用神偉大的愛來做每一件小事.
弟姐妹和大家分享的見證就是.
我們看見德蘭修女的生命的中心.
她聽到神的吩咐.
她照著神的吩咐去做.
並且在她最困難的時候.
她靈性經歷低潮的時候.
她仍然忠心地去完成神的吩咐.
去堅持關心照顧那些病人.
那些貧窮的人.
弟姐妹.
今天我不知道大家身處在一種什麼的境況.
若我們真的來到神的面前.
我們不單單追求外表的復興.
我們更加追求我們生命靈性的更新.
以致我們真的能夠見到.
不單單是我們生命被復興.
我們更加見到教會被復興.
上星期我去了港島.
我去了屯門.
那間教會在屯門志樂花園.
弟姐妹你認不認識志樂花園在哪裡?.
有沒有去過?.

$^{721}$我有去過.
我第一次去.
它在三聖村再入一點.
去到志樂花園的時候.
我看到是一個很低密度的屋苑.
傳道人去到一個新的屋苑.
第一時間會留意什麼呢?.
就是看看有多少間教會在那裡.
我一看就看到有四間教會在那裡.
我就想一個小小的屋苑.
就有四間教會在那裡.
明白就是傳道人第一件事都會想的就是.
四間教會很難做.
很多競爭.
雖然表面上沒有競爭.
但都可能暗地裡有些競爭.
當時我就有這些思維.
當我去到教會坐下來靜下來的時候.
我有一個感動.
其實我想到的就是.
其實神既然放了這四間教會在這裡.
在這個志樂花園這麼小小的屋苑.
其實我相信神真的一定有祂的心意.
所以我在港島的時候.
我鼓勵弟姐妹.
我說弟姐妹不要只是看到.
現在教會要去服事的困難.
要怎樣去祝福這個屋苑呢?.
好像都摸不著頭腦.
都不知道怎樣入手.
我鼓勵他們要擁抱這個屋苑.
不要只是想帶一兩個人信主就已經很開心.
我覺得要擁抱這個屋苑.
真的要祝福整個屋苑.
無論是教會自己有努力去做也好.
連同其他教會也好.
一起努力去做.
真的要去祝福這個屋苑.
所以我鼓勵他們去祝福志樂花園.
弟姐妹同樣.

$^{761}$我和你分享這個感受.
其實我很想鼓勵大家.
弟姐妹.
神放我們在榮光街榮輝商場裡.
神有祂的心意.
20年前.
哪有人想到.
土瓜灣會有這樣的慘劇.
會塌樓.
接著整個社區很多舊樓需要去遷拆.
重建.
以致我們現在望出去.
整條榮光街都是地盤.
弟姐妹.
我相信神一定有祂的心意.
我鼓勵我們整間教會的弟姐妹.
一起去擁抱這個社區.
不單單只是擁抱榮光街的新樓.
更加擁抱的是整個社區.
我們真的能夠祝福這個社區.
我們期望神真的去復興這個社區.
而這個復興是來自我們每一個.
我們生命願意去更新的人.
所以總結當我們今天要去講復興.
我們生命要去更新.
弟姐妹我們經常聽人說.
減肥是他畢生的事業.
弟姐妹我很盼望.
更新都成為我們畢生的事業.
我們都去追求我們生命不斷被神更新.
神今日都同樣向我們發出呼召.
不要只滿足過去的經歷.
也不要只停留在昔日那種感動.
也不要因為已經信了主很多年.
覺得自己已經很足夠.
神仍然要更新我們.
我們生命仍然要成長.
我們仍然要有一個好的見證去榮耀神.
所以盼望弟姐妹.
讓我們真的來到神的面前.

$^{801}$今日我鼓勵你.
你可以問神.
究竟你生命需要怎樣去更新.
被神去更新你.
讓你真的能夠成為一個分別為勝的生命.
你重新被神去更新你敬拜和安息的生活.
你更加能夠有一個順服神話語的心.
以致你真的能夠對神有忠心.
在極困難.
在你覺得最無從入手.
做不到的時候.
你仍然能夠靠著主.
可以踐行真理.
可以活出神美好的見證.
弟姐妹.
你願不願意呢.
請同起來.
今日和大家講這篇道.
可能是我回來土鎮這麼久.
最長氣的一篇道.
但今日這篇道.
其實花了大概九個字的時間.
很簡單.
其實講一樣東西.
就是弟姐妹的心要歸向神.
親愛的主耶穌.
主今日當我們來到你的面前.
我們懇求你.
懇求你幫助我們.
拿走我們很多的軟弱.
我們很多的疲乏.
甚至可能我們有很多的憤怒.
主我們來到你的面前.
願你的寶血再一次去潔淨我們.
一切的罪.
讓我們能夠可以坦然無懼的來到你的私人寶座前.
讓我們今日都開放我們自己的生命在你的面前.
我們願意去被你更新.
幫助我們.
讓我們不要被這個世界去影響.

$^{841}$不要被這個世界去同化.
我們能夠可以分別為聖.
我們更加像你.
我們亦願意懇求主.
你更新我們敬拜和安息日的生活.
幫助我們.
我們不要離棄你的殿.
讓我們顧念神的殿.
更加顧念神的事.
你都堅定我們的信心.
幫助我們.
讓我們在無論任何的環境.
任何的狀況.
我們都能夠忠心地去跟隨你.
讓我們能夠學會馴服你的話語.
以致我們真的能夠可以去聽道行道.
踐行真理.
能夠可以成為你美好的見證.
所以主我們今日就求你幫助我們.
祝福我們每一個弟兄姊妹.
我們的生命都可以被你更新.
以致我們真的能夠可以成為我們身邊人的祝福.
成為身邊人能夠從我們身上去看見這個美好的見證.
主我們就懇求你祝福幫助我們.
使用我們每一個.
主我們多謝讚美你.
我們這樣禱告.
是奉靠主耶穌.
基督的聖名自祈求.
阿們.
\newpage



\section{士師記 4:4}
\label{sec:032258}
\textbf{聖經人聖經事:第三位士師  底波拉 巴拉}~~~~黃克勤牧師博士
\newline
\newline
連結: \href{http://www.church.com.hk/acms/content.asp?site=cdc\&op=show\&type=product\&code=032258\&layout=sermon}{\texttt{032258}} ~~~~ 語音日期: 2026-06-06
\newline
\newline
\hyperref[sec:032257]{\small{< < < PREV SERMON < < <}}
~
\hyperref[sec:index_chronic]{\small{[返順時目]}}
~
\hyperref[sec:032259]{\small{> > > NEXT SERMON > > >}}
\newline
\newline
士師記 4:4
\newline
\begin{longtable}{cl}
\hline
\hline
章節 & 經文 (和合本修訂版)\\
\hline
4:4 & \begin{tabularx}{0.7\textwidth}{X} 有一位女先知底波拉，是拉比多的妻子，當時作以色列的士師。 \end{tabularx} \\ \\
[1ex]
\hline
\hline
\end{longtable}
$^{1}$各位電視節目大家好.
我們一起讀聖經.
生命樹宣教網絡很想鼓勵大家多看聖經.
來到聖經人聖經事這系列已經說了很久.
今次是講第三位事師.
其實第三位主角是底波拉.
她是一個助手.
巴拉.
底波拉是一位女士 女事師.
很厲害 得高望重 大家很尊敬.
她是管理士 文人.
要打仗.
她有一位助手.
神叫她找巴拉.
看看聖經人 事師和神怎樣做這件事.
打退了迦南人.
我們繼續.
士師記來到第四章.
大家知道.
士師記有一個惡性循環.
第一.
以前的事師死了.
太平了一輪之後.
以色列人又開始頑皮了.
第一.
以色列人又行耶和華.
眼中看為惡的事.
惡的事通常不是行為上惡.
最主要是他們去拜迦南地的神.
一拜巴力 神就很氣憤.
巴力統稱迦南地各個民族的神.
拜的時候會有廟祭一起淫亂.
最差的會獻上人.
或者獻上自己的兒子做祭.
神是很討厭這些.
行耶和華眼中看為惡的事.
於是有第二.
耶和華會將他們賦予其他外族人.
用皮鞭鞭他們.
這次在以弗左手的事師死後.

$^{41}$太平了一段時間.
又行耶和華眼中看為惡的事.
於是神就將他們交給夏所地的迦南王.
夏所地的迦南王.
迦南王只是一個統稱.
因為是迦南地一個頗為重要的本地王.
他統領了很多人.
很多部族聽他講.
所以統稱迦南王.
他自己住在夏所.
他的名字叫耶賓.
他的將軍叫西西拉.
西西拉對於耶賓很忠心.
西西拉才厲害.
叫士石拉.
也可以讀做士石拉.
士石拉的意思是戰鬥陣式.
這個士石拉很有名.
主要是他很會用陣式來打仗.
打仗很講戰陣.
能夠贏輸在乎戰陣.
特別以往在乎戰陣戰術.
住在外邦人的夏羅切這個地方.
有兩個人物.
一個是迦南王耶賓.
耶賓做王.
不過他手下有一個很厲害的將軍.
士石拉.
西西拉.
這個才重要.
戰陣很厲害.
看來他最厲害的是用戰車.
耶賓王.
耶賓是什麼意思呢?.
也賓.
也賓是神所看見的人.
這個名字很靈性化.
是不是耶和華看見的人?.
不是.
神所看見.

$^{81}$因為他信迦南帝的巴力神.
所以他的神是指巴力.
是神認識我的看見的人.
就是神都看著我.
所以叫耶賓.
他有鐵車九百輛.
這個士石拉最厲害的.
就是用鐵戰車的戰陣.
九百輛戰車.
過千匹馬.
很多步兵.
因此很厲害.
大大欺壓以色列人.
二十年.
是大大欺壓以色列人.
二十年這麼久.
一百多萬人都不夠以色列人來.
因為大家不集中.
以色列人這個時候又來了.
第三點就是.
呼求耶和華神.
救我們.
跟著你的應許來幫我們.
很需要你幫我們.
求了二十年.
神聽你祈禱.
你們是這樣的.
心又軟了.
於是又興起一位仙子.
今次這位女仙子.
這個士司叫底波拉.
這個底波拉本身是女仙子.
即是主要是講上帝的說話.
仙子是甚麼呢.
仙子應該叫代言人.
代表神講說話.
她是講神的說話.
講神的誡命.
帶領大家敬拜主.
叫止人們.

$^{121}$去敬拜神.
這個是女仙子.
神跟我們說甚麼.
這個是女仙子.
不是祭司.
祭司是代表人做.
向神.
仙子是講的.
講出神的說話.
所以仙子的意思是神的代言人.
就叫底波拉.
Debra這個好名字.
Debra.
Debra是蜜蜂.
很有趣.
她的名字叫蜜蜂.
可能出生的時候小小的.
所以爸爸媽媽就叫她蜜蜂仔.
所以就叫她蜜蜂.
而當時早期的人.
很多時候都用小動物的名字來做她的名字.
Debra是蜜蜂.
她丈夫叫拉比多.
當時是以色列的士司.
以色列的士司不是講拉比多.
是比底波拉.
又是仙子.
但她又是士司.
記得了.
士司是做甚麼的.
士司我們英文就翻譯為Judges.
士司就是Judges.
就是法官的判斷事情.
她就住在以法連山地拉馬和伯特利中間.
這個地方.
都在整個以色列由南到北的中間.
她就在底波拉的中樹下.
她就在自己家裡.
在中樹下幫別人做判斷.
這個就是士司的工作.

$^{161}$士司不是仙子.
仙子就是將整個說話講出來.
底波拉她本身又是仙子.
但她同時間又被推舉為士司.
大家都去找她.
不如你和我們做法官.
我們有甚麼事.
就來找你幫我們判決.
以色列人都想他那裡去聽判斷.
聽判斷之後.
將事情告訴她.
她就決定這樣不對.
這樣不對.
你就坐一步.
不如大家一人走一步.
這樣吧.
這件事你真的不對.
你要做這個賠償.
她就判斷.
其實都挺忙的.
這個時候.
她就打發人.
這個應該是上帝感動她.
你去幫我找一個人.
於是她就打發人.
從拿忽他利的基底師.
張亞比拿閹的兒子.
巴拉.
巴拉克.
巴拉克是雷電的意思.
可能那個人很神勇.
所以他的名字叫巴拉克.
雷電.
他就叫人.
叫雷電來.
照料來.
其實你要知道這位女士.
在當日是男性為主的裡面.
如果不是上帝給她.
她是做不到事事的.

$^{201}$所以很少女士師女先知.
你見到很少.
因為是男性為主.
她竟然打發人.
叫巴拉過來.
去拿忽他利另一個支派的基底師.
叫了巴拉過來.
為甚麼?.
對她說.
耶和華以色列的神吩咐你.
你率領一萬拿忽他利和西布倫人.
去他博山.
以他博山是後來.
到主耶穌時代登山變象的地方.
你是金人以色列地的中心點之一.
好,你去那裡.
帶一萬人去.
主要來自你自己拿忽他利和西布倫兩個支派.
今次神叫你帶領兩個支派的人.
一萬人去對付耶賓和西西拉.
耶和華說.
我必使耶賓的將軍西西拉.
率領他的車輛全軍去基順河.
去北面一點.
基順河比較北一點.
接近加密山那邊.
到你那裡去.
我必將他交在你手中.
你去那裡.
我會將他交給你.
你會贏他的.
留意紅色字.
你看不到為何這個南的能徵慣戰.
帶兵的名叫雷電的男士.
竟然自己不敢.
自己不敢是怎樣.
他就說你若同我去我就去.
你若不同我去我就不去.
扭計.
叫我帶一萬人去和耶賓西西拉打仗.

$^{241}$你跟我一起去.
其實巴拉在想.
我是懂得打仗的.
可以的.
不過你和我去我心定很多.
為甚麼呢.
因為如果你去.
我就更加知道神與我們同在.
耶和華與我們同在.
因為祂透過你說.
祂一定會應驗祂的同在.
祂一定會照自己所應許.
有時你跟我一起去.
我的心定很多.
不是說你很好打.
不是說你很懂得戰爭.
而是你跟我去.
就表示到耶和華與我同在.
我就很定了.
他就明白了.
帝波拉就說我必要你同去.
不過.
你這次的路得不著榮耀.
因為耶和華要將西西拉.
交給一個婦人守女.
我和你一起去沒問題.
不過你要知道.
這次既然你這樣說.
其實你會將你的榮耀.
分給另一個人.
另一個女的.
一個婦人守女.
這裡有些難明.
一個婦人守女.
究竟是指帝波拉.
你本來自己帶兵去.
你得到所有榮耀.
這次你叫我去.
你的榮耀就有部份給了我.
變了我帶.

$^{281}$你的榮耀就不太好了.
其實沒辦法.
你要我去我一起去.
但如果看經文.
一個婦人可能是另一個婦人.
這是預言.
帝波拉就說我和你一起去.
不過進一步.
我覺得神告訴我.
這次你去.
你是不得著所有榮耀.
是有另一個婦人會得著榮耀.
會分榮耀給她.
是嗎?.
不要緊.
於是帝波拉起來.
就和巴拉一起去基底斯.
巴拉招聚西布里人.
拿忽塔利人到基底斯.
和他一起的有一萬人.
帝波拉也一起去.
大家氣勢如虹.
摩西鄂府荷巴的後裔.
基尼人希伯曾離開基尼族.
靠近基底斯的.
撒拿欽的象船旁邊之搭象棚.
有人告訴西西拉.
阿比拿欽的兒子巴拉已經去他博山.
預備和你打仗.
西西拉說好.
既然如此.
我就聚集我的鐵車九百輛.
跟隨他的全軍.
從愛邦人的夏羅切出來.
到基純河迎接他打.
帝波拉對巴拉說.
你起來吧.
今天是耶和華把西西拉交付你的日子.
你要贏了.
耶和華豈不在你前頭行嗎?.

$^{321}$這句說話真是很安慰.
請你記得.
很多時候你要去應付一件很大很難的問題的時候.
你要知道.
今天耶和華將敵人交給你.
而且豈不在你前頭行嗎?.
你在前頭行的時候.
其實神就跟你一起去.
而且祂跟你開路.
放心吧.
神會幫你.
於是巴拉下了他博山.
跟隨他的一萬人.
下山.
預備打仗了.
耶和華使西西拉和他一切車輛全軍賄戀.
這個你看到神的工作了.
其實是誰好打?.
是不是西西拉好打?.
都可以打的,不過.
是耶和華使西西拉和他一切車輛全軍賄戀.
本來西西拉最出名的是什麼?.
他的名字叫戰陣.
我的名字就是戰陣.
我最擅長用陣式打仗的.
誰知道全軍賄戀就是他了.
你最擅長的那件事.
神一撥,你的東西就散了.
於是在巴拉面前被刀殺敗.
西西拉他那九百輛鐵戰車不要了.
亂了路,人們下車暴行.
你知道那些坐慣車的人下車暴行就不行了.
不能打了.
他是機械兵團.
現在機動的.
現在下車暴兵就很差了.
巴拉追趕戰車兩軍隊直到外邦人的下落前.
他們走回自己家.
巴拉帶著軍隊追趕到他家.
西西拉全軍覆沒.

$^{361}$刀刀在刀下沒有留下一人.
很慘了.
只有誰?.
只有西西拉自己一個人暴行逃跑.
他也洩露.
到了基尼人希伯之妻瓦益的帳篷.
留意這個瓦益.
剛才神說你的榮耀會分一部分給一個婦人.
原來是分給瓦益.
他走到瓦益的帳篷.
因為夏娑王爺賓和基尼人希伯家是好朋友.
走到希伯家的太太家.
瓦益出來迎接西西拉.
對她說請我進來不要怕.
西西拉就進了她的帳篷.
瓦益說我找些東西遮住你.
一會兒別人追來找你.
我會說你不在的.
遮住先進來.
西西拉對瓦益說.
瓦益我跑了很久.
很渴給我水喝.
瓦益說沒問題.
喝水不夠.
你要補充體力.
打開了皮袋給瓦益喝.
繼續躲起來我找些東西遮住你.
你休息一下.
睡一下.
補充體力.
西西拉對瓦益說.
你站在門口.
如果有人問你有沒有人在.
你就說沒有.
有沒有見過西西拉來過走過.
沒有.
他們相信你.
你說沒有.
誰知上帝工作了.
西西拉疲乏沉睡.

$^{401}$希伯的妻子瓦益怎樣做.
留意一下.
取了帳篷的缺紙.
什麼叫缺紙呢.
當你打開帳篷.
除了有些木柱頂住不塌之外.
你一定要將帆布拉緊.
然後釘在地上.
然後拉些長些的繩子.
拉些長些的繩子.
再釘在地上.
四圍拉著帳篷.
於是風吹.
無論哪方面吹來.
都可以扯著帳篷.
不會塌的.
那些我們叫做搭露營的營釘.
那些就叫做缺紙.
是有個勾或一個洞綁繩子.
勾著繩子釘在地上.
斜斜地倒轉地釘.
於是拉緊帳篷.
那叫做缺紙.
很長的.
沒有這麼上下一尺長.
是不夠插得實在地上.
而且那裡是沙的地方.
要很長.
不然就不夠實力.
誰知想不到這個啞翼.
手裡拿著鎚子.
拿著一條營釘.
然後靜靜地走到.
疲乏沉睡的西塞拉旁邊.
將那條營釘缺紙.
它睡了.
在它發笨的地方.
一鎚出力釘下去.
穿了進去釘進地裡.
它在地上睡著.

$^{441}$釘了一會兒.
可能釘一會兒會動一動.
接著它打了幾下.
腦袋受傷不能動了.
它就直接打穿了頭.
釘進地裡.
西塞拉就死了.
巴拉追趕西塞拉的時候.
啞翼就出來迎接它.
來吧.
我將你所要尋找的人.
我已經將他解決了.
巴拉馬上進去帳篷.
看到西塞拉已經被殺.
倒在地上.
缺紙還在它笨中.
釘了進去沒拔出來.
我已經將他解決了.
所以巴拉的榮耀.
分了部分給婦人啞翼.
當然也分了部分給狄底波拉.
記住是神使迦南王耶賓被以色列人制伏.
於是以色列人的守禦法有力.
勝了迦南王耶賓將他滅絕.
他們又開始一段和平的日子了.
有一個惡性循環.
恆爾說恆爾惡的事.
神就用外族人鞭打他.
很慘啊.
被奴役二十年.
於是有求神.
神就給他一個事司.
事司就將他解救了.
他又和平了.
過沒多久他又來過.
我們人都是這樣的.
會不停好一點.
我們又會離開神.
鬆弛了.
所以有時候神都要給我們一點困難.

$^{481}$我們就會練得信心更加好的.
最主要就是不論哪個時代有信心.
今次是幾個人一起合力.
將外族人打退了.
打退魔鬼的作為.
你想不到一個女的也可以做事司.
做先知.
做得很好.
而且神在他身上.
他就叫巴拉.
巴拉就說好啊.
我去啊.
其實神除了用巴拉.
用底波拉.
還用了瓦益三個人.
一起為以色列帶來和平.
讚美主.
向我們一起禱告.
基本上我們雖然見到事司記.
好像不停地重覆.
但我們知道其實是有心意的.
每一次的循環.
都表示著你對以色列人.
對我們那種愛那種關心.
雖然我們做了很多不好的事.
行爺禍華.
眼中漢為惡的事.
你已經氣到不得了了.
所以你用外邦人成為一條鞭子.
雞毛掃打他們.
令他們醒過來.
向你哀求再回轉.
回轉之後你就救他們.
每一次都是這樣.
一個慈父一個慈母.
但有時是炎父又是炎母.
在這中間我們見到.
其實對以色列人.
無論怎樣不離不棄.
是不會捨棄他們的.

$^{521}$就會為他們當中.
興起一些不見經傳的人.
成為他們的領袖.
成為他們的英雄.
帶領他們脫離外邦人的手.
讚美你.
你的手在聖經的人.
聖經的事後面.
不停地工作.
做了很多事.
彰顯你的榮耀與作為.
願你得著讚美.
祈禱奉耶穌基督的名字.
阿們.
\newpage



\section{詩篇 126:1-6}
\label{sec:032259}
\textbf{耶和華為我們行了大事}~~~~羅錫堅牧師
\newline
\newline
連結: \href{http://www.church.com.hk/acms/content.asp?site=cdc\&op=show\&type=product\&code=032259\&layout=sermon}{\texttt{032259}} ~~~~ 語音日期: 2026-06-09
\newline
\newline
\hyperref[sec:032258]{\small{< < < PREV SERMON < < <}}
~
\hyperref[sec:index_chronic]{\small{[返順時目]}}
~
\hyperref[sec:032260]{\small{> > > NEXT SERMON > > >}}
\newline
\newline
詩篇 126:1-6
\newline
\begin{longtable}{cl}
\hline
\hline
章節 & 經文 (和合本修訂版)\\
\hline
126:1 & \begin{tabularx}{0.7\textwidth}{X} 當耶和華使錫安被擄的人歸回的時候，我們好像做夢的人。 \end{tabularx} \\ \\ \relax
126:2 & \begin{tabularx}{0.7\textwidth}{X} 那時，我們滿口喜笑、滿舌歡呼；那時，列國中就有人說：「耶和華為他們行了大事！」 \end{tabularx} \\ \\ \relax
126:3 & \begin{tabularx}{0.7\textwidth}{X} 耶和華果然為我們行了大事，我們就歡喜。 \end{tabularx} \\ \\ \relax
126:4 & \begin{tabularx}{0.7\textwidth}{X} 耶和華啊，求你使我們這些被擄的人歸回，好像尼革夫的河水復流。 \end{tabularx} \\ \\ \relax
126:5 & \begin{tabularx}{0.7\textwidth}{X} 流淚撒種的，必歡呼收割！ \end{tabularx} \\ \\ \relax
126:6 & \begin{tabularx}{0.7\textwidth}{X} 那帶種流淚出去的，必歡呼地帶禾捆回來！ \end{tabularx} \\ \\
[1ex]
\hline
\hline
\end{longtable}
$^{1}$10年之前,一位基督徒有一個領袖.
就要開始一間教會,在港島東區聚焦宣教.
本地宣教和海外宣教.
幾個月之後,全因堂就開始了.
上帝為我們預備在台灣一個地方.
接著就發生了很多事.
2017年,大家記得.
到現在,我們有經歷過社會事件.
新冠疫情,移民潮.
上個月,我們播道總會,有人告訴我.
播道會的人數跌了超過三分之一.
整個播道會,60間堂會的人數跌了超過三分之一.
在這個情況之下,全因堂當然也會受到這些外圍環境的影響.
但是,靠著上帝的保守,我們屹立到如今.
不僅如此,去年九月,上帝為我們預備了北角這個這麼好的地方.
今天我回來就跟Kid說.
我們自己安裝了冷氣,星期天就不能用大廈的冷氣.
這麼熱的時候,夠不夠力呢?.
上帝為我們預備資金,為我們預備執事,為我們預備童工.
我們童心竭力,有這麼美好的地方.
今天下午,也會有一個聚會,很多其他堂會的來賓.
有些人來想聽這個故事,我們宣教的故事.
有些人甚至來想聽我們這個地方是怎樣的.
如果大家今天下午可以的話,出席.
還有,你是主人家,歡迎他們.
就算你沒有報名,現在都可以報名的,出席.
這是一個慶祝,我們的神為我們預備了大事.
我們知道,隨著時間過去,我們都會老的.
今天早上我們崇拜,整天崇拜,我們想到.
我們的爸爸媽媽,身體不是很好嗎?.
我們的爸爸媽媽年紀不小了,我的爸爸媽媽就不在了.
但Kid,小方,江牧師那些爸爸媽媽,個個都不小.
其實,當我們展望前面的時候,可能都有點擔心,那怎麼辦呢?.
但今天我要跟大家講一段很興奮的經文.
就是詩篇126篇,剛才惠雄弟兄跟我們讀.
上帝為我們預備了,做了大事.
我跟太太7月5號就不可以參加唐慶的會友大會,年夜會.
因為我們在台灣,所以如果你不是去台灣的話.
我鼓勵你不要去北上旅遊,來參加會友大會.
因為我不能夠去到,所以這篇信息也是我心裡的感受.

$^{41}$上帝其實是為全因堂做了大事.
我們看一看這篇詩篇,頭三節是講到感恩的.
為上帝所做的事感恩.
第四到第六節為前途盼望依靠.
就是這麼簡單.
我們全因堂面對的情況也是這樣.
這首詩是以色列人回歸被擄70年.
在巴比倫回歸之後頭20年左右.
以色列人回歸是分幾個時期.
第一次回歸是公元前537年.
所羅巴伯帶五萬人回歸.
他們的目的是重建聖殿.
第二次回歸就是以色列過了差不多80年之後.
就帶領五千人回來.
他們的目的就是重建敬拜.
因為被擄回歸之後散了,人人都不懂敬拜.
然後再過十幾年,黎希米和一小撮人的侍衛回來.
可能只有十來二十個.
目的就是重建城牆.
所以背景就是這樣.
這首詩就在所羅巴伯帶領五萬人回來之後.
頭十多二十年有寫.
你知道那時候的環境.
就可以感受到這首詩的情景.
我們全因堂的情況當然和他們不同.
但是也有這樣的感慨.
你為過去感恩.
也為前途有些擔心掛慮要將自己交給神.
不單止教會是這樣.
個人也是這樣.
你想想你個人的日子.
我想高太十多歲的時候一個美女一個.
現在也很漂亮.
但是又過了幾年.
我就更加.
將來怎樣呢.
第一節我們看看.
耶和華說.
使被擄的人歸回錫安的時候.
我們好像作夢的人.

$^{81}$以色列人被擄之後就很絕望.
雖然上帝一早就說了我會帶你回來.
但是你直接被人搬去巴比倫.
前途根本不知道怎樣.
好像不太順.
等到他真的回來的時候.
那五萬人就很開心.
我們好像作夢的人夢想成真.
開心嗎.
這個美夢成真的感覺.
但是他回到耶路撒冷.
見到聖殿就拆毀了.
城牆也沒有了.
房子也爛了.
退換敗瓦一片荒涼.
他們就面對一個很重要的工作.
就是要重新來過.
第一步就是建聖殿.
他們建了一個祭壇.
可以說是獻祭.
建了聖殿的地基.
但是接著就停了.
周圍其他人就在搞他們.
投訴給皇帝聽.
又阻礙他們建築.
他們就覺得有點累.
心悸意懶.
還是不要建了.
建自己的房子吧.
自己的房子就喜歡建了.
你建一間我建一間.
聖殿又沒有人管.
一路就停了16年的時間.
大約在這段時間.
這個詩人就寫了這首詩.
你想想.
本來你回來的使命是要重建聖殿.
現在十幾年都沒有建了.
你的心就很鬱.
以前那種好夢成真.

$^{121}$現在就變成一場驚夢.
好像散了一樣.
所以我們好像作夢的人.
這句話其實是很複雜的感情.
初初回來的時候美夢成真.
現在就好夢成空.
不知道大家的心是不是這樣.
初初你滿懷大志.
剛剛讀書.
說要徵忠報國.
現在就不知道怎麼辦.
自己年紀又大了.
有人告訴我.
我以前的老牧師說.
現在年紀這麼大.
眼又矇 耳又聾.
就好像打鐘一樣.
怎麼能做到什麼呢.
有一種渺茫的感覺.
我就想起李白.
唐朝李白一首詩.
是這樣的.
描寫到他自己的心境.
都是這樣的.
李白小時候是個劍客.
武俠小說的劍.
真的 他作千幾首詩.
其中有十分之一是講到劍.
他也很喜歡打.
他也做過官.
但做官做得不太好.
我猜他也不會做得好.
他常常作詩.
做得不太好就被人趕出朝廷.
趕出朝廷就很不濟.
他就寫一首詩.
描寫自己的心境.
那首詩的名字叫杏露蘭.
因為唐朝的時候.
那些官話跟我們廣東話很相近.

$^{161}$所以你讀唐詩宋朝.
用廣東話讀是特別順的.
國語是清朝的滿洲話譯過來.
其實就沒那麼順.
廣東話叫中原古音.
所以你不要覺得自己失威.
不懂普通話.
你那些才厲害.
雖然沒人聽就不能講.
他就作了一首詩.
他說我吃都沒心情吃.
停杯頭注不能食.
即是停一杯扔了個筷子.
沒心情吃東西.
拔劍四顧身亡疑.
我拔起劍出來.
周圍望下都不知怎麼辦.
玉渡黃河冰塞穿.
我要過這條黃河.
但那條河就被冰塞了.
我要上個太行山.
但整個山坡就雪滿山.
玉渡黃河冰塞穿.
將燈太行就雪滿山.
前路茫茫那種感覺.
這種就是詩篇作者的感覺.
回來十幾年.
但不知怎麼辦.
升殿未起好.
還回來五萬人.
但有二百萬人還未回來.
國家將來前途如何.
一至三節為過去事感恩.
但感恩的時候.
他想到現在好像塞住了.
有時會覺得自己也是這樣.
我十幾二十歲信耶穌.
那時很風光.
二十幾歲做執事.
為何現在六十幾歲都沒進展呢.

$^{201}$會否有這種感覺呢.
但我要大家留意第四至第六節.
因為那三節很興奮.
那三節:耶和華求你使我們被擄的人歸回,好像南地的河水伏流.
河水伏流需要神蹟.
因為以色列的南地是沙漠.
很剛好.
很久才有一次雨.
但雨一下就下很大.
就像河一樣的沙漠.
但一會兒就沒了.
他就回到上帝那裡.
第四節:耶和華求你使我們被擄的人歸回,好像南地的河水伏流.
上帝其實是聽他們的禱告.
上帝復興以色列.
上帝的應許是沒有假的.
以色列亡國之後.
公元前五百多年亡國.
他正式復國是什麼時候呢.
就是第二次世界大戰之後1948年.
那時候希特勒殺了六百萬猶太人之後.
他這樣的情況下復國.
1948年宣佈復國.
馬上被人打.
周圍的國家.
由埃及,阿拉伯,沙特阿拉伯全部一起去打他.
竟然他搞定打贏了.
這個是神蹟.
然後1956年又被人攻擊.
又打贏了.
1967年又打贏了.
1973年都打贏了.
我研究過這幾次戰爭.
因為我對戰爭都很有興趣.
全部都是神蹟.
不可能打得贏.
美國的軍事部.
將這幾場戰爭去研究.
去做科學研究.
他們都奇怪為什麼會打得贏.

$^{241}$沒可能.
這個看到上帝的應許.
賀水伏流.
所以弟兄姊妹.
如果你今天的心情不好.
我現在工又做得不好.
爸爸媽媽身體又不好.
我自己身體都不好.
好像又滯在這裡.
以前那種的大智.
現在好像消磨耗盡.
好像李白那樣.
不過我又不會加劍.
你要回到上帝那裡.
因為上帝的應許是真的.
耶穌基督說.
我來是要給你豐盛的生命.
祂要賜給你豐盛的生命.
祂給你的救恩的拯救是永遠的.
不單止今生.
不單止來世.
是從現在開始一直到永久.
有一節的經文.
我跟大家分享過.
其實我很喜歡.
就是上帝對你的應許.
《申命記》三十二章十字.
耶和華在曠野之地遇見了祂.
在荒涼的地方.
和野獸嘲叫的荒野遇見了祂.
就環繞祂,看顧祂,保護祂.
好像保護自己眼中的同人一樣.
所以弟兄姊妹.
今天我要將這節經文賜給你.
無論你的情況怎麼樣.
上帝環繞你,保護你.
你周圍可能有很多豺狼,獅子.
但是上帝與你同在.
剛才偉雄帶領我們唱的那首詩很好.
耶穌因由與我們同在.

$^{281}$我記得這首詩的音樂.
我第一次聽就是六零年代的《細蚊仔》.
那時候那些人出殯就有樂隊.
如果是基督教出殯.
我就聽到他們奏這首歌.
後來我聽到原來這是聖詩.
很好.
其實值得回家再聽.
《126篇》的詩人禱告之後.
他心裡面暢快.
他在第五,第六節說.
「那些流淚殺種的必歡呼修國,那帶著種子流著淚出去殺種的,必帶著禍菌歡呼快樂會來.」.
所以弟兄姊妹.
當我為傳恩堂來禱告守望的時候.
祂一定會保守我們傳恩堂.
一定會照顧弟兄姊妹.
流淚殺種的必歡呼修國.
上帝怎樣叫以色列人復國.
怎樣帶領他們經過過去那幾十年的很奇妙的生存.
上帝都會帶領我們傳恩堂前幾日子.
上帝會大大地使用你.
叫你成為這個教會.
我們十年前那個義將那個領袖.
要向香港少數民族傳福音.
要向海外的地方傳福音.
今天下午Steven Chandra.
美國播道警察會的亞洲總監.
會和我們分享亞洲很多奇妙的工作.
還有幾個地方是我們香港播道環球警察會和他合作的.
我鼓勵大家去聽.
我們看完這首詩.
我們想想怎樣應用在我們身上呢.
這首詩叫上行之詩.
上行之詩就是以色列人朝聖.
每一年有三次如意節,五秦節和鑄盤節.
他們會去耶路撒冷那裡朝聖.
耶路撒冷就在山頂上面.
他們就向上走才走到耶路撒冷.
這是上行之詩.
我的幻想是他們每走一步就覺得自己親近上帝一點.

$^{321}$來到這個聖所.
我都鼓勵大家親近上帝.
在你每日的靈修裡依靠神.
有幾樣東西我鼓勵大家可以做.
將這首詩的話語實踐在你的生命裡.
為過往感恩.
無論你過往怎樣坎坷都好.
你能夠活到今日就已經是奇蹟.
所以有很多東西你要為過往感恩.
為將來盼望.
將貧信心將你的前路教託給主.
毀身基督.
重建你自己的敬拜生活.
我剛才提到所羅巴巴帶了五萬人.
回去十幾二十年.
聖殿又建不好.
自己的生活又開始敗壞.
以斯拉尼和他們重建.
回到神的話語裡.
教他們讀聖經.
建立會堂.
然後重建敬拜生活.
你和我都要反思自己.
重建我們自己的敬拜生活.
每日讀經祈禱依靠神.
擁抱教會.
我上次提過很多人.
是很孤單的老年的時候.
甚至過一世都沒有人知道.
這個教會裡面不應該發生.
我們要擁抱教會.
擁抱弟兄姊妹.
將你的時間恩賜金錢.
在教會裡面大家能夠共享.
和能夠支持教會的工作.
或者你說.
羅木斯我一個人世孤力弱.
我做到什麼呢.
我一個人有什麼可以做到呢.
除了回來坐在這裡崇拜.

$^{361}$或者是收拾東西之外.
其實一個人可以做到很多事.
我要和大家講一個真實故事.
而賭數是怎樣發現的呢.
在加拿大1923年有一個醫生.
他是做研究的.
加拿大的醫生和香港的醫生不同.
加拿大醫生的入息是很低的.
香港的醫生可能是真的很好.
這個醫生的入息低是不要緊的.
做研究的更低.
這個班廷醫生.
他是做研究.
研究二賭數.
那時候不夠資金去研究.
他自己買了一輛車.
又叫家裡的父母.
做農夫的農莊接濟他.
一直做研究的工作.
他和其他三個人在研究.
他又帶領.
竟然在2023年研究成功.
然後他用一元將這個專利權.
賣了給多倫多醫院.
因為這樣盡快可以將這個藥物.
供應給全世界.
在1923年之前.
糖尿病是絕症.
不能醫.
我們當中有沒有人患糖尿病.
好在你現在是2026年.
1923年之前不能醫.
但因為他的緣故.
很快全世界的人都可以醫.
所以一個人.
就可以改變人類的命運.
今日全因堂.
亦都可以因為你一個人.
而得著興旺.
我們其實.

$^{401}$我們在全因堂弟兄姊妹.
我們可以全民皆兵.
很多事可以做.
譬如我今日第一日回來.
我今日做.
我和我太太早了回來.
我們原本坐電車.
後來坐巴士.
很快出去見到巴士.
返了來.
返了來我立即去做.
有一件事我做到.
有一件事做不到.
男廁那些廁紙.
那些洗手紙.
整開門.
我做到.
女廁做不到.
我不知道女廁有沒有人做.
這些事可以做.
我回來我都做.
所以如果你第一個回來.
可以做那些.
還有你可以參加.
八點九的祈禱會.
因為八點九.
你可以早幾分鐘回來.
一齊祈禱.
很好.
在關穆星的房裡.
我們有幾個人.
為教會祈禱.
為弟兄姊妹禱告.
這些都可以做.
散會之後.
不要太快走.
有很多弟兄姊妹.
新知舊友.
大家聊天.
甚至約出去吃飯.

$^{441}$這個很難得.
因為我們是一家人.
這個世界這麼孤單.
如果我們不一齊的話.
怎樣呢.
我鼓勵大家.
一會三點鐘有聚會.
如果同拜之後.
我想去參加聚會.
中間那段時間.
我可以一起出去吃飯.
我可以和你一起出去吃飯.
我已經約了邵芳一起出去吃.
你可以參加.
你可以祈禱.
你可以用經濟金錢.
來支持教會事工.
這是很讚大家的.
我們租了這個地方.
需要一百萬.
我們訂了預算一百二十萬.
真的很快就籌到.
其實我們教會人數不多.
不容易.
但這是上帝的恩典.
亦是你的參與.
才能成功.
大家知不知道.
我們希望請多一位傳道人.
發展更好的事工.
將來培訓他下一個梯隊.
但如果你看經濟邵芳的報告.
我們其實沒有錢請第二個傳道人.
現在剛剛好.
所以錢從何而來.
上帝公認.
他有兩個辦法.
一是透過烏鴉.
一是透過你.
以利亞有隻烏鴉養她.

$^{481}$但通常上帝是透過.
教會弟兄姊妹來公認.
其實我自己也是這樣想.
我們希望今年年底.
第二階段的台灣.
開第一間教會.
我計劃五年後他自立.
不需要在香港實施支持.
不需要在美國支持.
而且五年後這個教會.
就可以開直分堂.
開枝散葉.
所以你今天的奉獻.
如果你當中有為台灣宣教奉獻.
我們全恩堂也有為台灣宣教奉獻.
將來就會開枝散葉.
我記得我二十四歲的時候.
在加拿大Calgary.
帶領幾個弟兄姊妹開始.
第一間堂會現在有七間.
我帶領完之後.
雖然我已經搬到西人教會.
但是要繼續發展下去.
所以其實神的侍庚很好.
就好像種菜或種樹.
你將一粒種子放在地上.
上帝就祝福.
Calgary 1977年開了第一間.
現在有七間.
我也開過其他的.
一間浸信會.
一間獨立教會.
上帝都恩代.
所以弟兄姊妹.
今天全恩堂有你在.
不是偶然.
上帝把你放在這裡不是偶然.
他要我們學習彼此相愛.
學習依靠他.
經歷到上帝的恩典.

$^{521}$經歷到流淚殺種.
歡呼收割的恩典.
我在這裡特別祝福大家.
因為我們今個月底.
我和太太就會回台灣.
我們每個月都落門.
因為我們帶兩頂帽子.
兩個崗位一個在香港一個在台灣.
6月26日是我們結婚50週年.
我們就去台灣.
7月5日就還沒回來.
但我祝福大家流淚殺種的.
別歡呼收割.
要和為我們成就大事.
我們鼓勵大家.
你一個人可以為這個教會.
為上帝的家成就大事.
上帝會使用.
我們一起低頭祈禱.
親愛的主人.
我們謝謝你.
我們看到以色列人被擄這麼多年.
竟然在二千多年之後.
能夠再次復國.
而且過去那八十年.
能夠一次又一次的得以保存.
我們實在看到你的恩典.
不是以色列人自己可以做到什麼.
而是你的應許.
你成就你的應許.
今天我們十幾年.
十年前我們姐妹的領袖.
今天我們都在你面前依靠你.
我們不喪膽.
我們雖然知道很多的外圍環境.
很多的變動.
但是你與我們同在.
你的應許必定成功.
我們就恭敬你來到面前.
我們依靠你.

$^{561}$我們仰望你.
我們要在你面前歡呼.
我們知道將來我們會有美好的收穫.
為弟兄姊妹每一個禱告.
為每一個家人禱告.
求你祝福你的令.
引導我們.
安慰我們.
圍繞我們.
保護我們.
我們祈禱感恩.
奉耶穌基督的名.
阿們.
\newpage



\section{瑪拉基書 2:10-16}
\label{sec:032260}
\textbf{向行事詭詐說:不!}~~~~石偉漢牧師
\newline
\newline
連結: \href{http://www.church.com.hk/acms/content.asp?site=cdc\&op=show\&type=product\&code=032260\&layout=sermon}{\texttt{032260}} ~~~~ 語音日期: 2026-06-07
\newline
\newline
\hyperref[sec:032259]{\small{< < < PREV SERMON < < <}}
~
\hyperref[sec:index_chronic]{\small{[返順時目]}}
~
\hyperref[sec:032261]{\small{> > > NEXT SERMON > > >}}
\newline
\newline
瑪拉基書 2:10-16
\newline
\begin{longtable}{cl}
\hline
\hline
章節 & 經文 (和合本修訂版)\\
\hline
2:10 & \begin{tabularx}{0.7\textwidth}{X} 我們豈不都有一位父嗎？豈不是一位神創造了我們嗎？為何互相行詭詐，褻瀆了神與我們列祖所立的約呢？ \end{tabularx} \\ \\ \relax
2:11 & \begin{tabularx}{0.7\textwidth}{X} 猶大行事詭詐，在以色列和耶路撒冷中行了可憎的事；因為猶大褻瀆耶和華所喜愛的聖殿，娶外邦神明的女子為妻。 \end{tabularx} \\ \\ \relax
2:12 & \begin{tabularx}{0.7\textwidth}{X} 凡做這事的，無論是清醒的或回應的，即使獻供物給萬軍之耶和華，耶和華也要將他從雅各的帳棚中剪除。 \end{tabularx} \\ \\ \relax
2:13 & \begin{tabularx}{0.7\textwidth}{X} 你們又再做這樣的事，使哭泣和嘆息的眼淚遮蓋耶和華的祭壇，以致耶和華不再理會那供物，也不喜歡從你們的手中收納。 \end{tabularx} \\ \\ \relax
2:14 & \begin{tabularx}{0.7\textwidth}{X} 你們還說：「這是為甚麼呢？」因為耶和華在你和你年輕時所娶的妻之間作證。她雖是你的配偶，你誓約的妻，你卻背棄她。 \end{tabularx} \\ \\ \relax
2:15 & \begin{tabularx}{0.7\textwidth}{X} 一個人如果還剩下一點靈性，他不會這麼做。這人在尋找甚麼呢？神的後裔！當謹守你們的靈性，誰也不可背棄年輕時所娶的妻。 \end{tabularx} \\ \\ \relax
2:16 & \begin{tabularx}{0.7\textwidth}{X} 耶和華－以色列的神說：「我恨惡休妻的事和衣服外面披上暴力的人。所以當謹守你們的心，不可行詭詐。這是萬軍之耶和華說的。」 \end{tabularx} \\ \\
[1ex]
\hline
\hline
\end{longtable}
$^{1}$頂英姊妹主內平安.
各位新來的朋友.
歡迎你們來到我們教會當中.
也歡迎每一個在家中崇拜的頂英姊妹.
讓我們在崇拜講信息之前.
我們一起有個祈禱.
讓我們一起有聆聽你話語的時間.
還有多謝你的說話.
是跨越世代.
即使在馬拉基時代的教訓.
都對我們今天有提醒.
完全幫助我們在生命中活出主理喜悅的樣式.
聽我們祈禱奉耶穌名字而求.
阿們.
今天一開始先問一個問題.
我想問一個人物.
這個人物和什麼食物有關係?.
這個人物叫做秦奎.
秦奎.
我網上找到很多不同的圖.
有些秦奎是瘦的.
有些秦奎是肥的.
我不知道哪個才是他.
我想問一下.
秦奎這個人物和什麼食物有關係?.
為什麼你們都認識?.
為什麼你們覺得和油炸鬼有關係?.
其中一個傳聞.
確實是說油炸鬼和秦奎有關係.
其中一個說法是當時因為秦奎不好.
秦奎和哪一個人物有些相關?.
岳飛.
岳飛和什麼有關係?.
和刀背有關係.
背上有些字.
精忠報國.
事實上岳飛很忠心.
去打仗.
被秦奎說服皇帝發出十二道金牌.
結果令到他要回來.

$^{41}$還要承擔一個罪名.
因為這個秦奎的行為.
引致當時的人民很生氣.
很激動.
於是其中傳聞.
有一間食店.
說食物與秦奎有關係.
為什麼油炸鬼有兩條紋路?.
有沒有想過為什麼有兩條紋路?.
因為當時的平民百姓.
其中一個店主.
將兩條紋路拼在一起.
將兩條紋路黏在一起.
背靠背.
代表秦奎和他老婆.
一起拼在一起.
拿去炸.
所以為什麼要拼在一起呢?.
吃油炸鬼.
吃完一條紋路.
要吃另外一條紋路.
一次過吃兩條紋路.
才夠吃.
當時就寓意.
油炸奎油炸奎.
秦奎.
油炸奎油炸奎.
現在變成油炸鬼.
這就是那種說法.
所以.
原來.
秦奎這個人.
很詭詐.
有很多陰謀.
所以要用油炸他.
而今天我和大家說的訊息.
與油炸鬼有關係.
將它反轉.
即是什麼?.
對了.

$^{81}$就是鬼炸油.
因為有時候油字.
可以讀成油.
好油走油.
今天就說鬼炸油.
為什麼呢?.
因為事實上.
我們今天讀的聖經裡面說到.
原來.
馬拉基時代的.
一群回歸.
被擄回歸的百姓.
其實他們回到他們的故鄉.
回到以色列的地方.
耶路撒冷的地方.
似乎本來.
應該安居樂業.
和太平回復正常的生活.
但是.
他們在那個.
環境之下.
會出一種鬼炸的生活.
所以馬拉基.
特別寫了馬拉基書.
去說這個問題.
當然這個問題為什麼會產生呢?.
上次和上次.
都有提及.
整個馬拉基書裡面說.
兩群人.
上樑不正下樑歪.
就是那群祭司.
如果有印象的話.
我們說在以色列裡面.
特別回歸的時候.
那群祭司已經是一個.
核心的屬靈領袖.
帶領著百姓回歸.
的生活裡面.
能不能夠活出上帝的.

$^{121}$心意,期望和那種.
敬拜生活.
就是由祭司去做.
但是我們上一次也提及.
祭司他們.
和神本來有約的.
有一份契約.
但是他們背棄了這份約.
他們背棄了這份約.
不單止自己背棄了.
和上帝那份立約的關係.
更加引致下樑.
那群百姓.
原來都不知不覺.
和他們一樣背棄了.
上帝的心意.
所以我們上一次看了.
祭司們的情況.
就會看到那群百姓.
在這群祭司.
牧羊之下,帶領之下的那群百姓.
怎樣生活的呢?.
而今天這一段聖經.
就是《瑪麗基書》.
二章十至十六節.
其中.
重複得最多的是哪個字呢?.
猜不到嗎?.
說了這麼久.
就是「鬼咋」.
否則我們今天的題目.
也不會有那群「鬼咋友」.
那群「鬼咋」的人.
因為短短七節的聖經.
其實你會看到.
和合本的翻譯更加明顯.
用了五次.
說到「鬼咋」這件事.
五次說到「鬼咋」的事情.
所以「鬼咋」.

$^{161}$就是這次段落裡的重點.
所以我稍後也會分段.
也會用了.
其中這幾個重點裡的段落.
我們說了三個不同的「鬼咋」.
但是簡單來說.
我們要問一個問題.
什麼叫「鬼咋」呢?.
今天我們有一種對「鬼咋」的理解.
但是在希伯來人.
原文裡的理解.
原來「鬼咋」這個字.
原來的意思.
是指用衣裳遮蔽.
用一件衣服遮住一些東西.
就是「鬼咋」這個字.
最原本的字的「字根」.
就是這個字.
當用衣服遮蔽的「字根」.
去衍生出來.
就是說一些人用衣服.
去遮蔽了自己.
以致表面.
和裡面不一樣.
表裡不一.
這是其中一個.
衍生出來的意義.
就是說那些人.
表面一套 裡面另一套.
裡面真正的那套和外面的不一樣.
另外一個.
衍生出來的意義是什麼呢?.
就是奪了別人的外衣.
就是說.
本來的東西不屬於你的.
就用了一些方法.
包括了欺騙.
訛詐 欺負.
背信 這些方法.
去拿了一些不屬於他的東西.

$^{201}$所以今天.
我們就會看到.
在以色列裡面那群百姓.
如何去做了這些行為.
當中我嘗試用.
三個分題去講解.
包括三個詭詐指標.
指標是一個index.
幫我們去看到.
一個狀況.
三個詭詐的指標.
我們先看第一個詭詐的指標.
他們的詭詐是怎麼樣呢?.
就是「詭詐代帝卿」.
「詭詐代帝卿」.
這裡我已經附上了.
和合本的翻譯.
上面就是和修本的翻譯.
下面就是和合本的翻譯.
和修本就是說.
我們豈不都有一位父媽.
豈不是一位上帝.
創造了我們嗎?.
為何互相行詭詐?.
褻瀆了上帝.
與我們列祖所立的約呢?.
在這裡說「互相行詭詐」.
即是說不只是一個人.
是有一群人.
互相之間去行詭詐.
但這群人是什麼人呢?.
和合本就幫我們有更加明確的了解.
就是「詭詐代帝卿」.
說的是一些在以色列.
一些帝卿之間的詭詐.
所以這裡用上「我們」.
馬拉基說.
我們都是有同一位父.
有同一位上帝在我們當中.
但是這群人.

$^{241}$他們帝卿之間行了詭詐.
所以說.
「我們」是在說什麼呢?.
「本是同更新」.
這裡不是在說.
我們基督徒之間.
我們都信耶穌的那種帝卿的關係.
而是在說整個.
以色列十二支派當中.
他們那種的關係.
這個就是背景.
十二支派當中.
他們同樣是以色列國的子民.
他們是「我們」.
因為這群人.
十二支派當中整個以色列國.
都是和上帝有立約的.
都是在上帝的約裡有份的.
而且他們都是.
同一位的父.
因為他們的源頭就是亞伯拉罕的父神.
所以他們.
整個民族當中.
都有同一個的約.
在他們當中.
這個就是他們那種的狀況.
簡單來說.
即是他們是同胞.
他們是自己人.
他們是同是天涯淪落人.
因為他們經歷了.
被擄去到巴比倫之後.
現在終於回歸.
回到自己的故土.
這群一起回到故土安居樂業的人.
和這群人.
是淪落的.
去到外邦成為漁民.
漁民的意思就是.
剩漁的民.

$^{281}$回到巴比倫.
而留在耶路撒冷的那群人.
都不是好得多.
因為當他們被擄的時候.
整個國家都滅亡.
生活都很淒慘.
所以當這兩群人走在一起.
在以色列這個地方生活的時候.
看回大家的源頭.
原來大家都是同胞.
大家都是同更新的.
一起這樣.
理論上應該是怎樣呢?.
如果大家都是落難的.
悲慘的.
另外一群人也是這樣.
應該怎樣對待那些人呢?.
理論上大家都是經歷過這些艱苦.
應該互相幫助.
我想起.
圖畫.
左邊這幅圖.
不知道大家看過.
這好像是一套戲.
你認識的.
我不認識.
不坐飯的.
其實.
有些人看過這套戲.
這套戲有什麼特別呢?.
有兩個特別.
我自己觀察到.
第一.
有一個很重要的說話.
八個字.
人人為我.
我為人人.
這種精神.
我從小就聽這種精神.
有些戲都翻拍.

$^{321}$六七十年代.
七八十年代.
都會講這句說話.
經常都會聽這句說話.
代表什麼呢?.
大家都知道.
香港那些生活的日子.
不是很富貴.
不是很發達.
有時借錢.
有些人.
在中國走來香港.
你又會住在你家裡.
當時是.
互相幫助.
你幫我.
我幫你.
有些人分給別人.
我岳父也是這樣.
走來香港的時候.
他在香港住之後.
有些親戚來到.
他也是這樣照顧他.
大家走來香港.
生活的時候.
都是這樣的心態.
就像今天我們來到美國.
有些人早點來美國.
剛才說有些美國華僑.
遲一點來美國的人.
他初到.
貴境很陌生.
我們是不是應該幫他.
應該是這樣的.
那些同胞.
自己有來的.
這個精神很重要.
我為人人人為我.
但同時這幅海報有個特別的地方.
是什麼呢?.

$^{361}$就是不是彩色的.
不是彩色的.
代表什麼呢?.
就是黑白的年代.
黑白.
什麼叫黑白的年代?.
非黑即白.
很清楚的.
黑色就是黑色,白色就是白色.
石堅走出來一定是什麼?.
石堅一生人.
真的很差.
一走出來他就是奸的.
是不是.
為什麼呢?.
很黑白的.
中那個一直中下去.
就是那些角色.
就是這樣定型的.
但隨著時間的轉變.
我們看的顏色就不同了.
現在的海報是什麼顏色?.
彩色的.
有些黑,黑得來其實可能是白.
白得來有些是黑,其實有些是灰.
原來不知不覺.
變了顏色.
不是黑白兩種顏色,是彩色.
代表什麼呢?.
那種生活的狀況.
有時原來.
那些自己人.
原來都對自己人不好.
《馬拉基書》就在說.
這種狀況.
大家都在這個.
以色列的地方.
經歷了被擄歸回.
重生生活.
亡國這些,全部一起經歷了.

$^{401}$理論上大家一起.
彼此扶持,照顧大家.
但原來有些人.
掛著一些名號.
原來在你身上拿著數.
的.
即是你以為是黑白.
他應該幫你,但他不一定幫你.
你以為他應該.
為了你好,但原來他為了.
自己好.
《馬拉基書》就在說這樣的狀況.
以致到一起在.
以色列生活的人,特別當面對.
那種經濟環境.
不是那麼好,還有不同的人.
很多人窮人的.
艱難處境的時候,原來.
他們恨鬼咋.
所以.
鬼咋弟兄就說.
他們在社會上,用鬼咋.
去彼此對待.
自己同胞,同族的人.
以致到在對方.
身上拿很多的著數.
在.
聖經另外一本書叫.
《阿摩斯書》就這樣說.
他說.
當時那群以色列的人.
怎麼說呢?當然.
《阿摩斯書》在馬拉基書.
早大約三百多年.
的時間,但是他們那種.
行為表現,我們可以想像得到.
去到馬拉基時代.
仍然會發生.
《阿摩斯書》第八章五至六節.
說當時在.

$^{441}$阿摩斯的年代,那些人.
怎麼去恨鬼咋的呢?.
他說.
月朔何時過去,我們好賣糧.
安息日何時.
過去,我們好擺開.
物資.
日出要用小星斗.
收銀要用大凳子.
用鬼咋的天秤.
欺控人.
好用銀紙買貧寒人.
用一箱鞋換.
窮乏人.
將壞料的物資.
賣給人.
《阿摩斯》裡面其實是其中一個.
有意思的書卷.
我不能夠今天詳細地說.
《阿摩斯書》,但是要說一些狀況.
原來那些人恨鬼咋.
真是古古惑惑.
他說.
月朔何時過去,我們好賣糧.
說的是當時.
阿摩斯時代.
那些以色列人經常問.
何時月朔過去.
等於我們的十五.
何時過去.
即是說,慶祝上帝.
敬拜上帝的日子.
何時完結.
什麼意思呢?.
因為快點完結.
可以上班.
可以賺錢.
所以那群百姓.
敬拜上帝時都問.
夠鐘了嗎?.

$^{481}$何時完結?我現在出去玩.
我現在出去賺錢.
他們不斷想快點完結節期.
快點繼續做生意.
不要阻礙他們.
通常這些日子都是安息的日子.
所以後面說,安息是何時過去.
擺開墨紙,拿出來擺檔.
做生意.
很重要.
原來他們這樣去.
敬拜上帝.
還有什麼呢?賣出用小星斗.
即是說.
我們有時賣東西.
會用個梁器.
裝著.
用來秤多少東西.
都用個梁器兜著.
買墨那些都是這樣裝著.
賣出用小星斗.
即是賣東西給人.
裝給人的那個兜.
小了一個碼.
你明白嗎?.
即是本來一個小星斗.
大約22kg的.
重量size.
但可能到他.
做生意賣給人的時候.
其實他的籃是21kg.
不知你有沒有見過.
有時你去.
超市買東西都是這樣.
我明明見到那些日本食物.
寫著十枚入.
十枚入即是十包裝.
那個盒的設計.
是一樣的.
但不知不覺.

$^{521}$原來有個小字寫著九枚入.
即是同樣size的盒.
現在變了九包裝.
不知不覺少了一包.
他又沒有改小盒.
但他改少了一包.
就是這樣.
賣出用小星斗.
小了的東西.
收銀就用大凳子.
即是收銀有個東西.
也是這樣裝著.
有時他怎樣計算.
因為你知道.
古代不是拿現金.
古代是拿現金.
不過現金不是紙幣.
你拿多少張銀紙.
數就可以了.
但古代怎樣去買東西.
拿些東西去買.
而那些東西往往是計重量.
計大小.
當你去收人家的東西的時候.
就有一條線.
你去到那條線.
就夠那個數了.
但現在你整個盒子.
大了.
你整個盒子大了.
就算同一條線.
大了闊了.
你即是給多了.
買貴了.
有時也是這樣.
有時有很多古古惑惑的招數.
這裡還說.
用鬼扎的天秤欺控人.
怎樣叫鬼扎的天秤.
如果用今天的角度.

$^{561}$你想想你以前買豬肉.
香港.
香港買豬肉是怎樣的.
有個勾有個秤.
放個東西在那裡移前移後.
即是現在就說.
壓秤.
鬼扎的天秤.
如果用天秤的方式.
你放些東西在這裡.
這裡放一塊東西.
大家剛剛平衡.
就代表了對的重量.
在用這塊東西裡.
不同了.
所以你會發現.
原則就是鬼扎的人.
用了兩套的法馬.
兩套的標準.
給你看到的.
是一個很公平的標準.
暗地裡面用了一個.
不公平的標準.
阿摩斯時代是這樣.
去到馬拉基時代.
其實也是這樣.
所以今天說一個問題.
就是原來.
在意識裡面有一班人.
他們行鬼扎.
就是他們騙人.
但同時人家騙他.
你騙他他騙他.
騙到怎樣.
就個個都騙.
是不是.
有時都是這樣.
我不知道你在美國生活是不是這樣.
即是.
越來越多.

$^{601}$正正常常.
做生意或者買賣.
做一個星斗市民.
去做一些事.
不過原來有很多時候.
有很多人都不是很老實.
就算一些專業的人士.
即是一些專業的人士.
你聽到他專業.
他應該放心交給他做一些事.
原來有時候魔鬼在細節.
魔鬼在細節.
有一個姊妹.
她就說.
有一次她去賣樓.
她就跟經紀說.
買一個價錢.
我就當30萬.
買一個價錢的時候.
跟住經紀就說.
哎.
人家講價.
就講便宜一點.
希望能夠爭取好一點的價錢.
結果經紀就買了一個.
好一點的價錢.
但原來.
後來那個姊妹回家的時候.
再看那個合約.
看著看著.
原來不是.
你以為賣好一點的價錢.
原來有估劃.
姊妹分享就說.
原來那個合約裡面.
後來就說.
原來這個賣了好一點價錢之後.
業主之後要再給回佣.
給那個買家.
但那個經紀沒有跟她說.

$^{641}$即是說.
你賣貴了五萬.
原來你最後要回.
幾萬給那個買家.
即是說你沒有賺到.
沒有賣貴.
為何這個經紀這麼奇怪.
告訴我為何地產經紀會這樣做.
因為佣金多了.
是不是.
我聽到都覺得.
不是地產經紀是專業來的嗎.
原來他們才最會出估劃.
跟住.
他還告訴我.
是不是.
好像這樣的說話.
原來他騙你.
原來.
地產是這樣的.
原來有些專業也是這樣的.
今天很多律師.
我們香港就不同.
有錢才找律師.
這裡不用錢找律師.
是不是.
即是說.
你有案件.
如果一些訴訟案件.
是不是.
你只要找到他.
他會用盡方法幫你去爭取.
跟住分多少.
我都不知道分多少成.
有沒有分到一半.
所以這些訴訟.
其實有時本來.
那單案是不應該訴訟的.
但是.
原來在一起生活的環境裡.

$^{681}$他們為了自己的好處.
就發起訴訟.
是不是都告你.
十萬.
總之到時.
萬一講價有三十萬.
他都分到十五萬.
如果有五成的話.
所以原來.
今天那個環境就告訴我們.
詭詐.
是會蔓延的.
詭詐是會蔓延的.
當你詭詐人的時候.
你為了自己好一點.
可能你會又詭詐回另外一些人.
在這個時代裡.
一個詭詐的國度.
就是同胞之間.
互相詭詐.
帶來一種負面錯誤的社會風氣.
所以第一個指標.
就是我們怎樣去我們身邊的同胞.
我們的同鄉同族.
特別我來到美國.
我什麼都不懂.
我靠你們.
如果你們詭詐我.
我都被你騙了我的身家財產.
不知道有沒有被騙.
我找會計師幫我查一查.
所以相反.
其實我們要怎樣去對待弟兄.
怎樣去對待自己人同胞.
其實應該誠實.
做生意.
其實做長線.
是要君君珍珍.
你不是君珍的話.
你只是做一次的生意.

$^{721}$所以賣樓就容易被騙.
應該幾年都不會是下一次.
所以他就騙你.
但如果你做普通飲食生意.
就很難了.
你的態度不誠實.
你騙他一次那五元十元一個飯.
沒有十元那麼便宜.
十五元一個飯.
下次不來光顧.
上次我去某間酒樓.
喝茶無端端說要跟我查錢.
我已經不再光顧他.
我們要建立一種誠實的狀況.
最近我又有一次的經驗.
我去Pubix.
買東西.
Pubix買東西.
我和師母買的時候.
其中看到有包帶子.
有另外一個姐妹說.
這包帶子很划算.
我們看到帶子就想買.
我們一直看價錢.
我找了很久都找不到.
因為通常包裝在上面.
價錢在下面.
不知道為什麼.
放到亂了.
總之我們看了很久.
我看到那個價錢.
應該是我們想的價錢.
我們終於買了這包東西.
買完之後.
價錢牌不一樣.
貴了.
原來和我們在價錢牌下面.
那個牌是不一樣.
我們再看看.
不可以.

$^{761}$要不要退貨.
退貨又要排隊.
我覺得很煩.
走就算了.
兩三元都不划算.
過不了自己那關.
如果我回家我就不回去了.
但我在店鋪門口.
我馬上發現.
於是我就走去.
排隊退貨.
這個價錢和我們看到.
那個牌的價錢不一樣.
我拍了一張照片給她看.
看一看.
對喔.
退錢給我.
是不是應該要這樣的服務才可以.
退錢給我.
然後我就.
給她一包東西.
她還給我.
不愛.
什麼人來的不愛.
她還給我說.
我想買個包.
我退錢給你.
不愛她還給我.
我走了.
原來是這樣的.
是真的這樣.
我覺得.
這些真的值得去.
為什麼.
因為她打錯了價錢牌.
她知道自己不對.
而肯退錢給我.
代表什麼.
她不是故意騙我.
代表這間店.

$^{801}$她是用誠實的方法.
希望用一個正確的價錢.
心甘樂意買.
雖然是貴一點.
心甘樂意買.
買一個好的服務.
原來是這樣的.
我們今天也要想想.
怎樣對待我們身邊的人.
這不是說信不信耶穌.
這麼簡單.
身邊的人信不信耶穌.
特別我們在美國這個地方.
難得遇到華人.
我們華人.
應該要互相.
誠實相待.
好好幫助對方.
要人人為我我為人人.
不要像馬拉基書.
時代的人一樣.
落在一個.
既然是這樣的生活艱難.
我再詭詐一點.
小心謹慎.
所以有些反思和分享.
待會大家可以去.
行道團界裡分享.
你在生活裡有沒有.
真的經歷過一些被人詭詐的事件.
有時候在.
教會圈裡也會有.
一些詭詐的事件會發生.
因為有些人帶著一個.
不正確的動機來到神的教會裡.
所以我們也要謹慎.
同時間我們也要.
有一份智慧.
所以聖經說什麼?.
要靈巧像蛇.

$^{841}$順量與格致.
所以對人好是應該的.
但又要有智慧.
這是聖經教我們的.
所以我們也要去思考.
第一個.
詭詐待弟兄.
是一個詭詐的指標.
第二個詭詐的指標是什麼呢?.
就是詭詐.
待上帝.
原來我們有時候對上帝.
也會騙的.
上帝是全知,全能,全善的.
你應該沒有辦法騙到祂.
不過有些人仍然覺得.
你可以騙到神的.
可以騙到神的.
我們看看怎麼說.
十一,十二節.
猶大行事詭詐在以色列和耶路撒冷中.
行了可憎的事.
因為猶大人.
褻瀆耶和華所喜愛的.
聖殿取愛邦神明的女子為妻.
凡做者士的.
無論是清醒的或回應的.
即使獻供物給耶和華.
萬君子耶和華.
耶和華要將她從雅各的帳篷中剪除.
這裡說的是.
一種狀況.
這些以色列人.
明知故犯.
因為在他們.
祖先和神的立約.
也很清楚在摩西五經裡說.
他們和外邦神明的女子的關係.
但他們明知故犯.
在耶路撒冷.

$^{881}$以色列的地方.
娶了外邦神明的女子為妻.
在《出牙疾記》裡說.
你為你的兒子娶他們的女兒為妻.
他們的女兒.
因著隨從他們的神明行任.
就引誘你的兒子.
也隨從他們的神明行任.
生命句裡也說.
因為他必使你的兒女離棄我.
去事奉別神.
所以留意.
這段聖經裡說.
我們跟隨神的人.
小心不要選外邦女子為妻.
但這裡是什麼意思呢?.
不是說不可以選和你不同族的人結婚.
如果是這樣.
今天很多人都搞不定.
因為有些人喜歡娶外國人.
外國人.
特別是你的子女.
混在一起通常更漂亮.
更漂亮.
但這段聖經不是說.
今天不能娶金髮女子.
他說什麼?.
不能娶外邦神明的女子.
如果你知道她正在敬拜.
外邦神明.
敬拜另一個神.
這樣就會引致.
你敬拜一個神.
你自己也敬拜一個神.
大家敬拜不同的神.
就會出現一個問題.
要麼兩個神吵架.
要麼兩夫婦不和順.
另外一個情況是什麼呢?.
你為了和順.

$^{921}$要麼你跟他.
要麼他跟你.
而很大機會是你跟他.
這就是神指出的問題.
很大機會我們不知不覺就跟了別人.
跟了別人代表什麼呢?.
你不知不覺.
就是敬拜了別的神.
敬拜了另外一個神.
而不是敬拜耶和華.
你會隨從他們.
拜了外邦的神明.
外邦的神明裡面.
特別是在迦南地.
很多神明.
和淫亂可能也有些關係.
我們的信仰就會不知不覺的敗壞.
所以在《出埃及記》裡.
很清楚提醒我們要怎麼做.
所以神提醒我們.
原來有很多神明在這個世界存在.
事實上.
當我們讀到《出埃及記》的時候.
記不記得有十宅.
其實十宅裡面.
其實十個宅.
都是和埃及裡面.
不同的神明有關.
可能你會說.
神為什麼帶領以色列人出埃及.
不直接叫他們出來.
直接一開始殺光那些長子.
直接出埃及離開.
原來.
整個《出埃及記》的十宅裡面.
是在說神怎麼和.
十個埃及裡面的神.
鬥爭.
宣告神是得勝的.
一個個.

$^{961}$埃及所會拜的神.
去到神的面前.
他們都要輸.
以至到最後神就宣告.
你們就會知道.
哪一個是神.
《出埃及記》很清楚.
它的目的就是要讓人知道.
哪一個是在埃及的神.
哪一個才是真正在世界上的神.
就是神自己.
而這三幅圖正正好就是說.
在埃及裡面.
其中幾個.
他們拜的神明.
就是十宅裡面所引用.
衍生相關的神明.
原來就是這樣.
原來神要求我們.
不要和一些.
外邦神明的女子.
或者不要和一些.
外邦神明有交往的一些人.
他們是知道的.
但他們明知故犯.
明知故犯.
那你擺明車馬.
那也不是詭詐.
我擺明車馬.
我就是不聽.
或者我擺明車馬.
我娶一個這樣的女子.
明知故犯.
我們都看得出.
但他更加詭詐的地方是什麼呢?.
這裡說.
他還裝作沒事.
失敗了.
原來他裝作沒事.
去到神面前才可怕.

$^{1001}$他才欺詐上帝.
你擺明我是不跟你了.
我不愛神了.
那你就自己放棄.
但你卻說我很愛神.
我會來神的殿去敬拜祂.
我又要去讚美神.
Hallelujah.
其實我自己就知道.
我底牌裡.
和那些外邦神明的女子.
有一些外邦神明的關係.
所以他們說.
他們的欺詐行動是怎樣行動呢?.
他們繼續獻公物給耶和華.
在過程中.
褻瀆神的聖殿.
褻瀆神的地方.
你不知不覺.
其實我還很敬拜神.
當你這樣去到神的聖殿.
其實是在踩牆.
在踩牆宣告.
讓這些不潔.
來到神的聖殿.
是上帝不喜悅的.
原來他們就是這樣.
來到這個地方.
你會看到.
這個會幕就是上帝所要.
用來敬拜的場所.
但他們就帶著祭物.
來到這裡當中.
去獻祭.
神願不願意立這些獻祭呢?.
如果有這樣的情況.
就是去敬拜拜偶像.
敬拜一些外邦的神.
然後又回到神的殿.
去獻祭.

$^{1041}$還要獻最好的禮物.
那神願不願意呢?.
這麼貴 這麼好.
願不願意呢?.
不要啊 弟兄姊妹.
不要不是因為那些東西不好.
不是最上等 不是頭生.
而是那些東西.
是上帝知道你在開罪他.
弟兄姊妹你們要留意.
今天我們犯了罪.
是不是可以有贖罪祭.
大家記得.
贖罪祭是用來做什麼的?.
贖罪嘛.
如果你有罪.
那你拿祭物來贖罪.
那不就對了.
摩西五經.
我們學了五種祭.
其中第四種就是贖罪祭.
那獻贖罪祭就可以了.
每一次犯錯.
然後又回來教會.
贖罪 又出去犯錯.
又回去贖罪.
我們這樣循環生活可以嗎?.
不可以 為什麼不可以呢?.
聖經哪裡說不可以呢?.
我們要問一個問題.
聖經哪裡說不可以呢?.
聖經說不可以.
我們看一看.
贖罪祭的條件是什麼呢?.
我們看一看.
你留心看.
五個祭裡面.
四章就講到贖罪祭.
其實贖罪祭講得很清楚.
四章二十七節說.

$^{1081}$這地的百姓.
若果有人是無意中犯罪.
在任何事上犯的一條.
耶和華所吩咐的禁令.
就有了罪.
即是說.
無意中犯罪.
你才可以獻贖罪祭.
因為有時候我們得罪了.
我們不知道.
我們有時候失敗了.
我們不知道自己失敗了.
我們有意犯罪就不一樣.
有心犯的罪.
上帝是不會讓你平息的.
所以他們的欺詐是什麼呢?.
他們來到神的面前.
中獻祭的時候.
就說我不知道.
覺得自己是無意犯罪.
扮無知.
而在欺詐神.
這就是他們的問題.
他們扮無知.
扮到什麼地步呢?.
在瑪拉基書二章十二節說.
凡做者事的.
無論是清醒的還是回應的.
即使獻供物給耶和華.
神都不要.
要將他們剪除.
為什麼呢?.
因為神說.
他們在清醒.
或者回應代表什麼意思呢?.
清醒即是真的知道自己發生了什麼事.
一些真的知道自己發生了什麼事.
而現在走來獻祭的人.
又或者有些人.
回應說.

$^{1121}$我被逼的.
我今天犯罪得罪了.
我被逼的.
什麼叫被逼?.
犯罪很多時候就是.
自己心甘樂意.
和自己的罪相關.
所以神說這些罪.
神說都不閱立.
都不閱立.
所以今天.
第二個很重要的思考就是.
我們不要詭詐對上帝.
首先我們常常都要檢視.
我們生命裡面有沒有一些罪.
在我們生命裡面發生.
當我們發現.
我們生命裡面有些罪的時候.
我們要問一問究竟是有心還是無意呢?.
是否真是無意呢?.
當你無意的時候.
你醒察到 聖靈提醒了你.
你就要認罪.
但如果你有意的話.
你要立刻悔改.
神期望我們悔改.
而不是繼續犯罪.
而帶這些祭物.
去到神的面前.
神是不閱立的.
因為這是一種詭詐對上帝.
一種信仰生活裡面.
如果你行詭詐是一個問題.
《馬拉基書》時代裡面的百姓.
就是這樣去對待他們的信仰.
第三個詭詐的指標是甚麼呢?.
這問題待會我們有時間再討論.
第三個詭詐的指標是甚麼呢?.
詭詐待妻子或詭詐待家庭.
特別是說妻子.

$^{1161}$這裡說甚麼呢?.
「你們又做這樣的事.
使哭泣和嘆息的眼淚遮蓋耶和華的祭壇.
以致耶和華不再理會那供物.
也不喜歡從你們手中收納」.
你們說這是為甚麼呢?.
因為耶和華在你和你年輕時所娶的妻子之間作證.
她雖是你的配偶.
你誓約的妻.
你卻背棄她.
「一個人如果還剩下一點靈性.
他也不會這麼做」.
這人在尋找甚麼呢?.
「上帝的後裔.
當謹慎你們的靈性.
誰也不可背棄你年輕時所娶的妻子」.
耶和華說「以色列的上帝.
我恨惡丘妻的事.
和依附外面披上暴力的人.
所以當謹守你們的心.
不可行詭詐」.
這裡說甚麼呢?.
這裡說這班人在聖殿裡哭泣.
其實我想解釋一點.
在和學本和和修本有少許不同.
如果有時你喜歡看和學本.
你也要留意了.
和學本第十三節說.
「你們又行了一件這樣的事.
使前妻嘆息哭泣的眼淚遮蓋耶和華的痰」.
所以和學本的翻譯就是說.
誰在神的痰面前哭泣.
哭得一局局.
就是前妻.
和學本的翻譯就是這樣.
為什麼這件事去到十四至十六節.
神說他們在婚姻上的詭詐.
以致在十三節說是因為前妻.
前妻在神面前哭泣.
哭得一局局.

$^{1201}$以致神不收納他們的供物.
這就是和學本的翻譯.
但你看看和修本.
和修本其實沒有了「前妻」這個字.
「你們再這樣做.
使哭泣和嘆息的眼淚遮蓋耶和華的濟嘆」.
有少許不同.
我想說的不同是.
原文的嘆息哭泣主角.
根據他們的文法.
原來不是一個陰性的詞.
不是一個女性的詞.
不應該用「前妻」.
就是這樣.
所以應該說什麼呢.
就是那群憲制的百姓.
這裡說什麼呢.
這裡說原來是因為你們背弱.
所以神不會約納你們的祭物.
所以你們這群百姓.
就要在神面前哭得一局局.
不是前妻哭得一局局.
就是這個意思.
即是神說你們的詭詐臨到神面前.
祂看得到你們的問題.
所以你們來獻祭的時候.
你們都知道.
你們無論是前面剛才說的事.
或者現在跟著要說的事.
因為你們要承受你們的後果.
所以你們要面對神五月立.
神要剪除.
你們要面對著.
在神面前要哭泣流淚.
所以這群百姓要承受他們的後果.
為什麼要哭得這麼慘呢.
因為上帝不約納.
但這裡說的是.
特別是聖經裡面說.
他們之後要怎樣對待妻子的時候.

$^{1241}$特別是會帶出一個很重要的婚姻的意義.
婚姻的意義有三個重點.
今天我們要怎樣對待我們的妻子呢.
要回復我們婚姻的本質和意義.
經文裡面提及第一個婚姻的本質和意義.
就是成為一位配偶.
婚姻的意義就是.
上帝要為我們尋找一位合適我們的配偶.
而這個配偶將要跟我們有一個婚約.
所以後面就說.
是你的配偶 你誓約的妻.
所以馬拉基就說.
婚姻的意義就是這樣.
要成為一位配偶.
跟她有一個婚約 一個盟約.
第二就是這個盟約是有神做見證的.
為什麼我們不可以看輕我們的婚約.
不只是你們兩個的事.
我們經常說結婚是我們兩個人的事.
不關其他人的事.
是呀 不關其他人的事.
不過是關神的事.
因為神是在婚姻裡面的見證人.
神在你們的盟約裡面立誓.
所以我們教會在基督教裡面的婚姻的時候.
我們的婚姻是在神面前做的.
我們不是做給嘉賓看的.
我們在神面前宣告.
我們以我們的婚姻建立一個基督教的.
基督化的家庭.
這個是上帝喜悅的.
所以很多時候我們誤會了.
他不好 我們就離婚.
不是 因為這個婚約裡面還有神在當中.
神沒有說不批准 你要留心了.
這件事是不是還要做呢?.
第三 一個婚姻的意義就是延續你的後裔.
這裡說要得虔誠的後裔.
因為願你有更多的子孫在當中.
所以婚姻不只是兩個人的事.

$^{1281}$也牽涉到怎樣有我們的後裔.
怎樣教導我們的後裔成為上帝的子孫.
是重要的.
所以婚姻裡面有一個很重要的任務.
就是培育你們的下一代.
不只是培育他們得志體群美.
更加是靈和上帝的關係.
這才是上帝容讓我們有後裔的一個很重要的關鍵.
所以這裡說的是.
《馬拉基書》裡面說的就是.
這群百姓他們初初年輕的時候有妻子.
到自己生活好一點.
或者那個狀況不同了.
就有離婚.
他們興翰了年輕時所娶的妻子.
而這種狀況是什麼呢?.
就是令到妻子帶來受到傷害.
你知不知道離婚是很有害的?.
大家不知道有沒有看過.
其中一個的研究.
而他們的詭詐的行動是什麼呢?.
他們的詭詐就是.
這裡說衣服外面披上暴力的人.
《和學本》翻譯得比較好一點.
就是強暴代妻.
這群人既然有一個上帝為他們見證的婚姻.
但他們用什麼行動去處理面對他們的婚姻呢?.
第一就是家暴.
進行家庭暴力.
以至對自己的妻子殺人暴力.
暴力完之後還怎樣呢?.
就休妻.
就是不要這個妻子.
就像垃圾一樣扔掉.
然後就馬上另娶.
小心點.
當你婚姻完結之後.
馬上另娶是高危的.
因為你是一個不清楚自己心理狀態的狀況.
甚至很多時候就是為了用一物取代一物.

$^{1321}$或者用一個人取代另一個人的狀況.
當你用取締的態度去看婚姻.
你很容易就會再找下一個人取代另一個人.
所以今天很多時候社會裡面.
婚姻離婚率越來越高.
就是因為什麼?.
不合適就換.
iPhone不合適就換Samsung.
Samsung不合適就換華為.
我們當我們的婚姻就像手機一樣.
可以換換換換.
但不是的.
上帝眼裡的婚姻是一個人來的.
而這群百姓就是用暴力對待.
休妻.
接著換一個.
進行了這個鬼炸代妻三部曲.
帶來的傷害是令到前妻確實哭得一塌糊塗.
這裡說.
有一個研究說.
一個美國的整形和修復外科醫學的期刊裡.
做了一個這樣的研究.
他就說.
離婚生活.
離婚的生活壓力.
令到人有生活環境壓力.
異常消瘦.
抗抑鬱藥.
甚至會令人衰老.
所以離婚會令到整個身體帶來一種轉變.
而這種轉變是負面的.
以致對一些人來說.
其實這種婚姻狀況會破壞了很多人.
特別是女人老得快一點.
其實婚姻不如意也會老得快一點.
所以一定要抗這些衰老.
也要抵抗這些衰老.
明白我的意思嗎?.
剛才說的這個研究是怎樣的呢?.
研究員就找了186對同卵雙生的雙胞胎.

$^{1361}$雙胞胎其中一個特色是什麼呢?.
基本上很多基因都是一樣的.
而他們說對186對雙胞胎進行基因研究的一些設定.
假設他們完全相同.
結果發現離婚者的外貌.
看起來較已婚,喪偶,獨身的聯生兄弟姊妹老了近兩年.
即是會老了兩歲.
其中一個關鍵的原因是壓力.
壓力可以視為那些外貌較老雙胞胎的共同特點.
其實會帶來一些這樣的損害.
另外還有一個研究.
美國最喜歡做這樣的研究.
美國芝加哥大學和約翰霍普金斯大學.
研究了8652名51歲至61歲的男女.
發現離婚或者喪偶的人士.
腦癌,心臟病,糖尿病,癌病的慢性風險比已婚人士高出20\%.
走路氣喘的風險高出23\%.
即使再婚出現慢性病和走路困難的風險分別高出12\%和19\%.
研究員解釋離婚和關寡人士壓力很大.
傷心,收入減少,照顧兒女又心煩.
吃得不好,睡得不好,又欠缺做運動.
血液循環系統和免疫系統又受損.
心血管又發炎.
患癌的風險亦增加.
簡單來說 離婚死得比較快.
但如果能善待婚姻 就會長命.
所以有句說話叫Happy wife, happy life.
你老婆開心 你的婚姻開心 你的命長久.
所以要善待你的婚姻.
但重要的是不要詭詐待你的妻子.
不要詭詐待你的婚姻.
因為婚姻牽涉到整個系統 包括你這一代和你的下一代.
如果我們輕看婚姻 不知不覺意圖詭詐待妻三步谷.
只會帶來你自己,你的另一半,你的下一代很大的傷害.
所以神很想提醒我們要好好建立我們的婚姻.
營造我們的婚姻.
這才是上帝期待的生活.
最後總結部分.
有一個叫做脆弱國家指數.
脆弱國家指數是評估一些國家.

$^{1401}$透過四大領域12個指標去評估一個國家的脆弱程度.
即是說這個國家是很危機的 隨時會倒塌.
用這些指標去評估.
當中哪個國家是最高風險,最脆弱的呢?.
就是索馬里,蘇丹,敘利亞這些國家.
哪些國家是最不脆弱的呢?.
最後就是冰島,芬蘭和挪威.
這些指標就是在嘗試用幾個不同的指標去幫我們衡量一個國家.
會不會沒有了呢?會不會散了呢?.
今天我們也要衡量我們的生命.
會不會脆弱 會不會散了 會不會被破壞呢?.
所以我們思考到幾個重要的指標.
在社會層面裡面就是鬼咋待弟兄.
鬼咋待你的同胞的話 是會影響整個社會的風氣.
當你鬼咋待上帝的話 會影響你的信仰.
當你鬼咋待妻子的話 會影響你的家庭.
《瑪拉基書》裡面的百姓不知不覺影響了他們的社會,信仰,家庭.
整個國度不知不覺已經不在上帝的喜悅裡面.
所以《瑪拉基書》裡面帶來的那種對祭司和百姓的責備是真實的.
也是迫切的.
因為這個問題如果不認真處理的話 會帶來很大的破壞.
最後一句總結.
我們有一段聖經 大家一起讀一下這段聖經.
就是來自《尼加書》 另外一個先知書所講的說話.
我們一起讀一下這句經文 六章八節開始.
「西人啊!耶和華已指示你何謂善?他向你所要的是什麼呢?.
只要你行公義,好憐憫,存謙卑的心與你的上帝同行」.
求主幫助我們 我們能夠活出一個公義的社會.
一個憐憫的婚姻.
也能夠在上帝面前存一個謙卑的心.
以致我們能夠活出聖潔的國度.
我們一起祈禱.
最後幫助我們 因為我們知道你給我們不單止有君尊的祭司.
你也要求我們有聖潔的國度.
一個屬神理的國度出現.
願主幫助我們藉著《馬拉基書》的思考.
我們如何遠離逃避這個世代裡的詭詐行為和生活.
讓我們活出聖潔的生命.
以致能夠去榮耀你.
以致能夠在這個社會裡.

$^{1441}$更加讓人看見我們是誠實的一群.
我們是可信賴的一群.
以致更加多人願意來到神理面前去認識神理自己.
願主幫助我們.
也幫助我們一會兒下來有分享討論團契的時間.
祈禱奉耶穌的名字而求.
阿們.
\newpage



\section{哥林多後書 2:14}
\label{sec:032261}
\textbf{你是個得勝者 (華)}~~~~劉彤牧師
\newline
\newline
連結: \href{http://www.church.com.hk/acms/content.asp?site=cdc\&op=show\&type=product\&code=032261\&layout=sermon}{\texttt{032261}} ~~~~ 語音日期: 2026-06-07
\newline
\newline
\hyperref[sec:032260]{\small{< < < PREV SERMON < < <}}
~
\hyperref[sec:index_chronic]{\small{[返順時目]}}
~
\hyperref[sec:032262]{\small{> > > NEXT SERMON > > >}}
\newline
\newline
哥林多後書 2:14
\newline
\begin{longtable}{cl}
\hline
\hline
章節 & 經文 (和合本修訂版)\\
\hline
2:14 & \begin{tabularx}{0.7\textwidth}{X} 感謝神！他常率領我們在基督裡得勝，並藉著我們在各處顯揚那因認識基督而有的香氣。 \end{tabularx} \\ \\
[1ex]
\hline
\hline
\end{longtable}
$^{1}$讀一段聖經在哥林多厚書第二章第十四節.
只有一節的聖經我們一同來念.
並藉著我們在各處顯揚那因認識基督而有的香氣.
我們就讀到這裡.
今天我們真的生活在一個充滿危機的世代.
不但國際局勢動盪不安.
社會的次序價值體系越來越遠離聖經的原則.
而且在我們每天的生活當中.
我們也面對各式各樣的試探引誘挫折跟壓力.
做一個基督徒我不知道你怎麼樣來看.
你每天生活所遇見的挑戰呢.
其實在聖經裡面神早已經教導我們怎麼樣來面對.
這樣一個充滿危機的世代.
而且能夠靠主得勝.
首先我們需要明白.
神從來沒有應許我們在地上有一個平順安舒的日子.
約翰福音十六章三十三節.
你看見這個經文組很清楚的對我們說.
我將這些事告訴你們是要叫你們在我裡面有平安.
在世上你們有什麼.
在世上有苦難.
換句話說.
今天當我們遇見苦難的時候不應該很驚訝.
我們是會在地上遇到各式各樣的壓力跟危機的.
但雖然是如此.
主繼續怎麼說呢.
但你們可以放心.
為什麼呢.
因為我已經勝了世界.
Hallelujah.
今天無論環境怎麼樣的改變.
我們一定要確信我們的神.
祂不但是坐在寶座上的主宰.
祂更是勝過世界的主宰.
Amen.
那也許你會說那跟我有什麼關係呢.
大有關係.
因為今天只要你是神的兒女.
你就是走在神得勝的行列當中的.
Amen.

$^{41}$在剛才我們所讀的經文裡面講到.
我們可以向世界來誇聖.
向仇敵來誇聖.
這個誇聖.
這個原文就是走在得勝的行列當中的意思.
換句話說.
今天因著耶穌基督已經得勝世界.
所以我們每一個在基督耶穌裡面的人.
也可以天天向我們的環境來誇聖的.
Amen.
那你說我們今天怎麼樣.
能夠在每天的環境當中靠主來得勝呢.
今天從四方面跟大家來分享.
首先第一個.
一定要靠著主的寶血得勝.
跟你們要靠著主的寶血得勝.
啟示錄十二章第十一節這裡說.
弟兄勝過他.
就指著魔鬼撒旦.
是因為高揚的血.
和自己所見證的道.
他們雖至於死也不礙其性命.
這裡提到主耶穌的寶血.
就是我們能夠得勝仇敵的重要關鍵.
為什麼呢.
因為今天我們之所以在生活當中.
充滿了憂愁苦澀怨恨嫉妒.
甚至在紛爭當中.
甚至遇見天災人禍基本的原因.
都是因為今天我們活在罪惡的瑕疵跟捆綁當中.
第二點今天無論我們對罪的看法是如何.
我們作為一個人.
我們不得不承認.
在我們裡面是有個力量.
叫我們所願意做的善.
我們不去做.
但我們所不願意做的惡.
我們反倒去做不是嗎.
當我在這裡說惡的時候.
倒不見得是指什麼大壞事.

$^{81}$很多時候可能只是在生活中的小事.
甚至只是一些不好的習慣而已.
就好像發脾氣.
我們明明知道不好.
既得罪人又傷自己的身體.
可是脾氣一來就好像山洪爆發一樣.
一發不可收拾.
脾氣發完了之後就很後悔.
哎呀剛才怎麼回事情.
下次再也不無緣無故發脾氣了.
可是下次呢.
還有下次.
一次又一次.
好像我們永遠勝不過那在我們裡面的力量不是嗎.
不只是發脾氣如此而已.
如果我們肯察驗在我們裡面的生命.
我們就會發現有很多的事情.
一直到現在都在挾制我們.
讓我們無法勝過.
有些時候是一些不好的生活習慣.
長期的束縛我們.
有時候是我們心中一些怒氣.
特別在這樣一個動盪不安的世代.
多少的壓力多少的束縛.
讓我們心中好像有一種無名的怒火在裡面.
有些時候是過去偶像的敬拜.
那種偶像的詮釋到今天還挾制我們的生命.
有時候是一些過去所受的傷害.
那種傷害的陰影.
多少受過去的傷害到如今都挾制了我們的心靈.
讓我們心靈無法得以自由.
這為什麼主耶穌基督.
要在十字架上為我們成就偉大的救恩的原因.
因為唯有耶穌基督的寶血.
能夠赦免我們的罪.
而且斷開一切黑暗的詮釋.
約翰一書第一章第九節非常寶貴的經文.
我們一聽我們來唸.
我們若認自己的罪.
神似信使是公義的.

$^{121}$必要赦免我們的罪.
洗盡我們一切不義.
你要知道主寶血的赦免.
不是像我們所想的那種存檔式的赦免.
什麼叫存檔式的赦免.
就好像我們夫妻吵架之後的和好一樣.
那個都是存檔式的和好不是嗎.
這一次就饒了你.
下次給我記住.
一切都在檔案當中.
記得清清楚楚.
那種聖經上說主的寶血的赦免.
是一種塗抹式的赦免.
哥羅斯第二章十四節非常的重要.
一同來念來.
又塗抹在律例上所寫.
攻擊我們有礙於我們的自居.
把它撤去.
釘在十字架上.
換句話說.
當你我承認我們的罪跟軟弱的時候.
主就要完全的赦免我們.
神就要把我們一切的紀錄都除掉.
你可以想像.
當你在神面前認罪之後.
當魔鬼再一次到神的面前想控告你的時候.
你可以想像神怎麼對魔鬼來說嗎.
神會向魔鬼說什麼罪啊.
白紙一張.
Hallelujah.
白紙一張.
Oh Dear Ones.
你們多好受在我們生命當我們被這種罪跟罪惡感.
挾持捆綁.
卻不知道今天當我們在神面前承認我們罪的時候.
主的寶血就塗抹.
一切的過犯.
我們在神面前就是白紙一張.
這為什麼聖經說在基督裡面要成為.
新造的人.

$^{161}$New Creation.
就像重新被創造的人一樣.
是白紙一張.
Hallelujah.
這為什麼今天我們可以勝過一切的仇敵.
勝過一切環境的原因.
因為在耶穌基督裡面我們可以坦然無懼了.
主的寶血已經塗抹我們一切的過犯.
Hallelujah.
弟兄姊今天你可以不再被過去那種罪惡感.
羞愧感所挾持.
你也可以不再被仇敵所控告困擾.
每當仇敵再來控告你來攻擊來繳繞你的時候.
你能夠站立在主寶血的根基上對魔鬼說.
退去吧.
Hallelujah.
他就不能不離開你的.
Amen.
Hallelujah.
主的寶血能夠幫助我們勝過一切過去的捆綁跟殘餒.
並且帶給我們生命一個嶄新的開始.
Amen.
第二要靠著主的名來得勝.
跟你們說要靠著主的名來得勝.
腓立比書第二章第九到第十節這裡說.
所以神將他身為至高又賜給他超乎萬民之上的名.
叫一切在天上地上地底下因耶穌的名無不屈膝.
名字是非常重要的.
名字是代表權柄.
這裡提到我們主耶穌的名字.
是超乎萬民之上的名.
換句話說.
沒有一個權柄不是降伏在耶穌基督的權柄之下.
保羅在羅馬書裡面更是提到.
連這個世界上一切的權柄.
一切的政權也都是在神的權柄之下.
Amen.
所以在我們的生命當中要來高舉主的名字是太重要的事情.
你說那我們怎麼樣來高舉主的名呢.
那就要透過敬拜跟讚美來宣告耶穌是主.

$^{201}$當我們敬拜讚美的時候.
主的寶座就要在我們中間設立.
神的主權就能夠自由的運行在我們中間.
你要知道這是魔鬼撒旦最不喜歡我們做的事情.
魔鬼撒旦總是喜歡對我們說.
只要你自己知道就好了.
只要你知道耶穌是主就好了.
不要張揚了.
不要什麼宣告到處去宣告.
就好像我們中國人這種祖傳秘方的觀念一樣.
既不講又不傳留著自己用就好了.
結果最後什麼都沒有.
No.
我們看見今天我們擁有神給我們的恩典跟權柄.
在我們生活每一個層面裡面.
我們就要來宣告耶穌是主.
Hallelujah.
耶穌今天仍然掌權.
神的國必然降臨.
神的旨意必定成就.
Hallelujah.
你說牧師這樣的宣告有用嗎.
我們在魔鬼面前.
這麼好像手無抓機之力的.
還又喊又唱.
像個娃娃兵一樣有用嗎.
當然有用.
詩篇第八篇第二節非常寶貴的經文.
我們一同來念.
你因敵人的緣故.
從嬰孩和吃奶的口中建立了能力.
使仇敵和報仇的閉口無言.
請問誰是這裡所說的嬰孩和吃奶的.
就是我們.
我們怎麼去面對魔鬼撒旦.
我們這麼軟弱無力.
怎麼能勝過仇敵魔鬼撒旦.
但神說這就是我的美意.
我就是要藉著這些看起來非常軟弱的人的口中.
什麼意思什麼是口中.

$^{241}$口中就不是我們在神面前所發出的敬拜讚美跟宣告嗎.
當我們在神面前發出敬拜讚美宣告的時候.
神就要從我們口中建立能力.
藉著我們口中發出來的敬拜跟宣告.
除滅魔鬼撒旦一切的作為.
叫一切仇敵閉口無言.
換句話說本來我們是不可能勝過魔鬼撒旦不可能.
但什麼時候今天我們開口敬拜讚美高舉主的名的時候.
神就要做奇妙的事情.
神就要為我們興起為我們來征戰除滅魔鬼一切的作為.
Amen Hallelujah.
我最喜歡提到萬國主教佈道團他們同工的見證.
有一次他們在喜馬拉雅山的山區上傳福音.
他們來到一個村落.
挨家挨戶想要發單張的時候.
沒想到每一個當地的村民他們都拒絕接受這個單張.
看到這些同工來就把門通通關上.
所以他們沒有辦法他們就回到他們所居住的地方.
他們就在他們居住的地方在那敬拜讚美神.
在敬拜整個晚上禱告宣告耶穌是這個村莊的主宰.
結果到清晨的時候就聽到有人敲門的聲音.
他們就打開門一看.
所有的村民已經排隊一直排在他們門口.
一個一個對他們說昨天晚上我們在睡覺的時候.
主耶穌向我們顯現.
告訴我們耶穌是真神.
叫我們今天早上來到你們會告訴我們關於耶穌的真理.
Hallelujah.
我們看到我們敬拜讚美高舉主名的時候.
我們就能得勝一切捆縮我們的環境.
Amen.
所以同樣的弟兄姊妹.
今天當你在面對人生各樣困境的時候.
你一定不要先想盡自己的辦法.
不要一天到晚想用盡自己的方式來解決.
你一定要先來敬拜神.
你一定要宣告耶穌是主.
不但是為你的生命來宣告.
要為你的家庭來宣告.
也要為你的職場你的工作來宣告.

$^{281}$我們看見在七月份我們馬上要有耶利哥行動.
我們也特別要為我們的社區來宣告.
為我們的國家來宣告.
耶穌是主.
Amen.
什麼時候當我們開始高舉主的名的時候.
主能夠轉移我們的眼目.
不只是看見我們的環境而已.
更是要轉眼看見我們的主正坐在寶座上.
他的手正在工作.
神正在為我們征戰.
神正在為我們開道路.
Hallelujah.
Amen.
第三.
要靠著主的話語得勝.
各位靠著主的話語得勝.
以弗所說六章十七節.
我們一同來念這個經文.
並戴上救恩的頭盔.
拿著聖靈的寶劍.
就是神的道.
以弗所說第六章特別講到屬靈的軍裝.
如果你有在背聖經.
我們背經行從一月份.
我們就背這個屬靈的軍裝.
記得嗎.
在這個屬靈的軍裝當中.
你仔細看就發現很有意思.
其實保羅講到神的話的時候發現.
保羅所提到的軍裝多半是指的防衛性.
頭盔是保護頭.
護心鏡是來遮胸.
藤牌是擋火箭.
只有神的話是例外.
神的話就是聖靈的寶劍.
是用來攻擊仇敵的兵器.
這為什麼背誦神的話這麼重要.
你要把神的話存放在心裡.
你才能夠靈活的運用神的話.

$^{321}$來制敵能夠制勝.
所以我們看見背經行動是很重要.
六月份是我們教會背經行動到一半的時刻.
我不知道你跟上沒有.
你要知道.
你剛剛在看新聞裡面講到.
我們的獎品是一個榮譽紀念碑.
我看過那個sample 漂亮極了.
值得你背完12個月的經文.
能夠放在你的家裡.
最明顯的告訴大家說.
我背過聖經了.
背經非常漂亮.
非常重要不只是為那個獎品.
而是神的話存放在你心中的時候.
你就是能夠制勝仇敵.
為什麼神的話是一種.
得勝仇敵的兵器.
因為神的話是信實的.
神的話一出絕不突然返回.
神必定要親自成就他自己所說的話.
而且無論從人的角度來講是多麼不可能的事.
在神永遠沒有難成的事.
因為神所應許的他必成就.
這就兵器的意思.
你發現說神既然話已出來.
他一定會成就.
無論仇敵多麼厲害.
抵擋都沒有用的.
只要神說的.
神所應許的一定會發生.
一定會成就.
這就好像聖經裡面所多次神預言.
以色列人會回歸到神所應許之地的.
回到他們的土地上.
從人的角度來看.
一個亡國兩千多年的國家.
怎麼可能再復國.
但神的話一出絕不突然返回.
他就必成就他自己的話.

$^{361}$結果果然以色列人.
在1948年能夠重新復國了.
哈利路亞.
所以今天我們如果想在生活當中.
經歷神的得勝.
我們就必須在生活當中.
選擇來信靠神的話語.
你們要選擇信靠神的話語跟他說.
你知道我常聽到人說.
但是牧師我聽不見神的聲音.
其實根據我這麼多年的觀察.
問題從來不是你沒有聽見神的聲音.
其實你常常聽見神的聲音.
有時候神是藉著講道.
聽道的時候向你來說話.
有時候是藉著你在查經的時候.
神向你說話.
有時候是藉著弟兄姊妹分享的時候.
向你說話.
甚至中年以上的弟兄姊妹.
神常藉著你的身體向你說話.
血脂血糖血壓三高.
向你發出警告.
不要再虛度人生.
跟你們說不要再虛度人生.
好好為主而活.
跟他說.
但問題是當你聽見神的話語的時候.
卻不見了去選擇信靠神的話.
然後再次說很多時候真正的問題.
不是在於你聽不見神的聲音.
真正的問題是當你聽見了.
你卻沒有選擇信靠神的話語.
就像當年的以色列人.
他們要準備進神的應許之地的時候.
還記得派了十二個探子.
到迦南美地去窺探那地.
十二個探子回來兩種報告.
一種報告說那裡像巨人.
我們像蚱蜢.

$^{401}$我們不可能得到那地.
第二種報告說雖然那裡像巨人.
可是這是神應許之地.
我們一定能夠得地為業.
如果你是當時的以色列人.
你會相信哪一種報告.
信靠哪一種報告.
你要知道當時的以色列人.
不是沒有聽見神的聲音.
他們聽見了.
只是他們沒有選擇去信靠神的話語.
所以結果他們不肯進入迦南地.
結果就在曠野漂流了四十年.
你可知道以色列人.
其實是不需要在曠野漂流四十年的嗎.
只不過因為他們選擇.
相信了世界的聲音.
卻沒有選擇相信神的應許.
第二點今天在這樣一個充滿動盪不安的世代.
我告訴你如果你想經歷神的得勝.
你一定要選擇信靠神的話.
Amen.
Hallelujah.
最後第四你要靠著主的聖靈來得勝.
跟你們說靠著主的聖靈來得勝.
第二點今天這位聖靈.
不但是內住在我們生命當中.
與我們同在的神.
聖靈更是賜下能力給我們的神.
使徒信仰第一章第八節.
怎麼說我們一同來念來.
但聖靈降臨在你們身上.
你們就必得著能力.
便要在耶路撒冷猶太全地撒瑪利亞.
直到地極做我的見證.
我們可以講得很清楚.
聖靈降臨在我們身上就給我們能力嗎.
什麼樣的能力.
首先是一種生命改變的能力.
如果你仔細看整本使徒行傳.

$^{441}$你就會發現在使徒行傳當中.
最大的神蹟不是在於醫病趕鬼.
其實最大的神蹟.
來自於生命改變的見證.
你還記得使徒行傳第二章.
當聖靈降臨下來之後.
彼得跟其他幾個使徒就站起來.
高聲地向當時已逝的人來傳福音.
說老實話.
當時彼得是最不配站起來說話的人.
不是嗎.
他不久前才三次不認主.
從人的角度來講.
彼得應該站到後面去.
安靜點.
不要說話了.
但當聖靈降臨下來之後.
一切就不一樣了.
彼得勇敢地站起來.
高聲地說你們要聽我說.
結果他一開始講到.
三千人歸向耶穌基督.
Hallelujah.
到了使徒行傳第四章.
他被官府抓住.
而且恐嚇他們說.
不可以再傳講耶穌的名.
若是以前的彼得.
早就躲藏起來了.
可是這個時候的彼得不一樣了.
使徒行傳第四章19到20節.
這裡說彼得約翰說.
聽從你們.
不聽從神.
這在神面前合理不合理.
你們自己琢磨吧.
我們所看見所聽見的.
怎麼樣.
不能不說.
哇.

$^{481}$完全不一樣的彼得.
若是從前的彼得.
根本就躲在一個地方不敢出來.
這個時候彼得為神.
大膽勇敢地傳講神的道理.
我們看見聖靈降臨下來.
不但改變了彼得的生命.
聖靈降臨下來.
也改變了每一個.
聽見主福音人的生命.
那天彼得站起來傳講福音.
三千人歸向耶穌基督.
我們也看見三千個生命.
被改變人的生命.
以至於後來在整個耶路撒冷城.
帶來翻天覆地的影響力.
今天我們若是想翻轉我們的城市.
影響這個世代.
我們必須依靠聖靈的大能.
其次.
聖靈也帶給我們勝過環境的能力.
特別是今天我們面對.
這樣充滿危機的時代.
我們真的要呼求聖靈來幫助我們.
因為聖靈就是我們的幫助者.
跟你們說.
聖靈就是我們的幫助者.
因為祂是我們的Helper.
幫助者.
你說怎麼幫助.
反而首先.
聖靈能夠打開我們的眼睛.
跟你們說.
聖靈能夠打開我們的眼睛.
因為很多時候我們的眼目很狹窄的.
我們只看見我們的環境不是嗎.
很多時候當我們看見環境.
我們的驚慌失措.
我們覺得心灰意冷.
就像當時.

$^{521}$跟祂的僕人在多丹城一樣.
當整個多丹城被敵人四面圍攻的時候.
清早起來.
那個僕人看見敵人四面圍繞.
就匆匆忙忙到以利沙說.
以利沙怎麼辦.
以利沙就為她的僕人禱告.
求神開她僕人的眼睛.
讓她能夠看見.
神就開她僕人的眼睛.
她僕人眼睛一打開.
就看見原來天上天使天君四面圍繞.
數目比敵人還多.
Hallelujah.
同樣今天聖靈來.
就是要幫助我們打開我們的眼睛.
否則的話你所看見就是你的環境.
很多時候你就沒有辦法能夠看見神的同在.
只有聖靈降臨在你身上.
打開眼睛看見原來天使天君四面圍繞.
原來神正在與你同在.
原來神的手正在工作.
Hallelujah.
當你看見了就不一樣了.
當你看見了.
你就有信心往前走.
你就可以不再懼怕.
不再懼怕仇敵的攻擊跟攪擾.
因為你確信神正在與你同在.
Hallelujah.
聖靈來也給我們勇氣.
能夠踏出信心的行動.
很多時候我們不是不知道神的心意.
也不是不知道神的應許.
但很多時候就是沒有勇氣.
很多時候沒有勇氣做當行的事情.
我們也不是不知道神要我們做什麼.
但沒有勇氣去跟隨向前走.
但我告訴你聖靈一來.
祂就要賜給我們智慧與勇氣.

$^{561}$來得勝一切的環境.
就像當年的門徒一樣.
當年的門徒他們領受大使命.
他們知道了.
他們知道神的心意讓他們直到地極.
去傳揚主的福音.
但他們很徬徨.
怎麼去法.
也不知道怎麼能夠傳福音給萬民.
那些門徒他們一輩子沒有走離開過猶太權力.
什麼叫地極根本很難想像.
很徬徨怎麼去法.
很遲疑.
所以天使提醒他們回到耶路撒冷.
等候聖靈的降臨.
而聖靈一降臨下來的時候.
真的帶來翻轉的勇氣跟能力.
福音就真的從耶路撒冷.
猶太權力 撒瑪利亞傳到地極.
聖靈是我們的幫助者.
我們要歡迎聖靈.
要依靠聖靈的大人.
第二點今天在基督耶穌裡面.
你就是個得勝者.
給你們好處.
你就是個得勝者.
無論你怎麼看.
你自己說我這麼軟弱.
怎麼可能是個得勝者.
但不要忘記哥林多後書十二章第九節的經文.
太重要太重要.
這裡說他對我說我的恩典夠你用的.
因為我的能力是在人的什麼.
軟弱上顯得完全.
所以我更喜歡誇自己的軟弱.
好叫基督的能力復辟我.
第二點神的能力是要在人的軟弱上顯得完全.
但我真的提醒你.
是在人承認自己軟弱的地方上顯得完全的.
所以保羅說我喜歡誇我的軟弱.

$^{601}$什麼意思.
這個誇就是承認.
我喜歡承認我的軟弱.
好叫基督的能力才能復辟我.
所以今天當你承認你是軟弱無助的時候.
我告訴你.
你是最有資格來經歷神的得勝的.
今天無論你所面對的環境是充滿多少的危機.
你今天要來專心的來依靠神.
要靠祂的保險斷開一切罪惡的困鎖.
靠祂的名字勝過一些仇敵的攻擊.
靠著祂的話語能夠經歷神的信使.
神所應許的必定成就.
靠著祂的聖靈經歷神的大能.
得勝又得勝.
靠著主我們必然能夠勝過世界.
我們一起禱告.
親愛的主.
因為雖然我們人是這樣的軟弱.
你的能力就是在我們的軟弱上顯得完全的.
在我們承認我們軟弱的地方.
你的能力就能夠完全的彰顯.
主要今天早上讓我們重新謙卑的回到你面前來.
承認我們就是那個軟弱無助的.
也讓今天早上讓我們能夠經歷你得勝的大能.
我們謝謝主.
我們將一切榮耀都歸給主你自己.
奉耶穌得勝的名祈求.
阿們 阿們.
\newpage



\section{約書亞記 2:1-24}
\label{sec:032262}
\textbf{恩典的記號}~~~~李業基牧師
\newline
\newline
連結: \href{http://www.church.com.hk/acms/content.asp?site=cdc\&op=show\&type=product\&code=032262\&layout=sermon}{\texttt{032262}} ~~~~ 語音日期: 2026-06-07
\newline
\newline
\hyperref[sec:032261]{\small{< < < PREV SERMON < < <}}
~
\hyperref[sec:index_chronic]{\small{[返順時目]}}
~
\hyperref[sec:032263]{\small{> > > NEXT SERMON > > >}}
\newline
\newline
約書亞記 2:1-24
\newline
\begin{longtable}{cl}
\hline
\hline
章節 & 經文 (和合本修訂版)\\
\hline
2:1 & \begin{tabularx}{0.7\textwidth}{X} 嫩的兒子約書亞從什亭暗中派兩個人作探子，說：「你們去窺探那地和耶利哥。」於是二人去了，來到一個名叫喇合的妓女家裡，在那裡睡覺。 \end{tabularx} \\ \\ \relax
2:2 & \begin{tabularx}{0.7\textwidth}{X} 有人告訴耶利哥王說：「看哪，今夜有以色列人到這裡來窺探此地。」 \end{tabularx} \\ \\ \relax
2:3 & \begin{tabularx}{0.7\textwidth}{X} 耶利哥王派人到喇合那裡， 說：「你要交出那來到你這裡、進了你家的人，因為他們來是要窺探全地。」 \end{tabularx} \\ \\ \relax
2:4 & \begin{tabularx}{0.7\textwidth}{X} 但女人已把二人藏起來，卻說：「那兩個人確實到我這裡來過，他們從哪裡來，我卻不知道。 \end{tabularx} \\ \\ \relax
2:5 & \begin{tabularx}{0.7\textwidth}{X} 天黑、要關城門的時候，他們就出去了。他們往哪裡去我也不知道。你們趕快去追他們，就必追上。」 \end{tabularx} \\ \\ \relax
2:6 & \begin{tabularx}{0.7\textwidth}{X} 其實，這女人已經領二人上了屋頂，把他們藏在她擺列在屋頂的亞麻梗中。 \end{tabularx} \\ \\ \relax
2:7 & \begin{tabularx}{0.7\textwidth}{X} 那些人就往約旦河的路上追趕他們，直到渡口。追趕他們的人一出去，城門就關了。 \end{tabularx} \\ \\ \relax
2:8 & \begin{tabularx}{0.7\textwidth}{X} 二人還沒有睡之前，女人就上屋頂，到他們那裡， \end{tabularx} \\ \\ \relax
2:9 & \begin{tabularx}{0.7\textwidth}{X} 對他們說：「我知道耶和華已經把這地賜給你們了，並且我們也都懼怕你們。這地所有的居民在你們面前都融化了。 \end{tabularx} \\ \\ \relax
2:10 & \begin{tabularx}{0.7\textwidth}{X} 因為我們聽見你們出埃及的時候，耶和華怎樣在你們前面使紅海的水乾了，並且你們怎樣處置約旦河東的兩個亞摩利王，西宏和噩，把他們完全消滅。 \end{tabularx} \\ \\ \relax
2:11 & \begin{tabularx}{0.7\textwidth}{X} 我們一聽見就膽戰心驚，人人因你們的緣故勇氣全失。耶和華－你們的神是天上地下的神。 \end{tabularx} \\ \\ \relax
2:12 & \begin{tabularx}{0.7\textwidth}{X} 現在我既然恩待你們，求你們指著耶和華向我起誓，你們也要恩待我的父家。請你們給我一個確實的憑據， \end{tabularx} \\ \\ \relax
2:13 & \begin{tabularx}{0.7\textwidth}{X} 要救活我的父母、兄弟、姊妹，和所有屬他們的，拯救我們的性命脫離死亡。」 \end{tabularx} \\ \\ \relax
2:14 & \begin{tabularx}{0.7\textwidth}{X} 那二人對她說：「我們願意以性命來替你們死。你們若不洩漏我們這件事，當耶和華將這地賜給我們的時候，我們必以慈愛和誠信待你。」 \end{tabularx} \\ \\ \relax
2:15 & \begin{tabularx}{0.7\textwidth}{X} 於是女人用繩子把二人從窗戶縋下去，因為她的屋子是在城牆邊上，她也住在城牆上。 \end{tabularx} \\ \\ \relax
2:16 & \begin{tabularx}{0.7\textwidth}{X} 她對他們說：「你們暫且往山上去，免得追趕的人遇見你們。要在那裡躲藏三天，等追趕的人回來，你們才可以走自己的路。」 \end{tabularx} \\ \\ \relax
2:17 & \begin{tabularx}{0.7\textwidth}{X} 二人對她說：「你叫我們所起的誓與我們無關， \end{tabularx} \\ \\ \relax
2:18 & \begin{tabularx}{0.7\textwidth}{X} 除非，看哪，當我們來到這地的時候，你把這條朱紅線繩子繫在縋我們下去的窗戶上，並要叫你的父母、兄弟和你父的全家都聚集在你家中。 \end{tabularx} \\ \\ \relax
2:19 & \begin{tabularx}{0.7\textwidth}{X} 凡離開你家門往街上去的，他的血必歸到自己頭上，與我們無關；凡在你家裡的，若有人下手害他，他的血就歸到我們頭上。 \end{tabularx} \\ \\ \relax
2:20 & \begin{tabularx}{0.7\textwidth}{X} 你若洩漏我們這件事，你叫我們所起的誓就與我們無關了。」 \end{tabularx} \\ \\ \relax
2:21 & \begin{tabularx}{0.7\textwidth}{X} 女人說：「就照你們的話吧！」於是她送他們走了，就把朱紅繩子繫在窗戶上。 \end{tabularx} \\ \\ \relax
2:22 & \begin{tabularx}{0.7\textwidth}{X} 二人離開，到山上去，在那裡停留三天，直等到追趕的人回去。追趕的人一路尋找，卻找不著。 \end{tabularx} \\ \\ \relax
2:23 & \begin{tabularx}{0.7\textwidth}{X} 二人回來，下了山，過了河，來到嫩的兒子約書亞那裡，向他報告他們所遭遇的一切事。 \end{tabularx} \\ \\ \relax
2:24 & \begin{tabularx}{0.7\textwidth}{X} 他們對約書亞說：「耶和華果然將那全地交在我們手中了，並且那地所有的居民在我們面前都融化了。」 \end{tabularx} \\ \\
[1ex]
\hline
\hline
\end{longtable}
$^{1}$「Be a Saint Angel」的講道系列.
做一個獅子心人之心的聖徒.
《約書亞記》第二章.
拉合的故事.
很多我們都.
其實如果回到教會一段時間.
基督徒一段時間都一定聽過.
這個妓女拉合的故事.
耶利哥城.
她怎樣救回她家的故事.
但是為了幫助大家去重溫一下這個故事.
可能當中的確有些可能還沒聽過.
就好像上次那樣.
首先我們看一段短片.
再強調就是.
這些故事有時候有藝術創作的成分.
這段會是用約書亞的角度來說.
聖經大家仍然要回去看.
不過大家看看這個短片.
(畫面).
拉合的故事其實是一個充滿弔詭的傳奇.
拉合似乎是一個用人的角度來說.
出了最不幸的情況.
他不只是住在一個神即將毀滅的耶利哥城.
當中的一個外邦人.
他還是一個妓女.
這個故事講大話.
背叛了自己的民族.
是一個叛徒.
耶利哥城即將毀滅.
他不僅沒有阻止.
他還幫了一份.
幫了一些異色人.
是他的敵人.
按照聖經的描述.
拉合其實是順應天命.
成為了聖經當中最重要的一位信心英雄.
希伯來書11章說.
因著信妓女拉合曾友善地接待探子.
就沒有跟那些不信從的人一同滅亡.

$^{41}$拉合有的信心.
不是好像有些信條.
我願意信一些信條.
也不是做了一些什麼鎮禮.
割禮這些儀式的信.
雅各書第二章裡面說.
拉合的信心是透過行動.
透過勇敢活出來的.
雅各說.
同樣妓女拉合接待使者.
又放他們從另一條路出去.
不也是因行為稱義嗎.
所以就如身體沒有靈魂是死的.
信心沒有行為也是死的.
因為這樣的緣故.
在猶太人的傳統當中.
其實我們看到.
拉合雖然是一個外邦人.
但是他因為敬畏神而做的事情.
成為了信心的典範.
馬太科罕第一章第五節.
耶穌的家譜當中說到.
撒緩從拉合事生波亞斯.
波亞斯從路德事生俄彼得.
俄彼得生耶西.
耶西生大衛王.
大衛王從烏里亞的妻子生所羅門.
所以說到拉合這個名字的時候.
猶太讀者不僅看到一個信心的典範.
他們和拉合是有血緣的關係的.
但是值得留意的是.
拉合這個名字.
其實在舊約當中.
除了約書亞記第二章和第六章之外.
其他地方沒有提過.
新約當中.
也就是我們剛才讀過的三個地方.
馬太福音.
希伯來書.
和雅各書.

$^{81}$這三卷書是提到的.
三卷書.
大家如果熟悉聖經的人都知道.
是為猶太人的讀者所寫的.
所以我們可以推測.
在外邦人當中.
拉合這個故事.
可能比較沒有那麼廣為人知.
也都是可能容易被誤解的.
外邦人可能未必那麼能夠明白到.
這個故事背後的一些脈絡.
整個舊約.
我們叫做Metanarrative.
原序是故事的脈絡.
去理解這個故事.
所以有時候比較難解釋.
可能新約的作者沒有那麼多的引用.
缺少的可能性是什麼呢.
恩典.
拉合得救.
其實本乎恩典.
也是藉著信心的.
拉合之所以能夠進入耶穌的家譜裡面.
不是單純因為他的勇氣的行為.
好像贏取了一個Purple Heart.
一個勇氣勳章.
因為他其實是透過恩典.
來到耶穌的家譜當中.
什麼是恩典呢.
恩典就是得到一些本來你不配得.
白白得到的東西.
批重神白白得到而來.
所以我們有時候如果純粹是過分的.
去問一條問題.
譬如說拉合做了什麼好事.
或者拉合有什麼美德.
來看這個故事的時候.
來看他的《深淺凡人》的時候.
其實反而可能錯過了這個故事最重要的地方.
這是一個恩典的故事.

$^{121}$拉合生活在一個即將遭毀滅的民族當中.
拉合其實是一個職業性的罪人.
但是按理來說.
其實他應該好像他民族的其他任何一個成員一樣.
是一樣被毀滅的.
他獲得拯救.
只能夠說是一個白白得來.
神的恩典.
但其實恩典何嘗不是所有以色人的故事呢.
恩典何嘗不是所有教會我們每一個人的故事呢.
其實恩典就是人類的故事.
沒有人其實是配得得到這個永生的.
但是因為神的恩典.
我們不單止可以存活.
而且能夠得到生命的時候.
那個更豐盛更有意義的生活.
所以我們首先要開始的時候.
我覺得我們看這個約書亞記第二章.
這個拉合的故事的時候.
我們首先要有一點背景的眼光在當中.
因為如果你不用這個恩典的角度做這個前提的時候.
你未必可以明白到隱含的意義.
其實就算花很多時間研究這個故事的學者當中.
有時候也未必看得到其實恩典這件事在當中的重要性.
甚至有些學者可能先入為主.
有些問題帶入這個故事當中的時候.
就用一個不同的角度去看這個故事.
我最近在夏季在神學院教論理.
這段時間.
拉合這個故事.
如果你看很多基督教論理的書籍.
都會引用為一個兩難之間.
這個所謂Moral Dilemma的一個案例來研究.
上個星期五我和Focus組做了一個討論.
就是用這個故事研究一下這個道德兩難的議題.
很多時候的而且確我們不知不覺.
我們就用了這個角度去看拉合的故事.
我其實在金門神學院這一課.
我其實用Zoom的.
我錄了.

$^{161}$如果大家有興趣再研究一下這個道德兩難之間的問題的時候.
可以問我拿那個Zoom Link.
大家可以去看一下.
但按照這個倫理學的案例研究的時候.
很多時候大家就看拉合其實是夾在兩個絕對的道德命令之間.
一方面其實他有個義務對那些官員誠實.
但他又覺得有個義務要保護這些探子的性命.
做人誠實當然是我們一個最基本的道德要求.
是一個普遍性的道德原則.
所以耶穌說你們的話是就說是,不是就說不是.
若再多說就是出於惡者.
但問題是如果拉合來查問那些官員說真話的時候.
就會供那些探子出來犧牲了他們的性命.
因為這個緣故很多時候研究倫理學的學者.
用各樣的道德理論來解釋拉合的謊言為他辯護.
有些人認為拉合是兩害取其輕.
lesser of two evil.
他選擇了比較小的惡,就是說謊話.
雖然是有錯但是值得原諒.
另一些派別的人認為當人面臨一個保護性命.
和誠實兩項其實是相互衝突的絕對道德要求規則的時候.
人應該選擇更加高的原則保存性命.
如果你這樣做的時候不是犯罪的.
而且是對的,因為是規則之間你選擇一個更高的原則.
也有人說其實這個是錯的.
拉合其實沒有想清楚,他只是不有信心.
他如果想清楚的時候,或者他照樣講實話也好.
不講說話也好,神一定會幫他開一條出路的.
人在試探的當中神一定不會讓我們不能夠勝過的.
所以這些就是基督教倫理拿這個案例來做爭論的一個題目.
其實為什麼會這樣看呢?.
多多少少是因為第二次世界大戰的時候.
有很多很著名的故事.
在歐洲裡面有很多人為了救那些猶太人.
面對那些納粹的官員的時候說謊話.
譬如說Cory Ten Boom和Orca Schindler的故事.
兩個都是被改編成為電影的.
大家都可能聽過或者看過了.
探討一些道德兩難的問題的時候.
這些故事,尤其是第二次世界大戰當中的故事.

$^{201}$經常被作為案例的研究.
最近這幾年香港的社運都發生了很多類似的故事.
那些人藏起學生,藏起病人,記者隱瞞他們.
我想再過十年八年就會拿來做案例研究了.
所以其實很自然我們會用這三種不同的道德倫理的觀念.
去看這些現代的案例.
那些人想說完這個現代案例的時候.
想找一些聖經的支撐.
結果就去聖經裡面找一些人說謊.
但是神好像反而賜福他們.
拉俠和埃及那兩個修生婆就是最著名的例子.
作為一個聖經印證,似乎有時候是可以說謊的.
但是我覺得大家去聽這些解釋.
說一些倫理學家去說這些問題的時候.
要有一個原則.
當我們去按正義分解真道看聖經的時候.
我們不能夠太生硬去引用聖經.
其實很多時候你不一定要引用聖經.
如果你想討論一下這三個道德倫理的原則的時候.
如果我們拿一些現代的例子,例如Coryten Buma,Orcash Center來說.
拉俠的故事其實不是一個很恰當的對照.
其實拉俠這個故事實際上說的是.
他透過信心得到恩典.
不是意味著他在這個過程當中做的每一件事.
都是正確而且是值得效法的.
如果我們太過努力去將別人信心的故事每一個環節都美化.
或者叫做whitewash.
將它變成一個什麼事情都做得對的二人的時候.
可能反而是曲解了,誤解了恩典的精髓.
約書亞記第二章.
恩典才是貫徹這個故事的主題.
剛剛過了星期五晚.
我們邀請了盛曉梅姊妹和泥土音樂.
做這個2026年幸福的音樂會.
我們多加了幾十張椅子.
估計我們有450人參加.
將近20人走到台前.
決志接受耶穌基督的福音.
其中一首歌是盛曉梅姊妹一首非常著名的歌曲.
恩典的記號.

$^{241}$剛才我們唱的第三首歌就是這首歌.
這首歌第一段的歌詞是這麼說的.
站在大海邊才發現自己是多渺小.
登上最高山才發現天有多高.
浩瀚的宇宙中我真的微不足道.
像灰塵消失也沒人知道.
用國語唱比較有這個味道.
但大家聽到是押韻非常好的歌詞.
但這就是很多人的醒覺.
我們人生很多的努力.
三分幸運七分的努力.
人生當中去追求成功.
考到某些學校.
得到某些職業.
買了某些房子.
銀行帳戶裡有多少個零.
到達了人生的高峰.
但接下來又如何.
你已經走到這個高峰又如何.
這個就是所羅門王.
晚年所發現的事情.
他追求成功的過程當中.
或者很有趣的.
很值得的.
這些流的血汗都是很有意義的.
而且幫到人.
但最終所羅門看回頭的時候.
他說虛空的虛空.
凡事都是虛空.
都不過是暴風捉影.
因為人呢.
始終是有一個生命的終點.
到你未預備好的時候.
這個生命的終點就來到你那裡.
到這一刻的時候.
你走到盡頭.
你看回頭.
這一切又有什麼分別呢.
之後還有誰記得你做的事呢.
人總是有些時候.

$^{281}$未去到那一點.
都有機會可能會停下腳步.
想下思考下.
花了幾十年的寒暑.
幾十年的精力.
究竟你人生得到些什麼呢.
其中一些的時候.
就是當人走到死亡邊緣的時候.
耶利哥城.
當時其實是一個.
在迦南地非常富庶.
而且是無堅不摧的城市.
但是當耶利哥城的人.
聽到以色列人過約旦河.
差不多兵臨城下.
他們好像現在的人在醫生那裡聽到.
有一個末期的癌症一樣.
晴天霹靂.
心都融化了.
就正如拉合和那兩個探子所說.
我知道耶和華已經把這地賜給你們了.
並且我們也都懼怕你們.
你們所有的居民.
在你們面前都融化了.
因為我聽見你們出埃及的時候.
耶和華怎樣在你們前面使兇害的水乾了.
並且你們怎樣處置約旦河東的兩個.
阿摩利王西弘王國.
把他們完全消滅.
我們一聽見就膽戰心驚.
人人因你們的緣故.
勇氣全失.
耶和華.
你們的神是天上地下的神.
很多時候我們從以色列人的角度去看這個故事.
但你們有沒有試過從耶利哥人的角度去看這個故事呢.
那些以色列人是癌症.
要吞食你.
當人們知道末日將到的時候.
會做些什麼呢.

$^{321}$最自然的反應就是.
反過來嘗試去反抗.
這叫做求生本能.
這顯然就是耶利哥王想做的事.
現在快要到場了.
軍隊都不夠你打了.
於是拼死一搏.
在不同的方式當中.
我們都在拉合身上看到這個恐懼.
甚至可能有些人認為.
其實拉合都算不上什麼信心英雄.
他只不過是設法保住自己的性命.
勇氣很多時候是滴出來的.
只不過是求生本能的另一個名字.
正如我上次講到的時候提到.
六夜仙蹤裡面有隻沒有膽的獅子.
勇氣不是沒有恐懼.
勇氣是願意面對你所恐懼的事物.
有時候危險殺到你面前的時候.
你沒有願意.
你都沒有辦法.
你都要去面對.
不過如果我們停下來想一想的時候.
這種恐懼.
這種求生本能是從哪裡來的呢?.
相信有神的人就明白.
求生本能其實是來自神的.
如果不是神將這個求生的本能.
放在動物當中的時候.
萬物當中.
人就不可能有足夠的恐懼去求生.
所以對信神的人來說.
這個求生本能是一種神的恩典的開始.
我們在拉合身上所看到的第一個原則就是.
人是因敬畏懇求恩典.
拉合不是因為覺得自己很佩得這個恩典而向神求.
他沒有跟這些探子說.
我是一個好人.
我幫過很多人.
我好像這個城裡面其他人死有餘辜.

$^{361}$他沒有這樣說.
我懷疑即使是一個耶利哥城的標準來說.
拉合都不是特別算什麼好人.
其實我真的很不明白.
為什麼那些學者研究拉合的故事.
總是要說到拉合好像很誠實.
是他的矛盾.
他做什麼職業呢?.
寬敞無真愛.
說謊對他人有什麼困難呢?.
拉合懇求這個恩典.
不是因為他覺得他自己是好人.
能夠為自己所作所為去辯護.
他懇求恩典.
唯一的原因就是他害怕他怕死.
他敬畏神.
第十二節他繼續說.
現在我既然恩待你們.
求你們指著耶和華向我起誓.
你們也要恩待我的父家.
請你們給我一個確實的憑據.
要救活我的父母兄弟姊妹和所有屬他們的.
拯救我們的性命脫離死亡.
拉合懇求恩典當然是怕死.
但是他也承認懂得用常理來說話.
倫理學當中我們有所謂的黃金論.
這個Golden Rule.
是什麼呢?.
只要豈所不欲,勿施於人.
或者好像摩西所說.
耶穌所說要愛倫社羽和愛自己.
這個就是拉合的道德推理.
很簡單.
我不想死.
我知道你們也不想死.
不如我們談談有沒有辦法可以一起不死.
你如果善待我,我又善待你.
你對我好,我對你好.
如果我們川普總統說的話.
這個其實就是.

$^{401}$引出了其實恩典的第二個要點.
因慈愛承諾恩典.
在人際關係當中.
其實慈愛這個憐憫.
是很大部分建立在一個同理心當中.
拉合只是出於這個憐憫.
向這些探子承諾這些恩典.
他可以選擇不這樣.
他可以選擇攻出他們.
但他多多少少是因為憐憫他們.
當然他害怕.
但也是有憐憫他們.
他也是盼望探子能夠以憐憫去回報他們.
而那兩位探子也確實是這樣.
因此他們說.
我們願意以性命來替你們死.
你們若不洩漏我們這件事.
當耶和華將這地賜給我們的時候.
我們必以慈愛和誠信待你.
如同拉合一樣.
那兩名探子都是因慈愛承諾恩典.
但這兩位探子其實不只是進行一筆互惠互益的雙贏交易.
他們真的對拉合有同情.
拉合不是完全不懼怕耶利哥城的士兵.
但是因為他對耶和華神有更大的敬畏.
他就願意去做這個勇敢的行為.
其實這兩個探子都一樣.
他們都害怕.
但他們更加敬畏的是天上耶和華的神.
所以要打入耶利哥城.
好吧.
硬碰進去.
可以選擇.
我想他們可能選擇不去.
但他們因為敬畏耶和華神.
耶和華神的命令他們不能遵從.
所以他們都要打入去.
所以這就是他們有同理心.
可以互相有共鳴.
能夠自然產生憐憫的原因.

$^{441}$想想其實神也是.
神也是用這句話來形容神.
因此愛承諾恩典.
也是對神一個很恰當的描寫.
當然神不需要我們的憐憫.
不需要我們的恩惠.
但神自己賜恩典給我們的時候.
就是因為他看見我們的軟弱憐憫我們.
甚至不會因為我們的軟弱而輕視我們.
因為我們其實有這個軟弱的肉體.
我們這個講道系列裡說人子身.
本來就是神的旨意.
甚至我們可以想.
神其實容讓我們有這樣的軟弱.
是可以使祂更加憐憫我們.
所以大衛跟我們說.
那個憂傷痛悔的心.
神是不會輕看的.
又好像聖小培姐妹的歌詞.
夜空的星星彷彿在對我微微笑.
輕聲告訴我一切他都看見了.
我所有掙扎所有軟弱和跌倒.
將成為主恩典的記號.
神對這些弱者的恩典.
對某些人來說.
對中國人來說.
其實的確有些難理解.
因為我們從我們的文化角度來說.
在華人傳統我們受儒家思想的影響.
是一種人民主義的思想.
人本主義的思想.
而且也是精英主義的思想.
這種思想認為.
越是優越的人.
是越值得獲得最好的東西.
只是強者才是配得幫助的.
如果資源有限的時候.
你當然是投資在最頂尖的10\%的cream crop那裡.
你怎會將這些資源投資在10\%最不成氣候的羅底產那裡呢.
如果你從人本主義精英主義的角度去看這個故事的時候.

$^{481}$可能你就沒辦法去看準這個故事.
雖然你一定沒辦法去否認神的主權.
神的恩典.
但是如果你拿著這個精英主義人本主義不放的時候.
你很容易仍然是很努力的嘗試在拉闔身上.
找一些美德.
找一些善行.
證明他是配得這個恩典的.
但是你看這個故事當中.
除了拉闔的憐憫之外.
其實拉闔在他身上真的沒有什麼善行好說的.
這個其實就是人民主義和基督教思想最大的分別的地方.
一個最飢餓的分歧.
人本主義或者我們叫humanism.
對人性是一個很樂觀的態度.
無論是中國儒家思想.
或者是西方的人民主義者.
都認為人性是本善的.
如果有人變壞只不過是因為後天的時候.
受到一個不良的影響教壞了他.
因此人民主義者經常有一個思想障礙.
就是他不明白神的公義.
他認為如果用神的恩典公義.
來視為解決人類問題的方法的時候是行不通的.
人本主義者認為人類問題解決的方法在哪裡呢.
答案就在人那裡.
因為他們認為人性是正義的來源.
人性是正義的來源.
拉合的故事如果換成是中國人的讀者.
或者他直覺聯想到的不是恩典.
如果你看拉合的故事的時候.
很容易我們就會拉合一些民間傳奇當中.
那些學識淺薄.
但仍然能夠做一些好事的民間英雄人物.
明代有一個學者.
有一個做官的叫曹學全.
他有一句名句.
仗義每多好九倍,負心當是讀書人.
這個就是反映這種人性本善的觀點.
人為什麼會變壞呢.

$^{521}$讀得書多讀壞腦.
多讀書本來是應該使人更加有知識.
更加知書識禮的.
做人更加有智慧.
秉公行義的.
但是也要看你讀什麼書.
尤其是你讀書的動機.
是一種紙賄讀書功利主義的.
只是為了謀取官職.
書中自有黃金屋.
書中自有顏如玉的時候.
有可能不單止不能夠成為好人.
而且看得書多對你有個不良影響.
一朝赴京考上功名.
拿到一官半職.
連租康之妻都不管.
所以按照這個觀點.
雖然拉闔只不過是一個不學無術的女流之輩.
但是她反而對辨別是非有一個更加好的直覺.
這個就是人性本善總有四個善端的看法.
不過這是不是聖經的看法呢.
我覺得不是.
聖經一開始就反對人性是正義來源的看法.
自從亞當夏娃墮落之後.
人性就不再是正義.
人墮落之後就不知道什麼是何為善.
就算你知道什麼是好事的時候.
也做不出來.
人性既是軟弱也充滿罪惡.
因此人生既是空虛也充滿痛苦.
為什麼神又給我們這個人性呢.
聖經的答案是.
其實引申到聖小梅的詞彙.
人性其實就是恩典的記號.
所以同樣人生也是恩典的記號.
人性的軟弱使每個人都陷入一個真正的兩難當中.
我們以為好的事我們做不到.
我們以為惡的事沒辦法停止去做.
而且這個兩難的困境是死路一條你改不了的.
沒辦法走的.

$^{561}$在絕望的當中絕路的當中我們可以走的是什麼.
唯一可以做的.
也是拉俠做的.
去哭求幫助.
無語問蒼天人能夠再做些什麼呢.
就是向天下求助.
所以正是因為這個緣故.
人性的軟弱反而成為了一個恩典的記號.
人性的軟弱人生的困難.
迫使我們向神尋求憐憫.
成為一個路牌.
明白人處於兩難之間.
我們才去找這個恩典的記號.
第二章的結尾給我們一個很好的意象.
給我們看到一個恩典的記號.
就是那條珠紅繩.
作為這個故事提醒和教訓.
那裡說.
於是女人用繩子把二人從窗戶除下去.
因為她的屋子是在城牆邊上.
她也住在城牆上.
跳幾節.
二人對她說.
看啊,當我們來到這地的時候.
你把這條珠紅線繩子繫在除我們下去的窗戶上.
並要叫你的父母兄弟和你父的全家都聚集在你家中.
女人說.
就照你們的話吧.
於是她送他們走了.
就把珠紅繩子繫在窗戶上.
那條將兩個探子放下去的繩子.
就是成為了這個恩典的記號.
當以色列人再捲土重來來到耶利哥的時候.
見到這條繩子的時候.
就會想起他們對拉俠的應許.
對拉俠來說.
這條繩子也是一個恩典的記號.
提醒他去信靠的恩典.
這條繩子其實是將拉俠的姓名.
兩名探子的姓名.

$^{601}$以及全以色列人的姓名連在一起的一個慈性哀索.
這個故事其實真的幫我們解釋了.
三個人性和恩典的關係.
人是因為敬畏而懇求恩典.
因慈愛去承諾恩典.
是藉信心去領受這個恩典.
那條珠紅繩子某個程度上其實是一個十字架的預表.
十字架一方面向我們信徒應許生命.
但是也提醒那些不信的人.
死亡正在等著他們.
一條繩子能夠帶來生命.
也能夠結束生命.
剛才我們說到曹學泉.
曹學泉其實還有另一句大家都可能聽過的名句.
後來他吊頸自盡之前死的.
他說生前丹管不死後一條繩.
雖然曹學泉他一輩子就希望做一個清官.
做一個好官.
用他的筆桿去批趕這些案件去捍衛正義.
但是他最終都沒有辦法去阻止明朝的腐敗.
明朝的敗亡.
當最後一位明朝的皇帝去世的時候.
他有什麼可以做.
除了上吊自盡之外.
就是沒有另一條路了.
古往今來其實大家看看.
他人性過於樂觀的人.
最終在絕望當中結束他的生命.
其實又何止曹學泉一個人呢.
一條繩子能夠結束一個生命.
一條繩子都能夠帶來一條生命.
在耶穌基督裡面.
我們有盼望的理由.
拉合的故事當中.
後來再提出來的約書亞Joshua.
其實就是耶穌的一個儒表.
約書亞後來到第六章就提到.
約書亞的確是信守了兩個探子所作的承諾.
所以拉合和拉合的全家性命是得到保存.
好像歌詞所說.

$^{641}$當我呼求耶穌聽見我的禱告.
千萬人中他竟關心我的需要.
走過的路有歡笑有淚水.
都留下著恩典的記號.
我總結一下.
其實在我們面前.
每一個人面前.
尤其是如果你仍然是沒有.
真是透過水利.
透過承認耶穌.
是加合教會成為一個基督徒的.
你面前有一個選擇.
我們可以選擇走人本主義的觀點.
堅持人性是正義的來源.
人性有光明的一面.
只要我們盡力活出我們的人性.
我們就能夠拯救改變這個世界.
或者我們可以走另一條路.
恩典之路.
相信人可以因敬畏懇求恩典.
因慈愛承諾恩典.
藉信心去領受恩典.
當我們選擇了相信.
神其實是創造人類的.
神也是創造人性的.
我們不是很難相信.
人性其實是一個指標.
讓我們看到人性的軟弱.
人生的困難.
其實都是一些恩典的路牌.
指著我們去找神的恩典.
所以當我們見到這些患難的時候.
其實只不過是可以看為一個提醒.
我們其實是要走上這個恩典之路.
幾個禮拜前.
我跟大家提過.
我和Irene回了兩天去波士頓.
其實是見我姑媽最後一面.
其實昨天收到消息.
她已經去世了.

$^{681}$我其實唱這首歌.
《恩典的記號》的時候.
其實心裡有很多的記憶.
每次都不知道怎樣可以唱到今上下流眼淚.
所以我今天講到.
都是忍著自己去講.
如果不懂祝福大家.
不要介意.
但是很多時候.
我們去到這些人生的轉折點的時候.
真的讓我們停下來想一想.
究竟我們走了這麼多的路.
最終我們前面是一條怎樣的路呢?.
我很開心見到我姑媽.
我想她見到我們幾個回去的時候.
真的很精神.
我們那時候的焦點.
是不是回光返照呢?.
但是我們後來那些人探望完她之後.
慢慢慢慢就見到她很快的離開.
不過當我想起過去和她這麼多年的相處.
其實她教我鋼琴.
也教了我女兒鋼琴.
我女兒也不是很成才.
我女兒比我好一點.
當人活在神的恩典當中.
將他人生全部投資在神的恩典當中.
回看的時候是沒有任何的後悔的.
而且向前望的時候.
是沒有任何的恐懼的.
他知道他要走上的路是什麼路.
他也可以回看的時候.
每一個人生的轉折點.
都看到恩典的記號.
這個就是我們基督徒生活的分別.
所以如果當中還沒有完全明白這個道理.
多想一想.
停一停.
想一想.
其實神給我們有很多恩典的記號.

$^{721}$在我們身邊.
你和你身邊的人聊一聊.
你會發現其實人生不只是靠自己.
有神在我們當中.
帶領我們走向一個更美好的家.
讓我們低頭禱告.
\newpage



\section{士師記 6:1}
\label{sec:032263}
\textbf{聖經人聖經事:米甸人壓制 以色列人}~~~~黃克勤牧師博士
\newline
\newline
連結: \href{http://www.church.com.hk/acms/content.asp?site=cdc\&op=show\&type=product\&code=032263\&layout=sermon}{\texttt{032263}} ~~~~ 語音日期: 2026-06-07
\newline
\newline
\hyperref[sec:032262]{\small{< < < PREV SERMON < < <}}
~
\hyperref[sec:index_chronic]{\small{[返順時目]}}
~
\hyperref[sec:032264]{\small{> > > NEXT SERMON > > >}}
\newline
\newline
士師記 6:1
\newline
\begin{longtable}{cl}
\hline
\hline
章節 & 經文 (和合本修訂版)\\
\hline
6:1 & \begin{tabularx}{0.7\textwidth}{X} 以色列人又行耶和華眼中看為惡的事，耶和華就把他們交在米甸手裡七年。 \end{tabularx} \\ \\
[1ex]
\hline
\hline
\end{longtable}
$^{1}$各位弟兄姊妹,大家好.
生命樹又和大家讀聖經了.
聖經人,聖經事.
看到聖經背後上帝的工作.
在事,事,事記.
是一個很清楚的循環,說了很多次了.
每一次都是這樣,每一次都是這樣.
今次見到米殿人壓制以色列人.
米殿人壓制以色列人.
今次是以色列人和米殿人對抗.
或者是米殿人來攻擊以色列人.
以色列人頑皮了.
於是神就用外族的人成為一條皮鞭鞭他.
鞭十年八載,很慘的,我們看看.
以色列十二支派當中的迦南人.
這個圖看了很多次了.
不過有時我會移動它,會有些變化的.
主要就是我們見到.
其實當他們進入迦南地的時候.
迦南地主要是約旦河的西面.
東面也給了他們.
有三個支派拿到這裡.
約旦河東,今日的國家是約旦.
約旦河西岸就是以色列.
留下了很多迦南地的人.
當時神是叫他們趕走他們,殺死他們.
牛羊也不要,免得.
最重要是不要和他們一起拜他們的神.
要將他們的祭壇全部拆掉.
但很可惜,以色列人沒有聽他們說.
留下很多人.
其實是分了不同支派的地之後.
剩下一些未得之地就自己去拿了.
很多支派已經定居.
不太有意欲去趕那些人走.
或者有些人怕了他們,就來幫他們打工.
於是就叫他們留下來.
結果正如上帝所說.
你留下來的人將會成為你肋上的荊棘.
即是你的身上繞著荊棘.

$^{41}$你會覺得很痛,很慘,不能動.
而且會成為你們的網絡.
網絡著你動不了.
這是真的.
今次是米殿人壓制了以色列人.
米殿在哪裡呢?.
這個真的很難說.
看看地圖.
地中海,今天的以色列.
下面是阿拉伯.
再到這邊遠一點就是兩河流域.
即是底格里斯河和幼發拉底河.
再到南邊就是波斯灣那邊.
就是底格里斯河和幼發拉底河的末後.
即是出海那邊.
那個地方就是巴比倫.
巴比倫也就是今天的伊拉克.
其實中間就是阿拉伯曠野.
是一個沙漠地區或曠野地區.
其實米殿人在哪裡呢?.
眾說紛紜.
其實米殿很有可能在這個位置.
因為我參考了之後.
發覺原來米殿並不是在這裡.
米殿是比較下一點.
不過我的圖不夠位置.
我畫到最右下角.
米殿應該是在這邊下面.
和一些更低一點的沙漠或曠野地方.
他們在那裡是游牧的民族.
但由於周圍去.
他們習慣了游牧.
於是他們就周圍去.
其實在今天說的.
史詩機殿要和他們打仗的時候.
他們來壓制以色列人的時候.
他們是周圍走.
他們來和以東聯絡.
和阿瑪力聯絡.
今次是米殿和阿瑪力聯絡最多.

$^{81}$來到整個以色列境內周圍去.
又找了其他不同的民族.
或者他們裡面的迦南人民族.
一起對抗以色列人.
由於以色列人是進了去他們那裡住.
所以一叫被叫喚起對抗以色列人.
中間那些剩下的民族個個都響應.
原因就是我本來住在那裡.
現在被他們進來住了我們的地方.
如果說要和他們對抗.
個個都肯.
而且有大的勢力來到.
很好好於是就一起對抗.
這樣就是米殿.
如果今天我們所說的.
第四個士司機殿.
機殿其實是馬來西人來的.
在馬來西即是說.
馬來西在約旦河這裡.
東邊和西邊黃色這一塊.
就是馬來西最大的地域.
馬來西這裡有河東和河西.
而這個機殿就是馬來西支派.
你見到米殿人進攻進來.
連同其他人一齊來打.
或者內部的人一齊裡面裡應外合.
以色列就被他壓制住.
整個以色列地都被他壓制住.
不只是一個支派.
整個以色列地.
所以很慘.
米殿是比較低.
可能應該去到我所指的低一點.
由於他游牧.
帶著他的帳篷牛羊一起到處去.
這樣就很厲害.
我們先看看經文.
士司機第六章.
以色列人又行耶和華.
眼中看為惡的事.

$^{121}$耶和華就把他們交在米殿人手裡七年.
慘了.
米殿人壓制以色列人.
以色列人因為米殿人在山中挖掘挖洞.
建造營寨.
先解釋一下.
惡性循環大家看過很多次.
已經熟悉了.
到這裡第六章你已經可以知道.
整個士司機是怎樣的.
第一就是以色列人.
又行耶和華眼中看為惡的事.
這個惡的事多數是什麼呢.
就是他們有人去拜迦南地人的神.
迦南地的神全部都叫什麼呢.
巴黎.
這麼有趣嗎.
其實每一個他們本地的民族.
都有拜不同的神.
不過他們全部都用一個迦南字去描述.
Bael,Bael,巴黎,巴黎,巴黎.
全部都拜巴黎.
其實就是在說他們在拜不同的神靈.
不過都叫巴黎.
巴黎漸漸都共通了很多.
都是農業的神.
就是生殖的神.
女的又管著別的東西.
這裡有兩夫婦.
阿斯塔綠.
有些是阿瑟拉.
不同的.
又行耶和華眼中看為惡的事.
最惡是這個.
然後又不跟著神的說話去做.
還拜其他的神就很不歡喜.
於是耶和華是自己將他們.
交在外族人的手裡面.
好像一條皮鞭鞭他.
米殿人.

$^{161}$這次輪到.
米殿如果你看士師記裡面.
最長的紀錄就是這件事.
就是機殿和米殿的對抗.
在這裡有很多東西學習機殿這件事.
我們分了很多次說.
都在米殿人手裡.
多久呢?七年.
所以第一點.
循環的第一步.
又行耶和華眼中看為惡的事.
第二就是將它交給外族人.
很長,七年.
米殿人壓制以色列人.
以色列人因為米殿人壓制的緣故.
第一次是那麼慘.
未試過那麼慘.
在山中挖掘挖洞建造營寨.
即是說.
那些人都走上山了.
為什麼走上山呢?.
避了那些米殿人和外族人一齊進來.
那怎麼辦呢?.
因為爬了上山.
沒那麼容易攻他.
於是連平地的城鎮都不敢住.
平原都不敢住.
住上山.
住上山就麻煩了.
麻煩在哪裡呢?.
沒地方種田.
或者田在山下.
平時住就走上山.
那就很麻煩.
又上上落落.
或者落了還未想住.
落了去耕田.
耕幾日.
一個星期才上來.
幾日然後又住下去.

$^{201}$總之就躲在山裡.
山中挖掘挖洞就躲在那裡.
其實都不需要挖掘.
因為以色列很多山都是石.
有很多山洞.
泥土很少.
沙很多.
泥土很少.
而在山上有很多洞.
住在那裡不夠住怎麼辦呢?.
就挖洞.
即是擴大它.
整深它.
就在那裡連結起來建造營寨.
一起住.
以色列人每逢殺種之後.
米殿人,阿瑪力人,東方人.
都上來攻打他們.
所以你見到最主要的人是甚麼呢?.
最主要的人就是米殿人和阿瑪力人.
米殿人就在這邊.
其實已經進來周圍.
阿瑪力人就在這裡.
一起夾擊.
很麻煩.
很難對抗.
和其他東方人.
即是河的這邊的人.
東方人就包括阿滿的老婆.
一齊來搞事.
他們在殺種之後.
即是他們在附近敲著敲著.
附近有很多其他人.
敲到你殺種.
未殺種他不來.
一殺種他來做甚麼呢?.
你已經殺種了.
他就來搞亂你.
挖起地,踩爛他.
讓種子無法生長.

$^{241}$沒有收成.
其實這樣都算好.
原因是他們急著.
如果我是阿瑪力人,米殿人.
我先不來.
最好是等你收割我才來.
收割就順手把麥子拿走.
現在弄爛他.
只不過是你耕種不到.
你耕種不到我也收不到.
搶不到你的糧食.
最好是在你收割的時候我來.
其實他們應該都有.
一齊來攻打.
對著他們安營.
哪裡有人就在那裡.
整個營地預備跟你打.
他們做甚麼呢?.
想害死你.
毀壞土產.
直到家殺.
又要講講了.
真是不好意思.
他們走到上來.
是把遍地的東西毀壞.
整到家殺.
家殺就在這裡.
菲力士地.
其實菲力士地已經很慘.
因為菲力士人住在這裡.
這裡都沒有甚麼以色列人.
不過他們整到那裡.
為甚麼直到家殺呢?.
家殺不整.
因為菲力士人在這裡.
他們不會整爛菲力士人的東西.
所以去到猶大西面這些地方.
直到菲力士地就算了.
所以這些.
什爾茲派的地都被他搞爛了.

$^{281}$但是家殺不整.
所以他就直到家殺.
沒有給以色列人留下食物.
即是說.
整爛了他的耕種.
如果有羊,有牛,有驢.
搶走.
慘了.
羊,牛,驢這些.
在當時是很貴重的東西.
搶了.
好像你的車被人搶了.
哎呀,沒有了那輛車.
羊,牛,驢這些耕種.
運輸是很重要的東西.
是家當重要的財產.
搶了.
沒有了.
沒有留下.
沒有留下食物.
沒有留下牛羊.
所以.
唉,不夠他們走上山避.
很慘.
因為那些人帶了甚麼來呢?.
為甚麼這麼大.
這麼厲害.
毀壞這麼多.
因為他帶著他們自己.
米地人帶著.
阿瑪利人帶著.
他們的生畜和東方人.
帶著生畜和象棚一齊來.
不是進宮一次就走.
進來之後.
哇,帶著工具來.
住在這裡.
牛羊在這裡住.
在這裡吃.
自己的牛羊在這裡.

$^{321}$打爛你的東西.
搶走你的東西.
住在這裡.
多如蝗蟲.
人和駱駝都無數.
為甚麼有駱駝呢?.
這樣就說得對了.
因為.
米地是比較低的.
他們要走沙漠地.
走沙漠路.
就是要駱駝.
所以他們.
是駱駝無數.
人和駱駝.
和象棚帶著生畜.
這樣就很厲害.
有少少像.
中國以前那些.
元朝時候.
蒙古.
他們一路遠征.
除了騎馬之外.
帶著他們的.
駱駝各樣.
牛羊,家眷.
一齊去進宮.
就是這樣.
多如蝗蟲.
都進入國內.
當時以色列是未立國的.
所以這個國.
這個字希伯來文.
就是土地.
地帶裡面.
都進入他們的地帶裡面.
是毀壞全地.
哎呀.
真的很慘.
於是以色列人就因為.

$^{361}$米地人的緣故.
極其窮乏.
很窮.
極其窮乏.
用這四個字.
所以我紅色了一句.
極其窮乏.
在這個窮乏裡面.
他們就會.
惡性循環去到第三步.
就是呼求耶和華.
早些就不呼求.
早些就不理耶和華.
去拜那些巴力.
現在呢.
就呼求耶和華.
以色列人因為.
米地人的緣故.
就呼求耶和華.
嗯.
但是神呢.
本來就很憤怒他們的.
哎呀.
都敬拜我.
我不跟從.
我說了你們又不聽.
又去拜那些.
很憤怒.
那你知道呢.
在初一及二已經說了.
耶和華是.
忌邪的上帝.
什麼叫忌邪呢.
如果你看英文翻譯.
叫做jealous.
啊.
妒忌.
嫉妒.
jealous.
這麼小.

$^{401}$這麼窄.
胸襟的嗎.
上帝這樣.
其實不是胸襟窄.
我經常都用來做例子.
就是什麼叫忌邪呢.
就不是胸襟窄.
因為.
唉.
作為父親.
跟兒子說.
你只可以叫我做爸爸.
因為你爸爸只有幾個.
一個就是我.
那你不要緊的.
你個個都叫爸爸.
有蘇應做爸爸.
不可以的.
如果你叫別人做爸爸.
我生氣的.
不歡喜的.
你叫別人做媽媽可以嗎.
生氣的.
你試試叫別人的女兒做太太.
叫別人的男兒做老公.
你看看你的配偶怎麼樣.
你不要胸襟這麼窄.
jealous的.
不然呢.
這個jealous的意思就是說.
只可以有一個.
不可以叫別人的.
有時候.
神很生氣.
他們就呼求耶和華.
不過神都好.
在用外族人鞭打他們的時候.
都回心轉意.
於是就第四了.
在這裡就未曾猜顯示私.

$^{441}$不過就先猜顯了.
先知去跟我們講清楚.
喂.
哎呀神啊.
為什麼你這樣對我們啊.
為什麼你不記得你那些應許.
說要給這塊地我們嗎.
為什麼又會有其他外族人.
被他侵入來.
是你不對了.
那神就派了先知去跟他們講.
其實是你們的問題.
耶和華就猜顯先知.
到以色列人那裡.
對他們說.
耶和華以色列的神.
你們的神.
跟你們這樣說.
我就是曾經令你們從埃及上來.
出了遺牢之家的那個.
你記得吧.
就是我帶你們出來的.
然後呢.
是救你們脫離埃及人的手那個.
並且脫離一切欺壓你們的人的手.
將他們從你面前趕出.
將他們的地賜給你.
我就是那個.
帶你出埃及遺牢之地.
然後將這個地賜給你.
將當地的人在你面前趕出.
你看看還記不記得.
你試想想.
你圍著城這樣叫.
耶利哥就倒了.
然後很多的城都被你們全部攻陷了.
你知道.
其實就是我在你面前趕他出去的.
比冰雹,比黃蜂殺死的人還多.
嘩這樣.

$^{481}$是我將他們趕走.
是我的手將地賜給你們的.
其實我又跟你說過.
不過你記得留意了.
我是耶和華你們的神.
我還這樣說.
我又對你說.
我聲明了很多次了.
你們住在阿摩利人的地方.
記得.
不可敬畏他們的神.
你們竟不聽從我的話.
所以我將他迎光了.
我已經說了.
我是帶你們出來那個.
幫你們入迦南打贏了那個.
我叫了你們在當地人那裡住的時候.
不要拜他們的神.
特別說了很多次.
你們竟不聽從我的話.
所以有今天.
是你們不聽從我的話.
這樣講得很清楚的.
我想大家都知道了.
他們的惡性循環就是這樣.
又行耶和華眼睛開是惡的事.
又敬畏別人的神.
所以神這個忌邪的上帝就發火了.
就用外邦人去鞭打擊打他們.
於是他們就開始.
很慘啊.
被人搶了東西.
慘的時候又呼求耶和華.
呼求耶和華.
這次呢.
神跟他說清楚了.
這次米殿人.
不是一個小地方.
是全境裡面來搞呢.
其實是我叫他們來的.

$^{521}$不過是因為你們不聽從我的話.
所以.
見不見到這條藤條.
這條米殿藤條就來了.
弟兄姊妹今天跟大家分享到這裡.
你會發覺到.
最主要的就是.
為什麼會遭遇藤條呢.
其實就是因為.
不聽耶和華的話.
很多時候我們會遭遇很多困難.
遭遇很多攻擊.
遭遇很多的敵對.
很多時候我們是有一些地方.
是沒有聽上帝的說話.
於是神就鞭打我們.
小心.
這個鞭打其實對我們不是壞的.
是好的.
令到我們可以警醒.
哎呀 原來這樣是不行的.
有時候有些痛也好的話.
我小時候都講過很多次了.
怎樣知道牙痛慘過大病呢.
要牙痛過.
明明叫了你很多次.
不要吃糖粉.
要刷牙的.
刷牙那麼麻煩.
你一個小孩都不喜歡刷牙.
然後呢.
我就喜歡含著一粒糖粉.
一直吃著吃著睡著.
睡得特別香甜.
結果就蛀牙了.
真的.
牙有洞.
就去剝牙.
道牙根.
剝牙.

$^{561}$痛得不得了.
開始學到了.
牙痛的時候真的慘過大病.
好痛啊.
就去買些牙痛水.
當時沒有那麼多牙醫的.
很多都是以前那些.
學師回來的牙醫.
手腳又不太行.
哎呀 好慘.
就是這樣.
當時也沒有那麼多止痛的.
所以真的痛足七年.
這個米電的皮鞭.
是再次提醒我們.
要聽耶和華的話.
請你記得.
最重要記得.
要聽耶和華的話.
我們一起祈禱吧.
主我們向你認罪.
向你道歉.
很多時候我們是沒有聽你的話.
有時候我們覺得自己都幾好.
都幾強.
就會聽自己的話.
就會覺得我能夠掌控一切.
結果就去到最差.
有些知識的人這次.
是遭遇這麼大的困難.
是史詩裡面最長的篇幅記載.
他們竟然要走到山上住.
都不敢住平地了.
要住在山上.
在一個洞裡住.
山下眼光光.
看著他們的田被破壞.
看著好像蝗蟲那麼多的外族人.
帶著家眷帳幕.
他們的牛羊牲口樂陀無數.

$^{601}$霸佔了他們的地方.
令他們心裡覺得很悲傷.
他們回想.
他們質問神.
為什麼你不開顧我們.
多謝我們的主.
你是差派先知.
去跟我們說.
是你們不肯聽耶和華的話.
以至這條米電的皮鞭就來了.
幫助我們更加醒目.
不等到皮鞭來到.
我們就聽你的話.
吸引我們更加歡喜聽你的話.
隨聽我們的禱告.
奉耶穌基督的名求.
阿們.
\newpage



\section{創世記 4:10-16}
\label{sec:032264}
\textbf{活出召命——生命的覺醒與承擔 II (華)}~~~~梁明財牧師
\newline
\newline
連結: \href{http://www.church.com.hk/acms/content.asp?site=cdc\&op=show\&type=product\&code=032264\&layout=sermon}{\texttt{032264}} ~~~~ 語音日期: 2026-06-07
\newline
\newline
\hyperref[sec:032263]{\small{< < < PREV SERMON < < <}}
~
\hyperref[sec:index_chronic]{\small{[返順時目]}}
~
\hyperref[sec:032265]{\small{> > > NEXT SERMON > > >}}
\newline
\newline
創世記 4:10-16
\newline
\begin{longtable}{cl}
\hline
\hline
章節 & 經文 (和合本修訂版)\\
\hline
4:10 & \begin{tabularx}{0.7\textwidth}{X} 耶和華說：「你做了甚麼事呢？你弟弟血的聲音從地裡向我哀號。 \end{tabularx} \\ \\ \relax
4:11 & \begin{tabularx}{0.7\textwidth}{X} 現在你必從這地受詛咒，這地開了口，從你手裡接受你弟弟的血。 \end{tabularx} \\ \\ \relax
4:12 & \begin{tabularx}{0.7\textwidth}{X} 你耕種土地，它不再給你效力；你必流離飄蕩在地上。」 \end{tabularx} \\ \\ \relax
4:13 & \begin{tabularx}{0.7\textwidth}{X} 該隱對耶和華說：「我的懲罰太重，過於我所能承當的。 \end{tabularx} \\ \\ \relax
4:14 & \begin{tabularx}{0.7\textwidth}{X} 看哪，今日你趕我離開這塊土地，不能見你的面；我必流離飄蕩在地上，凡遇見我的必殺我。」 \end{tabularx} \\ \\ \relax
4:15 & \begin{tabularx}{0.7\textwidth}{X} 耶和華對他說：「既然如此 ，凡殺該隱的，必遭報七倍。」耶和華就給該隱立一個記號，免得人遇見他就殺他。 \end{tabularx} \\ \\ \relax
4:16 & \begin{tabularx}{0.7\textwidth}{X} 於是該隱離開了耶和華的面，去住在伊甸東邊挪得之地。 \end{tabularx} \\ \\
[1ex]
\hline
\hline
\end{longtable}
$^{1}$主持人:弟兄姊妹平安!今天我們還是繼續在兩個禮拜之前,跟各位來看創世記的該隱這個聖經的人物,希望大家還有點印象,到底該隱是一個怎麼樣的人..
劉仲敬:我再幫助各位來複習一下.從該隱身上,最主要是他不是一個很出名的人.當然,對於一些人來說,你也會認識他,因為他是創世記裡面一開始就出來的人,除了他的父母亞當夏娃,該隱和亞伯..
但是,我是盼望能夠從這個小人物身上,他的一些生命的反省,我們做一個反面的教材.所以,這個是從他一個做得不好的身上,我們來作為我們很重要的提醒..
跟各位來看,從該隱的生命當中,我們上期有提到的是從創世記第四章的三竇九節,講到該隱怎麼樣跟他的兄弟亞伯和一號阿神對話.這就是我們上期有提到的..
最後,該隱的哥哥把他的弟弟殺掉,然後神又再跟他有一些對話.這個大概就是我們上期講的..
所以,我們上期講的重點是講到該隱的沒有覺醒,雖然他有跟一號阿神對話的時候,但他仍然不知道,也不承認,也不願意去反省他所作所為..
這個我們上期有提到,就是因為其中一個影響可能是他覺得他是一個大哥,是家裡面的長子.所以,很容易他的面子問題,我們上期提到他的臉色變了,就是好像我們那個沒面子的樣子..
他引起他心裡面的怒火,他的生氣,讓他的怒火沒有控制的話,他就成為他做更大的惡的一個原因..
也不是說我們不能夠生氣,但我們要懂得怎麼樣去好好的控制我們生氣的時候那個怒氣或是那個脾氣..
所以,這個怒火中燒,以致他覺得神很不公平,他就覺得神偏心.我們說他在我們廣東話裡面是savage,他沒有負上他應該負的責任,他是一個不負責任的人..
然後他做了一些是選人不利己的,甚至說有時候可能是兩敗俱傷,或者說他根本就沒有獲到任何的好處.他只是覺得,既然他弟弟有,他就不賞他有這個神的月納他的貢物..
雖然他自己沒有,他也不賞別人有,他就把他殺掉.這些選人也不利己.這大概是我們上次講的.但我跟各位有提到,在新約聖經裡面有提到該隱這個人,在約翰一書和尤大書裡面..
所以,讓我們看到一個很重要的,該隱回應耶和華的時候說,我豈是看守我弟兄的那個人..
我想跟各位說,在我們跟人相處裡面,不單只是家庭裡面的關係,在教會里面也是.也不是說,我的要看守什麼人..
其實,在約翰一書告訴我們說,我們應該因為神是愛,我們應該在愛裡面彼此相愛.愛比看守更大..
其實,神不是責備該隱單單是因為他的貢物不好,或者說沒有看守弟兄,其實是因為他根本就沒有愛.所以,在約翰一書用到該隱作為一個例子的話,是提醒我們神是愛.我們應當在神的愛裡面,跟我們的家人,跟我們的弟兄姊妹彼此相愛..
這就是用該隱作為一個例子,一個反面的教材,因為這是該隱沒有愛..
然後就是說,要提醒我們不要走該隱的路,這是他的選人不利己.他既沒有的,也不想別人有,然後變成一個兩敗俱傷,甚至說犯了更大的罪惡..
所以,我們不要跟該隱走同一條路,這是我上次跟各位有從聖經當中我們去領受的..
不要跟該隱走同一條路,就是說,我們不要常常,有時候我們中國人是比較容易覺得面子很重要,所以如果有些事情讓我沒有面子的話,我就會很不高興,或許做一些很不好的事情..
面子對你我不曉得有多大的影響,或者說多重要,但如果你真的是一個很看重面子的人,要小心,有時候你會變成走該隱的路..
這些都是該隱不好的地方,我們要引以為戒.而且該隱是一個缺少自動自覺去省略自己生命的心和力.他沒有省略反省的心,而且他沒有那種能力..
那種能力是怎麼樣的能力呢?我想上次也跟各位看,聖經很清楚告訴我們,每一個人的心都是很複雜的,我只能夠說,他這裡形容是鬼子..
但就是說我們的心,對每一個人來說,可能你自己也不太明白你自己的心到底在想什麼,在籌劃什麼,在期望什麼.我們的心需要好好的去更多地省略我們的心,是一個怎麼樣的心..
我們的能力,能夠去省略我們生命的能力,其實不是靠我自己.我能夠去面對我每一天的生活,也不是靠我們自己的能力,是要依靠神..
約翰一書一個很好的提醒,就是說,縱然你覺得你不太能夠去愛的話,如果是這樣子,其實也不是你能夠去做,是神在我們裡面.我們纔有這種能力去愛我身邊的人,無論是愛自己也好,愛家人也好,愛弟兄姊妹也好,愛一些你不喜歡的人也好..
不是靠我個人有多少的愛,多少的愛心,是靠神在我們裡面幫助我們.所以這個是我上次跟各位聊從蓋茵生命當中的第一個所領受的信息..
今天我想跟各位來多一點放在,我們說到蓋茵是一個不負責任的,我們就需要多反省我們是不是能夠有承擔責任的勇氣,我們的生命的素質有沒有這樣子的承擔責任的勇氣在我們的生命裡面..
蓋茵是沒多少的了.蓋茵犯罪的結果是怎麼樣的呢?它破壞了它跟神,跟人,跟地的關係.這個是在蓋茵《創世記》第四章裡面,我們可以多看一點..
所以我今天接續上次我們看的經文,我們來看《創世記》第四章十到十六節.這個就是講到蓋茵犯罪的一個結果,她要承擔她這個犯罪的結果,其實她不願意,她覺得神對她還是不公平,為什麼要受那麼嚴重的懲罰..
那我們先看看那個經文,《創世記》第四章十到十六節.看到蓋茵,我請弟兄姊妹我們一起來讀這段的經文,預備!.
你被流離飄蕩在地上.蓋茵對耶和華說,我的刑罰太重,過於我所能當的.你如今趕赴我離開這地,以至不見你面.我被流離飄蕩在地上,凡遇見我的必殺我..
耶和華對她說,凡殺蓋茵的必遭報其輩.耶和華就給蓋茵立一個記號,免得人遇見她就殺她.於是蓋茵離開耶和華的面,去住在伊甸東邊羅的之地..
這裡是神跟蓋茵給她的一個懲罰,所以蓋茵覺得這個懲罰太重了,她好像承擔不起..
但其實這就是神要蓋茵承擔她所做的事的責任的那種勇氣.無論她是不願意也好,願意也好,她就是要做一個承擔責任的人.所以我們也要有多這一份承擔責任的勇氣在我們每一個人的生命當中..
這裡其實有一個比較接近剛纔我們所講到蓋茵為了她自己的面子的問題.如果你留意,這裡講到"如今敢逐我離開這地,以至不見你面"..
其實它在經文裡面用的詞,它是一個比較跟面也有關係的.它在當中講到"敢逐我離開這地",其實它有一個面在裡面..
如果用英文來說就是"the face of the land",它是離開那個地的面,然後它是讓它不能夠見到誰的面?神的面,the face of God.所以這個其實,寫聖經的人真的是很在這個面上面,那個face的面子,如果有用的話..
蓋茵終於自在神給她的那個結果裡面,她就知道她已經因為她自面子的問題,她就不能再留在這個地面,就要離開這個地,然後她不能再在這個地面,跟這個地的關係就已經破壞了..
然後她現在有點醒覺,她醒覺是什麼?她要離開神的面,是另外一種的沒有面子.如果她之前她的臉色變了,她覺得她的面子掉下來,她沒有面子..

$^{41}$現在她覺得一個更重要的沒有面子是什麼?她離開了神的面,神的面不再是跟她有關係.她另外一種的沒面子的話,就是神的面離開了她.這對她來說,她醒覺了,她察覺了這個是更重要..
沒有神的面,或者說神的面她不能再見到神的面,代表什麼呢?就代表她再沒有神的保護,神沒有再跟她同在.她就沒有神的保護了,離開地的面就是她再沒有一個家了,她已經要離開這個地面的話,就代表她沒有一個家,她要流離,要飄蕩到其他的地方去..
所以當人犯罪以後才醒察,或者說醒覺她那個嚴重的後果的時候,有點太遲了.她已經不能再沒有那個face of the land,沒有那個face of God,沒有神的面也沒有地的面,她沒有家了,也沒有神的保護..
所以這種沒有面子對該因來說是更嚴重的,她才醒覺,她不能再有神跟她同在,當然神也是很憐憫她的..
這後面跟她說是什麼呢?這有點如果簡單來說就是說,給她一點面子.其實有點是給她一個mark,一個記號在她的臉上,在她的面上給她一個記號..
當她覺得她自己會很無助,會被人追殺的時候,她還是有神給她一個記號,給回她一點點面子,讓她不是被殺.這還是神的憐憫,神讓她仍然可以繼續在那個地上來生活..
所以,如果從該因這個反面教材來說的話,我真的深賞弟兄姊妹記得,中國人真的很受面子的影響,首先我們真的不要,最重要不是我自己個人的面子..
其實如果真的當你去到一個人生的階段是神的面子,神的面你不能再見到,你離開了神的面,神的面已經沒有了,那纔是更嚴重的一個沒面子..
如果我們看第五到第七節的話,有一個講到是說,如果你行得好的話,神就會悅納他;只是他行得不好,他就不蒙悅納..
所以,行得好那個意思跟悅納是很有關聯的.悅納的意思是什麼呢?就是講到揹負,承擔.如果他是一個願意有承擔的話,他行得好,願意承擔的話,有承擔,有擔當的話,他就是神所悅納的..
所以在這段經文裡面,除了有關面子的話,另外一個就是承擔該人到底有沒有願意承擔他的責任,他是不是有擔當的認同..
通常我們形容一些男生,他跟女生交往的時候,就是說要多認清楚這個男生是不是有擔當的男人.如果他不會有擔當的話,那你要好好的考慮,你要跟這個男生繼續下去的話..
所以,該人的犯罪,他有時候覺得我們可以問到底是誰的錯呢?上次我們有講到,是不是亞伯的錯?是不是亞當·薩瓦為什麼多生一個兒子?只生一個獨子就不是比較好嗎?那就沒有那個比較,沒有那個增進..
是不是撒旦的錯?好像撒旦·薩瓦把那個責任都推給撒旦,他喫了那個不可喫的果子的話,就是不是我的錯,是撒旦的錯..
然後就是,還是神的錯了.神,你為什麼造了這個伊甸園,為什麼要給那些不可喫的果子在裡面?所以,人很容易走這條路,就是說不太願意承認或者說擔當自己所犯的錯..
所以,其實這個最重要的是,我們都認識是該人的錯,不是其他人,是該人他自己的錯,但只是該人他沒有那種承擔的勇氣..
神不悅納代表的是,他沒有那個行得不好,他沒有那種承擔,沒有那種勇氣在當中.所以,我們盼望從該人的身上我們要學會忍忍責任,我們應該做一個有擔當的人,要學會怎樣承擔..
我們要反省,有時候在我們的工作裡面比較容易,就是有時候在你工作,跟你的夥伴或者說你的同事合作的話,有時候你會發現你是給人家把責任推給你的,就是要卸給你,你就覺得被窩了,那你就很慘..
但有時候你會不會想到,其實你自己也有一些責任,會不會是你把責任推給別人的.這就是我們要學會多一點反省..
其實在我們的人生當中,神會不斷給我們有機會,讓我們去願意做一個有承擔的人,只是該人他沒有抓到這個機會..
有時候我們也要反省,我們看別人是怎樣.當我們跟別人一起來相處或者說一起來工作,一起來在教會里面有教會生活的時候,其實我們有時候要看我們身邊的人是一個怎樣的人,你是怎樣去看人,這個也會影響我們怎樣去做一個有承擔,有擔當的人..
該隱把亞伯看作不是他的兄弟,雖然他說他的兄弟了,但是該隱沒有看亞伯是他的兄弟.你知道該隱怎樣看亞伯?他大概是看亞伯是一個搶奪者,他覺得亞伯是搶了他長子的祝福,搶了神給亞伯的獻祭以後的祝福..
他認為亞伯是他的敵人,是他的對手,是搶奪他的東西的人..
所以有時候我們也要反省,我身邊的人到底我怎樣看他?會不會我們不會把他看作是在基督裡面的弟兄姊妹,沒有把他看作一個一家人,沒有用愛來去愛他們?.
所以有時候我們的責任不單止要珍惜神給我們的機會,也要反省我們有沒有看別人看錯了..
所以在當中我們要學會有承擔的勇氣,我們不要像該隱,我們要做一個有擔當的人,要願意承擔我們有做過的事情,那種承擔責任的勇氣..
主持人:最近大概是很多的大學都在舉行畢業禮.我不知道您還記得在大學的時候,畢業禮的時候通常那些領導都會出來講一些話嗎?我不知道您對他們的講話還記得有多少..
最近我看到一些大學,也有不少的報道去講他們那些領導在大學畢業典禮當中給學生,給畢業生的一些訓練,怎樣在他們準備要離開,進到社會上工作的時候..
劉仲敬:我想跟各位分享其中一個,我覺得蠻有意思的.有一個大學不在香港的,在臺灣,我不知道您有沒有聽過,有一個大學叫中原大學,在臺灣的中原大學..
它比較特別的是什麼呢?它是一個基督教的大學,在臺灣其中有幾家基督教大學裡面.這個大學其實算是一個私立的大學,沒有政府的很多資助,它是一個私立的大學..
最近它被選上了,是跟美國的企業來合作的.那個企業是誰呢?就是NVIDIA.我猜你們都應該大概有聽過了,那個回答是NVIDIA,那個黃仁勳..
中原大學就被挑選為其中在臺灣一個私立大學第一家,跟NVIDIA有合作的一個大學,所以這個很難得.所以這個大學其實在臺灣還算是有點影響力..
最近它畢業典禮的時候,它有一位畢業生在臺上去分享.這個畢業生大概是每個大學都可能有的,大概算是全校拿第一名的,這是所有畢業生.他們有一個全人標綱的大招,是特別優秀的畢業生,上臺跟同學一起來分享..
他有一句話很有意思,我跟各位來講一下.其實他是一個馬來西亞的僑身,到臺灣去讀書,很難得他真的是在讀完四年畢業以後就拿到最優秀的獎..
他在分享當中說,他很欣賞在中原大學有一個叫全人教育的訓練.他在中原大學裡面覺得不僅能夠得到一個專業的知識,更透過全人教育培養他們學生的責任感跟關懷的精神,讓他們在成長的過程當中不斷地突破自我,破綻他的戲也..
我覺得這是真正在該隱身上沒找到的,也可能在我們身上比較缺少的..
同學他能夠在四年的大學從馬來西亞到臺灣的時候,他覺得在一個基督教大學裡面能夠得到這個訓練,得到這個成長,重要的不是學到多少的知識,拿到多少的GPA和高分的分數,拿到什麼大獎..
他是覺得他能夠在這個培育當中,他的責任感跟關懷精神越來越好,越來越高,越來越多..
我相信這真正是一個很正面的分享.該隱真正就是沒有那種承擔責任的勇氣.我們盼望在我們還年輕的時候,我們在我們的生活當中能夠多有這份責任感,多有一個承擔的勇氣,讓我們的生命的樹枝不會像該隱一樣..
在《聖經》裡面的《舊約新約》當中,我想用兩個人物給大家作為另外一個的看法,有點像該隱的,但也不太像..

$^{81}$一個叫月拉.我不曉得你們有沒有聽過《舊約》《聖經》裡面有一個人物叫月拉.如果這個月是差傳月的話,月拉大概有時候在差傳的經文裡面都常用的一個人物..
月拉的覺醒和承擔,大家可以回去看《聖經》裡面的月拉書,它其實就是一個在當中所經歷到的,我們覺得是一個沉浸在船主坐船的時候所形容的沉睡的人..
他沒有覺醒,他仍然像人生當中在做夢的樣子.後來當然他是愚夢初醒,他最後還是跟從神的帶領,要去做神要吩咐他做的事情,但他還是一個沉睡的人..
所以有時候我們真的是會不會也提醒自己,我們不要做一個沉睡的人,讓我們真的是多有一點自動自覺的醒察的話,我們就能夠更做一個承擔責任的人..
月拉的沉睡,他後來愚夢初醒,他知道他自己做了什麼事情,他覺得在船上風浪或者說船快要沉了,很危險的時候,他知道是他的問題,他終於醒覺了,所以他還算有負責任..
其實這個負責任應該在於,我們有時候常用"早晚要早熟",你不能夠說你不負責任或者說做了事情就不管,早晚你要早熟了,還是要承擔這個責任..
所以他自己要早熟了,結果他叫人把他丟到海里面,風浪就會平靜..
但他好像有點覺醒,最後他在經歷到他要去城市裡面勳章的時候,他還是對神發怒,他的怒火還是有..
他對神有點是罵神,就是說"你幹嘛怎麼樣,我要死掉就好了",所以他還是不太覺醒,因為他怎麼樣看人,他怎麼樣看他要去城市裡面的人,他把他們看作是他的仇人,是他的敵人,所以他就不想去..
雖然他好像有點像認識神,是一個有憐憫的神,有慈愛的神,他覺得我去也沒用,神也不會把他們殺掉,所以他就好像是認識神,其實他也是一知半解,也不太完全認識神是一個怎麼樣的神..
他就是這樣子,有他的怒火,也為他的面子說"我去了以後,你宣告了神要滅城的信息,然後神又沒有把他們滅掉,他又好像沒有面子的人.".
所以就是這樣子,他有時候好像覺醒了一點點,但是後來還是回到夢中,還是有可能他還要回到一條大魚肚裡面,再一次地進到黑暗當中,他纔會有醒覺的機會..
所以這個是願,他有時候真的需要有更多的覺醒,他有時候會負責任,有時候也不太負責任..
另外一個我想跟各位來看是《新約聖經》的另外一個人物,他也是經歷到怎麼樣的在覺醒的過程當中,然後他願意承擔神給他的責任..
這個是保羅,他以前有一個名字叫掃羅,後來他改名叫做保羅.我用他其中一段的經文是《使徒行傳》的二十六章九到二十二節,如果大家有時間可以回家再慢慢去看這一段經文,可能你對他有點認識,如果不太認識的話,你可以回去再看《使徒行傳》的一些經文,就多認識保羅這個人..
他就在做一個見證的時候,他講到他自己所經歷到的.他以為信耶穌的人都是壞人,都是不好的人,所以他曾經追殺他們,抓他們,把他們帶到監牢裡面去..
他們那時候真的是很瘋狂,但最後他經歷到什麼呢?當他追殺他們的時候,他經歷到一個黑暗的時候,他瞎眼了,看不見,進到一個黑暗當中,然後他又開始願意承認自己做錯了,他開始覺醒了,然後他願意悔改,願意領受神給他的使命..
所以,後來他就經歷到上主給他的幫助,給他的接納,他願意為耶穌給他的使命,他願意承擔起來..
所以,他跟這個王有一句話我想跟各位來看,這句話是什麼呢?如果看到就是黃色的那個字.他跟那個王說:"我固此沒有違背那從天上來的意義上.".
他從來沒有違背那個從天上來的意義上.當他醒覺的時候,他知道他自己做錯了,他願意悔改,走回正路,他願意依靠神,他持守神給他那個意義上,他願意承擔他已經做錯了,他承認錯了,他也願意承擔神交託給他的使命..
所以,他最後做一個見證,就是:"我從來沒有違背那從天上來給我的意義上.".
我相信這就是從保羅的見證身上,他的覺醒跟他的承擔..
其實,甘影也好,月拿也好,保羅也好,他們大概有不同程度的覺醒.我不知道對於你來說,你曾經經歷過哪一個階段,跟哪一位比較像?當然,我希望你是最像是可以跟保羅像..
在我們今天,我們覺得是一個比較混亂,很多的戰爭,很多的變化的時代當中,我們有哪一些的覺醒呢?在你作為一個基督徒,你是神所愛的兒女,你生命當中有什麼覺醒?.
有時候我們會,特別在今天這個六月的時候,剛纔大家有唱不同的詩歌,在我們跟各位來分享拆船使命的時候,大家有沒有一些覺醒,你要去幫助其他人在你的生命當中?.
拆船樂大概就是提醒我們要在近處也好,在遠處也好,我們都能夠去把福音帶給不同的人.所以,有一首詩歌叫《人們需要主》.這裡它的副歌有一個歌詞,我覺得也很有意思,就是說我們何時醒覺?.
在最後一句話,人們需要主.有時候我們還是在像月娜做夢也好,在該引祂自己為了自己的命旨,也沒有那種醒覺的時候,我們什麼時候會醒覺到其實別人很需要主?.
我們不要單單只是把眼光,目光放在自己的身上,我的需要上.其實很多人的需要,我們有沒有醒覺到他們需要我們的主耶穌,需要我們已經有的福音,需要神的愛?這就是我們需要醒覺的地方..
當然,拆船樂是提醒我們為別人的需要醒覺的,但我還是覺得,從我個人,我自己開始,我的生命醒覺是怎麼樣的?我們要覺醒,不要走該演的路..
我們有提醒我自己,不要為命旨,不要太多的怒火控制不了,不要做一個不負責任的人,不要選人不利己.在我們所做的事情上,我願意做一個有承擔的人..
這個有承擔的人是讓我們去依靠神的恩典,不是靠我自己.我的能力很有限,我們要做一個真誠,有真愛,願意敬畏上主,靠主的恩典幫助的人,我們才能夠更多的覺醒,也有更多承擔能力,做一個有擔當的人..
所以我盼望能夠在該演身上這些幫助我們,就我們每天有機會的話,更多既動,既覺的覺醒,我們也能更多的在我們的生命的樹極裡面有這種擔當,承擔的那種勇氣,不是靠我們自己能力有多少,是靠著神的恩典..
最後我想跟各位分享,視頻139篇,這段經文其中一個提醒是,神很清楚知道我們每一個人的好的地方,不好的地方,知道我們的強項,知道我們的弱項,知道我們的優點,知道我們的缺點,神很清楚明白我們,認識我們..
所以,我們儘管很放心地來到神面前,繼續地跟隨祂,信靠祂.有時候你真的是像一些該影也好,越拿也好,你很想逃避神的話,我告訴你,你逃到天沼底,你都不能夠逃避神..
就算你有翅膀會飛的話,你也飛不到哪裡去.所以,我們不能逃避神給我們的責任,不能逃避我們自己所種下來的不好的一些後果,我們要承擔責任..
我們不能夠逃避,我們不能夠逃避神,也不能夠逃避我們應該承擔的責任.但神其實是很想扶持我們,祂的右手必扶持我們,無論在哪裡,祂最想是帶領我們走在正路上..
所以,詩篇139篇的最後兩節經文我覺得是很好的,常常成為我個人的提醒,我也盼望能夠成為弟兄姊妹們的提醒,讓求神不斷地檢查我們.當然神知道我們,神可以檢查我們,但求神幫助我們更多地自我地省覺,自我地省察我到底是走在什麼樣的路上..
其實有時候,試煉的領導,神要試煉我們的話,就是要提醒我們,幫助我們走在正路上,不要走在該引的路上..
所以我要請弟兄姊妹們,我們一起來讀這一篇139篇的最後兩節經文,讓它成為我們在生命當中,生命醒覺和承擔的一個很好的提醒,讓我們不要走該引的路,走在神帶領我們走的正路上..
我們一看到的話,請大家一起來讀這段經文,預備來!.

$^{121}$神啊!求禰檢查我,知道我的心思,試煉我,知道我的意念,看在我裡面有什麼惡行沒有,引導我走永生的道路..
我們一起來禱告!.
求禰保守我們,不要走該引的道路,讓我們繼續地依靠禰,依靠禰的恩典,願禰的手扶持我們每一個,保護我們,引導我們,以及我們能夠走在一條神力所喜悅的正路上..
主,謝謝禰決定我們的禱告,靠耶穌救救民生!Amen!.
\newpage



\section{使徒行傳 18:1-3}
\label{sec:032265}
\textbf{宣教的神國俠侶}~~~~劉銳鋒牧師
\newline
\newline
連結: \href{http://www.church.com.hk/acms/content.asp?site=cdc\&op=show\&type=product\&code=032265\&layout=sermon}{\texttt{032265}} ~~~~ 語音日期: 2026-06-07
\newline
\newline
\hyperref[sec:032264]{\small{< < < PREV SERMON < < <}}
~
\hyperref[sec:index_chronic]{\small{[返順時目]}}
~
\hyperref[sec:032266]{\small{> > > NEXT SERMON > > >}}
\newline
\newline
使徒行傳 18:1-3
\newline
\begin{longtable}{cl}
\hline
\hline
章節 & 經文 (和合本修訂版)\\
\hline
18:1 & \begin{tabularx}{0.7\textwidth}{X} 這些事以後，保羅離開雅典，來到哥林多。 \end{tabularx} \\ \\ \relax
18:2 & \begin{tabularx}{0.7\textwidth}{X} 他遇見一個生在本都的猶太人，名叫亞居拉。不久前，他帶著妻子百基拉從意大利來，因為克勞第命令所有的猶太人都離開羅馬。保羅去投靠他們。 \end{tabularx} \\ \\ \relax
18:3 & \begin{tabularx}{0.7\textwidth}{X} 他們本是製造帳棚為業。保羅因與他們同業，就和他們同住，一同做工。 \end{tabularx} \\ \\
[1ex]
\hline
\hline
\end{longtable}
$^{1}$自上次會議後發出的資料文件.
當你看到這個題目的時候.
不知道你腦海當中浮現的第一個畫面是甚麼.
不是神鵰俠侶 是神國俠侶.
你會不會想起一些武俠小說裡面的畫面呢.
每當我們讀武俠小說的時候.
最令我們心潮澎湃的.
往往不是那些孤高自傲.
獨來獨往的孤膽英雄.
而是那些能夠並肩作戰.
快意恩仇的神仙倦女.
我們會不會想起郭靖和黃蓉.
他們守衛襄陽的合骨柔情呢.
會不會也想到楊過和小龍女.
他們相見合璧的絕代風華呢.
他們心意相通.
在刀光劍影的江湖裡面.
成為彼此堅實的後盾.
我們弟兄姊妹有沒有曾經想過.
在兩千年前的風起雲湧.
危機四伏的福音江湖裡面.
也曾經出現過一對名震天下的神國俠侶.
在新約裡面這對夫妻的名字.
一共出現過六次.
每一次出現他們都是一起的.
是緊緊相連從來不分開的.
他們沒有絕世的武功.
但是他們手握聖靈的寶劍.
他們沒有門派的利祿.
卻是為了神國的大業浪跡天涯.
他們不是全職的使徒.
而是和在座各位一樣.
是一個有自己生計.
有自己工作的普通信徒.
然而他們卻是用自己的婚姻.
他們的家庭和事業.
建立起一個堅不可摧的宣教的基地.
他們是誰呢?.
他們就是亞居拉和伯基拉.
我稱呼他們為神國俠侶.

$^{41}$今天神要在這個猜拳的主日.
對我們每一個人,每一個家庭來說話.
宣教從來不只是一個人的事情.
而是上帝給你整個家庭榮耀的呼召.
神的國度不需要你做一個.
在安書區裡面的宅男宅女.
而是神正在尋找這個世代.
能夠為主去到.
勇於宣教的神國的俠侶.
我們今天就看看這對神國俠侶.
給我們有什麼可以學習的地方.
首先第一方面.
他們是一對相見合璧.
是同心同行的宣教伴侶.
在《使徒行傳》第十八章.
一到三節這樣說.
我們一起來讀一讀.
「這是以後保羅來了雅典.
來到哥林多.
遇見一個猶太人名叫亞居拉.
他生在本都.
因為格魯雕命猶太人都離開羅馬.
新近帶著妻伯基拉從意大利來.
保羅就投奔了他們.
他們本是製造帳篷為業.
保羅與他們同業.
就和他們同處作工」.
《使徒行傳》第十七章就提到.
保羅在雅典向外邦人宣教.
如果你還有印象.
我去年六月份猜拳月的時候.
就是講《使徒行傳》第十七章的經文.
題目就是「雅典報道會:跨文化宣教的智慧」.
那段是講保羅怎樣在雅典向外邦人宣教.
接著第十八章就提到.
保羅離開了雅典.
抵達希臘一個很重要的城市.
就是哥林多.
他在哥林多遇見了改變他宣教和生活模式的.
很重要的同工夫妻.

$^{81}$就是阿居拉和伯基拉.
在十八章第二節提到.
因為甲老丟命猶太人都離開羅馬.
有些歷史學家研究說.
在大概主後四十九到五十年期間.
當時的羅馬皇帝甲老丟下令.
將羅馬城的所有猶太人驅逐出境.
很多學者認為.
這很可能就是早期羅馬教會裡.
猶太基督徒和反對福音的猶太人之間.
針對耶穌是基督的信仰.
引起了很大的爭論甚至是紛爭.
而導致了在政治和宗教上的一些動盪.
這個事件迫使阿居拉和伯基拉夫婦逃離羅馬.
轉而來到哥林多.
這就促成了他們和保羅的相遇.
這一對神國俠女.
他們在信仰上能夠與生活事業一起相見合璧.
他們最令人佩服的地方.
首先第一點就是.
他們能夠在信仰生活和事奉上.
全然同心.
聖經說他們是以製造帳篷為業.
在當時的羅馬帝國.
這是很需要體力的工作.
但有好處就是.
可以經常接觸到很多不同的角色人等.
在市井裡.
他們不是全職的使徒.
但他們在世俗的職場當中能夠亮劍.
他們夫妻一起拿著剪刀.
拿著針線去一起縫這些皮革.
做帳篷.
他們雙手有可能已經佈滿了腦間.
但他們心中是有神國的藍圖的.
他們將每日賴以為生的帳篷的生意.
轉化為傳教的補給站和聯絡處.
他們同作帳篷.
亦都同開福音的門.
他們能夠同心去傳教.

$^{121}$在當時的羅馬社會.
夫妻能夠有共同的信仰和共同的事奉依仗.
其實是很罕見的.
但這裡的夫妻.
他們將他們的墳墓建立在同一個神聖的根基上.
這個根基就是神的國度.
在當今的傳教和尚裡.
同樣亦都有不少這樣的相見合璧的家庭.
我認識一對大概三十多歲的年輕夫妻.
丈夫是一個高級的工程師.
妻子是一個很優秀的室內設計師.
當他們看到中東創起地區有需要的時候.
他們就全家人移居到當地.
她先生就帶著一身的技能.
應聘到當地一間跨國的科技公司上班.
而他太太就利用她的設計的專長.
在當地開了一間文創的工作室.
他們都帶著兩個年幼的小朋友.
就好像一對欠著小合拍的江湖夫妻.
丈夫就在公司裡面用他的專業的技能.
贏得了他的同事的敬重.
太太就在他的工作室裡面.
和當地的婦女建立了很好的關係.
而她的兩個小朋友就在當地的小學裡面.
和其他本地的學生玩成一片.
所以這個家庭是當古代職場的.
一個神國俠侶的見證.
他們在職場上有很好的見證.
他們經常邀請他們的同事.
朋友和小朋友的同學.
包括同學的父母來到他們家裡吃飯.
就是這樣.
在那個沒辦法公開可以傳福音的地方.
他們已經帶領了差不多有十個家庭來信主.
弟兄姊妹.
你們的家庭和事業.
是否只是為了賺錢養家呢?.
還是可以成為你們夫妻聯手為神的國征戰的地方呢?.
當夫妻,父母和孩子在家庭目標上一致的時候.
將主耶穌基督頒佈的大使命作為終極的使命.

$^{161}$你們的家庭生活就會不再平庸.
而是一個充滿屬靈傳奇色彩的家庭.
我盼望我們弟兄姊妹的家庭也是一樣.
好像阿基拉和伯基拉一樣.
是一個雙劍合璧的家庭.
第二.
這一對的神國俠侶.
他們建立了福音的客棧.
他們讓他們的家成為了宣教的陣地.
在《使徒行傳》第十八章二到三節講到.
保羅就投奔了他們.
保羅因為與他們同業.
就和他們同住作供.
連使徒保羅也住在他們的家裡.
這裡講到「投奔」.
「投奔」的意思其實就是「投靠」.
即是說這個家庭是供應了保羅一些的需要.
甚至是供應他日常的所需.
另外在《羅馬書》第十六章的三至五節這樣說.
我們一起讀.
「問伯基拉和阿基拉安.
他們在基督耶穌裡與我同供.
也為我的命將自己的頸鏝自知到愛.
又問在他們家中的教會安.
又問在他們家中的教會安」.
這個就告訴我們.
他們的家就是一個小的教會.
弟兄姊妹就是他們家中聚會去敬拜神.
這個就是最早期的真正的家庭教會.
很厲害.
他們開放他們的家.
成為了福音的客棧.
也成為了福音的陣地.
他們是打開家門的宣教.
我們知道在武俠的世界裡.
客棧往往是各路英雄豪傑他們會聚的地方.
甚至是改變歷史命運的地方.
而阿基拉和伯基拉他們的家.
就成為了至關重要的福音的客棧.
無論在哥林多.

$^{201}$還是在以弗所.
還是在羅馬.
他們每到一個地方.
第一件事就是打開家門.
設立家中的教會.
當保羅孤單很疲乏的時候.
來到哥林多.
就是這個家庭向他們敞開了大門.
他們給了保羅的不單單只是一張床.
更加是家庭的溫暖和屬靈同伴的支持.
所以有時候宣教不需要你馬上遠赴異邦.
有時候只可以打開你的家門.
開放你的客廳.
就已經可以做得到了.
我認識另外一對很平凡的基督徒的夫婦.
他們住在大學城附近.
這十幾年來.
他們的家就好像一間永遠都不打烊的客棧一樣.
每個週末.
太太就會煮兩大鍋熱辣辣的滷肉飯.
先生就會去校園裡面搜羅那些孤單的.
上家的國際留學生和異鄉的青年.
那些來到他們家的年輕人.
很多是從來沒有聽過福音的.
但是這對夫妻他們沒有在家裡.
很生硬地跟他們傳教福音和做神學辯論.
而是有家的溫暖.
有無條件的接待.
在十幾年下來.
這間福音客棧.
竟然陸陸續續帶領了差不多有一百位的年輕人.
來信主.
在當中後來也有人成為了全職的傳道人.
這對夫婦就用很簡單的柴米油鹽.
打贏了一場很美麗的福音聖像.
他們開放他們的家.
他們就已經在做宣教.
另外一個故事.
也是我認識的朋友.
他們是一對美國的夫婦.

$^{241}$他們千山萬水來到中國.
收養了兩個孤兒.
這個故事其實也有很多的報章.
之前也有報導過.
這個報章就用了這個題目.
跨越國境的愛 八萬中國孤兒和他們的美國父母.
這個報道是說在中國開放以來.
有超過八萬的孤兒.
被外國的家庭收養.
而在這些家庭裡絕大多數是來自美國的家庭.
我所認識的那對美國夫婦就是其中一對夫婦.
她們領養了兩個孩子 兩個都是女孩子.
而這兩個女孩子的身體是有殘障的.
我不知道是不是因為有殘障的緣故.
所以被人遺棄.
也可能因為她們是女性等等.
總之就是被人遺棄了.
她們就領養了她們回去.
帶到美國之後 她們是用心去愛她們.
她們的生命見證 她們的愛.
就已經讓那兩個孩子很深刻地感受到耶穌基督的愛.
所以很快她們都順了主.
而且從小她們都在信仰的環境裡長大.
這對夫婦還將她們培養成運動的健將.
雖然她們殘疾 雙腳不能走路.
但她們坐著輪椅參加了很多殘障人士的比賽.
而且獲得了獎項.
有時候我在Facebook看到她們發佈出來的照片.
都很替她們開心.
因為看到她們很開心地去比賽和獲得獎項.
這個家庭都很平凡.
但她們用很簡單的領養和接待.
就將福音傳開去.
弟兄姊妹你不用怕 不是叫你去領養.
但是我們教會都有類似的服事.
親善家庭計劃 不知道大家有沒有聽過.
可能你都已經有參與過.
因為這個計劃我們都做了很多年.
接下來六月份我們都會繼續招募義工.
我盼望弟兄姊妹能夠積極去參與.

$^{281}$透過接待 透過開放你的家庭.
能夠去做到宣教的工作.
將福音帶給更多的人.
另外一位不得不說的人就是艾偉德.
其實他呢 我去年都有舉過他做例子.
但為什麼今年再說他呢.
因為說到客棧就不能夠不說他.
這個女士其實她是出生在英國.
她小時候是出生在一個很普通的家庭.
她沒有受過高等教育.
而且她曾經是做過使女 做女僕人.
在1930年的時候.
她就以傳教士的身份來到中國的山西.
在山西做傳教和慈善的活動.
後來她還加入了中國籍.
但是在1940年日軍侵華 日軍掃蕩的時候.
她就需要在她當地的客棧.
她在當地開了一個叫做百福客棧.
招待很多的人 也收養了很多的孤兒.
她就需要將差不多一百個孤兒.
由山西陽城這個地方.
要徒步翻山越嶺地送到陝西西安.
是應該超過一百公里的.
要徒步帶著差不多一百個小孩去.
是很艱鉅的任務.
而且當時日軍還在轟炸.
要避開日軍的轟炸都不容易.
但是她最後成功將差不多一百個孤兒.
平安送到陝西西安.
所以她這個壯舉.
後來荷里活就拍成了電影.
叫做六福客棧.
有機會大家可以繼續去了解一下.
她的故事是激勵了很多人.
她是一個很平凡的女子.
但是她做了很多其他人做不到的事情.
她藉著八福客棧將福音傳遞開去.
有一個弟兄曾經在網上.
去搜集艾偉德的生平資料.
在搜集的時候他偶然發現.

$^{321}$在相關的資料網頁上有網友留言.
有一個網友就說.
我的父親就是其中一位孤兒.
然後這個弟兄就很好奇.
他就順著這個留言.
去到那個網友的社交平台主頁裡面.
他就發現原來那個網友是山西籍的.
而現在他正正在非洲做醫療的服務.
他是支援非洲醫療的一個隊員.
這個是很震撼的事情.
這個是愛的傳遞.
和一百年前艾偉德來到中國一樣.
他將愛帶來中國.
他可能都想不到他帶的其中一個孤兒的兒子.
今天在非洲繼續做相同的事情.
他將愛帶到非洲這個偏遠落後的地區.
這個就是愛的傳遞.
我們回到經文那裡.
阿居拉和伯基拉這對神國的俠侶.
他們都是開放他們的家庭.
讓他們的家庭成為福音的客棧.
來做宣教的工作.
他們是有跨文化的適應.
就是從羅馬到哥林多.
從猶太到外邦.
他們因為政治驅逐而遷徙.
及在每一個新的城市.
他們能夠建立家中的教會.
我們是否都能夠像他們一樣.
願意為福音去接納新移民.
難民,外籍勞工.
讓我們的家能夠成為他們的家呢?.
你們是否願意為神的國創開福音的客棧呢?.
你們的家是否願意為那些孤單的留學生.
新移民創開呢?.
你們的客廳是否願意成為小組聚會.
福音聚會的一個據點呢?.
願我們都能夠效法這對神國的俠侶.
讓我們的家能夠成為福音的客棧.
成為宣教的陣地.

$^{361}$第三方面,他們是有信仰的傳承.
在信仰傳承的過程當中.
他們注重愛與真理的教導.
這是愛與真理並重的教導.
在《士郎行傳》第十八章二十四到二十六節這樣說.
我們一起讀一讀.
有一個猶太人名叫阿波羅來到以弗所.
他生在亞歷山大,是有學問的.
就連能講解聖經.
這人已經在主的道上受了教訓.
心裡火熱,將耶穌的事詳細講論教訓人.
只是他單曉得若翰的朕禮.
他在會堂裡放膽講道.
伯基拉雅居拉聽見,就接他來.
將神的道給他講解更加詳細.
如果我們講回武林的話.
武林真正的宗師,一代的宗師.
他們從來都不是那些道賢疾能.
和那些很保守的人.
而是很樂意去引導後輩的人.
當時以弗所來了一位很年輕的奇才.
或者容許我形容他是一個武學的奇才.
骨骼精奇,萬中無一.
他的武功底子非常好.
經文說他有口才和最能講解聖經.
但是他在真理的內心心法上還不夠全面.
他只曉得若翰的朕禮.
所以雅居拉和伯基拉聽了之後.
他們沒有在公開的場合去反對或批評他.
而是他們這對神國俠侶.
展現了很高的屬靈情懷.
他們將阿波羅拉他來.
即是請他去到他們的福音客棧.
去到他們的家庭當中.
在家庭的氛圍裡.
為這個年輕的領袖補上了一堂很關鍵的神學課.
將更加深奧的神國的心法傾囊相授給他.
這個過程就好像一個私屬門訓.
他們將阿波羅邀請到他們的家裡.
對他有些門訓去教導他,帶領他,栽培他.

$^{401}$在這裡我們見到經盟約伯記拉的名字放在前面.
新約聖經有六次提到這對夫婦的名字.
其中有四次是伯基拉太太的名字放在阿波羅拉先生的名字前.
這裡顯示他在教導的過程當中有主導的作用.
他起了很重要的角色.
這告訴我們當時他是一個很重要的女性教師.
所以我們今天也一樣.
我們的姐妹們也可以承擔很重要的教導的職事.
他們是用愛和真理去教導.
是愛與真理並重的.
他們的謙卑和智慧為早期的教會栽培出來一個很重要的神國的戰將.
就是阿波羅.
我們作為別人的父母.
在座每一位為人父母的弟兄姊妹.
我們可以反思一下.
我們是不是花很多的精力和金錢.
想要將我們的孩子培養成世俗眼中的武林高手呢?.
送他去名校.
希望他將來能夠高薪.
也有高的成就.
但是我們有沒有傳授他們神國的真正的內功心法呢?.
我們是不是常常花很多的資金送他去補習班.
送他去學一些才藝.
希望他將來能夠出人頭地.
但是我們沒有花時間全家人一起去讀經.
一起去為遠方的宣教士祈禱呢?.
我們警會有祈禱的計劃.
其中一個就是給小朋友的.
叫做「猜拳小天使」.
我相信也有很多小朋友有參與過.
這個「猜拳小天使」的計劃.
不單單是培養小朋友去為猜拳,為宣教士去祈禱.
也是培養他們從小學會奉獻金錢去支持宣教的工作.
他們所奉獻的金錢我們會集中在一起去支持有關宣教的士工.
如果我們可以和他們一起這樣做的時候.
其實我們就是在他們的心裡埋下了宣教的種子.
所以信仰的傳承.
最重要的不是外在的那些東西.
就好像學武功那樣.
學武功如果我要傳授武功給下一個年輕人的話.

$^{441}$不是那個武功的招式最重要.
而是那個內功心法是最重要.
所以將神國的真理這個內功心法要傳授給他們.
這個是這對夫婦很特別的地方.
他們雖然很平凡.
他們也不是很出名.
但是他們居然培養了一個非常之厲害的阿波羅.
我相信我們在座每一位弟兄姊妹都可以做得到.
你試一下開放你的家庭.
你試一下幫助一些年輕人.
神會使用你們.
第四方面.
這對神國俠侶是肝膽相照.
是將生死能夠自知道愛的夫妻.
他們有自知道愛的合骨柔情.
在羅馬書第十六章三到四節.
我讀給大家聽.
「問伯基拉和雅居拉安.
他們在基督耶穌裡與我同供.
也為我的命將自己的頸項自知道愛.
他們是有將頸項自知道愛的勇氣.
我不知道你們有沒有聽過一句話.
「為國為民,俠之大者」.
這句話是金庸所寫的.
在神鵰俠侶裡面.
是郭靖對他的後輩楊過所說的話.
如果我接著他這句話.
我就會說「神國俠侶為主寫命」.
使徒保羅在羅馬書中.
給了這對夫妻很高的評價.
他們是為保羅的生命.
將自己的頸項自知道愛.
自知道愛在希臘文原文的意思是.
準備好將頭放在斬首台上.
雅居拉和伯基拉是為了保護保羅.
為了福音的緣故.
他們願意承擔極大的生命危險.
當保羅被猶太人控告的時候.
被官府追捕的時候.
他們沒有退縮.

$^{481}$而是以家為盾牌收留他.
以帳棚為掩護藏匿使徒.
傳遞書信來聯絡教會.
當羅馬皇帝驅逐猶太人的時候.
他們被迫放棄他們的財產.
他們流離失所.
但他們從來沒有抱怨.
也沒有退縮.
他們收拾好行李.
背上帳棚.
就飄泊到那裡.
就在那裡建立起他們家中的教會.
建立起福音的基地.
這種隨時準備出發.
視死如歸的豪情.
正正就是神國俠侶最高的境界.
在歷史上都有很多視死如歸.
為福音勇於將生死置諸到外的人.
在這本書叫做「回首百年殉道血」.
這本書裡面有很多的見證.
其中一個就是講到顧姑娘.
顧姑娘是英國瘦羊傳道團的一個教士.
在英國的愛丁堡出世.
她是一個非常聽話和活潑的女孩.
從小就信主.
十歲的時候就被人送去瑞士讀小學和中學.
讀完之後就回到英國.
在牛津大學讀了四年的文學.
後來就拿了文學的學士學位.
之後又在中學教書教了六年.
教完書之後.
她就在1898年加入了瘦羊的宣道會.
就被派到中國山西省的太原來工作.
她的志願就是為耶穌基督竭心盡力.
她當時是負責教會女子學校的工作.
她是致力於消除對婦女兒童的歧視和偏見.
但是在1900年6月的一個晚上.
義和團.
在太原的義和團和群眾.
就包圍和燒毀了顧姑娘所在的.

$^{521}$由教會所建立的紀念醫院.
據當時的目擊者所說.
她兩度掙扎從火海當中走出來.
但是同樣兩次被人推回到火海裡面.
所以最後她在火海當中殉道.
她當時只有38歲.
她在最後一封家書裡面這樣說.
我們一起讀一讀.
到頭來 死亡對我們來說不是悲傷的.
而是神現今使我們得著生命之途.
受苦也不是完全悲傷的.
因我們知道從中學會了祝福的事.
這是一項無價的功課.
她是一個很好的見證.
她是能夠為福音的緣故.
將生死置之度外.
今日有什麼阻礙我們的家庭參與宣教呢.
通常不是我們缺乏才華 缺乏能力 缺乏資源.
而是我們對安逸和安全的過渡的執著.
我們害怕失去穩定的收入.
我們害怕我們的孩子受苦.
也害怕打破了現有的平靜.
但是亞居拉和伯基拉的生命告訴我們.
一個真正被基督的愛點燃的家庭.
是敢於為福音冒險的家庭.
當我們願意為了神的國度.
稍微放下世俗的保障的時候.
我們就會看見神蹟的發生.
最後有一封神國的英雄帖要發給大家.
邀請弟兄姊妹能夠接下這封神國的英雄帖.
這個世界其實正處於屬靈的黑暗當中.
有無數的迷茫痛苦的靈魂.
正在黑暗的江湖當中摸索著.
神今日去建立你的家庭.
賞賜你的婚姻.
祂將孩子托付給你.
不是僅僅為了讓我們能夠關起門去享受所謂的小確幸.
而是要讓你的家庭成為祂手中的一支利箭.
要射向這個黑暗的世代.
所以今日猜拳的主日.

$^{561}$有一封神國的英雄帖要發給全教會的家庭.
這封神國英雄帖帶著三個具體的挑戰.
第一個挑戰就是全家同心足談.
我邀請弟兄姊妹在接下來的這個星期.
你選擇一天的晚上.
全家人一起坐下為一個未得知民的國家.
或者為我們教會所支持的宣教士家庭.
來祈禱.
第二個挑戰就是重開福音客棧.
在接下來的這個月裡.
你邀請一位需要關懷的鄰舍.
或者新朋友.
或者教會的弟兄姊妹.
來到你家裡吃一餐家常便飯.
你和他們分享.
分享神在你生命當中的作為和見證.
第三方面的挑戰就是從小學習奉獻.
鼓勵你的孩子一起去參與猜拳的奉獻.
讓孩子學習為神國的擴張獻上自己一份的力量.
剛才我提到有猜拳朝天使計劃.
鼓勵弟兄姊妹邀請你的孩子一同去參與.
雖然我們不是宣教士.
但是我們可做宣教士.
我請弟兄姊妹一起來讀這句話.
不是宣教士.
可做宣教士.
在崇拜週刊.
在電子版的週刊裡.
有警察會的宣傳單張和回應表.
所以鼓勵弟兄姊妹積極去掃QR Code去回應.
願我們這群的神國俠侶.
我們靠著耶穌基督.
我們能夠去奪回這個世界.
我們一起同心去祈禱.
慈愛天父.
讓我們看見一對平凡的夫妻.
如何能夠發揮如此驚人的神國影響力.
主啊.
願你同樣興起我們的家庭.
興起我們當中的神國俠侶.

$^{601}$讓我們的婚姻家庭.
不再只是看見自己的需要.
而是看見天父你宣教的意象.
讓我們的職場.
我們的客廳.
我們的餐桌.
都奉獻給你.
作為宣教的祭壇.
求你賜給我們全家同心.
冒險跟隨的勇氣.
奉主耶穌基督的聖名祈求.
阿們.
\newpage



\section{羅馬書 5:17-21}
\label{sec:032266}
\textbf{靠著耶穌得勝}~~~~蕭壽華牧師
\newline
\newline
連結: \href{http://www.church.com.hk/acms/content.asp?site=cdc\&op=show\&type=product\&code=032266\&layout=sermon}{\texttt{032266}} ~~~~ 語音日期: 2025-11-09
\newline
\newline
\hyperref[sec:032265]{\small{< < < PREV SERMON < < <}}
~
\hyperref[sec:index_chronic]{\small{[返順時目]}}
~
\hyperref[sec:032267]{\small{> > > NEXT SERMON > > >}}
\newline
\newline
羅馬書 5:17-21
\newline
\begin{longtable}{cl}
\hline
\hline
章節 & 經文 (和合本修訂版)\\
\hline
5:17 & \begin{tabularx}{0.7\textwidth}{X} 若因一人的過犯，死就因這一人掌權，那些受洪恩又蒙所賜之義的，豈不更要因耶穌基督一人在他們生命中掌權嗎？ \end{tabularx} \\ \\ \relax
5:18 & \begin{tabularx}{0.7\textwidth}{X} 這樣看來，因一次的過犯，所有的人都被定罪；照樣，因一次的義行，所有的人也就被稱義而得生命了。 \end{tabularx} \\ \\ \relax
5:19 & \begin{tabularx}{0.7\textwidth}{X} 因一人的悖逆，眾人成為罪人；照樣，因一人的順從，眾人也成為義了。 \end{tabularx} \\ \\ \relax
5:20 & \begin{tabularx}{0.7\textwidth}{X} 而且加添了律法，使得過犯增加，只是罪在哪裡增加，恩典就在哪裡越發豐盛了。 \end{tabularx} \\ \\ \relax
5:21 & \begin{tabularx}{0.7\textwidth}{X} 所以，正如罪藉著死掌權；照樣，恩典也藉著義掌權，使人因我們的主耶穌基督得永生。 \end{tabularx} \\ \\
[1ex]
\hline
\hline
\end{longtable}
$^{1}$這錄音或錄影的版權是屬於基督教宣導會不各堂所有.
我們的屬靈傳統之一是以真理教導為本.
歡迎您將此錄音或錄影的連結傳給親友.
懇請注意,未經授權,請勿轉錄或剪輯本錄音或錄影.
甚願好似使徒行傳二十章三十二節下所講.
這道能建立你們,叫你們和一切誠誠的人,同得基業.
各位弟兄姊妹平安.
神賜給我們每個人都有自己地上的家庭.
神也賜給我們每一位有屬靈的家庭.
在屬靈的家庭裡面我們可以稱為弟兄姊妹.
他們是你的幫助,是你的關懷,是你的幫助.
所以我很想在今天開始講到之前.
請大家確認一下你身邊的家人.
你問問他,你姓什麼?.
我見到大家都是差不多坐在那個位置.
所以你這次問了,你下次知道了他是哪位.
那都是可以一步步更多認識你的家人.
好,今天神給我們一齊的時間.
我們再一起同心祈禱,我們仰望神.
祈禱神保佑,祈禱神保佑.
祈禱神保佑.
神保佑,祈禱神保佑.
祈禱神保佑,祈禱神保佑.
神保佑,祈禱神保佑.
當他一直看的時候,這棵樹可能不行了.
可能要死了.
但後來他隔了一段長時間之後,再走回樹林.
偶然又回到這棵樹的旁邊.
他一看的時候,很驚訝.
發現那些樹皮,有某一些纖維.
再重新生出來.
亦即交錯著,遮在.
這些傷口的上面.
是很穩固地,有新的樹皮.
蓋在受傷的.
傷口的上面.
而且他感覺,被受傷的地方.
所蓋上的新的樹皮.
比起其他地方,好像更堅固.
更穩固.

$^{41}$然後他就說,他想起聖經說.
罪在那裡顯多.
恩典就必更顯多了.
沒錯.
有時候你覺得自己軟弱到一個地步.
大概神也幫不了.
有時候你會感覺自己的問題太大了.
信徒主這麼多年,還是這樣的.
問題太大了.
但你記得勝經的應許.
罪在哪裡顯多,恩典就更顯多.
能夠大過你所面對的.
一切的錯失過犯.
能夠給你力量可以勝過.
一切的軟弱錯失.
這是神的恩典,是神的工作.
在過程當中.
你會留意到一件事情.
因為你的生命不再是罪掌權.
而是恩典掌權.
所以你再讀二十一節的時候.
聖經提到過去罪靠了死掌權.
但是今天,恩典藉著義掌權.
英文所提.
Grace reigns through righteousness.
因為神將義賜給你.
你只要繼續留在義裡面.
保留自己在光明的裡面.
你就會看到恩典在其中.
掌權,叫你有力量.
可以過平安靜寂的生活.
也因此你知道.
藉著我們的主耶穌基督.
神將永遠的生命賜給你.
各位靈姐妹,在你軟弱的時候.
不要太多看自己.
將你的眼睛轉向神.
去思想祂的恩典,祂的奇妙的能力.
你會發覺在祂的裡面.
你經歷的那種恩惠和能力.

$^{81}$不斷的彰顯出來.
我想回到.
大家常常會唱一首詩歌,奇異恩典.
這首詩歌,奇異恩典何等甘甜.
我罪已得赦免.
這首詩歌你知道是John Newton.
約翰牛頓先生所寫的.
其實當他寫的時候,他用的原來的英文.
He saved a wretch like me.
他用wretch這個字來描述.
自己那種大惡不赦的狀況.
你一定知道他過去的一些歷史.
你看他傳記的時候,你會知道他自少年期開始.
已經嫖賭飲吹,什麼都懂.
在他船上的時候.
他自己做很多的惡事.
在非洲自己做了奴隸.
在非洲做奴隸,面對很多的傷害.
但竟然他後來自己做了販賣奴隸的商人.
他曾經承認.
他說在過程當中,我殺了很多人.
未必是直接殺.
但當他運送這些奴隸的時候.
一些很細小的船隻.
很多的黑人運進去.
過程當中很多病死,餓死.
然後掉進大海.
然後去到目的地.
但過程當中死了很多人.
他說照他自己所估計.
可能有二萬人間接死在他手中.
但當他經歷到神自己的救恩臨到他的時候.
他的生命有一個很奇妙的改變.
你讀他的傳記,你會知道.
當他說奇妙恩典.
Amazing Grace的時候.
不是一些口號的說話.
是他深刻的體會和經歷.
他說這個大罪人.
神都能夠改變.

$^{121}$因此他將自己的生命奉獻了.
在他一生裡面,他將自己奉獻.
三十九歲成為一個牧師.
他會傳揚主的福音.
一生帶領數萬人歸向耶穌.
你會見到.
誰?約翰牛頓的改變.
不是,你會見到神的恩典.
如何在這個罪人身上彰顯出來.
你會更加知道.
不再是依靠自己.
而是依靠那位奇妙改變我們生命的主.
再讀聖經.
你會發覺聖經下面.
講出一個很重要的方向.
領受恩典,因此一生向罪死.
向神活出來.
六章十二節說.
所以不要讓罪在你們必死的身體上掌權.
使你們順從身體的私慾.
也不要把你們的肢體獻給罪.
當你領受了恩典之餘.
你要記得,你仍然面對著危險.
因為你仍然在身體之內.
曾經稱身體為必死的身體.
在這個世界裡仍然面對著各種不同的試探.
你的身體裡仍有盈餘的私慾.
這些私慾.
何時你自由的.
憑著你已經得到的自由.
將那種主權交回給罪.
給它掌權.
罪就會引誘你身體.
去按你的私慾行事.
叫你們順從身體的私慾.
以致你盡心落在罪裡面.
或者叫你裡面的罪.
死灰復燃.
本來已經死了.
不警覺著.

$^{161}$讓罪再次掌權.
罪就呼喚你裡面的私慾.
你聽從私慾的意思.
去犯罪,以致得罪上帝.
你記得.
當主耶穌要釘十字架.
其實是面對著一個很大的挑戰.
當時彼得在旁邊.
知道這件事之後.
彼得對主耶穌說.
這事必不能到你.
你不可以死.
你不可以被釘十字架.
其實彼得也有一份體貼耶穌的心.
不想祂死,不想祂受苦.
但是你會記得主耶穌怎樣回答他.
祂叫他做撒旦.
去吧,因為你只是體貼人的意思.
不體貼神的意思.
何時你不體貼神的意思.
你可以變成撒旦.
但是你發覺主耶穌很清楚地跟他說.
撒旦,退我後邊去吧.
我上星期聽了一位姐妹分享.
她說沒運氣的時候.
有些試探來到.
她自己知道有一種掙扎.
很想做,但又覺得不應該做.
然後她就跟自己說.
撒旦,退我後邊去吧.
你不一定要說這句話.
但是你內心有一個很清楚的方向.
耶穌說,撒旦,你沒有權柄來主宰我.
不容許罪在你身上再次掌權.
撒旦就要逃跑了.
就如淫亂,污穢,邪情,惡慾和貪婪的事情.
還有你要去棄絕一切的事情.
是什麼事情呢?.
包括了老恨,憤怒,惡毒,誹謗的說話.
和口中污穢的言語.

$^{201}$因為這些都是屬於舊的人.
聖經羅馬書第六章說到.
你舊的人已經死了.
你舊的生命已經與耶穌基督釘在十字架上.
死了,不再有作用.
你不需要理會他,你不需要回應他.
我們太快會忘記了.
我們信了主很多年.
忘記了我們的舊人其實已經死了.
所以法國聖經裡面一路提醒我們.
你都要留意死了的人已經擺脫了罪.
我們已經擺脫了罪.
因此你看自己向罪是死的.
你向罪對你的試探,對你的引誘的時候.
你要知道你已經死了.
因此你不用回應他.
你不用說,是嗎?我以前做慣了.
好吧,這次再做吧.
我想用一個比喻,大家會明白我的意思.
如果你參加一個安息禮拜.
你看到有一幅棺木放在台中間.
安息禮拜完結之後,大家出去瞻仰為榮.
大家帶著公敬的心去瞻仰這位過去的人.
走著走著,突然看到一位男士按著棺木.
看著裡面的人.
喂,你這麼快就死了嗎?.
你答應要幫我找一份好工作.
到現在還沒找到就死了嗎?.
你快點回來吧.
你猜那個死人會不會推開蓋子說.
對不起,我還沒幫你找到.
我現在搞不定,幫不了你.
大家明白我的意思嗎?.
罪人已經死了.
當罪人向你呼喚的時候,你不用理會他.
他在你身上沒有權柄.
他不會因為你不聽他的話而死.
你就會損失.
如果有試探在你裡面出現,告訴你.
你做一點點動作,你就會得到一筆大金錢.

$^{241}$你過去都做過了,為什麼不做?.
很有好處.
你的心告訴他,我向你已經死了.
我不需要反應,不需要理會你.
第一次會抓住你的權柄去拒絕罪.
去致死你的身體的私慾.
你能夠一直走在光明的路上,一直跟隨主.
在過程當中,剛才你記得勝經十三節的一句話.
不要把肢體獻給罪,作不義的工具.
聖經裡面提到,有些時候我們會留意.
我們的身體會帶我們有些行動,有些事情會做出來.
但請告訴我們,要留心不要讓你的肢體獻給罪.
成為了不義的工具.
我想回到一件事情,就是過去一段時間.
我們很多時候都會看到一些關於騙詐園區.
因為我剛剛從泰國回來.
有些人說,這裡很接近KK園.
我就想起很多事情.
你會記得很多人被欺騙去到這些地方.
但這些被欺騙的人,他們會說.
其實那些人是我認識的,是我的舊同學.
舊朋友,很久沒見,突然打電話給我.
有份筍工,你過來做吧,就被騙了.
但後來他們發現一些事情,報導繼續說.
這些人在這些園區被虐待,被毒打.
身上出現很多傷痕.
有些情況說,有些曾經欺騙他們的人.
回看這些照片,他們很後悔.
為了自己少少的利益,竟然做出這樣的事情.
令自己的舊朋友,竟然受到這麼大的創傷.
弟兄姊妹你會體會,罪是帶來傷害的.
是帶來死亡的.
何時你不為意,你就讓罪在你身上掌權的時候.
讓自己的身體成為易的工具.
不一定是你直接犯罪.
但你帶來了死亡,帶來了傷害.
聖經提醒,不要把你的肢體獻給罪作不義的工具.
倒要像從死中復活過來的人一樣.
怎樣做?.
將自己獻給神,並把你們的肢體獻給神.

$^{281}$你自己整個人,包括你的身體,獻給神.
作義的工具.
聖經所說的獻,是大家熟悉的.
《南無書》十二章第一節.
「所以弟兄們,我以神的慈悲勸你們,將身體獻上」.
是同一個字.
這個獻其實有一個很重要的含義在裡面.
就是旁邊站立的意思.
什麼叫做旁邊站立呢?.
你在神的面前旁邊站立.
你準備好隨時回應祂對你的吩咐.
是一種很積極的態度.
隨時有主的提醒,有主的引導.
就走在這條光明的路上.
隨時奉獻自己做光明的事情.
積極地將自己的身體心靈奉獻給主.
作義的工具.
不是作不義的工具.
工具在很多其他經文裡面也叫做兵器.
是在說你作神公義的徵兵.
為義來征戰.
你知道在這個罪惡的世界裡面.
有很多罪惡的權勢在不斷地去威嚇很多不同的人.
但是你要作義的徵兵.
去驅逐黑暗的事情.
去帶來光明的事情.
因此你領受恩典之餘.
你知道神不單單要讓恩典停留在你生命當中.
要讓你的生命成為了義的器皿.
當你將自己的生命獻呈在神手當中.
神要用你將更多的光明帶進這個黑暗的世界.
在那時候你會心裡面很快樂.
因為你會見到神的能力在你生命當中顯出來.
你會真的知道什麼是能力.
什麼是恩典.
來到第六章的十四節,這一段最後一節.
聖經就說.
罪必無法克制你.
再一次重複地提,罪必無法克制你.
因為你們不在律法之下.

$^{321}$意味著我們不再憑自己去行義.
憑著自己去對抗罪.
你不是,而是在恩典之下.
各位師妹,我們一起說.
我們在恩典之下.
你可以放膽面對一切的挑戰.
知道神的恩典成為你的能力.
你可以向罪死.
可以向神活著.
獻呈自己成為義的器皿.
我們開始時說,By grace I stand.
所以,詩歌提醒我們.
剛才師班,不要怕,不要動一下就走路.
你放心,你放心,走前面的路.
跟隨主耶穌,仰望著恩典.
你可以一步一步走在神給你的蒙恩光明的路上.
這個時候,我請你祈禱.
你為自己祈禱.
求神讓你在前面更深切地信靠神的恩典.
你努力地,積極地行義.
你為自己祈禱.
愛我們的主,我們何等的蒙恩.
一個在死亡中間被罪所控制的人.
竟然今天可以在主耶穌的復活裡有份.
主讓我們,竟然可以願意銅釘十字架.
救人死去,領受你給我們新生命.
主讓我們在你面前恭敬的感恩.
來仰望親愛的主.
在我們人生的路上靠著你.
我們繼續能勝過一切盈餘的邪惡.
讓我們靠著你的恩典,放膽無懼來面對.
主讓我們知道,你的恩典就是能力.
我們有時候會跌倒,有時候會軟弱.
但主我們仍然依靠你.
讓我們知道,你永遠是我們的盼望.
我們祈求天上的主,感動我們的心靈.
知道將自己的人,自己的身體.
你都完全奉獻在你的手當中.
能夠被你所用,能夠成為義的器皿.
讓我們蒙恩之餘,我們能夠見證你的恩典.

$^{361}$讓這個世界,因為你的功績,會更多一點的光明.
聽我們祈禱,奉主耶穌的聖名祈求.
\newpage



\section{馬可福音 1:1-45}
\label{sec:032267}
\textbf{不再無間道?}~~~~秦潤祥牧師
\newline
\newline
連結: \href{http://www.church.com.hk/acms/content.asp?site=cdc\&op=show\&type=product\&code=032267\&layout=sermon}{\texttt{032267}} ~~~~ 語音日期: 2026-06-07
\newline
\newline
\hyperref[sec:032266]{\small{< < < PREV SERMON < < <}}
~
\hyperref[sec:index_chronic]{\small{[返順時目]}}
~
\hyperref[sec:032268]{\small{> > > NEXT SERMON > > >}}
\newline
\newline
馬可福音 1:1-45
\newline
\begin{longtable}{cl}
\hline
\hline
章節 & 經文 (和合本修訂版)\\
\hline
1:1 & \begin{tabularx}{0.7\textwidth}{X} 神的兒子，耶穌基督福音的起頭。 \end{tabularx} \\ \\ \relax
1:2 & \begin{tabularx}{0.7\textwidth}{X} 正如以賽亞先知書上記著：「看哪，我要差遣我的使者在你面前，他要為你預備道路。 \end{tabularx} \\ \\ \relax
1:3 & \begin{tabularx}{0.7\textwidth}{X} 在曠野有聲音呼喊著：預備主的道，修直他的路。」 \end{tabularx} \\ \\ \relax
1:4 & \begin{tabularx}{0.7\textwidth}{X} 照這話，施洗約翰來到曠野，宣講悔改的洗禮，使罪得赦。 \end{tabularx} \\ \\ \relax
1:5 & \begin{tabularx}{0.7\textwidth}{X} 猶太全地和全耶路撒冷的人都出去，到約翰那裡，承認他們的罪，在約旦河裡受他的洗。 \end{tabularx} \\ \\ \relax
1:6 & \begin{tabularx}{0.7\textwidth}{X} 約翰穿駱駝毛的衣服，腰束皮帶，吃的是蝗蟲和野蜜。 \end{tabularx} \\ \\ \relax
1:7 & \begin{tabularx}{0.7\textwidth}{X} 他宣講，說：「有一位在我以後來的，能力比我更大，我就是彎腰給他解鞋帶也不配。 \end{tabularx} \\ \\ \relax
1:8 & \begin{tabularx}{0.7\textwidth}{X} 我用水給你們施洗，他卻要用聖靈給你們施洗。」 \end{tabularx} \\ \\ \relax
1:9 & \begin{tabularx}{0.7\textwidth}{X} 那時，耶穌從加利利的拿撒勒來，在約旦河裡受了約翰的洗。 \end{tabularx} \\ \\ \relax
1:10 & \begin{tabularx}{0.7\textwidth}{X} 他從水裡一上來，就看見天裂開了，聖靈彷彿鴿子降在他身上。 \end{tabularx} \\ \\ \relax
1:11 & \begin{tabularx}{0.7\textwidth}{X} 又有聲音從天上來，說：「你是我的愛子，我喜愛你。」 \end{tabularx} \\ \\ \relax
1:12 & \begin{tabularx}{0.7\textwidth}{X} 聖靈立刻把耶穌催促到曠野裡去。 \end{tabularx} \\ \\ \relax
1:13 & \begin{tabularx}{0.7\textwidth}{X} 他在曠野四十天，受撒但的試探，並與野獸同在一起，且有天使來伺候他。 \end{tabularx} \\ \\ \relax
1:14 & \begin{tabularx}{0.7\textwidth}{X} 約翰下監以後，耶穌來到加利利，宣講神的福音， \end{tabularx} \\ \\ \relax
1:15 & \begin{tabularx}{0.7\textwidth}{X} 說：「日期滿了，神的國近了。你們要悔改，信福音！」 \end{tabularx} \\ \\ \relax
1:16 & \begin{tabularx}{0.7\textwidth}{X} 耶穌沿著加利利的海邊走，看見西門和西門的弟弟安得烈在海上撒網；他們本是打魚的。 \end{tabularx} \\ \\ \relax
1:17 & \begin{tabularx}{0.7\textwidth}{X} 耶穌對他們說：「來跟從我，我要叫你們得人如得魚一樣。」 \end{tabularx} \\ \\ \relax
1:18 & \begin{tabularx}{0.7\textwidth}{X} 他們立刻捨了網，跟從他。 \end{tabularx} \\ \\ \relax
1:19 & \begin{tabularx}{0.7\textwidth}{X} 耶穌稍往前走，又見西庇太的兒子雅各和他弟弟約翰在船上補網。 \end{tabularx} \\ \\ \relax
1:20 & \begin{tabularx}{0.7\textwidth}{X} 耶穌隨即呼召他們，他們就把父親西庇太和雇工留在船上，跟從了耶穌。 \end{tabularx} \\ \\ \relax
1:21 & \begin{tabularx}{0.7\textwidth}{X} 他們到了迦百農，耶穌就在安息日進了會堂教導人。 \end{tabularx} \\ \\ \relax
1:22 & \begin{tabularx}{0.7\textwidth}{X} 他們對他的教導感到很驚奇，因為他教導他們正像有權柄的人，不像文士。 \end{tabularx} \\ \\ \relax
1:23 & \begin{tabularx}{0.7\textwidth}{X} 當時，會堂裡有一個污靈附身的人，他在喊叫， \end{tabularx} \\ \\ \relax
1:24 & \begin{tabularx}{0.7\textwidth}{X} 說：「拿撒勒人耶穌，你為甚麼干擾我們？你來消滅我們嗎？我知道你是誰，你是神的聖者。」 \end{tabularx} \\ \\ \relax
1:25 & \begin{tabularx}{0.7\textwidth}{X} 耶穌斥責他說：「不要作聲，從這人身上出來吧！」 \end{tabularx} \\ \\ \relax
1:26 & \begin{tabularx}{0.7\textwidth}{X} 污靈使那人抽了一陣風，大聲喊叫，就出來了。 \end{tabularx} \\ \\ \relax
1:27 & \begin{tabularx}{0.7\textwidth}{X} 眾人都驚訝，以致彼此對問：「這是甚麼事？是個新的教導啊！他用權柄命令污靈，連污靈也聽從了他。」 \end{tabularx} \\ \\ \relax
1:28 & \begin{tabularx}{0.7\textwidth}{X} 於是耶穌的名聲立刻傳遍了全加利利周圍地區。 \end{tabularx} \\ \\ \relax
1:29 & \begin{tabularx}{0.7\textwidth}{X} 他們一出會堂，就同雅各和約翰進了西門和安得烈的家。 \end{tabularx} \\ \\ \relax
1:30 & \begin{tabularx}{0.7\textwidth}{X} 西門的岳母正發燒躺著，就有人告訴耶穌。 \end{tabularx} \\ \\ \relax
1:31 & \begin{tabularx}{0.7\textwidth}{X} 耶穌進前拉著她的手，扶她起來，燒就退了，於是她服事他們。 \end{tabularx} \\ \\ \relax
1:32 & \begin{tabularx}{0.7\textwidth}{X} 傍晚日落的時候，有人帶著一切害病的和被鬼附的，來到耶穌跟前。 \end{tabularx} \\ \\ \relax
1:33 & \begin{tabularx}{0.7\textwidth}{X} 全城的人都聚集在門前。 \end{tabularx} \\ \\ \relax
1:34 & \begin{tabularx}{0.7\textwidth}{X} 耶穌治好了許多害各樣病的人，又趕出許多鬼，不許鬼說話，因為鬼認識他。 \end{tabularx} \\ \\ \relax
1:35 & \begin{tabularx}{0.7\textwidth}{X} 次日早晨，天未亮的時候，耶穌起來，到曠野地方去，在那裡禱告。 \end{tabularx} \\ \\ \relax
1:36 & \begin{tabularx}{0.7\textwidth}{X} 西門和同伴出去找他， \end{tabularx} \\ \\ \relax
1:37 & \begin{tabularx}{0.7\textwidth}{X} 找到了就對他說：「眾人都在找你！」 \end{tabularx} \\ \\ \relax
1:38 & \begin{tabularx}{0.7\textwidth}{X} 耶穌對他們說：「讓我們往別處去，到鄰近的鄉村，我也好在那裡傳道，因為我是為這事出來的。」 \end{tabularx} \\ \\ \relax
1:39 & \begin{tabularx}{0.7\textwidth}{X} 於是他走遍全加利利，在他們的會堂傳道，並且趕鬼。 \end{tabularx} \\ \\ \relax
1:40 & \begin{tabularx}{0.7\textwidth}{X} 有一個痲瘋病人來求耶穌，向他跪下，說：「你若肯，你能使我潔淨。」 \end{tabularx} \\ \\ \relax
1:41 & \begin{tabularx}{0.7\textwidth}{X} 耶穌動了慈心，就伸手摸他，說：「我肯，你潔淨了吧！」 \end{tabularx} \\ \\ \relax
1:42 & \begin{tabularx}{0.7\textwidth}{X} 痲瘋病立刻離開他，他就潔淨了。 \end{tabularx} \\ \\ \relax
1:43 & \begin{tabularx}{0.7\textwidth}{X} 耶穌嚴嚴地叮囑他，立刻打發他走， \end{tabularx} \\ \\ \relax
1:44 & \begin{tabularx}{0.7\textwidth}{X} 對他說：「你要注意，千萬不可告訴任何人，只要去，讓祭司為你檢查，又因為你已經潔淨，獻上摩西所吩咐的祭物，作為證據給眾人看。」 \end{tabularx} \\ \\ \relax
1:45 & \begin{tabularx}{0.7\textwidth}{X} 那人出去，倒說許多的話，把這件事傳揚開了，使耶穌不能再公開進城，只好留在外邊曠野地方，人從各處都到他跟前來。 \end{tabularx} \\ \\
[1ex]
\hline
\hline
\end{longtable}
$^{1}$剛剛我們已經說了為了這個奉獻祈禱.
我對最後那首回應詩挺有感動的.
大家喜不喜歡聽那首回應詩.
我們一會兒那首回應詩叫做《靠主因》.
那個名字叫什麼名字 我忘記了.
《一生不走錯》.
為什麼我們要唱《一生不走錯》呢.
因為我們一生裡面經常都走錯.
我們經常都要靠主因.
大家看到這幅圖嗎.
這幅圖是一個海.
看到海面挺平靜的嗎.
但看到裡面有很多暗湧嗎.
我們一生裡面有很多暗湧.
在我們的所謂的潛意識裡面.
心底裡面有很多解釋不了的.
但是包圍著我們.
例如我們原生家庭.
我們的創傷.
我們的所有的罪.
各方面等等的事情.
你有沒有看到有些探測器在這裡.
是不是.
我們很多時候行事為人的時候.
我們裡面雖然有很多暗湧.
但我們是會探測著自己在別人面前.
我們表面上是比較平穩的.
這個如果你看心理學.
這個應該叫做人格的面具.
又或者面具.
有面具不代表不好.
因為適當的面具.
是可以令我們在人際關係裡面.
是不走錯的.
但是如果你不處理裡面的暗湧.
你就會很辛苦了.
有沒有聽過.
為什麼回家就發脾氣.
為什麼對別人比對我好.
就是因為回家.

$^{41}$回家的時候你就放鬆.
看到那些探測器沒有了嗎.
放鬆之後就決堤了.
很多時候我們裡面有很多暗湧.
用家庭這個單位來做一個例子.
就是你看心理醫生也好.
你去探測什麼也好.
你回去的時候.
他給了你一個像夫的面具.
讓你會好一陣子.
但是如果你裡面的暗湧不處理好.
當那個面具你都受不了.
你要決堤的時候.
一樣是有同樣的問題.
所以我們要一生不走錯.
我們真的要靠著神的恩典.
去慢慢醫治好我們裡面每一個暗湧.
或者我們可以說的是傷口.
我們很多時候戴著面具做人.
如果太多的時候就會很累.
相信大家都看過一套戲叫做《無間道》.
看過沒有.
都看過是吧.
希望我們今天一起不再無間道.
無間道的主角是哪兩個.
梁朝偉和劉德華.
就是一個面具裡面的真實寫照.
劉德華本身是一個什麼.
賊來的.
但是他偏偏做警察.
是不是.
是很辛苦的.
又怕一樣又怕那樣.
梁朝偉本來是什麼.
警察來的.
他走去做賊.
做臥底.
也是很辛苦的.
如果你看《無間道》.
一二三.

$^{81}$你會看到他們兩個到最終都睡不著.
因為有很多事要處理 他們是受不了的.
記不記得在天台那場經典的戲.
文戲.
劉德華說.
可不可以給我一次機會.
我想做回一個好人.
你記不記得梁朝偉怎麼說.
跟法官說.
看他給不給你做個好人.
記不記得.
是吧.
想做但他又做不到.
因為很多東西纏累著.
很多暗湧在裡面.
但是如果你今天.
你問旁邊的弟兄姊妹.
可不可以給我一次機會.
我想做回一個好人.
弟兄姊妹這樣說.
跟上帝說.
他一定可以給你做好人.
是不是.
可不可以給我一次機會.
可不可以讓我跳.
可不可以讓我痊癒.
跟上帝說.
上帝恩典救理容.
靠著他.
將會一生都不走錯.
我們要靠著神的權能.
什麼是權能.
他的authority.
他的權柄.
權柄就是能力.
他的authority.
他的權柄能夠平息風浪.
記不記得.
在福音書裡面說.
他可以平息風浪.

$^{121}$就是我們的神.
神今天都可以在我們心裡面.
平息我們所有的風浪.
所有的暗湧.
希望我們可以靠著今天的經文.
讓我們都可以得以慢慢地痊癒.
慢慢地成聖.
享受神給我們的恩典.
我們今天會讀馬可福音第一章.
幾節而已.
我們先一齊讀.
21至27節.
1,2,3.
「多了加伯隆.
耶穌就在安息日.
進了會堂教訓人.
眾人很稀奇他的教訓.
因為他教訓他們.
正像有權柄的人.
不像民事.
在會堂裡有一個人.
被污鬼附著.
他喊叫說.
那撒勒人耶穌.
我們與你有什麼相干.
你來滅我們嗎.
我知道你是誰.
乃是神的聖者.
耶穌責備他說.
不要作聲.
從這人身上出來吧.
污鬼叫那人抽了一陣風.
大聲喊叫就出來了.
眾人都驚訝.
以至彼此對問說.
這是什麼事.
是個新道理.
他用權柄.
吩咐污鬼.
連污鬼也聽從他」.

$^{161}$如果你看回這段經文裡面.
出現了兩次權柄.
在第一章裡面權柄出現了兩次.
成為一個開括號.
刪括號.
括著了驅魔.
有民事的驅魔.
我以前上一堂的教授.
就說他是驅過魔的.
他什麼時候驅魔.
就是當有人想決戰的時候.
突然間嘔吐.
嘔吐然後變聲.
然後驅魔.
用神的大能大力.
他真的驅魔了.
他這樣告訴我.
我就這樣轉述.
但是弟兄姊妹.
你知不知道.
對於我來說.
驅魔就容易.
拉他出來.
靠著耶穌基督的名字.
但是如果要醫治我們心裡面的暗湧.
是難很多的.
因為我們有很多時候.
心裡面有很多暗湧未醫治的東西.
比撞鬼還慘.
我自己知道.
我不知道你是否一樣.
大家都是一個有罪的人.
我們需要耶穌基督的權柄.
他就說權柄不像民事.
民事的權柄和耶穌基督.
都是說當時同一本聖經.
那有什麼分別呢.
弟兄姊妹.
權柄就是大能大力.
權柄裡面有能力.

$^{201}$這個權柄在《創世記》第一章一至三節.
都已經說了出來.
起初神創造天地.
地是空虛混沌.
淤泥黑暗.
神的靈渾騰在水面上.
神書要有光就怎樣.
神書珠水之間要有空氣.
就將水分為上下.
見不見到神的權柄和能力.
他可以將水分為上下.
令到中間有空氣.
窮瘡.
成為我們現在可以住的地方.
這個能力就是神的權柄和大能.
不單止這樣.
在這段經文裡面.
你見到三位一體的神在運行.
他有個fellowship在裡面.
他有關係在裡面.
這個關係裡面不單止有權柄.
能力還有愛.
弟兄姊妹.
我們要回到這個權能裡面.
享受三一真神裡面的fellowship.
享受祂給我們的愛.
我們才可以得以痊癒.
希望大家記得神的權柄就是大能大力.
剛才我們說將水和水分開這個能力.
不單止在創世記那裡.
是見到我們在耶穌基督.
即是馬可福音.
即是剛才我們所讀的經文上文.
我們都可以看到.
我們一起讀好嗎.
一二三.
他從水裡一上來.
就看見天裂開了.
精靈彷彿甲子降在他身上.
又有聲音從天上來說.

$^{241}$你是我的愛子.
我喜悅你.
你見不見到將水和水分開的能力.
在耶穌基督的浸泥裡面.
將天念裂開這個能力.
在耶穌基督浸泥的場景裡面.
發揮得淋漓盡致.
不單止這樣.
這個經文也彰顯了.
三一真神的fellowship.
有聖靈.
有父神.
有耶穌基督.
剛才說神說的時候.
我們今天成為基督徒.
我們就明白.
道成了肉身.
充充足足在我們中間有恩典.
有真理.
今天這個真理.
都在我們的聖經上告.
如果你想被傳喻.
被醫治.
我們每人都要看這本聖經.
加上敬虔.
俯伏的敬拜.
我們將會在人生裡面.
得到了改變.
記住將天裂開的那一位神.
是今天我們每一位基督徒.
的天父上第三一真神.
權柄就是等於什麼.
大能大力.
這個大能大力.
不單止在剛才的浸禮出現了.
還在馬口福音比較後面.
你記不記得耶穌上了十字架之後.
馬口福音怎樣登記呢.
我們一起讀好不好.
一二三.

$^{281}$耶穌大聲喊叫.
氣就斷了.
電裡的萬子從上到下.
裂為兩半.
對面站著的白夫長.
看見耶穌這樣喊叫斷氣.
就說.
這人真是神的兒子.
你見不見到.
萬子是從上到下裂為兩半.
你知不知道萬子大概是怎樣的.
和家裡的地氈.
你有沒有試過一扯.
扯爛地氈.
由上至下扯爛地氈.
有沒有人試過.
扯不扯到.
扯不到.
你知不知道萬子的厚度.
不是手掌的厚度.
它有五層樓那麼高.
兩層半樓那麼寬.
由上至下這樣撕開.
需要大能大力.
我有時很無聊.
我就和AI談一談.
究竟怎樣才可以.
將一個這樣的萬子撕開.
AI告訴我.
一定要用F35.
F35.
美國剪輯機.
F35C.
因為它有個渦輪推動器.
這個畫得不太好.
但不要緊.
他說這兩個剪輯機.
一定要用那些鋼絲.
很粗的那些.
來吊著萬子的兩邊.

$^{321}$你不可以直扯.
一直扯就無限大力就會斷.
所以它要V字形飛出去.
還要停空.
因為它有渦輪.
反作用力.
停住空慢慢扯出去.
如果那兩個剪輯機.
角度和力度剛剛好.
它就會慢慢將萬子撕裂出來.
弟兄姊妹.
你見不見到不可能的任務.
見不見到神的權能.
當年有沒有F35C.
沒有.
今天你會不會拿兩架F35C.
做這些事.
不會.
但神當見到耶穌成了之後.
上了十字架.
成就了救恩.
祂將聖殿裡面的萬子.
祂拿手撕開了這個大能大力.
讓我們今天都可以.
來到施恩寶座面前.
蒙恩衛.
德蓮率.
作隨時的幫助.
這個就是神的權能.
這個就是神的大能和大力.
所以你會見到.
當耶穌祂說不要作聖.
不是跟我們說.
跟鬼說不要作聖.
從這人身上出來吧.
他就出來了.
命辣就辣.
要有光就有光.
要將水和水分開.
就有窮滄.

$^{361}$這個就是祂的權柄和大能大力.
今天這個權柄和大能大力.
一樣可以平息我們心裡面.
所有的風浪.
所有的暗湧.
希望我們都可以.
不再做屬靈的臥底.
放下面具.
三年又三年.
要做多少年呢.
人生有多少個三年呢.
希望我們可以記得.
我們的第一點.
我們一起讀.
讓神的權柄.
進入心中的混沌.
我特意用這個字.
讓.
讓神的權柄.
先進入我們內心的混沌.
我們的關鍵是要怎樣.
慢慢打開.
我們的心.
是慢慢打開我們的心.
看到這個心有什麼做的嗎.
鑽石.
很硬.
是不是.
驅魔容易.
一扯一扯出來.
那個不是你.
好像吸塵機.
一吸一吸出來就搞定了.
但這個是你的心的時候.
我們裡面有很多的創傷.
就這樣打開.
是會扯傷我們.
神給我們自由意志.
我們要學習.
慢慢打開我們每一個人的心.

$^{401}$我們要學習謙卑俯伏.
放下自己的防禦機制.
我們有很多的防禦機制.
我們不想神搞我們.
我們不願意打開我們的心.
因為我們有很多很多的防禦機制.
例如.
我們要放下我們的否認.
我們很好的了.
OK了.
就好像彼得.
彼得三次被厭主.
他都是OK的.
應該的.
這麼危險.
又或者大衛.
有沒有人聽過大衛的故事.
贏在尼亞.
輸又在尼亞.
誰是尼亞.
烏尼亞.
他貪戀別人的太太.
犯了奸淫.
又把她殺了.
聖經登記了.
她是生了出來.
之前殺了別人的老公.
即是懷胎都要十月.
大概一年時間裡面.
他都不覺得自己錯.
他沒有否認.
或者神都沒有責備我.
或者OK.
很多的否認.
我們要放下我們的否認.
我們都要放下我們的.
合理化.
當時的環境很難.
我都是受害者.
如果不是慈禧大後.

$^{441}$我們會被人侵佔.
我們又不會輸毒.
又不會捱窮.
不會做賊.
大概是這樣.
人不為己.
天誅地滅.
那條蛇叫我這樣做.
如果不是你做那條蛇出來.
我都不會這樣做.
即是你都要負一些責任.
合理化.
又或者我們要放下我們的手指指.
拿丹對大衛說.
你就是那個人.
神就派拿丹去罵大衛.
說了故事.
有個有錢人有很窮的人.
窮的人只有一隻羊.
有錢的人有羊群牛群.
但有錢人招呼客人的時候.
他就拿了窮的那隻羊.
那隻羊別人當是如珠如寶.
一齊抱著睡.
他拿了殺了請人吃飯.
大衛說這個人我以最碎屍萬段.
跟著拿丹就說.
這個人就是你.
指著他.
但大衛他不覺得.
他在說他.
我們很多時候看聖經都是他在說他.
很多時候我們要看聖經怎樣指著我們去說話.
大衛都是一個成功和失敗的例子.
我們要放下什麼.
我們的面具.
我們在神面前的面具.
下面這個不是經文.
OK.
要不要唱.

$^{481}$明明可以晝夜無間.
踏著面前路.
夢想中的彼岸.
怎樣.
為何還未到.
你戴著面具.
你去不到彼岸.
OK.
你永遠都是一個無恥境的無間道裡面.
不停在這裡.
將自己暗湧.
湧出來.
又收回去.
湧出來.
又收回去.
我們要放下我們的面具.
即是說我們要放下我們的自卑.
OK.
Jesus knows the worst about you.
Nathanael is the one who loves you the most.
即是說神愛你.
就不是因為你戴了面具.
你做得有多好.
神是一早已經看見你裡面有多不好.
但祂都愛你.
行不行.
你真的要想清楚這句是為什麼.
我覺得講得很好.
即是耶穌基督.
祂知道你心底裡面最差是怎樣.
你有多奸詐.
多壞.
多無恥.
都知道的.
但祂就是這樣愛你.
不是因為我們知道自己很無恥.
即是講者.
下流賤格.
接著我就做好一點.
祂才愛你.

$^{521}$不是.
在恩典裡面.
神是看見你最殘的樣子.
祂都愛你.
行不行.
即是好像我太太那樣.
發聲.
是吧.
我說很殘.
我說你多殘.
我都是這樣.
我不會走的.
是吧.
祂都愛你.
神和我們的關係.
祂要帶我們回到三一神的團契裡面.
無論你是怎樣.
祂都怎樣.
愛你行不行.
記得了.
如果你明白剛才那一句說話.
你才真正可以放下自己的面具.
放下自己的自卑.
你才真的可以體會到神怎樣去光照你的內心.
否則你會被遮蓋.
你是一個錫心.
不是一個柔軟的心.
第二點我們一起讀.
在神的愛中.
重整並傾倒恩典.
你讓神的權力進入你的混沌裡面.
你就可以在神的愛裡面重整.
記不記得創世記.
重整.
讓你有秩序.
慢慢去重整你的內心的時候.
就可以將恩典倒下去.
我們上次剛剛教完創世記第一章.
神是要將所有客觀的事情做好.
才倒一些生物進去.

$^{561}$讓他享受這個恩典.
今天耶穌基督的救恩.
讓我們裡面的空虛混沌.
再一次在我們屬靈的世界裡面重新創造.
重新拯救.
我們一起去經歷平靜風浪的美好.
好不好.
我們裡面有很多不同的暗湧.
AI挺好用的.
方便很多.
一張圖可以抵擋我十個slide.
暗湧很多時候我們好像職場的精英.
劉德華是警察.
又和一哥打高爾夫球.
光鮮亮麗.
但他裡面其實有很多的罪.
有很多的東西纏繞著.
有很多的煩惱.
焦慮.
夫妻也是.
你有看八卦新聞.
明明很好.
他們不是剛剛出來一起拍照.
後來又不行.
戴很多面具.
不是他們錯.
他們是公關.
但頂不住的時候就爆.
甚至宗教也是.
文士的教導.
剛才說沒有權兵.
文士都可以戴著面具.
牧師都可以戴著面具.
把經文包裝得很漂亮.
叫大家這樣做.
以前舊約.
安息日你要守到的話.
神就要立你.
你又犯罪.
騙制.

$^{601}$你做所有的事.
都不代表你內心的所有暗湧可以被醫治.
所以作為一個基督徒.
在神的恩典裡面.
借著十字架的功能.
我們要讓神的光進入我們深層裡面.
你會見到我選了四樣.
原生家庭的創傷.
有時有原生家庭創傷.
你就會有焦慮.
小時候他不稱讚我.
例如媽媽不稱讚.
但這就是焦慮.
這個我很喜歡.
你見到這個是誰.
梁朝偉.
梁朝偉這一幕在哪裡做.
就是黃秋生在天台跌下來.
死了的那一幕.
那一幕真是影帝.
他有很複雜的情緒在裡面.
第一他就死了.
之後立即想.
我身份就沒有人知道.
我永遠都是做賊.
又想起大家的感情.
我們感情那麼好.
你就這樣死了.
又傷心.
又想死了怎麼辦.
他又內疚.
就這樣走了.
那個複雜性在那幾秒裡面.
發揮得淋漓盡致.
人就是那麼複雜.
我們人生裡面都很複雜.
有很多焦慮.
慢性焦慮 急事的焦慮.
讓我們很多時候犯罪.
要那些什麼.

$^{641}$安多芬.
推上去的時候.
你就feel better.
慢慢上癮 就好像吸毒.
但你每次犯罪之後.
你心靈又空虛.
有人試過嗎.
心靈空虛.
講得輕一點.
剛剛shopping完.
心靈又空虛.
心靈空虛又去shopping.
你永遠都無止境.
這就是暗湧.
夢想的彼岸永遠都未到.
弟子們.
這些都是藤掙瓜掙藤.
原生家庭的創傷.
令我想起一個歷史人物.
瓦丁路德.
Martin Luther.
他有很多原生家庭.
即是改教的那位 瓦丁路德.
他有很多原生家庭的傷害.
其中一樣就是他爸爸很嚴.
經常打他.
所以他經常覺得神都一樣那麼嚴.
所以他很害怕.
他有很多焦慮.
他以前做他的神父.
就很慘.
他告解你知道說了多久.
六個小時.
說來說去.
都還是那麼焦慮.
那個神父.
就開OT了.
我不知道怎麼聽.
六個小時.
還要隔著玻璃.

$^{681}$很慘.
因為他有很多焦慮.
有一次走雷閃電的時候.
瓦丁路德怕到要躲起來.
他怕真的被雷劈中.
因為他覺得自己太多罪.
就算他做多少的修行.
去禁食.
去做修道士.
他都覺得自己還不行.
神一定不喜歡我.
直至到.
他自己查經的時候.
神的光.
照落他的心淵裡面.
他看見羅馬書.
那裡說.
福音是恩順稱義的.
本符恩.
恩著信.
所以如經上所記載.
二人.
邊緣似得身.
他才明白.
神的光照遍他的內心.
他人生不再一樣.
他有勇氣.
不然我們今天都坐不在這裡.
我們就做不了基督教.
那就沒有Living Grace.
亮透了瓦丁路德.
還有一個奧古斯丁.
有沒有人聽過.
一個以前的教父.
他又相反.
他有很多罪.
雖然他學哲學.
他有很多罪.
他有去喝酒.
吸煙.

$^{721}$亂搞男女關係.
但他心裡面就很空虛.
剛剛才犯完罪.
又空虛.
讀完哲學又空虛.
來來來去的時候.
有一次在米典的後院.
他看到一本聖經.
突然有個小朋友打開來看.
然後一打開.
羅馬書十三章十三字.
你行事為人要端正.
保障在白晝.
不要喝酒.
不要淫蕩.
要披戴基督.
不要為身體.
來安排.
不要有那麼多情慾.
就是這樣.
他被神的光.
亮透了他心中的混沌.
他就得以改變.
成為了我們的教父.
弟子姐妹你看到.
我們個個都是一樣.
我們個個都.
多不少原生家庭的創傷.
有很多上網看的見證.
原生家庭.
覺得爸爸媽媽不愛他.
缺少了愛.
就會成為一個.
外面很想.
多很多愛.
或者你會.
轉換拍檔.
來得到多些愛.
證明自己.
成為一個反人格的個性.

$^{761}$各方面.
但如果.
我們得不到.
應該有的愛.
無論社會也好,父母也好.
公民提醒我們.
神向我們說.
你是我的愛子.
我喜悅你.
我們都是神的兒女.
神.
是無條件的愛我們.
無論你有多骯髒.
剛才說的.
你最差最差.
那一格他已經看到.
他都呼召你.
他都愛你.
我們很多時候.
有很多慢性的焦慮.
為什麼我們會有焦慮呢.
我們為什麼有焦慮.
有沒有人透氣有焦慮.
很焦慮.
你記不記得創世記.
神創造天地的時候.
他創造飛鳥.
馬太福音說.
不種也不收.
為什麼沒有焦慮.
因為神設計他是飛的.
他飛的時候.
為什麼沒有焦慮.
他很開心.
洋溢著自由的.
拍翼而飛的豪邁.
如果你叫一條魚.
走上去飛.
他就會焦慮.
或者你叫小鳥.

$^{801}$像魚一樣游.
他就會焦慮.
很多時候我們在人生當中.
因為我們是犯罪的世界.
我們上班.
不是很享受的工作.
希望回到教會的時候.
你會沒有了焦慮.
因為你.
回到神設計.
你應該做的事.
你就做了.
我見到有很多.
弟兄姊妹做的事.
你讓我做.
很焦慮.
例如煮飯.
有沒有人在這裡吃過飯.
嘩.
我就會縮小了.
有魚蛋有奶茶.
有蘿蔔糕.
我只能安排.
你叫我做魚飯.
我還要擔心不夠.
太多.
不夠的時候要很聰明.
打電話去訂東西.
有沒有人做到這些事.
姐妹們做到.
很累.
他OK的.
你OK的.
你沒有焦慮.
多厲害.
很多時候我們真的要回到.
神給我們做事.
你叫我在後面.
做什麼.
跟Terence聊天.

$^{841}$他覺得沒什麼.
他在行.
你叫我做這些事.
我已經很焦慮.
我不可以繼續看下去.
我要看著你.
每一個人神設計我們.
都是不同的.
我們的焦慮要慢慢被醫治.
但你要記得.
我們回到我們的calling的時候.
焦慮慢慢醫治.
我們要學習怎樣去俯伏.
怎樣去在神面前謙卑.
我人生有很多焦慮.
我經常提醒.
我讀神學的時候.
是焦慮的.
剛剛信主八個月走去神學院.
然後.
每個本聖經都爛了.
我本聖經是新的.
我真的不焦慮都不行.
慢慢頂到.
覺得都應該OK.
我最不怕哪一科.
就是港道學.
我小時候去做話劇.
我小時候都差不多.
這樣就糟了.
不行的那些神又OK.
覺得可以的那些.
神又不讓你OK.
他磨練我.
我本來在一間中文的神學院讀.
當年叫他讀英文神學院.
我就焦慮.
我不去.
神真的扯了我.
扯了我出去.

$^{881}$去了英文神學院.
我在那裡之後才考港道學.
這樣就非常焦慮.
我英文.
好像本來手打字.
要用腳打字.
你明不明白我的意思.
第一堂的時候.
他叫我講些什麼.
你知不知道.
我講馬丁魯德金.
I have a dream.
整篇.
叫我.
Preach出來.
好像撞鬼.
他叫我.
學習他的Passion.
我就看.
我真的沒有.
我不是那種.
I have a dream.
我差點講了I have a dream.
I have a dream.
要等他停定位.
My 42 children.
will one day live in a nation.
where there will be no judging.
我忘記了.
他說.
the color of their skin.
I have a dream.
聽到那些人都看著我.
就像現在這樣的反應.
我都不知道.
怎麼走下去.
後來神就.
教導我.
在呼召學會謙卑.
我都感恩.

$^{921}$經歷了這些.
之後.
哎呀都要講到英文.
否則真的講不到.
在呼召裡面.
都會有焦慮.
但你慢慢謙卑在神面前.
焦慮慢慢減退.
你就被抑制.
很多的罪.
為什麼有很多的罪.
焦慮就會.
找些事情做.
好事就很難做.
你沒理由很焦慮.
你去做義工撿垃圾.
焦慮的時候.
犯罪當然是最容易.
最爽.
但犯罪的時候.
你要治好焦慮.
神在聖經裡面告訴我們.
貧勞苦擔重擔.
到我姐女來.
我就使你安息.
你就少點罪.
罪的時候.
有時候我們很沒面子.
又戴個面具離開神.
算了.
我還是不回教會.
就算回教會.
都要洗個澡才回去.
但你要記得.
神愛我們.
不是因為我們外面的光鮮亮麗.
因為父神愛愛愛到.
差派他的兒子.
來到這個世界.
為我們死亡.

$^{961}$我們今天才可以站立.
所以不是因為我們做任何的事.
而是因為神愛我們.
我們才可以.
在神面前站立.
我們心靈才不會空虛.
聖經裡面說.
我要使你.
活得更加豐盛.
有時候做人.
心裡好像很亂.
好像少了塊肉.
少了塊肉心.
以西結書給你肉心.
你要在神的救恩裡面.
你真的得到的時候.
你才感覺完整.
才是完美.
所以在神的眼中.
我們每一個人.
都是完美.
可以嗎?.
你不想看旁邊那個人?.
你覺得對方.
是不是完美?.
好像不是.
但你再看他.
想想耶穌基督的救恩.
完美.
我們都要在.
這些挑戰裡面.
學會如何.
放下我們的面具.
就算做牧師也是.
裡面都是很多東西.
我們不是特別神聖.
我們都是沒那麼多.
有很多不好的東西在裡面.
很多時候都要戴著面具.
適當的面具是可以的.

$^{1001}$但如果太多面具.
就假的.
假的就是假冒為善.
所以我們其實每一個.
基督徒我們都要學習.
如何去被醫治.
這個就是Living Grace所說的.
Being before Doing.
當你Being.
開始痊癒的時候.
你可以Do一下.
Do的時候神又操練你又Heal一下.
Do到累的時候你可以Be一下.
Being, Doing, Being, Doing.
所以如果這裡有弟兄姊妹做到累了.
休息一下.
如果你覺得Being都可以.
好像有一點痊癒了.
又Do一下.
一起讀一次.
1,2,3.
你會記得.
我們一定要.
被神在愛中為我們重整.
傾倒恩典.
我們讀三點一起讀.
神的恩典.
會出主裡面的光芒.
我們其實是有光芒的.
但是我們被很多罪.
遮擋住.
變成沒有光芒.
我們藉著神的恩典.
可以活出主裡面的光芒.
你看看.
剛才的經文裡面.
馬福音15節說.
對面站著的白夫長看見.
是不是.
剛才我們看.

$^{1041}$耶穌浸禮的時候.
馬可用的文法.
只有耶穌看見.
旁邊的人.
沒有看到.
但是耶穌基督.
他豐富地.
活出神應該.
覺得人應該有的.
樣式之後上十字架.
為了我們流寶血.
死了之後.
對面站著的.
白夫長愛邦人都看見.
行不行.
神的豐盛.
在我們裡面可以被人看見.
不單止這樣.
剛才下文說驅魔.
驅魔之後.
西滿的岳母.
Mother-in-law.
西滿的岳母.
正害疫病躺著.
就有人告訴耶穌.
耶穌進前拉著她的手.
扶她起來.
熱就退了.
她就事奉他們.
天晚日落的時候.
有人帶著一切害病和被鬼附的.
來到耶穌.
跟前.
你看見單數.
接著是雙數.
她被醫好就來服侍.
弟兄姊妹.
我們今天都可以成為.
這個世界的一點光芒.
就好像.

$^{1081}$在爾賽亞書.
在馬可福音之前大概.
或者九百年.
還是多少年.
很多年前.
那受過痛苦的.
必不再見幼吟.
從前神使西布倫蒂和.
拿不他利地.
被藐視.
沒後卻使者沿海的路.
若大河外.
外邦人的加里利得著榮耀.
在黑暗中.
行走的百姓.
看見大光.
住在死刃之地的人有光.
照耀他們.
在馬太福音有講這個經文.
耶穌就是那光.
而今天我們.
怎樣可以應用到.
今天我們就是那光.
因為耶穌先照我們.
我們可以照遍.
這個世界.
我們本來是死刃之地的人.
就好像當年被亞述游輪.
沒有希望.
但耶穌基督來到的時候.
祂的國原來是天國.
我們知道神的國.
在我們心中顯明的時候.
活出來.
最後.
記住兩個字.
一二三.
立刻.
在馬太福音.
立刻這個字.

$^{1121}$原文裡面出現了三十六次.
在第一章裡面出現了九次.
即是出現了四分一次.
即是很重要.
你見到那個馬可.
在唸書的時候.
他不是用一些很精美的希臘文來寫.
他是用一些很粗糙的希臘文來寫.
來直達一個很立刻要做.
很立刻.
不是喝那個.
很立刻應該要做的事.
就要立刻去回應.
所以如果你.
有防禦機制.
你未信耶穌.
如果你真的知道神是真的.
立刻信耶穌.
好不好.
不要再等了.
那首歌怎麼唱.
夢想的彼岸 為何仍未到的時候.
神已經在愛著你.
又或者你已經信了.
但是你不浸禮住.
有些東西卡住.
或者這個是你其中一個防禦機制.
立刻在神面前.
放下面具.
用祂的光照遍你的內心.
很多時候是神先醫好你.
你先經歷祂.
行不行.
做一個見證給別人看.
又或者你已經信了主.
但是你還纏累在很多罪.
很多陰影裡面.
立刻打開自己的心.
我們都是鑽石心.
我們的心很硬.

$^{1161}$我們在聖經一起.
我們一起去敬拜 一起去團契的時候.
我們都會暴露我們的缺點出來.
暴露缺點不是一件壞事.
在教會裡面.
我們有缺點 我們互相接納的時候.
你看到神的愛.
會慢慢被醫治.
好不好.
如果你很久沒有回教會.
你覺得.
我洗了澡都不想回 都不夠聖潔的時候.
立刻回來.
因為信仰是要.
在人前面 在一個群體.
由我至到你們.
由你至到他們.
這個就是.
神的榮耀.
這個就是神的權柄.
和大能 透過我們將大光.
發放出去 所以希望.
大家都不再無間道.
先讓神的權柄.
進入我們心裡面的混沌.
再透過他的愛.
重整.
得到恩典之後 再將恩典.
撥出來 展現我們.
基督徒的光芒 好不好 我們一起祈禱.
愛我們愛到底.
我們都不是很底的.
給我們有這麼多的恩典.
讓我們在糜爛的.
內心世界裡面 可以.
一滴一滴地被醫治.
痊癒 改變.
活得越來越像你.
祈禱是奉主 耶穌基督名其求.
Amen.

\newpage



\section{羅馬書 8:18-27}
\label{sec:032268}
\textbf{苦難世界的盼望 (華)}~~~~葛國光牧師
\newline
\newline
連結: \href{http://www.church.com.hk/acms/content.asp?site=cdc\&op=show\&type=product\&code=032268\&layout=sermon}{\texttt{032268}} ~~~~ 語音日期: 2026-06-07
\newline
\newline
\hyperref[sec:032267]{\small{< < < PREV SERMON < < <}}
~
\hyperref[sec:index_chronic]{\small{[返順時目]}}
~
\hyperref[sec:032269]{\small{> > > NEXT SERMON > > >}}
\newline
\newline
羅馬書 8:18-27
\newline
\begin{longtable}{cl}
\hline
\hline
章節 & 經文 (和合本修訂版)\\
\hline
8:18 & \begin{tabularx}{0.7\textwidth}{X} 我認為，現在的苦楚，若比起將來要顯示給我們的榮耀，是不足介意的。 \end{tabularx} \\ \\ \relax
8:19 & \begin{tabularx}{0.7\textwidth}{X} 受造之物切望等候神的眾子顯出來。 \end{tabularx} \\ \\ \relax
8:20 & \begin{tabularx}{0.7\textwidth}{X} 因為受造之物屈服在虛空之下，不是自己願意，而是因那使它屈服的叫他如此。 \end{tabularx} \\ \\ \relax
8:21 & \begin{tabularx}{0.7\textwidth}{X} 但受造之物仍然指望從敗壞的轄制下得釋放，得享神兒女榮耀的自由。 \end{tabularx} \\ \\ \relax
8:22 & \begin{tabularx}{0.7\textwidth}{X} 我們知道，一切受造之物一同呻吟，一同忍受陣痛，直到如今。 \end{tabularx} \\ \\ \relax
8:23 & \begin{tabularx}{0.7\textwidth}{X} 不但如此，就是我們這有聖靈作初熟果子的，也是自己內心呻吟，等候得著兒子的名分，就是我們的身體得救贖。 \end{tabularx} \\ \\ \relax
8:24 & \begin{tabularx}{0.7\textwidth}{X} 我們得救是在於盼望；可是看得見的盼望就不是盼望。誰還去盼望他所看得見的呢？ \end{tabularx} \\ \\ \relax
8:25 & \begin{tabularx}{0.7\textwidth}{X} 但我們若盼望那看不見的，我們就耐心等候。 \end{tabularx} \\ \\ \relax
8:26 & \begin{tabularx}{0.7\textwidth}{X} 同樣，我們的軟弱有聖靈幫助。我們本不知道當怎樣禱告，但是聖靈親自用無可言喻的嘆息替我們祈求。 \end{tabularx} \\ \\ \relax
8:27 & \begin{tabularx}{0.7\textwidth}{X} 那鑒察人心的知道聖靈所體貼的，因為聖靈照著神的旨意替聖徒祈求。 \end{tabularx} \\ \\
[1ex]
\hline
\hline
\end{longtable}
$^{1}$其實這個跟接下來講苦難世界好像有點不太對稱.
但是其實這個世界就是如此.
你會發現一方面.
我們在這邊有一群歡的人.
但是在世界的另外一個角落.
就有一群在受苦的人.
美國跟伊朗的戰爭不知道什麼時候會結束.
是不是.
我們看見在波斯灣那邊的國家.
最近我們跟一些韓國宣教士.
特別有一位是在卡塔爾的金穆斯在講.
他說因為這一次的戰爭.
所以今年上半年.
甚至到下半年.
很多準備要在杜拜舉行的聚會都cancel掉了.
因為大家覺得太危險了.
飛彈不時會打過來.
其實他說也沒那麼嚴重.
他們在那邊的人沒有感覺那麼嚴重.
但是大家一聽到那邊.
這個情形不知道什麼時候纔會結束.
我們看見美國跟伊朗的談判.
所以大家現在已經有點麻木.
不知道到底最後談的怎麼樣.
也不曉得到底誰講的是對不對.
你可以看見這個世界到處充滿了苦難.
加沙纔剛剛結束.
可能大家已經忘掉他們了.
但是那兩百多萬的人.
其實現在很多時候.
當我們想到當初他們在那邊嗷嗷待哺.
等候著救援的這些糧食來到.
現在這些苦難還繼續在那個地方.
世界已經忘掉了他們.
但是神沒有忘掉.
當我們看見這些苦難的時候.
我們想這些發生在這個世界上面真的跟我沒有關係.
我在美國這個地方還是舒舒服服地喫我的hamburger.
還是在我這個地方過我自己的生活.
但丟子妹即使是今天在美國你也會發現.

$^{41}$也有許多人在受苦的中間.
是嗎.
昨天我們二女兒舉行baby shower在我們家裡面.
所以小女兒跟我到downtown去幫她.
拿花.
downtown我們都可能有去過.
但是去過的時候大概我們不會在這個時候.
早上大概八點左右的時候.
其實我們七點半就出發了.
我到那邊去看.
看到那一大片的那個堵都是賣花的.
可是一路上面你去看都是homeless.
一個一個的tamp在那個地方.
然後你就看見那些人就好像一群棄嬰一樣.
精神有問題.
有的還在那邊吸什麼東西.
好像是一個沒有人管的一個地方.
加州政府在做什麼.
美國在做什麼.
這許多的這些homeless的人.
我們看多了就在Torrance.
你天天經過你會看到這些人.
I'm hungry.
當然我們教會有在做homeless工作.
可是這個太複雜的一個問題.
今天不是要講怎麼做homeless的工作.
但是我要講的就是說.
今天就在我們所生活的環境當中.
你就看到有許多人.
他們其實是過著一個很苦的生活.
還不僅僅這種physical.
這個身體上面的.
這個肚子餓的這種苦.
無家可歸 沒有地方可以住的這樣的苦.
你我是不是在生命當中.
我們也會經歷苦難 會嗎.
從小到大.
你從來都沒有經歷過.
沒喫過苦的 舉下手.
除了你來講.

$^{81}$也是有了.
現在都準備.
在這個苦難的世界當中.
就像今天我們所讀的羅馬書第八章.
18節到27節.
其實整個羅馬書第八章.
就是在講這個苦難的事情.
那我們要來看看.
苦難的原因到底是什麼.
為什麼這個世界會有苦難.
首先第一個.
在聖經裡面很清楚的讓我們看見.
罪是從哪裡來的.
罪從哪裡來的.
羅馬書說罪跡因義人.
而進入了世界.
這個義人是誰.
亞當.
奇怪了 聖經為什麼不講夏娃.
因為是夏娃先喫了 對不對.
是不是.
對不對 聖經讀懂嗎.
是不是.
姐妹都不說 是不是夏娃先喫了金果.
那為什麼亞當.
變成他擔罪了.
罪經因他而領入世界.
應該因夏娃吧 對不對.
所以男人.
你要做頭.
是要擔責任的.
妻子犯罪.
你要當 最關鍵於是你.
這個罪.
因著.
人類的墮落 始祖的墮落.
所以罪就進入了這個世界.
好 那我們來看看.
在詩篇裡面.
107篇17節.

$^{121}$這裡說 愚妄人因自己的過犯.
和自己的罪孽.
便受苦楚.
彼得千數第二章20節 我們一起來唸.
好不好 來 預備 請.
你們若因犯罪受責打.
能忍耐.
有什麼可誇的呢.
但你們若因行善受苦.
能忍耐.
這在神看.
是.
可喜愛的.
好 在這個地方就讓我們很清楚看見有兩種的受苦.
對不對.
第一個就是因為犯罪而受苦.
另外一個是什麼 因為行善而受苦 是不是.
所以.
第一個原因.
這也是我們今天所需要的.
知道的.
就是因為犯罪而受苦.
跟這個犯罪受苦有關係的.
是.
神藉著.
他.
管教我們.
當我們犯錯的時候 小孩子犯錯的時候.
你會不會去管教他 會不會.
一個好父親好母親會不會管教孩子.
會不會 會.
真正的愛是一定要管教的.
不管教的愛叫做溺愛 是嗎.
所以他說.
他說.
主所愛的他必管教.
你如果看希伯來書12章這個上下文裡面.
他講到的是.
我們因為犯罪.
然後所以呢.

$^{161}$神要管教我們.
要讓我們.
結出義的果子來.
要讓我們聖潔.
所以.
當我們面臨到苦難的問題.
聖經裡面也是這樣子告訴我們.
記不記得耶穌醫治一個生來瞎眼的人 有沒有.
那麼.
最後耶穌說什麼.
這個人病是因為什麼.
耶穌說這個人病.
這個人的眼睛瞎.
不是因為他的罪.
是不是.
所以表明一件事情.
耶穌也提醒我們.
有一些時候.
這個疾病也好.
或者這些苦難也好.
是跟罪有關係的.
所以我常常跟弟兄姊妹們.
特別是在我們健康的時候提醒.
因為生病的時候就.
不好講這個話 懂我意思吧.
你今天去探訪說 第一個你先認罪.
這個牧師好無情.
所以健康的時候我跟你講.
當我們生病的時候 我們第一個先來反省.
我有沒有什麼地方.
得罪神 得罪人的.
今天我們是聖餐 聖餐裡面也這樣提醒我們.
他說.
我們要.
省察自己.
為什麼.
因為如果我們不省察自己.
我們就會喫喝自己的.
罪.
接下來說什麼呢.

$^{201}$他說所以在你們中間的.
軟弱的.
還有呢.
生病的.
也不少 看到沒有.
所以第一個.
苦難是因為罪所帶來的.
苦難是因為神要管教我們.
管教的目的是要讓我們能夠結出.
義的果子.
也就是.
能夠合乎神心意的和神標準的.
這一個生命的果子要出來.
神藉著.
當我們犯罪的時候 他藉著鞭打我們.
因為他愛我們.
凡是他所愛的他必管教.
只有私生子.
神就不管了 對不對.
這是希伯來書在這裡講的.
苦難使得我們能夠.
感覺到.
甚至是苦難使我們.
去覺悟到神存在.
今天很多的人不認識神.
但是當他碰到.
苦難的時候.
他會開始思想神.
在英國的曼切斯特.
有一箇中國大陸來的學者.
他碰到很大的困難.
他不知道該怎麼解決.
他的英文.
也不夠好.
功課沒有辦法做下去.
經濟又絕居.
有一天早上起來.
他就大聲喊著說.
"My God, where are you?".
"我的神啊,你在哪裡?".

$^{241}$住在隔壁的人聽見了.
心裡想.
怎麼隔壁這個無神論的人會喊我的神呢.
那天中午.
這個人也不知道為什麼.
他痛苦到一個地方.
他就喊出"My God".
顯然在美國呆了一段時間了.
不然的話他應該喊什麼.
"我的媽".
我們中國人都喊媽.
痛苦到一個地步說"我的媽".
他喊什麼.
"My God".
這才美國人才會這樣"My God".
就那一天他就喊出"My God".
無神論的人怎麼喊"My God"呢.
你看到沒有.
苦難.
叫人會去覺悟到神的存在.
他不一定是相信有神哦.
請記得.
但是他這句話裡面.
他在苦難的中間.
他開始思想.
這個世界.
應該要有神吧.
不然的話像我這種狀況該怎麼辦呢.
是不是.
親愛的聽眾姐妹.
很多時候.
我們不認識這位神.
所以我們會犯罪.
當我們犯罪之後呢.
我們經驗到神的管教.
你會發現.
很多人因此就來認識神了.
是不是.
那對我們基督徒來講也是如此.
你想想看.

$^{281}$如果沒有苦難.
我們會知道我們所走的路.
可能是越走越偏離神.
我們都不曉得.
是不是.
最近加州.
聽說加州.
最近一個消息.
聽說加州.
在洛杉磯.
不是隻有洛杉磯.
整個加州.
我知道的就是.
我上網去查了一下.
在我們洛杉磯的.
大概downtown有一些什麼Sunset Boulevard.
就是有一些比較交通繁忙的地區.
已經設了攝像頭.
所以你Speeding.
你就會接到.
這個單子罰單來告訴你要去交款.
所以大家小心點.
跟旁邊說你小心點.
那我的重點不是在.
告訴你這個交通的.
那聽眾姐妹請問.
如果.
沒有告訴你說.
這個地方速限65.
請問你會開多少.
啊.
80啊.
蠻客氣了.
假設這個整個Freeway上面都沒有人.
晚上.
12點.
你會開多少.
你懂我意思嗎.
如果沒有一個標準.
我們人.

$^{321}$就會為.
就是真的是為所欲為是不是.
所以上帝給了我們一個標準是什麼.
神的話.
告訴我們說什麼該做什麼不該做.
什麼是合神心意什麼是不合神心意.
但是.
當我們不明白神不認識神的時候呢.
我們就很容易去.
越了那個標準.
是不是我們就犯罪了.
那麼我們犯罪怎麼辦呢.
神就透過管教懂嗎.
透過苦難.
所以當你出車禍的時候.
你不要馬上就是怪都是他的錯.
你看那個人開車這麼快撞到我了.
其實當你講他開車這麼快的時候.
你平常開車多少.
懂我意思嗎.
丟死唄.
所以.
苦難第一個原因是因為罪所帶來的.
第二個.
第二個我們來看.
我們一起來讀.
馬太福音第五章第四節來預備請.
畏異受逼迫的人有福了.
因為天國是他們的.
這是耶穌在登山寶序裡面所說的.
畏異受逼迫的人.
異這個異就是神的標準.
畏著神的標準.
畏著神的話.
受逼迫的人有福了.
因為天國是他們的.
請注意啊.
登山寶訓.
他所指的.
針對的.

$^{361}$不是隻是對一個信.
一個不信的人.
他是對信徒說的.
這個標準為什麼這麼高.
因為.
這是對得勝者的一個要求.
所以這裡講天國是他們的.
不是指的說可以上天堂.
不要把我們的信仰變成是隻有這麼淺.
這裡說天國是他們的.
表明的.
就是.
這樣子的人.
畏著神的話而受逼迫的人.
天國是他們的.
他們在神的國度裡面.
要做得勝者.
要與主一同統掌王權.
瞭解這個意思嗎.
彼得前書第三章十四節.
他這裡也說你們就是為義受苦.
也是怎麼樣.
有福的.
跟你旁邊說.
為義受苦.
就有福了.
為義受苦.
所以.
當苦難來臨.
我們首先要來問.
你不是跟那些不認識神的人在講說.
為什麼這個苦難臨到.
你應當問的是問說主啊.
什麼原因造成這苦難呢.
第一個.
先反省.
有沒有什麼的罪.
可能是導致這個罪.
可能導致這個苦難來臨.
今天開車超速.

$^{401}$本來限速65麥.
開到90麥.
因為要去參加.
女兒的畢業典禮.
可不可以.
因為女兒的畢業典禮.
時間很緊張.
快到了快遲到了.
可不可以.
不可以.
這麼沒有人情啊.
是是是是.
非常好.
你們都很懂啊.
兄弟妹妹.
這也是我們在講加勒臺書一直在講的.
律法教人治罪.
你可以看見我們多軟弱.
你可以看見我們很容易.
就是我們可以給自己找很多的理由.
然後我們去做了不應該做的事情.
是不是.
可是.
如果我們是為著主的緣故.
為著神的話的緣故.
我們受這些苦.
這樣子.
這是有福的.
我們接下來.
我們來看一下.
這有一個圖表.
把整個羅馬書.
剛纔我們讀的這段經文.
一直到.
二十五節.
十八到二十五節.
整個做了一個對比.
我們來看看.
保羅在這個地方說.
我們現在的苦楚.

$^{441}$對比將來的.
榮耀.
所以當你在苦楚的時候.
你現在在苦楚的時候.
不要只看見現在的苦楚.
你要看見什麼.
你要看見將來的.
榮耀.
再來看.
他說我們現在是在.
空虛之下.
是不是.
那麼.
將來呢.
將來會享受自由的榮耀.
是不是.
好.
那麼我們現在是嘆息勞苦.
可是將來我們要怎麼樣.
我們要得著兒子的名分.
我們的身體要得俗.
好.
那我們稍微解釋一下.
因為.
我們平常比較.
平常比較在講現在.
我們經歷苦難也是現在.
我們在現在嘆息勞苦.
這些我們都很明白.
但是丁宏哲妹你注意看.
如果我們今天把整個.
從.
過去到永恆.
我們來看這個圖表.
這是一個末世圖.
我們來看.
好.
我們來看看.
在這個圖表裡面.
主耶穌降生.

$^{481}$然後到了末世的時候.
聖經上提到說.
你如果要看祁書裡面.
會有七連大災難.
對不對.
七印 七號.
這些的印 號.
這些的碗.
都是代表這個災難.
而且是越來越激烈的.
越來越厲害的.
那麼這些末世的大災難裡面.
當然有神學的不同的.
有災前被提的.
災後被提的.
有災中被提的.
有多次被提的.
那這些又不是我們今天要探討的.
但是我要講的就是說.
羅馬書保羅在這裡講說.
將來那個榮耀.
他在講什麼.
他在講到說.
我們這個身體.
我們現在這個身體.
最近我們去探訪爺媽媽.
我想我們當中有些人認識爺媽媽.
爺媽媽現在在hospice.
在家裡面.
她因為多年洗聖.
現在已經沒有辦法再洗了.
所以就放棄洗聖了.
那身體就開始一下下.
可能現在情況就是很不好的情況.
但我們去看她的時候.
我們看到她.
其實你如果看過爺媽媽健康的時候.
你就知道.
這個真的是.
就是面容 身體 形態.

$^{521}$都已經跟以前很不一樣了.
是吧.
但親愛的兄弟妹妹.
如果我們今天看到這樣子的一個情況.
我們真是覺得人生這麼多的苦難.
是吧.
好像我們周遭.
當我們進入年紀也比較大了之後.
我們現在去探訪.
去關心的.
最常去的都是看什麼人.
都是看生病的人.
最常參加的都是什麼禮.
是不是.
當我們看著這些的時候就覺得說.
人生一下子又走到那個地步了.
是嗎.
當我們想到.
身體漸漸衰殘.
寶寶在這邊形容說.
我們的身體漸漸衰殘.
是不是.
但是他說什麼.
但是我們的內心怎麼樣.
一天心是一天.
所以保羅在這個.
羅馬書第八章裡面.
他講到這個苦難的時候.
現在的苦楚將來的榮耀.
講到這些的時候.
他後期還講什麼.
他講到.
將來的榮耀.
然後他講他說.
他我們所盼望的.
看到沒有.
所盼望的.
我們盼望什麼.
他說我們得就是在乎盼望.
他講到我們的身體得受之後.

$^{561}$接著就說我們得就是在乎盼望.
只是所見的盼望不是盼望.
誰還盼望他所見的呢.
保羅在講什麼.
保羅意思就是說.
我們將來那個榮耀.
就是我們今天的盼望.
弟兄姊妹.
現在我們身體漸漸衰殘.
我們開始喫東西牙齒這個也掉.
那個也掉.
要植牙了.
要植幾顆牙.
他說要植八顆牙.
那一顆牙多少錢.
兩千塊美金.
那八顆牙多少錢.
我們現在算的都是這些東西.
所以到底要去哪邊去植呢.
回中國植比較便宜.
機票這些加起來也不見得便宜.
所以好了.
可不可以有一副牙齒是不會壞的呢.
上帝你造牙的時候為什麼還會壞呢.
有沒有想過這個問題.
眼睛開始白內障了.
清光眼了.
沒這些東西那就像我一樣.
就老花眼了.
總之我們的眼睛開始四茫茫.
然後發蒼蒼.
今天早上梳頭的時候突然發現.
這邊一塊這邊一塊特別白.
很奇怪.
它為什麼不擺在後面.
它擺在前面.
就是我們開始感覺到我們的身體漸漸衰殘了.
是不是.
以後也不好了.
椎間排突出了.

$^{601}$五十肩.
現在提早四十肩了.
我們身體不好漸漸衰殘.
到最後有一天你會發現躺在床上的就是我們.
是不是.
但是保羅要我們不要停在這個地方.
不要只看見現在的這個苦楚.
當你看見現在這些苦楚的時候.
你就會在那邊嘆息勞苦.
你要看見了將來的榮耀.
他講到說我們身體得俗.
指的是什麼.
就是指的我們身體復活.
保羅在羅馬書.
在哥林多前書裡面講到復活的時候.
他講什麼.
他講到說我們將來的復活.
是不是.
講的這是我們的盼望.
如果我們沒有復活.
耶穌基督沒有復活.
那我們就不能夠復活了.
它就是那個復活出熟的果子.
所以我們也要復活.
保羅講到這個地方的時候.
他最後發出了一個偉大的宣告.
他說死啊.
你的毒鉤在哪裡.
死啊.
你的權勢在哪裡.
罪的權勢在哪裡.
他好像在向他挑戰.
他說死也好 罪也好.
這些都不能夠.
因為靠著那愛我們的主.
我們在這.
要在這一切事上得勝有餘.
羅馬書最後的結論.
第八章最後的結論告訴我們說.
靠著愛我們的主.

$^{641}$在這一切的事情上都得勝有餘了.
那個將來的榮耀.
將來的盼望.
就是新天新地的盼望.
就是主耶穌基督要再來.
這也是為什麼每一次聖餐的時候.
弟兄姊妹 我要講跟主再來有關係的信息.
因為主要再來是我們的盼望.
一點不是消極.
主再來不代表說我們在地上.
就乾脆渾渾噩噩過日子.
變賣財產也不要工作 不是這個意思.
因為主耶穌要再來.
所以我們在地上是有託付的.
是不是.
所以想到主耶穌要再來.
當主再來的時候 新天新地降臨.
啟示錄告訴我們說.
在那裡不再有眼淚 不再有哀哭.
也不再有.
不再有這一切的.
一切都更新了 哈雷路亞.
一切都更新了 阿們.
所以當你想到那個日子的時候.
我們要問自己.
我將來是不是可以在這個榮耀的裡面.
第一個問題你要問的.
我是不是將來這個新天新地.
上帝的國度來臨的時候.
我是不是在裡面有這個把握嗎.
有嗎.
我要一再地問大家.
多年前有一個老弟兄過世之前.
我去看他.
有人就提醒我說.
這個老弟兄不知道到底信主沒有.
家教會就很多了.
那有人說.
看他聖餐的時候爺爺拿起來喝.
聖餐也喫的.

$^{681}$是不是表示說他有受險.
也不一定.
對不對.
有的人他也不管.
反正東西來就照喫吧.
所以我說好 那我去探訪他.
我要把這個事情搞清楚.
他平常昏昏睡睡的.
當我去探訪他的時候.
他人很驚醒.
我就跟他講.
我問他說.
我說你.
後來我確定他沒有受險.
在教會訓會十年.
也沒有受險.
一個老弟兄.
姓藍.
每次喫聖餐還喫.
根本搞不清楚狀況.
等那一天在他.
清醒的時候.
我帶領他清楚的接受了耶穌基督.
哈雷厄亞.
過了一個禮拜他走了.
但是我心裡面何等的歡喜.
在他走之前.
他清楚他自己的.
舊恩怨.
Amen.
丟著妹他有把握.
有一天他離開這一個衰退的身體.
他要進入哪裡去.
進入到新天新地神的國度裡面去.
聖經告訴我們說用什麼身體進去呢.
不是全身插著管.
那一副.
憔悴摔船的那個面膜.
不是.
是嶄新的身體身體上說是靈性的身體是榮耀的身體.

$^{721}$是復活的身體.
我們沒有辦法去形容.
因為我們不知道.
世界你用再多的言詞沒有辦法去形容.
這就是上帝所允許的.
所以我們要.
將來在這個榮耀的國度裡面身體得俗.
這是保羅在這裡告訴我們說我們所盼望的.
我們所盼望的.
雖然我們看不見.
雖然我們眼睛看不見.
但是我們信心的眼睛看見Amen.
那怎麼辦呢.
剛纔我們已經講了.
如果你今天犯罪.
如果這個苦難是因為犯罪的緣故.
犯罪的就應該怎麼樣.
認罪.
大衛的故事我們都很清楚.
那天我們早上六點鐘的宣讀聖經的時候.
我在聽到大衛的故事的時候.
其實很熟的故事.
但是我在聽的當中又有一些.
臉色.
我就想到我們最近在讀大衛的詩篇.
想想大衛是多麼一個讓我們覺得很羨慕.
有沒有.
一天禱告三次.
一天讚美幾次.
清晨的他就起來禱告神.
他的生命真的是值得我們羨慕.
可是這樣子的一個人怎麼會犯罪呢.
而且犯這種罪呢.
他自己的氣他看到這個八十八.
他心中就動了一念.
然後怎麼樣把他的先生.
就調回來.
他犯罪之後.
他為了掩蓋自己犯罪.
他把他的先生從戰場的調回來.

$^{761}$希望讓他的先生跟他太太同房.
你懂這意思嗎.
結果這一個先生.
他的義正是顯出大衛的不義.
他說前些的戰士正在那邊打仗.
我怎麼能夠這樣子呢.
他不要.
結果.
他沒有辦法.
他只有最後.
所以我聽的時候就聽得說.
你就可以看見那個人性的那種墮落.
人性的那種敗壞.
大衛就找了他的.
這個派了將軍約壓告訴他說他前線的時候.
你就怎麼樣子的.
把他派到最危險的地方去.
結果烏里亞到了最危險的地方.
就被點殺了.
被殺了之後接下來做什麼事情.
他就把烏里亞取進來了.
不不不 把八十八取進來了.
剛纔準備要睡的人突然醒過來了.
弟兄姊妹.
我們有比大衛更好嗎.
我們真是都是在罪孽當中的.
但是大衛經歷了.
大衛並不是像掃羅一樣.
掃羅犯罪不願意承認自己的罪.
但是他也不願意承認自己的罪.
他也不願意承認自己的罪.
但是大衛當他犯罪之後.
先知拿丹對他說話.
則指出他來.
他先講了一個例子.
講完之後.
然後他說那人就是你.
哇 講完之後大衛應該怎麼樣.
如果是你的話你會怎麼樣.
來人啊.

$^{801}$把他推出去斬了.
是不是.
大衛沒有.
大衛畢竟是一個敬畏神的人.
他聽到了之後.
立刻說我得罪耶和華了.
他以為他神不知鬼不覺.
他以為他借刀殺人.
他以為他所做的這些天衣無縫.
沒有想到神都知道.
先知來告訴他.
是他做的.
所以在詩篇51篇.
他這樣子說.
神啊.
求你為我造清潔的心.
使我裡面重新有正直的靈.
詩篇51篇就是大衛的懺悔詩.
整個詩篇裡面.
代表著一個.
真誠悔改在神面前的一個表達.
親愛的聽眾妹妹.
不怕我們犯罪.
只怕我們不認罪.
沒有一個人不犯罪.
如果有人說他可以不犯罪.
他就是在說謊.
但是.
當我們犯罪的時候.
我們知道有一位義者.
耶穌基督.
當我們向他承認我們的罪.
神是信實的.
是公義的.
必要赦免我們的罪.
洗淨我們一切的不義.
Amen.
第二個.
在這裡.
羅保書裡面.

$^{841}$25節說.
他說我們如果盼望那所不見的.
就必忍耐等候.
忍耐等候.
雅各書也告訴我們要忍耐.
我們起來讀好不好.
雅各書第五章第七節第八節.
預備請.
弟兄們.
雅各書裡面講忍耐講了很多.
雅各也說.
他說你們要忍耐.
但是雅各在這個地方提到.
也是跟主再來有關.
他說忍耐到什麼時候.
忍耐直到主來.
你說我還忍到什麼時候呢.
我結婚二十年受苦夠了.
忍耐.
跟你旁邊說.
忍耐.
忍耐真不容易.
是不是.
忍耐這個字.
在這個地方.
這個希臘文有兩個字.
沒有時間清楚的解釋.
簡單來說.
這忍耐呢.
不是消極的在那邊等.
好像我們講忍耐就是很消極.
忍耐.
我在那邊消極的等.
等候.
不是.
他這裡告訴我們說.
忍耐的等候是什麼呢.
他的忍耐等候.
就好像是承擔一個.
好像是一個舉重的選手.

$^{881}$承擔著那個舉重的時候.
那個重量很大.
這個一百多少.
我沒舉過重.
一百八.
還是一百公斤.
一百公斤.
是不是.
那麼.
你看那個舉重的人怎麼樣.
他是這樣子嗎.
不是.
他舉重是怎麼樣.
他兩腿在重壓之下.
他要把他站穩.
然後怎麼樣.
所以那是一個積極的.
他不是消極的在那邊被動等候.
他是積極的怎麼樣.
在重壓之下.
在困難之下.
他怎麼樣.
他能夠去堅忍.
他去忍耐到底.
這就是這裡所講的忍耐.
他說你們要忍耐.
兼顧你們的心.
所以有很多的學者把這個地方.
翻譯成說.
堅忍更好.
堅忍.
又有一個翻譯叫做持守.
把它對.
我們忍耐的時候常常會有.
這個.
在這個希臘文裡面講的兩種的忍耐.
第一個就是對人的忍耐.
第二個是對環境的忍耐.
對人的忍耐在婚姻當中需不需要忍耐.
管教孩子需不需要忍耐.

$^{921}$真巴不得給他喫歐羅菲.
早期在臺灣我們給豬喫的叫歐羅菲.
就是讓它長得快.
現在就是打生長激素.
這幾天帶我小外孫.
就覺得真是.
這把年紀帶孩子不太容易.
喫個東西.
他也沒耐性.
餵了一口.
馬上再喂.
就他哦哦就表示說你要再喂.
心想這什麼時候才喫飽呢.
喂個東西就搞半天了.
我們等這個孩子長大要等多久.
要等多久.
18歲.
等18年.
容不容易啊.
我跟你講你看你搖頭的都是那些有小小孩的.
我都在想我三個女兒小的時候.
一直到現在我都不知道怎麼過來的.
我那個日子真的是不知道怎麼過來的.
是不是師母.
師母更不容易啊.
全天候.
全天候都要照顧孩子.
親愛的同事們.
人對人.
我們今天生活當中有很多這些人的需要忍耐.
對事情對環境.
有時候這些大環境不是我們能夠決定的是不是.
人家打仗你有辦法我也不是國防部長.
我也沒辦法去決定.
好他們就打了打了之後我們怎麼辦呢.
我們大家只有等啊等吧等到什麼時候.
他說很快.
就下個禮拜就有結果了.
就一次兩次.
最後你就算了不等了.

$^{961}$不管你了.
日子照過是不是.
那當然了今天你能夠負擔起高油價.
你能負擔起這些高消費.
我們這些人都沒有關係是不是.
很多人沒有辦法.
是不是他可能沒有辦法等下去了.
親愛的同事們.
聖經裡面告訴我們要忍耐等候.
跟你旁邊說要忍耐等候.
龍夫等待龍產的出來.
他需要等待.
他不能說我今天撒一盆水下去.
馬上明天就可以果子出來.
有辦法這樣子嗎.
沒辦法.
所以我們要忍耐等候.
聖經告訴我們說像約伯一樣的忍耐.
約伯經歷了家破人亡.
經歷到身體都受到很大的創傷是不是.
全身長包.
這樣子一個情況.
請問如果臨到你的身上.
你要忍耐嗎.
你能忍耐嗎.
聖經告訴我們說.
就如同眾仙姿的榜樣.
特別講到約伯.
約伯他的忍耐.
直到主給他的結局.
明顯主是滿心憐憫大有恩賜.
約伯的忍耐.
約伯的約伯記是許多在苦難當中的人.
最.
可以說是最喜歡去看的一卷書.
中子妹.
如果你在苦難中.
鼓勵你好好把約伯記.
把它好好讀一遍.
約伯忍耐.

$^{1001}$但是約伯的忍耐不是消極的等候.
約伯的忍耐.
他是積極地在盼望.
好 我沒有時間講了.
最後我們要看.
我們中間有受苦的.
他就該禱告.
有喜樂的.
他就該歌頌.
羅馬書十二章十二節我們一起來讀.
在指望中要喜樂.
在患難中要忍耐.
禱告要恆且.
在這個地方告訴我們說.
我們中間如果有受苦的應該要怎麼樣.
要禱告.
跟你旁邊說要禱告.
你說牧師又來了又禱告.
三句話不離本行又禱告.
那不請問你.
聖經上是不是這樣說.
是不是.
聖經上是不是說.
你們中間有受苦的要怎麼樣.
是誰說的.
誰說的.
誰說的嘛.
雅各蘇說的.
神的話說的.
是不是.
你們中間有受苦的就應該要怎麼樣.
禱告.
保羅也說.
在患難中要忍耐.
接下來就說什麼.
禱告要恆且.
其實雅各所說的受苦跟禱告.
在原文裡面都是現在是什麼意思呢.
表明這是一個持續的狀態.
表明這不是一次性的.

$^{1041}$表明不是這是短暫的受苦.
表明這是一個長時間的受苦.
所以這個禱告必須要恆且.
持續不斷地禱告.
不住地禱告.
常常地禱告.
我們常常的問題是禱告.
我們禱告好像是做個文章.
有一次我在一個教會的網站.
因為別人叫我幫他們介紹.
介紹這個教會.
我就去看了這個網站.
覺得這個教會.
我就介紹給.
有弟兄姊妹問我.
我就介紹給他們.
我就上去看了他們祖宗崇拜.
就剛好看到他們祖宗崇拜的時候.
師會在那邊上臺.
像我們這樣子.
我們是牧者們上臺禱告.
師會主席就在那邊禱告.
因為從攝像機照下來是很清楚的.
所以我們看的人是很清楚地看見.
他是在對著一個禱告文.
寫好的禱告文.
他在禱告.
懂我意思嗎.
就是他說.
親愛的主耶穌啊.
我們現在在.
就那個文字非常的典雅.
你一般如果眼睛閉著才禱告.
你是沒有辦法講出這麼典雅的文字的.
你懂我意思嗎.
哎呀 我看了就很感慨.
其實今天有許多基督徒.
平常沒有在禱告.
到了公開場合.
可能上臺了.

$^{1081}$需要禱告的時候.
他不知道該怎麼禱告.
丟著沒.
我們的禱告.
是不是發自內心的禱告.
還是做一篇文章.
形式化的禱告呢.
所以聖經裡面保羅也好.
彼得也好.
提醒我們.
我們要恆切地禱告.
在受苦的時候要禱告.
要持久地禱告.
有一位.
寫了一個很著名的一本書.
不知道你有沒有看過.
天路歷程有看過嗎.
有看過嗎.
我瞭解一下好不好.
這本很老的書.
有看過天路歷程的舉手.
請放下.
沒有看過的.
鼓勵你一定要看.
這個可以說是經典的.
可以說是一個屬靈書籍.
天路歷程這本書.
這個作者叫做本人約翰.
有的把它翻譯叫做約翰·班揚.
臺灣是把它翻成本人約翰.
其實我不知道為什麼翻本人約翰.
John是約翰.
班揚本人.
Anyway 我們已經習慣了.
本人約翰.
他真的是神使用一個人.
他真的是刻意地在塑造這一個人.
本人約翰.
他曾經經歷到有12年的被關在監牢的生活.
這是他一生當中最黑暗的這個時期.

$^{1121}$可是也是他操練禱告最深的一個屬靈的實驗室.
當時候他被送到監獄裡面去的時候.
跟妻子跟女兒分開.
他自己形容說.
他說與妻兒分離.
就像把肉從骨頭上撕扯下來一樣的痛苦.
他沒有辦法 這12年當中在監獄裡面.
每一天他只有把自己的家人交託在上帝的手中.
不僅僅如此.
在當時候17世紀的英國.
你不服從國教的法規.
私自講道這是違法的重罪.
所以本人約翰他隨時要面臨被判處死刑.
或者是被流放.
在這種的情況之下.
他每一天都在面臨到死亡的掙扎.
他在許多的痛苦的中間.
他能做的是什麼.
他唯一能做的就只有禱告.
雖然沒有一個記載說他到底花多少時間禱告.
可是根據當時候跟他一起關在監牢裡面的.
有一些清教徒的弟兄.
還有一些的獄友.
他們的描述.
本人約翰他說他當時候每一天都要在牢獄裡面做工.
他一面在那邊做那個長筒絲襪.
就是要做那個絲襪 女人的絲襪.
他在那邊一面在編織 一面在做的時候.
旁邊的人就看著他的嘴巴一直不停地在微微地.
在那邊咬動 就這樣子.
他在做什麼 他在禱告.
他除了這個之外.
他每一天早上清晨的時候跟晚上.
都是別人在睡覺的時候.
別人還沒起來的時候.
他自己一個人跪在牢房的那個溼溼潮潮的地上面.
跪在那邊獨自的幾個小時在那個地方禱告.
清晨夜晚的禱告.
每一天時時刻刻不住地.
默默地在那邊禱告.

$^{1161}$苦難使得本人約翰塑造了他那個禱告的生命.
親愛的弟兄姊妹.
苦難可以囚禁一個人的身體.
但是卻沒有辦法囚禁一個雙膝跪下向主禱告的人.
Amen.
求主幫助我們.
也許你現在正在苦難的中間.
但是不要把苦難當作是一個怨天尤人的時間.
把你的怨天尤人變成是一個交託禱告.
把你的怨天尤人甚至變成是一個讚美歌頌.
Amen.
把你的怨天尤人變成是一個忍耐等候.
因為知道主接著這些苦難.
要使我們的生命越來越像誰.
耶穌基督.
那就是將來的榮耀.
那個不但只是身體被改變的.
而是整個人你可以看見.
經過在地上的磨練之後.
這個人就已經到主面前的時候.
是像耶穌基督的一個人.
讓這個成為我們的盼望.
Amen.
\newpage



\section{約翰福音 4:31-42}
\label{sec:032269}
\textbf{遵行神旨,作成主工}~~~~駱柱成牧師
\newline
\newline
連結: \href{http://www.church.com.hk/acms/content.asp?site=cdc\&op=show\&type=product\&code=032269\&layout=sermon}{\texttt{032269}} ~~~~ 語音日期: 2026-06-07
\newline
\newline
\hyperref[sec:032268]{\small{< < < PREV SERMON < < <}}
~
\hyperref[sec:index_chronic]{\small{[返順時目]}}
~
\hyperref[sec:032270]{\small{> > > NEXT SERMON > > >}}
\newline
\newline
約翰福音 4:31-42
\newline
\begin{longtable}{cl}
\hline
\hline
章節 & 經文 (和合本修訂版)\\
\hline
4:31 & \begin{tabularx}{0.7\textwidth}{X} 就在這個時候，門徒求耶穌說：「拉比，請吃吧。」 \end{tabularx} \\ \\ \relax
4:32 & \begin{tabularx}{0.7\textwidth}{X} 耶穌對他們說：「我有食物吃，是你們不知道的。」 \end{tabularx} \\ \\ \relax
4:33 & \begin{tabularx}{0.7\textwidth}{X} 門徒就彼此說：「難道有人拿甚麼給他吃了嗎？」 \end{tabularx} \\ \\ \relax
4:34 & \begin{tabularx}{0.7\textwidth}{X} 耶穌對他們說：「我的食物就是要遵行差我來那位的旨意，完成他的工作。 \end{tabularx} \\ \\ \relax
4:35 & \begin{tabularx}{0.7\textwidth}{X} 你們不是說『到收割的時候還有四個月』嗎？我告訴你們，舉目向田觀看，莊稼熟了，可以收割了。 \end{tabularx} \\ \\ \relax
4:36 & \begin{tabularx}{0.7\textwidth}{X} 收割的人已經得工錢，為永生儲存五穀，使撒種的和收割的一同快樂。 \end{tabularx} \\ \\ \relax
4:37 & \begin{tabularx}{0.7\textwidth}{X} 『那人撒種，這人收割』，這話可見是真的。 \end{tabularx} \\ \\ \relax
4:38 & \begin{tabularx}{0.7\textwidth}{X} 我差你們去收你們所沒有辛勞的；別人辛勞，你們享受他們辛勞的成果。」 \end{tabularx} \\ \\ \relax
4:39 & \begin{tabularx}{0.7\textwidth}{X} 那城裡有好些撒瑪利亞人信了耶穌，因為那婦人作見證，說：「他把我素來所做的一切事都說了出來。」 \end{tabularx} \\ \\ \relax
4:40 & \begin{tabularx}{0.7\textwidth}{X} 於是撒瑪利亞人來見耶穌，求他在他們那裡住下，他就在那裡住了兩天。 \end{tabularx} \\ \\ \relax
4:41 & \begin{tabularx}{0.7\textwidth}{X} 因為耶穌的話，信的人就更多了。 \end{tabularx} \\ \\ \relax
4:42 & \begin{tabularx}{0.7\textwidth}{X} 他們對那婦人說：「現在我們信，不再是因為你的話，而是我們親自聽見了，知道這人真是世界的救主。」 \end{tabularx} \\ \\
[1ex]
\hline
\hline
\end{longtable}
$^{1}$今日我正道的經文是.
《約翰福音》四章三十一到四十二節.
上兩週我已經講了兩堂關於薩瑪利亞婦人的.
今日我們繼續薩瑪利亞婦人的事跡和成事的事.
今日的主題是「遵行神的旨意和造成神的功」.
我們在這個地上每個人都要謀生.
廣東人更實際「餵兩餐就乜都肯濟」.
「兩餐」就是指我們一天的兩餐.
後來「兩餐」這個字就變成「謀生」的意思.
就是為了謀生.
古時候古代的人就是吃兩餐.
因為早上去耕田 晚上回來再吃一餐.
「餵兩餐」這是一般的解釋.
「兩餐」不是一種享受.
「兩餐」只是人生存最基本的需要.
「餵兩餐」所以很多人就頻頻撲撲地「餵兩餐」.
奔奔波波就「搵兩餐」.
所以人的一生就是忙忙碌碌.
就是為了生活 為了戶口.
為了我的家人 為了我自己的口福.
如果有人跟我們說「有兩餐」.
那已經是很低調的一種說法.
「有兩餐」就代表他過得比較好的.
可以說過得還是過得去的.
很多時候我們中午也好 夜晚也好.
都會問你吃了沒有.
所以吃是很重要的.
但是今天的社會不擔心沒得吃.
大家關心的是去哪裡吃 吃什麼.
還有去哪一家酒樓吃.
很多時候都在考慮這些.
甚至有些人要吃那些很好吃的.
可能就是約了朋友就去.
朋友的電話一來就說「海鮮的季節到了」.
「你快點飛過來吧」.
就可能是當天或者當星期結束了工作.
星期五晚出發 星期一再回來吃大餐.
人的生命就是為了生活.
生活又為了生命.
「唔信主」就是這種感覺.

$^{41}$生命為生活.
生活又好像是為了延續肉身的生命.
但是如果人生只是為了這些的話.
最後發現人一生就是為了謀生.
一步步就去到一個終點 死亡那裡.
最後才發覺很多忙忙碌碌做的事.
只不過是沒什麼意義的.
或者不是那麼重要的.
很多不必要做的事到最後都做了.
但是真正要做的事反而沒有去做.
這樣是很可惜的.
我們活在這個世上到底是為了什麼呢?.
很多人將生活的需要看成最重要的.
但是可以說是將神的旨意放在一旁.
將神賜給我們的恩賜和責任也放在一旁.
將屬靈的,心靈的,更高尚的那種需要也放在一旁.
反而將人最基本的.
可以說是延續肉身的生命.
當成為人生唯一的目標.
這樣就出問題了.
所以耶穌基督提醒我們.
人活著不是單靠食物.
而是靠神口裡所出一切的話.
我們需要食物是沒錯的.
也需要工作也是沒錯的.
人不只是為了食物最基本的東西而活著.
因為專行神的旨意.
得著屬靈的需要的那種滿足和喜樂.
人心靈最需要的東西.
可以說是遠遠超過物質表面上的東西.
我們今週開始學傳道書.
周長老稍後會跟我們說.
傳統上傳道書被認為是所羅門晚年的反思.
所羅門有智慧,有財富,也有權力.
享盡一切美食榮華富貴.
但是從日光之下的人生得出一個結論.
就是虛空的虛空.
就是如此虛空.
吃飽了第二餐又來.
都是如此虛空.

$^{81}$凡事都是虛空.
因為沒有神的人生.
他擁有最多,他吃得最多.
都依然是虛空.
傳道書不是一本消極的書.
所羅門的觀點是.
人莫強如吃漢.
且在勞碌中享福.
我看這是出於神的手.
工作沒有什麼錯.
享受神所賜的福.
也是很好,也不是罪.
問題是出於在哪裡.
如果人一生就是為了吃.
為了勞碌.
為了積蓄一些所謂的金錢.
享樂,最終就是暴風.
好像捉空氣一樣,捉不住.
因此經過一生的探索.
所羅門最後得出一個結論.
總之就是你要敬畏神.
守祂的誡命,聽祂的話.
這個就是我們人要做的一種本分.
你要聽神的話,你要敬畏神.
這個才是最重要的.
真正人生的意義不是在乎你擁有多少.
而是你有沒有活在神祂的旨意.
所以我們可以說.
只是為現世,為現在.
為兩餐而活,是很虛空的.
但你要為永恆而活.
這樣你才是有福的.
你可以積到永恆.
但是人憑著自己.
你是沒辦法脫離這些物質的虛空.
你受很多的限制.
耶穌基督就來到地上拯救我們.
保羅這樣說過.
基督既為我們受了咒詛.
祂救贖出我們.

$^{121}$脫離那種律法的咒詛.
將我們脫離這些咒詛.
在這個世界上的咒詛.
主耶穌親自來到日光之下.
在這個充滿虛空的世界.
承擔我們的罪惡和痛苦.
將我們從虛空救贖出來.
使我們能夠從永恆的角度.
永遠的角度.
去思想地上的人生是怎樣的.
祂讓我們看得更長遠.
可以說是永恆的.
不是有限的目光.
所以保羅就說.
所以我親愛的弟兄們.
你們務要堅固.
不可以動搖.
且要竭力多做主公.
因為知道你們的牢籠.
要多做主公的時候.
你們的牢籠是在主的裡面.
不會是逃園的.
不會是那麼虛空的.
一下就過了.
什麼都沒有.
不是.
不會是逃園的.
傳道書說虛空.
保羅在主的裡面.
勞苦不是逃園的.
因為基督使我們的人生.
是有永恆的意義的.
我們在主的裡面.
是可以傳到永遠的.
所以我們今天做在主身上的.
在主裡面的.
專行神的話的.
是可以有永恆價值的.
今天一個人如果能夠脫離.
物質的那種捆綁.

$^{161}$我們被很多的金錢.
很多的東西綁住.
我們的心很難投入.
神的永恆那種旨意那裡.
因為太多東西綁住了.
但如果你可以做得到的時候.
全心去投入主耶穌.
祂就會出一個美好的榜樣.
當門徒買食物回來的時候.
祂對他們說.
拉比.
請吃.
這個就是食物了.
請吃.
如果我們走了很長的路.
一定會說.
我很餓了.
快點給我.
因為我有食物的需要.
但耶穌基督呢.
耶穌基督沒有這樣說.
祂反而說.
我有食物.
是你們不知道的.
門徒不明白這句話的意思.
就在那裡討論.
思考去討論.
莫非有誰給了他食物.
他不餓的.
但他們想的是地上的.
短暫的.
耶穌基督所思所想.
祂的心所想的是天父的旨意.
耶穌基督祂想的是天父的旨意.
門徒想的是食物.
我記得有一個詞叫廢寢忘食.
有一次孔子.
有人問孔子的學生.
子路.
孔子是一個.

$^{201}$到底孔子是一個怎樣的人.
子路就說.
一時都想不到.
他很想說.
但他說不出.
後來孔子知道了.
就教他這樣說.
發奮忘食.
樂而忘憂.
不知老之將至.
就是說.
孔子就是一個很專心去追求學問的人.
經常因為追求那種學問.
去思想的時候.
他就不記得吃飯了.
甚至忘記自己的優秀.
很多的優秀.
甚至忘記自己已經老了.
這句話的意思就是這樣.
上個月.
我去區會.
有中國的傳道人跟我一起去學習.
我們都要去考區會的license.
我問他最近怎樣.
他就說.
說到前幾天.
基本上從早到晚一直在修改整理.
所以.
他的太太為他準備午餐的時候.
他說等一會再吃.
一等.
低頭一抬起的時候.
午餐就變成晚餐了.
有時晚餐就變成宵夜了.
所以他又是很專注.
我就說.
你好像孔子.
他說哪方面.
我說廢寢忘食.
你就像孔子一樣.

$^{241}$你做得很專注.
但孔子是為追求學問.
耶穌是為專行天父的旨意.
所以耶穌進一步解釋給門徒聽.
我的食物就是專行差我來者的旨意.
作成祂的功.
我就是要作成父神的功.
父神差我來的.
我要造成祂.
我專行的就是專行父神的旨意.
你會看到大四章.
耶穌向撒瑪利亞婦人要水喝.
又差派門徒去城裡買東西吃.
你可以想到耶穌基督.
祂真是脾渴的.
真是肚餓的.
但是祂就是.
當祂傳講真理.
見證福音.
尋找這些失喪靈魂的時候.
和這個失喪的撒瑪利亞婦人.
講道的時候.
祂就將喝水.
吃東西.
都可以說擺在一邊了.
不是祂不肚餓.
不是祂不脾渴.
祂是真是.
一心投入在神的旨意當中.
所以你會看到整章.
再沒有提及耶穌基督.
怎樣喝水.
怎樣吃東西.
其實這個不是第一次的.
你會看其他福音書.
特別是《路加福音》.
你會看到.
耶穌基督吃飯的時間都沒有.
而且祂的生活.
是靠一些婦人供養祂的.

$^{281}$無論是祂還是門徒.
主耶穌自己也說過.
「湖泥有洞,天空的飛鳥有窩」.
人質連枕頭的地方都沒有.
可見祂在世上是很簡樸的.
祂曾經向婦人說.
請你給我喝水.
也吩咐門徒.
吩咐門徒去買東西吃.
但當祂傳講真理的時候.
祂就將這些放在一邊.
祂是以神的工作為重.
以父神的心意為中心.
對耶穌來說.
專行神的旨意.
造成神的功.
比身體上的需要都更加重要.
所以祂就宣告.
我的食物就是專行差我女兒的旨意.
造成祂的功.
第二點.
接著祂又說.
「你們豈不說到收割的時候.
還有四個月嗎?.
我告訴你們.
莊稼已經熟了.
可以收割了」.
主耶穌說這句話的時候.
門徒剛剛從城裡買東西吃回來.
而另一方面.
撒瑪利亞婦人已經向城裡的人.
說了這個福音.
所以很多撒瑪利亞人就從城裡出來.
因為他們聽了這個婦人的話.
他們要去到耶穌那裡.
所以很可能耶穌就在這個時候.
指著這些快要來到的撒瑪利亞人說.
莊稼已經熟了.
可以收割了.
主耶穌要那些門徒看到的是.

$^{321}$福音在撒瑪利亞人當中已經有收成了.
已經近在眼前了.
祂將門徒的那種目光.
食物的那種目光.
轉向那個和祥.
神的和祥那裡.
從自己的需要.
轉向神的工作那裡.
但是主提醒他們神的工作.
不只是收割.
也殺種.
所以祂又說.
那些人殺種.
那些人是收割.
這句話是真的.
今天我們傳福音也是一樣.
有人殺種.
有人囂觀.
有人收割.
但是一切的成長都是在神那裡.
那個果效都是在神的手中.
正如保羅自己說.
我栽種了.
阿波羅囂觀了.
唯有神才可以使他們生長.
所以當我們傳福音.
有果效的時候.
你要記得.
不是我們可以使人得救.
是神能夠使人得救.
不是我們可以赦免人的罪.
是耶穌基督十字架的救恩成就.
這份救恩成就的時候才可以赦免人的罪.
不是我們可以改變人的心靈.
是聖靈感動人的時候.
他才會信主.
歸向主耶穌基督.
我們不是.
我們只不過是神手中的一個器皿.
我們要做的就是忠心去傳揚福音.

$^{361}$所以殺種的人.
不要灰心.
就算你當時看不到一個果效也好.
你都要繼續去忠心.
收割的人.
你有收成了.
你一傳一傳.
別人就相信了.
你也不要驕傲.
因為我們在主的面前.
我們知道我們只是做了當做的.
我們可以歡喜快樂.
但是要知道.
當我們有這份收成的時候.
有人信主的時候.
不只是我們的功勞.
早就有很多人殺了種.
早就有很多人向他們傳過福音.
他們已經默默地附上了.
或者我們為我們的家人禱告.
也有跟他們去探訪.
也有一次次跟他們傳福音.
但是他們一直都沒有信.
但是當神一做公的時候.
差派另外一位去的時候.
每個人就有收割了.
每個人一傳就相信了.
所以這一切都是傳言是神自己的差遣和作為.
我們教會每年都會舉辦不同的福音外展活動.
昨天有個包種.
我們記得.
如果是六七月份的時候.
如果有小朋友.
我們記得袁牧師又有個圍棋.
到中秋節前後有個月餅.
所以在這些活動當中.
我們的牧師袁牧師.
經常都會有間斷的福音分享.
昨天我們就舉辦了包種的活動.
但是昨天有沒有人相信主.

$^{401}$有沒有當場就有人決志信主.
沒有.
但是不代表我們沒有殺了那個種.
因為有人殺種.
亦都有人收割.
將來如果有人因此信主.
很可能就是因為這一次活動的殺種.
亦都很可能是因為弟兄姊妹長期的關懷.
禱告和傳福音.
積累的一個結果.
到了神所定的時間.
神就讓我們看到收割的果效.
所以平時基督徒怎樣去傳福音.
怎樣去活出見證的時候.
很多時候就影響了這些福音的外展.
是否有那個收成.
不只是靠那個外展活動.
亦都不是靠一些的報道.
而是靠平時弟兄姊妹跟人分享福音.
為主做見證.
所以主耶穌就說.
我差你們去.
收你們沒有勞苦的.
人家已經勞苦了.
你們享受他們的勞苦.
所以收割的人要感恩.
殺種的人亦都不需要灰心.
無論是殺種的,收割的.
都在參與著神的工作.
享受著神的恩典和喜樂.
我們的責任是什麼?.
就是忠心去傳揚.
果效就在神那裡.
中國的福音的傳播亦都不是一朝一夕的.
我們知道我們中國的文化.
不是一直都是信耶穌的.
在唐朝的時候有景教.
雖然景教的信仰不純正.
但是當時唐朝景教已經傳入中國了.
元朝亦都有天主教去傳揚.

$^{441}$再到明朝的時候.
有一些宣教士利瑪竇.
是耶穌會差派去的.
亦都向中國人去傳福音.
他不單止傳福音.
亦都將當時的文化.
和豐富的科學技術和知識.
都帶到中國那裡.
讓那些官和學者可以學習到很多新的知識.
所以福音亦都在那個時候進入到這裡.
近代我們現在一般所說的基督教.
從馬登路德開始脫離了天主教.
一般稱我們為新教.
近代的新教在中國宣教的工作.
通常就是在1807年.
有一個叫馬禮遜的開始去計算.
馬禮遜進入中國的時候.
花了七年的時間.
只有一個基督徒真正不單止信.
而且是仇使.
那個人叫做蔡高.
是七年的奴婦.
只有一個「果」字.
但是今天回頭看.
這七年是不是桃源呢?.
不是的.
後來有一個牧師出來.
叫做梁發.
廣東佛山人.
所以就有新教的傳道人.
是一代一代地去傳開的.
今天我們看到的就是.
神祂的工作往往不是一天就搞定的.
有人殺種.
有人削冠.
亦都有人收割.
正如耶穌說的.
那個人收割.
這個人殺種.
這句話是真的.

$^{481}$所以我們看到.
修成的時候.
往往就是建立了.
很多前面殺種的人的耕耘.
我們才有那個修割.
在見和祥的時候.
他創辦的.
因為宣信博士原先是長老會的.
後來呢.
他就成為近代最傑出的.
一個宣教推動者之一.
當他看到世界的民族.
有很多的民族.
很多的地方都未聽過福音的時候.
他就呼籲教會一同去回應主的大使命.
要將福音傳遍地極.
不只是看到我們這片小小的和祥.
不是看到我們當地.
不是看到我們同膚色的人.
也都不是看到英國這個地方.
也都不是看到美國這個地方.
你要看到的是全世界.
所以他創立了宣道會.
不只是為了建立一個新的宗派.
而是為了推動普世的宣教.
在他的影響之下.
很多的基督徒就出來回應.
去不同的地方傳揚福音.
直到今天宣道會.
依然承接這個意象.
將福音帶給萬民.
所以你會看到宣信.
他推動的是廣大的宣教時光.
因為他看到真正的和祥是何等的廣闊.
當主耶穌看到撒瑪尼亞人的時候.
他看到的是很廣闊的和祥.
宣信他看到的是萬國萬民.
他看到的也是和主一樣的和祥.
他們就是等待收割,莊稼.
所以主就說.

$^{521}$莊稼已經熟了.
已經可以收割了.
所以今天主也同樣呼召我們.
你要看.
你要看得更遠.
不只是看到自己的需要.
也不是看到自己身邊的一些人.
要看到神他愛靈魂.
他神的國的和祥是很廣闊的.
五月份的時候.
宣道會有一個會叫做「華聯會」.
它不是行政方面的.
它是一個很大的團契.
不同教會服侍的人.
一同聚會.
一同交流.
彼此連結.
互相支持.
一起推動福音的時光.
在這些團契當中.
在華聯會當中.
我看到有兩個是白色皮膚的.
兩位都是女宣教士.
當我和她聊天的時候.
我打算要用我不寒不淡的英文跟她說.
但她一過來就說.
「你好,我是在哪裡服侍的?」.
「我現在在北京工作」.
她說「我在北京服侍」.
她說得比我更標準.
我不懂得捲舌.
她懂得.
另外一個更厲害.
「什麼回事呢?」.
我以為她只不過跟另外一個一樣.
只懂得說普通話.
她說「你好,我現在不是在北京」.
「我是在雲南」.
後來我知道她.
原來她在香港住了幾年.

$^{561}$後來就去雲南.
去廣西傳福音.
所以她說「你好厲害」.
「你很厲害」.
所以.
他們就是這些不同膚色的.
沒有那種民族的偏見.
看到的是主的禍場.
不是你的.
不是我的.
不只是我自己的小地方.
她是做什麼工作的呢?.
她是一個音樂老師.
所以她就是教中國的.
廣西,雲南一帶的年輕人.
學吹笛,小提琴,彈鋼琴.
辦了一個音樂培訓中心.
一方面教音樂.
其他時間會讀聖詩,聖經.
教他們聖經.
所以.
她就是通過這些方法.
將福音帶到那些沒有人願意去的.
很少人去的.
雖然已經有一些宣教士有足跡.
但她就在那裡默默耕耘.
她做了一段時間之後.
她的培訓中心在疫情期間被查封了.
她就將樂器送了給別人.
繼續去找其他的服侍.
不單止是中國.
去到緬甸.
如果是我們呢.
可能我們就會說.
既然都會廣東話.
不如在香港深圳就好了.
這麼方便.
吃東西又方便又好吃.
舒服點就好了.
又是做主公.

$^{601}$但她沒有選擇自己喜歡的地方.
她是願意信服神.
神要她去的地方.
為什麼呢?.
因為她看到的是神的禍場.
她看到的是很多人還沒聽到福音.
她看到的是靈魂需要基督.
她看到的是神國正在等候修葛.
要更多的人去信主.
也提醒我們.
事奉主最重要的是什麼?.
不是我想在哪個崗位.
我想做什麼.
而是主要我做什麼.
不是我的計劃可不可以成功.
而是我有沒有忠心去遵行神的旨意.
真正的事奉.
不是說我決定做什麼就做什麼.
而是我先簽了名.
簽了一份不可以更改的合同.
我先簽了名.
神你自己去填.
你填的內容是什麼.
我就去做什麼.
這就是信服.
這就是事奉.
先要將自己交給神.
然後再跟主說.
主啊.
你是我的主.
你是我的神.
你要我去哪裡.
我就去哪裡.
你要我在貧窮的地方.
我就去貧窮的地方.
去市中心.
我就去市中心.
你要我在職場.
在各個領域當中去見證你.
我就在各個領域當中去見證你.

$^{641}$不是我.
不是體貼我.
所以主差我們去修葛.
我們要信服.
主差我們要殺種.
我們也要信服.
因為真正的奉獻.
不是要我們成全我們的心意.
而是成全神祂自己的心意.
耶穌說.
莊稼熟了.
可以修葛了.
不是體貼我.
是體貼神祂的心意.
神要我們去修葛.
祂給我們的使命.
給你的使命是什麼.
或者有人就說.
我年紀大了.
不要了.
走不得遠的了.
你們宣教的年輕人去吧.
那麼.
說得過去的.
但是也都.
因為真的年紀大了.
沒有足夠的精力.
走得太遠.
但是現在很多的工作.
不一定出國的.
不一定要走很長的水路的.
因為神祂的作為是很奇妙的.
今天來自不同的地方的人.
有難民.
有移民的.
有留學生.
也有那些過來這邊工作的.
他們都來到這一個地方.
可以說是洛杉磯.
是他們一個很聚腳的地方.

$^{681}$所以.
他們不單止是我們的靈社.
也是神將他們放在我們面前.
成為我們福音的一個對象.
不一定走得很遠.
他們已經來了.
神的作為很奇妙.
很多人剛剛來到美國.
他們對環境陌生.
文化也都不是很熟.
他們在這個時候.
特別是前兩三年.
他們特別需要幫助.
在這個時候.
他們需要幫助的時候.
往往也是最需要聽福音的時候.
因為每一個人都需要耶穌基督救主.
所以這個時候.
就是我們關懷他們的時候.
幫助他們的時候.
是我們傳福音最好的時候.
可能就是我們平時去吃飯.
餐館的服務人.
是展發的服務人.
就可能是你住的地方.
新來的一些人.
我們不一定都是收割的.
可能就是在他們當中去殺種.
我們就是向他們傳.
為主去做見證.
主耶穌說.
收割的人得公教.
積蓄不谷度永生.
叫殺種的和收割的是一同歡樂的.
殺種的和收割的是一同歡樂的.
我們最主要做的就是忠心去做.
當華聯會三天結束的時候.
我就在波蘭那邊搭飛機回來.
一個司機約了車.
他就來搭我了.

$^{721}$車到我的機場.
我就介紹了我是這裡的什麼的.
我們這裡教會有一些福音的聚會.
就跟他說.
後來呢.
我才知道他是一個穆斯林.
他是亞洲那邊的國家走過來的.
他只不過來了幾年.
他娶了一個墨西哥裔的太太.
但是他從未去過教會.
他就問了我一句.
可能他娶的是墨西哥的太太.
所以他又是穆斯林.
他就問了一句.
基督徒可不可以跟不同宗教的人結婚的.
我就跟他說信仰.
說婚姻.
說人生.
最後.
下車的時候我就跟他說.
你有機會就跟太太也好.
去接我的地方是教會.
你就可以去那裡聚會.
也都去你家附近那裡去聚會.
今天美國在美國向穆斯林傳福音.
比我們自己去走到穆斯林的國家傳福音.
是輕鬆很多的.
神已經是將不同的人帶到這裡了.
和尚其實就是我們的眼前.
我記得有一次.
我們實學院的一個教授.
有一次他在退休會.
參加完退休會.
他是退休會的講員.
講完之後.
突然電話就來了.
說穆斯林.
你的太太發生車禍.
現在昏迷在醫院.
於是他一下台.

$^{761}$馬上就什麼都不帶.
就開車走了.
因為他講課講了很久.
接著走山路的時候.
一下就跟別人撞了.
撞了別人.
自己就滾下山.
連車帶人滾下山.
怎樣救他呢.
就直升機飛過來.
將他吊起來.
當救護人員救醒他的時候.
他第一句話是說什麼呢.
他那時候牙都斷了.
但是他還是很清醒.
醫護人員問他.
你怎樣.
他第一句話.
他不是關心自己.
他說你信了耶穌沒有.
如果沒有的.
快快都去信耶穌.
所以他在生死當中.
依然是愛人的靈魂.
珍惜這些機會.
隨時去服侍人去傳講福音.
今天可以說是很少見到的.
所以不要以為以後才有機會去傳.
有時你以為.
有時是我們已經沒有時間了.
有時是對方沒有時間.
有時不是因為真的沒有機會.
而是我們沒有把握那個機會.
可能是我們的親人朋友.
也都是剛剛認識的人.
如果你錯過了那個時機.
一下子就過去了.
沒有了那個機會.
本來他要聽福音.
但是他走了.

$^{801}$所以就是要把握那個時機.
趁著神將人帶到面前的時候.
我們要趁著那個門還沒關的時候.
那個死亡之門還沒關的時候.
莊稼熟的時候.
我們就要去修葛了.
我們要忠心地去殺種.
也都去修葛.
要把握那個機會.
多做主公.
第三點.
撒瑪利亞婦人.
她以前.
撒瑪利亞婦人.
她以前是一個不敢見人的人.
她總是躲起來.
見不得光的可以說.
也都不敢和別人去聊天.
因為她知道她的生活.
真的不光彩.
她也都忍受很多的批評和指責.
被人可能.
背後就暗暗沉沉.
因為她的生活就是.
可以說.
墮落了.
所以她寧願在那個.
最熱的時候去打水.
一個人出來.
但是當她遇見耶穌之後.
一切都改變了.
她從躲起來的一個人.
變成敢去做見證的人.
逃避的一個人.
變成一個向全城.
去傳講神的福音的一個人.
她很勇敢地走出去.
走出去過去那種黑暗.
她對眾人說.
你們來看.

$^{841}$有一個人將我.
數來所行的一切事.
都和我說出來了.
莫非這個就是基督.
莫非這個就是我們一直等的那個基督.
所以她的生命已經是.
發生了很大的改變.
從前她怕人.
但是現在呢.
我不怕.
全城我都不怕.
我就是敢去做見證.
今天我們很多基督徒.
叫她禱告.
哎呀不要啦.
叫她做見證.
也都不多敢.
去傳福音.
哎呀不要啦.
免得別人不喜歡我.
保羅呢.
保羅他是這麼說的.
我不以福音為恥.
這個福音本來是神的大能.
要救一切相信的人.
當聖靈真正做工的時候.
人的生命就開始改變了.
當這個婦人認識耶穌基督的時候.
聖靈就在她生命當中去做工了.
她就發生一個變化了.
從前她是以罪為樂.
現在是以基督為樂.
從前呢.
怕承認主耶穌.
不敢承認自己的過去.
現在勇敢去說出來.
而且為主做見證.
撒瑪利亞人最大的改變.
不是承認自己的罪.
她要高舉耶穌基督.

$^{881}$認罪是.
但是重要的是要高舉耶穌基督.
她不是一直說自己的過去.
不是再停留在過去了.
她要將人的目光帶到基督的身上.
你要將目光帶到基督的身上.
真正的見證不是高舉自己.
而是高舉耶穌基督.
不是再關注過去那種失敗的生活.
是讓人看到.
耶穌基督就在我這個生命.
去做改變的工作.
有一句話.
很有意思的.
請你在這裡住下來.
耶穌基督本來就是要來住在人的當中.
祂道成了肉身.
就是要住在人的中間.
祂是創造天地萬物的主.
祂也是願意與人同住的主.
祂不僅與這些門徒同住.
也與撒瑪利亞婦人同住.
祂歡迎人來認識祂.
也親近祂.
但是撒瑪利亞人邀請耶穌基督住下來.
不是一件很容易的事.
為什麼?.
因為猶太人和撒瑪利亞人.
他們是長期的仇視.
他們雖然同是亞伯拉罕以撒雅各的後裔.
但是因為他們的歷史不同.
他們的民族.
很多的習慣已經不同了.
宗教已經好像.
撒瑪利亞人已經有偏差了.
所以他們有很多的衝突.
雙方一直都在這裡.
你指責我,我指責你.
而且是沒有來往的.
你做你的,我做我的.

$^{921}$接著人的眼光來看.
一個猶太人來到撒瑪利亞的人.
那裡,那個城那裡.
可以說他是不受歡迎的.
但是耶穌基督.
祂在他們當中的時.
來到他們當中的時候.
情況就完全改變了.
這些原先仇視猶太人的撒瑪利亞人.
竟然求耶穌說.
你來和我住.
我們都不適合,你來和我住.
為什麼?.
因為耶穌去到哪裡的時候.
就將改變帶到哪裡.
去改變了.
原本被全城所輕視的撒瑪利亞婦人.
而這個婦人.
為耶穌基督做見證的時候.
又影響了全城的人.
這個婦人很有天份.
在宣教的事上.
撒瑪利亞人開始發現.
他們真正需要的.
不是那種民族之間的對抗.
我看到的是猶太人.
我是撒瑪利亞人.
不是那種對抗.
也不是在宗教那裡爭來爭去.
真正重要的是什麼?.
是耶穌基督.
當耶穌基督被彰顯出來的時候.
很多原先看似不重要的事.
就變得像.
民族的仇恨.
人際的衝突.
過去的那種成見.
都失去了那種重要性.
不再重要了.
當耶穌基督出來的時候.

$^{961}$因為他們發現.
這個城市真正需要的.
就是耶穌基督.
當基督進入人生命的時候.
生命就改變了.
這個婦人的生命改變了.
當耶穌基督進入一個家庭的時候.
那個家庭會改變.
當進入一個城市的時候.
那個城市也都會改變.
保羅說過.
以前我以為與我有益的.
我現在因著基督當作是有損的.
原先是看似重要的事.
不再重要了.
以為這些是重要有益的.
反而不是這樣.
接著又說.
我要以認識主耶穌基督為至寶.
原來耶穌基督才是那個至寶.
一個人沒有基督.
就算他做總統.
他都不一定快樂.
一個人有了基督.
他在各樣的困難當中.
世界的一個改變.
是取決於一個人.
對耶穌基督的認識程度.
去到哪裡.
耶穌基督是來改變人.
改變人的家庭.
改變人聚居的城市.
改變這個世界的.
近年.
如果你有留意一些基督教的報章的時候.
有一段報章說到.
在越南西北部.
叫做蔡州省.
有一個地方名.
那條村叫做新水河村.

$^{1001}$那裡是做什麼的呢.
那裡是種鴉片的.
是提煉鴉片的.
整條村的人.
九成的人.
都是吃鴉片的.
所以整個世界.
就稱這條村叫做.
引君子村.
君子是君子.
要加個忍字.
引君子村.
賣鴉片呢.
好像菜市場賣菜一樣.
整條村就在那個毒品的那個綑綁當中.
我記得我小時候.
我媽說你不要晚上出去.
現在治安不好.
這裡隔壁屋有人吸毒.
所以這些吸毒的人在當中的時候.
環境是很惡劣的.
所以這條村的生活很貧困.
因為每個人都習慣吸毒.
不再去耕種.
後來呢.
因為福音廣播.
傳入那條村的時候.
那些人就開始聽福音廣播.
聽見耶穌基督的福音.
所以開始就有一些村民去悔改信主.
就戒毒除去那些毒癮.
後來就更多的人.
歸向耶穌基督.
隨著福音在那條村紮根的時候.
依賴毒品的那種衛生.
那些人就開始轉行.
去做其他的工作.
經過十幾年的改變.
那條引君子村.
現在就變成了東南亞知名的旅遊景點.

$^{1041}$耶穌基督的福音.
耶穌基督的話.
耶穌基督在那裡的時候.
整條村都改變了.
讓我們看到.
當基督進入一個人的生命的時候.
改變了.
家庭改變了.
那條村改變了.
所以撒瑪尼亞人發現.
他們真正需要的.
不是再爭論這些.
勾心鬥角的事.
不是爭論這些守律法.
或者是守這些傳統的事.
他們真正需要的是.
耶穌基督.
於是他們就說.
主啊.
請你住在這裡.
最寶貴的是.
耶穌基督沒有拒絕他們.
沒有拒絕這些撒瑪尼亞人的邀請.
起初.
耶穌基督和撒瑪尼亞婦人.
要一杯水的時候.
當時婦人戴著.
看不起或者是.
針對受事的目光去說.
你是猶太人.
為什麼要向我一個撒瑪尼亞婦人要水呢?.
現在是整個撒瑪尼亞城的人.
都邀請耶穌基督留下.
於是耶穌就在那裡住了兩天.
耶穌基督在世上.
33年半.
而在這33年半當中.
有兩天是在撒瑪尼亞的.
這不僅顯明.
祂願意拯救撒瑪尼亞人.

$^{1081}$也顯明祂願意與他們同住.
因為祂來不是要定人的罪.
而是要尋找.
拯救那些失喪的人.
那些墮落了的人.
那些迷失了的人.
那些沒有方向的人.
耶穌說過「人看見我就是看見父」.
祂所顯明的不是人的樣式.
更是天父的樣式.
聖經沒有記載祂在這兩天.
做了什麼,說了什麼.
但是在這兩天.
撒瑪尼亞人對那個婦人說.
現在我們信不是因為你的話.
是因為我們親自聽見.
亦親自看見耶穌基督.
知道祂是真正的救世主.
祂就是真正的救世主.
起初他們因著婦人的見證.
而來到耶穌這裡.
後來他們親自聽見,親自看見.
他們與耶穌基督有接觸.
親自去認識耶穌基督.
信仰不僅是別人說祂的見證.
而是成為我們自己的經歷.
不僅是間接的那種相信.
他相信,所以我也相信.
有些信.
而是成為我們親自去相信.
這是每一個基督徒的經歷.
不僅是聽聽身邊的朋友.
是因為我也與這位耶穌基督有生命的關係.
可能因為牧師,父母,朋友認識耶穌基督.
但最終還是要自己與耶穌基督建立這個關係.
親自去認識耶穌基督.
親自去思想祂的話.
思想真理.
別人的信仰不可以代替我們的信仰.
別人的經歷也不可以代替我們的經歷.

$^{1121}$撒瑪利亞二人經過親自認識耶穌之後.
得出一個很偉大的結論.
這個就是救世主了.
這個就是.
救世主.
這個稱呼不是出自猶太人的口.
不是出自門徒.
而是這一群被人看不起的撒瑪利亞人.
當豬一樣的人.
他們竟然宣告.
耶穌基督就是救世主.
為什麼?.
因為他們看到.
連撒瑪利亞婦人.
這個婦人.
全城都出醜了.
全城的人都在議論她.
這樣一個人都還有希望的.
耶穌都願意去拯救她.
她認識了耶穌.
竟然是改變了的.
既然拯救這樣一個人的時候.
我們呢.
我們應該是比她更有盼望之力的.
我們不是好像比她站得更高嗎?.
所以他們就知道.
這個真是救世主了.
連這樣一個人.
她都還有希望.
所以他們相信.
耶穌不單止是猶太人的救主.
也是他們的救主.
是全世界人的救主.
撒瑪利亞城的人.
就這樣歸鄉了.
耶穌基督已經完成了.
祂這兩天.
本來是去加利利的.
現在專登來到撒瑪利亞.
為了拯救這個婦人.

$^{1161}$原來不單止是這個婦人.
祂不只是看到這個人.
祂是看到一個更廣闊的和場.
祂看到一座城的救恩.
祂來這裡.
福音已經傳開了.
很多的人已經歸向基督了.
祂就在這個時候.
兩天之後.
祂就離開了.
繼續去加利利.
主不會無緣無故去一個地方.
祂每一步都有祂的計劃和祂的旨意.
祂是要尋找失喪的.
被人撇棄的.
已經是無可救藥的.
可能是一個吸毒的人.
全世界的人可能都沒有信心再挽回.
可能祂的父母都已經筋疲力盡了.
但耶穌基督就是要來尋找這些人.
祂今天依然在尋找.
祂要將救恩帶到人的生命當中.
祂要得著更多的人.
賜福給更多的人.
在網上的朋友.
願你得著耶穌基督.
願主賜福給你們.
願我們也好似耶穌基督一樣.
以專行神的旨意為我們的心志.
以修葛莊稼為我們的使命.
並且盼望更多的人.
歸向我們的主.
我們的神耶穌基督.
我們要專行神的旨意.
造成神祂的功.
阿們.
\newpage



\section{士師記 6:12}
\label{sec:032270}
\textbf{聖經人聖經事:耶和華的使者 呼召基甸}~~~~黃克勤牧師博士
\newline
\newline
連結: \href{http://www.church.com.hk/acms/content.asp?site=cdc\&op=show\&type=product\&code=032270\&layout=sermon}{\texttt{032270}} ~~~~ 語音日期: 2026-06-08
\newline
\newline
\hyperref[sec:032269]{\small{< < < PREV SERMON < < <}}
~
\hyperref[sec:index_chronic]{\small{[返順時目]}}
~
\hyperref[sec:032271]{\small{> > > NEXT SERMON > > >}}
\newline
\newline
士師記 6:12
\newline
\begin{longtable}{cl}
\hline
\hline
章節 & 經文 (和合本修訂版)\\
\hline
6:12 & \begin{tabularx}{0.7\textwidth}{X} 耶和華的使者向基甸顯現，對他說：「大能的勇士啊，耶和華與你同在！」 \end{tabularx} \\ \\
[1ex]
\hline
\hline
\end{longtable}
$^{1}$各位姐妹,我們又讀聖經了.
今天是耶和華的使者呼召機典.
講呼召,大家留意了.
看看神會不會呼召你.
聖經人聖經時,看看上帝背後在人與事之間做些什麼.
在生命樹宣教網絡為大家提供的.
耶和華的使者呼召機典.
士師記第六章,耶和華的使者到了訛佛拉.
坐定,坐在亞比爾者族人約阿斯的橡樹下.
講這些做什麼呢?.
就是來到機典的家.
也就是說,耶和華的使者來的時候.
是不是預先已經想定,選定了一個人了.
是的,所以他一去就去到哪裡呢?.
就去到訛佛拉這個地方.
就坐在亞比爾者族約阿斯的橡樹.
就是他自己住的地方的原地.
約阿斯的兒子,約阿斯的兒子就叫做機典.
機典正在酒駕那裡打麥子.
說要防備米電人.
這麼有趣,要解釋一下.
這一節有很多東西要說.
在這裡,耶和華的使者.
我們一看耶和華的使者就會想.
神的天使來到了.
很有可能特別在舊約聖經.
耶和華的使者很多時候不是天使.
是耶和華本身的一個代表.
是代表著耶和華本身來到這裡.
用一個叫做使者的身份出現的.
你知道上帝使靈是沒有一個肉身的.
祂來到跟人說話.
很多時候祂會現身.
祂用一個耶和華的使者.
這樣來現身跟我們說話的.
這個耶和華的使者在今次.
我比較從祂的下文看.
祂比較像,語氣等等.
比較像是耶和華本身.
神不如簡單說.

$^{41}$上帝來到俄弗拉坐定在這裡.
坐在機典的門口,橡樹那裡.
看著,看著.
那時候機典,爸爸就叫約阿斯.
約阿斯的兒子機典在做什麼呢?.
是打麥紙.
打麥紙就是說那些麥已經成熟收割了.
但是那些麥水還在上面.
有麥的殼這樣.
這樣黏著它,不能吃,不能用.
打麥於是就要將那些麥拿著.
打下去,讓它掉下來.
不然你這樣用手,手也會受傷.
這樣打下去.
讓那些麥紙掉下來.
麥紙掉下來之後,他們就要用工具磨.
磨了它,磨了它之後.
就等麥紙的殼離開了一些肉.
於是這些殼和肉就混在一起.
磨了之後,他們就要剷一剷它.
就叫做「播機揚正他的牆」.
叫什麼呢?新西翁說.
播機一兜下去剷了.
其實是用植物這樣織成的.
於是向天空拋.
拋的話,麥紙就掉下來比較快.
因為那些肉比較重.
那些殼,那些坑就比較輕一點.
剛剛找一個有風的季節才拋.
於是一拋起.
那些風吹了那些殼和坑.
那些麥紙掉下來.
它就可以掃一掃.
就這樣做了.
整個場地,整個麥田.
很多麥紙.
那怎麼做呢?.
它就要這樣做.
「揚正他的牆」是整個打麥紙的牆.
所有麥紙都要這樣做.

$^{81}$拋得一次.
其實這樣做打麥紙要很大地方的.
怎會在走站那裡做呢?.
走站其實如果今天來說.
大概都只不過會是三米乘三米.
不是很大的地方.
一個三米乘三米的地方.
而且多數在山上因為是斜斜的.
撞了石斜斜的.
於是有四條邊.
於是走站就把葡萄扔進去.
扔到差不多高.
所以我說有條邊就是這個意思.
就好像一個大的池.
把葡萄扔進去.
然後脫掉鞋子.
就踩爛葡萄.
讓汁流出來.
在最尾那裡斜斜的.
有一個洞把葡萄汁流到下面.
挖一個坑.
載著葡萄汁.
撇走所有渣.
舀葡萄汁上來.
就開始拿去釀酒.
所以走站就不是很大.
而且走站不會在路邊.
多數都在後山.
有自己的走站.
為什麼會在走站打麥呢?.
第一,麥不多.
多就不夠擺.
第二,就是分開來打.
少少地少少地.
為什麼呢?.
不要被米店人看到.
周圍的米店人.
走來走去巡邏.
嘩嘩嘩嘩.
如果見到有麥子.

$^{121}$他們會搶的.
那就慘了.
沒有了.
於是躲起來.
打都不敢那麼大力.
還要找人看著.
「細佬,看著有沒有米店人」.
「如果有,我們馬上收起來」.
「不打住」.
所以在走站打麥子.
機電一個年輕人.
大大隻的.
不過他都要躲起來.
原因就是米店人太強了.
周圍都是.
要在走站打麥子.
第一,麥子不多.
第二,要躲起來做.
因為怕米店人看到.
哎呀,很慘.
神在看著他.
要這樣躲起來打你.
你都算慘了.
耶和華的使者就向機電顯現.
其實耶和華就向他顯現.
就對他講.
大能的勇士.
耶和華與你同在.
叫他做大能的勇士.
耶和華與你同在.
叫他大能勇士.
忽然之間.
如果神跟我說.
大能的勇士.
我都不知道認不認好.
我又不接受.
叫他做大能勇士.
那怎麼認呢?.
耶和華與你同在.
耶和華是看他為大能勇士.

$^{161}$大能就不用躲起來.
在走站打麥子.
這樣都要躲起來.
即是說.
其實大能都變無能了.
很害怕,很怕,很怯.
但神叫他做大能的勇士.
給他一個膽量.
耶和華與你同在.
這一點最為重要.
耶和華與你同在.
你想不到.
基澱這個信教神竟然這樣說.
你看他的信心全沒了.
基澱就說.
主啊.
耶和華若與我們同在.
我們何至遭遇這一切事呢?.
嘩.
真的頂過來了.
哎呀.
大能的勇士.
耶和華與你同在.
他就說.
同什麼在?.
如果真的同在.
我們就不會慘成這樣了.
怎麼會遭遇這麼慘呢?.
這麼淒慘,淒涼.
現在躲起來.
嘩.
如果我是神.
真的一巴掌打過去.
這樣跟我說話.
不是吧.
你哪位啊?.
多大啊?.
他接著就解釋.
我們的列祖不是向我們說.
耶和華領我們從埃及上來嗎?.

$^{201}$耶和華領我們從埃及上來.
這個是列祖跟我們說的.
大家以後歷世歷代要記得一件事.
記得什麼?.
就是這句.
耶和華領我們從埃及上來的.
記得啊.
去到哪裡都不要忘記.
我們的祖先就這樣跟我說.
聽到現在是我.
只不過.
唉.
這麼奇妙的作為現在在哪裡啊?.
還有沒有領導啊?.
這麼厲害.
這麼奇妙的作為現在還有嗎?.
沒有啦.
事過境遷以前就有.
現在沒有了.
現在是怎樣啊?.
耶和華卻丟棄我們了.
耶和華已經沒有理會我們了.
我們自生自滅.
現在很淒涼.
將我們交在米殿人手裡.
差不多了.
耶和華要我們同在了.
那基甸怎麼回答啊?.
如果用今天來說.
不是吧.
有沒有說錯啊?.
你還說耶和華要我們同在.
我們的列祖祖先就說.
耶和華把我們埃及領出來.
多厲害啊.
不過今天我看不到啊.
神的作為在哪裡啊?.
如果祂還在就不會丟棄我們了.
其實祂已經沒有理會我們了.
哇.

$^{241}$唉.
如果我是上帝聽到這句話.
我的心真的會痛啊.
唉.
怎麼這麼說話啊?.
這麼嘴硬的說話.
這樣你都說得出.
其實我是一直帶領你們.
保護你們的.
這樣都說得出.
那麼米殿人呢.
如果你保護我.
怎麼會有今天呢?.
哦.
好啊.
這樣.
第一點要說的.
就你靠著你這能力.
去從米殿人手裡拯救以色列人.
不是我差遣你去的.
哇.
讀個詞啊.
不要看那兩個藍色字.
你靠著你這能力.
去從米殿人手裡拯救以色列人.
不是我差遣你去的.
啊.
這份這麼有趣的.
如果呢.
將那些比較不是等等那些拿起來.
我將它迎光了.
那麼你會知道.
其實神是用一個.
小小的反話這樣說.
唉.
你上帝已經丟棄了我們了.
我們怎麼會衰到現在被米殿人欺負成這樣.
又說大能.
又說保護我們.
不要說了.

$^{281}$其實神就說.
我是想你.
做什麼呢.
我是想你.
靠著你這能力.
去.
從米殿人手裡拯救以色列.
我想你第一.
靠著你現在的能力.
去.
這個靠著你這能力.
這能力是什麼.
你現在有的.
英文叫as is.
as is 就是.
照現在有的.
那你不教我十八道武功.
我怎麼去啊.
那你起碼讓我大隻一點.
或者這樣.
那不用了.
就照你現在呢.
靠著你的能力就去得了.
因為我和你同在嘛.
那麼這個as is的一個概念.
是我覺得很重要的.
譬如什麼叫as is呢.
好像現在這樣.
那麼有些時候我們租屋啦.
買屋啦.
啊看完了多好啊就現在這間啦.
我租這裡.
或者我買這間屋啦.
那麼我想你不要動那些裝修.
不要動那些所有東西.
as is.
就是現在看到的.
現在有的.
我買了下來.
我租了下來.

$^{321}$as is.
那麼.
上帝是說.
你就照著.
基地呢.
現在我就照著你.
現在as is有的能力.
我就差遣你.
去拯救以色列.
啊.
不是吧.
那麼就當然我們會覺得.
我現在這樣怎麼做啊.
靠著你的能力去啦.
去從委員會拯救以色列.
那麼不是我差遣你去的嗎.
問好啊.
不是我派你去嗎.
就是我派你去你就可以了.
但是其實呢.
應該是.
靠著你這個能力拯救以色列.
其實是我差遣你去的.
如果正面不用問句的話.
不用問句的話.
靠著你現在有的能力.
就去從委員會拯救以色列.
是我派你去的.
我派你去的.
那麼.
所以這兩點我常常跟他說就是.
靠你現在的能力.
第二.
是我派你.
那麼這兩點很重要了.
當有一件工作.
有一件事工臨到你.
常常叫你去做呢.
那麼我哪裡行啊.
常常說.

$^{361}$行的.
照你現在的能力就可以了.
啊.
不行啊.
那為什麼.
因為.
我有能力嘛.
是我和你同在嘛.
那麼.
還有第二你記得.
是誰叫你的.
是奉旨去的.
是神差遣我去的.
派我去的.
照我現在的能力.
派我去.
好.
那麼基建就說了.
啊.
主啊.
我有何拯救以色列的人呢.
我家在嗎哪西支派中.
是第三點.
知貧窮.
知微小的.
我哪裡行啊.
基建就說.
不可能啊.
as is.
照我現在有的.
我是不可以了.
我做不到.
為什麼呢.
因為.
我當然不行了.
怎麼拯救到以色列人呢.
因為我在嗎哪西支派裡面.
是最窮.
最少的.
微小.

$^{401}$最低微.
你看看.
我窮到.
打麥我都沒有一個和場.
我只有一點點麥子.
在家裡的酒窖躲起來打打打.
吃了這一餐一兩餐之後再怎樣我都不知道.
我窮到這樣.
那麼又知貧窮.
知微小.
這麼大的任務.
我怎麼做到啊.
這個就是基建.
是承認的.
這個是呼召臨到基建.
而基建的反應是這樣的.
我知貧窮知微小.
很多時候我們都是這樣.
我又試過你又試過.
常常叫我做一些事我就說.
哎呀.
我書又讀不夠.
樣又不夠帥.
家裡又沒搞得定.
那麼我自己問題軟弱多多.
做不到啦.
叫他吧.
那個厲害很多.
哇那些真的厲害.
真的厲害.
那麼就.
又不怕和你分享我個人的經歷見證.
我本身就是這個宣導會希伯倫堂出來的.
那麼希伯倫堂裡面以往是有個很厲害.
我跟著他成長的牧師李飛吾牧師.
他有一個很厲害的兒子.
李思敬博士.
哇做過中神院長的.
在英國讀神學的.
很厲害.

$^{441}$他在我們一起中學.
和一個團契長大的時候.
他已經傑出了.
他考會考成績很好.
英文很好.
鋼琴很好.
領導能力很好.
講聖經很清楚.
已經很厲害了.
差不多你說.
如果要他去做傳道人.
去做工作.
事奉李思敬了.
我們見到他多偉大的工作.
李思敬.
我不是吧.
我當時很瘦.
肥包骨.
讀書又不厲害.
會考要考兩年才夠五科.
考一次未包括數學.
那麼就不厲害了.
他一坐下就已經是狀元級.
我們坐下成績很差.
要去事奉.
不過你不知道.
在一次下令會裡.
我和他另外三個人.
是五個人同時蒙召的.
就是我了.
其中一個.
蒙召做神的工作.
做傳道人.
他啦.
思敬啦.
我這些怎麼行.
不過你叫到.
我不妨一試.
頂住先.
跟著跟著去.

$^{481}$其實神是呼召我們.
As is.
就照你現在來跟隨我.
就照現在.
我哪有這樣的能力.
我和你同在.
我最弱.
最貧窮.
最瘦.
最微小.
家裡可以說.
第一個認真信主的就是我.
家裡沒有人信耶穌.
這樣都可以.
這樣都呼召我.
你會不會叫錯人.
我在我家族中.
最貧窮.
最微小.
找那些厲害一點的.
教會裡面很多厲害的人.
不到我.
不過神是說.
我和你同在.
我去叫你.
那我就起來做這個工作.
也都見到今天.
神真的透過.
我祝福到一些人.
都很感恩.
所以都是不怕的.
大家知道.
最重要就是神說.
靠著你這個能力去.
和我差遣你就可以了.
也都最重要的是什麼呢.
大量的勇士.
耶和華與你同在.
神和你同在.
你就不怕去.

$^{521}$而且是靠著.
你現在有的能力.
是神派你去.
記住有更加剛強壯膽的心.
其實看下去.
下面再看多幾天.
你會發覺.
原來機電是信心很小的.
信心不大起碼.
不過他面對的挑戰很大.
所以認到他的信心.
顯然是真的.
嘩!行不行啊.
去不去到啊.
給點證據行不行啊.
給點憑據啊.
都幾大考驗的.
我們繼續看下去.
記得.
靠著你現在有的能力.
上帝差遣你.
哎呀!我好卑微啊.
好窮困.
不怕.
神可以的.
望在神的堅強的手裡面.
我們一起禱告.
多謝主.
原來你是使用了很多.
自卑微.
自貧窮的人.
讓他.
照著他手裡面所有的小小.
就成就了很偉大的事功.
你是欣賞那個.
兩個小時的奉獻.
你自己是用五餅二魚.
來自一個小朋友的手.
一個天真的奉獻.
就餵飽了五千人.

$^{561}$多謝我們的主讚美你.
我們在這個基面身上有所學習.
他竟然這麼大膽跟你說.
神的看顧在哪裡啊.
不要再說了.
但是神你就是揀選他這樣的人.
你揀選你的同在最重要.
給我們看重.
隨天祈禱.
奉耶穌基督名求.
阿們.
\newpage



\section{約翰福音 8:21-29}
\label{sec:032271}
\textbf{榮耀君王:死在罪中}~~~~林鍵牧師
\newline
\newline
連結: \href{http://www.church.com.hk/acms/content.asp?site=cdc\&op=show\&type=product\&code=032271\&layout=sermon}{\texttt{032271}} ~~~~ 語音日期: 2026-06-06
\newline
\newline
\hyperref[sec:032270]{\small{< < < PREV SERMON < < <}}
~
\hyperref[sec:index_chronic]{\small{[返順時目]}}
~
\hyperref[sec:032272]{\small{> > > NEXT SERMON > > >}}
\newline
\newline
約翰福音 8:21-29
\newline
\begin{longtable}{cl}
\hline
\hline
章節 & 經文 (和合本修訂版)\\
\hline
8:21 & \begin{tabularx}{0.7\textwidth}{X} 於是耶穌又對他們說：「我去了，你們會找我，而你們會死在自己的罪中；我所去的地方，你們不能去。」 \end{tabularx} \\ \\ \relax
8:22 & \begin{tabularx}{0.7\textwidth}{X} 猶太人說：「他說『我所去的地方，你們不能去』，難道他要自殺嗎？」 \end{tabularx} \\ \\ \relax
8:23 & \begin{tabularx}{0.7\textwidth}{X} 耶穌對他們說：「你們是從下面來的，我是從上面來的；你們是屬這世界的，我不是屬這世界的。 \end{tabularx} \\ \\ \relax
8:24 & \begin{tabularx}{0.7\textwidth}{X} 所以我對你們說，你們會死在自己的罪中，你們若不信我就是那位，就會死在自己的罪中。」 \end{tabularx} \\ \\ \relax
8:25 & \begin{tabularx}{0.7\textwidth}{X} 他們就問他：「你到底是誰？」耶穌對他們說：「我從起初就告訴你們了。 \end{tabularx} \\ \\ \relax
8:26 & \begin{tabularx}{0.7\textwidth}{X} 我有許多事要講論你們，判斷你們；但差我來的那位是真實的，我從他那裡所聽見的，就告訴世人。」 \end{tabularx} \\ \\ \relax
8:27 & \begin{tabularx}{0.7\textwidth}{X} 他們不明白耶穌是對他們講父的事。 \end{tabularx} \\ \\ \relax
8:28 & \begin{tabularx}{0.7\textwidth}{X} 所以耶穌說：「你們舉起人子以後就會知道我就是那位了，並且知道我沒有一件事是憑著自己做的。我說這些話是照著父所教導我的。 \end{tabularx} \\ \\ \relax
8:29 & \begin{tabularx}{0.7\textwidth}{X} 差我來的那位與我同在；他沒有撇下我獨自一人，因為我一直行他所喜悅的事。」 \end{tabularx} \\ \\
[1ex]
\hline
\hline
\end{longtable}
$^{1}$歡迎大家來到我們的周六崇拜.
上年七月 我們開始研讀《新約聖經》當中的《約翰福音》.
我們講到系列的題目是榮耀君王.
我們花了半年的時間.
完成了《約翰福音》第一章至第四章的宣講.
今年一月 我們開始研讀《約翰福音》第五章.
接著五個星期前 我們來到《約翰福音》第七章.
從故事的時序來說 《約翰福音》第七章可以分為三大段.
一開始是猶太人的駐彭捷 近遙.
還沒到 很近的了.
接著到中間 節日已經過了一半.
節日是七至八天.
所以應該是節期的第四天.
耶穌終於來到耶路撒冷.
接著那個戲玉是第三個部分.
節期的最後一天.
節期的最後一天 耶穌在耶路撒冷的聖殿裡.
向眾人宣告兩件事.
祂說 我是生命的活水.
我是世界的光.
耶穌做了兩個自我的啟示.
表面上這兩個是很不同的自我啟示.
但其實這兩個啟示.
好像是分了在兩張聖經.
一個是《約翰福音》第七章.
另一個是《約翰福音》第八章.
可不可以留意兩個啟示是被一個故事所分隔了.
那個故事是第八章第一節.
十一節一個犯奸淫事被抓到的女人.
然後耶穌原諒了她.
這個故事打斷了兩個自我啟示.
然而這個故事.
我們相信在《約翰福音》的原稿裡是不存在的.
是後來加進去的.
因為最古舊的原稿裡是沒有這個故事.
簡單來說.
其實兩個自我啟示是同一個自我啟示.
當你將中間的故事拿走的時候.
我們現在變成了一個二合一的自我宣告.
耶穌說三句話.

$^{41}$其實耶穌第一.
祂開播的時候祂說我是.
我是其實是上帝的名字.
耶穌說你知不知道我是誰.
我是舊約的上帝的化身.
我是舊約上帝裡面的獨一兒子.
祂的名字叫我是.
我的名字也叫我是.
我們是同姓的.
祂說我是.
然後祂說我是生命的活水.
我同時又是世界的光.
是一句話來的.
我是生命活水.
我又是世界的光.
是同一句話來的.
我想帶你回去看這兩個我是的自我宣告.
第七章第37節.
至期的最後一天.
駐彭捷的最後一天.
就是最隆重的一天.
耶穌站著喊著說.
這個是耶穌自我宣告的背景.
如果你記得我早幾個禮拜說過.
當時猶太的朝聖者已經經歷了七天.
駐彭捷的水禮.
每一天七天的每一天.
大祭司都會有一個金水壺.
從西羅亞池那裡取水.
然後一大群人跟著他.
來到耶路撒冷聖殿裡面的祭壇.
而當這個大祭司倒水到祭壇的時候.
同時之間也會獻上酒祭.
於是有水有酒.
一起倒到聖殿的祭壇.
聖殿的祭壇又有血.
於是你可以想到那些顏色的.
混合有水有酒從聖殿的祭壇裡面流出來.
這是一個何等榮耀的福音圖畫.
因為當耶穌被掛在十字架上.

$^{81}$最後有士兵去戳祂身體的時候.
有血和有水從祂的身體流出來.
就是這個圖畫.
就在這個意義深重的水禮的背景之下.
耶穌站起來喊著說.
「人若喝了,到我這裡來喝」.
「信我的人,就如經上所說的」.
「從他的福中將流出活水的乾涸力」.
其實耶穌是故意讓所有的水禮都完結了.
祂不是中間有一個水禮跟祂鬥嘈.
祂就出來說.
「不是,讓所有的東西都靜下來」.
耶穌說「我將要用我自己去取締」.
「去承傳著彭哲的水禮」.
約翰這樣解釋耶穌的宣告.
約翰這麼說.
第七章第三十九節.
「耶穌這話是指信他的人要受聖靈說的」.
「那時還沒有賜下聖靈」.
「因為耶穌還沒有得到榮耀」.
「所謂的活水乾涸是指向聖靈的降臨」.
「才是聖靈從耶穌的福中流出來」.
大家記得那時我和大家說過.
原文有兩個可能.
第一個福中是說耶穌的福中.
接著是說信徒的福中.
按我們的思想.
先是聖靈會從耶穌的福中流出.
耶穌的福中好像是末世的聖殿裡流出活水.
是以西結書第四十七章裡的一個意象.
接著那個活水要流向黑暗罪惡的世界裡.
接著回轉悔改的信徒.
前來領受聖靈的施洗.
因為那個活水.
耶穌流出來的活水是聖靈的代號.
聖靈的施洗可以潔淨我們的罪孽.
潔淨我們的污穢.
施洗約翰在馬可福音第一章第八節曾經這樣說.
「我是用水幫你施洗的.
但然而後來有一個比我強大得多.

$^{121}$他會用聖靈給你們施洗」.
於是耶穌在十字架得到榮耀之後.
他就有聖靈為我們施洗.
因著聖靈回轉的罪人會得到一個新的心.
也是以西結書的語言.
接著我們會被聖靈所重生.
怎樣重生?.
從天而生.
成為上帝的兒女.
稍後你繼續讀下去.
你會發現耶穌都很強調這件事.
是從天而生.
於是我們信主的人.
真正重生得救的人.
在他們的福中都會流出活水的江河.
而這個從信主的人.
真信主,真跟隨主的人流出的活水江河.
會在這個黑暗的世界裡.
為上帝明光見證.
上帝永恆的天國.
於是從耶穌的活水江河.
來到我們的活水江河.
不是單獨的活水江河.
而是集體的活水江河.
接著我們將這個福音傳到地極.
因而第七章第39節.
聖靈降臨之前耶穌必須要先得著榮耀.
所謂得著的榮耀不是在地上拿什麼獎.
不是地上的榮耀.
而是十字架的榮耀.
在聖靈未降臨之前.
耶穌必須要上十字架.
上十字架是一個最榮耀的地方.
唯有這樣.
聖靈活水才可以將耶穌救贖的寶血.
在十字架上耶穌流的寶血.
聖靈帶著耶穌的寶血帶到罪人的身上.
接著再帶進罪人的體內.
以致我們可以流出活水的江河.
這是第一幅圖.

$^{161}$第一個「注我宣告」或者Part A.
耶穌不單止要為我們成為生命的活水.
祂也要為我們成為世界的光.
於是我們跳到第八章十二節.
其實是Back to Back.
是Part A和Part B.
不是中間被故事打斷了.
其實是同一句說話.
耶穌說:我就是世界的光.
跟從我的必不在黑暗裡走.
卻要得著生命的光.
我又要你記得一件事.
這個「自我啟示和宣告」.
又是連接第七章第三十九節.
其實是未發生的.
不可以發生.
未發生要等到聖靈降臨.
但聖靈必須要有耶穌.
德榮要耶穌要上十字架.
聖靈降臨.
才可以成為世界的光.
才可以成為生命的光.
首先你留意一下.
所謂世界的光就等於生命的光照遍世界.
其實就是生命的光.
有生命的光.
接著這個生命的光照遍世界.
就變成了世界的光.
一個人要得著新的生命.
得著上帝從天而生的生命.
需要有十字架的寶血.
唯有十字架的寶血.
耶穌基督的寶血.
才可以救贖罪人.
賜人生命.
第二.
唯有聖靈的洗禮.
才可以重生一個罪人.
讓這個人走出黑暗.
走進光明.

$^{201}$我們需要聖靈的洗禮.
需要聖靈的重生.
以致我們可以走出黑暗.
以致我們走進光明.
接著.
唯有聖靈降臨.
才可以給教會大能大力.
將十字架的福音傳到地界.
讓生命的光變成世界的光.
就是一個照亮全世界的生命光.
《紹行苑》第一章第八節.
但聖靈降臨要在你們身上.
你們必得著能力.
並要在耶路撒冷 猶太 全地 撒瑪利亞 直到地極作.
我的見證.
耶穌說了很多話.
說了一句話.
但是一句話背後有很多很多很多很深的意思.
而這句話.
祂說的.
讓我告訴你.
第八章第二十節.
這些話.
耶穌是在聖殿裡面.
銀婦女教導人的時候所說的.
在某一個地方.
在聖殿裡面有銀婦的地方.
你記不記得我上星期帶你去看這個地方.
耶穌所在的銀婦是在婦人的庭院裡面.
The Court of Women.
在婦人的庭院裡面.
耶穌做了一個世界之光的宣告.
你記不記得這個地方有一個很特別的地方.
祂有四條很高高的柱子.
那個柱子大概有七十至七十五尺高.
十個NBA Player.
每條柱子上面有四個金碗.
是用來裝橄欖油.
一共是十六個金碗.
裝著八十加侖的橄欖油.

$^{241}$所有的一切是為了要預備每年駐柏林寺.
連續六天晚上所舉行的火禮.
如果你上星期.
我想大家都記得的.
火禮.
祂要著燈 點燈的時候需要有一個燈芯.
那個燈芯原來是用當時祭祀的舊衣服.
以致可以成就這個火禮.
當然所有的一切都是指向耶穌的.
這個就是當時耶穌宣告.
祂說我是世界的光.
我是生命的光.
這個宣告的場景和背景.
正如我所說.
耶穌不是和祂一起搶鏡.
耶穌等這個火禮完結之後.
接著祂站起來宣告自己的身份和使命.
於是在六天的火禮完成之後.
耶穌以自己取締了這個火禮.
就如祂用自己去取締去成全了水禮一樣.
沒有火了.
耶穌說我是世界的光.
有少許應該說我是世界的火.
在古代人光不是按個按鈕.
然後有電就著燈.
光是和火連貫在一起.
於是耶穌其實告訴他們.
最終取締著行節火禮和水禮的.
就是祂自己的榮耀十家.
祂被上帝的憤怒所燃燒的榮耀十家.
那個是行節火禮的成全.
還有聖靈的降臨.
是水禮的成全.
這個就是耶穌向當時在聖殿當中的猶太人.
跟著他聽他說話的猶太人.
所宣講的榮耀福音.
其實祂想和他們說什麼.
祂來是為他們上十字架.
祂來是要將聖靈帶進他們的人生.
讓他們可以重生.

$^{281}$讓他們從天而生.
為上帝的天國而活.
這個是耶穌的訊息.
我們問這個問題.
他們有沒有聽明白耶穌的訊息.
有沒有聽得入耳.
答案是聽不明白.
也都聽不入耳.
於是接下來耶穌就心急了.
就心急了第八章第21節.
於是耶穌又對他們說.
因為他們聽不明白.
耶穌說我去了.
而你們會死在自己的罪中.
我所去的地方.
你們不能去.
為什麼我說耶穌心急呢.
為什麼耶穌說他著急呢.
原來這一段說話在第七章裡面都有出現過.
我讀一讀第七章的版本給你們聽.
第七章33和34節.
於是耶穌說.
我跟你們在一起的時候.
不會太久了.
我要回到那差我來的那裡去.
你們要找我.
你們找不到我.
我所在的地方你們不能去.
是不是很相似.
耶穌說我的時候差不多了.
我要回去差我來的那裡去.
然後他說你要找我.
不過你找不到的了.
然後他說我所在的地方你們不能去.
但是在第八章裡面.
耶穌特意.
甚至我可以說他很激動地.
不隨意地加了一句從來沒說過的說話.
他加了哪句說話.
你們會死在你們的罪中.

$^{321}$好像一個牧師最後一次說道.
看到你們聽了這麼久的福音.
都不醒覺的.
然後我最後一次了.
下個星期不是我了.
我死了.
我最後一次.
我要你們醒覺.
我再說一句.
說什麼一句.
我走了.
You will die in your sin.
你會死在你的罪中.
那些猶太人聽不聽得明白.
聽不聽得入耳.
完全聽不明白.
完全聽不入耳.
還在取笑他第八章二十二節.
他就說.
他說他去的地方.
我們去不了.
難道他要自殺嗎.
在第七章裡面.
耶穌說.
我要去的地方你去不了.
他們那時候還在猜.
他是不是要去外邦人的地方.
去到這裡.
他們變本加厲.
他們其實是在取笑耶穌.
用不著這麼肉緊.
用不著這麼緊張.
我知道你最後說到.
說完就算了.
單指甲路.
不用再回來了.
又要說到要生要死.
又要說我們在罪裡面.
我坐在這裡都傻.
你是不是想自殺.

$^{361}$其實他們在詛咒耶穌.
在取笑耶穌.
他從來沒有想過.
十字架.
正是耶穌.
要返回天父那裡.
天上榮耀的必經之路.
但其實話說回頭.
十字架不是必經之路.
誰告訴你.
耶穌回家一定要經過十字架.
可以直飛的.
為什麼一定要去十字架.
耶穌之所以一定要去十字架.
是因為祂要救我們.
如果耶穌不要救你.
不要救我.
祂根本不需要上十字架.
十字架成為了祂的必經之路.
祂的中轉站.
是因為唯有這樣.
外邦人 猶太人.
才可以有機會.
有一個可能.
當我們回轉得救的時候.
我們可以得救.
我們可以得著生命.
於是耶穌說.
耶穌說.
你們會死在罪中.
那一節.
每一張slide都會有.
我真的告訴你.
這一節是整段聖經的核心.
耶穌告訴你.
我以下所說的所有事情.
都是為了要抗衡上面那一句說話.
那一句說話必定發生.
You will die in your sin.
你會死在罪中.

$^{401}$我這次走了之後.
以後不知道還有沒有機會.
再和你說第二篇道.
但是我告訴你.
如果你不留心聽.
你一定會死在罪中.
耶穌說.
如果你們心硬不回轉.
你就一定會死在罪中.
這個罪其實很有趣.
在原文裡面是單數.
祂不是死在很多的罪中.
而是死在一條罪中.
這個罪代表了一個罪惡叛逆的狀態.
這個狀態是什麼狀態.
下面的狀態.
下面人的狀態.
接著第23節耶穌這樣說.
耶穌對他們說.
你們是從下面來的.
我是從上面來的.
你看看上下.
你們是屬世界的.
我不是屬世界的.
下面的人有一種狀態.
世界的人有一種狀態.
我是屬上面.
你們屬下面.
下面人的狀態就是罪.
下面的人處在一個單數的罪.
在一個罪惡叛逆的狀態裡面.
下面的人叛逆上帝.
下面的人敬拜偶像.
下面的人貪愛世界.
你看不看到.
但我是屬於上面的.
我不屬於這個世界.
上面的人馴服上帝.
上面的人盡心愛主.
目前為止我是唯一在上面的.

$^{441}$你們都是在下面.
但耶穌呼召他們做什麼.
由下面升上去上面.
你明白嗎.
祂呼召你移民.
因為如果你移不到民.
你就死定了.
耶穌說.
我是上面的.
你是下面的.
你們是屬世界的.
我不是屬世界的.
一樣解釋.
如果你繼續屬世界.
你有沒有得救.
你沒有得救.
你一定要離開世界.
你們將死在最終.
你們單數的.
你們將死在罪惡叛逆當中.
你們將死在偶像假神裡面.
你們將死在地下的世界裡面.
除非你離開地下.
離開世界.
跟我一起.
上去.
耶穌這句話.
其實意義很重.
祂說.
你們是屬於世界的.
我不屬於世界.
你不要將祂想得這麼輕.
我們當然屬於世界.
人在江湖身不由己.
你再說多一次.
上面那句就是你的.
你們會死在最終.
你明白嗎.
你不要講得下巴輕輕.
我是這樣的.

$^{481}$是不是.
信耶穌也是隨便信一下.
但我們人當然是活在世界裡面.
難道活在天上嗎.
你一這樣想的時候.
上面那句就是你的.
耶穌這句話.
你們是屬於這個世界.
但我不屬於這個世界.
令我們想起祂在彼拉多面前的那句話.
讓福音裡面.
金句中的金句.
耶穌最後被彼拉多審訊的時候.
耶穌說過這句話.
祂說.
我的國不屬於世界.
彼拉多說你是不是來做君王的.
你是不是要跟我練過.
你是不是要取締我.
你是不是要跟我作敵人.
耶穌說.
我的國不屬於這個世界.
我不會跟你搶的.
你要的你自己去搶.
我不會跟你搶.
我的國不屬於這個世界.
我的國若屬於這個世界.
我的部下就會為我戰鬥.
只是我的國不屬於這個世界.
有沒有留意耶穌三次說.
我的國不屬於這個世界.
中間那次是若我的國屬於這個世界.
意思就是我的國不屬於這個世界.
耶穌三次.
我的國不屬於這個世界.
今天很多基督徒不明白這件事.
信耶穌會有什麼好處.
幫我打天下.
是不是.
有沒有分別.

$^{521}$世界艱難.
所以需要耶穌.
是不是.
你這樣想.
你就會完全死在世界當中.
你會死在你的隨從.
上面那節聖經就是送給你的金句.
我的國不屬於這個世界.
這句話令我想起.
寫約翰福音的約翰.
他當年年輕的時候.
未信耶穌的時候.
他那種愚昧無知.
那種大大的錯誤.
你記不記得這個故事.
我甚至會說他寫整本福音書.
就是要undo他當年所做錯的東西.
所說錯的說話.
他的無知錯誤.
在馬可福音第十章.
第35至38節裡面.
我讀給大家聽.
當年寫約翰福音的約翰.
年輕的時候.
西比塔的兒子雅各和約翰.
雅各是他的兄弟.
進前來對耶穌說.
老師.
我們無論求什麼.
願你為我們做.
耶穌問他們.
要我為你們做什麼.
他們對耶穌說.
在你的榮耀裡.
請賜給我們一個坐在你的右邊.
一個坐在你的左邊.
他們以為耶穌來.
會成為地上的一個君王.
說耶穌我們.
買定你的.

$^{561}$我們看好你.
你一定做王的.
我們會跟你出生入死.
不過有個要求.
到你做王的時候.
我們一個是左丞相.
一個是右丞相.
完全不明白.
My kingdom is not of this world.
他以為耶穌會幫他們打天下.
傻子一名.
耶穌怎麼回答他.
你們不知道所求的是什麼.
然後耶穌說.
我所喝的杯.
你們能喝嗎.
說的是十字架.
我所受的洗.
你們能受嗎.
原來.
年輕的約翰和當時的猶太人一樣.
完全不明白.
君王耶穌天國的本質.
今天.
香港基督徒都是一樣.
全世界基督徒都是這樣.
根本不明白天國的本質.
有什麼就求耶穌.
求耶穌的事和求王的事有沒有分別.
沒有分別.
為什麼.
都是神.
都是偶像.
在你的腦海裡.
都是偶像.
都是假神.
所以.
你們會死在其中.
當年的約翰所要的東西.
和他的弟兄雅各所要的東西.

$^{601}$是地上的榮耀.
世界的榮耀.
耶穌要和他們分享.
耶穌想不想和他們分享.
當然想和他們分享.
耶穌要和他們分享.
是十字架的榮耀.
耶穌要和他們分享.
是永恆天國裡面的榮耀.
所以.
耶穌最後會告訴你.
我所受的洗.
你也要受.
我喝的杯.
你也有份.
因為當他真正悔改相信耶穌的時候.
他會過一個背十字架的人生.
約翰之後的五六十年.
忠心愛主.
為耶穌受苦.
在他信主忠心服侍上帝五六十年之後.
約翰寫了一本書叫《約翰一書》.
在《約翰一書》裡面.
老約翰.
真的老了.
現在八九十歲了.
老約翰.
語重心長地去警告勸勉.
新一代的信徒.
警告他們什麼?.
Do not love the world.
《約翰一書》第二章第十五節.
不要愛世界和世上的東西.
若有人愛世界.
愛父的心就不在他裡面.
第十六節.
因為凡世上的東西.
好比肉體的情慾.
眼目的情慾.
今生的驕傲.

$^{641}$都不是從父來的.
而是從世界來的.
第十七節.
這世界和世上的情慾都要消逝.
唯獨那遵行上帝旨意的人.
全都永遠.
這個是他自己的自白.
是他年輕的時候.
不認識上帝.
不認識耶穌的天國.
現在他反過來.
他終於明白為什麼耶穌說.
他不屬於世界.
而你們如果要繼續屬於這個世界.
你必與世界一同死.
於是我帶你回去.
第八章第二十三節.
耶穌再三強調.
我們一定要逃離罪惡叛逆.
我們要逃離今生地下.
我們要逃離黑暗世界的種種誘惑.
我知道很多人不明白這件事.
他說不是耶穌派免死金牌給我.
我繼續在世界就能解決嗎.
你是錯誤的.
如果你說我不要上去.
我繼續會在下面.
我繼續會在世界.
不過耶穌派一個.
恩順清義免死金牌給我.
總之我就是靠它得救.
你不明白.
你必定死在你的罪中.
我可以跟你保證.
你不信當然你可以找其他黃綠醫生.
黃綠牧師.
總有人告訴你.
你一定得救.
但是你想死在罪中.
你就死在罪中吧.

$^{681}$因為耶穌開報道會的時候.
他也是這樣說的.
除非你離開下面.
除非你離開這個世界.
耶穌這裡所描述的天上地下.
令我想起C.S. Lewis那本名著.
The Great Divorce.
有沒有看過.
天地大分別.
起初天地是一個地方.
因為亞當的墮落.
天地大分別才有上帝.
才有上面下面.
不是上帝創造的時候.
上帝睡上帝.
下面我們睡.
不是這樣的.
這個世界只有一個地方.
就是上帝與我們同在.
是因為人的墮落的緣故.
現在天地大分別.
不過說真的.
這幅圖畫真的不太對.
這幅圖畫太過平靜了.
這幅圖畫好像給我們一個假象.
覺得我們很多時間去考量.
左邊還是右邊.
左邊是天堂.
右邊是地獄.
這幅圖畫又給我們一個假象.
好像覺得我們兩邊都不是.
好像我們現在這個地方.
都是很安全的.
去到臨死之前.
總之選左不選右就沒問題.
如果你是這樣想的你完全錯.
這個絕對不是我們今天的光景.
我們今天真實的光景.
是一個非常緊急的狀態.
如果我可以打999.

$^{721}$救命啊SOS.
911什麼都打了全世界.
我們在一個極其緊急的狀態.
聖經告訴我們.
起初當人類的始祖亞當.
選擇了相信撒旦的謊言.
不聽上帝的話.
叛逆上帝的時候.
他一次叛逆上帝.
他完全的叛逆上帝.
他以為自己可以掌控自己的生命.
做皇帝.
其實他成為了罪的奴僕.
於是在亞當裡面.
我們和亞當一起.
從天墮下.
彷彿踏上了一個地獄快車.
無法自拔.
完全的失控.
有很多人不明白.
他說林牧師我再想一下.
我遲一些過多兩年.
退休的時候或者.
take a break的時候我再考慮.
你以為你真的有個按鈕可以停一停嗎.
你停不了的.
誰可以停一停按鈕說暫時不計數.
多給我兩年時間.
誰可以說我選左邊信耶穌就信.
你有個按鈕可以按信耶穌嗎.
你根本沒辦法信.
聖經告訴我們.
我們每一個人都是在一個地獄的快車裡面.
我們需要立定志向停下來.
去思想聽耶穌對我們的呼召.
我想帶你去看羅馬書.
羅馬書道出了人類叛逆之後的狀況.
罪惡帶來的致命的後果.
羅馬書第五章第十二節這樣說.
追思從一人進入世界.

$^{761}$那個人是亞當.
死有是從罪而來.
於是死就臨到所有的人.
因為人人都犯了罪.
在亞當犯罪裡面.
我們既是亞當的後裔.
我們生出來帶著罪進入這個世界.
我們不需要有人教我們犯罪.
我們都曉犯罪.
而因為我們曉犯罪.
因為我們是帶著罪進入這個世界.
結果死亡會臨到我們.
聖經說世人都犯了罪.
虧欠了上帝的榮耀.
罪的功架乃是死.
為什麼我們每個人都會死.
是因為我們每個人都犯了罪.
什麼是罪.
聖經說罪就是虧欠了上帝的榮耀.
有神學家解釋.
所謂罪就是missing the mark.
未能達標.
或者是未能擊中那個目標.
我們問那個目標是什麼.
那個目標不是一些道德的標準.
那個目標是上帝的榮耀.
虧欠了上帝的榮耀.
你missing the mark.
什麼mark.
上帝的榮耀.
當我們不聽上帝的話.
當我們不信上帝的應許.
當我們不再為上帝而活.
為自己開始為自己而活的時候.
我們就虧欠了上帝的榮耀.
這個是亞當的故事.
人類的始祖亞當.
夏娃選擇了叛逆.
做什麼.
不聽上帝的話.

$^{801}$信靠他自己的智慧.
為自己而活.
做自己的君王.
於是亞當的後裔.
我們同一樣叛逆上帝.
在我們叛逆的世界裡面.
有人選擇做壞人.
亦都有人選擇做好人.
在我們教會裡面.
和我們查過經的弟兄姊妹就知道.
有人選擇做叛逆的小兒子.
亦都有人選擇做乖乖的大兒子.
但是兩個都是按照自己的想法.
做自己眼中看為對的事情.
兩種人都是為自己而活.
兩種人都當上帝武道.
我就是一個好人.
我就是一個壞人.
我喜歡怎樣都行.
兩個人都是在他們的人生規劃裡面.
在他們的圖畫裡面.
上帝是不存在.
這個叫做罪.
什麼是罪.
在你人生規劃裡面.
上帝不存在.
這個叫做罪.
不是你要殺人放火.
殺人放火是罪.
做一個好人都是罪.
因為你不是為上帝的緣故.
而去活那個生命.
這個是罪的本質.
你不相信.
你問一下撒旦你就知道.
撒旦不一定要你天下大亂.
撒旦亦都可以天下太平.
只要你不為上帝而活.
只要你繼續為你自己而活.
你就是在那個獨立的罪裡面.

$^{841}$你就是活在罪中.
你會死在罪中.
你有沒有想過.
為什麼這個世界裡面.
無論是好人壞人都會死.
因為原來無論好人壞人.
在上帝的眼光裡面都是罪人.
因為無論好人壞人.
在他們人生的規劃裡面.
上帝都不存在.
西人都犯了罪.
規缺了上帝的榮耀.
罪的功架.
乃是死.
還有一句沒有打出來的.
死後且有審判.
那個才令人害怕.
死後有紅紅的烈火.
等待著我們.
在整個聖經故事裡面.
耶穌是唯一一個不需要死的人.
因為祂是唯一一個為上帝而活的人.
祂是唯一一個為上帝榮耀而活的人.
耶穌是唯一一個不需要死的人.
耶穌是唯一一個不應該受審判的人.
然而耶穌來為我們死.
耶穌來為要幫我們接受審判.
為什麼.
楊福音告訴我們.
因為上帝愛世人.
甚至將祂的獨生子賜給他們.
叫一切信祂的人不致滅亡.
反得永生.
耶穌來替我們待死.
耶穌來拯救我們.
叫我們跟著祂從地下移民到天上.
叫我們過一個新的人生.
你不要想著.
給你一個免死金牌.
不是給你免死金牌.

$^{881}$不是給你保險單.
是要帶你一起去移民.
耶穌答應我們.
我們若願意跟著祂.
祂會給我們一個全新的人生.
我們要和祂一起離開這個世界.
為天國而活.
耶穌不願看見我們領受永恆地獄的審判.
於是耶穌語重心長.
重覆又重覆地警告我們.
第八章第二十一節.
你們會死在罪中.
第二十四節.
所以我說你們會死在罪中.
你們就會死在自己的罪中.
這真是金句中的金句.
你在《聖經》裡找不到耶穌.
三次跟你說你會死在罪中.
這裡第八章第二十一節的罪是單數.
我剛才跟你說過.
有趣的是第八章第二十四節的罪是複數.
所以第八章第二十四節裡面出現過兩個複數的罪.
我怎麼去理解這件事呢.
為什麼一個是單數一個是複數呢.
其實很簡單.
第八章第二十一節裡面單數的罪.
是代表了我們叛逆上帝.
以自己取締上帝的一個狀態.
是在說這個狀態.
八章二十四節裡面複數的罪.
是告訴我們當我們叛逆上帝.
以自己取締上帝之後所衍生出來的惡念和惡行.
一個是罪的根本.
一個是罪所生出來的惡念和惡行.
記得我們之前說十誡的時候.
我曾經說過十誡其實很簡單叫一臣九誡.
十誡其實只有一誡.
除我以愛便無他臣.
當你決定我不要上帝做我的上帝.
我要自己做上帝的時候.

$^{921}$你就進入了那個單數的罪.
接著因為你進入了單數的罪.
你就可以衍生之下九誡裡面所犯的罪.
從單數的罪變成複數的罪.
一臣九誡就好像太陽在裡面.
曾幾何時太陽是有九個行星現在只剩下八個.
九個行星是圍繞著太陽而行的.
整個十誡裡面的核心是第一誡.
除我以愛便無他臣.
不知道你還記不記得勝經裡面的十誡.
聖潔上帝頒佈給世界以色列人的十誡.
是哪十誡呢.
第一誡除我以愛便無他臣.
第二誡便無他將.
不可以做偶像.
不可以呼喊其他偶像的名字.
當紀念安息日.
當孝敬父母.
第一誡是所有誡命的總綱.
你犯了第一誡.
你就進入了《養父音》第八章第21節的單數的罪.
你進入了這個狀態.
於是其他的九誡反映了複數的罪.
就是因為我們離開了上帝之後.
我們所產生的惡念和罪行.
第一誡 肝胖第一誡就好像一個叛逆的母親.
其他九誡是叛逆母親所生出來的叛逆的兒女.
你也可以這樣說.
第一誡其實就是那個壞樹的樹根.
而其他的九誡是那個樹所生出來.
結出來的不同的壞果子.
第六誡至第十誡不可殺人,不可姦淫,不可偷渡,不可做假見證,不可貪戀.
有沒有留意第十誡和其他誡命很不一樣.
第十誡說不可貪戀的時候.
說的再不是外在的罪名.
說的是內在的東西.
第十誡告訴我們其實上帝所說的誡命.
其實不是單單跟你說外面.
而是跟你說裡面.
罪行是罪.

$^{961}$惡念都是罪.
上帝對罪的定義遠遠超過了地上法律的要求.
我們當中都不少律師.
你知道政府沒可能告你.
因為你有邪念.
只有你有罪行才算.
才可以告你.
但是上帝的國度裡面.
惡行是罪.
惡念都是罪.
所有你犯過的惡行,惡念.
上帝都記錄下來.
叫我們不目定口外.
恐懼賤經.
俯伏認罪.
我聽過你都記得耶穌在登山寶訓裡面.
曾經這樣說.
祂說你聽過古人說不可殺人.
但是我告訴你.
凡向弟兄動怒的.
必須要受審判犯.
罵弟兄為廢物的.
必須要受議會的審判犯.
罵弟兄是白癡的.
必須要遭地獄的火.
耶穌說不是你有沒有殺人.
是你心裡面有沒有恨惡你的弟兄.
接著下面那裡還絕.
你們聽說過不可姦淫.
然而我告訴你們.
凡看見婦人就動淫念的.
這人心裡已經和他犯了姦淫.
如果你有空睡不著覺的時候.
可以回家想一想.
數一下這十誡.
數到這十誡之後問自己.
你犯過多少條.
你可以再問自己.
每一條你犯了多少次.
我再三強調.

$^{1001}$身體犯罪.
心中犯罪.
都是犯罪.
上帝全部都見到.
上帝全部都記錄了下來.
今天我們周圍有很多的閉路電視.
或者你要犯罪的時候.
你要去一個地方沒有閉路電視.
然而在上帝的國裡面.
沒有一個沒有閉路電視的地方.
所有的東西全部被記錄下來.
每一條罪都是死罪.
每一條罪都可以將你.
送進地獄永遠不滅的烈火裡面.
那個烈火是上帝聖潔的憤怒.
我們相信的上帝.
聖經裡面啟示的上帝.
是一個聖潔的上帝.
所有的罪都是死罪.
於是聖潔的上帝.
律法給我們的警告.
地獄的烈火.
帶我們回到耶穌的身邊.
有人說信耶穌有什麼好處.
如果你不覺得上帝是聖潔的話.
如果你不覺得你干犯了上帝的律法.
如果你不覺得地獄的烈火是真實的.
耶穌真的跟你沒什麼幫助.
耶穌真的跟你沒什麼關係.
你都不需要耶穌.
然而如果你真的相信上帝.
是一個聖潔的上帝.
聖經裡面的律法給我們發出來的警告是真實的.
地獄的烈火是真實的話.
那我們絕對有理由要回到耶穌的身邊.
正如我之前所說.
耶穌是唯一不需要死的人.
耶穌是唯一不應該受審判的人.
然而耶穌來到這個世界.
為了要我們死.

$^{1041}$為我們去受審判.
耶穌不願意看見我們領受地獄.
永恆地獄的審判.
於是耶穌,我可以這麼說.
祂本著極大的憐憫.
祂抱著極大的愛心.
耶穌語重心長.
重複又重複地警告我們.
你們會死在自己的罪中.
祂不是咒你.
耶穌的敵人會咒他.
耶穌不咒我們.
耶穌看著我們就愛我們.
於是祂帶著極大的憐憫和愛心.
祂向我們說.
你會死在你的罪中.
那些人聽不聽得懂.
聽不懂.
第八章第25節.
他們就問他.
你到底是誰.
耶穌對他說.
我從起初就告訴你們了.
你到底是誰.
其實一個更好的翻譯是.
你算老幾.
你憑什麼挑戰我們.
他不是問耶穌是誰.
他是說你憑什麼這樣對我們說話.
耶穌對他說.
我從起初已經告訴你.
一年半前我已經告訴你.
我憑什麼這樣說.
一年半前在耶路撒冷裡面.
耶穌已經對你說.
我一副做事直到如今.
我也做事.
一年半前耶穌已經對他們說.
我憑著我是上帝的兒子的身份.
來到你們當中.

$^{1081}$耶穌有一個小名.
祂經常說.
我就是耶和華上帝所差來的.
但是你一想到.
耶和華上帝所差來的.
你會想起摩西.
你會想起先知.
而耶穌說.
我是耶和華上帝差來的.
不過我不是摩西.
我也不是先知.
我是上帝獨一的兒子.
上帝是我的父.
我是祂獨一的兒子.
上帝的名字叫做「我是」.
「I am」的名字也叫做「我是」.
「I am」因為我與上帝本父一原為一.
你到底是誰.
我是上帝的兒子.
我起初已經告訴你.
我就是上帝.
我來到要拯救你.
然而你不要被我拯救.
你只會死在你自己的罪中.
當他們聽到這一切的時候.
他們有什麼反應呢.
我相信他們的反應.
與一年半前的反應是一樣的.
就是他們越發想殺耶穌.
《楊福音》第五章第十八節.
一年半前的事情.
為了這個緣故.
猶太人越發想殺耶穌.
因為他稱上帝為他的父.
把自己和上帝看為同等.
一個這樣的人.
不可以留生路給他.
於是耶穌繼續說.
我有許多要講論你們 判斷你們的.
但猜我來的那位是真實的.

$^{1121}$我從他那裡所聽見的.
就告訴世人.
耶穌說猜我來的上帝.
祂是真實的.
祂的判斷是真實的.
而我的判斷是出於祂給我的判斷.
我的判斷也是真實的.
我說你們會死在罪中.
你們真的會死在罪中.
信不信由你.
但我告訴你們.
這是最真實的故事.
不知道你們有沒有想過.
在你人生中.
你聽過最真實的說話是什麼呢.
你很英俊.
你很有才華.
你以後一定什麼什麼.
還是你一定會死.
還是死後且有審判.
還是你一定死在罪中.
除非你迴轉信靠耶穌.
這群人不要相信.
第八章第27節.
他們被耶穌的終極回答.
他們不明白耶穌是對他們說過的事.
我覺得有兩個可能.
第一個可能就是.
這群猶太人和之前第五章第六章第七章的猶太人.
都是不同的群體.
可能有新人來了.
所以沒有聽過耶穌說.
所以聽不明白.
但我覺得更有可能的是.
他們不是不明白.
是選擇性失憶.
假裝不明白.
有說過嗎.
再聽聽吧.
今天先不相信.

$^{1161}$下次有機會再聽聽你說話.
他們到底是不明白還是不相信.
我們當中也有很多人聽了無數次.
聖經的講道.
福音你是明白的.
只不過你不願意接受.
不是明白不明白的問題.
是接不接受的問題.
是聽不聽得入耳的問題.
所以我們也要問我們自己.
我們對耶穌這個警告這個呼召.
又有什麼回應呢.
你們會死在自己的罪中.
耶穌給了我們一個有這麼清楚的警告.
我們會不明白.
不同意.
還是不聽從呢.
我們不斷說.
聖經裡面說的信耶穌.
不是買個保險.
聖經裡面所說的耶穌.
是有三重的意義.
Notitia, Assensus, and Fiducia.
Notitia是知識.
我們明白耶穌說的是什麼.
我們真的明白.
接著你問你自己是否同意.
Assensus是同意.
世人都犯了罪.
虧缺了上帝的榮耀.
罪的公家乃是死.
死後且有審判.
如果你不信我.
你不能.
你就一定死在罪中.
Notitia, Assensus.
但最後要聽從.
Fiducia.
要毀身.
對耶穌.

$^{1201}$我要跟隨你.
這是唯一可以逃離地獄的途徑.
我要跟隨你.
我要一生跟隨你.
如果我們不聽從耶穌.
如果我們要堅持自己作王.
我們堅持與上帝為敵的話.
我告訴你.
我們只好承受上帝對罪人.
所預備公平的審判.
請告訴我們.
當上帝送你到地獄的時候.
是絕對公平的審判.
不是上帝特別嚴謹.
不是我們特別不幸運.
而是因為罪的公家乃是死.
死後且有審判.
罪是不可想像.
不可饒恕的黑暗.
但耶穌想救我們.
我們回到第24節.
第21節這樣說.
你們會死在自己的罪中.
第24節又重覆.
所以我告訴你們.
你們會死在自己的罪中.
你們就是會死在自己的罪中.
但中間有一句說話.
你們若不信我就是.
換句說話來說.
我們信你就是就不用死.
是不是?.
這就是重點.
整個三文治.
三文治的中心就是.
你不用死.
為什麼一定要死?.
如果你信我就是.
你就不需要死.
信我就是就是什麼?.

$^{1241}$三點.
第一,我就是是上帝的名字.
耶穌跟猶太人說的話.
耶穌一直跟他們說.
我是父的兒子.
父差我來.
父的名字叫我是.
我的名字叫我是.
我與父本為一.
你要信這件事.
你要信我就是天父的兒子.
信耶穌就是信祂是上帝.
第二.
你要相信我就是生命的活水.
我就是世界的光.
這是耶穌在主彭次裡面所啟示的.
可能你已經不記得三文治之前說過什麼.
什麼水禮,火禮.
然後水禮當中.
耶穌說我是生命的活水.
在火禮裡面.
耶穌說我是世界的光.
你忘記了就回去聽.
但是我怎麼可以不死在罪中?.
我怎麼可以得救?.
我要明白什麼叫做生命的活水.
我要明白什麼叫做世界的光.
如果你說我都不太記得.
不要緊.
我幫你總結.
因為無論生命活水.
或者世界的光.
都需要兩樣東西.
第一是榮耀的十甲.
第二是聖靈的降臨.
第七章第三十九節.
耶穌不可以成為生命的活水.
耶穌不可以成為世界的光.
除非兩樣東西發生.
第一祂要被升上十字架.

$^{1281}$第二在祂被高舉的十字架上.
祂得榮耀的時候.
聖靈會降臨.
就是生命活水的開始.
所以今天我們要逃離地獄的審判.
我們要逃離死在罪中.
我們需要兩樣東西.
我們需要耶穌基督的十字架.
我們需要聖靈來到我們當中.
順我的人從他的福中將流出活水的江河來.
第八章第十二節.
我就是世界的光.
跟隨我的必不在黑暗裡走.
卻要得著生命的光.
我們若不要死在自己的罪中.
我們要來到耶穌的身邊.
來,來喝耶穌.
來信耶穌.
當然我剛才已經說了.
Notitia, Ascensis, Volitia.
要跟隨祂.
來學祂,來信祂,來跟隨祂,來追光.
我們要仰望十字架.
我們要領受聖靈的洗禮.
第八章第二十八節.
耶穌的編度要完了.
他們還聽不明白.
但是耶穌想他們走之前.
再給多一個最後的提示.
耶穌說.
你聽不明白不要緊.
你記得第二十七節.
他們說我聽不明白.
不知道你在說什麼.
耶穌說OK不要緊.
我給你一個提示.
到時考試的時候拿出來看.
記得拿出來看.
那個提示在哪裡.
有一天我會被舉起來.

$^{1321}$那天你再想一想.
耶穌說你們舉起人質以後就會知道.
我就是那位.
你現在聽不明白.
或者你故意不明白.
你不願意接受.
你選擇性失憶.
你可不可以記住這個提示.
到時大考出來的時候.
你做不出那條題目.
你要肥佬的時候.
你要死在最終的時候.
記得這句話.
你會看到人質被舉起來.
那一刻你就可以知道.
猜我來的那位與我同在.
我們在行祂所喜悅的事.
祂的喜悅事是什麼.
拯救世人.
於是耶穌說.
如果你到時仰望我的十字架.
你就會知道.
你就會曉得.
你就會重生.
你就會相信.
你就會明白.
而果真在聽這篇道裡面.
其中有一個人叫尼哥德穆.
如果你回到第七章尾的時候.
你會發現尼哥德穆在人群當中.
他有留心聽耶穌的講道.
他拿了耶穌的筆記.
他那天三章裡面.
黑夜來找耶穌.
耶穌都跟他說.
當摩西舉蛇的時候.
人質被舉起來.
這次耶穌出題.
接著我們聽不明白的時候.
耶穌再給我們提示.

$^{1361}$都是講同一樣東西.
於是尼哥德穆跟著耶穌去到十字架.
他終於明白.
終於知道.
第十九章第三十四節.
士兵拿槍擲他立房的時候.
有血和水聲令來了.
救贖來了.
尼哥德穆也來了.
他看著耶穌被舉起來.
帶著一百斤的活藥和沉香.
他們用西麻布.
他用藥塞.
加上了香料.
將耶穌的身體補好.
我們重複又重複.
講了這個故事好幾次.
因為要害每一個故事.
最後所有故事.
都要帶我們去尼哥德穆的重生.
他明白了.
他得到重生了.
他相信了.
他知道了.
於是他離開黑暗.
他追隨震光.
他看到十字架.
他奉上他所有.
一百斤.
七十五年薪水買回來的.
這個人很搭水.
幾百萬送給耶穌.
他親手幫耶穌埋葬.
他拋棄了他所有的東西.
他開始追求一個不一樣的國度.
你明白嗎.
你明白發生了什麼事嗎.
尼哥德穆從前是地下人.
現在變了上了天.
變了天上人.

$^{1401}$尼哥德穆從前是屬於這個世界.
他現在拋棄了這個世界.
跟隨了耶穌.
你見不見到.
以前是地下現在是天上.
以前是死在罪中.
現在活在耶穌基督裡面.
被聖靈所重生.
我們要結束了.
結束的時候我想用一張slide.
五個point.
我們總結耶穌在耶路撒冷.
聖殿宣講朱鵬節的報道訊息.
整個朱鵬節報道訊息的中心是什麼.
十字架和聖靈.
若你跟著我們一直聽著約翰福音.
你知道我們走不掉的了.
約翰不斷重複講十字架.
不斷重複講聖靈.
生命的活水.
世界的光會從哪裡來.
都是從十字架和聖靈來.
如果你仰望十字架.
如果你等待聖靈.
被聖靈所駛.
以下你會發生什麼事.
饑渴來到耶穌的根前.
渴耶穌信基督得清兒.
下面離開黑暗.
隨真光求天國得生命.
上面AA.
不再做罪的奴僕.
死在罪中.
單數.
不會死在自己罪中.
單數的罪.
因為單數的罪.
代表我們是罪的奴僕.
我們再不是罪的奴僕.
不屬世界.

$^{1441}$不再追求地上的榮耀.
求主幫助我們.
我們所有人一起去回應.
耶穌對我們的呼召.
怎樣回應.
仰望耶穌的十字架.
領受聖靈的洗禮.
盡心追求上帝的國.
上帝的義.
這個就是耶穌呼召我們.
跟隨祂最後的目的.
讓我們和祂一起得著生命.
一起祈禱.
我們是天父.
我們需要你幫助我們.
我們已經讀了《楊福音》第七章.
很久很久.
讀到第七章.
讀到第八章.
然後告訴我們.
原來第八章也是第七章的一部分.
我們看這個故事.
已經重複又重複看了很多次.
朱鵬哲.
接著生命的活水.
世界的光.
接著耶穌再加多一篇.
報道會的道.
接著下個禮拜還回來.
還有夏文.
很多人信祂.
結果又不信祂.
主人我們都很害怕.
我們很想信你.
但我們又害怕我們的信是假信.
我們多謝你.
約翰將所有的東西寫下來.
要告訴我們.
何謂信耶穌.
但是這裡.

$^{1481}$在這十節九節的聖經裡.
我們看到耶穌發狂地警告.
祂眼前的罪人.
其實這群人都不是什麼好人.
最後究竟有多少人信.
我也不知道.
這群人在詛咒祂.
這群人在譏笑祂.
但是耶穌卻對著他們.
不斷狂呼喊.
你們若不聽.
你必死在你的罪中.
死在你的罪中.
而耶穌告訴我們.
祂不想我們死在罪中.
祂要我們和祂一起得生.
而祂為我們承傳了這個救贖.
求主你幫助我們.
仰望祂的十字架.
等待渴慕聖靈.
使我們重生.
讓我們不再是地上的人.
讓我們是天上的人.
讓我們不再屬於這個世界.
讓我們完全屬於你的天國.
求主你幫助我們.
豁出一個完全不一樣的人生.
多謝你.
我們敢祈禱.
奉主耶穌基督名求.
\newpage



\section{民數記 13:1-14:45}
\label{sec:032272}
\textbf{信不信神應許的結果}~~~~林誠信牧師
\newline
\newline
連結: \href{http://www.church.com.hk/acms/content.asp?site=cdc\&op=show\&type=product\&code=032272\&layout=sermon}{\texttt{032272}} ~~~~ 語音日期: 2026-06-07
\newline
\newline
\hyperref[sec:032271]{\small{< < < PREV SERMON < < <}}
~
\hyperref[sec:index_chronic]{\small{[返順時目]}}
~
\hyperref[sec:032273]{\small{> > > NEXT SERMON > > >}}
\newline
\newline
民數記 13:1-14:45
\newline
\begin{longtable}{cl}
\hline
\hline
章節 & 經文 (和合本修訂版)\\
\hline
13:1 & \begin{tabularx}{0.7\textwidth}{X} 耶和華吩咐摩西說： \end{tabularx} \\ \\ \relax
13:2 & \begin{tabularx}{0.7\textwidth}{X} 「你要派人去窺探我所賜給以色列人的迦南地；每個父系支派要派一個人，是他們中間的族長。」 \end{tabularx} \\ \\ \relax
13:3 & \begin{tabularx}{0.7\textwidth}{X} 摩西就遵照耶和華的指示，從巴蘭曠野差派他們去；他們都是以色列人中的領袖。 \end{tabularx} \\ \\ \relax
13:4 & \begin{tabularx}{0.7\textwidth}{X} 這是他們的名字：屬呂便支派的，撒刻的兒子沙母亞。 \end{tabularx} \\ \\ \relax
13:5 & \begin{tabularx}{0.7\textwidth}{X} 屬西緬支派的，何利的兒子沙法。 \end{tabularx} \\ \\ \relax
13:6 & \begin{tabularx}{0.7\textwidth}{X} 屬猶大支派的，耶孚尼的兒子迦勒。 \end{tabularx} \\ \\ \relax
13:7 & \begin{tabularx}{0.7\textwidth}{X} 屬以薩迦支派的，約色的兒子以迦。 \end{tabularx} \\ \\ \relax
13:8 & \begin{tabularx}{0.7\textwidth}{X} 屬以法蓮支派的，嫩的兒子何西阿。 \end{tabularx} \\ \\ \relax
13:9 & \begin{tabularx}{0.7\textwidth}{X} 屬便雅憫支派的，拉孚的兒子帕提。 \end{tabularx} \\ \\ \relax
13:10 & \begin{tabularx}{0.7\textwidth}{X} 屬西布倫支派的，梭底的兒子迦疊。 \end{tabularx} \\ \\ \relax
13:11 & \begin{tabularx}{0.7\textwidth}{X} 屬約瑟支派，就是瑪拿西支派的，穌西的兒子迦底。 \end{tabularx} \\ \\ \relax
13:12 & \begin{tabularx}{0.7\textwidth}{X} 屬但支派的，基瑪利的兒子亞米利。 \end{tabularx} \\ \\ \relax
13:13 & \begin{tabularx}{0.7\textwidth}{X} 屬亞設支派的，米迦勒的兒子西帖。 \end{tabularx} \\ \\ \relax
13:14 & \begin{tabularx}{0.7\textwidth}{X} 屬拿弗他利支派的，縛西的兒子拿比。 \end{tabularx} \\ \\ \relax
13:15 & \begin{tabularx}{0.7\textwidth}{X} 屬迦得支派的，瑪基的兒子臼利。 \end{tabularx} \\ \\ \relax
13:16 & \begin{tabularx}{0.7\textwidth}{X} 這些是摩西差派去窺探那地之人的名字。摩西叫嫩的兒子何西阿為約書亞。 \end{tabularx} \\ \\ \relax
13:17 & \begin{tabularx}{0.7\textwidth}{X} 摩西差派他們去窺探迦南地，對他們說：「你們上到尼革夫那裡，上到山區去， \end{tabularx} \\ \\ \relax
13:18 & \begin{tabularx}{0.7\textwidth}{X} 看看那地如何：住那裡的百姓是強是弱，是多是少， \end{tabularx} \\ \\ \relax
13:19 & \begin{tabularx}{0.7\textwidth}{X} 他們所住的地是好是壞，所住的城鎮是營地還是堡壘， \end{tabularx} \\ \\ \relax
13:20 & \begin{tabularx}{0.7\textwidth}{X} 那地是肥沃還是貧瘠，當中有樹木沒有。你們要放膽，把那地的果子帶些回來。」那時正是葡萄初熟的季節。 \end{tabularx} \\ \\ \relax
13:21 & \begin{tabularx}{0.7\textwidth}{X} 他們上去窺探那地，從尋的曠野到利合，直到哈馬口。 \end{tabularx} \\ \\ \relax
13:22 & \begin{tabularx}{0.7\textwidth}{X} 他們從尼革夫上去，到了希伯崙。在那裡有亞衲族的亞希幔人、示篩人和撻買人。希伯崙的建造比埃及的瑣安早七年。 \end{tabularx} \\ \\ \relax
13:23 & \begin{tabularx}{0.7\textwidth}{X} 他們到了以實各谷，從那裡砍下葡萄樹枝，上面有一掛葡萄，兩個人用槓抬著，又帶了一些石榴和無花果。 \end{tabularx} \\ \\ \relax
13:24 & \begin{tabularx}{0.7\textwidth}{X} 以色列人從那裡砍下一掛葡萄，所以那地方就叫以實各谷。 \end{tabularx} \\ \\ \relax
13:25 & \begin{tabularx}{0.7\textwidth}{X} 他們窺探那地四十天之後，就回來了。 \end{tabularx} \\ \\ \relax
13:26 & \begin{tabularx}{0.7\textwidth}{X} 他們來到巴蘭曠野的加低斯，摩西、亞倫，以及以色列全會眾那裡，向他們和全會眾報告，又把那地的果子給他們看。 \end{tabularx} \\ \\ \relax
13:27 & \begin{tabularx}{0.7\textwidth}{X} 他們告訴摩西說：「我們到了你派我們去的那地，果然是流奶與蜜之地；這就是那地的果子。 \end{tabularx} \\ \\ \relax
13:28 & \begin{tabularx}{0.7\textwidth}{X} 但是住那地的百姓很強悍，城鎮又大又堅固，我們也在那裡看見了亞衲族人。 \end{tabularx} \\ \\ \relax
13:29 & \begin{tabularx}{0.7\textwidth}{X} 亞瑪力人住在尼革夫；赫人、耶布斯人和亞摩利人住在山區；迦南人住在沿海一帶和約旦河旁。」 \end{tabularx} \\ \\ \relax
13:30 & \begin{tabularx}{0.7\textwidth}{X} 迦勒在摩西面前安撫百姓，說：「我們立刻上去得那地吧！我們必能征服它。」 \end{tabularx} \\ \\ \relax
13:31 & \begin{tabularx}{0.7\textwidth}{X} 但那些和他同去的人卻說：「我們不能上去攻打那些百姓，因為他們比我們強大。」 \end{tabularx} \\ \\ \relax
13:32 & \begin{tabularx}{0.7\textwidth}{X} 於是探子中有人向以色列人散佈有關所窺探之地的謠言，說：「我們所走過、所窺探之地是吞沒居民之地，並且我們在那裡所看見的百姓都身材高大。 \end{tabularx} \\ \\ \relax
13:33 & \begin{tabularx}{0.7\textwidth}{X} 我們在那裡看見巨人，就是巨人中的亞衲族人。我們在自己眼中像蚱蜢一樣，而在他們眼中，我們也確是這樣。」 \end{tabularx} \\ \\ \relax
14:1 & \begin{tabularx}{0.7\textwidth}{X} 全會眾大聲喧嚷，那夜百姓哭號。 \end{tabularx} \\ \\ \relax
14:2 & \begin{tabularx}{0.7\textwidth}{X} 以色列眾人向摩西和亞倫發怨言，全會眾對他們說：「我們寧願死在埃及地，寧願死在這曠野！ \end{tabularx} \\ \\ \relax
14:3 & \begin{tabularx}{0.7\textwidth}{X} 耶和華為甚麼要把我們領到那地，讓我們倒在刀下呢？我們的妻子和孩子必成為擄物。我們回埃及去豈不更好嗎？」 \end{tabularx} \\ \\ \relax
14:4 & \begin{tabularx}{0.7\textwidth}{X} 他們彼此說：「我們不如選一個領袖，回埃及去吧！」 \end{tabularx} \\ \\ \relax
14:5 & \begin{tabularx}{0.7\textwidth}{X} 摩西和亞倫在以色列全會眾面前臉伏於地。 \end{tabularx} \\ \\ \relax
14:6 & \begin{tabularx}{0.7\textwidth}{X} 窺探那地的人中，嫩的兒子約書亞和耶孚尼的兒子迦勒撕裂衣服， \end{tabularx} \\ \\ \relax
14:7 & \begin{tabularx}{0.7\textwidth}{X} 對以色列全會眾說：「我們所走過、所窺探之地是極美之地。 \end{tabularx} \\ \\ \relax
14:8 & \begin{tabularx}{0.7\textwidth}{X} 耶和華若喜愛我們，就必領我們進入那地，把這流奶與蜜之地賜給我們。 \end{tabularx} \\ \\ \relax
14:9 & \begin{tabularx}{0.7\textwidth}{X} 但你們不可背叛耶和華，也不要怕那地的百姓，因為他們是我們的食物。保護他們的已經離開他們，耶和華卻與我們同在。不要怕他們！」 \end{tabularx} \\ \\ \relax
14:10 & \begin{tabularx}{0.7\textwidth}{X} 當全會眾正說著要拿石頭打死他們的時候，耶和華的榮光在會幕中向以色列眾人顯現。 \end{tabularx} \\ \\ \relax
14:11 & \begin{tabularx}{0.7\textwidth}{X} 耶和華對摩西說：「這百姓藐視我要到幾時呢？我在他們中間行了這一切神蹟，他們還不信我要到幾時呢？ \end{tabularx} \\ \\ \relax
14:12 & \begin{tabularx}{0.7\textwidth}{X} 我要用瘟疫擊殺他們，使他們不得承受那地。我要使你成為大國，比他們強大。」 \end{tabularx} \\ \\ \relax
14:13 & \begin{tabularx}{0.7\textwidth}{X} 摩西對耶和華說：「埃及人必聽見，因你曾施展大能，領這百姓從他們中間出來。 \end{tabularx} \\ \\ \relax
14:14 & \begin{tabularx}{0.7\textwidth}{X} 埃及人要告訴這地的居民，他們已經聽見你---耶和華是在這百姓中間，因為你---耶和華面對面顯示自己，你的雲彩停在他們以上。你日間在雲柱中，夜間在火柱中，在他們的前面行。 \end{tabularx} \\ \\ \relax
14:15 & \begin{tabularx}{0.7\textwidth}{X} 你若把這百姓殺了，好像殺一個人那樣，那聽見你名聲的列國必說： \end{tabularx} \\ \\ \relax
14:16 & \begin{tabularx}{0.7\textwidth}{X} 『耶和華因為不能把這百姓領進他向他們起誓應許之地，所以在曠野把他們殺了。』 \end{tabularx} \\ \\ \relax
14:17 & \begin{tabularx}{0.7\textwidth}{X} 現在求主顯出大能，照你說過的話說： \end{tabularx} \\ \\ \relax
14:18 & \begin{tabularx}{0.7\textwidth}{X} 『耶和華不輕易發怒，且有豐盛的慈愛。他赦免罪孽和過犯，萬不以有罪的為無罪，必懲罰人的罪，從父到子，直到三、四代。』 \end{tabularx} \\ \\ \relax
14:19 & \begin{tabularx}{0.7\textwidth}{X} 求你照你的大慈愛赦免這百姓的罪孽，好像你從埃及到如今饒恕這百姓一樣。」 \end{tabularx} \\ \\ \relax
14:20 & \begin{tabularx}{0.7\textwidth}{X} 耶和華說：「我照著你的話赦免他們。 \end{tabularx} \\ \\ \relax
14:21 & \begin{tabularx}{0.7\textwidth}{X} 然而，我指著我的永生與遍地充滿了耶和華的榮耀起誓： \end{tabularx} \\ \\ \relax
14:22 & \begin{tabularx}{0.7\textwidth}{X} 這些人雖然都看過我的榮耀和我在埃及與曠野所行的神蹟，仍然這十次試探我，不聽從我的話， \end{tabularx} \\ \\ \relax
14:23 & \begin{tabularx}{0.7\textwidth}{X} 他們絕不能看見我向他們祖宗所起誓應許之地；凡藐視我的，一個也不得看見。 \end{tabularx} \\ \\ \relax
14:24 & \begin{tabularx}{0.7\textwidth}{X} 惟獨我的僕人迦勒，因他另有一個心志，專心跟從我，我要領他進入他所去過的那地；他的後裔必得那地為業。 \end{tabularx} \\ \\ \relax
14:25 & \begin{tabularx}{0.7\textwidth}{X} 亞瑪力人和迦南人住在谷中，明天你們要轉回去，沿著紅海的路往曠野去。」 \end{tabularx} \\ \\ \relax
14:26 & \begin{tabularx}{0.7\textwidth}{X} 耶和華對摩西和亞倫說： \end{tabularx} \\ \\ \relax
14:27 & \begin{tabularx}{0.7\textwidth}{X} 「這邪惡的會眾向我發怨言要到幾時呢？以色列人向我發的怨言，我都聽見了。 \end{tabularx} \\ \\ \relax
14:28 & \begin{tabularx}{0.7\textwidth}{X} 你要告訴他們，耶和華說：『我指著我的永生起誓，我必照你們在我耳中所說的待你們。 \end{tabularx} \\ \\ \relax
14:29 & \begin{tabularx}{0.7\textwidth}{X} 你們的屍體必倒在這曠野中。你們中間被數點，凡二十歲以上向我發怨言的， \end{tabularx} \\ \\ \relax
14:30 & \begin{tabularx}{0.7\textwidth}{X} 必不得進我所起誓應許給你們居住的那地。惟有耶孚尼的兒子迦勒和嫩的兒子約書亞才能進去。 \end{tabularx} \\ \\ \relax
14:31 & \begin{tabularx}{0.7\textwidth}{X} 你們的孩子，就是你們說要成為擄物的，我必領他們進去，他們就得知你們所厭棄的那地。 \end{tabularx} \\ \\ \relax
14:32 & \begin{tabularx}{0.7\textwidth}{X} 至於你們，你們的屍體必倒在這曠野中； \end{tabularx} \\ \\ \relax
14:33 & \begin{tabularx}{0.7\textwidth}{X} 你們的兒女必在曠野遊牧四十年，擔當你們不信的罪，直到你們的屍體在曠野消滅為止。 \end{tabularx} \\ \\ \relax
14:34 & \begin{tabularx}{0.7\textwidth}{X} 按你們窺探那地的四十日，一年抵一日，你們要擔當你們的罪孽四十年，你們就知道我疏遠你們了。』 \end{tabularx} \\ \\ \relax
14:35 & \begin{tabularx}{0.7\textwidth}{X} 我－耶和華說過，我必這樣對待這一切聚集對抗我的邪惡會眾。他們必在這曠野中消滅，死在這裡。」 \end{tabularx} \\ \\ \relax
14:36 & \begin{tabularx}{0.7\textwidth}{X} 摩西所差派去窺探那地的人回來，散佈有關那地的謠言，使全會眾向摩西發怨言， \end{tabularx} \\ \\ \relax
14:37 & \begin{tabularx}{0.7\textwidth}{X} 這些散佈謠言的人都遭受瘟疫，死在耶和華面前。 \end{tabularx} \\ \\ \relax
14:38 & \begin{tabularx}{0.7\textwidth}{X} 窺探那地的人中，惟有嫩的兒子約書亞和耶孚尼的兒子迦勒得以存活。 \end{tabularx} \\ \\ \relax
14:39 & \begin{tabularx}{0.7\textwidth}{X} 摩西把這些話告訴以色列眾人，他們都極其悲哀。 \end{tabularx} \\ \\ \relax
14:40 & \begin{tabularx}{0.7\textwidth}{X} 他們清晨起來，上到山頂，說：「看哪，我們要上到耶和華所說的地方；因為我們犯了罪。」 \end{tabularx} \\ \\ \relax
14:41 & \begin{tabularx}{0.7\textwidth}{X} 摩西說：「你們為何要這樣違背耶和華的指示呢？這事必不能順利。 \end{tabularx} \\ \\ \relax
14:42 & \begin{tabularx}{0.7\textwidth}{X} 不要上去，因為耶和華不在你們中間，恐怕你們在仇敵面前被擊敗。 \end{tabularx} \\ \\ \relax
14:43 & \begin{tabularx}{0.7\textwidth}{X} 亞瑪力人和迦南人都在你們面前，你們必倒在刀下。因為你們背離不跟從耶和華，耶和華必不與你們同在。」 \end{tabularx} \\ \\ \relax
14:44 & \begin{tabularx}{0.7\textwidth}{X} 他們卻擅自上到山頂。但耶和華的約櫃和摩西都沒有離開營地。 \end{tabularx} \\ \\ \relax
14:45 & \begin{tabularx}{0.7\textwidth}{X} 於是亞瑪力人和住在那山區的迦南人下來，擊敗他們，追擊他們直到何珥瑪。 \end{tabularx} \\ \\
[1ex]
\hline
\hline
\end{longtable}
$^{1}$等姐妹早晨.
今日跟大家繼續思考民數記給我們的訊息.
今早我回來講九點鐘的崇拜.
之前發生了一個小小的插曲.
我明明記得我帶了手提電話回來.
跟他們祈願土 之後回到辦公室.
拿了講章之類的東西.
發覺到處都不見了那個手提電話.
同工又幫我在洗手間找 前台 到處都找不到.
接著又打電話給我的手機.
但我跟他們說我的手機應該調校了靜音.
因為崇拜 調校了靜音.
真是糟糕了 怎麼辦呢.
當時又要進來講道.
所以懷了一個很高度的焦慮.
但心裡默默的祈禱.
在辦公室一直祈禱 祈禱出來.
求主憐憫 因為我今天整天都有聚會.
直到晚上.
沒有手機怎麼辦呢.
現在沒有手機比沒有錢包還慘.
因為所有的資訊.
人家如果找你又找不到.
又不知道你去了哪裡.
這樣就很大件事.
這種的焦慮.
我覺得和我今天要講的經文.
也有一個很大的關連性.
其實我們信不信得過上帝的應許呢.
聖經有沒有應許告訴我們.
祈禱就得著 尋找就尋見.
叩門就給你門開門.
如果上帝的應許這麼真實.
你今天怎樣向神去依靠.
有時上帝給我們的回覆.
是臨時的.
好像是掉進水裡.
救我主耶穌 會不會說.
讓你多喝幾口鹹水.
我們看其實整件事.

$^{41}$當時我心裡頭.
雖然是焦慮.
但是我求神真的幫我.
再找我 因為我不是第一次.
留下電話.
我真的在教會洗手間也留下過.
去七仔買東西也留下過.
放手機在別人的桌子上.
給完八達通卡.
很安然的就離開了.
為什麼會這樣呢.
原來呢 我發覺原來都是.
人年紀大了可能是比較笨拙.
我們仍然是要繼續依靠神.
也是最近看到一個報告.
是香港心理衛生會.
三月 很新的.
全港抑鬱和焦慮的指數調查.
港人的情緒健康正處於關鍵時刻.
歷年的新高.
屬於中度至重度的焦慮受訪者比例.
顯著上升了22.6\%.
上升至22.6\%.
即是說每4至5個人就有一個人正默默承受著嚴重的焦慮.
在去年25年的時候.
香港樹人大學做了一個四代香港市民研究.
也都指出Z世代.
是說那些18到27歲的焦慮比率.
是高達30.5\%.
是所有年齡層壓力最大的群體.
大家有沒有想過.
為什麼今天我們年輕人的焦慮這麼大.
因為我相信.
可能就是他們面對讀書壓力.
戀愛的壓力.
或者要結婚.
結婚你說有什麼壓力.
我有啊.
現在的人結婚.
以前就說要買樓.

$^{81}$現在的年輕一輩就不一定要買樓.
他說躺平了現在.
所以我們看到他有不同的焦慮.
甚至現在AI是整個浪潮.
在接下來的5年間是很明顯的.
是我們整個世界的社會都會有改變.
有很多的統計都告訴我們.
AI是會取代很多的工作.
但同時間AI是會產生一些新的工作.
問題就在這裡.
你說沒有問題.
每一個時代都會變.
失業的就去找一些新的工作就行了.
講就很容易.
在數據來講.
你可以講得通.
但在人來講我們不是機器.
你可能做的工作.
要你去做另外一些工作.
其實是有壓力的.
有很多不同的困難.
我們看到今天你是為了明天的幽女.
為了你的家庭.
為了你的兒女.
或者為了你要成立的家庭.
你的幽女.
今天你信不信得過上帝.
就是民數記13到14章.
你信神的什麼呢.
是神的應許.
這兩章的聖經給我們一個很清楚的思路.
你不只是信有神.
你是信神的每一句話.
神的吩咐.
神的命令背後是他的應許.
他與你同在.
要帶領你進去得著這個迦南美地.
但問題是你信心要有行動.
你不可以只是說我信.
你一點行動都沒有.

$^{121}$你用行動去領受上帝的恩典.
但你不相信的時候.
你就變成很多的焦慮.
很多的懼怕.
很多的埋怨.
你就錯失神的恩典.
我們一起來看看.
13章開場的情況是怎樣呢.
我們一起來讀.
耶和華吩咐摩西說.
你要派人去窺探.
我所賜給以色列人的迦南地.
每個父系之派要派一個人.
是他們中間的族長.
我們在這裡可以看到.
聖經是很清楚告訴我們.
是神賜的給以色列.
這個是明明上帝說.
是給過你們的.
給過你們.
但聖經也在這裡說.
耶和華吩咐摩西.
你要派人去窺探.
去視察.
他們沒有去過那個地方.
等於你今天去買一層樓.
你不只是看圖跡.
起碼也看看示範單位.
或者起碼也看到電腦的模型.
裡面好像進去了.
昔日他們沒有這些東西.
但他也要實地去看一看.
總之問題來了.
在申命記裡面.
當摩西複述這段歷史的時候.
是給那些新生代.
那一輩那些年輕一輩.
那些是以色列人.
唯有他們才能夠進入這個迦勒美地.
我們這兩張聖經會交代.

$^{161}$為什麼這些新生代的爸爸媽媽.
個個都進不了呢.
除了約書亞和迦勒.
就是因為他們不信神的應許.
他們是十二個探子去.
十個探子都說不行.
結果他們就去支持那十個.
報這個壞消息的.
這十個探子.
我們看到究竟是哪一個對呢.
哪個版本對呢.
其實兩個版本都對的.
其實聖經有時不是衝突的一些地方.
不是前後矛盾.
其實很肯定就是說那件事.
如果是以色列人.
他們說我們要去得地為業.
我們都要先看一下.
他們來到要求摩西.
但摩西作為他們的領袖.
摩西不會就這樣就去答應他們.
摩西是凡事都去求問上帝.
而上帝也都吩咐摩西.
你就派人去.
所以兩者沒有衝突的.
沒有矛盾的.
我們去看看這十二個探子的名字.
我們看到有兩個特別highlight出來.
就是迦勒和約書亞.
唯獨這兩個探子.
他們是用信心去扣得上神的應許.
他們就是可以作出一個.
是完全正確的一個報告.
這個報告是什麼呢.
這個報告帶來是可以.
令人是很雀躍還是很沮喪.
同一個事件.
原來是很視乎那個人的眼光和解讀.
帶出來是有不同的一個結論.
摩西叫探子去做什麼呢.

$^{201}$我們又看看下文.
由十七節到二十節.
其實是一個線型的結構.
在文學上就是前呼後應.
好像吃漢堡包吃到最裡面那層.
那個才是餡.
頭尾是說你們要去這個藍地.
去到山地去.
由藍邊進入這個迦藍地.
而去到那些山地.
你要去看看那些地是怎樣.
你還要拿一些果子回來.
拿一些果子回來.
其中那些居民是怎樣呢.
是多是少是強是弱呢.
還有那個土地是肥美還是貧脊呢.
有沒有樹木的呢.
叫他們做完這些考察.
好了.
中間最核心的是什麼.
上帝說這個是應許之地.
是留乃與物之地.
所住之處是好事大.
你們說吧.
你們翻出來看看那個情況是怎樣.
探子的行蹤是由21節到26節.
下一個地圖可以幫一下大家.
我們在藍邊就是這個迦底斯巴尼亞.
是近著這個尋的曠野.
他們在那裡進入去.
經過希伯倫耳實角谷.
然後他們是會一路北上.
去到這個利俠和哈瑪口.
哈瑪口就是最北面那個地方.
是他們以色列舊約的時候.
都是以那裡做一個關口.
當他們去到那裡而視察所有.
所以其實基本上是由南到北.
去看完這個迦南地.
他們看完那些山地.

$^{241}$看完那些果子.
看完那些居民.
然後他們才折返的.
整個行程他們頭尾40天.
探子回來的報告是什麼呢.
就這樣在第30節.
在13章30節.
迦納的報告.
我們立刻上去得地吧.
我們必能夠征服他.
迦納為什麼這麼有信心.
因為他知道這個是神所應許的.
我們不用理會其他的情況.
這個是神給予我們的.
我們一定得著.
問題在這裡了.
為什麼只有他一個人出聲.
12個探子.
那約書亞為什麼沉默呢.
在這個時候.
這個其實有些你要明白.
約書亞他是12個探子之一.
但約書亞當時也是怎樣.
他是摩西的副手.
他是摩西親自提拔.
揀選要做他的幫手.
約書亞完全明白摩西那個心意.
他是跟隨著摩西.
他的立場是完全一致的.
他知道上帝吩咐摩西.
而他們這個地是必然會得著.
所以他基本上的立場是迦納一樣.
但這個時候他不需要出聲.
因為如果他出聲的時候.
就好像跟他的利益有衝突了.
是不是我偏幫呢.
好了.
接著你看那10個探子.
他們的報告是什麼呢.
31節到33節我們一起讀.

$^{281}$但那些和他同去的人卻說.
我們不能上去攻打那些百姓.
因為他們比我們強大.
於試探子中有人向以色列人.
散佈有關所窺探之地的謠言.
說我們所走過所窺探之地.
是吞沒居民之地.
並且我們在那裡所看見的百姓.
都身材高大.
我們在那裡看見巨人.
就是巨人中的迦納族人.
我們在自己眼中像蚱蜢一樣.
而在他們眼中我們也確是一樣.
這裡他們報告了什麼呢.
報告就是那個地是怎樣.
那裡的人是怎樣.
他們都拿了那些果子回來.
但是聖經很清楚說.
他們是實地視察過的.
看過的.
但是他們所說的那些是謠言.
什麼是謠言呢.
謠言就是有些可能是.
一些觀察到的東西.
是事實.
但是將那些事實放在一起.
拼在一起.
可以扭曲了那個全面的真相.
那個圖畫.
這個扭曲是很嚴重的.
這個扭曲是誤導了所有的百姓.
他們就是要游說百姓.
不要進入這個迦納地.
因為他看眼前那個環境就是這樣.
大家你們有沒有試過.
誤信一些所謂的謠言.
試過沒有.
為什麼會信呢.
因為現實.
特別現在AI有圖有片.

$^{321}$所謂有真相其實是沒有真相.
這件事沒有AI的時候都在發生.
你有沒有留意過.
有時你看.
我不開名說.
有些你在便利店買到一些雜誌.
每個星期都會出一本.
每個星期都會說一些歌星明星.
說他的寵物.
他彈一張圖片出來.
或者封面的.
很出名啊.
這個女明星啊.
啊 拍到她.
就和這個年少的.
這個富家子弟.
富家子弟結了婚.
不過原來那張圖片.
是富家子弟旁邊都有他老婆.
不過他就斬件.
將那個男的和那個女明星合照.
這樣就可以大做文章了.
我們有時一見到這些情況.
你信的話.
你就會對那個人的印象大打折扣.
怎樣怎樣.
可能會罵爆他.
我們有時就是不信一些謠言.
可以將整個的所謂真相.
就埋沒了甚至扭曲了.
這件事很危險極度危險.
人的信心是否接得上神的應許呢.
當我們離開了神的應許.
人的眼光就完全看不到.
上帝要我們看見的.
他所說.
這個是流奶與物之地.
他們看到的不是流奶與物之地.
是會吞沒我們居民的地方.
為何他們會有這樣的一個印象呢.

$^{361}$因為那裡的居民是身材高大的巨人.
但這個是有前後矛盾的.
他們自己的報告裡面說.
是不少人居住在那個地方.
這個是很好的地圖.
地圖肥沃很多的果子.
流奶與物這兩個字.
你要知道是以色列裡頭.
三個重要的農業和畜牧的一個出產.
這些奶就是羊奶.
他們養很多羊.
蜜不是蜂蜜.
是那些蜜棗.
我去過以色列已經四次.
你有沒有吃過那些.
很甜的那些蜜棗.
那些蜜棗也都很有營養.
不過如果你有糖尿病.
當然不要不停地吃.
這個就是要告訴我們.
神的供應.
神要帶你進去的那個.
是一個美地.
極好的一個地方.
但他們這些全部就抹殺他.
全部就說這個變成一個吞沒我們的地方.
嘩這麼恐怖.
跟著他還把眼光只看那些巨人.
他看不到.
如果是這麼多巨人.
為什麼還有這麼多人在那裡住呢.
他們全部吞沒他們.
這個是很矛盾的.
第三是看他們自己.
我們好像乍賣一樣.
我們站過去.
好像你今天如果站在姚明的旁邊.
你真的好像小朋友一樣.
嘩他是巨人.
姚明很高的.

$^{401}$起碼七呎.
七呎近八呎.
很高的.
我以前我曾經第一次.
看他打NBA的時候.
我真的在家裡找張梯子.
我把他的高度拉一條尺.
跟著我發覺原來要找梯子站上去.
跟著我站上去.
嘩原來這個是姚明的高度.
非常是高瞻遠矚的.
但問題就是.
你把你自己的眼光.
只是看這個地方.
你看不到是上帝.
上帝比你的敵人是更高大衛猛.
上帝是全能者.
是不是.
跟你隔壁說好不好.
上帝是全能者.
不要害怕.
你害怕的時候.
你看我自己只是裝模作樣.
裝模作樣.
沒得揮.
他的意思就是.
我被他一踩就.
這個我用一個簡單的圖像.
給大家知道.
我把一個意念.
我跟辦公室同工.
你可不可以在網上幫我找一張圖片.
那些去賣鞋的人.
去到印度.
要開鞋舖的話.
他一見到那些.
很多印度人不穿鞋.
他的反應是怎樣.
他的反應會是怎樣.
怕沒生意.

$^{441}$但是.
另一種人的反應是怎樣.
正.
他們沒有鞋穿.
我現在給鞋他們穿.
你如果要做這個給鞋穿.
就要好像這個銷售員那樣.
解釋給你聽.
穿鞋的好處.
穿鞋可以保護你的腳.
你走路.
遠足.
是可以叫你.
可以走得更遠.
他可以很.
但這個是一個信心的眼光.
但要有後著的行動跟進.
要教育那裡那些不穿鞋的那些人.
但如果你只是看眼前.
他們不穿鞋.
我們開些鞋舖就沒生意了.
你如果這樣.
這個就是你的結論.
你就永遠.
你就是原地踏步.
我們的人生都一樣.
神有時會不會要帶領我們.
就是要走出你的安恕區.
你要進入一個你沒去過的地方.
一個新的領域.
如果是出於神的帶領.
他就一定有他的應許.
你看看那些百姓.
他們的反應是怎樣.
這是十四章的開場白.
我們一起讀.
傳輩眾大聲宣揚.
那夜百姓哭號.
以色列眾人向摩西和法輪法院言.
傳輩眾對他們說.

$^{481}$我們寧願死在埃及地.
寧願死在這曠野.
你說真是這班百姓.
難怪我講完一章的時候.
摩西真的激到吐血.
現在聖經用四個動詞去形容他們.
在原文希伯來文.
大聲宣揚其實是兩個字.
大聲.
不是細聲討論.
大聲說.
個個在說.
而且是宣揚.
引起輿論.
整個以色列的所有會眾.
個個在抱怨.
那晚他們全部在哭.
原來以色列百姓很厲害.
哭.
殺手簡.
在十一章我講他們抱怨的時候.
前後十一章出現了五次.
他們哭號.
那些水.
不是沒有水.
他們覺得缺乏.
那些不是沒有東西吃.
神每天都給他們吃碼拿.
他們吃到厭.
他們不是厭那個碼拿.
他們是厭棄上帝.
做他們的主人.
他們是在拒絕.
不是只是說不夠東西吃.
或者沒有東西吃.
他們是拒絕上帝.
但是夜間我們看到.
神做很多奇妙的偉大的神蹟奇事.
神好像火般的雲彩.
一路不離不棄.

$^{521}$帶領著百姓在曠野裡面.
而且他們在睡覺的時候.
神每晚降下碼拿.
給他們早上就去撿.
去撿拾.
每天是救用的恩典.
神也在夜間叫紅海分開.
他們是安然走過這個巷地.
到天亮的時候.
那些法老大軍追進去紅海.
神就叫東風將紅海遇合.
結果他們埃及大軍全軍覆沒.
我們看到他們完全忘記上帝.
昔日怎樣來帶領直到如今.
已經一年多的時間.
他們看過很多神蹟奇事.
但是他們仍然向摩西和阿倫發怨言.
怨言在十一章和前文都.
甚至在出埃及記的時候都出現過.
但是這次的怨言是最嚴重的一次.
為什麼這樣說呢.
這次說的是.
我們寧願死在埃及地.
寧願死在曠野.
意思是什麼.
我們都不進去.
神要應許給我們地方.
哇 你說什麼.
神帶你來到邊緣.
你們如果有信心就直進去.
可以進去得地為業.
但是他們竟然在這個時候.
向摩西和阿倫發怨言.
我們再看下一文.
約書亞和迦勒很正面的.
他們兩個人很齊心的跟他們說.
我們要信得過神的應許.
他說我們所走過所窺探的地方是極美之地.
耶和華若喜愛我們.
就必令我們進入那地.

$^{561}$擺著牛奶與麥之地.
賜給我們.
但你們不可背叛耶和華.
也不要怕那地的百姓.
因為他們是我們的食物.
保護他們的已經離開他們.
耶和華卻與我們同在.
不要怕他們.
看到嗎.
重複叫他們不要背叛神.
不要怕他們.
不是他們吞嚥我們.
我們吞嚥他們.
我們必然會得勝.
很清楚.
但是百姓仍然反應.
他們說我們要求要另立一個首領.
帶我們回埃及地.
聽完迦勒和約書亞這個陳詞.
竟然他們的回覆.
不是我們要走回頭路.
我們要回埃及地.
他們寧願做奴隸.
也不願意做神的選民.
去敬拜事奉耶和華.
這是很嚴重的否定.
就是神帶他們出埃及.
是出於祂的拯救.
他否定神是透過摩西的帶領.
也是否定神賜迦南地這個英雄.
否定神對這個地圖的評價.
否定神的聖情是美善.
神是極大的恩慈憐憫.
我們看這兩張聖經的時候.
你繼續看下文的時候.
你會看到.
我覺得上帝的忍耐是極高.
所以最後的後果.
你看到這群百姓.
我只能用兩個字.

$^{601}$我們廣東人有時很喜歡用這兩個字.
抵死.
是真的抵死.
給了你們那麼多機會.
悔改,迴轉.
你們都不醒水.
還是要繼續堅持你們的偏行己路.
所以我們看見領袖的回應.
摩西當時和亞倫就俯伏在地.
他們祈求上帝.
他們對上帝.
他們在那裡禱告.
約書亞和迦納聽到百姓這樣的要求.
他們是撕裂衣服的.
他們強調神的應許的地方.
和那些人不用怕他們.
你們最重要的是不要背叛神.
有神的同在.
如果神喜悅我們.
神一定會帶領我們進去.
其實今天我們人生都一樣.
你有沒有有時你要轉工作.
或者你看到一些異象.
你知道神呼召你.
你勇敢地你要踏上去.
踏出去.
我在這裡和大家介紹兩個年輕人.
這兩個年輕人都是很有創意.
左邊那個叫阿謙.
右邊那個叫Gordon.
是一個網紅.
他們兩個都有自己本身的職業.
阿謙是拍電影的.
他讀電影的.
他自己有他的正職.
但他自己一直很想.
為什麼今天教會.
即是說香港和全世界很多教會都是.
沒有那些年輕人.
很多年輕人流失了.

$^{641}$就是因為教會和時代有脫節.
和年輕人的溝通.
亦都是有一種疏離.
這件事他們覺得.
我們需要在這個網絡世界.
媒體的福音事公裡面.
我們要做一些事.
努力.
結果他就找他的好朋友Gordon.
這個網紅.
他說你在網上都講這麼多節目.
不如你幫我.
我們要成立一個叫紮路燈.
他們就做了一個網台的節目.
叫紮路燈.
紮路就是要告訴別人.
我們信耶穌其實是很正的.
信耶穌不是代表一帆風順.
因為我們走紮路是要付代價的.
我們做主的門徒不用怕.
當他們發覺要做的時候.
一天有時拍幾條短片就放上去.
當他們拍到一天要放七條片的時候.
發覺連自己的正職那份工作都搞不定.
他們有很多壓力.
有很多關於這些的財政經濟.
怎麼辦呢.
自己要養自己都養不了.
如果這樣.
但他們後來憑信心堅持要繼續做下去的時候.
他們也有一間機構幫助他們.
怎樣來解決每個月的營運.
而他們也請了這個節目.
就是請了王君興上去.
他們一起跟那些人講.
怎樣信耶穌不要誤解那個意思.
他用一些很創新的手法.
其實今天不只是香港.
其實今天在其他一些地方.
我都看見一些很有創意.

$^{681}$很有心的年輕人.
他們願意用另類的方式去接觸.
這些時代的人.
但我們看見百姓的回應是怎樣呢.
他一看到.
他們的反應就是要用石頭.
打死約書亞和迦勒.
你們該死.
十個人都叫我們不要去.
為什麼你們兩個要叫我們去呢.
但在這個時候.
耶和華的榮光就在活動中.
向以色人顯現.
神要叫這些百姓.
知道他們在說的話.
他們是在得罪神.
他們對上帝所揀選的.
無論是摩西亞倫.
迦勒和約書亞.
你看見神宣告迎罰百姓.
在上一節他說.
這班百姓藐視我要到何時呢.
這班百姓不信我要到何時呢.
兩次重覆.
神已經忍無可忍.
所以神說要用瘟疫擊殺百姓.
在12節.
摩西一聽到.
他馬上.
所以摩西真的極其謙和.
他忍這班百姓之餘.
但他很愛這班百姓.
他當時一聽到神要擊殺百姓.
他馬上跪在神面前.
他說神千萬不要.
不可以這樣做.
為什麼呢.
摩西為百姓代求.
基於兩個原則.
他說因為你的明星.

$^{721}$你帶他們出埃及.
過去你和埃及人.
遣過十災.
你都要讓他們知道你多厲害.
你亦都現在帶我們一起.
在曠野路裡.
你都是保護我們.
供應我們.
而且你將來.
你還要帶我們進去迦南地.
如果來到這個時候.
你擊殺這班百姓.
連迦南地那些外邦人.
都會來恥笑我們.
神的明星.
神你說你很偉大.
神為什麼.
你現在連百姓都搞不定.
接著他亦都宣告.
就是神的本性.
神你是不輕易發怒.
且有豐盛的慈愛.
你赦免罪孽.
但你無不以有罪為無罪.
神的本性就是既有慈愛.
但同時有公義.
這個宣告.
就是在西乃山裡.
當摩西領受十戒板.
他發生一件事.
就是神跟他說.
你快些下山.
那班百姓已經作反.
鑄造金牛毒.
因為摩西上山了.
四十晝夜.
他就說這個摩西.
那這個摩西去了哪裡呢.
到摩西下來的時候.
見他們這樣做.

$^{761}$更加愚昧的就是.
阿倫作為大祭司.
竟然是順從民意.
做一件得罪神的地方.
結果神是很大的懲罰.
來擊殺當時這一班的百姓.
摩西又是用這個經文.
他就向神求憐憫赦罪.
這一段的經文同樣是出現在.
《初一及記》三十四章裡面.
但是十九節.
是這一段文訴記.
繼續加下去的.
求你照你的大慈愛.
赦免這個百姓的罪孽.
好像你從埃及到如今.
要輸這百姓一樣.
神的赦罪.
為了要符合.
除了祂慈愛的本性.
但同樣一定要符合.
祂公義的要求.
所以當上帝說.
赦免我們的罪的時候.
其實背後是甚麼.
背後是有一個很大的代價.
我們今天新約的信徒就明白.
是耶穌背負了我們每一個人的罪.
死在十字架上.
還清這個罪債.
所以赦免.
其實背後是一個代價.
這個代價上帝一定要解決.
我們試一下.
明白了這個意義的時候.
你就好像這段經文說.
從父到子直到三四代.
神不會以有罪為無罪.
祂要懲罰人的罪.
所以當我們要處理這個.

$^{801}$其實是一個神學的問題.
赦免的一個代價.
但同樣神在這個時候.
祂不可以擊殺所有的百姓.
祂要繼續跟那些曾經有錯誤的人.
要保持一個關係.
但祂要按祂的慈愛去行這件事.
就必須要符合另一個條件的要求.
這個就是祂必須為了祂的名聲.
祂要來到有公義的這個原則.
公義的原則就是.
在以色列當中.
確是有一些是不忠於上帝的人.
包括這十個的探子.
他們報一個謠言的壞消息.
誤導了神的百姓.
神是要懲罰他們的.
神也都雖然說赦免百姓.
但是上帝是拒絕那些活逆者進入應許之地.
所以這個後文是會跟我們交代的.
在後文23到24節.
神說這一班藐視我的人.
他們沒有一個可以見到這個應許地.
除了迦勒.
接著29到30節.
他也都說到凡20歲以上向我發怨言的.
他們必不得盡我所起誓應許他們居住的地方.
他們的屍體是要所謂倒在這個曠野.
他們接下來這38年多的時間.
他們不斷這一班那一代的人.
就一直死在曠野.
他們都會生而育裡.
真正入應許之地是他們的兒女.
但是他們的兒女33到34節.
是要在曠野遊牧40年.
是擔當你們不信的罪.
所以父母的罪.
原來都是牽連到兒女都要承擔.
承擔的時候這個牽連性真的很大.
就因為他們誤信靠神.

$^{841}$以致他們帶來的窺探那個地方40天.
一年就抵一日.
所以他們要有這40年的漂流在曠野.
最後在37到38節.
那些散佈謠言的人.
他們遭受瘟疫.
死在神的面前.
這個是說即死.
當日神就是施行瘟疫.
即時審判.
這10個探子全部要在神的面前死.
只有約書和迦勒他們可以存活.
所以12個人死了10個.
那你說那些百姓怕不怕.
醒了沒有.
他們醒了.
那他們是不是真的悔改.
你再讀下去文蘇記.
遲些繼續回來聽.
然後你看下去.
他們有些人就說.
我們現在要去得神給我們應許的地為業.
我們要衝進去.
我們要打仗.
和那些阿瑪利人,迦南人.
我們打仗.
我們一定可以拿到那個地方.
摩西怎樣跟我們說.
你們為何要違背耶和華的指示.
這件事不能夠順利.
你們不可以上去.
因為耶和華都不在你們中間.
又沒有吩咐你們.
只是吩咐你們去探地.
你們現在回來.
你們現在看到神懲罰那些惡的探子.
你們馬上就說.
我們要去拿地了.
這個不是真正的信服.
弟兄姊妹今天我們聽神的話.

$^{881}$可能你聽到都聽了幾十年.
但你有沒有試過.
你真正的選擇至死忠心去跟從神.
還是有時聽一點不聽一點.
還是選擇性的信服.
你終有一日你會發現.
你其實倒在應許的恩典的門之外.
原來上帝可能給你的家庭.
給你的個人.
甚至給你的事業.
不是只有這麼多.
不過你錯過了上帝給你的很多恩典.
是因為你不肯去相信.
你去知取神的應許.
我們再看的時候.
這個我們需要今天重新回轉.
重新聽神的話.
這個很重要.
神的應許放在我們每一個人面前.
你今天用信心去連上去.
我在這裡和大家分享.
就是李聖林牧師.
我昨天早上星期六.
就在聖安德魯這個教會.
其實場面非常感人.
我是其中一個受邀請的牧者.
因為那個場地有限.
所以他們那個實地的.
舉行的安息禮拜.
必須有通行證才能進去.
在整個聚會裡面.
其實在另外一個看直播.
那裡有八百人.
那個地方.
很多很多弟兄姊妹.
他們很懷念.
這一位我們所敬愛的.
李聖林牧師.
他一生.
牧養教會和神學院教導.

$^{921}$而他不單是很精於.
舊約聖經的教導.
他能夠活出神的話.
特別在三年半之前.
大家知道.
他的小兒子阿Mo.
就在那個紅館.
MIRROR的音樂會.
發生了一個很大的意外.
當時在高空那個LED.
一個大屏幕就掉了下來.
結果他的小兒子阿Mo.
正正就壓在他的背上.
他由脊椎一直全身是癱瘓.
但是牧師當時在加拿大.
他馬上趕回來香港.
接著是說三年多時間.
一千三百六十八天.
他真是年終無憂.
無怨無悔.
是陪伴自己的兒子.
不離不棄.
這種父愛.
正是彰顯了天父.
今天都是每一天.
等待每一個浪子回頭.
當時李牧師.
很愛惜這個兒子.
他一直用詩篇二十三篇鼓勵他.
他每一天用靈修神的說話.
安慰他.
叫他不要放棄.
當醫生判斷他的兒子.
基本上是不能醫治的時候.
他後來轉到中大的醫學院.
在一個很好的團隊.
後來他遇到一個很好的機會.
見到一個類似受傷的個案.
一個病人.
得到很大的康復.

$^{961}$他可以坐在輪椅.
但起碼可以活動.
他可以手腳用輪椅來控制.
比阿牧現在的情況.
由完全不能活動.
他現在開始.
可以活動到一些地方.
而且他是朝著這個方向.
這是一個目標.
宗教的目標就是就算不能起身.
起碼可以照顧到自己.
可以用輪椅來代步.
這才是神.
神其實在他們整個家庭最困難的時候.
他們都能夠挺過去.
是因為他們相信.
神沒有離棄他們.
所以在昨天這個聚會裡.
他三個兒女都有講見證.
阿牧就透過視像.
只是錄音去講.
他說.
我給爸爸的信.
我很願意將來在天下.
站起來抱著你.
這個年輕人.
他若然沒有一個這麼愛他的父親.
和一個充滿著上帝的愛的家庭.
他們三兄弟.
最大的就是這個大女兒Karen.
他們在爸爸媽媽的身上.
看到神的說話.
神的應許.
而看到在李牧師的身上.
他不只是教.
他不只是講.
他就走出來.
成為很多人生命的幫助.
這首詩是他的大女兒Karen.
作出來的.

$^{1001}$而且昨天在現場演唱出來.
我自己很被觸動.
這首詩是講了詩篇23篇.
也都講了詩篇127篇.
121篇.
我要向山舉目.
我要幫助從何而來.
我幫助從天地的弱壞而來.
主你就是我的方向.
引領我一生的路.
樹林中同行.
幽谷中細雨.
我仍記得你的應許.
你的應許.
主你是我的安慰.
使我安潔青草地.
雖然經過黑暗.
我仍然選擇是相信.
我深信你必帶領.
我舉目定睛仰望群山.
我幫助從何而來.
主你就是從創造天地的主而來.
你在我的面前擺設筵席.
賜平安與喜樂.
丁姐妹.
願以今日你和我.
我們每一個一生.
都不忘記神的應許.
神的說話.
我不知道你現在人生.
是在經歷什麼階段.
你前行.
你看到的是一個巨人.
你很害怕.
或者是一個城牆很寬廣.
你沒有辦法去攻打這個城牆.
是不是.
是不是真的要這麼害怕.
最重要.
你是知道.

$^{1041}$是不是上帝吩咐.
無論你要轉工作.
無論你是你的家庭.
你或者你和你的一些同事.
你有一些關係.
你要復和的話.
這些就好像巨人一樣.
魔鬼就整天在恐嚇我們.
不行的.
沒有辦法解決的.
這件事很複雜.
你一定輸定了.
我告訴你.
在神的領域裡面.
沒有難成的事.
你信嗎.
你信得過神的應許嗎.
不要怕.
只要信.
不要怕.
只要信.
和你旁邊的人說.
不要怕.
只要信.
信神的應許.
神與我們同在.
用信心連結在神的話語上.
我們一起祈禱.
天父我求你幫助我們每一個弟兄姊妹.
每一個家庭.
我們身處的.
無論在多大的難處裡面.
求主給我們有信心的眼光.
是看見你.
不只是看見那些困難.
主要讓我們用信心的眼睛.
我們有新的眼光.
新的亮光.
我們敏銳.
你的帶領.

$^{1081}$你自己的幫助.
你的幫助.
就是創造天地萬有的主.
聖多主.
以致我們今天.
不需要焦慮.
不需要徬徨.
我們不需要懼怕.
我們只是單單默然在你面前.
信服你.
仰望依靠你.
求主引領我們.
前行的路.
能夠打亮美好的像.
榮神益人.
同心祈禱.
奉主名求.
阿們.
\newpage



\section{希伯來書 1:1-3}
\label{sec:032273}
\textbf{希伯來書系列 (1) 進入聖經中的奇異世界}~~~~張健庭牧師博士
\newline
\newline
連結: \href{http://www.church.com.hk/acms/content.asp?site=cdc\&op=show\&type=product\&code=032273\&layout=sermon}{\texttt{032273}} ~~~~ 語音日期: 2026-06-07
\newline
\newline
\hyperref[sec:032272]{\small{< < < PREV SERMON < < <}}
~
\hyperref[sec:index_chronic]{\small{[返順時目]}}
~
\hyperref[sec:032274]{\small{> > > NEXT SERMON > > >}}
\newline
\newline
希伯來書 1:1-3
\newline
\begin{longtable}{cl}
\hline
\hline
章節 & 經文 (和合本修訂版)\\
\hline
1:1 & \begin{tabularx}{0.7\textwidth}{X} 古時候，神藉著眾先知多次多方向列祖說話， \end{tabularx} \\ \\ \relax
1:2 & \begin{tabularx}{0.7\textwidth}{X} 末世，藉著他兒子向我們說話，又立他為承受萬有的，也藉著他創造宇宙。 \end{tabularx} \\ \\ \relax
1:3 & \begin{tabularx}{0.7\textwidth}{X} 他是神榮耀的光輝，是神本體的真像，常用他大能的命令托住萬有。他洗淨了人的罪，就坐在高天至大者的右邊。 \end{tabularx} \\ \\
[1ex]
\hline
\hline
\end{longtable}
$^{1}$好 艇姊妹早安.
我們再一次低頭祈禱.
天父願你教導我們.
當我們今天開始思想希伯來書的時候.
事實上有很多我們都承認.
好像是一個很陌生的世界.
因為希伯來書所講的當中很多的內容.
包括怎樣獻祭.
甚麼聖所.
甚麼祭司.
這些對於我們現代人來講.
我們坦白說.
我們真的一竅不通.
但是我們願你的聖靈.
穿越我們的有限.
在我們今天的時候.
一起聆聽你的話語的時候.
你自己向我們講說話.
奉耶穌基督的名祈求.
阿們.
好了 接下來就會開始.
和大家來思想希伯來書.
因為我想這個希伯來書本身.
對於很多人來講.
都是一本挺特別的書.
因為當中講到很多的內容.
和舊約有關.
我相信很多艇姊妹.
其實你對舊約或多或少有些印象.
因為你一定會看過一些舊約的故事.
或者舊約一些基本的內容.
譬如憲制啊 回望啊那些.
你都應該是不會陌生的.
但是對比在希伯來書的時候.
就似乎又有一些.
似懂非懂的一種印象.
因為希伯來書本身又不單止是.
好像重新講回舊約那些東西.
它就好像再加多了一層的演繹下去.
將舊約的東西用另一個角度來演繹.

$^{41}$對於一些本來舊約你都不是很了解透徹的時候.
再疊多一層演繹上去.
你就會發覺有時候都好像摸不著頭腦.
所以就希望接下來講希伯來書.
事實上.
我相信很多艇姊妹.
如果你回到這裡教會一段時間.
你都必然會知道其實很多書.
我一卷一卷都講過了.
約翰福音講過兩次.
其實從頭到尾.
然後羅馬書又講過.
甚至乎啟示錄都講過.
就是以弗所啊.
菲勒比啊那些彼得前書啊那些亞國書.
差不多全部很多都講過.
舊約那些都講過.
你都會問一個問題.
為什麼牧師一直都沒有講希伯來書呢.
而去到這次我去想要講希伯來書的時候.
其實我都有一些的.
又不是叫做掙扎的.
都有一些不是那麼自然的.
一想起還沒講過不如講一下.
因為我發覺希伯來書當中.
正如剛才所講.
你要某程度對於舊約有相當的了解.
因為希伯來書又不是只講舊約.
他將舊約.
如果你以前舊式的那些投影機.
或者一張張膠片.
你原本有一張東西.
然後再疊多一張上去.
再疊多一張上去.
你就發覺那些東西就開始模糊了.
有一點點是這樣的印象.
我覺得講希伯來書.
因為如果你原底舊約那些不是那麼清楚.
然後再講希伯來書.
其實他不是直接講舊約的.

$^{81}$他引用舊約.
但是他有他自己對舊約的一種很特別的演繹.
這樣的時候.
如果你不是很清楚原先舊約.
所以我就發覺有一定程度的困難.
所以我一直都很猶豫.
有一點點猶豫.
究竟是不是要講希伯來書呢.
但是經過禱告又考慮過.
之後其實都很多都講過了.
那就試一下挑戰一下大家.
也挑戰一下自己.
因為講這些東西不只是我講完就可以了.
你們也要明白.
這個也是一個挑戰.
但是在這樣的思想過程當中.
我就發覺了一件事.
其實聖經是一個很特別或者很陌生的世界.
你對於今時今日的人來說,你會不會憲制呢.
你完全不會有想.
你直接想都不會想憲制的.
對於信仰的事情.
很多東西你都覺得是跟心靈有關的.
你怎麼會無端端.
不是無端端.
你怎麼會思想聯繫到.
拿一隻動物.
拿一隻新畜來殺了牠.
然後就會得到解放呢.
只是講憲制.
你細心來想想.
是一個很陌生的世界.
而不單只是憲制.
又有祭司.
又有回眸.
這些全部都是希伯來書的背景.
我就問一個問題.
其實是否需要這樣去講聖經呢.
當然我不是說聖經的內容不重要.
但我的意思是.

$^{121}$可不可以我不講那些憲制的內容.
而只是講那個精髓.
你明白我的意思嗎.
就好像收禮物.
你只是拿了盒子.
或者去超級市場見到有些人買粟米.
不停地撕.
撕掉衣服.
只要心就行了.
可不可以這樣呢.
我自己的反省.
如果講聖經.
譬如我們有罪.
耶穌為我們贖罪.
死於十字架.
我們與神和好.
這些觀念我們完全是明白的.
我們也需要明白才能做到基督徒.
但是我可以講這些觀念.
講的很明白.
我不需要跟你講憲制.
你明白我的意思嗎.
所以我也明白為什麼有些基督徒.
他們說我信耶穌我就信了.
我看聖經看不明白.
不太有興趣看.
其實背後的心態如出一轍.
就是因為他覺得.
核心那些我明白就行了.
為什麼要看很多.
現在對於現代人來說.
這些又不是很明白.
又好像沒什麼關係的東西呢.
這個就是.
我相信對於現代的基督徒.
是一個很大的困惑.
你未必會問這個問題.
但是在實踐上.
其實你是正在遇到這個問題.
我究竟需不需要再用這些.

$^{161}$希伯來書自己也說.
舊約的原事後事得映而形體卻是基督.
這個時候.
形體已經是基督.
就行了.
為什麼要去看.
那些憲制舊約那些東西呢.
事實上在教會歷史當中.
其實這個也是經過一個很激烈的討論.
在早期的教會歷史的時候.
那時候甚至在馬吉安主義興起.
這些細節我不說了.
但是那時候聖經新約和舊約的正典還沒確立.
那時候已經有一些基督徒.
就發覺.
我現在相信耶穌.
我覺得耶穌是絕對超越的.
那還需要看舊約聖經嗎.
其實有些基督徒也是這樣問.
那時候在早期教會歷史的時候.
當然經過了很激烈的討論的時候.
正統的教會其實已經有一個共識.
就是我們是需要.
我們是需要進入去聖經的一個奇異的世界.
我所說的奇異的世界當然就是舊約那些.
但是其實某個程度對於我們今天來說.
不單是舊約.
新約有些東西我們都覺得很奇異.
那時候去理解信仰.
事實上我們就是需要這樣.
去進入聖經一個奇異的世界去理解信仰.
為什麼呢.
其中有一個很重要的問題.
就是有些東西.
如果你要明白它真實所要表達的意義的時候.
你說我好像撕粟米一樣.
撕掉表面的衣服.
竟然也可以吃粟米.
沒有分別的.
但是如果不是粟米呢.

$^{201}$是不是啊.
另外一個例子你可以想到.
洋蔥呢.
你撕掉多少層衣服.
可以吃到洋蔥的核心啊.
是沒有的.
如果你真的以為洋蔥有個核心.
撕完衣服就可以了.
撕著撕著撕到沒有了.
原來是沒有的.
事實上有時候.
可能是介乎兩者中間.
聖經是粟米還是洋蔥呢.
我不能夠很講實地告訴你.
一定是洋蔥.
撕掉一層層之後沒有東西.
我不能夠講實地告訴你.
一定不是粟米.
因為如果你撕掉衣服之後.
我也不知道你撕剩下的是什麼.
而事實上.
所以在早期的教父.
他們去讀聖經的時候.
他們已經體會到一件事.
如果你不高呼不明白舊約的時候.
你不能夠明白新約.
而希伯來書的作者.
當然他寫作的時候有他的背景.
當時的信徒.
他們是一些猶太的基督徒.
我相信是.
所以他們是非常明白舊約.
希伯來書的作者.
用舊約那些當然是魚得水的.
跟我們解釋.
但問題是.
希伯來書留在聖經的時候.
他也不只是給那些猶太基督徒.
你要明白.
他也是給你和我.

$^{241}$因為他留在聖經那裡.
聖經「都是神所物是的」.
「都是」這個字可以翻譯成.
聖經所有都是神的話語.
所以要高呼進入聖經的奇異世界.
是必不或缺的.
很多時候有時候.
我也明白.
有時候看聖經.
你似懂非懂的時候.
你會覺得.
是不是有這樣的必要呢.
有時候不如聽一下.
一些人解釋聖經就好過.
自己看聖經.
因為聽人解釋的時候.
人家都把衣服脫光了.
你只是吃個粟米.
或者你看看書就行了.
但是重點還是一樣.
我們每個人都需要.
好像自己進入聖經的世界.
去經歷聖經的世界.
所帶給我們的那種體會.
你才能夠明白那種衝擊力.
簡單來說.
我舉個例子來說.
對於今天的人.
犯罪.
求神饒恕.
這個經歷本身.
我相信每個基督徒都有.
否則你也不會做基督徒.
你有罪.
你求神來去饒恕你.
但是.
那個的深度就有深有淺.
有些人犯完罪很輕描淡寫.
然後祈禱求神饒恕.
我記得我初初信耶穌的時候.

$^{281}$也不知道是不是真心信.
也不知道就是.
我平時讀經教會那些就不回家.
但是每晚睡前.
我都是睡在床上.
閉上眼睛睡前.
必然會說一句話.
神饒恕我今天的罪.
奉耶穌基督的名祈求.
阿們.
這樣睡覺.
這個成為了我的一個習慣.
在頭兩年做基督徒的時候.
你說我是不是認罪呢.
當然是認罪.
有多深刻呢.
坦白說沒什麼深刻的.
只不過是一個.
我們都知道有些事情.
今天一定是做錯了.
不過好像.
補個底.
買個保險.
或者兜兜轉轉.
在神面前.
這樣不是很深刻的.
但是如果你進入了聖經的世界.
你看到那個血淋淋的.
當你犯了罪.
根本沒辦法去跟神相交.
也沒辦法去.
有任何恢復的可能.
唯有你要付出一個代價.
而這個代價.
窮的人.
就要拿一隻八甲或者斑溝.
比較豐富的人.
可能就要拿一隻羊.
再有錢的人可能就要拿一隻牛.
但是這個代價其實很重的.

$^{321}$窮的人拿一隻斑溝這麼容易嗎.
那時候肉都沒得吃.
是不是.
你拿一隻斑溝.
然後你來獻祭.
你有錢.
都是這樣說的.
牛啊羊啊那些是很矜貴的.
那些是等於.
你宰一隻牛就等於找一輛車.
拿去廢鐵場一樣.
很矜貴.
不是說爛車.
是一輛好好地地.
可以有生產能力的車.
就拿去廢鐵場報銷了它.
你肉不肉赤.
很肉赤的.
就是要進入一個這樣的世界.
你體會到那種無奈.
那個的無能為力.
而即使是要付出很大的代價.
但是這個也是唯一要經過.
真正你去到獻祭的時候.
哇那種血淋淋.
你以為獻祭是一家大小兒童適宜嗎.
兒童不宜的場面.
真的很血腥.
我記得那時候.
中學的時候.
去香港長洲宿營.
在附近有個屠宰場.
哇晚上聽到那些牛羊在叫.
那些聲音嚇得整晚都睡不著.
真的很恐怖.
可能他們都知道自己死了.
要被人宰了.
那種淒厲.
而就是要進入一個這樣的世界.
原來就是犯罪得到回復.

$^{361}$而經歷到那種可以回復的喜樂.
那種的深刻.
現代人是沒有的.
是沒有的.
坦白說是沒有的.
當然.
所以讀聖經本身.
有時候你跳過那些.
不用理那些很複雜很繁瑣的細節.
那些舊約的儀式.
你只不過是掌握到觀念.
你掌握到觀念.
所以我們有時候是需要.
真正的進入聖經的世界.
而事實上.
當我們不進入聖經的世界的時候.
我們現代人如何去理解信仰呢?.
如何去理解信仰呢?.
因為我們沒有了聖經的世界.
那如何去理解信仰呢?.
因為聖經那裡又是講那些.
獻祭,贖罪諸如此類的東西的時候.
那我們如何去理解呢?.
就是用我們現在的生活去理解.
所以我們看罪的.
很容易就想到罪犯.
打劫,偷東西那些就是罪.
因為那些的標誌就是罪犯.
但是聖經講罪根本就不是這樣的.
聖經講罪基本上有三個層次.
你看利未記那裡.
當然了.
打劫,殺人那些是其中一個罪.
但是也有宗教的層次.
神那裡相交是另外一個層次.
另外一個是禮儀的層次.
就是在生活上如何分別為聖.
也是另外一個層次.
但是我們完全沒有的.
因為我們就是要用現在我們的生活.

$^{401}$反而來說套入去理解聖經的觀念那裡.
理解了多少呢?.
很多時候就變成一知半解.
所以我今早在國語堂那裡講.
講《仕途信經》的時候.
才講到一個對比.
就是說現代的基督徒很多時候.
想信耶穌是什麼一回事呢?.
就是死了之後信耶穌就上天堂.
有多少人是這樣想呢?.
我現在不想做調查舉手.
但是我希望你不是這樣想.
為什麼?.
因為我今天講《仕途信經》.
我相信身體復活.
就是新天新地都不是上天堂的.
你找遍整本聖經.
基本上你如果能夠找到一句聖經.
說信了耶穌死了之後上天堂.
我請你吃飯.
就是一定沒有的.
因為我知道.
你看遍整本聖經根本都沒有.
死了之後與神同在就有.
即是離世與主永遠同在就有.
但是那裡是不是天堂呢?.
聖經都沒有的.
聖經只有新天新地.
神的障礙在人間.
那裡不是天堂是人間.
所以很多觀念本身.
我們就用我們現代人的想法.
反而來代入聖經.
如果用現代人的想法代入聖經.
可能聖經有些對應得上.
有些對應不上.
對應不上的就硬要按進去.
所以上天堂的是硬要按進去.
因為你聽那些故事.
有地獄當然有天堂.

$^{441}$但是聖經是有地獄.
但是聖經有沒有天堂呢?.
你有沒有想過?.
未必有.
有陰間是有.
但是陰間是不是地獄.
這個也是值得討論.
這個不說了.
但是聖經肯定是沒有天堂.
而一個更加深入的問題.
如果我們將我們的生活代入聖經.
或者將聖經那些觀念.
其實將我們現在所有的.
現代人的生活塞進去.
就帶出一個很大的問題.
就是說.
我們的信仰是我們現在的生活主導.
還是信仰主導我們的生活呢?.
這兩個是完全不同的.
是我們現在的生活主導我們的信仰.
還是信仰主導我們的生活呢?.
我相信很多人是前者.
是生活主導信仰的.
所以你有問題的時候.
就親近神多一點.
因為神是在你的生活當中解決你的問題.
我記得我那時候讀書的時候.
一到考試季節.
就有一點點換人的狀況.
什麼叫做換人的狀況呢?.
就是有些平時不上教會的人.
就會上教會.
有些平時上了教會的人.
又不上教會.
有沒有想過.
有沒有想過為什麼會出現這樣的情況呢?.
其實很明顯的.
平時不上教會的人會上教會.
因為可能第一讀書又悶.
上教會比較正派.

$^{481}$家人就不會阻止.
如果是基督教家庭的話.
另外還會想著讀書也不行.
上教會.
起碼神幫幫我吧.
考試.
這樣就上教會了.
反過來.
另外那些上了教會的人.
為什麼不上教會呢?.
就因為現在很忙.
要考試.
要努力讀書.
爭取上了教會.
起碼星期天沒有了伴奏.
不行了.
星期一考了.
那就不上了.
但是兩種都是一樣的.
想法都是.
其實都是生活主導了信仰.
是不是?.
都是生活主導了信仰.
不是信仰主導生活.
是生活主導信仰.
而這個就是一個很嚴重的問題.
就是如果你不能夠進入聖經的世界.
去理解信仰的時候.
多少我們的信仰.
其實是偏差了.
我們都不知道.
而被我們現世的生活主導住呢?.
所以就是因為.
如果是這樣的時候.
其實信仰很容易變為重擔.
尤其是.
如果現在我們信仰.
跟各樣東西沒有什麼衝突的時候.
你也不需要為信仰付上代價的時候.
必然是OK的.

$^{521}$但是如果.
有碧白尼林又如何呢?.
如果譬如舉例來說.
信耶穌可能會令到你失業的又如何呢?.
如果你要持守你信仰的原則.
是令到你吃虧的時候.
甚至受大苦的時候.
又如何呢?.
這個就成為一個很嚴重的問題.
而希伯來書的作者.
不是,希伯來書的讀者.
其實正正就是面對一個這樣的問題.
希伯來書的人.
其實他們是信耶穌的人.
他們是猶太人.
但是他們已經信了耶穌.
這個是.
基本上都沒有什麼需要爭論的.
基本上都很顯著.
他們是一群信耶穌的人.
所以那些內容.
本身你看到他們是信耶穌.
但是我想.
跟剛才所說的都是很類似的.
他信了耶穌之後.
基本上各樣東西相安無事的時候.
其實是說生活指導信仰.
信仰指導生活.
有時候哪裡想的那麼多.
都不知道就是.
繼續生活這樣.
是不是?.
信耶穌似乎都OK.
都是多麼合理多麼好.
但是當挑戰來臨了.
有人為信仰需要付上代價.
有人可能因為信仰的緣故.
生意都做不成了.
家產都被人充公了.
甚至乎有些被人迫害了.

$^{561}$有些人因為這樣而被人排擠了.
有些人甚至乎是面對生命的危險了.
那又如何呢?.
那就成為了一個很大的挑戰.
那就不是隨隨便便.
信仰指導生活.
生活指導信仰都是差不多的.
那就不是這麼說了.
如果你的信仰搞不清楚.
它不能夠成為你的指導.
你的生活的時候.
你就要後退了.
是不是?.
那就唯有你的信仰妥協.
妥協就是後退了.
而希伯來書的作者.
正正就是提醒當時的讀者.
你想清楚你選擇哪一條路.
如果你的生活仍然要先行.
信仰只不過是其中一個附加品的時候.
你選擇一條後退的路.
你就知道代價是什麼.
希伯來書其實是一個這樣的警告.
所以在希伯來書的中間.
第二章,第四章,第六章,第十章.
講著講著就加插了這樣的警告在其中.
就因為他們正正就是面對著這樣的問題.
你搞清楚你的信仰是什麼的時候.
讓它成為你的指導的時候.
有時候你進入一個很.
坦白說都不是這麼容易的一個世界.
但是你搞清楚的時候.
你就發覺原來是這麼豐富的.
原來神的救贖計劃是這麼豐富的.
正正就是面對著一個這樣的情況.
希伯來書就是面對著一群.
當時的猶太教信了耶穌的一群讀者.
他們似乎信了耶穌之後.
是需要有一個選擇.
因為當他們面對著一個很嚴峻的.

$^{601}$可能最初信的時候沒有這麼大的困難.
最初信耶穌的人也不覺得.
最初有很大的迫不及待.
但是後來慢慢慢慢的時候就發覺.
真的要他們選擇了.
這個時候他們要怎麼選擇呢?.
希伯來書基本上就是講出了.
叫他們要小心選擇.
甚至你不要選擇錯誤.
不要選擇後退的路.
所以就交代了希伯來書的背景就是這樣.
而希伯來書的第一章第一至第三節.
基本上就是整本書開頭的一個很重要的引言.
這個引言本身每一句話.
就涵蓋了後面所講的主要內容.
我們很快看看他講什麼呢?.
神在古時直著眾先知多次多方.
曉遇列祖就在者末世.
直著他兒子曉遇我們.
特別是他用了兩個字.
我們平時不會用的.
第一就是曉遇.
第二就是末世.
曉遇是什麼意思呢?.
曉遇的意思就是說神遇.
就是說好像.
你想像就是神親自跟人說一些東西就叫做曉遇.
神親自.
那個是沒有詮釋空間的.
意思就是說譬如舉例來說.
牧師講道.
代表神去將神的心意直著聖經講給你聽.
那個有詮釋空間的.
有些人說這個牧師.
他講的那些東西.
都不一定是跟神的旨意或者神的曉遇相關的.
因為你仍然有詮釋空間的.
但是曉遇這個字就沒有詮釋空間的.
就等於神在西乃山跟以色列人說話.
沒有詮釋空間的話.

$^{641}$那個不是神講的.
神直著摩西跟以色列人在申命記跟他們說.
叫他們生與死.
禍與福擺在他們面前.
那個是代表神來說.
也沒有詮釋空間的.
沒有說摩西臨死的時候.
臨死老人交代遺囑.
跟神說沒有關係.
沒有人會這樣質疑的.
而事實上曉遇就是這個意思.
就是說其實神是一位說話的神.
他會發出他自己獨一無二的心意.
是清楚地來表達.
就是曉遇的意思.
神是一位獨一無二的神.
我們的神永遠都是說話的神.
基督教信仰所說的神是一個說話的神.
有其他的宗教.
當然世界上各種宗教形形色色.
什麼都有.
但是有些宗教.
譬如本身他們的神是不說話的.
是需要人自己去領會的.
你立地成佛.
就是你領會了怎樣可以得道.
可以成佛.
中國人那些成仙那些.
都不是有什麼說話的.
都是他們自己練仙的.
其實道教成仙的就是道教.
其實他們都是一種信仰.
有很多不同的宗教.
他們那些都未必需要說話.
意思就是說他們的神未必需要說話.
人可能每個人內在都有一些佛性.
有一些神性.
你自己發掘你內在.
打通你的神性的開關.
你就可以得到一種宗教上的昇華.

$^{681}$但是我們的信仰就不是這樣的.
永遠都是神先說話.
神將我們的心意.
很清楚.
不會被扭曲地表達出來.
這個就是曉喻.
這個是信仰當中的啟示成份.
就是神是說話的神.
我們的神是說話的神.
我們的信仰不是我們猜出來的.
不是一些很有智慧的人去揣摩到.
是神自己親自說.
他說神怎樣說呢.
藉著眾先知.
在古時藉著眾先知.
多次多方的曉喻列祖.
神是藉著人來說話.
當然在聖經的故事當中.
有神直接說話的.
比如在西乃山那裡就是了.
比如以利亞先知在荷列山那裡就是了.
但是那是少數的.
多數是藉著先知來說話.
神的說話是需要藉著人來傳遞的.
但是那些人是先知.
是不會錯的.
不是說先知說話不會錯.
而是說先知在要傳遞神的旨意的時刻.
他就不會錯.
但是接著他說.
「造在者沒勢曉喻他.
藉著他而止曉喻我們」.
剛才說的第一個字就是「曉喻」.
第二就是「沒勢」.
這裡所說的「沒勢」.
不是我們一般人想的世界末日.
是不是?.
這裡當然不是世界末日.
耶穌基督「曉喻」我們.
耶穌基督二千年前降生為人.

$^{721}$來將神的心意表達出來.
那就是「沒勢」.
所以在聖經當中.
說「沒勢」不是說世界末日.
聖經當中說「沒勢」的意思就是說.
神的救贖計劃進入了最後的時刻.
這就叫做「沒勢」.
神的救贖計劃進入了最後的時刻.
耶穌基督來臨的時候.
好像神的救贖計劃已經進入了最後一刻.
那就是稱為「沒勢」了.
所以現在我們是不是活在「沒勢」?.
是不是?.
當然是的.
因為耶穌基督來臨的時候.
好像已經將這個「沒勢」的時鐘打開了.
已經開動了.
現在這個時鐘在跳的時候.
就是這個「沒勢」的時鐘.
已經不斷地向前了.
這個「沒勢」的時鐘.
就是在耶穌基督來臨的時候已經打開了.
所以整個新約聖經的使徒.
他們理解他們是活在「沒勢」的.
我們也是活在一個「沒勢」當中.
而第二個元素就說到.
「他又早已立他為他承受萬有的」.
這裡基本上.
在第一至第三節.
希伯來書第一至第三節那裡.
其實他講出了一個很重要的.
就是整個的信仰.
是信仰指導生活.
還是說生活指導信仰呢?.
是有一個很重要的關鍵.
就是你怎麼去理解耶穌基督.
你怎麼去認識和跟從耶穌基督.
如果這件事你是搞不清的.
是沒有很清澈的去決定的.
你一定.

$^{761}$信仰一定是指導不了你的.
指導不了你的生活的.
而信仰能夠指導你的生活的.
唯一你就認清那位的主.
你是跟隨著哪位的主.
其實背後的邏輯就是這樣.
你不認清你是跟隨著哪位的主.
或者你以為你跟隨著他.
但其實你不認識他的時候.
你的信仰是很膚淺.
沒辦法成為你生活上遇到挑戰的一種指導.
這是一個很真實的事實.
當你成為了神的兒女的時候.
他說以前是奴僕現在成為兒女的時候.
他整個人對於死亡的看法完全不同.
後來他去到那些.
可以這麼說.
那些曠工其實是那些很野蠻的老粗.
一聽到聽得不對.
整個煤礦扔在他頭上.
隨時被人扔死的.
他有幾次是遍體鱗傷.
傳完道之後被人扔到.
他其實不怕死的.
那些隨時變成了暴徒.
你聽港獨說得不對的時候.
你差不多進入了一些暴徒那裡去講道.
後來他不怕死.
因為他對於耶穌基督那種認識和體會是完全不同了.
不用說約翰維斯.
現在很多認識神的信徒.
他們面對著死亡那種.
我們自己也見過離開了的長者.
老人家面對著死亡那種.
心裡面那種沒有懼怕.
而是因為認識耶穌基督.
而那種認識不是一種.
耶穌降生在馬槽.
出來傳道.
升了十字架.

$^{801}$然後復活.
兒童主義學的故事那種認識.
不是停留在那個層次.
所以你怎樣認識耶穌基督是很重要.
而耶穌基督就是神終極的啟示者.
第二點就是說.
「早已納他為承受萬有的」意思就是說.
我們認識的跟從的這位主.
是終極的一切存在的源頭.
或者是主人.
「早已納他為承受萬有」這個字當然是.
如果進入聖經的世界.
是說到其實他們是.
神的遺產是可以傳的.
亞伯拉罕就傳給了以撒.
以撒兩個兒子就只傳給了雅各.
雅各就傳給了十二個先祖.
但是十二個先祖基本上後來就.
傳給猶大和約瑟.
然後約瑟基本上就傳給了嗎哪西.
不是.
是傳給了以法連.
不是嗎哪西.
傳給了以法連.
猶大呢.
後來傳來傳去就傳給了大衛.
所以他們將神的英海的遺產.
是一代一代的傳下去.
而最終他們也都在創世記第三章的時候.
他們也都去問一個問題.
究竟最終那個女人的後裔.
可以完全的去領受著神的權柄.
去將一切的仇敵打敗的.
是誰成為這個女人的後裔.
可以承受這一切的英海呢?.
其實這裡就是正在回答的問題.
所以你沒有進入到聖經的世界.
你看這些就覺得好像.
觀念上你只是知道.
但他就是在說原來歷世歷代神的英海.

$^{841}$一代一代的傳下去的時候.
去到最後現在終極的主人.
現在就是耶穌基督就是終極的主人.
一切的存在.
宇宙萬物的存在.
你和我的生命.
現在活在世上.
我們都是存在的一部分.
但是我們的終極主人就是耶穌基督.
現在這個世界.
很多的混亂.
很多的不合神的心意.
很多的腐爛.
很多的罪惡.
好像已經失控了.
但是終極的主人.
其實仍然都是耶穌基督.
在當時希伯來書的信徒.
他們就面對著又要被迫害.
這個世界又好像把神當作沒有到的.
甚至反對的.
但是終極的主人就是耶穌基督.
接著他就再說.
也曾藉著他創造諸世界.
這個世界追溯到源頭的時候.
也是藉著耶穌基督而做的.
這個在新約聖經的信仰就很清楚.
「太初有道,道與神同在,道就是神」.
「萬物是藉著他造的,凡被造的物也一樣,不是藉著他造的」.
一切都是藉著耶穌基督而做的.
在保羅的書信也有.
在歌羅西書也有.
愛子是守身神的像.
一切都是藉著愛子而被造的.
是藉著耶穌基督而被造的.
但是在說被造的時候.
就有一個觀念.
我們今天可能比較少想.
就是你不知道一切的創造.
必然是有其目的的.

$^{881}$是不是?.
一切的創造必然是有其目的的.
那樣東西的存在如果是被創造出來.
不是說.
舉例來說,如果我的手有疤痕.
我的疤痕不是創造出來的.
我的疤痕是因為意外而造出來的.
這樣就可能沒有目的.
我也不想有疤痕.
我也有疤痕,無緣無故有疤痕.
也不想有目的.
因為創造出來的東西一定是有其目的的.
一定有.
而宇宙萬物如果是由耶穌基督而創造的時候.
就是宇宙萬物也有在基督裡面可以找到它的目的.
你有沒有想過?.
宇宙萬物也有在基督裡面可以找到它的目的.
人存在的目的,你存在的目的.
宇宙萬物存在的目的.
是唯有回到基督那裡才可以找得到.
人生存在的意義和目的.
唯有在基督裡面才可以找得到.
失去了基督的時候就好像.
一樣東西沒有了目的.
我們小朋友那時候.
拿那些間尺.
就要來打架.
那個間尺的目的不是要來打架的.
是要來劃直線的.
我們頑皮就要來打架.
打到凹了.
就變成劃直線就劃到一癌一窟.
其實失去了它的目的.
因為它本身的目的根本也不是這樣.
它也不是跟著它的目的.
而運用也對它本身帶來了很大的破壞.
到它要去完成它目的所帶給它的使命.
就是劃直線的時候也達不到.
人有多少也是這樣呢?.
人有多少都是這樣呢?.

$^{921}$所以這個也是第二方面.
所以他說萬物都是藉著耶穌基督而被造的.
第三就是說.
他是神榮耀所發的光輝.
他是神榮耀所發的光輝就是說.
宇宙萬物對於當時的人來說.
其實他看到當然是有它的敗壞的一面.
但是也有它榮耀的一面.
我想今天的人也一樣會看到的.
是不是?.
聖經說到.
「諸天說說神的榮耀,窮瘡全陽祂的手段」.
就是神的榮耀在大自然界.
仍然都有很奇妙的設計.
但是這一切的榮耀.
本身好像是在創造物上反映了一些.
但是那些的反映.
其實它不是本身的光.
就正如太陽和月亮兩種的光是不同的.
太陽和月亮的光有什麼不同?.
太陽是自己發光的.
月亮是反射太陽的光的.
所以宇宙萬物反映著神的榮美.
有時候你看到人世間的那些美善.
那些愛.
那些美麗.
其實都是反映著神的榮美.
但是那些本身不是神的榮美.
是反映出.
不是他們自己的榮美.
是反映出來的.
但是這裡就說到.
祂是神榮耀所發的光輝.
就是說耶穌基督自己就不是月亮的光.
是太陽的光.
祂自己就是發出神的榮光.
然後他就再說.
「是神本體的真相」.
這個「是神本體的真相」.
簡單來說就是說.

$^{961}$是神自己的「essence」.
英文「essence」這一個字就是說.
神最本性的那一部分.
跟耶穌基督是一樣的.
「本體」就是那個意思.
「essence」就是神最本性的那一部分.
所以耶穌基督絕對不是天使.
絕對不是低級一點的.
很多基督教的異端.
之所以稱為異端.
就是因為他們看耶穌基督是低一點.
覺得即使耶穌基督是多麼高.
但是祂還是低於神.
這個都是錯的.
因為耶穌基督是跟神同榮同尊.
祂就是神的本體就是怎樣.
所以在今天.
我們讀新城要理問答那裡.
唯有神自己才可以將人救贖.
因為人的尊貴.
人是按著神的形象而被造的.
而人如果沒有神自己的榮耀的本體.
救贖的時候.
人是恢復不了神的那種形象的.
所以下面當然會有說「赦罪」.
但是接著他就再說.
「常用他的權能的命令托住萬有」.
這裡說到他不只是創造一切.
他也是維持一切的.
有時候有些東西.
你好像賣東西那樣.
你給了錢拿走了.
就沒有….
為什麼?沒有售後服務的.
是不是?.
那你就自己想自己.
壞了就你自己想吧.
你能不能維修?你自己想吧.
是沒有售後服務的.
但是有些東西你買了.

$^{1001}$你是不斷要訂閱的.
譬如舉例來說.
網絡服務.
你在家裡不斷每個月要給錢.
因為你要繼續維持的.
你一天不給錢.
接著就斷網了.
接著就沒用了.
有些東西是這樣的.
或者手機.
你就算你有個多好的手機.
接著你不交電話費.
就沒有維持就沒有了.
但是原來在聖經的角度來說.
原來宇宙.
整個宇宙萬物本身.
就是靠著神不斷來維持的.
神不托住.
這個托住的意思就是說神頂住.
神不頂就沒有了.
就塌了.
其實意思是這樣的.
這個宇宙萬物本身是靠神.
當然這裡.
接著他就再說.
他洗淨了人的罪.
就做了高天賜大者的右邊.
就說得夠熟了.
人和神那種的罪.
是需要洗淨.
接著以至到人和神恢復關係.
而主耶穌也都被升高.
當然今天這裡所說的.
是很多東西.
已經說了第一節到第三節.
但是正如剛才所說.
這是一個很高層次的概念.
接下來希伯來書的每一段.
都是嘗試再一次去解釋.
去將它更加詳細的去說.

$^{1041}$所以你說讀下一連串說了很多東西.
很周詳.
不要緊.
之後會再說.
但是你看到.
希伯來書的作者就看到.
他的跟從神的兒子.
耶穌基督是多麼的超越.
以至到他的生命.
是可以值得去毀身上去.
如果你看不到耶穌基督的超越.
你的生命是不會毀身的.
那些東西都不值錢.
你為什麼還要拖他頭過去.
是不是.
父母教子女都是這樣教的.
但是現在耶穌基督.
你看不到它這麼超越的時候.
你的信仰可以毀身的深度.
是很有限的.
\newpage



\section{士師記 7:1-25}
\label{sec:032274}
\textbf{士師記析讀 (9) 以少勝多}~~~~梁永善牧師
\newline
\newline
連結: \href{http://www.church.com.hk/acms/content.asp?site=cdc\&op=show\&type=product\&code=032274\&layout=sermon}{\texttt{032274}} ~~~~ 語音日期: 2026-06-10
\newline
\newline
\hyperref[sec:032273]{\small{< < < PREV SERMON < < <}}
~
\hyperref[sec:index_chronic]{\small{[返順時目]}}
~
\hyperref[sec:032275]{\small{> > > NEXT SERMON > > >}}
\newline
\newline
士師記 7:1-25
\newline
\begin{longtable}{cl}
\hline
\hline
章節 & 經文 (和合本修訂版)\\
\hline
7:1 & \begin{tabularx}{0.7\textwidth}{X} 耶路巴力，就是基甸，和所有跟隨他的人早晨起來，在哈律泉旁安營。米甸營在他北邊，靠近摩利岡的平原。 \end{tabularx} \\ \\ \relax
7:2 & \begin{tabularx}{0.7\textwidth}{X} 耶和華對基甸說：「跟隨你的人太多，我不能把米甸交在他們手中，免得以色列向我自誇，說：『是我自己的手救了我。』 \end{tabularx} \\ \\ \relax
7:3 & \begin{tabularx}{0.7\textwidth}{X} 現在你要向這百姓宣告說：『凡懼怕戰兢的，可以離開基列山回去。』」於是有二萬二千人回去，只剩下一萬人。 \end{tabularx} \\ \\ \relax
7:4 & \begin{tabularx}{0.7\textwidth}{X} 耶和華對基甸說：「人還是太多。你要帶他們下到水旁，我好在那裡為你試試他們。我指著誰對你說：『這人可以跟你去』，他就可以跟你去；我指著誰對你說：『這人不可跟你去』，他就不可跟你去。」 \end{tabularx} \\ \\ \relax
7:5 & \begin{tabularx}{0.7\textwidth}{X} 基甸就帶百姓下到水旁。耶和華對基甸說：「凡用舌頭舔水像狗一樣舔的，要使他單獨站在一處；那些用雙膝跪下喝水的，也要使他單獨站在一處。」 \end{tabularx} \\ \\ \relax
7:6 & \begin{tabularx}{0.7\textwidth}{X} 用手捧到嘴邊舔水的數目有三百人，其餘的百姓都用雙膝跪下喝水。 \end{tabularx} \\ \\ \relax
7:7 & \begin{tabularx}{0.7\textwidth}{X} 耶和華對基甸說：「我要用這舔水的三百人拯救你們，把米甸交在你手中；其餘的百姓都可以各回自己的地方去。」 \end{tabularx} \\ \\ \relax
7:8 & \begin{tabularx}{0.7\textwidth}{X} 百姓手裡拿著食物和角；其餘的以色列人，基甸都打發他們各自回到自己的帳棚，只留下這三百人。米甸營在他下邊的平原上。 \end{tabularx} \\ \\ \relax
7:9 & \begin{tabularx}{0.7\textwidth}{X} 那夜，耶和華對基甸說：「起來，下去攻營，因我已把它交在你手中。 \end{tabularx} \\ \\ \relax
7:10 & \begin{tabularx}{0.7\textwidth}{X} 倘若你害怕下去，可以帶你的僕人普拉下到那營裡去， \end{tabularx} \\ \\ \relax
7:11 & \begin{tabularx}{0.7\textwidth}{X} 你必聽見他們所說的，這樣你的手就有力量下去攻營。」於是基甸帶著僕人普拉下到軍營裡帶著兵器的人邊上。 \end{tabularx} \\ \\ \relax
7:12 & \begin{tabularx}{0.7\textwidth}{X} 米甸人、亞瑪力人和所有東邊的人都散佈在平原，如同蝗蟲那樣多。他們的駱駝無數，多如海邊的沙。 \end{tabularx} \\ \\ \relax
7:13 & \begin{tabularx}{0.7\textwidth}{X} 基甸到了那裡，看哪，有一人把夢告訴同伴說：「看哪，我做了一個夢。看哪，一個大麥餅滾入米甸營中，來到帳幕，把帳幕撞倒，帳幕就翻轉倒塌了。」 \end{tabularx} \\ \\ \relax
7:14 & \begin{tabularx}{0.7\textwidth}{X} 同伴回答說：「這不是別的，而是以色列人約阿施的兒子基甸的刀。神已把米甸和全軍都交在他手中了。」 \end{tabularx} \\ \\ \relax
7:15 & \begin{tabularx}{0.7\textwidth}{X} 基甸聽見這夢的敘述和夢的解釋，就敬拜神。他回到以色列營中，說：「起來吧！耶和華已把米甸軍隊交在你們手中了。」 \end{tabularx} \\ \\ \relax
7:16 & \begin{tabularx}{0.7\textwidth}{X} 於是基甸將三百人分成三隊，把角和空瓶交在每個人手中，瓶內有火把。 \end{tabularx} \\ \\ \relax
7:17 & \begin{tabularx}{0.7\textwidth}{X} 他對他們說：「看著我，你們要照樣做。看哪，我來到營邊，我怎樣做，你們也要照樣做。 \end{tabularx} \\ \\ \relax
7:18 & \begin{tabularx}{0.7\textwidth}{X} 我和所有跟隨我的人吹角的時候，你們也要在營的四圍吹角，喊叫：『為耶和華！為基甸！』」 \end{tabularx} \\ \\ \relax
7:19 & \begin{tabularx}{0.7\textwidth}{X} 基甸和跟隨他的一百人，在半夜之初換崗哨的時候來到營旁。他們就吹角，打破手中的瓶； \end{tabularx} \\ \\ \relax
7:20 & \begin{tabularx}{0.7\textwidth}{X} 三隊的人都吹角，打破瓶子。他們左手拿著火把，右手拿著吹的角，喊叫：「耶和華和基甸的刀！」 \end{tabularx} \\ \\ \relax
7:21 & \begin{tabularx}{0.7\textwidth}{X} 他們圍著軍營，各人站在自己的地方；全營的人都逃竄，一面喊，一面逃跑。 \end{tabularx} \\ \\ \relax
7:22 & \begin{tabularx}{0.7\textwidth}{X} 三百人就吹角，耶和華使全營的人用刀自相擊殺。全營的人逃往西利拉的伯‧哈示他，一直逃到靠近他巴的亞伯‧米何拉。 \end{tabularx} \\ \\ \relax
7:23 & \begin{tabularx}{0.7\textwidth}{X} 從拿弗他利、亞設和瑪拿西全地來的以色列人被召來，追趕米甸人。 \end{tabularx} \\ \\ \relax
7:24 & \begin{tabularx}{0.7\textwidth}{X} 基甸也派人走遍以法蓮山區，說：「你們下來迎擊米甸人，在他們的前面沿著約旦河把守渡口，直到伯‧巴拉。」於是以法蓮眾人聚集，沿著約旦河把守渡口，直到伯‧巴拉。 \end{tabularx} \\ \\ \relax
7:25 & \begin{tabularx}{0.7\textwidth}{X} 他們捉住了米甸的兩個領袖，俄立和西伊伯。他們在俄立磐石上殺了俄立，在西伊伯酒池那裡殺了西伊伯。他們追趕米甸人，把俄立和西伊伯的首級帶到約旦河對岸，到基甸那裡。 \end{tabularx} \\ \\
[1ex]
\hline
\hline
\end{longtable}
$^{1}$我們上個星期就查到機電被神蒙照.
但是它就要有很多的確據,證據.
然後才肯踏上.
我上個星期都說如果有時間.
就會分享一下我自己被神呼召.
我踏上這一條奉獻的路的經歷.
也都說 甚至有一天我離開了世界的時候.
都立此存照.
但是上個星期因為太豐富了.
所以說不出.
我今天在這個查經的後半段.
完結之前.
如果有時間就會跟大家分享.
機電被神呼召.
他在酒寨裡打麥.
他怕米顛人來搶他的東西.
但是神就跟他說.
大能的勇士.
究竟是奉賜他還是看到他的質素呢?.
後來機電就要神給他一些確據.
第一.
就是讓那個製物在磐石上.
然後磐石出火.
將那些製物帶進去燒了.
後來它要再拿證據.
就是大家很熟悉的.
就是將羊毛放在螢幕外面.
羊毛要濕.
周圍的地下要乾.
這樣就成功了.
但是它仍然缺乏這份的信心.
又再要求什麼呢?.
就是要羊毛乾.
但是地下是濕的.
我已經講過了.
是較為困難的.
因為.
早上太陽一出來的時候.
霧水就已經消散了.
如果只是羊毛濕.

$^{41}$而地下乾.
是較為困難的.
但是神都來答應了.
然後他的疑慮消除了之後.
他就願意信服神的召命了.
我們看一看今次.
我們看機電如何用很少的人.
來戰勝多餘海沙的米電人.
所以今次的題目就叫《以少勝多》.
其實今次機電真的以兔撲獅.
不是以獅撲兔.
是微小的力量.
去對付龐大的軍力.
我們看看《士師記》第七章.
我們先看第一節.
我一直都鼓勵弟兄姊妹.
如果我查警之前.
你知道我下一堂是講什麼的時候.
你先看一次會對你幫助大一些的.
《士師記》第七章第一節.
耶路巴力就是機電.
他和一切跟隨的人早晨起來.
在哈律全旁安營.
米電仍在他們北邊的平原靠近摩利干.
留意到聖經作者刻意稱他為耶路巴力.
記不記得耶路巴力是什麼意思?.
就是機電去拆毀了巴力的祭壇.
後來他們要逼機電的父親將機電交出來.
但是機電的父親愛機電.
本來他是米電人那些異教的祭師.
但這次就說到如果巴力是真的.
讓貝爾巴力和他們爭拗.
這個意思是有人稱他為耶路巴力.
聖經作者刻意稱他為耶路巴力.
然後就稱他為機電.
這就是他蒙神呼召時藉著名字,行動,爭戰.
來彰顯神的能力.
神力工作是由機電個人的家庭開始.
要他清除那些偶像.
因為他父親是米電人,拜異教的祭師.

$^{81}$但很可惜.
他沒有引以為鑒.
結果後來他的家庭卻是極度失敗.
甚至引致自己的民族都受影響.
這就是下一節和下一節我們會探討的.
我們看看機電如何去面對這班米電人.
我們看《士師記》第七章第二,第三兩節.
耶和華對機電說:.
「跟隨你的人過多.
我不能將米電人交在他們手中.
免得以色列人向我誇大說.
「借我們自己的手救了我們」.
現在你要向這些人宣告說.
凡懼怕膽怯的可以離開機電山回去.
於是有二萬二千人回去.
只剩下一萬」.
當機電呼召人來迎敵的時候.
結果有三萬二千人來參與.
我們不要否定這三萬二千人.
這三萬二千人都是很有勇氣的.
為甚麼呢?.
因為雖然三萬二千人都不少.
但他面對的敵人絕不少數.
是大大超越了他們的.
我們看一看《士師記》第七章第十二節.
「米電人,亞馬力人和一切東方人都佈散在平原.
如同蝗蟲那樣多.
他們的樂壇無數多如海邊的沙」.
聖經裡面怎樣描繪這班人呢?.
就好像描繪亞伯拉罕的后裔多到天上一星海邊的沙那麼多.
「米電人,亞馬力人和一切東方人都佈散在平原.
如同蝗蟲那樣多.
他們的樂壇無數多如海邊的沙」.
稍後我們會看到.
你知不知道他面對的是多少人呢?.
是十三萬五千人.
所以你可以想了.
雖然現在他們有三萬二千人參與.
但要對付十三萬五千人.
仍然是一個蚊比和牛比的數目.

$^{121}$但是神要的是精兵.
所以神就說現在太多人了.
免得這班人將來贏了自誇.
是我們的手來贏這一場仗的.
於是.
機電就怎樣宣佈呢?.
機電就說.
「犯懼怕膽怯的可以離開」.
結果懼怕膽怯的人竟然為數不少.
有二萬二千人離開了.
我們會不會有一個質疑.
難道他們參戰的時候沒有勇氣嗎?.
那時他們沒有膽怯的.
如果我們了解.
很多時候我們人是被當時的氣氛影響的.
又或者被朋友帶來影響的.
也都被情緒來影響.
所以可能他們沒有清楚的時候就參與了.
所以中國人有一句說話是這樣說的.
「慷慨受死而從容就易難」.
聖經也都指出.
有些人是不用當兵的.
我們看一看《申命記》第20章第5到第8節.
有沒有留意到《申命記》裡面很清楚的說.
那些人要當兵.
有幾種人是不用當兵的.
就是建了一間新屋.
但是還沒有入住的.
那就不要去了.
第二就是他種了葡萄園.
還沒有第一次的收成.
這裡是說到的.
他都可以不用去打仗.
好了 他定了婚.
預備去迎娶了.
但是還沒有去娶妻子.
那就都不用打仗了.
還有第四點是什麼呢?.
誰怕的就不要去打.
為什麼呢?.

$^{161}$因為如果建了屋還沒有入住.
又種了很辛苦的葡萄.
但是還沒有收成.
還有如花似玉的妻子還沒有迎娶.
死了就給人家享受.
於是這些人就不會勇往直前了.
那就會有所保留的.
還有膽怯的人又怎麼樣呢?.
所以會對整隊的軍隊沒有幫助.
反而會帶來壞的影響.
所以這些退場的人.
是有較多時間來思考.
為了自己的利益.
於是他怎麼樣呢?.
他就退出這一場戰役.
正如今天有不少的人.
來湊高興.
在報道會裡面.
那些大型的報道會.
在醫館,紅館.
甚至政府大球場.
願意信主的人請你下來.
可能很多人就下去了.
但是很可能他們其實是沒有經過.
很仔細的考慮,決志.
及後究竟這麼多所謂信主的人.
有多少有真正的追求.
受到好的栽培.
然後就成為一個成熟的教徒呢?.
主耶穌在世上的時候.
有人要跟隨祂.
耶穌都叫這些人.
拜託你要考慮清楚.
跟隨我的話要付代價的.
所以主期望跟隨祂的人.
要有生命力.
肯為神.
教會推動一些事工的時候.
或者一開始的時候都很多人參與的.
但是很可能沒多久.

$^{201}$有些人就會退出了.
所以記住.
神要求的是徵兵.
而不是護俠之眾.
神用少數的人去打這場仗.
目的就是防止這些以色列人.
有自高自大的錯誤的心態.
以為是自己的能力.
去戰勝這些米電人.
神要他們有謙卑和信心依靠祂.
看見的是祂的能力.
而不是自己.
所以這場戰爭.
不單單要解救以色列人.
脫離米電人的手.
脫離他們的核制.
更加重要是讓他們清楚.
他們敬拜的對象是一個怎樣的神.
剩下來的那一萬人.
算是好東西了.
一萬對十三萬五千.
絕對是勇敢不怕犧牲的.
因為要得勝的機會是很難的.
但是他們仍然堅持.
是絕對值得我們去讚揚的.
但是神仍然說人數太多了.
於是要大幅削減以色列軍隊的數目.
這個是很不合理的方法.
但是正如我剛才所說.
神需要的不是勇員.
更加不需要按人的方法去作戰.
祂要的就是徵兵.
於是用另外一個方法.
將那些人帶去西走.
在《士師記》第七章第四到第七節.
耶穌說:對姬澱說.
人還是過多.
你要帶他們下到水泊.
我好在那裡為你涉事他們.
我只染誰說.

$^{241}$這人可以跟你去.
他就可以跟你去.
我只染誰說.
這人不可跟你去.
他就不可跟你去.
既然就帶他們下到水泊.
耶穌說對姬澱說.
凡用舌頭添水.
將狗添的.
要使它單站在一處.
凡跪下喝水的.
也要使它單站在一處.
於是用手捧著添水的有三百人.
其餘的都跪下喝水.
耶穌說對姬澱說.
我要用這添水的三百人拯救你們.
將米顛人交在你手中.
其餘的人都可以各歸各處去.
神是要DQ一些人.
好了 傳統的解釋就是什麼呢?.
一去到一口口汗.
就趴在那裡.
好像狗一樣.
噁噁噁噁的舔一些水.
這些人就不得警醒了.
因為敵人來到的時候.
它起身也慢了.
有些跪下的也同樣是這個道理.
但是如果站在那裡.
然後拿些水來.
去喝那些.
那些就行了.
這是傳統的教導.
但其實我們要看一看.
神要的是什麼呢?.
是要少數的人.
所以那少數的人只有三百.
如果你反過來想一想.
如果有九千多人.
都是站在那裡拿水喝.

$^{281}$而趴在那裡喝水只有三百人.
因為神是要少的數目.
所以可能趴在那裡喝水的那三百人.
就被神揀選.
所以記住.
原來不只是在乎他們怎麼喝水的方法.
焦點是哪些人是最少的.
結果就是只有三百人.
但是三百人會不會少得可憐呢?.
怎麼可以贏十三萬五千人呢?.
直到這樣才可以彰顯到神的能力.
好了.
接著.
幾年你猜他在那一刻裡面.
有沒有內心也有一些戰驚呢?.
行不行呢?.
三百人對這麼多人.
神滿有憐憫.
神讓他再得到憑據.
而且這次不是機電主動要求.
是神主動的來給這個憑據機電.
我們看看《時詩記》第七章第九到十一節.
「當那夜耶穌說吩咐機電說.
起來 下到米顛迎女去.
因我已將她們交在你手中.
倘若你怕下去.
就帶你的僕人普拉下到那迎女去.
你必聽見她們所說的.
然後你就有膽量下去攻迎.
於是機電帶著僕人普拉下到迎旁」.
你看到神真是滿有憐憫恩典的神.
上一課我們特別看到機電要求的那些確據.
其實是神憐憫機電.
難道沒有祂不行嗎?.
仍然來好好地訓練祂.
建立祂的信心.
所以在機電沒有要求的情況下.
神就讓祂得到憑據.
讓祂有信心去面對挑戰.
神一早已經宣告了.

$^{321}$已經將敵人交在你的手中.
但是你有沒有留意到.
祂沒有叫機電立刻去作戰.
祂容讓機電稍作遲疑.
然後叫祂做一些刺探的工作.
甚至祂說出.
如果你怕的話.
你可以帶你的僕人普拉和你一起去.
於是神就讓我得到額外的情報.
而加強了祂的勇氣和信心.
機電果然沒有被神看錯.
祂真的害怕.
於是祂就帶了普拉這個僕人.
一齊靜悄悄地去米電營來刺探軍情.
但是看到那個景象.
使他們感覺到非常懼怕.
剛才我讀過第十二節.
就是這幫人散佈在平原裡面.
如同蝗蟲那樣多.
而且他們落塘無數.
好像海邊的沙那樣多.
但是雖然人很多.
但是機電竟然偷聽到.
米電軍營裡面的人的對話.
甚麼對話呢?.
我們看看13到15節.
「機電到了.
就聽見一人將夢告訴同伴說.
我作了一夢夢見一個大麥餅滾入米電營中.
到了帳幕將帳幕撞倒帳幕就翻轉傾覆了.
那同伴說這不是別的.
乃是以色列人約納斯的兒子機電的刀.
神已將米電和全軍都交在他的手中.
機電聽見這夢和夢的講解就敬拜神.
回到以色列營中說起來吧.
也會說已將米電的軍隊交在你們手中了」.
機電和他的僕人去刺探軍情.
就聽到那些米電軍隊講解他們的夢.
夢見大麥餅搗破了他們的軍營.
原來大麥餅是甚麼呢?.

$^{361}$就是很粗糙的食物.
如果稍為有錢的人.
都會用小麥來做餅.
所以如果你看《河西亞書》.
河西亞娶淫婦的時候.
是要怎樣呢?.
他不但要給錢.
因為錢不夠還要用大麥一鵝眉以半.
即是加上這些麥.
然後去換取淫亂的妻子回來.
他背後是在彰顯神的愛.
河西亞是一個窮的人.
但是你有沒有留意到是很特別.
不是一個人做夢.
我和你都會做夢.
但我和你做的夢是不會相同的.
但現在是很多很多米電人做的夢.
但他們做的夢竟然是相同的.
這豈不是神的工作嗎?.
大家一起來分享.
這個夢境令他們覺得非比尋常.
甚至這個夢境已經令他們膽戰心驚.
因為他們說這個大麥餅代表了機電.
他們了解到機電是約翰斯的兒子.
而且他們更加知道將來會和他打仗.
甚至神已經將米電和全軍都加在他們手裡.
這個夢令他們膽戰心驚.
機電雖然缺乏勇氣.
在要求他之下來刺探軍營.
但敵人似乎是一早認識了機電.
更加知道他的行動.
即是敵人的情報工作比機電來得好.
他們早已監視著這班以色列人.
知道以色列人的神.
以及在軍營裡宣告的預言.
雖然米電人的情報比以色列人的情報來得好.
但都不能令他們逃避災難.
機電得到這些確據後.
就不再猶豫.
於是回去召集這三百人和他們說.

$^{401}$今日耶和已經將米電的軍隊放在我們手裡.
但我們聽不到一件事.
在人和神的關係當中.
其實懷疑神比懼怕敵人更加有破壞性.
所以求神幫助我們.
讓我們有信心依靠神.
正正經經我們見到機電沒有信心.
相信神親口的宣告.
他是否主動要求神給人印證呢?.
但神就好像引退了.
完全再沒有主動出現了.
直到這一次.
是神主動吩咐機電去取印證.
祂只是信服而已.
但不要忘記.
祂之前都拿了羊毛這些印證.
但並沒有帶給我們甚麼實際的幫助.
直到他聽到敵人的夢境.
他反而有更大的信心.
是不是很諷刺呢?.
直接從神聽到的應許.
竟然不能夠叫機電去信服.
但從米電士兵裡聽到這樣的應許.
卻令他信服.
這豈不是今天很多基督徒的寫照呢?.
有時我們向神有很多很多要求.
但即使神答應了我們的要求.
但我們往往有時都沒有完全的信服.
反而有時外人一兩句說話.
就叫到我們貼貼束縛.
所以很可悲.
但無論如何.
機電在這一刻有一點值得我們去欣賞.
就是以神為重.
在他開戰之前他就要去敬拜神.
因為之前他是膽怯.
經常要憑據.
但現在在強弱懸殊之下.
他是有這份勇氣.
以敬畏神的心去呼召那三百人作戰.

$^{441}$但在他作戰之前.
他先去敬拜神.
他就是以神為先.
這個原則是值得我們欣賞的.
不要忘記日後.
第一個王蘇羅的失敗.
就是當他見到形勢好的時候.
就輕忽了神.
我們看看當日蘇羅要和菲律賓人打仗.
他失敗了.
但他打算抬個藥櫃出來去求問神.
誰知後來約拿丹去打贏菲律賓人.
他一聽到菲律賓人敗亡的時候.
他就不等了.
我們看看這段經文.
《撒布爾記》上14章18到19節.
那時神的藥櫃在以色列人那裡.
蘇羅對阿希亞說.
你將神的藥櫃運了來.
蘇羅正與祭司說話的時候.
菲律賓營中的軒轅越發大了.
蘇羅就對祭司說停手吧.
原來蘇羅當日很想拿藥櫃出來去求問神.
誰知他一聽到敵方陣腳大亂的時候.
他就馬上說不用了.
他就不再理會倚靠神了.
就跟著去打仗.
更加他要乘勝追擊.
及後他沒有理會士兵那種疲乏.
不准他們吃東西.
下這個命令.
之前是怎樣呢?.
原來當他要打仗之前.
他應該要等候撒布爾來為他們祈禱.
但是他等撒布爾還未到的時候.
他又見到跟隨他的人逐漸離開.
於是他就要自己獻祭.
我再講君王獻祭其實是可以的.
所羅門天聖殿的時候他都有獻祭.
但是這一刻蘇羅獻祭的目的.

$^{481}$絕對不是為了神.
而是一個政治行動.
我們看一看.
《十二記》第十三章第八到第十二節.
蘇羅照著撒布爾所定的日期等了七天.
撒布爾還沒有來到吉格.
百姓也離開蘇羅散去了.
蘇羅說把梵濟和平安濟帶到我這裡來.
蘇羅就獻上梵濟.
剛獻完梵濟.
撒布爾就到了.
蘇羅出去迎接他要問他好.
撒布爾說你做的是什麼事呢?.
蘇羅說因為我見百姓離開我散去.
你也不照所定的日期來到.
而且非利人聚集在墨沒.
所以我心裡說恐怕我沒有禱告耶和華.
非利人下到吉格攻擊我.
我就勉強獻上梵濟.
你有沒有留意到蘇羅他要怎麼樣呢?.
他要政治上的考量.
因為有些人擔心於是離開了.
要挽留百姓.
他也把責任推在撒布爾身上.
因為你遲到了.
我怎麼知道要等到什麼時候呢?.
他就說到我勉強獻上梵濟.
我再說.
他獻祭目的不是真的為了神.
是一個政治上的行動.
所以被神放棄.
但這次基地不是.
如果我看完經文的時候.
基地聽到這個軍情之後.
他呼召三百人.
但他先去敬拜神.
還有我們有沒有留意到一些細節呢?.
基甸做事往往是黑夜的.
他拆毀巴力的偶像和亞瑟拉的墓柱.
是在晚上進行的.

$^{521}$因為他不敢在日間來行事.
好了.
所以因為這件事被人稱為耶路巴力.
他要求羊毛濕或乾.
都是在晚上進行的.
他要第二天早上出去看.
現在他得到憑據.
都是在黑夜裡面進行的.
他回去之後.
就向那三百勇士.
其實我曾經想過.
這個題目叫「戰狼三百」.
不知道大家有沒有看過這部電影.
他向這三百個勇士.
好像神一樣宣告.
耶和華已經將米電的軍隊交到他們手裡.
這時候以色列人的領導者.
就是以神的話語來幫助他們.
提醒他們.
基甸讓他跟隨者.
一起承擔神的使命.
也一起分享勝利的光榮.
好了.
接著我們看到二小聖多.
我們很熟悉的一段經文.
稍為長一點.
麻煩Yvonne幫我們讀出來.
「於是基甸將三百人分作三隊.
把國和空瓶交在各人手裡.
瓶內都藏著火把.
吩咐他們說.
你們要看我行事.
我到了營的旁邊怎樣行.
你們也要怎樣行.
我和一切跟隨我的人吹角的時候.
你們也要在營的四處吹角.
喊叫說也會說和基甸的都.
基甸和跟隨他的一百人在三更之初.
才半更的時候來到營旁.
就吹角.

$^{561}$打破手中的瓶.
三隊的人就都吹角.
打破瓶子.
左手拿著火把.
右手拿著角.
喊叫說也會說和基甸的都.
他們在營的四處各站各的地方.
全營的都亂數.
八百人吶喊使他們逃跑.
三百人就吹角.
也會說全營的人用刀互相擊殺.
逃到西尼拉的伯哈斯塔.
直逃到靠近塔巴的巴米荷拉.
我們看到他們兵分三路.
每隊各一百人.
首先吹號角.
製造一個很大的響聲.
隨即將藏在火把裡面的瓶打碎.
然後突然之間.
看到有很多火光出現的時候.
就引發敵軍惶恐驚亂.
隨即大聲喊叫.
耶和華與基甸的都.
代替了角聲.
讓我們想起當日.
記不記得.
他們入迦南地的時候.
繞著耶利哥城多少次?.
是十三次.
然後大聲呼喊就得勝.
所以基甸就活出了摩西和約書亞所承傳的信仰.
留意.
也不是這三百人來和去擊殺那些軍隊.
請大家再留意第二十二節.
耶和華使全營的人用刀互相擊殺.
原來耶和華與基甸的都.
變成了敵軍的刀.
米殿人與他們的同盟互相殘殺.
當然觸發這件事的是誰?.
就是神.

$^{601}$絕對不是基甸.
以色列人用的刀是做什麼?.
是追殺逃亡的敵人.
三百個以色列人.
竟然可以追殺十三萬五千人的軍隊.
是以小勝多.
以兔損失.
真是難以想像.
但越難以想像的事情.
就越顯示神的能力.
跟著我們看看.
那些人一直處走的時候.
怎麼去追呢?.
我們看看《時史記》第七章23到25節.
「以色列人就從那佛塔利阿徹和嗎哪西全地聚集來追趕米殿人.
米瑪基甸打發人走遍以法連山地.
說:你們下來攻擊米殿人.
爭先把守約但河的渡口直到伯巴拉.
於是以法連的眾人聚集.
把守約但河的渡口直到伯巴拉.
抓住了米殿人的兩個首領.
一名俄納一名西伊伯.
將俄納殺在俄納半石上.
將西伊伯殺在西伊伯酒鎮那裡.
又追趕米殿人.
將俄納和西伊伯的首級帶過約但河到基甸那裡」.
我們看到始終三百人很難完全追那些逃亡的敵人.
於是基甸呼召了.
拿伐塔利,阿徹,嗎哪西,以法連人.
參與追殺這些敵軍.
這些支派當然樂於參與.
因為那時勝負已經很明顯出來了.
所謂「執個尾彩」要去追殺他們的時候.
最後他們將米殿人一名首領西伊伯.
殺死在西伊伯酒鎮那裡.
另外一個就殺死在盤石上.
我們看到被殺在酒鎮那裡.
都有一個很特別的描繪.
不要忘記當日基甸是很害怕的.
在酒鎮那裡來來去去打墨.

$^{641}$但現在酒鎮就成為處死敵軍首領的公開地方.
這個酒鎮竟然成為基甸得榮耀的地方.
以色列人勢力弱,情報差.
但沒有阻止他們靠著神得到勝利.
米殿人在兵力,在情報上佔優.
但他們不能夠歇止被殺的忠告.
神的事工上是需要資源的.
甚至各項資訊的幫助.
但我們不要將我們的注意力,信心.
集中在這些事上.
因為這樣的時候就會偏離正路.
甚至我們沒有跟隨神的引導.
將信心放在一旁.
我下一年年初就要幫傑尼索羅院搞一個課程.
大家都為我祈禱教我怎樣講道.
我相信神給我有機會.
我第一樣提醒年青一代的讀者.
不要用人工智能去寫講章.
我不是很懂得用人工智能.
到現在我都沒有.
我只是來不久用查GPT去旅行.
可以怎樣編排.
但越懂電腦的人.
越懂得用這些就越危險.
因為很多人都說.
牧師傑不用你自己做PowerPoint.
現在AI可以幫你做.
我到現在都要自己做.
我不懂得做.
所以PowerPoint很窄.
但你用這些的時候.
你沒有好的預備.
沒有好的等候.
沒有好的依靠神.
往往你就會怎樣呢?.
只是靠這些數據.
你沒有辦法讓神在你身上工作.
可能會花少了很多時間.
我今天寫再下一篇的講章.
就是講機電的兒子怎樣互相殘殺.

$^{681}$我由今早10點鐘.
一直寫到晚上6點鐘.
中間停了一段很短時間.
吃下午飯.
只是寫完講章.
未做PowerPoint.
但我會不斷去看書.
等候.
自己去發掘.
所以記住.
我們有些事是需要幫助的.
需要別人幫助我們.
但不要將信心放在那裡.
在聖經裡.
我們終於看到機電那種領導能力.
他用神的話語.
作為召喚以色列人上戰場.
當我們向神祈求的時候.
神很可能用我們來完成.
並且應運了我們的禱告.
所以禱告被應運.
很多時是來自禱告者本身的行動參與.
正如我們經常埋怨教會沒有.
那樣不行.
你求神改變某些事.
但當你去祈禱求神的時候.
你會不會想到你就是被神揀選.
去改變這些事呢?.
我自己很深感受一句說話.
與其奏作黑暗.
不如焚燒自己.
正如今天有些人說教會沒有兒童工作.
沒有人來參與.
我剛才在查經之前.
跟我一班弟兄姊妹去吃飯.
因為我們明天同一位.
過去的女傳道離開了.
來了英國.
我們一起來見面.
在英國的肢體都發出一些怨聲.

$^{721}$說甚麼呢?.
聚會的人多.
但參與事奉的人少.
很多人都期望有兒童工作.
有青少年工作.
但邀請他們來參與幫忙的話.
個個都撒手擰頭.
這樣就很可悲.
所以求神來幫助我們.
與其奏作黑暗.
不如焚燒自己.
你去禱告.
但神可能用你這位禱告者去建立生命.
開這場戰爭.
我們看到神很奇妙的作為.
竟然可以以小勝多.
但最重要是整個事情對我們.
看到神長久的忍耐和恩賜.
神不斷的忍耐基甸的膽怯.
去要求取得憑據.
但神仍然滿足到他.
去栽培基甸成為拯救以色列的士師.
神一開始的時候.
稱基甸為大能的勇士.
那時候他膽怯.
要求憑據.
但磨練了之後.
他終於有所突破.
所以成功並非一就而成.
我們都要有忍耐.
面對一些人的軟弱成長.
當然我們要有智慧去分辨.
因為很可能那些軟弱或者不足.
是會影響大局.
所以同樣我們都要求神忍耐我們.
但我們自己都要積極地去追求.
當我們學習的時候.
不要單單注視那個的凶險.
我們要轉眼仰望神.
求神來幫助我們.

$^{761}$我們更加要學習.
從悲觀裡面找到樂觀.
從失望中找到希望.
在危機之中我們找到機會.
願意了解神的計算方程式.
用信心去面對.
今天我們看這場戰爭已經完結了.
但我答應了一些人.
有些人問我.
牧師又說你講你蒙召的經歷.
我回答.
我十幾歲的時候.
我想十六,七歲已經開始感受到.
神對我的呼召.
但是我一直都不理會.
直到有一次.
我去大會堂音樂廳.
參加一個聚會.
好像是建道神學院的聚會.
昨天聚會完的時候.
那些人就呼召了.
有牧者呼召.
有誰很深感受到.
呼召你獻身做傳道人的時候.
請你走出來.
那時候我就像機電一樣.
我都可以的.
平時覺得自己很厲害.
但是又覺得我又喜歡玩又喜歡開玩笑.
不如做傳道人吧.
神就叫隔壁那個.
隔壁那個比較好.
豆紅色隔壁那個.
最重要的是甚麼呢.
因為那時候已經出來工作了.
十八,九歲左右.
我只是幫一個人的補習.
是一家銀行董事長的女兒.
那份工資已經等於一個人.
平時的文職人員的一個月工資.

$^{801}$還有給我的印象就是做傳道人.
又沒有收入.
又被人看不起.
我是神我寧願奉獻.
那一刻我是沒有決志的.
沒有出去.
我特別為自己說.
可能是周圍的氣氛影響.
當我離開的時候.
離開大會堂.
如果你看回這些圖片.
大會堂音樂廳對出的時候.
還在皇下碼頭.
下著微思細語.
但我心裡面非常不安.
還在這一刻出現了一句話.
就是聖經的話語.
「人若知道行善卻不內行.
這就是他的罪了」.
不是那一刻聽到的.
也不是那一天靈修的.
神就用這句話提醒我.
初時以為一些很負面的東西.
去猥瑣,看色情的東西.
那些才是罪.
但原來你知道什麼是對的.
而不去做.
這就是罪.
我記得那天晚上我回到家.
我跪在床上.
來祈禱和神說.
神啊.
呼召我.
可不可以讓我看到.
傳福音的重要呢?.
這是我對神的要求.
但還有一件事是什麼呢?.
就是要很實際.
我要在大會堂音樂廳.
再一次的站出去.

$^{841}$讓你站出去.
奉獻我自己.
押門之後.
第一.
我會知道傳福音的重要.
傳福音的重要.
這些只是藉口.
但第二個就是很實際.
我要在大會堂音樂廳裡面.
來奉獻.
押門之後我就告訴自己.
我以後都不會再去大會堂音樂廳了.
好了 怎知道在那一年裡面.
我是第一次回到中國.
當我回到中國的時候.
我在香港出生 在香港長大.
那時候是文革的後期.
我記得我在深圳一坐車.
預備上廣州.
我看到綠油油的田地.
我到現在都記得.
我在那裡流眼淚.
不知道為甚麼會流眼淚.
這個就是我的國家.
但那時候國內是非常悲慘的.
灰色的衣服或者藍色的衣服.
要糧票.
要肉票.
那些人很窮困.
當我上到去的時候.
我就住在別人的家裡.
晚上我都撕了一本聖經.
一卷書就是《馬太福音》和《傳道書》.
我挑燈夜讀.
怎知道被我的會友.
他的家人問這是甚麼.
很害怕 說不說呢.
因為那時候破四九.
如果講信仰就很有問題.
但我又記得勝經說.

$^{881}$如果我們在人的面前不認神.
神都不認我們.
死就死吧 跟他說.
這個是聖經來的.
是講耶穌的.
原來我們沒有聽過.
有些人說聽過.
但不知道是甚麼來的.
但最令我感動的是.
他拿了我的聖經手抄.
繼續看.
然後在那次裡面.
我跟他們講了福音.
但其實我不知道他們是否明白.
我又不明白他們說甚麼.
他們未必明白我們說甚麼.
我們在海珠廣場那裡.
圍著一個圈來講話.
為甚麼呢.
可以知道有沒有人走過來.
走過來我們就去看.
剛剛最近.
就和我以前小年團的團友.
來大陸去分享.
這個是早幾天我們的WhatsApp.
剛才我不是說.
我和冥恩堂一位離開了女傳道.
來大陸去吃飯的.
他現在就在報道會服侍.
而我這個小年團的團友.
現在就在這個教會裡面.
他就這樣說.
我們剛剛談到.
我記得我們十幾歲的時候.
膽粗粗的一起去廣州.
想向親人傳福音.
他說前幾個月我去廣州.
探望最小的表妹.
她談起往事.
突然說起你的名字.

$^{921}$記不記得你啊?.
還說記得你啊.
不知道你記不記得我舅父和我老婆呢?.
我當然記得 多謝他們不記得我.
為甚麼呢?.
很深印象.
因為那時候.
你一去到一些人家裡.
即刻有關愛隊.
有人來看你要查你是甚麼人.
所以我們不可以在家裡面說話.
我那時候和他一起踩單車.
所以我的單車技術很高.
一邊踩單車一邊講福音.
所以這個經歷是很深的.
這個是美好年輕的回憶.
這是第一點.
甚至我看到中國.
我的同胞需要福音.
第二是甚麼呢?.
就是我參加了宣導員.
一些六年級的畢業營.
去營地的.
我的孫女今天剛去了畢業營.
那時候原來畢業營裡面.
那些小朋友是東頭村十八座.
天台學校叫智江小學的小朋友.
我是1973年第一次參加.
第二次就是1974年.
所以你見到這本聖經.
現在我正在用的《死海古卷》.
就是1974年6月5日.
宣導員送給我的.
我來帶去.
那些小朋友很喜歡聽福音.
我未起床已經敲門.
要我來帶去和他們講福音.
接著我就去到東頭村.
為他們建立團契.
但其實我還未讀神學.

$^{961}$但看到他們的需要.
但其實我不懂得做導師.
但勝在有錢.
又請他們吃東西.
又請他們看電影.
所以很多人來跟隨.
現在也有人做了傳道員.
他們當中有一個.
神答應了我一個很騙他的禱告.
就是我去看廚房有沒有緊要.
但他大會堂音樂廳.
我不去.
當時宣導會想聯合這麼多間堂會.
有些友誼的互動.
記得剛剛旺角球場開了沒多久.
我們租了旺角球場有運動會.
世運會.
說我為了要為自己的團體增光.
我早上去跑馬地.
去練跑.
走一千六.
因為最後是一千六米.
因為我是最後一個.
很多人走來走去.
我跑到抽筋都完了.
所以有些人鼓勵我.
但發現原來這樣幫助不了人們彼此交往.
於是當年宣導會.
舉行了一個音樂會.
請一些牧者去唱歌.
因為我是北宣的.
我在團體裡都活躍.
所以北宣的牧師就叫我.
梁永善 你今晚帶隊去大會堂.
這個聚會裡面去收奉獻.
大會堂?.
我問今晚的內容是怎樣.
那些牧師唱歌而已.
我有沒有講到.
他說沒有.

$^{1001}$那我就去吧.
誰知今晚的唱歌很差.
我不太懂得欣賞.
最後結束的時候.
一位叫容保羅先生的時候還沒有做牧師.
在完結之前他就說.
好 現在我們先唱一首歌.
就是《萬世戰爭》.
如果當中你心裡有感動的.
神父召你做傳道人的時候.
你就出來吧.
我馬上說 不是吧.
今晚哪有感動.
這麼難聽的歌不好玩.
但是這一首歌.
唱到這幾句.
寶貴光陰瞬即過去.
世界終局將臨.
是主機會快將消逝.
現在趁時救人.
神就說你現在還想怎樣?.
你所要求的我已經給了你.
你還不出來?.
我記得我流著眼淚出去.
奉獻自己給神.
接著一年之後 1975年.
我就進了建築神學院.
讀我的神學.
現在50年了已經過了.
1979年開始事奉神.
我很清楚神的呼召.
初時我覺得奉獻可能是一個損失.
但是現在很肯定一件事.
絕對是一個祝福.
這個就是我自己要奉獻的時候.
要向神拿的憑據.
神都憐憫我.
給我機會.
祂對我的忍耐我要感恩.
好了下個星期.

$^{1041}$我們就會去看第八章.
啊不是.
Sorry不是下個星期.
請大家為我祈禱.
因為接著我就會上蘇格蘭去廣道.
接著我有一個星期的休息.
蘇格蘭的網絡較為弱一些.
所以我們會停一個星期.
是再下個星期.
我們會看《時事記》第八章.
我們經常看《機電的成功》.
我們忽略了《機電的失敗》.
我們就會看《機電的失敗》.
失敗之父是誰呢?.
希望大家記住.
就是6月24日.
我們就會看《機電》.
我們會看第八章.
希望大家可以先看.
今晚結束的時候.
和大家祈禱之前.
我們都會聽一首歌.
就是剛才我所說的.
《萬世戰爭》.
現在我們都是在一個屬靈的爭戰裡面.
魔鬼都要誘惑我們.
但我們靠著神要得勝.
和事主機會快將消逝.
現在是趁時救人.
我們一起來去祈禱.
天父我們多謝你.
我們看到你奇妙的作為.
機電只是帶著三百人.
卻可以戰勝十三萬有多的人.
上天原諒我們知道不是依靠勢力.
也不是依靠才能.
而是依靠神你的靈方能成事.
幫助我們.
很深的感受到神你對我們的愛.
對我們的忍耐.

$^{1081}$其實在事奉的歷程上.
不是沒有我們就不行.
但神你仍然讓我們有機會.
去服侍你.
讓我們懂得用信心去擺上.
正如剛才我們所唱的那首歌.
我們所聽的.
原來寶貴光陰.
很快就會過去.
我們可以服侍你的機會.
越來越少.
讓我們珍惜不要再拖延.
在你面前去領受.
你華語的弟兄姊妹.
無論是現場的.
又或者在YouTube裡聽的.
都祝福我們.
讓我們能夠靠著你.
我們有這份信心勇氣.
去面對挑戰.
更加願意積極地服侍.
聽下者面前的祈禱交托.
來奉主德性的名.
\newpage



\section{列王記上 9:3-9}
\label{sec:032275}
\textbf{華麗殿宇背後:神看重的不是建築,而是內心}~~~~袁惠鈞牧師博士
\newline
\newline
連結: \href{http://www.church.com.hk/acms/content.asp?site=cdc\&op=show\&type=product\&code=032275\&layout=sermon}{\texttt{032275}} ~~~~ 語音日期: 2026-06-07
\newline
\newline
\hyperref[sec:032274]{\small{< < < PREV SERMON < < <}}
~
\hyperref[sec:index_chronic]{\small{[返順時目]}}
~
\hyperref[sec:032276]{\small{> > > NEXT SERMON > > >}}
\newline
\newline
列王記上 9:3-9
\newline
\begin{longtable}{cl}
\hline
\hline
章節 & 經文 (和合本修訂版)\\
\hline
9:3 & \begin{tabularx}{0.7\textwidth}{X} 耶和華對他說：「我已聽了你在我面前的禱告和祈求，將你所建的這殿分別為聖，使我的名永遠立在那裡；我的眼、我的心也必時常在那裡。 \end{tabularx} \\ \\ \relax
9:4 & \begin{tabularx}{0.7\textwidth}{X} 你若以純正的心和正直行在我面前，效法你父大衛所行的，遵行我一切所吩咐你的，謹守我的律例典章， \end{tabularx} \\ \\ \relax
9:5 & \begin{tabularx}{0.7\textwidth}{X} 我就必堅固你在以色列國度的王位，直到永遠，正如我應許你父大衛說：『你的子孫必不斷有人坐以色列的王位。』 \end{tabularx} \\ \\ \relax
9:6 & \begin{tabularx}{0.7\textwidth}{X} 倘若你們和你們的子孫轉去不跟從我，不守我擺在你們面前的誡命律例，去事奉別神，敬拜它們， \end{tabularx} \\ \\ \relax
9:7 & \begin{tabularx}{0.7\textwidth}{X} 我就必把以色列從我賜給他們的地上剪除，也必從我面前捨棄那為我名所分別為聖的殿，使以色列在萬民中成為笑柄，被人譏誚。 \end{tabularx} \\ \\ \relax
9:8 & \begin{tabularx}{0.7\textwidth}{X} 這殿雖然崇高，將來凡經過的人必驚訝，嗤笑，說：『耶和華為何向這地和這殿如此行呢？』 \end{tabularx} \\ \\ \relax
9:9 & \begin{tabularx}{0.7\textwidth}{X} 人必說：『因為此地的人離棄領他們祖先出埃及地的耶和華－他們的神，去親近別神，敬拜事奉它們，所以耶和華使這一切災禍臨到他們。』」 \end{tabularx} \\ \\
[1ex]
\hline
\hline
\end{longtable}
$^{1}$請姐妹平安.
很高興又有機會來到你們教會.
和你們分享神話語.
今天我要講的是.
列王記合一王朝的第六講.
是列王記上的第六章和第七章.
剛才主席也有讀到我的題目.
就是華麗建語背後.
神看重的不是建築 而是內心.
我在結束的時候會大家一齊回來看.
我們應該怎樣回應神呢?.
因為我們的身體都是神的電.
現在我們見到所羅門起聖殿.
我們又應該怎樣來回應呢?.
其實今天的經文.
在列王記的第一章至十一章.
是很中心的經文 是很重要的經文.
我時時都說聖經裡面有很多首尾的結構.
而首尾結構最中心的地方就是最重要的經文.
今天在列王記上的第一章至十一章.
都有一個首尾結構 一共有五層.
第一層就是講到所羅門.
其實整個結構都是有關所羅門.
一開始就是講所羅門的登基.
列王記上的一章至二章十二節.
而在十一章的尾就是講到所羅門去世.
所以有沒有見到一頭一尾都是講所羅門.
這是第一層.
第二層在列王記上的第二章尾.
講到所羅門鞏固他的國位.
而在樹下半部的首尾結構.
在列王記上的十一章十四至四十節.
是講到所羅門的皇朝的衰弱.
有人說不妥 牧師.
上面是講所羅門鞏固皇位.
下面是講他的國位衰退.
好像是相反的.
其實如果我們注意作者.
早在列王記上的第二章尾已經給我們看到.
神鞏固了所羅門的國位.

$^{41}$但所羅門仍然用他自己的手法.
自己的智慧排除異己.
這是不被神所喜悅.
所以去到十一章的尾.
神是差遣所羅門的敵人.
以色列的敵人來攻擊以色列國.
這個結構的第三層.
就是所羅門求智慧.
列王記上的第三章.
而在下半部是列王記上的十章二十六至十一章十三節.
是講到所羅門愚昧的行為.
他是追隨他的妃嬪去拜偶像.
以至是帶來神的審判.
而全部經文最中心的地方.
其實就是列王記上的五章至九章九節.
部分的也是我們今天所講.
所羅門起聖殿奉獻聖殿.
今天的大綱很簡單.
一盤扑兩點.
所羅門起聖殿和所羅門佈置聖殿.
當他起聖殿的時候.
這個聖殿是有外殿和內殿.
而他佈置聖殿.
聖殿裡面有很多的陳設.
都是祭司需要來敬拜神,事奉神.
我們會看這些陳設.
明眼的聽眾一聽就會知道.
咦?牧師你好像講漏了一些經文.
如果你注意的大綱.
我是沒有寫到列王記上的七章一至十二節.
是講甚麼的經文呢?.
是所羅門起自己的皇宮.
其實作者有意讓我們看到.
所羅門起皇宮對所羅門是很重要.
這就讓我們看到他其實優先次序出現了問題.
以致他最後跟隨他的妃嬪去拜牛像.
第一點我首先要講的就是所羅門建造聖殿.
在列王記上的六章一節.
「以色列人出埃及地後四百八十年.
所羅門作以色列王第四年西伐月.

$^{81}$就是二月開工建造耶和華的殿」.
有沒有看到作者將所羅門建造聖殿.
和出埃及連在一起?.
這個非常重要.
有兩個重要的點.
第一點就是讓我們知道以色列人出埃及是何時.
很多人都說以色列人出埃及沒有記載.
究竟是何時?.
經文講得很清楚.
是所羅門登基之後的第四年.
而之前的四百八十年他們出埃及.
也就是說以色列人出埃及.
是公元前的一千四百四十七年.
這是很重要的一個日子.
作者將以色列人出埃及和所羅門建造聖殿.
連在一起第二個原因.
就是讓我們看到所羅門的問題.
甚麼問題?.
埃及是代表甚麼?.
是代表罪的綑綁在聖經裡面.
記不記得神差派摩西帶領以色列人出埃及.
拯救他們脫離罪的綑綁.
神拯救了他們之後.
就和以色列人.
尤其是以色列的王說.
千萬不要回去這條路.
不要走回埃及.
這條是綑綁之路.
你不要走回去拜偶像的地方就是埃及.
但很可惜所羅門和跟著他的王.
不斷走回這條路.
他們不斷和埃及立約建文約.
他們不斷走回埃及招兵賣馬.
就好像所羅門買了很多名貴的馬匹.
就是靠這些馬匹來征戰.
而不是靠耶和華的力量.
以致最後他們被亞述國,巴比倫國亡國.
來將他們擄掠.
而真的走回那條路.
當他們被亞述國,巴比倫擄掠的時候.

$^{121}$就好像走回埃及遺老一樣.
今天我們說所羅門建聖殿.
其實是神的恩典.
神曾經應許在《申命記》的十二章.
以色列人將會在一個納神明的地方來敬拜神.
現在神就給了這個應許給所羅門.
准許他建聖殿.
你說建聖殿這個設計從何而來?.
其實是從神而來的.
在歷代志上的二十八章.
十一至二十一節講得很清楚.
神甚至將這個聖殿建造的藍圖給了大衛.
日後大衛就將這個設計給了所羅門.
所以所羅門才能夠建聖殿.
在希伯來書的八章五節.
在新約很清楚.
「他們供奉的是本是天上時的形狀和影像.
正如摩西將要做帳幕的時候.
蒙神警誡他說.
你要謹慎作各樣的物件.
都要照著在山上指示你的樣式」.
有沒有看到當摩西建會幕的時候.
全部是照著神的旨意來建會幕.
現在歷代志上的二十八章.
也是一樣告訴我們.
連聖殿的設計都是從神而來.
這個不是說我們聽.
將來去到天上就有個會幕.
有個聖殿將來的天堂.
因為啟示錄的二十一章,二十二節.
說在天堂是沒有聖殿的.
但是經文是告訴我們.
一切會幕也好,聖殿也好.
的設計就是照著天上的形象來做.
我也要講講有很多自由派的神學家會說.
根本以色列人的聖殿是抄那些俗世的廟宇.
因為他們找到有很多的廟宇.
都是有一個聖所.
而聖所裡面有更加聖潔的地方.
就好像致聖所一樣.

$^{161}$所以有很多的學者.
有很多的自由派的學者都說.
喂 你又說以色列人的聖殿是照著神的設計?.
不是啊 是照著俗世的廟宇來建.
但是你們等一下會看到.
聖殿的設計其實和會幕是一模一樣.
只是尺寸有些不同.
如果聖殿是照著會幕的設計.
會幕是甚麼時候有的?.
是所羅門建聖殿之前的480年.
我剛才有講過.
就是公元的1447年.
而以色列人出了埃及之後.
第二年就建會幕了.
也就是說公元前的1446年已經有會幕.
而是早過了所有這些俗世的廟宇.
所以絕對不可以說以色列人是抄俗世的廟宇.
我講所羅門建聖殿.
首先是要講外殿.
然後我會講內殿.
我平時在YouTube講到.
很多人就這樣聽到.
今天我真的很鼓勵大家.
去下載一個PowerPoint來看.
因為我將會有很多的圖片.
如果看到這些圖片.
就會更加明白這個聖殿的設計.
以色列人首先用甚麼標準來量度?.
是爪子.
一爪是多少?.
由一個人的手指尖到手肘就是一爪.
他們就是用這樣的方法來量度.
以色列人來說一爪就是18吋.
就是一尺半.
我長得很小.
我就去量一量我自己的一爪有多少.
原來只有16吋半.
所以找我來量度就多很多爪.
如果你找一個巨人.
而的確以色列是有巨人.

$^{201}$找一個巨人來量度.
這個長爪有紀錄就是21吋.
但是聖殿就是照著普通的爪.
就是一爪18吋一尺半.
那你說聖殿的建築有多大?.
一共是90呎長.
30呎闊.
45呎高.
在加州我們的房子.
一個鎖的尺寸是多少?.
大概是60呎乘120呎.
所以你要裝這個聖殿能不能裝?.
裝得到.
這個聖殿其實不是那麼大.
90呎乘30呎.
它有一個走廊是30呎闊.
是有15呎深.
在這個建築物的前面.
是有一個內院.
是給祭司來作供.
這個內院和外院是隔開.
是有石頭和木的牆來隔開內院和外院.
而在耶利米蘇的第36章第10節.
叫祭司的院做上院.
為什麼?.
因為祭司的院的地方是比較高一點.
是高過外院.
所以也叫做上院.
有一件很重要的就是無論.
昔日的回望或者是聖殿.
門都是向著東面.
這個很重要.
我們稍後會回來看這一點.
回望和聖殿不同的地方.
就是回望是用帳篷做的.
那些帳篷很薄的.
但是聖殿呢.
你看到它是有牆的.
在南邊,西邊和北邊的牆.
都是有暗格.

$^{241}$在暗格裡面是有房的.
還是有三層的房.
每一層的房都是七尺半那麼高.
最高那層的房是十尺半那麼闊.
第二層就沒有那麼闊了.
去到九尺.
接著最低那層是只有七尺闊.
這些房是用來儲藏物品.
在南邊的牆.
是有一條通道可以通到下面的地下室.
那個地下室.
然後由地下室.
你可以有一個旋轉的樓梯.
可以走回上去中層和上層.
房和房中間都有通道.
在三層房間最高的那個地方.
在北邊的牆和南邊的牆.
都有一些很窄的窗.
是能夠透過一點點的光線.
進入這個聖所裡面.
真的不夠光.
昔日的會幕呢.
簡直是沒有窗的.
那你說祭司怎麼能夠見到東西在會幕裡面.
記不記得有個金燈台.
就是靠著金燈台的光來事奉神.
而在這個聖所裡面呢.
一共是有十盞的金燈台.
五盞是靠著北邊的牆.
五盞是靠著南邊的牆.
我們看到這個聖殿很宏偉.
全部的牆都是用石頭做成了磚.
然後去到聖殿裡面是砌在一起變牆.
所以不是在聖殿裡面做這些磚.
當時砌這個聖殿的時候.
砌聖殿的牆.
根本是聽不到有鐵器的聲音.
因為全部已經做好了磚頭.
進去砌成牆.
當所羅門完成了這個建築物.

$^{281}$就在這個屋頂鋪上了香柏木.
成為了屋頂.
在聖殿的入口如果大家注意的話.
是有兩條的銅柱.
當所羅門說要建聖殿的時候.
不僅向推羅王希蘭拿了很多的香柏木.
還請了一個是考手的工匠.
這個就叫做護蘭.
這個人就設計了那兩條的銅柱.
這兩條的銅柱高二十七呎.
圓周是十八呎.
光是青銅的部分的厚度已經是四吋.
這兩條的銅柱是空心的.
在柱頂上如果大家留意的話.
是有杜鑽的花紋.
好像是一個杜鑽的碗.
上面是有蓮花,有石榴的雕刻.
非常之美觀.
所羅門叫這兩條銅柱是有名的.
在北面那條銅柱.
如果你對著聖殿入口的話.
右邊那條銅柱就叫做瓦根.
瓦根的意思就是神建造的意思.
左邊那條銅柱在南面.
就叫做波亞斯.
記不記得波亞斯這個人?.
是誰?.
是路德的丈夫.
波亞斯的名字的意思就是神裡面的力量.
這兩條銅柱都提醒以色列人.
神平安帶給他們力量.
但是這兩條銅柱都提醒以色列人.
就是神和大衛立的約.
你說為什麼這麼奇怪?.
那條叫做瓦根的柱的意思是什麼?.
是神建立.
建立什麼?.
大衛的王朝.
這就是神和大衛在撒母耳下第七章.
大衛之約.

$^{321}$神要建造大衛的王朝.
而波亞斯的名字.
神的力量.
以色列人就是靠神的力量.
來保持維持這個約.
直至到永遠.
所以這兩條柱都提醒了以色列人.
神和大衛所立的約.
我們剛才讀的經文告訴我們.
耶和華的話又臨到所羅門.
耶和華的話什麼時候第一次臨到所羅門?.
記不記得神問所羅門.
你希望得到什麼?.
所羅門就說希望得到智慧.
這個是第一次.
那你說為什麼神現在第二次又來找所羅門?.
有很多學者說.
起聖殿不是一件容易的事.
所羅門可能沮喪.
所以神要來鼓勵他.
但是可能有更加好的解釋.
就是當所羅門起聖殿的時候.
他沾沾自喜.
很驕傲.
所以神是要來提醒他要謙卑.
這個是比較合理.
因為兩次無論第一次神給智慧給所羅門.
這次又向所羅門顯現.
都是要所羅門馴服.
而馴服也是聖經裡面.
很重要給我們基督徒.
保羅時是在他的寫作都提醒我們要馴服.
例如在以弗所書的第六章第六節.
有很多人會說.
喂!大衛之約不是無條件的嗎?.
無條件的約為什麼神要所羅門服從?.
為什麼又要所羅門的子孫服從?.
無條件的約啊!.
喂!其實神和大衛立的約.
像不像神和我們立的約?.

$^{361}$是不是靠我們的行為得救?.
當然不是.
就好像亞伯拉罕的約.
神怎樣稱亞伯拉罕為二人?.
當亞伯拉罕相信.
神就稱他為二人.
不是因為他有行國禮.
不是因為他做了任何的事.
我們和神立的約都是一樣.
不是因為我們做了什麼事.
是因為我們相信.
但是你說是不是我們的行為完全不重要?.
不是的.
不要忘記還有摩西之約.
其實我們和神的關係都一樣.
我們就是亞伯拉罕屬靈的子孫.
我們就像亞伯拉罕一樣.
恩信稱義.
但是當我們不服從神的時候.
神會管教他所愛的兒女.
希伯來書的十二章.
這個就是摩西之約.
而保羅在舍加拉太書的三章十七節.
講得很清楚.
在四百三十年前神和亞伯拉罕立的約.
不會被四百三十年後摩西之約廢掉.
也就是說我們會不會因為我們的行為失去我們的救贖?.
絕對不會.
如果我們真的相信的話.
不會.
但是當我們犯罪的時候.
神就會管教.
這個也是原因在列王記裡面我們會看到.
神也是管教以色列人.
將他們驅逐出去應許之地.
聖殿被毀.
耶路撒冷被毀.
就是將他們驅逐.
這個就是摩西之約.
其實是很可惜的.

$^{401}$當所羅門建造好聖殿.
將聖殿獻給神的時候.
我們將會看到神的榮光照亮了這個聖殿.
淋到這個聖殿那裡.
但是最後也因為以色列人的反叛.
最後神的榮光離開了這個聖殿.
在以西結書的第十章.
最後聖殿也被毀.
說完了外殿的結構.
我們現在看看內殿的結構.
這個是非常非常重要.
因為祭司就在內殿來事奉神.
而作者描述內殿的時候.
是說得非常輝煌,非常宏偉.
無論天花板,地板都是用香柏木鋪在上面.
上面都挑黑了有野瓜和開了花朵.
非常美觀.
而且地板再用松木鋪在上面.
你說香柏木不好嗎?.
為什麼又要鋪松木?.
香柏木就比較硬.
松木就比較軟熟.
走上來是比較舒服.
而無論殿內的牆,地和天花板.
都是鋪了金,包了金.
有兩道扇門是摺疊門.
是容許祭司由祭司的院進入聖所.
這兩道門都是用橄欖木來製造.
上面是貼了金.
無論是門書.
即是那些hinges或者是門鎖.
都是用金做.
上宮還掛上了今年非常非常之的名貴.
去到聖所最盡頭.
我剛才說聖所的門是向著東邊.
你由東邊進入聖所.
走到最盡頭就是走到西邊.
西邊的盡頭是什麼?.
就是致聖所.
你未進入致聖所之前會看到有一個萬紙.

$^{441}$就是耶穌基督為我們釘十字架的時候.
這個萬紙由上至下是分開了.
致聖所的尺寸是30尺乘30尺乘30尺.
是一個立方形的一個cube.
昔日的會幕也有致聖所.
但是尺寸只有一半.
15尺乘15尺乘15尺.
我剛才說這個聖殿牆是90尺闊是30尺.
所以30尺就是致聖所的尺碼.
你能夠想像最後去到西邊的時候.
那間房30尺乘30尺乘30尺就是致聖所.
昔日的會幕它的長度是45尺.
闊是15尺.
這個就是昔日的會幕.
致聖所的牆壁的地和天花板都是包著金.
連釘子都是鍍金.
你說為什麼不用精金來做呢?.
金是很軟的.
嘗試一下用金的釘子插進一些硬的木頭會怎樣.
那個釘子也彎了.
所以鍍金就是比較好一點.
在致聖所裡面是有兩個基路柏的天使.
有很多人說我知道.
在藥櫃上面不是有兩個基路柏嗎?.
不是的.
原來在致聖所裡面.
是有兩個巨大的基路柏的天使的像.
高是15尺.
它是打開了它的翼.
由一個翼尖去到另外一個翼尖是15尺.
我剛才說聖所有多闊.
30尺.
兩個基路柏的天使每個佔15尺是怎樣的呢?.
在一邊在房間的中間.
大家的翼尖對著大家.
但是去到另外一邊.
一個基路柏的天使的翼就佔著南邊的牆.
一個基路柏的天使就佔著北邊的牆.
而在這兩個基路柏天使的中間.
就有一個藥櫃.

$^{481}$藥櫃非常重要.
是神的寶座.
一共是45吋長.
27吋高.
27吋闊.
藥櫃是代表神和以色列人納藥.
在藥櫃裡面最重要的是什麼?.
是有法板.
是有十戒.
所以藥櫃也叫做法櫃或者藥櫃.
藥櫃也代表神和以色列人同在.
所以記不記得在《述撒母耳記》裡面搞出了一件事.
二利兩個兒子拿著藥櫃去打仗.
以為帶著神去就有神同在.
結果藥櫃也被斐理士人擄掠了.
在藥櫃的上方也有一對基路柏的天使.
是橄欖木造.
也是鋪了金.
兩個基路柏天使的中間就是神的座位.
為什麼要用基路柏天使來放在藥櫃上?.
基路柏的天使其實是守護者.
我們稍後會看到.
在聖經裡面也有很多經文.
例如歷代至上的28章.
詩篇的90篇和132篇.
說藥櫃是神的腳墩.
不是神的寶座 是神的腳墩.
所以你不要以為神被這個至聖所.
被這個聖殿限制著.
神是無處不在.
所以詩人和聖經的作者都會被我們看到.
這個只不過是神的腳墩.
神是無處不在.
一年一度大祭司都會進入至聖所.
就是在藥櫃的上方.
這個詩恩寶座的上方.
把祭物的血塗在上方.
為什麼要這樣做?.
你想想.
神在詩恩寶座的上方.

$^{521}$望向藥櫃看到什麼?.
看到那兩塊白板.
白板做什麼?.
指出人的罪.
但是當大祭司把血塗在詩恩寶座的上方.
神還會看到這兩塊白板反映出人的罪嗎?.
看不到.
其實這個就是一個很好的圖畫.
讓我們看到耶穌基督為我們所做.
因為耶穌基督為我們釘十字架.
祂的寶血遮蓋我們.
以致神看到我們.
就看不到我們的罪.
我們能夠與神和好.
這是一個很美麗的圖畫.
還有一件事.
就是在網絡上有很多關於藥櫃的圖畫.
其實都是錯誤的.
如果大家留意的話.
大部分這些圖畫.
都是把那兩條金棍穿過闊玉那兩邊.
其實應該是穿過窄玉那兩邊.
為什麼呢?.
當你進入聖殿.
大祭司打開蠻子.
進入致聖所的蠻子.
他就會看到藥櫃.
是闊那邊向著自己.
耶和華應該是坐在這個上方.
我們怎麼知道是這樣的擺法呢?.
聖經有說.
在「列王記」上的八章八節.
和歷代志下的五章九節.
說原來那兩條棍的尾.
是頂著那個蠻子的.
所以你在外面會看到那個棍的尾.
是頂著那個蠻子.
所以我們絕對知道.
是這樣的擺法.
所以那個金棍是穿過窄那邊.

$^{561}$而不是穿過闊那邊.
如果是穿過闊那邊.
你能夠想像.
耶和華不是對著聖殿.
他是對著南邊的牆.
或者是對著北邊的牆.
這樣就出現問題.
無論內殿和外殿的門.
都是用橄欖木來做.
都是鋪了金.
上方都是有基魯伯的天使.
有棕樹.
有開花的花朵.
是很美麗的.
我沒有辦法可以找到.
昔日聖殿的門是怎樣.
但是原來是找到有一扇門.
是古代以色列人的門.
上方也是有棕樹的花紋.
是有花瓣的圖案.
當然就是沒有基魯伯的天使.
我相信也沒有人夠膽用.
這些象徵來給自己的門.
以西結先知在《以西結書》的41章.
他是描述將來在耶路撒冷的聖殿.
在千禧年國的聖殿.
無論牆,地和門.
都是有基魯伯的天使.
那你說為什麼所羅門這麼喜歡基魯伯的天使?.
記不記得在《創世記》的三章.
當亞當和夏娃犯罪之後.
被趕出去伊甸園.
神差派誰看守著新領事?.
就是基魯伯的天使.
所以我相信所羅門是希望基魯伯的天使.
同樣守護著聖殿.
所以你會看到有這麼多基魯伯天使的圖案.
七年之後.
所羅門就完成了聖殿的建築.
完成這個字.

$^{601}$和《創世記》的二章一節.
講神六天完成了祂的創造.
第七天休息是同一個字.
而在這裡作者用同一個字.
也讓我們看到極有可能是有一個意思.
就是說正如神創造了一切都是美好的.
《創世記》的一章三十一節.
現在所羅門建了聖殿.
一切都是美好.
直至到以色列人犯罪.
神的靈離開了這個殿.
這樣就出現了問題.
講完了聖殿的建築.
現在要講聖殿的佈置.
這個部份也是很重要.
因為我們知道聖殿的佈置是很重要.
大祭司是用很多這些陳設來敬拜神.
來事奉神.
當祭司進入聖殿.
還沒有進入這個建築物的時候.
他先要經過外院.
然後要經過祭司的院就是內院.
在那裡他首先會看到的.
就是有兩樣東西.
在南邊也是在他的左邊.
有一個銅鎧.
有一個洗鑿盤.
在右邊有一個銅造的祭壇.
我們先看這個銅鎧.
這個銅鎧是一個圓形的盤.
它的邊那裡是有百合花的圖案.
下面是有野瓜的樣式.
這個銅鎧可以裝一萬七千加侖的水.
所以是非常龐大的一個建築物.
它的直徑是十爪.
它的圓周是三十爪.
它的高度是五爪.
在這個會幕的洗鑿盤.
也有一個銅鎧有一個洗鑿盤.
但是尺碼是比較小很多.

$^{641}$這個銅鎧也在十二隻牛的上方.
十二隻銅牛的上方.
十二隻牛代表什麼?.
十二支牌.
十二隻牛是分開四組.
每組是三隻牛.
分開是向著東南西北.
我剛才說這個銅鎧的直徑是十爪.
圓周是三十爪.
聖經又被人罵了.
在我們當中如果有讀科學的就會知道.
一個圓形的物體.
如果直徑是十爪的話.
圓周應該是多少?.
乘以π你就知道了.
π是3.1416.
10乘以3.1416就應該是31.416.
聖經怎麼會說三十爪?.
分明聖經就是有誤.
所以你的神都不知道怎樣.
連計算都不會計算.
有一個拉比解釋這件事.
解釋得很有趣.
他怎麼說的?.
他說原來以色列人量度圓周.
是找一條量乘來量的.
這條量乘叫做QAV.
是由兩個希伯來文的字母組成.
是QAV和VAL兩個字母.
我們也知道.
希伯來文每個字母都有一個價值.
QAV的價值就是100.
VAL的價值就是6.
你加上這條量乘應該是106.
但是這個拉比說.
你有沒有注意?.
在列王記上七章二十三節.
這個量乘的串法出現了問題.
不是QAV 是QAV.
多了一個H的尾音.

$^{681}$是由三個希伯來文字母組成.
除了QAV和VAL還有H.
H是5的價值是5.
你將這些價值加在一起.
神用一條量乘這條價值是111的總值.
你試一下將111除106.
這個比例就是31.416.
除30的比例是一模一樣.
所以這個拉比怎麼說?.
神用了一條很特別的量乘.
量度到的圓周是30爪.
但是其實應該是31.416.
這個拉比的解釋精不精彩?.
我就不知道是不是真的這樣.
但是我相信另有解釋.
正解就是因為這個銅海是有厚度.
那個黃銅的厚度是整隻手掌那麼厚.
所以如果你量了直徑是10.
如果你計算了厚度.
你量圓周的時候.
是將那條量乘放在銅海的厚的上面.
是中間的部分.
所以不是最外面那裡.
所以量出來就是30爪.
不是31.416.
這個洗作盤很大.
你說怎麼拿水出來呢?.
這個畫家他就是有想像力.
在牛的口那裡噴出來.
可能當時是需要有些器皿.
是舀些水出來.
也有可能的確有個水龍頭.
但是是不是真的從牛的口那裡出來呢?.
我就真的不知道.
祭司他們在工作的時候.
都要保持守潔心清.
他們都要洗手洗腳.
否則隨時會招來殺身之禍.
而在聖經裡面告訴我們.
水是代表聖靈.

$^{721}$而在聖經裡面也不斷提醒我們.
要用神的話語來潔淨我們.
洗淨我們自己.
我們事奉的時候我們都要很小心.
因為我們都是不潔的人.
我們時時都要用神的話語潔淨自己.
否則後果可能是不幸設想.
而耶穌基督也在約翰福音第十三章.
給了我們很好的例子.
記不記得耶穌基督替門徒來洗腳.
是一個很好的榜樣.
彼得說這麼好啊.
彼得說你不如幫我全身都洗了.
耶穌基督說不需要.
你已經潔淨了.
你只需要偶然洗一下你的腳.
這就好像我們待會要領受聖餐.
在領受聖餐之前我們要做什麼?.
凝聚悔改.
這就等於我們不斷地在這裡.
潔淨我們的手,潔淨我們的腳.
這就是聖餐的一個很重要的意思.
還有我剛才說你進到祭司的內院.
左邊你會見到銅鎧.
右邊你會見到一個銅的祭壇.
但是很有趣.
在列王記上居然沒有說到這個銅祭壇的尺碼.
沒有說到銅祭壇的製造.
但是就有提到是有銅祭壇.
那你說為什麼這個作者不去提它?.
學者是有兩個不同的解釋.
第一個解釋就是我剛才說.
推羅王曾經差派一個考將.
叫做烏蘭來幫助以色列人製造這些器皿.
在聖殿的銅柱,銅鎧.
和等下我們會見到在聖殿裡面有很多的器皿.
都是這個烏蘭設計.
唯有這個銅祭壇就不是他設計.
所以有學者說這個就是原因.
列王記上的作者沒有提到它.

$^{761}$另外又有一個原因.
記不記得大衛曾經是數點人民.
不為神所喜悅.
後來他怎樣?.
買了阿羅那的一塊地.
然後在那塊地上建了一個祭壇.
來獻祭給神.
有可能這些學者說.
作者是希望我們仍然記得大衛的祭壇.
而不是說將注意力放在所羅門的祭壇.
這個銅祭壇是30呎乘30呎.
高是15呎.
嘩!有沒有想像一下你是大祭司.
你要獻一隻牛給神.
嘩!你托著那隻牛.
怎樣舉15呎?.
沒有什麼可能的.
所以應該有一條樓梯可以上去.
而且也有很多學者認為.
根本這個銅祭壇在石座上.
你能不能想像這個銅祭壇燒得很大聲.
嘩!一個人一時之間不經意.
好像我這樣彎彎抖抖.
很多時候一撞過去.
真的被他燒得體無完膚.
所以有很多人說是有個石的座.
昔日在這回望的銅祭壇就小很多.
只有4呎半.
那個祭壇我相信.
大祭司真的可以托一隻牛放進去都可以.
第二樣我們要說的.
第一樣就是銅海.
第二樣我們要說的.
就是十個活動的洗盤.
你能不能夠想像.
剛才那個銅海這麼大.
你想拿些水出來洗東西.
嘩!不是簡單的.
尤其是你洗的東西不是這麼近的時候.
很不方便的.

$^{801}$原來有十隻這些活動的洗盤.
五隻在祭司院的南邊.
五隻在北邊.
下面還有洞的.
所以你裝了這些水.
你可以推去不同的地方.
來洗祭物或者裝祭物的器皿.
這些水應該由銅海那裡來.
如果洗完之後骯髒.
就將那些骯髒水運到別處來清理.
其實如果你留意的話.
這些洗作盤很精美.
上面有基魯柏.
有獅子.
有棕樹.
又有花瓣.
讓我們看到什麼呢?.
神不僅著重聖潔.
每樣東西都要洗得乾淨.
還有神也很著重微觀.
第三樣的神設.
就是在聖殿裡面有很多很多的器皿.
這些器皿都是顧蘭他所設計.
包括什麼呢?.
剪燈芯的剪.
灑水的器具.
碗,盤.
還有湯匙.
還有煮祭物的鑊.
甚至清除灰燼的鏟.
這些全部都要做的.
還有在聖殿裡面有很多金和青銅做的家具.
這些都是很實用的器皿.
如果祭司沒有了這些器皿.
真的很難能夠順利地完成他們的工作.
在《列王記》上的七章41至45節.
總結了聖殿裡面所有青銅的器皿.
這些青銅器皿都在約旦河谷被鑄造.
而所羅門經文說他沒有量度這些器皿有多少.
為什麼不去量度?.

$^{841}$因為實在太多了.
不是那麼容易量度得到.
第四個層次就是金香壇.
金香壇是用棗甲木來製造.
是非常名貴的木.
還鋪了金.
在經文裡面沒有告訴我們.
這個金香壇的尺碼有多少.
只是告訴我們它在萬寺的前面.
我們剛才說萬寺是隔開致聖所和聖所.
如果在萬寺的前面的話.
就不是在致聖所裡面.
所以每一天當祭司去點燈的時候.
他也會在這個金香壇的上方燒香.
在聖經裡面.
燒香是代表什麼?.
是代表我們聖徒的禱告.
到了啟示錄你就會看到有提到.
摩西他曾經得到神的啟示.
他懂得製造香料.
香料也有特別的配方.
用在金香壇上.
這些香料你不可以用在其他的用途.
你也不可以私自去廣祭.
你說多香啊.
我用來當香水.
你不可以這樣.
是犯很重的罪.
在贖罪日.
大祭司會將祭物的血.
從金香壇四個角.
你有沒有看到?.
真的有角的.
從這些角的上方來潔淨金香壇.
為什麼要潔淨?.
我剛才說過.
燒香就是代表我們向神的禱告.
所以我們向神禱告的時候.
我們需要什麼?.
守潔心清.

$^{881}$我們才能夠親近神.
否則我們也不能祈求神.
垂聽我們的禱告.
所以我們禱告的時候.
也需要守潔心清.
最後我要說的.
就是金燈台和陳設餅的檯.
在會幕裡面只有一盞金燈台.
也只有一個陳設餅的檯.
但是在聖殿裡面.
十個金燈台.
十張陳設餅的檯.
五張是靠近北面的牆.
五張是靠近南面的牆.
現在我們的金燈台是怎樣的?.
有很多人就用七支蠟燭來代表.
古時的金燈台不是的.
一支銅管上去.
然後分開七支.
每一支的上方都有一個小碗.
小碗裡面就裝著橄欖油.
而橄欖油裡面就有燈芯.
當你一點著這個燈芯.
你就能夠點著這個金燈台.
而在聖所裡面.
我剛才說是有一些小窗.
但是不夠光.
所以是需要這些金燈台.
使祭司能夠見到光.
知道怎樣來事奉神.
你說陳設餅的檯呢?.
也是金做的.
也有十張這麼多.
陳設餅真正希伯來文的意思就是麵餅.
我不是在說麵包的麵.
不是在說麵粉的麵.
是什麼麵呢?.
一個人的臉.
是代表什麼呢?.
神的臉常與以色列人同在.

$^{921}$所以是代表神與以色列人同在的意思.
而在啟示錄的二十二章四節.
的確是將神的臉和神的名放在一起.
我們看到這個就是一個很親切的意思.
就是神與以色列人同在.
最後去到列王記上的七章第五十一節.
經文說:所羅門作完了耶和華殿的一切功.
就把他父大位分別為聖的金銀和器銘.
都帶來放在耶和華殿的府府裡.
這樣所羅門就完成了他建聖殿.
和佈置聖殿.
我今天結束的時候.
是想和大家看兩點.
第一點就是神的恩典.
神甚至可以將我們的罪.
透過我們的罪變成祝福.
很有趣.
說到大衛犯了兩項最大的罪是什麼呢?.
有沒有人想猜一下?.
姦淫是吧?.
和拔蝕花的姦淫.
還有第二項呢?.
謀殺就是.
但是數點人民都是不被神所喜悅.
搞出很大的事情.
死了很多的以色列人.
但是我們看神就是透過這兩件事.
來祝福以色列人.
為什麼?.
大衛雖然和拔蝕花犯姦淫.
但是最後娶了拔蝕花.
第一個嬰兒是死了.
但是生了誰?.
所羅門.
所羅門日後就是建聖殿.
還有第二項罪呢?.
他數點人民死了很多以色列人.
他最後買了阿羅那的那塊地.
做了祭壇獻祭給神.
為自己,為以色列人贖罪.

$^{961}$將來這塊地變成了什麼?.
變成了聖所的所在.
聖殿就在那裡.
有沒有看到?.
雖然大衛犯了兩項非常大的罪.
但是神就是透過這兩件事.
來祝福以色列人.
你說很棒啊.
讓我多犯一點罪.
多犯一點罪.
神的恩典就更多.
這正是羅馬書的五章二十節所說的.
罪在哪裡顯多.
恩典就更顯多.
但是保羅立刻勸我們.
你千萬不要這樣想.
不是說我犯多一點罪.
讓神祝福我多一點.
但是當我們犯了罪.
如果我們真的肯認罪悔改的話.
神是信神的.
甚至透過我們的罪.
完成祂的旨意.
來祝福我們.
這就是第一點.
第二點就是我今天的題目說.
神不是注重華麗的奠禦.
神是注重什麼?.
是注重我們的內心.
在哥林多前書的三章十六和六章十九.
都說我們的身體是什麼?.
是聖靈的奠.
是神的奠.
有很多學者都會說.
聖殿致聖所就是什麼?.
好像我們裡面的聖靈一樣.
外院呢?.
好像我們的身體.
羅馬書的八章和加拉太書的五章.
保羅也有說.

$^{1001}$我們時時我們的靈和我們的身體都會有鬥爭.
不要讓身體勝過我們的靈.
我們要我們的聖靈勝過我們的身體.
還有一樣東西我們要注意的.
就是當耶穌基督為我們釘十字架.
在至聖所前面的萬子一分為二.
由上至下裂開.
不是由下至下.
下至下就是人為的.
有很多人說萬子開了就代表什麼?.
我們可以進入至聖所了.
我們可以與神親近了.
但是你不要想這個方向.
不是我們進入至聖所.
至聖所的亮光是怎樣的?.
釋放出來.
讓別人能夠看到神的亮光.
聖經告訴我們.
聖靈藉著神的話語來到光照我們.
使我們能夠明白神的話語.
當我們明白神的話語的時候.
我們要走出來.
我時時都說有很多基督徒頭大腳小.
我們要走出來.
走出來別人才能夠看到神的光.
這個很重要.
而且在聖殿裡面.
所有的東西都是鋪滿了金.
金碧輝煌.
非常之美麗.
但是很多基督徒.
我們的生命是怎樣的?.
好像很貧乏一樣.
我們真的很需要活出光輝.
很需要活出神的榮耀.
而不是說單單做一個基督徒.
每個星期回來聽一下道.
我們真的很需要被聖靈光照.
我們明白神的話語.
活出那美麗的潔淨.

$^{1041}$願神祝福他們自己的話語.
讓他們低頭禱告.
列王記上的第六章和第七章.
見到所羅門怎樣來建聖殿.
來佈置聖殿.
我們見到聖殿的金漆輝煌.
非常之的榮耀.
非常之的美麗.
但是願意我們能夠活出那美好的見證.
因為我們知道在至聖座前面的幕子.
已經被拿掉.
拿開已經是開了.
願意我們都能夠將神的榮耀活出來.
願意我們都能夠透過聖靈光照我們.
能夠明白聖經.
以致我們能夠活出神的榮耀.
活出神的光輝.
願意我們都有美麗燦爛的生命.
而不是貧乏的生命.
從大衛的生平我們能夠看到.
你就是那位滿有恩典的神.
雖然大衛犯罪.
但是你仍然能夠透過他的罪.
走出你的旨意.
也是讓我們學習到.
我們每一個都是罪人.
我們仍然都有犯罪的傾向.
但是不要緊.
當我們犯罪的時候.
只要我們認罪悔改.
繼續走在你的旨意裡面.
你甚至可以用我們的罪.
能夠實現你的應許.
能夠實現你的旨意.
願意我們有生之年.
都能夠做那美好的見證.
將榮耀頌讚歸給你.
我們今天的禱告.
同主耶穌基督的名義求.
阿們.

\newpage



\section{尼希米記 1:1-4}
\label{sec:032276}
\textbf{Keep Burning}~~~~顏沛恩牧師
\newline
\newline
連結: \href{http://www.church.com.hk/acms/content.asp?site=cdc\&op=show\&type=product\&code=032276\&layout=sermon}{\texttt{032276}} ~~~~ 語音日期: 2026-06-07
\newline
\newline
\hyperref[sec:032275]{\small{< < < PREV SERMON < < <}}
~
\hyperref[sec:index_chronic]{\small{[返順時目]}}
~
\hyperref[sec:032277]{\small{> > > NEXT SERMON > > >}}
\newline
\newline
尼希米記 1:1-4
\newline
\begin{longtable}{cl}
\hline
\hline
章節 & 經文 (和合本修訂版)\\
\hline
1:1 & \begin{tabularx}{0.7\textwidth}{X} 哈迦利亞的兒子尼希米的言語如下：亞達薛西王二十年基斯流月，我在書珊城堡中。 \end{tabularx} \\ \\ \relax
1:2 & \begin{tabularx}{0.7\textwidth}{X} 那時，我有一個兄弟哈拿尼，同幾個人從猶大來。我問他們那些被擄歸回、剩下殘存的猶太人和耶路撒冷的情況。 \end{tabularx} \\ \\ \relax
1:3 & \begin{tabularx}{0.7\textwidth}{X} 他們對我說：「那些被擄歸回剩下的餘民在猶大省那裡遭大難，受凌辱；耶路撒冷的城牆被拆毀，城門被火焚燒。」 \end{tabularx} \\ \\ \relax
1:4 & \begin{tabularx}{0.7\textwidth}{X} 我聽見這話，就坐下哭泣，悲哀幾日，在天上的神面前禁食祈禱， \end{tabularx} \\ \\
[1ex]
\hline
\hline
\end{longtable}
$^{1}$鼎姐妹早晨.
早晨.
我想給你們看一幅圖.
你們還記得他們嗎?.
這隻粉紅色的兔子.
和黃金的小鴨子.
還記得嗎?.
知道這是什麼嗎?.
這是我兒時.
一個很出名的廣告.
在賣什麼呢?.
我都知道是電芯.
知道是什麼電芯嗎?.
就是...叫什麼名字呢?.
耶!就是金霸王!.
果然是同年代的!.
為什麼這隻兔子這麼出名呢?.
因為這隻兔子.
是永遠跑得比別人快.
又不會突然慣低.
這隻粉紅色的兔子.
所以在他的廣告裡.
有一句名言.
金霸王強勁發揮.
耐電比其他多至七倍.
有沒有回憶呢?.
金電比其他多至七倍.
所以他整個廣告的核心價值.
就是這兩個字.
就是持久.
他整個廣告裡都在說.
你用我這隻電.
你用我這隻電就會持久又電.
記得...不是剛才.
是早三個星期前.
曉怡傳出跟我們分享的訊息.
祭壇的火是不熄滅的.
因為祭壇上的火.
要在其上.
這八個字很重要.

$^{41}$要常常燒著.
然後做什麼呢?.
不可熄滅喔!.
跟你旁邊的人說一次好嗎?.
要常常燒著.
不可熄滅.
所以祭司很勤力的.
他們輪更祭.
晝夜不止息地做什麼呢?.
在聖殿裡守著.
要加柴.
確保神的殿宇裡要有火.
弟兄姊妹.
剛才說祭壇的火.
不熄滅.
是因為有個祭司.
日以繼夜在殿裡.
來加柴.
但是我們.
心裡的火.
都要不熄滅.
所以今天的講題叫什麼.
還記得嗎?.
跟你旁邊的人說一次.
我們要.
keep burning.
我們的心裡要keep burning.
不是靠我們的意志.
那些叫死撐.
那些叫什麼火.
那些叫虛火.
當一些大能來到的時候.
大困難的時候.
生命承載不了的時候.
一刺就會怎樣呢?.
會倒塌的.
所以我還記得第一篇.
講尼希米基.
生命無常.
還記得後面那幾個字嗎?.

$^{81}$嘩!你這麼記得.
會怎樣呢?.
收檔啊.
跟你旁邊的人說一次.
生命無常是怎樣呢?.
是啊.
所以這個是虛火.
所以聖靈而來的火.
那裡是什麼火呢?.
那裡是實火.
所以你看看你旁邊的人.
到底他是虛火還是實火呢?.
看一看他的樣子.
就已經知道了.
弟兄姊妹.
今天我們要學一件事.
我們要將我們.
生命裡面那顆電.
我們經常漏電.
經常跪下.
經常走下就沒電了.
我們那顆電.
要換上剛才那隻兔子.
給你看看剛才那隻兔子.
嘩!這是什麼金霸王.
我多加兩個字.
屬靈上的金霸王電芯.
要放在我們的心裡面.
剛才說過.
廣告的核心那兩個字是什麼?.
持久嘛.
所以我們都一樣.
我們要keep burning.
在舊約裡面就要彈起一個地方.
所以當時的以色列人.
他們要坐11號巴士.
知不知道11號巴士是什麼呢?.
就是你的腳.
就要去到聖殿那裡.
但是.

$^{121}$我今天回教會.
其實你怎麼回來的?.
大家都是.
我也是.
人人都是坐11號巴士.
你知道什麼叫11號巴士.
11號巴士來到教會裡面.
我們一起來去敬拜主.
但是你知不知道.
我們這個生命.
其實都是一個教會.
我們這個生命.
其實都是一個祭壇.
因為.
哥林多前書講到.
豈不知你們的.
身子.
就是什麼?.
就是.
聖靈的電.
我們的身子就是什麼?.
聖靈的電.
為什麼好像不會讀.
豈不知你們的身子.
就是什麼?.
聖靈的電.
所以這個body是什麼?.
是神的居所.
是神住在其中.
我們這個祭壇.
怎樣可以持久有電?.
好像這個屬靈的金霸王.
是不會跪下,不會漏電.
一直要持久.
是持久keep burning.
去到見主面那一刻.
所以今天我們繼續用.
尼希米記.
我們在尼希米的生命裡面.
去看看這個人.

$^{161}$他怎樣來.
keep住那裡屬靈的火.
他怎樣來keep burning.
繼續持久.
去到重建.
這個城牆.
他內裡面.
生命這個祭壇.
一定有一些.
持久力的燃料.
裡面會放一些.
持久力的燃料.
有燃料才可以.
才可以燒.
有火.
為什麼有火?.
因為有柴.
濕柴會怎樣?.
濕柴.
濕柴是燒不著的.
乾柴會怎樣?.
乾柴就會.
烈火.
又看看你旁邊那個.
究竟是濕柴還是乾柴.
檢視一下自己.
究竟是濕柴還是乾柴.
稍微溫習一下.
剛才已經讀了.
經文裡面有一個小小的背景.
一章三節這樣說.
他們對我說.
那些被擄歸回.
剩下的漁民.
他們第三批的以色列民.
回家了回家了.
他們可以回到耶路撒冷這個地方.
居住.
在猶大省那裡.
租大欄.

$^{201}$回去的時候租大欄.
受靈辱.
耶路撒冷的城牆.
做什麼事?.
這個很重要.
耶路撒冷的城牆被拆毀.
然後那些門.
城牆圍著有幾道門.
十幾道門.
全部被火.
被火燒.
當時尼希米看到這個很傷心.
他覺得簡直是以色列人.
一大恥辱.
也是神住的地方.
這個聖城會是一個荒涼.
聖殿在其中.
好像什麼都沒有.
城牆被拆.
城門被火燒.
國破家亡的裡面.
上帝的榮耀在哪裡.
反而看見是一個很大的.
羞辱神聖的羞恥.
所以你想一座城.
如果沒有了城牆.
就等於我們的身體.
沒有穿衣服.
沒有皮膚的話.
我們會怎樣.
沒有了防禦力.
沒有了保護.
所以那群百姓.
就算讓他們回到家裡.
仍然受敵在外面.
一攻就倒了.
可以隨時衝進耶路撒冷.
移為平地.
擄走他們的東西.
然後殺死他們都可以.

$^{241}$所以其實他們正在經歷很惶恐.
威脅 恐懼.
焦慮當中.
有家了 但其實不能夠.
在當中安居樂業.
所以尼希米當時帶著一份.
同胞一份血淋淋.
被受辱的痛楚.
所以尼希米說.
二章十七節最後那句.
奈巴讓我們做什麼.
重建耶路撒冷的城牆.
免得再受凌辱.
重建城牆城門.
即是重建神的祭壇.
而這個祭壇.
就是我們的生命.
而這個生命要.
繼續燃燒.
繼續燃燒不止息的靈火.
怎樣才有不止息的靈火.
一會兒講到裡面.
我會解釋三個希伯來文的字.
這三個希伯來文的字.
就是點燃繼續燃燒的原動力.
其實就是一節的經文.
就是一章四字.
準備一起讀.
我聽見者說.
第一是什麼.
第二是悲哀的日.
在天上的神面前.
第三.
你已經講了這三樣東西.
一句說話裡面.
已經有三種.
繼續燃燒的元素.
我逐個逐個講.
第一個是坐下.
你現在不是坐著.

$^{281}$只有我站著.
現在每一個人都坐著.
但原來這裡不是在講.
一種姿勢的描述.
原來原文裡面在字典.
有一個很深層的意思.
一個人由sit down.
去到另外一個.
生命的裡面.
settle down.
settle什麼去處理什麼.
在他的字典裡這樣講.
the meaning of mental.
還有一個emotional settlement.
就是說我全人的精神.
我心理的裡面.
所有的地方.
思想的裡面.
整個人要安靜下來.
不只是坐下來這麼簡單.
原來是整個人.
安靜.
當時尼希米哭得極度傷心.
很難過.
很悲哀.
他清楚這一刻.
我要安靜.
不要再看著那班同胞.
不要看著耶路撒冷發生什麼事.
我現在要坐下.
我要對準神.
我內心要對準神.
我要settle我內裡面的情緒.
我要將我整個人抽離.
處理他這一刻.
很難表達出來.
難以控制的情感.
有一種叫做.
有口難言的痛楚.
弟兄姊妹.

$^{321}$過去我們都有坐下.
每一天都有些時間會坐下.
不是站起來走.
就是坐下.
不是坐就是睡.
我們都有些時間會坐下.
但是我們又有.
我們又有多少時間.
會settle down.
會觸摸一下.
內心.
自己內在的需要.
釋放一下內裡面那種不斷壓抑的情緒.
弟兄姊妹.
你想像一下.
現在我和你拿著一個玻璃瓶.
玻璃瓶裡面.
放了一些很骯髒的東西.
一扔那些是泥水.
於是我很想設法.
整乾淨玻璃瓶.
然後我用手指.
不斷搖搖搖搖.
然後不斷搖搖搖搖.
那你的玻璃瓶會怎樣.
會怎樣.
乾不乾淨.
很骯髒.
比剛才還要骯髒.
因為你不斷搖搖搖搖.
因為你很心急很想快點.
弄乾淨它就好了.
不如試一下這樣.
將玻璃瓶.
放在這裡.
坐下.
不要動.
這個玻璃瓶坐下.
不要動.
你猜一會兒會發生什麼事.

$^{361}$一分鐘.
三分鐘.
五分鐘.
再看一下玻璃瓶.
你會看到什麼.
那些泥.
去了哪裡.
那些泥那些骯髒的東西.
就會.
沉澱了.
而上面的.
那些水.
就很清澈.
就乾淨了.
其實我們的心都好像.
這個玻璃瓶一樣.
可能一聽到.
負面壞消息.
大能題怎麼辦.
突然會怎樣.
心跳加快.
血壓高.
驚惶失措.
用自己的方法.
打電話.
找人救.
現在不用找人救.
拿電話來問誰.
問誰.
不是問我.
相信大家都很厲害.
一拿電話來.
隨時隨地.
問誰.
問AI.
你有沒有試過.
很多自己的方法 總之呢 要去解決 找人來幫 找AI幫.
其實說的是什麼 我們的心不斷在搖晃這個玻璃瓶.
剛才說一搖搖搖 搖搖搖的時候 那些污泥在瓶子裡 嘩 骯髒骯髒.
整個瓶子都骯髒骯髒 其實我們的生命都是一樣的.

$^{401}$你起碼第一個功課 第一個燃料 告訴我們知道 我們要坐下.
sit down 將我們的焦點放回神的裡面 不要動 當你安靜的時候.
原來settle down我內心的恐懼 裡面有一些莫名其妙的焦慮.
和一些東西 嘩 突然間浮出來 但因為你安靜 坐下.
去到settle down的時候 這個玻璃瓶 所有骯髒骯髒的東西.
就將你沉澱下來 回應上星期韋傑傳道的分享.
可能我們有些弟兄姊妹 已經去到一個burn out的狀態.
又或者可能是一個很疲累的時候 可能去到一個臨界點.
我就快爆了 弟兄姊妹 我們其實需要預備自己sit down.
我們需要預備自己坐下 為了什麼?.
為了讓我們的心對準神 可以settle down我們內在的情況.
這是第一個燃料 第二個燃料是什麼呢? 記得剛才讀過的兩個字嗎?.
第二個燃料是什麼? 哎呀 忘記了 你起碼坐下 接著做什麼?.
哭泣嘛 很容易解釋 什麼叫哭泣 弟兄姊妹.
就是哭泣 流淚 但原來在希伯來文的三個字母裡面.
原來他們是說水的流動 水不停不停湧出來.
水的流動 即是說現在很悲哀 很難過 很哭泣的時候.
內裡有一種沒辦法 壓抑的情感 我要做什麼?.
我要爆了 我的眼淚突然間好像決堤一樣.
眼淚好像全水 聽到噴出來那種水的流動.
湧出來 是一種身體和心靈已經沒辦法 可以再忍受了.
很強烈地要發洩出來 尼希米當時那種哭泣是放聲.
是豪哭那種 尼希米當時那種悲傷 它和神對齊的.
所以尼希米的流淚 當時是神的流淚 是神的難過.
是神的憂傷 尼希米當時的眼淚就好像一份禮物.
而這份禮物都是一個奉獻 是一個無言的祈禱.
我就是不斷哭 不斷哭 不斷哭.
原來這個已經說不了 說不了任何東西.
只是可以無言地去禱告 當一個人已經沒有說話.
可以表達的時候 我只能夠哭很傷心.
原來眼淚就成為人裡面最真實 也都最誠實的祈禱了.
而這個祈禱內容 上帝的憐憫.
上帝是聽到人眼淚背後最深的聲音.
所以大衛說 詩篇56篇第8節 我們一起讀.
準備起 我幾次流離 你都數算.
求你把我的眼淚 裝在你的皮袋裡.
這一切不都記在你的冊子上嗎?.
這個瓶袋也可以譯作一個瓶子.
我們的主一滴眼淚都會珍惜.
並且他用莊 莊就是收集.

$^{441}$放在哪裡?放在他的瓶子裡 放在他的皮袋裡.
當時大衛跟尼熙米一樣 爆哭 非常辛苦.
非常痛苦的情況下 流離的狀況.
為什麼會這樣?流下淚水.
在上帝眼裡看見 聽見 人的痛苦 那種委屈.
眼淚從來都不會白流的.
因為上帝他完完全全看見.
並且眼淚背後的意思.
他說都會記在哪裡?.
記在上帝那本寶寶那裡.
所以弟兄姊妹 你有沒有經歷過這種無言.
只是眼淚的祈禱?.
可能我們曾經經歷過失去親人.
又或者我們曾經經歷過.
跟我們一起成長的小寵物離開我們.
可能又會經歷過一些感情關係的破口.
可能又會經歷過一些極度擔心狀況的情況下.
我也不知道可以怎樣 總之很擔心很害怕.
又可能曾經試過被人出賣背叛的時候.
我也不知道可以說什麼.
只是那份痛在心裡面.
只能夠發出一份無言的禱告.
對主獻上最真的眼淚.
眼淚帶出說不出豐富的表達.
尼希米就是這樣.
哭乾了 哭盡了.
一種嬌心的傷痛.
所以第二個keep burning的燃料.
其實是什麼?.
尼希米的眼淚背後是在說.
傳言將他的生命傾倒在上帝的面前.
他已經沒有話說.
已經說不出.
我只有眼淚去呈現給你.
盡 一個人的盡.
是在說尼希米眼淚的背後.
我將我的生命傳言獻給全kingdom.
講個故事給大家聽.
有一位陶瓷大師.
你知道陶瓷 當然要做很漂亮的陶瓷.

$^{481}$有一天他做了一個很漂亮的陶瓷.
誒 徒弟 你快點過來.
你知道做完陶瓷要怎樣的.
要拿去做什麼.
要拿去燒 令到他可以見顧.
誒 你拿去燒吧.
點子就出事了.
然後出爐的時候可能控制不了那些溫度.
一拿出來看見陶瓷做什麼.
龜裂.
裂了幾條痕.
於是徒弟看到之後嚇死了.
哇 好害怕 怎麼辦.
然後他馬上就用自己的方法做什麼.
拿膠水 拿金粉 接著做什麼.
摸著修補那些破口 那些裂痕.
以為這樣可以當是一個藝術.
演繹當中的瑕疵.
接著這個陶瓷大師走過來.
突然之間一巴做什麼師.
做什麼.
一巴做什麼.
將整個陶瓷打爛.
接著陶瓷就掉在地上.
整個就碎了.
那徒弟怎麼樣.
嚇得呆住了.
呆住了 很害怕.
接著大師跟他說.
你每一天在修補這個缺口.
你每一天都在執著這個半製成品.
其實你看清楚.
它永遠都只是一個有瑕疵的作品.
唯有徹底做什麼.
打爛它.
變成粉了 還有沒有用 沒有用.
接著怎麼樣.
那我就加水 做些泥再搓一下.
再重新去建築.
重新去雕塑.

$^{521}$另外一個.
嘩 這個帥哥.
完美的瓷器.
弟兄姊妹 我們很多時候在神面前.
都會的.
可能祈禱的時候.
都會流幾滴眼淚.
戴著頭盔 可能又會用自己的方法.
去解決一些問題.
又或者再深層一點.
其實我們常常都戴著一些假的面具.
裝出來很好.
但其實裡面是掩飾著我裡面的洞.
我生命的洞.
我生命的缺口.
我生命破裂的地方.
尼希米剛才說什麼.
她哭的背後是她徹底的將自己粉碎.
去到神的面前.
傾倒給主.
「主啊!你重新再建造我這個瓦器」.
「我這個瓷器」.
在神面前再一次打爛它.
當人願意放棄.
當人願意放下那些執著.
特別是那些你沒辦法搞定的時候.
你經歷生命的拆毀,破碎.
那一刻你很清楚知道.
我生命只有主.
我生命只有主自己.
一種完全徹底的降伏.
我給兩個字.
就是我們的生命要放手.
跟你隔壁的人鼓勵他.
你要什麼?.
放手.
不是放棄.
因為神在這裡.
放手的時候.
你知不知道這個瓦器會怎樣?.

$^{561}$當人的主權.
我願意放下,放手.
神的恩典才可以介入在我們的生命裡.
不知道你有沒有看過一本書.
是我初信的時候看的.
「放手,讓神來做」.
放手,接著就怎樣?.
「神來做,讓神來做」.
為什麼?.
因為這八個字很重要.
一起讀吧,這八個字.
「準備起人的盡頭」.
有兩個頭都很重要.
人的盡頭.
徹底赤裸地粉碎.
已經破爛了這個瓦器.
「主啊,再一次重建我」.
「雕琢我生命」.
「神的起頭」.
好了,第三個燃料.
第三個燃料.
先問問你,記不記得?.
坐下哭泣,接著怎樣?.
是呀!禁食祈禱.
問問第一姐妹.
有沒有第一姐妹經驗過.
禁食祈禱的揮揮手來看看?.
可能一天裡面抽一餐出來禁食.
又或者你會喝一些流質的飲品.
總之取代了平時的食慾量.
其實純粹是幫助去補充.
我身體裡面那一刻的需要.
但是為什麼特意定一個時間出來.
還要禁食祈禱呢?.
平時祈禱祈就行了.
我為什麼要禁食去禱告呢?.
背後一定會有一些原因.
禁食只不過是方式.
當然重點是一個祈禱.
剛才曉怡主席報告的時候告訴我們.

$^{601}$咦,原來我們已經開始了.
很欣賞.
很欣賞第一姐妹.
每一次探訪的時候.
一坐下的時候.
第一姐妹已經40天裡面.
已經為下令會來禱告了.
小組裡面又為下令會來禱告了.
如果你再留意周刊.
其實有一個指引.
第一天可以減,第二天可以減,第三天可以減.
到了第七天,主日的時候.
就為了什麼什麼來祈禱了.
剛才還記不記得主席報告的時候.
這個星期開始.
我們又鼓勵你做什麼呢?.
就是禁食祈禱了.
咦,為什麼會是這樣的呢?.
這個開始又要加入禁食祈禱.
作為一個操練.
可以鼓勵各人.
可能你自己一星期可以找一餐來禱告了.
又或者平時已經有操練.
幾天的裡面.
你預備就這一餐來禁食禱告了.
又或者有經驗過.
一天的裡面.
禁止一餐來禱告了.
有一個很重要的目的.
第一姐妹.
你知不知道.
飽暖就會施什麼的?.
哎呀,我不是說這個.
飽暖就會施.
床褥.
明白為什麼嗎?.
你有沒有試過.
當你吃完飽.
突然有碗飯在你旁邊.
又有意粉.

$^{641}$不知道呢.
在樓下街口.
那個飽什麼呢?.
經常我們勁快隊.
就會去那裡吃飯.
哇,真的.
其實那餐太大了.
對於我來說.
為什麼吃完飽.
咦,為什麼會加上意粉呢?.
其實根本是太飽了.
所以吃得飽.
你就會怎樣?.
就會想睡覺了.
行不行?.
所以你明白.
飽暖就會施床褥.
所以你明白.
為什麼要用一個方式禁食呢?.
讓你專心.
讓你定精在聖經裡面.
讓你定精在禱告裡面.
同神來到對象.
你想想.
吃完飽的時候.
我怎麼祈禱呢?.
你還坐在那裡.
然後.
期期下.
你會怎樣?.
然後我聽到腳的聲音.
然後就會睡著了.
所以當你禁食的時候.
只是飲流質.
反而讓你精神集中.
去更加幫助你.
專心來尋求主.
你對準神.
預備自己進入下靈會裡面.
可能安靜的時候.

$^{681}$神去掏空你生命裡面.
一些需要.
或者隱而未現的東西.
就在禁食禱告的時候.
你要去處理.
你要去靠著去潔淨你自己.
你要去呼求神的幫助.
又說在過去的日子裡面.
很忙碌的時候.
沒有停過.
就在禁食禱告的時候.
預備你自己安靜.
安靜就帶來經歷.
這是什麼經歷?.
是能力.
神平靜安穩.
住在你裡面.
當你禁食禱告的時候.
專心聆聽.
你去尋求生命的照明.
和方向.
預備好自己進入下靈會.
我很相信.
每一個人.
每一個弟兄姊妹.
你越大學武.
神就會在下靈會裡面.
為我們預備不同的禮物.
而這份禮物需要你接.
需要你聽到.
需要你去拿.
所以為什麼我們鼓勵你.
要去預備自己呢?.
倒空你心裡面的東西.
當時看到同胞受難.
看到神的家.
神的地方荒涼.
讓我吃多少餐都好.
更重要我要去祈禱.
尼希米將他整個人傾倒奉獻.

$^{721}$在神的面前考考你.
你知不知道尼希米禁食禱告多少個月?.
哎呀!其實我講過的.
去到彌散月.
即是去到四個月.
一期禁食祈禱.
去到四個月的時間.
其實在這個禁食祈禱的時間.
尼希米特意為了重建這件事.
來自卑懇求.
你知不知道在這段時間裡面.
神已經預備了尼希米的心.
神亦都同時在波斯王的心裡面去動功.
就是剛才我們讀的經文.
當王見到尼希米.
你為何面帶愁容.
心裡愁煩.
你沒事嘛.
跟著之後就問他.
你想求什麼?.
於是尼希米做了一個動作.
我向天上的神祈禱.
其實有另外一個譯本譯做默禱.
即是隨時隨地都可以禱告.
你不是很清楚知道神的心.
他馬上.
當你看聖經的時候.
他馬上已經懂得回應.
有三個要求.
就是王啊.
你差點讓我去到猶大.
第二我就要召書.
因為我要去到哪裡的時候.
我有一書.
我就可以經過.
就可以順順利利走過.
第三件事當然要木材.
重建城牆需要一些木.
經文說.
出於神的.

$^{761}$王就允准了.
因為我神私人的手幫助我.
你看不看到呢.
由禁食祈禱四個月.
跟著去到儀仗的默禱.
就知道神的心.
尼希米的生命.
一直都是流露著一個祈禱的生命.
一直都在祈禱.
由坐下.
去到settle down.
我裡面的情感.
這是個人一個很大的破碎.
去哭盡了.
哭乾了.
處理好自己了.
然後怎樣.
禁食祈禱.
去求主.
去仰望主.
很奇妙地.
在這短短四個月時間.
神就藉著那一刻的安靜.
將屬靈的膽量.
屬靈的智慧.
屬靈的能力.
擺在了尼希米的心裡面.
以至於王問他.
你要什麼.
然後默禱馬上可以回答出來.
並且記得剛才讀的經文是.
他沒有懼怕.
去到王面前.
很坦然很坦白地告訴王.
我的需要.
並且他信心說.
我神私人的手幫助我.
就是因為這樣.
他就走上了國度的服侍.
他就重現了.

$^{801}$所以尼希米一直點燃的燃料.
就是我們今天要讀的金句.
貼上來的詞書.
五章十七節.
我只要中間那五個字.
不住的什麼.
是呀.
不住的禱告.
跟你旁邊的人說一次.
我們學.
不住的禱告.
弟兄姊妹.
祭壇的火要keep burning.
說的是我們的生命.
這個祭壇要有一個實火.
繼續的燒旺.
其實還原基本功.
一切都是由祈禱開始.
上上臂聖靈充滿的人.
他裡面那個是聖靈的火.
聖靈的火是實火.
一直keep burning.
因為他每一個禱告.
都發自他內心的真實.
他裡面的悼空.
他的誠實.
他降伏在神的面前.
所以聖靈的火.
也是上個禮拜.
韋傑傳道跟我們分享.
上帝的愛.
上帝的愛同樣.
充滿在這個真實.
悼空破碎的生命.
以致我們才可以走出.
阿傑跟我們分享的.
Phileo Sorge的愛.
一個真實的祈禱人.
不是有事才來求.
就是保羅說的那句話.

$^{841}$不住的祈禱.
你可能問我.
怎樣不住.
我什麼都說出來了.
我知道弟兄姊妹.
一早起床.
可能上班也好.
可能做家事也好.
一早起床就坐下.
馬上就會祈禱.
祈祈祈祈.
祈完了.
沒話說了.
那我還要說什麼.
祈什麼.
為什麼保羅會說.
不住的祈禱.
弟兄姊妹.
如果你說到.
話都沒得說.
就不祈禱了.
你知不知道.
我們不明白.
祈禱本身的意思.
我們很容易將祈禱.
變成一種交功課.
因為牧者叫你回月土會.
剛才已經叫了.
記著這個禮拜.
我提醒你一次.
之後牧者又叫你.
四十天為夏令會祈禱.
牧者剛才又提醒你.
這個禮拜開始又鼓勵你.
要禁食祈禱.
多說一句.
補多一句.
你要按著自己身體的情況.
才去禁食禱告.
因為我知道有些弟兄姊妹.

$^{881}$每一天都要看著自己的血糖.
所以留意自己身體的狀況.
就禁食禱告.
你有沒有想過.
其實我們每一天.
都有很多祈禱清單.
祈禱需要第一第二第三.
其實主全部都知道.
那為什麼還要祈禱.
其實神是不是想我們.
這樣祈禱.
祈禱其實是什麼.
祈禱的焦點.
是在說關係.
而不是告訴主.
第一這樣.
第二這樣.
第三這樣.
第四這樣.
那個叫資訊傳遞.
祈禱的焦點是在說關係.
你想想.
爸爸還是很想子女.
哎呀 回來了.
然後一起做什麼.
一起玩.
一起看電視.
一起聊天.
是不是.
很想我們的孩子.
可以跟父親和媽媽.
有一個親子的交流.
所以禱告.
其實就是一個聊天的交流.
我記得初信的時候.
其實我當時是傻乎乎的.
百思不得其解.
然後聽到神學院有一個老師.
分享了一個見證.
他說有一次他很想去吃披薩.

$^{921}$然後跑去披薩室.
突然間靈裡面.
有一種很大的阻攔.
把他拉出去.
阻攔他.
不准他進去披薩室吃.
然後他也不明白為什麼.
但是他馴服.
生命的操練.
神學院的老師.
他馴服就不吃了.
當他主日講完道之後.
他完全明白.
為什麼神要阻攔他.
當他講了三場崇拜之後.
他已經失聲了.
如果他再失敗的時候.
相信他未必可以講得完.
他更加知道.
原來當刻.
神的聲音.
是要來保護他.
當下這個初信的我.
我是拗了頭的.
原來神會講話的.
為什麼這麼神奇呢.
保羅所講.
不住的祈禱.
不是叫你24*7合手.
然後閉上眼不斷講話.
而是.
剛才講的焦點在哪裡.
關係.
是24*7.
和神.
在裡面.
父子.
父女的關係當中.
意識到神的同在.
原來祈禱是.

$^{961}$在關係裡面.
意識到神的同在.
你不會24*7.
不斷不斷和你父母講話.
你都會有聊天的時候.
但你都會安靜的時候.
但是全部的人.
都在同一屋簷下.
同一道理.
不住禱告.
不是叫你24*7不斷講.
而是當你安靜.
閉上嘴的時候.
我們都在祈禱.
為什麼?.
因為我們意識到.
神原來和我們在一起.
我們擁抱著聖靈.
聖靈也擁抱著我們.
Henry Lauen是.
勞運靈修學大師.
他解釋Paul.
保羅所講的一句話.
就是不住禱告.
他用Breathing.
Praying is like breathing.
來解釋什麼叫祈禱.
祈禱就好像呼吸一樣.
弟兄姊妹你一定要呼吸.
你不呼吸你會怎樣?.
你會怎樣?.
不呼吸你會怎樣?.
會死的.
他就是用這個比喻.
原來我們的禱告就好像呼吸一樣.
當我們不呼吸的時候.
我們這個人會死的.
當我們不祈禱.
我們的屬靈的生命.
就會越來越枯乾.

$^{1001}$所以他說.
我們每一秒.
其實我們都在呼吸.
你不會突然意識到自己在呼吸.
因為已經很自然地.
已經在呼吸.
即是我們和爸爸一樣.
我們很有意識地知道.
爸爸在我裡面.
我們一起已經在祈禱.
基督徒的生命.
無論你工作,吃飯,休息都好.
你整個人都好.
你都是同神連結在一起.
這個就是關係裡面.
24*7不住的祈禱.
我常常會用這個比喻.
去描述我和神之間的關係.
我是一條鐵線.
神也是一條鐵線.
但我們兩個是怎樣的呢?.
不是這樣的.
我們是怎樣的?.
不是黏在一起.
是怎樣的?.
是扭在一起.
扭到一個位置很深的.
是打風都打不掉的.
以致我們彼此之間的生命.
是生命相容.
我的心靈和內在的聖靈.
是契合的.
是相通的.
去到這樣的關係裡面.
我更加貼近神的心意.
祈禱不是牧者叫你做.
你就要去做.
也不是閉上眼參與一個聚會.
而是真實生命的連結.
是一個同神之間.

$^{1041}$一個很深的神同在的關係.
講到最後很想和您分享一個.
最近的一個經歷.
可能前兩星期我突然失蹤了.
不知去了哪裡.
因為我飛回神學院.
美國神學院我畢業了.
當我下飛機的時候.
其實是凌晨.
所以我早已預訂了凌晨.
每個女子都有間酒店會比較好.
我其實是想省錢的.
我本來是想睡在機場的.
然後那些弟兄姊妹說你不要啦你快點找間酒店住.
所以我早已預訂了機場的酒店.
可以一晚.
早已在香港找好了路.
因為你知道我有個特性.
知道是什麼嗎.
就是路痴.
東南西北都搞不定.
所以早已在機場走左走右走上走下.
然後不知道經過什麼.
然後就會去到這個酒店.
去到酒店你會給什麼.
給passport.
check in.
然後你知道職員跟我說什麼嗎.
你知道他跟我說什麼嗎.
不是這間.
然後呆了一刻.
這間叫Hilton.
你預訂的那間多了一個字.
叫做Hilton Garden.
他還教我他沒什麼.
我睜大眼睛看著他.
他說你回到機場.
然後叫車.
叫酒店的shuttle會載你過去.
當下想叫什麼.

$^{1081}$我的卡是打不到電話的.
我的卡是上網的.
不可以顧網的.
總之就是打不到.
剛才說.
黃文尼熙米你有什麼要的時候.
尼熙米做了什麼動作.
還記不記得.
擦肚子.
然後那一刻你找誰.
沒有人可以找的.
知不知道.
然後那一刻很自然地.
不知道為什麼看了一看.
酒店的門口.
你知道凌晨時分.
你知道我已經有腦花了.
然後是什麼.
什麼都看不到.
黑漆漆的.
就看到一些光.
就是馬路.
那一刻我可以call you吧.
我想說這個經歷.
是在說.
其實我很清楚知道.
當刻那個不是我.
因為真正的我.
認識我的人就知道.
我會怎樣.
我站在那裡.
哎呀死了.
那要怎樣.
然後就很害怕.
但我知道那一刻.
為什麼這麼鎮定.
內心我形容我自己.
那一刻平靜安穩.
然後慢慢按電話.
call you吧.

$^{1121}$慢慢等.
那個時候想起.
其實在飛機那裡.
在香港預備去的時候.
我一早已經祈禱了.
就正正就好像尼熙米.
禁食禱告的時候.
神注入了一份.
出人意外的平安.
在我心裡面.
另外還有一件事.
更嚴重的.
然後有一天.
你知道我要行畢業典禮.
來畢業禮的時間.
還有半個小時.
我是住在神學院的apartment.
但是畢業典禮.
就不在神學院裡面.
就要十分鐘車程.
去到他附屬的一間學校.
一間college那裡.
去行畢業禮.
然後我和另一個同學.
已經mark了.
因為大家都不是local人.
然後就.
哎呀你有車.
正啊.
然後他懂得開車.
然後我就.
喂記得叫我啊.
到時候我call你.
你記住接我啊.
因為大家都是附近apartment那裡住的.
然後.
記住啊.
畢業典禮前半個小時.
然後我打電話給他.
不停地打.

$^{1161}$打爆機.
你知道發生什麼事嗎.
沒有人聽.
然後就開始滴汗.
開始害怕.
搞什麼.
因為他不是local人.
我也擔心大家平不平安.
但是我已經不知道.
那些人在哪裡找你了.
我一直call都沒有人.
那一刻很震驚.
真的震驚.
裡面跳得乒鈴bang lang 的時候.
但是我聽到我內裡有把聲音.
跟我說.
你走出去吧.
你明不明白什麼叫走出去.
叫我收拾好東西.
走人了.
離開你住的apartment.
你走出去吧.
你走出去.
我就很聽話.
我就走了出去.
走了出去就看到.
一對couple很靚仔.
然後穿了西裝.
然後太太化了妝.
然後我就很自然地問他.
請問你們是不是去畢業禮.
請問你們是否參與畢業禮.
然後我就說是的.
然後我馬上可以坐順風車.
去到這個畢業的地方.
還有.
原來他是神學院的高層.
他還可以帶我去那地方的大廳.
因為其實地方太大.
我還要去找.

$^{1201}$哪一個位置呢.
他還帶我去.
其實我有再一次經歷.
我曾經跟大家說過的.
「閒步健步主抱弟」.
我有再一次經歷.
我想說的是.
祈禱的生命就是這樣.
爸爸常與我同在.
我是爸爸,我的父親是我親密的朋友.
二十四乘七無間斷在關係裡面.
扣在一起的關係.
爸爸知道我.
我知道爸爸.
我聽得到是爸爸的聲音.
你叫我就行.
雖然我都不是很明白.
有時候.
但是你叫我就行.
最後.
keep burning.
不止息的靈火.
當你不會漏電.
不會灌低的時候.
裡面要有持久力的燃料.
就是要有這三種的燃料.
坐下.
誠實面對真正的自己.
痛哭背後說.
你是全然將你的生命破碎.
傾倒在神的面前.
禁食祈禱.
禁食其實不是一個重點.
最重要的是你徹底倒空.
全然降伏.
但是你知道.
這個操練的最重要的根基.
是在與神之間的關係.
如果不是在這個關係裡面.
這個禱告只會淪為.

$^{1241}$教功課一個宗教的儀式.
你要知道那惡者.
只做一件事.
就是怎樣.
掩飾我們這裡聖靈的火.
所以我們要向祂說.
No way.
叫祂走開.
聖靈不斷在我們生命裡面.
來去加柴.
最後一版.
當你對主.
擁留這個極大的渴慕.
剛才說的.
好像呼吸一樣.
那種與神的連結.
你的生命自然就會被點燃.
教會的復興.
是因為我們的生命願意破碎.
是因為我們的生命願意被打爛.
是因為我們的生命願意傾倒.
擺在神的面前.
一切根基來自阿爸父.
阿爸父.
爸爸.
\newpage



\section{士師記 6:17}
\label{sec:032277}
\textbf{聖經人聖經事:基甸第一次求證據}~~~~黃克勤牧師博士
\newline
\newline
連結: \href{http://www.church.com.hk/acms/content.asp?site=cdc\&op=show\&type=product\&code=032277\&layout=sermon}{\texttt{032277}} ~~~~ 語音日期: 2026-06-12
\newline
\newline
\hyperref[sec:032276]{\small{< < < PREV SERMON < < <}}
~
\hyperref[sec:index_chronic]{\small{[返順時目]}}
~
\hyperref[sec:032278]{\small{> > > NEXT SERMON > > >}}
\newline
\newline
士師記 6:17
\newline
\begin{longtable}{cl}
\hline
\hline
章節 & 經文 (和合本修訂版)\\
\hline
6:17 & \begin{tabularx}{0.7\textwidth}{X} 基甸對他說：「我若在你眼前蒙恩，求你給我一個證據，證明是你在跟我說話。 \end{tabularx} \\ \\
[1ex]
\hline
\hline
\end{longtable}
$^{1}$各位定睇姐妹,我地又一齊睇聖經,今次繼續係生命樹同大家睇機電,睇下呢個聖經人物同埋一件故事,又睇下上帝ge 手點樣做法..
機電第一次求證據呢就係今日ge 主題黎ge ,第一次求證據,求咩證據呢就唔係證明D 乜ye ,其實呢個證據ge 意思呢就係求上帝比一個明證佢,就係你幫我架啦咁你比D ye 我zar 手咁ge 意思呀,證實一件事..
咁好多時呢D 人都覺得機電係一個好大能ge 人啦,好夠膽啦咁樣,能夠帶三百人就去攻擊到成個米電咁多ge 部族啦,不過其實呢睇落去老老實實機電呢都幾驚架,好怕架..
好了,士司機第六章十六節,耶和華對他說,我與你同在,你就必擊打米電人如擊打一人一樣.呢個呢就係第一點要講ge ,咁呢就係神呢就去搵機電,咁呢就睇就見到機電呢個人呢就係呢個就係我要搵ge 人啦,.
就係呢要佢出黎我聽到我ge 百姓ge 哀求比米電人害到好苦,而家我要搵一個人出黎就幫我打擊米電人,其實呢米電人就好似一條皮鞭咁樣鞭打呢個又看耶和華眼中看為惡ge 事ge 以色列人..
咁事呢,但係今次興起返米電人ge 對手機電,咁去打擊返米電人呢又用返以色列做一條皮鞭去鞭打返呢個米電人ge ,咁神就同機電講啦,我與你同在,咁呢個係一個好大ge 應許黎架喎..
如果你睇呢度呢係寫住耶和華對他說,如果睇上幾節呢係耶和華的使者咁樣去搵佢喎,咁黎到呢度呢就變jor 耶和華冇jor 的使者呢幾個字ge ,其實係舊約裡面耶和華的使者好多時都係耶和華本身黎架..
不過有時寫呢就好似神係猜探個使者其實佢自己,因為好有可能呢係因為神本來係冇一個形體架嘛,不過佢黎到呢個出現ge 時候呢係用一個人ge 樣,於是呢就叫做耶和華的使者,其實係耶和華本身,我地ge 神本身黎架..
我與你同在呢件事係一個最大ge 應許黎架喇,如果神真係出現同你講我與你同在,咁我想問下大家如果咁樣講你覺得神要同我同在要同我一齊去攻擊米電人gwo 個應許夠唔夠呀?.
其實好夠架喇,神自己親自出現話我與你同在,仲加埋添你會打擊米電人好似打擊一個人一樣,米電人多到千萬,嘩點打呀?.
係架喎,成個米電人係你面前只不過似一個人架咋,咁你打一萬人就好難打一個人就易好多架嘛,而且我與你同在一定要記得呢一句,係我地ge 人生裡面最緊要記得神與我同在,我時常俾人地ge 祝福都係主與你同在,咁其實神與我同在就乜ye 都得啦,都可以解決到啦,都過到架喇..
神咁同佢講喎,直頭出現同佢講..
使我知道與我說話的就是主,聖經係咁講架,不過呢度有D 字改一改俾大家知道架,使我知道呢幾個字藍色畫jor 呢希伯來文係冇架,如果我係你眼前無因,你俾一個證明我呀咁樣,等我知道你與我說話喇..
我睇今日呢一個經文呢,我係覺得我地ge 主耶和華呢係非常之忍耐,對於機電係好忍耐,因為明明你睇下上一節佢先講我與你同在,然後你睇下一節佢竟然機電話你俾個證據我,等我知道同我說話ge 就是主你自己呀咁樣..
即係奇怪啦,譬如呢個總司令將軍黎同我一個人講,你同我帶兵出去打啦,然後話你俾個證據我知道你就是個將軍,好不敬架喇,將軍可能係千萬人當中我唔識,但係佢見到ge 係神,神同佢講喎,佢都講我若在你眼前無因,俾一個證據我,等我知道你就是主..
咁樣喎,咁我覺得其實咁講呢,都如果我係神呢,都一巴星過去呢,即係有D 不敬架喇,覺得咁樣同我講ye 咁樣,俾個證據我,咁一來信我唔過啦,你俾張鏡頭zar 手呀咁樣呀,證據呢個字呢,哦,就係記號,你俾一個記號我,等我知道與我講話ge 就是主..
la ,呢個說話呢,就daba,daba呢,係呢一個ge 時態裡面呢,係應許咁解架,俾,俾我知道你俾ge 應許係你俾ge ,la ,呢個主字呢,就其實就唔係主黎ge ,他係你呀..
基建話,我若在你眼前無因,求你俾我一個記號,知道同我俾呢個應許我的就是你.la ,咁arm arm 你同我講,就話你同我同在,你俾個記號我,等我知道話與我同在ge 就係你..
咁其實呢,我覺得基建係gwo 度有少少唔係好知自己講緊乜ye ,好似知好似唔知,不過你又覺得你叫我去打擊米電人,得唔得架,真唔真架,係咪可以做到呀,點做法呀,係咪真架,有冇咩zar 手呀..
咁,咁就俾個記號我,求你唔好離開呢度,la ,咁個記號呢,仲得意喎,唔係求神俾架喎,係我做D ye 俾你簽紙咁樣架喎,好得意架喎,你簽個名,我俾張紙你簽個名..
求你不要離開這裡,la ,我要,我想要個記號係咁ge ,就呢,你等等我,求你不要離開這裡,等我歸回將禮物帶來,供在你面前,你俾個記號我知道你真係你呀,咁樣,咁我覺得已經好唔尊敬架喇,你俾個記號我,等我知道你真係你呀..
la ,而家你唔好走,等我返黎,我返黎lor D 禮物黎獻俾你,咁樣喎,咁如果我係神,我就話唔使喇,而家你想要乜ye 證據呀,咁樣,咁就,不過神呢,就好耐性喎,咁主就話喇,我必等你回來,呢個必字呢,又係冇ge ,我等你回來..
咁呢,呢個等呢,係Yahshab,Yahshabge 意思係,Yahshab Ad,係我留,我會留低,Yahshab係留低,Ad係直至,我會等你留低直至你返黎,好,我等你返黎..
咁明明gwo 個係耶和華神,竟然間你叫佢你等我返黎呀,我而家返去lor 禮物,lor 咩禮物呢,係lor 獻祭ge 公物喎,一陣再睇啦,你知道啦,好犀利架..
第十九次,基奠呢,就跑去預備咩ye 呢,一隻山羊羔,好重禮呀,好重禮,基奠屋企都有D 錢架喎,一隻山羊羔,又用一個伊法ge 細面做了無價餅..
咁一個伊法,Far,Far有幾多呢,又或者當時另一個ge 單位叫Bath,Bath,Bath,咁係咩意思呢,一個伊法,伊法同埋Bath呢,都係流體ge 單位黎架喎..
即係譬如我地今日用一升咁樣,一升水,一升牛奶啦,咁一個伊法ge 或者一個Bathge 細面有幾多呢,一個伊法有幾多呢,其實好多架喎,22公升架喎,22公升都好多架啦..
咁佢仲用22公升咁樣咁多ge 細面呢,即係靚ge 麵粉,就做jor 無價餅lor 黎..
而且gwo 隻山羊羔呢,係湯jor 之後,係煮jor 呀,就將D 肉,煮jor ge 肉就放係腔內,即係一個竹籃或者一個漬出黎ge 籃咁樣,然後將湯就盛在壺中,就帶到像樹下獻在侍者面前..
咁佢將湯呢個湯呢,係Maroc,Maroc係肉汁,即係話呢,佢湯一隻羊,煮好jor D 羊肉,咁煮ge 時候有汁架嘛,咁你就放係個壺gwo 度,咁我地今日就有暖壺呀,就有不鏽鋼壺拎黎呀,當日有架咩?冇架喎,其實係個鑊或者個鍋,咁你連個煮羊肉gwo 個鍋都拎埋黎喎..
咁你入面有D 羊肉,咁樣就分開搵個籃裝住D 羊肉,又拎住個鍋,入面有汁,咁然後呢仲要咩呢?拎住有一個依法二十二星細麵煮整出黎ge 無膠餅都好多個架啦,咁樣嘩!好似去Picnic咁樣拎過黎,咁都好有心ge ,黎到像樹下,將佢獻係侍者面前..
成日侍者就吩咐機電話,你將肉同埋無膠餅放在這盤石上啦,將湯即係D 肉汁倒出黎係個盤石gwo 度,咁佢就照住做啦,咁呢個盤石呢個字呢,希伯來文如果係見到佢譯做英文字母呢,C係讀Sge ,石疴,石疴,石疴呢就係山崖,唔係一舊石,唔係咁樣一舊石,係嘩!係一個好大舊石..
咁其實可能佢呢就已經係山上面架啦,咁呢就係個山崖巨石或者山寨,點解咁樣譯ge 呢?盤石係大石喎,好大舊喎,唔係一個咁樣咁叫盤石喎,係嘩!一座山咁樣..
所以聖經裡面話呢,你係我ge 盤石,係我ge 山寨,咁山寨即係話高ge ,好大舊石ge ,你呢擺gwo 度啦咁樣,咁好似一個山崖巨石咁樣,咁呢就呢度呢,當時ge 景況就咁睇字就唔係好明佢講緊乜ye 啦,咁呢就係你將佢擺係大舊ge 盤石上面,.
就呢將肉擺係度,將無教餅擺係度,將D 肉汁倒係個石gwo 度啦,咁佢就照住咁做啦.耶和華ge 使者就伸出手內ge 杖,杖頭呢嗌下D 肉同無教餅,一嗌,有火從盤石中出黎,燒盡了肉和無教餅,咁跟住耶和華的使者就不見了..
咁呢件事呢,好得意喎,就係神ge 使者將D 杖出黎,就咁掂下gwo D 肉,咁跟住個盤石就起火啦,於是個盤石就變jor 好似係一個祭壇咁樣,將gwo D 肉同埋餅燒jor ,咁呢個燒盡了呢,有火出黎燒盡了呢,係表示神ge 使者或者耶和華神本身係閱納jor gwo 個獻上ge 禮物ge ..
咁收jor 之後呢就不見了啦.咁幾年見她是耶和華的使者,點會而家先至知道ge 呢,即係話唔通咁樣講,佢開頭見到gwo 個佢唔信係耶和華ge 使者,你係咪真係耶和華ge 使者呀,咁呢就你比D 證據我,你真係耶和華ge 使者呀咁樣,咁都得,咁呢就,但係呢佢又會走返去獻禮物喎,即係話佢已經係接受呢位係耶和華ge 使者架啦..
咁所以gwo 度係有少少我呢度睇佢ge 經文呢,覺得佢有少少矛盾ge ,未係捉得好實ge .咁呢就嘩,竟然間燒jor 然後唔見jor ,佢就說,愛在主耶和華,因我,呢個讀迪ge ,因我迪面看見耶和華的使者..
咁得意ge ,佢講ye 好得意架,愛在阿拉,呢個拉長jor 音ge ,阿拉,阿拉呢就係慘囉,死囉,阿拉,慘囉,主耶和華,因我迪面看見耶華的使者..
我覺得基甸呢個,呢件事係有少少唔通ge ,講ge ye 係,佢同主耶華講,主耶華話慘啦,我見到耶華ge 使者,咁即係咩意思呢?.

$^{41}$呢個字呀,呢個中文字,我查過字典,叫迪,讀迪,粵語係讀迪,就係palim,palimel,palimel呢就係,palim係面呀,兩個字都係面,palim,palim,因我面向住耶和華使者的面,我同耶華使者面對面呀,所以呢,我死囉,唔好囉..
因為當時ge 人呢,係好怕神ge 使者,或者神自己向佢地顯現架,因為無人見過神架嘛,如果見過呢,佢係認為自己會死架啦,會死架啦..
咁因此呢,哎呀,原來真係耶華華使者呀,慘囉,耶華華,原來真係你ge 使者呀,呢次我死啦.咁都幾得意ge ,又會向耶華華講,耶華華,慘囉,我見到你ge 使者..
咁佢向耶華華講,咁其實呢,就係可能當時機電ge 心裡面呢,都有D 亂架,講ye 呢,唔係好清楚,有少少唔係好清楚自己想點樣.咁呢,不過呢,都有,有D 紋路喎..
咁神ge 使者向佢講,我與你同在,我揀jor 你去對付米電人,嘩,你真係耶華華ge 使者黎架,咁咪好呀,你俾D 證據我呀.咁佢又無講到咩證據,然後就話你等我,我返去lor D 公物返黎,lor D 製品黎比你..
咁佢返黎lor 喇喎,煮曬ye 咁拎黎喎,嘩,拎曬黎好多ye 喎.咁拎曬黎之後,咁神ge 使者就話你倒係gwo 度啦,跟住呢,燒jor .嘩,跟住,哎呀,你真係神ge 使者..
耶華華,我面見你ge 使者,我死jor .耶華華對佢話,呢個你字係無ge ,你放心,不要懼怕,你必不致死.第三天要講ge 呢,神就同佢講啦,放心..
放心呢個字呢,就好好喎,Shalom,希伯來文係Shalom,Shalom我諗大家聽過啦,Shalom就平安.咁而家ge 以色列人呢,猶太人呢,佢地見到面打招呼,佢地都唔係話Hello架,佢地係話Shalom,Shalom,平安,平安..
耶華華對佢話,平安,不要懼怕,你必不致死.呢個呢,就係,呢個呢,就係佢向神講ge 一個要求,比佢證據我啦咁樣..
咁呢,就於是好呀,但個證據係乜ye 呢?佢就做jor 祭物禮物獻比主,主又收jor ,然後佢就話死囉,我見到你,咁神就話唔使擔心..
於是,基澱在那裡為耶華華築了一座壇,名叫耶華華沙龍.耶華華沙龍呢,其實係譯音黎ge je ,沙龍係耶華華Shalom,耶華華Shalom就係耶華華是平安.咁今日好多人呢,就話耶華華沙龍係耶華華是平安,其實耶華華是平安,耶華華Shalom..
就呢,呢個壇就係亞比爾籍俄弗拉,俄弗拉意思就係又鹿,直到如今.呢個俄弗拉呢,應該係基澱出生或者佢成長ge 地方黎ge ,咁呢個意思就係又鹿,耶華華似平安..
讀到呢度呢,你會睇到,基澱係第一次求一個證據,神話係像樹下面,我揀你為我做一個人,去帶領全民勝過米殿人.咁基澱就話,我呀,如果真係你比D 證據我,第一次求證據..
佢呢邊呢,係好歡喜求證據架喎,你比佢證據我,未講完就話你等我,我返去煮D ye 黎你食咁樣,佢煮jor 拎到黎啦喎,咁耶華華使者就用D 象一點收jor ,火燒哂,咁佢就哎呀,原來你真係神ge 使者呀,佢就向耶華華神講,哎呀慘啦神呀,我見到你ge 使者,咁我呢次死啦,咁神就話唔會死,平安平安..
即係好似我地今日講,真定唔真定,唔使驚,唔會死,唔會死,於是呢,基本上gwo 個地方,改jor 個名,就係耶華華使平安,好啦,黎到呢度,我想問你,基澱第一次求證據,gwo 個證據係咩呢?讀完呢個故事,個證據就係,一句說話,主耶華華同佢講,.
我必與你同在,其實神一早就比jor 證據,我必與你同在,但佢又個心唔係最接受到,唔係最信,於是佢做jor D ye ,咁又無講做乜ye ,不過神就好啦,我用收jor 你ge 禮物,再比一個證據你睇下,就係我遇納你,得啦呱..
然後佢就話,哎呀,死啦,我見到神ge 使者,我一定冇命啦,神就話,平安,定D 黎,於是佢就改jor 個叫耶華華沙龍ge 地方,咁因此呢,其實基澱係開頭,我個人對呢個解釋呢,佢係有D 慌亂架,忽然之間,嘩,上帝話我揀你受寵若驚,係咪真係揀我架,得唔得架,你有冇咩證據比我呀?.
大家冇傾過證據,神只係話,我與你同在,我必與你同在,你會贏,米殿人好似同一個人打交咁樣,失贏,咁然後呢,就收jor 佢ge 禮物,然後佢話我死啦,佢就話平安,耶華華係平安,所以呢,係有神ge 同在,有平安,呢個先至係一個最重要ge 證據比基澱,.
就是月立佢,賜佢平安,與佢同在,呢個呢,係最大ge 證據啦,咁呢個呢,就係我地一齊要學習架喎,咁呢,就有時見到神特別叫我做一樣ye 呢,唔好好似基澱咁第一次呢,就嘩,慌亂晒呀,有少少唔知道自己想點,咁呢,就有D ,我唔敢話基澱咁偉大ge 人物,語無倫次啦,.
不過呢,佢就係有D 方寸,有D 亂呀,咁呢,就不過到最尾,上帝比佢ge 安係咩呢?Shalom,平安,定D 黎,平安,耶華華是平安,呢個呢,就係我地最寶貴架啦,我地一齊祈禱,主佑有時候我地面對住一個好大ge 挑戰,或者你比我地一個好大ge 責任,心裡面係驚ge ,.
而且諗返基澱當日實在係值得驚ge ,因為係講緊打仗,因為係講緊對付米電咁強大ge 一個部族,而且聯埋jor 好多個部族一齊黎同我地以色列人打,已經比佢地糟蹋jor 好多年,心裡面真係驚ge ,.
甚至基澱呢一個咁大能勇士,佢都要地係酒駕gwo 度打墨紙,甚至佢ge 信心都失jor ,話其實神族邊有應許同我地同在呢?所以見到佢有咁大ge 驚慌,咁慌竄失去,係對ge ,主佑有時候我地都覺得自己好強,好勁,不過係一個比我地重大ge 挑戰面前,我地都會覺得驚慌,.
而呢個時候神你就黎到,你同我地講安慰ge 說話,你比我地ge 證據就係你同在ge 平安,主佑幫助我地去尋求有你同在ge 平安,以致我地心裡面係帶住喜樂,係帶住平安去面對新ge 事奉,工作與挑戰,幫助我地,禱告奉耶穌名求,阿們..
\newpage



\section{希伯來書 9:22-28}
\label{sec:032278}
\textbf{向過去說再見,向未來說嗨}~~~~謝又生牧師
\newline
\newline
連結: \href{http://www.church.com.hk/acms/content.asp?site=cdc\&op=show\&type=product\&code=032278\&layout=sermon}{\texttt{032278}} ~~~~ 語音日期: 2026-06-07
\newline
\newline
\hyperref[sec:032277]{\small{< < < PREV SERMON < < <}}
~
\hyperref[sec:index_chronic]{\small{[返順時目]}}
~
\hyperref[sec:032279]{\small{> > > NEXT SERMON > > >}}
\newline
\newline
希伯來書 9:22-28
\newline
\begin{longtable}{cl}
\hline
\hline
章節 & 經文 (和合本修訂版)\\
\hline
9:22 & \begin{tabularx}{0.7\textwidth}{X} 按著律法，幾乎每樣東西都是用血潔淨的；沒有流血，就沒有赦罪。 \end{tabularx} \\ \\ \relax
9:23 & \begin{tabularx}{0.7\textwidth}{X} 這樣，照著天上樣式做的物件必須用這些禮儀去潔淨，但那天上的一切，自然當用更美的祭物去潔淨。 \end{tabularx} \\ \\ \relax
9:24 & \begin{tabularx}{0.7\textwidth}{X} 因為基督並沒有進了人手所造的聖所—這不過是真聖所的影像—而是進到天上，如今為我們出現在神面前。 \end{tabularx} \\ \\ \relax
9:25 & \begin{tabularx}{0.7\textwidth}{X} 他也無須多次將自己獻上，像大祭司每年帶著牛羊的血進入至聖所。 \end{tabularx} \\ \\ \relax
9:26 & \begin{tabularx}{0.7\textwidth}{X} 如果這樣，他從創世以來就必須多次受苦了。但如今，他在今世的末期顯現，僅一次把自己獻為祭，好除掉罪。 \end{tabularx} \\ \\ \relax
9:27 & \begin{tabularx}{0.7\textwidth}{X} 按著命定，人人都有一死，死後且有審判。 \end{tabularx} \\ \\ \relax
9:28 & \begin{tabularx}{0.7\textwidth}{X} 同樣，基督既然一次獻上，擔當了許多人的罪，將來要第二次顯現，與罪無關，而是為了拯救熱切等候他的人。 \end{tabularx} \\ \\
[1ex]
\hline
\hline
\end{longtable}
$^{1}$猶太人有一個這樣的故事.
Simon一直在一條路上走.
要尋找他的人生.
當他走到一條橋邊的時候.
有一個人拿著一條繩子向他走.
很奇怪.
這個人的樣貌和衣著竟然和Simon一模一樣.
這個人就跟Simon說.
先生 你可不可以幫我解開這條繩子.
其實很簡單.
你只要抓住這條繩子這一端.
一路向前走.
無論任何的情況.
你都不能夠放手.
這樣就可以了.
Simon覺得也是舉手之勞.
於是就答應了.
但當他一路拿著這條繩子走過橋的時候.
突然間有一個強大的拉力.
將他向後攀.
究竟發生了什麼事呢.
原來剛剛跟他說完話的這個人.
他自己跳到橋上了.
大家可以想像.
一個人拿著繩子的另一端.
當他一路跳到橋上.
Simon就一直被他拉下去.
Simon很想將他拉上來.
但是力氣不夠.
他既然拉不到他上來.
就希望將繩子的另一端綁在橋邊.
但是那個拉力實在太大了.
他沒有辦法綁住.
他想放手.
但掉下去的那個人就大叫.
你千萬不要放手.
如果你一放手你會自責一生.
Simon用盡一切的能力.
都沒有辦法將他拉上來.
自己又慢慢被他拉下去了.

$^{41}$於是他覺得必須要做一個決定.
他想了一會就說了以下的說話.
他說每一個人必須要為自己的人生作出抉擇.
說完這句話他就放手.
然後繼續走他的路.
弟兄姊妹,你公佈到Simon遇見這個人代表什麼呢?.
這個人就是代表我們的過去.
可能是過去的內疚.
過去的遺憾.
過去的失望.
過去的羞恥.
過去的恐懼.
甚至是過去的憤怒.
過去的苦毒.
隨著日子過去慢慢.
這些的過去就會成為我們的包袱.
如果我們不學習去放手.
讓過去真的成為過去.
我們學習跟過去說再見.
慢慢這些的過去就會.
跟我們的人生糾纏不清.
可能我們慢慢就會不斷被這些過去捆綁.
我曾經為一位完美主義者的弟兄禱告.
這位弟兄告訴我.
我無論做任何事情.
都定了一個最高的標準.
這個標準.
我如果做到.
我就很驕傲.
於是我就看不起其他人.
但如果我做不到.
我突然變得很自卑.
很沮喪.
一團糟.
這幾年.
他就不斷在這些自卑和沮喪徘徊.
他的心情和情緒就像過山車般大起大跌.
我聽完之後就帶他去聖經.
跟他分享一句聖經.
就是哥林多前書的十三章十一節.

$^{81}$他說什麼呢.
他說我們是小孩子的時候.
我們說話像小孩子.
想事情像小孩子.
心思念念都像小孩子.
如果我們長大的時候.
我們將這些小孩子的東西完全丟棄.
亭姐妹,今天神要和大家說.
我們已經在耶穌基督裡面.
已經長大成人了.
長大成人是什麼呢.
我們要學習.
丟下那些不好的過去.
如果有些過去纏擾我們.
我們要學習和它說再見.
放手.
讓過去真的成為過去.
同樣神要為我們預備新的將來.
新的盼望.
所以我們應該帶著期盼.
去擁抱我們未來的人生.
不僅是哥林多前書這樣說.
希伯來書也是這樣說.
我們一起看看希伯來書五章十三節至六章一節.
因為希伯來書.
原本和合本是比較失誤的字.
大家比較難明白.
所以我用了和修版.
我們一起讀一讀.
希伯來書五章十三節至六章一節.
一二三.
凡只能吃奶的.
就不熟練仁義的道理.
因為他是嬰孩.
唯獨長大成人的才能吃乾糧.
他們的心跳因練習而靈活.
能分辨善惡了.
所以我們應當離開基督道理的基礎.
竭力進到成熟的地步.
不必再納根基.

$^{121}$我解釋一下.
嬰孩的意思不僅是肉身的BB.
而是屬靈生命的BB.
BB只能吃奶.
你有沒有人見過.
將出前一丁麵餵給BB吃.
不會有的.
是吃不到的.
他只能夠喝奶.
喝奶代表什麼.
不熟練仁義的道理.
就是說.
屬靈生命的BB階段的人.
他是分辨不到什麼是好什麼是不好的.
就舉剛才那個原理主義者的弟兄為例.
他小時候.
他跟我說一定要做到最好.
一定要做到最好.
如果你做得不好.
就不會有人疼你.
不會有人再愛你.
他的成長.
他沒有放下過去.
他就被過去一直綁住.
他分辨不到什麼是不好什麼是好.
過去的不好他無法放下.
他反而被不好的東西控制他這個人生.
所以他是分辨不到.
什麼是好什麼是不好.
但是屬靈生命長大成人的人.
就是快高長大的人.
這些人是可以吃飯.
可以吃固體食物的.
他們代表什麼.
就是說這些人是分辨到好.
分辨什麼是不好.
分到兩樣東西.
所以他們能夠將他們的過去放下.
所以喜不來時的作者.
不斷鼓勵我們.

$^{161}$要離開基督道理的基礎.
也就是說我們要離開屬靈生命幼稚的階段.
就好像剛才那些小五小六的小朋友.
我不做那些幼稚的事.
因為我長大了.
我開始慢慢會成人了.
長大了.
今天我們快高長大的時候.
我們必須要放下那些過去.
擁抱我們的將來.
但是問題是我們憑著什麼.
去放下過去呢.
喜不來時的作者告訴我們.
耶穌已經做了一些事情.
耶穌基督已經做了一些事情.
我們能夠有能力放下過去.
是因為耶穌開始了一些新的事情.
為我們完成了一些事情的事實.
今天我想和大家講耶穌完成的工作有兩個事實.
第一個事實.
就是耶穌流血.
跌在十字架.
其實祂不只是跌在十字架.
祂為我們開始一個新的局面.
是一個新的方向.
我們一起繼續看希伯來書的九章二十二至二十四節.
等一會我們一起讀.
因為基督並沒有進了人手所做的聖所.
只不過是真聖所的影像.
而是進到天上.
如今為我們出現在神面前.
等一會我想解釋一下用血潔淨.
用血潔淨的意思是在舊約聖經摩西律法.
人犯罪必須要獻祭.
獻祭是將一些牛羊那些動物來到祭壇那裡.
割了牠們流血.
用那些血洗淨犯罪人的罪.
但是你想想.
其實你有沒有想過呢?.
你聽了那麼多.

$^{201}$那麼血是怎樣潔淨人的呢?.
血怎樣去潔淨我們呢?.
你會不會去紅十字會拿一包血來洗臉呢?.
你會潔淨了.
如果Eleanor這樣.
我拿一些血來洗臉.
我就不會說她潔淨.
我只會說她恐怖.
又或者你不小心弄傷手指.
手指流血.
你會不會覺得手指乾淨了.
不會的.
我們經常說耶穌的寶血洗淨我們.
那是什麼意思呢?.
我要解釋一下.
在聖經裡面說.
罪的代價就是死.
其實死代表什麼?.
死是人的步向毀滅.
今天這個世界充滿罪.
人的代價就是死.
步向毀滅.
整個世界都好像充滿毀滅.
如果大家同意.
這個世界這麼混亂.
人犯罪的代價就是步向毀滅.
血是代表生命.
所以大家看舊約聖經的時候.
他又說不要流人血.
就是不要摧毀別人的生命.
不要殺人.
血是代表生命.
其實是指向一個事實.
就是耶穌用祂自己的生命.
幫我們犯罪.
我們有很多代價要結帳的.
耶穌就幫我們結帳了.
用祂的生命幫我們結帳了.
結果是.
我們信了耶穌之後.

$^{241}$不是說我們個個原理.
而是在神的眼裡面.
當我們好像沒有犯過罪.
到最後審判的時候.
我們不致滅亡.
反得永生.
所以.
用血來洗淨就是.
耶穌用祂的生命代替了我們犯罪的代價.
結帳.
然後到最後大審判的時候.
我們不致滅亡.
反而有永生.
所以這個就是一個事實.
這個永生.
耶穌為我們結帳.
是開設一些新的局面的.
還有怎樣新法呢.
我們再解釋一下.
沒有流血.
就沒有赦罪.
我剛才說舊約聖經裡面.
人犯罪要獻祭.
但不是一次.
你犯一次罪就要獻一次祭.
所以大家可以想像.
永無止境.
你犯罪就要獻祭.
因為大家都知道我們會犯罪.
犯罪就要獻祭.
如是者繼續犯罪獻祭下去.
就好像.
舊約聖經裡面犯罪獻祭.
其實就等於我們生病.
我們沒有找出根源.
我們只是吃Panadol.
你吃一次就止痛.
然後藥力過了.
然後又再吃.
如此類推繼續下去.

$^{281}$這些是治標不治本.
耶穌就不同了.
耶穌就是特效藥.
一劑搞定.
吃一次就搞定.
換句話說.
耶穌一次過成為我們.
變為我們人類.
一次過來到這個世界.
我們叫做道成肉身.
一次過將自己獻祭.
代我們犯罪結束.
所以一次過就完成了工作.
而賜給我們永生的新約.
聖經的新約.
賜給我們永生的新約.
就馬上開始.
我舉個例子.
我相信大家有用過電腦.
大家用電腦的時候.
你覺得檔案舊了.
不用了.
你是不是要刪除它.
那你按了刪除鍵.
是不是真的刪除了.
不是的.
然後有個訊息彈出來了.
你是不是真的刪除了.
你說是的.
我是真的刪除了.
確定.
那就先刪除了.
耶穌就是這樣.
耶穌一定要死.
一定要流血.
祂一流血.
在十誡上說.
成了就是確定了這個刪除.
我們就馬上我們的罪被赦免.
我們賜給我們永生的新約.

$^{321}$就馬上開始.
新約開始代表什麼.
新約開始是代表.
神已經為我們創造了一個新的局面.
什麼是新的局面.
新的局面是新的一頁.
新的開始.
新的人生.
新的未來.
我們是否擁抱這個新的未來呢.
怎樣新化呢.
等一會我們繼續看.
希伯來書九章二十五至二十六節.
我們一起讀.
一二三.
他也無需多次將自己獻上.
像大祭司每年帶著牛羊的血.
進入至聖所.
如果這樣.
他從創世以來就必須多次受苦了.
但如今他在今世的末期顯現.
僅一次把自己獻上為祭.
好除掉罪.
等一會這個僅一次就是一個新的.
因為舊約聖經裡面不斷地發誓不斷地獻祭.
但耶穌一次過就成就了心意.
開創了一個新的局面.
好除掉罪.
就是那個結果了.
耶穌赦罪的結果是包括了.
全人類所有人.
由古至今直至將來.
所有覆蓋.
那個影響力是什麼呢.
我們的過去都被除去.
那些不好的過去都被除去.
好除掉罪的罪.
在希伯來書出現了二十五次.
在整本聖經出現一百七十三次.
除掉的意思.

$^{361}$就是我們等於禁Delete那樣東西.
而這個罪是代表我們的過去.
可能我們的過去今日.
直至今日我們仍然有很多人被過去都仍然綁住.
被過去的勢力仍然控制住.
可能我們今日仍然在那種內疚.
在那種遺憾.
在那種失望.
在那種苦獨.
在那種恐懼裡面.
今日仍然綁住.
但耶穌在十字架上面成就的工作.
就等於禁了那個Delete制.
因為有聖靈復活的大能幫助我們.
耶穌曾經說過.
若有人在基督裡就是新造的人.
舊事過都變成新的了.
等於今天有沒有一些過去仍然在控制著你的今日.
控制著你的人生.
控制著你的情緒.
控制著你的思考.
耶穌說若有人在基督裡就是新造的人.
舊事過都變成新的了.
耶穌要將這個說話宣告在你的生命裡面.
將這個真理宣告在你的生命裡面.
今日第一個的事實.
就是耶穌為我們流血.
成就開創一個新的局面.
這個事實.
我們可以基於這個事實放下我們的過去.
第二個事實.
就是耶穌會再回來的.
祂會再回來完成祂答應了我們的救恩.
我們繼續看希伯來書.
希伯來書的九章.
二十七至二十八節.
我們讀一二三.
安著定命人人都有一死.
死後且有審判.
同樣基督既然一次被獻.

$^{401}$擔擔了多人的罪.
將來要向那等候他的人第二次顯現.
與罪無關.
而是為拯救熱切等候他的人.
第二次顯現就是可能大家都知道的主再來的時候.
就是耶穌將來會再回來.
祂回來是做些什麼呢.
記住第一次耶穌來的時候.
是釘十字架.
犧牲自己.
流血為我們罪找一數.
這是第一次來的使命.
第二次來的使命就不是釘十架.
是祂答應我們的救恩.
是完全成就.
換句話說將來的日子裡面.
今天這個世界可能大家都看到.
步向毀滅這麼混亂.
不知道怎麼辦.
將來是沒有這件事的.
因為所有這個世界混亂根源的罪已經拿走了.
我們會和耶穌在一起.
但是耶穌還沒回來之前.
祂已經賜下聖靈給我們.
所以今天是聖靈的時代.
耶穌說主的靈在哪裡.
哪裡就得著自由.
所以今天聖靈是可以幫助我們.
可以放下我們的過去.
影響我們的將來.
等於我們基於這兩個事實.
我們怎樣和聖靈配合.
以致我們真的有能力去放下我們的過去.
人面對這件事有三種不同選擇的態度.
第一種態度.
有人回望過去.
他只是找一些為他今天的行為而解釋合理化的藉口.
包括什麼呢.
我今天這麼喜歡控制別人.
這麼喜歡作主.

$^{441}$這麼喜歡掌控別人.
不能怪我.
不是怪我.
最壞是我父母在我小時候就離婚.
搞到我沒有安全感.
我沒有安全感就要控制.
所以不關我的事.
最壞是他們.
這是藉口.
第一個態度.
第二個態度.
就是有人會否定過去.
當沒有事情發生過.
傷害?.
什麼傷害?.
沒有.
我一點感覺都沒有.
這是第二個態度.
第三種態度.
就是有人回望過去.
他找找我的生命有什麼需要改變呢.
我有什麼需要進步呢.
我有沒有什麼需要給聖靈去醫治呢.
我有沒有地方是要原諒人的呢.
我有沒有地方是需要求人的原諒呢.
我有沒有地方是盡力去與人和好呢.
這種態度是同感聖靈配合.
讓我們真正放下我們的過去.
你又要選擇什麼態度.
去配合聖靈的工作.
以至我們真正放下我們的過去呢.
有位弟兄今天已經主懷安息了.
在他回天家之前.
我問他.
我可不可以和教會的弟兄姊妹.
分享你的見證.
他說可以.
不過請你不要說我的真名.
所以我就習慣幫人家改這個綽號.
我叫他S弟兄.

$^{481}$S弟兄的下半生.
一直被他的內疚所折磨.
1949年.
S弟兄帶著他太太和四個兒子來香港定居.
到了1967年.
當時的內地沒有現在那麼先進那麼進步.
當年物質很缺乏.
很多人生活都有點艱難.
所以S弟兄見到他的親戚.
也是物質缺乏生活艱難.
於是就叫太太和大兒子帶了很多物資.
和少量現金回鄉.
大家現在看到的圖.
那個現在不存在的.
就是當年的尖沙咀火車站.
那個紅磡火車站其實當年是尖沙咀.
當時回鄉探親.
他們帶著一擔擔東西.
很多物資.
其實就是香港帶去的物資.
回到內地去接濟親戚.
太太和大兒子就這樣回到內地.
但是當他一回到內地的時候.
就馬上被人抓住了.
當年1967年是文化大革命最激烈的時候.
所以他們被人抓住.
因為他們以前在鄉下的背景是地主.
所以被打為地主階級.
於是五花大綁.
被批鬥.
很多的折磨和侮辱.
太太終於受不了那種侮辱.
在第二天就吊頸自殺.
S弟兄非常自責.
他沒辦法原諒自己做了一個錯誤的決定.
不但害死他太太.
甚至他的大兒子因為這麼混亂.
沒辦法回來香港和他其他家人團聚.
所以S弟兄決定終生不娶.
用自己的努力帶大三個兒子.

$^{521}$但是由那個時候開始.
S弟的性情有很大的改變.
他變成一個非常孤僻.
非常抑鬱.
並且經常發脾氣的人.
就因為這樣.
他這麼大的脾氣.
就和他三個兒子的成長關係一點也不好.
直至弟兄信耶穌的那一刻.
我有幸那一次有機會和S弟兄做決志祈禱.
我一路帶他去祈禱的時候.
我心裡面湧現一句聖經.
就是加泰書的三章十三節.
他說:基督既為我們受了詛咒.
救贖出我們脫離律法的詛咒.
我讀這段聖經給S弟兄聽.
接著為他祈禱.
S弟兄後來跟我說.
當他閉上眼的時候.
他彷彿覺得耶穌基督親自來到他身邊.
親自背負他一切的內疚在耶穌身上.
我從來沒見過一個幾十歲的大男人.
哭得好像小朋友一樣.
很感動.
我都沒辦法忍住我的眼淚.
所以兩個男人就抱著在那裡哭.
這位S弟兄的重擔真的離開了他.
那個內疚真的離開了他.
因為由那個時候開始.
他的性格又變得隨和.
慢慢變得平易近人.
連他之前的大脾氣都沒有了.
但最大的突破是什麼呢.
我相信不到自己的眼睛.
一個幾十歲的長輩.
親自走到他三個兒子的面前.
跟他說 兒子 對不起.
求你寬恕我.
求你寬恕爸爸.
弟兄姊妹.

$^{561}$過去再不可以在S弟兄身上形成枷鎖.
過去再不可以綁住他.
過去再不可以控制他的未來.
剛才那位長輩幾十歲 年紀大一點.
以下這位Clannon年紀小一點.
Clannon的童年是很不開心的.
因為他家裡都比較吵鬧.
第一 他媽媽長期遇酒.
第二 他爸爸長期吸毒.
連他爺爺奶奶的婚姻.
都在他小時候破裂.
所以大家可以想像Clannon的處境.
他的成長是很困難的.
所以他很小的時候就開始迷失了.
包括十二歲跟人打架.
可能跟人打架而已.
小朋友打架沒問題.
但他打架是打到去警署的.
十三歲開始偷東西.
十六歲跟女朋友擲籃子.
但他擲籃子的那年.
反而成為Clannon的人生轉捩點.
因為他有一群很愛他的朋友.
一手抓了他回教會.
並且在教會裡信了耶穌.
Clannon很小的時候都開始懂事.
他發覺他意識到.
如果他想真的放下他的過去.
他必須要放下他對家人的憎恨.
所以Clannon做了一個.
很清楚他心裡的決定.
以往他跟家人怎樣相處呢.
我剛才提過他媽媽是長期遇酒的.
遇酒喝大了不只是暈倒那麼簡單.
他是行為不檢亂的.
因為他喝大了行為不檢亂經常被人拉進警署.
以往的Clannon只是在哭.
在埋怨在哭.
但Clannon這次很不同.
他走去媽媽的房門.

$^{601}$跪下舉手.
為他媽媽的生命改變禱告.
為他跟媽媽的關係去祈禱.
慢慢由那時開始他發覺.
慢慢他跟媽媽的關係有點不同了.
很明顯的不同了.
是說話的時間多了.
以往他媽媽喝大了就躺在那裡.
或者被人抓了或者不理家人.
Clannon就說我用不著理他.
就不理他.
大家不理的.
沒什麼機會沒什麼兩句.
但由那時開始他發覺.
他跟媽媽的關係開始有少少改善.
開始聊天多了.
開始聊的時間又長了.
直至最近有一次.
他媽媽跟他說兒子你在家播那些詩歌.
不知為什麼.
你每次播我都心裡平安點.
你可不可以把詩歌寫在紙上.
把詩歌名寫給我.
他覺得有點不同了.
第二個不同.
就是他爸爸.
剛才提過他爸爸是長期吸毒的.
長期吸毒會犯事.
又會被人抓去坐牢.
他爸爸就抓去坐牢.
坐牢.
他爸爸最近出獄發現自己患了鼻咽癌.
於是Clannon就順理成章邀請他回教會.
但雖然是這樣.
Clannon覺得是應該做的事.
但他心裡對爸爸的憎恨是沒有消除過的.
直至有一次祈禱的時候.
在參加祈禱會.
他一開始祈禱.
他發覺不知為什麼好像有一種能力幫助他.

$^{641}$他慢慢覺得不太憎恨他爸爸.
開始他對爸爸的憎恨消失了.
那一年的聖誕節.
他有生以來第一次親手寫信給爸爸.
還有他有生以來親口第一次稱呼他爸爸是爸爸.
他信裡面大約的意思就是說.
爸爸我希望在這個聖誕節.
我很希望你真的明白.
你不需要為以往做錯過的事愧疚.
你也不需要為你對我們以往的態度去愧疚.
因為我已經在乎你.
我所信的耶穌改變了我的心很大很大.
我也相信我所信的耶穌一定可以改變你.
那年聖誕節.
第二年的一月他爸爸在教會接受浸禮.
接受浸禮之後他爸爸也在教會開始事奉.
後來我看Clannon的見證的時候他這樣說.
他說我一直很長時間為家人認識耶穌祈禱.
但我怎麼也想不到.
我爸爸是我家裡第一個信耶穌的.
是第一個生命改變的.
今天我的爸爸不單在肉身上得到醫治.
他在屬靈上也得到醫治.
以往我爸爸是被毒癮綁住脫不了的.
今天我的爸爸可以在教會大敬拜.
帶人進入神的同在.
弟兄姊妹.
過去是不可以綁住我們的.
我邀請大家停下幾秒.
你很鄭重地跟你父親弟兄姊妹說.
過去是不可以綁住我們的未來的.
今天耶穌要在我們當中宣佈.
過去是不可以控制我們的未來.
過去是不可以決定我們的未來.
過去是不可以攔阻我們領受.
神為我們預備的新的未來,新的人生的.
Amen.
弟兄姊妹.
我說的不單止是我們個人的屬靈生命可以放下過去.
迎向未來.

$^{681}$我們的群體,我們的教會.
也可以放下過去迎向未來.
弟兄姊妹,我們的教會叫什麼名字?.
這麼小聲的?.
對了,同福堂.
我作為主教牧師,我很長期為教會祈禱.
最近禱告,我心裡有一個很大的感受.
就是這一兩年是同福堂非常重要的兩年.
今年我們的主題是多代同行.
我剛才也提過,多代同行我們會有很多一起的機會.
如果大家有留意,我們今年多了很多聯合崇拜.
弟兄姊妹,聯合崇拜不是為了聯合崇拜而聯合崇拜.
其實聯合崇拜是很重要的屬靈意義.
因為教會是耶穌的身體,是屬靈的大家庭.
所以我們不同年紀走在一起去敬拜神.
是一個很重要的屬靈的意思.
代表我們宣告耶穌是我們的神.
也好像我們在灣景樓上的那塊牌寫什麼.
耶穌是主,是我們的主.
所以是一個非常重要的屬靈意義.
不單止是聯合崇拜.
我們這個暑假,我們的兒童和少年人會有更多的機會合作.
一起去事奉神.
我們的少年人也將會去到金靈群體.
跟他們有很多互動.
我們的職青群體也將會去到向日葵.
我們哥哥姐姐長輩的群體.
跟他們有很多交流.
我們同福堂是社區教會.
我們對大坑這個地區有很多的服侍.
我們的服侍這一年也有很多不同年紀的弟兄姊妹一起去參與.
所以你看到這個同福堂是有思耐的.
未有過的這麼大的可以一起的聯合.
也是我很相信.
也是神在預備我們不同年紀的弟兄姊妹.
進到我們教會發展的另一個階段.
就是明年的關心社區的主題.
這兩個主題,多對同行的關心社區主題.
中間是一個很重要的里程碑.
就是我們今年的聖誕崇拜.

$^{721}$弟兄姊妹今年的聖誕崇拜會在12月24日.
平安夜那天的黃昏開始.
我們會去到亞洲博覽館.
一場過崇拜.
讓所有的弟兄姊妹一起可以同一場地.
同一場,不用分幾場.
一起敬拜我們的神,好不好呀?.
弟兄姊妹,這個不單止是同福堂今年踏入35週年.
我們為了神這麼多年帶領我們去獻上感恩崇拜.
但不只是這樣.
不單止是感恩.
也是我們同福堂每一位弟兄姊妹一起有份敬拜.
有份邀請人.
有份傳福音.
有份事奉的機會.
所以是一個報道會.
我也會邀請不同地區的弱勢社群.
很大機會是用車載他們來亞洲博覽館.
因為那個場地不單止夠地方讓我們所有弟兄姊妹一起敬拜.
並且多了很多地方讓我們有很多的傳福音活動.
詳情我就交給今屆的籌委他們去決定.
但我想跟大家說一件事.
香港人,我們都是香港人.
我們很流行計劃.
我們很流行計劃聖誕的時候要怎樣做.
去外遊,要怎樣的活動.
所以我在半年前我跟大家說了一個預言.
這個預言是什麼呢?.
請大家在日曆上記住12月24日那一天.
如果大家可以的話.
請你預先那天如果你要上班的話.
請你考慮請半天假或者一天假.
在這半年裡面我們的童工會找大家.
因為我們盼望每位教會的弟妹都可以一起去服侍事奉傳福音.
所以你先記住.
如果有童工找你.
Joshua,那天可不可以一起事奉.
你就可以很有型的跟他說.
知道了,我早就知道了.
我記住了,我請假了.

$^{761}$我將那天奉獻給神.
那你就很有型了.
還有那一天會是黃昏開始.
我們的活動.
那天不會很晚結束.
所以如果有些弟兄姊妹.
你是計劃聖誕去外遊的話.
記住亞洲博覽館和機場只差一個站.
所以你可以預備.
你去服侍完,經擺完,傳完福音.
然後跟神說,成了.
然後你可以拖著你的行李.
平安地,喜樂地上機.
或者不要這麼趕.
我要十二二十五天,又可以.
你說我是萬馬斯斯人的執行禮.
平平安安地,那天飛機又成了.
所以弟兄姊妹.
盼望我們可以一起同心的事奉.
我們很盼望.
那個聖誕崇拜是我們.
一起可以事奉服侍神的機會.
同心合意的.
其實聖誕崇拜我剛才提到.
是一個里程碑.
由今年的多代同俠.
去到明年的關心社區.
所以弟兄姊妹我們會重啟東區異象.
我不知道有多少弟兄姊妹.
知道東區異象.
你見這個範圍.
就是港島東區.
由銅鑼灣去到柴灣和小西灣.
如果當年在同福堂的弟兄姊妹.
你會知道.
東區異象是怎麼來的.
是神興起我們這個群體.
我們同福堂這個群體.
讓我們今天有這個地方.
去崇拜去敬拜神.

$^{801}$就是當年神說我聽見東區孤兒寡婦哭聲.
我聽見東區弱勢社群貧窮人的哭聲.
所以他在撒旦的座位興起他的教會.
他會差遣這群人向東區的弱勢社群傳福音.
讓他們得著耶穌基督的救恩.
所以弟妹我們會重啟東區異象.
所以我們盼望全教會.
明年關心社區.
我們的社區不單是大廈.
我們的社區是整個港島東區.
弟妹我不是只說明年.
其實今年我們已經有行動了.
現在大家看的照片.
就是我們的童工.
已經和柴灣同學的童工.
去行了兩次的行區祈禱.
我們行區祈禱是做什麼的呢.
是為整個地區神所愛的人去祝福.
是為他們禱告.
讓神的救恩讓他們能夠得著.
我們也求神給我們一些新的計劃.
新的做法給我們.
令我們整個同福堂可以配合神的計劃去工作.
所以弟妹明年我們啟動重啟港島東區的異象.
我們需要大量的人力,物力,財力.
請弟妹為教會祈禱.
我們很盼望神可以使用我們的群體.
可以祝福整個地區.
還有明年.
明年我覺得.
明年的整個香港.
屬靈上面有一個很大的逆轉.
為什麼呢.
明年的十一月很有機會.
因為他們現在還沒有確定.
將會在政府大球場有六場的報道會.
是全香港的報道會.
這個是當日的發佈會.
但是早幾個月.
我去香港報道大會的異象分享會的時候.

$^{841}$我有一個很深刻的感動.
為什麼呢.
因為我去到的時候.
我發覺香港的教會很積極參與.
包括.
我不知道大家知不知道.
香港的基督教中派一共有三十個中派.
但是有二十五個中派.
中派代表很多教會.
一起參與.
香港所有的神學院.
都一起參與.
香港絕大部分的基督教機構.
都鼎力支持.
所以我那天去到的時候.
我有一個很大的感動.
我覺得同心合意.
興旺福音就是這樣來的.
回顧我們2027年19年.
可能大家都感受到.
香港的基督教圈子.
很多的混亂很多的決裂.
直至那天我感受到.
我們神在香港有更大的工作要做.
就是同心合意.
興旺福音.
你知道為什麼叫做同福堂嗎.
叫做同福堂是因為同心合意.
興旺福音.
所以很重要.
傳福音是我們的使命.
也是教會的使命.
也是同福堂的核心價值.
所以我心裡面很大感動.
於是我就走回來.
跟我們的同工分享.
他們全部都一致支持.
跟著我就跟執事會分享.
執事會一致支持.
所以我很感受到.

$^{881}$大家都有一個很大的心.
如果神在香港.
神在香港才是祂大能的手.
讓我們整個.
香港的信徒群體.
可以同心合意.
興旺福音.
是非常之美好的事.
所以明年我們一定會全力去參與.
我們盼望.
包括我們自己的社區.
港島東區.
我們的.
或者我們自己的家人.
我們認識的圈子.
能夠合作救人.
我相信頂智慧神.
不只是呼籲一些小絕的人.
神從來都不是.
只是說那些精英.
不是的,神是呼籲每一個.
所以在我們當中.
無論你年紀有多大.
小朋友也是.
年輕人也是.
大朋友也是.
長輩也是.
無論你是什麼年紀.
無論你是什麼背景.
無論你哪一年加入同福堂.
這個都不是最重要的.
因為神呼召你和我.
一起去參與.
讓我們為自己禱告.
讓神將這個全福音的熱忱.
這個服侍人的熱忱.
放在我們的心裡面.
因此我們教會真的可以.
同心合意,興旺福音.
頂智慧你願不願意這樣跟從神呢.

$^{921}$我們一起讀一次.
希伯來書.
十二章一至二節.
一二三.
我們既有許多的見證人.
如同雲彩圍著我們.
就當放下.
各樣的重擔.
脫去容易.
纏累我們的罪.
全心忍耐.
不那擺在我們.
前頭的路程.
仰望我們.
創始成眾的耶穌.
頂智慧的.
重擔和罪.
剛才說的就是我們的過去.
那些我們要放下.
的過去.
所以神說得很清楚,要放下.
要脫去,要我們離開.
那些過往的事.
那些一切.
神不喜悅的事,又或者.
不好的事,讓我們.
不要被這些東西綁住.
因此我們可以不那擺在.
我們前頭的路程.
我們一起同心去.
帶著期盼.
去迎接神為我們.
預備新的未來,新的方向.
新的一頁,新的人生.
同樣不單止是個人.
我們教會都一樣.
迎接新的一頁.
新的方向.
所以頂智慧,你又願不願意.
這樣去跟從神呢?.

\newpage



\section{馬太福音 3:1-3}
\label{sec:032279}
\textbf{天國的福音}~~~~王梯青長老
\newline
\newline
連結: \href{http://www.church.com.hk/acms/content.asp?site=cdc\&op=show\&type=product\&code=032279\&layout=sermon}{\texttt{032279}} ~~~~ 語音日期: 2026-06-07
\newline
\newline
\hyperref[sec:032278]{\small{< < < PREV SERMON < < <}}
~
\hyperref[sec:index_chronic]{\small{[返順時目]}}
~
\hyperref[sec:032280]{\small{> > > NEXT SERMON > > >}}
\newline
\newline
馬太福音 3:1-3
\newline
\begin{longtable}{cl}
\hline
\hline
章節 & 經文 (和合本修訂版)\\
\hline
3:1 & \begin{tabularx}{0.7\textwidth}{X} 在那些日子，施洗的約翰出來，在猶太的曠野宣講： \end{tabularx} \\ \\ \relax
3:2 & \begin{tabularx}{0.7\textwidth}{X} 「你們要悔改！因為天國近了。」 \end{tabularx} \\ \\ \relax
3:3 & \begin{tabularx}{0.7\textwidth}{X} 這人就是以賽亞先知所說的：「在曠野有聲音呼喊著：預備主的道，修直他的路。」 \end{tabularx} \\ \\
[1ex]
\hline
\hline
\end{longtable}
$^{1}$大家早晨.
2026年6月已經差不多去了一半.
半年很快.
2026年可以說是一個前所未有的一年對我來說.
因為2026年開始我又有第二次的退休.
23年前我是第一次.
所以2026年第二次是頗難過的.
但是接著退休下來.
今年頭三個月我遇到的事情.
我想也是前所未有的.
一直連著下來.
從一月到三月總共有六個人離開.
我的朋友.
我的親戚.
兩個我的中學一起畢業的同學.
到四月的時候還有老姐妹.
加起來剛剛有七個.
所以連著這麼多人離開的時候.
我的心有時候很難.
因為自己也是在退休年齡.
坐在那裡怎樣想呢.
老實說有時候也是有一些.
不知道怎樣.
當中有兩個同學.
我們同學是有一個群組的.
因為我們中學一起.
由頭讀起讀書讀到畢業.
所以我們有一個群組.
在WhatsApp群組上.
當這兩個同學離開的時候.
有很多人寫了一些節哀的話語在上面.
我看看想寫.
但是每個人都是寫那些東西.
基本上都是那幾個字.
他們去雲歸天國 安息主懷.
有時候會加多一句.
一路好走.
當我看完這幾個字.
這幾個字沒有什麼問題.
本身沒有什麼問題.

$^{41}$但問題是我們是一個生的人寫給一個死的人.
我不知道.
因為聖經也沒有這樣說.
我們是不是可以和死人溝通.
這是一個問題.
另外寫這幾句話的時候.
死者又不是信主的話.
我不知道他聽完這句話有什麼作用.
再加上寫的人.
我知道有些人也是未信主.
所以看完這幾個字的時候.
有時候不知道會怎樣.
跟別人說去天國或者去道主那裡.
在那個時候說的話.
似乎有些都不太妥.
有些都不知道怎樣.
這是我自己的感覺.
在二十年前.
我剛剛退休之後.
我上了馬太福音.
幾乎是第一課.
馬太福音是新約裡面第一本書.
我記得每次上課的時候.
老師要我們寫一篇不是很長的論文.
二十頁紙.
幾乎每一堂課都要寫.
通常那些題目都是有一些爭議.
或者有一些不同的解讀的題目.
在神學院裡面他們要我們寫.
寫的時候.
我們要交題目上去給老師看.
老師看完告訴我們可不可以寫.
還有一些題目不是那麼容易寫.
當時全班都交上去的時候.
老師就告訴我們.
結果發現大部分的人都是寫天國的題目.
這個是我們在神學院的學生寫的.
天國在馬太福音.
幾乎是一個專用的名詞.
在聖經裡面.

$^{81}$天國是在馬太福音裡面特別出現的一個字.
其他兩本福音書裡面不出現.
其他兩本福音書裡面.
兩本我不說第四本.
路加和馬可福音裡面用的都是神的國.
但是馬太福音是用天的國.
兩個字是一樣的.
因為天國在馬太福音裡面.
出現了三十二次或者三十三次.
我記得我們以前.
微信主持年輕的時候.
我們學了聽了天國這個字的時候.
我們就很喜歡用這個字來和信徒笑.
我們經常對基督徒嘲笑.
你們可以將來去天國.
信主他們可以去天國.
但事實上.
我們的心當時不是很善良的.
事實上是詛咒他們.
去天國是要死了才會去.
這是一個一樣.
所以當時我上馬太福音的時候.
我都不敢寫.
每個人都寫天國的事.
但我不敢寫.
因為有時候題目太大.
不是很容易.
因為當時如果你看馬太福音裡面.
耶穌都用了很多時間.
用比喻去解比要講天國的事情.
為什麼會不明白耶穌講的.
但是我們要想.
看到天國的時候.
我們怎樣去思考這個問題.
但我告訴你.
天國是一個很嚴肅的問題.
我們不可以好像我剛才說.
我們叫人去天國.
不可以亂這樣叫.
我很大膽問大家一個簡單的問題.

$^{121}$假如你今天死了之後.
你肯不肯定你是不是會去天國.
我只是一個差人來挑戰大家這個答案.
今天我想只是用一些時間.
大家回去我們再去看天國的事情.
因為這個和我們每一個人都息息相關.
我不知道大家記不記得.
上個禮拜的講道是講什麼.
你們不用舉手教我.
上個禮拜是講基督何等大的恩典.
再上一個禮拜.
不記得了 鄭牧師不在.
鄭牧師在Detroit跟我們說.
是講基督律.
再前一個禮拜.
Ken牧師在那裡經常見到.
我不記得那條題目.
那個禮拜是講耶穌升天.
實際上都是跟著我們一個主題.
但是有人走來.
好像是前個禮拜跟我說過.
他不是不記得.
聽不完上兩個禮拜的講道.
我不知道為什麼.
幾個禮拜都是講同樣的東西.
重複又重複.
但是有人會這樣說.
所以我說.
不如好好的我們回去看聖經裡面.
怎樣去教導我們.
當耶穌來的時候.
在祂來之前.
施洗約翰是祂的先鋒.
所以祂來的時候.
是一個很簡單的訊息.
在馬太福音三章二節裡說.
天國近要.
你們應該悔改.
我想我們今天守聖餐的時候.
都重複同樣一個題目.

$^{161}$後來耶穌來的時候.
祂亦都說一樣.
天國近要.
你們應當悔改.
所以當我們看天國的時候.
一件很特別的事情.
我們要看到就是.
我們要悔改.
這是耶穌說的.
我們今天時間不是很多.
我想從頭十五分鐘.
我想我們用一些很簡單的.
認識天國有三個很簡單的.
相關的一個組成元素.
這個不是.
實際上在我讀馬太福音的時候.
教導我們怎樣去明白天國.
第一個就是天國有一個王.
第二講到天國的子民.
第三講到天國的律法.
這是很簡單的三點.
一個王 一個子民 一個律法.
因為有一天我們去天國的時候.
我們要知道天國有一個王.
我們怎樣去天國.
如果我們是熟悉天國的話.
我們怎樣進入天國.
天國又有天國的律法.
這是三點.
我想從這三點我們去認識.
繼續去認識天國的事情.
天國的原文在希臘文的時候.
是不實的.
我不是很懂.
我不是很懂原文.
我可以抄出來.
我懂得字母.
意思就是神的統治.
在天上神的統治和權柄.
通常講到天的時候.

$^{201}$我們大家都知道天是神的代表.
猶太人不太敢叫神.
神的名字.
我知道我講了很多次神的名字.
所以馬太開始是用天國.
一個天的事情.
天是神的代表.
這裡我下面加了另外一個字.
和天國的事很相似.
天國是Bessalia.
下面那個我不是很懂.
後面的尾音是差一點點.
是Basilica.
有個K音.
Basilica是教堂.
特別是指建築物.
所以我們想到天國的時候.
我讀到教堂那個字.
我想大家可以聯想一下.
因為我們今天都在教堂裡面.
但我們今天不是特別要講教堂.
我們今天是注重天國.
在我們今天世代裡面.
我們不太喜歡黃.
如果聽到我們去到每個國家.
我們都不太喜歡黃的事情.
不太接受.
當然猶太人自己都是要立一個黃.
但是他們這個立黃不是神所喜歡的.
猶太人有四十幾個皇帝.
由蘇羅開始去到美國.
大概四十多個.
數不清楚.
跳來跳去很難數.
但是四十多個.
但是我們看完這些黃的最後的結局是怎樣.
這些是人的黃.
看回中國.
我們中國人是最多黃的.
我想我不去證明它.

$^{241}$我想大家知道.
因為我們有二千多年.
猶太人只有那幾百年.
中國人有五百.
數也是很難數.
五百多個黃.
隨時誰想做黃就殺了前面那個.
這是歷史上我們看到的.
但是很有趣的一件事.
中國最後兩個朝代之時.
當皇帝登位.
最後那兩次.
他們登位的時候.
開國的時候.
很喜歡說一句話.
他們下了一篇文書.
特別寫到.
奉天承運.
有時候我也不知道是怎樣.
但是我只認得前面兩個字.
奉天.
他們事實上是.
他們做皇帝之時.
就要用神的名字來讓他們做皇帝.
不是人們造反.
這是一個這樣的事情.
但是當我們看到天國的時候.
我們應該要想到一件事.
今天我們打開聖經的時候.
我們第一件事看到的是甚麼呢.
今天世界在天上的黃.
不是黃.
你叫他甚麼都好.
我們有一個神.
這個神是創造萬物一切的神.
我沒有辦法去證明.
但是我們聖經打開.
第一件事是告訴我們.
我們知道這個神.
不只是創造天.

$^{281}$不是創造地.
亦都是創造我們人.
雖然祂做人很簡單.
拿札幣吹口氣.
但是這個是聖經裡面告訴我們.
所以我們看到祂是一個全能的上帝.
祂是創造一切的上帝.
所以我們看天上的神的時候.
我們要去了解這件事.
所以當我們看到一件事的時候.
就是後來當神差耶穌基督來的時候.
耶穌基督是擁有一個神的形象.
神當初做人是照祂形象做.
但是我們看不到我們自己怎樣.
如果我們今天去照一下鏡子的時候.
我想我們都會頗令人害怕.
猶太人當日期待的一件事就是.
一個所謂的比賽.
是一個特別的人物出來.
去到他們當中拯救他們.
但是當他們看到是耶穌基督的時候.
他們不是很覺得耶穌基督是一個這樣的身份.
所以他們後來是拒絕.
第一個我們看到就是.
耶穌的門徒安德烈.
走去告訴他的兄弟知道.
我現在見到這個.
甚至後來施洗若翰自己都有少許懷疑.
叫人走去問耶穌.
究竟祂是不是神差來他們當中.
耶穌來祂是神的兒子.
祂是神.
我想我再加多一句.
所以祂是和神起初一起同在.
所以後來寫在歌羅西書的時候.
我們翻看一章.
很特別一樣.
保羅特別再提醒就是.
講到耶穌是神的愛子的時候.
是和神一起做.

$^{321}$特別提到一件事就是.
萬有去做.
因為耶穌是和神一起.
如果耶穌不是神的兒子.
祂只不過是另外一個.
所以當我們想到耶穌這個神的時候.
我們要想到祂是一個怎樣的神.
祂由開始.
這是我們要認識到.
所以保羅後面再加多一句.
在十八節的時候.
特別注重一件事就是教會.
因為耶穌是死而復生.
今天我們可以在祂的裡面.
第二個我們要講到天國的子民.
我想我們本來是神族.
但結果因為人犯了罪離開了.
我們離開了從亞當夏娃之後.
我想大家都知道.
我相信大家都不是在這裡土生土長.
作為一個國家的國民.
通常都是在那裡出生.
但如果我們不是出生的時候.
我們要申請也好.
或者要怎樣也好.
同樣做天國的子民的時候.
我想我們都是有同樣的一個.
這是一個對比.
我們本來是有罪.
我們不是去彌敦神的裡面.
但是如果我們有改變.
我們會轉.
所以耶穌的訊息就是.
你們應該悔改.
悔改的意思就是離開罪惡.
我們今天活在世界上就是活在罪裡面.
就算有時我們信了主.
但我們都要想一想.
我們今天有沒有活在罪裡面.
很多事情我們今天想.

$^{361}$我們是不是做很多事情.
是不是得罪神.
為什麼我們每個月要上首聖餐.
我們要紀念這件事.
這個我想很快就會具體.
有些其他的見證我想先跳過.
在彼得前書我們今天早上都提過.
我們要看到今天如果我們是被神揀選.
我們有一個不同的身份.
好像當天以色列人出埃及的時候.
但是以色列人沒有聽從神.
沒有轉向神.
他們的結果都是一條路.
但我們今天不是.
我們今天是有神.
我們得到拯救.
我們可以成為祂的子民.
所謂在以弗所說後面.
再加一件事就是.
我們今天同時有一個聖靈在我們當中.
我們在史達林傳看到日後.
當耶穌離開之後.
祂差遣了聖靈在我們當中.
所以我們今天紀念主餐的時候.
我們都紀念.
我們大家都是有一個聖靈.
將我們帶在一起.
第三是天外律法.
天外律法.
神制定這個律法.
由一開始.
由亞當夏娃開始.
祂給了他們一件事.
叫他們不要碰一些東西.
但是沒有聽到.
結果就離開了.
最後到了西乃山的時候.
神也給了摩西頒發律法給他們.
但是我們看到那個結果.
猶太人可不可以守律法.

$^{401}$沒有用.
結果仍然是不聽神的話.
那一代的人在曠野的死亡.
被我們看到.
這些圖畫沒有什麼問題.
所以在《羅馬書》那裡.
再一次讓我們知道.
我們今天沒有辦法.
因為在律法裡面.
我們可以行清義.
律法是叫人知罪.
叫人知道罪.
我們知道耶穌到最後告訴他們.
他們要重生.
所謂重生就是要迴轉.
他們不是可以憑著行律法.
就可以進入天國.
因為人沒有辦法可以守律法.
不是不可以.
他可以,但是他守不到.
很多人也走來問耶穌.
怎樣才可以進入.
耶穌給他們一個答案.
很特別.
他們最大的誡明是什麼呢.
是盡心盡義愛他的神.
剛才我們唱那首詩歌的時候.
也是特別提到.
耶穌最後在離開的時候.
再給他們加多一條律法.
我先跳了.
就是愛的律法.
我叫它愛的律法.
我們叫門徒彼此相愛.
這是最後的一條.
耶穌給他們的.
這條律法也是給我們今天.
我們今天做信徒的.
要明白來遵守.
剛才我們唱詩的時候也提到.

$^{441}$我們人自己不知道怎樣愛.
只有神愛.
感動我們.
耶穌上十字架.
為我們擔當的罪.
也戰勝死亡.
這是神給我們的恩典.
愛我們的一個結果.
今天因為我的時間不夠.
我們要好好地去認識天國有一個王.
這個王是怎樣.
整本聖經裡面也會告訴我們.
也要認識我們天國的子民是怎樣一回事.
我們如果是天國子民應該怎樣.
在《馬太福音》裡面耶穌用了很多時間講百福的事情.
講到天國子民有一個性格.
但是祂不是講什麼律法.
那條是天國子民應該有.
如果你進到天國就是這樣.
所以也講到天國的律法是怎樣.
天國就是神的國.
天國是代表神.
人今天是跟神分離.
我們的結局是死.
上一次陳太在這裡講了很多次.
關於死的.
但這個我們講到不單止是肉體的死.
也是以後跟神的遠遠的分離.
耶穌在地上很短時間.
三年的訊息就是天國近.
耶穌也帶了福音的消息.
因為祂到最後是上十字架.
讓人知道.
到最後我們能夠接受祂的人.
我們要悔改接受祂.
我們可以得到在天國的一個部分.
所以我們上到去的時候.
我們不需要怕究竟我們能不能進去.
這是一樣的.
所以我們到最後有一天.

$^{481}$我們會知道.
我們是跟基督在一起的時候.
我們在神的面前.
這個在地上.
我們可以一起在上面.
我們讚賞神.
我們會成為天國的子民.
今天大家坐在這裡.
這張論述箱不是你們的.
我想大家知道感受是什麼.
我希望藉著馬太福音的訊息.
這是我們應該知道我們要悔改.
今天雖然在地上.
但是我們要知道我們有一個永恆的盼望.
我們是可以回去.
回應第一次Ken在這裡說道的時候.
說到耶穌升起.
就是讓我們有一個永恆的盼望.
我們一起祈禱.
主啊,來到你給我們的恩典.
求你再一次加添力量給我們.
讓我們在我們今天行走天路的時候.
我們都能夠在你裡面得到力量.
我們願意同心.
我們都敬拜,尊敬你.
讓主你自己在我們心裡面做王,做主.
我們禱告,奉祭耶穌的名銜.
\newpage



\section{詩篇 78:1-8}
\label{sec:032280}
\textbf{信仰的接力}~~~~葉慧敏傳道
\newline
\newline
連結: \href{http://www.church.com.hk/acms/content.asp?site=cdc\&op=show\&type=product\&code=032280\&layout=sermon}{\texttt{032280}} ~~~~ 語音日期: 2026-06-07
\newline
\newline
\hyperref[sec:032279]{\small{< < < PREV SERMON < < <}}
~
\hyperref[sec:index_chronic]{\small{[返順時目]}}
~
\hyperref[sec:032281]{\small{> > > NEXT SERMON > > >}}
\newline
\newline
詩篇 78:1-8
\newline
\begin{longtable}{cl}
\hline
\hline
章節 & 經文 (和合本修訂版)\\
\hline
78:1 & \begin{tabularx}{0.7\textwidth}{X} 我的子民哪，要側耳聽我的訓誨，豎起耳朵聽我口中的言語。 \end{tabularx} \\ \\ \relax
78:2 & \begin{tabularx}{0.7\textwidth}{X} 我要開口說比喻，我要解開古時的謎語， \end{tabularx} \\ \\ \relax
78:3 & \begin{tabularx}{0.7\textwidth}{X} 是我們所聽見、所知道，我們的祖宗告訴我們的。 \end{tabularx} \\ \\ \relax
78:4 & \begin{tabularx}{0.7\textwidth}{X} 我們不要向子孫隱瞞這些事，而要將耶和華的美德和他的能力，並他所行的奇事，述說給後代聽。 \end{tabularx} \\ \\ \relax
78:5 & \begin{tabularx}{0.7\textwidth}{X} 他在雅各中立法度，在以色列中設律法；他吩咐我們的祖宗要傳給子孫， \end{tabularx} \\ \\ \relax
78:6 & \begin{tabularx}{0.7\textwidth}{X} 使將要生的後代子孫可以曉得。他們也要起來告訴他們的子孫， \end{tabularx} \\ \\ \relax
78:7 & \begin{tabularx}{0.7\textwidth}{X} 好讓他們仰望神，不忘記神的作為，惟遵守他的命令； \end{tabularx} \\ \\ \relax
78:8 & \begin{tabularx}{0.7\textwidth}{X} 不要像他們的祖宗，是頑梗悖逆、心不堅定，向神心不忠實之輩。 \end{tabularx} \\ \\
[1ex]
\hline
\hline
\end{longtable}
$^{1}$弟兄姊妹五芒.
今日我給她的講道題目叫做.
信仰的接力.
不知你看到這個題目時你有甚麼感覺.
昨晚臨睡前我兒子問我.
「明天你講道嗎?」.
我說「是呀」.
「你講甚麼?」.
我說「信仰的接力」.
然後他說.
「哦!即是耶穌講給他的門徒聽」.
「他的門徒又講給下一批門徒聽」.
「行了!講完」.
他都講到我想講的其中一些意思.
但今日我都很想.
我們可以再透過一篇詩篇一起思想.
我想問大家.
如果有一天你的小朋友.
你的下一代.
或者你說.
「我其實未有小朋友」.
「或者我未結婚」.
如果一些比你小的人跟你講.
「其實我不是很知神是誰」.
他一直在教會.
但他問一個這樣的問題.
你會有甚麼感受呢?.
是否都會有少少覺得心酸,難過?.
弟兄姊妹,這不是一個理論的問題.
而是一個很真實,很貼身的問題.
我們今日很多時為下一代做決定.
當初移民來英國.
我們去想住哪裡.
教育是怎樣.
或者未來怎樣對下一代最好.
完全是出於愛.
因為我們想給予他們一個.
最好而我們負擔得到的環境.
但有一件事.
我們有沒有同時都問.

$^{41}$我怎樣在這些新環境裡.
將這個信仰傳給我的下一代呢?.
因為來到英國之後.
我們會發現.
信仰不再是自然會發生的事.
在香港我們可能有一個很自然的環境.
小朋友在媽媽肚裡的時候.
已經回到那間教會.
接著出生之後.
每個星期都在那個地方走來走去.
教會裡面的人又認識你.
在教會,在群體,在語言裡面.
都幫我們下一代托住信仰.
但來到英國之後.
所有東西都不同了.
教育上他們變了是英文學習.
回教會.
我們現在是在講廣東話.
但他們接收已經可能只能得一半.
或者一半多一點.
信仰開始變成一件.
需要刻意維繫的事.
有一件事是真實發生中.
信仰可以在一代之間慢慢消失.
不是因為有人去反對神.
而是因為慢慢沒有人再傳.
士師記二章十節這樣說.
後來有另一代人興起.
「不知道耶和華,也不知道他為以色列所行的事」.
他們不是離棄神.
而是他們不再認識神.
因為沒有人再提起.
我想我們都不想見到.
士師記這樣的情況成為我們下一代的故事.
今天我想和大家透過詩篇七十八篇.
一篇講信仰接力的詩篇.
一起去探討.
詩人阿薩用這篇詩篇.
講兩個問題.
第一個問題是信仰怎樣傳下去.

$^{81}$第二個問題是如果不傳會發生什麼事.
求主幫我們一起祈禱.
讓我們一起去學習.
主我求你去幫助.
讓今天無論聽的講的.
我們一起去堅持.
也一起去學習.
怎樣將我們對你那份信仰.
你那份愛可以傳遞下去.
主求你去眷顧引領.
我們將到禱告交託.
奉主明求 阿們.
詩篇七十八篇是詩篇裡面第二長的一篇.
它是一篇訓誨詩.
訓誨詩的原文意思是.
帶來動見使人明白的詩歌.
是用來教導的詩歌.
不是用來唱的.
作者是阿薩.
他是大衛時代的敬拜帶領者.
是一位先知性的詩人.
他不是旁觀者.
而是一位深深認識神歷史作為的人.
他選擇用詩歌來記載歷史.
因為詩歌可以記住.
可以傳下去.
可以一代一代這樣唱.
整篇的詩篇七十八篇.
回顧了以色列人的歷史.
呈現了一個不斷重覆的循環.
神施恩給以色列民.
人叛逆 上帝管教.
人悔罪 迴轉.
神再次的憐憫.
給我們有恩典.
這個讓我們看見就是人是會失敗.
但是神是沒有放棄.
而一至八節.
是這篇詩篇的引言.
是阿薩堆信仰傳承.

$^{121}$所發出的一個呼召.
我們一起看看.
這個呼召對我們有甚麼提醒.
第一點.
這個呼召是提醒我們.
要說 要說出來.
傳承是需要開口.
我們一起讀這四節經文.
一二三 我的民啊.
你們要留心聽我的訓誨.
則以聽我後宗的話.
我要開口說比喻.
我要說出古時的謎語.
是我們所聽見 所知道的.
也是我們的祖宗告訴我們的.
我們不將這些事.
向他們的子孫隱瞞.
要將耶和華的美德和他的能力.
並他奇妙的作為.
述說給後代人.
第一節可以說是一個牧者的呼召.
詩人一開始.
稱呼這一班的以色列民是我的民.
不是一個冷冰冰的開場.
而是一種很深的連結的稱呼.
是一個被神託付的人.
看著自己最掛心的人去說話.
則以聽 留心聽.
連續用兩個動詞.
帶著一種迫切感.
意思是說 你一定要聽.
因為這件事不可以忽略.
他其實是說一句很重的說話.
你們一定要聽.
因為這件事.
關乎你們和你們的下一代.
第二節他說我要開口說比喻.
比喻愚民的意思不是說普通的故事.
而是一種有目的 有結構.
可以記住 可以傳遞的智慧說話.

$^{161}$《箴言整卷書》就是用比喻這個方式去表達.
阿發說我要說.
意思是我不是隨便說.
我是刻意用可以傳下去的方式去說.
我要說出古時的謎語.
謎語的意思是.
需要人停下來思考.
追問意義的說話.
不是叫人混亂.
而是叫人開始問.
為什麼會這樣.
神想我們明白什麼.
而信仰就是在這種提問開始.
傳承不是看我們說得有多好.
而是看我們有沒有開口說.
去到第三節經文.
好像一個接力棒的畫面.
是我們所聽見所知道.
也是我們的祖宗告訴我們.
阿發說的信仰不是他自己創造出來.
而是他接收回來的.
他先聽見 先知道.
然後再傳出去.
原文用一種群體的語氣.
我們所聽見 我們所知道.
即是說這個不是個人的信念.
而是一個群體一直盛載著的記憶.
這裡已經出現一個很重要的畫面.
信仰好像一條接力鏈.
有人傳給阿發.
阿發再傳給下一代.
他不單止是接收者.
而是傳遞者.
早幾天 當我預備這篇講道的時候.
我收到一個訊息.
是我一個中學的老師.
他跟我分享生活.
有問候我的一個訊息.
我在那個時候.
對天父說我真的很感動.

$^{201}$其實畢業了幾十年.
我和他仍然可以維繫著這個關係.
他就是當日挑戰我信仰.
希望我可以在畢業之前.
清楚做一個決定.
我是否在信耶穌的一個老師.
今天我們坐在這裡其實是一樣.
我們都是因為有人曾經將這一棒交給我們.
去鼓勵我們信耶穌.
去將耶穌分享給我們.
而今天我們都站在同一個位置.
我們是接收者.
同時我們是傳遞者.
這一棒輪到我們要交出去.
第四節他這樣說.
我們不將這些事向他們的子孫隱瞞.
隱瞞的意思不是不講這麼簡單.
而是令某些東西逐漸消失.
變得不再存在.
阿薩提這樣說.
我們不會讓信仰消失.
但反過來的是.
如果我們不講.
信仰會慢慢消失.
沉默從來不是中立.
沉默是選擇放棄這個成全.
阿薩所說要傳下去是甚麼.
經文繼續說.
要將耶和華的美德和他的能力.
並他奇妙的作為述說給後代聽.
三樣很具體的事.
神的美德,神的能力.
和他奇妙的作為.
不是抽象神學.
而是可以用第一身講出來的故事.
弟兄姊妹你和我都有我們自己的故事.
正如我剛才分享.
我怎樣去信耶穌.
就是因為有一位老師.
在我臨畢業之前.

$^{241}$很記掛我的信仰.
以致他約我去傾.
他問了我幾個問題.
我到今天仍然很深刻.
他當時問我.
你信不信這個世界有一個神創造出來.
我說我信.
你信不信人有罪.
我都信.
你信不信上帝預備了一個獨身指揮我們的罪.
寫在十字架上.
我都信.
你信不信耶穌復活.
我都信.
他問到第五個問題.
我就說我未信.
他問我.
問什麼?弟兄姊妹.
我們通常問別人信不信耶穌.
都是問這幾個問題.
第五個問題就是.
你願不願意相信耶穌是你的救主.
和你的生命的主.
你願不願意接受.
我在那個時候跟他說.
我就不是很願意接受.
我未覺得祂是我生命的主.
但是他就跟我說.
你就繼續去想吧.
這個問題我繼續讓你回家去思考.
如果你真的離開這間中學.
你人生裡面還有很長的年日去思考.
我後來就因為一些經歷.
想起他給我的挑戰.
我就相信了主.
同樣我們來英國.
其實都有我們的故事.
其實我們每一個人都可以.
將一個關於神.
怎樣帶領.

$^{281}$怎樣供應.
怎樣同行的見證.
和我們的下一代分享.
特別是和我們的子女分享.
為什麼我們可以來到這裡.
上帝怎樣去引領我們.
如果我們不說.
下一代會視為一個理所當然的事.
他們不會知道.
他們也不會知道.
上帝參與在其中.
信仰不會自然傳承.
只會慢慢消失.
我也知道在在有些父母.
你們的子女已經長大了.
可能已經不是天天坐在你身邊.
甚至他們已經有自己的生活.
我兒子最近就跟我說.
你一天問我三個問題就好了.
為什麼只可以問三個問題.
問得多覺得你煩.
我說那好吧.
你也問我三個問題就好了.
他就說不是的.
其實不是的.
你不是問我三個問題就好.
你問我多少個問題都好.
你快點問我吧.
我就說不要.
你覺得我煩.
我們的子女越來越大的時候.
可能都會覺得.
你不要說那麼多.
我們也許已經錯過了.
教導的黃金時間.
但是傳承從來不單只是一種方式.
對年幼的小朋友來說.
我們可以用說話來引導他們.
對已經長大的小朋友來說.
更重要的是用我們的生命見證.

$^{321}$他們可能不再即時回應.
我們所說的話.
但是他們仍然在看著.
我們怎樣去生活.
他們仍然在看著.
我們可以做的有兩件事.
第一件就是.
我們繼續活出我們真實的信仰.
不是表面而是真實.
讓我們的下一代看見.
信仰不單只是星期天在教會裡面發生.
而是每一天在我們人生裡面.
最真實的支撐著我們過每一天.
第二件事是持續為他們祈禱.
當我們沒有辦法再控制.
他們選擇上帝還是放棄上帝的時候.
我們仍然可以將他們交給神.
這是一個更深的信靠.
阿薩寫詩的時候.
他都不知道結果何時會出現.
但是他仍然寫.
因為他相信上帝會使用.
我們每一個禱告.
每一個真實活出信仰的時刻.
都是我們放在我們的小朋友.
我們的下一代.
生命裡面的種子.
種子不是在我們手裡面有結果.
但是種子一定要有人殺下去.
才會有機會有結果.
第二個呼召是.
詩人說要對.
我們要對甚麼呢?.
傳承的目的是仰望神.
第五節他這樣說.
因為他在雅各宗立法道.
在以色列宗設律法.
是他吩咐我們祖宗要傳給子孫的.
立法道設律法.
這個動詞的意思是建立.

$^{361}$令到他更加鞏固的意思.
不是臨時安排.
是神主動刻意建立的傳承機制.
這個和申命記錄.
將七節一節我們很熟悉的經文.
其實是一脈相承.
也要殷勤教訓你的兒女.
無論你坐在家裡 行在路上 躺下.
起來都要談論.
希伯來的文化裡面.
信仰教育是滲透在日常生活裡面.
而不是局限在某個場合當中.
而第六節他說.
使將要新的後代子孫知道.
他們也要起來告訴他們的子孫.
這一節我們看見有一個結構.
它不單止是一個一次的傳承.
而是一個連鎖式的傳承.
起來這個意思和第五節裡面說.
設立是同一個字根.
神設立了這個制度.
而人的回應是要主動站起來承擔.
神起來 人都起來.
一個不是巧合.
而是阿撒他刻意的描寫的方式.
其實阿撒描繪的圖畫就是.
祖宗傳給父母.
父母告訴兒女.
告訴我們.
我們告訴子女.
子女告訴孫兒.
每一代都是接力的一棒.
每一代都有責任主動的起來.
將這一棒交出去.
信仰的傳承不是某一代人的責任.
而是每一代人都有一份使命.
阿撒寫這篇詩的時候.
他自己都起來.
用詩歌的形式.
將神的作為傳給下一代.

$^{401}$對我們移民的家庭來說.
更加真實.
我們這一代正正站在一個很關鍵的位置.
銜接著香港的信仰根基.
和英國成長下一代之間.
這一棒交得好.
下一代才能夠有可能起來.
再傳給他們的子女.
這一棒如果斷了.
在我們這一刻.
不單止是我們的下一代失去.
而是整條鏈都斷了.
去到第七節.
詩人用了三個動詞.
他順序地表達.
我們一起讀讀第七節經文.
1,2,3.
好叫他們仰望上帝.
不忘記上帝的作為.
為要守他的命令.
三個動詞是順序刻意地寫出來.
首先要仰望上帝.
接著是一種信靠和關係的結合.
接著是紀念神的作為.
記住一些歷史和感恩.
最後是遵守神的命令.
行為加信服.
如果次序亂了.
信仰就會變成一個重擔而不是一個關係.
很多信仰教育的問題.
是用了第三步開始.
你要回教會,你要讀聖經,你要祈禱,你要乖.
但是沒有了第一步的仰望.
第三步的遵守.
就會變成律法主義.
或者只是一個家庭的文化習慣.
當文化改變了.
習慣自然消失.
但是關係不會.
仰望這個字的意思是.

$^{441}$將心放在神身上.
這是一個主動持續的轉向.
不是被動接受.
這是我們要在孩子心裡建立的.
一個主動尋求神的心.
不單止是一套宗教行為.
弟兄姊妹先關係後行為.
先仰望後遵守.
這裡有一個很真實的問題.
我們都要面對.
在我們的下一代裡.
無論是你的子女.
或者是教會的小朋友.
或者是教會的青少年.
直情,教會是甚麼?.
是一個群體?.
是一個家?.
還是只是眾多課外活動中的其中一項呢?.
游泳班,鋼琴班,足球的練習.
或者看一場戲.
教會排在第幾位呢?.
來到英國之後.
我兒子都經常收到生日會的邀請.
通常生日會是在星期六,日進行.
我們的原則.
我們要跟他分享.
我們要敬拜主.
所以我們是有一些原則.
生日會如果是在崇拜前.
崇拜開始前半小時結束.
我們會參加的.
或者生日會如果是在星期日發生.
生日會在崇拜開始前半小時結束.
我們會參加的.
或者生日會在崇拜結束之後才開始.
我們都可以參加.
爸爸辛苦一點.
要接送他們.
但是我們一直跟他分享.
一個很重要的訊息.

$^{481}$我們不要停止聚會.
我們解釋給他聽為何要回教會.
為何要敬拜神.
其實我講不講.
他明不明.
他從小到大都在教會生活.
從小到大我們都是這樣的模式.
去敬拜主.
我們的小朋友其實已經知道答案.
他從來都沒有跟我們掙扎過.
去到就去,去不到就跟我說去不到.
因為小朋友不是用耳朵去學.
他是用眼睛去看.
當信仰變成生活的其中一項.
我們下一代學到的結論只有一個.
上帝其實不是那麼重要.
第三個呼召是要活.
傳承是靠生命不是說教.
第八節他這樣說.
不要像他們的祖宗是還敢叛逆.
居心不正之輩.
向著上帝心不誠實.
詩人列出了四個描述層層遞進.
非常值得我們仔細去看.
新譯本更加平衡.
免得像他們的列祖成為還敢叛逆的一代.
心不堅定的一代.
他們的心對神不忠.
還敢的意思是固執拒絕被塑造.
不是說他這樣說.
而是說他的心眼不願被改變.
不願被神的話去承認.
叛逆是主動的抗拒.
不單是被動.
而是主動和神的方向對抗.
心志不堅定.
心不穩固.
意思是心沒被建立起來.
沒根基.
這不是道德的判斷.

$^{521}$而是靈命狀態的描述.
一個在信仰上沒根基的人.
對神不忠.
不忠信不可靠.
對神的承諾沒有持久性.
這四個描述不是描述一個大奸大惡的人.
而是描述一個漸漸偏離上帝的人.
不是突然的背棄.
是慢慢的冷淡.
慢慢的疏遠.
慢慢的心不在焉.
信仰的斷層.
大多數都不是突然性發生.
而是漸進的過程.
第八節他說「心」這個字出現了兩次.
在希伯來的人觀裡.
心不僅是情感.
而是意志,思想,品格的核心.
整個人的生命方向由心決定.
阿薩說問題不在於外在行為.
而是生命的方向.
我們的生命是不是朝向神.
這對我們作為基督徒父母一個很深的提醒.
或者我們面對著我們下一代.
我們都是一班成年人.
我們面對著下一代.
一個很深的提醒.
孩子最在意的不是你教他甚麼.
而是你活著甚麼.
我有一位神學院的同學.
他跟我一樣都是移民來英國.
但他移民的路是崎嶇很多.
他來到之後他沒有目會.
但他仍然參與一些教會的邀請.
邀請他去港獨.
他有一位今年考GCSE的兒子.
在這個兒子成長的過程裡.
他親眼看見爸爸在事奉上所經歷的挑戰和困難.
我的朋友他沒有刻意的隱瞞.
沒有裝假.

$^{561}$沒有告訴他教會就是很美好.
機構就是很美好.
基督教就是非常美好.
他的兒子看見有些問題.
但更重要的是.
他看見他爸爸怎樣面對這一切的問題.
怎樣去形現經歷上帝的愛.
經歷上帝的恩典.
就是因為這樣.
每次當我朋友有愛出港獨的時候.
他的兒子就會說.
我跟著你一起去聽.
我想聽你講道.
不是因為他沒有事做.
其實下一代很想在家裡.
叫他不出街其實完全沒有問題.
但他說我很想跟你一起去做.
我問我朋友是不是他教會生活不開心.
不是.
他教會生活很開心.
但這個兒子就是很喜歡聽爸爸講道.
他之所以這麼喜歡.
是因為他從他爸爸生命裡面.
他看見一位真實的神.
弟兄姊妹今天我們的下一代.
是否同樣看見.
上帝在我們生命裡面的真實.
求主去幫我們.
阿薩講完一至八節之後.
好像有一些原則.
他跟著在九至七十二節.
他就用了整段的歷史.
以色列人的歷史回答一個問題.
如果不傳承會發生甚麼事.
但由於時間關係.
我不能夠逐節跟大家看.
大家如果有興趣.
可以回家看回九至七十二節.
我嘗試用四個畫面.
去分享以色列人的一些縮影.

$^{601}$他們的歷史.
信仰斷層的四個警告.
第一個畫面.
忘記神.
九至十一節講的是.
以法蓮的後代.
以法蓮是以色列十二支派其中之一.
是一個強大的支派.
有裝備,有能力,有資源.
但是關鍵時候.
他們轉身退後.
為甚麼呢.
經文說因為他們忘記上帝所行的事.
信仰斷層.
通常不是公開背叛作為開始.
而是從忘記開始.
忘記上帝的恩典.
忘記神的供應.
忘記上帝的作為.
當記憶消失.
忠心亦都會流逝.
信仰斷裂從記憶開始.
第二個畫面.
是以色列人在曠野中的生活.
記載在十二至四十一節.
以色列人出埃及親眼看見紅海分開.
親身經歷上帝行的神蹟.
但是在曠野中每次遇到困難.
他們第一個反應不是祈禱.
不是信靠.
而是埋怨.
詩篇七十八篇十八節.
他們在曠野任意試探上帝.
在荒地求他.
按自己的心願賜食物.
他們知道上帝的能力.
卻沒有仰望上帝的心意.
他們只要上帝為他們服務.
卻沒有以上帝為主.
這是一個很真實的屬靈狀態.

$^{641}$知道關於上帝的事.
但是沒有真正的仰望關係.
信仰變成了功能性的工具.
我需要的時候才會想起你.
沒有仰望只剩要求.
第三個畫面是一個持續的叛逆.
以色列人進入應許地之後.
開始安定下來.
但是安定之後反而更加容易忘記神.
他們在應許的地方仍然拜偶像.
跟隨周圍的文化形式上的宗教還在.
但心已經離開.
在五十六至六十四節裡面記載.
這個畫面同樣都有可能發生.
我們的孩子仍然有宗教的標籤.
我是基督徒.
我是信二代.
我是信三代.
我在媽媽的肚裡已經在回教會.
但是這些標籤的背後.
有沒有真實的神.
有沒有真實的認識神呢.
沒有真實的信教只是一個文化身份.
宗教存在.
但是神已經被遺忘.
第四個畫面是神興起了大衛.
詩篇七十八篇不是以絕望作結束.
上帝棄掉了二百年.
揀選了猶大.
棄掉了皇族.
從陽圈呼召了大衛.
上帝的計劃沒有失敗.
但是那些叛逆的世代已經被上帝越過了.
上帝繼續前行.
但是他們不在其中.
上帝仍然掌權.
但是那一代錯過了.
信仰一旦不被刻意傳承.
就會被世界取代.
來到這裡.

$^{681}$可能你都會有一種感覺.
我們未必做得完美.
甚至有些地方好像已經錯過了.
阿薩用整篇詩篇告訴我們.
上帝最後的動作不是要審判.
而是要牧養.
他一次又一次被人辜負.
但是上帝仍然憐憫.
他興起了大衛.
不是因為以色列配得.
而是因為神的憐憫沒有盡頭.
詩篇七十八篇最後一節.
七十二節這樣說.
「於是他按心中的純正牧養他們.
用手中的巧妙引導他們.
他呼召大衛.
去用心中的純正牧養子民.
用手中的巧妙引導他們.
純正是神的心意.
是真實沒有扭曲.
巧妙是神知道怎樣帶領我們走每一步.
就算我們不懂得.
我們不是完美的父母.
不是完美的長輩.
我們都不是總是純正.
不是總是巧妙.
我們有時會說錯說話.
有時會做錯決定.
有時會讓我們的孩子.
看到我們的軟弱或者失敗.
但是神不同.
他的牧養不受我們的限制.
他的引導不因為我們的失敗而停止.
所以弟兄姊妹今天的訊息.
不是要用來叫我們感覺到罪疚.
而是用來呼召我們.
喚醒我們.
上帝仍然在牧養我們每一位.
我們只是被祂邀請成為這一棒接力者.
上帝邀請我們參與.

$^{721}$但是這支棒不會自己傳下去.
需要有人新出手.
信仰都一樣.
如果我們不傳.
世界會教.
我們的下一代.
每一天在學校.
在社交媒體.
在朋友當中.
已經有人告訴他.
什麼是重要.
如果家庭和教會是沉默.
他聽到的只是一種聲音.
今天我也想邀請你.
去思想一個問題.
我們可以做的一步是什麼呢?.
在家裡.
也許去分享上帝在你生命裡的一個見證.
你怎樣去經歷上帝.
你分享你的事.
他就算不相信.
他也會聽的.
也許在家裡.
邀請我們的子女.
做一個簡單的家庭祈禱.
在教會裡.
我們面對著我們的兒童主日學.
面對著我們的青少年.
面對著大專生.
即是他們都很需要有人去牧養.
甚至我們自己都很需要被牧養.
會不會參與在裡面.
去成為他們的一個同行者呢?.
一個小小的行動.
可能成為下一代.
信仰的轉捩點.
懸疑是詩記義.
詩記義將十節不會成為我們的故事.
讓我們將這一棒傳下去.
我們一起祈禱.

$^{761}$很想在這個時候.
我們都去有一段安靜的時間.
如果你有孩子的時候.
很想你先為自己祈禱.
同時也為你的孩子祈禱.
求主在裡面.
仍然在你們裡面當中.
求主讓我們自己.
讓我們的孩子能夠成為.
真的願意以耶穌作為我們的救主.
我們生命的主那一位.
正如當年.
我的一位很疼愛我的中學老師.
他那份緊張.
那份很擔心我離開了學校.
就未必再可以信耶穌那份心.
我們先為我們自己.
和我們下一代圖告.
天父上帝我們多謝你.
多謝你讓我們看見.
是我們人真的很有限.
我們真的不是完美的人.
我們都會有軟弱有出錯的時候.
但我們感恩.
上帝你一直沒有放棄我們.
我們可能很多次令你傷心難過.
但你的憐憫你的恩典.
仍然與我們同在.
上帝我求你.
你就去幫助.
幫助我們能夠看見.
我們領受了你這份救恩.
領受了你是我們生命的大靈.
你的改變之後.
我們怎樣去傳承給我們下一代的肢體.
無論是我們教會的弟兄姊妹.
無論是我們的兒女.
上主很需要你不斷的加力在裡面.
我們很需要有著你給我們的巧妙.
你給我們的智慧.

$^{801}$你給我們那份愛去傳承下去.
我求主你去鼓勵我們.
我們過去可能真的有甩漏.
上帝求你填補我們這些甩漏.
我們今天開始重新去建立.
我們在家庭裡那份屬靈的氣氛.
我們重新建立在教會裡那份屬靈的氣氛.
主我求你.
你就與我們同在.
幫助我們加力給我們每一位.
讓我們每一位弟兄姊妹都看見.
我們接了這支棒.
需要去交給下一梯隊.
主求你去眷顧人民.
我們多謝你.
我們將到禱告交托.
奉主明求.
阿們.
\newpage



\section{出埃及記 20:1-5}
\label{sec:032281}
\textbf{當A.I.成為偶像}~~~~招世超牧師
\newline
\newline
連結: \href{http://www.church.com.hk/acms/content.asp?site=cdc\&op=show\&type=product\&code=032281\&layout=sermon}{\texttt{032281}} ~~~~ 語音日期: 2026-06-07
\newline
\newline
\hyperref[sec:032280]{\small{< < < PREV SERMON < < <}}
~
\hyperref[sec:index_chronic]{\small{[返順時目]}}
~
\hyperref[sec:032282]{\small{> > > NEXT SERMON > > >}}
\newline
\newline
出埃及記 20:1-5
\newline
\begin{longtable}{cl}
\hline
\hline
章節 & 經文 (和合本修訂版)\\
\hline
20:1 & \begin{tabularx}{0.7\textwidth}{X} 神吩咐這一切的話，說： \end{tabularx} \\ \\ \relax
20:2 & \begin{tabularx}{0.7\textwidth}{X} 「我是耶和華－你的神，曾將你從埃及地為奴之家領出來。 \end{tabularx} \\ \\ \relax
20:3 & \begin{tabularx}{0.7\textwidth}{X} 「除了我以外，你不可有別的神。 \end{tabularx} \\ \\ \relax
20:4 & \begin{tabularx}{0.7\textwidth}{X} 「不可為自己雕刻偶像，也不可做甚麼形像，彷彿上天、下地和地底下水中的百物。 \end{tabularx} \\ \\ \relax
20:5 & \begin{tabularx}{0.7\textwidth}{X} 不可跪拜那些像，也不可事奉它們，因為我耶和華---你的神是忌邪的神。恨我的，我必懲罰他們的罪，自父及子，直到三、四代； \end{tabularx} \\ \\
[1ex]
\hline
\hline
\end{longtable}
$^{1}$我不知道大家的禱告生活是怎麼樣的.
通常平時我們很少禱告.
除了借飯祈禱 或者睡覺前有沒有祈禱.
我希望有啦.
否則什麼時候最容易祈禱呢.
就是當我們遇到問題的時候.
通常有問題 有困難.
我們就會祈禱.
這些是基督徒應該有的習慣.
但我發覺現在的世代變得太快了.
在科技發展的這個時間當中.
現在的人有問題 有難處.
以前就Google一下.
現在不Google了.
其實不是不Google的.
不過現在是什麼.
去問CHPT.
Germany.
或者是Copilot.
總之現在人工智能.
我們發覺到一個地步.
好像橫行無忌.
其實不是橫行無忌.
是我們使用它.
使用到很多時候沒有分寸.
所以很多時候我們發覺.
在不知不覺當中.
不知道我們會不會有些人.
有這樣的感覺.
覺得現在AI.
人工智能這樣東西.
好像是無所不知.
否則你就不會去找它.
有什麼問題.
它都可以給你答案.
有什麼需要.
它都會給很多的建議.
你可以一二三四五.
怎樣去做.
不知道會不會在不知不覺當中.

$^{41}$我們失去了對上帝.
那種單純敬畏和信服的心.
AI 人工智能.
大家都知道.
是人製造出來的東西.
但是現在我們發覺.
這個被人製造出來的東西.
應該是成為人的工具.
但是不知道在哪個時間.
可能我們都不Aware.
不警覺到.
是取代了神的位置.
甚至成為某些人的偶像.
所以有什麼問題.
有什麼事.
首先就去查看TVT.
去問一下AI.
到底是什麼一回事.
和我們以前的基督徒生活習慣.
是不是好像不同了.
另外一個人工智慧.
帶給我們影響的是什麼.
不知道會不會我們在不知不覺當中.
扭曲了對人生命價值.
的看法.
特別是我們中國人.
很強調品格.
基督徒很著重.
屬靈生命.
我們的靈性的情況是怎麼樣.
但是現在.
在過了數碼時代.
現在人工智能的時代.
我們的情況變成怎樣.
看數據.
譬如有沒有Facebook.
你在Facebook裡有多少個朋友.
這個很重要.
有沒有一些社交媒體的帳戶.
人工.

$^{81}$不是 不是.
社交媒體的帳戶.
有多少個跟隨者.
有沒有.
有沒有去想到這些問題.
還有現在很多人自己自拍.
做很多的影片.
TikTok也好.
或者Instagram也好.
或者很多時候放上Youtube.
那你很緊張.
有多少人看.
有多少人看我放上Youtube的影片.
還有多少人按讚.
說很棒.
很多時候.
會不會像我剛才所說.
不知不覺當中.
從前我們強調品格.
強調內在的靈性.
現在以數據作為人價值的一個指標.
於是很多時候.
我們就將我們的生命.
數據化了.
人在不知不覺當中.
很自然地.
落在一些無止境的網絡的焦慮.
有網絡焦慮.
或者怎麼說.
一個無止境的用數據.
去互相比較的一個環境當中.
而當我們落在這樣的處境.
上網上到真的瘋了.
好像一不上網就不自在.
然後又去比較我們的數據.
當我們這樣做的時候.
帶來的結果是人際關係.
和所謂人所需要的群體生活.
就會被一個虛擬的環境.
去隔離了我們.

$^{121}$甚至是取代了.
這就是另外一個AI帶來的情況.
我們大家都回到教會.
你不是上網去崇拜.
現在很多人就上慣了網.
不回到實體去敬拜.
就在家裡穿著睡衣也好.
在吃東西也好.
坐在那裡就說敬拜上帝.
那個到底是不是敬拜呢.
《聖經》裡面很強調.
基督徒之間要有群體的生活.
我們要彼此相愛.
這些互動.
我們現在是上網不能夠得到的.
但是智能手機.
人工智慧製造出來的虛擬環境.
再加上現在社交媒體的即時傳達訊息.
幾秒鐘就到你手上了.
就覺得這個很厲害.
我覺得自己無所不在.
雖然我現在站在這裡.
不過我去到香港 去到加拿大 去到哪裡也好.
有些東西傳過來 我馬上就看到了.
相片也看到了.
影片也看到了.
於是在現實的生活當中就缺席了.
所以我今天有這麼多空位在這裡.
這個不知道是不是其中一個原因.
I don't know.
這是我們面對一個AI時代當中.
我們需要面對的其中一個問題.
而另外一個比較嚴重的.
就是剛才我稍微提過的.
很容易我們會沉溺上網.
會沉溺在那個虛擬的世界當中.
特別是AI出來之後.
那些景象是越來越真實的.
以前是看卡通片.
現在不是了.

$^{161}$好像真人一樣.
看上去好像是一個真實的景況.
我們不知不覺被拉進去一個虛擬的世界.
虛擬的環境當中.
而在教會生活當中.
我是說愛的傳遞.
我是說要包容 要服侍 要牧養.
但是不知不覺當中.
就被那些很冰冷的資訊交易所取代.
今早我回家祈禱的另一件事.
我很感恩的是什麼呢.
有一個姊妹帶了一張白蘭花出來.
整個房間都香了.
你看著電腦或手機.
你就聞不到那些香味了.
你不會受到那種環境氣氛的感染.
沒有了.
但是現在很多時候我們享受.
我們沉醉在那些虛擬的環境當中.
為什麼呢.
那些虛擬的環境好像天馬行空.
你想什麼 想什麼.
你吩咐一下那些AI.
text就變成一個畫像.
甚至一個video就出來了.
我們沒有了那種.
怎麼說呢.
靈裡面的沉澱.
或者是安靜.
當是這樣的時候你會發覺.
似乎我們越來越難聽到.
我在祈禱當中要聽見.
上帝那種微小的聲音.
有沒有聽到呢.
當我們資訊成為一種.
好像毒癮一樣.
資訊成癮的時候.
就會取代了.
甚至是搶走了.
我們跟神個別親近的時間.

$^{201}$這是一個很嚴重的問題.
大家其實要記住.
今天怎麼跟大家說這個問題呢.
因為我們發覺現在AI這個世代.
對我們的生活的影響.
是越來越重要.
我們知道其實我剛才提過.
其實人工智慧我們可以怎麼說呢.
是上帝賜給人管理世界的智慧的結晶.
應該是我們的工具.
但當這些工具承擔了過多我們的安全感.
有什麼不妥當的時候就去問AI.
承擔了過多的成就感.
或者甚至是承擔了過多.
我們能夠解決問題.
得到拯救的盼望.
人工智慧到了那個地步.
就會成為我們心中的偶像.
我們將人手所創造的東西.
捧上了神壇.
大家覺不覺得呢.
我覺得 否則不會說這個題目.
我們的生命其實當我們將一樣東西.
捧上我們心中的神壇的時候.
我們就離開上帝 越來越遠.
屬靈的生命開始倒退 開始墮落.
所以諸大弟兄姊妹.
求神幫助我們.
我就去形容現在AI的人工智能.
加上大數據 沒有大數據就沒有那麼厲害.
這個是AI人工智慧橫行的一個世代.
我們怎樣能夠跟隨耶穌基督呢.
這是我們這次講道系列的主題.
我們怎樣能夠在這樣的環境下.
仍然可以警覺得到.
提醒我們自己.
神真是存在 神真是很偉大.
首先我會建議大家有一個曠野時刻.
我們看馬可福音第一章32到35節.
這裡講耶穌基督.

$^{241}$應該是我們學習的榜樣.
《此經》說 夜晚天晚日落的時候.
有人帶著一切害病的和被鬼附的.
來到耶穌的庚前.
合成的人都聚集在門前.
耶穌就治好了很多人各樣的病症.
又趕出許多鬼.
不許鬼說話.
因為鬼認識他.
35節 次日早晨.
多早呢 天未亮的時候.
耶穌起來到曠野的地方去.
在那裡禱告.
耶穌出來傳道的時候.
大家知道不單是四處去廣道傳福音.
在廣道傳福音的同時.
吸引了許多人來.
不單單是廣道.
跟會堂裡的律法師 法律出人不一樣.
耶穌來到的時候還會醫病趕鬼.
這個不是醫生.
但竟然能夠治病趕鬼.
奇難雜症 什麼污鬼.
都可以從人的身上趕出來.
所以耶穌沒有去到哪裡.
周圍都有很多人跟著他.
被群眾包圍著.
所以連休息 吃飯的時間都沒有.
說到耶穌這麼忙碌的時候.
今天大家坐在這裡是不是要感恩.
因為今天不用上班.
不是的 牧師.
我回來就是想透透氣.
禮拜一到禮拜五 甚至禮拜六.
現在都要開工.
特別是灣區.
我們都是大忙人.
這個應該怎麼說.
AI時代帶給我們其中一個好處是什麼.
就是可以上網做很多事情.

$^{281}$如果你做工就上網開會.
好像節省了一點時間.
為什麼 不用塞車在路上.
節省了很多時間.
但其實你想想.
可能你更加忙.
以前是上班在公司.
現在家裡也要上班.
因為有互聯網就行了.
在公司開會 在家裡也要開會.
而且不是中國晚午.
半夜三更.
因為現在是全球的.
全世界的.
所以在公司開會 在家裡開會.
要在不同的時段.
跟世界不同地區的同事一起說話.
其實會不會覺得比以前更加忙.
有沒有時間放假.
放假也要帶著電腦手機去.
為什麼 照樣做工.
所以現在公司可以很靈活.
你要回香港就回.
不過照樣做工.
你去台灣 去哪裡玩.
不要緊 不去.
帶著電腦去.
其實沒有假房.
根本是一邊做一邊忙到不得了.
耶穌很忙.
今天我們也很忙.
忙到怎麼樣呢.
忙到不得了 就要放假.
我們現在跟師母經常聊天.
慘了 現在找人開會很難.
因為以前一年放一次假.
現在一年放很多次假.
要遷就人開會很困難.
除非你真的在網上開會.
不然很難遷就所有人.

$^{321}$但是能夠走開一下也是好的.
能夠去旅行 放假也是好的.
但問題是耶穌的時代.
沒有我們今天那麼豪華.
耶穌沒有我們有這個機會.
帶著電腦到處走.
耶穌很忙碌.
但我們要歸根本.
耶穌很忙碌的時間.
祂有一個秘訣.
讓祂能夠很快重新得力.
剛才聖經說得很清楚.
是什麼.
是早上.
天還沒亮.
門徒還沒起床.
吃了早餐沒有.
應該還沒吧 我相信.
去哪裡.
一個人去狂野.
什麼意思.
是沒有人能夠騷擾他的一個地方.
在那裡和上帝聊天.
和上帝祈禱.
其實我今天煮的.
今天講這篇訊息.
因為當我預備這篇訊息的時候.
我看回我自己.
其實最近真的很忙.
很多會議要開.
教會會議 學校會議.
語言研究團隊.
然後傳道部.
聖經研究.
門徒訓練.
我每個星期不是要開會.
就是要做門徒訓練.
現在Tracy又開始了一個門徒訓練.
似乎是為事公.
為教會忙.

$^{361}$我不是抱怨的.
還可以做這些事.
但是為什麼會提醒我.
因為忽然之間我看到耶穌的故事的時候.
我想起從前我沒那麼忙的時候.
還沒做ISP的時候.
我每個星期都上山.
一個人.
開車不用五分鐘就有一個停車場.
可以去走一圈.
我走完回來才沒有味道.
我很喜歡自己的因.
每個星期起碼有一次去上山.
上山的時候其實不是做運動那麼簡單.
我喜歡擔天忘地.
因為天色每天都不一樣.
而每一個季節我所看到的花草樹木又不同.
有時候我走著走著.
我喜歡留在這裡留心看看每一朵野花.
看看野草.
似乎沒有人去兼顧.
但是這邊忽然之間讓我看到神創造萬物.
將生命力放在花草樹木.
所有動物人當中.
那是一個很清新的.
對我來說是一個很令我心靈的更新.
很多時候我在那裡神給我很多的啟發.
和神聊天.
那段時間是很開心的事.
所以星期四我又去了上山.
很久沒有自己去了上山.
還不是氣喘還可以走一圈下來.
當然不是很長時間.
但是那種感覺開始回來了.
你知道人忙起來會煩躁.
所以有時候師母都提醒我.
你是牧師來的.
在會友面前不要表現得這麼煩躁.
老人家是比較煩躁的.
再加上忙碌的時候有時候很難控制.

$^{401}$所以如果被大家看到的話.
請大家包容一下.
當我們忙到一個地步會煩躁.
而且會容易出錯.
容易有遺漏的地方.
所以我們搞這裡真的要大家去體諒和包容.
出去看看師妹.
我猜你可能附近有沒有公園.
或者有沒有一些道路可以去走.
其實我剛才說抗疫的地方是什麼.
只要找到一個地方不受人騷擾.
Only you and God.
一段安靜的時間.
我一睜開眼就去祈禱.
你要感覺到神在那裡才說話.
我與神同行.
讓你自己去走.
但你也不覺得神拖著你.
那是沒什麼用的.
那時候只是做運動.
找個時間找個地方.
沒有人騷擾你.
安靜下來和神聊天.
當你意識到神的同在的時候.
向上帝盡訴心中情.
然後聽.
讓那個微小的聲音.
你越來越聽得清楚.
我不知道你有沒有試過.
那種屬靈的體會.
我鼓勵你如果你沒有的話.
試一試.
第二件事.
我們要用真道取代演算法.
提摩泰的前書第四章.
第七節和六章二十節怎麼說呢.
祂指示要棄絕那些世俗的言語.
和老虎荒謬的話.
在經權上操練自己.
提摩泰.

$^{441}$你要保守所託付你的.
躲避世俗的虛談和那敵真道.
此是而非的學問.
保羅勸命提摩泰.
要放棄要拒絕那些世俗的虛談.
和此是而非的言論.
世俗的空談和此是而非的言論.
其實我們可以引申到.
現在我們說的電腦很厲害.
AI很厲害.
那些電腦程式的語言.
用來寫很多東西.
用來收集資料.
有推算有邏輯分析諸如此類的東西.
所以其實現在我們上網問AI的問題的時候.
它是用大數據推算出一些意見和觀點.
我不是很懂電腦也不是很懂這些東西.
如果當中有科學家幫我矯正一下我說錯的話.
但是我將來引申這些電腦程式或者AI程式.
然後幫我們收集很多資料.
因為它的資料附很多東西.
然後很快很短時間用推算的方式.
將一些答案一些建議給我們.
這些其實有沒有問題在這裡呢.
其實有問題在這裡.
上一次我在另一篇講過.
今天就不講了.
在這裡講到老虎荒謬的話.
其實在當時來說.
是一些特別貶低對方的一些說話.
這些是老虎荒謬的話.
因為當時那些紅面家用這樣的說話.
去諷刺他們對手所講的一派胡言.
沒有根據的或者是荒謬不合理的.
就會稱那些東西為老虎荒謬的話.
因為以前那些姐妹沒有什麼機會受教育.
但是人老了之後有很多東西可以講.
結果都出來講.
所以被人引用老虎荒謬的話.
去貶低對手所講的話是不合理的.

$^{481}$其實我剛才也講過.
人工智能其實是根據人輸入.
那些數據和演算的方程式.
來推算出一個答案給我們.
如果人輸入的資料或者數據有偏差.
或者不足夠或者漏了什麼東西.
那答案可能就有問題了.
可能會出錯了.
但是他要給你答案.
所以那些答案其實可能很有問題.
可能你會爭辯的.
數據越多越來越精確就行了.
但是仍然會有人為的錯誤參雜在其中.
所以當我們要用人工智慧.
問人工智能問題的時候.
我們要很小心.
求上帝給我們有警覺性.
不要閉上雙眼.
總之他出來就照樣搬過去.
應該沒問題的.
因為他那些資料快過我.
推算又快過我.
這些就以為覺得OK的.
其實是有問題的.
所以保羅勸變提摩塔說.
要保守所託付你的.
躲避世俗的虛談和拿捏真道.
似是而非的學問.
我們除了對人工智能的使用.
要有警惕性.
有警覺性之外.
這裡還要告訴我們.
要保守神所託付的真道.
你怎麼能夠有警覺性.
怎麼能夠知道是真是假.
如果我們都不熟悉聖經.
我們不知道什麼是真理.
你真的不知道怎麼評估.
你不是很容易聽起來很邏輯.
很合理.

$^{521}$那你就照單傳修.
但是保羅除了說要警覺之外.
我們更加要保守神託付我們的真道.
這個真道在哪裡.
在聖經裡面.
蘊含在神的話語當中.
所以我們要信一個.
聖經裡面所說.
所啟示的真理.
是永恆不變的真理.
特別是在今天這個時代.
今天這個時代人們是不講真理的.
是講什麼 我講過了.
記不記得.
講feeling.
講印象.
講關係.
不講真理.
這是我們今天世代所面對的問題.
所以當我去問AI的時候.
我就發覺這個問題.
其中一個例子.
我就上去問.
這個比較敏感的問題.
關於LGBT.
LGBT+.
OK.
那我就.
去找的時候就將那個Gallup Poll.
Gallup Poll這個我經常用的.
一個很出名的.
很popular的一個民調的一個機構.
它是最新.
今年2026年.
它說.
美國人有9\%.
說自己是LGBT+.
一個這樣的群體.
9\%.
即是多少人呢.

$^{561}$大概2300萬到2400萬的美國人.
說自己是屬於這個群體的.
其實這個Gallup Poll.
是到2012年.
才開始有系統性每年去追蹤.
搜集這些LGBT的身份認同.
在2012年的時候.
當時只有3.5\%.
說自己是屬於這個群體的.
OK.
到了2016年.
增加到4.1\%.
我剛才說2016年10年之後.
增加了多少\%.
變成9\%.
超出了一倍多.
所以我們看到.
現在這一類.
自己認同這群人的群體.
不單單是增長迅速.
差不多接近10\%.
10個有1個.
這個不得了.
所以已經成為一個似乎是.
勢不可擋的風氣.
所以我們說話要很小心.
是不是.
在教會裡面說不要緊.
因為我們有信仰的立場.
這個就是我要說的東西.
越來越多人現在.
包括基督徒在內.
是接受LGBTQ+的人的.
那我就去問AI.
我問了一個問題.
他說一般社會人士.
我不是說基督徒.
一般社會人士.
如何看待自稱是LGBTQ+的敵人.
我就將這個問題打進去.

$^{601}$他給我大堆的答案.
一個結論是什麼.
結論說大部分的社會人士.
在人際交往中.
通常會保持基本的禮貌和尊重.
並傾向以個人的工作表現.
與待人處事.
來評斷一個人.
而非單純以其性向.
或性別認同來作定義.
這個就是AI給我的答案.
看完之後.
有沒有問題呢.
這個答案是歸納了三類人的反應.
第一個是支持 接受.
OK 沒問題.
第二個是保守 中立.
有concern 有保留的.
不過都無所謂.
有得有不得 無論如何都無所謂.
第三是宗教人士的立場.
歸納了這個東西之後.
還要加上去標籤化這個概念.
然後推算了剛才螢光幕上說的結論.
當你看這個結論的時候.
我不知道你如何去分辨.
數據給你.
那麼聖經是如何看待的呢.
聖經是如何看待同性戀 LGBTQ.
大家知不知道.
很簡單的一集 利未記.
利未記有十八章和二十章都有說.
不可與男人苟合.
像與女人一樣.
這本事可憎惡的.
人若與男人苟合.
像與女人一樣.
他們二人行了可憎的事.
總要把他們賜死.
罪要歸到他們身上.

$^{641}$除了利未記之外.
如果大家熟悉聖經.
《羅馬書》第一章都有講得很清楚.
再肯定起碼同性戀的行為.
不要說變性轉性這些東西.
是神所厭惡的.
為什麼《舊約聖經》《新約聖經》都有這樣的講論呢.
因為創世記講得很清楚.
當神創造人的時候.
人只有一種.
不過神造人的時候.
有男性和女性.
這個大家應該記得清楚.
剛才AI給我們這個結論.
它只是人際關係.
人與人相處的角度去回答這個問題.
有沒有從一個倫理道德的角度去回答這個問題.
有沒有.
沒有.
但我們基督徒.
甚至那些比較保守的中國人.
或者西人也好.
我們不只是說大家相處.
相處也沒有問題.
一點問題也沒有.
我們不應該有偏見.
相處的角度.
因為你覺得他是犯罪嗎.
多少人犯罪.
這個也是他們的爭論點.
但很多時候我們不懂得根據《聖經》.
從倫理角度.
從道德的觀點去評估這個問題.
那些相信所謂進步神學.
我不是要說很深的東西.
Progressive Theology.
那些人相信什麼呢.
你記得羅馬書有關同性戀的記載.
其實有他們的歷史背景.
所以不應該套用在兩情相悅彼此相愛的關係上.

$^{681}$這是一般的答案.
他們兩人相愛.
兩個人喜歡在一起.
有什麼問題呢.
所以男可以和男 女可以和女.
甚至和自己的父母子女.
只要身心相愛.
都可以有夫妻一樣的關係.
這個就是危機.
這個就是問題.
在教會裡我們很少說這些東西.
但我覺得我們的立場要很清楚.
我們的觀點要很清楚.
從舊約到新約聖經裡.
聖經裡都是強調愛的.
這是不論的.
基督教的標誌就是愛.
不論的.
但是要定義什麼是愛.
神愛世人.
我曾經說過這句說話.
大家都知道.
神愛世人 神愛世界上每一個人.
包括了LGBTQ+的人.
I am confident神愛這些人.
但聖經很清楚地說.
神萬不以有罪都為無罪.
你信不信這件事.
神愛每一個人.
但神不包容罪惡.
不包容還會什麼.
會刑罰 會處罰那些犯罪.
得罪他 離開他的人.
所以我在講壇上可以大膽地說.
We love all those people.
我們希望他信耶穌.
讓聖靈可以調整他們的偏差.
愛這些人.
但是我們不接受這些行動.
為什麼不接受.

$^{721}$因為上帝不接受.
這不是上帝的心意.
上帝不喜歡.
這就是得罪神的事.
就好像那些偷偷拐騙的人.
我們都不喜歡他們.
為什麼 得罪上帝.
其實是一樣的.
所以不要被人說我們基督教只針對同性戀LGBT的人.
No.
我們針對所有犯罪得罪上帝的人.
但這些人都是上帝所愛的.
所以求神幫助我們.
如果我們要在數碼時代AI橫行的世代當中.
祂給我們意見.
但是我們要謹慎.
但是你怎麼謹慎.
如果你都不能夠掌握神的話語.
不明白聖經的真理.
你是不知如何去分辨的.
有一部聖經.
以上亞書40章第8節.
我們一起讀吧.
讀得很短而已.
以上亞書40章第8節.
一起讀讀.
草必枯幹.
花必凋謝.
惟有我們神的話.
永遠立定.
信不信.
再來一次.
草必枯幹.
花必凋謝.
惟有我們神的話.
永遠立定.
弟兄姊妹今天我們面對一個很大很大的挑戰.
因為我們.
不看重神的話.
甚至用我們自己的觀點去評估神的話.

$^{761}$好不好.
對不對.
應不應該去聽.
這個是.
我們的問題.
在這個AI世代裡面來說.
就是.
有很多由AI推算出來的言論影像觀點.
衝擊著.
我們的頭腦.
每天我們接收很多這樣的資訊.
所以我們真的要在神的面前.
真的把時間.
靜下來.
跟神去談去說.
那些令我們感到疑惑的地方.
賜給我們有那個.
分辨的能力.
並且相信聖經跟神的話是真理.
如果我們發覺.
我們不是太熟悉聖經.
那就很簡單.
OK.
參加教會的查經班.
一會兒下午我回來開會.
開什麼會呢 開Growth.
就是屬靈成長的.
我們教會有小組.
我們教會有門徒訓練.
不過門徒訓練放下了一段時間了.
怎麼去做門徒訓練呢.
我們教會有.
查經的系統.
但是現在不是很多弟兄姊妹能夠去掌握和利用這些資源.
我們又未能夠有效推動給大家知道.
所以我們覺得這個不行.
怎麼能夠幫助大家建立一個信仰的基礎.
讓我們掌握什麼是真道呢.
這是教會的責任我們需要做.
請你為我們祈禱.

$^{801}$希望能夠.
按部就班 提供一些資源.
和一些機會.
讓大家在神的話語上面.
紮根.
第三樣東西我發覺就是要.
回歸實體的團契.
這是Ada帶給我們的另一個問題.
剛才說到這個虛擬世界我們享受在裡面.
聖經說 不可停止聚會.
這個很適合聖經.
好象那些停止慣了的人 都要彼此.
勸勉.
既看見那日子臨近 就更當如此.
就是我們.
立定心智要回來.
不單單是回來我們聚會.
要參加我們的小組 參加團契.
就是.
有時我們.
不知道會不會成為我們的藉口呢 其實不是啊 有AI我們可以上網.
去參加聚會.
可能我們參加多了也說不定 因為上網可以崇拜 不單崇拜 牧師你說崇拜.
上網可以參加小組.
上網可以參加查經班 上網可以參加門徒訓練 上網還可以有神學課程.
只要你願意 其實現在提供給我們很多便利 我們可以這樣做.
你可以這樣去爭辯.
但不知道我們有沒有想過.
我們在網上.
可不可以互相勸勉呢.
我們在網上.
可不可以用短訊互相關心呢.
我騙你們的.
可以的.
為什麼不可以呢.
可以.
不過.
是有瑕疵的.
是不缺陷的.
對我來說.

$^{841}$那些社交媒體啊.
那些訊息啊 Whatsapp啊Whatever 這些這樣溝通的社交媒體來說.
是.
當我們不能夠.
面對面實體見面的時候呢.
用一些補救的方法.
絕對不能夠取代了我們實際參與敬拜.
團契和彼此相交的方法.
OK 是我們沒辦法.
不能夠做到一樣東西 OK 那好過沒有.
因為在網上 我剛才說那些白蘭花很香.
你聞不聞得到 你聞不聞得到.
你的body language 你的肢體語言.
可不可以觀察到呢 只見到樣子 甚至你樣子都見不到.
OK 你怎樣去知道他.
其實說那句話是真的還是假的.
出自內心還是.
你去觀察到 咦 不是喔 一邊說一邊笑.
不行囉 在網上.
你做不到這些東西囉 是不是.
我好像跟大家說過呢 有時候我去做輔導 去到家裡.
兩個人吵架 一個人在樓上一個人在樓下吵架 行不行.
不是大聲吵 用手機吵架.
可以斯文一點.
是吧 又不摔倒對方的惡樣 多好呢.
這是好處來的 可能是好處來的.
但是是不是真的能夠解決問題呢.
是不是真的能夠溝通呢.
是不是真的能夠大家彼此相交呢.
就是.
AI和社交媒體呢.
就是我們去.
彼此問候相交的時候 其實是.
有很多的缺陷 我們很難有.
就是生命那種互動和情感上.
那個的交流.
都是有缺陷的.
都是片面的 都是不足夠的.
你要和人建立一個真心的關係的話.
你必須要面對面去見面.

$^{881}$當我們.
聚在一起的時候呢.
其實就算你不說話.
就已經能夠生發出一種感染力了 大家知不知道.
上個星期五.
我們在Lupitas.
有一個.
誰啊 成小梅啊.
成小梅.
成小梅的音樂報道會 他們打算三百人 結果來了四百七十人.
我們沒試過爆棚的.
我們走到大堂 如果我們有去過呢.
還要加幾十張椅子在裡面.
坐到爆了.
完了之後 那些人就出來了.
那些人 很多人和我們反映 做檢討的時候.
都不用聽他講道 就是那種感染力.
突然之間我們錄了一個漢幕.
一起唱歌 舉手 讚美上帝.
然後.
呼召的時候 那些人走出去.
我們在裡面 怎麼去鼓勵 拍手掌.
就是那種氣氛.
I was lifted up.
我心靈被提升的那種感覺.
你覺得在裡面.
好震撼啊.
如果你在.
家裡對著電視機來敬拜.
或者電視機.
或者我們上網去參加這個音樂會.
你就看到裡面 黑漆漆 人頭湧湧.
你有沒有那種感受啊.
你沒可能.
人與人之間的相處呢.
是有感染力的 弟兄姊妹 你知道嗎.
不要小看你自己.
我們每個人.
都有感染力的.
問題是你感染些什麼東西.

$^{921}$好還是壞呢 正面還是反面呢.
但是如果你不見面 個個躲在家裡的時候.
我們就喪失了.
這些去感染人 和被別人感染我們.
那些的機會.
而我們會越來越身為一個.
就是孤獨.
空虛的世界當中.
你以為你什麼都可以知道.
你以為你到處都能去到 打開電視機你能看到的.
越來越真實嘛 是不是.
但是其實.
你心靈裡面.
有一種莫名的空虛.
我們最近幫了一個弟兄.
以前不是那麼煩躁 現在真的很煩躁 不是說我自己 是說別人.
就是太忙了.
由早到晚就是對著電腦.
所以跟他老婆說 你不要趕著跟我聊天.
我很忙 沒空 跟老婆說.
這句話可以反映出 其實已經到了那個地步.
去到那裡了.
但是不做也不行.
不做也不行.
跟家人的關係 因為這樣的緣故.
都受到虧損.
所以求神幫助我們.
早期教會我們看到 是怎麼樣的.
全體信徒.
《小人生》第二章.
是繼續在一起過團契的生活 所有的東西都是大家共用的.
他們同心合意天天在聖殿裡面聚會.
又分別在各人的家裡面.
分享愛言.
以歡喜純潔的心一切用法.
頌讚上帝.
跟人人保持和睦的關係 這個人人 不只是基督徒.
主.
天天把得救的人數加給他們.
早期教會那種團契生活.

$^{961}$當然那時候沒有互聯網.
但是我們看到那種效果.
在那個時間他們不單在家裡.
不單在聖殿裡聚會.
家裡也有家庭聚會.
我們現在有沒有.
疫情之後我們每個人都保留.
開放家裡有時候都很緊張.
阿彪說開放家裡叫做Outreach.
那個人我不認識的 帶陌生人來家裡.
行不行.
很多時候有很多顧慮.
所以我們避免跟別人接觸 以為是保護自己.
但是失去了很多.
建立真正友誼的機會.
所以求神幫助我們.
當基督徒能夠聚在一起.
其實是一個.
美好的見證.
簡單來說聖經帶來的結果是什麼.
主將得救的人數天天給他們.
所以求神幫助我們.
要製造更多這樣的機會 讓我們能夠真情相交.
多點見面.
回來團契回來小組.
可以在主教裡面.
互相提醒激勵.
將更多的愛 更多的溫暖 更多的鼓勵.
帶給那些.
沒有盼望.
孤單一個人生活的人.
最後.
我是不是說得太長了.
OK 最後一點.
我們要界定人工智慧 AI的角色是什麼.
其實我們現在每個人都暴露在AI的文化裡面.
有誰沒有用過AI 上次我問沒有人舉手.
現在每個人都在用.
所以其實OK的 用AI不是一個問題.
用AI不是犯罪 不要誤會那個意思.

$^{1001}$我們不需要排斥.
或者是.
故意不去用.
問題是.
我們要清楚知道AI扮演的角色是什麼.
剛才我提過的 AI人工智慧.
是人造出來的東西.
應該是給人用的.
工具.
OK.
但是如果我們不經意的將他打上神壇.
甚至將他變成一個偶像.
那我們就將他拉下來.
其實.
上帝裡面說給我們很清楚的指示.
上帝作了以下的指示.
除我以外不可敬拜別的神明.
不可為自己做任何的偶像.
也不可仿造天上地下地底水中的任何形象.
不可向任何的偶像跪拜.
因為我是上主你們的上帝.
我絕不容忍.
跟我對立的神明.
足夠記載的當然是那些有形的偶像.
AI人工的智能.
可以成為那些.
心裡面無形的偶像.
今年的世界.
現在包括青少年小朋友都有.
手裡面有一個智能手機能夠上網的時候.
我們忽然之間有一個錯覺.
以為AI是.
無所不知.
無所不在.
無所不能.
錯覺啦.
我們不是這麼相信的.
不過有時候我們真的有這樣的想法.
因為問你什麼問題都有答案.
你將什麼挑戰給他.

$^{1041}$他都會給你很多解決的方案.
所以難怪AI在很多人的心目當中.
奉之為神明.
但是神一早知道.
神一早知道.
所以.
不是現在.
在以賽亞的時代神已經透過以賽亞先知跟我們說了.
第二章第八節以賽亞說.
他們的地滿了偶像.
我們今天這個地方也是充滿了偶像.
他們跪拜自己手術做的.
包括AI.
就是自己指頭所做的.
所以我希望在結束的時間大家要記得AI.
人工智慧是.
人類發揮智慧和管理萬物.
因此的一個產物.
神讓我們能夠創造出AI人工智能.
人做出來的.
但是神是神做的.
神enable us.
讓我們有這樣的能力可以製造這樣的工具.
是幫助我們能夠更有效管理這個世界的.
而不是讓他管理我們.
所以我們不要把那些.
服侍我們的工具變成為我們的主人.
我們的軍師.
不知道怎麼做.
問一下AI.
OK.
應該問上帝嘛.
剛才我說過了.
對基督徒來說.
基督徒.
去用AI的時候.
他有龐大的資料庫.
和有高速運算的能力.
求神讓我們有智慧去用.
用來做什麼?.

$^{1081}$去傳福音.
去拓展.
神國的事工.
讓他成為我們.
完成我們使命的一個資源和幫助.
氣象藝術第二章十七節說.
那日子來臨的時候.
人的驕傲要被制止.
我們很proud of我們所創造的東西.
但是.
人的自大要破滅.
偶像要完全被粉碎.
只有上主.
受尊崇.
這是我們基督徒.
對AI應該有的態度.
我們不需要怕.
擔憂或者逃避使用AI.
只要我們能夠什麼?.
心中建立一個為上主.
受尊崇的祭壇.
我們心中有一個聖殿.
裡面的王是我們的救主.
是我們的上帝.
唯有祂.
值得我們尊崇.
當我們把我們的注意力反過來.
現在我們的注意力很分散.
我們注意力很分散.
不僅是AI.
很多東西都拿了我們的注意力.
把注意力拿回來.
我要訓練自己能夠安靜.
聆聽.
和分辨.
上帝.
在我們心中所說的那種微小的聲音.
夠不夠?.
不夠,你聽到.
現在我們很多時候聽不到.

$^{1121}$聽到之後.
憑信心.
去順從.
憑信心去順從.
當你這樣做的時候.
你經歷上帝.
那種真實和存在.
AI帶給我們那種迷惑.
自然就會從那個偶像的神壇裡面.
被拉下來.
所以求神幫助我們,可以繼續用AI的沒問題.
剛開始那幾幅圖畫我們叫AI做而已.
給他一個comment就出了幾幅圖畫出來.
就是這麼簡單.
但是讓我們去用.
不要讓我們去控制.
我們一起祈禱.
因為你賜給人的智慧能力.
竟然可以創造AI這麼特別.
這麼Amazing.
令我們難以相信竟然可以做到這樣的東西出來.
但是這個和你的創造相比是沒得比的.
求天賦裡面我們能夠懂得去分辨.
怎樣去善用人創造出來的這個工具.
去拓展神國的福音.
並且讓我們能夠經常要撥出時間.
和家人相處.
特別是要和你相處.
建立一個更深厚的友誼.
以致我們不會誤會了你的意思.
不會被這個世代裡面那些不正確的言論,思想,潮流.
去衝擊我們.
以致我們離棄你.
天賦我們需要你.
需要聖靈的幫助.
祈求你加添我們的力量.
祈求你天天與我們同在.
又讓我們能夠意識到你真的與我們同在.
我們這樣的祈禱.
靠天父的瑞典幫助我們.

$^{1161}$放耶穌基督的聖名.
是不是?.
\newpage



\section{士師記 6:32}
\label{sec:032282}
\textbf{聖經人聖經事:基甸第一次勝仗 耶路•巴力}~~~~黃克勤牧師博士
\newline
\newline
連結: \href{http://www.church.com.hk/acms/content.asp?site=cdc\&op=show\&type=product\&code=032282\&layout=sermon}{\texttt{032282}} ~~~~ 語音日期: 2026-06-10
\newline
\newline
\hyperref[sec:032281]{\small{< < < PREV SERMON < < <}}
~
\hyperref[sec:index_chronic]{\small{[返順時目]}}
~
\hyperref[sec:032283]{\small{> > > NEXT SERMON > > >}}
\newline
\newline
士師記 6:32
\newline
\begin{longtable}{cl}
\hline
\hline
章節 & 經文 (和合本修訂版)\\
\hline
6:32 & \begin{tabularx}{0.7\textwidth}{X} 所以那日人稱基甸為耶路巴力，意思是：「他拆毀了巴力的壇，讓巴力與他爭辯吧。」 \end{tabularx} \\ \\
[1ex]
\hline
\hline
\end{longtable}
$^{1}$各位姐妹,又一日了,眨下眼又讀聖經.
今天是基澱第一位聖像耶路巴黎.
耶路巴黎是基澱的名字.
第一次聖像.
《生命說話》你帶到聖經人聖經時.
看看上帝在聖經人士裡面.
祂怎樣彰顯祂奇妙的作為.
祂幫了基澱打贏第一場屬靈的戰爭.
是司祭六章二十五節,當那爺.
當那爺即是何時?.
「耶和華沙龍」即是耶和華是平安.
當那爺,耶和華就馬上有說話來到基澱.
就跟祂說:你娶你父親的牛來.
就是那七歲的第二隻牛.
並拆毀你父親為巴黎所築的壇.
撼下壇旁的木偶.
這裡很有趣.
其實這一節我研究了很久都未能夠明白.
或者說得不清楚,可能當日記載的字.
那一爺,耶和華是吩咐基澱做什麼呢?.
「你娶」即是「拿」你父親的牛.
很明白了,是不是?.
「拿」你父親的牛來.
這個「牛」字是用了兩個「牛」字.
一個是「趴」,一個是「show」.
「趴」就是牛獨,或者公牛.
而這個「show」這個字就是一隻牛.
那麼怎麼重覆兩個「牛」字呢?.
這個「牛」第一個「趴」.
其實多數聖經在其他字用「趴」.
是公牛的.
但是有幾段聖經是牛獨,即是小牛.
所以真的難明白.
「你拿你父親那隻公牛來」.
這樣就很明白了.
「那隻公牛牛來」.
這樣就明白了,是不是?.
「就是那」.
原來希伯來文是沒有的.
不過可能要令我們容易明白.

$^{41}$所以中文就加了「就是那」這幾個字.
「拿你父親那隻牛來」.
公牛.
「七歲的第二隻牛」.
也是用「公牛」這個字,或者「牛」這個字.
但是呢.
「七歲」.
「就是那七歲的第二隻」.
那就是哪一隻呢?.
難道這裡第一隻牛,這裡第二隻牛?.
但其實是一隻牛而已.
所以這裡我也要向大家承認.
就是我這一句不是很明白.
「就是拿你父親那隻公牛來」.
「七歲的第二隻」.
我不知道是第一隻還是第二隻了.
總之不要緊.
「拿了一隻牛來」.
可能他父親有幾隻牛,第二隻.
「並且拆毀你父親為巴力所築的塔」.
什麼?.
機電的父親原來是拜神佬?.
他的父親竟然是家裡開壇拜巴力的.
是壇主.
是帶頭拜巴力的人.
而且還不止這樣.
「還撼下壇旁的木偶」.
最有趣的是「木偶」這個字.
「木偶」這個字原來是阿舍拉.
阿舍拉就是巴力的太太.
阿舍拉這個神.
中文翻譯成「木偶」.
對不起.
我要說是翻譯錯了.
「並撼下壇旁的阿舍拉木杖木偶」.
這樣就對了.
阿舍拉是財富幸福.
所以巴巴.
不是巴巴.
丈夫巴力.

$^{81}$巴力就是男性的神.
是管農作和生育的.
而太太阿舍拉.
就是管財富與幸福.
所以他們一拜就拜兩個.
拜巴力和阿舍拉.
似乎機電的父親.
是拜巴力和阿舍拉的壇主.
不是拜神婆.
是拜神公.
真的想不到.
機電竟然是阿舍拉的兒子.
神就跟他說.
你拿你爸爸那隻牛.
七歲那隻.
很強壯的公牛.
然後拆了你爸爸拜巴力的壇.
拼在神壇旁邊.
阿舍拉太太的木偶.
然後再借盆石大塊石上.
整整齊齊地.
為耶和華神築一座壇.
這就是第二件事.
然後把第二隻牛獻為梵西.
撼下阿舍拉木偶做柴.
再看這件事.
第一件事就是拿你爸爸那隻牛.
拼在阿舍拉木偶上.
拆了他的壇.
即是把你爸爸拜神的東西全部拆掉.
把他的牛拿來做祭物.
我很欣賞這四個字.
整整齊齊.
神是說整整齊齊.
即是說把他砌得很漂亮.
那些假神不理他怎樣.
不過我們獻給主的壇.
是要整整齊齊.
牛放上去是整整齊齊.
柴火是整整齊齊.

$^{121}$做一座壇.
把那隻牛獻為祭.
然後用木偶做柴.
再從他僕人中挑了十個人.
很有趣.
於是你見到.
其實姬甸家裡有沒有錢.
有錢.
所以他一回去.
獻祭給耶和華的使者是怎樣.
是一隻羊羔.
二十二升這麼多的細麵做毛餃餅.
還要熬了汁這樣拿來.
很厲害.
拿不拿到這麼多東西給耶和華使者.
可能要找兩個僕人幫忙拿.
因為一隻羊也不少東西.
他有很多僕人.
找了十個僕人.
這十個僕人是甚麼呢.
應該說這十個僕人一定是以色列人.
而且還是根基奠.
而且對上帝還有心.
對耶和華還有心.
否則其他都在拜巴黎.
你叫他們來砍那個像.
他們不肯.
叫得很吵.
他們是選十個忠心.
又能夠和他一起.
又有相同比較好的信仰的人.
照著耶和華吩咐他做了.
不過怕爸爸和本城的人知道.
就不敢在白晝做.
就在夜間做了.
我都很奇怪.
因為不是一件小事.
砍了個壇 劈了個木偶.
帶一隻牛來獻祭.
燙了 燒了.

$^{161}$都很大件事.
沒理由不知道.
夜間就行了.
又沒有人來反對.
好像又沒有人知道.
所以我都認為.
這是神的手在保護他.
是沒有人知道的.
亦都有人跟他幫忙.
城裡面的人清早起來.
看見巴黎的壇拆毀.
這樣都行.
壇旁的亞瑟拉木偶又砍了.
有第二隻牛獻在新竹的壇上.
都是提第二隻牛.
這第二隻牛是甚麼呢.
我就不算最知道.
因為沒有提到第一隻牛是甚麼.
第二隻牛獻在新竹的壇上.
獻了 燒了.
他們就說是誰做的.
他們就查訪之後.
就知道是約阿斯的兒子基奠做的.
雖然那晚做的時候.
那麼大工程.
十個人搞了一整晚.
又拆了壇 又砍了木偶.
又整整齊齊鋪了新的壇.
把牛獻上.
牛有聲音.
有叫聲.
都很轟動.
要放火燒.
竟然沒有人知道.
第二天早上才知道.
而查訪之後.
就有人說.
好像昨晚蹦蹦鬆鬆見到基奠.
是啊是啊 我都見到.
於是查訪之後.

$^{201}$就發覺是約阿斯的兒子基奠.
約阿斯就是他爸爸.
約阿斯這個名字都挺好的.
以後的王都叫約阿斯.
就是基奠做的.
城裡的人就對約阿斯說.
把你兒子交出來.
讓我們殺死他.
因為他拆毀巴力的壇.
砍下壇旁的木偶.
嘩 拆巴力的壇.
又把阿舍拉的木偶砍了.
很生氣.
我們要弄死你兒子.
怎知忽然之間.
上帝的工作很奇妙.
一定要相信的.
就是什麼呢.
約阿斯就說.
回答站著攻擊他的眾人.
這個站著的意思就是站著說.
就是圍住了.
你們是為巴力爭論嗎.
你們為他出頭嗎.
你們要救巴力嗎.
誰為他爭論.
趁早將誰致死.
你要致死這個人.
不如致死你先吧.
為什麼.
如果巴力若果是神.
有人拆毀他的壇.
讓他為自己爭論吧.
很有道理.
真厲害.
我記得我小時候讀書的時候.
也有讀過.
孫中山是基督徒.
他認識了耶穌基督之後.
他將他家附近的神廟裡的.

$^{241}$偶像的手撕斷了.
嘩 有沒有搞錯.
竟然這樣.
你是什麼人.
他好像我印象中.
也是他爸爸出頭.
如果那個木偶真的是神.
他就不會被他撕斷他的手.
這樣東西挺有道理的.
很有趣.
在第三十一節.
若亞斯改一下.
反而兒子基澱做了這件事.
他忽然間醒悟.
回頭.
喂 是啊.
如果巴力是神.
讓他自己和兒子爭論吧.
這麼好東西.
讓他搞定吧.
真是一個真的神.
那你為什麼要救他呢.
為什麼要去救他呢.
為什麼要幫他出頭呢.
讓他自己爭論.
那些人就說.
哎 是啊.
那沒辦法了.
所以當日就叫基督是什麼呢.
耶路巴力.
Yerubal.
Yeru 就是爭辯.
讓巴力和他抗爭.
意思就是.
他拆毀巴力的壇.
讓巴力和他爭論.
所以他的名字叫耶路巴力.
那耶路巴力呢.
反而成為他的花名.
或者別人叫他這樣.

$^{281}$就是讓巴力和他爭辯.
意思就是說.
叫巴力和他爭辯.
就是說他勝過巴力.
這個就是基督的第一場勝仗.
耶路巴力.
那麼他夠膽這麼做.
真厲害.
我覺得.
再看下去.
我們基督是看到十次八次.
很多次.
基督其實.
你看完他的一生.
他是一個很膽小的人.
不算大膽的.
不過是神給他一些東西.
他要求證據.
神給他一些證據.
他就大膽了一點.
他又要多走一步.
神叫他多走一步的時候.
他又要證據.
他又大膽一點.
這樣的.
他一步步練出來的.
膽量都是這樣的.
好了.
神為他做了第一個神蹟.
我要你同在.
你放心.
平安不會死.
然後呢.
就叫他們幫我拆了巴力的壇.
於是基建就真的做了.
然後神保守他.
還有一個很厲害的名字.
叫耶路巴力.
那就厲害了.
等巴力自己跟他談.

$^{321}$這是他第一場屬靈的聖像.
他的聖像是.
神叫他做.
他就做.
他得勝的地方在這裡.
他依靠著上神.
叫我做.
我做.
竟然得了.
這個對他的信心有很大的幫助.
很多時候你會發覺你想不到.
這次呢.
我都是本死的.
可能在基建來講.
就是呢.
拿了爸爸的牛.
拆了我家裡的生意.
接著全城的人走來圍攻.
都想著沒命了.
我爸爸都罵死我了.
誰知道.
他爸爸就經過這樣改變了.
你這麼巴力真是真神的.
等他跟兒子講.
厲害了.
這個呢.
他爸爸又厲害.
他又厲害.
最厲害的就是上帝.
幫助他做了第一次屬靈的聖像.
我們一起禱告.
讚美主.
你是得勝的主.
是一件不是很大的事.
不是一場大仗.
一場小仗在他們的城裡面發生.
不過.
我們的主又先勝一仗.
是透過你的部人.
基建.

$^{361}$本來是一個心驚驚的人.
都躲起來的人.
不過.
你是透過他做了這件事.
竟然間為他打開了一條.
更加勇敢前行的路.
也都令到他爸爸悔改了.
真是想不到你有這麼奇妙的作為.
有這麼大的能力.
主我們一起在這裡就這件事.
基建第一場聖像.
名字都改了.
耶路巴黎.
我們來讚美你.
願以你得著榮耀.
我們讚美主的大臨.
我們再看你自己奇妙的作為.
怎樣透過這個.
本來心驚驚的基建.
去勝過米田人.
求主祝福.
禱告奉耶穌名求.
阿們.
\newpage



\section{馬太福音 9:35-10:15}
\label{sec:032283}
\textbf{門徒的天國使命}~~~~陳傳華牧師博士
\newline
\newline
連結: \href{http://www.church.com.hk/acms/content.asp?site=cdc\&op=show\&type=product\&code=032283\&layout=sermon}{\texttt{032283}} ~~~~ 語音日期: 2026-06-07
\newline
\newline
\hyperref[sec:032282]{\small{< < < PREV SERMON < < <}}
~
\hyperref[sec:index_chronic]{\small{[返順時目]}}
~
\hyperref[sec:032284]{\small{> > > NEXT SERMON > > >}}
\newline
\newline
馬太福音 9:35-10:15
\newline
\begin{longtable}{cl}
\hline
\hline
章節 & 經文 (和合本修訂版)\\
\hline
9:35 & \begin{tabularx}{0.7\textwidth}{X} 耶穌走遍各城各鄉，在他們的會堂裡教導人，宣講天國的福音，又醫治各樣的病症。 \end{tabularx} \\ \\ \relax
9:36 & \begin{tabularx}{0.7\textwidth}{X} 他看見一大群人，就憐憫他們；因為他們困苦無助，如同羊沒有牧人一樣。 \end{tabularx} \\ \\ \relax
9:37 & \begin{tabularx}{0.7\textwidth}{X} 於是他對門徒說：「要收的莊稼多，做工的人少。 \end{tabularx} \\ \\ \relax
9:38 & \begin{tabularx}{0.7\textwidth}{X} 所以，你們要求莊稼的主差遣做工的人出去收他的莊稼。」 \end{tabularx} \\ \\ \relax
10:1 & \begin{tabularx}{0.7\textwidth}{X} 耶穌叫了十二個門徒來，給他們權柄，能驅趕污靈和醫治各樣的疾病。 \end{tabularx} \\ \\ \relax
10:2 & \begin{tabularx}{0.7\textwidth}{X} 這十二使徒的名字如下：頭一個叫西門（又稱彼得），還有他弟弟安得烈，西庇太的兒子雅各和雅各的弟弟約翰， \end{tabularx} \\ \\ \relax
10:3 & \begin{tabularx}{0.7\textwidth}{X} 腓力和巴多羅買，多馬和稅吏馬太，亞勒腓的兒子雅各，和達太， \end{tabularx} \\ \\ \relax
10:4 & \begin{tabularx}{0.7\textwidth}{X} 激進黨的西門，還有出賣耶穌的加略人猶大。 \end{tabularx} \\ \\ \relax
10:5 & \begin{tabularx}{0.7\textwidth}{X} 耶穌差遣這十二個人出去，吩咐他們說：「外邦人的路，你們不要走；撒瑪利亞人的城，你們不要進； \end{tabularx} \\ \\ \relax
10:6 & \begin{tabularx}{0.7\textwidth}{X} 寧可往以色列家迷失的羊那裡去。 \end{tabularx} \\ \\ \relax
10:7 & \begin{tabularx}{0.7\textwidth}{X} 要邊走邊傳，說『天國近了』。 \end{tabularx} \\ \\ \relax
10:8 & \begin{tabularx}{0.7\textwidth}{X} 要醫治病人，使死人復活，使痲瘋病人潔淨，把鬼趕出去。你們白白地得來，也要白白地給人。 \end{tabularx} \\ \\ \relax
10:9 & \begin{tabularx}{0.7\textwidth}{X} 腰袋裡不要帶金銀銅錢； \end{tabularx} \\ \\ \relax
10:10 & \begin{tabularx}{0.7\textwidth}{X} 途中不要帶行囊，不要帶兩件內衣，也不要帶鞋子和手杖，因為工人得飲食是應當的。 \end{tabularx} \\ \\ \relax
10:11 & \begin{tabularx}{0.7\textwidth}{X} 你們無論進哪一城、哪一村，要打聽那裡誰是合適的人，就住在他家，直住到離開的時候。 \end{tabularx} \\ \\ \relax
10:12 & \begin{tabularx}{0.7\textwidth}{X} 進他家時，要向那家請安。 \end{tabularx} \\ \\ \relax
10:13 & \begin{tabularx}{0.7\textwidth}{X} 那家若配得平安，你們所求的平安就臨到那家；若不配得，你們所求的平安仍歸你們。 \end{tabularx} \\ \\ \relax
10:14 & \begin{tabularx}{0.7\textwidth}{X} 凡不接待你們，不聽你們話的人，你們離開那家，或是那城的時候，要跺掉你們腳上的塵土。 \end{tabularx} \\ \\ \relax
10:15 & \begin{tabularx}{0.7\textwidth}{X} 我實在告訴你們，在審判的日子，所多瑪和蛾摩拉地方所受的，比那城還容易受呢！」 \end{tabularx} \\ \\
[1ex]
\hline
\hline
\end{longtable}
$^{1}$我們知道世界追求都是最大的企業.
但是作為神的家.
我們建立就是追求大的致命.
我讀大學的時候.
在我同學的家裡經常做功課.
原因是當時開始有電腦了.
不用打字機.
不知道什麼年代才知道.
但是我的電腦可能要等一段時間.
打印很久.
所以有些比較有錢的同學.
他已經買了電腦.
因為我打印機.
我經常在這裡工作.
經常通宵大彈工作.
他有個同學的爸爸看到.
哇,這小子很勤奮.
所以他就跟同學說.
介紹我認識.
認識了之後.
不如你大學畢業下來幫我做事.
其實我也不想晚上的.
為什麼?.
因為白天的電腦當然是同學先用.
晚上我用.
結果我真的去了公司工作.
公司做什麼的呢?.
就是做眼珠的.
你看看我們旁邊的眼珠.
是誰做的?.
其實是做公仔眼珠.
全球百分之十的眼珠都是他做的.
他的企業有一千多人.
哇,很大的.
真的不得了.
大企業.
做了一年多.
就想也未必最適合我.
做回我在金融上面.
在他的會計師樓叫做羅冰咸.

$^{41}$當時有幾千人.
現在不知道會不會倒閉了.
被人告了十幾億.
真的慘了.
哇,很大.
所以有幾千人.
國內又有人.
這樣很厲害.
全球也很多人.
慢慢做著做著.
在金融企業裡面.
又幫一家公司上市.
公司是做玻璃的.
全球百分之十幾的頭髮玻璃去做.
哇,很多人.
公司一萬多二萬人.
真的是很大的企業.
其實也不是那麼大.
原來全世界有五百強的企業.
是請很多人的.
行不行?.
例如排第441位.
這間公司你一定認識的.
就叫做Star什麼?.
八隻勁.
猜到了.
可以玩那些動作.
猜一下.
你猜他請了多少人呢?.
三萬?.
五萬?.
全球?.
十萬?.
是三十六萬人.
你每次喝咖啡.
你會發現.
你還在三十六萬之一.
你不要欺負他.
後面很多人的.
真的是很多人的.

$^{81}$我不知道你知不知道.
全香港的勞動人口.
其實只有三百多萬.
就是他請了十分之一.
好了.
不要數那麼久.
四百零.
還是四百三百九.
這樣數.
數一百三十五.
你一定認識的這間公司.
這間公司你應該用過他.
就是叫做Fenix.
有沒有人用過Fenix?.
舉一下手.
哇,真的很多.
你猜他請了多少人呢?.
全球?.
是.
叫一下.
八零八十.
這樣是不是有六百多人?.
快到了,慢慢上.
五十萬而已.
都很大了.
好了,誰入油的.
可能你用過那些油.
第五位馬上.
很快.
是中國石油天然氣.
你猜他請了多少人呢?.
全球?.
是.
猜中.
梁牧師一會兒請你.
上去喝黑牛.
多少人呢?.
八十萬.
差不多了.
八十萬,多數一百萬.

$^{121}$好了,第二位你一定用過他的.
第二位你猜是什麼公司呢?.
是第二大.
就是Amazon,行不行?.
好了.
Amazon,OK.
就是人最多的.
差不多第二了.
有沒有人用過Amazon?.
有了,舉一下手.
是,大部分人用過了.
他請了多少人呢?.
猜一下.
是,差不多了.
八二多少,一百五十萬.
一半是香港人不到的.
好,最後.
第一位請人最多的是.
等等等等等等.
是.
麥當勞.
麥當勞的朋友.
肉二碼.
很厲害.
請了多少人呢?.
二百萬.
嘩,很大.
真是很大的企業.
真是不得.
你想一下,很難想像.
香港整個勞動有三分二.
你想一下,一出街.
嘩,你每三個有一個是他的員工.
如果是.
都是他全球.
不過,這麼大的企業.
都不夠這間企業大.
這間企業叫什麼名字?.
叫做天國企業.
我不知道大家知不知道.

$^{161}$全球有多少間教會.
大概統計超過三千萬間.
對不對?.
如果.
每間的童工.
我們教會童工三十多.
我當每個平均童三個人而已.
三千多萬間已經請了一億人.
對不對?.
還沒算義工.
我們有童工一個至少有十個義工.
無論敬拜隊.
無論迎新.
咖啡.
SV.
很多的.
最少估計了十億.
你看一看你旁邊.
你是天國企業的其中一環.
你不覺得很感動嗎?.
你有沒有事奉?.
對不對?.
你跟你旁邊的人說.
嘩.
你很厲害.
真的很厲害.
天國企業.
以億計了,肯定.
就是有事奉.
當然童工可能也有一億.
好了.
我想今天說這段經文有什麼特別?.
有什麼關係呢?.
就是.
天國企業今天以億計的人.
第一次開始去招聘.
是這段經文.
第一次他聘請天國企業的員工.
是從這段經文開始.
這段經文裡面.

$^{201}$包括我們的心態.
心情.
包括我們要仰望的.
我們的權柄等等等等在裡面.
所以這段經文對我們來說非常重要.
也是因為我們十二週年.
今天我多送一點給大家.
總共有幾點?.
四點.
再加六個秘笈.
嘩.
十全十美才離開.
所以跟我們旁邊的人說.
你今天一定吃飽了.
好不好?.
一定吃飽了.
我們立刻進入.
第一個.
第一點.
我們學習的第一點.
看看.
在這段經文裡面.
怎麼開始.
怎麼天國企業第一次做了什麼?.
背景是什麼?.
我們一起說.
一二三.
耶穌走遍各城各鄉.
載會堂女教訓人.
宣講天國的福音.
由伊斯華人的長城.
花園的多人.
走遍了花園.
因為花園很受歡迎.
與同樣會有福人一般.
當時耶穌走遍什麼?.
各城各鄉.
會堂女教訓人.
又講福音.
又醫治.

$^{241}$他見了很多人.
然後.
原來我們知道.
如果我們要做天國事公.
第一樣不是要人多.
也不是要計劃.
那些都需要計劃.
他最重要的是什麼?.
憐憫.
他說見到人憐憫.
見到他們的需要.
而裡面說.
因為他們.
困苦.
流離.
意思是說.
原本的意思是說.
他們被騷擾.
裡面不平安.
無助.
無奈分散.
愚弱沒有牧羊人.
外表都很活躍.
不過他沒有了方向.
我們不要覺得.
我們身邊人個個很正常.
你看看你旁邊的正不正常.
或者你公司人很正常.
你的家人很正常.
裡面都有很深層的需要.
和缺乏.
我就用這幾天的報道會.
作為例子.
梁牧師在牛津畢業.
還是律師.
厲害嗎?.
對我來說夢寐以求.
牛津畢業當然.
但他分享的東西很真實.
他當日說.

$^{281}$那裡的人很厲害.
經常怕自己不夠別人給.
他怕有東西會失去.
外表什麼都有.
他經常說他父親很疼他.
這真是真的.
他說他太太很愛他.
這真是真的.
他經常都說.
那些童工很好.
這真是真的.
但是他當日未信主.
他裡面發現有一種空.
一種恐懼.
於是乎他就用賭.
來遮掩自己裡面的空虛.
很真實的.
我們很多時候都會.
到處去找一些填補.
可以嗎?.
第二天牙醫.
他說牙醫什麼都有.
就一個浪.
金融危機.
稍微什麼都沒有.
發現人可以有東西.
可以這麼快沒有.
究竟上主又是怎樣.
在那一刻.
他聆聽上主.
在《梁牧師》裡面說.
當日耶穌跟他說.
你認識那個耶穌.
但是你跟他未有關係.
建立.
我們見到他.
我們周圍的朋友.
不覺得他的樣子.
都是兩隻眼.
一個鼻一個口兩個耳.

$^{321}$其實他裡面可能是破損的.
耶穌有這種眼睛.
如果我們被邀請成為天國的員工.
我們一定要見到耶穌見的東西.
而當耶穌見到.
於是乎在三十七,三十八節.
他這樣說.
我們一起讀.
一,二,三.
「於是對門徒說.
要收的莊稼多.
作工的人少.
所以你們當求莊稼的主.
打發工人出去.
收他的莊稼」.
好了.
他說莊稼多很多.
但是你看不看到.
其實很多的.
作工的人怎樣.
少比對比.
多其實是豐富.
周圍都是.
很有趣.
然後他就叫門徒.
你求我差你.
不是我差門徒.
其實耶穌有很多方法.
處理這個問題.
很多失散的靈魂.
但是不夠人做事.
第一個他可以做的是怎樣呢?.
他就不要一天做十個鐘那麼少.
過去真的做得很忙.
做十二個鐘.
十四個鐘.
十八個鐘.
神來的不睡也可以.
一天.
千年如一日.

$^{361}$一分鐘如一千分鐘.
睡一分鐘立刻睡了八個鐘.
他沒有.
他又可以這樣說.
神做七天的.
不如做八天,九天做多點事.
做足九天才有安息.
第十天進入安息就可以了.
沒有.
又或者他不要三十三對四.
做到八十三.
再不是八百三也可以.
沒有.
他就是要招聘你們.
他就是要邀請我們成為天國軟弓.
這段經文很有趣.
不單是這樣.
他不是差你.
他要你求他差你的意思是.
你準備好了嗎?.
你是不是見到他們的需要.
還是經常見到自己.
當你開始經歷上主的改變.
你就會發現.
我開始見到神見到的東西.
不同了.
今天也很有趣.
門徒是誰呢?.
不是每一個信主的人.
是耶穌學生.
是跟隨者.
常與耶穌同在人.
耶穌就邀請他求耶穌差他.
所以你看看你旁邊那個.
你問他.
你信了主沒有?.
信了是不是?.
你也有份被神邀請的.
你跟他說.
神也邀請你.

$^{401}$好了.
你說門徒.
當然要很厲害才行.
錯.
為什麼呢?.
第九章.
你看看第九節就開始.
呼召門徒.
十二門徒都還沒呼召完.
第十章才呼召.
他呼召一個就差一個.
呼召一個就差一個.
對不對?.
我先讀了神學.
讀了博士才說錯.
如果等到那一刻才傳就不會傳了.
立刻就傳可以了.
我們看看第九章第九節.
讀一讀這一節.
一二三.
耶穌從那裡往前走.
看見一個人名叫馬太.
坐在隧棺上.
就對他說.
你跟從我來.
他就起來.
跟從了耶穌.
對不對?.
跟從我就差了.
可以嗎?.
很真實.
那你會說為什麼會.
這樣也不夠工人呢?.
為什麼這麼迫切用宗教做比喻呢?.
因為你知道如果宗教熟了.
你不收他會怎樣?.
死.
他不會等你的.
我不會掉下來的.
我繼續等你.

$^{441}$明天來.
下個星期.
下年.
不會的.
他掉下來就死了.
就是說人生準備著福音.
如果你不去收割.
他會怎樣呢?.
他不等你的了.
他會給別的東西填滿的.
如果用梁牧師為例子.
怎樣?.
倒填滿他.
是不是?.
牙醫怎樣?.
用工作,用金錢.
還有很多東西可以填滿這個世界的需要.
偶像.
物質.
娛樂.
旅行.
日本食物.
是不是?.
福光.
東京.
是不是?.
北極.
你遲些接下來可以去.
火星.
很快就沒了神.
接著下一個循環可能要十幾年後了.
所以我們要立刻見到就要.
但是為什麼也不夠工人呢?.
因為可能有幾個原因的.
第一.
I don't care.
關我什麼事呢?.
那些你做你做你做你做.
我記得2024年奧運我也很喜歡看.
就是乒乓球.

$^{481}$喜歡看游泳.
我覺得何詩蓓真的很厲害.
打劍.
張家朗六尺高.
多羨慕.
這一輩子都不行了.
可能不行.
我記得有一次跟第一姐妹聊天.
這個項目很棒.
什麼時候做?.
這個第一姐妹說我喜歡看什麼?.
再上一屆說我喜歡打空手道.
很有型的.
誰知道那次聊天有個第一姐妹說.
我都不看了.
沒興趣.
說到興起.
說幾句而已.
不過沒什麼問題的.
奧運的確有人喜歡看.
也有人不喜歡看.
可以選擇.
不過大使命是沒得選擇的.
你信主就必然要猜出去的.
所以這個很真實.
不是I don't care.
你要說I care.
待會我會解答這些問題的.
知道有的.
但是好像看不到.
很忙.
很多事情要做.
家裡有事要做.
又要照顧小孩.
然後工作又要做.
又要進修.
然後這個朋友又罵我去玩.
什麼都優先的.
耶穌淪為最後的乞丐.
我們有時間給你.

$^{521}$我們叫消費主義.
不是的.
沒有抵觸的.
你可以做好你的工作.
我可以關心我的家人.
更加需要將福音告訴他們.
待會我會提到.
第三.
我看到也想.
不過我不行.
我沒能力.
因為我覺得我沒口才.
我不像梁牧師的口才.
我不像阿J這麼帥.
又說得好.
很乖.
不好意思.
我看到什麼就玩.
對不起.
但是.
不是.
這裡告訴我們.
待會我會提到.
福音不是口才.
是關乎生命和權柄.
我們會回應.
如果我們經歷過舊恩.
不是這樣的.
經歷舊恩我們會充滿著力量.
是不會的.
你裡面澎湃的.
就好像.
Gregg Ordon.
莫夫勒老師說.
我們今天世界妥協了.
如果我們成為門徒就不會妥協的.
我們會經歷聖經.
經歷神的擁護和舊恩.
我們會效法基督.
我們會做基督所做的事情.

$^{561}$我們會看到基督所做的事情.
我們的價值觀和世界價值觀.
會轉換到天國價值觀的.
我們慢慢去到那裡很重要.
並且我們開始會經歷一個恩賜的建立.
我們會發揮屬靈的恩賜.
我們有不同的品格.
所以第一點我們要學.
不是避了差.
而是求差.
是避神差.
很重要.
我記得我這兩天看到一個組長.
星期四晚帶了他的老闆.
他也是很高級的老闆來.
聽完梁牧師很有感動.
信了主.
星期五又看到他帶另一個部門的上司來.
聽完陳意很感動.
又信了主.
我就WhatsApp他.
哇.
很感恩啊.
你帶了這麼多同事來.
真的很感動.
我就跟他說.
你記不記得兩年前.
我跟你聊天的時候.
你說我不會帶人信主的.
我都幾十歲了.
我只是做教導而已.
我什麼都不會的.
那些福音交給別人.
他就記得了.
但原來不是的.
所以我們每個人都可以有這個力量.
他說今天很開心.
可以這樣帶.
並且他說.
是因為神的大能.

$^{601}$The power of Jesus.
在他生命中可以彰顯.
所以第一點我們每一個人都可以的.
我們不要被撒旦騙了.
不要被我們的世界騙了.
我們在說.
第一個.
我們學到不是避開差.
是著重在主那裡.
主啊求你差遣我.
主啊求你讓我看到你見的東西.
主啊求你讓我看到憐憫的心.
主啊求你讓我看到哪一代的莊稼.
主啊求你差遣我.
所以這個大使命關於每一個人.
所以你會要去的.
作為傳道.
作為組長.
作為頂禪會也是.
我們都要思想.
究竟我們是不是在做主的工作.
教會不是行政.
教會不只是崇拜.
教會是使命.
教會是栽培.
教會是生命.
第二點我們繼續下去.
「不在乎能力,在乎權兵」.
剛才我說我都沒有口才.
我都不懂得說.
不是.
是在乎生命和耶穌的權兵.
剛才我早上說.
我看到耶穌的權兵.
耶穌的大能彰顯.
我們看下去你就知道了.
去到第十章.
我稍後會讀第一節.
第五.
第六.

$^{641}$第七節.
咦?你說為什麼不讀第二.
第三.
第四節呢?.
其實第二.
第三.
第四節是在說.
耶穌在叫十二門徒的時候才開始的.
可以嗎?.
剛剛好.
我們讀讀第一章.
第五.
第六.
第七節.
一起讀.
一.
二.
三.
「耶穌叫了十二門徒來」.
「吸他們的權兵」.
「讓他們到廣州總會去」.
「去見諸多人的聖聖」.
「耶穌派了十二萬人過去」.
「吩咐他們說」.
「海外人的路你們不要走」.
「撒瑪利亞人的城你們不要進」.
「聽我講」.
「以色列加上新疆拿到去」.
「隨走隨存」.
「說天國近利」.
是呀.
這裡說.
「十二門徒來」.
「吸他們什麼」.
「權兵」.
所以我們可以傳福音.
可以講述這一節不是我們的能力.
我們要禱告.
我們要尋求.
我們要交給主.

$^{681}$而裡面說天國近了.
就是說我們要悔改.
跟從主.
去不同的地方.
裡面我們說.
為什麼有些人可以.
帶這麼多人信主.
經歷了神的改變.
為什麼我沒有呢?.
如果你這樣想.
可能是被撒冷騙了.
其實你有權兵.
我記得.
有幾個例子.
我也很感動.
記得二十多年前.
我在一個小組.
在帶茶經.
有一個組員.
他也是信主.
沒多久.
接著他跟我說.
他帶了姐姐信主.
接著他說.
他要跟姐姐家拆偶像.
他就邀請我.
我覺得他是組長.
我還沒有讀神學.
去到他家拆偶像.
看到他姐姐家有一堆藥.
我就問什麼事.
原來他骨折了很多年.
什麼醫生看完.
吃了什麼藥.
都很痛.
我就說.
這樣.
不過他說現在不用吃了.
我說為什麼.
他說弟弟將風音帶給我.

$^{721}$為我的腳祈禱.
康復了.
我說.
康復了.
我剛剛跟他弟弟確認一下.
今天已經二十多年前.
到現在完全康復了.
當時拆完偶像.
認識完他.
我就跟他弟弟聊天.
我說你怎樣祈禱.
他說怎樣祈禱.
按照你之前.
茶經.
我們在查《路加福音》.
他說耶穌有權.
並有能力.
可以醫治.
我就照著聖經所說的祈禱.
他就康復了.
我張開嘴巴就可以了.
原來真的可以.
有時候我都忘記了.
當然神何時醫治.
不是我們的掌權.
我們只是責任將風音帶給人.
並且按照我們的感動祈禱.
神怎樣工作.
不工作我們不知道.
好像昨晚.
黃明慧.
一個瞎了八歲.
藥物敏感.
她神沒有醫好她的眼睛.
不過不要怕只要信神給她另一個能力.
去渡過一個瞎的人做不到的事.
結果拿了六個獎學金.
三個大學畢業碩士.
這個神有另一個恩怨.
我們不知道神給了甚麼恩怨.

$^{761}$我又很記得.
在星期五早上講見證的嘉傑.
這個能力是很奇妙.
十四年前我們第一次開嘉年華.
在寶安路.
他當日進來他是街坊.
本來他拖著兩個兒子.
他想了結自己的生命.
但他公開講過我就不公開講.
因為他見到兒子很困難.
有自閉有過度活躍.
小巴坐不剛剛那個位.
他會踢不開心大叫.
在學校當他上課.
拿一把硬石不停打到玻璃碎.
差點插到另一位同學.
很困難.
他老公又很麻煩.
又賺不到錢最後都被離婚.
他這麼艱難的時候.
他覺得人生沒有盼望.
沒有意義沒有力量.
他當日想了結生命的時候.
就收到三張門票.
就是嘉年華門票.
進來玩完.
並且見到大家很有愛心.
你都是天國的員工.
因為這樣他信了主.
他分享.
晚上回到家.
耶穌就跟他說話.
這是我帶你小學的教會學校.
你認識過耶穌.
你沒有跟我沒有來找我.
今天我差派教會來找你.
哇 耶穌說話的.
耶穌不單止跟梁牧師說話.
跟陳醫說話.
跟嘉傑說話.

$^{801}$他那一刻很不同.
這個故事實在太感動.
但我現在不去講.
因為時間不多.
也沒有錄音.
所以下次再來.
你要把握機會聽.
把故事講短了.
今天他不單止不需要被服侍.
他已經參加了服侍了一幅詩人.
今天兩個小朋友.
在兒童少年成長了.
今天已經有工作去做.
上帝是拯救生命.
將死的都可以變復活.
是嗎?.
將沒有盼望的.
都可以成為有盼望的嗎?.
Hallelujah.
很真實的.
為什麼這麼多人正在行使權柄.
但是這麼多人也沒有力量去傳道.
因為你被撒旦騙了.
你被這個世界的工作繁忙騙了.
你被自己的議程騙了.
上帝告訴我們.
不在乎我們的能力.
在乎神的權柄.
秘笈的第四 你要行使.
只是聽到行使不了.
你出去了可能感動完之後.
你經歷不了.
但是你一旦經歷過.
你人生就永遠會記得.
神給我們的權柄不單止越來越多.
更加鼓勵我們要寫出去.
第三點.
不是得更多這麼簡單.
是要捨棄更多.
接著我們讀第三點的經文.

$^{841}$裡面有兩個分點.
我們先讀第八到第十三字.
一二三.
醫治病人 叫死人復活.
叫想大麻風的潔淨.
把鬼趕出去.
你們白白的得來.
也要白白的捨去.
腰袋裡不要帶金銀銅錢.
行路不要帶口袋.
不要帶兩件掛紙.
也不要帶太多話杖.
因為眾人不認識自己講的.
你們無論進大城還是進村.
都要打聽考慮誰是好人.
就算在客家.
經濟做不好的時候.
都要進大家裡去.
要請大家平安.
大家若不匯得平安.
你們所求的平安.
就未能到大家.
若不匯得你們所求的平安.
乃歸你們.
有兩件事.
第一.
福音是白白來白白去的.
第二.
就是我們要尋找願意聽的.
有些人不聽的.
不聽不要緊.
我們繼續尋找.
平安.
你是有能力賜給他的.
你見不見到.
你去到請得平安.
配得平安之後.
你將這個平安的福音帶給他.
你是有力量.
為什麼嘉傑.

$^{881}$為什麼梁牧師.
為什麼陳毅會改變.
因為.
上帝用了我們卑微不完美的人.
可以完成他完美的救恩.
這個是瘋的.
一個完美的救恩可以用著.
是有瑕疵的你和我.
你望一望你旁邊那個有沒有瑕疵.
有就對了.
有神才用得到.
很奇妙的.
第一樣我想今天特別講的.
得飲食是神供應.
不用追.
我想這部分特別想和我們的傳道.
或者將會在想全世界讀傳道的人.
分享.
我們的確行上的路.
是需要信心的.
我記得兩年前十週年的時候.
我請了Pastor 勁.
請了我老婆分享一件事情.
都是講道.
做傳道.
我偷了一點點來講一下.
當時Pastor 勁他講.
有一晚他講.
他是科大畢業.
他的同學已經找到錢了.
很厲害.
錢方面.
他就說他作為傳道人.
他的工資是很穩定地低.
他這樣分享.
我寫了.
這麼多年了.
的確他的同學又多又成.
不過他覺得他服侍主.
非常感動滿足.

$^{921}$簡單生活就夠了.
我老婆又講.
這個是我請不到她分享的.
是梁小姐請到的.
因為我老婆不理睬我.
他又說因為我老婆在國內十幾歲才來.
在國內.
吃每個星期.
都是吃一件事就是辣椒.
辣椒炒飯.
一個星期只有一次吃雞蛋.
所以我老婆吃辣椒吃到她很辣.
一個月只吃一次肉.
那塊肉只有蛋麻雀這麼小顆.
所以來到香港.
很努力讀書找工作.
並且在讀大學的時候.
認識了一個很好的男生.
我想問一下為甚麼你笑呢.
你是不是不同意呢.
牧師你看不看到你的回應.
重點不對.
嚇到他哭了.
對了.
他說.
然後等等.
好了.
然後他又這樣分享.
當我呼召.
當時我老公事業上攀向高峰.
高級會計師和分析師.
也拿了獎.
當時成為公司CEO接班人.
他收入佔了家裡的95\%.
而在那個時刻他被呼召.
他說他當然順服神.
也支持他老公.
但這代表甚麼呢.
代表負代價.
甚麼代價呢.

$^{961}$他說以前他上班是賺錢買花.
但現在真的要努力工作.
然後他說.
兒子四歲女兒才一歲多.
他裡面盤旋了很多東西.
他說做教會的傳道人的子女.
通常都要收拾別人的衣服.
吃東西也沒錢.
然後要問人拿食物.
他說不知為何這麼多的圖畫.
在腦袋裡浮出來.
這些都沒有發生過.
不用怕.
沒有發生過.
神對我們很好.
然後呢.
他在分享二十年過去了.
很感恩.
一樣都沒有發生.
我兒女都長大了.
同工對我很好.
教會對我很好.
電視節目對我很好.
這二十多年他的收入升了幾倍.
彌補了缺口.
當時我媽媽也很擔心.
也不會說出聲.
所以現在我經常看著我老婆說.
She earns, I spend.
She works, I serve.
到了今天.
回頭看.
我始終覺得.
足夠就好了.
生命就滿足.
我想問一下.
這一堂的傳道.
梁牧師.
做了哪個傳道金的彩樓.
你覺得.

$^{1001}$阿姐在這裡.
你同不同意嗎?.
梁牧師同不同意嗎?.
同意.
對了.
我每堂見到的傳道都是.
還有Brian.
恭喜你.
Brian.
你會見到.
你會同意嗎?.
我深信.
二,三十年後.
到我差不多退休.
你回頭看.
恩典救人.
這是真的.
不用怕.
網上在聽的父母不用怕.
他們隔壁在聽.
夠的不用怕.
還有不僅如此.
我們要尋求願意聽.
剛才我早就想說.
你不要以為這麼容易.
那個老闆又來.
那個老闆說.
星期五個個都不來.
去到最後.
他打到隔壁部門的主管.
問你有沒有興趣來聽.
他突然說肯.
那天晚上只有一個人肯.
那天晚上就信了.
好嗎?.
我也是.
星期四晚放飛機的.
沒有的.
在不同的天也放飛機.
到星期六.

$^{1041}$我就找到四個人來.
有兩個信了.
兩個還沒信.
竟然星期六兩個都來.
接著來了.
還沒信的那兩個星期六晚.
都信了耶穌.
是啊.
很多時候被人拒絕的.
不要緊.
你拒絕我.
明年又會見到我的.
你越瞄我.
我就會用福音來愛你.
行嗎?.
你怎麼瞄我都沒用.
你越瞄我.
你會見到我.
一年後兩年後.
又是我.
我繼續會帶他聊天.
關心他.
肯定會吃飯.
很重要的.
第五個秘訣是不斷尋找的.
不用怕.
行嗎?.
不怕被人拒絕.
沒問題的.
我現在很興奮.
昨天那些人信了.
哇.
大家這個爸爸.
又那個.
這幾天又那個媽媽.
又這個朋友.
昨天帶來信了.
我走去.
哇.
大家都抱著.

$^{1081}$哇.
很開心.
你信主.
通常都是這樣的.
初信主那些不是很明白的.
你是不是瘋了.
比我更開心.
我跟他說.
新信主.
你遲些你會明白.
當你經歷福音的改變.
當你帶人信主.
你就知道.
那種力量.
那種渴望.
是多麼真實.
弟兄姊妹.
如果你沒有經歷這個豐富.
有時我們會有惰性.
有困難.
不好.
那是撒旦的蒙蔽.
當你經歷這種力量.
你就知道.
更新的力量.
必定是.
從誰來的?.
從耶穌.
我們給贊成主.
真的很辛苦.
最後.
其中一個難阻的就是.
我們覺得時間很多.
但不是.
其實時間不多.
第四十五節我們一起讀.
一二三.
你們無論盡.
讀錯了.
對不起.

$^{1121}$再來.
剛才那一段.
兩秒刪掉了.
再來一次.
一二三.
凡不攝待你們.
不聽你們話的人.
你們會離開大家.
困於難性的時候.
受罰約瑞的塵土.
會躲到花卉.
我實在要告訴你們.
當審判的日子.
所躲藏的禍.
我帶在左手裡.
給難性的環境.
容易受益.
這裡有背景的.
為什麼塵土躲下去.
就是說當日其實.
信奉猶太教宗的耶和華神.
那些在外邦地.
離開的時候.
他會除去那些腳.
衣服那些塵土.
代表異教的風俗.
對不對?.
他們拜偶像.
或者你走一些不合意的東西.
不關事.
傳播音也是.
就是去完之後.
你給個平安.
你有力量給他.
他接他有平安.
他會改變.
不接還給你.
而且他說.
有審判的日子的.
的確是有的.

$^{1161}$而審判的日子.
比所多瑪俄摩拉所受的更加嚴厲.
意思是說.
當日所多瑪俄摩拉.
都沒見到耶穌的神蹟.
今天你看到神蹟都不信.
那個後果很嚴重.
我們今天很怕告訴別人.
你不信主會落地獄.
我都知道.
其實很難說的.
但是這個是真實的.
是永遠都分隔的.
我記得當我女朋友.
當日我們中五,六拍拖.
去到大學的時候.
我的導師跟我說.
我拍拖的時候大家沒信主.
我拍拖我就信主.
我不斷傳播音.
我女朋友都不信.
我導師跟我說.
這樣不行的.
因為跟不信主是很難走的.
你想想我們瘋癲的.
我們又會回去崇拜.
又會受到奉獻.
又會插街坊.
搞一大堆事情.
對不信主是匪夷所思.
我很尊重不信主的人.
也有很多人跟我做好朋友.
不過的確很難.
可以明白我們做的事情.
長期走不了.
所以我就定了.
那一年.
就是兩年後拍拖的那個暑假.
如果他不信主.
之後就分手了.

$^{1201}$他不知道的.
行不行?.
去到八月中都沒信主.
我又傳播音.
又帶報道去都不信.
我記得那一次.
是在我爸爸的茶餐廳九龍城.
我得知.
應該跟爸爸說.
咦 九龍城.
潮語浸信會有個報道會.
嘩 把握最後時刻.
但我老婆.
不是 不是 不是.
我女朋友不知道的.
我又很想我爸爸去.
我又很想我女朋友去.
但我爸爸又要收錢.
在茶餐廳當日.
沒有現在的這樣才進的.
如果我爸爸繼續收錢.
他又去不了.
結果怎樣就是.
我在茶餐廳收錢.
就是我爸爸和我女朋友的報道會.
是不是很奇怪這件事.
這個組合.
他們都OK的朋友.
那次其實.
我想我爸爸都不知道這個故事.
他已經在聽.
那時其實是最後機會.
結果去完之後.
我都沒什麼盼望.
我都問.
怎樣呀今晚.
然後我女朋友說.
我信了.
不是吧.
原來我以前不信主的難處是我.

$^{1241}$行不行.
你們拍手的.
你們這麼愛心.
真是太離譜了.
都是梁牧師的朋友.
你們很愛心.
都是我行不行.
我問我女朋友.
你聽到什麼.
牧師說什麼這麼感動.
你信主嗎.
我和你說這麼多次都不信.
她說我不知道牧師說什麼.
不過她呼召說一件事.
就是如果今天就是世界末日.
愛你的基督徒的人.
就要和你永遠分開.
所以那一刻她很感動.
我不想和愛我的基督徒分開.
我就問是不是我.
你們又笑了.
太離譜了.
但是真的.
是永遠分開.
當然她現在很愛主.
但很多人信主都服侍.
這個是真實的.
弟子們.
我們都會面對完結的.
也都會到期的.
人是會過期的.
是會面對審判的.
我記得在新冠肺炎的時候之前說過一次.
就是有一次我在吃辣椒醬.
在吃蘿蔔糕.
但那些辣椒醬味怪怪的.
我就看一看辣椒醬.
嘩.
過了一期年紀.
我和我老婆不過來看一看.

$^{1281}$我有沒有看錯.
是不是過了一期年紀.
她說不是.
過了最佳食用日期年紀.
沒事的繼續吃吧.
有事的.
你去你冰箱看看.
你開一下.
我包你十枝八枝的過期.
有時過期你都不知道.
每次過年都是.
我們和很多過年禮品.
都去了Pantry給大家吃.
會過期的.
你看一看你旁邊那個.
他都會過期的.
不要那麼好意思.
不過都是真的.
即是什麼呢.
什麼是愛呢.
當我們對家人的愛.
和他玩去旅行吃東西.
看是對的.
是好的.
買號碼給他.
不過真正對他永恆的愛.
就是永恆的孝順.
將福音給他.
明白.
真的.
他怎樣不要都好.
你不斷在那兜著他.
挪著他.
五個月又見到我.
又說這件事了.
同時一樣.
你和他看球賽很好.
去旅行很好.
但是最好的友情.
是永恆的友情.

$^{1321}$明白.
真的.
街坊都是.
我們探訪他有月餅.
總之這個真是真實.
不過最重要的愛.
是永恆的盼望.
明白.
不可以沒有的.
我們都是愛他的.
在地上.
如果不是我們都愛他.
愛在地上都幾年而已.
但是很重要.
我們要見到他永恆.
所以最後這個秘笈.
第六.
是要迫切.
你知道人是會過期的.
隨時都會過期的.
當然.
如果我們中間有些正在考慮讀神學的.
大家多過第七.
剛才我說過.
就是不用怕.
神供應足夠.
白白的來白白的去.
讓我再選.
我告訴你一定選回現在.
所以Brian不用怕.
快要踏入了.
很開心的.
但是我告訴你.
裡面真的很多眼淚.
不要說到天堂一樣.
因為的確辛苦.
不過.
有主有對公.
我們裡面就經歷上帝的供應.
我們不是等候差遣.

$^{1361}$而是回應差遣.
神已經差我們了.
我們從心出發.
那個憐憫.
被主差遣.
我們真的被福音改變.
我們一定會.
成為別人.
遇見福音的原因.
我再講一次.
如果我們先被福音改變.
而是真正的.
我們必定會成為別人.
遇見福音的原因.
所以不要問你隔壁的我.
跟他說.
你是可以的.
最後我請靜脈隊來總結.
去到《馬太福音》28章.
亦都再提一次.
這個權柄他要快要走了.
他說天上地下所有的.
權柄都至給我.
所以你要去.
各位弟兄姊妹親愛的弟兄姊妹朋友.
我們信了主是有權柄的.
很多奇妙的事你想都沒想過.
醫都醫不到.
心理都幫你不了.
剛剛我家一個醫生.
一個心理學家.
他做事是有限的.
我們說包了.
屬靈上我來.
心理上我女.
身體上我仔.
但的確我告訴你很多東西.
醫學是做不到的.
真的.
直到世界的末了.

$^{1401}$最後這個重點.
假如基督是王.
所有的事情.
按字面所說的.
所有的人和事.
都需要重新想像.
你重新去突破你的想像就可以了.
重新調整.
重新轉向一個絕對順從.
基督的生活方式.
而我告訴你.
你這樣將會經歷一個.
奇妙的人生.
我們一起祈禱.
天王我們多謝你.
你讓我們.
站立了.
我們站立祈禱.
多謝你.
很積極.
我們站立再祈禱.
如果今天這篇經文.
你知道上去對你說.
你說主啊.
我求你答應我.
讓我見到你所見的.
讓我憐憫你所憐憫.
讓我們見到祂們背後內心的需要.
讓我們知道我們有權柄.
讓我們知道我們的迫切.
如果你是其中一個.
這一刻將你的手放在你的心.
我和你一起禱告.
每一次我都會.
因為我都需要主燃點我的眼光.
亦都藉著不同弟兄姊妹去服侍.
我見到主.
神今天我們禱告.
無論在這裡.
在網上的.

$^{1441}$我們求你都向我們說話.
讓我們見到你的真實.
你的恩典是夠用的.
在不同的狀況.
用不同的方法.
燃點我們.
改變我們.
讓我們更加像你.
用耶穌的生活方式.
去活出我們美好豐盛的一生.
環境是困難.
但沒有什麼比我們裡面那種的光難.
耶穌說祂就是世界的光.
願這個光光照我們.
直到我們見主面.
我們多謝你.
Hallelujah.
\newpage



\section{馬太福音 25:14-30}
\label{sec:032284}
\textbf{善用今生,拓展神國}~~~~李文康傳道
\newline
\newline
連結: \href{http://www.church.com.hk/acms/content.asp?site=cdc\&op=show\&type=product\&code=032284\&layout=sermon}{\texttt{032284}} ~~~~ 語音日期: 2026-06-06
\newline
\newline
\hyperref[sec:032283]{\small{< < < PREV SERMON < < <}}
~
\hyperref[sec:index_chronic]{\small{[返順時目]}}
~
\hyperref[sec:032285]{\small{> > > NEXT SERMON > > >}}
\newline
\newline
馬太福音 25:14-30
\newline
\begin{longtable}{cl}
\hline
\hline
章節 & 經文 (和合本修訂版)\\
\hline
25:14 & \begin{tabularx}{0.7\textwidth}{X} 「天國又好比一個人要出外遠行，就叫了僕人來，把他的家業交給他們。 \end{tabularx} \\ \\ \relax
25:15 & \begin{tabularx}{0.7\textwidth}{X} 他按著各人的才幹，給他們銀子：一個給了五千，一個給了二千，一個給了一千，就出外遠行去了。 \end{tabularx} \\ \\ \relax
25:16 & \begin{tabularx}{0.7\textwidth}{X} 那領五千的立刻拿去做買賣，另外賺了五千。 \end{tabularx} \\ \\ \relax
25:17 & \begin{tabularx}{0.7\textwidth}{X} 那領二千的也照樣另賺了二千。 \end{tabularx} \\ \\ \relax
25:18 & \begin{tabularx}{0.7\textwidth}{X} 但那領一千的去掘開地，把主人的銀子埋藏了。 \end{tabularx} \\ \\ \relax
25:19 & \begin{tabularx}{0.7\textwidth}{X} 過了許久，那些僕人的主人來了，和他們算賬。 \end{tabularx} \\ \\ \relax
25:20 & \begin{tabularx}{0.7\textwidth}{X} 那領五千的又帶著另外的五千來，說：『主啊，你交給我五千。請看，我又賺了五千。』 \end{tabularx} \\ \\ \relax
25:21 & \begin{tabularx}{0.7\textwidth}{X} 主人說：『好，你這又善良又忠心的僕人，你在少許的事上忠心，我要派你管理許多的事，進來享受你主人的快樂吧！』 \end{tabularx} \\ \\ \relax
25:22 & \begin{tabularx}{0.7\textwidth}{X} 那領二千的也進前來，說：『主啊，你交給我二千。請看，我又賺了二千。』 \end{tabularx} \\ \\ \relax
25:23 & \begin{tabularx}{0.7\textwidth}{X} 主人說：『好，你這又善良又忠心的僕人，你在少許的事上忠心，我要派你管理許多的事，進來享受你主人的快樂吧！』 \end{tabularx} \\ \\ \relax
25:24 & \begin{tabularx}{0.7\textwidth}{X} 那領一千的也進前來，說：『主啊，我知道你，你是個嚴厲的人：沒有種的地方也要收割，沒有播的地方也要收穫， \end{tabularx} \\ \\ \relax
25:25 & \begin{tabularx}{0.7\textwidth}{X} 我就害怕，去把你的一千銀子埋藏在地裡。請看，你的銀子在這裡。』 \end{tabularx} \\ \\ \relax
25:26 & \begin{tabularx}{0.7\textwidth}{X} 他的主人回答他說：『你這又惡又懶的僕人，你既知道我沒有種的地方也要收割，沒有播的地方也要收穫， \end{tabularx} \\ \\ \relax
25:27 & \begin{tabularx}{0.7\textwidth}{X} 就該把我的銀子放給兌換銀錢的人，到我來的時候可以連本帶利收回。 \end{tabularx} \\ \\ \relax
25:28 & \begin{tabularx}{0.7\textwidth}{X} 把他這一千奪過來，給那有一萬的。 \end{tabularx} \\ \\ \relax
25:29 & \begin{tabularx}{0.7\textwidth}{X} 因為凡有的，還要加給他，叫他有餘；沒有的，連他所有的也要奪過來。 \end{tabularx} \\ \\ \relax
25:30 & \begin{tabularx}{0.7\textwidth}{X} 把這無用的僕人丟在外面黑暗裡，在那裡他要哀哭切齒了。』」 \end{tabularx} \\ \\
[1ex]
\hline
\hline
\end{longtable}
$^{1}$冬季慢慢過去 天氣回暖 夏天終於要來了.
我和太太就有一個決定.
在家裡找一些夏季衣服出來穿 準備.
不知道大家有沒有相同的經驗.
有時衣服雜物太多.
就算家裡有衣櫃也好 可能也穿不完.
不知大家有什麼其他方法.
我們的方法是放一些衣服在行李箱上.
因為我家裡有很多行李箱.
特別是過去在泰國服侍生活七年的時間.
所以在泰國七年的東西就拿回香港.
有一個行李箱放在我們窗台最邊緣.
也有很多雜物壓著它.
所以平時我們比較少把它拿出來.
有一天我們終於決定把它拿出來.
一打開 我和太太都很驚訝.
為什麼呢?.
原來這個行李箱裡放了很多冬天可以穿的衣服.
我們一件一件拿出來.
我們都異口同聲地說.
原來這件衣服在這裡.
如果早前冷的時候可以拿出來穿.
那就最合適了.
誰知道整個冬季過去了.
其實我們也有其他厚衣服.
因為我們完全沒有.
也有長袖衣服.
但原來不止這麼少.
有很多可以穿的冬季衣服.
竟然埋藏在行李箱裡.
我們覺得太浪費了.
沒辦法 因為夏天來了.
沒理由拿出來穿.
接著我們收拾心情.
我有另一個念頭.
因為我們的睡床.
其實也有一個收納的空間.
我們的睡床的設計比較高身.
一邊放很多櫃桶可以放衣服.
另一邊打開床板可以放很多雜物.

$^{41}$因為床褥很大很重.
平常我們很少把床褥移開.
打開床板去看.
因為印象中有很多雜物放在裡面.
接著我們很辛苦地把很大的床褥移開.
接著打開床板.
大家猜猜我們看到什麼.
我們竟然看到一個很寬闊的收納空間.
當我們一直以為放了很多東西.
誰知道原來這個空間一直都浪費了.
我們和太太都彼此在笑.
當中我心裡浮現了兩句說話.
這兩句說話是什麼呢?.
首先第一句是「埋在地裡」.
接著第二句是「白佔土地」.
為什麼呢?.
因為很多冬天的衣服應該可以穿的.
誰知道塞在行李箱裡.
接著就好像埋在地裡一樣.
接著我們家裡一直有很多東西.
很費神怎麼可以把物件收納得好一點.
誰知道原來這個床褥裡面有這麼大的空間.
我們都沒有好好運用.
同樣好像白佔土地一樣浪費了.
其實當我們一些衣服沒有穿一季.
或者一些收納空間沒有運用.
其實都不算是很大問題.
我們有得補救.
做好一點就可以了.
但是如果這個問題.
是發生在我們的屬靈生命上.
神賦予了我們每一位神的兒女.
都給我們有很寶貴的屬靈資源.
我們沒有好好運用.
我們忽略了的話.
這就不是一個小問題了.
並且我們可能會不明白.
神在我們整個人生當中.
究竟有什麼計劃和旨意.
更甚者我們可能失去了.

$^{81}$能夠為神拓展他的國度的機會.
原來「埋在地裡」和「白佔土地」.
這兩句話都是出自主耶穌基督.
祂自己的一個比喻.
一個是說到僕人將銀子埋在地裡.
另一個是說到沒有花果樹.
祂沒有結果子.
祂們的共通點不是沒有資源.
而是沒有將這些資源好好運用出來.
生命結束果子.
沒有產生果效.
今天很想和大家分享的題目.
藉著耶穌基督所說的這兩個比喻.
我們分享一個題目就是.
善用今生 拓展神國.
盼望我們一起思想.
我們可以怎樣去明白神的心意.
讓我們今生能夠活得合神的心意.
讓我們一起同心再次禱告.
讓我們有去敬拜你的時間.
去親近你.
願意當我們每一次去研讀你的話.
親愛聖靈親自的引導.
教導我們.
幫助我們.
不單止明白.
並且能夠有能力.
有決心的去遵行主導.
願你親自的帶領.
禱告奉靠耶穌基督的名字而求.
阿們.
我們一起先看第一個比喻.
第一個比喻就是.
主人按著柴幹分給僕人有家業.
記載在《馬太福音》25章14-30節.
在這個比喻裡提到.
有一個主人要去遠行.
他將他的家業分給了他的僕人.
按著個人的柴幹.
分別有給五千的.

$^{121}$有給二千的.
也有給一千的銀子.
之後這個主人就離開了.
接著這個領五千銀子的僕人.
隨即立刻去做買賣.
他積極行動.
他沒有浪費時間.
接著他另外賺了五千.
另外領二千的僕人.
他沒有比較.
沒有說為什麼五千這麼多.
他自己也同樣地.
將這個二千拿出去運用.
結果也賺了二千.
接著唯獨那個領一千銀子的僕人.
因為他害怕.
他將銀子埋在地裡.
過了一段很長的時間.
最後主人回來了.
每個僕人都要向主人交帳.
這個僕人就將他原本.
從主人所得的一千銀子.
原封不動地交給主人.
主人回來的時候.
稱讚第一個和第二個僕人.
稱讚他們什麼呢?.
因為他們將銀子運用.
他說:你這個又良善又忠心的僕人.
當中善良有一個美好的意思.
指的是一個好的存心和好的動機.
去體貼主人的心意.
而忠心就是代表著.
有一個好的態度和行為的表現.
是可靠的.
是值得信賴和殷勤的.
雖然第一個和第二個僕人.
他們所領的銀子是不同的數目.
但在主人所得來的稱讚是一樣的.
主人說:你們在不多的事上有忠心.
我要派你管理許多的事.

$^{161}$並且可以盡來享受主人的快樂.
這個快樂我們可以理解為.
當主人回來的時候.
他所帶來的福氣.
可以和主人一起去同享他的滿足.
亦可以理解為彌賽亞賢侄的喜樂.
另一千銀子的這個僕人.
他究竟發生了什麼事呢?.
有什麼問題呢?.
原來他向主人辯護.
他說:「主啊!我知道你是一個嚴厲的人.
沒有種的地方你都要有收穫.
沒有播的地方你都要有收穫.
我就害怕你把你的一千銀子賣藏在地裡.
請看你的銀子在這裡」.
其實在正常的情況下.
耕作需要先有播種才可以有收成.
打和的時候.
先要去迎風讓穀糠可以去播散.
才可以剩下一些籽粒能夠聚集起來.
所以這句說話.
其實這個僕人暗示主人.
都沒有工作.
他是想有收穫.
但如果這個僕人很清楚自己的角色的時候.
他就會明白.
其實僕人正正就是要為主人去工作.
這是他的遺訓.
他的問題是.
主人是很嚴厲.
於是他說他害怕.
然後也不願意去承擔.
他將責任拋給主人.
用害怕來掩飾自己的不忠心.
他選擇了一個最安全的方法.
把銀子埋藏的時候.
不做就不錯了.
但是主人對他有什麼評價呢?.
我們看到原來主人對他有一個責備.
主人說.

$^{201}$「你這又惡又懶的僕人.
你既知道我嚴厲.
為什麼不把銀子放給兌換銀錢的人.
好叫我回來的時候.
可以廉本大利收回呢?」.
主人對他有一個要求.
最低的要求是.
你可以把銀錢放給銀行.
當我回來的時候.
就可以連息收回.
但是這個僕人沒有這樣做.
他連最基本的動作也沒有做.
這個僕人說主人是嚴厲的.
這個僕人對主人的看法.
究竟是對還是錯呢?.
其實當我們看到.
當主人稱讚.
第一個和第二個僕人的時候.
稱呼他們什麼呢?.
「又善良又忠心的僕人」.
所以僕人得到五千.
其實主人沒有要求他賺夠二萬元.
或者拿二千元的.
沒有要求他賺夠一萬元.
所以我們看到主人同樣稱讚這些僕人.
是善良又忠心.
如果這個主人.
好像這個僕人所說這麼嚴厲的話.
相信他又怎會能夠向僕人.
作出一些又良善又忠心的稱讚呢?.
「又善良又忠心」.
和在二十六節裡面.
「又惡又懶」.
成為了一個很強烈的對比.
「又惡又懶」這個形容.
凸顯了這個僕人.
他沒有一個全正確的心態.
去承擔主人所交託給他的.
他也沒有體貼到主人那份心意.
並且對主人有一個錯誤的認識.

$^{241}$這個僕人的問題就是說.
他並不是沒有能力.
他有能力.
但他選擇不行動.
並且透過害怕.
掩飾了他對主人的不忠心.
這個比喻帶給我們.
有什麼屬靈學習的地方呢?.
主人遠行.
其實代表著主耶穌基督.
祂完成了救恩之後.
祂回到天父那裡.
主人將銀子交託給僕人.
也預表著主耶穌.
祂將很獨特的資源.
交給我們每一位.
願意跟隨主,相信主的人.
而今天的銀子對我們來說.
我們相信不會單單只是金錢.
也包括到神所賜給我們.
今生裡面一切寶貴的資源.
可能是時間,恩賜,智慧.
一切的能力.
主人回來是會算帳的.
有一天我們要相信.
主耶穌基督會再一次回到世上.
每一個信徒都要向主交帳.
去表明我們怎樣去用我們的人生.
其實神是量過.
我們每一個人一生裡所應得的那一份.
我們不過是一個管家.
我們忠心地運用.
神所賦予給我們的一切.
我們並不需要跟其他的信徒比較.
正如聖經裡說.
多給誰就向誰多取.
多託誰就向誰多要.
主人將這個僕人僅有的這一千奪過來.
給那個有一萬的.
因為聖經裡這樣描述.

$^{281}$「因為凡有的還要加給他.
叫他有餘.
沒有的連他所有的也要奪過來.
將這個無用的僕人丟在外面.
黑暗裡在那裡他要哀哭切齒了」.
為何要將這一千交給那個有一萬的呢?.
其實當我們看到.
主人知道當有僕人願意善用.
去運用,願意發揮的時候.
神會更加大大地去加添給他.
但是沒有忠心地運用.
神所交付的僕人.
就要面對神的懲罰了.
在這個比喻裡.
我們其實學習到甚麼呢?.
我們要重申.
在這個比喻中.
並不是要教導我們如何好好學習投資.
也並不是鼓吹成功神學.
人努力工作.
用合法,合理的情況下.
去賺取金錢是沒有錯的.
但是這段經文的核心.
就是說我們願不願意.
忠心地將主人所交託給我們的.
我們去拿出來,去運用,去發揮.
零五千和零二千的僕人.
他們所得的數量是不同的.
但是得到主人的稱讚是完全一樣的.
所以我們看到主人所看重的.
並不是回報有多少.
那是對神要有忠心.
在這裡我們看到.
這個零一千的僕人.
他將銀子埋藏在地裡.
好像是寶本.
起碼不會虧損了.
但是卻埋沒了神給他的託付.
既然我們知道.
神所給我們的一切的資源,能力.

$^{321}$都是神賜予的.
我們可以放心.
我們的能力是神知道的.
祂不會要求我們.
做一些過於我們能力範圍所能及的工作.
親愛的弟兄姊妹.
今天不知道我們對神.
有沒有像這位零一千的僕人的心態.
對主人,對神有一個偏差的認知呢?.
我們如果產生了錯誤的認知的時候.
可能我們心裡也會產生恐懼.
我們會害怕.
但是求主去幫助我們.
我們知道神所託付給我們的.
是祂相信我們.
相信我們有能力去發揮,去運用.
雖然主今天好像遲遲還未回來.
但是在主說還未回來的這個時間.
這個空間裡.
神交託了給我們每一個弟兄姊妹.
都有很獨特的資源.
不單止是金錢.
我們擁有著我們人生裡的時間.
我們有不同的恩賜.
有些弟兄姊妹可能很容易.
與人打開話題去傳福音.
可能有人很容易去關心.
有人去奉獻.
有人在總務方面能夠有美好的發揮.
有人有能力,有憐憫的心腸.
可以來幫助很多人.
我們都可以在不同的恩賜上.
去發揮神所賜予的.
還有就是我們每一位弟兄姊妹.
其實我們都有很獨特的人生的經驗.
哪怕我們過去可能有愉快的.
甚至可能有艱難的.
一些痛苦的經歷也好.
只要我們繼續走在神的心意當中.
我們繼續衷心地跟隨主的時候.

$^{361}$我們深信我們所走過的每一段路.
神都可以去使用.
不單止去祝福我們.
也可以將來成為別人的祝福.
記得之前在工場的時候.
發覺很多弟兄姊妹很有心.
來到工場探訪宣教士.
探訪關愛使者.
他們都有帶著不同的技能,專長.
有些是懂得維修汽車的.
我最記得的是.
當我在工場開始第四年的事奉.
開始落實投入在教會很多的服事當中的時候.
要買一手二手車.
我自己看了很多次.
看了很多款式.
但因為是外行.
不知道這架車究竟有沒有水浸過.
有沒有一些大的問題.
不懂得.
誰知道有一位弟兄.
他願意恆常每一年都來到工場.
因為他很了解汽車的機械.
於是我們正正約了一個時間.
透過他來到的時候.
一起去看我準備想買的那一部二手車.
最終他看過後.
說應該沒有大問題.
應該是合格的.
於是我就很放心地買了那部車.
在餘下的年頭裡.
都不斷有好的運用.
甚至這位弟兄.
都不斷幫我們每一位宣教士的單位.
捐了車底.
去檢查我們不同的車.
究竟有什麼問題.
接著也記得.
有一些可能是建築師.
可能在工場上.

$^{401}$我們有一些不同的項目.
需要去興建.
要有不同的建築物.
有些建築師.
專長的弟兄姊妹.
有時可能駐守在香港的.
可能提供有不同的專業意見.
幫助在前線工場的弟兄姊妹.
同工可以怎樣去跟進.
怎樣可以做得更好.
有些可能是醫生,醫師.
來到工場幫我們關愛使者打脈.
或者幫我們成為一些顧問.
了解我們身體的情況.
很難忘有一次在斷宣隊裡面.
曾經有一個隊伍當中.
有髮型師.
也有美容師.
這個髮型師在一個主日的下午.
他就開檔.
連續四個小時.
幫我們不單止弟兄姊妹.
甚至有一些新來賓.
一些微信主義的朋友去剪頭髮.
可能因為是他的專長.
如果給我們剪四個小時一定很累.
但這位弟兄很樂意剪了很多頭髮.
但仍然沒有疲倦.
很樂意來服侍我們.
有一位姊妹是美容師.
她幫我們提出一個服侍.
可以有手部美容.
怎樣做呢?.
原來可以用一些很簡單的物料.
就是用一些砂糖和橄欖油.
和一些不同的物料.
混合在一起之後.
就可以在手上捉.
一邊捉一邊捉.
然後發現原來捉完之後.

$^{441}$真的好像變白了.
在一個聚會當中.
我們設計了一些家長.
有些可能是微信的.
有些可能是教會的弟兄姊妹.
來到一起在房間裡面.
大家都是姊妹.
然後彼此在捉.
捉著對方的手.
然後不斷地捉.
在捉的過程當中又可以分享.
這位美容師的姊妹.
在當中分享見證.
我們發覺有些微信的家長.
都很留心去聆聽.
從而建立了好感.
亦幫助他們慢慢地.
更加投入在教會的生活當中.
弟兄姊妹.
我們發覺原來我們所獨特的專長.
我們的背景不單止在本土教會.
本土的社群當中可以有發揮.
在不同的工場環境當中.
同樣的可以給神使用.
只要我們願意的時候.
神就能夠去使用.
我們深信今日我們一切所需要的.
神都會來賜予給我們.
弟兄姊妹.
不知道你有沒有更多的空間.
願意給神更加多的賜予給你呢?.
我們深信只要你願意有更多的.
放手給神的時候.
深信神是願意更加大大的.
來祝福你.
去加添你更多的恩賜.
更多的可能性來服侍主.
當我們在第一個比喻當中.
祂提醒我們.
不要將神所託付的埋藏起來.

$^{481}$在第二個比喻當中.
我們會見到.
其實這個比喻更加直接地問.
如果神來到我們當中.
在我們生命裡面的時候.
祂想找到些什麼呢?.
在第二個比喻.
《路加福音》第十三章六至九節.
當中這個比喻提到.
耶穌說.
有一個人在葡萄園裡栽種了一棵無花果樹.
他來找果子.
但是他找不到.
他就說.
看看.
我這三年來.
來到這棵樹面前去找果子.
但都竟然找不到.
不如將它砍了.
即是說將它砍了下來.
何必白占土地呢?.
元丁就求情.
就說.
多給他一年時間.
讓我可以去花開土.
去施肥.
如果再不結果子的話.
才將這棵無花果樹砍了下來.
這個比喻是記載在.
《路加福音》十三章一至五節之後.
一至五節發生了什麼事情呢?.
原來當時有些人.
去到耶穌面前.
去告訴耶穌.
彼拉多殺害了一些獻祭的加尼尼人.
亦有提及.
有一個樓房.
叫做西羅亞樓.
倒塌.
竟然在當中壓死了十八人.

$^{521}$當時很多人以為.
這些人的死亡.
是否因為他們犯罪.
而得到這樣的結果.
於是耶穌就回應了.
耶穌說.
「你們以為這些人比其他人更有罪嗎?」.
「不是的」.
「你們若不悔改.
都同樣要滅亡」.
接著耶穌就用無花果樹的比喻.
藉此帶出了悔改和審判.
要警告的訊息.
在這個比喻當中.
我們會知道.
原主是象徵著神.
而無花果樹是象徵著以色列民.
就是神所一直栽種,看顧的人.
耶穌藉此要提醒,警告以色列民.
神已經揀選了他們.
栽種他們.
賜福給他們.
他們多年來都是蒙恩的.
神賦予他們一個獨特的位置.
要他們去結果止.
但可惜他們屢次背棄了神.
現在又不接受耶穌基督.
就是神所差來的這位彌賽亞.
如果他們再不悔改的話.
就將要面對神的審判.
在這個比喻當中.
重點是同時彰顯了.
神的寬容和審判.
神期待贖祂的人.
是需要結果止的.
即是說.
生命要有回應.
要有改變.
元主並不是沒有耐性.
其實當元丁去求情的時候.

$^{561}$「今年且留著.
讓我翻開肚,施肥.
一年之後.
看看祂會不會結果止」.
原來這也顯示出.
神是有恩典.
神願意去給予機會.
但再等一年.
這個恩典.
不可以成為我們再拖延的藉口.
乃是提醒我們.
恩典是有期限的.
人是要把握機會.
不要白佔土地.
要趕快去結果止.
今天神也同樣藉著這個比喻.
去呼召我們每一位神的兒女.
我們要把握神賜給我們的人生.
我們要為主去結出果子.
究竟神要在我們身上.
找到甚麼果子呢?.
我們深信可以包含三方面.
首先第一就是一個.
真實悔改行動的果子.
神可能不是問我們有沒有信主.
而是問我們有沒有真實的悔改.
有沒有回轉,歸向給神.
生命有沒有常常在神面前有反思.
我最記得當我信耶穌的時候.
最初信主那一刻.
當然要有悔改認罪.
但之後最深刻的悔改是甚麼時候呢?.
當我二十多歲的時候.
當時因為年輕力壯.
對人生很有抱負.
很有期望,很有計劃.
當時我做印刷.
轉換不同的公司.
轉換了不同的公司的時候.
希望學習不同的新領域,範疇.

$^{601}$於是希望不斷不斷提升自己.
最終能否做到經理呢?.
會否將來回國內發展.
甚麼時候可以買樓,買車,娶老婆呢?.
人生有很多自己的打算.
但是在一直做的過程當中.
其實也發現了自己放了很多時間在工作上.
甚至影響了自己的教會生活.
接著價值觀就慢慢改變了.
這些都是人之常情.
「我年輕時不博,何時博呢?」.
但是發覺在過程當中.
其實慢慢自己所失去的.
是失去了很多.
不單止是回教會的時間.
或者沒有回團契.
失去了甚至是和家人相處的時間.
有時做到OT很晚.
回到家裡.
只見到家人打聲招呼.
接著家人就睡覺了.
自己接著才吃飯.
在過程當中.
神給我一個很大的提醒.
出現了一個突發的事件.
有一天的晚上.
我自己發現了一個很不舒服的情況.
長話短說.
是一個危疾.
要即時送到急症室.
透過這個病的過程當中.
我發現原來一下子.
一隻不知道甚麼名字的細菌.
已經可以將整個人打倒.
要躺在病床上.
並且不知道自己是否能康復.
透過這個危疾.
我在過程當中.
再一次將自己的生命重新交給神.
神說:我將生命交給你.

$^{641}$你才是主人.
之後透過悔改.
慢慢地回轉.
回到神身邊.
我知道自己過去是去掌舵.
但原來耶穌基督才是人生的主人.
在祂的引導當中.
人生才是最寶貴,最穩妥的.
所以弟兄姊妹都彼此勉勵.
我們不單只是在信主的那一刻.
我們需要向神悔改認罪.
可能整個人生當中.
我們都要隨時檢視自己.
是否偏離了神的方向.
是否需要在主的面前.
去結出悔改的果子.
接著第二個果子就是生命的果子.
因為我們每位弟兄姊妹.
信主之後都有聖靈內在我們裡面.
神會藉著聖靈.
提醒,引導和幫助我們.
去結出一些屬靈的果子.
大家記不記得屬靈的果子有什麼?.
會有人愛,喜樂,恩慈,良善.
信實,溫柔,節制.
這些不單是聖經裡面的經文.
而是可以活生生地.
出現在我們的生命當中.
只要我們願意將生命交出來.
我們深信神會塗做我們.
不知道我們的家人,同事.
會否在我們身上看到基督的香氣呢?.
最後我們會深信.
神要我們去結出使命的果子.
神揀選以色列民.
其實都很想藉著他們去祝福萬民.
所以我們深信今天.
神都期待在我們的生命當中.
我們都能夠為主作見證.
無論是遠方或是近處.

$^{681}$在家裡面.
在我們家樓下.
在教會,職場,社區.
其實我們隨時都可以為主耶穌.
作美好的見證.
和大家認識一位弟兄.
他的生命的見證.
這位弟兄在一間茶餐廳工作.
他也分享到他作為侍應.
在一間茶餐廳工作.
有時候都很艱難.
為什麼呢?.
原來在廚房裡面.
有幾位同事.
他們的脾氣都很差.
不知道是否因為做廚房很熱.
而影響情緒.
接著這幾位廚師.
其實都很無理地整天罵人.
在這間餐廳裡面.
很多同事都被他們罵.
甚至被他們罵到哭.
這位弟兄很坦白地說.
他其實是一個成年人.
一個大男士.
同樣都試過被這幾位廚師罵到哭.
在這個時候.
他怎樣回應呢?.
他說.
其實當然會很不開心.
但是他在心態裡面說了一句話.
他說他更加看到.
這個艱難的工作場景.
更加是能夠為基督作見證的機會.
於是他沒有用同樣的咒罵方式.
去回應對方.
或者和他理論.
大家不要誤會.
不要誤會以為這位弟兄是懦弱.
原來這位弟兄在未信主之前.

$^{721}$同樣性情都很暴躁.
甚至容易和人有衝突.
試過有衝突的時候.
握起對方的胸口.
很想揮拳去打架.
但是他信主之後.
悔改.
生命得到改變.
所以他不用別人的方式去回應.
因為他深信.
雖然這個人暴躁.
但是在整個餐廳裡面.
每一位同事.
每一位茶客.
他們的眼睛都是雪亮的.
一定會知道誰對誰錯.
這就是他的心態.
於是他就用這個心態.
繼續忠心地去工作.
很感恩.
因為他這個與眾不同的生命表現.
吸引了一位同事.
主動有一天告訴他.
「我也很想跟你回教會」.
這個見證真是令這位弟兄很開心.
在這樣的氣氛下.
一個一個同事.
慢慢地一個一個跟這位弟兄回教會.
在當中明白福音.
領受信仰.
神交託給我們.
原來並不是單單一個.
在教會裡面服侍的崗位.
一個角色.
原來神所交託的是我們整個人生.
是在我們工作的場景.
生活的場景裡面.
原來我們都可以在當中事奉神.
都可以在當中為主作美好的見證.
弟兄姊妹.

$^{761}$在剛才我們.
無論是埋在地裡.
或者是白佔土地的這個比喻當中.
其實主耶穌基督.
都是在問同一個問題.
神給你的 你有沒有忠心耿去運用出來呢?.
有沒有願意為主去作結果子呢?.
當將兩個比喻放在一起.
我們嘗試一個整合.
我們會發覺.
僕人的比喻是神託付.
要我們去使用.
在無花果樹的比喻當中.
我們見到神給生命.
是要我們去結出果子.
神給我們有今生.
並不是單單要求我們.
好好地去過每天的生活.
神更加希望我們.
善用我們每一刻的時間.
去善用我們每一份很獨特的關係.
機會整過人生.
去回應主.
去榮耀神.
去拓展神的國度.
我們可以有一個生活的實踐.
也許我們不會一下子.
需要做一個很大的事情.
我們可以從一小步做起.
在這個星期裡.
我鼓勵大家一起去做.
看看我們會否在這個星期當中.
很認真地撥出一個時間.
為四個世界.
為神的國度去禱告.
盼望世界上更多的人.
更多未認識主的人.
可以有機會認識主.
我們可以在這個星期當中.
更加主動地關心我們一位未信的朋友.

$^{801}$或者在我們家庭裡面.
如果有時有誤會,衝突.
我們會否先做一個主動.
去調解,去包容,去寬恕呢?.
在我們教會當中.
我們也會見到很多年長的弟兄姊妹.
特別是很多的成人.
當然在我們當中有青少年.
但在當中我們成人的弟兄姊妹.
其實我們也要小心.
我們不要以為.
當我們慢慢年紀越來越大的時候.
會覺得我們也要退下來.
留給年輕人去做.
我們深信神使用我們每一位.
無論我們是什麼年紀也好.
神所體重的就是我們那份忠心.
年長的弟兄姊妹.
可以繼續做很穩定的代禱者,陪伴者.
做忠心的奉獻者.
甚至是支持者.
我們可以鼓勵我們的下一代.
去回應神的呼召.
在過去我們看到教會當中.
有很多年長的弟兄姊妹.
他們很用心去服侍.
很努力去事奉神,服侍人.
這也是我們值得去學校的地方.
有些是去探訪的.
有些是去傳福音的.
有些在教會裡當中.
有不同的事奉崗位.
都是默默地為主去作供.
沒錯,我們教會需要有接棒.
我們需要去傳承.
但我們更多需要彼此有合作的機會.
以致我們年長的,我們青年的.
大家能夠彼此交流.
能夠發揮美好的果效.
我們可以有更多的機會去嘗試.

$^{841}$更多的包容,更多的發展.
這個願意我們一起去學習.
當我們願意這樣去回應的時候.
就會發現.
善用今生並不只是單位為了自己.
而是為了神的國度.
為了下一代.
為了未得之民.
當我們說要拓展神國的時候.
我們最後一定不能夠離開.
回應主的大使命.
主耶穌基督在大使命裡面說.
我們要去使萬民作主的門徒.
這個呼召並不單單只是對宣教士.
而是對每一位門徒,作主門徒的人.
都要去回應主親自的呼召.
我們都是在主的差遣當中.
而在羅馬書裡面同樣是教導我們.
如果我們沒有差遣的.
怎樣能夠有去的呢?.
沒有去傳的.
怎樣能夠有人可以聽到福音呢?.
沒有聽到福音.
怎樣能夠信主呢?.
所以差傳並不只是單單宣教士的事.
而是教會整體的群體.
亦都需要去參與的.
我們可以參與的.
如果在我們在座當中.
已經很清晰神的呼召.
準備自己醞釀.
將來要去踏上宣教的時候.
我們就可以憑著這份感動.
不斷地禱告,等候.
好好地來裝備自己.
給神將來的人使用.
如果我們知道自己.
我們也不覺得神是呼召我們去遠方的.
其實我們也可以關注在本土裡面.
也可以關心幫助一些來港的.

$^{881}$一些不同種族的人士.
在我們香港社會當中.
亦都有很多不同種族的群體.
來到我們香港裡面生活.
我們可以繼續以生命的見證.
來感化他們.
剛才說到這位在餐廳裡面工作的弟兄.
他分享在三年裡面.
他帶領了八個人回教會.
這八個人裡面.
有他在茶餐廳裡面的同事.
有幾位是茶客.
分別在回到教會的過程當中.
都能夠信主了.
他傳福音的方法.
原來是與人聊天.
跟著與人分享他的得救見證.
最重要的是在他職場裡面.
去活出一個生命的見證.
他分享到曾經有一位茶客.
常常來餐廳吃東西.
但原來他患了癌症.
有一天這位弟兄主動邀請他回教會.
但這位茶客竟然怎樣回應呢?.
用粗言穢語去回應他.
但是沒有打垮這位弟兄的信心.
這位弟兄繼續努力地關心他.
最後真的帶領這位茶客回到教會.
並且不單止信主.
這位患癌病的茶客.
也帶領了其他的朋友回到教會.
甚至這位弟兄.
他不單單在茶餐廳裡面傳福音.
他分享到就算他去釣魚.
釣魚可能是弟兄姊妹的嗜好.
在釣魚的時候.
原來都認識到別人.
真的帶了兩個女孩回到教會.
並且信了主.
釣一些在海裡的魚.

$^{921}$最後釣到別人回到教會.
甚至這位弟兄.
他回鄉參加一些喪禮的時候.
都跟一些親人傳福音.
同樣帶領到一些親人信主.
所以我們看到.
原來他的傳福音.
他為神作供.
是融入在他的生命裡面.
他並不只是在什麼時刻.
什麼崗位.
什麼角色上面.
才知道自己是傳福音.
原來在他的生命裡面.
他都願意隨時去體貼神的心意.
記掛著神的國度.
所以弟兄姊妹.
我們每位弟兄姊妹.
同樣都是在神領受了呼召.
領受了生命.
我們同樣可以在主面前一起努力.
我們都深信神會使用我們每一位.
只要我們願意.
讓我們一起低頭去禱告.
我們有獨特的屬靈的恩賜.
有不同的裝備.
有不同的背景.
亦都有很獨特的.
不同的人際的關係.
我們深信這些都不是偶然.
都是神要使用我們.
要成為你的祝福途徑.
所以主啊.
在你面前.
我們就將我們自己的心交託給你.
如果在我們裡面.
過去曾經有膽怯的.
神啊 求你將這份膽怯挪開.
讓我們有信心地踏前一步.
去行在神的心意裡面.

$^{961}$神啊 如果過去我們有不同資源的時候.
求你聖靈親自感動我們.
以致我們知道有沒有哪些資源.
我們是埋藏了的.
是我們沒有發現的.
我們能夠拿出來.
讓神使用.
用在我們的國度當中.
神啊 求你引導我們.
幫助我們.
主啊 亦都懇求你.
更加開拓我們的空間.
我們知道.
只要我們願意的時候.
我們所領受的.
不單止是這麼少的.
是可以有更多的廣闊的機會.
更多的能力.
更多的機會可以給神使用.
神啊 求你親自臨在.
去祝福我們每一位.
以致我們能夠將這個祝福.
傳揚開去.
帶給更多的人.
成為你祝福的器皿.
求你親自使用和幫助.
我們多謝你.
誰聽我們在你面前同心禱告.
是奉靠主耶穌基督的名字而求.
阿們.
\newpage



\section{使徒行傳 2:37-47}
\label{sec:032285}
\textbf{屬乎耶穌基督的群體}~~~~翁曉怡傳道
\newline
\newline
連結: \href{http://www.church.com.hk/acms/content.asp?site=cdc\&op=show\&type=product\&code=032285\&layout=sermon}{\texttt{032285}} ~~~~ 語音日期: 2026-06-07
\newline
\newline
\hyperref[sec:032284]{\small{< < < PREV SERMON < < <}}
~
\hyperref[sec:index_chronic]{\small{[返順時目]}}
~
\hyperref[sec:032286]{\small{> > > NEXT SERMON > > >}}
\newline
\newline
使徒行傳 2:37-47
\newline
\begin{longtable}{cl}
\hline
\hline
章節 & 經文 (和合本修訂版)\\
\hline
2:37 & \begin{tabularx}{0.7\textwidth}{X} 眾人聽見這話，覺得扎心，就對彼得和其餘的使徒說：「諸位弟兄，我們該怎樣做呢？」 \end{tabularx} \\ \\ \relax
2:38 & \begin{tabularx}{0.7\textwidth}{X} 彼得對他們說：「你們各人要悔改，奉耶穌基督的名受洗，使你們的罪得赦免，就會領受所賜的聖靈。 \end{tabularx} \\ \\ \relax
2:39 & \begin{tabularx}{0.7\textwidth}{X} 因為這應許是給你們和你們的兒女，並一切在遠方的人，就是給所有主—我們的神所召來的人。」 \end{tabularx} \\ \\ \relax
2:40 & \begin{tabularx}{0.7\textwidth}{X} 彼得還用更多別的話作見證，勸勉他們說：「你們當救自己脫離這彎曲的世代。」 \end{tabularx} \\ \\ \relax
2:41 & \begin{tabularx}{0.7\textwidth}{X} 於是領受他話的人，都受了洗；那一天，門徒約添了三千人。 \end{tabularx} \\ \\ \relax
2:42 & \begin{tabularx}{0.7\textwidth}{X} 他們都專注於使徒的教導和彼此的團契，擘餅和祈禱。 \end{tabularx} \\ \\ \relax
2:43 & \begin{tabularx}{0.7\textwidth}{X} 眾人都心存敬畏；使徒們又行了許多奇事神蹟。 \end{tabularx} \\ \\ \relax
2:44 & \begin{tabularx}{0.7\textwidth}{X} 信的人都聚在一處，凡物公用， \end{tabularx} \\ \\ \relax
2:45 & \begin{tabularx}{0.7\textwidth}{X} 又賣了田產和家業，照每一個人所需要的分給他們。 \end{tabularx} \\ \\ \relax
2:46 & \begin{tabularx}{0.7\textwidth}{X} 他們天天同心合意恆切地在聖殿裡敬拜，且在家中擘餅，存著歡喜坦誠的心用飯， \end{tabularx} \\ \\ \relax
2:47 & \begin{tabularx}{0.7\textwidth}{X} 讚美神，得全體百姓的喜愛。主將得救的人天天加給他們。 \end{tabularx} \\ \\
[1ex]
\hline
\hline
\end{longtable}
$^{1}$聖嬰姐妹平安.
我們一起祈禱.
教會的主.
你是次下聖靈的主.
都可以被聖靈充滿的時代.
由揀選我們.
成為有份見證你的群體.
願主幫助我們在今天的訊息裡.
再次被聖靈觸動.
以致我們可以勇敢闖開自己的生命.
被聖靈繼續透過我們的敬拜.
禱告,團契,經歷裡面.
塑造我們成為流露基督同在的群體.
我們的同心禱告.
奉主耶穌基督的聖名祈求.
阿們.
聖嬰姐妹.
不知道大家有沒有在社交媒體.
或者地鐵的廣告上面.
有聽聞過.
或者你也有玩一個遊戲叫皮克文呢.
好像沒有什麼.
因為沒有人見過這樣的東西.
總之地鐵可能會見到.
你們下次留意一下吧.
根據官方的說法.
這一堆的生物是半生物半植物.
在皮克文這個遊戲裡面.
它會陪著你去走路和探險.
會去收集果實,種花.
在累積足夠的友好度的時候.
它們也會回到自己的出生地.
去得到一些裝飾.
這個遊戲透過不同的設計.
想鼓勵玩的人多走路,多走動.
在遊戲裡面有一些可以個人努力的部分.
例如你每天走多少路.
你去了多少地方.
你怎樣跟自己的皮克文互動.
你有沒有餵它們吃東西.

$^{41}$或者去收集裝飾,活動徽章等等.
同時也有一些有社群性的元素.
例如你可以寄明信片給你的朋友.
你可以跟你的朋友組成五個人一隊.
去每個星期一起走十萬步.
或者在遊戲裡面的戰鬥的時候.
你就呼叫你的朋友幫忙.
我本身兩年前已經有下載這個遊戲.
在我的電話裡面.
但就是一種有空拿出來玩一下.
有空就沒有玩的一個狀態.
當時我身邊也不多人有玩這個遊戲.
但是去到今年二,三月的時候.
就突然由台灣開始掀起了一輪皮克文的風潮.
就熱到來香港.
突然有一次在我服侍的初職團體裡面.
聽到有些人聊天的時候提起皮克文這三個字.
我就衝進去加入這個對話.
然後我們就聯繫了.
我們開始在皮克文裡面做朋友.
開始叫其他人一起下載.
我們就可以一起走路.
當我們越來越多人玩的時候.
我們聽到或者自己也有經歷到一些生活模式上的改變.
有些人去上班放午餐的時候.
他不再留在座位上看YouTube或者睡覺.
他會選擇出去散步.
有些人在回家之前會選擇去繞路.
就是為了多走一些路.
我也有一個朋友.
本身有一些狀況.
不是很喜歡出門口.
又很怕一些人多的地方.
但他為了玩皮克文去儲這些裝飾.
他得到足夠的力量去過海.
是很奇妙的.
在社交媒體也看到有些人邀請了他的父母一起玩.
因此跟父母多了對話互動.
以及共同用的術語.
這個遊戲令我發現.

$^{81}$原來皮克文沒有說不可以自己玩.
但原來有人一起玩的體驗是完整得多.
有趣得多.
而且也是這樣一起玩.
才帶來真正對日常生活有很大幅度的改變.
在這幾個月裡經歷很多人一起玩.
很多人一起研究不同的攻略的時候.
也有不斷想起我們的信仰.
我覺得我們的信仰帶來的改變也是這樣.
我們信仰的群體也一樣這麼重要.
我們今天說的「士徒人轉移章」37-47節.
是說五旬節後的群體如何形成.
他們又是有什麼樣的面貌.
之前已經說了.
門徒一起禱告等候聖靈.
聖靈降臨.
他們旁邊的人的反應.
到今天這個部分就是五旬節.
聖靈降臨的故事發展到尾聲.
一個被聖靈充滿和改變的群體就形成了.
我們先來一些上回提要.
在耶穌設立聖餐的余節之後.
第五十天來到五旬節.
在這五十天裡.
跟隨耶穌的信徒已經經歷了.
耶穌被捕,受苦,受死,復活,升天.
到他們現在沒有了耶穌.
跟他們肉身上同在的時候.
他們等候聖靈.
當他們聚在一起祈禱.
然後聖靈降臨的時候.
每個人都欣著聖靈充滿.
而可以說起別國的話.
他們就用別國的話來講論神的作為.
而有一群被聖靈降臨的時候.
那個大響聖所吸引來聚集的猶太人.
見到他們這樣的情景既疑惑又驚奇.
有些人選擇去笑他們只是聚酒.
彼得就發表了《士徒行傳》裡的第一篇《道》.
透過去詮釋約珥書和詩篇的預言.

$^{121}$去解釋聖靈降臨是什麼一回事.
並且用這件事去論述.
拿撒勒人耶穌就是猶太人.
由先祖以來一直等候來臨的那個尼賽亞.
縱然猶太人將耶穌釘在十字架上.
但耶穌的死而復活.
卻顯明了祂是救主,是基督.
到37節那些因著想守節旗而來到耶路撒冷.
去到聖殿朝聖的經乾猶太人.
在聽到彼得的訊息後覺得「宅心」.
就問「士徒」他們應該做些什麼.
「宅心」的原文直譯為「刺透心臟」的用語.
是用來表達一種深度的憂傷.
我覺得中文是絕對足夠博大精深去譯得很生動.
「宅心」可以表達那種聖靈使用彼得的訊息.
直擊了這班敬虔猶太人的心.
以至他們進入了一份認訊.
相信耶穌是救主.
相信自己的罪需要耶穌去拯救.
而他們會急切渴望自己應該怎樣做.
唯有有一份這樣的認訊.
才可以帶來真正的悔改和轉化.
我們現在經常說「願聖靈在人心動工」.
剛剛報道會前我們不斷祈禱時.
也會說「願聖靈在人心動工」.
讓人「宅心」是聖靈其中一種具體的工作.
當人在聖靈的工作下穿透自我和驕傲.
而深刻體認到自己的罪.
為著自己遠離神,背叛神,傷害神而感到悲痛.
才會由心而發地渴望做一些事來改變局面.
弟兄姊妹,你還記不記得自己初次感到「宅心」時?.
當初為著自己的罪感到哀傷.
又很感激上主的救恩感受.
有沒有仍然推動著你將生命獻給神?.
你信主後到最近還有沒有經歷「宅心」?.
我是一個信二代.
所以沒有很記得有「宅心」而信主的畫面.
但卻很記得因若「宅心」而蒙回應呼召的畫面.
當時是我中二升中三的暑期營會.
當中的詩歌及訊息令我領會到神有很大的愛.

$^{161}$以及祂很體重我的生命和價值.
但我覺得自己實在不夠好.
然後就覺得很「宅心」.
以致講完去呼召人將來全職事奉時.
我找不到理由跟聖靈說「不如不要吧」.
當日這一份對神的愛很廣大.
而我無以為報的深刻體會.
成為了我日後在軟弱,掙扎,不太想做的感受中.
一次次選擇信服的原因.
是當初的那份「宅心」令我日後也有其他的「宅心」.
亦是這些「宅心」令我,令很多其他人.
或曾經也讓你在生命中作出不一樣的選擇.
彼得以一個悔改的呼召回應這群人的「宅心」.
他叫他們各人要悔改,奉耶穌基督的名受洗,使你們的罪得赦免,就會領受所賜的聖靈.
要相信並承認人只能透過耶穌基督得潔淨和赦罪.
就可以領受耶穌應許此下的聖靈,住在人的心中,只因人按著神的心意而活.
這對當時的群體應該頗為重要,因為他們已不再有肉身的耶穌與他們同在.
但當人奉耶穌基督的名受洗,領受聖靈的時候.
仍然是應驗間施洗約翰所說的.
「耶穌要以聖靈和火吸引施洗」的語言.
復活後的耶穌成為了一個不可見的施洗者.
而這句應許是「給你們和你們的兒女,並一切在遠方的人,就是給所有住於我們的神所召來的人」.
五重節的事件標誌著一個新的時代開始.
是所有相信耶穌的人都會領受聖靈,有聖靈內著的時代.
五重節前和耶穌前,其實舊約開始一直有記載有人被聖靈充滿.
但往往那些只是少數人的經歷,並且是單次發生.
但由五重節開始,聖靈不再僅僅充滿某些人.
而是更加普遍和廣泛地持續充滿每一個願意悔改,相信耶穌是救主的人.
聖靈持續和信徒同在,幫助教會成為耶穌復活的一個有力和堅韌的見證.
雖然我們經常這樣說,但不知大家在你的生活中有沒有深刻地相信和經歷到聖靈持續同在呢?.
之前有和一些弟兄姊妹討論的時候.
剛說到想像一下耶穌在進入曠野時被聖靈充滿.
我們就想大家想像一下被聖靈充滿的狀態是怎樣.
有很多人不約而同地戴著頭盔表示.
「我不敢說自己有經歷過聖靈充滿,但我覺得耶穌那時應該怎樣怎樣吧」.
那次之後我回想了很久,為何我們一直這樣教導每個人都有聖靈,但我們又好像不太感受到呢?.
我們怎樣才會覺得有聖靈同在呢?.
我們會否感受到現在是靈命高峰,世界充滿盼望,才會有聖靈同在呢?.
其實在我們每日的生活中,我們仍然找到原因相信上主,堅持我們的信仰.
我們見到有需要的人時有憐憫的心,我們在生活中見到上帝的恩典.

$^{201}$可能這些都是聖靈一直充滿著我們,與我們同在的一些記號.
希望我們可以確信我們有聖靈同在的一群人.
並且在我們生活中也能見到聖靈與我們同在的一些記號.
40節給予血面人,他們要救自己脫離彎曲的世代,我覺得很有意思.
因為他說的是「耶穌成為了肉身來到世上,揹起所有人的罪,死而復活,升埋天」.
這個世代仍然是彎曲的,仍然是羅馬帝國統領的世界.
猶太人仍然有很多散居,但當彼得延續耶穌的教導時.
救自己脫離這個世代,應該不是當時很多猶太人所期盼的.
有一個人帶來強大的戰力幫他們復國.
彼得更有可能是透過悔改受洗,歸入耶穌基督的名下.
以一個新的世界觀去生活,在真理中重新得著自己的生命.
於是雖然帝國可以以權勢暫時統治列國,令列國伏在他的腳下.
甚至嘗試控制人民的語言及文化,但聖靈的福音才是真正可以改變人的心靈及頭腦.
才是到終末終極得勝的那一位.
當我們帶著這樣的認信時,信徒雖然仍然活在彎曲的世代.
但在價值上,生活實踐上,我們都可以救自己脫離這個彎曲的世代.
我們可以有另一份信念,有另一個世界觀去抗衡這些文化.
這種可以脫離或用另一種信念去抗衡的可能性.
今日仍然有透過聖靈賦予我們.
無論你看這個世界時,特別見到甚麼問題.
可能是消費主義,可能是社會上資源財富分配不均.
可能是某些位置的權力不受制衡.
可能是一些人濫權,一些人自我中心等等.
我們見到這個世代和二千多年前一樣彎曲.
但我們因著聖靈,都有一切努力去抗衡的可能性.
我們都會繼續思考,究竟我們信仰天國的世界觀.
怎樣幫助我們去看我們生活的世代,我們應該怎樣去回應.
回到當時,那些責心的人去問自己可以做甚麼.
彼得叫他們「悔改受洗」,他們就這樣做了.
當日這樣悔改受洗的人,約有三千.
路加沒有描述很多這班人的背景.
亦沒有分析究竟散居在甚麼地方的猶太人特別容易信主.
路加的焦點在於五旬節後這群群體的面貌特徵.
他們專注於使徒的教導和彼此的團契,抹餅和祈禱.
這群群體專注在四件事.
第一是聽使徒去複述和解釋耶穌的教訓.
學習去活出一個悔改的生命.
第二是彼此的團契.
團契這字帶有緊密連結彼此分享的意思.
福音書沒有用過,是去到《使徒行傳》這個位才第一次用.

$^{241}$顯示了耶穌的復活,升天和聖靈降臨.
為教會帶出一種互相參與在大家生活中的分享.
一種新的群體連結.
第三是抹餅.
抹餅本身在猶太人用餐禮儀中也有.
但因著基督設立了聖餐.
抹餅本身和教會群體一起抹餅這個動作的意思重要了許多.
而一起抹餅除了一起紀念設立聖餐的主.
記著主先是現職的主人及我們這個群體是主復活的見證外.
當時抹餅也很有可能在他們正式吃飯前進行.
每個人都可以得到一份聖餐.
也展現出當時他們不需要受自己有錢沒錢.
有多少資源影響他們有沒有東西吃.
每人都可以得到一份.
展現出被聖靈塑造的教會有慷慨及群體的特徵.
有份開始去應驗天國的價值及盼望.
第四是42節最後說到的禱告.
禱告是一個眾訴.
禱告不是單次發生的行動 而是教會群體持續的敬虔生活.
至43節有提到歧視神蹟.
最大的作用應該是預告之後會發生的神蹟故事.
同時也以一種和耶穌平衡的敘事.
象徵著教會作為耶穌在地上的代表延續耶穌當初的工作.
當這群會眾專注在聽道,團契分享,抹餅及祈禱中.
他們就成為一群心存敬畏的群體.
而這些心存敬畏的人 在五旬節當日的三千人當中.
應該有一部份不會留在耶路撒冷.
因為他們本身是散居各地的猶太人.
他們住在當時形容的天下各國.
與港地極差不多的描述.
所以這群人帶著他們的看見,被更新的信仰.
以及這四個群體生活的焦點.
帶回他們所散居的地方.
以至在使徒將福音傳到地極之前.
神在五旬節已經揀選了一批在地極來的人.
經歷悔改,更新.
再回去成為耶穌是復活主的見證人.
所以我們見到由使徒首先成為主的見證.
到其他信主的人見到他們的見證後.
都成為了聖靈能力的見證人和傳播者.

$^{281}$按著耶穌所說的大使命.
由聖靈降臨開始一路應驗著.
44至47節更加仔細地描述他們的群體生活.
信的人聚在一處繁物公用.
又賣了田產和家業.
照每一個人所需要的分給他們.
路加描繪了一個很街尾的教會群體畫面.
本來在社會中分配不均的資源和財富.
在這個教會群體中得以合宜地再分配.
以致有錢的人不需要為他的財富而自傲.
而本身貧困的人也不需要為他的缺乏而羞恥.
在教會群體中.
所有人都是因著聖靈的引導而聚集在這裡.
他們無論是吃飯的時候還是其他相處互動.
都不像當時的羅馬人這樣以社會階級主導.
他們在聽道,團契,抹餅,禱告中.
就在聖靈的幫助下脫離羅馬社會對人區分的淚鏡.
而轉向用救主的眼光看重每一個人的生命.
見到每一個人都是同等地被上主看重.
因為他們都是同樣被神以愛創造的人.
當一群被聖靈充滿的信徒不分階級彼此相助.
人與神與人的關係可以得以恢復的福音.
首先展現在教會群體中.
見證著社會和宗末的盼望.
今天我們身處的香港社會和羅馬社會都有點相似.
都會因應那個人的成就所掌握的資源.
或者可以帶來的利益.
將人分門別類或暗暗地分等級.
希望我們都可以在聽道,團契,抹餅,禱告中.
重新校正我們的眼光.
見到每一個人都是同等地被上主所看重.
成為社會和宗末的盼望的一種見證.
46至47節上.
講到這群信徒天天同心合意.
行徹地在聖殿裡敬拜.
且在家中抹餅.
全著歡喜坦誠的心用飯讚美神.
敬拜,抹餅,分享用飯,讚美.
都是一些現在式的分詞.
表示這些是持續地發生的行動.

$^{321}$成為了這群教會群體的日常生活方式.
所以他們的日常生活包括會公開到聖殿去敬拜神.
也包括私底下會在家中去敬拜神.
也包括他們會一起去抹餅.
在食飯中的平等相待.
都成為了他們的日常生活.
耶穌基督的「道成肉身」.
死而復活,聖靈的降臨.
令到連簡單日常的食飯.
都有了更深層的意義.
都可以成為天國價值的展現.
成為五旬節當日見證的延續.
對今天的教會來說.
重現聖靈降臨的畫面從來不是我們的實踐重點.
因為在五旬節發生的那一次聖靈降臨.
已經刻劃和肯定了神已經將聖靈賜給教會.
教會有著藉著聖靈在地上延續耶穌基督工作的身份和角色.
所以我們的重點是我們如何帶著這種身份.
帶著聖靈的同在.
在現今的社會裡.
在很多不同的主流文化裡.
我們如何學似耶穌.
展現出天國的價值.
我們一起在教會裡的生活.
平時的日常生活.
群體生活都是一些很重要的範疇.
剛才說了那麼多節經文.
都不斷提到他們的群體展現.
顯示出基督是會藉著群體而流露.
一個信徒如果只有個人的屬靈,操練而沒有群體.
他無法完整地展露基督給人看到.
而那種藉著群體包括了教會群體之間的互動.
亦包括了教會群體和身處社會的互動.
四十節下亦描述了五重節後的教會群體.
因著他們的群體生活見證.
得全體百姓的喜愛.
主張得救的人天天加給他們.
得全體百姓的喜愛是一個結果.
但未必是他們當初群體生活見證最著想的東西.
而他們會有新加入的人.

$^{361}$最主要的原因是上主的工作.
和四十一節一樣.
四十一節提到有三千人信主的時候.
原文是「若有三千人被加入」.
同樣是指有人加入教會群體是上主的工作.
真正使人信主的是主耶穌自己.
就如福音書中耶穌差派十二個門徒.
七十二個門徒去傳福音.
他是差派和持能力的那一位.
當那些門徒回來說「我們做到了,很開心」的時候.
這些成果並不是表明他們的能力.
而是表明他們屬於神.
教會宣講而有人信主.
我們去做見證,然後有人信主.
其實都是一樣.
並不是去證明我們做得很好.
而是證明我們屬於將人加入的神.
所以彼得的宣講很重要.
我們的分享和見證都很有意義.
不過當我們這樣做的時候.
希望我們都是想為主作見證.
我們都是帶著敬畏神,敬拜神的心去做.
因為唯有聖靈可以在人心動功.
我們都不太知道聖靈會如何使用我們所做的.
唯有聖靈可以令人有那份責心.
我們除了做好自己的見證外.
都應該要準備好作為一個群體.
將神新加入的人納入在我們的群體之中.
而這個得全體百姓的喜愛.
我都覺得是一個頗吊詭的一點.
因為我們都知道.
他在這裡說得全體百姓的喜愛.
但我們都知道在教會歷史上.
初代教會以至今天.
都不斷受到各種的逼迫.
可能去到今天都好像那麼吊詭.
有時人們會稱讚基督徒的不一樣.
他們會見到基督徒的群體會善良一點.
可能會有愛一點.
但有時同樣的人會挑戰基督徒的不一樣.

$^{401}$可能挑戰我們的信仰.
或者挑戰我們虔誠上的不一致.
我們固然仍然要做好見證.
就像在不同的書信裡.
這些使徒都有教會要做好自己.
嘗試去達至無可指責.
但其實基督教的天國價值.
和這個世界的文化.
注定是無法完全融合的.
因為這個世界是一個有罪的世界.
這個世界有很多帶有罪而扭曲的世界觀.
和意識形態.
我們總會去到一個位置.
要堅持告訴他們.
天國價值不是這樣.
我們就是有那份不一樣.
這條線很難在這裡用三言兩語去講清楚.
但我猜最重要的是.
我們有一個敬畏神的心.
當我們感受到我們所持守的價值.
和社會有衝突的時候.
我們以敬畏神的心去持守天國的價值.
亦記住是主去掌管一切.
因為是主會將得救的人加給教會的群體.
說起群體見證和時機.
我上星期六幫一個團體的姊妹做姊妹.
幫她處理了一個婚禮.
也挺特別的.
因為她的另一半是微信主.
我們覺得很接近.
不過她就是微信主.
而她的家人也有很多不是信徒.
所以當我們有份在婚禮裡帶詩歌敬拜.
我們的團契在他們簽婚書有獻詩的時候.
我們一直準備將這班微信的親友放在心上.
希望可以將我們敬畏敬拜主的心唱給他們聽.
讓他們可以感受到一些.
最後我們的姊妹說她老公覺得很感動.
我們也覺得很感恩.
但我們也知道她會不會感動.

$^{441}$其實與我們沒有太大關係.
也與我無關.
誰決定的 其實都是聖靈去決定.
聖靈決定她這個位置會感動.
但她也未必去到決止.
但這樣不代表我們覺得這樣就不用做.
甚至她去到當天這個時刻會感動.
與我們之前這麼多年一直有認識她.
她來到團契當中與我們互動.
認識我們這一班人究竟在信什麼.
平時怎樣去展現.
她才可以更加被我們分享信仰的時候所感染和感動.
希望我們做見證 群體要持守不同價值和信念的時候.
我們記得做好自己去敬畏神.
上帝會按著他的時間去動功.
信仰的成長和展現離不開群體.
但也不能沒有個人努力的追求.
就好像玩皮克問 如果你完全不玩的話.
其實你很難與朋友一起交流.
甚至你的朋友可能開始不想與你一起玩.
因為你會拖他後腳.
但你要自己玩才可以與朋友一起交流 一起努力.
而如果你只是自己玩 沒有群體與你一起玩的話.
你的體驗又不完整.
甚至我本身只是自己玩.
我去到有人一起玩的時候 我才知道原來我曾經損失了什麼.
可能我們信仰上也是這樣.
我上次也很喜歡說打排球 打排球也是這樣.
你可以自己在場邊與自己搓球.
你可以想很多東西.
但你一天不下場與隊友合作 你也不見真章.
團契也是這樣.
如果我們平時只是週末回來交換任務進度.
其實我們未必能考驗到我們能否在生命中同行.
我們唯有慢慢參與在大家的日常生活中.
才能真正認識我們的弟兄姊妹更多.
我們才可以以我們的群體生活一起展現給人看.
一群因著基督而聚集起來的人.
一群學習基督彼此相愛的人.
是怎樣的不一樣.

$^{481}$信仰群體可以擴闊我們的眼界.
因為我們在同一個信仰裡 可能也會有一些不同的焦點.
好像玩皮克文 有些人想儲多些步數.
有些人想儲多些不同的皮克文裝飾.
有些人想多撿些不同的東西.
但其實他們最後的大目標是一樣的.
他們最後成就的東西都是他們走多了.
就好像我們在信仰裡面.
可能有些人對聖經研究比較有興趣.
有些人很喜歡研究怎樣將福音講得清楚明白.
有些人希望將自己的經歷成為與其他人同行的一種力量.
我們可能有不同的焦點.
但這些都沒有辦法完整地反映我們的信仰.
唯有當我們放在一起.
將我們不同的焦點中的得著彼此分享.
互相調整大家對信仰的認識 對上帝的認識.
我們在信仰裡面的體會.
對神的體會才會更加闊 更加豐富.
也會同時豐富了我們一起展現給其他人的見證.
人可以由我們群體裡面不同人.
怎樣有不同經歷神的模式.
去知道神是一個很有創意.
很豐富地與人互動的一個神.
我們在團體,抹餅,讀經,祈禱.
不同的操練裡面.
都去思考更多我們應該怎樣去生活.
有更大的勇氣在我們的信仰群體裡面.
可能有一些碰撞 一起去理解其他人.
怎樣去認識,經歷上帝.
以致我們一起可以有更加豐盛的信仰經歷.
也更加讓人看到一個屬於耶穌基督的群體.
是有怎樣的生命力,怎樣的豐盛.
\newpage



\section{哥林多後書 10:1-5}
\label{sec:032286}
\textbf{勝利不在乎刀劍}~~~~梁錫強醫生
\newline
\newline
連結: \href{http://www.church.com.hk/acms/content.asp?site=cdc\&op=show\&type=product\&code=032286\&layout=sermon}{\texttt{032286}} ~~~~ 語音日期: 2026-06-14
\newline
\newline
\hyperref[sec:032285]{\small{< < < PREV SERMON < < <}}
~
\hyperref[sec:index_chronic]{\small{[返順時目]}}
~
\hyperref[sec:032287]{\small{> > > NEXT SERMON > > >}}
\newline
\newline
哥林多後書 10:1-5
\newline
\begin{longtable}{cl}
\hline
\hline
章節 & 經文 (和合本修訂版)\\
\hline
10:1 & \begin{tabularx}{0.7\textwidth}{X} 我－保羅與你們見面的時候是溫和的，不在你們那裡的時候向你們是勇敢的，如今親自藉著基督的溫柔和慈祥勸你們。 \end{tabularx} \\ \\ \relax
10:2 & \begin{tabularx}{0.7\textwidth}{X} 有人認為我們是憑著血氣行事，我認為必須敢於對付這等人；我但求在那裡的時候，不必這樣勇敢。 \end{tabularx} \\ \\ \relax
10:3 & \begin{tabularx}{0.7\textwidth}{X} 我們雖然在血氣中行事，卻不憑著血氣爭戰。 \end{tabularx} \\ \\ \relax
10:4 & \begin{tabularx}{0.7\textwidth}{X} 因為我們爭戰的兵器本不是屬血氣的，而是憑著神的能力，能夠攻破堅固的營壘。我們攻破各樣的計謀， \end{tabularx} \\ \\ \relax
10:5 & \begin{tabularx}{0.7\textwidth}{X} 和各樣攔阻人認識神的高壘，又奪回人心來順服基督。 \end{tabularx} \\ \\
[1ex]
\hline
\hline
\end{longtable}
$^{1}$剛才那首歌 最後一首歌.
「我要耶穌 我要耶穌」.
什麼感覺都回來了.
這首歌我從小在教會唱到大的.
是一首非常好的歌.
我上網查了一下.
作者是George Webster.
他來自一個牧師的家庭.
他很明白.
他從小到大都是信主的.
但他小時候是幫他家做木工 農場的工作.
一些刻苦的工作.
因為他覺得當時的教會.
很多人的誤解就是以為.
教會只是給一些中產階級的.
有錢人 白人的地方.
當時我說的是美國在十九世紀的年代.
後來他在神學院畢業出來後.
他做了牧師.
有一次他感動寫了一首歌.
「I need Jesus 我要耶穌」.
為什麼要寫這首歌呢.
他的見證就是無論你是什麼種族.
無論你是來自什麼階層.
無論你是有錢還是沒錢.
陶容克今天無論你是在英國住了多久.
或者你是新在香港過來.
無論你是來自哪個不同的教會的背景 中派.
我們都有一個共同點.
我們共同點有一個所謂英文所說的.
common denominator.
就是我們都需要耶穌.
就是這麼簡單.
很簡單的一首歌.
我們很多時候看回《舊約聖經》.
我們就覺得有一樣東西.
為什麼整天都打仗.
我們從《出埃及奇一路》講.
看下去 看到慈思記 撒母耳基.
仍然是在打仗的.

$^{41}$為什麼真的有這麼多仗去打呢?.
實在真的有很多仗去打.
今天在世界上.
就算今時今日.
我們都在中東打仗.
俄羅斯 烏克蘭在打仗.
有很多地方在打內戰.
戰爭是沒有停過.
其實從他們出埃及之前.
你看回《創世記》的時候.
其實雖然沒有說國與國之間的打仗.
但是其實已經有人與人之間的打仗.
甚至兩兄弟的打仗.
該隱和阿伯.
《創世記》一開始已經說了.
結果阿伯被他大哥該隱殺死了.
有以素和雅各的紛爭.
結果因為這樣.
雅各要逃離他地方.
很多年的地方.
十幾二十年的地方.
要逃避他哥哥.
有約瑟和他哥哥的紛爭.
結果約瑟被賣了去埃及的地方.
這全部都是人與人之間的戰爭.
後來去到曠野的時代.
一出到曠野.
第一件事是什麼?.
就是打仗.
不是以色列人去撩交打.
是有些叫阿瑪力人來跟他打仗.
聖經形容阿瑪力人.
喜歡打仗.
善於打仗.
很奇怪.
今天我們年輕人喜歡打遊戲機.
不是打仗.
但是阿瑪力人天生就是喜歡打仗.
你不搞他他也會搞你.
等於走過來看他.

$^{81}$他看什麼.
接著打你兩拳.
那種人來的.
後來去到迦南地.
不停的打仗.
在還沒去到迦南地的時候.
在《申命記》七章一節.
摩西已經說了.
他說你們耶和華上帝帶領你們進去德維葉的地方.
從裡面要趕出許多國家.
就是比你更加強大的七個國家.
所謂迦南的七族.
包括赫人.
甲加薩人.
阿摩利人.
迦南人.
比利時人.
希美人.
和耶布斯人.
這些都是.
含的後裔 即是羅亞的兒子含的後裔.
迦南的後裔.
是成為當時的迦南的原居民.
他說你要趕走這些人.
滅絕這些人.
為什麼會說得這麼嚴重呢.
我看下去就知道了.
原來迦南人他們是用嬰兒去獻祭的.
他們是神像拜假神.
他們是發洩他們的情慾 肉體.
那種所謂的宗教.
你會發覺原來.
聖經所說的戰爭.
有少許是正與邪的戰爭.
有少許是上帝和魔鬼的戰爭.
去到士師記.
非利士人成為主角.
成為以色列人最強大的敵人.
之前說參宣 已經說過非利士人.
薩普爾基的年代.

$^{121}$也是和非利士人不斷的打仗.
到最後大衛統治了江山.
才可以徹底降服了非利士人.
以色列國才得到一個太平的盛世.
打仗.
為什麼經常說打仗呢.
原來聖經是有它的意思在裡面.
這個也是歷史的真實.
但有個零裡頭的意思就是因為人犯了罪.
罪的緣故.
令到這個世界充滿紛爭.
說得沒錯的.
其實在該隱和阿伯之前.
我們看回當人類犯罪的時候.
當夏娃犯罪的時候.
已經很清楚了.
這個是零裡頭的戰爭.
是上帝和魔鬼的爭戰.
上帝說女人的後裔.
要傷你的頭 要傷你的腳筋.
這個就是所謂的原始的福音.
是預言女人的後裔耶穌基督.
最後是會戰勝了魔鬼.
但在這段時間.
魔鬼是將耶穌帶到十字架打斷了他的腳筋.
這個仗還沒有完.
到今時今日這個二十一世紀我們仍然是在打仗.
直至到有一天上帝再回來.
耶穌再回來的時候就新鮮新地.
所以聖經經常用打仗來說.
不是說要我們去打仗.
不是說喜歡打仗.
而是反映了人的罪性.
人生是如戰場.
你說我沒有去打過仗.
但是其實我們一生.
其實每一天我們都在打仗 你知不知道?.
很多困擾圍住我們.
我們不斷地打仗.
我們這樣做對 就是這樣做不對.

$^{161}$這樣做是上帝.
這樣做是魔鬼.
我們其實在靈裡不斷地在鬥爭.
特別是我們基督徒.
我們屬靈的人 有更加大的挑戰.
保羅說得很清楚.
我們應該做的做 我們反而不去做.
不應該做 我們就去做了.
耶穌說得很清楚.
我們心靈願意想去做.
但是肉體就去軟弱下來.
其實每一天我們在爭戰.
其實很簡單 大家看到的.
嬰兒出生 我之前說過.
是很痛苦地出生的 知不知道?.
有沒有見過嬰兒出生是笑著出來的?.
那個就是怪嬰 如果笑著出來.
不喜歡這個世界的.
媽媽的肚子這麼舒服.
出來之後怎麼樣?.
醫院從剛猛到死.
一群人圍著看著你.
空氣又污濁 又多細菌.
他來到這個世界上 他是很不想的.
一生在打仗.
你沒打之前 你媽媽幫你打了.
去超市搶奶粉 是不是?.
搶尿片 那時候不足夠.
COVID那時候是不是這樣?.
排隊入學 申請表 是不是?.
要搶回來的 要霸佔的.
要打架的 那些人在那裡.
去到學校 上學了.
成績每天都是打仗 是不是?.
比較 為什麼姐姐讀書比你厲害?.
你是不是蠢?.
為什麼隔壁的B仔比你高?.
你是不是吃東西不夠營養?.
其實我們每天都在打仗.
為什麼有這樣的情形?.

$^{201}$就是因為反映了人類的罪在裡面.
我們要比較 我們比較就有驕傲.
有比較就會自私 有比較就會妒忌.
就中了魔鬼的詭計.
魔鬼就利用我們這些軟弱的事物.
去不停的攻擊我們.
直到我們去找工作.
我們要踩著別人去上位的.
我們要戳別人的背部 可以領功的.
一模一樣的.
我們出來結識男女朋友的時候.
我們要爭的 聽說要爭女.
原來要爭的.
說得沒錯的 知不知道?.
因為我們是要鬥爭 要戰爭.
所以聖經戰爭其實告訴我們.
有一個意思 屬靈的背後意思就是.
我們基督徒每一天都是屬靈的戰爭.
我們要知道什麼是對什麼是錯.
什麼是喜悅 什麼是不喜悅.
我們要分辨出來.
我們會掙扎 我們有軟弱的時候.
我們要在教會的生活.
有牧者 有弟兄姊妹.
要不斷的提醒我們.
不然的話我們一定會跌在谷底.
任我們怎麼做 我們走的路一定是歪路.
就像小朋友 我們不教好他.
他一定是學壞的.
BB仔 很小的時候.
我的孫女已經懂得扭計.
扭計 知不知道扭計?.
發脾氣 無理取鬧那些 知不知道?.
動物會不會?.
你的貓 你的狗會不會扭計?.
我沒聽過.
我江魚不懂扭計.
我走過去牠們很開心.
我餵牠吃 牠非常開心.
扭什麼計 假裝不餓?.

$^{241}$不會的 這是人.
耶利米蘇說人比萬物更詭詐.
說得沒錯的 人比萬物更詭詐.
在屬靈之間當中 我們要認清兩件事.
第一件事就是誰是我們的敵人.
我們很多時候賴什麼呢?.
賴魔鬼 是不是?.
鬼作弄我.
我們不信主的人就說是鬼作弄.
信主的人就說是魔鬼作弄我們.
我們賴了魔鬼 是.
最終的敵人 終極的靈界惡魔是魔鬼.
但魔鬼是利用人的軟弱.
所以敵人其實很多時候是我們自己的罪.
我們自己的罪.
這個就很清楚了.
以弗所說 保羅說得很清楚.
很多時候 以弗所說的六章.
大家很熟悉的就是十章至十二節那一段.
是說到六章十四至十七節那一段.
是說到我們要穿戴全副的軍裝.
大家都很熟悉了 是不是?.
我們遲些有搞一個背聖經的計劃.
可能都要我們背一節的聖經.
就是要公義做護心鏡.
和平福音當作走路的鞋.
信德做盾牌.
平安福音 差不多吧.
接著有救恩的頭盔.
聖靈的寶劍 我都不記得了.
但我們很多時候都是想著穿些什麼軍裝.
但其實這節聖經之前是有說過.
究竟我們的敵人是誰.
先認清了敵人.
認清了敵人.
因為我們征戰並不是對抗有血有肉的人.
而是對抗那些執政的掌權的管轄這個幽暗世界.
以及天空靈界的惡魔.
說得很清楚.
原來我們今天打仗的真正的敵人.

$^{281}$是天空靈界的惡魔.
很清楚 很清楚.
魔鬼令我們犯罪.
罪令我們窺決上帝的榮耀.
罪令我們沉淪.
罪令我們有驕傲 有自私 有妒忌 有仇恨.
我們放不下.
今天我們首先要打贏這個仗.
要打贏我們自己的罪.
我們要面對這個罪.
不容易打的.
說到打仗 其實我們身體裡面每一天都在打仗.
就是我們的免疫能力.
其實很豐富的話題可以說的.
今天早上都沒有說 因為沒有時間.
我們是在打仗的.
為什麼他會有這種病毒.
他有事我沒事呢.
因為你打贏了 他打不贏 他打敗了.
所以傳統我們廣東話說.
病了就是打敗仗 就是這個意思.
所以古老的中國人都很有哲理.
他知道是打仗 是打輸了.
很多病其實都和你的免疫能力有關.
我讀大學的時候 上世紀.
N那麼多年前.
那時候免疫學剛剛起飛.
但今天免疫這個治療學.
是醫了很多很多的病.
很多以前是絕症的.
現在是可以靠免疫的治療 是可以控制的.
為什麼有癌病呢.
因為身體細胞複製的時候出錯了.
一定有出錯的.
我試一下叫你罰抄.
以前我被人罰抄一百句 二百句 三百句.
有些寫字寫一塊就一定出錯.
身體有時候就是複製細胞的時候出錯了.
但通常你的免疫能力是可以對抗的.
你找到出錯的.

$^{321}$但當有一天你的免疫能力不足.
你看不到那個錯處.
容許它繼續滋生的話 就成為癌細胞.
很簡單 為什麼會有敏感.
有些人帶了口罩來 今天又說因為他有Hey fever.
就是你的免疫能力過強.
那些小小的花粉.
你當作是很大的敵人去攻擊.
結果就流鼻涕 打黑痴.
很多病都是這樣.
為什麼會有風濕關節炎.
不是因為有細菌進入你的關節.
而是你的免疫看錯了 認錯人了.
以為那些是敵人走去打.
有人來探望我 我開門認不到你.
我就打你兩拳 以為你是賊.
害怕吧.
真是的 這個就是免疫能力.
所以其實打仗是說得沒錯的.
無論是肉身 靈裡頭 每天都是打仗.
但今天主要是說靈 今天不是醫學講座.
我們說靈裡頭最大的敵人就是魔鬼.
但魔鬼是利用人的軟弱的罪.
所以最大的敵人其實是自己.
在薩姆意記上有一個很大的例子就是蘇羅.
這個蘇羅是舊約的蘇羅王.
不是新約的保羅 還沒改名的蘇羅.
聖經怎樣形容這個蘇羅.
薩姆意記上九章二節.
他說蘇羅又健壯又英俊.
在以色列人中沒有一個可以與他相比.
如果用這句說話來形容你 你會怎樣.
飄飄然.
你這樣說很難不會驕傲.
以色列人中沒有一個可以與他相比.
又健壯又英俊.
因為蘇羅真是一個很突出的人.
你很遠已經見到他 為何見到他.
還有一句 他比眾百姓高出一個頭.
高大衛猛.

$^{361}$高大是比較有優勢的.
真的比較有優勢 有人做過統計.
高大的人搵工成功率是大很多.
都理不順.
沒辦法 這個數據是這樣寫.
我不是說矮人蠢 矮人沒用.
不是 是人的眼中就是用這些來衡量.
蘇羅比人高大一個頭.
所以我在英國搭巴士 搭地鐵.
我很自卑.
那些師奶都高過我一個頭.
那些鬼妹.
沒辦法.
這個蘇羅.
薩姆爾後來就高立了他做猶太人的王.
以色列人的王.
蘇羅有些戰績.
他不是真的沒用.
他做到事的.
他都聽上帝的說話.
他不是完全不聽.
因為他是猶太人 一定聽的.
他一開始去打的時候.
去打阿滿人 阿滿人都是很厲害的.
大家知道摩訶和阿滿兩個族.
都是遊牧民族 經常四處搞事.
他去打贏了他們.
拿到戰績 人們覺得他是不錯的.
但接著要打菲律賓人.
菲律賓人就真的不同了.
菲律賓人就厲害很多.
如果阿滿是英格蘭.
菲律賓人就是巴西阿根廷.
應該是厲害很多.
照道理就厲害很多.
聖經十五意記上十三字五字.
他說菲律賓人集合.
要與以色列人作戰.
他們有戰車三萬兩.
騎兵六千.

$^{401}$士兵在海邊的沙攬揚多.
不只是多這麼簡單.
他們善於作戰.
見到都害怕.
蘇羅都不敢.
他們知道要靠上帝.
撒末爾說你們聚在前線.
但不要走去打.
等我來.
我們祈禱.
我們獻祭.
我們在上帝那裡.
秩序作戰的能力.
結果在撒末爾十三章說到.
蘇羅在那個地方等了七天.
等撒爾姆爾.
撒末爾沒有來.
撒爾姆爾去完麗池夜總會.
遲到.
等了七天沒有來.
人們就開始亂了.
很多蘇羅的軍隊和百姓就走了.
不行.
撒末爾是我們的精神領袖.
他是我們的先知.
他是我們的士師.
怎麼辦.
保羅.
不要等了.
等我來.
把蘇羅和平安祭帶到我這裡來.
蘇羅就獻上飯祭.
不知道是不是撒爾姆爾作弄他.
剛獻完飯祭.
撒爾姆爾就來了.
突然出現.
那有什麼不對呢.
你不來我就幫你做你的那份.
有很嚴重的錯誤.
第一.

$^{441}$蘇羅是不可以獻祭的.
我們知道撒爾姆爾.
他們當時的士師是他們的領袖.
他們是精神領袖.
他們是皇帝.
都是他按立的幾個皇帝.
他也是利未人.
他有祭司的職份.
但他不是大祭司.
他不是亞倫的後代.
但他當時是有獻祭的責任和職份.
容許他做.
所以他是可以.
但蘇羅就一定不可以.
蘇羅是越了他的權.
這是很違反了上帝的律法.
是很嚴重的錯誤.
都是做.
你等不到他.
他先做.
好過軍心崩潰.
但不是違反上帝的律法.
不是上帝抓得那麼嚴.
其實最重要的是.
你對上帝沒有信心.
你驕傲.
你以為可以靠自己的能力.
可以解決這件事.
你不等候命令.
你不聽撒爾姆爾的話.
撒爾姆爾是上帝的代言人.
他是先知.
你不聽撒爾姆爾的話.
就是不聽上帝的話.
信心不夠.
信心不夠是怎樣呢.
我們祈禱不單只七天.
有時我們祈七年.
上帝還未應允.
繼續祈.

$^{481}$上帝有他的時間.
撒爾姆爾有他的時間.
不要衝動.
不要靠自己的方法.
去解決問題.
一定要靠上帝給你的方法.
我們怎知道哪些是上帝的方法.
哪些是我自己的方法.
有得知的.
要了解上帝的心意.
一定知道.
但很多時候我們急於.
用自己的方法去解決.
蘇羅就是這樣.
他內心.
他是一個不信服的人.
他是一個驕傲的人.
他要告訴別人.
我都行 我都行.
那個是精神領袖.
我才是真正的領袖.
我是皇帝.
我帶你上前線出生入死.
我是皇帝 帶著你打仗.
不是坐在皇宮享樂.
撒爾姆爾.
來祈禱獻祭.
拿了功 拿了彩.
他不想有這些事.
他忍不住 要做出來.
所以.
撒爾姆爾就真的.
去告訴他.
你做了什麼來.
撒爾姆爾說.
等你來 等你來 等你來.
等不到你來 我就獻上繁祭.
告訴布羅廷 你做了糊塗的事.
沒有遵守耶和華.
你上帝吩咐你的命令.

$^{521}$今天我們做了多少糊塗事.
我們過去人生當中.
信主或未信主之前.
我們做了多少糊塗事.
令我們後悔.
很多時候為什麼我們不看清楚.
不想清楚 才去做.
一句說話為什麼不想清楚.
就說出去.
我們做了很多糊塗事.
就因為我們不得信服.
上帝的說話.
所以.
索羅最大的敵人.
就是他自己內心的罪.
又是回到之前.
我所說的三大罪性.
就是眼目的情慾.
肉體的情慾 今生的驕傲.
今天我們的主題.
勝利不在乎 刀劍.
就告訴我們 要贏.
不是靠這把刀 靠這把劍.
為什麼不是.
他們索羅.
他們打贏仗.
他是靠上帝的能力.
見到非利士人.
空馬 他們真的要.
在那裡祈禱.
要獻祭.
才能夠打贏.
這個是很重要的.
今天我們打屬靈戰爭的時候.
很多時候我們想起.
傳福音 傳福音真的是屬靈戰爭.
你想一個未信主的人.
去傳福音.
是有一個戰爭在裡面.
你要夠膽去做.

$^{561}$他可能會有很多的反駁.
很多的藉口 甚至是辱罵你教會的人.
一定會.
你怎麼能站穩呢.
靠著自己的能力.
可以嗎 不可以.
我們要靠禱告.
每一次報道會.
或者我們出隊 或者我們去探訪的時候.
我們一定要禱告.
禱告求循給我們的能力.
所以他們.
打非利士人之前就要禱告.
要獻祭 當然今天我們不是獻祭.
因為耶穌基督一次過已經獻祭.
擺上了救贖.
但我們仍然是要禱告 我們求問神.
求上帝的旨意 求上帝給我們能力.
這個就是.
告訴我們 勝利不在乎.
我們有多厲害.
勝利不在乎.
我們的刀有多鋒利.
或者我們有多大的權力.
有多少錢.
完全沒有關係.
靠著自己.
什麼都做不到 靠著自己.
生命有多長都不知道.
這個星期 香港娛樂圈.
大家有沒有留意 有三個人離開了世界.
一個星期.
都是有點名氣.
有很大年紀的.
也有很年輕的.
人生有多長.
多短 我們是沒有辦法.
可以明白.
那你如何打仗.
你的刀多漂亮.

$^{601}$你的履歷表寫得多好.
又如何.
你有多少錢.
又如何呢.
好 那我們.
在勝利不在乎的刀劍.
最.
大家最熟悉的.
一個例子.
也是我們自小.
主義學都熟悉的故事.
就是大衛打倒哥利亞.
記不記得.
這是一個很大.
很大的對比.
哥利亞是什麼人.
哥利亞是一個巨人.
聖經形容.
他是巨人.
我相信這個.
可能就再高過.
數落兩個頭的巨人.
聖經有說.
多少爪加個虎口.
這個哥利亞.
不只是這樣.
他是很有裝備的.
頭戴銅造的.
頭盔.
身穿的盔甲.
重多少千兩.
腿上又有.
護膝.
兩肩有銅的.
矛.
銅是當時很堅硬的金屬.
所以.
無論是哥利亞天生.
出來的.
強大氣勢.

$^{641}$和他的裝備.
所以他是.
菲律賓人第一勇士.
他是帶頭去打仗的.
大衛.
是什麼呢.
大衛當時.
都是數落王的.
他是一個普通的.
年輕人.
聖經怎樣形容大衛呢.
他說這個哥利亞.
大衛那時候.
他覺得大衛.
他藐視他.
因為他年輕.
因為他面色紅潤.
容貌盡美.
奇怪了.
他帥還可以藐視他.
你怎可以藐視個帥呢.
他說他不像打仗.
他像一個書生多.
像一個在溫室長大的年輕人.
面色紅潤.
不是曬到黑.
家慣田的那些.
打慣仗的那些.
你明白嗎 容貌盡美.
不是像刀疤強有幾條疤痕.
在臉上的那些.
他覺得他.
真的小看他.
和年紀.
年輕 經驗有多少.
打過多少場仗.
哥利亞就是這樣想.
他就和大衛說.
萊巴 我要把你的肉.
給吸空中的飛鳥.

$^{681}$和田野的走獸吃.
以前打仗.
殺了人 殺了他之後.
你不會和他撿屍走去撞火.
沒有這樣的東西.
放在那裡 等野獸 飛鳥來吃了你.
我會打死你.
大衛.
怎樣回答他呢.
你來攻擊我是靠著.
刀槍和銅矛.
但我來攻擊你是靠著.
萬軍之野華的名.
就是你所辱罵大力意識軍隊的上帝.
我靠著的.
是野華華的名.
剛才早堂之後.
陳全和我們分享一首歌.
是靠著耶穌得勝.
其實很多歌.
他講這首歌之後.
我就想起另一首.
靠我靈成仕.
拜My Spirit.
撒迦利亞書.
靠上帝的靈成仕.
萬軍的野華華.
是被非利士人辱罵.
但是這個野華華.
就是會帶他們.
去打贏仗.
怎樣打大家都很熟悉的故事.
大衛伸手從廊中拿出一個石頭.
用機緣甩出去.
擊中非利士人的.
哥利亞的前額.
石子進入額內 他就撲倒.
跌倒了 跟著上前.
就用刀殺了他.
爆冷門.

$^{721}$知道嗎?.
爆冷門.
大家有沒有聽過.
腐賴什麼國家.
腐賴奇.
腐賴索.
聽過嗎?.
你沒聽過吧?.
聽過?.
這麼厲害.
聽過是因為有留意世界盃.
他是第一次參加.
入選.
去到世界盃.
去到世界盃最後.
42強.
48強.
今天下午會對德國.
你想想.
如果這個腐賴索贏他5比0.
就好像大衛打死哥利亞一樣.
知道嗎?.
我形容就是.
在這個強弱懸殊當中.
他竟然可以贏.
等於一個.
你從來沒聽過的國家.
可以贏得這麼強勁.
贏過幾次的德國.
所以其實.
告訴我們什麼呢?.
勝利是不在乎刀劍.
沒錯.
你的刀劍利的話.
你可能真的會打得好一點.
那段時間我去搞.
牛腩跟蚊.
切來切去.
切到手都軟.
磨下把刀.

$^{761}$真的爽很多.
所以中國人的說話.
功欲善其事必先利其器.
是對的.
不是錯.
你磨你的刀是對的.
你搞一個CV是好的.
我不是說這樣不好.
但不要靠這些.
因為你磨你的刀.
煮出來的牛腩可能不好吃.
我們知道的.
我小時候.
要寫毛筆字.
不知道現在的小朋友.
寫什麼毛筆.
練字.
練到寫得漂亮.
我就不覺得了.
現在的字都這麼醜.
有人說買一支漂亮的毛筆.
就寫得漂亮.
我求我爸爸買一支漂亮的毛筆.
給我.
狼毛.
那個年代.
狼的毛.
有些馬毛 羊毛.
狼毛.
都寫得這麼醜.
沒有分別.
這個世界就是這樣.
讀書的時候.
有個同學.
有個有錢的香港人同學.
他很厲害.
打網球.
以前都說過.
買個球拍.
很貴.

$^{801}$他說這個是波格.
贏了溫布頓五次.
那個波格打的那個款.
用什麼金屬做的.
結果輸給我.
糟了.
我的金身驕傲出來了.
不好意思.
他強調什麼呢.
我強調就是.
你不可以靠這些刀劍的東西.
屬靈戰爭更加不是.
不是靠你有多少權力.
有多少勢力.
不是靠你有多少錢財.
不是 絕對不是.
是靠上帝的能力.
剛才說開.
他真正的敵人.
是他自己.
為什麼呢 他是不輸的 他是驕傲.
全世界稱讚他帥的時候.
他就驕傲 很簡單.
當他打贏了仗.
回來 菲律賓人其實都是靠大衛.
他打贏了回來之後.
那些女人都走出來了.
擁抱 不是擁抱.
走出來歡歡喜喜.
打鼓奏樂.
唱歌跳舞迎接.
他們的軍隊回來.
那些婦女又唱什麼.
蘇羅殺死千千.
好聽啊 真是 哇 蘇羅勁.
接著一句呢.
大衛殺死萬萬.
就是說 蘇羅勁.
不夠大衛勁.
蘇羅就很憤怒.

$^{841}$一 我帥過你.
我高大過你 二 我是皇帝.
你是幫我打仗.
居然功勞.
給了大衛.
蘇羅就很憤怒.
很憤怒.
就說 你給萬萬給大衛.
千千給我 那又怎樣.
我都是皇帝 我王國沒有給他.
他還是憎恨這件事.
我是皇帝.
今天那些執政權力.
那些人 都是這樣想.
做錯事.
不用負責.
因為他覺得自己是對的.
沒錯的 他高高在上.
沒有人能動他.
驕傲就會出來.
驕傲就會出來 這是人的罪.
所以最大的敵人就是罪.
第二個就是.
最重要打贏仗.
的武器.
不是物質上的.
武器 而是.
屬靈上的武器.
這才能幫助我們.
今天我們面對不同的戰爭.
無論是我們在.
外面 職場上.
家庭上 或者國與國之間.
我們每天都有新聞.
其實回到我們.
教會當中.
也有很多的戰爭.
在教會經歷了.
很多年的事奉.
也有不同的教會.

$^{881}$經歷過很多不同的戰爭.
真的不想用這些來形容.
有些人就說.
為什麼基督徒都是這樣.
我就告訴他 因為基督徒都有罪.
所以我們更加需要耶穌.
從來沒有一個道德光環掛在我頭上.
我是特別好過你們.
或者你們好過外面的人 沒有的.
只不過我們靠著聖經的說話.
不斷地提醒我們.
但是很多時候我們做不到.
這個問題很大.
今天魔鬼在我們中間.
每天都在中間.
你不要以為我們的十字架可以擋住.
魔鬼不能進來.
沒有這樣的東西 我們信的基督教.
是沒有的.
魔鬼會進來.
它會霸佔我們的空間.
它會迷惑我們.
它會引誘我們.
它會令我們分開.
分裂.
令我們不再同心.
最近和一些弟兄姊妹.
最近兩個禮拜.
不約而同有幾個弟兄姊妹和我說.
他說教會很多是非.
他不是說.
教會很多是非.
他說這間教會很多是非.
這句真的揶揄我.
他說他都不想再回去.
不止一個人 有幾個都是這樣說.
藉口也好 什麼也好.
是不是我做得不夠好.
是不是很多時候我們被魔鬼利用了.
甚至有姐妹和我說.

$^{921}$如果大一點的地方.
看現在教會一早已經分裂了.
我當然是很不想.
也不希望會發生.
但是看到人的軟弱.
為什麼聖經不斷提醒我們.
為什麼不斷講戰爭.
不是要我們去打仗 不是要我們去威風.
是提醒我們.
你每一天都在打仗.
你是不可以靠你自己.
你靠自己一定是.
你那些壞事就會走出來.
你就會用這些東西去解決問題.
就會造成很多的紛爭.
無論是教會還是教會以外.
教會不是特別多.
外面都有.
不過教會發生的事我們更加心痛.
因為大家是弟兄姊妹.
是我們的骨肉.
是我們的親人.
更加不想見到.
令我們流淚流得更加多.
一定的 這是一定的.
所以今天我們來到神的面前.
我們就要明白.
就好像剛才大家看的哥林多後書.
他告訴我們我們征戰的兵器.
不是屬於血氣的.
而是憑著上帝的能力.
我們雖然在血氣中行事.
卻不憑著血氣征戰.
這句話.
之前保羅講的東西.
是很妙的.
他說我來到你們那裡.
我用很溫柔的方法去跟你講.
但是我出去的時候.
我又要很勇敢.

$^{961}$所以保羅說我精神分裂.
我對著屬靈的世界的時候.
我一定要靠著上帝的能力.
我要很勇敢面對.
但是我對著弟兄姊妹的時候.
我又要很溫柔對著你們.
今天做牧者的困難.
就是這樣.
以前我事奉了這麼多年.
我沒有這樣想過.
最近這幾年.
我真的有這樣想.
不容易處理這個衝擊.
所以.
保羅說.
你們今天仍然在血氣中.
但是你不要用血氣.
來去征戰.
這是保羅很深入的神學.
Now but not yet.
今天已經是末世時代.
但不是末日.
今天你們基督徒特別多挑戰.
特別多掙扎.
因為你是一個屬靈的人.
在一個屬世的地方.
你要面對屬世問題的時候.
是靠著上帝的聖靈還是靠著血氣.
這就是我們掙扎.
世人沒有這個問題.
想到想做就去做.
不犯法就可以了.
他們沒有這個枷鎖.
但我們屬上帝的人.
我們有這個枷鎖.
我們要明白.
我們要天天提醒我們什麼是對什麼是錯.
沒錯很多時候我們都犯錯.
昨天有個姐妹問我.
為什麼出了耶穌就要回教會.

$^{1001}$我很怕是非不想回.
我家裡敬拜.
這個不是合上帝的心意.
絕對不是.
因為你很快就會跌下來.
魔鬼能力你一個人是撐不住.
就是這麼簡單.
所以我們要憑著上帝的能力.
就可以攻破堅固的淫慮.
再堅固的淫慮.
魔鬼最堅固的淫慮.
魔鬼下去也只有幾招.
其實很容易攻破.
不過我們經常都失敗了.
在這幾招之下.
但靠著上帝的能力.
我們就必然可以做到.
上帝的教導就是告訴我們.
今天我們要回到上帝那裡.
回到上帝的角度那裡.
我們要伶侶頭我們要堅強.
我們教會要有一個合一的團體.
我們才可以抵抗外敵.
否則魔鬼就.
很容易就攻破到我們.
面對自己的罪.
面對自己的敵人.
就是自己 就是魔鬼.
每一天我們都要面對.
求主繼續帶領我們.
大家低頭.
當我們來到你面前的時候.
我們想到我們真的很軟弱.
我們沒有能力.
可以延長自己一刻中的壽命.
更加沒有能力.
去靠著我們的肉體.
我們的血氣.
可以面對這個世界這麼複雜的環境.
無論是教會以內教會以外.

$^{1041}$我們都知道我們.
不是那麼容易.
靠著自己的能力去做.
只會做錯 只會自己的罪性出來.
只會被魔鬼利用.
被魔鬼迷惑.
求主讓我們靠著聖靈來行事.
靠著耶穌.
我們每一個人都需要耶穌.
我們靠著上帝的能力.
我們要打贏.
不在乎有刀劍.
不在乎人多人少.
是在乎萬軍耶和華的權能.
求主繼續帶領我們.
無論我們在未來的日子當中.
遇到什麼困難的時候.
我們都要穿戴全副的軍裝.
屬靈的軍裝.
我們是要對抗.
這個惡者.
直到再回來的日子.
鬥爭才可以平息.
才來到一個世界.
是沒有罪的.
沒有痛苦的 沒有死亡的.
求主繼續帶領我們今天.
在這個動盪的時代當中.
我們仍然緊緊握著你的手.
我們不要放棄.
我們祈禱奉耶穌得勝名之祈求 阿們.
\newpage



\section{約翰福音 13:31-38}
\label{sec:032287}
\textbf{主的榮光往哪裡去了?}~~~~黃紹權牧師
\newline
\newline
連結: \href{http://www.church.com.hk/acms/content.asp?site=cdc\&op=show\&type=product\&code=032287\&layout=sermon}{\texttt{032287}} ~~~~ 語音日期: 2026-06-14
\newline
\newline
\hyperref[sec:032286]{\small{< < < PREV SERMON < < <}}
~
\hyperref[sec:index_chronic]{\small{[返順時目]}}
~
\hyperref[sec:032288]{\small{> > > NEXT SERMON > > >}}
\newline
\newline
約翰福音 13:31-38
\newline
\begin{longtable}{cl}
\hline
\hline
章節 & 經文 (和合本修訂版)\\
\hline
13:31 & \begin{tabularx}{0.7\textwidth}{X} 猶大出去後，耶穌說：「如今人子得了榮耀，神在人子身上也得了榮耀。 \end{tabularx} \\ \\ \relax
13:32 & \begin{tabularx}{0.7\textwidth}{X} 如果神因人子得了榮耀，神也要因自己榮耀人子，並且要立刻榮耀他。 \end{tabularx} \\ \\ \relax
13:33 & \begin{tabularx}{0.7\textwidth}{X} 孩子們！我與你們同在的時候不多了；你們會找我，但我所去的地方，你們不能去。這話我曾對猶太人說過，現在也照樣對你們說。 \end{tabularx} \\ \\ \relax
13:34 & \begin{tabularx}{0.7\textwidth}{X} 我賜給你們一條新命令，乃是叫你們彼此相愛；我怎樣愛你們，你們也要怎樣彼此相愛。 \end{tabularx} \\ \\ \relax
13:35 & \begin{tabularx}{0.7\textwidth}{X} 你們若彼此相愛，眾人因此就認出你們是我的門徒了。」 \end{tabularx} \\ \\ \relax
13:36 & \begin{tabularx}{0.7\textwidth}{X} 西門‧彼得問耶穌：「主啊，你去哪裡？」耶穌回答：「我所去的地方，你現在不能跟我去，以後卻要跟我去。」 \end{tabularx} \\ \\ \relax
13:37 & \begin{tabularx}{0.7\textwidth}{X} 彼得對他說：「主啊，為甚麼我現在不能跟你去？我願意為你捨命。」 \end{tabularx} \\ \\ \relax
13:38 & \begin{tabularx}{0.7\textwidth}{X} 耶穌回答：「你願意為我捨命嗎？我實實在在地告訴你，雞叫以前，你要三次不認我。」 \end{tabularx} \\ \\
[1ex]
\hline
\hline
\end{longtable}
$^{1}$好,今天我想和大家一起去看的就是約翰福音第十三章.
其實如果大家來退休會由第一晚開始.
其實我們老早已經和大家說了.
在退休會裡面我們看的是集中於約翰福音第十三章裡面的經文.
想討論到一個題目就是《上帝你的榮耀去了哪裡》.
在約翰福音的最後的一個部份.
亦都指出是耶穌基督在耶路撒冷的最後的一段日子.
祂向很多的門徒都談及很多的題目.
但是其中一個最突出的題目就是叫做《榮耀》.
現在個個都說找AI來做一些聖功或者工作.
對於我來計我其實對AI一直都不是不用的.
不過我就覺得它的東西你不要盡信.
第二樣東西它就對於一些題目就未必是那麼透徹理解.
因為它經常都用教會一向的慣性.
大家喜歡怎樣解讀聖經.
它就給回解釋給他.
但是隨著最近有一個科技發展.
我就和他聊聊天.
我就和他說你可不可以和我講一下約翰福音第十三章的內容.
他說你想講什麼約翰福音第十三章.
我說不如這樣你和我講這段聖經和榮耀神有什麼關係.
或者和榮耀這個題目的關係.
他就不斷在那裡只懂得擦鞋.
他說這個真是很好的題目.
我說那其實是怎樣.
你可不可以和我將這段經文解拆講榮耀.
他就只說這是一個很好的題目.
我見到在第31節以後講過好幾次的題目.
他很蠢其實我發覺AI.
跟著我就再和他講多兩句.
我就說你可不可以按照前文.
為什麼要講有關於耶穌幫門徒洗腳的題目.
又要講有關於彼得要談論到要洗完腳之後又要洗身.
但是又跟著又馬上要談及到猶大.
他是說準備出去出賣主.
他就開始猶疑那個AI.
我發覺我又不要較精AI了.
不如我在講道上一次過和大家討論這個題目.
所以就選了這個題目和大家一起去看.
如果我們仔細去看.

$^{41}$在《約翰福音》裡面.
從第13章至第17章的經文裡面.
這個談論到榮耀的題目是多次出現的.
如果大家熟悉一點元文聖經.
希臘文的聖經.
講到有關榮耀這個字眼.
通常是不出這個字眼的.
叫做Dosa.
Dosa是一個名詞來的.
榮耀.
如果是動詞來算的話.
就叫做Dosaso.
原來在從第13章.
特別是《約翰福音》第13章開始.
是出現過很多次的.
譬如在第13章的第31節裡面.
出現過兩次.
在第32節裡面出現過三次.
在第14章的第13節裡面又出現過一次.
在第15章的第8節又出現過一次.
在第16章的第14節又出現過一次.
在第17章的第1節又出現過兩次.
在第17章的第4,第5,第10,第22節和第24節.
都談到榮耀這個問題.
不期而我們就會問.
其實為什麼要在《約翰福音》的最後一部分.
約翰要重重複複記載到這個題目.
和耶穌受死有什麼關係呢?.
和出賣耶穌的猶大有什麼關係呢?.
甚至我們要說到那個很熱心的追隨主的彼得.
我們一會兒會看的.
他很厲害的.
他願意和耶穌一起去死都可以.
他講出.
他其實和這個有什麼獨特的關係呢?.
這個正正就是我想透過這段聖經.
去看一下榮耀其實是想指什麼東西.
和榮耀究竟是不是只是.
談論到耶穌和聖子之間的關係.
和我們有什麼關係呢?.

$^{81}$特別我們今天作為信徒.
榮耀神這件事其實和我們是不是完全.
是毫無一個直接的關聯呢?.
這一方面我希望和大家透過這段聖經.
一起探討.
重新發掘榮耀和我們的關係.
尤其是神的榮耀和我們的關係.
我們一起看看《約翰福音》的第十三章.
大概說些什麼.
如果大家有聖經的.
我就和大家好快地去看一看.
如果大家去看《約翰福音》的第十三章第一節.
一開始就說「逾越節以前」.
接著什麼呢?.
「耶穌就知道自己離世歸附的時候到了.
祂既然愛世間屬自己的人.
就愛他們到底」.
這句話如果就這樣看.
就好像沒有什麼特別的.
就是耶穌死我又知道.
耶穌接著會有如月餐我又知道.
但是那句很奇怪的說話就是.
祂既然愛世界屬自己的人.
祂多了一個屬自己.
就是耶穌原來有些自己人.
而耶穌說祂對於這些祂的自己人.
祂愛到底.
其實反映了幾樣東西.
第一樣東西就是.
這個世界是不是有些耶穌不愛的人呢?.
或者不屬於祂的人呢?.
還有如果是這樣的話.
為什麼要說愛到底呢?.
那不就是愛了嗎?.
祂也和祂在一起.
其實這幅圖畫首先說了一個.
耶穌的一個性情.
就是去到最後的事奉階段.
其實祂很擔心祂的門徒.
他們接下來的屬靈生命.

$^{121}$能不能夠繼續在祂不存在.
就是耶穌不存在的時候.
仍然是可以向著上帝.
那個屬性的方向去發展.
祂擔心他們會倒退.
但這是屬祂自己的人.
所以耶穌一定要想一個方法.
保證可以使得這班信徒.
是要一直存在在.
耶穌所看的這個救恩裡面.
為什麼這麼重要呢?.
因為原來當信徒.
他是能夠真真正正活出一個.
在救恩裡面的生命.
就能夠真正能夠成為.
上帝的榮耀的一份子.
我們一起去.
首先看到這一個.
是似乎讓很多基督徒.
今天成為一個很大的一個隱憂.
我放不下它,可能過熱了.
我忘了插一把風扇.
不用找嘴吹了,不會吃的.
我應該有把風扇的,不過不要緊.
其實我們會有一個疑問.
究竟救恩是不是給所有的人.
所有世界的人都一定得到呢?.
在一些教派.
尤其是天主教裡面就是這樣想的.
只要你跟教會有接觸.
那你就會得救的了.
Whatever type of contact or connection.
你都可以得救的.
因為他們認為在聖經裡面.
曾經上帝在提摩太前書第二章第四節說到.
「不願有一人沉淪,乃願意人人都悔改」.
其實這句話是一個combination.
前者就是提摩太前書所談及到的.
「神願意萬人得救」.
就組合了彼得後書的三章第九節裡面的.

$^{161}$就是「願意人人悔改」.
我們就不知道是誰不知道在什麼時候.
就說到這句話「願意人人得救」.
其實是「願意人人悔改」而已.
不是「得救」,這是上帝的心願.
「得不得救」是那個人「悔不悔改」的一個之後的發展.
但是我們很多時候,你知道教會的事情.
聽一點不聽一點,說一點不說一點.
然後將它分開來組合.
他就憑著這句話去指導.
因此救恩應該是每個人都得到的.
但是當我們從聖經裡面看見.
耶穌其實反映出的不是每個人都一定得救的.
尤其是剛剛我們所看的第十三章的第一節裡面.
說得很清楚.
其實真真正正願意去接受耶穌基督.
可以悔改進入到去成為耶穌所愛的人.
不是全世界的人.
所以約翰福音的三章十六節.
是指出上帝祂的屬性的心願.
來到十三章裡面的第一節.
其實就反映出那個實際的救恩的成就的發展.
從這個種種的背景之下.
我們就會了解到.
為什麼耶穌一開始就在這裡約翰提及到.
他離世的日子很近,歸去父的時候.
門徒如果沒有了耶穌照顧著,看著他們.
或者教導他們,糾正他們.
那他們會不會走樣的?.
如果走了樣的話.
這個會不會就成為了上帝的一個羞辱呢?.
榮耀的相反詞其實就是一種羞辱.
今天,我想作為信徒的.
我們活在這個地上面.
其實我們最怕就是使我們做了一些事.
使到上帝羞辱.
但是什麼使到人能夠在羞辱裡面.
都能夠扭轉成為上帝榮耀呢?.
在我們這個時代.
就剛才我們所講到的彼得後書.

$^{201}$我們願意悔改.
但是沒有耶穌存在,是不是就不會悔改呢?.
或者悔改就會倒退呢?.
甚至會用癌症的陳述叫做Relapse.
翻法,回歸這個罪的狀態裡面.
如果我們回歸到.
回歸到那個罪的狀態裡面的話.
我們就會失去了上帝的榮耀.
虧損了上帝的榮耀.
所以我在這裡再講多一點的就是.
原來當我們看從約翰福音13章至17節的時候.
我們就會發現到一個事實.
就是原來.
耶穌每次講到他榮耀父神的行動.
都是因為他能夠按照父神他的指示而行.
換句話來說.
如果一個屬神的子民.
同樣肯願意按照父神的指意而行事的話.
他所作的事就會成為神的榮耀的一部分.
你不需要去高舉.
因為我是基督徒所以我才這樣做.
做了之後有上帝的心思意念.
別人就說.
是啊,所以我們的父神很榮美的.
你要歸榮耀給神.
不是這樣做的.
是由人去發現到我們原來是肯聽上帝的話.
所以榮耀這個字眼除了指導做了一些事情.
去能夠成為一個光榮的事情之外.
另一個意思就是翻譯為值得別人讚賞.
所以英文就稱之為praise.
別人見到你好的時候.
他會回想起你的神.
就好像做子女一樣.
如果我們做得好的時候.
別人很自然透過我們就會發現到.
是啊,王媽王爸的兒子真的挺聽話的.
就是這個意思.
我很記得我小時候.
我媽媽有時候有些朋友來家裡探訪.

$^{241}$探訪其實有些娛樂.
我們不詳細說了.
但是有一次我很記得我出街.
出夜街的時候.
我很少出夜街,我不喜歡出夜街.
師母知道的.
我去旅行一入夜我就回家了.
但是有一次我出夜街的時候.
我臨出去的時候.
我就見到我媽媽和她的朋友在那裡聊天的時候.
我就和我媽媽說.
媽媽我出街幾點回家了.
那我就出去了.
後來我回家我媽媽就和我說.
她說你知不知道你出街的時候乜乜太走來和我說.
王太,原來你兒子出街會告訴你的.
我媽媽說有什麼奇怪.
她說你兒子真是生性.
我媽媽說.
哇,真的別人稱讚你的時候.
我也挺開心.
你知不知道榮耀神不是你打鑼打鼓.
然後招惹做了一些很燦爛的事.
然後還要告訴別人是神做的.
不需要的.
你和我身邊在做的事.
這個群體在做的事.
就會顯出上帝的榮譽.
而這個就是我們今天去看的.
第31節至到第38節的主題.
講的內容.
我們一起去看看.
這部份我分開兩個部份和大家一起去看.
第一個部份就是第31,32節.
我們跳一節然後34,35節.
找到之後我們自己又讀多一次.
31,32和34,35.
我們去了解第一個部份.
怎樣是透過信服神.
能夠歸入到神的榮耀的部份.

$^{281}$我們一起找到之後一起讀.
他既出去耶穌就說.
如今人子得了榮耀.
神在人子身上也得了榮耀.
神要因自己榮耀人子.
並且要快快地榮耀他.
我們跳一跳第33節.
然後第34節.
我賜給你們一條新命令.
乃是叫你們彼此相愛.
我怎樣愛你們.
你們也要怎樣相愛.
你們若有彼此相愛的心.
眾人因此就認出你們是我的門徒了.
如果我們看回第13章裡面.
所有的內容.
我們就會發現到耶穌的擔心不無道理.
擔心信徒會走樣.
譬如最簡單第13章.
耶穌最開頭就是為了所有門徒.
在筵席間去洗腳這件事.
首先這一段經文我們會發現到.
耶穌洗腳的包括了加略人尤大.
加略人尤大沒有因為耶穌早知道他是一個會出賣他的人.
他就不洗他的腳踢他出局.
其實耶穌愛人.
這就是我剛才所說的三章十六節的上帝普及的性情.
愛人的性情.
想人歸回他的性情.
上帝用盡了所有他可以做的方法.
要讓你知上帝其實是想包羅你在他的拯救他的永恆國度裡面.
所以這個洗腳這一段經文其實是透過.
給我們看耶穌甘願作為一個老師.
都願意去以一個職責去服事.
去盡心去做.
他不會去有所取捨.
選擇去篩選一些人.
我不給你一個救恩.
我不讓你知道上帝拯救你.
他沒有這樣做.

$^{321}$做了.
但是接著我們去看的是什麼呢.
從第二十一節開始就馬上跳進去一個現象.
就是談論到耶穌面對這位的尤大.
這位尤大其實他的心仍然是想離開耶穌.
然後去出賣他.
藉此為自己謀取一些利益.
因為他也知道耶穌說他快要死了.
要上耶路撒冷上面.
如果這樣那條財路沒多久.
這種叫做什麼.
叫做最後割禾青.
屬靈上的最後割禾青.
賺得多少得多少.
跟著主.
他反正都說他自己會死.
那不如趁他死之前.
我撈一點都不差.
不過我們去看看耶穌第二十一節這裡說.
說了這話就心裡憂傷說些什麼.
其實談論到先前有關於洗腳的問題.
跟彼得的一番的爭論.
是談論到彼得一路想說.
主啊你洗了我的身體.
全身洗淨它.
但是他不明白耶穌其實說的主題.
是要去忠誠地去實踐到上帝的說話.
就是榮耀神的基礎.
他看不到這一點.
他看的仍然是一個外觀上的行為.
一個包裝.
一些的動作.
來證明自己是很委身在耶穌的身上.
但是這位門徒.
彼得既然這麼盲目的時候.
當耶穌看到猶大準備出賣他.
耶穌其實也很坦白地去點出了.
給包括門徒知道.
其實猶大你是會出賣我的.
如果大家去誠實地去看第二十一節.

$^{361}$至到去第三十節裡面.
你會很清楚看見一個事實.
當耶穌說有人要出賣他的時候.
門徒問他:誰會出賣你?.
耶穌其實說了給他聽.
誰啊?.
那個我一會給麵包,沾醋給他吃的那個.
那個就是了.
任何人都聽到了.
連那個會出賣他的那個都聽到了.
但是奇在奇在.
那個聽到的猶大都去接麵包來吃.
你應該做賊你都不要這樣.
收起來一點好不好?.
不是啊,照吃啊.
為什麼?.
因為他明白到信徒.
那些恨人於不義的人.
其實很看得到其他人那些軟弱的.
其實他看得明白.
這班門徒其實都是蒙的.
耶穌做了這個展示.
猶大已經做了.
彼得你們應該看得到了吧.
是不是啊?.
但是這段聖經卻是很諷刺地告訴我們知道.
這班人是怎麼想的.
第29節.
有人因為猶大帶著錢囊.
以為耶穌是對他說.
你去買我們過節所應用的東西.
或是叫他什麼周濟窮人.
他們想事是用自己那套去想事的.
啟示做了給你看你都不會明白.
那令不令耶穌擔心呢?.
弟兄姊妹.
我做牧者這麼多年.
不是沒有向弟兄姊妹.
有時弟兄姊妹犯錯,犯罪.
其實我之前都和你討論過這個問題.

$^{401}$他說我忍不住.
這些就是耶穌擔心的問題.
因為我們仍然在這個世界上.
很多時候用我自己那套思想.
去想耶穌他想的事.
不是耶穌所看得明明白白的事.
更諷刺的是什麼?.
我們不要漏最後第30節.
猶大受了那點餅立刻就出去.
那時候是夜間.
其實如果原文聖經這句話的翻譯就很有趣.
當猶大接受了麵包之後.
他還拿著這些麵包一起走.
耶穌才說得清楚過無比.
連他都做了.
摺了麵包,立了麵包.
走了一起去做什麼?.
走出去.
半夜走.
他進入到夜間裡.
其實我想說.
其實這一點就反映了今天很多信徒.
其實他的罪不是不知道.
是甘願跑下去的.
或者我這樣說都明白.
這些教徒明知道有些事不合上帝心意.
都照樣衝下去.
立刻做.
怕做得遲.
教會猶大做了一個見證給我們看.
就是一個人如果他的心完全是困在自己的裡面.
他的目的.
其實他是不理一切地衝下去.
他是在上帝恩典的源頭的身邊.
你不要留意.
這幅是什麼圖畫?.
是由上帝恩典的源頭走下去黑暗的裡面.
耶穌在這樣的狀態.
擔不擔心留下那十一個門徒呢?.
我不在的時候.

$^{441}$我的源頭不在你身邊.
你會不會衝下去?.
是不是很危險?.
所以第十三章我們讀的時候.
我們應該有這份迫切感.
當門徒一個又一個.
如果在耶穌不在他身邊.
他們都會走下去這個黑暗的裡面的時候.
剛才我們說過什麼叫榮耀神.
是會聽上帝說話的時候才是榮耀神.
好像耶穌從第十三章至十七章裡面.
每一次說.
父啊我照了你做的指示去做了.
所以我能夠去榮耀你.
譬如特別是第十四章裡面.
就談論到.
耶穌說我已經照足你的說話.
你的指示去做.
而你得到榮耀.
門徒的榮耀是什麼?.
耶穌得到榮耀.
耶穌和父神得到榮耀.
那不就好了.
其實搞定了.
不需要我們有.
為什麼耶穌要有這麼擔心.
因為耶穌想囊括我們在上帝的榮耀裡面.
不過我們在一個危機的狀態當中.
因此在這段聖經裡面.
當我們去到第十三章第三十一節.
剛才我們所讀的經文.
這裡所講到.
如今人子得了榮耀.
神在人子身上也得了榮耀.
這個好像只是和上帝和上帝之間.
和聖子之間得到互相榮耀.
但是有extension.
extension是什麼?.
第三十二節.
神要因自己榮耀人子.

$^{481}$並且要快快地榮耀他.
在希臘文裡面.
原來做了一個文化上的改變.
由上半節三十一節所講到的榮耀.
是一個過去已在完成式.
就是已經得到榮耀.
轉入一個未來式.
future tense講到榮耀的問題.
上帝永恆榮耀是人都知道的.
但是為什麼要在這裡向門徒講這一點呢?.
因為耶穌想給的是門徒們.
你們都能夠在這個榮耀當中.
所以這是一個未來繼續要發展的事情.
上帝彼此得到榮耀.
人子得到榮耀.
既定事實.
但是他是想將這個榮耀.
持續在我們身上.
都繼續發生.
因此在這個文化上面.
讓我們看到.
當加勒人猶大正在趕急地.
向著這個榮耀的源頭離開的時候.
耶穌就在這個時候.
要更加趕緊地去將這幫門徒.
去將他們穩固地.
去將他們建立在一個.
不會失去神榮耀的可行.
可行的一個我們叫framework.
一個架構裡面.
所以就有給了剛才我們所看的.
第34和35節.
這個新的命令.
這個新的命令有個特點.
我們再讀多一次.
34和35節.
我賜給你們一條新命令.
乃是叫你們彼此相愛.
我怎樣愛你們.
你們也要怎樣相愛.

$^{521}$你們若有彼此相愛的心.
眾人因此就認出.
你們是我們的門徒了.
這裡有個可行的實踐的框架.
告訴我們.
就是你們要照我的definition的相愛.
來在你們之間當中去施展.
這一點是今天我們教會最多人說.
但是最多誤會的一個題目.
我們經常說愛.
沒錯.
不過這個愛的定義.
不是你定不是我定.
而是主耶穌所定.
所以這裡說.
我怎樣愛你們.
你們也要怎樣相愛.
人間的定義的愛.
很多時候是你怎樣在對方上得到一些東西.
但是耶穌的愛.
是永遠針對著我們屬靈上最大的益處而付出.
如果大家記得我在歌羅西書的時候.
也說到一點教會的成就.
教會群體的建立.
不是說只是牧師在做一些聖禮.
你搞一些聖功.
然後做一些愛展.
這樣就完結了.
而是每一個信徒彼此之間.
為對方最大的屬靈的益處而去付出.
有時這些付出.
對方可能未必一下子是明白的.
但是這個付出是對你屬靈上最大的益處.
所以你會去回心去想並且因而去接受.
所以基督徒.
我們對於愛的定義.
很多時候不是純粹利益上的一種好處.
譬如我回了這間教會.
接著我買車有汽車經紀.
我買燕梳有燕梳經紀.

$^{561}$我買樓有買樓經紀.
我有法律問題.
有律師可以的信徒在我裡面.
我們以為是這些東西.
然後我以為去實踐這些東西.
互相互惠互利的時候.
就等同於教會彼此相愛.
不是.
耶穌在這裡不是說這些東西.
是講及於當你面對著.
你在這個世界有任何挑戰的時候.
會拉你們離開上帝的榮耀.
使你不遵從上帝的事事去行事.
因而虧損了上帝的榮耀的時候.
你們要作出彼此的關懷.
所以當我們去看回其他的書卷.
譬如保羅的書信.
他都談論到很多時候教會的建立是什麼.
是一種提醒.
是一種挽回.
是一種一起去面對挑戰的一個經歷.
而不是大家有什麼好處.
大家跟著互相在這裡只是利益輸送.
當我們看見耶穌在這個時候.
給了這個框架他們的時候.
那個效果的最終.
不只是這個群體.
我們再看第35節.
「你們若要彼此相愛的心」.
這個說話就已經是.
假設上面所講到你們已經相愛了.
當你們真的做到.
能夠彼此對對方的靈明益處去著想的時候.
最終塑造了的.
不只是我們.
是眾人因此認出你們是我的門徒了.
這句說話要斬得很仔細去看的.
首先.
眾人.
如果大家去看回這句說話.

$^{601}$眾人這裡所講是一個普羅大眾.
Everyone.
所有人.
在這裡.
我們會發覺.
原來那個影響是跳出了.
這個教會的空間.
第二點.
是他們認出.
是他們留意得到的.
如果大家做信徒做了這麼多年.
你們都會明白的.
這個世界留意你最多的那個人.
就有兩個.
一個.
一定是敵人.
一個就是你的親人.
尤其是你的父母.
他最留意你的了.
譬如你媽媽你爸爸.
最怕你吃不飽.
穿不暖.
去了哪裡會不會交到一些壞朋友.
去害你.
你有傷害.
第二個就是剛才Sandy所講的.
就是敵人.
世人最喜歡看基督徒的.
為什麼.
因為看基督徒可以跳出.
基督徒的不好的地方的時候.
他們就會最開心.
他開心在什麼.
證明我不信都不是很差.
但是如果我們做到.
好像耶穌基督一樣.
彼此相愛的時候.
這些人.
就會認出.
第四樣東西.

$^{641}$在這段聖經除了眾人認出.
第三樣就是.
我們是基督的門徒.
這個說話是什麼意思.
原來我們做得好.
被這個世界所有的人看到的.
無論是關心你的人.
想你死的那個人都好.
所有人都好.
他們會見到我們的身份.
是屬於上帝.
屬於耶穌基督.
這句說話.
和剛才我所講.
我晚上出街.
我和我媽說.
我現在準備出街.
幾點鐘會回來.
然後她的朋友跟她說.
王太.
你兒子很乖.
這就是我的兒子.
有沒有以前看過一套.
電視的廣告.
說有一個男孩去到哪裡都好.
他的弟弟.
譬如畢業典禮.
他遲到.
不過他進場都說.
This is my brother.
以自己的兄弟為榮.
這種不是自己得到什麼好處.
你不是台上得獎的那個人.
你不是台上得到讚許的人.
你只是台下的.
其中這個人的家人一份子.
不過你會分享到那份喜悅.
人家看見你好的時候.
就曉得讚賞你的家人.
耶穌是想將門徒.

$^{681}$納入一個這樣的結構裡面.
讓所有的人.
在這個世界上.
讓你和我能夠.
作一個這樣的引導的角色.
所以這句聖經其實.
不單只是說到一個信徒的教會.
他們之間這個群體.
他們怎樣建立合神心意這麼簡單.
而是我們所說的真正的.
報道的動力.
而這個動力是要讓我們納入到.
他在主的榮耀當中.
這個很有趣的一段經文.
但是相影那個.
現在正在進入黑暗的猶大.
他卻是沒有這份參與的機會.
造成了一個很大的反差.
一個對比.
這個對比也能夠讓我們.
看回我們今天自己的信徒的生命.
我們是走一個什麼樣的方向呢?.
上帝的榮耀.
怎樣在藉著我們建立的時候.
被人認得到.
成為一個報道的動力呢?.
第二方面我想和大家透過這段聖經所看的.
第三十三節和第三十六至三十八節.
在這裡我們可以透過.
在屬靈的操練的過程當中.
去避免我們被感性的事情.
去操控了我們.
為什麼我要提出這一點.
也為什麼在這段聖經.
《約翰福音》要提出這一點呢?.
因為當我們很容易被感性去操控的時候.
我們就會變得.
很容易離開上帝那份榮耀.
我們一起看看這段聖經說些什麼.
我們一起讀第三十三節和第三十六至三十八節.

$^{721}$小子們,我還有不多的時候與你們同在.
後來你們要找我.
但我所去的地方你們不能到.
這話我曾對猶大人說過.
如今也照樣對你們說.
第三十六節.
西門彼得問耶穌說.
主往哪裡去?.
耶穌回答說.
我所去的地方.
你現在不能跟我去.
後來卻要跟我去.
彼得說.
主啊,我為什麼現在不能跟你去.
我願意為你寫明.
耶穌說.
你願意為我寫明嗎?.
我實實在在地告訴你們.
稽叫而先.
你要三次不認我.
在聖經裡面再次強調.
耶穌說我要離開你們.
和你們很快會離開.
哪裡說過的話?.
十三章第一節剛剛說過.
說完又說.
其實耶穌.
在這裡反映出.
祂有一種很強的迫切感.
所以大家要明白.
這裡不在今天說到說.
但為什麼耶穌要去到.
赫西瑪利園的時候禱告.
祂的禱告流出來的汗.
如血一樣.
是不是真的血.
不是這樣的意思.
是說到那種迫切感.
那種掙扎是很激烈的.
在這個情況之下.

$^{761}$我們看到插入了一段插曲.
在一個新誡命說完之後.
彼得就走來跟耶穌說.
你去哪裡?.
因為他前文說.
耶穌說你要走.
這麼迫切.
但耶穌說多了一點.
給我們知道.
後來你們要找我.
但我所去的你們不能到.
這話我曾經對猶太人說過.
如今也照樣對你們說.
這個說話是回應.
在約翰福音第七章.
第三十三至三十五節的經文.
如果有我們打開看看.
約翰福音第七章第三十三至第三十五節.
我們由第三十二節讀起.
我們一起讀.
法利賽人聽見眾人為耶穌這樣紛紛議論.
就師長和法利賽人就打發差役去捉拿他.
於是耶穌說.
我人還有不多的時候和你們同在.
以後就回到差我來的那裡去.
繼續讀.
你們要找我卻找不著.
我所在的地方你們不能到.
猶太人就彼此對問說.
這人要往那裡去.
叫我們找不著呢.
難道他要往散著希臘中的猶太人那裡.
去教訓希臘人嗎.
他說你們要找我卻找不著.
我所在的地方你們不能到.
這話是什麼意思呢.
我不知道當你如果零修讀到這段聖經.
你會覺得是不是一頭霧水.
你一定不知道他講什麼.
你又去那裡你又不知道去那裡.

$^{801}$然後猜耶穌就去希臘那裡去教訓希臘人.
你根本不知道他想講什麼.
不過讀完剛剛我們這段聖經的時候.
我們就會明白到.
做了一個很清楚的對比.
同樣類似的說話.
耶穌對著門徒說.
耶穌怎樣對著門徒說.
他是小子.
小子這個字.
如果你看英文翻譯叫做Children.
希臘文叫做Teklon.
是什麼意思呢.
唉!小子.
我爸爸也是這樣叫我的.
唉!小子.
你不要這樣叫我.
就是你未成熟.
不是不行.
我看起來的不過.
不太行.
所以我爸爸這次過來加拿大.
我看我女兒結婚.
然後我和他去旅行.
他都沒有叫過我一句叫小子拳.
長大了.
像樣了一點.
不過父母永遠看子女都是小孩子.
小女孩.
未成熟.
但是.
原來這番說話在先前和猶太人說過.
但是那班人怎麼想耶穌呢.
剛才讀過了.
讀的時候.
我不知道你會不會感受到那班.
祭司長,文士,法利塞人說那番說話.
說耶穌都不知道和什麼人.
是不是他想走去希臘那裡.
教訓那些希臘人.

$^{841}$希臘人是什麼.
當時認為是學術知識最豐富的文化代名詞.
其實是充滿諷刺的義氣.
猶太人.
即是先前那班祭司長.
其實對耶穌完全不客氣.
他認為什麼呢.
耶穌你很厲害嗎.
你去那裡和一些厲害的人訓練吧.
他其實完全不明白.
耶穌所說到他要離開的去哪一個地方.
耶穌在這段聖經就要說得明白給門徒知道.
所以我們要有這段聖經.
我們能夠更加清晰了解.
讀到第七章的時候.
我們只會看見這班人是懵的.
不知道耶穌想說什麼.
你如果讀到那段聖經.
你得到的答案就是.
其實這班人都不知道.
耶穌說要說什麼的地方.
他們是懵的.
你得到這個答案.
你就是讀到那段聖經的意思.
但當你讀到第十三章的時候.
你讀到這裡.
你就會知道這班人是懵的.
不過耶穌不想那班門徒懵.
所以耶穌就會和他解釋.
但怎樣介紹.
就用了一個經常都很聰明的門徒.
叫做彼得.
讓彼得來發問一條問題.
而耶穌就用這個機會來帶他們去看.
耶穌怎樣帶他們去看.
耶穌回答.
我所去的地方你現在不能跟我去.
後來卻要跟我去.
不是不能去.
不是現在去.

$^{881}$聽清楚多一次.
不是不能去.
不過不是現在去.
有什麼重要的牧師這個說法.
記不記得剛剛走了那個猶大.
留意到他走的時候.
聖經說了一句什麼話形容他.
看多一次.
第十三章第三十節.
猶大受了那點餅.
立刻就出去.
聽清楚.
猶大受了點餅.
立刻就出去.
猶大做任何事要立刻見果孝.
耶穌叫門徒做事.
不是沒有果孝.
不過果孝是按照耶穌的時間表來發展.
換言之都是回到那句話.
父啊我已按照你的旨意行事.
你我榮耀了你.
你也榮耀我.
今天很多基督徒.
不按上帝說話做事.
就要按照他自己的時間表來做事.
要按照他自己的結果來做事.
並且我們會留意到.
今天很多的信徒做事.
都要立刻見果孝.
我對靈因派不是討厭.
不過我是對於靈因派那種.
凡事都要立刻有感覺.
立刻收果孝.
是很存疑的.
聖經裡你留意一下.
做所有事無論先知也好.
無論後來的仕徒門徒也好.
以至我自己屬靈的經驗也好.
我們從來有沒有試過.
要強迫上帝要立刻見效.

$^{921}$如果你要立刻.
往往那些都不是好的結果.
這裡做了一個很清楚的對比.
西門彼得他要想好似.
就快了.
跌入好似猶太一樣.
你現在要去.
你就快要去.
而且講到很大.
怎麼大.
他說.
主啊.
我為什麼現在不能跟你去.
我願意為你捨命.
我們先不要理前頭那句.
他講到現在.
誓神癖願都會和上帝和耶穌一起去.
我和你赴湯蹈火.
在所不計.
我捨命陪君子.
這個好像是耶穌在這個時候見到的彼得.
但這個正正就是彼得最大的缺陷.
在時間的上面.
他要現在.
他在表達的上面.
他要立刻和耶穌有一個緊密接觸.
後者不是大問題.
我們喜歡和上帝和耶穌一起.
尤其是你喜歡的人是走在一起.
不是一件大的問題.
不過.
你有沒有聽那個人給你的指示.
是重要的.
有一次有個網上的會友.
和我這樣講.
他說.
你的會友在二馬內環浸算會和你過個禮拜聚會就好了.
我就不行了.
我遠在很遠很遠的地方.
我就問有什麼好呢.

$^{961}$他說我也想.
因為他立刻問到你的事.
那你現在也可以問到我的事.
不同的.
我不知道有什麼不同.
不過我自己覺得至少一件事.
他想和我有一個緊密的接觸是一件好事.
但是.
你如果和我接觸.
麻煩一起按照上帝的心思教導來做事.
因為我是按這種東西.
按這些基礎來做事.
不是只有嘴巴.
只有一種熱心來去講.
還有.
他說我為什麼現在不能跟你去.
這個問題.
比起.
加勒人猶大是最截然不同的地方.
加勒人猶大的特點就是.
他不用和主討論.
他已經落了conclusion.
彼得不同.
彼得會聽完主的說話會問多一個問題.
為什麼不行.
或者為什麼.
why or why not.
然後他等耶穌給答案.
弟兄姊妹如果我們做一個基督徒.
你經常說想有些屬靈操練.
我想這個就是一個好的操練.
當你知道上帝講一些東西給你知道的時候.
尤其是在讀聖經靈修的時候.
而你有所啟發對於你自己人生有所啟發.
你應該問為什麼.
或者為什麼不行.
在這個思想的過程當中會幫助到你去.
摸索清楚上帝的心思.
俗語都有句.
停一停想一想穩陣很多.

$^{1001}$你問多一條問題.
對於你的屬靈生命的影響.
可能深度完全不同.
而這個就是彼得和.
加勒人猶大最大不同的地方.
耶穌給答案.
彼得.
耶穌給的答案是什麼.
願意為我寫明嗎.
這條問題的答案是什麼呢.
當有人問到這條問題.
回應你所說的話的時候.
應該的答案就是.
我是不是真的會寫明呢.
這個不是答案.
這是一個繼續反思的循環問題.
然後耶穌向彼得這樣說.
我實實在在地告訴你.
阿們阿們.
雞叫已鮮.
你三次不認我.
當時正正就是在這個.
晚上的時間.
正在吃這個真是如意節餐還是預備餐.
這些我們不討論了.
在這個餐的時候.
已經入夜.
耶穌被捕的前夕.
耶穌科斯前面的事情.
彼得應該在這個時候會問.
耶穌說得出一些東西.
不會是一定會落空的.
耶穌說得出的東西.
沒有試過落空的.
你看看整個約翰福音書裡面.
所有的神蹟.
耶穌說過的事情.
做過的事情.
從來沒有落空的.
那就是我會在雞叫已鮮.

$^{1041}$三次不認主.
那我是不是真的寫明.
會願意和主寫明呢.
我相信當彼得在這個時候.
最大的難題就是.
他感情上和主所決定的那種矛盾.
他很難融合.
他很想和耶穌在一起.
但是耶穌又說他不行.
而耶穌也沒有落空過的時候.
既然如此.
那我是不是真的會寫明呢.
似乎就不會了.
不會的時候.
那我是什麼價值的人呢.
是什麼靈性的門徒呢.
我是不是一個真正相信主.
跟隨主的門徒呢.
耶穌真的要他問這條問題.
我們作為基督徒.
我問你,你是不是真的相信主呢.
尤其我們做牧師.
你是不是相信主呢.
弟兄,姐妹,你是不是相信主呢.
你會怎麼想呢.
是的,我是相信主的.
還是牧師問我這句問題.
是不是我有些事情不妥呢.
弟兄,姐妹.
對自己的信不斷地去作出這條問題.
不是不信.
這條問題是你作為信徒.
應該長期要去問的問題.
因為當你曉問這條問題的時候.
你就會不斷去監察.
讓上帝的靈在你裡面.
去參與你的生命的更新進程.
教會很多時候說所謂的更新變化.
又說事奉主.
流於這些聖經字眼的引述.

$^{1081}$但是如果真的要說更新變化.
你的靈性真的會成長的.
其中一個問題就是.
你對你自己的靈性不斷要有一個減省.
大家知道你們牧師是做 Auditor 做核數員.
核數的目的不是要去湊成.
核數的目的是想自己知道哪些地方有缺陷而去補足.
核數的目的不是要去罰款.
核數的目的是要將漏洞堵塞.
靈性的成長其實就是一個長期的核數的狀態之下.
不過這個核就不是數目字.
是你的靈性的狀態.
一個人決定由信主.
表白是信主那一刻開始.
其實就是進入了這個運作的裡面.
而這個運作就能夠保證到我們去學會怎樣去滿足.
是耶穌剛剛給我們的新的誡命.
你是不是真的能夠做到彼此相愛呢?.
不做到的,為什麼做不到呢?.
做不到是不是很丟臉呢?.
還是做不到,你想竭力去改進去做到呢?.
當你肯這樣做的時候.
包括弟兄一,姊妹二,弟兄三,姊妹四.
大家都曉這樣想的時候.
這個群體就是願意彼此成長.
而這個成長就要在之間大家互相提醒.
付出最好的為他的靈性著想.
這就是耶穌所指的照我愛的定義去定.
耶穌對人的愛和上帝對人的愛.
從來都是先賦予的,先付出的.
為你的靈性最好的地方而願意騰出,付出.
看似好似損失,其實卻是一個一同建造的過程.
我希望這段聖經引起大家去想想.
到最後,誰是最大的輸家?.
加勒人猶大.
他最蠢的就是去到捉到一隻鹿都不會脫殼.
他竟然覺得要去找一匹驢會好一點.
就是這樣的狀態.
所以他願意自己去奔跑他自己的人生.
最後他的下場就好像他一直的心態.

$^{1121}$他馬上出賣了耶穌,馬上收到錢.
馬上就吊頸死,他的生命馬上.
但是基督徒的生命是在耶穌基督裡面細水長流.
當時候到,耶穌就帶我們去到他要去的地方.
現在不行,因為這個時候就是給你和我這個屬靈舞台上面要演好這個角色的時間.
所以聖經有一句說話.
耶穌這麼快回來不是擔言.
而是一種對於我們的一種恩賜,忍耐.
讓我們在這個忍耐裡面成長起來.
讓我們在屬靈生命上面有這個方向.
我們就可以配得起在新天新地裡面同主同活.
未必你說去到見主被提那天.
我當被提,我們說這些神學性的說法.
你做得已經是圓滿,沒有一個人圓滿.
每一個人做的程度成熟程度都可以不同.
但是上帝看的就是你有同一個生命.
這樣就足夠進入到天國裡面.
所以你說牧師你堂堂道道說得這麼厲害.
我不行的,我進不了天堂的.
不是,你做到你做到的地步就夠了.
那你的地步是什麼?不做嗎?不是.
你要做到你能夠上帝給你的屬靈能力.
恩賜達到的東西,這樣就足夠了.
至少我們在聖經裡面看到的.
耶穌是這樣幫助門徒.
在最後一個時間裡面去塑造一個這樣的群體.
我希望以馬萊尼環浸信會.
無論是先上的會友或是線下的會友.
不是線下,是實體的會友.
我們都是有這份生命,能夠一起去建造.
希望我們能夠在主的裡面朝向一個將來.
耶穌也說,現在你不能到,但卻要最後跟他去到.
讓我們一起從心低頭度過.
父神多謝你讓我們在這段聖經裡面學習.
讓我們看到你對我們的關懷.
讓我們看到一個墮落的生命.
是你給足的機會,祂怎樣仍然錯過.
然後求你讓我們每一個聽到這個訊息.
讀到這段聖經的弟兄姊妹.
我們都能夠好好的把握.

$^{1161}$能夠一同在這份救恩的上面.
能夠參與在這個榮耀的當中.
真正的榮耀不是一個只是一件禮物.
送給上帝為你去沾花,拉紅的一個狀態.
而是我們真真正正在我們的生命上面.
能夠按照你所指出彼此相愛的建立.
讓我們在地上讓人認出我們是你的門徒.
從而讓上帝你的名可以讓人認識得到.
我們同心的禱告.
奉靠主耶穌基督得勝名至於靠.
阿們.
\newpage



\section{帖撒羅尼迦前書 1:1-10}
\label{sec:032288}
\textbf{培育生命,建立教會:感恩有你!(粵/英)}~~~~梁康民牧師
\newline
\newline
連結: \href{http://www.church.com.hk/acms/content.asp?site=cdc\&op=show\&type=product\&code=032288\&layout=sermon}{\texttt{032288}} ~~~~ 語音日期: 2026-06-14
\newline
\newline
\hyperref[sec:032287]{\small{< < < PREV SERMON < < <}}
~
\hyperref[sec:index_chronic]{\small{[返順時目]}}
~
\hyperref[sec:032289]{\small{> > > NEXT SERMON > > >}}
\newline
\newline
帖撒羅尼迦前書 1:1-10
\newline
\begin{longtable}{cl}
\hline
\hline
章節 & 經文 (和合本修訂版)\\
\hline
1:1 & \begin{tabularx}{0.7\textwidth}{X} 保羅、西拉、提摩太寫信給帖撒羅尼迦在父神和主耶穌基督裡的教會。願恩惠、平安歸給你們！ \end{tabularx} \\ \\ \relax
1:2 & \begin{tabularx}{0.7\textwidth}{X} 我們為你們眾人常常感謝神，禱告的時候提到你們， \end{tabularx} \\ \\ \relax
1:3 & \begin{tabularx}{0.7\textwidth}{X} 在我們的父神面前，不住地記念你們因信心所做的工作，因愛心所受的勞苦，因盼望我們主耶穌基督所存的堅忍。 \end{tabularx} \\ \\ \relax
1:4 & \begin{tabularx}{0.7\textwidth}{X} 神所愛的弟兄啊，我知道你們是蒙揀選的； \end{tabularx} \\ \\ \relax
1:5 & \begin{tabularx}{0.7\textwidth}{X} 因為我們的福音傳到你們那裡，不僅在言語，也在能力，也在聖靈和充足的確信。你們知道，我們在你們那裡，為你們的緣故是怎樣為人。 \end{tabularx} \\ \\ \relax
1:6 & \begin{tabularx}{0.7\textwidth}{X} 你們成為效法我們，更效法主的人，因聖靈所激發的喜樂，在大患難中領受了真道， \end{tabularx} \\ \\ \relax
1:7 & \begin{tabularx}{0.7\textwidth}{X} 從此你們作了馬其頓和亞該亞所有信主的人的榜樣。 \end{tabularx} \\ \\ \relax
1:8 & \begin{tabularx}{0.7\textwidth}{X} 因為主的道已經從你們那裡傳播出去，你們向神的信心不只在馬其頓和亞該亞，就是在各處也都傳開了，所以不用我們說甚麼話。 \end{tabularx} \\ \\ \relax
1:9 & \begin{tabularx}{0.7\textwidth}{X} 因為他們自己已經傳講我們是怎樣進到你們那裡，你們是怎樣離棄偶像，歸向神來服侍那又真又活的神， \end{tabularx} \\ \\ \relax
1:10 & \begin{tabularx}{0.7\textwidth}{X} 等候他兒子從天降臨，就是神使他從死人中復活的那位救我們脫離將來憤怒的耶穌。 \end{tabularx} \\ \\
[1ex]
\hline
\hline
\end{longtable}
$^{1}$頂尖輩平安!.
今日我們紀念恩傳浸信會十九週年的紀念.
培育生命,建立教會.
最緊要是什麼?.
感恩有你啊!.
我想問一下,恩傳十九年了.
有幾多頂尖輩在這裡呢?.
在這裡的,已經在這裡十九年了.
嘩!在這裡已經十九年了.
意思是什麼?.
他們是建立恩傳的第一批弟兄姊妹輩.
可不可以請你們站起來給我們看看.
第一批的.
感恩有你.
請坐.
要建立教會.
當我們,譬如說我們一會兒都會切蛋糕,唱生日歌.
我們祝福恩傳浸信會.
恩傳浸信會是什麼來的?.
是祝福有你啊!.
Happy Birthday 是每一個一起有份的.
耶穌基督的大使命.
大家都很熟悉了.
有兩個重點.
第一,就是使萬民作耶穌基督的門徒.
使萬民建立教會.
還有就是要教導門徒遵守耶穌基督的教訓.
就是培育生命.
所以耶穌基督的大使命.
不是只是一個口號.
是要去做的.
我們要有行動的.
我們過去教會的歷史.
都是講到耶穌基督的大使命是要去實踐的.
什麼是教會呢?.
建立教會.
和建立教堂.
在英文,church又可以是教會.
但亦都是一個教堂.
什麼是教會呢?.

$^{41}$教會如果,舉一個例子.
這個地方.
我們今天都坐在這裡.
參加完崇拜.
那就散會了.
大家都出去了.
關了燈.
關了門.
這個地方是不是教會呢?.
這個地方不是教會.
這個是教堂.
明白嗎?.
你看看聖經怎樣講到教會.
我們要建立教會.
在提摩太前書三章十五節裡就說了.
這家就是永生神的教會.
這家,這個是這家.
我喜歡一個講法.
當我們講恩泉浸信會的時候.
我喜歡講是恩泉家.
恩泉家.
就是永生神的教會.
這裡有兩幅圖是很清楚講解教會是什麼意思.
左邊那幅圖.
就是我們再一起就成為教會.
我們一起就是教會.
所以在聖經的原文.
這個教會這個字.
就是被神呼召聚集一起的意思.
有一個教會的名稱.
叫做神召會.
這個翻譯是很正確的.
就是被神召聚的人.
所以我們一起就是教會.
右邊那幅圖.
去教會或者是回來教會.
我們今天是不是回來教會.
但我們不是去了一個地方.
我們回家.
懂嗎.

$^{81}$這家就是永生神的教會.
我們在一起才是教會.
所以是一個教會的成長.
在剛才吳牧師的短片裡面的介紹.
由七十人開始到今天接近四百人.
這個家一路成長.
感恩有理.
十九週年.
我們值得感恩的.
一起關於去建造.
所以耶穌基督的大使命.
從保羅開始.
他領受從天上來的意象.
他就不斷的去傳福音.
去建立教會.
這個地圖就表示了他很多次宣教的旅程所去過的地方.
今天你和我.
都是領受著神的大使命.
我們看一看保羅的經歷.
和我們今天所讀的這段經文.
那裡是帖撒羅尼迦這個教會.
是怎樣開始的.
這個是《試圖文宣》第十七章.
是保羅第二次的宣教旅程.
順便他們說保羅和塞拉經過了很多很多的地方.
就去到帖撒羅尼迦這個地方.
這裡就看到.
沒有了.
這個看不到.
我們看一看這個地圖.
圓圈那裡就是帖撒羅尼迦所在的地方.
圓圈那裡就是帖撒羅尼迦.
在右邊就看到馬其頓的異象.
他們由亞洲去到歐洲.
就是第一個地方去到菲律賓.
他們記住好幾個城市就去到帖撒羅尼迦.
帖撒羅尼迦是馬其頓省首府的首都.
那裡去到猶太人的會堂.
在當中傳道.
一年三個安息日.

$^{121}$就傳耶穌基督從死復活.
我們所傳的這個耶穌就是基督.
這個教會是保羅離開了之後.
因為在帖撒羅尼迦那裡.
那裡的猶太人要逼迫他.
捉了他要拉他去結官.
宣教都不是容易的.
去到那裡就遇到好多的反對.
所以門徒就說你快點走.
他們去到一個地方.
初立的教會.
有多少人.
很少人的.
初初幾個月.
大家知道初期教會都是家庭教會.
不會好像一個這麼大的建築.
坐得幾百人在那裡.
所以最初都是十幾二十人的信徒.
他們在哪裡聚會呢.
就去到其中一個信主的人的家裡聚會.
那個人的名字叫做耶孫.
那些人去到耶孫的家裡要捉保羅.
他們就說你快點走.
保羅就離開了帖撒羅尼加.
到南走.
就去到希臘.
然後就去那裡繼續傳福音.
然後就去到哥林多.
所以離開已經好幾個月了.
那究竟帖撒羅尼加.
我離開了幾個月.
發生什麼事呢.
在這麼受逼迫的環境裡.
他們可不可以生存呢.
他們心裡面好擔心.
他們就不斷地祈禱.
剛才的經文就說他們不斷地祈禱.
後來尼迪提摩泰就去到哥林多.
就告訴保羅.
帖撒羅尼加教會現在的狀況.

$^{161}$那保羅就很開心了.
他感恩神的保守帶領.
但這個教會是在困難中成長的.
在反對逼迫當中不斷地成長.
這個教會有多大呢.
都是十幾二十人開始的.
恩傳多少人開始的.
剛才看完短片有沒有印象.
大概多少人開始的.
大概七十人開始的.
比帖撒羅尼加教會還要多.
但是帖撒羅尼加的教會.
他們存在的環境是怎樣呢.
是受逼迫的.
是受反對的.
但是他們仍然在困難當中去見證神的恩典.
不斷地傳福音.
所以剛才所讀的經文.
他們傳到馬其頓阿戈雅各地方的人都知道你們的信心.
他們聽到你們信主之後生命的改變.
雖然我離開了這麼久.
這些消息都傳到我的耳中.
是不是啊.
教會在社區裡面的見證.
不是帖撒羅尼加這麼短的範圍.
馬其頓傳承.
阿戈亞傳承.
這麼廣闊的地方.
所以帖撒羅尼加教會可以成為我們的榜樣.
它是一個模特兒教會.
是一個示範的教會.
所以波羅就提到幾件事.
他為了他們的感恩.
有三件事.
有沒有印象.
為他們感恩有三件事.
第一.
就是為他們.
因信心所做的工作.
第二就是.

$^{201}$因愛心所受的勞苦.
第三就是因盼望我們主耶穌基督所傳的堅忍.
如果用三個字.
代表這三樣東西.
是什麼.
信.
愛.
望.
是不是.
在聖經裡面的保羅的書信.
信望愛這三個字.
你一定找到的.
基督徒生命的特質.
有信心.
有愛心.
有盼望.
信望愛.
大家最熟的哥林多前書十三章.
現在不常傳的又盲又愛.
最大的是什麼.
信望愛.
貼索了教會的特質.
成為我們的榜樣.
因全家.
我們的建立.
都需要靠信愛望愛.
我們要看這三方面.
究竟我們有沒有信.
信不是一個靜止的字.
而是有行動的.
他說你是因信心所做的工作.
所以信心是要加工作的.
雅各書就是這樣說.
我們只是說信但是沒有行為.
那就是死了.
沒有用的.
要有一個活的信心.
教會的建立需要用信心.
信靠神的供應.
因全家是怎樣來的呢.

$^{241}$有多少人是這五年之內才來的.
這五年之內才來的.
有很多是這五年之內才來的.
剛才吳律師也說了一部分的片段.
重要的事說三次.
不厭其嘗讓大家去回顧一下.
因傳浸信會的誕生是信心的工作.
中間的三段.
北約懷仁浸信會.
黎明山懷仁浸信會.
還有因傳浸信會.
我們是第三代.
我們的母會就是黎明山浸信會.
我們的祖母.
祖母就是北約懷仁浸信會.
這是教會的誕生.
左右兩個字你翻譯一下.
示明教會的見證.
左邊這裡.
就是憑信展行大示明.
我們建立教會.
這幅圖我們先看看.
北約懷仁浸信會.
那個年份是1970年.
是第一次主日崇拜的年份.
當然是經過籌備.
幾個家庭開始祈禱.
預備設立教會.
一路發展的時候成長很快.
所以去到1984年的時候.
即是14年之後.
教會一路成長成長到坐不下.
開始準備要席堂.
當時有80個成年人願意參與.
剛才也有一些弟兄姊妹說.
你有參與.
由北約懷仁浸信會出來席堂.
有哪些弟兄姊妹.
是由北約出來的.
Pastor Leo.

$^{281}$Pastor Karen.
開拓先鋒.
第一代席堂.
第二代席堂.
去了泰明山之後再席堂.
他們又成為席堂的先鋒.
他們的長輩也參與事奉.
Julian.
Iris.
他們的一家也事奉.
最初我們在OBC.
在道社聚會.
當時是劉達榮傳道成為席堂領導的牧者.
由1988年開始第一次聚會.
當時他們要邀請一位主任牧師.
他們就回香港找主任牧師.
找了誰.
找了我.
我就成為.
黎明山懷仁浸信會的第一個主任牧師.
我1991年開始在黎明山.
當時我們借一個神學院聚會.
我們也買了一幅地預備建堂.
那間教會黎明山.
在1996年入伙.
有一個很特別的日子.
第一代在這個新堂行婚禮.
是雷哥.
Auntie Grace.
在這個新堂第一代行婚禮.
我和他們證婚.
他們的子女也長大了.
Jessica.
感恩.
看到生命的成長.
培育生命.
去建立教會.
到黎明山建立了之後.
開始慢慢人數有增長.
你看到那幅圖片.

$^{321}$這是黎明山浸信會十週年紀念的照片.
十週年.
如果是1988年開始.
但是十週年也是2001年.
因為2001年是自立的日子.
所以我們算黎明山自立那一年.
是那個開始.
所以那時候你看圖片.
已經有五百人.
信徒增長得很快.
所以去到2003年的時候.
就要成立一個直堂小組.
要準備再建立另一個教會.
六百人禮堂坐不完.
我們曾經要去到禮堂後面有一個大的課室.
坐不下就在後面那個課室坐著看螢光幕.
停車位又不夠.
我們就去附近找些地方去泊車.
甚至有一段時間.
我們鼓勵我們自己人.
就要泊去那個離開教會三公里的停車場.
然後我們就找人去接他們回來教會聚會.
所以一定要搬走一些人出去直堂.
我1991年開始在黎明山受訪.
到2006年我就退休了.
我做了十五年退休.
因為我1991年去的時候.
我大概五十歲.
我做了十五年退休.
再做多五年顧問務之.
即是多少年.
五十加二十.
那我退休多久.
退休由2010年退休.
退休了十六年.
你心算計不計到我現在幾多歲.
五十加二十加十六.
知道吧.
能夠繼續去事奉祂.
又真人看到恩泉的誕生.

$^{361}$所以恩泉在2007年.
就開始第一次的主要崇拜.
所以我們看到神的恩典.
幾代的教會.
都是實踐耶穌基督的大使命.
是建立健康的使命教會.
亦都是培育生命建立教會.
信心.
信心開始.
因愛心所受的勞苦.
愛心亦都是一個行動.
愛心是要勞苦的.
在一個英文的翻譯.
是因愛的激勵而付出的辛勞.
教會的建立.
培育信徒的生命.
是要出於愛心.
要勞苦.
在哥林多後書五章第十四節.
保羅說因為基督的愛激勵了我們.
要被激勵.
弟兄姊妹.
你的愛要被激勵.
要被基督的愛激勵.
要激勵.
要激活.
我們一起去付出這個愛.
培育生命建立教會.
感恩有你.
我們都要每一位弟兄姊妹的參與.
大家提到我們植糖的歷史.
植了糖之後又怎樣.
我種了些東西在這裡.
很多弟兄姊妹是種瓜種菜.
等它自生自滅.
你要用一口心機.
要淋水.
要除草.
保羅也說我是去栽種.
但是誰去淋水.

$^{401}$他說阿波羅去淋水.
其實是供了一切的配搭.
當我們建立了.
開始了一個教會.
無論你開始.
我們開始由七十人開始.
那怎樣.
今天我們坐在這裡都很舒服.
但你知道哪裡十九年來.
恩泉經歷了什麼.
我們開始借了一間西人的教會.
就好像我們今天借了一間西人的教會.
但是當時不是一個禮堂.
我們是用了他們的球場.
所以我們要搬桌搬椅.
那些音響教那些.
例如麥克風 喇叭 電腦那些.
教會沒有供應.
頂展會要搬回教會那裡用.
搬回家裡.
每個禮拜就這樣做.
辛苦嗎.
麻雀雖小 五臟俱全.
我們又要有主日崇拜.
又要主日學.
又要有小組.
又要有些外展工作.
需要很多人手.
場地的需要.
經費的需要.
我們靈元三席堂的頭三年.
母會有津貼.
開始時需要很多錢.
第四年就要成立.
一切經費都是頂展會愛心的奉獻.
一起的事奉.
恩賜的奉獻.
我們過去幾個禮拜的講題.
都提醒頂展會要立志去事奉神.
我們要發揮我們的恩賜.

$^{441}$我們要做恩賜的好管家.
所以教會的發展.
我們要同心建造.
培育生命.
建立教會.
感恩有禮.
所以每一位頂展會的參與.
我們現在進行建堂.
建立一個教堂.
我們需要很多很多頂展會的參與.
真愛心所受的勞苦.
很多頂展會都很辛苦.
建堂委員會.
很多事情要考慮.
工程.
經費的來源.
要和政府溝通.
要申請許可證.
工地的巡查.
建堂委員會頂展會.
在這裡.
為你們感恩.
感恩有你們.
恩愛心所受的勞苦.
在去年籌款晚宴.
很多頂展會都在.
很勞苦的工作.
特別是吳牧師和吳師母.
他們好幾次回香港.
分享我們建堂的意象.
有一次我遇到吳師母.
我說吳師母辛苦你.
辛苦你了.
她說不辛苦.
怎會不辛苦.
因為愛就不辛苦.
頂展會.
建堂委員會的頂展會.
很辛苦.
不辛苦.

$^{481}$因為出於愛心.
愛神的心.
不辛苦.
頂展會我們要同心建造.
培育公平.
教會.
感恩有你.
我們彼此感恩.
感恩上帝.
亦都是彼此感恩.
感恩有你們.
最後一點.
盼望我們主耶穌基督.
所傳的堅忍.
等候他御子從天降臨.
就是神使他從死人中復活的那一位.
我們的脫離將來憤怒的耶穌.
全心盼望.
盼望甚麼.
主耶穌快回來了.
我們有個迫切感.
因為我們將來都要向主交徵.
我們能夠面對主耶穌基督的時候.
主耶穌會稱讚我們.
有忠心有良善的僕人.
我們一起去努力.
靠著神的恩典.
培育生命去建立教會.
我們一起用信心去工作.
我們用出於愛心的勞苦.
我們因為有盼望.
所以一定要去堅忍.
感恩有你.
有每一位頂主妹在.
我們明年.
主若願意.
我們明年還會在新堂.
舉行二十週年.
的感恩.
好不好.

$^{521}$我們繼續一起努力.
新堂你看到嗎.
還有那幅圖畫.
是楊恩泉拍的.
我們今天都坐滿了.
我們感恩.
我們一起祈禱.
多謝天父恩典.
因為你自己的信實.
和你一切豐富的供應.
讓我們全家能夠一路成長.
十九年走來.
一路都看到.
豐盛的看顧.
我們願意將感恩獻上.
願你得到一切的榮耀.
以聖靈繼續引導我們走前面的路.
主若願意.
我們明年可以在新堂.
舉行感恩崇拜.
奉耶穌基督的名.
Amen.
\newpage



\section{士師記 6:39}
\label{sec:032289}
\textbf{聖經人聖經事:基甸用第二,第三次求證}~~~~黃克勤牧師博士
\newline
\newline
連結: \href{http://www.church.com.hk/acms/content.asp?site=cdc\&op=show\&type=product\&code=032289\&layout=sermon}{\texttt{032289}} ~~~~ 語音日期: 2026-06-11
\newline
\newline
\hyperref[sec:032288]{\small{< < < PREV SERMON < < <}}
~
\hyperref[sec:index_chronic]{\small{[返順時目]}}
~
\hyperref[sec:032290]{\small{> > > NEXT SERMON > > >}}
\newline
\newline
士師記 6:39
\newline
\begin{longtable}{cl}
\hline
\hline
章節 & 經文 (和合本修訂版)\\
\hline
6:39 & \begin{tabularx}{0.7\textwidth}{X} 基甸又對神說：「求你不要向我發怒，我再說一次，讓我用羊毛再試一次，但願羊毛是乾的，遍地都有露水。」 \end{tabularx} \\ \\
[1ex]
\hline
\hline
\end{longtable}
$^{1}$各位姐妹平安,又讀聖經啦.
生命樹鼓勵大家讀聖經.
看聖經人物,聖經故事.
然後看看上帝奇妙大能的手在背後.
和祂的心意.
這次我自己看下去就是.
機電本來小時候一直聽聖經故事.
都說機電很厲害.
帶著三百人打破瓶子大叫.
贏了米電人.
但原來之前上帝是一個人.
我稱呼機電一開始是一個膽小的人.
竟然令他做了很大能的工作.
不過機電為何說他膽小呢?.
因為他經常都害怕,經常都求.
你給我一個證據.
經常求神給他證據.
給他證據做少少,給他證據做些.
都幾好.
不過我覺得上帝在他身上有很多功夫.
對於他很忍耐.
其實我都不給機電大膽.
不過神給我學些學些學些.
一直為他工作.
機電第一次求證據.
然後今次第二次又第三次.
我覺得他第三次是求得最差.
我們看看他求什麼證據.
《史詩紀錄》章33節.
你看看米電人已經很多人了.
這是沙漠民族.
而且米電人的意思就是競爭者.
米電出了名是跟你鬥的.
愈強愈強.
阿瑪利人也出了名.
很厲害,跟摩西打過仗.
東方人,什麼叫東方人?.
約旦河東的所有民族.
包括摩爾阿曼,以東的就是東方人.
都聚集過.

$^{41}$或者由東方過約旦河過來河西.
跟以色列人打仗.
去到河西的北耶穌列平原.
哇,耶穌列平原.
說著說著大家都有機會知道在哪裡.
或者去看看.
耶穌列平原就是.
如果你有看聖經.
最尾的一場重要戰爭之一.
哈米吉多頓.
就在耶穌列平原那裡決鬥了.
你會知道耶穌列平原是很大的平原.
去到那撒勒山上看到耶穌列平原.
哇,很厲害.
就已經聚集了.
安營,即是預備打了.
這個時候第一點.
最重要的是什麼?.
我們靠著非勢力,非才能.
是靠神的靈方能成事.
我中學的時候.
很有幸,很感恩.
在中學的團契就開始學這首歌.
非勢力,非才能,乃靠耶和華的靈方能成事.
耶和華的靈降在基甸身上.
哇,這一點第一點最厲害了.
降在基甸身上他就很有信心了.
他就吹角.
吹角的意思就是呼召.
叫人來集合.
這個集結號.
吹集結號,立刻.
集合.
一吹就「阿彼爾謝族都聚集跟隨他」.
阿彼爾謝就是他本身的族.
他只是族裡面的其中一個人.
年輕人.
竟然一吹號,一召集,一召馬.
整個族都來聚集跟隨他.
這個基礎在哪裡呢?.

$^{81}$這個人很勇敢,厲害.
竟然夠膽叫耶路巴力.
改了一個新名,叫做耶路巴力基甸.
巴力他都夠膽動.
這個人真是跟隨神的.
去將我們那些巴力神要.
好啊,我們跟隨耶和華.
於是整個家族最先響應了.
因為耶和華的靈在他身上.
記住這一句.
不是他號召靈,不是他靚仔,不是號角吹得夠響.
是因為耶和華的靈降在基甸身上.
於是他一吹號,神就叫了很多人來.
記不記得這個羅亞做方舟的時候.
他用不用去找那些動物?.
不用的,那些動物自己來的.
那為什麼會自己來的?.
神叫那些動物來.
神叫他整個家族.
最先我們集合.
由基甸帶我們去跟米殿人.
各個族我們要跟他對抗.
他打發人走遍馬來西地.
馬來西人都聚集跟隨他.
整個馬來西支派又來了.
接著又打發人去阿徹.
西布倫,拿忽他利.
幾個支派的人都一齊來了.
其他支派不叫的?猶大那些不叫的?.
遠啊.
一叫了附近幾個支派.
他們都出來與他會合.
哇,這個集結號一吹就厲害了.
馬上大家都出來了.
很多人.
因為大家都受米殿人壓迫了很久.
很慘.
再加上有耶和靈在他身上.
一聽到他叫.
大家就說好,我們出去吧.

$^{121}$在這個時候你看到基甸對神說.
你若果照著所說的話.
藉我手去拯救以色列人.
我就想你再給我一些證據.
唉,又來?.
這個時候神的靈在他身上.
已經叫他吹號.
幾個支派和自己的支派都來了.
來了之後預備要打.
忽然之間基甸的心是怎樣呢?.
怯不怯?.
我自己的解釋推測,怯不怯?.
接著他就跟神說.
如果真的你藉我手拯救以色列人.
不是吧?如果我跟神說.
你不是吧?現在才說這些?.
叫這麼多人來才跟我說.
好像叫了所有的夥計來.
然後跟老闆說話.
是不是真的請我們先?.
這些說話很有趣.
如果我覺得有一點點要脅.
那些人都來了.
給個證據來.
是不是先出自己的糧?.
這樣都可以?.
如果我覺得基甸真的在這裡.
不是很好.
我的解釋而已.
好了,如果是這樣.
你照我這樣做.
還自己開了證據要求跟著自己做.
要上帝跟自己.
那究竟你是主還是我是主呢?.
我就把一團羊毛放在和牆上.
若單是羊毛上有露水.
別的地方是乾的呢?.
我就知道是你會照著所說的話.
藉我手拯救以色列人.
就是說我給一個要求你.

$^{161}$你做個神蹟讓我看看.
如果有這個神蹟.
我就相信是你叫我.
人都叫了幾個支派.
半個以色列都來了.
然後你說你做個神蹟讓我看看.
我才出發.
不是吧?.
如果我是耶和華神.
真的有東西打下去.
我叫別人不用你了.
我叫別人都可以.
你也只有你.
但是神對他真的很體恤.
很憐憫.
很溫柔.
很體寵.
大家明白嗎?.
我放一團羊毛在野外和牆.
羊毛上我要明天早上有露水.
這個神蹟很大.
早上有露水.
不過只是羊毛是濕的.
整個地都是乾的.
有露水整個地都是濕的.
為什麼只有羊毛是濕的?.
所以只有羊毛是濕的.
所有地方都要乾.
這個神蹟很大.
明天早上要沒有露水.
除了我的羊毛.
這樣我就知道.
真的是你叫我拯救以色列人.
這個神蹟的要求.
我認為實在太大了.
如果我是神.
我就忍不住.
我不用你了.
我叫別人.
不用機電我叫電機來.

$^{201}$叫別人不用你了.
結果明天早上機電起來.
見果然是這樣.
神父跟他說話就做給他.
將羊毛捏一捏.
嘩!很多水.
其他地方是乾的.
這個神蹟超級大.
誰知機電又對神說.
又來?.
求你不如向我發露.
不要生我的氣.
為什麼不生氣?.
聽到也生氣.
第二個神蹟.
第二個證據我已經生氣.
我再說.
反過來.
這次讓我.
第三點要看.
將羊毛再試一試.
但願羊毛是乾的.
別的地方都有露水.
嘩!這件事真是.
明天早上我再看看.
我放進來的羊毛藥.
是沒有露水的.
其他全部地方都是濕的.
嘩!怎麼做?.
結果.
這夜神也如此行.
讀羊毛尚是乾的.
別的地方都有露水.
我真的沒辦法明白.
神為什麼這麼低聲下氣.
會肯跟機電這樣做?.
是不是沒有他不行?.
找別人也不行?.
找機電的弟弟也不行?.
找別人的支派也不行?.

$^{241}$其實我敢跟他爸爸做也行.
為什麼要他?.
要我玩這樣的兩個神蹟.
兩個神蹟都是超大.
於是.
就是這樣發生.
又照做.
第二第三次.
最離譜我認為是第三次.
大家記得嗎?.
因為已經做了第二次.
全部地方是乾的.
只是羊毛是濕的.
這個神蹟已經超級大.
然後你還敢做一次相反的.
第三次求證明.
我想.
神真是很寬容.
禱告.
多謝我們的主.
我們在這裡看.
很不歡喜機電.
為什麼這麼不敬.
可以對你有這樣的要求.
其實第一次.
已經向祂顯現了.
然後.
支像掃一掃羊肉和細麵餅.
毛餃餅.
已經有火在石上燒了.
然後神的使者就不見了.
這個神蹟已經很大.
一塊石怎麼會出火呢?.
怎麼會將那些祭物燒盡呢?.
這個使者怎麼會不見了呢?.
很大的神蹟.
然後又來第二個神蹟.
一吹號.
那些人就來集合了.
幾個支派.

$^{281}$半個以色列都來了.
但是他仍然不信.
還要再問一次.
再給我一次證據.
一個很偉大的羊毛與露水的神蹟.
到這裡夠了.
不過他竟然說.
神呀,你不要生氣.
再做一次相反的給我看.
主人,我只能夠說.
就是機電面對著很大的對手.
米電人心裡面很害怕.
一個神蹟不夠.
兩個神蹟不夠.
給多一個證據來.
主人,我們求你幫助.
再想深一層.
其實我們並沒有比機電厲害很多.
事事的心裡面就很害怕.
不行,不行.
最大的敵人就是自己的恐懼.
因為很害怕.
主人多謝你.
你是體恤機電.
用你自己大能的手.
用你的神蹟告訴他.
不用害怕.
有你在這裡.
你的同在就是平安.
就是Shalom.
所以可以坦然無懼.
面對著那些敵人.
真的好像沙那樣多的米電人.
和各個部族.
主人你自己是幫助了他.
繼續下去.
我們看到你要做更大的神蹟.
其實這些一步一步.
原來我們再看.
你是預備著.

$^{321}$培訓著.
栽培著機電的信心.
之後他再沒有求什麼試驗.
再沒有求什麼證據.
他就一心相信跟隨你.
行在神蹟之中.
看到你怎樣去幫助他.
看到你自己怎樣擊敗米電人.
主我們都需要.
有這些人生裡面的學習.
在很害怕的時候.
主啊.
你與我們同在給我們平安.
多謝我們的主.
你體諒機電.
同樣求你體諒我們.
求你不要向我們發怒.
當我們失去信心的時候.
求你同樣去加添我們的信心.
垂聽祈禱.
奉耶穌名求.
阿們.
\newpage



\section{羅馬書 5:1-8}
\label{sec:032290}
\textbf{看透我,依然愛我}~~~~林玉英牧師
\newline
\newline
連結: \href{http://www.church.com.hk/acms/content.asp?site=cdc\&op=show\&type=product\&code=032290\&layout=sermon}{\texttt{032290}} ~~~~ 語音日期: 2026-06-14
\newline
\newline
\hyperref[sec:032289]{\small{< < < PREV SERMON < < <}}
~
\hyperref[sec:index_chronic]{\small{[返順時目]}}
~
\hyperref[sec:032291]{\small{> > > NEXT SERMON > > >}}
\newline
\newline
羅馬書 5:1-8
\newline
\begin{longtable}{cl}
\hline
\hline
章節 & 經文 (和合本修訂版)\\
\hline
5:1 & \begin{tabularx}{0.7\textwidth}{X} 所以，我們既因信稱義，就藉著我們的主耶穌基督得以與神和好。 \end{tabularx} \\ \\ \relax
5:2 & \begin{tabularx}{0.7\textwidth}{X} 我們又藉著他，因信得以進入現在所站立的這恩典中，並且歡歡喜喜盼望神的榮耀。 \end{tabularx} \\ \\ \relax
5:3 & \begin{tabularx}{0.7\textwidth}{X} 不但如此，就是在患難中也是歡歡喜喜的，因為知道患難生忍耐， \end{tabularx} \\ \\ \relax
5:4 & \begin{tabularx}{0.7\textwidth}{X} 忍耐生老練，老練生盼望， \end{tabularx} \\ \\ \relax
5:5 & \begin{tabularx}{0.7\textwidth}{X} 盼望不至於落空，因為神的愛，已藉著所賜給我們的聖靈，澆灌在我們心裡。 \end{tabularx} \\ \\ \relax
5:6 & \begin{tabularx}{0.7\textwidth}{X} 我們還軟弱的時候，基督就在特定的時刻為不敬虔之人死。 \end{tabularx} \\ \\ \relax
5:7 & \begin{tabularx}{0.7\textwidth}{X} 為義人死，是少有的；為仁人死，或者有敢做的。 \end{tabularx} \\ \\ \relax
5:8 & \begin{tabularx}{0.7\textwidth}{X} 惟有基督在我們還作罪人的時候為我們死，神的愛就在此向我們顯明了。 \end{tabularx} \\ \\
[1ex]
\hline
\hline
\end{longtable}
$^{1}$各位弟子朋友安安.
請大家再一次低頭祈禱.
因為你揀選我們成為你的孩子.
我們知道你很深愛我們.
想引誘聖靈去打開我們的心.
讓我們能夠明白你對我們的說話.
你知道我們的光景.
你知道我們的需要.
求天父你就將你的話去深植我們的心.
也願我們能夠去進行.
因為這是我們生命要走的路.
多謝你聽我們的祈禱.
藉著耶穌基督的命求.
阿們.
今天我想跟大家分享一個.
我們每一天都會面對的實況.
不論你是學生哥.
還是剛剛進入職場的年輕人.
還是只是三四五十歲.
正是面對家庭,事業.
兩邊都要背重的成年人.
甚至你好像我這樣.
已經進入身體的功能漸漸衰退.
面對晚年生活的長者.
我們都活在一個.
一起活在一個很巨大的社會遊戲規則裡面.
這個規則告訴我們.
你的價值有沒有用.
取決於你擁有多少東西.
年積的時候.
我們的學歷.
我們在哪間公司裡面上班.
我們的薪金有多少.
成為我們自己和其他人.
認為我們有沒有價值的指標.
好了 到了中年.
我們的職位.
我們的月薪有多少.
我們擁有多少的資產.
駕駛著什麼車.

$^{41}$住在哪一區.
家裡有多大.
甚至是我們兒女的成就.
都會成為我們怎樣看待自己.
別人覺得我們是否成功.
有沒有價值的指標.
好了 到了年紀大了.
雖然退休了.
我們依然在比較.
誰的退休金夠豐厚.
誰的身體保養得更加好.
誰的生活更加有體面.
我們會用這些來衡量自己.
衡量其他人.
他有沒有價值.
我們太習慣這套.
市場經濟的評分標準.
我們確實經歷過.
當你擁有這些東西的時候.
身邊的人不單止會羨慕你.
會稱讚你.
會尊重你.
不過我也想誠實地問大家一個問題.
這種尊重.
這種讚美.
到底是尊重你這個人.
抑或是尊重,讚美.
你背後所擁有的資源呢?.
現代人最深層的焦慮就在這裡.
因為我們知道這個愛.
這個社會的愛是有條件的.
我們每一個人都活在.
都生活在一個很大的壓力底下.
我們要好像一隻倉鼠一樣.
在這個輪子裡面.
不斷地走.
拼命地走.
拼命地證明自己.
我很好.
我有用.

$^{81}$我有價值.
於是我們職場裡面表現得很專業.
在家裡面表現得很堅強.
在社交媒體裡面表現得很幸福.
我們極度地渴望被愛.
被別人肯定.
不過我們又極度地不相信.
這個世界有真愛.
因為我們都很害怕.
只要有一天.
我們停下來.
又或者我失敗了.
我搞禍了.
那些讚美,那些肯定,那些接納.
就會一下子清零.
這就是兩千年之前.
保羅寫信給羅馬教會的時候.
當時一班信徒所面對的真實境況.
當時羅馬教會的信徒在社會上.
被主流文化排擠.
他們被貼上一個標籤.
就是他們不愛國.
他們面臨經濟和政治的重壓.
在教會裡面有兩班信徒.
一班是外邦信徒.
一班是猶太信徒.
他們因為彼此之間的文化,血統.
飲食條例不同.
他們彼此歧視.
他們爭做大佬.
每一個人都想證明自己.
才是上帝面前那個正統.
那個更加純正的子民.
在這種內憂外患.
每一個人都被壓得透不過氣來的時候.
保羅在《羅馬書》第五章.
寫下了這一段宣告.
這一段是驚天動地的宣告.
他要為那群活在重壓下的靈魂.
送上一份永遠不落空的定心丸.

$^{121}$我們一起來看看經文的第一節.
去到第二節.
我們一起讀出這兩節經文.
預備起.
「所以,我們既因信稱義.
就藉著我們的主耶穌基督.
得以與上帝和好.
我們又藉著祂.
因信得以進入現在所站立的恩典中.
並且歡歡喜喜.
盼望上帝的榮耀」.
如果用我們現代人的說話來說.
保羅的意思是.
各位.
你們可以停下來.
停在內心那一台.
拼命要證明自己的跑步機.
這裡出現了兩個極其重要的神學詞彙.
不過.
他們一點都不抽象.
一個叫「稱義」.
一個叫「和好」.
什麼叫「稱義」?.
「稱義」不是指你變成了一個.
道德完美.
從來都不犯錯.
不是.
不是這個意思.
「稱義」指的是在全宇宙.
最高的審判者.
創造你和我的上帝面前.
去看著你.
他就敲下這個法官的定音槌.
去告訴你.
他也告訴所有的人.
你已經沒有罪了.
你是我所接納的.
現代流行心理學.
每天在電視上.
在書本上告訴我們.

$^{161}$我們要自我肯定.
我們要愛自己.
這個沒有錯.
很有安慰.
不過自我肯定.
也有一個致命傷.
也有一個弱點.
這個弱點是什麼呢?.
它很脆弱.
以及不穩定.
因為自我肯定.
是由我自己.
做我自己的審判官.
是我去評論我自己.
今天你工作很順利.
你減肥達標.
你得到別人的稱讚.
你內心的法庭就會宣判.
你很厲害.
你很有價值.
不過當明天你搞砸了一個項目.
搞砸了上頭給你的項目.
又或者搞砸了一個關係.
你面臨中年危機.
或者你年紀大了.
覺得為什麼這麼力不從心.
你內心的法庭.
就會馬上跟你說.
你真的失敗.
不過聖經的恩順清義.
是完全不同的.
順利到你的接納.
不是由你自己的感覺.
由你自己擁有什麼來決定.
也不是由這個社會告訴你.
你在這個社會的表現是怎樣.
去決定你有沒有價值.
那是什麼呢?.
是由耶穌基督的表現來決定的.
是祂所做的救贖計劃.

$^{201}$救贖工作來決定你有沒有價值.
因為主耶穌基督在十字架上面.
替我們還清了代價.
所以上帝接納每一個相信祂的人.
這個意味著.
你在宇宙中最重要的審判者面前.
你已經拿到這張合格證書了.
這張合格證書不會因為你明天失業.
就會過期.
不會因為你今天考試搞砸了.
就沒收了這張合格證書.
也不因為你年老體弱.
喪失了對這個世界的生產力.
就被上帝看低你.
當你拿到宣判這份合格證書.
你和上帝的關係就進入了和好的狀態.
你不需要再戰戰兢兢地在上帝面前.
在人面前去演戲.
基督信仰給現代人最大的解放.
就是容許你放下你的面具.
大方地承認我自己是不完美的.
其實聖經也說了.
沒有一個人完美.
因為你已經被造物主無條件地接納了.
而且你也被他增殖了.
我相信大家聽到這裡的時候.
我都會覺得上帝這份愛太美妙了.
恩上清二所帶來的平安實在太好了.
問一問大家.
知不知道教會這三年的主題是什麼名字?.
如果你忘記了.
打開你的秩序表程序表.
去到報告那一欄.
寫在這裡.
我們看看我們這三年的全年主題是什麼.
如果看到的話要大聲說.
預備起.
「深化生命 令人歸主 榮耀上帝」.
我們期望在三年裡.
希望大家可以深化生命.

$^{241}$可以去令人歸主.
可以去榮耀上帝.
特別我今天跟大家看「令人歸主」.
當你聽到這個主題的時候.
你心裡面第一個是興奮.
抑或是有一份壓力呢?.
尤其我看到很多大型節目都跟我一樣.
比較內向的.
不是很擅長用說話去說服別人.
於是大家看到我每一次講到一定是寫完的.
我們都不是很擅長用說話去說服別人.
所以聽到「令人歸主」的時候.
可能你會覺得.
嘩 教會為什麼給我這麼大壓力.
又給我KPI.
又給我業績壓力.
今天在這裡我想跟大家.
我想代表大家向上帝做一個最誠實的靈魂拷問.
鼓勵大家對上帝誠實一點.
不需要回答出來.
你在心裡面回答上來就可以了.
我們當中有多少人.
雖然每一個禮拜來崇拜.
參與團契.
參與事奉.
不過其實我們自己都已經很久很久.
沒有經歷過第五節.
第五節經文所說.
我們被上帝的愛所囂觀的呢?.
因為當你被上帝的愛所囂觀的時候.
你就會有真正的平安.
當我們面對生活的成本不斷上漲的壓力.
當我們面對職場的惡性競爭.
報紙說的.
又說今年大學畢業的同學.
他們投身職場的時候.
對比起以前.
有65\%的職位是沒有了.
職場的惡性競爭很大.
或者你面對身體有慢性疾病.

$^{281}$或者你面對家庭關係很冰冷.
我們都跟一些未信的人一樣.
我們很害怕.
很憂慮.
如果這個時候.
我們自己的心都好像一口沽乾了的井.
裡面一滴水都沒有.
我們到底要拿些什麼去跟別人分享呢?.
難道我們要變成一個推銷員.
去跟別人推銷一個連我自己都感受不到.
都不相信的產品嗎?.
如果我們自己都沒有被上帝愛過.
沒有經歷過上帝的愛.
我們怎樣跟別人分享信耶穌呢?.
跟自己欺人和自圓其說其實沒有分別.
多謝大家的盛失.
如果你現在心裡都在經歷這種沽乾.
保羅接下來的經文.
要為你解開這個困境.
我們要重新來看.
《羅馬書》第五章的屬靈次序是怎樣.
保羅不是一開口就叫羅馬信徒去傳福音.
去帶人信耶穌.
不是這樣的.
保羅說.
這份從上帝而來的愛.
我們要在我們的現實生活裡面.
特別是要在我們的患難裡面.
能夠驗證得到.
我們來看看第三至第五節.
第三至第五節.
我們一起讀.
預備起.
「不但如此.
就是在患難中也是歡歡喜喜的.
因為知道患難生忍耐.
忍耐生老年.
老年生盼望.
盼望不至於落空.
因為上帝的愛.

$^{321}$以藉著所賜給我們的聖靈.
囂冠在我們心裡」.
請大家注意.
保羅不是給我們一帖廉價的心靈雞湯.
他沒有說信耶穌一帆風順沒有痛苦.
因為當時羅馬教會信徒的患難.
是真的很大.
而「患難」這個字原文的意思就是.
他們在很窄縫的地方.
在很強烈的重壓之下.
當時的信徒.
願意信耶穌.
去做一個信徒.
是一點都不容易.
保羅給了一個非常反主流文化的.
生命提煉的方式.
去告訴我們.
原來患難會生忍耐.
忍耐會生老年.
老年會生盼望.
這個不是叫你咬緊牙關.
死頂死撐.
在心理學上給了我們的知識.
他說當人面對無法改變的打擊的時候.
我們要學會去接納那個實況.
這樣做是會能夠止血.
止蝕.
止住我們內心的虛耗.
我們不再消耗自己裡面的能量.
聖經神學帶給我們走得更加遠.
當我們在具體生活的壓力和挫折當中.
我們不去問為甚麼.
我們是要找一個地方下來.
去讓自己安靜.
向上帝承認.
我們是多麼無力.
多麼枯乾.
你記不記得以利亞.
以利亞她面對很大的壓力的時候.
她和上帝呼吸.

$^{361}$她覺得自己不好.
她甚至向上帝求死.
這個就是忍耐的開始.
我們在重壓之下沒有逃避.
當然要逃離那個殺你的地方.
也沒有向世界妥協.
而是選擇去找上帝.
站在上帝的面前.
我們記得以利亞的故事.
天使來接待她.
讓她在休息和飲食當中得到飽足.
在這種安靜的堅持當中.
上帝要在我們裡面的戰場.
去熬煉出老年靈.
老年圓滿的意思就是.
金屬經過烈火的提煉之後.
出了純淨的品質.
我們會在這個挫折當中突然發現.
過去我們以為最實在最寶貴的學歷.
職位和財產.
其實根本撐不住我們的靈魂.
就在我們發現的這一刻.
我們來到上帝面前尋求平安的那一刻.
我們體會到一種超乎理性.
從神而來的安靜和平安.
我們發現大環境.
雖然仍然很糟糕.
但我們竟然沒有崩潰.
我這樣跟大家說不是在書本內.
是我曾經都經歷過.
當我面對很大的難題.
心裡又害怕又憤怒又擔心的時候.
我不知道怎樣去面對.
簡直是真的.
覺得心裡好像有個洞,逃空了,很虛弱.
不過我就是選擇坐在那裡,安靜下來.
整整三個小時裡面.
我只跟上帝說了兩個字.
不停重複又重複說了兩個字.
「上帝,請你賜我平安」.

$^{401}$「請你將你的平安賜給我們」.
然後經過這三個小時的安靜.
坐下來之後.
慢慢我覺得心裡變得很實.
這口枯乾的井.
就在我們願意去對上帝誠實.
在患難裡面去仰望祂的時刻.
就被聖靈的愛重新挑灌,滿到寫.
今年我們的教會主題是「令人歸主」.
上帝對你第一個的心意.
不是要你勉強去拉兩個人回來教會.
上帝今年對你第一個的心意.
是請你將你自己先帶到上帝的面前.
讓祂的愛先填滿你一口.
你內心裡面很枯乾的井.
當我們願意回到上帝面前的時候.
我們會看得到.
看得到是一幅什麼圖畫呢?.
原來我們所相信.
所要跟人分享的這位上帝.
祂的愛不是我們想像出來的.
是真真實實我們會經歷得到.
我們一起讀第六節第八節吧.
預備起.
我們患軟弱的時候.
基督就在特定的時刻.
為不見情之人死.
為人死是少有的.
為人死或者有感做的.
唯有基督在我們還作罪人的時候.
為我們死.
上帝的愛就在此向我們顯明.
今天的講題是.
看透我.
依然愛我.
人世間最美好的愛.
不論是父母對我們的愛.
還是戀人之間的愛.
在現實當中.
往往帶著一些隱形的條件.

$^{441}$隱形的條件.
而且能力很有限.
人間的愛往往在說.
因為你表現得很好.
當然不會說出口.
因為你表現得很好.
因為你很優秀.
因為你很可愛.
所以我愛你.
這個導致我們在最親密的人面前.
有時候也要小心翼翼.
去藏起自己的缺點.
很怕他看穿我的時候.
愛就消失了.
不過上帝的愛完全是顛覆了人類的邏輯.
保羅用了三個詞語來形容.
我們最真實,最不堪的底牌.
我們還軟弱的時候.
我們還罪人去對抗神的時候.
我們根本無力自救的時候.
這個意味著什麼呢?.
上帝在看透了我們.
看穿了我們生命裡面所有的黑暗.
去看穿了我們那些.
連我們的伴侶.
連我們的父母都不知道的私密罪惡.
去看穿了我們內心那份的自私.
懦弱,虛榮和失敗之後.
他完全看穿了我們的情況底下.
去依然愛我們.
他現在無限不是只有一個「講」字.
他不是叫我們在腦海裡面去想像.
上帝很愛我們.
不是.
保羅給了一個歷史的事實告訴我們.
記錄就是在特定的時刻為我們死.
在公元一世紀的羅馬歷史上清楚記錄.
耶穌基督實實在在被釘在十字架上為我們死.
十字架的意思就是上帝看過我們最差.
最不可愛的一面.

$^{481}$他沒有轉身走開.
而是選擇走上這個刑場.
去代替我們承受了審判.
他付上了自己的生命作為代價.
去救贖你,救贖我.
靚姐妹.
這個才是呼音的震撼的地方.
在一個每一天都跟我們打分.
跟我們比分數.
每天都在看我們表演的這個.
很可怕的世界裡面.
有位造物主對我們說.
「我知道你今天過得一團糟」.
「我知道你心裡面充滿了懷疑和灰光」.
「不過我已經在十字架上證明了我對你的愛」.
「你不用再戴面具的」.
「你過來我這裡」.
這份愛是我們靈魂唯一的永遠不會沉沒的藍.
當我們自己被這份看透我依然愛我的.
客觀歷史事實所震撼的時候.
被聖靈的平安重新去囂觀我們的時候.
帶人信耶穌.
就不再是一件強人所難的KPR任務了.
傳呼音領人歸主的本質是什麼呢?.
它不是用我們絕頂的口才去跟別人辯論.
也不是強拉人去進我們的教會.
傳呼音只不過是一個很可憐的人.
發現了哪裡有麵包.
於是很安靜地.
很真誠地.
去帶另一個同樣可憐的人的手.
帶他一起去吃麵包.
在我們身邊有太多同樣生活在重壓之下的靈魂.
你的同學.
可能他只是為他的前途.
覺得很大的壓力很迷茫.
搞到整晚都睡不著.
有可能是你的同事.
可能他正面對中年危機.
面對著家庭的破碎.

$^{521}$面對著在辦公室裡面要辦得很堅強.
不過他內心卻是很荒涼.
可能那位是你的鄰居.
他正默默地面對身體的衰退.
和死亡的恐懼.
覺得自己再沒有任何的價值.
這些人都在這套殘酷的社會標準底下.
生活得很疲累.
很累.
累得只剩下半條人命.
他們都在尋找一個不需要再裝扮得完美的避風港.
在講到最後我想邀請大家.
無論你是多少歲都好.
我們一起安靜下來.
閉上我們的眼睛.
請你在心裡誠實地回答我以下簡單的三個問題.
如果今天拿走你的學歷.
拿走你的職位.
拿走你的財產.
你還知不知道自己是誰呢?.
請你誠實地去面對上帝.
你現在的心是一口很枯乾的井.
抑或是充滿了上帝的愛呢?.
在你的身邊.
有哪一個人正是活在重壓之下.
你記不記得他的名字呢?.
又或者有誰的名字只是浮現在你的腦海當中呢?.
你願不願意在今年.
陪他一起回到這位看透我們.
卻依然愛我們的上主面前呢?.
讓我們同心祈禱.
在我們還在做罪人去抗拒你.
去叛逆你.
還是滿身破碎的時候.
你就為我們釘死在十字架上.
求你今天並且以後的每一天.
親自去安慰我們的心.
用你那個源源不絕的愛.
去囂冠在我們中間每一個枯乾的心靈.
讓我們帶著這份不會被搖動的安全感.

$^{561}$在今年.
將我們自己.
也將我們身邊那些疲累不堪的靈魂.
可以一步一步地帶到你的身邊.
在你的懷抱裡面享受真正的安息.
我們這樣祈禱過.
藉主耶穌基督的名而求.
阿們.
\newpage



\section{羅馬書 7:13-25}
\label{sec:032291}
\textbf{別害怕,祂知你難受!}~~~~周曉暉牧師
\newline
\newline
連結: \href{http://www.church.com.hk/acms/content.asp?site=cdc\&op=show\&type=product\&code=032291\&layout=sermon}{\texttt{032291}} ~~~~ 語音日期: 2025-11-15
\newline
\newline
\hyperref[sec:032290]{\small{< < < PREV SERMON < < <}}
~
\hyperref[sec:index_chronic]{\small{[返順時目]}}
~
\hyperref[sec:032292]{\small{> > > NEXT SERMON > > >}}
\newline
\newline
羅馬書 7:13-25
\newline
\begin{longtable}{cl}
\hline
\hline
章節 & 經文 (和合本修訂版)\\
\hline
7:13 & \begin{tabularx}{0.7\textwidth}{X} 那麼，那善的是叫我死嗎？絕對不是！叫我死的是罪。罪藉著那善的叫我死，為要顯出這真是罪，以致罪藉著誡命更顯出是惡極了。 \end{tabularx} \\ \\ \relax
7:14 & \begin{tabularx}{0.7\textwidth}{X} 我們原知道律法是屬靈的，我卻是屬肉體的，是已經賣給罪了。 \end{tabularx} \\ \\ \relax
7:15 & \begin{tabularx}{0.7\textwidth}{X} 因為我所做的，我自己不明白。我所願意的，我並不做；我所恨惡的，我反而去做。 \end{tabularx} \\ \\ \relax
7:16 & \begin{tabularx}{0.7\textwidth}{X} 如果我所做的是我所不願意的，我得承認律法是善的。 \end{tabularx} \\ \\ \relax
7:17 & \begin{tabularx}{0.7\textwidth}{X} 事實上，這不是我做的，而是住在我裡面的罪做的。 \end{tabularx} \\ \\ \relax
7:18 & \begin{tabularx}{0.7\textwidth}{X} 我也知道，住在我裡面的，就是我肉體之中，沒有善。因為立志為善由得我，只是行出來由不得我。 \end{tabularx} \\ \\ \relax
7:19 & \begin{tabularx}{0.7\textwidth}{X} 我所願意的善，我不去做；我所不願意的惡，我反而去做。 \end{tabularx} \\ \\ \relax
7:20 & \begin{tabularx}{0.7\textwidth}{X} 如果我去做我不願意做的，就不是我做的，而是住在我裡面的罪做的。 \end{tabularx} \\ \\ \relax
7:21 & \begin{tabularx}{0.7\textwidth}{X} 我覺得有個律，就是我願意行善的時候，就有惡纏著我。 \end{tabularx} \\ \\ \relax
7:22 & \begin{tabularx}{0.7\textwidth}{X} 因為，按著我裡面的人，我喜歡神的律， \end{tabularx} \\ \\ \relax
7:23 & \begin{tabularx}{0.7\textwidth}{X} 但我看出肢體中另有個律和我內心的律交戰，把我擄去，使我附從那肢體中罪的律。 \end{tabularx} \\ \\ \relax
7:24 & \begin{tabularx}{0.7\textwidth}{X} 我真苦啊！誰能救我脫離這必死的身體呢？ \end{tabularx} \\ \\ \relax
7:25 & \begin{tabularx}{0.7\textwidth}{X} 感謝神，靠著我們的主耶穌基督就能！這樣看來，一方面，我內心順服神的律，另一方面，肉體卻順服罪的律了。 \end{tabularx} \\ \\
[1ex]
\hline
\hline
\end{longtable}
$^{1}$這錄音或錄影的版權是屬於基督教宣導會不各堂所有.
我們的屬靈傳統之一是以真理教導為本.
歡迎您將此錄音或錄影的連結傳給親友.
懇請注意,未經授權,請勿轉錄或剪輯本錄音或錄影.
甚願好似使徒行傳二十章三十二節下所講.
這道能建立你們,叫你們和一切誠誠的人同得基業.
明明應該這樣做,卻在常常變成相反的事.
然後無心睡眠,腦交戰.
明明心裡想做對的事,但實際做出來的是相反.
回過頭來很懊悔,很無奈.
這種明明和常常拉扯的感覺,我想幾乎是所有人都經驗過.
明明你下定決心一定要早睡.
結果呢,就常常凌晨兩點還在滾手機.
我不知道你今早凌晨是不是這樣.
明明告訴自己,這一次一定要心平氣和地和對方溝通.
但常常一出口,就變成最尖銳的指責,急向對方.
明明心裡很清楚什麼是對的.
但身體裡好像有另一個意志.
常常就把你拉去相反的方向.
你以為是關乎意志力薄弱的事情嗎?.
今日經文告訴我們,其實是觸碰到人的一個內心非常特別的狀態.
而聖經的羅馬書裡有一段將這個內心的衝突描寫得最淋漓盡致.
就是今日我們讀的第七章後半段,十三節到二十五節.
大概剛才你聽主席去讀,或者你以前看過羅馬書裡面幾節七章的經文.
你可能都會很心痛感受,這個是我,是我裡面那種沮喪的狀態.
今天我們會用幾個角度去切入,看看這一場的腦交戰.
去討論一下今日我們有的,不明白的地方.
並且到底真正的出路在哪裡.
我們進入經文之前,一起祈禱吧.
主啊,你知道我們裡面的難受,害怕,擔心,羞恥.
主啊,幫助我們,無論我們已經信了你,或者我們仍然在尋求你,認識你.
都讓我們得著你聖靈的指引,明白羅馬書七章下半段的經文.
幫助我們每一個,多謝你,我們祈禱在奉耶穌基督得勝靈之祈求,阿們.
保羅一開始在第七章先講的是律法.
他說律法本身是神聖的,是公義的,是良善的.
這個律法是神給的,應該是這樣的狀態.
但是為什麼這個律法常常都跟死拉上關係呢?.
或者這樣說,上帝給的好的東西,律法是他頒布的.
為什麼會帶來好像壞的結果?.
這個就是第十三節,保羅開始要處理的問題.

$^{41}$他直接這樣說,我請大家讀這四行.
保羅一開始就在表達,好像在做一些事情.
令到律法看來是跟死拉上關係,是律法導致人有這個問題.
導致我們會死,沒有律法可能好一點.
不過保羅立刻很堅決地否定,不是,絕對不是.
原本文字是很強烈很強烈的語氣.
保羅這樣說,問題的根源不是律法本身,而是罪.
這個罪很狡猾,但是律法本身是好的.
不過罪卻是他利用了律法,去透過律法,去使用了律法.
扭曲了,變成產生更多更多的東西出來.
以至是罪大惡極必死,必須要死.
這個就是罪的狀態.
他藉著良善的律法,在前面第七章,早些的經文他這樣說.
譬如說不可以貪心,這個是誡命來的,是神頒布的律法.
但是當罪去使用他,就變成一個他生產的基地.
產生了每個人都有的死亡.
在這裡說,罪捉住了機會,第二行,最後那幾個字.
或者在這裡可以翻譯為,罪成為了一個根據地.
在那個地方去不斷產生更多,這個本來說不好做的事情出來.
結果呢,就產生了死亡.
也都讓人和神隔絕.
所以其實不是律法令到人死.
而是罪利用了律法,去達成讓人必死無疑的目的.
其實小時候你都明白的,我們所身處的城市有很多標語.
我不知道你有沒有去過維多利亞公園,看到那個牌.
不准採摘花朵,不准踐踏草地,不准在這裡踩滑板車.
不准,不准,不准,不准,很多的.
有時候我們小時候就覺得,其實我都沒有想過在這裡.
真的要做這件事,踩草,你說不准,沒有人,踩兩下先.
好像你本來不為意,但是因為提了出來.
反而罪就在你裡面,將它成為一個生產的基地.
產生這些不好的行為.
又譬如,我那個年代的中學生,常常都初中就有個辯題.
中學生不應談戀愛,你那個時候有沒有?.
本來沒有想的,拿出來討論一下,可能談戀愛都不錯.
是啊,這個就是我們人性的一個弱點.
而且,這種狀態令我們進入一個更深的無力感.
搞清楚這個罪和律法的關係之後.
我們就可以看一下,本來接著在第十四節要我們想的東西.
它那裡說,事實上我們知道律法是屬靈的,即是好的.

$^{81}$但是我呢,這個我呢,是屬肉體的.
這個肉體在保羅的書信裡面常常都用到的字詞.
不是說我們這個身體的物質這麼簡單.
更多的時候是表達我們墮落之後.
我們人常常都自我中心.
傾向做不好的事,很軟弱,那種人性的狀態.
用肉體這個字去表達出來.
而且這一節更加說,我是屬於這樣的人性狀態.
並且已經賣給罪了.
好像一個奴隸一樣,失去了自由,被罪這個主人牢牢控制著.
從這一點開始,或者從第十四節開始.
如果從原本聖經的文字.
就會突然發現保羅從一堆過去式討論法律.
討論耶穌基督彌帶給我們恩訊稱義的東西.
來到這一段特別的,突然轉了.
轉了現在式,保羅用第一人稱.
我現在非常非常非常痛苦.
我現在內心很掙扎去說一下的話.
我卻是屬肉體,我已經賣給罪作奴僕了.
然後就是一連串,剛才你聽經文,很經典的掙扎.
我所願意的,我不去做.
我所憎的,我反而去做.
明明好的東西我是願意的.
但常常做出來的是相反,是錯的.
這個描述的痛苦程度.
一路去表達,其實這個人很矛盾.
並且在二十二節他這樣說.
他內心是很想做好的.
明明是很喜愛神的律法.
但是二十三節說.
我請大家讀下面那三行.
但我發覺.
將人那種人性.
並且他內心應該明明向好.
但是常常做出來壞的東西.
他內心的掙扎表達到極點.
他現在感覺到他內心有兩個律.
有兩股力量將他拉扯.
而且有一股力量要他服從罪的律.
這一場的爭戰到了極點,很疲倦.

$^{121}$所以他最後在第二十四節.
他發出一個呼喊.
我,我這個人真是苦啊.
我不知道你有沒有試過.
人生掙扎到一個位置.
你覺得我控制不了.
我搞不定,我很苦.
於是他問.
誰能救我脫離這個必死的身體呢?.
你今天第七章單獨抽出來看.
很簡單直接.
但如果從羅馬書整個脈絡.
由第一章講到去.
我們前兩星期第四章第五章.
上星期第六章.
我們本來已經讀到.
因信稱義然後信徒已經怎樣.
是向罪死了,得了自由.
並且我們的金句是這樣說.
你們也要看自己對罪是怎樣.
死的.
在耶穌基督裡面我們對神是活的.
所以不要再給罪在我們必死的身上.
掌權使我們順從身體的肆慾.
得救了.
為什麼第七章又講自己這麼厲害.
這麼痛苦的掙扎呢?.
這個位置就帶來了在過去數百年.
在羅馬書這一章.
特別是第七章十四節.
開始這一大段.
《釋經學》不斷在爭論.
到底這裡說的「我」是誰?.
或者是哪一種人?.
他講到那個人好像是未信主.
因為他仍然在律法之下.
發現自己無能為力.
並且有那種矛盾.
到一個位置很絕望.
「我真是苦」.

$^{161}$另一邊就覺得.
這裡應該是在說一個重新得救的基督徒.
在這個成聖路上.
和他內裡的罪性去搏鬥.
我會和你從這兩方面都一起思想.
因為可能這兩方面.
你過去在讀這段經文的時候.
或者聽以前不同的人去分享講道.
都可能採取了其中一邊.
第一邊.
這個「我」描述的是一個未信主的人.
他在律法面前掙扎.
無助的狀態.
即使他做得多好.
他都去不到神的標準.
所以很有道理.
這也是我們認識福音時.
最基本的就是.
我們向善.
我們做得好.
行得正 站得正.
但是我們都去不到神那裡.
所以他說「我是屬肉體的 賣給罪」.
在這段痛苦的對白裡.
你有沒有留意到.
剛才主席讀的時候.
有沒有提過聖靈?.
應該沒有吧?.
沒有吧?.
因為第七章這裡沒有這樣表達.
好像在鋪排第八章.
聖靈的幫助.
好像在一個黎明之前.
最黑暗的時候.
然後迎來曙光.
聖靈來幫助你.
單靠自己面對律法.
人就徹底失敗了.
因為到了第八章.
那個對比很大.

$^{201}$當這個我去面對到.
沒有辦法 完全無助.
第八章就是歡慶了.
第七章最後預備鋪排到第八章.
「即在我們的主耶穌基督就能救了」.
第八章更加說.
「因此如今在基督耶穌裡的人就不被定罪」.
將那個人唯一得救的路說了出來.
只有耶穌基督.
你只能夠靠耶穌基督.
然後耶穌基督施下聖靈.
讓你可以從罪中得釋放.
這種的絕望.
到了很有盼望.
很有說服力.
我相信在這個位置.
就算這一段有些思考.
在我們在座當中.
大部份我相信都已經順了.
不過我亦都認為.
通常都有一些朋友.
你是回來聽一下.
你認識一下耶穌.
或者你來坐一下.
很想認識這個真理.
認識這個信仰.
甚至有些人聽了很久.
你仍然覺得我不是很需要跟耶穌.
因為我都覺得自己是一個挺好的人.
我是一個很好的人.
行得正 站得正.
我不知道你身邊的弟兄姊妹有沒有這樣的朋友.
他真是一個大好人.
比很多基督徒還要好.
這個大好人他喜歡一些對的東西.
甚至他聽到耶穌基督的道理.
聖經講的說話.
他都覺得很正面.
是對的 這樣做.
不過羅馬書第七章這裡.

$^{241}$是對這些覺得自己是大好人的人.
發出一個提問.
無論你行得多正 站得多正.
我想你總是會有遺憾的.
而且那些遺憾.
很可能都是因為你做錯了選擇.
做錯了.
你本身不是這樣的.
但偏偏做了相反的事.
不是你有意.
不過你卻害了別人.
或者苦了自己.
每一個大好人都有這些片段.
如果是這樣的話.
這段經文是繼續推你去想.
你誠實面對自己的時候.
你知道就算我做得多好.
總有些地方我都沒有辦法.
中國有一句古語叫做.
「聖人都有錯」.
「知聖賢的人都知道自己會有錯」.
那誰可以救他呢?.
如果放在羅馬書第七章.
就是對這批聖賢人發出提問.
推你去想一想.
唯一出路在哪裡呢?.
就是福音.
因為聖賢人都要想.
他都會遇到.
「我所願意行的事我沒有行」.
「我所不願意作的惡我反而去作」.
真的有些時候.
你無法靠你那種做人的原則.
那些意志那些良心高尚的品格.
繼續做好的.
反而有一些.
你都不想再提.
不想再想.
因為你知道那些是不好的東西.
很怕提起.

$^{281}$你情願忘記.
你害怕面對.
甚至你想.
總不知道我帶入棺材算了.
不會再提了.
但你知道你忘記不了.
羅馬書的真理讓你知道.
是的.
這是你人性的狀態.
你無法幫自己.
如果是這樣的話.
怎麼辦呢?.
今日經文說.
不要怕.
神知道我們難受.
神明白你這個大好人的難受.
所以各位朋友.
不要等.
再聽下去都是好的道理.
要接受這個福音.
要信耶穌.
完了崇拜.
找同工幫你決志.
聽到這裡.
你開始會不會想.
如果這真是講世人.
講那些未信者的經驗.
跟我們基督徒日常的掙扎很有脫節.
難道信了主之後.
這種內心的膠戰就沒有了嗎?.
一下子就洗白了.
以後都做得好了嗎?.
顯然這不是我們大部份人的經驗.
你的經驗可能就是.
信主的時候很認真.
去上洗禮班都很認真.
洗完禮之後.
第三天.
你又死了.
開始犯罪了.

$^{321}$你那時候可能有一個疑問.
我不是一個蒙恩基督教的人.
為什麼我會這樣呢?.
如果這一段不是對著這班基督徒講的話.
那會不會就是我根本就不合格?.
怎麼辦呢?.
上帝會不會不要我?.
這可能是你曾經有的疑問.
甚至你不敢問的問題.
不過弟兄姊妹不要慌.
不要亂.
不要怕.
剛才講了第二個觀點.
就是基督徒.
誠信路上的掙扎.
關鍵在第22節.
我請大家讀.
因為喜愛神的律.
你想想一個未信主的人.
同神為敵的人.
會發展內心很喜愛神的律法嗎?.
這一派的解經家就以此為一個根據.
一個未重生的人.
他會怕刑罰.
或者為了社會認同的標準去遵守律法.
但是他內心深處不是真正喜愛這些標準.
留意這個喜愛.
不只是理智上的同意.
不只是認同.
是那種裡面的喜悅和渴慕.
如果這樣的話.
聽起來.
不應該更加像詩篇所講的詩人嗎?.
喜愛神的話語就夜思想.
喜愛神的話語好像蜜糖那麼甜.
真金那麼棒.
那種狀態嗎?.
如果這樣的話.
這一段不像是說一個想做壞事的我.
和律法之間的爭戰.

$^{361}$更加像一個想討神喜悅的我.
和我盤踞在我身體裡面.
那種很頑固的罪.
不斷的爭鬥.
是呀.
保羅在第15節.
又回到那一節.
我們又看到.
他就是這樣說.
因為我所做的我並不知道.
我所願意的我不去做.
我所憎惡的我反而去做.
甚至他在第17節繼續說.
既然是這樣.
這事就不再是我做的.
事主就是我裡面的罪做的.
這個區分實在太重要了.
為什麼這樣說呢?.
他不是說不關他的事.
不是我的意思.
而是在表達.
根本跟他的意願和他本身.
成為一個得救的基督徒.
完全相反的.
但是他就不知道為什麼.
就是這樣.
常常做了神喜悅的事.
他將自己認同那個願意行善的自己.
將那個作惡的規律.
成為一個入侵的.
住在他裡面的.
不應該在這裡的東西.
然後形成爭戰.
丙姐妹,我想這份的痛苦掙扎.
你很明白的.
你信得越久.
你越明白.
我信了50年了.
我還在犯罪.
那怎麼辦呢?.

$^{401}$還有沒有得救呢?.
我做了導師.
我做了執事.
我做了牧師.
還要犯罪.
那就很差了.
沒有盼望了.
這個是會令我們害怕的地方.
那個痛苦就是.
得救了.
但是還沒有完全得救.
第23節將這個戰場描寫得更清楚.
這個冥冥和常常的戰場是怎麼樣的?.
就是心中的律和肢體裡面的罪的律爭鬥.
23節,我請大家讀.
「但我快樂,自己的肢體中仍有一個律.
但我愛神的律,我愛神,我愛我的律.
是我愛神的律,我愛神,我愛我的律」.
這個我發現自己裡面有兩個律在打仗.
心中的律和肢體中罪的律.
結果他被罪的律抓了他.
擄了他.
拉著走.
不能自拔.
這個律.
其實新漢業館翻譯得很好.
不是理解為波西律法的那些條文.
這裡是說一種運行的原則.
定律.
信主之後我們裡面有一個新的屬神的運行原則和定律.
那個就是我們內心的律.
同時那個舊有的原屬於罪的原則.
那種運行的方式,那種定律.
沒有完全消失.
所以在我們的生命當中.
在我們信了到將來永遠得救的這個過程裡面.
我們就常常在這裡爭戰.
而且這裡說好像罪打贏了.
保祿說將我擄去,成為俘虜.
這個是很多信徒真實的感受.

$^{441}$但是保祿在講完這一節之後就發出.
我好苦啊.
怎麼辦啊.
誰誰可以來救我啊.
這個苦表達的一個極度的悲慘.
這個必死的身體就是被罪所捆綁.
充滿了內在的衝突.
也都令我們情緒低落.
的確我剛才問你信了多少年.
或者這樣說.
你有多久沒有犯過罪.
可能你剛才主席帶你去認罪祈禱.
你剛剛在認剛才在進來之前犯的罪.
我們都是在這麼慘的狀態裡面.
好苦啊.
但是經文在下一秒.
馬上給你一個曙光.
因為上帝知道你難受.
祂馬上給你一個答案.
要你不要害怕.
一起讀吧.
你認真一點讀.
由問誰可以.
到耶穌基督能救.
中間沒有其他的了.
馬上告訴你這個答案.
這個是我們作為基督徒仍然要捉住的真理.
所以有了這個答案.
我們覺得一天光了.
是的,我們會有掙扎.
但是耶穌繼續幫我們.
不過你的困難就是.
是的,耶穌幫了我.
但是很快又失敗了.
二十五節都是這樣表達的.
你看清楚一點.
下面那三行.
一方面我是服事神的律.
另一方面肉體卻服事罪的律.
你看到嗎?.

$^{481}$好像本來應該去到二十五節上半那兩句就完了.
大禹團圓結局.
但是他下半節馬上告訴你.
不是的.
這個情況是兜兜轉轉.
兜兜轉轉.
弟兄姊妹.
這個得救大微得哂的狀態.
是我們的生命.
是的,打贏了耶穌基督的十字架.
我們屬於耶穌.
不過戰鬥還在.
就好像本來罪是作王的.
上一兩章跟你說作王.
但是耶穌基督將罪拉下來.
現在上帝透過耶穌基督作王.
那罪去哪裡呢?.
繼續躲藏在你的身體裡.
成為了恐怖份子.
找機會出來偷襲你.
影響你.
他沒有散去.
我們這一生就要在這個爭戰裡面.
不斷拉扯.
在這個張力當中.
來到這裡我們要問.
到底是哪一個觀點?.
這邊是人的.
還是未信的世人,還是信了基督徒?.
因為這個觀點.
很多大師級的解經家都是這樣說的.
甚至連我們本地有一位.
專門研究保羅,研究羅馬書的大臣學家.
他寫了很厚很厚的羅馬書著釋.
他對這一章的看法.
都是在說世人.
未信主的狀態.
但同時第二個觀點.
亦都有近年.
很著名的港道家常常都有的論點.

$^{521}$甚至我們說改革的時候.
宗教改革的時候.
馬丁路德加爾文.
都認為是已經得救.
認為是已經得救的.
重新得救基督徒的狀態.
那誰是對的?.
我認為這一章其實就是一顆寶石,鑽石.
很多的切割面.
你在不同的位置.
看到不同的折射.
不同的光彩.
沒有哪一個位置.
折射出來的光不是假的,不是真的.
每一個點看都是不同的.
而每一個方位所看到的.
通通都是鑽石折射出來的光彩.
都是真的.
或者這段經文的章理.
本身就是他想給世人的寶藏.
是要逼每一個讀這段經文的人.
你未信的,你面對自己.
你說你自己是好人.
實際你自己知道.
你沒辦法幫自己.
對於基督徒來說.
你也是要面對的.
你沒辦法靠你努力返教會.
努力去做好行為.
去改變你自己太多.
因為我們的確仍然是和罪增賤.
在這個矛盾裡面得救了.
但還未得哂.
頂尖這個也是真實的.
這個第七章是一粒寶石.
讓我們每一個人都可以看到.
福音的光彩.
是有效的對每一個人.
所以我認為兩邊都是對的.
不過這件事很灰.

$^{561}$怎麼辦?.
經常失敗怎麼辦?.
冥冥常常.
這個掙扎.
這個時候.
我想引用騰金輝牧師.
他曾經說過一個比喻.
騰牧師曾經提出一個.
對基督徒常常失敗.
他去解釋這件是怎樣的事.
他說基督徒的一生.
就像在一個靈性的學院裡面.
持續地學習.
在一個獨特的學院裡面.
失敗不是終點.
不是遊戲.
而是一個重要的學習過程.
你每天的失敗.
你每個星期的失敗.
你常常的認罪.
應該成為你一個學習的機會.
這是上帝給我們一個屬靈的考試.
並且當我們在某方面跌倒之後.
騰牧師說.
神不是立即宣判你.
踢你出校.
而是讓你再重新補考.
不是補考一次.
是不斷可以補考.
不過我們一想到補考.
就不太想了.
騰牧師繼續說.
這個補考的分數.
是不會因著你補考多少次.
打折扣的.
反而你越補考越加分.
為什麼?.
因為在經歷補考的過程裡.
你已經不斷付出悔改.
反思和掙扎.

$^{601}$令到你的人.
你的機會和生命更加成熟.
更加真.
當你每一次在軟弱當中重新站立.
這個更能夠塑造我們的靈性.
彰顯神真實的恩典.
騰牧師更加說.
聖經裡面有很多這些補考的人.
但是這兩個你都明白.
大衛.
為什麼大衛寫這麼多詩篇?.
經常要補考.
還有新約的彼得.
由他做門徒第一年.
到他做門徒的第三年.
都經常要補考.
不過他們根實耶穌基督.
擁抱這個真神的話語.
我們可以在補考裡加分成長.
騰牧師說.
在補考的過程裡.
我們的責任是什麼?.
我們在那個位置要通過補考.
就來自四個步驟.
知罪.
悔罪.
認罪.
離罪.
知道自己做錯了.
然後那個是罪的律.
很痛苦啊.
悔改.
認了.
是啊.
在我裡面的.
我沒有辦法了.
然後要離開.
頂智慧如果你信了很久.
你依然在同一種的罪裡不斷補考.
大概你沒有知罪.

$^{641}$認罪.
悔罪.
離罪.
這樣是考不到的.
我們依然要配合上帝給的恩典.
採取實際的行動.
我們說是不是又要回去呢?.
要靠我自己的舊腦上?.
不是啊.
而是你要明白.
採取實際的行動.
是讓你領受神的恩典並且運用出來.
可能我曾經過去都講過一個很小的片段.
最初我爸爸媽媽.
帶我去完台灣宣教.
回來香港.
拆班.
在台灣的時候我是沒有英文學的.
一拆班就拆五年班.
第一次的墨書我是零分.
第二次墨書是零分.
第三次老師給我機會.
他還給我那張Unseen Dictation的紙.
叫我拿回去讀.
我有23分.
現在上帝給你加分的機會.
你也要拿去讀.
有恩典.
但是我們要接納並且運用.
這個是成長.
弟兄姊妹.
上帝為我們預備一個又一個報效的機會.
是我們承認失敗.
但是不會停留在失敗裡面.
我們會認罪悔改.
我們跟著採取具體離開罪的行動.
並且在這個過程裡面.
一生在我們人性依然要爭鬥的軟弱裡面.
依靠神的恩典.
所以弟兄姊妹不要怕.

$^{681}$今天你依然覺得自己很失敗.
上帝知道你難受.
在尋找人生出路的朋友.
不要怕.
上帝就是給了福音.
成為你的出路.
羅馬書七章十三到二十五節這段經文.
不是給我們為自己的失敗找藉口.
反而是提醒你正確認識自己的掙扎.
如果上上在保羅描述有那種內心的交戰.
腦交戰.
你要知道這不是在說你信仰有問題.
反而更加可能對基督徒來說.
是在證明你的新生命不斷奮力成長.
突破.
努力報效中.
第二.
要培育對耶穌基督的倚靠.
要保持謙卑.
當你上上向神喊出我真的苦.
你就是知道裡面的罪的頑強.
你沒辦法靠自己.
搞得定.
謙卑下來.
倚靠基督.
下一章就說聖靈的幫助.
第三.
這段經文激勵你持續同罪爭戰.
不要怕.
主戰你難受.
不要怕.
雖然是苦.
苦不是結局.
結局永遠都是那一句.
靠著我們主耶穌基督.
這個最終的勝利.
是屬於每一個願意.
跟從耶穌基督的人.
我們一起祈禱.
主啊.

$^{721}$你知道我們裡面.
想起我們自己很多的狀態.
甚至我們有時候覺得羞愧是.
我們好像擺脫不了.
一些常常令我們跌倒的事情.
神啊你知道我們的軟弱.
我們只能夠倚靠你.
我們要領受你的恩典.
求主幫助我們.
不單止知道悔罪.
我們都願意認罪和離開罪.
給我們有這個力量.
求聖靈幫助我們.
將耶穌基督的福音.
活現在我們的生命當中.
走每一天的道路.
多謝你.
全數基督德盛名之祈求.
阿們.
\newpage



\section{創世記 18:1-15}
\label{sec:032292}
\textbf{在主前蒙恩}~~~~高國雄牧師
\newline
\newline
連結: \href{http://www.church.com.hk/acms/content.asp?site=cdc\&op=show\&type=product\&code=032292\&layout=sermon}{\texttt{032292}} ~~~~ 語音日期: 2026-06-14
\newline
\newline
\hyperref[sec:032291]{\small{< < < PREV SERMON < < <}}
~
\hyperref[sec:index_chronic]{\small{[返順時目]}}
~
\hyperref[sec:032293]{\small{> > > NEXT SERMON > > >}}
\newline
\newline
創世記 18:1-15
\newline
\begin{longtable}{cl}
\hline
\hline
章節 & 經文 (和合本修訂版)\\
\hline
18:1 & \begin{tabularx}{0.7\textwidth}{X} 耶和華在幔利橡樹那裡向亞伯拉罕顯現。天正熱的時候，亞伯拉罕坐在帳棚門口。 \end{tabularx} \\ \\ \relax
18:2 & \begin{tabularx}{0.7\textwidth}{X} 他舉目觀看，看哪，有三個人站在他附近。他一看見，就從帳棚門口跑去迎接他們，俯伏在地， \end{tabularx} \\ \\ \relax
18:3 & \begin{tabularx}{0.7\textwidth}{X} 說：「我主，我若在你眼前蒙恩，請不要離開你的僕人走過去。 \end{tabularx} \\ \\ \relax
18:4 & \begin{tabularx}{0.7\textwidth}{X} 容我拿點水來，請你們洗腳，在樹下休息。 \end{tabularx} \\ \\ \relax
18:5 & \begin{tabularx}{0.7\textwidth}{X} 既然你們來到僕人這裡了，我再拿點餅來，讓你們恢復心力，然後再走。」他們說：「就照你說的去做吧。」 \end{tabularx} \\ \\ \relax
18:6 & \begin{tabularx}{0.7\textwidth}{X} 亞伯拉罕急忙進帳棚到撒拉那裡，說：「你趕快拿三細亞細麵，揉麵做餅。」 \end{tabularx} \\ \\ \relax
18:7 & \begin{tabularx}{0.7\textwidth}{X} 亞伯拉罕又跑到牛群裡，牽了一頭又嫩又好的牛犢來，交給僕人，僕人就急忙去預備。 \end{tabularx} \\ \\ \relax
18:8 & \begin{tabularx}{0.7\textwidth}{X} 亞伯拉罕取了乳酪和奶，以及預備好了的牛犢來，擺在他們面前，自己在樹下站在旁邊，他們就吃了。 \end{tabularx} \\ \\ \relax
18:9 & \begin{tabularx}{0.7\textwidth}{X} 他們對亞伯拉罕說：「你妻子撒拉在哪裡？」他說：「看哪，在帳棚裡。」 \end{tabularx} \\ \\ \relax
18:10 & \begin{tabularx}{0.7\textwidth}{X} 有一位說：「明年這時候 ，我一定會回到你這裡。看哪，你的妻子撒拉會生一個兒子。」撒拉在那人後面的帳棚門口也聽見了。 \end{tabularx} \\ \\ \relax
18:11 & \begin{tabularx}{0.7\textwidth}{X} 亞伯拉罕和撒拉都年紀老邁，撒拉的月經已停了。 \end{tabularx} \\ \\ \relax
18:12 & \begin{tabularx}{0.7\textwidth}{X} 撒拉心裡竊笑，說：「我已衰老，我的主也老了，怎能有這喜事呢？」 \end{tabularx} \\ \\ \relax
18:13 & \begin{tabularx}{0.7\textwidth}{X} 耶和華對亞伯拉罕說：「撒拉為甚麼竊笑，說：『我已年老，果真能生育嗎？』 \end{tabularx} \\ \\ \relax
18:14 & \begin{tabularx}{0.7\textwidth}{X} 耶和華豈有難成的事嗎？到了所定的時候，我必回到你這裡。明年這時候，撒拉會生一個兒子。」 \end{tabularx} \\ \\ \relax
18:15 & \begin{tabularx}{0.7\textwidth}{X} 撒拉因為害怕，就不承認，說：「我沒有笑。」那人說：「不，你的確笑了。」 \end{tabularx} \\ \\
[1ex]
\hline
\hline
\end{longtable}
$^{1}$頂住會平安.
讓我們先做一次同心祈禱.
天父上帝我們獻上感恩.
是你的恩典讓我們成為蒙恩的人.
是你的恩典讓我們今天.
在這個環境狀態當中.
能夠在你的殿裡面.
安靜,舒服地去享受你所賜的一切.
並且專心地將榮耀頌讚歸給你.
求你恩代我們眾人.
開我們耳朵讓我們細心去聆聽.
開我們眼睛讓我們能夠沾上你的榮光.
並且以我們的手腳去回應.
你的吩咐,你的大愛.
由於僕人口中的言語.
頂住會正是心思意念.
蒙聖靈的啟迪蒙主的月立.
祈禱奉耶穌基督的名靠.
阿們.
頂住會今天和大家分享一個題目.
叫做蒙恩.
你是不是一個蒙恩的人.
又或者這麼說.
蒙恩這個字在你的生命裡面.
是怎樣去展現的.
很多人說基督教講很多愛.
同樣地在教會裡面.
教會群體裡面都常常講蒙恩這件事.
回想一下.
過去這段日子你有多少講到.
我真是蒙恩.
又或者和旁邊的人講.
真是好了 你真是蒙恩.
有沒有.
過去這段日子有沒有這樣去分享.
或者有沒有這樣去思考.
蒙恩這個字其實我們都很熟悉.
常常都掛在嘴邊.
我們見到人生活得很順利.
家庭很幸福就會說你真是蒙恩.

$^{41}$很好 很順利.
當我自己可能是不舒服 病了.
好了又或者升職 加薪等等.
很順利 很滿足.
我蒙恩.
但在聖經裡面所講的.
在主前蒙恩.
是不是正正在講我們.
很順利 順風順水的日子.
當我們面對生命裡面的荒涼.
面對身體的軟弱.
甚至罪惡的捆綁的時候.
我們是不是還是一個蒙恩的人呢.
今日的經課.
四段經文自然不同.
很有趣地.
上帝的靈在當中就將它串連.
成為了讓我們去思考.
從上帝的應許到實踐.
一起去思考.
我們作為一個在主前蒙恩的人.
應該擁有怎樣的生命的確具呢.
在詩篇裡面.
116章.
我們聽到一個很特別的訊息.
這個描述就是.
原來在主前蒙恩.
第一樣就是意味著.
我們的聲音被上帝聽見.
因為詩篇說.
我愛耶和華.
因為祂聽了我的聲音和我的懇求.
祂既向我質疑.
我一生要求告祂.
詩人在這裡.
用一個很親密的.
關係的方法去表達.
他和上帝的狀況.
上帝是靠在他身邊.
去到耳邊和他聊天.

$^{81}$當你看完這段聖經.
這段經文之後.
你就知道當時究竟.
詩人面對什麼狀況 什麼遭遇呢.
在聖經裡面後面一路看你就會發覺.
他講到在死亡的困所裡.
遭遇痛苦.
愁悶 在最絕望的時候.
他就只有呼求上帝.
原來這段經文裡面.
告訴我們.
這個蒙恩的開始.
是由上帝的質疑開始.
中文字很特別.
用質疑這個描述.
就好像我們這麼偉大的.
創造萬物的主宰.
竟然為我們這麼微小的人.
用他的耳朵貼近我們.
可能就好像爸爸媽媽.
聽到晚上小朋友哭.
就走過去聽聽.
又或者看到他呼吸慢了.
用一個很留心.
並且很親密的關係.
去為了捕捉一個很弱小.
可能是小小的哭聲.
可能是小小的呼喊.
原來蒙恩的開始.
恩典的起點.
不是我們要多偉大.
或者做多精彩豐富的事.
或者我們要成為多偉大的人.
而是上帝親自願意.
去垂聽.
上帝親自願意來到我們當中.
一切開始就是上帝要質疑.
去聽我們.
當我們在軟弱困苦絕望.
受煩.

$^{121}$甚至只懂得呼求的時候.
上帝就質疑,聽我們的聲音.
所以我們的回應很簡單.
既然上帝向我們質疑.
我們一生的回應.
就是應該求告他.
在他面前蒙恩的人.
真正在主裡面蒙恩的人.
就是一個知道自己軟弱.
但都願意隨時向上帝呼求的人.
頂尖的有時候.
我們太擔心.
我們經常擔心就是上帝不聽我們的.
我那麼多罪.
上帝那麼忙,那麼多事去做.
哪有空理我.
我記得很多年前看過一套電影.
有一個人突然間.
他站在上帝的位置去做事.
發覺很忙,忙到沒辦法處理.
然後他就很驚訝,為何上帝可以那麼忙.
所以有人會這樣說.
上帝那麼忙,沒空理我.
又或者我這麼大問題.
他看不到我.
我們最軟弱.
當我們在一個最差的環境.
我們最需要的是呼求上帝.
所以頂尖姐妹,如果你想成為一個蒙恩的人.
第一件事你要記住.
任何時候,特別是最軟弱的時候.
請你記得,有勇氣.
向上帝呼求.
創世記十八章.
剛才聽了一段故事,很精彩,很豐富.
這個故事我們都可能熟悉.
就是亞伯拉罕的故事.
裡面講到當亞伯拉罕.
在萬利杖樹旁.
接待了三個神秘客人.

$^{161}$然後.
很有趣地.
他就對那位領頭的人.
這樣說,我主,我若在你眼前.
蒙恩,求你不要離開.
僕人,往前去.
這個故事講的是.
有三個人出現,死者在後面講.
用揚華描述,是上帝的侍者.
這三位侍者出現的時候.
就和亞伯拉罕有種相遇.
而亞伯拉罕就很快地.
盡力去接待.
然後他們就發出預言.
就應許.
他們就說,在明年這個時候.
我必會回到你這裡.
你的妻子撒拉必生一個兒子.
在最後最後.
聖經裡面這樣描述.
揚華說,豈有難成的事呢.
亞伯拉罕.
求在主面前蒙恩.
他求蒙恩是蒙什麼恩呢.
當時聖經沒有特別講.
蒙恩就是.
陪多一會.
我蒙恩,你上來陪我多一會.
上帝給的恩典竟然.
打破了我們的想像.
甚至打破了.
當時亞伯拉罕的期盼.
當時聖經裡面講到.
亞伯拉罕和撒拉已經年紀很大.
大家都知道當.
兒子依撒出世的時候.
亞伯拉罕已經一百歲,年紀很大.
撒拉也生育斷了.
聖經裡面講得很清楚.
當揚華的使者.

$^{201}$在講這句說話.
撒拉要生一個兒子的時候.
撒拉裡面就笑.
這個可能是心裡面暗笑.
那種想法就是.
不要說笑了,哪有可能的.
又或者是一種無奈.
一種苦笑.
究竟我們在生命裡面.
有沒有像撒拉那樣.
有時候也有這種暗笑.
或者苦笑呢.
主啊,我的壞習慣改不了.
算了,不行的.
說笑什麼,幾十年了.
我和家人的關係已經.
冰凍凍結了這麼多年.
和好?.
不要說笑了.
這個世界這麼混亂.
上帝真的重掌權的嗎?.
不要說笑了.
人很多時候.
面對不同的狀況,甚至覺得.
沒辦法解決.
最大的失望.
有時候就是絕望.
最無奈的就是我們覺得有些事情沒辦法.
甚至有人覺得不是的.
不要說笑了.
當撒拉.
對上帝的說話.
暗笑的時候.
上帝對撒拉這個暗笑.
發出了.
很重要的一個靈魂的敲問.
你不知道耶和華.
豈有難成的事嗎?.
你知不知道我們作為信徒.
我們作為追求信仰的人.

$^{241}$我們知道上帝全能,上帝很厲害.
在上帝裡面,當然沒有難成的事.
一在上帝,凡事都能.
金句也會背.
不過我們總覺得.
上帝能力可以移山填海.
不過我這件事上帝就搞不定.
我這個關係上帝就搞不定.
我這個狀態沒辦法.
死了一樣.
死了也要說,如同死人一樣.
在上帝面前蒙恩.
不是因為我們有什麼能力.
或者我們有多厲害.
是因為上帝是大能的上帝.
恩典就是我們在已經絕望.
毫無辦法的時候.
上帝的應許.
依然能夠成就.
這個就是恩典.
上帝的恩典永遠都比我們想像的大.
我經常說這一句.
上帝的恩典永遠比我們想像的大.
我們想像的那麼大.
上帝的恩典就更大.
因為我們以為不能的地方.
在上帝,凡事都能.
頂智慧你要知道.
恩典,蒙恩的人.
是要看見上帝.
在不可能裡面成就.
他的應許.
如果說舊約.
剛才的兩段經文是將恩典.
成為一個預表的圖畫.
新約聖經裡面.
羅馬書就將蒙恩的核心.
最清楚地,很震撼地.
告訴我們.
羅馬書第五章,第六節和第八節.

$^{281}$這樣說,因為我們還軟弱的時候.
基督就按時候.
為罪人死.
唯有基督在我們還作罪人的時候.
為我們死.
上帝的愛就在這裡.
向我們顯明.
你試圖保羅在聖經裡面.
很多的書信,很多的教導.
很多精彩的吩咐.
但是保羅在這裡.
很清晰地在羅馬書裡說.
當我們未蒙上帝恩典之前.
我們有幾個狀況,一個叫軟弱.
一個叫罪人.
一個叫仇敵.
在地上的邏輯很簡單.
你做得好,你表現得好.
就愛你,你有價值.
我就回報你,你工作做得好.
就讚賞你,升職加薪.
你聽說你乖,我就疼你.
很多父母可能這樣的邏輯.
去對小朋友,你不乖我就罰你.
你乖就獎你,當然獎罰.
給小朋友學習當然重要.
但在上帝裡面.
這不是他的恩典的表達.
這是交換.
這個世界很喜歡用交換的概念.
去想所有的事情.
包括待人好.
對人有恩賜.
這個好像是交換.
因為你對我好,我對你好.
因為你聽說我對你好.
因為你做得好,所以我對你好.
但在上帝的恩典裡面.
正正要打破.
這個交換.

$^{321}$交易,上帝的恩典.
就是那個狀況.
當我們還是一團糟,很混亂.
甚至背叛他.
同時可能甚至扮演.
仇敵的角色的時候.
耶穌基督.
就已經為我們釘在十字架上.
這個是最精彩.
最豐富,最大最大的恩典.
很多人問.
為什麼基督教的信仰這麼特別.
有些信仰有很多不同的要求.
基督教信仰的要求是什麼呢?.
你信就行,這樣真的行嗎?.
不是應該做些事情嗎?.
我們很想做些事情去換取救恩.
真的有弟兄姊妹.
有朋友來我面前問.
牧師,我背多幾句金句.
是不是會多些獎賞.
信耶穌這麼簡單.
是不是一定要看聖經一次.
背幾句金句,回十次八次教會.
才等於信了耶穌.
當然鼓勵大家讀聖經.
背金句回教會.
但聖經裡面很清楚講到.
上帝的恩典正正就是那種.
當我們還是罪人.
還是不行的時候,上帝就愛了我們.
耶穌基督為我們死.
我們很乖,我們肯追隨祂.
我們肯盡力在宗教裡面.
向前走,所以祂就給你多些恩典.
不是的.
當我們還是罪人,還是受敵的時候.
基督就為我們釘十字架.
基督是為罪人釘十字架.
而罪人就是我們.

$^{361}$這個描述是講.
一個叫白白的恩典.
為什麼是一個白白的恩典.
就是因為其實我們在上帝面前.
沒有東西可以跟祂交換.
你憑什麼去到上帝面前.
跟祂交換.
你的罪有什麼辦法自己解決.
你的和上帝的關係.
是決裂離開.
你有什麼辦法自己走近.
所以原來我們從來沒有辦法.
因為我們沒有辦法,上帝很明白.
所以上帝的愛.
大到一個地步就是.
當我們最軟弱.
最負面,最離開上帝的時候.
甚至是祂的敵人的時候.
祂已經為我們死,救恩已經放在那裡.
你只要願意去回應.
你願意去用你的行動.
用你的心,用你的信心.
擺上,是呀,上帝.
我沒有東西可以回應你.
我沒有能力去賺取救恩.
上帝我就是這麼軟弱.
但是多謝你的愛.
當你這樣回應的時候,白白的恩典.
我們就能夠跟上帝去相通.
能夠有生命得到指望.
能夠.
找到我們的方向.
記住這些完全不是我們的功勞.
不是我們做到什麼.
在患難裡面我們要歡歡喜喜.
因為這個是恩典的根基.
是基督.
為我們死和復活.
恩典的根基是基督.
為我們死和復活,而不是我們.

$^{401}$有多大的能力或者做了多大的事情.
所以本來這麼說.
我們甚至在患難裡面.
都要歡歡喜喜的.
因為知道上帝的愛已經藉著聖靈.
囂貫我們的心.
在主前蒙恩.
是給我們在苦難裡面最大的盼望.
弟兄姊妹.
原來上帝要施恩給你.
是在你最軟弱.
最有難處.
甚至患難當中的時候.
已經臨在了.
所以這一刻你需要的話請你去呼求.
在.
馬達科罕裡面.
很清楚的描述.
耶穌走遍各城各鄉.
看見許多的人.
就憐憫他們.
因為他們困苦流離.
如同羊.
沒有牧人一樣.
接著就數了.
十二個門徒的名字.
弟兄姊妹有時經常都.
記不到很多門徒.
但聖經裡面.
刻意要數門徒的名字.
告訴我們.
這些就是耶穌吩咐.
他們要去的.
重點不只是名字.
而是他們要去做些什麼.
他們要醫治病人.
死人復活.
把鬼趕出去.
裡面很清楚講到.
他們得到上帝的能力.

$^{441}$門徒當日得到的恩典.
不只是感受上.
和耶穌一起美好的恩典.
而是他們得到一些能力.
而那些能力就像耶穌一樣.
可以醫病趕鬼.
可能弟兄姊妹說.
如果我們有門徒這些能力.
能夠醫病趕鬼死人復活就好了.
我們好像今天沒有這些能力.
當然我們今天的能力.
和昔日門徒經歷的有不同.
但我們同樣一樣.
有能力是什麼呢.
有能力去回應上帝的愛.
和他的恩典.
這段經文裡面最後講的.
你們白白地得來.
就要白白地捨去.
很獨特這段經文.
講一個很重要的.
重心.
這個重心就是耶穌.
他自己的一個示範.
他動了憐憫的心.
他見到很多困苦流離的人.
他動了慈心呼召門徒.
你去吧.
你幫我一起去完成這些的使命.
這十二位門徒.
本身就是那些跟著耶穌.
他可能在他的難處裡面.
你看到聖經裡面講耶穌呼召門徒有不同的狀況.
有些可能本身是罪人.
人體來的罪人.
所以放下工作跟著耶穌.
有些人忙碌著自己的生計.
捉魚生活.
耶穌放下忙就跟隨耶穌.
每一個都是在他的狀況流離.

$^{481}$在他的難處限制裡面.
耶穌就是這樣憐憫他們.
呼召他們.
讓他們成為門徒.
所以耶穌就呼召這一班.
被愛被憐憫的門徒.
他去出去.
見到有需要的人就憐憫.
和愛.
所以今天.
耶穌和我們都講同樣的說話.
你們白白的得來.
都要白白的捨去.
是我們今天.
可能沒有當日門徒那種.
即時見到神蹟奇事.
醫病趕鬼的能力.
未必有 上帝許可當然可以有.
但更重要的是我們有沒有看見.
原來上帝的恩典.
是白白的臨到我們身上.
而我們需要白白的.
去使捨出去.
去分享出去.
在主前蒙恩.
為了赦免的恩.
醫治的恩 安慰的恩等等.
不是為了讓我們很舒服.
上帝醫治我.
赦免我 安慰我.
我就很舒服地享受一個.
屬靈狀況很好 很安逸.
耶穌在這裡親自教導.
蒙恩的人.
必須要成為.
施恩的人.
真正蒙恩的人.
不只是坐在那裡.
蒙恩真是好了.
是要將上帝的恩典.

$^{521}$白白的跟別人分享.
當我們面對.
這個世代 很多人.
同樣都很困苦流離.
我們很多時候.
都很冷眼旁觀.
在生命裡我們可能.
已經很怕去幫人.
今時今日的教導就是很多人欺騙你.
你要小心.
所有跟你講需要的人都是欺騙你的.
但當我們見到.
不同的人需要的時候.
上帝有沒有憐憫的心.
在你當中.
上帝有沒有提醒你.
有需要的人需要你去伸手去幫助.
當然我們要醒覺.
不要被人欺騙.
但同樣地見到有需要的人.
特別是心靈需要的人.
靈魂需要的人.
我們有沒有想過.
將上帝給我們白白的恩典.
這個福音傳給.
身邊同樣需要的人.
各位弟兄姊妹.
說到蒙恩.
一開始可能你們在想.
我身上經歷了多少恩典.
我身邊的人有多少經歷.
我感恩.
今天很想和你們將幾段經文串在一起.
再看一段很重要的概念.
第一.
我們知道我們蒙恩.
是因為我們生命不再孤單.
因為我們這位上帝.
是側耳聽我們的上帝.
弟兄姊妹你們要記住.

$^{561}$上帝是側耳聽我們的上帝.
祂常常都很想.
聽我們的聲音.
在我們的信仰追求裡.
很多時候都問.
很多弟兄姊妹都問.
牧師我怎樣知道上帝的聲音.
我想知道上帝的聲音.
第一就是要熟悉祂.
看聖經.
看教會.
多祈禱.
同樣地你都要將你的事跟祂說.
我們上帝想和我們互動.
當然你跟祂說.
祂不是說你要什麼.
而是說你的心事.
更重要的是.
上帝願意側耳聽每一個人.
很多時候我們到了一個地步.
有時我們尋問上帝的心意.
甚至我們禱告.
去到不知怎樣說.
因為我覺得上帝都不聽.
又或者是我自己一個人說.
但我們深信上帝是側耳聽我們的上帝.
上帝在我們有難處.
困惑困苦.
甚至遭迫害的時候.
我們任何呼求上帝都聽到.
我們就好好的呼求.
我常常和頂尖妹分享.
很多人覺得禱告好像沒有什麼可以做.
或者做不到很多事的時候.
最後就禱告.
很多頂尖妹都問.
每一次崇拜我們都公討很多內容.
其實我們都未必真的做到很多事.
是不是做不到所以我們唯有禱告.
當然啦.

$^{601}$很多時候我們真的做不到很多事.
但同樣地我常常提醒自己.
也和頂尖妹分享.
不是因為我們做不到而我們禱告.
而是因為我們深信我們和上帝溝通.
告訴上帝我們是心聲.
和上帝的呼求.
是我們心中最重要的事.
在一切我們要先做.
不是想我們做到.
未做到以先我們先要向上帝呼求.
去尋問.
去求告.
去告訴他我們是心事.
不是因為他不知道.
而是因為我們好想和他有更親密的關係.
這個側耳聽我們的上帝.
他想聽的是你的心跳.
你的心聲.
所以你就將你的心事向上帝講.
禱告最大的權柄.
就是我們可以尋求最大的幫助.
所以不是沒事做才去祈禱.
而是我們先去到上帝面前.
將最重要的.
將最想分享的去和他講.
他會側耳聆聽.
他也會回應我們的需要.
因為蒙因我們生命不再絕望.
因為上帝沒有難成的事.
亞伯拉罕的例子告訴我們.
什麼是上帝的心意.
上帝的心意是要讓他成為萬國之父.
成為最重要的以色列民的一個很重要的列祖.
就要透過以撒亞國.
所以上帝的心意就是這樣.
上帝的定義就是這樣.
對亞伯拉罕來說.
不可能發生的.
我已經年紀老了.

$^{641}$撒拉還要偷笑.
還要苦笑.
上帝是會覺得絕望的.
但弟兄姊妹記住.
我們從來都不在絕望當中.
因為上帝仍然會看顧.
上帝仍然會保守.
當你最不想祈禱的時候.
你記住這句話.
你就去祈禱.
當你最想逃避.
你覺得上帝都不會理我.
我都是絕望.
我都是躲起來.
甚至你覺得滿身醉污.
不配去見上帝的時候.
上帝顯得最大.
我常常這樣說.
當最黑暗的時候.
微弱的光線都能夠引你走向光明.
所以不要沉溺在自己的黑暗當中.
不要躲起來覺得已經絕望.
沒有辦法.
因為在上帝沒有難成的事.
在人不能的.
在上帝凡事都能.
因為蒙恩.
不需要再絕望.
因為蒙恩我們生命會有滿足感.
安全感.
因為我們知道.
當我們在座罪人的時候.
上帝就愛我們.
不是要我們好好表現.
上帝才愛我們.
當然當我們回應上帝的愛.
成為上帝的兒女追求的時候.
我們生命應該有不同.
我們應該有新的生活.
新的生命.

$^{681}$但不是因為這些努力.
去換取上帝的恩典.
不是因為我們做了多少好事.
讀多少學位.
上帝的恩在我們在座罪人的時候.
已經臨在.
基督在我們在座罪人的時候.
為我們死.
所以每一位朋友,弟兄姊妹請你謹記.
當我們還是罪人的時候.
當我們還是.
基督的敵人的時候.
他已經為我們死了.
所以這個愛就是這樣顯明了.
讓我們能夠知道.
怎樣去回應.
因為蒙恩我們生命有了方向.
弟兄姊妹記住.
白白的禮,白白的捨去.
我們每一個人都要好好去見證.
上帝給我們身上的恩典.
真正蒙恩的人,是施予的人.
越施予得多.
你越能夠蒙恩更深.
我很喜歡聖經裡面一句說話.
很多人都會說.
施比受更為有福.
為什麼呢?因為當你施予的時候.
你才擁有,你才可以施予.
你完全沒有的時候,你只可以默默的接受.
就好像我們在座罪人的時候.
默默接受上帝的恩典.
你才能夠施予.
所以當我們施予的時候,就會更為有福.
弟兄姊妹讓我們回應這個宣教的使命.
將白白得來的恩典.
白白的去捨去.
願我們每一個人.
每天都心心在主前.
去經歷他的豐富.

$^{721}$真正成為蒙恩的人.
能夠昂首挺胸.
走合人群.
去見證,去分享.
去傳遞.
每一個人都需要的福音.
請同心祈禱.
天父上帝我們獻上感恩.
你的恩典豐厚.
你賜給我們每天的生命氣息.
你賜給我們能夠.
有豐富,有生活.
你賜給我們有教會.
有崇拜,多謝你.
因為你無為不值的.
看顧,讓我們體會.
生命的豐盛和美好.
基督我們多謝你.
你愛我們,你愛我們到底.
你甚至.
將我們去拯救.
在我們軟弱.
還是罪人,甚至可能敵擋你的時候.
基督你就為我們死.
你為我們受死.
以致我們生命.
可以得到改變.
你的復活成為我們生命的盼望.
和我們活在一樣.
基督我們多謝你.
求你賦予大能在我們身上顯大.
以致我們能夠去經歷豐富.
聖靈我們多謝你的提醒.
繼續幫助我們.
讓我們被提醒.
我們是蒙恩的人.
我們這班蒙恩的罪人.
今天領受你豐富的恩典.
就提醒我們要去.
洗萬民做你的門徒.

$^{761}$要去讓更多人和我們一起.
經歷豐富的恩典.
和我們一起來到你的殿當中.
去體會明白.
生命的精彩和豐富.
更知道我們應該怎樣處.
面對這個可能.
不容易身處的世代.
更有力量去經歷.
你所知的一切.
去見證你的榮美.
幫助我們眾人在未來的路上.
無論是多窄多難走都要努力堅守.
以致有一天.
我們去到你面前經歷你的豐富.
能夠蒙你的月立.
為你的恩成為.
你所愛良善終生的僕人.
多謝主在我們當中.
多謝主你的恩由.
以上感恩加諾禱告奉耶穌基督的名禱.
\newpage



\section{瑪拉基書 2:17-3:5}
\label{sec:032293}
\textbf{你都有今日}~~~~石偉漢牧師
\newline
\newline
連結: \href{http://www.church.com.hk/acms/content.asp?site=cdc\&op=show\&type=product\&code=032293\&layout=sermon}{\texttt{032293}} ~~~~ 語音日期: 2026-06-14
\newline
\newline
\hyperref[sec:032292]{\small{< < < PREV SERMON < < <}}
~
\hyperref[sec:index_chronic]{\small{[返順時目]}}
~
\hyperref[sec:032294]{\small{> > > NEXT SERMON > > >}}
\newline
\newline
瑪拉基書 2:17-3:5
\newline
\begin{longtable}{cl}
\hline
\hline
章節 & 經文 (和合本修訂版)\\
\hline
2:17 & \begin{tabularx}{0.7\textwidth}{X} 你們用言語使耶和華厭煩，卻說：「我們在何事上使他厭煩呢？」因為你們說：「凡行惡的，耶和華看為善，並且喜愛他們；」又說：「公平的神在哪裡呢？」 \end{tabularx} \\ \\ \relax
3:1 & \begin{tabularx}{0.7\textwidth}{X} 萬軍之耶和華說：「看哪，我要差遣我的使者在我前面預備道路。你們所尋求的主必忽然來到他的殿；立約的使者，就是你們所仰慕的，看哪，快要來到。」 \end{tabularx} \\ \\ \relax
3:2 & \begin{tabularx}{0.7\textwidth}{X} 他來的日子，誰能當得起呢？他顯現的時候，誰能立得住呢？因為他如煉金匠的火，如漂洗者的鹼。 \end{tabularx} \\ \\ \relax
3:3 & \begin{tabularx}{0.7\textwidth}{X} 他必坐下如煉淨銀子的人，必潔淨利未人，熬煉他們像金銀一樣；他們就憑公義獻供物給耶和華。 \end{tabularx} \\ \\ \relax
3:4 & \begin{tabularx}{0.7\textwidth}{X} 那時，猶大和耶路撒冷所獻的供物必蒙耶和華悅納，彷彿古時之日、上古之年。 \end{tabularx} \\ \\ \relax
3:5 & \begin{tabularx}{0.7\textwidth}{X} 萬軍之耶和華說：「我必臨近你們，施行審判。我必速速作見證，警戒那些行邪術的、犯姦淫的、起假誓的、剝削雇工工錢的、欺壓孤兒寡婦的、屈枉寄居者的和不敬畏我的人。」 \end{tabularx} \\ \\
[1ex]
\hline
\hline
\end{longtable}
$^{1}$弟兄姊妹主 平安.
很高興見到大家一齊來崇拜.
一起敬拜主.
也問候家中的弟兄姊妹.
願你們平安.
我們一起有個祈禱.
親愛的主 多謝你.
因為你讓我們有生命的氣息.
讓我們有生命的氣息去敬拜你.
去聽你的話 聽你的吩咐.
願主幫助我們.
讓我們生命裡面有你與我們同在.
我們都能活出主你喜悅的生命.
求你帶領我們的學習.
我們的思考.
我們的聆聽.
祈禱奉耶穌名字而求.
阿們.
今天和大家分享一個題目.
叫做「你都有今日」.
我不知道你有沒有說過這句說話.
你什麼時候會說這句說話的呢.
大家心裡都笑了 也不敢說出來.
「你都有今日」.
很多時候我們都想起畫面就是這樣.
有班這樣的人.
很壞壞地跟你說.
「你都有今日」.
我有沒有說過呢.
我也有說過這句說話.
我什麼時候說過這句說話呢.
原來在大約三年多前.
我說過這句說話.
是什麼情境說呢.
話說在大約三年多前.
我還沒有來美國之前.
我都要安頓好那邊香港教會的各樣事情.
其中除了要處理教會的堂主任的接班.
和教會一些轉制的工作.
另外一個很重要的任務.

$^{41}$就是安納我們的同工.
因為事實上.
現在你眼前見到的這兩位同工.
分別在網上崇拜.
在網上都跟我們說過道.
一位就是劉牧師.
一位就是黃牧師.
事實上他們過去的日子.
他們都忠心耿耿.
在教會裡面都一直參與.
教會的青少年的牧養.
婦女的牧養.
和整體教會的一些關顧.
探訪等等不同的工作.
所以經歷了幾年.
亦都通過我們的總會.
開放了多一些.
容易安納牧師.
因為以往你都知道.
華人教會安納一個牧師.
是重重艱辛.
十幾二十年都不會安得到.
這裡可能有些不同.
但在香港如果你能夠.
十幾二十年都等得到.
安得到都是一個很大的恩典.
但我們經過轉制之後.
我們這兩位同工.
很用心去牧養教會的時候.
我肯定他們.
亦都得到我們教會的肯定.
所以終於在我離開之前.
可以順利在2023年3月26日.
為他們安納牧師.
所以當時我和他們說的一句說話.
就是「你都有今日了」.
這句說話是甚麼呢?.
是一個祝福.
因為牧養辛苦了很多年.
其實牧養不是為了命.

$^{81}$不是為了利.
不過這個鹹頭.
這個積鹹.
是對他們生命事奉的肯定.
所以「你都有今日」.
對他們是一個很大的祝福.
因為他們能夠成為牧師.
今日仍然在香港.
我之前牧養的教會裡面.
繼續擔當著教會的領導.
事奉等等.
為甚麼他們有今日呢?.
因為他們的昨日做得好.
他們之前用了很多的勞苦.
所以他們有今日.
如果在幾年之後.
再問你.
幾年之後的今日會是怎樣呢?.
結果會是怎樣呢?.
其實就很視乎你.
今日如何去對待你的明日.
所以今天我們都去.
看這個題目的時候.
正好思考著.
我們昔日所做的事.
會不會影響著我們今日所做的事.
今日的結果.
同樣的.
今日我們都是一個抉擇的時間.
去決定我們的明日如何.
所以在你的明日.
將會聽到一句怎樣的說話.
究竟是「你都有今日」是正面.
還是負面呢?.
就真的看你今天如何.
就好像聖經裡面所說.
你的日子如何.
你的力量也必如何.
你的日子怎樣過呢?.
究竟你將來的日子會怎樣過呢?.

$^{121}$所以今天我們都會去.
藉著馬拉基書裡面.
二章十七至三章五節的部分.
去看看那一天.
即是我們說今天那一天.
那一天是一個怎樣的一天呢?.
我們就透過三個分題.
去逐個段落.
去了解一下這個部分.
第一個段落就是來自.
二章十七節的段落.
讓我再讀一次.
「你們用言語使耶和華厭煩」.
「卻說我們在何事上使他厭煩呢?」.
「因為你們說反恨惡的」.
「耶和華看為善」.
「並且喜愛他們」.
「又說公平的上帝在哪裡呢?」.
這裡其實承接著我們一直.
在看馬拉基書的部分.
如果你還記得有印象.
我們在馬拉基書裡面的第二章.
的部分是說什麼呢?.
上梁和下梁.
有幾個印象.
上梁不正.
下梁歪.
上梁就是我們上次才看到的.
就是祭司.
在當時來說.
就是一個回歸的國土裡面.
那些祭司有一個屬靈的.
帶領的職份.
帶領百姓怎樣如何呢?.
從過聖見生活.
去獻祭的一班人.
不過他們卻不是獻上帝喜悅的祭物.
當他們是這樣的時候.
也都帶來下梁就是那些眾百姓.
也都活在一個欺詐.

$^{161}$一個詭詐的角度裡面.
他們活出一種詭詐的生活.
這個是我們上次所看到的.
所以你見得到.
就是在整個以色列的回歸的角度裡面.
他們就見到那個土地裡面.
並不是他們所期望的.
因為理論上就是.
聖殿已經重建.
應該聖潔的生活.
也都應該是會被重建.
但是事實上卻不是這樣.
卻不是這樣.
當時有很多人見到這樣的社會環境.
他們都會帶來很多的疑問.
特別他們心裡面有一些想法.
他們很期待著一個日子.
特別在這裡說一種什麼狀況呢?.
就是一種彌賽亞降臨的日子.
很想這個國家能夠再一次.
不被其他國家壓制.
重新建立一個自己的角度.
重新去有一個他們期待的救主.
他們彌賽亞來幫助他們.
以致到他們真的不醉一躍.
但在這個等待的日子裡面.
剛才面對了那種上樑下樑的一些問題.
所以他們心裡面就出現了一些疑問.
這種疑問怎樣去表達呢?.
就在乎他們本身的信念.
不知道你們有沒有看過報紙.
現在沒有人買報紙的.
現在沒有人買報紙了.
很多報紙都倒閉了.
不過雖然我們不買報紙.
但我們對某些報紙是極度相信的.
我介紹兩份報紙.
很多人都會常常去留意和相信的.
就是一份叫善報一份叫惡報.
因為有些人會相信善有善報.

$^{201}$惡有惡報.
所以這兩份報紙大家都期望.
做善事應該有善報做惡事應該有惡報.
這也是以色列百姓在耶路撒冷重歸的故土裡面.
他們深深期待的.
很多人都相信這兩種報紙.
善報和惡報.
不過原來還有一種報紙.
在這個社會裡面更多出現的.
就是微報.
因為若然微報時辰未到.
而事實上微報真的常常發生.
微報當然不是指一份報紙.
我只是想借這個意思說一種微報未出現的結果.
你期望善有善報.
你期望惡有惡報.
但很多時候善報和惡報都不是即時在你面前發生的.
往往你會見到很多應該有報應的人.
見不到他們的報應.
甚至乎那些正在做不好事的人.
他們繼續活得好.
甚至活得比你更好.
這就是那種狀況.
以致到人去問一個問題.
就是這個世界是否有這個字.
這個字有四個字.
就是顛倒是非.
即是人們都會問一個問題.
是否是是非已經被顛倒被反轉.
以致到惡人有好報.
善人面對惡報.
所以好好的在耶路撒冷生活的人.
他們卻被人欺壓.
上次說到的是什麼.
壓情.
壓寵.
很多詭詐.
婚姻裡面用一些古怪的方法.
希望離婚然後重新找另一段婚姻.
那些是非不斷顛倒的時候.

$^{241}$就帶來很多人心裡的疑問.
所以我們今天讀到的17節裡面.
正正好就提出了三個疑問.
其中一個他們都會常常問到.
所以他們說用言語去不斷問神.
問到的是什麼呢.
使耶和華厭煩 包括什麼呢.
我們在何事上使他厭煩呢.
他們覺得不斷去質問神.
不斷去問神的時候.
覺得不是令神厭煩.
我們什麼時候有做到這件事呢.
第二就是犯行惡的耶和華看為善.
原來做壞事的人.
耶和華看他們是善的.
還很喜愛他們.
所以令他們有好報.
這個在眼裡面.
在現實生活裡面.
真的好像是這樣的.
行惡的人竟然有好報.
所以第三個向神質問的問題是什麼呢.
公平的上帝在哪裡呢.
究竟有沒有公平的上帝呢.
不知道你有沒有問過一些信仰上.
有時候在你信神的過程裡面.
有些生活令你問一些問題.
究竟神你想怎樣呢.
究竟神你是怎樣呢.
究竟神你是否我心目中所想的那個神呢.
事實上在聖經裡面.
有時候有些先知.
都會同樣有一些對神的發問.
包括在這裡見到.
左邊的代表耶利米.
右邊的代表哈巴谷.
他們都向神發出問題.
事實上我們信神的人.
我們可不可以向神發出一些問題呢.
是可以的.

$^{281}$你不明白就去問神.
是可以的.
不過有時候我們去到一個地步.
就是我們問神問到什麼呢.
對神失去了信心.
失去了信心的時候.
以致見到那些環境和你的期待不一樣的時候.
就開始產生疑問.
而疑問之後就是懷疑.
懷疑之後是什麼呢.
就是抱怨.
然後是什麼呢.
言語的否定.
甚至乎言語的暴力.
是呀我們問了神一些東西.
好像答不到.
越來越懷疑.
越來越不斷反覆在懷疑的困局裡.
以致到什麼呢.
開始抱怨神.
否定了神是一個善的神.
甚至乎跪神觀賞一個污名.
就是神是不行善的.
甚至乎喜愛行惡的.
甚至乎什麼.
剛才說到言語上.
對神作出一些懷疑抱怨.
如果你還記得我們看二章十三至十六節的時候.
說到百姓對妻子行詭詐的時候.
都是什麼呢.
慢慢就是那種暴力.
遠離那個妻子.
然後就另外再娶一個妻子.
所以如果我們在言語上.
在懷疑裡面.
如果去到一種言語的否定.
慢慢我們都會遠離了神.
我們不知不覺都不會相信神.
這個也是一些人現在在信仰生命裡的處境.
我不知道你有沒有經歷過這樣的狀況.

$^{321}$好像信仰裡面有些問題.
看聖經又好像看不到答案.
或者問牧師又好像拿不到答案.
怎麼辦呢.
會不會就這樣去懷疑了.
這位你所相信的上帝呢.
事實上我們今天都習慣了.
有什麼不明白就問誰.
問AI.
AI你們用什麼軟件.
有Gemini,ChatGBT等等.
不同的軟件.
我們今天買手機.
第一時間買有AI功能的.
就第一時間可以問到答案.
但有時候我們知道信仰的疑問.
是不是一定等於即時可以問到答案呢.
有時候都不一定.
我們信仰面對困惑的時候.
有時候不一定找到答案.
甚至乎有時候隨著時代的改變.
我們的疑問會越來越多.
問題是你有沒有因為這樣而懷疑神.
遠離神疏遠神呢.
你知不知道最近五月份的時候.
其中有一單新聞.
就令到我們很多基督徒都要面對一種新的挑戰.
知不知道是什麼新聞呢.
由美國政府頒布的新聞.
有沒有看過這單新聞.
美國政府就說.
在2026年5月份的時候.
透過官方的網站.
發佈了一個叫不明異常現象的解密檔案.
當中發佈了一些很奇怪的相片.
講一些影片.
在天空上可能有些不知道是什麼的物體.
但卻突然間有些光的物體出現.
這個成為了基督徒都要面對的一個問題.
當中講什麼呢.

$^{361}$這個神秘的片段裡面.
講很多奇異的光體的時候.
其實那些人就開始問.
是不是有UFO.
更加問的不是單單是.
因為UFO其實是代表了不明飛行物體.
Unidentified Flying Object.
這些不明飛行物體裡面.
他們最關心的問題.
這個物體裡面會不會有一些不明的生物.
於是當這個公佈了之後.
也有一個科學家說.
他有份參與一些物體的研究.
參與政府的研究.
他說在這些不明的物體裡面.
裡面找到四種不一樣的人物.
所以現在網上就說有這四種.
不知道是不是外星人.
原來真的有羅茲威爾.
原來這裡說的是什麼呢.
包括有昆蟲人.
當中說有些像昆蟲那樣的人.
有手有腳的一些物體.
昆蟲人.
有小灰人.
就是大大眼睛的灰色的.
另外還有蜥蜴人.
最後還有一個比較像人的.
就叫北歐人.
就像北歐人那樣大大隻的.
那些人物.
當有人問你.
你會怎麼回答他呢.
人家說這個世界上真的有UFO的.
有外星人的.
所以你們信耶穌是假的.
你會不會有這樣的疑問呢.
或者你怎麼去回答這些疑問呢.
弟兄姊妹.
有時候我們面對著一些.

$^{401}$突如其來的事情.
心裡解釋不了的時候.
不知不覺會引誘我們.
拖累我們.
好像信神信得越來越不真實.
越來越不堅韌.
當然這裡要說一點.
如果真的有UFO.
如果真的有外星人.
我們信耶穌是不是真的呢.
你認為還是真的請你舉手.
那你認為因此.
不可以再信耶穌請你舉手.
那你想想.
發現了又怎樣呢.
正如我們有很多事情.
我們都說不斷發現的.
在這個社會裡面.
其中一個就是.
例如我們說尼芝湖水怪.
大家聽過尼芝湖水怪.
不斷都在傳.
尼芝湖有隻水怪在那裡.
長長頸那樣的.
偶爾有些照片就說有這些.
我想問你.
就算被你發現了尼芝湖.
真的有隻水怪.
是不是代表上帝是假的呢.
只不過代表我們終於發現了.
一種我們不難以命名的生物.
又再出現了.
如果你發現了一些昆蟲人.
是不是代表你信耶穌是假的呢.
只不過發現了這個宇宙裡面.
有另外一種.
都是上帝所創造的生物.
並不能否定上帝是那位.
創造萬有的主.
反而就能夠透過.

$^{441}$原來上帝創造了.
很多不同的東西的時候.
我們就明白.
什麼叫做諸天說說神的榮耀.
窮倉全揚他的手段.
在這個宇宙萬物裡面.
原來我們都可以透過.
看到上帝創造的奇妙.
更加重要的.
無論你找到多少新的生物也好.
你都不能否定.
神對你的愛是真實的.
神願意為你犧牲了他的愛子耶穌.
去帶來我們生命的拯救.
所以當我們面對.
信仰移民的時候.
我們就真的要想一想.
就算解答不了.
我們要明白.
我們的信仰其實都是真實的.
所以我們之前都說過.
在香港大火裡面.
這位牧師.
他其中有三位的親人.
就在大火裡面離世.
他也同樣面對一個問題.
為什麼上帝讓他這三位親人.
在大火裡面離世呢?.
可能心裡面也有很多疑問.
特別他作為一個牧師.
他真的面對著親人離開的時候.
他也有很多信仰的疑問.
不過他得出來的結論是什麼呢?.
他就說他明白.
他能夠藉這個機會.
去和這班大火的居民.
去面對哀傷和盼望.
確實面對著一份很真實的傷痛.
心裡面都好像信仰解釋不了.
不過那份盼望就讓我們知道.

$^{481}$就是在神的裡面.
有一天我們會再一次相見.
有一天會再一次相見.
所以即使我們見到.
好像那個環境裡面.
很多懷疑.
我們仍然要相信神是能夠.
你相不相信你這個神.
是可以做得到的呢?.
所以剛剛最近.
我們祈禱紙裡面.
我都分享給大家.
就是最近有一個訪問.
就是訪問了一個著名的.
無知的報導家.
無知就是無手無腳的報導家.
就是Nick.
其中他就說.
很多人都懷疑他是disabled.
因為他無手無腳.
他無手無腳什麼都做不到.
不過他說不是.
因為他明白.
別人看他是disabled.
不過他說.
只要在前面加一個go.
這個字的話.
就會變成了.
God is able.
神能夠.
他在生命裡面.
他找到一個答案.
就是他雖然是四肢.
沒有一個健康的四肢.
不過他知道.
神能夠使用他.
成為一個出色的報導家.
所以弟兄姊妹.
第一個我們要思考的問題.
就是我們信仰裡面.

$^{521}$會不會都像這群意識的人那樣.
有時候我們不知不覺.
在一個信仰的懷疑裡面.
特別當見到.
善未能夠有善報.
惡未能夠有惡報.
甚至乎顛倒是非的時候.
你相不相信.
在終極裡面.
上帝就是那位.
掌管一切的主.
他有能力.
他在最後裡面.
作出他最好的審判.
我們這位神.
確實在我們生命裡面.
很多的疑問裡面.
他並不是讓我們.
在一個混亂當中.
很多時候他更加希望.
給到一些方向和答案.
讓我們知道.
他初步的方向.
所以在第二個部分.
B裡面.
除了第一部分.
未報之日.
好像見到什麼都沒有.
得到報應的日子裡面.
神就告訴我們什麼呢?.
B就是要期待之日.
我們要對將來的日子.
有一份清楚的期待期盼.
而這份期盼.
就是神最終會拯救.
經文就在第三章.
一至四節的部分.
萬君子耶和華說.
漢雅我要差遣我的使者.
在我前面預備道路.

$^{561}$你們所尋求的主.
必忽然來到他的殿.
納入的使者.
就是你們所仰慕的漢雅.
快要來到.
他來的日子.
誰能當得起呢?.
他顯現的時候.
誰能納得住呢?.
因為他如煉金匠的火.
如漂洗者鹼.
他必坐下.
如煉證銀子的人.
必結證利未人.
奧煉他們匠金銀一躍.
他們就憑公義獻供物.
耶和華.
那時猶大和耶路撒冷.
所獻的供物.
必蒙耶和華月立.
彷彿古時之日.
上古之列.
在這裡三章一節.
就說一個狀況.
神就回答.
你們覺得好像現在解答不了.
你還未見到.
那種報應.
臨到在一些人身上的時候.
神說.
是的.
你等候.
有些事必然發生.
原來我們有時需要去等待.
等候.
因為有些事必然發生.
特別是神的拯救.
必然在之後會發生.
所以在這裡說.
他們這些以色列的百姓.

$^{601}$面對著他們眼裡的困局的時候.
神就說.
有三種人.
你要去等待.
有些人要等待.
等待什麼呢?.
包括了你們所尋求的主.
必忽然來到他的殿.
就是神自己.
當你無望的時候.
神自己.
有一天會親自臨到當中.
臨到他的殿裡.
所以你尋求的主就是神自己.
但同時間說.
當神說他要來臨之先的時候.
亦都說到.
萬君的耶和華.
要差遣他的使者.
在萬君的耶和華面前.
預備道路.
我們知道.
一個國家的元首去外地的時候.
先有使者去之前鋪排一切.
神亦都說.
他都會派遣定一個使者.
在之前預備.
在他前面鋪排一切.
除了神這個.
預備道路的使者之外.
原來後面亦都說到.
有一個立約的使者.
就是你們所仰慕的.
這個立約使者很不一樣.
是令人很仰慕.
很期待的那一位.
他們分別是誰呢?.
如果我們再看下去.
《瑪利基書》之後開展的.
就是《馬達福音》.

$^{641}$《福音書》等等的話.
我們就會明白.
這個更加指向的是什麼呢?.
施洗約翰和耶穌.
因為施洗約翰就要.
預備住的道.
修植他的路.
去預備住耶穌基督的降臨.
而立約的使者.
就是和我們納一個新約的那一位.
就是主耶穌自己.
所以你看得到.
當耶穌在一帶河.
接受洗禮的時候.
就好像三個同在的畫面一樣.
天上有聲音說.
「善我的愛之仇」.
就是那位的神在當中.
同時間施洗約翰.
亦都在這個使者旁邊.
為著耶穌去施洗.
而這個立約使者.
就是那位立約者.
耶穌就在當中.
原來那天會來臨.
當他們在等待的時刻裡面.
有一天會來臨.
今天我們都是.
我們要等待住主耶穌.
再一次回來.
第二次再來的時候.
我們同樣需要面對住.
我們在等待的過程裡面.
要做的事.
當你期待有一天.
某些事情會發生.
但發生之前.
你中間要做一些事.
是些什麼呢?.
這裡說有兩件很重要的事情.

$^{681}$但都歸納在同一個意思之下.
就是你要練.
練正.
奧練.
即是說.
我們需要面對練正和奧練.
同時間我們都.
好像要漂洗潔淨一樣.
要整乾淨.
整體的意思就是說.
好像一些金器銀器一樣.
要經歷火去燒.
經歷這些火去燃燒之後.
就能夠將裡面的雜質.
能夠去清除.
又好像我們將一些的衣服.
經過一些漂染的過程.
去洗淨的過程.
就將一些污糟去除去.
所以原來.
當我們在等待拯救的日子之先.
我們自己需要去練正自己.
去漂淨自己.
而這個過程裡面.
整乾淨.
是一件不容易的過程.
所以經文就說.
他來的日子誰能當得起呢?.
他顯現的時候誰能立得住呢?.
誰能夠準備得好.
去面對這個過程呢?.
因為我們要避火去經煉.
我們要不斷去漂染.
去洗乾淨我們所有的不潔.
我們怎樣去等待著.
神拯救的日子呢?.
我們要預備好我們自己.
成為一個怎樣的器皿呢?.
各位姐妹.
原來我們今天在這個地上裡面.

$^{721}$有時生活裡面.
很多的挑戰.
都是為了去熬練我們.
去練成我們.
將我們提升.
本來我們是雜質的金.
你也是金.
不過你是雜質的金.
但經過這個火燒了之後.
你又變成純淨的金.
你可能會問.
那成為純淨的金又如何呢?.
是有分別的.
大家讀經計劃.
在讀五經.
其中五經裡面.
說的一個畫面就是見聖典.
不知大家有沒有見回眸.
有沒有想過回眸.
其實回眸的鋪排.
建設是很有層次的.
因為裡面有很多意思.
我不能夠一一去說.
不過其實你留意一下.
回眸裡面所用到的金屬.
有哪幾種金屬?.
金銀銅.
有沒有鐵?.
沒有鐵的.
金銀銅沒有鐵的.
為什麼呢?.
因為那時候.
未是一個鐵器的時代.
後來菲律賓人就有很多鐵器.
所以你見到.
菲律賓人的武器是很厲害的.
因為他們擁有很多鐵的物資.
但在出埃及的時候.
其中他們主要的金屬就是金銀銅.
金銀銅哪種比較貴?.

$^{761}$金比較貴.
你送金給我就可以了.
銀紙那些紙的銀不要.
要金就可以了.
為什麼呢?.
因為原來在回眸設計裡面.
你認真看一看.
其實是分了等級的.
你想想在回眸裡面的燈台.
叫什麼燈台?.
金燈台.
那個藥櫃用什麼包的?.
用金包的.
在回眸裡面.
原來這裡說很多器皿.
獻給神或擺設的物品.
都是用金的.
越近致勝所.
回眸最核心的地方.
原來都是用金.
才可以放進致勝所裡面.
而銀呢.
這裡說銀就是那些外面的柱.
根基的位置.
用了銀的器皿.
就是銀的物質.
而還沒去到回眸之先.
有個什麼?.
有個祭壇.
什麼祭壇?.
我讀了這麼久的五經.
就是銅祭壇.
所以你見得到.
原來進入回眸的一層一層的過程裡面.
是先去經過銅.
再去到銀.
再去到金.
再去到上帝的致勝所面前.
所以以色列人就將他們不同的金屬.
按照這些次序擺設獻給上帝.

$^{801}$所以今天.
我們都要成為一個被上帝奧煉的生命.
我們有很多生命的雜質.
我們要經歷一路提煉一路提升.
可能我們真的雜得不順.
慢慢我們要堅壯如銅.
慢慢我們要成為銀.
更重要的是.
我們好像金一樣.
成為金.
煉成了一些雜質的金.
可以去到神的面前.
去獻給神.
今天你想獻給神的時候.
你想想如果你的生命裡面有很多污穢.
很多雜質.
很多邪惡.
你怎樣能夠去到回眸裡面呢?.
不能夠的.
但正正好我們能夠.
讓我們自己去成為金的話.
我們就能夠去到神的面前.
而這個過程是需要奧煉.
是需要漂靜.
即是一個提煉一個磨煉的過程.
兄弟姐妹你們想想.
今天你們信了耶穌那麼久.
有沒有想過.
怎樣去磨煉自己的生命.
變得更加純正呢?.
我們很期望.
回到教會安安靜靜.
如常每一個星期.
好像例行公事.
很穩定地.
星期日某個鐘數就是那個鐘數.
做完就完.
不是的.
是在我們每一天裡面.
經歷那種奧煉.

$^{841}$提煉.
鍛煉.
才能夠讓我們提升成為精金.
上帝讓我們的生命裡面.
都是不斷要作出一些挑戰.
一些磨煉.
幫我們生命去提升.
剛才我們說第一個重點裡面.
就是說.
我們怎樣面對.
我們的懷疑.
對不對?.
剛才我們說了懷疑.
很多問號令我們遠離神.
但接下來我要介紹的是.
一個人物.
他叫懷志.
不是懷疑.
他叫懷.
不是懷疑.
他叫懷志.
其實這個懷志.
有一個哥哥.
叫懷義.
懷義.
就是義人的義.
懷義和懷志.
不是懷疑.
懷疑是不喜悅的.
不要質疑上帝.
但我們今天和大家去看一個.
年輕人懷志.
這個懷志有什麼特別呢?.
原來他在小時候.
已經決心要被上帝去使用.
以致他被上帝不斷去改變他的生命.
提升他的生命.
大家知道他是誰嗎?.
有沒有見過這個人?.
沒有就算了.

$^{881}$那我繼續說吧.
這個懷志是怎樣呢?.
他說的是.
其中一些他生命的經歷.
就是他說他高中一年級的時候.
他立志去追求一個聖潔的生活.
是不是?.
而當時發生什麼事呢?.
發生什麼事呢?.
他有一個挑戰.
原來當時他.
即是他的年代.
那個髮型.
那時候流行什麼髮型呢?.
光頭.
這個時代.
你見到每個人都有頭髮.
光頭是被迫的.
原來他們喜歡.
頭髮是分界的.
原來那個時候.
頭髮分界就是那個朝代.
那個時代的尖端.
好看的就是.
頭髮會分界.
他也是這樣.
後來他發覺.
在神教裡面發覺.
原來他.
那麼注重頭髮分界.
是一種貪慕虛榮.
因為為什麼呢?.
每個人都是這樣.
做年輕人.
追得上時尚的年輕人.
就要這樣.
即使他也覺得.
我要過一個聖潔生活.
我不應該追求分界.
結果怎樣呢?.

$^{921}$他怎樣去.
在神的面前.
立志成為聖潔呢?.
和參孫相反.
參孫留長髮.
他又怎樣呢?.
他就剃光頭髮.
因為他說.
我沒有頭髮.
我就在神面前.
不會再想貪慕的虛榮.
就是分界.
當時時代分界的虛榮.
雖然今天我們覺得.
好像很奇怪.
分不分界.
但當時引誘他的就是這樣.
就好像我們早一輩.
有些教會都說.
屬神的基督徒.
就不會進戲院.
有這樣的時代.
因為當時戲院.
有很多不同的戲.
有黑暗.
但重點是.
他在神面前.
有一種體會.
他要不要跟著這個世界.
他就是這樣去立志.
要追求聖潔.
就是這樣剃光頭髮.
這是其中一件他發生的事.
他就是這樣立志追求神.
而他怎樣呢?.
後來神不斷去塗造他.
去誣衊他.
其中有一次.
他很年輕的時候.
他已經開始講道.

$^{961}$你知不知道他幾多歲開始.
在一些教會聚會裡面.
開始講道?.
這個懷智.
17歲.
19歲.
19歲就是大學生.
所以我們教會.
有些學習講道的.
那些人不說幾多歲.
但你已經遲了.
因為早些訓練.
原來他在年輕十多歲的時候.
原來已經被牧師委派去訓練他.
培育他成為將來神使用的人.
他就是這樣.
由19歲開始.
第一次講道.
一路一路去講.
講了很多年.
所以今天我們要.
早些訓練我們的年青人.
因為越早訓練.
他越能經歷多些生命的惡念.
接著到他讀大學的時候.
他讀的大學叫西北大學.
他其中一個大學的活動是什麼呢?.
就是他走60里路.
走60里路去西北醫學院.
是一間醫學院.
走60里路去幫團契聚會.
即是帶一個團契周會.
走60里路.
帶團契周會.
所以如果邀請你帶團契周會.
很短路程.
你開車一會就到了.
但他走了很多個小時.
因為60里路.
很多個小時.

$^{1001}$去到就是為了負責一個周會.
到第二天才走60里路.
回到大學宿舍.
就是這樣.
他一路一路去接受不同的惡念.
甚至乎他親眼見到.
他爸爸離開世界的時候.
是什麼時候離開.
成為他生命中一個.
歷歷在目的一個體驗呢?.
就是他爸爸原來是一個牧師.
他爸爸是一個牧師.
他爸爸有一天在主日崇拜的時候.
他爸爸講道.
可能覺得太激烈.
腦充血.
結果在台上暈倒.
結果五六天後就離開了這個世界.
我有時候都想.
如果講道都不可以太激動.
所以你看我今天節奏比較慢.
其實不可以太激動.
事實上我也認識另外一個牧師.
他也是在台上講道的時候.
離開這個世界.
但當中他們離開世界的時候.
其中留下了.
離開世界的時候.
有讀到一段經文.
就是啥呢?.
就是啟示錄十四章十三節.
他說我聽見.
從天上有聲音說.
你要寫下.
從今以後在主裡死去的人有福了.
聖靈說.
是的.
他們要從自己的奴婦中得安息.
因為工作的成果永隨著他們.
見證到爸爸原來就是為主.

$^{1041}$在台上離開這個世界.
他很小的時候就是這樣.
經歷著生命的熬煉.
一路熬煉.
一路將他的雜質清除.
以致他慢慢越來越調教到.
自己可以去專心去事奉主.
所以他用了六十年的時間去服侍主.
你認識他嗎?.
你知道他是誰嗎?.
給點提示吧.
他姓藤的.
是不是?.
他就是藤近分牧師.
原來他小時候的名字叫懷智.
藤近分牧師.
如果你在香港或世界各地.
你都認識他.
他參與了很多的福音工作.
他起碼用了六十年的時間.
用他的生命去做了一個見證神的時刻.
牧羊教會.
特別宣道會北角堂.
整個宣道會區聯會.
還有見道神學院.
全部都和他有一定的關係.
見道神學院的院長.
他一生用了很多的時間去為主宣講.
原來他今天成為一個眼中的大牧.
一個出名的牧者.
是因為他從小就開始用生命去奧戀自己.
已經決定要去跟隨主.
這就是生命的奧戀.
經歷了這種的奧戀之後.
他能夠成為上帝所期待的.
因為他們就憑公義獻公物給耶和華.
猶太和耶路撒冷所獻的公物.
必蒙耶和華月立.
因為當你經歷了這些奧戀之後.
你這個祭物是上帝月立的.

$^{1081}$你成為一個純淨的金.
上帝會安置你在一個至聖所裡.
或者聖所裡.
很不一樣.
所以馬拉基先知對這群百姓的回覆就是.
沒錯,那一天還沒到.
不過神應許的拯救必然會發生.
我們今天在等待的過程裡.
我們需要不斷去提升奧戀自己.
以致我們可以成為一個上帝月立的一個供物.
一個禮物給神.
弟兄姊妹,一份美麗的禮物不是在於包裝.
包到自己外面有多美.
而是你實質的價值.
真的活出那種價值.
今天我們怎樣重視我們的生命.
我們有沒有那份價值呢.
這是我們值得去思想的.
第三,失望之日.
面對的是審判.
第五節說,萬軍的耶和華說.
我必臨近你們施行審判.
速速作見證.
警戒那些行邪術的,犯奸淫的,.
起假誓的,剝削故宮工錢的,.
欺壓孤兒寡婦的,屈往寄居者的.
我不敬畏我的人.
這裡說的失望之日是什麼呢.
是說有些人很期待著上帝有一天會來到.
有些人很期待著彌賽亞有一天會來到.
但是這一天來到對你是不是好事呢.
就要視乎你是一個什麼樣的人.
所以當時上帝對這群百姓說.
他來到的時候.
他要來做一個很重要的工作.
就是施行審判.
如果你是沒有資格的.
你是一個罪人的話.
你神來到那天.
你是面對審判那天.

$^{1121}$不是歡天喜地那天.
如果你成為一個很純正的金獻給上帝.
這一天當然是歡天喜地.
但是如果你是一個被上帝審判的人.
這一天來到對你來說是一件不好過的事.
這裡正正好就是說.
神就說審判的人出現.
誰是審判的人呢.
是我必臨近你們施行審判.
就是那個萬軍的耶和華這一天來到的時候.
他會親自來到施行審判.
就是說那個法官就是神自己.
一件案子除了有法官之外還有證人.
誰是證人.
他說我必速速作見證.
神說證人都是我.
因為我經歷了歷世歷代.
我見證了你們每一個人所發生的一切的事.
所以原來在我們終極的日子裡面.
我們要面對上帝親自來到審判.
而且他也用他見過我們所做的一切.
包括我們以為別人看不到的一切.
他都見到的時候.
這些就成為我們被上帝去見證的一些行為表現.
而我們面對著這個審判的結果.
在這個審判裡面.
神特別說有些在這個以色列地裡面.
一群人他們就是這樣活在這個回歸的國土裡面.
他們回歸的國土裡面有些什麼問題.
為什麼一直成為神要責備他們的原因呢.
這裡說有七宗的罪.
七宗的罪包括了什麼呢.
行邪術.
就是說用一些靈異的方法.
一些交軌的方法去希望讓自己有一些好的結果.
所以今天有些人會養鬼仔.
不是認識美國人的那些鬼仔.
是靈異的那些鬼仔.
所以有些人一方面敬拜上帝.
一方面跟邪靈交軌有往來.

$^{1161}$所以要謹慎.
假如你家裡有個敬拜神.
又放了一些偶像在家裡.
小心了.
第二就是犯奸淫.
代表了一些婚姻的不忠.
就是在當時的以色列地裡面.
他們在做的問題.
在家庭裡面的婚姻裡面.
竟然不能夠活出一份的約.
就好像我們之前在講義章裡面.
在講婚姻的約一樣.
婚姻是兩夫妻的盟約.
但如果當你犯奸淫的時候.
就破壞盟約.
這個是上帝不喜悅的.
起假誓.
你見得到嗎.
起假誓代表什麼.
就是藐視法庭一種表現.
藐視一種你在別人面前.
忠誠的可信的一種表現.
你將你自己那個可信的信靠的程度.
由一十分可能變成零分.
還有什麼呢.
剝削故宮工錢的.
代表什麼.
你藉著你是一個有錢人.
就去在錢財上面壓榨一些窮人.
上次我們沒有講到的.
阿摩斯書其中一樣在講什麼.
就是用一對拖鞋的價錢.
去買一個窮人回來.
你想一下拖鞋值多少錢.
但是他們就壓榨那些窮人.
在他們不行的時候.
最困苦的時候.
就用這個價錢.
就去買了他回來做他的奴隸.
一種財富的壓榨.

$^{1201}$還有什麼.
欺壓孤兒寡婦.
對社會不公.
就是說在社會裡面.
去欺負一些沒有能力保護自己的人.
帶來一種這樣不公的狀況.
屈往寄居者.
就是什麼.
一種民族的歧視.
就是說和你同胞的.
你就好聲好氣.
但當和你是不同胞.
不同族不同種的人.
你就盡量去什麼.
去屈往他們.
去屈他們.
去對他們歧視.
最後就是什麼.
就是不敬畏神.
對神不敬.
所以馬拉基就講一個狀況.
就是當我們在這個社會裡面.
在不同的範疇裡面.
我們做一些得罪上帝的事情.
其實上帝再回來的日子.
其實你就慘了.
因為上帝再回來的日子.
正正是我們面對審判的日子.
上帝就在高高在上的寶座裡面.
審判我們一群犯不同罪的人.
這真是值得我們去警醒.
特別我們更加要去警惕.
我們不要因為現在這些人.
做一些上帝不喜歡的事情.
我們就靠邊站靠過去那邊.
站過去同一陣線.
不是啊.
我們要持守我們的信仰.
持守我們的陣線.
特別這裡最近有很多不同的.

$^{1241}$一些新聞一些報導.
包括了我們都要面對的.
一個社會的問題.
這裡說最近加拿大.
應該如果沒有記錯.
已經通過了取消一個.
宗教豁免的仇恨言論條款.
不知道你們知不知道這是什麼.
簡單來說就是.
譬如我在教會裡面.
我是宗教員我就宣講聖經.
就說同性戀是聖經不容許的.
按著聖經說.
我講的是聖經裡面說的東西.
不是說我不同意.
我只是說聖經有這樣說.
不同意同性戀這樣說.
這種情況應該在宗教上.
在宗教自由之下.
我可以說的一些說話.
但是加拿大已經取消了.
這種宗教豁免.
無論你是什麼宗教人士都好.
只要你就算掛著.
你是聖經裡面這樣說的.
而這些言行是引致了.
一些仇恨的言論.
你都是要坐牢.
就是說.
剛剛我說了那些東西.
請幫我刪除.
被網上聽到.
我就死定了.
如果我在加拿大.
要坐牢的了.
甚至我看回資料.
現在其實英國也是這樣做.
英國加拿大都是這樣.
那你想想美國.
可能都不願意.

$^{1281}$可能不願意.
就是說一種.
傾左的一種狀況.
為了要表示大愛.
表示尊重.
你的信仰裡面.
就算某些行為是不對的.
你都不可以開聲說不對.
就算你是有一個宗教的原因.
是不可以的.
反過來.
他們說他們大愛.
他們就可以這樣說.
還有另外一單新聞.
最近世界盃.
世界盃也有一單新聞.
世界盃裡面.
你知道原來世界盃裡面.
在足協裡面.
說明所有的足球員.
都不可以公開.
表明你的信仰.
什麼意思呢.
就是說.
有時候有些球員.
進了球之後.
脫衣服.
撕開衣服.
有些信耶穌的人就寫.
I love Jesus.
會這樣的.
顯示你的信仰的狀況.
有時候他們會把衣服展示出來.
就像說了一集.
不可以.
總之在任何賽事裡面.
在任何球賽裡面的人.
都不可以去表明.
你有任何的信仰.
不過呢.

$^{1321}$這裡說.
就是什麼呢.
但你是可以掛著綠色彩虹旗的.
所以剛剛最近.
其中在.
世界盃裡面.
在美國其中一個場地裡面.
他們.
那個.
比賽裡面.
就掛了綠色旗.
還有說了四個字.
什麼什麼球賽.
就是說那場比賽裡面.
是要尊重同性人士的.
我再發給大家看.
就是說.
這些是可以的.
但你可以告訴別人.
我信耶穌.
耶穌萬歲.
主佑榮耀.
這些全部.
都不可以在球賽裡面.
這個社會裡面.
就是這樣用雙重的標準.
去判斷我們信神的人.
反而不可以告訴別人信神.
不信神的人.
他們有些問題.
他們繼續可以說出的問題.
這些是會令我們面對著.
更加多的疑惑.
就好像.
以色列地裡面那群的百姓一樣.
但是神說.
任何在神的裡面犯了不同.
上帝不喜悅的事情的時候.
的人.
他們都要面對他們自己.

$^{1361}$不同的罪.
在那一天.
神就會回來去審判這個世代.
所以我們.
不需要自己做審判官.
但你卻又要堅持你的信仰.
在這個社會裡面.
面對著不同的奧煉.
不同的困難.
不同的煉證的時候.
我們仍然要提升自己.
成為今天一樣.
可以獻給上帝.
所以今天和大家思考的就是.
今天.
確實可能我們是一個未報之日.
見不到一些報應.
在某些地方臨到.
但是我們要知道.
有一天.
要出現.
那一天.
可以是一個拯救的日子.
當你預備好的時候.
但當你.
如果想著.
既然.
這些的錯誤.
這些的惡報都沒有報的話.
你接著成為惡人的話.
到那一天.
審判來到.
你就成為了.
失望之日.
最後介紹一本書給大家.
不知道大家有沒有看過.
我自己有時候經常看完再翻看.
這本書叫.
與成功有約.
The Seven Habits of Highly Effective People.

$^{1401}$這本書裡面.
是一本很出名的書.
其中它說.
怎樣能夠讓你成功.
成為一個有效的人.
其中說了七個的.
七個的習慣.
其中.
第二個習慣.
就是說怎樣成功.
就是以終為始.
以終為始.
即是說.
你期待著明天.
你是面對審判.
還是面對上帝拯救.
當你決定了.
你選擇哪個終結的時候.
你就今天.
你就要怎樣去生活.
來令到這件事發生.
令到這件事發生.
以終為始.
我們不是看現在多惡劣.
而是看終極.
我們究竟想蒙神的讚賞.
懸納我們成為一個可喜悅的祭物.
還是我們成為一個犯人.
在神的面前面對審判.
今天的你.
可以決定明天的你.
求主幫助我們.
今天要立志.
立志最後還有幾句說話.
我們一起讀這幾句說話.
好 開始.
思想決定行動.
行動決定習慣.
習慣決定品德.
品德決定命運.

$^{1441}$究竟在很多年之後.
神跟你說.
你都有今天了.
究竟是什麼日子呢.
求主幫助我們.
我們一起祈禱.
因為你讓馬拉基先知.
在他們面對困惑的世代裡.
作出審判的訊息.
一個挑戰的訊息.
以致我們真的能夠知道.
我們等待彌撞亞的來臨的時候.
究竟對我們是真的好還是壞.
要視乎我們怎樣活出我們的生命.
求主你幫助我們教會.
每一個聆聽你話語的人.
我們都有份心志.
有份悖念的心志.
我們要成為精金獻給上帝.
或許我們今天.
我們習慣了很悠閒.
很自在的生活.
但求主你幫助我們.
更加知道我們要有一種韌力.
去面對上帝你自己.
以致我們成為一個最好的祭物.
去獻給你.
幫助我們在這個世代裡.
堅守我們的信仰.
聽我們祈禱奉耶穌的名字而求.
阿們.
\newpage



\section{以弗所書 5:1-20}
\label{sec:032294}
\textbf{活出你是誰 (華)}~~~~孫雷雷牧師
\newline
\newline
連結: \href{http://www.church.com.hk/acms/content.asp?site=cdc\&op=show\&type=product\&code=032294\&layout=sermon}{\texttt{032294}} ~~~~ 語音日期: 2026-06-14
\newline
\newline
\hyperref[sec:032293]{\small{< < < PREV SERMON < < <}}
~
\hyperref[sec:index_chronic]{\small{[返順時目]}}
~
\hyperref[sec:032295]{\small{> > > NEXT SERMON > > >}}
\newline
\newline
以弗所書 5:1-20
\newline
\begin{longtable}{cl}
\hline
\hline
章節 & 經文 (和合本修訂版)\\
\hline
5:1 & \begin{tabularx}{0.7\textwidth}{X} 所以，作為蒙慈愛的兒女，你們該效法神。 \end{tabularx} \\ \\ \relax
5:2 & \begin{tabularx}{0.7\textwidth}{X} 要憑愛心行事，正如基督愛我們，為我們捨了自己，當作馨香的供物和祭物獻給神。 \end{tabularx} \\ \\ \relax
5:3 & \begin{tabularx}{0.7\textwidth}{X} 至於淫亂和一切污穢，或是貪婪，在你們中間連提都不可，這才合乎聖徒的體統。 \end{tabularx} \\ \\ \relax
5:4 & \begin{tabularx}{0.7\textwidth}{X} 淫詞、妄語和粗俗的俏皮話都不合宜；總要說感謝的話。 \end{tabularx} \\ \\ \relax
5:5 & \begin{tabularx}{0.7\textwidth}{X} 要確實知道，無論是淫亂的，是污穢的，是貪心的（貪心的就是拜偶像的），在基督和神的國裡都得不到基業。 \end{tabularx} \\ \\ \relax
5:6 & \begin{tabularx}{0.7\textwidth}{X} 不要被人虛浮的話欺騙了，因這些事，神的憤怒必臨到那些悖逆的人。 \end{tabularx} \\ \\ \relax
5:7 & \begin{tabularx}{0.7\textwidth}{X} 所以，不要與他們同夥。 \end{tabularx} \\ \\ \relax
5:8 & \begin{tabularx}{0.7\textwidth}{X} 從前你們是暗昧的，但如今在主裡面是光明的，行事為人要像光明的子女— \end{tabularx} \\ \\ \relax
5:9 & \begin{tabularx}{0.7\textwidth}{X} 光明所結的果子就是一切的良善、公義、誠實。 \end{tabularx} \\ \\ \relax
5:10 & \begin{tabularx}{0.7\textwidth}{X} 總要察驗甚麼是主所喜悅的事。 \end{tabularx} \\ \\ \relax
5:11 & \begin{tabularx}{0.7\textwidth}{X} 那暗昧無益的事，不可參與，倒要把這種事揭發出來。 \end{tabularx} \\ \\ \relax
5:12 & \begin{tabularx}{0.7\textwidth}{X} 因為，他們暗中所做的，就是連提起來都是可恥的。 \end{tabularx} \\ \\ \relax
5:13 & \begin{tabularx}{0.7\textwidth}{X} 凡被光所照明的都顯露出來， \end{tabularx} \\ \\ \relax
5:14 & \begin{tabularx}{0.7\textwidth}{X} 因為使一切顯露出來的就是光。所以有話說：「你這睡著的人醒過來吧！要從死人中復活，基督要光照你了。」 \end{tabularx} \\ \\ \relax
5:15 & \begin{tabularx}{0.7\textwidth}{X} 你們要謹慎行事，不要像無知的人，要像智慧的人。 \end{tabularx} \\ \\ \relax
5:16 & \begin{tabularx}{0.7\textwidth}{X} 要把握時機，因為現今的世代邪惡。 \end{tabularx} \\ \\ \relax
5:17 & \begin{tabularx}{0.7\textwidth}{X} 不要作糊塗人，要明白主的旨意如何。 \end{tabularx} \\ \\ \relax
5:18 & \begin{tabularx}{0.7\textwidth}{X} 不要醉酒，酒能使人放蕩；要被聖靈充滿。 \end{tabularx} \\ \\ \relax
5:19 & \begin{tabularx}{0.7\textwidth}{X} 要用詩篇、讚美詩、靈歌彼此對說，口唱心和地讚美主。 \end{tabularx} \\ \\ \relax
5:20 & \begin{tabularx}{0.7\textwidth}{X} 凡事要奉我們主耶穌基督的名常常感謝父神。 \end{tabularx} \\ \\
[1ex]
\hline
\hline
\end{longtable}
$^{1}$粵語堂的弟兄姊妹們 屬內平安.
很感恩今天有機會來到大家當中.
和你們分享神的話.
我們在開始的時候一同做一個禱告.
天父.
當我們打開你話語的時候.
求你親自的向我們的心說話.
光照我們的心思意念.
引導我們與你同行.
禱告奉主耶穌的名.
阿們.
有一位.
有一位父親.
在他女兒出嫁的前一天的晚上.
然後送給他三面鏡子對他說.
第一面鏡子是讓你記住你是誰.
你是我深深愛著的女兒.
第二面鏡子是要讓你看清楚周圍是什麼.
哪些是光.
哪些是黑暗.
第三面鏡子讓你看清楚你要走的路.
不是憑著感覺走.
而是憑著智慧走.
今天我們所讀的以弗所書的這段經文.
在原文當中正是有三個詞.
應該是一個詞.
三次的出現就是行走.
也正像是天父賜給神兒女的三面鏡子.
五章二節就告訴我們說要在愛中行走.
五章八節讓我們在光中行走.
五章十五節讓我們謹慎行事.
當向智慧人在智慧中行走.
所以三次的行走.
三面的鏡子就構成了今天經文的骨架.
那麼接下來就讓我們來看一看這三面鏡子.
首先第一面鏡子.
保羅說告訴我們要在愛中行走.
我們應當要效法神.
好像那蒙慈愛的兒女一樣.
保羅在告訴我們.

$^{41}$我們是神的兒女.
所以我們就應該要像我們的父親.
就好像是一個孩子.
從小耳濡目染.
自然而然的就會學到他父親的走路的樣子.
說話的腔調.
做事情的那些風格方式.
還有性情.
這是很自然的事情.
那我們呢.
我們是蒙神所愛的兒女.
我們也應當要活出天父上帝的樣式來.
那麼神是什麼樣式呢.
第二節保羅就說.
我們也要憑愛心行事.
正如基督愛我們.
為我們舍了自己.
當做心相的供物和祭物獻於神.
基督的愛.
就是這樣一個捨棄自己的愛.
是給出去的犧牲的愛.
我們常常把愛看得太表面.
我們會說我愛喫這道菜.
我愛看這部電影.
甚至我們覺得愛一個人的時候.
也是更多的就是喜歡那個人給我們的感覺而已.
但這種愛有一個特點.
都是向內的.
就是向著自己的.
都是圍著自己轉的.
但是基督的愛卻是向外的.
他愛我們.
不是因為我們讓他高興.
讓他快樂.
而是在我們還做罪人的時候.
他就為我們死了.
這才叫捨己的愛.
保羅說我們要來效法他的愛.
為人付出.
犧牲捨己的愛.

$^{81}$本身就是我們可以獻給神的祭.
正好像基督他捨棄了自己.
獻給神做新香的供物一樣.
那我們也同樣要去反省.
在我們的日常生活中.
在我們的家庭中.
婚姻關係中.
在教會和弟兄姊妹的關係中.
我們的愛更像是哪一種的.
我們是更多向外給出去的愛.
還是更多向內的索取的.
而接著保羅他的筆鋒一轉.
就從這一個至高的愛.
直接跳到了愛的反面.
他說至於淫亂並一切的污穢.
或是貪婪在你們中間.
連提都不可.
他說的這些淫亂貪婪污穢.
其實都有一個共同點.
就是都是向內的.
都是索取的.
性方面的犯罪.
其實就是貪圖別人的身體.
貪婪就是貪圖別人的財物.
所以保羅把這些都並列的放在一起.
因為他們都是個同一類型的東西.
都是那種自我放縱的.
與基督那種捨棄的愛.
是完全不同的方向.
他還說.
淫詞妄語和嬉笑的話都不相宜.
那些粗俗的笑話.
就是把神的美好的恩賜.
給當作那種消遣品.
那性 身體 我們的生命.
都是神賜給我們的寶貴的禮物.
是鄭重的 是美好的.
但是當我們用那些粗俗的笑話.
來消遣他們的時候.
實際上就是把這一個貴重的東西.

$^{121}$給踩在腳底下來開玩笑.
所以保羅提醒我們感恩.
感恩就是讓我們去承認.
這些事物都是神美好的恩賜.
我們需要鄭重的捧在手心裡面.
來向神獻上感恩.
而保羅在第五節也說.
有貪心的就與拜偶像的一樣.
那麼偶像其實就是在我們心裡面.
那一個佔據了本該神所佔據的位置的.
所有的事物.
我們可以是金錢.
可以是地位.
也可以是別人的認同.
別人的認可.
也可以是我們對別人的掌控.
又或是我們戒不掉的某一些的隱.
任何一種東西.
都可能成為我們生命當中的偶像.
有一個富商.
在他臨終前把他的兒子.
兩個兒子叫到床前.
他對他大兒子說.
我把我所有的財產都給你.
對小兒子說.
我把我的名字留給你.
大兒子拿了財產.
不到十年的時間.
揮霍一空.
一無所有.
小兒子沒有拿到一分錢.
但是他繼承了他父親的名聲.
誠信 慷慨.
他最後在商界裡面.
他的成就比他的父親還要大.
保羅在這裡就說.
實際上就是同樣的道理.
我們是神的兒女.
我們繼承的是神的名字.
是神的性情.

$^{161}$是神的那種自我犧牲的捨己的愛.
這纔是神給我們的最寶貴的家業.
所以當我們來看著第一面鏡子的時候.
我們也需要來反思.
我們的生命是否有真的活出了神的性情呢.
我們的愛更多是向外的還是向內的呢.
在我們的所有的關係當中.
我們更多是在給予 付出還是在索取呢.
而接下來我們來看著第二面鏡子.
就是要在光中行走.
保羅說從前你們是闇昧的.
但如今在主裡面你們是光明的.
行事為人就當像光明的子女.
我們可以注意保羅在這裡他是怎麼說的.
他沒有說你們從前在黑暗裡.
他是說你們從前就是黑暗.
他也沒有說你們現在在光明裡.
他說的是你們現在就是光明.
這就不是一個處境的改變.
而是一個我們存在的本質的改變.
我們不是說從這個黑暗的地方搬到了光明的地方.
而是我們這個人 整個人被神徹底的重新造過了.
我們的本質已經被神改變了.
這對我們的屬靈生活來說也有很大的很好的含義.
就是當我們去面對試探.
當我們面對誘惑.
當我們感到軟弱.
或是感覺到自己好像什麼都沒有做好做不到的時候.
我們需要提醒自己的.
不是說我要更加的努力.
靠著我自己的努力做一個好人.
而是提醒我們自己.
我是光明之子.
這是我的身份.
身份決定了我的行為.
而不是行為決定了身份.
美國的著名作家馬克·吐溫.
他曾經寫過一個故事.
叫做《乞丐王子》.
講的就是說一個真正的王子.

$^{201}$因為他換上了乞丐的衣服.
換了處境.
被所有人當作乞丐來對待.
最後連他自己都開始動搖了.
我真的是王子嗎.
但是他的身份從來都沒有變過.
他的衣服變了.
環境變了.
別人對他的看法變了.
但是他是誰這件事情.
從始至終沒有變過.
今天作為基督徒.
作為天國的子民.
天父上帝的兒女.
我們生活在這個世界上.
這個世界其實用各種不同的方式和聲音.
來PUA我們.
就是想要來操控我們 影響我們.
讓我們去忘記我們自己的身份.
保羅就是在提醒我們一件事情.
我們的身份是改變不了的.
我們在基督裡面就是光明的.
感覺不能定義你.
處境也不能定義你.
只有神的話語能夠定義你.
光明之子的身份是不會改變的.
而光明之子所要結的果子.
就是一切的良善 公義 誠實.
這三樣東西就是光明之子的生命的特徵.
然後第十節保羅說.
要查驗何為主所喜悅的事.
這個查驗不是我們定期去做一個.
好像身體檢查.
我們也定期地做一個屬靈生命的檢測.
而是說一種日常的生活方式.
就是在我們的生活當中.
每一個處境面前.
每一個選擇面前.
我們都不斷地問.
這件事情是光明的果子嗎.

$^{241}$這件事情是否是曖昧無意的事.
這件事主喜不喜悅呢.
同時保羅還說.
那曖昧無意的事不要與人同行.
倒要責備行這些事的人.
就是不要去參與那些黑暗的事情.
而不參與黑暗.
不等於說我們就要把自己.
關在一個基督教的圈子裡.
不和這個世界有任何的接觸.
保羅的意思不是這樣.
光要真正地成為光.
就是要發出光來.
就是要照到黑暗裡面去.
但是我們知道.
這樣做其實說起來容易.
做起來難.
因為在我們的真實生活當中遇到的處境.
是很複雜的.
往往可沒有那麼簡單.
比如在職場裡面.
幾個人開始聚在一起.
說某一個同事的閒話.
然後開始營造出一種.
我們都共同很討厭一個人的氛圍.
那你會如何顯得.
又不像一個道德衛士那樣去說教別人.
巧妙地轉移話題.
或者是溫和地去提醒在場的人.
為那個不在場的人說一句公道話呢.
當你拒絕這種有毒的社交綁定的時候.
其實你也要做好心理的準備.
會可能會被邊緣化.
同樣在我們的這些朋友圈子裡面.
當你知道了某些人.
他們正在做某件事情.
是不符合聖經教導的.
那你一個人要站出來的時候.
也是要付代價的.
所以在這些處境裡面.

$^{281}$不參與就是要付代價的.
在我們的現實生活當中.
這種做光.
往往不是那種轟轟烈烈的英雄主義.
而是在這種充滿著人情世故.
灰色地帶.
還有那些利益的糾葛的日常瑣碎當中.
去經歷我們內心掙扎之後.
所做出的艱難的選擇.
但是同時我們也要小心.
這種不參與的本身.
還有可能走向另一個極端.
那就是冷漠論斷高高在上.
用一種令人窒息的方式來劃清界限.
把這種我不參與.
變成了我比你高級.
滿臉的寫著就是宗教的優越感.
這種自以為是聖潔的.
這種自以為是的聖潔.
其實就和法利賽人相差不遠了.
所以我們如何能夠拒絕在不參與罪的同時.
又不變得那種道貌岸然反應過度論斷人呢.
那我們就需要回到.
保羅在已夫所說四章二節給我們奠定的根基.
那就是謙卑 溫柔 忍耐 用愛心互相寬容.
這四樣東西不是說我們拒絕參與罪之後.
用來做的某種的裝飾.
而是我們拒絕參與罪的這種方式的本身.
我們可以拒絕參與罪的和錯誤的行為.
但是同時也可以真誠地去愛那些陷在其中的人.
我們也可以用愛心去說真理去行真理.
同時讓對方感受到的是愛而不是論斷.
這兩件事情不是矛盾的.
我們要同時做到需要的.
就正是真正的謙卑 溫柔 忍耐和愛.
而保羅也讓我們看到.
聖潔並不等於自義.
區別也並不等於隔離.
光它真正要成為光.
發揮光的功用.

$^{321}$那最終的問題還是要回到我們的身份.
我們是誰.
誰在定義我們.
光明之子.
不應該讓黑暗來定義我們.
十四節保羅就說.
你這睡著的人當醒過來.
從死裡復活.
基督就光照你了.
有的時候我們雖然已經是光明之子了.
但是卻活得好像還在睡覺一樣.
所以保羅呼籲我們說.
是時候該醒過來去發出光來了.
所以第二面鏡子提醒我們.
我們要記住我們的身份.
我們不是一個掙扎著要靠自己的努力.
做一個好人的這樣的一個罪人.
我們是一個偶爾可能還會犯錯的光明之子.
這兩種身份的認同.
會帶來截然不同的生命狀態.
前者就會讓我們反覆地.
不斷地在我們的罪疚當中來打轉.
而後者讓我們可以真正地.
在恩典裡面站起來.
行在光中.
同時我們也需要去照亮.
而不是躲在黑暗裡面.
我們可以拒絕參與罪.
同時也深切地愛那些陷在罪中的人.
聖潔不等於自義.
區別也不等於隔離.
而最後讓我們來看這第三面鏡子.
第十五節.
保羅說:你們要謹慎行事.
不要像愚昧人.
當像智慧人.
要愛惜光陰.
因為現今的世代邪惡.
愛惜光陰.
原文執意就是要贖回時機.

$^{361}$把時間買回來.
我們每一個人都希望自己做出的每一筆投資.
都是精明的.
而保羅在這裡鼓勵我們.
要去搶購每一個機會.
這不是在說一個時間管理的問題.
而是我們要有的那種屬靈的緊迫感.
根據統計.
全球平均每人每天花在屏幕上的時間是六個小時四十分鐘.
就將近於我們每天清醒的時間的一半.
當然了.
我們現在的工作生活和學習交流都離不開電子設備.
但是我們同樣也要問自己的是.
這一天下來.
你花在屏幕上面的時間有多少是真正有益處的呢.
有多少是真正能夠使你自己也使他人得益處的呢.
時間是我們這一生最不可再生的資源.
我們的錢花光了.
我們可以再賺.
我們的關係破裂了.
我們還可以去修復.
但是時間一旦過去了就永遠回不來了.
保羅說.
現今的世代是邪惡的.
時間不多了.
不要浪費.
在我們的生命當中.
有哪些時間正浪費在那無益的事情上呢.
我們需要靠著神的光照.
需要去靠著神智慧的引導.
去不斷的去求問神.
在這個時刻.
神要我做什麼.
保羅說.
不要做糊塗人.
要明白主的旨意如何.
那麼我們怎麼樣去明白神的旨意呢.
我們讀經.
禱告.
崇拜.

$^{401}$參加各樣的聚會.
團慶.
這些當然都非常重要.
但是保羅的意思不是說.
當我們把這些屬靈的操練都打卡完成了.
就算能夠明白神的旨意了.
並不是.
我們做這一切的屬靈的操練.
最終的目的是什麼呢.
就是要在我們的生活中每一個處境.
在我們生活的每一個處境裡面.
清楚的去辨認出.
神讓我們做的事.
並且照著去做.
接著保羅就來到了這段經文的核心.
在第十八節.
他說.
不要醉酒酒能使人放蕩.
乃要被聖靈充滿.
醉酒和聖靈充滿.
保羅把這兩樣東西給並列的放在一起.
來做一個對比.
那絕對不是偶然.
在以弗索這個城市.
當時也是拜偶像聖行的.
他們當中就有一個偶像.
就是酒神.
很多人對於這個酒神來崇拜.
那麼崇拜的人就相信.
當他喝到大量的酒之後.
進入到這種癲狂.
迷狂的狀態裡.
他才能夠更加的來接近神明.
保羅說.
不是這樣.
你們要的那種和與神同在的那種真實的感覺.
那種超越自我的滿足感.
不是酒精能夠給得了的.
只有聖靈能夠給.
其實也不只是酒了.

$^{441}$我們今天每一個人在我們的生活中.
我們也會用各樣的東西來充滿自己.
我們可能會用工作.
用我們的消費.
用我們的娛樂.
用尋求別人對我們的認可.
或是用對別人的掌控.
這些都可能會充滿你.
但是都帶來的是短暫的滿足.
而且會讓我們感到上癮.
而且最終給我們的感覺會是越來越空.
只有聖靈的充滿.
是那種真實的.
是永恆的.
是持久的.
是會讓我們感到越來越滿的.
被聖靈充滿了.
也不是那種一次性的體驗而已.
而是一種持續性的生命的狀態.
我們持續地向聖靈敞開.
持續地讓他來掌管.
讓聖靈充滿我們的生命.
那保羅說.
被聖靈充滿的生命.
就是當用詩章 誦詞 臨歌 彼此對說.
口唱心和地讚美主.
凡是要奉我們主耶穌基督的名.
被聖靈充滿的生命.
是歌唱 是感恩.
是彼此的建造.
其實.
早期的教會就是一個歌唱的教會.
他們喜樂的根源不是建立在處境好.
而是處境非常糟糕.
受逼迫 被排斥.
他們的感恩是因為他們知道神是誰.
也知道神為他們做了什麼.
所以他們感恩是成了他們生命的一種基調.
不是表演出來的.
也不是建立在這些外在的環境之上的.

$^{481}$而是完全從他們的內心湧流出來的.
凡是常常感恩.
不是有的時候.
也不是環境好 遇到好事的時候.
是凡是常常感恩.
我們也知道這很難.
當你失去工作的時候.
當你失去健康的時候.
當你家裡面出事的時候.
當你禱告了很久.
為某件事情持續地禱告.
沒有得到答案的時候.
感恩這兩個字是很難說出口的.
但保羅所說的感恩.
不是因為一切都好才感恩.
而是因為知道一切都在神的手中才感恩.
這兩種感恩是完全不同的.
前者是建立在環境上面.
會隨著這個環境的變化而消失的.
而後者是建立在神的信實上面.
是任何的環境都撼動不了的.
所以這第三面的鏡子.
就讓我們被提醒.
我們需要認真的來對待我們的時間.
在智慧當中來行走.
在我們的生命當中.
去找出那些白白被我們浪費的時間.
並且去問神.
你要我在這段時間裡做什麼.
不要讓這個世界的東西來充滿我們的心.
而是要讓那永恆的聖靈來充滿我們.
我們要不斷地停下來轉向聖靈.
求聖靈來引導我們的心思意念.
更新我們的心.
而當我們遇到困難.
遇到困境的時候.
我們也是否願意選擇感恩.
還是開始抱怨呢.
你是否有把敬拜帶入到你的日常生活中去呢.
敬拜不只是每週日我們來到教會.

$^{521}$在這崇拜當中發生的.
敬拜是發生在我們日常生活中的每一個時刻的.
在我們出門的時候.
我們開車的時候.
洗碗的時候.
工作的時候.
我們和家人在一起的時候.
任何的時間.
任何的時刻.
都可以成為我們敬拜的時刻.
我們對神的那種感恩.
敬拜和讚美.
是需要成為我們日常生活的背景音樂.
所以讓我們再回到開頭.
那位父親送給他女兒的三面鏡子.
這第一面鏡子就告訴我們.
你是誰.
你是蒙神所愛的兒女.
基督為你捨棄了自己.
這愛不是你贈來的.
而是在你死在過犯罪惡中的時候.
他就已經甘心情願為你付出的.
所以我們要在愛中來行走.
效法我們的主耶穌基督.
效法我們的天父.
用那種向外給出去的捨棄犧牲的愛.
來活出在我們所有的關係裡.
第二面鏡子告訴我們.
你在哪裡.
你曾經就是黑暗.
現在在主裡面.
你就是光明.
你不屬於黑暗.
你屬於光.
所以我們要在光中來行走.
結出那良善公義誠實的果子.
然後用我們的光.
去照亮我們周圍的黑暗.
而第三面鏡子告訴我們.
你要往哪裡走.

$^{561}$現今的世代是邪惡的.
時間是寶貴的.
我們不能憑著感覺.
憑著衝動去行事.
我們需要聖靈的充滿和引導.
我們需要在智慧中行走.
不要被這個世界的東西所充滿.
而要被那永恆的聖靈所充滿.
活出那歌唱感恩的生命.
三面鏡子.
都是來自於同一位天父上帝.
他愛你.
他光照你.
他以智慧引導你.
他也已經在基督裡為你成就了這一切.
現在他也正在邀請你.
作為蒙神所愛的兒女.
去活出這一切吧.
我們一同來禱告.
使我們成為蒙慈愛的兒女.
主啊.
我們承認我們從前是曖昧的.
但如今在主裡面是光明的.
求主幫助我們真正的效法你.
活出犧牲與捨己的愛.
懂得愛惜光陰.
明白什麼纔是你的旨意.
願你的聖靈在我們心中掌權.
無論順境逆境.
都能凡是奉我們主耶穌基督的名.
去活出一個不斷感恩與蒙昭的恩.
相襯的生命.
禱告奉主耶穌基督的名.
阿們.
\newpage



\section{啟示錄 3:7-13}
\label{sec:032295}
\textbf{非拉鐵非教會:軟弱卻又強壯的教會}~~~~林穎豪傳道
\newline
\newline
連結: \href{http://www.church.com.hk/acms/content.asp?site=cdc\&op=show\&type=product\&code=032295\&layout=sermon}{\texttt{032295}} ~~~~ 語音日期: 2026-06-14
\newline
\newline
\hyperref[sec:032294]{\small{< < < PREV SERMON < < <}}
~
\hyperref[sec:index_chronic]{\small{[返順時目]}}
~
\hyperref[sec:032296]{\small{> > > NEXT SERMON > > >}}
\newline
\newline
啟示錄 3:7-13
\newline
\begin{longtable}{cl}
\hline
\hline
章節 & 經文 (和合本修訂版)\\
\hline
3:7 & \begin{tabularx}{0.7\textwidth}{X} 「你要寫信給非拉鐵非教會的使者，說：『那神聖、真實的，拿著大衛的鑰匙，開了就沒有人能關，關了就沒有人能開的這樣說： \end{tabularx} \\ \\ \relax
3:8 & \begin{tabularx}{0.7\textwidth}{X} 我知道你的行為。看哪，我在你面前給你一個敞開的門，是沒有人能關的。我知道你有一點力量，也遵守我的道，沒有否認我的名。 \end{tabularx} \\ \\ \relax
3:9 & \begin{tabularx}{0.7\textwidth}{X} 那屬撒但會堂的，自稱是猶太人，其實不是猶太人，而是說謊話的，我要使他們來到你腳前下拜，使他們知道我已經愛你了。 \end{tabularx} \\ \\ \relax
3:10 & \begin{tabularx}{0.7\textwidth}{X} 因為你遵守了我堅忍的道，我也必在普天下人受試煉的時候保守你免受試煉。 \end{tabularx} \\ \\ \relax
3:11 & \begin{tabularx}{0.7\textwidth}{X} 我必快來，你要持守你所有的，免得人奪去你的冠冕。 \end{tabularx} \\ \\ \relax
3:12 & \begin{tabularx}{0.7\textwidth}{X} 得勝的，我要使他在我神的殿中作柱子，他必不再從那裡出去。我又要把我神的名和我神城的名—從天上我神那裡降下來的新耶路撒冷，和我的新名，都寫在他上面。 \end{tabularx} \\ \\ \relax
3:13 & \begin{tabularx}{0.7\textwidth}{X} 凡有耳朵的都應當聽聖靈向眾教會所說的話。』」 \end{tabularx} \\ \\
[1ex]
\hline
\hline
\end{longtable}
$^{1}$今天是六月十四日.
我有一個很喜歡的歌手.
他寫了一首歌叫做《六月》.
歌詞是這樣說的.
別恨自己生於這悲情世代.
怎麼永垂不老 轉眼亦已不再.
像六月天空都會忽然飛霜.
想天天笑容燦爛.
可惜世情冷酷.
當中喜與怒 樂與哀 未到我揀.
萬種悲歡離合 逐一化為灰燼.
連同白雪花遍地瀰漫.
這個歌詞描述到生命裡有很多生離死別.
悲歡離合.
以至社會或世界發生的事情.
都是不盡人意.
是人所無法去掌握去控制.
甚至人在當中連喜怒哀樂都好像無得揀.
人在一切面前.
彷彿只可以無奈地去接受.
苟然殘存.
昔日有一間教會同樣都是軟弱無力.
面對逼近患難 只能逆來順受.
這間就是今天我們來到六月.
看到耶穌耶穌基督寫信給第六間教會.
菲拉鐵菲.
菲拉鐵菲這個地方的名字很繞口.
真的要練習一下怎麼讀.
菲拉鐵菲.
讀快一點可能讀錯了其他東西.
菲拉鐵菲.
但是當大家聽到它的英文名字.
可能就會比較覺得熟悉一點.
因為菲拉鐵菲的英文名字就是菲律賓.
現在美國賓夕法尼亞州人口最多面積最大的城市.
費城就是以菲拉鐵菲為它的名字.
今天就讓我們藉著這封信.
我們看看耶穌教導軟弱的教會.
這間菲拉鐵菲的教會應該怎樣去生存.
這封信我們今天看是很特別的.

$^{41}$這間放在七間教會裡排名第二的教會.
和排在第二的仕姆那教會.
一樣都是一個名第二一個第二.
這兩間教會是七間教會裡.
唯二耶穌沒有在信裡去責備他們的教會.
不過有一樣東西是很相同的.
就是耶穌寫給教會的這七封信.
都是有一個很相同的格式.
一開始耶穌先一個自我介紹.
然後耶穌就說祂知道.
祂看得見這間教會一些特質,一些情況.
這封寫給菲拉鐵菲教會的信.
主耶穌說祂知道教會有什麼特質,有什麼情況.
主耶穌說祂知道菲拉鐵菲教會有好的行為.
並且略有一點力量.
如果今天有人走到你面前跟你說.
你煮東西也挺好吃的.
你也挺厲害的.
你也有點力的.
按照我們華人的文化.
我們華人謙遜的文化.
凡事不要說得太過的文化.
我想大幾率每個人都走到你面前是想稱讚你的.
不過耶穌不是華人.
按照當時的語境.
按照當時的處境他們說的話.
略有一點力量原來不是一個讚美的詞彙.
耶穌是想很客觀的指出一個現實.
原來菲拉鐵菲教會的力量.
真的只有一點點.
真的很微弱,很有限.
甚至是瀕臨耗盡的狀態.
論人數和事工.
他們比不上推也推拉的教會.
論財力.
他們比不上老弟家的教會.
論牧者和神學家的牌面.
他們比不上已不所的教會.
他們力量單薄.
欺弱,尤施.

$^{81}$勢孤,力弱.
而為何他們那麼軟弱呢?.
他們那麼軟弱不是沒有原因的.
菲拉鐵菲教會的軟弱主要來自兩個主因.
第一個是因為人為的敵視和逼迫.
第九節說到他們要面對一群猶太人.
這群猶太人只是相信《舊約聖經》.
或者叫做《希伯來聖經》.
他們不相信耶穌是從死裡復活的基督.
所以他們會視基督徒是猶太教的異端或異教徒.
而這種從猶太人來的敵視.
就為這群基督徒.
相信主耶穌基督的一群人.
帶來幾個層面的打擊.
從信仰的層面.
信耶穌之後的猶太基督徒.
他們被趕出了猶太會堂.
昔日這個地方是他們在這裡聚集.
在這裡社交.
他們一起去學《舊約聖經》.
學摩西的《五經》.
但是信耶穌之後.
會堂的大門就對他們關上.
猶太教的信徒也不斷地去攻擊.
去批評他們.
說他們是不純正.
不正統的一些教徒.
去剝奪他們對信仰的歸屬感.
第二.
猶太教的信仰和他們日常生活息息相關.
環環相扣.
所以當這群猶太教拒絕基督徒的時候.
其實這群基督徒都失去了.
他們過去有的一些社會地位.
社會的身份.
在社會的層面.
他們也好像被隔絕了一樣.
他們的家族.
可能在知道他們信了耶穌之後.
就和他們斷絕關係.

$^{121}$昔日他們一起吃飯.
一起相聚的一些朋友.
一些論社.
知道他們信了耶穌之後.
不再理會他們.
甚或者.
他們每天去買賣的市集.
一些猶太人的市集.
也不再歡迎.
不再接待他們.
第三.
從一個政治的層面.
因為基督徒不被猶太人所承認.
他們就被拒於猶太教的門外.
這樣對於一個政治來說有什麼影響呢?.
原來當時羅馬政府.
他們會將當地不同的宗教區分.
一些是合法的.
一些是非法的.
當猶太教不承認基督徒.
這些基督徒就在羅馬政府的眼裡.
成為了一群非法的宗教組織.
基督徒要去傳福音嗎?.
他們要去聚會敬拜神嗎?.
這些就是進行非法的活動.
基督徒不可以很高調.
去出席一些公開的場合或者活動.
他們要聚會.
只可以在暗裡進行.
默默的才可以去敬拜主耶穌.
這就是人為的敵視.
逼迫帶來不同層面的影響.
第二個令到飛拉鐵飛教會很軟弱的主因.
就是來自周圍環境的創傷.
又或者叫做他們身份的不安.
飛拉鐵飛是一個什麼地方呢?.
飛拉鐵飛是一個位於火山大邊緣.
或者叫做地震大上面的一個城市.
在主後的十七年.
這裡曾經發生了一場毀滅性的大地震.

$^{161}$整個飛拉鐵飛城差不多移為平地.
即使他們經歷過重建也好.
但是餘震在接下來的很多年繼續發生.
這就是飛拉鐵飛教會身處的環境.
一個充滿著不安,不確定性.
隨時都會震,會動的一個環境.
我們可以想像.
住在飛拉鐵飛這裡的人.
他們一生長年都活在一個創傷,一個恐懼之中.
只要一震.
人們就馬上扔下所有東西.
馬上要跑去外面.
跑去城外面的曠野裡面.
甚至有歷史學家說.
其實當時有很多飛拉鐵飛的居民.
自從大地震之後.
他們不敢再回到城裡面去住.
只是夠膽在曠野,在城外面.
找個地方搭個帳篷就住在那裡.
不敢再回到城裡面.
他們隨時都要去逃跑.
另外說到飛拉鐵飛這個名字.
我們覺得很蹺口中文.
這個名字是怎麼來的呢?.
原來飛拉鐵飛這個地方的名字.
是來自主前大約140年.
當時建這個城的王.
一個別加摩王.
他去定這個名字.
這個名字是什麼意思呢?.
原來就是他為了紀念他一個很佩服,很欣賞,很喜歡的一個卿長.
為了這個城定了叫做飛拉鐵飛的名字.
剛才說到這個地方發生過大地震.
大地震之後又重建過.
原來有一個王當時發生過大地震之後.
他重建這個城的時候.
他就免除這個城要給的稅收.
並且又給資助他們去重建.
所以有一段時間為了紀念這個王.
紀念這個提比利亞王.

$^{201}$去幫這個城重建.
這個城改了名字.
叫做新凱撒利亞.
紀念這個王為他們做的一切恩情.
再後來.
這個城又再一次為了紀念另一個皇帝.
就用這個皇帝家族的姓氏再度改名.
不是叫飛拉鐵飛.
不是叫新凱撒利亞.
叫做飛拉維亞.
三個名字都很巧合.
但是原來這個城出現過的三個名字.
都是出自昔日的當權者.
彷彿表達這個城就是屬於這個政權.
這個城的人民都是他們所管治擁有的資產.
當權者可以肆意使用他們.
皇帝一聲令下他們就要服從.
當權者可以按心意去定義.
哪些人是非法 哪些人是合法.
在政治上被排擠.
在社會上被孤立.
在天災裡他們失去安全感.
在強權之下他們失去了自主.
這個就是飛拉鐵飛教會很真實的軟弱.
他們很勢孤力弱.
沒有控制權.
並且充滿著無力感的一群邊緣人.
不過耶穌在這封信裡.
祂說祂不單知道他們的軟弱.
耶穌更加知道.
更加看得見他們如何使用他們僅餘一點點的力量.
耶穌說他們也曾遵守我的道.
沒有棄絕我的名.
其實可能我們也會想.
只要這群人稍微妥協一下.
向羅馬皇帝上一炷香.
或者向猶太會堂低頭.
可能他們的日子會好過一點.
但是他們沒有這樣做.
他們卻是將最後一絲的力氣.

$^{241}$用來緊緊地抓住耶穌的旗幟.
對著這群很軟弱.
但是很忠心的信徒.
耶穌在信的開頭.
這樣去介紹自己.
祂說祂是聖潔.
是真實的.
祂又引用《倚創亞書》22章22節.
祂說自己是拿著大衛的約時.
開了的門沒有人可以關.
關了的門沒有人可以開.
祂很想表達祂才是真正掌有天國決定權的那一位.
然後這位真實的主.
這位大君王.
真正的大君王.
他針對著菲拉鐵妃他們的苦況.
他們的軟弱.
這位擁有最高權柄的主.
賜下了三個很有針對性的應許.
第一個應許.
主要在他們面前.
開一個敞開無人能關的門.
剛才我們記得.
猶太人要將他們拒於門外.
但是耶穌說.
沒有人可以關上這裡.
由主親自去打開.
由主親自去允許的門.
社會要圍堵他們嗎?.
政權要邊緣化他們嗎?.
耶穌說.
菲拉鐵妃教會面前的門是敞開的.
沒有人可以阻擋.
沒有人可以封殺他們.
而那些要去攻擊.
去針對.
處處去壓迫他們的猶太人.
主要將他們帶離.
帶離到這間菲拉鐵妃教會面前.
打敗.

$^{281}$耶穌要他們清楚明白知道.
要承認原來教會的基督徒.
才是真正的選民.
剛才我們有說到.
菲拉鐵妃教會.
菲拉鐵妃的人.
他們長期都很害怕地震.
他們覺得沒有任何建築物是穩固的.
但是.
耶穌第二個對菲拉鐵妃教會的應許.
就是要叫他們在神的殿裡作柱子.
成為一條柱.
代表的是永不動搖的安穩.
而主要親自將人.
建立在神永恆的殿裡.
作一條永遠都不會斷裂,不會倒塌的柱.
對於菲拉鐵妃的人來說.
剛才說他們長期處於一個.
他們長期對於他們身處的地方.
都很感到不安.
隨時都準備要逃命.
但是.
但是耶穌有應許他們.
他們必不再從那裡出去.
你們以後不用再這樣走出去.
不用再出去代表永恆的歸屬和安全.
菲拉鐵妃人可能一生都在逃跑.
但是耶穌基督承諾他們.
你們終於可以安頓下來.
不需要再恐懼.
不需要再流浪.
你們在神的家裡是永恆.
是一個安全的居所.
第三.
對於一次,兩次,三次.
被冠上羅馬皇帝名的城市居民.
耶穌應許他們說.
我要將我神的名.
我神承的名.
這個新耶路撒冷.

$^{321}$並且我的新名.
都寫在菲拉鐵妃身上.
耶穌說不再是甚麼菲拉鐵妃.
不再是甚麼新凱撒尼亞.
不是甚麼菲拉維亞.
是我神的名.
是我神的承.
新耶路撒冷的名.
是我的新名.
耶穌去到宣告.
他們不再屬於任何西族的政權.
他們是單單屬於主耶穌基督的天國子民.
他們是耶穌基督用重價買回來的寺寶.
對於軟弱的菲拉鐵妃教會.
耶穌賜下了三個.
一個好像出自永恆目光的英雄.
但是當我們從永恆回到地面.
當我們從天上回到地面.
我們經常說不要太離地.
當我們回到地面 回到現實.
我們又如何面對現實環境上的苦況呢?.
當猶太人仍然在逼迫他們.
當地震還在出現的時候.
當羅馬政權皇帝的影響力仍然存在的時候.
當現實是當我們仍然身處在軟弱,無力之中.
除了有剛才所說的永恆的盼望,應許之外.
我們人總會很渴望.
我們也很想可以脫離軟弱.
我們也很想可以成為強壯.
叫我們現在的生活不要活得那麼苦.
主耶穌有沒有答允這種很人性的心願?.
主耶穌有沒有叫菲拉鐵飛的教會從此不再受苦?.
第十節 耶穌說:.
「你既遵守我忍耐的道.
我必在普天下人受試煉的時候.
保守你免去你的試煉」.
不過 原來我們所看的這個和合本的翻譯.
對於聖經的意思,經文的意思有些微的簡化.
是一個較好的翻譯 我覺得是來自新漢語譯本.
「你既遵守了我忍耐的道.

$^{361}$在試煉將要來到全世界的時候.
就是試煉住在地上的人的時候.
我必保守你」.
原來耶穌真正的意思.
不是要完全免去菲拉鐵飛教會的試煉.
當菲拉鐵飛教會仍然生活在地上.
這一天 他們依舊要面對現實生活中種種的試煉.
當菲拉鐵飛教會不想再繼續下去.
不想再繼續軟弱.
不想再繼續承受這些苦難患難.
耶穌在第十一節竟然叫他們.
「你們持守你們所有的」.
耶穌沒有叫他們去爭取.
或者變成一些他們所沒有的東西.
耶穌沒有說你們這麼少人.
你們去變成一間大教會 變成一間千人的教會.
耶穌沒有這樣說.
當他們很軟弱的時候.
沒什麼錢 沒什麼影響力的時候.
耶穌沒有叫他們去變成一間有錢的教會.
有影響力的教會.
祂只是說「持守你所有的」.
「所有」有什麼?.
就是略有一點力量 再加上對主的忠心.
耶穌說這樣就夠了.
即使在軟弱無力之中.
只要我們持守對主的忠心就夠了.
今天我們未必像菲拉鐵飛教會一樣.
有猶太人去拒絕或排擠我們.
但或者我們同樣面對著很多道向我們關上的門.
可能是你發出過無數的申請或履歷.
但始終都得不到面試的機會.
可能是你很渴望建立家庭.
卻總是遇不到合適的對象.
可能是你很想你的家人順主.
你已經祈禱超過十年.
但他們的心仍然去到緊閉.
我們今天未必像菲拉鐵飛教會一樣.
經歷地震,經歷很多天災的震動.
但我們經歷的震動可能是經濟的震動.

$^{401}$可能是政權的更替.
可能我們經歷的震動是我們自己或身邊人健康的衰退.
雖然沒有倒塌在我們身上的建築物.
但倒在我們身上的可能是破裂的關係.
可能是家庭的重擔.
我們的名和身份不是被羅馬政府去定義的.
但是我們生命的價值.
彷彿卻是被這個社會,被公司的業績,被上司的評價評語.
政權的決策或銀行戶口的數字去定義我們.
今天可能我們都在各自的軟弱無力當中.
面對著不同的試煉.
而耶穌基督很想跟我們說.
在我們的軟弱裡面.
持守你所有的,持守你對我的忠心就夠了.
雖然我們竭盡全力.
等了很久都找不到工作.
但或許我們仍然可以持守對神的忠心,對神的信心.
相信祂是公認的主.
雖然我們可能在感情上遇不到合適的對象.
但或許我們可以不因為寂寞,不因為年齡的焦慮.
或者社會的壓力去降低我們的道德標準.
叫我們仍然去持守神對婚姻,神對感情,神對性的教導.
從而我們成為一個見證.
雖然我們的家人可能我們祈禱很久了.
他們仍然未信主.
但或許我們仍然可以在他們面前.
讓他們見到從主而來的那份忍耐,那份愛,那份包容.
並且那份不灰心.
雖然我們遭遇經濟的打擊.
但或許我們可以學習保羅如何處悲戰,處缺乏的秘訣.
我們可以重新審視,重整人生的優先次序.
雖然我們在社會上看到很多不公義的局面.
我們沒有辦法去推翻.
但或許我們仍然可以在日常遇到每一個接待的人.
最少在弟兄身上我們可以去接待他.
是讓神的公義在我們身邊逐點的微光去展現.
雖然我們和身邊的人關係破裂.
但或許我們可以選擇不用審判的目光去看對方或自己.
而是我們可以用主的恩慈,憐憫.
對待對方,對待自己.

$^{441}$雖然我們被工作,被職場的壓力壓到我們喘不過氣.
但或許我們可以將生命的終極價值放回創造主的身上.
持守主為工作,為安息所設立的界線.
雖然我們只剩下一點點的力量.
但或許我們仍然可以持守.
只剩下一點點的時間.
但或許我們可以不要停止聚會.
我們可以不要停止靈修.
只剩下一點點的愛心.
但或許我們可以不要吝嗇去關懷身邊每一個有需要的人.
雖然我們在職場上可能只是一個微小的齒輪.
只是一顆很小的螺絲.
但或許我們可以不要放棄活出耶穌基督的誠實和光明.
雖然工作,家庭,日常的瑣事已經令我們只剩下一點點的精力.
但或許我們仍然可以在睡前親近主和主禱告.
耶穌說一點點的力量再加上對著耶穌基督的忠心就足夠了.
足夠嗎?.
如何足夠呢?.
如何在試煉中繼續生活呢?.
繼續下去就甚麼都沒有了.
原來耶穌不是說甚麼都沒有.
原來耶穌說有些事是當你繼續持守著軟弱中的信心.
是你才有而別人沒有.
這就是祂所賜的冠冕.
有甚麼冠冕呢?.
可以當飯吃嗎?.
沒錯,冠冕不是可以當飯吃.
但卻是比有飯吃更加重要.
這就是主耶穌基督在試煉中對我們的保守.
這個保守原來與第十節上的遵守.
第八節出現的遵守是同一個字.
意思是耶穌對我們說.
「你持守著我,我同樣也持守著你」.
不是沒有艱難,沒有苦難.
因為世界有很多苦難.
是因為魔鬼撒旦可能會攻擊我們.
亦因為我們自己會犯罪.
我們的罪性會帶來苦難.
更加因為除了有魔鬼撒旦的攻擊.
除了因為有自己的軟弱.

$^{481}$更加因為生命本身就是一個經歷生老病死.
高低起跌的歷程.
聖經說生有時死有時.
禍有時,笑有時.
但主耶穌基督要保守你.
祂要親自去持守你,陪伴你,與你同在.
並且將那份可以承受生命中各樣患難,試煉的平安賜給你.
原來主耶穌基督自己.
祂就是我們在現實中,在現世中最需要的保護.
主耶穌基督自己就是我們現世中最安穩的處子.
我們現世中可以安身的神.
主耶穌基督寫給非拉帖非教會的信.
徹底顛覆了神國度中對強壯和軟弱的定義.
原來真正的強壯不是追求世俗眼中的龐大和強勢.
真正的強壯也不是憑藉一己之力.
去推翻現實中的困境.
真正的強壯是在充滿無力感,充滿不安的環境中.
仍然去到持守那一點點的力量,去到緊緊的抓住.
降伏於主耶穌基督.
這一份在軟弱中的強壯.
正正是我們經歷主親自保守的關鍵.
當我們承認自己的渺小,微小.
並且將主權交託給那位掌握著大位弱勢的主.
我們就不再是在現實洪流中任人擺佈的邊緣人.
而是在神永不震動的國度中作堅固不倒的處置.
開始時提到我有一位很喜歡的歌手.
他寫了那首六月的歌.
他同樣也寫了一首叫做《強弱》的歌.
歌詞是這樣的.
願你記得軟弱沒有錯.
痛亦沒有錯.
要投降方知可以被援助.
如此際想變更強嗎.
請你允許你真心痛過.
可能我們都曾經遭遇過不同程度的創傷.
以致我們活在恐懼之中.
以致我們很缺乏安全感.
可能這些是我們個人的經歷.
又可能好像飛拉鐵飛一樣.
來自遭遇天災的苦痛.

$^{521}$亦可能好像飛拉鐵飛一樣.
是整座城市整個地方我們一同經歷過集體的創傷.
叫我們不再覺得地上有一個地方是穩固.
讓我們可以安心.
但是主耶穌很想用祂大能的話語.
用祂的應許去安慰我們不安的內心.
很想醫治我們心裡的傷口.
包裹撫平我們心裡的傷痕.
主耶穌說只要我們願意在怨陽裡仍然忠心於祂.
祂就賜給我們一個永恆.
可以安全棲身歸屬的住處.
而在困苦充滿著思念的現世.
祂也要與我們繼續同在.
保守我們心中有一個被環境不被世界動搖的平安.
願我們在這個充滿著動盪不安的世代.
都能善用我們手上一點點的力量.
去緊緊抓住主耶穌.
我們同心低頭祈禱.
親愛的主耶穌.
那位拿著大位的約士.
聖潔而真實的主.
我們帶著真實的自己來到你的面前.
我們承認在這個充滿著震動不安.
各樣思念的世代裡.
我們就好像菲拉鐵非教會一樣.
我們只是略有一點力量.
我們面對著生活裡的難處.
面對著那些緊閉的門.
甚至面對著各樣的失去.
各樣的創傷.
我們心裡會感到無力.
會感覺到疲乏.
甚至我們會懼怕.
因為你沒有苛求我們要假裝剛強.
你也沒有要求我們要憑一己之力.
去成就甚麼的豐功偉業.
你只是呼召我們在軟弱裡.
去持守我們僅有的中心.
求聖靈幫助我們.
即使在職場的重擔.

$^{561}$即使在家庭的難處.
即使在健康衰退.
或者對未來迷茫的思念裡.
幫助我們仍然緊緊抓住你.
永不棄絕你.
你應許我們當我們若持守你的時候.
你就在思念當中保守持守我們.
求你與我們同在.
醫治我們心裡的傷口.
也將那份不被環境動搖的平安賜給我們.
我們確信在我們承認軟弱的那一刻.
就是我們經歷你剛強的開始.
我們獻上我們的禱告.
是奉旨的,所以我覺得他明智其靠.
\newpage



\section{士師記 7:2}
\label{sec:032296}
\textbf{聖經人聖經事:免得以色列人向我自誇}~~~~黃克勤牧師博士
\newline
\newline
連結: \href{http://www.church.com.hk/acms/content.asp?site=cdc\&op=show\&type=product\&code=032296\&layout=sermon}{\texttt{032296}} ~~~~ 語音日期: 2026-06-12
\newline
\newline
\hyperref[sec:032295]{\small{< < < PREV SERMON < < <}}
~
\hyperref[sec:index_chronic]{\small{[返順時目]}}
~
\hyperref[sec:032297]{\small{> > > NEXT SERMON > > >}}
\newline
\newline
士師記 7:2
\newline
\begin{longtable}{cl}
\hline
\hline
章節 & 經文 (和合本修訂版)\\
\hline
7:2 & \begin{tabularx}{0.7\textwidth}{X} 耶和華對基甸說：「跟隨你的人太多，我不能把米甸交在他們手中，免得以色列向我自誇，說：『是我自己的手救了我。』 \end{tabularx} \\ \\
[1ex]
\hline
\hline
\end{longtable}
$^{1}$各位弟兄姊妹大家好.
我們又一齊看聖經.
生命樹和大家看看聖經的人.
聖經的事.
認識一下聖經裡面.
上帝對我們講些什麼.
今天的題目繼續是講人物.
就是這個事司機殿.
當然還有他自己的同胞.
以色列人.
題目就是.
免得以色列人向我自誇.
這句說話是引自聖經的.
就是免得以色列人向我自誇.
這個我是上帝說.
向我自誇.
不用他們.
不想他們對著我誇口.
認叻.
說自己厲害.
這個也是要很小心的.
因為很多時候我們人.
贏啊勝利啊.
厲害啊.
我們廣東話叫做好得戚.
鬆毛鬆逆.
是我們厲害.
就忘記了上帝的恩典.
這個是.
神說免得以色列人向我自誇.
是我自己打贏敵人的.
事司紀錄章第一節.
耶路巴力就是基奠.
大家知道基奠的名字叫耶路巴力.
因為他就是把巴力.
即是他爸爸.
原來基奠.
早前兩次講了.
基奠家裡自己有一個拜巴力的壇.
原來他爸爸是拜巴力的.

$^{41}$而且家裡開了壇.
設了壇.
向巴力獻祭.
在壇的旁邊還有阿瑟拉.
就是巴力的太太的神像.
神叫他把那個壇拆掉.
整整齊齊.
四個字.
建造一個新的壇.
把阿瑟拉這個木頭的像.
砍成柴之後.
就做火燒.
然後把爸爸第二隻牛.
七年的公牛.
把牠tong 了獻在上面.
於是那些人就說.
嘩!查出原來晚上是基奠做的.
就找他爸爸了.
這個開神壇的巴力壇的爸爸.
約斯竟然就說.
如果巴力是神的.
就讓他自己和基奠爭辯.
於是基奠的名字就叫耶路巴力.
就是讓巴力和他爭辯.
就是說.
他爸爸在那一夜的兒子做了這件事.
他醒悟了.
悔改了.
基奠就是耶路巴力.
和一切跟隨的人早晨起來.
在哈律全旁安寧.
當然你知道安寧.
當然有水在附近才重要.
因為很多人.
基奠厲害.
上帝就透過他叫.
上面那一集說了.
有很多人.
很多個支派的人.
都一起來跟隨他.

$^{81}$要和米殿人.
我們被他壓制了這麼多年.
我們要和他爭戰.
米殿營就在他們北邊的平原.
有原因的.
就是安寧要麼在泉旁.
要麼在平原.
原因就是如果軍隊多就在平原.
軍隊少就找個地方躲起來.
軍隊超多你找個平原才能放得完.
才容易通傳和聚集.
但平原的不好處就是.
軍隊多你不夠水喝.
因此基奠就找了個泉.
就在他旁邊安寧.
而米殿營就在他北邊的一個平原.
他就靠近摩尼岡.
耶和華對基奠說.
「跟隨你的人過多」.
看看這個螢光的字.
「跟隨你的人過多」.
什麼叫過多?.
如果我是基奠就說.
「不是吧?是你帶領人來跟隨我」.
「然後去和米殿人打仗」.
「米殿人比我們還要多」.
怎麼會多?.
當然比少一個還要多.
「跟隨你的人過多」.
「我不能夠將米殿人交在他們手裡」.
即是跟隨你的那些以色列人.
「免得以色列人向我誇大」.
說什麼呢?.
「是我們自己的手救了我們」.
「免得以色列人向我誇大」.
「說自己的手救了自己」.
這裡有兩點跟大家去思考,靈修.
第一就是.
「跟隨你的人過多」.
這一點如果我是基奠.

$^{121}$我就會抗議.
我就會爭論.
怎麼會多?打仗嘛.
多的人就好辦事.
人太少不行.
人家有幾萬人.
我們現在數完也不夠.
怎麼會過多呢?.
「過多到以致我不能讓你勝過米殿人」.
「我不幫你」.
這樣也行?.
為什麼?.
「免得以色列人向我誇大」.
誇大誰呢?.
誇大自己手的能力.
這就是第二點.
在這裡我們學到一些熟練的功課.
就是有時候面對一個很大的項目要做.
教會很大的事工.
或者上班有一個很大的工作.
或者是一件大事又大又難擺在前面.
通常我們都會祈禱.
「主啊,求你幫助我」.
「你給我智慧力量」.
「求你給我合適的人來幫我忙」.
哇,祈含肚,憤外迫切.
終於戰戰兢兢,小小心心.
很小心地把事情做好.
「啊,你真是厲害了」.
「這件事這樣也讓你做到」.
「找了這麼多人幫忙」.
「這樣的計謀你也想到」.
「啊,真是厲害」.
「有領導才能,很厲害」.
「你的辦事能力真是高了」.
「甚至比你老闆還要高」.
「哇,很厲害」.
於是我們很多時候會自大驕傲.
當被人讚了兩讚後,失魂站不穩.
飄飄然飄上天.

$^{161}$飄上天就差不多.
「幸好我厲害」.
「幸好我認識的人多」.
「幸好我也有些思想能力」.
「不然這件事哪裡能搞定」.
「你看看旁邊的組就搞不定」.
「幸好我搞得定」.
但其實心裡自高自大.
覺得是「我得」.
很多時候我們會.
一開始禱告很懇切.
不過後來卻是搶奪了上帝的榮耀.
真的不好.
CPA說「耶和華,耶和華」.
「榮耀不要歸於我們」.
「不要歸於我們」.
「要因你的慈愛和誠實歸在你的名下」.
這樣就對了.
「不要歸於我們」.
「榮耀不要歸於我們」.
「要因你的慈愛和誠實歸在你的名下」.
少一點要記得.
「免得」.
「我厲害」.
「紅啊」.
「那個人紅啊」.
「不好意思」.
「那個人紅呢」.
「跟我的人多」.
「於是我救了我自己」.
「我有計謀,我自己的手」.
「記住是上帝的手救了我們」.
「而不是自己的手」.
現在你要向這些人宣告.
現在上帝跟大家說.
「凡懼怕膽怯的」.
「可以離開基烈山回去」.
直接這樣說.
如果我是基甸.
我又不會這樣說.

$^{201}$甚至你叫我說這些話.
很羞辱人.
直接說.
「你們各位誰怕的」.
「現在可以退出回家」.
說這些話.
那你什麼意思.
如果我真的有些事想回家.
那些人就說我怕了.
「縮沙了」.
「現在領算身走了」.
「現在走的全部沒用了」.
「驚慌了」.
「懼怕了」.
「沒膽子了」.
「就跟著來拿威了」.
「現在一說可以回家就走了」.
「凡懼怕膽怯的」.
「可以離開基烈山回去」.
其實我不想說這句.
如果你感動得他們來.
不如你感動他們走.
為什麼要我說呢?.
跟隨你真的麻煩了.
要我宣告.
要你向這些人宣告.
「凡膽怯的回家吧」.
基甸也算是好東西.
算了.
我都說了.
說啊.
於是二萬二千人回去.
只剩下一萬.
就是說.
算起來本來是三萬二千人.
三萬二千人.
不打也看得到.
三萬二千人一起拿一塊石頭.
扔過去那三萬塊石頭.
嘩!真的.

$^{241}$填平了你也行.
起碼跟米電人.
阿瑪利人和所有東部的人.
來的民族.
即是摩爾阿曼那些.
以東那些一起來的.
都可以撐一會兒.
三萬二千人.
只是擺起一個陣你也嚇一驚.
誰知道.
如果心裡有疑慮的.
膽怯的.
回家吧.
不要緊的.
神在這裡.
我在這裡.
我們搞定了.
於是有二萬二千人會去.
只剩下一萬人.
當你看著三萬人的時候.
心裡面很有信心.
但是你看到.
很多人起來收拾東西離開的時候.
嘩!你想一想.
二萬二千人起來收拾東西離開.
那個隊伍大不大.
嘩!真的被所有場合的散場燈你開.
走囉.
嘩!那條隊真的厲害.
嘩!好像螞蟻搬家那樣.
我走了.
哎呀!糟糕了.
這麼多人離開.
怎麼辦.
還有人在離開.
哎呀!怎麼辦.
你的心裡面真的很慌.
我想機電經歷了這些.
真的沒有聲音出.
怕得不得了.

$^{281}$如果我是機電.
我都走一份.
哎呀!神呀!我都怕.
我先走.
不過.
他唯有站在這裡.
只剩下一萬人.
神就說.
不用這麼多.
要心裡面有膽怯的就走吧.
留下一些我不怕的.
這個就是一個考驗.
當你是那三萬二千人裡面.
其中一個的時候.
打仗怕不怕.
當然怕了.
我武藝很好.
我勤練武功.
很能打的.
這樣的幾萬人之間.
亂殺之間.
到處都有人在打.
你是沒有辦法保護到自己的.
我有戰陣.
有陣式.
有陣式有時都沒有用.
陣式會衝亂了.
然後你發覺到處都是敵人.
剛剛幾個在你那裡.
三,四個人夾擊你的時候.
你多有陣式都沒有用.
多好武藝都不夠人來.
因為三,四個人一齊找長矛刺你.
你都不知道躲到哪裡去.
你以為拍戲嗎?.
跳跳起.
然後四,五支矛.
然後你站在上面.
跳又一個關刀下來.
就殺了五個.

$^{321}$一個圈.
你以為拍戲嗎?.
是不行的.
在這些戰陣亂殺之間是很可怕的.
你想像起來都會想走.
於是剩下來的一萬人.
是真的夠膽的人.
無論他是對自己個人有信心.
或者對神有信心.
但他總是有信心.
不怕留下來.
所以其實都好的.
因為就算看回這件事.
從現在的角度來看都是對的.
譬如三萬多人好看.
不過不好用.
打了.
馬上有人在後面喊.
撤退.
那些人就亂跑了.
我以前認識一些.
本來做強盜後來做傳道的人.
其中有些是在黑社會裡面打的.
不過他們說.
我們出去打架.
出去斬人.
大家鬥斬.
你以為真的像拍電影那樣.
把西瓜刀拿出來.
打了幾個字.
不會啦.
是這樣拿出來.
大家打了幾下之後.
打啊打啊打啊殺啊斬啊.
沒斬了多少下.
後面就有個小子在後面喊.
喊什麼.
警察啊.
王熙啊.
馬上就走了.

$^{361}$大家就說有王熙啊.
大家是走得很有理由的.
打了幾下就斬了一個.
斬不到斬不到.
就走了.
有些受傷了.
有些沒受傷.
如果真的斬死的不是很多.
都是被暗算到才會.
你見到其實怕不怕.
怕.
因為你去到那些位置就怕了.
一打的時候後面就有人喊.
走啊.
這樣就有人走了.
真的,如果三萬多人一打.
有二萬人裡面有這麼多人喊.
走啊.
就真的有很多人走了.
接著走啊走啊.
有二萬多人走.
常障輸的了.
不如還是剩下不怕的一萬.
阿神真的有智慧的.
對的.
因為我想三萬多人.
不打得看得.
看得你會站在那裡嗎.
做立仗嗎.
是要真的打才行的.
夠膽的是不會走的.
那剩下一萬人.
神是對的.
如果怕就回去吧.
太多了.
另外.
祂會搞亂你的當的.
而且.
神就說.
搞不搞亂當都不要緊.

$^{401}$因為我會搞定那個當的.
不過最重要的就是.
我不想他們贏了之後自誇.
當自誇之後.
他們以後就不懂得依靠我.
如果依靠自己.
已經注定不行了.
我想這個是上帝的苦心.
這個免得以色列人向我自誇.
其實不是說.
你不給一些財給以色列人.
給你的兒女.
不是這個問題.
而是.
如果只是他們自誇.
以後他們就不懂得依靠神.
用他的手.
他們是不會救世界打的.
就是這樣的.
正如聖經說.
你們來了我就不能作甚麼.
所以都是要和神一起.
祂保護我們.
我們依靠跟隨祂就最好了.
好,我們一起祈禱吧.
天護上帝我們很感恩.
開始看這段聖經的時候.
心裡面有很多不明白的.
三萬二千人不好嗎?.
不能打看燈.
又是你招聚這麼多人來的.
這件好事.
為甚麼你會說人太多.
士兵怎會怕太多呢?.
其實再重新想過.
真的怕太多.
太多很難管理.
太多又要食物.
太多又要這麼多照顧.
又怕一會他們自己叛變.

$^{441}$或者他們懼怕逃走.
影響了整個士氣.
不過實際就是這樣.
這麼聰明有智慧.
就叫怕的可以回家.
而你知道如果太多人贏了.
會自誇.
這個才是最致命的.
這個是屬靈上最致命的地方.
主要幫助我們.
更多依靠你.
更多榮耀你.
要將榮耀歸在你名下.
知道我們來了你就不能作甚麼.
求主幫助我們.
更多依靠.
不會自誇.
搶奪你的榮耀.
祈禱奉靠耶穌基督的名求.
阿們.
\newpage



\section{約翰福音 8:30-36}
\label{sec:032297}
\textbf{榮耀君王:上帝兒女}~~~~林鍵牧師
\newline
\newline
連結: \href{http://www.church.com.hk/acms/content.asp?site=cdc\&op=show\&type=product\&code=032297\&layout=sermon}{\texttt{032297}} ~~~~ 語音日期: 2026-06-13
\newline
\newline
\hyperref[sec:032296]{\small{< < < PREV SERMON < < <}}
~
\hyperref[sec:index_chronic]{\small{[返順時目]}}
~
\hyperref[sec:032298]{\small{> > > NEXT SERMON > > >}}
\newline
\newline
約翰福音 8:30-36
\newline
\begin{longtable}{cl}
\hline
\hline
章節 & 經文 (和合本修訂版)\\
\hline
8:30 & \begin{tabularx}{0.7\textwidth}{X} 耶穌說這些話的時候，有許多人信了他。 \end{tabularx} \\ \\ \relax
8:31 & \begin{tabularx}{0.7\textwidth}{X} 耶穌對信他的猶太人說：「你們若繼續遵守我的道，就真是我的門徒了。 \end{tabularx} \\ \\ \relax
8:32 & \begin{tabularx}{0.7\textwidth}{X} 你們將認識真理，真理會使你們自由。」 \end{tabularx} \\ \\ \relax
8:33 & \begin{tabularx}{0.7\textwidth}{X} 他們回答他：「我們是亞伯拉罕的後裔，從來沒有作過誰的奴隸，你怎麼說『會使你們自由』呢？」 \end{tabularx} \\ \\ \relax
8:34 & \begin{tabularx}{0.7\textwidth}{X} 耶穌回答他們：「我實實在在地告訴你們，所有犯罪的人就是罪的奴隸。 \end{tabularx} \\ \\ \relax
8:35 & \begin{tabularx}{0.7\textwidth}{X} 奴隸不能永遠住在家裡；兒子才永遠住在家裡。 \end{tabularx} \\ \\ \relax
8:36 & \begin{tabularx}{0.7\textwidth}{X} 所以，神的兒子若使你們自由，你們就真正自由了。」 \end{tabularx} \\ \\
[1ex]
\hline
\hline
\end{longtable}
$^{1}$歡迎大家來到我們的周六崇拜.
上年七月我們開始研讀《新約聖經》當中的《約翰福音》.
講道系列的題目是「榮耀君王」.
我們花了半年的時間完成《約翰福音》第一至第四章的宣講.
接著今年一月我們開始研讀《約翰福音》第五章.
六個星期前我們來到《約翰福音》第七章.
從故事的時序來看.
《約翰福音》第七章其實是一直延伸到《約翰福音》第八章.
因為大家知道在希臘文的原文裡面其實是沒有章節的.
可以分為三大段.
第七章至第八章.
第一是猶太人的鑽棚節近了.
第二是節期已經過了一半.
大概是第四天的時候耶穌來到耶路撒冷.
最後是節期的最後一天.
在那天耶穌在耶路撒冷聖殿裡面向眾人宣告.
祂說「我就是生命的活水」.
「我就是世界的光」.
這樣看來好像是兩個不同的自我宣告.
但其實他們是一個二合一的自我宣告.
耶穌說「我是」.
「我是」是耶和華上帝在舊約裡面的名字.
「我是上帝」「我是上帝獨生的兒子」.
所以祂的名字叫「我是」.
「我又是我是」.
「我是」「我就是生命的活水」.
和同時之間二合一世界的光.
我們可以說這個二合一的自我宣言.
其實是約翰福音第七章至第八章的核心.
就好像一個核彈爆炸.
輻射是擴散至影響周圍.
來到這裡我們必須要謹記.
作者約翰給我們的溫馨提示.
第七章第三十九節約翰強調又強調又強調.
耶穌這話是指信他的人要受聖靈所說的.
那時還沒有次下聖靈.
因為耶穌還沒有得榮耀.
由此可見真正的福音核彈.
是耶穌基督的榮耀十家.
再加上聖靈的降臨.

$^{41}$但當你想到聖靈的時候.
你必須要記得勝靈其實是耶穌的靈.
所以其實這兩件事是一樣的.
耶穌在十字架上得榮耀.
隨著這件事發生的時候.
他升天他參拜他自己的靈.
來到這裡當中替罪人眾人洗禮.
聖靈的洗禮.
這是福音的核心.
於是榮耀福音核彈爆發.
輻射擴散覆蓋了前和後.
在前面我們有兩個故事.
一個故事是不信的弟兄.
接著下面那個故事是安息的葛禮.
如果大家還記得.
舊約裡面的葛禮.
最後其實是指向耶穌基督在十字架上.
替罪人所承受的審判葛禮.
耶穌受了兩個葛禮.
第一個是他第八天出生之後.
他行了一個葛禮.
他就是上帝應許的救主.
第二個葛禮是為我們割的.
是在十字架上.
他被上帝切割割除.
以致我們不需要受到那個刑罰.
為甚麼?因為世人都犯了罪.
虧缺了上帝的榮耀.
罪的公鑒乃是死.
死後且有審判.
如果耶穌不死不受審判.
就是我們要死我們要受審判.
於是我們可以說.
耶穌為著他的兄弟釘了十字架.
被上帝割除.
承受了上帝的憤怒.
於是最終成為他兄弟的救贖.
這是前面兩個故事.
現在到後面兩個故事.
後面兩個故事是甚麼呢?.

$^{81}$第一是上星期的死在罪中.
到今天使你自由.
不單止使你自由.
你見到今天的題目不是使你自由.
使你自由以致你可以成為上帝的兒女.
使你自由不是那個高潮.
上帝兒女才是那個高潮.
但我們進入今天的經文.
《約翰福音》第八章第三十節.
接著我們再讀之前的第八章二十一和二十四節.
耶穌說這些話的時候.
有許多人信了祂.
說甚麼話的時候.
說前面那些話的時候.
你們會死在自己的罪中.
第八章第二十一節.
第八章二十四節.
你們會死在自己的罪中.
你們若不信我就是.
就會死在自己的罪中.
這是一個不可思議的結局和回應.
試問誰可以連出三拳.
三拳是你們會死在罪中.
而結果竟然是許多人信了主.
如果上戀是你們死在自己的罪中.
下戀怎麼可能是許多人信了主呢.
於是你說這真是一個神蹟奇事.
然而英文有句說話.
Too good to be true.
有些事如果好像很得令人難以置信的話.
那樣東西很可能根本就不是真的.
所以你橫看橫看這件事情發生.
原來它不是一個神蹟奇事.
它也不是一個所謂屬靈的大復興.
依我預見.
所有一切都是一場誤會.
第八章第28,29節.
解釋給我們知道.
為什麼突然間會有這麼多人.
許多人信了耶穌.

$^{121}$第八章28節這麼說.
你們舉起人子以後就會知道.
我就是那位.
並且知道我沒有憑自己做的是照著.
我父所教導我的.
差我來的那位與我同在.
我想很有可能是有人誤會了.
舉起人子這句說話的意思.
他們以為耶穌將要在地上作王.
於是祂被舉起.
然後祂被舉起的時候.
祂的名字叫做人子.
但我們回到《大義地主》第七章13-14節.
舊約裡面.
人子將要從上帝那裡得著一個永恆的國度.
於是那些人都熟讀舊約聖經.
於是想起祂要被舉起.
祂要在地上作王.
然後還要領受一個永恆的國度.
於是有許多人信了祂.
其實他們全部都不明白.
耶穌的國度是不屬於這個世界.
他們也不知道耶穌的榮耀不是地上的榮耀.
舉起這個動詞在《約翰福音》裡面曾經出現過三次.
最終的所謂舉起.
就是耶穌上升.
被舉上升.
回到天上的榮耀.
回到天父那裡.
那個叫做被舉起.
舉起all the way to heaven.
但是為了拯救罪人.
耶穌選擇經過榮耀的十字架.
祂不需要經過榮耀.
祂可以馬上回家.
直飛.
但是祂不直飛.
祂要經過十字架.
這個是舉起人子的真正意義.
很可惜.

$^{161}$這一群耶穌的新粉絲.
他們完全聽不明白耶穌的說話.
其實他們應該好好請教猶太人當中的前輩.
如果你回去看.
你會發現周圍有些人早就不信耶穌了.
那些前輩可能曾幾何時都是耶穌的粉絲.
但是現在他再不是了.
因為他們後來發現耶穌的國度不屬於這個世界.
他發現耶穌的榮耀不是地上的榮耀.
我想帶大家去看看前輩是怎樣回應耶穌的.
前輩之前已經在回應耶穌了.
耶穌一直說你們死在最終的時候.
前輩說難道他要自殺嗎?.
很明顯地在笑他.
你是哪一位?.
你到底是他?.
憑什麼告訴我們我們要死在最終?.
接著又裝作不知道.
他們不明白耶穌.
但是楊寶欽第一章第十節已經告訴我們.
他在世界世界是藉著他做的.
然而這個世界不認識他.
這個世界對於耶穌是完全不認識的.
對於他所說的話是懵然不知一頭霧水的.
他們不認識耶穌.
他們不明白耶穌.
他說到最後是他們不接納耶穌.
說到這裡.
其實有不少的聖經學者會說.
耶穌的教導將當時的猶太人一分為兩半.
當然不是平等的.
耶穌的教導將他們分為兩組人.
第一組是信他的人.
第二組是不信他的人.
當然不信他的人比較多.
信他的人是比較少的.
但是依我預見.
兩組猶太人其實都不信耶穌.
一個一開始已經拒絕不信.
另一個自以為信.

$^{201}$但是最後有沒有想過.
這兩條路最後會回到同一條路.
好像有人信有人不信.
結果兩條路兜兜轉轉.
回去殊途同歸.
這兩群人最後都是半逆不信.
會接受上帝嚴重嚴重終極的審判.
真正的相信耶穌.
是跟隨耶穌逆水行舟逆風飛行.
追求十字架的榮耀.
追求天上的國度.
很有趣.
追求天上的國度.
但是又向下走.
在地上所有東西都是倒轉了.
向下等於向上.
向上等於向下.
你越是向上就越是向下.
你越向下卻要向上.
我們要跟隨耶穌.
逆水行舟逆風而行.
追求十字架的榮耀.
追求天上的國度.
這群新的粉絲.
跟耶穌說真的.
就是因誤會而結合.
那下一句是什麼呢?.
當然是因了解而分手.
所以你會發現到了第37節.
我們今天沒有正式去看.
你已經發覺他們已經開始要分手了.
第8章第37節.
要向耶穌這樣說.
你們想要殺我.
耶穌看得出已經知道.
因為你們心裡容不下我的道.
第8章43節.
耶穌再說.
你們聽不進我的道.
你們容不下我的道.

$^{241}$你們聽不進我的道.
你們出於你們的父魔鬼.
他說謊是說謊者的骨.
他們的心容不下耶穌的道.
他們聽不進耶穌的道.
原因是什麼?.
因為他們滿腦子都是魔鬼的謊言.
你不要太快開心.
不過我聽得中.
我容得下耶穌的道.
我經常在說.
你容得下.
聽得進的意思.
就不是你可不可以正確說出來.
我剛才說過的那篇道是什麼.
我們現在很多弟兄姊妹在做功課.
他聽完那篇道.
他用三點來總結.
可能我會說總結得很好.
接著下面還有一句.
Go and do and live it out.
你要活出來.
你不活出來.
你是容不下耶穌的道.
耶穌的道是用來活出來的.
不是用來考試.
接著說我達到滿分.
這群人連耳朵都容不下.
他們根本從來沒有相信耶穌.
真正的相信耶穌.
我們不斷重複地思想.
真正的相信是notitia.
清楚地明白耶穌的福音.
第二.
ascensus.
心裡確認確信.
耶穌說的是對的.
耶穌說世人都犯了罪.
沒錯.
我就是那個世人.

$^{281}$就不是說隔壁街的那個人是那個罪人.
有了notitia.
有了ascensus.
最後是要你信靠鬼神Fabrizio.
這群人.
大概聽完之後.
有了notitia.
有了清楚地明白.
但是他們沒有辦法心裡認信.
他們覺得耶穌是瘋的.
覺得耶穌看錯了他們.
更加不用說.
他們完全不會信靠鬼神.
來到這裡我們要問一個問題.
耶穌究竟說了些什麼.
那麼刺激他們.
令他們如此激動.
一節而已.
耶穌說完之後.
不知為什麼他突然間容不下他的道.
耶穌究竟說了些什麼.
去刺激他們.
令他們如此激動.
還想要殺他.
OK.
我想帶你去看回第八章的31和32節.
從信他到想要殺他.
問題出在哪裡.
就是出在這兩節聖經.
第八章第31節.
耶穌對信他的猶太人說.
不要說真信假信.
總之當他信了.
你們若繼續遵守我的道.
你就是我的門徒了.
好像沒什麼問題.
下一句.
你們將認識真理.
真理會使你們自由.
如果你問我.

$^{321}$我會說第八章第31節是OK的.
不過第八章第32節就壞了大事.
會令人有點反胃.
如果我當時在場.
我會建議耶穌.
可不可以改一改.
改一改第八章第32節.
與其說你們將要認識真理.
真理會使你們自由.
倒不如這樣寫.
你們將要看見天國.
我會使你們得釋放.
你這樣寫就正面很多了.
可能有人會問.
分別在哪裡.
分別的地方就是.
Who is to blame.
是誰的錯.
你現在這樣的情況是誰的錯.
如果我們說.
你們將看見天國.
我會使你們得釋放.
意思是錯在外面.
我們今天有什麼問題.
我們的問題在外面.
我們的問題很明顯在外面.
所以你去教會報道會.
教會報道會有些人信主很簡單.
將問題放在外面.
你有什麼問題.
耶穌會幫你解決.
錯在外面.
我來幫你解決.
你們看見天國.
我會使你們得釋放.
這樣就沒事了.
但耶穌不是這樣說.
耶穌說你們將要認識真理.
真理會使你們得到自由.
什麼意思.

$^{361}$錯在裡頭.
錯在你的裡頭.
而你現在要承擔.
要負責你裡頭出了的問題.
原來你現在.
是做了別人的奴役.
奴僕.
所以你要通過.
我給你的真道使你.
裡面得到自由.
你要承擔你自己的後果.
你要承擔你要負責.
於是令他們很不開心.
耶穌反過來.
他拒絕了說外在的問題.
他選擇了說內在的問題.
我不是說耶穌覺得外在沒有問題.
我們活在一個墮落的世界裡面.
我們外在有沒有問題.
有問題.
政治有沒有問題.
有問題.
社會有沒有問題.
有問題.
很多問題.
但是.
耶穌的福音.
不是說外面.
是說裡面.
耶穌拒絕和你討論外在的問題.
選擇和你說你內裡的問題.
而這群人竟然聽明白了.
一句話他聽明白.
於是引來一個極大的恨惡和決裂.
他一聽這句話.
他聽到我知道你說什麼.
我知道你戳我背.
我來是想你幫我解決我外面的問題.
我們這群都是好人.
如果你真的是上帝派來的.

$^{401}$是不是應該上帝派來幫我們這些好人申冤.
你這樣說我沒有問題的.
但是你現在說我們是壞人.
你說話.
於是他們很大很大的反應和耶穌決裂.
他說我們是亞伯拉罕的後裔.
從來沒有作過誰的奴隸.
你怎麼說我會使你們自由呢.
這個是亞伯拉罕在約翰福音第一次的出現.
亞伯拉罕在約翰福音裡面一共出現過11次.
11次都是約翰福音第八章.
所以由這一節開始.
第30單節一直到第58節裡面.
亞伯拉罕會不斷重複不斷重複11次的出現.
所以未來的兩三篇裡面都離不開亞伯拉罕了.
你已經知道了.
甚至今天這一篇.
稍後有一節聖經.
就這樣看是沒有他的名字.
但是是關於亞伯拉罕的.
他們說我們從來都不是奴隸.
他們說的是心裡面的奴隸.
他說的是屬靈裡面靈女的奴隸.
因為在外在而言.
以色列人其實是多次成為別人的奴隸.
由埃及出來.
後來被亞瑞人攻破他們的北國.
接著是被擄巴比倫.
接著是波斯希臘.
現在是羅馬.
所以從外在而言.
他們其實是在做別人.
在看別人的臉色.
在做別人的奴隸.
他們是心裡面不服.
但是耶穌說的不是外面.
耶穌說的是裡面.
耶穌說你裡面有個問題.
但是我可以幫你解決.
他想都不想.

$^{441}$你瘋了.
我沒有問題.
我的問題是外面.
你幫我解決問題.
外面我就叫你做我的神.
我的救主.
但是你不肯.
走.
他們說既然.
我們從來都沒有在靈裡面做奴隸.
耶穌你也不需要用你的道.
用你的真理來浪費時間.
釋放我們獲得靈女的自由.
耶穌回答他們說.
還不放棄.
我實實在在的告訴你們.
所有犯罪的人.
就是罪的奴隸.
我想你們已經知道.
耶穌不是嫌棄他們.
所以破口大罵.
吃碗底飯碗麵.
我現在說實實在在.
他不是這樣的.
他不是要定他們的罪.
而是他帶著無盡的憐憫和慈愛.
看見他們就愛他們.
於是他呼籲他們要信靠跟隨他.
就可以得到自由得到拯救.
他不是告訴你.
你知道你是罪的奴隸.
他是想告訴你.
我有辦法.
我可以幫你搞定.
只要你知道你是罪的奴隸.
我就可以將你從罪裡面釋放出來.
在這裡我們已經在背景裡面聽到.
讓福音的主題.
上帝愛世人.
甚至將他的獨生子賜給他們.

$^{481}$叫一切相信的.
因為耶穌的轉捩點就是.
你相信我.
你如果相信我.
你就不需要做罪的奴隸.
叫一切相信他的.
不致滅亡.
反得永生.
因為上帝差他的兒子.
來到世上.
不是要定世人的罪.
而是要使世人因他得救.
信他的不被定罪.
不信他的罪已經定了.
因為他不信上帝.
獨一的名字.
我想帶你回去.
第八章二十一二十四節.
你們會死在自己的罪中.
你們若不信我.
就必會死在自己的罪中.
如果你不信我.
你就會繼續做罪的奴隸.
這個是耶穌在說的.
如果你繼續做罪的奴隸.
你就必死在自己的罪中.
但是不一定要這樣的.
結局可以改寫的.
如果你信我.
當然我們說的是真心信.
不是假意信.
如果你真心全心信靠跟隨我.
就必從罪中的奴役和審判中.
拯救你出來.
我會使你從奴役中得自由.
我會使你從審判中得拯救.
你明白嗎.
兩件事是一樣的.
你是罪的奴役.
你需要釋放.

$^{521}$如果你信我.
我可以釋放你.
如果你繼續做罪的奴隸.
你會死在自己的罪中.
但是如果你信我.
我可以拯救你.
我可以連你罪的審判.
都可以幫你得到救贖.
但是我們來到最後兩節.
很快已經去到最後兩節.
最後兩節是《約翰福音》第八章.
第三十五節和三十六節.
耶穌這樣說.
奴隸不能永遠住在家裡.
兒子才永遠住在家裡.
這是第一節.
接著第二節.
所以上帝的兒子.
若使你得自由.
使你自由.
你們就真正自由了.
你明白耶穌說了些什麼嗎.
原來在第三十五和三十六節的背後.
藏著兩個舊約的故事.
你可以開始猜一下.
第一個是《創世記》裡面的故事.
藏在《約翰福音》第八章第三十五節.
三十五節的背後有一個《創世記》的故事.
三十六節的背後.
藏著一個出埃及記的故事.
我想先和你看出埃及記.
接著再回去看《創世記》.
我們先看出埃及記的故事.
上帝的兒子耶穌基督說.
他若使你自由.
你就真正自由了.
耶穌正在帶領我們回去出埃及記.
因為出埃及記是一個講自由.
北島自由的一個故事.
我帶你去看的是出埃及記第十二章.

$^{561}$第二十一至二十三節.
大家很熟悉的經文.
你們要為家人取羔羊.
把愉悅的羔羊採了.
要拿一把牛膝草.
招盤裡的血.
把盤裡的血.
塗在門尾上和兩邊的門框上.
直到早晨.
你們誰也不可以出自己家裡的門.
二十三節.
因為耶和華要走遍埃及.
私行擊殺.
他看見血在門尾上和兩邊的門框上.
耶和華就必逾越冷門.
我想你想想上帝是怎樣拯救以色列人的.
上帝拯救以色列人的方法是.
通過十災.
上帝要帶以色列人出埃及.
而埃及的法老和上帝拗手瓜.
於是上帝通過十災.
尤其是第十災.
去擊殺埃及人的頭生.
但是如果沒有逾越羔羊的母血.
以色列人也難逃上帝的審判.
於是上帝設立了逾越節.
羔羊和羔羊的血.
以致可以保住以色列人的命.
以致可以帶他們出埃及.
於是我們會說.
通過逾越節的羔羊.
以色列人逃離了埃及的奴役.
獲得了肉體上的自由.
大家明白嗎.
上帝的兒子使你自由.
你就真自由.
我們說的是埃及的.
他們現在得到的是甚麼自由.
是肉體上的自由.
是外在的自由.

$^{601}$你記得耶穌說的是甚麼.
內在的自由.
可怕的事.
可惡的是.
他們外面得到自由.
他們裡面有沒有自由.
You shall know the truth.
the truth shall set you free.
你要認識真理.
真理會使你得自由.
那個自由是在裡面.
他不是在說羅馬人.
他們要的是外面.
現在以色列人在《舊約聖經》裡面.
他們得到了身體上的自由.
但是可怕的事.
他們沒有逃離偶像之家.
埃及偶像之家.
屬靈的奴役.
他們仍然深想埃及.
他們仍然貪愛世界.
他們俯伏敬拜金牛族.
他們不應該拜金牛族.
他們應該拜的是上帝的羔羊.
所以你明白為何耶穌回來告訴你.
你需要自由.
因為儘管你出了埃及.
但是你心裡仍然在埃及.
我救了你出羅馬.
你猜你會怎樣.
你會繼續去敬拜羅馬的偶像.
世界的偶像.
耶穌說犯罪的人是罪的奴隸.
拜偶像的人是偶像的奴隸.
你聽聽.
愛世界的人是世界的奴隸.
追求地上榮耀的人是地上的奴隸.
所有的東西都是同一樣東西.
罪的奴隸.
犯罪的人是罪的奴隸.

$^{641}$拜偶像的人是偶像的奴隸.
愛世界的人是世界的奴隸.
追求地上榮耀的人是地上的奴隸.
於是耶穌說.
你們會死在自己的罪中.
他們.
舊約的世界.
死在自己的罪中.
我們.
也會死在我們自己的罪中.
你可以替代.
你不是單單在說罪.
你也在說這個世界.
你在說地上的榮耀.
所有的東西你扔進去就對了.
因為你愛這個世界.
你就是地上的奴隸.
你會死在你的罪中.
死在你的奴役當中.
耶穌說.
除非你們.
他們.
我們.
全心盡心.
信靠跟隨耶穌.
上帝的羔羊.
你要做什麼.
要捨棄地上世界的偶像.
和榮耀.
你一定要捨棄金牛族.
如果不是.
你一定會死在.
自己的罪中.
這是第一個故事.
第二個故事.
有關的.
所有的東西有關的.
然後我們有第二個故事.
是創世記的故事.
耶穌這麼說.

$^{681}$他說奴隸不能永遠住在家裡.
兒子才永遠住在自己的家裡.
這裡所指的是什麼故事呢.
你猜不猜得到.
主要是學的題目.
這個背後原來是一個關於亞伯拉罕的故事.
我們現在要跳過舊約聖經的亞伯拉罕的故事.
直接帶你去新約.
看看保羅怎樣幫我們做這個總結.
我們現在要shortcut了.
我帶你去看.
佳太書第四章.
保羅怎樣幫我們總結這個故事.
他說亞伯拉罕有兩個兒子.
一個結果是奴隸的兒子.
另一個才是真兒子.
亞伯拉罕有兩個兒子.
佳太書第四章第二十二節.
一個是使女生的.
一個是自由的婦人生的.
那使女所生的是按著肉體生的.
那自由的婦人所生的是憑著嬰許生的.
然而經上是怎麼說呢.
是把使女和她的兒子趕出去.
不可以長期留在家裡.
你記不記得他講的那句.
因為使女的兒子絕不能與自由的兒子一同承受產業.
亞伯拉罕有兩個兒子.
第一個叫做以撒瑪利.
是他還年輕的時候.
多年輕啊.
86歲.
是他還年輕的時候和下家使女所生的.
是按著自己意願所生的.
按肉體所生的.
因為上帝說他會有個兒子.
他等不及了.
於是自己想自己的辦法.
就生了以撒瑪利.
第二個兒子是以撒.

$^{721}$是他年紀老邁100歲的時候生的.
是和撒拉生的.
是按著上帝的嬰許生的.
因為他99歲的時候.
上帝差遣使者跟他說明年會抱兒子.
不是抱孫子.
抱兒子.
他在那裡笑.
以撒的名字叫做He Laughs.
他在那裡笑.
按照上帝的嬰許.
按照上帝的旨意所生的.
奴隸不能永遠住在家中.
在講這個故事.
使女的兒子不能承受產業.
在講這個故事.
唯有真正的兒子才可以長住家中.
唯有真正的兒子才可以承受產業.
我們將這個故事引進去約翰福音.
你就會明白耶穌在講什麼.
因為約翰福音一開始的時候.
已經在講兩種不同的兒子.
約翰福音第一章第十節.
他在世界.
世界是藉著他做的.
世界卻不認識他.
他來到自己的地方.
自己的人並不接納他.
第十二節.
凡接納他的人.
就是信他名的人.
他就賜他們權柄作上帝的兒女.
上帝的兒女.
不單單是自由.
《釵及記》在講他會自由.
這裡在講的.
剛剛的故事是講做上帝的兒女.
這些人不是從血生的.
不是從情慾生的.
也不是從人的意願生的.

$^{761}$以撒瑪利就是從人的意願生的.
不算數.
而是從上帝的心.
是憑上帝的應許.
從上帝而生的.
於是你去到這裡.
你明白耶穌在講什麼.
你要問這個問題.
人.
罪人.
墮落的罪人.
從人而所生的墮落罪人.
怎樣才可以從上帝那裡重生.
你明白嗎.
耶穌在講的是一個重生的問題.
你信我就得重生.
你不信我你沒有重生.
信我重生成為上帝的兒女.
於是耶穌說.
你一定要得自由.
因為奴僕是沒有辦法繼承上帝的天國.
於是你必須要重生.
才可以成為上帝的兒女.
我們問這個問題.
人.
如何可以從上帝那裡重生.
成為上帝的兒女.
答案是.
羔羊的寶血.
聖靈的洗禮.
全部都在第一章裡面已經出來了.
第一章.
「約翰看見上帝」第一章二十九節.
「看!上帝的羔羊除去世人的罪」.
你必死在罪中.
我怎樣可以救你出罪.
「上帝的羔羊除去世間的罪」.
你明白嗎.
除非你相信我.
你會死在罪中.

$^{801}$如果我相信你的話.
你怎樣幫我解決這個罪.
我會成為上帝的羔羊.
「上帝的羔羊除去世間的罪」.
接著.
你不可以繼續做一個奴隸.
你要被釋放.
你要被改造成為上帝的兒女.
你才可以繼承.
承繼天國永恆的產業.
第一章第三十三節.
「我先前不認識你.
可是那差我來的用水施洗的對我說.
你看見聖靈降下來.
停在水的身上.
誰就是用聖靈施洗的.
我看見了.
所以作證.
這一位是上帝的兒子」.
如果上帝的兒子令你自由.
誰會令我們自由.
上帝的兒子.
耶穌基督.
耶穌基督怎樣令我們自由.
怎樣可以將我們從奴役變成上帝的兒女.
通過聖靈的洗禮.
你看到嗎.
所有的一切.
其實帶我們去《楊博恩第一章》.
這裡有兩件事情.
兩件事情都和你無關的.
兩件事情都是和我無關的.
兩件事情都是和耶穌有關的.
羔羊的寶血.
聖靈的洗禮.
全部都是救主耶穌的功德和作為.
OK.
那你問.
他做了之後.
我們要做些什麼.

$^{841}$我們必須要做些什麼.
才可以得救.
我們必須要做些什麼.
才可以重生.
他說他.
我說我的.
說真的.
你經常看他那張考卷.
沒有用的.
他考他的.
你考你的.
根本都不是同等級的.
你要問.
上帝告訴我們.
他做完他的那份卷.
他考了一百分.
他可以帶我們進天國.
那我們需要做些什麼.
答案就是信耶穌.
但是問題來了.
什麼叫做信耶穌.
在楊寶欽裡面有太多人信耶穌了.
太多人信了.
第一節信完之後.
下一節就不信了.
有太多人說他信耶穌.
耶穌看著他們卻不信他們.
所以我們必須要問.
什麼叫信耶穌.
不是隨隨便便說一聲.
我信.
而是全人投入.
盡心竭力.
一生一世.
信靠跟隨.
Notitia.
記不記得.
清楚明白福音的核心.
福音的核心是什麼.
高揚的寶血.

$^{881}$聖靈的勢力.
這個很重要.
也是教會經常會忘記了.
你不說高揚的寶血.
你不說十字架.
是沒辦法救人的.
Notitia.
清楚明白.
第二.
心裡要認同.
心裡要確認.
確認什麼.
耶穌給我們的控告.
耶穌說什麼.
你是罪的奴隸.
耶穌說.
你會死在罪中.
如果你不行.
我接納不了耶穌的控告.
我沒辦法心裡確認的話.
這個福音與你無關.
清楚明白福音的核心.
第二 心裡確認福音的控告指控是真實的.
第三.
毀身履行信靠福音的呼召.
福音的呼召是什麼.
捨棄地上世界的偶像.
你不可以留在埃及.
你不可以貪戀世界.
因為耶穌說.
我不屬於世界.
你屬於世界.
耶穌說.
我是天上的.
你是地上的.
你跟不跟我.
你如果要跟我.
你沒辦法留在地上.
你要跟我.
你不可以屬於這個世界.

$^{921}$你要跟我.
你一定要追求我所追求的榮耀.
十字架的榮耀.
你若不愛十字架的榮耀.
我與你有何干呢.
所以清楚明白福音的核心.
心裡確認福音的指控是真實的.
接著我們回轉悔改.
毀身履行福音的呼召.
捨棄地上世界的偶像.
跟隨上帝而至.
信靠高揚的寶血.
追求天上永恆的國度和榮耀.
我們回到第八章三十一,三十二節.
耶穌說這些話的時候.
很多人信了祂.
接著祂繼續.
很可愛.
祂知道他們不是真信.
但是祂繼續說.
先當他們是真信.
接著解釋一下真信的人應該會怎樣.
接著繼續第三十一節.
耶穌對信他的猶太人說.
你們若繼續遵守我的道.
就是真是我的門徒了.
你們將認識真理.
真理會使你們得自由.
這裡說的是一個真信心的實踐和經驗.
一個真信心會怎樣實踐呢.
如果是真的話.
第一你會繼續遵守我的道.
今天沒有時間講.
如果你看到英文是.
Abide in my word.
Abide這個字.
在讓福音重複又重複出來.
我是葡萄樹你是枝子.
你要藏在我裡面.
我都藏在裡面.

$^{961}$那個是Abide.
所以這裡中文翻譯是遵守.
但是其實是藏在我的道裡面.
是親密很多的.
不是遵守OK我們今天怎樣怎樣.
不是你外在的遵守.
是裡面我活在他的道裡面.
我常在這道裡面.
接著你就會認識.
但是我們一見到認識.
我們知道講的是愛慕.
愛慕耶穌的真道.
這個是真信心的實踐.
接著真信心會得到一個經驗.
那個經驗是什麼.
得自由.
行不行.
什麼是真信心.
真信心是怎樣實踐出來.
真信心的實踐就是藏在耶穌的道裡面.
真信心的實踐是認識愛慕真理.
真信心的結果出來的效應.
那個經驗是什麼.
得著自由.
在這裡我們見到兩個希臘文.
一個是Logos.
另一個是Elysium.
兩個關鍵主題同時出現.
真道與真理.
這個組合是罪奴獲得自由的關鍵.
一個罪奴怎可以獲得自由.
需要Logos.
需要上帝耶穌的道.
需要真理.
但是這個組合其實早在《約翰福音》第一章裡面.
已經出來了.
有一個另外一節.
大家都很熟悉的.
道成了肉身.
我幫他改了翻譯.

$^{1001}$將莫我們中間充充滿滿地有恩典.
有真理.
你可以將它翻譯成有慈愛有信實.
真理就是Elysium.
道就是Logos.
這個組合早就出來了.
這個組合可以令我們得自由.
如果你可以這樣記住.
然後他說我們也見過他的榮光.
就是負獨一兒子的榮光.
這一節聖經其實是更加豐富的.
可以說第一章第十四節是整本《約翰福音》最核心的主旨.
我們沒有時間重複看完所有的東西.
但是我和你指出幾個重點.
第一有Logos上帝的道.
耶穌就是道成了肉身住在我們當中.
道成了肉身住在我們.
第二個字.
帳篷在我們當中.
它不是住在我們當中.
它是帳篷.
帳篷和住有什麼分別.
帳篷是會動的.
帳篷的功能就是要引導我們.
帳篷是會動的.
然後有恩典.
有真理.
或者有慈愛.
有信實.
慈愛信實.
舊約裡面的Hesed和Emmet.
是代表了上帝救贖的認許和承傳.
在這個帳篷裡面.
在這個道成了肉身裡面.
我們看到上帝救贖的認許.
救贖的承傳.
在誰的身上.
在副獨一兒子的身上.
所有一切的東西.
都是副獨一兒子耶穌基督的身上.

$^{1041}$所有的榮耀.
如果可以這麼說.
所有的恩典.
真理.
所有的慈愛信實.
全部在耶穌基督最榮耀的地方.
在哪裡.
耶穌基督的一生最榮耀的地方.
在十字架上.
耶穌替我們承傳了救恩.
耶穌也是上帝的榮耀的帳幕.
我再說一次.
帳幕是用來做什麼的.
帳幕不是住的.
帳幕是引領我們的.
耶穌在十字架上承傳了救贖.
然後.
祂成為這個榮耀的帳幕.
要帶我們升空.
離開地上的世界.
追求天上的國度.
這是不能少的東西.
這是自由的定義.
你明白嗎.
什麼叫自由.
自由就是我可以和耶穌一起升空.
所以我這裡畫了.
唯有背十字架的人才可以跟耶穌上飛船.
什麼叫背十字架.
背十字架在聖經新約裡面的其中一個意義.
就是向世界死.
向上帝活.
只有很想升空.
不想再在地球的人.
才可以和耶穌一起上飛機.
做太空人到庫那裡去.
如果你好像以色列人那樣.
貪愛這個世界.
不捨得埃及.
經常回頭的話.

$^{1081}$你是想不到.
耶穌的飛船.
祂的帳幕.
如果我們和祂同死同活.
我們背起十字架.
我們就想和祂一起上飛船.
到庫那裡去.
追求上帝的國和上帝的義.
於是我要問你.
你願意嚮往跟隨上帝的兒子.
獨一的救主耶穌基督.
和祂一起飛天.
一起到庫那裡去.
那裡嗎?.
耶穌說.
如果你不去的話.
你們會死在罪中.
我再想你留心一下.
耶穌怎麼定義罪.
因為我們的定義.
和耶穌的定義是很不一樣的.
你們會死在罪中.
什麼是罪?.
耶穌對他們說.
你們是從下面來的.
我是從上面來的.
下面就是罪的世界.
下面沒有一點一滴不是罪.
你很難明白的.
我們看到這個世界.
我們很難說這個世界.
所有人都是罪人.
我們有很多好人.
我有很多好朋友.
我也是好人.
為什麼你經常說我是罪人.
不.
在上帝的眼光.
你們會死在罪中.
這個罪等同於下面.

$^{1121}$我是上面.
上面沒有罪.
下面有罪.
下面都是罪.
上面一點罪也沒有.
你們是屬世界.
世界是黑暗.
世界是罪孽.
我不屬於這個世界.
我們要問這個問題.
誰是罪的奴隸.
誰會死在罪中.
很多人一想起這兩條問題.
覺得這兩條問題跟他們沒有關係.
誰是罪的奴隸.
誰會死在他自己的罪中.
這些說的是很極端的罪犯.
殺人犯.
罪的奴隸.
連環殺人犯.
行不行.
死在自己的罪中.
強姦犯.
連環強姦犯.
罪的奴隸.
死在罪中.
毒友.
戒幾個戒不掉.
罪的奴隸.
死在罪中.
我們一想起罪的奴隸.
我們就在想這些東西.
賭徒.
切了左手.
切了右手.
都繼續去賭.
上網又賭.
過海又賭.
偷竊上癮的人.
喜歡偷.

$^{1161}$怎樣都要偷.
習慣了.
來教會偷聖經.
色情上癮的人.
我們這些人.
我一想起他們.
罪的奴隸.
死在罪中.
手機上癮.
那就慘了.
個個都是.
上癮的人.
罪的奴隸.
或者你已經是其中一個.
你知道這個故事跟你有關係.
但你說不是的.
我都不是的.
對比他們.
我.
我自由的.
我不會完美.
不過我不是那種.
我自由的.
我正常的.
我陽光的.
我正氣的.
但有沒有想過.
當聖潔上帝從天上.
看下來地下的時候.
這個世界裡面.
一個人.
一個自由人都沒有.
沒有一個是正常的.
沒有一個陽光的.
沒有一個正氣的.
全部都是陰氣.
陰氣很重.
全部都是黑暗的.
每個人都是不正常的.
每個人都是被罪捆綁.

$^{1201}$被世界捆綁.
被地上的榮耀所捆綁.
原來被地上榮耀所捆綁.
都是罪的捆綁.
你明不明白.
有些人很喜歡讀書.
有些人很喜歡考試.
考得高分.
有些人很想升職加薪.
有些人很想七人投隊.
有人很想為家人爭一口氣.
原來這些都是罪的奴役的一部份.
耶穌說.
聖經告訴我們.
地上世界所有的人.
所有.
你想最好的.
你認識的最好的.
你想你自己最好的狀態的時候.
如果他還沒有重生的話.
他們全部都是罪的奴隸.
你會發現在這個故事裡面.
有些人是擺明想殺耶穌.
你說這些人是罪的奴隸.
你說這些人是會死在罪中.
是對的.
但我想你知道.
在約翰福音第七章至第八章裡面.
約翰特意擺了另外一個正氣人在那裡.
這個正氣人.
還要正氣到為耶穌說話.
而約翰想告訴我們.
儘管我將這個正氣人擺在這裡為耶穌說話.
這個正氣人其實是罪的奴隸.
這個正氣人.
其實也是耶穌教導拯救的對象.
這個正氣人是誰呢.
這個正氣人其實是尼哥德姆.
他在第七章的結尾出現過.
唯有他有個名字.

$^{1241}$你可以想像.
當耶穌在說這些.
你會不會死在罪中.
其中一個在後面.
在人群中看著的就是尼哥德姆.
好像薩戈一樣.
不關我的事.
我是好人.
不過又關我事.
他又想聽.
尼哥德姆可以說是約翰福音裡面特選的配角.
最重要的配角.
他在第三章出現過.
在第七章露了露面.
說了一句話.
為耶穌說的.
正氣人.
然後第十九章再出來.
我們可以說.
尼哥德姆比耶穌所有的門徒都更加重要.
我想你真的要很認真去想想尼哥德姆這個人.
在第三章裡面.
我們問誰是黑夜人.
答案是尼哥德姆.
誰進不了天國.
尼哥德姆.
約翰專門幫我們打這個小報告.
有一個人進不了天國.
他不是連環殺人犯.
他是尼哥德姆.
誰需要重生.
因為耶穌對著他.
第一個說話.
You must be born again.
誰很需要重生.
你很不濟.
尼哥德姆.
去到第七章和第八章當中.
我們問.
誰是下面的人.

$^{1281}$耶穌是上面的人.
誰是下面的人.
尼哥德姆.
尼哥德姆是屬於下面的.
其實第三章耶穌已經跟他說.
他是屬於下面的.
誰是屬於世界的.
尼哥德姆.
誰是罪的奴隸.
尼哥德姆.
誰會死在自己的罪中.
尼哥德姆.
這是一個不可思議的發現.
我們也可以說.
這是約翰被他讀者的驚訝.
或者驚喜.
尼哥德姆是一個怎樣的人.
照我們所知.
他是尋找真道.
尋找基督的人.
所以夜裡來找耶穌.
他是一個敬虔公義.
道德高尚.
不是假冒為善的猶太人.
我想.
我這樣說不是錯的.
雖然我沒有一句說話是這樣說.
聖經沒有一句說話.
但是聖經從頭到尾沒有打過他任何一個小報告.
說他是假冒為善的.
他不是的.
他是好人.
第三.
他是以色列的老師.
D-theologian.
以色列的神學家.
他熟讀他的聖經.
他不會和耶穌吵聖經意外的東西.
他就是和耶穌一起翻本聖經去思想.
第四.

$^{1321}$他謙卑受教.
耶穌摑他兩巴.
是不是我摑你兩巴.
你下個星期就不回來了.
耶穌摑了他不知多少巴那一晚.
繼續回來.
默默的跟隨耶穌.
第五.
他公平公正.
他在第七章為耶穌發聲.
我們的律法是這樣的嗎?.
東尼審過他就已經說他是罪不正確的.
接著他被人罰兩下.
這個尼哥德姆.
這樣的尼哥德姆.
不是壞人.
是好人.
這個比我們這裡.
可以說是每一個人.
無人可比的尼哥德姆.
居然是罪的奴隸.
除非他信主.
除非他重生.
否則他會死在罪中.
你可不可以這樣想.
你將尼哥德姆.
他的名字放進去.
罪的奴隸.
死在自己的罪中.
接著你將你自己跟他比一比.
你問.
你需不需要.
做同一樣的動作.
這個發現.
尼哥德姆的出現.
是一個地獄般的壞消息.
同時又是一個天堂般的好消息.
你聽我說.
壞消息是.
如果尼哥德姆是罪的奴隸.

$^{1361}$如果尼哥德姆除非他重生.
不能.
會死在自己的罪中.
我們這裡每一個人.
無論我們.
道德有多高尚.
無論我們多熟悉聖經.
無論我們在教會服侍多少年.
我們都很有可能.
未重生.
你需要面對這件事.
如果他都未重生.
你有沒有可能未重生.
絕對有可能.
這是一個地獄般的壞消息.
但有另一個天堂般的好消息.
林牧師有什麼好消息.
好消息就是.
如果你對你自己的基督徒人生.
感到有些不滿足.
做了這麼多年.
做來做去.
越來越沒有人.
但你又心裡面有一種期盼.
希望有一個更豐盛的人生.
更大的喜樂.
更大的盼望.
更大的滿足.
尼哥德姆的故事告訴我們.
是有的.
有的是什麼.
你需要重生.
如果你有重生.
你就立刻變成一個完全不一樣的人生.
不一樣的世界.
令我想起C.S. Lewis的那句名言.
時間關係.
我直接讀中文的翻譯.
上帝成為人.
要將他的造物改變成為他的兒子.

$^{1401}$不僅是將舊的人改良.
變成一個比較好一點的人.
而是改喚我們成為新人.
有翅膀懂得飛的那種.
這和教一匹馬跳欄.
要它跳得更高.
更好不同.
而是將這匹馬變成能飛的造物.
你看不看到分別.
不是令那匹馬更加好.
是令那匹馬可以長翅膀.
曉飛.
尼哥德姆的重生經驗.
其實就是一個飛行的經驗.
不是他信了主之後.
他有什麼陋習改了他.
而是他信了主之後.
他真的可以最終起飛.
你們會死在最終.
除非你們仰望耶穌.
信他.
怎樣信呢.
《人子被舉起》.
第十九章第三十四節.
士兵拿槍擲他的立旁立刻有血和水流出來.
水代表了聖靈的降臨.
於是有羔羊的寶血.
有聖靈的洗禮.
而這個接下來的就是重生的模樣.
我知道我說了不知道多少次.
但我想沒有人真的接受我這個解釋.
我看了很多書.
大家都不敢說這件事.
因為如果這樣才是重生.
那誰是重生.
連寫的人都不敢說.
所以他不敢寫.
但是你怎麼看.
約翰就是要寫這個故事.
他夠膽寫為什麼.

$^{1441}$因為他也是.
因為約翰曾幾何時都是在追求地上的榮耀.
他現在不再追求地上的榮耀.
他就夠膽寫這個故事.
其他人就不敢.
其他人就恩信稱義.
這個人夠膽寫重生.
尼哥德姆帶著一百斤的抹藥.
尋向用細麻布把耶穌的身體裹好.
這個是重生的模樣.
這個是地獄般的壞消息.
卻又是天堂般的好消息.
原來基督徒的人生是可以這樣活的.
你明不明白.
反過來說.
只有這樣的人生才是基督徒的人生.
其他的就是你自己創造出來的基督徒人生.
不是聖經所說的基督徒人生.
於是我們說.
尼哥德姆全心信靠耶穌.
仰望榮耀的十字架.
他領受了聖靈的施洗.
於是他那天離開了他的黑暗.
他進入光明.
冒險信耶穌.
真的冒險.
接著他帶了一百斤七十五磅值七十五年人工的末藥.
從沉香去敬拜耶穌.
接著他親手埋葬耶穌.
於是他等於以後都不用再做事了.
他要拋棄他這個世界.
他要求上帝的國度.
衝刺.
他還沒做這個動作之前.
你可以想像嗎.
他還沒做這個動作之前.
他是罪的奴隸.
你說什麼罪的奴隸.
他哪有罪.
你用人的角度來說是沒有罪的.

$^{1481}$在上帝的眼光.
在地上的人.
在世界人周圍有罪.
他還沒做這個動作之前.
他是一個罪的奴隸.
他是地上世界的奴隸.
這個人必然死在罪中.
但是現在他自由了.
你看不到嗎.
他這個動作告訴你.
他再不是罪的奴隸.
他自由了.
他成為上帝的愛子.
他不單止自由.
他再不是罪的奴隸.
他成為了上帝的愛子.
他再沒有束縛.
沒有任何的束縛.
他可以單單的為主而活.
你的人生是怎樣的呢.
我不敢問.
你的人生是不是這樣的呢.
因為聖經約翰告訴我們.
這樣的人生才是上帝兒女的人生.
我想帶大家回到很熟悉的聖經.
「世人都犯了罪,規缺了上帝的榮耀」.
什麼是罪?.
規缺了上帝的榮耀.
聖經學者這樣解釋.
所謂規缺了上帝的榮耀.
罪的意思就是missing the mark.
未能達標.
大家聽過了.
未能擊中目標.
我們問目標是什麼.
我們的目標是上帝的榮耀.
通常我們一想到罪就是這樣.
十戒.
或者上帝給我們的道德標準.
但我們經常都是射不到的.

$^{1521}$missing the mark.
但是你有沒有想過.
其實這幅圖畫不是真的.
真正的圖畫不是這樣的.
規缺了上帝的榮耀.
是我們不及上帝的箭靶.
是當上帝呼召我們.
去以祂為我們人生的中心的時候.
我們經常看另外一邊.
另外一邊是世界.
我們每一個人都是自己瞄準自己的箭靶.
我們坐在那裡都傻.
你有沒有想過.
什麼叫世人都犯了罪.
規缺了上帝的榮耀.
因為我們每個人都是背向世界.
背向上帝.
面向世界.
我們自己有自己的箭靶.
難怪我們規缺上帝的榮耀.
因為你都不是對準上帝的榮耀.
你從來沒有想過上帝的榮耀.
你從來沒有用上帝的榮耀的眼光去看人生.
在我人生.
返教會.
慢慢成為基督徒的過程裡面.
其中有一樣東西最.
危害我很深的.
就是有人告訴我們.
我們和上帝的關係可以和上帝寫計劃.
可以向上帝提案.
你想讀什麼書.
告訴上帝.
希望他簽名.
如果你能夠進去.
那就是他簽名.
你想要什麼工作.
你可以和上帝提一個東西.
你說上帝這樣.
然後只要你後面加一句話.

$^{1561}$「隨你的旨意不求我」就可以了.
有沒有人這麼傻呢.
是不是我們這樣做了之後.
然後我找到的工作.
做到的東西就是上帝的旨意.
有沒有人這麼傻.
沒有理由是這樣的.
我們不應該是這樣的.
這樣其實和拜偶像有什麼分別呢.
我和黃大仙一樣.
求了支籤.
大仙我現在申請這份工作.
不知道大仙可不可以高抬貴手.
簽一簽名.
你有什麼分別.
這個不是基督的人生.
基督的人生應該是全心全意.
回轉.
靶在這邊.
你在這邊.
你的靶可以不斷錯失的.
是不是.
我們永遠不會射中心的.
耶穌幫我們射中心的.
你射不到.
但是我對著哪個靶.
我再不是對著這個靶.
我是對著那個靶.
這個就是一個基督徒的人生.
耶穌來將我們的舊靶全部毀掉.
我們現在只有上帝是我們的靶.
虧缺上帝榮耀的反面是什麼.
獻上作活祭.
虧缺神榮耀.
獻上作活祭.
上戀是這樣.
下戀是這樣.
我們一定要成為一張白紙.
然後拿著白紙跟上帝說.
不求我的意願.

$^{1601}$只求你的旨意.
行不行.
不要拿著你自己的計劃.
跟上帝說.
是啊.
都是.
那你有什麼辦法.
上帝批就批.
不批就不批.
你有什麼辦法.
王大仙也是這樣.
沒有人會去王大仙拿著白紙.
大仙.
你告訴我.
我人生應該怎樣.
你也不應該去這麼做.
但是上帝不一樣.
我們結束的時候.
我想給你一個John Calvin的名言.
我們將中文翻出來.
Calvin這麼說.
他說人生最重要最偉大的事就是這件事.
我們被上帝分別為聖.
獻給上帝.
為的是讓我們從今以後.
一切的思想言語默想行為.
一切.
真的很誇張.
都只為祂的榮耀.
我們舊時的東西可以抹掉.
我們被上帝分別為聖.
獻給上帝.
讓我們從此以後.
一切的思想言語默想行為.
全都只為祂的榮耀.
你不會完全的.
你不會完美的.
但是起碼你的靶是射中.
是正確的方向.
然後他說.

$^{1641}$因為屬神之物.
若用於不合神心意的用途上.
如果我們用屬神.
我們已經是屬神的之物.
如果我們用來追求這個世界的話.
必將嚴重損害上帝的榮耀.
求主幫助我們.
你需要一個新的標靶.
轉向上帝.
你怎麼射.
射不到.
射不中.
不完美.
No problem.
總比你以前的人生好得千倍萬倍.
因為你終於瞄向一個正確的靶.
求主讓我們.
將我們的人生獻給他們.
保羅所說的獻上活祭是眾數的.
他不是說一個人獻上活祭.
是整間教會獻上活祭.
求主幫助我們.
讓我們可以為祂作見證.
我們去祈禱.
我們是天父.
一直帶領我們去思想.
耶穌是世界的光.
是生命的活水.
我們都覺得耶穌真是.
是一個很奇妙的救主.
當這班這樣的人.
擺明根本不信祂.
他們說他們信祂的時候.
耶穌竟然這麼有耐性地繼續說.
好 如果你真的信我們就繼續這樣.
耶穌其實是告訴我們.
有一天當我們真的回轉去信靠的時候.
我們應該怎樣去實踐我們這個真的信仰.
我們要藏在祂的話語裡面.
Abide in Him.

$^{1681}$我們要藏在祂裡面.
我們要好像我們的生命.
和祂的生命是二合為一.
祂的導向就是我們的導向.
我們不需要胎盤.
祂是我們的胎盤.
我們要藏在耶穌基督裡面.
我們要認識祂的真道.
我們要愛慕祂的真道.
這樣我們就可以得自由.
那個自由不是單單說道德上的進步.
正如C.S. Lewis說.
不是一匹馬跳得更高.
跨欄跨得更好.
而是說一匹馬曉飛.
主要這個是一個不可想像的東西.
我們從來沒想過罪的奴役.
是將我們困在地球上,世界上的東西.
就算我們很有道德,很好人.
很有慈善家的風範.
被世界景仰讚美.
如果我們的人生就被鋸在地上.
飛不起,沒辦法做天國奇妙不可思議的東西的時候.
原來我們一直都是罪的奴役,罪的奴隸.
更加嚴重的是.
這等人必死在罪中.
《尼哥德姆》是一個不可思議的故事.
因為一個這樣的人仍然需要重生.
但一個這樣的人仍然可以重生.
而當他重生的時候.
他真是改頭換面的.
一個已經這麼好的人.
有時候我們會說你的人都已經這麼好.
你做了基督徒我想都是加多兩個小時的教會而已.
不是的.
原來我們的人生是需要被大大的改造.
而這個改造已經再不是說道德倫理上的改造.
而是究竟我們的箭靶的方向是在哪裡.
這個是我們這個時代需要聽的福音.
求主你憐憫我們的教會.

$^{1721}$讓我們去思想耶穌基督的呼召.
真正順靠祂改造我們的人生.
我們這樣祈禱.
奉祖耶穌基督明求 阿們.
\newpage



\section{歷代志上 22:1-13}
\label{sec:032298}
\textbf{上帝的居所:從聖殿到子民}~~~~林于謙宣教士
\newline
\newline
連結: \href{http://www.church.com.hk/acms/content.asp?site=cdc\&op=show\&type=product\&code=032298\&layout=sermon}{\texttt{032298}} ~~~~ 語音日期: 2026-06-14
\newline
\newline
\hyperref[sec:032297]{\small{< < < PREV SERMON < < <}}
~
\hyperref[sec:index_chronic]{\small{[返順時目]}}
~
\hyperref[sec:032299]{\small{> > > NEXT SERMON > > >}}
\newline
\newline
歷代志上 22:1-13
\newline
\begin{longtable}{cl}
\hline
\hline
章節 & 經文 (和合本修訂版)\\
\hline
22:1 & \begin{tabularx}{0.7\textwidth}{X} 大衛說：「這是耶和華神的殿，這是以色列獻燔祭的壇。」 \end{tabularx} \\ \\ \relax
22:2 & \begin{tabularx}{0.7\textwidth}{X} 大衛吩咐人召集住以色列地的寄居者，又派石匠鑿石頭，要建造神的殿。 \end{tabularx} \\ \\ \relax
22:3 & \begin{tabularx}{0.7\textwidth}{X} 大衛預備許多鐵，要做門上的釘子和鉤子，又預備許多銅，多得無法可秤； \end{tabularx} \\ \\ \relax
22:4 & \begin{tabularx}{0.7\textwidth}{X} 還有無數的香柏木，因為西頓人和推羅人給大衛運了許多香柏木來。 \end{tabularx} \\ \\ \relax
22:5 & \begin{tabularx}{0.7\textwidth}{X} 大衛說：「我兒子所羅門還年幼脆弱，要為耶和華建造的殿宇必須高大輝煌，使名聲榮耀傳遍萬國，所以我要為殿預備。」於是，大衛在未死之前預備了許多材料。 \end{tabularx} \\ \\ \relax
22:6 & \begin{tabularx}{0.7\textwidth}{X} 大衛召了他兒子所羅門來，吩咐他為耶和華－以色列的神建造殿宇。 \end{tabularx} \\ \\ \relax
22:7 & \begin{tabularx}{0.7\textwidth}{X} 大衛對所羅門說：「我兒啊，我心裡本想為耶和華－我神的名建造殿宇， \end{tabularx} \\ \\ \relax
22:8 & \begin{tabularx}{0.7\textwidth}{X} 可是耶和華的話臨到我說：『你流了許多的血，打了多次大仗；你不可為我的名建造殿宇，因為你在我面前使許多血流在地上。 \end{tabularx} \\ \\ \relax
22:9 & \begin{tabularx}{0.7\textwidth}{X} 看哪，你要生一個兒子，他必成為安寧的人；我必使他得享安寧，不被四圍仇敵擾亂。他的名字要叫所羅門，在他的日子，我必使以色列平安康泰。 \end{tabularx} \\ \\ \relax
22:10 & \begin{tabularx}{0.7\textwidth}{X} 他必為我的名建造殿宇。他要作我的子，我要作他的父。我必堅定他國度的王位，使他治理以色列，直到永遠。』 \end{tabularx} \\ \\ \relax
22:11 & \begin{tabularx}{0.7\textwidth}{X} 我兒啊，現今願耶和華與你同在，使你亨通，建造耶和華－你神的殿，正如他指著你所說的。 \end{tabularx} \\ \\ \relax
22:12 & \begin{tabularx}{0.7\textwidth}{X} 但願耶和華賜你聰明智慧，好按著他吩咐你的去治理以色列，遵行耶和華－你神的律法。 \end{tabularx} \\ \\ \relax
22:13 & \begin{tabularx}{0.7\textwidth}{X} 那時候，你若謹守遵行耶和華藉摩西吩咐以色列的律例典章，就得亨通。你當剛強壯膽，不要懼怕，也不要驚惶。 \end{tabularx} \\ \\
[1ex]
\hline
\hline
\end{longtable}
$^{1}$各位弟兄姊妹.
主佑平安.
我今天要分享的題目就是.
上帝的居所.
從聖殿到子民.
不知道大家想起上帝的居所.
會想起什麼呢?.
我小時候有看《孫悟空》.
我的年代看的《孫悟空》是張偉健的時代.
我不知道現在應該翻拍了很多新版本.
但是我記得.
有一幕我記得很深刻的.
就是《孫悟空大鬧天宮》.
可能這一幕還在我腦海裡.
我經常想起.
上帝的居所.
神住在哪裡.
可能有時候就受這個影響.
覺得神在天上有個天宮.
祂在上面.
我們在下面.
我記得《孫悟空》.
我很喜歡那個唐三藏.
每次《孫悟空》搗亂的時候.
他念經.
《孫悟空》就會頭痛.
就會變得很乖.
有時候我也想有這個超能力.
有些小朋友跑來跑去發瘋的時候.
念經他就乖乖坐在那裡.
回到我們的經文裡面.
《歷代至上》第二十二章一至十三節.
先說說自己的背景.
我爸爸是美國人.
我媽媽是菲律賓人.
所以我的背景有一點.
跟在座可能很不同.
但是一樣東西應該相同.
就是小學和中學.
都是在香港.

$^{41}$我在隔壁出生.
小學和中學都是在這裡讀書的.
我四年前去了柬埔寨.
在那邊宣教.
剛剛回到香港.
兩個禮拜多一點.
我現在太太是三十九州華人.
她今天也在這裡.
可不可以站起來.
她應該是禮拜二就生孩子了.
所以.
如果現在生你揮手.
我早點完.
我趕快一點.
OK.
好.
我們先開始.
我請各位同學一起朗讀今天的經文.
一二三.
大衛說.
這是耶和華神的殿.
這是以色列建制主義的殿.
大衛吩咐人.
要集合以色列地居者.
又派石像石頭.
要建實行的殿.
大衛準備裹帖.
又做門上的口字.
他預備許多堂.
多得無法可停.
並有無數的賓客來.
因為西頓人和吐羅人.
大衛說.
我兒子所留門.
貧民空出.
以為耶和華建制主義.
必須擴大安樂.
使民生物榮耀傳遍各國.
所以為了慰勞眾人.
於是大衛在臨死之前.

$^{81}$準備了許多材料.
大衛千里遍野以所留門.
吩咐他為以色列地居者建造殿宇.
大衛對所留門說.
我兒啊.
我心裡本想為耶和華實在未建造殿宇.
可想而知大衛講了那麼多話說.
比我們每個大衛.
大衛的好事大得多.
為甚麼大衛不願意建造殿宇.
因為建造殿宇是最多人能做的事.
安娜.
我要生一個兒子.
他必成為安寧的人.
我必使他得到安寧.
不必死為受敵擾亂.
他的名字要為所愛我.
在他的日子.
我必使以色列平安安泰.
他的愛我的名字就是殿宇.
大衛作我的字.
我要作他的父.
我必奠定他的國度皇位.
使他使以色列直到永遠.
我兒啊.
願公願耶和華與你同坐.
使你健康.
願祝耶和華你神的健.
積聚他至在理所.
大願耶和華使你聰明智慧.
我奉請他吩咐你.
要建立以色列.
宣揚耶和華與以色列的合法.
那時候.
你若謹守宣宣耶和華.
俗外科的豐富儀式的基督教儀式.
就獲得安康.
你安安穩穩.
不必懼怕.
也不必驚慌.

$^{121}$多謝各位.
今天這個故事裡面主要有兩個角色.
大衛和所羅門.
不知大家今個星期看新聞.
知不知道這個是誰?.
安東尼高帝.
大家可能未必認得他的樣子.
因為這幅相片也非常舊.
但可能認得出他的成就.
這個星期.
巴塞隆拿的聖加堂.
經歷了一百四十四年.
終於封了頂.
有個十字架放在上面.
高帝在一百八十三年動功.
不幸地他在一九二六年離世.
他生前未能親眼看見教會的完成.
他的子孫.
甚至子孫的子孫.
都未能見證他人的願望.
跨越了多個世代.
才今天得到實現.
你看到相片裡面.
一九零五一九零七.
還是一個很小的塔.
這個就是高帝用了一百四十四年前的意象.
而我相信他一百八十三年動功的時候.
有沒有人一百八十三年在這裡.
應該沒有人這麼老吧.
有的我都怕.
我們見到大衛今天這個故事裡面.
大衛應該都很相信.
他清楚或者他明白高帝這種感受.
因為大衛領受了建造聖殿的意象.
他準備了所有材料.
他聚集了資金.
他組織了工人.
但是大衛最後他不可以親眼去見到聖殿員工.
相反地他就告訴他的兒子所羅門斯.
不可以啊.

$^{161}$你放心去建造上帝放在你心裡面的那座殿.
我相信上帝最偉大的工作.
有時因為太過宏大.
沒有辦法由一個世代去完成.
在聖經裡面.
我們見到亞伯拉罕.
他從來未見到以色列成為一個國家.
我們看看摩西.
他亦都不可以進入應許之地.
大衛他亦都未見到聖殿的員工.
好像在裡面的高帝.
他見不到聖家堂的完成.
大衛和高帝.
他們都明白一件我們這個世代經常忘記的事.
就是忠心與否.
不是取決於我們是否完成了那個工作.
反而是取決於我們在上帝宏大的故事裡面.
有沒有忠心去扮演自己那份角色.
有時候忠心就是準備到老.
我們今天看看一個關於神居所的故事.
我們看看以色列歷史.
我們今天的故事歷代至第22章的背景是甚麼.
在第10章裡面.
大衛登基作王.
10章的下半段.
他就攻取了耶路撒冷.
在第15至16章.
大衛就將一個藥櫃搬進了耶路撒冷.
在第17章裡面.
這裡就是大衛有意象去建一個聖殿.
在第17章裡面.
大衛剛剛入房他自己的新屋.
他和先知拿丹說.
你看看我現在住在香柏木的棺中.
但是耶和華的藥櫃.
反而在一個300至400年的萬寺裡面.
所以大衛心裡面.
就很想為神去建這個殿.
第18至20章.
記述了大衛一連串軍事的勝利.

$^{201}$他擴張了以色列的邊界.
在第21章我們見到.
大衛為建造聖殿的那塊地他買了.
而第22至29章.
亦都是今天的經文所記載.
就是大衛為聖殿所準備的各種材料.
在下一本書《歷代志》下2至7章.
我們就見到他的兒子.
所羅門所建造的聖殿.
所以《歷代志》第22章.
已經是記載大衛的晚年.
準備去傳承王位的事.
他已經是一位連萬的君王.
但是他心裡面還有一個夢想.
我們見到今天的經文裡面.
第1至4章我們見到.
大衛他用愛邦人的材料.
愛邦人的勞工.
因為聖殿不是單純是給猶太人.
而是為了萬國萬民.
神起初就有這個心意.
在第5章裡面.
大衛就說所羅門年幼脆弱.
在英文年幼脆弱就寫.
Young and experienced.
又年輕又沒有經驗.
所以他就說我要準備.
準備到多少都要做到最好.
上帝跟大衛說.
你不可以做起建的人.
但是大衛沒有因為這樣去生氣上帝.
他心裡面可能在想.
起就不讓我起.
但是上帝你沒有說過.
我不可以準備材料.
你也沒有說.
我不可以去準備買那塊地.
你也沒有說.
我不可以去準備工人去建這棟建築物.
所以大衛就準備了.

$^{241}$他可以準備的都做到最好.
在第8章.
我們知道.
為什麼大衛不可以去建這個聖殿.
上帝就透過先知說.
你流了很多血.
打了多次的大仗.
你不可以為我的名去建造殿宇.
因為你在我面前.
使許多血流在地上.
聖殿是一個和平的殿.
所以需要一個帶來和平的人.
所羅門(Solomon).
他的名字本身就和平安(Shalom)相關.
他就是這位和平之子.
和平的人.
歷代自稱他為得享安息的人.
他生出來不是來打仗的.
而是為了完成上帝給他的意象.
我們看到大衛召集了所羅門.
把建造殿的心願和神的應許.
託付給他.
他就叮囑所羅門.
你要肝強,忠心.
因為真正的成功在於馴服神.
大衛把所有的材料和工匠交給所羅門.
讓他知道所羅門不是孤單的.
聖殿真正的根基不是金,銀或石頭.
而是一火向神擺上的心.
大衛最後的吩咐就是.
遵行神的命令.
高過一切.
這是一個父親可以留給孩子.
最大的傳承.
忠心去跟從神.
作合神心意的領袖.
大衛的意象他準備了很多東西.
土地,資源,一切事務.
甚至他也有跟其他領袖說.
因為一個偉大的領袖所羅門是不夠的.

$^{281}$要把所羅門周圍的領袖都召起來包圍.
幫他完成這件事.
最後也透過軍事勝利為國家帶來和平.
這些根基就讓所羅門之後可以順利建造聖殿.
我不知道我們有沒有留意到.
或者知不知道原來大衛在背後做了這麼多.
很多時候我們把所羅門的殿.
大家都聽過所羅門的殿.
沒有人聽過大衛的殿.
但是原來他背後做了這麼多.
你覺得大衛介不介意呢?.
我們其實在聖經裡面看到.
大衛是不介意的.
他從來沒有介意這個聖殿沒有自己的名字.
他不在乎功勞.
他不在乎榮耀.
只要看到神的殿被建造.
就算自己完全隱身在幕後都沒有關係.
有一個原則對我們來說很重要.
如果我們甘願默默無聞.
只做神交給我們的事.
不在乎誰得到明星.
只要神的工作得到完成.
祂的家園被建造.
我相信上帝都對我們很滿足.
所羅門就負責興建.
而聖殿就代表了神的同在.
但是在將聖殿獻給神那一天.
所羅門都明白一件事.
他心裡面問.
神果然真的住在地上嗎?.
我們知道聖殿的建造確實是有意義.
但是沒有一個殿譽能夠真正勝在神的同在.
所羅門他什麼優勢都有.
所有資源,所有需要的事情都有.
但是他最後都沒有一個完美的結局.
今天想在這個時刻轉一轉話題.
大家今天走過來的時候.
有沒有留意走過多少路牌呢?.
可能大家不是第一次回教會.

$^{321}$有時候自行駕駛.
就這樣走回教會.
都沒有留意身邊有沒有什麼路牌.
在網上看到這幾幅照片.
有時候路牌都很誤導性的.
左邊這個禁止停車.
但是又寫著停車場.
旁邊幸福通道.
但是嚴禁入內.
另外有一幅照片有個老師問一個學生.
你為什麼今天遲到?.
那個學生就說因為路上有標誌.
那個老師就問.
那個標誌跟你遲到有什麼關係呢?.
他說看到這個.
前面寫著前方學校減速慢行.
所以他遲到.
有時候我上班走過M記.
可能買個包.
我自己上班都遲到.
好像這個小朋友那樣.
我為什麼說路標呢?.
就是因為我有時候在柬埔寨教神學的時候.
柬埔寨有時候沒有一些詞語.
例如舊約去預表新約的東西.
沒有專有typology這樣的詞語.
所以在英國一個著名的聖經學家和主教.
Antirite拉特.
他就用了一個比喻.
就說很多時候舊約.
好像一個向前展望的路牌.
或者一個路標.
我們在舊約裡面.
看到以撒.
在亞伯拉罕犧牲了他的獨生子.
就好像新約裡面.
上帝準備了他的獨生子耶穌.
若撒就是被賣為奴隸.
但是最後成為拯救者.
好像耶穌那樣.

$^{361}$大衛也是指向將來的孟王之王.
聖經裡面也有事件.
好像如意節的羔羊.
指向一個路牌.
指向新約.
神的寶血.
過紅海就是神在埃及.
拯救我們出一個綑綁我們的地方.
有時即訓例如先知.
祭司.
君王.
都是路牌.
所以我經常教書.
很喜歡用這幅相.
這幅相就連著了.
聖經裡面很多的金句.
例如.
舊約有一個金句.
和新約有關係.
有一個研究家畫了一條線.
將聖經的故事連在一起.
最後你看到聖經.
我們很多時候看聖經.
只看到小的故事.
但是你看不看到大的故事.
也指向有連結性.
我們看到左邊是創世記.
右邊是啟示錄.
下面黃色的就是今天的經文.
詩篇119就是那條很長的線.
因為每條線都和經文的長度有正比.
我希望大家有這個概念.
對一個路牌.
或者聖經的連貫性有一個概念之後.
我們去看看今天這個路牌.
就是上帝的居所.
在聖經裡面描述了.
上帝學武和他的子民同住的經歷.
上帝的居所除了聖殿.
大家有沒有想到.

$^{401}$在聖經裡面有沒有什麼故事.
也被描述為上帝的居所.
早在聖殿被規劃.
或者被建造之前.
就已經開始了.
第一個就是上帝起初的願望.
就是在伊甸園.
住在人中間.
《創世記》一至三章說.
伊甸園被描述是地上第一個聖所或者聖殿.
上帝和亞當夏娃同行.
但是人就因為墮落犯罪.
失去了與神的同在.
在這個時間.
關鍵問題就是.
上帝如何重新和他的子民同住呢?.
我們向前看幾百年.
就見到在會幕裡面上帝居住在以色列民中間.
會幕.
上帝就居住在以色列中間.
他就說.
他們要為我建造聖所.
使我可以在他們中間居住.
會幕用了一段時間.
今天的故事就是大衛看到的會幕.
已經三百五十至四百年.
他說我要建一個新的聖殿給上帝.
今天的故事就是大衛的意象.
和所羅門的聖殿.
大衛準備所羅門興建.
聖殿的意義在於神的同在.
耶和華的榮光充滿了他的殿.
聖殿就成為一個永久的象徵.
象徵神就住在這裡.
神和以色列同在.
神的心意就是當萬國抬頭.
見到以色列的時候.
他們會見到和明白一件事.
神真的住在這些人中間.
舊約有兩個故事.

$^{441}$清清楚楚說明了這一點.
按猶太人的理解.
舊約裡面是萬國要來到我們以色列這邊.
其中一個故事.
萊曼是一個愛邦人.
他有皮膚有疾病.
他去以色列尋求以利沙先知的醫治.
因為他聽到他的僕女跟他說.
以色列的神可以醫治他.
以利沙跟他說.
你去那條河裡面浸七次.
他一開始不想.
但最後都拒絕了.
他浸了一次.
兩次.
三次.
四次.
五次.
六次.
第七次上來.
聖經說他的皮膚好像嬰兒一樣.
不用SK-II或任何東西.
他出來完完全全的醫治.
他就說了一句.
我知道除了以色列.
全地沒有別的神.
所羅門的時代.
另外一個斯巴女王.
他原本不相信所羅門所聞的故事.
他說我本來不相信這些話.
但我來了之後.
我親眼看見了.
我以前聽的一半.
都不及我現在親眼所見到.
聖殿的建造就是以色列歷史的高峰.
很可惜地.
這個不可以維持很久.
好像伊甸園一樣.
人選擇了自己的路.
領袖偏行己路.

$^{481}$結局就是分裂和毀滅.
今天的故事.
今天的經文.
歷代字.
最初讀者的對象.
就是剛剛一班在巴比倫.
被擄了之後.
返回去的猶太人.
他們這個時候.
就是想重建他們的國家.
回復他們的民族身份認同.
重建他們對神的敬拜.
試試想像這班人走了七十年.
回來見到自己耶路撒冷.
聖殿不在.
重建了都沒有以前的規模.
甘碧輝煌.
他們心裡面有很多問題.
其實我們還是不是神的子民呢?.
神還是不是住在我們中間?.
聖殿現在又扮演什麼角色呢?.
神對大衛的後代.
對大衛的王朝.
還有什麼計劃?.
當我們向前看.
在《約翰福音》就解答了這些問題.
耶穌就成為真正的聖殿.
《約翰福音》第一章十四節寫.
祂獨渡成了肉身.
住在我們中間.
居住這個詞語.
本身有露營的意思.
指的就是住在帳篷裡面.
或者是紮營.
我不知道在座各位.
有沒有人喜歡露營呢?.
耶穌真正的聖殿.
在之後的經文第二章.
耶穌回答對他們說.
你們拆毀這殿.

$^{521}$我三天之內將它重建.
猶太人就說.
用了四十六年去建.
你三天真的可以重建?.
但是門徒就知道.
耶穌所說的殿就是祂的身體.
在這個故事的背景.
其實我們知道.
也是耶穌一個很有趣的故事.
因為耶穌去了愛邦人院外面.
見到動物聲.
見到那些人買賣.
整個市集很亂.
祂知道愛邦人院.
是應該給愛邦人來敬拜禱告的地方.
但是現在人家做買賣.
聖殿成為了嘈雜的市場.
祭司逃離.
窮人被剝削.
耶穌就憤怒.
祂不是一個暴力的憤怒.
而是義憤.
祂拿起鞭子.
趕走那些換錢的人.
將父的殿徹底弄乾淨.
因為聖殿本來就應該是萬國禱告的殿.
耶穌對父神的殿.
猶太人的聖殿的純潔.
充滿了內心的火熱.
就是因為這個故事.
祂生氣.
可能大家都不知道.
祂除了反抬之外.
祂還自己織了一條鞭子.
去趕走那些人.
我有一個朋友很喜歡這個故事.
因為他說這是我最喜歡的耶穌.
拿一條鞭子出來和祂反抬.
他說如果耶穌都這樣.
我也可以這樣.

$^{561}$我就跟他說.
但是耶穌也是一次而已.
你天天這樣不是很好.
我們再向前看.
耶穌是真正的聖殿.
耶穌親自就宣告.
祂自己是聖殿.
神的同在.
不是在一個建築物裡面.
而是在一個人身上.
所以在這個時刻.
我們看到.
大衛和所羅門所做的一切建造工作.
都是指向將來的基督.
耶穌死了之後.
祂就將使命交託給門徒.
祂的同在就要住在我們每一個人裡面.
祂要我們每一個人.
成為一個小聖殿.
神同在的群體.
這個就是教會的意象.
不是一個建築物.
而是一群人.
彼得前書兩章五節說.
「你們作為活石要被建造成屬靈的殿」.
教會如今.
就成為神與人同住的地方.
以弗所書兩章二十二節說.
「你們也靠祂同被建造.
成為神藉著聖靈居住的所在」.
所以.
神就將門徒建在一起.
聖靈在神指紋當中.
教會不是一個人去參加的場所.
而教會就成為神居住的家.
當我們在最後的路標.
或者最後一個終點.
就是新天新地.
神最終的居所.
我們在啟示錄二十一至二十二章裡面看到.

$^{601}$結局回到起初的設計.
啟示錄二十一章三節和二十二節說.
「看吶 神的帳篷在人間」.
啟示錄的作者說.
「我看不到 我沒有看見城內有殿」.
「因為神親自臨在」.
「整個受造之物就成為了祂的居所」.
聖經是結束.
不是一個新石頭的聖殿.
而是一個全新的宇宙.
一個宇宙性的聖殿.
將來不需要特定的建築.
因為藉著耶穌.
整個世界都要成為.
神和祂子民同住同治理的地方.
我們看舊約和新約.
對比在舊約.
萬國都要來到神的殿裡面.
今天神就差遣祂的子民去萬國.
我們在故事裡面.
如果我們想想我們的角色像不像大衛.
神是否呼召我們做一個預備者.
你會留下什麼給下一代去傳承呢?.
可能你是一個鎖門.
一個建造者.
有時我們將上一代人所付出的犧牲.
視為理所當然.
我們記不記得.
為了讓我們今天.
可以享受今天所擁有的一切.
上一代的人付出了多少.
我自己都不認識整群創辦人.
甚至警察會的創辦人.
八九年是誰我都不知道.
但是很多人默默打下了根基.
我們今天才可以在這個根基上.
繼續去建造.
我太太Nana.
她爸爸七十年代偷渡來香港.
她試了三次.

$^{641}$她第一次被人抓了.
被人抓回去.
第二次也被人抓了.
被人抓回去.
她心裡面想.
我再試一次.
看看能不能過去香港.
幸好她們三次來到香港.
就是因為她七十年代來到香港.
我太太的家庭.
她哥哥和她妹妹和她媽媽.
都現在在香港這邊.
我爸爸也是八十年代來香港.
他要住滿七年才拿到香港身份證.
我記得我和我太太結婚的時候.
我太太的爸爸和我爸爸.
美國人.
他就問我爸爸.
你是不是也偷渡來香港的?.
我爸爸就說.
不是 我八十年代是打工來香港的.
他不明白.
他可能以為全部人都是這樣過來的.
所以我們看今天的故事.
聖經就記述了.
上帝渴望他子民同住的經歷.
由伊甸園到新天新地.
我們也可以這樣說.
聖經是一個關於神同在,失落,恢復.
和最終達到完美的故事.
如果我們方言去看一個完整的聖經故事.
不只是單看每一個小故事.
我們看到神在伊甸園.
和亞當夏娃同行.
之後神藉著會幕和聖殿.
住在以色列人中間.
之後神度成肉身.
在基督裡住在我們中間.
神如今就藉著聖靈住在他的教會中.
而將來有一天.

$^{681}$各國各族的人會來到萬民.
將要永遠和他同在.
有一天這個世界不會再需要宣教士.
也不需要傳道人.
因為大家都是敬拜神.
美國有個牧師約翰·派博.
他有這一句.
宣教之所以存在.
因為敬拜不存在.
我這四年在柬埔寨服侍.
因為柬埔寨百分之九十五的人.
都從來沒聽過福音.
不知道耶穌的喜訊.
所以我去宣教.
就是因為大部分的人都不敬拜上帝.
如果有一天柬埔寨百分之百的人.
都敬拜上帝.
那不需要宣教士去傳福音.
我們今天之所以在這裡.
都是因為百多二百年前.
有宣教士來香港.
我們領受了神的祝福.
所以我們今天都要成為祝福其他人.
宣教之所以存在.
就是因為敬拜不存在.
總結今天.
大衛準備了一座.
他永遠無法親眼見到的聖殿.
所羅門建造了一個.
最終都是消失的聖殿.
耶穌今天成為真正的聖殿.
而今天聖靈住在我們每一個人裡面.
我們被呼召要走出去.
將神的同在帶到萬國.
我去柬埔寨不是帶著一個建築物.
而是帶著神的同在去.
要總結今天.
雖然有很多點.
但只有三點我想大家記得.
第一.

$^{721}$神是否呼召你.
或者在你現在這個時刻.
你是否一個準備者.
要準備神怎樣去傳承下一代.
可能你是一個所羅門.
是一個建造者.
神有什麼任務擺在你面前.
和尚未完成的故事.
我們見不見到聖經故事的連貫性.
由伊甸園 會幕 聖殿 耶穌.
現在教會我們每一位弟兄姊妹.
帶著神的同在.
這個任務未完.
直至到新天新地.
我們怎樣可以成為神工作.
幫他完成他這個故事.
大衛和所羅門所做的一切.
的準備和他們建造的工作.
最後都是指向耶穌基督.
問問在座各位.
你的工作和你的生活.
最終是指去那裡.
當我們思考神的居所的時候.
重要的問題不是我們會留下一個怎樣的建築物.
更重要的應該是問自己.
我們會成為一個怎樣的一群人.
和下一代能不能夠透過我們.
經歷到神的同在.
弟兄姊妹.
大衛和所羅門的故事.
都是指向一個更大的現實.
就是耶穌基督.
大衛渴望去建這個殿.
所羅門最終建了這個聖殿.
都是預告著那位真正的大衛子孫.
和平君王耶穌的降臨.
所羅門用石頭去建殿.
耶穌就用你和我的心靈和我們的生命.
去建造教會成為神的居所.
大衛為兒子準備一切.

$^{761}$正如父神為耶穌準備了道路.
真正的聖殿不是用石頭去建築.
而是一個跨越世代.
萬國因為信基督合一的大家庭.
真正的傳承不是紀念碑.
而是一些被耶穌所愛的指紋.
大衛最終沒有親手去建成這個聖殿.
但大衛是不是一個失敗者?.
我相信大衛不是一個失敗的人.
他完成了神所容許他做的一切.
弟兄姊妹.
我們的人生也是這樣.
如果神呼召你去做一件事情.
這就是你需要完成的事情.
大衛投入一個自己看不見的現實和應許.
為下一代打下根基.
今天為神建造神的夢想.
要藉著你和我延續下去.
願我們可以像大衛一樣.
衷心完成神給我們的工作.
將信心傳給下一代.
最後總結想用一個故事去總結.
我爸爸八十年代來香港.
他其實很想去創啟地方去宣教.
但他在一個浸水會教會結了婚.
我們幾個小朋友要照顧我們.
一直留在香港.
他原本心裡想去宣教的地方.
他自己去不到.
他沒有想過我們幾個小朋友在香港大學中文.
學不同語言.
竟然二十年之後.
他想去宣教的地方.
我們完成了他的心願.
去到他想去的地方.
他也跟我說.
原來神的信實.
神的應許.
最後也是得到Full Fulfillment.
所以我希望藉著這個故事.

$^{801}$藉著大衛,所羅門和未完成的故事.
我們大家可以想想.
神有什麼任務給我們每一位.
最後我們一起禱告.
大衛做一個預備者.
和所羅門做一個建造者.
我們知道我們在你裡面.
我們都是你的活石.
你想將我們建造一個新的點.
求你讓我們每一個人可以體驗你.
知道你怎樣去呼召我們.
去完成你在我們心裡每一個人的任務.
願你保守我們在座每一位活石.
我們以上禱告.
奉主耶穌基督的名祈求.
\newpage



\section{以弗所書 2:11-22}
\label{sec:032299}
\textbf{合一和睦ge 代價唔簡單}~~~~鄺偉志牧師
\newline
\newline
連結: \href{http://www.church.com.hk/acms/content.asp?site=cdc\&op=show\&type=product\&code=032299\&layout=sermon}{\texttt{032299}} ~~~~ 語音日期: 2026-06-14
\newline
\newline
\hyperref[sec:032298]{\small{< < < PREV SERMON < < <}}
~
\hyperref[sec:index_chronic]{\small{[返順時目]}}
~
\hyperref[sec:032300]{\small{> > > NEXT SERMON > > >}}
\newline
\newline
以弗所書 2:11-22
\newline
\begin{longtable}{cl}
\hline
\hline
章節 & 經文 (和合本修訂版)\\
\hline
2:11 & \begin{tabularx}{0.7\textwidth}{X} 所以，你們要記得：從前你們按肉體是外邦人，是「沒受割禮的」；這名字是那些憑人手在肉身上「受割禮的人」所取的。 \end{tabularx} \\ \\ \relax
2:12 & \begin{tabularx}{0.7\textwidth}{X} 要記得那時候，你們與基督無關，與以色列選民團體隔絕，在所應許的約上是局外人，而且在世上沒有指望，沒有神。 \end{tabularx} \\ \\ \relax
2:13 & \begin{tabularx}{0.7\textwidth}{X} 從前你們是遠離神的人，如今卻在基督耶穌裡，靠著他的血，已經得以親近了。 \end{tabularx} \\ \\ \relax
2:14 & \begin{tabularx}{0.7\textwidth}{X} 因為他自己是我們的和平，使雙方合而為一，拆毀了中間隔絕的牆，而且以自己的身體終止了冤仇， \end{tabularx} \\ \\ \relax
2:15 & \begin{tabularx}{0.7\textwidth}{X} 廢掉那記在律法上的規條，為要使兩方藉著自己造成一個新人，促成了和平； \end{tabularx} \\ \\ \relax
2:16 & \begin{tabularx}{0.7\textwidth}{X} 既在十字架上消滅了冤仇，就藉這十字架使雙方歸為一體，與神和好， \end{tabularx} \\ \\ \relax
2:17 & \begin{tabularx}{0.7\textwidth}{X} 並且來傳和平的福音給你們遠處的人，也傳和平給那些近處的人， \end{tabularx} \\ \\ \relax
2:18 & \begin{tabularx}{0.7\textwidth}{X} 因為我們雙方藉著他，在同一位聖靈裡得以進到父面前。 \end{tabularx} \\ \\ \relax
2:19 & \begin{tabularx}{0.7\textwidth}{X} 這樣，你們不再是外人或客旅，是與聖徒同國，是神家裡的人了， \end{tabularx} \\ \\ \relax
2:20 & \begin{tabularx}{0.7\textwidth}{X} 被建造在使徒和先知的根基上，而基督耶穌自己為房角石， \end{tabularx} \\ \\ \relax
2:21 & \begin{tabularx}{0.7\textwidth}{X} 靠著他整座房子連接得緊湊，漸漸成為在主裡的聖殿。 \end{tabularx} \\ \\ \relax
2:22 & \begin{tabularx}{0.7\textwidth}{X} 你們也靠他同被建造，成為神藉著聖靈居住的所在。 \end{tabularx} \\ \\
[1ex]
\hline
\hline
\end{longtable}
$^{1}$各位早晨,我們一起祈禱.
藉著思想主義華語,聖靈開導我們有效的敬拜.
因為我們更加明白你,我們更加敬畏你.
更加敬拜你,這樣來服侍你.
主,求你再一次聖靈幫助我講的孩子.
可以將你的華語講得明白.
也幫助我們姐妹用心聆聽.
主,求你聖靈在我們中間,因為你就是那位那麼奇妙.
使我們聖靈在我們中間,冠乎在我們當中的這位奇妙的上帝.
祈禱,奉耶穌基督名志,阿們.
好,合一和睦的代價不簡單.
其實我們為什麼要講合一的呢,教會.
我們是怎樣講合一的呢.
有時候我們講合一,好像說到大家和睦共處.
和和氣氣,大家沒有意見,沒有爭執是最好的.
這些就叫做合一.
但是想一想,是不是我們要追求這種合一呢.
這種才叫做合一呢.
有時候我們的經驗,合一都未必是這樣才是一個健康的合一.
很簡單,我舉出兩點.
第一,有時候一個合一的群體,大家想一想.
無論是一個群體,一個組織,或者一個家庭.
是不是只有一個意見就叫做合一呢.
可不可以有不同的意見都可以合一呢.
當然我們的經驗是這樣的.
當然有時候也不是反過來講.
你的意見完全相反,就是合一.
不是這個意思,但是透過健康的討論,辯論,協調,磋商.
我們可能最後會有接近的意見.
這種都可能是合一.
而且有時候大家都是一個意見.
這些未必是合一,為什麼呢.
因為可能有些人是透過操控的手法.
使得大家,現在我講,說了算.
到時候每個人都同一個意見,這是不是合一呢.
再嚴重一點,如果我們了解那些教派.
有些異端的教派.
很多時候就是將這件事用得更加厲害.
所謂合一其實是洗腦.
大家都是同一個想法.

$^{41}$但其實大家那些所謂合一,只是一個表面的合一.
另一方面,我們講合一.
是不是說沒有意見,甚至是沒有爭執呢.
我們很多時候講小組的動力.
team building或者一個組織的發展.
我們都會學到.
如果你,坊間那些.
forming, storming, norming, performing, joining.
再簡單一點.
一個小組開始的時候,形成的時候.
就蜜月期,大家多開心.
一個新的群體,大家不是很熟.
因著一下,又有些新的刺激.
大家就很開心在一起.
婚姻都是這樣,蜜月期.
誰知道,到差不多就開始磨合.
大家意見不同,可能吵得很厲害.
甚至組裡面都有很多想法.
一起工作,很多不同的意見.
但慢慢磨合了,就去到成熟期.
那個時候就合一,更加合一.
不是說大家意見不同.
而是說我們摸通了大家的想法.
或者大家懂得相就.
這種都是合一.
所以今次我們講聖經裡面的合一.
其實是有很多層次的.
在已忽所書講的角度.
亦都有不同的角度.
這個題目合一和睦的代價不簡單.
集中在代價.
即是說原來合一是有個代價的.
基礎得來不容易.
但其實合一這件事亦都不簡單.
我先講大綱.
第一,已忽所書第二章經文是很著名的經文.
待會我們會集中在十一至二十二節.
但我們都會提少少前文.
第一段,即是十一節開始講.
或者前文少少,就是合一和睦的基礎.

$^{81}$第二段,就是如何表達合一.
有三個教會的類比,或者叫表象.
比較難明的表象.
即是imagery.
你一想到一個形象.
你又想起去吐上去.
例如教會是神的家.
一個家就是教會.
這叫做表象.
這三個表象就在第二段講.
讓我們搞清楚.
誰和誰合一,如何合一.
並且對我們有什麼啟發和應用.
所以第一方面就是合一和睦的基礎.
我們看看經文.
希望大家看得到.
剛才講了,十一節是我們今天讀的開始.
但有所以,我們經常說一定有前面那堆.
所以你們應當紀念.
你們從前看肉體是愛邦人.
如果你看十二節highlight紅色那裡.
有時你們與基督無關.
在以色列國民以外.
然後你再看前面.
其實是講第二章開始.
已忽所書.
我們一班人是怎樣的.
本來我們是一班在罪惡過犯裡的人.
這個我們很容易明白.
不詳細去解釋.
第八節那裡講.
我們其實是有罪的.
我們需要主的拯救.
而我們得救是因為神的恩典.
得救是本福恩.
我們變成一班甚麼人.
第十節那裡講,我們原是他的工作.
我們經常提及,原是他的工作.
就是我們成為神的masterpiece.
即是一個最偉大的創造.

$^{121}$我們每一個都是masterpiece.
是神的masterpiece.
所以使我們可以行善.
神預備我們行善.
神救我們脫離罪惡.
要我們去行善.
所以我們本身是一班這樣的人.
如果你讀多一點.
你可以說我用一個說法.
我們之前是甚麼.
五毛的人士.
甚麼叫五毛呢.
不要想其他東西.
現在講的五毛.
第一就是武國禮.
即是他是非以色列國民.
以色列人,大家知道.
有神的應許.
當然其實.
他有神的應許.
在舊約不是只救以色列人.
他是選民.
但神是想透過以色列這班選民.
成為萬國的榜樣.
並且那些人可以相信他.
這個本身在背後.
是有一個這樣的想法.
不過以色列人.
將他被揀選的身份.
變成了獨有的尊榮.
你們是外邦人.
我們是以色列民.
不過事實上.
聖經要講我們沒有國禮.
我們不是以色列民.
當然如果有些初一班教會.
沒多久就說國禮是甚麼.
那你自己問隔壁那個.
問不到隔壁那個再問隔壁那個.
不想很詳細的解釋.

$^{161}$一個記號.
表示你是以色列民.
一個這樣的意思.
所以第一毛就是無國禮.
第二你留意他的經文的說法.
第十二節.
你們與基督無關.
我們說無關基督.
不關事的完全是無關的.
跟耶穌基督.
跟彌賽亞.
跟著在以色列國民義.
我剛才講了.
在應許的諸約上.
是局外人.
即是沒有應許沒有諸約.
當然那個應許是單數的.
特別是指彌賽亞可以有解釋.
不過要不詳細講.
應許沒有的.
沒有指望沒有神.
即是五毛.
有的是甚麼? 犯罪.
只有這些外邦人.
所以我們的身份.
其實以前是這樣的.
一個對比現在是怎樣.
現在是蒙恩的基督徒身份.
以前遠離神.
你看到在第十三節.
我們再看下面.
你們以前是遠離神的.
現在因為耶穌基督靠著他的血.
已經得親近.
即是遠近.
以前猶太人是近.
內邦人是遠.
我們本身都是遠.
但現在變成近.
為甚麼? 是因為耶穌基督的血.

$^{201}$其實這段經文.
如果講快一點.
你見到雖然很熟悉.
不過你見到很密集式在幾節裡.
講到耶穌基督的救恩.
那些關鍵字都出來了.
靠著耶穌的血.
第十五節.
靠著耶穌基督的身體.
第十六節.
在十字架上待死.
這些主要的字眼.
來表示耶穌基督的救恩.
是因著耶穌基督救恩的工作.
我們才能夠.
成為了.
以前是遠的變成近的.
成為了基督徒.
好像前面的經文.
曾經說我們得救是本乎恩.
是因為恩典的緣故.
我們成為了得救的基督徒.
真是差很遠.
有什麼分別呢?.
以前和神很遠.
現在就很近了.
與及之前有仇.
後來.
第十五節.
成就了和睦.
第十六節.
特別講.
與神和好.
兩下歸為一體.
成為一個新人.
這裡他講了很多次.
基礎是什麼呢?.
原來我們因為.
信了耶穌.
因為耶穌的救恩.

$^{241}$我們與神和好了.
以致我們都能夠.
與其他信耶穌的人.
他們都與神和好.
我們就可以和好.
我們與神和好.
因為救恩的緣故.
我們聚得赦免成為了基督徒.
有永生一切一切.
大家都能夠.
有這種救恩的.
效果在我們身上.
我們就能夠和睦.
這個基礎.
其實很清楚.
我們講一些背景.
講了兩下歸為一體.
你想清楚.
是哪兩下.
哪兩個派對.
是以色列人和愛邦人.
兩個不同種族.
兩個根本就.
勢成水火.
但是因著救恩的緣故.
可以走在一起.
這個其實道理.
不是很複雜.
不過有些背景要講.
這裡第二章.
第十四節.
有一個中間隔斷的牆.
都破了.
其實這是什麼呢?.
考古學家.
找到一些.
古蹟.
你看到這兩塊板是什麼呢?.
其實據研究.
共有十三塊.

$^{281}$講什麼呢?.
這個其實就是.
中間隔斷的牆.
隔什麼呢?.
中間隔斷的牆.
我們經常想.
我們和猶太人有牆.
或者人與人之間有牆.
隔斷不只是這樣.
我們看一看.
這是一個.
聖殿警告的碑文.
待會再看.
那個位置大概在哪裡.
那個碑文就是.
外邦人不能夠越過.
這道牆進入聖殿的範圍.
為者是怎樣呢?.
面臨死亡.
很嚴重.
還要再加一句.
如果這個人死了,他自己越過了.
只能歸咎他本人.
而他的死亡必然臨到.
什麼情況呢?.
這塊板在哪裡呢?.
我們看看聖殿的圖.
這些網上圖片.
這個很好用.
中間隔斷的牆.
好像很小.
其實這是.
《新約時代聖經》叫第二聖殿.
就是大希律造的.
在那時候還沒建完.
但是你會看到.
這個聖殿敬拜神的地方.
在這個外邦人院.
裡面.
在聖殿.

$^{321}$裡面的範圍中間.
其實就有一個矮牆.
大概這麼高.
但是有13塊.
據研究有13塊.
那些板就在各個入口.
提醒人們.
你不要過那裡.
假設我們在當時.
在這裡,我們不小心進去了.
你死你死.
不要怪.
其他人是你自己越過了.
是這麼嚴重的.
一個這樣的牆.
但是這個牆.
因著耶穌基督.
就破了.
如果這樣說.
外邦人能不能進去?.
可以進去.
因為耶穌基督可以進去.
當然當時仍然不能進去.
你要記住.
保羅寫的書信的時候.
聖殿還在.
做這些祭祀的事情.
不能進去.
但是在神學上.
在理論上.
我們每人因著耶穌基督.
其實連這個牆都可以廢掉.
你明白嗎?.
你實質都沒事.
不過.
其實理論上.
因著耶穌基督的救贖.
猶太人和外邦人因為同信一個神.
與神和好.
這個牆其實破了.

$^{361}$用聖殿的比喻.
那些例子.
來講這件事.
接著.
這幅圖就更加清楚.
有些人就畫.
當時的時間.
你看到有些牆.
一個個的人在走.
我想起什麼呢?.
給大家不要太睏.
大家有沒有看過.
這套戲.
精武門.
華人不得入內.
有些人真的去研究有沒有這個牌.
原來沒有這個牌.
只有兩個牌.
一個牌.
就是華人不能進去.
和其他人不能進去的租界.
背景不講了.
另外一個就是狗不能進去.
將張力.
放在一起.
耶穌基督.
是因為救恩將猶太人.
和外邦人.
可以合一.
我要強調很多次.
這個不簡單.
我們.
再看下去.
這裡就是說.
強調幾次.
因神使我們和睦.
使兩下合而為一.
講了幾次.
十五節黃色那幾個.
為了將兩下.

$^{401}$藉著自己 藉著耶穌.
造成一個新人.
接著就.
藉著十字架使兩下歸為一體.
十八節.
因為我們兩下.
藉著一個聖靈所感.
得到進度符的面前.
你明白他的強調.
兩個可以走在一起.
是不簡單的.
因為這兩個是猶太人和外邦人.
你試想一下.
我們這個世界.
很亂的其中一個原因.
就是很多族群的政治.
很多的差異.
你說那些人.
怎樣可以走在一起.
甚至有些人會說基督徒.
你們那些宗教人士.
都會搞到有些差異.
我們都承認.
因為大家是有不同的想法.
會有差異.
不過最難搞的.
就是那些族群的政治的手刷.
你見很多國家很多問題.
Genocide.
你要將兩個.
完全不同的.
猶太人和外邦人.
走在一起.
絕對是很難.
不過.
因著耶穌基督的救恩.
大家都相信了耶穌.
成為教會的一個群體.
合一問題就能搞.
不過當時可能有很大的衝擊.

$^{441}$因為當時新的社會.
新的聖經的背景.
是一個多元文化.
一個多元文化的情況.
尤其是初信耶穌.
是猶太人.
接著就有很多外邦人.
信耶穌就湧入了教會.
那就很大差異.
但是本來要提他們的.
就是這個差異.
是可以搞的.
因為我們的基礎已經是合一.
我們的生命已經是屬於耶穌基督.
我們和神和好了.
因著這樣我們都可以和好.
記住要強調.
基礎一定在神那裡.
你說兩個族群.
怎樣可以和好?一般來說不行.
但是因著耶穌基督.
一起信了主成為教會的群體.
就有得搞.
這個衝擊.
保羅拿來提他們.
就是在教會裡面.
也會有很大的挑戰.
就等於.
如果你認識經過的教會.
我去年也去過.
突然之間多了.
數十萬的移民進去.
香港人在環教會本身.
死氣沉沉.
不是說歧視大家.
那些教會老化得很厲害.
突然之間.
很多新的家庭進去了.
每一個教會都搞不定.
這個衝擊.

$^{481}$保羅就再提.
有得搞的.
因為我們一起信了主.
還有一個關鍵字.
不過我們現代人很明白.
就是我們.
同一個聖靈所敢.
得以進到父面前.
進到父面前.
如果你看本身的字或英文.
access to god.
大家現在能不能上網?.
不是叫你上網.
假設.
如果我們沒有.
沒有數據的電話.
我們要怎樣.
上網呢?在這裡.
要上網的wifi.
問題是你有沒有access code.
你能不能access到?.
這裡說什麼呢?.
我們都因為信耶穌.
我們可以access to god.
ok.
我們一起.
我們本身是沒有數據的.
但我們有wifi.
所以我們能夠上網.
access to internet.
因著這樣我們又能夠.
給予access.
即是說我要傳訊息給你.
因為大家都有wifi.
如果不是呢.
我們沒有數據.
我們怎樣去溝通呢?溝通不到.
所以大家不要到處寫whatsapp.
不要搞了.
更加我們沒有access.

$^{521}$但我們因為可以信耶穌與神和好.
可以access to god.
進到神面前.
祈禱和神溝通.
所以我們都能夠溝通.
我們都能夠和好.
這個就是.
說這件事.
所以呢.
我們再.
看下去.
這裡會進一步.
去說.
不合一的.
都合一了.
而我們.
這裡我待會再說多一點.
但回到前面.
如果說兩下合一.
這麼難的都可以合一.
何況在教會裡面.
有不同的差異.
是不是可以合一呢?.
在我們教會講三語的處境.
可不可以合一呢?.
事實上我們有三個不同語言.
三個不同文化.
甚至國語堂有四個文化.
國語堂大家都講國語.
不過我們講的國語是普通話.
究竟我們用什麼字眼呢?都討論了一大段時間.
國語堂有馬來西亞來的.
或者東南亞那邊的.
有台灣來的.
有內地來的.
有香港來的.
國語堂有四種.
差異怎麼搞呢?.
是很難搞的.
不過因為我們信耶穌.

$^{561}$我們在基督裡面.
是可以合一的.
這是一個基礎.
接著就說三個類比.
這三個類比.
都不是不熟的.
剛才已經讀了.
我們就在這裡.
討論了三個類比.
在黃色那裡.
與聖徒同國.
第一個類比.
原本的文字是沒有國字的.
其實他是在說citizenship.
即是說公民的身份.
我們同是公民.
大家現在就在討論.
PR在澳洲就是福利級.
但是公民就沒有問題.
是有分別的.
個個都是公民.
我們在基督裡面.
是公民.
這是第一個類比.
第二個就是.
神家裡的人.
這個很熟悉.
因為我們有一個團契叫ICOS.
ICOS是神的家.
應該這麼說.
ICOS本身的字是房屋.
或者是家.
也可以指聖殿.
所以ICOS其實就是.
神裡面的家.
我們就是一家人.
這個我們也明白.
也不是很不熟悉.
接著黃色那裡就說.
漸漸成為.

$^{601}$主的聖殿.
這個我們也會有點熟悉.
聖殿.
剛才說到.
我會看整段的經文.
其實他在說.
雖然有不同的類比.
不同的表象.
但其實主要的表象是什麼.
聖殿.
為什麼呢.
中間隔奪的牆.
也是在聖殿那裡.
接著這裡說.
耶穌基督為防國石.
這個是在說.
哪一個建築物的.
防國石.
是主的聖殿.
所以他又用神的家.
當然.
也是對的.
在家裡面.
教會就像一個家.
跟我們有一個親屬的關係.
但是.
特別是說到主的聖殿.
聖殿的意思是什麼呢.
一起去敬拜神.
但這裡有多些不同的意思.
逐樣可以再highlight.
其實仍然是說基礎.
兩下合而為一.
因著耶穌基督的救恩.
但是這裡特別說.
第20節.
並且建造在使徒和先知的根基上.
有基督耶穌自己.
為防國石.
特別去提這個根基.

$^{641}$留意不是使徒和先知.
是根基.
我們作為使徒先知.
傳道一些領袖.
是透過神的話語去建造.
但是我們是建造者.
我們不是那個根基.
根基是什麼呢.
耶穌基督防國石.
又看看.
這個防國石在聖殿裡面.
一讀就知道.
這些資料.
我自己前幾年再去.
這個就是當時的聖殿.
聖殿是向東的.
大家知道.
它的面向東.
這裡向西.
所以耶路撒冷所謂的曲牆.
現在的人不叫曲牆了.
叫西牆.
就是這個位置.
你會看到.
這是畫回當時的情況.
還有.
這個位置有很多石.
原來聖殿的建築.
厲害.
你看這個位置.
我們有一個西牆的.
旅程.
晚上去看西牆.
因為現在.
這些看不到.
經過很多年.
我們要去隧道.
才看得到這幅牆.
的情況.
總之要付錢.

$^{681}$有一個大的石.
是這幾塊.
其實我有些照片.
但找不及.
超巨型.
你可以想像一下.
13.6米.
即是45呎.
這麼闊的一塊.
是一塊石.
估計是500多至600多噸.
你想像一下.
那什麼叫噸呢.
一塊石.
原來它是.
古代世界建築.
最大.
的一塊石.
最大那塊是這麼巨型.
這是聖殿的根基.
大過什麼呢.
要想一下.
復活島那些石.
我聽過有人說想去復活島.
有沒有人去過.
那些石很巨型.
大過復活島那些石.
金字塔去過沒有.
大過金字塔最大那塊.
你說真的很厲害.
當時.
那個建築.
是很厲害.
希律.
他為了討好猶太人.
他建了.
聖殿共建了八十多年.
而耶穌時代.
指著聖殿.
聖殿是正在建的.

$^{721}$是46年.
所以其實大家是知道的.
一直在建.
後來.
很傷心說起來.
主後63年.
肯定寫.
已忽所誅未建好.
63年建好了.
建了八十多年.
很感恩.
但聖殿被毀什麼時候.
主後70年.
建好了七年.
我希望現在.
西班牙聖家堂.
剛剛封頂.
那個很厲害.
建築師建了144年.
還未建完.
還有十多年.
是建好了封頂.
高地很厲害.
希望不要幾年後.
不知道有什麼問題.
炸了.
所以.
是正在建.
但是.
保羅要用.
這個聖殿的.
防國石這麼厲害.
大家都知道.
來說耶穌基督.
作為防國石.
是一個怎樣的.
耶穌基督.
這是一個禁忌.
而我們一班.
先知仕徒.

$^{761}$或者傳道.
透過神的話語.
來建造.
所以你讀.
經文的時候.
為什麼在.
這個禁忌建造.
其實暗裡面.
是說你用什麼建造.
是要透過先知仕徒.
神的話語去建造.
所以說.
這件事我們怎樣合一.
就走不開.
在這一方面.
來被建造.
所以好像老生常談.
但是從這個角度.
如果我們的信仰.
或者我們的聖經.
是有問題的.
信仰又搖搖噩噩.
你不在根基去紮.
你說合一.
很難.
另外一件事.
暗指一件事.
這裡有兩件事.
有一個很清楚的.
國防靠他聯絡得合適.
漸漸成為主的聖殿.
很有趣.
漸漸變了主的聖殿.
ING來的.
他還未完成.
聖殿還未建好.
這個概念一會再說.
但是他特別強調.
國防靠他聯絡得合適.
事實上.

$^{801}$看資料.
那些石頭那麼巨型.
根基原來不用交水.
大家剛剛好.
靠力.
防地震.
他靠力.
根基就這樣建好.
巧妙地.
聯絡得合適.
所以教會.
在根基上.
都是這樣建造.
強調彼此的聯繫.
但是這個聯繫.
是有一個進程的.
並且成為.
神藉著聖靈居住的所在.
這個神的殿.
是聖靈居住的所在.
所以我們能夠合一.
所以我說.
暗指是什麼呢?.
我們有神的聖靈.
我們信了主.
但是我們怎樣去接觸神呢?.
我們靠著聖靈祈禱.
因為本身聖殿.
就是禱告的殿.
耶穌說這是禱告的殿.
不要變成一個賊窩.
聖殿的本質就是禱告.
所以怎樣合一呢?.
因著我們都是同一個聖靈.
有聖靈在.
有神在我們心裡.
所以我們可以去祈禱.
可以透過這樣.
使大家去聯絡得合適.
有一個合一.

$^{841}$以及這個進程.
這個聖殿還沒有建好.
是不斷在建的.
所以這個聖殿.
或者神的家.
這個比喻有時候我們想到.
這個聖殿是一個建築物.
不是,是一個有機的.
是不斷的在建.
甚至彼得前書說.
他另外說了一個說法.
這個聖殿.
教會的聖殿.
用靈工.
我們是活石,所以我們有活石堂.
怎樣叫Living Stone?.
因為是一個有機的.
這個不是一個建築物.
這個聖殿是不斷的.
會有一個擴展性.
有一個發展性.
不斷有人去成為.
聖殿的一員.
所以我們就因為這樣.
能夠合一.
所以.
說了這麼多之後.
我們想說.
那個結語了.
五分鐘就出了.
那麼.
我們為什麼要說合一?.
為什麼?.
因為我們得到.
那個救恩.
這個合一是基礎.
耶穌基督用祂最重的代價.
來救了我們.
有些人可能.
不明白,但是耶穌救了我們.

$^{881}$我們只不過認識了神.
認識了耶穌基督的救恩.
我們願意相信主.
所以我們成為教會的一員.
而.
以分所施給我們看到的.
不單止這個合一是人際關係的問題.
其實這樣說.
也包含了.
所有東西都合一.
不同種族,猶太人.
外邦人都合一.
包括很多不同的外邦人.
都可以合一.
你說人際關係那些的差異.
可不可以合一呢?甚至.
說到.
他其實用了一個字.
他剛才漏了去說.
他特別去.
強調一種.
他用這個字.
去說到我們.
有冤仇.
那些字.
其實是說一種.
兩個的關係已經是水深火熱.
敵對的關係.
敵對的關係都可以合一.
那有什麼不可以合一呢?所以.
包含了人際關係.
但不單止是解決人際關係.
不單止是說我們要合一.
目標一起去做一些事.
所以合一,不是.
不是不是,而是說包含了.
更加是一個救恩的問題.
你說怎樣去解決.
人與人之間的隔閡.
種族之間的隔閡.

$^{921}$人類的問題.
就是.
透過這個救恩.
以致人來去.
同神和好,以致人與人和好.
以致我們有一個合一的基礎.
連這些所有的衝突.
都可以化解.
去到這個程度.
所以其實.
背後有宣教的,你留意到.
剛才說的那段經文.
第.
先看看.
第17節,並且.
傳和平的福音給你們.
遠處的人和近處的人.
講傳福音,這個聖殿.
是不斷增長的,是有傳福音的角度.
所以不是說小組.
大家很開心,很合一.
不傳福音.
都不是合一,而這個合一.
是不斷有人來信主.
他本身跟我們有差異的.
信仰我們都能合一.
是這樣的一個局面.
想起.
很多年前.
在短宣中心事奉的時候.
大家知道有一個福音船.
叫望渡號,世界最大的海上圖書館.
有四百個船員.
很多年輕人.
來自不同的國家,幾十個國家.
有一次.
我很多次跟他們.
去接待,以及帶他們去傳福音.
有幾個年輕人.
有一個團隊.

$^{961}$他們派一些團隊去當地教會學習.
有一個團隊.
真是經典.
有匈牙利人.
有.
華人.
有美國人.
但是有兩個民族.
真是令我們很感動.
當然有很多故事,有以色列人.
有巴勒斯坦人.
那個團隊.
跟他們相交.
真是他的福音船.
在活現.
在基督裡面合一.
因為信了主.
所以.
以色列人和巴勒斯坦人.
這兩個年輕人.
成為死黨.
可以一起.
傳福音.
我們講合一,那些差異.
是因為我們的信仰.
我們要問的就是我們的信仰是怎樣.
我們有沒有建立在一個正確的根基.
我們有沒有去依靠著神.
靠祂.
聯絡得合適.
而沒有了這件事.
什麼都無話可說.
但願我們一起去追求合一.
在一個正確的路上,我們一起祈禱.
保守我們體驗經脈的時候.
實在是很宏大很奇妙.
主,你自己的救恩.
無可測度.
耶穌基督是防國石,也是很偉大.
求你讓我們.

$^{1001}$在一個好的基礎上.
去討論合一.
去依靠著你,成為合一的群體.
祈禱奉耶穌基督的名字.
阿們.
\newpage



\section{希伯來書 10:19-20}
\label{sec:032300}
\textbf{每一次你禱告 (華)}~~~~曾興才牧師
\newline
\newline
連結: \href{http://www.church.com.hk/acms/content.asp?site=cdc\&op=show\&type=product\&code=032300\&layout=sermon}{\texttt{032300}} ~~~~ 語音日期: 2026-06-14
\newline
\newline
\hyperref[sec:032299]{\small{< < < PREV SERMON < < <}}
~
\hyperref[sec:index_chronic]{\small{[返順時目]}}
~
\hyperref[sec:032301]{\small{> > > NEXT SERMON > > >}}
\newline
\newline
希伯來書 10:19-20
\newline
\begin{longtable}{cl}
\hline
\hline
章節 & 經文 (和合本修訂版)\\
\hline
10:19 & \begin{tabularx}{0.7\textwidth}{X} 所以，弟兄們，既然我們靠著耶穌的血得以坦然進入至聖所， \end{tabularx} \\ \\ \relax
10:20 & \begin{tabularx}{0.7\textwidth}{X} 是藉著他給我們開了一條又新又活的路，從幔子經過，這幔子就是他的身體。 \end{tabularx} \\ \\
[1ex]
\hline
\hline
\end{longtable}
$^{1}$一個人的成功,不但要靠著努力,還是要堅持,而且堅持到底..
因為一件事情的成功不是靠你一次的努力,乃是要繼續不斷的努力,並且要堅持到底..
禱告也是如此..
所以耶穌記錄在《路加福音》第十八章第一節,鼓勵我們說,要常常禱告,不可灰心..
要常常禱告,不可灰心.其實追求女朋友也是如此..
老張的朋友很想知道他怎麼娶到一個這麼漂亮的太太..
老張就很自豪地說,.
"其實我娶這個太太,我堅持了三年,每天都去買彩券,.
中的最多的只不過是兩百塊錢.朋友們都笑我愛做百日夢,.
結果我把那賣彩票的女孩娶回家了.".
告訴朋友們說,要常常禱告,不可灰心..
為什麼耶穌要我們常常禱告,不可灰心呢?.
因為耶穌記錄來了以後,祂真的為我們開了一條直接通往復神的一個道路..
因為在舊約聖經裡面看到舊約的信徒們,他們必須要透過祭司,包括祭司的獻祭,.
才能夠把他們的需要帶到神的面前..
而且大祭司只能夠一年一次進到資聖所裡面與神會面,而且是戰戰兢兢的..
那耶穌記錄來了以後,祂不但為我們捨身流血在十字架上贖我們的罪,.
我們透過信靠祂,我們得以被稱為神的兒女..
因此祂也恢復了我們和祂天父美好的關係..
這是為什麼在希伯來書作者非常清楚告訴我們,一起來看第十章的十九二十節..
這裡說,弟兄們,我們既因耶穌的血得以坦然進入資聖所,.
是藉助祂給我們開了一條又新又活的路,.
從曼子經過,這曼子就是祂的身體..
這裡講到,因為耶穌基督的血,我們可以進入到資聖所..
資聖所乃是舊約聖經裡面所提到大祭司一年一次才能夠進去的資聖所,神同在的所在..
而且是藉助我們,藉助祂給我們開了一條又新又活的路,從這個曼子經過..
這個曼子指的是什麼呢?.
乃是聖殿裡面資聖所和聖所之間非常厚的一個曼子,.
從上到下沒有人可以用他的手力可以撕開的曼子..
當聖靈告訴我們說,當耶穌基督釘死在十字架上,當他說成了將他的靈魂交付給天父的時候,.
聖靈告訴我們說,殿裡的曼子從上到下裂為兩半,就是這個曼子..
換句話說,神與人之間這樣一個隔閡的牆已經被斷開了,.
我們可以坦然無懼地進到父神的寶座面前,阿們..
所以換句話說,禱告不但只是一個我們把需要帶到神的面前,.
禱告乃是神為我們預備一條又新又活的路,.
好讓我們藉著禱告,藉著耶穌的寶血,能夠進入到神的同在裡面..
因此我們說,禱告乃是主與你親密的關係,阿們..
一起講一遍來,禱告是你和主的親密關係..
是透過耶穌基督的寶血,我們可以進到神的寶座面前..
因此每次你來到神面前禱告的時候,.

$^{41}$你不但只是把你的需要,禱告祈求交到給神,.
你說,你是可以進到神祂寶座的面前..
這是一個生命的關係,耶穌開了一個直接通往父神的道路..
所以耶穌基督自己在路加福音第十五章第七節很清楚告訴我們,.
你們若藏在我裡面,我的話也藏在你們裡面..
凡你們所願意的祈求就給你們成就..
你們若藏在我裡面,耶穌基督說,我是葡萄樹,你們是枝子..
我們與祂有這樣一個生命的連結,.
不再透過祭物,不再透過祭祀,.
而是直接能夠住在主裡面,主也住在我們裡面..
祂的話也在我們裡面的時候,我們願意什麼?.
只要祈求就給你們成就..
親愛的網友姐妹,你可知道禱告是何等大的一個權柄,.
何等大的一個恩典嗎?.
這乃是舊約聖經裡面這些祭祀先知他們默密以求的事情..
如今,因為耶穌基督的保血,撕開了這個曼子,.
好讓我們直接,貪婪無懼地進到祂的執聖所裡面去,阿們..
不但是如此,在約翰福音第十五章第十六節,.
耶穌繼續告訴我們說,不是你們揀選了我,.
是我揀選了你們,並且分派你們去結果子,.
叫你們的果子蠶存,使你們奉我的名,.
無論相父求什麼,他就是給你們..
看到這個生命的連結,也是我們差的同工,.
他要差一派我們每一個人去為他多結果子..
換句話說,當你在跟隨主的道路上,.
當你在服侍神的道路上,有任何的需要,.
只要你向父神來祈求,.
因為耶穌基督的明智的緣故,就會為我們成就,阿們..
換句話說,我們依靠神的人,服侍神的人,可以一無所缺,.
因為耶穌基督就是我們通往天父的一個唯一的道路..
因此,禱告,你要記得,不再只是一個儀式,.
不再只是一個形式,而是你與主非常親密的一個關係,.
是一個生命的連結..
所以每次你禱告的時候,耶穌說我們在天上的父,哇!.
當你依靠耶穌基督的那一刻,聖靈就內住在你裡面,.
你可以呼叫阿爸父..
想一想,這位創造宇宙天地萬物的主宰,.
如今成為你在天上的阿爸父..
這是何等寶貴的一個運續,何等寶貴的一個生命的連結,.
因此你要知道,每次你禱告的時候,.

$^{81}$你還是與主有這個親密的一個關係,阿們..
那你說,既然是如此的話,.
我們必須要相信耶穌基督所應許我們的,必定會成就..
換句話說,主必會隨聽和應允你的禱告,阿們..
一起講一遍來,主必隨聽應允你的禱告..
是的,我們主既然為我們開了一條又新又活的路,.
因此祂必隨聽應允我們的禱告..
你可否知道,其實聖經裡面很清楚告訴我們至少五方面,.
我們什麼樣的禱告肯定是蒙主隨聽和應允的..
你想知道嗎?.
不少人啊..
五個非常重要的原則,主必隨聽應允我們的禱告..
首先第一,要奉主耶穌的名求..
一起講一遍來,要奉主耶穌的名禱告..
奉主耶穌基督的名禱告..
剛纔已經看過希伯來稅第十章告訴我們說,.
是耶穌基督的血讓我們能夠坦然無懼地進到至聖所..
所以耶穌基督也親自應許我們,.
在約翰福音第十四章十三節十四節說,.
你們奉我的名,無論求什麼,我可能,我或許,不是嗎?.
他說,我必成就,教父應兒子得榮耀..
耶穌再重複恐怕你聽不清楚..
你們若奉我的名,求什麼,我必成就..
哇,耶穌的名字為什麼那麼寶貴呢?.
因為耶穌是那位神的兒子,道成肉身,.
住在我們中間,匆匆忙忙有恩典有真理..
乃是透過耶穌基督我們可以來到神的面前,.
直接進到他的至聖所..
所以你們奉我的名,無論求什麼,我必成就..
因為耶穌的名是大有權柄,大有能力..
這是為什麼我們基督徒每次禱告結束之前都要說什麼?.
奉主耶穌基督的名,求..
這個不但只是一個形式,.
或者是一個我們必須要講的一句話,.
每次你講這句話的時候,.
你必須想到耶穌基督的名字是大有權柄的,.
是因為祂的血使我們能夠直接進到父神的寶座面前..
所以我們可以看到在五星節之後,.
使徒們經歷了聖靈的充滿,.
然後第一批的門徒彼得跟約翰,.

$^{121}$他們的確經歷了耶穌的名所帶來的能力,.
所帶來的權柄..
有一天早上他們到聖殿裡面去禱告的時候,.
經過每門看到一個伸來切腿的人坐在那裡討飯..
彼得就看著他說,你看著我..
這個人以為彼得要給他施捨什麼州紀..
他就望著彼得,彼得說什麼呢?.
我們來看看第三章第六節..
彼得說,金銀我都沒有,那幹嘛看我呢?.
他說,只是我把我所有的給你..
彼得有什麼呢?非常寶貴的..
我奉拿撒勒人耶穌基督的名,叫你起來行走..
他們過去跟隨了耶穌基督三年半,.
如今耶穌基督復活昇天了..
他們第一次經歷耶穌的名是何等的有能力,.
何等的有權柄..
他說,金和銀我都沒有,我將我所有的都給你..
我所有的是何等的寶貴..
我奉拿撒勒人耶穌基督的名,叫你起來行走..
拿著他的手,這個伸來切腿的人又站又跳,.
站在那裡唱詩,讚美我們的主..
他們經歷了耶穌的名是大有能力,大有權柄的..
所以每一次當你禱告的時候,為病人禱告的時候,.
你要奉耶穌基督的名..
每次當你要醫病趕軌的時候,要奉耶穌基督的名..
為人祝福禱告的時候,都要奉耶穌基督的名..
因為祂的名是大有權柄,大有能力的..
所以第一,要奉主耶穌的名禱告..
第二,要奉主的旨意祈求..
一起來講,要造主的旨意求..
另外一個,垂聽蒙應允的禱告,是要按照主的旨意求..
很多時候你說,我求什麼沒有得到呢?.
雅各蘇告訴我們說,因為你萬求,按照你自己的私心,.
按照你自己的私慾來祈求,並不是按照主的旨意而求..
所以在約翰義書第五章,第五章裡面第十四節,十五節,.
使徒約翰很清楚告訴我們,.
我們若照祂的旨意求什麼?.
記得啊,神的旨意求什麼?.
祂就聽我們..
這是我們向祂存什麼?.

$^{161}$坦然無懼的心..
希伯來書第十章是同樣的字眼,.
坦然無懼的心,因為耶穌的飽血..
既然知道祂聽我們一切所求的,.
就知道我們所求於祂的無不得著..
換句話說,當你按照神的旨意求的時候,.
祂就聽我們..
而且我們所祈求的無不得著.阿們..
那你說牧師啊,什麼是主的旨意呢?.
我怎麼知道如何按主的旨意求呢?.
親愛的兄弟姐妹,主的旨意都在哪裡?.
都在聖經的話語裡面..
這是為什麼我們今年教會特別鼓勵你,.
要來背誦聖經,好好地默想神的話語,.
好好地背記神的話語,思想神的話語,應用神的話語,.
因為神的話語裡面就是神的旨意..
那你按照神的旨意來祈求的時候,.
無不得著.阿們..
因為神的旨意都在祂的話語裡面,.
所以你越熟悉神的話語,你的禱告越有能力,.
你的禱告越蒙神垂聽.阿們..
那你說有時候,當你遇到特殊的情況,.
你不知道該如何禱告,也不知道神的旨意是如何的時候,.
你怎麼禱告呢?.
在羅馬十一八章,二十六節,二十七節,.
我們才能記得二十八節,沒有看前面兩節..
二十六節說什麼呢?.
況且我們的軟弱,有誰幫助啊?.
有聖靈幫助..
我們本不曉得誕怎樣禱告,.
只是聖靈親自用說不出的探喜替我們禱告..
二十七節說什麼呢?.
見察人心的,曉得勝靈的意識..
見察者說什麼呢?.
因為聖靈造作什麼?.
連聖靈都造作神的旨意,替聖徒祈求..
Amen!.
連聖靈都會按照神的旨意為你祈求..
所以聖靈為你禱告,肯定神一定垂聽的..
為什麼呢?.

$^{201}$他是按照神的旨意求..
我這裡親身經歷了一次,.
真的不知道該如何禱告,不知道神的旨意是如何..
二零零六年大海嘯之後,.
我帶短線隊到印尼去做一些救濟的工作..
當地的牧者就帶我們去印尼的另外一個島叫雅琪省,.
必須經過南中國海..
南中國海著名的是常常有大風浪的..
這個牧師告訴我們說,.
上個禮拜剛好一艘漁船大風浪一打的話,.
船翻了,五個人都死掉..
我說,幹嘛告訴我這些事情?.
我們當天早上五點鐘就要坐船,也是坐個漁船,到對面去..
我就說我不怕,因為從小在一個島上唸書,.
六年船來來回回坐好幾次,大風浪都經過,沒問題..
驕傲的人,有好戲可看..
後來船還不到一半,大風大浪就打在船上..
我從來沒有如此暈船過的..
我見過風浪,沒有見過如此大的風浪..
被風浪打得我暈得想吐,吐不出來,想呼吸,不能夠呼吸..
真的覺得幾乎要死了,就趕緊爬到船的,.
我坐在船尾,爬到船的正中央,抓住柱子裡面..
心裡想著,主啊,我不能死,我兩個女兒還小..
我的工作還沒有做完,.
心中只有這個念頭,不知道該如何禱告..
真的是以為自己真的是沒命了,.
真的不能呼吸,也想吐,吐不出來..
沒想到不久以後,另外一個短裙子的姐妹也爬了進來..
她也暈船..
她爬進來做什麼?要找牧師為她禱告..
我從來服侍那麼多年,從來沒有那麼尷尬過,.
第一次有地方姐妹要我為她禱告,.
我沒有辦法禱告,我說不出話來,.
我就根本快暈死在那邊了..
她看到我比她更可憐..
結果是她為我禱告..
聖靈用說不出話來的嘆息為我們禱告,.
按照神的旨意,結果我沒事..
很多時候當你不能禱告的時候,.
有聖靈幫助你啊..

$^{241}$因為當你按照主的旨意的時候,.
特別聖靈按照主的旨意為你求的時候,.
就必成就,阿們..
第三方面,要憑信心地禱告..
大家都知道我們是因性稱義,.
因為相信耶穌基督成為天父的兒女..
我們做基督徒,跟隨耶穌基督的道路,.
服事神的道路,也是要繼續憑信心,.
但是禱告更是要憑信心..
所以耶穌基督告訴我們一句非常寶貴的應許,.
也是非常喜歡的一句話,.
在馬可福音第十一到二十四節,.
耶穌說:所以我告訴你們,.
凡你們禱告祈禱,.
凡你們禱告祈求的,.
無論是什麼事情,.
家庭,婚姻,各樣的事情,.
只要信是得著的,就必得著..
一起講一遍來,.
只要信是得著的,就必得著..
換句話說,.
耶穌說當你們禱告祈求的時候,.
無論是什麼事,.
你可以想到的為你的孩子,.
為懷孕,為醫治,各樣的事情禱告,.
無論是什麼,你只要信是得著,就必得著..
換句話說,.
出於信心的禱告,.
主一定會垂聽應於你的禱告,阿們..
這是耶穌基督親口告訴我們的應許,.
你信是得著就必得著,.
為什麼信心那麼重要呢?.
雅各蘇就告訴我們了,.
雅各蘇第一章我們看第六節到第七節,.
他說只要憑信心求,一點都不疑惑,.
因為疑惑的人,.
不相信神能夠做,不相信神夠好,.
就像海中的波浪被風吹動翻騰,.
這樣的人不要想從主那裡得到什麼..
哇!非常清楚!.

$^{281}$憑信心祈求,一點都不可以疑惑,.
因為疑惑的人好像海中的波浪翻騰,.
這樣的人你不要想從神那裡得到什麼..
那你說,不小心有疑惑怎麼辦?.
就奉耶穌的名,斥責疑惑,離開我吧!.
就斥責疑惑,退去吧,不要再回來!.
那你有完全的信心,依靠神,.
相信你信是得著就必得著..
上上禮拜我剛剛從瓜裡馬拉醫療轉學回來,.
在那裡服侍了450位當地的居民,.
然後我們帶領了42位弟兄姊妹慶祝,.
下午的時間我們就到非常貧困的家庭,.
分派食物給貧困的家庭,.
一些油,米給一些日常用品..
在回程的途中我們看到一位姐妹,.
她叫阿拉塞裡的姐妹,.
跟她一個六七歲的孩子剛走回來,.
我們就看到她就問她有什麼需要,.
其實這位姐妹她就得了腦癌,.
一兩年的時間,.
但是她是非常有信心的一個姐妹,.
我們就問她說你是否相信耶穌基督會醫治你呢?.
她說我相信,.
我相信神已經在醫治我了,.
我們整個醫療團隊跟我們短訓隊,.
我們就一直按手為她禱告,.
看著相片,.
她的腦瘤已經影響到她的嘴巴,.
影響到她的眼睛,.
她的六七歲的孩子跟她一起,.
我們就按手為她禱告,.
我們大家同心合一地奉主的名禱告,.
我們問她說只要我們信是得著就必得著,.
咱們一起為她禱告宣告的時候,.
哇,這個媽媽跟這個孩子就被聖靈充滿,.
去那裡哭,在那裡顫抖,.
我們禱告完她我們就非常有信心說,.
耶穌基督告訴我們說我們禱告祈求,.
無論是什麼只要信是得著就必得著,.
你相信耶穌醫治你嗎?.

$^{321}$她說我相信耶穌醫治我,.
非常有信心,.
我們相信這是禱告耶穌給我們的應許,.
我相信神會醫治她,.
雖然不是馬上醫治她,.
但是神會醫治她,.
因為她是大有信心,.
她相信她禱告祈求的時候,.
信已經得著,.
她就已經得著了,.
所以第二方面要平信心來禱告,.
第四,清心的禱告,.
告訴我們要清心的禱告,.
清心的禱告也是蒙神垂聽應允的禱告,.
在馬太福音三個月的經文,.
第六章第五節說,.
清心的人有福了,.
因為他們必得見神,.
所以當你禱告的時候,.
先求神的靈來檢查你是否心中有一些隱藏,.
還未認罪,.
悔改的罪,.
因為當你沒有一個清潔的心來到神面前禱告的時候,.
你的禱告會被攔住,.
怎麼知道呢?.
聖經告訴我們,.
在以賽亞書第五十九章第一節第二節,.
耶和華的板背並非縮短,.
不能拯救,.
耳朵並非發沉,.
不能聽見,.
是因為什麼呢?.
但你們的罪孽使你們與神隔絕,.
你們的罪惡使祂掩面不聽你們..
換句話說,.
耶和華並非沒有能力,.
祂的板背並非縮短,.
祂的耳朵並非發沉,.
沒有空聽你的禱告,.
乃是因為你的罪孽,.

$^{361}$乃是因為你的罪惡使祂掩面不聽你們..
換句話說,.
每次來到神面前禱告的時候,.
一起先來認罪悔改..
約翰不是告訴我們,.
約翰一書一章九節說,.
我們落認自己的罪,.
神是信實的,.
是公義的,.
必要赦免我們的罪,.
洗淨我們一切的不義..
這乃是對信徒講的,.
親愛的弟兄姊妹,.
不是對非信徒講的,.
約翰一書..
我們落認自己的罪,.
神是信實的,.
是公義的,.
必要赦免你的罪,.
洗淨你一切的不義..
因此你的禱告就毫無攔阻的,.
神的膀背就沒有縮短,.
耳朵就沒有發沉,.
就可以聽你的禱告..
阿們..
雅各書我們大家非常喜歡的一節聖經,.
第五章第十六節怎麼說呢?.
他說,.
所以你們要彼此什麼呀?.
認罪,.
互相代求,.
使你們可以得一致..
義人祈禱所發的力量是大有功效的..
我們常常只講義人祈禱所發的力量是大有功效的,.
所以前面那幾句話,.
要怎麼樣?.
彼此認罪呀,.
互相代求呀..
換句話說,.
當你要尋求一致的時候,.

$^{401}$在神面前,.
問你自己,.
我是否有得罪神,.
得罪人的地方?.
我是否心中有怨恨,.
苦毒,.
不饒恕?.
求神來赦免我..
彼此認罪,.
然後才互相代求,.
你纔會得一致..
這纔是義人祈禱所發出來的力量,.
是大有功效的..
是要你彼此認罪,.
對付你的罪,.
包括引而未泄的罪,.
求神來檢查,.
好讓你的禱告,.
毫無攔阻的,.
達到神的私人寶座面前..
這是非常重要的,.
清心的禱告,.
避謀神,.
除一定陰雲..
第五,.
同心合一的禱告,.
各位同心合一的禱告..
《聖經》裡面非常看重我們弟兄姊妹同心合一,.
因為合一非常非常的重要,.
教會的合一,.
弟兄姊妹的合一,.
家庭的合一,.
代詞的力量是大,.
而且大得非常非常有能力的..
所以耶穌基督在他大祭司的禱告的時候,.
約翰福音第十七章,.
他多次為門徒禱告,.
至少有五次,.
也為後來信靠他的人,.
就是我們,.

$^{441}$就是我們每一個人,.
至少三次為我們禱告說,.
叫他們合而為一,.
像我們合而為一一樣,.
叫他們合而為一,.
像我和天父合而為一一樣,.
叫他們合而為一,.
像我和天父合而為一一樣..
為什麼合一那麼重要呢?.
同心是非常非常寶貴的,.
同心合一的禱告,.
神避除一定陰雲,.
怎麼知道呢?.
耶穌在馬太福音第十八章十九節,.
給我們很清楚的應許,.
我又告訴你們,.
若是你們中間有兩個人,.
在地上同心合一地求什麼事,.
我在天上的父必,.
可能或許嗎?.
沒有,.
必為他們成全..
所以我告訴你們,.
若是你們中間這兩個條件,.
有兩個人沒有說,.
一個人一個人禱告好不好?.
好!.
但是更好的禱告,.
更有能力的禱告,.
更蒙神垂聽應允的禱告是,.
至少兩個人在地上同心合一,.
這個要求是兩個人,.
兩個人同心合一,.
在地上無論求什麼事情..
這是我常常用來鼓勵夫妻,.
一起禱告的應許的經文..
丈夫一個人禱告可以嗎?.
可以啊!.
太太一個人禱告可以嗎?.
可以啊!.

$^{481}$但是夫妻二人的禱告,.
就像馬太福音第十八章十九節所講的,.
你們在地上若有兩個人,.
最基本的單位是誰?.
就是丈夫跟太太兩個人..
所以鼓勵你謙卑一下跟你太太一起禱告,.
謙卑一下跟你丈夫一起禱告,.
兩個人同心合一地禱告,.
在天上的父必為你成就..
包括同物也是如此,.
跟你的同工一起同心合一地禱告,.
跟你的小組員一起同心合一地禱告,.
跟你同事的弟兄姊妹一起同心合一地禱告,.
如此地禱告,.
在天上的父必為他們成全..
阿們!.
這些瓜裡麻拉短行為,.
也是親身經歷了同心合一禱告所帶來的力量..
我們在瓜裡麻拉分堂,.
準備有五位弟兄姊妹受洗..
因為他們是租用俱樂部來做逐日崇拜,.
因此沒有自己的教堂,.
就借用另外一間教堂來進行洗禮..
而且洗禮的禁詞是在教會外面,.
你看看這個圖畫,.
在教會的外面,.
左邊底下藍色那個地方是禁詞,.
旁邊那個比較高的建築就是他的,.
樓上是教室,.
底下是廁所,.
就在旁邊那個禁詞..
我們要事效禮拜六四點鐘要舉行這個敬禮,.
然後我們就在那裡看到,.
我們在開始敬拜讚美,.
預備要舉行敬禮的時候呢,.
忽然間,.
天烏雲滿布,.
然後清盆大雨倒下來,.
又打雷又閃電,.
整個地方完全覺得好像黑暗籠罩整個大地,.

$^{521}$我們大家都嚇了一跳,.
怎麼辦,四點鐘,.
那時候已經四點鐘,.
我們在那裡準備要洗禮了,.
四點鐘在那裡烏雲滿布,.
然後大雨大雷又閃電,.
真的很可怕,.
非常非常大的雨不斷的下來,.
我們的薩伯多牧師,.
我們看看下一張,.
我們在那裡敬拜讚美的時候,.
我們的分堂牧者薩伯多牧師大有信心,.
他就起來宣告說,.
我們一起同心禱告,.
禱告神要把太陽再一次伸出來,.
這個雨要停止,.
我聽了以後蠻有信心的,.
我們就一起禱告,.
我就想說豬啊,.
雅各蘇告訴我們說,.
以利亞跟我們一樣性情的人,.
禱告以後三年半不下雨,.
我說豬啊,.
我們現在同心禱告,.
只要三個小時不下雨就好了,.
我們就禱告,.
禱告,.
拼命地禱告,.
沒想到不到十分鐘,.
天開了,.
你看天開了,.
我們就可以在外面為五位弟兄姊妹洗禮,.
你可以看到當時弟兄姊妹的心情好嗨,.
我們的神真棒,.
他們都沒有想到,.
大家禱告還不到十分鐘,.
雨停了,.
閃電也停了,.
雷也不打了,.
大家能夠順利地在外面為五位弟兄姊妹受洗,.

$^{561}$我想這一天是我們非常非常難忘的一天,.
第二天逐日崇拜的時候,.
他們還在講起這個事情,.
我們的神是聽禱告的神,.
我說阿們,.
因為聖經說你們在地上,.
若有兩個人同心合一,.
無論求什麼事,.
我們在天上的父,.
必為你們成就,.
千萬不要輕看你們夫妻的禱告,.
為你孩子的婚姻禱告,.
為那個叛逆的青少年禱告,.
為你的工作來禱告,.
為你的服飾禱告,.
大有能力的禱告是兩個人同心合一的禱告,.
她非常堅持,.
我們每個禮拜至少三四天早上是一起同心,.
為孩子禱告,.
為兒孫女禱告,.
求神來幫助我們,.
讓我們領受這樣一個應許,.
能夠同心合一的禱告,阿們..
不論是如此,.
我們知道禱告既然是這麼多寶貴的應許,.
你可不知道禱告是你最有能力的武器嗎?.
各位朋友們說,.
禱告是你最有能力的武器..
禱告是你最有能力的武器,.
保羅在許多的書信告訴我們,.
不論是以夫所書,.
更連多書,.
我們在血契中行事,卻不靠血契來征戰..
因為我們征戰的惡魔是空中掌權屬靈契的惡魔,.
你靠你自己血肉之軀,.
肯定是粉身碎骨..
但是神給你一個最有能力的武器是什麼武器呢?.
就是禱告..
怎麼知道?.
看哥林多後書第十章,第四節跟第五節..

$^{601}$這裡說我們征戰的兵器本不是屬血契的,.
Hallelujah!.
乃是在神面前有能力,.
可以攻破堅固的引領..
什麼是在神面前呢?.
是禱告進到神的面前嘛!.
耶穌我們看到一條又新又活的基督,.
禱告進到神的面前,.
有能力可以攻破堅固的引領,.
將各樣的巨謀,.
各樣攔阻人認識神的那些至高之事,.
一概攻破了,.
又將人所有的心意奪回,.
使他都順服..
基督,Hallelujah!.
你可不知道你手中有的是一個最厲害,.
最有能力的武器嗎?.
這個武器讓你能夠在神面前有能力,.
可以攻破一切堅固的引領..
親愛的弟兄姊妹,.
魔鬼撒旦不害怕你逐日來參加崇拜,.
魔鬼撒旦也不害怕你待會到行外去參加AVL的服飾,.
魔鬼撒旦最害怕的是你在神面前禱告祈求,.
特別是同心合一的禱告祈求..
魔鬼最害怕的是當主的兒女們,.
憑信心奉主的名,.
照主的旨意,.
同心合一祈求的時候..
這是為什麼我們在神面前有能力,.
可以攻破一切堅固的引領,.
一概攻破了,.
可以將人的心意,.
可以想象,.
人的心意怎麼能夠奪回?.
是靠著禱告..
你的孩子很叛逆嗎?.
你的丈夫很叛逆嗎?.
禱告!.
將他的心意奪回,.
使他能夠順服基督..

$^{641}$這是你最有能力的武器,.
你知道嗎?.
不是靠著你嘴巴嘮叨,.
你越嘮叨他跑得越遠..
你還是在神面前,.
好好的禱告..
在神面前你是有能力的..
這個能力不但使我們能夠攻破堅固的引領,.
這個能力也能讓我們經歷神蹟啟示的祕訣..
所以經歷神蹟的祕訣在於什麼?.
在於禱告..
你想經歷神蹟啟示嗎?.
禱告啊!.
這是剛纔所講的五個非常重要的神的應許的這樣的禱告..
在神面前禱告就能夠經歷神蹟啟示..
使徒保羅聽到了馬其頓的呼聲,.
就到菲律賓去要把福音傳給菲律賓的人..
沒想到到那裡的時候遇到一個被誣毀所使用的一個使女,.
她就把誣毀趕出去以後,.
後來就被當地人抓拿,.
被打了以後,就把她和她的同伴希臘關在監獄裡面..
你可以想象保羅說,.
我領受了馬其頓的呼聲,.
來到馬其頓的菲律賓要把福音傳給你們,.
沒想到被關在監獄裡面,好倒黴啊!.
她和希臘在做什麼呢?.
非常熟悉的經文在《實錄生》第16章第25節說什麼呢?.
約在半夜,保羅和希臘做什麼?.
禱告!.
在哪裡禱告啊?.
監獄裡面禱告,唱詩讚美神,.
眾囚犯也徹而聽,忽然!.
哇,非常喜歡這個字,.
禱告以後忽然地大震動,.
甚至監牢的地基都搖動了,.
監門立刻全開,眾囚犯的鎖鏈也都鬆開了,.
這裡面至少有三個神蹟,.
首先第一,保羅他們禱告,.
唱讚美詩的時候,.
這裡的囚犯都不是好搞的,.

$^{681}$竟然沒有咒罵他們,叫他們住手不要再吵,.
反而怎麼樣?.
徹而聽啊!.
哇,現在保羅和希臘的讚美詩肯定唱得特別好,.
他們徹而聽,.
然後第二,他們一聽拜讚美禱告的時候,.
忽然間地大震動,.
哇,監門打開了,.
囚犯的鎖鏈竟然都已經鬆了,.
這個第二大神蹟啊!.
第三神蹟是什麼?.
這御主看到以後想要自殺糟糕,.
囚犯要跑了,肯定要被砍頭的,.
保羅說,不不不,不要害你自己,.
我們都在這裡,.
竟然這些囚犯鎖鏈都脫了,.
監門也打開了,.
竟然不走啊!.
你可以想象嗎?.
這是不是神蹟啊?.
是神蹟啊!.
這裡御主看到以後,.
不但御主信靠了耶穌基督,.
全家都受洗歸羅族的名下,.
成為菲律賓教會最重要的成員,.
哈利路亞,.
禱告是你經歷神蹟的祕訣,.
阿們!.
你在你的婚姻家庭裡面,.
需要神蹟嗎?.
目前你工作好像在監獄裡面,.
久久都沒有一條出路,.
你需要神蹟嗎?.
在你的福祉上,.
好像遇到史無同,.
沒有什麼果效,.
你需要神蹟,.
這什麼?.
禱告啊!.
禱告是使你在神的面前有能力,.

$^{721}$可以攻破所有魔鬼撒旦所興起的堅固引領,.
一旦攻破了,.
可以將人的心,.
包括你孩子的心,.
包括你丈夫的心,.
包括你老闆的心,.
你員工的心都可以奪回,.
使他們都聖服耶穌基督..
這是主根本的應許啊!.
所以禱告是你最有能力的武器,.
不但如此,.
禱告詐美也能夠轉化成事,.
影響世代..
一講一遍來,.
以禱告詐美,.
轉化成事,.
影響世代..
禱告詐美可以轉化成事,.
影響世代..
這是為什麼教會今年做了一件,.
我們過去三十年沒有做過的事情,.
耶利哥行動,.
我們來推動,.
來鼓勵弟兄姊妹,.
一起來,.
透過禱告,.
透過敬拜詐美,.
透過請以盧,.
來轉化成事,.
來影響我們的世代..
阿們..
大家今天進來的時候,.
都已經拿到這一份,.
耶利哥行動手冊..
非常謝謝劉師母,.
她真的是做手冊是一流的,.
非常非常精彩,.
裡面所有的資料都在裡面,.
花了很多的心血,.
花了很多的心思,.

$^{761}$特別為大家預備了一個手冊,.
一共有三個階段,.
第一階段是二十一天的首萬禱告,.
今天是幾號?.
六月十四號..
到七月五號一共幾天?.
二十一天..
二十一天裡面首萬禱告,.
每天都有禱告的經文,.
禱告的重點,.
甚至禱告詞都給你了..
二十一天預備迎接耶利哥行動,.
從今天開始,.
一直到七月五號,.
這個第一階段是什麼?.
禱告..
第二階段是做什麼呢?.
就是耶利哥行動..
為什麼要有耶利哥行動?.
我們來看看,.
約書亞記第一章第三節..
我們來是領受了神給約書亞的應許,.
做我們的應許,.
祂說,.
"凡你們腳掌所踏之地,.
我都照做我所應許摩西的話,.
賜給你們了.".
神應許賜給摩西的,.
就是對約書亞說,.
我都已經賜給你了..
這是為什麼我們今年要把這個行動.
做為耶利哥行動..
你知道當神呼召約書亞,.
要帶領百姓進入到迦南美地,.
要收取神所賜給他們的迦南美地的時候,.
給他一個非常非常有趣的策略,.
特別是要攻打耶利哥城的時候..
耶利哥城是當時舊約聖經裡面所記載,.
最堅固的城牆,.
非常堅固到一種地步,.

$^{801}$沒有任何的軍隊能夠攻破的城牆..
當神賜給他們一個特別的策略,.
什麼策略呢?.
要他們吹腳的祭司和軍隊,.
每天繞成一次,.
什麼都不做,繞成一次..
繞了六天,第七天繞成七次,.
然後神就做功了..
請問如果你是其中一個士兵,.
你會怎麼想?.
拜託神啊,這怎麼打法?.
孫子兵法裡面沒有這一套,.
哪裡有這種打法?.
不動一刀一槍,.
每天帶著軍隊,然後祭司吹腳,.
繞著城走一圈,.
而且每天都走六天,.
第七天做什麼呢?.
看看經文怎麼講,.
第六章十五節,.
到第七日清晨黎明的時候,.
他們起來照樣繞城多少次?.
七次,唯獨這日把城繞了七次,.
到了第七次,.
祭司吹腳的時候,.
嶽叔要吩咐百姓說呼喊吧,.
因為嶽後華已經把城交給你們了..
你相信嗎?.
只有一個人相信,.
可以叫他嶽叔爻了..
嶽叔爻是信靠順福神給他的,.
unconventional,.
不平凡的一種打法,.
不動一刀一槍,.
每天繞著城走,.
第七天的時候還繞七次,.
最後一次怎麼樣?.
大聲地呼喊神啊,.
一呼喊,.
什麼事情發生了?.

$^{841}$我們看看第六章二十節,.
於是百姓呼喊,.
祭司也吹腳,.
你看這聲音大得不得了,.
百姓聽見叫聲便大聲呼喊,.
Hallelujah!.
城牆就怎麼樣?.
坍陷,.
百姓便散去進城,.
個人往前直上,.
將城奪取,.
Hallelujah!.
哇,.
神的道路畢竟是高過我們的道路,.
所以永遠高過我們的意念,.
耶律格行動就是一個處於完全信靠順服的行動,.
不是一般的行動,.
乃是一個屬靈的行動,.
要求我們信靠順服,.
所以我們耶律格行動是從七月五號禮拜天,.
七點半開始,.
一直到禮拜六,.
七月十一號的晚上,.
有復興之夜,.
每天晚上我們會用不同的城市,.
每天晚上為不同的城市來禱告,來宣告,.
來繞圈,.
然後也為影響世代的七座山頭,.
來為這七座山頭來禱告,.
求神能夠將這些堅固的陰雷都拉下來,.
使人的心被奪回,.
順服耶穌基督,.
Amen!.
七天的時間,.
我想禮拜六晚上一定會大聲的,.
呼喊啊,.
Hallelujah!.
大聲的呼喊,.
Hallelujah!.
耶律格的城要坍塌了,.

$^{881}$要坍陷了,.
彎曲的,.
堅固的陰雷都要來坍陷了,.
Hallelujah!.
你興奮嗎?.
你期待嗎?.
請不要忘記,.
七月五號到十一號,.
七月十二號我們做什麼呢?.
逐日崇拜結束以後,.
每個墓區都已經安排了你們要去禱告行經的地點,.
就在這個首次的後半段,.
每個地方都安排了禱告行經的地方,.
你就帶領你們的弟兄姊妹一起到七座山頭裡面,.
政府機構,.
學校,.
藝術,.
宗教,.
家庭,.
娛樂,.
傳媒,.
各個機構裡面就在那裡宣告禱告,.
就像耶律格城一樣,.
當耶穌帶領他的軍隊,.
帶領祭祀去禱告的時候,.
他們就是如此做,.
我們相信,.
這麼樣的信心的宣告,.
信靠順福,.
是要來翻轉我們戲谷的城市,.
要轉化我們每一個所居住的城市,.
要影響這個時代,.
為什麼呢?.
因為耶穌在馬太福音第十六章,.
第十九節,.
給我們這樣的應許,.
我把天國的鑰匙給你,.
天國鑰匙是什麼?.
就是禱告啊,.
禱告宣告啊,.

$^{921}$禱告宣告就是天國的鑰匙,.
罰你們在地上所捆綁的,.
在天上也已經捆綁了,.
在原文意思是說,.
已經捆綁了,.
也要捆綁,.
罰你在地上所示範的,.
在天上也要示範,.
當我們出去到各個城市的,.
每一個不同的山頭禱告的時候,.
我們宣告諸啊,.
帶走你的權柄,.
帶走你的能力,.
我們宣告諸啊,.
在這個城市裡面,.
堅固的陰謀也都要破除,.
哈利路亞,.
所有男主人認識神的至高之色都要被除掉,.
讓你的心能夠被奪回,.
使他們順服嫉妒,.
我相信你的城市一定會經歷轉化,.
你的小主要復興,.
你的墓區要復興,.
整個灣區要被神大大的得罩,.
阿們,.
這是神給我們的應許,.
因為神已經把最有能力的武器給你了,.
那就是禱告,.
特別是同心合一的禱告,.
阿們,.
我們一起禱告,.
因為祝你愛我們,.
不但救贖了我們,.
也把這最寶貴,.
最有能力的武器賜給我們,.
讓我們能夠貪婪無懼的,.
來到你的至聖所,.
透過禱告,.
我們能夠經歷,.
祝你自己的大能,.

$^{961}$祝你的同在,.
我們恭敬姜定文姐妹,.
再次到主人的面前,.
那麼今天明白,.
每一次我們禱告的時候,.
主啊,.
禰就在工作,.
教堂當年戴德森宣教士說,.
到我們工作的時候,.
我們在工作,.
到我們禱告的時候,.
是禰在工作,.
When we work, we work..
When we pray, God works..
幫助我們,讓我們明白,.
禱告的真理,.
讓我們願意,.
真心主義的話語,.
同心合一,.
在你面前使用,.
你給我們最有能力的武器,.
來轉化城市,.
影響時代..
謝謝我們的主,.
聽我們的禱告,.
奉考主耶穌基督,.
寶貴了民窮..
\newpage



\section{約翰福音 4:1-26}
\label{sec:032301}
\textbf{放下水桶}~~~~梁紫健傳道
\newline
\newline
連結: \href{http://www.church.com.hk/acms/content.asp?site=cdc\&op=show\&type=product\&code=032301\&layout=sermon}{\texttt{032301}} ~~~~ 語音日期: 2026-06-14
\newline
\newline
\hyperref[sec:032300]{\small{< < < PREV SERMON < < <}}
~
\hyperref[sec:index_chronic]{\small{[返順時目]}}
~
\hyperref[sec:032302]{\small{> > > NEXT SERMON > > >}}
\newline
\newline
約翰福音 4:1-26
\newline
\begin{longtable}{cl}
\hline
\hline
章節 & 經文 (和合本修訂版)\\
\hline
4:1 & \begin{tabularx}{0.7\textwidth}{X} 耶穌知道法利賽人聽見他收門徒和施洗比約翰還多， （ \end{tabularx} \\ \\ \relax
4:2 & \begin{tabularx}{0.7\textwidth}{X} 其實不是耶穌親自施洗，而是他的門徒施洗，） \end{tabularx} \\ \\ \relax
4:3 & \begin{tabularx}{0.7\textwidth}{X} 他就離開猶太，又回加利利去。 \end{tabularx} \\ \\ \relax
4:4 & \begin{tabularx}{0.7\textwidth}{X} 他必須經過撒瑪利亞， \end{tabularx} \\ \\ \relax
4:5 & \begin{tabularx}{0.7\textwidth}{X} 於是到了撒瑪利亞的一座城，名叫敘加，靠近雅各給他兒子約瑟的那塊地。 \end{tabularx} \\ \\ \relax
4:6 & \begin{tabularx}{0.7\textwidth}{X} 雅各井就在那裡；耶穌因旅途疲乏，坐在井旁。那時約是正午。 \end{tabularx} \\ \\ \relax
4:7 & \begin{tabularx}{0.7\textwidth}{X} 有一個撒瑪利亞婦人來打水。耶穌對她說：「請給我水喝。」 \end{tabularx} \\ \\ \relax
4:8 & \begin{tabularx}{0.7\textwidth}{X} 因為那時門徒進城買食物去了。 \end{tabularx} \\ \\ \relax
4:9 & \begin{tabularx}{0.7\textwidth}{X} 撒瑪利亞婦人對他說：「你是猶太人，怎麼向我一個撒瑪利亞女人要水喝呢？」因為猶太人和撒瑪利亞人沒有來往。 \end{tabularx} \\ \\ \relax
4:10 & \begin{tabularx}{0.7\textwidth}{X} 耶穌回答她說：「你若知道神的恩賜，和對你說『請給我水喝』的是誰，你早就會求他，他也早就會給了你活水。」 \end{tabularx} \\ \\ \relax
4:11 & \begin{tabularx}{0.7\textwidth}{X} 婦人對耶穌說：「先生，你沒有打水的器具，井又深，哪裡去取活水呢？ \end{tabularx} \\ \\ \relax
4:12 & \begin{tabularx}{0.7\textwidth}{X} 我們的祖宗雅各把這井留給我們，他自己和兒女以及牲畜都喝這井裡的水，難道你比他還大嗎？」 \end{tabularx} \\ \\ \relax
4:13 & \begin{tabularx}{0.7\textwidth}{X} 耶穌回答，對她說：「凡喝這水的，還要再渴； \end{tabularx} \\ \\ \relax
4:14 & \begin{tabularx}{0.7\textwidth}{X} 誰喝我所賜的水，就永遠不渴。我所賜的水要在他裡面成為泉源，直湧到永生。」 \end{tabularx} \\ \\ \relax
4:15 & \begin{tabularx}{0.7\textwidth}{X} 婦人對他說：「先生，請把這水賜給我，使我不渴，也不用到這裡來打水。」 \end{tabularx} \\ \\ \relax
4:16 & \begin{tabularx}{0.7\textwidth}{X} 耶穌對她說：「你去，叫你的丈夫，再到這裡來。」 \end{tabularx} \\ \\ \relax
4:17 & \begin{tabularx}{0.7\textwidth}{X} 婦人回答，對耶穌說：「我沒有丈夫。」耶穌說：「你說沒有丈夫是對的。 \end{tabularx} \\ \\ \relax
4:18 & \begin{tabularx}{0.7\textwidth}{X} 你已經有過五個丈夫，你現在有的並不是你的丈夫。你這話是真的。」 \end{tabularx} \\ \\ \relax
4:19 & \begin{tabularx}{0.7\textwidth}{X} 婦人對他說：「先生，我看你是一位先知。 \end{tabularx} \\ \\ \relax
4:20 & \begin{tabularx}{0.7\textwidth}{X} 我們的祖宗在這山上敬拜神，你們倒說，應當敬拜的地方是在耶路撒冷。」 \end{tabularx} \\ \\ \relax
4:21 & \begin{tabularx}{0.7\textwidth}{X} 耶穌對她說：「婦人，你要信我。時候將到，你們敬拜父，既不在這山上，也不在耶路撒冷。 \end{tabularx} \\ \\ \relax
4:22 & \begin{tabularx}{0.7\textwidth}{X} 你們所敬拜的，你們不知道；我們所敬拜的，我們知道，因為救恩是從猶太人出來的。 \end{tabularx} \\ \\ \relax
4:23 & \begin{tabularx}{0.7\textwidth}{X} 時候將到，現在就是了，那真正敬拜父的，要用心靈和誠實敬拜他，因為父要這樣的人敬拜他。 \end{tabularx} \\ \\ \relax
4:24 & \begin{tabularx}{0.7\textwidth}{X} 神是靈，所以敬拜他的必須用心靈和誠實敬拜他。」 \end{tabularx} \\ \\ \relax
4:25 & \begin{tabularx}{0.7\textwidth}{X} 婦人對他說：「我知道彌賽亞—就是那稱為基督的—要來；他來了，會把一切的事都告訴我們。」 \end{tabularx} \\ \\ \relax
4:26 & \begin{tabularx}{0.7\textwidth}{X} 耶穌對她說：「我就是，正在跟你說話呢！」 \end{tabularx} \\ \\
[1ex]
\hline
\hline
\end{longtable}
$^{1}$各位早安.
不知道大家早上起來第一件事是做什麼呢.
嘩 這麼開心 這麼好.
第一件事不是張開眼睛.
第一件事 除了張開眼睛之外.
有沒有人第一件事是拿你的電話出來.
或者ipad.
或者鬧鐘.
如果你第一件事起身.
除了張開眼睛之外.
你是拿ipad或者拿電話的 請舉手.
ok 都不少.
不知道大家有沒有試過.
在深夜本身有些睏.
結果你打算有自己的me time.
正呀 打算看一陣手機或者ipad.
看了半個小時.
看了一個小時.
經常都會覺得 為什麼自己這麼晚睡.
明明十點坐在床上 看到11點.
有這樣的情況.
我們經常都是這樣.
看message 看很多不同的社交媒體.
看短片.
其實短片都給我們一些安慰.
或者給我們分散日常的注意力.
覺得我們工作很辛苦.
照顧完小朋友 終於有自己的時間.
那就看一下.
電話就好像一個無底黑洞.
你不斷去loop.
去看新的post 短片.
透過很多你自己的時間.
希望我們拿回一些屬於自己的空間.
屬於自己的滿足回來.
但問題是.
大家不知道有沒有發現.
你會覺得好像永遠都不夠.
loop完一次 看完一次.
應該都開始覺得滿足.

$^{41}$但又想loop第二次.
心裡面的那份渴望.
始終都沒有被解決.
loop到累 就ok 我終於可以睡覺了.
但第二天 早上又來.
每一天重複重複跌入這個循環.
我相信大家都很熟悉這個感覺.
用電話 用電子世界去填補這份空虛.
那份渴望 好像沒有真正被滿足.
接下來 我們今天會看一下經文裡面.
有一個婦人.
她一生的人都在尋找不同的滿足.
希望去到能夠fulfill滿足自己.
直到有一天 她在井旁.
遇到一個比她有滿足的人.
我知道FF家庭團契的弟兄姊妹.
下星期會用這段經文去查經.
所以請你特別留心.
今天的經文會帶給我們有不同的思考.
我們一起來.
讀一下今天的經文.
《約翰福音》四章.
耶穌離開猶太 回到加利利去.
祂必須經過撒瑪利亞.
於是到了撒瑪利亞的一座城叫敘迦.
接近雅各給他的兒子約瑟的地.
雅各井在這裡.
耶穌因為累 坐在井旁.
那時是正午.
有一個撒瑪利亞婦人來打水.
耶穌就跟她說.
請給水我飲.
大家可以看到我highlight了一些字.
有些bold和underlined.
值得留意的是.
經文寫耶穌必須經過撒瑪利亞.
是否因為這條路近一點 所以是必須呢?.
經文提到雅各給他的兒子約瑟的地.
雅各井.
這也是一個重要的地方.

$^{81}$接著繼續讀.
第八節.
因為那是門徒進城馬食蜜.
撒瑪利亞婦人對他說.
你是猶太人.
為何要向我一個撒瑪利亞女人要水飲呢?.
因為猶太人和撒瑪利亞人沒有來往.
為何耶穌會親自開口.
主動去問一個撒瑪利亞婦人說話.
亦和她借水飲.
這裡的沒有來往其實很值得留意.
寫得很隱晦.
究竟實際上代表甚麼呢?.
第十節.
耶穌回答.
如果你知道上帝的恩賜.
跟你說給水你飲的是誰.
你早就會求祂.
祂也會早就給你活水.
這裡提到活水.
究竟活水和這個婦人飲的井水有甚麼分別?.
看完經文之後.
我們會有三個問題.
就照教我們收集這個問題.
就是聖經作者帶著我們去理解他的思路.
第一就是.
第四節說.
耶穌必須經過撒瑪利亞.
必須因為甚麼呢?.
第二.
耶穌作為一個猶太男人.
為甚麼會跟一個撒瑪利亞婦人說話.
而且問祂借水飲呢?.
第三.
甚麼是活水?.
跟婦人飲的井水有甚麼分別?.
我們就帶著這幾條問題去看我們今天的經文.
回答第一個問題.
耶穌必須經過撒瑪利亞.
就要明白這個必須有多出奇.

$^{121}$這幅是地圖.
在說這幅地圖之前.
要說一下猶太人和撒瑪利亞人的淵源.
在北國以色列國亡國之後.
亞述王就將外族人遷入撒瑪利亞.
和當地人通婚.
他們的後代就是撒瑪利亞人.
他們只接受摩西吳敬.
還要在基里森山建了自己的聖殿.
在主前大概一百二十幾年.
這個聖殿就被猶太人大兵毀滅了.
另一邊.
北國以色列就是撒瑪利亞人的祖先.
曾經聯合外邦人剿滅殘殺猶太國.
大家可以想像.
這兩族人雖然同源.
但是他們的仇恨很深.
因為他們互相殘殺.
以至在虔誠的猶太人.
在猶太地他要去到加利利的話.
他寧願繞路.
多走幾天都不願意踏入撒瑪利亞.
大家看這幅地圖就知道.
最南面橙色的那部分是猶太地.
橙色的那部分.
有耶路撒冷 有伯利恆.
中間那塊就是撒瑪利亞 有雅各井.
然後之後北邊.
綠色那塊是加利利.
有拿撒勒.
大家看到地圖中間有紫色的箭咀和一條線.
大家看到嗎.
很可能就是耶穌當時走的路.
由耶路撒冷經過雅各井.
撒瑪利亞再上到北邊的加利利.
當然你可以想像到.
猶太人是不會走這條路的.
猶太人他們是有另一條路繞過撒瑪利亞.
就是紅色那條路.
紅色那條線.

$^{161}$由約旦河東邊直上繞過了中間那塊.
就是黃色那塊地.
然後才到北邊.
綠色那塊是加利利.
猶太人寧願走很遠的路.
都不想和撒瑪利亞人有接觸.
所以耶穌的必須.
是不是因為趕時間.
走捷徑呢?.
這裡的必須不是講地理.
是講上帝的救恩計劃.
耶穌是必須走進去.
因為那裡有一個人.
甚至有一個民族.
必須經歷這個救恩.
就好像我們平時用GPS找路.
有時有些區域我們是不想走的.
就好像例如我這樣.
我上年撞車.
有時我都會不想再走回那條Highway.
有些陰影.
你會想繞過那裡.
不要再走那裡.
但耶穌的導航就不是計時間.
是計人的.
他專登選了一條路.
為了一個避開全村人的婦女.
因為當時的婦女.
通常打水的時間是什麼時間?.
是早上或是下午黃昏的時間.
很少會在中午的時間去打水.
因為中午.
你想一下在中東.
嘩!熱到什麼地步.
曬到太陽最猛烈的時候.
又乾又熱.
你可以想像到.
那位婦女.
很明顯是一個孤獨的人.
她連打水都要避開人.

$^{201}$耶穌就專登在最曬的中午.
走入撒瑪利亞.
找這個婦人.
解決了第一個問題.
我們就要問.
第二個問題就是.
為什麼耶穌會開口.
借水喝呢?.
回答這個問題.
第九節有一個很關鍵的句子.
就是猶太人和撒瑪利亞人沒有來往.
原文的意思還要具體.
就是撒瑪利亞人和猶太人之間沒有來往.
沒有來往是什麼?.
不是說他們沒有來往.
是說在井旁的場景裡.
他們沒有共用器皿.
猶太人覺得.
撒瑪利亞人是什麼?.
是混種人士.
整個民族都是不潔的.
他用他的杯.
他用他的水桶.
他碰一碰.
都會玷污了整個猶太人的習俗.
所以有些似就是.
大家經歷了疫情之後.
就會明白.
那幾年我們和別人吃飯.
我們用公筷.
戴口罩.
搓手液.
每個人都會用.
別人的杯子你會不會拿來喝呢?.
不會吧?.
如果有一個陌生人.
突然走過來.
說可不可以借你的杯子來喝水?.
你會怎樣?.
他用完之後.

$^{241}$我當然不會用.
是不是?.
所以經歷了疫情之後.
每個人的防線.
都特別明顯.
當時的猶太人和撒瑪利亞人.
經歷了打仗.
互相殘殺之後.
更加是見面都不會見面.
還遑論喝他手碰過的水.
是不是?.
所以耶穌這一句.
請給我水飲.
已經打破了兩重界限.
第一就是猶太人和撒瑪利亞人之間不潔的界限.
第二就是男人和女人之間的界限.
當時在耶穌的年代.
特別是猶太拉比.
通常會在公眾場合.
避免和非親屬的婦女說話.
這是一個社會禁忌.
所以當耶穌.
專登進入撒瑪利亞的時候.
走近這個婦人.
還要開口問她喝水.
這個舉動是一個很大的突破.
沒有什麼猶太人會這樣做.
但是當這個井旁婦人.
用兩軸人的界線拒絕耶穌的時候.
耶穌的回答聽起來有點九不搭八.
就說如果你知道上帝的恩賜.
和你說給我水飲的是誰.
你早就會求他.
而且他會給你活水.
當婦人還以為.
這個活水就是流動的水.
耶穌就說了一句金句.
凡喝者水的還要再喝.
誰喝我所賜的水就永遠不喝.
我所賜的水要在他裡面成為泉源.

$^{281}$直湧到榮生.
那麼耶穌所指的活水是什麼.
這是我們第三條的問題.
婦人一開始聽錯了.
活水的詞語其實是一語相關.
相對於井底的死水.
自免可以解作活動的山泉水.
所以她才會這樣說.
井這麼深.
你又沒有打水的器皿.
你憑什麼拿活水.
難道你比雅各更厲害嗎.
她這樣說.
當婦人還在考慮.
怎樣實際上打水的時候.
耶穌已經將層次提升到生命的層次.
留意約翰的用字很仔細.
他說雅各井的時候是祖宗留下的井水.
但是耶穌所賜的水.
他用了另一個字.
泉源.
我所賜的水要在他裡面成為泉源.
直到永生.
一個是靠人力一桶桶地打出來.
一個是不斷湧出來.
不用你打的.
所以我們進一步問.
究竟活水.
泉源是什麼.
大家如果將聖經翻後幾章.
約翰福音第七章.
約翰自己解了這個謎底.
人若喝了.
道了我者裡來喝.
信我的人.
就如同經上所說.
從他腹中流出活水江河.
耶穌指著.
這話是指著他.
信他的人.

$^{321}$要受聖靈所說.
原來活水就是聖靈.
是上帝親自鑄在我們裡面的生命.
如果我們帶著這個敘事的背景.
我們才會明白.
其實耶穌為什麼在下文提到.
婦人五個丈夫的對話.
祂不是因為想八卦唱衰這個婦人.
而是因為這個婦人代表著.
整個撒瑪利亞的民族.
一直在尋找水源.
一直在尋找泉源.
一直都沒有找到.
所以聊聊活水.
耶穌就說.
叫你的丈夫來吧.
那個婦人就說.
我沒有丈夫.
耶穌就說.
你說得對.
你已經有五個丈夫.
而且現在那個也不是你的丈夫.
很多人讀到這裡.
很快就會將這個婦人定罪.
覺得她的私生活不檢點.
但是大家慢慢看一下這段經文.
耶穌有沒有說過一句.
她是罪人.
耶穌有沒有說過.
你悔改吧.
一句都沒有.
對比起約翰福音第八章行淫的婦人.
耶穌起碼有一句說.
從此不要再犯罪了.
但這裡一句都沒有.
在那個時代.
其實婦女基本上是沒有離婚的權利.
優妻的主動權是在男人手上.
一個女人嫁五次.
很可能不是她選的.

$^{361}$很可能是因為她五次被離婚.
或者五次守寡.
無論是被離婚或者守寡.
都是代表著一次又一次的被遺棄.
一次又一次的失去了要求.
大家還記得她在中午打睡.
很可能不是因為她生活不檢點.
是因為她一直找的依靠.
到頭來沒有一個靠得住.
是因為她一直找的依靠.
到頭來沒有一個靠得住.
更加令人驚喜的是.
這五個丈夫其實是在映射整個撒瑪利亞的民族.
給大家看一段經文.
如果大家有看列王記夏十七章.
你就會發現在北角亡國之後.
亞述王就將五族人遷入了撒瑪利亞.
這五族人帶著五族的神明.
而撒瑪利亞人又拜了一份.
經文是這樣說的.
亞述王在這五個地方.
巴比倫古他等等的地方.
帶了很多人來撒瑪利亞.
各族的人住在城裡.
為自己製造神像.
安置在撒瑪利亞人所做的幽壇殿裡.
而撒瑪利亞人都有拜.
所以在撒瑪利亞的歷史.
正正就是一段嫁過五個神的歷史.
同一時間拜耶和華.
同一時間都拜當地五族的神.
是一個混雜的敬拜.
所以五個的丈夫這句話.
其實他不單只是在說一個婦人的故事.
他是在回響整個撒瑪利亞民族的故事.
他就說你們找了很多依靠.
但沒有一個靠得住.
所以耶穌說這句.
他不是說要踢爆婦人.
說她生活不檢點.

$^{401}$而是告訴你.
我認識你.
我知道你喝過多少口井的水.
我知道你水桶裝都不滿.
留意婦人的回應很有趣.
她說先生我看得出你是善智.
反映著耶穌講中了.
她生命中最深處的需要.
言下之意.
耶穌就說我才是賜活水聖靈的那位.
是生命的傳源.
是依靠你們以前的族人不斷地找.
都未找到.
現在耶穌親自來到你的面前.
都是因為這個原因.
所以如果大家發現第四章.
有講到敬拜的問題.
在下文就緊接著有敬拜的問題.
今天我們沒有時間講這裡.
留待下次再講.
你會明白到.
其實就是因為突然間.
講到生命的需要.
講到敬拜的問題.
她才會說我是這座山派.
你是那座山派.
這個婦人拿著手上的水桶.
就是她一直畏為生.
生命的傳源.
以為有很靠得住的關係.
但無奈五段的婚姻.
五次的成災.
多次的失敗.
反映著很多人生不如意的事.
你以為那件事順利.
我們就會找到倚靠.
滿足到我們生命的空虛.
其實我們每一個人.
手上都有自己的水桶.
一些是我們以為.

$^{441}$可以滿足到我們的來源.
有人的水桶是事業.
有人的水桶是關係.
有人是子女的成就.
有人是別人怎樣看我.
這些事本身都不是壞事.
但請留意.
如果我們每天拿著水桶.
以為裝回來的水.
可以滿足我們內心深處的需要.
但我們最終會發現.
真正的滿足.
在耶穌裡面才找到.
所以第28節的動作.
成為了整段經文最美的鏡頭.
就是這個婦人放下水桶.
進了城裡.
她放下水桶是.
你想一下.
一個中午熱死辣辣的時候.
專登拿著水桶來打水.
但沒有打到水.
她放下水桶.
然後她不要了.
因為她已經找到存源.
她已經找到真正.
滿足她生命需要的位置.
所以她不需要水桶.
所以來個總結.
說回剛剛開頭的三條問題.
我猜我們都應該回答得到.
耶穌為何要經過撒瑪利亞.
就是因為祂知道.
人最深層的需要.
是耶穌親自跨越界限.
去到井旁婦人當中.
尋找她.
她說 先生.
我看得出你是先知.
祂是一個真正敬拜的人.

$^{481}$第二.
為何耶穌要開口借水喝.
這個問題其實是引導人.
人以為真正的需要.
是我面前的水.
是我們面前的電子世界.
是我面前的關係.
是我的事業.
是我賺到多少錢.
水桶就是我們日日以為.
我們可以滿足到自己的倚靠.
耶穌問這條問題.
就是令我們覺醒.
真正滿足人生命的需要.
就是面前這位彌彩人.
第三.
甚麼是活水.
活水就是聖靈.
是在我們裡面.
湧流的生命泉源.
水桶是我們日日都以為.
可以滿足到我們的倚靠.
耶穌就要邀請我們今天.
放下這些虛假的倚靠.
面對這個主動越界.
似活水的耶穌.
我們可以怎樣回應呢.
放下我們手上的水桶.
是一些我們以為.
可以滿足到我們心靈.
最深處的倚靠.
這個世界不斷說服你.
那樣東西很重要.
這樣東西很重要.
這個世界的敘事.
非常之強大.
令我們從心裡面相信.
我們休息是怎樣.
我們省下時間去看電話.
看短片.

$^{521}$Netflix and chill.
令我們覺得我們心裡面最需要的是.
電子世界帶給我們的安慰.
給我們的感官刺激.
是一些現實世界裡面.
給不到我們的東西.
但最後其實你聽一聽.
內心深處你的身體.
告訴你這些東西.
從來都滿足不到我們.
我們可以做的事.
是我們回到聖經裡面的敘事.
約翰剛剛就和我們說.
從來滿足我們心靈的東西.
就唯有是活水.
唯有是聖靈.
唯有是耶穌自己.
換言之.
這個世界都不單止和你說.
電子世界才滿足到你.
還有其他東西的.
關係,成就,期望.
人家的認同.
不要誤會我的意思.
我不是說那些東西不好.
我是說如果我們將人生的剛確.
將我們的人生完全.
依賴在這些東西身上的時候.
我們就會成為一個有洞的水桶.
裝多少都不夠.
飲多少都不夠.
所以這個星期.
希望大家找一個安靜的時間.
問一下自己.
究竟我現在的人生.
最渴望的是什麼.
然後呢.
我想很誠實地和大家說一句.
沒有了那句.
不要緊.

$^{561}$就是耶穌請你把水賜給我.
耶穌請把這水賜給我.
可不可以大家和我說這句.
我們一二三說這句.
一二三 耶穌請把這水賜給我.
我希望這句成為我們今個星期的祈禱.
好嗎.
第二就是面對著活水的上帝.
我們可以走過去我們經常避開的人那裡.
耶穌必須經過撒瑪利亞.
其實我們人生裡面.
其實有沒有一些撒瑪利亞.
有沒有一些敵區 有沒有一些人我們正在避開.
可能是一些很久沒有和他講話的家人.
或者是在教會裡面曾經和他吵過架.
不是很合口味的一些肢體.
可能是一些你心靈裡面已經認定.
這傢伙應該不會相信.
那些的朋友.
今個月.
可不可以選一個你一直避開的人.
你問他主動借一口水來喝.
我的意思不是真的借一口水.
一個訊息 一餐飯.
一個問候.
從一件最小的事開始.
而第三.
就是.
由拿水的人變成分享傳言的人.
這個婦人的報道很簡單.
她沒有讀過神學.
她只是說你們看吧.
這個人將我之前所行的事一切都說出來了.
她得到活水 她沒有自己享受.
她馬上走進城裡和一群人說.
你們看看.
誰知整個撒瑪利亞城都相信了.
活水的特性就是湧出來的.
是收不住的.
當你從你心裡面相信.

$^{601}$耶穌才是你最需要的那個位置.
你就會很自然和別人說.
耶穌在你生命有多重要 有多滿足.
你的故事 你的以往.
你被尋找的經歷.
是你最有力的見證.
我們希望在這個最撕裂的世代.
如果我們肯成為跨族群.
跨階層 跨立場的教會.
我們教會就會成為一個井塘的教會.
讓那些乾渴的人.
在這裡感受到活水的真實.
最後.
思考一下.
你的水桶是什麼呢.
你願意放下它嗎.
你的撒瑪利亞在哪裡呢.
當我們思考下一步的時候.
下一步的時候我們想.
是怎樣呢.
我想我們可以.
從第一步開始.
就是由我們睡醒 張開眼睛那一刻開始.
我們就問.
耶穌.
可不可以賜活水給我.
我們問一下耶穌.
今天我們的生命最需要的是什麼.
我很希望你的答案是.
原來我最需要的是你.
耶穌.
我們一起祈禱.
耶穌多謝你.
多謝你必須經過我們的撒瑪利亞.
在我們最乾渴.
最避開人的時刻.
你坐在井旁等我們.
你將你的活水賜給我們.
讓我們放開我們手上的水桶.
讓你的靈成為我們.

$^{641}$生命中永留的泉源.
帶領我們.
讓我們走進我們一直避開的人.
讓我們的教會成為一個井旁的教會.
讓一些乾渴的人在這裡遇見你.
我們在禱告耶穌的聖面前.
\newpage



\section{士師記 7:7}
\label{sec:032302}
\textbf{聖經人聖經事:用手捧著舔水的三百人}~~~~黃克勤牧師博士
\newline
\newline
連結: \href{http://www.church.com.hk/acms/content.asp?site=cdc\&op=show\&type=product\&code=032302\&layout=sermon}{\texttt{032302}} ~~~~ 語音日期: 2026-06-13
\newline
\newline
\hyperref[sec:032301]{\small{< < < PREV SERMON < < <}}
~
\hyperref[sec:index_chronic]{\small{[返順時目]}}
~
\hyperref[sec:032303]{\small{> > > NEXT SERMON > > >}}
\newline
\newline
士師記 7:7
\newline
\begin{longtable}{cl}
\hline
\hline
章節 & 經文 (和合本修訂版)\\
\hline
7:7 & \begin{tabularx}{0.7\textwidth}{X} 耶和華對基甸說：「我要用這舔水的三百人拯救你們，把米甸交在你手中；其餘的百姓都可以各回自己的地方去。」 \end{tabularx} \\ \\
[1ex]
\hline
\hline
\end{longtable}
$^{1}$弟子們今天又繼續講士師記.
就是《機殿》.
《機殿》來到這裡.
我覺得《機殿》終於煉成了一個.
跟隨神很信服的人.
變成有一個很信服的心,而且有信心.
看看《生命樹》怎樣跟大家講.
聖經人聖經士看到上帝的心意,怎樣做.
今天是用手捧著舔水的三百人.
有三百人.
是用手捧起水來喝的.
當然捧起水來舔,舔的舌頭是怎樣舔的呢?.
不太清楚,因為捧起水來舔,多數都是這樣喝的.
他說捧著舔水,意思就是捧起水來喝.
又或者跟機殿說,人還是過多.
不是吧?.
唉,不如你告訴我,你想要多少人?.
叫了,你感動了,我又叫了三萬二千人來.
然後你跟他說,誰怕的回去.
又來又去,那些人怨我的.
好了,然後三萬多人,人多勢好,你又說不好.
好,剩下一萬人.
一萬人,你竟然說人還是過多.
喂,上帝不如你說,你說要多少?.
你說要多少人,我就找多少個人.
而且你最好說人數多少人.
還想他們怎樣的,用刀還是用矛的.
跳得多高的,他的身體怎樣的.
跑一百公尺要多快的.
體力怎樣的,你說完就好了.
人還是過多,帶他去水泉旁邊.
我會逐個試他.
唉,如果我是機殿就只有深呼吸.
嘆一口氣,不是吧.
你堂堂上帝,你給我名字吧.
你就不用叫那些人去喝水,在那裡試.
我要在水邊試他.
我就會逐個指點你.
這個可以去,這個不可以跟你去.
我指點誰,這個不可以去.

$^{41}$我指點誰,這個可以去.
我逐個點,你就給我起個名字叫他就行了.
上帝,你這個辦法會不會煩一點.
不如一次過,你叫那些人自己分開站.
機殿就帶他們走到水邊.
神沒這樣做,機殿就神吩咐他.
你叫他們喝水,全部人喝水.
凡用舌頭舔水,好像狗一樣的.
就使他們站在一處.
跪下喝水的,要使他們單站在一處.
那些是單站在一處,這裡是單站在一處.
我講一下,分別人還是太多.
現在我要選哪些,用舌頭舔水那些.
無論他跪下喝水,都是說.
我選哪些呢,我要用舔水的三百人.
因此你會發覺到,大家都舔水的.
舔的,這麼有趣.
舔水,舔就好像狗一樣,喝水.
舔舌頭,舔舌頭舔水.
好像狗一樣,那些人站在一處.
跪下喝,跪下怎麼喝?跪下喝不到,離開水很遠.
其實下面就講了,用手捧著舔水的有三百人.
其他都跪下喝水,其實應該這樣說.
在飲水過程裡面,誰把水放在手上舔上來.
然後人不是趴在河邊.
而是把水這樣,然後這樣喝呢?.
我就選那些,趴下去,好像狗一樣.
伸頭下去喝呢?我不選那些.
我選那些跪在水邊,或者去到水邊.
不過無論如何,他們是把水舀在手上拿上來喝.
我就選那些,我就選這些人.
這些用手捧著的有三百人.
好,這三百人站在一邊.
不是的,就這樣趴到水邊這樣.
全部又站在另一邊.
那些是跪下喝水的.
即是跪在水邊,然後人這樣彎身下去.
那些去到水邊這樣拿上來喝呢?.
那些我就選了.
其實這裡有一個隱喻.

$^{81}$就是那些跪下的我不選.
跪下的我不選.
因為他去到跪,隱喻他會跪拜權勢.
他會跪拜巴力,他會跪拜假神.
容易跪的,其實未必是.
只不過說這個跪下,就是他會跪在水邊.
當然他不會趴在水邊.
他會跪在水邊,然後頭這樣下去.
這樣拿來喝,這樣拿來喝.
但有些人可能是比較不小心站在那裡.
他選那三百個人.
其實分別就是.
有些人就說,如果他是拿酒瓶來喝.
他喝的時候,他會留意到周圍有沒有敵人.
比較警覺的.
那些看到水就趴下去喝.
那些就不夠警覺了.
終於他找那些用手捧起水來喝.
這樣可以觀察住那些人.
就三百人.
第七節耶和華對基甸說.
我要用這添水的三百人.
來拯救你們.
將米電人交在你們手中.
其餘的人都可以各歸各處去.
神在這裡說,我要用三百人.
如果我是基甸.
我就會大力手腳都舉起抗議.
或者倒給幾口冷氣.
唉,搖頭.
為什麼這樣,神啊.
你叫那三萬二千人來跟隨我,是你感動的.
我叫他們來,每個人都冒著生命危險.
終於有些人害怕.
打仗有沒有人不害怕的.
然後呢,害怕的回家吧.
不要緊的,神有工作的.
我會看著.
然後有二萬人走了,剩下一萬人.
對手五六萬人,我推測.

$^{121}$哇,我一萬人已經很難打了.
竟然你就說,人還是太多.
全部帶去水邊.
然後我要試一下,他終於剩下三百人.
三百人是三萬人的多少啊.
百分之一,我只是用你來的人的百分之一就可以了.
三百人,如果對方是和我們平手.
三萬人,就是說我只是他百分之一的人.
來打這麼多人.
有什麼可能呢.
我要用這三百人.
對手的百分之一或者還小一點.
就可以了.
其他人叫他回家.
這三百人就跟著機電了.
神啊,這次好了.
我們很多武器,弓箭,刀劍,盾牌太多了.
因為幾萬人丟下了.
現在可以了.
神說不用帶了.
帶什麼去打仗啊.
帶食物和豪閣.
就可以了.
劍呢?不用.
刀呢?不用.
弓箭呢?沒有用.
盾牌都要背幾個吧.
前一個背後一個都有保護.
不用,可以走了.
帶著食物和豪閣.
其餘的機電打發他們回自己的帳篷.
只留下這三百人.
而米電人就在他們下面的平原裡面.
當這三百人看著下面的平原裡面滿了人之後.
你猜你會怎樣.
害怕得不得了.
就帶著這三百人.
沒辦法相信.
這個就是用手捧起水來喝的那三百人.
其實神選用手捧起水來喝的那三百人.

$^{161}$為什麼呢.
其實是神沒有解釋的.
為什麼呢.
於是多數解釋的只能夠用現在的人.
套進去當時來想.
或者套進去打仗的時候想.
這三百人如果是跪下喝水.
那一萬人如果是跪下喝水.
就是他趴下去.
而且跪下.
所以那個屬靈意義可能不好.
這三百人是怎樣的.
可能他們站著彎身拿水上來喝.
於是他們不跪.
拿上來喝.
這些站得直了.
很有信心有志氣.
而且我覺得最為重要的就是這樣.
他們可能是會對周圍的事物敵人的情況.
周圍的情況比較留意.
所以他們是屬於比較有素質些.
上乘些的戰士都不一定.
不過無論怎樣.
這個令我們想起那套舊的電影.
叫戰狼三百.
只有三百人.
戰狼三百是說以前的希臘的人.
去跟別人打仗只有三百個戰士好好打.
其實他們都是全軍覆沒到最後.
因為三百人實在是打不過別人的.
當時跟波斯大軍對抗.
不過這裡三百人是贏了戰爭.
而且他們沒有拿武器.
是拿食物和豪華.
這是一個很大的考驗與挑戰.
不過我們學機電在這裡看到的.
就是他對上帝的信心和信服.
機電信心本來不是很好.
是懼怕的.
是害怕的.

$^{201}$是躲在酒駕打麥的.
然後神叫他怨氣的.
然後開始出來試驗.
第一個試驗是什麼.
你等我回來.
跟神說你等我回來.
我拿些祭物.
真的來了.
然後神就收了他的祭物.
用岩石盤石起火.
燒了他的祭物.
第二次就是曠野地上的.
我擺執羊毛在那裡.
整個野地是乾的.
神祖沒有露水的野地.
是不會的.
是很大的神蹟.
不過整個野地是要乾的.
然後我的羊毛要濕.
這樣都為他成就.
然後不行.
我再試多一次.
明天全部野地有露水.
是我擺執羊毛要乾.
這樣三次試探上帝.
試驗上帝.
要求證據.
神真的聽他很忍耐.
然後又在這裡做這件事.
你見到機電開始有信心.
而且有信心得來.
學到的功課是信服.
信心與信服要在一起.
就在這三百人那裡表現到了.
三萬二千人不要.
要三百.
百分之一.
真的戰狼三百對著米電大軍.
一個很大的信心.
有信服神的安排.

$^{241}$最大的信服之一.
只是帶食物和豪菊去打仗.
這個真的超級難得.
今天跟大家分享到這裡.
我們禱告之後明天再見.
我們一起祈禱.
天父來到這裡的時候.
心裡面真的全然的佩服.
寫一個服字.
服的是機電對你來到這裡的信心與信服.
你叫三萬多人來.
終於剩下三百.
要去打仗的時候.
武器都沒有.
是拿著自己需要的食物與豪菊就出發了.
平原下面滿了好像沙一樣的米電人.
就這三萬人要跟他們打仗.
主人怎麼可以得呢?.
機電心裡面真的不停跳.
可以怎麼做呢?.
但是他有這個信心與信服.
是你自己的說話他就跟隨.
更加感恩的是有三百個人.
是夠膽又同樣的信心與信服.
就跟他一起.
主人幫助我們.
其實我們真的需要一些.
跟我們有相同分享到這樣的信心與信服的人一起同行.
求主幫助我們.
為主打美好的仗.
誰聽祈禱.
奉耶穌基督名求.
阿們.
\newpage



\section{馬可福音 4:35-41}
\label{sec:032303}
\textbf{真正的平安  (見證:張瑪莉姊妹)}~~~~倫國平牧師
\newline
\newline
連結: \href{http://www.church.com.hk/acms/content.asp?site=cdc\&op=show\&type=product\&code=032303\&layout=sermon}{\texttt{032303}} ~~~~ 語音日期: 2026-06-14
\newline
\newline
\hyperref[sec:032302]{\small{< < < PREV SERMON < < <}}
~
\hyperref[sec:index_chronic]{\small{[返順時目]}}
~
\hyperref[sec:032304]{\small{> > > NEXT SERMON > > >}}
\newline
\newline
馬可福音 4:35-41
\newline
\begin{longtable}{cl}
\hline
\hline
章節 & 經文 (和合本修訂版)\\
\hline
4:35 & \begin{tabularx}{0.7\textwidth}{X} 那天晚上，耶穌對門徒說：「我們渡到對岸去吧。」 \end{tabularx} \\ \\ \relax
4:36 & \begin{tabularx}{0.7\textwidth}{X} 門徒離開眾人，耶穌已在船上，他們就請他一同去；也有別的船和他同行。 \end{tabularx} \\ \\ \relax
4:37 & \begin{tabularx}{0.7\textwidth}{X} 忽然狂風大作，波浪打入船內，以致船灌滿了水。 \end{tabularx} \\ \\ \relax
4:38 & \begin{tabularx}{0.7\textwidth}{X} 耶穌在船尾上，枕著枕頭睡覺。門徒叫醒他，說：「老師！我們快沒命了，你不管嗎？」 \end{tabularx} \\ \\ \relax
4:39 & \begin{tabularx}{0.7\textwidth}{X} 耶穌醒了，斥責那風，向海說：「住了吧！靜了吧！」風就止住，大大平靜了。 \end{tabularx} \\ \\ \relax
4:40 & \begin{tabularx}{0.7\textwidth}{X} 耶穌對他們說：「為甚麼膽怯？你們還沒有信心嗎？」 \end{tabularx} \\ \\ \relax
4:41 & \begin{tabularx}{0.7\textwidth}{X} 他們就非常懼怕，彼此說：「這到底是誰？連風和海都聽從他。」 \end{tabularx} \\ \\
[1ex]
\hline
\hline
\end{longtable}
$^{1}$弟兄姊妹們 早晨.
嘩 今日好感恩.
一黎到一個咁有力ge C班.
真係為你地感恩.
你知唔知我點睇C班ge  弟兄姊妹.
我其實唔係聽佢地ge 音樂ge .
因為冇理由係舊廳聽音樂架 同大家講.
係舊廳音樂不如你去大會堂好聽D .
係呢度聽乜ye 呢.
係聽佢地生命ge 見證.
係聽佢地點樣能夠藉住佢地ge 生命.
將上帝ge 榮耀彈出黎.
啊 我又幾鍾意你地ge 詩琴啊.
啊.
咁當然唔會忘記呀陳sir啦.
我地ge 指揮呀真係.
真係好耐 好多年朋友啊.
為你感恩 上黎就祝福佢啦.
俾佢好返一手心工.
哎呀 指揮又冇你手你點指.
冇搵你腳你指架嘛先.
我同大家講啊.
我可以聯想到 收到聽天國ge 時候.
一萬四千四百個C班員啊.
一齊係度見證神.
點解咁樣諗呢.
佢地 佢地今日呢個 呢個頌讚呢.
令我諗起.
1984年 我係全球gwo 個Urbana.
係全世界ge mission conference.
每三年開一次ge .
哇 gwo 度.
全部弟兄姊妹.
係咪.
三萬幾人.
一齊黎唱詩歌讚美神.
你大家諗o下.
勁唔勁啊.
我先明白咩叫天上ge 勁拜.
哇 你冇話指啊.

$^{41}$gwo 個震撼力啊.
真係 我都鼓勵大家有機會呢.
特別你鍾意missiongwo D .
你一定要去呢個Urbana.
我而家唔記得幾多年.
因為我去jor 一次已經.
比佢帶jor 我 呼召jor 我.
做上帝ge 僕人啊.
一次已經夠啦.
所以我唔使去第二次 第二次上天國得架啦.
所以 弟兄姊妹 C班門.
所以我為你地感恩.
係度能夠將上帝ge 榮耀帶出黎.
不過我都想鼓勵o下.
係我地當中ge 弟兄姊妹們.
其實C班呢.
係度將上帝ge 榮耀帶出黎.
一方面讚頌上帝.
另一方面.
都係比到我地作為回族.
弟兄姊妹 你記住啊.
我地.
都係C班黎架.
你唔好當你.
今日星期日黎到呢度.
你係純粹坐係度.
係度做咩啊.
坐係度做禮拜.
知唔知 因為.
我少同大家講.
今日上帝比我地每一個星期.
有一日.
大家見面.
一齊係度做咩事呢.
我話個故事比大家聽.
上帝係舊約就興起.
以色列國.
係眾國當中做咩啊.
去榮耀.
佢自己.

$^{81}$外邦人見到佢地gwo 種.
對上帝敬拜gwo 種莊嚴.
嘩 佢地話.
嘩 你個神真係勁啊.
我都要相信你地ge 神啊.
係舊約就興起.
以色列國.
我問大家啦.
係新約興起咩啊.
你同我都有份.
就係興起咩啊.
教會.
去榮耀上帝.
你睇o下.
幾蒙福 原來我地今日.
第一次會係度啊.
點樣榮耀上帝 我問大家啦.
係敬拜當中.
最榮耀上帝ge 時候.
係邊一個moment.
諗o下.
唔單止係師班.
讚頌啊.
師班係帶領我地.
俾我地一個example.
係點樣去讚你神啊.
繼續最重要.
你知唔知你同我.
就係上帝榮耀ge .
器皿.
平時ge 時候.
我地對上帝ge 敬拜跟隨.
係咪 信靠係我地ge 工作環境.
係我地ge 家 係我地ge .
社區去用我地.
所做ge 一切去榮耀神.
星期日呢個moment.
原來我地.
每一個人都有上帝榮耀ge 器皿.
咩黎架.

$^{121}$就係我地ge 生命.
同埋我地ge 嘴.
你知唔知我地嘴.
就係榮耀上帝ge 器皿.
當我地一起ge 時候.
讚頌神ge 時候.
上帝ge 榮耀唔係淨係充滿jor .
係呢個祭壇上.
係充滿整個會堂.
Hallelujah.
詩琴.
詩琴妹.
或者詩琴姐.
睇你大唔大得過我啦.
我今年71歲架啦.
你試比個D Major.
彈個D Major.
考我地D .
D 和聲.
大家試一齊.
唱Hallelujah.
Hallelujah.
不過要真係開口唱.
試下要睇下點樣.
上帝點樣接住我地嘴.
充滿整個榮耀.
充滿整個會堂.
特別係山上面gwo D .
你gwo 度.
我見到我D friend係度.
無錯啦.
唔該你鼓勵下大家.
一齊去榮耀.
我地ge 上帝.
Hallelujah.
來.
Hallelujah.
你睇下係咪充滿整個會堂.
頂住個門.
你呢就唔可以下面咁樣.

$^{161}$今日詩班都幾好.
跟住你有D .
你咩表現呀睇住囉.
就拎個銀包出黎唱得唔錯.
比多少少奉獻.
你地消費者.
ge 態度.
係唔要得我同你講.
頂住門你小心呀.
我成日同大家講呢.
你坐係度你唔好以為好安心呀.
其實有人坐係你旁邊.
特別你留心呀.
有D 吉位gwo D 地方呢.
gwo D 吉位呀.
就係我地ge 敵人魔鬼.
撒旦就坐係度.
呢個做咩係度.
梗係唔係好ye 啦.
你知魔鬼好多個名.
係咪.
又殺方知父啦係咪.
但係係度星期日.
佢做乜ye 呀同大家講.
佢就拿控訴者.
The Accuser.
咁你諗下佢係度控訴咩呢.
哎呀.
你睇下上帝.
你話佢地黎讚美敬拜你.
你睇下佢地咁樣ge .
都唔係讚美你ge .
你睇下佢逐個係度點名呀.
係咪.
我地話黎呢度去榮耀神.
la .
gwo 頭就中招啦係咪.
有D 就踢踢呃就.
所以.
你留心.

$^{201}$你側邊gwo 個位呀係咪.
有D 牧師提醒我地.
千祈唔好比D 吉位比魔鬼坐jor .
比邊個坐呀.
比耶穌坐.
係咪但係你比耶穌坐.
有D 好提醒囉.
有耶穌同住我一齊.
arm arm 呀你好唔比.
比個魔鬼坐係度.
你就真係傻啦係咪.
知唔知你話我要訓覺.
你梗係唔敢比耶穌坐係度架.
你知唔知點解呢.
佢成日提住你咪冇得訓.
弟兄姊妹所以敬拜呢.
唔係隨隨便便架.
敬拜之前gwo 晚.
就早D 訓覺.
因為要gwo 一剎那就show比神睇.
我點樣prepare myself.
準備我自己.
聽日去敬拜我D 神.
你記得嘛.
神睇我地敬拜.
佢唔係睇我地係度做乜.
佢睇我地ge 心呀.
所以大家留心D 呀.
我地ge 心好重要架.
耶穌話我地ge 心諗乜.
我地就會點做ye 架啦.
所以我想鼓勵弟兄姊妹們.
今日你話冇寫過對我咁嚴謹.
唔係呀.
我係唔想你sai jor 上帝.
比你ge 生命.
氣息呀.
你問我好朋友呀.
阿晨晃上大哥呀.
唱歌要幾多戲呀.

$^{241}$係咪先.
佢正話係度練ge 時候都已經.
一二三四嘩.
我聽到之後嘩勁呀.
邊個比佢呢D 戲呀.
四篇一百五十篇.
最尾一次點講呀.
凡有戲.
你睇下你地講.
幾陰公.
凡有戲.
你已經就快瓜得架啦.
我地聽到之後.
嘩我地咁樣同時講.
凡有咩呀.
戲 冇錯呢.
你咁你先能夠都要黎咩呀.
讚美上帝呀.
所以我淨係開始同大家講呢.
我同大家講.
離教會聚會.
唔係係度唱歌.
only.
not for the sake of singing.
大家搞錯呀.
不過我地用singing呢個skill.
將我地ge 讚美帶出黎.
我比例子大家.
而家我地最鍾意唱一首.
國歌係咪 點唱呀.
好啦 示範 聽下啦.
遲來 遲來.
遲來.
我問大家呢個係.
讚美呀.
因為係度唱歌.
正式點唱呀.
係呀.
我鍾意此回.
冇錯 呢個先係讚美.

$^{281}$所以呢.
How great art 個首歌 你真偉大.
個副歌點講呀.
我寧歌昌.
為咩呀.
讚美.
教主我神.
係咁.
但係我地今日成日.
比個衰鬼.
魔鬼呀.
呃jor 我地.
佢地叫我做咩黎呢度.
唱唱歌啦.
唱歌你又開心又鍾意.
係咪.
我問大家 你京班實在係你鍾意我鍾意.
ge 咩.
你係咪流失咁係度.
我好鍾意gwo 個女仔唱歌.
佢點會咁樣講ge 呢.
所以我地需要反思.
聲樂主義就提醒返我地呢樣ye .
第一次會.
我地黎京班實首先.
sai jor 上去比我地gwo D .
生命氣息.
係咪 唔用佢.
仲係度聽音樂.
上帝話既然你唔想.
唱歌係咪.
好易je  佢點做呀.
我咪收返佢囉.
咁你就大鑊囉.
咁你呀.
你又唔使點樣ge .
你有永生ge 你信耶穌ge 人係咪.
不過你上到天國話.
上帝點解收我返黎呀.
我比你去讚美我你都.

$^{321}$唔讚 我比你戲你又唔.
唱歌讚美我 咁咪收你上黎囉.
反正上黎你都要唱架囉.
係咪先 所以兄姐們.
有時都諗下呢D 神學.
反思呢 會幫助.
你糾正我地黎敬拜.
gwo 個心態.
明唔明呀.
詩歌就係咁幫助.
我地 將我地ge 生命.
藉著詩歌帶出黎.
去榮耀神.
咁你又點解 你就咁講得啦.
乜詩歌咁麻煩je  弟兄姊妹.
上帝用得詩歌.
音樂 梗有D 原因ge .
你真係當佢.
戇居架.
係咪 我話因為好聽呀.
好聽.
因為你當jor 音樂 詩歌.
係一種藝術去欣賞.
但係 音樂有.
另一個名叫做.
Musical Language.
我問大家 語言.
做咩架.
你開始明呀.
你返去睇一本聖經呀 成本聖經.
從來冇將音樂 係當jor 一種.
藝術去睇呀.
你話係咪呀 林牧師.
我問你 以色列人.
戴美完之後 你地人戴完唱歌.
有冇話.
嘩 今日D 利未人唱歌.
好好聽呀.
你有冇見過咁個comment呀.
你記得話比我聽.

$^{361}$我比你一萬蚊.
因為我即刻可以開工教聖樂.
乜都唔教.
一邊教音樂.
無呀 弟兄姊妹.
我幾辛苦返去找成本聖經.
我從來未見過 聖經裡面.
上帝將音樂當一種娛樂咁聽.
有人來串我喎.
老婆.
上帝將音樂當做.
music therapy呀.
音樂治療呀.
你記唔記得呀 素羅唔開心果時 做咩事呀.
大衛呢.
我問大家啦.
大衛彈琴係邊個叫佢彈呀.
係神呀.
好多時候呢.
我地斷章取義.
以為呢.
佢同聖經都有講呀.
我同大家講.
gwo 個係世人ge 做法呀.
如果係arm ge 話.
耶穌仲唔回應佢.
凡勞苦擔重擔的.
都當謝我.
到我者來做咩呀.
聽音樂.
係咪 你有冇見到咁ge 教導.
但係世人要喎 點解呀.
因為世人無神呀.
所以佢地唯有靠音樂.
去做咩呀.
慰藉下佢地gwo 種不平安呀.
因為音樂有另一個名呢.
你一講你就明架啦.
音樂唔係淨係叫music咁簡單.
聽住啦.

$^{401}$叫做soul music.
因為music.
係fit我地個soul呀.
係咪呀.
你問下陳大哥呀.
一講音樂第一件事.
講咩呀.
有冇feel呀.
feeling.
記得我七五年.
寫呢首歌好出名架.
feeling.
係咁多咋嘛.
但係你吹佢唔脹.
親愛的弟兄姊妹.
係廿一世紀.
我地親愛ge 蔡元為醫生.
佢點樣形容呢.
佢話依家係後現代ge 世界.
著做咩呀.
問大家.
feeling.
I don't feel like doing it.
arm 唔arm 呀.
而且仲要一個feeling.
仲要同一個好match.
一種個character.
就非常咩呀.
自我呀.
所以我成日講笑呢個咩generation呀.
呢個係me.
generation.
你睇o下.
幾危險.
音樂點解咁容易去呃jor 我地呢.
就因為音樂講咩呀.
feeling.
知唔知.
音樂好危險呀.
弟兄姊妹們.

$^{441}$因為音樂.
你一個唔小心.
佢就可以令我地.
聽住啦.
飲酒你就會酒醉係咪.
你個音樂會唔會醉呢.
記得個咩字呀.
陶醉呀.
所以你敬拜啦.
你一個唔小心呢.
音樂就令你咩呀.
陶醉在其間.
唱唱歌仲黑埋隻眼添.
我問大家啦.
如果你有神奇係度.
睇住我地.
我地個個係度咁呀.
Oh I love you lord.
And I lift my voice.
係咁一唱.
你ok唔ok呀.
所以gwo D 人呢就以為.
我係度敬拜呀.
佢地唔知因為佢地唔熟音樂D 厲害呀.
我一夾呀.
我就知佢地係度被音樂.
陶醉jor 係當中.
佢仲以為敬拜.
我問大家啦 邊個最鍾意呃我地呀.
The father of lie.
你知唔知點解咁厲害呀.
因為那惡者.
曾幾何時.
係天上甲邊雞呀.
潮州話 即係做乜ye 呀.
佢就係.
係度搞性惡.
佢個位呀.
帶領敬拜上去.
比佢精美ge 鼓底去做.

$^{481}$點知佢就因為咁樣驕傲.
而被跌jor 跌落地上.
咁你諗o下地上佢做乜ye 呀.
大家諗o下.
你返去睇o下.
你ge 以世傑書.
gwo 邊就講呢個問題.
大家見到 嘩 原來今日.
我地黎敬拜上帝ge 時候.
你一個唔小心就跌落D 陷阱裡面.
雖然世界企圖.
將音樂據為己有.
但係上帝絕對無.
容養呢件事發生.
讓教裡面繼續有.
聖詩黎幫助我地.
然後聖詩裡面D 歌詞.
就將我地.
對上帝ge 一種讚嘆.
gwo 種經歷寫落去.
你睇o下 剛剛唱gwo 首詩歌.
It is well with my soul.
係一個好好ge 詩歌.
你知唔知gwo 首歌講咩.
你知唔知 gwo 個作者.
十九世紀末ge 一個.
好出名 好有錢ge 地產家.
係芝加哥.
點知佢人生就遇著呢大風暴.
一個大火燒哂佢D 產.
最不幸ge .
其中佢唯一一個兒子.
就係呢個火災裡面.
被燒死.
所以佢好傷心之餘.
就讓佢ge 太太.
帶著三個女.
坐船去英國.
諗住行開o下.
點知真係.

$^{521}$人算不如天算.
去到英倫海峽.
遇著咩.
gwo 頭大火災 gwo 頭大風暴.
結果船沉jor .
剩得佢老婆得救.
去到倫敦.
等個電報比佢.
佢話三個女兒.
走jor 啦.
嘩.
人生係咪成日發生呢D 事.
唔係講笑.
係好可悲.
係好慘ge .
結果佢自己就即刻坐船.
去英國會合佢太太.
經過個海峽ge 時候.
佢就係gwo 度.
禱告.
上帝.
ge 安慰.
就淋到佢生命裡面.
然之後佢就寫出.
呢首詩歌.
上帝ge 平安.
你諗唔到就係呢.
咁危險ge 時候.
最咁悲哀ge 時候.
就成為jor 佢.
心裡面gwo 種.
平穩ge .
平安.
弟兄姊妹.
今日馬可福音.
其實第五章一個好簡單ge 故事.
就係講耶穌.
記得同埋同一齊過海.
係當中遇著大風浪.
你大家留心.

$^{561}$我覺得好有意思.
我跟隨耶穌.
你而家諗住咩.
新年最鍾意講gwo 句.
你平平安安跟住神都仲唔平安.
點知跟住佢.
仲遇著大風浪.
係咪.
而且呢個大風浪.
唔係一般ge 大風浪.
係好.
得人驚.
係生命.
你面對死亡ge .
佢話船入滿水.
弟兄姊妹諗船入滿水係咪沉jor .
大家諗下.
你撐艇仔.
入jor 水你會點做.
你繼續撐.
你做咩.
淨係我第一堂.
有人同我講一句潑水.
我大佬.
gwo 度廿幾呎船.
我當你入jor 半水.
你可以潑水潑出去.
梗係啦漁夫.
佢好醒目.
你而家做咩.
走人.
但你留心.
呢班門徒仲有少少.
疑似.
佢搵耶穌.
因為耶穌同佢一條船.
係咪.
你睇下.
少少信靠救jor 佢.
大家留心.

$^{601}$雖然之後耶穌都話.
你點解咁緊張.
我就好鍾意聖經呢個地方.
佢將人性.
真白話出黎.
當你生意失敗.
情意失敗.
當你人生失敗.
你第一句做咩事.
OH MY GOD.
就黎啦.
天啊點解咁對我.
人性黎架.
所以好多人話.
聖經好屬靈.
我同你講.
屬靈唔係等於無人性.
佢相對就好欣賞我地呢種人性.
佢比我地睇到人性.
gwo 種厲害.
但係更重要ge .
耶穌係佢地當中.
有無留意到.
我好鍾意呢段經文.
點解六節經文裡面有五節.
都有耶穌係度.
你留心下.
所以當我跟隨耶穌.
行走我地人生.
ge 道路ge 時候.
有乜ye 發生.
都唔緊要.
最緊要耶穌係咪到.
只係意思係我地有無信靠我地耶穌.
我地有無去跟隨佢.
跟隨佢你就定啦.
我同大家講.
所以你睇下點樣.
耶穌一起身就即刻.
平靜風浪.

$^{641}$我又好鍾意呢段聖經.
耶穌第一件事話佢.
風來啦佢地點解咁膽怯.
佢地膽怯住.
因為呢個生死問題.
但係大家睇下.
仲勁D ge 膽怯.
佢咪就大大巨攀.
點解呀?見到耶穌.
點解佢一講.
事就咩呀?成了呢?.
我問大家.
如果你係以色列人.
希伯來人.
猶太人.
你好熟呢個摩西.
悟經ge .
你見到有人咁樣做.
gwo 把口一講事就成啦.
你話gwo 個係邊個黎ge 呢?.
根據你ge 明白.
一定咪.
鄰國平牧師啦.
個邊個黎架?.
你睇睇創世記第一章.
上帝點樣創造呢個.
宇宙.
Let there be light.
係咪?只要口一講.
跟住點呀?事就成了.
所以佢地見到耶穌咁樣講.
佢地有D 驚呀.
佢知邊個黎ge .
所以佢咪驚.
連風和海都聽佢.
到底係邊個呢?.
弟兄姊妹們.
好多時候當我地遇見人生ge 慘事.
我地都會問.
嘩!D 神幫唔幫我地架?.

$^{681}$因為我地對佢唔認識呀.
弟兄姊妹們.
所以我感恩耶穌帶佢地過海.
就讓佢認識佢係邊一個.
但係耶穌好幽默呀.
呢個chapter就停jor .
第一節冇ye 啦.
所以你一定要係第五章.
睇下發生D 咩事呢.
你先知耶穌ge 厲害.
你聽住喇喎.
你睇返第35節.
你睇下耶穌講咩.
你俾35節我地睇睇.
係喇.
大家睇o下喇.
我們到到哪邊去.
耶穌記得話講緊呢句ye .
我地過去gwo 邊.
就咁多je .
第五章呀第一節.
你知唔知佢講咩呀.
他們就埋岸呀.
我問大家呢句說話咩意思呀.
即係架船到Po jor 喇喎.
有冇可能呀.
弟兄姊妹們.
架船有冇可能到Po 喎.
跟住你就明白.
你記住喎.
入滿了水喎.
居然可以埋岸喎.
有好多神蹟係度喎.
聖經裡面有好多神蹟呀.
你慢慢睇你覺得嘩咁都得ge .
就係咁就得jor 喇.
但係呢你知唔知呀.
一上岸.
耶穌見到咩知唔知.
第五章講咩.

$^{721}$一上岸就撞鬼喇.
係咪呀.
隻鬼一見耶穌即刻黎喇.
佢即刻撲到咁黎到耶穌面前.
去.
跪拜佢.
聽住喎親愛的弟兄姊妹.
我當你而家個個都係門徒.
一齊問佢個句.
一齊講呀.
睇返最後gwo 節.
四十一節.
登登登登登登.
到.
呢個藍色gwo 個字一齊讀.
這到底是誰.
係咪.
chapter就完jor 喇.
chapter 5 第五篇.
一開始就撞鬼喇.
跟住隻鬼就黎到佢面前.
睇住喇喎.
馬可點記載呀.
隻鬼黎到耶穌面前.
拜佢就大聲呼喊.
記住喎.
大聲呼喊.
你係門徒喎.
arm arm 問完咩話.
這到底是誰.
跟住隻鬼講呀.
至光神的兒子耶穌.
連龍都聽到.
點解知唔知呀.
因為佢鬼咁大聲.
原來宗教講ye 都幾好意思.
佢鬼咁靚.
我問大家.
鬼靚唔靚呀.
你覺得佢恐怖其實佢好靚.

$^{761}$鬼靚唔靚呀.
鬼咁靚.
你知呀.
帕克尼拉娃烏.
你點形容佢呀.
一隻魔鬼仔幫佢一隻拉.
真係出神入化.
見到無呀.
所以你開始明白.
中國人哩節學都幾叻.
好近乎聖經.
真係鬼咁靚.
鬼咁靚就鬼咁大聲呀.
所以呢.
你睇下.
邊個答jor 門徒呢個問號.
你驚呀.
鬼.
你都虛大聲呀.
我地有上帝ge 權柄.
鬼.
咁大聲係咪.
所以睇下.
佢終於明白到原來佢地.
所認識ge .
呢個耶穌.
唔係單單係夫子.
唔係單單係醫生.
更重要ge .
佢係神呀.
因為只有神.
先有呢種能力.
口講就咩呀.
事成.
所以弟兄姊妹們.
係咪好精彩呀.
上帝真係好厲害架.
我好感恩.
原來當我跟隨神ge 時候.
第一耶穌ge 保證係咩呀.

$^{801}$佢就與我地同在.
聖經舊約新約都咁講啦.
耶穌個別名.
叫咩名呀.
(唱).
呢首歌叫咩名呀.
(唱).
咩意思呀.
神與我.
同在.
幾精彩.
詩歌唔係淨係比我地唱好聽.
弟兄姊妹.
係提醒緊你.
弟兄姊妹們.
佢知道咩事.
神與我地同在呀.
佢就知我地係人世間.
好多悲慘事.
所以佢特登黎到世界上.
俾我地再一次同上帝和好之後.
點呀.
佢就一生一世與我地同在.
你老爺呀.
我今年71歲.
我睇下轉頭.
如果過去上帝唔係咁樣同我一齊.
我死jor 唔知N咁多次架啦.
趁我係弟兄姊妹.
正話一首詩歌你知唔知呀.
It is well with my soul.
我心裡得安康.
係1992年ge 時候.
我仲係Texas.
gwo 間教會裡面.
發生jor 一件好.
全美國都知道ge 事.
淨係summer.
暑假.
三架十五人van.

$^{841}$裝住一班arm arm 畢業ge 年輕人.
去retreat.
去退休會.
點知行到半個時候.
其中一部van.
俾一架水泥車係車路直撞落黎.
撞正中間.
好不幸.
個油箱就擺.
淨係一架車ge 中間.
點解本來擺後面架嘛.
佢就係驚人撞後面.
所以點呀擺jor 係中間.
你諗都諗唔到.
佢就係係中間撞埋黎.
一死死jor 五個人.
嘩.
人生ge 意外.
原來唔到我地話事.
但係我好感恩.
感恩咩呢.
雖然事情發生jor .
但係上帝.
又差派佢ge 僕人.
蘇榮志牧師.
佢arm arm 咁巧.
係阿靈屯距離.
三十里ge 地方.
佢收到呢個消息.
即刻zar 車過黎我地教會.
同我一齊去面對呢件事.
你知呀.
我當時係做緊傳道.
係咪.
好似仲有Mark傳道咁.
arm arm 開始.
鬼識做ye 咩.
好在蘇榮志牧師.
同我一齊.
的晚開jor 培靈會.

$^{881}$就係唱呢首詩歌.
It is well with my soul.
哇.
成功會中.
好傷心.
但係.
就詩歌做一次話比我地聽.
無論點樣咁.
辛苦幾艱難.
上帝與我地.
存在我地ge 靈魂.
仍然可以得到安慰.
鼓勵就感受.
詩歌帶領jor 我地之後.
gwo 幾個禮拜.
我諗唔到詩歌係咁厲害.
聽師妹們.
點解呀.
係因為主耶穌.
祂是王.
係泛濫洪水泛濫當中.
祂仍然是王.
所以當我地面對.
人生艱苦時候.
上帝仍然幫助我地.
我就係咁得到呢個.
得到呢個幫助得到呢個鼓勵.
更認定好似門徒咁.
你睇下大家留心呀門徒.
當你認識到呢個神之後呢.
佢生命都改變過黎.
我as a result.
經歷呢D ye ge 時候.
我先知我主神係咁厲害.
所以返去香港.
有機會仲要係.
宋牧師教會服侍.
一次邂逅可以咁多.
ge 機會發生.
甚至乎係 COVID ge 時候.

$^{921}$都好多艱難.
但係上帝又好好.
將我呢個消息呢.
借住 Mary 呢.
我認識呢.
我地有個禱告會.
音樂禱告會.
就將呢個訊息帶出去.
去祝福好多係 COVID 裡面.
困難ge 人.
歡迎 Mary 去同我地分享.
多謝教會.
比我返黎自己ge 屋企.
分享.
又好多謝利明學院院長.
倫國平牧師.
係我 COVID gwo 幾年裡面.
一路即係用音樂.
令我度過jor gwo 幾年.
係非常感恩.
而且係真係見到天賦ge 恩典.
咁我簡簡單單講一講.
其實世界上.
好多個衝擊.
我地睇到戰爭.
連續幾個二巴ge 戰爭.
一路落黎.
但聖經係講jor .
民要攻打民 國要攻打國.
到處飢風地震 呢只係災難ge 開始.
究竟經濟.
醉和.
突然其來ge 一D 疾病.
我地點樣去面對呢.
咁係疫情.
開始ge 時候.
我係係 2020 年係紐約.
係紐約ge 時候.
就有D 培訓係度做.
突然間就係個.

$^{961}$lobby 度食早餐.
不斷聽到全部都係咳聲黎ge .
咁原來有一種病.
就係呢個冠狀肺炎.
已經係度傳播緊.
咁當時好多人都好擔心.
咁我係gwo 時間.
同其他人一樣.
因為我地係上jor 年紀gwo D .
60 以上gwo D  我諗如果你睇返新聞.
即係係意大利係講緊.
掘一個窿就擺一個屍體.
不斷聽到唔同ge 人.
係離世.
我自己係疫情ge 時候.
有四十幾個朋友.
返jor 天教.
好難受.
我係紐約知道佢地又冇呼吸機.
又冇.
醫療gwo D ye 比我地呢D .
所謂遊客.
所以我就返黎香港.
返jor 香港之後呢 孤獨老人一個.
咁係點呢 就係.
去D 會所度食飯.
大家都識ge .
Hi Bye 即係大家都唔敢坐埋一齊.
過jor 好多ge 中秋節.
感恩節 聖誕節.
新年 都係一個人.
我相信好多朋友gwo 陣時都係.
有時我一路有做D 長者ge 服務.
所以當去探D 長者ge 時候.
大家只係隔住塊膠.
隔住個玻璃 掂唔可以掂.
拎住面之後.
只係有淚水.
係好難ge .
但你而家見到呢D 咁靚ge 相.

$^{1001}$無端端我唔係去.
登台做歌聲.
係因為gwo 段時間實在太難受.
咁我就一路都.
冇話即係用gwo D 咩.
社交媒體gwo D ye .
咁我就忽然間唔知.
幾時有一次食完早餐.
係一個會所度 鳥語花香.
咁我就覺得 咦 原來有一樣ye .
叫Facebook 喎 不如lor 出黎.
唱下啦 咁就自己.
因為我背歌詞ge  好鍾意 尤其聖詩.
我真係全部背得到ge 我都會背.
因為我覺得係好多弟兄姊妹.
或者係D 屬靈ge 長者.
佢地寫低jor 好多ye 係.
淨係讀歌詞.
唔使唱 都已經好多安慰.
好多ge 感恩.
好多ge 意志咁樣.
咁於是乎呢 我就.
開始就係個Facebook度唱.
咁係唱個過程裡面呢 就.
感恩ge 一樣ye 就係.
見到好多字出黎.
咁我初初就好驚.
佢地係鬧鬼我 即係唱成咁ge 樣.
收聲gwo D  但返去睇.
睇返gwo 陣時候 咦原來D Comment.
係好好喎 D Comment呢就咩呢.
你幫我祈禱啦.
我呢就COVID中jor .
我呢就 呀仲好呀.
雖然gwo 間公司執jor  我有另外一間.
公司就請jor 我啦 咁即係.
有感恩ge 又有難受ge .
咁gwo 陣時我就覺得 咦.
原來真係得架喎 一路唱.
聖詩 然後一路幫人祈禱.

$^{1041}$分享分擔 真係好喎.
咁我自己又得到好多個平安.
好多個喜樂架喎 譬如你唱.
即係譬如唱好似會出外呀 呢D 歌呀.
誰曾應許呀 嘩真係好正呀.
陪我走過春夏秋冬呀.
即係呢D 全部都係我會背ge 歌黎架.
當我每次唱gwo 陣時候我就覺得.
今早我自己先arm 係車度又唱jor 一次.
擺jor 上Facebook 咁原來.
唔係淨係我自己得到安慰.
我自己得到平安 得到恩典.
原來好多睇ge 朋友.
佢地都有 咁你而家淨係見到.
呢D 相呢 就係輪廓屏幕師啦.
同埋師母啦 我地係戴住口罩 係2021年開始.
係我公司.
亦都係利美學院.
一路呢 就係唱聖詩.
然之後祈禱 原來我地.
帶到俾好多朋友gwo 段時間.
好多ge 安慰 今早第一堂.
尋日ge 五點半gwo 堂.
都有好多朋友去Facebook特登過黎.
話俾我聽.
佢地係gwo 陣時開始返教會.
開始聽福音.
開始用聖詩.
我地叫聖樂.
俾佢地好多ge 安慰.
亦都可以帶到俾佢地好多人生ge .
意義同埋方向.
因為聖經裡面真係教jor 我地好多ge 智慧.
咁所以我自己好聽話ge .
我成日覺得呢.
聖經裡面有好多.
我地講香港小姐就美麗與智慧.
我好注重智慧ge .
認識耶和華就係智慧ge 開端.
聖經話俾我地聽.

$^{1081}$我地要去傳福音 傳到去地極.
我地要先求神ge 國.
同佢ge 義 呢D ye 都會加給我地.
當我地走過死刃ge 幽谷.
我唔會遭害.
因為神應許與我同在.
佢ge 像 佢ge 光都會安慰我.
咁所以你見到我就一路唱一路唱.
呢個就係我地facebookgwo 陣時ge ye .
仲感恩到呢.
竟然係係疫情.
唔知第三波定幾時呢.
就有幾日 即係大約四十日.
係零感染ge .
咁我地同利美學院ge .
弟兄姊妹 牧師.
去到恩福gwo 度做jor 個千幾人ge .
利美學院ge 三周年紀念.
咁一路用詩歌.
一路用.
聖經裡面ge 話語去分享.
聖誕節好多人都出唔到去街食飯.
我地自己班朋友呢.
就走去唱 係聖誕節裡面.
唱好多好多ge 聖詩.
然後同影音使團ge 合作啦.
同好多好多 咁係gwo 段時間.
我開始就.
唱過 冇停過到而家為止.
其實上年ge 十一個月.
我都係海外.
係唔同ge 國家啦 英國啦.
法國啦.
澳紐啦 美加啦.
我淨係飛紐約.
飛jor 兩次 上年喎.
然後仲有一次係飛係Toronto.
係當中每一次.
去到唔同ge 教會裡面.
係分享神ge 愛同埋.

$^{1121}$福音ge 時候 見到每一次.
係尤其是呼召.
原來聖詩好有力.
因為裡面都係.
聖經裡面ge 話語.
都係好多人經歷jor .
神ge 醫治 神ge 憐憫.
神ge 寬恕.
神ge 施恩.
所以呢我地先至可以.
咁所以你見到呢D 全部都係.
我再要多謝倫牧師ge .
因為gwo 段時間一路唱歌.
咁多年我都冇唱.
我淨係做jor 張.
福音ge CD 做jor 一個.
咁就而家做緊第二張添.
咁我淨係覺得感恩ge .
我會將所有天父比我地ge .
恩典 我地ge 聲音.
我地能夠經歷到神ge .
呢一份福氣.
係應該要廣傳 唔使驚.
總之每一次出去.
天父就會帶領我地 咁呢D 都係上年.
ge 一D 報道會.
咁我去jor 11個月 我發覺到我而家.
好似傳道人咁樣 但我係.
好感恩ge  因為我自己係.
一個受惠者 我經歷jor 好多好多ye .
咁可能朋友如果你地.
上網睇都知道.
有D 見證係講緊50幾萬人.
睇 叫活出外 咁因為.
今日ge 時間 我就唔可以講太多.
我只係想話比你聽.
我ge 人生只有感恩 冇有.
奏作.
所以弟兄姊妹們.
原來我地睇到.

$^{1161}$聖言與聖學.
係似咩呢.
就好似馬丁路特.
所講 佢話.
除jor 神學.
第二樣最重要ge .
就係聖惡.
原來兩者.
個關係.
係咁密切.
我稱佢之叫天地劍.
人生一定有.
兩本 一本聖言.
就聖經一本就聖詩.
聖詩就係我地經歷.
個教導 聖言ge 教導.
係我地人生裡面.
活出黎 嘩 原來咁精彩.
有音樂呢個語言.
將我地帶到去.
同聲同氣.
同心合拍.
一齊將呢個榮耀帶出黎.
Hallelujah.
完全因為上帝.
比我地有呢個平安.
冇錯呀.
記住呀.
唔係平平安安.
而係唔平安裡面.
先至能夠得到基督裡面.
真正ge 平安.
我地一齊禱告.
天父我地好感恩你.
每個星期日.
你比我地有一日.
黎到你面前.
我地唔係係度做禮拜.
我地係度.
係將過去gwo 個禮拜.

$^{1201}$牧師教導我地ge .
點樣依靠 順靠你.
去活出黎.
睇到上帝你ge 真實.
就好似門徒跟住你一樣.
佢地係當中見到個個真實ge .
嘩 上帝.
真係厲害.
我深信.
如果我地當日每日都係咁睇.
將我眼目放係你身上.
我地同門徒ge 經歷.
都係一模一樣.
我地一定係大大懼怕.
因為.
你係我地生命裡面.
日日都屢創神跡.
我地就係咁.
能夠活到今日咁解咋嘛.
所以讓我地感恩.
要繼續ge .
係永恆裡面.
跟隨上帝.
祝福每一位今日係你面前.
榮耀.
頌讚歸比你ge 弟兄姊妹.
奉耶穌基督.
聖名而求.
阿們.
\newpage



\section{俄巴底亞書 1:1-9}
\label{sec:032304}
\textbf{驕傲的以東}~~~~鄭兆奇牧師
\newline
\newline
連結: \href{http://www.church.com.hk/acms/content.asp?site=cdc\&op=show\&type=product\&code=032304\&layout=sermon}{\texttt{032304}} ~~~~ 語音日期: 2026-06-14
\newline
\newline
\hyperref[sec:032303]{\small{< < < PREV SERMON < < <}}
~
\hyperref[sec:index_chronic]{\small{[返順時目]}}
~
\hyperref[sec:032305]{\small{> > > NEXT SERMON > > >}}
\newline
\newline
俄巴底亞書 1:1-9
\newline
\begin{longtable}{cl}
\hline
\hline
章節 & 經文 (和合本修訂版)\\
\hline
1:1 & \begin{tabularx}{0.7\textwidth}{X} 俄巴底亞所見的異象。論以東說：我從耶和華那裡聽見信息，並有使者被差往列國去，說：起來吧，一同起來與以東爭戰！ \end{tabularx} \\ \\ \relax
1:2 & \begin{tabularx}{0.7\textwidth}{X} 看哪，我要使你在列國中為最小，被人大大藐視。 \end{tabularx} \\ \\ \relax
1:3 & \begin{tabularx}{0.7\textwidth}{X} 你狂傲的心欺騙了你，你住在巖穴，居所在高處，心裡說：「誰能把我拉下來到地上呢？」 \end{tabularx} \\ \\ \relax
1:4 & \begin{tabularx}{0.7\textwidth}{X} 你雖如鷹高飛，在星宿之間搭窩，我必從那裡拉你下來。這是耶和華說的。 \end{tabularx} \\ \\ \relax
1:5 & \begin{tabularx}{0.7\textwidth}{X} 盜賊若來到你那裡，小偷夜間來到，豈不是只偷他們所需要的嗎？摘葡萄的若來到你那裡，豈不留下幾串嗎？你竟全然滅絕！ \end{tabularx} \\ \\ \relax
1:6 & \begin{tabularx}{0.7\textwidth}{X} 以掃遭到搜查，他隱藏的寶物竟被尋出！ \end{tabularx} \\ \\ \relax
1:7 & \begin{tabularx}{0.7\textwidth}{X} 與你結盟的都驅趕你，直到邊界，與你和好的欺騙你，勝過你，吃你飯的人設下圈套陷害你---他卻毫無聰明。 \end{tabularx} \\ \\ \relax
1:8 & \begin{tabularx}{0.7\textwidth}{X} 到那日，我豈不從以東除滅智慧人？從以掃山除滅聰明人？這是耶和華說的。 \end{tabularx} \\ \\ \relax
1:9 & \begin{tabularx}{0.7\textwidth}{X} 提幔哪，你的勇士必驚惶，以致以掃山的人都被殺戮剪除。 \end{tabularx} \\ \\
[1ex]
\hline
\hline
\end{longtable}
$^{1}$我們一起禱告好嗎?.
讓我們再一次能夠聽到.
上帝禰自己在日本的工作.
能夠有姐妹們去到當中.
我們真的要將剛纔每一個生命.
交到禰的手裡面.
我們離開了那個地方.
但上帝我深信禰仍然在那裡工作.
我們特別為當地的宣教士小慧去交賞.
求主禰的光照進每一個人的心裡.
讓他們每一個人的家庭.
都能夠感受到上帝禰的平安.
他們很需要醫治.
因為他們很多寡婦,很多孤兒.
他們都孤單,他們都哀傷.
願主禰再一次.
重新讓他們有生命的力量盼望.
更加知道日本的教會有很多很多的需要.
但願我們主禰能夠去興起他們.
好像興起我們一樣.
讓他們的信徒無論是多還是少.
但在禰的愛裡.
他們就知道上帝是在他們當中.
我們再一次求主禰.
建立當地的宣教士.
讓我們的短宣隊.
讓我們的教會.
能夠再有機會去服侍他們.
禱告奉主明求.
阿們.
OK 再一次謝謝你們.
其實我經常都在想.
其實日本人.
我不知道你們有沒有日本人朋友.
他們其實很多年前.
會不會都是中國人來的呢.
會不會其實大家都是同一個祖先呢.
為什麼我會這樣想呢.
因為過去我相信.
最近去日本週圍去看.

$^{41}$很多神社 很多寺廟.
其實很多時候都好像我們中國人的信仰傳統.
但因為過去歷史的緣故.
我以前經常覺得日本人是很驕傲的.
當然現在最驕傲的國家.
大家都知道是誰.
但每一個國家.
我又看歷史.
如果他是驕傲的.
其實他最終是得不到上帝的祝福.
教會一樣.
哪個教會驕傲.
哪個人是驕傲.
最終聖經告訴你.
他得不到上帝的祝福.
今天其實我也想和大家分享.
聖經裡面有一個驕傲的民族.
這個民族叫以東.
可能我們對猶太人和以色列都聽過.
但我們對以東就比較陌生.
但聖經告訴我們.
這個以東的民族是一個非常驕傲的民族.
以致在今天來說.
你都聽不到有以東人.
或者有他的文化在.
以東在舊約來說.
是一個非常重要的民族和國家.
在舊約的時代.
為什麼呢.
大家知不知道以東人的祖先是誰.
是的,他叫以掃.
以掃是誰呢.
以掃就是雅各的雙胞胎.
就是他的哥哥.
為什麼以掃的後裔會叫以東呢.
當然了,雅各的後裔我們知道是以色列.
所以他們兩個其實是兄弟.
都是同一個父親生的.
他父親叫Isaac.
他父親叫Isaac,其實是同一個祖宗.

$^{81}$但為什麼以東後來結束了呢.
本身以東的字的意思其實是紅色的意思.
為什麼會叫紅色呢.
在創世記告訴我們.
在二十五章的時候.
他說到了生產的時候.
即是李伯嘉要生產了.
即是他的媽媽.
他的肚子裡果然是一對雙生子.
先生出來的那個.
聖經說他是全身赤紅有毛.
像毛衣一樣.
他們就給他起名叫以掃.
所以以東的意思就是紅色.
因為他生出來的.
大家接生的時候.
不單止是一個毛鞋.
是一個紅毛鞋.
紅色的,全身紅色的.
當然,這個紅字對這個以掃.
在歷史上又有另一個事件出現.
記不記得?.
那個紅豆湯事件.
紅豆湯事件就是.
他用一碗紅豆湯出賣了自己長子的名分.
聖經又怎樣說呢.
他這樣說,我們一起讀一二三.
以掃對雅各說.
求你把這紅豆湯給我喝吧.
因為我疲乏得很.
因此,以掃的名字又叫以東.
為什麼他要叫以東呢.
就是因為那個紅豆湯的紅字.
紅豆湯改變了他的一生.
改變了他和他的雙胞胎兄弟的關係.
但其實你看看創世記.
他們兩個的民族.
兩個的後人都得到上帝很大的祝福.
很多牛羊,他們很富裕.
到了三十六章的時候.

$^{121}$創世記跟我們說.
因為兩個人的財物群畜甚多.
寄居的地方容不下他們.
其實他們都住在一起.
後來相遇了,然後住在一起.
所以不能同居.
不是說磨擦.
聖經說他們太多財寶了.
所以以掃去了西乙山.
以掃就是以東.
聖經說以掃的後裔為什麼叫以東呢.
完全是這個意思.
因為他們後期去了西乙山.
雅各就是以色列.
以掃就是以東.
以東去了西乙山.
以色列就沒有去到.
西乙山是個什麼地方呢.
他們去了西乙山的地方呢.
西乙可以說是現在死海的南方.
死海的南方.
就是以色列的地理.
死海的南方.
下面我寫了西乙山的地方.
那是以東的地方.
如果相對今天來說那是什麼地方呢.
在座我相信也有不少弟兄姊妹去過.
約旦嘛.
約旦很近當時的首都.
就是Petra.
貝特拉.
有去過貝特拉的.
有些人去過.
有些進不了以色列.
我們進了約旦.
我們也去過聖地.
這個Petra的地方.
大家如果好幾年前我們和弟兄姊妹去過的地方.
它叫做玫瑰紅城.
Rose Red City.

$^{161}$是一個很漂亮的城市.
是一個世界七大奇景之一.
如果你看那些奪寶奇兵.
變形金剛.
在那裡取景.
你會見到在這個這麼宏偉的地方.
有些Transformers走出來.
然後又有奪寶奇兵.
那個石頭滾來滾去.
因為這個你要走一個峽谷才能見到這個.
這個你見到的叫做卡芝利神廟.
所謂神廟其實和日本一樣.
都是那些陵墓.
都是皇帝的陵墓.
很壯觀的.
在這裡你經過那個峽谷.
旁邊那幅很幼的.
慢慢走進去就變成很大的環境之外.
那個只是頭塞.
那個西義山是山來的.
要上山的.
在那裡就要騎駱駝.
騎完駱駝之後又會騎小駱駝.
騎駱騎駱上山.
上到山上面就要見到什麼呢.
如果大家有些不是騎駱的.
有些是走上去的.
我知道.
Rowena有些是走上去的.
我是騎駱上去的.
Arena有些是走上去的.
有沒有人坐飛機上去的.
我不知道她什麼時候會坐升降機上去.
上到那裡就會見到另一個環境.
那個環境其實是一個修道院.
我記得.
那個叫修道院.
走到很高.
見到有個這樣的建築物.
修道院.

$^{201}$當我一路去看的時候.
我想想.
哪有人住到這麼高的.
住在山上面的.
原來這個就是西義山.
就是以素的後裔.
他帶著那群以東人.
就是去到這個地方居住.
居住的時候.
在古代.
現在來說.
你覺得很不行.
但是當時古代來說.
其實給他們一個非常安全.
好像堡壘一樣.
他們是靠著這些山.
去居住有個安全感.
敵人來到.
他是易守難攻的.
敵人來到的時候.
攻他不陷的.
一個這麼好的戰略據點.
為什麼最終都是不攻而破呢?.
為什麼以東人在歷史上.
真是蕭生的跡呢?.
剛纔我們看《創世記》.
很多他的足跡.
但是如果我們看一個小先知書.
叫做《俄巴底亞書》.
其實就給了我們很多的答案.
俄巴底亞本身是神的僕人的意思.
他是一個敬拜耶和華的人的意思.
他這本書在聖經來說.
是很短的.
只有一章.
21節聖經.
很值得看的.
你看完就說看完一本書了.
俄巴底亞書在舊約.
可能你也找不到.

$^{241}$因為只有一版.
不過現在不用了.
現在是用電腦看的.
在頭一兩節.
我們一起讀好嗎?.
他說俄巴底亞所見的異象.
主耶和華指示以東者揚說.
我們從耶和華那裡聽見了一個信息.
有一位使者被派往列國去說.
我們起來與以東爭戰.
看啊.
我必使你在列國當中成為弱小的.
你必大受藐視.
在頭兩節聖經告訴我們.
俄巴底亞所見的異象.
就是上帝要興起很多的列國.
與以東爭戰.
剛纔我們看過地理環境.
有沒有人敢與以東爭戰呢?.
沒有.
就住在高山上.
你怎麼打呢?.
你走上去.
不小心把石頭滾下來.
你也避不了.
你也玩完了.
你基本上應該是打不了他.
但是現在上帝興起所有的列國.
說要一起打你.
為什麼?.
做什麼?.
他在山上.
應該挺好的.
他不是有些什麼.
因為他在俄巴底亞書的第三第四節.
給我們看到一個原因.
我們一起讀.
上帝說.
你心中的慪氣欺騙了你.
你這住在岩石的隱祕處.

$^{281}$居所的高處的.
你心裡說.
誰能把我拉下地呢?.
你雖如鷹高飛.
在聲速之間答我.
我也必從那裡把你拉下來.
這是瑤華的宣告.
在經文當中.
你會見到.
以東在上帝眼中是一個什麼民族?.
是驕傲的民族.
他當時有什麼驕傲呢?.
就是因為他住在西義山.
他覺得地勢險要.
高山環繞.
易守難攻.
但是上帝告訴他.
當你再說.
我是無敵的.
沒有人能夠攻下我的時候.
其實上帝是可以的.
當然他背後也做了很多.
上帝眼中不好的事情.
稍後荷包底亞就會再告訴我們.
為什麼他這麼驕傲呢?.
他不只是款式驕傲.
他做的東西很驕傲很霸道.
我們沒看之前我想問一下.
有沒有人覺得自己是很驕傲的?.
有人可以舉手嗎?.
沒有.
只有一個.
其餘都覺得自己不驕傲.
那你們是最驕傲的了.
這個問題很tricky.
為什麼承認自己驕傲是那麼困難呢?.
叫大家剛纔舉手.
我怕醜而已.
其實我們是需要承認我們是驕傲的.
我再問多一次.

$^{321}$有沒有人覺得自己不驕傲的舉手呢?.
很聰明.
真的不舉手.
驕不驕傲呢?.
有沒有人覺得只是驕傲呢?.
因為驕傲本身很tricky.
他就是讓你見不到自己驕傲.
明白嗎?.
驕傲和其他罪很不一樣的.
譬如你說我憤怒.
人人都知道你憤怒.
走過來都要你著火.
貪婪,你知道的,你只知道自己貪婪.
說謊,說人家壞話.
你知道的,你做完這些罪是知道的.
那些情慾,諸如此類的.
但是驕傲呢.
是你很難察覺到自己是驕傲.
這樣纔是驕傲.
驕傲本身就是讓你見不到自己的問題.
高估了自己的質素.
等等.
夫妻間很多時候都會有這些磨擦發生.
因為人的驕傲.
是不是?.
有一個丈夫就和太太說.
他說你真是太好了.
娶到一個這麼好的丈夫.
太太立刻說.
真的?.
什麼事?為什麼?.
他說.
你看看我.
我又不吸菸,又不喝酒.
又不賭錢,又不去嫖.
絕世好男人,是不是?.
你嫁到這樣的老公.
那太太.
和大家在座的姊妹一樣.
目定口呆,都喝醉了.

$^{361}$跟著跟她說一句.
我覺得你很驕傲.
是啊.
男性有時候是很驕傲.
為什麼?.
那男性問.
但是驕傲又不知道.
跟著就很不明白地問那位太太.
為什麼我驕傲啊.
我剛纔又不.
喝酒又不成,多好啊.
她說你驕傲是因為你不做家務.
是啊,很多時候是看不到自己的問題.
你看到自己有很多好東西.
但是你沒有做家務.
這個是問題,你看到嗎?.
看不到.
這個是一個很驕傲的男人.
我們希望我們LG的弟兄.
繼續不要這麼驕傲.
不過我見到你們,我很驕傲.
一個教會如是.
不要說夫妻間.
教會也是.
如果教會驕傲的時候.
很難得到上帝的祝福.
在啟示錄裡告訴你.
有個教會叫老底家.
這個教會怎麼說呢?.
他說你說.
我是富足的.
已經發了財.
毫無缺乏.
卻不知你是困苦的.
可憐的.
貧窮的.
黑眼的.
赤身的.
當老底家覺得自己發達了.
我很行.

$^{401}$我一無所缺.
OK的,有這個感覺.
但是他看不到自己的缺點.
你不喫不喝就OK.
但是你看不到自己.
原來上帝告訴我們.
我們其實是可憐的.
我們是貧窮的.
我們是黑眼的.
我們是赤身的.
驕傲最大的特徵.
就是讓人看不到自己的光景.
而依賴了一些.
我們不應該去依賴的東西.
以東人的安全感.
我們剛纔說了.
他就是靠這個山.
但是上帝就說.
我可以拉你下來.
經文告訴我們.
到底他們驕傲的是什麼呢?.
他也這麼說.
他說以素的隱祕處.
何竟被索取.
他隱藏的寶物.
何竟被查出.
第一件事讓我們看到.
原來以東在高山上有很多財寶.
很多寶物.
怎麼會一個人在高山上有那麼多財寶.
他一定是怎樣的.
他肯定是搶的.
他肯定是去拿別人的東西.
他說曾與你結盟的.
都把你送到邊界.
本來與你有好的.
卻欺騙你勝過你.
與你同籍的在你腳下.
暗設網羅.
你卻毫不曉得.

$^{441}$另一件事讓我們看到.
以東不單靠這些財寶.
又住高山又有錢.
半山區那些很貴的地方.
他也有些盟友.
但是他靠著這些關係.
他以為有很多的支持和保障.
但是這裡說.
驕傲的人很容易.
判錯了他和人的關係.
他也覺得大家都是站在我這邊的.
但是神提醒他們.
其實人際關係不是你真的可以去依靠的.
你要依靠的是上帝的關係.
有個姐妹.
跟我們分享在國語部.
她說她有個師姐.
她師姐本來和她兒子的關係很好.
但是發生了一件事.
很離奇.
他們的關係決裂了.
崩潰了她的師姐.
什麼事呢.
她兒子十八歲.
他十八歲的時候.
他有個女朋友.
他女朋友就這樣拍拖.
因為他十八歲.
他十八歲他就跟他媽媽說.
他說接下來.
我會去做一個變性手術.
我要做變女人.
他媽媽覺得很奇怪.
為什麼呢.
你現在不是有個女朋友嗎.
你要拍拖嗎.
但是因為這個關係的緣故.
他女朋友就告訴他.
其實我是男性.
你是女性.

$^{481}$我們應該去交換我們的性別.
你應該去做那個手術.
變女人.
我是男人.
他媽媽聽到的時候很愕然.
一個十八歲的小孩.
為什麼會做出這種衝動的行為呢.
他又為什麼會告訴他媽媽呢.
因為他跟他媽媽說.
我很尊重你的.
我現在十八歲.
我會變女人了.
但是你想不想保留我的精子呢.
將來你有後裔.
我有個兒子.
他媽媽聽完之後更加氣憤.
更加荒謬.
這麼好的家庭有一個好的關係.
因為這件事.
他和他兒子的關係很差.
而他的丈夫是支持他兒子的.
他和丈夫的關係也不是很好.
他人生當中是一個非常低落的時候.
當中當人去驕傲的時候.
在一個關係裡面.
有時候可能就被人騙了.
他如果變成女人的時候.
那怎麼辦呢.
如果他這個關係到最後走不下去了.
他真的做了手術了.
那他變不變得回來呢.
他變不回來的.
很多很多事情.
他媽媽跟他說.
但是他聽不到.
他覺得他媽媽是傻的.
當人去到驕傲.
想做什麼就做什麼的時候.
最後是一個災禍.
當我們教會覺得我們.

$^{521}$都是一無所缺的時候.
我們都好像很好的時候.
我們見不見到我們教會.
其實都有一個光景.
以致我們不可以去驕傲.
聖經裡面說.
耳東耶和華宣告說.
到哪一天.
我豈不從耳東滅絕智慧人.
豈不從耳塞除掉聰明人嗎?.
你們的勇士要驚慌.
他們當中有智慧人.
有聰明人.
有勇士.
耳東人不是一般的民族.
他們很精壯.
很厲害.
當日一個這麼小的日本.
能夠有這樣的破壞力.
他們是有勇士,有智慧人.
大家都知道東京,大阪.
幾十個這麼好的城市.
要征服其他國家.
真的有可能.
但當人去驕傲的時候.
聖經告訴我們.
上帝要除掉他們.
今天我相信我們要學習謙卑.
絕對不是一個我們人可以去努力的事情.
我覺得這個一定是聖靈的工作.
當我們再看回聖經.
聖經告訴我們.
耳東人他那個極度的驕傲.
其實是來自他內心一種很大的不安和自卑感.
而這個後面的自卑感和不安.
我們很多時候是察覺不到的.
今天如果你在上帝面前.
你都發覺.
你都會有驕傲的情況.
今天我們一起去學習.

$^{561}$我們讀一次聖經.
一起讀一次真言.
驕傲在敗壞之先.
狂心在跌倒之前.
滅亡之先.
人心高傲.
尊榮之先.
必有謙卑.
驕傲的人有什麼結果.
大家都知道.
人靠的是上帝的工作.
我們再不能靠自己.
全方一樣.
我們盡心盡意.
但是我們將一切榮耀一切的結果交給上帝.
上兩個星期當我去到我們GC.
我們那個崇拜的時候.
大家都有去到.
我們有安納典禮.
在那一刻我真的覺得.
哇!教會是真的很好.
真的很好.
去到那個地方.
我真的有點驕傲.
我見到大家其實我很驕傲.
我引以為傲見到你們.
但當我很多時候想的時候.
其實我們教會是什麼的光景.
我覺得這個應該是我們的出發點.
過去我們真的看到上帝比我們好多很多.
但我們不要錯判了我們的關係.
我們不是靠我們的勇士們.
我們有很多勇士.
我們不是靠我們的offering.
我們不是靠我們弟兄姊妹的關係.
我們真的靠我們和上帝的一切事情.
我們不要有錯誤的focus.
願意我們過去一年上帝的帶領.
我們接下來的一年.
我們更加一起謙卑.

$^{601}$我們去跟從我們上帝.
我們一起禱告.
當我們來到您面前的時候我們再一次願意.
去承認我們自己.
有衣無衣服我們都會有驕傲.
當很多人都說.
婚姻教會LG很好.
弟兄姊妹又有恩賜.
一無所缺.
但上帝叫我們不要走老底家的路.
然後我們都看到我們在哪一方面.
我們真的貧窮.
哪一方面我們真的赤身.
願意我們不要被魔鬼欺騙.
願意我們在這個世代.
我們看到仍然有很多的需要.
無論是日本剛纔我們說的名字.
我們身邊有很多的人.
他們struggle著他們的身份.
他們struggle著他們的身體.
主啊 在這些struggle當中.
讓我們去靠著禰謙卑的去跟隨禰.
求禰給我們一份謙卑的心.
禱告奉主名求 阿們.
\newpage



\section{加拉太書 5:7-12}
\label{sec:032305}
\textbf{直奔天路 (華)}~~~~葛國光牧師
\newline
\newline
連結: \href{http://www.church.com.hk/acms/content.asp?site=cdc\&op=show\&type=product\&code=032305\&layout=sermon}{\texttt{032305}} ~~~~ 語音日期: 2026-06-14
\newline
\newline
\hyperref[sec:032304]{\small{< < < PREV SERMON < < <}}
~
\hyperref[sec:index_chronic]{\small{[返順時目]}}
~
\hyperref[sec:032306]{\small{> > > NEXT SERMON > > >}}
\newline
\newline
加拉太書 5:7-12
\newline
\begin{longtable}{cl}
\hline
\hline
章節 & 經文 (和合本修訂版)\\
\hline
5:7 & \begin{tabularx}{0.7\textwidth}{X} 你們向來跑得好，誰攔阻了你們，使你們不順從真理呢？ \end{tabularx} \\ \\ \relax
5:8 & \begin{tabularx}{0.7\textwidth}{X} 這樣的勸導不是出於那召你們的。 \end{tabularx} \\ \\ \relax
5:9 & \begin{tabularx}{0.7\textwidth}{X} 一點麵酵能使全團都發起來。 \end{tabularx} \\ \\ \relax
5:10 & \begin{tabularx}{0.7\textwidth}{X} 我在主裡深信你們必不懷別樣的心；但騷擾你們的，無論是誰，必須承受懲罰。 \end{tabularx} \\ \\ \relax
5:11 & \begin{tabularx}{0.7\textwidth}{X} 弟兄們，我若仍舊傳割禮，為甚麼還受迫害呢？若是這樣，十字架絆倒人的地方就沒有了。 \end{tabularx} \\ \\ \relax
5:12 & \begin{tabularx}{0.7\textwidth}{X} 恨不得那騷擾你們的人把自己閹割了。 \end{tabularx} \\ \\
[1ex]
\hline
\hline
\end{longtable}
$^{1}$保羅在加拿大書上次我們講到說.
保羅又再提.
你看整個加拿大書六章的聖經裡面.
他花了五章的聖經裡面一直在反覆的講.
反覆的講這件事情.
也許對你我來講.
今天陰性稱義這件事情是一個好像是理所當然的事情.
但是我想今天我也會告訴你.
他已經不是理所當然.
因為在整個的特別是西方基督教世界裡面.
這已經是被衝擊的到一個地步.
現在很多的教會要怎麼樣能夠守得住.
這都變成是一個挑戰.
今天我們華人教會面對在特別是美國.
這個所謂的基督教的大本營.
我們更是需要能夠在真理上面能夠站得住.
所以保羅在這邊一再的講.
因為當時候這些猶太的基督徒.
他們為了受到他們自己的猶太教的影響.
所以他就勸那些外邦的信徒.
外邦就是非猶太人.
勸這些外邦的信徒說.
你們必須要受割離才能夠得救.
割離是猶太人男孩子生出來第八天之後.
就必須接受的一個割離.
瞭解這意思嗎.
懂嗎.
可能有一些還剛來的或不太清楚.
就是以色列的男孩生出來第八天就要割包皮.
這個叫做割離.
這是上帝對以色列人所立的一個約.
用這個約來代表著神跟他們之間的關係.
所以猶太人有這種很強烈的觀念.
認為這是上帝在舊約裡面跟他們所立的這一個約.
因此世世代代我們都要去遵守這個約.
可是他不曉得.
因為他們今天也不相信.
耶穌基督這一位他們所盼望的彌賽亞已經來了.
這就是今天我們所謂的基督教.
我們相信耶穌基督使我們能夠得救.

$^{41}$Amen.
好.
那麼保羅基亞也講到這件事情之後.
接下來第五章第七到十二節.
保羅基亞提到了.
他說你們向來跑得很好啊.
有誰攔阻你們 叫你們不順從真理呢.
這樣的勸導不是出於那招你們的.
在第七節第八節保羅就繼續的講到.
他說你們過去向來跑得好.
什麼意思.
保羅在這裡把我們走這條天路.
好像是我們在跑天路一樣.
還記得保羅也曾經用賽跑來形容我們今天在人生道路當中.
是不是.
我們要向著標杆怎麼樣.
直跑.
是不是.
所以保羅提醒我們.
他說你這些.
他說這一家的基督徒.
他說你們過去向來跑得好.
跑得好.
這個字在原文裡面的詞態.
它是所謂的過去不完成式.
什麼意思.
在希臘文裡面這個意思就是講說.
在過去的一段時間裡面.
他持續的一個動作.
所以保羅講說你們過去跑得好的意思就是說.
他說你們向來跑得好的意思就是說.
你們過去一直跑得很好.
懂這意思吧.
你們過去信仰.
你們過去從受洗之後.
你們信耶穌之後.
你們一直都跑得很好.
他說加勒泰人他鼓勵他們.
但是接著他說什麼.
他說有誰攔阻你們.

$^{81}$叫你們不順從真理呢.
所以在這個地方他提醒.
他說有誰攔阻你們呢.
你們過去一直跑得很好.
是誰攔阻了你們呢.
剛纔我們講到美國的教會.
美國的教會 西方的教會.
從歐洲的教會開始.
受到新派神學 自由派神學的影響.
帶來的影響非常的大.
直到今天.
我們來看看.
自由派的神學.
我等一下會跟大家一起解釋.
什麼叫做自由派的神學.
這邊是一把自由派的神學.
跟福音派的神學做了一個對比.
一個有七項對比.
我沒有時間整個講.
我就highlight幾個.
第一個 比如說關於聖經的權威.
自由派的神學呢.
他們認為這個聖經是一個記錄.
記載古代的這些人的一個記錄.
你看 當這樣講的時候.
那麼所以聖經就不是一個絕對性的東西了.
懂我意思嗎.
那跟我們今天所謂福音派的教會.
我們所強調的是什麼.
聖經是無誤的.
聖經是神的啟示.
是上帝的話語.
聖經不會改變.
從過去一直到永遠.
神的話永遠不改變.
你相信這件事情嗎 弟兄姊妹.
當你在讀聖經的時候.
你是把它當作一本歷史的記載.
還是把它當作是上帝的話語.
是神的啟示.

$^{121}$這是非常非常重要一件事情 對不對.
對我們福音派來講.
這是沒有改變的餘地.
馬丁路德改教的時候.
就把這件事情清楚地講下來.
是不是 改變天主教會.
教宗也好 或者這些天主的主教也好.
他們用自己的方式來解釋聖經.
甚至不按照聖經的本意.
那所以馬丁路德改教.
讓我們看見聖經是沒有錯誤的.
OK 好 那弟兄姊妹.
第二個 關於耶穌基督的神性.
自由派他們認為什麼.
耶穌是一個在歷史上面.
所出現的一個完美的人.
所以耶穌是不是神.
不是神.
所以在他們的觀念當中.
耶穌不是神.
親愛的弟兄姊妹.
如果耶穌不是神.
你我信的是枉然的.
因為耶穌定在實相上面.
祂死了 對不對.
祂是人啊.
祂沒有再復活.
所以講到這些東西的時候.
那耶穌基督的復活.
耶穌基督的再來這些通通都沒有了.
瞭解這個意思嗎.
因此那我們所相信的這位耶穌是誰.
祂是道成了肉身.
充充滿滿 有恩典 有真理 住在我們中間.
好了 第五項 我們要講.
那麼基督徒.
今天我們信了耶穌 我們做什麼.
我們在地上的使命是什麼.
我們存在這個地上的目的是什麼.
信了耶穌之後就坐在教堂裡面.

$^{161}$每個禮拜聽牧師講道.
然後聚完會之後.
我們不唱三一誦了 我們唱什麼.
唱祝福誦 祖導文.
唱完之後 回去之後.
從此再過另外一個生活.
教堂的那段一個半小時沒了.
不是的.
可是自由派的神學認為什麼呢.
我們基督徒在地上.
我們要積極地參與在地上社會的工作.
去幫助弱勢的 好不好.
但是去幫助弱勢 去解決貧窮問題.
去積極地參與政治的目的.
是為了要在地上建立一個天國.
懂這個差別嗎.
但是身心上告訴我們.
特別你看看我們啟示那邊.
地上所有的一切是會越來越好.
還是越來越差.
明天會更好嗎.
明天的股市市場會更好.
會嗎.
可能好了一下 又沒了.
但是我告訴你 身心上讓我們看見.
這個世界有一天.
最後統統要銷燬.
這個世界會越來越糟糕.
這個世界會越來越墮落 是不是.
所以 不是說基督徒不要管這個世界.
脫離這個世界 好像過一個與世無關的生活.
不是 但是不是像自由神學一樣.
他們認為說我們要積極地在地上.
甚至他們贊成要透過革命.
透過這些所謂的解放.
這個自由神學跟解放神學.
很多的地方受到他們影響.
所以教會也會去參與這些的.
怎麼樣去參與這些戰爭 參與這些事情.
好 再給我準備.

$^{201}$最後一個.
那麼 講到這裡.
我們基督徒的使命是什麼.
來吧 佛地28章49節20節會背吧.
來 預備 請.
所以 你們要去使萬民做我的門徒.
奉父子聖靈的名給他們施血.
凡我所吩咐的你們的都教訓.
他們遵守 我就長與你們同在.
直到世界的末了.
好 丟準備.
大使命是我們基督徒在地上的使命.
這些世界怎麼樣子.
有一天都要改變的.
不是 again 我再次強調.
不是要變成另外一個極端.
基督徒都是一群與世隔絕的人.
不是.
但是基督徒你不要忘記.
我們關心社會.
我們今天關心貧窮的人.
我們去做homeless的工作.
我們今天去發放這些物資給這些有需要的人.
我們做這一切的目的是為了什麼.
使萬民做主的門徒.
社會關懷不傳福音就變成社會福音.
這就是自由神學所帶來的結果.
但是我們社會關懷.
我們目的是為.
最後是為了達到福音的目的.
對不對.
好 那麼最後一個.
這個是跟今天美國整個趨勢很有關係的.
就自由神學強調宗教多元與真理.
多元.
你聽這個多元這名稱.
六月是所謂的驕傲月 懂嗎.
什麼叫驕傲月.
就是我以我自己是同性的.
同性戀 我驕傲 是吧.

$^{241}$我以我自己的 我覺得.
我雖然身體是一個男的.
但是我覺得我是女的.
所以我去做了變性手術.
我就變成女了.
I'm proud of myself.
驕傲月.
最近在美國各地舉行同志驕傲月.
哎呀 我在禮拜天 禱告會的時候跟大家放.
你看見特別是在紐約的那個.
那個多少人啊 幾十萬人啊.
那麼在那個地方遊行的在旁邊圍觀的啦.
然後你就看見那個主持人帶著大家一起一直講.
Love is love.
什麼意思啊 愛就是愛.
什麼叫愛就是愛.
講到後來.
也就是說我跟誰愛 只要我們愛就可以了.
那什麼叫做愛.
愛可以同性愛.
愛也可以跟我的狗狗愛 你懂意思嗎.
歐洲已經有些國家走到這個地步.
就是從接納多元 接納同性戀 接納跨性別.
到最後就接納人壽交 懂嗎.
可悲 當人離開了真理之後.
人就連畜生怎麼樣都不如.
這個意思不是說我們基督徒.
我們恨惡同性戀的人 懂這意思嗎.
我們恨惡罪.
但是我們愛罪人 懂這個差別嗎.
神愛世人 但是神不愛罪.
這是要很清楚分別出來的.
所以很多一些不明白的人.
他們定罪基督徒.
他們認為基督徒說.
你們都是不接納 我們不愛 我們不是.
我們愛 但是我們不能夠容許罪 懂這意思嗎.
所以這個是真理.
真理就是真理.
很多時候真理是隻有一個 不可能有多個.

$^{281}$瞭解意思嗎.
可是自由神學多元主義的一個結果就是.
只要我喜歡 沒有什麼不可以.
好了 那丟什麼.
所以在基督教裡面.
長老會 衛理公會 還有信義會.
這幾個大的宗派.
全部都因為這一個自由神學的影響.
在他們總會里面最後形成一些人.
這些很厲害.
他們花了幾十年的工作.
最後進入總會里面掌握了多數.
最後投票 決定我們總會的立場.
是決定接納同性戀 案例同性戀的牧師.
這一決定一下分裂.
凡是不贊成的 最後就離開.
那這是今天美國的這些主流教會.
一個一個的衰退的原因.
很多的基督徒離開了這些的主流教會.
有的離開了教派.
他們自己成立新的教派.
這是為什麼我在這邊語重心長地要講.
想想看當年約翰·魏斯里.
英國大復興的時候.
衛理公會是多麼樣子的被上帝使用.
想想看 加爾文.
這也是其中 改教的其中一個祖師.
他在改教的時候帶來後面長老會的影響.
長老會一般來講是非常看重聖經的.
但是今天的長老會 今天的衛理公會.
信義會 馬丁·路德帶出來的信義會.
全都被自由神學怎麼樣 腐敗了.
所以講到這裡 他說有誰攔住你們呢.
丟子妹 那你說牧師這些東西太高大上了.
我這個平民小族跟我有什麼關係呢.
丟子妹 那就是我們要從聖經裡面去想到.
有誰攔住我們在這條信仰的道路上奔跑呢.
可能不是主有主義.
但是可能是誰.
我們當中還沒有結婚的.

$^{321}$年輕人.
我記得以前我在帶年輕人的時候.
好些很愛主的姐妹.
很可惜.
後來碰到了不信主的男生.
最後結婚之後.
能夠保住信仰的已經不錯了.
很多最後就跟著先生就走了 離開主了.
那就是信仰上面的攔住.
丟子妹 你瞭解我意思嗎.
在我們的家庭裡面 在我們的婚姻當中.
會不會有這種情形呢.
哎呀 一個很愛主.
另外一個就在後面扯後腿 有沒有.
當然 不要誤會我意思.
那好 你既然攔住我.
那我覺得怎麼樣.
牧師今天講了要認清這個攔住 是不是.
那我就隔絕了.
不是 不是 千萬不要亂解釋.
我的意思就是說 今天可能你會覺得你的利益一敗.
在你的信仰上面成為你的攔住.
你更有責任 你要堅持的往前奔跑.
同時呢 你要積極的用你的生命見證.
來帶領他 影響他跟你一起跑 懂這意思嗎.
你不要說 你看你都不讀經禱告.
每天早上起來 你都在那邊睡大覺.
我參加教會的六點鐘 六點半的晨導.
你就在旁邊睡覺 十年了.
你都是這個樣子.
哎呀 算了 我不理你了.
我今天自己追求吧.
親愛的劉子梅 不要放棄.
不要放棄 持續下去.
繼續的為他禱告 繼續為你的另外一半禱告.
繼續的你自己不要受到影響.
你仍然要保持追求主的信仰.
Amen.
同樣的劉子梅.
這個攔住可能是誰呢.

$^{361}$這個攔住可能是你的兒女.
會不會.
很多弟兄姊妹.
沒有孩子的時候很積極很追求.
有了孩子之後呢.
全世界的忠心都在我的兒子身上.
一切的忠心都在女兒身上.
劉子梅 不可以如此.
繼續的抓住神.
那個攔住我們跟隨神的.
很多在我們跑天路上面都會來影響我們.
有一個在臺灣根據統計.
他是在臺灣宣教的歷史最長的宣教士.
你們猜是誰.
他早期的時候曾經到中國大陸去宣教.
12歲的時候.
他就聽到這個護照要他去中國.
他就預備心了.
就21歲的那年.
那麼他就單身一個人就到了中國去.
可是剛好碰到國共戰爭.
那麼最後他就跟著國民黨的政府.
就逃到了臺灣.
這個人是誰.
這個人叫做彭蒙慧.
臺灣來的人肯定都認識他.
他在臺灣創辦了空中英語雜誌.
就是傳播協會.
他曾經 他不是 你看他長得很漂亮吧.
有點反應嘛.
不是我 他長得很漂亮吧.
不是那個嫁不出去的吧.
可是他不是沒有機會.
他曾經在臺灣碰到一個門諾會.
從美國來的 在門諾會的一個醫生.
在那邊做一診 做短期宣教.
結果兩個人就相戀了.
那這個男的 這個醫生.
他後來就是要調到越南.
去越南也是一個短期的去那邊做宣教.

$^{401}$做服侍.
那麼他就跟彭蒙慧講.
他說你一定要等我.
他說等我去完越南服務完之後.
我一定會回來娶你.
男人的話可以相信嗎.
OK 出乎你意料之外.
他果真短期之後他回來了.
回來之後照著他所應許的.
他就跟Doris跟他求婚 跟彭蒙慧求婚.
可是他希望彭蒙慧跟他結婚之後.
跟他一起回美國.
開始他們自己的兩個人的婚姻的生活.
但是Doris剛纔講說.
她說不 神給我很清楚的負擔 我要留在臺灣.
但是這個醫生他不願意.
最後他割捨了這一段的戀情.
繼續地留在臺灣.
一直到他98歲的時候.
好像是2024年還是2024年.
他蒙祖崇照到祖裡面去了.
不知道有多少人.
不僅僅是聽到他用英語教英語.
更是有多少的人受到他的影響.
而甚至是耶穌.
所以他最後講一段話.
他說我跟隨上帝帶領我走的路.
看起來我捨棄了許多.
可是我得到的卻是超過我所能想象的.
我這一生不後悔.
親愛的同志們.
選擇來跟從上帝.
不是代表就是一帆風順.
這個過程當中會有人攔住我們.
會有人會有一些思想觀念.
社會環境或者任何的人事物.
會來攔住我們叫我們繼續的跑.
繼續的跑好.
丟著沒.
這個人事物是誰呢.

$^{441}$最後接著.
第九節第四節我們起來唸.
一點面教能使全團發起來.
我們有沒有在座會做饅頭做包子的.
凡是要發的面的會嗎.
舉下手.
我看看有多少新好來人.
都是姐妹.
有有有有有有有人舉得很高.
你看這個世光他也舉得很高.
表示他是會做的.
請問弟兄姊妹.
你做只要要發的面.
你需要放什麼.
酵粉是不是.
放多少.
這纔是真正的好樣.
放多少.
Average 多少.
1\%.
我經常大概就是四杯麵粉.
One teaspoon and a half.
的意思.
一匙半的意粉.
不放呢.
剛開始學做的時候.
忘了放酵粉.
結果四個小時之後.
出來的麵包變什麼.
自動麵包機.
就是一團死麵.
懂這意思嗎.
沒有發.
所以他說一點面教.
就能夠使全團發起來.
我不知道保羅有沒有做過麵包.
但是這一個形容真是太好.
是不是.
保羅在這裡講的意思就是說.
一點點的面教.

$^{481}$這個面教是有這個擴散滲透的能力的.
這個面教可能就是一點點的罪.
在我們的生命當中有一個小罪.
這個小罪你剛開始你不注意它.
你輕忽它.
你覺得沒什麼關係.
哎呀我只不過偶爾不太就是偶爾會撒個謊.
這也不是什麼.
哪有人不撒謊嗎.
我偶爾而已.
我也不是常常對不對.
現在就準備很多時候我們會這樣子是不是.
好一點面教就能夠使全團發起來.
在哥林多前書第五章第六節.
我們一起來讀好不好來.
你們這自誇是不好的.
豈不知一點面教能使全團發起來.
保羅在這地方講這個一點面教.
他就講了一次.
他這人講什麼呢.
當時候哥林多的教會有人娶繼母為妻.
請告訴我可不可以娶繼母為妻.
為什麼.
為什麼不可以.
今天那個自由派神學家就問你.
why not.
她也不是我mother.
她是我繼母.
你懂這意思嗎.
可以嗎.
why.
為什麼.
亂來了.
亂倫了.
是不是.
現在就這麼一看到一件事沒有.
保羅說你們當中竟然有人娶繼母為妻.
這是淫亂.
然後你做教會的.
你知道這件事情.

$^{521}$你不處理.
一點面教能使全團發起來.
丟著美.
今天的美國教會.
最大的問題就是如此.
好像覺得這一個不是什麼大問題.
今天也許不一定是所娶繼母為妻.
但是在美國的教會里面.
這種不合聖經的離婚.
越來越多.
多到一個地步.
很多牧師也是如此.
因此你講臺牧者不敢講這個道.
你懂意思嗎.
你不敢講這個道.
你不敢講這方面聖經的規定.
你不敢用耶穌的話說.
聖經怎麼說.
主耶穌說什麼.
若不是為淫亂的關係.
對不對.
神所配合的人不可以分開.
但是今天聖經的話.
真理擺在一邊.
因為多數覺少數覺.
瞭解這意思嗎.
因為現在多數人都如此.
所以牧師你不要講得罪人的話.
對不對.
你這樣講下去的話.
可能教會里面很多人要離開教會.
可能換教會.
丟著美.
這就是今天很多的傳道人.
不敢傳講真理.
那我這樣講的意思.
不是說我們今天只站在公義的一面.
我們不能夠體恤人的軟弱.
可是親愛的同事們.
我要套剛纔這個同性戀講的.

$^{561}$Love is love 是不是.
Sin is sin.
罪就是罪.
聖經上告訴我們說.
罪要怎麼解決呢.
怎麼解決.
在這個公義的聖潔的神面前.
罪使得我們沒有辦法來他面前.
那罪怎麼解決呢.
我們落任自己的罪.
神是信實的.
是公義的.
必要赦免我們的罪.
洗盡我們一切的不義.
以色列因為亞甘一個人的犯罪.
全軍覆沒.
提醒我們.
教會是一個整體.
保羅在這邊講的意思就是說.
教會是一個整體.
你一點的面教會使全團發起來.
你一個小組.
你一個團體裡面.
你不把這個罪的問題解決.
你包容這些的事情.
最後是使你整個的小組.
整個團體.
你通通都會受虧損.
所以教會必須要聖潔.
教會不聖潔.
今天在美國這個社會.
你可以看見外邦人.
世界上面不信主的人.
那亂七八糟太多了.
是不是.
現在年輕人.
很多都是婚前同居.
你問他說.
你們要不要結婚啊.
他說.

$^{601}$再看唄.
再看什麼.
為什麼再看.
因為看看我們適不適應.
合不合.
最近我女兒的以前的高中同學.
他們兩個終於決定結婚了.
那這只是其中一個例子.
很多.
如果有機會你問問看你的孩子.
他們周邊的人都是這樣子.
很多人覺得就是.
我們為什麼一定要.
那婚前同居已經不是一個問題.
那請問婚前有關係可不可以啊.
可不可以.
劉哲明你要斬釘截錘講啊.
是不是.
不僅僅你自己.
你對你自己的下一代.
下下代.
你要清楚地告訴他們這些事情.
今天這個世代已經沒有標準科研.
所以教會不能夠有神的同在.
因為教會不聖潔.
神不會跟教會同在.
所以這是今天整個美國社會的問題.
教會軟弱.
因為教會自己沒有力量.
教會自己犯罪.
教會自己有這些問題存在的時候.
教會的講話就講不出來.
你就沒有辦法講出真理來.
所以保羅提醒我們.
一點的面教全談都會發出來.
在我們生命當中也是.
有沒有什麼一點的面教.
這些是不好的.
不合神心意的.
可能是一個誤會的思想.

$^{641}$可能你想說.
我只是偶爾看一下那些.
網絡上面的一些色情的東西.
我只是偶爾.
我沒有變成是經常.
親愛的同志們.
不要容許它存在.
你要跟罪惡征戰到一個地方.
像聖經上說.
征戰到流血.
那是什麼意思.
流血會怎麼樣.
流血的征戰會是什麼一個征戰.
很激烈的.
我們今天對付罪要這樣激烈的態度.
靠著主我們一定可以得勝.
對沒.
那麼.
保祿說如果我們這些.
他說.
我在主裡很信你們必不還別的心.
但繳饒你們的無論是誰.
必擔當他的罪名.
所以保祿在這邊講的意思就是說.
這些傳這種靠著遵守立法而得救的這些人.
他說他們必擔當他的罪名.
那麼.
我現在要講到.
可能今天不僅僅是在美國.
可能在全世界的教會里面.
受到一個很大的影響.
就是世俗化.
英國十九世紀著名的福音派主教.
J.C. Ryan.
中文把它翻成叫萊爾.
他說.
在今天.
在今天普遍有一種世俗化的基督教.
這個是什麼叫做世俗化的基督教呢.
他說這是一種不會冒犯任何人.

$^{681}$不需要付代價.
不求犧牲的廉價基督教.
因為缺乏真實的悔改與實質家.
這樣的信仰在神眼中毫無價值.
世俗化.
剛纔我們所提到的叫做世俗化.
世俗化就是以世界變成是.
世界變成我們的標準.
世界的風氣.
世界的做法.
世界的這些觀念變成我們的追求.
或者說我們學習的對象.
最近我們談到.
有跟姐妹有談到有關穿著的問題.
那麼就談到說.
我們今天請問基督徒姐妹們.
可不可以穿露胸的衣服.
怎麼露法.
你說露多少\%是不是.
請問弟兄.
我們問弟兄好不好.
請問弟兄.
你看到女孩子露胸要露多少\%.
會讓你想露飛飛.
請告訴我.
Kevin.
你比較老實莊重.
我看看你的程度.
多少\%.
好.
最近在Freeway上面有一個廣告.
每次我遠遠地看到的時候.
我都要趕快眼睛閉開.
我不知道在你們走的Freeway裡.
在710吧還是.
一個女的.
我不知道她賣什麼東西.
賣鞋啊.
你知道我在講什麼吧.
OK 你們都看到了.

$^{721}$你看姐妹都很誠實地承認.
弟兄都裝著很聖潔的樣子.
請問弟兄你看到時候.
你會不會想什麼.
會不會.
你不會你就不太正常.
對不對.
會不會.
會 會.
那我請問你.
如果今天教堂裡面.
姐妹你穿了露胸的到教堂裡來.
你是在做什麼呢.
你在試探弟兄們.
是嗎.
是不是.
所以.
我們的一舉一動.
我們的穿著.
都應該要榮耀主.
Amen.
三個原則.
榮神一人立己.
你說天氣很熱啊.
所以我穿的比較少一點.
清涼一點.
對你自己清涼.
但是呢.
你半導別人.
耶穌說半導人的有禍了.
OK 好嗎.
我不是要給大家一個dress code.
像早期.
美國宣教士來臺灣宣這個時候規定的.
姐妹主日崇拜.
需要穿過膝的長裙.
什麼叫過膝的長裙.
我們今天穿過膝的長裙.
來秀一下.
來來來.

$^{761}$這是弟兄們建議的.
來來來嘛.
來.
有沒有這個.
非常女性.
非常lady.
她已經是兩個孩子的媽媽了.
現在整個世界.
男不男 女不女.
我不是反對姐妹不可以穿長褲.
有時候你工作或者有需要的時候.
是不是.
但是現在看到姐妹穿裙子的已經不太多了.
反而現在好像覺得穿裙子的不太正常.
有沒有.
這個就叫做.
這就是叫做文化.
當majority都不這樣做的時候.
你會覺得奇怪.
當majority都那樣做的時候.
你就會覺得那是正常的.
像有什麼.
請你記得一件事情.
真理是不受環境影響的.
天地都會改變.
但是神的話一點一話都不會改變.
好.
那麼最後.
最後保羅在這11節12節.
我們起來唸好不好.
來請弟兄們.
我若仍舊傳歌裡.
為什麼還要受逼迫呢.
若是這樣.
那十字架討厭的地方就沒有了.
恨不得那攪擾你們的人把自己割掘了.
保羅在這邊講了一個很生動的.
他講那個攪擾你們就是那些傳歌裡的.
他說恨不得他把自己割掘了.
懂這意思嗎.

$^{801}$保羅用了同樣的字根來表明一件事情.
就是說這些人我真巴不得你把你自己割掘了.
親愛的同志們.
保羅受逼迫的原因是因為他不傳講這一個錯誤的福音.
不傳講歌裡.
他單單傳講耶穌基督病他十字架.
哥林多前書第二章保羅在那邊說.
他說他不知道別的.
他只知道耶穌基督病他怎麼樣.
定十字架.
因為哥林多人強調哲學.
他們很懂哲學.
他們很有知識.
他們很有這些的學問.
他們講了一大堆這些東西.
但是保羅告訴他們講的意思就是我不講.
我不知道別的.
請問保羅知不知道哲學.
保羅懂不懂希臘哲學.
保羅懂不懂這些天文地理.
保羅非常懂.
但是保羅講的意思就是說.
我懂這一切的事實上的學問.
可是我不去講這些東西.
我單單傳講我的信仰的核心.
我傳講的中心就是耶穌基督病他定十字架.
Amen.
我不傳講成功神學.
雖然成功神學很popular.
讓人聽了很舒服.
教會的人數也可能會增加.
可是耶穌基督病他定十字架.
這是我們今天要傳講的.
各位同學們.
我們會不會為信仰付代價呢.
你們喜歡看電影的人認識這個人嗎.
他曾經演過美國隊長.
星際爭霸戰.
還有很多片子.
是一個很著名的美國演員.

$^{841}$他有五個孩子.
在好萊塢.
而且他有一個妻子.
從頭一直到現在.
這是很難的.
他是一個基督徒.
可是呢.
他是一個虔誠基督徒.
但是他定下一個原則.
他說.
我這雙嘴脣.
是給我自己的妻子.
我絕對不去親吻其他的女人.
做演員的可以這樣子嗎.
你看看好萊塢的影片.
有沒有接吻的鏡頭嗎.
沒有.
幾乎找不到.
所以因為這樣子的原則.
最後他說.
他被fire了.
好幾年的時間.
沒有人找他拍片.
因為他不能.
他這樣等於就是自己把自己的路都斷絕了.
後來神給他開了一條路.
讓他去演反派的角色.
反派的角色.
就不需要去搞這些什麼東西了.
所以什麼美國警長.
什麼這些都.
他通通演反派的角色.
神就繼續的使用他.
一直到如今.
親愛的同志們.
為了信仰.
像他這樣子.
我真是說在好萊塢是一個奇葩.
不僅僅是他.
今天基督徒在各個的行業裡面.

$^{881}$甚至我們自己在我們自己的弟兄姊妹.
我們的家人中間.
在我們的這個是我們的職場裡面.
我們都可能會面臨到要為我們的信仰.
為真理的緣故要受到逼迫.
你願意嗎.
最後給你看一個人.
這個大家肯定就比較熟悉了.
這是一個高中美式足球的教練.
Joseph Kennedy.
他在帶領高中的足球隊.
每一次當他們打球的時候.
他都會帶領這些足球隊員.
一起跪在中場50碼線的地方.
單膝跪下.
在那邊禱告.
後來被告了.
校區的學校告他.
告他說.
他這樣子是違反所謂的.
Separation of Religion and State.
就是宗教跟政教分離.
OK 好.
那麼他也上訴.
從地區的法院.
一路上面打到哪裡呢.
打到最高法院.
最後經過幾年的時間.
總算最高法院判決6比3勝訴.
可是那幾年的時間.
他沒有工作 被學校fire.
這是很辛苦的一件事情.
等到最高法院把他判決勝訴了.
學校雖然只有接受他回去工作.
可是過了沒多久.
他就搬到佛羅里達去了.
我看到他最後的情況就是說.
他就從教練的職場退休下來.
他現在就是出來專心地在speaking.
在教導.

$^{921}$就是說當你面對這樣子的情況的時候.
為了你的信仰的時候.
你應當怎麼樣子來defense.
你要怎麼樣子來堅持.
所以他現在在做這些事情.
但是我要講的就是說.
你可以看見這些基督徒.
這些都是比較算是現在是知名的.
對不對.
可是他也可能就發生在我們身上.
我們在信仰的時候.
為了基督的緣故.
為了生理的緣故.
為了遵守神的話.
你可能要付代價.
你願意付嗎.
這是今天保祿對我們的勉勵.
他不傳別的.
只傳耶穌基督並他定是之家.
他不傳講其他偏離福音的.
偏離聖經的.
他單單的遵守神的道.
求主幫助我們.
在這個末後的時代.
成為一群遵守神旨意的得勝者.
\newpage



\section{撒母耳記上 26:1-25}
\label{sec:032306}
\textbf{兩種人生,兩種死亡}~~~~李思聰牧師
\newline
\newline
連結: \href{http://www.church.com.hk/acms/content.asp?site=cdc\&op=show\&type=product\&code=032306\&layout=sermon}{\texttt{032306}} ~~~~ 語音日期: 2026-06-14
\newline
\newline
\hyperref[sec:032305]{\small{< < < PREV SERMON < < <}}
~
\hyperref[sec:index_chronic]{\small{[返順時目]}}
~
\hyperref[sec:032307]{\small{> > > NEXT SERMON > > >}}
\newline
\newline
撒母耳記上 26:1-25
\newline
\begin{longtable}{cl}
\hline
\hline
章節 & 經文 (和合本修訂版)\\
\hline
26:1 & \begin{tabularx}{0.7\textwidth}{X} 西弗人來到基比亞，到掃羅那裡，說：「大衛不是在荒野東邊的哈基拉山藏著嗎？」 \end{tabularx} \\ \\ \relax
26:2 & \begin{tabularx}{0.7\textwidth}{X} 掃羅動身，帶領以色列人中挑選的三千精兵下到西弗的曠野，要在那裡尋索大衛。 \end{tabularx} \\ \\ \relax
26:3 & \begin{tabularx}{0.7\textwidth}{X} 掃羅在荒野東邊的哈基拉山，在路旁安營。那時大衛住在曠野，看見掃羅到曠野來追趕他， \end{tabularx} \\ \\ \relax
26:4 & \begin{tabularx}{0.7\textwidth}{X} 大衛就派人去探聽，知道掃羅果然來了。 \end{tabularx} \\ \\ \relax
26:5 & \begin{tabularx}{0.7\textwidth}{X} 大衛起來，到掃羅安營的地方，看見掃羅和尼珥的兒子押尼珥元帥躺臥之處；掃羅睡在軍營裡，士兵安營在他周圍。 \end{tabularx} \\ \\ \relax
26:6 & \begin{tabularx}{0.7\textwidth}{X} 大衛對赫人亞希米勒和洗魯雅的兒子約押的兄弟亞比篩說：「誰同我下到掃羅營裡去？」亞比篩說：「我同你下去。」 \end{tabularx} \\ \\ \relax
26:7 & \begin{tabularx}{0.7\textwidth}{X} 於是大衛和亞比篩夜間到了士兵那裡；看哪，掃羅睡在軍營裡，他的槍在頭旁，插在地上。押尼珥和士兵睡在他周圍。 \end{tabularx} \\ \\ \relax
26:8 & \begin{tabularx}{0.7\textwidth}{X} 亞比篩對大衛說：「神將你的仇敵交在你手裡，現在讓我拿槍把他刺透在地上，一刺就成，不用再刺他了。」 \end{tabularx} \\ \\ \relax
26:9 & \begin{tabularx}{0.7\textwidth}{X} 大衛對亞比篩說：「不可殺害他！有誰伸手害耶和華的受膏者而無罪呢？」 \end{tabularx} \\ \\ \relax
26:10 & \begin{tabularx}{0.7\textwidth}{X} 大衛又說：「我指著永生的耶和華起誓，他或被耶和華擊殺，或死期到了，或出戰陣亡， \end{tabularx} \\ \\ \relax
26:11 & \begin{tabularx}{0.7\textwidth}{X} 耶和華絕不允許我伸手害耶和華的受膏者。現在你可以把他頭旁的槍和水壺拿來，我們就走。」 \end{tabularx} \\ \\ \relax
26:12 & \begin{tabularx}{0.7\textwidth}{X} 大衛從掃羅的頭旁拿了槍和水壺，他們就走了。沒有人看見，沒有人知道，也沒有人醒過來。他們都睡著了，因為耶和華使他們沉睡了。 \end{tabularx} \\ \\ \relax
26:13 & \begin{tabularx}{0.7\textwidth}{X} 大衛過到另一邊去，遠遠地站在山頂上，與他們相離很遠。 \end{tabularx} \\ \\ \relax
26:14 & \begin{tabularx}{0.7\textwidth}{X} 大衛呼叫百姓和尼珥的兒子押尼珥說：「押尼珥，你為何不回答呢？」押尼珥回答說：「你是誰？竟敢呼喚王呢？」 \end{tabularx} \\ \\ \relax
26:15 & \begin{tabularx}{0.7\textwidth}{X} 大衛對押尼珥說：「你不是個大丈夫嗎？以色列中誰能比你呢？百姓中有一個人進來要害死你主你王，你為何沒有保護你主你王呢？ \end{tabularx} \\ \\ \relax
26:16 & \begin{tabularx}{0.7\textwidth}{X} 你做的這件事不好！我指著永生的耶和華起誓，你們都是該死的，因為你們沒有保護你們的主，就是耶和華的受膏者。現在你看，王頭旁的槍和水壺在哪裡？」 \end{tabularx} \\ \\ \relax
26:17 & \begin{tabularx}{0.7\textwidth}{X} 掃羅認出大衛的聲音，就說：「我兒大衛，這是你的聲音嗎？」大衛說：「我主我王啊，是我的聲音。」 \end{tabularx} \\ \\ \relax
26:18 & \begin{tabularx}{0.7\textwidth}{X} 又說：「我主為何要追趕僕人呢？我做了甚麼？我手做了甚麼惡事呢？ \end{tabularx} \\ \\ \relax
26:19 & \begin{tabularx}{0.7\textwidth}{X} 現在求我主我王聽僕人的話：若是耶和華激發你來攻擊我，願耶和華悅納供物；若是出於人，願他們在耶和華面前受詛咒，因為他們今日趕逐我，不讓我在耶和華的產業中有分，說：『你去事奉別神吧！』 \end{tabularx} \\ \\ \relax
26:20 & \begin{tabularx}{0.7\textwidth}{X} 現在不要使我的血流在遠離耶和華面的地上。因為以色列王出來，只不過是尋找一隻跳蚤，如同人在山上獵取一隻鷓鴣。」 \end{tabularx} \\ \\ \relax
26:21 & \begin{tabularx}{0.7\textwidth}{X} 掃羅說：「我有罪了！我兒大衛，回來吧！我必不再加害於你，因為你今日看我的性命為寶貴。看哪，我是個糊塗人，大大錯了。」 \end{tabularx} \\ \\ \relax
26:22 & \begin{tabularx}{0.7\textwidth}{X} 大衛回答說：「看哪，這是王的槍，可以吩咐一個僕人過來拿去。 \end{tabularx} \\ \\ \relax
26:23 & \begin{tabularx}{0.7\textwidth}{X} 今日耶和華將王交在我手裡，我卻不肯伸手害耶和華的受膏者。耶和華必照各人的公義誠實報應他。 \end{tabularx} \\ \\ \relax
26:24 & \begin{tabularx}{0.7\textwidth}{X} 看哪，我今日看重你的性命，願耶和華也照樣看重我的性命，並且拯救我脫離一切患難。」 \end{tabularx} \\ \\ \relax
26:25 & \begin{tabularx}{0.7\textwidth}{X} 掃羅對大衛說：「我兒大衛，願你得福！你必做大事，也必得勝。」於是大衛上路，掃羅也回自己的地方去了。 \end{tabularx} \\ \\
[1ex]
\hline
\hline
\end{longtable}
$^{1}$我們今天要看的題目是兩種人生兩種死亡.
上一次跟大家分享的題目是三個地名三種提醒.
還記得嗎?.
所以今天我們要講兩種人生兩種死亡.
為什麼要講這個呢?因為在三國義記上我們已經來到結束了.
在三國義記上26章一直到31章最後一章.
記載了蘇羅他人生的結束.
所以蘇羅就死了 31章記載關乎他死亡的事.
所以這裡可以說是回顧一下蘇羅的人生.
蘇羅的人生也有最後帶出來他的死亡.
接著我們要看另一種人生.
另一種死亡 我們先買個關子.
所以第一我們先來看一看.
在三國義記上最後.
看到蘇羅他人因為打非利士人的時候受了重傷.
結果他要自盡而死.
結束了一個王朝的時代.
所以我們看到這個可以說是一個朝代的結束.
也是他的死亡帶來他人生的結束.
所以第一我們來看看蘇羅之死.
第一種人生第一種死亡.
蘇羅之死我們發覺.
在這裡看到的就是他糊塗一生而結束.
糊塗一生而結束 為什麼呢?.
這個在三國義記上26章21節.
我們一起讀一讀這節經文好嗎?請大家同聲讀出.
蘇羅說:「我有罪了,我疑大衛你可以回來,因你今日看我的性命為寶貴,我必不再加害於你.我是糊塗人,大大錯了.」.
這裡我們看到是蘇羅追殺大衛的時候.
大衛放他一命.
不止一次 兩次之多.
所以他就覺得自己一次再一次的來做錯事.
或者是對大衛有過錯.
所以他就跟大衛說這段話.
他說:「我有罪了,我疑大衛你可以回來,因我今日看你的性命為寶貴.」.
這句話是真還是假的?.
那麼大衛有沒有回去?.
當然沒有了.
因為這個蘇羅是而又非,出爾反爾的.
絕對是不值得信任的.
很有可能是騙他回去的時候殺了他.

$^{41}$所以當然沒有了.
不過呢.
他自己就套了一些真言.
這個真言就是說.
「我是糊塗人,大大錯了.」.
所以在這裡我們看到.
這句話無論他是真心假意也好.
不過真的是他的描寫.
套了一句真言出來.
他說他是一個糊塗人.
回顧大衛和蘇羅之間的這段記載.
看到蘇羅真的是一個糊塗人.
他糊塗一生而結束.
所以這個結束也是很悲慘的事.
希望我們的人生不會像蘇羅那樣.
到最後的時候說.
我真的是糊塗地過了一生.
實在是糊塗人.
人人都會有一死.
人只有一生的機會.
一生完結就會結束了.
就沒有了.
但有些人死的時候.
是叫做不枉此生.
就是說他很有價值很有內容.
覺得他過了一生的人生.
無論長短是不枉此生的.
但亦有些人是糊塗一生.
糊塗一生就是空白一片.
乏善可陳.
即是沒有什麼值得紀念.
兩種不同的人生.
亦有兩種不同的死的方法.
究竟我們要不枉此生.
還是要糊塗一生.
你要怎樣?.
你要糊塗一生過了.
還是覺得回過一生的時候.
真的不枉此生.
如果這樣的話.

$^{81}$我們人生就會很有價值.
蘇羅良心一時發現.
在他年老的時候檢討自己的人生.
是糊塗的.
是犯錯的.
而且是大錯.
所以有一位學者叫Sidlo Bextor.
他分析蘇羅人生的糊塗.
有很多方面.
綜合跟大家分享.
可以有以下幾點.
第一.
他是.
不聽勸.
糊塗在乎什麼?糊塗在乎他忽視別人的勸告.
其實有一位.
敬虔的恩師.
可以說是他的朋友.
來忠誠的勸告他.
薩姆爾勸告他一次兩次.
亦都給他機會提醒他.
一個人將良師亦有摒棄.
固執.
不肯聽勸.
這個人就是糊塗了.
真是糊塗.
所以我們看到蘇羅糊塗在乎什麼呢?.
其實他有機會回轉的.
也有機會改變的.
不過很可惜他不肯聽人勸.
這個就是一個很大的.
問題.
今天我們都會有時候.
會有這些毛病.
在上一次我們團體聚會的時候發覺.
越老就越什麼?.
越老就越不肯聽人說.
這個都是很糟糕的事.
蘇羅來到他年紀.
又差不多了,不過我們發覺他仍然.

$^{121}$不肯聽人的勸告.
箴言十二章這樣說.
我們一起讀一讀.
十五節,箴言章十二章的十五節,請讀.
愚妄人所行的.
在自己眼中看為正直.
唯知威人肯聽人的勸告.
所以這裡將.
兩個不同的人比較,一個叫做愚妄人.
一個叫做智慧人.
智慧人就肯聽人的勸告.
但是愚妄人就以為自己一切所做的都對的.
所以意思跟後面對比很明顯就是說.
他是不想聽別人說.
所以愚妄人所行的.
自己以為自己對,但是智慧人就聽.
因為他知道有可能不對.
所以他願意聽別人的勸告.
這個就是一個很不同的心態,影響我們.
願不願意聽勸.
箴言十三章第一節,亦都這樣說,我們一起讀,請讀.
智慧子聽父親的教訓.
竊萬人不聽責備.
所以這裡當然是指到家庭裡面.
但其實在希伯來的背景.
他們都是以一些親屬,甚至一些老師為他們的父.
所以這裡就指到他不聽這些.
比較長輩一點,有些經驗一點的人的教訓.
但是呢.
竊萬人,我們看到就是這種的表徵.
所以聖經裡面提醒我們不要做竊萬人,要做智慧的子.
智慧子就聽父親的勸導.
聽長輩,聽別人,聽別人的教訓.
耶利米書十七章,這裡亦有另一節經典,我們一起讀,第23節.
是耶利米先知的勸告,請讀.
他們卻不聽從,不執以而聽.
竟硬著頸看不聽,不受教訓.
這個真的是先知的警告.
所以我們看到無論是真言提醒我們要成為智慧人.
不要做一些愚妄人,或者不要做一些竊萬人.

$^{161}$我們看到先知同友勸告,因為這種是忠誠的勸告.
真的是出於心來勸告他們不要不聽.
但是他們卻沒有執以而聽.
相反硬著頸看不聽,不受教訓.
所以勸歸勸,聽就不聽.
所以不要真的勸,很多人勸你.
你有你說,我行我素.
實在是受虧損自己.
不能夠得到這套益處.
如果能夠虛心接受教訓.
接受教導,往往是得益的開始.
所以我們求神幫助他們.
不要做糊塗人,為什麼呢?.
因為我們看到素羅第一,就是一種不聽勸的人.
第二,我們看到他的問題也是糊塗的,就是他藐視神.
因為他越過了神的旨意.
神明明吩咐他要滅絕那些阿瑪利人.
但是他沒有照辦,他沒有聽神的話.
明明是說你要等撒母耳來才獻祭.
但是他越過界,又不願意聽從.
所以有一個我們很熟悉的千古名言.
就是神透過撒母耳提醒素羅.
聽命勝於獻祭.
順從勝於供養的自由.
他重複一次,供養的自由代表獻祭的意思.
所以什麼是神所喜悅的?就是聽命和順從.
勝過了外表的一些宗教行動.
甚至可能是祈禱,或者獻祭,或者讚美,或者奉獻.
因為這些不是不重要.
不過如果我們一方面祈禱,一方面讚美,一方面奉獻,一方面事奉.
不過是不聽神的命的話.
就失去了最終的意義.
真正神要的是透過這些來表達出我們對他的敬拜.
對他的信命,對他的服侍,對他的順從.
這就是神所喜悅的.
不過很可惜我們見到素羅一而再,再而三.
沒有將神當受敬畏的心放在他的心中.
所以在《箴言十三章》提醒我們.
第七節我們一起讀,好嗎?請讀.
「敬畏耶和華是知識的開端,如往人藐視智慧和憤悔.」.

$^{201}$所以這裡給我們看見如往人另一種情況.
就是他沒有敬畏耶和華.
他藐視神的智慧和神的憤悔.
但是原來真正的知識,真正生命的指引.
就是從敬畏耶和華開始.
這是我們很熟悉的「敬畏耶和華是智慧的開端」.
第十三章再提醒一次「敬畏耶和華是知識的開端」.
《箴言二十八章》亦是智慧書.
我們再讀一次《箴言》的書.
第十二十八章十四節,請讀.
「常存敬畏的,便為有福.
心存光明的,必陷在禍患裡.」.
所以這裡提醒我們要敬畏神,不要藐視神.
上星期我們講了三個地名,三種提醒.
其中最後我們要記得的是「彼列斯烏薩」.
還記得這個地名嗎?.
一定記得,是關乎烏薩的事.
「彼列斯烏薩」提醒我們要尊重神.
尊重神,敬畏神就是蒙福的源頭.
我們習以為常,失了敬畏,尊重神的心.
一定受虧損.
第三,在數來我們見到他的糊塗在於他遮掩自己的罪.
他做錯只去遮掩,不肯承認.
不僅不去承認,還用很多理由去遮掩他的錯失.
所以他跟薩姆爾說,因為你遲了來,所以我來不及了.
因為我知道上好的羊是神所喜歡的.
所以我留下來,沒有照上帝的命令,全部殺了.
他用很多理由去遮掩他的罪.
其實這正是他糊塗的地方.
在《箴言二十八章》,又是《箴言二十八章》,十三節.
我們再讀二十八章,這次是十三節,請大家讀.
「遮掩自己罪過的,必不亨通.
承認利器罪過的,必無連述」.
如果能夠早些承認自己的錯,去改過.
神一定給他機會,很可惜.
他一錯再錯,而且一遮掩再遮掩.
一不承認,繼續不承認.
所以我們見到這實在是很可惜的事.
他糊塗在這裡,因為人人都有錯.
不過為何不去認錯呢?為何不去改變去認錯呢?.

$^{241}$有一句西方諺語是「肯承認自己有缺點的人,他的缺點已經少了一個」.
聽不聽得懂?.
「肯承認自己有缺點的人,他的缺點已經少了一個」.
即是當我們不肯承認自己有缺點的時候.
其實這就是我們一個很大的缺點.
人人都有錯的,如枉人就堅持錯下去.
蘇羅和大衛不同的地方,其實我們嚴格講.
看散佈意見上,蘇羅犯的罪沒有大衛的罪那麼重.
後來,即是大衛.
我們知道大衛犯錯,這種錯是謀財害命那種,很嚴重.
蘇羅,即是獻祭遲些,沒有殺光那些羊.
但為何神不越納蘇羅,廢棄他.
揀選大衛,揀選他成為以色列王.
後來神以他為量度所有以色列王和猶大王的標準.
因為大衛肯悔改,大衛肯承認他的錯.
大衛肯回歸到神的面前,這就是分別.
所以我們看到大衛得著智慧.
因為他能夠回到神的面前承認自己的錯.
蘇羅卻作糊塗人到底,因為他遮掩自己的罪.
彼得和猶大兩個都是有犯過錯的門徒,耶穌的門徒.
但彼得三次不認主,結果怎樣.
結果他悔改,痛哭,結果重新起來.
後來成為神大大使用的仕途.
但是猶大卻沒有回到神的面前.
他死硬到底,情願至盡.
都沒有去到耶穌的面前,或者去到神的面前悔改.
所以我們看到這就是很大的分別.
在約翰一書第一章告訴我們.
這是和教會說的,不是和未信的人說的.
是和懂得認識神的人說的.
我們看看聖經上的經文英序.
約翰一書一章九節,我們一起讀,請讀.
「我們若認自己的罪,神是信實的,是公義的,.
必赦免我們的罪,洗淨我們一切的不義.」.
這就是上帝的應許,祂要洗淨我們的不義,赦免我們的罪.
但是我們需要什麼?祂說「我們若認自己的罪」.
所以這是很重要的一個關鍵.
求神幫助我們不要學蘇羅的糊塗.
因為他不斷地來遮掩罪.
第四我們看到他懷忌恨,這就是他另一個糊塗.

$^{281}$我們知道他一而再,再而三.
主要要殺大衛的原因就是他心裡的忌恨.
他妒忌大衛,要把王位給他.
人們擁戴他,所以他心裡就產生這種忌恨.
這種嫉妒和仇恨就控制了他.
控制到一個程度他被魔鬼入侵.
由惡魔來到他的心中影響他.
所以有時失去理智,情緒有時失控.
更加當他恢復清醒,有時會認錯的時候.
又重新再一次地用這個忌恨來控制他自己.
所以這實在是一個很大的糊塗.
在現代的醫學裡面都有很清楚的證據.
如果人經常被這些腐毒或者是老奴.
傳在我們心裡,不去解決,不去處理的時候.
他有時會影響我們的情緒,影響我們的健康.
影響了我們判斷事情的準確,甚至是身體的疾病.
所以很多的醫生都發現,其實他沒有什麼身體大的毛病.
不過總是頭痛,又悲痛,又過敏,又有潰瘍.
有時血壓高,甚至有心臟問題.
當然都有生理方面的因素.
但有些不是生理因素,而是他心理的因素影響他.
就是他心裡充滿著很多這些腐毒,很多這些疾病.
令到他心靈影響他的身體.
真是一個今天很多現代人所面對的問題.
在聖經裡面,雅各書提醒我們.
我們一起讀讀這幾節,好嗎?.
雅各書第三章,這是新約的智慧書.
三章十四到十八節,請讀.
「你們心裡若懷著腐毒的疾度和紛爭,就不可自誇.
也不可說謊話抵擋真毒.
這樣的智慧不是從上頭來的,乃是屬地的,屬情慾的,屬毀魔的.
在何處有疾度紛爭,就在何處有擾亂和各樣的壞事」.
十七節.
「唯獨從上頭來的智慧先是清潔,後是和平.
溫柔柔順,滿有憐憫,多見善果,滿有偏見,滿有假冒.
並且使人和平的是用和平所栽種的異果」.
這兩個將屬地的智慧和屬上頭來的智慧作一個比較.
世人的智慧就是爭奪,有很多紛爭,雜惰,擾亂的事就不斷增加.
上頭來的智慧就是和平的,清潔的.
溫柔柔順,滿有憐憫,多見善果,滿有偏見,滿有假冒.

$^{321}$而且能夠使人和平,心裡面能夠有平安.
這種就是在聖靈所結的平安的果子,用和平所栽種的異果.
人生到尾如果回顧一下,仍然充滿著很多苦毒,很多嫉妒,很多埋怨.
沒有好好處理和釋放.
沒有在神的面前求上帝來更新,潔淨我們.
將這些好像神饒恕我們一樣的放開.
這個人生一定不會快樂.
這個人生一定是仍然帶著這些來結束.
沒有靠著神的恩典能夠去釋放這些的時候.
是一種糊塗的人生.
所以你看素羅大部分的時間都是用來追大位.
為什麼不用這些時間去做一些有益的事呢?.
有很多人一生就是帶著這些來結束人生.
為什麼不用精力去做一些建設社會造就別人的好事呢?.
不如放開這些,讓我們能夠有貢獻的人生.
最後我們看到素羅的糊塗最後結鬼魔.
他最後死的原因是因為他沒有照神的吩咐.
他也知道頒佈給以色列人要守的命令.
他自己卻去交鬼.
所以如果看回這段記載,你會發覺如果他沒有去巫婆那裡交鬼.
沒有自己嚇自己.
沒有去給這些通靈的人嚇他的時候.
他不會死得這麼快.
他就是因為聽完之後.
聖經形容他猛然瀑倒.
挺身在地,非常軟弱,極其恐懼.
所以這個就令到他完全失了方.
打仗的時候一定輸.
所以我們看到這是他最後的下場.
在這裡他心理和身體差不多全年崩潰.
未打就先輸了.
所以聖經裡面提醒我們.
以別神代替耶和華的,他們的受苦必加增.
求神幫助我們,不要中了魔鬼的詭計.
千萬要小心.
我最近聽到一位傳道人的家人說.
他說爸爸從小就這樣提醒我們,今天神有什麼要跟你說.
然後他說今天魔鬼有什麼要跟你說.
他說如果你聽到魔鬼跟你說的話.
你記得馬上趕走他,一腳踢走他.

$^{361}$神跟你說的話,你記得要聽神的話.
有沒有啊,有沒有有時候魔鬼跟我們說話.
有時候上帝也跟我們說話,究竟我們聽誰的.
更加千萬不要有邪靈交鬼的行動.
如果過往未信主之前有這些.
我們切切實實將它斷絕.
在《新命紀》第十八章,這是素羅也知道的.
他說你到了耶和華,你所去的地方.
那些國民所行可憎惡的事,你不可學行.
你們中間不可人洗兒女經火.
也不可有占卜的,觀兆的,用法術的,行邪術的,用迷術的,交鬼的,行邪術.
過陰的,凡行這些事,都為耶和華所憎惡.
所以我們看到,總括起來,素羅的糊塗是非常自我.
他不聽神的命令,也不肯悔改.
不聽人的勸告,糊塗一生,這是第一方面.
第二,我們看到他的死是一大諷刺的死亡.
這就是《新命紀》上最後的結束,形容他怎樣死.
我們很快讀一讀,我們一起讀這幾節.
「非利士人與以色列人爭戰,以色列人在非利士人面前逃跑.
在基地坡有被殺戮到的非利士人.
緊追素羅和他兒子們,就殺了素羅的兒子.
約拿達·阿比拿達·麥基舒瓦.
西派甚大,素羅被弓箭手追上,射傷甚重.
就吩咐拿他兵器的人說,你不出刀來將我刺死.
免得那些未受過禮的人來刺我,凌辱我.
但拿兵器的人甚懼怕,不肯刺他.
素羅就自己伏在刀上死了」.
所以他最後是怎樣死的?.
第八節請讀:「此日,非利士人來沒那未殺之人的衣服.
看見素羅和他三個兒子佈倒在基地坡山.
就割下他的首級,剝了他的軍裝.
打發人到非利士地的四景.
報信予他們廟裡的偶像和眾物.
又將素羅的軍裝放在阿斯塔祿廟裡.
將他的屍身釘在白山的城牆上」.
這樣死,真是很可惜.
當然我們要知道,上帝是公義的.
真是為大衛親冤.
所以素羅真是如此下場.
他的死亡是一大諷刺.

$^{401}$為什麼呢?因為他一生最怕的是沒面子.
但結果死的時候,衣服也沒有,頭也沒有.
被人高掛於木架上,被人嘲笑.
真是很可惜.
他的戰衣就放在偶像的廟裡.
就是非利士人的廟.
屍身就被釘在白山的牆上.
這種是不是一種諷刺呢?.
有人研究白山城其實是非常近.
素羅被薩摩爾高臨,就職為王的地方.
可以說是繞了一圈,一生就回到原點.
So back to square one.
可以說是得個零.
所以這真是很可惜的一個人生.
真是叫我們警惕.
這種素羅自我中心,死也要面子.
結果最後成為一大諷刺,被人笑柄.
我曾經去過一個靈命更生講座.
講完之後他說.
很多人都是這樣,寧願要面子不要神的恩典.
因為他要面子,死也不肯認.
所以不斷遮掩,不斷不願意接受神的赦免和恩典.
但有些人願意放下自己.
坦然來到神面前,求神去施恩.
得著神的恩典.
所以有句說話跟大家分享.
這叫做Gaining Faith,Losing Grace.
中文有個翻譯.
只顧面子,失去恩典.
不如失去面子,得著恩典.
即是我們來到神的面前,承認自己的有限.
需要你的恩典.
所以不要硬要面子.
因為就像素羅一樣,最後只有個零.
而且一大諷刺的死亡.
今天跟大家說兩種人生,兩種死亡.
究竟第二種是什麼呢?.
這個主要是我讀了一本書.
是由前Dallas神學院院長寫的.
就是Charles Wendel.

$^{441}$他有一個很特別的比較.
他對於素羅的死有個比較.
所以跟大家分享最後做結束.
這個就是基督的死.
基督的死和素羅有什麼分別呢?.
素羅的死就代表國家的希望完了.
他本來就是以色列國第一個皇帝.
結果就此完畢.
但基督的死就帶來什麼?.
基督的死就帶來人類新的希望.
這個死亡很不同.
素羅的死就代表.
他被逆境邪惡而得勝.
所以被邪惡而勝.
但基督的死就是戰勝了罪惡和死亡.
是直著死亡敗壞了掌死權的魔鬼.
所以這是基督的死.
素羅就是引進了大衛王朝的開始.
大衛王朝因此開始.
而基督的死就是引進了大衛之約.
所應許的永恆的國度的開始.
所以這是基督的死帶進一個新的國度.
直著大衛之約所延伸出來的.
所以基督就開始了一個永恆新的國度.
素羅的死就是展示了一個愚昧的人生.
正如我們剛才所說.
基督的死就是展示了神的大愛和大能.
讓我們能夠活出一個豐盛的人生.
所以耶穌基督的死亡.
他換來千萬人得著豐盛的生命.
包括你和我.
素羅的死是被釘在白山城.
罪有應得.
但基督的死是被釘在葛國他山上.
而異代替了不異.
讓我們能夠承受神的恩典.
所以我們看到素羅最後的身體被分解和毀壞.
基督就是死人復活.
而他復活的身體顯出他是神的兒子.
得著榮耀.

$^{481}$所以這兩個比較很特別.
我和大家分享一下Charles Wendel的比較.
所以在這裡我做一個結語.
我們看到素羅糊塗的一生.
糊塗的一生是什麼?.
他不聽勸.
所以我們應該反過來學習.
應當謙虛受教.
這就是智慧了.
求神幫助我們.
第二 他是什麼?.
他是藐視神.
所以我們學習如何?.
學習謹守 遵守神的命令.
這是很重要的.
我們要持守.
不要學素羅的糊塗.
因為他遮掩罪.
所以我們應該如何?.
我們應該肯認罪.
肯悔改.
這就是智慧的人生.
素羅的人生是壞忌恨的.
所以我們求神幫助我們.
求神釋放心裡的苦毒.
這就是我們能夠得著智慧.
不要學素羅的糊塗.
最後 素羅是結鬼魔.
所以我們應該如何?.
我們應該尚存忌邪的心.
這就是神所希悅的.
要我們能夠有一個忌邪的心.
要我們能夠守住神給我們的聖潔.
最後 素羅是假面具.
因為只顧面子.
我們應該如何?.
我們應該到主前.
活出真我 得著釋放.
來到神的面前 有一個真的自由.
因為神預納我們 接納我們.

$^{521}$我們能夠活出在主裡面的真正的我們.
這就是我們從素羅.
這個壞的榜樣 糊塗的人生.
我們學習一個智慧的人生.
不單止這樣 我們更進一步.
追求一個豐盛的人生.
我們不要忘記耶穌基督.
所為我們所作的一切.
耶穌為我們所成就的.
所以我們要歸主 歸信救主.
如果還沒有信耶穌.
快點信靠主耶穌.
我們需要來學習基督.
學習耶穌基督的人生.
追求一種仿似基督的生命.
活出在我們裡面神賜給我們新的生命.
我們更加靠主活出豐盛.
結出聖靈所結的果子.
活出豐盛的生命.
然後讓我們兩種人生.
兩種死亡 你要選擇哪一種.
不得我們死而無憾.
更加不枉此生.
讓我們兩種人生.
兩種死亡.
你選擇哪一種.
你要糊塗一生這樣過.
還是死而無憾.
不枉此生.
願意讓我們得到神賜給我們的智慧.
不要學習數落.
應該思想.
學習主耶穌基督.
我們一起同心圖告.
請大家低頭祈禱.
而且悔改.
叫我們能夠靠著神.
釋放我們心裡面一切的苦毒.
一切的罪惡.
又讓我們有個忌邪的心.

$^{561}$來遵行你.
守住你給我們的聖潔.
叫我們能夠坦然的來到你面前.
活出我們真正的惡.
在基督裡面我們的自由.
因著信靠耶穌.
活出豐盛.
好讓我們一生無憾.
告一個榮神益人.
能夠榮耀你.
能夠去建立自己.
也能夠造就別人的人生.
讓我們不枉此生.
因為主耶穌活在我們裡面.
因為信靠主.
是我們一生最大的福氣.
多謝你聽我們祈禱.
奉耶穌基督的名.
阿們.
\newpage



\section{士師記 7:13}
\label{sec:032307}
\textbf{聖經人聖經事:一個夢:一個大麥餅}~~~~黃克勤牧師博士
\newline
\newline
連結: \href{http://www.church.com.hk/acms/content.asp?site=cdc\&op=show\&type=product\&code=032307\&layout=sermon}{\texttt{032307}} ~~~~ 語音日期: 2026-06-14
\newline
\newline
\hyperref[sec:032306]{\small{< < < PREV SERMON < < <}}
~
\hyperref[sec:index_chronic]{\small{[返順時目]}}
~
\hyperref[sec:032308]{\small{> > > NEXT SERMON > > >}}
\newline
\newline
士師記 7:13
\newline
\begin{longtable}{cl}
\hline
\hline
章節 & 經文 (和合本修訂版)\\
\hline
7:13 & \begin{tabularx}{0.7\textwidth}{X} 基甸到了那裡，看哪，有一人把夢告訴同伴說：「看哪，我做了一個夢。看哪，一個大麥餅滾入米甸營中，來到帳幕，把帳幕撞倒，帳幕就翻轉倒塌了。」 \end{tabularx} \\ \\
[1ex]
\hline
\hline
\end{longtable}
$^{1}$各位電影姐妹大家好.
生命算書的網絡.
又和大家讀聖經.
聖經人,聖經事.
看看背後上帝的手有什麼作為.
今天的題目是一個夢,一個大麥餅.
你知道我們一直在說時事機點.
來到時事記.
時事記的事情,故事.
事件實在非常非常多.
而且是很好說的.
原因就是它是一個故事.
每一個是一個民族英雄.
神用他來打垮外國人.
也是用外面的民族來鞭打以色列.
懲罰他們.
然後神又用以色列作為皮鞭.
鞭打這些邪惡的外族人.
聖經人,聖經事.
今次繼續說機點.
一個夢,一個大麥餅.
這個夢和大麥餅就拉上了關係.
我們看看《斷聖經》.
時事記第七章.
晚上.
其實這裡是當那晚.
就是那一個晚上.
那晚是什麼呢?.
如果看回昨天.
你看到機點找了三百個人.
本來三萬二千人一起來.
有兩萬二千人回去了.
剩下一萬人.
神就說仍然太多.
你帶他們去水邊喝水.
大家記得昨天這件事.
喝了水之後.
神就說那些趴下跪在地上喝水的.
就叫他們回去.
是那些.

$^{41}$很有趣.
他說用手捧著水來舔.
當然我們不會像狗舔.
這樣喝的人.
這樣吸進去的人.
那些就是了.
有三百個這樣的人.
其實為什麼呢?.
有些人就解釋.
跪在地上.
一個屬靈的意義.
屬靈的象徵性.
象徵性,靈修性的解釋就是.
你是會容易跪下的.
看到有水喝就跪下.
然後甚至趴下去喝.
完全沒有志氣的.
這樣不行.
如果其他人是怎樣呢?.
去到河邊.
水離開這麼遠.
那當然是跪下的.
就是說這些人.
兩種人都是用手舔的.
不過有些跪下.
有些自然沒有跪.
沒有跪是怎樣呢?.
可能他彎下身.
拿了河水之後.
就這樣喝.
這樣喝的時候.
他站著看到周圍的東西.
也反應會比較快一點.
很多人是這樣.
從這個角度去看的.
但是其實我覺得.
不可以.
可以接受.
不過這個有少少是用.
現在我們的想法.

$^{81}$套進去《亂聖經》裡面.
如果直接地說.
不知道為什麼.
因為那裡喝水.
都可以周圍看的.
這樣喝水.
跪下都可以周圍看的.
那就不一定說.
站著的人警覺一點.
反應容易快一點.
可能是可能不是.
不知道.
不過.
神的揀選就是那些.
應該不是跪下.
而用手拿著水這樣喝的人.
只有300個.
如果我是機電.
即是暈倒.
300個.
本來有3萬人.
我的心一定很多.
本來我都很怕的了.
現在剩下1萬.
我覺得都勉強可以.
有1萬人.
那你有3萬多人.
看個陣式擺出來都嚇一驚.
現在只有1萬.
都算了.
都不行剩下多少.
300.
在這些橋段拍戲才有.
300.
以前有一套舊戲叫《戰狼300》.
300個斯巴達人.
怎樣去對抗敵軍數萬人.
那些希臘時代的勇士.
《戰狼300》.
好了.

$^{121}$結果其實那套電影.
那300個都死了.
如果我是機電.
300個真的很怕.
米電都無數這麼多人.
不過就是這樣選完之後.
當那夜.
神都知道機電心裡面.
步步跳.
唉呀很怕.
於是你說那300人怕不怕.
都怕的.
本來就說膽怯的可以回家.
走了2萬多.
剩下1萬.
其實我不怕的.
不過當你剩下300之後.
那300人不怕就搞了.
好.
耶和華就吩咐機電.
當那夜.
起來下到米電營裡去.
即是說.
起身.
你進去米電營地窺探下.
做下這個窺探.
看看怎樣.
唉.
其實我想你知道多些.
其實我已經將他們交在你手裡.
你去看看.
這些人全部都將要被你打敗的了.
不過.
你要知道我已經將他交在你手裡.
下去看看.
深入滴陣一個人去窺探很危險的.
神一早就知道機電是這樣的.
好像神是知道我們內心怎樣想的.
神就說.
倘若你怕下去.

$^{161}$如果你怕.
帶上你的僕人普拉夏道那營去.
和你的僕人叫普拉.
我就不特別說他的名字什麼意思了.
你和普拉一起去.
那麼.
真是上帝厲害.
這樣都知道他害怕的.
其實我也知道他害怕.
這樣不害怕是假的.
一個人進去千萬的米電軍隊的營中.
被人捉到怎麼辦.
有沒有少少化妝.
穿起他們的衣服.
可能你都沒有準備.
於是.
怕就找僕人普拉去.
壯壯膽.
你必聽見他們所說的話.
你聽到他們.
留意聽敵軍那些士兵說的話.
聽了之後.
然後.
第一點了.
你就有膽量下去攻營的了.
於是機電就帶著僕人普拉.
一起去營房了.
你去聽一下對方的士兵說什麼.
然後你就有膽量下去攻營的了.
你去聽一下.
不要只在那裡猜.
其實很多時候我們都會被自己.
最大的敵人.
就是自己的恐懼所嚇到的.
一看.
哇.
無千無萬這麼多敵軍.
比王重環多.
駱駝無數.
比海邊的沙一樣.

$^{201}$這樣你真的害怕的.
不過其實神是不斷跟他說.
我和你一起.
我和你同在.
你鎮定一點.
放心.
可以的了.
然後你下去看看.
你害怕就看清楚一點.
哇.
這個時候.
真的不害怕的.
你看完之後.
你就會有膽量的了.
於是.
機電就好了.
我就下去了.
其實機電都學到信服了.
和有信心了很多.
第十二次.
米殿人.
阿瑪力人.
阿瑪力人是都出名.
戰士.
戰爭很厲害的.
米殿人.
米殿人就是.
競爭的意思.
競爭.
是一個.
在約旦河的東岸.
稍為南一點的.
沙漠曠野地帶.
很闊的.
很多個細小步驟.
不過他們是一起的.
這樣就叫統稱米殿人.
米殿人.
阿瑪力人.
就是在地中海南邊那邊的人.

$^{241}$一切東方人.
東方人是指約旦河的東邊.
有很多民族.
那些更遠的.
就是阿滿.
摩亞.
以東這些.
散步在平原那裡.
哇.
如同蝗蟲那樣多.
駱駝無數.
多如海邊的沙.
為什麼用駱駝.
不是馬的呢.
因為他們全部都在沙漠游牧的.
因此.
他們很善於.
沙漠的行動.
鬥爭.
而且他們騎駱駝的.
駱駝怎麼夠快.
駱駝不快.
駱駝如果跑起來.
很快的.
特別在沙裡面.
跑得很快.
馬是不能比的.
馬在沙裡面是很難跑的.
駱駝就可以了.
駱駝趕起路上來很快的.
第十三節.
幾年到了營地裡面.
當然是躲起來了.
聽見有一個士兵.
發了一個夢.
將夢告訴同伴.
第二點了.
什麼夢呢.
哎呀我做了一個夢.
夢見一個大麥餅.

$^{281}$滾入米電營中.
一直滾到帳篷.
哎呀帳篷被他撞爛了.
帳幕就是翻轉.
平復了.
留意再說一次.
那個士兵就說.
我發了一個夢.
夢見一個大麥餅.
滾入米電營中.
一直滾到帳篷.
帳幕那裡.
就是我們住的地方.
將帳幕撞倒.
帳幕就翻轉.
平復了.
你知道帳幕是這樣拉開的.
一撞帳幕就倒塌了.
布就翻開了.
大麥餅是什麼呢.
大麥餅就是一般當地人吃的.
用大麥粉.
有大麥有小麥.
小麥貴一點.
就是我們今天吃的麵粉.
大麥粉是粗一點的麥子.
是大粒一點.
便宜一點.
粗一點.
做出來的纖維是更加粗一點.
吃下去比較耐吃.
因為纖維多.
有部分不是很直接消化到.
吸收到.
所以呢.
它是有這個腸臟.
固然這麼說.
但是它好在是耐吃.
有時窮的人.
做勞力的人.

$^{321}$要吃大麥餅.
五餅魚裡面有兩條小魚.
五個麥餅.
五個麥餅就是五個大麥餅.
加上兩條小魚.
即是少少的蛋白質.
五個麥餅用來充飢.
可能是兩餐.
小朋友拿著兩餐.
一條小魚一餐.
加上兩三個大麥餅.
大麥餅我想都是差不多大.
如果今天去吃.
在當地沙灘吃的都是差不多大.
大麥餅其實是粉搓開之後.
將它放進一個熱的.
像鑊一樣的上面.
中間是火.
鑊圓拱形.
才放進去.
它燒焦了.
不是燒焦了.
燒熟了一邊.
它就反轉.
再反轉.
就熟得均勻,好吃一點.
軟身的.
這麥餅.
放久了,硬了.
就沒那麼好吃.
但如果帶出野外.
都可以放很多天.
而且地方乾燥.
所以大麥餅其實不是一個.
一個很大的木的輪.
車輪.
不是.
它是一塊麥餅.
這個大麥餅能夠將橡木撞lum .
我們以為麥餅那麼大.

$^{361}$好像車輪那樣滾過來.
其實.
你會看到一個印象.
就是一個很大的,軟的餅.
不過這麼大,軟的餅.
滾過來,其實餅是軟軟的.
怎樣能夠直滾都說不了.
總之它滾過來,將橡木撞lum .
所有的橡木撞lum .
一走過之後,全都撞lum .
傾覆,即是說全部倒地.
發出一個這樣的夢.
那麼有趣.
所以其實.
也有人問,為什麼用大麥餅.
不用大木車輪呢?.
不用一塊大石輪呢?.
要用麥餅呢?.
其實可能有個意義.
就是上帝讓他知道.
讓這個基甸聽到這個夢的時候.
麥餅是軟的.
而且麥要來攻擊人.
麥餅怎樣攻擊到人呢?.
麥餅攻擊人,攻擊不了的.
軟的.
而且不是一個貴重的東西.
大麥餅.
但是上帝用一個這樣的.
大麥餅,看來不能夠做到武器.
攻擊不了人的大麥餅.
將米殿的營地沖反了.
當他的同伴聽到.
這個同伴不一定是一個人.
同伴聽到就說.
哎呀,這樣啊.
這個不是別的東西了.
一定是這個了.
是以色列人約阿斯的兒子.
基甸的刀.

$^{401}$這個基甸的刀的意思.
就是說是基甸的武力.
是基甸的刀.
就是說基甸的軍隊.
基甸的軍隊與武力.
神已將米殿和全軍都交在他手中.
哎呀,我們米殿已經被他們的上帝戰勝了.
我們將它好像傾覆的帳篷.
而基甸就變成一個又不值錢.
又軟的麥餅.
大麥餅,即是大麥做的餅.
不是很大個.
那麥餅有多大他沒有說.
如果大麥餅可能是差不多大.
這樣的麥餅那麼小就已經搞定了.
那些帳篷都可以的.
又或者變成很大,所以滾倒了都可以.
那這個餅究竟是直滾過去還是橫向飛.
好像飛碟那樣,沒有人知道.
因為他沒有說.
一個大麥餅就進去米殿帳篷就全部推倒.
奇了,這樣都可以?.
這個就是基甸的軍隊和他的刀.
基甸的軍隊可以做這樣的事.
大家都很害怕.
大家看到一個夢一個大麥餅.
你看到整件事有沒有感覺.
當很害怕的時候.
神就說不用害怕.
你下去看看吧.
怎樣看都好.
我現在給你選剩三百人.
你還叫我下去看萬一捉了我.
真的不知道怎麼辦.
不怕,我會保護你的.
找你的同輩一起去吧.
下去後閃來閃去.
在帳篷旁邊蹲下的時候.
聽到有士兵聊天.
神叫你下去聽他們說話.

$^{441}$你就知道我已經將他們交在你手裡.
聽聽吧.
一個大麥餅飛過來.
將所有帳篷全部傾覆推倒.
這就是基甸的軍隊了.
看來軍隊好像軟軟的.
不太值錢.
不太厲害.
裝備又沒有很特別.
不過他將我們得勝了.
我們全輸了.
發過這樣的夢.
最奇妙的地方是有人發夢.
然後有人解夢.
其實最有趣的地方在這裡.
就是你發這樣的蠢夢.
傻的.
怎會有這樣一個麥餅推倒我們.
你為什麼這麼害怕.
我們這麼多人.
千萬樂陀.
千萬軍隊.
怕他什麼.
只有麥餅一口就吃了.
怎會推倒我們.
傻的.
再發一個好一點的.
但是有人解他的夢.
解了之後大家都害怕.
如果你再看.
很有可能這個夢.
很快就全遍整個米殿的軍隊.
大家都害怕.
基甸的軍隊快要來了.
他們最怕什麼呢.
他們對神的信心比基甸還要大.
神已經將我們的軍隊.
交在基甸的手裡了.
所以這個是基甸的刀.
反而敵人對神的印象尊敬還大一點.

$^{481}$基甸親歷其境者.
反而.
如果我是基甸.
三百人.
神說你別玩了.
不過你走到這一步.
死就死了.
不過神說怎會死呢.
我已經將他們交在你手裡了.
你聽聽他們說什麼.
我已經預備好了.
所以神用一個夢和解夢最厲害了.
就是關於一個大餅的夢.
加強了基甸的信心.
這個是神體恤的.
也是神有這樣的能力做的.
祈禱.
多謝主.
你真是滿有大能.
你不單止能夠控制每一件事.
你更加能夠控制我們發的夢.
你令到一個基甸的.
細小的士兵發了這樣的夢.
和同伴講解.
竟然有同伴解到這個夢.
剛好也是給基甸.
同一個晚上.
最害怕的時候.
只剩下三百人在心裡.
很多疑問在心中的時候.
你就帶基甸去聽到這個夢.
和這個夢在米殿人口中的解釋.
這個真是上帝無比的能力.
讚美你.
願你得到榮耀.
願你被高舉.
戰勝這個米殿人的上帝.
願你得到榮耀.
讚美你.
祈禱奉耶穌基督的名.

$^{521}$阿們.
\newpage



\section{以弗所書 6:1-4}
\label{sec:032308}
\textbf{把孩子帶到主面前}~~~~杜崇基牧師
\newline
\newline
連結: \href{http://www.church.com.hk/acms/content.asp?site=cdc\&op=show\&type=product\&code=032308\&layout=sermon}{\texttt{032308}} ~~~~ 語音日期: 2026-06-14
\newline
\newline
\hyperref[sec:032307]{\small{< < < PREV SERMON < < <}}
~
\hyperref[sec:index_chronic]{\small{[返順時目]}}
~
\hyperref[sec:032309]{\small{> > > NEXT SERMON > > >}}
\newline
\newline
以弗所書 6:1-4
\newline
\begin{longtable}{cl}
\hline
\hline
章節 & 經文 (和合本修訂版)\\
\hline
6:1 & \begin{tabularx}{0.7\textwidth}{X} 作兒女的，你們要在主裡聽從父母，這是理所當然的。 \end{tabularx} \\ \\ \relax
6:2 & \begin{tabularx}{0.7\textwidth}{X} 當孝敬父母，使你得福，在世長壽。這是第一條帶應許的誡命。 \end{tabularx} \\ \\ \relax
6:3 & \begin{tabularx}{0.7\textwidth}{X} 見上節 \end{tabularx} \\ \\ \relax
6:4 & \begin{tabularx}{0.7\textwidth}{X} 作父親的，你們不要激怒兒女，但要照著主的教導和勸戒養育他們。 \end{tabularx} \\ \\
[1ex]
\hline
\hline
\end{longtable}
$^{1}$各位電影姐妹平安.
各位家長平安.
因為我們都知道今天.
是一會兒會有孩子的.
詩歌班.
所以今天有一些.
父母在當中.
都是特別支持他們的孩子.
其實我都想一下.
自己小時候都做過這些.
兒童詩班的.
我們一舉一動.
那些爸爸媽媽.
都是非常興奮的.
我預先講定.
一會兒不要太誇張.
都是含蓄一點.
應該我們都會有人幫他們拍照.
很開心.
一會兒我們都會為這班小朋友.
祝福的祈禱.
今天要和大家.
分享神的說話.
在以弗所書六章一到四節.
就是剛才我們讀的那一段經文.
你們作兒女的.
要在主裡聽從父母.
這是理所當然的.
要孝敬父母.
使你得福在世長壽.
這是第一條大應許的誡命.
作父親的.
你們不要激怒兒女.
但要照著主的教導.
和勸誡養育他們.
在這一段經文裡.
就講到作兒女的本份.
亦都講到作父母的本份.
而當中在主裡.
就是貫穿這兩個本份.

$^{41}$最重要的地方.
我們知道作為父母.
我們都很愛自己的孩子.
由孩子出世開始.
我們已經為他們操心.
吃什麼 穿什麼衣服.
讀什麼學校 學什麼才藝.
將來怎樣發展.
特別在香港.
作父母實在不容易.
很多家長白天工作忙碌.
但到晚上仍要照顧孩子的功課.
情緒 生活習慣.
有時我們很想做一個好的爸爸.
好的媽媽.
但現實是我們亦都會疲倦.
焦慮 發脾氣.
甚至覺得自己做得不夠好.
最近太古城發生了一宗慘劇.
就是母親和女兒先後墮樓.
據說這個母親也是一個輔導員.
有時我們都有個感受.
就是能醫不自醫.
當我們聽到這些消息的時候.
我們覺得很傷心 很難過.
因為我們看到.
在今天這個教育很緊迫.
很有令人窒息的裡面.
甚至令到我們喪失了我們的生命.
但是我們在主的裡面.
我們不自絕望.
我們在主的裡面.
我們能夠常存盼望.
這就是為什麼我們這麼著緊.
要每一位家長都認識主耶穌.
在我們生命裡面的重要.
我們沒有一個人能夠逃避人生的風浪.
但是我們能夠依靠神的時候.
我們能夠經歷和經歷人生風浪的裡面.
神給我們的安慰和同在.

$^{81}$今天神的話不單止要讓我們知道怎樣去做.
祂更加提醒我們.
我們應該怎樣去追求我們生命的意義.
家庭不單單是一個單位.
家庭是神塑造人生命的地方.
我們的孩子不單止需要食物 衣服 教育和活動.
他們更加需要愛 真理 榜樣和信仰的根基.
所以願意今天和大家分享的時候.
我們先將一切交給主.
我們同心禱告.
你的聖經 你寶貴的話語.
就成為我們的指引.
就成為我們的幫助.
以致我們在這個彎曲末路的世代.
在一個充滿困難 焦慮 抑鬱的世代當中.
我們仍然能夠靠主站立得住.
求天父你的話親自與我們眾人同在.
帶領我們明白你的真理.
更加活出你的真理.
我們這樣在你面前祈禱感恩.
奉主名及 阿們.
當我們再看二弗所書六章一節.
作兒女的你們要在主裡聽從父母.
這是理所當然的.
第二節又說.
我們要當孝敬父母使你得福在世長壽.
這是第一條大應許的誡命.
這裡所說到的我們聽從.
這裡說到我們的孝敬.
這個都是聽從 是外在的行動.
孝敬是我們內心的態度.
孩子不單要學習按父母的教導去做.
更加要學習尊重父母 愛父母.
在中國文化裡孝道是一個很重要的教訓.
古人說 大孝尊親 其次不欲 其下能養.
這是出自於禮記祭儀.
是曾子提出孝道的三個層次.
原來最高層次的孝就是尊敬父母.
使父母得到真正的尊榮.
其次不欲 意思就是不要辱罵父母.

$^{121}$羞辱父母 令父母蒙羞.
而最低的層次 其下能養.
就是能夠供養父母 使父母生活有依靠.
這些中國民間的智慧.
和聖經裡孝敬父母的精神是有相通之處.
其下能養提醒我們.
孝順不只是口頭說的.
也包括實際生活的照顧.
當父母年老的時候.
兒女應該關心他們生活的需要.
聖經也很重視實際照顧家人.
提摩太前書五章八節說.
人若不看顧親屬就是背了真道.
信仰不只是在禮拜堂裡說敬虔.
信仰也要在家庭中活出責任.
中國的智慧也指出.
只是供養父母仍然不是最高的孝.
其次不欲 意思就是兒女的品格.
行為 做人 處事不應該羞辱父母.
這一點對孩子的教育很重要.
我們不僅希望孩子成績好.
更希望孩子做人正直.
誠實 有禮 懂得愛人.
如果孩子將來有很多成就卻沒有了品格.
甚至用欺騙 驕傲 自私的方式生活.
這就不是父母真正的喜樂.
我不知道大家有沒有留意.
最近一位實習醫生的事情.
我也不多說.
但是我真的很感受到.
這是對社會 對他個人一個很大的損失.
社會用了很多資源栽培他成為一位醫生.
但是我們卻覺得他不懂得分輕重.
以致他現在喪失了成為醫生的資格.
今天我們看到我們的孩子成績好.
但是如果我們的思想不夠成熟的時候.
我們仍然不能夠對社會有所貢獻.
而最高的孝就是大孝尊親.
尊親不僅是外表稱呼父母.
不僅是節日送禮物.

$^{161}$而是從心裡敬重父母.
體重父母的角色.
珍惜父母的付出.
這與聖經所說的孝敬父母是非常接近的.
聖經所說的孝敬不僅是物質的供養.
也不僅是避免做壞事.
而是內心承認父母是神所賜的權柄.
是神在家庭中設立的引導者.
曾經聽過一個法人深省的故事.
就是一位年長的母親因為她手震.
吃飯的時候經常失手打爛碗.
她的媳婦她的兒子感到心煩.
而且抱怨去斥責這個長者.
結果為了不打爛這個碗.
她的兒子就挑剔了一個木碗給母親專用.
即使掉了在地上也不會破損.
有一天這個兒子看到自己的小朋友在剋木頭.
就問他做什麼.
他的兒子就這樣說.
我正在剋這個木碗.
等你們老了的時候.
我也要用這個木碗栽飯給你們吃.
免得打爛.
這個兒子聽了之後很震驚很羞愧.
他體會到自己今天怎樣對待他的父母.
明天他的孩子就會怎樣對待自己.
因而趕緊改善對待母親的態度.
孝順父母是為人子女的責任.
更是一種新教.
我們對父母的態度.
會直接影響下一代對我們的態度.
我是出生於一個牧師的家庭.
所以我很有深的感受.
雖然我父親是牧師.
但其實牧師也是常人.
做父親的當然是望子成龍.
望子成龍的方法就是他要成績好.
但偏偏我在人生的一個階段.
成績就不是那麼理想.
所以爸爸曾經對我說過一些很失望的話.

$^{201}$你快點去學車.
將來你開起重機.
其實我想現在看不是差.
你現在AI暫時都未取代到這份工.
很多工都是你現在越高人工那些.
越取代得快.
這些基層工作沒有什麼人會取代.
但你明白當時我就覺得很慘.
爸爸這樣看我.
覺得我沒有前途.
覺得我讀不到書.
其實讀不到書是那個考試制度令我讀不到書.
我現在都不知道多喜歡讀書.
我現在都不知道多喜歡去學習.
但當時父親就因為望子成龍的心態.
他們喜歡用激將法.
他就說要激你 等你發奮.
所以這個就成為我心中常常的一個遺憾.
但爸爸其實在他生命裡面都有成長.
我記得當我長大出來工作的時候.
我就有一個心.
就是希望能夠成為一個傳道人.
你想想如果爸爸是牧師.
兒子想做傳道人 當然很開心.
誰知道他又不開心.
他說什麼呢 兒子做傳道人很慘的.
你最好做執事 做執事是最好的.
爸爸這樣說又令我又傷心.
我爸爸說不要緊.
我們一起祈禱尋求主好不好.
我每天上班之前有半個小時.
還沒有出門之前就和爸爸禱告.
真的很奇怪.
禱告將我們的心教正.
因為如果我這個只是一時衝動.
禱告已經消滅了這個感動.
但是越禱告 神又改變我和我爸爸.
我越來越確認神是揀選我成為他的僕人.
我爸爸就明白.
兒子 我都知道你要走傳道的路.

$^{241}$我都不阻你了 為什麼.
我就是亞伯拉罕.
我要獻耳撒.
獻耳撒就是要放手.
現在是神負責你的將來.
我供養你讀完書.
你自己決定是怎樣.
就這樣說.
但是我爸爸其實心裡很開心.
神一直改變他.
後來到我去讀神學的時候.
他經常都很開心.
這個兒子真是一個奉獻給神的兒子.
其實我覺得做爸爸都是這樣.
其實我們一直在學習怎樣做爸爸.
沒有一個爸爸是完美的.
沒有一個爸爸一出生就懂得做爸爸.
很多學問都有讀書有學位拿.
但是做父母就沒有.
而且做父母你都還沒讀好那樣東西.
你就已經要做了.
所以你說邊做邊學.
我這樣說不是對我爸爸不敬.
其實我是很尊敬我爸爸.
因為我爸爸真是在神裡面.
很認真地讀事實.
在他生命裡面.
我從他裡面看到.
他怎樣對神話語的重視.
其實做爸爸有一樣東西.
你猜猜子女最記得的是什麼.
就是帶我去玩和我下棋.
我和別人做安息禮拜很多.
通常寫家庭或者歷史.
通常那些子女經常記得爸爸.
就是記得爸爸帶她去玩帶她去吃東西.
不知道你會不會記得這件事呢.
這個做爸爸真是一個秘訣.
有時多些時間和你的子女有獨處的時間.
你的子女就一定很記得你.

$^{281}$聖經這個很重要的在主裡.
就說到孩子聽從他的父母.
是在主裡的聽從.
意思就是父母的權柄不是絕對的.
而是在神的真理和愛之下.
父母不是用一個威權去壓迫他的孩子.
孩子也不是盲目跟從錯誤的要求.
家庭的忠心不是父母的脾氣.
也不是孩子的慾望.
而是主耶穌基督的真理.
在在當中有很多是幼稚園的家長.
孝敬父母可以從很小的事開始.
願意聽爸爸媽媽的話.
學習說多謝對不起.
吃飯的時候不任性.
玩完玩具願意收拾.
對長輩有禮貌.
睡覺前向父母說晚安.
這些都是生命品格的開始.
但是各位家長我們也要記得.
孩子學習孝敬.
不單只是靠我們的吩咐.
更加是藉著我們的新教示範.
孩子看見我們怎樣對待自己的父母.
將來很可能也會照樣對待我們.
如果我們尊敬上輩.
關心父母.
常常感恩的話.
孩子就能夠在生活中學習什麼是孝敬.
以弗所書六章四節說.
作父親的你們不要激怒兒女.
但要照著主的教導和勸誡養育他們.
聖經不單要求孩子要聽父母.
也要求父母按神的心意對待孩子.
原來在當時的社會裡.
父母是操生殺大權的.
如果那個爸爸說.
孩子我不要了.
那這個孩子馬上就會沒命了.
父親擁有這些權柄.

$^{321}$但是聖經提醒我們.
你有權柄不代表你可以任意發怒.
你有責任不代表你可以用羞辱.
恐嚇比較的方式來控制你的孩子.
相反父母是孩子生命的牧者.
可能你也聽過詩篇二十三篇.
耶和華是我的牧者.
我必不自缺乏.
原來我們是孩子生命的牧者.
我們要帶他們走去青草地.
溪水旁得著餵養.
所以我們又發現.
孩子最終不是屬於我們的.
那是神給我們託管的產業.
原來真正愛孩子的是我們的天父.
是我們的神.
所以今天我們更加要學習.
不要激怒我們的孩子.
原來聖經的管教.
不單是怒氣的發洩.
而是一個愛中的引導.
什麼叫做令孩子發怒呢.
就是父母要求往往是不一致.
今天可以是這樣.
明天又不可以.
孩子覺得無所適從.
又或者父母只是責罵不去聆聽.
看見錯處又看不到孩子的努力.
又或者常常將孩子和別人比較.
這些說話可能令孩子一時間會聽話.
但是長遠其實會傷害孩子的心.
反過來我們怎樣在主裡面養育他們呢.
我們在教養裡面.
我們要有清晰的界線有規範.
我們願意聆聽和引導.
我們願意教導人去認錯和道歉.
指向神的赦免.
我不知道大家覺不覺得.
現在這個社會覺得認錯是很難.
我們發現我們做了一些令人傷害這麼大的事的時候.

$^{361}$我們都不願意去認錯.
認錯讓人覺得我們是弱者.
但是真正的強者就是他肯承認錯誤.
今天我們以為嚴厲就是教育.
其實聖經教導的管教絕對不是怒氣的發洩.
而是最終的引導.
神自己也是這樣待我們.
祂不縱容罪但是祂滿有憐憫.
祂會管教我們.
但是祂的管教是為了使我們得生命.
經文說「要著主的教導和勸戒養育他們」.
這句話提醒我們父母不單止是孩子的供應者.
也是孩子的牧者.
父母不單止是負責孩子讀書,吃飯,睡覺.
更要引導孩子認識真理,分辨對錯,建立敬畏神的心.
猶太人的教育理念幫助我們明白這一點.
在《申命記》六章四到七節這樣說.
「以色列,你要聽,耶和華,我們的上帝是獨一的主.
你要盡心盡性盡力,愛耶和華,你的上帝.
我今日吩咐你的這些話都要記在心上.
也要殷勤教導你的兒女.
無論坐在家裡,走在路上,躺下起來都要淫訟」.
這六章四至五節,紅色這個字就是特別說到猶太人.
那個最重要每天禱告的禱文.
就是讓孩子每天都重念這些禱文.
就是說到耶和華,我們的上帝是獨一的主.
如果你記得耶穌,人家問什麼誡命是最大的時候.
耶穌就是引用這個《申命記》六章五節.
你要盡心盡性盡力,愛耶和華,你的上帝.
這個就是最大的誡命.
其次也相仿,就是要愛人如己,或者愛倫如己.
這個經文特別提醒我們.
今日我們的生命是建基於認識這位獨一的主.
和我們要全心全意的愛我們的上帝.
而我們也要留意,這個教導是貫徹在我們的生活當中.
猶太人日常信仰的節奏,不單是說到宗教的活動.
而是在我們每一天生活的細節.
當我們起身的時候,我們就教導我們的孩子迎接新的一天.
常存感恩的心.
當我們行在道上的時候.

$^{401}$我們就是觀察這個神所創造的大自然,向神發出讚歎.
當我們在家裡,我們吃飯很重要.
我之前也說過,吃飯是一家人能夠一起相聚的時間.
所以吃飯千萬不要拿開手機來看.
吃飯我們就專注,不是吃飯.
是專注在吃飯前後我們的聊天.
當我們願意去分享,當我們願意去祈禱.
將我們生命裡的東西交託神的時候.
家人的關係就能夠更加密切.
我們常常遇見一個問題.
就是當青少年長大了的時候.
越來越不肯跟我們說話.
所以我們要找機會跟他們一起去說話.
去溝通,彼此明白了解.
所以在飯桌上的溝通很重要.
我們一起去聊天,說話.
我們一起為對方祈禱.
甚至在我們小組裡推動家庭崇拜.
我們能夠藉著詩歌,藉著經文,藉著代禱.
我們彼此的建立.
這個是很重要的一個生命的練力.
其實有些什麼你能夠給孩子呢?.
就是能夠給你的心,能夠給你的時間.
親自培育你的孩子,讓他在愛中去成長.
猶太人重視的家庭教育.
他相信孩子是神所託付的.
父母是有責任將神的話傳給下一代.
我們留意今天孩子是受世界價值觀的影響.
我們發現今天我們的孩子.
是很容易產生一個與人比較.
和重視物質的生活,重視成就.
這些都是世界價值觀對我們的影響.
而我們發現青少年抑鬱焦慮的情況.
也是非常的普遍.
有些人告訴我們.
在成年的群體裡面.
有15-30\%都有情緒的問題.
其中抑鬱焦慮的問題是很嚴重的.
有些人說我們屬靈,我們愛主就不會有抑鬱的.
我告訴你,我就是抑鬱過了.

$^{441}$我是做傳道人,做牧師的,我都抑鬱過了.
但是抑鬱做什麼?除了祈禱之外.
更加重要的是吃藥.
就是你一定要吃對的藥.
調好你身體裡面的化學物質的比例.
還有你的屬靈生命靠著神去克服.
還有我覺得有時候這個情緒病是有些嚴重的.
有些嚴重的時候我們真的要求主幫助我們的孩子.
陪伴我們的孩子經歷這段艱難的時刻.
我們知道信仰能夠帶給我們的出路.
不過我們說我們要去影響我們的孩子的時候.
第一步我們要先將神的話記在心上.
父母自己要先去渴慕這個真理.
然後才可以透過生命的示範.
讓孩子都能夠體會認識這些真理.
猶太人教育還有一個很重要的特色.
就是重視故事和紀念.
每逢二月節猶太人的家庭會一代一代講述.
神怎樣去拯救以色列人離開埃及.
那些小朋友就會問.
為什麼我們今晚和其他的晚上不同呢?.
父母就會講述神的作為.
原來孩子不單要學習規條.
而是要進入神救贖的故事.
今天我們也要將福音的故事告訴孩子.
神創造我們愛我們.
人犯罪離開神.
但是神沒有放棄過我們.
祂差遣耶穌基督為我們死在十字架上.
第三天復活.
凡相信祂的人就罪得赦免.
成為神的兒女.
我們的孩子需要知道.
祂的價值不單是來自成績.
不單是來自祂的表現.
更加是因為祂是神所創造的.
是神所愛的.
所有的宗教人都不敢說.
耶穌是道路,真理,生命.
若不是藉著祂.

$^{481}$沒有人能夠去到負那裡去.
耶穌為人也沒有一個宗教的人說.
祂為人的罪怕身在十字架上.
更加在第三天復活.
所以今天什麼是福音呢?.
福音就是我們要接受自己是一個罪人.
但是同時接受耶穌基督.
作為拯救我們的主.
和統領我們生命的主.
今天我們怎樣去回歸神呢?.
我們就是要回歸神的懷抱.
耶穌不是等到我們變好才接納我們.
祂邀請那些不完美的父母.
先來領受神完美的愛.
所以我有幾樣東西可以鼓勵我們在座的父母.
可以去實踐的.
第一樣就是我們每天去祝福我們的孩子.
我們怎樣能夠.
不要整天說.
催他.
快點啦.
不要這樣.
你要多點說.
天父是很愛你的.
你要多點去欣賞和鼓勵他.
另一方面.
我們每天都要為孩子做一個簡短的禱告.
特別是他上學前或是臨睡前.
讓孩子知道他害怕的時候.
是可以去求神.
感恩的時候.
也可以去讚美神.
第三就是要有穩定教會的生活.
我們作為家長.
不要只想著把孩子送來教會就算.
我們都一起在教會裡面.
領受神的真理真道.
以致我們知道怎樣去教養我們的兒女.
又或者作為父母.
我們自己都要持續地去學習真理.

$^{521}$有些家長可能會說.
我都未信主.
我不是太明白聖經.
但是不要緊的.
信仰的開始不是說你已經完全明白.
而是你願意追求.
你可以由回教會開始.
由聽聖經開始.
由祈禱開始.
神很樂意帶領我們每一個人.
願意尋求祂的人.
聖經教導.
孝敬父母.
但是更深一層.
就是帶我們去認識天父.
地上的父母有愛.
但是也有限制.
地上的父母會疲倦.
會犯錯.
會有情緒.
但是天父的愛是完全的.
他創造我們.
認識我們.
愛我們.
也都願意拯救我們.
所有父母都想給孩子最好的.
但是我們不能夠給孩子永恆的生命.
我們不能夠完全保護孩子.
一生不遇到風浪.
也都不能夠代替孩子面對罪和死亡.
只有耶穌基督可以拯救我們.
聖經說我們都犯了罪離開了神.
罪不單止是做壞事.
更加是以自我為中心.
不按神的心意生活.
罪使我們與神隔絕.
也都破壞了我們與神的關係.
但是神愛我們差遣耶穌來到世上.
耶穌並沒有犯罪.
卻為我們的罪死在十字架上.

$^{561}$承擔我們的刑罰.
他第三日復活.
勝過罪和死亡.
凡願意相信他.
接受他作救主和生命之主的人.
就可以罪得赦免.
成為神的兒女.
得著新的生命.
今天耶穌不單止是邀請他的孩子來認識他.
也都邀請父母去認識他.
神不單止要改善我們教育的方法.
更加要更新我們的生命.
當父母先被神的愛觸摸.
家庭就會開始改變.
當父母學習在主裡面生活.
孩子亦都會在父母身上.
看見信仰的真實.
或者你會覺得自己都不是一個好的父母.
甚至你會發脾氣.
會焦慮.
甚至不懂得怎樣教.
但是耶穌不是要等你變好.
才接納你.
他今天就邀請你來到他的面前.
他願意赦免你.
醫治你.
給你新的力量.
亦都帶領你的家庭.
以下我們一起禱告.
親愛的主耶穌.
多謝你愛我.
多謝你愛我的家庭.
我承認我有罪.
常常靠自己生活.
沒有按你的心意而行.
多謝你為我的罪死在十字架上.
並且從死裡復活.
今日我願意打開心門.
接受你作我的救主和生命的主.
求你赦免我的罪.

$^{601}$賜我新的生命.
亦都求你幫助我學習.
作合你心意的父母.
帶領我們的孩子在你的愛和真理中成長.
我們這樣在你面前奉主耶穌基督的聖名祈求.
如果剛才這個禱告是你的心願的時候.
我很樂意鼓勵大家.
繼續的去認識這個真理.
繼續的去相信.
唯有神的話才能夠改變我們的生命.
改變我們孩子的未來.
如果你剛才是真心有這個祈禱.
我相信這個是你人生一個很重要的開始.
盼望你繼續的不是今天回來教會.
你繼續的去認識神的話語.
與楊弟兄姊妹陪您同行.
\newpage



\section{申命記 9:1-29}
\label{sec:032309}
\textbf{感恩 MODE (2) 恩典中的感恩}~~~~王啟恩牧師
\newline
\newline
連結: \href{http://www.church.com.hk/acms/content.asp?site=cdc\&op=show\&type=product\&code=032309\&layout=sermon}{\texttt{032309}} ~~~~ 語音日期: 2026-06-13
\newline
\newline
\hyperref[sec:032308]{\small{< < < PREV SERMON < < <}}
~
\hyperref[sec:index_chronic]{\small{[返順時目]}}
~
\hyperref[sec:032310]{\small{> > > NEXT SERMON > > >}}
\newline
\newline
申命記 9:1-29
\newline
\begin{longtable}{cl}
\hline
\hline
章節 & 經文 (和合本修訂版)\\
\hline
9:1 & \begin{tabularx}{0.7\textwidth}{X} 「以色列啊，你要聽！你今日要過約旦河，進去佔領比你更強大的列國，那裡的城鎮又大，城牆又堅固，如天一樣高。 \end{tabularx} \\ \\ \relax
9:2 & \begin{tabularx}{0.7\textwidth}{X} 那裡的百姓是亞衲族人，又高又壯，是你所知道的；你也聽說過：『誰能在亞衲族人面前站立得住呢？』 \end{tabularx} \\ \\ \relax
9:3 & \begin{tabularx}{0.7\textwidth}{X} 你今日應當知道，耶和華---你的神在你前面渡過去，如同吞噬的火，要除滅他們，並要在你面前將他們制伏，使你可以趕出他們，速速消滅他們，正如耶和華向你所說的。 \end{tabularx} \\ \\ \relax
9:4 & \begin{tabularx}{0.7\textwidth}{X} 「耶和華---你的神將他們從你面前趕出以後，你心裡不可說：『耶和華領我得這地是因我的義。』其實，耶和華將這些國家從你面前趕出去是因他們的惡。 \end{tabularx} \\ \\ \relax
9:5 & \begin{tabularx}{0.7\textwidth}{X} 你能進去得他們的地，並不是因你的義，也不是因你心裡正直，而是因這些國家的惡，耶和華---你的神才把他們從你面前趕出去，為了應驗耶和華向你列祖亞伯拉罕、以撒、雅各起誓應許的話。」 \end{tabularx} \\ \\ \relax
9:6 & \begin{tabularx}{0.7\textwidth}{X} 「你當知道，耶和華---你的神將這美地賜你為業，並不是因你的義；你本是硬著頸項的百姓。 \end{tabularx} \\ \\ \relax
9:7 & \begin{tabularx}{0.7\textwidth}{X} 你要記得，不要忘記，你在曠野怎樣惹耶和華－你的神發怒。自從你出了埃及地的那日，直到你們來到這地方，你們常常悖逆耶和華。 \end{tabularx} \\ \\ \relax
9:8 & \begin{tabularx}{0.7\textwidth}{X} 你們在何烈山惹耶和華發怒，耶和華對你們動怒，甚至要除滅你們。 \end{tabularx} \\ \\ \relax
9:9 & \begin{tabularx}{0.7\textwidth}{X} 我上了山，要領受兩塊石版，就是耶和華與你們立約的版。那時我在山上住了四十晝夜，沒有吃飯，也沒有喝水。 \end{tabularx} \\ \\ \relax
9:10 & \begin{tabularx}{0.7\textwidth}{X} 耶和華把那兩塊石版交給我，是神用指頭寫成的；版上是耶和華在大會的那一天，在山上從火中對你們所說的一切話。 \end{tabularx} \\ \\ \relax
9:11 & \begin{tabularx}{0.7\textwidth}{X} 過了四十晝夜，耶和華把那兩塊石版，就是約版，交給我。 \end{tabularx} \\ \\ \relax
9:12 & \begin{tabularx}{0.7\textwidth}{X} 耶和華對我說：『起來，趕快從這裡下去！因為你從埃及領出來的百姓已經敗壞了；他們這麼快偏離了我所吩咐的道，為自己鑄造偶像。』 \end{tabularx} \\ \\ \relax
9:13 & \begin{tabularx}{0.7\textwidth}{X} 「耶和華對我說：『我看這百姓，看哪，他們是硬著頸項的百姓。 \end{tabularx} \\ \\ \relax
9:14 & \begin{tabularx}{0.7\textwidth}{X} 你且由著我，我要除滅他們，從天下塗去他們的名，我要使你成為比他們更大更強的國。』 \end{tabularx} \\ \\ \relax
9:15 & \begin{tabularx}{0.7\textwidth}{X} 於是我轉身下山，山上有火燃燒，兩塊約版在我雙手中。 \end{tabularx} \\ \\ \relax
9:16 & \begin{tabularx}{0.7\textwidth}{X} 我觀看，看哪，你們得罪了耶和華－你們的神，為自己鑄成了一頭牛犢，迅速偏離了耶和華所吩咐你們的道， \end{tabularx} \\ \\ \relax
9:17 & \begin{tabularx}{0.7\textwidth}{X} 我拿著那兩塊石版，從我雙手中扔出去，在你們眼前把它們摔碎了。 \end{tabularx} \\ \\ \relax
9:18 & \begin{tabularx}{0.7\textwidth}{X} 因為你們所犯的一切罪，做了耶和華眼中看為惡的事，惹他發怒，我就像從前一樣俯伏在耶和華面前四十晝夜，沒有吃飯，沒有喝水。 \end{tabularx} \\ \\ \relax
9:19 & \begin{tabularx}{0.7\textwidth}{X} 我很害怕，因為耶和華向你們大發烈怒，要除滅你們。但那一次耶和華又應允了我。 \end{tabularx} \\ \\ \relax
9:20 & \begin{tabularx}{0.7\textwidth}{X} 耶和華也向亞倫非常生氣，甚至要除滅他；那時我也為亞倫祈禱。 \end{tabularx} \\ \\ \relax
9:21 & \begin{tabularx}{0.7\textwidth}{X} 我把那使你們犯罪所鑄的牛犢拿來，用火焚燒，搗碎後再磨成粉末，好像灰塵。我把這灰塵撒在從山上流下來的溪水中。 \end{tabularx} \\ \\ \relax
9:22 & \begin{tabularx}{0.7\textwidth}{X} 「你們在他備拉、瑪撒、基博羅‧哈他瓦又惹耶和華發怒。 \end{tabularx} \\ \\ \relax
9:23 & \begin{tabularx}{0.7\textwidth}{X} 耶和華叫你們離開加低斯‧巴尼亞，說：『你們上去得我所賜給你們的地。』那時，你們違背了耶和華－你們神的指示，不信服他，不聽從他的話。 \end{tabularx} \\ \\ \relax
9:24 & \begin{tabularx}{0.7\textwidth}{X} 自從我認識你們的日子以來，你們常常悖逆耶和華。 \end{tabularx} \\ \\ \relax
9:25 & \begin{tabularx}{0.7\textwidth}{X} 「我因耶和華說要除滅你們，就在耶和華面前俯伏四十晝夜，像我以前俯伏一樣。 \end{tabularx} \\ \\ \relax
9:26 & \begin{tabularx}{0.7\textwidth}{X} 我向耶和華祈禱，說：『主耶和華啊，求你不要滅絕你的百姓，你的產業。他們是你用大能救贖，用你強有力的手從埃及領出來的。 \end{tabularx} \\ \\ \relax
9:27 & \begin{tabularx}{0.7\textwidth}{X} 求你記念你的僕人亞伯拉罕、以撒、雅各，不看這百姓的頑梗、邪惡、罪愆， \end{tabularx} \\ \\ \relax
9:28 & \begin{tabularx}{0.7\textwidth}{X} 免得你領我們出來的那地之人說：耶和華因為不能將這百姓領進他所應許之地，又因恨他們，所以領他們出去，要在曠野殺他們。 \end{tabularx} \\ \\ \relax
9:29 & \begin{tabularx}{0.7\textwidth}{X} 其實他們是你的百姓，你的產業，是你用大能和伸出的膀臂領出來的。』」 \end{tabularx} \\ \\
[1ex]
\hline
\hline
\end{longtable}
$^{1}$今天的講題是恩典中的感恩.
經文是在《申命記》第九章.
今天有三個重點跟大家分享.
第一個重點就是感恩的起點.
我們要先明白為何我們需要感恩.
首先我們要認清楚自己的本相.
當我們不認清自己的本相的時候.
或許我們會忘記了.
我們一起看看經文.
先看第九章四至六節.
我們一起讀.
.
弟兄姊妹,這裡說到摩西對以色列人發出一個很嚴厲的警告.
這個嚴厲的警告是什麼呢?.
就是不要把恩典當作自己努力得來的獎賞.
因為恩典是神賜予的,是白白的.
我們所做的事不是我們做了什麼有多厲害,而是得來的.
所以如果我們要活出感恩的生命.
我們必須要很誠實地面對自己的本相.
我們的本相是什麼呢?.
大家有沒有想過這個問題?.
如果看回經文,我們的本相有兩件事.
第一件事是說到我們是不義的,是不正直的.
剛才大家讀了經文.
這裡有三節的經文都說到.
不可以說因為我的義.
不是因為你心裡正直.
也不是因為你的義.
即是神令他們得勝.
不是因為你們做了什麼,不是因為你們公義.
義是說一些什麼呢?.
義是說一些善良的行為.
即是一些義行.
而心裡正直是說全心的純正.
高尚.
給大家一個例子.
有時候我也覺得自己很公義.
去到地鐵的時候,我永遠不會坐關愛座.
大家會不會坐關愛座呢?.
不會吧?.

$^{41}$話說有一次,去到一個站.
開了,我知道一個已經是小學生.
跳踏踏的樣子.
走進來,一得就得了去關愛座.
坐在關愛座那裡.
坐在關愛座前面有個老人家.
他不讓給他.
我看到,你猜我做過什麼動作呢?.
我就跟那個小學生說.
喂,小朋友,你剛才走進來跳踏踏.
你也挺年青力壯的.
不如你讓個位給這個老人家吧.
那你猜他讓不讓呢?.
喂,讓了,他聽我說.
我真的很高興.
驚訝地聽我說.
這件事情完了.
回想我自己.
有時候我在地鐵,坐地鐵.
我不會坐關愛座.
但有座位我會坐.
有時候如果我很累的時候.
我就會跟自己說.
最好不要讓位出去.
最好不要有個老人家站在我前面.
或者有個懷孕的.
或者有些戴著拐杖的.
站在我前面.
有時候也會的.
覺得這個小朋友.
應該要讓位給別人.
我就想,最好不要了.
有時候是這樣的.
當我累的時候.
所以人有時候也會很有趣.
做些事或者想些事.
另一方面我們要留意一個重點.
人有罪性.
不知道大家認不認同.
這裡說的是什麼呢?.

$^{81}$說的是內心深處與生俱來的.
一種自我中心的傾向.
是一切犯罪行為的根源.
人並非單純犯罪才成為罪人.
大家留意一點.
而是因為內心帶有一個罪性.
因此我們必然會表現出各種的罪.
有什麼例子呢?.
我看到有兩個小朋友.
說回小朋友.
你會不會教你的兒子自私或貪心呢?.
兒子你最好貪這些.
兒子你最好自私.
你自己有的東西不要跟別人分開和玩.
有時候小朋友是這樣的.
看到一些東西.
你會說,爸爸,我又要,我又要.
有些小小的貪念.
另一方面如果擁有一些東西.
有些小朋友跟你說.
你讓我玩吧.
有時候就很捨不得.
去給別人玩.
有時候可能也會給別人玩.
但最後我也要收回.
是我擁有的東西.
我們看看羅馬書怎樣說.
羅馬書三章十至十二節.
我們一起讀.
就如經上所記.
沒有二人.
連一個也沒有.
沒有明白的.
沒有尋求神的.
都是偏離正路.
一同變為無用.
沒有行善的.
連一個也沒有.
我們的本相就是這樣.
弟兄姊妹.

$^{121}$另外一個是甚麼呢.
剛才也有說.
硬著頸汗的百姓.
硬著頸汗的意思是.
還緊.
活逆.
好像這兩隻牛.
這兩隻公牛原本是耕田的.
公牛耕田就要套上鴨.
我們有時候也會這樣.
我們不願意套上神給我們的鴨.
我們不肯負這個鴨.
大家清楚這個鴨是甚麼嗎.
這個鴨就是要跟著神.
這個鴨就是神的使命.
我們去履行神的使命.
剛才也有說.
高宗生有三個反省.
很深刻的反省.
他第一個反省是甚麼呢.
第一個反省就是這個.
他就說我輸了輕飄飄的虛榮.
卻是贏得沉甸甸的懂得.
開始懂得.
大家聽了真人真事的這篇文章.
虛榮即是說.
我需要一部甚麼的漂亮的電話.
拿出來令隻散鏡的那種虛榮.
但他沒有買到.
最後還給他錢.
他就說我贏得沉甸甸的懂得.
究竟懂得甚麼呢.
大家回顧我剛才所說的字跡.
即是這個文章.
我相信是這樣.
他懂得爸爸媽媽賺的錢.
全部都是血汗錢.
辛勞和當中的愛.
他懂得我覺得是醒覺.
即是醒覺自己身在不福中不知福.

$^{161}$原來是這麼有福氣.
爸爸媽媽其實這麼愛他.
他懂得認清楚自己的本相.
他自私自我中心的那種本相.
還有有名牌電話不是必然的.
他懂得去留意和欣賞他爸爸媽媽.
他以前從來沒有這樣做過.
所以其實他一開始已經明明.
他是活在一個恩典的裡面.
但是他看到的只是自己沒有的東西.
我們今天不知道會是怎樣.
有沒有一些東西你覺得沒有恩典.
有沒有看到你有的東西呢.
弟兄姊妹.
當我們認清楚自己的本相和無有的時候.
我們才能夠明白.
能夠蒙恩不是因為我們自己配得.
而是全部都是神的恩典.
而當我們深深感受到神的恩典的時候.
我們自然就懂得感恩.
即是相反來說.
當我們將恩典當作理所當然的時候.
這些東西一定有的.
神一定給我的.
我們的感恩的心就會停止.
什麼叫做真正的感恩呢.
真正的感恩就應該是這個.
就是當我們站在恩典裡面的時候.
低頭望向你自己.
然後你很驚訝地說.
我們一起讀吧.
一二三.
主啊,我其實是個硬著頸汗的人.
我不配得這一切.
但你竟然將最好的給了我.
突然間神給了你什麼.
你接收了神的東西.
這就是最好的東西.
這一種感恩會帶來很大的謙卑.
同時之間也是神的喜悅.

$^{201}$第二個重點.
盼望的感恩.
因為我們有盼望.
所以我們感恩.
為什麼呢.
因為神已經為我們跨過生命的巨人.
我們又要看九章一至三節.
我們一起讀吧.
好,一二三.
以色列啊,你當聽.
你今日要過約旦河.
進去講出比利強大的國民.
得著廣大堅固高得頂天的誠意.
那民是阿納族的人.
又大又高是你所知道的.
也曾聽見有人指著他們說.
誰能在阿納族人面前站立得住呢.
你今日當知道.
耶和華你的神在你面前過去.
如同烈火要滅絕他們.
將他們制伏在你面前.
這樣你就要照耶和華所說的.
趕出他們.
使他們速速滅亡.
這裡說到神要以色列民去過約旦河.
趕出比利強大的阿納族人.
回想當年.
即是三十八年前.
大概,三十八年前.
當年有十二個探子.
大家應該都聽過.
十二個探子去窺探迦南地的時候.
其中有十個.
即是除了迦納和約書亞之外.
其他十個.
嘩,看到裡面的人又神高神大.
嘩,真是發震了.
因為自己好像眨眼.
那些巨人一捏就捏死你.
即是那一種.

$^{241}$因而就和摩西說.
不行不行.
我們不可以去攻打這個地方的人.
就是因為這件事情.
令到以色列民要在曠野去漂流多.
三十八年了.
這是在文述的十三章.
說到出來的.
三十八年後,神知道他們的過去.
在人的眼光裡面.
看來這一場好像絕對打不贏的仗.
但是這一段我們剛才所讀的經文.
告訴我們.
在智慧行間裡面.
我們會意識到.
神是如今.
他是保證.
會帶領以色列民得到勝利.
而他們要做的事情就是.
很簡單.
要照耶和華所說的.
最後highlight的紅色一段經文.
我們一起讀.
要照耶和華所說的.
這個就是.
神在你面前過去.
即是之前的紅色字.
如同烈火的意思.
等一等,姐妹,如果你今天只是面臨人生的約旦河.
前面看來好像沒有辦法跨過的困難.
甚至好像有很強大的屏障擋住你.
不知道是什麼.
可能是你的財政壓力.
財務壓力.
家庭的關係.
工作上的瓶頸.
等等.
但是神跟你說.
你今天要過約旦河.
盡量說出給你強大的國民.

$^{281}$因為.
神也會在你面前過去.
如同烈火一樣.
所以我們人生的.
在我們人生的巨人當中.
神已經在我們的前頭為我們去爭戰.
大家相不相信.
神已經在我們的前頭為我們去爭戰.
如果相信的話就說阿們吧.
Amen.
感恩這兩個字.
不只是為了我們在手頭已經得到一些物質的好處.
更加為了神在我們生命裡面所跨出的每一步.
就好像神就像烈火一樣.
在我們面前走過.
燒盡在我們面前.
爛了我們屬靈的生命的成長.
爛了我們蒙福的障礙.
我們不要看輕.
這個其實是一個確據.
這個確據的意思就是說.
是來自耶穌基督已經復活了.
耶穌基督已經復活了.
祂已經戰勝魔鬼仇敵.
而我們真正的聖像不在乎我們.
我們剛才的知覺叫在乎.
你就是在乎神.
我們的聖像不是在乎我們.
現在在乎神.
我們不是靠自己去打仗.
而我們只是跟隨一位已經得勝的元帥.
進入他的得勝的裡面.
就是說我們正在打一場必勝的仗.
必勝的仗.
來到這裡大家會不會感覺到.
突然間好像有個遠徑.
這個代表什麼呢.
彩虹.
是不是人生換希望.
雖然老套一點.

$^{321}$人生換希望也要說.
人生換希望.
我們不單止有永生的盼望.
而我們今生也得到倚靠.
同時之間也可以得到一個得勝的生命.
得勝的人生.
只要.
當然有前提.
有前設.
只要你.
question mark.
你覺得是什麼呢.
剛才讀的是經文.
這裡也有的.
兩個字.
只要你做過什麼動作呢.
好像聽到了.
隱約.
當然是動詞.
是跟隨.
這裡有兩個字.
跟隨.
我們要跟隨.
跟隨的意思是按照主所說的去做.
也是剛才經文所說.
我們要跟隨去做.
跟隨去做就可以了.
跟隨.
弟兄姊妹.
大家有沒有盼望.
精神一點好不好.
有沒有盼望.
有了.
我們已經有了耶穌.
怎麼可以沒有盼望呢.
所以其實大家已經很有盼望.
有了盼望其實我們已經很感恩.
我們的盼望不單止是回天家那種盼望.
永生的盼望.
在世的日子其實神有看顧我們.

$^{361}$在世的日子我們有倚靠.
神走在我們前頭.
看不看得到經文.
走在我們前頭.
好像烈火一樣.
燒盡我們前面那種障礙.
所以大家有盼望的.
沒有盼望的人會是怎樣呢.
會出現這些現象.
第一.
他會失去前進的動力.
接著.
會帶來心軟的心理和生理的危害.
第三.
往往會陷入這個.
習得性無助.
我有沒有聽過這個詞語.
我解釋給大家聽.
這是一種精神的狀況.
是指一個人.
當他重複經歷挫折.
和沒有辦法去控制的負面事件之後.
即使後來會有機會去改變現狀.
他仍然選擇放棄嘗試.
而變得消極,被動.
以為自己沒有辦法脫離這個困境.
就叫做習得性無助.
即使有機會翻身.
他也放棄了.
因為他覺得不行了.
擺脫不了這個困境.
當我們沒有盼望的時候.
就會是這樣.
接著就是.
長期處於負面情緒當中.
同時可能會削弱一些免疫力.
誘發一些抑鬱症.
甚至可能在絕望當中.
放棄生存的意志.
當我們有盼望的時候.

$^{401}$我們應該不會這樣.
當我們反省心想.
有些事情在我生命中正在發生.
大家要回到神那裡.
你要看看盼望是什麼.
神在你的生命中是什麼.
第三點.
感恩的傳言.
是耶穌基督中堡的恩典.
接著我們看第九章七至二十四節.
這裡說到什麼呢.
說到摩西帶領百姓進行一次.
很深刻的歷史回顧.
一段充滿活躍的歷史.
摩西很詳細地數算.
以色列人的失敗.
主要有三方面.
在經文中有三方面.
第一件事情是金豪毒事件.
有沒有人沒聽過金豪毒事件.
應該有聽過吧.
摩西上山領受十誡和律例典章四十日.
四十晝夜.
以色列人看到摩西去了哪裡.
會不會再回來呢.
接著向阿倫申訴.
為什麼不下來.
我們需要怎麼做.
阿倫要了他們的金器.
注了一種所謂的金豪毒.
這就是金豪毒事件.
期盼金豪毒會帶領他們走前面的路.
大家都知道這是一個拜偶像的表現.
惹得耶和華發怒.
要撕掉他們.
第二件事件是向神發預言.
在三個地方.
在二十二節說到.
這三個地方發生了什麼事情呢.
在塔比拉, 以色列人很愕然地去抱怨神.

$^{441}$大家要留意.
在瑪薩, 因為沒有水飲而發預言.
在基博羅哈他亞, 很懷念那時在埃及的日子.
很有肉吃, 出到來.
他那些恩賜的嗎哪不行, 我要吃肉.
這三件事件.
為什麼摩西要數呢.
還是三件事件, 他只是數三件.
其實不止三件.
帶出一個很重要的事情.
第三件, 先說完.
就是十二個探子的事件.
剛才說了, 十二個探子.
在這裡, 申命記也有提到這件事件.
這三件事件在一起.
說到一個重要的事情.
事實就是說, 一連串的事件.
證實人不是偶然為過犯所幸.
而是連貫性的犯罪損益.
不是一兩件事.
而是一連串的事件.
加在一起.
其實說的是兩件事.
第一件事, 他們不要神.
人要為自己去做假神, 偶像去代替真神.
今後毒事件就是其中一個.
就是剛才所說.
我們心中也會有偶像.
當然我們不會弄一個像去拜.
大家都會明白.
將一些事情, 比神重要的時候.
沒了他不行的時候.
其實也是你的偶像.
第二個就是不信神.
不信神的意思是, 以自己的心思意念為一切行事.
或者行動的依據.
這樣就等同自己是做神.
即是我要作主.
神, 你走開吧.
我要自己作主.

$^{481}$所以剛才所說的十二探子的事件.
賣怨神的事件.
都是歸類於這個不信神的表現.
好了, 去到申命記.
九章七節.
我們讀一讀, 九章七節.
綠色的字, 123.
你當記念不忘.
即是摩西叫他們.
或者神叫他們.
不要忘記你過往失敗的經歷.
不要忘記以往失敗的經歷.
雖然這樣做, 我不知道大家會不會.
有時會令自己很痛苦.
即是我要記起以前那些痛苦的事.
有時會很憤怒自己.
可能是自己做錯的事.
或者是別人對你不好的事.
一般都是.
失敗的經歷.
但是我們要想.
為什麼神要他們這樣呢?.
你不要忘記以往的經歷.
就是更重要的事情.
就是我們要清楚.
記得其實自己是活躍的.
以色列民你們是活躍的.
還有你們犯錯.
這個是你們的本相.
你不要忘記你的本相是這樣.
接著, 當他們.
即是以色列民認清楚的時候.
他們就可以去學習.
不要驕傲自義.
同時不要自以為是配得.
我們所有的事都是恩典.
是白白的, 雖然.
但是我們不是配得的.
因為我們每一個人都是.
曾經犯罪, 我們在罪的當中.

$^{521}$有時是心思意念.
我們會得罪了神.
這就是為什麼我們在.
這個時候, 這個是什麼?.
什麼日子?.
聖餐, 對了.
聖餐通常我們要怎樣?.
我們要審查自己.
特別是審查自己有沒有犯罪.
得罪了神.
這是其中一個原因.
當然我們不應該只是聖餐.
一個月一次才審查自己.
我自己的習慣是怎樣呢?.
我自己的習慣是這樣的.
我每天都會做這個動作.
睡覺之前.
我會回想那天.
我做了什麼呢?.
或者我祈禱.
神有什麼要提醒我呢?.
特別是我做得不對的事.
特別是有些事我誤會了你.
當然有時想著想著.
就會睡著了.
但這是我的習慣.
很幫助到我.
同時之間.
我也會謹記.
這個我覺得不錯.
謹記自己容易犯錯.
或者容易犯罪的地方.
我會謹記的.
分享給大家.
我最容易犯錯.
或者最容易犯罪的地方是什麼呢?.
大家不要看我斯斯文文.
溫文爾雅.
不會暴躁.
不會摒人.

$^{561}$我有時也會摒人.
我會爆人.
就是當我被人冤枉的時候.
被人逼得很厲害的時候.
所以你千萬不要逼我.
不過如果你想我爆你.
你就逼我吧.
你要逼我埋牆的時候.
我就會爆你.
我不知道大家是什麼.
就是你要提醒自己.
去到那些位置我也會提醒自己.
這些是我的位置.
暫時不要爆.
話說兩年前.
一年半前.
差不多兩年前.
我曾經爆一個童工.
在開會當中.
當然不只是他在現場.
還有其他童工在現場.
當我爆了他之後.
所有童工呆了.
這個是Larry來的嗎?.
就是這樣.
我也會爆人.
提醒自己.
這件事情叫我不要忘記.
我真的不要忘記這件事情.
在什麼情況下.
我不要這樣去做.
摩西在當中.
接著就說了.
為他們代求.
成為關鍵.
摩西提醒他們.
你們不要忘記.
同時之間.
他也為以色列人去代求.
成為他們不被滅絕的關鍵.

$^{601}$他做了什麼動作呢?.
在二十五節說.
我因耶和華.
要滅絕你們.
就在耶和華面前.
照舊伏伏四十晝夜去求情.
令到以色列人.
不被神滅絕.
仍然可以站在神的面前.
摩西為百姓代求的理據有三個.
第一個.
在二十六節.
大家看就行.
這裡說的是什麼呢?.
以色列人是神.
是你所拯救的.
是你分別出來的產業.
為什麼你要滅絕他們呢?.
第二個.
在申命記九章二十七節.
這裡說到.
神顧念你的僕人.
亞伯拉罕以撒亞國.
大家都應該聽到.
神應許亞國是什麼呢?.
他的後裔多於天上的聖聖.
地上是沙.
是神給他們的.
所以他求神.
不要因為以色列民的過犯.
而不顧對列祖的應許.
第三個.
就是說.
這節經文我覺得很有趣.
就是說.
免得在那地方的人.
說你閒話.
你將這些百姓.
領出來.
你又憎恨他們.

$^{641}$領了他們出來.
又要在曠野殺他們.
因為神說要殺他們.
因為他們是活役.
活役的後果就是死.
所以他就說.
你不要因此被外邦人嘲笑你.
或者褻瀆你.
摩西.
為那時候的以色列人.
代求.
我們每一個人在神的面前.
都可能是.
好像當年的以色列人一樣.
會不會還緊.
會不會活役呢.
會不會偏向世界那種.
金牛毒呢.
真的.
如果神是按公義去審判.
其實我們沒有一個人可以站立得住.
其實我們早已經滅亡.
耶穌基督跟他們做了.
其實他給了應許我們.
就好像當年耶和.
給了應許亞伯拉罕.
和基薩亞國一樣.
今天也給了應許我們.
和他為我們代求.
經文在哪裡呢.
主應許的聖靈.
為印記.
既然信他.
就受了所應許的聖靈.
為印記.
在.
以法所述一章.
十三節說得很清楚.
為何這麼重要.
在一章十四節.

$^{681}$之後.
他就說.
因為這個聖靈.
是我們得基業的憑據.
大家想不想得到基業.
這個基業即是說.
神的東西是我們可以擁有.
神是我們的產業.
同時我們也是神的產業.
是互相的.
不單止這樣.
主也為聖徒獻上代禱.
其實在約翰福音.
很多地方.
耶穌基督也是為他們的門徒.
而是說到今時今日.
他也為我們代求.
當我們信了他的時候.
所以我們憑著耶穌基督的代求.
我們終不至滅亡.
反得拯救.
這一節經文說得很清楚.
我們一起讀.
凡靠著他.
進到神面前的人.
他都能拯救.
到底.
他是長遠活著.
替他們祈求.
弟兄姊妹.
昔日由摩西.
為到以色列民人祈求.
今日我們有誰為我們祈求.
耶穌基督.
神.
親自為我們祈求.
給了我們一個嬰兒.
是不是很大.
我剛才答應了.
高宗生有三個深切的反省.

$^{721}$說了一個.
第二個反省是什麼.
第二個反省.
愛不是理直氣壯.
索取的底氣.
媽媽你要愛我.
所以我要索取.
你應該要愛我的.
而是.
我們能夠看見.
看見父母藏在.
很不容易的愛.
裡面.
發自內心的心痛和體諒.
這個高宗生.
我覺得他很聰明.
他懂得反省.
他看見.
媽媽的手.
很粗糙的老人家的手.
滿泥漿.
很多污漬.
裂開的手.
他看見了.
今日我們看見什麼.
我們應該看見這些.
我們看不看見耶穌基督的丁行手.
因他受的懲罰.
我們得平安.
因他受的鞭傷.
我們得醫治.
我們看不看見呢.
我們會不會像高宗生一樣.
去醒覺呢.
醒覺.
我們今天能夠.
活在恩憐的裡面.
大家覺得.
你有沒有平安.
你有沒有盼望.

$^{761}$我們可以有永生.
全憑我們的中寶.
耶穌基督.
全憑他那種.
無條件犧牲的愛.
和剛才所說的代求.
他的應許.
所以當我們明白.
本相是什麼.
本身我們其實是一個罪人.
我們白白蒙受.
但卻是白白蒙受.
基督的寶血.
大家看到丁行手.
會不會聯想到寶血呢.
聯想到這件事情的時候.
怎麼.
能夠.
不從我們心底裡面.
發出那種很深層的.
感恩呢.
這就是我們感恩的源頭.
感恩之源.
耶穌基督是我們的中寶.
他有很大的應許給了我們.
而這一個是永遠不會消失的.
所以我們感恩.
不是因為我們生活裡.
我們拿到一些什麼.
的東西.
而是一切的順利.
而是因為我們.
靈魂得救.
和在世上的日子.
我們是有.
主的眷顧.
有主的同行.
為我們去代求.
這個恩典是永遠不會被奪去.
你跟隔壁左口說.

$^{801}$這個恩典是永遠不會被奪去.
永遠在這裡.
因為神是信實的.
祂說的話.
不會不實現.
好了,高宗新.
第三個反省.
第三個反省是什麼呢.
就是這個,今天的主題.
是什麼.
想起那兩個字.
感恩,恩典中的感恩.
我們一起讀.
發現以前的我.
是多麼的自私.
及虛榮.
能夠寫出漂亮的文章.
卻說不出.
謝謝爸媽的感恩.
說話.
好了,來到這裡.
不知道大家有什麼感應.
有什麼感受.
我覺得.
這個高宗新.
他的醒覺.
他的醒悟很大.
他開始動得感恩.
他不單止開始動得感恩.
這種感恩是.
有附帶的行動.
這個附帶的行動是什麼.
我不要這個電話.
這個電話.
他不捨得要.
不捨得買這一部電話.
弟兄姊妹.
我都曾經經歷過你們的年紀.
我相信你們大概都.
三十出頭.

$^{841}$大概吧,或者二十尾三十出頭.
我告訴大家.
我那時候渾渾噩噩的.
什麼渾渾噩噩.
那時候信了耶穌.
信了耶穌.
所以大家都很乖.
當然有時候.
都會有些罪在裡面.
當然.
不會很壞.
但都知道有些行為.
是神所不起悅.
我那時候就少回教會.
因為已經出來.
工作了,二十八尾.
九十初的時候.
二十尾.
二十尾或者九十初的時候.
二十尾.
三十初的時候.
都在工作.
年輕力壯.
世間上實在太多誘惑.
在我身邊的人.
沒有回教會,全部都不是基督徒.
所以大家都會想像到.
一個渾渾噩噩的青年.
大家好好了.
我很欣賞大家.
但是.
我很感恩,神把我帶到加拿大.
我應該跟大家分享過.
把我帶到加拿大.
在我一個很特別的處境下.
我回到教會.
我從那時候.
才開始.
對我的信仰認真.
我那時候才醒覺.

$^{881}$三十初.
才醒覺.
我要對我的信仰認真.
我聽說.
我聽得入耳.
睡著了.
不是睡著了,是睡著了.
很用心的聽牧師的講道.
翻煮,學習等等.
那時候才醒覺.
弟兄姊妹.
我不可以說大家不醒覺.
但如果你仍然覺得自己.
都不是很醒覺的話.
我真的很提醒大家.
醒了.
當我們面對未來.
面對生活裡面.
真的,我知道有很多挑戰.
但大家仍然年輕.
三十初.
當你有八十多歲的人生.
你都有五十多歲.
即是五十五年.
還有很長的人生路要走.
但是.
很想鼓勵大家.
不要忘記回顧.
神給你的恩典.
這個恩典是白白的.
這個恩典永遠在裡面.
不會奪去.
所以我們可不可以做這個行動.
不要常常以為.
自我誇耀.
我們將這個放下.
成為什麼呢.
對了.
感恩,讚美.
因為我們擁有的一切.

$^{921}$我們有的氣息.
我們有的生命.
我們有的力量.
全部都是出於.
恩典.
所以今天的講大是什麼.
恩典中的感恩.
所以在.
這個申命記第九章.
主要有三個.
的重點.
都是圍繞感恩.
我們要看清楚自己的本相.
我們不是異.
但是神給我們恩典.
我們要看清楚.
其實神已經在我們前面.
為我們去打勝這個仗.
我們要跟隨他.
第三個就是.
神給了我們這麼大的恩典.
和我們做中堡.
耶穌基督在我們中間.
做我們的中堡.
為我們禱告.
為我們祈求.
和他上十字架.
我們可以得到.
這裡寶血的洗淨.
這些全部都是我們的感恩.
的源頭.
所以.
鼓勵大家.
也是鼓勵我自己.
成為生活的一部分.
好不好.
因為當我們這樣做的時候.
還有這個.
by products.
即是bonus.

$^{961}$這些bonus是什麼呢.
自己去看.
當我們這樣做的時候.
我們可以建立更謙卑的品格.
我們可以.
降低我們的焦慮.
可以帶來快樂.
我們可以面對困難的時候.
幫助我們不同的角度.
去看問題.
並且提升我們那種抗壓的能力.
也可以促進我們.
心理的健康.
讓我們容易感到.
有幸福感.
和我們可以激勵.
利他的行為.
即是行善.
是神所喜悅的.
\newpage



\section{腓立比書 1:3-6}
\label{sec:032310}
\textbf{多代同行•與同輩同行:追求共同目標}~~~~唐郭麗婷牧師
\newline
\newline
連結: \href{http://www.church.com.hk/acms/content.asp?site=cdc\&op=show\&type=product\&code=032310\&layout=sermon}{\texttt{032310}} ~~~~ 語音日期: 2026-06-14
\newline
\newline
\hyperref[sec:032309]{\small{< < < PREV SERMON < < <}}
~
\hyperref[sec:index_chronic]{\small{[返順時目]}}
~
\hyperref[sec:032311]{\small{> > > NEXT SERMON > > >}}
\newline
\newline
腓立比書 1:3-6
\newline
\begin{longtable}{cl}
\hline
\hline
章節 & 經文 (和合本修訂版)\\
\hline
1:3 & \begin{tabularx}{0.7\textwidth}{X} 我每逢想念你們，就感謝我的神， \end{tabularx} \\ \\ \relax
1:4 & \begin{tabularx}{0.7\textwidth}{X} 每逢為你們眾人祈求的時候，總是歡歡喜喜地祈求， \end{tabularx} \\ \\ \relax
1:5 & \begin{tabularx}{0.7\textwidth}{X} 因為從第一天直到如今，你們都同心合意興旺福音。 \end{tabularx} \\ \\ \relax
1:6 & \begin{tabularx}{0.7\textwidth}{X} 我深信，那在你們心裡動了美好工作的，到了耶穌基督的日子必完成這工作。 \end{tabularx} \\ \\
[1ex]
\hline
\hline
\end{longtable}
$^{1}$好,我們說教會需要有一個共同一齊的目標.
我們今天在同福堂這個教會.
大家知不知道同福堂的目標是什麼?.
大家記不記得上個星期.
謝牧師不知道,大家有沒有在這裡.
如果在這裡,不知道你回過堂有沒有說.
就是我們同福堂的名是什麼意思?.
「同心合意,興旺福佑」,所以「同福」.
其實這個「同心合意,興旺福佑」.
不是因為我們要叫「同福堂」.
所以我們就做這個口號出來.
不是的,其實這句話是出自《肥臘比書》.
保羅是對肥臘比教會的信徒,一個勉勵.
當我們看到華語寫在聖經裡面的時候.
一定就不只是保羅神指定說給當時的肥臘比教會.
而是說給眾教會.
今天我們所有的教會.
其實神都盼望我們能夠「同心合意,興旺福佑」.
好,我們一起來讀一下這幾節的聖經.
就是肥臘比書一章三至六節,我們一起讀.
每逢為你們眾人祈求的時候.
嘗試歡歡喜喜地祈求.
因為從頭一天直到如今.
你們是同心合意的,興旺福佑.
我深信那在你們心裡動了善功的.
必成全者功,直到那遭忌毒的日子.
這裡就是保羅和肥臘比教會信徒一開始的信.
哇,就很溫暖了,是不是?.
他說:我想起你們就很開心,很感恩.
我想起你們祈禱的時候都很歡喜.
其實肥臘比教會就是保羅在歐洲建立的第一間教會.
他寫信給他們說:我很開心.
因為從第一天我去到肥臘比建立這間教會.
直到今天,我們看資料,其實已經過了十二年.
這個時候他已經因為福音的緣故被抓去坐牢.
所以他在監牢裡.
但是他說:今天雖然我不在你們當中.
但是你們仍然同心合意,興旺福佑.
他說到同心合意這個事的時候.
如果我們看回這個聖經的原文.

$^{41}$就用了Cornelia這個字.
這個字是什麼意思呢?.
就是逆作一個團契,相交,神學上也有共同分享.
一起參與這個意思.
所以保羅說:除了因為我坐牢.
你們為我祈禱,供應我這個經濟的需要之外.
我更加開心的是.
我去建立這個教會就是想將福音傳給你們.
傳開去,今天你們能夠跟我一起去參與,共同分享這個使命.
所以他覺得很開心.
然後稱讚他們,你們是同心合意.
由我在這裡的那天到今天.
你們都是跟我一起同心合意,興旺福佑.
這個就是這樣的意思.
所以我們說,他們不當其實傳福音.
保羅你來建立教會,是你的使命,你傳給我.
我接受了就接受了.
但是記住,他們是能夠承擔教會的使命.
我想問大家,你清不清楚教會的使命是什麼?.
我們今天在這個教會.
你知不知道教會是什麼目的下建立的?.
是怎樣成立的?.
大家知不知道?.
你看回聖經就會發覺.
其實這個教會就是當天源於耶穌的一群門徒.
他們在地上,我們跟著耶穌.
耶穌在他們在地上的時候,對他們有一些吩咐.
我們一起來讀這些吩咐.
其實你也知道是很熟悉的.
如果你已經信奉了一段時間,你就知道這個大使命.
我們一起來讀,《馬太福音》28章19節.
1 2 3 「所以你們要去使萬民作我的門徒.
奉父子聖靈的名,給他們施洗」.
這裡說得很清楚,你們要去.
你們要出去,然後將我這個福音告訴他們.
然後讓他們成為我的門徒.
他們要知道這個福音.
然後到耶穌升天之前.
他又再一次有一個吩咐給他的門徒.
是說什麼呢?.

$^{81}$就是「使徒行傳一張白紙」.
這個大家可能都熟悉的,我們一起來讀.
「但聖靈降臨在你們身上.
並要在耶路撒冷猶太傳遞.
直到地極作我的見證」.
嘩,大家都知道.
這個是耶穌吩咐你們要去.
你們得著能力要等待聖靈.
然後你們要藉著聖靈的能力.
將這個福音傳到地極.
所以今天,你看看旁邊那個.
我們聽到福音.
是因為當日他們領受到上帝的吩咐.
神的吩咐,耶穌的吩咐.
然後他們將這個福音不斷傳.
傳到今天我們每個人都聽到了.
嘩,所以就是這樣.
所以教會就是在當耶穌升天之後.
這個五旬節聖靈降臨之後.
就開始,門徒就開始聚集.
一起敬拜,一起禱告.
然後他們有一個共同的目標就是.
我們要將這個耶穌吩咐的.
我們今天得救,我們這個最好的福音.
我們將這個福音傳開.
所以教會就是這樣成立.
所以弟兄姊妹,如果我們要擁抱.
同心合意,興旺福音這個使命的時候.
我們必定有一樣東西要很清楚.
就是我們要知道.
教會的目的.
建立的目的是傳福音.
因為今天我們在裡面做很多事情.
我們在這裡敬拜,我們會禱告,查經.
小組,你就要以為我們只是一個團體.
聚集一下人,做一下這些事情.
其實都一定要做這些事情.
這些都是我們要做的.
但是最終,你要記得.
我們教會最終的使命.

$^{121}$是傳福音.
如果教會不傳福音,你期望誰傳福音?.
你期望什麼團體傳福音?.
百貨公司?餐廳?.
是教會.
所以我們要認清.
今天我們教會,譬如我們說同福堂.
宗教會存在在地上.
我們的意義是什麼?.
不是只是聚集一班人.
來做一些裡面的事情.
而是神說,你們要去.
所以是一個內向的.
然後我們知道教會是傳福音.
所以我們大家要擁抱這個使命.
不是說,只有牧師的責任.
哦,這個教會領導站在台上,是他的責任.
不是的.
其實當今天我們是信徒的時候.
我們就領受這個這麼美好的福音.
今天我們承受了.
我們要告訴其他人.
你一起來得到這個這麼好的救恩.
這個這麼好的福音.
你要得到,是不是?.
都要跟別人說.
然後我們就一起向著這個目標去邁進.
我們盼望.
我們這樣去認清教會的目的的時候.
我們就能夠大家手挽手地.
向著這個方向去邁進.
所以教會其實不是為.
我們今天得到這個福音.
我們已經得到這個救恩.
所以教會不只是為我們自己而存在.
當然在裡面我們也做很多事情.
但是教會是為了這個世界而存在的.
是為了很多仍然未得救.
未聽過福音的人而存在的.
大家明不明白?.

$^{161}$是啊,所以不只是為了我們一群人.
當然我們裡面.
我們去敬拜,我們去禱告親近神.
我們來有信仰的教導.
讓大家知道聖經說什麼.
我們都是在說.
我們很盼望裝備我們.
讓我們更加認識神.
更加生命被神去改變.
能夠得著這個豐盛的生命.
然後我們能夠裝備好.
來去見證神.
所以福音能夠讓更多人得到.
去分享給他們.
其實我們也會關懷社區.
明年我們會關懷社區.
我們關懷社區的原因.
是因為神叫我們愛地上的人.
但是我們從來都不用關懷.
來做一個傳福音的手段.
即是說我關懷你.
其實是因為我想傳福音給你.
其實平時我不是很關心你.
即是說我想傳福音給你.
我請你吃飯.
其實不是的.
我們是真心去愛每一個人.
神說在小子上面.
最少一個就是在他身上.
所以我們是無條件去愛任何一個人.
但是你們要知道.
有時候我和一些社福機構合作的時候.
我們很欣賞他們.
即是關懷有需要的人.
但是我和他們說.
我說其實你可不可以讓我們傳福音.
因為有些他們會說你們不要傳福音.
你們不要傳教.
但是我開頭的時候.
我都想要和你合作.

$^{201}$好啊 我們都要關懷人.
但是到最後我都和他們說.
請你讓我們能夠講福音.
因為我真的和你有些分別.
我是教會.
我希望人得到幫助.
得到耶穌的愛.
但是我有一個終極的使命.
我希望人能夠.
不只是得到幫助.
而是能夠得救.
Amen.
我們是要人得救.
不只是停在關懷.
關懷是必須的.
是真心的.
但是你不可以不和他講這個福音.
因為唯有耶穌才能夠拯救人.
我們真的拯救不了.
我把我的財產都給他.
都未必會改善他的生活.
但是唯有耶穌讓人有盼望.
讓人在困難當中.
能夠得著醫治安慰.
走出困境.
Amen.
所以我們要很清晰.
很清晰這件事情.
否則我們只是照顧我們裡面的需要.
整天集中在我們裡面.
要怎樣怎樣.
照顧這些失去了福音的使命.
個個都不願意出來傳福音.
只是自己為內一班人.
其實我們真是一個.
失去了最重要的靈魂.
的一個宗教組織.
別人看著這個教會.
哦.
是不是.

$^{241}$完全不知道他們在做什麼.
可能是一班人.
他們喜歡這樣.
但是教會真是有一個使命.
好了.
我們說我們要一起擁抱這個使命.
實在都不容易的.
因為有人的地方就不同了.
所以其實同心合意.
一定我們會遇到挑戰.
挑戰是什麼呢.
最簡單最直接就是意見不一樣.
分歧.
然後因為大家意見不同了.
說著說著就.
你不妥我 我不妥我.
可能表面都妥妥的.
但是裡面就不太妥.
即是張力.
見到面阻阻.
都很辛苦.
其實是很正常.
有人的地方.
如果教會越大.
越難.
是不是.
要更加多人去同心.
是不是比較難.
但是.
難道我們說沒有嗎.
不是的.
有些人說.
我們沒有.
沒有不同意見的.
我都想請他們講一個講座.
因為.
我都很想學一下.
怎樣整個這麼大的群體.
會沒有不同的意見.
但是一般都會有.

$^{281}$其實強如這一間.
被保羅稱讚為同心合意的教會.
其實他在肥辣秘書.
都說到裡面的人都吵架.
都有紛爭.
他講到這間教會裡面.
特別提到有兩個女童工.
叫做呂雅蝶和秦都姬.
兩個吵架.
當然他在這裡提出來.
他亦都想勸勉.
教會是有這些分歧的.
不過你們記住要在主裡面.
記住同心.
你們記住要謙卑.
你們記住要有愛.
你們記住要放下自己的想法.
意見 立場.
你們要顧著 守著.
教會的使命.
當然.
如果是你.
你都可能會這樣勸勉別人.
你見到別人吵架.
你記住要放下自己.
大家聽聽意見.
但其實當你火遮眼的時候.
兩個人都是火遮眼的時候.
其實你這樣講.
都未必聽你的.
所以我就想與其這麼辛苦.
要經常勸勉.
不如我將這些勸勉.
做一點點的原則給大家.
來看看當我們遇到不同意見的時候.
我們要拿些什麼來做標準.
例如我和你的意見不一樣.
我們怎樣去解決呢.
用什麼來做標準呢.
因為我們在講.

$^{321}$我們是一個信主的群體.
福音是我們最大的使命.
所以福音是大前提.
然後人的意見.
不能夠大過神的心意.
我們是純屬於神的心意底下.
所以人你怎樣想都好.
你一定要想福音是大前提.
神怎樣想.
神最想我們做什麼.
你不可以不去做那件事.
你只能夠放下你自己.
如果你自己的意見和神一樣.
所以我們說個人的意見要讓步.
然後最終我們不要忘記.
我們不要做一些事情.
令我們離開了神整個教會.
對基督的使命.
因為有時候我們會忘記.
可能我是敬拜部的.
教會只敬拜就可以了.
全部都是敬拜.
永遠都是敬拜.
就不用傳福音了.
這個意見很強.
那些人都跟著我.
那就錯了.
一點福音都不傳.
所以不能夠離開基督的使命.
說到這幾件事.
我們就說.
其實做真的不容易.
不過在聖經裡面.
我們見到保羅雖然這樣讚佩勒比教會的信徒.
但其實他經常這樣想.
他有沒有試過和別人吵架.
你在聖經裡面有沒有見到保羅和別人吵架.
有的.
都說他和其他屬靈領袖吵架.
意見不合.

$^{361}$今天我想和大家講兩個例子.
看看他和當中參與了爭拗的一些領袖.
他們怎樣來用這個原則.
剛才我們說.
福音為大前提來解決問題.
第一個就是保羅和巴拿巴拉.
如果大家熟悉《律行傳》.
就知道他們一起去宣教.
第一次的時候他們是很好的.
第二次的時候.
巴拿巴說想繼續帶他的表弟馬可去.
保羅就說不行.
給大家一些背景.
在第一次宣教旅程.
馬可也有去 但中途就走了.
聖經沒有說什麼原因.
但可能估計他覺得辛苦.
或者怎樣就走了.
所以第二次.
巴拿巴又說要帶他去.
保羅就說不行.
現在我不是兒戲.
你看這個人這麼沒有紀律.
無緣無故就走了.
但巴拿巴就說.
不是啊 年輕人多給他一次機會.
就讓他再去一次.
好了 又問大家.
如果你在中途看到他們爭拗.
你會站在哪一邊.
你會幫誰.
幫保羅的舉手.
幫巴拿巴的舉手.
巴拿巴比較多.
其實不舉手的那些是怎樣呢.
不過我知道的 我有時也不舉.
其實我也不知道.
好了 其實如果你幫保羅.
你可能會這樣想.
我們現在去宣教.

$^{401}$我們現在不是去遊樂園.
其實是很重要的事.
我還要去建立教會傳福音.
一個這樣的人突然又來又走.
完全沒有紀律.
到時我怎樣搞他.
搞到我們整隊散了.
所以你看到.
他也有理由的.
真的有可能發生.
不過你見到他的出發點.
他視功為重.
我一定要有紀律.
搞好視功 建立教會.
他也對的.
如果你幫巴拿巴.
為什麼你會這樣想.
你可能會想.
年輕人給一次機會吧.
那保羅 你想想你年輕時怎樣.
你逼白教會 別說走掉了.
是不是.
如果耶穌不給你第二次機會.
你又今天.
是不是很敢說.
你可能這樣想.
那怎樣.
那你又對 你又對.
那怎樣.
那要不要再舉一次.
你會不會改變你的想法.
真的不知怎樣.
聖經沒有說誰對誰錯.
不過我們見到出發點不一樣.
巴拿巴是以人為重.
講恩典.
耀樹 恩典 寬容.
其實真的談不攏.
在我們現今的人可能就翻臉了.
那就不去了.

$^{441}$我們不要去了.
談不攏.
我們搞旅行.
不搞了 不去了.
整家人都不去了.
有沒有試過.
我聽過.
真的這樣翻臉.
結果就是要取消第二次宣教旅程.
可能是這樣 很糟糕.
但是保羅一位宣教士.
一位很認識神心意的人.
他說不可以讓任何東西攔阻福音的工作.
福音我們是教會.
不可以因為吵架而不傳福音.
不可以因為意見不一樣而不傳福音.
所以最後他的決定就是.
他不堅持不讓馬可去.
然後他想了一個方法.
馬可 好啊.
巴拿巴 我尊重你的意見.
給一次機會.
所以最後他讓巴拿巴帶著馬可.
走海路去傳福音.
然後他就和西拉走陸路去傳福音.
所以其實他這個決定.
願意放下自己.
以福音為大前提.
其實我相信神很喜悅.
他沒有再爭吵下去.
堅持 然後得到神極大的祝福.
你看到那兩隊宣教隊.
一分為二之後.
福音去到不同的地方.
不僅沒有被影響停止.
更加是拓展了.
我們看到如果我們做事.
適合神的心意.
其實可以有分歧意見.
但當我們解決的方法.

$^{481}$是適合神的心意的時候.
你見到是吃一個祝福.
變成一個極大的祝福.
其實我們見到最後.
保羅和巴拿巴.
其實他們是冰釋前嫌.
不是因為這樣而吵架.
因為保羅之後在歌羅西書.
和提摩太后書.
我們見到他稱呼馬可為同宮.
又說馬可做的福音工作.
對他有益處.
所以你見到最終他們的關係.
沒有因為之前的意見分歧而有影響.
之後還有馬可的福音.
嘩 你見 其實真是太美麗了.
第二個例子就說.
保羅和彼得.
彼得其實是保羅的前輩.
當聖經說.
聖靈降臨.
有幾千人信著保羅的講道.
彼得的講道後來才有保羅出來.
所以彼得是保羅的前輩.
但這個事情就發生.
有一次彼得和一群外邦人在吃飯.
吃飯之前外邦人覺得他們不潔淨.
我們猶太人怎會和外邦人一起吃飯.
但當時耶路撒冷會議.
教會已經協議.
外邦人不需要受國禮.
他們不需要遵守所有猶太人的律法.
只要守幾個戒條.
他們就可以和猶太人一起吃飯.
他們就得救了.
他們是恩順清儀.
不需要守所有的律法.
雖然這樣 彼得當然知道.
他是教會的柱石.
但當他和外邦人吃飯時.

$^{521}$他看到有猶太人進來.
他突然覺得.
怕這群猶太人.
不知怎樣想我和外邦人一起吃飯.
他就靜悄悄跑到猶太人那裡.
坐在猶太人的桌子上.
他是領袖.
其他巴布巴那巴猶太人看到他.
就跟著他走走走.
全部走回猶太人那邊.
好了.
保羅看到之後.
非常不高興.
向外邦人說.
我仍然覺得你不潔淨.
你又怕那群猶太人.
說你為什麼和外邦人一起吃飯.
你就走過去.
保羅說.
你這樣不行.
你是教會的柱石.
這個真理關乎外邦人是否能夠得救.
福音傳不傳給他們.
為什麼你會這樣做.
多面罵他.
如果是給你會怎樣.
被人這樣散.
對著那麼多人.
後輩.
都很難.
但是.
長話短說.
是時間關係.
這兩個巨頭相爭.
最後彼得接納了保羅的責備.
因為他想起了什麼.
對 為了福音的緣故.
其實我真的不應該這樣做.
但可能他那一刻想.
我坐過去.

$^{561}$因為他們不明白.
我怕他們不明白.
但其實都是為了自己的面子.
所以他接受了這個責備.
他體諒到.
其實這個責備是來自保羅對福音的堅持.
對外邦人能夠得救的堅持.
所以之後.
在彼得後書都見到他稱呼保羅是.
親愛的弟兄.
所以你見到.
我覺得我很想學習.
我們記在聖經裡面.
一些屬靈.
所謂的偉人.
他們的氣量.
他們不記仇.
他們真的以福音為大前提.
他們就可以繼續彼此相愛.
所以這兩個經歷.
保羅的經歷.
讓我們看到.
其實在一個群體裡面.
包括教會.
其實分歧是可以存在的.
是不是.
一定是有分歧的.
張力可以出現.
但是我們怎樣來到.
應該怎樣走呢.
你永遠走在福音的方向.
福音是大前提.
是不是.
然後有一樣東西很重要.
就是.
福音是大前提.
雖然我們意見分歧.
但是最後.
弟兄姊妹我們都要.
彼此相愛.

$^{601}$你記住.
那個人.
就算弟兄姊妹跟你意見多不一樣.
你要記住.
他不是你的敵人.
他不是你的仇人.
他是跟你一起向同一個方向的信徒.
我們只有一個敵人.
是誰.
撒旦.
撒旦偷竊毀壞.
要破壞我們的關係.
我們只有一個敵人.
所以記住.
到最終記住.
要和好.
要彼此相愛.
這就是神對我們的教導.
當我們來到講道.
我們能夠同心合意的時候.
為我們這個使命.
帶來什麼祝福呢 教會.
我們就會說.
我們不需要時常掙扎在彼此相拗.
你覺得吵架累不累.
累不累.
我很累.
吵架令我們只集中在吵架.
因為要贏.
要贏多辛苦.
要想很多原因來說別人輸.
很累.
然後心情不好.
整個人很沮喪.
整天陷在一個很不開心的情況底下.
沒有全福音的力量.
所以當我們願意同心.
我們以福音為大前提的時候.
我們記住要彼此相愛.
我們就少了很多內耗.

$^{641}$你知不知道什麼叫內耗.
一個家庭都可以很內耗.
一個教會都可以很內耗.
你不能夠將所有精力集中在應該要做的事情.
而是做了一些不值得的事.
就是內耗.
神說我們是福音的軍隊.
我真的很盼望.
我們能夠同心合意.
希望福音將教會的注意力.
心力.
真的回到全福音.
我們是有活力有動力的.
我們不是時常被挨打的一隊軍隊.
我盼望是這樣.
但是不是靠一個人 不是靠鄭牧師 不是靠你.
是靠我們整個教會.
同心合意.
好不好.
其實有個前輩跟我說過一句話.
很久了 我都忘記是什麼時候.
就是同心不是一個情緒的融合.
而是一個使命的對焦.
我們不是要求所有人很開心.
你明白嗎.
而是我們在很多艱難不同的想法裡面.
我們是對焦.
一起去對焦.
這就是真正的同心.
我正在想.
最近大家都有聽到.
我們下年2027年.
有一個叫一臨香江的福音聚會.
我記得鄭牧師都提過幾次.
其實這個報道大會.
我都聽到一些資料.
我去了兩個聚會.
第二個聚會就是這個聚會的啟動禮.
在5月28日.
我都去聽聽這個聚會究竟是怎樣.

$^{681}$因為我們教會都會有份去做一些服事.
去參加鄭牧師會.
後來就知道其中一個講完了.
這個聚會其實現在都沒有場.
是打算在政府大球場舉行.
我去之前 去之後都聽到不同的聲音.
有人就問.
其實在現在這個世代.
我們還需要一個這麼大型的報道會嗎.
有沒有聽過.
你有沒有聽過別人在說.
有些人就說不要再搞這些了.
個人報道更好.
做一些小型聚會更好.
你怎樣想法呢.
你有什麼想法呢 對於這個聚會.
你的想法是什麼呢.
我自己就思考一下.
我就想 其實對於我來說.
哪個為之最有效的傳福音方法呢.
是沒有一個定論的.
因為我想起我當年.
我都有去國培理報道大會.
我有很多聚會.
最終都是別人來探訪我 我信得住.
其實不同人接收福音的渠道是不一樣的.
所以對我來說.
我覺得討論的方向不是在形式上.
而是這個聚會所想帶來的就是.
我們很想全香港的人都聽到福音.
你想不想啊.
哇 我真的很想啊.
我們得救 我當然想所有人都得救.
每個人都聽到福音啦.
所以我覺得討論的方向.
對我來說不是在形式上.
因為香港有多少人呢.
750萬人 我昨天再查過一次.
有多少人信了耶穌.
20萬?.

$^{721}$太多了吧.
還因為社會事件.
過去幾年就跌了幾萬人.
我心想 還有700多萬人還沒信耶穌.
那這個大球場你坐進去不就20多萬這麼多場.
還有那些進不去的又怎樣.
所以我覺得.
如果我們需要擁抱的.
不是在那裡搞.
我們要擁抱的是.
你想不想全香港的人都聽到福音.
這才是我們的目標.
這才是我們的方向.
如果我們真的想一起擁抱這個目標的時候.
其實每個人.
你就可以用你覺得合適的方法去傳福音.
是不是啊.
那你說我不答應去大球場.
我覺得不適合他.
那你就用適合他.
你覺得請他吃飯請20次.
哇 你就請20次囉.
你說請他看電影.
下次有個報道會.
或者你自己個人報道.
其實真的可以的.
因為實在太多人還沒得救.
所以我覺得討論的方向.
其實對於我來說形式不重要.
而是我們自己有沒有做好我們自己.
每一個人同心合意.
興旺福音.
你有沒有被神用你自己是一個流通的管子.
將福音傳給其他人.
讓他能夠得救.
這個才是重點是不是啊.
各適其適.
可以就好像保羅在哥林多前書裡說.
無論如何怎麼樣都好.
你說什麼都好做什麼都好.

$^{761}$總要盡力去救一些人.
你跟隔壁的說.
我們總要盡力去救一些人.
是啊 我們真的盼望.
個個都得救是不是啊.
人人都得救.
這是上帝的心意.
其實到了啟動禮5月28號.
我最大的感動.
我是見到一班前輩.
其實我90年代去參加國賠禮大會.
我還沒信主.
但到今天我看到他們仍然堅持.
福音幾十年.
從年輕到今天年老.
我覺得對於我們這些後輩來說.
不要說別人搞什麼.
我有沒有這樣的傳福音的心智.
是不是啊.
我真的很感動.
我真的覺得我真的需要學習.
我看到那天很多人來.
跨教會 跨宗派.
所有的神學院 跨機構.
哇 這麼多人一齊上來傳福音.
我覺得很感動.
然後忽然間我覺得.
我很想參與.
因為我希望在我有生之年.
真的能夠真真正正經歷一次.
同心合意 興旺福音.
是香港的信徒.
Amen.
你想不想啊.
我真的很想.
所以我相信神會閱納這個心.
然後就好像《肥納比書》一章六節這樣說.
我們一起讀.
我深信那在你們心裡動了善功的.
必成全者功.

$^{801}$直到基督耶穌的日子.
叫人能夠信耶穌的.
不是我們說的.
好像表面上是我們向人傳福音.
讓人聽到這個福音.
但他能不能夠去信主的時候.
其實是神在他心裡動了善功.
是聖靈感動了他的心.
讓他能夠接受這個救恩.
接受這個福音.
所以我覺得神叫我們的角色就是.
我們去傳.
當他看到我們同心來傳福音的時候.
他必完成者功.
一個報道會有沒有果效.
不是說你在哪裡搞 搞什麼.
而是那班人是不是真的有心.
同心去傳福音.
所以我們來到神的面前.
我真的很想去經歷.
我相信當面對一些事情.
我們又有不同的意見.
好像這次福音聚會也是這樣的時候.
我們會想一下在福音的大前提下.
我們會怎樣去回應.
我相信這個神給我們一個第二次機會.
來學習什麼叫做不同而和.
Amen.
好了 其實到最後.
在這個聚會裡面.
我有一個經歷.
他有一本祈禱手冊.
他叫我們你們安心隨便揭.
揭一頁 你們就跟著那頁來自己禱告.
我就揭了那些頁.
就寫著這些東西.
我覺得 啊 神真的叫我去這個聚會.
也是有目的的.
他跟我說了這些.
就是那本祈禱手冊的訊息.

$^{841}$我們一起來讀.
籌辦全港性報道會前.
不是場地也不是人數.
是宗教會要有合而為一的心.
在彼此相愛中向著標桿直跑.
在彼此的不同中.
讓我們有共同的目標.
向著傳福音這個標桿.
一起去直跑.
Amen.
我們真的需要.
我們真的需要校正我們的目光.
以至我們能夠同心同行.
我們是教會.
我們要傳福音.
我們要人得救.
Amen.
好 我們最後起來.
我們讀一句聖經.
然後我們就唱回剛才的詩歌.
結束我們今天的崇拜.
我們一起讀羅馬書一章十六節.
我不以福音為恥.
這福音本是神的大能.
要救一切相信的人.
\newpage



\section{列王記上 8:1-29}
\label{sec:032311}
\textbf{耶和華的榮耀充滿聖殿:神同在的高峰,人心順服的考驗}~~~~袁惠鈞牧師博士
\newline
\newline
連結: \href{http://www.church.com.hk/acms/content.asp?site=cdc\&op=show\&type=product\&code=032311\&layout=sermon}{\texttt{032311}} ~~~~ 語音日期: 2026-06-14
\newline
\newline
\hyperref[sec:032310]{\small{< < < PREV SERMON < < <}}
~
\hyperref[sec:index_chronic]{\small{[返順時目]}}
~
\hyperref[sec:032312]{\small{> > > NEXT SERMON > > >}}
\newline
\newline
列王記上 8:1-29
\newline
\begin{longtable}{cl}
\hline
\hline
章節 & 經文 (和合本修訂版)\\
\hline
8:1 & \begin{tabularx}{0.7\textwidth}{X} 那時，所羅門召集以色列的長老、各支派的領袖和以色列人的族長到耶路撒冷，所羅門王那裡，要把耶和華的約櫃從大衛城，就是錫安，接上來。 \end{tabularx} \\ \\ \relax
8:2 & \begin{tabularx}{0.7\textwidth}{X} 以他念月，就是七月，在節期時，所有的以色列人都聚集到所羅門王那裡。 \end{tabularx} \\ \\ \relax
8:3 & \begin{tabularx}{0.7\textwidth}{X} 以色列眾長老一來到，祭司就抬起約櫃。 \end{tabularx} \\ \\ \relax
8:4 & \begin{tabularx}{0.7\textwidth}{X} 祭司和利未人將耶和華的約櫃請上來，又把會幕和會幕一切的聖器皿都帶上來。 \end{tabularx} \\ \\ \relax
8:5 & \begin{tabularx}{0.7\textwidth}{X} 所羅門王和聚集到他那裡的以色列全會眾一同在約櫃前獻牛羊為祭，多得不可勝數，無法計算。 \end{tabularx} \\ \\ \relax
8:6 & \begin{tabularx}{0.7\textwidth}{X} 祭司將耶和華的約櫃請進內殿，就是至聖所，安置在兩個基路伯的翅膀底下約櫃的地方。 \end{tabularx} \\ \\ \relax
8:7 & \begin{tabularx}{0.7\textwidth}{X} 基路伯張開翅膀在約櫃上面的地方，從上面遮住約櫃和抬櫃的槓。 \end{tabularx} \\ \\ \relax
8:8 & \begin{tabularx}{0.7\textwidth}{X} 這槓很長，從內殿前的聖所可以看見槓頭，從外面卻看不見。這槓直到今日還在那裡。 \end{tabularx} \\ \\ \relax
8:9 & \begin{tabularx}{0.7\textwidth}{X} 約櫃裡沒有別的，只有兩塊石版，就是以色列人出埃及地，耶和華與他們立約的時候，摩西在何烈山放在那裡的。 \end{tabularx} \\ \\ \relax
8:10 & \begin{tabularx}{0.7\textwidth}{X} 祭司從聖所出來的時候，有雲充滿耶和華的殿， \end{tabularx} \\ \\ \relax
8:11 & \begin{tabularx}{0.7\textwidth}{X} 祭司因雲彩的緣故不能站立供職，因為耶和華的榮光充滿了耶和華的殿。 \end{tabularx} \\ \\ \relax
8:12 & \begin{tabularx}{0.7\textwidth}{X} 那時，所羅門說：「耶和華曾說要住在幽暗之處。 \end{tabularx} \\ \\ \relax
8:13 & \begin{tabularx}{0.7\textwidth}{X} 我的確為你建了一座雄偉的殿宇，作為你永遠居住的地方。」 \end{tabularx} \\ \\ \relax
8:14 & \begin{tabularx}{0.7\textwidth}{X} 王轉過臉來為以色列全會眾祝福，以色列全會眾都站立。 \end{tabularx} \\ \\ \relax
8:15 & \begin{tabularx}{0.7\textwidth}{X} 所羅門說：「耶和華－以色列的神是應當稱頌的！因他親口向我父大衛應許的，也親手成就了；他曾說： \end{tabularx} \\ \\ \relax
8:16 & \begin{tabularx}{0.7\textwidth}{X} 『自從那日我領我百姓以色列出埃及以來，我未曾在以色列各支派中選擇一城，在那裡為我的名建造殿宇，但我揀選大衛治理我的百姓以色列。』 \end{tabularx} \\ \\ \relax
8:17 & \begin{tabularx}{0.7\textwidth}{X} 我父大衛的心意是要為耶和華－以色列神的名建殿。 \end{tabularx} \\ \\ \relax
8:18 & \begin{tabularx}{0.7\textwidth}{X} 耶和華卻對我父大衛說：『你有心為我的名建殿，這心意是好的； \end{tabularx} \\ \\ \relax
8:19 & \begin{tabularx}{0.7\textwidth}{X} 但你不可建殿，惟有你親生的兒子才可為我的名建殿。』 \end{tabularx} \\ \\ \relax
8:20 & \begin{tabularx}{0.7\textwidth}{X} 現在耶和華實現了他所應許的話，使我接續我父大衛坐以色列的王位，正如耶和華所說的，我也為耶和華－以色列神的名建造了這殿。 \end{tabularx} \\ \\ \relax
8:21 & \begin{tabularx}{0.7\textwidth}{X} 我也在那裡為約櫃預備一處。約櫃那裡有耶和華的約，就是他領我們列祖出埃及地的時候，與他們所立的約。」 \end{tabularx} \\ \\ \relax
8:22 & \begin{tabularx}{0.7\textwidth}{X} 所羅門當著以色列全會眾，站在耶和華的壇前，向天舉手， \end{tabularx} \\ \\ \relax
8:23 & \begin{tabularx}{0.7\textwidth}{X} 說：「耶和華－以色列的神啊，天上地下沒有神明可與你相比！你向那些盡心行在你面前的僕人守約施慈愛， \end{tabularx} \\ \\ \relax
8:24 & \begin{tabularx}{0.7\textwidth}{X} 這約是你向你僕人大衛守的，是你應許他的。你親口應許，親手成就，正如今日一樣。 \end{tabularx} \\ \\ \relax
8:25 & \begin{tabularx}{0.7\textwidth}{X} 耶和華－以色列的神啊，你向你僕人我父大衛應許說：『你的子孫若謹慎自己的行為，在我面前行事像你所行的一樣，就不斷有人在我面前坐以色列的王位。』現在求你信守這話。 \end{tabularx} \\ \\ \relax
8:26 & \begin{tabularx}{0.7\textwidth}{X} 以色列的神啊，現在求你成就向你僕人我父大衛所應許的話。 \end{tabularx} \\ \\ \relax
8:27 & \begin{tabularx}{0.7\textwidth}{X} 「神果真住在地上嗎？看哪，天和天上的天尚且不足容納你，何況我所建的這殿呢？ \end{tabularx} \\ \\ \relax
8:28 & \begin{tabularx}{0.7\textwidth}{X} 惟求耶和華－我的神垂顧僕人的禱告祈求，俯聽僕人今日在你面前的祈禱呼求。 \end{tabularx} \\ \\ \relax
8:29 & \begin{tabularx}{0.7\textwidth}{X} 願你的眼目晝夜看顧這殿，就是你說要作為你名的居所；求你垂聽禱告，你僕人向此處的禱告。 \end{tabularx} \\ \\
[1ex]
\hline
\hline
\end{longtable}
$^{1}$列王記合一王朝系列的第七講.
我今天要講列王記上的第八章.
一共有66節的經文.
所以我沒可能能夠講完.
能夠講多少我今天就講多少.
下一次會繼續完成列王記上的第八章.
今天的題目叫做.
耶和華的榮耀充滿聖殿 神同在的高峰 人心馴服的考驗.
為何有這樣的題目?.
因為今天要講所羅門奉獻聖殿.
耶和華的榮光臨到聖殿.
耶和華與以色列人同在.
但很可惜 因為以色列人沒有馴服.
最後聖殿被毀 耶路撒冷被毀 人民被擄去.
所以雖然我們的神願意與我們同在.
但很多時候我們都需要馴服.
今天結束的時候 我們會回來看這一點.
其實所羅門奉獻聖殿.
在聖殿禱告祈求.
是在列王記當中非常重要的經文.
而我時時都說.
列王記的作者有很多神學的思想.
在第六章 當所羅門起聖殿的時候.
作者就將起聖殿的事和出埃及記連在一起.
讓我們看到不單止知道.
究竟以色列人是何時出埃及.
亦都是透過埃及讓我們看到以色列人的罪.
埃及就是代表罪的捆綁.
神說你們不要回去這條路.
但很可惜 以色列的君王就不斷走回這條路.
和埃及立約,起文約.
向埃及招兵買馬.
所以最後是怎樣?.
就好像回到埃及.
人民被擄再一次在這裡捆綁的當中.
就好像回到埃及一樣.
現在所羅門奉獻聖殿.
作者又將這件事和大衛之約連在一起.
讓我們看到雖然神應許大衛之約.
這個王朝是會永遠的.

$^{41}$但就因為以色列人的反叛.
最後都是聖殿被毀.
作者描述所羅門的一生.
有很多的伏筆我們已經說了.
今次作者描述獻聖殿的事.
亦都有很多的伏筆.
其實列王記作者的文筆非常精彩.
在第八章就有一個非常精彩的首尾結構.
我沒有時間一點一點地跟大家分析首尾結構.
但最重要的部分.
就是在首尾結構的最中心.
一共所羅門在聖殿裡祈求耶和華七件事.
有六件是具體的.
三件在前面是列王記上的八章31至36節.
三件在後面是列王記上的八章41至51節.
這六件事都是具體的要求.
求神赦免以色列人的罪.
但經文最中心的地方.
就是說當以色列人不服從神降災難.
這些災難都是一般的災難.
不是那麼具體的 是一般的.
這些經文就在列王記上的八章37至40節.
所以由以所羅門向神祈求的七件事.
成為第八章最中心的經文.
我會把這個PowerPoint留在網上.
如果大家有興趣的話可以下載來看.
今天的大綱非常簡單.
一共是三件事.
第一件事是講到聖殿是神榮耀的殿.
在列王記上的八章1至11節.
第二件事是聖殿是神同在的殿.
列王記上的八章12至21節.
第三件事是聖殿是向神祈求的殿.
就是在列王記上的八章22至53節.
我以前也講過.
我講很多系列都是希望和大家一起讀經.
所以當你聽這篇道的時候.
你就要讀列王記上的第八章.
但不要忘記.
今天的經文也在歷代志下的五章至六章.

$^{81}$甚至第七章.
所以也應該讀這些經文.
我首先講第一點.
就是這個聖殿是神榮耀的殿.
所羅門在耶路撒冷.
他招聚了所有十二支派的領袖.
和以色列人從北面到南面的居民.
來和他一起奉獻聖殿.
所以聖殿這個字在經文裡非常重要.
在今天的經文列王記上第八章出現26次.
但在歷代志下的第五至第七章.
一共出現37次.
無論是大衛.
無論是所羅門.
都用了最好的人力物力來建造聖殿.
現在所羅門就奉獻聖殿.
聖殿亦是代表神和以色列人同在.
這是禱告的殿.
所以今天最中心的經文.
就是所羅門向神祈求的七件事.
其實當以色列人起聖殿奉獻聖殿的時候.
他們已經變得和外方人和他們的鄰居很相似.
記不記得在《說撒姆》以上的第八章.
他們要求有一位王.
就是因為他們的鄰居有一位王.
他們亦都希望有一位王.
他們的鄰居亦都有很多的廟宇.
來敬拜他們的偶像.
以色列人亦都希望有一座聖殿.
能夠敬拜他們的神.
我上次都講過.
其實有很多學者都認為.
以色列的聖殿的結構是抄外方人的廟宇.
因為他們說在這裡.
以色列聖殿之前已經有很多這些廟宇存在.
而這些廟宇都有一個聖所.
而裡面有一個致聖所.
是最聖潔的地方.
但是我們要知道.
其實以色列的聖殿是按照甚麼來設計.

$^{121}$是按照回望.
在所羅門起聖殿之前的480年.
回望已經存在.
所以是比這些所有外方人的廟宇還要早.
所以不可能以色列的聖殿抄外方人.
這個絕對是神榮耀的殿.
所羅門就將約瑰從耶路撒冷的帳篷.
搬到聖殿裡面.
我們剛才說.
這些異教國家都有他們的廟宇.
都有祭壇,祭司,祭品.
就好像我們剛才所唱的詩歌.
那些人是向石頭來敬拜.
這些偶像都是石頭造的.
都是假的.
只有唯一的真神又真又活的耶和華.
就在摩利亞山的聖殿裡面.
這個就是原因.
所羅門是將約瑰搬到聖殿裡面.
他也是將回望裡面的陳設裝備.
從機片搬到耶路撒冷的聖殿裡面.
他將約瑰放在至聖所裡面.
兩個基魯伯天使的翅膀下.
這兩個基魯伯的天使.
不是在約瑰上高的兩個天使.
原來在至聖所裡面.
有兩個巨大的基魯伯天使的像.
每個像是15呎高.
兩個天使每個展開的翼.
由一個翼尖去到另外一個翼尖.
也是15呎.
而至聖所的闊度就是30呎.
所以兩個基魯伯的天使.
可以遮蓋整個至聖所最後面的牆.
兩個天使的翼尖在至聖所的房間中彼此相接.
而在另外一邊.
一個天使翅膀的尖是碰著南邊的牆.
一個是碰著北邊的牆.
兩個基魯伯的天使都是面向著萬子.
就是東邊耶路撒冷聖殿的入口.

$^{161}$當然在約瑰的上高.
也有兩個基魯伯的天使.
但是這兩個基魯伯的天使彼此相望.
在聖殿裡面有很多天使的雕刻.
而我也剛剛說了有關基魯伯天使.
為什麼有這麼多天使?.
我們說基魯伯天使就是守護者.
所以所羅門將這麼多基魯伯的天使雕刻在牆上.
也在至聖所裡面.
其實聖經告訴我們.
天使是不能被救贖.
在啟示錄裡面告訴我們.
有三分之一的天使墮落了.
但是這些天使都不能被救贖.
只有人類能夠被救贖.
所以天使他們都很好奇.
他們都觀察神恩典的奧秘.
就好像彼得前書的一章十二節所說.
他們也透過以色列人的事奉.
希望能夠了解這個恩典是怎樣的.
當聶王紀的作者寫聶王紀的時候.
很明顯聖殿仍然存在.
耶路撒冷還未淪陷.
因為他見到聖殿裡面的情形.
他見到什麼?.
這個作者不是大祭司.
他不能夠進入至聖所.
但他見到在這個萬寺有兩點.
這兩點就是抬藥櫃的金棍穿過了藥櫃.
這兩條棍的尾部是碰著萬寺.
相信在外面在聖所是能夠見得到.
你見不到那條棍.
但你見到這兩條棍的尾部是碰著萬寺.
你就知道藥櫃在至聖所裡面.
如果你打開萬寺的話.
你就會見到神的寶座向著你.
也就是說那兩條金棍是穿過藥櫃.
短的那兩邊 窄的那兩邊.
而不是穿過闊的那兩邊.
有很多在網上有關藥櫃的圖畫都是錯的.

$^{201}$他們將這條金棍穿過闊的那一邊.
如果穿過闊的那一邊就出現問題了.
因為如果你打開萬寺的時候.
而和華的寶座是向著哪一邊.
是向著南邊或是北邊.
而不是向著東邊的入口.
所以這個是很重要.
我們也知道在藥櫃裡面有些甚麼.
是有三樣東西本來.
是有裝著嗎哪的罐.
是有裝著亞倫開花的像.
和當然是有藥板或法板.
就是十戒.
其實無論是嗎哪或是亞倫開花的像.
都是紀念以色列人的反叛.
嗎哪為甚麼會有呢.
因為以色列人在《民數記》的第十一章.
他們是埋怨.
現在來到曠野出了埃及.
甚麼都沒得吃.
在埃及多和味.
又有蒜頭又有洋蔥.
你不要說甚麼.
你煮甚麼都好.
你找個鍋爆香了洋蔥和蒜頭.
你煮甚麼都好吃.
現在來到曠野.
真的甚麼都沒得吃.
不斷地埋怨.
於是神就降下嗎哪給他們.
為甚麼會有亞倫開花的像.
又是以色列人的領袖.
在《民數記》第十七章十二支派的領袖.
走來找亞倫.
為甚麼你做領袖.
我們都可以做.
神和摩西說好啊.
你叫他們每人拿他們的像.
在上面寫了他們的名字.
擺在會幕裡面.

$^{241}$藥櫃的前面.
你日後早上來看看.
只有亞倫的像是開了花.
其他領袖的像都沒開花.
所以又是把這支像擺在這個藥櫃裡面.
現在這個藥櫃已經去到聖殿神立祂名的地方.
不再需要提這些事.
是一個新的開始.
所以經文在列王記上的八章九節告訴我們.
現在藥櫃裡面只剩下八板.
神仍然是要以色列人來遵守祂的律例殿章.
但是嗎哪的冠和開花的像已經沒有了.
當藥櫃來到聖殿至聖所.
神的榮耀降臨在聖殿.
藥櫃就是代表神與以色列人同在.
在聖列王記上的八章十至十一節.
經文說祭司從聖所出來的時候.
有雲充滿耶和華的殿.
甚至祭司不能站立供職.
因為耶和華的榮光充滿了殿.
有沒有看到?.
有經文說雲充滿了殿.
亦有經文說神的榮光充滿了殿.
雲就是代表神的榮光.
記不記得昔日在這裡曠野.
神是以雲柱和煙柱來帶領以色列人.
現在神的榮耀在聖殿與以色列人同在.
於是祭司就用歌聲和樂器來讚美神.
就好像在詩篇二十二篇第三節.
耶和華住在他子民的讚美中.
這也是原因為何我們每個星期進來.
我們就用詩歌來讚美神.
在聖經裡不斷告訴我們.
神的榮耀來來去去.
在《著列王記》上第八章.
作者又一次將經文和出埃及記連在一起.
提到雲的榮耀.
就是提醒我們昔日以色列人出埃及之後.
神以雲柱煙柱的榮耀來帶領以色列人.
直至到神叫摩西起了回眸.

$^{281}$在《出埃及記》裡.
神的榮光就淋到回眸裡.
但也因為以色列人的反叛.
記不記得以色列人拿著神的藥櫃.
去和比利時人爭戰.
連藥櫃也被擄.
所以在《撒母耳上》第四章告訴我們.
神的榮耀就離開了回眸.
現在神的榮耀又回到聖殿裡.
但很可惜不會長久.
也會因為以色列人的罪.
神的榮耀會再一次離開聖殿.
就是在《以西結書》裡.
第八章,九章和第十章裡的預言.
神的榮耀離開聖殿.
而以塞瓦先知.
他預言神的榮耀將會再一次臨到人間.
就是在《以塞瓦書》第四十章.
是何時?.
是耶穌基督道成肉身來到世上.
就是《約翰福音》第一章.
第十四至二十三節所說.
神的榮耀再一次臨到人間.
但很可惜.
以色列人將榮耀的主.
釘死在十字架的上方.
而當耶穌基督復活升天的時候.
就是架著榮耀的文采回到天上.
而相對來說.
神的聖殿就變得荒涼.
你看看現在耶路撒冷的聖殿.
所在的地方.
就是有一個金色頂的清真寺.
神的榮耀絕對不在那裡離開了.
神給以西結先知有一個意象.
能夠見到神的榮耀.
將會又一次臨到耶路撒冷的聖殿.
這個就在以西結書的第四十一章至四十三章.
也是千禧年的時候.
耶穌基督將會帶領以色列人進入千禧年國.

$^{321}$神的榮耀再一次在聖殿出現.
所以這個是很重要.
沒有了千禧年國.
神的應許就會落空.
現在神的榮耀在哪裡?.
在使徒行傳的第二章告訴我們.
當那些門徒在敬拜神的時候.
聖靈降臨在他們的身上.
現在神的榮耀就在每一個兒女的身上.
也都在地方教會和普世的教會那裡.
我們都知道每次當我們用心靈誠實來敬拜神.
我們就將榮耀歸給神.
每一次我們是透過我們的事奉.
來榮耀神的時候.
神的榮耀也都被彰顯.
但我們都要小心.
我們不只是星期日或者是事奉的時候.
才將榮耀歸給神.
而是我們每一天的生命當中.
我們都要活出那美好的見證來榮耀神.
在將藥櫃帶到至聖所的儀式當中.
以色列人是獻祭了無數的祭物.
經文告訴我們.
根本是沒有辦法可以計算得到有多少祭物.
讓我們見到神的榮耀有多大.
讓我們見到以色列人是多麼尊敬神.
而經文也都告訴我們.
所羅門起好了聖殿之後.
他等了十一個月才將藥櫃搬去至聖所.
他起好聖殿的時候應該是八月.
而他將藥櫃搬去至聖所.
所以就是七月.
七月是什麼時候?.
正是鑄盤節的時候.
這個就令我們想起.
《申命記》的十六章第十五節.
經文說:你要在耶和華你神所揀選的地點.
向耶和華你的神守節.
以色列人起好了聖殿.
第一個節期就是鑄盤節.

$^{361}$鑄盤節代表什麼?.
神與以色列人同在.
啟示錄.
神要鑄在人間.
神的帳盤在人間.
就是與人同在的意思.
現在以色列人就能夠在耶路撒冷的聖殿來敬拜神.
因為神曾經允許以色列人.
可以在神納祂名的地方來敬拜神.
《申命記》第十二章.
但我們如果注意希伯來文的原文的話.
經文是用Mahon這個字.
來形容天上神的居所.
這個字出現在列王記上的八章三十九,四十三和四十九節.
為什麼用這個字?.
Mahon這個字其實是出於另外一個希伯來文的字.
就是Kun這個字.
Kun這個字通常在經文形容到.
大魏王朝被鞏固.
大魏的王朝被建立.
就是Kun這個字.
亦都是說到所羅門王位被鞏固.
例如在列王記上的第二章.
我們說所羅門用自己的方法來鞏固自己的王位.
就是Kun這個字.
為什麼作者要用這兩個字?.
同出一個字根.
就是讓我們看到有這個連帶的關係.
首先是大魏的王位,所羅門的王位.
而一直連帶到神天上的王位.
我時時都說.
千萬不要說大魏王位的應許.
是在天上才是應驗.
因為天上的王位不是大魏的王位.
所以我們看到這裡帶出了一個訊息.
是由大魏和所羅門的王位一直帶到天上.
但中間還有一個步驟.
就是耶穌基督會回來帶領以色列人進入千禧年國.
然後進入永遠的國度.
這個就是列王記上八章十三節告訴我們.

$^{401}$我們說完了神的殿是神榮耀的殿.
我們現在要說這個亦都是神同在的殿.
剛才我說這個作者不斷地將.
起聖殿,奉獻聖殿和出埃及的故事連在一起.
讓我們看到以色列人的歷史.
神就是透過以色列人的歷史.
彰顯出祂自己的歷史.
有人說History這個字是大草H.
History是誰的Story?是誰的歷史?.
是神的歷史.
所以這個字是很有意思.
我亦時時說在系統神學有不足的地方.
就是沒有了有關以色列人的系統神學.
聖經裡面三分之二的經文.
都是有關以色列人居然系統神學沒有了這一範.
出現問題.
神就是透過以色列人的歷史.
彰顯出祂自己的榮耀.
神要以色列人透過禮儀.
透過節日來紀念自己.
例如是安息日.
現在建了聖殿.
神又叫以色列人在聖殿獻祭.
來敬拜.
為自己贖罪.
都是要紀念耶和華.
記得耶和華.
所以在《述申命記》裡面.
「記得」和「記住」這個字最少出現十四次.
但是在《述申命記》第六章都有另外一個字是很重要.
就是「Shema」這個字.
要聽見經文在《述申命記》第六章第三節.
說以色列雷耀廷.
所以神不單止要以色列人記得祂.
記住祂.
亦都要以色列人聽祂的話語.
行出祂的話語.
當所羅門獻聖殿的時候.
他就在同台的上高.
他面對著聖所.

$^{441}$祭司在祭壇的面前.
百姓朝聚在會場當中.
所有人都見到神的榮耀彰顯出來.
而所羅門就開始他的演講.
他一開始講什麼呢?.
耶和華陳說祂必住在幽暗之處.
列王記上的八章十二節.
說大件事神是光.
在《若翰一書》的一章五節說神是光.
在《提摩太前書》的六章十六節.
說神住在光中.
為什麼現在所羅門突然說神住在幽暗裡面.
剛剛奉獻聖殿.
所有人在場都見到神的榮耀.
哪裡來幽暗呢?.
其實這就是作者文筆精彩的地方.
為什麼幽暗呢?.
又一次把獻聖殿和出埃及連在一起.
記不記得以色列人出埃及在西乃山.
神頒十戒給摩西.
他怎麼說呢?.
耶和華對摩西說.
我要在默魂中臨到你那裡.
叫百姓在我與你說話的時候.
可以聽見也可以永遠順利了.
出埃及記的十九章第九節.
注意這裡兩個字.
在默魂中和百姓可以聽見.
這個就是讓我們見到神的奧秘.
在幽暗的裡面.
在默魂的裡面.
當時在西乃山.
神頒發律法的時候.
有雷軍有閃電.
是有默魂.
摩西是戰戰兢兢地上了西乃山.
以色列的人民都非常害怕.
不敢觸摸這座山.
任何人碰到這座山都會沒命.
人不能夠見到神.

$^{481}$在出埃及記的三十三章二十節說.
你見到神就不能存活.
其實經文的重點就是聽見神的聲音.
神是不以形象來和以色列人會面.
因為神不要以色列人製造任何的偶像.
為自己雕刻任何的偶像來敬拜神.
所以神是奧秘.
其實我們都應該好像以色列人一樣.
要成為神話語的子民.
聽從神的話語來服從神.
知道這位神是興邁不得.
是住在幽暗當中的奧秘的神.
為何奧秘那麼重要?.
為何神在幽暗當中那麼重要?.
讓我們見到神的奧秘.
我們就會懂得謙卑.
我們不能夠完全明白神的道路.
我們不能夠完全甚至明白神所有的旨意.
透過神的奧秘就會鼓勵我們去相信祂.
去服從祂.
而不是靠自己的智慧.
希望去了解 希望去明白.
希望用自己的想法來完成任何的事.
而是完全依靠神 聽從神的話語.
很可惜 現在有很多人希望做甚麼?.
成為神的好朋友.
Buddy Buddy.
當然神是愛我們.
但神是興邁不得的.
就好像哈巴谷書第二章第20節提醒我們.
主在他聖殿裡.
願全地在他面前肅靜.
我們都很需要在神的面前肅靜.
我們都很需要來敬畏這位奧秘的神.
看著列王記上的八章十二節.
我剛才說神是住在幽暗之處.
住在這個字就是原文希伯來文的Chahan這個字.
是一個很特別的字.
在列王記上只出現了兩次.
為甚麼作者要用一個這麼特別的字來說.

$^{521}$神住在幽暗之中.
其實有另外一個希伯來文是更加普遍的住在.
或者甚至坐在皇位上這個字.
就是Yashgab這個字.
在列王記裡一共出現91次.
為甚麼作者不用這個普遍的字.
而要用這麼獨特的Chahan這個字.
原來Chahan這個字和另外一個希伯來文.
Mishkan這個字是同出一個字根.
Mishkan就是帳棚 帳幕.
剛才我們說以色列人守甚麼節.
是守鑄盤節.
鑄盤節就是代表神與人的同在.
所以我相信作者用了Chahan這個字.
雖然神住在幽暗當中.
但他仍然與以色列人同在.
就是透過Mishkan這個字 帳棚.
在《啟示錄》的二十一章第三節.
也有告訴我們.
神的帳幕在人間.
就是神要與我們永遠同在.
所以我們看到這兩個字的連帶關係.
聶皇紀的作者是否文筆非常精彩呢.
絕對是.
我們說完神與以色列人同在.
他是奧秘.
第二點他也是信實的神.
作者又一次將所羅門獻聖殿.
和大衛之約連在一起.
提醒會眾大衛之約.
就是神允許大衛有一個兒子.
能夠為耶和華起聖殿.
這就是所羅門.
現在所羅門在聶皇紀上的八章十五節說.
神親口所說的.
他就能夠親手成就了.
這就讓我們看到神的信實.
神向大衛允許.
所羅門一定坐在他的王位上.
這樣就成就了.

$^{561}$而且所羅門真的為神.
不僅起聖殿.
也奉獻聖殿.
你說神這樣做.
是否要將榮耀歸給大衛.
歸給所羅門呢.
絕對不是.
神是要榮耀自己的名.
所以在所羅門的講話和禱告裡面.
至少十四次提到神的名.
就讓我們看到這個重要性.
所羅門小心翼翼地將榮耀歸給神.
而在耶路撒冷.
每一個以色列人來到聖殿敬拜.
他們都會記得神的恩典.
神的信實.
神講出來的話就算數.
所以他們去到聖殿祈求神.
希望神能夠祝福他們.
注意.
經文的文法在聶皇紀上的八章十三節至二十一節.
突然間文法由第三者轉到第一者.
所羅門說我為你起聖殿.
這裡就讓我們看到所羅門的驕傲.
所羅門的驕傲也會引致他將來.
跟隨他的妃嬪去拜偶像.
這個就是所羅門的問題.
我們講完兩點了.
我說我今天要講三點.
第一這個聖殿是神榮耀的殿.
第二這個聖殿是神同在的殿.
第三點就是一個向神祈求的殿.
所羅門貴在台上.
在祭壇附近.
雙手舉向天開始禱告.
而天或天上這個字.
在聶皇紀上的八章二十二至五十四節.
一共出現不下於十次.
所羅門舉起兩隻手來禱告.
這和我們現在向神禱告的姿勢很不同.

$^{601}$我們怎樣禱告?.
合上雙手 閉上眼睛.
但以色列人舉起兩隻手.
仰望天來祈求.
你說為什麼?.
舉起兩隻手就是代表自己的空虛.
代表自己的貧乏.
神啊 你快些答應我.
你快些應允我禱告.
就是很期待地向神禱告.
其實這個姿勢在聖經裡.
在早期的教會也是很普遍.
在提摩太前書的二章八節補羅說.
當那些人禱告的時候.
舉起聖傑的手來向神禱告.
在這裡所羅門又提到.
他說這些全部都是神的恩典.
以至我能夠為你來起聖殿.
他說是你守了你的約.
在世上沒有神能夠比起你.
聶王紀上第八章第十三節.
而的確神是為那些盡心.
行在祂面前事奉祂的人.
是會思慈愛.
大衛就是一個最好的例子.
當大衛成為合臣心意的人.
他是願意來事奉神的時候.
神就實驗了他的應許.
所羅門所以在這裡.
也提出了另外一句神的應許.
他說如果大衛的子孫.
謹慎自己的行為.
在神面前行事像大衛所行的一樣.
就不斷有人會坐以色列的國位上.
在撒母耳夏第七章.
神和大衛立約的時候也說過這句話.
就是說這個王朝是會永遠的.
但是如果以色列的王不聽話.
背弱的話.
你分分鐘不能夠坐在這個王位上.

$^{641}$現在大衛的王位在地上的王位.
誰坐在上面?.
沒有人了.
大衛的王朝已經完了在地上.
而將來耶穌基督會回來.
帶領以色列人進入千禧年國.
這個就是實驗了神的應許.
如果沒有了這個千禧年國的話.
神的應許就落空.
我也說過在天上的王位.
不是大衛的王位.
我們是要搞清楚.
所以我們這裡看到有兩個部分.
一個部分就是神的應許是無條件的.
直到永遠.
但是每一個的王都要為自己的行為負責.
我也說過這個好像我們的信心一樣.
我們的信心是不是因為我們的行為得救?.
不是.
是神的恩典算我們為二人.
就好像亞伯拉罕之約是無條件.
但是每一個信徒都要小心自己的行為.
因為如果我們犯罪的話.
神是會管教.
這個不是亞伯拉罕的約.
這個是摩西之約.
而保羅說430年後摩西之約.
絕對不能夠廢去430年前亞伯拉罕之約.
所羅門的禱告裡面.
也說到自己所建造的殿是有不足的地方.
就是說到天上的神.
全能的神.
怎麼可能住在人手所做的建築物裡面.
其實當所羅門邀請推羅王希蘭.
來幫他建聖殿的時候.
他已經說過這句話.
在《歷代志下》的二章第六節.
而在《以賽亞書》的66章一至二節也說.
耶和說:如此說.
天是我的座位.

$^{681}$地是我的腳墩.
你們要為我做何等的奠宇.
那裡是我安息的地方呢.
耶和說這一切都是我手所做的.
你不要以為是所羅門為神來建聖殿.
神怎麼說這一切都是我手所做.
是祂實現了祂自己的應許.
此提反在猶太人公會要將祂處決之前.
也引用了所羅門和以賽亞先知的話.
在《使徒行傳》的第十七章.
而當保羅向外邦人傳福音的時候.
他也強調了這一點.
在《使徒行傳》的第十七章.
我們見到這一切都是出於神的恩典.
我們說是祂准許以色列人.
在那裡納祂名的地方來敬拜.
就是在《申命記》的第十二章第五節.
耶和說:你們的神從你們各支派中.
選擇何處為納祂名的居所.
你們就當往那裡去求問.
這節的經文你有沒有看到.
將兩件事放在一起.
神的聖殿和向神求問.
神的聖殿就是萬國禱告的殿.
所羅門就是拿著這個應許.
來向神祈求.
當我們的子民犯了罪.
但是當他們向聖殿祈求.
所羅門說求你垂聽他們.
求你看顧這個聖殿.
求你垂聽他們的禱告.
在《少羽列王記》上的第八章第五十六節.
經文說:願耶和說:.
我們的神晝夜垂念.
每日為祂的僕人與他民以色列.
這就是申冤.
所羅門知道有一件事就是.
《新命紀》第28章和第29章.
是神與以色列人立約的條款.
當以色列人走約的時候.

$^{721}$神會祝福.
但是當以色列人毀約.
神就會詛咒.
神甚至會差派敵人來攻擊他們.
也容許災難臨到他們身上.
這就是神的管教.
但是無論如何.
在《新命紀》第30章.
神還允許以色列人的子民.
就算你們背了約.
我是差派你的敵人來攻擊你.
我是有災難臨到你們的身上.
但是只要你們悔改.
只要你們迴轉.
就怎樣?.
神一定會赦免.
一定會復興.
這個也在《述若拿先知》的身上.
我們會見到.
若拿仰望神的聖殿來禱告.
神就赦免了他.
在《述若拿》書的第二章第四節.
但以利先知.
也是遙望耶路撒冷的聖殿.
因為他當時已經不再在耶路撒冷.
他甚至遙望耶路撒冷的聖殿.
為百姓來祈求.
但以利先知的第六章第十節.
我剛才也說過.
無論新約舊約都告訴我們.
神的聖殿就是萬國禱告的殿.
所以所羅門在第八章最中心的地方.
說出了有七個祈求.
希望神能夠答允.
這就是當以色列人毀了約.
神來到責罰他們.
管教他們的時候.
所羅門說求你赦免他們.
其實這是在第八章最中心的經文.
我剛才也說過.

$^{761}$三件是具體的祈求.
在首尾結構的上方.
三件亦是具體的祈求.
在首尾結構下面的部分.
而最中心的經文.
就是一些普通天然的災難.
我們來看看首尾結構.
這個首尾結構一共有四層.
第一層是說到.
當以色列人得罪他們的鄰居.
列王記上的八章三十一至三十二節.
或者是以色列人得罪了神.
有沒有看到?.
列王記上的八章四十六至五十一節.
無論是得罪鄰居.
無論是得罪神.
另外第二層.
當他們敗在敵人的手上.
列王記上的八章三十三至三十四節.
就是有敵人來攻擊他們.
他們敗在敵人的手上.
而列王記上的八章四十四至四十五節.
當他們去攻打他們的敵人.
無論敵人來攻打他們.
或者他們去攻打敵人.
他們打敗仗.
因為他們的罪.
求神寬恕他們.
還有第三層.
就是當你的民得罪你.
這裡是以色列人.
列王記上的八章三十五至三十六節.
又或者不屬你民的外邦人.
當他們禱告的時候.
求你聽他們的禱告.
列王記上的八章四十一至四十三節.
有沒有看到這裡包括所有人.
以色列人或者不屬於以色列人.
他們向你禱告.
求你聽他們的禱告.

$^{801}$而經文最中心的.
就是一些一般的災難.
天然的災難.
包括有饑荒,瘟疫,寒風,梅瀾,蝗蟲,馬乍.
或者是遭遇任何其他的疾病.
或者災難所羅門說.
當我們面臨這些災難.
因為我們的罪的話.
求你赦免以色列人.
我今天都說我沒可能.
講完全部列王記上的第八章.
也都沒可能可以講完這七點.
我今天只可以講兩點大概.
第一點就是人若果得罪了鄰居.
我們知道所羅門剛剛登基為王的時候.
他是為兩個妓女來破案.
所羅門絕對不可能為所有的人民來破案.
所以他很需要有審判官來幫助他審案.
也都祭司可以透過律法來作出裁判.
這些都能夠幫到所羅門.
經文說如果有一個人被他的鄰居控告.
說他得罪了鄰居.
這個鄰居去到聖殿裡面.
指著壇來起誓.
說這個人是得罪了自己.
神就要宣佈這個人的罪.
得罪鄰居這個人的罪.
那麼你說神怎樣宣佈?.
怎樣講給這些祭司或者領導人聽?.
我相信是透過污名和土名.
就是祭司身上的污名土名.
神叫以色列人要維持公義.
才能夠安居樂業在這片土地居住.
但很可惜很多以色列人都犯了罪.
不單得罪鄰居也得罪了神.
他們始終讓不義進入了這片土地.
流了無辜人的血.
所以最後遭受神的管教被趕出這個應許之地.
講完人如果得罪鄰居還有第二呢?.
因為他們的罪.

$^{841}$如果他們敗在敵人面前的時候又會怎樣呢?.
我剛才講過.
如果以色列人守約.
他們能夠安居樂業.
平安居住在應許之地.
但如果他們不守約背約.
神就會甚至差派他們的敵人來攻擊他們.
他們會被他們的敵人擊敗.
如果以色列人打敗了之後.
他們願意悔改願意歸向神呢?.
神說就算你們被擄.
我都能夠使你們回到自己的土地.
這裡作者又有一個非常精彩的描述.
就是當以色列人知道自己犯罪.
他們願意悔改.
悔改這個字.
原文希伯來文是哪個字?.
是Shuf這個字.
我講過悔改Repent是怎樣?.
180度轉回.
就是希伯來文的Shuf.
當以色列人Shuf悔改的時候.
悔轉的時候.
神會怎樣呢?.
使他們能夠回歸到自己的土地.
而回歸這個字又是Shuf這個字.
有沒有看到作者文筆的精彩地方?.
這個列王記的信息.
也是被很多在這裡流亡的以色列人所讀到.
因為列王記這本成書之後.
以色列人已經是聖殿被毀.
他們被擄去巴比倫.
他們已經在流亡.
當他們讀到神的這個信息.
只要他們悔改.
轉向神.
他們Shuf.
神就會怎樣呢?.
就會讓他們回到自己的土地.
你說對他們是否一個很大的鼓勵呢?.

$^{881}$今天結束的時候.
我要和大家再回去看.
我今天的這個題目.
神同在的高峰.
人心馴服的考驗.
為什麼有這個題目呢?.
首先神怎樣與以色列人同在?.
是透過所羅門所建的聖殿.
神的榮耀臨到聖殿.
所以神是與以色列人同在.
那你說我們呢?.
耶穌基督渡乘肉身來到世上.
與我們同在.
還不僅如此.
還為我們釘十字架.
赦免了我們的罪.
每一個接受祂的人都赦免我們的罪.
我剛才說有很多審判官.
幫助所羅門來審案.
他們是怎樣審案的?.
是定罪人有罪.
定義人為義.
但是耶穌基督在十字架上為我們所成就的.
是否定罪人有罪?.
這就糟糕了.
如果所有罪人願意認罪悔改的話.
願意赦免的話.
神也是給我們有機會接受祂.
這就是耶穌基督的救贖.
不是定罪人有罪.
那你說那些義人又怎樣?.
羅馬書的三章二十三節.
告訴我們世上有沒有義人.
完全沒有義人.
所以每一個人都需要神的救贖.
我們現在看到神已經與我們同在.
我們應該要怎樣回應?.
我這個題目的下半部分就是人心馴服的考驗.
所羅門有沒有馴服?.
外表好像頗為馴服.

$^{921}$剛才其實經文作者已經有伏筆.
在列王記上的八章十三至二十一節.
突然間經文用到第一位.
「我為神起聖殿」.
有沒有看到所羅門的驕傲?.
有六件事是神所憎恨的.
加上神所最憎恨的一共七件.
哪一件排頭?.
高傲的眼睛.
箴言的第六章.
所羅門就是犯了驕傲.
以致他日後向十一個支派徵收重稅.
為什麼十一個支派不是有十二個支派?.
這正是問題.
他沒有向自己的猶大支派徵收重稅.
而是向其他十一個支派.
而引來十一個支派的不滿.
最後還被他的一千個妻妾帶領去拜偶像.
你被一個妻妾來引領你拜偶像已經很大件事.
你還有一千個妻妾.
你想想有多少偶像?.
還有一個問題.
不單止所羅門不夠謙卑他驕傲.
以色列人又怎樣?.
都忘記了神的恩典.
是神救以色列人出埃及.
與以色列人立約.
神亦起了聖殿與以色列人同在.
但是以色列人因為他們背若.
以致最後聖殿被毀.
這位作者在列王記下二十五章.
十三至十七節.
非常之簡短的一筆就給我們見到結果.
就是所羅門為耶和華所起的殿.
所造的兩根柱子.
海和可移動支架的青銅.
被打碎並運往巴比倫.
我在這裡說.
所羅門起聖殿的時候.
說了很多這些青銅和金造的家具.

$^{961}$那些器皿.
被巴比倫全部打破.
將這些物品都運往巴比倫.
就完結了.
所以我們也要很小心.
我們應該要怎樣?.
我們應該有好的回應.
這就像所羅門一樣.
當他起聖殿的時候.
他走去向神宣佈.
他說神啊.
我完成了你給我的使命.
摩西亦完成了起會幕的使命.
耶穌基督亦完成了.
救贖就是為人類釘十字架.
而施洗約翰和保羅.
亦都完成了神的任務.
有一天我們每一個人.
每一個真正相信的基督徒.
都會去到天家.
神就會問我們.
喂,Sidney.
我給你的任務你完成了沒有?.
不要去到當時才突然間很驚訝.
不懂怎樣回答.
其實我們已經在地上.
已經要上.
我們怎樣能夠真真正正.
為神服務.
完成我們的任務.
將榮耀歸給神.
願神祝福我們自己的話語.
讓我們低頭禱告.
讓我們見到所羅門.
不單是起聖殿.
亦都是奉獻聖殿.
讓我們見到我們不要學習所羅門.
在經文告訴我們.
他開始驕傲.
以至日後他帶領以色列人去拜偶像.

$^{1001}$我們不單止不要驕傲.
亦都要記得你的恩典.
就好像你帶領以色列人出埃及.
和他們立約.
起聖殿與他們同在.
我們都滿有你的恩典.
你差派你的獨身子來到世上.
來與我們同在.
為我們贖罪.
願意當我們將來去到天上與你見面的時候.
都能夠有美好的報告.
都能夠被你稱為一個忠心良善的僕人.
將榮耀頌讚歸給你.
我們禱告.
奉主耶穌基督的名而求.
阿們.
\newpage



\section{希伯來書 5:7-10}
\label{sec:032312}
\textbf{在基督裡效法基督 (華/粵)}~~~~戴永富博士
\newline
\newline
連結: \href{http://www.church.com.hk/acms/content.asp?site=cdc\&op=show\&type=product\&code=032312\&layout=sermon}{\texttt{032312}} ~~~~ 語音日期: 2026-06-14
\newline
\newline
\hyperref[sec:032311]{\small{< < < PREV SERMON < < <}}
~
\hyperref[sec:index_chronic]{\small{[返順時目]}}
~
\hyperref[sec:032313]{\small{> > > NEXT SERMON > > >}}
\newline
\newline
希伯來書 5:7-10
\newline
\begin{longtable}{cl}
\hline
\hline
章節 & 經文 (和合本修訂版)\\
\hline
5:7 & \begin{tabularx}{0.7\textwidth}{X} 基督在他肉身的日子，曾大聲哀哭，流淚禱告，懇求那能救他免死的神，就因他的虔誠蒙了應允。 \end{tabularx} \\ \\ \relax
5:8 & \begin{tabularx}{0.7\textwidth}{X} 他雖然為兒子，還是因所受的苦難學了順從。 \end{tabularx} \\ \\ \relax
5:9 & \begin{tabularx}{0.7\textwidth}{X} 既然他得以完全，就為凡順從他的人成了永遠得救的根源， \end{tabularx} \\ \\ \relax
5:10 & \begin{tabularx}{0.7\textwidth}{X} 並蒙神照著麥基洗德的體系宣稱他為大祭司。 \end{tabularx} \\ \\
[1ex]
\hline
\hline
\end{longtable}
$^{1}$我們先低頭禱告 禱告我們就不需要反應.
天父我們真感恩.
因為今天我們可以在你面前.
可以敬拜你 聆聽你的話語.
求你幫助我們 打開我們的心竅.
使我們可以明白你的話.
消化你的話 遵守你的話.
我們把你的兩位僕人都交託仰望.
求你與我們同在 阿們.
所以今天我們的題目就是在.
「基督離校法基督」.
(記者:今天我們的題目是「基督離校法基督」).
我們要先從比較大的角度看這個問題.
(記者:我們從大的角度來看這個問題).
我們就是圍繞上帝的創造和上帝的救贖.
(記者:圍繞著什麼呢?就是圍繞著上帝的創造和上帝的救贖).
上帝創造賜給我們什麼呢?賜給我們生命.
(記者:上帝創造賜給我們什麼呢?就是賜給我們生命).
如果我們看《聖經》裡頭.
我們就發現生命跟愛是息息相關的.
因為生命就是愛的表現.
(記者:我們看《聖經》就看到生命跟愛是息息相關的).
(記者:因為生命就是愛的表現).
而且特別有意思的是在《聖經》裡頭.
當上帝創造生命的時候.
上帝也是給人氣息 對嗎?.
(記者:當《聖經》裡說神創造人的時候).
(記者:也都給我們人有氣息的 對不對呢?).
所以當我們呼吸對生命很重要.
呼吸是有吸然後呼 對嗎?.
(記者:我們呼吸是很重要的).
(記者:在我們生命之中 我們是要呼和吸的).
所以你不能只是吸而不呼.
(記者:我們不可以只是吸而不呼出去的).
或者呼出來但是沒有氣也不行.
(記者:或者我們不可以只是呼 但我們不吸氣進去的).
所以這個呼吸是可以象徵什麼呢?可以代表什麼?.
(記者:這個呼吸是可以象徵什麼呢?).
(記者:吸就是等於得 領受 呼就是等於給 分享).
吸就可以得 領受 呼就可以好像分享一樣.

$^{41}$(記者:那我剛才說這個生命是愛的表現).
(記者:愛也是生命的表現).
剛剛我講到生命就是愛的表現.
而愛也是生命的表現.
(記者:那我們就是知道這個愛呢).
(記者:離不開給 離不開得 離不開受 離不開奉獻).
(記者:我們知道愛是離不開給 離不開受 離不開奉獻的).
(記者:那現在比方說有兩個人在教會結婚).
(記者:我就是不管是富貴貧窮身體健康疾病到死).
(記者:我都是輸於你的 什麼意思呢?).
(記者:兩個人在一個教會如果在這裡結婚的時候).
(記者:好像說我們無論貧窮或者富貴 或者身體健康).
(記者:都是雙手的 究竟是什麼意思呢?).
(記者:我們要注意 我這個是等於說我把自己給你).
(記者:但是很有意思 當我把自己給你的時候).
(記者:這種給也是領受的條件).
(記者:這裡要注意大家 我是將我自己給你的時候).
(記者:這種給也是一種領受的意思).
(記者:所以相愛就是相給 相互充滿 相互奉獻).
(記者:相愛就是相給和相互的分享).
(記者:所以你們可以了解為什麼耶穌說).
(記者:非我失去性命的會得到性命).
(記者:所以這裡就證明好像耶穌所說的).
(記者:非為我失去生命的就會得到生命).
(記者:這個好像跟我們的邏輯不一樣).
(記者:因為我們不斷地要激烈 是不是不斷地要得).
(記者:我們跟我們自己的邏輯好像不一樣的).
(記者:因為我們永遠都是想我們自己).
(記者:永遠都要得很多的東西).
(記者:因為通過倒空自己 我們才能夠領受).
(記者:因為通過倒空了自己).
(記者:我們才能夠真真正正領受).
(記者:所以愛相愛就是忽給忽得).
(記者:愛相愛就是忽給忽得).
(記者:那愛 愛也是一種聯合 對不對).
(記者:愛 愛也是一種聯合 對不對).
(記者:那我跟對方 我跟我的愛人 我跟我的孩子).
(記者:怎麼樣聯合呢).
(記者:我跟我的對方 我的愛人 我的細蚊仔).
(記者:怎麼樣聯合呢).

$^{81}$(記者:不僅是坐在一起 不僅是空間上).
(記者:我們是緊愛著).
(記者:不是 但是我們坐在一起).
(記者:或者在空間之中 我們是緊緊地).
(記者:是靠著大家).
(記者:靈魂之間的聯合就是忽給忽得).
(記者:靈魂之間的結合就是忽給和忽給).
(記者:我把自己給你).
(記者:我把我自己給了你).
(記者:但是當我給你的時候 我也從你那裡得到你).
(記者:當我自己給你的時候 我也從你那裡得到東西).
(記者:過去50年代 60年代有一首歌叫《不了情》).
(記者:大家有沒有知道這首歌).
(記者:50年代 60年代有一首歌叫《不了情》).
(記者:《不了情》).
(記者:你們有沒有聽過呢).
(記者:《不了情》).
(記者:所以她說忘不了 忘不了 忘不了雨中的散步).
(記者:她說忘不了 忘不了 忘不了在雨中的散步).
(記者:為什麼呢).
(記者:因為女朋友散步的時候下雨).
(記者:因為這個女朋友在散步的時候下雨).
(記者:男朋友在裡面 也忘記帶雨傘).
(記者:但是她已經在裡面).
(記者:男朋友在裡面 也忘記帶雨傘).
(記者:但是她已經在裡面).
(記者:所以男朋友就跑到外面).
(記者:寧可跟他的女朋友一起在雨中散步).
(記者:男朋友就跑出去).
(記者:寧願在外面跟他女朋友一起在雨中散步).
(記者:第二天男朋友感冒了).
(記者:第二天男朋友就感冒了).
(記者:感冒了 他很感動 女朋友很感動).
(記者:那女朋友就會很感動).
(記者:我怎樣都不會嫁給你).
(記者:為什麼我一定要嫁給你).
(記者:所以給跟得相應).
(記者:現在聖經說 上帝就是生命 上帝就是愛).
(記者:聖經告訴我們 上帝就是生命 上帝就是愛).
(記者:所以在上帝的生命裡).

$^{121}$(記者:有這種忽給忽得的關係).
(記者:所以在上帝的生命裡).
(記者:有呼吸和呼吸的愛).
(記者:但是我們現在先不講這三個議題).
(記者:我們現在先不講這三個議題).
(記者:所以我的意思是 上帝本身是相愛的關係).
(記者:我們在講的是 祂本身就是相愛的關係).
(記者:所以我們現在就開始講創造).
(記者:所以我在這裡講創造).
(記者:那愛 我剛才說 愛是聯合的力量).
(記者:剛剛我講到 愛就是那個聯合的力量).
(記者:那上帝創造 不是因為上帝需要我們).
(記者:就覺得好像如果沒有我們的話 上帝就感到很孤單).
(記者:上帝創造不是因為沒有我們 上帝就會覺得很孤單).
(記者:上帝創造是出於祂的自由).
(記者:上帝創造是因為出於祂自己的自由).
(記者:因為上帝想跟其他的存在 祂以外的存在聯合).
(記者:是上帝想與祂以外的東西聯合一起).
(記者:所以上帝的愛的表現是什麼呢).
(記者:上帝要創造一個離祂距離遠的存在).
(記者:所以上帝創造是一個離祂很遠的物質的創造).
(記者:那離祂距離更遠的是什麼東西呢).
(記者:離祂距離更遠究竟是什麼呢).
(記者:是物質).
(記者:就是物質了).
(記者:為什麼物質離祂遠).
(記者:為什麼物質會離神很遠呢).
(記者:上帝無所不在).
(記者:因為上帝是無所不在).
(記者:但是物質必須有固定的地方 固定的空間).
(記者:因為物質一定要有固定的空間 固定的地方).
(記者:上帝是永恆的).
(記者:但是物質呢 物質是短暫的 會容易分解).
(記者:我們是從很多原子元素組成的 容易分解).
(記者:物質是很短暫的 很容易分解).
(記者:我們是從很多的原子結合 所以很容易分解的).
(記者:所以上帝離物質距離遠).
(記者:所以上帝離物質距離是最遠的).
(記者:我們的身體是從塵埃塵土造成的).
(記者:我們的身體是從塵埃塵土所造成的).

$^{161}$(記者:這個就是因為我們是從物質來的).
(記者:因為我們是從物質而來的).
(記者:但是物質沒有生命).
(記者:但是物質是沒有生命的).
(記者:所以上帝就護了它的氣息).
(記者:所以上帝呼了它的空氣).
(記者:我們有靈魂了).
(記者:我們人就是能夠說 我愛你的一坨塵土).
(記者:所以上帝的創造就是上帝要收招一種聯合).
(記者:所以上帝的創造是要創造一種聯合的).
(記者:所以上帝要給).
(記者:所以上帝要給予).
(記者:給予什麼呢?給予亞當夏娃能力,智慧,聰明,感性,理性都給).
(記者:給予亞當夏娃能力,智慧,聰明,感性,理性都給).
(記者:上帝也是給予亞當夏娃永遠的生命).
(記者:所以如果沒有上帝的維持,沒有上帝的恩典).
(記者:我們人很快就死了).
(記者:所以如果沒有上帝的恩典,沒有上帝的維持).
(記者:我們人很快就會死亡的).
(記者:所以創造可以被理解為一種婚姻,一種聯姻).
(記者:造物主跟受造物).
(記者:所以上帝是靈,然後我們是有身體的靈魂).
(記者:上帝就是靈,我們就是有身體的靈).
(記者:聯合).
(記者:這個就是生命).
(記者:這個就是生命了).
(記者:所以聖經說永生就是認識你).
(記者:所以聖經說有生就是認識耶穌).
(記者:因為永生就是跟上帝聯合).
(記者:所以永生就是與上帝聯合的意思).
(記者:我們剛才已經說生命,愛,聯合).
(記者:我們所以說生命,愛,聯合).
(記者:那跟聯合需要什麼呢?).
(記者:聯合需要什麼呢?).
(記者:互給,對不對?).
(記者:是互相的給予,對不對?).
(記者:所以上帝把一切給亞當夏娃).
(記者:所以上帝將一切都給了亞當夏娃).
(記者:但是亞當夏娃也需要給啊).
(記者:但是亞當夏娃也需要給予的).

$^{201}$(記者:他們要給什麼呢?).
(記者:什麼最重要的呢?).
(記者:最重要的就是他們的自由).
(記者:最重要的是他們的自由).
(記者:所以他們把最好的給上帝).
(記者:所以他們將最好的給上帝).
(記者:這個就是順服的意思).
(記者:把自由給上帝).
(記者:把自己的自由給上帝).
(記者:但是你能夠把你的東西給上帝嗎?).
(記者:但是你想想).
(記者:你可以將你自己有的東西給上帝嗎?).
(記者:你的一切都是從上帝來的).
(記者:所以我受造物能夠把自己的東西給上帝).
(記者:這個神學上不太正確).
(記者:我受造之物可以將我自己的東西給上帝).
(記者:在神學裡好像不太合邏輯).
(記者:比方說有一個人欠我一百萬美金).
(記者:比方說).
(記者:比如來講有個人欠我一百萬美金).
(記者:當然我的錢不是這麼多).
(記者:但是比方說有人欠我一百萬美金).
(記者:這是一個例子).
(記者:所以有一天我可能忘記了).
(記者:不帶信用卡不帶現金).
(記者:有一天我忘記了).
(記者:有一天我忘記了沒帶自己的信用卡或現金).
(記者:我要買東西先借你十塊錢可不可以).
(記者:所以這個人說我就借給你十塊錢).
(記者:然後記住).
(記者:我說笑話一百萬你的錢都是我的).
(記者:那個人說可以我借你十塊錢).
(記者:但是你要記住).
(記者:但是我在想笑話).
(記者:一百萬我都給了你).
(記者:亞當夏娃要給上帝怎麼給).
(記者:所有的都是給上帝來的).
(記者:所以亞當和夏娃有什麼可以給神).
(記者:其實所有的東西都是從上帝而來的).
(記者:所以其實).

$^{241}$(記者:給是得的方式).
(記者:其實給就是得到的方式).
(記者:亞當夏娃給上帝自由).
(記者:亞當和夏娃給上帝他們的自由).
(記者:他領受上帝更大的祝福).
(記者:他會領受上帝更大的祝福).
(記者:這個就是聯合).
(記者:這個是創造).
(記者:這個也是創造).
(記者:但是亞當夏娃墮落了).
(記者:但是亞當和夏娃他們墮落了).
(記者:每個婚姻都需要磨合).
(記者:需要適應期 對嗎).
(記者:每個婚姻都需要磨合和適應的期 對不對).
(記者:大家要適應彼此的脾氣 彼此的習慣).
(記者:大家要適應大家的脾氣 大家的習慣).
(記者:然後呢).
(記者:我們需要跟對方讓步 不斷地給).
(記者:我們需要向對方的讓步 不斷地給予).
(記者:但是可惜的是什麼呢).
(記者:亞當夏娃跟上帝的聯合呢).
(記者:是以失敗告終的).
(記者:但是可惜的是亞當和夏娃).
(記者:與上帝的聯合是以失敗為告終的).
(記者:為何呢).
(記者:為什麼呢).
(記者:因為亞當夏娃不願意適應上帝).
(記者:因為亞當和夏娃不願意適應上帝).
(記者:他們只要得不要給).
(記者:他們只想得到 但是不想付出).
(記者:只要得不要給 就是跟生命分開).
(記者:只要得到而不要付出的話).
(記者:就好像跟生命分開了一樣).
(記者:就比方說你只想吸 不願意呼).
(記者:可以嗎).
(記者:你好像比例來講).
(記者:你只想吸入 但是不想付出).
(記者:這樣做行不行呢).
(記者:所以亞當跟夏娃跟上帝隔絕了).
(記者:痛苦死亡就出現了).

$^{281}$(記者:痛苦和死亡就出現了).
(記者:外體不斷毀壞).
(記者:外體不斷的毀壞).
(記者:所以人就沒有安全感).
(記者:所以人就開始悲觀 如果不是悲觀那是驕傲 自我欺騙).
(記者:如果不悲觀就變成驕傲 自我欺騙).
(記者:但是聖經的信息很偉大).
(記者:但是聖經裡面的信息是很偉大的).
(記者:這個是關係到第二部分 創造).
(記者:我們講到第二部分).
(記者:不是 救贖 救贖).
(記者:就是救贖了 第二部分).
(記者:人不願意適應上帝 上帝願意適應人).
(記者:人願意適應上帝 但是上帝願意適應人).
(記者:人要變成上帝 上帝要變成人).
(記者:人不願意適應上帝).
(記者:人願意適應上帝).
(記者:所以上帝成為人).
(記者:所以上帝就成為了人).
(記者:道成了肉身).
(記者:那耶穌復活前的身體是像我們的身體).
(記者:在耶穌復活之前的身體就和我們的身體一樣).
(記者:希伯來書告訴我們 他在一切事情上受試探 只是不犯罪).
(記者:所以他會捱餓 甚至會死).
(記者:他可以捱餓 甚至可以死).
(記者:他能夠難過 能夠流淚).
(記者:他可以感受到難過和流淚).
(記者:他經歷了人生的痛苦).
(記者:但是特別有意思的是 他和我們不一樣的是什麼呢?).
(記者:但最有意思的是 他和我們不一樣的是什麼呢?).
(記者:他信服到底).
(記者:耶穌基督他從上帝那裡吸了什麼呢?).
(記者:吸了他上帝的恩典 上帝的生命).
(記者:耶穌基督在上帝那裡吸了什麼呢?).
(記者:他吸了上帝的恩典).
(記者:所以他領受了上帝自己的能力 上帝的偉大).
(記者:他領受了上帝的能力 上帝的偉大).
(記者:所以他被釘在十字架的時候).
(記者:他把氣息呼出去).
(記者:神啊 接受我的靈魂 我給你我的靈魂).

$^{321}$(記者:神啊 求你接受我的靈 你創造的靈).
(記者:所以耶穌基督向天父說阿們).
(記者:我要成全你的旨意).
(記者:聖經說 你看 你預備了我的身體).
(記者:聖經告訴我們 你看 你預備了我的身體).
(記者:反忌和救贖忌不是你所喜歡的).
(記者:但是你給我預備身體 什麼意思呢?).
(記者:凡你所做的 你給我身體是什麼意思呢?).
(記者:如果我自己有翻譯一下).
(記者:我自己有翻譯一下 大體上是這樣子).
(記者:真正的反忌 真正的贖罪忌就是身體).
(記者:真正的贖罪忌 凡忌或者贖罪忌就是自己的身體).
(記者:所以耶穌基督把身體獻給天父).
(記者:所以耶穌基督把祂的身體獻給天父).
(記者:但是耶穌基督也跟我們說 是).
(記者:耶穌基督也跟我們說 是).
(記者:說阿們 阿們 阿們).
(記者:他說阿們).
(記者:因為祂願意成為我們的贖罪忌).
(記者:因為祂願意成為我們的贖罪祭).
(記者:所以耶穌基督給).
(記者:所以耶穌基督祂是吸出了).
(記者:那你看 當耶穌基督給的時候 發生了什麼).
(記者:看看當耶穌基督吸的時候 吸出的時候 發生了什麼事呢).
(記者:愛的力量就反擊).
(記者:愛的力量就反擊了).
(記者:顛倒了自然規律).
(記者:過去生而死).
(記者:過去是生然後死).
(記者:人生去古誰無死 對不對 過去生都是死的).
(記者:人生自古都是生之後就會死的).
(記者:現在死就是生).
(記者:現在就是死而復生).
(記者:過去多年前 我看過一個影片).
(記者:我覺得大家也看過 就是 Jurassic Park 知道嗎 這個影片).
(記者:很多年前我看過一套電影 叫做 侏羅柏公園).
(記者:大家也看過這套電影的).
(記者:所以有很多科學 有一些科學家 剋了恐龍).
(記者:然後他們想控制恐龍 給恐龍很多限制).
(記者:這個恐龍不會害人).

$^{361}$(記者:科學家就是 發掘了一些恐龍).
(記者:產生一些恐龍 他對這些恐龍有很多的限制).
(記者:以致他們不會傷害到人).
(記者:然後有一個科學家告訴這個老闆).
(記者:但是有一個科學家告訴這個老闆).
(記者:告訴 Jurassic Park 這個公園的老闆).
(記者:侏羅紀公園的老闆聽).
(記者:Life will find a way).
(記者:生命會找到辦法).
(記者:生命會找到方法).
(記者:你限制它沒用).
(記者:你限制它是沒用的).
(記者:我剛才說 生命就是愛的表現).
(記者:我剛剛講到 生命就是愛的表現).
(記者:愛會找到方法).
(記者:愛就是能夠改變的).
(記者:所以在死亡中 這個愛會反擊 耶穌復活).
(記者:失去生命會得到生命).
(記者:失去了生命會得到生命).
(記者:你把生命獻給上帝).
(記者:這個奉獻本身是領受的方式).
(記者:其實這個奉獻是本身領受的方式).
(記者:所以耶穌復活了).
(記者:耶穌勝利了).
(記者:所以這個死而復活告訴我們愛的奧秘).
(記者:這個死而復活告訴我們愛的奧秘).
(記者:你看在聖經裡 愛一方面).
(記者:聖經說愛是永不止息 愛是永遠長存).
(記者:在聖經裡告訴我們).
(記者:愛是永不止息 永不長存).
(記者:但是另一方面 聖經說).
(記者:沒有比這個更大的愛).
(記者:一個人把自己的性命犧牲給他的朋友).
(記者:但是聖經也告訴我們).
(記者:沒有其他愛比一個人).
(記者:將自己的生命獻出更加重要).
(記者:這個很矛盾).
(記者:一方面愛是永遠的).
(記者:一方面說愛是永遠的).
(記者:但是另一方面 愛是離不開死的).

$^{401}$(記者:另一方面就是愛你不應該死的).
(記者:這個矛盾怎麼解決呢?).
(記者:就是在耶穌基督的死而復活).
(記者:給才能夠得).
(記者:吸出去才能夠有所得著).
(記者:你把自己倒空給上帝).
(記者:你將你自己倒空給上帝).
(記者:上帝就充滿你).
(記者:上帝就會充滿你).
(記者:被上帝充滿).
(記者:上帝就是生命本身).
(記者:被上帝所充滿).
(記者:上帝就是生命的本身).
(記者:死亡不能夠殺生命).
(記者:生命能夠弄死死亡).
(記者:兩種死 弄死死亡).
(記者:但是生命能夠弄 整死死亡).
(記者:整死死亡).
(記者:所以這個是耶穌基督給我們的).
(記者:這個是耶穌基督給我們的).
(記者:那這樣子連最悲見的器皿).
(記者:最大的痛苦不能夠妨礙).
(記者:不能夠攔阻上帝的工作).
(記者:所以最悲的器皿).
(記者:最大的痛苦都不攔阻).
(記者:我們從上帝的愛).
(記者:其實我的孩子小的時候).
(記者:他們很喜歡看這個迪士尼漫畫).
(記者:迪士尼).
(記者:我小時候很喜歡看迪士尼那些戲的).
(記者:卡通片的).
(記者:那個英國的文豪CS Lewis).
(記者:他說過很多童話).
(記者:其實有呼音的種子).
(記者:他說CS Lewis提過).
(記者:其實細蚊仔的童話裡面).
(記者:都是有呼音的種子的).
(記者:所以過去有一個童話叫).
(記者:Beauty and the Beast).
(記者:美女與野獸).

$^{441}$(記者:大家知道嗎?).
(記者:有一個兒童的戲).
(記者:叫做美女與野獸).
(記者:大家都聽過).
(記者:所以這個野獸).
(記者:本來是英俊的王子).
(記者:因為他驕傲是).
(記者:遭到詛咒變成野獸了).
(記者:那他變成野獸).
(記者:只要有一個漂亮的姑娘愛上他).
(記者:他就能夠擺脫他的詛咒).
(記者:就回到過去的樣子).
(記者:他是變成了野獸).
(記者:只要有一個漂亮的姑娘愛上了他).
(記者:這個詛咒就會除滅了).
(記者:他就可以變回一個英俊的王子).
(記者:這個怎麼可能呢?).
(記者:猴年馬月呀).
(記者:什麼時候才能夠找到).
(記者:漂亮的姑娘愛上我?).
(記者:怎麼可能呢?).
(記者:猴年馬月才能找到).
(記者:一個美麗的女生喜歡上他呢?).
(記者:奇怪的).
(記者:有一個姑娘).
(記者:有一個漂亮的姑娘).
(記者:來到他的地方).
(記者:愛上了他).
(記者:很奇怪的).
(記者:有一個美麗的姑娘).
(記者:來到他那裡).
(記者:竟然愛上了他).
(記者:當他們要開始結婚之前).
(記者:當他們開始要在一起之前).
(記者:當他們要開始結婚).
(記者:或者在一起之前).
(記者:那個姑娘聽到).
(記者:她的爸爸有苦難).
(記者:需要回到自己的家).
(記者:這個野獸說放開他).

$^{481}$(記者:那個野獸就說放了他).
(記者:你去拯救你的爸爸).
(記者:你去拯救你爸爸).
(記者:不要管我).
(記者:不要理我).
(記者:他學到愛).
(記者:他學到了愛).
(記者:剛剛呢?).
(記者:他給).
(記者:他吸出去).
(記者:當他給以後他就死了).
(記者:但是他給了之後).
(記者:他自己就死亡了).
(記者:他復活了).
(記者:但是在這個童話裡).
(記者:他就得到復活).
(記者:為什麼?).
(記者:愛反擊).
(記者:因為愛反擊).
(記者:所以這個愛的力量).
(記者:所以上帝通過耶穌基督).
(記者:適應我們).
(記者:所以上帝通過耶穌基督).
(記者:來適應我們).
(記者:亞當不願意適應上帝).
(記者:上帝願意適應人).
(記者:亞當不願意適應上帝).
(記者:但是上帝願意適應我們人類).
(記者:所以我們在基督裡).
(記者:我們要通過適應基督).
(記者:效法基督適應上帝).
(記者:我們所以在基督裡面).
(記者:能夠適應上帝).
(記者:是因為我們效法基督的).
(記者:所以什麼叫效法基督?).
(記者:什麼叫效法基督呢?).
(記者:效法基督就是適應上帝).
(記者:適應上帝是什麼意思呢?).
(記者:變得像上帝一樣).
(記者:所以魔鬼告訴亞當).

$^{521}$(記者:你吃了那個果子).
(記者:你就變成像上帝一樣).
(記者:這個其實是上帝要給的).
(記者:魔鬼就跟亞當說).
(記者:你吃了這個果子之後).
(記者:你就像上帝一樣).
(記者:其實這個是上帝).
(記者:已經給他的東西來的).
(記者:但是不像撒旦的方式).
(記者:但是他沒有防到撒旦的方法).
(記者:聖經告訴我們).
(記者:上帝是由此愛的).
(記者:上帝告訴我們).
(記者:上帝是愛的).
(記者:這個告訴我們).
(記者:上帝是不斷的給).
(記者:不斷的分享).
(記者:不斷的給).
(記者:這個告訴我們).
(記者:上帝是不斷的愛).
(記者:不斷的分享).
(記者:亞當犯罪了).
(記者:他給他們做衣服).
(記者:亞當犯了罪).
(記者:他還為亞當做衣服).
(記者:以色列人背叛上帝).
(記者:上帝還釋放他們).
(記者:以色列人背叛了上帝).
(記者:繼續釋放他們).
(記者:耶穌的門徒離開他們).
然後他彷彿聽到主的聲音.
他彷彿聽到神的聲音.
那你要給我你的心.
那你也要給我你的心.
我要最大的禮物.
我最大的禮物就是上帝的心.
那上帝告訴他.
那你要給我最大的禮物.
那上帝告訴他.
那你也要給我最大的禮物.

$^{561}$最大的禮物是什麼.
最大的禮物是什麼呢.
我剛才說最大的禮物就是我們的自由.
最大的禮物我剛剛說過了.
就是我們的自由.
也就是我們的心.
也就是代表我們的心.
所以最大的禮物需要最大的代價.
為什麼.
為什麼呢.
因為領受上帝的心需要我們的心.
所以領受上帝的心需要我們自己的心.
能夠感受基督的感受.
能夠感受到基督的感受.
以基督的心為心.
能夠思考基督的思想.
通過什麼.
通過什麼呢.
通過心.
通過我們的心.
如果我要給你幾百年的葡萄酒.
如果我要給你一個幾百年醞釀過的葡萄酒.
你準備杯子.
你要準備一個杯.
空的杯子.
一個空的杯子.
心就是這個空的杯子.
心就好像這個空的杯子一樣.
最大的給.
最大的給予.
得到最大的禮物.
所以得到最大.
之後就得到最大的禮物.
所以我們能夠感受基督的感受.
我們要感受基督的感受.
我們過去不知道這個基督是什麼樣子.
我們過去不知道基督究竟是怎麼樣的.
現在我們大體上隱隱約約能夠感受到.
現在我們大體上隱隱約約都可以感受得到了.
當我們在痛苦的時候說.

$^{601}$不是跟著我的意思.
是跟著你的意思.
當我們在痛苦的時候.
我們能夠說.
不是跟著我自己的意思.
是跟著你的意思.
所以我們要給.
所以我們要吸出.
怎麼樣給呢.
我們要把我們的生命.
我們的心.
我們的生命給基督.
我們要將我們的心.
我們的生命給基督.
就好像五餅二魚.
好像五餅二魚.
你向上帝禱告.
你去向上帝禱告.
上帝啊.
我的生命.
我的心有很多軟弱.
有很多罪惡.
有很多痛苦.
有很多掙扎.
但是我願意先給你.
上帝啊.
我的身體.
我的心裡有很多軟弱.
有很多掙扎.
但是我都願意將我的身體獻給你.
我什麼都沒有.
我就是有這個瓦器.
我就是一個器皿.
我就給你.
求你使用.
求你適應.
上帝就會喜歡.
上帝就會適應了.
上帝看你這樣禱告會喜歡.
上帝看到你這樣禱告會很高興.

$^{641}$為什麼.
為什麼呢.
因為上帝說.
因為上帝是什麼呢.
這個孩子像我.
他就會說.
這個小孩子像我.
所以我們就把這個生命給上帝讓他使用.
所以我們就將我們的生命獻給上帝.
讓他使用.
所以我們很多時候.
我們痛苦的時候.
我們所想的.
怎麼樣得.
得.
怎麼樣索取.
有很多時候我們痛苦的時候.
我們想怎麼樣得.
怎麼樣索取.
但是現在我們奉獻.
怎麼樣給呢.
怎麼樣去吸出去呢.
我怎麼樣事奉呢.
我怎麼樣去施予呢.
所以這個就是事奉.
這個就是事奉了.
以基督的心為心.
我們適應上帝.
我們去適應上帝.
適應上帝.
上帝會充滿我們這個塵土.
適應了上帝.
上帝會充滿我們整個身體.
所以通過我們的身體.
很多人會受祝福.
所以通過我們的身體.
很多人會受到祝福.
所以上帝怎麼樣幫助你.
怎麼樣拯救你.
所以上帝怎麼樣幫助你.

$^{681}$拯救你呢.
你的問題得到解決是額外獎金.
你的問題得到解決是額外的獎金.
真正的報酬是什麼.
得到的報酬是什麼呢.
當你有痛苦的時候.
你把自己給上帝.
你把自己交給上帝.
上帝通過你工作.
當祂會通過你.
你變成像祂一樣.
就會好像祂一樣.
最後我還是想講迪士尼的童話.
最後我也想講一個童話.
但是這個童話其實不是迪士尼的.
這個是意大利幾百年前的童話.
不是迪士尼的童話.
其實是意大利幾千年前的一個童話.
大家都知道這個Pinocchio.
大家都聽過木偶奇遇記裡面那個木偶Pinocchio.
所以一個木匠很想有孩子.
但是沒有孩子.
所以有一天有一個天使可憐他.
就把他的木偶改成有生命的木偶.
就變成好像小孩子一樣.
一個木匠他沒有小孩子.
他很想有小孩子.
所以天使有一天將他所造的木偶.
就變成有生命.
好像小孩子一樣.
所以這個天使跟木偶說.
你要學習適應.
適應一下.
變成適應人.
適應這個好人的生活方式.
天使跟這個木偶說.
你要嘗試去適應人.
做好人的生活方式.
變成好孩子.
才能夠變成一個好的小孩子.

$^{721}$通過了以後.
考試通過了以後.
考試通過了之後呢.
才給你小孩子的身體.
我會給你一個真真正正的小孩子的身體.
但是現在還是木偶還是木頭.
現在暫時還是木頭造的.
這個孩子很調皮很調皮.
就不聽爸爸的話.
花天酒地.
這個小孩子很調皮很頑皮.
不聽爸爸的話.
整天花天酒地.
所以爸爸就想辦法找他.
結果爸爸被鯨魚吞掉了.
但是爸爸還活著在海洋裡頭.
那個爸爸去拯救他.
但是他發生了意外.
被一條魚吞掉了.
還在水裡面.
所以皮諾超就想辦法去海洋游泳.
好不容易救他的爸爸出來.
這個木偶呢.
皮諾奇他就很不容易地跳進海裡去救他爸爸.
他救爸爸.
但是他失去性命.
他為了救他爸爸.
這個皮諾奇自己就失去了生命.
他在海邊死了.
他在海裡面就死了.
這個天使出現了.
他說你學到.
他說你學到了.
什麼叫做好的孩子.
怎樣去維持一個好的小孩子.
所以給他新的身體.
所以給他一個新的身體.
小孩子的身體.
一個小孩子的身體.
我們的身體還是能夠分解.

$^{761}$能夠死亡的灰塵.
我們的身體也會繼續分解和死亡的灰塵.
效法基督是要變成一個上帝的兒女.
我們效法基督要成為上帝的兒女.
什麼叫上帝的兒女.
怎樣維持上帝的兒女呢.
向上帝的人.
要好像上帝的人.
不斷地給.
通過給能夠得.
通過我們所施與和給的人所得著.
以後我們就會得到復活的身體.
以後我們就會得到一個復活的身體.
Amen.
\newpage



\section{約翰福音 1:29-36}
\label{sec:032313}
\textbf{神的羔羊,以靈施洗}~~~~李應新牧師
\newline
\newline
連結: \href{http://www.church.com.hk/acms/content.asp?site=cdc\&op=show\&type=product\&code=032313\&layout=sermon}{\texttt{032313}} ~~~~ 語音日期: 2026-06-14
\newline
\newline
\hyperref[sec:032312]{\small{< < < PREV SERMON < < <}}
~
\hyperref[sec:index_chronic]{\small{[返順時目]}}
~
\hyperref[sec:032314]{\small{> > > NEXT SERMON > > >}}
\newline
\newline
約翰福音 1:29-36
\newline
\begin{longtable}{cl}
\hline
\hline
章節 & 經文 (和合本修訂版)\\
\hline
1:29 & \begin{tabularx}{0.7\textwidth}{X} 第二天，約翰看見耶穌來到他那裡，就說：「看哪，神的羔羊，除去世人的罪的！ \end{tabularx} \\ \\ \relax
1:30 & \begin{tabularx}{0.7\textwidth}{X} 這就是我曾說『那在我以後來的先於我，因為在我以前，他已經存在』的那一位。 \end{tabularx} \\ \\ \relax
1:31 & \begin{tabularx}{0.7\textwidth}{X} 我先前不認識他，如今我來用水施洗，為要使他顯明給以色列人。」 \end{tabularx} \\ \\ \relax
1:32 & \begin{tabularx}{0.7\textwidth}{X} 約翰又作見證說：「我曾看見聖靈彷彿鴿子從天降下，停留在他的身上。 \end{tabularx} \\ \\ \relax
1:33 & \begin{tabularx}{0.7\textwidth}{X} 我先前不認識他，可是那差我來用水施洗的對我說：『你看見聖靈降下來，停留在誰的身上，誰就是用聖靈施洗的。』 \end{tabularx} \\ \\ \relax
1:34 & \begin{tabularx}{0.7\textwidth}{X} 我看見了，所以作證：這一位是神的兒子。」 \end{tabularx} \\ \\ \relax
1:35 & \begin{tabularx}{0.7\textwidth}{X} 又過了一天，約翰同兩個門徒站在那裡。 \end{tabularx} \\ \\ \relax
1:36 & \begin{tabularx}{0.7\textwidth}{X} 他見耶穌走過，就說：「看哪，神的羔羊！」 \end{tabularx} \\ \\
[1ex]
\hline
\hline
\end{longtable}
$^{1}$弟兄姊妹平安.
今個星期我們很多弟兄姊妹都在憂傷之中.
特別我們在太魯閣城居民.
我們在一個很大的難過裡面.
我們求神憐憫.
亦幫助我們可以走出死暗憂谷.
我們有少少安靜的時間.
然後我們一齊同心祈禱.
親愛的天父求你幫助我們.
我們在極大的憂傷當中.
求你幫助我們.
聖靈幫助我們轉眼仰望主耶穌.
我們知道我們很軟弱.
求你恩典臨到我們眾人身上.
讓我們成為你合用的器皿.
流通的管子.
將盼望,將恩典,將慈愛.
帶給我們身邊受驚嚇的人.
亦都求你幫我們先飽得你的慈愛.
以致我們今天可以先在你面前蒙恩.
撫平我們心靈裡面的創傷.
我們繼續仰望你.
拖著你的恩手.
繼續走前面有生有活的路.
祈禱交托奉耶穌名靠.
阿們.
我們一起讀出約翰福音第八章第四十四節.
預備一二三.
因為心裡沒有真理.
他說謊是出於自己.
因他本是說謊的.
也是說謊人之父.
在這裡講到魔鬼撒旦是說謊人之父.
魔鬼撒旦將謊話放在我們心中.
將假的,不真實的東西放在我們心裡面.
所以我們提防,警惕.
免得這些謊話改變我們.
或者使我們陷入網奈.
謊話使人死亡.
在這裡說是殺人的.

$^{41}$因為謊話殺人.
我們今天牧師特意要講這件事情.
因為我們這個世界裡面.
被很多的謊話去改變我們自己.
擺錯了價值觀.
甚至殺害了我們的生命.
這個世界有什麼謊話呢.
很多的謊話.
例如如果讀不到這間小學.
你就死定了.
所以使我們很多家長都陷在網奈裡面.
讀不到這間中學.
你將來沒有希望了.
你讀不到書就死定了.
很多的謊話在我們心中.
使我們情不自覺地.
將不必要,不應該有的壓力.
放在我們身邊的人身上.
其實不單是家長.
其實很多人都有很多謊話.
丈夫很多時候.
或者太太也是.
覺得我老公很壞.
全世界最壞的他.
這是不是謊話.
很多價值觀都是謊話.
當然不是.
你旁邊的老公雖然是壞.
但是和旁邊的老公差不多壞.
不比的壞.
只要你回到那些夫婦團契.
你就知道每個老公都很壞.
對不起.
很多老公都很壞.
你老公算是這樣的.
很多謊話.
放了錯誤的價值觀在我們身上.
叫我們不能自拔.
所以求神提醒我們.
分辨真理和謊話.

$^{81}$這個社會裡面有很多謊話.
也讓我們陷在網絡之中.
在這段時間裡.
我心裡也覺得.
特別是為了那個家庭.
我們心裡很難過.
因為很多時候我們被這些謊話.
改變了我們生命的方向.
以致我們不能夠活得美好.
神的心意是讓真理讓我們活得美好.
在這裡提多一兩件事情.
我們為了這個家庭來禱告.
這個家庭.
其實是主裡面的家庭.
他們有回教會.
我們求神憐憫.
幫助他們.
特別是他們的教會.
也陷在一個很大的憂傷當中.
我們被謊話偶爾.
我們有時即使是信徒.
我們也會被謊話.
或者被惡者所勝過.
兩件事情.
我們昨天祈禱會有分享.
我們弟兄姊妹知道.
謊話是殺人的.
讓我們不單止自己受害.
很多其他人也受害.
昨天我們分享會的時候.
也有弟兄姊妹說到.
他們自己個人內心的經歷.
一個人選擇了.
用這個方法離開世界的時候.
其實全家受害.
和賢子孫.
他們跟著幾十年後.
他們的兒女.
他們的家人.
他們身邊的人.

$^{121}$都仍然在苦難當中.
因為那惡者的謊話.
使我們做了錯誤的決定.
求主幫我們.
我還想分享兩件事情.
我們有時也會心裡想.
為什麼會這樣發生呢?.
或者當中有些什麼.
為什麼謊話會令我們這樣呢?.
兩件事情.
有時有些叫做「調製」.
突然間「調製」了.
「調製」什麼呢?.
其實就是衝動「調製」.
有時也會衝動做些不合理的事情.
後來回想.
自己也覺得自己怎麼會做這種事.
為什麼會說出這種話.
或者做出這種事.
說話比較多.
做事也有時會有.
「調製」嘛.
你看我們吵架的時候.
也有時會「調製」.
說一些話後來會發覺.
為什麼會說出這種話呢?.
可以嗎?.
分享兩件事情.
一個是我自己身上.
就不要告訴我兒子.
我上個星期回家.
其實我回家是.
回家的時候.
其實我兒子和.
我女兒是在讀A-Level考試.
他們是全世界最遲的.
DSE考完很久還在考.
我兒子其實去了倫敦讀書.
但是被媽媽召回來.
看著妹妹.

$^{161}$幫她讀書.
但是怎麼說呢?.
我的女兒其實不太喜歡讀書.
其實地球人也不喜歡讀書.
總之就是不喜歡.
哥哥奉命來看著她讀書.
又要鞭策她.
要她溫習.
要她做練習.
但是弄著弄著.
弄著弄著.
妹妹也覺得很煩.
哥哥就覺得她擺爛.
於是妹妹就發脾氣.
擺爛.
不做練習.
我死她好過.
就這樣說.
突然哥哥跳掣.
我這麼辛苦來教你讀書.
你死我和你一起跳樓.
說這些東西嚇得.
他媽媽就嚇了一跳.
你們兩個都跳樓.
我怎麼辦呢?.
這叫做什麼?.
跳掣.
說完之後.
我問我女兒做了什麼?.
說了什麼?.
我說說而已.
但是我不知道哥哥跳掣這麼厲害.
不明白.
她說說而已.
我女兒其實我覺得她是擺爛躺平的人.
她怎會什麼呢?.
不過無論如何.
跳掣.
弄到哥哥跳掣.
行不行?.

$^{201}$有些時候.
也有些是心裡面有個漩渦.
那叫漩渦.
那漩渦是什麼呢?.
就是你想一遍.
或者情緒抑溫.
有個漩渦會一直沉下去.
就是個漩渦.
跳掣是一種.
漩渦是另一種.
漩渦一直沉下去.
整天都想著要這樣.
再說一個.
是我自己的經歷.
我有位教友.
不是我經歷.
是我的教友的經歷.
我的教友.
其實在幾十年前的時候.
他三姐妹的媽媽.
他媽媽有抑鬱病.
也很嚴重.
整天都想著要死.
也試過.
就是嘗試過.
那次嘗試.
救得回來.
進了醫院.
住了一陣子.
他們三姐妹下定決心.
下定決心輪著.
每人都陪著他.
輪著來陪.
隔了一陣子.
可能一個多兩個月.
那天是我教會的那位姐妹.
當更.
是不是叫當更?.
那更是她.
總之是她陪媽媽.

$^{241}$她被媽媽耍開了.
媽媽耍開了.
你整天黏著我都沒用.
我沒事.
你不要整天黏著我.
你幫我去買些什麼.
那一類.
那天她耍開了一個多小時.
她自己做了一些不能挽回的決定.
姐妹非常難過.
其實這件事已經隔了十幾二十年.
當然不會每天都想起.
但你可以想像中.
她被其他兩個姐姐.
說成怎樣.
叫了你.
什麼事不看緊媽媽.
就是那些.
但她也明白一件事.
如果她有心要這樣做的話.
難道我能夠看到她一生嗎?.
也不能夠.
但那一刻的難過.
這種傷痛.
差不多叫一生之久.
間中又會彈出來.
我們舊神都憐憫我們.
所以我們也有安排一些活動.
我們希望可以釋放惡者的謊話.
求主也幫助我們釋放惡者的謊話.
我們連程序表也等不及.
全部都是剛剛約好的.
下個星期六.
我們兩點半.
昨天就有祈禱會.
下個星期六.
我們邀請了一位醫生.
跟大家分享一下.
那些情緒.
你突然間跳掣那些.

$^{281}$有些漩渦是有徵狀的.
可以及早尋求幫助.
另外我們也有正向家長心理的專家.
所以在後面坐著.
所以幫我們說一下.
其實我很感恩.
過往的一年裡.
我們也有一起合作.
去幫助家長.
去分辨好一些謊話.
正向來教導子女.
但同樣地我也希望.
家長也好.
不是家長也好.
也要學習.
希望你也邀請身邊的人.
一起來學習.
能夠身心靈得著康健.
也能夠分辨到一些情緒上.
有警號.
需要怎樣去尋求幫助.
可以嗎?.
下個星期六.
我們一起邀請身邊的人.
我們自己也一起來學習.
求神幫助我們分辨哪些是謊話.
可以在當中.
求神使用我們.
我們繼續看下去.
今天我們要看經文.
剛才叫你揭曰漢.
我們剛才讀了一些.
但我讀不完.
我想今天你也一起看.
我們揭下去.
聖經曰漢福音第一章.
我們從十九節開始看.
這是曰漢所作的見證.
一章十九節.
一章十九節大家可以看到.

$^{321}$我們一起看.
我讀你也讀.
有些時候你讀.
有些時候我讀.
我先讀十九節.
這裡說到.
曰漢所作的見證.
記在下面.
猶太人從耶路撒冷差.
祭司和利未人到曰漢那裡.
問他們算.
你是誰.
我們一起讀二十一節.
他就明算.
並不隱瞞.
明算.
我不是基督.
他們又問他算.
這樣你是誰呢.
是以利亞嗎.
他算.
我不是.
是那先知嗎.
他回答算.
不是.
這裡說到.
他們問問你是誰.
常常在我們心裡.
或者在人身邊.
問你是誰.
你究竟是誰.
你是否先知.
是否那個先知.
是否以利亞.
他說我不是.
你不是的話.
於是他們再問.
第二十二節.
於是他們就說.
你到底是誰.

$^{361}$叫我們好回覆那猜我們來的人.
你自己說你是誰呢.
我們又讀二十三二十四節.
他說.
我就是那在曠野.
有人聲喊著說.
修直主的道路.
正如先知以賽亞所說的.
那些人是發利賽人.
猜來的.
我們讀到這裡.
原來他明說.
我就是那把聲.
而且那些人是發利賽人.
猜來的.
剛才說到.
是什麼人來的.
是祭司和那些利未人.
那些利未人.
祭司是從哪裡來的.
耶路撒冷來的.
有些人派他們來.
不是他們自己來.
他們派.
不是自己想來.
是有人派他們來.
派他們來的是什麼人.
發利賽人.
我們繼續讀下去.
我先讀二十五節.
使他們就問他說.
既不是基督.
又不是理令.
那也不是那.
先知為什麼施洗呢.
由二十六節一起讀.
約翰回答說.
我是用水施洗.
但有一位站在你們中間.
是你們不認識的.

$^{401}$就是在我耳後來的.
我給他解鞋帶也不配.
這是在約旦河外.
伯大利約翰施洗的地方.
作的見證.
原來他就說.
我是用水施洗.
在你們中間.
有一個你不認識的.
在我耳後來.
我是為他解鞋帶不配.
也為他作見證.
我又說回來.
這個問問我是誰.
你是誰.
為什麼做這些事.
是不是大字這樣問.
總合來說.
你是誰.
為什麼做這些事.
其實這段時間.
香港也有這些醞釀.
你是誰.
為什麼做這些事.
譬如有個媽媽在家裡生小孩.
你是誰.
你又不是醫生.
你又不是護士.
又不是助產士.
你為什麼在家裡.
明白.
你是誰.
為什麼做這些事.
明白.
大字這樣的意思.
幾個重複出現.
或者有個.
怎麼說呢.
有時候我們有個身份.
有個位份.

$^{441}$才做這些事.
不是一定不行.
有時候也是可以的.
但是給我們一個提醒.
挑戰.
你有什麼特質.
你為什麼會做到這些事.
施洗約翰其實是誰.
是天父派他做一個很特別的事情.
他原本是什麼人.
施洗約翰其實應該是一個祭司.
因為他爸爸就是祭司.
撒迦利亞是一個祭司.
他應該按照班次來電女公職.
如果這樣說.
當時是世襲.
這個祭司是世襲.
他如果安安分分.
應該在哪裡.
在耶路撒冷等候公職.
然後就怎樣.
然後就做祭司.
他現在去了哪裡.
去了曠野.
去了為人施洗.
你是誰.
你是不是才知道.
為什麼要做這些事.
說回頭.
每個人都有個特質.
有時候這些質問.
其實給我們檢視自己.
究竟我是誰.
我要做些什麼.
為什麼要做這些事.
每個人都有些特質.
人人都有些特質.
教會有沒有特質.
每間教會都有些特質.
我們教會有什麼特質.

$^{481}$我們教會的特質是什麼.
有什麼特質.
我們教會的特質.
是很認真的.
你知不知道.
你不知覺不覺.
我們教會的特質.
玩也很認真玩.
知不知道為什麼玩認真.
我見到部長在那裡說.
我們玩.
我們那些.
反過來有些人很認真的.
都好像玩一樣.
譬如阿叻.
他們很認真的.
都很兒戲玩.
但我們反過來.
我們玩的都很認真的.
你見到我們亂世遊戲.
時間都計算得很準確.
我們有我們的特質.
但我們是誰.
我們教會的特質.
是這方面.
但同樣地.
我們教會還有特質.
我們教會的特質.
是希望我們成為一間社區的教會.
不知你覺不覺得.
我們希望可以將神的愛.
帶給我們身邊的街坊.
未必很成功.
要我們特質.
我們希望可以做.
可以嗎.
我們希望可以做.
究竟我們是誰.
我們究竟的目標是怎樣.
我們希望可以成為神使用的器皿.

$^{521}$約翰是誰.
他究竟是誰.
他的目標.
很明顯.
我們再看一次經文.
第23節.
原來他是那把聲.
我就是那在曠野有人聲喊著說.
修直主的道路.
是那把聲.
修直主道路的那把聲.
希望神引領主耶穌.
來到眾人的面前.
我們都希望做到那把聲.
約翰做的我們都要做.
因為地上最大的是約翰.
主耶穌說的.
但是天國最小的都大過他.
我們同樣要做.
我們將修直主的道路.
引領人去到主的面前.
還有約翰要做甚麼呢.
帶領人受洗,施洗.
我是試用水施洗.
我們看經文.
我們繼續看.
跳過去看.
第29節我們看.
次日.
我們一起讀.
次日29,30,31.
次日約翰看見耶穌來到他那裡.
就說.
看那神的羔羊.
去世人聚業的.
這就是我曾說.
有一位在我以後來的.
反成了在我以前的.
因他本來在我以前.
我先前不認識他.

$^{561}$如今我來用水施洗.
為要叫他顯明給以色列人.
所以還有第二個.
約翰要做的就是.
用水來和人施洗.
亦宣告主耶穌是神的羔羊.
神的羔羊.
這裡遙遙指向甚麼呢.
指向在禹節裡面.
神的羔羊.
將血灑在門欄上.
神的審判不臨到當中.
神的羔羊.
他去宣告.
亦用水來施洗.
引領我們明白主的救恩.
約翰要做的事.
只能夠做到甚麼呢.
是用人去用水來施洗.
讓我們知道我們需要除去罪孽.
但等候著甚麼呢.
等候著用聖靈為我們施洗的主耶穌.
我們繼續看下去.
我們看三十二節.
我們讀到三十六節.
預備,一二三.
約翰又作見證說.
我曾看見聖靈.
彷彿甲子從天降下.
住在他身上.
我先前不認識他.
只是那差我來的用水施洗.
對我說.
你看見聖靈降下來.
住在誰的身上.
誰就是用聖靈施洗的.
我看見了.
就證明這是神的兒子.
再次人.
約翰和兩個門徒在那裡.

$^{601}$他看見耶穌行走.
就算看那這是神的羔羊.
我們讀到這裡.
這裡再次宣告.
再次即是第二天.
又再說.
這是神的羔羊.
向著兩個門徒說.
因為他也看到了.
他去證明.
施洗約翰.
用水和人施洗.
提醒我們是有罪人.
要來到主耶穌的面前.
來到神的羔羊面前.
領受聖靈.
因為我是用水施洗.
但是我看見聖靈降在誰的身上.
誰就是用聖靈施洗.
我看見了.
就為此作證明.
證明這是神的兒子.
所以約翰做了三件事.
約翰是施水.
做了三件事.
第一件事就是那把聲.
修籍主的道路那把聲.
向世界宣告.
第二.
行動.
用水向人施洗.
證明我們是.
或者指示我們是罪人.
第三.
去證明這是神的兒子.
去作見證.
施洗約翰所做的三件事.
我們也提醒自己.
讓我們知道.
因為聖經提到.

$^{641}$我先前不知道.
不是完全不知道.
是不是很知道.
似乎知道.
他是不是不認識耶穌.
不是.
他是耶穌的親戚.
但是後來他才知道.
這是神的兒子.
因為看見了一個很明確的證據.
那明確的證據是什麼.
那明確的證據是.
因為聖靈降臨在他身上.
我今天還想提醒一件事.
今天我們究竟是誰.
我們今天是上帝所使用的器皿.
去宣告主耶穌是神的兒子.
引導人去到神的面前.
我們要做那把聲.
修織主的道路.
我們要做什麼呢.
引領人去得著主的救恩.
我們提醒人.
我們是罪人.
需要主的救恩.
同樣證明耶穌基督是神的兒子.
在我們身上.
我今天還想提一兩件事.
如果我們這樣的時候.
我們也提醒自己.
我們究竟是不是在做這樣的事.
是不是在做修織主的道路.
引領人去到神的面前的事.
今天我還想提一件事.
我是誰.
我是神所使用的.
你也是.
我們是弟兄姊妹.
我們都在一起.
「同打那美好的仗」.

$^{681}$在我上一次.
在我還沒放假前就提到.
提到一篇信息.
有位弟兄來提醒我.
我看看.
轉到PowerPoint.
一邊看一邊看.
在經文那邊一邊看.
在5月17日.
有弟兄姊妹回應我.
那天我說大使命.
我有提到.
其實我沒有展示經文.
《馬太福音28章19-20節》.
大使命.
我相信你也記得.
大使命叫我們往普天下去傳福音.
給萬民聽.
凡我所教訓的.
都教訓他們遵守.
我就常與你們同在.
記不記得.
那段經文.
我有提及前一部分.
但弟兄提醒我.
你沒有提及後面的部分.
因為兩件事.
第一件事是傳福音.
第二件事是教訓門徒.
遵守主的教訓.
繼續看下一頁.
今天我所說的信息.
在5月17日提醒大家傳福音.
但是沒有強調.
信徒要成為門徒.
遵守主的教訓.
各位弟兄提醒我.
原文中.
這個N的意思.
是連接詞.

$^{721}$將19-20節.
寫到連在一起.
那麼重要.
整傳大使命.
傳福音.
但要做門徒.
可以嗎.
大致的說法.
當然有人提醒我.
我第一個感覺就是.
我是誰.
我是牧師.
你提我講到講漏的東西.
甚麼人.
不是.
是弟兄.
那又何妨呢.
提醒我就告訴你.
所以我覺得.
好啊.
因為我們都是需要.
互相提醒的.
同意嗎.
互相提醒.
我就告訴你.
所以肯定.
你看這些不是我的風格.
我拋一個重點.
弄出來給你看.
清楚一點.
繼續看下去.
不用看了.
下一個.
因此一定要一起強調.
避免肢體.
主教強調.
真實實踐主的教導.
所以我認同肢體.
給我愛心的回應.
這裡特別宣告.

$^{761}$互相勸勉.
我是誰.
我們是神所使用.
我們一切都是上帝所使用.
提醒人要來到主的面前.
互相教導.
實踐真理.
不過我還想繼續看下去.
我們先前知道.
先前不知道.
現在知道.
現在開始實踐.
不過我還有一件事.
想繼續說下去.
重點.
這裡提到.
神的羔羊主耶穌.
是用聖靈和我們洗禮的.
可以嗎.
今天你和我.
有沒有聖靈呢.
我們是否受了聖靈的洗禮呢.
不要懷疑.
我們能夠接受主耶穌的救恩.
已經領受聖靈.
聖靈一定在我們當中.
不需要懷疑.
聖靈在我們心裡面.
三件工作.
我們已經因為主耶穌.
已經賜下聖靈.
因為主耶穌離開世界之前.
肉身離開世界之前.
祂已經賜下保衛師.
我要離開你們.
但我們賜了一位保衛師.
在你們當中.
是永遠與我們同在.
所以我們不需要擔心.
我們能夠信主耶穌.

$^{801}$是因為聖靈.
如果不是聖靈.
沒有人能夠稱耶穌是主.
如果不是聖靈.
我們也不能夠接受主耶穌的救恩.
是因為聖靈.
我承認主是神的兒子.
我承認主耶穌是我的救主和生命主.
聖靈在我心中是證明,證據.
如果有聖靈在心中.
沒有人會說耶穌是可救助的.
羅馬書告訴我們.
還有第二件事.
我們有聖靈在心中.
我們必定會怎樣呢.
我們對罪有一個醒覺.
會為罪,為義,為審判.
自己責備自己.
我不詳細說.
但我們對罪有敏感.
聖靈在我們心中提醒.
我們可以拒絕聖靈.
我們會的.
有時也會的.
聖靈提醒我.
但我抗拒消滅聖靈的聲音是可以的.
但聖靈的聲音在我心中是證據.
聖靈在我心中提醒.
我在犯罪.
在虧缺上帝的榮耀.
提醒我.
一定有.
我們有聖靈在心中.
第三件事.
我們會結出聖靈的果子.
「仁愛,喜樂,和平,忍耐,溫柔,良善,淑女,節制」.
這些是一定有的.
雖然未必結得很靚仔,很成熟.
但起碼有結的.
求神也幫助我們.

$^{841}$不需要懷疑.
聖靈已經在我們心中.
求主幫我們.
聖靈在我們心中的時候.
我們就歡喜地在地上走天路.
我還想和你分享一件事.
有時我們也會覺得患得患失.
有沒有聖靈呢?.
我們一定有聖靈.
但有些弟兄姊妹跟我說.
有些教會是比較明確的.
會講方言.
我們教會是不講方言的.
有些教會又有神蹟奇事.
我們又比較少.
聖靈不同的表達.
在我們宗教會有獨特的遺訓.
表達的方式.
但是聖靈一定在我們心中.
我想分享兩個見證.
在我們教會裡.
聖靈有時比較.
天父會讓我突出一點.
讓我感受到聖靈與我們同在.
在這裡的.
不過突出一點.
他害怕我不知道的時候.
他會突出一點.
幾年前.
沒有幾年前.
兩年前.
有一位弟兄.
其實回教會很久了.
但沒有洗禮.
因為他身體比較軟弱.
經常回不了.
他有回.
但又回不了洗禮班.
經常坐不了那麼久.
沒有回.

$^{881}$但是他病情突然急轉直下.
進了ICU.
他想洗禮.
他有跟他老婆說.
我想洗禮.
遲些痊愈.
就上洗禮班.
但一直都沒有.
去到那天沒有機會.
回到洗禮班.
他太太跟我說.
我可不可以幫我丈夫.
就是牧師幫我掌洗禮.
這樣.
他進了ICU.
表達不了.
我可不可以幫他洗禮.
我真的要做決定.
有人會挑戰牧師.
他又不能告訴你.
我做決定.
按照聖靈給我的心.
不一定對.
幫他洗禮.
但其實他太太跟我有些疑惑.
都有一點點疑惑.
究竟是不是他的意願.
有些意願.
我跟他洗完禮之後.
那星期我又去探望他.
一起唱詩歌.
一起祈禱.
那天從來都沒有.
我平時都這樣做.
但那天很特別.
在ICU裡面.
我感受到聖靈所賜下的平安.
我覺得是天父額外給我的印證.
給我的姊妹.
我們一起去印證.

$^{921}$聖靈在這裡.
我問他.
平時不是這樣.
我感受到出人意外的平安.
心裡的溫暖.
很特別.
聖靈在此.
看我.
也給我們證據.
我們沒有冤枉他.
沒有冤枉那位弟兄.
他是上了.
神現在如立他.
聖靈在我們當中.
神在我們當中.
還想說一個見證.
我本來.
這位弟兄.
我認識他幾十年.
這也是幾十年前的見證.
聖靈.
有時很明顯在我們信主人心中.
這位弟兄.
一個聖靈的經歷.
見證就是.
有一天他睡覺.
聖靈提醒他.
看聖經.
揭聖經.
看,看,看,看.
看完之後.
聖靈提醒他.
聽.
他就停下來.
他就祈禱.
我接下來怎麼辦.
聖靈給他聲音.
沖涼.
他就去沖涼.
沖涼之後.

$^{961}$他又祈禱.
我接下來怎麼辦.
聖靈有聲音跟他說.
我那時候的綽號.
現在也有使用.
叫做布魯圖.
聖靈的聲音叫做尹布魯圖.
六點多.
那天早上六點.
搞了這麼多事.
到了六點.
打給我.
嚇得我有什麼事.
有個人六點鐘打給我.
他就把這件事告訴我.
他提了一段經文.
告訴我.
我問他做什麼.
但不知道為什麼.
聖靈有跟他說.
他沒有跟我說.
他就這樣跟我說.
我就跟他說.
這段經文是這個意思.
大概意思是怎麼樣.
我正在讀神學.
不知不覺.
然後呢.
然後就沒有然後了.
也不是完全沒有然後.
就完了.
因為他移民到英國.
我見到他.
這次我去旅行的時候.
我說你記不記得這件事.
他說.
很清晰.
又走出來.
那天的脈絡是很清晰.
不過你不說.

$^{1001}$我又放下了.
沒有記得.
我今晚回去再看那段經文.
他還記得那段經文.
我再看.
剛才是他跟我說的.
他就跟我說.
他說.
我終於明白.
少少明白.
為什麼神要叫我找你.
就是找我.
原因是什麼.
原因是我很記性.
記得幾十年後告訴你.
今天再看.
然後呢.
然後我在等然後.
現在剛剛開始.
繼續等然後.
神叫我們總不孤單.
叫我們在地上的時候.
互相作字體.
聖靈一定在我們心中.
有時候我們不察覺.
神有祂的心意.
還想多說一個故事.
唱首詩歌之前.
預備唱詩歌.
我們多說一個故事.
聖靈.
讓我們.
明白神的心意.
同樣.
神也賜我們有同伴.
剛才我們讀經文.
再次一日.
約翰和兩個.
你看看聖經說的.
三十六節.

$^{1041}$三十五節.
再次一日.
約翰和兩個門徒.
站在那裡.
他見耶穌行走.
就說看啊.
這是神的羔羊.
記不記得.
這裡是兩個門徒.
一個門徒叫安德烈.
另外一個是誰.
不知道.
因為接著沒有說.
不知道是誰.
總之一個是安德烈.
另一個是神.
叫我們不孤單.
兩個門徒在其中.
第二個門徒是什麼呢.
不知道是誰.
不過第二個門徒有沒有用呢.
不知道.
不過神讓這個門徒.
陪安德烈一起去見耶穌.
讓我們長路變近.
什麼事呢.
我還想起一個經歷和見證.
想和你說.
這個經歷.
又是幾十年前.
我讀中學的時候.
其實有一次去踩單車.
那時候踩單車.
其實我不太會踩單車.
不過總之別人叫我不敢不去.
一起去.
當年的時候.
我讀中學的時候.
踩單車最高級數是什麼呢.
平時是活道.

$^{1081}$灣仔租單車.
最高級數的人大衛租單車踩去.
新娘探.
這個最高級數.
有沒有人踩過.
最高級數.
厲害.
踩過這些.
很遠的.
很辛苦.
當年租單車.
要放下身份證或者學生證.
做抵押.
租單車在那裡踩.
浩浩蕩蕩踩去.
三個同學去.
我就記得三個同學.
三個同學踩單車的時候.
最厲害踩單車那個.
炒車.
不會踩單車那個.
沒事.
然後呢.
嚴重到要進醫院.
救傷車撿了那個.
進醫院的那個.
那剩下什麼呢.
剩下兩個.
但是呢.
有三輛單車.
那怎麼辦呢.
你說怎麼辦呢.
我們要還單車.
我們想了個方法.
我們兩個.
我就踩.
兩個人一踩.
踩兩輛車.
然後呢.
你看著單車.

$^{1121}$一個走回頭.
就踩那個回來.
然後又踩一段.
然後又放下單車.
你看著.
這次輪到你.
你走回去.
又踩一段.
行不行呢.
最後就還了單車.
我覺得很感恩.
一開始覺得沒可能.
不過.
只要有人同行.
長路便近.
我想起這件事的時候.
我一直在想.
那兩個同學是誰呢.
我忘記了.
我用很大的力氣都想不到.
那兩個是誰呢.
不過他是陪我一起踩單車的人.
有他的美意放在當中.
就正如我們今天在太古城.
一起相聚.
同行.
有他的美意.
可能幾十年後你不記得是誰.
不過這段路是有價值的.
這兩個門徒.
一個是安特烈.
一個呢.
不知道是誰.
不過一起同行.
\newpage



\section{士師記 7:18}
\label{sec:032314}
\textbf{聖經人聖經事:耶和華和基甸的刀}~~~~黃克勤牧師博士
\newline
\newline
連結: \href{http://www.church.com.hk/acms/content.asp?site=cdc\&op=show\&type=product\&code=032314\&layout=sermon}{\texttt{032314}} ~~~~ 語音日期: 2026-06-15
\newline
\newline
\hyperref[sec:032313]{\small{< < < PREV SERMON < < <}}
~
\hyperref[sec:index_chronic]{\small{[返順時目]}}
~
\hyperref[sec:032315]{\small{> > > NEXT SERMON > > >}}
\newline
\newline
士師記 7:18
\newline
\begin{longtable}{cl}
\hline
\hline
章節 & 經文 (和合本修訂版)\\
\hline
7:18 & \begin{tabularx}{0.7\textwidth}{X} 我和所有跟隨我的人吹角的時候，你們也要在營的四圍吹角，喊叫：『為耶和華！為基甸！』」 \end{tabularx} \\ \\
[1ex]
\hline
\hline
\end{longtable}
$^{1}$時間過得很快,又一天了.
又看看今天聖經人,聖經事,生命樹為你預備了.
耶和華和機電的刀.
昨天看的就是那個大墨餅.
一個軟軟的墨餅飛過來,我們所有帳幕都倒下了.
這個就是約阿斯的兒子機電的刀.
神已經將我們米電人交在機電的手裡.
哎呀,慘了.
其實那些米電人都挺有趣的.
其實我真的想不到他會有這樣的想法.
為什麼他會做這樣的夢.
這個夢是表示他心底裡很怕機電.
很怕機電,很怕對方是元帥機電.
當然他們可能聽到機電找到三萬多人.
後來有二萬多人走了.
不知道是不是那些人怕呢?.
剩下一萬人而已.
後來又走了,全都走了.
現在還有沒有人?只有機電一個不是.
還有我們算了大概三百人的.
三百人他都夠膽還留在那裡?.
啊,明天他就知道了.
最應該是這樣的說法.
這樣的想法.
為什麼會是做夢?哎呀,好害怕.
為什麼有墨餅飛過來,我們的帳幕都倒下了.
是神已經將我們交在他的手裡了.
還有人懂得解釋,這樣才厲害.
你會看到上帝在工作,他們的心已經很慌了.
剩下三百人,我怎麼會怕他呢?.
就算一萬人我都不怕.
因為我們不止.
耶和華和機電的刀.
耶和華和機電的刀,厲害過機電的刀嗎?.
當然厲害過了.
機電的刀已經像墨餅一樣飛到帳幕上了.
何況加上耶和華和機電的刀.
哎呀,那真是通殺,沒辦法輸了.
士師記七章第十五節.
機電在窺探的時候聽見.

$^{41}$這夢和夢的講解.
他就立刻敬拜神.
厲害了.
聽見了,聽見了,走了.
阿普拉走了,回去了.
回去就敬拜神.
他一聽見這個夢和夢的講解.
就敬拜神.
哇,機電成不成功呢?.
醒了,大膽到不得了.
因為依靠神,看到上帝的工作.
就在那個帳篷旁邊.
就敬拜神.
讚美你,你真是一個偉大全能的上帝.
讚美你,就敬拜神.
然後他就步行普拉.
回到以色列營中.
就跟所有人說了.
起來,起來吧.
耶和華已將米電的軍隊.
交在你們手中了.
起身了.
出發.
喂,當晚就出發了.
厲害了,沒有拖.
一回去就說,起來.
耶和華已將米電的軍隊.
交在你手中了.
於是機電將三百人分成三隊.
看到這個夢和夢的講解.
之後信心大增.
敬拜讚美神之後.
出發.
馬上出發.
那三百人呢.
分成三隊,用三國包抄的方法.
就將國和空瓶.
交給國人手裡.
不要問.
怎麼會忽然間有三百個空瓶這麼多.

$^{81}$絕對有可能,因為有幾萬人來過.
走了二萬二.
然後一萬都走了.
剩下三百.
我想他們一定拿了很多空的水罐.
這些瓶子應該是放水的.
一定是,或者放食物.
或者放水.
食物多數放在籃子裡.
用布蓋住.
而空的瓶子多數裝水.
有沒有三百個瓶子這麼多.
我想一萬個一定有.
扔在礦上.
那些人走了.
空的瓶子可能不拿了.
將豪閣.
吹的豪閣.
和空瓶就交給國人.
大家來拿武器.
哇,厲害厲害.
會不會劍很重.
不用劍,會用刀的.
會不會很重,不用想.
因為我們的武器.
就是一個吹的豪閣.
還有牛閣,羊閣,羚羊閣之類.
空的瓶子.
大家一人拿一個.
平內都藏著火把.
將點著的火把.
放在瓶子裡.
那怎樣固定火把?.
拿的時候怎麼發來發去?.
動來動去怎麼拿?.
我想他們有他們的辦法.
可能用泥把火把插進去.
插進去.
怎樣插?燒著.
不用想那麼多.

$^{121}$一個瓶子,下面放些泥.
一半左右,將火把插進去.
固定了,在瓶子裡.
然後才點火都可以.
不用點了才放進去.
臨出發就點.
火把在裡面.
一定要先點,不然去到那裡.
啪啪啪,點不著,沒有打火機.
還要用火石.
一定是點著了才去.
所以有三樣武器.
就是一個吹的.
蠔角,羊角,牛角.
空的瓶子裡.
放著燒著的火把.
我認為火把要先點著.
正如我們要做一件事.
事先都要有火在心裡.
沒有火在心裡.
去到見到了.
想點火已經遲了.
就吩咐他們.
你們要看我行事.
我去到營裡的旁邊點一做.
你們就點一做.
這件事是有些有趣.
原因就是.
我們分了三隊人,每隊一百人.
去到我們分散.
然後你聽到我一聲令下.
你打爛瓶子.
然後拿著火把.
然後大聲吹蠔角.
叫爺和華和機電的刀.
全部叫了.
一會兒你見到他們就這樣做.
全部叫了.
但是他們沒有.
機電都很有趣.

$^{161}$去到的時候看著我怎樣做.
你們就怎樣做.
為什麼不事先做一下這個briefing.
有個簡報先.
有個簡說先.
去到有心理準備.
沒有的,你這麼心驚驚.
拿著一個水瓶裡面.
有個火把拿著蠔角.
出發了,武器就是這樣.
你以為我急人嗎.
千萬軍隊.
樂途無數.
你們要看我行事.
我到了營的旁邊.
旁邊而已.
你們就要這樣行.
我跟隨我的人.
吹角的時候.
你們都要在營的四周吹角.
聽到我那隊吹角.
你們就在營的四周吹角.
然後一起喊叫.
爺和華和機電的刀.
其實有的.
我們去到.
我就會吹角.
你們聽到吹角就跟我一起吹角.
還有喊一句口號.
是我們的武器.
那句口號就是.
爺和華和機電的刀.
Sir,那我們.
武器弓箭在哪裡拿.
是不是搶他們那些.
沒有.
那我們只是.
喊的.
Sir,不是吧.
喊什麼?爺和華和機電的刀.

$^{201}$就有力量.
能夠殺到他們.
你會聽到.
當日的士兵怎麼想呢.
莫非又來個.
爺彌哥之役.
喊叫.
就會整個倒.
都可能有這樣的背景.
他們一定知道.
這麼偉大的事.
總之照著我這樣做.
一起喊叫.
爺和華和機電的刀.
什麼的刀是指.
他的軍隊.
這是爺和華和機電的軍隊.
的意思.
爺和華和機電的軍隊.
刀啦.
明白嗎.
舉個例子.
今天.
有犯罪的事情.
警察來到就說.
不要動,警察,警察,警察.
全部不要動,站好.
走到牆邊,警察.
就是說.
警察,警察.
這裡是說爺和華和機電的軍隊.
爺和華和機電的軍隊.
來了,大家不要動.
這裡是這樣的意思.
機電和跟隨.
他的一百人在.
三更之初.
剛剛才換更的時候.
你看看那三更.
三更其實比較都.

$^{241}$過了半夜了,原因就是.
當晚.
找完這三百人.
那三百人其實都很睏.
然後夜晚機電就走.
下去他們的營地.
看完,哇,爺和華和機電的刀.
有這麼大一塊餅,這個夢.
然後走回來,就開始帶著他們.
拿著自己的瓶子就開始下去.
這個時間.
已經過了半夜了.
到三更半夜.
三更半夜.
就開始.
他們日夜換更.
上一隊很睏,這一隊也很睏.
剛剛半夜換更.
來到營房旁邊.
就吹角,打破手中的瓶.
三隊人都吹角.
第三隊打破瓶子.
拿著火把.
右手拿著角,喊叫說.
爺和華和機電的刀.
你看到他們就開始這樣.
打破瓶子有什麼作用呢?.
啪一聲,周圍都聽到啪啪啪啪啪啪啪啪啪啪啪啪啪啪啪啪啪啪啪啪啪啪啪啪啪啪啪啪啪.
然後一驚醒一看,啪啪聲驚醒,嚇醒,聽到嘩!.
這麼大火光,周圍都有火,什麼事?.
然後右手拿著角吹號角,吹號角吹什麼呢?.
吹雜,什麼叫?叫做結雜號,大家來啦,集合打,吹號角就是進攻,嘩!.
聽到砰砰聲,即是打爛東西,然後火光紅紅,然後聽到吹起角,即是來進攻了.
然後聽到很多人叫什麼?.
耶和華和機電的軍隊,耶和華和機電的軍隊,嘩!機電的軍隊,機電的刀已經嚇到臉青了.
再加上耶和華和機電的刀,大家都知道以色列人的神耶和華出了名好打.
由埃及一直曠野,直打到迦南,打到耶利哥,然後打了整個應許之地,迦南地.
都是耶和華神那麼厲害,所以耶和華和機電的軍隊的刀,嘩!死咯.
他們在營的四圍,各站各的地方,全營的人都亂數.
米電軍隊聽到立馬跑出來,三百人不停的吶喊,耶和華和機電的刀,打到你了,嘩!他們逃跑.

$^{281}$三百人就摧割,耶和華和機電全營的人用刀互相擊殺.
唉!就算這樣跑,怕一怕跑出來,會不會見到自己軍隊就會用刀斬他?.
當然未必會啦,當然未必會啦,為什麼要穿制服?就是一看就知道自己人.
晚上睡覺會不會脫了制服,所以難認一點,可能都會啦,暗啦,那些火光影不清楚.
不過都未必會用刀互相擊殺的,最重要就是這幾個字.
耶和華使道人用刀互相擊殺,是神令他們互相用刀擊殺.
那他們就逃命,可能當時有驚慌啦,有人跑出來阻一阻,不是用刀先砍了他,這樣都不一定啦.
不過是耶和華使人用刀互相擊殺,他們就逃到西尼拉的白哈士他.
直到靠近哈巴的阿伯米荷拉,這些名字我就沒有翻譯了.
他們就逃跑呀,到處跑呀,於是很多軍隊,一定是幾萬軍隊就亂了籠.
自己見人就砍,不是啦,黑暗中不知道是誰,見到有人出現就先砍他.
於是以色列就從拿忽他里亞撤往嗎哪西,全地聚集追趕米殿人.
馬上派人回去再召集人馬出來,嘩,本來一兩萬人可能馬上聽到馬上回頭拿著工具來追殺.
召集了他們,另外機電再做了一件事,就是打發人走遍.
走遍,我用紫色紙都是經文來的,只不過我想大家要有另一個重要性,打發人走遍以法連山地.
一會明天就講以法連山地,以法連山地就說你們都一齊下來,在山上下來跟我們一齊攻擊米殿人.
我想你們的分工是什麼,在戰爭裡面,有,機電說你們下來的分工就是爭先大家跑去約旦河的渡口.
直到白拉巴,白拉巴的意思就是淺灘的房子,機電請人叫以法連人在山上下來.
你們快點跑去守住約旦河的渡口,不要讓他們過河逃跑,因為他們已經進來了.
現在一打輸仗,逃跑當然是逃回家裡那邊,他們就會過河回到東邊.
守住就是不讓他們過約旦河走回他們的家鄉,他們這樣進來,誰知道現在逃跑他們就要過約旦河.
守住不讓他們過,整條河守住不讓他們過,直守到哪裡呢?守到一個淺灘.
就叫做白拉巴,巴拉就是白巴拉,白巴拉就是淺灘,白就是房子,或者屋,或者家之類.
譬如伯特利,伯利恆等等,這些都是房子,白巴拉,可能去到約旦河某個地方是淺灘的地方.
於是很多人在那裡建淺灘的房子,守住那裡,淺灘容易過河,整條約旦河守住了.
特別是碼頭不讓過,淺灘沒得去,於是以法連人很聽話,馬上聚集.
把手約旦河的道口,直到白巴拉,就在那裡守住了.
並且真的有用,機電真的有辦法,只要他們逃跑,未必那麼多人,我們只有三百人追他.
再加上幾個支派的人,再走出來追他,都比他跑得早很多.
誰知就是叫以法連人快點過來,守住那條河,不讓他撤退,不讓他過河回家.
就抓了米電人的兩個首領,叫做俄勒,俄勒是搶奪的意思.
一個叫西伊伯,石伯,石伯,不是,一個叫西伊伯,石伯,石伯,石伯就是野狼的意思.
把這個首領抓了,他們就把俄勒殺在俄勒盆石上.
即是說抓他之後,你哪裡人?我叫俄勒,你家在哪裡?那邊.
就押他回家,然後在家裡當做殺他,在家鄉.
西伯,西伊伯就殺在西伊伯的酒寨那裡,即是說回到他家的酒寨那裡將他殺死.
即是說全屋的人都知道,全村人,全眾人都知道,哎呀,死了我們的王.
兩個王來的,兩個首領,又追趕米電人,將俄勒和西伊伯的首級.
然後再追趕約旦河,他們守住,那兩個首領來到被他們抓了.
抓了押回家,殺了,然後將首級帶回去過河,交給姬甸.

$^{321}$兩個被我們殺死了,首級在這裡,嘩,真厲害.
所以你見到姬甸這一次大勝,以色列人大勝,這個上帝很奇妙的作為.
請你記得,是耶和華使米電人互相用刀擊殺他們的同胞,弟兄.
是不是姬甸那三百人那麼厲害呢?是,不過不是那麼厲害.
是耶和華的刀,耶和華的能力更厲害.
不過,耶和華和姬甸的軍隊,成為他們的一個武器,一個手段.
就是他們信靠耶和華,而耶華神幫助了他們.
好了,今天跟大家講耶和華和姬甸的刀.
請你記得,他們是依靠著神去的.
第一點,拿什麼武器?水瓶,火把在裡面.
耗過三百個人,進攻幾萬的米電聯軍,唉,沒可能贏.
是耶和華神使他們一驚醒,聽到瓶子打碎的聲音,火光閃閃.
接著聽到吹角,呼喊聲音,耶和華和姬甸的軍隊來到了.
進攻,唉,個個驚到散了,見人就斬,殺死了自己人,然後逃跑.
接著再叫,一叫,那些人就馬上下來幫忙,大獲全勝.
這個是誰做的?耶和華做的,也是姬甸做的.
不過這個叫做耶和華的刀是更加合適.
不過,耶和華都讓姬甸可以和他一起參與.
所以叫做耶和華和姬甸的刀,整個故事你看到的是什麼呢?.
是耶和華的刀,用姬甸這個麥餅,不是小麥餅,是大麥餅軟軟的.
去贏了米電人,其實是耶和華的刀的工作.
債美主,我們將榮耀歸給他,不要搶奪上帝的榮耀.
我們一起祈禱,債美主,你是一個德性的上帝.
你的刀,你的軍隊,你的能力才是我們最要依靠的.
如果以在戰爭當中,很清楚,好像一個夢一樣.
我們只不過是一個大麥餅,從來沒有人拿過大麥餅做武器.
他一塊餅拿在手上,好像一個披薩一樣,怎麼可以做武器呢?.
也不是一個貴重的物質,只不過是人日常吃的東西.
不過實際上你可以用他,好像姬甸一樣,一個這麼害怕的人.
這麼要求上帝證據的人,你用他這個不值錢的麥餅.
又不是武器,去勝過了千萬的姬甸軍隊.
成為那些這麼不好的,這麼罪惡的人的一個審判的鞭.
將他們的鞭打回頭,主要是求你自己幫助我們.
讓我們更加信靠你,如果是你的揀選,是你自己的使用.
是你的同在,與你自己的作為,我們是放心的.
是必定德性的,因為是你,耶和華的刀最為重要.
願你得罪 榮耀讚美.
祈禱奉耶穌基督的名.
阿們.
\newpage



\section{羅馬書 8:1-17}
\label{sec:032315}
\textbf{惟獨在基督裡}~~~~周依娜牧師
\newline
\newline
連結: \href{http://www.church.com.hk/acms/content.asp?site=cdc\&op=show\&type=product\&code=032315\&layout=sermon}{\texttt{032315}} ~~~~ 語音日期: 2025-11-23
\newline
\newline
\hyperref[sec:032314]{\small{< < < PREV SERMON < < <}}
~
\hyperref[sec:index_chronic]{\small{[返順時目]}}
~
\hyperref[sec:032316]{\small{> > > NEXT SERMON > > >}}
\newline
\newline
羅馬書 8:1-17
\newline
\begin{longtable}{cl}
\hline
\hline
章節 & 經文 (和合本修訂版)\\
\hline
8:1 & \begin{tabularx}{0.7\textwidth}{X} 如今，那些在基督耶穌裡的人就不被定罪了。 \end{tabularx} \\ \\ \relax
8:2 & \begin{tabularx}{0.7\textwidth}{X} 因為賜生命的聖靈的律，在基督耶穌裡從罪和死的律中把你釋放出來。 \end{tabularx} \\ \\ \relax
8:3 & \begin{tabularx}{0.7\textwidth}{X} 律法既因肉體軟弱而無能為力，神就差遣自己的兒子成為罪身的樣子，為了對付罪，在肉體中定了罪， \end{tabularx} \\ \\ \relax
8:4 & \begin{tabularx}{0.7\textwidth}{X} 為要使律法要求的義，實現在我們這不隨從肉體、只隨從聖靈去行的人身上。 \end{tabularx} \\ \\ \relax
8:5 & \begin{tabularx}{0.7\textwidth}{X} 因為，隨從肉體的人體貼肉體的事；隨從聖靈的人體貼聖靈的事。 \end{tabularx} \\ \\ \relax
8:6 & \begin{tabularx}{0.7\textwidth}{X} 體貼肉體就是死；體貼聖靈就是生命和平安。 \end{tabularx} \\ \\ \relax
8:7 & \begin{tabularx}{0.7\textwidth}{X} 因為體貼肉體就是與神為敵，對神的律法不順服，事實上也無法順服。 \end{tabularx} \\ \\ \relax
8:8 & \begin{tabularx}{0.7\textwidth}{X} 屬肉體的人無法使神喜悅。 \end{tabularx} \\ \\ \relax
8:9 & \begin{tabularx}{0.7\textwidth}{X} 如果神的靈住在你們裡面，你們就不屬肉體，而是屬聖靈了。人若沒有基督的靈，就不是屬基督的。 \end{tabularx} \\ \\ \relax
8:10 & \begin{tabularx}{0.7\textwidth}{X} 基督若在你們裡面，身體就因罪而死，靈卻因義而活。 \end{tabularx} \\ \\ \relax
8:11 & \begin{tabularx}{0.7\textwidth}{X} 然而，使耶穌從死人中復活的神的靈若住在你們裡面，那使基督從死人中復活的，也必藉著住在你們裡面的聖靈使你們必死的身體又活過來。 \end{tabularx} \\ \\ \relax
8:12 & \begin{tabularx}{0.7\textwidth}{X} 弟兄們，這樣看來，我們不是欠肉體的債去順從肉體而活。 \end{tabularx} \\ \\ \relax
8:13 & \begin{tabularx}{0.7\textwidth}{X} 你們若順從肉體活著，必定會死；若靠著聖靈把身體的惡行處死，就必存活。 \end{tabularx} \\ \\ \relax
8:14 & \begin{tabularx}{0.7\textwidth}{X} 因為凡被神的靈引導的都是神的兒子。 \end{tabularx} \\ \\ \relax
8:15 & \begin{tabularx}{0.7\textwidth}{X} 你們所領受的不是奴僕的靈，仍舊害怕；所領受的是兒子名分的靈，因此我們呼叫：「阿爸，父！」 \end{tabularx} \\ \\ \relax
8:16 & \begin{tabularx}{0.7\textwidth}{X} 聖靈自己與我們的靈一同見證我們是神的兒女。 \end{tabularx} \\ \\ \relax
8:17 & \begin{tabularx}{0.7\textwidth}{X} 若是兒女，就是後嗣，是神的後嗣，和基督同作後嗣。如果我們和他一同受苦，是要我們和他一同得榮耀。 \end{tabularx} \\ \\
[1ex]
\hline
\hline
\end{longtable}
$^{1}$(錄音/錄影).
此錄音/錄影的版權是屬於基督教宣導會北角堂所有.
我們的屬靈傳統之一,是以真理教導為本.
歡迎您將此錄音/錄影的連結傳給親友.
懇請注意,未經授權請勿轉錄或剪輯本錄音/錄影.
甚願好像《使徒行傳》二十章三十二節下所說.
「這道能建立你們,叫你們和一切誠誠的人同得基業」.
第三堂的弟兄姊妹平安.
就說好像很久沒見過我講道,是不是不再講呢?.
其實我轉去第二堂服事.
所以不能常常來第三堂見到大家.
但是今天很感恩,再一次見到大家.
心裡面都很牽掛著,亦都是很有親切的.
我們都有一個祈禱.
阿爸父啊,今天我們是你親愛的兒女.
我們可以坦然無懼來到你的面前去敬拜你,去愛你.
你思念我們不過是塵土.
但是你竟然去紀念我們這些不配有罪的微臣.
差耶穌基督下來陳刑去拯救我們.
並且為我們寫生在十字架上.
並且身體復活.
叫我們將來都可以有屬天的盼望.
這一切一切的賞賜都是你豐沛的愛.
聖靈啊,求你指教我們並充滿我們.
叫我們在每天的日子裡面.
常常依靠你的能力去行出新造的人的樣式.
禱告奉耶穌基督的聖名而求.
阿們.
有一次有位宣教士跟我說.
「Ina,我們要跟基督徒再傳一次福音」.
明不明白他的意思?.
有時信久了,忘記了福音的內容和福音的大能.
所以藉著這一季的羅馬書.
我們重新去兼顧我們起初的信仰.
我覺得讀羅馬書.
只要你很專心去研讀.
都不難的.
不過最難的是甚麼呢?.
去實踐出來.
所以今天讓我們立定心智.

$^{41}$靠著聖靈的能力.
來活出新造的人的責任.
剛才主席已經幫助我們讀了羅馬書八章一至十七節.
是比較長一點的.
我有一個撮要讓大家比較容易掌握.
一會兒中心的內容.
福音不只是基督的救贖.
福音也是基督改變生命的能力.
叫輕易的罪人不再定罪.
從靈性的死亡中活過來.
從此成為神的兒女.
在基督裡靠著聖靈.
追求公義正直的成聖生活.
這是一個好消息.
在基督裡我們被稱義.
一個很簡單的定義.
稱義就是不被定罪.
你今天來到神的殿裡.
可以在神的面前.
不是因為你有甚麼好.
而是神的眼睛看著誰.
是耶穌基督.
祂成為我們的贖罪祭.
成為贖罪的羔羊.
以致祂所流出的寶血.
遮蓋我們的罪.
以致上帝的眼睛再臨到我們.
所以今天我們來到敬拜.
我們來到祂的生活.
完全是耶穌基督為我們去成就.
而在基督裡.
我們就可以靠著聖靈的能力.
目的是第四節.
讓我們律法所要求的義.
就是去罪的生活.
可以完全表現出來.
記不記得上星期我們說到.
律法本是美善.
但因為肉體而變質.
被罪利用.

$^{81}$以致我們行不出上帝的心意.
而罪的意思就是永遠都打不中.
上帝的要求.
就像射箭一樣.
達不到的雄心.
所以保羅說.
「我真是苦」.
而當你感覺真是苦.
也好一點.
因為你好像有努力去淺行神的話.
當你連「我是苦」這種感受都沒有的時候.
你就要小心.
你是否成為了一個肉體的基督徒.
而不是一個屬靈的基督徒.
「清義不被定罪」.
本來罪和死亡.
這個咒詛的鎖鏈.
已經斷開了.
但弟兄姊妹.
你有沒有意識到.
綑綁你的罪的鎖鏈.
已經釋放了呢?.
有一個心理學.
裡面有一個.
這樣的.
稱呼.
大笨象與小木樁.
知不知道是甚麼來的?.
你看到這個大笨象.
其實他是力大無窮.
他有足夠的力量.
來向前步行.
但他竟然.
被鎖實了.
他不敢前行.
這麼奇怪.
原來在他小時候.
這個小笨象.
被綁在.
這個小木樁裡.

$^{121}$可能是馬戲團.
因為他太小.
以致他爭奪不到.
這個鎖鏈.
所以越來越久的時候.
他就宣佈放棄了.
反正我怎樣爭奪.
都爭奪不到這個鎖鏈.
我一定不能夠.
向前走.
我們會覺得.
這個小笨象實在可笑.
但是.
你回想一下.
自己有沒有被罪的鎖鏈.
綑綁著你呢?.
剛才大笨象.
他學到一樣東西.
就是.
Learned helplessness.
他覺得我永遠都是無助.
我永遠都做不到.
我永遠都爭奪不到這條鎖鏈.
你會不會都一樣.
面對你的罪的時候.
貪婪,嫉妒,怨怒.
情慾等等.
這些試探,這些罪的時候.
你會覺得.
我做不到.
我永遠都改變不了.
於是你就像.
大笨象一樣.
會有一種習慣性的無助.
而這種的習慣.
是讓你更習慣犯罪.
以致你不願意.
再去意識到.
原來我的罪的鎖鏈.
已經是.

$^{161}$靠著耶穌基督.
斷開了.
弟兄姊妹你要問自己.
曾經綁住.
你的鎖鏈.
其實已經鬆了綁.
鬆了綁.
無論是原生家庭.
無論是你自己的成長.
無論是什麼原因.
讓你身出.
這些罪的時候.
或者是別人的罪.
去傷害你的時候.
每個人都捆著不同的鎖鏈.
但是根源.
都是那個的罪.
被自己的罪所傷.
或者是被別人的罪所傷.
但這條鎖鏈.
捆綁著你的鎖鏈已經鬆綁了.
你知不知道.
如今那些在基督耶穌裡的人.
就不被定罪.
這條鎖鏈已經鬆開了.
你是一個獲釋者.
今天.
我們就是一個獲釋的新造的人.
是一個新的身份.
來去活出天父對神的兒女.
他有新的一種生活方式.
但這種生活方式.
當你真的去踐行的時候.
你就會發覺到.
有一些矛盾.
或者是那種鬥爭.
你覺得很痛苦.
為什麼呢?.
我們就要來去認識.
一些保羅.

$^{201}$所說的教義.
在五至八節裡.
他提醒我們.
在基督裡的那種拔河.
肉體,聖靈,肉體,聖靈.
他指到我們要.
專注思想,聖靈.
你會在這裡看到.
五至八節.
就好像勢不兩立的兩股勢力.
一個是肉體.
一個是聖靈.
如果你隨著.
隨從肉體的時候.
會帶來什麼呢?.
死亡.
隨從聖靈呢?.
生命和平安.
如果你隨從肉體的思想.
原來是與神為敵.
是勢不兩立的.
我們看到.
我們如果還是活在肉體的基督徒的時候.
第八節說.
屬肉體的人不能討神的喜歡.
因為你在最終與神為敵.
又如何討神的喜歡呢?.
神對你罪的鎖鏈的捆綁.
已經解開了.
我們應該向前行.
向一條光明生命的道路.
大堅姐妹,我們首先來了解一下.
什麼叫「略先」.
你就會知道.
如何依靠聖靈.
解脫這條鎖鏈.
「略」的意思是什麼呢?.
是定律.
什麼叫定律呢?.
例如我手上的物件.

$^{241}$如果我向下扔的時候.
會發生什麼事呢?.
一定會掉下去.
請問是什麼定律呢?.
地心吸力.
地心吸力會不會改變呢?.
除了你上了太空.
上了太空.
還有一些律.
所以我們要認定.
在你有生之年.
未去見主耶穌之面之前.
有一個律在.
就是.
罪和死的律.
其實耶穌不是已經.
解開了我們的罪.
沒錯是.
但我們的罪性仍然在.
所以.
有罪的律.
仍然影響著我們.
同時我們基督.
賜了聖靈給我們.
就是生命之靈.
聖靈的律.
同時也存在我們的生命裡.
所以這兩個律.
就好像拔河一樣.
而在中間.
拔河的遊戲.
怎樣才算是勝利呢?.
就是看著中間.
那條紅色的線.
那條繩子.
向左邊邊就是勝利.
而這條紅色的線.
也可以代表著.
我們屬肉體的心思.
的一種律.

$^{281}$當我們這個心思的律.
充滿著罪性的時候.
我們很容易會偏向邊邊的律.
罪的律.
我問你.
罪的律和聖靈的律.
哪個律比較強大?.
聖靈的律.
但在現實裡面呢?.
在現實裡面.
是罪的律.
我們來看一下我們充滿罪性的心思.
其實是.
什麼樣的解釋呢?.
在舊約和新約.
希伯來文.
這個Soul.
即是人的心思或是生命.
它的意思是.
人的情感.
意志,思想.
希伯來文叫Nephesh.
希臘文叫Psyche.
它通常形容的.
就是人的意願.
或者是意欲.
但在舊約裡面.
特別是曠野之旅裡面.
不停出現這種Soul.
發生了一個問題.
就是那些以色列民.
忘記神的愛.
忘記神的恩情.
以致隨從自己的意欲.
特別是那種的Desire.
和那種的Want.
特別是.
生活上的欲望.
是隨心所欲.
隨己而行的.

$^{321}$特別提到的是口服之欲.
記不記得沒有食物.
沒有水喝.
沒有吃火鍋.
全部都是這些口服之欲.
我們香港人都很懂得吃.
很懂得享樂.
這些都是沒有問題的.
但是如果.
我們每天.
中日思想的.
這些口服之欲.
或者是享樂之欲.
這些貪戀世俗的思想.
就會讓我們傾向了.
成為一個肉體的基督徒.
所以思想很重要.
思想思想.
不斷重複了五次.
思想.
在和合本裡面.
是體貼的意思.
即是Focus.
你專注一些什麼的事情.
而在原文裡面.
如果你專注肉體的事.
你的Outlook.
會Shaped.
By the things of the flesh.
就是你整個的行為.
你的外觀.
都是一直被肉體影響.
我們說.
思想影響行為.
行為影響習慣.
習慣就成為.
你的品格和生命.
所以如果.
我們不斷體貼.
肉體的事的時候.

$^{361}$我們好像不能夠.
去順服神的律法.
我們也無能.
去順服神.
我們要知道我們的無能.
所以我們就要思想.
聖靈的事.
到底我們的思想.
讓我們偏向了罪的律.
還是偏向了.
屬靈的心思呢?.
我記得有一個.
菲律賓姐姐.
她每次晚上.
她做完自己所有的.
家務之後.
她就會和那位僱主.
她就說.
我現在可不可以去公園.
初時的僱主就以為.
她去公園和同鄉.
一起聊天.
但是原來.
她去公園知道做什麼事嗎?.
她是去和同鄉.
傳福音.
她也是常常讀聖經的.
一位菲律賓姐姐.
她的心思意念.
是很充滿著上帝的話.
上帝的愛.
和上帝的使命.
所以當我們在日常生活裡面.
其實不妨.
可以想一想.
在我現在的崗位.
在我現在的地步裡面.
其實我可以怎樣更多.
改變別人的心思.
我記得也有一個職場.

$^{401}$有一位弟兄.
他是上師.
他曾經見到一個同事.
不是基督徒.
見到他深陷一個很低谷.
很低谷的一個痛苦幽谷.
這一位上師.
他是基督徒.
但是他很低調.
他也是選擇了.
一張金句卡.
在後面寫了一些鼓勵的說話.
然後給下屬.
十五年之後.
這位上師退休了.
在他退休的宴會裡面.
或者是farewell party裡面.
這一位職員.
他就來到他面前.
在他銀包裡面.
拿出了這一張.
差不多發黃的一張卡片.
是十五年之前.
他說老闆.
多謝你十五年前.
給了這張卡片我.
我收藏在.
我的銀包裡面.
足足有十五年.
是這張卡片鼓勵我.
可以站得起來.
我們可以在.
不同的場景裡面.
都可以讓聖靈的律.
去充滿我們.
幫助我們有聖靈的心思.
所以控管思想.
是很重要的.
你可能會問.
為什麼要控制.

$^{441}$的確是.
有一個思想就這樣出現.
好像現在說的.
你可以在想喝茶.
那你怎樣去控管呢.
我不會怪你的.
因為思想就是這樣出現.
但是我們有一個方法.
我教你.
其實只有一個方法.
叫做替代法.
替代法.
大家看一下這兩杯水.
一杯.
其實是漂白水.
我可不可以.
將我旁邊那杯.
純淨的水.
混在漂白水裡面去喝.
你就算.
怎麼稀釋到1\%.
你都會出傷你的喉嚨.
你的胃甚至你的生命.
是不是.
是勢不兩立的.
聚的律和聖靈的律.
你一定是換過杯水.
你不能混在一起.
所以肉體的基督徒和.
屬靈的基督徒其實.
不能混在一起.
我們要向著這個神聖的道路.
成為一個屬靈的基督徒.
所以.
你是被什麼充滿呢.
是被漂白水還是.
被純淨的水呢.
聖靈的充滿的意思是.
他的工作其中之一.
就是叫我們想起.

$^{481}$耶穌.
祂說的一切說話.
就是真理.
就是聖言.
所以如果你不靈修.
你不讀經的時候.
你怎麼想起.
神的話.
你不會被聖言.
聖靈充滿.
所以一個屬靈的基督徒.
他最基本的充滿.
就是充滿著聖言.
和充滿著.
聖靈的愛和引導.
我們大家都是屬靈的基督徒.
要記住.
因為神的靈.
他就住在我們裡面.
請問第三堂弟兄姊妹.
誰有聖靈在你內心的.
舉手來看看.
要很肯定.
即使你未受洗的話.
你都是有聖靈住在這裡.
有聖靈住在這裡.
你就有新的生命.
記不記得在受浸的時候.
你要填這個犯罪書.
有沒有填過.
有吧.
其實我覺得.
信了主無論五年.
十年,二十年,四十年.
都要再去填一填.
因為這裡.
其實寫了很多.
聖經所指出的一些罪.
我會請弟兄姊妹.
用藍色筆.

$^{521}$來圈出.
你信主之前.
所犯過的罪.
然後用紅色筆.
來圈出.
你現在還會有一個試探.
是會犯這些罪的.
而這些罪性.
會引導你.
行出一個罪行.
例如你貪念.
就會偷竊.
所以.
這種罪性的思想.
我們可以怎樣用替代法.
來控管呢.
我就教弟兄姊妹一個方法.
如果你每天.
都在吃這些罪果.
你一定是屬肉體的.
那些行為.
怎樣辦呢.
我們要伸出聖靈.
九個特質的一個果子.
例如.
考考大家.
如果你要摘.
仇恨的罪果下來.
你要栽種哪一個果子.
仁愛.
如果是真正的一個罪果.
和平.
然後貪戀世俗.
節制.
我聽到了.
你們比第二課慢一點.
他們答得很快.
你們要再多練習.
再更多更多的練習.
所以如果你只是.

$^{561}$摘了一個罪果下來.
你可以種一個聖靈果子.
但你不能摘出聖靈果子.
我舉一個例子.
如果你在職場很憤怒.
例如Mary.
你很仇恨她.
經常欺負她.
你祈禱.
靈修的時候.
神啊求你幫助我不要再仇恨Mary.
其實是無效的.
因為你越是這樣祈禱.
你越是仇恨多一點.
一見到她可能會兜路走.
我們換一換替代法.
祈禱法是怎樣的呢?.
神啊.
你明白我真的做不到.
我真的很憎恨她.
但求你幫助我.
將人愛的種子.
種落在我的心田.
的好土壤裡面.
以致她靠著聖靈的幫助.
可以慢慢地成長.
當你這樣禱告的時候.
你會發現聖靈就開始做事.
擴闊你的心胸.
然後慢慢一點一滴.
人愛的果子.
阿加皮.
就會這樣成長.
到最終有一天.
聖靈會幫助你.
可能再次在公司見到她.
你跟她說.
Hello 早晨 耶穌愛你.
你要相信是做到的.
聖靈可以幫助你做到.

$^{601}$八章十一節.
是很重要的一節.
我們一起來讀一讀.
(聖經).
必死的身體.
和活過來.
是十一節至十七節.
很重要的主題.
所以我需要去講一講.
就是必死的身體.
必死的身體.
就是因為罪的咒詛.
每人都有一死.
但是耶穌基督.
祂從死人中復活.
是很厲害的.
一個死人為什麼會復活呢.
原來就是靠著神的靈.
就是聖靈.
所以祂的能力是很大的.
大到可以叫.
已經死了的屍體.
復活過來.
我們見證過拉薩路.
也見證過耶穌基督.
所以在安息禮拜裡.
當我們見到一些遺體.
在棺木裡的時候.
我們相信.
這位弟兄或姊妹.
有一天.
祂的身體會怎樣.
是復活的.
無論是土葬或火葬.
祂的身體是會復活的.
因為我們相信.
裡面的聖靈原來是這麼大力的.
所以.
這個是指著.
將來身體復活的泡沫.

$^{641}$但是今天.
今天上帝復活.
叫人在死人中.
復活的大能.
今天可不可以.
一樣.
去.
執行得到.
是可以的.
你不用等死.
而是你今天已經可以.
活過來.
因為聖靈.
是住在.
你的裡面.
保羅用的字眼.
是很優美.
他常常會用一些不同的字.
或者不同變化的動詞.
來讓讀書信的人.
變得有優美.
有節奏.
所以我就去查一查「住在」.
我們翻譯裡面.
是同樣這個「住在」.
是沒有錯的.
看回原文的時候很有意思.
第一次的住在.
保羅他用.
ok這個字.
ok這個字就是單單的居住.
住家.
在你決志的那天.
聖靈就住在你裡面.
第二次他講住在的時候.
enokutos.
enokutos的意思.
是這個聖靈.
是永遠不會離開你的.
永遠不會離開.

$^{681}$住在你的裡面.
是定居的.
所以這個定居.
你未來身體會復活.
但今天.
他可以幫助你.
將身體的惡行.
治於死地.
基督復活的大能.
讓我能夠治死身體的惡行.
我們讀一次.
基督復活的大能.
讓我能夠治死身體的惡行.
講一下身體.
其實這裡保羅.
他有一個很深刻的意思.
就是身體的惡行.
其實就是來自我們剛才的思想.
是不是.
如果將剛才那些.
罪性.
成為罪行.
我可以簡單的歸納概括.
其實你的身體.
你的五官.
手手腳腳.
我們都會犯罪.
而歸納之後.
你會看到耶和華.
恨惡身體的惡行.
總共有多少樣.
七樣.
高傲的眼.
說謊的舌頭.
流無辜人血的手.
圖無惡事的心.
快跑行惡的腳.
說謊的假證人.
這些全部.
都是罪的律那邊.

$^{721}$罪的思想那邊.
以致我們.
如果有這些.
思想的時候.
我們就會行出.
身體的惡行.
保羅說要治死他.
怎樣治死.
就是靠聖靈的能力.
你要想起保羅.
所處的世代是什麼世代.
是一個.
羅馬的時代.
當時對基督徒的逼迫.
還未算是很激烈.
但在他寫完羅馬書之後.
那幾年就開始變得激烈.
而且他要坐牢.
最後他死在哪一個.
王的底下.
羅馬王.
尼祿.
現在AI很厲害.
將尼祿復刻成這樣.
這些樣子.
在大英博物館裡拍攝.
歷代的一些.
羅馬的皇帝.
極殘酷的皇帝.
來迫迫基督徒.
迫迫教會.
他們犯了.
流無辜人的血.
大家都會聽過.
尼祿很喜歡.
有些紀錄.
說他將基督徒綁在.
後花園的樹裡.
一把火燒.
好好娛樂.

$^{761}$好好看.
放煙花.
是這麼殘酷的.
其他歷任的羅馬皇帝.
起碼在初期教會歷史裡.
是有十次教會的大迫害.
這些就是身體的惡行.
流無辜人的血.
但主耶穌.
他鼓勵我們.
那些殺身體不能殺靈魂的.
不要怕他們.
因為當時.
羅馬他們是要.
基督徒去拜他們是神.
沒可能的.
我們只是拜.
獨一的真神.
所以他們就會受到這些殺害.
但主耶穌說.
唯有把身體和靈魂都滅在地獄裡.
怎樣?.
淨要怕他.
就是神.
我們面對.
很多的試煉.
逼迫.
特別當時的一些信徒.
當時羅馬書的信徒.
他還沒遭遇那麼多逼迫的時候.
其實保羅已經.
預備他們.
一種受苦的心智在裡面.
將來你們會遇到.
很多的逼迫.
將來你們會有很多的.
讓你很害怕的一些試煉.
去發生.
但保羅這裡就為他們打預防針.
就是說.

$^{801}$你們的身份是什麼?.
十四至十七節不停重複.
是神的兒女.
你有神兒子的名份.
有聖靈的.
內注的意思是.
我們不再懼怕.
而且稱呼神.
叫做什麼?.
是很親密的關係.
好像一個爸爸.
牽著一個小朋友.
小朋友很多事情做不到的時候.
親自叫爸爸.
或者是老爸.
或者是爸爸.
來求助於爸爸.
我做不到.
但爸爸是做到的.
最後我分享一個見證.
就是我去德國教會.
去服事的時候.
我最後一天服事完的時候.
他們是.
動起了這個十字架.
我很開心.
在最後一天能夠見到.
因為我也要走了.
我就拍下這張照片.
在橫額上寫著.
Why Jesus Lived.
就是因耶穌.
活著.
並且下面寫著.
我就是.
道路真理生命.
耶穌就是道路真理生命.
祂到今天仍然活著.
德國的教會.
他們很著重.

$^{841}$是.
將耶穌基督的說話.
怎樣去實踐出來.
我去他們教會.
去講道的時候.
我就見到那面牆.
十字架.
寫的經文.
就是關於聖賢的.
因為耶和華的言語正直.
凡他所做的.
盡都誠實.
耶和華有一個榜樣.
就是他的言語怎樣.
是正直.
以至耶和華的作為.
是很正直的.
同樣今天.
我們都是來到.
行上帝的話的時候.
尊稱祂為我們父親的時候.
我們就好像祂的樣式.
我們去.
尋求祂的言語.
然後尋求聖靈的能力.
然後我們的行為.
乘聖的路上.
我們都是正正直直.
公公義義的.
然後我走出街的時候.
你猜一下.
我就會踩到那個牌.
我經常在德國都會半跌.
那些磚.
經常都凸起.
這一塊你知不知道是什麼來的.
它就是叫做.
半腳石.
但是它是.
德國藝術家發起的.

$^{881}$一個紀念的項目.
是由1995年開始.
全歐洲鋪設超過10萬塊.
是用來紀念.
納粹時期的受害者.
特別是猶太人.
他們不是隨便放的.
而是陷入人行道上.
或者最特別是.
陷在受害者生前.
他最後住所的門口.
而牌上面.
通常是包括著.
他的名字.
出生年日.
被捕和處決的日子.
和在哪一個監獄裡面.
是象徵著納粹的迫害.
為什麼要擺出來.
這些牌.
是由當地的居民.
或者家屬去資助的.
和鋪設之前.
是需要研究.
受害者的一個生平的.
這個牌的意義.
其實.
不是半腳石讓人摔倒.
而是.
藍色字那裡.
是半倒良心.
讓行路的人.
看到的時候.
會想起歷史.
避免我們去遺忘.
習慣了去遺忘.
我們不去遺忘.
我們不去遺忘的時候.
我們就不會再犯了.
在德國裡面.

$^{921}$他們有一個習慣.
就是會邀請學校的.
老師 同學.
或者是社會各界的.
律師.
社工 醫療.
消防.
警察等等.
全部人他們要走這條.
半腳石路線.
每一個人都要學習.
去記住.
這些歷史.
當他們邊看邊讀故事的時候.
或者我們自己去走的時候.
下次去德國.
你也不妨去走走.
這些牌子提醒我們.
歷史不是書本.
而是腳下的石頭.
這些石頭.
是有它的教訓.
叫我們不遺忘.
主耶穌說.
凡不因我半倒的是有福的.
意思是將來.
或者我們.
都會受苦.
身體上受苦.
或者別人給你的傷痛.
或者是你面對.
任何的一些試煉.
弟兄姊妹很容易.
會放棄.
不相信耶穌.
我在浸禮班裡面.
都會有一些受了浸.
最後離開耶穌.
因為勝不過那些苦難的試煉.
所以耶穌說.

$^{961}$凡不因我半倒.
是有福的.
見到這些試煉的時候.
就好像羅馬的信徒一樣.
要有受苦的心智.
因為羅馬信徒.
他們以信心著稱.
我們雖然沒有.
這些黃色的牌子.
但是我們有浸禮牌.
你受浸的時候.
有沒有收過這些賞架呢.
我就收過一個.
我到現在都保留著.
他說我已經與基督同釘十字架.
現在活著的.
不再是我.
乃是基督在我裡面活著.
並且我如今在肉身活著.
是因信神的兒子而活.
他是愛我.
為我捨己.
在剛才十一節.
到十七節裡面.
致死身體.
然後還有一個字.
是不停重複.
是活過來 活過來 活過來.
你已經.
在罪的鎖鏈.
綑綁裡面釋放了出來.
你是一個活的.
新造的人.
你會不會將耶穌基督.
活出來.
活出他的形象來呢.
保羅.
他戴著鎖鏈.
但他不是戴著.
他自己罪的鎖鏈.

$^{1001}$他是為主.
受苦.
他是戴著這條鎖鏈.
但是.
戴鎖鏈福音使者.
保羅他生命的宣言.
是因我活著.
就是基督.
在我教肥勒比書的時候.
我和弟兄姊妹分享了.
這個圖畫.
保羅之前是追捕基督徒.
去敵對耶穌.
當他剛信耶穌的時候.
他要重新學習.
怎樣成為一個屬靈的基督徒.
有一個開始.
好像我們這樣.
我們自我很大.
什麼都是me me me.
什麼都要去找.
很多自己的東西.
但是慢慢慢慢.
當基督在我們生命裡面.
顯大顯大.
其實顯大的意思.
有點像我們觸摸手機.
不停顯大.
有時我看一本聖經.
我覺得老花眼.
我都好像習慣性.
要將它擦大擦大.
直至基督完全.
在我的生命裡面.
去彰顯出來.
這個就是.
活著就是基督的意思.
而當我們.
活著為了基督.
原來可以成為我們.

$^{1041}$喜樂.
深層的原因.
今天我不知道你被.
什麼肉體的思想佔據.
而這些肉體的思想.
是讓你更喜樂.
還是更空虛呢.
你會不會已經.
很習慣很習慣.
去做一些事情.
或者是很習慣.
想都沒有想.
就是這樣去做.
但下一次你再想多一步.
我裡面有聖靈的律.
好像一個手把一樣.
當你有習慣性地.
有聚性的思想的時候.
你拉一拉這個手把.
聖靈的律.
可以幫助你.
是將耶穌基督顯為大.
我們一起來.
讀.
這個主旨.
以致我們可以記住.
上帝.
祂對我們的託付.
在基督裡.
我是神的兒女.
專注聖靈的思想.
致死身體的惡行.
為主而活.
顯主為大.
同得榮耀.
阿們.
將榮耀歸給我們的主耶穌基督.
\newpage



\section{馬可福音 5:21-43}
\label{sec:032316}
\textbf{兩極連結於基督,絕望中的信心轉向}~~~~李成傑傳道
\newline
\newline
連結: \href{http://www.church.com.hk/acms/content.asp?site=cdc\&op=show\&type=product\&code=032316\&layout=sermon}{\texttt{032316}} ~~~~ 語音日期: 2026-06-14
\newline
\newline
\hyperref[sec:032315]{\small{< < < PREV SERMON < < <}}
~
\hyperref[sec:index_chronic]{\small{[返順時目]}}
~
\hyperref[sec:032317]{\small{> > > NEXT SERMON > > >}}
\newline
\newline
馬可福音 5:21-43
\newline
\begin{longtable}{cl}
\hline
\hline
章節 & 經文 (和合本修訂版)\\
\hline
5:21 & \begin{tabularx}{0.7\textwidth}{X} 耶穌又坐船渡到對岸，有一大群人聚集到他身邊；他正在海邊。 \end{tabularx} \\ \\ \relax
5:22 & \begin{tabularx}{0.7\textwidth}{X} 有一個會堂主管，名叫葉魯，也來了，一見到耶穌，就俯伏在他腳前， \end{tabularx} \\ \\ \relax
5:23 & \begin{tabularx}{0.7\textwidth}{X} 再三求他，說：「我的小女兒快要死了，求你去為她按手，使她痊癒，可以活下去。」 \end{tabularx} \\ \\ \relax
5:24 & \begin{tabularx}{0.7\textwidth}{X} 耶穌就和他同去。有一大群人跟隨他，擁擠著他。 \end{tabularx} \\ \\ \relax
5:25 & \begin{tabularx}{0.7\textwidth}{X} 有一個女人，患了經血不止的病有十二年， \end{tabularx} \\ \\ \relax
5:26 & \begin{tabularx}{0.7\textwidth}{X} 在好多醫生手裡受了許多苦，又花盡了她所有的，一點也不見好，反而更重了。 \end{tabularx} \\ \\ \relax
5:27 & \begin{tabularx}{0.7\textwidth}{X} 她聽見耶穌的事，就夾在眾人中間，從後面來摸耶穌的衣裳， \end{tabularx} \\ \\ \relax
5:28 & \begin{tabularx}{0.7\textwidth}{X} 因她想：「我只摸到他的衣裳，就會痊癒。」 \end{tabularx} \\ \\ \relax
5:29 & \begin{tabularx}{0.7\textwidth}{X} 於是她的流血立刻止住，她覺得身上的疾病好了。 \end{tabularx} \\ \\ \relax
5:30 & \begin{tabularx}{0.7\textwidth}{X} 耶穌頓時心裡覺得有能力從自己身上出去，就在眾人中間轉過來，說：「誰摸我的衣裳？」 \end{tabularx} \\ \\ \relax
5:31 & \begin{tabularx}{0.7\textwidth}{X} 門徒對他說：「你看眾人擁擠著你，還說『誰摸我』呢？」 \end{tabularx} \\ \\ \relax
5:32 & \begin{tabularx}{0.7\textwidth}{X} 耶穌周圍觀看，要見做這事的女人。 \end{tabularx} \\ \\ \relax
5:33 & \begin{tabularx}{0.7\textwidth}{X} 那女人知道在自己身上所成的事，就恐懼戰兢，來俯伏在耶穌跟前，將實情全告訴他。 \end{tabularx} \\ \\ \relax
5:34 & \begin{tabularx}{0.7\textwidth}{X} 耶穌對她說：「女兒，你的信救了你，平安地回去吧！你的疾病痊癒了。」 \end{tabularx} \\ \\ \relax
5:35 & \begin{tabularx}{0.7\textwidth}{X} 耶穌還在說話的時候，有人從會堂主管的家裡來，說：「你的女兒死了，何必還勞駕老師呢？」 \end{tabularx} \\ \\ \relax
5:36 & \begin{tabularx}{0.7\textwidth}{X} 耶穌不理會他們所說的話，就對會堂主管說：「不要怕，只要信！」 \end{tabularx} \\ \\ \relax
5:37 & \begin{tabularx}{0.7\textwidth}{X} 於是他帶著彼得、雅各和雅各的弟弟約翰同去，不許別人跟著他。 \end{tabularx} \\ \\ \relax
5:38 & \begin{tabularx}{0.7\textwidth}{X} 他們來到會堂主管的家裡，耶穌看到一片吵鬧，並有人大聲哭泣哀號， \end{tabularx} \\ \\ \relax
5:39 & \begin{tabularx}{0.7\textwidth}{X} 就進到裡面，對他們說：「為甚麼大吵大哭呢？孩子不是死了，是睡著了。」 \end{tabularx} \\ \\ \relax
5:40 & \begin{tabularx}{0.7\textwidth}{X} 他們就嘲笑耶穌。耶穌把他們都趕出去，帶著孩子的父母和跟隨的人進了孩子所在的地方， \end{tabularx} \\ \\ \relax
5:41 & \begin{tabularx}{0.7\textwidth}{X} 就拉著孩子的手，對她說：「大利大，古米！」翻出來就是說：「女孩，我吩咐你，起來！」 \end{tabularx} \\ \\ \relax
5:42 & \begin{tabularx}{0.7\textwidth}{X} 那女孩子立刻起來走動—她已經十二歲了；他們就非常驚奇。 \end{tabularx} \\ \\ \relax
5:43 & \begin{tabularx}{0.7\textwidth}{X} 耶穌切切地囑咐他們，不要讓人知道這事，又吩咐給她東西吃。 \end{tabularx} \\ \\
[1ex]
\hline
\hline
\end{longtable}
$^{1}$今天分享的這段經文是在馬可福音5:21-43.
但是如果我們很熟悉聖經的時候.
我們會知道這段經文的其中一個故事.
艾魯的女兒快要離開的時候.
這段故事也會在路加福音記載.
但是為什麼在福音書裡面.
雖然是同一段經文.
但是馬可的技術或多或少有些許不同.
並且要將12年患血流病的女人這件事.
很黏連地放在這裡.
是有它的意思在當中.
我在這裡起了一個題目.
兩極黏於基督就是一個方華正茂.
不應該說是方華正茂.
應該說是一個錢包待放的少女.
她去著她父母的希望.
一直進化的時候突如其來.
人生裡有一個風暴.
她也不知道要面臨死亡.
同樣地一個本來是方華正茂的女性.
因著12年患血流病.
她就圓華盡逝去到絕望當中.
但是無論身處在什麼處境也好.
她們都在同一個點遇見耶穌.
其實我們的生命真正需要什麼呢?.
我們以為自己很清楚.
我們生命需要什麼.
或者信仰當我們祈求上帝的時候.
我們以為自己很清楚我們需要什麼東西.
但其實我們真正的需要或者生命需要的東西.
我們未必知道.
但是透過禱告.
透過有很多種種的事件.
與耶穌基督相遇的時候.
我們本來只是去到這個地步.
但是當我們與耶穌相會的時候.
耶穌就會帶我們進入不同的地步.
同樣地在聖經裡面.
這兩個主角可能他們的生命只是想著好起來.
又或者能夠讓他回復某些地位.

$^{41}$但是耶穌希望他的生命.
不只是限制於他所想的.
而是上帝讓他可以開闊他的眼界.
這段經文《馬可福音》第五章.
與之前第四章.
耶穌到理海,平靜風和海.
和去到格拉森人的地方.
趕走群鬼的一段經文息息相關.
因為有很多動詞,有很多字句.
都是很連接的.
為什麼我這樣說呢?.
第一就是祈求,央求這個字.
因為當耶穌去到格拉森人的地方.
鬼就央求耶穌離開.
同樣地那群格拉森人也央求耶穌離開他們.
所以「求」這個字.
就是串連在第四章,第五章.
而那個祈求.
他們也不知道為什麼耶穌基督在他們身上作出好事.
這群格拉森人要求耶穌離開.
所以鬼也求耶穌離開.
人也求耶穌離開.
你說是不是很有趣呢?.
同樣地.
那群人看到格拉森人被鬼附的人.
他們心裡明白了.
即是回復正常.
而這個被耶穌治好的人.
這個從前被鬼附的人.
他也求耶穌.
他求耶穌做什麼呢?.
求耶穌讓他可以跟從祂.
但是耶穌也不依照著祂自己的心願.
要讓他有更大的使命.
是怎樣呢?.
回到他家裡.
因為他家裡已經很久沒有擁有這個兒子了.
並且這個常住在格拉森的墳營的地方.
後來聖經的《馬可福音》.
怎樣記載呢?.

$^{81}$他從這個墳營這個好笑的地方.
走遍底加玻利.
底加玻利原本的翻譯就是十座城.
所以你會看到一樣很重要的事情.
就是整段聖經有一個訊息.
就是從一個很小的地界.
從一片很小的信心的土地.
可能可以擴展自己信心的地界.
所以這種無論是字句.
或者是祈求裡面的詞彙.
來到這裡第五章.
我們發覺一個很重要的事情.
第二十二節說.
有一個管會堂的名叫艾魯來見耶穌.
就伏伏在他的腳前再三求他.
再三求他.
就是說他好像那些祈求的上一段一樣.
他再三求耶穌.
在我們的生命裡面.
我們也祈求上帝很多很多次.
或者我們心裡面有些願望.
也希望求上帝可以改變.
但是這個再三呼求.
耶穌基督沒有立時去應允艾魯.
究竟是在耶穌離開平靜峰和海之前.
艾魯已經求他.
還是耶穌從格拉森那邊.
從此岸回到彼岸.
從彼岸回到此岸的時候.
艾魯再三求他.
究竟是A還是B.
但無論如何.
艾魯再三求他.
他說:我的小女兒快要死了.
求你按手在她身上使得痊癒.
得以活了.
整段經文給我們看到有一個很重要的背景.
因為艾魯這一刻也打定輸數.
為什麼呢?.
剛才我們讀經文的時候.

$^{121}$不是有一群人在當中嗎?.
其實那群是特別請回來.
為已死的人哭泣,哀號.
所以其實艾魯本著一個希望與絕望之間.
求耶穌到他家裡讓他的女兒得以痊癒.
在艾魯的祈願裡.
他將自己或者對耶穌的身份.
定性為一個能夠使人痊癒的醫生.
並不是一個使人復活的主.
而我們看到一個很重要的事實.
艾魯在這個地位.
或者他生命的處境裡.
他比稍後一瞬間出現血漏病的女人優越些.
有一個社會學家叫做內布爾.
他講的是一種道德幸運.
他告訴我們人之所以能夠領取道德的責任.
或者他有一些優越的主權.
他們以為他們自己很take for granted.
很知道這些是當然的.
但是如果細心一想.
我們會看到明顯地艾魯在這裡.
是有一種道德的幸運或者自己的優越感.
第一 他是管會堂的.
管會堂不是我們平時想像的那種.
在大廈的管理人.
而是在會堂裡的一個很重要的主管.
地位是相當之高.
所以他能夠去到這個海邊.
你要注意 聖經說的海邊.
能夠去到耶穌的前頭的位置.
是因為他自己本身.
無論在社會上.
無論在宗教地位上.
甚至在出身上.
都能夠去到那個位置.
可以接近耶穌.
等於有時候我們要反思.
我們今天站在這個恩典當中.
或多或少都是因為.
在上帝未揀選我們以前.

$^{161}$已經有某一種的社會給我們.
無論是學歷或者環境因素.
給予我們一些位置.
令我們可以很容易得到一些東西.
我早前說了一樣的經歷.
我想一個月前有林憶蓮演唱會.
大家有沒有聽過.
我在Facebook看到我以前服侍的.
那班弟兄姊妹在呼喚.
今晚有林憶蓮演唱會的門票.
有人要嗎? 舉手說要.
挺好的 我就想要.
我想弟兄姊妹當我回應時.
他們會送門票.
其實不是送門票的.
當你要的時候.
他們說Sam 我私信你.
以前教會 以前在港福堂有幕會.
門票大概是二千元的價錢.
我一聽到那個價錢.
我都想著給錢的.
不會想著送門票那麼簡單.
我就嚇了一跳.
我就問幾點鐘開始.
因為我真的不太知道.
好像六點鐘就開始.
哎呀 不好意思.
我真的趕不到.
因為我在聯絡部要做頒獎嘉賓.
那時剛好五點九就要去沙田頒獎.
就去不了了.
我就輸了一口氣.
不是那麼丟臉.
因為二千多元看一場演唱會.
真的被我老婆罵嗎.
又要帶她去看.
應該是雙倍的價錢 價錢很貴.
她就說不好意思 因為這些全部都是真的.
但上帝又讓我講到這些說話.
頂智媒 其實很有趣.

$^{201}$頂智媒在那個層次裡面.
她是不察覺的.
她不察覺這是怎麼回事.
因為那缸水泡了很久.
沒什麼 很平常.
我記得我以前服侍她的教會.
有個阿姨就跟我說.
Sam 我前陣子又撞了一輛賓士.
現在的賓士很便宜 幾十萬一輛.
幾十萬一輛很便宜.
不是的 其實被謝霆鋒撞的法拉利.
都便宜了一截.
幾十萬撞一輛都會便宜.
但我開頭都不是很習慣.
但她在那缸水泡上 習慣了去講.
如果切身處地去講.
我們有時候講一下換電話.
或者我們講一下吃一下兩sung 飯.
四十八塊 四十八塊都可以.
但你在某一個區域.
這個價錢對某些人來說是難的.
所以我們要知道.
為什麼我講這麼多的例子.
因為艾佬他不知道他自己本身天賦.
或者上帝賜予他一些道德優越的地位.
以致他能夠在這麼前列的位置.
再三央求耶穌.
同樣地 在經文那裡.
後頭來的那一個女人.
和前頭能夠求耶穌那裡.
成為了一個很強烈的對比.
因為宗教的環境.
因為社會的地位的緣故.
能夠給我們某些的地位.
以致我們能夠履行一些社會的責任.
道德的責任很治愈.
但相反地.
一會兒在經文裡面所呈現的那個患十二年血漏病的女人.
不能夠像艾佬那樣.
可以有這個位置.

$^{241}$和可以有這個優越的位置.
可以央求到耶穌.
讓我們看到一件事.
弟兄姊妹.
無論我們身處什麼什麼的境況的時候.
我們都需要從一個視角.
讓我們看到.
總有感恩的緣故.
總有上帝給我們的恩典的緣故.
而艾佬.
就是這個在生命裡面.
處於優越的地位.
而能夠可以在心求耶穌.
但是在耶穌答應他.
去到他的家裡.
去醫治他的女兒的時候.
在第二十四節.
耶穌就和他同去.
有很多人跟隨.
擁擠他.
這個擁擠.
是阻礙了耶穌.
去到他的家裡.
可以及時.
可以醫治他的小女兒.
復原的.
因為這個擁擠.
成為了他的生命裡面.
一個困局.
他是沒有辦法預計到的.
很多時候我們都希望.
能夠作出一個.
很好的計劃.
但是擁擠.
甚至突然間殺出一個人.
成為了一個.
停止的.
一個記號.
或者一個立腳點.
是艾佬千千萬萬想不到的.

$^{281}$頂智媒如果你讀這本經文的時候.
最心急的是誰?.
是那個患十二年血流病的女人?.
還是艾佬心急?.
應該是艾佬最心急.
為什麼?.
不要再擁擠了.
不要再阻礙耶穌了.
為什麼?.
因為我的小女兒.
那個時間不夠了.
但是這個時候.
殺出一個程咬金.
說十二年患血流病的經歷.
很奇怪.
門徒說的那些話.
是很合理的.
等同我們.
去旺角想著很心急.
趕那班車.
逼來逼去.
去到一個地步趕不到那班車.
有時候你等巴士很有趣.
你平時在巴士站那裡沒有人.
剛剛巴士來的時候.
突然間那些人殺出來.
可以排隊等巴士.
特別是在夏天的時候.
你有沒有試過這樣的經歷?.
不知道為什麼.
站沒有人.
剛剛巴士來的時候.
突然間那些人殺出來.
在車站那裡.
很有趣.
不過可能現在有APP.
計算好時間出去.
總而言之你很急的時候.
特別覺得是擁擠.
但這個不是你感覺的情況.

$^{321}$是一個很實際的情況.
但是耶穌基督.
他走的時候.
還要.
要患12年血漏病的女人.
出來說她患12年.
血漏病的見證.
我們有時候說見證.
我們都需要很好的時間去管理.
但她患12年.
血漏病的女人.
一年兩年三年四年這樣說.
艾魯就是告訴你.
不行你找那個醫生.
然後告訴他.
找那個醫生12年.
怎樣去虎方的時候.
拜託你快點說完吧.
艾魯就在這個心情.
有生命裡面的.
一些疾谷.
點滴.
弟子們.
在這個.
時候裡面.
上帝給我們.
有學一個功課.
原來當我們.
祈求的時候.
在這個時刻.
成為了.
我們生命裡面.
一個很神聖的停頓.
艾魯是不會察覺.
並且門徒是不知道的.
為什麼?.
因為在這一刻.
門徒的生命.
和患12年.
血漏病的女人的生命.

$^{361}$是很有.
不同的改變.
他這個停頓.
能夠述說.
他過往的日子裡.
那種不堪.
也可以在公開的場合.
讓.
所有眾人看到.
耶穌基督的能力.
因為在.
聖經的記述.
門徒是從來都不知道.
原來.
祈求的時候.
耶穌的默然.
能力可以從.
他的身上去到.
這個女人的生命當中.
有時候.
在我們的生命裡面.
我們以為.
一些大鑼大鼓.
一些很鏗鏘的說話.
可以幫助我們.
生命的成長.
但是原來在這個.
神聖停頓的時刻裡面.
耶穌基督的.
另一面.
或者耶穌基督的能力.
可以讓門徒知道.
門徒那一刻才知道.
原來默然的禱告.
上帝是會透過.
這種的靜默.
這種的默然.
是可以進到.
祈求者的生命裡面.
我有個老師.

$^{401}$叫做Rachel Miller.
她寫了.
一份博士論文.
只是研究聖經.
其中的一個字.
那個是什麼字呢?.
就是寂靜這個字.
寂靜寫一份博士論文.
是在乎你如何去堆砌.
是不是?.
她說這個寂靜在哪裡出現呢?.
她說在她年少的時候.
年齡少女的時候.
她讀聖經讀到啟示錄.
她說那些天使.
剪開書卷.
揭開神的話語的時候.
天上寂靜若有.
若有片刻.
她就很奇怪.
究竟這個寂靜片刻.
是創造了出來呢?.
還是.
是因為被造物.
有些動作.
而令到那一刻.
是寂靜呢?.
她得出的答案.
原來是.
被造物.
在那一刻.
當上帝的話語揭開的時候.
進行一個敬禮.
指示到天上.
寂靜若有片刻.
所以.
在這個.
「永濟誰摸我的衣裳」.
在這個停頓裡面.
是讓我們的生命停頓.

$^{441}$讓上帝的說話.
透過不同的方式.
使我們的生命.
有所不同.
有所提升.
這個提升是甚麼呢?.
從前.
這個.
患十二年血漏病的女人.
因為宗教的緣故.
與很多很多.
上帝的東西.
有很多很多.
與神的東西.
都隔絕的.
因為在舊約的聖經裡面.
告訴她一件事.
血漏病是一個不潔的.
不潔的災病.
要找祭司察看.
痊癒後才可以履行宗教任務.
所以.
這個女人的生命.
是一個隱匿的生命.
不能夠公開.
一些場合裡面.
亦不能夠.
好像我們的生命.
那麼自如地去到任何地方.
去行走.
但是耶穌基督.
透過呼喚.
呼召這個女人.
三十三節.
那女人知道自己身上所成就的事.
就恐懼戰驚.
來苦伏在耶穌面前.
將實情全告訴他.
加拉太書這樣說.
當你接觸了耶穌基督.

$^{481}$從此就不是.
受奴僕的心仍夠害怕.
你是有.
這個兒女的心.
所以這個女人的生命.
本來她的祈願.
只是希望.
能夠很安靜地.
或者很生命地.
來尋找那個信仰.
只要摸耶穌的衣裳.
她就可以慢慢地.
去恢復.
她自己.
原本的地位.
但是耶穌基督.
讓她的生命的提升.
不是自怨自艾的女人.
在三十四節這樣說.
耶穌對她說.
女兒,你的信救了你.
從女人.
去到女兒.
從這個女人.
這個被拒.
被拒諸門外.
甚至乎生命隱匿的.
這個女人.
驚嚇之間因為這個信心的緣故.
被提升到女兒.
弟兄姊妹,我們要想.
究竟我們.
在這個信仰裡面.
我們仍是.
跟耶穌基督.
很遠嗎?.
還是我們已經知道.
我們的生命已經被.
主所提升成為.
神家裡面的一個人.

$^{521}$從一個女人進到.
女兒的地步呢?.
好.
當這個十二年脫漏病的女人.
作完她的見證.
我們都不知道.
花了多少時間.
經文裡面.
有句說話.
三十五節.
還說話的時候.
有人從.
管會堂的家裡來說.
你的女兒死了.
何必還勞動先生.
很有趣的一段經文.
亦都是馬可將心.
獨蘊地說.
還說話的時候.
對應一直.
形容艾魯是一個.
管會堂的人.
管理時間.
管理會堂.
一切一切的程序.
甚至乎管理.
整個宗教的制度.
管會堂.
跟我們香港人.
本身的生活作息.
很息息相關.
因為我們很懂得管理時間.
懂得管理一些東西.
但是.
三十五節.
還說話的時候.
對我們這班懂得管理的人.
一個很大的諷刺.
無論你怎樣懂得.
管理都好.

$^{561}$時間都並不是你能夠掌控的.
一環.
而還說話的時候.
對一個管會堂的人來說.
是生命裡面一個.
很重要的一個.
孝敏一個重擊.
我如何管理得好.
我如何可能安插很多人都好.
總而言之.
時間和計劃.
都不被我生命所管理.
另外一樣.
三十五節.
你女兒死了.
何必還勞動先生.
勞動希臘文這個字.
在路加福音.
同樣都出現.
為什麼你會問我.
Sam 你在說馬可福音.
為什麼出現同一個字.
又要在路加福音說.
因為.
我女兒這個.
記載.
同樣地.
在路加福音有記載.
而路加都是用勞動那個字.
而勞動那個字.
在路加福音都同樣地出現過.
這個勞動.
這個字怎麼出現呢.
是在.
伯夫長.
去求耶穌.
伯夫長在路加福音說了什麼.
他說 耶穌.
我家裡的僕人.
病了.

$^{601}$求你.
醫治他.
然後耶穌就說.
好啊 我去到你家裡醫治你的僕人.
你記不記得伯夫長做了什麼事.
不用.
我不用勞動.
耶穌你.
去到我家裡去醫治.
我的伯夫長 我的僕人.
他說.
你只要吩咐就可以了.
你吩咐就可以.
然後耶穌就讚伯夫長的信心大.
但是.
如果這個勞動的字在.
伯夫長說不用勞動先生.
和這個.
第35節裡面.
勞動先生那一拼.
你就很明顯看到.
艾魯他的生命.
他裡面.
和伯夫長的生命裡面.
有天淵之別.
因為伯夫長知道.
上帝一吩咐他的僕人.
就會立事好了.
而伯夫長.
他將他自己看耶穌基督的地位.
是更加比他高一層次.
因為我都做到了.
你更加是.
我的主.
但是.
管會堂 艾魯 管時間.
管程序.
管什麼.
管什麼東西.
甚至乎在他的生命裡面.

$^{641}$勞動這個字.
連上帝都可以.
管理.
我們值得反思.
是.
我們這位主.
管理我們.
還是.
上帝要我們.
再去提升一個層次.
我才是.
你的主.
我是時間管理的主.
所以很明顯地.
不要怕.
只要信.
艾魯在他的生命裡面.
他將耶穌基督.
定性為一個醫者.
很像是.
要醫治.
血流病女人的醫者一樣.
但是耶穌基督.
透過不同的經歷.
將這個.
十二年患血流病的女人.
的層次提升.
從女人成為女兒.
他的生命得到改變.
得到提升.
同樣地艾魯也透過這件事.
他的生命能夠得到提升.
他再不是.
操控或者管理.
時間的一人.
而是真正他相遇的.
是告訴他今天.
他相遇的不是在他生命裡面.
一個醫治人的耶穌.
而是一個復會主.

$^{681}$並且是管理時間掌控的.
這位上帝.
所以兩個.
不同的生命.
透過接觸耶穌基督.
連結於耶穌基督同歸於一.
在這個.
馬可福音裡面.
在第.
四十二節那裡.
剛才我們讀經文.
「那閨女立時起來,她們就大大的驚奇,.
閨女已經十二歲了.」.
在這個.
《路加福音》裡面記載.
很簡單.
一開始就介紹了.
這個女孩.
這個閨女是十二歲.
但是為什麼馬可將.
十二歲那裡擺到最後.
是成為了整段經文的首尾呼應.
為什麼呢?.
因為.
一個含苞待放.
或者將會.
亭亭玉立的十二歲的.
女孩子.
帶著希望.
帶著快樂.
正如弟兄姊妹有時候在我們生命裡面.
我們有很多愉快的時刻.
無論十年,二十年,三十年.
四十年都好.
但是馬可.
希望透過這個十二年.
擺到最後.
讓我們突然間醒覺.
當我們愉快的那十二年.
原來在平行時空底下.

$^{721}$有一個哀傷的十二年.
是我們不知道的.
哀傷的十二年在哪裡出現呢?.
就是在十二年血漏病的女人裡出現.
讓我們知道.
在教會.
這扇門裡面.
我們可能.
是很感恩地.
愉快地過著.
但是原來.
在這個教會外面.
當我們很快樂.
或者覺得很信心的同時.
也會在外面.
有很多不稱心的人在當中.
正所謂.
幸福的是一樣的.
但是不幸的是.
各有各的不幸.
到了最後.
四十二節.
那裡同樣地.
告訴我們一個很重要的事實.
為什麼一直叫做女孩.
後來就叫到.
四十二節那裡說.
那閨女立時起來.
她就大大的驚奇.
閨女已經十二歲.
然後就說.
閨女我吩咐你起來.
原來耶穌基督.
透過.
馬可.
這個寫作.
是叫那一群.
或者在艾路.
隱蔽底下的.
在所有.

$^{761}$宗教保護.
底下的信徒.
能夠走出去.
起來.
去看.
這個世界裡面.
一些需要.
正如波福堂的.
一個主題.
要走出去.
傳遞外.
就是閨女要起來的一個訊息.
在結束的時候.
讓我有一點點滴.
如果我們熟悉.
馬可這個人.
他為什麼要這麼有心思.
或者聖靈為什麼要感動他.
要去寫.
這段經文.
讓我們透過他的寫作.
透過聖靈的安排.
讓我們看到生命.
不同層次的感動.
馬可.
本身他是一個.
有錢仔.
曾經被保羅拒絕.
因為某些緣故可能宣教的時候.
很辛苦他離開了行列.
後來他的表哥.
巴拿巴都因為.
馬可他的表弟.
和保羅爭吵.
但是到了後來.
其實馬可他的生命.
不是像保羅一樣.
進行宣教.
原來馬可.
他本身是一個神學家.

$^{801}$一個福音書的作者.
是要將上帝的.
說話很匠心毒運.
去傳講.
到了後來在保羅的書信裡面.
當他臨走的時候.
回到上帝的身邊.
都要吩咐提摩泰.
叫馬可.
來到他的身邊.
並且要帶一些書卷.
所以你會明白到.
馬可為什麼要.
透過這段經文.
去講生命的不同層次.
因為馬可慢慢都猜不到.
在他生命的失敗裡面.
上帝.
這位耶穌基督.
提拔他.
要他正視原來他自己的生命.
未曾發掘的恩賜.
以致到他今天.
成就了他自己的生命素質.
去將這個馬可福音.
留下.
並且要我們注視到一件事.
原來在我們的生命裡面.
有很多未發掘不知的東西.
不只是停留在我們的腦裡.
我們所需要.
原來是耶穌基督透過不同的事件.
發掘在需要背後的.
上帝要我們知道的.
一些事情.
是我們不知道的.
並且他在創世.
已經欣賞了我們一些東西.
就好像馬可一樣.
就好像這個.

$^{841}$十二年患血漏病的女人一樣.
從女人提升到.
為女兒.
並且也要像艾魯一樣.
這個女孩.
這個閨女要走出去.
要看到這個世界不同的東西.
上帝要他所要看到的.
讓我們一起有一個禱告.
才主耶穌.
幫助我們.
有時候在我們的生命裡面.
我們都因為某些事情.
限制了我們自己的想像.
我們覺得可能過去日子裡面.
我們的生命有很多限制.
甚至乎在我們的.
生命裡面我們會覺得.
有很多被高估.
或者自視甚高的狀況.
但是上帝你透過不同的事件.
人事物去改變我們.
求你幫助我們.
好像昔日.
你怎樣去幫助馬可.
讓他能夠發掘到他自己的生命裡面.
那些很重要的養分.
以致到他能夠.
將你的說話.
好好地留給我們.
以致到我們能夠透過主耶穌基督.
你的話語.
去成為我們生命的養分.
提升我們自己.
幫助我們.
很勇敢地走出去.
看到上帝你所創造的世界.
是多麼美好.
也有多麼多的需要.
多謝主你.

$^{881}$禱告是奉教主耶穌基督的名祈求.
Amen.
\newpage



\section{以弗所書 4:1-6}
\label{sec:032317}
\textbf{在基督裡合一}~~~~歐陽萬廷弟兄
\newline
\newline
連結: \href{http://www.church.com.hk/acms/content.asp?site=cdc\&op=show\&type=product\&code=032317\&layout=sermon}{\texttt{032317}} ~~~~ 語音日期: 2026-06-14
\newline
\newline
\hyperref[sec:032316]{\small{< < < PREV SERMON < < <}}
~
\hyperref[sec:index_chronic]{\small{[返順時目]}}
~
\hyperref[sec:032318]{\small{> > > NEXT SERMON > > >}}
\newline
\newline
以弗所書 4:1-6
\newline
\begin{longtable}{cl}
\hline
\hline
章節 & 經文 (和合本修訂版)\\
\hline
4:1 & \begin{tabularx}{0.7\textwidth}{X} 我為主作囚徒的勸你們，既然蒙召，行事為人就要與你們所蒙的呼召相稱。 \end{tabularx} \\ \\ \relax
4:2 & \begin{tabularx}{0.7\textwidth}{X} 凡事要謙虛、溫柔、忍耐，用愛心互相寬容， \end{tabularx} \\ \\ \relax
4:3 & \begin{tabularx}{0.7\textwidth}{X} 以和平彼此聯繫，竭力保持聖靈所賜的合一。 \end{tabularx} \\ \\ \relax
4:4 & \begin{tabularx}{0.7\textwidth}{X} 身體只有一個，聖靈只有一位，正如你們蒙召，是為同有一個指望而蒙召， \end{tabularx} \\ \\ \relax
4:5 & \begin{tabularx}{0.7\textwidth}{X} 一主，一信，一洗， \end{tabularx} \\ \\ \relax
4:6 & \begin{tabularx}{0.7\textwidth}{X} 一神－就是萬人之父，超越萬有之上，貫通萬有，在萬有之中。 \end{tabularx} \\ \\
[1ex]
\hline
\hline
\end{longtable}
$^{1}$多謝神給我有機會繼續和大家分享.
是已忽所書.
第四章.
我的祈求是聖靈親自作為我們的導師.
而不是大家聽我的講道.
讓我們再次低頭祈禱.
祈求你安靜我們的心.
教導我們怎樣可以 怎樣來到敬拜你.
求聖靈教導我們.
做我們的導師.
讓我們明白你自己的話語.
祈禱是奉主耶穌的名字.
今天我繼續講到已忽所書.
剛才讀的經文是從第四章第一節到第六節.
今天我覆蓋的內容一直是到第十六節.
但我的中心就在第一節到第六節.
我希望大家能夠打開你自己的聖經.
如果你有帶聖經的話.
或者打開你的手機.
打開已忽所書第四章.
我們會希望大家一邊聽一邊讀.
亦都一路能夠來思考.
我將今天的題目分為四部份.
第一部份是我們想一想教會合一的意義是甚麼.
第二件事就是講到如何維繫教會的合一.
第三方面是教會合一在有甚麼不變的本質.
然後第四是講到教會如何在合一裡能夠增長.
剛才所讀的經文大家是用和合盤來讀.
在我準備今天的講道的時候.
我亦都是參考其他的聖經版本.
其中一個版本是呂鎮忠的譯本.
我覺得是非常合適在這裡跟大家講.
讓我讀一次.
「所以我為主的緣故做囚犯的勸你們.
行事為人要對得起你們所蒙的呼召」.
為何我喜歡這個譯本其中有幾個原因.
在英文的聖經裡面有「therefore」即是「所以」.
「所以」這兩個字將第四章連接起以上的經文.
就是說從第一章到第三章是一個大分段.
從第四章到第六章是另一個分段.

$^{41}$另外他提醒了我一件事.
這個第一節.
請你回到剛才的講道.
這裡講到神對我們有一個呼召在此.
而你們需要在呼召和我們行事為人.
有一個平衡在此.
「你們行事為人要對得起你們所蒙的呼召」.
即是你一邊放在神的救恩在左手邊.
另一方面你考慮一下.
你自己的行事為人是否真的對得起你所得到的救恩.
因此我希望能夠重複一下.
重點性地重複一下.
從第一章到第三章的一些重點.
因為我相信如果我們更加了解明白.
救恩對我們的重要性的話.
我們亦都會能夠相應地回應神對救恩的那種寶貴.
我們相對會懂得怎樣去回應祂.
「行事為人」是中文的一個很長的翻譯.
英文就是Walk.
就是我們信了甚麼我們應該怎樣去走.
在呼召裡面我們想到幾樣東西.
從第一章到第三章.
第一件事就是首先神在創世以初.
神就在基督裡面揀選了我們.
叫我們成為聖潔無有瑕疵.
第二件事.
就著祂的愛子耶穌基督的字.
我們得蒙救贖.
並且洗淨我們一切的不義.
第三件事.
我們每一個信了耶穌的人.
我們都會受到聖靈的印記.
這個印記是我們能夠將來得到基業的憑據.
另外一件事.
聖靈在我們做了甚麼工作.
聖靈能夠照亮我們的心靈.
讓我們知道祂的呼召.
我們的呼召所得的指望是甚麼.
而且這個基業是何等的榮耀和豐盛.
很多時候一天過一天.

$^{81}$我們似乎對我們的生命.
我們的生活都覺得很平淡.
沒有甚麼喜氣樂樂.
但當如果我們細心思考神的救恩的時候.
我們很難同時想到我們的救恩是有豐富的榮耀.
原來神給我們的救恩不是單單只是救了我們.
叫我們成為一個基督.
我們的指望,我們將來要得到的基業.
是有豐富的榮耀.
大家去思考一下.
下一個.
「神在信的人所顯的能力是何等浩大」.
這個「何等浩大」是甚麼意義.
「神的能力何等浩大」的話.
這個同樣的能力是叫主耶穌基督從死裡復活的那種浩大.
即是我們在座每一位信耶穌的人.
神叫你從死裡復活.
那種能力是和神叫耶穌基督從死裡復活的那種能力是一模一樣.
沒有改變.
得救是本福恩也是一種出生入死的經歷.
你會發覺我做了很久基督徒.
或者是想著要信耶穌的.
我經常都提醒.
重複提到一樣東西就是.
我們的救恩是出於恩典.
而不是因為我們個人的努力.
因為只有神才擁有這種權力.
將我們從一個被定罪的罪人.
轉變成為一個有永恆生命的人.
「求一母種族之分」.
剛才我也讀了一些經文.
因為主耶穌基督的緣故.
藉著祂在十字架上所附上的.
祂的身體拆毀了猶太人和外邦人.
拆毀了猶太人和外邦人相隔的牆壁.
使得無論你是猶太人還是外邦人.
在耶穌基督裡面都合二為一成為一.
所有基督徒都在耶穌基督裡面.
是一個大家庭的成員.
為什麼重要性這麼重要呢?.

$^{121}$為什麼要經常提到.
再沒有猶太人,再沒有外邦人.
為什麼這麼重要?.
因為今天我們所讀的聖經裡面.
都是很多有關猶太人的歷史.
我們得到救恩之後.
亞伯拉罕的經歷.
以撒瓦各的經歷.
他們怎樣向著神.
都成為我們今天的經歷.
是完全沒有分別的.
因為主耶穌基督在我們裡面.
我們能夠充分領悟神的愛.
只有建立一個唯一的根基就是耶穌基督.
我們能夠彼此來分享.
那種神的愛的長,闊,高深.
我們今天的問題就是說.
你們有沒有體驗到神這個大愛?.
無論我們當中得救的經歷是怎樣.
有些人很安安靜靜.
有些人空空裂裂.
一個吸毒的變成一個信耶穌.
無論是怎樣.
這個不是很重要.
最重要的是每個人都能夠認知一個同一個的訊息.
就是說.
神為了愛你付出了極大的代價.
所以保羅在第四章.
向我們有一個呼籲.
他的呼籲是甚麼呢?.
你要竭力保持聖靈所賜的合一.
很奇怪.
很多不同的書信裡面.
我們都發覺保羅會很重點性地提到一些問題.
但是以弗所書第四章.
他提到一個很重要的訊息.
就是合一的訊息.
而且他要求我們的是甚麼呢?.
他要求我們是要竭力去保持這個合一.
合一是甚麼意思呢?.

$^{161}$這個合一的意思.
並不是一個新的提出的題目.
因為在以弗所書第二章.
第一章和第二章都有提到.
有關這個合一的訊息.
在第一章那裡.
他說在日子滿足的時候.
神使天上地下和一切的.
都叫他們在基督裡面同歸於一.
在第二章十九節他說到.
他說猶太人和外邦人裡面.
在基督裡面都是同為一家人.
他用一個例子.
就是房屋的例子來說明這件事.
所以合一的概念並不是一個新的概念.
但是在第四章的開始.
他提到合一這個討論的時候.
一個勸告的時候.
帶來一種新的體驗.
合一就是根據希臘的牧師提到.
除非建立在第一章到第三章的教導上.
否則就沒有在基督裡面的合一.
既然基督徒的合一是源於.
神在基督裡面的作為.
祂就不是為了合一這個目標.
而刻意追求一些的使命.
簡單來說是怎樣說呢.
就說合一在教會裡面的合一.
和外面的團體的合一有很大的分別.
可能你參加過團體.
大家很死黨的兄弟.
盡心來做一件事.
做一些慈善的工作.
但這個不是教會裡面的合一.
我不是為了做一些事而做一些事.
而互相很親密地做一件事.
而是說這個合一是.
我們對神救恩的一種回應.
不是為了工作而工作.
而是回應神自己對我們救恩的一種回應.

$^{201}$怎樣去保持聖靈所賜的合一.
在華本裡面是用「竭力保守」.
在修靈本裡面是說到「竭力保持」.
我這裡寫了兩點.
第一點是說到屬靈的合一.
是有一個原則性的存在.
這個需要解釋一下是什麼意思.
我用一個例子.
有時候我們打工的要出差.
有時候我出差去到不同的地方的時候.
如果遇到週末的時候.
我會找當地一間教會.
參加當地教會的崇拜.
原則上那間教會的信徒是我的兄弟姐妹.
應該是這樣.
他們都是信耶穌的.
他們都是我的兄弟.
也是我的姐妹.
但是我只是一個過客跟他們一起敬拜.
我不能夠幫助他們做任何合一的工作.
或者幫助他們教會有什麼成長方面的貢獻.
這個只是一個原則性.
一個原則性是叫我們.
我們沒有分別在神的家裡面.
在不同的地方,不同的言語,不同的種族裡面.
我們都是一家人.
但是我們不能夠.
因為只是一個過客.
我們不能夠參與當地的事奉.
或者是任何的工作.
但如果你是T4C的成員.
你經常來T4C聚會.
你就必須要參與其中的工作.
因為你是其中的一部分.
你有義務盡力去維護在基督裡面的合一.
今天我帶了一樣東西給大家看看.
一支小旗子.
很簡單的小旗子.
我包著它以免弄髒.
這支小旗子是我在亞利桑那讀書的時候.

$^{241}$畢業的時候.
在團契裡面給每個畢業的人有一支小旗子在那裡.
大家都可以忽略ICC.
Illinois Chinese Christian Fellowship.
在這裡有個標誌寫什麼呢?.
Bloom where you are planted.
就是說在你所栽種的地方你要開花.
當然要結果.
但是沒有地方寫這麼多東西.
這支小旗子我已經帶了40年了.
在我家裡.
我每次搬家我都捨不得丟掉.
因為這是一個很好的提醒告訴我們.
如果今天有一天因為工作或者其他原因.
你搬到別的地方.
你必須要在那個地方找一個你能夠服侍的教會.
你能夠參與的教會.
你能夠成長的教會.
你要成為當地教會的一份子.
你才能夠在那個教會裡面.
You are planted 種在那裡.
你要成為那個教會能夠保持合一的人.
合一是需要當地的信徒.
竭力,盡力來實踐出來.
而不是一個口號.
在和合本裡面.
如果大家看清楚.
打開你自己的聖經.
在第三節.
它講到用和平來彼此聯絡.
讓我慢慢解釋這個字.
為什麼要解釋.
聯絡這個字在希臘文來說就是Synthesma.
這個是什麼呢?是一個連接,一個關係,一個connection.
這個字可能是因為你是某某人的親戚.
或者你結婚.
你的太太和你成為一家人.
或者其他的利益的聯繫.
某些事情你和她很buddy.
這個是那個字.

$^{281}$但在聖經裡面我們要留意在聯絡之前有兩個字.
是什麼呢?和平.
在英文裡面幾個翻譯本都用同一個幾個字.
就是the bond of peace.
所以我們講到這種的關係.
這種的聯繫是跟我們的平安有關係.
不是簡簡單單一個聯繫.
我們喜歡打乒乓球.
打乒乓球的會.
不是一個團體.
我們是一個因為和平的緣故所建立的聯繫.
因為我們信耶穌.
因為神給我們有平安.
我們要維持這種的平安.
而所建立那種的親密的關係.
在《雷俊宗的譯本》.
我很喜歡.
因為這個很特別.
以和平之承索.
鐵力保守聖靈所賜合而為一的心.
他用一個很生動的圖畫給我們看到.
和平的承索.
意思是什麼呢?.
今天我們沒有帶一條長的繩子.
所以不能夠做一個demo.
如果有一條繩子叫做和平的話.
我們將聚會的人全部綁在一起.
大概是這個意思.
就是說當我們是信耶穌那幫人.
我們要維持緊緊地.
將我們所有信耶穌的人綁在一起.
不要放鬆.
要維持那種親密的經歷.
是那種親密的關係.
而將我們全部因為這個信仰.
我們信同一樣東西的時候.
將我們綁在一起.
怎樣去維持合一.
我們怎樣在行事為人的時候.
保持那種合一的態度.

$^{321}$保羅提出四樣東西.
就是凡事謙虛,溫柔,忍耐,用安心,互相寬容.
如果大家看《和合本》的時候.
其實從第一節到第六節.
只有兩句話.
就是第一節到第三節是一句話.
第四節到第六節是另外一句話.
所以只有兩句話在六節的經文裡面.
所以全部是連在一起的.
沒有分開的.
在《和合本》裡面是將它切開兩段.
我們講一講.
這些是基督徒的一種行事為人.
一種態度.
是怎樣去令到教會能夠緊緊地綁在一起.
是互相有那種和平的存在.
謙虛.
什麼叫謙虛.
謙虛是要求我們放下我們個人那種驕傲.
不是用一種高高在上的態度.
跟一個初信的人講一些大道理.
不是要做一個主人.
而是做一個僕人.
我們最好的例子就是在菲勒比書講到.
描述耶穌怎樣去謙卑自己.
祂可以跟神有同樣的位置.
但是謙卑成為人.
更加成為一個奴僕的樣式.
這是需要一個很大的定力.
很大的認知.
就是說在什麼情況我應該做到謙虛的狀態.
另外一樣就是溫柔.
溫柔通常來說.
是一種我們覺得很溫和很柔和.
是一種人的個性和特質.
在新約的書信裡面.
有兩個地方給我們有更加清楚的一種.
對溫和這兩個字有不同的認知.
在加拉太書六章一節.
請聽我讀.

$^{361}$「若有人偶然被過犯所成.
你們屬靈的人.
就要用溫柔的心把他挽回過來」.
我再讀一次.
「若有人偶然被過犯所成.
你們屬靈的人就要用溫柔的心.
把他挽回過來」.
這是什麼意思.
有人犯罪了.
在教會裡面有人犯罪了.
你認識他是一個弟兄.
你怎樣處理這件事.
你是不是拿著他.
你做錯了.
你不應該講話.
你不應該做什麼.
在這裡的教導是什麼.
我們應該用一個溫柔的.
我不是認同他們所犯的罪.
而是說我們用什麼方法.
將他們挽回回來.
將他們帶回來.
我們很親密的和平關係.
怎樣將他們帶回來.
另外一段經文.
在提摩泰後書二章二十三節.
這裡說.
「要用溫柔勸導反對的人.
也許神會給他們悔改的心.
能明白真理」.
用溫柔勸導反對的人.
也許神會給他們悔改的心.
能明白真理.
意思是說什麼.
有些人走歪了.
他走歪了.
和真理有所分別.
他不願意悔改.
反而他用一種對立的態度.
和你吵架.

$^{401}$然而在這樣的情況下.
保羅勸告提醒提摩泰.
你要用一種溫柔的態度來勸勉他們.
這是很難做到的.
人家錯了.
你錯了.
你信的錯了.
我們怎樣去提醒自己.
我們遇到這樣的情況的時候.
我們第一件事不是發怒.
是發疑怒.
要和你講實話.
然後想一想.
我怎樣能夠將他們帶回來.
什麼是最好的方法.
我們能否用溫柔的態度.
將他們帶回來.
忍耐.
忍耐是一種不放棄的精神.
在逆境裡面堅持到底.
在《雅各書》第五章第十節.
講到弟兄們.
你們要把先前奉主明說話的.
當作能受苦能忍耐的榜樣.
然後跟著那一節的經文提到.
用「約伯」作為一個榜樣.
忍耐不容易.
我們說.
特別是在一個快速的社會裡面.
我做一件事.
我就要得到另一件事.
我們不懂得去等待.
是一種克制.
是一種自我的克制.
而不是說.
當有一些對你不利或者不公平的事件的時候.
我就很急急去報復對方.
要忍耐.
等待神自己給你的時間.
最後一點是講到.

$^{441}$愛心互相來寬容.
在愛裡面是學習包容的功課.
當你認為自己.
我行事為人和我所蒙的恩相稱的時候.
你有沒有帶著愛去做這件事.
你有沒有帶著愛去走.
當你選擇謙卑來退讓的時候.
你是只是一個外表的表達那種謙卑.
還是你真真正正為了愛那個弟兄姊妹的時候.
你願意用愛來包容他.
當人們陷在罪裡面.
或者對真理有些不同的看法的時候.
你能不能夠用一種溫和來跟他交談.
用一種愛來包容.
希望將他從罪裡面引渡回來.
當你不能夠忍耐.
有些事我要看到某件事發生要做得好.
但見到的就是一種遲遲沒有改變的狀態.
你願不願意等待.
願不願意用一種愛心去包容你的弟兄.
在信徒裡面的群體去維持合一.
是充滿挑戰的.
有時候甚至令人覺得有一種窒息的感覺.
是不是?.
從第四節到第六節.
在保羅是提出幾樣東西.
是在合一的過程裡面是不能夠改變的.
他們有同一個身體.
在哥林多前說十二章十二節.
他說身體是一個卻有很多的肢體.
身體的肢體雖然多仍然是一個身體.
基督也是這樣.
所以剛才提到一樣東西.
無論我們去到哪個國家哪個民族.
講哪種的語言.
但我們同樣都是屬於基督身體的一部分.
同一位聖靈.
剛才我們讀到關於猶太人是希臘人.
在哥林多前書十二章十三節.
他說我們無論是猶太人是希臘人.

$^{481}$是圍爐的是自主的.
都是從一位的聖靈成為一個身體.
並且共享這個聖靈.
所以其他教會的聖靈怎樣感動他的話.
同一個聖靈可以同樣在教會裡面感動你.
同一位主 依然是主耶穌基督.
他說猶太人和希臘人.
在羅馬書十章第十二節那裡.
猶太人和希臘人是沒有分別的.
因為人人都有同一位主.
他亦都厚待求告他的每一個人.
當我們祈禱.
當我們用耶穌的名字.
說奉主耶穌基督求的時候.
我們是求同一個主.
並沒有分別.
同一個信仰.
這個是說到建立我們的信仰.
是建立在神的根基.
神自己話語的身上.
這個真理無論是昨日.
今日和將來都不會改變的.
同一個洗禮.
受洗的形式.
我們都在次要神學.
我們都有很多的討論.
受洗的形式固然是重要.
但更加重要的是什麼呢?.
受洗所謂標誌著一些什麼?.
是標誌著我們向這個世界宣告.
我再不屬於世界.
我是屬於耶穌基督.
是宣告我們已經和耶穌基督一同死.
一同復活.
亦都願意跟隨耶穌過一個新生的生活.
來配得神所賜給我們的永生.
同一位父神.
那位父神是誰?.
在第六節裡說得非常漂亮.
「祂就是萬人之父.

$^{521}$超乎萬有之上.
貫通萬有.
在萬有之中」.
保羅所提到的這一點.
為什麼這麼重要?.
就是說無論在什麼時代.
無論在什麼環境.
在什麼國家,什麼民族裡面.
我們教會能夠團結在一起的時候.
而不分開的時候.
是因為這些原因.
當你找到一間教會的時候.
你要注意它所標榜的是什麼.
它所走出來的是什麼.
是否符合聖經裡面的教導.
我用一個很簡短的.
來講第七節到十六節.
只有一張slide.
在教會合一的教會.
基督的身體怎樣才能夠成長?.
在教會裡面保羅提到一件事.
耶穌基督是聖靈和聖靈是所賜的恩賜.
這裡所提到的有使徒,有先知,有傳福音.
有牧者和教師.
但是他們和所有無論是什麼職位.
他們都有同樣的目標.
他們為了做什麼?.
是要裝備,是要事奉.
是要建立耶穌基督的身體.
經文裡面提到.
「靠主全身都連接得緊湊」.
「百折各按其跡」.
「照著身體的功能彼此相助」.
「令到身體漸漸增長」.
「在愛中建立自己」.
這就是神自己的心意.
不單止是要你團結.
而在這個團結當中.
因為神所賜給你的恩賜.
令到這個團體.

$^{561}$我們所被種在這個團體T4C裡面.
讓T4C能夠慢慢地增長.
而不是停留.
就這樣結束.
我只有三點.
如果我們明白神拯救的大愛.
就必須要在我們行事為人上.
作出一個回應.
我們不能夠採取.
成為基督的身體的一部份.
我們就必須要參與.
讓教會才能夠增長.
你多一分鐘坐在你的位置.
完全沒有動作.
連祈禱都不能為教會祈禱的話.
這個教會不會增長.
教會需要你.
基督的身體需要你.
教會的合一是需要.
每一位信徒有所付出.
所以我今天的講道.
是對大家問一個問題.
你願意付出一些什麼.
令到教會能夠有合一的身體.
能夠有增長的身體.
我們低頭做一個祈禱.
真是顛覆我們多一次.
多謝你.
也多謝寶樂經過以弗所殊.
來向我們的教導.
最後讓我們檢測自己.
問問自己的問題.
我們愛你有多少.
我們認識你有多少.
我們願意付出有多少.
主要是求你的話語與我們同在.
不斷提醒我們.
共主耶穌領教.
阿們.
\newpage



\section{箴言 25:11}
\label{sec:032318}
\textbf{金蘋果}~~~~曹子光牧師
\newline
\newline
連結: \href{http://www.church.com.hk/acms/content.asp?site=cdc\&op=show\&type=product\&code=032318\&layout=sermon}{\texttt{032318}} ~~~~ 語音日期: 2026-06-14
\newline
\newline
\hyperref[sec:032317]{\small{< < < PREV SERMON < < <}}
~
\hyperref[sec:index_chronic]{\small{[返順時目]}}
~
\hyperref[sec:032319]{\small{> > > NEXT SERMON > > >}}
\newline
\newline
箴言 25:11
\newline
\begin{longtable}{cl}
\hline
\hline
章節 & 經文 (和合本修訂版)\\
\hline
25:11 & \begin{tabularx}{0.7\textwidth}{X} 一句話說得合宜，就如金蘋果在銀網子裡。 \end{tabularx} \\ \\
[1ex]
\hline
\hline
\end{longtable}
$^{1}$好 大家望一望這個做得很好的宣傳的海報.
大家不要以為這個是世界盃.
這個不是世界盃的宣傳.
其實 蘋果大家會不會知道是什麼顏色.
你見過什麼顏色的蘋果.
綠色 紅色 有沒有見過金色的蘋果.
今天介紹給你認識金色的蘋果.
根據最新的科學的研究告訴我們.
有一樣東西正在逐漸消失.
就是我們實體講說話的能力.
在2007年每一個人平均講的說話.
大概是一萬六千個字.
到了2019年的時候.
這個數字跌到一萬二千八百個字.
最令人擔憂的原來就是25歲或以下.
最喜歡上網的那些年輕人.
最喜歡去社交平台的年輕人.
最喜歡玩數碼科技的年輕人.
他們越來越少講說話.
大家做家長的同不同意.
在家裡面是不是發覺.
他們幾乎在家裡面.
都不知道有沒有講多過十句說話.
甚至他們將自己個人的感受.
轉化成為符號去代表.
甚至是按一個like就算數了.
另一方面.
過去我們總有一個黑板的印象.
就是覺得女人講話總是比男人多.
同不同意嗎.
以前我們總是這樣想的.
但是原來在2007年的報告告訴我們.
不是這樣的.
男和女講話的數字.
其實是差不多的.
不過最近有一個最新的調查.
就是訪問了這四個國家.
跨越了十歲到九十四歲.
大型的一個研究.
就發現喜歡不喜歡講說話.

$^{41}$原來是取決於人生階段的.
而在二十五歲到六十四歲.
這個成家立業的人生黃金的階段.
女人平均是每天講.
你猜多少個字.
是二萬二千個字.
差不多是二萬二千個字.
從你起床到今天.
不知道你講了多少個字.
如果你是這個年數裡面的姐妹們.
男人竟然只有一萬八千五百個字左右.
所以女性多了差不多三千個字或以上.
我們講說話的減少.
其實正正是代表著.
可能在這個時代.
年輕人真的很不喜歡講話.
女人講說話多了.
在這個年歲裡面的女士多了.
其實大概是因為她們那個年齡層.
可能真的又要照顧家庭.
又要工作.
她們講說話真的多一點是明白的.
但我們要明白一件事.
說話多不等於講得對.
說話少也不等於合宜.
有時候還要配合你的身份.
有沒有合乎聖經的教導.
看看你的身份是否願意講這個說話.
還有你講的話是否合理.
前兩三天是否太古城.
發生了一宗很不開心的事件.
媽媽和女兒都一起跳樓.
據聞是和說話爭吵有關的.
可能是為了某些原因.
而近期香港的教育界風波不斷.
其中一個大家都可能聽過.
屯門一個中學的校長帶隊去新加坡.
然後在旅遊巴士上跟別人爆粗.
結果現在應該是被人解僱了.
他本來想辭職的.

$^{81}$後來辭都不用辭了.
被人解僱了.
但最近有一個不知哪個中學的.
孔什麼中學的周什麼校長.
被指用一些很侮辱女性的字眼.
說那些女學生.
而且用一些不當的言辭來對待教師.
引發出極大的爭議.
還有一個很出名的籃球教練翁什麼.
他在教的時候說的話非常不恰當.
這三件事件告訴我們.
作為老師在不適當的場合.
說不適當的話.
後果可以非常之嚴重.
我問過今早有一個教育界的朋友.
弟兄他就說.
老師原來的規則是.
上課或者學校範圍之內.
是不能講粗口的.
總之在學生面前是不能夠講粗口的.
今天《箴言書》25章第11節.
要告訴我們什麼叫做真正說得合宜的話.
今天我們一起看看一節的經文.
《箴言書》25章第11節.
我們一起讀一下這一節的經文.
請讀.
我們再讀多一次.
一句話說得合宜.
就如金蘋果才望穩.
這一句話其實可以從三個角度去理解.
第一就是合宜的話要帶著神聖的份量.
合宜的話要帶著神聖的份量.
國際上很著名的溝通技術裡面.
其中一個黃金法則叫做Think.
就是思考這個字.
T這個字就是代表True.
真實的東西.
這個是門檻來的.
到底你講的東西是不是真的呢.
H是Helpful.

$^{121}$就是說你面對當時的情況.
你跟對方講的這件事有沒有幫助呢.
I是Inspiring.
就是說你能不能夠帶來一些.
是不是一些燈光效果.
帶來一些正面的影響呢.
或者純粹是宣洩你的感情呢.
N是Necessary.
現在是不是真的要我這樣講出來呢.
K是Kind.
我的說話 我的動機.
是不是出於上帝的善意.
或者是對人的尊重呢.
溝通技巧當然是重要.
這個Think提供了技術的方法.
提供了由外到內的操練.
好的.
但是講到底其實都是包裝來的.
就是包裝你怎樣講話.
但聖經在箴言這裡25章第11節這裡.
意思是剛剛相反.
學再多的技巧.
但生命如果沒有智慧.
沒有正直.
問題更多.
在箴言書25章.
其實這章的經文主要是跟權柄.
跟公義 跟智慧是有關的.
特別提到人倫的倫理.
還有自我控制.
去到第11節到28節更加是集中在言語上面.
那個不同的格言.
所以是特別針對那些宮廷未來的領袖.
所以針對的不是技巧.
而是你的身份本身.
跟生命關係塑造出來的品格.
這節經文原文的上半句講到金蘋果.
這個絕對不是水果牌子.
玫瑰金的電子產品裡面的那個脈頭.
我想你知道首先蘋果這個字.

$^{161}$Tapioca這個字.
其實這個字在原文.
Tapioca這個字在舊約只是出現了很少次數.
最多是牙膏.
牙膏是出得最多的.
在那裡是描繪一種很吸引人.
很香的 很美麗的果.
令人很想要的.
不知道大家喜不喜歡買蘋果.
有誰喜歡吃蘋果 你有幾個舉手.
都不少人.
買哪個牌子.
今早早堂都是這個答案.
識食.
Pink Lady大家知不知道多少錢.
真是你看看哪間超市.
起碼都差不多逼近三磅.
有些還要三磅幾.
如果你去那些名牌的店舖買的話.
當然如果比起AZDA牌.
TESCO牌那些.
當然貴一兩磅.
大佬咬下去的時候.
很好吃的.
所以Paul姐答得對.
你們是不是真的買Pink Lady.
真是識食.
不過你要知道.
在以色列這個地方.
是很難找到蘋果的.
特別是古代.
今時今日如果你去哥蘭高地.
即是以色列的北邊哥蘭高地這個地方.
是有的 有蘋果的種植場.
為什麼.
是因為蘋果需要極凍的氣候.
令到它那個氣候的改變.
逼到它可以生了蘋果出來.
可以咬下去爽脆到不得了.
是有氣候的因素在裡面的.

$^{201}$所以以色列在古時.
基本上是找不到蘋果場的.
所以這裡所說的蘋果.
極有可能不是我們認識的蘋果.
而是很有可能是杏.
又或者是木梨.
或者叫溫缽這個名字.
木梨就是我們認知的名字.
而不是今日所說的蘋果.
然後經文馬上就告訴我們.
這個蘋果是金.
這裡絕對不是說金色的蘋果.
試問你何時見到金色的蘋果.
見到你都不敢吃 是不是.
這裡是在說一種食物.
不是在說顏色.
而是在說一個貴金屬的蘋果.
而下半句因為用銀這個字.
所以很明顯這裡是在說金屬.
而這裡下半句.
25章第11節真是下半句.
那裡一開頭的字就是講的說話.
所以其實金蘋果是等於講的說話.
所以用貴金屬的金.
去描繪講的說話.
大家知不知道金是很寶貴.
很值錢 是不是.
第一樣東西.
金是什麼東西呢.
為什麼用金成為一個講說話的特質呢.
就是金是真實的.
是不會造假的.
大家知道999金是不會生鏽的.
有沒有.
當然說沒有的.
999金是不會生鏽.
亦都不會褪色的.
不能夠以假亂真.
杜金隔了一段時間之後.
你雖然是表面閃亮亮.

$^{241}$但是一段長時間之後.
經過水 經過空氣等等因素.
很快就會露出馬腳出來.
可能褪色 可能生鏽.
但是唯獨真正的999金.
才可以經歷過這個時間的考驗.
永恆不變 真正有價值.
所以大家結婚的時候.
是不是用金.
就是結婚的時候用什麼金器.
999那些戒指 金鏈那些.
因為保值嘛.
很多人都會用這些.
但是如果是假的.
表面好像很甜蜜的.
好像用回說話那樣.
表面好像很順義的.
就好像那種薄薄的杜金.
但是日子久了.
遇到考驗的時候.
它就會褪色.
它就會露出馬腳出來.
金亦都是有份量.
它不是空洞的.
用真正的黃金.
999黃金.
代表著重量的份量.
是貴重 是有價值的.
大家知不知道在五年前.
用金實的代表.
五年前的黃金和今天五年後的黃金.
知不知道升了多少倍.
兩倍 是兩倍.
是兩倍 足足是.
當然我只是高位.
俠義的說話的價值.
的而且確是很重要的份量.
它在聽者的心裡面.
留下一個非常之重要的諾印.
不是說完就飄過就算數.

$^{281}$不是簡單一個讚一個like.
發表幾個符號.
轉接幾個貼文.
去代替真實的說話.
是代替不到的.
有份量的說話是很重要.
所以金是代表有份量.
不是空洞的.
而且金第三代表永恆.
是不會過時的.
一句真實的說話.
出口的時候.
可以在人心裡面.
逗留了很久很久的年日.
大家有沒有試過.
人家跟你講一句說話.
你記他很久呢.
很厲害的 有沒有試過.
我試過.
我讀神學的時候.
可能大家不知道.
我告訴大家一個秘密.
寫錄影了.
這個秘密是什麼呢.
我見到時候就亂讀.
我們有天台.
我們每一晚.
跟幾個男生就上去.
我們就叫做頹廢三人組.
包括我在內.
我們都是平日不讀書的.
不太讀書的.
我們就很頹廢.
上到去就馬上風花樹葉.
講一輪東西.
講一下這個講一下那個.
這樣就上完一天課了.
然後.
其中一個我們的同學.
他今天已經離開了世界.

$^{321}$他在一本書.
他臨畢業的時候.
送了一本書給我.
他在這本書裡面.
寫了一句說話.
到今時今日都記得.
他說不要做小人.
要做大人.
金就好像代表永恆.
是不會過時的.
那種價值的一句話.
是那種真實性.
是你永遠都會記得的.
有沒有留意.
這三樣東西.
沒有一樣東西是技巧來的.
全部都是一個人.
內在生命的狀態.
是不是用think.
可以用這些法則.
可以全部代表.
因為講出來的.
依然是道金的假貨都未定.
而不是999的真金.
美國南北戰爭的時候.
有一位媽咪叫做北克斯比.
她就失去了.
五個戰場上陣亡的兒子.
林肯總統在白忙裡面.
打算親手去寫一封信.
去慰問她.
信裡面沒有任何官方的說話.
也沒有任何官方的文件的印鑑在裡面.
她真的寫了一封個人的信給她.
在信裡面她就說.
她感受到.
那種失去那麼多孩子的悲傷.
是多麼沉重.
是多麼痛苦.
她祈求天父.

$^{361}$以全能的安慰.
減輕你的哀傷.
又讓你記住.
你所摯愛的人.
為了保衛這個國家的自由而犧牲.
這一份是一件光榮的事件.
這封信對媽咪來說.
是一封非常珍藏的信.
她保留了這一封信一生.
在原文裡面.
這一節的格言.
二十五章的第十一節.
有八個希伯來文字.
上半句有四個.
下半句有四個.
一開始那兩個字是金蘋果.
而下半句開始那兩個字.
剛才也說了.
是說話.
即是說的話.
是平衡的.
所以說的說話.
就好像金蘋果那麼寶貴.
那麼有價值.
說的說話.
其實在原文的字義裡面.
有聖經學者說.
其實不是說口齒伶俐那麼簡單.
而是說當那個人.
看清楚所有複雜的環境之後.
他衡量過了.
然後考慮得很清楚.
才去說出那句說話.
是一個協商的決定.
即是說他是考慮得很清楚.
才去說出那句說話.
換句話說.
這件事是真實的.
是公義.
就好像聖經裡面.

$^{401}$那種很有神聖份量的說話.
神聖的愛.
就算你處於什麼狀況.
矛盾壓力之下.
經過深思熟慮.
仍然都能夠令到我們.
去說一些有智慧.
有神聖的份量的說話出街.
第二件事.
合而的話.
是能夠承載恩慈的美麗.
有聖經學者就說.
純金製造的蘋果.
雖然是極其貴重.
但如果沒有一個合適的底座.
去襯托它的時候.
它的美感和價值.
就會蕩然無存.
為了讓大家方便理解這件事.
我就拿了我家裡的一件寶物出來.
我就很努力.
找家裡有沒有金的東西.
是沒有的.
銀的東西都沒有.
對不起.
你說.
對不起.
這個是甚麼來的.
有沒有人知道這個是甚麼來的.
如果你從遠處看.
可能是發布膠球或者球.
不知道是甚麼來的.
可能是一件裝飾物.
如果純粹扔在地上.
你真的不知道它是甚麼來的.
但是.
當如果有個座盤.
當這個是銀色.
扔上去之後.
是不是變漂亮了.

$^{441}$喂.
給一點.
臨時演員都是演員來的.
大哥.
還有.
其實它是會發光的.
發光了之後.
才是它最漂亮的東西.
明白了嗎.
OK.
沒錯.
所以你要知道.
金蘋果.
它需要放在一個座盤那裡.
這樣才可以襯托它的美麗.
如果這個金蘋果.
它是扔在泥瓣裡面的時候.
它雖然本質是金蘋果.
你說是不是很噁心.
整件事就噁心了.
所以它需要.
就好像.
如果真的這樣放在泥瓣裡面.
就好像箴言所說.
金環插在豬鼻裡面.
遊街一樣.
非但不漂亮.
更加是很噁心的災難.
它類比出說話.
如果缺乏愛.
說真的東西.
沒有用力去承載它的時候.
金蘋果再美麗.
只不過是.
好像孤零零地下.
你剛才見到的.
這個圓滾滾的東西.
不漂亮.
雖然你都說它現在不漂亮.
這一次的經文說到.

$^{481}$金蘋果是放在銀製的網紙裡面.
首先我要你知道.
銀都是一個金屬.
大家又知不知道.
五年前如果你買了銀.
現在.
2026年銀最高的時候.
你知不知道升了多少倍.
是四倍.
我沒有叫你買.
是四倍.
這裡是在說.
金銅銀.
一個最貴的金屬.
最名貴的東西.
你說出說話.
那種美麗.
那種高貴.
那種神聖.
這裡說銀網紙.
不過我要指出.
網紙這個字.
Maskive這個字.
其實更有可能是指雕刻.
或者是展品.
原文這個字在舊約裡面.
是多數用來指雕刻.
塑像.
或者是影像.
或者是展品.
所以有學者就說.
藝術家.
很有可能是說.
這裡是說.
藝術家做一個.
雕刻的藝術品.
將金的蘋果.
鑲上去.
又或者是將金蘋果.
放上去.

$^{521}$這個藝術品的上面.
襯托起它.
展示出一個很漂亮的.
一個展品出來.
我們知道.
金和銀都是很貴重的金屬.
在古代裡面.
金紙雖然很貴重.
但是因為它的質地比較偏軟.
如果沒有了其他的金屬.
例如銀.
結構支撐住它.
它很難雕刻到.
更加細緻複雜的.
一些工藝品出來.
這節的經文.
將兩個很珍貴的東西.
放在一起.
所以這句格言.
告訴我們.
似乎暗示.
金是象徵著說話的內容.
銀是象徵著方式.
當然這裡不是說.
說話的內容比較重要.
還是方式比較重要.
不是.
而是兩樣都這麼重要.
才是完美.
才是perfect.
金再好.
都需要銀器來承托它.
才成為一件藝術品.
用價值連城的金.
做了蘋果.
但是也用銀.
做了一個承載物.
將這個藝術品.
完完全全的公諸於世.
這個才是美麗.

$^{561}$才是越人欣賞.
用金屬製的器皿.
彷彿將金屬的蘋果.
將本質進一步提升.
這個才是這段.
這個格言的重要提醒.
我們的說話.
我們說話的方式.
態度和動機.
有沒有想過.
你的語氣是溫柔.
還是像人的呢.
你的態度謙卑.
還是高高在上的呢.
你的關係.
你站在對方那邊去想.
還是你站在自己的角度去想.
你的動機.
你說話的這句說話.
是為對方好.
還是為自己舒服.
還是想贏對方呢.
今天在這個世界裡面.
出現了兩種極端的問題.
第一.
就是很喜歡讚.
不講真心話.
很空泛.
很客套.
沒有真實的金的東西在裡面.
第二.
就是在網絡裡面.
一言不發.
對不起.
一言不合就是集體攻擊.
說可能是真的.
你說的話可能是真心話.
是真的.
但沒有任何的憐憫去承托住它.
結果傷痕累累.

$^{601}$兩種都是失敗.
有一位非常出色的外科醫生.
他叫做保羅布蘭特.
在印度行醫很多年.
專門治麻風病.
有一次有一個壞消息.
要跟一個年輕人說的時候.
他想想怎樣可以將這個年輕人.
讓他能夠明白.
又接受到呢.
布蘭特醫生於是乎.
他就走到這個年輕人的旁邊.
握住他的手.
跟他說.
讓他感受一下他的溫暖.
然後我告訴你.
有一個很困難的消息.
但是我希望你明白.
你知道我們是會一起面對這件事.
那一句話其實是很殘酷的真相.
那個是金.
但是握住他的雙手.
那個熱度.
溫暖.
貼近我們一起的這句話.
那個是銀.
就是承托住那個實相.
那個雕刻.
那個消息其實是非常令人難受.
但是他的表達的方法.
令到這個年輕人感覺到很舒服.
結果.
他後期成為了布蘭特醫生的助手.
他甚至說.
為什麼我很喜歡這個醫生.
甚至我欣賞他的原因.
不是因為他說了什麼.
是因為他用什麼方式來說話.
今天可能我們擔心傷害對方.
怕說了之後不知道是什麼後果.

$^{641}$所以我們選擇沉默不說話.
但沉默本身是逃避.
我們希望神今天藉著這句金真言來幫助我們.
多些用體恤多些用恩慈.
去承托那個蘋果.
金蘋果.
最後第三樣東西.
俠兒的說話.
是要說得精妙.
也要聽得受教.
這裡下半句的真言.
最後那兩個字是最困難的.
大概是中文翻譯.
俠兒這個字就是最後的這個字.
奧番這個字.
原文是最後的一個字.
其實在《聖經》裡面出現過一次.
有語言學家就發現.
這個字在古代.
其實很有可能和face.
或者面是相關的.
如果是這樣演變出來的時候.
根據這個字的字源.
似乎就是在說.
以他自己的方式.
用對的態度.
妥當.
合乎禮貌.
這個阿拉伯的意思.
阿拉伯語的意思.
套入這個格言裡面.
他就會體會到一樣東西.
這裡可能說的是.
向著他的臉.
顧及他的面子.
用最尊重他的態度.
說那句話出來.
一句穩妥的話.
是因為說話的那個.
到底有沒有看著對方.

$^{681}$有沒有顧及對方的感受.
有沒有看過他的自尊.
看過他的感情是怎麼樣.
才傳達那個訊息出來.
當我們如果有做到的話.
這個是一個生命的連結.
弟兄姊妹.
這個是徹底打破了.
我們現代人說話的習慣.
今天我們最喜歡的是.
什麼情況去批判人呢.
就是在網絡上面罵人.
大家有沒有試過看那些片子.
在那些片子你有沒有見過.
在下面那些留言那裡.
很多人就罵不同的東西.
我們最近.
我自己在.
不要緊.
說出來.
我是見到神樂園畢業的.
最近我們見到.
神樂園的後任院長.
高博士.
我先離身.
本身我不是嚴格來說.
我並不認識他.
我只是曾經在Facebook留言.
在訊息裡面.
和他溝通了幾句.
最近他有給兩個主媒的網媒.
叫公X門徒.
和WeXchurch.
在網絡裡面罵他.
隔空.
隔著空氣.
用一些很難聽的字眼說他.
我不是說誰對誰錯.
而是這些字眼本身.
很明顯是有一些.

$^{721}$很貶低一個人.
很不尊重一個人.
我們很流行在社交平台.
在WhatsApp留言.
其實都是同樣的東西.
在鍵盤打字的時候.
我們是看不到對方的臉.
我們不知道他的感受是怎樣.
他在流淚.
他痛苦.
他心碎的臉.
我們看不到.
網絡的延遲.
但是越來越刻薄.
越來越肆無忌憚.
我們以為自己維護真理.
但其實我們在發洩.
當然.
可能那些人會覺得.
我在做一個真理的事情.
所說的東西.
我都是尊重他才這樣說.
但是一個客觀的人.
看他的說話.
一點都不是尊重的說話.
我真是利申.
我並不認識高博士.
不是完全認識他.
所以我可以很客觀地.
這樣說他們的說話.
是非常不尊重.
我們需要用恩慈.
成為了真理的承托.
令到整個藝術品有價值.
另外有些學者研究了很久.
發現這個字.
這個字其實可以指輪子.
奧番這個字其實可以指輪子.
在《初開及記》第十四章.
曾經出現過一次.

$^{761}$是用輪子這個字.
為什麼呢.
原來奧番這個字.
其實在這個真言裡面.
是一個複數的表達.
即是兩個輪子的意思.
更有趣的是.
學者說兩個輪子.
正好是希伯來詩歌裡面.
一個比喻.
上面那句.
下面那句詩.
又或者阿拉伯語裡面.
同一個字可以說到.
一線門的左邊門和右邊門.
我用一個例子.
讓大家明白.
大家有沒有試過換車胎.
有啦.
換車胎一定有的.
開車的人.
通常你去換車胎的時候.
很多人都會告訴你.
全部換了一對.
很少叫你換一隻.
因為不久之後.
你一定會換另一隻.
所以那些汽車公司.
通常告訴你.
全部換一對.
因為一對是很重要的.
現在這個兩個輪子.
其實是類似這樣的意思.
或者門都是這樣的意思.
我們首先用門去理解一下.
古時代宮殿的大門.
左邊和右邊的門.
都是刻了很多圖案的.
而且關門的時候.
是沒有縫隙的.

$^{801}$但是英國的人.
大家有沒有留意到.
英國的特色.
做房子的時候.
做門的時候.
你會發覺有很多縫隙.
下面有縫隙.
旁邊有縫隙.
上面有縫隙.
基本上你是關不住的.
在香港的門.
多數都可以的.
這個就是英國的文化.
我估計.
這是一個文化.
如果你普通做一間屋的時候.
大部分的屋企都是這樣.
除非你那間屋是很豪華.
否則的話他們都是這樣.
所以英國人見怪得怪.
但是現在是在說皇宮的門.
是不會有縫隙的.
所以這兩道門.
應該是非常之一模一樣.
是關得很密的.
是不會這邊大這邊小的.
等於車輪一樣.
你不會這邊大輪那邊小輪的.
這樣是會翻車的.
你的口一說一句話出來的時候.
是會翻車的.
是會有意外的.
意外遲早都會出現的.
不能夠因為一時衝動.
流於情緒化.
就爆了那些說話出來.
這樣是失去了真理的.
這樣是失去了真理的金.
因為你流露的情緒.
沒有說那個.

$^{841}$只是純粹流露情緒.
但是不能夠過度的懼怕.
而成為一種虛偽.
這樣是失去了恩慈的銀.
所以我們要知道.
在《真言二十五章》第十一節.
他說其實是一個藝術的產品.
金蘋果在銀器的裡面.
在銀紙的網裡面.
或者易造.
下半句其實是一個現實的時空.
就是說你的說話.
在一個適當的時間.
適當的環境.
來講出適當的說話.
這個是重要的.
另外如果你用一個圖像的相輪去講的時候.
有另一幅圖畫出現.
大家有沒有見過圖像.
用泥來去模造那個瓷器.
或者那個產品出來.
即是那個輪是在轉的.
你用雙手去扶著它.
一路去模造它出來.
這個就是相齒輪的一個解釋.
如果你用力太猛會怎樣呢.
整個東西會散.
如果你太輕又弄不出一個形狀.
所以你必須要很有力量.
很平均去弄它出來.
講一句俠義的說話就是這樣.
你要在兩個輪子之間找一個平衡.
既要表達真理.
真實.
亦都要有恩慈配合.
這個就是這裡的意思.
有一個婦人.
剛剛失去了她剛出世的小孩.
痛不欲生.
她的朋友都不知道怎樣好.

$^{881}$有很多人都嘗試去安慰她.
例如跟她說.
你要堅強一點.
孩子是天家.
他們很安好放心.
有些就說你還年輕.
可以生多幾個.
但是這些說話對她來說.
越聽越心痛.
因為每一句說話都好像叫她.
快點收拾心情.
繼續前行.
後來德蘭修女來探望她.
那個婦人心裡面都挺抗拒的.
以為又是坐在那裡教訓我.
講一些這樣的說話.
但德蘭修女走進來的時候.
只在她身邊靜靜的坐在那裡.
握著她的手.
靜默了一段比較長的時間.
然後她就講了一句說話.
上帝知道你的痛.
就是這一句.
那個婦人後來就說.
這一句說話是我一生裡面.
覺得是最合宜的一句說話.
是對的時間.
講出對的說話.
用對的方式.
神眼中的完美的藝術品.
不單止是講得精妙.
更加是要聽的人聽得受教.
這個也是兩個論詞.
雙論的意思.
因為第十一節之後就是第十二節.
第十二節講什麼呢?.
經文我不知道有沒有打出來.
經文在第十二節好像有.
不如我們讀一讀.
自慰人的勸誡.

$^{921}$在聽從的人.
好比金耳環和精金的首飾.
講的人講的話是金蘋果.
聽的人如果你接收的話.
就是一個的金耳環.
雙論同時轉動.
這個才是最完美的狀態.
你的話可能是一件.
是一件金的蘋果.
但必須有個底盤承托它.
這個承托可能是你的態度.
因此憐憫.
但與此同時.
我們也需要有一個.
對方願意聽一個謙卑的心.
聽.
這樣聽完之後.
這樣這個就是最完美的圖畫.
三合一.
才是一件珍藏的藝術品.
就好像前面這個.
亮了燈就真的是三合一.
你要知道.
我們這個時代最缺乏的.
正正是這種說話的藝術.
大家知不知道.
英國人被稱為.
紳士風度的民族.
他們說話真的很有趣.
舉個例子給大家聽.
當一個英國人跟你說.
嘩 你說話好有趣.
你以為很讚你嗎.
你以為在讚你嗎.
其實不是.
他的意思就是說.
嘩 實在太荒謬了.
我根本不在乎.
又或者當一個英國人跟你說.
That's very brave idea.

$^{961}$就是說一個勇敢的想法.
你以為他在讚你嗎.
錯了.
其實他說.
你瘋的 絕對行不通.
當你鄰居敲你的門.
告訴你.
喂 你的小朋友真的勤力練琴.
練得這麼好.
練琴的.
有時候也要放鬆一下.
你以為他緊張你嗎.
其實他告訴你.
你的小朋友的琴聲很大聲.
吵著他.
英國人說話就是這樣.
皇室重敬.
查理斯三世曾經在全國講話裡面.
他說.
我要對Harry.
即是他的兒子.
和他老婆叫什麼.
Maddon 好像是.
表達愛意.
希望他們兩個在海外建立愉快的生活.
英國傳媒就懂得看了.
懂得讀了.
英國傳媒就馬上說.
其實你的意思是遠距離祝福.
不要耍我.
其實你是在說一件事.
你們兩夫妻在美國好好生活.
不要再回來了.
其實他是在說這些.
你們明白嗎.
這些不是叫做俠義.
這些叫做口黑學.
或者是語言藝術.
不是金蘋果.
這個是外表的鍍金.

$^{1001}$裡面裝著地雷的假蘋果.
當我們語言心口不一的時候.
充滿著暗示.
諷刺.
委婉的譴責的時候.
關係只會越來越疏離.
神所要的是合宜的.
是按著真理的.
用愛心說誠實話.
真面話一句的說話說得合宜.
就如金蘋果在銀網紙裡面.
核心是合宜.
不單只是真.
更加是要tiny.
appropriate.
sincere.
不單只是表面好聽.
更加是裡外都一致.
帶有能力和意志的說話.
最後也要聆聽的人用心去接受.
在應用的部分的時候.
我想說三個很重要的應用.
給大家考慮一下.
第一就是家庭上.
我們對家人是特別喜歡說一些.
很傷害性的說話.
不知為何對外面的朋友不會這樣說.
但是家人就肆無忌憚.
什麼都可以說.
我們是很流行這樣.
而且我們留意.
同一句說話.
有些人聽了感受是被愛.
但是某些人聽了想逃跑.
差別往往不是在於內容.
是在於我們說話之前.
我們人的靈魂的狀態是怎樣.
在開口之前.
有沒有先在心裡面停一下.
問自己.

$^{1041}$我要說的這句說話是真實的嗎.
現在我說出來對他有沒有幫助.
我的動機是什麼.
是不是想要做低他.
贏他呢.
這個不是技巧.
這個是說話之前.
我們的心靈狀態.
如果我們的心靈狀態是充滿批判.
不耐煩的話.
就算學多少方式都是沒用的.
都不是金蘋果.
其實我們小朋友就知道.
你知不知道我們小朋友.
都記得我們所說的話.
我記得我和其中一個小朋友說.
很多年前他和我說過一句話.
你記不記得你小時候這樣罵過我.
其實是說好幾歲而已.
很小的好幾歲.
他還記得這麼久.
可想而知.
一句說話隨便開口的時候.
諷刺否定帶出來的後果.
是超過你所想的.
第二樣東西就是職場和網絡上面.
我知道我們做工的時候.
我們在職場可能要做很多的決定.
要說很多的說話.
有時候我們明知道人家做錯了.
我們都不敢去得罪人.
我們可能怕尷尬.
可能覺得關係影響了.
於是就出現了極端的兩種人.
第一種人就是什麼都不說.
明知做錯了就不理他.
什麼都不想什麼都不說.
第二是會說.
什麼都說出來.
其實這個都是問題.

$^{1081}$兩者都是不值得你們學習的.
我們說的話是要幫助對方.
還是要證明自己對呢.
如果是前者想幫助人.
找一個適合的時間私底下.
帶著尊重面對面跟他說.
那個金蘋果就放在合適的銀盤上.
對方收的就會是禮物.
如果是後者你想贏人的.
那就不要說了.
先禱告.
讓神處理好你心裡面想贏的觀念.
然後想下第二個時間再說出來.
最後第三樣東西是最困難的.
就是在教會裡面.
前面那兩個應用呢.
都是叫你注意你說什麼.
但這個應用就是你注意你接收到什麼.
第十二節告訴我們.
智慧人的勸戒在聽從的人的耳中.
是化為金耳環和精金的裝飾.
留意一下這裡.
其實他講到一句很重要的東西.
就是原來我們在教會裡面.
我們很流行被人讚.
但是我們又很抗拒被人勸戒.
當有人指出我們的盲點的時候.
第一個反應就是反駁.
防衛.
就好像昨天我們在一個特別的座會的時候.
我們講弱失的時候.
我們都會講到有些人特別說他們的反應是防衛的.
先駁了你.
我們很多時候都是這樣.
因為大家都是事奉而已.
我們都不收人工的.
他是什麼人.
憑什麼這樣說我.
他自己也是這樣.
他根本不了解我的處境.

$^{1121}$這些的反應.
人之常情.
但是真言告訴我們.
真言告訴我們.
這樣子就好像將一句對的話.
扔進了泥裡.
金蘋果扔進了泥裡.
真正的尊榮.
從來都不是犯錯.
與不犯錯的可能.
而是來自當你面對善意的勸戒的時候.
你沒有去逃跑.
我再講一次.
是你面對真正善意勸戒的時候.
你沒有去逃跑.
你沒有去反擊.
而是用謙卑的心跟他說.
多謝你告訴我.
我會想一想.
那一句本來很刺耳的話.
因為你接受了.
就在你的生命裡變成了一副閃亮的耳環.
耳環.
或者金耳環.
所以.
弟兄姊妹.
非常重要的就是.
如果神差遣你生命裡有屬靈的朋友.
在教會裡有你的牧者.
有你的朋友.
甚至你的團契裡的導師.
等等的人.
或者你的家人.
對你說一句不太舒服的話.
但可能那句話是真實的.
試一下.
不要馬上反駁.
讓那句話停一停在你的心裡.
然後祈禱.
主啊.

$^{1161}$這句話裡有沒有你要我聽的東西.
如果有的.
就讓它在我的生命裡.
從一句刺耳的話.
變成一副金耳環.
我們一起禱告.
主我們知道.
在我們人世間裡有很多東西.
我們很多時候都要我們即時去反應.
特別是人們說很多的說話.
或者我們自己本身說很多的說話.
有時候是別人傷害我們.
有時候是我們傷害別人.
但是原來我們都需要去.
在你的裡面學習.
因為.
金蘋果的說話.
是你想我們做到的.
就放在一個銀的器皿上面.
因為這個擺放.
既是美麗.
既是一個非常完美的組合.
說話.
說真的說話.
能夠有好的態度.
好的時間.
合適的方式.
來去配合的話.
這個當然是最理想.
但更重要的就是對方.
願不願意去聆聽.
如果當對方聆聽.
又願意接受的時候.
這件事就好像是.
第十五章第十一節和第十二節一樣.
這麼寶貴.
求你放在我們的心裡面.
特別是在教會裡面和弟兄姊妹相處.
在家庭裡面和家人相處.
在朋友的圈子.

$^{1201}$在工作的圈子裡面和其他人相處.
我們求你將這個真言放在我們的心當中.
阿們.
\newpage



\section{路加福音 24:13-35}
\label{sec:032319}
\textbf{原來主一直同在}~~~~李適清博士
\newline
\newline
連結: \href{http://www.church.com.hk/acms/content.asp?site=cdc\&op=show\&type=product\&code=032319\&layout=sermon}{\texttt{032319}} ~~~~ 語音日期: 2026-05-03
\newline
\newline
\hyperref[sec:032318]{\small{< < < PREV SERMON < < <}}
~
\hyperref[sec:index_chronic]{\small{[返順時目]}}
~
\hyperref[sec:032320]{\small{> > > NEXT SERMON > > >}}
\newline
\newline
路加福音 24:13-35
\newline
\begin{longtable}{cl}
\hline
\hline
章節 & 經文 (和合本修訂版)\\
\hline
24:13 & \begin{tabularx}{0.7\textwidth}{X} 同一天，門徒中有兩個人往一個村子去；這村子名叫以馬忤斯，離耶路撒冷約有二十五里。 \end{tabularx} \\ \\ \relax
24:14 & \begin{tabularx}{0.7\textwidth}{X} 他們彼此談論所發生的這一切事。 \end{tabularx} \\ \\ \relax
24:15 & \begin{tabularx}{0.7\textwidth}{X} 正交談議論的時候，耶穌親自走近他們，和他們同行， \end{tabularx} \\ \\ \relax
24:16 & \begin{tabularx}{0.7\textwidth}{X} 可是他們的眼睛模糊了，沒認出他。 \end{tabularx} \\ \\ \relax
24:17 & \begin{tabularx}{0.7\textwidth}{X} 耶穌對他們說：「你們一邊走一邊談，彼此談論的是甚麼事呢？」他們就站住，臉上帶著愁容。 \end{tabularx} \\ \\ \relax
24:18 & \begin{tabularx}{0.7\textwidth}{X} 兩人中有一個名叫革流巴的回答：「你是在耶路撒冷的旅客中，惟一還不知道這幾天在那裡發生了甚麼事的人嗎？」 \end{tabularx} \\ \\ \relax
24:19 & \begin{tabularx}{0.7\textwidth}{X} 耶穌對他們說：「甚麼事呢？」他們對他說：「就是拿撒勒人耶穌的事。他是個先知，在神和眾百姓面前，說話行事都大有能力。 \end{tabularx} \\ \\ \relax
24:20 & \begin{tabularx}{0.7\textwidth}{X} 祭司長們和我們的官長竟把他解去，定了死罪，釘在十字架上。 \end{tabularx} \\ \\ \relax
24:21 & \begin{tabularx}{0.7\textwidth}{X} 但我們素來所盼望要救贖以色列民的就是他。不但如此，這些事發生到現在已經三天了。 \end{tabularx} \\ \\ \relax
24:22 & \begin{tabularx}{0.7\textwidth}{X} 還有，我們中間的幾個婦女使我們驚奇：她們清早去了墳墓， \end{tabularx} \\ \\ \relax
24:23 & \begin{tabularx}{0.7\textwidth}{X} 不見他的身體，就回來告訴我們，說她們看見了天使顯現，說他活了。 \end{tabularx} \\ \\ \relax
24:24 & \begin{tabularx}{0.7\textwidth}{X} 又有我們的幾個人往墳墓那裡去，所發現的正如婦女們所說的，只是沒有看見他。」 \end{tabularx} \\ \\ \relax
24:25 & \begin{tabularx}{0.7\textwidth}{X} 耶穌對他們說：「無知的人哪，先知所說的一切話，你們的心信得太遲鈍了。 \end{tabularx} \\ \\ \relax
24:26 & \begin{tabularx}{0.7\textwidth}{X} 基督不是必須受這些苦難，然後進入他的榮耀嗎？」 \end{tabularx} \\ \\ \relax
24:27 & \begin{tabularx}{0.7\textwidth}{X} 於是，他從摩西和眾先知起，凡經上所指著自己的話都給他們作了解釋。 \end{tabularx} \\ \\ \relax
24:28 & \begin{tabularx}{0.7\textwidth}{X} 他們走近所要去的村子，耶穌好像還要往前走， \end{tabularx} \\ \\ \relax
24:29 & \begin{tabularx}{0.7\textwidth}{X} 他們卻強留他說：「時候晚了，天快黑了，請你同我們住下吧。」耶穌就進去，要同他們住下。 \end{tabularx} \\ \\ \relax
24:30 & \begin{tabularx}{0.7\textwidth}{X} 坐下來和他們用餐的時候，耶穌拿起餅來，祝福了，擘開，遞給他們。 \end{tabularx} \\ \\ \relax
24:31 & \begin{tabularx}{0.7\textwidth}{X} 他們的眼睛開了，這才認出他來。耶穌卻從他們眼前消失了。 \end{tabularx} \\ \\ \relax
24:32 & \begin{tabularx}{0.7\textwidth}{X} 他們彼此說：「在路上他和我們說話，給我們講解聖經的時候，我們的心在我們裡面豈不是火熱的嗎？」 \end{tabularx} \\ \\ \relax
24:33 & \begin{tabularx}{0.7\textwidth}{X} 於是他們立刻起身，回耶路撒冷去，看見十一個使徒和與他們正在一起的人聚集在一處， \end{tabularx} \\ \\ \relax
24:34 & \begin{tabularx}{0.7\textwidth}{X} 說：「主果然復活了，已經顯現給西門看了。」 \end{tabularx} \\ \\ \relax
24:35 & \begin{tabularx}{0.7\textwidth}{X} 於是，兩個人把路上所遇到，和耶穌擘餅的時候怎麼被他們認出來的事，都述說了一遍。 \end{tabularx} \\ \\
[1ex]
\hline
\hline
\end{longtable}
$^{1}$大英姊妹早晨.
我不知道你有沒有經歷過以下這些時刻.
一段關係結束了.
一個機會錯過了.
前路茫茫 好像沒有出路.
又或者醫生的一句話.
令你好像突然間沉入谷底.
那種感覺.
好像我們的人生走到最後一頁.
但是 耶穌基督的復活提醒我們.
人以為不可能的.
上帝可以令到它可能.
在人以為已經結束的地方.
往往就是上帝開始工作的重要時刻.
昔日的門徒.
他們跟了耶穌三年多.
以為耶穌會在地上作王.
誰知道他們就見到耶穌被門徒出賣.
被人抓了 還釘在十字架上.
你可想而知門徒有多麼害怕和多麼失落.
因為他們還沒明白耶穌一直在說.
祂死後要復活的這件事.
我們今天剛剛讀了經文.
發生在耶穌被釘十字架之後.
經文記載在《路加福音》二十四章十三至三十五節.
說的是門徒在以馬五師的路上遇見耶穌.
在四福音裡只有《路加福音》記載了這兩個門徒的經歷.
當中耶穌出現的時候.
和其他地方都很相似.
就是祂很簡短地出現了.
然後又不見了.
都是初時見不到或者不知道那是耶穌.
然後就發現耶穌或者認得耶穌.
其實我想真一點.
如果耶穌一出現在這兩個門徒旁邊.
門徒就認得了.
你猜有什麼感覺呢.
見鬼會不會很害怕.
因為他們從來都沒有預料過耶穌會復活.
但是耶穌就是用了一個很溫柔的方法.

$^{41}$一路走跟他們一路談.
然後門徒心裡面好像覺得火熱似的.
後來他們才認得耶穌.
我想我們也有不認得耶穌的日子.
曾幾何時.
你以為他不聽祂的祈禱,不在,沒有保守你.
但是原來其實耶穌一直都在.
一直都和你一起.
我們會不會有時給自己的一些想法.
或者想要的東西來阻擋我們見到上主的帶領.
以為祂不在呢.
你有沒有試過呢.
騎行徒.
不過都好像沒有結果.
耶穌就正正不是那種有求必應的神明.
祂會聽我們禱告,和我們一起.
但是祂不一定會照我們所想的方法來回應我們的祈禱.
因為我們連什麼對我們自己最好都不知道.
就好像你家裡有小朋友.
他不斷問你拿糖果喫.
但是你不會不斷給糖果他喫.
因為你知道什麼纔會對他最好.
有一個古老的傳說.
講到印度有一個部族.
當年輕人成長的時候.
他們到一定年紀就要行成年禮.
少年人準備要做成年人.
他們的爸爸會帶他們去森林.
然後他們要蒙著眼睛單獨過一夜.
不可以打開眼睛,直到第二天早上.
他們又不可以喊救命.
如果他們過了這一晚都沒有事.
這個少年人就真正成為成年人了.
他們經過了這一晚就要承諾.
不要將這個經歷告訴其他少年人.
因為每一個少年人要自己去經歷這個成年禮.
這一晚有一個年輕人.
他按照成年禮的要求蒙著眼睛.
坐在漆黑的森林裡面.
四周肯定有很多猛獸.

$^{81}$但是他不僅聽見風吹草蟲的聲音.
他還覺得他坐的樹幹好像有點動.
雖然他很害怕,全身都在繃緊.
不過他沒有打開蒙著的眼睛.
他只是動也不敢動地坐著.
因為他知道只有這樣纔會成為成年人.
終於他經過了一個很恐怖的夜晚.
他感覺到陽光照在他的臉上.
他可以打開蒙著的眼睛.
這個時候他一睜大眼.
他發現自己的父親.
原來一整晚都坐在樹幹上.
在他的旁邊守護著自己.
保護著自己.
原來愛我們的天父.
都是在我們不知不覺或不明的時候.
他都與我們同在.
昔日的門徒按照他們自己的理解.
他們很希望耶穌這位尼賽亞要在地上作王.
但是這個時候耶穌死了.
他們的希望好像破滅了.
當時人人都知道耶穌被處死的消息.
但為什麼眼前這個陌生人居然不知道呢?.
他們不認得眼前這位就是耶穌.
是什麼阻止他們辨認祂呢?.
或許門徒太沉醉於自己的感受和想法.
他們根本不會認為真的會見到耶穌.
就好像門徒一直跟著耶穌.
耶穌一直說自己會死後三日復活.
他們都不明白.
我想很多時候我們都以為自己去找神.
我們自己去決定信不信耶穌.
但其實我們連怎麼去找神都不知道.
我們都沒有能力.
是耶穌親自去找回我們.
在我們身邊等著我們歸向祂.
路加對於整段故事的描述.
核心就在第23節.
那裡說「天使顯現,說他活了」.
復活的耶穌就像一個隱言的形態.

$^{121}$跟他們相遇.
然後將聖經的形容.
基督會受苦走上榮耀的路.
去告訴他們.
在路加福音24章25節.
耶穌對他們說.
「無知的人先知所說的一切話.
你們的心信得太遲鈍了」.
可能我們都信了主一段日子.
就像門徒一樣.
跟耶穌一起生活了幾年.
但耶穌說他們信得太遲鈍了.
其實這兩個門徒在做什麼呢?.
他們正正要走這條二馬五師的路.
將有人見到耶穌復活的消息.
帶回耶路撒冷.
但似乎他們的心情一點都不興奮.
因為耶穌的死讓門徒陷入深沉的悲哀.
他們好像覺得所有都沒有希望.
那些探訪墳墓帶來空墳墓消息的男女.
只不過是給他們一個很疑惑不明的觀念.
耶穌如何回應這個情況呢?.
祂說門徒信得太遲鈍了.
耶穌就將聖經告訴他們.
令到門徒的內心都火熱.
我不知道我們讀聖經的時候.
有沒有都會感動.
有沒有會覺得火熱,受觸動.
這個時候他們終於都準備好了.
他們可以憑著聖經見證耶穌是死而復活的神.
今天我們不單止知道耶穌是死而復活.
我們更加有聖經的說話.
有歷世歷代信徒的見證.
讓我們可以學習和成長.
所以你有沒有都為自己設計生命成長的時間和空間呢?.
無論我們在人生什麼階段.
我們都可以繼續尋求,認識主.
你有沒有經歷你初初信主的時候.
發現真理的那種喜悅呢?.
有一個年輕人他在木廠裡面工作.

$^{161}$工作第一天他斬了十棵樹.
實在是很厲害.
自己都很高興,老闆都很滿意.
第二天他同樣這麼努力斬了八棵樹.
成績都不錯.
一個星期之後他一樣這麼勤力.
但是生產力越來越少.
老闆都覺得不知道哪裡出現了問題.
這一天他只斬了一棵樹.
老闆都覺得很難接受.
就叫他去辦公室瞭解.
年輕人和老闆說.
我每天都很努力地工作,我已經盡力了.
老闆就說.
你其實有沒有去磨礪你斬樹的斧頭.
然後再去斬呢?.
年輕人說.
我哪有時間去磨礪斧頭呢?.
我實在太忙了.
原來有時候.
我們實在太過沉醉於自己做的事情上.
我們很容易失去了焦點或者事情的重心.
其實是有選擇的.
你要怎樣去度過你餘下的日子.
不失去重心.
而在最重要的生命事情上.
繼續去成長,繼續學習追求呢?.
那些門徒何時才知道.
其實和他們一起走過這段路的是耶穌呢?.
這個陌生人.
後來受到邀請留宿在他們家裡.
這個陌生人反客為主.
他用主人的身份為他們抹餅.
款待他們.
這兩個舉動.
就令到門徒心裡感動.
他們終於認出耶穌的身份.
所以到最後.
其實也不是他們有能力去認出耶穌.
而是耶穌親自行動.

$^{201}$成為主人.
讓他們能夠認得出來.
我想你也有.
你和耶穌經歷相遇的故事.
我自己最難忘的一個屬靈經歷.
其實也發生在多倫多這裡.
就是我第一次發現原來上帝是這麼真實的.
當時我在多倫多出來工作沒多久.
都三十多年前了.
由我讀大學的城市搬了出來.
我很想找一間教會.
可以成為我屬靈的家.
不過呢.
其實我是一個慢熱的人.
由初來步到.
於是我也不知道去哪裡好.
我就自己去了家附近一間華人教會.
回了幾次.
還遇到一些我平時上班的同事在那裡.
不過回了好幾次之後.
我又覺得很難投入.
於是過了沒多久.
我就想.
我不如去家對面那間西人的教會試試.
那裡也不錯.
人少少靜.
於是我回了好幾次之後.
我又覺得有些問題.
因為沒有相交生活.
於是心裡想.
我應該找一間真正適合我的教會.
去哪裡好呢.
記得有一個星期日.
我回到那間西人教會.
通常我都喜歡早一點去到.
坐在長的木凳一邊安靜.
我心裡就在想.
不如下次去那邊.
另外還有一間華人教會試試.
那裡好像也不錯.

$^{241}$突然之間我心裡有一種聲音.
或者是一種感動.
是說.
你在這裡做什麼.
其實你來到多倫多的時候.
我已經帶了你去你應該去的教會.
為什麼你不留在那裡.
在那一刻我四周看.
因為感動是很清楚的.
不過周圍都沒有人.
我心裡想.
回到之前那間華人教會.
就是我最不想的.
因為我去過了.
都投入不了.
又回到那裡.
可能有些人又認得我.
又離開了.
停一輪又回去.
我覺得很尷尬.
當日崇拜開始之後.
我就覺得很特別.
因為當時牧師是講一系列的講道.
每一次就用主土文其中一句.
就是講他的講道.
主土文應該還有很多句才完.
但是很特別的是.
那天他一次過講了這麼多節經文.
然後還加上了總結.
給我的感覺.
就好像告一段落.
在一個階段.
下次要去另一個地方了.
崇拜之後.
我走出去坐在車上.
看著太盤.
我心裡想.
是真的嗎?.
因為在我的生命裡.
祈禱都是我向神求和講的.

$^{281}$上帝是會感動我回應我的嗎?.
是我第一次遇到.
其實這些經歷都不多.
但是.
我就將這個經驗.
那一個星期.
我和一位比較成熟.
我也信任的基督徒分享.
我和他分享完我的經歷.
我以為他會給我一些意見.
問下我的問題.
誰知道他只說了一句話.
他說.
你講出來也好.
會幫到你去做的.
我就想.
那就是要去做了.
我最不願意的.
不過帶著萬分的不願意.
覺得尷尬也好.
我第二個星期.
就回到我原本的那間華人教會.
我很記得.
在那個星期天.
和之後的主日的講道.
是非常觸動我的.
牧師是講到.
《初一級記》裡面.
上帝如何如鷹揹負以色列人.
來歸向他.
然後那個星期就說.
上帝如何像鷹搞動鳥巢.
鳥掉下去就要逼他們學飛.
這兩篇道給我很大的感動.
我覺得.
天父就是這樣看顧我.
是他親自背我回去.
來歸向他.
教導我成長.
於是後來我就留在那間教會里面.

$^{321}$原來上帝和我們的距離是這麼親密的.
聖經告訴我們.
聖靈是住在我們每一個人裡面.
不是外在的.
而是他在我們生命裡面.
每一天和我們同行.
就好像門徒在二馬五師的路上一樣.
有時候我們不知道.
不認得.
或者我們會懷疑.
但其實他一直都在.
大英姐妹.
你也有你很獨特的.
與神相遇的屬靈經歷.
關鍵不是我們知不知道他在.
而是他一直都在.
有時候可能我們會覺得.
未信主或者困難的時候.
特別覺得他在.
但其實只是因為那時候.
我們會特別去追求.
所以會很敏銳.
上帝和我們的互動.
這個互動.
每一天在我們的生活當中.
就可以出現.
每一天我們都可以經歷上帝.
在《路加福音》24章32節.
門徒回想起耶穌和他們一起的時候.
他們發現.
「他們彼此說.
在路上他和我們說話.
給我們講解聖經的時候.
我們的心豈不是火熱的嗎?」.
心裡火熱.
有時候我們說我們要服侍神.
要會應神.
其實不是硬要為主做些什麼.
或者我們怎麼去想象.
而是當我們與神同行的時候.

$^{361}$在人生路上的經歷裡面.
我們發現.
原來耶穌一直都在.
我們心裡就火熱了.
以致無論環境好壞.
我們按照聖經的話語.
可以充滿喜樂.
有能力去面對我們每一天的挑戰.
有一位非洲的皇帝.
他有一個從小玩到大的好朋友.
這個好朋友有一個習慣.
就是無論遇到什麼事.
好不好.
他都會說.
「這件事真是好好」.
有一天皇帝與這個朋友一起去森林打獵.
通常朋友都會為皇帝預備好獵槍.
上子彈.
這次朋友也一樣準備槍.
不過他準備的時候出了錯.
皇帝拿上手的時候.
發射.
突然間走火.
哎呀.
射中了皇帝的手指弓.
皇帝就是這樣.
沒有了手指弓.
肯定是痛得死死的.
還流了很多血.
朋友見到這樣的情況.
他說了和平時說的同樣一句話.
他說.
「這件是好事,好好」.
哇,皇帝火都來了.
當然不好,這件事不好.
非常之不好.
於是他很生氣.
就抓了這個朋友去坐牢.
大約一年之後.
皇帝又去打獵了.

$^{401}$這次他去了一個他沒去過的地方.
不知不覺就迷路了.
他進入了一個食人族的地盤.
那些土人見到有外人來到.
就抓住他.
綁住他.
然後搭上柴準備.
開火燒了他來喫.
正當那些土人上前點火的時候.
他們發現這個人的手少了一隻手指.
原來那些土人很迷信.
他們不會喫身體不完整的人.
所以最終他們放了皇帝送走了他.
皇帝回到家裡.
他就記起一年前為什麼他失去了一隻手指.
他就記起他當時很生氣.
把朋友扔進了監獄.
他很後悔.
於是他就去監獄跟朋友說.
原來你對.
幸好我失去了一隻手指.
否則我就被食人族喫了.
然後他就告訴朋友發生了什麼事.
皇帝就說.
我真的很抱歉.
讓你坐了這麼久的牢.
這真的不是一件好事.
朋友就說.
不是,這件是好事,很好.
朋友說,當然是好事.
如果我不是坐牢.
我剛纔就跟你一起去食人族的森林.
現在就被人喫了.
弟兄姊妹.
其實這個故事是講一個重點.
在所有事情發生的時候.
可能都有些原因.
也有一個指向.
表面上不是好的事.
我們不知道它將會怎麼樣發展.

$^{441}$也可能會變成美好.
我們就在人生的環節面上.
我們看不到未來.
我們經歷過去.
我們以為有很多事我們懷疑.
不知道神在不在.
不知道為什麼會這麼差.
為什麼會遇到這些困難.
但是眼前發生的.
雖然好像是壞事.
但原來這些事情是未結束的.
我們仍然不能落定論.
我會說.
當我們發現.
而且確信.
原來在人生的路上.
主耶穌一直與我們同在的時候.
我們知道是祂親自去掌管.
是祂親自去找回我們.
以至我們雖然有不明.
但我們願意去相信祂.
祂每一天就像在二馬五司路上.
與我們同行.
以至我們可以在讀祂的話語.
去尋求祂的時候.
我們心裡火熱.
不斷地經歷上主奇妙的恩典和大靈.
如果你也想經歷.
其實很簡單.
你只需要直接向上帝祈禱.
你也不用告訴我.
也不用告訴其他人.
耶穌在等著我們.
祂就在我們身邊.
一直陪伴著我們.
兩年前2024年.
上帝將我從香港帶到多倫多.
我其實是年輕三十多年前.
在多倫多成長.
後來我回到香港工作.

$^{481}$讀神學.
和在神學院教書.
我以為我會在香港退休.
誰知道上帝突然帶我來到多倫多.
上任加拿大華人神學院多倫多院長.
過去十幾年.
我研究的都是職場神學和系統神學.
職場神學其實是幫助.
所有的弟兄姊妹.
在我們的生活裡來活出信仰.
即使你說我沒有一份出糧的功.
但是其實上帝在人生不同的階段.
都給了我們有時間.
因此能力.
我們接觸到很多身邊的人.
不是你叫牧師去.
他就能夠接觸到.
因為那些在你生活範疇裡.
你的朋友.
你的家人.
你的同事.
其實他都只是認識你.
上帝就是這樣.
擺我們在人群當中.
以至我們的生命.
可以自然地流露出基督徒的特質.
讓人去預見我們生命裡的主.
所以我們的神學院.
除了培育未來的牧者和宣教士.
我們也有一些職場的課程.
有一些生活化的學習.
例如我們如何去看財富.
我們如何在生活裡.
去整理我們的屬靈的生命等等.
我們也有一些免費的網上講座.
大家都可以去參加.
今天我們有同工在外面擺了一個攤位.
如果大家崇拜之後有興趣.
歡迎你參觀我們嘉臣多倫多的攤位.
去拿一些資料.

$^{521}$我相信上帝在這個時代.
祂要興起所有的弟兄姊妹.
我們都說我們要做主的門徒.
上帝是呼召每一個願意跟隨祂的弟兄姊妹.
猜測我們去到生活的場景位置.
你的社區,你的家庭,你的朋友的身邊.
來回應祂給我們的使命.
不是我們想象或者想一些所謂偉大的事.
去為主而做.
而是很簡單地在我們的生命裡面.
我們告訴別人一個基督徒是怎樣的.
讓別人去遇見我們生命裡面的主.
所以祝福大家在你的生活當中.
不斷地經歷上帝.
我們同心禱告.
天父我們多謝你.
因為你不單猜測你的愛子耶穌基督.
為我們釘身在十字架上.
你也留下你的話語.
也讓我們有聖靈住在我們的心裡面.
與我們同在.
以至我們能夠繼續經歷你.
繼續認識你.
繼續學習和成長.
求你賜給我們信心和勇氣.
在這個過程中.
無論我們在人生的什麼階段.
我們是願意回應你.
繼續在你裡面成長.
求你幫助我們.
我們禱告奉耶穌基督的名求.
阿姆.
\newpage



\section{羅馬書 8:22-39}
\label{sec:032320}
\textbf{靠著愛我們的那一位}~~~~周曉暉牧師
\newline
\newline
連結: \href{http://www.church.com.hk/acms/content.asp?site=cdc\&op=show\&type=product\&code=032320\&layout=sermon}{\texttt{032320}} ~~~~ 語音日期: 2025-11-30
\newline
\newline
\hyperref[sec:032319]{\small{< < < PREV SERMON < < <}}
~
\hyperref[sec:index_chronic]{\small{[返順時目]}}
~
\hyperref[sec:032321]{\small{> > > NEXT SERMON > > >}}
\newline
\newline
羅馬書 8:22-39
\newline
\begin{longtable}{cl}
\hline
\hline
章節 & 經文 (和合本修訂版)\\
\hline
8:22 & \begin{tabularx}{0.7\textwidth}{X} 我們知道，一切受造之物一同呻吟，一同忍受陣痛，直到如今。 \end{tabularx} \\ \\ \relax
8:23 & \begin{tabularx}{0.7\textwidth}{X} 不但如此，就是我們這有聖靈作初熟果子的，也是自己內心呻吟，等候得著兒子的名分，就是我們的身體得救贖。 \end{tabularx} \\ \\ \relax
8:24 & \begin{tabularx}{0.7\textwidth}{X} 我們得救是在於盼望；可是看得見的盼望就不是盼望。誰還去盼望他所看得見的呢？ \end{tabularx} \\ \\ \relax
8:25 & \begin{tabularx}{0.7\textwidth}{X} 但我們若盼望那看不見的，我們就耐心等候。 \end{tabularx} \\ \\ \relax
8:26 & \begin{tabularx}{0.7\textwidth}{X} 同樣，我們的軟弱有聖靈幫助。我們本不知道當怎樣禱告，但是聖靈親自用無可言喻的嘆息替我們祈求。 \end{tabularx} \\ \\ \relax
8:27 & \begin{tabularx}{0.7\textwidth}{X} 那鑒察人心的知道聖靈所體貼的，因為聖靈照著神的旨意替聖徒祈求。 \end{tabularx} \\ \\ \relax
8:28 & \begin{tabularx}{0.7\textwidth}{X} 我們知道，萬事都互相效力，叫愛神的人得益處，就是按他旨意被召的人。 \end{tabularx} \\ \\ \relax
8:29 & \begin{tabularx}{0.7\textwidth}{X} 因為他所預知的人，他也預定他們效法他兒子的榜樣，使他兒子在許多弟兄中作長子。 \end{tabularx} \\ \\ \relax
8:30 & \begin{tabularx}{0.7\textwidth}{X} 他所預定的人，他又召他們來；所召來的人，他又稱他們為義；所稱為義的人，他又叫他們得榮耀。 \end{tabularx} \\ \\ \relax
8:31 & \begin{tabularx}{0.7\textwidth}{X} 既是這樣，我們對這些事還要怎麼說呢？神若幫助我們，誰能抵擋我們呢？ \end{tabularx} \\ \\ \relax
8:32 & \begin{tabularx}{0.7\textwidth}{X} 神既不顧惜自己的兒子，為我們眾人捨了他，豈不也把萬物和他一同白白地賜給我們嗎？ \end{tabularx} \\ \\ \relax
8:33 & \begin{tabularx}{0.7\textwidth}{X} 誰能控告神所揀選的人呢？有神稱他們為義了。 \end{tabularx} \\ \\ \relax
8:34 & \begin{tabularx}{0.7\textwidth}{X} 誰能定他們的罪呢？有基督耶穌 已經死了，而且復活了，現今在神的右邊，也替我們祈求。 \end{tabularx} \\ \\ \relax
8:35 & \begin{tabularx}{0.7\textwidth}{X} 誰能使我們與基督的愛隔絕呢？難道是患難嗎？是困苦嗎？是迫害嗎？是飢餓嗎？是赤身露體嗎？是危險嗎？是刀劍嗎？ \end{tabularx} \\ \\ \relax
8:36 & \begin{tabularx}{0.7\textwidth}{X} 如經上所記：「我們為你的緣故終日被殺；人看我們如將宰的羊。」 \end{tabularx} \\ \\ \relax
8:37 & \begin{tabularx}{0.7\textwidth}{X} 然而，靠著愛我們的主，在這一切的事上，我們已經得勝有餘了。 \end{tabularx} \\ \\ \relax
8:38 & \begin{tabularx}{0.7\textwidth}{X} 因為我深信，無論是死，是活，是天使，是掌權的，是有權能的 ，是現在的事，是將來的事， \end{tabularx} \\ \\ \relax
8:39 & \begin{tabularx}{0.7\textwidth}{X} 是高處的，是深處的，是別的受造之物，都不能使我們與神的愛隔絕，這愛是在我們的主基督耶穌裡的。 \end{tabularx} \\ \\
[1ex]
\hline
\hline
\end{longtable}
$^{1}$基督教宣道會.
此錄音或錄影的版權是屬於基督教宣道會不各堂所有.
我們的屬靈傳統之一,是以真理教導為本.
歡迎您將此錄音或錄影的連結傳給親友.
懇請注意,未經授權,請勿轉錄或剪輯本錄音或錄影.
甚願好似使徒行傳二十章三十二節下所講.
這道能建立你們,叫你們和一切誠誠的人同得基業.
的確在困難當中我們會感覺到很低沉,很失落.
尤其是突而其來的困苦.
這些我們無能為力的事,那種感覺.
可以好像心裡塞了一個很大的石,紮住了.
你拿不走,它在那裡,但塞住了.
你有很多事你都無法去反應.
你只能夠發出一聲嘆氣.
是呀,人生很難.
所以羅馬書第八章下半,繼續和我們講現世的苦難.
是那種不單是個人,是面對整個受造的世界.
受造物所承受的災難,痛苦.
我們要去承認是有這些災難.
我們並且要承認我們有需要哀悼的感受.
頂著當我們覺得很難,甚至感覺黑暗.
今日這段經文,福音給我們有怎樣的亮光.
我們先看二十二節.
二十二節裡提到的嘆息或是新淫.
其實是表達那種哀嘆的感受.
是從心底最深處的無奈.
所以待會你讀的時候,你將這種感覺讀下去.
二十二節請大家一起讀.
事實上,我們知道一切受造之物.
至今仍然不同新淫,同受陣痛.
那種嘆息悲哀哀鳴.
不是我們個人的,是普遍而且是整個受造界.
曾經有個喪失家人的人這樣寫.
在那一天,天不再藍,葉不再綠.
生命變得很遲緩,連呼吸都覺得很累.
無力嘆息.
保羅說我們是一同面對這件事.
而且他用另一個字詞說我們是同受陣痛.
同受陣痛是表達用一個懷孕的術語.
說媽媽在生孩子的經歷當中會有劇烈的疼痛.

$^{41}$如果痛是分十級的話.
陣痛就是十級頂級的痛.
因為這個痛楚是帶來承受的人.
一個無法忘記的苦.
一下一下一下.
不同的時候又再出現,直到生產的那一刻.
但這個痛楚卻是有一種特質.
保羅用這個字詞的時候.
是提醒我們這種苦難在世上我們所經歷的.
有一個特點,就好像陣痛一樣.
是很難,十級難,但是會停止.
十級難是有盡頭.
世上的災難是時有發生,但是終有一天不再有.
因為到那一天,有新天新地.
神所應許的,這個世界得著更生.
一起在裡面的人,得著更生的生命.
如果是這樣的話,這個盡頭是有盼望的.
有這個盼望就能夠忍受這一刻的十級痛.
如果你這一刻很痛苦.
但你記得,有盡頭.
並且這一年我們感受到.
不同的地方都有災難,災禍,戰爭.
有人的問題,有社群的問題.
有我們個人的困難.
一浪接一浪,很難受.
但你記得,是會有盡頭的.
這是福音給我們的亮光.
屬神的人要記得.
你繼續等待那一天,不要放棄.
當你受苦的時候,歎息是有限的.
但神應許的終點是會來的.
保羅繼續在二十八節提出了一句.
我們都很熟悉的金句.
不如我們一起讀.
萬事都互相效力,叫愛神的人得益處.
我不知道你有沒有送過這句金句給別人.
或者你自己有抄下來.
遇到一些事的時候特別困難,災難.
你會安慰自己.
萬事都互相效力,叫愛神的人得益處.

$^{81}$在網上留意到一些資料.
說我們熟悉的一些基督徒的球星.
其中一個是NBA的林書豪.
他就最喜歡帶著這句金句.
上NBA的球場打比賽.
他相信無論他贏球還是輸球.
都是互相效力的.
對他是有益有幫助的.
但如果將這一節的經文.
放在保羅要討論福音的大能.
普世的苦難.
似乎不是說我們一個人.
看清楚二十八節.
我們抄就抄這兩行.
其實整句是這樣的.
不如我們只讀二十八節.
並且我們知道萬事都互相效力.
叫愛神的人.
就是按著他旨意被召的人得益處.
看到嗎?.
前面有四個字.
我們知道.
不是說你個人.
是一群人.
大家都會知道這個真理.
這個真理是給一群人的.
是我們知道.
不是我知道.
是我們這一群人要知道.
知道萬事都互相效力.
有困難有災難.
好的.
神會轉化他.
這個層次沒有問題.
但不要停在這個層次.
這句金句更加讓我們看到.
無論是好事壞事.
神都是用來讓這一群人.
屬於他的子民.
能夠走向一個方向.

$^{121}$知道萬事都互相效力.
叫愛神的人.
按著他旨意被召的人.
就是我們.
得益處.
然後二十九節說.
因為.
說了原因.
方向.
是甚麼呢?.
要你像耶穌基督.
要這群人像耶穌基督.
這個就是保羅這樣寫的原因.
放在災難的處境.
若果災難發生在我們當中的一個.
困難在我們當中的一個.
這一群人的其他人.
會不會好像耶穌所做的事.
與哀哭的人同哭?.
會不會去安慰,輔助他?.
我想會.
因為你知道耶穌會這樣做.
你會跟與神的兒子形象相似.
若果這件事是發生在我們以外的.
我們會不會.
好像耶穌基督一樣.
走進人群當中.
將愛與醫治與福音的亮光.
帶給這些受創的人.
安慰那些已經失去家人的人呢?.
會.
這個是與神的兒子形象相似.
所以萬事都互相效力.
是在說我們這群信主的人的團隊合作.
不是說你一個人.
好事壞事都幫你成長得很好.
若果你要拿一隊人來說.
也不是在說輸贏.
而是說這群人一起去比賽繼續講球隊.
有球員受傷.

$^{161}$這些球員怎樣做?.
隊友怎樣做?.
大家在更衣室吵架.
你會怎樣做?.
輸了比賽.
遇到不同的事情.
判罰不公平.
你給什麼反應?.
一個屬主的人.
他要像耶穌的話.
這群人會不會有不同的選擇呢?.
當經歷這些萬事.
這些高低的起伏.
會塑造這群人.
應該變成更像耶穌.
能夠彼此包容.
互相保衛.
不再是互相指責.
這是在主裡面的成長.
所以萬事互相效力.
是我們的團隊合作.
最終的益處不是贏一個長輩.
而是贏我們這群人.
是否越來越像主.
乘勝.
苦難在神的手中.
當然是轉化來成就祂的旨意.
也成就這群被祂召喚出來.
作祂兒女的人.
走在祂的計劃當中.
繼續像耶穌一樣.
走進祂的心意裡.
不過弟兄姊妹.
你可能會說.
我知道萬事都互相效力.
這群人知道苦難有盡頭.
但我這一刻很苦.
我這一刻很痛苦.
有些事我不明白.
實際有時候.

$^{201}$我們都未必能夠立刻接受.
你會繼續問.
難道世上的苦難悲劇是好的嗎?.
難道絕症讓我失去摯愛親人.
是有益的嗎?.
火災是有益的嗎?.
我們受到各式各樣的折磨.
挫折是有益的嗎?.
為什麼神會容許我們遭遇這些事?.
這些問題是揮之不去的.
因為痛苦就像一塊大石.
將你一直沉到底.
沉到一個地步.
我們不懂得祈禱.
尤其是一浪接一浪.
一波未平一波又起.
在這麼低沉的底下.
保羅在這段經文說.
有聖靈支持你,幫助你.
我們之所以可以堅持下去.
因為我們擁有一個看不見.
但必會出現的將來.
並且是有證據的.
我再說一次.
我們之所以能夠堅持下去.
是因為我們擁有一個看不見.
但必會出現的將來.
而且這個出現一定會有證據證明.
是可以的.
當保羅在二十三節,二十四節提到.
聖靈就讓你知道.
你盼望著的,神所應許的將來.
是會出現的.
這是真的盼望.
有一個看不見,真的盼望.
相對一個假的盼望.
是隨處可見的.
二十四,二十五節請大家讀.
當時弟兄姊妹收到保羅這封信.
就知道保羅在鼓勵他什麼.

$^{241}$他所說的漢德建的盼望.
就不是盼望.
是什麼意思呢?.
是在表達在羅馬的帝國裡面.
羅馬帝國都會宣傳.
你在我們帝國很有盼望.
我這個帝國給你的盼望.
就是你會有太平盛世.
但實際在裡面的人.
他們面對的是無止境的.
征戰,征服,擴張,奴役.
這個不是盼望.
而且這個盼望.
是被人製作成一個可以看得見的東西.
即是每天羅馬城裡面的居民都看得見.
到處都是.
就是那些雕像.
那些號,那些領導的樣子.
到處都有.
讓你望見.
你會知道羅馬有盼望.
其實歷史上很多政權都做這件事.
國王元首,首領,將軍,領袖,總統.
都很喜歡將他的頭像,雕像.
後來是畫像.
現在就是照片.
甚至是影片.
到處貼,到處登,到處播.
家家戶戶有些都要貼.
讓你望見.
你當那個是盼望.
一個可見的盼望.
弟兄姊妹呼音讓我們看見的真相就是.
這個世界給你的.
是一個一瞬即逝的保證.
可能那個王,那個領袖會做得好的.
但是有多久呢?.
一瞬即逝.
一個很脆弱的梵華.
甚至在當時羅馬城的信徒.

$^{281}$他就是面對著羅馬城勢.
但奢華背後他們知道.
是建築在很多人被奴役,被剝削.
被統治的殘酷結構裡面.
這個不是盼望.
經文說我們有看不見真的盼望.
上主給我們的是真的盼望.
而且是有證據的.
那個看不見必會出現的將來.
是有證據的.
證據在哪裡呢?.
一開始這裡說:就連我們這些有聖靈為初結果子的人.
甚麼是初結果子呢?.
就是當一個農夫很辛苦地種植了一季,兩季.
有一天他發現.
他的果樹上有第一顆果子熟了.
可以摘下來.
這個就是初結的果子.
這顆果子代表著樹上其他的果子.
接下來的那幾天.
都會一直成熟成長.
他有很長時間的忙.
有收成.
這一顆的證據拿在手上.
他就知道前面會有收成.
這個是初結果子的意思.
在英文本有些人都會將希臘字詞翻譯為.
down payment.
首期.
或者不可以退的訂金.
聖靈就好像一個不可以退的訂金.
神已經給了你.
那筆錢在你手上.
將來還有一大筆給你.
是這樣的意思.
你已經拿著了.
你已經率先因為聖靈成為這個down payment.
你已經率先體驗神的同在.
不是等到將來.
你去見到上帝面對面的時候.

$^{321}$才有神的同在.
這一刻因為聖靈內著.
你有神的同在.
你正在體驗.
你已經在今天生活.
有證據,有證明,有渣拿.
不需要等到最後才有.
不是純粹等.
頂智慧信耶穌的人.
是經歷生命的改變.
這個實證讓我們知道.
我們存有的是看不見的盼望.
是一個活的盼望.
我們的生命有改變.
我們隨著信耶穌的年日增長.
你都會長進的.
今天你看到你重生以後的形態.
你的生命多多少少都會改變.
並且應該.
你跟隨著實神的話.
一路一路會成長的.
多少歲都可以成長.
所以你是有渣拿的.
聖靈讓你知道你心裡存著.
這個看不見的盼望是真的.
以致你對應在現實你看見的困難.
你因著你有這個成長的痕跡.
你就知道上帝一定會帶你去.
最後他答應給你的結局.
在我年輕的時候.
我跟太太結婚沒多久已經很想有小朋友.
我們後來很快就懷孕了.
很開心很高興.
不過兩三個月後就小產了.
在那一次的經歷裡.
由一個高點跌到最低點.
心情非常低沉.
並且往後的年日.
我們沒有再有成功懷孕一段時間.
在那段日子裡.

$^{361}$我們不斷地去問.
為什麼?.
為什麼我會遭遇這些事?.
是不是我做錯了什麼?.
是不是我吃錯了什麼?.
是不是我以前有什麼問題?.
是不是我不夠好?.
很多問題不明白.
很低沉.
甚至連祈禱都不懂得祈禱.
直到兩年後.
有一天我們突然間好像開竅一樣.
我們向神祈禱說.
神啊.
我們明白生命是屬於你的.
是你所賜的.
無論你給或不給.
我們都願意信服.
就在那一天.
我們開始看到身邊.
原來在我成長的歷程裡.
有不同的弟兄姊妹,團友,朋友.
他們可能小時候已經遭遇.
不同的挫折困難,苦痛.
但是我的感受不大.
有些可能是學業的問題.
有些可能是工作出來工作的困難.
有些可能是交友戀愛婚姻的困難.
有些是身體疾病.
甚至面對死亡.
這一切在我自己所經歷的苦之後.
我開竅了.
我明白多一點.
多了一點同理心.
我可以學耶穌做一點.
去關心這些人多一點.
明白多一點.
弟兄姊妹.
這件事是不好的.
但是有神的幫助.

$^{401}$這一個是因為上主給我們有聖靈.
祂幫助我能夠度過那個日子.
萬事都互相效力.
叫我屬於這一群的人.
能夠越來越像耶穌基督.
在二十六節.
保羅繼續說.
我們在這個漫長的等待裡.
有軟弱,有艱難,有挫折.
甚至連祈禱都不會祈禱.
軟弱.
我想你也試過傷心到連禱告都不會祈禱.
躲在被窩裡哭.
或者連哭都哭不了.
長期的壓力.
那種心力交瘁.
說不出話的狀態.
可能你也經驗過患病.
你很明白的.
每一次去醫院覆診.
每一次去進行療程.
反反覆覆很多的等待.
做完療程了.
還不舒服.
又嘔吐.
又吃不了東西.
可能有些時候有副作用.
好了,弄了很久.
終於醫好了.
但是又不安心.
果然幾年後又復發.
又重複這個過程.
要到什麼時候呢?.
還怎麼祈禱呢?.
根本已經不懂得怎麼祈禱了.
保羅說:就算在這樣的軟弱裡.
聖靈也會幫助我們.
祂會幫助我們.
經文的「幫助」這個字.
不是普通的help.

$^{441}$也不是assist.
不是輕輕的逗你.
這個原本在希臘文的字詞.
是形容一個狀況.
就好像一個人.
背著一個二十斤重的背包.
一路走.
走到跌倒了.
聖靈就不是在旁邊幫你加油.
打氣或者扶你一把.
站起來繼續走.
不是這樣的.
而是聖靈二話不說.
就將你這個二十斤重的背包.
直接拿過來.
放在自己肩膀上.
陪你負重前行.
這個是聖靈的「幫助」.
很不同.
不是我努力加油.
是完全由聖靈去幫助我們.
繼續走這麼艱難的下一步.
聖靈怎樣幫助我們?.
就是將這麼難的大石.
無法言喻的痛.
幫助我們心底深處最辛苦的事情.
祈禱祈不到的事情.
轉化為一個完全符合天父旨意的完美禱告.
祂代我們向天父說.
「父神啊,你的兒女在這樣的情況下.
求你憐憫他們.
引領他們所作的選擇.
讓他們所走的下一步.
越來越像你的兒子耶穌基督.
走在你的旨意裡.
得著你所應許的福分」.
而天父就是那位監察人心的.
他又完全明白聖靈的心.
所以你那種以為沒有人明白的.
哀痛和破碎.

$^{481}$扎在大石下面的那一件事.
天父完全收到.
靈姐妹,聖靈就是這麼貼近你和我的心.
去幫你揹著.
代你求告.
祂將你的歎氣揹起.
傳達給父神.
祂就好像一個最明白你的人.
望你一眼就知道你要怎樣.
你需要什麼.
二話不說就幫你.
就祈禱求告.
這是聖靈的幫助.
不過世上太多苦難了.
所以經文不單叫我們.
看回我們裡面有神的同在.
有聖靈.
更叫我們從眼前的患難.
轉向望向耶穌基督的十字架.
所以三十四節繼續說.
耶穌基督已經死了.
也復活了.
現在祂在哪裡?.
如今就在神的右邊.
我們每個月都背的使徒信經.
復活升天.
坐在全能父上帝的右邊.
羅馬書裡加了一個.
祂在那裡坐在右邊做什麼?.
並且為我們祈求.
你看到這個很厲害的圖畫嗎?.
裡面你心底最深處.
有聖靈在為你祈求.
完全明白你.
在高天之上.
在神的寶座.
全宇宙最大的權力核心.
耶穌基督為你祈禱.
將你的軟弱難處.
祂就直接跟父神說.

$^{521}$祂在我旁邊.
祂為你祈禱引領你.
我記得有一次.
我在工作.
很多年前.
突然收到一個團友的短訊.
短訊是這樣寫的.
我突然想起你.
為你祈禱.
知道你忙 不用回覆.
他沒有問我什麼事.
我沒有說什麼.
他就給了我這幾個字.
那一刻我很感動.
因為我正可以遇到一個很大的事.
我也不知道該怎麼辦.
我還沒有祈禱.
我覺得很苦惱.
還在想怎麼辦.
我就收到這個短訊.
那一刻我知道上帝知道有人為我祈禱.
但是朋友也不是常常為你祈禱.
你有沒有收過這些短訊?.
可能一次兩次 一生人.
但是基督是坐在上帝的旁邊.
不斷為你祈禱.
祂是全方位為你祈求.
不只是全天候.
全時間.
隨時.
剛才我們唱了一首詩歌.
這是一件很厲害的事.
你明白嗎?.
當苦困環繞著你.
但是你能夠有的幫助.
不僅有聖靈.
有耶穌基督向父神禱告.
不過你說苦難這麼大.
我依然這麼苦,這麼困難.
那件事沒有走.

$^{561}$有時候我們信心動搖.
那怎麼辦?.
我們會不會失去了信心?.
保羅列出了很多可能切斷我們和神關係的患難.
他問:這七種苦難通通都是突然而來的.
患難,困苦,刀劍.
都是我們不要的.
即使你遇到這七種.
他最終的結論是.
這些都不能夠讓我們和神的愛隔絕.
是不會讓神所揀選的兒女和神永遠地隔絕.
我用手機的圖畫跟你說一件事.
有時候我們在教會.
正堂是很容易響的.
因為修得很好.
進去後面的課室.
不知道你有沒有進去過.
最進去的房間是一格.
經常都不通.
有些人就很擔心.
害怕家人找我.
經常擔心斷聯,斷線,沒有網絡.
有一個焦慮在.
保羅說這些患難,困苦.
就好像這些極端的環境.
令我們有一種斷網的焦慮.
好像你去了一個深山.
去了一個荒島.
沒有網絡.
怎麼辦呢?.
但是保羅說.
我們人世所遇到的患難,困苦.
信神的人遭遇到.
我們不怕.
因為正是有耶穌基督.
透過十字架上的死亡.
復活,升天.
打通了這個無阻礙的網絡.
或者這樣說.
他已經幫每一個相信他的人.

$^{601}$綁定了一個終身頂級無限數據太空漫遊計劃.
你去到外太空都收到的.
不會斷網的.
神和你的連線.
沒有任何的東西可以切斷的.
這條電纜是不會被人剪斷的.
我們害怕失聯.
有時候是因為怕我們自己做錯的事.
我們會焦慮地想.
神會不會封鎖我們.
但是羅馬書告訴我們.
只要你是靠主.
你無能為力.
但是你靠主.
這一切都不會.
因為我們的福音是建立在一個不可以動搖的基礎上.
第八章不斷跟我們說.
耶穌基督在十字架上的犧牲是非常重要.
打通了這個網絡.
而且父神已經永恆定義要建立這個網絡.
又讓你有這個宇宙太空漫遊計劃.
永遠不會拿走.
不會收回你的SIM卡.
如果這樣的話.
還有聖靈內置在你裡面.
內置的一個電話.
不會沒有電.
一路你知道打得通.
有經驗是有證據.
還有耶穌基督確保這條通道一路打通.
祂坐在父上帝的右邊.
全權地去管理.
為你去代求.
不止這樣.
祂坐在父上帝的右邊表達.
按照聖經的說法.
右邊代表掌權.
國王把他的權交給坐在他右邊的人.
全權管理.
耶穌基督是全權執行對你的救恩請求.

$^{641}$確保你完全不屬.
你的失聯,失腳.
不會失去主的愛.
這條網絡是存在的.
主耶穌不但憐憫同情我們.
祂還全權執行對我們福音的請求.
給了救恩你得救了.
但還未得到的中間這條人生路.
一路帶著你走到底.
一定去到目的地.
所以弟兄姊妹.
因著有主耶穌.
你要記得神的愛是不斷往的.
不是五支.
是十支.
神的愛那種力量.
無限.
是超越你和我的腦袋.
所以當你遇到困苦.
你感受不到上帝怎樣幫助你.
你要看回十字架.
看回我們這個十字架.
我們這個十字架沒有耶穌雕像.
因為耶穌不在十字架上.
不在墳墓中.
祂坐在父上帝的右邊.
全權執行我們這個救恩計劃.
一直帶領我們走到底.
不止剛才那七樣.
可能截斷連線的擔心.
保羅還數了十項的距離.
人的生死.
靈界的勢力.
時間.
空間.
各樣的團勢.
各種的物質.
各種的受造物.
即使是這一切這麼有距離.
這麼宏大的東西.

$^{681}$這麼闊的東西.
他們出絕招.
死亡.
令你突然死亡.
都不能夠令你和神的愛隔絕.
我小時候.
團契導師問我.
那時候是中一.
他說你每晚睡覺之前都會祈禱.
你會為那一天感恩.
也會為那一天所做錯的事情.
求主赦免.
如果一天你到晚上祈禱.
你就死了.
那還有罪未認清.
那你能不能回到天下?.
能不能回到天下?.
能回到天下.
你懂得回答了.
同樣信主的人.
若果因著各種原因不幸突然離世.
連肚子都站不住了.
他們能不能回到天下?.
能回到.
是的.
死亡這個絕招施行了.
在這些弟兄姊妹身上.
但是.
他們不會因為這件事而不能回到天下.
有時候有些人因著自己錯誤的選擇.
喪失生命.
他不會因為這件事.
只要他是靠主.
他不會因為這件事.
而不能回到天下.
經文提醒我們.
神已經稱我們為異.
祂不會收回.
耶穌基督十字架的拯救.
讓我們最後依然可以靠著.

$^{721}$得到完全的救贖.
被死亡擊倒了.
是的.
不過血疤又再重新出現.
德勝有餘.
所以我們記住這一節.
經文一起讀.
但靠著愛我們的那一位.
我們在這一切事上都得勝有餘.
這次的火災我們知道.
在我們宣導會.
也有一些堂會的弟兄姊妹.
他們有家人.
在其中有信主的弟兄姊妹喪命.
我相信在火場裡面.
有神的言語.
很哀傷.
但這個火.
不會令他和神的愛隔絕.
他已經得到完全的拯救.
耶穌會保證這一切.
所以弟兄姊妹.
你不要害怕死亡.
你不要害怕自己的錯謀.
只要你看十字架.
靠著主.
即使身處深淵的人生低谷.
這條基督連著你的網絡連線.
永遠都在.
讓我們得勝有餘.
死亡也不能將我們這份神的愛奪走.
弟兄姊妹我們無法預知明天.
正如我們沒有想過.
這個星期會發生這樣的事.
在這種極致的無常裡.
正好今天我們讀到八章下半的經文.
讓我們看見我們最大的盼望和安慰.
就是神所應許的.
這個世上是有苦難.
但神的應許讓我們有著耶穌的平安.

$^{761}$今生的痛苦,掙扎,失敗,軟弱.
我們無力.
但羅馬書讓我們記得.
我們信靠祂.
我們將來可以完全得救.
在最後和你分享一個叫Rev. John Alpert的人.
他說過這句說話.
他說可能有時候我們身處黑暗.
很難受.
但不要忘記你在光明裡已經學過的東西.
所以當下次你覺得人生很難.
身處黑暗.
感到無助.
覺得沒有人明白.
心被大石壓著.
連肚子也不懂得站.
你記得羅馬書八章.
你翻去看.
找回這七個讓你撐過逆境.
福音的亮光.
幫助你和我.
再去面對這個世上仍然有的陣痛.
相信這個為我們揹重擔的聖靈.
讓我們永不同神的愛失聯.
的耶穌基督.
我們一起祈禱.
主啊,我們心裡知道你的話語鼓勵我們.
當我們今天離開禮堂.
回到我們的生活裡面.
面對苦難的現實.
主啊,讓我們記得.
我們可以長大了.
因為萬事都互相效力.
叫我們這一班被你呼召出來.
又愛你的人.
能夠得著益處.
越來越似耶穌基督.
在這個神聖路上.
一路走到最後.
求主幫助我們北宣家的弟兄姊妹.

$^{801}$多謝你.
我們祈禱是奉主耶穌基督得勝名之祈求.
阿們.
\newpage



\section{士師記 8:1}
\label{sec:032321}
\textbf{聖經人聖經事:為甚麼沒有招我們同去}~~~~黃克勤牧師博士
\newline
\newline
連結: \href{http://www.church.com.hk/acms/content.asp?site=cdc\&op=show\&type=product\&code=032321\&layout=sermon}{\texttt{032321}} ~~~~ 語音日期: 2026-06-16
\newline
\newline
\hyperref[sec:032320]{\small{< < < PREV SERMON < < <}}
~
\hyperref[sec:index_chronic]{\small{[返順時目]}}
~
\hyperref[sec:032322]{\small{> > > NEXT SERMON > > >}}
\newline
\newline
士師記 8:1
\newline
\begin{longtable}{cl}
\hline
\hline
章節 & 經文 (和合本修訂版)\\
\hline
8:1 & \begin{tabularx}{0.7\textwidth}{X} 以法蓮人對基甸說：「你去與米甸爭戰，沒有召我們同去，你為甚麼這樣待我們呢？」他們就和基甸激烈地爭吵。 \end{tabularx} \\ \\
[1ex]
\hline
\hline
\end{longtable}
$^{1}$各位電影節目,大家平安.
我們又一齊讀聖經了.
生命樹宣教網絡為大家預備了聖經人,聖經事.
講了很多集都是在講機電.
如果你看時詩記.
你會發覺機電是佔了很長的篇幅.
可以說是重要的時詩之一.
我們今天繼續講機電.
機電我想也有十次八次.
一個這麼重要的人物.
今天的聖經人物機電.
事情就是機電對抗米電人.
為什麼沒有召我們同去是問號.
機電遭遇了一個問題.
人家問他.
為什麼你不叫我們一起去.
被人怪責的.
打了聖仗搏命之後還被人怪責.
這個也是我們今天很多時候會遭遇的問題.
跟大家看看這段聖經.
這段聖經很短的.
不過我們一起研究一下.
這兩節聖經上一次讀過了.
時詩記第七章.
他已經打贏了仗了.
機電打發人走遍以法連山地.
叫人去以法連.
以法連的十二支派分給他的地方在山上.
高一點的.
就說你們下來攻擊米電人.
再補充一下.
米電人是一個在約旦河的東面.
是屬於阿滿.
摩押阿滿.
以東這些地方.
再下一點的阿拉伯沙漠那裡.
是很多個大世民族的.
一起就叫米電人.
這次就來攻擊以色列人.
就在他們收割的時候.

$^{41}$就說弄爛他們的田地令他們很痛苦.
哀求耶和華.
這次就叫了神派機電跟他們對抗.
機電就是一個本來很弱信心.
很害怕的人.
經過很多次問神.
你給我一個證據.
知道你和我同在一定可以.
如果我是上帝.
就已經沒有耐性一巴掌打過去.
要找其他人.
不過上帝對他們很有耐性.
終於將他們練成一個很有信心的人.
只是帶著三百人去和米電人.
如果讀下去是有十二萬人的米電.
但是三百人好像戰狼三百那套電影.
可以和別人打.
但是機電有這個信心.
終於他竟然贏了.
打到米電人逃跑.
馬上就叫人去以法連山地.
快點在山上下來.
請你們以法連人幫我守住約旦河的渡口.
就是碼頭.
不要讓他們從西岸逃回東岸.
他們自己的本地本鄉.
以法連人就馬上走下去截住河.
不讓他們走.
結果真的有兩個米電人的首領過河.
一個叫俄納.
一個叫西伊伯過河.
俄納的名字意思就是搶奪.
西伊伯的名字就是野狼.
上次講了.
將他們兩個人殺死了.
而且還在俄納的盆石上.
即是說這兩個人逃跑回到他們的家鄉.
結果被以法連人追到那裡.
將他們兩個人殺了.
將他們的首級斬了.

$^{81}$帶去給機電.
接著就來到今天講的.
來到以法連人拿著那兩個首級.
來到見機電就說.
你去與米電人爭戰.
沒有招我們同去.
為什麼不叫我們?.
為什麼這樣待我們?.
這個待阿沙就是做.
為什麼對我們這樣做?.
於是他們就在罵機電.
與機電大大地爭吵.
爭論.
爭論為什麼不叫我們?不是兄弟嗎?.
其實原因機電沒得說的.
主要原因就是他是馬來西斯派.
他就叫了馬來西斯派.
西布倫,阿徹那些就一起來.
幾個派可能是更加熟的.
一起來因為你叫人家.
大家一起出來打仗.
你不是吧?.
為什麼叫我不叫他?.
這些出去沒命的東西你叫我?.
其實你不熟你不叫到的.
因為大家出去打仗找命搏的.
是不是保證一定不死的?.
連割傷手腳都沒有.
很辛苦很危險的打仗.
當日的機電就沒叫到以法連之派的人.
而且他在山地上.
你叫了附近他熟的那些就去打.
來了多少人呢?.
三萬二千人.
來了三萬二千人.
誰知上帝.
如果你讀上次的經文.
上帝說人太多了.
叫些怕的回去.
各位如果大家怕的你先回家.

$^{121}$三萬二千人就走了.
二萬二千人剩下一萬幾人.
如果我是機電真的暈倒.
人家是十幾萬人.
那些駱駝無數好像山那麼多.
這樣也好.
有三萬二千人不看得.
不打得看得.
誰知兩萬幾人走了.
剩下三分一.
誰知上帝又說.
一萬人都太多.
去水邊.
我看他喝水的姿勢我就會選他.
終於選了三百個.
是站著將水潑上來.
在手上站著喝的.
跪下趴下去這樣喝的不要.
全部要這樣拿上來站著的.
跪下去然後在水邊喝都不行.
當然就是.
那麼苛刻.
剩下三百個.
那三百個真的特別.
如果我都多數會跪在水邊喝的.
那三百個是站著.
拿了水之後站著喝.
那就說那些可能警惕一點.
可能在那邊喝到處看到.
其實跪在那裡都可以一邊看.
所以上帝的方法我們不知道.
選了三百個.
結果那三百個就大獲全勝.
就對付米電人贏了.
不讓他們走.
於是就叫以法連人.
我們這裡打緊.
你們在山上快點下來.
去到河邊指著他們.
不讓他們回頭.

$^{161}$不讓他們回鄉下.
於是以法連人馬上走過來.
就守住.
就將兩個領袖追到殺了.
接著就罵機電.
為什麼你不早點叫我們去的.
如果我是機電.
真的把工具都扔了.
還好說你們住到那麼遠.
其實老老實實.
如果一早叫他去.
可能又被他罵的.
這些東西叫我們?.
想我死嗎?.
我這裡有雞爺仔.
有八十歲的媽媽在家.
其實很難.
不過我們今天都會說.
有些困難的地方是什麼呢?.
贏了你又說一早不叫我.
如果輸了又怎樣呢?.
輸了你又說一早為什麼叫我.
這個很難.
就大大地罵他.
為什麼你一早打的時候不叫我們.
其實只是說我要三百個.
當然我叫他.
是贏了叫他幫我.
拿三百人分不開身.
你們幫我守住河.
機電很好.
有智有勇有謀有信心.
機電對他們說.
我所行的.
豈能比你們所行的.
以法連十取剩下的葡萄.
不強過阿彼以謝所摘的葡萄嗎?.
紫色紙.
紅色紙黑色紫色紙.
都是聖經.

$^{201}$只不過我想吸引你的眼球.
紅色紙特別要吸引.
紫色紙都重要.
機電很厲害.
起碼第一件事.
他的EQ情商都很高.
如果我打完仗三百人辛苦到死.
叫你幫幫忙.
殺了兩個領袖.
又罵我豈有此理.
真是打了一萬多人.
你不見你禮.
那就算了.
他就說我所做的.
和你們以法連無得比.
我們做什麼呢.
拿著一個水瓶.
裡面大概放些泥.
插住一支燒著的火把.
還有一個豪閣.
我們刀都沒有.
我們就去打了.
又沒有那麼多工具.
和你們精銳的部隊與兵器無得比.
接著他很厲害.
他沉得住氣.
他說了一句話.
這句話令到以法連人所有怒氣消散.
是什麼呢.
以法連自己在山上.
執政的葡萄.
收割葡萄.
掉到地上剩下的不收拾的.
都比我們以亞比爾者所摘的漂亮.
摘的就是我們很漂亮的那些.
你多到掉到地上的.
都比我掛在樹上的漂亮.
你們厲害很多.
用了葡萄做例子.
你掉下來剩下的.

$^{241}$不收拾的都比我家鄉以亞比爾者漂亮.
亞比爾者是幫助的父.
father of help的意思.
亞比爾者是什麼呢.
亞比爾者是一個地方的名字.
亞比爾者是father of help.
幫助我們的上帝.
父親的意思.
亞比爾者是什麼呢.
其實就是基地的家鄉.
你們以法連又高又強又大.
連葡萄執政的都比我家鄉的漂亮.
一聽到這些.
當堂以法連人都信了.
當然你知道.
士師記早過列王記.
薩姆爾那些很多.
因此其實以後大衛.
有一個很厲害的勇士叫做亞比爾者.
這個亞比爾者是基地的家鄉地帶.
神已經將米殿人的領袖.
就是俄納和西伊伯交在你手裡.
你們殺死他們首級都拿到了.
你們做了這麼偉大的戰士.
我所行的豈能比你們所行的呢.
基地說道以法連人的怒氣就燒了.
燒了一道氣.
這麼會說話.
是好的.
我這裡是很欣賞基地.
基地是什麼呢.
他這麼辛苦打完仗回來贏了.
只不過是跟著還被人罵.
為什麼不早點叫我們.
我們那些小兒科.
哪裡夠你們那麼厲害.
神將兩個首領交給你們.
你殺了他們比我們厲害很多.
我們殺的那些很少.
我們都沒有武器.

$^{281}$你掉到地上不撿那些葡萄.
都比我們掛在樹上的那些好.
我做的豈能彼得起你那麼多.
當場燒了氣.
各位電影姐妹有時真的這樣.
做一些事情自己辛苦到不得了.
誰知道會被人罵.
當然你心想現在做得好.
贏了仗你又說為什麼不早點叫你.
如果輸了今天我還有命嗎.
一定被你冤到死.
為什麼叫我們.
我沒有奸爺仔.
其實你自己出兵就可以了.
當然很多這樣的事情你是說不到的.
不過我見到如果我是基地.
我就忍不住氣了.
不過他竟然這樣說.
很委婉.
大國為眾.
這個基地果然由本來一個人.
變成有大將風範.
有信心,有智,有勇,有謀.
有變才,有急才.
很感恩的.
見到基地就很感恩.
其實老老實實.
基地本來是一個很怕死的人.
很膽怯的.
神叫他的時候.
他是躲在酒鎮打墨紙.
躲在家裡最偏僻的地方打墨紙.
是怕被人見到.
接著不停地問上帝要三次的證據.
你真的幫我.
給我看看神跡.
問了三次.
而且那些神跡的內容是非常離譜的.
但是神都認他.
然後神就叫他.

$^{321}$人太多人太多.
剩下三百.
如果我是基地.
我的信心已經失去了.
神已經令到他.
由沒有信心的人.
已經變成很有信心的人.
很勇敢.
再加上神給他智慧.
就是情商.
EQ這麼高.
我們很多時候都會被人怨的.
為什麼沒有招我們和去.
被人說的.
你要知道.
在這個時候我希望你記得基地.
基地沒有將豐功偉績放在自己那裡.
我去打的時候不見你.
現在贏了你又說.
我給機會你殺了兩個首領.
你又向我邀功.
又賴我.
又說為什麼不叫我.
你這些人都不知道什麼心.
如果這樣就沒有了.
大家失去了同心.
請你記得.
學到更好的情商.
EQ.
另外自己謙卑.
不要緊.
你知道上帝說什麼嗎?.
凡自卑的必升為高.
自高的必降為卑.
你要知道.
神是對於驕傲的人.
是很看不過眼的.
聖經講得很清楚.
驕傲的是神.
眼中看為幾件很不妥的事.

$^{361}$所以更謙卑.
其實你一定要知道.
相信上帝在這裡.
祂知道.
大局為重.
所以上帝選擇了基地.
做當代以色列人的事實.
而沒有選擇一個以法列人.
是對的.
我們一起祈禱吧.
天皇上帝我們在你面前禱告的是.
很多時候我們做了一點事.
獲得一次小的大的勝利.
也好.
我們心裡也會沾沾自喜.
是我.
做完這件事.
是我.
對整個大局有幫助.
是我做了這件事.
對整個教會都有幫助.
對整個人類都有幫助.
好像全世界都是自己救回來一樣.
基地確實是.
帶著三百人.
勇敢地去勝利了.
打贏了十幾萬人的米田人.
他也很懂得運用戰略.
運用戰術.
哪個時候做哪些事.
自己在這裡去不到.
就請人叫以法連人.
在山上最下來最近.
馬上截住條河.
見到他有勇有謀.
不過完了被人罵.
又更加感恩.
見到他這麼謙卑.
他就說出你們做得真好.
比我們厲害很多.

$^{401}$這樣的謙卑與情傷.
其實我們需要學習的.
多謝我們的主.
在今天也需要學.
給我們有謙卑的心.
給我們柔和地回應別人的質問.
求主幫助我們.
讚美你.
你實在是大能的上帝.
願你得著榮耀.
祈禱奉耶穌.
基督名求.
阿們.
\newpage



\section{以弗所書 5:21-6:9}
\label{sec:032322}
\textbf{飛翔需要兩隻翅膀}~~~~禤凱榮牧師
\newline
\newline
連結: \href{http://www.church.com.hk/acms/content.asp?site=cdc\&op=show\&type=product\&code=032322\&layout=sermon}{\texttt{032322}} ~~~~ 語音日期: 2026-06-21
\newline
\newline
\hyperref[sec:032321]{\small{< < < PREV SERMON < < <}}
~
\hyperref[sec:index_chronic]{\small{[返順時目]}}
~
\hyperref[sec:032323]{\small{> > > NEXT SERMON > > >}}
\newline
\newline
以弗所書 5:21-6:9
\newline
\begin{longtable}{cl}
\hline
\hline
章節 & 經文 (和合本修訂版)\\
\hline
5:21 & \begin{tabularx}{0.7\textwidth}{X} 要存敬畏基督的心彼此順服。 \end{tabularx} \\ \\ \relax
5:22 & \begin{tabularx}{0.7\textwidth}{X} 作妻子的，你們要順服自己的丈夫，如同順服主。 \end{tabularx} \\ \\ \relax
5:23 & \begin{tabularx}{0.7\textwidth}{X} 因為丈夫是妻子的頭，如同基督是教會的頭；他又是這身體的救主。 \end{tabularx} \\ \\ \relax
5:24 & \begin{tabularx}{0.7\textwidth}{X} 教會怎樣順服基督，妻子也要怎樣凡事順服丈夫。 \end{tabularx} \\ \\ \relax
5:25 & \begin{tabularx}{0.7\textwidth}{X} 作丈夫的，你們要愛自己的妻子，正如基督愛教會，為教會捨己， \end{tabularx} \\ \\ \relax
5:26 & \begin{tabularx}{0.7\textwidth}{X} 以水藉著道把教會洗淨，使她成為聖潔， \end{tabularx} \\ \\ \relax
5:27 & \begin{tabularx}{0.7\textwidth}{X} 好獻給自己，作榮耀的教會，毫無玷污、皺紋等類的缺陷，而是聖潔沒有瑕疵的。 \end{tabularx} \\ \\ \relax
5:28 & \begin{tabularx}{0.7\textwidth}{X} 丈夫也應當照樣愛妻子，如同愛自己的身體；愛妻子就是愛自己了。 \end{tabularx} \\ \\ \relax
5:29 & \begin{tabularx}{0.7\textwidth}{X} 從來沒有人恨惡自己的身體，總是保養愛惜，正像基督待教會一樣， \end{tabularx} \\ \\ \relax
5:30 & \begin{tabularx}{0.7\textwidth}{X} 因我們是他身體的肢體。 \end{tabularx} \\ \\ \relax
5:31 & \begin{tabularx}{0.7\textwidth}{X} 「為這個緣故，人要離開父母，與妻子結合，二人成為一體。」 \end{tabularx} \\ \\ \relax
5:32 & \begin{tabularx}{0.7\textwidth}{X} 這是極大的奧祕，而我是指基督和教會說的。 \end{tabularx} \\ \\ \relax
5:33 & \begin{tabularx}{0.7\textwidth}{X} 然而，你們每個人都要愛妻子，如同愛自己一樣；妻子也要敬重她的丈夫。 \end{tabularx} \\ \\ \relax
6:1 & \begin{tabularx}{0.7\textwidth}{X} 作兒女的，你們要在主裡聽從父母，這是理所當然的。 \end{tabularx} \\ \\ \relax
6:2 & \begin{tabularx}{0.7\textwidth}{X} 當孝敬父母，使你得福，在世長壽。這是第一條帶應許的誡命。 \end{tabularx} \\ \\ \relax
6:3 & \begin{tabularx}{0.7\textwidth}{X} 見上節 \end{tabularx} \\ \\ \relax
6:4 & \begin{tabularx}{0.7\textwidth}{X} 作父親的，你們不要激怒兒女，但要照著主的教導和勸戒養育他們。 \end{tabularx} \\ \\ \relax
6:5 & \begin{tabularx}{0.7\textwidth}{X} 作僕人的，你們要懼怕戰兢，用誠實的心聽從你們肉身的主人，好像聽從基督一般； \end{tabularx} \\ \\ \relax
6:6 & \begin{tabularx}{0.7\textwidth}{X} 不要只在人的眼前這樣做，像僅是討人的喜歡，而是作基督的僕人，從心裡遵行神的旨意， \end{tabularx} \\ \\ \relax
6:7 & \begin{tabularx}{0.7\textwidth}{X} 甘心服侍，好像服侍主，不像服侍人， \end{tabularx} \\ \\ \relax
6:8 & \begin{tabularx}{0.7\textwidth}{X} 因為知道每個人所做的善事，不論是為奴的或是自主的，都必按所做的從主得到賞賜。 \end{tabularx} \\ \\ \relax
6:9 & \begin{tabularx}{0.7\textwidth}{X} 作主人的，你們待僕人也是一樣，不要威嚇他們，因為知道他們和你們在天上同有一位主，他並不偏待人。 \end{tabularx} \\ \\
[1ex]
\hline
\hline
\end{longtable}
$^{1}$大家聽完剛才所讀的經文之後.
弟兄們會不會覺得有期盼和雀躍的感覺呢?.
而太太們會不會覺得.
今天在家裡看崇拜比較好.
為什麼呢?.
因為剛才我們所讀的第24節.
那裡說教會怎樣信奉基督.
妻子也要怎樣信奉丈夫.
那麼男士們可能馬上說.
「穆斯你講多點這類題目吧」.
女士們可能會說.
「穆斯你看著說啊」.
「不要說錯話」.
會有很大壓力.
我其實真的聽過一些.
真實的弟兄的分享.
他說有一次.
可能也不止一次.
跟太太爭執有不同的意見.
這個丈夫就引用這一節的經文.
來要求他的太太信奉.
太太聽完之後說.
「信耶穌是這樣的嗎?」.
「如果是這樣的我寧願不信」.
事實上我最近聽到這對夫婦的近況.
覺得太太真的離開了教會.
不知道是不是因為這個原因.
很多人都相信剛才我們所讀的那段經文.
似乎是討論了一個很重要.
它是很爭拗性的問題.
什麼問題呢?.
妻子是否完全信奉她的丈夫呢?.
有不少的丈夫.
用這段經文成為一個理據.
要求他們的太太信奉.
而有些太太為了要做一個虔誠.
好的基督徒.
她就被迫.
就算是不自願.
都要被迫信奉他們的丈夫.

$^{41}$究竟這個看法.
是不是真實聖經的教導呢?.
讓我們一同通過這段經文.
去嘗試探討.
理解究竟神建立地上不同類型的關係.
祂的心意是怎麼樣的呢?.
我們都跟大家說過.
要理解某一段經文.
我們必須要按照上文下理來理解.
我們不可以抽離斷章取義.
這段經文最重要的.
我們需要理解的就是第十八節.
那是整段經文的理解基礎.
第十八節說什麼呢?.
祂說我們不要醉酒.
酒能使人放蕩.
乃要被聖靈充滿.
而當一個人被聖靈充滿之後.
我們就應該有個表現.
接著保羅就說.
就會用詩章頌詞靈歌.
不僅如此.
我們也能夠彼此信服.
一個被聖靈充滿的基督徒.
是應該有一個能夠信服的心.
但是大家要留意.
不代表基督徒就像門前的地毯.
人人走過都可以踩兩下.
搓兩下.
基督徒不要因為這段經文就覺得我們是一個弱者.
這段經文告訴我們.
我們反而是因為主耶穌的緣故.
我們自願的犧牲一些我們的權利.
來讓對方得到好處.
接著保羅在接下來幾節的經文.
就分享用三個不同的關係.
來分享我們作為神兒女.
應該怎麼活出這些地上的關係.
這三個關係雖然是不同的.
但是他們共同有一個相同的基礎原則.

$^{81}$原則記載在第21節.
全敬畏基督的心.
彼此信服.
告訴我們這不是一個單方面的命令.
而是建基於我們對耶穌基督的敬畏和尊重.
彼此謙卑服侍的態度.
在今天這個世代裡面.
這個世界是講權力的.
我們經常說.
誰大誰惡誰正確.
所以我們說著說著.
想贏就越來越大聲.
我們從小到大的教育就是.
我們一定要懂得維護我們自己的權益.
千萬不要被人欺負.
要比其他人厲害.
這是我們普遍在世的教育.
我想大家想像一下一個遊戲.
大家都玩過了.
扯大欄,有玩過扯大欄嗎?.
扯大欄是怎麼玩的?.
出盡你的力.
將對方拉過來.
拉到對方.
這就是遊戲的勝負規則.
這個世界其實看關係.
很多時候都是這樣.
誰有力一點.
那邊就贏了.
而這個關係的建立方式.
不只是人與人之間.
國與國之間都是這樣.
有些國家有些政治家就說.
我以實力出發.
就是說我有實力的時候.
我就怎麼做都行.
這個世界是這樣的似乎.
不過教會卻不一樣.
教會不是這樣.
教會玩的遊戲不是扯大欄的遊戲.

$^{121}$教會玩的遊戲是.
一起成為一個對攻去做一些事情.
例如托起一條大木頭.
這個就是教會的遊戲.
彼此信服.
並不是雙方面去爭勁去拿權力.
誰大力就誰贏.
不是這樣.
我們玩的遊戲.
是在耶穌基督的主權底下.
共同努力去達到一些目的.
所以我們所追求的.
不是自身的權益.
我有些什麼.
我可以得到些什麼.
我們所說的不是這件事.
教會我們所追求的.
就是可以去考慮到.
其他人的需要和好處.
為的是神的榮耀.
換言之.
當我們說要信服給其他人的時候.
我們首先要明白一件事.
我們首先有一個基本態度.
就是我們看神的榮耀.
重要過我們自身的利益.
敬畏神可以說是彼此信服的先決條件.
所以大家一定要不要搞錯.
我們願意和能夠去信服.
不是因為對方配得我們信服.
不是這個原因.
我們願意和能夠信服是因為我們對神的敬畏.
我們盼望信服神.
我們盼望事奉神.
為什麼?.
耶穌基督為我們成就救恩.
祂配得我們的事奉.
好了.
保羅用的第一個例子來講解彼此信服.
就是丈夫與妻子彼此信服.

$^{161}$保羅一口一句就說.
「你們作妻子的,當信服自己的丈夫」.
我們理解這節經文.
我們首先第一件事要搞清楚.
保羅這裡並不是說妻子是低等一點.
丈夫是高等一點.
威一點,叻一點.
這一句話保羅並不是這樣的意思.
因為保羅在另外一個地方.
加拿大書講得很清楚.
那裡講什麼?.
那裡講說不分你是什麼人.
只要我們在耶穌基督裡面.
我們就披戴基督.
然後怎樣?.
我們就與在基督裡面都成為一了.
就是說每一個罪人來到耶穌基督面前的時候.
我們都是罪人.
我們都是蒙恩.
通過我們的信心而得著這個救恩.
所以在救恩面前人人平等.
所以各位弟兄聽得很清楚.
保羅沒有在這裡給了弟兄們一個做暴君的許可證.
沒有.
不要搞錯.
所以當我們與妻子有爭執的時候.
不要拿這一節的經文來強迫我們的妻子信服.
好了.
我們看一看第22節.
所以說一說.
「你們作妻子當信服自己的丈夫」.
這裡的信服.
大家要很留意.
是自願的.
不是強制的.
就好像基督徒信服耶穌基督一樣.
大家覺得耶穌基督是我們的主.
祂要求我們作為信徒去信服祂.
你覺得這個信服是自願還是強迫的?.
我相信大家都很清楚.

$^{201}$耶穌基督從來沒有強迫我們去信服祂.
因為如果這是神的做法的時候.
事情可以很簡單.
神做我們全部都像機械人一樣.
大家今天都見過機械人表演.
按個鍵.
一起跳舞.
一起打功夫.
一起打關斗.
不會有其他的反對.
不會有其他的不同的動作.
這個是一致性.
是絕對的信服.
但是耶穌基督從來沒有這樣對我們.
大家記不記得我剛才所舉的那個例子.
那個例子就是說我們基督徒不是在玩扯大欄的遊戲.
基督徒的關係的建立就好像兩個人一起做一些事情.
大家同坐一條船.
大家是同工來做同一件事情.
所以對基督徒來說.
信服並不是說某一個人好一點.
或者某一個人重要一點.
而是兩個人雙方一同的合力努力.
扮演著不同的角色.
但是要達到同一個的目標.
所以可以這麼說.
丈夫與妻子是一個對工.
一同經營同一個的家庭.
彼此信服在婚姻裡面.
就是要信服對方的需要.
我想問大家.
你覺得丈夫與妻子的需要是什麼?.
最重要的需要是什麼?.
我相信對太太來說.
最需要的是愛.
對丈夫來說.
最需要的是尊重.
對太太來說.
對妻子來說.
信服的意思是指將自己放在丈夫的保護和關懷之下.

$^{241}$尊重他的領導.
和他一起努力建立經營同樣屬於他們的家.
這就是信服對他們的意思.
這絕對是一種出於對基督的敬畏.
而不是出於丈夫給他的恐嚇.
或者覺得是一種義務.
對丈夫來說.
信服意味著愛他的妻子.
有什麼例子呢?.
正如基督愛教會.
耶穌基督.
大家留意.
從來沒有透過強權.
控制或者恐嚇.
來達到愛.
耶穌基督給我們的榜樣是什麼?.
耶穌基督愛的實現在哪裡?.
他對人的服侍,犧牲,保護,給予.
甚至於死.
所以丈夫對妻子的愛.
如果我們要仿效教會對基督的愛的時候.
我們需要明白的就是.
如同愛自己的身體.
滋養他.
珍惜他.
將他的需要.
置於自己的需要之上.
這就是丈夫對妻子信服的意思.
當兩方面.
丈夫與妻子都做好他們在婚姻裡面的角色的時候.
當太太將她的生命.
投放在丈夫的愛與關懷裡面的時候.
妻子會感到安心.
願意將他繼續放在丈夫的帶領和愛裡面.
而當丈夫感受到.
從太太而來的尊重的時候.
他就會再一步,更進一步地愛他的妻子.
珍惜他的妻子.
這是一個良性不斷向上的循環.
第二個的關係.

$^{281}$本來就是父母與兒女的關係.
在今天的世界裡面.
父母很多時候都不理智地.
溺愛他們的兒女.
所以很多時候兒女會變成在家裡做小霸王.
現在也有另一個極端的例子.
今天我們聽到很多的父母.
要絕對地控制他們的兒女.
他們的兒女沒有任何的話事權.
什麼都是我說了算.
父母都是另一個極端.
而他們給的原因很簡單.
是我為了他們好.
所以父母們留心聽清楚.
神從來沒有要我們養育出一些絕對聽從父母的兒女.
聖經沒有這麼說.
聖經反而是要父母教導兒女,培育兒女.
成為一個愛耶穌,信服上帝的人.
故此保羅在這裡的教導就是說我們需要按著聖經的教導.
來養育我們的兒女.
不要惹兒女的氣.
很簡單的.
因為你有管教.
但沒有愛的時候.
你只會製造兒女們當中的苦毒.
如果你很多的規矩但卻沒有關係作為基礎的時候.
你只會引發兒女們的背叛.
那對於兒女呢.
對於今天的兒女來講我覺得有兩個極端.
第一個極端就是今天我們所說的.
啃老一族.
渣渣渣.
渣到他們的父母一滴都沒有了.
另外一個極端就是說.
完全放棄他們.
拋棄他們,離開他們.
他們自己有他們的世界.
我有我的生活.
似乎在兩個極端裡面.
保羅告訴我們.

$^{321}$我們作為兒女需要孝敬父母.
我問大家什麼是孝敬.
每個人的孝敬可能都不同.
聖經裡面所說的孝敬.
意味著我們在言語和態度上.
我們一定要有對父母的尊重.
我們一定要給他做父母應有的尊嚴.
好了,時代不斷地變.
可能父母對某些東西落後了.
給予意見的時候.
你懂什麼.
好了,父母動作慢一點.
你不要阻礙地球轉.
這些說話千萬不能講.
保羅甚至用.
一個十誡裡面的第五誡來強調.
神如何看重父母兒女之間的關係.
在《創埃及記》第20章第12節說到.
我們應當孝敬父母.
因為使你的日子.
在耶和華你神所賜你的地上.
得以長久.
這個很明顯說到.
神是會賜福給一些進行心意的人.
好了,接著保羅再說.
第三個的關係.
主人,僕人的關係.
當然今天我們不會有盧博.
所以我們大致可以用這個關係.
來思考我們今天在地上.
僱主與僱員的關係.
我只是打份工.
這個關係和其他兩個關係有點不同.
剛才所說的兩個關係.
都有很大愛的因素在裡面.
但是僱主與僱員.
除非發生一些.
另外的關係.
否則僱主與僱員.
是利益的關係.

$^{361}$利益的關係.
僱主希望僱員.
你做多些.
幫我賺多些錢.
僱員就說.
你給多些工我做.
讓我可以賺多些錢.
由於這個關係的基礎.
建立在利益那裡.
所以很自然.
僱主當然會想盡辦法.
讓僱員.
多為他們工作.
多為他們賺錢.
僱員就要想辦法.
怎樣可以升職加薪.
搖搖直上.
不過保羅在這裡.
告訴我們.
我們不是單是想.
自己的益處.
保羅在這裡說.
我們需要用誠實的心.
聽從你們肉身的主人.
監心事奉.
好像服侍人.
不像服侍人.
不知道大家有沒有在工作的地方.
安裝了這個.
那些叫什麼.
CCTV.
大家喜不喜歡.
當然不喜歡.
經常被人監視.
經常被人看著我工作.
不過無論你喜不喜歡.
事實就告訴我們.
很多時候在工作的場所裡.
如果安裝了這個東西.
的話.

$^{401}$那個效率就高很多.
很容易理解.
當你知道你的老闆.
天天看著你的時候.
當然不會過份偷懶.
大家是不是這樣.
有人看著.
然後才會.
落力工作.
沒有人看著的時候.
偷懶,人都要鬆一鬆.
是不是這樣.
保羅這個教導很有趣.
基督徒在我們.
事奉我們的主人也好.
僱主也好.
心態要如同.
服侍主.
主是24小時的.
主不需要CCTV.
主24小時都看著你.
有沒有偷懶,有沒有做一些不應該做的事.
有沒有努力地.
忠心地做一些應該做的事.
所以基督徒.
不需要CCTV.
因為我們知道.
我們事奉地上的.
工作.
或者僱主.
其實就是事奉我們的神.
如果心態.
擺得準確的話.
對於僱主.
來說.
所以大家要留意.
如果有人在幫你打工的話.
千萬不要.
運用你那個.
僱主的身份與地位.

$^{441}$來去威嚇.
壓榨,欺騙,苦待.
你的員工.
為什麼?因為保羅說.
每一個人.
無論你是僱主還是僱員.
在天上都有同一個主.
每一個人,有一天.
都要向這個主交仗.
三種的關係.
其實都告訴我們.
一樣同一個事情.
如果我們需要.
在這三種不同的關係裡.
活出神的榮耀的時候.
第一件事我們需要什麼?.
需要知道我們在這個關係裡.
我們要扮演的角色.
第二件事就是.
我們應該要怎麼樣.
去做好我們的角色.
以致這個.
我們所付出的.
能夠帶來.
雙方在關係裡.
雙方的好處.
同一個時間.
將榮耀歸給神.
我想大家想一想.
一隻鳥,鳥有兩隻翅膀.
一隻左,一隻右.
都有不同的角度.
有不同的.
一個左,一個右.
不過,兩個翅膀.
雖然有不同.
他們卻是同樣的重要性.
不需要我說,隨便一隻翅膀.
出了事,或者不想動.
不配合的時候.

$^{481}$鳥就飛不起來.
不但飛不起來,還會跌在地上.
跌得很厲害.
今天的經文講到這裡,我可以結束了.
不過我覺得.
最後這一點.
一定需要交代清楚.
否則的話.
當我們面對一些極端.
不過不是罕見.
並非罕見的情況.
出現的時候.
我們會形容錯了這段經文.
沒錯,一隻鳥需要兩個翅膀.
聖經的教導.
永遠都是對的.
當在一個關係的雙方.
都扮演著他們應該要扮演的角色.
都是按照心意來付出的時候.
這個關係.
必定是.
強壯與健康.
問題在於.
如果一方盡了自己的能力.
做好本分.
另外一方.
卻是將這一邊的努力.
看為他們.
拿著數的機會.
那怎麼辦?牧師.
不會有這樣的事情發生吧.
我做了牧師之後.
我發覺這個世界什麼人都有.
這個世界什麼關係都有.
首先我們需要明白.
彼此馴服.
並不是.
我們在一個遊戲裡面.
去爭管轄權.
所以做太太.

$^{521}$千萬不要坐在一起.
看看有什麼好的策略.
治好我的老公.
老公坐在一起的時候.
你有什麼方法可以控制你的太太.
讓我學一下.
當你有這樣的心態的時候.
你對這個婚姻.
完全搞錯了.
這個關係彼此馴服.
並不是去爭取管轄權.
而是一個夥伴的關係.
是一個同工.
一個夥伴.
在這個夥伴裡面.
我們需要明白.
你好我好大家好.
如果我好你不好的話.
不能有大家好.
這個就是同工的意思.
所以如果有一個人.
利用對方的謙卑努力.
成為一個自己.
提升自己的機會的時候.
問題不是在於.
願不願意馴服.
問題在於.
有一方.
根本沒有敬畏神.
而有一方.
根本對另外一方沒有真正的愛.
根據保羅的教導.
如果將它做一個剛才我們所講的總結的話.
妻子要馴服丈夫.
丈夫要愛妻子.
孩子要孝敬父母.
父親不要惹兒女的氣.
下屬要忠誠服侍.
上司要公平的對待下屬.
好了我問你.

$^{561}$這一大堆的教導.
是不是任何處境.
都要嚴格遵守的.
是不是.
剛才你所講.
另外一方沒有做什麼.
還在拿我著數.
是不是都要馴服.
我想大家留意一件事.
保羅從來沒有講過.
我們做我們關係的那一份.
是因為對方做了.
所以我們做.
沒有講過.
剛才我一開始的時候.
如果大家有留意聽清楚的話.
我們願意彼此馴服.
並不是因為.
對方配得我們的付出.
而因為什麼.
而是因為.
神的命令.
神希望我們這樣做.
神要我們這樣做.
所以.
保羅在第八節.
這樣講的.
我們每一個人.
都是.
按我們.
個人的情況.
行善事.
而有一天.
神會按著我們所行的.
得主的.
賞賜.
每一個人到某一天的時候.
必須向耶穌基督.
交賬.
交我們今生所做過一切的賬.

$^{601}$然後按我們的所行.
來去接受.
神所為我們預備的獎賞.
所以大家.
要留意一件事.
就是對方沒有做好.
他在關係裡面的本份.
讓我們不遵守.
神的心意.
以致我們失去了.
神為我們留下的獎賞.
我們要看耶穌基督的榜樣.
耶穌基督的榜樣是怎麼樣的.
耶穌基督領受.
天父上帝的使命.
他的心意來到這個世界.
為人施行救恩.
不過在整個地上的.
事公裡面.
世人是怎樣的.
世人是輕視他的愛.
宗教領袖是控告他.
群眾嘲笑他.
士兵怎麼樣.
士兵鞭打他.
猶大怎麼樣.
猶大出賣他.
不過.
耶穌基督從不.
從來沒有放棄.
仍然的.
堅守.
神給他的使命.
不過同一個時間.
大家也要留意.
耶穌基督也沒有完全讓人.
控制他的.
有的時候.
面對很憤怒.
很充滿敵意的群眾.

$^{641}$他選擇離開.
對於虛偽的善行.
他選擇.
嚴厲的抨擊.
對於罪惡.
毫不妥協.
他責罵罪惡.
他都畫定一些界線.
從耶穌基督的.
例子裡面我們看到.
信服千萬不要.
以為.
是門前的地毯.
我們不是doormat.
是人走過都可以踩兩腳.
我們不需要逆來順受.
那你說究竟怎樣.
如果對方濫用我的.
信服我應該怎樣做.
在這樣的處境裡面.
都要考慮.
在不傷害自己的情況下.
保持信服.
剛才我所說的.
有些很極端的例子.
你越信服的時候.
對方就會越濫用你的信服.
在這樣的.
變質的關係裡面.
反而信服不是帶來.
雙方的好處.
而是帶來雙方的傷害.
這個傷害.
不單只是身體.
也可以是情緒,精神,財務上的傷害.
在這樣的處境的時候.
我們需要.
我們不只是.
死忍.
我們需要嘗試.

$^{681}$去跟對方溝通.
因為如果我們死忍.
我們只是想忍耐一下.
我們忍耐一下.
有一天他會改變.
他會改變的.
這種單方面的盼望.
在現實上.
對不起了,很難修成正果.
我們需要考慮的就是.
雙方坐下來.
認真的,真誠的溝通.
把問題的所在交出來.
好了.
第三.
我們需要按著聖經的原則.
來畫一些界線.
信服不等於什麼都做.
對方要我做這個,我就做那個.
在聖經的原則下.
就是主女信服.
什麼意思呢?.
有些東西不在我們信服的清單裡面.
例如什麼?.
不道德的東西.
罪惡的東西.
犯法的東西.
這些不能信服.
如果我做完所有的事情.
我們之間的關係仍然是這樣.
那怎麼辦呢?.
我們試試找第三者.
聖經教導我們.
當我們在生活上遇到不同類型的困難的時候.
教會是我們一個資源.
一個找幫助的地方.
教會裡面.
應該.
應該有一些智者.
應該有一些熟齡的人.

$^{721}$可以幫我們去處理一些.
我們處理不到的問題.
弟兄姊妹.
所有.
健康的關係.
都需要雙方的.
付出和努力.
一隻翅膀.
是飛不了的.
雖然我們很多時候.
在不同的關係裡面.
我們都會覺得.
這一方是有話事權比我多.
也許是這樣.
不過.
話事權是那個角色裡面.
要有的東西.
不等於.
有話事權或多一點話事權的那個.
就是高一級.
低一級.
我希望大家離開這個崇拜之後.
一定要澄清一個觀念.
觀念就是.
這段經文從來沒有說過.
在關係裡面的一方.
是高級過.
上等過另外一方.
付出.
這一個關係.
我們有時候.
可能在關係裡面.
有一方話事權比較多.
但是千萬要注意.
你有話事權比較多的時候.
不要用這個東西.
去得到自己的好處.
而傷害到另外一方.
反而是將你.
在關係裡面所有的東西.

$^{761}$用出來.
讓這個關係更加的強壯.
更加的健康.
為的就是讓神得著榮耀.
不容易的.
因為所有涉及人的東西.
我做了牧師這麼多年.
所有涉及人的東西都不容易的.
但我告訴大家.
當我們願意攜手一起的時候.
那個回報相當的豐厚.
讓你在人生的路上.
得著很大的鼓舞.
那個關鍵.
正如保羅所說.
我們必須要.
被聖靈充滿.
這是第一點.
什麼意思被聖靈充滿.
按聖靈行事.
聖靈會告訴你什麼是對.
什麼是不對.
我們需要對基督敬畏的心.
這樣.
有了之前的東西的時候.
我們才能夠有機會.
健康地彼此馴服.
享受神賜給我們地上.
所有關係的好處.
能夠活出.
這個人生的精彩.
能夠將這個精彩化為對神的榮耀.
我們都能夠做到.
願意我們能夠彼此鼓勵.
請我們低頭祈禱.
天父上帝我們承認.
我們很多時候看權力多於信服.
很多時候我們看權利.
而不是看服事.
神也求你饒恕我們.

$^{801}$神也求你.
讓聖靈充滿我們.
教導我們.
引領我們.
使我們能夠效法耶穌基督.
學習祂的虛己.
得著敬畏基督的心.
以致我們能夠彼此馴服.
享受你所賜給我們地上.
每一種的關係.
求你恩待.
祈禱奉主名求.
\newpage



\section{俄巴底亞書 1:1-21}
\label{sec:032323}
\textbf{不要沾沾自喜?}~~~~黃紹權牧師
\newline
\newline
連結: \href{http://www.church.com.hk/acms/content.asp?site=cdc\&op=show\&type=product\&code=032323\&layout=sermon}{\texttt{032323}} ~~~~ 語音日期: 2026-06-21
\newline
\newline
\hyperref[sec:032322]{\small{< < < PREV SERMON < < <}}
~
\hyperref[sec:index_chronic]{\small{[返順時目]}}
~
\hyperref[sec:032324]{\small{> > > NEXT SERMON > > >}}
\newline
\newline
俄巴底亞書 1:1-21
\newline
\begin{longtable}{cl}
\hline
\hline
章節 & 經文 (和合本修訂版)\\
\hline
1:1 & \begin{tabularx}{0.7\textwidth}{X} 俄巴底亞所見的異象。論以東說：我從耶和華那裡聽見信息，並有使者被差往列國去，說：起來吧，一同起來與以東爭戰！ \end{tabularx} \\ \\ \relax
1:2 & \begin{tabularx}{0.7\textwidth}{X} 看哪，我要使你在列國中為最小，被人大大藐視。 \end{tabularx} \\ \\ \relax
1:3 & \begin{tabularx}{0.7\textwidth}{X} 你狂傲的心欺騙了你，你住在巖穴，居所在高處，心裡說：「誰能把我拉下來到地上呢？」 \end{tabularx} \\ \\ \relax
1:4 & \begin{tabularx}{0.7\textwidth}{X} 你雖如鷹高飛，在星宿之間搭窩，我必從那裡拉你下來。這是耶和華說的。 \end{tabularx} \\ \\ \relax
1:5 & \begin{tabularx}{0.7\textwidth}{X} 盜賊若來到你那裡，小偷夜間來到，豈不是只偷他們所需要的嗎？摘葡萄的若來到你那裡，豈不留下幾串嗎？你竟全然滅絕！ \end{tabularx} \\ \\ \relax
1:6 & \begin{tabularx}{0.7\textwidth}{X} 以掃遭到搜查，他隱藏的寶物竟被尋出！ \end{tabularx} \\ \\ \relax
1:7 & \begin{tabularx}{0.7\textwidth}{X} 與你結盟的都驅趕你，直到邊界，與你和好的欺騙你，勝過你，吃你飯的人設下圈套陷害你---他卻毫無聰明。 \end{tabularx} \\ \\ \relax
1:8 & \begin{tabularx}{0.7\textwidth}{X} 到那日，我豈不從以東除滅智慧人？從以掃山除滅聰明人？這是耶和華說的。 \end{tabularx} \\ \\ \relax
1:9 & \begin{tabularx}{0.7\textwidth}{X} 提幔哪，你的勇士必驚惶，以致以掃山的人都被殺戮剪除。 \end{tabularx} \\ \\ \relax
1:10 & \begin{tabularx}{0.7\textwidth}{X} 因你向兄弟雅各施暴，你必蒙羞，永被剪除。 \end{tabularx} \\ \\ \relax
1:11 & \begin{tabularx}{0.7\textwidth}{X} 當陌生人擄掠雅各的財物，當外邦人進入他的城門，為耶路撒冷抽籤分取財物的日子，你竟站在一旁，像與他們同夥。 \end{tabularx} \\ \\ \relax
1:12 & \begin{tabularx}{0.7\textwidth}{X} 你兄弟遭難的日子，你不該瞪著眼看；猶大人被滅的日子，你不該幸災樂禍；他們遭難的日子，你不該說狂傲的話。 \end{tabularx} \\ \\ \relax
1:13 & \begin{tabularx}{0.7\textwidth}{X} 我子民遭災的日子，你不該進他們的城門；他們遭災的日子，你不該瞪著眼看他們受苦；他們遭災的日子，你不該伸手搶他們的財物。 \end{tabularx} \\ \\ \relax
1:14 & \begin{tabularx}{0.7\textwidth}{X} 他們遭難的日子，你不該站在岔路口剪除他們逃脫的人，你不該交出他們的倖存者。 \end{tabularx} \\ \\ \relax
1:15 & \begin{tabularx}{0.7\textwidth}{X} 耶和華的日子臨近萬國；你所做的，人也必向你照樣做，你的報應必歸到自己頭上。 \end{tabularx} \\ \\ \relax
1:16 & \begin{tabularx}{0.7\textwidth}{X} 你們在我聖山怎樣喝了苦杯，萬國必照樣不停地喝，且喝且吞，他們就必歸於無有。 \end{tabularx} \\ \\ \relax
1:17 & \begin{tabularx}{0.7\textwidth}{X} 但在錫安山必有逃脫的人，那山必成為聖；雅各家必得原有的產業。 \end{tabularx} \\ \\ \relax
1:18 & \begin{tabularx}{0.7\textwidth}{X} 雅各家必成為大火，約瑟家成為火焰；以掃家必如碎秸，遭燃燒，被吞滅，以掃家必無倖存者。這是耶和華說的。 \end{tabularx} \\ \\ \relax
1:19 & \begin{tabularx}{0.7\textwidth}{X} 他們必得尼革夫和以掃山，得謝非拉，非利士人之地，他們必得以法蓮地和撒瑪利亞地，得便雅憫和基列； \end{tabularx} \\ \\ \relax
1:20 & \begin{tabularx}{0.7\textwidth}{X} 被擄的以色列大軍必得迦南人的地，直到撒勒法，在西法拉被擄的耶路撒冷人必得尼革夫的城鎮。 \end{tabularx} \\ \\ \relax
1:21 & \begin{tabularx}{0.7\textwidth}{X} 必有一些解救者上到錫安山，審判以掃山，國度就歸耶和華了。 \end{tabularx} \\ \\
[1ex]
\hline
\hline
\end{longtable}
$^{1}$大家好,經過退休會後,我們又回到我們的教會聚會點,一起去敬拜.
亦是一個很好的轉折,我們開始新一卷書卷《俄巴底亞書》.
如果你對俄巴底亞書很陌生,絕對不出奇.
因為俄巴底亞書,你如果打開聖經是很難找的.
因為它只有二十一節,但是這卷書卻是藏在十二卷.
或者用基督教的分類叫做小先知書的其中一部分.
我們能夠在這卷書裡面有什麼學習.
能夠透過它能夠透視到在這個被擄之間的時間.
以色列民,尤其是猶大國的子民,怎樣在期間去面對他們的種種信仰的挑戰.
以及上帝的恩典怎樣打開去到普世.
都能夠給我們一個很不一樣的啟示.
讓我們一起再次在聆聽主的說話一起研讀以先.
加上禱告,祈求主給我們有智慧能夠在當中明白俄巴底亞書的啟示.
父神多謝你讓我們再次回來.
讓我們能夠在真理上不斷去學習並且去尋求.
多謝你讓我們能夠有這份生命能夠再繼續走在這條屬靈的道路上.
蒙受聖靈的帶領和去造就我們為此獻上感恩.
又是能夠在新的一卷書卷裡打開有更加透徹明白上主你的心思意念.
為此在這個時候加上禱告.
祈求你的帶領讓每一個無論是聽的或者說的.
都能夠在這一卷俄巴底亞書的當中.
能夠更加將我們的屬靈生命深化.
而不是留於浮面膚淺的狀態.
求你也幫助我們每一個弟兄姊妹.
能夠在這個學習的過程當中.
我們能夠彼此守望.
透過大家在認識這本書的過程當中的領受.
透過分享.
透過鼓勵和提醒.
我們能夠更加實質地去將主你的教訓.
融合進到我們生命每一個的環節.
求你的帶領.
禱告奉靠主耶穌基督得勝名字而求.
阿們.
如果當我們去看回在新約書卷的約翰福音裡面的時候.
大家會留意到去到約翰福音的最後.
彼得被耶穌三次問祂.
你是不是愛我.
如果當我們仔細去看回約翰福音那段經文的時候.
其實耶穌的問題不單只是問愛不愛耶穌這個問題.

$^{41}$一方面當然是指到彼得究竟是不是愛主耶穌基督.
但其實同時間也都展示出彼得對於身邊他一起同伴的門徒.
或者甚至這班門徒的所作所為.
是不是彼得都帶著愛心來去接受和面對呢.
這個約翰福音的記載其實是一個很有深度的介紹.
也都讓我們去看到在往後日子甚至去到仕途行傳的當中.
彼得怎樣去看待身邊的門徒.
甚至是一些我們稱之為仕途的保羅.
怎樣去面對彼得或者彼得怎樣面對他們.
但是在這段聖經裡面.
尤其是在第21章約翰福音第21章第20節至到23節裡面.
這段經文是非常之有趣.
但是也都非常之難解釋的.
如果大家有機會的或者和我一起去打開.
約翰福音的第21章第21節至到23節.
我讀出給大家聽聽.
這段聖經其實我們明不明白他當中想講些什麼.
這段聖經這樣講.
他說彼得轉過來看見耶穌所愛的那門徒跟著.
就是在晚飯的時候靠著耶穌胸膛說.
主啊賣你的是誰的那門徒.
彼得看見他就問耶穌說.
主啊這人將來如何.
耶穌對他說.
我若要他等到我來的時候與你何干.
你跟從我吧.
於是這話傳在弟兄中間.
說那門徒不死.
其實耶穌不是說他不死.
而是說我若要他等到我來的時候與你何干.
這段聖經看似很簡單.
甚至在基督教會的傳道之間.
有時甚至將彼得對身邊這位耶穌所愛的門徒.
有一種不客氣的態度.
似乎彼得有一種增進的心智.
究竟是否這樣呢?.
我在對於約翰福音的詮釋的時候.
來到這一段聖經.
當然有幾個比較難處理的問題.
第一這段聖經是否真的在原文聖經裡存在.

$^{81}$因為約翰福音的最後這一部分.
稱之為一個增長版本的結束.
究竟是否原本耶穌在聖經新約啟示裡的一部分呢?.
這其中在學術上其中一個掙扎.
不過我不在這裡和大家去看.
我們當這段聖經都是上帝的啟示的時候.
彼得是否出於與這位耶穌所愛的門徒.
或者我們所知道的這位約翰.
是有一個比較呢?.
還是彼得是出於因為這位是耶穌所愛的門徒.
所以彼得同樣關心到這位門徒.
他的未來發展是怎樣一回事呢?.
在這種解釋上其實當然你問我.
我有我個人的看法.
但是不少我們活在這個地上作為基督徒.
其實不多不少我們都會可能掉在這種矛盾裡面.
一方面我們似乎是需要按照耶穌基督的教訓.
我們要去彼此相愛.
去彼此接待.
我們很想做到的.
作為教會的一份子.
我們通常都會說弟兄姊妹.
要大家彼此照顧關懷守望.
但是也都有同時一種矛盾.
就是我們會有一種比較.
我曾經和一位弟兄分享到的時候.
他分享到一個他在教會裡面的一個難處.
他跟我講到他說在教會當中他回到一間教會已經接近十年的時間.
和一對夫婦非常熟悉.
兩家人一起經常去出遊.
甚至大家一起坐下來傾談.
在信仰上彼此信心上支援.
但是當一去到他們自己的子女的層面上.
他突然打了一個凸.
他說有一次當兩個小朋友參加一個比賽的時候.
那一次不幸也好榮譽也好.
他的小朋友得獎.
贏出這個比賽.
但是同樣參賽的這位很好的小朋友.
因為輸了的緣故.

$^{121}$就覺得很不開心.
而這對夫婦就突然間有一次和他這樣講.
他說你的女兒或者你的小朋友就厲害了.
他突然間覺得為何會有一個這樣的比較呢.
既然大家在主裡面都有這麼親密的關係.
為何我要引用約翰福音這段聖經.
就是因為彼得這段聖經裡面的表達.
可以有兩種很不同的詮釋.
他究竟是一種比較.
和耶穌這位所愛的門徒.
是不是一種在屬靈上都要爭耶穌的愛.
還是他真的體貼耶穌基督他的教導叫彼此守望.
尤其是彼得是耶穌門徒當中的領袖.
他是不是真的很關懷這位耶穌所愛的門徒的未來呢.
我不在這裡給一個答案給大家.
我想讓大家看到你和我.
都會可能跌在一種這樣的屬靈的矛盾裡面.
你說牧師你又說講俄巴底亞書.
為何變成講新約呢.
因為這種矛盾就正正是在.
俄巴底亞書裡面去陳述了出來.
在聖經當中我們會發現到.
原來俄巴底亞書裡面所提及到的.
以東以及猶大國.
原來是源自於舊約聖經裡面一對親生兄弟.
以素和雅各.
這一對兄弟從出生.
之前到出生的那一刻開始.
以至於往後他們兩兄弟.
以及他們所衍生出來的兩個邦國.
他們都不斷地沉醉在一種鬥爭或者比較裡面.
結果就發展成為一種互相攻擊.
甚至互相陷害的狀態.
當見到對方失敗的時候.
在俄巴底亞書我們見到以東人.
當看到猶大國被擄滅亡這件事情發生的時候.
他們充滿沾沾自喜.
甚至可以用熟稔一點的說法.
這位的以東民族他們覺得.
好了,上帝伸張正義為我們都出一口氣.

$^{161}$出些什麼氣呢?.
這個正正就是在整段我們去看俄巴底亞書的基礎觀念.
也都讓我們重新想一下.
在聖經的本意當中.
屬靈人在他們之間.
或者屬神的子民之間.
或者甚至親兄弟之間.
是不是真的有一份比較.
或者一份作為.
成為我們在地上要去掙扎的題目呢?.
所以讓我們一起去打開俄巴底亞書.
我們看看在這本書裡面.
耶和華神怎樣向以東人的一連串意象宣佈之後.
猶大國的國民是不是應該讀了俄巴底亞書.
他們就應該大開慶祝會呢?.
還是在這段聖經有更深一層的意義呢?.
讓我們一起去看看這本書.
所以為什麼今天我給了這段聖經的題目.
叫做不要沾沾自喜.
我們會在這段聖經裡面去看到.
有很多人都在見到對方失敗的時候.
似乎好像一步自己的恩仇.
或者可以說是見到對方被上帝懲罰.
就好像一口冤屈氣被紓緩一樣.
我們一起去打開這本書.
去看看俄巴底亞書.
我相信現在這個世界.
應該沒有人找不到俄巴底亞書的了.
如果你用手機電子版一戳就中了.
但是如果我們下次有機會玩個遊戲.
我們就不如拿本聖經出來.
看看你能不能夾到俄巴底亞書.
又看看你能不能夾到猶大書.
保證這個遊戲蠻好玩的.
回去有空就練習一下.
我們一起打開俄巴底亞書.
我們今天會和大家一起去看的.
就是第一節的經文.
沒有哪一章的了.
只有一章而已.

$^{201}$所以不需要哪一章的了.
我們只是看第一節.
有弟兄姊妹我們都打開聖經的.
我們一起讀一次好不好.
找到之後我們一起讀.
這個聖經想講什麼呢.
一般就這樣看下去.
就很簡單.
開頭講到俄巴底亞.
一位先知他得了耶華的啟示.
然後就講一大堆有關於以東的事情.
究竟是講些什麼呢.
很簡單而已就是這麼多.
上帝有話要講.
當我們去了解這段聖經的時候.
第一個我們要考慮的問題就是.
上帝為什麼在一本希伯來人聖經裡面.
要講一些關於另一個族群.
叫做以東的事情呢.
為什麼呢.
講一些背景.
當我們去讀俄巴底亞書.
第一個問題就是.
究竟俄巴底亞書寫作的歷史背景是什麼時間.
在研究聖經裡面.
學者很厲害的.
很多時候永遠給你兩個答案.
又不告訴你實質答案是什麼.
每個人學院都是這樣的.
很多時候我們的同學.
我們叫做功力未深的時候.
都會有這個問題的.
就是說為什麼教授講的東西無伶兩可.
宗教不是很實在的東西嗎.
尤其是我們將來出來募會.
或者有弟兄姊妹問.
俄巴底亞書是什麼時候寫的.
學者通常給兩個答案.
一個叫做早期寫作.
一個叫做末期寫作.

$^{241}$兩者之間的分別.
其實就是告訴我們.
這本書究竟他想成書的時間是什麼時候.
還是他寫的時候是哪一段時間.
其中早期理解的.
或者很多時候講到早期理解.
就是按照俄巴底亞書的內文.
來嘗試去斷定他的歷史背景.
很大部分的學者.
都認為俄巴底亞書所寫的.
正正就是新巴比倫的時代興起.
而進攻猶大國的時間.
如果真的用陽曆來講.
大約是主前的第六百年至六百零二年.
這段時間.
因為在這段時間裡.
我們知道猶大國是被新巴比倫國.
從北往南進攻.
所以這段聖經的直根點是.
被擄的過程.
尤其是被擄剛剛完成的時間.
相信是主前第六百零二年.
這是很多學者的解釋.
另一些學者提供後期成書的時間表.
後期成書的時間表.
往往是講到這本書寫了之後.
寄了下來.
然後真正成為一個書卷.
傳遞給人去讀的時間.
因此很多時候就講到.
這不是在被擄期間.
或者剛剛被擄完結的時間.
是在被擄已經過了好一段時日.
甚至猶大國不但不再存在.
是在羅馬世界的時間.
而寫成的一本書.
去重提這段歷史的因素.
成為在後世歷代的.
讀這本書的讀者的一個境界.
究竟是哪一個的詮釋比較好呢?.

$^{281}$或者才是對的呢?.
很多時候信徒就要求一個答案.
我就將這兩個答案告訴你.
因為我相信這兩個的作用.
是同時在一本書裡面發生的.
一方面要告訴我們.
在當寫俄巴底亞書的時候.
作為一位先知.
他將這一大堆的內容寫下來.
其實他的心路歷程.
以及當時猶大國被滅的時候.
他們的掙扎.
以及這段聖經裡面所講的主人翁.
或者叫做以東人.
啟示的對象.
他們究竟他們的心情.
靈性以及他們的下場是怎樣.
這兩個不同的解釋.
講到後期的話.
就會好像是一種.
可以值得被人滿心歡喜.
因為終於有人行惡被懲治.
而且是被上帝懲治.
在這樣的情況之下.
那不就是彰顯上帝公義了.
究竟是不是這個意思呢.
我們用兩方面來看看.
第一部分.
第一節的上半節.
上半節講什麼呢.
俄巴底亞得了耶和華的默示.
在一開始.
這裡要讓我們去理解.
這本書.
是要從上帝的意象的角度來去詮釋.
換句話來說.
我們不是將這本書套入.
我覺得它看完之後.
我覺得它是什麼意義.
反而我們要去問的是.

$^{321}$上帝想講過什麼訊息呢.
這個詮釋.
就是涉及於我們很多時候解釋聖經的兩個原則.
一個是對的,一個是錯的.
對的原則叫Exegesis.
我們是要從上帝的內容.
去找回上帝講解的動機和目的.
所以我們是從聖經去告訴我們.
從聖經向外來告訴我們.
另一個錯誤的詮釋叫Exegesis.
Exegesis是從我的角度套入聖經.
覺得勝經怎樣解釋.
為什麼我要在這裡強調一點呢.
首先可以幫助大家讀聖經的時候.
尤其是.
我們要做一些聖經分享的時候.
我們很容易跌入Exegesis的局面.
我覺得一段聖經怎樣適合我的情況.
按照我的理解和我的歷史背景.
我人世的遭遇來詮釋上帝應該是這樣的意思.
但是上帝往往他的解釋不是按照你和我的狀態來陳述的.
所以一開始俄巴底亞書要講得很清楚.
這個必須要從神的角度來詮釋.
神以下一連串的異象陳述.
為什麼說異象呢.
這裡明明中文寫明是默示的.
如果大家有聖經其他的譯本.
尤其是英文.
這裡用到一個英文的翻譯.
希伯來文叫做Vision.
我可以說著緊90\%的翻譯譯本.
都是用Vision這個字眼.
Vision和默示有一些意義上的深度不同.
默示就好像某一件事.
我靜靜地跟你說了.
上帝是不是靜靜地躲起來和俄巴底亞這位先知說呢.
似乎不只是這樣.
因為Vision或者在希伯來文上這個字眼.
異象是一幅活活的圖畫.
而先知俄巴底亞是經歷過親歷其境.

$^{361}$或者甚至我可以這樣說.
是親眼目睹一幅圖像.
來了解上帝的啟示.
所以這個不是自己自發地去一種詮釋.
而是上帝有展示給人去看的一本書.
所以不是只是俄巴底亞先知的個人經歷.
然後將他自說自話寫了下來.
反而是一連串的讓上帝帶他所看的景象.
來將他記下來.
所以不是一種只是吸眼睛的一種觀望之後.
寫下來的一種個人的一種經歷靈性幻想.
而是上帝帶他去行看過這幅圖畫.
我相信每一個信徒都很期望有一個這樣的經歷.
弟兄姊妹 俄巴底亞未死 未在天國裡面.
但是能夠和上主一起在事奉的裡面.
看上帝所要我們看的事物.
是一份極大的光榮.
今天我們能夠打開俄巴底亞書去看.
俄巴底亞怎樣和神一起讓他看見這幅圖畫.
我們應該有一份很大的感恩.
所以當我們讀俄巴底亞書.
第一樣東西原來是帶我們去看東西.
而看東西看完之後究竟怎樣回應這個下一步.
但至少我們要去看的是上帝想我們看的東西.
而不是我們覺得那段聖經想說甚麼.
為甚麼我要這樣說.
因為作為以下所看的論道以東人所說的內容.
如果你是猶太人來計或者猶大國的子民.
你就會覺得 嘩 真是好了.
還不看你有幾折朵.
是不是這樣的意思.
上帝想帶我們去看他有幾折朵.
還是上帝另有原因.
所以我們去看上帝所要看的東西.
第二 在這段聖經說到.
這個異象是由耶和華神去帶他去看.
所以這個異象是由主耶和華帶來的.
是屬於主耶和華的.
請留意在希伯來文裡面.
這裡不是用一般子神.

$^{401}$我們和合館都譯得準確.
直接指到是耶和華的墨使.
在希伯來文裡面是用到Adonai Yahweh.
或者我不敢讀Yahweh.
如果有猶太人在這裡.
不過 Adonai Yahweh是指我們的主耶和華神.
這是一個很親密的名字.
上帝要介紹他的尊姓大名的標誌.
而這個異象是由一個這麼個人化的神.
帶著俄巴底亞去看的.
不是一個很有距離的神的啟示.
而是一個有名有姓.
一個想你和他知道是怎樣的一位神.
個人的神耶和華.
去帶你去看.
就好像你去旅行一樣.
如果你去旅行.
你只是說有個導遊帶我去看某個風景.
和我那個老公某某某.
或者我的太太某某某帶我去看.
完全兩回事.
因為他作為導遊帶你去看.
就是一份工作.
有回報的.
但是你如果和一個你的親密人士.
你的丈夫你的太太你的好友.
他有名有姓帶你去看的話.
這個就是和你的人生的經歷.
上帝是一個Personal God.
是一個很貼我們身的一位上主.
他不是遠遠地和我們保持距離.
他是願意將他自己.
很親密地和我們一起.
去看我們他想我們所看到的事情.
這個異象的內容是關於哪一個呢.
這裡講得很清楚.
論以東說.
以東是一個什麼的民族呢.
在之前我們講過了.
在引言裡面.

$^{441}$原來他是和猶太國的國民.
有歷史及血緣緊密關係的一個民族.
從聖經神學而言.
當先知在以西結書裡面.
就是以西結這位先知.
去宣告了上帝要審判猶太國這個訊息之後.
分別就在小先知書裡面.
包括馬拉基或者我們所看的西班牙哈巴谷.
和兩位大先知書.
以賽亞和耶利米先知書裡面.
就看見上帝這個審判.
正正式式地實現在猶太國的身上.
在這裡我想停一停.
讓大家先搞清楚這個訊息.
我們所看的先知書本身其實是有它的結構意義的.
不是純粹說誰是大先知誰是小先知書.
誰是有寫作先知有些沒有寫作先知.
這些全部是神學性系統神學性的一個詮釋.
但是如果在聖經神學裡面.
因為聖經神學的核心是要將聖經自己.
它想告訴我們一個整全訊息.
它本身是一個活體來的.
聖經本身是一個在新約說是道.
差不多是耶穌來的.
它直接來跟我們說話的.
它想說些什麼呢.
原來先知書裡面有這樣的結構.
是隱藏在裡面.
猶太國被審判懲罰.
或者我們說他被滅亡.
是先由以西結這位先知將它宣告出來.
所以以西結先知可以說是在聖經裡面.
打響上帝懲治猶太國的第一號宣告人士.
接著在聖經裡面.
又自己透過不同的先知的記載.
告訴我們這個審判的過程是怎樣進行的.
剛才我列舉了那些先知書.
包括小先知書和大先知書.
我們很多時候說大先知書很重要很重要很重要.
是沒錯它很重要.

$^{481}$但是原來小先知書裡面.
都原來有三卷是介紹到猶太國怎樣滅亡的過程.
為什麼要這樣記載呢.
記載了是不是代表上帝的憤怒得到釋放呢.
原來還有一個層面.
就是我們包括今天我們所看的俄巴底亞書.
以及我們先前已經看過的約珥書.
這兩位先知都同時要介紹出一個.
我們少有所看的題目.
就是在亡國之下的疼痛.
痛是被懲罰的其中一個過程.
但是痛是有一個目的的.
很多人將基督教包裝成為一個很世外桃源的地方.
我很清楚告訴大家教會不是一個世外桃源.
可能是一個更加多掩藏著很多包著羊皮的狼的場地的好地方.
所以真正在聖經的教導裡面.
不是要將教會塑造成一個天堂.
而是要告訴大家教會是一起去面對自己的罪的時候.
所種種遇到的疼痛.
而在這個疼痛的過程當中.
能夠彼此提醒他的疼痛的意思是什麼.
為什麼我們會痛因為我們犯罪.
我們有罪在我們當中.
我們就會有很多種不同層次的痛在我們裡面.
而這些痛可能是解決不了的.
尤其是當你做了一些事情而得罪了那個人.
你是再沒有辦法跟他說對不起的時候.
那種痛是永恆地在你生命裡面刻在當中的.
我曾經做過一件事.
是現在我知道很難去修補的.
就是我在讀中學的時候.
因為我自己是屬於領袖生在學校.
領袖生有一個歡送會.
我們離開學校的時候玩得很瘋狂.
在玩的過程當中.
我不小心弄傷了一個同學的鼻子.
雖然不是一件什麼嚴重的事.
只是損了一點點而已.
但對於我來計算.
是我一生其中一個最大的一個的愧怯.

$^{521}$可以這麼說.
得到這個同學我相信都很難了.
除非有一天上帝帶到我可以跟他面對面見得到.
說一聲對不起.
但是大兄姐們.
當有一些比我這件事這麼小的事.
更大的抱歉你做了的時候.
而那個人你沒有辦法說對不起.
你一生的痛苦就在你裡面.
俄巴底亞書就是說這種痛隱藏在猶大國人裡面.
為什麼它會痛呢?.
因為以東人正正就是猶大國的.
我可以這樣說.
親兄弟種族的一個倫邦.
當你讀到俄巴底亞書的時候.
輪到以東人的遭遇.
當你看到什麼華臣要這樣處理以東的時候.
你就會覺得.
我應該的態度處理是開心還是不開心好呢?.
以東人對不起過我.
但是如果我們看回歷史.
因為以東人是屬於耳塑的後裔.
而猶大國是屬於雅各的後裔.
在這個發展歷史當中.
就讓我們去看到.
一個親兄弟如果昔日都由爭開始在這個地上出現.
到現在他仍然在爭.
爭究竟國度進入到你死我活的狀態.
而你看到對方現在失敗了.
你看到上帝現在要懲罰他了.
你是不是應該開心呢?.
還是你覺得有種難言之痛呢?.
這個很視乎於你自己對於和對方的關係.
究竟以東人和猶大國.
是否仍然視之為兩個兄弟邦呢?.
還是好像很多人說.
當一個家族大的時候.
什麼一代親兩代表.
三代咀妙妙.
是不是這樣呢?.

$^{561}$我昨天見到我的表妹失散多年.
終於約到了,昨天見到面.
親表的.
是不是這樣呢?.
還是突然間某一天我在某個商場裡面.
譬如first market place的parking lot.
和人爭車位的那個就是我的第三代表呢?.
是哪一個呢?是什麼關係呢?.
是不是這樣呢?.
這段聖經其實俄巴底下書.
其實是講一個這樣的背景.
因為如果我們去看耶利米埃哥書第4章第21節.
以及我們所提到的以西結書第25章第12節.
或者35章第10節當中.
都指到當時的新巴比倫帝國.
踐踏來到巴勒斯坦地區.
而在這個時候準備攻打猶大國的時候.
處於南方的以東國趕都趕不及地.
或者應該這樣說.
就是很急地去趕快向這位從北面而來的帝國的王.
巴比倫的尼布甲尼沙王.
向他投降.
並且怎樣協助巴比倫王尼布甲尼沙.
來攻擊耶路撒冷.
這些種種的事情.
就記載了在我們的俄巴底下書第11節至第14節的內容.
我們之後會看的了.
但是當我們看到以東人是這樣去對待一個親兄弟邦國.
現在作為讀俄巴底下書的猶大國子民或者神國的子民.
你的心態是怎樣呢.
是不是仇人見面憤外眼紅.
於是乎見到他跌倒.
你就樂此不疲呢.
甚至見到上帝現在動這個以東的時候.
你沾沾自喜呢.
或是你在這個時候.
你會嘗試回去用耶和華臣的角度.
耶和華臣為何要向以東人進行一個審判的宣告.
我們有沒有從這個角度去看呢.
所以讀俄巴底下書看似很短的一段經文.

$^{601}$但是卻勾畫出我們作為信徒.
在這個世界的上面.
我們怎樣去面對一些屬靈的掙扎.
不是代表你不想見到上帝的公義得到彰顯.
你不是不想見到.
你想見到.
但是上帝彰顯他的公義的時候.
是否想那些被懲罰的人變壞了.
我們就開心快樂呢.
還是在這個時候.
回頭去看看.
是否自以數雅各的年代.
那份一直掙扎.
去爭取自己個人利益.
這種思想是錯誤的呢.
所以當我們去看看這段書的時候.
我們會留意到有一連串的鋪排.
告訴我們知道的.
剛才我們略略講過了.
第一我們看到這段書第一個要討論的人物是耶和華.
主耶和華.
第二提及到的是這位將上帝訊息傳來的.
是一位先知.
俄巴底亞.
他是一位忠於神的僕人.
他在傳這個訊息的時候.
他的位置有沒有將上帝的意思扭曲了呢.
還是準確地將上帝的動機介紹了給我們.
接下來我們從這個角度去看.
第三我們看到以東人.
以東人在接著整卷的俄巴底亞書裡面.
我們會留意到一個現象.
你看一看就可以了.
除了見以東人之外.
你會見到一個名字叫做耳塑.
為什麼要在這個時候.
耳塑已經死了多少年了.
重提這個關係呢.
上帝透過俄巴底亞先知不斷提耳塑與雅各.
耳塑與雅各.

$^{641}$將耳塑和耳東這兩個字眼互相.
英文好說.
有時在西方世界久了用英文好說.
interchangeably.
不斷地互相交替使用.
其實想塑造一種什麼樣的訊息.
給讀俄巴底亞書的猶大國子民.
和後來的猶太人呢.
在這三者之間我們看到的.
就是這兩位兄弟的關係.
原來上帝一直沒有忘記過.
而俄巴底亞這位先知.
就將他準確地將這個目的介紹了出來.
不過你問我.
牧師講了這麼久.
其實你想講些什麼呢.
我都挺沉氣.
講這麼久才講.
你想講些什麼呢.
其實我想印證到聖經在.
彼得前書第三章第九節裡面.
有一句這樣的說話.
彼得前書第三章第九節這樣說.
「不以惡報惡.
以辱罵還辱罵.
都要祝福.
因為你們是為此蒙召.
好叫你們承受這福氣」.
這番說話是彼得.
向著當時一大群的猶太人去寫的.
誰得罪過我.
我就動他.
是不是這個意思.
不是.
原來在聖經裡面.
有一個很恆常的一個.
我們叫 counter balance force.
一個平衡我們會走錯的.
屬靈的心態.
不要以惡報惡.

$^{681}$以辱罵還辱罵.
我們不要這樣做.
聖經從來不鼓勵我們做這件事.
所以說上帝好像向我們伸冤.
很多人很喜歡講.
引用聖經那句.
「伸冤在我,我必報應」.
很喜歡將這句話套在上帝身上.
給我們生一口氣.
但是上帝其實不是伸冤.
上帝是在挽回.
將每一個屬祂的子民.
和上帝子民有關係的人.
都帶到一個被挽回的位置上.
而我們有沒有明白上帝這樣的心思呢.
俄巴底亞書正正就摸得準這個心思.
讓我們在上帝的啟示裡面.
透過上帝的詮釋角度.
不受人的感受的操控之下.
來明白一套上帝喜悅祂的子民.
當活的生命樣式.
聽清楚.
是合神心意的一份心意的樣式.
不只是行為.
是整個人的取態.
不只是取態沒有行為.
是兩者並肩.
所以當大家還記得我在引言的時候.
為什麼引用彼得向著那位耶穌所愛的門徒.
然後問耶穌這位你所愛的門徒.
他將來會怎樣.
很多人解釋那段聖經就說.
耶穌說關你什麼事.
不是關你什麼事.
而是耶穌的角度.
想這位將來的門徒.
他所愛的門徒.
一路怎樣去等候耶穌基督.
又是主的心思意念來的.
不是我們值得去關心的事.

$^{721}$你關心弟兄姊妹是需要的.
你想他的將來是好的是好的.
這個最重要的就是.
那段聖經裡面有一句話.
來跟從我.
什麼叫來跟從我.
你不明白上主的心意.
你就不知道怎樣跟從他.
有時如果你有叫一個你的同事.
你是他上司.
你來跟著我看看我怎樣做.
他不明白你教的東西的動機.
為什麼你會教他.
他來跟著你是沒什麼意思的.
今天很多信徒就是會在這樣的狀態.
你明白了上帝的心思.
你就知道怎樣去跟著主行.
當你讀到俄巴底亞書的時候.
你有沒有去跟著耶和華神的心思.
來去詮釋他看整件事.
那個目的是什麼.
如果你跟我是明白的話.
你的屬靈生命.
就可以藉著俄巴底亞書這塊跳板.
進入一個更加深的屬靈的境在.
可以跟上帝更加合一的心智.
來走面前的屬神的道路.
好了第二件事.
第一節的下半節.
在這裡神是藉著介入列國.
來實踐他的計劃.
這個也是很多時候.
曉講我們不曉看的一個經文.
我們看看下半節講什麼.
輪到以冬雪那裡.
跟著我們看了.
我為大家讀出來.
大家一起去看好不好.
我從耶和華那裡聽見訊息.
並有使者被差往列國去說.

$^{761}$起來吧.
一同起來與以冬爭戰.
我讀一次.
New Revised Standard Version的聖經給大家聽.
看看怎麼樣.
Thus says the Lord God concerning Edom.
We have heard a report from the Lord.
and a messenger has been sent among the nations.
Rise up.
Let us rise against it for battle.
如果你想看多一個譯本的話.
我可以讀New English Translation給大家聽.
他是這樣說的.
We have heard a report from the Lord.
and every was sent among the nations.
saying arise.
Let us make war against Edom.
在這裡首先以先知耶和巴底亞指出.
他聽到耶和華來到有一個報告.
一個報告.
即是這件事已經做了.
已經做了.
不做就沒有報告.
做了就有報告了.
有一個由眾多國家差來的使者團.
原來有一個組織組織了出來.
包含了很多國家一起成為一個使者團.
而這個使者團要作出一個宣佈.
從和合本聖經的讀法.
我們就好像是耶和華猜了一個人去.
然後去到向所有國家講.
一起來興起與以冬的爭戰.
不過在詮釋這段聖經的時候.
講到這個使者.
我們可以看到的是一個.
從多國組成一個使者的團體.
他們的組成是包括所有周邊的列國.
這些周邊的列國被聚集在一起.
他們有一個共同的宣言.
好像現在很多國家經常開什麼會議.

$^{801}$經常開什麼國際會議.
開什麼國際會議呢.
有很多宣言的.
好像上個禮拜有一個七國G7會議.
有沒有看.
大家熟不熟悉新聞的.
聽聞好像我們的總理講了一些話.
有沒有聽過.
他講了些什麼呢.
這些大家可以看一下.
但是在這裡不是私底下.
哪個總理和哪個總理哪個國家和哪個國家講話.
不是誰買誰怕.
不是誰討好誰國家.
是一起作出一個共同宣言.
他們的共同宣言是什麼呢.
起來吧.
一同與以冬爭戰.
原來當我們聽到這裡的時候.
我們就發覺到上帝有一個報告.
原來這個報告是講及周邊的邦國.
他們一起來好像.
撐猶大國來去處理以冬國.
你作為猶大國聽到這個時候.
開不開心.
假設而已不是上帝的工作.
好像全世界的人都來幫烏克蘭.
來去反抗俄羅斯入侵.
你作為烏克蘭人開不開心.
我上個禮拜才和一個烏克蘭的女士聊天.
她也知道多開心.
開心不是她的國家被人打.
而是她開心其他國家來幫她們守.
猶大國在這一刻聽見耶和華臣.
有一個這樣的報告.
而且是先知親口告訴我們.
他親身聽到耶和華臣有這個報告給他的.
可以說是全世界的人都來撐猶大國.
是不是這一波真是可以一洩心頭脂粉呢.
停一停.

$^{841}$以冬又和你不妥又想陷害你.
現在猶大國你見到其他國家來.
好像興起來去打以冬國.
好像撐你一樣.
你又很開心.
這是不是神的目的呢.
是不是俄巴底亞書想講的訊息.
不過在未回答這個問題之前.
我相信從一般人性.
你就會覺得好事.
是不是.
有人和你出頭.
我被人欺負.
你就多給我兩百支導彈.
我可以隨時射過去.
不用我動手.
你動手.
多好.
但是你有沒有想過.
為什麼這些列邦的人會願意一起興起來去.
好像.
我再重複.
好像支持猶大國.
好像很支持猶大國就是去打以冬人.
為什麼.
他們.
是不是很大意凜然.
弟兄姊妹.
在這個世界.
你猜有沒有人會做蝕本的生意.
有沒有人做過生意.
我爸爸做生意的.
怎麼會叫蝕本生意呢.
蝕本生意.
通常有個特色.
你根本不知道自己.
想做這盤生意是為了什麼.
在加拿大.
怎樣斷定一盤生意.
大家知不知道.

$^{881}$在法律上.
生意只得一個定義.
就是你會期望得到利潤.
什麼人會支持你.
是不是他們因公義來支持你.
誰想去賣軍火給那些被欺負的國家裡面.
因為那些人想在過程當中去賺錢的人.
弟兄姊妹.
我們看事情要懂得這樣看.
當眾人都來到有一個這樣的宣言.
來到去支持好像很合猶大國人的心意的時候.
強調很合猶大人國的心意的時候.
這是不是神的心意呢.
不是的.
牧師明明說上帝這樣宣告.
請聽清楚.
這是一個報告.
上帝說有一個報告.
有這樣的背景.
上帝不是說他在這裡同意這群人這樣做.
上帝是將這個報告講給這群.
要聽這個訊息的人知道.
但上帝就在這個時候做事了.
這群人是滿有自己的心思籌算.
用英文說.
This group of countries are opportunities.
他們全部都是利益機會主義者.
當見到你兩個國家打對了.
好了.
我們就從中當中去得到一些利益囉.
這個世界在歷史上從來沒有改變過.
這個做法.
哪個國家打哪個國家.
甚至你隔壁屋那個人和你爭執.
都是為了某些利益之小.
譬如以我家為例.
我後院那個人不是很happy我.
為什麼.
因為他見到我為什麼草沒有野草.
我不告訴你我為什麼後院沒有野草.

$^{921}$現在都有了.
愛莫能助搞不了那麼多了.
但是都比後院少很多.
他就不服氣.
為什麼你這樣呢.
你搞過什麼呢.
我見到他不搞我又很怒氣.
我想告訴大家.
每一個人其實在這個世界.
如果按照你個人意願去行事.
就有很多的東西在爭執.
為的都是你自己個人的利益.
而猶大國在這個時候.
是不是因此而開心見到好囉.
這麼多人來支持我.
去動一動.
我都想他死.
因為他害我嘛.
他怎麼害我啊.
歷史已經告訴你了.
我被巴比倫欺負的時候.
入侵他趕快去投降.
然後教他怎麼來攻打耶路撒冷.
還不是他害我.
這次我有機會報仇.
但是有沒有看清楚.
在聖經裡面.
這幫人其實是從中得利.
上帝很清楚這幫人.
是從中去索取利益.
但是俄巴底亞.
以為擅自將這件事講了出來.
當猶大國子民去看的時候.
他有沒有留意到.
這幫人其實未必是值得他們信賴的人呢.
為什麼很多基督徒順上下主不順主.
因為好像在這個世界.
對我的幫助多過.
很實質上的幫助.
是多過上帝對我的幫助.

$^{961}$於是教會就變質了開始.
要搞多點幫助囉.
將教會變成一個.
只是幫助人的機構.
從來對這個靈性的深度的建立.
是沒有需要去考慮.
為什麼.
建立了這麼有深度有什麼意思.
沒有什麼意思.
可能要付出.
但是在這種情況之下.
俄巴底亞就要講給猶大國.
或者甚至今天我們每一個讀者.
你要明白.
這幫人這麼厲害.
講得這麼宏大.
起來吧 一同來與以東征戰.
看起來好像是一件有很宏觀的幫助.
但卻原來是一份隱藏的陷阱.
作為神的子民猶大國.
縱使你被擄.
你有沒有想過你為什麼會被擄.
你在疼痛的當中.
有沒有想過上帝怎樣讓你回轉上來.
而當在俄巴底亞書裡面講到以東.
你是被上帝.
使用一些外邦人來使你們.
去到有一連串的糟禍.
其實上帝是不是也都.
在這個時候.
想你回轉呢.
在這裡我們可以學習到的就是.
我們真的不要按照我們自己的意欲.
來去活在地上.
這個世界有很多東西很吸引我們.
很多東西很容易使我們誤會上帝.
他行事的目的.
很多時候我們在被人間接地操縱了之後.
我們還以為自己好像在行公義.
我們參與了這一幫烈國來一起陷害.

$^{1001}$不要陷害.
一起來攻打以色列的仇人.
以東人.
我們就好像在行公義.
是不是呢.
在這裡我想講一點點.
一點點涉及一些政治的東西.
我不詳細講了.
不過有很多人經常問我.
牧師你對於過去在你的家人的原居地.
他們的情況.
你究竟支持還是反對什麼人呢.
我可以告訴大家.
很多人以為基督徒在過去這些日子裡面.
參與了一些運動是在行公義的.
很多包括基督徒.
他有個預設的理解就是.
總言之我們是被欺壓.
所以我們要反抗.
但是如果你想深一層.
事實是不是真的這樣.
你支持哪一邊.
其實他們都有個人的動機在背後.
最蠢的是沒有經過一對熟齡的眼睛和智慧.
去看這些事情.
然後隨便站邊.
在過去很多人經常以為我講.
你是什麼顏色的.
因為你姓黃的所以你是黃色的.
你不要再冤枉我了.
我不是黃色的.
我什麼顏色都不是.
我看的是事實的全面.
作為基督徒.
我們看的是上帝的眼睛看下來的時候.
看到什麼呢.
當看到人肉橫流的時候.
我們有沒有制止自己.
錯誤地投在某一些人肉橫流的狀態之下呢.
當你以為你正在行公義的時候.

$^{1041}$其實可能你已經在不經不覺之下被人利用了.
參與在一場你不應該參與的活動裡面.
在仕途行轉.
甚至我們所看見的仕途保羅.
他就是曾幾何時被人煽動的一個人.
他不是沒有聖經.
他是一個法利塞人.
他在加瑪列門下讀滿聖經的.
聽清楚.
他讀滿聖經的.
舊約的聖經.
他是非常熟悉的.
但是當我們在仕途行轉裡面的時候.
當他被當時猶太教煽動的時候.
去攻擊基督徒信耶穌基督的人的時候.
他甚至以此為喜樂.
他大發熱心.
但是當他的信主啟蒙了他之後.
他在羅馬書裡面這樣說.
我們一起打開羅馬書第十章的一至二第四節.
羅馬書第十章.
第十章一至二第四節.
找到之後我們一起讀好不好.
弟兄們我心裡所願的.
向神所求的.
是要以色列人得救.
我可以證明.
他們向神有熱心.
但不是按著真知識.
因為不知道神的義.
想要立自己的義.
就不服神的意了.
律法的總結就是.
基督使凡信他的都得著義.
看到沒有.
保羅去到他寫羅馬書.
的時候是經驗了他年輕的時候.
以一個很有聖經知識的人.
所經歷過的一連串的錯誤.
而寫出的反思.

$^{1081}$在這裡可以很清楚地告訴他們.
他說尤其是第二節.
我可以證明.
用自己的生命整個人來做一個證明.
他們向神有熱心.
因為他自己也曾經好像向神有熱心.
以為迫害基督徒就是一份好事.
但是誰不知.
但不是按著真知識.
什麼叫真知識.
有些人又說知識.
知識使人自高自大.
不要一下子這樣下定論.
知識就是你的眼睛如何看見真相.
你的眼睛如何看見真相.
剛才我們讀完羅馬書裡面.
就告訴我們當看不到真相的時候.
我們就用自己來定什麼叫真相.
就用了一個字眼.
義.
我們自定的義.
我要為我什麼來爭取我的權益.
我要為這個世界什麼來爭取權益.
因為用我的角度來看.
但上帝的義是什麼.
你有沒有去看.
你沒有看到的話.
那你就趕快去看.
而保羅他就看見這個真的知識.
是源自於耶穌基督.
當你明白耶穌基督的時候.
你就會明白什麼是真正上帝的義義.
今天有很多的基督徒.
卻沒有花心思時間去洞悉上帝在歷史當中.
他介入的時候的目的是什麼.
純粹就這樣看這段聖經.
俄巴底亞書.
就是等同於上帝懲罰以東來報.
他們如何向猶大國不善的仇.
不是這樣.

$^{1121}$所以當你讀俄巴底亞書的時候.
你聽我說到.
我肯定我100\%告訴你.
你沒有一本聖經書找到像我這樣解釋的.
十本就算九本.
他都一定告訴你.
俄巴底亞書就是耶華臣.
如何為猶太國向以東人報仇.
因為他們曾經如何傷害猶大國人.
但是當我們仔細去看.
俄巴底亞書的用字用詞的時候.
就告訴你原來不是報仇.
而是如何將這兩個國家的人的屬靈生命.
帶回上帝原本做的屬靈生命上.
透過兩兄弟彼此相愛來建立起來.
剛剛退休會才講完.
約翰福音13章講什麼.
去到最後耶穌給了一條新的命令你.
你們要彼此相愛.
人就認出你是我的門徒了.
兩兄弟在俄巴底亞書打生打死.
我們以為好事了.
好啊 公義得到彰顯了.
不是啊.
這段聖經是上帝很痛地去看見.
這兩兄弟他如何互相殘殺.
而如何去挽回他們.
記不記得我剛才講什麼.
聖經神學裡有兩本書表達了這種痛.
一本書你們現在知道了.
俄巴底亞書.
另一本書是約珥書.
你說牧師 是嗎.
是啊.
不記得 不知道.
那你回去看我講約珥書吧.
求主幫助我們.
在今天這個世界裡有很多人在欺騙不少人.
包括基督徒.
有些人好像KOL.

$^{1161}$Key Opinion Leaders.
在不同的層次叫什麼.
社交媒體裡面不斷煽動你們去看一些東西.
有些人甚至搞政治.
甚至有些基督徒給那些穩食的牧者.
你知道什麼叫穩食的牧者嗎.
哄得你開開心心.
只要你奉獻和他一起去玩.
這些騙得他們神魂顛倒.
滿足他們的需要.
你要搞什麼抗爭我陪你去搞.
你要搞什麼事我陪你去搞.
真章上帝公義.
扭曲上帝的意思.
甚至有很多成為一些政治的工具.
今天分為很多黨.
仇侵黨.
支持侵黨.
還有什麼叫擁侵黨.
有些叫民族主義黨.
我們做基督徒活在這個世界裡面.
看多些俄巴底亞書的時候.
當你的屬靈生命深化的時候.
我相信你就越曉得.
看得透這些人的動機是什麼.
他們全部都是列邦當中一起叫你興起.
一起去動而動.
用今天的說法.
誰是我們的敵人.
我們一起和你去動那些敵人.
你是不是隨波逐流呢.
如果是這樣的.
弟兄姊妹.
愧我作為一個你的牧者.
給這樣的教導你.
我愧對上帝.
我可能解釋俄巴底亞書不清楚.
讓你們不能夠深化你的屬靈生命.
因此這個只不過是俄巴底亞書的開始.
俄巴底亞書在這段聖經裡面.

$^{1201}$是叫所有的聽眾.
能夠好好把握上帝這個意象.
是看到上帝在歷史當中怎樣作為.
也都同時叫我們小心提防.
所謂大眾的主流的思想.
大眾不一定是對的.
人人說做的時候.
未必是對的.
我說的是從上帝的角度開始.
很多時候我們看到.
當這個世界越亂的時候.
這種的風氣.
就越容易引誘基督徒沉淪下去.
有沒有我們回到上帝的說話裡面.
聖經有一個教導叫做辨別諸靈.
諸靈.
不是豬的靈.
是文雅的說法.
Various spirits.
眾多的靈.
不是外面的靈.
不是邪靈這麼簡單.
不是那些仇人的靈性是怎樣.
包括你和我自己的靈.
你有沒有常常拿一個屬靈的X光機出來.
照照你自己的靈性.
什麼樣呢.
我們一起和心底人討論.
恩主多謝你給我們在這段聖經裡面的學習.
我們相信俄巴底亞這位先知.
當他看過上帝你給他的一連串的報告和意象之後.
他所帶來的訊息.
是讓我們怎樣深化去面對這個世界的時候.
有上主你的角度而去測驗.
求你的幫助.
讓我們每一個聽到今天訊息的弟兄姊妹.
都能夠從這段聖經裡面.
再次學會怎樣將我們的生命.
透過你的角度來去透視.
做到在保羅所說的.

$^{1241}$測驗上帝你的善良,純全,可喜悅的旨意.
求你的聖靈幫助.
我真的需要你的幫助.
禱告奉靠主耶穌基督得勝靈治球.
\newpage



\section{啟示錄 3:14-22}
\label{sec:032324}
\textbf{老底嘉:富足卻又貧窮的教會}~~~~譚綺敏傳道
\newline
\newline
連結: \href{http://www.church.com.hk/acms/content.asp?site=cdc\&op=show\&type=product\&code=032324\&layout=sermon}{\texttt{032324}} ~~~~ 語音日期: 2026-06-21
\newline
\newline
\hyperref[sec:032323]{\small{< < < PREV SERMON < < <}}
~
\hyperref[sec:index_chronic]{\small{[返順時目]}}
~
\hyperref[sec:032325]{\small{> > > NEXT SERMON > > >}}
\newline
\newline
啟示錄 3:14-22
\newline
\begin{longtable}{cl}
\hline
\hline
章節 & 經文 (和合本修訂版)\\
\hline
3:14 & \begin{tabularx}{0.7\textwidth}{X} 「你要寫信給老底嘉教會的使者，說：『那位阿們、誠信真實的見證者、神創造的根源這樣說： \end{tabularx} \\ \\ \relax
3:15 & \begin{tabularx}{0.7\textwidth}{X} 我知道你的行為，你也不冷也不熱；我巴不得你或冷或熱。 \end{tabularx} \\ \\ \relax
3:16 & \begin{tabularx}{0.7\textwidth}{X} 既然你如溫水，也不冷也不熱，我要從我口中把你吐出去。 \end{tabularx} \\ \\ \relax
3:17 & \begin{tabularx}{0.7\textwidth}{X} 你說：我是富足的，已經發了財，一樣都不缺，卻不知道你是困苦、可憐、貧窮、瞎眼、赤身的。 \end{tabularx} \\ \\ \relax
3:18 & \begin{tabularx}{0.7\textwidth}{X} 我勸你向我買從火中鍛鍊出來的金子，使你富足；又買白衣穿上，使你赤身的羞恥不露出來；又買眼藥抹你的眼睛，使你能看見。 \end{tabularx} \\ \\ \relax
3:19 & \begin{tabularx}{0.7\textwidth}{X} 凡我所疼愛的，我就責備管教。所以，你要發熱心，也要悔改。 \end{tabularx} \\ \\ \relax
3:20 & \begin{tabularx}{0.7\textwidth}{X} 看哪，我站在門外叩門，若有聽見我聲音而開門的，我要進到他那裡去，我與他，他與我一起吃飯。 \end{tabularx} \\ \\ \relax
3:21 & \begin{tabularx}{0.7\textwidth}{X} 得勝的，我要賜他在我寶座上與我同坐，就如我得了勝，在我父的寶座上與他同坐一般。 \end{tabularx} \\ \\ \relax
3:22 & \begin{tabularx}{0.7\textwidth}{X} 凡有耳朵的都應當聽聖靈向眾教會所說的話。』」 \end{tabularx} \\ \\
[1ex]
\hline
\hline
\end{longtable}
$^{1}$今日繼續講腦底家教會.
富足又貧窮的教會.
首先講一個小故事.
有位老先生第一次來澳洲旅行.
出發前興奮自滿地說.
我帶了一萬元去澳洲豪玩一個月.
羨慕嗎? 跟他的朋友說.
大家聽了都很羨慕.
有少少嫉妒的眼光望著他.
一萬元?.
這趟旅程真的很舒服很豪玩.
於是他一落飛機.
先去找一間華麗的餐廳去醫肚.
點了一些名菜.
吃完之後.
侍應就拿了張帳單來.
他一看 九十八元.
很便宜.
他瀟灑地從錢包裡拿出一張大鈔.
心裡說 九十八元一餐.
我可以豪吃很多餐.
誰知侍應一看.
就微笑地對他說.
先生 你這個是日元來.
不是澳元.
先生就大聲說.
我們這些貨幣是全球通用.
你這麼大的餐廳難道不收嗎?.
侍應就說.
先生 你這張一千元的日本大鈔.
兌換率只有澳幣九元左右.
你不夠找數.
老先生就很驚訝.
一萬元夠不夠找數?.
侍應就說 哎呀 都不夠.
我們看見這位老先生自以為自己很富足.
原來他很貧窮地去旅行.
問題不是他沒有錢.
是他對自己的財富有錯誤的認知.
所以今天說到老弟家的教會也是這樣.

$^{41}$你發覺剛才所讀的經文.
看見他自以為富足.
擁有很多.
一樣都不缺.
但主耶穌怎麼說呢?.
祂說 你卻不知道自己是困苦.
可憐 貧窮 黑眼 赤身嗎?.
為什麼呢?.
他們忽略了些什麼呢?.
今天我們就來探討一下.
我們先來祈禱.
撒華天父 求你打開我們的眼睛 讓你的話語再次感動我們.
讓我們知道什麼是我們生命上所追求的.
不要像老弟家教會一樣.
他自以為富足 其實是貧窮.
求主在這個時刻進到我們的心裡.
將我們心中所有的掛慮都清除出來.
以致我們可以安靜地來到神面前.
單單聆聽神的話語.
好使我們的生命可以改變.
得著豐盛的生命.
願主帶領我們.
聽我們祈禱 奉主耶穌基督得勝一命祈求 阿們.
今天我們繼續看七個教會的書信.
是約翰寫的.
我們先來複習一下.
因為是最後一個教會.
我們來看看這七個教會.
約翰在拔摩島裡對望.
其實就看到這七個地方.
以弗所教會的問題就是.
起初的愛心離棄了.
善梅拿教會就是面對逼迫的時候.
耶穌鼓勵他要至死忠心.
到別加摩教會.
他們的問題就是服從了巴蘭的教訓.
意思就是假先知的教導.
推瓦推拉教會的問題.
就是容讓耶駛別的淫行.
意思就是容讓了罪.

$^{81}$到撒迪教會.
就是耶穌吩咐他們.
在雖微中仍然要去警醒.
到上次說的飛拉鐵飛教會.
就是耶穌說.
在軟弱裡他們要持守他們的強壯.
所以今天來到第七個教會.
最後一個教會.
魯迪加教會的特點就是.
不冷也不熱.
不知道像不像你們現在的狀況呢.
想一想.
到底當中主耶穌有什麼教導呢?.
我們先來看看魯迪加這個字是什麼意思.
魯迪加這個字是.
人民(Laos)和大權利.
兩個希臘文組成的.
意思就是說.
魯迪加這個字.
是一個很好的字.
顯示這個地方很重視人民.
更加重視權利.
很人民至上.
權力至上.
民主自由的國家.
所以這麼高舉人權自由的地方.
但其中的教會.
你發覺剛才讀.
很令我們驚訝.
為什麼呢?不知道大家有沒有留意.
因為在七個教會裡.
只有魯迪加教會.
沒有神一句的稱讚.
一句都沒有的稱讚.
而且是很嚴厲的責備.
神更加說.
必從我口中把你套出去.
這麼厲害.
所以更加令人驚奇的就是.
魯迪加教會自己不知道.

$^{121}$自己覺得自己是富足.
一樣都不缺.
但主耶穌卻說.
你是困苦可憐貧窮.
黑眼赤身.
我經常提過這個句.
你發覺很強烈的反差.
所以我們更加要正視.
其實當中有什麼原因.
令到這樣呢?.
首先我們要認識魯迪加教會.
這個樣樣都有的城市.
魯迪加其實是.
今天土耳其的西南部.
可能很多人想去旅行.
都會喜歡去這個地方.
是一個非常重要的.
考古遺跡的旅遊地點.
所以很多人都喜歡去.
蓮花堡那邊的地方.
十幾年前.
發掘出來的古蹟很大.
你看這個面積多大.
在2013年更加列入了.
世界文化遺產的預備名單.
魯迪加的古城.
其實在公元前三世紀.
即是距離現在二千多年.
已經建立了.
當時是小亞細亞很富有的商業中心.
因為它在一個三條大道的交匯處.
所以其實令到城內的人很驕傲.
他們那麼富有.
所以有三大驕傲.
他們很自豪.
第一就是金融業.
好像我們的香港.
魯迪加是銀行商業的中心.
他累積了很驚人的財富.
在主後六十年.

$^{161}$他們建立了三百多四百年之後.
主後六十年有一場大地震.
城內的人民.
他們拒絕了羅馬皇帝的救濟金.
他們覺得自己可以重建這個城市.
所以他們自己拿錢出來建立.
再重建這個城市.
所以你看看他們多有錢.
還有就是紡織業.
魯迪加的土地很肥沃.
放羊就很好.
養羊是最好的地方.
他們有一隻黑綿羊.
羊毛很光滑.
所以他們製造出來的黑毛衣.
是世界馳名的.
到他們的醫學院.
有一間聞名的醫學院.
其中以鼻喉科最著名.
學院發明了一些耳藥粉.
或者一些眼藥膏.
都是世界馳名的.
所以當時的人覺得自己很富足.
一樣都不缺乏.
但在這個樣樣都有的城市裡.
有一個尷尬的問題.
就是他們缺乏了好的水.
魯迪加的城市沒有天然的水源.
所以要透過管道將遠處的水引入魯迪加.
在途中會受到影響.
魯迪加的水已經變質.
不像魯迪加的北面.
有溫泉可以有醫療功效.
也不像魯迪加的東面.
哥羅西有冷泉.
泉水可以解渴.
所以魯迪加的水.
充滿了沉澱的礦物質.
變了溫溫的水.
令人作嘔作悶.

$^{201}$又不可以治病.
治病又不可以解渴.
所以魯迪加這座富裕的城市.
慢慢地像泉水一樣.
不知不覺被感染了.
他們沒有撒旦的攻擊.
不像之前的教會有撒旦攻擊.
也沒有宗教政治的迫害.
因為他們距離比較遠.
他們應該是很安逸的.
但他們的教會被屬勢影響.
魯迪加的教會被屬勢影響.
好像不溫不熱一樣.
他們很忙碌去賺錢.
成為富貴的教會.
他們很自豪自己是富貴的教會.
主耶穌如何評價這個教會呢?.
剛才讀的15-17節我們看到.
主耶穌說:我知道你的行為.
你不冷也不熱.
但我很想將你們套出去.
這個不是好的現象.
不是溫暖的現象.
不要以為不冷不熱是溫暖.
不是這樣的現象.
主耶穌更加責備魯迪加教會.
魯迪加教會自己是不知道的.
他說:我是富足的,發財了,一樣都不缺乏.
但主耶穌說:你知不知道你是困苦.
可憐貧窮,核硬,赤身.
一直提著主耶穌所說的話.
因為這正正是魯迪加的問題.
他不知道.
所以到第二點我們要看.
魯迪加教會致命的問題就是.
他不知道.
當我們知道的時候我們可以改.
當我們不知道的時候我們可以改.
我們不知道魯迪加教會具體做了什麼事.
在城市的背景裡我們推測到.

$^{241}$可能裡面有不少商人,銀行家,有產階級.
也有猶太人在裡面.
他們將城市的價值觀搬進了教會裡.
以致活動很多,節目很多.
可能人數也不少.
有這些資源.
大家多活動的時候.
可能吸引越來越多人來教會.
大家嘻嘻哈哈去參加活動.
也歡迎有產階級的人士來教會.
他們看著教會的帳幕.
看著教會的帳戶.
華麗的教堂.
多姿多彩的節目.
他們覺得我們什麼都有了.
我們是蒙恩的教會.
蒙福的教會.
如果是這樣,你覺得這個教會有沒有問題呢?.
今天你看見這些教會.
如果是這樣,你覺得有沒有問題呢?.
如果我們教會是這樣,有沒有問題呢?.
但是不要忘記.
主耶穌怎麼說他們呢?.
祂說:你卻不知道.
你是勒困苦,可憐,貧窮.
狠硬,赤膝的.
不知道你聽到這句話.
有沒有一點害怕呢?.
你卻不知道.
就是說教會錯了.
你以為是對的.
教會是貧窮的.
你以為是富足.
瞎了眼的.
你看不到自己的問題.
卻以為看得很清楚.
好像耶穌去責備法利賽人一樣.
今天如果我們知道自己是不對的.
我們都會去改.
如果你知道自己婚姻的狀況有什麼不對的.

$^{281}$你自己都會去改.
如果你知道自己缺乏了.
你會去積極去追求.
但是如果你覺得自己一樣都沒有缺乏.
你就連動力去改都沒有.
這正正是老弟家教會的致命問題.
自以為富足.
卻不知道在屬靈上是貧窮的.
也沒有了屬天的財寶在天上.
在神永恆的角度是貧窮的.
具體來說.
教會出現了什麼問題.
他不知道呢?.
我們可以在第十四節.
去找回主耶穌說到的一些答案.
因為如果你有留意.
第十四節.
主耶穌怎樣稱自己呢?.
那為阿們的.
為誠信真實見證的.
在神創造萬物之上為元首的.
主耶穌這樣介紹自己.
有什麼意思呢?.
如果你有留意.
每封書信一開始之前.
耶穌都先介紹自己.
他針對教會的一些問題.
或者一些屬靈的危機去介紹自己.
好像是士梅拿教會.
他們面對壁壩的時候.
耶穌介紹自己一開始就說.
那首先的,末後的.
曾經死去而又活過來的.
為什麼耶穌要這樣介紹自己呢?.
就是讓士梅拿教會認識.
原來耶穌一直在掌權.
以致他們面對壁壩的時候不害怕.
但當耶穌這樣跟老弟家說.
自己是阿們的.
阿們是什麼意思呢?.

$^{321}$是一個希伯來語.
意思就是神是真真實實的.
我是真真實實的.
我是誠信真實見證的.
意思就是耶穌的見證.
是很全然誠實.
也是很真實.
他見證了自己是神.
他愛我們.
更加是神創造萬物之上為元首的.
這個意思也就是說.
耶穌是所有東西的元首.
是一切萬物之先.
是最高所有創造物的統治者.
物質的東西.
他都是一個統治者.
耶穌為什麼要這樣去介紹自己呢?.
其實是針對老弟家教會的問題.
如果老弟家教會有見到.
耶穌是那位真實萬物之上的神.
他應該是怎樣呢?.
他應該是高舉神.
但今天老弟家教會不是.
他高舉的是自己的財富.
他們自滿的是說自己已經發財.
一樣都不缺.
所以他們倚靠的是物質的東西.
他們沒有真實地活出主耶穌基督的見證.
如果我們真實地看見神是萬物元首.
我們自然會高舉祂.
自然會信靠祂.
有困難的時候我們會信靠祂.
也會活出神的見證.
耶穌想我們怎樣呢?.
我們就去做.
以致我們去活出.
「恆公二號憐憫傳謙卑的心」.
與神同行的見證.
我們應該是這樣的.
不是以世界的財富或物質享受來自誇.

$^{361}$所以如果大家記得.
我們上星期的程序表.
禱告的事項其實很好.
我們怎樣禱告呢?.
如果大家記得.
其實在現在的世界裡.
都是盛行物質主義的.
你發覺我們的世界也是.
有物質主義.
所以當我們以前強調人本主義.
什麼叫人本主義呢?.
就是強調人的價值.
人的尊嚴.
人是有潛能和理性的.
亦都主張在人類的世界裡.
我們是以人為福祉的.
這個就是人本主義.
但當人本主義慢慢去淡薄的時候.
其實這個世界是高舉物質的.
你發覺打仗各樣的世界的國家元首.
都是去搶奪這個物質.
但當這個物質強盛的時候.
我們要怎樣呢?.
上一次的禱告事項很好.
求主幫助我們信徒可以活出.
基督信仰歷來所重視的德行.
這是傳統上我們基督信仰有的德行.
就是我們有理智,有明智,有智慧.
有公義,有堅毅,有節制.
有仁愛.
都要在我們身上彰顯到耶穌基督的見證.
所以當教會失去這個見證的時候.
神就要責備和提醒老弟家教會.
所以主耶穌用老弟家最熟悉的事物.
去提醒他們的信仰.
在第十六節.
祂說:你既如溫水也不冷也不熱.
所以我必從我口中把你吐出去.
溫水剛才也提及過.
是指老弟家去到他們的泉水.

$^{401}$其實他們的泉水不可以治病.
也不可以像冷的泉水那樣解渴.
因為它充滿了一些雜質.
所以他們沒有功用可言.
所以其實是說到一個功用的問題.
所以當耶穌用溫水作比喻.
不是說我們有溫度.
其實是指這個溫度不是一個溫度的問題.
而是一個你有沒有效的問題.
有沒有功用的問題.
好像一杯溫咖啡.
你想想你喜不喜歡喝.
我們很喜歡喝熱咖啡.
因為有熱度.
但一杯溫咖啡一喝下去.
好像什麼味道都沒有.
很想吐出來.
一點吸引都沒有.
所以不冷不熱的信仰.
其實是指什麼信仰呢?.
就是指我們失去了對神的熱心.
更加是說到我們的信仰失去了功用.
所以今天教會有沒有失去功用呢?.
其實教會存在的目的是要向世界見證.
耶穌基督是拯救人的主.
使人可以有新的生命.
但當世界混合到教會.
要享受財富帶來的好處.
不是再要信靠神.
而是沒有去信靠神.
也沒有為教會的罪擔憂的時候.
也不再去傳福音.
不再關心周圍的社團.
不再為失散的人去祈禱.
只顧著自己的享樂.
只顧著來教會得到什麼好處.
其實這個教會已經變質了.
變成一個不冷不熱的溫水.
只是維持了聚會.
其實在神的眼中失去了他的功能.

$^{441}$所以說了這麼久.
其實你知道不冷不熱的教會出了什麼問題嗎?.
知不知道為什麼耶穌要責備老弟家教會這麼重要呢?.
大家知不知道出現了什麼問題呢?.
我講了很久了.
經常提到又富足又貧窮的教會.
其實出了什麼問題呢?.
其實老弟家教會出了一個致命的問題.
就是他不知道.
他不知道什麼呢?.
他不知道今天教會已經失去了他們的功用.
他不知道在教會裡已經沒有了神.
教會已經被世界同化了.
以成功取代了神.
以忙碌取代了親近神.
沒有了傳福音的使命.
我想最令主耶穌擔憂的情況.
不是教會的軟弱.
是他自己都不知道.
以為自己很好.
其實快要死了.
你覺不覺得有點可怕.
快要死了.
因為軟弱的人會呼求神.
自滿的人連呼求神都沒有.
因為他心裡沒有了神.
他不會找神.
所以今天可能是我們基督徒信仰的危機.
我們很自滿覺得自己不需要神.
我們不在乎神.
我們可能沒有否認神.
但我們不在乎神.
不在意神的教導.
講什麼都好.
聽一下就行了.
沒有活出見證.
如果你的信仰是不冷不熱.
你要小心.
剛才說到主必從我口中把你吐出去.
真的很嚴厲.

$^{481}$意思是不冷不熱的時候.
可能神對你的事情都不在意.
因為你都不在乎神.
神很在意你.
但你都不在乎神.
於是神在你生命上做的事.
你根本不在意.
根本看不到.
所以如果整個教會都是這樣.
其實教會就站立不穩.
所以你想想.
一個民主自由的地方.
一個繁榮富有的城市.
一個發財了.
一樣都不缺的教會.
最後怎麼收場呢.
教會就站立不住.
因為沒有了傳福音的使命.
慢慢吸引不到人來.
活動多.
不如去圖書館搞活動.
不如去俱樂部.
那些更多資源.
政府辦.
所以教會要這樣吸引人.
已經吸引不到人.
沒有了傳福音的使命.
沒有了福音的使命.
沒有了作門徒生命的承傳.
教會就會關門.
所以今天我看見不少西方教會.
教會賣了作商業用途.
被拆了.
或者被賣.
不知道我們有什麼感想.
縱然教會有不少問題.
但主沒有宣判教會要關門.
聖經上從來沒有叫教會要關門.
所以19節.
主耶穌怎麼說呢.

$^{521}$「凡我所疼愛的.
我就責備管教他」.
雖然老弟家教會一樣.
值得稱讚的東西都沒有.
但主仍然用疼愛.
去關心這個教會.
主耶穌其實深愛著教會.
祂的責備出於愛.
祂很想人去迴轉.
迴轉到去神面前.
接受主的管教.
不要落在一個自己都不知道的自欺裡.
所以第三點我們要看.
就是基督的挽回.
就是愛女的責備和醫治.
唯有挽回.
唯有醫治.
教會才會復興.
所以主對老弟家教會.
說要套他出來.
其實不是說他不要教會.
而是希望教會在衝擊裡.
去反省.
去醒悟過來.
所以主耶穌提出了三樣買的勸告.
第十八節.
「我勸你.
向我買火煉的金紙.
叫你付足.
又買白衣穿上.
叫你赤身的羞恥.
不露出來.
又買眼藥.
擦你的眼睛.
使你能看見」.
是什麼意思呢?.
你會發覺這三樣東西.
都是老弟家教會驕傲的三樣東西.
引以為榮的三樣東西.
第一樣就是金融業.

$^{561}$所以現在說到.
火煉的金紙.
意思就是.
你要用火煉的金紙.
去取代銀行的金紙.
不要以為自己有錢就很強.
真正的屬靈財富.
是經過用火試煉的信心.
然後對神有信心.
不是依靠金錢.
而是要依靠神.
去經歷神的實在.
然後也說到.
要穿上白衣.
白衣其實是取代了黑毛衣.
不要以為衣著上.
用黑毛衣.
會令到身體和暖.
所以更加是說到.
人犯罪之後會害怕.
害怕之後就會躲藏起來.
好像創世記.
說到考阿當犯罪之後.
唯有神賜予白衣穿上.
你才可以遮蓋罪.
今日教會雖然有罪.
但神是願意去遮蓋你的罪.
只要你是穿著基督的衣袍.
你再回到神那裡.
你不會羞恥的.
因為神是赦免我們的罪.
但我們更加要活出.
因信稱義的信仰.
第三就是眼若擦眼.
這就更加明白了.
被聖靈的光光照我們的心眼.
要我們看見自己現在的本相是怎樣.
以致更加看見神的榮耀.
所以今日主耶穌告訴你.
我在門外叩門.

$^{601}$若聽見我聲音就開門的.
我要進到他那裡.
他與我一同坐直.
這句說話很多時候.
我們傳福音會用這節經文.
耶穌叩你的心門.
向未信主的人傳福音.
我們說耶穌會進入你的心門.
但如果你看回上下文.
在這裡出現.
其實不是說未信主的人.
剛剛相反.
它是說一些我們信了主的人.
這個門是指教會的門.
是指老弟家教會的門.
所以老弟家教會自以為甚麼都有.
其實它將耶穌丟了出去門外.
擠了出去門外.
所以主耶穌他叩教會的門.
他很關心教會.
他主動去叩教會的門.
他對教會說.
所以這句福圖是很感動的.
即使現在很多教會.
或者現今教會很危險的光景.
就是很多節目將耶穌擠了出去.
其實今天教會很危險的光景.
就是物質主義.
和對人的冷漠.
剛剛在這個世代.
真是最危險的兩件事.
所以只顧自己關上門.
在教會的團契.
將主耶穌逼了出去門外.
所以今天如果我們教會是這樣.
我們真的要去反省.
但教會即使是這樣.
但主耶穌最感動的是.
我沒有放棄教會.
我很主動去教會的門.

$^{641}$去叩門.
可能你們太忙.
忘記開門給我.
可能你們太忙.
以致忘記將我放進去.
先親近主耶穌.
可能事奉已經變得很形式化.
可能教會有很多的爭執.
是因為大家不同的意見去事奉.
不要緊.
主耶穌說我已經在門外叩門.
如果有一個信徒聽到主的呼喚去開門.
耶穌很歡喜快樂地進到教會裡.
所以今天我們反省一下.
在我們生命裡會不會有很多的東西.
以致多到我們失去了最重要的東西.
就是和主耶穌七問的關係.
我們不知不覺地變成了溫水.
我們不再負神的話語.
我們更沒有渴慕讀經祈禱的心.
或者我們信了主多年.
我們的品格沒有改.
我們很容易發怒.
很容易不饒恕人.
也不願意去改.
沒有神的生命在我們裡.
我們今天去問問自己.
我們會不會是這樣的光景呢?.
要醫治這些不冷不熱的良藥.
關鍵是我們有沒有重新回到神那裡呢?.
我們有沒有和神一起呢?.
我們的生命有沒有和神連線呢?.
我們的生命有沒有去歡迎主耶穌和我們一起呢?.
如果要我們火熱起來.
我們要更加去接上和主同在的關係上.
就像電話沒電的時候失去功能.
就像教會失去功能一樣.
但怎樣可以恢復功能呢?.
就是一接電源就可以恢復功能.
所以今天我們只要回到主耶穌那裡.

$^{681}$我們熱切去追求祂.
我們教會的功能就會恢復.
所以十九節講到.
你要發熱心也要悔改.
什麼叫做發熱心呢?.
原文有重新點燃我們對主熱誠的意思.
也不是一次性的感動一下就沒有了.
而是持續的生命狀態.
就是你要一直有神在你心裡.
而悔改我之前牧師也講過.
悔改就是我們每天都承認自己的錯誤.
每天都相信主耶穌有赦免.
我們不斷犯錯不斷有神的赦免是真實的.
但我們也要承認每天我們都需要神.
願意回轉回到主耶穌的面前.
所以盼望我們重新有讀經祈禱.
盼望我們重新記得回來崇拜.
盼望我們重新回到小祖裡.
彼此去鼓勵去支持.
被聖靈的火重新點燃.
我記得羅馬書十二章十一節.
有一句很好的經文.
「因勤不可懶惰,要心裡火熱,常常服侍主」.
怎樣可以常常服侍主心裡火熱呢?.
其實就是要我們因勤不可懶惰.
我們要記得去服侍神.
如果你只是星期天回來坐下來聽.
享受別人給你的美好聚會.
你會發覺你的生命是不會有很大的成長.
這是我自己在屬靈旅程裡很深深的體驗.
我覺得我要去服侍主.
我才有生命的成長.
所以其實有時候牧者請你去服侍.
或者弟兄姊妹有需要的時候.
我們的服侍其實是令我們的生命很有成長.
其實服侍有很多層面.
你很積極地參加小祖也是一個服侍.
你來不來其實生命是沒有成長的.
你很用心地為別人代禱也是一個服侍.
我記得有一個很年長的老人家.

$^{721}$他已經八十多歲了.
他一個人獨居.
他信了主很多年.
但是他的身世也很可憐.
他信主沒多久.
丈夫就去世了.
他是一個宣教士的太太.
丈夫在越南宣教的時候去世了.
是一個美國的宣教士.
然後他就回到澳洲因為家人在那裡.
但是他放不下他丈夫的一些事情.
因為突然間被共產黨開門衝進屋裡射死了.
所以其實他是很久也擺脫不了的.
但是當他去祈禱的時候.
他經歷了很深的恩典.
他試過去露宿街頭.
因為他在美國的時候去夫家.
他以為回到夫家.
但是他丈夫的家人不太能接受他.
因為他是華人.
他家人回來澳洲.
所以他在美國試過去露宿街頭.
所以其實是很可憐的身世.
所以他有很多的精神問題.
但是當他深深經歷神的時候.
他現在已經八十多歲了.
快要九十歲了.
我經常跟他祈禱.
是我以前教的一個姐妹.
在我認識他的時候.
北宣也有很多人認識他.
因為他是北宣出來的.
他一天有二十多個人打電話.
可能沒有二十多.
二十個左右.
打電話去他家跟他祈禱.
他也會打電話去別人那裡跟人祈禱.
所以其實他雖然很多事情都做不到.
但是他去服侍神就是透過祈禱.
為國家為團契.

$^{761}$為不同的人去祈禱.
他心裡面很火熱.
就是因為他不斷地為人祈禱.
所以盼望弟兄姊妹.
很少的事奉,很少的服侍.
其實都是一個服侍.
令到你生命可以成長.
只要我們用心去做.
我們接待新朋友也是一個成長.
我們願意為家人去改變.
去服侍也是一個服侍.
所以神設立教會很盼望.
信徒裡面可以彼此服侍.
肢體可以連接起來.
跟神,跟人連接起來.
這個是很有價值的.
今日的世代你會發覺.
人很冷漠很疏離.
有AI就可以了.
有電腦就可以了.
如果你看見有些人.
跟電話去聊天.
就好像跟人去聊天.
你會發覺有些輔導的app都出來了.
所以其實你會發覺.
很多事情都不同了這個世界.
但是我覺得教會是很難能可貴的.
我們可以真實地見到面.
有不舒服的時候.
人可以牽著你的手.
給你一些溫暖.
有一個情感上的支持.
其實是很難得的.
所以盼望我們也要回到教會去崇拜.
去經歷神,去事奉神.
也說一些例子.
我們婦女組我很欣賞.
其實我們這班婦女組已經很大.
年紀很大.
但我越多分享的時候.

$^{801}$每次都說不停.
每次都有很多的分享.
還有一些姐妹.
都不想煮東西吃了.
都年紀大了.
有一個八十多歲的都不想煮東西吃了.
但她每次都帶東西回來吃.
問起她上來.
她就說看見有位姐妹帶兩三個菜回來.
每個菜都十人份量.
我們不用煮都夠吃.
所以我們每次都有東西帶回家.
那些姐妹就覺得.
沒理由那位姐妹都煮得到.
我煮不到的.
我比她還年輕.
所以大家都很迫切.
很殷勤地去煮東西吃.
所以其實這個就是互相感染的火熱.
也很欣賞教會的一些長輩.
他們盡心竭力地為教會.
出錢出力去建立教會.
不論在哪個地方.
在廚房,煮飯,清潔.
兒童主教學幫忙.
甚至維修,廁所,鏟草.
購買什麼日用品.
教會有很多事情做.
都有長輩的身影.
所以今天我們坐得這麼舒服.
其實很多長輩.
他們勞心勞力.
出錢出力去建立教會.
為什麼我們教會可以.
不用再有貸款.
都是長輩的一些出錢出力.
所以最後鼓勵弟兄姊妹.
第21節.
得勝的,我要賜他在我寶座上與我同坐.
就如我得了勝在我父的寶座上與他同坐一樣.

$^{841}$這是得勝的認許.
一個沒有稱讚只有責備的教會.
他都可以得勝.
今天我們都可以得勝.
只要我們重新讓主耶穌去掌權.
重新恢復我們與神的關係.
重新有美好的見證.
教會再吸引人來到教會當中.
我們就可以得勝.
同樣家庭的.
我們仍然可以得勝.
只要我們有神在我們當中.
所以盼望弟兄姊妹.
我們一起去讓神去掌管我們的人生.
第22節.
聖靈向宗教會所說的.
「凡有意的就應當聽也」.
我們一起去祈禱.
神話天父.
我們很容易像老弟家教會一樣.
去追求擁有更多的外庫.
更多的物質.
以致我們可以有豐足富裕的生活.
主耶穌這也是你給我們的恩典.
但當我們自滿去追求的時候.
我們知不知道自己在神面前是什麼光景呢?.
主耶穌伯伯.
我們在你面前都不是貧窮的.
不是困苦的.
求主開我們的眼睛.
讓我們看見自己的本相.
讓我們重新願意去重視與主耶穌基督的關係.
更加追求去愛主.
求主在這一刻進入我們內心的深處.
用你的話語再去提醒我們.
教導我們.
重點我們好像初初信主.
昔日愛你的心.
以致我們上上都有那個幹勁.
那個熱心.

$^{881}$願意去彼此服侍.
活出生命的見證.
吸引更多的人.
更多我們的朋友.
更多我們的親戚好友.
願意去相信主.
求主去帶領.
多謝主聽我們祈禱.
奉主耶穌基督得勝的命祈求.
阿們.
\newpage



\section{創世記 21:8-20}
\label{sec:032325}
\textbf{天父與人同渡時艱}~~~~蘇善美宣教師
\newline
\newline
連結: \href{http://www.church.com.hk/acms/content.asp?site=cdc\&op=show\&type=product\&code=032325\&layout=sermon}{\texttt{032325}} ~~~~ 語音日期: 2026-06-21
\newline
\newline
\hyperref[sec:032324]{\small{< < < PREV SERMON < < <}}
~
\hyperref[sec:index_chronic]{\small{[返順時目]}}
~
\hyperref[sec:032326]{\small{> > > NEXT SERMON > > >}}
\newline
\newline
創世記 21:8-20
\newline
\begin{longtable}{cl}
\hline
\hline
章節 & 經文 (和合本修訂版)\\
\hline
21:8 & \begin{tabularx}{0.7\textwidth}{X} 孩子漸漸長大，就斷了奶。以撒斷奶的那一天，亞伯拉罕擺設豐盛的宴席。 \end{tabularx} \\ \\ \relax
21:9 & \begin{tabularx}{0.7\textwidth}{X} 那時，撒拉看見埃及人夏甲為亞伯拉罕所生的兒子戲笑， \end{tabularx} \\ \\ \relax
21:10 & \begin{tabularx}{0.7\textwidth}{X} 就對亞伯拉罕說：「你把這使女和她兒子趕出去！因為這使女的兒子不可與我的兒子以撒一同承受產業。」 \end{tabularx} \\ \\ \relax
21:11 & \begin{tabularx}{0.7\textwidth}{X} 亞伯拉罕為這事非常憂愁，因為關乎他的兒子。 \end{tabularx} \\ \\ \relax
21:12 & \begin{tabularx}{0.7\textwidth}{X} 神對亞伯拉罕說：「你不必為這孩子和你的使女憂愁。撒拉對你說的話，你都要聽從；因為從以撒生的，才要稱為你的後裔。 \end{tabularx} \\ \\ \relax
21:13 & \begin{tabularx}{0.7\textwidth}{X} 至於使女的兒子，我也必使他成為一國，因為他是你的後裔。」 \end{tabularx} \\ \\ \relax
21:14 & \begin{tabularx}{0.7\textwidth}{X} 亞伯拉罕清早起來，拿餅和一皮袋水，給了夏甲，搭在她肩上，把她和孩子一起送走。夏甲就走了，但她卻在別是巴的曠野流浪。 \end{tabularx} \\ \\ \relax
21:15 & \begin{tabularx}{0.7\textwidth}{X} 皮袋的水用完了，夏甲就把孩子放在一棵小樹下， \end{tabularx} \\ \\ \relax
21:16 & \begin{tabularx}{0.7\textwidth}{X} 自己走開約有一箭之遠，相對而坐，說：「我不忍心看見孩子死」。她就坐在對面，放聲大哭。 \end{tabularx} \\ \\ \relax
21:17 & \begin{tabularx}{0.7\textwidth}{X} 神聽見孩子的聲音，神的使者就從天上呼叫夏甲說：「夏甲，你為何這樣呢？不要害怕，神已經聽見孩子在那裡的聲音了。 \end{tabularx} \\ \\ \relax
21:18 & \begin{tabularx}{0.7\textwidth}{X} 起來！把孩子扶起來，用你的手握住他，因我必使他成為大國。」 \end{tabularx} \\ \\ \relax
21:19 & \begin{tabularx}{0.7\textwidth}{X} 神開了夏甲的眼睛，她就看見一口水井。她就去，把皮袋裝滿了水，給孩子喝。 \end{tabularx} \\ \\ \relax
21:20 & \begin{tabularx}{0.7\textwidth}{X} 神與這孩子同在，他就漸漸長大，住在曠野，成了一個弓箭手。 \end{tabularx} \\ \\
[1ex]
\hline
\hline
\end{longtable}
$^{1}$今天是父親節.
我祝願在在的.
我們如果是爸爸的.
或者是我們本身的爸爸.
我們自己的爸爸.
到父親節快樂.
不知道大家對父親這個形象.
有些什麼感覺.
我記得當年我生第一個小朋友的時候.
有件很重要的東西要買.
生小朋友要買什麼.
尿片是要買的.
床也要買.
我還要買一條背帶.
在在應該很多人都會用過背帶.
如果你要照顧小朋友.
在德國的時候有一條背帶.
它有得獎的.
德國很多背帶都得獎.
我就看看這些得獎的背帶.
有什麼特別之處.
在德文背帶就叫Target.
它這個背帶的特性.
就是那幾條得獎的背帶的特性.
是同一個特性.
都是寫著Papa und Mama Target.
Papa和Mama的Target.
很有趣.
多數我們買背帶的時候.
你去買都是媽媽試的.
但是德國拿這幾條得獎的背帶.
是以爸爸為名的.
Papa und Mama.
為什麼呢.
這個背帶得獎的原因.
其中一個就是.
它的腰圍特別闊.
做了爸爸之後.
很多時候都會有幸福肥.
特別闊的.

$^{41}$粗一點的勒著腰.
肩膀的位置就闊一點.
讓嬰兒在背帶裡面的時候.
就不會夾死它.
所以有幾條背帶都是得獎的.
它們的功能都叫做.
Papa und Mama Target.
當年我們都買了一條.
我當然叫我先生試一下.
過了幾個月之後.
小朋友長大後回來香港探親.
他就很開心.
因為那個背帶很貼身.
貼著自己的胸口.
經常背著小朋友.
有一天跟我親戚吃飯.
我親戚就跟我先生說.
她說喂.
你真是.
她這樣說.
男人大丈夫.
在街上.
我必須要說.
那條背帶是綠色的.
棕色的.
背著這樣的東西.
背著一個嬰兒在裡面.
整個體統.
怎麼做男人啊你.
接著就說.
我先生很好.
兒子我有份的.
她說我們這些.
男主外女主內.
你給別人知道.
你不要給別人看.
背著還大搖大擺.
後來我先生當然很好.
沒有告訴我.
親戚說這句話.

$^{81}$他都不需要說.
因為我全部聽到.
他只不過在牆後面.
那塊板後面說我.
我心裡就想.
原來在香港.
當時的香港.
或者一些中國人的社會裡面.
我們看爸爸的角色.
都是可能不一樣的.
在西方來說.
很多西方的國家.
他們都很習慣.
爸爸一齊來到身兼務職地照顧小朋友.
不知道大家的感覺.
爸爸在你心目中.
你究竟是像西方那種感覺.
還是像我們中國人傳統的形象.
在今天的三代經課裡面.
其中一段經文.
就記載在.
《暢世紀》的第21章8到20節.
今天的主題是.
天父與人共度時艱.
父親節在每一個人.
人生裡面.
都是很不一樣.
我們的父親在我們的感覺.
在我們的經歷裡面.
都是很不一樣.
有人會經歷父親那種.
好像詩歌那種.
被包圍被保護的父愛.
甚至乎很肯定他.
但亦都有人會因為.
可能父親一些不完美的地方.
甚至乎這個父親是一位缺席的父親.
或者父親講過一些說話.
而感到傷害.
今日我們必須要面對一個.

$^{121}$真實的問題.
有時地上的父親.
或許會令人感到失望.
這樣.
我們還能夠仍然相信.
我們天上的父媽.
我們還能夠仍然相信.
我們天上的父媽.
在《暢世紀》第21章裡面.
描述到一個父親.
我不知道大家對他有一些什麼感覺.
但正正告訴我們知道.
地上的父親雖然有限制.
但並不阻礙天上的父親.
看顧我們每一個.
在剛才《生禮獨告》那段經文裡面.
講到一個背景.
就是有一個媽媽.
我們稱之為妾侍.
因為她是一個外邦人.
叫夏甲.
她生了一個兒子叫爾什瑪利.
如果根據那個時程.
這個兒子是出生先的.
之後就到一個正室.
我們廣東話可能叫大婆.
撒拉生了第二個兒子.
叫爾什瑪利.
就是兩個的媽媽.
兩個的兒子裡面.
開始了今天故事的這個場景.
我們看經文第八節那裡.
它說孩子漸漸長大.
就斷了奶了.
當爾什斷奶的那一天.
亞伯拉罕擺設豐盛的筵席.
在經文裡面我們覺得很普通.
孩子漸漸長大.
好像我們近幾十年都很流行.
擺百日宴 滿月酒.

$^{161}$在這裡爾什所說的漸漸長大.
其實不只是他慢慢長大成人那麼簡單.
在這裡所說的長大.
原來意思就是.
這個孩子慢慢長大的時候.
他慢慢要開始接受一個獨特的身份.
就是他要承繼爸爸的產業.
這個人將來是一個承繼人.
所以慢慢在家庭裡面.
那個身份和地位.
開始會有少少的不同.
甚至乎在這裡經文說到.
亞伯拉罕這個爸爸.
為這個兒子擺設了豐盛的筵席.
豐盛的筵席.
故事在這裡開始的時候.
發生了一件很重要.
也都影響心願的插曲.
第九節.
那時撒拉見到這個埃及人夏甲.
為她的丈夫亞伯拉罕所生的兒子.
嬉笑 取笑他.
這件事沒有什麼特別的.
可能就是說笑.
說一些話來刺激他.
或者嘲笑他.
但很可惜這個嘲笑.
觸動了正室撒拉的心.
於是他們家庭裡面.
起了一個很大的變化.
到經文下去的時候.
這個正室撒拉非常之憤怒.
你竟然笑我這個正室.
我這個漸漸長大.
開始去承擔去繼承的兒子.
於是他就有一個很憤怒的決定.
他就要亞伯拉罕.
要將這個侍女和兒子趕出去.
簡直是一套粵語長片的劇.
亞伯拉罕有什麼感覺呢.

$^{201}$十一字這樣形容.
亞伯拉罕當時為這件事非常憂愁.
因為這個關乎他兒子的事.
在這裡經文所說的兩個時間.
第一個時間是下甲的艱難時期.
同樣是亞伯拉罕的時間.
當以色列這個地方.
或者這個大國慢慢來到興起之前.
我們當然有很多不同的偉人.
在當中裡面一個很重要.
一個我們叫做.
一種孩子去成長的時候.
是將整個的一件事來掀起.
在這裡的主角.
這個孩子以撒被看重.
但在同時間.
他在一個比較.
這個正室之下的兒子.
被看重的時候.
這個妾侍的兒子以撒瑪麗.
他好像被看低一樣.
當有人歡喜快樂.
家庭慶祝的時候.
這個慶祝.
或者這件事.
正正是將某些人推到家庭的邊緣.
我不知道大家小時候.
怎樣去坐.
我記得我三年級的.
我們開始兩張桌子.
一起坐.
即是我旁邊有個女同學一起坐.
這個女同學第一次和他坐.
我覺得怪怪的.
但我小時候不太清楚.
總覺得他身上有一股味道.
我小時候不知道什麼叫體味.
又不知道什麼叫臭.
但我又覺得應該不是這些味道.
但覺得不知道為什麼.

$^{241}$有時候吃完早餐.
回學校坐在旁邊.
很想吐.
坐在旁邊.
過了三個月之後.
大概是冬天的時候.
那種味道更厲害.
更臭.
我就想應該不是體味.
體味應該是流完汗之後.
才會有味道.
冬天那種味道更酸.
更熟更臭.
後來我和這個同學熟了之後.
我就問這個女生.
我就說.
為什麼你的身體臭臭的.
她很尷尬的跟我說.
「是嗎?連你也聞到嗎?」.
我說「我聞到啊」.
「我坐在你旁邊的」.
她就說.
「是啊 是我的校服」.
我說「你夏天又臭」.
「冬天為什麼會更臭?」.
「你的校服有什麼問題?」.
這個女生就告訴我.
因為她家裡是一個.
比較重視男性的家庭.
她的校服裙 她的衣服.
爸爸不讓她晾高在公屋的騎樓.
公屋的騎樓是有些窗的.
衣服晾在那裡容易曬乾.
她說爸爸覺得她的裙 她的褲.
她的衣服晾高在上面的時候.
弟弟走過 爸爸走過.
就好像被姐姐 被同學騎著一樣.
這樣一輩子就不會發財.
更嚴重.
弟弟可能將來老婆也娶不到.

$^{281}$她這樣跟我說的時候.
我心裡覺得很難過.
我說「你的衣服怎麼辦?」.
她說「我的衣服就是放在廁所裡」.
「或是放在櫃子裡晾乾」.
「你媽媽的小女孩」.
「她的衣服長期都沒有被晾乾」.
「怎麼會呢?不會有一股臭味嗎?」.
在這個家庭裡 她被比較.
她被看下去.
這個感覺其實很難受.
這個家庭為這個弟弟.
來到這裡著想的時候.
大分裂了這個姐姐的需要.
在同時下隔在這個處境裡面.
她經歷了第一個時間.
就是好像看到以撒被重用.
但她的兒子好像被焊在一邊.
在這個時候.
一個非常艱難的決定就是.
亞伯拉罕清早起床.
就將這個餅和一瓶水.
放在她的肩膊上.
將她和孩子一起送走.
這個是下隔第一個時間被拋棄.
一個女人被拋棄是一件很慘的事.
如果你的丈夫死了的話.
對這個女士來說還好一點.
因為你是寡婦.
這個丈夫的親戚都會照顧你.
但這個女子她不是被稱為寡婦.
她被拋棄.
而且將她和孩子一起送走.
下隔面對的時間.
和亞伯拉罕一樣.
是一件很可惜.
也是一件很淒涼的事.
到第二個下隔要面對的時間.
當她被趕走的時候.
她內心充滿了那種孤單和絕望.

$^{321}$或者我們這樣想.
當人在失望的時候.
才會見到上帝的臨在.
當下隔帶這個小孩.
去到一棵樹的下面.
他自己走開大概有一箭之遠.
其實聖經在第十五到第十六節去描述.
下隔怎麼來將這個小孩放在曠野裡.
其實是充滿愛.
他不是將小孩放在曠野裡曬死.
他將他放在一棵小樹的下面.
而且自己也不能走得遠.
和他有一箭之隔.
即是很近的.
他望著這個小孩.
他就說.
我不忍心見到這個孩子死.
於是就坐在對面放星的大雄.
當下隔在這個荒野裡.
水用盡了.
皮袋的水也用盡了.
他知道沒有水就大件事.
於是下隔就放星的來到大雄.
他不忍心孩子死去.
當他的資源用盡了.
沒有人來幫忙.
前路好像沒有答案的時候.
人很容易就好像以為自己.
一定是走上絕路.
這是當時下隔面對的時間.
不知道大家有沒有曾經經歷過.
你很需要喝水.
但你沒有水喝的時候.
我記得幾年前.
我和一個姐妹去韓國旅行.
我們上下看耶穌山.
我們就去問當地旅遊中心.
我們要帶些什麼去.
他們說這是很簡單的.
我看你們兩個走30分鐘就差不多了.

$^{361}$這裡走上去.
拍幾張照片下來就可以了.
於是我們就想.
好啊 半個小時可以走走山.
我們上去吧.
怎知道我們上去的時候.
我們發覺山很大很美.
我們和途人一起聊天拍照.
慢慢我們大概10點鐘去到這個山.
11點就下來吃午餐.
後來我們不知為何.
不知不覺間我們在這個山上.
我們就越走越遠.
我們身邊的途人都沒有了.
只剩下我們兩個.
於是覺得很熱很渴.
我們10點鐘上山.
應該10點半下來.
但我們看時鐘已經12點半.
沒有人在這個山上.
因為太熱了.
我們就想.
糟了 我們昨天太自信了.
我們沒有帶水上山.
於是很感恩.
真的很感恩.
我這個同學竟然帶了一瓶150ml的水.
我昨天坐飛機的時候.
沒有喝到一瓶水.
還在背囊.
快點喝吧.
你想想150ml的水.
夠誰喝.
當然大家都不夠.
於是我的同學就說.
我有個方法.
我拿一條毛巾.
把水弄濕.
塞進你的嘴裡.
你嚇我一下.

$^{401}$這樣就可以了.
我哪有教你.
這是什麼方法.
我看電視學的.
那些戲劇都是這樣的.
塞進嘴裡就不會死的.
你覺得我有沒有這樣做.
當然沒有這樣做.
我說你75ml.
我75ml.
喝了它就算了.
一邊喝一邊走.
差不多走到1點半.
我們10點走到1點半.
那裡沒有水機.
沒有水源.
也沒有地方.
洗手間也沒有.
我們想喝生水也沒有.
後來我們就想.
糟了 怎麼辦.
那麼熱.
在山上真的沒有人.
後來我們就在遠處.
終於見到耶穌的像.
我和姐妹就說.
糟了 我們怎麼辦.
求耶穌幫助我們.
在這個時候.
我們真的沒有留意.
我們身邊有一輛巴士經過.
我們就說不要理了.
就上這輛巴士.
於是我們兩個就跳上了這輛巴士.
因為對方不會說英文.
我們就問對方.
是這樣嗎 還是這樣.
他就這樣.
我們就說是這樣嗎 還是這樣.
其實你明白我們說什麼.

$^{441}$是上去還是下去.
他就知道是這樣.
即是下山了.
不要再打他們上去.
然後我們就下山了.
下去了.
我記得一下去的時候.
其實我們都沒有坐下.
我們就直直在門外.
開關門等.
我們一見到有下了市區.
就立刻想著出去.
我很記得那一刻.
下了市區的時候.
剛剛見到有間餐廳.
你會怎樣.
真的很口渴.
差不多四五個小時沒有喝水.
又熱.
於是我們兩個女孩.
一人一邊徒手把門撬開了.
我們忘記了.
原來韓國的巴士.
和香港一樣都可以按鐘.
因為我們沒有按鐘.
對方沒有開門.
我們就撬開了門.
就立刻衝去喝水.
我想說的是.
如果人在一個非常乾渴的環境裡.
缺水的話.
他真的會有一點出現幻覺.
當時我也是有一點.
所以我們已經忘記了.
坐巴士要按鐘.
和要叫人開門.
不是自己撬開門.
我想說在這裡的下甲.
他在水源裡耗盡.
你想想當時.

$^{481}$他有多麼的失望.
正正當資源耗盡.
沒有人幫忙.
前路沒有答案的時候.
人容易就會哭.
放棄.
或者感到非常之失望.
但是在21章第17節裡面.
是一個非常之關鍵的經文.
他說上帝聽見孩子的聲音.
上帝的使者從天上呼叫下甲.
如果你看回經文前一點的話.
當時誰在哭.
經文沒有形容小朋友在哭.
下甲的兒子二十瑪利.
沒有形容孩子在哭.
反而在第16節的下半節.
他就坐在對面.
放聲大哭.
放聲大哭的是誰?.
是下甲.
沒有聲音的是誰?.
應該是他的兒子二十瑪利.
但可能有一點點嬰兒的聲音.
但這裡上帝是聽見孩子的聲音.
上帝是先聽見的.
是孩子的聲音.
你知道.
下甲很愛這個小朋友.
他將他放在一個樹下.
在一支箭不遠的地方.
來坐著看著這個孩子.
他哭不是因為自己被拋棄這麼簡單.
他哭是因為他不忍心看到這個孩子的死.
這個母親很愛這個孩子.
上帝知道.
上帝是聽見孩子的聲音.
他不是聽見下甲的聲音.
他先是聽見孩子的聲音.
於是他就說.

$^{521}$下甲下甲你為何這樣?.
不需要害怕.
如果我們作為父母.
或者我們可以想一個父母.
自己是怎樣也不要緊.
最重要是什麼?.
自己的孩子好.
有一年.
幾年前我帶一些家庭去中大旅行.
其實沒有什麼好去的.
不過就帶他們上去.
看看那個山很漂亮.
我們自己沒有訂車.
因為我們當時人數本來是40人去.
最後是120人.
我們又訂不到車.
不要緊.
我們就坐中大的校巴上去.
一輛校巴就上去.
導師就上去.
第二輛校巴又上去.
到第三輛校巴了.
慘了.
司機認得我們了.
又是你們這班人.
我們校巴不可以讓你們整輛坐上去.
我們要留給學生和中大的職員.
我們就說.
不好意思.
不行.
夠了.
我們要關門了.
司機很絕情.
馬上就關門了.
那怎麼辦呢?.
當司機說要關門那一刻.
我們還有十多個家庭.
還有孩子在那裡.
那些父母.
特別是那些父母就說.

$^{561}$你不用理我.
蘇雪瑩.
你帶我走.
我有兩個兒子交給你.
於是我和你那個導師.
在三四分鐘裡面.
我們突然在手裡接了十個小孩.
交給你.
我自己走上山.
你幫我好好照顧他.
帶他上去.
很熱.
不要讓我兒子在那裡.
這樣就走了.
那輛車就走了.
兒子就在窗裡哭了.
爸爸.
於是我們就上了山.
你知道我想說的是甚麼嗎?.
原來在一些很艱難的絕境或處境裡面.
我們做父母的.
其實是不忍心看到孩子受苦.
上帝所顧念的.
不單只是夏甲和這個孩子.
上帝所顧念的是夏甲所顧念.
因為他知道夏甲愛這個孩子.
上帝的回應並不是先分對錯.
亦都不是先看見痛苦.
上帝是告訴他知道.
我聽到你的聲音.
我知道你的需要.
這個對孩子的呼喚.
對夏甲當時來說.
上帝的聽見是一個直接的安慰.
亦都是一個直接的介入.
上帝聽見孩子的聲音.
在裡面是最軟弱.
最微小.
最無助的一位.
但上帝所看見的.

$^{601}$不單只是孩子.
更加是這位孩子的母親夏甲.
上帝的幫助除了聽見之外.
還有甚麼幫助呢?.
在第十九節裡說.
「上帝使他的眼睛明亮」.
「他」的意思是夏甲.
使夏甲的眼睛明亮.
我們心想.
是否曬到很厲害.
沒有辦法看東西呢?.
是否眼睛有事要做手術.
這樣就明亮了?.
上帝使他的眼睛明亮.
眼睛明亮這件事.
不是指擦乾他的眼.
或者不是指擦亮他的眼睛.
使他的眼睛明亮.
這個原來意思就是說.
上帝開啟他的眼界.
上帝開啟夏甲的眼界.
使他以新的角度去看.
他身邊的人與事.
這個很重要.
上帝使他的眼睛明亮.
是代表上帝開啟夏甲的眼界.
使他用新的角度去看.
他身邊的事和人.
這個開啟在當時.
第一個時間.
第二個時間裡面.
對夏甲來說.
是一個很重要的開始.
好像一個人在絕望.
在沒有水喝.
在盡頭的時候.
重現那種對人的信任.
和那種盼望.
人在絕望的地方.
仍然可以找到出路.

$^{641}$好了 上帝聽見.
上帝使他的眼睛明亮.
那上帝怎樣實質來幫助夏甲呢?.
我跟他說.
他說這裡有一個井.
你起來去看那井.
上帝不是說你起來去那井.
經文裡說.
上帝說你起來去看那井.
看是一個很重要的意思.
剛才上帝是聽見孩子的聲音.
聽到了.
使夏甲的眼睛能夠明亮.
所以能夠看見.
看見那一口的井.
為他提供出路.
井是有水.
亦都是代表生命之源.
人在絕處.
人在艱難的時候.
我們最想的是看到什麼?.
是看到出路和盼望.
在這裡.
夏甲是一個孤單被拋棄的人.
他竟然在這個曠野裡面.
他能夠看到這一口井.
是生命當中的盼望.
上帝聽見孩子的聲音.
起來叫夏甲來到去走.
開了夏甲的眼睛.
叫他看見一口井.
去拿他的皮袋裝水.
然後給孩子喝.
在這裡我們可以納為七個動作.
這七個動作正是代表上帝的愛.
來怎樣改變夏甲當下的處境.
用以身的角度來看.
曠野不等於沒有末路.
曠野也不等於沒有井.
上帝的英文.

$^{681}$上帝的救恩.
並不是短暫的解困.
而是為這個家庭.
為這對母子帶來未來和方向.
在這裡到經文最後的時候.
第20節那裡這樣說.
上帝與這個孩子同在.
他就漸漸長大.
住在曠野.
成為了一個弓箭手.
剛才讀第8節的時候就說.
孩子漸漸長大.
這個孩子是子爵二十.
但第20節這個孩子是指著.
夏甲的兒子二十瑪利.
經文剛剛在第8節和第20節.
做了一個非常好的首尾呼應.
上帝不單止叫以撒在裡面慢慢長大.
去承繼.
到這個故事即將要結束的時候.
上帝沒有偏心.
同樣這個非正室的兒子二十瑪利.
上帝也說.
上帝.
是有上帝這個字.
上帝與這個孩子同在.
上帝要祝福他.
他就漸漸長大.
承繼,承擔.
住在曠野裡面.
成為一個弓箭手.
上帝沒有偏心.
上帝是愛人.
有時候我們未必能夠.
馬上知道.
但今天這段經文給我們有兩個結論.
上帝的愛.
是超乎人的限制.
地上的父親.
可能是不夠完美.

$^{721}$或者好像亞伯拉罕那樣.
在困獄裡面他自己也沒有選擇.
他只可以看著他的愛而離開這個家庭.
亞伯拉罕愛孩子.
但他仍然有他家庭的限制.
仍然有他人的限制.
父親可能是失焦.
甚至失去控制的能力.
但上帝在這裡.
第二點.
人的錯失不等於上帝的離棄.
在這裡.
下甲一定是做錯.
他為什麼要來嬉笑這個耳殺呢?.
但是在這裡.
上帝告訴我們.
縱然有人做得不對的地方.
甚至人因為不對.
而產生一個可能不開心的結局.
下甲在曠野裡面.
但他從來沒有被上帝來遺棄.
這個正要彰顯.
上帝永遠不變的憐憫和愛.
甚至願意為一個做錯事.
做得不夠完美的人.
開出一條出路.
父親節提醒我們.
雖然我們不是每一個人的童年.
都曾經被愛或者溫柔地對待.
但聖經今天藉著這個訊息.
要我們轉向.
天父的愛勝過地上人的限制.
天父的愛不會因為人的過失.
而等於要離棄你.
雖然我們在我們的家庭裡面.
可能我們的處境都是不一樣.
但上帝告訴我們.
祂並不是等人完美.
才來看顧我們.
可能我們在家庭裡面.

$^{761}$或者在我們父的形狀裡面.
我們不多不少都有一些不滿意.
甚至乎缺失的地方.
在我們人生裡面.
我們總會有機會跌倒在荒野.
但我們要相信我們要堅信.
上帝的看顧是隨時都在.
祂看顧你.
聽見你.
呼叫你.
帶你走到人生一個新的處境.
開你的眼界.
叫你看到生命的正旁.
豐富的賢侄.
和孩子的同在.
並不止於一個開心快樂成的家庭.
同樣是一個似乎被遺忘的家庭裡面.
上帝都依然看顧.
何處幫助我們.
我們一起祈禱.
我們親愛的天父.
你還在人失落.
軟弱.
時間.
甚至是被傷害的時候.
你仍然聽見.
你仍然呼叫.
你仍然看顧.
在今天父親節裡面去提醒我們.
在我們地上的父愛可能不夠完全.
但主啊.
你透過下甲的故事.
你如何在他的荒野裡面.
來帶領他和孩子走過.
主啊.
求你都看到我們每一個人的需要.
可能我們的荒野.
我們的時間都不一樣.
主啊.
求你看到我們.

$^{801}$求你成為我們的安慰.
更重要的.
主啊.
求你在我們生命當中.
看到這個生命之源.
看到這一口井.
這個新的盼望.
求主使我們.
讓我們在家庭.
在教會裡面.
活出更美好的見證.
我們將我們未信諸的家人.
特別是我們的父親.
我們都希望能夠將他帶到上帝你的面前.
我們今天不只紀念父親節的情誼.
我們更加紀念上帝你永恆不變的愛.
你補足我們的不完美.
你補足我們的不完全.
多謝你.
聽我們在你面前的禱告.
是奉主耶穌基督的名祈求.
阿們.
\newpage



\section{士師記 8:11}
\label{sec:032326}
\textbf{聖經人聖經事:因為他們坦然無懼!}~~~~黃克勤牧師博士
\newline
\newline
連結: \href{http://www.church.com.hk/acms/content.asp?site=cdc\&op=show\&type=product\&code=032326\&layout=sermon}{\texttt{032326}} ~~~~ 語音日期: 2026-06-17
\newline
\newline
\hyperref[sec:032325]{\small{< < < PREV SERMON < < <}}
~
\hyperref[sec:index_chronic]{\small{[返順時目]}}
~
\hyperref[sec:032327]{\small{> > > NEXT SERMON > > >}}
\newline
\newline
士師記 8:11
\newline
\begin{longtable}{cl}
\hline
\hline
章節 & 經文 (和合本修訂版)\\
\hline
8:11 & \begin{tabularx}{0.7\textwidth}{X} 基甸從挪巴和約比哈的東邊，從住帳棚人的路上去，趁米甸的軍兵以為安全的時候攻擊他們。 \end{tabularx} \\ \\
[1ex]
\hline
\hline
\end{longtable}
$^{1}$各位大英姐妹,眨眼又一日啦,又一齊睇聖經啦.
今日ge 題目就《因為他們坦然無懼》.
呢個題目我諗jor 好耐,因為gwo 段經文就有好多好ge 內容.
咁呢就有第二個題目我都好想揀ge .
不過呢我就最後用呢個《因為他們坦然無懼》.
呢四個字呢今日就送比大家.
因為坦然無懼呢做人黎講都真係好重要ge .
有坦然無懼呢幾個字會令到我地發揮到更加大ge 力量.
更加好ge ,gwo D 潛質盡量釋放,而且壓力比較細.
希望大家都能夠坦然無懼ge .
生命樹同大家一齊讀聖經,今日聖經人,聖經事.
睇下聖經背後上帝ge 作為.
先睇呢一個地圖先,呢個地圖呢.
我就著時機搵唔到,不過呢就搵返瓦角同耳塑ge 地圖比大家睇下.
略為呢就介紹下兩個地方.
一個叫梭角,一個叫比奴伊勒.
比奴伊勒呢大家可能聽得多D 熟D 啦.
比奴伊勒記唔記得呢,就係瓦角係呢個瓦博道口.
就預備一陣遲D 明天就會見到gwo 個耳塑架啦.
咁就大家分別jor 幾十年.
咁好想見返哥哥向佢賠罪ge 耳塑真係.
係瓦角真係好想同佢哥哥認罪.
係我當日呃jor 你個長子名分又真係唔好呀咁樣.
咁呢就不過好驚.
因為而家我已經有老有嫩有曬我ge 家當牛羊無數.
咁而哥哥耳塑呢亦都係.
係呢個所以ge 地方呢佢已經係好有錢.
亦都好大族.
咁呢就一派人去話哎呀哥哥你ge 弟弟瓦角想搵你.
就見返面喎.
哎呀好呀我即刻帶人去迎接佢.
帶幾多呀.
四百人嘩瓦角聽到真係暈低.
起返身死囉帶四百人黎.
唔通佢想報仇咁樣.
咁呢個呢心底你面前意識ge 投射呢就令佢覺得.
嘩四百人梗係想報仇呀如果唔係點使帶咁多人呢.
其實佢哥哥耳塑係想黎歡迎佢最盛大的歡迎.
咁呢個藍色線呢就係阿爾耳塑呢就一路黎喇喎.
咁呢條河呢呢條河呢就叫做瓦博河.

$^{41}$咁係呢個有好多地方係過去架啦過河啦所以叫渡口.
咁呢就係一個過河ge 地方.
咁就係瓦博河ge 瓦博渡口gwo 度.
咁瓦角呢就自己過jor 河.
早一晚就係gwo 度呢就擺低晒D 仔女.
所有族人就一個人過夜喎.
想安靜下啊聽日就見到阿哥啦.
如果佢一殺我我全家人就冇架啦.
咁gwo 晚冷靜下咁樣自己點知gwo 晚夜晚一個人黎同我摔交.
打摔角咁終於呢就打到天光.
咁gwo 人就話我要走啦記唔記得呢件事呀我要走啦.
咁就唔得唔比你走除非你為我祝福啦其實佢知道gwo 個係上帝.
你為我祝福啦咁gwo 人呢就為佢祝福你改個名啦叫以色列.
咁以色列人呢就係由瓦角開始架喎.
瓦角ge 重要性叫以色列改名做以色列.
其實呢你見到佢完全不比亞伯拉罕失色架喎.
gwo 個係信仰上上帝ge 心目中.
因為你冇見到叫做亞伯拉罕人架喎.
今日見到叫以色列人架喎而家個國家名都叫以色列架喎.
都係由瓦角開始架咁呢就係呢個地方呢就瓦博道口gwo 度摔交.
咁然後gwo 個天使呢就最尾就捏一捏佢隻腳咁佢呢就跛jor 啦.
瘸子啦變成咁佢呢就話呢個地方.
我要改個新ge 名叫比魯伊啦咁樣咁呢就係咁樣啦.
咁係度見返哥哥囉喎跟住呢就咁呀.
阿哥呢就話我地一齊返去我住gwo 度啦我已經預備jor 一個大地方比你住架啦咁樣.
咁佢好好阿哥你先行啦我呢人多gwo D 細路仔又細D 羊呢又細.
我地行得慢D 架你先返去先返去咁呢哥哥就先返去囉喎呢個藍色線先返去.
呢度呢就係約旦河黎架有條有條河流入去呢就叫瓦博河.
咁呢哥哥先返去我跟住黎跟住黎點知瓦角gwo D 古惑ye 呢係都冇改到.
見阿哥行遠jor 呢佢即刻掉頭就去jor 唔行下面喎就去jor 疏角呢個地方.
原來疏角呢係一個當時ge 市集黎架係有好多疏角呢係gwo D booth.
咁介紹先我地今日呢係有gwo D gwo D ge 嘉年華會呀咁好多個攤位ge .
疏角呢其實就好多個攤位咁解ge 就係gwo 度呢買賣喎咁佢呢就走jor 去gwo 度就變jor 個富戶係gwo 度啦.
呢就疏角同埋比魯伊啦如果以佢地十二支派入去迦南地黎講呢.
呢個地方呢都係河東呢邊呢都係佢地幾個支派lor jor ge .
所以疏角同埋比魯伊啦人呢都有相當部分係佢地以色列自己人黎ge .
即係以色列十二支派ge D 人黎ge 夾集jor D 當地居民啦咁樣呢個位.
就知道係約旦河ge 東邊好近ge 地方自私ge 第九章機電和跟隨他的人.
三百人將來呢係天上真係好感恩呀點解將來如果我地係天上面呢會我一定會搵機電機電機電你gwo 三百人.
係邊呀我介紹大家識同我一齊出生入死ge 三百人.

$^{81}$呢三百個真係勇士呀佢勇呢係來自信心ge 而唔係來自佢地ge 武藝高強.
呢三百人去對付米電人第一仗係拎住一個陶瓷器造ge 水瓶裡面擺把火把.
另一隻手就拎住個豪國係咁啦冇帶刀咩冇講得好清楚.
佢地只係去吹豪國佢地會呼喊耶和華和機電的刀就係咁樣ge je .
然後打佢地打爛個瓶D 拎火把出黎聽到啪一聲咁然後呢就得為都係火閃閃咁呀.
嚇到D 米電人亂曬係神使米電人潰散.
三百人贏jor 第一仗之後呢一直跟住佢呢就去到約旦河追埋佢地呢gwo D 米電人要逃返去.
約旦河ge 東邊然後呢就向南走返自己家鄉呀嘛.
咁而米電人黎到呢除jor 米電人南邊阿瑪力人呢仲帶jor 好多東方人東方人咩意思呀.
約旦河東ge 人呀嘛包括jor 摩爾阿曼以東等等gwo D 人.
帶住三百人去到約旦河呢要渡過喇喎過約旦河.
留意第一點.
雖然疲乏還是追趕.
咁呢一句說話呢句聖經呢.
係我地讀神學ge 時候呢好多時我地同學之間學院都勉勵我地呢.
知道係好辛苦呀雖然疲乏.
還是追趕.
就呢做傳道人呢都好多時gui ge 壓力又大ge 咁其實.
疲乏ge 不過咁囉其實我都想講下做媽媽壓力大唔大呀疲唔疲乏呢.
超級疲乏.
但係點呀有仔都係度你還是要追趕呀係咪呀咁你返工係唔疲乏呢超級疲乏呢壓力好大呀還是追趕人生就咁喇.
雖然呢好多時候係好疲倦呀不過冇辦法呢我仍然要趕路呀人生ge 路就係咁.
好唔容易.
gwo 300人點解會疲乏呢打仗呀嘛.
跑呀嘛追呀嘛.
廝殺呀嘛會唔會呀梗係會啦做一陣15分鐘20分鐘運動都會到坐係度啦.
疲乏不過.
會都好啦我地繼續追.
但係個問題就係呢300人追得太耐呢係其實佢地去ge 時間去打米電人呢係冇帶食物架喎.
有冇帶住好多人呢去跟你煮飯喎.
所以gwo 300人呢又gui 仲係追緊.
就追到我地淨係講ge 疏國gwo 度啦.
過河嘛.
機電就對疏國人說求你地比D 餅我跟隨gwo D 人食啦.
因為佢地疲乏而且好餓呀可唔可以比D 餅我地呀.
問疏國人話我地追趕米電人ge 兩個王就係呢叫做西巴同埋薩摩那而家我地呢追緊gwo 兩個王gwo 兩個首領.
咁呢就追緊佢地而家呢過jor 河佢地走緊返鄉下我地繼續追緊.
其實當時ge 疏國人裡面呢好有可能住係疏國呢個城je 呢個地區je 但其實好有可能佢地都係以色列人黎架喎.
點解咁樣同兄弟講ye 呀求你比D 餅我跟隨我ge 人食啦因為佢地好gui 啦好耐冇食ye 又打jor 仗咁耐.
我地追緊米電ge 王呀大家同一個目標架喎點知.

$^{121}$疏國人ge 首領就回答西巴和薩摩那已經在你手裡你使我地將餅給你ge 軍兵嗎.
睇多次呀其實呢句呢好得意喎不過只不過係中文ge 翻譯呢令到我地有D 唔係好明je .
其實呢個係一個比較尖刻D ge 挑剔ge 問題黎架各位疏國ge 領袖可唔可以比D 餅我地三百人食呀.
我地打jor 好耐好gui 好肚餓啦我地追緊西巴同埋薩摩那呢兩個米電首領.
咁巴閉西巴和薩摩那你地捉到啦咩已經在你手裡啦咩咁你走黎問我lor 餅食咁巴閉係咪已經捉到呀捉到我比你.
咁喎其實咁ge 意思呀你都未捉到就黎lor 餅咁快咁樣喎.
真係機電真係有D 嬲呀機電就話啦有一日如果耶和華將西巴和薩摩那交係我手之後.
我就會返黎用野地ge 荊條荊棘一條條嘛當做邊咁樣同埋紫棘紫棘即係話一條gwo D 有刺ge 木材呀.
咁樣我就會刺棍咁樣我就會打你地因為佢唔肯比呀嘛仲笑佢嘛lor 餅食捉到啦咩捉到先算啦.
機電就話好有一日如果機電都好謙卑架喎好有信心喎機電就話如果耶和華比我捉到gwo 兩個領袖返黎ge 時候呢.
我就會對付你點對付你呢我就會用荊棘一條條嘛我呢好似鞭咁樣鞭你同埋呢我會用刺棍打傷你.
打傷又冇呀打死呀我打死你咁樣咁就話我打傷你點解呢自己人呀嘛咁呢就機電呢就跟住去啦去到隔離附近.
當時大家睇個地圖係咪疏國隔離咪比魯伊勒啦去到比魯伊勒啦又問當地ge 人可唔可以比D 餅我地呀我地呢三百人好餓啦即係打jor 好耐.
而家追緊兩個首領呀米田人都係咁講點知比魯伊勒ge 人同疏國人回答ge 話一樣捉到啦咩捉到先算啦捉到咪比你囉咁o下.
咁機電真係好嬲咁佢就話如果我有一日能夠平平安安冇事又唔餓死返到黎呢我會拆jor 你呢個樓房拆jor 你呢個食咁樣囉喎咁呢就都都好和平ge 竟然咁講.
就咁肚餓ge 人行過路人你都比餅佢啦何況我地以色列一齊去打緊仗哪時西班和薩姆拿呢就並跟隨ge 軍隊都在加國約有幾多人呢留意有一萬五千人走緊佬ge 呢兩個王米建王都有一萬五千兵馬.
其實呢萬五千人呢係東方各族一齊黎攻打以色列全軍淨返一萬五千人淨返一萬五千人gwo 個機電戰狼三百已經被殺ge 呢有十二萬軍兵係拿刀ge 所以呢即係話起碼米電人黎攻擊佢ge 軍隊呢攻擊以色列呢係十三萬五千已經比機電.
追殺jor 十二萬啦其實機電唔係好殺佢地因為佢地冇乜武器呢要樽咁樣打咋嘛不過呢.
耶和華市米電ge 軍隊與自己用刀擊殺弟兄呀擊殺佢地ge 同胞佢呢就話聯曬會見到就斬癲jor 咁樣.
咁而家淨返走ge 呢係一萬五千人.
追佢ge 三百就好啦機電就係拿巴和約比哈ge 東邊就往著張鵬ge 人路上去就殺敗jor 最後gwo D ge 米電軍兵.
咁呢因為佢地咁好打呢最後gwo 萬幾都比佢追殺埋呢係因為佢地.
六色字係解釋黎架係因為他們坦然無懼呀呢四個字呢就要送比大家架啦呢個係我要講ge 第三點黎架.
返返上去先.
第一點呢同大家講鼓勵大家雖然疲乏.
還是要追趕有好多ye 呢係咁架啦.
好累啦好累唔係理由我呢要完成佢ge 咁呢就gwo D 人唔比佢食喎咁雖然疲乏.
唔比ye 食呢係只不過係呢段聖經ge 插曲黎架je 其實佢地係三百人追緊一萬五千人.
嘩好大個數目然後點解佢地咁厲害呢能夠到最後擊敗埋一萬五千人呢係因為佢地坦然無懼.
最大ge 力量唔係佢地人多佢得三百人人地有一萬五千人一萬五千人對三百嘩個比例真係實在不成比例.
但佢仍然能夠贏到係因為乜ye 係因為上帝幫助佢地之餘呢係佢地坦然無懼.
但我要同大家特別講呢四個字呢因為他們坦然無懼就係我地中文聖經有一個希伯來文ge 字呢係冇係度出現ge .
因為佢地應該係佢地是坦然無懼咁呢有一個hia係是is咁解ge 呀或者佢地呀咁佢地佢地是坦然無懼.
坦然無懼呢係一個字je 喎hia betach係chkge 呀hia betach係是安然或者叫是平穩或者叫安然平穩.
大家知道咩叫安然平穩嗎.
就係呢面對住一個好大ge 考驗.
一陣返工有ye 要present啦一陣要見老闆啦要比佢zar 啦或者一陣呢就要去睇醫生啦醫生就話頭零零咁樣睇住個file好耐唔講ye 頭零零咁樣或者要去見牙醫啦.
好多人去見牙醫一坐落張床gwo 度呢成個薑硬曬好驚或者一陣要做D 咩事啦.
好心慌意亂咁呢個呢就不是安然不是平穩架喎點解佢地能夠變得咁好呢.
如果呢你現在gwo D 運動員啦gwo D 世界級運動員其實如果去跑步賽跑個個ge 速度差唔多gwo D 去到奧運gwo D 差唔多.

$^{161}$入到奧運決賽裡面個個速度差唔多點解有D 人會贏呢就係佢個臨場臨場ge 心理素質好好.
人就叫素質士就叫質素佢地ge 心理質素好好佢會贏其他人零點幾秒係因為佢可能心理ge 質素比佢好D .
咁呢班戰狼三百呢三百個人呢佢地ge 心裡面ge 質素好好呀係咩呢安然平穩所以佢地表現出黎呢係好好架.
音樂又係武術又係打打戰鬥又係打交又係考試又係演講又係當佢安然平穩ge 時候佢地好好表現ge .
gwo D 人雖然餓不過呢佢地ge 心裡面應該講死都唔驚仲驚乜ye 呢.
嘩如果連死都唔怕都冇咩怕囉所以佢地打起黎仍然好奮勇心裡面好平安.
原來平安穩妥係可以有一個作戰ge 力量ge 令到gwo 三百人呢係戰勝jor 一萬五千人.
gwo 一萬五千人同呢三百人ge 唔同係邊度呀gwo 一萬五千人個心驚到急急震呀見路就跑呀見人阻住佢就已經拔刀斬佢啦.
咁於是呢gwo 一萬五千個驚恐ge 人反而比呢一個三百個安穩平安然平穩ge 人呢係勝過jor .
咁gwo 三百個人已經冇飯食好耐架喇喎哎呀腳都軟曬就靠佢自己淨返ge 體力嘩繼續追趕.
而呢獲得全勝ge 咁呢最緊要呢就係第一有意志力雖然疲乏仍然要趕路要追趕.
第二呢雖然人少不過唔緊要我地仍然得ge 最緊要就係坦然無懼坦然無懼呢個最重要啦.
好咁呢就西巴和撒姆拿逃跑啦激烈追趕佢捉住gwo 兩個王咁呢就連兩個王都捉埋gwo 萬五人呢驚散全軍走曬啦.
咁約拿丹就對喇頓住要忽然間黎係度咁全軍驚散咁呢就事先記到呢度跟住我有一節呢係參考比大家ge .
呢一節呢就係撒母耳記上ge 十四章第六節呢個經文好緊要啦約拿丹約拿丹係邊個呢就係以色列第一個王.
素羅ge 兒子當時呢佢ge 好朋友就叫大衛咁呢個以後ge 事黎架喇去到撒母耳記架喇約拿丹Jonathan呢就同拿兵器ge 少年人.
當時打仗呢係有好多少年人跟住gwo D 戰士架喎gwo 個少年人做咩架係一個學習ge 戰士ge 學習ge 小戰士後生仔咁呢係比佢地細D ge 做咩架.
就係幫手服侍呢個戰士啦飲水啦著衫啦帶住D 兵器架倉啦把刀跌jor 話有第二把啦即係gwo D 咁樣咁你去到邊個帳棚gwo 度處呢佢就服侍佢啦咁gwo D 打鬥ge 人如果係高級D ge 軍事呢佢有個少年人跟班ge .
咁約拿丹就同佢拿兵器即係拎劍黎咁gwo 個咁呢就講啦不如我地去到未受國禮ge 人ge 防刑gwo 度啦未受國禮當時佢地ge 對手就係非利時人.
佢就話gwo D 未受國禮即係唔屬於神ge 人或者呢個或者呢係相信應該咁講如果相信耶和華為我地施展能力就得架啦.
咁不如我地直接去呀嘩呢個約拿丹呢個傻仔同埋gwo 個拎兵器ge 少年人少年人識打架喎可能拎劍架喎lor 刀黎la 有D 把咁樣咁呢個人就話我地不如直接我地我跟你一齊去進攻佢地ge 營呀咁樣喎.
非利時人好好打架喎佢進攻非利時ge 營地呀咁個少年人話好呀我地去呀如果神幫助我地就得架啦真係咁大膽呀因為耶和華使人得勝不在乎人多人少架呢句螢光掌呢送比大家啦心裡面呢有一個記得下有個記住耶和華使人得勝不在乎人多人少架.
gwo 三百個人係戰勝jor 十五萬幾ge 米電兵靶.
米電係個統稱姐包括埋阿瑪利人同埋好多東方ge 民族一齊結果失敗係因為耶和華使人得勝係不在乎人多人少架.
咁我同大家分享到呢度啦再次睇返個題目佢地得勝除jor 因為神ge 幫助必然啦係因為佢地本身坦然無懼呀心裡面係好釋放好安穩好平靜.
是死如歸嘩個力量真係非常之大啦咁就祝福大家呀生活上呢靈性上呢雖然疲乏仍然追趕.
面對好大ge 挑戰但係心裡面係坦然無懼因為上帝使你得勝係不在乎人多人少在乎耶和華施展大能祝福大家一齊祈禱.
主幫助我地好多時我地遇到事情馬上第一ge 就係諗我地手上有幾多籌碼我地ge 手臂有幾多能力.
馬上諗ge 我地識幾多朋友有邊個熟ge 人可以幫到手有邊個人可以比到我地指導有邊個人幫我地出面擺平.
我地馬上就諗呢一個但係我地係事司記係薩姆爾記我地見到ge 呢神使人得勝原來不在乎人多人少係在乎施展大能的上帝.
更加幫助我地係疲憊ge 身軀裡面仍然追趕係倚靠你ge 時候心裡面係坦然無懼.
以致我地獲得更多ge 恩典與能力比我地更多仰望以求更有信心誰聽其求奉耶穌基督的名.
阿們.
\newpage



\section{約翰福音 10:1-18}
\label{sec:032327}
\textbf{從一扇門,到心深處}~~~~鄧瑞強牧師博士
\newline
\newline
連結: \href{http://www.church.com.hk/acms/content.asp?site=cdc\&op=show\&type=product\&code=032327\&layout=sermon}{\texttt{032327}} ~~~~ 語音日期: 2026-06-21
\newline
\newline
\hyperref[sec:032326]{\small{< < < PREV SERMON < < <}}
~
\hyperref[sec:index_chronic]{\small{[返順時目]}}
~
\hyperref[sec:032328]{\small{> > > NEXT SERMON > > >}}
\newline
\newline
約翰福音 10:1-18
\newline
\begin{longtable}{cl}
\hline
\hline
章節 & 經文 (和合本修訂版)\\
\hline
10:1 & \begin{tabularx}{0.7\textwidth}{X} 「我實實在在地告訴你們，那不從門進羊圈，倒從別處爬進去的，就是賊，就是強盜。 \end{tabularx} \\ \\ \relax
10:2 & \begin{tabularx}{0.7\textwidth}{X} 那從門進去的才是羊的牧人。 \end{tabularx} \\ \\ \relax
10:3 & \begin{tabularx}{0.7\textwidth}{X} 看門的給他開門，羊也聽他的聲音。他按著名叫自己的羊，把羊領出來。 \end{tabularx} \\ \\ \relax
10:4 & \begin{tabularx}{0.7\textwidth}{X} 當他把自己的羊都放出來，就走在前面，羊也跟著他，因為牠們認得他的聲音。 \end{tabularx} \\ \\ \relax
10:5 & \begin{tabularx}{0.7\textwidth}{X} 羊絕不跟陌生人，反而會逃走，因為不認得陌生人的聲音。」 \end{tabularx} \\ \\ \relax
10:6 & \begin{tabularx}{0.7\textwidth}{X} 耶穌把這比方告訴他們，但他們不明白他所說的是甚麼。 \end{tabularx} \\ \\ \relax
10:7 & \begin{tabularx}{0.7\textwidth}{X} 所以，耶穌又對他們說：「我實實在在地告訴你們，我就是羊的門。 \end{tabularx} \\ \\ \relax
10:8 & \begin{tabularx}{0.7\textwidth}{X} 凡在我以前來的都是賊，是強盜；羊沒有聽從他們。 \end{tabularx} \\ \\ \relax
10:9 & \begin{tabularx}{0.7\textwidth}{X} 我就是門，凡從我進來的，必得安全，並且可進出，找到草吃。 \end{tabularx} \\ \\ \relax
10:10 & \begin{tabularx}{0.7\textwidth}{X} 盜賊來，無非要偷竊、殺害、毀壞；我來了，是要羊得生命，並且得的更豐盛。 \end{tabularx} \\ \\ \relax
10:11 & \begin{tabularx}{0.7\textwidth}{X} 「我是好牧人，好牧人為羊捨命。 \end{tabularx} \\ \\ \relax
10:12 & \begin{tabularx}{0.7\textwidth}{X} 雇工不是牧人，羊不是他自己的，他一看見狼來，就撇下羊群逃跑；狼抓住羊，把牠們趕散。 \end{tabularx} \\ \\ \relax
10:13 & \begin{tabularx}{0.7\textwidth}{X} 雇工逃走，因為他是雇工，對羊毫不關心。 \end{tabularx} \\ \\ \relax
10:14 & \begin{tabularx}{0.7\textwidth}{X} 我是好牧人；我認識我的羊，我的羊也認識我， \end{tabularx} \\ \\ \relax
10:15 & \begin{tabularx}{0.7\textwidth}{X} 正如父認識我，我也認識父一樣；並且我為羊捨命。 \end{tabularx} \\ \\ \relax
10:16 & \begin{tabularx}{0.7\textwidth}{X} 我另外有羊，不屬這圈裡的，我必須領牠們來，牠們也要聽我的聲音，並且要合成一群，歸一個牧人。 \end{tabularx} \\ \\ \relax
10:17 & \begin{tabularx}{0.7\textwidth}{X} 為此，我父愛我，因為我把命捨去，好再取回來。 \end{tabularx} \\ \\ \relax
10:18 & \begin{tabularx}{0.7\textwidth}{X} 沒有人奪去我的命，是我自己捨的；我有權捨棄，也有權再取回。這是我從我父所受的命令。」 \end{tabularx} \\ \\
[1ex]
\hline
\hline
\end{longtable}
$^{1}$各位弟兄姊妹 早晨.
願神的恩惠常與你們同在.
今天是父親節.
願神祝福每一位父親.
每一位父親都是天父在人間的代表.
付出了自己 守護了家庭.
成就了兒女的生命.
保羅在哥林多前書第四章第十五節這樣說.
「師父雖有一萬 為父的卻是不多」.
所謂師父 類乎教練.
重點是訓練一個小朋友成功.
但是保羅口中的父親卻是不同.
父親的重點在於以高貴的人格.
引領孩子歸向神.
提升孩子生命的神聖來函.
近來有一則很熱門的新聞.
一位很出名的教練.
以公開羞辱的方式來訓練球員.
效果可能很成功.
令球隊的成績進步.
但這肯定不是保羅口中的父親.
為父的教孩子尊重其他人.
教孩子領悟生命的尊嚴.
其實比成功更加重要.
多一些這些父親.
將來的世界才不會變為格鬥的場所.
元天父引領每一位父親成為真正的父親.
之前我已經講過頭三卷福音書.
每卷書都很不同.
馬太福音用心表明.
耶穌在什麼意義上超越猶太教.
以激發耶穌的同胞.
用心在耶穌身上領悟神的新作為.
馬可福音表明.
耶穌十字架的人生的沉重.
讓我們每一個人用嚴肅的態度.
來對待這種生命.
盧卡福音卻藉著耶穌.
去顯明天父的大愛在人間.
人生被天父的大愛承托著.

$^{41}$而約翰福音卻表明.
耶穌是宇宙的大道.
這個大道灌注人心.
讓人成為流露大道的器皿.
用約翰福音的話來說.
這就是永生.
真正豐盛的生命.
今天我們看約翰福音一段.
都很熟悉的經文.
就是那一段耶穌是好牧人的經文.
在耶穌的大道裡.
讓我們細味一下.
我們應該有的信仰的生命.
這段經文開頭是說一道門.
最後卻是說到內心最高的領悟.
就好像引領我們從一扇門進入.
然後一步一步進到內心的深處.
我們看約翰福音第十章一至第十八節.
經文都頗長.
我把它分成五個段落.
一個段落,一個段落來說.
第一個段落是約翰福音第十章一至第六節.
經文都印在我們的程序表較後的部分.
約翰福音第十章一至第六節.
耶穌說:我實實在在告訴你們.
人進羊圈,不從門進去.
倒從別處爬進去.
那人就是賊,就是強盜.
從門進去的,才是羊的牧人.
開門的就給他開門.
羊也聽他的聲音.
他按著名叫自己的羊.
把羊領出來.
既放出自己的羊來.
就在前頭走.
羊也跟著他.
因為認得他的聲音.
羊不跟生人.
因為不認得他的聲音.
必要逃跑.

$^{81}$耶穌將這個比喻告訴他們.
但他們不明白所說的是甚麼意思.
第十章第一節.
耶穌說:我實實在在告訴你們.
實實在在.
耶穌很認真地說.
希望我們每一個人都認真地去聽.
在這裡用了舊約.
牧羊人和羊的這個意念.
這是舊約描述.
神和祂的子民關係的一個經典圖畫.
譬如我們很熟悉的詩篇第23篇.
耶和華是我們的牧者.
我們的牧羊人.
耶穌在這裡很認真地說.
我實實在在地告訴你們.
神和祂的子民之間的關係.
有一些很重要的道理.
這段經文的主要動作就是.
「進羊圈」.
羊當然就是子臣的子民.
就是我們.
「進羊圈」就好像帶領我們.
一步一步地去了解.
信仰生命的真正含義.
這裡分開兩種人.
兩種進羊圈的方式.
第一種就是.
不從門進去的就是賊.
就是強盜.
他們是屬於黑暗的.
從門進去的就是牧羊人.
是屬於光明的.
黑暗和光明的對立.
這是《約翰福音》一貫的主題.
唯有從正門進羊圈的.
那個才是真正的牧人.
羊認得牧羊人的聲音.
會跟著牠走.
我們這裡應該沒有人比羊大.

$^{121}$但是很多人應該比狗大.
你叫狗的名字.
通常叫「Bobby」.
牠很遠地聽到你的聲音.
牠就會跑過來.
搖搖尾巴.
接著你摸牠 「Good Boy」.
如果你叫牠.
牠不回應你.
卻跟著一個給牠薯條吃的陌生人.
並且越走越遠.
這隻狗一定是「Bad Boy」.
這其實就是信仰生命的起點.
是很基本的.
認定耶穌.
認定祂是你的牧人.
認定之後緊緊地跟從祂.
不要再跟其他人了.
還說這邊有人允許你成功.
那邊有人允許你富貴.
你不要再跟了.
認定耶穌.
緊緊地跟著祂.
這是信仰生命的起點.
接著是第十章第七至第十節.
「所以耶穌又對他們說.
我實實在在地告訴你們.
我就是羊的門.
凡在我以先來的.
都是賊是強盜.
羊卻不聽他們.
我就是門.
凡從我進來的必然得救.
並且出入得草吃.
盜賊來無非就是盜竊殺害毀壞.
果乃了是要叫羊得生命.
並且得的更豐盛」.
耶穌再一次說.
我實實在在地告訴你們.
耶穌很認真.

$^{161}$剛才提到門.
現在耶穌開展關於門的討論.
門是什麼?.
門在哪裡?.
耶穌說.
「我就是門」.
明顯.
這道門與豐盛的人生有關.
這是生命的門.
第十章第八節說.
「凡在我以先來的.
都是賊是強盜」.
所謂「以先」.
基本有兩種解釋.
一個很簡單.
就是時間上的「以先」.
按照上下文.
第十章指第九章.
那些在耶穌出道以先.
已經在教導猶太人的法利賽人.
這些是不是耶穌所說的「以先」.
他們就是賊呢?.
但「以先」有第二個意思.
「以先」也可以指.
優次上的排位.
排在前頭的意思.
即是說.
那些將自己擺在比耶穌更優先次序的人.
那些就是「以先」的那些人.
我覺得這個解釋更加有意思.
若果按照這個解釋來說.
可以是指法利賽人.
也可以是指我們任何一個人.
即是說我們可以將自己擺在耶穌「以先」.
如果是這樣的話.
耶穌所說的「以先」就是指我們.
這些人是賊.
是強盜來的.
將自己擺在耶穌「以先」.
這種人對所有人的生命.

$^{201}$都會造成很大的破壞.
經文說.
這些人就是盜賊.
盜賊來無非就是偷竊,殺害,毀壞.
將自己放在耶穌「以先」的人.
他只會利用耶穌來滿足自己.
這些人在群體裡.
他們會引導其他人去注意他.
注意他的成功,他的厲害.
他不會教人去仰望耶穌.
他只會去想自己的心意的成就.
而不會願意神的旨意,成就.
這類人充滿私意.
令人看不到神在那裡.
耶穌說他是「門」.
要從這道門進入才是正途.
經文說.
從門進去的必然得救.
並且出入得草吃.
生命要得到豐盛的供應.
一定要從這道門進入.
才能有真正的養分.
第二就是從正門進入的.
他就叫揚得生命.
並且得得更豐盛.
當然這裡不是指物質財富.
是指生命意義的圓滿.
孔子說:朝聞道,夕死可矣.
朝早領悟到大道.
晚上死去都無憾了.
這些聖人都很明白.
了悟生命的真正意義.
比生命的長短更加重要.
這是信仰生命的第二個層次.
你要學會放下自己.
將耶穌放在最高的位置.
不單單聽從耶穌的聲音.
更加以耶穌為進入圓滿生命的大門.
這是第二個段落.
第三個段落是第十章第十一至第十五節.

$^{241}$耶穌說:我是好牧人.
好牧人為羊捨命.
若果是故宮,他不是牧人.
羊也不是他自己.
他見到狼來了,就撇下羊逃走.
狼抓住羊,趕散羊群.
故宮逃走,因為他就是故宮.
他不顧念羊.
我是好牧人.
我認識我的羊.
我的羊也認識我.
正如父認識我.
我也認識父一樣.
並且我為羊捨命.
耶穌進一步用牧羊人和羊的關係.
來表明合乎正道的生活態度.
這裡有一個對比.
牧人和故宮.
兩種很不同的存活方式.
先說故宮.
故宮的重點當然是賺錢.
他是故宮,當然賺錢.
與人相處.
他只想到利益.
身邊的人對一個故宮來說.
只是土生活的工具.
沒事的時候.
工作看起來很正常.
但是狼來了.
危險來了.
要找人背鍋了.
他當然不會理會身邊的人.
一走了之.
經文說故宮逃走.
因為他就是故宮.
他並不顧念養.
他根本不愛身邊的人.
他來到這個世界.
他的生活只不過是謀生.
虧本的事他不會做.

$^{281}$好牧人.
完全不同.
是另外一種存活的方式.
就正正因為狼來了.
危險來了.
他更加守著崗位.
他知道狼是會來的.
而他明白.
他的存在就是為了這個時刻.
他努力守著崗位.
因為他很想守護他的羊.
他知道這就是他的天職.
狼來了.
並不是他逃走的理由.
反而是他站立在崗位上的理由.
他認識他的羊.
他的羊亦認識他.
一種生命的連繫.
一種愛的連繫.
而這種愛的連繫.
耶穌說到.
深刻到好像聖父和聖子的連繫一樣.
就因為這種愛的連繫.
好牧人是不會放棄他的羊.
他甚至願意為他們捨命.
信仰生命來到這個層次.
就要求我們每一個人.
對自己的存活方式.
有一種深刻的反省.
我們究竟是故宮.
還是好牧人.
我們其實是如何對待身邊的人呢?.
我們如何看待他們呢?.
我們的生活其實追求些甚麼呢?.
我們在人類悠長的歷史裡面.
出現在這個特定的時間點.
究竟我們的天命在哪裡呢?.
時代的艱難.
是我們逃避責任的藉口.
抑或是我們背負十字架的機遇.

$^{321}$人世間有無數.
對生活在謀生的故宮.
人世間有無數故宮.
但是若果多一些好的牧人.
人世間將會有多大的不一樣呢?.
故宮是很多的.
牧羊人是罕有的.
第四個段落.
第十章第十六節.
耶穌說:我另外有羊.
不是在這個圈裡面的.
我必須領牠們來.
牠們也要聽我的聲音.
並且要合成一群.
歸一個牧人了.
耶穌在這裡說.
牠不單單有眼前的羊.
眼前的羊即是眼前的猶太人.
牠不單單有牠們.
牠還有另外一些羊.
當然這就指全世界的非猶太人了.
牠們都是神所愛的羊.
這個觀念突破了.
耶穌當代的傳統猶太教.
限定神的子民的範圍.
傳統的猶太教規定.
要成為神的子民.
你必須先要成為一個猶太人.
猶太人才有資格成為神的子民.
你要守猶太教的一切規矩.
這其實就是一種排外的思想.
自己人就是自己人.
不是自己人就不是神的子民.
耶穌在這裡說.
是一種心靈的擴闊.
要突破人為的界限.
要擁抱界外的人.
現代的世界.
你每天看新聞.
所有的國家.

$^{361}$都是為了守護自己的利益.
越來越排外.
合作越來越難.
耶穌在這裡想的.
是要打破這樣的藩籬.
促成不同社群之間的和解和合一.
我們每個禮拜應信的時候.
都宣讀《仕途信經》.
其實《仕途信經》裡面有一句是這樣的.
「我信聖而公的教會」.
「公」.
Catholic.
所謂「公」其實是一種大公精神.
甚麼叫做大公精神呢?.
「我信聖而公」甚麼叫做「公」呢?.
其實是指一種有容乃大的包容的精神.
包括民族,文化,經濟,地域的大公精神.
即是擁抱不同的民族,文化,經濟,背景.
不同的居住地域的人.
這是同時代的橫面.
其實是包括時間上的大公精神.
即是我們思念過去的人.
關心當下的人.
亦都顧及將來的人.
這是時間直線的一種大公.
耶穌的心靈是一種確而充之的心靈.
突破人為的界限.
使人愛及於未及之處.
信仰不應該使人越來越封閉.
而應該使人越來越打開心門.
向他人開放.
越多這樣的信徒.
基督教的信仰才有將來.
最後第五段.
《楊福音》第十章第十七節第十八節.
耶穌說:我父愛我.
因我將命捨去.
好在取回來.
無人奪我的命.
是我自己捨的.

$^{401}$我有權柄捨去.
亦都有權柄取回來.
這是我從父所授的命令.
他能夠理解的一種天命.
耶穌在此表達他對生命的一種最高領悟.
面對生死.
他不訴諸於命運.
而訴諸於自由.
面對十字架.
這樣不是意味著遇到甚麼厄運.
而是他自己權能下的選擇.
耶穌的一生.
浸淫在天父的愛裡面.
若果活著.
能夠流露天父對人的愛.
他就選擇活著.
若果死去.
能夠讓人驚訝於天父的愛的闖入.
他就選擇死去.
生的慾望.
死的恐懼.
並不主宰耶穌.
主宰耶穌的是天父的心.
是他在天父那裡領受的天命.
他面向天父.
展現真正的自由.
自由的出入生死.
自由的去愛.
自由的活出生命豐盛的內容.
這就是信仰生命的最高體會.
活著.
就是實踐天命.
流露天心.
沒有人強迫他捨命.
也不是因為要履行甚麼律法的規矩而捨命.
捨命是出於他自己的自由的行動.
他不會問這樣公不公平.
他也不會想怎樣可以逃避.
他捨命是因為愛.
他真心去愛.

$^{441}$因為他體會到這就是天父內心的那種想法.
而他的想法與天父的想法完全一致.
活著就是天父的心的流露.
2010年有一齣法語電影上映.
按照法語片名直譯為「人與神」.
英文片名是「Of Gods and Men」.
香港一般正規片名翻譯為「人神之間」.
根據《真人真事改編》.
我在神學院教書,教牧神學或教牧倫理時.
這齣電影是必看的.
是堂上拿來討論的.
話說在北非阿爾及利亞.
有一條回教的村莊.
那裡有一間天主教的修道院.
那裡住了九名法國的修士.
這些修士在那裡修道,宣教,醫病.
與那裡的回教村民和諧共處.
相處得很好.
忽然間有一天.
一群極端的回教武裝分子出現.
殺害那些外國人.
之後整套電影的內容.
整個主調都改變了.
之後的內容就是.
在死亡的威脅下.
這九名修士不斷討論.
究竟去或留呢?.
整套電影其實是一個漫長的討論.
很快可能會死了.
是應該離開還是繼續留在這裡呢?.
當初不少的修士恐懼死亡.
他們的想法就是回到法國.
這也是很正常的.
那些修士對村民說.
「我們是鳥仔,有一天我們會飛走」.
那些村民對他們說.
「甚麼你們是鳥仔?」.
「我們才是鳥仔」.
「你們其實是樹」.
「你們樹立在這裡」.

$^{481}$「我們才能有安生立命的地方」.
電影一幕又一幕地呈現修院的生活.
一幕一幕地呈現修院的生活.
其實就是一幕一幕地展示.
基督信仰對人生命的轉化.
他講到聖誕節的崇拜.
那些修士就在思考.
甚麼叫做基督卑微的臨在?.
甚麼叫做真正的和平?.
他們每一天的修會都有聖餐的環節.
那些修士就一步一步地去明白.
甚麼叫做基督的犧牲?.
甚麼叫做基督在受苦的群體裡面愛的臨在?.
在一個受難節最後的晚餐裡面.
他們這班修士慢慢去體會.
基督的受難和基督在死人裡面復活.
所帶來的希望.
最後所有的修士.
都與基督的生命融成一體.
他們以謙和的生命.
與那些受苦的人同在.
他們用自由選擇.
至生死於道愛.
繼續在那個貧窮的社區.
靜靜地流露天賦的人愛.
電影結束的時候.
恐怖份子襲擊修院.
擄走了七位修士.
有兩位修士因為躲起來逃過一劫.
最後這七位修士都被殺害.
他們陪伴了基督活出豐盛的人生.
普羅曾經說:「我們若活著是為主而活.
若死了是為主而死.
所以我們或活或死總是主的人」.
《羅馬書》第十四章八節.
普羅又說:.
「我已經與基督同釘十字架.
現在活著的不再是我.
那是基督在我裡面活著」.
《解泰書》二章二十節.

$^{521}$原來一個人真的可以或活或死.
都是神的天心的流露.
我希望每個人都能夠體會這種生命.
活出這種生命.
這是豐盛,最豐盛的生命.
這就是永生.
願神賜福你們.
願榮耀歸於聖父,聖子,聖靈的名.
阿們.
\newpage



\section{歌羅西書 3:18-25}
\label{sec:032328}
\textbf{不要惹兒女的氣 (英/粵)}~~~~嚴建平博士
\newline
\newline
連結: \href{http://www.church.com.hk/acms/content.asp?site=cdc\&op=show\&type=product\&code=032328\&layout=sermon}{\texttt{032328}} ~~~~ 語音日期: 2026-06-23
\newline
\newline
\hyperref[sec:032327]{\small{< < < PREV SERMON < < <}}
~
\hyperref[sec:index_chronic]{\small{[返順時目]}}
~
\hyperref[sec:032329]{\small{> > > NEXT SERMON > > >}}
\newline
\newline
歌羅西書 3:18-25
\newline
\begin{longtable}{cl}
\hline
\hline
章節 & 經文 (和合本修訂版)\\
\hline
3:18 & \begin{tabularx}{0.7\textwidth}{X} 你們作妻子的，要順服自己的丈夫，這在主裡面是合宜的。 \end{tabularx} \\ \\ \relax
3:19 & \begin{tabularx}{0.7\textwidth}{X} 你們作丈夫的，要愛你們的妻子，不可虐待她們。 \end{tabularx} \\ \\ \relax
3:20 & \begin{tabularx}{0.7\textwidth}{X} 你們作兒女的，要凡事聽從父母，因為這是主所喜悅的。 \end{tabularx} \\ \\ \relax
3:21 & \begin{tabularx}{0.7\textwidth}{X} 你們作父親的，不要惹兒女生氣，恐怕他們會灰心。 \end{tabularx} \\ \\ \relax
3:22 & \begin{tabularx}{0.7\textwidth}{X} 你們作僕人的，要凡事聽從你們肉身的主人，不要只在眼前服事，像是討人喜歡的，總要心存誠實，因為你們敬畏主。 \end{tabularx} \\ \\ \relax
3:23 & \begin{tabularx}{0.7\textwidth}{X} 你們無論做甚麼，都要從心裡做，像是為主做的，不是為人做的； \end{tabularx} \\ \\ \relax
3:24 & \begin{tabularx}{0.7\textwidth}{X} 因為你們知道，從主那裡必得著基業作為賞賜。你們要服侍的是主基督。 \end{tabularx} \\ \\ \relax
3:25 & \begin{tabularx}{0.7\textwidth}{X} 行不義的人必受不義的報應；主並不偏待人。 \end{tabularx} \\ \\
[1ex]
\hline
\hline
\end{longtable}
$^{1}$各位弟兄姊妹,我好開心可以來到這裡..
上次我應該可以來到這裡,但沒辦法,這次就來補償了..
我們一起去禱告..
上帝,平靜我們的心,讓我們有聆聽的耳朵,讓我們有信服你的心..
奉耶穌名求,阿們..
你們有沒有從來想過為什麼父親節會在母親節之後呢?.
為什麼會是母親節然後到父親節呢?.
有人說,因為母親節的禮物那張單,剛剛在父親節來到的時候,就可以足夠給錢了..
所以就有足夠的錢了..
多好啊,有得可以給錢..
父親有足夠的錢財可以去給,是何等的好呢?.
我們今天所讀的經文,我們專注在國際書三章二十一節..
這一節的經文給我們父親是一個嚴肅的指示..
他說,父親不要惹兒女的氣,恐怕他們失了志氣..
在同一個教導裡面,已不鎖書六章四節也都有出現..
在那裡,保羅警告父親說,不要惹兒女的氣..
意思是不要讓你們的孩子們變成記恨不滿..
否則整件事會變得更加的擴大,就會失了志氣了..
父母是孩子人生裡面第一個的屬靈者..
就是你們了,爸爸媽媽就是第一個屬靈者..
父母也要作為有好的父母的責任..
爸爸和媽媽各自有自己的英俊的本份,也有好的責任..
但父親也有他自己需要做的職責..
作為父親,你不可以忽略或者推卸自己的職責..
不可以將這些責任推卸到其他的母親或者是屬靈的導師身上..
當然,母親和屬靈長輩可以去做這個支持或者去幫助..
在這段經文裡面,是用一個嚴肅警告的方式去鼓勵這些父親..
他說,父親啊,不要惹兒女的氣..
讓我們想一下,究竟怎樣可以不激怒我們自己的子女呢?.
或者不是去令他們心裡面變得很苦澀,不要傷害他們的志氣呢?.
以至於他們不會失去了成長甚至活下去的盼望..
如果我可以的話,我今天會給三個實際的教導給各位的父親..
有很多的小朋友將閻牧師講為一個牧師爺爺..
(主持人:牧師爺爺).
牧師爺爺,他也是啦,牧師爺爺..
我今天像牧師爺爺那樣去告訴所有的父親,這三個的教導..
第一個,.
不要令他們心裡面有苦澀,就是不要輕視你們的孩子..
究竟是怎樣解釋呢?.
輕視他們就等於羞辱上帝的禮物,因為小朋友是神所賜的產業..

$^{41}$是我們上帝所賜的禮物..
你是羞辱上帝的禮物..
同時你也不相信上帝在這班孩子身上有工作..
要記得神確實在我們的兒女身上是在動工的..
如果不贊同上帝將這個孩子交給我們照顧的安排,也是一樣輕視他們..
意思是G-Leg今天是他的,給他鼓掌..
輕視孩子就是羞辱上帝的禮物,不相信上帝在他們身上有工作,也不贊同上帝給我們照顧的安排..
每一個孩子都是獨特的..
每一個都是被上帝特別的關愛..
既然孩子被託付在我們的照顧之下,我們更加應該體重重視他們..
我們越重視他們,就越能將盼望灌輸在這班孩子的生命裡面..
這班孩子就會有很強的品德和有上帝的榮美..
我將這個故事講述給大家聽..
這個故事講述在德國有一位教十至十一歲的男子的一個老師..
在每次上課之前,他會將他的高帽脫下..
然後鞠躬致意,要尊重..
就是尊重所有的孩子..
有人就問這個老師為什麼要這樣做呢?.
他的答案不單單很有幫助,也很有說服力..
他說不知道這班孩子將來會變成怎樣的人..
而當中其中一個是馬丁路德..
他之後成為了一個很出名的大學教授,也是一個被尊敬的教會領袖..
我們從來都不可以只看待我們的孩子是現在這個狀況..
要看他們預期之後幾年究竟會變成怎樣的呢?.
讓我再講另一個故事,是講關於一個領袖叫菲利普·麥蘭頓..
他和六歲的小朋友講話..
他和這班小朋友講話是會用很禮貌去問案的..
他和我們講教授們..
醫生們..
律師們..
老師們..
作家們..
農夫們..
藝術家們..
家庭主婦..
還有畫家們..
父母,畫家..
他和他們講話是好像很有成就的大人一樣的態度..
因為他沒有輕看他們,他看他們很貴重..
他很愛這班孩子,因為他知道他們是尊貴有上帝的形象..

$^{81}$而這個老師更加將價值觀好好地教導這班孩子..
在這樣做的時候,這班孩子就不會失去志氣,失去希望,或者失去方向..
當我七歲的時候,我很沒有安全感..
我完全是沒有自信,很沒有信心..
我是有很多的掙扎在裡面..
每一晚我都不能夠去進睡..
因為我的父親在我早年已經過世了..
我是很害怕..
每一晚我都要抱著我母親的腳去進睡..
在教會裡面的執事,其後也成為了我的屬靈的導師..
他真的將我變成一個兒子..
他花很多的時間去陪伴我,去鼓勵我..
他跟我說:你有一天會在公眾場合跟幾千人去說話..
他真的看得到我以後,很多年之後會變成這樣的樣子..
而我在這幾年裡面真的跟很多很多的人去說話..
我當時很害怕,很沒有自信,甚至有自殺傾向,.
都依然可以成為這一刻的牧者去講這個福音,講這個信息..
今時今日不同了..
你不能說我沒有進步..
我沒有任何的成就..
可能我真的沒有覺得有很多的成就,.
但是因為透過那個屬靈的父親鼓勵我之後,.
我現在可以成為現在的我..
有很多的屬靈父親是很愛我們的..
第二點,.
怎樣可以避免他們不再激怒我們的子女呢?.
就是不要輕視他們的感受..
有時有些的父親真的非常理性,.
那個理性去到他們自己掙扎都不知道,沒有感覺的..
我們作為成年人有時候真的會忘記了原來這班小朋友是有感受,.
也很快就可能會不理會他們..
很多時候我們很容易去忘記其他東西,但就忽略了我們小朋友的感覺和傷痛..
青少年或者一些成年的孩子經常都和我抱怨..
我多麼希望我的父母能夠聽我說多一點這些疑問,.
因為他們可能是很重要的..
是的,可能他們真的很重要..
有時候我覺得我很被壓垮,很有挫敗感..
因為我的父親經常都將這些質問變成誤會,覺得我在挑戰他的權威..
其實一點都不是,我只是希望我們的話題可以再深入一點,所以我才去問我的父親..
希望可以多一點,講多一點,談多一點..

$^{121}$就算是意見不同的時候,都可能可以講多一點,談多一點..
我不是要去挑戰男士或者爸爸父親的權威或者尊嚴..
是的,最厲害的就是他了,今天..
我只是希望他能夠聽我說多一點,我們可以一起學習,一起成長..
是的,父親,我們可以在孩子身上學到一些東西..
就算只是一個心對心的交談,真是心連心的交談..
一定要有一個更加深的連結在父親和孩子之間..
對方都要知道對方的感受是怎樣..
父親和孩子之間的關係要更加有聯繫..
父親和孩子之間的連結要變得更深..
父親和孩子之間的聯繫要變得更加深..
否則就會變成很多束縛..
很多束縛到一個地步你不能夠分享..
那種束縛是你無法跟對方分享的..
如果沒有這個連結,就更加會有束縛..
當有束縛的時候,就不只是束縛,是會有轟炸,好像感覺被壓碎了..
有束縛,有被壓碎和被轟炸..
作為父親的我,我的兒子在小時候成長的時候,我就試著不去忽視他任何的感受..
因為我作為一個教授,我可以很理性..
有時候真的完全沒有情感的表達..
理性和感性是要在一起的..
情傷和智商都要在一起的..
兩種都是要合在一起的..
總之那種情感和智商一定要全部在一起的..
有時候我們真的聽不夠..
我們沒辦法聽得更多,所以我們也不明白對方..
所以有時候我們真的沒辦法有心連心的那種交談..
我作為父親在成長中幫助我的兒子,真的由一個不理性的抒發,變成一個合理的宣洩..
又再講一個故事給大家聽..
當我的兒子是八歲的時候,他在我的後座裡突然大叫..
又哭又大叫,我就問他,什麼事呢?.
他說,有人偷了我的紅筆..
紅筆?你沒有紅筆..
他說,有啊,我在學校的班房的地上撿到的..
就是那支紅筆,我在地上撿到的那支紅筆不見了..
所以他就哭啊哭啊哭,又大叫..
作為父親,我知道這是一個不合理的哭..
對於一個很理性的人來說,我快要罵他一頓了..
他說,為什麼要哭呢?.
又不是你的..

$^{161}$你哭什麼呢?.
又不是你的,為什麼要哭呢?.
哭什麼呢?.
但是那一刻,我控制住自己..
我知道縱使那個原因可能是錯的,.
我也嘗試去理解我的兒子為什麼有這個傷痛和失望..
我的兒子當時不夠成熟去知道究竟什麼是合理和不合理的淚水..
究竟為什麼他不應該哭,因為那些東西根本不屬於他..
他就要教他怎樣才有合理抒發的淚水..
我就慢慢嘗試引導他由不合理的淚水變成合理抒發的情緒..
我就和我兒子說,Hansel,這個是他的名字..
我就說,我會幫你多買兩支紅筆..
這兩支紅筆是比你在地上撿到被人偷走的更加漂亮..
他就很高興了..
他很高興第二天拿了一支紅色的紅筆上學..
他就告訴所有的同學,很高興地告訴大家,他有一支新的紅色的紅筆是爸爸買給他的..
但很不幸,那支紅筆也被人偷走..
這次他哭得更淒涼,更加慘,更加控制不了自己..
但你知道我怎樣做嗎?.
由得他哭,繼續哭.因為這次是你自己的紅筆被人偷走,所以讓他哭..
就繼續讓他哭..
當他沒有再哭之後,我就去解釋給他聽..
那你有沒有發現第一次和第二次哭的分別呢?.
第一次是一個不合理淚水的那種抒發..
但作為父親的我,我依然確認他當時有那種失望和不開心..
但我嘗試引導他不再有一個不合理,而慢慢變成一個合理的情況..
他就很開心地解釋給他聽這兩者的分別..
幾年之後,我見到我兒子有所改變..
他就開始有合適的淚水抒發,無論他是很傷心還是很憤怒的時候..
父親,你們一定要好好有耐心地對待你的孩子..
一定要慢慢引導他們出來,由不合理變成合理..
父親,容許我去講這句話..
在我們的積分裡面,我們必須要有情緒表達的榜樣..
父親可能有很多很多的壓力在裡面,.
但我們一定要學習怎樣去控制這些情緒..
我們究竟要怎樣去應付自己的情緒呢?.
我問他們,為甚麼他們不回家吃聖誕晚餐呢?.
他們說,因為他們有很多的專注在學校..
有一個老虎在家裡..
他說,誰是老虎?.

$^{201}$他說,母親,好像老虎一樣..
他說,但是老虎之外,還有就是獅子..
家裡有母老虎,也有獅子..
我說,你這是甚麼意思?.
他說,因為我父親在家裡經常有不合理的怒氣..
所以父親,你們在家裡,你們在你的積分裡面,.
一定要有表達情感上的榜樣..
如果你容讓你自己這些無理的怒氣,.
或者是一些不開心的情緒,.
真的完全發洩出來的時候,.
你可能會知道你的家裡會變成怎樣..
其他人可能會失去志氣,失去了不同的盼望..
第三點就是我們要花時間和我們的孩子一起..
我們不想錯過任何在他們成長的機會裡面的任何一個時刻..
小朋友最需要的其實就是我們的陪伴和時間..
我們永遠都給不夠這些時間他們..
無論是質量還是數量..
但是只有我們花更多的時間,.
那個關係才有質素,才能夠提升..
如果沒有數量,就會沒有質量..
很多小朋友在成長的時候,都會有他們的掙扎..
他們一直在成長的過程中,去嘗試了解究竟中間在發生什麼事..
他們問這個問題,.
當我們在學習上,在教育上,真的沒有辦法去符合的時候,究竟我們還有沒有盼望呢?.
當我沒有辦法去符合父母的期望的時候,究竟我們還有沒有盼望呢?.
如果同輩和靈人給我們壓力的時候,我們有沒有辦法去面對或超越他們呢?.
孩子只需要我們一直在那裡陪伴..
要一起黏著..
一起哭..
一起笑..
一起去打枕頭戰..
可能是說一些無聊的笑話..
有時候那些無聊的笑話才會和我們一起分享..
每日瑣碎平凡的事,能夠分享的時候,才能夠彼此了解對方..
而這些共享一起分享的回憶,才能夠成為喜樂和鼓勵的傳言..
有一個很聰明的大學生,在所有的科目他都很聰明,但是就只有一科不合格,就是物理科..
她是一個很大膽的女孩,變成一個很消沉的女孩..
她說我不知道怎樣面對..
她說在那個時刻,她父親打電話給她..
她爸爸由多倫多開六個小時的車去到麥嬌大學找她..

$^{241}$當她見到他時,她就瘋了..
真的跑過去..
在她的肩膀上,她就哭了..
她說我的父親一句話都沒有說..
爸爸只是好好的抱著她一起哭..
他沒有改變我任何東西..
沒有有任何的怒氣或者責罵我,因為我不合格..
她說爸爸的出現,當刻陪伴我一起,就是最大的醫治..
她說那些回憶成為我可以面對生活壓力的一個很好的動力..
現在她已經完成了她的學位,也成為一個藥劑師..
現在她有五個小朋友要照顧..
超級有效率,很會生,生三五個..
很會生,一生生五個..
.
閻牧師說希望大家可以生多一點..
.
我的母親有十個..
.
我母親真的照顧了所有十個,每一個都被照顧..
.
當我是七歲的時候,我的父親已經離世,但有十個..
.
但我的母親除了母親的職責之外,她還要擔當父親的職責..
.
可能有些在座的弟兄姊妹朋友,你的父親未必在這裡..
.
可能母親真的要做完所有的職責..
.
我們真的需要明白背後的故事,背後他們的苦衷..
.
另外一個大學生又有故事了..
.
他說我的女朋友離開了我..
.
他說她和其他同學去拍拖..
.
就把我氣死了..
.
令我很憤怒..
.

$^{281}$總之我們很憤怒,很傷心..
.
他說我不想再在這個世界上活..
.
有沒有生活的意義呢?.
.
他現在的掙扎就是究竟應該怎樣去面對這件事呢?.
.
怎樣去解釋呢?.
.
父親就應該在這個時候,進來這個位置去幫忙去解釋..
.
小朋友其實不是壞的..
.
他們只是在掙扎..
.
就算好像父母一樣,其實他們不是很差的..
.
他們依然是好像小朋友一樣,都是掙扎..
.
我掙扎去成為一個父親..
.
我有ADHD專注力的影響..
.
我的兒子也有ADHD,有其他的問題..
.
那你可以怎樣呢?.
.
那你就要花很多時間,繼續花時間去花時間..
.
那個,那個那麼多前書..
.
十三章..
.
要繼續去閱讀,去讀給自己聽..
.
去花時間..
.
父親是很重要的職責,就是要陪伴著你的孩子一起..
.

$^{321}$那些繼續有成長的,有年紀開始慢慢增長的..
.
我們也都要好好地去看著這一班的..
.
如果你的父親是在其他的國家,不是加拿大的,打電話給他們..
.
或者飛過去,跟他們好好地相處..
.
時間都跑到哪裡去了?.
.
眨一下眼,我的兒子現在已經28歲了..
.
時間繼續在推緊我們..
.
作為一個屬靈的父親,很多人會來找我去幫忙..
.
是一個很好的職位,去幫助其他人..
.
他們真的,他們所有的那些,都是call,去挖我們..
.
屬靈的爸爸,屬靈的父親..
.
一開始是一個覺得,好一個,真是害怕的呼召..
.
可能你真的會激怒他們的..
.
讓我去講一下關於Pastor Ada..
.
他在香港,現在應該在看..
.
當我第一次見到Pastor Ada的時候,我知道一定有些事情會發生..
.
他說,Pastor Ada,你知道我的眼睛是大過你的..
.
他的樣子都顯露出,他是很擔心..
.
他就說,我就會祈禱了..
.
他說,我就會多些花時間在這個女孩身上..
.

$^{361}$我的出現,我和她一起陪伴..
.
繼續好好地去培養,培育她..
.
然後教導她..
.
這是一個過程..
.
你必須在那裡,閉著嘴..
.
在那裡聽..
.
你會感受到他們的傷痛..
.
然後你會成為一個個人的那種培育的個人..
.
然後我就教導她,怎樣去講道..
.
我去改她的講章..
.
但我是很友善,很溫柔地對待她..
.
所以你見到為什麼Pastor Ada是這麼好的..
.
我講這個東西,我也要你們來誇我一下..
.
我今天講給大家聽,希望大家能夠誇我..
.
尤其我做老師也很慘..
.
個個人都疼愛學生,不疼愛老師..
.
我們有牧者,有執事,有領袖,都是作為父親或男士的..
.
很好的對待他們..
.
他們真的為我們好,花很多的時間..
.
時間..
.

$^{401}$陪伴..
.
陪伴..
.
培育..
.
教導..
.
這三樣是你如果懂得做的話,你會有很多的追蹤者..
.
如果你是懂得去做這三樣東西的話,你就會知道怎樣去做一個好的父親..
.
所以我們一起去做咖啡,去花時間..
.
謝謝..
.
謝謝,謝謝,太好了..
.
我還想說我們三個一起去做咖啡..
.
你們付錢,我祈禱..
(笑聲).
\newpage



\section{瑪拉基書 3:6-12}
\label{sec:032329}
\textbf{福浸眼眉的驚喜}~~~~石偉漢牧師
\newline
\newline
連結: \href{http://www.church.com.hk/acms/content.asp?site=cdc\&op=show\&type=product\&code=032329\&layout=sermon}{\texttt{032329}} ~~~~ 語音日期: 2026-06-14
\newline
\newline
\hyperref[sec:032328]{\small{< < < PREV SERMON < < <}}
~
\hyperref[sec:index_chronic]{\small{[返順時目]}}
~
\hyperref[sec:032330]{\small{> > > NEXT SERMON > > >}}
\newline
\newline
瑪拉基書 3:6-12
\newline
\begin{longtable}{cl}
\hline
\hline
章節 & 經文 (和合本修訂版)\\
\hline
3:6 & \begin{tabularx}{0.7\textwidth}{X} 「我－耶和華是不改變的；所以，雅各的子孫啊，你們不致滅亡。 \end{tabularx} \\ \\ \relax
3:7 & \begin{tabularx}{0.7\textwidth}{X} 從你們祖先的日子以來，你們就偏離我的律例而不遵守。現在你們要轉向我，我就轉向你們。這是萬軍之耶和華說的。你們卻說：『我們如何轉向呢？』 \end{tabularx} \\ \\ \relax
3:8 & \begin{tabularx}{0.7\textwidth}{X} 人豈可搶奪神呢？你們竟搶奪我！你們卻說：『我們在何事上搶奪你呢？』其實就是在你們當納的十分之一奉獻和當獻的供物上。 \end{tabularx} \\ \\ \relax
3:9 & \begin{tabularx}{0.7\textwidth}{X} 因你們全國上下都搶奪我的供物，詛咒就臨到你們身上。 \end{tabularx} \\ \\ \relax
3:10 & \begin{tabularx}{0.7\textwidth}{X} 你們要將當納的十分之一全然送入倉庫，使我家有糧，以此試試我，是否為你們敞開天上的窗戶，傾福與你們，甚至無處可容。這是萬軍之耶和華說的。 \end{tabularx} \\ \\ \relax
3:11 & \begin{tabularx}{0.7\textwidth}{X} 我必為你們斥責蝗蟲，不容牠毀壞你們的土產。你們田間的葡萄樹，果實未熟以先也不會掉落。這是萬軍之耶和華說的。 \end{tabularx} \\ \\ \relax
3:12 & \begin{tabularx}{0.7\textwidth}{X} 萬國必稱你們為有福的，因你們必成為喜樂之地。這是萬軍之耶和華說的。」 \end{tabularx} \\ \\
[1ex]
\hline
\hline
\end{longtable}
$^{1}$好 弟兄姊妹主 內平安.
願大家都有開心的日子.
特別是父親節的日子.
無論你是父親的.
或者你不是父親的.
但是我們盼望每天都開開心心的.
這樣去過我們的每一天.
好 我們一起先有個祈禱.
讓我們一起可以在你的面前.
我們去敬拜天父你.
因為你就是我們最好的父親.
讓今天父親節.
你就得到最好的榮耀.
也讓你自己先喜悅.
我們在你面前的敬拜.
也讓父親你自己親自對我們說話.
讓我們聽到父親你的說話.
我們就成為一個乖孩子一樣.
我們是遵從主你的吩咐.
求你帶領我們.
聽我們祈禱奉耶穌名字而求.
阿們.
今天和大家分享的題目.
叫做《覆浸眼眉的驚喜》.
就是一個這樣的畫面.
在未去到裡面的內容之先.
首先和大家說一個故事.
大家都有喝中國茶.
有喝過中國茶.
中國茶裡面.
我們有時候也有一些叫茶道.
就是茶友茶道.
其中我們的茶道裡面.
有一個叫寇子禮.
不知道大家知不知道這個寇子禮.
為什麼會有這個寇子禮這件事發生呢.
原來按一些資料顯示.
其實這個寇子禮是和乾隆王有關的.
原來話說乾隆有一次微服出巡.
微服出巡就要出到裡面.

$^{41}$去了解一下民間的疾苦或者狀況.
於是就去到民間裡面其中有一次.
就去茶樓喝茶吃東西.
當時他就突然興致勃勃.
就幫他的侍從斟茶.
皇帝為你斟茶.
怎麼辦呢.
皇帝為你斟茶.
平時皇帝在皇宮上賜你要怎樣.
要謝主隆恩.
現在在光天化日之下.
全世界都看到你的時候.
如果他們謝主隆恩怎麼辦.
微服出巡就要露餡.
被人發現了微服出巡是不行的.
但是皇帝斟茶你不禮貌.
也不行的.
所以當時其中有一個智者侍從.
他就做了一個動作.
就是我們今天的購茶禮.
去表達了皇帝為你斟茶.
所以問題來了.
你知道我喜歡講道.
團體都要大家用一下腦.
問題來了.
以下有三種的購茶方式.
請問你哪一種是乾隆手下.
他的購茶方式呢.
第一個就是左邊.
用一隻手指這樣購茶.
第二個就是兩隻手指購茶.
第三個就是這樣手指購茶.
請問你這幾個侍從.
皇帝幫你斟茶的時候.
要用哪一種的方式去購呢.
第一種的舉手.
第二種的.
第三個的.
請問廣哥為什麼是第三個呢.
你會覺得是撞的.

$^{81}$謝謝你承認是撞的.
有沒有人猜到.
有些原因你會支持這種想法呢.
因為當時的大臣要五體投地.
因為皇帝平時在皇宮要五體投地.
暗地裡五體投地就用五隻手指.
向皇帝叩頭.
所以規矩是要叩三下.
表示對皇帝的叩頭.
所以今時今日我們的購茶裡.
都是這樣的.
如果你是晚輩.
向著你的長輩.
如果長輩斟茶給你.
就記得這樣輕輕叩三下.
你不要大大力.
質問人才是大大力.
輕輕叩三下就代表你是晚輩.
向長輩多謝他為你斟茶.
如果大家平輩.
平輩是怎樣.
就用兩隻手指.
叩三下.
平輩叩三下.
為什麼這樣呢.
原來兩隻手指代表.
像我們打功夫那樣.
抱權利.
兩隻手這樣.
表示一種恭敬.
平輩互相切磋的態度.
兩隻手.
所以如果晚輩斟茶給長輩.
我斟茶給你們呢.
你們要怎樣回應我.
就用一隻手指.
一隻手指.
叩多少下.
錯了.
你是長輩.

$^{121}$通常晚輩送禮物給你.
這樣就行了.
叩一下.
因為你是長輩.
表示輕輕點頭.
這樣就行了.
點一下頭.
所以叩一下就行了.
如果你很欣賞晚輩.
都可以拍三下.
你很欣賞晚輩就拍三下.
就是這樣.
所以原來.
喝茶有這些禮貌.
不過不知道你們有沒有發現.
我剛剛和弟兄姊妹去喝茶.
發現無論我叩多少下.
他們都繼續斟茶給我.
有時候我們理解.
叩幾下就代表夠了.
不要再斟了.
有時候我們這樣理解.
但有些人不斷斟.
喝多兩口又斟.
有時候喝到膀胱都有病.
但代表什麼呢.
你的杯子.
從來都不會不夠茶喝.
今天我們思考這個題目.
都是在說.
茶有茶道.
我們的茶杯.
都會滿舍.
因為有人幫我們斟茶.
而神有神道.
神的道是什麼呢.
讓我們的福杯滿日.
是因為什麼呢.
三章六字說.
因為我耶和華是不改變的.

$^{161}$所以雅各的子孫.
你們不自滅.
因為原來神的心意.
就是愛世人.
神愛世人.
甚至將他的獨生子.
賜給他們.
接著呢.
把自滅亡.
凡得永生.
證明你們有運輸.
一生人有幾句經文.
在心裡一生有幸.
不自滅亡.
凡得永生.
所以你記得前面三章五節.
上次我們說到.
三章五節裡說的七宗罪.
很多以色列人背棄上帝.
那些罪.
神說有一天要審判.
但第六節就說到.
一個很重要的問題.
雖然神會審判你們的罪.
但其實神的本意.
是不自滅亡.
所以神對我們的愛.
都是不改變的.
所以神就告訴我們.
其實神想在我們的生命裡.
是怎樣呢.
就是帶來祝福.
神其實本來就是愛我們.
本來就是想祝福我們.
但當你要經歷這種祝福的時候.
有時候都需要做一些事情.
去體驗,去經驗.
才能夠蒙受這種祝福.
原來我們的信仰就是這樣.
上帝就是要我們不自滅亡.

$^{201}$是凡得永生.
凡得上帝的愛.
凡得上帝的祝福.
不過問題是.
你願不願意去接受上帝的祝福呢.
上帝要祝福你.
但你要做一些事情.
或者這樣說.
你要不做一些事情.
你要調整你的生命.
才能夠承受到這種.
福杯滿日的狀況.
我們今天會看到.
《馬拉基書》後面的部分.
我們會分開兩個分題.
去看這個部分.
今天你們會發現.
我是有點不同的.
有什麼不同呢.
就是說慢了.
減輕了碧姐的壓力.
希望可以控制得舒服一點.
不用翻譯這麼辛苦.
我們先看第一個分題.
七至九節.
我想請大家讀一讀.
因為剛才主席讀了.
你們沒有機會讀.
他不可以搶奪你們.
讀聖經的福分.
我請大家讀一讀.
七至九節開始.
你們卻說.
我們在何事上搶奪你呢.
其實就是在你們.
當立的十分之一奉獻和.
當獻的供物上.
因你們全國上下.
都搶奪我的供物.
詛咒就臨到你們身上.

$^{241}$七至九節裡面.
我給大家一個分題.
叫做萬君之夜和所要求的獻賞.
原來神很想我們.
在蒙受福氣之前.
我們去學習去獻賞.
去擺一些東西給神.
擺些什麼給神呢.
我嘗試在這裡分開了.
兩個分題的重點.
第一個是什麼呢.
第七節.
就是要全人的獻賞.
全人的獻賞.
是整個人去獻賞給神.
怎樣獻賞呢.
這裡說.
他們要求這班.
以色列人要轉向神.
怎樣轉向呢.
如何轉向呢.
裡面是要求.
轉向的一種變化.
今天我們.
我剛才說了很多世界盃.
世界盃裡面其實有很多.
踢足球的技術.
其中一種技術.
美心很喜歡的.
就是插花.
足球有一種叫插花的動作.
即是走來走去.
插花.
不過你不明白不要緊.
但還有一種動作叫.
走左望右.
知不知道什麼叫走左望右.
走左邊.
但其實是望右邊.
引人以為是望那邊.

$^{281}$其實是望這邊.
就是騙了一個人.
整個身體就過了那個人.
原來有時候我們.
都會走左望右.
好像心在這邊.
但其實人已經在這邊.
這段經文告訴我們.
我們要轉向的時候.
就是什麼呢.
當你決定要轉向神的時候.
轉向一種悔改.
轉向神的時候.
不只是你想轉向.
而你的身體是不配合.
因為身體是很誠實的.
你的身體都要配合.
一同轉向.
這樣才能夠配合.
你的心裡想的動作.
所以神就吩咐.
當你要去決定.
要去獻上給神的時候.
首先第一個就是.
你要全人獻上.
將你整個人.
不只是想一些東西.
去給神.
而是整個人.
都真的擺上去給神.
所以在羅馬書裡面.
都很清楚.
我們經常讀到的經文.
就是說.
所以弟兄們.
上帝的慈悲勸你們.
將身體獻上當作活祭.
是聖潔的.
是上帝所希願的.
你們如此事奉乃是理所當然的.

$^{321}$當你要成為一個.
祭物獻給神的時候.
神期望你獻上的.
是一個將身體獻上.
當作活祭.
你不可以說.
我的心我要獻給上帝.
但身體就繼續留在祭壇以下.
是不可以的.
所以在這段聖經裡面.
說的是期待著.
那種以色列人的轉向.
就是180度的改變方向.
整個人向著上帝.
但其實.
當時的以色列人有什麼反應呢?.
其實當時的以色列人.
他們就是覺得很委屈.
他們心裡面就說.
我們不是已經留在耶路撒冷了嗎?.
我們千里迢迢.
由被擄之地回到耶路撒冷.
我們今天就特意回到耶路撒冷嗎?.
神大力要求有敬拜的生活.
我們不是有一些宗教儀式.
去配合上帝的期望嗎?.
你還要我們轉什麼方向?.
你要我做的功夫.
我都做足了.
以色列人心裡面這樣想.
所以你看到我們上次說的.
上樑不正下樑歪.
上樑和下樑.
都做了很多表面的功夫.
但裡面沒有實質的生命的轉向.
整個人轉向神裡面.
所以神就期望他們.
在偏離的路裡面.
再一次回來.
不要走了忘憂.

$^{361}$要全心全意向神奔跑.
神就期望人.
今天上帝的子民百姓都是這樣.
我們不只是在某些場合.
某些地方去說.
我很愛神.
不是這樣.
神期望我們整個人.
在我們生命裡面.
活出那種全然愛神的表現.
是時候又要說世界盃了.
世界盃裡面.
在以往很多不同的賽事.
今年有些不同的賽制.
有很多隊.
四十多隊的隊伍.
但今年很多花邊的消息裡面.
其中一個新聞是說什麼呢.
就是說美國隊.
美國隊已經成為了第二支.
必然篤定出勝的.
打入三十二強的隊伍.
其中原來美國隊裡面.
在整個球賽裡面.
做了一樣東西.
不知道你有沒有看新聞.
新聞裡面就這樣說.
6月12號.
在加州那場比賽裡面.
D組賽事裡面.
美國隊擊敗了烏拉圭的球隊.
但他們打勝了這場比賽之後.
他們做的是什麼呢.
他們就是圍埋一起祈禱.
圍埋一起祈禱.
他們在整個的賽事裡面.
很清楚就是.
他們打勝這場比賽.
他們看到上帝對他們的帶領.
他們願意將這件賽事交回給神.

$^{401}$他們勇敢地展示了.
他們是信耶穌的.
他們勇敢地在美國的國民面前.
顯示了他們和其他人不一樣.
他們是一齊玩意祈禱的人.
當中帶領著祈禱的就是後裔馬克麥肯齊.
他就是這樣帶領著球員去祈禱.
事實上他們不只是在.
球賽比賽結束的時候這樣祈禱.
在球賽比賽之前.
他們已經在一起禱告.
當中其中在這隊的隊伍裡面.
除了後裔帶領著這班隊員一起祈禱.
原來還有一個邊封.
就是克里斯蒂安·普利西奇.
他說原來他自己不只是在球場上.
一起參與祈禱.
原來他自己在平時的生活裡面.
他常常在自己的一些公開的頻道.
一些Instagram裡面.
都分享聖經和和別人談論.
他自己生活裡面.
他會有一個叫Bible Time.
聖經時間.
就是說他不只是在人面前.
或者在球場上一個樣子.
回到教會另一個樣子.
不是.
他願意整個人的生活習慣.
都是這樣去顯示出.
他有一種讀經親近神的生活習慣.
是整個人的.
他曾經在2021年這樣說過.
他說我必須一次又一次地證明自己.
但一如既往.
我向上帝祈禱.
祂給了我力量.
有了這份力量.
真的沒有什麼能阻止我了.
他自己就是這樣去深切的祈禱.

$^{441}$事實上在整個世界盃的賽事裡面.
記錄著不單止美國隊.
是這樣在賽後裡面祈禱.
原來.
報導裡面發現原來有很多的隊伍.
一些基督教的國家或者團體.
他們都是這樣.
另外一場的賽事.
就是這一場的賽事.
這場賽事裡面.
也是在說他們.
Prayer before the whole world.
就是說.
他們在賽事裡面.
向世界裡面一起去祈禱.
而且.
說的是巴拿馬.
德國隊.
應該是德國隊的隊員.
他們就是在這個賽事裡面.
就是這樣圍著.
圍著這班富拉索的球員.
在賽後裡面一起祈禱.
這幅圖畫特別的地方是什麼呢?.
不只是德國球員打完球就祈禱.
而他們圍著對方富拉索的球員.
圍著圈子一起祈禱.
表明了他們的信仰.
在別人面前.
而且是跨越了國家.
雖然這個比賽是為國家而戰.
但是在比賽之後.
他們仍然願意.
大家認為.
認定大家是弟兄.
其中帶領著這個祈禱.
其中一個球員.
就是剛才我們看到右邊.
這位球員.
他就說什麼呢?.

$^{481}$他說在這個比賽裡面.
雖然我們是對手.
但是在比賽完了之後.
其實我們都是基督徒.
所以我們都是兄弟.
所以我們就會一起去祈禱.
因為我們相信.
Jesus is glorified.
耶穌基督是要被榮耀的.
要被榮耀.
甚至在這位球員裡面.
他說過一句話.
他說在歐洲.
做基督徒最困難的是什麼呢?.
歐洲做基督徒.
不是一些回教國家.
要被逼迫.
不過他說在歐洲.
做基督徒最困難的.
就是什麼呢?.
就是隱性的壓力.
隱性的壓力.
什麼叫隱性的壓力呢?.
就是說.
有一種想法.
就是民間的想法.
就是信仰可以存在.
但最好不要太認真.
宗教可以保留.
但最好留在私人領域裡面.
即是你自己躲在一角.
做認認真真.
怎樣都好.
但公開場合.
最好不要這樣.
這個就是他自己.
公開表示的見證.
就是他願意.
站在別人面前.
去表明了.

$^{521}$他要成為一個活祭一樣.
去獻給神.
站在萬人見到的台前.
去顯明了.
他是甘心為主拜上.
所以第一個要求就是.
神要求的獻上.
就是一種全人.
整個生命的生活裡面.
那種獻上.
而第二種獻上是什麼呢?.
就是在後面的第八第九節.
我們見到.
就是一種財物上的獻上.
這裡特別說.
一種狀況就是.
這群人.
第九節說.
你們全國上下.
都搶奪我的貢物.
就是說.
我們獻上禮物給神.
其實是應份的.
不過這裡.
正好說一個.
很特別的處境.
這裡是全國上下.
都一起去搶奪.
屬於上帝的貢物.
搶奪這個字.
在原文裡面的意思就是.
一種以暴力.
強行奪取的狀況.
也有一種.
偷走的意思.
知道什麼叫偷東西嗎?.
偷東西是什麼?.
那個東西不是你的.
你拿了就叫偷.
如果那個東西是我的.

$^{561}$我想拿回.
那是不是偷?.
不是.
現在神就說.
我想拿回自己的東西.
因為這個是我的貢物.
神說這些東西是我的.
我想拿回自己的東西.
神不是在你身上拿東西.
神是想拿回.
一些屬於他自己的東西.
但是他反而遭到什麼呢?.
那群以色列人.
就是他們上下一心.
去搶奪了屬於上帝的東西.
其實上帝要求我們的獻賞.
這裡特別說.
一些財物上的獻賞.
其實這些本來就是神自己的.
所以第八節不斷強調.
人豈可搶奪上帝呢?.
豈可搶奪我呢?.
就是豈可搶奪上帝.
當他們還說.
我什麼時候.
有搶奪過上帝的東西呢?.
那些是我的.
不是你的.
上帝說不是的.
這些當立的十分之一的奉獻和貢物.
都是我的.
這裡正正好說.
一種狀況就是.
上帝不單止要用人獻賞.
我們也用一些物質的表現.
去顯明我們願意為主去擺上.
事實上我們見到.
在聖經裡提過一些奉獻的功課.
希望我們可以了解一下.
其中關於奉獻的功課.

$^{601}$其中可以簡單來說分為三大類.
第一類是什麼呢?.
就是對利未人的奉獻.
就是第一個十一奉獻的概念.
就是在利未記裡說到.
要對利未人.
就是一群他們分別出來.
專心服侍神的人.
十二個支派.
其中有十一個支派.
都可以有土地.
有自己的農作物.
有自己的產業.
有自己的牛牛羊羊.
有自己的生意.
什麼都可以.
但其中有一個支派.
他不可以有自己的土地.
當你剛剛最近在讀文書的時候.
你發現不斷強調什麼呢?.
利未人是分不了土地的.
他們一直都沒有分土地.
他們甚至乎就是.
他們代表著一群以色列人.
成為了代表他們的一份子.
去將他們的藍丁獻上給神.
專心服侍神.
以致十一個支派不用派人.
來在聖殿或會幕裡服侍神.
所以對利未人要奉獻.
因為要照顧這群.
代表他們去服侍神的人.
第二個十一奉獻.
就是在靜夏裡.
要獻上一種殿捐或會幕.
因為你知道那裡說.
在會幕裡或聖殿的建立裡.
都有很多燈油火蠟.
很多維修 很多物資.
很多物料.

$^{641}$你要換油換燈.
很多東西.
以色列人都要對聖殿裡.
或會幕裡都要獻上.
以致能夠維持物資上的需要.
第三類的十一.
就是在有需要的人身上.
獻上一些詞彙.
一些幫助.
特別幫助窮困的孤兒寡婦.
讓他們在生命裡.
能夠得到一些愛心的關懷.
所以其實聖經裡說.
好幾個不同的十一奉獻.
但無論說哪一個都好.
重點就是說.
這些本來就是神.
本來就應該物歸原主.
就要去給神自己.
所以其實神就是要求我們.
去學習將一些屬神的東西.
給神.
物歸原主.
就是這麼簡單.
一個這樣的功課.
但很多時候我們知道.
實況是什麼呢.
在這裡說.
Nice people rob God.
很多時候我們都會.
在生活裡不知不覺.
做好人之餘.
我們會搶奪上帝.
我們掛著一些很好聽的理由.
去搶奪上帝的東西.
而這些理由聽起來都很舒服.
在我信主年輕的日子裡.
我都聽過一個叫.
十個蘋果的故事.
大家聽過十個蘋果的故事嗎.

$^{681}$沒聽過吧.
十個蘋果的故事裡.
就是說有一個人.
他擁有了十個蘋果.
上帝祝福他.
讓他擁有十個蘋果.
上帝跟他說.
你獻一個給我就行了.
而這個人當時說好啊.
接著他望一望這十個蘋果.
他想一想.
我要租地方住.
我都要花蘋果去換這個地方.
於是他就想一想.
他將三個蘋果.
拿去交租供樓.
接著他又想一想.
我都要吃的.
我都有生活需要.
我都要確保不會餓死.
於是他又想.
我就將三個蘋果.
拿去生活裡飲飲食食.
接著再望一望.
不是啊.
我還要買保險.
還有很多要去旅行.
又要買漂亮的衣服.
有空又要出去娛樂一下.
這三個蘋果.
我就拿去學玩去玩.
當他想完之後.
然後分配好了.
接著他再望一望.
只有一個蘋果剩下.
現在神就要我一個蘋果.
神不就是拿了我八分之八.
於是他就說.
神說不行.
你拿我八分之八.

$^{721}$你拿我全部剩下.
所以我不能夠給你.
就是這樣.
他本來答應了神.
給其中一個蘋果給他.
於是最後.
他一個蘋果都沒有給神.
是不是有時候都是這樣.
我們的優先次序.
放在這些東西.
分了在外生活.
分了這個分了那個.
最後才望回.
本來要獻給神的那個蘋果.
才望回.
神你好像拿了我很多.
其實有時候.
我們的優先次序的排列錯誤.
正好令我們有時候.
去學做奉獻很艱難.
但其實一開始.
上帝說你有十個.
你拿一個獻給我.
其實我們應該怎樣做呢.
應該先獻這個蘋果給上帝.
接著想想.
只有九個蘋果.
其實我少看一場戲都可以.
一年都去幾次旅行.
少去一次都可以.
家裡那麼多衣服.
還要買衣服.
少買幾件都可以.
其實你這樣.
調節了優先次序.
你會發現.
剩下的蘋果也夠用.
有時候就是這樣.
是我們奪取上帝.
用一個很好的理由.

$^{761}$不過原來我們的優先次序錯了.
令到我們本來是上帝的東西.
我們給不回上帝.
所以我自己奉獻的習慣.
就是月頭已經做奉獻.
因為我知道月頭做了奉獻之後.
我就將我剩下的錢.
在剩下的錢裡.
想想怎樣去花.
怎樣去渡過.
就不是去到月尾.
才想想再做不做奉獻.
有時候我們就是這樣.
去奪取了上帝的東西.
而且這種奪取.
甚至成為了一種問題.
一種我們信仰根基的問題.
這本書叫為何教會不像樣.
是陳希希牧師所寫的一本書.
當中其中有一課.
講到有關金錢觀的兩個基本真理.
兩個真理說什麼呢.
第一 神就是萬有的主宰.
問題是我們是否相信一切是屬於神的呢.
第二 祂擁有一切.
祂必然看顧我們.
既然你相信神是創造萬有的.
祂又擁有一切.
一切的需要祂都會看顧你.
但是為什麼我們走不到信心那一步.
去將當立的獻給神呢.
陳希希牧師說.
這可能源於我們對神的不相信.
我們不相信神真是萬有的主宰.
我們也不相信神真是看顧你.
照顧你生活所需.
所以今天我們會對著神怎麼說呢.
我們說 神啊 我的收入微薄.
我還要供養父母 養育兒女.
神啊 你不會不足夠的.

$^{801}$但是我不足夠.
神啊 你不需要衣食住行.
但是我需要.
神啊 你不是負資產.
但是我是負資產.
因此暫時我不能奉獻.
當我安頓好一切.
安排好生活之後.
有多的我再給你.
陳希希牧師說.
我們有時信仰裡面就是對神不信的狀況.
以致我們今天在物質上獻給神.
都成為了我們那種困難.
甚至成為了信仰裡面的糾結.
令我們和神顯得格格不入.
有一個美國浸信會的牧師.
不知道你有沒有聽過他.
這個牧師叫Isaac Barton.
這位浸信會的牧師.
是一位美國的牧師.
是1826年至1902年的人物.
已經過了西班牙天家.
他說很多年前.
他為神學院募捐.
很多神學院其實都沒有錢的.
而經費或者照顧神學生的錢.
都是像剛才說的那些第三個十分之一.
一些慈惠的方式.
靠別人捐助的方式.
才能夠營運到一些神學院.
這位美國牧師為了神學院去募捐.
於是他去到不同的地方.
其中有一次在半路中途.
他遇到兩個持槍匪徒.
持槍匪徒叫他將所有錢交出來.
而牧師從馬背上下來.
因為當時不是駕駛私家車.
當時的年代騎馬.
於是牧師從馬背上下來.
將一個大的和一個小的錢包放在路邊.

$^{841}$對這兩個匪徒說.
先生我是浸信會的牧師.
我的任務是來走遍田納西州.
為了訓練年輕的傳道人的神學院.
去募集一些款項.
小包的錢是我的.
你儘管拿去.
大包的錢是為了有困難的神學生準備的生活費.
你敢不敢接受.
就這樣問他.
於是那兩個匪徒聽完之後.
他敢不敢接受呢.
當時這兩個匪徒和牧師聊天.
開始問他為何神學院需要這些經費.
關於這些神學院籌募經費的事情.
而牧師就和他解釋.
和他解釋那些錢是怎樣.
接著聽完解釋之後.
那兩個匪徒就決定甚麼呢.
他們決定不會去拿那兩包錢.
兩包錢都不拿.
故事是不是完結了.
還未完結.
接著牧師就說你們都是善良的人.
多謝你們的好意.
既然你們都聽了這麼久.
你們都知道我所做的工作是甚麼.
那你們現在是否要做些奉獻呢.
接著那兩個匪徒.
每人各自奉獻了五元美金.
牧師就放他們走.
這個是很真實的.
你上網找到這個見證.
連匪徒都不敢奪神學院這些錢.
因為這些是屬神的.
但這群以色列的百姓.
他們竟然上下一心.
上梁下梁一起欺騙搶奪屬於上帝的物品.
所以這裡正正好說一個很重要的事情.
就是神的要求很簡單.

$^{881}$我們身心財物都是屬於上帝的.
我們只需要將這些屬於上帝的物品獻給上帝.
當然有時候我們信耶穌很可愛.
我們覺得好像神要這樣要那樣.
但我們不要誤會.
剛才說過這些物品本身是屬神的.
而這個獻上的目的只有一個.
就是要祝福你.
不自滅亡反得永生.
不自滅亡要得祝福.
所以第十至十二節就說出.
上帝要我們獻上的目的.
要讓我們蒙福.
我們一起讀讀第十至十二節.
開始.
你們要將當立的十分之一.
全然送入倉庫.
使我家有糧.
以此試試我.
是否為你們闖開天上的窗戶.
傾福於你們.
甚至無處可用.
這是耶和華說的.
我必為你們斥責蝗蟲.
不用它毀壞你們的土產.
你們田間的桃桃樹果實未熟已鮮.
也不會掉落.
這是萬君之耶和華說的.
萬國必稱你們為有福的.
因為你們必成為喜樂之地.
這是萬君之耶和華說的.
這裡十至十二節正好說出.
上帝的目的心意.
就是要讓我們成為一個蒙福的生命.
這裡的福氣有三種.
哪三種呢?.
第一就是屬天的福氣.
這裡說到為你們闖開天上的窗戶.
傾福於你們無處可用.
剛剛最近香港有黑雨.

$^{921}$雨多得水浸 路堵塞.
地庫整層都沒了.
原來神可以降雨.
有不同的雨可以降下來.
所以你記得在創世記裡.
是不是有一次下很大雨.
是不是下到不停.
是不是下到滿社.
是不是下到連方舟都載得起.
原來上帝肯開天上的窗.
雨可以下到滿社.
神就藉著這個片段.
回憶給他們知道.
如果上帝開的不是雨水.
如果上帝開的窗是開了福的窗.
那怎麼辦?.
掉下來.
你一開始以為是綠福.
接著變成黃福.
接著紅福.
接著黑福.
就是滿到社.
可以滿到社的福氣.
神就在這裡說.
原來真是水浸眼眉.
多到浸了你整個人都可以.
上帝的福氣.
但是這個福氣從哪裡來呢?.
從天上敞開.
所以這裡正正好是一個畫面.
原來神要給我們的.
是一種屬天的福氣.
神而來的福氣.
以至我們每一個願意獻上的人.
都是成為蒙福之人.
看到的不是一般的景象.
而是一種不是從天上來.
都不會發生的景象.
那種福氣是超越了.
你所能夠理解.

$^{961}$你所期待.
你所計算的.
問題是你相不相信.
可以真的大到這樣呢?.
原來神就說祂要大大的賜福給我們.
不只是天上要很多福氣降臨.
而這些福氣.
當來到地上的時候.
成為我們能夠肉眼見得到的.
地上的福氣.
地上的福氣是什麼呢?.
我必為你們斥責蝗蟲.
不容他毀壞你們的土產.
你們田間的葡萄樹果實未熟.
以先也不會掉落.
當時他們的困難.
就是可能面對一些饑荒.
但上帝說不是的.
上帝容讓你有充足的雨水.
讓你的土產能夠生長.
你們不需要擔心.
你們吃不夠,著不暖.
所以當這裡說.
如果有蝗蟲上來襲擊你們的話.
因為事實上你見得到.
聖經裡面說了很多蝗蟲成為災害.
其中《若爾書》就說到蝗蟲.
他說神都會為你斥責這些蝗蟲.
責備他們.
以至這些蝗蟲都要轉頭走.
不會毀壞你們的土地.
你們的財產得到保護.
接著你田間的葡萄.
真的要成熟的時候.
才會掉下來.
成熟代表什麼?.
就能夠賣到很價錢的葡萄.
最大,最漂亮,最合味道的.
已經成熟,甜的.
才會收割.

$^{1001}$不會掉下來,浪費.
原來神給我們地上的福氣是這樣的.
今天不知道你有沒有這樣的經歷.
有時候無端端沒有了錢.
有時候也會這樣.
有時候無端端壞這個壞那個.
又要花錢.
又要弄這個弄那個又要花錢.
無端端有什麼單又發給你又要花錢.
無端端掉了東西又不見.
又要買東西又要花錢.
你發現我們有時候.
可以隨時隨地都有一些.
所謂的蝗蟲在我們身邊經過.
侵蝕我們的財富.
但如果上帝保護你.
你不用擔心.
上帝保護你.
那個賊在你旁邊.
上帝都斥責他要離開.
地上的福氣是滿滿的.
不單止這樣.
第三個他說的是什麼呢?.
他說你們是稱為有福的.
為什麼呢?.
你們的福氣還會淋到.
你們裡面成為喜樂之地.
喜樂之地.
喜樂是什麼?.
喜樂不是一種物質.
是在心靈裡面的.
原來我們的心靈裡面.
都得到一份滿足.
一份喜樂讓我們生命裡面.
不是一個憂愁裡面.
所以你聽了很久.
《肥辣比書》很多課.
應該都記得什麼叫喜樂.
所以你們要喜樂.
你們要經歷喜樂.

$^{1041}$你們要經歷心靈那種福氣.
才由你們去學習.
去放上海錢.
事實上喜樂.
這種狀況並不一定.
需要和錢有關係.
今天我們誤會了.
我們要經歷喜樂的生命.
一定要很多財產.
很多物質的豐富.
其實不是.
而是我們在神的裡面.
我們可以經歷那種喜樂.
剛剛在星期三.
我們又有我們的團契.
星期三我們的迦勒團契.
我們都有很好的時光.
當中黃太跟我說了一些話.
黃太捉住我跟我說.
我真的很開心.
每個星期三的團契.
他說很開心.
他玩得很開心.
我們每個星期三的團契.
沒有賭錢.
不會賭錢.
不用有錢人瓜葛.
但我們一起聽聖經.
吃東西.
一起玩遊戲.
很開心.
如果你打麻將.
輸了,整晚都睡不著.
不服氣.
但如果一群弟兄姊妹相交.
神給你一種福氣.
一種相交的生命.
你走在一起.
輸也開心.
你看到下面的幅圖.

$^{1081}$就是黃太和幾位姊妹.
一起玩潛烏龜.
玩潛烏龜.
師母輸得最多.
成為了龜王.
龜后.
成為了龜后.
輸得最多.
輸了三次.
但他們四個人都不停地笑.
輸了都可以笑得出聲.
因為我們有心靈的喜樂和滿足.
弟兄姊妹.
我們不要被物質限制.
但同時上帝也要祝福我們物質.
我們不要只想著物質.
我們想著天上的福氣.
所以很重要.
原來上帝本意就是要我們去蒙福.
所以今天我想去澄清.
很多人都說.
基督徒要去做十一奉獻.
因為事實上這段聖經.
是和十一奉獻有關.
很多時候都說.
十一奉獻是一種責任.
是一種逼迫.
上道的逼迫.
去拿我的東西等等.
其實很多都是錯誤的觀念.
這些東西是屬於上帝的.
上帝要你奉獻十分之一.
去考驗你的信心.
考驗你信不信得過上帝.
上帝不是不夠東西.
很多時候說.
是不是教會不夠錢.
要奉獻.
不是的.
如果上帝不祝.

$^{1121}$我自己的看法.
如果上帝祝福的話.
這間教會是合上帝的心意.
就算人們不奉獻.
上帝都會令這間教會有錢.
這個就看回.
慕勒他們開孤兒院那件事.
稍後有機會再說.
上帝就會開路.
所以上帝不會不夠錢.
上帝會感動一些.
願意承擔的人去做.
但問題是.
上帝很想每一個信徒.
都去經歷信心的功課.
學習奉獻的功課.
讓他們都經歷蒙福.
弟兄姊妹.
不是教會要.
是上帝想你們每一個都經歷蒙福.
體會到天上的福氣是很多.
體會到地上的福氣.
你肯付出.
上帝祝福你.
看到你心靈裡.
不再被金錢所限制.
上帝要讓你有喜樂的生命.
最後總結.
馬拉基怎麼說.
馬拉基很簡單.
這是他體會到真實的情況.
但你需要做一個事情.
就是以此示善.
以此示善.
神說當你去學習.
將上帝的東西.
送入倉庫.
使神的加油量.
你用這件事試一試.
你就會體會到上帝是真實.

$^{1161}$這裡的試一試.
並不是試探神的那種試.
而是你要做一次.
神抵不住.
你總是覺得.
在你那個位置裡.
卡住了走不了.
你試試做一次.
你做一次就會體會到.
體會了你就明白了.
牧師說這麼久都沒有用.
因為你沒有做過.
你體會不到原來那種福氣.
我信耶穌還有三十多年.
已經三十五年了.
三十五年裡我都有做.
而且是超過十一奉獻的.
我自己做.
是不是.
我來到佛羅倫多.
我胖了.
我沒有瘦.
我不會因為餓而瘦了.
你信不信.
上帝要求我們去試一試.
你去做一次.
你去經歷上帝給你的祝福.
怎樣去試呢.
其中有兩個很重要的.
向度.
第一.
一定要轉向.
你記住一定要向神.
你不要向著錢.
錢會吸引你很多方法.
金光閃閃吸引你.
但經文要求我們要轉向.
真的要整個身轉向神.
放上信心.
答案還不知道.

$^{1201}$但你試了就知道.
你去做一次.
沒事的.
安心的.
你去做一次.
你就會經歷上帝那種生命的祝福.
滿滿.
最後說一個見證.
是我們阿加勒團.
這個星期有看的一個見證.
因為在座當中.
只有兩個是當天在場.
我再說一次.
當中的見證是.
大家聽過中年好聲音嗎.
中年好聲音大家有聽過.
其中的主持.
就是車婉婉.
車婉婉就是中年好聲音的主持.
事實上.
我最近看她的生命分享見證經歷.
她自己都經歷過.
怎樣去信耶穌.
其中一個信耶穌的經歷.
就是面對著.
她媽媽有一個餐館.
就是一個餐廳.
其實就是因為.
有些伙計穿貴通.
偷錢.
特別是那個會.
不知不覺間偷了很多錢.
公司不夠錢 公司要錢等等.
甚至有一次.
她跟車婉婉說.
你今天要拿到四十萬回來公司.
公司才能夠過到今天.
四十萬.
即是五萬美金.
老實說.

$^{1241}$誰借五萬美金給我.
沒有吧.
我跟你這麼熟你都不借給我.
車婉婉的見證是怎樣呢.
其實她當時都有少少接觸.
因為她自小都有一些基督教學校.
天主教學校的背景.
有些信仰的背景.
但她就未說信耶穌.
但她自己在過程中.
因為她身邊有些基督徒的藝人朋友.
亦給予她很多鼓勵.
於是她就向神做了一個呼求祈禱.
她說.
神啊我現在要四十萬.
真的應付不來.
如果神今天讓我籌到這四十萬.
我就順利了.
就是這樣.
這四十萬.
她就跟神說.
其實她在試神.
如果你今天做到我就順利.
她就是這樣去做的時候.
神就知道她的需要.
知道這個女兒.
她怎樣要經歷神.
而這種經歷她不是要試探神.
而是她在困境中的呼求.
以致上帝回應她的禱告神是怎樣呢.
她就說剛好那天.
她就跟一個認識了兩三個月的阿姨.
一個認識了兩三個月的阿姨.
一個做生意的阿姨.
就聽到她說.
她需要四十萬.
這個阿姨就叫她.
跟職員去樓下銀行.
馬上把四十萬現金給她.
我認識了你三年了.

$^{1281}$剛才問你借五萬都沒有人肯.
但這個認識了兩三個月的阿姨.
她就這樣白白地借了四十萬給她.
是不是因為這個阿姨是好人呢.
可能她是一個好人.
但我更加相信.
是因為她向神發出求救的禱告.
所以神就幫助她.
讓她真的那一天.
不是遲一天.
是那一天就見到這四十萬.
來到她需要裡面.
她就這樣更加認定了.
原來這個神是真的.
她願意信靠神.
她經歷了很多不明白的事情.
但她就走了信心的步伐.
信神之後.
一直去很多不明白的地方求神.
去引領她的路.
其中第二個的事是什麼呢.
第二個事就是她的工作.
其中一個工作是在電台做的.
之前信耶穌的時候.
她有一個工作在電台做.
其中做什麼節目呢.
她做星座 講星座的節目.
她心裡很疑問.
她問她的牧者.
她自己都覺得有點不妥.
她問我在電台講星座.
星座會不會不是一個很好的記錄圖.
因為星座是什麼呢.
相信那些星座.
你今天見到什麼就很黑仔.
你今天要近些什麼.
她說會不會有問題呢.
她自己都在想.
她就問牧師.
牧師說什麼.

$^{1321}$牧師怎樣回答她.
不行.
可以.
牧師說我們為這件事祈禱.
在這件事祈禱.
結果.
轉頭那個星座節目就倒閉了.
神知道她的需要.
就倒閉了那個節目.
做其他.
現在中年好聲音.
第三個.
她怎樣去經歷.
試到神是真實呢.
她自己都體會到.
她願意讓她兒子在教會裡成長.
以致她體會到.
神祝福她的兒子.
有一天她兒子問她.
媽媽.
你知不知道誰是超級英雄.
她媽媽說.
是不是超級英雄.
不是.
是不是蜘蛛人.
不是.
是不是蜘蛛人.
不是.
那是誰呢.
Jesus.
是的.
這個小孩就是這樣說.
Jesus.
當她願意將她的小孩.
信心地交給神.
她就能夠在神的裡面.
這是最神的一種回應.
你試一下.
你將你的子女交給神.
有時我們拿著子女都不給神.

$^{1361}$我自己用我的方法去管教.
但你信得過神.
就交給神去管教.
求主幫助我們.
學習信心的一步.
今天說的不單止是.
一種十一奉獻.
對這些以色列人不遵從上帝的吩咐.
更加是一種挑戰我們信仰.
信心的功課.
求主幫助我們.
更加在神的裡面去堅穩.
踏上信心的一步.
以致真的面對上帝.
給我們大大的祝福.
我們一起祈禱.
因為你給我們的心意.
原是要賜福.
不是降災的心意.
而是要賜福的心意.
你不是要我們滅亡.
而是要我們永生永福.
你本意就是對我們這麼好.
你就是那位愛我們的天上爸爸.
你這位爸爸對我們這些子女.
無論我們多麼不聽話.
你仍然顯明了這份愛.
主我們作為你的兒女.
我們學懂將有的糖請你吃一粒.
我們有的金錢獻給你.
讓你自己得著榮耀.
幫助我們去學習生命的功課.
實踐的功課.
聽我們祈禱奉耶穌名字而求.
今天的你可以決定明天的你.
求主幫助我們今天要立志.
最後還有幾句說話.
我們一起讀這幾句說話.
好 開始.
思想決定行動.

$^{1401}$行動決定習慣.
習慣決定品德.
品德決定命運.
究竟在很多年之後.
神跟你說你也有今天.
究竟是什麼日子呢?.
求主幫助我們.
我們一起祈禱.
因為你讓馬拉基先知.
在他們面對困惑的世代裡.
作出審判的訊息.
一個挑戰的訊息.
以致我們真的能夠知道.
我們等待彌賽亞的來臨的時候.
究竟對我們是真的好還是壞.
要視乎我們怎樣去活出我們的生命.
求主你幫助我們教會.
每一個聆聽你話語的人.
我們都有份心志.
有份悖念的心志.
我們要成為精心獻給上帝.
我們今天習慣了很悠閒.
很自在的生活.
求主你幫助我們更加知道.
我們要有一種韌力.
去面對上帝你自己.
以致我們成為一個最好的祭物獻給你.
幫助我們在這個世代裡.
堅守我們的信仰.
聽我們祈禱奉耶穌的名字而求.
阿們.
馬拉基書裡面那班百姓.
不知不覺影響了他們的社會.
信仰 家庭.
整個角度不知不覺.
已經不在上帝的喜悅裡面.
所以神體馬拉基書裡面.
帶來那種對祭司和對百姓的責備.
是真實的.
亦都是迫切的.

$^{1441}$因為這個問題如果不認真處理的話.
會帶來很大的破壞.
最後一句總結.
我們有一段聖經.
大家一起讀這段聖經.
就是來自尼加書.
另外一個先知書所講的說話.
我們一起讀這句經文.
六章八節開始.
西人啊.
耶和華已指示你何為善.
他向你所要的是什麼呢.
只要你行公義.
好憐憫.
全謙卑的心與你的上帝同行.
求主幫助我們.
我們能夠活出一個公義的社會.
一個憐憫的婚姻.
亦都能夠在上帝面前.
全一個謙卑的心.
亦都能夠活出聖潔的國度.
我們一起祈禱.
最後幫助我們.
因為我們知道.
你給我們不單止有君尊的祭司.
好像上次所講一樣.
亦都要求我們有聖潔的國度.
一個屬神理的國度出現.
願主你幫助我們.
記著《馬拉基書》裡面的思考.
我們怎樣去遠離.
逃避這個世代裡面.
那種詭詐的行為和生活.
讓我們活出聖潔的生命.
你只能夠去榮耀你.
你只能夠在這個社會裡面.
更加被人見得到.
我們是誠實的一群.
我們是可信賴的一群.
以致更加多人願意來到神理面前.

$^{1481}$去認識神理自己.
願主你幫助我們.
亦都幫助我們.
一會兒下來有分享討論團契的時間.
祈禱奉耶穌的名字而求.
阿們.
\newpage



\section{以弗所書 4:1-16}
\label{sec:032330}
\textbf{合一的關鍵}~~~~吳銘淇傳道
\newline
\newline
連結: \href{http://www.church.com.hk/acms/content.asp?site=cdc\&op=show\&type=product\&code=032330\&layout=sermon}{\texttt{032330}} ~~~~ 語音日期: 2026-06-21
\newline
\newline
\hyperref[sec:032329]{\small{< < < PREV SERMON < < <}}
~
\hyperref[sec:index_chronic]{\small{[返順時目]}}
~
\hyperref[sec:032331]{\small{> > > NEXT SERMON > > >}}
\newline
\newline
以弗所書 4:1-16
\newline
\begin{longtable}{cl}
\hline
\hline
章節 & 經文 (和合本修訂版)\\
\hline
4:1 & \begin{tabularx}{0.7\textwidth}{X} 我為主作囚徒的勸你們，既然蒙召，行事為人就要與你們所蒙的呼召相稱。 \end{tabularx} \\ \\ \relax
4:2 & \begin{tabularx}{0.7\textwidth}{X} 凡事要謙虛、溫柔、忍耐，用愛心互相寬容， \end{tabularx} \\ \\ \relax
4:3 & \begin{tabularx}{0.7\textwidth}{X} 以和平彼此聯繫，竭力保持聖靈所賜的合一。 \end{tabularx} \\ \\ \relax
4:4 & \begin{tabularx}{0.7\textwidth}{X} 身體只有一個，聖靈只有一位，正如你們蒙召，是為同有一個指望而蒙召， \end{tabularx} \\ \\ \relax
4:5 & \begin{tabularx}{0.7\textwidth}{X} 一主，一信，一洗， \end{tabularx} \\ \\ \relax
4:6 & \begin{tabularx}{0.7\textwidth}{X} 一神－就是萬人之父，超越萬有之上，貫通萬有，在萬有之中。 \end{tabularx} \\ \\ \relax
4:7 & \begin{tabularx}{0.7\textwidth}{X} 我們每個人蒙恩都是照基督所量給每個人的恩賜。 \end{tabularx} \\ \\ \relax
4:8 & \begin{tabularx}{0.7\textwidth}{X} 所以有話說：「他升上高天的時候，擄掠了俘虜，將各樣的恩賜賞給人。」 \end{tabularx} \\ \\ \relax
4:9 & \begin{tabularx}{0.7\textwidth}{X} 既說「他升上」，豈不是指他曾降到地底下嗎？ \end{tabularx} \\ \\ \relax
4:10 & \begin{tabularx}{0.7\textwidth}{X} 那降下的，就是高升遠超越諸天之上的，為要充滿萬有。 \end{tabularx} \\ \\ \relax
4:11 & \begin{tabularx}{0.7\textwidth}{X} 他所賜的有使徒，有先知，有傳福音的，有牧者和教師， \end{tabularx} \\ \\ \relax
4:12 & \begin{tabularx}{0.7\textwidth}{X} 為要裝備聖徒，做事奉的工作，建立基督的身體， \end{tabularx} \\ \\ \relax
4:13 & \begin{tabularx}{0.7\textwidth}{X} 直等到我們眾人在信仰上同歸於一，認識神的兒子，得以長大成人，達到基督完全長成的身量。 \end{tabularx} \\ \\ \relax
4:14 & \begin{tabularx}{0.7\textwidth}{X} 這樣，我們不再作小孩子，中了人的詭計和欺騙的法術，被一切邪說之風搖動，飄來飄去。 \end{tabularx} \\ \\ \relax
4:15 & \begin{tabularx}{0.7\textwidth}{X} 我們反而要用愛心說誠實話，各方面向著基督長進，連於元首基督， \end{tabularx} \\ \\ \relax
4:16 & \begin{tabularx}{0.7\textwidth}{X} 靠著他全身都連接得緊湊，百節各按各職，照著各體的功用彼此相助，使身體漸漸增長，在愛中建立自己。 \end{tabularx} \\ \\
[1ex]
\hline
\hline
\end{longtable}
$^{1}$靈姐妹早晨.
讓我們為著講到和聽到的福事一同禱告.
是因著我們在基督和聖靈的當中.
叫我們能夠去得著這個身份.
得著這個屬靈的家.
求主在當中讓我們都去一同領受上帝的話語.
求主在當中祝福我們.
讓我們能夠聽到而行到.
奉主耶穌基督的聖名祈求 阿彌陀佛.
認真去研讀聖經.
可能去參與一些網上的查經研經的聚會.
甚至和幾個志同道合的基督徒的朋友定期聚會.
這就是他的信仰生活.
而美國甚至有一個研究報告說.
在2023年的時候.
有25\%的美國成年人經常透過網絡和電視.
去參與宗教的聚會.
他們覺得這個狀態是很滿意的.
而這個數字其實也在不斷上升.
也有不少人說信仰依然是存在的.
只不過是不再需要依賴教會.
他們為什麼不再在教會裡面或沒有教會的生活呢?.
研究有很多不同的原因.
但可以歸類成兩個.
第一個是關於教會的問題.
就是對教會的失望.
有不少受訪者認為教會不能夠回應當下時代的苦難.
心理健康的問題.
又或者一些社會的議題.
也有人覺得教會只是說一套做一套.
所以面對著有時看到一些領袖的道德失誤.
又或者在一些權力鬥爭上的問題.
他們就很失望離開.
第二個問題是來自於信仰漸漸變得很個人化.
這個也是一個很重要的文化趨勢.
就是越來越多人覺得信仰只不過是一個很個人的事.
這是我和上帝之間的關係.
我的屬靈生命是否健康.
我和神的關係是否親密.
其實這個最終都是值得我和上帝知道.

$^{41}$其實其他人你們根本沒有辦法去評論.
也不能夠因著我有沒有參與什麼聚會.
而去決定我和上帝的關係是怎樣.
所以群體不再是信仰裡面不可或缺的一個部份.
而只是其中一種的幫助.
人仍然可以透過我自己有讀經.
我有祈禱.
我甚至有聽講道.
和我幾個朋友一起查經.
這個都是我的信仰生活.
所以即使沒有教會.
我仍然可以做一個好的基督徒.
在這種觀念底下.
教會不再是信徒身份一定要有的一個部份.
而只是信徒可以選擇的其中一種屬靈的資源.
就好像你可以選擇這本靈修書.
也可以選擇那本靈修書一樣.
當教會有幫助.
覺得這件事不錯的時候.
我就參與.
當覺得好像不太適合我的茶的時候.
我就可以離開.
但是我的信仰本身從來都不是依賴在群體當中.
於是我們就會問一個問題.
就是我仍然信耶穌.
我有讀經.
我仍然和我的基督徒朋友一起查經.
其實我是否已經不需要教會呢?.
可能你都不需要去到離開了教會的人才會問這個問題.
可能現在很多人都說自己是埃人.
或者很內向.
可能我們都會問這個問題.
是否需要教會呢?.
是否需要那麼多的社交.
那麼多的人.
經常都要和那麼多的人去相處呢?.
這個正正就是今天我們都會問的問題.
而保羅在二福所書裡.
似乎都給了我們一些思考的方向.
當今天我們會認為.

$^{81}$首先應該是由個人的.
然後才是我和上帝的關係.
然後才去決定.
是否需要用屬靈資源.
是否需要進入教會當中.
但是保羅是否這樣說呢?.
保羅卻是覺得.
當一個人被基督拯救的時候.
他同時應該要被帶入一個群體當中.
基督徒的身份從來都不是自己一個.
而是群體的.
所以二福所書的對象.
是對著每一個信徒.
但是他的焦點.
一直都是整個的群體.
就是神使我們成為了他的指紋.
他的家.
甚至是屬靈的宮殿.
全部都是一個群體的描述.
甚至之後四到六章.
那些道德的教導.
都不是說信徒各自要做些什麼.
而是說整個群體裡面.
怎樣互相提醒.
互相督責.
去活出夢照的身份.
說到二福所書.
不如我們嘗試重溫一下.
二福所書的結構.
其實它和其他的保羅書信都很相似.
二福所書可以分為兩個部分.
前半部是說一些神學.
說上帝是誰.
他在做些什麼.
後半部分是一些應用.
就是因著我們看到上帝的心意是這樣的時候.
教會當下處境要怎樣去回應.
二福所書二到三章裡面.
告訴我們一個很重要的奧秘.
就是以前你不知道的.

$^{121}$現在你應該知道的.
就是藉著基督的拯救.
然後叫猶太人和外邦人.
能夠一同去建立教會.
又因著聖靈的印記.
叫教會的群體.
有一個新的身份.
就是上帝子民的身份.
這是上帝的救贖計劃的一個很重要的見證.
因為一章十節說.
上帝要藉著基督使天上地下萬有都要同歸於一.
因此原本兩班人.
猶太人和外邦人.
有很多的文化衝突.
甚至有很多的國仇家怨.
但都能夠在同一個群體裡面.
去合一的生活.
就是這個見證.
去見證展示半生福音的能力.
來到第四章.
就是開展下半部的一個很重要的開始.
就是由神學上帝是誰.
他在做什麼.
到我們要做什麼.
我們應該要怎樣.
保羅說四章一節.
所以既然行事為人.
就當與蒙召的因相稱.
進入到第四章.
保羅開始講解.
既然神已經創造了.
或建立了這個新的群體.
上帝的子民應該怎樣去活出.
和蒙召相稱的生命.
而四章十七節之後.
就開始講各種很多道德的教導.
例如怎樣寬恕人.
怎樣不要貪心.
怎樣不要行淫.
講了很多這些問題.

$^{161}$其實這些都是建立在哪裡呢.
就是四章一到十六節.
裡面講群體是什麼.
你們為什麼會成為一個群體.
就要講這個部分.
所以我們今天會處理兩個問題.
在四章一到十六節.
第一個就是這群原本彼此敵對的人.
他們怎樣可以維持合一呢.
他們靠些什麼呢.
第二就是神為什麼這麼看重.
一定要他們合一.
一定要有一個合一的群體呢.
第一個四到十六節.
是講維持合一的原因或基礎.
保羅提醒異忽所有信徒.
你們要保守聖靈所賜合而為一的心.
之後他引用了一段信仰的宣言.
就是一個是身體.
一個是聖靈.
一個是指望.
一主一信一洗一神.
不少學者都認為.
這個可能是初期教會.
他們在洗禮時候.
會用的信仰宣言.
值得留意的是.
保羅在這裡不是要求這群信徒.
去努力創造合一.
而是要提醒他們.
這個合一是存在的.
是存在了的.
你們本來已經是一個群體.
很難明白.
舉個例子.
我剛剛剝完智慧刺.
今天可以講到原來已經很感恩.
之前已經腫了.
完成了手術的時候.
醫生開始坐起來.

$^{201}$朗一下口的時候.
醫生就說你剛剛做手術的時候.
有一個位置不斷出了很多血.
然後他就這樣說.
我就很緊張.
心裡面其實都知道.
應該處理了.
不然沒有辦法讓你坐起來.
但心裡面就這樣.
那我現在應該做些什麼.
我應該怎樣才可以治到血.
醫生就靜了兩秒.
就很冷靜的說.
都搞定了.
其實你想一想.
一個專業的醫生.
他能夠完成一個手術.
一定是要處理的東西應該處理了.
或者是如果處理不了.
都要告訴你要怎樣去做.
我們要做的是什麼呢.
可能就是要怎樣再好好地.
定時地吃藥.
敷一下冰讓它沒那麼腫.
然後做好一些清潔.
吃一下流質的食物.
總之就是不要.
盡量不要讓傷口再刺激它.
令它出血.
或者令它發炎.
同一樣的道理.
要猶太人和外邦人.
他們要處理傷口.
這個多年來的恩恩怨怨.
要他們.
我們就要抱頭抱頸.
彼此接納.
根本是不可能的.
他們根本靠自己沒辦法做到.
他是因著基督的十字架.

$^{241}$就是原本彼此敵對的猶太人和外邦人.
已經在身體裡.
因著基督和上帝和好.
又因著同一位的聖靈.
讓他們能夠成為神家裡的人.
這就是二章裡講給我們聽的那件事.
因此保羅就要呼籲.
他們要保守這份的合約.
而不是要怎樣去創造合約的機會.
因為醫生已經做完手術.
我們只要好好照顧好傷口就足夠了.
合約已經出現.
神已經在基督裡成就了這個合約.
所以換句話來說.
四章的四到六節.
這個信仰宣言.
是一個合約的命令.
而為了這個命令.
提供了一個很重要的根據.
就是你們憑什麼.
怎樣可以合約呢.
信徒之所以合約.
是因為他們在基督裡是同一個身體.
我們有著同一個信.
有同一個上帝.
今天我們每隔個月.
就會唸一次《仕途信經》.
唸的時候不知道你們在想什麼呢.
可能就是追著字讀.
又覺得有點快.
可不可以讀慢一點.
可能想很多不同的東西.
但是你們有沒有想過這是上帝給教會的一個證據.
就是上帝已經將合約賜給我們了.
然後因為我們都相信同一個上帝.
相信祂為我們帶來的救贖工作.
就叫我們雖然很多不同.
但因著在基督裡都已經成為同一個身體.
所以我們才能分享同一個《仕途信經》.
一齊讀 一齊說阿們呢.

$^{281}$不過可能有人會說.
Kevin 不是這樣.
因為不同宗派都有不同的《信經》.
你唸這個 我又唸那個.
但可能未必大家認同.
不要說到《信經》那麼大.
不同的教會都有在不同的神學議題上.
有不同的看法.
這樣我們是不是還是合一呢.
這個上帝給我們的根據.
是不是還存在呢.
答案是可以的.
因為保羅不是說.
那個一信就是你們所有信的東西是完全一致的.
沒有任何不同的答案.
而是說我們共同承認福音最核心的真理.
例如我們都相信父子 聖靈.
都是神 有位格的.
但是在歷代裡面.
都會有這樣的.
都會在想怎樣表達三一呢.
又或者他們在位格上的關係是甚麼呢.
這些都有很多不同的討論.
但會不會影響我們去承認.
所以祂不是三一的上帝.
應該不會.
這個東西我們必然要這樣去承認.
同樣我們相信耶穌基督是真神.
也是完全的人.
但祂也為了我們的罪而死.
死在十字架上.
不過神學家仍然都會在很多細節裡面.
有不同的理解.
例如是祂是完全的人.
也是完全的神.
究竟耶穌在世的時候.
祂是否仍然擁有著很多神的能力呢.
包括祂是不是全知.
是不是全能呢.
這些問題神學家仍然有很多不同的見解.

$^{321}$但是福音的核心.
耶穌帶來的拯救.
這個從來都沒有人能夠否認.
再來就是.
我們相信基督會再來.
死人會復活.
神會施行最終的審判.
不過千禧年是甚麼呢.
怎樣去理解呢.
不同的教會都可能有不同的看法.
但同樣不會影響教會要去承認.
就是我們一起去等候主再來的時候.
所以原來合一都不是說每一個意見.
完全我們就是一致的.
大家都是完全的相同.
而是我們共同去承擔了福音的核心的真理.
沒有人能夠否認的.
因為否認了你就不是基督徒了.
所以我們去誦讀信經的時候.
我們其實都是在見證這件事.
就是神已經將我們放在同一份的救恩裡面.
是因著同一位的主.
同一位的聖靈.
同一位的父神.
我們相信的同一位的基督.
承受著同一個的福音.
所以讓我們有合一的憑據.
這也是作為基督徒一個很重要的根基.
一個基礎.
我們知道保羅說完.
原來我們合一是因為同一位上帝.
已經將我們放在這份救恩.
放在基督的身體裡面.
然後去建立教會.
所以之後他就要再一步去說.
就是神為什麼要去建立這個合一的群體呢?.
保羅引用了詩篇六十八篇的十八節.
裡面他要證明的是.
基督已經將屬靈的恩賜賜給他的子民.
就好像詩篇裡面.

$^{361}$他本身是想歌頌那個古代德姓的君王.
他凱旋回朝.
那些百姓全部都很歡慶.
然後他要做什麼呢?.
就是這個君王將那些戰利品派給他的百姓.
他是一個這樣的圖像.
這裡也一樣.
他想說的是基督戰勝了死亡.
也戰勝了罪惡和一切屬靈的權勢.
這是一章裡面告訴我們.
之後他就將這些恩賜帶給教會.
他所示下的恩賜包括什麼呢?.
就是使徒,先知,傳福音,牧者和教師.
裡面的角色包括有人負責開拓教會.
要去建立教會.
有人負責要去勸勉.
也要造就信徒.
也有人要去傳揚福音.
有人負責牧養.
有人負責教導的工作.
這些的職份都是特定的.
因為是基督賜給你的.
所以不是人說要做就可以做的.
但是這些的職份也不是為了讓少數的人去承擔所有的職份.
或者所有的事奉.
而是為了要裝備整個的教會.
所以我們看第十二節.
十二節他不是說.
所以這班領袖就是這班使徒,先知,傳福音,牧師,教師.
他們就要負責事奉.
那其他你們沒有這些恩賜的人怎麼辦?.
你就接受服侍.
坐在那裡按按腳.
然後就有飯吃了.
他不是說這樣.
而是他說神賜下這些的恩賜.
為的是要裝備這些的聖徒.
聖徒就是我們.
使得我們這班的聖徒都能夠各盡其職.
每一個人都一起去建立基督的身體.

$^{401}$所以換句話來說.
神賜下的這份救恩.
不是為了要建立領袖.
而是要建立整個的教會.
也是全體的信徒都要一同成長.
當信徒在信仰和認識基督上都同歸於一.
意思是什麼呢?.
就是說他們都了解了.
都認識到這個上帝了.
然後去接納這個上帝給他們的價值觀.
接著這個價值觀去活的時候.
有兩樣東西.
一個步驟會發生.
就是整個群體會慢慢成熟.
然後會長大.
而最終就會達成基督長成的新量.
就越來越像基督.
其實這個就好像和在四福音所說的.
追隨基督要去作門徒的道理是一樣的.
就是作門徒這件事要怎樣發生呢?.
原來是在群體裡面發生.
保羅說這件事是需要在群體裡面.
我們一起互相提醒.
然後一起去接納上帝的價值觀的時候.
整個群體就會一起長大.
然後越來越像基督.
保羅用一個人去形容這種成長成熟成為一個例子.
所以這個也是教會建立的過程裡面.
告訴我們它是需要時間的.
亦都會經歷起跌的.
我們都做了二十三十幾年人.
我們可能都知道.
作為一個小朋友變成一個大人.
在一個過程裡面必然經歷很多挫敗.
可能有很多我們知道要做的事.
但又不想做.
可能會經歷扭計.
經歷發脾氣.
經歷很多禮頑節曲的事都可以說出來.
這個就是小朋友.

$^{441}$然後慢慢經歷成長.
慢慢變得成熟.
我們都是由這樣變出來的.
所以一樣.
當教會朝向這個方向去成長的時候.
經過時間.
經過這群人一起成長.
會有一個結果.
保羅就形容.
這群信徒不再像小朋友一樣.
容易受各種不同的錯誤思想.
或者價值觀所動搖.
這些錯誤思想不是來自於.
那些非信徒的生活方式.
所以保羅在第五章裡面告訴他們.
你們要怎樣小心不要在淫亂裡面.
不要在貪婪和拜偶像等等的事情上.
因為他在說的正正就是當時.
異忽所那群非信徒過的生活.
信徒要互相提醒.
不要重蹈覆轍.
更加不要得罪他們的上帝.
當信徒在真理上成長.
這個群體就不再隨波逐流.
他們的思想亦不會那麼容易被人引走.
好像小朋友一樣.
因為我估計我們今天.
作為一個成熟的大人.
我相信你不會為一粒糖而在荒野打架.
第一個小朋友就會.
之後保羅就在十五到十六節.
最後用身體成為了一個比喻.
這個比喻成為了一個總結.
他說教會就好像一個身體.
基督是個頭.
每個肢體都有自己的功用.
有人是教導的.
有人是學習的.
有人是勸勉的時候.
就有人需要他被勸勉.

$^{481}$有人去扶持軟弱的人.
正正是因為有軟弱的人需要被扶持.
當每一個人都按著上帝給他的恩賜服侍.
彼此的配搭.
彼此都認真做好上帝給他的恩賜的時候.
就能夠彼此的建立.
然後整個身體就會在愛裡面.
一路一路地成長.
這個亦是保羅對教會的一種想像.
就是一個成熟的群體.
從來不是靠著幾個領袖.
你們自己搞定它.
你們要撐起.
你們不可以跌下.
不是這樣.
而是每一位信徒都要去參與其中.
不只是一齊來聚會.
然後坐下聽完我就走.
不是這樣.
而是在真理裡面.
我們有責任要彼此的提醒.
彼此的扶持.
彼此的建立.
這樣才會讓教會逐漸由小朋友變成大人.
不再為了一顆糖而打架.
弟兄姊妹.
今天我們願不願意成為彼此提醒.
彼此去認罪.
甚至當我們有軟弱的時候.
能夠彼此代討的群體呢?.
我們能不能夠一起和弟兄姊妹.
和你旁邊的弟兄姊妹走人生不同的階段.
去和他們同行.
又或者當我們偏離了真理的時候.
你覺得在教會裡面有沒有人願意用愛心去提醒你呢?.
當跌倒的時候.
又有沒有人願意伸手去扶持呢?.
今天我們教會有沒有發揮這個功能呢?.
在保羅的眼中.
教會存在的目的不是一群信徒聚會.

$^{521}$然後大家坐下來完成一個崇拜.
大家離開就可以.
不是這樣的.
而是信徒要一同成長.
然後滿有基督的新娘.
在《以弗所書》四章166節裡.
保羅很清楚要我們知道一件事.
就是教會的合一不是人的計劃.
而是神的計劃.
神的工作.
首先,神已經藉著基督在十字架裡.
將原本敵對的猶太人和外邦人帶進基督的身體裡.
去建立這個教會.
所以信徒的責任從來都不是創造這份合一.
而是要保守上帝賜下的這份合一.
同時,神也不單單賜下了合一.
祂也賜下了恩賜.
所以基督將不同的恩賜和積分賜給教會.
為的就是讓信徒可以裝備.
讓每一個信徒都發揮他的能力,恩賜.
然後去建立基督的身體.
最後,當信徒彼此去配搭.
彼此去建立一同成長的時候.
教會就慢慢成熟.
不再像小朋友一樣.
在各種錯誤中容易被動搖.
如如此基督.
所以保羅提醒教會要合一,成熟,彼此建立.
要成為一個這樣的群體.
這樣才能見證上帝讓世人看見.
上帝要藉著基督.
在天上地下要同歸於一這麼重要的上帝旨意.
這正好回應了我們今天要問的問題.
如果我們仍然相信耶穌.
我們有讀經.
甚至可能我也有和幾個我自己熟悉的基督徒朋友.
一起查經,祈禱.
其實我是否就不需要再回教會呢?.
二不熟說告訴我們.
神拯救我們不只是要我們和上帝建立關係.

$^{561}$神拯救我們是要將我們帶進一個新的群體裡.
是因為信徒無法單打獨鬥.
無法自己一個.
我們需要學習基督的教導.
所以我們需要被教導.
可能你說這個也可以在網上或其他資源裡做到.
但我們仍然需要群體.
不只是因為我們在群體裡學習.
也是因為這個學習不是一學就懂得.
正如你做運動.
現在很流行哈羅斯.
不是一開始就能做到.
原來當中有很多技巧.
當中也要訓練.
信徒一樣.
我們學習謙卑.
學習基督的樣式.
我們仍然有很多東西要操練.
而這些操練需要在群體裡進行.
因為我們需要互相提醒.
亦需要彼此的鼓勵.
我們需要在我們知異難行的時候.
我們需要守望.
需要幫助.
同時這個群體不是我們自己選擇的.
是上帝給我們的.
因為我們無法選擇誰來建立基督的身體.
甚至我們需要跟自己不同的人.
才能夠提醒我們的問題究竟在哪裡.
所以神設立教會不只是讓信徒聚集.
而是讓信徒一同成長.
整個群體都成為門徒.
學習基督.
既然我們都知道信仰群體很重要.
因為這是上帝所賜的.
祂賜下了合一.
也賜下了恩賜.
然後讓我們能夠各按其職.
一同成長.
那麼我們又怎樣去面對.

$^{601}$或者處理這個群體呢.
深水清的你應該知道.
有兩隻經文我們很熟悉.
但剛剛一直沒有說.
就是四章的二到三節.
凡事要謙虛 溫柔 忍耐.
用愛心互相寬容.
以和平彼此聯絡.
竭力保守聖靈所賜.
合而為一的心.
謙虛 溫柔 忍耐.
這幾個字的原本意思.
或者是很直白地表達出來.
都是在說一些要放下自己的東西.
包括放下自己的驕傲.
放下自己的憤怒.
或者是報復的心.
要自己去控制自己.
是一個這樣的意義.
所以事實上.
保羅之所以要求信徒.
要去謙虛 溫柔 忍耐.
用愛心互相寬容.
正正就是因為教會裡面.
本身就是充滿衝突.
有很多的摩擦.
甚至裡面仍然有很多的軟弱.
如果教會裡面的人.
全部都已經很成熟.
甚至到了大家的性格很相近.
背景很相似.
想法亦都很相近的時候.
我相信保羅應該不會需要去提醒.
你們仍然要忍耐.
你們仍然要包容.
但正正是教會.
是一群不同的人.
被基督帶來一起.
所以可能會有衝突.
可能會有誤會.

$^{641}$甚至經歷到失望.
今天我們經常在網上看到.
教會病了 教會病了.
其實這句說話.
多多少少都有真實的.
有些人可能因為在教會裡面經歷受了傷.
所以他要離開.
有人可能是見到領袖.
在道德上的失誤.
可能有人經歷了.
在權力上的挑戰.
覺得在教會裡面有這麼多政治.
亦有人會因為誤解.
被忽略或被論斷.
而感到失望.
這些傷害.
不一定是虛構出來.
是假的.
有時可能教會都有很多問題.
這些失望.
這個感受亦都很合理.
不過教會其實不是今天才開始病.
當我們去讀保羅書信的時候.
你會發現每一間教會.
都有自己的問題.
哥林多教會分黨分派.
有很多淫亂亦有互相的控訴.
加泰甚至偏離了福音.
加上加教會.
有人是硬硬腳不肯做.
只是飯來張口.
以不所教會.
今天都在被人提醒如何去合一.
換句話說.
教會從第一世紀裡面.
已經不是一個完美的群體.
保羅沒有說要逃避教會的問題.
沒有說教會沒有問題.
相反他只是要去知道.
教會根本就是充滿很多問題.

$^{681}$充滿很多軟弱的人.
所以他就是這樣去呼籲信徒.
要如何互相寬容.
彼此的接納.
因為我們都是蒙恩的罪人.
通常一說到要接納教會.
就是蒙恩的罪人.
就會有一種感受.
又是教會有問題.
所以大家就忍讓一下.
大家就當睜一隻眼閉一隻眼.
但不是這樣.
伊溫所說不是在說這件事.
因為第四章裡面.
他在說神不僅次下合一.
祂次下的是有恩賜.
所以祂不僅是要我們彼此接納.
是要彼此建立.
換句話來說.
就是神將我們放盡在教會裡.
不是要我們去麻木.
去忍大家的軟弱.
然後去逃避那些問題.
而是我們要去一同成長.
就是當有人偏離真理的時候.
我們需要提醒.
當有人軟弱跌倒的時候.
我們需要去扶持.
當有人受傷的時候.
我們需要安慰.
當有人犯罪的時候.
我們需要用愛心去挽回.
這個正正就是保羅在四章十六節裡說.
八節都要各按其職.
照著各體的功用去彼此輔助.
使身體漸漸增長.
在愛裡面建立自己.
所以當我們見到教會的問題的時候.
其實我們不如試一下.
不要先問的是.

$^{721}$這間教會有這麼多問題.
我還是要走了.
我們更加不如試一下問.
究竟上帝給了我們什麼恩賜.
讓我們可以一起和弟兄姊妹一同成長.
怎樣八節各按其職呢?.
在基督裡同心前行.
這個不是靠少數的領袖.
你們努力,你們做吧.
去完成.
而是靠著整個群體彼此的建立.
神不是呼召我們要去.
一來就是一個完全沒有問題的教會.
而呼召我們在這間教會裡.
一同慢慢成熟.
一同成長.
讓我們一同禱告.
因為你是師恩的主.
你看見我們需要群體.
你看見我們需要互相提醒.
叫我們一同成長.
但可能在群體當中.
我們有時候都會經歷失望.
有時候我們都很怕在人際關係當中.
但上帝求你在當中.
讓我們看見群體的重要性.
又在當中叫我們感到孤單.
感到失望的時候.
願上帝藉著群體去重新鼓勵我們.
求主祂繼續祝福和帶領.
祝福我們在靈福堂裡.
無論在粵語堂當中.
或者在三語我們一起的聯合裡.
叫我們一同看見上帝賜給我們的群體.
讓我們一同成長的群體.
奉主耶穌基督的聖名祈求.
阿們.
\newpage



\section{士師記 8:21}
\label{sec:032331}
\textbf{聖經人聖經事:人如何,力量也如何}~~~~黃克勤牧師博士
\newline
\newline
連結: \href{http://www.church.com.hk/acms/content.asp?site=cdc\&op=show\&type=product\&code=032331\&layout=sermon}{\texttt{032331}} ~~~~ 語音日期: 2026-06-18
\newline
\newline
\hyperref[sec:032330]{\small{< < < PREV SERMON < < <}}
~
\hyperref[sec:index_chronic]{\small{[返順時目]}}
~
\hyperref[sec:032332]{\small{> > > NEXT SERMON > > >}}
\newline
\newline
士師記 8:21
\newline
\begin{longtable}{cl}
\hline
\hline
章節 & 經文 (和合本修訂版)\\
\hline
8:21 & \begin{tabularx}{0.7\textwidth}{X} 西巴和撒慕拿說：「你自己起來殺我們吧！因為人如何，力量也如何。」基甸就起來，殺了西巴和撒慕拿，取了他們駱駝頸項上的月牙圈。 \end{tabularx} \\ \\
[1ex]
\hline
\hline
\end{longtable}
$^{1}$大家好,一齊看聖經.
聖經人聖經時在生命樹宣教網絡為大家提供.
今次繼續講機電.
今次的題目是「人如何,力量也是如何」.
是聖經裡面的一句說話.
我們看看「人如何,力量也是如何」.
看看它說什麼.
時詩第八章十三節.
「約阿斯的兒子機電由希烈斯坡從陣上回來」.
機電在希烈斯坡這個地方.
名字很有趣.
因為他一直追米電的領袖.
追到很遠的地方.
這些名字只出現一次.
都懶得考究它在哪裡.
不過應該是去到米電人比較深入的地方.
從陣上回來.
他抓了他們追趕的兩個領袖.
記得上一集說什麼呢?.
去到疏國這個地方跟疏國的人說.
其實是自己人.
就是我們三百人追趕米電人.
為我們全個以色列民爭戰.
上帝要祝福我們.
不過我們已經很餓又很累了.
所以你記得昨天說.
雖然疲乏仍要追趕.
雖然累了我們繼續趕.
要追趕.
可不可以給我們吃點餅.
誰知道疏國的人說.
要餅?你抓到他們了嗎?.
抓到才算吧.
不肯給.
其實疏國人真是頑皮一點.
本來就這樣一個路人很餓.
你都要給他吃點餅.
誰知何況是自己人.
當中可能部分不是自己人.
不過照地理位置.

$^{41}$是屬於他們的支派裡面的一部分.
是以色列人.
不過大他.
你抓到了嗎?這麼厲害.
抓到才算吧.
有時要嘲笑他們.
看你三百人追趕別人.
我怕你沒命回頭.
這些我說的.
基甸就說.
如果我抓到那兩個領袖回來.
我就會用荊棘做鞭來.
用刺棍打傷你.
你等我.
基甸就說要走了.
去到另一個地方.
就是比魯伊納.
那裡的人都是這樣嘲笑他們.
如果我回來.
我會燒了你這座樓.
拆毀你這座樓.
約阿斯的兒子基甸.
由希列斯坡從鎮上回來.
他們在那裡抓住梭葛.
梭葛這個貨攤是甚麼呢?.
貨攤是綠色字.
是我的解釋.
梭葛這個字.
其實就是貨攤.
我上次也解釋過.
那裡有很多攤檔賣東西.
有名的市集.
在那裡一檔一檔.
好像嘉年華那樣.
在那裡可以買到很多不同的東西.
去到那裡.
梭葛就抓住梭葛的一個少年人.
就問他.
去到見到那個少年人.
抓住他.

$^{81}$馬上把他拉過來.
如果我們今天好像拍戲那樣.
或者叫做話劇.
抓住他.
喂!小伙子.
就說給我聽.
你們梭葛的首領長老是誰?.
你們梭葛的首領小伙子.
講給我聽,長老首領什麼名?.
他就把首領長老的77個人的名字都寫了出來.
這個少年人就害怕.
好吧好吧.
照我記得的寫給你.
寫完不要漏.
他就寫了出來.
那就是說.
在當時來說也算是少有的.
因為一個少年人.
他讀過書才能寫出這麼多名字.
這些名字都不是容易記的.
而且很記性.
記得自己首領的名字.
多少個?.
77個.
記得了.
好了.
我解釋一下這個少年人是什麼意思.
少年人不是走過路邊.
就抓住一個人.
喂!小伙子.
說給我聽你們有哪個領袖的名字.
不一定是這樣的.
因為通常叫少年人.
在戰爭打仗裡面.
少年人其實是一些.
跟著一個戰士服侍他的人.
通常都是年輕很多.
就遞劍,遞食物,遞水.
幫他守.
所以其實就是說.

$^{121}$是他的助手.
打仗的助手.
不過去到臨陣交鋒.
真的出去衝殺的時候.
那些少年人會找二哥躲起來.
可能一班少年人在一起.
躲在那裡.
看著情勢.
受傷就出去救主人.
或者要什麼就遞給他.
這些是幫忙打仗的人.
很有可能.
抓住一個這樣的人.
就問他.
你領袖什麼名字?.
長老什麼名字?.
全部寫出來.
他真的寫了77個.
不過我要說的就是.
抓住.
我用螢光的.
不知為何.
我有少少覺得.
特別這裡寫是抓住.
其實有點不禮貌.
當然今天我不知道我說的話.
會不會太過將我今天的思想.
套在當日的場景裡面.
可能當日的場景.
這個抓住是沒有禮貌的.
不過今天抓住一個少年人.
嚇死他嗎?.
有點害怕.
咬他.
有點惡.
機電到了疏閣.
未到疏閣就抓到少年人.
這個少年人可能出來做什麼事.
或者他的主人猜探他做什麼事.
他就被他抓住.

$^{161}$抓到嚇得他半死.
去到疏閣就對那裡的人說.
你們疏閣從前譏笑我.
就說西巴和薩姆拿已經在你手裡了嗎?.
所以你叫我們將餅給你跟隨的人.
被罰的人.
是不是你已經抓到?.
抓到這麼厲害.
我就給你餅.
你都沒抓到就拿餅?.
不是吧?.
抓到才算吧.
你這樣譏笑我.
又不給餅.
我們繼續.
我們雖然肚餓疲乏.
繼續追趕.
現在西巴和薩姆拿在這裡.
你看我抓到他了.
於是就抓住那城內的長老.
第二點要說.
用野地的荊條和紫荊紮打疏閣人.
於是抓住他.
那七十七人是疏閣的長老,領袖.
綁住他們.
可能多數這樣紮打.
就用一棵樹綁住他.
背部露出來.
然後就用荊荊.
像泥田般掃他背部.
然後用樹上的針打他.
用刺激和荊棘條鞭他.
紮打.
看死.
笑我當日.
還拿個餅吃.
你們抓到嗎?.
這麼厲害.
捉一拿,笑一拿就打死他們.
沒有打死.

$^{201}$照他昨天所說.
只是打傷他們.
不過也很傷.
如果像泥田掃月般掃他背部.
被荊棘和刺棍打到.
很傷.
入肉的.
背部沒有肉.
是入骨的.
又去到比魯伊納.
比魯伊納的意思.
就是面對神.
亞國當日在比魯伊納.
和神摔跤面對面見到神.
他們就改地方名叫比魯伊納.
面對神.
其實我們時時都會面對神.
我們都有我們比魯伊納的地方.
在那裡面對神.
心裡有很多事情和神摔跤.
或者心裡有很多不忿.
或者心裡有很多難解的事.
很多憂悶.
很多心痛.
心碎的事.
在那裡面對神摔跤.
神你為何這樣.
神就說我是這樣.
於是大家摔跤.
心靈上摔跤.
比魯伊納就是面對神.
有你的比魯伊納.
後來那裡發展成為一個地方.
很多人在那裡住.
而且有一個樓.
比魯伊納樓.
基田就說.
你給我們餅我們很餓.
又拿餅.
他們就像疏割那樣.

$^{241}$拿餅?捉到了嗎?.
捉到了我領袖.
我就給你餅.
你還沒捉到.
其實你追著去.
都不知道有沒有命回來.
不要浪費我的餅.
你們先追吧.
基田很生氣.
他就說如果我捉到他們回來.
我會拆了你們比魯伊納的樓.
這個樓其實是一個塔.
即是說比魯伊納發展了一個塔.
塔是用來做甚麼呢?.
高一點.
可以在裡面看得遠一點.
和變成一個中心點.
很遠就可以看到那個地方.
如果在頂部放煙.
更加遠就可以看到你.
如果晚上放火.
走夜路的人都知道.
那裡叫做比魯伊納.
就可以去到.
所以是一個比較發展的地方.
他就拆了比魯伊納的塔.
就是他上一次說的.
我現在這麼辛苦為以色列人打仗.
都為你們去打仗.
去追擊米電的兩個首領.
你竟然餅又不給.
還笑我.
我回來就會拆了你的塔.
拆了你的樓.
他真的回來拆了比魯伊納的塔.
不過還做了一件事.
就是殺了那城裡的人.
殺了那城裡的人.
是上次沒有說的.
沒有說你不給我餅.

$^{281}$笑我.
好.
算了.
我們餓著.
累著.
繼續追.
不過我回來.
我找你算帳.
我會拆了你的樓.
這是上次說的.
但沒有說殺了城裡的人.
這次機電回來.
餓起來.
一不做二不休.
拆了你的城樓.
拆了你的塔.
殺了城裡的人.
城裡的人可能不是以色列人.
但可能是自己同胞以色列人.
可能殺紅了眼.
殺了城裡的人.
殺人.
第二點.
我想說的第一第二點.
捉住一個少年人.
搖他.
喂.
把名字寫出來.
第二點要說的.
就是他鞭打那些蘇國人.
又拆了比魯伊納的塔.
還殺了城裡的人.
在這裡我個人.
我怕我用了我今日.
我個人的思想套入當時的場景.
可能機電這樣做.
在當日來說是沒有什麼問題的.
在當日的文化與場景裡面.
是沒有什麼問題的.
被羈絲.

$^{321}$所以被殺死是很通常的.
被責打是很通常的.
但今天.
你想想.
給我一些餅.
不給.
你打贏了回來我才給.
才算.
你一定贏到的.
你可能沒命回來.
米電的王那麼容易追.
你只有三百人就餓成這樣.
哎呀.
這樣笑你真的很生氣.
講到明我回來我會鞭打你.
我回來我會拆了你那座樓.
那塔.
但都不用殺裡面的人.
捉住少年人不用惡到搖他.
寫了名字出來.
七十七個全部.
找當日那幾個.
跟你說的那幾個夠了.
有七十幾個長老.
可能老人家都不知道發生什麼事.
就七十七人被你捉住走.
鞭一輪.
很傷的.
很痛的.
流血的.
入骨的.
被那些人勾到.
哎呀.
很慘.
我個人在這裡第一第二點.
我是覺得機電.
他可能殺紅了眼.
而且過分嬲怒了.
他沒給餅而已.
不用打他.

$^{361}$沒給餅而已.
打倒他那座樓算了.
不用殺裡面的人.
為什麼不用?.
想餓死我?.
現在我要報仇.
殺他.
我覺得可能殺紅了眼.
可能是他嬲氣很厲害.
我個人覺得.
捉住那少年人搖.
其實我覺得很不禮貌.
很不尊重.
兼夾有些惡.
如果你說機電有些什麼.
如果我個人覺得.
是不夠好的.
就是他在這件事上.
第一和第二點.
他是惡了一點.
很惡.
變成惡死.
沒給餅他又打又殺.
我覺得是有點惡了.
過多了.
過分了.
嬲過分了.
火遮了眼.
令我反省.
會不會有時都是這樣.
當自己功勞很大的時候.
很辛苦做事的時候.
會對那些不肯付出.
沒有付出的人.
是很不開心很惡的.
會想要他的命.
收入都很大.
不過要不要又打又殺.
又捉那麼厲害呢.
好像做多了一點.

$^{401}$我覺得機電在這裡.
有些少是他生命裡.
曾經做過的過分.
我想叫錯了.
或者過分了的憤怒.
造成他一個缺點.
或者叫污點.
是我個人的看法.
可能我太多用今天的想法.
套入當日.
當日可能他沒事.
因為你看到.
上下文沒有人罵他.
沒有人說他沒有.
上帝沒有責備他.
哎呀不如你這樣做.
我個人看.
就有少少覺得是這樣.
好了.
另一件事.
第十八節.
機電問西巴和薩姆拿雪.
那些人搞定了.
打打殺殺.
當日在他博山殺了人.
是什麼樣式.
當時他們聚集在他博山.
然後這兩個米電領袖.
都有份攻打他博山.
殺了以色列人.
他就問他們.
當日在他博山殺我們以色列人.
他們是什麼樣式.
樣式的意思我看過了.
就是他們像什麼.
他們是什麼樣子.
是什麼人.
就是說.
你一看就知道.
這些是平民的裝束.

$^{441}$一看這個是達官貴人.
貴就的裝束.
一看這個是將軍.
這個是普通士兵.
他們是什麼級別.
是什麼樣子的.
這兩個領袖就回答他.
大家都很直接.
他們很像你.
他們是什麼樣子.
像你.
很厲害.
領袖來的.
各人都有王子的樣式.
看上去就知道他是尊貴的人物.
都是重要的人.
當日我們在他博山上.
殺的都像你一樣.
這麼有型,這麼尊貴的人.
像王子,像領袖的.
這句話是老實還是說反話?.
你在他博山殺了什麼人?.
他們是怎樣的?.
殺了很多人.
他們的裝束樣貌是怎樣的?.
看上去就像你一樣.
都是這麼英俊,這麼厲害.
都是做領袖級的.
究竟是真是假不知道.
不過這樣說.
這兩個人就害死自己了.
基建說.
他們是我同母的弟兄.
是不是我親家的弟兄呢?.
不是的.
他們應該是我的同胞.
都是我們的支派.
因此我要指著永生的耶和華起誓.
我一定會這樣做.
向永生的耶和華起誓.

$^{481}$誓言是不能修改的.
你們從前虐存留他們的性命.
我如今就不殺了你們.
如果你當日留守沒有殺他們.
我就會今天都不殺你.
但你當日殺死了他們.
所以我今天就要殺你了.
基建這樣說.
於是對他的長子奕帖說.
你起來殺他們.
跟他的大兒子說.
奕帖.
我的大兒子.
你殺了這兩個領袖.
叫你殺了我的兄弟.
現在我叫我兒子殺你.
如果你當日不是殺了我的兄弟.
我今天可能會放你.
不過.
既然你殺了他們.
我現在向永生神起誓.
你必定要死.
誰知奕帖.
耶和華.
耶和華在中文譯為奕帖.
耶和華的意思是風俗.
他的長子奕帖.
奕帖的意思是風俗.
也就是說.
奕帖出生的時候.
其實基建家庭不差.
有錢的.
你怎知道.
因為他爸爸有一段時間是做什麼?.
是離棄了上帝.
在家裡建了巴力的祭壇.
他把壇拆了.
然後把旁邊的亞瑟拉木像拆掉.
然後做柴.
於是就說.

$^{521}$你怎麼這樣做?.
然後把爸爸七歲公牛.
放在壇上.
然後那些人就說.
喂.
約亞斯.
你兒子做的.
把他交出來.
約亞斯就說.
如果你們保護巴力.
如果他真是神.
讓巴力和他爭辯.
讓巴力算他帳.
那些人說是.
沒有聲音.
於是就算了.
嬲嬲地走了.
約亞斯就回覆信仰.
我痊癒了.
他做了這件事情.
大家記得嗎?.
機電就叫耶路巴力.
讓巴力和他爭辯.
機電就叫耶路巴力.
機電在他爸爸任下生活.
爸爸做這些拜巴力的事情.
一定有很多錢.
因此機電的兒子.
益帖出生的時候.
應該生活環境都很好.
所以見到益帖出生.
就給他取了一個名字.
叫Yetha.
Yetha就是風燭.
兒子你出生就金鎖匙含著.
厲害了.
你就有很多錢用了.
誰知道益帖因為是童子.
Nah.
Nah這個字就是我們剛才說.

$^{561}$找到一個疏覺的少年人.
相同的字.
這裡譯作童子.
那裡就譯作少年人.
少年人我特別說了.
少年人很多時候.
是幫勇士拿兵器.
服侍他.
幫他守的人.
是年輕比較小的.
幫忙一個好像服侍他的工人.
不過是童工.
Nah意思譯作童子.
即是他還小.
小時候害怕.
其實Nah是做什麼呢?.
可能跟著爸爸去打仗.
跟著幫忙.
遞東西給爸爸.
跟著幫忙.
爸爸小心那邊有人來.
之類的.
他就幫忙.
小時候叫童子.
頂多叫少年人.
這個益帖.
益帖就害怕.
Yet.
Yetpha.
Yetpha是益.
Yet.
Yetpha就害怕.
不敢拔刀.
雖然他都拿著刀.
不過不敢拔刀.
他真的捅死一個人.
他鋸他的頸.
他不敢.
哎呀爸爸好害怕.
出不到手.

$^{601}$拔不到刀.
拿著刀很震.
於是.
西巴和薩姆娜就說.
紀典不如你自己殺我們吧.
這個.
對不起我用粗俗一點.
這個小子哪裡行.
他還這麼小.
因為第三點.
人如何力量也是如何.
就是我第三點要說的.
他這麼小.
力量很有限.
叫他割我的頸.
他都不夠力.
人如何力量也如何.
不如你自己這麼好.
你自己殺我們吧.
我們的命在這裡.
紀典二話不說.
就起來.
殺了西巴和薩姆娜.
奪獲他們樂陀.
項上大的月牙圈.
原來他們每一個人.
自然都有自己的坐騎.
他們不是騎馬.
因為米電是沙漠遊牧民族.
就是樂陀.
樂陀頂著很珍貴的東西.
每人有一隻野獸的牙.
叫月牙圈.
即是彎的.
綁在那裡.
可能是吉祥.
或是表示尊貴階級.
彎的牙很多時候.
都是野豬的牙.
不會象牙嗎?.

$^{641}$應該不會.
野豬的牙比較多.
野豬的牙比較大.
獲得奪獲他們樂陀.
項上大的月牙圈.
即是獲得他們樂陀.
項上的月牙圈.
即是已經征服了他們.
而且獲得他們的樂陀.
就殺死他們.
人如何力量也是如何.
所以你叫一個小朋友.
做一些比他能力化的事情.
是不可以的.
我勉強要求.
做我的兒子.
可以的.
你拿著刀.
我捉住你的手.
我們一起去殺死他.
不要嚇壞小朋友.
他還未到.
你是太早pre-mature.
叫他做這些事情.
其實未必好的.
你自己做吧.
這個小朋友沒有力量.
你這麼大.
你才有這樣的能力.
因此當我們一路成長的時候.
會有跟著我們成長而出來的力量.
用這個力量才可以做到事.
否則做不到要做的事.
當然另一句.
力量也自如何.
另一句也出現過.
在聖經裡面.
你日子如何.
力量也必如何.
這是什麼意思.

$^{681}$隨著你年日的增長.
你的力量會一路增加.
當然.
年老了.
力量低了.
年老了才厲害.
老奸巨猾.
哪需要自己動手.
年老了.
但他有自己的力量.
可能叫別人做.
可能有計謀.
可能更厲害.
日子如何.
你的力量也會相同地增多.
日子一路增多.
你的力量也會增多.
這是祝福.
這是事實.
人如何.
力量也如何.
人去到哪裡.
力量去到哪裡.
所以如果你要大力量.
你這個人要大量一點.
力量會大一點.
你的訓練.
你的鍛鍊要好一點.
於是力量會好一點.
如果你這個人不好.
力量就會差.
譬如我這樣.
現在身體差很多.
有糖尿.
心臟有事.
力量很差.
瘦了.
沒有肌肉.
為甚麼.
要還債.

$^{721}$過去二三十年都沒做運動.
你以前不是這麼高.
不是打籃球嗎.
以前打籃球很好.
做中鋒.
打中鋒.
很快.
人家一交球.
我在中途空中.
就拿了他的球.
然後就上了.
人家一跳籃.
我的手要高過他.
擋住他的球.
很快.
很叻.
很高.
不過上次打籃球.
甚麼時候.
都不是多久.
四十年前.
而我記得.
四十年.
我是欠了四十年的運動量.
坐四十年的運動量.
於是今天身體不好.
糖尿了.
心臟不好.
這樣.
這個人的力量如何.
人幾十年沒做運動.
力量就衰微.
所以要把自己這個人.
做好一點.
你的靈性不好.
怎樣有屬靈的能力做事.
聖經都沒甚麼看.
我要教聖經.
做不到.
書都沒甚麼看.

$^{761}$怎樣有新的思想.
世界都沒望客.
怎樣有突破.
所以人如何.
力量也是如何.
所以要把自己這個人.
鍛鍊好.
做好一點.
機電二話不說.
殺了他.
我想和大家看看.
這個就是今天想和大家講的.
人如何.
力量也是如何.
講兩點.
結束.
總結.
我覺得機電去到這個位.
最後勝利的時候.
他是惡了.
他的仇恨.
是大了.
為甚麼.
人如何.
他可能沒把自己那些.
溫柔那些.
那些仇恨.
沒過濾也好.
變得很躁.
很惡.
其實如果他沒那麼躁.
沒那麼惡.
我覺得更加好.
又沒壞.
起碼我見到沒有神說.
你做得過份了.
或者沒有當時的人起來.
就說機電你不應該這樣做.
就殺了城內的人.
他沒給你餅.

$^{801}$不用死吧.
我現在覺得是這樣.
不過當時可能沒事.
覺得他這樣做是對的.
不過我覺得.
好像有點過份惡了.
他沒把這個人.
自己修養得好.
在仇恨裡面的時候.
當然你會見到.
殺了他兩個.
叫我兒子殺.
其實讓他練一下刀.
練一下膽.
可能是這樣的.
兒子將來也要像我這樣.
夠膽.
其實他兒子可能是文人.
不過被你帶到這樣.
現在就沒辦法了.
要這些.
我覺得我們都要做好.
自己個人的修養.
在神面前的等候.
成為一個合神心意的人.
然後我們的力量.
也要同樣合神的心意.
今天跟大家分享到這裡.
我們彼此祝福.
彼此勸勉.
同心禱告.
求主教導我們.
在基地這個位置裡.
他帶著三百人.
非常辛苦.
餓得不得了.
但是被人家譏笑.
他完成任務.
回來他將人鞭打.
將人的城樓拆了.

$^{841}$高塔拆了.
還殺死裡面的人.
還抓住少年人.
其實這些我們今天看.
個人看未必需要.
可能有些不必.
不過他餓了.
那個人餓死了.
主要求你自己幫助.
讓我們有更好的.
內心的修為.
讓我們有更加溫和的處理.
可能將來更好.
主要求你自己帶領.
教導我們.
誰聽祈禱.
奉耶穌基督名求.
阿們.
\newpage



\section{路加福音 19:1-10}
\label{sec:032332}
\textbf{從大起大落到大喜大樂的人生}~~~~麥偉祺牧師
\newline
\newline
連結: \href{http://www.church.com.hk/acms/content.asp?site=cdc\&op=show\&type=product\&code=032332\&layout=sermon}{\texttt{032332}} ~~~~ 語音日期: 2026-06-21
\newline
\newline
\hyperref[sec:032331]{\small{< < < PREV SERMON < < <}}
~
\hyperref[sec:index_chronic]{\small{[返順時目]}}
~
\hyperref[sec:032333]{\small{> > > NEXT SERMON > > >}}
\newline
\newline
路加福音 19:1-10
\newline
\begin{longtable}{cl}
\hline
\hline
章節 & 經文 (和合本修訂版)\\
\hline
19:1 & \begin{tabularx}{0.7\textwidth}{X} 耶穌進了耶利哥，要從那裡經過。 \end{tabularx} \\ \\ \relax
19:2 & \begin{tabularx}{0.7\textwidth}{X} 有一個人名叫撒該，作稅吏長，是個財主。 \end{tabularx} \\ \\ \relax
19:3 & \begin{tabularx}{0.7\textwidth}{X} 他要看看耶穌是怎樣的人，只因人多，他的身材又矮，所以看不見。 \end{tabularx} \\ \\ \relax
19:4 & \begin{tabularx}{0.7\textwidth}{X} 於是他跑到前頭，爬上桑樹，要看耶穌，因為耶穌要從那裡經過。 \end{tabularx} \\ \\ \relax
19:5 & \begin{tabularx}{0.7\textwidth}{X} 耶穌到了那裡，抬頭一看，對他說：「撒該，快下來！今天我必須住在你家裡。」 \end{tabularx} \\ \\ \relax
19:6 & \begin{tabularx}{0.7\textwidth}{X} 他就急忙下來，歡歡喜喜地接待耶穌。 \end{tabularx} \\ \\ \relax
19:7 & \begin{tabularx}{0.7\textwidth}{X} 眾人看見，都私下議論說：「他竟然到罪人家裡去住宿。」 \end{tabularx} \\ \\ \relax
19:8 & \begin{tabularx}{0.7\textwidth}{X} 撒該站著對主說：「主啊，我把所有的一半給窮人；我若勒索了誰，就還他四倍。」 \end{tabularx} \\ \\ \relax
19:9 & \begin{tabularx}{0.7\textwidth}{X} 耶穌對他說：「今天救恩到了這家，因為他也是亞伯拉罕的子孫。 \end{tabularx} \\ \\ \relax
19:10 & \begin{tabularx}{0.7\textwidth}{X} 人子來是要尋找和拯救失喪的人。」 \end{tabularx} \\ \\
[1ex]
\hline
\hline
\end{longtable}
$^{1}$但我們今天要講道的時候.
一定要知道聖經怎樣教我們.
從大喜大樂到大喜大樂的人生當中.
剛才我們選讀的經文是在《路加福音》第十九章一到十節.
是講一個人的.
人生當中有時我們都知道.
有時我們行差踏錯.
好像文森弟兄說的.
初初不知道自己做錯.
覺得是理所當然的.
出去風花雪月 交朋結友.
有時就是這樣.
有個校長說錯一句說話.
在巴士上說一句不知是不是髒話.
就髒話了.
就已經搞到要馬上離職.
這些很熟悉的個案.
全程哄動.
本來是高薪厚職做校長的.
現在變成炒魷魚.
還要即時去炒魷魚.
不可以先做暑假.
真的很麻煩.
亦有很多不同的人都有很多不同的際遇.
可能有時一失足成千古恨.
總之是這樣就有種子無乞兒.
但有時又很難說.
幸福不是必然的事情.
好像馮福遠.
大家以為都安居樂業.
一個中產家庭.
誰不知一個煙頭.
一些不合規格的建築材料.
或者一些不守規矩的建築箱.
或者是消防的設備.
搞到這麼多棟樓都出事.
本來說可以有些重建或者原地的選址等等.
現在可能亦都另外一些的說法.
沒人知道下一步怎樣.
現在已經開始有些人被追究責任.

$^{41}$開始ICAC已經告一些工程界人士.
但除了工程界之外.
還有沒有其他更多的一些人.
是需要去負責任的呢.
最重要不只是追究責任.
而是怎樣去改善那個措施.
以至其他類似的建築.
或者類似的工程.
都能夠有多一點的保障.
我們都想搞好那個制度.
當然我們亦都知道.
前線很多的工作人員都好辛苦.
有些消防員都是無私的奉獻.
為了救火 建築自己的生命.
所以幸福不是必然.
所以今天我們能夠坐到這裡.
能夠有機會去分享上帝的一些工作.
其實很感恩.
是不是 弟兄姊妹.
很感恩.
但是我們今天講一個經文.
耶利哥裡面的地方.
有一個故事發生.
耶利哥海洋地球上面.
最低的一個水平線.
它是低過海拔.
地中海是負258米.
即是地中海在這裡.
它在下面那裡.
那些水就水向低流.
所以所有的水落到地中海.
即是落到這個死海之後.
是不會走出來的.
那個耶利哥這個城市.
就在耶路撒冷亞東北的23公里.
這也是當時我們知道聖經裡面.
第一個以色列人進入的城邑.
就是由約旦過河進入了耶利哥.
這是一個很重要的歷史故事.
歷史是會重寫的.

$^{81}$即是最初進入耶利哥.
現在耶穌最後這個站又去京興耶利哥.
這個地方我們叫做「低處未算低」.
地理上的低是一個肯定的.
那裡的水各方面的氣壓等等都不同.
但那裡有一個人是更加低.
就是有一個叫做薩蓋.
薩蓋是個甚麼人來的.
是個稅吏長.
是個財主.
這幾個這樣的稱呼就很明顯了.
因為剛才講到那個地理環境的時候.
是約旦過來這一邊.
以色列是必經之路來的.
即是這條是一個很豐富的運輸的中樞來的.
那你說那裡做稅官發不發達.
還是做關長.
是不是很發達.
所以這個薩蓋做了一個關長.
做了一個稅吏長的時候.
他是管所有耶利哥裡面其他的稅吏.
即是這個是很高的一個職位來的.
是一個很有錢的人來的.
但在當時的猶太文化來講.
做稅吏是不受歡迎的.
因為當時以色列是被羅馬統治.
你所有的納稅是納給誰.
納給上級領導.
納給羅馬.
所以在猶太人來講.
覺得你的稅吏是賣國賊來的.
是我們那些的叛徒來的.
還有當時的稅制來講是怎樣的呢.
就是誰想做稅吏.
誰想做關長.
你就開標.
投標.
你就承包到這個地區裡面.
你可以上繳多少.
那你只要你的價高者得.

$^{121}$那你說這些多少的圍標.
多少的既得利益者在裡面中間做事呢.
但薩戈是一個很有能力或者智慧的人來的.
他是做到這個稅吏長.
管理整個耶利哥的所有的稅務.
所以這個人來講.
在同胞來講.
就覺得這個人是最討人厭的了.
討人厭到相當討厭的了.
行不行啊.
我們現在的老人家.
我們知道的.
老人家最近要交醫療保險那些.
看醫生又貴了那些.
我當然罵那個了.
搞到我們現在的看醫生那麼貴.
還要給藥錢.
就是當我們要有一些負擔的時候.
我們當然覺得不公道的.
是不是啊.
同意嗎.
很多時候做這個稅吏或者稅吏長的時候.
因為他要去交稅給上面的話.
有時他都要做一些不光彩的手段.
他才能夠令到那些下面的百姓去服.
無論是欺詐也好,訛詐也好,威逼也好,利誘也好.
都是做一些不是很光彩的事情.
所以基本來講.
就是說做得稅吏,做得稅吏長.
是會斷綠親的.
同意嗎.
人生就是這樣.
可能你職位上是很高.
但是品格上別人覺得你很低.
很低俗.
明白我的意思嗎.
我們都是一樣的.
有時我們.
就好像文森剛才講的那些故事那樣.
他以前做銀行的.

$^{161}$我們覺得都是受人尊重的.
但是我們背後知道.
那些銀行裡面很多的生活習慣都是很萎靡的.
是不是.
很腐敗的.
就是兩種不同的那個落差當中.
人生就是這樣.
在這個時候他說.
這個薩埃很想去見耶穌.
很想去見耶穌.
為什麼呢.
其實他不是第一次聽到耶穌的故事.
做得稅吏長的.
不可能沒有那些情報機構.
跟他報那些附近發生的事情.
是不是.
耶穌剛剛經過耶利哥.
就醫了一個盲仔.
那個盲仔就說病得醫治.
是不是很神奇.
當時就說這是神蹟來的.
如果不是上帝與這個耶穌同在的話.
不可能這個盲仔可以得醫治.
還有過程當中也醫了很多其他人.
全都趕鬼治病.
耶穌的工作對薩埃來說是鳥語紙掌的.
但是他就真的很想知道.
這個是什麼人來呢.
他就很想見耶穌.
看看耶穌是誰.
還有另一個可能性就是.
因為耶穌的門徒當中.
也有一個叫做馬太.
馬太就是利美.
以前也是做加拔隆的稅務的稅吏.
同行當中也有一個是跟耶穌的.
會不會有些消息傳給他.
是呀 你以前的拍檔也曾經有一個跟了耶穌.
所以我們也知道.
每一個人來想見耶穌.

$^{201}$也有些人介紹的.
會不會突然之間.
今天來這裡當中.
有沒有人介紹你來的.
誰介紹人來這裡的.
有誰 你有介紹誰來.
你旁邊那個.
是呀 是呀.
他也是去介紹你來的.
為什麼他信耶穌這麼開心.
這樣來看看.
我們今天信耶穌的人開心嗎.
當然開心.
這裡有平安有喜樂.
所以很多時候.
有些人是做了預公的.
我們都相信薩埃能夠想見耶穌.
不是一件偶然的事情.
也是他想了很久.
還這麼近.
耶穌專程來到耶利哥.
當然想見他.
但經文講到很有趣的地方.
只因為他人多.
他又是身材矮小.
所以見不到耶穌.
人矮是沒有辦法的.
前面有人擋著你.
我們看放煙花.
看其他的表演.
人矮就要騎駁馬.
如果有人疼愛你.
就揹你上膊頭.
我跟孫子去看表演也是這樣.
我就揹他.
是不是很開心.
但這個矮的薩埃竟然怎樣.
他很努力很努力的希望見耶穌.
但卻是不能.
是沒有能力見耶穌.

$^{241}$不是他想見就見.
因為前面太多人擋著他.
人們見到他不想讓位給他.
你可以想像當中.
他被人排斥.
那種無力感.
那種心酸.
你明白我的意思嗎.
江湖地位你有.
但去到這些場合沒有人給你面子.
慘不慘.
都很慘.
剛才我在街口已經見到弟兄姊妹.
姊妹跟我打招呼.
然後他問你認識他們嗎.
我說是的我認識他們.
很多年了.
我們合作了十年.
八九年了.
都見得了認得了.
你明白我弟兄姊妹嗎.
但竟然因為他的身份.
因為做稅吏稅吏長.
而人們排斥他.
甚至不讓他走到前面.
其實那種被排斥的感覺很強烈.
正如文森所說.
他做到一些囚犯的時候.
已經沒有社會的關係.
那些社交網絡.
那些朋輩不跟你接觸了.
甚至不會讓你.
是一個無力感很強的.
無論他的身形.
或者他的工作.
或者他違反了上一代對他的期望.
最終他只剩下錢.
聖經形容他財主有錢.
但在這時候他矮仔多計.
他說我命由我不由天.

$^{281}$雖然你們擋著我不讓我走.
但我可以急步走上前面.
走上前面對一個成年人來說.
是一個很失禮的行動.
我們經常都覺得要穩定.
自由淡定.
爺爺在大廳大家不用怕.
有什麼可能你跑的呢.
一跑就證明你沒有實力.
你一是奴隸才會走來走去.
一是你沒有信心.
沒有那種道行才會走來走去.
應該是剩餘生.
但這時候他沒有辦法.
他都願意走去前面.
不單止走去前面.
他還知道走去前面別人都會擋你.
還怎樣.
他爬上一棵桑樹.
一些矮的樹木能夠爬上去.
摘一些梳果來.
那些是桑子果的樹木.
他說去到那棵樹上面.
就可以看見耶穌要見耶穌.
因為這是耶穌必經之路.
他違反了人的常態.
違反了他們的禮儀.
違反了那些人對他的期望.
他甚至要冒一些風險.
爬樹裡面可能會勾到衣服褲.
弄傷手腳.
會不會啊.
我們都見到小朋友爬樹.
我們都害怕.
行不行啊.
你讓不讓他爬啊你.
當然叫他不要爬.
但是這個撒該就是這樣.
不怕羞恥.
不怕被人辱罵.

$^{321}$因為反正都是乳慣了.
於是他爬上去.
爬上去的時候.
很有趣的是.
他很想去見耶穌.
剛才說的.
他的心態是想見耶穌.
誰不知當他見到耶穌的時候.
是耶穌到了那裡.
然後耶穌是怎樣.
抬頭一看.
兩個動作.
耶穌是突然停在那裡.
掃描上去上面.
然後就找到撒該.
你想見耶穌嗎.
原來是耶穌來見你.
行不行啊.
我們都想來找耶穌.
其實你沒進來之前.
耶穌都見到你了.
這個就是動與被動的關係.
你想見耶穌嗎.
你們尋找的必然是尋見.
因為神就會讓我們被尋見.
並且神是主動來尋見他.
耶穌到那裡見到.
抬頭一看.
就跟他說.
撒該快下來.
今天我必須住在你家裡.
沒有人介紹撒該給耶穌認識.
但原來耶穌一早就知道.
誰是撒該.
也知道他的樣子是怎樣.
也知道他多高多矮.
多肥多瘦.
他能夠叫得出撒該這個名.
Amen.
等於我們每一個人的名字.

$^{361}$上帝都知道.
祂記掛著我們每一個人.
然後祂說.
你下來吧.
你不用在上面.
下來之後.
祂說今天我必須住在你的家裡.
我今天要跟你一起生活.
這個住在你家裡.
有一句相似的經文出現過.
就是在耶穌最後的十字架上.
跟囚犯說.
今天我要和你在樂園裡.
是同一個寫法.
就是說能夠和耶穌一起生活的話.
就是樂園.
Amen.
剛才文心說到.
在裡面的時候.
竟然可以和上帝對話.
就算坐牢.
心裡竟然得到自由.
因為耶穌說我要和你一起.
你就在樂園當中.
而耶穌說我要在你家裡住下來.
即是怎樣.
我要去你家裡作客.
作客是一件什麼事情.
今天我們經常北上旅遊.
我們洗兩口水已經是住五星級酒店.
吃一些豪華的自助餐.
很便宜的.
九十八元都可以做.
但是在以前二千年前的時候.
作客.
即是那些旅館.
條件是沒有現在那麼高級的.
行不行.
作客的原文來說.
就是說我要屈就.

$^{401}$即是去到別人家裡.
你會不會住別人的主人房.
你會不會住別人的套房.
你說隨便吧.
綠架床我們那時是帆布床.
或者摺床.
或者打地鋪都不一定.
我們那時很多時候.
鄉下來的人都是這樣.
有什麼就住什麼.
你明白我的意思嗎.
耶穌說我要去你家裡居住.
其實說到耶穌是要來到你家裡屈就.
耶穌是很願意來到我們人的當中.
耶穌來到是尋找拯救失常的人.
祂道成肉身來到我們的世界裡.
祂不介意和我們一起生活.
屈不屈就.
頗屈就的.
並且祂說要用另一個方法.
就是聖靈要住在什麼.
住在門徒的身體裡.
豈不知我們的身體就是聖靈的電.
你想想如果上帝住在我這個身體裡.
屈不屈就.
又三高又糖尿又血壓高.
你明白我的意思嗎.
其實上帝住在我們當中是很屈就的.
不過祂很樂意和我們一起生活.
祂說今天我必須住在你的家中.
於是撒戈就立刻急急下來.
很歡喜快樂.
歡歡喜喜的招待耶穌.
接待耶穌.
讓耶穌可以成為他家中的一份子.
在當時的文化能夠住在他家裡.
就成為家中的一份子.
受他保護.
是不是.
行不行.

$^{441}$我們都明白.
不是很熟悉的都不會留縮.
但是耶穌會成為他家裡的一份子.
很開心.
和耶穌接觸就是這樣.
很簡單.
好像耶穌又沒有說過話.
又沒有叫他信耶穌.
又沒有叫他做決志.
但是在這個時候.
有些人就問.
為什麼耶穌去罪人家裡住宿呢.
所有人都知道這個撒戈是個討人厭的人.
是個稅吏.
是無良的官府的人.
也是賣國賊.
是我們所認識的當中的罪人.
和他一起去收債的都是罪人.
那些都是不見得綱的東西.
但是為什麼耶穌應該是個聖人.
為什麼不去住一些高尚的地方.
要住在這些敗類裡面呢.
是不是耶穌是這樣.
不知道的呢.
是不是這麼糊塗.
有些人就澀澀的施言.
在這個時候.
耶穌是不解釋的.
任何人罵耶穌.
耶穌是不會和你對話.
答嘴的.
不過在這個時候.
事實就勝於什麼.
雄辯.
當撒戈和耶穌吃完飯之後.
在吃飯的時候.
他感覺自己被接納.
感覺耶穌那種公平公開公正.
感到耶穌那個慈悲憐憫.
竟然他受到耶穌的感動.

$^{481}$然後站起來對耶穌說.
他現在稱耶穌為主.
因為他知道生命主權.
不再在他自己掌握當中.
而是有個生命的主和他接觸了.
當你的生命和生命之主接觸的時候.
你的生命就會被改變.
他站起來說什麼呢.
他說主啊.
我會將我所有的一半分給窮人.
我如果勒索過誰.
我就會還他四倍.
這個時候我們覺得很奇怪.
在撒戈來說應該怎樣.
每樣事情想什麼都是想錢.
竟然今天他將他最緊張的事情.
視錢財如糞土.
是不是很違反原理.
很違反常理.
但他知道當他認識這個生命的主的時候.
他的眼目不再在錢那裡.
而是怎樣去能夠進入上帝的國裡面.
將他一半分給窮人.
這個是在說他的價值觀.
和他生命的意義已經重整了.
他用上帝那個天國的觀念去活在人間.
神就是要我們有一個知足的心.
我們今天夠不夠用.
我們都夠的,我們都相當富裕.
在香港人來說不錯的了.
比起其他貧窮國家我們已經是很富有的人.
我們每一天都有一個知足的生活.
並且我們要更加懂得跟別人分享.
分享我們所擁有的.
同不同意嗎?.
當我們能夠有奉獻的能力.
或者有其他做慈惠的能力的時候.
我們是不要吝嗇.
這個是在說上帝的慈愛.
他見到耶穌的慈愛.

$^{521}$他就學到耶穌的慈愛.
將他所有的一半分給窮人.
給他更加有需要的人.
這是一種慈悲的行為.
第二個就是說.
如果我勒索過誰.
我還他四倍.
這是出埃及記裡面二十二章一節裡面.
對一些不公平的交易裡面的賠償原則.
一賠四.
因為如果你拿了人家的牛.
那隻牛會生小牛的.
拿了人家的羊會生小羊的.
所以你拿了人家一隻的話.
要賠四隻給人家.
所以你就不要亂去斂財.
應該拿的就拿.
不應該拿的就不要拿.
還四倍.
這是在說什麼?.
公義的原則.
一方面他有憐憫.
另一方面有什麼?.
公義.
原來當你和耶穌接觸了一餐飯那麼多.
被接納之後.
你是會將自己的價值觀改變了.
有回上帝那種的屬性.
公義和憐憫在當中.
不再敲詐人不得已.
還要將自己曾經敲詐過的.
還回四倍.
將功贖罪.
在這個時候.
我們都知道.
在外國來說.
很多的富翁都會將他自己很重要的財產.
去成為一些基金.
或者做一些善事.
巴菲特,蓋茨這些都是.

$^{561}$每年都捐幾十個億美金去做這些慈惠工作.
幫一些貧窮國家去做一些科技上的發展.
或者一些學習上的求學.
我不多說了.
其實西方比我們中國文化會看得闊一點.
金錢觀是不同的.
香港可能都在學習.
都在做一些慈惠的事情.
我們都知道在上次的大埔事件裡面.
香港人的捐款都是相當多的.
本來愛心滿東華就變成慈心滿東華.
那個善款都是相當多.
所以我們要明白.
有一些事情發生都是去刺激我們.
應不應該做一些比較合宜的事情.
有公義和有憐憫的事情.
在這個時候.
在殺戒的改變裡面.
我們知道人是有很多不同的訴求的.
起碼我們開始有生存的訴求.
接著有些安全的訴求.
然後社會上的訴求.
更加重要有沒有尊嚴.
然後有沒有自我實現.
或者超越我們自我的一些實現.
這個就是每人生的向度.
不是在說吃多少穿多少花多少錢.
而是我們找不找到那個地位.
安全感和那個自我的超越.
現在在這個殺戒裡面真的找到了.
他已經超越了他自己原來的目光.
當然在這個過程當中.
耶利哥的人最開心了.
那些窮人馬上有個大富翁.
將他的身家分一半給我們.
如果是李嘉誠先生.
將他分一半給我們香港人.
每人起碼拿20萬.
發達了.
你明白我的意思嗎.

$^{601}$我比他還要窮.
耶利哥的貧窮人真的最開心了.
馬上有即時的津貼.
平常的老百姓都開心.
為什麼呢.
因為文化變了之後.
大家沒那麼辛苦.
消費能力高了.
或者其他方面都會提升了.
而那些稅吏都開心了.
為什麼.
如果耶利哥的總稅吏長.
這樣改邪歸正的話.
我們做他下面跟班的.
都不用背那麼多費.
不要訛詐那麼多.
可以嗎.
我雖然不是做得那麼好的職業.
但起碼不用背那麼多重費.
都是舒服一點的.
同意嗎.
還有就是坐直的人都開心.
原來我跟耶穌吃一餐飯.
原來耶穌的感染力那麼強.
你明白我的意思嗎.
可以感染到一個那麼有錢的人.
可以放棄他一半的資產.
亦都願意去如數付還.
給那些訛詐了的對象.
原來耶穌是那麼有能力的.
那些門徒開心嗎.
我跟師傅沒有跟錯.
還有當然就是撒該的本人和家人.
他的家人可能都不想告訴別人.
我爸爸做稅吏長.
現在告訴別人.
我爸爸就是分錢給別人的那個.
是不是很開心.
就是這樣.
當然我們做得耶穌的門徒都開心.

$^{641}$耶穌真是生命的主.
能夠改變人生命的主.
當一個人願意回轉的話.
天上的天使都為他鼓掌.
為他開心.
你開心嗎.
聽到你的故事開心嗎.
聽到民心的故事開心嗎.
一會兒聽你的故事可能更開心.
耶穌就向他宣告.
今天救恩到這個家庭.
因為他都是阿伯的子孫.
以前的人覺得這個撒該.
不應該是我們以色列人的同一個部落.
因為他是剝削我們的人.
但耶穌說他都是阿伯的子孫.
我來到尋找拯救失喪的人.
他曾經失喪了.
但我現在尋找他.
讓他成為神國的一份子.
Amen.
救恩就來到這個家庭.
不只是撒該一個人得救.
而是撒該整個家族.
都因為撒該的改變而得救.
將這個關係復和了.
一個人悔改.
天上都為他歡喜快樂.
因為人子尼是要尋找和拯救失喪的人.
今天我們都知道.
耶穌來尋找耶利哥的稅理長撒該.
並且他的全家因此而得救.
並且他的同行都很多人被改變.
他說繼續這個大使命交給我們.
我們都是被耶穌尋找拯救那班人.
繼續尋找拯救其他需要的人.
所以大使命都給了我們.
我們要將福音傳開去.
聖經沒有說撒該後來是怎樣.
但如果撒該的名字能夠記錄在聖經裡面.

$^{681}$就證明他在這幾十年裡.
人們都認識他.
教會的門徒都知道他什麼存在.
所以有名字可以叫他撒該.
原來我後來看回教會歷史裡面.
撒該就成為一個忠心跟隨主耶穌的人.
因為他不能做稅理長.
不能回頭了.
既得利益已經沒有了.
其他人已經做了職位.
但他就忠心成為耶穌的門徒.
並且成為使徒.
去繼續傳福音的工作.
當耶路撒冷被擊打的時候.
門徒四散的時候.
撒該竟然留在當時的凱撒尼亞.
當時耶穌問彼得.
你們說我是誰的那個地方.
那個是很大的軍區.
凱撒尼亞是一個很大的軍區.
是管領著耶路撒冷的.
可以嗎?.
他做了凱撒尼亞區的第一任主教.
本來是做稅理.
就是做官方的工作.
現在成為耶穌的門徒之後.
成為教會的領袖.
是整個凱撒尼亞區的主教.
Bishop of凱撒尼亞.
在第一到第二世紀之間.
在最艱難的日子當中.
他留下在這個地區裡面.
去鞏固當地的教會.
感不感動?.
這是真人真事.
在聖經歷史裡面.
他初初被人厭惡,被人排斥.
後來被人尊重.
我們的主教就是凱撒尼亞.
以前的撒該稅理長.

$^{721}$今天成為我們所尊重的教會領袖.
真正的喜樂不是你有多少錢.
不是你有多少江湖地位.
而是你能夠在上帝的國裡面.
找回你自己的位份.
找回你自己的可能性.
成為神中心的僕人.
剛才文森的故事.
我們今天的主題都是講大喜大樂.
他也是大喜大樂的人.
我們曾經都經歷過很多大喜大樂.
在座的弟兄姊妹朋友.
你們有沒有大喜大樂的故事?.
都有的 一定有的.
想當年也有大喜大樂的.
但今天我們想怎樣呢?.
你想不想得到上帝給我們的真正的平安.
和真正的喜樂呢?.
是不是?.
想不想?.
想就舉手好嗎?.
舉手 真是有平安有喜樂的.
文森的故事是很清楚的.
他那些案情.
你上網搜尋一下.
陳文森的那些全部上街了.
那些照片.
裁判在2023年就說他是誤殺罪名成立.
那些很多不同的貼文都有的.
後來呢.
剛才他講了.
當他肯認罪.
肯認自己是誤殺的時候.
不是謀殺的時候.
法官是會用聖經作為結案陳詞.
擺在現在的法院裡面的.
他說 我所做的我們不明白.
我所願意的我們不做.
我恨惡的我反去做.
立志為善的由得我.

$^{761}$行出來由不得我.
因此我所願意的善我反而不做.
我不願意的惡我反都去做.
但我們若認自己的罪神是信實的.
是公義的.
必要赦免我們的罪勢.
即是我們一切不義.
這個就是法官寫的判詞的其中一部分.
亦都是對文森來講一個最大的鼓勵.
他肯認自己的罪.
就是按理應的罪.
受回應有的刑罰.
那就坐滿了 坐夠了.
就可以被釋放.
然後他就願意奉獻自己.
行在上帝的心意.
不是再行自己想怎樣的心意.
Amen?.
當然上帝亦都給他很多回報.
人生可能會有做錯.
但亦都有知錯能改的機會.
這個就是他那兩個老師.
一個就是白牧師.
一個就是鐵盾牧師.
這個是他人生大喜大樂的其中一個大喜.
小冬.
即是他能夠上帝挽回他另外一份婚姻.
給他一個好的歸鎖.
他亦都有很多人在當中支持他.
令到他不離不棄.
然後他因為得到減刑的時候.
他有一個更加積極的人生.
當他重獲自由的時候.
他願意奉獻自己.
去分享神給他當中的領袖.
並且願意在信仰上面自己去真正學習.
繼續這麼艱難.
都這麼大年紀還去讀神學.
都不容易的事情.
是不是? 弟兄姊妹.

$^{801}$但他讀得很開心.
我經常見到他讀得很開心.
一會兒他讀完之後.
就要回去改功課了.
但無論如何.
我們知道.
絕處逢生.
我們可以蒙救恩.
人去到最低谷的時候.
最絕境的時候.
竟然是上帝會為我們開路.
同不同意嗎? 弟兄姊妹.
同不同意嗎?.
有沒有覺得你是有困難.
有絕境的呢?.
有的話你可以跟上帝說.
上帝啊.
你幫我吧.
好像耶利哥這個撒該.
幫我吧.
好像文心.
他跟上帝說.
你幫我吧.
是不是?.
上帝會聽我們禱告.
是不是?.
在最低處的時候.
上帝會來尋找拯救我們.
因為我們有很多信主的弟兄姊妹.
都可以為我們祈禱.
為我們去禱告.
求神去幫助我們.
有沒有?.
跟我們的親戚朋友說.
我有困難.
求你幫助我們.
求神開我們眼睛.
我們需要救主.
每一個人都需要救主.
不是說我不用祂.

$^{841}$我們面對生命的問題.
我們沒有能力去改變我們的生命.
我們面對死亡的問題.
我們沒有能力去面對死亡.
但我們的主是復活的主.
祂為我們的罪被釘十字架.
洗淨我們的罪.
並且祂流血捨身被埋葬.
三天之後祂從死裡復活.
告訴我們祂是滿有權柄.
有生命的主.
我們認祂為主的時候.
我們生命得益保證.
並且我們活得更加精彩.
求主憐憫我們.
我們希望能夠教會開放.
接納更加多不同的弟兄姊妹的家人.
親戚朋友來到.
也讓我們一起去讚美我們的主.
\newpage



\section{馬太福音 11:20-30}
\label{sec:032333}
\textbf{一同找到安息的智慧}~~~~魏文俊傳道
\newline
\newline
連結: \href{http://www.church.com.hk/acms/content.asp?site=cdc\&op=show\&type=product\&code=032333\&layout=sermon}{\texttt{032333}} ~~~~ 語音日期: 2026-06-14
\newline
\newline
\hyperref[sec:032332]{\small{< < < PREV SERMON < < <}}
~
\hyperref[sec:index_chronic]{\small{[返順時目]}}
~
\hyperref[sec:032334]{\small{> > > NEXT SERMON > > >}}
\newline
\newline
馬太福音 11:20-30
\newline
\begin{longtable}{cl}
\hline
\hline
章節 & 經文 (和合本修訂版)\\
\hline
11:20 & \begin{tabularx}{0.7\textwidth}{X} 那時，耶穌在一些城行了許多異能。因為城裡的人不肯悔改，他就責備那些城說： \end{tabularx} \\ \\ \relax
11:21 & \begin{tabularx}{0.7\textwidth}{X} 「哥拉汛哪，你有禍了！伯賽大啊，你有禍了！因為在你們中間所行的異能若行在推羅、西頓，他們早已披麻蒙灰悔改了。 \end{tabularx} \\ \\ \relax
11:22 & \begin{tabularx}{0.7\textwidth}{X} 但我告訴你們，在審判的日子，推羅和西頓所受的，比你們還容易受呢！ \end{tabularx} \\ \\ \relax
11:23 & \begin{tabularx}{0.7\textwidth}{X} 迦百農啊，你以為要被舉到天上嗎？你要被推下陰間！因為在你那裡所行的異能，若行在所多瑪，它還可以存留到今日。 \end{tabularx} \\ \\ \relax
11:24 & \begin{tabularx}{0.7\textwidth}{X} 但我告訴你們，在審判的日子，所多瑪地方所受的，比你們還容易受呢！」 \end{tabularx} \\ \\ \relax
11:25 & \begin{tabularx}{0.7\textwidth}{X} 那時，耶穌說：「父啊，天地的主，我感謝你！因為你把這些事向聰明智慧的人隱藏起來，而向嬰孩啟示出來。 \end{tabularx} \\ \\ \relax
11:26 & \begin{tabularx}{0.7\textwidth}{X} 父啊，是的，因為你的美意本是如此。 \end{tabularx} \\ \\ \relax
11:27 & \begin{tabularx}{0.7\textwidth}{X} 一切都是我父交給我的；除了父，沒有人知道子；除了子和子所願意啟示的人，沒有人知道父。 \end{tabularx} \\ \\ \relax
11:28 & \begin{tabularx}{0.7\textwidth}{X} 凡勞苦擔重擔的人都到我這裡來，我要使你們得安息。 \end{tabularx} \\ \\ \relax
11:29 & \begin{tabularx}{0.7\textwidth}{X} 我心裡柔和謙卑，你們當負我的軛，向我學習；這樣，你們的心靈就必得安息。 \end{tabularx} \\ \\ \relax
11:30 & \begin{tabularx}{0.7\textwidth}{X} 因為我的軛是容易的，我的擔子是輕省的。」 \end{tabularx} \\ \\
[1ex]
\hline
\hline
\end{longtable}
$^{1}$各位弟兄姊妹平安.
今天是父親節.
首先祝福各位叔寧的父親.
以及在家教會不同崗位.
默默無聲付出的弟兄和男士.
說一聲父親節快樂.
願神報答你們的勞苦.
今天揀選的經文是.
《馬太福音》十一章二十至三十節.
內容可能大家耳熟能詳.
是耶穌說怎樣才能找到安息.
亦是誰才能得著安息.
今天嘗試回到還原基本步.
回到我們信仰的核心當中.
透過這段經文.
我們嘗試在這個很忙碌.
或者很躺平的世代裡.
一起領受這個訊息.
希望我們作為耶穌基督的門徒.
可以回到小孩子那樣.
單單去倚靠耶穌找到安息.
好讓我們作父母的.
作叔寧父母的.
甚至我們出去傳講福音的.
能夠將悔改平安的福音.
承傳給我們下一代.
讓人得著那份喜樂的福音.
我會從兩個部分去看這段經文.
第一個部分是二十至二十四節.
耶穌去責備不信世代.
那個不信世代的原因和提醒.
看看耶穌對那個世代的人的提醒.
又對我們有沒有什麼特別.
我們要留意.
第二個部分是二十五至三十節.
是耶穌去呼籲所有人去傳言相信祂.
以至能夠在裡面找到安息.
透過這個部分我們一起檢視一下.
看看平日我們常常說的安息.
究竟是不是和耶穌說的一樣呢.

$^{41}$我們今天應該是我們的安息日.
我們今天過的安息日.
和耶穌說的有沒有不同呢.
這兩個部分我用兩個英文字去分出來.
希望兩個英文字可以幫助大家記住.
第一個英文字是"Believe".
一個簡單的英文字.
我們要相信耶穌所行的神蹟.
祂行神蹟的目的是要呼籲人.
呼召人去相信這個悔改的福音.
我們一起讀出二十至二十四節.
如果能看到我們一起.
「耶穌在諸城中行了許多義能.
那些城的人終不悔改.
就在那時則被他們說.
「約拉遜啊,你有禍了.
白菜大啊,你有禍了.
因為在你們中間所行的義能.
若行在推羅西頓.
他們早已披霸蒙灰悔改了.
但我告訴你們.
當審判的日子.
推羅西頓所受的.
比你們還容易受呢?.
加伯隆啊,你已經升到天上.
將來必墜落陰間.
因為在你那裡所行的義能.
若行在所多瑪.
他還可以存到今日.
但我告訴你們.
當審判的日子.
所多瑪所受的.
比你還容易受呢?」.
謝謝大家.
剛才讀出那幾節經文有否發現到.
有一個重複的結構.
耶穌去說某一個城.
他們怎樣.
然後他就用另一個舊約的城.
去做一個比較.

$^{81}$耶穌去透過特別提到.
哥拉遜,白菜大和加伯隆這三個城市.
他強烈去譴責當時那個世代的人.
看到神蹟,但是不願意悔改.
耶穌說,末後他們就會受到審判.
這三個是什麼地方呢?.
哥拉遜,白菜大和加伯隆.
都是很接近加利利湖邊的地方.
換句話說.
就是很接近耶穌事奉的起點.
也是耶穌事奉的基地附近.
這幾個地方.
這裡的附近.
是耶穌去行神蹟,行他的義能最多.
也應該是那些人看到耶穌行神蹟最多的地方.
理論上應該是最親近耶穌.
應該是越看到神蹟就越相信.
不過如果大家打開聖經會發現.
其實那些人好像不是這樣.
耶穌在剛才大家讀的那段經文.
去到強調看到神蹟.
悔改和承受審判的關係.
在這幾節裡面.
耶穌說如果他的神蹟行在推羅西頓那裡.
裡面的百姓一早就已經披麻蒙灰.
即是他們會悔改.
如果耶穌的神蹟.
他行的義能在索多瑪這個地方.
索多瑪的百姓看到神蹟.
這個地方到耶穌的時代.
可能甚至是今天.
都能夠存活.
可以留到今天.
這三個是什麼城市呢?.
其實就是舊約裡面作為搏翼神.
一個代表的城市.
耶穌說這三個在舊約裡面這麼差的城市.
只要他們看到耶穌的神蹟.
也會悔改.
可能大家坐在這裡就心想.

$^{121}$我也經歷過耶穌基督.
應該是因為經歷過耶穌才會信耶穌.
應該一看到就悔改.
為什麼這三個城市不悔改呢?.
耶穌就用舊約這幾個城市.
去突出了哥拉遜,白塞大和加伯隆的惡.
他們的惡是什麼呢?.
在二十三節去講加伯隆的時候.
就很明顯地用不同的聖經去講給當時的人聽.
在二十三節裡面有一句很特別的.
他提到加伯隆就說.
加伯隆已經升到天上面.
將來必墜落到陰間.
不知道大家讀的時候有沒有像我一樣.
有個畫面就是有個城市升到天上.
然後又跌到地上讀一讀的時候.
會不會是這樣呢?.
原來耶穌是在引用《以賽亞書》其中一段經文.
他將這幾個地方.
包括加伯隆,哥拉遜和白塞大.
去比喻舊約的巴比倫王.
他說這幾個地方為什麼不悔改.
就是他們很驕傲自大.
和巴比倫王一樣.
將來就好像巴比倫一樣會被神離棄.
這三個城市的惡.
是在於他們驕傲自大覺得不需要神.
原來他們的罪.
這三個城市的罪.
是代表當時猶太人的驕傲.
當時猶太人的驕傲.
就是在於他們覺得我們自己可以守好律法.
就可以取悅到神.
我就可以去到神那裡.
耶穌基督不是神的獨生子.
不是祂來的.
我們還要等待另外一個神派下來的人.
去拯救我們.
我們守好律法就可以.
所以這段經文.

$^{161}$整段今天的經文核心的信息是講舊因的.
是關乎舊因的經文.
如果在座有誤導未信主的.
你可以這段經文告訴我們.
原來我們看到耶穌基督的神蹟.
我們去悔改.
去禱告,認罪.
承認耶穌基督是生命的主.
就可以得著舊因.
如果當中真的有這樣的朋友.
你可以在崇拜聚會之後.
找我們白色制服的傳道同工為你祈禱.
也可以找我為你祈禱.
其他的弟兄姊妹呢?.
會不會今天的報道信息.
其實和大家好像沒有什麼關係呢?.
不是的.
想說即使信了主.
我們也要持續地相信神所做的神蹟.
每天認罪悔改.
不過問題是我們有沒有覺得.
有看到神蹟在我們身邊發生呢?.
弟兄姊妹你有沒有見過神蹟發生呢?.
有沒有見過聖經裡面的神蹟發生呢?.
例如盲眼的得醫治,瘸腿的能夠醫治,站起來.
有沒有見過這些呢?.
我看到有人搖頭.
但是你想清楚有沒有見過神蹟.
還是我們看到有些事發生了.
我們覺得那些不是神蹟.
覺得這些事其實是每天都在發生.
也不關神事.
是我自己去做或是其他人去做的事.
我有不同的問題.
所以我自己解決.
神只是會介入病患要醫治的事情.
或者有其他的事.
我有困難.
我自己就可以解決.
所以我想我們在當中去學習.

$^{201}$其中一個很重要的信息是.
我們要承認自己的軟弱.
如果我們願意承認自己的軟弱.
我們才發現原來我們身邊大大小小的事情.
都是那位似人出人意外平安的神所施行的神蹟.
如果我們能夠這樣.
我們的屬靈眼睛能夠打開.
我們才發現每個星期回來.
能夠崇拜,唱詩,聚會,坐在那裡聽信息.
能夠參與教會聚會.
打開成長班的課本去查經.
甚至去參與服事,唱詩.
在成長班裡面做主席祈禱.
甚至去做奉獻祈禱.
每件日常發生的事.
其實都是神所賜下的神蹟.
不是我們能力能夠控制的.
是神自己親自讓我們能夠成就.
也做到這些事情.
所以第一個學習.
我們要在耶穌面前承認自己的軟弱.
將主權歸還給神.
對付自己心裡的一些軟弱,一些罪.
拆毀我們心裡的偶像.
律法是猶太人的偶像.
他們覺得不需要了.
他們做得非常好.
不需要悔改.
他們做得很好.
我們今天擺了什麼偶像在我們心裡面?.
小時候我脾氣很大.
我的家人非常熟悉我.
所以都會幫我改了一個花名.
我的師長或一些屬靈前輩更加明白我.
或者都會知道我是這樣.
在教會99.9\%的人都沒有見過我發脾氣.
是真那種發脾氣.
不是裝那種.
所以都是成長裡面不同人提醒.
所以長大後學習去壓住一些怒氣.

$^{241}$那種脾氣都不發出來.
又不滿意都放在心裡就算了.
直到我自己小朋友出生.
發現照顧小朋友的壓力真的很大.
壓力大在哪裡呢?.
如果他還未懂得走.
放在嬰兒床上都可以.
他不懂得說話.
你看他都很開心和他玩.
但到他懂得走的時候.
是父母噩夢的開始.
因為他會走.
會跑.
然後你叫他停.
他就會走了.
到他會說話.
就更加難搞.
有人會心未笑就很明白.
小朋友出生我的脾氣就再大一點.
有時太太都在家裡說.
你不出聲不說話沒有罵人.
但好像全世界都得罪你似的.
小朋友入幼稚園之後.
那個挑戰就越來越大.
我自己裡面的脾氣都越來越多.
有時按都按不住會發出來.
發脾氣之後.
不知道大家有沒有試過有這樣的感受.
就是覺得為何要發脾氣.
其實很小事.
遲了一點點出門.
很小事一分鐘.
為何要發脾氣.
又會覺得自己很無能.
很無奈.
為何要發脾氣令到大家都不開心.
不過感恩是.
小朋友入了幼稚園之後.
神讓我們遇到很有愛心的老師和社工.
他們提醒和幫助.

$^{281}$不是小朋友而是我們作為家長.
知道原來小朋友的成長.
不是捉實去控制他成長.
是陪伴去關心去引導.
放手讓他在神的恩典底下成長.
是不是聽起來都很容易呢.
我見到很多人都會心微笑.
說就很容易.
做就很困難.
很多掙扎.
的確過去兩年我自己都很多掙扎.
要放下那種執著很難.
不過我又回望自己成長路上.
神其實安排了很多人讓我能夠成長.
我自己的家人就接納我很多.
知道我是很大脾氣.
我自己的太太.
或者在教會裡面的屬靈導師.
我的舊同事.
也是一些神放在我身邊成長的地方.
這些是我自己的軟弱.
時常我都會想.
怎樣才能夠平復到內心那種不開心.
或者憤怒的脾氣呢.
究竟神又放了些什麼.
其實要拆毀心裡面一些捉緊了的東西.
你要放下給神.
其實很不容易的.
因為跟我們的性格成長是很不同.
但是原來我們看到神的神蹟.
在我們身上一直發生.
我們有很多人都是神放在我們身邊.
那些神蹟.
原來神就是呼籲.
呼召我們回去祂的身邊.
一起去學習.
降伏在祂十字架的恩典底下.
放手交回給神.
去處理我們心裡面的掙扎和執著.
這是第一步.

$^{321}$悔改是我們回到神裡面.
找到安息的第一步.
所以第二個英文字也很簡單.
Belief之後第二個英文字就是Find.
一起找到安息.
我們要一起在主裡面找到安息.
那種安息不是我們開玩笑的那種安息.
也不是在安息禮拜時候才有的安息.
而是忙碌或者自己獨處.
很多時間我們都能夠相信神的恩手一直在作供.
這才是在聖經裡面所說的安息.
現在想大家跟我一起讀出第二部分的經文.
看到我們一起讀吧.
預備起.
那時耶穌說.
父啊 天地的主.
因為你將這些事向聰明通達人就藏起來.
向嬰孩就顯出來.
父啊 是的.
因為你的美意本是如此.
一切所有的.
都是我父交付我的.
除了父沒有人知道子.
除了子和子所願意指示的.
沒有人知道父.
凡奴婦擔重擔的人.
可以到我這裡來.
我就使你們得安息.
我心裡柔和謙卑.
你們當付我的額.
學我的樣式.
這樣你們心裡就必得向安息.
因為我的額是容易的.
我的擔子是輕散的.
耶穌在剛才的二十至二十四節說.
那個世代的人不悔改有什麼後果.
下一段在這個部分就說.
其實怎樣才能找到安息.
我們跳過二十五至二十七節比較難明白的部分.
先去我們最熟悉的那一節經文.

$^{361}$如果不打出來可能大家都會懂得背那段經文.
是二十八節.
凡奴婦擔重擔的人.
可以到我這裡來.
我就使你們得安息.
如果我們只看這一句金句.
我們可能對「救恩」有很大的誤解.
因為剛才開頭有說.
這一段經文其實都是在說.
在耶穌基督裡面的「救恩」.
我們只看這一句會以為.
將那些奴婦重擔給了耶穌.
我有「救恩」就行了.
事情我就不用做了.
如果我有家裡很大的重擔就不用照顧了.
我健康隨時都會好起來.
所以可能有些人誤解之後會怎樣呢?.
他信了耶穌.
家裡雖然關係很複雜.
也跟家人關係不好.
但我不管了.
我回來事奉就好了.
如果你誤解了.
可能身體都很差.
不過我不管了.
不用治療了.
我回來事奉就夠了.
不過原來經文不是在說這件事.
這個經文是在說那種「救恩」.
其實在耶穌基督裡面怎能夠得著安息.
不是說我們全給了祂就可以.
我們如果看二十九節往後的部分.
甚至看完十二章我們會更加清楚.
不過今天我們沒時間看十二章了.
我們只看二十九至三十節.
我們都會明白.
那種耶穌基督在二十八節這裡的邀請.
其實是什麼意思.
在二十九節裡面.
耶穌說:我心裡柔和謙卑.

$^{401}$你們當乎我的額.
學我的樣式.
這樣你們心裡就必得享安息.
下面那兩點其實就是.
根據原文翻譯擺好次序.
耶穌原來是說.
給兩個命令當時的人.
祂說你們要當乎我的額.
學我的樣式.
為什麼要學?.
耶穌說:因為我心裡柔和謙卑.
你們就學吧.
兩個命令.
第一個「乎額」.
就是回應當時那群猶太人.
他們覺得要守好所有律法.
才能夠承受到救恩的錯謀.
耶穌叫他們放下.
放下那些律法.
放下心裡那種執著.
你乎我的額.
這一句告訴我們.
救恩是否需要背負一些事情.
也需要吧.
我聽過有弟兄姊妹說.
沒有的.
你放下所有律法.
就可以跟隨耶穌.
但好像說漏了.
二十九節.
耶穌說:你們當乎我的額.
原來救恩也是由「額」要負的.
不是什麼都沒有的.
不過在三十節.
耶穌基督給了很大的安慰.
我們一群耶穌的門徒.
耶穌在三十節說.
祂強調祂的「額」是容易和輕散.
這是對比起法利賽人那種.
要做好所有律法.

$^{441}$是一個很大的對比.
是容易和輕散的.
所以這也是信仰浸禮班學了很多年.
原來「乎額」是相信耶穌基督.
是神所差來的微塞亞.
相信祂有能力拯救我們.
不再需要依靠律法.
這是第一個命令「乎額」.
第二個命令.
學祂的樣式和學本.
就將這句經文翻譯了.
說「向我學習」.
耶穌基督直接叫人向祂學習.
為什麼要向祂學習呢?.
又要學習什麼呢?.
我們又回到剛才跳過的幾節.
比較難解的經文.
二十五至二十七節.
其實有兩句耶穌這樣說的.
什麼叫做有智慧呢?.
耶穌說那些智慧.
在二十五節說.
因為你(神)將這些事.
將救恩向聰明通達人藏起來.
向嬰孩顯出來.
二十七節說.
除了子所願意指示的.
也就是說救恩.
沒有人知道乎.
這個顯出來和指示的.
其實都是在說啟示的意思.
耶穌去提醒那些人.
特別是跟從猶太人.
要守律法的人.
祂才是天父所啟示祂.
能夠去成就救恩的那一位.
祂要打破猶太人的社會.
那種要靠律法.
靠行為得救的習慣.
耶穌說那種智慧.

$^{481}$不是像法利賽人那樣.
也不是像我們今天這樣說.
要做好所有事情.
才能得著救恩.
而是單單好像嬰孩一樣.
去承認自己沒有能力.
必須要像倚賴父母那樣.
全然信靠耶穌.
弟兄姊妹.
你近來有沒有試過.
真的很無助地求救?.
是怎樣的呢?.
很無助地求救.
或者你是否能夠回想.
你是嬰兒的時候.
真的攤開雙手.
要父母牽手才能走的時候.
可能很難回想.
但是其實.
原來耶穌在聖經裡說.
我們是攤開雙手.
找耶穌基督.
承認自己其實沒有那個能力.
必須要倚賴耶穌基督.
才是真正的安息.
我們能夠倚賴耶穌.
是因為耶穌在二十五至二十七節.
都說得很清楚.
裡面剛才大家都讀了很多次.
他提了很多次父.
他在強調他和天父很緊密的關係.
天父就差派了耶穌成就救恩.
所以我們能夠全然相信他.
所以回答剛才的問題.
猶太人要學習什麼呢?.
就是學習耶穌的柔和謙卑.
不要以為自己已經守好律法.
就不需要救恩.
耶穌在這段經文二十五至二十七節.
其實就像舊約的律法師傅.

$^{521}$舊約《箴言傳道書》的智慧老人.
這樣告訴他.
什麼叫做有智慧呢?.
就是相信耶穌的神蹟.
他所成就的救恩就可以悔改.
這樣才能夠有安息.
所以我們其實都是這樣相信的.
我們不再需要律法.
而是我們單單相信耶穌.
和祂建立關係.
祂就會聆聽我們的禱告和需要.
就能夠和神復和.
我們都是這樣相信的.
大家聽到這裡.
會覺得好像今天.
和大家重溫信仰浸禮班的初信的事情.
我們今天還原基本部有問.
我都很想問一下.
其實我都是這樣問自己.
其實我們和神.
是不是有這樣的信任關係呢?.
有還是沒有在我們生活上反映出來.
我們生活上發生一些事情.
我們是不是真的在耶穌裡面有安息呢?.
的確我們不會再打開摩西五經.
去跟著逐條逐條手的.
不過可能今個世代有一種律法.
我們都跟了.
亦都是無意識地去跟.
我作了一個名字給祂.
這個律法叫做抵抗焦慮的律法.
什麼意思呢?.
做父母的.
怎樣才可以不用擔心小朋友.
幫他安排好.
入一間好的學校.
幼稚園,小學,中學,大學.
入一個好的科.
找到一份好的工作.
看著我要安排好他的成長.

$^{561}$如果他到一間不好的學校,我就很擔心了.
我的工作有很多事情要做.
下班都做不完,那怎麼辦呢?.
帶回家做吧.
半夜都做不完,那怎麼辦呢?.
那就繼續做吧.
或者都做不完,那怎麼辦呢?.
很擔心,我睡不著.
我的健康.
有些病患.
我看醫生吧.
看醫生解決不了,那怎麼辦呢?.
很擔心,我睡不著.
我們其實是不是相信.
跟神有這樣的信任關係.
在我們生活中反映.
可能你不是這一種.
不是要安排好所有東西.
是今個世代的另一種.
就是很多事情要做,很煩.
那怎麼辦呢?.
不管了.
神就搞定了.
不用安排那麼多.
兒女很感恩.
神都讓他入了一間好的學校.
搞定了,我不用想.
在這兩種好像很不同的做法下.
其實都是同一種原因.
就是我們很想控制所有事情發生.
不過一種是我極力去控制.
一種是算了,我都控制不了.
隨便吧.
那種隨便的其實都未必是信靠.
在我們生活上各種的控制慾.
其實實際上.
就是我們不願意向耶穌基督悔改的硬心.
事奉上可能都會的.
我們要安排所有事奉.
每一分鐘怎樣運用周會.

$^{601}$每一分鐘要怎樣做.
其實是很好的.
不過遲了一分鐘就快到了.
我們就很心急.
祂說完沒有?.
或者可能崇拜的當中.
大家都會看一看鐘.
哎呀,阿彌陀佛.
祂說完沒有?我一會兒又要怎樣?.
還未完的,為甚麼?.
事奉上可能都有這樣.
為甚麼那個人還未交那樣東西給我?.
不單止自己焦慮.
當我們去催迫別人的時候.
我們自己都會焦慮.
不單止是我們放手.
不用計劃,交給聖靈就可以了.
不單止是我們這樣放手就算了.
我們令到別人都會焦慮.
這樣的信任關係.
是怎樣呢?其實我們同生.
傳言倚靠耶穌的安息.
其實是在不確定之中.
我們仍然知道祂會掌權.
按祂的旨意去成就救恩.
這種信任關係.
其實就好像和我們最心愛的人去溝通.
或者最信任的人溝通.
不知你有沒有和最信任的人去WhatsApp.
我和太太有時上班都會交代.
我們在做甚麼,甚麼時候吃飯.
有時傳完訊息出去.
幾個小時都不回覆.
大家會怎樣呢?.
為甚麼你不回覆我?.
你在做甚麼?.
是不是很空閒?.
其實沒有這樣問過.
不過可能會有這樣想.
我和太太都認識了很多年.

$^{641}$所以我們都明白.
大家不回覆是甚麼原因.
可能在開會,可能忙碌中.
未有空回覆.
如果忙完開完會就會馬上回覆.
你是否這樣相信耶穌基督?.
你知道耶穌祂很清楚你的需要.
你都知道祈完禱之後.
或者甚至不用祈禱,祂都知道你的需要.
你願意耐心等候祂的旨意成就.
還是你心裡反過來想.
我祈禱了很久,為甚麼祂都不回覆我呢?.
當我們信任主耶穌基督.
相信祂的時候.
我們有這種和神的深厚關係.
在耶穌裡面負的壓才容易和輕散.
所以我們問的問題是.
不是做多還是做少.
而是放下我們要控制所有事的那種控制慾.
安心地相信耶穌基督.
其實在我們生命裡面工作.
祂其實是會施行神蹟.
問題是我們當不當那些是神蹟.
我們覺得那些是普通人去做一些事.
但是可以告訴大家其實不是.
神是一直在掌管我們的生命這個世界.
祂仍然在掌管這個世界.
我想用最近的經歷來作結束.
在六月初,我正在預備這篇訊息.
我和太太在家裡發生了一些事.
就是我和太太要帶小朋友去做一些平衡.
在平衡的前一晚安排了小朋友睡覺.
我們都很擔心.
我和太太談了很久.
大約談到十二時左右.
很睏,太太睡覺.
我也很睏,不過我很擔心睡不著.
然後我做了一件事.
就是現在的人很喜歡做的.
即是年輕人特別喜歡做的.

$^{681}$我不知道大家有沒有做.
就是我打開了電腦.
然後問一下AI.
我將所有的情況打在AI上.
然後就問,那是甚麼意思.
那天晚上我和AI辯論了一小時.
我和它辯論了一小時.
它說了一些情況給我聽.
我說不是的,你這樣好像不是這樣.
可能是這樣,然後它也很好.
給了我很多感受.
它說你觀察得很好.
不過可能都是這樣.
談了一小時,它勞列了ABCDEF很多的情況給我.
然後我就覺得兩點多,很睏.
明天要早起和小朋友出去.
所以我就睡覺.
第二天早起,其實也很擔心.
大家都預備好了.
然後我坐在那裡.
太太和小朋友在預備.
我很擔心,不敢告訴他們.
因為說了他們會更加擔心.
小朋友也擔心,無謂了.
影響了大家的心情.
我坐在那裡擔心,打開那些AI的字.
因為真的不行.
然後我突然間想到一首.
我以前在某個聚會聽的一首詩歌.
這首詩歌叫做《平靜討問》.
是英文的歌詞.
有人去做中文翻譯.
然後我就拿起家裡的結他在那裡唱.
今天沒有拿結他.
要請詩琴幫我彈音.
謝謝詩琴.
由此勇氣扭轉可變改的事.
願給我智慧看得清.
主就讓我的心安穩.
來承受不可扭轉難事.

$^{721}$由此勇氣扭轉可變改的事.
兩者的區分靠睿智.
其實我在那天在早上唱了很多次.
一邊唱一邊聽歌詞提醒我.
其實有些事我是不能夠改變的.
因為我是一個人.
我有我的界限.
我不是神.
其實我很多事都做不到.
其實有些事我不能夠控制.
那個結果我不能控制.
不過我可以做一件事.
就是我依靠著耶穌基督去面對那個結果.
我相信其實祂在每一件事上都有祂的心意.
AI可以幫到我們很多.
如果你有用過AI.
或者你上網問過問題.
其實你就知道AI可以幫到我們一些事.
因為就算你沒有和AI談過.
你上網找東西.
它都幫你整合了一些資料給你.
它幫助我們很多資料搜集.
不過AI做不到一件事.
就是它不能夠令我們心裡面那種恐懼.
我們那種焦慮不安.
它消除不了.
它給多少資訊.
說你觀察力很好.
你真的分析很好.
其實心裡面那種焦慮它幫不了我們.
唯有耶穌基督那種安穩.
在祂裡面的安息.
才能夠平靜到我們.
因為我們相信其實祂是我們的救主.
祂真的能夠安排和掌權.
我們相信祂.
所以唱完詩歌.
我的心就平靜了.
不需要再和AI談話.
我和太太帶了小朋友出去.

$^{761}$不知你一直聽著.
你有沒有一些事情.
你都想著放在心裡.
其實真的解決不了.
你做了很多事.
都沒有什麼能力去再處理.
今天藉著經文.
這個真是AI幫我做的.
藉著經文告訴我們.
如果我們真的相信耶穌基督.
祂行的神蹟.
我們就看到神蹟.
我們能夠回到耶穌基督那裡.
我們一齊就能夠找到安息.
我們一齊都要學習.
特別是在香港這個這麼忙碌的地方.
這麼焦慮的地方.
我們更加要學習.
在耶穌裡面有那份安息.
不是做多或者做少.
而是在做多或者做少的時候.
我們在耶穌裡面.
有沒有那份安息.
如果你未來有一次突然間.
你又對自己的子女.
或者屬靈的子女.
或者在事奉上.
甚至是你的健康.
你的工作.
你突然間都很焦慮.
你試試停一分鐘.
你停一分鐘.
你跟耶穌說.
主啊.
這件事我其實管不了.
我處理不了.
現在我放手.
交回在你手上.
我願意依靠你去面對面前的挑戰.
我們一齊學習在耶穌裡面有安息.

$^{801}$這樣今天的傳道主日和父親節才有意義.
因為當我們傳承信仰的時候.
如果我們真的有安息.
我們就將我們自己裡面那種焦慮不安.
給了一些未信的人.
或者我們的子女.
我們學習在耶穌裡面有安息.
我們才能夠將悔改.
平安的福音傳承給下一代.
我們一起祈禱.
我們要求你賜下勇氣和智慧.
讓我們能夠分辨.
在哪些時候要學習.
謹耐和陪伴.
在哪些時候要開口宣揚.
使人悔改.
歸向你平安的福音.
求你差遣我們.
使用我們.
禱告乃奉我主耶穌基督.
得勝之名祈求.
阿們.
\newpage



\section{申命記 6:6-7}
\label{sec:032334}
\textbf{信仰的傳承 (華)}~~~~劉彤牧師
\newline
\newline
連結: \href{http://www.church.com.hk/acms/content.asp?site=cdc\&op=show\&type=product\&code=032334\&layout=sermon}{\texttt{032334}} ~~~~ 語音日期: 2026-06-21
\newline
\newline
\hyperref[sec:032333]{\small{< < < PREV SERMON < < <}}
~
\hyperref[sec:index_chronic]{\small{[返順時目]}}
~
\hyperref[sec:032335]{\small{> > > NEXT SERMON > > >}}
\newline
\newline
申命記 6:6-7
\newline
\begin{longtable}{cl}
\hline
\hline
章節 & 經文 (和合本修訂版)\\
\hline
6:6 & \begin{tabularx}{0.7\textwidth}{X} 我今日吩咐你的這些話都要記在心上， \end{tabularx} \\ \\ \relax
6:7 & \begin{tabularx}{0.7\textwidth}{X} 也要殷勤教導你的兒女。無論你坐在家裡，走在路上，躺下，起來，都要吟誦。 \end{tabularx} \\ \\
[1ex]
\hline
\hline
\end{longtable}
$^{1}$讀一節聖經在申命記第六章第六到第七節.
我們一同來念這兩節的聖經來.
我今日所吩咐你的話都要記在心上.
也要殷勤教訓你的兒女.
無論你坐在家裡,行在路上,躺下,起來都要談論.
我們就讀到這裡.
今天是父親節.
首先在這裡恭賀每一位偉大的父親.
父親節非常非常快樂.
Amen, Hallelujah.
你知道每年的父親節.
我們教會都會為父親們安排一個全家一起照相的機會.
希望你今天三會之後把握機會去照相.
這些年我都會把我們每一年跟孩子們照的照片.
放在我辦公室的這個bulletin board上.
當然現在他們都長大了.
不見得每一年我們都能湊得在一起.
但是當我回頭看.
這些年跟孩子們一起成長的照片的時候.
雖然我常常覺得自己做一個父親.
常常在軟弱無助當中.
但是每一次當我在軟弱當中的時候.
就看見神總是帶領我們走過人生每一個過程.
成為我們的供應跟力量.
Amen.
我就看見其實在一個孩子成長的過程當中.
父親的形象是非常重要.
而這個父親的形象不是因為我們的教導的緣故.
而是好像一種在生活上的潛移默化的一種影響力.
是我們的生活行為.
是我們的喜好.
是我們的厭惡.
是我們的習慣.
甚至是我們的憂慮跟恐懼.
所以你看見每一件我們生活的事情.
其實都會直接或間接的影響到我們兒女.
所以有時候當你看見你兒女的時候.
就好像看見你自己不是嗎.
所以做一個父親的人.
我們一點不要輕忽神所託付給我們教養兒女的責任.

$^{41}$信仰是需要傳承的.
Amen.
跟你們說信仰是需要傳承的.
我們要求神幫助我們.
在我們兒女身上能夠帶來屬靈的影響力.
那你說我們怎麼樣才能夠在家中.
將我們的信仰能夠傳承下去呢.
今天早上從幾方面跟大家來分享.
首先第一個.
一定要主動來跟隨神.
跟你們說要主動來跟隨神.
我想做一個基督徒.
我們一定是要為我們自己的靈命來負責任.
我們一定要主動的來建立一個與神的關係.
特別是作為基督徒父親.
更是要為著我們兒女的緣故.
我們真的是要建立一個與神美好的關係.
為什麼呢.
因為信仰是需要傳承的.
一代又一代的傳承.
當這個父親主動來跟隨神的時候.
這個一代一代.
我們的兒女也會跟著來跟從神的.
在聖經當中我們看見亞伯拉罕就是個最好的例子.
你知道在聖經裡面看見亞伯拉罕跟羅德兩個人.
他們都是從吾爾出來的.
兩個人都要進入神的應許之地.
這兩個人其實都是認識神的人.
這兩個人也是都是屬於神的人.
但當我們看見這兩個人的結局的時候.
那真是何等的不同啊.
亞伯拉罕他的信仰真的是一代又一代的傳承.
所以在聖經裡面每次講到我們所信的神的時候.
都說我們所信的神是亞伯拉罕 以撒 雅各的神.
這是何等美好的見證.
你看見信仰從亞伯拉罕傳到以撒.
從以撒傳到雅各.
真的一代一代的傳承.
可是羅德呢.
他的妻子因為貪念索多瑪.

$^{81}$在逃走的時候捨不得索多瑪.
就回頭看.
而一回頭就變成一個炎柱.
他的女婿不肯從索多瑪出來.
後來就跟索多瑪一同滅亡.
他的女兒不敬畏神.
與他犯了亂倫之罪.
生出兩個兒子摩亞跟雅門.
也成為神咒詛的指名.
我們看見這兩個人亞伯拉罕跟羅德.
同樣是從吾爾出來.
同樣是屬於神.
為什麼兩個人卻有這樣不同的結局呢.
最主要的原因.
我們看見亞伯拉罕.
他是一個主動跟隨神的人.
你看見亞伯拉罕無論走到哪裡.
聖經說他都為神燭壇獻祭.
創世記第十二章第七到第八節.
這裡告訴我們.
耶華向亞伯拉罕顯現.
說我要把這地賜給你的後裔.
亞伯拉罕就在那裡.
為向他顯現的耶華燭了一座壇.
從那裡他又遷到.
伯特利東邊的山支搭帳篷.
西邊是伯特利東邊是愛.
他在那裡又為耶華燭了一座壇.
又求告耶和華.
你看見亞伯拉罕無論走到哪裡.
他都燭壇獻祭.
為什麼要燭壇.
這段經文很清楚告訴我們.
壇是他求告耶華的地方.
這個祭壇是他向神獻上感恩的地方.
這個祭壇也是他獻上他自己生命.
獻上自己給神的一個地方.
在聖經裡面我們也看見.
這個祭壇更是神與人相會的地方.
神在這裡向亞伯拉罕說話.

$^{121}$神在這裡向亞伯拉罕立約.
也就在這個壇的所在.
我們看見亞伯拉罕與神建立了一個非常美好.
非常親密的關係.
所以在創世記十八章十七節你看見.
神怎麼說亞伯拉罕呢.
創世記十八章十七節耶華說.
我所要做的事豈可怎麼樣.
瞞著亞伯拉罕.
我常說如果今天神能夠對我們說這句話.
我們死而無憾了不是嗎.
換句話說神在說我跟亞伯拉罕沒有什麼秘密.
我所要做的事我通通會告訴亞伯拉罕.
你看見亞伯拉罕與神之間有多麼密切美好的關係.
其實這不是因為神偏心.
也不是神特別愛忠義.
這個亞伯拉罕不是.
是因為亞伯拉罕他是一個肯花時間.
主動跟隨神的人.
反過來你看羅德呢.
你看在聖經裡面講到羅德的時候.
你從來沒有看見羅德為神築壇獻祭過.
當然你會說亞伯拉罕已經築了壇.
他幹嘛再築另一座壇呢沒錯.
但是後來當羅德離開亞伯拉罕.
他自己帶他家人到約旦河谷去的時候.
那你也沒看見他在約旦河谷的時候.
為神建築壇來獻祭沒有啊.
所以我們看見亞伯拉罕.
他是一個主動跟隨神的人.
那羅德呢.
從其量你只能說他是一個跟隨亞伯拉罕的人.
他從來沒有自己跟神建立一個親密的關係.
他只是跟隨亞伯拉罕.
亞伯拉罕到哪裡他跟到哪裡.
亞伯拉罕築壇獻祭敬拜神.
他就順便跟著亞伯拉罕敬拜一下.
我發現今天有很多基督徒.
也是這種我叫做順便敬拜一下的基督徒.
你們這裡敬拜真好耶.

$^{161}$講道也很好耶.
小主也很好.
那我順便來跟你們敬拜一下吧.
今天很多的基督徒是這種我叫做順便型的基督徒.
他們好像羅德一樣.
他們沒有跟神建立一個直接的關係.
所以說不但他們不能從神那裡得著祝福.
他們也無法為他們的家庭帶來真正的影響力.
今天做一個基督徒.
你一定要在基本的事情上下功夫.
你一定要花時間來親近神.
要做一個主動跟隨神的人.
我常說在人的一生當中.
總有一個地方是能夠決定你一生成敗的.
總有一個地方.
這個地方是什麼地方呢.
這個地方就是你雙膝跪下與神面對面的地方.
所以不要再做一個順便型的基督徒.
我跟你們講說不要再做一個順便型的基督徒.
今天讓我再一次呼召你.
在每天的生活當中做三件事情.
第一個每天讀神的話.
第二個每天花時間禱告.
第三個按照神的話語去遵行.
主動來跟隨神.
每天三件事情.
花時間讀神的話.
花時間禱告.
並且按照神的話.
遵行神的話.
當你這樣做的時候.
你不但能夠為你自己的生活當中帶來突破.
你更是能夠為你的家庭帶來祝福的.
所以第一個主動跟隨神.
第二個你一定要開口祝福你的家人.
給你們要開口祝福你的家人.
你要知道在聖經當中.
一個做父親的人.
他的話語是很重要.
一個做父親的人.

$^{201}$他是擁有神對家庭的一個權柄.
所以一個父親的話.
是可以帶來祝福.
但是另一方面父親的話.
也可以帶來咒詛的.
在舊約裡面的雅各.
他就非常了解這個真理.
所以當他的父親以撒將要離開世界的時候.
你看見雅各花了多少功夫.
要盡各樣的計謀.
就是要騙取他的父親原來要給他哥哥長子的祝福.
你知道我常常在想.
為什麼要這麼大費周章呢.
難道那不就是一個祝福的禱告而已嗎.
幹嘛要花這麼多功夫還要欺騙他的父親.
要搶奪這個長子.
為什麼要大費周章.
其實並不那麼簡單.
只是祝福禱告而已.
雅各非常的明白.
一個父親的祝福是可以影響他一生的生命的.
所以你看見當雅各自己在晚年的時候.
他還為他的孩子一個一個按手祝福禱告.
所以特別讀到創世記49章的時候.
其實這個49章整章的聖經.
就是雅各為他的孩子一個一個祝福禱告的話語.
所以28節裡面在結束的時候這裡說.
這一切是以色列的十二個支派.
這是他們的父親就是雅各.
對他們所說的話.
為他們所祝的福.
都是按照個人的福分為他們祝福.
整章的聖經雅各為他的孩子一個一個祝福禱告.
我們看見雅各非常明白.
一個父親的祝福是很重要.
今天在我們中間做父親的.
你有沒有常常為你的孩子做祝福的禱告.
向他們說祝福的話語呢.
也許你說牧師我雖然沒有常常祝福.
但起碼我沒有咒詛他們沒有啊.

$^{241}$你要明白啊.
當你不開口祝福他們的時候.
你就是剝奪他們得蒙神祝福的權利.
因為在神的心意當中做一個父親的你.
就是神祝福家庭的管道.
Amen.
你要看重你的這個位分的父親們.
你就是神祝福家庭的管道.
如果你不祝福他們.
他們就無從得著神的祝福.
但什麼時候你一開口祝福.
神的祝福就臨到在你的家庭當中.
Hallelujah.
Amen.
我發現我們做華人的父親們.
在這方面都常常很欠缺.
我們多半都是愛在心裡.
口難開.
我們都是暗地裡欣賞啊.
卻很少主動的去鼓勵.
No no no no.
你要知道你的孩子是非常需要你的鼓勵的.
無論年齡多大多小甚至長大成人了.
都仍然需要你的鼓勵.
Amen.
因為你的鼓勵代表了你的接納.
當一個孩子知道他自己完全被父母接納的時候.
他才可能有一個正確的自我形象.
才能夠不斷地往積極的方向去發展自己.
Amen.
我還記得我小女兒在上大學的時候.
有一段時間她功課非常不好.
也沒有什麼興趣繼續讀下去.
但當時她人在紐約我們在這裡好像幫不上什麼忙.
我只有不斷地在禱告當中為她來祝福.
結果新的學期來到了.
有一位新的老師.
在上課的時候就拿起他的作業念了一遍.
念完之後對全班同學說.
這是我所見過最有創意的作業.

$^{281}$從此就改變了她整個學習的態度.
她的功課開始突飛猛進.
從這件事我就開始明白.
其實很多時候只要一點點的鼓勵.
就能夠帶給人多大的改變.
Amen.
反過來說如果今天我們對待我們的孩子.
總是很苛刻很嚴厲.
那我們就是好像傳達一個不接納的一個信息.
我們等於是為仇敵打開了一個門.
讓他把自卑缺乏安全感放在我們小孩的心中.
所以說父親我們口中的話語太重要了.
你口中說出祝福.
你的孩子就能夠領受祝福.
你口中說出咒詛.
他們便受咒詛.
我們的話語真的是可以影響他們一生的道路.
所以非常重要常常開口祝福你的家人.
第三要常為家人禱告.
各位同學要常常為家人禱告.
特別是我們今天活在這樣充滿危險不安的世代.
我想我們的心也都常常憂慮.
到底有誰能夠保護我們的家庭.
我就想到當年神怎麼樣呼召尼希米起來.
為耶路撒冷建造一個保護的城牆.
同樣我相信神今天也呼召我們要起來.
為我們的家庭建造一個保護的城牆.
那就是一個禱告的城牆.
特別是在我們中心做父親的.
你要知道你就是在這個家庭當中.
神所設立一個屬靈的頭.
這個屬靈的頭的意思不是叫你做暴君的.
而是一個屬靈的領導者.
所以你要常常為你的家人禱告.
也要常常帶領你的家人一起禱告.
你知道當時尼希米.
當他領受神的呼召.
回到耶路撒冷要建造耶路撒冷城牆的時候.
他面對當時仇敵厲害的攻擊.
他就召集以色列百姓.

$^{321}$在尼希米基地四當地十四節.
他對他們這樣說.
我們一同來唸來.
我查看了就起來對貴州官長和其他人說.
不要怕他們.
當紀念主是大而可畏的.
一定要為弟兄兒女妻子家產征戰.
他說不要怕他們.
因為我們的神是什麼樣.
大而可畏.
第五點你可知道今天我們所信的神是大而可畏的神嗎.
無論我們今天所遇到的環境多麼的艱難.
無論我們所遇到的所生活的環境充滿險惡動盪.
我們的神永遠是坐在寶座上掌權的神.
而且他已經把得勝的權柄賞賜給我們了.
這為什麼今天可以不必害怕.
這樣現今世代的黑暗權勢.
我們可以靠著神我們必然得勝.
但另一方面神也在這裡吩咐我們說.
要起來為我們的家庭征戰.
這段經文講得很清楚.
要起來為弟兄兒女妻子家產而征戰.
一方面要認識我們的神大而可畏.
另一方面要看見我們的責任.
我們需要起來為我們的家庭征戰禱告征戰.
換句話說不要再埋怨了.
不要再埋怨世界太黑暗.
不要再埋怨環境太險惡.
孩子又不聽話你要起來禱告.
Amen.
跟你們說你要起來禱告.
你一定要抓住神所給你的權柄.
為你的家庭來做宣告.
做父親的要非常明白你是屬靈的頭.
換句話說神給你一個權柄.
是個家庭的權柄.
今天只有你能夠站立在神所給你的權柄上.
斥責仇敵撒旦一切的作為.
能夠向仇敵宣告說這個家庭是屬於神的.
仇敵你退去吧.

$^{361}$Amen.
不但要為你的家庭禱告.
而且要帶領你的家人一起禱告.
你知道帶領你的家人一起禱告非常的重要.
我曾經讀到一篇的統計調查.
他特別問一個問題問很多基督徒一個問題.
有多少人你曾經記得你的父親與你一起禱告過.
有多少人你記得你的父親曾經與你一起禱告過.
結果不到10\%的人還記得父親曾經為他們禱告過.
我相信這個比例在華人信徒當中更低了更低了.
難怪很多時候我們在兒女的身上.
我們失去了屬靈的影響力.
保羅在前述第二章第八節怎麼勸勉我們的弟兄們呢.
這裡說保羅說我願男人怎麼樣.
無憤怒無爭論.
換句話說不要無緣無故亂發脾氣.
而是要怎麼樣.
舉起聖潔的手隨處禱告.
Hallelujah Amen.
也許說牧師我不夠聖潔所以不能禱告了.
No No No No.
今天你是主耶穌用寶血所買俗回來的.
你在神的眼前就是聖潔的.
在基督裡面我們能夠因性稱義.
神說當我們在神面平信心.
舉起我們的手的時候.
這個手就是聖潔的手.
所以說父親的人你要每天舉起你的手.
為家庭來禱告.
你要大聲宣告為你的家庭宣告.
願神的國降臨.
願神的旨意成就在我們家庭當中.
Amen.
當你這樣宣告的時候.
神就要為你興起.
神就要為你征戰.
仇敵就不能不退去的.
Amen.
這是為什麼大衛在四篇二十三篇第五節裡面.
他說這句話.

$^{401}$他說在我敵人面前你為我擺設宴席.
很特別在敵人面前打仗都來不及了.
怎麼神居然為我們擺設宴席呢.
而是什麼時候我們不再依靠自己去征戰呢.
當我們在這裡依靠神的時候.
舉起手禱告來仰望我們的主的時候.
神自己要為我們征戰.
不在我們需要去征戰.
而神要為我們征戰.
當神為我們征戰的時候.
我們就是能享受神為我們擺設的宴席.
Amen.
所以第三個很重要.
要起來為家人禱告.
最後第四.
要珍惜現在的時光.
各位要珍惜現在的時光.
我發現我們人啊.
常常就是活在對過去的懊悔當中.
特別是做父母.
當孩子小的時候.
你抱著他嫌累啊.
可是等孩子長大了.
你想再抱抱他.
他都不要.
你就懊悔了.
為什麼當初沒有多抱一點呢.
當孩子吵的時候.
你就嫌他煩啊.
可是當孩子上大學之後.
你天天盼他回家.
他都不回來.
你又懊悔了.
為什麼以前沒有花多點時間陪他呢.
可是現在好不容易他回了家.
你就開始嫌他懶喔.
又不起床.
又不整理.
你看我們做父母多煩啊.
多嘮叨啊.

$^{441}$你看看你.
這為什麼很多年前.
我自己做了一個非常重要的決定.
我永遠不再讓我自己.
活在對過去的懊悔當中了.
我要珍惜現在每一個時刻.
Amen.
你要知道很多時候我們做父母的人.
面對兒女的成長.
常常有一種.
只要熬到他長大就好了的心態.
就好像爬山的人一樣.
當你聽到上面的人說.
趕快上來.
這山林景色很美.
你就很急.
很想趕快上去.
也催大家趕快上去.
卻忽略了一件事情.
那就是沿途也很美咧.
的確孩子都會長大的.
事實上轉眼就長大了.
但沿途卻是你與孩子建立關係的時刻.
沿途也是你與他們一同成長的時刻.
Amen.
所以不要等到明天.
不要等到有一天.
今天只要神給我們機會.
今天就花時間陪伴他們.
Amen.
前一陣子我跟一位弟兄在聊天的時候.
他的孩子已經大學畢業.
做事很多年了.
他就告訴我.
有一件事情到現在.
他都感覺非常愧疚.
很遺憾.
那他兒子從小就喜歡打棒球.
他孩子一直很渴望.
能夠去看一場職業的棒球賽.

$^{481}$可是他自己一直忙碌.
從來沒有帶他兒子去看過一場職業棒球賽.
直到如今他都覺得非常的愧疚.
弟兄姊妹.
其實我告訴你很多事情.
其實不需要我們花很多時間.
但當我們肯花時間陪伴我們孩子的時候.
常常就把我們跟孩子的關係.
能夠軋得更近一點.
珍惜現在的時光太重要了.
夫妻之間也是一樣.
很多人說不像從前.
當然不像了.
隨著時間的過去.
他不同了.
你也不同了.
你要知道.
我有個好朋友.
他是心理醫生.
他曾經寫過一篇文章.
形容中年夫妻.
那個文章的題目就是.
禿頭王子胖仙女.
大家都不一樣了.
不但體態不同.
心態也都不同了.
但無論有多少不同.
最重要的是珍惜現在.
所以不要再埋怨了.
要學習為神所給你的每一個恩典.
我告訴你.
若想家庭和諧.
就不要等將來有一天.
現在你就要學習.
能夠營造一個充滿溫馨的.
安樂的一個環境.
如果你希望你家庭充滿喜樂.
就從你開始吧.
先喜樂起來你自己.
不要一天到晚愁眉苦臉滿口埋怨.

$^{521}$從你開始先喜樂起來吧.
如果你希望家庭充滿神的愛.
就從你自己開始.
去服侍你的家人.
不要一天到晚等你的家人服侍你.
如果你希望你家庭充滿了溫暖.
就從你開始吧.
不要一天到晚堅持己見.
什麼都要聽你的.
現在就去營造一個和諧喜樂的環境吧.
從現在開始.
珍惜每一個你能跟家人在一起的時間.
Amen.
Hallelujah.
第二點.
今天當我們生活在充滿混亂不安的世代當中.
我們常常會感覺教養兒女真的是不容易.
越來越不容易.
我告訴你.
特別是做父親的.
我想我們常常會呼喊.
我的幫助從何而來.
但今天你做一個神的兒女.
你一定要記得.
無論你多麼軟弱無助.
你是有一位愛你的天父.
你有一位愛你的主耶穌.
祂永遠與你同在.
而且祂就是你隨時的幫助.
Amen.
所以當你在教養兒女的事情上.
你感覺徬徨無助的時候.
不要灰心.
不要失望.
堅持的信靠神.
並且把兒女交託給神.
始終信靠神的人.
必不致羞愧的.
Amen.
我們一起禱告.

$^{561}$謝謝主.
你是那一位一生與我們同在的神.
你永遠不離開我們.
無論我們今天所遇到的環境多麼艱難.
我們所遇到的困境多麼難以解決.
我們永遠可以舉目仰望你.
你就是我們的依靠.
你就是我們的幫助.
特別為我們中間做父親的人來禱告.
很多時候在這樣忙碌.
充滿危險挑戰.
充滿各樣壓力的世代當中.
我們常不知道怎麼樣做一個父親.
你從來沒有撇下我們.
你是那永遠與我們同在的神.
我們可以隨時仰望你.
隨時來到你面前.
你就成為我們的力量跟幫助.
今天求你的聖靈充滿教官我們.
再次在每個人心中帶來安慰與力量.
Amen.
\newpage



\section{傳道書 5:8-6:9}
\label{sec:032335}
\textbf{錢值幾錢?}~~~~譚文倩莊傳道
\newline
\newline
連結: \href{http://www.church.com.hk/acms/content.asp?site=cdc\&op=show\&type=product\&code=032335\&layout=sermon}{\texttt{032335}} ~~~~ 語音日期: 2026-06-21
\newline
\newline
\hyperref[sec:032334]{\small{< < < PREV SERMON < < <}}
~
\hyperref[sec:index_chronic]{\small{[返順時目]}}
~
\hyperref[sec:032336]{\small{> > > NEXT SERMON > > >}}
\newline
\newline
傳道書 5:8-6:9
\newline
\begin{longtable}{cl}
\hline
\hline
章節 & 經文 (和合本修訂版)\\
\hline
5:8 & \begin{tabularx}{0.7\textwidth}{X} 你若在一個地區看見窮人受欺壓，公義公平被掠奪，不要因此驚奇；有一位高過居高位的在鑒察，在他們之上還有更高的。 \end{tabularx} \\ \\ \relax
5:9 & \begin{tabularx}{0.7\textwidth}{X} 況且地的益處歸眾人，就是君王也受田地的供應。 \end{tabularx} \\ \\ \relax
5:10 & \begin{tabularx}{0.7\textwidth}{X} 喜愛銀子的，不因得銀子滿足；喜愛財富的，也不因得利益知足。這也是虛空。 \end{tabularx} \\ \\ \relax
5:11 & \begin{tabularx}{0.7\textwidth}{X} 貨物增添，吃的人也增添，物主得甚麼益處呢？不過眼看而已！ \end{tabularx} \\ \\ \relax
5:12 & \begin{tabularx}{0.7\textwidth}{X} 勞碌的人不拘吃多吃少，睡得香甜；富人的豐足卻不容他睡覺。 \end{tabularx} \\ \\ \relax
5:13 & \begin{tabularx}{0.7\textwidth}{X} 我見日光之下有一件令人憂傷的禍患，就是財主積存財富，反害自己。 \end{tabularx} \\ \\ \relax
5:14 & \begin{tabularx}{0.7\textwidth}{X} 他因遭遇不幸，財產盡失；他生了兒子，手裡卻一無所有。 \end{tabularx} \\ \\ \relax
5:15 & \begin{tabularx}{0.7\textwidth}{X} 他怎樣從母胎赤身而來，也必照樣赤身而去；他所勞碌得來的，手中分毫不能帶去。 \end{tabularx} \\ \\ \relax
5:16 & \begin{tabularx}{0.7\textwidth}{X} 這是一件令人憂傷的禍患。他來的時候怎樣，去的時候也必怎樣。他為風勞碌有甚麼益處呢？ \end{tabularx} \\ \\ \relax
5:17 & \begin{tabularx}{0.7\textwidth}{X} 並且他終身在黑暗中吃喝，多有煩惱、病痛和怒氣。 \end{tabularx} \\ \\ \relax
5:18 & \begin{tabularx}{0.7\textwidth}{X} 看哪，我所見為善為美的，就是人在神賜他一生的日子吃喝，享受日光之下勞碌得來的好處，因為這是他應得的報償。 \end{tabularx} \\ \\ \relax
5:19 & \begin{tabularx}{0.7\textwidth}{X} 而且，一個人蒙神賞賜財富與資產，又使他能享用，能獲取自己當有的報償，在他的勞碌中喜樂，這是神的賞賜。 \end{tabularx} \\ \\ \relax
5:20 & \begin{tabularx}{0.7\textwidth}{X} 他不多思念自己一生的日子，因為神使他的心充滿喜樂。 \end{tabularx} \\ \\ \relax
6:1 & \begin{tabularx}{0.7\textwidth}{X} 我見日光之下有一件禍患重壓在人身上， \end{tabularx} \\ \\ \relax
6:2 & \begin{tabularx}{0.7\textwidth}{X} 就是人蒙神賜他財富、資產和尊榮，以致他心裡所願的一樣都不缺，只是神使他不能享用，反被外人享用。這是虛空，也是禍患。 \end{tabularx} \\ \\ \relax
6:3 & \begin{tabularx}{0.7\textwidth}{X} 人若生一百個兒子，活許多歲數；他即使壽命很長，心裡卻不因福樂而滿足，又不得埋葬；我說，那流掉的胎比他倒好。 \end{tabularx} \\ \\ \relax
6:4 & \begin{tabularx}{0.7\textwidth}{X} 因為這胎虛虛而來，暗暗而去，名字被黑暗遮蔽， \end{tabularx} \\ \\ \relax
6:5 & \begin{tabularx}{0.7\textwidth}{X} 而且沒有見過天日，甚麼都不知道，這胎比那人倒享安息。 \end{tabularx} \\ \\ \relax
6:6 & \begin{tabularx}{0.7\textwidth}{X} 那人雖然活千年，再活千年，卻不能享福；眾人豈不都歸同一個地方去嗎？ \end{tabularx} \\ \\ \relax
6:7 & \begin{tabularx}{0.7\textwidth}{X} 人的勞碌都為口腹，心裡卻不知足。 \end{tabularx} \\ \\ \relax
6:8 & \begin{tabularx}{0.7\textwidth}{X} 智慧人比愚昧人有甚麼益處呢？困苦人在眾人面前知道如何行，有甚麼益處呢？ \end{tabularx} \\ \\ \relax
6:9 & \begin{tabularx}{0.7\textwidth}{X} 眼睛所看的比心裡妄想的倒好。這也是虛空，也是捕風。 \end{tabularx} \\ \\
[1ex]
\hline
\hline
\end{longtable}
$^{1}$終於來到傳道書的中間.
猜猜今天這段經文.
傳道者要跟我們講什麼.
我在練習的時候發覺.
這東西已經講給大家聽今天要講什麼.
By the way 如果你今天是第一次在我們教會.
來到我們教會那裡.
我想代表教會歡迎你.
完了崇拜之後歡迎你留步.
讓我們大家可以有些互相認識的時間.
OK 回到傳道書.
傳道書是聖經舊約裡面.
其中一本智慧文學.
探討人生的終極意義.
道德的實踐.
苦難的奧秘.
和我們人和做物主之間的關係.
今天的經文.
在傳道書的五章八節.
去到六章九節.
傳道書的作者今天要跟我們講錢.
你聽到這個主題的時候.
是不是馬上醒了.
還是馬上在想.
哇 講錢是傷感情的.
我不是很怕傷感情.
因為不關我的事.
不是我講的.
是傳道者講的.
我只是依經直說.
不如你又想想錢有什麼好說的.
在這段經文裡面.
傳道者不是站在道德高地那裡.
去跟我們說錢是好的.
或者錢是不好的.
而是他要藉著講錢這件事.
要我們思考一下.
喜歡錢的人不會因為錢而滿足.
錢多又如何.
你知不知道可以是分分鐘得個體字.

$^{41}$接著還要慢慢失眠.
為什麼呢.
有賺錢的能力.
又懂得存錢.
但可以無神神發生一些事.
一鋪清袋.
有財有勢一無所缺.
但卻沒有享受的能力.
很殘忍.
就好像回到第一章那裡說.
虛空的虛空.
一直在聽.
不知道你會不會覺得很depressing.
什麼這麼慘.
其實傳道者在這段經文裡面.
有什麼好說的.
有沒有好的東西在裡面說的呢.
是有的.
妹妹就在五章十八召二十節那裡說過.
「看那我所看見為善為美的.
就是人在神賜他一生的日子吃喝.
享受日光之下勞碌得來的好處.
因為這是他應得的報償.
而且一個人蒙神賞賜財富與資產.
又使他能享用.
能獲取自己當有的報償.
在他的勞碌中喜樂.
這是神的賞賜.
他不多思念自己一生的日子.
因為神使他的心充滿喜樂」.
意思是叫我們不用想太多.
又不用計算太多.
總之有得吃你就吃.
有得喝你就喝.
有工返就返工.
開開心心過日子.
存著感恩的心去領受.
就叫做懂得享受上帝給你的恩典.
如果我告訴你.
不是.

$^{81}$至少這個不是重點.
那你猜猜甚麼是重點呢?.
不如我們先一起低頭祈禱.
我們求聖靈幫助我們去聆聽神的話語.
我們人算得是甚麼.
你竟然眷顧我們.
求你幫助我們.
我們渴望聽到你的聲音.
奉主耶穌基督成名祈求.
1980年.
林子祥就出了一首歌.
叫做「分分鐘需要你」.
懂不懂唱?.
懂嗎?.
「願你會zar 火箭帶我到天空去」.
「在太空中兩人住」.
以前聽可能就想.
沒甚麼可能.
現在去太空住好像真的可能有可能.
然後後面說甚麼?.
「有了你開心D  乜都清心滿意」.
「鹹魚白菜又好好味」.
是嗎?.
每天吃會不會很好味?.
現在經常說鹹魚是一個processed food.
吃得多會致癌.
這首歌說甚麼呢?.
它是說不需要生生海味.
不需要豪華的生活.
只要和你一起.
我就滿足.
五章十字十二字.
「喜愛銀紙的不因得銀紙滿足」.
「喜愛財富的也不因得利益知足」.
「這也是虛空」.
「貨物增添吃的人也增添」.
「物主得甚麼益處呢?不過眼看而且」.
「勞碌的人不拘吃多吃少」.
「睡得香甜 婦人的豐足卻不容他睡覺」.
「錢無問題」.

$^{121}$「鍾意錢都無問題」.
「唔滿足先至係個問題」.
Jandith Rockefeller.
美國有史以來其中一個最有錢的人.
曾經有人問她.
「仲要賺多幾多錢先叫夠?」.
她說「再多D D 就夠」.
所以喜愛銀紙的人.
不會因銀紙滿足的問題.
是在於人那種「多D D 先夠」的心態.
永遠都不滿足.
還記不記得史陶寶洛在提摩泰前書6章10節.
說貪錢財是萬惡之根.
為甚麼.
原因是我們人和錢建立了一個錯誤的關係.
這裡說的喜愛銀紙.
有一種朋友的概念在裡面.
記不記得之前我說過「ahalf」這個字.
一個深層忠誠跟你同行患難的關係.
陪你 撐你.
在你有困難的時候不會離開你.
問題是人會否當錢是自己最好的朋友.
就是那個你可以依靠.
在你困難的時候幫助你的關係.
你期望錢對你忠心.
滿足你的需要.
在你最需要人陪你 撐你的時候不離開你.
一種錯誤的關係.
經文繼續說.
「貨物增添 吃的人也增添.
物主得甚麼益處.
不過眼看如此」.
換句話說.
當你的財富增加的時候.
圍著你 消耗你的資源的人也會增加.
曾賢也有一個很類似的提醒.
他說「財物使朋友增多.
窮人連朋友也遠離他.
人人都討好慷慨的人.
送禮的人 人人都與他為友」.

$^{161}$現實是當一個人沒有資源的時候.
可以連最親密的關係也疏離.
為甚麼.
因為他已經沒有被需要的價值.
所以財富帶來的是朋友.
還是期望和索取呢.
當你有資源的時候.
你要照顧更多的人.
分配更多的資源.
應付更多的需求.
甚至維持更多的關係.
相對如果你沒有或沒甚麼的時候.
其他人對你沒甚麼期望.
可能你跟對方的關係還來得簡單一點.
「勞碌的人不拘吃多吃少 睡得香甜.
富人的風築卻不用他睡覺」.
心理學上有一個很簡單的觀察.
一個人的壓力大不大.
你看他晚上睡得好不好你就知道.
傳道者在這裡做了一個對比.
他說一個辛勤工作.
但是生活簡單的人.
他不是很富有.
但他有夠好睡.
擁有多的人.
因為背著很多掛慮.
就算回到家的時候.
他的腦袋仍然不停在運作.
那個交易最後成不成交.
有沒有漏了甚麼.
成不成交.
如果成不成交怎麼辦.
對手會出甚麼招數.
換句話來說.
財富不會令你沒問題.
而是令你不停處理問題.
喜歡錢的人不會因為此而滿足.
分分鐘還慢慢失眠.
13至17節那裡說.
「我見日光之下有一件令人憂傷的禍患.

$^{201}$就是財主積存財富反害自己.
他因遭遇不幸財產盡失.
他生了兒子手裡卻一無所有.
他怎樣從母胎赤身而來.
也必照樣赤身而去.
他所勞碌得來的.
手中薰毫不能帶去.
這是一件令人憂傷的禍患.
他來的時候怎樣.
去的時候又怎樣.
他為勞碌有甚麼益處.
並且他終身在黑暗中吃喝.
多有煩惱病痛和怒氣」.
不知道大家會否像我一樣.
還經常留意香港的新聞.
香港最近有幾宗很大的財經新聞.
震撼彈.
其中一個就是抱黃波叔.
他在21年離開之後.
他的兒子接管了他的家族生意.
他在疫情之前買了很多舊酒店.
打算翻身打造酒店的王國.
但遇上疫情的關係.
大量的資金被扣住.
加上樓市大跌.
他們家族用了四五十年.
建立回來的八百億的王國遺產.
不斷在縮水.
今年2026年.
他被兩間公司申請破產.
發生這件事大約是四五年.
用了四五十年建立回來的.
八百億公司的王國.
四五年被申請破產.
另一個鄭志剛.
曾經是香港第三大富豪.
鄭裕彤的孫子.
他主政期間大撒金錢搞藝術.
興建你們可能有些人都去過的K11 MUSE.
結果新世界發展在2024年.

$^{241}$蝕了171億.
是20年來首次虧損.
市值蒸發了720億.
他因遭遇不幸財產盡失.
投資失敗遇上疫情.
生意蝕本.
八百億 什麼錢.
八百萬我都沒見過.
八百億好難想像.
天災人禍.
被人騙 被人搶.
我們這個世界可以這麼無厘頭.
可以發生一些你估都估不到.
好像沒什麼可能會發生的事.
《馬太福音六章》有一句經文.
可能我們當中有些人都很熟.
他說不要為自己積攢財寶在地上.
答真一點這句經文在想的時候.
思考一下這些新聞.
你都不是說不害怕.
幾百億.
什麼都沒了.
我在想.
他們沒了東西會怎樣 找個四百呎單位住.
好諷刺 怎麼接受.
就是這麼瘋狂.
十六字 這是一件令人憂傷的禍患.
可能你們要幫我 我不知道為何不動.
這是一件令人憂傷的禍患.
他來的時候怎樣 去的時候也怎樣.
他為勞碌有什麼益處.
傳道者批判的不是反對努力.
工作或者追求卓越.
事實上如果你記得的時候.
傳道者之前也多次在這裡肯定.
勞碌本身是神給我們的其中一個恩賜.
我們人也都應該去享受.
那個勞碌之後所得回來的份.
但在這裡他真正是想指出.
我們人生充滿很多我們控制不了的變數.

$^{281}$很多超出我們可以控制的事情.
悲劇是會發生的.
損失是會出現的.
計劃都可以被打亂.
It happens.
問題不是在說我們有沒有努力.
而是我們要Better think twice.
我們的安全感.
是否建立在這些東西上面.
在這些財富或者成就.
我們想著可以建立回來.
給我安全的這些東西上面.
6章1-9字.
我見日光之下有一件禍患重壓在人身上.
就是人夢神賜他財富資產和尊榮.
以致他心裡所願的一樣都不缺.
只是神賜他不能享用.
反被愛人享用.
這是虛空也是禍患.
人若生了一百個兒子.
活許多歲數.
他即使壽命很長.
心裡卻不因福樂而滿足.
有不得埋葬可說那流掉的胎被他討好.
因為這胎虛虛而來暗暗而去.
明知被黑暗掩蔽.
而且沒有見過天日.
什麼都不知道.
這胎被那人倒向安息.
那人雖然活千年再活千年.
卻不能享福.
眾人豈不都歸同一個地方去嗎.
人的勞碌都為後福.
心裡卻不知足.
智慧人比愚昧人有什麼益處.
困苦人在眾人面前.
知道如何行有什麼益處.
眼睛所看的比心裡妄想的都好.
這也是虛空.
也是暴風.

$^{321}$一個人可以擁有很多很多.
但若神沒有給他享受的能力.
他其實與沒有沒有分別.
之前我跟大家說的John D. Rockefeller.
美國有史以來其中一個最有錢的人.
在他大約五十歲的時候.
因為他長期工作的壓力.
他的健康開始出現問題.
他經常周身痛.
脫髮.
消化系統出現問題.
他去看醫生.
醫生預約他什麼.
就是喝奶和吃餅乾.
就算身體和腸胃康復.
他只可以吃很小的青豆.
生菜.
魚.
Brown rice, big potato.
其中一個最有錢的人.
想吃什麼都可以.
卻沒有能力享受.
那個年代的人大約五十歲.
但John D. Rockefeller.
活到九十七歲.
很多人都很渴望長命.
但他身體一直有很多狀況.
在他身邊跟他很親的人說.
John D. Rockefeller是長期失眠.
長期沒有笑容.
對任何事情都提不起勁.
財富資產一樣都不缺.
但卻沒有享受的能力.
這也是虛空.
也是暴風.
回到剛才說的五章十八至二十節.
我們再思考一下這節經文.
(看那我所見為善為美的,就是人在神賜他一生的日子吃喝,享受日光之下勞碌得來的好處,因為這是他應得的補償,而且一個人蒙神賞賜財富與資產,又使他能享用能獲取自己當有的補償,在他的勞碌中喜樂,這是神的賞賜,他不多思念補償.).
不知道你有沒有留意這裡.
有些很重要的教導在這裡.

$^{361}$(看那我所見為善為美的,就是人在神賜他一生的日子吃喝,享受日光之下勞碌得來的好處,因為這是他應得的補償,而且一個人蒙神賞賜財富與資產,又使他能享用能獲取自己當有的補償,在他的勞碌中喜樂,這是神的賞賜.).
(還有,你若在一個地區看見窮人受欺壓,公義公平被掠奪,不要因此驚奇有一位居高位的在監察,在他們之上還有更高的,況且地的益處歸眾人,就是君王也受田地的供應.).
就是君王也受田地的供應,都是蒙神的賞賜,都是神賜他的結果..
結束之前,我想跟大家分享一個我在新聞裡看到的故事..
我之前跟大家說過,香港雲福縣在去年十一月發生一場大火災,八十七歲的周婆婆,她就住在其中一棟大廈裡面的十四樓,一場大火,燒毀了她住了二十多年的家,大半生人的財物..
上個月,周婆婆終於可以在她的孫子陪同之下,回到家裡,沒有升降機,撐了十四層樓..
撐到上去的時候滿頭大汗,她周圍看一下自己的家,已經被燒到黑了,被煙燻到黑漆漆,門窗都破爛了,火災之後留下一陣糯味..
她走去客廳的角落,她自己的老朋友那部琴的前面,戴著手套坐下,就開始彈那首她自己很喜歡的詩歌,《耶和華是愛》..
她一邊彈一邊唱,她還問孫子,你有沒有跟著我唱呀?孫子說,有呀有呀..
原來周婆婆一早已經帶好歌詞,準備回家的時候跟她的家人一起唱..
她說《耶和華是愛》這首詩歌,是她最喜歡,最熟悉的詩歌..
上樓之前,有記者問周婆婆,她說大火已經燒了你所有的東西,你有什麼感受?.
婆婆就說,是呀是呀,很多東西都燒了,不過我現在所需要的東西,我都有,上帝的恩典夠我用..
所以我的心很平靜,我很感恩..
她又說,苦難是魏莊的祝福,因為這場大火之後,有很多人去關心她,.
甚至她一個以前的學生,特意從上海飛來香港去探望她,給她錢..
婆婆在房間裡找找找,其中一樣她說她很想找到的是什麼呢?.
是她藏在房間裡一堆印著耶穌愛你的毛巾,她說我很想找到這堆毛巾,.
因為每一次我出去傳福音的時候,我就會將這些毛巾送給那些福音對象..
錢,值幾錢?我們現在是活在人類歷史上最富裕的世代,.
但為什麼我們不是一個最滿足的世代?.
如果我們人生所有問題都是可以用錢去解決的,.
那似乎我們有錢的人就應該活得很快樂..
為什麼有人什麼都有卻是享受不到?.
反過來有些人沒有什麼,卻活得很開心很滿足呢?.
傳道書五章八節到六章九節,傳道者要教我們認清一個事實,.
我們所擁有的全部都是神給的,是神賜予的,.
甚至是神給我們的能力去享受,.
真正的滿足不是擁有更多更多,.
而是我們要學習怎樣從神的手中去領受..
一個八十七歲的老人家,在失去所有日光之下的財物之後,.
你能不能想像他仍然可以笑著坐在琴面前唱「耶和華是愛」?.
神保守我,我什麼都沒有了,但我很感恩,.
神的恩典夠我用,他心裡記掛的是那些還沒聽到福音,.
他很想派一條毛巾給他的人..
周婆婆的一生沒有建立在十四樓那層樓上的家裡,.
她的家當裡,沒有在她的錢和金器上,她的人生卻是建立在一個大火怎樣燒都燒不到的事實上,.
「耶和華是愛」,夢神賞賜的滿足,神使人心充滿喜樂,Amen?.
\newpage



\section{馬可福音 2:1-12}
\label{sec:032336}
\textbf{爸爸的心 (粵/華)}~~~~梁廷益牧師
\newline
\newline
連結: \href{http://www.church.com.hk/acms/content.asp?site=cdc\&op=show\&type=product\&code=032336\&layout=sermon}{\texttt{032336}} ~~~~ 語音日期: 2026-06-21
\newline
\newline
\hyperref[sec:032335]{\small{< < < PREV SERMON < < <}}
~
\hyperref[sec:index_chronic]{\small{[返順時目]}}
~
\hyperref[sec:032337]{\small{> > > NEXT SERMON > > >}}
\newline
\newline
馬可福音 2:1-12
\newline
\begin{longtable}{cl}
\hline
\hline
章節 & 經文 (和合本修訂版)\\
\hline
2:1 & \begin{tabularx}{0.7\textwidth}{X} 過了些日子，耶穌又進了迦百農。人聽說他在屋裡， \end{tabularx} \\ \\ \relax
2:2 & \begin{tabularx}{0.7\textwidth}{X} 於是許多人聚集，甚至連門前都沒有空地；耶穌就對他們講道。 \end{tabularx} \\ \\ \relax
2:3 & \begin{tabularx}{0.7\textwidth}{X} 有人帶著一個癱子來見耶穌，是由四個人抬來的； \end{tabularx} \\ \\ \relax
2:4 & \begin{tabularx}{0.7\textwidth}{X} 因為人多，無法抬到耶穌跟前，就把他所在那房子的屋頂拆了，既拆通了，就把癱子連所躺臥的褥子都縋下去。 \end{tabularx} \\ \\ \relax
2:5 & \begin{tabularx}{0.7\textwidth}{X} 耶穌見他們的信心，就對癱子說：「孩子，你的罪赦了。」 \end{tabularx} \\ \\ \relax
2:6 & \begin{tabularx}{0.7\textwidth}{X} 有幾個文士坐在那裡，心裡議論，說： \end{tabularx} \\ \\ \relax
2:7 & \begin{tabularx}{0.7\textwidth}{X} 「這個人為甚麼這樣說呢？他說褻瀆的話了。除了神一位之外，誰能赦罪呢？」 \end{tabularx} \\ \\ \relax
2:8 & \begin{tabularx}{0.7\textwidth}{X} 耶穌心中立刻知道他們心裡這樣議論，就說：「你們心裡為甚麼這樣議論呢？ \end{tabularx} \\ \\ \relax
2:9 & \begin{tabularx}{0.7\textwidth}{X} 對癱子說『你的罪赦了』，或說『起來！拿你的褥子行走』，哪一樣容易呢？ \end{tabularx} \\ \\ \relax
2:10 & \begin{tabularx}{0.7\textwidth}{X} 但要讓你們知道，人子在地上有赦罪的權柄。」就對癱子說： \end{tabularx} \\ \\ \relax
2:11 & \begin{tabularx}{0.7\textwidth}{X} 「我吩咐你，起來！拿你的褥子回家去吧。」 \end{tabularx} \\ \\ \relax
2:12 & \begin{tabularx}{0.7\textwidth}{X} 那人就起來，立刻拿著褥子，當著眾人面前出去了，以致眾人都驚奇，歸榮耀給神，說：「我們從來沒有見過這樣的事！」 \end{tabularx} \\ \\
[1ex]
\hline
\hline
\end{longtable}
$^{1}$大英姐妹早上好.
首先我多謝Carmen幫我翻譯.
我講的不是廣東話.
我的是香港話.
所以Carmen我都跟她說.
她盡量翻譯到意思就可以了.
不需要逐句翻譯.
我盡量翻譯出來.
我們今天慶祝父親節.
我們慶祝父親節.
其實有得慶祝.
就好把握這個機會.
有得慶祝的話.
我們就要把握這個機會.
可能很快就沒有了這個日子.
可能很快就沒有這個日子了.
那你說大家勵士嗎.
什麼意思呢.
你要咒我爸爸死是不是.
你詛咒我們嗎.
但其實現在已經開始.
幾年前已經開始一個叫做.
反對父親節母親節的運動.
幾年前有一個開始反對父親節.
跟母親節的運動.
覺得這些所謂的父親節母親節.
這些有性別的立場的節日.
是不夠包容性.
不夠inclusive.
因為他們覺得父親節母親節.
這個比較有性別的分別的節日的話.
不夠包容性.
所以就應該改為這個叫做family day.
或者家庭節或者叫做parents day.
就是家長節.
所以想把它改成了像是一個家庭.
就是family day或者父母節.
大家都知道在加拿大和美國.
有些學校是不鼓勵學生.
叫自己的爸爸媽媽叫做父母.

$^{41}$在加拿大或者美國.
有些學校已經不鼓勵他們的孩子.
叫爸爸媽媽父母.
只可以叫他們做家長.
只叫他們家長.
因為你知道現在有一些.
同性婚姻的雙方.
都是可以領養小朋友的.
因為現在很多同性婚姻的雙方.
他們都可以領養小孩子.
所以如果他的小朋友.
是有兩個男性做家長的.
那你怎麼慶祝母親節呢.
所以如果他們一個家裡.
有兩個父親的話.
他怎麼樣慶祝呢.
所以就開始不鼓勵做這件事.
所以他們就有些人.
開始不鼓勵做這件事.
因為做這些事就不夠包容性.
因為這樣子的話就不夠包容性.
所以我也見過一些.
所謂包容性的神學家.
在聖經裡面找出很多.
上帝是女人的這些形象.
所以很多說自己.
比較有包容性的神學家的話.
他也會嘗試在聖經裡面找出.
聖經是女人的.
那究竟上帝是男人還是女人呢.
那到底上帝是男人還是女人呢.
我只知道聖經告訴我們.
上帝做人是做男做女.
我只知道聖經告訴我們.
上帝造人是造男造女.
所以希望大家繼續把握機會.
去慶祝父親節.
希望大家把握機會去慶祝父親節.
那父親節Jackie就交代我.
你們開始說馬可福音的訊息.

$^{81}$那父親節Jackie就交代我.
你們現在在說馬可福音的訊息.
那今天是說到馬可福音第二章一至十二節.
今天是說到馬可福音二章一到十二節.
當我拿到經文的時候.
我就在思考跟父親節有什麼關係呢.
所以當我拿到經文的時候.
我就在思考這跟父親節有什麼關係呢.
我們首先請Carmen幫我們用國語.
先讀出整段的經文.
就是馬可福音第二章第一節開始.
Carmen麻煩你.
我現在把它讀出來.
過了些日子耶穌又進了加白濃.
人聽說祂在屋裡.
於是許多人聚集.
甚至連門前都沒有空地.
耶穌就對他們講道.
有人帶著一個攤子來見耶穌.
是有四個人抬來的.
因為人多無法抬到耶穌跟前.
就把祂所在的那房子的房頂拆了.
一拆通了就把攤子連鎖.
躺臥的褥子都吹下去.
下一章.
耶穌見他們的信息.
就對攤子說.
孩子你的罪赦了.
有幾個文士坐在那裡.
心裡議論說.
這人為什麼這樣說呢.
他說褻渎的話了.
除了上帝一位之外.
誰能赦罪呢.
耶穌心中立刻知道他們心裡這樣議論.
就說你們心裡為什麼這樣議論呢.
對攤子說.
你的罪赦了.
話說起來.
拿你的褥子行走.

$^{121}$哪一樣容易呢.
但要讓你們知道.
人子在地上有赦罪的權柄.
就對攤子說.
我吩咐你.
起來.
拿你的褥子回家去吧.
那人就起來.
立刻拿起褥子.
當著眾人面前出去了.
以致眾人都驚奇.
歸容要給上帝.
說我們從來沒有見過這樣的事.
當你看完經文的時候.
我們首先去看一下.
今天我們的父親.
對子女的心是怎麼樣.
當我看到這個經文的時候.
我們首先想要來看一下.
父親對孩子的心是怎麼樣子的.
有人說養兒一百歲長憂九十九.
憂就是憂慮的憂.
不是休息的憂.
憂是憂慮的憂.
不是休息的憂.
究竟你最憂心你的子女是什麼呢.
你最憂心的你的子女是什麼呢.
如果當中還沒有父親.
希望有一天你都成為父親的時候.
你都了解.
今天很多人對子女的憂慮憂心是什麼.
當中有人還不成為父親的話.
你可以了解一下.
現在的人父親對孩子的憂慮是什麼.
Peer Research是一個很出名的研究機構.
2023年的時候發表了一個報告.
說父母最憂慮,擔憂子女的最多是什麼.
在2022年的時候.
他就發出了一個這樣的報告.
說父母最擔憂子女的是什麼.

$^{161}$我不知道你們可能最憂心的.
就是孩子考不到一間好的學校.
可能你們最憂心的就是.
你的子女考不上一個好學校.
或者年輕的孩子晚上開車出去.
怕他發生意外.
或者年輕的孩子晚上開車出去的時候會發生意外.
其實Peer Research告訴我們.
最擔憂子女的.
是最怕孩子有depression,憂鬱.
這個研究出來的時候.
最擔憂的子女他們最擔心的是他會憂鬱.
接著是被bully,欺凌,被霸凌.
第二個就是被霸凌.
或者有些可能比較有錢的.
怕被人綁架.
或者有一些比較有錢的話就擔心被綁架.
酒精問題,被槍擊.
有20\%的父母怕子女在學校被槍擊.
還有22\%的怕子女在學校被槍擊.
接著就是被人弄大肚子.
接著就是懷孕或者試別人懷孕.
或者跟警察發生麻煩.
究竟你憂慮的子女是什麼呢?.
究竟你憂慮子女的是什麼呢?.
作為一個爸爸.
我想我都有30年的做爸爸的經驗.
作為一個父親的話.
我也有30年做父親的經驗.
我一會兒可以跟你分享一下.
我兩個女兒的一些往事.
我也可以跟你分享一下.
我有兩個女兒的往事.
第一件事我們要明白的就是.
我們要看見這段經文.
第一件事我們要明白的就是.
看到這段經文.
看見這段經文的時候.
我們看到這四個朋友.
四個朋友帶著個毯子將它摘下.

$^{201}$給耶穌去醫治.
這段經文裡面我們看到.
四個朋友把毯子摘下去.
給耶穌去醫治.
當時的房子就不是現在這些石屎.
你怎麼拆那個屋頂.
當時是一些很簡陋的.
可能是一些草的天綿.
很容易就拆了下來.
當時的房子就不像現在的水泥房子.
當時可能是用草或者這些泥做的.
所以就很容易可以把它的屋頂拆下來.
但是我們要看的是.
耶穌祂所做的行動和祂所說的話.
怎麼去表達天父對我們的心.
還有我們應該怎麼去學習.
天父對我們的心.
所以我們要從耶穌當時的行為來學習.
天父對我們的心是什麼.
然後我們怎麼樣去學習天父的心.
我想我們自己也有很多對孩子的夢想.
我們自己也有很多對孩子的夢想.
我們很想他們活出我們的夢想.
很多家長希望他們的子女.
做一些什麼職業.
進哪個學校等等.
很多家長希望他們的兒子.
做什麼職業.
進什麼學校等等.
又或者跟哪一個的結婚.
哪一個類型的人結婚.
或者跟誰結婚.
跟什麼類型的人結婚.
我以前當我岳父岳母未信主的時候.
當我岳父岳母沒有信主的時候.
就不准我太太嫁給我.
不是因為我人格不行.
第一樣是因為我窮.
因為他想他女兒過得好一點.
因為他想他女兒過好一點.

$^{241}$第二因為我是做傳道人.
第二是因為我做傳道人.
因為我岳父以前在大陸文革的時候是黑五類.
因為我岳父在文革的時候是.
黑五類.
所以他教書的.
所以經常被人抓去監房.
因為他是教書的時候.
他經常被人抓去坐牢.
他說他坐牢的時候.
大部分和他一起坐的都是傳道人牧師.
他坐牢的時候.
大部分人跟他坐牢的都是傳道人.
所以他說你做傳道人很危險.
所以他就覺得做傳道人是很危險的.
第三個原因就是因為我要回香港.
第三個原因就是因為我要回去香港.
我們這麼辛苦才從大陸走出來.
你這麼辛苦才從大陸走出來.
你又要把我女兒帶回香港.
又要把我女兒帶回香港.
當然不願意.
所以他就不願意.
他甚至這樣說.
他甚至這樣講.
他說我家裡有一支槍.
如果你硬要和我女兒結婚的話.
我就射你.
如果你硬要和我女兒結婚的話.
我就一槍斃了你.
但是最後我岳父信了主.
最後我岳父信了主.
而且公然說我是他最好的女婿.
然後公然的講我是他最好的女婿.
因為他只有一個女兒.
我當然是最好的.
因為他只有一個女兒.
所以他是最好的一個.
但是其實我也對女兒有很多自己個人的夢想.
但是其實我對我自己女兒的話.

$^{281}$也有很多個人的夢想.
我爸爸也對他兩個孫女有很多的夢想.
我爸爸對他兩個孫女也有很多的夢想.
香港的時候就不斷逼那兩個女兒去彈琴.
在香港的時候一直逼他們去彈琴.
因為彈琴就是可以服侍.
香港很多的小朋友都是去彈琴.
因為彈琴就是可以服侍.
香港很多的小朋友都是去彈琴.
他們一直考到六級琴都已經彈得挺好的.
但是當我爸爸死了之後.
就當我爸爸死了之後.
沒有再碰過那個琴.
他們就再也沒有碰過琴了.
我琴在香港運貨櫃搬了來美國.
碰都沒有碰過.
然後我把琴從香港搬到美國來.
他們碰都沒有碰過.
我就問他們兩個為什麼不彈琴呢.
我大女兒現在就跟我說.
我大女兒就跟我說.
我現在喜歡爬龍舟.
她是這個南灣龍舟隊的隊長.
她是南灣龍舟隊的隊長.
現在練得一身肌肉.
現在是一身肌肉.
經常出去比賽.
有時候去Fenix 有時候去西雅圖爬龍舟.
端午節龍舟賽她又去爬龍舟.
經常飛各地去爬龍舟.
我小女兒我看她彈琴都是喜歡文靜的東西.
但她去到中學之後不再彈琴.
她去到中學之後不再彈琴.
她喜歡打籃球.
她做了學校的校隊.
現在長大了在洛杉磯工作.
我叫她拾起彈琴可以在教會服侍.
她說她現在不喜歡彈琴.
但我學到一樣東西就是陪伴.
她會記得你是陪伴她走過人生的路.

$^{321}$我剛來美國的時候面對一個很大的困難.
就是小女兒接送的問題.
當時搬到San Jose.
因為在三藩市塞車回來.
差不多三點鐘.
因為在三藩市塞車回來.
差不多三點鐘.
因為在三藩市塞車回來.
差不多三點鐘.
我差不多每天.
就是這樣來回.
我差不多每天.
就是這樣來回.
去接送我女兒.
但經過這個陪伴.
我覺得我和女兒的感情.
是好很多.
所以經過這個陪伴.
我覺得我和女兒的感情.
是好很多.
等一下.
孩子最大的問題.
孩子最大的祝福.
不是有一個完美的父親.
這種姐妹的孩子.
最大的祝福.
不是應該有一個完美的父親.
而是一個願意參與.
他生命的父親.
而是一個願意參與.
他生命的父親.
我鼓勵大家.
多點陪伴你們自己的子女.
我鼓勵大家.
多點陪伴你們自己的子女.
我鼓勵大家.
多點陪伴你們的子女.
因為這些是錢買不到的感情.
這些感情是錢都買不到的.
但是還有一樣更重要的就是與神同工不單止是做爸爸的與神同工.

$^{361}$另外更重要的是與神同工不單止是爸爸與神同工.
更加我們需要有其他的肢體.
基督徒的肢體去和我們的子女同工.
我們更需要的是其他基督徒的肢體也跟我們的子女同工.
有一個神學家在庫拉的青少年的學院.
是負責導師的主任.
他叫卡拉豹.
有一個神學家在這個庫拉的學院裡面做一個院長.
他寫了很多本書做了很多研究.
他說有很多年輕人的小朋友.
為什麼他們沒有放棄信仰呢?.
為什麼有很多年輕的小朋友他們沒有放棄信仰呢?.
或者為什麼他們會放棄信仰呢?.
很多是繼續成長沒有放棄信仰.
在信仰裡面很堅持的一班年輕人.
有一班很多是在信仰裡面很堅持的.
他就發現至少有五個屬靈的長者和他們一同的同行.
他就發現在這班年輕人裡面.
至少有五個屬靈的長者在他們人生裡面跟他們同行.
不單止是一個使徒贊穆斯.
不只是一個Pastor Zen.
也都是有另外多四個.
總共有五個可能是長者或者是師兄等等.
是陪伴他同行.
還至少有另外四個屬靈的長者陪他們同行.
所以這個就是他所說的.
很多本書你自己都可以回去看看.
叫五比一的原則.
這就是他很多本書裡面講到的五比一的原則.
五比一的原則就是.
一個小朋友至少要有五個屬靈長者.
陪他一同走這個信仰的路.
就是一個小孩子至少有五個屬靈的長者陪他走他的路.
我想問一下各位的爸爸.
我想問各位的爸爸.
你的小朋友在教會裡面.
或者在他人生裡面有沒有五個基督徒.
可以陪伴他一同走他們信仰的路呢.
你的小朋友呢.
有沒有至少在教會裡面有五個屬靈的長者陪伴他走呢.

$^{401}$如果沒有的話.
我覺得你就需要去做這個配合了.
如果沒有的話.
我就需要你去做一個配合.
我的盼望就是每一個小朋友都有這五個人去幫助他們成長.
我們盼望這些小朋友都有這五個人可以幫助他們成長.
這個壇子有四個朋友幫助他走到耶穌面前.
這個壇子是有四個朋友幫他走到耶穌面前.
我們都需要有信仰的同行者陪伴我們的小朋友.
我們也需要信仰的同行者來陪伴我們的孩子.
因為聖經告訴我們我們與神同工.
神是會幫助我們小朋友成長.
但是人有自己的責任.
但是人也有自己的職任.
我們做教會也是應該擴闊我們小朋友的.
年輕人的屬靈的同行者的比例.
我們做教會也是應該擴闊我們小朋友的屬靈的同行者的比例.
除了斯魯贊之外還有其他的導師.
希望這些導師能夠幫助我們小朋友.
還有我們作為父母的一同幫助小朋友的成長.
希望其他的導師和我們一起作為父母的一起成長.
最後我們看到一件很特別的事情.
這群人想將他的炭治朋友帶到耶穌面前得到醫治.
帶來醫治的不只是他的肉身更是他的靈魂.
他們從未想過耶穌會說這句話並且會寬恕這個炭治的罪.
他們從未想過耶穌會講這些話並且會寬恕這個炭治的罪.
這就是我們在人生裡面對子女必須要對上帝有信心.
有時我們鋪排了很多路給我們的子女.
但上帝的心意未必是這樣.
有時我們安排很多路給我們的子女.
但上帝的心意並不是這樣.
我認識一個很好的朋友,在CLC的兒子在一間名校的中學5.0畢業.
他讀了AP,成績是最優異的.
然後進了UC Berkeley,第一年就落伍了.
憂鬱,因為他有抑鬱.
但是上帝就帶他去另一間學校讀書.
然後他蒙神呼召做了宣教士.
結了婚繼續在宣教工廠服侍.
結了婚之後繼續在宣教工廠工作.
結了婚之後繼續在宣教工廠服侍.

$^{441}$其實我們真的不知道上帝對我們子女的心意是怎樣.
其實我們真的不知道上帝對我們子女的心意是什麼.
但是你信不信得過上帝會看著你們的子女呢?.
但是你信不信得過上帝會看著你的子女呢?.
只要你盡了自己的本分.
只要我們盡了自己的本分.
我們就要把他們交託給上帝.
每次說到我兩個女兒.
我一定會分享這件神奇的事件.
每次說到我兩個女兒.
我都會分享這件神奇的事件.
就是我大女兒來美國讀書的事蹟.
有些節目都聽過了,再聽一次吧.
當我大女兒13歲的時候.
她就跟我說她不喜歡香港的教育.
她想過來美國讀書.
我的大女兒13歲的時候.
她就說她不喜歡香港的教育.
她要來美國讀書.
她就說她不喜歡香港的教育.
(主持人說話).
\newpage



\section{士師記 8:23}
\label{sec:032337}
\textbf{聖經人聖經事:唯有耶和華管理你們}~~~~黃克勤牧師博士
\newline
\newline
連結: \href{http://www.church.com.hk/acms/content.asp?site=cdc\&op=show\&type=product\&code=032337\&layout=sermon}{\texttt{032337}} ~~~~ 語音日期: 2026-06-19
\newline
\newline
\hyperref[sec:032336]{\small{< < < PREV SERMON < < <}}
~
\hyperref[sec:index_chronic]{\small{[返順時目]}}
~
\hyperref[sec:032338]{\small{> > > NEXT SERMON > > >}}
\newline
\newline
士師記 8:23
\newline
\begin{longtable}{cl}
\hline
\hline
章節 & 經文 (和合本修訂版)\\
\hline
8:23 & \begin{tabularx}{0.7\textwidth}{X} 基甸對他們說：「我不治理你們，我的兒子也不治理你們，耶和華會治理你們。」 \end{tabularx} \\ \\
[1ex]
\hline
\hline
\end{longtable}
$^{1}$各位聽眾大家好.
和大家一起讀聖經.
《生命說書》的網絡.
《聖經人》《聖經事》.
看看上帝在人與事背後有什麼動作.
「惟有耶和華管理你們」.
「耶和華管理你們」.
這是基甸所說的.
今天聖經人物繼續在基甸.
故事就是他和以色列人說的話.
是詩記八章二十二節.
以色列人對基甸說.
「你既救我們脫離米殿人的手」.
「願你和你的兒孫管理我們」.
基甸說「我不管理你們」.
「我的兒子也不管理你們」.
「惟有耶和華管理你們」.
基甸又對他們說.
「我有一件事求你們」.
「請你們各人將所奪得的耳環給我」.
「原來仇敵是以實瑪利人」.
「都是戴金耳環的」.
他們說「我們情願給你」.
「鋪開一件外衣」.
「各人將所奪得的耳環」.
「雕在其上」.
廿六節「基甸所要出來的金耳環」.
「是重一千七百射赫勒金紙」.
「此外還有米殿王所戴的月環」.
「以誰和所穿的紫色衣服」.
「並落駝項上的金鏈紙」.
「基甸以此製造了一個耳筏德」.
「設立在本城俄弗拉」.
「後來以色列人拜那爾弗德」.
「行了邪淫」.
「就作了基甸和他全家的網囊」.
今天跟大家讀這段經文.
我發覺先讀完一次.
然後再講它的重點.
可能會更加清晰.

$^{41}$所以我回去跟大家讀聖經.
「惟耶和華管理你們」.
「基甸贏了米殿人」.
「最後那兩個王在沉日體都殺死了」.
「沒有什麼再有米殿人」.
「米殿人是一個沙漠遊牧的民族」.
「他們其實包含了很多很多族人」.
「當然以大的種族來看」.
「他們都是其中一個分支」.
「不過裡面有更多小分支」.
「米殿人非常豪大」.
「上次他們來攻擊是十萬軍隊」.
「組合在一起追攻以色列」.
「很厲害的」.
「以色列人跟基甸說」.
「你既救我們脫離了米殿人」.
「差不多這次永遠都打不過他們」.
「既然如此」.
「就願你和你的兒孫管理我們」.
「基甸一口拒絕」.
「我不管理你們」.
「他們已經想過了」.
「如果他們叫我做領袖」.
「我不做」.
「管理」的意思是「統治」.
「統治」聽起來好像很高壓.
「我不會治理你們」.
「管理」的意思是「治理」.
「我不會做治理你們的人」.
「我家裡都不會」.
「我兒子都不會」.
「其實最好的」.
他說出了我想說的第一點.
「其實最好的就是」.
「交給耶和華來管理你們」.
「讓上帝自己工作」.
「我不管理」.
「我兒子都不管理」.
「唯有Yahweh」.
「Yahweh」是上帝的名字.

$^{81}$他們不敢讀這個字.
「唯有阿都拿」.
「唯有主管理你們」.
「耶和華自己管理」.
「交回這件事給主」.
「上帝親自管理」.
「大家尊敬祂就行了」.
其實基典講出了一個.
很重要的靈性上的訊息.
很喜歡找人來管理.
我們全部都找人來管理.
「你那麼厲害不如你管我們」.
「要和我們打勝敵人」.
「不如你管我們」.
「你那麼厲害不如你管理我們」.
基典說人的管理不行.
沒用不夠好.
最重要是耶和華上帝.
上主親自做我們的管理者.
我們凡事求問祂.
跟著祂給我們的律例典章.
跟著祂走就好了.
為耶和華管理我們.
無論誰做領袖都好.
最重要祂都是.
跟著主耶和華的管理就行了.
這一點是今天我想講的重點.
我們很多時候都看著人.
牧師怎麼說.
執事會主席怎麼說.
執事會怎麼決定.
我們的組長,家長,長老怎麼決定.
看他決定.
一看人的決定.
很多時候人之間就會爭.
而且他的決定都未必特別好.
當然必須要所有決定都是.
基於聖經基於上帝的帶領與管理.
所以是為耶和華管理我們最好.
所以聖經你會見到.

$^{121}$聖經講不講民主呢.
聖經老老實實.
不是很民主的.
因為上帝一創造.
我六天創造要創造.
這些就創造了.
創造人是這樣的就創造了.
創造人就有一雙眼在前面.
後面人來搏頭擋.
我看不到.
為什麼不做一雙眼在後面呢.
左右的車過來看不到.
不如兩邊又有兩雙眼.
最好頭頂多一雙.
因為有高空雜物我看到.
又可能下排有一雙我看到地下.
那就好了.
全身都是眼看得清楚了.
要全身都是手才很方便.
全身都是腳跑得快.
什麼都要多.
其實上帝無論怎樣.
他有沒有創造的時候問我們呢.
他是沒有的.
一做就做了.
十誡的時候上山.
兩塊法板已經寫好給他拿下去了.
有沒有叫那些.
十二支派的長老上來山上我們圓桌呢.
為什麼要上來山上.
應該神都下來和我們一起坐在那裡聊.
這樣才叫民主.
我們是同等的.
對等的談判.
這樣就叫民主了.
對等.
上帝怎樣跟你對等呢.
其實聖經是不民主的.
民主從聖經出來.
因為大家都是上帝創造.

$^{161}$大家擁有相同的人性.
相同的靈.
所以其實我們是同等的.
一個人都應該有一個聲音.
有一個聲音應該被尊重.
是的.
但是整個聖經的管治是什麼呢.
是神主的.
是神話的.
不是人主嗎.
不是.
不是民主嗎.
不是的.
是神帶領的.
這個很有趣的.
民主雖然有聖經背後.
或者間接的支持.
但是聖經並不民主.
聖經是神主.
祂是創造者.
祂是審判者.
祂是修復者.
祂是決定到永遠的一位主宰.
你無端端要找一個主來在你頭上.
無論我喜歡不喜歡.
祂都是.
已經是我的主了.
我是被祂創造的.
世界都是祂創造的.
這是我們的信念與承認.
有時候是唯有耶和華管理你們.
無論你服還是不服.
這是基澱帶出來說的.
他經歷到這位極有能力的上帝.
是其他人沒有經歷到的.
所以他完全信靠耶和華神的管理.
當人們還在看.
看不到神的管理這麼好.
看不到神的管理真是實在.
他說人管吧.

$^{201}$基澱說我不管.
耶和華管.
大家要學到這件事.
不過.
他為了令到耶和華管我們.
他想落實一點.
不過這反而成為了他後來的網絡.
我有一件事求你們.
這個求綠色字是解釋來的.
求Chael.
Chael是詢問或者請求.
這個我求你們.
好像說得很重.
其實這裡.
我有一件事想問大家.
可不可以將你們.
搶到回來.
贏敵人搶到回來.
討略回來的金耳環給我.
我想問大家.
可不可以將金耳環給我.
不是請你們將金耳環給我.
不是這樣的.
他們是請大家給我.
將奪得的耳環給我們.
原來受敵都是以實馬利人.
都是戴金耳環的.
以實馬利人不是剛剛說米電人嗎.
什麼是以實馬利人.
其實米電人.
其他等等的人.
米電人是什麼後裔呢.
就是以實馬利後裔的分支.
所以說他們其實是以實馬利人.
以實馬利是誰呢.
Ishmael.
以實馬利就是.
阿巴拉罕的太太.
薩拉的一個婢女.
就是薩拉給阿巴拉罕做妾的下甲的兒子.

$^{241}$他理論上.
以實馬利大家還記得.
應該是以撒的同父異母哥哥.
不過他也是出自阿巴拉罕.
所以神就說.
以實馬利.
我都認.
我都認他是你所生的.
所以我都會給他成為一個大族.
好像你阿巴拉罕.
我給你的以撒一樣.
成為大族.
所以以實馬利人.
想不到由於這樣.
以實馬利人將來會和以色列作對.
正如以蘇是雅各的哥哥.
以蘇的後裔以東人.
後來都是和以色列作對.
其實今天我們統稱他們為現代的名字.
都是叫阿拉伯人.
由於以實馬利是米電在他裡面的其中一個分支.
以實馬利人在這裡看到.
原來他們都有一個習慣.
就是戴金耳環.
男的戴金耳環.
這個金耳環.
男的用來做飾物.
女的都用來做飾物.
有些是戴兩隻耳.
有些是戴一隻耳.
女的很多時候都加多一個.
有錢的漂亮的就加多一個鼻環.
就覺得時髦了.
加個鼻環漂亮了.
男的就不會鼻環.
耳環一隻或者兩隻.
今天很多人都會這樣.
將那些耳環給我行不行.
將那些環給我.
其實不是想要那個環.

$^{281}$是要那些金.
給我.
他們就說好啊.
我們是很甘心留意給你的.
既然我們領袖都想找你做了.
你要金耳環.
給你吧.
我真是情願的給你了.
他們就打開外衣.
外衣鋪得很大.
各人就將金耳環放在衣服裡.
有一大包.
有多少呢.
基地所要出來的金耳環.
是1700赦克納金紙.
我用藍色而加多一劃.
在電腦裡用的.
其實是中文聖經和合本是有的.
但是原文希伯來文是沒有的.
是重1700金紙.
沒有赦克納的.
基地所要出來的金耳環.
是重1700金紙.
有多少金呢.
1700.
大家就知道.
單位自然就是赦克納.
赦克納是多少呢.
一個赦克納Shackles.
一個是不會S的.
一個是Shackles.
這個S就不要了.
對不起.
赦克納是11.4克.
大概.
11.4克.
記住不要這個S.
一個赦克納一個Shackles.
就是11.4克.
這個是重量.

$^{321}$所以1700個Shackles.
就一個S.
1700個Shackles的金紙.
就大概是20公斤.
20公斤即是40磅.
44磅.
都很重.
一大塊金是用了.
還有什麼呢.
全部放進去.
這些有米建王所戴的月環.
耳垂.
穿的紫色衣服.
和下台頸上的金鏈.
長一點.
這兩句不要緊.
只是一個補充.
人們很樂意獻上他們的金耳環.
或者金環.
就有大概20公斤左右.
1700個Shackles.
機電就以此製造了一個爾忽德.
這個故事的轉折就在這裡.
第三點.
機電是做一個爾忽德.
做什麼呢.
其實爾忽德是大祭司的標準禮服.
其中一件.
有一件外袍.
有一件內衣.
一件外袍.
然後就穿一件爾忽德.
爾忽德就好像背心形一樣.
這裡有帶.
一件中等去到差不多長的.
一件藍色為主的背心.
然後再在上面就有一個金的牌.
一個空牌.
如果沒有了金.
有很多都想辦法用布也好.

$^{361}$都是雙12塊的寶石.
12種顏色的寶石.
象徵大祭司是胸懷著以色列12個支派.
12塊寶石.
大家看看空牌.
12塊寶石.
以色列12支牌.
底部就是一件爾忽德.
爾忽德是一件背心形.
藍色為主的衣服.
裡面有紫色線.
有金色線.
有紅色線穿插著.
主要是藍色為主.
所以有四種顏色.
四種線織成的.
藍色為主.
有金有紫有紅.
另外就有祭司各樣的冠冕.
我不說那些了.
一件祭司的爾忽德.
掛在城裡.
他的城就是俄弗拉.
掛在城裡.
表示去到俄弗拉.
看到有祭司的代表.
外面要有一件爾忽德和空牌才是祭司.
見到一件爾忽德.
空牌不能隨便做.
爾忽德做一件就掛在那裡.
這個爾忽德.
就是一件在祭司禮袍前面.
一件好像背心的衣服.
上面有扣連著.
穿過掛在前面.
後面就包著.
這件爾忽德就好像背心.
上面有個空牌.
空牌有以色列十二支派的代表.
代表用十二塊寶石.

$^{401}$不同顏色來代表.
表示掛在胸口.
是胸懷以色列十二支派.
是大祭司的工作.
是他揹著十二支派.
空牌後面的背景.
是一件爾忽德.
好像一件背心.
爾忽德主要是紫色,藍色,金色,紅色的線織起.
紫色為主.
即是專櫃顏色.
而在這件空牌後面.
爾忽德後面是一件.
去到一半身上下五吋.
或者膝蓋上下.
是一件藍色的爾忽德袍.
爾忽德就在爾忽德袍的上面.
爾忽德上面有個空牌.
三層的.
然後到外袍白色.
你會見到.
就是這樣.
都幾專櫃.
爾忽德掛在機電本城俄弗拉的門口.
或者掛在裡面有一間屋再住.
反而後來那些以色列人在機電在的時候.
他們就尊敬.
禱告.
但機電不在之後.
沒有什麼人有這樣的影響力與權力.
那些以色列人就開始用了他們當時拜本地的神.
就是巴爾巴力的方法.
巴爾很多個巴力很多個的.
所以叫珠巴力.
很多個不同的巴力什麼的.
拜的時候.
他們就習慣了後來就發覺到.
想與巴力有更多的深入的來往.
就會與當時的廟祭發生關係.
於是就行了邪淫.

$^{441}$誰知去拜這個爾忽德的時候.
逐漸都用了這些當時加男人拜巴力的方法.
於是就搞到.
什麼骯髒垃圾.
本來是去思想仰望上帝的.
就變成了拜巴力這樣的辦法.
於是很骯髒垃圾.
就成為了機電和他全家的網羅.
網羅的意思就不是那個網看著.
而是變成他們全家的污點.
什麼這樣.
那這件事是誰做的.
起初是機電做的.
機電會是這樣的嗎.
是呀,很好,他很厲害,很有信心.
不過後來做了這件事.
搞到我們失敗了.
罵了他.
其實是以色列人自己拜慣了巴力的方法.
套用在敬拜神那裡.
神是很不歡喜的.
你見到是說這些是邪淫.
很壞的事來的.
神是不歡喜的.
變成了好事就叫做好心做壞事.
好心做壞事.
本來應該是今天的主題.
機電好心做壞事.
我不想這樣說.
所以我就用了正面的.
其實機電是說我不管你們.
唯有耶和華管理你們.
不過他做了一件事.
反而被魔鬼利用了.
被以色列人利用了.
反而滿足他們拜巴力的方法.
真是好心做壞事.
好事多磨.
不要上帝的管理.
而要巴力的管理.

$^{481}$逐漸做了出來.
真是很可惜的.
其實很多時候.
為什麼人這麼差.
靈性這麼低落.
這樣也會的.
人就是這麼軟弱的.
覺得自己最厲害的人就是最差的.
很多時候都會這樣.
如果有些時候.
我很想為香港的.
基督徒教會的靈性復興禱告.
猛唱復興的歌.
猛開一份慶會有沒有用呢.
好過沒有的.
不過最大的復興.
我發覺在聖經裡面也好.
在人類教會歷史裡面也好.
那種復興的來臨.
我是知道的.
但是也很怕這樣祈禱的.
就是復興很多時候.
在什麼時候會出現呢.
就是在我們信徒被逼迫的時候.
甚至是很大逼迫的時候.
很艱難.
甚至會流血的時候.
監禁的時候.
那時候反而靈性好.
好像火一樣煉成了教會.
哪些是草木和街.
哪些是金銀寶石.
馬上知道.
很多人面對著這些災難逼迫.
站得很直.
反而成為很多人信仰的榜樣.
這個是真的信仰了.
反而教會復興.
譬如中國經過很艱難的文化大革命.
十年浩劫裡面.

$^{521}$是沒有了教會.
沒有了傳道人.
甚至聖經都沒有了.
當時的信徒怎樣記住聖經呢.
少少人.
不如我們分開.
你唸希伯來書.
我唸羅馬書.
你唸詩篇.
怎麼唸到150篇.
好了我們分開五個人唸.
一人唸50篇.
或者一人唸一章聖經.
羅馬書十章是誰.
我.
他唸出來了.
於是我們今天讀經到羅馬書十章.
是那個人唸出來的.
沒有了聖經.
大家記住.
讀在心裡面.
嘩多精煉啊.
那時候禱告多行刺啊.
多大的水都蓋不住.
那些火不熄.
因為苦難.
苦難是會令到我們復興.
好好的享受反而會令我們衰弱.
有時候主要是讓我們平平安安.
順順利利.
靈性有時會低的.
讓我們有些艱難.
靈性會高.
我發覺很多人在病的裡面.
嘩靈性多好.
當然我當然不想.
神啊我的靈性好.
讓我生病.
讓我意外.
我不想.

$^{561}$不過很多時候.
都是有一根刺.
反而那個人的信心更加好.
求主幫助我們.
不用刺不用火.
仍然會知道耶和華管理是最好的.
這樣就最好了.
我們動心禱告.
天父求你幫助我們.
讓我們自己會歡喜你的管理.
會信服你的管理.
於是我們真的好像聖經所講.
我們是得到平安的.
得到喜樂的.
而無需要去火的煉證.
主我們很多時候都為到平安禱告.
為到平順順利禱告.
為到富裕禱告.
主我們又怕痛又想好.
想好的時候越不要痛.
我們就越退縮.
不敢面對艱難的事情.
信心會低落.
求主幫助我們在這個當中.
滿有你自己的恩典與憐憫.
幫助我們更多順求你主動的跟隨你.
主在需要一些鞭策的時候.
我們都願意接受.
求你幫助我們.
給我們更加緊的跟隨你.
你管理我們.
這是我們的禱告.
奉靠耶穌基督命求.
阿們.
\newpage



\section{以弗所書 6:10-24}
\label{sec:032338}
\textbf{門徒當自強}~~~~呂君望牧師
\newline
\newline
連結: \href{http://www.church.com.hk/acms/content.asp?site=cdc\&op=show\&type=product\&code=032338\&layout=sermon}{\texttt{032338}} ~~~~ 語音日期: 2026-06-21
\newline
\newline
\hyperref[sec:032337]{\small{< < < PREV SERMON < < <}}
~
\hyperref[sec:index_chronic]{\small{[返順時目]}}
~
\hyperref[sec:032339]{\small{> > > NEXT SERMON > > >}}
\newline
\newline
以弗所書 6:10-24
\newline
\begin{longtable}{cl}
\hline
\hline
章節 & 經文 (和合本修訂版)\\
\hline
6:10 & \begin{tabularx}{0.7\textwidth}{X} 最後，你們要靠著主，依賴他的大能大力作剛強的人。 \end{tabularx} \\ \\ \relax
6:11 & \begin{tabularx}{0.7\textwidth}{X} 要穿戴神所賜的全副軍裝，好抵擋魔鬼的詭計。 \end{tabularx} \\ \\ \relax
6:12 & \begin{tabularx}{0.7\textwidth}{X} 因為我們的爭戰並不是對抗有血有肉的人，而是對抗那些執政的、掌權的、管轄這幽暗世界的，以及天空靈界的惡魔。 \end{tabularx} \\ \\ \relax
6:13 & \begin{tabularx}{0.7\textwidth}{X} 所以，要拿起神所賜的全副軍裝，好在邪惡的日子能抵擋仇敵，並且完成了一切後還能站立得住。 \end{tabularx} \\ \\ \relax
6:14 & \begin{tabularx}{0.7\textwidth}{X} 所以，要站穩了，用真理當作帶子束腰，用公義當作護心鏡遮胸， \end{tabularx} \\ \\ \relax
6:15 & \begin{tabularx}{0.7\textwidth}{X} 又用和平的福音當作預備走路的鞋穿在腳上。 \end{tabularx} \\ \\ \relax
6:16 & \begin{tabularx}{0.7\textwidth}{X} 此外，要拿信德當作盾牌，用來撲滅那惡者一切燒著的箭。 \end{tabularx} \\ \\ \relax
6:17 & \begin{tabularx}{0.7\textwidth}{X} 要戴上救恩的頭盔，拿著聖靈的寶劍—就是神的道。 \end{tabularx} \\ \\ \relax
6:18 & \begin{tabularx}{0.7\textwidth}{X} 要靠著聖靈，隨時多方禱告祈求，並要為此警醒不倦，為眾聖徒祈求。 \end{tabularx} \\ \\ \relax
6:19 & \begin{tabularx}{0.7\textwidth}{X} 也要為我祈求，讓我有口才，能放膽開口講明福音的奧祕， \end{tabularx} \\ \\ \relax
6:20 & \begin{tabularx}{0.7\textwidth}{X} 我為這福音的奧祕作了帶鐵鏈的使者，讓我能照著當盡的本分放膽宣講。 \end{tabularx} \\ \\ \relax
6:21 & \begin{tabularx}{0.7\textwidth}{X} 今有親愛、忠心服事主的弟兄推基古，為了你們也明白我的事情和我的景況，他會讓你們知道一切的事。 \end{tabularx} \\ \\ \relax
6:22 & \begin{tabularx}{0.7\textwidth}{X} 我特意打發他到你們那裡去，好讓你們知道我們的情況，又讓他安慰你們的心。 \end{tabularx} \\ \\ \relax
6:23 & \begin{tabularx}{0.7\textwidth}{X} 願平安、慈愛、信心從父神和主耶穌基督歸給弟兄們。 \end{tabularx} \\ \\ \relax
6:24 & \begin{tabularx}{0.7\textwidth}{X} 願所有恆心愛我們主耶穌基督的人都蒙恩惠。 \end{tabularx} \\ \\
[1ex]
\hline
\hline
\end{longtable}
$^{1}$領姊妹平安.
如果說我們個個都身經百戰.
你會想起什麼呢?.
又說人生充滿了戰爭.
所謂人生如戰場.
無論是生活還是職場.
回到家裡好像打仗一樣.
家人的關係很緊張.
為了子女 兩夫婦可能吵架.
為了子女不讀書 不做功課.
跟子女吵架.
我們未必會動手動腳.
但那種罵戰 那種冷戰.
是經常發生在我們和家人的當中.
多年前 我和媽媽冷戰了三天.
當對方透明 最主要是我.
但最後我降服 我認輸.
媽媽永遠在我心中是我放不下的母親.
無論如何冷戰 最後我們都會妥協.
但其中一個戰 不知道你有沒有想過.
我們經常有內戰.
什麼內戰 是你和自己的爭戰.
你昨晚有沒有打過 今早有沒有打過.
我不知道今早有沒有世界盃的比賽.
可能有一把聲音告訴你.
看過去吧 這麼好看.
當然看現場 為什麼看錄影帶.
崇拜 上網可以看.
但你那把聲音告訴你.
你選擇我 還是選世界盃.
老我和新的我在爭戰.
所以弟兄姊妹 其實我們經常打內戰.
究竟你是代表上帝的那邊性.
還是代表魔鬼的老我性.
信耶穌 你以為天下太平.
什麼都好了 天下太平沒事.
不是的 弟兄姊妹.
信耶穌戰爭不但沒有減.
反而更加激烈.
因為撒旦最喜歡打仗.

$^{41}$他挑起他和上帝 和上帝子民的爭戰.
他一個最勁的敵人.
地上的敵人可能只不過是我們的家人.
是我身邊的弟兄姊妹.
是我們認識的人 是血氣之爭.
但魔鬼是屬靈的.
他的靈魂是由地上帶到天上爭戰.
今年到現在.
我們一直由本堂的教務同工思想.
和大家一起去研讀 已忽所書.
今天是最後了.
在過去的日子裡面.
我們會看到第一章至五章二十節裡面.
保羅告訴我們.
我們在基督裡面.
我們領受了上帝很多很多恩典.
我們是蒙恩的人.
我們蒙愛的兒女.
我們在基督裡面已經成為新造的人.
所以我們應該活在光明的日子.
我們要捨棄那些污穢.
穿上新人.
這是在經文裡面不斷地告訴我們.
而彭恩教師在上個主日.
來到了五章二十一節至六章九節.
告訴我們.
我們人是活在不同的關係裡面.
夫婦,父母子女,主人,僕人.
在不同的關係當中.
我們都要活出這個福音的大能.
而這一切都是建基在我們和主的關係上.
如果我們和主沒有一個緊密的關係.
任憑我們如何努力建立我們夫婦,父子,父女,主人,僕人的關係.
我們都會出現很多問題.
所以有關係就沒有關係.
只要我們和耶穌有一個緊密的關係.
人際的關係.
就算有些磨擦.
都是很少的問題.
是很容易的解決.

$^{81}$如果大家有留意.
在二弗所書前面不斷講到在基督裡面.
或者相繼字.
在基督裡面等等.
或者在基督裡面互著這些字眼.
是呀.
是保羅在二弗所書告訴我們.
我們在基督裡面互相為肢體.
建立基督的身體.
耶穌基督是頭.
我們作為主的門徒.
然後信了耶穌.
我們就成為一個新的創造.
我們是新造的人.
我們是新的群體.
這些新的群體.
這些美好的生活.
照理來說.
如果我們每個人都是新的群體.
回來聖經堂應該是一個天堂.
但事實上.
我想會比外面好很多.
但是會不會去到天堂呢.
又未至於.
每次回來我都很大掙扎.
我講可能你.
又見到他了.
所以是有爭戰的.
我們這個新的生命.
不是用來閱兵的.
閱兵是用來看的.
告訴別人我們有多麼的裝備.
多麼的先進.
基督徒的生命不是用來給人看的.
基督徒的生命.
這個新的生命要在我們不斷的爭戰當中.
在人際關係的磨練當中.
去實踐出來.
今天你和我.
作為基督徒的精兵.

$^{121}$我們去面對這個終極的爭戰.
這個敵人不是小事.
是屬靈的惡魔.
保羅提醒我們.
吩咐我們一件很重要的事.
一起讀這節聖經.
預備.
軟弱但有大力.
作剛強的人.
作剛強的人是保羅的吩咐.
是命令來的.
不是你的選擇.
可有可無.
我要選擇軟弱的不行.
你既然是門徒的兒女.
作為主要門徒.
你就要剛強.
最後.
或者話本叫做《沒了的話》.
可以說是保羅在葉忽所書的一個總結.
在總結裡面保羅語重心長的.
去命令弟兄姊妹.
要做一個剛強的人.
這一句可以說是整個葉忽所書的總結.
是一個很重要的標題.
基於這樣.
今天的講題叫做.
門徒當自強.
保羅勸勉我們要做剛強的門徒.
為什麼.
是因為.
我們是一個真正強大的敵人.
所以你不剛強.
你輸定了.
我們就讀.
十一,十二節.
預備起.
對了.
剛才說過.
我們真正的敵人不是地上的人.

$^{161}$不是你的家人.
不是你教會的弟兄姊妹.
不是你公司的同事.
學校的同學.
這些都是屬於血氣的.
我們大家都是屬於血氣.
可能討論的東西.
爭議的東西.
吵架的東西.
都是義氣的事.
放不下的意思.
執著的意思.
但是我們的敵人.
原來屬於靈界邪惡的勢力.
有權有勢.
去掌管一切.
他甚至可以告訴你.
你要榮華富貴.
我給你.
你只要拜我.
你要的東西.
不用那麼辛苦.
不用做牛做馬.
即時給你.
你只要拜我.
放棄耶穌就可以了.
他真的有這個權有勢.
不過他的領域是一個黑暗.
邪惡的領域.
而這個領袖.
就是撒旦魔鬼.
如果上帝是真實.
撒旦魔鬼也是真實的存在.
不是幻覺.
不是你和我想的東西.
是真實的存在.
所以我們作為主門徒.
我們和魔鬼爭戰.
是必經的階段.
或者是有機會遇到的事情.

$^{201}$過去端午節.
我們同意會遊退修會.
我們邀請了一位傳道人.
不是我們教會傳道人.
是我們教會聚會.
現在在我們當中的傳道人.
帶領整個聚會.
讓我們去退到主的面前.
就在剛剛過去星期一開始.
她先生就發高燒.
很不舒服.
因為只有他們兩夫婦.
作為太太的.
面對一個患重病.
有年紀那麼上下的一位丈夫.
總會擔心.
總會有掙扎.
是不是應該推遲這個負責呢?.
撒旦會告訴你.
當然推遲了.
照顧丈夫天經地義.
你怎麼會那麼不負責任呢?.
但是這把聲音.
這次是堂而會很難得的機會.
是上帝為你預備.
讓堂而會邀請你.
來幫助他們.
一起去到上帝的面前.
這就是神.
這是一個宿命的爭戰.
當我知道這位講完壞話的時候.
我得到堂而會的同意之後.
在堂而會群組當中.
我邀請每個童工.
因為我們有份的.
為港元禱告.
為她丈夫禱告.
雖然出發的時候.
丈夫仍然是發燒.
仍然很不舒服.

$^{241}$但是我們憑了信心.
讓港元來到我們當中.
他沒有任何掛慮.
他都告訴你.
電話都不聽了.
讓我們專心的去到神的面前.
這是宿命的爭戰.
是真實的.
稍後我們會為L國的.
短孫的團友來禱告.
全部現在病了.
培育教師病了.
不過如果延遲一星期.
就會有些淪盡.
神都在恩典念文.
我不肯定這是否宿命的爭戰.
不過不多不少都會有.
但是我們知道.
這是我們必須面對的事實.
有人說 撒旦的座右銘是甚麼呢.
無論何處.
他都要抵擋神.
簡單來說.
撒旦和上帝是世仇.
是對敵的.
所以撒旦也會抵擋.
屬於上帝的事情.
攔阻上帝的事.
來順利的完成.
也會攔阻屬於上帝而來.
要做的事.
你想立旨.
你立旨要去跟從上帝.
你立旨要親近上帝.
我就不給你機會.
魔鬼要做這些事.
保羅說.
我們要剛強.
我們要靠主.
目的是抵擋魔鬼的詭計.

$^{281}$我穿著全副軍裝.
是要抵擋魔鬼的詭計.
說明詭計是負面的.
詭計的目的.
就是要動搖我們屬神的兒女.
他們的門徒對上帝的信心.
扭曲了我們在神面前應有的價值.
他會令我們不相信.
神已經原諒了我們.
我做了這麼多事.
神真的會原諒我嗎.
不會不會不會.
我都不會原諒自己.
上帝更加不會原諒我.
所以活得不開心.
沒有喜樂.
很痛苦.
魔鬼會叫他相信.
你是沒得救了.
你正在進入地獄.
也會令人相信.
在基督裡面.
他沒辦法成為新造的人.
你賭了一輩子.
那一輩子的白粉.
你變了舊.
你只是垃圾.
他看不到.
在基督裡面.
神已經饒恕了我的罪.
他已經改變了我的人.
就算人都覺得我們生命改變了.
但我自己不相信.
我在基督裡面已經成為新造的人.
所以我們不能向世人見證.
因為我還是一個很老我的人.
這是魔鬼要蒙蔽我們的心.
使我們人失去了平安和盼望.
魔鬼很自然做這些事.
因為聖經說魔鬼是說放人之父.

$^{321}$最近太古城有一宗慘案.
母女同一天早晚分別跳樓.
沒有人知道女孩子在想什麼.
是很痛苦的事.
很遺憾的事.
不知道爸爸能不能承受這些壓力.
不過有很大機會.
是因為這個十二歲的女兒.
回到家裡.
看到媽媽自盡的地方.
勾起了她很深層的悲哀.
不捨和那種自責.
如果不是我今早頂著媽媽.
媽媽就不會跳下去.
她將整個責任歸咎於自己.
其實爸爸和媽媽已經分居了.
不一定是因為女兒喊死.
是累積了的問題.
但在那一刻.
女兒會覺得是我讓媽媽跳下去.
這是我們估計的.
我們沒辦法知道.
魔鬼會用這些方法.
叫我們一生都在自責的裡面.
而且我們失去在主裡面.
應有的平安喜樂和寬恕.
感受不到那種新的生命.
繼續自怨自艾的捱過一生.
所以保羅要再次重新提醒我們.
一起讀十三節預備起.
所以每當媽媽致上帝.
所賜予的羽絨和棍裝.
好在她是我們人之能力的教授.
並且是人生的一切好運.
和三萬萬的尊嚴.
拿起和十一節穿上有點不同.
意思是我們要拿著這套東西.
或者穿上這套棍裝.
目的是什麼.
抵擋魔鬼撒旦.

$^{361}$以至我們仍然能夠站立得住.
我們要自強.
是因為我們需要面對的是.
強勁的仇敵.
既然這樣的時候.
我們就要穿上這個全副的棍裝.
如果大家留意.
在以弗所書裡面不斷有.
「活在基督裡面」剛才提過.
或者基督裡面.
可以說是整個以弗所書.
不斷強調的觀念.
穿上就是要告訴別人.
既然我們活在基督裡面.
我們就要穿上在基督裡面應有的表現.
包括真理,公義,和平,信德,救恩等等.
這是我們屬於上帝在基督裡面的.
新的創造應有的表現.
保羅要告訴我們.
在現實世界裡面是殘酷的.
在現實世界裡面.
是會攻擊我們的弱點的.
要過一個得勢的生活.
只有穿上這個棍裝.
這個棍裝不是我們電影裡面.
看到的萬俠的裝備.
是很獨特的.
有六個裝備.
讓我們簡單稍微看看.
第一個裝備請大家讀白色的經文.
所以.
第一個裝備就是那條腰帶.
要站穩了是保羅的命令.
又再一次命令.
你要站穩.
第一件事.
最重要找到一條合適的腰帶.
腰帶對於瘦了的人來說很倫循.
一定要有的不然會脫掉.
但是腰帶如果大家留意.

$^{401}$第一就是和當時的人的衣著.
普通人的衣著是長袍長衣服.
如果你走路.
打仗什麼都好.
被人追你走路.
如果穿著鬆鬆的袍.
就走得不好.
做士兵更加倫循.
所以腰帶是等你.
束起的衣服.
走也好打仗也好.
行裝輕便.
但是腰帶更重要的是.
將我們人體的力.
可以集中在某一點裡面.
這個你從舉重裡面就看到.
舉重的人.
一舉重之後放下啞鈴就鬆了腰帶.
原來腰帶是讓大力士更加有力.
保羅說用真理做腰帶.
真理的意思是什麼呢.
以真理來約束我們的心思.
而不是讓世界的潮流.
人的說話意見.
世界價值觀去影響我們.
我們要讓神的真理引導我們向前行.
將焦點放在我們集中的一點裡面.
而不是讓腰的肥膏和脂肪.
令我們沒有力量.
這就是真理作腰帶的意思.
因為有力量的時候.
你就能夠站立得穩.
原來綁著腰帶.
你的腰也可以直一點.
保護你的腰.
第二方面.
一起讀預備起.
心很重要.
就連真言也說你要保守你的心.
勝過保守一切.

$^{441}$冰釘當然要護心.
因為如果你傷了我的手手腳腳.
多少次都不會致命.
但一個心一傷就會死人.
所以心很重要.
而聖經用公義當作護心鏡.
在御塞夏書也有類似的意思.
在御塞夏書59章17節.
已經記載神他穿上公義的鎧甲.
神也用公義作為一個保護身體的裝甲.
不過不是用護心鏡.
其實也是保護身體.
即是避彈衣.
神用公義作避彈衣.
去保護自己.
為何上帝用公義.
或者先知引用上帝的時候.
用公義或保留在這裡用公義呢.
原來公義可以解道神自己.
公義也可以講道神的標準.
神的價值觀.
我們以神的標準和價值觀.
去保守我們的心.
當你有上帝的標準.
當你有上帝在你裡面.
你知道你的心不會遠離上帝.
未信耶穌的人會說.
我問心無愧.
我對得起天地良心.
意思是甚麼呢.
就算你怎樣說我.
你怎樣誣告我.
你怎樣詆毀我.
我對得起自己.
有天造證.
同樣當我有上帝的公義在我內心裡.
當我以上帝的公義來行事為人的時候.
我就知道.
我不需要介意任何的攻擊.
包括魔鬼的誣蔑.

$^{481}$雅各也說.
所以要馴服神.
要抵擋魔鬼.
魔鬼就必逃避你們.
魔鬼見到你們走.
不是因為我們這麼兇.
我們絕對沒有魔鬼這麼兇.
我們絕對沒有魔鬼這麼有能力.
但魔鬼怎會逃跑.
雅各怎樣說.
是因為甚麼.
因為你和我馴服神.
當你馴服神去過生活的時候.
你不需要害怕任何的說話.
不需要害怕任何的攻擊.
因為連魔鬼都離開了.
因為你馴服神 它害怕你.
如果你是一個只是馴服人.
而不馴服神的人.
魔鬼說.
行 你是我好友.
我們是朋友.
弟兄姊妹.
你願意當上帝的朋友.
還是願意當撒旦的朋友.
第三個裝備是甚麼.
一起讀白色字 預備起.
以前做夢時試過不要穿鞋.
很尷尬.
鞋很重要.
第一 保護腳不用弄髒.
第二 鞋對於人來說很重要.
如果鞋不適合腳形.
要不就穿得很辛苦.
要不就走路很不自然.
搖搖腳 腳zoung 那些.
搖搖腳 你走不了.
你走不快.
一個適合鞋很重要.
一個士兵很重要.

$^{521}$士兵穿對鞋是幫助他打仗.
幫助他追士兵或逃跑.
保羅說以和平的福音作為鞋.
是有意思的.
因為在戰爭環境.
在現實世界充滿了不安.
如果我們的鞋不是平安和平的鞋.
我們基督徒活在這個世界.
就無法繼續向前走.
我們無法到處去見證耶穌的福音.
唯有這個和平的福音.
一個平安的福音.
使我們能夠到處向人分享.
耶穌多少福音的大能.
正如聖經所說.
報福音傳喜上人.
他的腳zoung 是何等皆美.
頂智慕我們需要的是一個平安的鞋.
使我們去到任何地方.
我們自己走得很開心.
我們都能夠留下平安和平.
給我們所接觸的人.
究竟人因為我的存在.
令他們從然盼望重拾喜樂.
還是因為我的出現.
帶來紛爭帶來的不安.
這是值得你和我去反省.
羅馬兵的鞋是有釘.
使我們走在不同的環境裡面.
今天我們這個和平的鞋都應該有釘.
這個釘不是釘人.
不是踩人.
而是能夠我們在腳踏實地.
以致不會被任何東西衝擊我們.
頂智慕你有沒有平安作為你的鞋.
以致你去到任何地方裡面.
都能夠播下平安.
第四的裝備一起讀八式.
預備起.
使我們了解所有當中的盾牌.

$^{561}$如何防禦馬國旗的實際使用和指示.
其實盾牌是高一點的.
不會像以前我們舊式的警察的盾牌.
我記得那時候.
有一段時間警方的盾牌.
原來是透明的.
顧得上又顧不了下.
但是中東的時候.
古代的.
即是羅馬時代的盾牌是不高的.
可以有盾牌陣.
但是敵人都很厲害.
這裡告訴你為什麼說.
我要撲滅那惡者.
一切燒著的箭.
箭不是箭都必有燒著.
是的.
他們打仗很厲害.
不單要射你的盾牌.
射死你.
他還會加上火.
就算我射不中也好.
那些火都會燒著你的盾牌.
盾牌不是那些現在先進的物料做的.
會著火的.
所以說.
保羅說.
我的盾牌是以信德做的.
什麼叫信德.
就是信心.
可以說是我們信仰的核心.
是我們依靠的信念.
當我們有信心.
我們有依靠的時候.
我們不需要怕任何的火箭.
就算有火的箭我也不怕.
甚至好像和學本所說.
我要滅盡他們.
其實和學本比和修本更加好.
和修本說撲滅.

$^{601}$撲滅好像很笨.
感覺上.
但是他的意思是滅盡.
來一支我滅盡一支.
當你有一個信靠.
有一個信念的時候.
真的可以這樣.
曾經聽過一個提醒.
小心魔鬼三支箭.
哪三支.
明利.
明.
識.
利.
其實不止三支.
魔鬼有很多箭.
魔鬼有一支箭是你和我都不留意的.
就是暗箭.
不是寒箭.
是暗箭.
暗暗射了你和我都不知道.
他將我們的信仰世俗化.
信耶穌不要緊.
講福音都不要緊.
不要得罪來賓.
是不是.
不要講罪.
很難聽.
很不舒服.
講多一點.
講多一點他們喜歡聽.
講給他們聽.
信耶穌有什麼好處.
他們才會信.
你講不講那些講客.
剛剛最近我去了葵福的一個音樂聚會.
是報道會.
當然他名稱不能叫報道會.
因為在學校裡面舉行.
但很欣賞.

$^{641}$無論是Utility唱的詩歌.
當中一直貫穿那些詩歌的訊息.
和他們的話劇.
和講完陳奕迷牧師的訊息.
都很清楚將福音完完整整的去演繹出來.
這個就是真正的福音.
不隱藏的.
不保留的.
這個是人真正的需要.
但很多時候無論我們自己.
或者從商業角度.
你要請多些名人來.
這個報道會才有人來的.
將信仰世俗化.
將信仰商業化.
信仰或者聚會 教會活動.
成為滿足人.
無論弟兄姊妹滿足不同的人.
的一個需要的宗教的時候.
這個就是魔鬼時時射給我們的暗箭.
可能寫了很多枝.
有些被我們發現了.
有些可能我們都以為是這樣的.
所以弟兄姊妹.
在一點裡面.
順道告訴你 你要持守這個信仰核心.
什麼是你所信.
你敢信什麼.
你信耶穌就蒙福.
你信耶穌單單可以上天堂.
所以你信.
你信耶穌知道以後我沒有病沒有痛.
什麼問題都可以解決.
這個是不是你的信念呢.
為什麼你相信.
你本來是一個罪人.
但因為耶穌基督的救贖.
你和我的罪得到祂的赦免.
並且你的生命被改變.
從此你那個豐盛的生命不再一樣.

$^{681}$然後你為主而活.
作主的門徒.
等候主榮耀的再來.
將來在天下裡面.
能夠與耶穌一起得著榮耀.
這個是不是我們的信仰呢.
當你有這個真正的信仰的時候.
當你有信靠的時候.
你就可以防備魔鬼任何的火箭.
第五個裝備 一起讀白色的字 預備起.
很短的幾個字.
除了心之外.
頭也是很重要的身體部分.
如果踩單車.
雖然沒有要求一定要戴頭盔.
但對於自己的命很重要的人來說.
戴一戴都是好的 是不是.
這個是真的.
這個是保命的裝備.
關乎生命的存留.
我經常聽到一句很俗的說話.
其實是來自電影的.
他說 你信不信我用槍打爆你的頭.
很俗吧.
他想置你於死命.
你不聽話.
我就會置你於死命.
如果你不想置人於死命.
你再不讀書.
你信不信我打斷你的腳.
打斷你的手 最多是這樣.
當人說打爆你的頭的時候.
即是他完全不給你機會.
所以保護頭部.
正正是我們基督徒很重要.
頭部是我們中央司令的地方.
魔鬼就算不打死我們.
他也會打到我們昏昏沉沉.
叫我們朦朦朧朧去過基督徒的生活.
頭可以說是撒旦最賣力攻擊的部位.

$^{721}$他會叫我們產生.
打到我們昏昏沉沉的時候.
產生對舊因的懷疑.
什麼叫對舊因的懷疑.
貼壽羅家前書五章八字說得清楚一點.
就是把得救的盼望當作頭盔戴上.
什麼意思.
就是我們產生對得救的盼望 懷疑.
哎呀 我偶爾回教會我得不得救.
我不是做到那麼完美.
我還有些軟弱我得不得救.
當然我不是鼓勵你不回教會.
我不是鼓勵你繼續放縱.
正如保羅所說.
我仍然有肉體軟弱.
我心靈軟弱 我仍然有肉體軟弱.
約翰很清楚告訴我們.
人有神的兒子又有生命.
沒有神的兒子還有沒有生命.
正正是因為你有神的兒子.
有聖靈在裡面你才有那種掙扎.
如果你不相信耶穌.
你不覺得你有問題.
正正有聖靈在裡面.
你才覺得自己還有很多不足.
你有一種內疚感.
但如果內疚感.
魔鬼就會轉化成為對你的一種動搖.
你不行了 你們有多差.
人家有事奉你不事奉.
人家勤力回教會你不回教會.
人家努力讀聖經你不讀聖經.
人家有祈禱你不祈禱.
我講這些不是告訴你.
我以後不讀聖經不祈禱也行.
不是這個意思.
不過魔鬼會利用這些.
令我不單止沒有動力去努力.
我甚至覺得我是徹底的失敗.
我完全躺平.

$^{761}$死就死罷了.
但是我們要讓救恩成為我們的頭盔.
這個救恩不是因為我做這些做那些.
這個救恩是因為耶穌的寶血的大能.
祂的大能大力已經拯救了我.
保護了我的生命.
所以用救恩做頭盔是名正言順.
很重要的.
今天能夠給你生命.
只有耶穌基督的救贖.
不是你做什麼做那樣東西.
最後一樣東西一起讀.
預備起.
(大聖靈的寶劍就似上帝的座位).
是的.
聖靈的寶劍.
這個武器大家都喜歡.
但是這把劍不是長劍.
是短劍.
好像這個裝備有些料到.
那些全部都是防備的.
這個起碼都有把小刀.
或者小劍.
短劍.
無論怎樣也好.
保羅這裡不是強調.
這個寶劍是要我們攻擊.
我們沒有能力憑這把寶劍.
去對付撒旦.
也就是去消滅撒旦.
只能夠用這個寶劍.
去保護自己.
去擊退撒旦而已.
在保羅在爾忽所書一直說的時候.
這些裝備目的是要我們剛強.
站立得住.
就是叫我們不要跌倒.
包括這個最後.
神聖靈的寶劍.
都是這樣的功能.

$^{801}$為什麼?.
因為在其他書信裡面.
提到基文告訴你.
基督已經聖了世界.
聖上告天.
遠超過一切執政的掌權的.
耶穌已經得勝了.
祂已經掌權了.
所以你不需要來到很偉大.
我要消滅撒旦.
我們沒有本事消滅撒旦.
耶穌會做這件事.
你和我今天作為主觀門徒.
只需要努力的.
剛強的.
去完成這個屬靈的征戰.
這是一個防禦戰.
不要跟別人搶功.
努力死守就行了.
其他事情讓我來搞.
神的道.
就是這把劍.
而這把劍.
是需要聖靈的感動.
來運作.
因為聖經說.
聖靈是真理靈.
是叫我們進入真理的裡面.
今天如果有聖靈的寶劍.
即是有神的道.
如果沒有聖靈的感動.
我們不懂得運用.
我們也不懂得運用神的說話.
去對抗魔鬼.
所以聖靈在當中很重要.
所以被稱為聖靈的寶劍.
就是這個意思.
我們時時都有寶劍練習.
不要驕傲.
就算你每次寶劍練習.

$^{841}$你拿第一也好.
不代表你在日常生活裡面.
你可以勝過撒旦魔鬼.
唯有你將神的話語.
透過聖靈.
好好地馴服.
魔鬼才會遠離我們離開.
我們需要好像師父一樣.
時時保持練習.
保持操練.
同樣我們好像一般市民一樣.
時時要採取防火演習.
這是今天你和我的態度.
但如我們好像一位.
偉大的英國宣教士李文斯敦所說.
無論何處.
我們都要勇敢.
勇往直前.
不需要害怕.
因為這個全副軍裝.
足夠叫我們去面對撒旦.
要自強.
不可恥的元素.
就是倚靠聖靈的大能.
我們知道今天最先進.
第五代的隱形戰機.
就算你買了回來.
你都要懂得駕駛.
有人教你駕駛.
要有那個創造的國家.
提供密碼也好.
技巧也好.
正如美國有最先進的實質系統.
防禦系統.
有些西方國家擁有.
但要得到美國的專業訓練.
和技術轉移.
否則就算有實質系統.
只會霸佔地方.
同樣幾號的全副軍裝.

$^{881}$都需要依賴聖靈才能發揮功能.
我們來看看什麼叫依賴聖靈.
我們來看看經文保羅給我們什麼講解.
一起讀給他聽.
(依賴聖靈).
最強的元素就是依靠聖靈.
就像保羅這裡所說.
在靈裡面禱告.
我們需要用禱告.
而在聖靈裡面的禱告去遮蓋我們.
今天有很多爭戰.
沒有禱告.
我們沒辦法解決.
好像前幾天的退休會.
如果不是頂智會不斷的禱告.
我相信港元都很為難.
很掙扎.
整個退休會都不會這麼順利.
這麼豐富的得著.
所以禱告不單只是一個人的禱告.
是彼此的代禱.
而且是警醒不倦的禱告.
自強是需要不斷的禱告.
在聖靈裡面的禱告.
其實在《耳筆上書》3到16節.
保羅都說.
藉著他的靈.
就是聖靈叫你們心裡的力量.
剛強起來.
我們都是軟弱.
但是因為神的靈.
我們就可以剛強起來.
好像培恩教師上星期都說.
吾將十八節提到.
在這麼多人際關係裡面.
多麼惡劣多麼困難都好.
我們要先被聖靈充滿.
當你被聖靈充滿的時候.
你願意順服聖靈的時候.
一些關係你會主動的和好.

$^{921}$你會靠著聖靈可以和好.
所以頂智會不是你努力哄老婆.
哄丈夫.
哄子女.
哄爸爸媽媽.
或者妥協地算了.
不是這樣.
你將這些關係.
這些問題.
在禱告當中切切的教堂.
當你有聖靈充滿的時候.
你就能夠成為.
你就體會到神的工作.
我們教會是一個屬靈的群體.
我們的自強不是孤軍的作戰.
我們透過彼此代求.
守望連結整個的群體.
6月7日.
李玉明教師說的題目叫.
「點,線,面」.
很難讀的.
最後那個「面」記得嗎?.
是不同的點加上一些線.
連繫就成為一個球體.
這個球體可以成為一個.
好做的見證.
可以為上主做很多的事情.
頂智妹.
今天你的信仰生活.
只是你一點嗎?.
還是已經在一個大的群體裡面運作?.
這個球體能夠可以運行.
是因為彼此的代求.
可以這麼說.
聖靈的禱告.
是這麼多裝備裡面.
最重要的武器.
不是那把劍那麼簡單.
就算有神的道.
如果沒有聖靈的禱告.

$^{961}$神的道只不過是留於知識.
留於頭腦裡面.
只不過是閱兵的菜式而已.
最後一件事.
自強真正的意義.
就是群體而不是個人.
所以「門徒當自強」這個題目.
不是告訴大家.
我們靠自己可以剛強.
我們不可以.
第十節告訴你.
你們要靠主.
依賴祂的大能大力.
而福瑟書保羅告訴我們.
我們在基督裡面已經互為肢體.
建立基督的身體.
不再是個人主義.
各自修行.
而是連在一起.
向目標前進.
彼此交往.
向世界見證.
向撒旦宣戰.
弟兄姊妹.
敵人不在教會裡面.
敵人不在家裡面.
敵人在屬靈的領域裡面.
他是我們真正的敵人.
在二十一至二十四節.
是以福瑟書的結語也好.
祝福語也好.
是一個整段結語的祝福.
但這個祝福保羅.
提醒了一個很重要的訊息.
正如前面六章十八至十九節提醒我們.
現在聖靈裡面彼此待求一樣.
保羅要告訴我們.
在這個祝福當中.
群體的觀念.
互為肢體.

$^{1001}$合一.
是一個很重要的訊息.
基督的身體.
是要合一的.
基督的身體不是個人的.
我們是互為肢體的.
相關你看回這段經文.
其實是說一個人叫推基古.
他是保羅的同工.
他是附一張信.
大道以福瑟也好 哥羅西也好.
是大信人.
保羅在二十二節告訴我們.
透過推基古.
讓你們知道我們這班同工.
包括我的警方.
現在在監獄裡面.
又讓他安慰你們的心.
不用擔心.
沒事的.
不要為我擔心.
原來在我們群體裡面.
我們可以互相待求.
透過群體.
透過面.
我們互相待求.
你成為我安慰.
你成為我祝福.
彼此關心.
今天我們到教會.
除了來崇拜聚會之外.
我們有沒有互動呢.
所以為什麼我們教會.
既然科技這麼好.
還用回來嗎.
燈油火蠟都曬了.
個個在家裡用自己的電.
崇拜就好了.
省水省力.
連車錢都省了.

$^{1041}$是因為這個不是我們的信仰.
我們互為之體.
我們要回來.
這是一個見證.
我們的互動是一個見證.
我們的互動是可以具體.
向撒旦宣告.
我們是一班合一的精兵.
我要對付你.
魔鬼可以逐一逐一打破我們.
但只要我們同心.
魔鬼都會逃跑.
因為我們個個都成為.
信服上帝的子民.
挺子妹.
鼓勵你投入教會的生活.
今年我們要彼此在靈裡面團契相交.
不只是搞個活動.
搞個口號.
因為這是我們的信仰.
是我們聖經的教導.
真正的敵人是撒旦.
希望我們面對撒旦任何的挑戰.
包括暗戰.
我們能夠恩著上帝的話語.
因為在聖靈裡面.
我們彼此的代求.
我們不讓他得逞.
我們仍然可以站立得穩.
要能夠站立得穩.
我們相信聖靈的能力之外.
同時間我們自己都要自強不息.
是不是.
沒錯.
一隊足球隊打贏.
除了教練的好的佈陣.
好的策略.
球員自己都要努力.
你都不練習.
你都不去跑步.

$^{1081}$你都不去健身.
怎麼能夠跑得動呢.
所以挺子妹.
自強既有個人的責任.
也有群體的需要.
但願我們彼此的努力.
說真心話.
在我們當中.
雖然我們信主都不少年入.
但我相信仍然有很多弟兄姊妹.
你們的生命.
仍然未讓主耶穌基督去更新.
你仍然未讓神的話語去改變你.
你知上帝的話語.
但你未讓上帝的話語.
你不信服上帝的話語.
所以你未被上帝的話語去改變.
你怎麼知道 阿牧師.
你不讀聖經.
你怎麼有上帝的話語.
你一個星期是靠聽崇拜 講道.
如果講完那一次聖經.
你就有一次聖經.
如果講完那一場.
你嫌它那麼長.
長又不行 短又不行 是不是.
好了 你有讀了.
但讀得很少.
來來去去都是那一本書.
幾十年的信仰生活.
一次都未看完整本聖經.
你怎麼能讓上帝去改變我們.
更新我們的生命.
我們就算有讀.
讀了很多次聖經也好.
很驕傲地告訴別人.
我讀了很多次聖經.
但原來我將聖經.
只不過成為我的頭腦知識.
我沒有和別人分享.

$^{1121}$我在神話語的領袖.
你和家人的相聚是講什麼.
有沒有分享.
我們今天在靈裡面的得著.
有沒有分享神對你的經歷.
我們門徒同行小組.
我們有一個很重要的.
就是真誠的分享.
不要緊 壞也不要緊.
大家鬥壞而已.
但我們一起在神面前.
領受恩典.
這個就是門徒同行小組的精髓.
我們是同行.
我知道你是怎樣的.
我知道大家是什麼.
知道大家有什麼困難.
呂牧師也在我們當中.
沒有保留.
我也是成年之一.
我們一起同行.
一起勉勵.
我們大家立志做主要門徒.
壞就一起壞.
起身一起起身.
衝就一起衝.
努力一起做.
鼓勵.
弟兄姊妹.
教會需要你開始自強.
只有每個人都自強的時候.
我們群體才能有真正的改變.
明年是我們一起來完成一個使命.
我們需要有剛強的門徒才能做到.
就如今天是父親節.
勉勵弟兄們.
姊妹們都幫忙聽.
弟兄們.
無論你是否爸爸.
我希望你知道.

$^{1161}$上帝做我們.
是做一家之主.
我希望你剛強.
承擔家庭的需要.
不要逃避自己的責任.
在信仰,在屬靈上.
在生活上.
要剛強,站穩.
不要每次都要太太.
姐姐,妹妹,媽媽提醒你.
我們是弟兄.
我們做兒女的榜樣.
如果我不是爸爸.
我都做其他人的榜樣.
可能你有外甥.
有姪仔,姪女.
就成為他們的榜樣.
沒錯.
我們是不容易的.
沒錯.
教會圈子裡面.
通常都是姊妹生.
爭氣一點.
弟兄們.
父親節沒有甚麼送給大家.
只是勸勉大家.
做一個剛強.
仁愛.
謹守的一個弟兄.
好好為主作見證.
這個家需要有姊妹.
更加需要有剛強的弟兄.
求主幫忙.
同心祈禱.
天父上帝求你的話.
傳到我們的內心裡面.
使我們知道有立志去行動.
因為今天這個世界.
有一個很大的守敵.
蠢蠢欲動.

$^{1201}$時時攻擊我們.
我們沒有辦法靠自己去對付.
唯有你的道.
唯有聖靈的寶劍.
唯有我們彼此的代求守望.
求主的話幫助我們.
奉耶穌基督的聖名.
阿們.
\newpage



\section{出埃及記 32:1-10}
\label{sec:032339}
\textbf{偷天換日}~~~~羅泰然牧師
\newline
\newline
連結: \href{http://www.church.com.hk/acms/content.asp?site=cdc\&op=show\&type=product\&code=032339\&layout=sermon}{\texttt{032339}} ~~~~ 語音日期: 2026-06-21
\newline
\newline
\hyperref[sec:032338]{\small{< < < PREV SERMON < < <}}
~
\hyperref[sec:index_chronic]{\small{[返順時目]}}
~
\hyperref[sec:032340]{\small{> > > NEXT SERMON > > >}}
\newline
\newline
出埃及記 32:1-10
\newline
\begin{longtable}{cl}
\hline
\hline
章節 & 經文 (和合本修訂版)\\
\hline
32:1 & \begin{tabularx}{0.7\textwidth}{X} 百姓見摩西遲遲不下山，就聚集到亞倫那裡，對他說：「起來！為我們造神明，在我們前面引路，因為領我們出埃及地的那個摩西，我們不知道他遭遇了甚麼事。」 \end{tabularx} \\ \\ \relax
32:2 & \begin{tabularx}{0.7\textwidth}{X} 亞倫對他們說：「你們去摘下你們妻子、兒女耳上的金環，拿來給我。」 \end{tabularx} \\ \\ \relax
32:3 & \begin{tabularx}{0.7\textwidth}{X} 眾百姓就摘下他們耳上的金環，拿來給亞倫。 \end{tabularx} \\ \\ \relax
32:4 & \begin{tabularx}{0.7\textwidth}{X} 亞倫從他們手裡接過來，用模子塑造它，把它鑄成一頭牛犢。他們就說：「以色列啊，這是領你出埃及地的神明！」 \end{tabularx} \\ \\ \relax
32:5 & \begin{tabularx}{0.7\textwidth}{X} 亞倫看見，就在牛犢面前築壇。亞倫宣告說：「明日要向耶和華守節。」 \end{tabularx} \\ \\ \relax
32:6 & \begin{tabularx}{0.7\textwidth}{X} 次日清早，百姓起來獻燔祭和平安祭，就坐下吃喝，起來玩樂。 \end{tabularx} \\ \\ \relax
32:7 & \begin{tabularx}{0.7\textwidth}{X} 耶和華吩咐摩西：「下去吧，因為你從埃及領上來的百姓已經敗壞了。 \end{tabularx} \\ \\ \relax
32:8 & \begin{tabularx}{0.7\textwidth}{X} 他們這麼快偏離了我所吩咐的道，為自己鑄了一頭牛犢，向它跪拜，向它獻祭，說：『以色列啊，這就是領你出埃及地的神明。』」 \end{tabularx} \\ \\ \relax
32:9 & \begin{tabularx}{0.7\textwidth}{X} 耶和華對摩西說：「我看這百姓，看哪，他們真是硬著頸項的百姓。 \end{tabularx} \\ \\ \relax
32:10 & \begin{tabularx}{0.7\textwidth}{X} 現在，你且由著我，我要向他們發烈怒，滅絕他們，但我要使你成為大國。」 \end{tabularx} \\ \\
[1ex]
\hline
\hline
\end{longtable}
$^{1}$我想問問大家.
當我說到瑞士這個國家.
你會想到些什麼?.
想到朱古力?.
想到美好的風景?.
雪山?.
還有什麼奢侈品?.
手錶?誰在戴瑞士金錶?.
大家都知道瑞士.
特別是日內瓦.
是一個鐘錶業的中心.
即使現在我們不流行戴錶.
有手機.
或者之前是用石英錶.
但它仍然保留著一個很名貴.
一個很奢華的鐘錶.
大家都知道鐘錶可以賣到很貴.
幾百萬.
一千萬都視乎裡面有多少鑽石.
有多少卡在裡面.
但原來它未成為鐘錶業中心之前.
瑞士日內瓦是由珠寶中心開始演變.
當時在歐洲的情況演變下來.
說到這件事情.
是一件很久很久之前的歷史事件.
可以說是一個宗教改革.
帶來一個誤打誤撞.
結合出這麼特別的改變.
由一個珠寶中心.
歐洲的珠寶中心.
成為一個鐘錶業中心.
後來成為世界現在奢華的鐘錶業中心.
它是因為一道禁令.
一班難民.
和不同的手藝一起結合出來.
造成今天的情況.
這個禁令是什麼呢?.
當時是十六世紀初.
大家知道是宗教改革.
在德國亞丁路德開始宗教改革.

$^{41}$由天主教引伸出來.
有新教.
我們今天說的基督教.
而瑞士漢家二汶是一個宗教改革家.
他在當時整個瑞士.
同樣一個大公教會裡面.
演變一個新教.
基督教.
當中瑞士受到家二汶很大的影響.
當時家二汶看到瑞士的珠寶業很盛行.
但珠寶是用來做什麼的呢?.
大家買到的珠寶.
當然最好有機會帶上身.
很方便可以看到我有這些鑽戒.
有這些耳環.
有這些很漂亮的項鍊.
一來可以炫耀一下自己有這樣的身份地位.
或者這樣的財富.
在家二汶的影響下.
當時政府出了一個法例.
叫「限制奢侈法」.
這個法例是反對任何奢華虛榮的東西.
曾經也教導過.
我們不是要帶很多的裝飾.
很多的財富去裝身.
他們當時是崇尚一種清教徒的生活.
樸素的生活.
這樣的做法.
所以就限制他們.
所以這個珠寶業突然間在限制之下.
當然就式微了.
珠寶就式微了.
但是同一時間.
到了當時的1541年.
過了幾年之後.
法國又有宗教改革.
但是法國的宗教改革就是.
一班新教徒被迫要去受到迫害.
唯有逃跑.
跑到瑞士這個地方.

$^{81}$日內瓦的地方.
而當中有很多的新教徒.
他們特別是一些工匠.
是哪些工匠呢?.
是鐘錶業的工匠.
他們對於鐘錶,齒輪或者那些機械.
是很好的手工.
當時說是歐洲最好的手工這班人.
來到一個瑞士裡面.
那班珠寶商就有些頭腦想想.
不如將這些珠寶放在鐘錶上.
變成我們就不會受到限制.
因為鐘錶我們是要帶的.
要看到的.
會有準時的.
我們有時間觀念的.
於是就這樣結合了.
所以你看到當時一拍即合.
日內瓦的珠寶業中心雖然停頓了.
很快就變成一個鐘錶業的中心.
但是那個鐘錶業中心更加叫做.
奢華的鐘錶業的中心.
合法地炫耀他們的奢華.
所以他們以前是戴一些項鍊.
戴一些首飾.
炫耀自己的財富.
現在他們不是這樣做.
是戴一隻手錶來炫耀自己的財富.
但是美其名他們說我們是要準時.
我們要看時間.
所以這個歷史裡面.
我們看到同班剛才所說的一段出埃及記.
摩西準備下山.
而山下的百姓.
竟然去做一隻金牛獨去敬拜.
很類似的地方.
因為始終人是需要.
在那個人的心裡面.
仍然有一個罪性.
是喜歡去炫耀.

$^{121}$終於有一個敬拜偶像.
終於有物質的崇拜.
可能多年都是不變.
對於我們今日來說.
仍然是成為我們一個很重要的.
提醒學習的地方.
所以我們有幾個題目叫做偷天換日.
偷天換日大家明白.
即是比喻是用一個很大膽.
很巧妙很隱蔽的一個手段.
暗中去操縱一些玩弄的手法.
改變一個事物的真相.
所以變了對他們的鐘錶業.
美其名是我們要有一個時間的觀念.
但是中間仍然是一種奢華物質的崇拜.
而且人的信仰屬靈的生命裡面.
是在這件事看到他們暗中企圖改變的真理.
是去滿足到他們一己的私慾.
題目我們是分開三方面.
第一是放肆的文.
第二是偷換概念.
第三是沉重代價.
為什麼是放肆.
你會看到第一節裡面所講.
百姓見摩西未見人.
就開始有些衣服想自己來.
不需要摩西帶領.
不需要神帶領.
要自己引路.
講回這段經文的背景.
大家會留意到知道.
當然神已經宣告了十誡.
神宣告十誡出在他們二十章裡.
但是祂宣告很特別.
大家知不知道神是怎樣宣告十誡的.
當時祂是親自講給百姓聽的.
跨越了時空直接向百姓講出來宣告.
不是藉著摩西抄寫.
摩西自己拿著法板去講出來.
當時百姓聽.

$^{161}$聽著聽著感覺怎樣呢.
感覺很害怕.
這些來自天上的聲音.
令他們很害怕.
這些聲音很威嚴.
嚇得他們呆在那裡.
當時他們就跟摩西說.
不如摩西你跟我們說好了.
不用由神親自來講.
因為我們真的怕那些聲音.
隨時好像大聲一點.
我們都會嚇到心臟病.
後來神就說.
摩西你上山吧.
去將神所講的說話.
跟百姓講的說話.
透過摩西.
所以摩西上上落落.
很多次上山落山.
最長一次是在二十四章記載.
有四十晝夜.
就是這段時間.
摩西主要當然上山有兩件事.
第一個就是剛才所講.
神親自將那十誡.
摩西抄寫在一個法板.
一個木板裡面.
帶回來.
作為一個藤居.
作為一個納藥.
另外就是講了一大段.
很繁複和很詳細的.
背幕的設計和祭司的條例等等.
講給摩西聽.
所以出埃及記去到一些位置.
我們覺得很多這些條例.
好像很難讀下去.
然後申命記又講一次.
所以你會覺得很困難的時候.
我們讀聖經.

$^{201}$出不到埃及.
然後在山下的人.
看到摩西這麼久都沒有回來的時候.
心裡面.
你說他們很驚恐也好.
他們好像等得不耐煩也好.
或者他們信心沒辦法維持到.
去跟隨一個看不見.
但曾經和祂講過話的神.
他們這麼多的猜測.
甚至說摩西會不會死了.
這樣的想法.
不怕萬一.
不怕一萬最怕萬一.
死了又怎樣.
所以就催促阿倫.
不如我們自己做一隻.
看得見摸得到.
不用等了.
馬上可以有求必應.
滿足到我們自己的金牛毒.
這個假神.
這裡看到這班百姓不想被帶領.
放肆的民.
所以他們說起來吧.
我們做一個神明.
在我們前面引路.
表示我們仍然有所跟從.
但是我們會看到.
他剛剛聽完神所講的十誡.
第一誡不可信靠別的神.
第二誡不可為自己挑剔偶像.
甚至第三誡.
不要妄稱神的名字.
都會犯了這個情況.
他們當然就用了自己身上的金銀.
去鑄造了一隻金牛毒出來.
用來敬拜.
明知十誡的律法下.
是不容許的.

$^{241}$但是他們公然犯罪.
當然他們會有自己的藉口去做這件事.
為什麼他們知道有十誡.
聽過神的聲音.
不可以拜偶像.
十天後就破戒就犯罪.
他們是不是覺得等得不耐煩.
失去了內心的安全感.
所以他們要有金牛毒.
有時等待是不容易的.
等待是痛苦的.
看不見的東西更加是焦慮的.
看不到結果.
又是令我們很煩躁的.
所以今天我們現代人.
特別基督徒.
有時在我們信仰的路程裡.
需要等待.
需要在看不見裡.
你要有一個信心.
一個憑據在裡.
有一個實質在裡.
我們有時不能夠這樣做.
我們有時禮拜日回到教會裡.
我們敬拜.
我們去聽道.
我們去班主學.
我們都知道聖經的說話.
給我們一時的感動.
一時的提醒.
但是去到禮拜一.
開始上班.
很多事情.
生活上的節奏很快.
很多事情我們發覺.
不能夠很容易在我們的手掌握裡.
我們又不知道前面的路怎樣走.
有時在一些很多的決定裡.
我們心裡很想很快即時.
有個答案給我們.

$^{281}$但是有時我們又不願意多讀聖經.
不願意在禱告裡等候神的聲音.
於是我們覺得生活裡有這樣的情況.
有時很不實在.
有時很空虛.
有時聽不到神的回應.
好像離我們很遠.
所以我們很難忍受這樣的情況裡.
我們等待不到神在這裡的聲音.
於是我們很容易受到世界很多的影響.
很多的誘惑.
現在世界盃期間.
最大的誘惑是甚麼?.
賭球.
誘惑是不睡覺看球賽.
不是的.
是賭球.
所以聽到人家說幾比幾之後.
然後說我昨天買了多少.
或者昨天我輸了多少.
然後他還有另一個結果.
球賽的結果.
是另外他自己有一個結果.
這個是賭球的結果.
所以今天我們都要留意.
很多時候我們都會藉著一些事情去滿足自己.
因為一場球賽很快.
90分鐘.
久一點120分鐘.
完了就知道結果了.
知道結果.
你就知道你自己做球賽的結果.
所以很多時候我們就會這樣.
當我們遇到一些這樣的生活壓力.
這些等待的時刻.
無聊的時候.
我們會找一些刺激尋找.
或者不是說世界盃.
我們平時有些時間.
我們又會拿著手機.

$^{321}$不斷在刷手機.
手機裡面有些甚麼東西呢?.
現在有很多人看那些短劇.
短訊.
不斷一路看一路看.
看多久都看不完.
會瘋狂的去看手機.
或者是打遊戲.
或者是圖.
網上購物.
圖一下有甚麼好東西.
這個真的挺好的.
馬上帶來滿足.
買到最便宜的東西.
很開心.
然後很快又會送到家裡.
一種滿足感.
這個是我們今天去填補我們的滿足.
一種金牛毒.
代替了一位要用我們信心等待的神.
這個就是我們覺得.
這些這麼遲這麼晚的.
來到我們屬靈生命.
神的話語聲音.
實在是幫不到我們.
這個是我們今天屬靈的危機.
另一方面當然.
他們除了不是信服神的帶領.
等候不到神的帶領.
亦不甘心在神的僕人裡面的帶領.
你看他們後面跟著說一句.
因為你我們出來埃及那個摩西.
我都不知道他遭遇甚麼事.
說得粗俗一點.
或者應該更加貼切一點.
我都不知道帶我們出來埃及那隻東西.
我都不知道他現在死了去哪裡.
是這樣的意思.
那個摩西.
是一種完全不尊重的說法.

$^{361}$因為的確百姓他們是對神的僕人不信任.
還有一種急躁的想法.
當然他們背後不甘心是神的僕人帶領.
為甚麼摩西帶領我們呢?.
為甚麼神是揀選摩西還是領袖呢?.
一時間又見不到他.
一時又出了去.
一時又上了山的情況.
找不到他就缺乏信心.
記得有一位牧師跟我說.
他說他去了一間教會做主任牧師.
每天早上九時到正的時候.
教會一位堂委就打電話給他.
早晨呀牧師你怎樣呀?.
幾好嗎?.
開頭覺得幾好.
開始回來教會理直理新.
每天都來問候.
原來不是.
每天都繼續打來.
講一些不太特別的東西.
一時又問這樣問那樣.
其實簡單來說.
就好像打卡一樣.
看你是否肯讀.
後來有位同工跟他說.
如果覺得不需要那麼煩打給你.
我教你.
你每天回到教會.
早一點吧.
八點九.
你打電話給他.
打給那個人.
弟兄姊妹.
好嗎?.
你起了身沒有?.
你今天怎樣?.
有什麼為你祈禱等等.
打了幾天之後.
沒有打又沒有打給你.

$^{401}$一個不信任.
一個這樣的想法.
當然大家都很煩躁.
這個故事當然是以前的長輩告訴我的.
我記得我做清朝人工作的時候.
樓下是辦公室.
清朝人很喜歡進我的房間.
失敬無神進來.
我們沒有特別需要違約.
我們通常幾間房間打對面.
不會關門.
突然間進來.
有時我走開他們進來坐下.
進來聊天講很久.
當然我也很樂意.
為什麼剛才我來不見你的?.
後來我發覺這樣的情況.
我就跟他們說.
記住.
你明白.
我這一刻不在的時候.
一定有重要的事情在做.
在另一個地方.
你們等待.
放心.
當然今天我們人.
要對神的僕人有信任.
對神的僕人.
他們知道自己是神自己所揀選,所猜派.
他們都要向神交代.
不單單是向一個組織交代.
更加重要的是向神交代.
因為我們向神的交代.
不是工作上的交代.
是整個生命的交代.
所以我自己也有很大的反省.
有時整個事奉的過程.
我們不是看我在教會裡.
我做了幾件工作.
完成了多少事情.

$^{441}$是整個生命裡的交代.
就算不是在退休那一刻的交代.
是整個生命在天國.
在神的面前的交代.
其實我們每一個人.
都要知道我們的生命.
是需要向神交代.
所以當神派遣你在不同的角色.
不同的崗位裡.
你都要做好你自己.
當然作為神的僕人.
摩西他要神的猜派上山來回.
你想想他都很辛苦的.
上落落.
講起上山.
我最近上山.
我發覺我有高山症.
我初頭以為走路很舒服.
但開始有些頭脹.
最後作嘔.
完全是高山症的反應.
去了三四千公尺的地方.
我初頭沒想過那麼高的地方.
不過呢.
完了之後下山就完全沒事.
我就感受到一個好像八九十的老人家.
上斜路的辛苦的感覺.
完全沒有氣.
你想想摩西上山.
這樣上上落落.
他都很辛苦.
我們要體諒.
我們要明白.
我們要等待.
我們要信任.
但始終摩西都是神的器皿.
如果真是百姓去信靠神.
他就不會因為摩西一時的缺席.
心裡面動搖.
第二樣就是.

$^{481}$偷換概念在第二節第六節.
雖然當時瑞士有限制奢侈法.
但他們能夠成為奢侈品中心.
是完全用偷換概念.
偷換概念的意思是.
在同一個論證裡面.
表面上維持同一個詞語名義.
名義名字.
暗地裡慢慢改變了它的實質內涵.
他們將一個計時的工具.
偷換成為一個裝飾的珠寶.
我戴的不是奢華的珠寶.
我只不過是戴一隻.
好像珠寶一樣的計時工具.
一隻錶.
所以昂貴的鐘錶.
這樣可以讓他們繼續有奢華的物質.
在當時宗教改革裡面.
他們繼續過奢華的生活.
可以說他們繼續有物質的罪惡.
但他們用屬靈的表達.
去洗白了他們內心.
這種炫耀奢華的罪惡.
為什麼這樣說呢.
因為實際上本身是一個.
宗教信仰裡面的偷換概念.
最初信仰的邏輯是這樣的.
清教徒是一班全敬的人.
敬虔的人.
敬虔的人要怎樣呢.
要去愛惜光陰.
因為現今世代邪惡.
所以愛惜光陰的時候.
特別在嘉義文的理論.
預定論裡面.
我們要努力去傳揚福音.
雖然有預定.
誰最終會得救.
但這個要在神的手裡才知道.
我們永遠都不知道.

$^{521}$但我們不可以說.
我們不做任何事.
所以我們得不得救.
神會決定的.
我們不用做事.
但他們的想法是相反的.
就是你不知道.
你每個人都要傳福音.
就要努力去傳福音.
所以你要愛惜光陰.
努力工作.
做得井然有序.
但是他們偷換概念.
我現在戴了一個很精準的手錶.
我是一個懂得運用時間的人.
代表著我能夠有很好的管理時間.
我會有很好的工作.
代表著我能夠蒙神的祝福.
概念上已經偷換了.
但這些以色列百姓.
他們將看不見的神.
換成一個看得見的金牛毒.
當時因為嘉義文覺得.
什麼都不可以戴.
手錶可以戴.
因為手錶本來幫我們把握時間.
但他們竟然變成了一個奢華的東西.
所以你看到這班百姓.
顯示了人類的特性.
最厲害的地方就是.
當我們不願意去順服一個.
或者更高的條例的時候.
一個的律法之後.
我們又不想不敢公然去反抗.
或者我們去公開反叛它.
但我們會用一些曲解偷換的想法.
能夠去遮蓋.
用冠冕堂皇的肅靈的藉口.
包裝我們自己內心裡面.
這些的罪惡,情慾,肉體的生活.

$^{561}$這件事裡你會看到.
亞倫也有自己的問題.
雖然是百姓.
但亞倫的罪也可以看到.
他為什麼會容許這班百姓.
甚至說出來.
我們就注導金牛毒出來.
你會想到其實.
當我們今天去敬拜神的時候.
我們都好像很容易為自己挑剔偶像.
亞倫的事奉是獻祭.
是為這班百姓獻祭.
當百姓說我看不到神.
我們不知道怎麼獻祭.
我不要獻祭的時候.
亞倫怎麼辦呢?.
亞倫好像似乎很熱愛的工作.
不能夠停頓.
不如我也好.
做隻金牛毒.
大家就向著他來獻祭.
我又繼續有工作.
我繼續能夠有服侍.
又能夠滿足我的身份.
多好啊.
可能這也展示了亞倫自己的問題.
我們會不會讓我們的事奉成為了我們的偶像.
只不過我們是希望.
很多人會跟從我們.
會不會因為我們的事奉.
能夠給我自己的滿足感呢?.
金牛毒代表了代替了一個真神.
但是竟然說這個名字是耶和華.
假借了神的名字為自己得到的益處.
這個也可以叫妄稱神的名.
平日的僕人我們很留意.
我們事奉的人很留意.
我們可能會做一件工作.
我們會做很多很多周邊更多的事情.
為了更多人的參與.

$^{601}$為了更多人可以去知道.
為了更多人能夠去擁戴.
越做越大.
但是那個本質上.
希望更多人聽到我的指揮.
聽到我的帶領.
這個是很小心的事情.
或者我們會假借了神的名字.
我們會分配工作.
我們會去做好一件事情.
有時我們都會留意.
這個事奉是否會帶來一種.
只不過是一種堆砌我們自己的功績呢?.
所以亞倫這個事奉.
某情況也是崇拜著自己.
其實他們也是.
為什麼呢?.
他們是在用自己的金銀.
他們的金屬和金鏈.
這些首飾.
去鑄造了金牛.
更加重要的是什麼呢?.
他們學習到埃及一門手藝.
就是這些金屬的鑄造.
他們能夠鑄造一隻.
好像藝術品一樣的金牛.
金牛的學習鑄造是從哪裡來的呢?.
是從埃及的.
因為聖經學者說.
那隻金牛其實就是埃及的牛神.
那個形象.
雖然他不是叫那個名字.
仍然是叫耶和華.
但是他看到我們自己的能力.
我們自己去做神.
我們自己可以做他的偶像去敬拜.
我們不需要看一個看不到的神.
我們看到有渣拿和我們自己建築出來的.
更加有滿足感.
更加能夠自我陶醉.

$^{641}$這個就是偷換了概念.
偷換了神的名義.
當然他們在這件事情裡面.
繼續包裝.
繼續讓他們的情慾放縱出來.
聖經裡面說.
做好金牛獨.
擺在蟬前.
亞倫就宣告.
明天我們就向耶和華守節.
指著那隻牛向耶和華守節.
於是第一日清早.
第六節說.
百姓出來獻梵祭和平安祭.
接著就吃喝.
接著就玩樂.
這個玩耍的字.
在平行經文裡面說.
就是他們的一種放蕩.
一種狂飲作樂.
脫去衣裳.
眾情淫亂.
一個玩耍的情況.
所以他們宣佈向耶和華守節.
是一個偷換概念.
但是裡面卻不是真正的敬拜耶和華.
所以今天我們看到.
用很多宗教的外表.
去掩飾了自己渴望的玩樂.
放縱了肉體情慾.
當然玩樂的本身.
是空虛或者誘惑所造成.
所以今天我們真的要留意.
很多事情我們會因著.
我們不能夠在神的面前.
看不到的神去尋求到.
這種神的說話.
神的溝通.
神的臨到我們當中的滿足.
我們就會用其他的東西去影響.

$^{681}$尤其是在社會生活上.
你們很多的誘惑.
剛才說世界盃開羅華.
全城抗議.
然後可能會心思思.
賭球.
我看球賽這麼久.
都猜到哪隊會贏哪隊輸.
賭一個贏輸.
我不懂得賭.
不過我就是這樣想.
起碼都會贏到一些錢吧.
這個很容易受到誘惑.
但是這些是無底深潭.
這些會令我們泥足深陷.
剛才我所說.
我們就會用很多的想法.
去開脫我們自己的想法.
找藉口.
我不是貪錢.
我只不過是支持那個足球隊.
所以我就去樂著.
因為這樣就增加了.
我投入在這個賽事裡面.
我們沉迷手機短劇.
或者我們打遊戲機.
為什麼呢.
通常什麼時候做呢.
如果你的工作有空.
你有得做.
你的工作真的很有空.
但一般就未必這麼有空.
有空打遊戲機都會被人解僱.
通常是什麼呢.
有時候我看手機都會覺得.
忙完一天.
吃完飯.
但是真的不要坐.
坐身體不夠健康.
現在我吃完飯會洗碗.

$^{721}$洗完碗都好.
都有點靜的時候.
做什麼呢.
打手機.
很容易一看.
不知不覺.
快一個小時.
慢一個小時.
慢慢就洗了兩個三個小時.
都不奇怪.
人就很容易.
好像失去了時間.
去到手機裡面.
所以這種情況我們可能說.
因為我忙我累.
所以我就想輕鬆一下.
看看手機.
看一下東西.
我們就會很合理的放鬆.
神都明白我的.
另一個情況可能我們會.
網上購物.
安慰我們自己.
為什麼網上購物呢.
因為網上購物方便.
還有可以找到.
可以格價.
現在按個鍵就看到.
誰最便宜.
還有很划算.
所以買完之後覺得.
心態上好像應該這樣做.
因為可以省錢.
但你會發覺.
你不斷不斷買.
越來越多累積了.
我發覺可能有些人.
家裡很多包裹都還沒拆.
因為他知道要買什麼.
不過東西不等用.

$^{761}$放在那裡.
很多包裹還沒拆.
放在那裡.
一邊堆積一邊堆積.
其實是想填補.
內心空虛.
對物欲的那種空虛.
所以這個是.
向耶和華守節.
其實背後.
有一個不能滿足的.
心情在裡面.
要讓這些玩樂去取代了.
我們能夠真正聆聽.
和神的關係.
所以這是我們要留意的地方.
所以對受.
他們那個是付出沉重的代價.
在第七到第十節那裡.
耶和華跟摩西說.
你下去吧.
他們已經敗壞了.
他們是偏離了我吩咐的道.
接著就說.
這群是硬著頸汗的百姓.
一個很負面的評價給他們.
接著還說.
你不要抓住我.
我現在生氣得不得了.
要滅絕他們.
這個就是耶和華神說的.
這個金牛毒的陷阱.
是帶他們來到這樣的情況.
他們還以為自己沒有犯罪.
還覺得自己能夠可以更快更及時.
去解決自己面前的憂慮.
或者是解決了很多.
摩西不用怎麼帶領我們這麼麻煩.
我們自己都會走的這樣的情況.
偏離了神吩咐的道.

$^{801}$就等於失去了神的同在.
神吩咐這個十誡的律法.
他們不是去限制我們的自由.
而是希望我們在危機四伏的路程裡.
這個人生路程裡.
保護我們.
這個很像一個欄杆.
但這個欄杆是圍欄.
是保護我們.
不是困住我們.
但我們會覺得困住了我們.
所以我們很想突破這些聖經的教導.
突破神的道理在裡面.
但是耶和華神就說太快偏離我們的道.
跟神建立不到一個親密的關係.
一點也不想聽神的聲音.
又不想聽到神那麼多的說話.
想自己走前面的路程.
他們失去了神的保護.
失去了神的引領.
這個是他們今天的問題.
當然另外一班硬著頸汗的人.
這個是他們會受災的原因.
為什麼呢?.
當時他們在說牛.
最硬頸的牛.
牠不聽主人的拉扯帶著.
主人拉牠的方向.
牠就走去另一個方向.
甚至寧願跳下山崖.
也不願意被主人帶領.
這個是硬著頸汗.
這樣形容這班百姓硬著頸汗.
是一個很負面.
屬靈的標記.
本身是一種羞辱的說法.
今天當然我們仍然很硬頸.
我們對於我們的沉迷.
我們一個地方.
剛才說可能賭球.

$^{841}$可能是去屠東西.
可能是通宵打機.
看短劇.
或者是我們很多自己想做的事情.
一直很沉迷在裡面.
甚至家人提醒我們不要這樣做.
我們都聽不到.
好像硬著頸不肯去聽.
甚至聖靈很微小的聲音跟我們說.
夠了.
停了.
我們都好像沒辦法回頭一樣.
這個是罪.
因為剛才所說前面的情況.
一直在說放肆的百姓.
不願意被神的帶領.
不願意被神的僕人的帶領.
然後他們在當中放縱情慾.
他們假借敬拜的名義.
去用其他的方法滿足他們自己.
這些在神的眼中是一個罪.
罪當然來到懲罰.
但你覺得懲罰神是特意懲罰.
後續大家會知道.
在32章後面.
他們摩西下來之後.
看到這樣的情況就很憤怒.
打碎了法板.
焚燒了金牛毒.
然後吩咐利未人.
用刀殺死了三千個大頭的百姓.
甚至後來百姓遭到瘟疫等等.
這些很血腥或很殘忍的懲罰在裡面.
你會覺得神為何這麼憤怒.
這麼大的懲罰.
但我想跟你說.
罪的毀滅性不是神的毀滅.
大家要分得清楚.
不是神毀滅我們.
是罪帶來我們自己的毀滅.

$^{881}$為何這樣說呢?.
因為神的烈怒.
當然是神的憤怒.
但你要明白.
神的憤怒是因為祂屬性.
是因為神是聖潔的神.
這個聖潔必定要彰顯.
一旦彰顯出來.
人是沒有辦法承受的.
唯有是一些守潔心清的人.
清潔的人.
我們會體會.
光進到黑暗裡.
不會再有黑暗.
沒有可能大家會同時間共存.
光進來後.
黑暗就驅散了.
就等於聖潔的神.
面對罪惡的時候.
是沒有辦法容許罪惡在神的面前.
如果你說罪惡.
這些人可能就在神的面前.
去站立不住,滅絕.
但是.
是不是這樣我們每一個人.
都沒有辦法在神面前.
去站立得住呢?.
但是感恩的神.
是讓耶穌基督.
給我們一個中補.
是救贖我們.
耶穌基督的釘心十價流出寶血.
但是一個救恩.
是給我們一個機會.
向主耶穌基督去禱告.
去認罪,悔改.
因為今天每一個人成為基督徒.
你都要去向主耶穌基督認罪,悔改.
又要求赦免.
得到祂寶貴的救恩,救贖.

$^{921}$我們才能夠坦然無懼.
回到神的面前.
回到施恩寶座的面前.
所以神已經給我們一個.
很好的憐憫.
因為聖經說.
神是施憐憫的神.
神是施慈愛,好憐憫.
祂不願意我們被滅絕.
我們仍然有一個機會.
回到神的面前.
所以今天我們看一段經文.
最終我們明白.
罪一直纏繞著我們.
但是罪是可以解決的.
只要我們在神的面前.
主耶穌基督的面前認罪,悔改.
你就能夠坦然無懼.
在神的面前.
所以今天我們如果.
教會追求復興的時候.
我們都很樂意去認罪,悔改.
但是如果我們心裡面.
仍然很硬著脖子.
仍然不願意去聽神微小的聲音.
這個復興的路.
是很難很難看到曙光的.
一定要全體的百姓.
全體的人都明白.
神的威嚴,神的聖潔.
是值得我們知道恐懼的.
但是同時間我們要知道.
神的恩典,神的慈愛救我們.
我們很簡單.
我們有做錯的事情.
我們要去認罪,悔改.
神一定是幫助我們.
你說那些認不到的.
我們改不了的.
那些認不到的改不了的.

$^{961}$是不是這樣呢?.
你看不到神的能力.
實際上神的能力是很大的.
你認為那些上了賭球癮.
或者是手機癮.
或者甚至性癮.
或者其他不同的什麼癮都好.
神的能力是夠我們用的.
是能夠斷絕我們的罪惡的.
要相信我們要認罪,悔改.
回轉向神那裡.
讓百姓再有機會被神帶領.
所以求主幫助我們.
在今天這段經文裡面.
能夠坦誠地面對我們自己.
送來的生命.
不要遮遮掩掩.
透天瞞日.
瞞得過神.
瞞得過自己.
瞞不到神.
我們求主幫助我們.
坦然無懼來到神的面前.
是求神讓我們能夠有清潔的心.
正直的靈.
繼續走前面的路程.
是同心祈禱.
因為你自己知道我們在罪中打滾.
但是你知道我們是軟弱無力.
你是願意用你自己一個救贖.
救恩的力量.
成為我們自己的幫助.
好似我們能夠可以.
繼續直著認罪,悔改來到面前.
潔淨我們自己.
但是再一次在神的面前.
跟你一同走前面的路程.
你的同在.
你的一起.
我們去感受到這種無比的力量關懷.

$^{1001}$求主幫助我們.
可能我們今天仍然有些情況.
我們不敢面對面.
在神的面前.
或者我們有羞愧.
或者我們有膽怯.
或者我們有罪性.
但是我們知道你是願意寬恕.
原諒我們.
更加給力量我們改變我們自己.
讓我們能夠在你面前.
肯誠地躺開我們自己.
被你所清理.
被你所修正.
我們禱告交託.
奉主耶穌明求.
阿們.
\newpage



\section{歌羅西書 1:23}
\label{sec:032340}
\textbf{A.I.機器人襲地球}~~~~招世超牧師
\newline
\newline
連結: \href{http://www.church.com.hk/acms/content.asp?site=cdc\&op=show\&type=product\&code=032340\&layout=sermon}{\texttt{032340}} ~~~~ 語音日期: 2026-06-21
\newline
\newline
\hyperref[sec:032339]{\small{< < < PREV SERMON < < <}}
~
\hyperref[sec:index_chronic]{\small{[返順時目]}}
~
\hyperref[sec:032341]{\small{> > > NEXT SERMON > > >}}
\newline
\newline
歌羅西書 1:23
\newline
\begin{longtable}{cl}
\hline
\hline
章節 & 經文 (和合本修訂版)\\
\hline
1:23 & \begin{tabularx}{0.7\textwidth}{X} 只要你們持守信仰，根基穩固，堅定不移，不致動搖，離開了你們從前所聽見的福音的盼望；這福音也是傳給天下一切被造之物的，我—保羅作了這福音的僕役。 \end{tabularx} \\ \\
[1ex]
\hline
\hline
\end{longtable}
$^{1}$父親節通常是被人忽略的一個節日.
Mother's Day就有更多的食物,更多的活動.
父親節的禮拜天,有沒有收過你子女給你的禮物呢?.
我記得我兒子有一次在讀大學的時候.
開了兩小時車,連夜回來給我做早餐.
然後驚喜,誰知爸爸不在家,去了港島.
一個很難得的經驗,我最近跟他講的時候.
一早四點起床,開車回來家裡.
我們應該學習怎樣去多謝我們的父親.
有沒有想過買什麼禮物給父親呢?.
如果我們在座,你的父親還在.
有沒有想過買這個呢?.
這個,還沒出來.
這個是什麼來的,這個是去年出的.
現在開始量產了,是你的Human Rights Companion.
即是說可以陪伴你的.
你會發覺這個叫做Area.
我看到這件事的時候,我心裡想過.
現在一個趨勢是什麼呢?.
以後單身一族,或者年老孤寡,一個人住的長輩.
以後就有人形的機器人相伴.
男的也有,女的也有.
那媽媽做家務呢?.
媽媽做家務呢,現在有AI機器人幫忙.
我看看下一張.
AI機器人,下一張.
OK,幫你吸塵,打掃.
還有,連做飯煮菜都可以.
這個是Leon Musk的機器人.
不但如此,甚至連照顧嬰孩.
都可以請AI幫忙.
下一張,再下一張.
以後我相信教會不用請Babysitter了.
買兩個AI機器人回來就搞定.
現在出街,我最近才坐Pastor黃靖.
他租了一輛Tesla.
他載我去Tracy.
我們一起聊天,聊了一個多小時.
車就自動開到Tracy了.
出街,有自動控制的車去代步.

$^{41}$可能有一天,到了一個時間.
只要我們有錢的話,我們可以完全躺平.
什麼都叫AI機器人幫我們做.
我們什麼都不用做.
但不知道大家有沒有想過.
AI機器人襲地球這個日子.
侵襲地球,充滿了地球.
我們想不想過這樣的日子呢.
全部都是AI機器人在我們身邊.
可能變成了生人勿近.
其實呢,我要做一點研究.
早在1950年,科幻小說家Isaac Asmore.
不知道大家熟不熟悉這個人.
其實他已經開始探討人和人工智能.
AI互動裡面所產生的倫理道德問題.
以及它對人類未來造成的影響.
1950年代,都很先進.
到了2004年的時候.
這個背景不知道大家記不記得是什麼.
2004年的時候.
人們就將他的小說改編成一個電影.
搬上大銀幕.
這套電影叫什麼名字.
是Will Smith主演的.
I Robot,我機器人.
電影的背景真的很先進.
說到2035年.
地球上很多工作都被人工機器人AI取代了.
Will Smith是一個偵探,被邀請去調查一個.
人工機器人設計師自殺身亡的一個案件.
後來發現出了什麼問題.
原來公司的中央人工智能中樞.
叫做Virtual Interactive Kinetic Intelligence.
很複雜,簡單叫做VIKI.
出了問題,人工智能中樞出了問題.
出了什麼問題呢.
他說本來機器人做出來的人.
是用來服侍人,幫助人工作的.
但是VIKI在進化的解讀當中.
扭曲了機器人要遵守的三大定律.

$^{81}$那是什麼就不說了.
就確認了一件事,是什麼呢.
人類最大的威脅是什麼.
就是人自己.
人最大衛脅就是人自己.
如果不介意控制的話.
最終就會導致人類自我毀滅.
於是VIKI為了保護人類的緣故.
重新編寫了AI機器人NS5.
這個電腦的程式.
要它監管控制人類的自由.
以免人類走上一個自我滅亡的道路.
這就是一個電影故事.
這是科幻小說的劇情.
但是我們想不到到了今天.
今天是什麼時代.
我們叫做AI時代.
人工機器人開始量產了.
剛才說的那個Area.
其實也是2萬多塊就搞定了.
打一下不用2萬就可以買一個回家了.
在我們被這個趨勢.
智能發展的影響下.
不知道我們有沒有想過.
這個發展的趨勢對我們基督徒.
跟隨耶穌基督的未來.
會帶來什麼影響呢.
或者我們說得好一點.
會給我們什麼鼓勵.
或者會給我們什麼提醒呢.
今天我跟大家分享幾件事.
首先第一件事我們不要忘記.
或者我們要知道.
AI人工機器人.
即是智能機器人.
有一個終極的弱點.
大家猜猜這是什麼.
AI機器人終極的弱點.
從我的角度來看.
他們是被不完美的人創造出來的.

$^{121}$我再說一次.
這些AI機器人.
人工智能.
是被不完美的人創造出來的.
所以AI機器人裡面.
我上次也說過.
他們的資料庫裡面有很多不同的數據.
很多不同推算的方法.
都是人放進去的.
當是人放進去的時候.
就包含了人所有的偏差,誤判.
從基督徒裡面來看.
甚至是有罪性的.
因為AI機器人的演算方法.
就是透過分析大數據.
一大堆的進去.
從歷史到現在.
收集了很多的資料放在裡面.
如果在歷史數據裡面來說.
就是包含了人過去的偏見,歧視.
譬如很簡單,以前我在香港投標做事.
是有潛規則的.
買房買屋.
如果你認識到人.
又會好一點的.
所以無論在招聘,住屋,貸款,投標工程上.
那些歧視和潛規則.
AI機器人去學.
怎樣去學.
可能當成是一個規矩.
是一個正當的方法去做.
所以如果你問他一些問題的時候.
可能AI機器人給你的答案.
就會包含了這些偏差在裡面.
將那個歧視,偏見甚至是不公義.
就自動化了.
我不知道大家聽得明白多少.
不過我盡量幫大家去了解一下.
這個可不可能的.
有可能的.

$^{161}$我相信其實.
我們在座大家都是成年人.
特別我們今天坐在這裡的弟兄姊妹.
我相信我們都不會否認.
人是很軟弱的.
我最近在口裡經常掛念人性的弱點.
因為見到太多人性的弱點.
我們人真是很軟弱的.
很多事情真是很不濟.
有時候我們連自己都不喜歡.
自己現在還有的一定缺點.
而聖經是更加清楚的告訴我們知道.
什麼東西呢.
羅馬書第三章大家已經很熟悉了.
沒有二人.
聖經說.
沒有二人連一個也沒有.
沒有明智的.
也沒有尋求上帝的.
每個人是人人都是背離了上帝.
一起進入祈禱.
沒有行善的連一個也沒有.
這兩三個聖經我相信大家不陌生的.
大家都聽過.
對我們來說到底有什麼意義.
或者我們心裡面可能有人會這樣想.
我們身邊有很多朋友.
他們心地良善.
熱心公益.
雖然不信耶穌.
他們是好人.
最低限度他們比那些自私自利.
選人利己.
甚至是犯法做壞事的人好很多.
沒錯的.
大家的觀察是沒錯的.
但是這些人的好.
這些人的善.
是相對的好.
是相對的善.

$^{201}$人就算是做了很多善事.
動機很多時候.
都是夾著一些不是很純正的心思.
比如你雖然認同你是一個好人.
好人就做這些事.
好人做好事.
我們獲得人的認同.
我們希望怎樣.
好心有好報.
做好事天比佑.
道德上的優越感.
我是一個正人君子.
而且我見到有人有需要我去幫.
還有些人是怕被人罰.
或者以前那些犯毒的人.
犯完毒之後賺了很多錢.
他又捐很多錢給慈善機構.
怎樣?.
他知道自己不好.
贖罪.
自己和自己贖罪.
我隨口說一些例子.
很多時候我們做好事.
問心講句.
很多時候多多少少.
有些不是很純正的心思在當中.
我不是想否定人類行善的社會價值.
這不是我想說的話.
我想是指出甚麼呢.
沒有任何一個人的善行.
是能夠完全純粹.
到完全不自私.
完全只是為了幫助人.
或者對我的基督徒來說.
是為了要榮耀上帝.
我發覺今天我們人.
其中一個最大的問題是甚麼呢.
就是故意不去理會神.
包括基督徒在內.
雖然我們信得主.

$^{241}$沒事我都不去找神.
我都講過.
不信耶穌的人.
故意不去理會神.
很少會主動去尋求認識神.
這就是為何聖經裡面說.
人都是背離上帝.
一起走進去祈禱.
這就是我們的問題.
因為我們不理會神.
所以我們一起走進去祈禱.
我們每個人都是想走自己的路.
我們每個人都是想做自己想做的事.
每個人都以為自己做的事是最好的.
是最對的.
誰不知原來我們偏離了神的心意.
我們得罪了神.
我們自己都不知道.
就好像I, Robot電影裡面.
那個人工智慧中樞.
叫做Viki.
不知是否這樣讀.
她錯誤地解讀了三大定律.
導致成為人類一個最大的危機.
傷害了創造她的人.
親愛的姐妹.
當我們自以為是.
當我們不警覺到自己.
在道德行為上的缺陷和失敗.
認為我們都不錯.
我們就沒有了上進心.
我們就不會想著怎樣去求神幫助我們.
去尋求認識神多一點.
因為我們覺得我們OK.
我們覺得我們可以.
這個就是我們人性的弱點.
所以我剛才說AI機器人.
都承襲了這些人性的弱點.
因為這些東西是我們給他們的.
灌輸在裡面的.

$^{281}$所以裡面有偏差有誤判.
一點都不出奇.
所以AI機器人可能會帶來.
對人潛伏性的危機.
現在有很多人正在研究.
應該怎樣可以好好管理.
怎樣可以使用.
對我們Gertrude Lee來說.
如果我們有這樣的認知和明白.
我們都要小心謹慎.
當以後我們越來越多用.
這些人工智能.
甚至家裡買幾個機器人回家的時候.
應該怎樣做呢.
不過第二樣想跟大家說的是.
AI機器人或者人工智能.
是不能夠操控人的.
《聖經》裡面怎樣說呢.
耶利米蘇十三章.
二十三字裡面說.
古實人豈能改變皮膚呢.
豈能改變斑點呢.
若能,你們者習慣行惡的.
便能行善了.
這個比喻就很簡單.
有些事情是不能改變的.
不是說你想改便改.
你生出來就是這樣.
AI機器人的出現.
即AI人工智能和機器人的出現.
改變了我們社會的生活方式.
會不會呢.
會.
亦都改變了我們每日的生活習慣.
會不會呢.
會.
但是.
改變不了我們的心.
可不可以改變呢.
AI人工智能不能夠完全操控人.

$^{321}$除非我們想讓它操控.
但它不能夠完全操控人.
因為好像我剛才說.
它不能夠操控我們的心.
除非好像上星期我和大家分享的時候.
我們讓AI人工智能.
甚至讓這些機器人.
成為我們心裡面的偶像.
沒有它不行.
它很厲害.
甚麼都靠它.
將它奉約神明.
這個我們自己拿回來.
如果不是的話.
AI機器人,人工智慧.
不能夠去改變我們的心.
所以為甚麼聖經特別提醒我們.
要保守你的心.
勝過保守一切.
要保守我們的心勝過保守一切.
AI智能和人工機器人來說.
我剛才已經說過不能夠操控人的心.
另外一個問題.
我們自己呢.
我們自己能不能夠操控我們自己的心呢.
可以嗎.
另外有兩句聖經大家都很熟悉的.
保羅說的.
你猜到是甚麼嗎.
保羅自己的反省.
保羅說我也知道.
在我的裡頭.
就是在我肉體之中.
沒有良善.
保羅怎麼知道.
他怎麼發覺.
因為他發覺.
立志為善由得我.
只是行出來.
由不得我.

$^{361}$故此我所願意的善我就不去做.
我所不願意的惡我就去做.
兄弟姐妹.
如果我們對自己誠實.
如果我們有機會真是安靜地.
自己人坐下來想一下.
我相信我們絕大多數的人.
都會有保羅類似這樣的經驗.
會不會呢.
我們曾經有多少次.
無論有心也好無意也好.
說了一些後來發覺我不應該這樣說.
我不知道我這樣說原來得罪了人.
傷害了人我們都不知道.
不過我們說了.
我們又曾經有多少次.
想做一些事.
當你想著要這樣做的時候.
你心裡覺得不太妥當.
覺得不太舒服.
心裡告訴你你不應該這樣做.
不過你照做.
有沒有試過.
我們又曾經有感動.
去幫助一個人.
或者想做一件事.
但到後來.
為了種種原因.
我們放下.
沒有去做.
甚至是拒絕去做.
所以保羅.
我有行善的意願.
我們都有.
我們都想做好人.
但是我們沒有行善的能力.
我們所願意的事我們不去做.
我們所不願意的惡.
我們反而做了出來.
有沒有這樣的經驗.

$^{401}$有沒有去舉手.
沒有人舉手.
心裡面舉吧.
各位聽眾.
如果我們有類似的經驗.
是因為在我們心靈的深處.
有些問題還未解決.
我不是看你表面的行為.
神也不是看我們表面的行為.
神是監察你的內心.
不知道會否在我們心靈的深處.
有些問題還未解決.
這個問題是甚麼呢.
可能就是我們的靈魂.
如果你相信人有靈魂的話.
我們的靈魂仍然在屬靈的黑洞裡流浪.
雖然說想耶穌每個星期回來教會.
但是靈魂的歸宿找緊了沒有.
找著了沒有.
很多時候我們發覺.
我們基督徒.
甚至是特別不信耶穌的人.
如果我們警覺到這樣的情況.
我們發覺我們不能夠在屬靈的黑洞裡爬出來.
AI 人工智能幫不了我們.
我們自己也幫不了自己.
所以我們的心思意念.
我們的生活行為.
就會隨著環境的轉變.
就作出相對自以為是的行動反應.
仍然要生活的.
環境變得這麼厲害.
隨悅而安.
這就是我們中國人經常說的道理.
我們找不到讓自己生命靈魂可以復甦的方法.
耶佛所說第二章第三節.
你們死在過犯罪惡之中.
祂叫你們活過來.
那時你們在其中行事為人.
是隨從今世的風俗.

$^{441}$信服空中掌權者的首領.
本為可怒之子和別人一樣.
然而神既有豐盛的憐憫.
因祂愛我們的大愛.
當我們水在過犯中的時候.
便叫我們與基督一同活過來.
很明顯聖經裡的「他」是指甚麼.
這個「他」是指創造宇宙萬物的主宰道.
一模一樣的真實.
我們所相信的上帝.
只有神只有祂才可以叫我們的靈魂肅成.
只有上帝才能讓我們從屬靈的黑洞裡爬出來.
所以現在有很多人在講AI.
現在大力量產AI機器人.
可能會造成一個風潮.
現在就比較貴一點.
都要萬多二萬元買一個次養.
便宜的也有的 一二千都有的.
當我們面對這些風潮湧到來的時候.
我們應該怎樣去應付呢.
我覺得有幾樣東西.
首先不要害怕.
有時候一覺得世界有甚麼轉變.
有甚麼大變動的時候.
人的反應就是怎樣說就怎樣說.
不知道怎樣去應付.
因為太多不清楚的東西會引來懼怕.
我們跟隨上帝不用害怕.
因為我們跟隨上帝.
我剛才講過甚麼呢.
我們要記住人工智能.
人造出來的AI機器人.
是透過人類的智慧創造出來的工具.
應該是用來幫助人管理世界的.
所以我們要學怎樣去善用他們.
第二樣我發覺應付的方法或態度是甚麼呢.
剛才也講過要保守我們的心.
不要本末倒置.
將他們變成我們的偶像.
變成我們的神.

$^{481}$變成反客為主.
除此之外我覺得要讓我們的靈魂蘇醒.
如果我們相信神.
我們相信有靈魂.
我們跟神的關係是如何呢.
到底我們醒了沒有.
因為仍然是沉迷在我們自己的自我世界當中.
求神讓我們能夠有一個清晰的眼目可以看到.
現在我們人最大的盼望.
無論那個世界變成怎樣.
亂成怎樣.
我們的盼望是在基督耶穌的身上.
是在天國的呼音上面.
所以我們要牢牢將福音的盼望抓住.
剛才修詩歌靈魂的鬧.
應該放在耶穌基督的身上.
這是一個有意識的抉擇.
世界不斷在改變.
很多事情我們快到不知道如何去想如何去看.
如何去吸收如何去消化.
明白 沒辦法.
但如果我們相信上帝.
抓住上帝讓我們去應許.
哥羅遜書第一章說.
只要你們在所信的道上恆心.
根基穩固.
堅定不移.
不至被引動.
失去了或離開了福音的盼望.
這個福音就是你們所聽見.
有時傳於普天下萬人聽敵.
在教會裡面你經常聽到福音這個字.
大家知不知道這個福音是在說甚麼.
昨天我在Tracy給他們訓練的時候.
我都在問他們這個問題.
聖經說要傳福音.
耶穌叫我們要將這個福音傳給萬民.
耶穌還沒死.
你心裡面的福音是在說甚麼呢.
其實我們看福音書是很清楚的.

$^{521}$當日耶穌告訴我們要傳福音.
是傳甚麼福音呢.
天國近了.
所以你們要悔改.
這是耶穌所說的福音.
而耶穌的出現.
就是要為人打通這條進入天國的道路.
所以耶穌說道路真理生命.
若不藉著我不能到乎那裡去.
天國近了,你怎能去天國.
靠著我.
透過相信我的說話.
透過相信我死和復活.
你就可以去到神那裡.
進入天國裡面生活.
如果我們是相信這個福音.
就要牢牢抓住堅定不移.
不要放棄這個福音帶給我們的盼望.
這個盼望我剛才已經說過.
無論我們活在甚麼世代.
無論我們周圍的環境.
我們的世界怎樣改變.
我們都可以在裡面.
找到我們生存的價值和意義.
這個就重要了.
今日覺得很多人生活得很無聊.
我會問為甚麼.
太亂了.
有些被迫躺平.
有些做不同的事.
但很多時候我們會發覺.
我們做人沒有意義.
為了生活而去生活.
這個盼望就不同.
這個福音的盼望.
讓我們無論在任何的境況底下.
都可以找到生存的價值和意義.
最後提醒小妹.
如果我們明白神的聖徒和教導的時候.
就應該在這個萬變的世代當中.

$^{561}$抓住那個永遠不會改變的神.
然後尋找為神做見證.
傳福音的空間和機會.
讓我們不會因為世俗的浪潮.
讓我們迷失在變幻無常的世界裡.
就好像我們生活在極端國家.
即是極端無神論的國家.
裡面有很多弟兄姊妹.
怎樣傳福音.
你為了做見證.
你傳福音會被打壓.
不單止被打壓.
會被抓去坐牢.
甚至會有生命的危險.
被人活摘器官.
但我們看到這些弟兄姊妹.
沒有失去了他們的信心.
沒有離開福音的盼望.
他們仍然在他們的處境裡面.
想方法去不單單找生存的空間.
是找為主做見證傳福音的空間.
難不難.
難.
危險嗎.
危險.
挑戰大不大.
很大.
但為何要這樣做.
因為這就是生存的意義.
去找到做人的意義和價值.
只要能夠帶一個人順而疏上天堂.
天上天君天使都會歡喜快樂.
所以就算環境更惡劣.
他們仍然找到一個.
事奉神的空間和機會.
活出一個有意義的人生.
可能神知道我們軟弱.
我們生活在自由民主的國家裡面.
我們生活的環境和質素.
很明顯比我剛才說過的地方.

$^{601}$國家裡面的人好很多.
而且我們還有法律的保障.
有法律的保障可以自由傳道.
可以和別人分享信仰.
人工智能機器人這個浪潮.
我們應該怎樣去面對.
這些新的發明不應該成為.
我們基督徒跟隨主服侍神的障礙.
相反我們應該想方法.
這就是我們面對的事.
想方法好好利用人工智能.
利用這些機器人去做什麼.
去為主做見證去傳福音.
人工智能機器人可以模仿人做很多事.
而且比人做得更好更快.
在某些事情上.
以後剛才你們看到那個區域.
越來越像人的樣子.
不過如果你看那些影片.
行動都是有些snuggish的.
但以後你怎麼知道.
三,五年之後.
但讓我再提醒大家.
人工智能機器人沒有靈魂.
人工智能機器人是不可以完全被信賴的.
為什麼呢.
我開始的時候說.
因為他是被不完美的人.
不完全的人.
被創造出來的.
而今天你和我是有靈魂的活人.
雖然我們人有軟弱.
甚至我們有罪性.
但是我們的靈魂.
卻會在我們感到孤單,寂寞,恐懼.
無奈,徬徨的時候.
懂得去尋求至高上帝的幫助.
人工智能機器人就不會了.
所以主流電視媒.
人工智能的迅速發展.

$^{641}$現在讓很多人失業.
找不到工作.
再加上政治經濟不穩定.
令很多人現在不知道要走哪條路.
人生的路應該怎麼走呢.
我們怎麼能夠在這樣的環境裡.
仍然找到做人的價值和意義呢.
有時候好像很困難.
我不知道在這兩個星期.
剛剛聽到中國大陸.
現在有三代同時躺平的情況出現.
有沒有聽過.
三代都躺平.
為什麼呢.
祖父母那一輩就退休躺平.
爸爸媽媽那一輩被離職要躺平.
新一代剛剛畢業找不到工作.
又要躺平.
以後當這個人工智能機器人大行其道的時候.
只要你有錢.
我剛才一開始說.
你有錢就多買幾個機器.
我們都可以躺平.
所以無論我們的生活環境.
在資源充足的情況下.
或者有其他的缺陷.
其實我們都會同樣有機會找不到.
或者失去了做人的意義.
靈魂.
我說人的靈魂.
人的心靈.
那種空虛難耐的情況.
一定會出現.
耶穌來.
就是要告訴我們.
神的國度已經藉著他的誕生,死和復活臨到這個世界.
祂可以讓我們已死的靈魂復甦.
並且賜給我們一個全新的生命.
這就是我們今天福音的主題.
讓我們可以體會到.

$^{681}$當我們用耶穌的生命的時候.
讓我們體會到神的愛.
和經驗到神的同在.
而且我們有足夠的力量.
去面對人生的挑戰.
去過一個有意義的人生.
但我們必須做一個抉擇.
我不知道我們在座有沒有誤導的朋友.
還沒有決志信主.
就是要相信耶穌真是救主.
耶穌是神的兒子.
然後將我們靈魂的牢.
牢牢地存在耶穌基督的身上.
如果你還沒有決志信耶穌.
你今天願不願意做一個抉擇呢.
而我們信耶穌的人.
求神讓我們能夠看到.
當我們和神之間有了隔膜.
那是什麼呢 大家看不看得到.
我也相信耶穌的.
但有時候我們好像看不清楚.
神在前面帶領我們.
阻礙我們看著神的是什麼呢.
求神的靈幫我們能夠知道.
然後將這個障礙拿走.
如果今天你坐在這裡.
你發覺你屬靈的生命好像沒什麼長進.
可以停在這裡多一段時間.
跟隨耶穌基督.
我也相信耶穌的.
不過好像越來越沒力.
事奉上帝好像越來越不起勁.
就算服侍神也沒什麼心機.
弟兄姊妹我們要小心.
求神的靈幫助我們.
讓我們從那個懶洋洋沉睡的靈裡蘇醒.
千萬不要藉著工作.
藉著娛樂或者其他的東西.
去逃避上帝.
去逃避面對這個萬變世代.

$^{721}$帶給我們的挑戰.
無論我們的世界變成怎樣.
我不是第一次說.
今天早上是吧.
其實聖經都有一個提醒.
教我們應該怎樣做的.
尼家先知第六章第八節怎樣說.
世人啊.
耶和華已指示你何為善.
他向你所要的是什麼呢.
只要你行公義好憐憫.
團謙卑的心與你的神同行.
弟兄姊妹無論我們生活在什麼世代當中.
無論我們的環境是怎樣.
我們都要遵從上帝給我們的指示.
就是要什麼.
伸張正義.
對人實踐耶穌的愛.
並且要恭敬謙卑的與神同行.
跟在神的身邊.
看著神做什麼.
跟著去做.
今年五月十七日.
在佛羅里達州Jacksonville.
Hannah Park的海灘.
有二千五百五十二人.
集體受浸.
嘩一個很震撼的場面.
可能大家不知道.
其實在二零二三年在加州.
Private Cove.
Private海盜.
Private Cove的海灘.
有多少人集體受浸.
四千一百人.
然後去年二零二五年.
在南加州的Huntington Beach.
有七千七百人.
一起在海灘受浸.
有時我們會印象.

$^{761}$美國的教會不行 一直沒落.
主要崇拜實體聚會.
我上星期也說了.
越來越少.
不足三分一的人.
未能浸水.
但原來最近這三年來.
在美國不同的地方.
出現了靈性復興的跡象.
而在全世界.
特別是非洲和中南美洲.
信耶穌的人數是Exponential Growth.
是指數級的飆升.
大家知不知道 神正在用.
近代發明的通訊科技.
互聯網 社交媒體.
還有其他傳播的工具.
將天國的福音和聖經的譯本.
傳到全世界各地.
從前那些福音被封鎖的地方.
現在神又藉著政治的動盪.
經濟的恐慌.
激發起很多人.
尋求物質世界.
無法提供的終極真理.
生命的價值和意義.
所以最近的聚會.
AI機器人越來越多.
我家裡也有一個吸塵的機器.
但聖經告訴我們.
用這東西激發我們.
莊稼已經熟了.
可以收割了.
我們沒想過.
我們可以怎樣使用AI.
人工智能和機器人.
作為工具.
去將天國的盼望和榮耀.
帶給人.
最後一首胡應的詩歌.

$^{801}$你看看那個場面.
AI做出來的話.
可不可以幫助我們有新的啟發.
不單單聖經這樣說.
剛才AI做出來的視頻.
也提醒我們.
耶穌基督昨日今日.
直到永遠.
不單單仍然存在.
而且已經坐著為王.
我們今天來見到一個.
這麼偉大的主宰.
AI雖然能夠做出來這些影像.
都不能夠代表天上的情況的一絲一毫.
但已經足以令我們.
能夠在你面前.
謙卑恭敬.
仰望祈求你的憐憫.
幫助我們天父上帝能夠抓住.
你自己.
抓住你被人福音的盼望.
將我們的靈魂交託在你的手中.
以致我們能夠在今世活著的時候.
學習怎樣跟隨耶穌基督.
找到每一個應有的機會.
去為主做見證傳福音.
因為很多人.
很多人.
還有很多人.
他們沒有這個盼望.
既然上帝已經在這個世代.
做了這麼多的工作.
我們說跟在你的身邊.
我們說要與你同行.
賜給我們力量.
做你現在要做的事.
我們祈求天國早日降臨.
但願耶穌基督的恩典.
天父上帝.
你向我們所彰顯的大愛.

$^{841}$聖靈的感動帶領和幫助.
只因我們前面的道路.
使我們成為福音的出口.
榮耀聖名.
直到我們歡喜快樂.
迎接基督耶穌的再來.
阿們.
\newpage



\section{羅馬書 4:17}
\label{sec:032341}
\textbf{模範的父親 (華)}~~~~葛國光牧師
\newline
\newline
連結: \href{http://www.church.com.hk/acms/content.asp?site=cdc\&op=show\&type=product\&code=032341\&layout=sermon}{\texttt{032341}} ~~~~ 語音日期: 2026-06-21
\newline
\newline
\hyperref[sec:032340]{\small{< < < PREV SERMON < < <}}
~
\hyperref[sec:index_chronic]{\small{[返順時目]}}
~
\hyperref[sec:032342]{\small{> > > NEXT SERMON > > >}}
\newline
\newline
羅馬書 4:17
\newline
\begin{longtable}{cl}
\hline
\hline
章節 & 經文 (和合本修訂版)\\
\hline
4:17 & \begin{tabularx}{0.7\textwidth}{X} 亞伯拉罕所信的是那叫死人復活、使無變為有的神，在這位神面前亞伯拉罕成為我們眾人的父，如經上所記：「我已經立你作多國之父。」 \end{tabularx} \\ \\
[1ex]
\hline
\hline
\end{longtable}
$^{1}$今天是父親節,是嗎?.
父親節,今天父親們都比他生意大一點..
我不曉得前面那些席位是不是表示absent father..
你知道,很多年的時間,當我在講父親節信息的時候,.
當我去查考一些資料,我常常看到的資料都是父親缺席,統計數字..
今天我就決定不講這個了..
雖然那是一個社會的現象,.
但是我們相信這就是為什麼我們需要耶穌基督..
因為沒有耶穌基督,你可以看見這個世界就是如此..
整個世界的走向就是如此..
今天我們要來跟大家分享,.
其實我們當中,我今天這個不是隻是跟父親講的,.
我是跟所有做父親,做母親的,.
甚至做祖父母,外祖父母的,.
我們都需要來看看今天這篇信息..
我們中國人講說,再苦也不能夠苦孩子,.
再窮也不能夠窮教育..
請問這是誰說的話?.
有聽過嗎?.
這個話最早是中國大陸在文化大革命結束之後,.
改革開放,政府為了推及九年義務教育,.
解決農村失學問題,就提出的一個宣傳的標語..
後來就在很多的牆壁上面都貼了這個標語..
鄧小平後來在1970年代到80年代推動科教新國的時候,.
他也強調,千方百計在別的方面忍耐一些,.
甚至犧牲一點速度,把教育問題解決好..
那不是隻有中國大陸講這些話,.
臺灣在前任的行政院長蘇貞昌,.
也曾經引用過..
馬英九總統也講過,.
窮不能窮教育,苦不能夠苦孩子..
所以你會發現我們中國人對於孩子的教育非常看重,對不對?.
我在墨西哥宣講,我告訴你,.
墨西哥人想法跟我們不一樣..
中國人很看重教育,對不對?.
中國父母親看重孩子,.
這是一方面是好,.
但是另外一方面,今天我要問大家一個問題..
父母親,你想,.
如果有一天你離開這個世界,.

$^{41}$你希望把什麼傳承給你自己的孩子?.
上個禮拜我去一個老闆是基督徒的一個公司去講道,.
禮拜一的時候,.
帶敬拜的姐妹就用了一個破冰遊戲,.
那時候剛好要準備過端午節..
她就問說,.
你們最想跟哪一個人,如果可以的話,.
你最想跟哪一個人一起喫粽子?.
她就用了這個破冰遊戲,.
每個人都講了,我也講了,等等這些..
結果他們老闆就講了,.
他說他最想要跟馬斯克一起喫粽子..
所以那一天他就.
其實馬斯克我們都認識吧,認不認識?.
不認識啊,好,不認識..
不認識,我給你看看他是誰好不好?.
這個照片比較小一點..
馬斯克現在可以說是全世界的首富,.
億萬富翁,.
兆萬,不得了..
所以你想跟他一起喫粽子嗎?.
馬斯克,我想很多人很羨慕他,.
很羨慕他這麼有錢,.
很羨慕他事業這麼樣子成功..
可是其實,我今天沒有時間多講他,.
但是我要問大家一個問題,.
如果我們今天可以,.
我們希望把什麼留給孩子?.
馬斯克他一共有四個,.
我要怎麼稱呼他們?.
有四個,他跟四個女人生了14個孩子,.
除了第一任的妻子之外,.
其他統統沒有結婚,.
甚至有一個是他的科技高管,.
他替他生孩子是因為他認為聰明的人應該多生孩子,.
所以他就跟這個女的,.
這個叫做IVF,.
就是人工受孕,.
透過這個方式去幫他生孩子..
所以他想要建一個大的莊園,.

$^{81}$可以把這14個孩子統統能夠跟他住在一起..
但是unfortunately,.
他的大女兒,也就是第一任的妻子生的這個女兒叫做Vivian,.
在16歲的時候,.
她覺得自己應該要變性,.
所以她就去變性了,.
所以她本來是男的,.
現在是女的..
到了18歲的時候,.
她甚至跟她的父親公開地決裂,.
她宣稱她不再有這一個父親,.
也不希望跟他有任何的關係..
她把她自己的last name改成了母親的這個last name..
那這一個女孩呢,你現在看到這個上面,.
她成為Gucci,.
Gucci知道嗎?.
我念錯了嗎?.
Gucci.
對不起,.
所以你就知道我對名牌不太瞭解..
Gucci.
她替Gucci這一類的名牌的包包做時裝的模特兒..
其實想到她的家庭,我們不是在看笑話,.
我們是覺得真的是蠻難過的..
其實我進一步瞭解,我發現說她的老爸也是一個很糟糕的老爸..
好像變成是一代一代地在把這樣子的一個好像咒詛一樣..
親愛的幼子們,.
今天我們在這裡慶祝父親節,.
那我們看看,神要我們做什麼樣子的一個父親呢?.
聖經裡面哪一個父親你認為是模範父親的?.
有沒有想過?.
哪一個?.
好,我們今天讀的這個經文,我們來看,.
我們再來讀好不好?來,預備,請..
亞伯拉罕所信的是那叫使人復活,使無變為有的神,.
他在主面前做我們世人的父,.
如經上所記,我已經立你做多國的父..
在這裡他說亞伯拉罕在主面前做我們世人的父,.
後面又提到說他是多國之父,.
所以亞伯拉罕他不僅僅是他自己這一個家族的父親,.

$^{121}$他也不僅僅是以色列的父親,.
他這裡說他作為歷世歷代,他作為我們的父親..
今天要跟大家用亞伯拉罕來作為我們今天要學習的模範的父親..
亞伯拉罕是一個什麼樣子的父親呢?.
第一個,亞伯拉罕是一個竹壇的父親..
我們一起來讀這段聖經,.
我們讀下面的那一節,來,預備,請..
來到西部人曼利的橡樹那裡居住,.
在那裡為耶和華竹了一座壇..
聖經當中記載亞伯拉罕一生當中有四次竹壇,.
從他一開始蒙昭,.
他一路一路一直到最後他陷以上,.
每一次在這些重要的關口的時候,他都竹壇..
現在我們準備竹壇是什麼意思呢?.
竹壇代表著與神的關係,.
竹壇代表著禱告,敬拜,親近神..
所以亞伯拉罕是一個竹壇的父親,.
他無論到哪個地方,.
他第一件事情一定下來,.
他第一做這件事情是什麼?就是竹壇..
親愛的同志們,.
亞伯拉罕沒有說今天到任何的地方去買房地產,.
就好像溫州的商人一樣,.
啊,一到一個地方買啊,炒房啊..
我們做父親的,我們做父母親的,.
很多時候有能力的話,.
我們都想要替孩子能夠多喫一點東西,是不是?.
最好呢,.
以前我記得我在年輕讀書的時候,.
我一個同學,.
他的父親給他們每一個孩子都預備了一棟房子..
哇,40年前的臺灣哦!.
我們當時覺得好像,哇,真的,這個父親太好了!.
但是親愛的同志們,.
我們是不是到處想要怎麼樣子給兒女.
存留更多的經濟上面,物質上面的幫助?.
但是我們有沒有想過,.
其實更寶貴的是我們怎麼樣子來建立孩子,.
把那個屬靈的傳承,.
把他跟神的關係,這是最重要的..

$^{161}$所以亞伯拉罕無論到哪個地方去,他都足談..
做父親的是家中的屬靈的祭祀,.
祭祀是要獻祭的,對不對?.
所以每一個做父親的,.
在家中有一個很重要的責任,.
就是要帶領全家一起來敬拜神,.
一起來親近神..
我們看見在聖經上當中,.
有一個成為我們的榜樣的是誰呢?.
約伯..
我們一起來讀約伯記第一章第五節,一起來,.
預備,請..
豔豔的日子過了,.
約伯打發人去叫他們自結,.
他清早起來,按著他們眾人的數目獻番祭,.
因為他說,.
他說按著他們眾人的數目,.
這個眾人就指他的子女,.
按著他們眾人的數目獻番祭,.
獻番祭代表的就是他天天的,.
他每一天早上起來,.
他為他的兒女們禱告,.
他下面說,恐怕我兒女犯了罪,.
心中離棄神..
我覺得今天我們這個世代,.
我們更需要這樣子的禱告,.
是嗎?做父母親的..
我們不知道孩子去了學校,.
或者他看手機或者上網,.
他會看到什麼,他會受到什麼影響?.
今天這個世代已經比亞伯拉罕那個時代,.
比約伯那個時代更是厲害,.
罪惡無孔不入,.
你根本不需要走出門來,.
坐在家裡面,手機一上去,.
罪惡就可以進入他的思想..
所以約伯在那個時代,.
他每一天早上都需要起來禱告,.
他不是隻是為自己禱告,.
他也為著兒女們天天在神的面前禱告..

$^{201}$我記得我們三個女兒還小的時候,.
那時候我們在前一個教會,幕會,.
最小的安娜是兩歲..
兩歲,四歲,六歲,大概是這個年紀..
我們這裡大概很多喜樂的夫妻,孩子都差不多這樣年紀..
這個年紀真是太多事情,是不是?.
我們雙方也沒有父母親,也沒有祖父母,.
什麼都沒有,我們就是我們夫妻兩個人..
所以你說,如果師母一個人搞三個,搞得來嗎?.
搞不來怎麼辦?.
搞得來啊!.
你看師母不講道,她得到的支持最多..
但是我的女兒可以做見證..
我在這種情況之下,.
又要幕會,教會很多的一些事工..
晚上要出去探訪,禱告會,聚會..
可是你知道嗎?.
當時候教會沒有禱禱會..
我們從臺灣背景來說,.
早上沒有禱告會,這是很奇怪的..
所以後來我就想說,要怎麼樣開始成島?.
因為小孩子小,所以也需要照顧..
所以後來我們就決定說,就從我們自己家開始..
每天早上七點鐘有一場,然後八點鐘有一場..
兩場禱告會,成島會在我們家..
那麼七點鐘不到,我們預備等等這些,.
弟兄姊妹來了,成島完了之後,他們還喫點早餐,.
完了之後走了,去上班了,去什麼樣子的..
接著八點又來一批,.
這樣子一直搞到差不多要九點鐘..
每天早上九點以前,我們的家全都是人..
所以那一段時間,三個孩子他們長大的時候,.
眼睛一睜開,家裡面都是人..
親愛的同志們,是不是傳道人就應該如此呢?.
教會沒有要求我們..
但是我相信這個對他們以後一生,.
這是我們希望留給他們的一個屬靈的產業..
讓他們看見跟神的關係是最重要的一件事情..
你做父母親的你躲不了,你也隱藏不了..
孩子知道你是什麼樣子的..

$^{241}$你可以在教會里面服侍很多,.
但是你在家中,可能根本你跟神沒有什麼關係..
每一天的daily life如何呢?.
那個東西騙不了人的..
所以做父母親的,不單單是做父親的你要起來,.
我知道我們當中有一些可能先生不在,.
或者是長期是absent father,.
那你做母親的你就起來扮演這個角色..
總之家中要人起來,做父母親的做祖父母要人起來,.
來帶領這個家,使這個家建立一個跟神的關係,.
這是非常重要的事情..
跟你旁邊說這是很重要的事情..
講得有點好像不太重要的感覺..
雖然你可能你覺得你自己都做不到,.
但是親愛的兄弟姐妹,沒有關係,.
我一開始也做不到的,.
我也是掙扎慢慢慢慢的一步一步的才建立這個習慣..
你建立起這一個每一天跟神的關係,.
你建立每一天有好的靈脩生活,.
你纔可以去影響你的下一代,是吧?.
對不對?.
如果一個父母親他自己跟神每一天親近的時候,.
他就會知道這一個太重要的,.
怎麼樣子能夠pass on給他自己的孩子..
所以,鼓勵我們當中做父母親的,.
特別是做父親的,要做一個足壇的父親,好嗎?.
你旁邊如果是弟兄的話,跟他說你要做一個足壇的父親..
第二個,亞伯拉罕他也是一個帶領的父親..
我們來看看,在《創世記》第二十二章的時候,.
這時候是亞伯拉罕上帝要他把以薩奉獻給獻祭,對不對?.
那我請問大家,你們都知道亞伯拉罕的故事,.
以薩是亞伯拉罕他年老的時候上帝應許他給他生的孩子,對不對?.
上帝給他,這個以薩應許給他,.
這個孩子已經是這麼樣不容易得到了..
請問上帝有一天,假設是你的話,上帝跟你說,.
把你的兒子放到祭壇上面,放祭壇上面是什麼意思?.
要殺掉他的,要作為祭屋的..
請問對亞伯拉罕來講,容不容易?.
不容易啊,是不是?.
好,那你看看,這個以薩呢,這個孩子我們來看看他,.

$^{281}$他那一天對亞伯拉罕他發出了一個問題,.
他說父親啊,請看,火與財都有了,.
但番祭的羊羔在哪裡呢?.
以薩問了這個問題之後,這個孩子很單純的,.
他不知道他就是那個祭屋,是不是?.
他就問說羊羔在哪裡呢?.
你做父親的你怎麼回答,如果是你的話你會怎麼回答?.
不容易啊,結果亞伯拉罕就回答了這句話,.
他說我兒神必自己預備做番祭的羊羔,.
亞伯拉罕說神必自己預備做番祭的羊羔,.
這句話彷彿是亞伯拉罕從信心裡面講出來的這麼一句話,.
所以這一個做父親的,他也是一個信靠神的父親,.
所以當他的孩子在人生當中發出這樣子的問題,.
有很多時候當孩子們在問我們說這些人生的一些各式各樣的,.
將來我應該去讀什麼啊,我讀哪一個科系啊,.
我上哪個學校啦或者是這些的問題的時候,.
親愛的同志們,你做父親的你會如何回答呢?.
亞伯拉罕他是一個信靠神的父親,.
所以他說神必自己預備,神必自己預備,.
後面接著就講出了一句話,.
結果後面發生什麼事情?.
當亞伯拉罕講了這句話之後,.
結果有一隻羊羔小被樹纏住了,.
所以那個羊羔就代替了伊薩,.
所以後來就稱那個地方叫做耶和華以勒,.
意思就是耶和華必預備..
親愛的同志們,你相信耶和華以勒嗎?.
可能你在擔心,.
不曉得你將來你的孩子他到底在這個環境這麼惡劣的情況之下,.
他能不能夠找到好的工作?.
甚至現在的孩子,我們有時候會擔心他,.
他能不能夠買個房啊?.
親愛的同志們,你相不相信上帝會預備?.
耶和華以勒的神是不是也是你的神呢?.
戴德森,我們都很認識他,.
這個內地會的創辦人,.
在戴德森他的臨坡辦公室,.
還有他自己的家中,.
他貼著兩副的對聯,.
一邊是以便以謝,.

$^{321}$這個是他用的另外一個文號,.
就是以便以謝的意思,.
另外一邊是耶和華以勒..
以便以謝是什麼意思?.
以便以謝就是到如今耶和華都幫助我們..
那耶和華以勒剛纔我們講什麼意思?.
神必預備,是不是?.
所以戴德森他整個一生的宣教,.
這兩句話成為他的座右銘,.
因為神過去一直到如今,.
他都幫助我們,.
所以我相信神一定會供應我們一切的需要,.
這是他放這兩副對聯的目的..
你知道嗎?.
這個戴德森家族不僅僅是戴德森如此,.
而且他自己親身的也經歷這樣子的一個對神的信靠所帶來的經歷..
有一次戴德森在倫敦預備要前往中國宣教的時候,.
有一次他的口袋裡面只剩下最後的,.
他全身只剩下,可以說他全部的家當,.
只剩下一個五仙令的硬幣..
可是當天晚上他去探望一個非常貧困,.
一個太太在重病的一個平民家庭,.
當他為他們禱告的時候,.
聖靈就感動他,.
要把他身上剩下的這隻有這五塊錢,五仙令,.
要把它奉獻給這一對家庭..
哇!他心裡面很真假,這是他全部的家當啊!.
但是聖靈在裡面很強烈地感動,.
最後他就順服把這個錢就奉獻了給這個家庭..
你知道嗎?第二天的清晨,.
他就收到了一封利民信,.
裡面就有一筆足足是他奉獻的大概幾倍的金額..
所以從那一次的經歷之後,.
他就相信上帝絕對負責他所差派之人的需要..
因此戴德森就講了一句話,.
可以說是一句在整個中國的內地會,.
後來的服飾當中,大家一直遵守的一個原則是什麼呢?.
用神的方法做神的工作,.
斷然不會缺乏神的供應..
Amen?.

$^{361}$親愛的兄弟姐妹,戴德森的神是你的神嗎?.
是嗎?.
戴德森所經歷的神蹟,耶和華以的必供應,.
它也可能發生在你身上嗎?.
我只聽到有一個人說是..
或者你不覺得需要,.
因為我退休金都足夠我下半輩子了..
我覺得我們在南灣,.
我不能講說所有的人,.
我覺得我們在南灣可能很多人會有一個問題,.
就是我們不太需要去經歷神的供應..
因為我幹嘛需要神的供應呀?.
我如果我的經濟能力都很好,.
我幹嘛需要禱告上帝你供應我,你懂我意思嗎?.
所以永遠我們聽到的見證都是宣教士的見證,傳道人的見證,.
好像我們自己覺得那不是我..
我甚至有聽到有時候有些弟兄子會說,.
哎呀,因為你們是傳道人嘛..
好像傳道人的神跟我們基督徒一般信徒的神不一樣哦,是嗎?.
耶和華以樂是上帝跟亞伯拉罕說的,.
還是上帝跟你說的,跟我說的?.
所有人說的,你相信嗎?.
你相信嗎?.
你相信的跟你旁邊說,阿們!.
好了,那你說,啊,牧師啊,道理我都懂啊..
你要去經歷,.
不一定是金錢上面的,各式各樣的,各式各樣的..
讓你的孩子從你做父母親的身上,.
來經歷他是這位信使的神,阿們!.
好,接著我們看第三個,.
亞伯拉罕他也是一個放手的父親,怎麼講?.
剛纔講出來《創世記》22章第9節到第10節,.
他講完了耶和華以樂,上帝必有預備,是不是?.
好了,其實當他講這句話的時候,.
他是不是知道上帝會用一隻羊羔來代替他?.
不知道,他如果知道就不會有這一個圖片出現了..
他不知道,他講完之後兩個人繼續同行,.
然後到了祭壇的地方之後,他就把以撒放上去..
你要記得那時候他已經不是一個baby了,.
這是一個童子了..

$^{401}$好了,他把他綁起來..
我在想,大家都是在講亞伯拉罕,.
但是我也在想,.
你看這個孩子,如果是你的話,你幹嗎?.
想盡辦法,就算不是跟爸爸打架,.
你也會要跳下來,趕快跑,對不對?.
這個孩子乖乖的怎麼樣?.
就父親把他綁在上面..
當亞伯拉罕正拿起刀來準備把他的孩子要殺的時候呢?.
結果怎麼樣?.
天使攔住他..
所以後來纔有那個羊羔的事情..
親愛的同志們,我們看這個故事的時候,.
我們看見這件事情,亞伯拉罕,.
他要把自己最心愛的以撒奉獻出去的時候,.
你自己想想看,這是你老年所得的子,.
這是上帝曾經應許要給你的孩子..
然後你可以有太多的理由跟上帝爭論,.
上帝呀,doesn't make sense..
你既然應許給了我,你為什麼還要我把他收回去呢,是不是?.
但是亞伯拉罕沒有,亞伯拉罕順服..
亞伯拉罕順服,他學習把他自己的孩子放手,.
交託給主..
放手這一個用詞是今天在我們談到家庭,.
所謂的親子教育當中常常用到的一句話,是不是?.
我們中國父母親最難放手了..
根據曾經有人對美,日,韓,中這四個國家的高中生的父母做了一個調查..
那麼在左邊的這個調查出來的結果,他們就說,做中國的父母親是在這四個國家當中最低的,也就是說最難放手的..
做父親的大概是64.7\%,他們會把中國父母親,這邊是講把孩子當成成年人來看待的,這個比例是在四個國家當中最低..
那麼母親是隻有50.7\%,所以表明什麼?母親更難放手,中國的母親更難放手..
好,那另外右邊這一個圖,這是中國大陸自己他們在做的一個survey..
那麼他們做的這個調查裡面,結果顯示說,高達94.6\%的受訪者認為中國家長對孩子大包大攬..
聽懂了嗎?大包大攬是什麼意思?就是事事帶勞,不願放手..
他說這個現象是非常的普遍,也就是說接受訪問的人他們都說,.
95\%的人都說這些中國父母親幫孩子代辦一切,做的一切所做的事情大包大攬,這種現象非常的普遍..
甚至到一個地步,父母親包辦到最後,小孩子的婚姻他也要包辦..
在上海的人民公園去過嗎?有一個相親角,給大家看看..
我看了一個視頻,那個人他說他進去上海人民公園相親角走一圈,他說他估計肯定有兩三百號人..
全都是這些爸爸媽媽在那邊拿著自己的兒子女兒相親的資料..
這邊有個經典的上海上進女生,94年戶籍,身高1.66米,身材佳,美國私立大學本科,.

$^{441}$現就職外企知名公司數據分析師,愛好健身,瑜伽,聽音樂,父親是德企銷售總監,母親退休,尋覓男生28到38歲,1.72米以上,本科以上,北美加拿大留學更佳..
家庭和睦,上海有房,家長微信,非誠勿擾..
上海人民公園有相親角,北京中山公園也有相親角,廣州全中國的這些城市的公園都有這麼一個相親角..
所以有一些老外到上海,到北京去,他還專門跑去這個相親角當作旅遊點去看..
這實在是一個很特別的現象..
你不要笑他,像劉子梅,你會不會對自己的兒女也是把他抓在手中不想放掉,擔心他,怕他不安全,是吧,是不是,會不會?.
我們需要學習放手,放手需要學習,兒女是神所賜的產業,我們需要學習用禱告把他們交託放在神的手中..
我們放手,讓神接手,我們不放手,神很難接手,兒女也很難成長的..
你聽過媽寶嗎,有沒有,媽寶是怎麼樣,去讀大學了,家裡面爸爸媽媽,祖父祖母,外祖父外祖母,全家出動去幫他打理一切,有沒有,這在中國是很普遍的現象..
我們這邊大概不可能有六個人,但是父母親也是一樣幫忙上大學去幫他打理一切,有沒有,有沒有,有沒有..
不曉得是他們幸福呢,還是我們這一代,我想想看我們成長過程當中,去讀大學,考大學,讀書,什麼都沒有,父母親也管不了,六個孩子你還叫他管多少呢..
可是往往就是在這種環境當中,我是我們家中老妖,只要我長大,我需要裝勁自強,處變不羈,我需要自力更生,我不自力更生,我就沒有辦法活了,自然就養成了一個從小就要獨立..
當然我們今天講這個東西,跟你下一代講,你千萬不要跟他講這些東西,講這些東西他會怎麼說,他會告訴你時代不一樣了,你不要跟我講那些..
但其他東西有什麼,屬靈的原則不會改變,有沒有,我們需要學習像亞伯拉罕一樣,我們需要學習把我們的孩子交託在主的手中,.
我們需要按照聖經的原則帶領他們,我們也需要從聖經的角度來操練他們,預備他們,讓他們能夠長大成人,滿有基督長成的神量..
作為父母親的,請記得一件事情,你的兒女不是你的產業,他是神所賜給你的,你只是管家,所以你需要做一個忠心的管家,最後你所做的這一切,你忠心地去做,那麼神就負責他一切..
我記得今年的一月,我大女兒跟女婿,他們就搬到休斯頓去參加他們在聖地亞哥的教會,一個年輕的教會的牧師,他們有負擔要到休斯頓開拓教會..
那這個事情其實是在去年的,應該講是前年的年底就已經開始了,那麼首先是我的女婿,他有負擔想要跟著教會的牧師去支堂,所以他就跟我女兒分享..
那跟我女兒分享的時候,他們還不知道女兒懷孕,等到去年的時候,就是他們女婿跟女兒分享完了之後,結果隔了一段時間,發現說原來女兒懷孕了..
這時候他們心中很掙扎,因為想到去年的七月,教會就準備要全部支堂的人都要出去了,所以教會已經在分享異象,牧師也是在這邊傳遞異象,呼召如果有大家願意要跟著他一起搬去的人跟他一起了..
所以當那個時候他知道自己懷孕的時候,其實女兒也很掙扎,連女婿聽了之後也很掙扎..
因為心想說,你先看算算,那時候大概預產期是九月,所以七月要去支堂.那我請問你,他是怎麼樣去?就等於是大概七八月的時候,就要搬過去了..
所以就覺得說,哎呀,這個其實蠻掙扎的..
事實上他們在San Diego那個地方也纔剛結婚兩年,也剛在那邊settle down,教會了,生活,各方面都是..
可是我今天早上還跟我女兒在電話裡面視頻當中在問他,我本來想說叫他錄一下來,他就叫我自己講了..
他就說其實當他知道自己懷孕,然後就在考慮要不要去支堂的時候,.
他說有一次他們聽到有他們前面的一個支堂的一些弟兄姊妹們分享他們的見證,神怎麼樣帶領他們,他心裡面很感動..
那加上那時候他又在教會參加了一個課程叫做不再一樣,英文就叫做Experiencing God,.
那個課程就是在教你學習說,在你每一天的生活當中你要去留意,神可以透過聖經,透過環境,透過聖靈微小的聲音,.
你要去發現神他怎麼樣來引導你.所以他上了那個課的時候他就開始操練,.
他說就在有一天他們聽到了這個見證之後,他自己在禱告的時候,他就覺得說聖靈很清楚地告訴他要去..
這個我們常常講就好像突然出點閃電或者是怎麼樣子,.
可是他跟我講的就是說,他說其實他就是在這一個過程當中越來越禱告,越來越想要明白神的聲音的時候,.
就在那個時候神很清楚地讓他知道他要去.所以他們夫妻兩個人都有同樣的感動,.
清楚神的引導之後,他們就決定去了.在生孩子之前,他們同時也就在那邊預備要去祠堂..
後來就在今年的一月,也就是孩子出生之後的三個月,他們全家就搬到休斯頓去..
在走之前把San Diego的房子要處理,要出租,等等這些很多的事情,要打包,要運輸等等這些事情,然後又要照顧這個寶寶..
所以這一路的過程當中,可是當我聽到有一天晚上我們家庭在一起禱告的時候,.
我們固定,雖然他們都已經結婚,但是我們固定每個月至少有一次會讓大家一起來禱告..
那一天剛好我們是可以在二女兒的家,我們就一起在那邊為禱告..

$^{481}$那她就跟我們分享這個事情.說實在的,做父母親的,聽到自己的孩子在這種狀況要出去,確實會捨不得..
但是當我聽到他們兩個人在分享神的代理,特別女兒在一面分享的時候,她一面流著眼淚說,.
她說我不願意錯過神的代理..
我自己聽了之後,我自己也流淚,我說主啊,是的,這是我從他們出生就一直希望他們做一個順服你,跟隨你的孩子,.
所以把他交託在天父的手中..
親愛的,我們做父母親的,不僅僅只是把孩子養育長大,讓他們有好的生活,.
更重要的是讓他們能夠認識神,讓他們一生愛神,事奉神,這是我們最大的責任..
你不扮演這個角色,沒有人可以扮演這個角色..
求主使用我們,做父母親的,甚至我們做祖父母的,.
我們一同來養育我們的下一代,建立盡情的下一代..
\newpage



\section{列王記上 3:3}
\label{sec:032342}
\textbf{從大衛家庭看父親對兒女的影響}~~~~李思聰牧師
\newline
\newline
連結: \href{http://www.church.com.hk/acms/content.asp?site=cdc\&op=show\&type=product\&code=032342\&layout=sermon}{\texttt{032342}} ~~~~ 語音日期: 2026-06-21
\newline
\newline
\hyperref[sec:032341]{\small{< < < PREV SERMON < < <}}
~
\hyperref[sec:index_chronic]{\small{[返順時目]}}
~
\hyperref[sec:032343]{\small{> > > NEXT SERMON > > >}}
\newline
\newline
列王記上 3:3
\newline
\begin{longtable}{cl}
\hline
\hline
章節 & 經文 (和合本修訂版)\\
\hline
3:3 & \begin{tabularx}{0.7\textwidth}{X} 所羅門愛耶和華，遵行他父親大衛的律例，只是還在丘壇獻祭燒香。 \end{tabularx} \\ \\
[1ex]
\hline
\hline
\end{longtable}
$^{1}$今年就是我的絕作.
寫了這本書20週年.
所以今天是父親節.
既然這本書是講大衛的家庭.
主要是講家庭的.
相信可能都應該再講一講.
在書中寫關乎父親的一些研究.
以及在大衛給我們看見他的家庭中.
有值得尷尬的地方.
亦都值得學習的地方.
所以今天的主題是從大衛家庭.
看父親對兒女的影響.
講到大衛家庭.
在我的研究中包括了他的爸爸.
因為他是原生家庭.
所以今天我會從大衛的爸爸.
講一下他和子女之間.
特別是大委身上的一些影響和關係.
跟著第二方面就從大衛.
講一下他自己和子女之間的關係.
兩個都是講一個缺點一個優點.
缺點我們就尷尬.
優點就學習.
先來看看大衛的爸爸.
大衛的爸爸是誰.
講到耶西.
我們知道她是波亞斯和路德的孫.
所以大衛就是波亞斯的曾孫.
所以在路得記記載.
當波亞斯和路德結婚之後就生.
然後一直數下去.
他就數耶西跟著大衛.
所以我們見到耶西在聖經裡面.
是出現了有超過20次之多.
所以很多時候都是形容.
關乎和大衛有關的記載.
有些地方就是被人奚落的時候.
如果大家記得那一班惡男人.
去藐視大衛的時候.
就是耶西的兒子.

$^{41}$但大部分都是講耶西的兒子是正面的.
而事實上在以賽亞書.
有記載關乎耶西的本.
講到微塞亞.
救主出於大衛的後裔.
他才是講耶西的本.
羅馬書就講到他是耶西的根.
即是耶穌基督從大衛家族而出.
就講他是耶西的根.
所以大衛和耶西之間.
在聖經裡面的確有不少地方提到.
我們知道我們讀過的經文.
在三國義記上十六章,十七章.
有關耶西和大衛的記載.
我們可以看見第一方面.
我們都曾經提過.
就是大衛不被父親重視.
在神吩咐撒母耳去到耶西的家裡.
去選照將來要高立的以色列王.
當時撒母耳不知道是誰.
只知道他是耶西的其中一個兒子.
所以他就通知他們有一個大宴席.
一起獻平安祭.
所以是一個很豐富的場合.
很專榮,很特殊的日子.
當撒母耳以為耶西是第一個兒子.
又不是第二個兒子.
一直數下去.
發現原來神說不是.
都不是頭七個.
結果撒母耳問耶西你還有沒有兒子.
耶西說還有一個.
這個是最小的.
在哪裡呢?.
他說現在在放羊.
我們已經說過.
在耶西的心目中.
不單止這些重要場合.
這些重要的家庭宴會.
沒有他的小兒子大衛.

$^{81}$他根本沒可能想到.
他將來可能成為以色列的皇帝.
根本不是這個材料.
所以我們看見他一直以來.
跟我們很多家庭不同.
通常「拉仔」是什麼?.
他沒有「拉心肝」.
不過就是藐視或者不重視他比較多.
其實放羊我們說過.
他們家庭都不是很窮.
其實可以找僕人放羊.
這些重要場合.
每個人都有份.
沒理由小兒子沒有份.
詹正義牧師在大衛的研究.
可以說是一個華人學者的專家.
在他的著作裡.
在這方面他都同意.
就是耶西基本上認為.
大衛沒有什麼特別的地方.
也沒有做皇帝的材料.
更加在家庭裡.
沒有什麼很重要的地位.
所以我們見到大衛一直以來.
是被父親比較輕視.
但其實耶西很注重跟兒子的聯繫.
所以全部的兒子都來了.
還有一段經文就在撒米爾記上.
17章 大家記得.
他幾個兒子其實就被選召.
做素羅的軍士.
甚至是兵士長.
所以他都快要跟哥利亞打那場仗之前.
17章就記載一件事情.
17章18 17 17 8 兩節.
我們一起讀一讀這兩節經文.
一日耶西對他兒子大衛說.
你拿一伊發控了的穗子和十個餅.
速速地送到營裡去.
交給你哥哥們.

$^{121}$再拿這十塊奶餅.
送給他們的千夫長.
且問你們哥哥們好.
你向他們要一封信來.
所以這個爸爸跟兒子有沒有聯繫.
有 今天會寄信.
或者寄電郵 或者發信息.
跟他們有聯繫.
而且很想念他們.
希望能夠報平安.
那叫誰去呢.
叫他做走腿.
就是做大衛.
這個小兒子去.
叫大衛去做走腿.
叫他做信差.
做跑腿 做傳遞.
希望能夠得到這三個大兒子的平安消息.
因為在打仗期間.
所以我們知道他很想念兒子.
想念幾個兒子.
但是大衛呢.
大衛就是做跑腿.
可見他真的很偏心.
父親們如果多過一個兒子或一個女兒.
很容易有些偏心.
某些兒子的樣子特別像你.
性格特別像你.
可能你是偏心多些.
或者他特別聰明.
或者某些特別才幹.
你覺得很好.
所以我們見到耶西跟兒子有聯繫.
不過很可惜跟大衛似乎不太好.
不止這樣.
這種父親影響子女.
有偏心.
一個比較重要一個沒那麼重要.
影響哥哥弟弟之間的感情.
所以我們見到第17章和第28節.

$^{161}$當這些哥哥知道大衛將這些東西送給他們.
要問他們安.
然後拿一封信告訴爸爸平安.
這個哥哥.
這個以利亞就是長兄.
第28節.
我們一起讀一讀.
大衛的長兄.
以利亞聽見大衛與他們所說的話.
就向他發怒說.
你下來做什麼呢?.
在曠野的那幾隻羊.
你交託了誰呢?.
我知道你的驕傲和你心裡的惡意.
你下來要看征戰.
即是看熱鬧.
這個大哥真的很不客氣.
弟弟那麼遠來.
特意拿了一醫法控了的燧紙和十個餅.
所以帶著好東西來.
好東西有東西吃.
好東西爸爸想念我們.
好東西弟弟你那麼遠來.
向我們問安.
不是的.
在幼稚園你來做什麼呢?.
你家裡的事不用做嗎?.
你的羊不用看嗎?.
我知道你無端端去誣蔑和控訴他.
我知道你的驕傲和心裡的惡意.
真是無辜.
無端端被他罵一場.
他不但誤解還藐視和攻擊.
所以這個真是很不值.
我們看到原來父親看重這個輕那個.
有時會引起這些家庭的問題.
所以這個都是我們父親們需要留意的.
看到這些子女不要輕看某一個.
我們發覺後來全靠這個小兒子大衛.
結果他們全家都能夠蒙恩.

$^{201}$因為後來素羅免了他們的壞.
給他們尊榮.
讓他們這些子女有機會做一些重要的崗位.
所以大衛大協了成家.
但卻一早被父親表示.
所以我們看到千萬不要輕看某一個.
因為我們發覺真的有很多事要方言才知道.
所以這個就是一個比較負面的影響.
你說耶西有沒有正面的影響呢?.
所以第二方面我們看到正面的影響.
父親對他的正面的影響.
雖然耶西不是很看重這個最小的兒子.
但是大衛的信仰和他的為人.
據我看的確是有一些正面的影響.
從耶西而來的.
我們看到在《撒米爾記》上20章第6節.
當大衛給約拿丹和當時的親信講的時候.
他提到一件事很特別的.
在我讀經的時候留意到20章第6節.
這節經文是約拿丹講的.
即是大衛和約拿丹對話的時候講的.
大衛怎樣講呢?我們一起讀一讀.
「你父親若見我不在席上.
你就說大衛切求我許他回本城伯利行去.
因為他全家在那裡獻蓮仔」.
「你父親」即是約拿丹的「你爸爸」.
即是誰呢?.
是蘇羅.
他說如果你爸爸蘇羅見我不在.
其實當時大衛就開始知道他有危險.
蘇羅開始去逼迫他.
又找一些矛飛他.
然後就知道他的女兒很喜歡他.
開始越來越嫉妒,越來越對付他.
所以他要避一避.
所以平時這些主將等等的人.
都和皇帝坐直.
那為什麼會不見了他呢?.
他說如果你爸爸問我在哪裡.
你就告訴他.

$^{241}$他說大衛原來求我.
即是約拿丹講.
我回到我家就是伯利行.
因為全家在那裡獻蓮仔.
所以這件事我們見到是他們經常有的事.
他們的家庭很著重這些宗教禮儀.
這些蓮仔.
他們敬拜神的時候要專心敬拜.
所以我們發覺其實在大衛的家庭裡.
他們也有一種氣氛.
事實上是敬拜神.
所以當撒母耳尼說他們一起守平安祭的時候.
他們很開心.
一起去舉行節旗.
現在這裡也見到他們經常有這些獻祭的事.
特別是這些大的日子.
所以我們相信他從小就在宗教環境裡受培育.
我們又發覺大衛其實他的靈性的基礎.
就是他早期的建立.
他很多時候自己去到曠野牧羊.
他作詩和神交通.
經歷神的力量.
幫助他擊退獅子和熊.
他年幼的時候就有這種跟神的關係的建立.
奠定了這個基礎.
所以後來我們發覺他人生經過很多風波也好.
經過很多磨難也好.
甚至經過試探和失敗也好.
他都懂得回歸到神那裡.
這個就是大衛的特點.
就是他自小的時候有的基礎.
所以這個自小的基礎.
相信就是我們見到從大衛的家庭裡受的影響.
我們參照大衛這幾個兒子.
他們改的名字.
我特別做了一個研究.
將大衛這麼多兒子的名字.
原文的意思找出來.
所以我們見到這麼多兒子裡.
我們見到有以利亞就是大兒子.

$^{281}$英文的意思就是.
God is father.
中文就是神是父親.
所以我們見到他改那些兒子的名字.
都反映了一些他自己的信念和信仰.
第二個兒子叫阿彼拿達.
就叫做Noble.
他很著重這些德行和尊貴的事情.
第三個兒子就叫做Chimera.
這個就是我們很出名的.
就是《新聞記》第六章.
你們要聽.
你們要盡心盡性盡義盡力.
愛主你的神.
這個Chimera就是猶太人很出名的一個.
這個律法裡面提醒我們愛神的意思.
現在Chimera就是說.
神也就是說聽我的祈禱.
然後我們也看到拿丹葉.
我們很熟悉的.
在耶穌時代有一個拿丹葉.
真理人 耶穌稱他為.
就是他說God is given.
就是神給予的.
然後Radar這個拉袋.
就是說去踐踏仇敵.
意思就是將不好的東西放在他們腳下.
還有一個叫阿仙.
就是代表力量.
而他幫大衛改這個名字.
其實這個名字很好.
叫做Beloved.
就是蒙愛的意思.
所以大衛就是蒙愛David.
所以我們看到其實.
這個耶西他都很著重這些.
和神 和有好的德行.
和著重在聖經裡面.
給我們的一些的應許.
將仇敵踐踏.

$^{321}$賜予我們力量等等.
所以我們看到.
剛才我們看到.
他叫大衛送一些東西給兩個哥.
他又沒有忘記.
他說你記得.
你將那些餅送給誰.
剛才大家有沒有留意.
將十塊奶餅送給誰.
千夫長.
為什麼.
千夫長就是誰.
大官.
他管一千個人.
千夫長下面有八夫長.
還有十夫長.
千夫長他很高級.
所以現在.
你記得去到那裡.
是要有禮貌.
就是不單止送東西給你的家人.
送給你幾個哥.
還記得什麼.
還記得有禮貌.
做一些禮.
送給他的上司.
他的上司 他的高官.
所以這個我們看到.
這個耶西是很懂禮節的.
很懂禮貌.
很懂禮節.
很能夠和家人聯繫.
很能夠和這些家人所處的處境.
他很明白要有禮貌.
要有禮節.
這個很重要的.
而家庭裡面的子女.
全部都是改一些.
是關乎宗教情操.
或者性格塑造.

$^{361}$這些名字去提醒他們.
所以我們發覺.
大衛也都從耶西的身上.
是受到一些感染的.
我們看到後來.
他都很懂得禮節.
他雖然要求喇叭.
但都很有禮貌.
請一些人去祝福他們.
所以我們看到.
這個實在是一個正面的影響.
我不知道大家.
影子妹有沒有回顧一番.
昨天我們獎執退休會.
我們有一些Ice Bricket的時間.
我就帶領其中一組.
我就給一條題目.
給這些組員問一下他們.
每人分享一下.
你爸爸給你.
即是人生當中.
記得有些什麼.
你記得爸爸給你.
一些重要的對你的提醒.
或者一些建議.
你現在還記得的.
所以我們每個人都分享一下.
我不知道你有沒有.
去記得爸爸曾經.
怎樣去影響過你.
爸爸怎樣培育過你.
Peter都跟我說.
聽了我講了很多次.
我爸爸經常講一句成語.
叫什麼.
他說.
學如逆水行舟.
不進則退.
心如平原走馬.
易放難收.

$^{401}$所以這句說話.
娃娃時常都提醒我們.
學如逆水行舟.
不進則退.
心如平原走馬.
易放難收.
所以要約束自己的心.
所以我爸爸都曾經這些.
很多時候.
他有時候古語成語會提醒我們.
暑期的時候.
要我跟弟弟默寫.
諸子自家格言.
所以黎明即起.
山蘇定徐.
要內外整潔.
所以暑假的時候.
要做暑期作業.
就是他給我們默書.
所以這些都影響我們.
我記得爸爸很影響我.
就是他待人有禮.
他非常有禮誠懇.
對人有恩.
很多親戚都讚他.
五叔經常照顧我的表弟.
所以我自小都聽到很多.
對我爸爸正面的.
親戚而來的好感評語.
都影響我們.
所以我們父母教.
見到人要叫.
你有沒有教你們.
有嗎.
不是每個人.
見到人要叫人.
所以這些自小學的.
叫人打招呼叫人.
昨天我們拍攝DAN還說.
見到人要叫人的名字.

$^{441}$這些相信自小都影響我們.
這些叫家教.
又或者是家庭的影響.
未必是教你.
不過是在他身上的影響.
所以這些都很重要.
想一想.
父親給我們一些影響.
我們一些教導.
當然有些不好的.
我們會避免.
但很多好的.
我們都需要去自首學習.
同時可以繼續傳承下去.
所以傳承下去.
我們教導我們的兒女.
都可以.
你知不知道爺爺怎樣教我的.
你告訴他.
所以這是一個很好的.
以身示範.
以身作則.
示範給下一輩.
有時想說一下.
上一輩怎樣影響我.
這是一個很好的教導方法.
你看到爸爸這樣.
因為爸爸這樣跟我說.
接著你跟他說的就有力很多.
所以這也是一個很好的方法.
去幫助我們的子女甚至孫輩.
所以這是很好的.
所以我們見到.
從爺妻身上我們避免的.
不要輕視每一個.
如果你多過一個兒女.
我們知道每一個都重要.
更加不要讓他們自小就感受到.
你那種偏心.
因為它會很影響他們的心理健康和發展.

$^{481}$和他自己跟朋友.
或者是他們的兄弟姐妹之間的關係的建立.
更加影響他在人面前的信心.
所以這也是我們很重要的.
作為父親.
在靈命和德行上培育我們的子女.
主世代人做好的榜樣.
他見到我們的榜樣都會學習.
誠實的榜樣會學習.
我女兒自小她在學校裡邊.
一直不會出貓.
一直不說謊.
她在小時候曾經有幾次看見我們的榜樣.
所以她就學了.
一直以來她都很誠實.
所以這也是很重要的.
對子女兒女的德行的培養.
在Dorothy Notte講到一個很重要的教養兒女.
她說有什麼呢.
我將她選列.
其實很長的.
她說成長在批評當中的孩子.
她學會責難別人.
當她在很多批評當中的時候.
她爸爸媽媽經常批評她.
所以她很容易責難別人.
在成長的時候有很多嘲諷當中的孩子.
看你多骯髒 看你多笨.
所以在這些成長嘲諷當中.
她就慢慢開始越來越萎縮.
覺得自己都是這樣.
我都是懶 我都是讀不成書.
因為爸爸媽媽都這樣跟我說.
所以這也是很重要的.
如果我們稱讚她.
有些事著重她.
譬如她有些德行.
很好的 建立她.
她就會有自信.
成長在羞辱當中.

$^{521}$她就很容易有內疚感.
因為經常被人羞辱.
覺得自己都是自己不好.
都是自己不對.
所以就會有這種內疚感.
反過來在成長當中.
在寬容當中.
不要緊 錯了 錯你了.
鼓勵她.
她就學會在挫折當中仍然忍耐堅持.
或者更加沉著氣繼續努力.
因為知道我還有機會.
成長當中有激勵.
她就會充滿著信心.
因為鼓勵她.
當我們說不要緊.
這次不可以下次.
又或者其實你已經很好了.
你下次會更好.
所以這個就是.
培育她 激勵她.
讓她更加有信心.
成長當中.
其實都是很公平的.
主事很公正公平的.
她就學會公正.
這些是德行.
所以我們如果能夠.
我們分享和神的關係.
見到你祈禱.
你能夠懂得祈禱這麼好.
那你就鼓勵她祈禱.
所以我們在靈性上去培養她.
鼓勵她.
不要所有事情都稱讚她.
只是稱讚她那些.
世人認為好的東西.
你跳舞很厲害.
成績很好.
入到名校.

$^{561}$這些都是稱讚的.
但是你去稱讚她其他的.
德行.
譬如你這麼誠實.
真的很好.
你能夠處事這麼精緻.
很好.
你能夠倚靠神.
神很喜歡.
我們要在這些上面去鼓勵她.
不要像世人一般.
她成績優異.
只稱讚她一頓.
只是這樣.
所以我們都要在其他方面.
在靈性.
特別在德行上面.
去培育她.
這也是很重要的.
所以我們見到在大衛.
後來因為他認識神.
和神的關係.
他和父母之間.
和兄弟之間.
能夠處得很好.
我們都講過.
如果大家記得.
在三位的二十二章.
我們講到.
當時大衛逃避蘇羅的時候.
他就去到阿道蘭洞.
我們讀一讀這幾節經文.
第一節,第三節,四節請讀.
大衛就離開那裡.
逃到阿道蘭洞.
他的弟兄和他父親的全家聽見了.
就都下到他那裡.
大衛從那裡往摩押的米斯巴去.
對摩押王說.
求你容我父母搬來.

$^{601}$住在你們這裡.
讓我知道上帝要為我怎樣行.
大衛領父母到摩押王面前.
大衛住山寨多少日子.
他父母也住摩押王那裡多少日子.
所以你看見第一.
他的弟兄和父親全部來投靠他.
當然當時大衛已經成為有名有望.
但是我們看到不同了.
和以前真的不同了.
他的哥哥,爸爸,媽媽.
全部都來投靠他.
他就照顧他的父母.
特意把他安置在摩押王.
不用和他一起住山洞.
也不用和他一起逃難.
所以結果大衛住山寨幾場.
父母就安置在摩押王那裡幾場.
所以我們看到大衛能夠表達他的孝心,愛心.
和他和父母之間的關係.
弟兄弟之間的關係.
因為他認識上帝.
因為他能夠用神的愛去愛他的父母.
所以親愛的弟兄姊妹.
今天我們成為基督徒.
可能我們的父母還未信主.
或者我們回想起.
他們不在也好.
有些事情不是那麼好.
我們也可以用神的愛去描述.
用神的愛去和好.
用神的愛去建立.
用神的愛去做復和的工作.
所以作為我們基督徒的弟兄姊妹.
我們也能夠對我們的父親.
有這種美好的見證.
透過認識神.
透過和神和好.
我們也能夠和父母和家人和好.
所以這是一個很寶貴的.

$^{641}$耶西和大衛之間的關係.
我們有的一些提醒.
接著我們看看大衛自己的家庭.
大衛的家庭.
總括來說.
有很多很多的失敗.
所以我們看到.
從他好幾方面.
第一就是他可能太著重國事.
所以忽略了家庭.
這也是很重要的提醒.
所以工作很重要.
事奉很重要.
但是家庭也不可以忽略.
事實上他在私下的家庭生活和管教兒女上.
有很多都是很被動的.
很疏忽的.
甚至逃避的.
逃避就是avoidance.
疏忽就是negligence.
所以有時很passive.
不能夠主動去做父親.
應該是有的.
不過事實上今天我們看到很多父親都有這種問題.
就是很passive.
什麼都說媽媽搞定.
有時我也有這種危機.
我是女兒.
女兒還說媽媽搞定.
但是很多時候父親不能夠做一個領導.
和主持整個家庭的屬靈的方向.
和去做公道.
做公正.
是一個很重要的柱石.
所以這個也給我們一些提醒.
就是我們都需要主動.
不要逃避.
有問題需要處理.
還有更加不要疏忽了我們一些的責任.
所以這個事實上大衛給我們看到.

$^{681}$是一種很深的提醒.
不過他在所羅門身上.
我們看到有一個很好的屬靈的影響.
所以我們先來看看父親壞的榜樣.
大衛和他的子女之間的影響的問題.
剛剛上個禮拜六.
喜樂團契請我講了一個主題.
專題就是如何與成年子女相處.
所以這個我用大衛的例子來分享.
其中有一兩點是有點相同的.
不過很多在上一次都沒有時間講.
首先我們第一次看到.
在撒米爾記下十三章發生一件事.
是聖經裡面很慘痛的一件事.
也很真實的一件事.
也很少見的一件事.
就是家庭亂倫的問題.
就是大衛的大兒子暗亂.
他強暴了同父異母的妹妹塔瑪.
所以這件事是一個很大的黑點.
一個很大的遺憾和恐怖的一件事.
但是學者安奴在撒米爾記下十三章.
這件事發生其實原因就是近因.
原因就是大衛一律之利的問題.
近因就是因為大衛犯罪的問題.
所以我們看到上兩章.
十一章十二章就是講到大衛的犯罪.
大衛剛剛犯完姦淫和殺人罪.
所以這種情況一落到十三章.
在這個罪事文學裡面.
將這幾章連在一起描寫是有意思的.
所以我們看到一開始十三章.
講完大衛這件事之後.
馬上就開始形容暗亂.
強暴了同父異母的妹妹塔瑪這件事.
是一件淫亂的罪.
所以這兩者是有很大的相關性.
所以大衛以他的權力去凌辱婦孺.
同樣暗亂是用他的權力.
即是哥哥和王子的權力.

$^{721}$去凌辱了弱質女流的妹子.
但問題是這樣.
當三米爾記下發生這件事的時候.
大衛的反應是什麼呢.
大衛的反應就是大衛王聽見這事.
馬上照暗亂來罵他一頓.
而且罰他.
而且可能要彌補這件事.
有沒有呢?沒有.
聖經就形容他大衛王聽見這事就甚發怒.
就甚發怒.
所以他只是一種情緒上的表現.
只是自己生氣.
他沒有正面的半點質疑犯罪.
或者是不合國家的律法.
不合上帝的律法更加.
他沒有主持公道.
又沒有去安慰他的女兒.
他的女兒無論如何都受了這麼大的衝擊.
這麼大的折磨.
應該做父親的都主動去關懷她.
安慰她,為她想出一些出路.
和為她主持公道.
我們發覺當他媽受了極大的羞辱.
他連衣服都扯爛了.
蒙頭流淚,關閉自己.
有一個很大的情緒的憂鬱開始的時候.
這個大衛一點都沒有做.
所以我們發覺這證明了.
他逃避忽視和很被動的情況.
有學者認為這個問題出於他自己良心有愧.
因為大衛自己都做了這些虧心事.
他又沒有能夠有正義的表態和力量去處理.
在家庭裡的淫亂和違反神的心意的事.
所以我們自己作為父母,特別是父親.
提醒我們以身作則是很重要.
或者我們自己去有一種好的榜樣.
是很重要的.
才能夠去處理家庭裡的一些問題.
當然無論如何我們都是以神的律法為首.

$^{761}$雖然我都是不完美的.
我們都可以向子女說.
雖然我都是不完美的.
但我向神認罪,神幫助我.
我們都可以幫助兒女在錯失裡回到神那裡.
雖然我們自己不完美.
不過我們都可以做這個榜樣.
就是我們能夠認罪,悔改,回歸到神那裡.
所以去鼓勵我們的兒女都同樣做這件事.
這個亦都很重要.
所以一位學者培恩說.
大衛沒有樹立一個好的榜樣.
相反是一個壞的榜樣.
暗戀因此看到爸爸也是這樣.
用權力去強暴其他人.
他也可以用權力去欺負他媽.
所以一連串的悲劇.
正如拿丹先知預告.
刀劍不離開他家.
我們看到後來那些暴力.
互相仇殺.
結果阿沙隆這個暗戀的弟弟.
其實是他媽的親哥哥.
就殺了暗戀.
阿沙隆就殺了他哥哥.
後來還要殺了他爸爸.
因為開始逼迫大衛.
所以一連串的悲劇就開始.
Dr. Donald Joy一位基督徒學者.
他說這個叫第一課.
形容父母如果行為榜樣.
我們影響子女的時候.
是所有他人生當中上課的第一課.
因為他未去學校之前.
就是我們教他.
未去學校看到這個世界之前.
就在我們身上看到天父是怎樣的.
或者愛是怎樣的.
或者道德是怎樣的.
或者婚姻是怎樣的.

$^{801}$或者家庭是怎樣的.
他未上課之前.
這一課就是第一課.
所以著名的邱清泰博士.
他曾經來過這裡.
他是家庭輔導機構的創辦人.
他的師母竹君.
他亦都引用Gary Collins.
一個基督徒心理學者.
他說就算一個細蚊仔看見.
關乎天堂地獄天使.
或者上帝的色彩的圖片.
但是有一個殘酷沒有愛心的父親.
你告訴他天父是多麼充滿慈愛.
多麼愛護我們.
都是一個愚蠢的比喻.
意思就是說.
你講得多好都好.
在你身上卻相反.
他就對天父的形象.
有一種非常錯誤的看法.
所以今天天父是怎樣.
從我們父親的角色.
是很重要的.
甚至勝過母親.
因為我們稱我們的神為我們的父.
所以這實在是給我們一個.
很重要的提醒.
所以當我們能夠被兒女尊重的時候.
他都懂得尊重天父.
所以我們怎樣能夠被兒女尊重呢.
我相信第一.
講過就不要不算數.
經常都食言.
你講過又不算數.
答應你的事又不遵從.
他就覺得天父都不可信.
因為爸爸都不可信.
或者我們有些道德不好的地方.
他看見你都做些不對的事.

$^{841}$你告訴他要做對的事.
其實都是多餘的.
所以這些都是我們很重要的去建立.
他對於父親的形象.
反射對天父尊重和敬愛的形象.
當然所有的事都是出於愛.
當我們告訴他我們很愛他.
我們為他祈禱.
我們鼓勵他.
好像剛才我們所講.
去讚賞他.
或者在一些道德和靈性上去鼓勵他.
這個實在能夠表達到我們對他的愛.
當我們說天父很愛你的時候.
他們很深深感受到天父對他們的愛.
我的女兒差不多二十多歲就開始有個習慣.
以前沒有的.
她就開始長大了.
對神的認識越來越清楚.
她就自己有個習慣跪下祈禱.
所以她每逢現在很多時候.
早上起床她都會跪在床邊祈禱.
我相信有很大的原因.
就是她從小就看見.
我或者師母都有這個習慣.
就是跪下祈禱.
所以我們沒有教她.
其實我們沒有教她.
祈禱一定要怎樣.
坐下還是躺下還是跪下.
但是她自己就去跪下祈禱.
這個我們沒有教她.
只是她覺得她想這樣或者她看見有這個榜樣.
所以今天其實我們.
上一個禮拜畢業生主日.
畢業生主日我們很多時候慶祝.
我們的兒子女兒孫子.
他們學業有成.
這個是很好的事.
我們當然希望他們學業有成.

$^{881}$做一個好好對社會有建設的人.
但是其實我們畢業生的主日.
為什麼我們為他祈禱呢.
其實也很重要就是提醒我們自己.
我們不單止注重他們的學業.
我們也要注重他們的德行和信仰.
所以如果我們誇耀.
我們的兒女讀書有多好.
入哪間名校.
這個不是犯罪.
但是有沒有同時間.
我們的心裡都覺得.
其實更重要的就是他們的德行和他們的靈性呢.
這個是值得我們反省的.
所以我們求神幫助我們.
希望我們的兒女都能夠在愛主.
在他們的德行.
在他們對神的關係上.
能夠得到建立.
很多時候真的可以說.
從我們父母開始.
其實今天橫觀.
昨天我們在獎執會也有談這些東西.
也不只是在獎執會分享的時候談.
有很多我們平時聊天的時候也是.
現在好像越來越難找到好的弟兄.
很多時候姐妹愛主的弟兄.
現在再看一看.
我們弟兄少一點.
姐妹多一點.
信耶穌的心軟一點.
願意信主的多一點.
我們如果未信主的配偶.
我們有一些也是未信主的太太.
但是未信主的丈夫多很多.
今天我們看到.
我們作為弟兄父親.
我們也需要影響我們下一代.
特別是兒子.
現在越來越難找到好的弟兄.

$^{921}$包括下一代.
年輕的很難找到長進一點.
愛神的男孩子.
所以真的很慘.
好的姐妹很難找到好的丈夫.
這是真的.
所以我們發覺.
如果我們有兒子的話.
我們作為父親.
父母當然.
但是父親的影響對兒子更大.
我們更加需要培育一些愛主的弟兄.
所以希望我們彼此互勉更加努力.
我沒有兒子只有女兒.
女兒也一樣.
也要培育她.
所以我們更加彼此互勉.
否則基督教會越來越.
那些叫什麼.
女多過男有句成語叫什麼.
陰性陽衰.
第二我們看到.
大衛值得.
我們學習的在他的家庭裡.
他很看重建立子女與神的關係.
特別是我們以所羅門為例.
大衛和所羅門有一個很獨特的關係.
因為神指示他.
你不可以建聖殿.
但是你的兒子就要負起這個責任.
所以大衛對神很有心.
他不僅沒有忘記約櫃漂流在外邦.
或者其他不應該在的地方.
約櫃應該在回望.
所以他很想把約櫃運回耶路撒冷回望.
不僅如此.
他還想起上帝的殿.
所以他很希望要建耶和華的殿.
他說我住在皇宮那麼漂亮.
但是耶和華卻住在回望那裡.

$^{961}$怎麼行呢?.
所以他就在心意要為耶和華建殿.
但是神說我不要你建.
我要你的兒子.
他就把這個重任.
也把所有的負擔.
來培育他的兒子.
去預備做這件屬靈的很重要的工作.
所以我們知道大衛雖然不是完全人.
真的不是.
不過在靈性上.
他實際上是一個很好的.
給我們去效法的榜樣.
就是無論多失敗也好.
都懂得靠神站立回來.
所以被神稱為與神同行.
重心能夠去心與神追隨.
所以神最後評語他.
雖然他犯過不少幾次過失大罪.
在聶王記上第九章第四節.
神的評語說.
他是全誠實正直的心行在上帝面前的一個人.
所以這是一個很重要的評語.
這個評語不單是神的評語.
也是他的兒子對他的評語.
這個才重要.
因為他實在看見在他爸爸身上.
有這些質素.
不只是聽回來.
我們看到後來所有以色列諸王.
無論南國北國也好.
特別是南國.
恆耀華眼中看為正.
還是恆耀華眼中看為惡.
就以大衛作標準.
所以這個實在有很多經文記載.
恆亮聖經的記載.
而最重要我們看到大衛.
在籌建聖殿的過程中.
他是被聖靈感動.

$^{1001}$指示他應該怎樣做.
照著各樣的樣式.
我們知道最初摩西上西乃山.
就是神指示他照著各樣的樣式.
怎樣做聖文.
怎樣回望.
現在到大衛.
神指示他各樣的樣式怎樣做.
所以你看到神使用大衛.
好像使用摩西那麼厲害.
我們看見他在.
宣導所羅門.
怎樣宣導他呢.
在歷代志上.
歷代志上第二十二章.
我們讀一讀這幾節.
就是大衛和所羅門說的一段說話.
請大家讀.
我而言今願耶和華與你同在.
使你亨通.
照他指著你說的話.
建造耶和華你上帝的殿.
但願耶和華賜你聰明智慧.
好治你以西列國.
遵行耶和華你上帝的律法.
你若謹守遵行耶和華藉摩西吩咐.
以西列的律例典章.
就得亨通.
你當剛強壯膽.
不要懼怕也不要驚慌.
你看到大衛勸所羅門.
記得什麼.
記得就是要遵行耶和華.
你神的律法.
所以他說如果你謹守遵行.
說了兩次.
謹守遵行耶和華的律法.
神就要你同在.
你就亨通.
你就可以為他建造聖殿.

$^{1041}$所以你要剛強壯膽.
你要完成這件事.
所以我們看到大衛在屬靈上指導他.
真誠向他祝福.
一定影響所羅門和神的關係.
所以我們看到所羅門最初的時候.
受大衛的影響.
事實上記載他是非常愛神.
在聶王記上第三章.
這樣形容所羅門.
我們一起讀.
所羅門愛耶和華.
遵行他父親大衛的律例.
你留意到嗎.
所以他說.
所羅門很愛神.
因為他遵行爸爸大衛的律例.
即是說大衛教他的律例.
所以你看到.
我們看到他有一個結果.
就是他愛耶和華.
所以你看到這種的.
當時就是獻殿.
即是殿被建築起來.
一個很大的宗教儀式.
在機片那裡.
來敬拜神.
所以很大的表達.
所羅門對耶和華的敬畏.
以數千的生畜來獻祭.
向神禱告.
而且他的祈禱.
成為我們今天經常所說的.
這個祈禱就是說.
求神怎樣給他智慧聰明.
來治理神的國家.
所以他不是求富足.
不是求兵強.
不是求有錢.
所以神說.

$^{1081}$既然你不求富不求秀.
我就將智慧聰明.
和你不求那些都給你.
所以這個就是神對他的賞賜.
所以你看到.
他禱告以神的國為先.
而且他在祈禱裡很特別.
他的祈禱怎樣說呢.
這個《尼環紀》上第六.
三章這裡的祈禱第六節.
我們一起讀.
好嗎 一起讀.
所羅門說.
你僕人我父親大衛.
用誠實公義正直的心.
行在你面前.
你就向他大師恩典.
又為他存留大恩.
賜他一個兒子.
就是我.
坐在他的位上.
值正如今日一樣.
所以在他的祈禱裡.
這個父親是怎樣的.
他是用誠實公義正直的心.
行在你面前.
這個很寶貴.
正如我剛才所說.
他不只是聽到.
他看到.
看到他父親是誠實公義正直.
行在上帝面前.
所以他就說.
因此你向他施恩.
包括謝罪之恩.
為他存留大恩.
更加為他預備了我.
來坐在這個位上.
為你完成這個.
屬靈的豐功偉績.

$^{1121}$所以今日我們作為父親.
真是.
原來最重要的影響.
給我們看見的就是.
當然我們沒有辦法.
擔保我們的子女日後將來.
我們死了之後.
或者他長大之後.
他究竟還愛不愛主.
事實上後來見到所羅門就一般般.
但是無論如何.
在他開始的時候.
我們怎樣影響他.
我們就看見他怎樣被建立.
是從我們身上的靈性.
是會影響他們的.
所以今日我們見到.
如果他都在我們身上看見.
上帝是值得敬佩的.
值得愛的.
我們亦都用誠實公義的心.
來追求神的.
他們都會受影響.
好像所羅門.
去稱仰他的父親.
和學校的父親的榜樣.
所以在《箴言》裡有很多處提見.
他說.
你們要敬畏上帝.
我兒啊.
用心良苦.
我要敬畏神.
這些提到很多次.
在《箴言》裡面提到18次.
提到敬畏神.
18次提到尊榮和信教神.
所以有三十多次.
都是這種對上帝的勸導.
是對後輩提醒的勸導.
所以我們見到這個真是.

$^{1161}$今日我們再講.
我們很著重我們子女的才藝.
學術.
自小就要很多.
甚麼都懂.
彈琴又叻.
打球又叻.
因為現在的教育制度.
害到我們雞毛鴨血.
每個星期都沒空.
我們盡量.
就算他未必入到.
靠這些入到最好的名校都好.
但如果他.
逢星期六星期日都不回教會.
只顧著打球.
畫畫.
做其他的事.
一過了這個階段.
沒有機會了.
過了這個階段.
你再要他們回教會.
沒有機會了.
所以如果我們.
不把握著這些機會.
來培育他.
除了這些.
學術.
才藝等等.
其實這些.
很叻又得.
一般叻又好.
都不是死得的.
事實上今天很多那些.
出來那些.
做得事很叻的那些.
未必是名校畢業的.
但反而他的人.
他的品格.
他的處事.

$^{1201}$他的智商.
他的情商.
這些可能是更加重要.
特別是他的靈相.
即是屬靈的相數.
所以我們發覺這個都是一個.
我們今天作為基督徒父母.
特別父親們.
對我們的一個提醒.
在最後我們培育.
我們的兒女.
甚至孫子.
現在與神的關係.
有什麼建議呢.
這些小小參考.
第一就是著重教導.
神的話語.
真的很需要.
從小幫他讀.
神的話語.
聖經故事.
不要只是講故事.
你要將那個屬靈的教訓.
提出來告訴他.
然後培養兒女禱告的習慣.
你和他一起祈禱.
慢慢他自己也會祈禱.
分享自己屬靈的經驗.
剛剛前兩天.
我們教務同工.
靈修.
我們每個星期二.
都有靈修的.
教務同工.
張雷牧師.
國語堂牧師.
他分享.
他以前在結婚初期.
教會建堂.
他憑信心.

$^{1241}$奉獻了一筆錢.
其實當時他沒有錢.
跟太太商量.
我們憑信心.
結果真的成功.
一百萬.
日本錢.
都很多錢.
結果他和我們分享.
最近他的兒子.
收了他.
UC Berkeley.
小兒子.
全面獎學金.
很厲害.
他和他看學校.
看學校的時候.
他剛剛星期一找教會聚會.
剛剛那個教會聚會的時候.
就是建堂奉獻.
叫頂尖會為建堂.
來奉獻金錢.
支持教會這個項目.
他就有機會.
他說從來沒有跟他的兒子.
講過他這個見證.
他就有機會跟他的兒子分享.
二十幾年前.
他們自己的見證.
這個很powerful.
所以我想.
我們這些屬靈的經驗.
跟我們這些子女分享.
勝過你讀多少次聖經.
當然聖經都要讀.
但是這個.
就是這些屬靈的經驗.
讓他們體會到神的真實.
和我們父母愛主的心.
他們一定一生都會受用.

$^{1281}$所以這個我們要多一點.
跟我們子女的孫輩.
分享一下你們的屬靈經驗.
這個.
他沒有聽過.
很開心聽你說.
你怎樣經歷神.
怎樣信耶穌.
怎樣去在一些困難當中.
經歷到神的幫助.
怎樣去做你人生的決定.
這些我相信很寶貴.
跟你的子女分享你的屬靈經驗.
然後鼓勵他們去祈禱經歷神.
好像剛才所說.
然後幫助他們盡早清楚信主.
這個很重要.
很小的時候就可以信主.
你不要放過這個機會.
五六歲就可以清楚信主.
所以從小幫他們講清楚福音.
教他們認罪悔改.
教他們信耶穌.
我們不斷為他們祈禱.
這個很重要.
不斷為他們祈禱.
然後如果他是青少年的.
跟我們青少年的同工合作.
怎樣幫助他.
然後在家裡提醒他.
你記不記得.
哪個minister教你怎樣.
然後他回到教會的時候.
他又記得.
爸爸告訴他.
要怎樣配合在教會的生活.
所以這個相輔相成.
在教會時間很短.
在家裡時間很長.
彼此配合.

$^{1321}$然後參與青少年侍工.
如果可以的話.
你親身去到.
我們英文堂有很多父母.
他們自己去做青少年的主學老師.
青少年的輔導員.
所以跟子女一起在青少年侍工.
看著他們成長.
這個更加寶貴.
當然我們語言有些限制.
不過英文堂那些就更加好.
如果可以的都可以考慮.
然後考慮跟他們一起參與事奉.
都很好.
參加短孫.
譬如.
每個人都去郵埠.
賺錢去參加短孫.
所以參加短孫.
帶著你的孫子和子女去.
都很好.
我們有些都是這樣.
小時候帶他們去短孫.
現在他們長大了.
慢慢自己事奉.
真的很好.
我記得.
我女兒現在是讀公共衛生的.
她四年undergrad.
master.
六年都是讀公共衛生.
為什麼從來沒有轉過呢.
就因為她在十一,二歲的時候.
和Anne Mo和我們那幾個.
去巴拿馬短孫.
去到那條村那裡.
看見那些環境那麼差.
她就開始覺得.
我們都不做醫生.
又不做老師.

$^{1361}$本來她想做這兩樣.
我要做這些衛生健康的東西.
你不相信.
竟然在小學的時候.
這個印象在她的心目中.
到她後來申請大學的時候.
她一直以來就做這件事.
現在在政府那裡.
服務這些under-served的社區.
做公共衛生的工作.
原來我們從小跟她一起事奉.
這些經驗.
是很深影響她的人生的.
最後我們見到.
她建立他們的靈敏.
品格和人際關係.
這個剛才我們說過.
所以我們見到.
這個我們都可以去參考學習.
我們了解我們的子女.
個個都不同的.
但我們為她祈禱.
神給我們智慧.
繼續主動.
我們去樹立好榜樣.
以聖經的原則.
自誓引導他們.
無論我們子女孫輩.
所以如果你自己不回教會.
你怎樣可以叫你的子女回教會呢?.
你自己躲在家裡上網.
永遠不回教會.
你怎樣帶他們回來呢?.
你怎樣帶你的子孫回來呢?.
讓我們重新反省一下.
父親很重要的使命.
最後在這裡結束.
斯亞斯醫生.
他說父親最大的人生焦點.
就是做一個自己回歸神.

$^{1401}$和帶領子女回歸神的人.
而關鍵就是全然毀身.
做這件事.
所以這位基督徒的醫生.
他作為一個反省,一個提醒.
叫我們作為父親們.
繼續努力一生.
毀身.
做我們子女甚至孫輩的.
美好的榜樣.
自己先回歸上帝.
帶領他們回歸上帝.
我們低頭祈禱.
請大家同心禱告.
求主承蒙我們的幫助,勉勵.
叫我們好像大衛的提醒.
我們要剛強.
繼續努力.
不斷學習.
改變我們自己.
總不會太遲.
因為只要我們再世一日.
仍然會影響我們的兒女.
求神幫助我們.
就算我們年老發白的時候.
我們怎樣信靠神這個榜樣.
都會影響他們日後對於你的認識.
主求你今天.
讓我們對於我們的兒女.
我們的孫輩.
都能夠產生正面的影響.
幫助他們愛你,愛人.
希望我們都能夠建立.
一群大能的子民.
無論弟兄或姊妹.
更需要弟兄.
讓你的國能夠被推薦.
你的名得到榮耀.
讓更多基督的精兵被興起.
求主隨聽我們禱告.

$^{1441}$奉耶穌基督的名.
\newpage



\section{士師記 8:33}
\label{sec:032343}
\textbf{聖經人聖經事:忘恩負義的以色列人}~~~~黃克勤牧師博士
\newline
\newline
連結: \href{http://www.church.com.hk/acms/content.asp?site=cdc\&op=show\&type=product\&code=032343\&layout=sermon}{\texttt{032343}} ~~~~ 語音日期: 2026-06-20
\newline
\newline
\hyperref[sec:032342]{\small{< < < PREV SERMON < < <}}
~
\hyperref[sec:index_chronic]{\small{[返順時目]}}
~
\hyperref[sec:032344]{\small{> > > NEXT SERMON > > >}}
\newline
\newline
士師記 8:33
\newline
\begin{longtable}{cl}
\hline
\hline
章節 & 經文 (和合本修訂版)\\
\hline
8:33 & \begin{tabularx}{0.7\textwidth}{X} 基甸死後，以色列人又去隨從諸巴力而行淫，以巴力‧比利土為他們的神明。 \end{tabularx} \\ \\
[1ex]
\hline
\hline
\end{longtable}
$^{1}$各位電影姐妹大家好.
生命樹宣教網絡又和大家讀聖經了.
聖經人,聖經事,人物.
今次繼續是機電.
事情就是機電退休了.
那麼就看看上帝是如何祝福他.
以及今天的主題.
忘恩負義的以色列人.
以色列人如何忘恩負義呢?.
本來我想叫以色列人忘恩負義.
不過我想將忘恩負義擺前一點.
強調一點.
就是忘恩負義的以色列人.
《士師記八章》二十八節.
先講機電幫助天下太平.
以及神祝福他.
這樣米電人被以色列人制服了.
不敢再抬頭.
抬頭都不敢.
機電還在的日子.
國中太平四十年.
就是說機電在的時候.
米電人四十年都不敢再來.
而其他的民族也都不敢搞事.
因為機電竟然將米電人.
是一個很多的民族的組合.
而且遙勇善戰.
在沙漠曠野地方的集合.
再加上約旦河東的很多民族.
東方民族一起打.
都被他三百人戰勝了.
殺了很多個米電的王.
大家提起都打了冷顫.
平安了四十年.
很厲害.
大家以色列人都很歡迎他.
就說機電不如你做我們的首領.
你管理我們.
你的兒子都繼續管理我們.
但是記不記得機電說不用了.

$^{41}$不好我不管你.
惠友耶和華管理你們.
不過就拿一些金來做一個金的伊弗德.
伊弗德本來是布的.
用藍線,紫線,金線,紅線.
做一個背心形.
後面是包著的.
然後這裡有一條帶.
上面掛一個胸牌.
胸牌是金做的.
有十二塊寶石.
象徵以色列十二個支派.
大祭司胸懷十二支派在胸口.
那件叫伊弗德就好像是一件背心.
就是掛著胸牌.
以色列十二支派的胸牌.
用布做的.
用了二十公斤黃金做什麼呢?.
應該後面再想想.
今天跟大家應該說.
我推測機電做這個伊弗德不是一塊布.
應該是一個金的伊弗德.
吊在這裡.
大個的.
嘩.
後來結果就搞到以色列人用了.
他們在迦南地.
諸巴力.
很多個巴力.
他們拜神的方式.
都是跪拜.
以及會跟廟祭發生關係.
神是很不開心.
這件事就反而成為機電的弱點.
以及他被指責.
網羅.
就是說他的缺失.
是他人生的缺點.
說回機電這個人.
那些人很快又走去拜了.

$^{81}$那件伊弗德偶像化了.
這件事機電都有少少責任的.
起碼他都沒有阻止他.
這樣就真的不好.
可能阻止不來.
好了.
太平機電還在.
米電人不敢動.
接著機電退休了.
就回到自己家裡.
機電的名字叫耶路巴力.
讓巴力跟他爭辯.
因為他拆了巴力的壇.
讓他跟他爭辯.
回到自己家裡.
就是說退休了.
神很祝福機電.
機電有70個親生的兒子.
很多妻.
咦?.
那不是一夫一妻的嗎?.
上帝伊甸園一夫一妻.
是的.
不過呢.
其實他們傳統一直下來都不是.
你看亞伯拉罕都不是.
亞伯拉罕給次拜下甲他.
神都給的.
當時還沒去到一夫一妻.
不過都應該傾向是這樣的.
不過那些人很厲害.
有什麼.
於是很多要求摩西.
給他們娶多一個.
或者休妻之類.
好了.
這裡下去.
祂有70個兒子.
我的名字都記不清.
我的教會少少人.

$^{121}$幾十人.
有時候名字都記不清.
老了記不清楚.
很多的妻子.
還有妾.
這些70個兒子.
很多妻子.
其實象徵著上帝很大大大地祝福他們.
子女多.
兒子.
女還沒寫下去.
哇.
百多人.
子女百多人.
都生了一個設在士劍.
士劍就是當時的.
以色列的首都.
還沒有去耶路撒冷.
還沒有耶路撒冷.
就在士劍.
會幕都是放在士劍的.
在那裡有一個妾.
因為她經常去士劍.
有一個妾為她生了一個兒子.
基澱為這個在士劍的兒子.
改了一個名叫阿比米娜.
好了.
約阿斯的兒子基澱年紀老邁而死.
這裡想第一點.
這個藍色螢光.
我想說.
神是大大地祝福基澱的.
第一就是子女很多.
太太多.
第二就是年紀老邁壽終正寢而死.
年紀老邁才死.
厲害.
而且還在哪裡呢.
葬在阿比爾族的俄弗拉.
即是她自己出生的城.

$^{161}$而且是自己族的城裡面.
即是說葬在自己祖宗的墳墓.
也葬在她父親約阿斯的墓墓裡.
為什麼自己沒有一個墓.
而葬在父親那裡.
因為你看到舊約.
他們經常說.
瑞鳥歸到祖宗那裡與列祖同睡.
是什麼意思呢.
就是他們會將屍體放一年.
然後用香料包住.
沉香抹藥.
然後讓它在墳墓裡面石.
放在那裡露空氣.
一年.
那個肉就不見了.
剩下的骨頭.
就將它放在盒裡面.
將骨頭放進去.
在山洞裡面的石再鑿深一點.
放一個骨盒進去.
然後兒子死了.
又這樣處理了他的肉之後.
又將他的骨汁弄乾淨.
放在骨盒裡面.
再鑿闊一點的小山洞.
將它放在爸爸的骨盒旁邊.
鑿在父親阿司的墳墓裡.
兄弟也是這樣.
所以墳墓會越來越大.
有些孫子也是放在那裡.
基澱死後.
以色列人又去隨從.
朱巴力行邪淫.
又去.
史經是這樣記載的.
基澱死後以色列人又冇往管.
這就是時事記經常發生的惡性循環.
就是以色列人去拜其他的神.
於是神用外族人成為邊打皮邊打他.

$^{201}$很慘啊很痛啊神啊.
為什麼你不顧念我們啊.
你答應帶我們的.
由埃及帶到現在.
神就好了.
見你這麼慘.
好了找個時事幫你走.
時事又平定了之後.
一會兒他又再拜.
當然時事記有一個說話.
那時以色列中沒有王國人任意而行.
就停下來走這個循環.
以色列人又忘恩負義.
又去隨從.
朱巴力.
什麼叫做朱巴力.
就是很多個巴力.
迦南地所有神都叫巴力.
就是Bael.
Bael有很多個.
這次是拜誰呢.
就是巴力比利圖.
Bael Berith.
Bael Berith 就是主.
Berith 就是誓約.
誓約之主就是你向他發誓一定靈的.
於是他們這樣想.
而且都是敬拜的時候.
會跟廟祭行淫.
這個神是很不歡喜的.
就拜巴力比利圖.
Berith 誓約之主.
做他們的神.
上帝又很氣憤很生氣了.
好了以色列人留意.
第三樣經文這麼說.
以色列人是不紀念耶和華他們的神.
也都不照著.
我們特別是螢光燈給大家看.
就是不紀念和不照著.

$^{241}$不記起他們的神.
就是曾經救他們脫離四維受敵的手.
忘記了.
不理.
不想那麼多.
很多恩典神很多恩典給我們嗎.
沒見到.
我這族沒經過.
什麼的.
不記得.
不紀念.
其實應該紀念就是由很遠古.
就是由埃及帶他們出來的.
不紀念.
也都不照著耶路巴黎基奠.
當時向他們所施的恩惠.
來厚待他們的家.
有兩件事的.
第一不紀念他們的神.
為他們做過什麼工作.
帶他們脫離受敵的手.
第二.
不照著耶路巴黎基奠.
向他們所施的恩惠.
因為基奠是救過他們.
他們當時也說.
你做我們的管理人.
你跟基奠說.
而且我們會厚厚地待你的家.
你放心.
你的後人我們會照顧的.
你退休就可以了.
既然你要退休.
好,沒問題.
他們是不照著.
他們曾經應許過會厚待基奠的家.
忘恩負義.
餓了,窮了.
另外就是不紀念神.
這是第三點.

$^{281}$他們是忘恩負義的以色列.
你知道忘恩負義.
對神忘恩負義.
對人忘恩負義.
神是會動手的.
你要知道我們的神.
是又真又活.
祂會動手的.
一定要小心.
記得.
除非你認為.
你的神是巴力.
由你召祂.
我們的神是耶和華上主.
是祂會召你的.
祂會主持的.
祂會出手的.
你要記得.
忘恩負義對神對人.
神是很不歡喜的.
好了.
我今天要跟大家講到這裡.
希望你會記得.
不要做一個忘恩負義的以色列人.
我們一起禱告吧.
天護上帝.
我們在這個地方是要求你幫助我們.
很多時候我們只看眼前.
後面的東西忘記了.
阿神你是帶領我們經過了這麼多的歷練.
這麼多苦難.
你都與我們同行.
幫助了我們.
開解了我們.
救了我們.
又是你差遣了人.
好像你的使者那樣.
幫助了我們脫離很多危險.
祂是身先士卒.
甚至為我們去捨棄生命.

$^{321}$去作戰.
去奮戰.
照顧.
但是我們又忘記了祂當日.
我們不懂得回報.
不懂得回報你的恩典.
不懂得回報人給我們那種的幫助.
主,這就是以色列人當日在這段經文裡面看到的.
所以我叫祂做忘恩負義的以色列人.
主幫助我們.
不做忘恩負義的人.
紀念你自己歷練給我們所施的大恩大德.
紀念到人曾經祝福了我們.
我們都可以回應.
可以回報.
求主幫助我們.
紀念你.
多謝你自己的恩情.
我們感恩祈禱.
奉耶穌基督名求.
阿們.
\newpage



\section{路加福音 15:11-12}
\label{sec:032344}
\textbf{一位父親的心}~~~~李慧珊傳道
\newline
\newline
連結: \href{http://www.church.com.hk/acms/content.asp?site=cdc\&op=show\&type=product\&code=032344\&layout=sermon}{\texttt{032344}} ~~~~ 語音日期: 2026-06-21
\newline
\newline
\hyperref[sec:032343]{\small{< < < PREV SERMON < < <}}
~
\hyperref[sec:index_chronic]{\small{[返順時目]}}
~
\hyperref[sec:032345]{\small{> > > NEXT SERMON > > >}}
\newline
\newline
路加福音 15:11-12
\newline
\begin{longtable}{cl}
\hline
\hline
章節 & 經文 (和合本修訂版)\\
\hline
15:11 & \begin{tabularx}{0.7\textwidth}{X} 耶穌又說：「一個人有兩個兒子。 \end{tabularx} \\ \\ \relax
15:12 & \begin{tabularx}{0.7\textwidth}{X} 小兒子對父親說：『父親，請你把我應得的家業分給我。』他父親就把財產分給他們。 \end{tabularx} \\ \\
[1ex]
\hline
\hline
\end{longtable}
$^{1}$弟兄姊妹,父親節快樂,父親節蒙科.
我們在神的裡面很感恩.
讓我們經歷到地上有我們的父親.
亦有天上的父親.
今日和大家分享的主題.
我選擇了「一位父親的心」.
剛才大家讀的經文很短.
稍後我準備和大家分享耶穌所說的故事.
是一個序幕.
在說這個故事之前.
想看看你們心目中的父親.
在你們從小到大的時候.
你的父親是怎樣的父親呢?.
我又… 請看一看.
今日是父親節,所以想起父親.
有些父親… 看看第一幅圖.
很能幹的,很懂得維修.
你的父親從你小時候到現在.
你的印象,你的父親是怎樣的?.
是不是這樣?.
我們在當中也有爸爸,舉手,誰是爸爸?.
你呢?你覺得自己也是在家裡做維修員的,舉手.
嘩!很多,換燈泡.
有時電掣壞了.
洗衣機的按鈕不懂得按.
看一看… 其實打開關掉就行了.
我們很經典的父親形象.
很懂得工作,很懂得維修.
有些爸爸很懂得烹飪.
誰是我爸爸是很懂得烹飪的?.
我爸爸是比我媽媽的烹飪更好吃.
很經典,爸爸是一個很好的廚師.
現在的年代會多一點.
爸爸會和小朋友看故事,說故事.
在我的年代,沒有這些親子活動.
但是爸爸也會說家族的歷史.
說他小時候的辛酸史.
大家在爸爸身上有沒有聽過這些故事?.
或者你做爸爸時也會和你的子女說.
爸爸也很懂得教導.

$^{41}$教你去看這個世界.
叫你做人最重要是有禮貌.
大家也會經歷到爸爸是一個很好的導師.
教你欣賞這個世界和看事物.
爸爸也很友善.
誰和爸爸很友善?.
少一點,有啊.
爸爸除了會教導.
也願意和我們做朋友,去玩.
有些親子的時間.
還有些爸爸好像超級英雄.
我爸爸可以有一個技能.
徒手捉蟑螂.
我覺得很厲害.
不知道爸爸有沒有一些特別的技能.
你覺得他很厲害嗎?.
我們父親的印象.
也有很多很美好的印象,經驗.
我們也有一些爸爸.
甚至我自己的爸爸.
也會有人的缺陷,人的軟弱.
很需要做子女的我們來疼愛他們,關懷他們.
也會有這樣的情況.
甚至在地上.
可能爸爸要工作,賺錢.
有不同的原因.
變成有缺席的父親.
但今天他說的比喻.
是耶穌所說的.
他說的這個父親.
他是沒有缺席過的.
但在說這個比喻之前.
原來耶穌已經說了另外兩個比喻.
在《路加福音》十五章的時候.
耶穌說了兩個比喻.
就是說失楊和失淺的比喻.
當中帶出很重要的訊息.
很簡單和大家說一下.
失楊.
有沒有記得耶穌說過的比喻.

$^{81}$若果你是做牧羊人,養羊的.
你有一百隻羊.
但走失了一隻羊.
因為這也是你的財產.
暗地裡.
所以你怎樣都要放下那九十九隻羊.
去找那一隻羊.
因為是很寶貴,很矜貴.
對你的價值連城.
那一隻羊.
所以找到的時候.
很歡天喜地.
然後就把這隻羊放在你的肩膀上.
然後放在你的羊圈裡面.
帶出了第十五章七節.
「一個罪人會改.
在天上也要這樣為他歡喜」.
這個就是失楊的比喻.
失淺的比喻就是說.
那些婦女.
我們在座當中有很多婦女.
精打細算的.
婦女要打理頭家,買菜.
都要省著錢.
如果那些婦女.
她有十塊錢.
當作是十塊錢.
一個個十塊錢放在那裡.
你只掉了一塊錢.
那怎麼辦?.
一塊錢掉在地上.
很倒霉.
你都要點燈.
那時候是油燈.
要點燈.
怎樣都要掃到家裡.
怎樣都要找回那一塊錢.
找到之後就很開心.
然後跟朋友說.
「我那天找回那一塊錢」.

$^{121}$這個是重視.
第十五章十節.
他帶出去說.
「一個罪人會改.
神的使者也是這樣為他歡喜」.
接下來的比喻.
就是到現在準備說故事.
就是耶穌接著這兩個比喻.
那個心情.
來到那種不見了.
而找回那種歡喜快樂.
是說生命的悔改.
那種迴轉.
在這個比喻裡.
耶穌帶出一個家庭.
這個家庭很簡單.
祂說了三個人物.
但這三個人物.
有沒有名字呢?.
有它的角色.
就是有爸爸.
爸爸有兩個兒子.
一個是大兒子和一個是小兒子.
在聖經裡說.
父親,大兒子,小兒子.
就是這麼簡單.
但裡面看到他們當中的關係.
看到他們雙方的互動.
來帶出耶穌想說的話.
讓我們一起去祈禱.
主耶穌,多謝你.
你去愛我們.
你去寶貴我們每一個生命.
也讓我們知道.
有一位天上的父親.
有一位父親很愛我們.
並且你今天要帶我們.
去進到這個父親的心裡面.
這位父親的心.
你顯明給我們每一個人知道.

$^{161}$無論我今天在台上說的.
和在台下的每一個眾弟兄姊妹.
都能夠打開我們的心.
能夠和你心對心.
來領受今天的恩典.
多謝你.
願聖靈親自帶動我們每一個.
讓我們每一個都蒙你賜福.
奉主耶穌基督的聖名自祈求.
阿們.
好,說這個故事.
這個比喻是說那兩個兒子.
小兒子叫小兒子.
小兒子跟父親說.
他說父親將你的家業分了.
即是分身家.
但分身家原來在當時來說.
猶太人是不可能這樣做的.
因為他們有些經典的教導.
其中一個經典叫做變西拉自訓.
是他們那些智慧的緣分.
教導說只要你再生.
就不要給任何的權給任何人.
即是你的權力不可以給任何人.
無論是你的妻兒.
你的朋友,兄弟都不要給.
你的財產也不可以這樣分給任何人.
以免你會後悔.
要你的心意要他取回.
為何會這樣說呢.
可以想到.
因為若果你太早分了你的家業.
分了甚麼人.
是在說關係的.
那時候你跟他關係很好.
就分給你.
但如果有些事發生了.
你對我突然變了.
我跟你的關係好像有些事發生.
我真的後悔了.

$^{201}$我想期待你沒做.
有些事發生.
人之常情.
這也是智慧.
若果分了給別人.
然後有一天.
有甚麼事發生.
突然有些事發生.
你需要用錢.
需要用你本身有的物業財物的時候.
分了不夠.
所以這個在猶太人的傳統裡.
是不可以做到的.
但耶穌竟然用這件事.
來帶出一些訊息.
這個父親竟然怎樣呢.
他父親就把財產分給他們.
即是說他父親答應了這個小兒子.
很奇妙.
這個小兒子所說的是非傳統的事.
而這個小兒子.
你猜爸爸是否知道這個小兒子的心呢.
一定會認識這個兒子.
這個兒子吉起條尾.
都知道你想怎樣.
是不是.
助助各位爸爸.
會不會的.
你都會了解你兒子所想的.
他做這件事.
究竟是想怎樣呢.
接著.
我相信這個爸爸.
都心中有數.
這個兒子怎樣.
後來我們繼續看.
他說隔了沒有多少天.
這個兒子便執包袱.
遠走他方.
經文十三節告訴我們.

$^{241}$他怎樣.
任意放蕩.
還有浪費錢財.
這個兒子其實拿了一些錢.
去花天酒地.
去頌呈他的慾望.
玩樂享樂.
我相信他當時.
都會結識到一群人.
那群人很喜歡和他在一起.
為甚麼.
又可以吃又可以玩.
認識這個朋友.
真是好得無比.
吃東西又不用錢.
玩東西又娛樂.
又是他出了很多錢.
所以這個小兒子.
應該都頗受歡迎的.
在當地.
但是錢都拿走了.
那時他分了多少錢.
就是分了爸爸三分之一的財產.
這個爸爸.
想像是很有錢的.
一個當地的富翁的時候.
他拿了三分之一.
為甚麼拿了三分之一.
不是分大兒子和小兒子的.
因為原來當時分財產.
大兒子是分兩份的.
小兒子是分一份.
大家現在知道了.
大兒子那兩份.
即是三份來分兩份.
分兩份.
兩份一份.
所以他三分之一的財產.
應該都是很豐富.
可以夠他過一輩子.

$^{281}$基本上.
正正常常生活.
但是他選擇.
浪費所有他爸爸給他的東西.
他不珍惜.
他沒有計劃.
他沒有打算.
就是這樣.
十四節說他.
耗盡一切.
還要那裡有饑荒.
沒有錢.
更加有饑荒.
整個經濟環境都非常差的時候.
他非常窮.
窮到連買東西吃的錢都沒有.
接著很不容易才找到工作.
經濟差.
環境差.
都影響生活.
影響找工作.
這個有時我們打工仔.
有時見到的時候.
或者覺得工作上都不太順利.
都很擔心那份工作.
都不知道穩不穩陣.
有份工作都很感恩.
而當時候這個小兒子.
他只可以找到一份工作.
那份工作是放豬.
聖經說.
他是去養豬.
而當時候養豬來說.
在他們猶太人文化.
是極度之低賤的工作.
因為豬.
猶太人會怎樣?.
覺得是不潔的.
因為豬的腳是分開的.
分體.

$^{321}$但是他又不可以去反戳.
聖經會跟我們這樣說.
你還沒記嗎?.
新命紀.
所以這個豬是不潔的.
而他只有去放豬.
去養豬.
去度活.
並且是不夠食物.
他當時.
「恨不得拿豬所吃的豆鴿充饑.
也沒有人給他什麼吃的」.
什麼叫豆鴿呢?.
就是這種在地中海.
有的植物叫「割豆樹」.
這個割豆樹有些豆.
它就是中間的種子.
那些豆.
去成為一些豬吃的東西.
看到這個豆.
都黑黑的.
不知道好不好吃.
可能人都不會想吃它.
都不是很青綠色.
其實它沒有東西吃.
到現在我們說吃什麼?.
吃豬哨.
沒有東西吃到.
想吃豬哨.
是多麼淒涼.
悲涼.
絕境.
窮途末路.
人的絕境.
接著他就想到.
爸爸.
因為他的爸爸.
是豐富的.
有足夠不同的資源.
所以他就醒悟過來了.

$^{361}$「我爸爸有很多工人」.
「有很多糧食」.
他想著想著.
「很多精美的食物」.
「我用不著在這裡吃豬哨」.
「我用不著這樣餓死嗎?」.
接著他就想著.
他爸爸知道他現在的情況.
真是敗家子.
一定不原諒他.
作為爸爸的你.
原諒他嗎?.
可能都原諒.
難一點.
心裡要罵他一頓.
「你這個壞蛋」.
他的憤怒.
他都想像到.
他爸爸一定會生氣.
生氣到發瘋.
所以他就預演一個小劇場給自己.
他說什麼?.
「我要起身」.
去到爸爸那裡.
接著就跟他說.
想一想對白怎麼說.
「父親 我得罪了天又得罪了你」.
「從今以後我不配稱為你的兒子」.
「就把我當做一個工人吧」.
這個對白好不好?.
你猜這個兒子.
是否真心誠意悔改呢?.
覺得自己做得很差呢?.
在我自己覺得裡面.
有可能是 有可能不是.
不是因為他的處境.
和他的品性的問題.
因為他不喜歡這個家.
他想自己自由地亂用他擁有的資源.
更加開心地放蕩很多人的自由.

$^{401}$做什麼都可以.
去到爸爸可能管理得很乾淨.
但現在因為他肚子餓.
沒東西吃.
又饑荒.
又沒東西做.
爸爸那裡有東西吃.
做工人也好.
所以我相信因為他肚子的問題.
他肚子餓 要找東西吃.
所以要找爸爸.
是有他自己生命上.
他身體上的那份需要.
那份想望.
那份渴求.
才去找回這個爸爸.
有沒有錯?.
都知道有個爸爸去靠.
都知道有個爸爸可以依傍.
去拉著他.
這個爸爸你猜知不知道.
這個爸爸他沒有計較.
我們繼續看下去.
他就真的去找他爸爸.
但原來這個爸爸.
這個父親.
他日思夜想這個兒子.
不見了他.
很掛念他.
也知道這個兒子的品性.
沒什麼好東西.
他去到他的大宅.
很空曠的地方.
以前的豪宅或者大富人家.
不是像現在那樣住半山.
某些地方住半山是窮人才住半山.
但香港很多時候半山區.
偏遠一點.
就是富有一點的地方的感覺.
但當時是平地.

$^{441}$他可能在村落裡.
某處是他的家.
但這個爸爸選擇走出屋外.
在一個很空曠的平原來仰望.
他就是看到遠處.
遠處已經看到他的兒子.
這個兒子當時是怎樣的呢?.
都可以想像一個養豬的工人.
沒東西吃.
什麼家當都沒有的時候.
他所穿的是什麼衣服呢?.
他當時的身世.
可能離開家的時候.
會不會有些資源.
讓他不用走得那麼辛苦.
可能坐小鑼.
可能小鑼也不會讓他坐.
他走回來.
徒步的回家.
那種的落寞.
那種的爛身爛勢.
但爸爸認得他.
當然他的兒子.
發了灰都認得他.
遠遠看到他的兒子.
跟著很遠.
然後這個爸爸.
聖經說什麼.
二十節一起讀.
就動了慈心.
這個慈心.
除了是kindness.
那種的恩慈.
是一種compassion.
那種的憐憫.
遠處看到他的兒子弄成這樣.
回來了.
他的心很受觸動.
因為太想念這個兒子.
那份憐愛出來.

$^{481}$他的心帶動他的腳.
然後跑去.
並且擁抱他的兒子.
他當時的兒子.
會不會身碎身亢.
骯髒骯髒.
哭哭啼啼.
不斷連連親他.
不斷親他.
不斷捶捶捶.
捶他.
那種的愛.
那種的擁抱.
是多麼的大.
這個兒子.
很記得.
他來要講對白.
這個兒子跟他的父親說.
寓言的對白.
他說.
父親.
我得罪了天.
又得罪了你.
從今以後.
我不配稱為你的兒子.
好 講到這裡.
不用說下去了.
因為父親.
他立刻叫工人來.
人來人來.
你快點把最好的袍.
給我兒子穿.
把你的戒指.
戴在他手上.
把最好的鞋穿在他腳上.
把一切最好的.
都放在他身上.
換掉他的骯髒骯髒的衣服.
我要證明.
這是我的兒子.

$^{521}$這個兒子.
好像死了回生一樣.
所以我們一起來開派對.
我們快點幫我們.
在那些肥牛毒.
很久沒有宰了.
宰一隻肥牛毒來.
我們一起高興歡喜快樂.
一起去慶祝.
Let's party.
我們一起去開派對.
這個就是父親.
這個父親.
在等這個兒子.
很高興歡喜.
他沒有計較這個兒子.
是什麼原因來到這裡.
這個是他的兒子.
這個兒子.
也不需要說要當他的傭人.
因為他是爸爸的兒子.
你不配就知道.
你是不配.
父親的愛.
很大的淋到這個兒子身上.
是傾注下來.
我相信這個兒子.
他經歷這份擁抱的時候.
會很不同.
這個兒子不認識他爸爸.
這個兒子以為他爸爸一定氣到擦天.
甚至不原諒他.
和他脫離父子關係.
但這個爸爸不是這樣的.
不是他所想的爸爸.
不是一般我們想的爸爸.
因為這個是天父的心.
耶穌所說的.
是天父的心那份愛.
接下來就看大兒子.

$^{561}$大兒子拿了多少身家.
三分之二.
即是說他現在有的身家.
全部是他的.
很正常.
但他很勤力工作.
他沒有吩咐傭人.
去落田.
他自己落田去工作.
然後他回家.
沒有多遠.
他聽到奏樂跳舞的聲音.
很遠距離都聽到.
跳舞聽到嗎?.
說笑而已.
可能是那個音樂就是跳舞.
很歡樂的聲音.
整體形容都很貼地.
因為他那時候很歡喜跳舞.
奏樂跳舞的聲音.
然後他就叫了一個童僕來.
見到有一個比較小的童僕.
「什麼事啊?什麼事啊?」.
這麼高興歡喜的奏樂聲音.
然後他說「你的弟弟回來了」.
「你爸爸說因為他無災無病回來」.
「所以幫他宰了一隻肥牛」.
「和他慶祝現在開派對」.
然後這個大兒子很生氣.
怒火就從心裡湧出來頭頂.
非常生氣.
不肯進家裡.
然後傭人見到他這麼生氣.
像殺人一樣就通知他的主人.
他爸爸出來勸他.
然後他和他爸爸說了些什麼呢?.
他就說.
他和他爸爸說.
「你看看,我服侍了你這麼多年」.
「從來都沒有違背過你不聽話的」.

$^{601}$「我一直很服從你的命令的」.
「但你連一隻小山羊都沒有去宰給我」.
「叫我和我的朋友去歡喜開派對」.
「合不合理啊?」.
但是他說.
「你這個兒子」.
「這個兒子是和昌記去吃了你的家產」.
「他一回來你就幫他宰肥牛」.
原來當時這個大兒子.
他怎樣看自己的身份呢?.
原來他是做傭人的.
他是家裡其中一個傭人.
幫他爸爸服侍.
原來服侍希臘民的原民.
是有奴隸和服從的意思.
這個大兒子.
他明白父親的心嗎?.
他看自己.
他是傭人.
一直以來這麼辛苦去服侍這個爸爸.
得到些什麼?.
「神你怎樣對我?」.
「我服侍了你這麼多年」.
有沒有可能會這樣對神說話呢?.
這個爸爸怎樣說呢?.
「兒啊」.
「你常和我同在」.
「我所有的都是你的」.
因為這個兒子想不通.
爸爸和他在一起了.
他沒有享受過爸爸和他在一起的恩典.
和美好.
和爸爸的那種關係.
那種親近.
在一起的那種親近.
原來這個兒子沒有去經歷.
他的關係.
始於是神僕的關係.
主人和僕人的關係.
爸爸對他說.

$^{641}$「我常和你同在」.
「我所有的一切都是你的」.
「你知不知道?」.
爸爸說.
「可視這個弟弟死而復活」.
「失而又得」.
「所以我們要一起歡喜快樂」.
「和他慶祝」.
這個兒子是不知道.
爸爸的心.
其實是透過這個小兒子.
很清楚地,很具體地去彰顯出來.
就是那份擁抱.
那份擁抱.
在原文是說「樂在」或「臨到」.
這是聖經所說.
原來在.
我們先看第二版.
因為在坊間.
我們也會找到一些資料.
就是關於心理學.
原來心理學上.
擁抱帶來很大的力量.
其實你想去擁抱.
找AI就會回答這些問題.
我嘗試去看一看.
原來擁抱不僅是身體的接觸.
更加是很強大的情感結連.
在心理上.
在醫學和心理上也證實了.
適當的擁抱能夠帶給身體的分泌.
有一種叫催產素和多巴胺.
能夠刺激到快樂的感覺.
並且有效降低壓力寒耳蒙.
穩定情緒和提升免疫力.
若果每天只擁抱數十秒.
就可以紓解壓力.
大家有沒有和身邊的人擁抱?.
和老公,好好榜樣.
大家有沒有和老公擁抱?.

$^{681}$數十秒.
大家壓力全消.
還有你的感情一定會繼續升溫.
還有有些焦慮感.
因為透過擁抱而釋放,解放了出來.
還有關係的促進.
調節情緒.
那種孤獨感,憤怒感.
都會減低,甚至是沒有的.
所以我鼓勵女組員和女組員之間.
多些擁抱.
不用有突發情況.
「主耶穌愛你,我也愛你」.
和身邊的人擁抱.
讓我們彼此心靈得到很大的祝福.
還有在第二次世界大戰.
有個心理學家叫做René Spitz.
去孤兒院觀察.
發現孤兒太多了.
在孤兒床上.
有些寶寶會照顧得到.
會抱得多一點.
有些很少會抱得那麼多.
原來都會影響他們的成長.
生長遲緩.
原來抱得少,碰得少的小朋友.
他的成長都很有障礙和問題.
就知道坊間是怎樣.
但是回到聖經.
擁抱就是「樂在臨到」.
一個原文.
耶穌說的比喻.
父親去找小兒子.
飛奔去擁抱那一刻.
原來原文是有「樂在臨到」.
而在《使徒行傳》十章提到一個事件.
就是彼得去探訪哥尼劉.
在這個事件中.
他們經歷到聖靈的臨在.
就是彼得宣講耶穌基督的十字架的福音時.

$^{721}$聖靈就在當中充滿覺人.
而在第十章的四十四節.
我們一起讀這兩節.
和十一章十五節一起讀.
一,二,三.
彼得還說這話的時候.
聖靈降在一起聽到的人身上.
我一開講.
聖靈便降在他們身上.
正像當初降在我們身上一樣.
剛才讀的時候有否留意.
有兩個字重覆了.
除了「聖靈」.
還有「降在」.
原來「聖靈降在」.
就是天父的擁抱.
這是我之前沒有想過的.
原來「聖靈臨在降在」.
就是天父藉著聖靈顯靈給我們聽.
在羅馬書五章五節說.
聖靈就是挑冠神的愛.
這是一個很重要的真理.
神的愛藉著聖靈去顯出.
顯在我們身上.
大家有沒有經歷過聖靈的挑冠.
那種喜樂.
那種能力.
那種釋放.
帶來的開心快樂.
有沒有?.
很記得很久之前的照片.
是我.
不用想.
我肥了三十磅.
我的朋友穿的T恤.
看了一年份的.
看不看到多少年.
眼力非常好.
看不看得到.
20.

$^{761}$看不看得到.
前排的看不看得到.
原來這麼小.
04.
2004年.
這裡是韓國的祈禱山.
有沒有人2004年去訪韓的.
有沒有?.
我們這場沒有.
我們長青有個領袖.
明珠和我同一年去訪韓.
我穿這件衣服應該都有.
當年訪韓.
我們有一個很渴求的心.
很想經歷耶穌.
很想經歷聖靈.
我們做了什麼很渴望.
現在我們鼓勵大家去禁食.
有沒有人在禁食祈禱?.
有四十天.
當時我們很乖.
因為大會叫我們四十天禁食祈禱.
我們就四十天禁食祈禱.
不過我每天禁一餐.
禁五餐很餓.
不過有同伴就會過得到.
因為有人陪.
當時有我在福音正主協會服侍.
說了機構的名字.
這個就是我正主的同工.
我們有幾個人一起去禁食祈禱.
然後去到祈禱山.
其實應該是這個旅程.
是有個要求再要禁食的.
除了去之前的四十天禁食一餐.
我們選一餐.
還有去到連續禁七餐.
即是兩天多一餐.
就是這樣.
當時為什麼這麼瘦.

$^{801}$可能是沒得吃.
說笑的 其實還是瘦的 年輕.
我想和大家分享一個很特別的經歷.
是我人生裡面.
第一次經歷耶穌基督十字架的愛.
而很真實地感受到這一份的愛.
當時的我.
2004年.
我1985年信主.
即是有幾多年了 當時的距離.
2004 減回1985.
幾多年?.
19年 差不多20年.
我不說幾多歲.
當時我相信耶穌在我的頭腦裡多一點.
當然有人生少少的經歷.
流淚的經歷在他之前也有.
不詳細說 不是個人見證的時間.
但我想說.
我是沒有因為耶穌釘十字架而流過一滴眼淚.
而在當時.
這份學無原來帶我.
不是要找神奇的經歷.
去找到耶穌釘十字架.
這份愛令我在那幾天裡不斷地哭.
後面是一個最重點的位置.
是甚麼來的 後面那東西.
其實它不是平面.
它是一個長方體.
門可以進去.
因為裡面有人.
所以我選擇了另一個.
開了門 進去後.
有一個地 像床墊一樣.
很容易稱為「滑」.
像膠面一樣 床墊軟軟的.
你可以坐進去 你可以跪 可以坐 或者可以睡.
站立不了 因為它只有這麼高.
差不多吧.
所以在裡面有燈 還有十字架.

$^{841}$還有少許通風 門有通風的位置.
就是這麼簡單的一個空間.
自己關在裡面 獨處.
就是大概六呎乘四呎.
你可以坐 可以睡 可以跪的地方.
有一個小的黃燈 是燈膽.
吊著在那裡.
當時我就進去了其中一格.
這個叫「祈禱洞」.
在裡面一跪在那裡祈禱的時候.
又想起十字架.
不斷地哭得很厲害.
那種是決堤式.
接著是肉身感到一份很大的重量.
有個頭 淋到我的肩膀.
整個身體很重.
好像看不到的物質.
有東西傾倒在我身上.
這是一種很神秘.
一種奧秘的屬靈經驗.
我是唯一一次去經驗過.
那一刻大到淋到我身上.
大到我承受不了.
「主啊!好多啊!好多啊!好多啊!好多啊!」.
而當時我也在尋求.
我會不會全世界去事奉.
有一個很卡我的位.
就是進去這個祈禱洞之前有一個聚會.
不知道大家有沒有聽過蓮嘉茵醫生.
當時去了現場分享.
一個青年的場.
我去參加了.
聽到她說了一句.
撫了我的心.
她說蓮嘉茵.
她怎樣回應神的呼召呢?.
她說「神啊!我命就有一條.
你隨便拿去吧」.
很震撼我.
我命就有一條.

$^{881}$我很害怕.
你拿去會怎樣?.
神你會怎樣用我的命?.
會死的.
我害怕死.
就是這個祈禱神的我.
真實臨在那一刻出來.
人還未回魂.
俗語所說.
還未定神氣.
聽到聖靈跟我說.
很清楚.
聽得很清楚.
特別清楚.
我的愛.
還夠不夠你為我去?.
很大的愛.
神的愛.
那一刻經歷到.
原來神的愛繼續帶我下山.
有血有肉地經歷生命.
不等於那一刻.
領受到這份超自然的經歷後.
我一帆風順.
反而更加落地.
上山是為了下山.
下山去經歷到人間疾苦.
不要再那麼天真.
很真實.
人是有血有肉.
有淚有痛.
之後的人生不斷去經歷這份擁抱.
不是一刻傾注.
一鋪過.
原來這份傾注.
在人生裡面不斷去經歷.
不斷去渴求神這份愛.
這份充滿.
就像小兒子一樣.
軟弱到一個地步.

$^{921}$呼求.
軟弱到一個地步.
學武.
軟弱到一個地步.
去抓住.
這份很真實的經歷.
神怎樣去赦免.
神的愛怎樣去擁抱.
去依持.
去釋放.
去更新.
去恢復.
去恢復我們.
我不知道大兒子小兒子.
你覺得自己像哪一個兒子.
我們明不明白天父的心.
我們不住要傾注.
去經歷這份擁抱.
我不知道夏令會.
是否一個神給你.
在你人生裡面.
一份很獨特.
一份屬靈的經驗去賜給你.
我不肯定.
或者是另一個.
或者是你自己的禱告裡面.
你自己和神的關係裡面.
神帶領你不斷去走.
一個新的旅程.
但是一個不可以缺少.
預備自己的心.
一份學武的心.
一份很期待.
天父更大的擁抱臨到你身上.
這一份的心.
Restore.
恢復.
再次起來.
更深的旅程.
更深的經歷.

$^{961}$打下我們的雙手.
去接耶穌的愛.
天父的心.
今天跟我們說.
很愛我們.
天父很愛你.
很愛我.
\newpage



\section{路加福音 11:9-14}
\label{sec:032345}
\textbf{懂愛敢愛的爸爸}~~~~王安洋牧師
\newline
\newline
連結: \href{http://www.church.com.hk/acms/content.asp?site=cdc\&op=show\&type=product\&code=032345\&layout=sermon}{\texttt{032345}} ~~~~ 語音日期: 2026-06-21
\newline
\newline
\hyperref[sec:032344]{\small{< < < PREV SERMON < < <}}
~
\hyperref[sec:index_chronic]{\small{[返順時目]}}
~
\hyperref[sec:032346]{\small{> > > NEXT SERMON > > >}}
\newline
\newline
路加福音 11:9-14
\newline
\begin{longtable}{cl}
\hline
\hline
章節 & 經文 (和合本修訂版)\\
\hline
11:9 & \begin{tabularx}{0.7\textwidth}{X} 我又告訴你們，祈求，就給你們；尋找，就找到；叩門，就給你們開門。 \end{tabularx} \\ \\ \relax
11:10 & \begin{tabularx}{0.7\textwidth}{X} 因為凡祈求的，就得著；尋找的，就找到；叩門的，就給他開門。 \end{tabularx} \\ \\ \relax
11:11 & \begin{tabularx}{0.7\textwidth}{X} 你們中間作父親的，誰有兒子求魚，反拿蛇當魚給他呢？ \end{tabularx} \\ \\ \relax
11:12 & \begin{tabularx}{0.7\textwidth}{X} 求雞蛋，反給他蠍子呢？ \end{tabularx} \\ \\ \relax
11:13 & \begin{tabularx}{0.7\textwidth}{X} 你們雖然不好，尚且知道拿好東西給兒女，何況天父，他豈不更要把聖靈賜給求他的人嗎？」 \end{tabularx} \\ \\ \relax
11:14 & \begin{tabularx}{0.7\textwidth}{X} 耶穌趕出一個使人成為啞巴的鬼，鬼出去了，啞巴就說出話來；眾人都很驚訝。 \end{tabularx} \\ \\
[1ex]
\hline
\hline
\end{longtable}
$^{1}$各位電影節目來賓,早晨.
聽到小朋友的歌聲很雀躍.
小朋友帶給我們無限的活力和動感.
讓他們的歌聲帶著祝福.
向在座每一位的父親加油,努力.
所以在電腦簡報也找了一個這樣的圖片.
一個爸爸牽著兩個兒女,一個背影.
相當溫馨,挺高的溫度.
今天和大家分享的題目是.
懂愛,敢愛的爸爸.
對於愛這個課題,不知道大家心裡怎麼想.
有些難度,但是也可以去互動.
兩段經文,最初這段是羅家夫音11章9至14節.
講到禱告的課題.
裡面提到祈求,尋找,求門.
也有一段是相類似關乎爸爸的一些比喻的特性.
所以羅家夫音15章11至32節.
一會兒我會和大家再繼續努力.
在這個比喻當中.
這位父親給我們有些什麼學習.
也有一些幫助.
對在座每一位尋找信仰的朋友.
成為你的一個鼓勵.
也祝福不怕多.
辛苦大家了.
作為家中的男士.
我們的奴婦都想家人個個開開心心,快快樂樂.
保持家裡的溫度和暖度.
也想一想父親節的來源.
其實多得母親節.
所以母親很偉大.
如果沒有他們就沒有父親.
1905年父親節就開始了.
五年之後大約1909年至1910年.
第一個父親節就在美國出現.
多謝這位道德夫人.
她的媽媽比較早逝.
所以她的爸爸是一個退伍的軍人.
像老虎一樣養大六個子女包括她.
所以她都很願意將爸爸這份敬意.

$^{41}$帶給地上的每一個家庭.
一起努力鼓舞每一個家庭的父親.
所以她在1909年就發動一些設立父親節的舉動.
1910年就正式設立了.
不過當時還不是全國或者一個法定的節日.
要去到1972年.
李克信總統就簽訂了那個法案.
就立定了六月的第三個主日就是父親節.
當然很多地方都有不同父親節慶祝的日子.
如果台灣的地方或者某些地方.
他叫做爸爸.
所以八月八號是父親節的日子.
所以各式各樣各處一個例.
亦都想表達對爸爸的敬意.
在在的父親.
你可能想一下你自己是屬於哪一類型的人.
想一下這張圖.
簡單說有五類.
但是你試一下找一下網上.
八樣九樣都有的.
簡單分類一下.
或者一些特性就是.
會不會有些爸爸都比較嚴肅.
比較霸氣.
一出來就大家要恭恭敬敬.
收聲呢.
我小時候爸爸有時候一出來我們就收聲了.
很恭敬的.
這一類的爸爸不是等於他不關心家裡.
他只是較少用言語去表達.
但是內心是很想家人個個都平平安安的.
另外一類就是愛理不理.
但是一個好像旁觀者這樣.
但是他事事都能夠細心觀察.
一到適合的時候.
一句適合的說話或者行動就起了果效.
所以有時候這些叫智者.
又或者一些暖男的爸爸.
比媽媽更厲害.
煮得兩手的.

$^{81}$所以那些小朋友或者家人就有一個飽福了.
或者一些爸爸很和子女打成一片的.
好玩得什麼都沒什麼所謂的.
子女很喜歡和他玩的.
或者一類爸爸很專心他的工作.
希望家庭有一個穩固的經濟基礎.
無論你是哪一類型的人.
或者你在想你爸爸屬於哪一類也好.
每一類型的爸爸都有時候有這種的表現.
一種叫做外顯式的父職.
或者叫做內隱式的父職.
外顯式就是用說話用行動去表達那份的愛.
會關心你會著重溝通.
會彼此互動.
又用他各樣可以用到的資源.
能夠帶給家人有那份的溫飽安全.
以致到有些時候如果子女需要教.
這些爸爸都會出聲說的.
但是有一類爸爸就是內隱一點.
就是內斂一點.
他用情緒去做一個回應的.
內心常常都想著子女.
其實都掛心每一個子女的.
好想每一個子女都好好.
其實也都很願意支持子女.
不過礙於他內斂.
他不是很能夠表現出來的.
隱藏不露.
但是無論外顯還是內顯.
無論是哪一類型的爸爸.
都是一樣的.
有哪一個爸爸不想家裡好.
我家裡就太太兩個女兒.
我常常都有這句話.
弟兄姊妹常常都學我的.
我家裡女人好我就好了.
不知道在座有沒有同樣類型的爸爸.
你的家庭背景都一樣.
你會不會共鳴這句話.
家人好我們就好了.

$^{121}$特別在這個時候.
我兩個女兒都已經大學畢業.
出來工作了.
其中一樣我好開心.
就是有家用收.
每個月都買一個日子.
我戶口就會多一些錢.
收到的時候.
心裡面都是值得的.
但是你要知道.
有時候你要在適當的時候.
講適當的東西是高難度動作的.
你不適當的時候.
講適當的東西又死的.
你不適當的時候.
講適當的東西也是困難.
適當的時候.
你不懂得講話也是困難.
所以我常常都會這樣想.
三個女人訓練我.
應該我的腦退化會比較遲一些出現.
因為常常都考驗我的IQ和EQ.
連政府部門都教我們.
做父親需要留心.
做父親不要只想著賺錢養家.
我小時候爸爸是做海員的.
他就是船上的大廚.
我小時候常常都想.
為什麼我爸爸會是這樣的呢.
久不久才見他一次.
短則幾個月.
長則兩年.
一次回來的時候.
心裡面五歲左右的時候.
心裡面不忿.
鬧彆扭.
爸爸媽媽為什麼要這樣.
我就心裡面很想知道.
誰搶了我爸爸.
就很想去我爸爸工作的地方.

$^{161}$去示威一下.
就鬧彆扭.
我爸爸說你不要去.
那艘船泊在大海那裡.
我受不了.
但你知道小朋友不聽.
只用眼淚回應他.
只要我爸爸搞我不行.
第二天一大早就把我帶去大船那裡.
大家知道要坐那些叫做博艇.
我們叫瀑布仔.
就瀑瀑瀑瀑瀑瀑瀑.
瀑一個多小時.
那一程是吐了幾次.
去到那裡都未至於.
那艘大船真的很大.
那些是大郵輪的船.
我看著我爸爸.
怎麼上去.
就走上去.
怎麼上去就走那條承梯.
中間那個洞是穿得下去的.
你知道我五歲有多大個人.
但走這裡樓底都高.
比這裡樓底高幾倍.
兩三層樓.
我都印象不記得走了多久.
終於走完.
但走到上去腳都軟了.
跟著我爸爸說.
爸爸跟我說.
你又說來.
就看你壞了.
但去完那次.
了解了爸爸工作的地方.
我就覺得原來爸爸很辛苦.
真的飄洋過海.
為了家裡有餐溫飽.
所以從此以後.
我就不夠膽再發爸爸媽媽脾氣.

$^{201}$而且很尊敬爸爸.
終於我爸爸.
順了主.
也安息主懷了.
也想到父親節.
一些特性.
鼓勵一下在座的父親.
我們都要著重那種.
跟子女彼此互動的關係.
特別那種陪伴.
那種留心子女的需要.
跟他們共同培養一些興趣.
能夠幫助他們.
在人生裡有一個美好的回憶.
也都需要耐性.
你知道有時候小朋友.
一時是天使.
一時是魔鬼.
總是摸我們的底線.
和挑戰我們的容忍程度.
但依著這份忍耐.
你要把握很適當的時間.
他需要你的時候.
你就在那裡.
世上很可惜.
有些叫做absent father.
曾經何時.
我都以為我爸爸absent了.
但原來present father.
是很重要的.
讓子女能夠有一個.
更好的環境.
孕育他們的生命.
能夠有更好的成長.
當然作為爸爸.
都要幫忙家裡的家務.
有時候都要欣賞太太.
那你試一下.
一個太太.
她要照顧家庭.

$^{241}$如果是在職的太太.
就更加厲害.
我太太這幾年.
就不需要再工作了.
因為她教了三十年書.
她說夠了.
到現在她是家庭主婦.
有時候都要幫幫忙.
就不是奉旨回到家裡.
斟茶遞水.
一起努力.
這樣的互動.
能夠讓這個家的溫度.
更加提升.
以下有一條片.
我在網上見過.
挺好.
這條片是得獎的.
但這條片不說話的.
大家一起欣賞一下.
(電車聲).
看到父愛.
這條片沒什麼好說的.
但你看到.
不知道大家心裡會不會觸動.
當然可能我們.
如果套落我們現實爸爸.
可能還有很多欠缺.
但我想每一個爸爸.
心裡都很想這樣.
遮風擋雨.
最近香港天氣都很變幻.
莫測突然間黑雨.
打風打雨的時候.
這件大衣.
扯緊了.
濕了.
很努力.
陽光普照之後.
放鬆一下.

$^{281}$拿出來的.
有太太的衣服.
有三個兒女的衣服.
維保他們的安全.
很窩心.
作為父親的我.
有時候看這條片.
都很得提醒.
其實爸爸內心.
都是想這樣.
愛這個字.
如果用簡體.
是明白不到裡面的精髓.
是繁體才明白到.
愛的中間是什麼.
心.
你拿走的心就是仇.
愛是要用心去領受.
去感受.
愛是一種觸動.
愛不是一種計算.
愛不是一種.
你可以計算出來的東西.
愛不是用你的腦袋去接收.
愛不是你用分析去處理.
不是用分數去打.
這個究竟愛我多少.
愛其實是要用你的心去領受.
你才知道愛為何物.
剛剛你們程序表.
這篇文章都是講述.
李成林牧師.
是一位我非常尊敬的牧者.
他的安息禮上.
三位兒女對這位爸爸.
有一些的悼詞表達.
其中一封是來自阿Mo.
啟賢弟兄.
他寫給爸爸一封信.
三十一歲半的他.

$^{321}$經過一千三百六十多日的難處.
但這位爸爸的陪伴.
無聲但默默的支持.
那種無微不至的付出.
那種真誠的表達.
用這個兒子.
曾經兒子在最困難的時間.
發爸爸脾氣.
很不容易.
病的人很容易.
而且發當然是最近那個發下去.
但爸爸沒有反應.
沒有回應.
也沒有罵他.
只是默默地用他.
強有力溫暖的雙手.
拍一拍他的兒子啟賢.
你知道男人不喜歡說話.
有時候拍幾下.
就收攏了.
所以啟賢這位兒子.
他馬上跟爸爸說對不起.
在啟賢的文章裡提及.
爸爸給了他一張信.
就是用廚房的手紙寫下.
字體也有些震動.
寫下啟賢加油.
這個兒子體會爸爸.
聖林是豐盛的森林.
是講到他那份對兒子.
那種不離不棄的愛.
那份堅持勉勵這個兒子.
要繼續努力.
所以他答應爸爸.
無論多艱難.
他都願意繼續努力.
剛才跟大家說的.
路加福音是一章七字.
九至十三節的經文.
是講到一個祈求.

$^{361}$經文這樣說.
我又告訴你們.
我是耶穌.
告訴那些人聽.
你們祈求就給你們尋找.
尋見叩門就給你們開門.
一位犯祈求的就得罪.
尋找的就尋見.
叩門的就給你們開門.
你們中間作父親的.
誰有兒子求餅.
反給他石頭.
求魚反拿蛇.
當魚給他.
求雞蛋反給他蠍子.
你們雖然不好.
尚且知道拿好東西給兒女.
何況天父豈不更將聖靈給.
求他的人嗎.
這裡說的是一個祈求.
一個求.
一個訴求.
人什麼時候去求.
就是當你知道自己有需要.
當你知道自己沒辦法.
在其他方面.
或者自己搞得定.
你就會求.
所以宗教有句說話.
在家靠父母.
出外就靠朋友.
所以在家就會問父母.
我小女兒當她入大學的時候.
不知為什麼.
她又找了一個門路.
她說爸爸.
我有機會過兩年西班牙讀書.
她真的問了我一句.
那你支不支持.
其實大家明白.

$^{401}$支不支持.
就是付不付錢去花.
我說好啊.
對你的人生有一個好的發展.
那就去吧.
作為爸爸每個人都是這樣的.
你求你生活所需.
可以有能力做得到的.
那就給吧.
這個聖經所說.
是一個生活所需.
這個聖經不是說求超級豐裕.
超級的成就.
發財發達.
不是這些.
是說人生的基本需要.
所以從這個角度.
Marslow在1943年已經說了.
人的需要的五個層面.
當然這個分析不錯.
不過這個分析不是夠詳盡.
人的宿命方面這裡沒有太多說.
雖然有安全感.
有些成就感這些事.
但是靈裡面.
與神的關係這裡沒有提.
但更重要的是與神的關係.
所以這個金字塔.
從底層開始.
就是一些生理的需要.
從一個安全的需要.
從一個人際互動的需要.
從一個被尊重被愛的需要.
到自己能夠實現理想的需要.
都是一個很好的轉識.
和一個指向.
這個基礎.
也都剛才聖經所說.
求的那些東西.
就是餅,魚,雞蛋.

$^{441}$是你日常所需.
父親雖然都有限.
父親都未必知道子女所求.
不過父親都盡力.
將這些他們所需的.
為他們預備.
父親可能自己都有限制.
自己都不知道自己需要什麼.
但是子女當要求的時間.
提出他們的需求的時候.
父親都會繼續努力.
特別這裡映照了一個.
作為天父上帝.
他將最好.
人最需要的就是聖靈.
讓人在生命裡面得著.
為什麼聖靈不是其他呢?.
突然間層次.
由生活所需.
去到一個屬靈的層面.
這麼有趣.
那麼聖靈是什麼呢?.
等一下我們可以透過.
《路加福音》十五章.
另外一個家庭的比喻.
我們可以再了解更多.
如果從這樣的訴求.
這樣的互動裡面.
讓我們也能夠留心.
這個家庭.
是聖經《路加福音》.
跟著十五章說的一個比喻.
是耶穌特別說給那些.
宗教領袖和對律法很有研究的人.
那些人就在指責耶穌.
為什麼你要黏著那些罪人.
尤其是那些稅利.
幫羅馬政權收款項的.
也都多收我們同胞.
鍾寶思郎.

$^{481}$那些是壞人.
那些是賣國賊.
那些是叛徒.
你反而接觸那些人.
你不潔了.
你被抓給你教人遵守律法.
你不應該接觸.
就質疑耶穌.
也都指責耶穌.
耶穌就用幾個比喻對應.
其中一個比喻就是.
說到一個家庭.
這個父親有兩個兒子.
一般我們用浪子比喻.
去描述這個比喻.
不過坦白說.
不單止浪子是個精髓.
更重要是個父親是個精髓.
所以某個程度.
這個比喻的內文是說.
這個爸爸有兩個兒子.
小兒子突然有一天就問爸爸.
將我應得的家業分給我.
爸爸真的分給他.
在這個層面爸爸沒有做錯.
一會兒再解釋.
但是這個兒子得到那份家業.
他就自己走去變賣了.
並且要離開這個家庭遠走高飛.
他覺得這個家庭束縛他.
他沒有什麼自由.
他想闖自己的天下.
能夠得到一個自由.
得到一個豐富.
然後做自己的生意.
這個大約十八到二十歲的青少年人.
走出去自然都有那種滿足.
情慾快樂的追求.
接著這些壞朋友就開始花天手地.
開始用錢在一些女人的身上.

$^{521}$但是那地方有饑荒然後生意失敗.
就淪落到做工人.
還要做最差的工人.
那個工人就是要養豬.
大家明白猶太人是不接觸豬的.
因為不潔.
但是他連最不潔的工作都要做.
還要吃豬手.
他心裡面就有一個很大的反省.
人去到這個境地.
就開始反省自己的問題.
呼求就是這樣.
他就想一下自己.
我真的得罪了天.
得罪了爸爸.
我真的不配去做兒子.
但是想到爸爸的家裡的工人.
都有餐溫飽吃.
好過我現在連豬手都沒得吃.
那不如就回爸爸家.
不做兒子了.
做工人都比沒有好.
那他就立意回家.
聖經就描述他回家的過程.
爸爸其實每天都在家裡對開.
隔居的腳尖.
然後遠遠看著.
小孩回來了沒有呢?.
果然見到小孩回來了.
這個爸爸二話不說就撲牆.
小孩想好對白.
我得罪了天.
得罪你.
我不配你做兒子.
只是配你做你工人.
都還沒說完.
爸爸已經連連親嘴.
跟著吩咐工人.
恢復他兒子的身份.
包括給他有外袍.

$^{561}$給他戒指.
給他鞋子.
這些全部是兒子身份的表徵.
而且還宰了一隻肥油毒.
去慶祝開派對.
二話不說在這樣的過程裡.
兒子在田裡回來了.
做完事.
突然間聽到有開派對的聲音.
兒子就問僕人.
發生什麼事?.
僕人就跟他說.
你弟弟回來了.
爸爸為此開心.
就慶祝.
大兒子一聽到就心不妙.
不滿.
不快.
也拒絕進入派對.
爸爸知道大哥這樣.
特意跑出來.
跟大哥互動.
這個大兒子就跟他說.
有沒有搞錯?.
弟弟敗了你的家產.
用了在妓女身上.
你還要這樣接待他.
還要恢復他做兒子的身份.
不可能.
做公人也要想一想.
我幫你這麼久.
一隻肥羊仔.
你都不給我和朋友享用.
開始計較了.
爸爸很耐心地聽完之後.
就跟他說.
你一直和我一起.
我的就是你的.
而且你的弟弟.
因為回來.

$^{601}$失而復得.
死而復活.
這個才值得高興歡喜.
整個互動.
其實這個比喻.
就是耶穌用來對應.
那些猶太的領袖.
那些宗教的領袖.
而大兒子就是他們.
因為他們不願意.
接納那些罪人.
信主和親近耶穌.
但是在整個比喻裡.
都值得我們反思.
其實誰才是浪子.
很明顯小兒子是.
但是其實大兒子是不是?.
小兒子.
他要求分家產.
在當仇的做法.
這件事不一定錯.
當仇作為父親.
如果他知道自己百年歸老.
快要中老.
他可以納為足.
讓他的產業分給兒女.
個個都一份.
長子就兩份.
所以小兒子就先問爸爸.
分他應得那份.
其實OK的.
但是這個分家產.
在爸爸還在生的時候.
有兩款的家產.
一種就是全部給小兒子.
可以掌管的.
但是另外一款就是.
只是給他使用.
但是他沒有擁有.
那種轉換的權.

$^{641}$小兒子的問題就是.
這一款他自己主動.
去變賣家產.
是屬於爸爸的.
所以侵吞了爸爸的家產.
也都奪了權.
這裡是他犯錯的恩由.
所以他要走.
他走就是希望自己.
有條生路.
有條出路.
也都自己去拼搏.
反映著什麼呢?.
小兒子是一個浪子.
他的恩由是.
他只想著自己.
他不理爸爸.
不理家裡.
只是為了自己的好處.
他著重物質.
著重自己.
爭取自己的自由.
當然也都用在.
一些情慾的事上.
他要求的就是.
我要被人尊重.
他要求的就是.
我未有愛.
我要被愛.
所以整個過程.
我用我的字眼就是.
他做了很多.
不當做的事.
這個就是小兒子.
他作為浪子的恩由.
但是大兒子.
你不要以為他很公正.
但是他的內心.
有一種掙扎和訴求.
他就和爸爸說.

$^{681}$你要公正地對待.
這位小兒子才行.
你不罰他都說了.
但是要接納他.
恢復他做兒子.
不行.
所以心裡他有把尺.
他覺得自己做得很好.
他也都覺得自己.
做得很對.
他的數算得很準.
但是他卻在計較得失.
他不肯接納那些.
他心中認為差的人.
他不願意和那些人分享.
分享本來是人應該有的特色.
所以用我的字眼就是.
他應做的不走去做.
所以小兒子有他的訴求.
爸爸給了他.
但是小兒子.
潛越了爸爸的領域.
大兒子有他的訴求.
去要求爸爸公正地對待.
一個犯錯的弟弟.
但是他缺乏了那種愛和接納.
你作為這位爸爸.
面對兩位兒子.
一個可以說是浪子.
一個是走開去.
一個就算在自己身邊.
但是心卻遠離爸爸.
你作為一個爸爸.
你面對這種處境.
你覺得難不難?.
但是爸爸用什麼態度回應?.
有沒有責備?.
有沒有敦起爸爸的款?.
閉嘴.
有沒有這樣表現?.

$^{721}$但是這位爸爸只是憐憫.
只是寬恕.
只是忍耐等候.
接納受傷的.
接納有罪自私的.
大概給他們有再生的機會.
人需要take two的.
有些時候.
生命很多時候以為.
只是一次機會.
想起市面上.
我們聽到不少的訊息.
有時有些倫常的慘疾.
我們覺得很可惜.
因為內心只是想一次機會而已.
我小時候考試制度比較難.
就是因為一次機會定生死.
你讀不到書.
考不到試就出來工作.
現在可能機會豐富一點.
選擇多一點.
或者你考多幾次都無拘.
但是我們小時候讀書.
在爸爸媽媽的艱辛裡面.
讀到你就讀.
讀不到你就出來工作.
所以壓力很大.
只有一次機會.
但是我們有沒有想過.
人生可以有第二次機會?.
這個庇育父親.
他是能夠恩愛給兒子有兩次機會.
不單只給小孩子.
不要忘記還給大兒子.
他和大兒子出來互動.
是給機會大兒子去表達.
接納弟弟這個可能.
這個爸爸願意放下自己的身段.
放開自己的雙手.
不要緊抓一些東西.

$^{761}$你都願意去聆聽.
願意去給一個寬恕.
你都願意賜人一個機會.
一個很好的爸爸.
小孩子為什麼回來.
是因為當他想到自己的問題.
我剛才說祈求是因為你自訂.
祈求是因為你知道自己不行.
你沒有.
你才會去愛求.
一個生命如果知道自己已經非常滿足.
你不會求任何東西.
不過奈何我們怎樣可以確保自己.
長年累月我們都沒有缺乏呢?.
我們都知道不可能.
所以人去到低谷.
有沒有第二次的機會呢?.
這個浪子的比喻就是說.
其實兩位兒子都有機會迴轉.
開車的你如果走錯路.
你唯一的選擇就是U-turn.
走錯路.
如果你死都說那條正路.
你就越走越遠.
謙卑的人才知道.
走錯路即掉頭.
因為你知道再走下去.
越走越差.
所以某程度知訂.
是一個我們的關鍵.
信仰於我們就是.
人要明白自己真正的需要.
你才思考信仰的問題.
你不肯面對自己真正的需要.
你就不會想信仰的問題.
我小時候問媽媽.
為什麼你生我?.
我小時候問媽媽.
為什麼你是我媽媽?.
小時候我問媽媽.

$^{801}$為什麼我生存在香港?.
沒有問第四個問題.
我臉上已經有兩個五巴掌的印.
但我問一些很深的問題.
小時候我媽媽也問我.
你死了去哪?.
我就跟媽媽說.
我真的不知道人死了去哪.
罪的定義就是.
不在箭靶中心的狀態.
亦都是我們沒有活出神.
在我們生命裡已經定定.
那個生命的意義.
能夠榮耀神.
用剛才那個比喻的特性.
就是不應該做的我們去做.
應該做的我們偏不做.
這個就是不在箭靶的狀態.
所以某程度上.
罪不是沒有東西.
罪帶來滅亡的.
會死的.
罪不是自彈彈.
中文有句話.
聖人都有錯.
何謂聖人都有錯?.
即是有機會說錯話.
想錯事和做錯事.
聖人即是很少犯錯.
如果用聖人的標準.
我當一個人.
他一日只想錯一件事.
說錯一句說話.
做錯一些事的時候.
一日三件.
你說不要緊.
三件累積起來看看.
一個月就九十件.
一年就一零八零件.
我當四四五一.

$^{841}$一千件.
當小朋友五歲不識世界.
給多一點.
十歲不識世界.
如果我現在六十歲.
有五十年的日子.
我已經累積了五萬件罪.
你不要以為小事.
其實罪引動生命.
亦都污染生命.
亦都扭轉了生命.
越行越差.
世人都犯了罪是聖經講的.
世界如果個個都搞得定.
就不用信仰.
耶穌亦都不需要為我們死.
但是因為我們的緣故.
我們犯了罪.
虧缺了神的榮耀.
亦都死路一條.
因為犯罪有功架.
人最大的問題.
或者最大的敵人.
就是罪和死亡.
做牧者一族團隊.
我們很多時候去醫院探病.
很多朋友面對死亡.
都有很真實的表現.
一些未有信仰的朋友.
他們心裡面很懼怕.
我問他們怕些甚麼.
第一他們不知道自己去了哪裡.
第二他們不明白死亡的過程.
第三他們很擔心.
他們還在生的家人.
相反我們面對一些.
信主的弟兄姊妹.
雖然疼痛的疾病.
令到他們不容易.
有些時候他們說神在念我等等.

$^{881}$但是他們知道自己將會去哪裡.
亦都很能夠放心.
因為他們將家人交託給.
他們所信的主耶穌.
兩碼子的事.
就是這段經文.
唯有耶穌給我們的恩典.
在基督裡.
不是要我們死路一條.
而是給我們生路.
曾經何時交通工具最快的.
其中有一個廣告.
就在一個牛角尖裡製作出路.
我覺得這個廣告非常好.
亦都非常熟悉.
人生在牛角尖裡裡.
如果有個出路.
是非常之有恩典的.
因為神愛世人.
神愛的世人是包括你我.
世人是甚麼.
就是還在罪裡面.
而不認識神的時候.
那些世人.
然後神做了一個很大的舉動.
就是讓他獨生愛子耶穌.
為我們負上了罪的贖價.
流出寶血.
目的是想我們乾淨.
想我們不是死路一條.
所以這個是.
當我們中了罪人的時候.
耶穌已經做了這個舉動.
在座.
你信耶穌好.
你暫時未信耶穌好.
耶穌都不會走過來問你.
你信不信我.
你信我就死給你看.
沒有.

$^{921}$耶穌是死了給你看.
然後才走來問你的.
有誰愛人.
我死給你看.
然後告訴你.
你愛不愛我.
沒有.
我們是計算.
我值不值得為你死.
耶穌說我是道路真理生命.
你想一想世界上信仰的對象.
有誰可以生勾勾的時候.
說這句話.
我就是道路真理生命.
有誰說自己是神.
沒有.
唯有耶穌.
在整個活動.
這位耶穌流出寶血幫助我們.
和神有一個復和關係.
恢復關係是很重要.
家庭的溫度在於關係的緊密度.
家庭的溫度不在於家庭的豐原.
不在於你住在哪裡有多大.
那張床有多好.
家庭的溫度在於關係.
這份關係就是聖經所說.
人與神那份關係的復和.
然後有平安.
我們用和平.
這是我剛剛去完馬來西亞怡保訪孫的時候.
有一個福音扭扭波的訊息.
這是我的作品.
就是兩個心中間有十字架連繫著.
心連心.
用人的努力很困難.
不過透過耶穌基督十字架我們就OK了.
以至到因為這位耶穌成為代做的羔羊.
為我們付上了罪的代價.
本來被釘死的是你和我.

$^{961}$但是耶穌愛的緣故.
祂願意為我們付上這個罪的刑罰.
然後問我們你想不想得生命.
所以當我們知道罪的功架是死.
耶穌裡面然後有出路的時候.
耶穌就成為了父神和人這個關係裡面一個很重要的中堡.
Mediator.
今天的Mediator叫調解員.
但其實在聖經的角度這個叫做中間人.
是成為一個橋樑上有罪的人.
可以透過耶穌成為一個義人.
去到父神的根前恢復這份關係.
就耶穌選擇為我們犧牲流出寶血.
被釘死.
你知道那些釘釘下去那一下.
知道被釘死死亡的因素是什麼嗎.
是窒息而死.
很困難的一個過程.
但是耶穌甘願為我們擺上.
祂已經為我們預備了這個救恩.
然後再來問我們.
你經不經歷我這份給你的愛與恩典.
你願不願意相信.
耶穌的犧牲是帶來平安.
耶穌的受鞭傷給我們有醫治.
耶穌的復活給我們也有盼望.
有德性的把握.
有生存的動力.
我小時候已經知道自己挺艱險的.
為什麼呢.
我讀書都OK的.
所以都不容易的時間.
都會拿第一第二那些.
通常拿第一第二都給你拿到品行獎.
中六中七都拿的.
但是那時候其實心裡面.
就是有一磨磨那些所謂的黑料.
是出不了街的.
我五歲會打麻將.
所以是高手.

$^{1001}$我贏到一個地步.
就是不要要私贏人家.
而且因為贏人所以經常吵架.
我打到一個地步就是.
中六七的時候我們學校有一個folding class.
就是沒有固定班房的.
下午不用上課.
我買通班長.
你和我take到.
然後我就和三位同學.
走去一個同學的家裡打麻將.
打風打雨我最開心.
一知道打風打雨明天不上學.
就那晚第一個開台的就是我.
近乎今上用各樣的方法.
都處理不了這個賭的問題.
打麻將的問題.
只差在尾端一手指.
但是碎了柱很奇妙.
聖靈的幫助.
我聽到麻將聲我就覺得很厭煩.
是從來都沒有發生過.
但是信了耶穌之後.
我求耶穌幫我搞定打麻將這個心癮.
神就有私人奇妙的作為.
就是一聽到麻將聲就戒.
就是按得了.
到一個地步.
我現在有一天忍不住和我說.
阿楊你現在不打麻將這麼奇怪.
那時候每次開台都是你來的.
我說是耶穌幫我不打麻將.
你耶穌這麼厲害.
我說是真的很厲害.
神給我們有復活的盼望.
不單止給我們有得勝罪惡的能力.
扭轉了我們生命的路徑.
以致我們越走越好.
所以當這位耶穌在問我們.
究竟我們承認自己的問題.

$^{1041}$呼求之前承認自己有一個限制.
如果你說沒有.
那你就自欺了.
聖經說其實每個人都需要.
因為世人都犯了罪.
世人包括你和我.
但耶穌的寶血流出.
是願意讓我們的生命得潔淨.
所以能夠得赦免是恩典.
但得赦免的關鍵是你肯不肯認錯.
小孩能夠被恢復兒子的關係.
是因為他意識到自己的罪.
他願意回頭已經知道自己錯了.
能夠認錯其實是難得的.
今日人對罪對自己的錯.
是否認拒絕逃避推卸.
始終犯罪偷食禁果.
都是推卸責任.
今日有誰犯罪.
會感應自己是我的問題.
很難得.
但信耶穌必須要經過這個過程.
你肯勇敢地面對你所犯的錯誤.
無論內心無論行徑無論言語無論心思意念.
肯回轉就是關鍵.
聖經有一個很濃烈的神學觀念叫做回轉.
或者叫歸回.
Return.
關係的復活是一種修復.
是一種歸回最基本關係的那種互動.
家人之間別人說沒有隔夜仇.
不過坦白說就是肯不肯低低地灌鋪.
無論你是什麼身份.
你是兒女還是父母還是更高的長輩.
其實能夠面對錯而認錯.
這個就是關鍵.
信仰之所以精彩.
就是讓每一個本來犯錯的人.
能夠有機會認錯.
而得到一個不一樣的結局.

$^{1081}$以致當一個人肯認.
耶穌是主.
心裡信這位神從死裡復活.
是讓我們得生命的基礎.
這個生命就得救.
信耶穌有多難.
其實很直接.
耶穌從來都沒有因為你曾犯的錯而不愛你.
人會因為別人犯的錯而不愛他.
耶穌也沒有因為你過往是多軟弱.
但是耶穌很想你生命有改變.
你想不想.
因為耶穌愛你.
所以如果你想信耶穌.
今天可以有個機會讓你去想一想.
信耶穌是用一個祈禱去接受.
這個祈禱的內容簡單.
就是你用信心.
一個悔罪的心.
接受耶穌基督進入你生命裡.
成為主和救主.
然後你的罪得赦.
你就被算為義.
成為神的兒女.
跟著你的一生就不再一樣.
因為有耶穌與你一同.
面對你人生的種種.
而你是有永活的盼望.
不需要再懼怕罪和死亡的陰影.
你想想這個生命多精彩.
在在每一個信耶穌的生命.
都可以與你作見證.
為什麼我們相信耶穌.
是因為如此這麼深厚的一份恩典和愛的領受.
所以很想現在請大家可以閉上眼.
我們都能夠面對一下自己的內心.
無論你今天是父親.
你的大日子.
還是其他每一個人生階段的你.
我們來望自己的人生.

$^{1121}$我們走過什麼路程.
當然有些很開心的路程.
不過可能都有一些很艱難的時候.
有些見得光.
有些不見得光.
但如果你很真誠面對自己.
總有些幽暗面.
總有些地方.
連你自己都未必可以接納到自己.
都不出奇.
不過耶穌說不要緊.
只要你肯面對.
你願意在神的面前謙卑下來.
你肯叩門.
神就給你開門.
你肯尋求.
神就給你尋見.
你肯呼求.
神就進到你生命.
所以我先問每一位基督徒的朋友.
你已經肯定耶穌在你的生命作主作王嗎?.
你願不願意繼續用你的生命為主作見證?.
如果願意的基督徒朋友.
我鼓勵你.
舉手示意你這份納祭.
見到相信神必定使用大家.
在在有未信主的.
這個可能是你第一次聽到福音.
或者你已經聽過很多.
不過很想鼓勵你.
當你靜下來來望你自己心靈的深處.
其實我們是否可以這樣說.
我有一個無愧的良心.
是我沒有做壞事.
不過所有好的事.
應該的事.
你有沒有做足呢?.
當我們面對自己有限制的時候.
耶穌跟你說.
不要緊.

$^{1161}$耶穌來就是幫助你.
只要你願意信祂.
能夠解決你罪和死亡的問題.
洗淨你的生命.
至你有永遠更新的生命的朋友.
這個過程.
這個處境.
你願意接受請舉起你的手示意.
還有沒有一些朋友很勇敢去表達.
比如你的父親.
他懂得愛.
他也敢愛.
耶穌懂得愛.
他也敢愛.
他很想你也懂得愛.
敢去愛.
有沒有這樣的朋友願意接受耶穌?.
我們一起去祈禱.
特別請已經舉手.
願意相信耶穌的.
我們用這個祈禱表達這個心聲.
因為你的愛長闊高深.
你的愛觸動我.
你給我勇氣面對我自己的本相.
我都是一個罪人.
我很有限制.
我不能夠幫助自己.
也不能夠自救.
耶穌你為我已經付上了罪的贖價.
流出你的寶血洗淨我的生命.
賜我永遠的那個生命的榜樣.
所以今天我願意真誠地打開我的心.
接受耶穌進到我生命的內處.
成為我的主和生命的主.
讓我能夠得清為宜.
成為神的兒女.
也恩著耶穌給我的力量恩典.
給我有平安之餘.
給我有動力.
能夠勝過罪.

$^{1201}$勝過黑暗.
在生活裡面能夠榮耀主的名.
多謝主.
感恩祈禱.
讓我奉主耶穌基督名而求.
讓我們繼續努力.
將主的名,主的愛繼續傳揚.
以致更多生命得著益處.
\newpage



\section{士師記 8:1-35}
\label{sec:032346}
\textbf{士師記析讀 (10) 失敗之父}~~~~梁永善牧師
\newline
\newline
連結: \href{http://www.church.com.hk/acms/content.asp?site=cdc\&op=show\&type=product\&code=032346\&layout=sermon}{\texttt{032346}} ~~~~ 語音日期: 2026-06-24
\newline
\newline
\hyperref[sec:032345]{\small{< < < PREV SERMON < < <}}
~
\hyperref[sec:index_chronic]{\small{[返順時目]}}
~
\hyperref[sec:032347]{\small{> > > NEXT SERMON > > >}}
\newline
\newline
士師記 8:1-35
\newline
\begin{longtable}{cl}
\hline
\hline
章節 & 經文 (和合本修訂版)\\
\hline
8:1 & \begin{tabularx}{0.7\textwidth}{X} 以法蓮人對基甸說：「你去與米甸爭戰，沒有召我們同去，你為甚麼這樣待我們呢？」他們就和基甸激烈地爭吵。 \end{tabularx} \\ \\ \relax
8:2 & \begin{tabularx}{0.7\textwidth}{X} 基甸對他們說：「我現在所做的怎麼與你們所做的相比呢？以法蓮拾取剩下的葡萄不強過亞比以謝族所摘的葡萄嗎？ \end{tabularx} \\ \\ \relax
8:3 & \begin{tabularx}{0.7\textwidth}{X} 神已把米甸的兩個領袖俄立和西伊伯交在你們手中；我所做的怎能與你們所做的相比呢？」基甸說了這話，他們對他的怒氣就消了。 \end{tabularx} \\ \\ \relax
8:4 & \begin{tabularx}{0.7\textwidth}{X} 基甸和跟隨他的三百人來到約旦河，渡了過去；他們雖然疲乏，還是追趕。 \end{tabularx} \\ \\ \relax
8:5 & \begin{tabularx}{0.7\textwidth}{X} 基甸對疏割人說：「請你們拿幾塊餅來給跟隨我的百姓，因為他們疲乏了。我正在追擊米甸王西巴和撒慕拿。」 \end{tabularx} \\ \\ \relax
8:6 & \begin{tabularx}{0.7\textwidth}{X} 疏割人的領袖回答說：「西巴和撒慕拿的手掌現在已經在你手裡，因此我們該將餅送給你的軍隊嗎？」 \end{tabularx} \\ \\ \relax
8:7 & \begin{tabularx}{0.7\textwidth}{X} 基甸說：「好吧！耶和華將西巴和撒慕拿交在我手之後，我必用曠野的荊棘和枳條鞭打你們。」 \end{tabularx} \\ \\ \relax
8:8 & \begin{tabularx}{0.7\textwidth}{X} 基甸從那裡上到毗努伊勒，對那裡的人也提出同樣的請求；毗努伊勒人給他的答覆跟疏割人的答覆一樣。 \end{tabularx} \\ \\ \relax
8:9 & \begin{tabularx}{0.7\textwidth}{X} 他也對毗努伊勒人說：「我平平安安回來的時候，必拆毀這城樓。」 \end{tabularx} \\ \\ \relax
8:10 & \begin{tabularx}{0.7\textwidth}{X} 那時西巴和撒慕拿，以及跟隨他們的軍隊都在加各，約有一萬五千人，是東邊的人全軍所剩下的，因為拿刀戰死的約有十二萬人。 \end{tabularx} \\ \\ \relax
8:11 & \begin{tabularx}{0.7\textwidth}{X} 基甸從挪巴和約比哈的東邊，從住帳棚人的路上去，趁米甸的軍兵以為安全的時候攻擊他們。 \end{tabularx} \\ \\ \relax
8:12 & \begin{tabularx}{0.7\textwidth}{X} 西巴和撒慕拿逃跑；基甸追趕他們，捉住米甸的兩個王西巴和撒慕拿，使他們全軍潰散。 \end{tabularx} \\ \\ \relax
8:13 & \begin{tabularx}{0.7\textwidth}{X} 約阿施的兒子基甸從戰場，沿著希列斯斜坡回來， \end{tabularx} \\ \\ \relax
8:14 & \begin{tabularx}{0.7\textwidth}{X} 捉住疏割人的一個少年，查問他。他就為基甸寫下疏割的領袖和長老的名字，共七十七人。 \end{tabularx} \\ \\ \relax
8:15 & \begin{tabularx}{0.7\textwidth}{X} 基甸到了疏割人那裡，說：「你們從前譏笑我說：『西巴和撒慕拿的手掌現在已經在你手裡，因此我們該將餅送給跟隨你的疲乏的人嗎？』看哪，西巴和撒慕拿在這裡。」 \end{tabularx} \\ \\ \relax
8:16 & \begin{tabularx}{0.7\textwidth}{X} 於是他拿住城內的長老，用曠野的荊棘和枳條責打疏割人。 \end{tabularx} \\ \\ \relax
8:17 & \begin{tabularx}{0.7\textwidth}{X} 他又拆了毗努伊勒的城樓，殺了城裡的人。 \end{tabularx} \\ \\ \relax
8:18 & \begin{tabularx}{0.7\textwidth}{X} 基甸對西巴和撒慕拿說：「你們在他泊山所殺的人是甚麼樣子的？」他們說：「他們很像你，個個都有王子的樣子。」 \end{tabularx} \\ \\ \relax
8:19 & \begin{tabularx}{0.7\textwidth}{X} 基甸說：「他們都是我的兄弟，我母親的兒子。我指著永生的耶和華起誓，你們若存留他們的性命，我就不殺你們了。」 \end{tabularx} \\ \\ \relax
8:20 & \begin{tabularx}{0.7\textwidth}{X} 他對他的長子益帖說：「你起來殺他們！」但是這少年害怕，不敢拔刀，因為他還是個少年。 \end{tabularx} \\ \\ \relax
8:21 & \begin{tabularx}{0.7\textwidth}{X} 西巴和撒慕拿說：「你自己起來殺我們吧！因為人如何，力量也如何。」基甸就起來，殺了西巴和撒慕拿，取了他們駱駝頸項上的月牙圈。 \end{tabularx} \\ \\ \relax
8:22 & \begin{tabularx}{0.7\textwidth}{X} 以色列人對基甸說：「你既然救我們脫離米甸的手，願你治理我們，你的兒子孫子也治理我們。」 \end{tabularx} \\ \\ \relax
8:23 & \begin{tabularx}{0.7\textwidth}{X} 基甸對他們說：「我不治理你們，我的兒子也不治理你們，耶和華會治理你們。」 \end{tabularx} \\ \\ \relax
8:24 & \begin{tabularx}{0.7\textwidth}{X} 基甸又對他們說：「我有一件事求你們，請你們各人把所奪的耳環給我。」因敵人都戴金耳環，他們是以實瑪利人。 \end{tabularx} \\ \\ \relax
8:25 & \begin{tabularx}{0.7\textwidth}{X} 以色列人說：「我們情願送給你！」他們就鋪開一件外衣，各人將所奪的耳環丟在上面。 \end{tabularx} \\ \\ \relax
8:26 & \begin{tabularx}{0.7\textwidth}{X} 基甸所要求的金耳環，重一千七百舍客勒金子。此外還有米甸王所戴的月牙圈、耳環，和所穿的紫色衣服，以及駱駝頸項上的鏈子。 \end{tabularx} \\ \\ \relax
8:27 & \begin{tabularx}{0.7\textwidth}{X} 基甸以此造了一個以弗得，設立在他的本城俄弗拉。全以色列就在那裡拜這以弗得行淫，這就成了基甸和他全家的圈套。 \end{tabularx} \\ \\ \relax
8:28 & \begin{tabularx}{0.7\textwidth}{X} 這樣，米甸就被以色列人制伏了，再也不能抬頭。基甸還在的日子，這地太平四十年。 \end{tabularx} \\ \\ \relax
8:29 & \begin{tabularx}{0.7\textwidth}{X} 約阿施的兒子耶路巴力回去，住在自己家裡。 \end{tabularx} \\ \\ \relax
8:30 & \begin{tabularx}{0.7\textwidth}{X} 基甸有七十個親生的兒子，因為他有許多妻子。 \end{tabularx} \\ \\ \relax
8:31 & \begin{tabularx}{0.7\textwidth}{X} 他在示劍的妾也為他生了一個兒子，基甸給他起名叫亞比米勒。 \end{tabularx} \\ \\ \relax
8:32 & \begin{tabularx}{0.7\textwidth}{X} 約阿施的兒子基甸年紀老邁而死，葬在亞比以謝族的俄弗拉，他父親約阿施的墳墓裡。 \end{tabularx} \\ \\ \relax
8:33 & \begin{tabularx}{0.7\textwidth}{X} 基甸死後，以色列人又去隨從諸巴力而行淫，以巴力‧比利土為他們的神明。 \end{tabularx} \\ \\ \relax
8:34 & \begin{tabularx}{0.7\textwidth}{X} 以色列人不記得耶和華－他們的神，就是那位拯救他們脫離四圍仇敵之手的， \end{tabularx} \\ \\ \relax
8:35 & \begin{tabularx}{0.7\textwidth}{X} 也不照著耶路巴力，就是基甸向以色列所施的恩惠善待他的家。 \end{tabularx} \\ \\
[1ex]
\hline
\hline
\end{longtable}
$^{1}$今晚和大家看《事實記》第十講.
就是失敗之父.
是講參宣的.
不知道大家有沒有看《事實記》第八章.
我們很熟悉的一句話是.
「失敗乃成功之母」.
經過很多次失敗.
你吸收到經驗之後.
然後你自己有改變.
這時候可能你就會成功了.
但是有另外一句說話.
我直接說就是.
「成功是失敗之父」.
為什麼呢?.
因為很多時候我們的失敗.
是因為我們的成功.
當我們成功的時候.
我們就會沾沾自喜.
我們放鬆自己.
於是我們就慢慢墮落了.
正如有錢的人就會飽斷私淫慾.
蘇羅是被神揀選第一個的王.
但是後來當他有權有勢的時候.
他著重物質.
結果就被神帶入去放棄.
大衛作王的時候他不需要逃亡.
結果他的生活糜爛.
他睡到日上三竿.
日頭平西斜起身.
他在皇宮的平台悠悠閒來閒去.
但是其實那時候要打仗.
他看到人其實是很遠的.
他出獄於是動了人臉.
結果就陷自己在這個沉獄當中.
所以原來成功往往就是帶來失敗.
問題就是我們有沒有警醒.
這封聖經就讓我們看到.
機電如何逐步逐步地走下坡.
而且他的人生的失敗不僅影響他.
是影響到他的家族.

$^{41}$還有一件事是什麼呢?.
就是機電是所有的事司裡面.
帶來危害性最大的一個事司.
因為很多時候我們只是集中在機電.
只是以小贏大.
就是這個事蹟.
但是我們沒有跟進下去.
就算我們很熟悉的參孫去放縱.
但是他的影響.
壞的影響都不及機電那麼厲害.
今天我們會看下去.
在這裡我們看到機電打贏仗之後.
有三類的人是對機電不滿意的.
而我們看到機電對這三種人.
至少有兩種是不同的對待.
其實就看到機電的政治手腕和他的皮性.
這三類人就是以法戀人.
梳割人和比奴以納人.
我們先來看看.
機電是游釋紛爭.
我記得在海報裡面講了這個題目.
但當我繼續研究下去的時候.
就不會用這個題目了.
這個只是一個分題.
以往都問我今晚的題目是什麼.
機電要追殺米電人的時候.
於是他就呼喚一些支派來幫助他.
以法戀就是其中一個支派.
當以法戀追殺了那些敵軍之後.
他就向機電提出質問.
《普氏詩記》第八章第一節.
以法戀人對機電說.
你去與米電人爭戰.
沒有招我們同去.
為什麼這樣待我們呢?.
他們就與機電大大地爭吵.
米電人他們來去打敗.
機電就要叫那些人來幫忙抓人.
但以法戀後來就跟他告.
為什麼你沒有叫我呢?.

$^{81}$一開始的時候.
還跟他大大爭吵.
我們先了解一下那個地勢.
以法戀山地是最早期以色列歷史裡面.
很少有外族入侵的地方.
所以他們得天獨厚.
他們有充分的時間去鞏固自己的地位.
時期的以法戀人最重要就是成為宗教的中心.
因為士羅和伯特利.
為什麼是伯特利呢?.
就是雅各與神相遇的地方.
就是他們宗教的重要地方.
所以我們覺得自己很好.
他們的爭吵都是有利益關係的.
因為他們大敗米電人之後.
一定有很多的老物.
掠奪了不少的財物.
一會兒我們會看到機電也要擁有這些財物.
我們很想分一杯羹.
以法戀某個程度其實是什麼呢?.
就是無理取鬧.
我們可以估計.
如果機電當時只有300人.
接著要對十幾萬的米電人.
就算當時他呼召以法戀人.
「你來幫忙」.
你估計以法戀人會不會幫忙呢?.
「你不要搞我」.
「我哪有膽量」.
你也要留意一件事.
做領袖你經常都會被人批評的.
叫他他不理你.
你不叫他他會說「為什麼你不叫我呢?」.
還有就是以法戀人.
如果你留意到他一直都覺得自己.
是比其他支派優勝很多的.
因為十二支派之中.
他覺得自己是最尊貴的.
以法戀人是誰呢?.
十二支派.

$^{121}$如果你留意到有十個支派.
是雅各的兒子.
但是以法戀嗎哪西.
不是雅各的兒子.
是雅各的孫子.
他是約瑟的兩個兒子.
而以法戀就是次子.
是第二個兒子.
但是雅各本身是第二個兒子.
所以他是很特別的.
他要將特別的祝福給次子.
即是以法戀.
結果以法戀他就能夠與白富平起平坐.
及後北角以色列.
都用以法戀作為一個代表.
所以你就可以想像到以法戀.
自以為是.
他自己是很厲害的.
所以他就要向諸多不滿.
向機電來挑戰.
稍後我們會看.
應該是第三次.
我們會看到耶佛他打贏了阿們人之後.
以法戀又照樣主管.
向耶佛他挑戰.
為什麼你們不叫我呢?.
但是結果是截然不同的.
這些我們稍後會講.
為什麼他們會這樣呢?.
我經常都這樣講.
一個清潔工人.
可能經常被人罵的.
普通的民賢就罵他.
可能洗東西洗得不乾淨.
廁所那麼髒.
老闆就罵他.
他不敢駁嘴的.
但是如果一個掌權者.
一個老闆.
如果被下屬罵了一兩句.

$^{161}$死定了.
就馬上解僱.
所以你就留意到.
一些人在高位的時候.
報案是不能夠被人指出他的錯.
在這裡我們看到.
機電就用一個很溫柔的語氣.
很圓滑的答案.
來跟米電人聊天.
然後平息他們的怒氣.
我們看看《史詩紀》第八章第二到第三節.
「機電對他們說.
『我所行的豈能比你們所行的呢?』.
以法戀拾取剩下的葡萄.
不強過阿彼爾謝所摘的葡萄嗎?.
朕已將米電人的兩個首領.
俄勒和西伊伯交在你們手中.
『我所行的豈能比你們所行的呢?』.
機電說了者說.
『以法戀人的怒氣就消了』」.
你看看機電其實剛剛打贏仗.
但是他是很忍氣吞聲.
很謙卑,很謙讓的.
我做了些甚麼給你呢?.
以法戀摘剩下的葡萄.
即不是最好的那些.
都比我阿彼爾謝家族裡面所摘的葡萄好.
即是你最摘的葡萄比我最好的葡萄好.
而且兩個米電人的首領都已經在你手裡面.
那我何必及你呢?.
結果這個這麼溫柔的回答.
又使他們的怒氣消退了.
這就是正正《箴言》裡所說的.
我們一起來看看《箴言》第十五章的第一節.
「回答又和,使怒消退,言語暴戾,觸動怒氣」.
現在大家在很熱的天氣裡面.
有時會心火燥的.
特別都要很小心我們的言語.
特別可能太太在這個時候很寶貴.
很熱的時候都要煮飯.

$^{201}$可能心火燥罵你兩句.
特別丈夫都鼓勵你.
突然吹一聲「忍氣吞聲」.
先忍了它.
千萬不要勞言相對.
否則就糟糕了.
但不要忘記.
忍讓都要有底線的.
如果不理會對或錯.
甚麼都說忍讓.
甚至對方是無理要求,指責.
我們都忍讓的時候.
只會越縱容對方.
令到他越來越放肆.
最終其實會令到對方失敗和吃虧.
及後耶佛他.
剛才我已經說過了.
面對以法連同樣地挑戰他的時候.
但耶佛他就完全不是這樣.
但是可以對以法連之派的人.
結果以法連之派很慘.
差點被耶佛他剿滅了.
我們遲些會看.
所以容忍的態度.
是面對挑戰這個原則是不變的.
但是這次是否用在基甸身上呢?.
未必一定是很正確.
因為神曾經說過.
你的人數太多.
所以就一直削減到三百人.
但基甸因為要乘勝追擊.
於是就召喚了不同的支派.
來帶入去.
叫他們來參與.
這究竟是否合乎神的心意呢?.
教會有時候有些事工.
是覺得力量不夠.
要找些人來幫忙也是對的.
但是會不會找一些未信的人來幫忙呢?.
我見過有些人真的是這樣的.

$^{241}$這是否合乎神的心意呢?.
絕對不是.
基甸甚至將兩個米殿的首領.
交在以法連的手上.
讓祂發落.
讓祂們得到榮耀.
是否合乎神的心意呢?.
祂似乎重視人.
多於重視神.
因為神很清楚說.
要將米殿人交在基甸的手中.
神沒有說要將米殿人交在以法連的手中.
我們看看《事實紀第七章》第七節.
特別紅色的字.
我已經將米殿人交在你手中.
所以要很小心.
原來祂取悅人.
是多於取悅神.
祂用祂的智慧考量.
去討好以法連人.
多於去討好神.
祂絕對不想得罪這個支派.
免得祂成為領袖的攔阻.
因為也要買祂的派.
正如香港如果你選特首.
現在有幾百人要提名.
如果那個要選特首的人敢不敢得罪那些.
有份提名祂的人.
有能力提名祂的人.
當然不敢.
當然要討祂歡心.
怎麼可以得罪祂呢?.
所以我們看到.
基甸及後對另外兩個支派.
都是自己人.
反應是截然不同的.
這就是為了政治因素.
而捨棄了真理.
所以不少基督徒的政治人物都是這樣.
面對以法連人的質詢.

$^{281}$祂沒有將自己蒙召的過程.
神怎麼引領祂.
讓祂有人證.
祂怎麼用三百人來打贏這場仗跟我們分享.
如果祂將這些經歷跟我們分享的事實.
擺在眼前.
以法連人也無話可說.
所以這一場聖戰之後的基甸.
已經慢慢失去了祂的意象和使命.
祂體重人多於體重神.
正正是這種心態.
讓祂一步一步墮落下去.
接著面對另外有兩批人反對祂.
這兩個宗派的人怎麼來反對祂呢?.
我們看看第四節到第七節.
「也和跟隨他的三百人到約旦河過道,雖然疲乏,還是追廣..
基甸對所覺人說,求你們拿餅來給跟隨我的人吃..
因為他們疲乏了,我們追廣米堅人的兩個和西巴和撒姆拿..
所覺人的首領回答說,西巴和撒姆拿已經在你手裡,.
你使我們將餅給你的軍兵嗎?.
基甸說:也和說將西巴和撒姆拿交在我手之後.
我就用野地的荊條和節稽打傷你們.」.
如果你留意到的時候.
所覺和彼羅伊勒這個地區.
其實是屬於哪個支派呢?.
是屬於迦德支派的分地.
迦德支派在哪裡呢?.
在約旦河分為河東和河西.
原來有兩個支派的人.
其中包括迦德支派的人.
因為當時他們的財產很多.
牛羊很多.
他跟摩西說我們不想去過河.
如果按照約書記十三章二十七節記載.
所覺在約旦河東屬於迦德支派.
而迦德支派因為非常富有.
所以他們就不和其他人去河西那邊.
所以其他人已經對他們不滿了.
我們看一看.
當時摩西說一樣很重要.

$^{321}$我都讓你留在河東.
但如果要打仗的時候.
你的兄弟有需要你就要過來幫忙了.
我們看回這段歷史.
民數記三十二章二十五到二十七節.
迦德子孫和劉弁子孫對摩西說.
僕人要照我主所吩咐的去行.
我們的妻子,孩子,羊群和所有的生畜.
都要留在激烈的國城.
但你的僕人凡拿兵器的.
都要照我主所示的話.
在日和華面前過去打仗.
你留意到當時劉弁和迦德就求摩西讓你留在河東.
摩西說沒問題.
但如果你的兄弟要打仗你就來幫忙了.
所以有某個責任.
他們是應該要給機電一些援助的.
機電不是要他們參戰.
只是要他們拿一些食物而已.
但這些領袖很可能他們都妒忌機電.
因為第六節回答的意思是甚麼呢?.
你們已經勝利在望了.
為甚麼還需要我來幫忙呢?.
但機電的憤怒回應.
沒有指出你當日曾經答應過摩西的.
你的誓言沒有來得及去承諾.
更加沒有引用約書亞當年發出的神聖召命.
當年約書亞要面對戰爭的時候.
都要求河東的劉弁和迦德之派.
要來打這場仗.
但他用神的召命叫他們來參戰.
我們看看約書亞第一章的十二到十四節.
約書亞對劉弁人,迦德人和馬來西本之派的人說.
你們要追獵耶護華的僕人摩西所吩咐你們的話.
說耶護華你們的臣使你們得享平安.
也必將這地賜給你們.
你們的妻子,太子和生畜.
都可以留在約旦河東摩西所給你們的地.
但你們中間一切大能的勇士.
都要帶著兵器在你們的弟兄面前過去.

$^{361}$幫助他們.
看見嗎?.
當日約書亞要求河東的人.
就是劉弁,迦德和馬來西本之派的人.
要過去幫忙.
他們有幫忙.
但這次機電沒有.
沒有引用神聖的召命.
他很生氣.
因為這群疏國人來拒絕他的要求.
說到「我會報復你」.
他的憤怒充滿了他的心.
他沒有引用神的說話.
他甚至要報復的時候.
這場仗已經不是聖戰.
即不是神聖的戰爭.
是私人的戰爭.
而最重要在這時候.
機電覺得自己是受到挑戰的.
而不是神的命要被告他.
他覺得自己的地位,名譽受到虧損.
於是他向疏國的首領要求糧食是合情合理.
但處理這些反對者的方法.
他是值得商榷.
待會我們會看他怎樣處理這些反對者.
第三批反對他的.
是比勞爾納的居民.
他們要求同樣的被拒絕.
機電更生氣.
你前一班人也是這樣對我.
現在也是這樣對我.
於是他要宣告做甚麼呢?.
更加嚴厲.
要拆毀他們的城池.
我們看第八,第九兩節.
機電從那裡上到比勞爾納.
對那裡的人也是這樣說.
比勞爾納也與疏國人回答科里伯一樣.
他向比勞爾納人說.
我平平安安回來的時候.

$^{401}$我必拆毀這樓.
他對迦德之派的報復.
就是要打傷他們.
但這次更加嚴厲.
去到比勞爾納的人就說我要拆你的城.
要將你們的人殺戮.
你不要忘記.
機電被興起.
就是要解救自己的子民.
不要被米顛人欺負.
但現在呢.
機電由拯救者的身份.
轉變成為克制自己人的壓迫者.
正如一些搞革命的人.
就說要解放人民.
讓人民得到自由.
不用再被剝削.
但你看到這些奪權的人.
往往使自己在江山穩固的時候.
更加來壓迫之前的人.
比之前更加的深.
很可悲的歷史不斷重複.
最簡單舉個例子.
伊朗以前就叫巴列維政府.
就說它腐敗.
於是高美尼用神權政治.
推翻了巴列維王朝.
但誰知它對人民的壓迫.
是更加更加的厲害.
正如近代的中國.
大家都可以想像到.
或者知道那個情況.
其實都類似.
機電遭受到同胞的拒絕.
它的請求受到挫敗.
但它沒有向神尋求幫助.
沒有問神我應該怎樣做.
反而是用自己的軍事力量.
來鞏固自己的領導地位.
去忠於自己的承諾.

$^{441}$它勝利之後.
就回到蘇國.
鞭打城內77個領導.
大家回去看這些經文.
我們就不詳細看了.
後來又拆毀了比盧伊勒的城樓.
使他們沒有辦法來防範外敵的入侵.
並且殺了城裡的人.
使這裡的城永遠得不到重建.
不要忘記這些是他們自己的同胞.
不是外人.
大家都是亞伯拉罕的後裔.
他這個做法.
是絕對不合神的心意.
他為甚麼會這樣做.
就是基於自己的地位被輕視.
而作出這個審判.
所以他又怎樣做呢?.
我們看看16,17兩節.
於是捉拿.
捉住那城內的長老.
用野地的荊條.
和截擊.
襲打蘇國人.
又拆了比盧伊勒的城樓.
殺了那城內的人.
原來古時以色列人打樁架的工具.
是兩塊連在一起的木板.
下面有很多石塊.
打樁架的時候.
將木板一面壓在谷上磨.
磨了那些殼出來.
身上就拖著這些.
磨這些谷.
墨水就會脫落.
所以他就將蘇國的首領.
用這些荊棘條.
放在一起.
來鞭打他們.
讓他們不單單受傷.

$^{481}$甚至可能因為這個緣故而死.
但他都是對城內的長老.
這個懲罰是非常非常殘酷.
是不人道的.
他們曾經都用這種方法.
去懲罰亞蘭人.
在《阿摩斯書》第一章第三章裡面有記載.
不過我們不看了.
但基澱跟著就殺了比盧伊勒城內的人.
他就要報家族的血仇.
他其實就是取了神原本應有的權利.
他過了約旦河之後.
他的表現已經不再是為神為人民打仗.
不再是聖戰.
他將比盧伊勒城內的男人全部都殺了.
使城內的女人暴露在外敵當中.
還有一件事你有沒有留意到呢?.
他對以法連人.
因為以法連人是有權有勢的.
會影響他做領袖的地位.
所以他就討好以法連人.
但他就看不起比盧伊勒和其他的領袖.
所以你留意到的時候.
大家都是同胞.
但是他對比盧伊勒有不同的看法.
所以帶來的問題就是令到他們之間分化.
同樣都是.
你想想如果一些國內來的人亂拋垃圾.
只是說我們要忍耐接納他.
一些紀律部隊的人.
特別是警察犯罪.
但往往得到輕判.
但普通的市民.
香港的市民就不是這樣.
你這樣的時候有兩把尺.
對著同樣犯錯的人.
但是有完全不同的對待.
只會觸動那些人對其他人不滿.
其實歷史就是這樣.
所以基煉他遭受同胞的拒絕.

$^{521}$他的請求受到挫敗.
但他沒有向神尋求他的幫助.
反而用軍事力量來鞏固自己的領導地位.
所以他就很敗了.
結果被神來懲罰.
我們留意到原來做領袖.
最大的危機是什麼呢?.
就是自我中心很重.
當地位越高就越不願意被人挑戰.
往往就令他變成一個獨裁者.
基煉由拯救者變成壓迫者.
他被呼召的時候.
他是膽怯的.
是懼怕的.
但當他獲得成功之後.
他就變成一個政客.
為自己來籌算.
所以他用很委婉的言語.
來平息以法連支派的怒氣.
因為以法連支派是一個強大的支派.
是他日後要成為領袖.
會被人幫助.
不要忘記.
底波拉記不記得?.
這位女事師.
都是來自這個支派的.
即使他們有出色的領導者.
而且也有不錯的軍事力量.
可以幫助基煉追殺米電人.
所以他的政治觸覺.
就叫他一定不可以開罪以法連支派.
所以就用間接的手段.
去平息他們的憤怒.
但另一方面.
就用無情和鐵腕的手段.
去對付較弱的人.
這種前後不一的作風.
我們見到的就是欺善怕惡的人.
見高就敗 見低就踩.
他輕忽了神的律法.

$^{561}$用自己來替代了神.
他們這些戰役.
不再是神聖的戰役.
是私人的鬥爭.
你有沒有留意到基煉.
他對神有懷疑.
但神容忍他.
一次又一次地去給他一些的確據.
但他現在成功.
就不能夠容忍別人對他的疑惑.
所以這是很大的問題.
還有.
他就是很難去接受挑戰的.
好 我們看看.
接下來那場戰爭剩下來是怎樣.
《士詩》第八章第十節.
那時西巴和撒姆拿.
並跟隨他的軍隊都在加國.
約有一萬五千人.
就是東方人全軍所剩下的.
已經被殺約有十二萬拿刀的.
其實他們已經殺了十二萬多人.
但有一萬多人正在逃亡.
有兩個人就是西巴.
我經常弄錯了.
我怕巴西.
特別看足球.
西巴和撒姆拿.
但是我們在第八章裡面就看到.
這一萬五千人是不斷走的.
那時我們就看到.
他們走的時候是怎樣呢?.
就被基煉抓了.
我們看看第十二節好不好.
「西巴和撒姆拿逃跑.
基煉追趕他們.
抓住米顛的二王西巴和撒姆拿.
驚散全軍」.
好了 基煉厲害.
抓了這兩個王.

$^{601}$其他人就來不及了.
於是基煉就帶他們去詢問一些事情.
好 我們看看第八章第十八節.
「基煉問西巴和撒姆拿說.
你們在他珀山所殺的人是什麼樣子?.
回答說他們好像你.
各人都有王子的樣子」.
如果你留意到.
原來之前他們在他珀山裡面.
在他珀山裡面殺了一些人.
基煉就問他們那時殺的是什麼人?.
他就揶揄他們.
這兩個王都厲害.
「好像你這樣.
好像王一樣樣子」.
其實神夫召他們打這場仗.
是沒有叫基煉翻舊帳.
基煉蒙召就是拯救以色列脫離苦難.
但他就很想挖苦這兩個王.
反而給這兩個王.
反正都死了.
他們真是硬漢子.
他挖苦回去.
知道自己都難逃一死.
所以臨死擺出最後一次機會.
要羞辱基煉.
相信他之前曾經在他珀山.
和以色列人有爭戰.
並且殺了不少以色列人.
所以另一個很諷刺的口吻.
就說「你好像你推王子一樣.
但是他都死在我手上」.
基煉聽到之後就很憤怒.
馬上怎樣來回應呢?.
我們看看《時詩記》第八章第十九節.
「基煉說:他們是我同母的弟兄.
我指著永生的耶和華起誓.
你們從前若存留他們的性命.
我如今就不殺你們了」.
他留意到.

$^{641}$他是我同母的弟兄.
記住不是同父的弟兄.
是同母.
因為當時很可能.
他父親不止一個妻子.
所以他是同母的弟兄.
這時候一些同母的兄弟.
是比較親近的.
正如雅各有十二個兒子.
但是雅各最愛的妻子所生的兩個.
就是約瑟和辯雅敏.
所以約瑟和辯雅敏的關係特別親近.
約瑟特別疼愛辯雅敏.
但是對哥哥沒有那麼親近.
所以你留意到.
他現在是公報私仇.
到了當日.
你不是殺了我的兄弟.
我今天就留你的性命.
是不是?.
他甚至高舉神的名.
我指著耶華來起誓.
如果你當日留下他們的命.
我就不殺你.
其實他已經犯了.
世界裡不可忘稱神的名.
於是基澱又怎樣呢?.
他就要奪財.
叫他的大兒子益帖.
殺這兩個王.
因為這兩個人都已經被抓住了.
可能已經綁著手綁著腳了.
藉此去建立自己的兒子的地位.
更加就是說.
你兩個人不配讓我殺你.
就讓我兒子來殺你吧.
想來去羞辱他.
誰知他這個兒子.
真的不是很長進.
因為可能太小了.

$^{681}$竟然不敢拿刀出來殺西巴和撒姆蘭.
於是又再被這兩個人羞辱一次.
我們看看第四詩.
第八章二十到二十一節.
於是對他的長子益帖說.
你起來殺他們.
但益帖因為是童子.
害怕不敢拔刀.
西巴和撒姆蘭說.
你自己起來殺我明白.
因為人如何力量也是如何.
既然就起來殺了西巴和撒姆蘭.
奪獲他們樂佛頸上大的月牙圈.
他的兒子是一個童子.
但要記住.
聖經裡面的童子不是十一,二歲以下.
是一個少年人.
甚至是一個青年.
四十歲都是少年人請你記住.
但這個兒子害怕不敢去殺他.
這兩個王子說.
你不行的.
要死就殺我吧.
你有力量就殺我吧.
於是姬澱被他這樣挑戰.
意思是你有勇氣.
你來殺我吧.
何須叫小孩來殺我行兇呢?.
你現在是否沒有膽量呢?.
原來就是這樣挑戰姬澱.
這是引用當時人的諺語.
人如何力量也如何.
就是這一句說話.
姬澱原本是用這句說話來羞辱這兩個王.
就是一個小孩子都可以殺你.
但現在事情是出乎意料之外.
是倒轉過來的.
所以姬澱必定親手來殺這兩個王子.
否則他們反過來嘲笑姬澱的力量.
只不過和他相等.

$^{721}$即是他殺這兩個王子已經滲入了.
要維護自己的明星的心態.
所以用最無情的方法來進行殺戮.
接著我們留意待會我們會看這一句.
就是奪獲了他們樂陀項上大的月牙圈.
我們看到原來人是會慢慢走下坡.
請你記住一個人的失敗.
絕對不是剎那之間.
一個地方越來越衰敗都不是剎那之間.
是不斷地走下坡.
是慢慢漸進慢慢惡化.
所以聖經提醒我們要時刻警醒反省.
如果我們有問題不早些解決.
後果是難以想像的.
你看弘福院大火.
現在聽證會一路一路揭露很多的事.
但也有問題是甚麼呢?.
沒有這個法律效力.
即是今天的聽證會已經在說.
他們已經做了實驗.
如果那些彭網可以阻延那些火洩的話.
是不會燒成這樣的.
現在就是燒的時候還拉上結他.
但是你想一下.
民建聯的區議員黃碧嬌.
或者一些政府官員有沒有出來道歉認錯.
有沒有問題呢?.
沒有啊.
積緊民憤.
我自己最近都有一個感受.
因為我沒怎麼除草.
過了一段時間.
我的女兒已經跟我說要快些剪草.
否則草會越來越長.
草很長.
終於有一天我就剪了.
怎知道我只剪了一行.
我這一部剪草機.
就已經馬上腎了.
壞了.

$^{761}$救不了.
因為草太長了.
所以用了幾年的機.
就沒有了.
結果怎樣呢?.
我就要浪費多些錢.
買一部新的剪草機.
都是很大的.
提醒原來你不早些去清除的話.
到你要去清除.
是會很難的.
我最多是這樣.
當機電戰勝了米電人之後.
我就看見它慢慢慢慢敗壞了.
而且最大的一個不幸就是.
它個人的敗壞.
是影響到整個民族.
我們看看士斯基尼伯章22到23節.
好像很好.
我們看看.
「以色列人對機電說.
你既救我們脫離米電人的手.
願你和你的兒孫管理我們.
機電說我不管理你們.
我的兒子也不管理你們.
唯有耶和華管理你們」.
機電由一個膽怯的人.
變成一個民族領袖.
以色列人就請求機電管理他們.
其實不要忘記.
拯救他們的不是機電.
是神的作為.
我們在這裡看到機電都似乎不錯.
因為他沒有接受他們的邀請.
並且指出.
「唯有耶和華管理你們」.
但是他沒有詳細說出.
神怎樣拯救他們.
以色列人很想立王.
因為他們見到隔離的國家都有王.

$^{801}$而且他們的國勢很強盛.
所以這是一個很現實的考慮.
因為當時的君王主要是帶領百姓去作戰.
而且他們過去一直都是以耶華為王.
故事他們很想有一個實在的王帶領他們.
機電解救了他們脫離了米電人.
那七年的管轄.
所以百姓很想立他為王.
機電表面拒絕了百姓的建議.
但我們看到他其實已經實質.
接納了他們的提議.
他沒有真正拯救以色列人.
脫離了壓迫.
因為他除滅了米電王之後.
他卻是以自己成為王.
有甚麼表現呢?.
第一.
他體重物質.
剛才我看了一篇經文.
機電向以色列人提出一個要求.
就是要求他們將那些物質給他.
這個要求已經顯示出他有問題了.
他就是要將所奪的耳環給他.
這是他們部分的老物.
我們看看《士師記》第八章二十四到二十六節.
機電又對他們說:我有一件事求你們.
請你們各人將所奪的耳環給我.
原來受的是以色列人都是戴金耳環的.
他們說我們情願給你.
就鋪開一件外衣.
各人將所奪的耳環吊在其上.
機電所要出來的金耳環.
重一千七百克錫克勒金紙.
此外還有米電王所戴的月環.
瑞和所穿的紫色衣服並樂堂頸上的金領子.
你有沒有留意到機電要求一些財物?.
以色列人就戴金耳環.
米電王就戴玉環,耳兌這些珠寶.
月牙圈就是彎月形的金屬裝飾.
古代近東的遊牧民族.

$^{841}$他們多數是去拜天的.
向星體來拜的.
所以他們會拜月光.
這時這些就已經是宗教色彩的物品.
這是一個很大的問題.
金耳環多數是婦女貴重的飾物.
但顯然你看到.
這些金耳環不是從女人耳朵拔下來的.
是米顛士兵所擁有的.
即是說米顛人都搶了很多其他人的財富.
所以這一場戰爭已經不再是聖戰.
他們現在拿金耳環這些金飾.
是1700射赫勒的金紙.
即是大概多少呢?.
就是20公斤的金.
機電只是要百姓給金耳環.
但他們將更多的飾物交給機電.
剛才我們看過.
他還要奪取他項上戴的月牙圈.
迦南王就有收集金銀財寶的習慣.
這一刻機電就要將他們的財物據為己有.
這就顯示他已經被外邦人的生活和宗教影響.
正如我們都會受西方的一些儀式的影響.
譬如韓國的男士都很女性化.
現在發覺中國很多年輕人都扮韓星.
都是很女性化.
化妝也是.
即是說我們往往會受外族影響.
機電就會受著這些米電外族人影響.
越來越重視物質.
如果我們再看回原來當日摩西都曾經贏了米電人.
但神如何吩咐他們去面對那些老物呢?.
較為長的.
我們大家看PowerPoint.
我們只是來帶入去看第27節.
我們一齊來帶入去看.
請耳光堂讀27節.
把所老來的分作兩半.
一半歸於出去打仗的精兵.
另一半歸於傳媒中.

$^{881}$沒錯.
他拿了禮物是分給打仗的人.
另外一半就要來交給一些平民.
還有一樣東西很重要.
但機電沒有做到的.
我們看29節到31節.
從他們一半之中要取出來高級祭司以利亞薩.
作為耶和華的舉祭.
從以色人的一半之中.
就是從人口,牛,驢,羊群各樣三畜中.
每五十除一.
交給所守,看守.
耶和華將莫的利未人.
於是摩西和祭司以利亞薩.
照耶和華所吩咐摩西的行了.
這一段經文.
如果你再自己去看就知道.
就是摩西打贏了米電人.
努力了很多的東西.
但神吩咐他們.
不是一個人要的.
分給一些有需要的人.
分給平民.
分給打仗的人.
但更加重要的是要怎樣?.
要來帶去.
抽出來.
奉獻給神.
所以你看到機電所要的.
與適合摩西所做的完全不同.
機電就將收集了金銀珠寶.
做了一個伊弗特.
正正是因為這個伊弗特.
就陷以色列全國.
進入了一個深淵裡面.
所以剛才我一開始說.
機電是這麼多個祭司裡面.
是令到他們有最壞影響的一個祭司.
神要求以色列人做的伊弗特.
根據《出埃及》的28章6到30節.

$^{921}$我們看到這個伊弗特是神所要求去做.
但絕對不需要用20公斤的金紙.
所以這裡機電所做的伊弗特.
不是神所指示的伊弗特.
是一個偶像.
將它樹立起來.
及後我們見到米加加中的伊弗特.
都是與其他的偶像並列.
這就是去到第17章.
第14和第18節.
就是我們去到《士師記》差不多最後的時候.
會看到的.
他製造這個伊弗特的目的不是要拜偶像.
記住他不是拜偶像.
但讓這個伊弗特好像是神的神諭.
很讓他可以治理以色列人.
但因為他做了這個偶像.
就引領了以色列人去拜偶像之路.
所以當日亞倫和耶羅波安做金牛獨.
他都不是說金牛獨就是神的.
他只不過是有他自己的政治恩訊.
但卻令到以色列人陷入了這個罪.
機電是他自己做這個伊弗特.
是個人的戰利品.
代表了他獲得重大勝利.
但很可惜這個戰利品成為了以色列人敬拜的對象.
詩詩記第八章第27節.
「機電以此製造了一個伊弗特.
接納在本城俄弗拉.
後來以色列人拜拿伊弗特.
行了邪淫.
這就作了機電和他全家的網羅」.
我們這樣看的時候.
就是在俄弗拉這個地獄裡面.
做這個伊弗特.
你有沒有留意到這個字出現了兩次.
第一次是甚麼?.
大家記不記得?.
就是機電在俄弗拉這個地方裡面.
拆除了爸爸拜巴力的偶像.

$^{961}$他爸爸就是巴力的祭司.
他在俄弗拉裡面拆除了巴力的偶像.
現在他就在俄弗拉這個地方裡面.
建立了這個伊弗特.
讓自己的子民陷入拜偶像的處境當中.
多麼可悲.
成也蕭何?.
敗也斯何?.
機電死了之後.
甚至他們設立了巴力比利圖.
作為敬拜假神的地方.
《士師記》第8章第33節.
機電死後.
以色列人又去隨從諸巴力行邪淫.
以巴力比利圖為他們的神.
所以你看到因為機電的貪婪.
高舉自己.
令到整個民族陷入墮落的情況.
機電從一個拯救者.
變成了克制自己的人.
我們已經看到了.
但因為更加看到機電.
成為一個粉碎偶像的人.
但現在竟然變成一個製造偶像的人.
所以剛才我說成也蕭何?.
敗也斯何?.
他贏得的勝利結果交還給巴力.
是很可悲的一件事.
接著我們看到他生活的墮落.
機電去做以伏德.
他的墮落更加可以從他家裡有反應出來.
我們看看他的家庭.
《士師記》第8章第29到32節.
約瓦斯的兒子耶路巴力回去.
住在自己家裡.
機電有七十個親生的兒子.
因為他有許多妻.
他的妾住在士堅.
也給他生了一個孩子.
機電與他起名叫阿比米勒.

$^{1001}$約瓦斯的兒子機電年紀老邁而死.
葬在阿比爾謝族的阿弗拉.
在他父親約瓦斯的墳墓裡.
這裡很清楚講到機電有眾多妻妾.
而且為他生了七十個兒子.
成為一個大家庭絕對不是一個短時間.
都是一段漫長的時間.
當他大勝米電軍隊的時候.
他得他的長子.
長子只不過是一個孩童.
幼童.
當然可能是十多二十歲.
但當伊斯蘭人要求他和他的兒孫管理他的時候.
他只是一個兒子.
因為在愚民裡面只是單數.
所以機電已經被家男人那種生活影響.
所以有很多妻妾.
我今天跟他們吃飯.
都是在講一些年青人.
香港來的年青人.
覺得一夜情是沒有所謂的.
男男女女.
有時候有一晚之間發生親密的關係.
沒有所謂的.
這裡的人都是這樣的.
只要大鼻人肚便可以了.
這裡有多麼可悲呢?.
所以我們都要很小心.
他雖然沒有自稱為王.
但他卻過著王一般的生活.
還有機電為自己的兒子取名是甚麼呢?.
阿比米納.
意思就是我父親是王.
所以他和這個兒子取這個名字的時候.
就是叫甚麼呢?.
就是看出他的狂傲.
他更加以自己來代替神的地位.
他沒有以公義和敬虔的心管理以色列人.
反而用政治手腕弄權.
玩政治手段.

$^{1041}$在這張聖經裡面.
似乎都是看不到神的名字出現.
這張聖經就以悲劇作為結束.
34到35節.
「以色列人不記念耶穌說他們的神.
就是拯救他們脫離四處受敵之受敵.
也不照著耶路巴黎就是機電向他們所施的恩惠.
好待他的家」.
機電成為領袖之後.
他以個人的好處作為焦點.
他期望以色列人會敬重他和他自己的家族.
但是當他死了之後.
以色列人就完全不再照機電所做過的事.
好待他的家人.
這個時候就很可悲.
所以你留意到這句說話.
「也不照著耶路巴黎就是機電向他們所施的恩惠.
好待他們」.
我們看到一個人越想追求別人對他的敬重.
但很可能就會令到其他人對他更加輕視.
歷史裡的人物就是這樣.
他們很想甚至立法.
你要尊重我.
但怎可以立法來尊重你呢?.
就算表面對你有尊敬.
心裡也鄙視你.
原來一個人只是追求自我榮譽.
只會帶來無盡的毀滅和羞辱.
所以為何今天是說失敗之父呢?.
作為一個領袖一定會被批評.
當受到批評時.
如何去回應才是問題所在.
這一章聖經特別提醒我們.
我們不可以靠賴錯誤的政治動機.
作為領導的終極手段.
信服神的帶領比玩弄政治手段勝過千萬倍.
機電也缺乏了這種堅持不懈的心.
更缺乏了長期依賴仰望神的心.
他本來有很大的潛力.
可以像摩西那樣.

$^{1081}$但我們看不到他的謙卑,敬虔的表現.
他雖然獲得成功.
但很可惜只是短暫.
雖然在咀唇上他去信服神.
甚至以神的名起誓.
但與他所做的絕不相符.
機電追求的榮耀.
但最大的敵人反而是甚麼呢?.
就是不斷羞辱他的人.
這張聖經特別提醒我們.
要如何面對成功.
不要忘記失敗乃成功之母.
但成功往往是失敗之父.
隨著成功而來的海台生活.
再加上我們的自我膨脹.
以及自我中心.
就會讓領袖陷於失敗之地.
成功背後所隱藏的誘惑.
不知不覺之中.
是攻克了我作為領袖的心智.
得完其來的成功.
我們更加不可以掉以輕心.
如果機電在得勝之後.
能夠繼續有謙卑的心跟隨神的時候.
或許那些子民所享受的平安會更大.
所以記住.
屬靈的領袖應該視自己的成功.
但是他覺得要將榮耀歸給神.
受人尊敬的虛榮.
只是浮雲和幻夢.
也要記住.
一時之間的凱旋.
未必是長遠的勝利.
剎那一夕之間的成功.
也不會是永遠的勝利.
所以有一段經文成為我的提醒.
我們用這節經文作為結束.
第十章第十二節.
所以自己以為站得穩的.
需要謹慎.

$^{1121}$免得跌倒.
這個就是我自己不斷提醒自己的.
因為過去都曾經試過以為自己很厲害就失敗.
哥羅多前書十章十二節.
大家都要謹記.
所以求神幫助我們.
這一刻的信服.
不代表我們能夠恆久同神也好的關係.
當我們有警醒的心依靠神的時候.
持久的信靠.
先就是真正的勝利.
我們一起來祈禱吧.
天王多謝你讓孩子有機會.
和弟兄姊妹一起來.
去查考神的話語.
神祈求你親自來眷顧.
讓我們透過你的話語讓我們有反思.
讓我們追求的不是自己的榮耀.
我要看到原來我們要高舉主理.
我們自然就會得到主理的祝福.
更要讓我們看到一次的勝利.
不是永遠的勝利.
如果我們滿足於目前.
我們就會有失敗.
讓機電成為我們的提醒.
也透過你的話語成為我們的幫助.
願意顛覆你自己.
彌大人去祝福每一個學習的人.
特別在這麼熱的天氣.
還有世界盃的一個直播.
但我們仍然有這麼多弟兄姊妹.
來直接的彌大人去參與.
我們彌大人可以感恩的.
祈求主你自己親自的彌大人去眷顧.
也祝福在YouTube裡看到的弟兄姊妹.
也給他們有領受和幫助.
聽下這次的基督教課.
大奉主的德性的名.
阿們.
\newpage



\section{列王記上 8:1-66}
\label{sec:032347}
\textbf{萬國禱告的殿:悔改與復興的真正道路}~~~~袁惠鈞牧師博士
\newline
\newline
連結: \href{http://www.church.com.hk/acms/content.asp?site=cdc\&op=show\&type=product\&code=032347\&layout=sermon}{\texttt{032347}} ~~~~ 語音日期: 2026-06-21
\newline
\newline
\hyperref[sec:032346]{\small{< < < PREV SERMON < < <}}
~
\hyperref[sec:index_chronic]{\small{[返順時目]}}
~
\hyperref[sec:032348]{\small{> > > NEXT SERMON > > >}}
\newline
\newline
列王記上 8:1-66
\newline
\begin{longtable}{cl}
\hline
\hline
章節 & 經文 (和合本修訂版)\\
\hline
8:1 & \begin{tabularx}{0.7\textwidth}{X} 那時，所羅門召集以色列的長老、各支派的領袖和以色列人的族長到耶路撒冷，所羅門王那裡，要把耶和華的約櫃從大衛城，就是錫安，接上來。 \end{tabularx} \\ \\ \relax
8:2 & \begin{tabularx}{0.7\textwidth}{X} 以他念月，就是七月，在節期時，所有的以色列人都聚集到所羅門王那裡。 \end{tabularx} \\ \\ \relax
8:3 & \begin{tabularx}{0.7\textwidth}{X} 以色列眾長老一來到，祭司就抬起約櫃。 \end{tabularx} \\ \\ \relax
8:4 & \begin{tabularx}{0.7\textwidth}{X} 祭司和利未人將耶和華的約櫃請上來，又把會幕和會幕一切的聖器皿都帶上來。 \end{tabularx} \\ \\ \relax
8:5 & \begin{tabularx}{0.7\textwidth}{X} 所羅門王和聚集到他那裡的以色列全會眾一同在約櫃前獻牛羊為祭，多得不可勝數，無法計算。 \end{tabularx} \\ \\ \relax
8:6 & \begin{tabularx}{0.7\textwidth}{X} 祭司將耶和華的約櫃請進內殿，就是至聖所，安置在兩個基路伯的翅膀底下約櫃的地方。 \end{tabularx} \\ \\ \relax
8:7 & \begin{tabularx}{0.7\textwidth}{X} 基路伯張開翅膀在約櫃上面的地方，從上面遮住約櫃和抬櫃的槓。 \end{tabularx} \\ \\ \relax
8:8 & \begin{tabularx}{0.7\textwidth}{X} 這槓很長，從內殿前的聖所可以看見槓頭，從外面卻看不見。這槓直到今日還在那裡。 \end{tabularx} \\ \\ \relax
8:9 & \begin{tabularx}{0.7\textwidth}{X} 約櫃裡沒有別的，只有兩塊石版，就是以色列人出埃及地，耶和華與他們立約的時候，摩西在何烈山放在那裡的。 \end{tabularx} \\ \\ \relax
8:10 & \begin{tabularx}{0.7\textwidth}{X} 祭司從聖所出來的時候，有雲充滿耶和華的殿， \end{tabularx} \\ \\ \relax
8:11 & \begin{tabularx}{0.7\textwidth}{X} 祭司因雲彩的緣故不能站立供職，因為耶和華的榮光充滿了耶和華的殿。 \end{tabularx} \\ \\ \relax
8:12 & \begin{tabularx}{0.7\textwidth}{X} 那時，所羅門說：「耶和華曾說要住在幽暗之處。 \end{tabularx} \\ \\ \relax
8:13 & \begin{tabularx}{0.7\textwidth}{X} 我的確為你建了一座雄偉的殿宇，作為你永遠居住的地方。」 \end{tabularx} \\ \\ \relax
8:14 & \begin{tabularx}{0.7\textwidth}{X} 王轉過臉來為以色列全會眾祝福，以色列全會眾都站立。 \end{tabularx} \\ \\ \relax
8:15 & \begin{tabularx}{0.7\textwidth}{X} 所羅門說：「耶和華－以色列的神是應當稱頌的！因他親口向我父大衛應許的，也親手成就了；他曾說： \end{tabularx} \\ \\ \relax
8:16 & \begin{tabularx}{0.7\textwidth}{X} 『自從那日我領我百姓以色列出埃及以來，我未曾在以色列各支派中選擇一城，在那裡為我的名建造殿宇，但我揀選大衛治理我的百姓以色列。』 \end{tabularx} \\ \\ \relax
8:17 & \begin{tabularx}{0.7\textwidth}{X} 我父大衛的心意是要為耶和華－以色列神的名建殿。 \end{tabularx} \\ \\ \relax
8:18 & \begin{tabularx}{0.7\textwidth}{X} 耶和華卻對我父大衛說：『你有心為我的名建殿，這心意是好的； \end{tabularx} \\ \\ \relax
8:19 & \begin{tabularx}{0.7\textwidth}{X} 但你不可建殿，惟有你親生的兒子才可為我的名建殿。』 \end{tabularx} \\ \\ \relax
8:20 & \begin{tabularx}{0.7\textwidth}{X} 現在耶和華實現了他所應許的話，使我接續我父大衛坐以色列的王位，正如耶和華所說的，我也為耶和華－以色列神的名建造了這殿。 \end{tabularx} \\ \\ \relax
8:21 & \begin{tabularx}{0.7\textwidth}{X} 我也在那裡為約櫃預備一處。約櫃那裡有耶和華的約，就是他領我們列祖出埃及地的時候，與他們所立的約。」 \end{tabularx} \\ \\ \relax
8:22 & \begin{tabularx}{0.7\textwidth}{X} 所羅門當著以色列全會眾，站在耶和華的壇前，向天舉手， \end{tabularx} \\ \\ \relax
8:23 & \begin{tabularx}{0.7\textwidth}{X} 說：「耶和華－以色列的神啊，天上地下沒有神明可與你相比！你向那些盡心行在你面前的僕人守約施慈愛， \end{tabularx} \\ \\ \relax
8:24 & \begin{tabularx}{0.7\textwidth}{X} 這約是你向你僕人大衛守的，是你應許他的。你親口應許，親手成就，正如今日一樣。 \end{tabularx} \\ \\ \relax
8:25 & \begin{tabularx}{0.7\textwidth}{X} 耶和華－以色列的神啊，你向你僕人我父大衛應許說：『你的子孫若謹慎自己的行為，在我面前行事像你所行的一樣，就不斷有人在我面前坐以色列的王位。』現在求你信守這話。 \end{tabularx} \\ \\ \relax
8:26 & \begin{tabularx}{0.7\textwidth}{X} 以色列的神啊，現在求你成就向你僕人我父大衛所應許的話。 \end{tabularx} \\ \\ \relax
8:27 & \begin{tabularx}{0.7\textwidth}{X} 「神果真住在地上嗎？看哪，天和天上的天尚且不足容納你，何況我所建的這殿呢？ \end{tabularx} \\ \\ \relax
8:28 & \begin{tabularx}{0.7\textwidth}{X} 惟求耶和華－我的神垂顧僕人的禱告祈求，俯聽僕人今日在你面前的祈禱呼求。 \end{tabularx} \\ \\ \relax
8:29 & \begin{tabularx}{0.7\textwidth}{X} 願你的眼目晝夜看顧這殿，就是你說要作為你名的居所；求你垂聽禱告，你僕人向此處的禱告。 \end{tabularx} \\ \\ \relax
8:30 & \begin{tabularx}{0.7\textwidth}{X} 你僕人和你百姓以色列向此處祈禱的時候，求你在你天上的居所垂聽，垂聽而赦免。 \end{tabularx} \\ \\ \relax
8:31 & \begin{tabularx}{0.7\textwidth}{X} 「人若得罪鄰舍，有人強迫他，要他起誓，他來到這殿，在你的壇前起誓， \end{tabularx} \\ \\ \relax
8:32 & \begin{tabularx}{0.7\textwidth}{X} 求你在天上垂聽、處理，向你的僕人施行審判，定惡人有罪，照他所行的報應在他頭上；定義人為義，照他的義賞賜他。 \end{tabularx} \\ \\ \relax
8:33 & \begin{tabularx}{0.7\textwidth}{X} 「你的百姓以色列若得罪你，敗在仇敵面前，卻又歸向你，宣認你的名，在這殿裡向你祈求禱告， \end{tabularx} \\ \\ \relax
8:34 & \begin{tabularx}{0.7\textwidth}{X} 求你在天上垂聽，赦免你百姓以色列的罪，使他們歸回你賜給他們列祖的地。 \end{tabularx} \\ \\ \relax
8:35 & \begin{tabularx}{0.7\textwidth}{X} 「你的百姓若得罪了你，你使天閉塞不下雨；他們若向此處禱告，宣認你的名，因你的懲罰而離開他們的罪， \end{tabularx} \\ \\ \relax
8:36 & \begin{tabularx}{0.7\textwidth}{X} 求你在天上垂聽，赦免你僕人你百姓以色列的罪，將當行的善道教導他們，並降雨在你的地，就是你賜給你百姓為業之地。 \end{tabularx} \\ \\ \relax
8:37 & \begin{tabularx}{0.7\textwidth}{X} 「這地若有饑荒、瘟疫、焚風、霉爛、蝗蟲、螞蚱，或有仇敵圍困這地的城門，無論遭遇甚麼災禍疾病， \end{tabularx} \\ \\ \relax
8:38 & \begin{tabularx}{0.7\textwidth}{X} 你的百姓以色列，或眾人或一人，內心知道有禍，向這殿舉手，無論祈求甚麼，禱告甚麼， \end{tabularx} \\ \\ \relax
8:39 & \begin{tabularx}{0.7\textwidth}{X} 求你在天上你的居所垂聽、赦免、處理。因為你知道人心，惟有你知道世人的心，求你照各人所行的一切待他們， \end{tabularx} \\ \\ \relax
8:40 & \begin{tabularx}{0.7\textwidth}{X} 使他們在你賜給我們列祖的土地上一生一世敬畏你。 \end{tabularx} \\ \\ \relax
8:41 & \begin{tabularx}{0.7\textwidth}{X} 「論到不屬你百姓以色列的外邦人，若為你的名從遠方而來， \end{tabularx} \\ \\ \relax
8:42 & \begin{tabularx}{0.7\textwidth}{X} 他們因聽見你的大名和大能的手，以及伸出來的膀臂，來向這殿禱告， \end{tabularx} \\ \\ \relax
8:43 & \begin{tabularx}{0.7\textwidth}{X} 求你在天上你的居所垂聽，照著外邦人向你所求的一切而行，使地上萬民都認識你的名，敬畏你，像你的百姓以色列一樣，又使他們知道我所建造的是稱為你名下的殿。 \end{tabularx} \\ \\ \relax
8:44 & \begin{tabularx}{0.7\textwidth}{X} 「你的百姓若奉你的派遣出去，無論往何處與仇敵爭戰，他們若向耶和華所選擇的城，以及我為你名所建造的這殿禱告， \end{tabularx} \\ \\ \relax
8:45 & \begin{tabularx}{0.7\textwidth}{X} 求你在天上垂聽他們的禱告祈求，為他們伸張正義。 \end{tabularx} \\ \\ \relax
8:46 & \begin{tabularx}{0.7\textwidth}{X} 「你的百姓若得罪你，因為沒有人不犯罪，你向他們發怒，把他們交在仇敵面前，擄他們的人把他們帶到仇敵之地，或遠或近， \end{tabularx} \\ \\ \relax
8:47 & \begin{tabularx}{0.7\textwidth}{X} 他們若在被擄之地那裡回心轉意，在擄掠者之地悔改，向你懇求說：『我們有罪了，我們悖逆了，我們作惡了』； \end{tabularx} \\ \\ \relax
8:48 & \begin{tabularx}{0.7\textwidth}{X} 他們若在擄他們的仇敵之地盡心盡性歸向你，又向自己的地，就是你賜給他們列祖的地和你所選擇的城，以及我為你名所建造的這殿禱告， \end{tabularx} \\ \\ \relax
8:49 & \begin{tabularx}{0.7\textwidth}{X} 求你在天上你的居所垂聽他們的禱告祈求，為他們伸張正義， \end{tabularx} \\ \\ \relax
8:50 & \begin{tabularx}{0.7\textwidth}{X} 饒恕得罪你的子民，赦免他們向你所犯一切的過犯，使他們在擄他們的人面前蒙憐憫。 \end{tabularx} \\ \\ \relax
8:51 & \begin{tabularx}{0.7\textwidth}{X} 因為他們是你的子民，你的產業，是你從埃及，從鐵爐中領出來的。 \end{tabularx} \\ \\ \relax
8:52 & \begin{tabularx}{0.7\textwidth}{X} 願你的眼目看顧僕人和你百姓以色列的祈求；他們無論何時向你呼求，願你垂聽。 \end{tabularx} \\ \\ \relax
8:53 & \begin{tabularx}{0.7\textwidth}{X} 主耶和華啊，你將他們從地上萬民中分別出來作你的產業，是照著你領我們列祖出埃及的時候，藉你僕人摩西所應許的。」 \end{tabularx} \\ \\ \relax
8:54 & \begin{tabularx}{0.7\textwidth}{X} 所羅門在耶和華的壇前屈膝跪著，向天舉手；他在耶和華面前禱告祈求完畢的時候，就起來， \end{tabularx} \\ \\ \relax
8:55 & \begin{tabularx}{0.7\textwidth}{X} 站著，大聲為以色列全會眾祝福，說： \end{tabularx} \\ \\ \relax
8:56 & \begin{tabularx}{0.7\textwidth}{X} 「耶和華是應當稱頌的！因為他照著一切所應許的賜平安給他的百姓以色列，凡藉他僕人摩西應許賜福的話，一句都沒有落空。 \end{tabularx} \\ \\ \relax
8:57 & \begin{tabularx}{0.7\textwidth}{X} 願耶和華－我們的神與我們同在，像與我們列祖同在一樣，不撇下我們，不丟棄我們， \end{tabularx} \\ \\ \relax
8:58 & \begin{tabularx}{0.7\textwidth}{X} 使我們的心歸向他，遵行他一切的道，謹守他吩咐我們列祖的誡命、律例、典章。 \end{tabularx} \\ \\ \relax
8:59 & \begin{tabularx}{0.7\textwidth}{X} 願我在耶和華面前祈求的這些話，晝夜靠近耶和華－我們的神，好讓他每日為他僕人和他百姓以色列伸張正義， \end{tabularx} \\ \\ \relax
8:60 & \begin{tabularx}{0.7\textwidth}{X} 使地上的萬民都知道惟獨耶和華是神，沒有別的了。 \end{tabularx} \\ \\ \relax
8:61 & \begin{tabularx}{0.7\textwidth}{X} 所以你們當向耶和華－我們的神存純正的心，遵行他的律例，謹守他的誡命，如同今日一樣。」 \end{tabularx} \\ \\ \relax
8:62 & \begin{tabularx}{0.7\textwidth}{X} 王和全以色列一同在耶和華面前獻祭。 \end{tabularx} \\ \\ \relax
8:63 & \begin{tabularx}{0.7\textwidth}{X} 所羅門向耶和華獻平安祭，二萬二千頭牛，十二萬隻羊。這樣，王和全以色列為耶和華的殿行了奉獻之禮。 \end{tabularx} \\ \\ \relax
8:64 & \begin{tabularx}{0.7\textwidth}{X} 當日，王因耶和華殿前的銅壇太小，容不下燔祭、素祭和平安祭牲的脂肪，就將耶和華殿前院子的中間分別為聖，在那裡獻燔祭、素祭和平安祭牲的脂肪。 \end{tabularx} \\ \\ \relax
8:65 & \begin{tabularx}{0.7\textwidth}{X} 那時所羅門守節，從哈馬口直到埃及溪谷的以色列眾人都與他同在一起，成了一個盛大的會，在耶和華－我們的神面前七日又七日，共十四日。 \end{tabularx} \\ \\ \relax
8:66 & \begin{tabularx}{0.7\textwidth}{X} 第八日，王遣散百姓；他們都為王祝福。他們為耶和華向他僕人大衛和他百姓以色列所施的一切恩惠都心中喜樂，愉快地各回自己的帳棚去了。 \end{tabularx} \\ \\
[1ex]
\hline
\hline
\end{longtable}
$^{1}$列王記合一王朝系列的第八章.
我們上次已經開始了列王記上第八章.
但因為有66節經文.
所以我說不完.
我們今天就會結束列王記上的第八章.
我給的題目是.
萬國禱告的殿.
悔改如復興真正的道路.
我們說所羅門建了聖殿.
他將聖殿獻給神.
他鼓勵以色列人在殿內甚至向著殿禱告.
尤其是當他們犯了罪.
得罪了神之後.
求神來赦免他們.
以至神能夠復興他們.
這就是我們今天要說的重點.
我們也說過.
其實列王記上的第八章.
神的名是非常重要.
聖殿就是神的名所居住的地方.
神是要以色列人在他納祂名的地方.
來敬拜自己.
向自己祈求.
列王記上的第八章.
我們也已經說過.
其實整張書都是一個首尾結構.
尤其是在經文最中心的地方.
所羅門是說出了七個祈求出來.
這七個祈求本身都是一個首尾結構.
一共有四層.
第一層就是所羅門說.
無論我們是得罪了我們的鄰居.
在列王記上的八章.
31至32節.
或者是得罪了神你自己.
46至51節.
有沒有看到這個首尾結構.
得罪了人 得罪了神.
我們得罪了我們的鄰居.
我們就得罪了神.

$^{41}$第二層.
無論當敵人來攻擊以色列人.
因為他們犯罪.
他們打敗仗在33至34節.
或者他們去攻打他們的敵人.
又是打敗仗因為他們的罪.
44至45節.
求神垂聽他們的禱告.
第三層.
無論是以色列民神的民.
或者是外邦人不屬神的子民.
向祂禱告 求神垂聽.
其實求神在天上垂聽.
一共出現七次.
我說七個祈求.
每一個祈求.
所羅門都說.
求你在天上垂聽.
而經文最中心的地方.
就是在列王記上的八章.
37至40節.
是一些一般天然的災難.
一些天然的疾病.
當他們不聽話.
被神管教的時候.
所羅門說.
當以色列人向你認罪悔改.
求你赦免他們.
我上一次已經和大家開始了.
其中兩個祈求.
如果你記得的話.
我已經說了.
以色列人如果得罪他們的鄰居.
如果又因為他們的罪.
被敵人打敗的時候.
求神憐憫.
今天我就要開始說第三個祈求.
就是當以色列人得罪了他們的神.
被神懲罰.
其實.

$^{81}$神和亞伯拉罕所立的約.
就給了應許之地給以色列人.
這個應許之地是無條件的.
不是說以色列人做得不好.
他們就不能夠得到這個應許之地.
是無條件給以色列人.
但是在這個應許之地上的祝福.
是和以色列人的行為息息相關.
當他們毀約的時候.
神會管教.
但當他們服從的時候.
神就會祝福他們.
而其中一個最嚴重的詛咒.
最嚴重的管教.
就是旱災.
天閉塞不下雨.
為什麼這麼多天然的災難.
我們稍後會見到.
在經文最中心的地方.
是所有天然的災難.
為什麼現在又會神.
突然間拿出旱災.
來懲罰以色列人.
為什麼這麼重要.
因為以色列人是耕種的民族.
和中國古代都是一樣.
如果你是一個耕種的民族.
雨水就非常重要.
所以當以色列人馴服.
神就給他們.
風調雨順.
國泰民安.
五穀豐收.
牛羊眾多.
但若果他們不服從.
神就給他們旱災.
來警告他們.
管教他們.
列王記上的十八章.
以利亞先知的時候.

$^{121}$我們會見到.
就是一個例子.
神閉塞天不下雨.
一共有三年這麼長.
左羅門第四個祈求.
就是我剛才所說.
其他天然的災難.
其實神在利未記的二十六章.
和申命記的二十八章.
都講得很清楚.
當以色列人不服從的時候.
可能會有饑荒.
有枯萎病.
有昆蟲入侵.
有各種疾病,瘟疫.
但當他們服從神.
他們能夠避免這一切的災難.
左羅門在這裡求神.
如果神真的管教以色列人.
如果以色列人知錯了.
認罪悔改.
祂求神憐憫.
復興以色列人.
讓他們能夠安居樂業.
在所羅門的禱告裡.
「那片土地」這個字重複出現.
那片土地是什麼?.
就是「應許之地」.
這是以色列人從神繼承的產業之一的一部分.
當以色列人開始犯罪.
神就降災在這一片土地的上方.
如果以色列人再不聽話.
執迷不悟.
繼續叛逆,忤逆神的時候.
神甚至會准許他們的敵人來攻打他們.
將他們從這一片土地擄去.
我們知道以色列人的歷史.
在公元前的722年.
神容許亞述人征服北國以色列.
將以色列國的人民擄走.

$^{161}$在公元前的586年.
神再次容許巴比倫人.
來征服南國猶大.
也將人民擄走.
聖殿被毀.
耶路撒冷被燒毀.
但是就是因為神這樣懲罰以色列人.
終於治好了他們拜偶像的病.
將他們擄走之後.
他們就不再去拜偶像.
第五個所羅門的祈求.
就是那些不屬神的人民.
那些外方人.
當他們向神祈求的時候.
所羅門求神.
垂聽他們的禱告.
這些不是以色列人進入迦南地的時候.
所征服的敵人.
這些是享有律法特權的外方人.
他們可能住在以色列人當中.
也可能是因為冒神的名.
來到聖殿敬拜.
來向神禱告.
他們也可能是起聖殿其中的外方人.
記不記得.
所羅門曾經是給很多外方人來起聖殿.
其實以色列人是有責任.
在這些外方人的面前彰顯神的榮耀.
神要他們成為外方人的光.
所以所羅門在《樹列王記》上的八章.
四十三節和六十節.
所羅門求神垂聽外方人的禱告.
祂說:使天下萬民都認識你的名.
敬畏你像你的民以色列一樣.
又使他們知道我建造的這殿.
是稱為你名下的殿.
所羅門可能認為.
當這些外方人向神禱告的時候.
他們就是邁出了第一步.
來信靠神或者是敬拜神.

$^{201}$其實早在神父趙亞伯拉罕.
離開吾爾的時候.
神已經告訴以色列人.
你們要成為萬國的祝福.
就是在《創世記》的十二章一至三節.
神在外方人的面前.
不斷彰顯祂自己的榮耀.
記不記得以色列人在埃及的時候.
神降下十災給埃及.
神最後打開了紅海給以色列人經過.
神就在埃及人的面前.
讓世人見到祂的大能.
見到祂的榮耀.
記不記得在《約書亞記》裡面.
第二章有一個妓女叫做拉罕.
她聽見神在外方人的面前.
怎樣彰顯祂的榮耀.
她相信了.
她願意接受這位以色列人的神.
後來神又打開約旦河.
讓以色列人經過進入迦南地.
神甚至幫助約書亞.
戰敗所有迦南地的國家.
也在外方人的面前.
彰顯祂自己的力量.
彰顯祂自己的榮耀.
還有大衛.
他戰勝哥利亞.
哥利亞是甚麼人?.
是非利士人.
又是外方人.
有沒有見到神不斷在外方人的面前.
彰顯祂的榮耀.
神叫以色列人要成為萬國的祝福.
這個也是給我們的命令.
我們都要成為萬國的祝福.
就好像詩篇的六十七篇一至二節所說.
願神憐憫我們.
賜福給我們.
用念光照我們.

$^{241}$使地上得知你的道路.
萬國得知你的救恩.
我們都要令到萬國得知神的救恩.
這是甚麼?.
這是大使命.
很多人認為.
得救了之後就行了.
我自己得救了.
不是的.
如果你自己得救.
那就算神隨時可以接走你上天堂.
神留我們在這裡祭上做甚麼?.
就是要成為萬國的祝福.
就是要將這個好消息傳給.
未聽過福音的人.
要執行大使命.
還有第六個所羅門的祈求.
就是當以色列人.
去與仇敵真戰打敗仗.
求神隨聽他們的禱告.
我說過七個祈求都是首尾結構.
這個其實和上面那個.
被敵人打敗的時候.
是有一個首尾的呼應.
上面說當仇敵來攻打他們.
而現在這個經文.
當他們去攻打他們的仇敵.
有沒有看到是一來一去.
首尾呼應.
這裡其實是說.
當神差派祂的子民去真戰.
即是說打一場聖戰的時候.
要以色列人靠著神的力量.
祂的智慧去得勝.
祭司打仗的時候就會吹響號角.
甚至會幫助軍隊明白神的旨意.
祭司都要鼓勵士兵來信靠神.
而在真戰當中.
如果士兵打敗仗.
或者是有危難的時候.

$^{281}$他們也可以向聖殿來禱告.
向神來禱告.
保羅在新約教我們要穿上全副軍裝.
其中禱告是非常重要的武器.
在《說耳忽》所書的六章第十八和第十九節.
保羅說我們要拿起聖靈的寶劍.
就是神的話語.
還有靠著聖靈隨時多方祈求.
有沒有看到.
無論舊約新約禱告都是非常重要的武器.
有人說神永遠站在最強大軍隊的那一邊.
當然是錯的.
其實是甚麼.
其實神永遠站在按照祂的旨意祈求的那一邊.
其實我也應該再具體一點說.
不是神站在你那邊.
而是向神祈求.
按照祂的旨意祈求的人站在神那邊.
這個才是正確.
還有第七點.
第七個祈求.
就是當神的子民得罪了神.
在列王記上第八章第四十六節.
經文說「他們」這個代名詞.
其實「他們」就是指以色列人.
在以色列人的歷史當中.
我們見到以色列人頻頻犯罪.
而經文亦說甚麼.
因為沒有人不犯罪.
今天的經文在舊約.
很多人以為保羅所說的.
眾人都犯了罪.
是新約才有.
其實不是.
早在舊約已經告訴我們.
沒有人不犯罪.
這個亦包括了所羅門.
所以他沒有辦法逃彼得罪耶和華.
就是因為所羅門跟隨他的妻妾去拜偶像.
以致導致以色列國分為南北國.

$^{321}$但是神非常有祂的恩典.
因為神記得和大衛所納的約.
祂姑息所羅門.
所以沒有在所羅門在位在世的時候.
使得以色列國分為南北國.
是所羅門去世之後才發生.
以賽亞先知和尼加先知都有預言.
巴比倫將會征服以色列的預言.
早在以賽亞先知和尼加先知的時候.
已經有預言.
而耶利未先知亦有預言.
但是這個預言更加具體.
他直接說以色列人將會被擄七十年.
你說為何七十年?.
原來以色列人有四百九十年.
沒有守神的安息年.
我們知道每七年就是一個安息年.
你要不耕種給一塊地來安息.
但是以色列人很貪心.
六年之後為何第七年不去耕種.
這樣不是徹底嗎?.
所以他們繼續耕種.
有四百九十年沒有守安息年.
四百九十年有多少個安息年?.
每七年就是一個.
四百九十除七就是七十.
所以神一次過將安息年收回來.
就是要以色列人被擄七十年.
當他們被擄的時候離開了應許之地.
應許之地就能夠得到安息七十年.
所以神所說的話就是算數.
但以利先知當他明白到.
耶利未先知的預言.
他就開始向神祈求.
希望神早些實驗他的應許.
以致以色列人能夠早些回到自己的故鄉.
他們的應許之地.
最後神就是激勵波斯王居魯士.
准許猶太人回到以色列地.
重建他們的聖殿.

$^{361}$所羅門告訴我們.
在聶皇記上的八章五十一節.
為何當神的子民願意認罪悔改.
歸向神的時候.
神一定會寬恕他們.
為甚麼?.
因為這是神的應許.
因為以色列人是神的子民.
神應許當他們認罪悔改.
神一定會寬恕他們.
是神當年帶他們出埃及地.
圍爐之家的地.
將他們釋放.
神是信實的.
在申命記的第三十章第二節.
神說如果以色列人盡心盡性.
歸向耶和華他們的神.
祂是會寬恕他們.
就算他們被擄離開了應許之地.
神說我都會讓你們回到你們的本國.
當所羅門結束他的禱告的時候.
他求神垂聽所有向聖殿禱告的人.
其實在經文裡面有一個首尾結構.
如果你注意到在聶皇記上的.
第八章二十二和八章五十四節.
無論所羅門禱告的姿勢.
無論他所講的話.
都是一模一樣.
而在聶皇記上的第八章.
和歷代志下的第六章.
和詩篇一百三十二篇的經文.
有很多相似的地方.
如果大家讀詩篇的時候.
可能會留意得到.
還有一樣很有趣.
所羅門禱告結束的時候.
求神不要厭棄受告者.
祂說神啊.
你不要忘記了受告者.
你要實驗你的應許.

$^{401}$你說你會差派你的受告者來到世上.
這就是神和大衛所納的約.
大衛之約講到什麼呢?.
將會有一位大衛的後人.
將會來到世上.
這個就是受告者.
是誰呢?.
是耶穌基督.
不是所羅門.
所以你看到舊約新約全部都是一致.
你讀聖經真的讀得很開心.
因為你發覺.
一位的聖靈在背後寫出聖經.
以至無論新約舊約都是一致.
大衛每一個的子孫.
都要為自己的行為來負責.
連所羅門也不例外.
我們知道所羅門最後.
他是被他的妻妾引誘去扮偶像.
所以很多人認為所羅門是不得救.
但是我認為他最後都要歸向神.
為什麼我這樣說呢?.
因為所羅門在聖經裡面.
作了三部書.
《傳道書》,《箴言》和《牙歌》.
說起來.
現在我們的周長老在說.
《傳道書》,《主日學》.
非常精彩.
所以我勸大家留步.
去聽周長老教《傳道書》.
我講了這麼久.
才來到今天的大綱.
大件事.
要講很久.
記不記得我上次已經講了三點.
我講了神榮耀的殿.
八章一至十一節.
我也講了這個聖殿是神與以色列人同在的殿.
十二至二十一節.

$^{441}$我也講過這個聖殿是向神祈求的殿.
二十二節到五十三節.
我今天要講另外那三點是什麼?.
這個聖殿都是神應許的殿.
在五十四節到六十一節.
這個殿也是以色列人一起慶祝的殿.
六十二節到六十六節.
第三點就是神話語的殿.
九章一至九節.
有人說牧師你又說第八章.
為什麼講到第九章?.
其實聖經的原文是沒有分章節.
是後人加上去.
我相信九章一至九節.
其實是應該與八章的經文放在一起.
有沒有看到這裡有一個首尾結構.
神榮耀的殿與神話語的殿是一致的.
神的榮耀與祂的話語是分不開的.
這是一個首尾呼應.
第二個首尾結構.
這個是神與祂子民同在的殿.
但是也是祂子民一起慶祝的殿.
有沒有看到這兩個也是有一個首尾呼應.
還有最中心的經文是講什麼?.
這個是向神祈求的殿.
但是也是神應許禱告的殿.
這個也是一個首尾結構.
而且在經文最中心的地方.
而經文最重要.
就是所羅門提出的那七個祈求.
就是我們剛才已經講了.
所以我們現在要講的.
這個殿就是神應許的殿.
還有上次我曾經講過.
這個是神榮耀的殿.
是一個首尾的呼應.
然後所羅門就替神向人民祝福.
你說什麼王可以祝福人的嗎?.
不是祭司才做的嗎?.
通常是祭司祝福人民.

$^{481}$但是在特殊的場合.
君王都可以祝福人民.
是有先例的.
如果我們記得大衛在《撒母耳夏》的第六章.
十八節和二十節.
都有同樣祝福人民.
所羅門回顧猶太人的歷史.
他有一個結論.
這個結論是什麼?.
神的應許一句都沒有落空.
神的指紋.
以色列人時時都令到神失望.
但是神從來沒有令到以色列人失望.
那你說.
神給以色列人有什麼的應許?.
實驗了嗎?.
起碼有三樣.
第一就是應許之地.
應許之地就是迦南地.
流乃與物之地.
神說你們要去那裡.
神就的確給約書亞.
有力量能夠戰勝迦南人.
給以色列人進入迦南地.
第二樣神應許以色列人是有安息.
在摩西的時候.
神已經承諾這個國家將會有一段時間有安息.
出埃及記的三十三章十四節.
約書亞記裡面都有講到.
而在所羅門的統治之下.
以色列人的確是有短暫的安息.
那你說這個安息要來做什麼?.
就是應驗了神第三個應許.
起聖殿.
神曾經應許以色列人.
在一個納祂名的地方來敬拜自己.
真命記的十二章.
現在所羅門真的起好聖殿.
以色列人不單止可以在那裡敬拜.
還可以在那裡向神祈求.

$^{521}$約書亞在他告別領袖的時候.
告別以色列人的時候.
在約書亞記的二十三章十四至十五節.
都有講同樣的話.
就是說神的應許沒有一件落空.
但是約書亞也提醒以色列人.
神的應許不單止只有祝福.
神的應許也有詛咒.
而祝福和詛咒的應許.
只不過是一善之差.
很多時候以色列人.
希望得到神的祝福.
但是最後他們都是離棄神.
而被神來詛咒.
他們被趕出應許之地.
就是一個好的例子.
那你說神還有什麼應許?.
有一個很重要的應許.
就是神應許以色列人和他們同在.
這個應許不單止是給亞伯拉罕.
也是給以撒,雅各,摩西.
約書亞,基澱,撒母耳,大衛,所羅門.
神不斷和以色列人說.
我會與你們同在.
太多經文我不會一一讀出來.
我會留在PowerPoint那裡.
而我們也要知道.
這些經文不單止說.
神會和以色列人同在.
有些經文也說.
神不會撇下以色列人.
神不會刁棄以色列人.
所羅門之後.
以賽亞先知也有重申這個應許.
就是去安慰那些被巴比倫人.
所擄去的猶太人.
神透過以賽亞書來鼓勵他們.
神也有用同樣的話來鼓勵耶利米先知.
在耶穌基督升天之前.
也有用同樣的話來安慰門徒.

$^{561}$就是我們很熟悉的大使命.
馬太福音的二十八章.
「去使萬民作我的門徒.
奉父子聖靈的名為他們施洗.
凡我教你們的都要教他們遵守」.
「我尚與你們同在直到世界的末了」.
所以這個應許也是給教會.
如果要神與我們同在的話.
我們一樣都是要去執行大使命.
而不是說接受了神之後.
就翹起雙手坐在教會.
希望被人服侍.
而是我們真的要出去.
落手落腳來執行大使命.
所羅門跟著祈求神.
求神能夠給以色列人.
一個歸向神.
遵行祂道的心.
有這樣的心.
歸向神的心.
所羅門是熟讀《申命記》.
他知道《申命記》的五章二十九節所說.
「神或敬畏我的.
常遵守我一切誡明的.
我會使他們和他們的子孫永遠特服」.
這一節的經文其實跟《列王記》上的八章.
六十一節是很相似.
所羅門要以色列人全誠實的心.
遵行神的律例.
謹守祂的誡明.
最後.
所羅門求神紀念祂口中.
和心中的禱告.
《列王記》上的八章五十九至六十節.
科學家告訴我們.
我們說出來的話是什麼.
是聲波.
我現在說出來.
你們就聽到.
如果我說完之後.

$^{601}$聲波就會消失.
在空中消失.
因為音波一直消失.
但是.
我們的禱告去到神那裡.
會不會消失?會不會不見了?.
絕對不會.
神是記得我們的禱告.
神是會回應我們的禱告.
一定會回應.
但是未必是照著我們的時間.
未必是照著我們的方式.
來回應我們的禱告.
但是祂永遠會比我們最好.
所羅門向神的禱告一點都不自私.
祂說希望以色列人能夠忠於神.
祂希望地上的國家都能夠認識神.
信靠神.
就是以色列人的神.
祂希望這個聖殿能夠成為萬國禱告的殿.
這個就是我們今天的主題.
萬國禱告的殿.
神的確是回應了.
所羅門的請求.
我們怎麼知道神應許了所羅門的請求?.
當時有一把火從天而降.
是將所有的祭物完全燒了.
然後神的榮耀就充滿了聖殿.
在歷代志下的七章一至三節.
其實在以色列人的歷史當中.
同樣的事都有發生過.
記不記得亞倫祝福百姓的時候.
在利未記的九章二十三至二十四節.
都是一把火從天而降燒了所有的祭物.
另外以利亞先知.
遲些我們會看這件事.
當他和巴力先知鬥法的時候.
神怎樣?.
降下火.
將以利亞先知的祭物全部燒了.

$^{641}$巴力先知的祭物原封不動.
如果我們記得的話.
其實以利亞先知的祭物還在淋水.
都能夠一把火將他們燒了.
現在百姓見到在這裡壇上的祭物.
一把火完全燒了.
他們怎樣?.
立刻俯伏在地上.
他們同聲說.
「耶和華本為善,他的慈愛永遠長存」.
歷代志下的七章第三節.
能不能夠想像成千上萬的.
以色列人同聲來讚美神.
讓我們見到神完全接受所羅門的祈求.
一把火燒了祭物.
我們說完了第一點.
就是這個聖殿是神應許的殿.
第二點我們要說的.
就是這個殿是以色列人一同慶祝的殿.
這個我又要再提一提大家.
還有上高我們曾經說過.
聖殿是神與以色列人同在的殿.
是有一個首尾結構.
這裡就是一同慶祝的殿.
所有的會眾.
從南邊埃及的小河或者是溪谷.
到北面的哈瑪口.
很多的以色列人都來到一起獻這個殿給神.
帶了很多的祭物.
只是所羅門自己的祭物有多少?.
二萬二千隻牛.
十二萬隻綿羊和山羊.
你想想有多少祭物.
所以經文告訴我們.
所羅門新建的銅祭壇.
根本沒有辦法可以獻這麼多的祭物.
於是所羅門就將聖殿前面的院子的中間部分.
將它分別為聖.
然後在那裡獻梵祭,掃祭和平安祭.
注意「平安祭」這個字.

$^{681}$「平安」這個字.
就是所羅門的名字.
Shalom.
就是有連帶關係的.
獻殿禮我們以前說過.
是在鑄盤節來舉行.
鑄盤節是紀念以色列人昔日在曠野.
當他們在曠野貧走的時候.
神以雲柱,煙柱來帶領他們.
他們住在帳盤那裡.
就是紀念他們住在帳盤.
有神與他們同在來照顧他們.
現在聖殿已經建好了.
而且他們將祭物獻給神.
當中是有平安祭或者和平祭.
平安祭是燒給神.
但是其中有一部分的肉是交給祭司.
另外一部分是交給奉獻者.
以致他可以和他的家人和他的朋友來分享.
猶太人是有養很多的動物.
但是通常他們要這些動物.
是要來拿牛奶,羊毛.
或者是使他們的身畜能夠繁殖.
而不是要來吃的.
只有當他們獻祭給神.
是可以拿回一部分.
他們就會和家人和朋友一起來慶祝.
一起來享受.
這個獻殿禮舉行了一個星期.
另外一個星期是一個盛宴.
以色列人和中國人很相似.
我們中國人都很喜歡吃東西.
好像我們敬拜完時時都一起去飯堂吃飯.
以色列人都是一樣.
獻完聖殿又怎樣?.
一個星期有盛宴.
然後最後一天就是嚴肅會.
來結束這個獻殿禮.
這個平安祭我們要講講.
平安祭的律例.

$^{721}$就是你只可以吃兩天.
我們剛才說和家人來分享.
去到第三天你吃剩的要全部燒掉.
為什麼?.
其實這裡有一個耶穌基督的預表.
祂就是我們的平安祭.
預表耶穌基督第三天復活.
所以祭肉你只能夠吃兩天.
第三天你吃不完.
一定要全部燒掉.
在聯王紀上的八章六十六節.
經文說.
第八天王遣散眾民.
他們都為王祝福.
都心中喜樂.
各歸各家去了.
所羅門有生之年.
的確是給了以色列人平安喜樂.
但這只不過是短暫.
只有耶穌基督給我們的平安喜樂.
才是永恆.
我講完了兩點.
就是這個殿是神應許的殿.
這個殿亦是以色列人一同慶祝的殿.
第三點這個殿亦是神話語的殿.
這個和我講第一點.
神榮耀的殿是有一個首尾的結構.
當所羅門完成了獻殿禮之後.
耶和華第二次向所羅門顯現.
第一次是什麼時候?.
當神來到問所羅門.
你要些什麼?.
所羅門就說我希望得到智慧.
那次在基片就是第一次.
神向所羅門顯現.
在列王記上的第三章.
剛才我們才見到神的榮耀降臨.
將祭物全部燒了.
有沒有見到?.
神的榮耀和神的話語是分不開的.

$^{761}$在以塞瓦書的四十章第八節.
新約的彼得前書一章二十五節.
都告訴我們.
只有神的話語是永遠長存.
神來到所羅門那裡.
就是跟祂說了神的話語.
你說有些什麼話語?.
一共有三點.
第一點就是神給所羅門的英文.
在列王記上的九章三節.
神對所羅門說.
你向我所禱告即祈求.
我都英文了.
我已將你所見的這殿分別為聖.
使我的名永遠在其中.
我的眼我的心也必尚在那裡.
這個經文和歷代志下的七章十一至十六節.
是一模一樣.
神英文了所羅門.
主要是祂七個祈求.
當以色列人犯了罪.
神說我一定會寬恕他們.
只要他們謙卑自己.
向我禱告.
尋求我的面.
遠離他們的惡行.
神說我一定會聽他們的禱告.
我們要知道在歷史當中.
神從來沒有跟其他國家立約.
神只有跟以色列國立約.
今日這個英文亦是給我們每一個神的子民.
每一個基督徒.
新約也有教導.
在約翰一書.
約翰告訴我們.
神是信實的.
只要我們犯了罪之後.
願意認罪悔改.
神是一定會寬恕我們.
這個無論新舊約都是這樣說.

$^{801}$這個亦是我們今天的主題.
今天的主題是什麼?.
就是悔改和復興的真正途徑.
當我們是願意悔改的時候.
神一定會寬恕我們.
所以第一點.
神跟所羅門說的話.
就是應允了所羅門的禱告.
第二點.
神是要所羅門來服從他.
在《樹列王記》上的九章四至五節.
「你若效法你父大位.
全誠實正直的心.
行在我面前.
遵行我一切吩咐你的.
謹守我的律例典章.
我就必兼顧你的國位.
在以色列中直到永遠.
正如我應許你父大位說.
你的子孫必不斷人.
坐以色列的國位」.
非常個人化的華語.
是跟所羅門說的.
你要得到我的祝福.
你要怎樣?.
要好像你父大位一樣.
成為一個合神心意的人.
有很多人讀到這裡會怎樣?.
會瘋瘋癲癲.
大位真的這麼好?.
不是,你說大位犯很多罪.
很嚴重的罪.
他犯奸淫.
他還殺了人.
現在要所羅門學他?.
不是學他殺人.
不是學他犯奸淫.
學大位甚麼?.
有一個悔改的心.
這才是神所體重.

$^{841}$正如詩篇的五十一篇十七節.
大位犯了罪之後.
他說「憂傷痛悔的心,你必不慶悍」.
不是我們有沒有犯罪.
剛才我已經讀了經文.
世人都犯了罪.
無論舊約新約都是這樣告訴我們.
我們都是會犯罪.
但最重要的是甚麼?.
知道認罪悔改.
這才是神所體重.
這亦是我們今天的主題.
悔改和復興真正的道路.
還有第三點.
就是神給以色列人的警告.
所羅門的後人.
神說「你和你的後人.
如果轉去祀奉別神.
我就會使以色列人.
從我賜給他們的地上剪除.
並且我為幾名所分別為聖的殿.
也必捨棄不顧」.
這裡說得很清楚.
神告訴以色列人.
是有祂的管教.
我之前也說過這個作者.
聶王紀的作者.
他的文筆是非常非常精彩.
這裡其實他用了一個字.
他說「你們以色列人.
如果不聽話.
你們將會成為萬國的笑柄」.
笑柄這個字.
原文希伯來文.
是Marshall這個字.
Marshall這個字.
也是解釋「真言」.
你會說為何這麼奇怪.
真言是智慧的言語.
笑柄又是Marshall.

$^{881}$有沒有看到.
都是一線之差.
你要成為笑柄.
你要成為智慧.
都是一線之差.
還有一個字.
又是Marshall這個字.
在希伯來文.
但是同音.
只不過這個字的解釋是甚麼.
是統治.
其實這個作者真的很有智慧.
文筆非常精彩.
他在聶王紀給我們看到甚麼.
是神將所羅門放在大衛的王位上.
給他統治以色列國.
統治就是Marshall這個字.
也給所羅門智慧.
以至所羅門能夠作出真言.
就是有智慧的言語.
這個又是Marshall.
但是最終就是因為所羅門.
和他的後裔.
他的子孫.
沒有聽到神的話語.
以至以色列人最終成為萬國的笑柄.
笑柄這個字.
又是Marshall這個字.
你有沒有看到.
這個聶王紀的作者的文筆.
真的非常精彩.
以一個神給他智慧.
被譽為世界上最聰明的人.
最後怎樣.
他的國家成為萬國的笑柄.
你說這是不是一個很悲劇的收場.
非常之可悲的一個收場.
在結束的時候.
我們一起來看看.
我今天的題目是甚麼.

$^{921}$就是神的殿就是萬國禱告的殿.
悔改與復興真正的道路.
經文很清楚.
我一直在說的時候.
不是我們有沒有犯罪.
我們每個人都會犯罪.
而是當我們犯了罪的時候.
我們怎樣.
如果我們願意認罪悔改.
回到神的身邊的時候.
歸向神的時候.
神一定會赦免.
這個就是我們今天的訊息.
但是也有另外一些提醒.
當我們去嘲笑所羅門.
成為萬國的笑柄的時候.
我們自己都要自省.
其實有很多的教會.
基督教的機構.
宗派.
甚至很多基督教的事工.
都已經放棄了真正基督教的信仰.
真的很可惜.
他們不再將榮耀歸給神.
只有個「做」字.
但不是真的將榮耀歸給神.
他們這些機構曾經做過很多好事.
曾經將福音傳給很多未聽過福音的人.
但是現在.
你可以在這些機構上寫三個字.
就是「以加博」.
「以加博」是什麼意思.
記不記得在撒母耳上的第四章.
「以利」這個大祭司.
記不記得.
在撒母耳的時候.
他有兩個五翼仔.
兩個五翼仔做什麼.
拿著藥櫃和斐理士人爭戰.
結果藥櫃被擄.

$^{961}$當時其中一個兒子的太太懷孕了.
當這兩個兒子戰敗死在沙場的時候.
他就生了這個嬰兒.
他幫他起名做什麼.
就是「以加博」.
「以加博」的名字的意思是什麼.
榮耀已經離開了.
所以我們都要很小心.
如果我們學這些機構一樣.
我們一樣都會變成「以加博」.
願神憐憫.
願神繼續祝福我們.
是有去做大使命的心.
去將祂的福音繼續廣傳.
我說「列王記」是有很多的伏筆.
在「列王記」上的八章五十八節.
所羅門是要以色列人的心歸向神.
「歸向」這個字原文希伯來文.
是「Natah」這個字.
作者在「列王記」上的十一章第四節.
說所羅門被他的妻妾誘惑去敬拜偶像.
「誘惑」這個字又是「Natah」這個字.
有沒有看到?又是一線之差.
你說為什麼這麼奇怪?.
你一是就歸向神.
一是就被引誘去歸向假神.
有沒有看到?是一個方向的問題.
這是其中之一個伏筆.
還有作者說.
所羅門向以色列人說.
你們要向耶和華傳誠實的心.
遵行他律例.
謹守他的誡命.
這裡誠實的心的誠實.
這個字原文希伯來文.
就是「Shalem」這個字.
「Shalem」這個字其實是.
所羅門名字的另外一個寫法.
有沒有看到這個作者的文筆.
真的很有趣,很精彩.

$^{1001}$這個呼籲以色列人一心歸向神.
傳誠實的心的君王居然怎樣?.
自己被妻妾誘惑去拜偶像.
這個也是我們要很小心.
成為一個傳道人.
不要只懂得呼籲回眾來歸向神.
自己的心也是很重要.
還有最後以色列人遭到亡國被流放.
在聖經裡說要等到耶和華.
使他們的心歸向他自己.
這個其實是在摩西的申命記.
在申命記的三十章第六節.
摩西說「耶和華你的神.
必將你心裡和你後裔心裡的污穢除掉.
好叫你盡心盡性愛耶和華你的神.
使你可以存活」.
以色列人一看到這節經文就會想.
行國禮.
將身體污穢的一部分切除了.
他們以為行了身體上的國禮.
那就聖潔了.
但不是的 神要他們行心的國禮.
心的國禮在舊約已經出現了.
在以西結詩的三十六章第二十六節.
神說我要給你們有一個新的心.
在耶利米蘇的三十一章三十一至三十四節.
神說我要給你們立新約.
這個律法是放在你們的裡面.
這個律法是放在你們的心上.
不再在石板上.
神要成為我們的神.
我們要成為神的子民.
這個就是神希望看到.
神給我們立新約.
所以我們看到以色列人的慘劇.
是給我們一個借鏡.
我們要小心.
我們要接受神.
願意未接受神的能夠接受神.
願意已經接受神的能夠更加愛神.

$^{1041}$而最後我也想再說的是平安祭.
我剛才說平安祭.
就是預表耶穌基督的救贖.
祂在十字架上為我們犧牲.
所有願意接受祂的人都能夠與神和好.
你說為什麼?.
因為接受了祂之後.
你就能夠得到從神而來的平安.
但是不要停留在那裡.
很多基督徒就是停留在那裡.
成為基督徒之後就怎樣?.
哎呀!我又平安了.
這樣就夠了是嗎?.
不是的.
我們要明白神的話語.
我們要遵守神的誡命.
才能夠活出耶穌基督的生命.
耶穌基督才能夠成為我們生命之糧.
這樣祂才是我們真真正正的平安祭.
而不是說只是.
得救了我就平安.
但是真的要活出美好的戰勝.
願神祝福祂自己的話語.
讓我們低頭禱告.
讓我們看到所羅門是起聖殿.
不單是起聖殿.
也是奉獻聖殿.
我們也看到今天的經文告訴我們.
這個聖殿就是萬國禱告的殿.
我們每一個人都是罪人.
正如無論新約舊約都提醒我們.
眾人都犯了罪.
我們沒有辦法可以防止不去犯罪.
但是我們知道.
如果我們犯罪的時候.
我們願意認罪悔改.
你就是那位信實的神.
你一定會赦免我們.
願意我們都有這個心.
來成為一個真正的基督徒.

$^{1081}$在你的話語上下功夫.
遵守你的誡命.
透過你的話語.
能夠將耶穌基督的生命活出來.
以致我們能夠真真正正接受耶穌基督.
成為我們的平安祭.
願意我們都能夠繼續有這個心.
來執行大使命.
而不是說單單我們自己得救就算.
願意我們都是在你回來的時候.
能夠被你稱為一個忠心良善的僕人.
將那榮耀還賺得貴給你.
我們這樣禱告.
奉主耶穌基督的名而求.
阿們.
\newpage



\section{路加福音 15:11-32}
\label{sec:032348}
\textbf{慈愛天父心:從被愛到愛人的呼召}~~~~蕭海波牧師博士
\newline
\newline
連結: \href{http://www.church.com.hk/acms/content.asp?site=cdc\&op=show\&type=product\&code=032348\&layout=sermon}{\texttt{032348}} ~~~~ 語音日期: 2026-06-21
\newline
\newline
\hyperref[sec:032347]{\small{< < < PREV SERMON < < <}}
~
\hyperref[sec:index_chronic]{\small{[返順時目]}}
~
\hyperref[sec:032349]{\small{> > > NEXT SERMON > > >}}
\newline
\newline
路加福音 15:11-32
\newline
\begin{longtable}{cl}
\hline
\hline
章節 & 經文 (和合本修訂版)\\
\hline
15:11 & \begin{tabularx}{0.7\textwidth}{X} 耶穌又說：「一個人有兩個兒子。 \end{tabularx} \\ \\ \relax
15:12 & \begin{tabularx}{0.7\textwidth}{X} 小兒子對父親說：『父親，請你把我應得的家業分給我。』他父親就把財產分給他們。 \end{tabularx} \\ \\ \relax
15:13 & \begin{tabularx}{0.7\textwidth}{X} 過了不多幾天，小兒子把他一切所有的都收拾起來，往遠方去了。在那裡，他任意放蕩，浪費錢財。 \end{tabularx} \\ \\ \relax
15:14 & \begin{tabularx}{0.7\textwidth}{X} 他耗盡了一切所有的，又恰逢那地方有大饑荒，就窮困起來。 \end{tabularx} \\ \\ \relax
15:15 & \begin{tabularx}{0.7\textwidth}{X} 於是他去投靠當地的一個居民，那人打發他到田裡去放豬。 \end{tabularx} \\ \\ \relax
15:16 & \begin{tabularx}{0.7\textwidth}{X} 他恨不得拿豬所吃的豆莢充飢，也沒有人給他甚麼吃的。 \end{tabularx} \\ \\ \relax
15:17 & \begin{tabularx}{0.7\textwidth}{X} 他醒悟過來，就說：『我父親有多少雇工，糧食有餘，我倒在這裡餓死嗎？ \end{tabularx} \\ \\ \relax
15:18 & \begin{tabularx}{0.7\textwidth}{X} 我要起來，到我父親那裡去，對他說：父親！我得罪了天，又得罪了你， \end{tabularx} \\ \\ \relax
15:19 & \begin{tabularx}{0.7\textwidth}{X} 從今以後，我不配稱為你的兒子，把我當作一個雇工吧。』 \end{tabularx} \\ \\ \relax
15:20 & \begin{tabularx}{0.7\textwidth}{X} 於是他起來，往他父親那裡去。相離還遠，他父親看見，就動了慈心，跑去擁抱著他，連連親他。 \end{tabularx} \\ \\ \relax
15:21 & \begin{tabularx}{0.7\textwidth}{X} 兒子對他說：『父親！我得罪了天，又得罪了你，從今以後，我不配稱為你的兒子。』 \end{tabularx} \\ \\ \relax
15:22 & \begin{tabularx}{0.7\textwidth}{X} 父親卻吩咐僕人：『快把那上好的袍子拿出來給他穿，把戒指戴在他指頭上，把鞋穿在他腳上， \end{tabularx} \\ \\ \relax
15:23 & \begin{tabularx}{0.7\textwidth}{X} 把那肥牛犢牽來宰了，我們來吃喝慶祝； \end{tabularx} \\ \\ \relax
15:24 & \begin{tabularx}{0.7\textwidth}{X} 因為我這個兒子是死而復活，失而復得的。』他們就開始慶祝。 \end{tabularx} \\ \\ \relax
15:25 & \begin{tabularx}{0.7\textwidth}{X} 「那時，大兒子正在田裡。他回來，離家不遠時，聽見奏樂跳舞的聲音， \end{tabularx} \\ \\ \relax
15:26 & \begin{tabularx}{0.7\textwidth}{X} 就叫一個僮僕來，問是甚麼事。 \end{tabularx} \\ \\ \relax
15:27 & \begin{tabularx}{0.7\textwidth}{X} 僮僕對他說：『你弟弟回來了，你父親因為他無災無病地回來，把肥牛犢宰了。』 \end{tabularx} \\ \\ \relax
15:28 & \begin{tabularx}{0.7\textwidth}{X} 大兒子就生氣，不肯進去，他父親出來勸他。 \end{tabularx} \\ \\ \relax
15:29 & \begin{tabularx}{0.7\textwidth}{X} 他對父親說：『你看，我服侍你這麼多年，從來沒有違背過你的命令，而你從來沒有給我一隻小山羊，叫我和朋友們一同快樂。 \end{tabularx} \\ \\ \relax
15:30 & \begin{tabularx}{0.7\textwidth}{X} 但你這個兒子和娼妓吃光了你的財產，他一回來，你倒為他宰了肥牛犢。』 \end{tabularx} \\ \\ \relax
15:31 & \begin{tabularx}{0.7\textwidth}{X} 父親對他說：『兒啊！你常和我同在，我所有的一切都是你的； \end{tabularx} \\ \\ \relax
15:32 & \begin{tabularx}{0.7\textwidth}{X} 可是你這個弟弟是死而復活，失而復得的，所以我們理當歡喜慶祝。』」 \end{tabularx} \\ \\
[1ex]
\hline
\hline
\end{longtable}
$^{1}$弟兄姊妹平安.
多謝張牧師的邀請.
讓我在這裡有機會.
透過分享聖經.
和大家一齊敬拜神.
今天是父親節.
很開心收到父親節的禮物.
小朋友很可愛.
但我都恭祝各位當中的父親.
借用約翰所講的.
恭祝大家願你們凡事興盛.
身體壯健.
正如你的靈魂興盛一樣.
恭祝大家事常主因滿人.
大家有沒有聽過查理士迪根斯這個名字.
我記得我小時候在香港.
經常聽香港電台有個查理士迪根斯.
原著的長篇小說.
很喜歡聽的.
這個查理士迪根斯.
他講過浪子.
這個浪子的比喻.
稱為有史以來最美好的短篇小說.
因為他用了很短的篇幅.
剛才我讀了好像很長的聖經.
讀了21節.
但這21節的聖經.
在這個短篇小說裡面.
就導進了人性的冷暖.
又講了人性的墮落.
講到人怎樣有機會回歸.
有盼望有出路.
所以查理士迪根斯說.
這個真是一個很了不起的短篇小說.
問題是這樣.
問題是究竟這個耶穌基督浪子的比喻.
裡面的主角是不是真的講這個浪子呢.
會不會是這樣呢.
會不會一個焦點其實是講這個父親.
因為是父親節.

$^{41}$其實會不會一個焦點是講這個父親呢.
當我們看這個浪子的故事的時候.
這個浪子的比喻的時候.
當我們看到中間這個慈愛的父親.
剛才大家讀聖經都讀到了.
很容易就讓我們想起.
我們慈愛的天父.
他那種寬容.
他那種對我們的愛.
是多麼的長闊高深.
在這個路加福音15章第11節.
耶穌就講了.
聖經裡面講過.
一個人有兩個的兒子.
就因為這一句話.
很多人就認為這個中間的主角就是浪子.
因為他有兩個兒子.
其實如果我們根據上下文我們看下去.
我們不難找出.
當中的主角.
不可能真的不是這兩個兒子.
很有可能就是那個父親.
反正我們看這個比喻的時候.
我們來看看這個比喻的背景.
和當時的聽眾是怎樣.
很簡單的.
當時是這樣的.
稅吏和罪人靠近耶穌.
要聽他講道.
但是那些法利賽人和文士.
就覺得這個人接待罪人.
他又和他們吃飯.
當時這些稅吏.
為什麼和罪人一起.
混為一談成為一個群體.
大家如果讀聖經都知道.
當時在羅馬帝國的管治裡面.
這些稅吏就是幫羅馬政府服務.
來收稅的.
說起收稅很多東西可以講.

$^{81}$原來他們隨便可以收多少.
不是像我們這樣.
在我們稅務局.
限定有一個標籤.
他喜歡收多少就收多少.
但是他們要錢買這個稅來做.
所以那些猶太人.
就很不喜歡他們.
當他們罪人一樣.
但是耶穌歡迎他們.
所以耶穌和他們一起吃飯.
所以那些法利賽人和文士就說.
你看這個人.
這個人就是誰.
就是耶穌.
這個人接待罪人.
又和他們吃飯.
聽見他們洩洩之語.
明白他們的心聲.
所以耶穌又怎樣.
就開始講比喻.
開始講比喻.
這個比喻是和當代耶穌的時候.
那些罪人.
那些稅吏.
那些法利賽人.
那些文士那些律法師講的.
當時二千多年前.
但是這個比喻.
是和我們今天講的.
為什麼?.
因為當中有可能.
我們都好像他們一樣.
我們從這裡看一下.
我們來看一下.
在耶穌講這個比喻的時候.
他一年講了三個比喻.
大家都熟悉.
今天我們講浪子的比喻.
但是我們不能夠.

$^{121}$很快地看一看.
之前那兩個比喻是講什麼的.
很短的.
就是他講羊的比喻.
羊的比喻.
就在路加福音.
十五章四到七節那裡講.
講到一個人有一百隻羊.
大家都熟悉.
就不見了一隻.
他就去找,放下九十九隻去找.
九十九隻.
不見了一隻,就是百分之一.
他不見了百分之一.
但是他都去找,去找這些羊.
找這些羊.
我在神學院裡面.
做事的時候有一個童工.
他就去.
以色列旅行.
他就學人家叫那些羊.
那些羊其實原來都有.
很聰明的一面的.
有一群羊.
那個童工就學那個主人.
那個主人有個名叫那些羊的.
不知道叫什麼名.
叫那些羊來.
那些羊理不理他.
不理他.
但是呢.
那個羊群的主人.
有一叫,那些羊.
跑去那個主人那裡.
所以聖經說,羊認識.
他的主人.
真的這樣的,認識主人.
但是那些羊.
在豬般的動物.
裡面來講,又可以說.

$^{161}$聰明,但又可以說.
很蠢.
他迷路了.
不會回家.
流浪在.
流浪在曠野.
或者在草叢裡面.
不會回家.
如果我們喜歡貓和狗.
那些狗出了街.
或者貓走了出去.
他都會迷路回家.
但是那些羊就不會.
那些羊一邊聰明.
認得主人的聲音.
一邊又蠢蠢的,傻傻的.
不會回家.
在曠野流浪.
等待黑夜的來臨.
豺狼的來臨.
這個主人就去找他.
主人找到他怎樣.
找到他回來就很開心.
那些羊迷路.
是不會迷路回自己家的.
好了.
第二個比喻.
第二個比喻講到什麼.
不見了錢,有個婦人.
有十元雞.
當她十元雞.
就不見了一元雞.
那又怎樣,翻箱倒籠去找.
去找.
結果就找回了.
那些羊.
剛才我們講的比喻.
那個主角是不是羊.
是找羊那個人.
那個牧羊人.

$^{201}$那這個錢呢.
錢,羊就不會回羊群.
錢當然.
更加不知道自己去哪裡.
更加講不上.
回到原本跌的位置.
所以講到這個.
失錢的人.
其實失錢的這個婦人.
才是這個比喻的主角.
是不是.
而且我們見到那個遞進.
那三個比喻的鋪陳.
第一個是不見了百分之一.
第二個是不見了十分之一.
第三個.
大家考一下你的IQ.
多少分之一.
不見了二分之一.
這個父親.
有兩個兒子.
就不見了50\%.
不見了二分之一.
大家都可以想到.
三個比喻,耶穌講了一年三個比喻.
三個比喻的內容.
都差不多.
但是主題其實很一樣.
大家認為這個浪子的比喻.
主角是誰.
父親.
當然是父親.
當然有時比喻.
有時有些分析.
可能是不一樣的.
似乎這個父親是主角.
都很合理.
所以這個浪子的比喻的主角.
就是這位慈愛的父親.
我們進入今天的經文.

$^{241}$十一到三十二節.
一共二十一節.
比較長一點.
我們當看故事.
我們看到這個爸爸.
在經文裡面.
我們看到他失去了一個很心愛的兒子的父親.
一個很傷心的父親.
這個經文告訴我們.
這個爸爸的小兒子.
他和爸爸要分家產.
最近如果你看一些娛樂新聞.
分家產這件事真是滿天飛.
有時都看到很可悲.
講回聖經.
這個小兒子要分家產.
他分到家產之後.
他就遠走他向.
遠走高飛.
在這個比喻裡面.
主耶穌形容這位父親.
我們雖然沒有很清楚.
一個字寫出來.
但是我們看到這個父親.
經歷了失子之痛.
如果我們看到後面的聖經.
我們看到這個爸爸.
這個兒子我差不多當他死了.
死而復得.
當他沒有了.
你說多傷心多痛.
弟兄姊妹.
當時在猶太人的文化裡面.
這個兒子和爸爸分家產.
是可以的.
是可以拿的.
為甚麼呢.
因為申命記十一章到十七節裡面說.
是可以爸爸在生的時候.
兒子是可以要求拿家產的.

$^{281}$還有一個規矩是怎樣的呢.
大兒子是得到兩份.
分豬肉分兩份.
中國人都有這樣.
其他的兒子就分一份.
是可以這樣拿的.
是可以要求的.
但是這個兒子要求他爸爸分家產.
如果我們看清楚一點.
這個小兒子其實有一個.
隱藏的意義在裡面.
就是.
父親我寧願你現在離世了.
我寧願你現在不存在了.
說得再作一點.
我寧願你死.
大家知不知道有這麼嚴重的事.
他要求分身家.
其實是有這樣的含義在裡面的.
為甚麼這樣說呢.
因為按照猶太人的傳統和文化.
在父親仍然在生的時候.
兒子是可以要求分身家.
但是有一個條件.
這個條件是甚麼呢.
兒子分了身家之後.
是不能夠拿這些財產去變賣的.
只可以在當中.
來代理爸爸的生意.
現代的說話就是這樣.
可以管理爸爸的財產.
但是我們在聖經裡面看到甚麼呢.
我們在聖經裡面看到.
第十三節.
過了不多幾日.
小兒子就把他一切所有的.
都收拾起來.
往遠方去了.
那即是怎樣呢.
即是他賣光了.

$^{321}$他不能夠帶著羊欄.
帶著田.
這樣拖走的.
即是他怎樣呢.
快點去拍賣行.
看看誰隔離屋賣了他.
拿著一大袋錢.
走人.
即是怎樣呢.
你當爸爸是死的嗎.
真是這樣.
行動上其實就等於這樣.
大家又知不知道.
這個小兒子.
年齡是多少歲呢.
不知道 聖經沒有說.
但是根據.
歷史學家 學者.
和聖經學者.
估計這個小兒子.
是由十七到二十歲這段距離.
因為他當時還未結婚.
又可以分身家的年紀.
又可以賣了離走了.
大概十八十九歲.
這樣的年齡不是絕對.
總之是十七二十一歲這段.
就是一個年青人.
剛剛才.
成人.
十八歲.
就分身家.
就拿走 你說爸爸傷不傷心.
很傷心.
這個小兒子的行為.
我們看到.
差不多成為了.
爸爸一種羞辱.
成為了爸爸一種.
真是.

$^{361}$心是撕裂的.
我不知道你可不可以想像到.
爸爸小兒子被走的時候.
心多撕裂.
試想一想.
從父親的角度來說.
這個兒子是沒有了.
想一想.
如果我們有個.
兒子或女兒.
兒女也好.
我們做父母的 做爸爸的.
那麼辛苦的.
幫他們完成了中學.
預備入大學.
打算去到外州.
一間長春藤.
將來.
出人頭地.
將來就會做牧師.
不是 可能做律師.
做律師.
或者醫師.
怎知道.
有一晚看到新聞.
電話響起來.
鈴鈴鈴鈴.
響起來 看到電視.
看到原來警察.
在校舍裡面.
滿了警車 紅燈綠燈.
閃甩閃甩.
有些人.
被人抓了 反著手銬.
原來在中間吸毒的其中一個.
是你的兒子.
傷不傷心.
傷心.
傷心欲絕.
本來滿心盼望.

$^{401}$望子成龍.
怎知道現在變成了蟲.
現在變成這樣.
這個傷心.
是不可以用筆墨來形容.
是不是.
心好像撕裂了.
所以弟兄姊妹.
我們從這個看法想到.
我們天父的心腸.
如果看到我們.
沉溺在這個世界上.
我們離開天父.
我們放縱.
我們的自由.
我們在世界上遊蕩的時候.
天父的心.
傷不傷痛.
因為主耶穌說的這個比喻.
是一個比喻.
想比喻什麼.
告訴我們天父的心腸.
是怎樣的.
弟兄姊妹給我們一個提醒.
既然說這個望子的比喻.
我們又不能夠不說這個小孩.
他雖然有些氣氛在裡面.
爸爸是主角 他也是配角.
說一下這個小孩.
這個小孩離開之後.
是怎樣的.
他放蕩.
浪費了所有的資材.
然後將自己一生.
放在煙裡.
不顧後果.
只是希望得到自己的滿足.
這個小孩的行為.
有時候我們想到.
就好像罪的一種寫照.

$^{441}$罪就是令我們離開天父.
我們好像在這個世界上.
遊蕩.
覺得好像更加的自由.
滿足自己的私慾.
這個小孩更慘.
他不單只傾家蕩產.
最終還要怎樣.
去養豬.
那個年代.
養豬是不可以的.
為什麼.
猶太人認為豬是不潔淨的.
所以後面有說.
我得罪了天.
這個就不詳細解釋了.
他養豬.
餓到連豬吃的豆蝦.
他想吃怎樣.
人家都不給他.
不過說回今天.
廣東人叫豬拿菜.
以前養豬的菜.
叫豬拿菜.
現在的豬拿菜很貴.
現在健康.
減肥去了脂肪.
豬拿菜原來很貴.
不過在那個年代.
是很低賤的東西.
豆蝦.
他想吃豆蝦.
人家都不給他.
我們作為小小的反思靈姊妹.
今天我們會不會想到.
有時候我們人生.
某一段路上.
我們得了神給我們的.
很厚厚的恩典.
我們揮霍在.

$^{481}$我們自己的屍肉上面.
其實這些都是神給我們的.
無論我們的時間.
我們的金錢.
我們的才幹.
我們的學問.
一切我們的家庭都是神給的.
會不會有時候我們浪費在.
我們自己的資材上面.
我們覺得會.
好像遠離了神的時候.
我們更加自由.
回教會又要小組.
又要查經又要祈禱.
很多東西捆綁.
會不會自由一點.
其實離開神沒有真正的自由.
只有失去了我們真正的.
平安和喜樂.
雖然我們離開神.
神是不會.
忘記我們的.
同不同意嗎.
聖經裡面怎麼說.
聖經裡面在耶利米斯三章.
二十節那裡.
二百年是我的愛子嗎.
是我可喜悅的孩子嗎.
我每逢責備他.
還是迎心顧念他.
神責備迎顧念.
所以我心腸戀慕他.
我必要憐憫他.
雖然我們離開神.
但是神的愛.
環繞著我們.
神的愛 神不會忘記我們.
另外.
在二賽亞書這段聖經我很喜歡.
我想做媽媽的更加會清楚.

$^{521}$那裡說什麼.
婦人怎麼能夠忘記.
她吃奶的寶寶呢.
婦人應能忘記.
她吃奶的嬰孩呢.
不連恤她所生的兒子呢.
即或有忘記.
神說.
即或有忘記.
我卻不忘記你.
神在這個比喻裡面.
讓我們思想到.
這個慈愛的父親.
也想起了我們.
天父的慈愛.
天父的慈愛讓我們知道.
如果我們稍稍遠離他.
他的心會很傷痛.
我們要記住.
我們離開神.
不單止我們受苦.
但是我們的神 我們的天父的心.
是會傷痛的.
但是.
雖然我們離開.
或者我們遊蕩.
但是我們的天父.
不會忘記我們.
我們記住.
我們離開 天父的心會傷痛.
我們遊蕩.
天父的心會傷痛.
但是無論如何.
天父不會忘記我們.
來到第二段經文.
講到浪子.
就快回頭了.
講故事的開始的時候.
當然有很多悲慘.
慢慢開始進入高潮.

$^{561}$這裡有少少高潮.
還沒有疊起.
少少高潮.
講到這個小孩.
他醒不過來.
經文裡面告訴我們.
這個小孩在異鄉落魄.
落魄他鄉.
他慢慢醒不過來.
換句話說.
他以前是沉迷在私慾.
和自己的享樂裡面.
他醒不過來的時候.
他突然間想起.
我爸爸家裡面.
工人的食都吃不完.
難道我在這裡餓死嗎.
為什麼要吃豬的豆蝦.
於是他決定.
回家.
他準備了一段說話.
他跟他爸爸說.
他要認罪.
他爸爸我得罪了天.
我又得罪了你.
我現在回來你不要當我是兒子.
就當我和你打工.
想著這樣說.
不知道你們有沒有跟爸爸.
說過對不起.
有沒有.
我有.
我相信主沒多久.
我和我爸爸的關係不是很好.
怎樣不好呢.
就是我離開家.
我離開家關鐵閘.
他就在外面.
回家見到面.
是不打招呼的.

$^{601}$關係很差.
這些家庭的問題.
真是家家有本.
難唸的經.
後來信了主之後.
我覺得不行.
無論誰對還是不對.
我都要和我爸爸認錯.
於是我想了很久.
那時候我已經25歲.
25歲血氣方剛.
不是那麼容易認錯.
和爸爸認錯.
很難的.
真的很難.
我怎樣呢.
我買了六瓶.
白蘭氏雞精.
哭著爸爸在床上睡覺.
就去.
這些白蘭氏雞精給你的.
對不起.
我見到你都不會打招呼.
我忘了說了什麼.
我爸爸看見他.
眼.
男人大丈夫.
眼都有一點濕濕的.
又不知道怎樣好.
突然和我認錯.
不過無論如何.
總之做過這件事.
如果得罪爸爸快點去認錯.
道歉.
爸爸的愛真是很難得.
這個兒子就想和.
父親認錯.
我們看到這個.
爸爸.
其實他一直在等兒子回來.

$^{641}$他知道.
就算他抓兒子回來.
一得到他的人.
我們經常看愛情小說都有這樣說.
得到人得不到心是沒有用的.
你綁住他是沒有用的.
要他的心來回轉.
所謂浪子回頭.
是什麼.
金不換.
浪子.
現在的金很貴.
但是金都不行.
貴都不可以換.
浪子回頭金不換.
這個爸爸等這個浪子回來.
就等得很辛苦.
雖然他沒有去找.
沒有像尹陽那樣去找.
沒有像尹淺那樣去找.
他沒有出去找.
但是他的心一直在等.
很急很渴望在等他回來.
因為他希望他的心回轉.
重點在這裡.
怎樣說呢.
重點在二十節.
相離還遠.
還在很遠的時候.
那又怎樣.
他父親看到就動了慈心.
然後這兩個字很重要.
跑去.
跑去抱住他的頸.
連連與他親嘴.
他回來的時候.
我不知道他有沒有換西裝.
他養豬的.
豬的豆蝦都不吃的.
不給他吃的.

$^{681}$抱住他連連親嘴.
弟兄姊妹其實這節聖經是很震撼的.
怎樣震撼.
跑去當然形容到.
這個爸爸很心急的.
衝出去抱住小孩回來.
是不是.
但是在猶太的社會傳統裡面.
年長的父親.
或者那些領袖.
家庭的長輩.
通常是不會這樣跑去的.
為什麼呢.
因為當時他們穿的衣服是穿什麼的.
是穿袍的.
穿袍的.
我們中國人都會穿什麼.
長衫馬褂.
你想一下那些悲情長篇劇.
那些嘉春秋.
坐在正堂.
正襟危坐.
穿著拐杖戰長袍.
如果他跑的時候會怎樣.
跌死他.
因為他蓋過他的膝蓋.
如果他要跑的時候.
會怎樣.
要拿高那件袍.
那件袍超過他的膝蓋.
他才可以跑.
是不是.
拿高膝蓋的時候會怎樣.
看到飛毛腿.
是怎樣.
很沒體面的.
是很.
當時視為一種羞辱.
是不好的.
但是我們看到這個父親怎樣.

$^{721}$跑去.
他遠遠望見.
其實他在這裡等了很久.
香港有個山.
叫什麼山.
望夫山.
變成石頭.
望仔山.
等了很久.
站在家門口.
等.
周圍可能有些鄉親父老.
看到.
你兒子這麼五逆.
回來養豬.
豬的豆蝦.
都不給你吃.
你不是.
廣東話.
你個五逆仔.
一回來不是這樣罵的.
你去了哪裡.
不是這樣罵的.
或者起碼都蹲起款.
你跟我說你的錢去了哪裡.
他不是.
他連連跑去抱著他.
這是什麼.
父親的愛是超越了什麼.
超越了.
他的尊嚴.
超越了.
當時文化的限制.
超越了.
表面.
人間的一般普通的愛.
是超越.
尊嚴.
放下尊嚴.
超越文化.

$^{761}$在愛的深處的一種表達.
剛才姊妹說到.
愛的表達.
我們又想到我們的天父.
對我們的愛.
我們天父對我們的愛是主動的.
是藉著他的恩典.
來找我們.
剛才詩歌所唱的.
這個奇異的恩典.
我信主的時候.
這首歌是我很喜歡的詩歌.
我很喜歡的詩歌.
我不懂彈琴的.
但我買了這麼小的電子琴.
連了一二三四進去.
來表達.
奇異的恩典.
我以前看不見的.
現在看見了.
弟兄姊妹.
神的愛是主動的.
天父的愛是主動的.
他用恩典來尋找我們.
在我們未曾知道有需要的時候.
他已經安排了環境.
透過朋友.
透過聖靈來感動我們.
我們的回頭.
其實不是我們單方面的.
是神安排了的.
因此我們怎樣.
可不可以說阿們.
奇異的恩典.
阿們.
第三點,真的來到終場的結局.
還不是大結局.
終場結局.
講完這一場之後.
就落幕,然後中場休息一下.

$^{801}$再有下一幕.
有時間再講.
講一下這裡,很快的.
這裡說到什麼呢?.
說到浪子終於回到家門口.
他預備好講認罪的說話.
他講了一半.
他的爸爸有沒有讓他說?.
沒有讓他說完.
爸爸沒有用說話來回應.
他的父親.
用行動來回應.
行動回應他.
你不用說.
你不是故宮,你也不是奴僕.
你是我的兒子.
他用什麼行動來回應呢?.
首先.
他爸爸的赦免是超過.
浪子的認罪.
兒子怎麼說?.
他說我得罪了天.
又得罪了你.
想做你的故宮.
爸爸說不是,你是我的兒子.
這個爸爸.
將上好的袍子給他穿上.
上好的袍子.
我們可以這樣想.
當時的僕人或奴僕.
在那個時候.
穿的衣服是很粗糙的.
上好的袍子就是.
兒子的身份.
用行動讓他知道.
你仍然是我的兒子.
你是我的兒子.
你不是奴僕.
看到這裡我就想起.
加拉太書三章26節裡說的.

$^{841}$你們因信基督耶穌.
都是神的兒子.
今天我們身上.
穿上兒袍.
是不是?.
你身上穿上兒袍.
你知道嗎?.
還有你額上有聖靈的印記.
你知道嗎?.
今天我們是兒子的身份.
這個爸爸讓他穿上兒袍.
穿上這件上好的袍子.
還有什麼?.
讓他戴上戒指.
戒指代表什麼?.
在當時.
戒指就代表家族的權柄.
讓他成為兒子.
恢復他的身份.
讓他權柄.
還有什麼?.
穿上鞋子.
穿鞋.
穿鞋有什麼重要?.
當時的奴僕是沒有鞋穿的.
當時的奴僕.
被賣到家中是沒有鞋穿的.
但是有鞋穿.
代表什麼?.
你在家裡是自由的.
你喜歡去哪裡.
你都可以去哪裡.
你可以吩咐家人.
做你想要他做的事.
還有什麼?.
宰一隻肥牛仔.
歡喜快樂.
從3月到4月.
我去了九寨溝.
去那裡看到壯族的人.

$^{881}$他們養了很多毛牛.
有沒有去過.
九寨溝看過那些毛牛?.
他們說那些牛.
就是那些家人的財產.
有多少隻牛.
他有多富有.
他說他最多問導遊.
最多有多少隻牛.
他說最多有幾十隻.
他不是這麼隨便.
宰一隻牛仔.
這裡是宰一隻肥牛族.
在當時.
宰一隻肥牛族.
不是小事.
是大事.
歡喜快樂的大事.
才會這樣宰.
所以你看到爸爸.
那種歡喜快樂.
不是不脈可以形容的.
浪子回頭.
金又不可以換.
戴上戒指.
穿鞋 宰肥牛仔.
歡喜快樂 抱著他.
跑過去 親他幾口.
我們的天父對我們也是這樣.
很多年前.
我在參加牧會的時候.
有個姊妹.
初來的時候很熱心.
她上主一學 上團契.
上崇拜 上小組查經.
很乖.
過了兩三年.
失蹤了.
不見了.
沒回來教會.

$^{921}$我打電話問候她.
沒接電話.
找了很多次.
都不知道.
吩咐弟兄姊妹你們去關心她.
弟兄姊妹打電話去.
不接電話.
過了若干年我都忘記了.
有一次團契.
完了之後.
沒看到她在.
團契完了之後.
我去問她最近好嗎.
她想哭.
牧師.
對不起.
我知道我真是很頑皮.
令你傷心.
要你擔心.
我說你不用說.
你回來就好了.
我讀起這段聖經的時候.
我就想起這位姊妹.
稍稍.
稍稍體會一下.
這位父親.
更加稍稍體會一下.
天父的心腸.
我們回到天父那裡.
就好了.
是不是.
我們認識一些弟兄姊妹.
可能現在.
已經同神相離甚遠.
有沒有這些人你認識.
或者以前有回來.
慕道園浸信會崇拜.
現在不見了.
會不會去了哪裡.
會不會他離開了神.

$^{961}$或者在世界路上.
遊蕩.
我們可不可以邀請他回來.
回來天父的身邊.
求主給我們體會.
天父的心情.
或者我們可不可以去.
做一些探訪.
將他們帶回來.
天父的身邊.
甚至有時候我們想起.
可能我們自己都曾經遊蕩去哪裡.
是不是.
得罪天父 叫天父的心傷痛.
神學家約翰加爾文.
他這樣說.
悔改不是一次性的事件.
而是一個持續.
直到行完人生路程的.
屬靈成長過程.
多謝張牧師.
介紹我們三十天的靈修.
我每天都有看.
其中有一個chapter是不是說到.
我們都要時時去到神的面前.
來反省.
想想我們有什麼做得不好的地方.
求神赦免.
有一首詩歌我很喜歡.
這首詩歌叫.
Mercy is more.
憐憫.
他憐憫更多.
其中一句說話是這樣說的.
我們的罪雖然眾多.
但是他的憐憫更多.
有時候我們不知不覺遊蕩.
或者有弟兄姊妹在網上崇拜.
很久沒有回到這裡.
我相信.

$^{1001}$教會弟兄姊妹都很掛念你們.
我們的罪雖然眾多.
但是他的憐憫更多.
主幫助我們.
讓我們知道.
我們的悔改雖然稍稍轉頭.
但是天父的慈愛.
天父赦罪的恩典.
遠超過我們所求所想.
我們都要怎樣.
好了.
這個故事說到這裡就稍為要落幕.
但是未完.
鏡頭一轉.
好像拍電影.
鏡頭一轉去那邊.
突然間有個大頭出來.
這個是誰.
是個大仔.
大仔怎樣呢.
這個大仔在田裡開工.
做得很辛苦.
回來就遇到經文裡面所說.
離家不遠.
聽見很多人跳舞.
又有跳舞的音樂.
叫了個僕人走過來.
什麼事.
你的兄弟你的弟弟.
因為無災無病回來.
你爸爸宰了一隻肥牛仔.
來跟他慶祝.
大仔就很激氣.
很激氣.
怎麼辦呢.
他很憤怒.
他要不要進屋裡慶祝.
他拒絕進去.
裡面很開心.
他拒絕進去.

$^{1041}$大仔裡面有一句聖經.
反映大仔裡面的心腸.
心情是說什麼.
十五章二十九節.
我服侍你這麼多年.
從來沒有違背過你所吩咐的命令.
你沒有給我一隻山羊仔.
叫我跟朋友一起快樂.
這句說話流露.
他心情裡面有什麼想法.
我可以想到幾樣東西.
第一樣東西他看到.
大仔和他爸爸之間.
他不是父子那種親密的關係.
他看到他自己.
好像僱主和工人.
我服侍你.
你吩咐我的事我都做.
這裡看到.
好像不是一個.
互相親愛的父子關係.
然後他心裡面累積了.
很多年的不滿.
好像覺得不被重視.
然後他就將那個關係.
放在得失上面.
你看看你有什麼比我好.
你看看你有什麼比別人好.
表面看來.
這個大仔有沒有離家出走.
沒有.
還在裡面做很勤力的.
在田間辛苦工作.
但是他其實內心有沒有離家出走.
有.
他的內心.
就好像離家出走的一個.
大仔那樣.
他心裡面有一幅圍牆圍住.
和他爸爸的關係好像很疏離.

$^{1081}$好像一個工人那樣.
但是我覺得.
當我們讀聖經的時候.
我們稍稍看一看這個.
好像給我們一個提醒.
我們的信仰不只是看外表的行為.
而是看我們內心的關係.
是怎麼樣.
弟兄姊妹.
這麼多年我們在當中.
來服侍很忙碌.
用了很多時間.
但是我們會不會覺得.
好像我們不是太認真享受.
作為神兒女.
在當中有這個平台.
給我們很自由服侍主.
會不會.
或者以前有現在忘記了.
如果現在有.
那就好了.
會不會這樣呢.
求主都幫助我們.
但是大仔他的憤怒.
還有怎麼樣.
他不是說我弟弟回來了.
他說什麼.
他說你這個兒子.
你這個兒子.
他有沒有當他弟弟是弟弟.
沒有他是你這個兒子.
不是我弟弟.
你這個兒子散盡了家財.
和那些妓女在一起.
你還對他這麼好.
是不是.
所以我們看到這個大仔.
他沒有離家.
但是仍然好像一個浪子.
我們.

$^{1121}$看到這裡我們很簡單.
覺得我們要.
反省一下我們在教會裡面.
在神的面前的身份是什麼.
我記得張牧師.
有一次講到.
他說我們要麼做罪的奴僕.
要麼做什麼.
要麼做義的奴僕.
沒得選擇.
沒得中間騎牆.
這個是一個真理.
在做義的奴僕裡面.
但是我們還是什麼.
我們是兒子.
我們一起讀聖經.
加拿大書四章.
六到七節.
請一起讀.
求你不是奴僕.
乃是兒子了.
既是兒子.
就靠著神為後子.
弟兄姊妹我們服侍神.
我們不是奴僕.
我們是僕人.
我們謙卑在主的面前.
但是不要忘記我們是兒子.
你可以心裡面叫.
阿爸父嗎.
可以.
我們心裡面因為有神的靈在我們裡面.
我們心裡面都可以叫阿爸父.
我們是兒子.
所以在裡面服侍主的時候.
我們滿心歡喜快樂.
說到這裡.
我覺得弟兄姊妹在這個教會服侍.
都是很歡喜快樂的.
為什麼我這樣說呢.

$^{1161}$因為我初初第一次來這個教會.
聚會的時候.
不知不覺我已經來了六個月了.
眨一下眼.
初初來的時候.
弟兄姊妹很熱情的.
很歡迎.
我和張牧師.
教會真是有福.
有群弟兄姊妹這麼熱情的.
這麼歡迎新來賓.
我和.
德勝傳道.
也是這樣說的.
叫呂牧師.
德勝傳道.
上次我和他吃飯的時候.
教會弟兄姊妹真的很熱心.
真的很熱情服侍主.
很喜樂.
我都很被吸引.
留到現在.
有這麼好的弟兄姊妹你們都是很喜樂.
但是.
如果有服侍主.
稍微失去喜樂.
求主幫助我們.
透過今天的聖經.
讓我們知道我們在當中可以很自由.
很歡喜快樂服侍主.
不是因為我們要做什麼.
是因為神的愛.
激勵我們.
因為天父愛我們.
我們的愛的回應.
我們能夠和神有很親密的關係.
願意我們.
對天父的愛.
關係密切.
遠遠多過我們對事奉的要求.

$^{1201}$當我們愛主.
和主有好的關係.
我們的服侍就會很喜樂.
很平安.
多說一件事.
我雖然年紀都不小.
不要揶揄.
好幾年前在教會講道.
那間教會的創辦人.
九十幾歲.
牧師多謝你來今天講道.
都是靠著恩典.
做得就做.
神給你服侍.
就是神的恩典.
神給我們服侍.
就是我們的恩典.
是主給我們的自由.
和喜樂.
我們心裡都叫.
阿爸父.
可以講阿們嗎.
阿們.
今天是父親節.
父親節我們一定要講一講.
父親的勸勉.
真是大結局.
這兩節聖經是很震撼的.
父親對他說.
你常和我同在.
我一切所有的都是你的.
只是你的兄弟.
是死而復活.
失而又得的.
所以我們理當歡喜快樂.
這兩節聖經很精彩.
這個父親說.
我和你常在一起.
我的都是你的.
是不是.

$^{1241}$所以我想到.
我們的天父.
時時都和我們在一起.
我們的心裡有天父在裡面.
天父的心.
有我們在裡面.
我們裡面有天父.
天父裡面有我們.
真是分都分不開.
但是是不是很奇妙呢.
神的愛是不是很奇妙.
當我看到這段時候.
我就想起二賽亞書第四十章.
四十章十五節.
漢萬民都像縱水桶的一滴.
又算如天平上的微塵.
弟兄姊妹你和我都是天平上的微塵.
但是天父對我們的愛.
不離不棄.
他的慈情愛索.
一路追著我們.
一路跟著我們.
一路都不放棄.
是不是很奇妙.
今早我吃早餐的時候.
我當溫習這一篇講章的時候.
我還想起有詩歌.
何能如此呢.
像我這樣的罪人.
主耶穌竟然流出寶血救贖我.
何能如此.
像我這樣的罪人.
是不是.
阿們.
真是不可以這樣.
弟兄姊妹.
我們當然很盼望.
我們勤力的事奉.
忠心的事奉.
但是我想提醒.

$^{1281}$我們彼此的勉勵.
當我們跟天父的關係.
很密切的時候.
天父寧願我們跟他的關係很親密.
遠遠超過我們的事奉.
是不是.
是不是這樣.
我們在當中享樂.
服侍主那種歡樂的自由.
真是最後.
要完結了.
今天看這個浪子的比喻.
其實主角.
我覺得是這位.
慈愛的父親.
也比喻我們的天父.
我們慈愛的天父.
永遠不會忘記我們.
我們的慈愛的天父.
無論我們走了多遠.
他在當中等待我們.
他的恩典超過我們所求所想.
我們慈愛的天父.
寧願我們跟他的關係密切.
遠遠超過我們事奉的一切.
所作所行.
我們有天父的同在.
就有一切.
是最寶貴的.
教會是神的家.
教會是天父的家.
求天父使我們以.
天父的心為心.
我們不單止領受.
天父的慈愛.
我們去回應,我們去愛人.
我們去尋找,我們去接納.
我們去歡迎,我們去恢復.
我們去關心那些有需要的人.
在你身邊的有誰.

$^{1321}$你想起的,我們去關心他.
我們去探訪他,我們去幫助他.
我們帶他回來,或者.
無論他是信主的,還是未信主的.
求天父幫助我們.
在我們心中對我們的說話.
帶著今天的訊息.
讓我們在當中來回應.
神的愛的呼召.
我們領受愛.
我們在當中去愛人.
我們一起來祈禱.
讓我們今天可以一同來思想.
這個訊息,求你將你的話.
在我們的心中,成為我們腳前的燈.
路上的光,帶我們前行.
讓我們享受和天父的慈愛.
在當中我們有服侍主的喜樂.
和平安.
求主恩代我們,我們再一次多謝主.
心裡面,心靈裡面呼叫.
阿爸父,多謝主.
奉主耶穌基督的名祈禱,阿們.
\newpage



\section{士師記 9:5}
\label{sec:032349}
\textbf{聖經人聖經事:亞比米勒用計殺兄弟作王}~~~~黃克勤牧師博士
\newline
\newline
連結: \href{http://www.church.com.hk/acms/content.asp?site=cdc\&op=show\&type=product\&code=032349\&layout=sermon}{\texttt{032349}} ~~~~ 語音日期: 2026-06-21
\newline
\newline
\hyperref[sec:032348]{\small{< < < PREV SERMON < < <}}
~
\hyperref[sec:index_chronic]{\small{[返順時目]}}
~
\hyperref[sec:032350]{\small{> > > NEXT SERMON > > >}}
\newline
\newline
士師記 9:5
\newline
\begin{longtable}{cl}
\hline
\hline
章節 & 經文 (和合本修訂版)\\
\hline
9:5 & \begin{tabularx}{0.7\textwidth}{X} 他來到俄弗拉他父親的家，在一塊磐石上把他的兄弟，就是耶路巴力的七十個兒子都殺了，只剩下耶路巴力的小兒子約坦，因為他躲了起來。 \end{tabularx} \\ \\
[1ex]
\hline
\hline
\end{longtable}
$^{1}$各位電影小妹,好開心又可以一齊讀聖經.
我自己日日預備這些經文對於我自己很有幫助.
希望對你有幫助.
因為變了日日都讀聖經.
有時一日讀好幾段.
以前都忘記了,現在又重溫.
真是如蜜甘甜.
生命是宣教網絡帶大家一齊看聖經人物,聖經故事.
最重要是見到上帝其實不停地跟我們說.
我的手在後面,在做事,我看著.
讀了這些經文是很安慰的.
今天仍然繼續看《士師記》.
機殿,不過機殿已經垂了,歸到列祖那裡了.
今次看到機殿的兒子.
機殿有一個兒子叫阿比米勒.
這個阿比米勒真的不是好人,是一個很差的人.
是用妓殺了他的兄弟而做王.
他是用鬼妓.
而且你知道一殺就是很殘忍的了.
用鬼妓和殺自己的兄弟.
就是上帝都很想動他的了.
不過都是.
我想都要多少給他機會.
不過上帝的手在後面,在做事.
有些人就失去生命就很不值.
是的,在世界上的生命有長有短.
當然長就好一點.
但長其實就辛苦多一點.
其中經過的不過是勞苦愁煩.
轉眼成空,我們便如飛而去.
所以其實長一點短一點不要緊的.
最重要是敬畏神,覺得有意義開心.
就好了,有意義,服侍主.
好,看看阿比米勒做什麼.
事實記第九章,耶路巴力.
耶路巴力在以色列人當日當中.
耶路巴力比基奠本身的名字更出名.
有時候人們叫我高佬.
叫得多,一說高佬就記得是我了.
一想起我,他就說高佬,或者光頭佬.

$^{41}$高佬和光頭佬這兩個名字.
是不敢在我面前叫的.
但其他人聊天時,全部都這樣說.
當然很多人都講老實話,說我英俊.
但本身的名字反而不知道.
因為耶路巴力這個名字.
反而在基奠是更出名的.
大家都叫他耶路巴力.
耶路巴力大家記得嗎?.
就是讓巴力和他爭論.
讓巴力和他爭辯.
如果巴力是真神,巴力和他談判.
因為他拆了巴力的壇.
把他的太太還是女朋友.
亞瑟拉的木匠斬碎以來做木柴燒.
獻上他爸爸那隻七歲的公牛做獻祭.
於是人們就來尋仇.
把兒子交出來.
他爸爸忽然之間竟然是巴力壇的壇主.
就有這個兒子,就悔改了.
如果巴力是真神,讓他自己和他爭論.
和他出頭,用得著你出頭救巴力嗎?.
於是人們就報世出.
基奠有個新名叫做耶路巴力.
爭辯,巴力和他爭辯.
在上一集看到耶路巴力說他有七十個兒子.
女兒就沒有說.
七十個兒子是說很多妻子.
其實是說上帝很祝福他.
令他兒孫滿滿滿,滿堂.
最重要是養得起.
其中有一個兒子叫阿比米勒.
七十個兒子裡有一個叫阿比米勒.
阿比米勒不是和他們一起住在俄弗拉.
是住在士劍.
由於基奠時時都去士劍.
士劍是當時以色列人的首都.
經常去的,他是士師.
在那裡就娶了一個妾.
妾通常都是服侍婢女.

$^{81}$就娶了妾,不是正妻.
妾生了一個兒子叫阿比米勒.
阿比米勒不是正妻所生.
是妾,就是婢女所生.
由於以色列人問耶路巴力基奠.
你管理我們,你帶領我們.
以色列太平了四十年.
不如你和我們戰勝米殿人.
不如你做我們的王,管理我們.
當時王沒有戰勝米殿人的王國.
只不過是一個地方的首領.
管理者,統治者,管理者.
基奠說我不會管理你們.
我的兒子也不管理你們.
唯有耶和華管理你們就最好了.
他就做了一件金的伊芙特.
可能是掛在那裡,放在那裡.
誰知那些人又拜.
就變成基奠和他家庭的一個缺點.
來到這裡.
他七十個兒子都不做王.
好啊,其實做王都不是那麼好.
做首領真的很大壓力.
其中一個兒子叫阿比米勒.
就是住在他妾時劍的妾所生的兒子.
阿比米勒的意思就是阿比米勒.
阿比米勒,阿比是一個字來的.
米勒是一個字來的.
所以是阿比米勒.
阿比你會發覺有一點像阿爸.
阿爸就是阿比.
阿比就是父親.
米勒就是米勒.
如果米勒是父親.
那麼米勒就是米勒的父親.
又或者父親是王.
米勒就是王的意思.
父親是王.
很有可能阿比米勒.
我不知道,這是我自己猜的.

$^{121}$自己的名字叫父親是王.
父親不在,當然到我了.
阿比米勒有興趣做王.
於是他就去妾劍這個地方.
去到首都.
即是說他到了妾劍.
即是說阿比米勒.
應該不是和媽媽一起住在妾劍.
可能回到娥忽那裡.
是和家人一起住的.
即是和姬甸一起住的.
七十個哥哥一起的.
在娥忽拉可能回去了.
現在他就從那裡出來.
回到他出生的城市妾劍.
在娥忽拉那裡出來.
那麼來到妾劍.
妾劍我們經常聽到很多次.
妾劍,妾劍.
其實就是當時以色列地.
迦南地的他們以色列.
十二個支派的首都.
那麼妾劍.
Shek-em.
可能就由於.
K-E-M.
其實是Shek-em.
Shek-em.
或者有些人叫.
Shek-em.
妾劍就是背脊.
或者兩個肩膀的意思.
見她的.
眾母舅.
即是說她回到.
她媽媽那裡.
不是自己直系那邊.
回到母系那邊.
因為她媽媽是妾劍人.
她回去見她的.

$^{161}$眾母舅.
即是她叫做外家那邊.
對她們和她外祖.
全家的人說了一段話.
即是說.
她刻意去說.
請你們問士劍的眾人.
就是耶路巴力的.
眾子.
七十人一起管理你.
好一點.
還是一個人.
管理你們好一點呢.
還有你要紀念.
我是你們的骨肉.
這兩個幾個問題.
她說什麼呢.
我再解釋.
她回去找到.
她外家媽媽那邊.
的鄉親父老.
我告訴他們.
其實我和你們.
是骨肉.
因為.
我是你這邊出生.
和我媽媽是這邊的人.
所以.
我是你們的骨肉.
她說得.
強調一點.
如果是中文.
其實我是你們的骨.
與肉.
骨.
這個字比較難讀.
骨.
肉.
我是你們的骨.
骨中的骨,肉中的肉.

$^{201}$好像.
始祖.
骨中的骨,肉中的肉.
即是說.
好親好親好親.
和我自己.
這個人身體一樣.
因為我在這裡出生.
我是你們的家人.
我和你們是骨肉,骨與肉.
你幫我去.
問一下事件的人.
究竟是.
如果要人管理.
整個以色列.
究竟是70個人.
即是耶路巴黎70個兒子.
一起管理好一點.
還是一個人管理好一點.
還用問嗎?.
當然是一個人管理好一點.
70個人怎麼聽,開會都來不及.
而且記得.
我是事件人.
而且你們幫我去說話.
因為我是你們的骨與肉.
骨中的骨,肉中的肉.
於是眾母舅就說.
原來是這樣.
於是.
他們就去和事件人說.
其實這段說話是說什麼呢?.
我想.
做你們的王.
做你們的管理人.
其實你們以前就拍照.
因此你們是和.
我爸爸機電說.
你不如管理我們.
不過爸爸說不好.

$^{241}$然後他說我的兒子.
都不會管理,但是.
我作為他的兒子,也是事件人.
我想管.
我想做王.
好不好,行不行?.
請你記得,其實我不是外人.
我是你們骨與肉的人.
哦.
而且我是耶路巴黎機電的兒子.
這樣.
於是你和我說,你問一下.
如果一個人管我們,機電有70個兒子.
那麼你就會反駁他.
70個人管理好還是一個人管理好.
當然是一個.
那麼在70個裡面,選一個都可以.
怎麼選,70個裡面.
我就是這個本地事件人.
你們就是事件人.
那即是怎樣,選我吧.
我自薦.
用得著的.
於是他的眾母舅.
就和事件人說.
起碼和領袖說.
事件的人就說.
也好,和.
阿比米勒談得來.
而且還說.
有原因的,他那70個.
也是好.
也是機電的兒子,不過.
這個阿比米勒.
是我們兄弟.
弟兄骨肉來的.
所以這個最適合.
其實我們也有一個人管好,不是很亂.
好了.
那麼他們就由巴黎.

$^{281}$比利時的廟裡面.
拿了70赦赫勒銀紙.
就交給.
阿比米勒.
那麼阿比米勒就用這些.
70赦赫勒的銀紙.
就雇了一些匪類.
來跟他,所以這個阿比米勒.
也不是好人來的,你看到.
他很多那些匪徒.
朋友的話,如果你叫我.
找個匪徒朋友.
我們這些正人君子沒有.
匪徒朋友的嘛,好了.
而且那些錢,我們給.
錢他,不如你做我們的王.
那麼阿比米勒可能就說.
好啊,不過我也要.
多少洗費,你們就不能夠.
幫幫我,給些洗費我先.
好啊,那麼我們最有.
錢的地方,公家錢就是.
巴黎比利土的廟.
裡面的銀庫.
大家記得上次說.
比利土是什麼嗎?巴黎.
巴L是主.
比利土就是.
細約,細約的主.
那麼到處整個.
迦南都有很多巴黎,珠巴黎.
那麼在他們那裡拜的.
巴黎就叫巴黎比利土.
那麼意思就是細約之主.
那麼在廟裡面.
拿七十赤赫勒.
那麼大家記得嗎?一赤赫勒.
就大概是十.
十四.
十四克銀子.

$^{321}$十四克,那麼就.
大概都有.
一個重量的銀就給了阿比米勒.
阿比米勒就用這筆錢.
是沽了些匪類,所以.
都不是一個很少的錢.
那麼就找了一些歹徒跟他.
跟他賺錢.
有馬的.
他就往俄弗拉.
這個俄弗拉就是他.
爸爸那裡,自己.
自己爸爸家裡.
正話事劍是媽媽家裡.
就去了他父親家裡.
就叫齊兄弟.
各位弟兄,我們開一個.
家庭會議.
兄弟會啦,我們有些.
重要的事要談,就叫.
他們出來,就去哪裡?.
去一塊盆石那裡.
一起談,一塊大石.
在上面開會.
奇奇錯錯,忽然之間.
有很多人走出來.
就將耶路巴力的七十個兒子.
即是紀念的七十個兒子.
七十人都.
殺在一塊盆石上,殺死了他.
有一個兒子.
是躲藏了.
就剩下耶路巴力.
最小的一個小兒子.
叫約坦.
就逃脫了,逃脫了.
為什麼呢?是因為.
他躲藏了.
是不是殺人的時候.
他躲起來呢?不是.

$^{361}$一塊石都不能躲起來.
而且會數人的,七十個殺光了.
其實應該這樣看.
和看廈文.
約坦應該是沒有去.
很有可能.
他沒有去.
小兒子你不要想.
他是小孩子,不過他可能沒有去.
剩下最小那個.
沒有去,約坦.
又分.
又分呢?.
又分的意思是耶和華.
是完美的.
又分,耶和華是完美的.
約坦.
當然你說.
廣東話可能變成約番.
又分,那就沒有去.
他的名字的意思.
就是耶和華是完美的.
如果改名.
改這個英文名都不錯.
如果又呢.
這些字如果翻譯成英文.
就變成J.O..
好了,不要緊,不說那些了.
約坦就.
沒有在.
士劍人和米羅人.
都一同聚集.
士劍人就是我剛才說.
以色列當時的首都.
那裡的人就是.
士劍人,由於他們是首都的領袖.
於是他們就一起.
做這件事.
米羅人其實就住在耶路撒冷.
當日已經有耶路撒冷.

$^{401}$不過就不是首都.
離開都比較遠.
就是在那裡.
就一起約.
不過可能耶路撒冷都開始逐漸.
變成一個都幾重要的地方.
上帝就祝福你在這個地方.
城市的人就開始突出.
兩個城市的人一同聚集.
就說我們不如.
在士劍一棵橡樹旁邊.
的柱子.
那柱子是什麼意思呢.
其實很有可能.
就是一個廟的柱子.
就在那裡.
這個廟會不會是會幕.
當然不是.
有這樣的柱子,即是說附近有神像.
都可能是巴黎廟.
很差的.
就在那裡一同.
納瓦比米納造王.
其實士劍人.
自己納瓦比米納造王.
可不可以呢,都可以的.
不過拉多些人一起.
都很多人贊成,當然有人承認.
就造王.
這個王就不是.
指一個王國是造哥利安王.
所以你看到機殿對付米殿人.
都殺了很多個王.
是一個統領.
一個領袖.
造了王.
這個就是納瓦比米納.
用計殺兄弟造王.
這件事今天就到這裡.
你看一下.

$^{441}$納瓦比米納是個什麼人.
這麼差的.
自己想造王.
爸爸.
他們請過爸爸造王.
爸爸不肯.
就說我的兒子都不會造王.
很聰明,機殿真的聰明.
因為知道這些位置不好造.
自己造過,做過領袖.
我不會造.
不如讓耶和華管理你們.
這樣很好,即是說他很敬畏神.
讓神自己造.
我都講過,其實聖經裡面.
說不是人主.
是神主.
是神作主.
凡事求民作主.
但是納瓦比米納不是這樣的人.
很有可能.
納瓦比米納本身就知道.
耶和華是神,不過.
他都很傾向巴力那邊.
所以他找的人.
就算按納.
高納去造王.
都去處子那邊.
很有可能是巴力的廟.
如果我是敬畏神.
我都想造王.
我都不要在那邊,找另一塊盆石.
去處子那邊.
有代表性的地方.
就覺得在那邊好.
而且他用了毒計.
請匪徒一起陰他七十個哥哥.
一次殺了.
七十個兄弟.
不一定是哥哥,有弟弟.

$^{481}$七十個,很殘忍.
很奸.
很殘忍,他會用計.
用說話,會說話.
會說話.
就是七十個人管理你好還是一個人好.
當然是一個.
那個就是我,因為我就是你們.
使劍人來的,是骨與肉.
骨中的骨,肉中的肉.
所以是我最好.
而且我造王.
你們使劍人都光彩.
使劍人出來造王.
不是,是俄弗拉人.
好啊,於是就這樣做.
於是大家串起來.
再拉一些.
米羅人一起串起來.
就殺了那些兄弟.
自己造王.
很差很差很差的人.
有時候.
今天我們都會覺得.
為什麼我們.
國家裡面,我們社會裡面.
有這麼差的人做管理呢?.
為什麼.
我們教會裡面有這麼差的人.
做管理呢?.
為什麼裡面有一些.
好像王克勤這麼差的人.
做的呢?.
這麼差,這麼卑劣.
的人呢?當然.
我們都會說其實我都很卑劣.
個個都是,但是.
有些人自己會.
這樣做了.
管理者.

$^{521}$又不自己反省.
不過.
上帝就讓他做了.
上帝計數的.
你一直下去,上帝是計數的.
這個亞比米勒的結果.
都很不光彩.
很差的,很慘的.
神是會知道的.
有工作的,在聖經的.
日與時的背後.
上帝仍在掌管.
希望我們都有這個信心.
雖然有一朵黑雲.
叫亞比米勒遮了以色列.
不過.
神是在的,祂沒有離開.
讓我們一起禱告.
多謝我們的主.
當聖經裡面.
烏西亞王死的時候.
當時是一個很好的皇帝.
希望以色列國在最慘的時候有救.
誰知道很早.
短命駕崩.
以賽亞先知都非常.
慘痛.
進到聖殿之中.
誰知道見到上帝高高坐在.
寶座上.
衣裳垂下遮滿聖殿.
表示你仍然.
眷顧屬於你的.
奠禦百姓.
在這裡見到亞比米勒.
當日他用妓女.
殺了70個兄弟而造王.
拿到.
使劍以色列.
這個領袖的位置.

$^{561}$神就讓他做一段時間.
看看會不會有什麼.
給個機會,但是他沒有改.
仍然做得很差.
結果上帝.
沒有就這樣放過他.
我們知道你仍然作王.
你仍然掌管,讚美你.
自己的大能,讚美你色在.
金在,以後永在.
又真又活的神.
願你得著讚美敬拜.
祈禱奉耶穌名求.
阿們.
\newpage



\section{申命記 10:12-22}
\label{sec:032350}
\textbf{感恩 MODE ON (3) 喂,係時候「落場」喇!}~~~~謝又生牧師
\newline
\newline
連結: \href{http://www.church.com.hk/acms/content.asp?site=cdc\&op=show\&type=product\&code=032350\&layout=sermon}{\texttt{032350}} ~~~~ 語音日期: 2026-06-20
\newline
\newline
\hyperref[sec:032349]{\small{< < < PREV SERMON < < <}}
~
\hyperref[sec:index_chronic]{\small{[返順時目]}}
~
\hyperref[sec:032351]{\small{> > > NEXT SERMON > > >}}
\newline
\newline
申命記 10:12-22
\newline
\begin{longtable}{cl}
\hline
\hline
章節 & 經文 (和合本修訂版)\\
\hline
10:12 & \begin{tabularx}{0.7\textwidth}{X} 「以色列啊，現在耶和華－你的神向你要的是甚麼呢？只要你敬畏耶和華－你的神，遵行他一切的道，愛他，盡心盡性事奉耶和華－你的神， \end{tabularx} \\ \\ \relax
10:13 & \begin{tabularx}{0.7\textwidth}{X} 遵守耶和華的誡命律例，就是我今日所吩咐你的，為要使你得福。 \end{tabularx} \\ \\ \relax
10:14 & \begin{tabularx}{0.7\textwidth}{X} 看哪，天和天上的天，地和地上所有的，都屬耶和華－你的神。 \end{tabularx} \\ \\ \relax
10:15 & \begin{tabularx}{0.7\textwidth}{X} 然而，耶和華專愛你的列祖，愛他們，從萬民中揀選你們，就是他們的後裔，像今日一樣。 \end{tabularx} \\ \\ \relax
10:16 & \begin{tabularx}{0.7\textwidth}{X} 所以你們的心要受割禮，不可再硬著頸項。 \end{tabularx} \\ \\ \relax
10:17 & \begin{tabularx}{0.7\textwidth}{X} 因為耶和華－你們的神是萬神之神，萬主之主，是偉大、強有力、可畏的神，不看人的情面，也不受賄賂。 \end{tabularx} \\ \\ \relax
10:18 & \begin{tabularx}{0.7\textwidth}{X} 他為孤兒寡婦伸冤，愛護寄居的，賜給他衣食。 \end{tabularx} \\ \\ \relax
10:19 & \begin{tabularx}{0.7\textwidth}{X} 所以你們要愛護寄居的，因為你們在埃及地也作過寄居的。 \end{tabularx} \\ \\ \relax
10:20 & \begin{tabularx}{0.7\textwidth}{X} 你要敬畏耶和華－你的神，事奉他，緊緊跟隨他，奉他的名起誓。 \end{tabularx} \\ \\ \relax
10:21 & \begin{tabularx}{0.7\textwidth}{X} 他是你當讚美的，是你的神，為你做了大而可畏的事，這些是你親眼見過的。 \end{tabularx} \\ \\ \relax
10:22 & \begin{tabularx}{0.7\textwidth}{X} 你的列祖七十人下埃及，現在耶和華－你的神卻使你如同天上的星那樣多。」 \end{tabularx} \\ \\
[1ex]
\hline
\hline
\end{longtable}
$^{1}$弟兄姊妹晚安.
晚安.
老人家聽歌不是太好.
可能聽不到,所以大聲一點.
好,我們先祈禱.
主一讓你今晚將你的心意向我們每一位弟兄姊妹顯明.
是神你自己親自來到我們當中和我們講說話.
讓我們更加明白到並且擁抱到你的心腸.
你對我們所有人,你對你的群體.
你對你所愛的兒女這一份的心.
讓我們真的擁抱明白到並且去回應你.
求你的靈繼續在我們當中親自和我們每一位講說話.
我們每一個都是你所塑造,都是不同的.
但我們卻是有一個心,我們很想很想認識你.
你監察我們的心,帶領我們,讓我們更加明白你.
在我們當中,我相信也有相當的弟兄姊妹.
你回到教會聚會也有不短的日子.
甚至有些可能是你很小的時候,爸爸媽媽帶你回來聚會.
可能你聽到的也比在YouTube上看影片更多.
但坦白說,不知道大家有沒有試過覺得你的信仰.
就好像看Netflix一樣.
看滿片也沒辦法,但就算你看了一套很精彩的電影.
過程你也覺得很爽.
但看完之後,你會發覺生活一樣照舊.
上班一樣拼搏,一樣掙扎.
拍拖的時候還是繼續人選你,你選人.
又或者阿男叫你快點多點事奉,投入教會事奉.
你會仍然繼續說,先看清楚一點,先看看怎樣.
如果當中真的有這樣的情況,又或者好像有人這樣形容.
如果有這樣的心態的話,這種就是一種屬靈上的消費主義的心態.
我們今晚說申命記第十章.
申命記第十章就是摩西對著一群曠野第二代的以色列人說的話.
又或者摩西說的話,今晚值得我們有深刻的反思.
所以今晚說到的語調有一點點不像感恩mode.
可能會有一點點挑戰mode.
所以可能大家要預備一下.
如果挑戰得太多就不要介意,不要怪我.
這個也是聖靈感動我所說的.
所以感恩得來,我們也要挑戰一下大家.
神帝執著摩西和當時的曠野第二代以色列人這樣說.

$^{41}$他說你不要再做聽道KOL了.
現在是世界盃,每個人都想著看球賽.
不過神叫我們不要光看,要下場踢真球了.
各位頂智妹,你準備好下場踢真球了嗎?.
剛才跟弟兄說,他又打羽毛球又踢球,他又踢球.
但是神叫我們在我們的信仰上真的要下場踢球.
今天你準備好了嗎?.
我首先說申命記是怎樣來的呢?.
申命記其實是劇情發生的時候.
以色列人已經在曠野裡度過了40年.
舊一代已經過去,新一代興起.
這群人準備進入神應許之地,也是迦南地.
怎樣預備呢?.
自己的心態要預備,自己的心靈也要預備.
跟神的關係也要預備.
摩西對著這群人說你要重新.
重新將神當初賜給以色列人的律法再說一次.
邀請他們重新跟神立約.
重新肯定對神的毀身.
以致他們能夠預備將來的挑戰.
頂智妹,人最面對的挑戰跟當時的以色列人一樣.
最大的挑戰未必是困難.
當時以色列人那40年面對的挑戰就是那些困難.
就是過紅海,在路途中遇到強大的外族要打仗.
又或者經歷沙漠那些嚴峻的氣候的挑戰.
但最大的挑戰不是這些.
最大的挑戰是當這群以色列人進入了迦南地應許之地之後.
定居了,一切看似安好,無風無浪.
而帶來的種種試探,種種誘惑.
如果我們要看,要講整個申命記.
我們必須要講兩節的聖經.
就是申命記的六章四節五節.
他說什麼,以色列啊,你要聽.
耶和華我們的神是獨一的主.
你要盡心盡性盡愛耶和華你的神.
我相信這裡很多頂智妹都講過N次了.
有很多人你送給人的時候,寫的金句都是寫這句送給人.
但是頂智妹,這個不止金句這麼簡單.
其實他是整個申命記裡面的地基和核心晶片.
是核心中之的核心.

$^{81}$為什麼呢,因為他一開始就劈頭一句.
他宣告神對以色列人有獨家代理權.
什麼叫獨家代理權,那個關係啊.
獨家代理權就代表.
神叫那些以色列人進了迦牟尼之後不要跟別的大佬.
你不要拜那些巴力,不要拜那些亞瑟拉的偶像.
這個世界的大佬就只有一個,就是耶和華自己.
當他宣告了這個獨家代理權之後.
後面的那些律法,那些警告.
就全部是懷疑我們怎樣才是向這個唯一的大佬效忠呢.
我們怎樣才是對這個真神毀身呢.
用這個世界的term是怎樣才是對這個唯一的大佬真的效忠呢.
根據這個地基,根據這個核心的晶片.
以下申命記的第十章,十二至二十二節.
他提出了三點.
這三點是在說我們怎樣才是向這個唯一的神效忠.
第一點,神說我不是要你like瑤慧,而是要你的life.
我再說一次,祂不是要你like她,祂是要你的life.
第一段,申命記的十章十二至十三節.
他說以色列現在耶和華,你的神要向你要些什麼呢.
就是要敬畏耶和華你的神,遵行祂的道,愛祂,盡心盡性,事奉祂.
遵守祂的誡命律例,就是我今天所吩咐你的,為要叫你得福.
神一開始就說,我要些什麼呢.
我不是要你like我,我不是要你在互聯網多給幾個like我.
我也不是要你在小組崇拜的出席率爆燈.
我也不是要你做到perfect,我要什麼呢,就是要你整個的生命.
整個的生命是什麼呢,那段聖經說了,是敬畏神和愛神.
什麼是敬畏神呢,敬畏神很簡單,八個字而已.
以神為神,以人為人,這個世很多人都混亂了.
人現在最大的問題是什麼呢,以人為神.
什麼叫以人為神呢,明明是人,我當我是神.
如果我是神的話,我喜歡做什麼就做什麼,喜歡怎樣就怎樣.
我喜歡任意妄為,這個是以人為神.
以神為人是什麼呢,又是相反的.
將神降格成為我操控的工具.
神在我生命裡面唯一的目的呢,就是令我榮華富貴,心想事成.
大家見過那些偶像就是這樣.
觀音是借父,他是令我發達而已.
他的目的,他存在的意義,是成就我的心意.
所以以神為人和以人為神,相反的,最後的終極的意思是.

$^{121}$是神聽我的話,不是我聽神的話.
所以是相反的.
所以神要我們真心敬畏他,以神為神,以人為人.
第二樣東西,什麼叫愛神呢.
聖經裡面的愛不是一種感受性的.
什麼叫感受性,我有feel就愛,我沒有feel就不愛.
這種是感受性的.
但是聖經裡面講的愛神是命令式的.
就是我們憑意志去愛.
可能大家也試過,你的際遇不太順利,頭頭碰著黑.
沒運行那些時候,你叫愛神先想一想.
那些就是了,你feel不到就不去愛.
但是神要求我們那種命令式的愛就是.
我們雖然有時候不太順利,有時候心情不太好.
我們仍然決定去愛.
但是我們講回頭,敬畏神和愛神,其實是有一點抽象的.
什麼叫愛呢?什麼叫敬畏呢?.
所以可能到最後實踐的時候,很多人的理解都不同.
所以神定了一個客觀的標準.
什麼叫客觀的標準呢?.
客觀的標準就是你用具體行動去表示你的愛.
就是遵守神的命令.
這個不只是舊約聖經說的,記得耶穌說過什麼嗎?.
你若愛我,就必遵守我的命令.
所以神的吩咐是一致的.
什麼叫敬畏神,什麼叫愛神呢?.
就是用具體行動去遵守神的命令.
我舉個例子,你抱著B,大家都結婚了.
有些可能拍拖了.
你結了婚,你弟兄姊妹,我問你,你愛不愛你老婆?.
真的愛?.
真的愛吧?.
有些拍拖,我不知道你女朋友男朋友有沒有來,你愛不愛他們?.
不敢說,因為沒有拍拖.
你愛不愛啊?.
你記住了,現在是在說什麼呢?.
愛是很抽象的,什麼叫愛,I love you, I love you too.
那些是很抽象的.
但是神說愛是什麼?具體行動.
老實說,你記不記得結婚週年是哪一天?.

$^{161}$你先結婚老婆生日,你先結婚老公生日.
這些是具體行動,還有.
你在實際的日子,你會送花,你會請她吃飯.
你會送禮物給她.
大家可能很多都是新婚.
我今年結婚30週年.
我們習慣有時老婆都會說.
年年都是送那些,年年都花那麼多錢.
今年不要送了.
你信不信?.
當然不信了,你不信了.
如果你信了,你就糟糕了.
其實她在說什麼呢?.
具體行動是什麼?.
你既然是說愛,你不是說I love you那麼簡單.
當然這個是很簡單的例子.
人生裡面有更多的行動去表達你的愛,用具體行動.
所以神都是這樣,用具體行動去表達你的愛.
為什麼我們不要忘記這些具體行動?.
因為這一切的行動都是指向一樣東西.
我把我整個心都給了你.
當太太收到你的花的時候,你記住.
其實是在跟你說話,不是摘花.
我整個心都給了你.
我整個人生都給了你.
這是我對你的盟約.
所以我永遠都記住.
神就是在說這件事.
為什麼要具體行動,遵守祂的命令呢?.
就是要用具體行動告訴神.
我將我的人生交給你.
將我的心交給你.
就是這樣的意思.
其實呢.
我也養了一段時間,我曾經是YA.
我以前認識很多弟兄姊妹一直這樣成長.
我之前認識一個三十多歲的弟兄.
他在一些大機構做經理.
經常都說他的工作很忙.
但是他很乖的,每個禮拜日都回來崇拜.

$^{201}$他真的坐足整座崇拜.
不過他一直聽廣道的時候,他就在手機裡開著BOOM BOOM.
在看著市場.
他就覺得我回來都坐足了,都交到功課了.
但是呢.
我很相信神跟他說,我不是要你的attendance.
我是要你的心.
是要你的敬畏.
什麼叫敬畏都用回我的terms.
是尊敬神,尊敬到一個地步.
你每一個決定.
時間怎麼過,你的錢怎麼花.
都問過神先,主啊,你想我怎樣呢?.
大家明白我說什麼吧,什麼叫敬畏神.
尊敬神到一個地步就是.
在你人生裡每一個決定,是要學的.
每一個決定,時間怎麼過,錢怎麼花.
你都先問過神先,主啊,你想我怎樣做決定呢?.
這個就是敬畏神.
所以今天神都對大家有一個挑戰.
神對你的挑戰就是,你現在的生活裡有些什麼是你做主.
有些什麼是不給神做主的呢?.
會不會是你的工作呢?.
你現在的relationship呢?.
或者你的錢怎麼花呢?.
神是要你的life,他不只是要你like他.
like他可能就是我禮拜六回來坐兩個小時.
禮拜日回來坐兩個小時.
但是life就是你的人生,你的方向.
你是不是以神為先呢?.
這是第一個,神不是要你的like,是要你的life.
第二,第二點就是.
葛禮是換個心,不是做個樣.
我們繼續看申命記,在這裡看下去.
十章的十四至十六節.
他說什麼看啊.
天和天上的天,地和地上所有的.
都屬耶和華你的神.
耶和華誕喜悅你的列祖愛他們.
從萬民中揀選他們的後裔.

$^{241}$就是你們,像今天一樣.
所以你們要將心裡污穢除掉.
不可再硬著頸看.
摩西要意思是將心裡污穢除掉.
那是什麼呢?.
用原文呢,如果我.
將他的意思逐個希伯來文直譯是什麼呢?.
你們閣下包庇你們的心.
大家都知道葛禮是什麼吧,我不詳細說了.
你們閣下你的包庇你們的心.
你拼在一起的意思.
如果大家記得在新約聖經保羅說過相類似的東西.
就是心裡面的葛禮.
不是肉身的葛禮.
是心靈裡面的葛禮.
其實以色列人可能大家知道歷史了.
以色列人每個BB出生第八天都要行葛禮.
就是割了包庇.
原因就是神和亞伯拉罕有一個亞伯拉罕的盟約.
就是你的子孫世世代代.
就是我的指紋,神的指紋.
那個標記就是這個葛禮.
但神說不是.
不只是你的身體要行葛禮.
是你的心要行葛禮.
對於我們今天說.
肉身的葛禮對你可能沒什麼意義.
但對於我們今天可能有很多的標記.
我們都有的.
不如你受一浸吧.
可能你的教會有些事奉.
但是會不會有機會當我們正在做這些事.
擁有了這些標記的時候.
我們的心裡面都是很冷.
很冷漠.
很剛硬.
冷到可能你自己都震驚自己.
我剛才說我都牧養了很久的時間.
我都見過很多很多的基督徒.
我以下說的是負面例子.

$^{281}$我相信不是在大家身上發生.
但是我覺得我應該說的是.
去表達一件事.
我們不是要樣貌的.
我們要真實.
神要我們那個真實的生命.
我曾經見過.
我想你都見過那些基督徒借飯.
吃飯的時候借飯.
你見到那個人吃飯之前祈禱.
那是基督徒來的.
但是有很多人這些都是一些浮面的.
表面的東西.
我認識的就是他一樣借飯.
很殷勤的樣子.
但是回家之後.
就對那個工人姐姐呼呼喝喝.
他一樣在群組chat的背後.
在狂插他的同事.
又或者他見到一些貧窮人的時候.
他沒有什麼反應.
他說不關我的事.
都不是我製造出來的.
不關我的事.
這些我好感覺.
他的心好像被皮包著.
他的真正的心好像被冰封了.
好多症狀.
譬如聽到滑了牙.
又或者可能思覺很感動.
他覺得很感動.
不過沒有眼淚流.
又或者見到身邊的人很有需要的時候.
無動於衷.
我感覺有些人的心是僵硬到一個地步.
不知道會不會有些人會是這樣呢.
有時候發覺自己的心僵硬到.
你對人是沒有什麼敏銳感的.
你對人沒有什麼敏銳感的時候.
對神都應該沒有什麼敏銳感.

$^{321}$等子妹.
我剛才說的.
將心裡的污穢除掉.
就是這一堆東西.
其實神的心是很想我們.
行心靈的國禮.
什麼意思呢.
破碎我們冰冷的心.
破碎我們那些冷漠.
讓神幫助我們.
那個心不要被皮包著.
神要幫我們將它變為肉心.
如果大家記得舊約聖經.
聖靈來到是將我們的錫心變為肉心.
今天我們都要檢視一下自己.
會不會我們在信仰日子裡面走這麼久.
我對身邊的人.
我對這個社會.
我對有需要的人.
好像沒有了那份感覺.
等子妹.
神要我們除掉這些污穢的心.
不要讓這些污穢.
難了我們去給神使用.
可以幫助.
可以服侍到很多人.
這是第二個.
心靈的國禮.
是讓我們改喚心.
將我們的心變得柔軟.
變為肉心.
不是將樣子扮出來.
第三點.
申命記繼續下去十章.
十七至十九節.
意思是.
神是特別關顧弱勢社群的.
今天神說我很關顧弱勢社群.
你跟不跟神呢.
你跟不跟神這個心呢.

$^{361}$你跟不跟神這個呼召呢.
我們繼續看.
心靈記十章十七至十九節.
他說什麼呢.
因為耶和華你們的神.
他是萬神之神.
萬主之主.
至大的神.
大有能力.
大而可畏.
不以貌取人.
也不受賄賂.
他為孤兒寡婦伸冤.
又憐愛寄居的.
賜給他衣食.
所以你們要憐愛寄居的.
因為你們在埃及地也作過寄居的.
這段我覺得有一點瘋狂.
為什麼呢.
因為開頭說神不以貌取人.
也不受賄賂.
接著一下子就跳出去了.
但是神為孤兒寡婦伸冤.
又憐愛寄居的.
接著就一句縫起來了.
所以你要憐愛寄居的.
因為你們在埃及也作過寄居的.
其實他在說神不是大世超.
不是玩偏心.
但是神真的特別愛.
特別憐憫那些弱勢社群.
他是呼召我們服侍弱勢社群.
因為我相信在我們人生裡面.
你也有試過被欺負.
有試過被欺負.
有試過處於弱勢.
所以神的心是憐憫這些人.
我們每一個人的人生都有機會是這樣.
所以神的心是呼召我們.
同一樣擁抱他的心.

$^{401}$去關顧一些弱勢社群.
好了 今天我相信.
我們當中有一些弟兄姊妹.
可能你在供樓.
可能爸爸媽媽資助你.
又或者可能你還是住在爸爸媽媽那裡.
可能有些住公屋.
有些我不知道有沒有人在供納米樓.
很多很多原因.
但是始終我們也有一種幸福感.
因為有瓦遮頭.
大家知不知道.
當我們在這個比較上幸福的環境裡面.
這個世界還有很多比我們差的人.
有露宿者.
我想大家都見過了.
不過當然不是在這個地區集中.
大家知不知道.
我們明年會重啟東區異象.
大家如果沿著英皇道一直走下去.
在北角那些舊樓那些舊區.
去到西灣河 筲箕灣那些舊樓.
有無數的tong 房.
大家有沒有去過tong 房.
tong 房呢.
它不只是小那麼簡單.
現在呢 今天開始又熱了.
陽光普照35 36 37度.
這些tong 房是40度以上.
所以不懂得開冷氣.
省電費.
那些人吃貧.
電費是很貴的.
所以他們不會開冷氣.
可能晚上開一點點.
在那裡涼了又關掉.
是這樣省的.
所以白天呢.
可能超過40度 43 44度.
那些人是不會回家的.

$^{441}$在外面蒲.
是有些這樣的人.
大家知不知道香港還有.
沒有屋區的.
沒想過是不是.
我一會兒那個例子再說.
一樣有沒有屋區的.
大家知不知道香港有很多.
商居 獨居長者.
今天才有一單新聞.
有一個長者在家裡死了.
沒有人知道.
特別呢 可能大家都知道.
這幾年移民潮.
幾十萬人移民.
大家知不知道呢.
移民了一些人.
放下了多少老人家.
上一代在香港.
是沒有人管的.
所以呢.
有很多獨居 商居長者.
其實同福都服侍很多.
這個地區.
大家知不知道那些.
老人家在家裡.
如果真的沒有人理他.
他真的死了也沒有人知道.
所以呢.
當我們好像很幸福的時候.
神就是憐憫這些人.
他憐憫這些弱勢的社群.
沒有人看到的.
大家有沒有見過那些少數屬裔.
少數族裔被人歧視的.
我在柴灣的時候.
我見過一個.
南亞裔小朋友那麼高的.
無端端被一個伯伯罵了十分鐘.
原因不是他頑皮.

$^{481}$他沒有頑皮.
因為他是不同種類.
被一個伯伯.
我不知道那個伯伯是不是傻傻的.
罵了十分鐘我才去干涉.
我也有一點害怕.
害怕他被我的伯伯打.
但是.
你真的搞不定.
你看到這樣的情況.
你有沒有見過那些清潔姐姐.
那些漢奸伯伯被人當透明的.
很多.
其實很多這些的遭遇.
我講我自己的一些經歷給大家聽.
我有一次去一個公屋單位.
去探望一個獨居婆婆.
那你們猜一猜.
整程一個人去探望一個獨居婆婆.
大家究竟在談什麼呢.
你猜不到我們在談什麼.
我進去那個屋裡.
那個婆婆住在公屋.
我發覺她的牆.
那個司令那裡.
掛了很多個膠袋.
那婆婆你知不知道她的膠袋.
她載著什麼.
載藥.
婆婆不是賣藥的.
大家知道老人家會去醫院看病.
全部都是膠袋.
婆婆做什麼呢.
逐個膠袋拿下來.
開了藥.
問我這是什麼藥.
那我想一想.
老實說我要承認自己的軟弱不足.
因為除了Panadol之外.
其他藥我完全不知道要喝什麼.

$^{521}$可能如果大家有些醫務界的翹楚.
可能我帶你去好一點.
但我完全不知道是什麼.
但我心裡想一想.
你問我做什麼.
我又不是醫生.
又不是什麼專業人士.
但我一邊想一邊突然間明白了一件事.
其實她問什麼藥 我回答什麼藥.
究竟說什麼不重要.
因為婆婆很孤單.
其實她拆東西下來.
給藥物問我.
其實談話的內容不重要.
是有人和她一起.
是有人和她聊天.
那一刻我特別體會到.
神真是很看到人那種心靈的需要.
婆婆需要什麼.
不是談什麼.
內容不重要.
而是有人聽她說話.
或者和她說話.
大家有人與她同在.
我覺得神真是很明白.
很看到人生命的需要.
我剛才問大家知不知道香港還有木屋區.
有沒有人知道.
沒有人知道香港還有木屋區.
大家知道亞高岩在哪裡嗎.
沿著英皇道一直去到筲箕灣道尾.
有一條大斜路.
叫柴灣道大斜路上去.
沒多久轉左去到以前的明華大廈.
走到盡頭.
有一間養和醫院的診所.
幾十層樓那麼高.
你再走進去.
就是那個叫亞貢岩村.
是木屋區.

$^{561}$去年我做門道訓練一年多.
差不多完結後的兩個多月.
我就挑戰門道訓練的弟兄姊妹.
我們一起去探望亞貢岩村.
亞貢岩村.
你見到.
這班弟兄姊妹後面的就是.
有些石 有些木屋區.
現在不是很流行木屋區.
不過叫木屋區.
鐵皮和晶片疊在一起.
沿著山一直上去.
如果有高一點的弟兄姊妹.
可能Hugo那些.
分分鐘你和我去.
有機會會撞頭.
有些洞一定會撞到頭.
就是這樣.
我就和門訓的弟兄姊妹.
我brief他們當地的情況.
有些門訓的姊妹.
她說 哎呀.
我有潔癖.
我去那裡又怕老鼠又怕小強.
那裡又骯髒.
哎呀 我不知習慣不習慣.
那我說.
有時人生的事.
你在安書區都要走出來.
都要突破.
然後你才經驗到神與你同在.
舉個例子 我之前.
都和他一樣.
我自己也有潔癖.
我都一樣走出來.
突破我的安書區.
我記得第一次去探訪的時候.
耶穌吩咐.
路加福音說你進去探人.
傳福音.

$^{601}$主人給你什麼食物.
你都要吃.
我就literally跟著聖經.
進去後.
那個阿姨給了我一個蛋糕.
我拿著西餅.
下面有隻螞蟻.
你吃不吃.
我聽聖經吩咐.
我照吃了.
我又沒事.
我的point是什麼.
突破你的安書區.
你覺得你這個圈裡面comfortable.
但你走出來.
你才經驗到神.
you take the risk.
你才經驗到神.
如是者.
大家一起去.
剛才見到的照片.
就是我們門訓的弟兄姊妹.
和我們有些外展的弟兄姊妹.
一起去探訪.
那天就很特別.
剛才那位姊妹.
潔癖的那位姊妹.
我們探訪完之後.
就去附近坐下.
有個deep breathing.
那位姊妹說.
哎.
我沒想到.
我這麼潔癖的人.
那天走得很舒服.
走得很開心.
我覺得.
原來我是可以突破安書區的.
原來祝福人是這麼開心的.
記住.

$^{641}$探人不是很容易.
給大家看一張照片.
見不見到我背部.
我不是特意站在這裡.
是因為那間屋是小到.
沒位置坐.
沒有地方.
就好像我們招呼人進來.
你又坐一個位.
人家又坐一個位.
就算物質質都不要緊.
我們幾個和一位姊妹去探訪.
我們對著主人.
我們只是站在那裡.
一個小時.
因為沒位置坐.
就是這樣的情勢.
但是那群弟兄姊妹探訪完之後.
發覺.
我們被神洗完很開心.
另一位門訓的弟兄.
他整個門訓他都說.
我是I仔.
不和別人說話.
對著我一個小時.
我可以不和你說話.
沒問題的.
我很舒服的.
那天.
他和一個拍檔去探訪.
結果.
他講兩句.
大家講兩句.
那個屋主就信了耶穌.
他很開心.
然後debriefing的時候.
原來我也可以.
我也可以給世人使用.
我這麼I也帶了人信耶穌.
所以整件事是在說什麼呢.

$^{681}$原來我們配合到神的計劃.
原來神可以使用我們.
我們都一樣可以成為祝福人的出口.
神很特別眷顧.
那些弱勢社群.
今天他也呼召我們.
今天神也給大家一個挑戰.
在你們的小組裡面.
或者你的生活圈子裡面.
是不是整體來說是同聲同氣的呢.
你的social economic status.
是不是差不多一樣的呢.
你會不會有時間去認識一下.
或者去愛一些和你不同的人.
你會不會接受這個挑戰呢.
不是叫你做救世主.
而是讓你學習神怎樣看待人.
神是怎樣看待一些被遺忘的人.
一些可能大家忘記了他們的人.
神真的看到這些人的需要.
神讓我們學習去關顧他們.
耶穌基督的心腸就是這樣.
我們為自己祈禱.
求神幫助我們.
去看到這些人的需要.
這個挑戰我會跟大家說 你會不會去祈禱呢.
或者問神.
主啊 我的心裡面有哪一部分是硬了的呢.
我的心有哪一部分是硬了的呢.
我有沒有對一些人或者某些事已經麻木了呢.
神給大家的挑戰就是.
今天是禮拜六.
下個禮拜你會不會特意去為一件你平常不會去理會的事去祈禱.
行不行呢.
你平常一定不會去理會.
你特意去為那件事去祈禱.
或者你以前是不會去理會那些人的.
你特意去找那些人去關心一下.
是一個而已.
一件事.

$^{721}$一個而已.
你會不會這樣做呢.
這些是在過程裡面.
我相信神會幫我們將我們本來比較硬的心.
比較冷漠的心變為肉心.
是一個歷程來的.
今天神呼召大家.
下個禮拜不妨去試一下這件事.
說了這麼久.
神回應了今天神的話題.
喂 你下場了嗎.
你不要只是看球賽.
你下場了.
你真的踢球了.
你真的不是聽到.
你行到了.
我講一個見證給大家聽.
我有一個都幾熟的弟兄.
都經常t的.
t就吹水.
有時候講一下信仰的東西.
這位弟兄有一次跟我說.
他說其實他中學時間已經信耶穌了.
因為他生長在一個經常喜歡跟別人比較的家庭.
他說他媽媽小時候經常將他跟鄰居輪射比較.
Ruby 你95分很厲害嗎.
你隔壁的B仔拿了100分.
其實我覺得拿95分已經很好.
還要比別人 不夠別人.
所以他從小到大.
他就是這樣比較.
但他的媽媽不是不好心腸.
不是壞人.
不過有些中國人都很希望.
你這樣不要滿足.
你驕傲 95分.
你看人家100分多好.
其實是這樣的意思.
不過對於這位弟兄就不好的影響.
所以他從小到大就要跟別人比較.

$^{761}$Hugo 幾多分.
80分而已.
我98分.
這樣的環境.
後來他中學的時候信了耶穌.
你記住這幅戒這個設定就這樣帶入教會.
教會有什麼可以突出自己的呢.
你唱歌不及Max那麼好聲.
SV又不及這幾位姐妹那麼好.
那些迎新又不及後面那幾位姐妹那麼好.
他厲害在什麼呢.
不過他厲害在讀書.
他厲害在讀書有什麼competitive advantage在教會.
聖經.
他讀得書嘛.
就精研聖經.
所以他沒多久就出人頭地了.
在教會裡面大查經.
有料到.
那些弟兄姊妹藍的那些又另眼相.
有前途啊這個靚仔.
有沒有玩過什麼聖經常識問答比賽.
那些知識問答比賽.
每次都贏.
接著有人提拔他.
藍就找他做主任學教師.
自然就很型.
在教會裡面出人頭地很型.
如是者.
他就由中學升上大學.
是我師弟.
發覺香港大學有一個大學的基督徒團契.
他就自然我是基督徒.
我入了大學.
當然是入基督徒團契.
有基督徒團契查經的時候.
一直在那裡查經.
阿松就帶著查經.
整班在那裡.
自然這位弟兄就精研聖經.

$^{801}$阿松信了很久.
我就當阿松.
他問信了多久.
阿松就說.
兩年左右.
因為兩年都暴露了帶查經的技巧.
要多厲害.
他看著看著.
唉.
無聊到.
阿松就看不起那個團契.
接著就沒有回教會.
沒有回那個團契.
如是者過了一年.
他就說很記得.
有一年在SWY那裡.
他見到SWYHALL.
即是香港大學有個叫做.
太古樓SWYHALL的宿舍.
有個貼紙貼著.
基督徒團契有早討會.
大家記住.
有些人做過大學生.
香港的大學生要早起都很困難.
七點鐘起床.
早討會.
老實說我當時做大學生沒那麼早起床.
我記得有一次.
要上早lecture.
我醒了.
上了lecture theater.
去到門口聽見教授.
在說話.
突然覺得很睏.
還是走回去hall裡睡覺.
不要告訴別人.
那時候是這樣.
你叫我早很難.
但是竟然有班基督徒.
那麼早起床祈禱.

$^{841}$他就好奇也好.
被他吸引也好.
就走了去.
參加那個早討會.
參加完之後.
其中有一個他們組織的成員就約他.
不如去樓下吃早餐.
好呀.
那就下去.
在那裡吃.
那個成員就問他.
你有沒有回團契.
我這位弟兄就很囂張.
眼睜睜.
嘴瞄瞄.
哦.
你帶茶經的年資就相當淺.
他的知識又未足夠.
很難牧養到我的.
他雖然這樣很謙卑的語氣.
但他的樣子很驕傲.
那位弟兄聽完.
哦.
那就繼續吃早餐.
吃到差不多就跟他說.
啊.
你聖經知識都很豐富.
不過呢.
你知道這些真理.
但是卻沒有走出來.
是不是等於不知道呢.
這位弟兄他覺得.
人生首次.
聖靈在他心裡面工作.
他突然間覺得整個身體是暖.
有一股熱力出來.
不是生氣的那種.
他突然覺得勝靈在他心裡面工作.
令他突然間好像.
神用一些.

$^{881}$physical的effect.
去提醒他.
這個就是真理.
這個就是聖靈跟你說話.
那位成羽的弟兄就問.
那你會回團嗎.
回!回!回!.
自此之後他就回了.
這個經歷告訴我們什麼呢.
我們信耶穌.
我們的信仰真的要走出來.
所以今天如果大家.
聽完之後聖靈有提醒你.
有一些行動.
大家必須要回應神.
不如我請大家一起讀.
今天最後的一段聖經.
《新命紀》十章二十至二十二節.
是我們今天講的聖經最後那三節.
我們一起讀一二三.
你要敬畏耶和華你的神.
事奉他 專靠他.
也要指著他的名起誓.
他是你所欣慰的.
是你的神.
為你做了那大而可畏的事.
是要你親眼所看見的.
你的列祖七十人下埃及.
現在耶和華你的神.
使你如同天上的星那樣多.
我想講的是.
專靠他.
以及指著神的名起誓.
其實摩西最後的總結是.
你要專一信靠我們的神.
並且要指著他的名起誓.
他呼召我們.
對神真的有一個life commitment.
一個生命的衛生.
生命不是說話出來的.

$^{921}$生命不是說出來的.
生命是行出來的.
是活出來的.
以致當你這樣行.
神的話的時候.
你的生命慢慢是向人講故事.
如果大家看那些音樂比賽.
那些評判說.
你唱這首歌是在講一個故事.
他就是在講這件事.
你如何待人.
就等於在講一個生命故事.
神要我們的生命就是這樣.
你如何為人.
你如何對神.
你如何對其他的人.
你怎樣對待.
你用什麼態度對待你的信仰.
就等於你生命在講一個故事給人聽.
慢慢當你的故事講得動聽.
神會使你的故事動聽的時候.
你就會在人生命裡面有權柄.
接著你叫人信耶穌.
那些人會回應信耶穌.
你不信呢.
你試一下這樣走.
我相信神會成就他的心意.
聽道和行道中間還差一個字.
就是決定.
等於可能.
我不知道大家有沒有在youtube看過那些.
健身的那些.
教育的短片.
可能你正在做健身.
我認識一些姐妹.
她看過很多這些健身.
動一動又教你怎樣.
methodology是怎樣.
但如果你拼命看那些短片.
你自己不去gym room.

$^{961}$你沒有舉過鐵.
你是不會fit的.
你不會strong的.
我們的信仰都是一樣.
我們只是聽.
你要愛人.
你要毀身.
你要謙卑.
你就這樣聽.
是沒有辦法改變你的生命.
是沒有辦法讓你的生命可以講故事出來.
你必須要下場踢球.
你必須要下場practice.
你才能做到.
生命的第十章其實是在說什麼呢.
是摩西對著一群.
聽了他說話很久的人.
這群人就快要進入那個應許之地了.
摩西要提醒他要挑戰他的生命.
當你接受挑戰.
當你願意走出來的時候.
你就得著神所給你的應許.
今天其實我們都是應許之地的邊緣.
神的心很想賜給我們豐盛的生命.
和神的親密的關係.
和有永恆價值觀的生命.
但這些神賜給我們的東西.
不是聽就得到的.
是要走出來的.
今天你是否願意走出那一步呢.
有一個生命的轉化.
是讓你從聽到變為行到.
將你本來可能有一點冷的心.
將它變為肉心.
將大家可能吃得很飽的.
習慣了吃花生.
變為真的下場踢球了.
你有沒有聽過廣播的人是多過踢球的.
因為很多人聽.
喜歡聽 喜歡說.

$^{1001}$但是聽和說不是行動.
你必須要下場踢球.
最近有幾個同工問我.
聽說你踢不踢球.
我踢 我幾十年前踢.
他說 你踢不踢.
我踢不要緊.
不過我怕扭傷.
令到你們羞辱就不太好了.
但問題是你要下場踢才行.
你要下場踢才能經歷.
鄧小慧 你會不會.
最後一個挑戰.
你會不會禱告時問神.
我應該怎樣開始跟從你.
我怎樣才開始.
我從哪裡開始.
我怎樣開始跟從你.
走出這一步.
等於我挑戰大家.
神也挑戰大家.
你願不願意真的下場踢球.
我們祈禱.
七位主人.
當我們來到你的先祖寶座面前.
神也給我們這個挑戰.
這個挑戰很簡單.
在我們聽完你的教導的時候.
我們聽完聖靈跟我們說的話的時候.
我們是否真的回應.
我們是否用行動去回應你.
所以我們當中這麼多人.
可能每個人的回應模式都不同.
每個人的回應方式都不同.
因為我們生命光景是不同.
你挑戰我們都是不同的位置.
所以我們當中.
今天聖靈也邀請大家.
如果你願意回應神.
我願意啟動.

$^{1041}$我願意起行.
我願意祈禱.
求神主教我怎樣開始.
我從哪裡開始.
我願意開始.
但我願意禱告.
我願意開始.
我願意回應你.
我邀請弟兄姊妹.
如果願意這樣做.
你把你的手放在你的心上.
我很想為你祈禱.
因為我們這樣走.
我們要決定.
我們要意志.
我們要有一個開放的心.
但是聖靈會給我們能力.
讓我們足夠.
讓我們有豐富的恩典.
因為神今日呼召我們.
祂永遠不會撇下我們.
祂一定會與我們同在.
我們回應祂的時候.
祂是喜悅的.
祂是開心的.
祂會祝福我們.
所以如果你願意這樣做.
請你把你的手放在你的心上.
主要我奉耶穌的名去祝福弟兄姊妹.
當他們很想真心回應你的時候.
這個真心是你很喜歡.
你很開心.
你很喜悅.
所以我奉你的名祝福弟兄姊妹.
打開他們的心眼.
打開他們的耳朵.
讓他們能夠感受到你的心.
主要是因為你答應過我們.
你應許過我們.
當我們願意親近你.

$^{1081}$你就必親近我們.
所以願你的知情就在弟兄姊妹的生命裡面.
所以無論今日他們的出發點在哪裡.
他們的起步點在哪裡.
神呀 你就讓他們有爆炸力.
可以衝出他們的安全區.
讓他們突破這個安全區的時候.
他們要經歷神你是何等的真實.
何等的活著在他們的生命裡面.
主呀 我相信你加能致力給我們.
主呀 我亦相信你必定不會撇下我們.
所以我奉你的名祝福弟兄姊妹.
給他們那種勇氣,能力,勇敢.
讓他們可以回應你.
你差遣他們做些什麼.
你差遣他們關心些什麼人.
你差遣他們去做些什麼事.
他們怎樣回應你.
神呀 你就跟弟兄姊妹講說話.
讓他們的禱告臨到你的施恩寶座面前.
讓他們的心很明亮地見到耶穌基督.
主呀 我多謝你.
我將弟兄姊妹交給你.
求主你恩高他們.
恩典滿滿地充滿他們.
奉耶穌基督得勝人名.
Amen.
\newpage



\section{路加福音 12:13-34}
\label{sec:032351}
\textbf{跳出填倉人生}~~~~李觀南傳道
\newline
\newline
連結: \href{http://www.church.com.hk/acms/content.asp?site=cdc\&op=show\&type=product\&code=032351\&layout=sermon}{\texttt{032351}} ~~~~ 語音日期: 2026-06-21
\newline
\newline
\hyperref[sec:032350]{\small{< < < PREV SERMON < < <}}
~
\hyperref[sec:index_chronic]{\small{[返順時目]}}
~
\hyperref[sec:032352]{\small{> > > NEXT SERMON > > >}}
\newline
\newline
路加福音 12:13-34
\newline
\begin{longtable}{cl}
\hline
\hline
章節 & 經文 (和合本修訂版)\\
\hline
12:13 & \begin{tabularx}{0.7\textwidth}{X} 人群中有一個人對耶穌說：「老師！請你吩咐我的兄弟和我分家產。」 \end{tabularx} \\ \\ \relax
12:14 & \begin{tabularx}{0.7\textwidth}{X} 耶穌對他說：「你這個人！誰立我作你們的判官，或給你們分家產的呢？」 \end{tabularx} \\ \\ \relax
12:15 & \begin{tabularx}{0.7\textwidth}{X} 於是他對他們說：「你們要謹慎自守，躲避一切的貪心，因為人的生命不在於家道豐富。」 \end{tabularx} \\ \\ \relax
12:16 & \begin{tabularx}{0.7\textwidth}{X} 然後他用比喻對他們說：「有一個財主，田地出產豐富。 \end{tabularx} \\ \\ \relax
12:17 & \begin{tabularx}{0.7\textwidth}{X} 他自己心裡想：『我的出產沒有地方儲藏，怎麼辦呢？』 \end{tabularx} \\ \\ \relax
12:18 & \begin{tabularx}{0.7\textwidth}{X} 就說：『我要這麼辦：要把我的倉庫拆了，另蓋更大的，在那裡好儲藏我一切的糧食和財物， \end{tabularx} \\ \\ \relax
12:19 & \begin{tabularx}{0.7\textwidth}{X} 然後要對我自己說：你這個人哪，你有許多財物積存，可供多年享用，只管安安逸逸吃喝快樂吧！』 \end{tabularx} \\ \\ \relax
12:20 & \begin{tabularx}{0.7\textwidth}{X} 神卻對他說：『無知的人哪！今夜就要你的性命，你所預備的要歸誰呢？』 \end{tabularx} \\ \\ \relax
12:21 & \begin{tabularx}{0.7\textwidth}{X} 凡為自己積財，在神面前卻不富足的，也是這樣。」 \end{tabularx} \\ \\ \relax
12:22 & \begin{tabularx}{0.7\textwidth}{X} 耶穌又對門徒說：「所以，我告訴你們，不要為生命憂慮吃甚麼，為身體憂慮穿甚麼。 \end{tabularx} \\ \\ \relax
12:23 & \begin{tabularx}{0.7\textwidth}{X} 因為生命勝於飲食，身體勝於衣裳。 \end{tabularx} \\ \\ \relax
12:24 & \begin{tabularx}{0.7\textwidth}{X} 你們想一想烏鴉：牠們既不種也不收，既沒有倉又沒有庫，神尚且養活牠們。你們比飛鳥要貴重得多呢！ \end{tabularx} \\ \\ \relax
12:25 & \begin{tabularx}{0.7\textwidth}{X} 你們哪一個能藉著憂慮使壽數多加一刻呢？ \end{tabularx} \\ \\ \relax
12:26 & \begin{tabularx}{0.7\textwidth}{X} 這最小的事你們尚且不能做，何必憂慮其餘的事呢？ \end{tabularx} \\ \\ \relax
12:27 & \begin{tabularx}{0.7\textwidth}{X} 你們想一想百合花是怎麼長起來的：它也不勞動，也不紡線。然而我告訴你們，就是所羅門極榮華的時候，他所穿戴的還不如這些花的一朵呢！ \end{tabularx} \\ \\ \relax
12:28 & \begin{tabularx}{0.7\textwidth}{X} 你們這小信的人哪！野地裡的草今天還在，明天就丟在爐裡，神還給它這樣的妝飾，何況你們呢？ \end{tabularx} \\ \\ \relax
12:29 & \begin{tabularx}{0.7\textwidth}{X} 你們不要求吃甚麼，喝甚麼，也不要掛慮。 \end{tabularx} \\ \\ \relax
12:30 & \begin{tabularx}{0.7\textwidth}{X} 這都是世上的外邦人所求的；你們需要這些東西，你們的父都知道。 \end{tabularx} \\ \\ \relax
12:31 & \begin{tabularx}{0.7\textwidth}{X} 你們只要求他的國，這些東西就必加給你們了。 \end{tabularx} \\ \\ \relax
12:32 & \begin{tabularx}{0.7\textwidth}{X} 你們這小群，不要懼怕，因為你們的父樂意把國賜給你們。 \end{tabularx} \\ \\ \relax
12:33 & \begin{tabularx}{0.7\textwidth}{X} 你們要變賣財產賙濟人，為自己預備永不壞的錢囊和用不盡的財寶在天上，就是賊不能近，蟲不能蛀的地方。 \end{tabularx} \\ \\ \relax
12:34 & \begin{tabularx}{0.7\textwidth}{X} 因為你們的財寶在哪裡，你們的心也在哪裡。」 \end{tabularx} \\ \\
[1ex]
\hline
\hline
\end{longtable}
$^{1}$經文有些長,多謝大家讀了我這麼久.
正如剛才所說,這段經文其實就是耶穌在地上傳道時.
那麼多個比喻中的其中一個.
如果大家有看聖經的話,大家說去都知道.
耶穌很喜歡將一些天國的道理.
用一些日常生活的例子,用比喻的方式去表達出來.
祂的目的是為了讓聽到的人可以更加容易明白.
更加容易理解到.
而每當耶穌去講比喻的時候.
祂當然不會說失敬無神,不如先講個比喻.
正常人都不會這樣做,是嗎?.
通常都是有起因的.
起因就是可能祂在群眾裡面有人向祂提出疑問.
又或者可能在現場裡面.
有些事情發生了,需要耶穌去回應.
然後耶穌就會用比喻的方式去回應這些問題.
又或者這些事情.
而且很多時候,在比喻之後會有一個更加深入的解釋.
不過,這部分的講的對象通常都只是限於祂的門徒.
而我們今天剛剛一起讀的經文.
我同樣會按著起因,比喻內容和今日的講解這三部分.
作為我三個要點去跟大家分享.
我們先看一下這個比喻的起因.
故事一開始就說.
當時有一個人來到耶穌面前.
他要求耶穌為他們兄弟二人去分家業.
其實在經文裡面沒有交代過.
究竟這個人和他的兄長之間.
他們究竟爭拗的是什麼.
是不是因為爸爸過世之前.
他未立好遺囑.
抑或究竟是不是他的兄長沒有按照他爸爸的遺囑.
按照遺囑所說的.
將他應得的那份分給他.
是不是這些呢?.
我們不知道.
但是聖經裡面提到的是.
他來到耶穌的面前.
他要求耶穌為他主持公道.
可能大家都會想.

$^{41}$為什麼要分家產這件事.
都要來找耶穌呢?.
如果我們看文素和生命的一些經文.
你會發現其實處理家業這件事.
正正是摩西律法裡面其中一個很大的課題.
而在耶穌時代.
詮釋摩西律法也正正是那些猶太拉比要做的一件事.
是他們的責任.
所以我們在經文裡面看到的就是.
這個人一來到耶穌面前.
他正在稱呼耶穌為夫子.
也就是老師的意思.
他正正期望的是.
耶穌可以像那些一貫的拉比一樣.
會按照摩西律法去為他和兄長.
去解決他們之間的紛爭.
不過事原原為的是.
耶穌沒有順著他所想的去做.
反而是帶有一點點斥責意味的去反問這個人.
他說究竟是誰立我.
去為你們兩個人去搞家業這件事呢?.
不知道大家聽到耶穌.
當一個人來找耶穌.
耶穌這樣回答他的時候.
大家會覺得怎樣?.
可能你也會覺得.
耶穌其實是不是也挺嚴厲的?.
他那個人其實也是在按照猶太人一般的做法.
去找拉比而已.
也就像剛剛所說.
無論是詮釋摩西律法.
還是處理家業.
本身就關猶太拉比的事.
所以當這個人看到耶穌周圍跟別人說.
天父的心意是怎樣的時候.
他會找耶穌幫忙.
其實看起來好像不是很錯.
雖則我們今天的信徒.
我們會知道的是.
這個人將耶穌誤當成為一個普通的猶太拉比.

$^{81}$這件事本身很錯.
但是細心想想.
當時那個人他跟我們不同.
他不知道原來耶穌是神的兒子.
所以他會找耶穌這個拉比去幫忙的時候.
其實看起來也挺合理.
不過我想無論怎樣也好.
我覺得在這個起因裡面.
最值得我們留意的.
不是這個人為什麼要找耶穌去分家產.
而是這個人當他來到耶穌面前的時候.
他說的那句話.
他來到耶穌面前就說.
夫子,請你吩咐我的兄長和我分家業.
如果我嘗試用這句話套進我們今天教會的setting.
話說有一天有個弟兄走去謝牧師面前.
他就跟牧師說.
牧師啊牧師,我拜託你幫我和某位弟兄去拆小組吧.
那是怎樣呢?.
大家有些笑聲.
其實就是說這個人已經帶著一個問題的解決方案.
他來到耶穌面前.
他來到謝牧師面前.
他不是要謝牧師也好,耶穌也好.
去投訴那個弟兄,投訴他哥哥有什麼做得不好.
有什麼不公義的地方.
他不是要耶穌來幫他評評理.
他也不是來到耶穌面前尋求耶穌的意見.
而是直接吩咐,直接叫耶穌去吩咐他的兄長.
去幫他搞分家產這件事.
如果難聽點這樣說的話.
我會說這個人其實是直接吩咐耶穌去做事.
可能大家聽到的時候.
你的反應都會想.
哇,這個人其實都很大膽,吩咐耶穌.
但是當我認真去想到這個時候.
我會問一個問題.
我們自己今天跟這個人有很大分別嗎?.
我們平時的祈禱何嘗不是吩咐耶穌做事呢?.
我們求神給我們很多時候見公順利.

$^{121}$求神給我們遇見生命中命中注定的那個.
我們會求神給我們的兒女聽話.
我們連巴士都求神.
我們一走到站上,拜託馬上有.
因為天氣真的很熱.
但是聖經明明在說的是.
它在教我們要先求上帝的國和上帝的義.
但是為什麼我們的禱告裡面.
大部分都是在求一些我們日常的飲食.
求一些我們盼望的事能夠心上事成呢?.
我們甚至會更壞的想的是.
當上帝沒有按著我們所祈求的事去成就的時候.
我們甚至會埋怨上帝.
我們會怪祂,為什麼你不聽我們祈禱.
又或者你沒有那麼大膽.
你說你不敢埋怨上帝.
但是你心裡面卻慢慢地覺得祈禱都沒有用.
我們就把祈禱都放下.
不再去找神,不再去祈禱.
所以我會說其實我們今天.
都不比要求耶穌分家產的那個人特別為好.
我們很多時候同樣都把耶穌當成了解決問題的一個工具.
我們的分別可能是那個人親口跟耶穌說.
而我們把我們吩咐耶穌的事情.
放在我們的祈禱裡面.
放在我們的心裡面.
所以今天這段經文這個氣音.
我想第一個對我們的提醒就是.
我們要好好反思我們自己.
我們究竟今天我們說我們很想親近耶穌.
其實我們背後是不是不過是為了一些自己個人的好處.
為了那個所謂有用.
我們才祈禱,才去親近耶穌呢?.
我們會不會其實都無形間跟故事那個人一樣.
把耶穌當成了一個普通的猶太拉比.
把耶穌當成了一個解決問題的工具.
所以弟兄姊妹,我盼望我們不要成為了.
聖經我們很多都熟悉的,金句經常說的.
不要只是看到別人眼中的刺.
看不到自己眼中的亮目.

$^{161}$盼望我們自己不要成為一個這樣的人.
常常去反思自己.
常常去想一下你自己今天跟耶穌的關係.
是不是純粹來到了為利益的緣故呢?.
這個是第一點.
我們繼續看故事的發展.
故事中其實沒有交代過究竟這個人.
當他聽完耶穌的反問之後.
他有什麼反應.
在聖經裡面只是記載了.
耶穌當時沒有直接回應這個人的要求.
他沒有選擇去處理那個人想解決的事情.
他反而跟手就將他對話的對象.
從這個人身上轉去其他周圍的人.
耶穌乘著這件事情的發生.
他乘著那個勢就用比喻去教其他的人.
在比喻裡面,耶穌說有一個財主.
因為他田產豐盛.
他就想了如何可以好好收藏自己田產的出處呢?.
於是他就想到一條好計.
就是打算拆掉他本身有的倉房.
繼而再建一個再大一點的.
以致他能夠有足夠的空間.
可以將他所擁有一切的糧食一切的財物.
都可以好好地收好.
叫他將來的人生可以有豐足的供應.
可以無憂無慮安逸地過他的生活.
如果我們單看這一部分的話.
其實我相信我們普遍的我們.
都不會覺得有什麼大的問題.
就好像我們今天.
假如你上司你老闆突然間告訴你.
我加你薪水.
甚或因為公司生意很好.
我今天突然間想派一個額外的額外獎勵給你.
人之常情你都會想.
那我怎樣用好呢?.
我會想究竟有沒有一些我喜歡的東西.
不如買一兩件給自己哄自己開心.
又或者可能你疼愛家裡的.

$^{201}$跟家裡人吃一餐是好的.
甚或你可能好像故事裡面的財主一樣.
這一刻還沒有用.
但是我先收好它.
保存在銀行.
我不知道大家現在還放不放在銀行.
放在銀行.
等將來有用的時候再拿出來用.
可能我們都是這樣.
所以我會說其實.
這個財主當下因為他的田產豐收.
他需要建一個更大的倉庫.
去為自己做好準備.
去為自己的將來打算.
其實有什麼錯呢?.
為什麼耶穌要在比喻裡面說這個人.
這個會為自己將來打算的人是無知呢?.
如果我們細心看經文的話.
其實不難發現.
真正的問題其實不是這個財主會為將來.
為自己的將來打算.
而是這個財主究竟他為將來打算.
他的背後的動機是什麼呢?.
在第十九節的經文.
當這個財主想好了這個計劃之後.
他就跟自己的靈魂說.
他說靈魂啊 你有許多財物積存.
可以作多年的費用.
只管安安逸逸吃喝快樂吧.
這節經文其實指出了.
這位財主他要建一個更大的倉庫.
來收好自己的糧食和財物.
他為的就是為了自己將來可以安安逸逸吃喝玩樂.
這才是耶穌形容他無知的原因.
可能你又會想.
其實為了吃喝玩樂有什麼問題呢?.
我們很多人上班.
上班那麼辛苦為了什麼?.
到最終其實何嘗不是為了要吃好一點.
住好一點.

$^{241}$旅行去久一點.
去得遠一點.
都是為了這些而已.
如果大家都是這樣想的話.
我們需要先理解一下當時猶太社會的一些文化.
在利未記的十九章九至十字.
和申命記的二十四章十九至二十一字.
我不讀了 大家可以看.
這兩節經文說的是什麼呢?.
說的是原來當時的社會.
即是這些財主即是有田地的人.
在摩西的律法裡面其實很清楚地說明.
財主是有相應的社會責任是他們要負責的.
他們是有責任要照顧窮人.
寄居的和那些孤兒寡婦.
如果大家對路得記有多少印象的話.
裡面有位財主叫波亞斯.
他容讓路德在他自己的田裡面去灑墨水.
他正正就是在履行猶太社會裡面.
作為一個財主的一個社會責任.
故事的主角這個財主有沒有呢?.
經文沒有寫.
沒有說他有沒有履行這個責任去照顧窮人或者孤兒寡婦.
但是在他自言自語的說話裡面.
我們看到有一個字常常都見的.
就是「我」字.
在三節的經文裡面出現了五次.
其實反映這個人一直最緊張的是他自己.
他將上帝賜給他的田地,財產.
一切的豐盛.
他視為滿足自己慾望的一件事.
上帝稱呼他是一個無知的人.
而無知的人在舊約聖經裡面.
常常都會用來形容一些不敬畏實的人.
比喻裡面的這位財主.
他不明白到財富是上帝賜給人的一個恩典.
他也沒有按照上帝的吩咐.
很正確地用這些財富去周濟窮人.
去幫助那些有需要的人.
反而將自己的需要,自己的安逸.

$^{281}$自己將來的生活的保障.
看得大於上帝的教導.
我會說這個才是他真正的無知.
是他不敬畏神的地方.
亦都是經文說.
為何他在神面前也不會負足的一個原因.
因為他將焦點放錯了.
他將為了自己.
將滿足自己看為最重要.
而不將神放在他生命的首位.
丁姐妹我們呢?.
我們有沒有將我們的生命有時也放錯焦點?.
剛才經文第十五節提醒我們要謹慎自受.
要免去一切的貪心.
人會有貪心的念頭出現.
其實正正就是源自於我們很想滿足自己.
在今天的社會裡面.
滿足自己.
不只是說錢.
可以有很多其他的事情.
可能是說我們對一些權力的追求.
我們要在一個群體裡面.
我們要有話事權.
我們要在對人裡面.
或者對那個社會身份地位.
我們同樣都會有很多的追求.
我們很盼望別人去認可自己.
很想成功 很想有成就.
我們又甚至滿足自己.
可以說我們跟某些人的關係.
跟一些不同的人的關係.
可能我們會覺得.
我們認識到某些達官貴人.
認識到某些出名的人.
認識到某些出名的牧師.
我們就會覺得自己厲害一點.
甚至最簡單的是.
我們在網上.
大家自己的社交平台裡面.
你出一個帖子.

$^{321}$有沒有人Like.
你的朋友會不會回應你.
都可以成為我們滿足自己的一些小事情.
在教會裡面又怎樣呢?.
我會說教會裡面都不例外.
特別是一班常常參與服事的弟兄姊妹.
我們會不會因為經常都在人前.
經常在別人的目光底下.
甚至有時候有些人可能很好的弟兄姊妹.
甚至讚你.
你會不會因為這些事情.
你就忘記了你最初參與服事的那個心智.
忘記了你參與服事.
最初只不過是為了你要回應上帝對你的愛.
的一個很小的選擇.
會不會我們慢慢地變得追求了這些人的目光.
這些人的讚賞.
就忘記了我們的上帝是怎樣看待我們.
我們會不會在意那些事情.
多於在意實在怎樣看待.
所以我會說其實教會一樣.
我們都很容易就擺錯Focus.
但是大家有沒有留意到的是.
在經文裡面當耶穌聽到要求分家產的弟弟.
聽到這個要求的時候.
他其實是沒有去處理過這件事.
耶穌選擇的是他不處理人的私務.
他反而選擇處理的是人的內心.
他選擇處理的是這個人.
他心底裡面為了滿足自己的貪念.
去處理他這一份無知.
正正這件事告訴大家的是.
耶穌其實就是不在乎那些外在的東西.
他真正在乎的是我們每個人的心.
而且這個情況是去到故事裡面的分家產的弟弟.
還是比喻裡面的財主.
他裡面都告訴我們.
就算你要求分家產.
你要求建立一個更大的倉庫給自己.
你在做的這些事情.

$^{361}$在人的角度裡面看為多正確也好.
有多麼合乎常理也好.
只要我們將Focus擺錯.
我們同樣都會成為經文裡面所說的無知的人.
所以第二點很想大家記得的是.
上帝他不只是看人怎樣做事.
怎樣處理好他們的事情.
處理好他們的人生.
他是在看我們每一個人的內心.
究竟我們內心想的東西是否純正.
我們所做的事情.
背後的動機是否正確.
我們仍然在做每件事裡面.
有沒有專注為大.
我們有沒有將我們的Focus放在上帝身上.
這是第二點.
去到最後一點.
亦都是去到耶穌和門徒講解的部分.
耶穌進一步去解釋.
人之所以要為自己去積財.
原因很多時是人不斷地為自己的生命憂慮.
吃什麼 為我們的身體穿什麼.
的確這兩個問題對我們現代人來說.
同樣都是經常困擾我們.
特別是我們早上起床打開衣櫃的時候.
和我們快要吃午餐要想吃什麼的時候.
我們都不斷每天花時間去處理這兩個問題.
但是耶穌在這裡告訴我們.
我們不用憂慮.
因為我們所掛心的飲食.
我們經常想著的穿什麼衣服.
這些外在的東西.
其實是不能夠令我們的壽數多加一克.
耶穌還分別舉了.
烏鴉,百合花和野對壘的草為例子.
祂帶出的是連舊約被視為不潔淨的烏鴉.
和好像百合花和野草這些微小的植物.
上帝也會養活.
上帝也會顧念.
更何況我們是去搜索創造.

$^{401}$去眼看為貴重的我們這群人類.
所以耶穌在跟門徒說的是.
吃什麼喝什麼這些問題.
其實根本就不值得我們花時間去想.
因為這些肉身上的需要.
只要我們把焦點放在上帝身上.
去追求祂的國.
去按祂的教導去過我們的生活的時候.
上帝是會必然供應我們.
換言之.
耶穌所說的重點就是.
真正能夠令人得生命的.
不是我們不斷地追求一些外在的事物.
也不是我們不斷地追求我們到退休的日子.
我們有一個安逸吃喝的生活.
而是我們要保持一個不斷追求神.
依靠神的心.
是敬畏上帝.
凸起個個的身外之物.
這個才是我們能夠真正得著生命的方法.
這個才是我們真正能夠在神面前.
得到真正的富足.
所以弟兄姊妹們.
不知道大家今天來的時候有沒有帶著憂慮.
你會不會都經常想很多東西.
很掛心很多東西.
你會不會一樣又憂慮著吃什麼喝什麼.
憂慮著你的錢怎麼辦.
憂慮著你那層樓升和跌.
你家產因而有改變.
又會不會憂慮著你的事業.
憂慮著可能是你的家.
可能家中有人有患病的需要.
甚至可能是我們剛好有很多父母在這裡.
憂慮著你的子女的成長如何.
是的 我們很多憂慮.
當我在預備這篇信息的時候.
由於我知道我要講貪婪這個主題.
所以我也上網去特意找一下.
有沒有一些官方的解釋.

$^{441}$究竟貪婪這個字眼是怎麼解釋的呢.
當然找官方就沒有了 哪有官方的.
不過就找到其中一個解釋是這樣說的.
就是說貪婪是指一個人永遠不滿足.
貪圖可能屬於其他人的事.
就是他們的東西.
當我看完這個解釋之後.
我再看回經文裡面的那個弟弟和財主.
我會想其實這兩個人.
不過都是為了拿回自己的家產.
為了想安置好自己有的財富.
看來他們兩個都只是想做好處理好自己的事.
他們其實跟貪婪這個主題有什麼關係呢.
為什麼在這個故事裡面.
會說他們是無知的時候.
叫他們有貪心這個字出現過呢.
我會說是的.
如果我們這樣看的話.
他們其實表面上真的沒有貪過其他人的東西.
因為就算弟弟叫哥哥分回家產也好.
他也沒有說我拿了哥哥的那份.
他不是貪別人的那份.
所以其實他們貪的是什麼呢.
他們真正貪的是.
他們對物質所帶給他們的那份安全感.
他們想有錢.
想保管好自己的錢財.
是因為錢財的確是可以叫他們的生活有所保障.
可以叫到他們的生命能夠.
對他們自己將來的生命有多一份的掌握.
但是這一份的對安全感的追求.
這一份我們很想生命可以自己能夠掌握多一點.
對這件事的渴望.
如果我們用另一個角度來看.
其實在變相說的是.
他們在減少他們自己對神的一份依靠.
我們也是的.
我們常常都覺得一些事情我們可以自己處理.
我們很盼望我們人生能夠掌握到.
明天做些什麼 後天做些什麼.

$^{481}$我們計劃了所有東西.
但是我們會不會因為這一切的事情.
慢慢地就變得所有東西都要我控制.
所有東西都要在我的控制下.
而忽略了原來我們已經對神的依靠.
慢慢地覺得可有可無.
我剛剛都有提到.
耶穌要處理的真正的問題是人的內心.
是根源的問題.
祂在重申真正的保障是從神而來的.
真正的生命是在乎我們有沒有敬畏上帝.
不是說做了這兩件事.
就代表我們可以不用理會我們仍然掛心的事.
那些事情可能我們仍然需要處理.
但是我想最起碼.
當我們去到校正我們的思想的時候.
我們知道原來神是顧念我們.
知道神是那個供應者.
祂會供應我們 祂會叫我們不會缺乏的時候.
我們不會將那些我們掛心的問題.
我們擔心的事情看得過大.
我們不會讓那一切的憂慮困住我們的生命.
不會令我們一生不斷地只會追求物質的需要.
不過你聽到這裡你當然會覺得.
說就容易 做就很難.
因為我們自小成長.
其實一直以來世界也好 家人也好.
在教我們的是什麼呢?.
我們要靠自己.
我們要相信自己.
我們要做好自己.
要顧好自己.
這個是我們自小的教導.
我和大家分享一個小故事.
這個小故事是話說有一個小朋友.
有一天和爸爸一起外出.
要出遠門去參加一個朋友的生日派對.
由於路途遙遠.
爸爸一出門口的時候已經準備好兩顆糖.
因為怕他中途覺得太遠 又悶 準備好.

$^{521}$給了一個小朋友.
那個小朋友一直拿著.
在整個車程裡他都拿著兩顆糖.
沒事的.
也很安全地到達了他們要參加派對的地方.
但是去到這個派對的時候.
他就發現.
邀請他們來派對的生日小朋友和他們的家人.
很好 因為要招呼人的緣故.
預備了很多禮物.
預備了很多食物在那個舉辦派對的地方.
但是這個小朋友.
由於他兩隻手拿著兩顆糖.
但是由於年紀相若.
可能只有兩三歲.
其實他不懂得拿著兩顆糖的同時.
其實他的手臂是可以繼續拿東西的.
他還沒有了解到這件事.
所以呢 爸爸也知道他這樣.
就和他說.
不如你放下兩顆糖吧.
你兩顆糖放下之後.
你就有手去拿你想要的禮物.
拿你想要的東西吃.
但是礙於這個小朋友.
一向都很喜歡吃這兩顆糖.
所以他一直不想放下.
結果他就沒有手去多餘的去拿他想要的禮物和食物.
這個小故事.
其實想告訴大家.
我們就是一個小朋友.
當我們常常以為那兩顆糖很重要的時候.
我們看不到原來上帝是在告訴我們.
祂預備了更多更大的事情.
更多更大的恩典給我們每一個.
但是我們永遠就是很想捉緊那兩顆糖.
因為我們覺得這兩顆糖是我們生命的保障.
是我們的依靠.
而當我們說那個小朋友.
他如果當天.

$^{561}$就是那個時候.
他聽到他爸爸說.
你先放下兩顆糖吧.
的時候.
他可能已經得到更多.
但是為什麼他不會這樣選擇呢?.
是因為他不夠相信他爸爸.
他不相信他爸爸會幫他保管好那兩顆糖.
他不相信他爸爸所說的那件事是有效的.
所以弟兄姊妹.
剛剛都說我們就是一個小朋友.
你信不信你天上的爸爸所說的話呢?.
你願不願意放下你那兩顆糖.
繼而你可以張開雙手.
去迎接那份更大的祝福呢?.
盼望我們同樣學習有一個小朋友的心.
去看我們的天賦.
去相信他的教導.
好 終結了.
我今天將我的講題改了個名.
叫做「跳出填倉的人生」.
如果大家有買股票.
都心知應該填倉都不是一件好事.
我改這個講題的題目.
原因當然其中一個.
就是為了要配合我那個比喻裡面的財主.
他想建一個更大的倉.
這個故事.
而其次我想說的是.
其實我們每個人生命裡面都一樣.
我們每個人都有一個倉.
這個倉其實是我們的內心.
也都是代表著我們生命裡面.
很需要的一份安全感.
但是今天究竟.
我們在用一些什麼去填這個倉呢?.
是不是我們的財富.
是不是我們有一份穩定的工作.
穩定的收入.
還是我們一早已經計劃好的那個退休大計?.

$^{601}$可能我們人生裡面真的有很多東西預備了.
但是我盼望今天這篇信息.
讓我們能夠學懂的是.
我們要常常都有一個自省的心.
我們要好好地去反思自己.
要常常反思自己所做的每一件事情.
究竟我們的動機.
會不會有些時候很不小心.
你不是故意的.
很不小心的.
就把焦點放錯了.
會不會有些時候.
我們為了解決我們那個所謂的憂慮.
我們連神都放下了.
繼而選擇的是我們選擇相信自己.
多過選擇相信上帝.
很想鼓勵大家.
今天嘗試再一次將你這個倉.
將你的心.
將你的人生.
再次交回給上帝.
再一次去信靠祂.
再一次去依靠祂的時候.
你能夠將你的憂慮.
完完全全地給祂.
讓你的生命.
不需要每一天都再次.
為了填滿這份安全感.
因為我們所相信的天父.
所相信的耶穌.
祂是我們那個生命最大的滿足.
祂是我們每一個人的供應者.
祂是我們每一個人的保護者.
最後想分享一張圖片給大家.
這張圖片在我預備這篇讀的時候.
在網上的社交平台看到.
而它所說的.
正正是我整篇讀的重點.
所以其實前面那半個小時.
你們聽不到什麼.

$^{641}$不要緊的.
你這一刻拿出你的手機.
拍下這張圖片.
這張圖片說什麼呢?.
我們很喜歡的就是.
靠自己掌握自己的人生.
但是我們的上帝.
在神裡面項為富足.
項為成功.
是我們要學懂怎樣去依靠上帝.
和上帝去走.
去過好我們的生命.
\newpage



\section{路加福音 17:20-37}
\label{sec:032352}
\textbf{天國的來臨}~~~~俞大鳴牧師
\newline
\newline
連結: \href{http://www.church.com.hk/acms/content.asp?site=cdc\&op=show\&type=product\&code=032352\&layout=sermon}{\texttt{032352}} ~~~~ 語音日期: 2026-06-07
\newline
\newline
\hyperref[sec:032351]{\small{< < < PREV SERMON < < <}}
~
\hyperref[sec:index_chronic]{\small{[返順時目]}}
~
\hyperref[sec:032353]{\small{> > > NEXT SERMON > > >}}
\newline
\newline
路加福音 17:20-37
\newline
\begin{longtable}{cl}
\hline
\hline
章節 & 經文 (和合本修訂版)\\
\hline
17:20 & \begin{tabularx}{0.7\textwidth}{X} 法利賽人問：「神的國幾時來到？」耶穌回答：「神的國來到，不是眼睛看得見的。 \end{tabularx} \\ \\ \relax
17:21 & \begin{tabularx}{0.7\textwidth}{X} 人也不能說：『看哪，在這裡！』或說：『在那裡！』因為神的國就在你們心裡。」 \end{tabularx} \\ \\ \relax
17:22 & \begin{tabularx}{0.7\textwidth}{X} 他又對門徒說：「那些日子將到，你們渴望能看見人子的一個日子，卻看不見。 \end{tabularx} \\ \\ \relax
17:23 & \begin{tabularx}{0.7\textwidth}{X} 有人要對你們說：『看哪，在那裡！』或說：『看哪，在這裡！』你們不要出去，也不要追隨他們。 \end{tabularx} \\ \\ \relax
17:24 & \begin{tabularx}{0.7\textwidth}{X} 好像閃電從天這邊一閃直照到天那邊，人子在他的日子也要這樣。 \end{tabularx} \\ \\ \relax
17:25 & \begin{tabularx}{0.7\textwidth}{X} 可是他必須先受許多苦，又被這世代所棄絕。 \end{tabularx} \\ \\ \relax
17:26 & \begin{tabularx}{0.7\textwidth}{X} 挪亞的日子怎樣，人子的日子也要怎樣。 \end{tabularx} \\ \\ \relax
17:27 & \begin{tabularx}{0.7\textwidth}{X} 那時，人又吃又喝，又娶又嫁，直到挪亞進方舟的那日，洪水就來，把他們全都滅了。 \end{tabularx} \\ \\ \relax
17:28 & \begin{tabularx}{0.7\textwidth}{X} 同樣，就像在羅得的日子，人又吃又喝，又買又賣，又耕種又建造， \end{tabularx} \\ \\ \relax
17:29 & \begin{tabularx}{0.7\textwidth}{X} 到羅得離開所多瑪的那日，有火與硫磺從天上降下來，把他們全都滅了。 \end{tabularx} \\ \\ \relax
17:30 & \begin{tabularx}{0.7\textwidth}{X} 人子顯現的日子也要這樣。 \end{tabularx} \\ \\ \relax
17:31 & \begin{tabularx}{0.7\textwidth}{X} 在那日，人在屋頂上，東西在屋裡，不要下來拿；人在田裡，也不要回家。 \end{tabularx} \\ \\ \relax
17:32 & \begin{tabularx}{0.7\textwidth}{X} 你們想想羅得的妻子吧！ \end{tabularx} \\ \\ \relax
17:33 & \begin{tabularx}{0.7\textwidth}{X} 凡想保全性命的，要喪失性命；凡喪失性命的，要保存性命。 \end{tabularx} \\ \\ \relax
17:34 & \begin{tabularx}{0.7\textwidth}{X} 我告訴你們，在那一夜，兩個人在一張床上，一個被接去，一個被撇下。 \end{tabularx} \\ \\ \relax
17:35 & \begin{tabularx}{0.7\textwidth}{X} 兩個女人一同推磨，一個被接去，一個被撇下。」 \end{tabularx} \\ \\ \relax
17:36 & \begin{tabularx}{0.7\textwidth}{X} 兩個人在田裡，一個被接去，一個被撇下。 \end{tabularx} \\ \\ \relax
17:37 & \begin{tabularx}{0.7\textwidth}{X} 門徒回答他說：「主啊，在哪裡呢？」耶穌對他們說：「屍首在哪裡，鷹也會聚在哪裡。」 \end{tabularx} \\ \\
[1ex]
\hline
\hline
\end{longtable}
$^{1}$我不知大家在世上最關心的事是甚麼.
我相信今天大部份弟兄姊妹.
都不會為柴米油鹽擔憂.
我們最關心的是甚麼呢?.
會否是自己的工作,職業.
會否關心何時買屋買屋好呢?.
或者何時退休好呢?.
下輩子換架車呢?.
換架甚麼車好些呢?.
電車還是混合車好呢?.
有些年輕兒女的家長.
可能要考慮搬到哪裡住.
要考慮兒女的教育問題.
又或者已經退休的.
可能考慮下次去旅行去哪裡好呢?.
我有位已經退休的朋友.
我不知道他自己是否喜歡去旅行.
但他太太是很喜歡去旅行的.
時時約一班志同道合的朋友到處去.
環遊世界.
去年去了Greenland.
今年剛剛去完印度回來.
這些事可能是我們現代人關心的事.
但在主耶穌的時代.
猶太人特別關心的事.
可能是天國的降臨.
我們怎知道呢?.
在《馬可福音》一章十四到十五節.
記載了主耶穌一出來傳道的時候.
祂是這樣說的.
耶穌來到加利利村講「上帝的福音」.
甚麼是「上帝的福音」呢?.
「日期滿了,上帝的國近了」.
「你們要悔改順福音」.
很多時候我讀到這節經文的時候.
我就在想.
「既然我效法主耶穌」.
「我走到人多的地方,去到市中心」.
「站在那裡,帶著過音機大聲說」.
「天國近了,你們要悔改順福音」.

$^{41}$大家認為這樣傳福音會不會有果效呢?.
會不會吸引到很多人來聽福音呢?.
我相信不會有果效.
還可能惹人討厭.
在耶穌的時代.
這個方法似乎很有效果.
耶穌就是這樣傳福音的.
招收了很多門徒跟隨祂.
在《馬太福音》第五章.
記載了有一大群人.
是跟隨著耶穌要聽祂的講道.
在我們很熟悉的經文.
八福裡面.
耶穌一開始提到什麼呢?.
祂一開始就提到.
「什麼人可以進入天國?」.
可想而知.
當時的人對天國的來臨.
是非常有興趣的.
另一段我們很熟悉的主禱文.
第三句話.
就提到「願你的國降臨」.
因此我們知道.
不但當代猶太人關心天國.
主耶穌也同樣關心天國.
並且教導門徒也要關心天國.
並且祈禱求上帝的國降臨.
主耶穌來世的目的.
其實就是要重建一條路.
一條可以回家的路.
回到天父的家.
所以主禱文的第一句話就是.
「我們在天上的父」.
耶穌用家的觀念.
來比喻上帝和人的關係.
但對猶太人來說.
他們對上帝和人的關係的觀念.
就不是家而是國.
上帝是他們的君王.
掌管一切.

$^{81}$猶太人就是上帝的特別選民.
要在地上建立一個國家.
來宣揚上帝.
這個觀念在上帝呼召亞伯拉罕的時候.
就已經提出了.
記載在《創世記》十二章一至三節.
「我要使你多子多孫」.
他們要怎樣?.
「形成大國」.
我要賜福給你.
我要藉著你賜福給萬民.
當摩西帶領以色列人出埃及之後.
上帝的目的.
也是要帶領他們進入迦南.
形成一個大國.
因此當人民在曠野拜金牛的時候.
上帝對摩西說什麼?.
「我知道自己人民多麼頑劣」.
「現在你不要攔阻我」.
「我要向他們發烈怒滅絕他們」.
然後做什麼?.
「我要使你和你的子孫成為大國」.
上帝要用以色列人建立大國的計劃.
要經過差不多四百年的時間.
在大衛的帶領之下.
十二支派才能夠統一.
才算得上是建立了一個國家.
有國王.
有整個迦南地作為他們的國土.
也有國民.
不過以色列國最強盛的時候.
是在大衛的兒子.
所羅門作王的時代.
但是非常可惜.
以色列人的確是頑劣的.
在建國四百多年之後.
由於他們背叛上帝.
敬拜偶像賈臣.
上帝就使他們的國家.
被其他強國毀滅了.

$^{121}$但是上帝也藉著那些先知.
去安慰那些愚民.
正愚的愚.
愚民或者說愚種.
講俗一點就是死正種.
國家很多人都打仗死了.
有死正種.
那些沒有死的以色列人.
告訴他們.
即使他們是失敗.
以色列人失敗.
但是上帝沒有失敗.
上帝的計劃裡面.
是有下一步棋的.
就是要派一位彌賽.
救世主.
來去拯救他們.
並且在《先知書》裡面.
有很多關於彌賽.
各道的榮耀預言.
這就是為什麼.
當施洗約翰和耶穌.
出來傳道的時候都說.
「天國近了」.
「天國」也可以翻譯成.
「上帝的國」.
同義詞在這裡.
「天國」這個字眼很有趣.
對日本人來說.
他們認為是人死後去的地方.
但是對猶太人來說.
是以色列重振聲威.
統管天下萬國的時刻.
我不知道大家對天國有什麼觀念.
但是如果從國家的角度來說.
一定會有國土,國民,國軍和國法.
從聖經的記載.
天國似乎更著重的是國法.
和神與人的關係.
當屬神的人進行上帝的法典的時候.

$^{161}$那就是天國的實現.
因此上帝和以色列人在抗夜立約的時候.
除了十誡.
還給了他們很多其他的條例.
對猶太人來說.
摩西在《申命記錄》.
將有二十五節.
律法的功用.
這裡特別提到是律法的功用.
「我們約照耶和華我們上帝所吩咐的.
在他面前謹守信行的一切誡明.
這就是我們的意料」.
我用大白話說一下.
其實這句話的意思就是說.
我們如果謹守信行上帝所立的法例.
我們就是天國的良好公民.
說了聖經對天國的一些背景之後.
讓我們一起思想今天的經文.
今天的經文是談論到天國何時到.
和從何時開始.
在今天的經文第二十節.
我們看到是法利賽人來見耶穌.
問祂有關上帝的國.
何時到來.
這是一個時間的問題.
英文就是「When will it come?」.
法利賽人是研究舊約的專家.
他們知道舊約的先知預言彌賽亞是會來.
是會來要建立一個上帝的國.
他們一直在觀察彌賽亞的先鋒.
和彌賽亞的出現.
可惜在他們的眼裡.
他們看不到先使約翰就是彌賽亞的先鋒.
而耶穌就是彌賽亞.
他們真的比盲人更加盲.
盲人只是眼睛有問題.
法利賽人就是腦袋有問題.
在《約福音》第七章.
記載一些群眾說過一句話.
那些普通群眾.

$^{201}$他們說.
「基督來的時候會比這個人行更多的神蹟嗎?」.
耶穌所行的神蹟是非一般的神蹟.
只不過由於耶穌出身韓美.
那些法利賽人堅決不肯相信耶穌就是彌賽亞.
因為他們認為.
彌賽亞是一位軍事和政治上的強人.
彌賽亞是會帶來以色列國的復興.
雖然他們看到耶穌所行的神蹟.
但是他們不喜歡耶穌的背景.
但是他們其中一些的法利賽人.
很喜歡來和耶穌討論神學問題.
在福音書裡面.
我們見過有好幾位法利賽人.
來問耶穌問題.
有些是來試試耶穌.
看看他究竟有沒有實力.
有些是真誠地來求問耶穌.
例如在《約翰福音》第三章.
我們見到尼哥德姆夜晚去找耶穌.
今天經文二十節的法利賽人來問耶穌.
在中文聖經我們看不出.
究竟一個法利賽人還是有幾個呢?.
原文那一節其實是眾數.
即是有一群法利賽人提出這個問題.
至於是出於甚麼動機.
他們為甚麼要問這個問題呢?.
我們不知道.
可能是法利賽人自己有爭論.
來試試看耶穌的看法究竟是怎樣.
所以耶穌回答他們.
只用了一句說話.
是很簡單的.
「上帝的國來到.
不是眼睛看得見的.
人也不能說.
看那在這裡或說在那裡.
因為上帝的國就在你們心裡」.
甚麼意思呢?.
就是說神的國是在乎關係的建立.

$^{241}$過於看得見的榮華富貴.
如果我記得耶穌曾經在曠野.
遭遇魔鬼的試探.
魔鬼的其中一個試探是怎樣呢?.
在《馬太福音》第五章記載.
「魔鬼又帶他上了一座很高的山.
將世上的萬國和萬國的榮華都賜給他.
對他說:你若苦復敗我.
我就把這一切賜給你.
耶穌說:打翻!退去!.
因為經上記載:要敗主,利益上帝.
但要事奉他」.
所以我們看到人或魔鬼看角度.
看榮華富貴.
耶穌看上帝的角度.
是看關係.
上帝要求他的子民和他建立.
一個忠貞又充滿愛的家人關係.
如果我們全心全意愛主.
上帝的角就建立在我們的心裡.
未出來的時候.
不是像一位國家元首.
乘搭元首專用飛機.
預早通知別人.
等其他人預備禮炮,紅地毯去迎接他.
而是在人的心裡.
預備好地方去接納耶穌做我們生命的主.
詩體約翰和主耶穌都說:.
「天國近了,你們當悔改順福音」.
上帝的主權,上帝的角就在你們的心裡.
意思是如果我們願意接受上帝的角法.
接受上帝的說話來管理我們的人生.
上帝的角就會在我們心裡發生果下.
而上帝的角就已經臨到我們心中.
上帝的角是超越各界,超越民族的.
所以我們華人都可以成為上帝的國民.
只要我們悔改順福音.
上帝的角就已經降臨在我們心裡.
不過這只是上帝國權彰顯的第一階段.
上帝國權彰顯的第一階段是要將人心.

$^{281}$從撒旦魔鬼的權勢下搶回來.
這是一場屬靈的戰爭.
因此不是肉眼看得到的.
要憑信心去看.
法理賽人時時從外表去觀察.
所以他們看不到上帝直著主耶穌.
正在實現祂作為元首的主權.
我們生長在耶穌之後二千年的人.
可以看到上帝的主權在歷史中.
是戰勝了羅馬帝國.
是戰勝了前蘇聯的共產政權.
也戰勝了中共的四人幫高壓政權.
在四人幫時要將中國教會的建築物改變用途.
又要將聖經打入歷史博物館.
現在四人幫已經消失了.
教會仍然在中國存在.
現在全球每個國家都有信主耶穌的人.
這就是上帝國的實現.
也是天國已經降臨了.
但聖經的啟示.
天國的全面降臨還有第二步.
還有第二階段.
那就是上帝要重整地球.
審判世界.
在羅馬書八章二十至二十一節這樣說.
「整個被造的變成虛空.
不是出於本意.
而是出於上帝的旨意.
然而被造的仍然盼望著有一天能擺脫.
那會扭壞的枷鎖.
得以跟上帝的兒女分享光榮的自由」.
而天使其實也曾經對門徒說過.
「且離開你們被接升上天的耶穌.
你們見他怎樣升上天去.
又要怎樣來臨?」.
我不知有沒有人是期望著耶穌早點再回來.
每次當我聽到那些自然災害殺了不少人的時候.
就覺得主來的日子真是近了.
我也快望祂早日再回來.
為這個烏煙瘴氣.

$^{321}$還有很多地方都有戰爭的地方.
收拾殘局.
把那些罪惡滔天的人.
把那些犯了種族滅絕的人.
扔進地獄就對了.
但主耶穌在二十二至三十七節對門徒有三個提醒.
無論我們多渴望主耶穌早日再回來.
我們都要留意.
「免於被人迷惑」.
第一個提醒就是二十二至二十四節.
這個提醒可以叫做「不要受迷惑」.
當我們不斷渴望主耶穌早日再來的時候.
二十二節說「卻見不到」.
見不到的時候怎麼辦呢?.
人會有什麼反應呢?.
人的心理會怎麼反應呢?.
很可能會焦慮.
很可能會心急.
於是就看聖經.
計算日子.
試圖找出耶穌何時會再回來.
這些事情在歷史上發生過.
有很多異端.
基督教的異端就是這樣產生的.
舉個例子.
在1816年.
美國有一位浸信會的傳道人.
叫William Miller.
開始對基督復臨的道理產生興趣.
他研究聖經.
他以《但義理書》八章十四節.
認為這一節的寓言快要應驗了.
在1831年的時候.
Miller開始遊行報導.
寓言耶穌復臨的日期.
勸人應當悔改.
但是他幾次的寓言.
基督復臨的日子.
都沒有應驗.
怎麼辦呢?.

$^{361}$面子掛不住了.
於是就作個學說出來.
自圓其說.
只要說.
基督已經來了.
不過秘密再來.
很秘密的.
一部分的人已經被提走了.
你不懂的.
他也曾經推動過一個運動.
叫做「再臨運動」.
後來就影響了1863年成立的.
「基督復臨安息日會」.
我們今天也有.
叫做「Seven Day Adventist」.
就是這個教派.
在澳洲還有一個很小的教派.
叫做「Philadelphian」.
現在說起來.
可能三十多年前.
我工作的時候.
有一位同事屬於這一派.
他們也很有興趣.
研究聖經的寓言.
他當時幾乎肯定.
主耶穌二千年的時候就會再回來了.
我跟他有個辯論.
另外有一位韓國的牧師.
也曾經預言主耶穌在某一天會回來.
日期我不太記得了.
也是大約幾十年前.
在雪梨的.
因為我在雪梨的電視新聞報道看到這件事.
當時他預測耶穌哪一晚會再回來.
於是他鼓勵會友.
「耶穌再回來吧」.
「你們買好賣積了,沒有用的」.
那些賣積過期無效了.
結果那個日期到了.
那些弟兄姊妹.

$^{401}$如果有屋頂就走上屋頂.
等耶穌接他們.
結果沒有.
沒有發生.
那些韓國人喜歡打人.
比牧師還小.
所以就上了新聞.
所以主耶穌提醒門徒.
要小心.
無論什麼人告訴你.
在這裡或在那裡.
我相信無論是在某年某日.
「你們不要出去看也不要追隨他們」.
因為主耶穌再來的時候.
所有人都會看到.
都會知道.
不會只有一部分人看到和知道.
第二個提醒.
在二十五節至三十節.
第一個提醒是不要受迷惑.
第二個提醒是不要愛世界.
有些人等主耶穌第二次來.
等到不耐煩了.
自己預計祂再來的日子.
但是主耶穌何時再來呢?.
是上帝的主權.
其實我們沒有辦法知道.
主耶穌也說.
「可使他必須先受許多苦」.
「尤其是世代所棄絕」.
「世代」這個詞對我們來說.
可能是一百幾十年的時間.
但聖經對時代的觀念.
只有三個世代.
第一是舊約的時代.
第二是末世和主耶穌再來的.
三個主要的世代.
舊約時代,耶穌時代和主耶穌再來的時代.
耶穌誕生,末世就是耶穌來了的時代.
耶穌出生的時代就算是末世.

$^{441}$如果你在舊約聖經讀到末世.
就不要以為是最後結束的時代.
只是指耶穌會來的時代.
那就是末世的時代的開始.
如果耶穌再來就是第三個世代.
耶穌誕生之後就是末世的開始.
熟悉聖經和慣用和諧本的弟兄姊妹要留意.
末世和末期是不同的.
主耶穌在《瑪達福音》二十四章六節.
所提到的另一個詞叫「末期」.
「末期」不等於「末世」.
「末期」是指主耶穌再來之前一段短的時期.
「末世」是很長的.
「末世」是指主耶穌來到世上.
一直到第二次再來的時候.
都是末世.
這是從猶太人的角度去看.
很多基督徒看到主耶穌過了這麼久還沒有回來.
他們有另一個反應.
就是開始不太相信耶穌會再回來.
他們開始做什麼?.
「探外世界」.
就是只顧今生的享受.
主耶穌用了兩個歷史事件去提醒門徒.
不要「探外世界」.
第一件事件就是「羅亞洪水」的事件.
上帝通知羅亞告訴他.
他將會用洪水來毀滅世界.
羅亞相信上帝.
他就開始做方舟.
羅亞也將這件事告訴其他人.
但是其他人不相信他.
朝上吃喝嫁娶.
結果當洪水來到的時候.
只有羅亞一家八口得救.
另一件事發生在羅德時代.
當時所多瑪城的人沒有得到預先通知.
但是他們罪惡滔天.
惹得上帝要派天使去視察.
上帝執行審判之前一般都會先搞清楚情況.

$^{481}$然後用火燒所多瑪城.
羅德有一天的時間.
可以帶他的妻子,女兒.
和未來的女婿離開所多瑪城.
但是他的未來女婿不相信羅德的話.
他們就死在所多瑪城裡面.
而羅德的妻子逃難的時候.
違背了天使的吩咐.
回頭一看.
看看所多瑪城.
結果變成了一支鹽柱.
貪愛世界的人.
就會死在上帝的審判之下.
所以我們要警醒.
不要貪愛世界.
第三十至三十七節是第三個提醒.
主耶穌第二次再來的時候.
也會帶來審判.
主耶穌第一次來世是為了拯救.
第二次就是為審判.
因此第三個提醒就是.
不要只愛現世的生命.
現世的生命是寶貴的.
一個人如果沒有今生肉體的生命.
也不會有來世的生命.
但是我們要知道上帝給我們今生的生命.
是要我們把握機會.
用今生的生命去換取永恆的生命.
如果我們只顧著蓋房子.
顧著工作.
顧著吃喝玩樂.
當主耶穌第二次再來的時候.
我們就會被帶走.
去接受審判.
讓我總結一下.
今天經文有關天國的教導.
天國的降臨是分兩個階段.
第一個就是主耶穌第一次來世.
是先建立天國的國民.
拯救一批真心接受耶穌為主.

$^{521}$管理他們人生的人.
第二個階段就是主耶穌第二次再來的時候.
祂要審判世界.
而祂提醒門徒.
等祂再來的時候.
要注意三件事.
第一不要受迷惑.
隨便相信一些人對主再來的預測.
第二不要由於主耶穌遲遲未到.
就改變價值觀.
貪愛世界.
第三要記住.
主再來的時候是會私人審判.
到時我們就沒有悔改的機會.
已經信主的第一次門.
我們要反省自己的生命.
是否真正接受主耶穌的話.
作為我們人生的指南呢?.
是否一方面信耶穌.
另一方面就像生活到像微信的人一樣.
只記掛著安定的生活.
買屋,賺錢,工作,飲食,旅行,玩樂呢?.
主耶穌說.
「凡想保存聖名的,要喪失聖名」.
「凡喪失聖名的,要保存聖名」.
這裡說的是兩種的生命.
一種是現今的生命.
另一種是靈性的.
永恆的生命.
我們要把握機會.
用現今有限的生命.
去保證自己獲得那個永恆無限的生命.
我們一起祈禱.
至於被任何事物迷惑而失去了永恆.
奉耶穌基督的名祈禱.
阿們.
\newpage



\section{尼希米記 1:1-7}
\label{sec:032353}
\textbf{尼希米:以神的異象引領,以神的使命前行}~~~~劉振惠弟兄
\newline
\newline
連結: \href{http://www.church.com.hk/acms/content.asp?site=cdc\&op=show\&type=product\&code=032353\&layout=sermon}{\texttt{032353}} ~~~~ 語音日期: 2026-06-21
\newline
\newline
\hyperref[sec:032352]{\small{< < < PREV SERMON < < <}}
~
\hyperref[sec:index_chronic]{\small{[返順時目]}}
~
\hyperref[sec:032354]{\small{> > > NEXT SERMON > > >}}
\newline
\newline
尼希米記 1:1-7
\newline
\begin{longtable}{cl}
\hline
\hline
章節 & 經文 (和合本修訂版)\\
\hline
1:1 & \begin{tabularx}{0.7\textwidth}{X} 哈迦利亞的兒子尼希米的言語如下：亞達薛西王二十年基斯流月，我在書珊城堡中。 \end{tabularx} \\ \\ \relax
1:2 & \begin{tabularx}{0.7\textwidth}{X} 那時，我有一個兄弟哈拿尼，同幾個人從猶大來。我問他們那些被擄歸回、剩下殘存的猶太人和耶路撒冷的情況。 \end{tabularx} \\ \\ \relax
1:3 & \begin{tabularx}{0.7\textwidth}{X} 他們對我說：「那些被擄歸回剩下的餘民在猶大省那裡遭大難，受凌辱；耶路撒冷的城牆被拆毀，城門被火焚燒。」 \end{tabularx} \\ \\ \relax
1:4 & \begin{tabularx}{0.7\textwidth}{X} 我聽見這話，就坐下哭泣，悲哀幾日，在天上的神面前禁食祈禱， \end{tabularx} \\ \\ \relax
1:5 & \begin{tabularx}{0.7\textwidth}{X} 說：「唉，耶和華－天上大而可畏的神，向愛你、守你誡命的人守約施慈愛的神啊， \end{tabularx} \\ \\ \relax
1:6 & \begin{tabularx}{0.7\textwidth}{X} 願你睜眼看，側耳聽你僕人今日晝夜在你面前，為你眾僕人以色列人的祈禱，承認我們以色列人向你所犯的罪；我與我父家都犯了罪。 \end{tabularx} \\ \\ \relax
1:7 & \begin{tabularx}{0.7\textwidth}{X} 我們向你所行的非常敗壞，沒有遵守你吩咐你僕人摩西的誡命、律例、典章。 \end{tabularx} \\ \\
[1ex]
\hline
\hline
\end{longtable}
$^{1}$很久沒跟大家一起在講台上分享神的話語訊息.
很高興今天有機會.
尼希米以神的異象引領.
以神的使命前行.
弟兄姊妹在座當中的朋友.
我們線上的弟兄姊妹朋友.
早晨.
給我有機會再一次跟大家分享神的話語.
在預備今天的訊息裡面.
其實也是一個不小的挑戰.
大家可能都知道.
我在三月份和四月份.
都去了兩個頗長的海外旅遊.
祂親自預備了我的心.
在今天讓我可以跟大家一起思想祂的話語.
今天我會替松尼希米記.
跟大家一起思想.
我們事奉生命裡面一些重要的功課.
當然我們不可能在一篇的講道裡面.
講完整卷尼希米記.
全部十三章.
今天我主要是講第一章和第二章.
我相信透過這兩章的經文.
我們可以從尼希米的生命和祂的記錄作當中.
學習到什麼是神的異象.
什麼是神的使命.
今天的訊息.
不會是逐章逐節地詳細地解經.
我只是想將從尼希米記.
第一章第二章我所領受的一些重點.
跟大家一起思想一起放在心裡面.
也可以學習怎樣應用在我們日常生活裡面.
跟隨主事奉主的生命當中.
如果主許可的話.
將來有機會再分享尼希米記餘下的部分.
我可以分享關於屬靈的復興.
以及神子問怎樣在持守了.
和延續了屬靈復興的課題.
多少年來不少聖經學者.
都是從屬靈領袖的角度去研究尼希米記.

$^{41}$坦白說這些年日.
我見過不少很多曾經熱心事奉主的人.
他們後來都離開了事奉的崗位.
有些甚至離開了他原本聚會的教會.
每次我見到這些情況.
我心裡面其實都很難過.
因為我親眼看到他們起初很多都是滿有熱誠.
滿有愛心去服侍弟兄姊妹.
他們願意回應神的呼召.
而放下了很多東西.
有些花了很多時間去接受聖經的訓練.
有些甚至接受了全時間神學院程度的裝備.
預備自己更好地去服侍教會.
但事奉了幾年之後.
有什麼事發生呢?.
他們離開了事奉的崗位.
有些人也慢慢在信仰上冷淡了.
背後的內心歷程.
我相信一定是很複雜.
是很痛苦的.
其實這個問題不單止發生在.
牧者,長老,執事或者是教會的領袖身上.
很多時候弟兄姊妹都可能曾經熱心事奉主.
有些做過導師.
有些做過主學老師.
有些做過侍工組長.
有些很積極地參與教會的事奉.
有些願意承擔責任.
背負教會的需要.
但因為一些個人的掙扎.
家庭的壓力.
人際關係.
或者有時候是教會裡面的一些小摩擦.
他們就選擇放下了事奉.
有些人不再回應神起初給他們在生命裡面的呼召.
有些人轉到較大的教會裡面.
在崇拜人群裡面很安靜地聚會.
成為一個不再想參與事奉的普通基督徒.
你覺得這樣的人多不多?.
你有沒有認識這樣的弟兄姊妹?.

$^{81}$或者我們可以誠實地問一問自己.
誠實地問一問自己.
我們會不會是其中一個呢?.
我這樣問是因為我們現在生活在一個.
非常快,非常容易分心的年代.
我們很忙碌.
我們有時候忙到筋疲力盡.
我們不停地看手機,查信息,回信息.
忙著工作,忙著照顧孩子.
當然還要忙著去賺錢.
有時候忙著看劇集.
有時候忙著各樣的責任.
也承受了生活上很多的壓力.
我們今天看看聖經裡面的一個人物.
他很清楚自己是在服侍神.
他也很清楚自己在做什麼.
他更加知道為什麼要服侍神.
他的名字就是尼希米.
他的故事就是一個人被神意象緊緊抓住.
並被神的使命差遣出去的故事.
在我分享之前.
我想大家一起低頭禱告.
讓我們一起禱告.
親愛的天父,我們來到你面前.
求你引導我們.
一起進入尼希米記的經文.
求你幫助我們.
不單止聽見你的話語.
也能夠在我們的生活裡面.
活出你的話語.
我們這樣的禱告是奉主耶穌基督的名禱.
阿們.
弟兄姊妹,當我正式進入尼希米記之前.
我想先分享一下我自己.
屬靈旅程裡面我個人對神的意象.
使命的一些經歷.
我記得我是在高中最後一年順著的.
跟很多年輕人一樣.
當時其實我還沒有真正完全降伏在神的面前.
很多時候自己有需要的時候.

$^{121}$才會去神面前禱告一下.
可能隔幾個星期就回到教會崇拜.
但神對我的生命有很大的計劃.
好幾年之後.
我在三藩市灣區讀研究院的時候.
神接著不同的人.
不同的弟兄姊妹.
和一些環境.
喚醒了我的沉睡的靈魂.
聖靈覆興了我的心.
不久以後.
我開始在帕魯阿圖的主恩教會聚會.
也在那裡我認識了我敬愛的屬靈導師.
是堂會的主任牧師張唐牧師.
張牧師是一位很難得的牧者.
他有深厚的神學根基.
也有真誠的牧者心腸.
他牧養了我和我太太Winnie.
差不多十幾年的時間.
他來加州牧會之前.
張牧師其實是在加拿大Winnipeg.
牧養一間華人教會.
他影響了很多後來成為基督教領袖的信徒.
包括了香港突破的蔡元雲醫生.
前中信的羅曼華博士.
以及很多愛主愛神的僕人.
對我和Winnie來說.
主恩教會是我們事奉最長.
也是最深刻的屬靈之家.
當我預備這篇講道的時候.
尼希米基不斷地湧入我的心裡面.
尼希米是一位領受神意象.
並且全心投入神使命的僕人.
而聖靈也讓我看到.
其實尼希米和張牧師之間.
其實有很多美好相似的地方.
在1992年張牧師曾經邀請我考慮擔任.
當時教會的執事.
但我很有禮貌地婉拒了他這個邀請.
並且告訴牧師我和Winnie.

$^{161}$其實感到神可能很快就呼召我們回亞洲.
參與Tag Making.
大職宣教事奉.
張牧師不單沒有因此覺得被我拒絕了.
反而鼓勵我們要忠心地跟隨神的帶領.
差不多同一時間.
原來張牧師他自己也是在神的呼召中尋求.
透過他與香港建道神學院的親密關係.
他收到當時院長張牧海牧師的邀請.
請他回去加入神學院的任教.
這個邀請是由1992年.
來我們教會講道的梁家倫博士親自傳達給他.
1993年年初.
張牧師就向教會宣佈.
他相信神正是呼召他回到香港.
在神學院裡面訓練將來的牧者和領袖.
因此他會在幾個月之後就離開教會.
好好趕上神學院1993年秋季的開學.
老實講教會裡面有很多人感到驚訝.
首先他從來沒有強調過自己有博士的學位.
對我們來說他就是張牧師.
一位很謙卑,忠心,愛人,愛主的牧者.
有些教會領袖更加不明白.
為什麼牧師要離開一間他牧養了.
從很細小到成長的教會.
而轉而去到海外神學事奉.
當時他受到不少很友善的壓力.
極力希望他留下.
但張牧師仍然保持著謙卑,溫柔.
並且他對使命堅定不移.
他相信神已經將這個意象放在他心裡面.
也知道真正的信服一定要帶來實際的行動.
從人的角度來看.
牧師真的有很多很多理由可以留下.
但他願意放下這些理由.
為了要跟隨神的呼召.
當神賜下意象的時候.
信服就必須帶來行動.
即使這個行動會使信徒可能付上一些.
舒適,安全,經濟利益等方面的代價.

$^{201}$1993年我們教會參拜張牧師和他一家人.
以宣教士的身份前往香港.
參與神學的教育.
然後在1994年的下半年.
我和Winnie也憑信心踏出了一步.
我們離開了三藩市灣區的工作.
我們回到香港開展了我們在亞洲.
帶職宣教旅程的第一站.
回頭看我可以誠實地體會到.
張牧師的榜樣給我們.
我和Winnie有了特別的勇氣.
他願意信服在神的意象.
兼顧了我們的信心.
也幫助我們踏出自己信服神的一步.
所以我而今天更加明白.
為何聖靈在我預備這篇講章的時候.
將尼希米基放在我心裡面.
我覺得這些全部都不是偶然.
神會透過歷世歷代.
很多忠心信服.
愛祂的人的榜樣.
來一步一步塑造我們屬靈生命的方向.
對我來說.
尼希米的故事不單止是一個.
很熟悉的聖經歷史故事.
這個故事也告訴我們.
神今天仍然賜給他的子民神的意象.
仍然呼召他們進入神的使命.
仍然藉著那些願意回應.
主啊,我在這裡.
請差遣我的普通信徒.
去完成神的事工.
開始看經文的時候.
我們一起看看.
認識一下尼希米和他的時代背景.
尼希米是屬於舊約歷史後期的人物.
他生活在主前第五世紀.
即是離開今天大約有二千五百年.
在聖經裡面記載了三個比較特別的近東帝國.
第一個就是亞述帝國.

$^{241}$第二個就是巴比倫帝國.
這兩個帝國我們都較為熟悉.
因為他們導致了以色列國.
和猶大國滅亡的主要勢力.
而尼希米生活的時代.
就是第三個.
也是最後一個近東的帝國已經興起了.
就是波斯帝國.
也是今天大家也知道的伊朗.
由主前的第六世紀到主前的第五世紀.
波斯帝國曾經統治了一百二十七個杭省.
以斯帖記裡面也很清楚記載了這個數字.
他的版圖由亞洲的印度河流域.
一直伸延到今天的土耳其,埃及.
和北非一些的地區.
尼希米就是生活在波斯帝國裡面.
他就是住在波斯帝國時代的一個猶太人.
雖然他遠離耶路撒冷.
住在波斯的書山城裡面.
但他從來沒有忘記自己的身份.
他是神立約子民裡面的一份子.
聖經告訴我們尼希米是哈加利亞的兒子.
他有一個弟兄叫哈拿尼.
正就是哈拿尼將一個令人很難過的消息.
帶給了耶利米.
除了這些資料.
聖經其實記載他的家庭背景不是很多.
不過從尼希米禱告裡面.
我們能夠很清楚看到.
他應該是來自一個相當保存對著神.
對著神話與律戒律.
對以色列歷史很有深刻認識的家庭和群體.
那個時候猶太國已經滅亡了150年.
尼希米從小應該就很開始努力學習.
勤奮工作.
跟我們差不多的.
一步一步一步一步這樣爬.
最後他成為波斯帝國權力中心一個重要官員之一.
他成為了王的酒席.
首先我們看到尼希米不是祭司.

$^{281}$他不是先知.
就是我們今天說的那些神職人員.
他是王的酒席.
是波斯國裡面服侍君王的一個人.
我們不要小看酒席這個官職.
酒席不單止是那些斟茶遞水的僕人.
君王所有的食物和飲品.
都要由他最信任的僕人來試飲試食.
所以酒席其實是君王身邊一個非常親近.
非常受信任的人.
尼希米作為一個被征服的人.
民族的人.
竟然能夠爬到波斯帝國這個權力中心這個官職.
享受王宮裡面的尊榮和安穩.
真的絕對不是很簡單,很容易的事.
他一定是努力工作.
也似乎已經爬到在官場上的人生的高峰.
還有一點值得我們留意的地方.
尼希米的寫作是很特別的.
他讀起來好像我們正在看一本日記的寫作.
裡面有很多經文,很多語氣.
尼希米都不斷用「我做了什麼」.
「我想了什麼」.
「我要怎麼樣」.
聖經引用了尼希米自己寫下的一些記錄.
幫我們進入他的內心世界裡面.
尼希米是一位站在帝國權力中心的人.
他每天都住在王宮.
看見帝國的榮華富貴.
然後有一天.
他的弟兄哈瑪尼從耶路撒冷回來找他.
尼希米當然會問.
那些被擄歸回的猶大人情況怎麼樣?.
他的回答是.
「他們遭大能」.
「受凌辱」.
「耶路撒冷的城牆被拆毀」.
「城門被火焚燒」.
沒有城牆.
耶路撒冷就是沒有防衛.

$^{321}$沒有防衛是什麼呢?.
什麼仇敵都可以進來.
那座城就是暴露的,脆弱的.
也是被敵人隨時可以羞辱的.
如果是今天的人.
或者我們這樣.
可能聽完就算了.
因為人人忙於工作.
忙於家庭.
忙於眼前的生活.
但尼希米不是這樣.
因為耶路撒冷那裡是他列祖的家園.
聽見這個消息之後.
尼希米的心裡面失去了平安.
所以尼希米自己.
他自己親自寫下來.
「哦」聽見者說.
就坐下哭泣.
悲哀幾日.
在天上的神面前禁食.
經文告訴我們.
尼希米他悲哀.
他禁食.
他禱告.
是持續了好幾日.
這些肯定不是幾分鐘的難過.
三分鐘的熱度.
而是深層而持續的傷痛.
他悲哀他禁食.
他因為心裡沉重而無法吃喝.
他在禱告中將這個重擔.
帶到神的面前.
尼希米眼前所見到的.
每天看到的就是波斯王宮的榮華.
但他心裡面看到的.
卻是耶路撒冷的荒涼.
當時尼希米應該不是.
很年輕的一個人.
要爬到酒鎮這個位置.
是不太容易的.

$^{361}$他一定要有相當人生的歷練.
我們說過.
那些人說打無人巷過很多關.
不少的聖經學者都認為.
他很可能已經是超過了五十歲了.
超過了五十歲.
用今天的說法是什麼呢?.
尼希米應該是差不多要開始.
為退休作準備了.
這個時候他應該想什麼呢?.
如何安享晚年呢?.
想想將來如何安排退休計劃.
安排去哪裡大大的旅行.
走走呢?.
但正正在這個時候.
神.
是神自己.
將耶路撒冷的荒涼放在了他的心裡面.
所以他不能夠安心吃喝.
不能夠睡覺.
他的心一直都沒有平安.
我相信神是用了幾十年的時間.
去預備尼希米的.
他的經驗.
他的資源.
人脈.
地位.
都不是偶然得來的.
這些都是.
有神幫助他.
這些有些像是神在摩西的生命裡面.
所做的事.
記得摩西嗎?.
神用各樣的經歷和處境.
預備摩西.
將來帶領以色列人離開埃及.
如今神亦要大大的使用耶利米.
我們看看耶利米接下來做了些什麼.
耶利米等候機會.
差不多等了四五個月.

$^{401}$從基斯留月.
即是現在我們的十一,十二月左右.
一直到離散月.
即是我們現在的三,四月.
終於有一天.
很特別.
王問他.
你既然沒有病.
為什麼面帶愁容呢?.
這句話對尼希米來說.
其實是不是很危險呢?.
因為酒精.
事奉王的.
怎麼可以在王面前露出愁眉苦臉呢?.
用香港人的話.
不想撈嗎?.
是不是?.
如果真的有病的話.
根據現在的 COVID-19.
根本就不應該回來工作了.
王的真正意思是什麼呢?.
我知道你心裡面有事.
尼希米非常的害怕.
因為如果他回答錯了一句話.
他隨時就會失去工作.
而且在君王面前.
隨時也會失去一條命.
但他仍然鼓起勇氣.
尼希米沒有直接提到.
「耶路撒冷」這個字.
而是說.
「我列祖墳墓所在的那城荒涼」.
「城門被火焚燒」.
「我怎能不愁容滿面呢?」.
然後王就問他.
「你要求什麼?」.
聖經記錄這一刻非常有意思.
因為尼希米接著記載了.
「於是我默討天上的神」.
大家看到了嗎?.

$^{441}$看到了嗎?.
當王問他這麼重要的問題的時候.
尼希米先做了什麼事情?.
在他開口回答王的請求之前.
他先禱告神.
然後尼希米才提出他的請求.
結果很特別.
王全部都批准了.
為什麼呢?.
今天作為一個基督徒.
我們知道為什麼.
因為神已經為尼希米預備到老.
這是美麗而真實的故事.
可以說是真實的故事.
很清楚記載在聖經裡面.
成為我們今日後世信徒的見證.
在很多方面.
我們和尼希米都有點相似.
我們都在異鄉美國.
打拼努力工作.
經過多年的打拼.
我們都在某種程度上.
好像有舒適和安穩.
我們有家庭要照顧.
有些人快要退休.
有些人已經退休.
好像我這樣.
我們都在事業和家庭生活裡面.
追求安穩.
坦白說.
每個人都喜歡舒適安穩.
不喜歡遇到一些生活上的困難.
特別是在華人的文化裡面.
我們常常傾向於避開一些風險.
追求穩定.
意義上是藉著困難磨練去塑造他的僕人.
很多人對困難的第一個反應就是逃避.
不要管.
今天的文化充滿了一些聲音就是.
如果這麼辛苦.

$^{481}$不如嘗試選擇另一條路走走.
慢慢地人就開始走一條比較容易走的人生道路.
但服侍臣很多時候不是這樣的.
尼希米原本在波斯帝國的書山城.
他有穩定的職位.
亦有安全的前途.
但神呼召他離開他的舒適圈.
回到耶路撒冷.
回到一個充滿廢墟和破牆的城市.
跟著一步.
尼希米聽見耶路撒冷城牆被拆毀.
城門被火燒.
他的心深感憂傷.
他不是聽完就算了.
他沒有說到.
這個都不是我的問題.
總會有.
或者這麼說.
總會有其他人去處理.
相反地.
他為神的誠哭泣.
他為神的百姓禁食禱告.
我想再提醒大家一個很重要的細節.
第一章開始的時候是什麼?.
是基斯留月.
到第二章的時候已經是彌撒月.
中間大約過了四五個月的時間.
換句話來說.
尼希米不是一時衝動.
不是一時衝動.
因為我們很多時候.
因為一種意一陣孽心.
一陣熱心就立刻行動.
但他首先來到神面前.
讓神塑造了他的心.
尋求神確認了他的呼召.
很多時候.
當神將一個負擔放在我們的心裡面.
我們以為我們要立刻行動.
立刻起來做些事情.

$^{521}$但神很多時候是要我們先等候一下.
等一等.
等候的目的不是拖延.
而是潔淨我們的心.
神要潔淨我們的動機.
我們的心,愛心.
很多信徒當意象來到的時候.
就大發熱心.
但過一段時間就冷卻了.
所以學習等候不是浪費時間.
等候本身就是神塑造他僕人的過程之一.
尼希米在皇宮中做酒精.
但那天他的憂傷再藏不住的時候.
王看到.
王就對他說.
你既然沒有病.
為什麼面帶愁容呢?.
這不是別的.
必是你心中受煩.
尼希米回答一句.
我甚懼怕.
尼希米這句話其實非常真實.
我們常常以為肅靈成熟的人就是.
沒有怕的人.
他沒有怕.
但其實不是這樣的.
真正的勇氣不是沒有懼怕.
真正的勇氣是即使懼怕.
我們仍然選擇信服神.
尼希米的害怕.
雖然害怕.
但他仍然拿著神賜予他的機會.
我們必須有肅靈的敏銳.
學習捉住神為我們打開門的時刻.
很多時候.
有些人說.
如果神要我做一些事奉.
他自然就會開一條路讓我走.
這句話聽起來也挺對的.
但真正的挑戰是.

$^{561}$當神打開門的時候.
我們準備好了沒有?.
我們敢不敢走出去?.
四個月之前.
尼希米已經開始禱告了.
四個月.
如今神終於打開機會.
直接問他你要求些什麼.
這是不是難逢的機會?.
如果尼希米沒有預備好的話.
他可能當場就呆了.
讓我想想.
說不出他的要求.
但我們看見尼希米其實早已經預備好了.
他知道方向.
他知道使命.
他知道神放在他心裡的負擔.
所以當機會一來到的時候.
他就能夠很清楚地回應君王的問話.
很多時候人不是沒有機會.
而是機會來到的時候.
人還沒有預備好.
我們每天的禱告.
我們每天的靈修.
我們每天的讀經.
正正就是神.
要我們在生命裡面.
要我們做好預備的鍛鍊.
所以不要小看這些靈修,讀經,禱告.
接著我想再請大家留意一下.
尼希米第二章一句非常寶貴的話.
「於是我默討天上的神」.
這可能是聖經裡面較為最短的禱告之一.
為什麼呢?.
因為尼希米根本沒有時間做很長的禱告.
王問了這個問題.
他在等你回答他.
你有多少時間?.
尼希米是沒有時間禁食.
沒有時間退休.

$^{601}$也沒有時間急拖幾個弟兄姊妹.
一起開一個禱告會.
就在短短幾秒的時候.
尼希米先向神禱告.
然後他才向王說話.
一個真正敬虔的人.
不單止是有固定靈修的生活.
更加重要的是他懂得每個時刻.
特別是一些關鍵時刻.
都依靠神.
那個時候根本沒有時間做很長的禱告.
但我們仍然可以學習.
好像尼希米那樣禱告.
如果遇到困難的時候.
我們可以說.
「主啊!主啊!主啊!幫助我!幫助我!」.
這些禱告可能往往比很多長篇的禱告.
更加真實.
尼希米對王提出了幾個請求.
很令人驚訝的就是.
王全部都答應了他.
尼希米沒有將功勞歸給自己的口才.
也沒有歸給自己.
「我和王很特別,很友善的」.
王只是說了一句話.
「王就允許我」.
為甚麼呢?.
「因我神施恩的手幫助我」.
尼希米知道.
成功不是靠他自己.
不是靠他的能力.
不是靠他的地位有多威風.
如果神不成就的話.
再好的計劃都不會成功.
但如果神工作施恩典的時候.
沒有一個困難是我們不可以跨越的.
接著我會將餘下的時間.
轉去看看第二章的記載.
講述了很多尼希米和王在書山城.
尼希米終於到了耶路撒冷.

$^{641}$他在那裡住了三天.
然後在夜間他帶著幾個助手.
出去看一看城牆.
他親眼看見那座城的破敗.
有些地方損毀得非常非常非常嚴重.
甚至他所騎的束手.
牲口.
都過不到去.
他只好原地折返.
聖經記載了尼希米到達耶路撒冷之後.
沒有急於立刻召開會議.
工作會議.
沒有立刻宣佈他已想好的計劃.
也沒有急於展示王賜給他的權柄.
地位.
相反的他先做了一件非常重要的事.
是甚麼事?.
他親自去了解真實的情況.
尼希米沒有急著說.
他先去看.
很多時候我們在事奉中有一個問題.
就是在事奉中我們可能說得太多.
太快.
卻看得太少.
真正屬靈的領袖應該願意花時間.
去多了解真實的情況.
尼希米發現情況比想像中更嚴重.
城牆倒塌.
城門被燒.
有些地方爛得完全不能通行.
那一刻他看見的不單是人.
所謂的土木建築工程問題.
而是一個屬神子民的大傷口.
因為城牆代表了安全.
代表了身份.
代表了神子民的見證.
如今城牆倒塌了.
也象徵了神的子民長久活在羞辱當中.
最令人驚訝的.
是這些廢墟不是剛剛才出現的.

$^{681}$這些廢墟已經存在超過一百年了.
很多人每天都經過.
每天都看見.
但是什麼都沒有做過.
什麼都沒有改變過.
時間久了人就開始習慣問題了.
人就開始接受失敗.
人開始相信事情本來就是這樣的了.
是這樣的了.
這是非常危險的.
最大的危機不單止是有問題的存在.
而是人已經對問題麻木了.
今天可能有些教會信徒.
可能他是這樣.
有些問題可能已經存在了很多年.
比如我們缺乏了禱告的生活.
缺乏了門徒訓練.
缺乏了傳福音的熱誠.
缺乏了認真地對我們下一代的信仰傳承.
心裡面都知道這些是問題.
但是如果心裡面不相信改變.
仍然可以通過大能的神發生的話.
問題就會變成正常.
廣東話一句話.
是這樣的了.
但看尼希米.
他有沒有退縮?.
他沒有退縮.
反而他召集了百姓.
對他們說.
你們所遭的難.
耶路撒冷怎樣荒涼.
城門被火焚燒你們都看見了.
內罷.
我們重建耶路撒冷的城牆.
免得再受凌辱.
而你要明白.
這一切都不是偶然的事情.
他那天.
站在那裡.

$^{721}$跟百姓說.
不是因為他自己有什麼大能力.
有什麼糾正的地位.
是因為神施恩的手.
一步一步地引導他回到耶路撒冷.
他一生在波斯所做的工作.
他所有的成就.
最終都是神為尼希米.
這個使命作的準備.
重建城牆.
使百姓在屬靈上重新被建立.
才是神賜給尼希米的一生最大的使命.
今天很多信徒都說神的意象和使命.
但我們知道神的意象不是個人的理想.
神的使命也不是NATO.
不是那個NATO.
是NO ACTION.
TALK ONLY.
只說不做.
真正神的意象和使命.
是在神的恩典綱照中.
我們清楚看見神打算要成就什麼事情.
也是憑信心認定.
神一定會成就.
確實實現在我們眼前.
這些不是空中流角.
更加不是只說不做.
更重要就是我們不可以將自己的東西.
比如我們的理想.
我們的想法.
特別是我們的喜好.
加入了神的意象裡面.
意象是發自於神自己.
也單單屬於神.
我們只是領受了神在我們生命中所賜的意象和使命.
在尼希米公開呼召百姓行動之前.
我們看到他已經經歷了.
哭泣.
禁食.
禱告.

$^{761}$等候.
和加上很仔細的觀察.
我們要看到帶領事奉是一種屬靈的責任.
尼希米沒有責怪任何人.
沒有翻舊帳.
也沒有找代罪羔羊.
他只是說.
你們所遭的難.
耶路撒冷怎樣荒涼.
城門被火焚燒.
你們都看見了.
來吧.
我們重建耶路撒冷的城牆.
免得再受凌辱.
尼希米接著分享.
他分享什麼呢?.
我告訴他們.
我神施恩的手怎樣幫助我.
並和對我所說的話.
百姓就回應說.
我們起來建造吧.
我們請大家留意一下.
尼希米鼓勵百姓的方法.
不是我們今天的營銷口號.
也不是我們現在所說的高漲情緒.
情緒價值.
也不是靠 peer pressure.
群眾的壓力.
尼希米是藉著見證.
藉著見證他自己真實經歷過的神的作為.
來鼓勵百姓.
真正能夠激勵神的子民.
永遠都是神自己.
當人看見神在工作的時候.
信心就被挑撥.
當人看見神為我們開路的時候.
我們的勇氣就被建立.
尼希米非常明白.
一個人不可能重建城牆.
無論他多有能力.

$^{801}$都不可能一個人完成.
很多事工可能的失敗.
不是因為缺乏好的計劃.
而是缺乏愛.
缺乏團隊.
也缺乏一切事奉的夥伴關係.
所以聖經記載了.
尼希米第一個任務.
不單止去找一些資源.
也需要去召集同工.
召集同工.
神的侍工從來都不是要我們單打獨鬥.
從摩西到保羅.
從耶穌到初期教會.
神一直都是藉著團隊的工作.
來完成神的使命.
今天的教會也是一樣.
如果一切都集中在少數個人身上的話.
burn out 疲乏一定會出現.
一間真正健康的教會.
是需要每一個弟兄姊妹.
都參與建造其中.
弟兄姊妹記不記得這個?.
我們T4C教會從去年到今年.
2016年的主題是什麼?.
就是神的使命.
我們的參與.
我也想問一問會眾一個問題.
你有沒有在你自己的生命裡面.
看見異象.
看見使命.
以及想想怎樣參與?.
有時候有些教會的挑戰.
其他人都可能認為.
解決不了.
有些事工其他人可能認為.
沒有能力了.
沒有人做了.
但是神就偏偏讓你見到.
也許正正就是別人看不到的.

$^{841}$也許正正就是別人沒有的呼召.
別人沒有的經驗.
別人沒有的能力.
你起碼看見的不單止是廢堆.
廢墟.
他看見神的使命.
你起碼一生最大的不是在皇宮作酒精.
而是回流到耶路撒冷重建城牆.
今天.
神為什麼要使用我.
使用你呢?.
也許你只是在事業打拼當中.
也許你快要退休了.
也許你已經退休了.
好像我這樣.
或者你已經覺得人生最好的年日.
已經過去了.
在後面.
或者你看見前面.
期待明天會更好.
更精彩.
但是聖經告訴我們.
神一直都是不停地呼召.
預備,準備,差遣.
那些信靠祂的子民.
兒女和門徒.
不分年齡的.
不分地位.
不分學識.
神一步一步帶領我們.
要到達祂要我們站立站穩的位置.
我祈求神會讓我們每一個跟隨祂的弟兄姊妹.
有祂從祂而來的意象和使命.
我今天的訊息是用一個個人的見證開始.
我希望也想用我自己.
在屬靈旅程中一個深刻的提醒.
最近的.
來結束今天神話語的分享.
可能很多弟兄姊妹.
會提到我太太.

$^{881}$坐在前面.
知道Winnie和我在五月中.
剛剛從英國回來.
Winnie就是被英國Birmingham.
一間新成立的教會.
邀請去帶領兩堂.
關於基督徒的戀愛和婚姻的小組聚會.
起初我沒有打算去.
因為剛剛從亞洲回來.
我兩個星期就要離開.
我時差還沒倒過來.
她就說要走.
我說不想去.
但Winnie對我說.
很有說服力.
你一定要跟我去.
這個講座是關於基督徒的婚姻和戀愛.
如果一個人去的話.
好像沒有那麼有說服力.
最擅長的那句話就是.
是不是呢?.
是.
於是我就去了英國.
我以為自己去支持這個講座.
大錯特錯.
因為我很快發現.
我的角色變成了一個司機.
開租回來的車.
接送人.
帶朋友去不同的地方旅遊.
英國.
去牛津大學.
劍橋大學.
莎士比亞小鎮等等.
其實我享受這個悠遊休息的時間.
但我心裡面仍然會問神.
主啊.
你叫我來了.
又要倒時差一點.
你到底要我在這次的英國旅程中.

$^{921}$得著什麼的屬靈領受呢?.
有一天我們計劃去.
布里斯多的.
布里斯多這個城市.
出發前一晚有一位姐妹.
建議我們去參觀布里斯多的.
Georgetown.
Mueller Museum.
佐治·慕勒的紀念館.
我以前聽過慕勒的事蹟.
特別是他的孤兒院事工.
在十九世紀.
神使用了他.
在英國照顧了超過一萬名的孤兒.
奇怪的就是我參觀紀念館的時候.
吸引我最大注意力的.
是神在他七十歲.
七十歲以後.
呼召他做的新的事工.
按世界的看法.
七十歲了.
七十歲是什麼?.
應該是退休的年齡.
但是慕勒在七十歲.
才開始一個全新的使命.
他在三十多年.
三十多年的事奉.
孤兒院的事奉之後.
他成為一位海外巡迴.
福音報道的宣教士.
慕勒是有充足的屬靈準備的.
剛才說了屬靈準備.
他五十七年不斷地讀經.
他研讀神的話語.
他一生讀完.
我都希望有他十分之一.
整本聖經是超過二百多次.
所以他能夠成為宣教士.
做福音報道工作.
由七十歲到八十七歲.

$^{961}$他行走超過二十萬英里.
去過我們熟悉的中國上海.
日本東京.
印度.
美國.
墨西哥.
澳洲.
紐西蘭.
東歐西歐.
不同的城市.
不同的國家.
八十七歲他才停止了他的宣教旅程.
神好恩待他.
九十二歲他就安息主懷.
我看見的時候心很有感動.
這是何等美好的見證.
我今天也到了六十六歲.
很多人經常聽著我提.
我還會不斷地問神主.
你是否仍然有新的使命給我.
是否會使用我和Winnie.
去服侍教會.
去服侍弟兄姊妹.
去服侍人.
我們必須為神的遺像.
使命和參與.
繼續不斷地去禱告尋求.
願我們的生命被神的遺像引領.
願我們也受到神使命的差遣.
也願我們在神的呼召我們的時候.
好像黎希敏和百姓一樣.
謙卑而堅定地回應神.
我們一起低頭作結束的禱告.
好嗎?.
親愛的天父.
藉著今天的訊息.
你再一次提醒我們.
關於神的遺像.
使命.
和我們的參與.

$^{1001}$願我們好像黎希敏和耶路撒冷的百姓一樣回應說.
我們起來建造吧.
願我們今天的禱告.
我們的禱告是奉主耶穌基督的名禱.
阿們.
\newpage



\section{約翰福音 4:43-54}
\label{sec:032354}
\textbf{慈父憑信見神蹟,得人間三寶}~~~~許開明牧師
\newline
\newline
連結: \href{http://www.church.com.hk/acms/content.asp?site=cdc\&op=show\&type=product\&code=032354\&layout=sermon}{\texttt{032354}} ~~~~ 語音日期: 2026-06-21
\newline
\newline
\hyperref[sec:032353]{\small{< < < PREV SERMON < < <}}
~
\hyperref[sec:index_chronic]{\small{[返順時目]}}
~
\hyperref[sec:032355]{\small{> > > NEXT SERMON > > >}}
\newline
\newline
約翰福音 4:43-54
\newline
\begin{longtable}{cl}
\hline
\hline
章節 & 經文 (和合本修訂版)\\
\hline
4:43 & \begin{tabularx}{0.7\textwidth}{X} 過了那兩天，耶穌離開那地方，往加利利去。 \end{tabularx} \\ \\ \relax
4:44 & \begin{tabularx}{0.7\textwidth}{X} 因為耶穌自己作過見證說：「先知在自己的家鄉是沒有人尊敬的。」 \end{tabularx} \\ \\ \relax
4:45 & \begin{tabularx}{0.7\textwidth}{X} 到了加利利，加利利人都歡迎他，因為他們也上耶路撒冷去過節，曾經看過他在節期間所做的一切事。 \end{tabularx} \\ \\ \relax
4:46 & \begin{tabularx}{0.7\textwidth}{X} 耶穌又到了加利利的迦拿，就是他從前變水為酒的地方。有一個大臣，他的兒子在迦百農病了。 \end{tabularx} \\ \\ \relax
4:47 & \begin{tabularx}{0.7\textwidth}{X} 他聽見耶穌從猶太到了加利利，就來見他，求他下去醫治他的兒子，因為他兒子快要死了。 \end{tabularx} \\ \\ \relax
4:48 & \begin{tabularx}{0.7\textwidth}{X} 耶穌對他說：「若不看見神蹟奇事，你們總是不信。」 \end{tabularx} \\ \\ \relax
4:49 & \begin{tabularx}{0.7\textwidth}{X} 那大臣對他說：「先生，求你趁著我的孩子還沒有死就下去吧。」 \end{tabularx} \\ \\ \relax
4:50 & \begin{tabularx}{0.7\textwidth}{X} 耶穌對他說：「回去吧，你的兒子會活！」那人信耶穌所說的話，就回去了。 \end{tabularx} \\ \\ \relax
4:51 & \begin{tabularx}{0.7\textwidth}{X} 正下去的時候，他的僕人迎面而來，說他的兒子活了。 \end{tabularx} \\ \\ \relax
4:52 & \begin{tabularx}{0.7\textwidth}{X} 他就問甚麼時候見好的。他們對他說：「昨天下午一點鐘熱就退了。」 \end{tabularx} \\ \\ \relax
4:53 & \begin{tabularx}{0.7\textwidth}{X} 他就知道這正是耶穌對他說「你的兒子會活」的時候；他自己和全家就都信了。 \end{tabularx} \\ \\ \relax
4:54 & \begin{tabularx}{0.7\textwidth}{X} 這是耶穌從猶太回到加利利後所行的第二個神蹟。 \end{tabularx} \\ \\
[1ex]
\hline
\hline
\end{longtable}
$^{1}$今日是父親節主日.
我們的題目就是.
慈父憑信見神蹟.
得人間三寶.
這個題目我會說什麼呢?.
第一個重點我會說.
憑信會見到神蹟.
接著會說人間有三個寶.
這三個寶是什麼呢?.
現在先不說,一說就很長了.
這一篇道是很適合今天洗禮的.
幫你去知道你到底信耶穌是信什麼.
或者你說我知道許牧師.
你知道的幫你重溫.
福音主日也不單止是對未信.
或者是已經信,初信的.
對每一個信徒都會有提醒.
到底我們是不是有憑信見神蹟.
得人間的三寶.
我們講道之前一起禱告.
主耶穌我們來到要敬拜你,讚美你.
之後我們要來聆聽主你的話語.
願你幫助我們能夠聽到你的道.
明白你的道.
也都跟到你的道.
以致我們的生命得到成長.
奉主名求,阿們.
今天的經文.
就是在約翰福音四章四十三到五十四節.
經文我不讀了.
因為四章四十三到五十四節.
Benny牧師好像跟我同感一靈.
他五月講了,叫做信心的建立.
如果我沒有記錯.
所以這段經文大部份人看了.
我只是想大家讀跟著那裡.
請退一章,第五章.
第五章.
第四章十四到六至七節.
我們一起讀.

$^{41}$耶穌對他說,我就是道路,真理生命.
若不直著我,沒有人能道負那兩句.
今天我們比較重點是這一句.
好,我們跟著.
憑信見神跡.
意思是從神跡認識主耶穌.
自古到今.
都是有很多人認為.
要看到神跡才信.
證明神的存在才接受.
但是聖經的回應是什麼呢.
我們信了神.
就先看到,見到神跡.
如果你未信神.
你是見不到神跡的.
你需要先信了神.
所以不信神的人.
你是看不到神跡的.
這句說話背後.
可能大家都有點懷疑.
到底耶穌想說什麼呢.
你聽不明白.
一會兒希望你明白.
這一段經文的背景.
就是大臣的人生困境.
這位大臣,聖經告訴我們.
這個大臣的兒子在加伯隆患了病.
這個大臣不是普通人.
他可能是希律王下面的官員.
甚至可能是分封的王.
他掌握有軍權和政權.
所以他是一品大官.
他有地位,有權力,有財富.
但是當他的兒子有病的時候.
病得快要死了.
他的權力沒辦法救到兒子.
他的財富也不能買到兒子的生命.
他在束手無策的時候.
找耶穌.
人生最大的問題.

$^{81}$往往不是因為你沒有資源.
或者沒有地位,沒有身份.
人的生命的問題.
是你生命有限.
你的生命會面對很多很多的限制.
今天英國社會.
你看它好像科技都進步.
醫療發達.
但是英國.
我們又發現它經濟衰退.
失業率偏高.
政府管治能力越來越薄弱.
種族主義興起.
在英國焦慮增加.
抑鬱增加.
孤獨增加.
家庭破碎又增加.
人生終極的問題.
是有關個人生命的問題.
「憑信見神蹟」.
這句說話是誰說的呢?.
是耶穌說的.
在這段經文裡面.
四十八節.
耶穌對他們說.
「若不看見其事.
你們總是不信」.
這句說話是當大神要求耶穌.
你快點下去救我的兒子.
耶穌就對他說.
其實不只是對他.
因為耶穌是用你們的.
若不看見神蹟.
你們那裡那麼多位.
總是不信.
神蹟其事到底是什麼呢?.
神蹟在《約翰福音》的了解.
是一個「Sign」.
「Sign」的意思是什麼呢?.
就是神的蹤跡.

$^{121}$就是說透過一些神蹟.
神蹟一般人理解.
就是「奇事」.
「Miracle」.
令人驚訝的.
但是《約翰福音》的重點是.
神蹟是要顯示有神.
和神的大能作為.
這樣去理解神蹟才是對的.
因為「奇事」,「大能」.
有些術士都做到的.
先知也可以做到的.
但是神蹟就不是他們可以做到的.
因為這個大神能夠.
「兒子得醫治」.
不是普通的.
我們深入去看他.
耶穌是要透過大神.
信他.
如果他沒有信他.
這個神蹟就不會出現了.
要信了耶穌.
他才會見到這個神蹟.
「因信」是可以得到三種的寶.
人生的三寶.
是什麼呢?.
道路,真理,生命.
我們一起說.
道路,真理,生命.
這三個我們先來看第一個.
耶穌說「我是道路」.
意思是憑相信有位神.
有位主.
你就會見到神蹟.
這個神蹟就是要讓你看到.
第一樣.
耶穌不是一般的人.
耶穌是神來的.
各位,時至今日有很多人都說信耶穌.
但他不是信耶穌是神的.

$^{161}$摩門教那個耶穌不是神來的.
是先知來的.
甚至法輪功那個.
就算是神.
都是其中一個.
而且好像說是假的神.
不是一個真的神.
基督教那個的神.
你信了有神之後.
你的人生有出路.
這個是聖經說我們知道的.
耶穌說他是道路.
道路是什麼?.
道路就是出路.
而且能夠去到神那裡.
當人生是無路可走的時候.
耶穌就是你的道路.
留意英文或者原文.
The way, the way.
唯一,唯一的路.
各位很多人是要求神拜佛.
很多時候都是去遇到人生.
已經去到束手無策.
無路可走的時候.
他找不到出路.
神要告訴我們.
如果你是相信有神.
相信主耶穌.
你的生命就有出路.
耶穌是道路.
對人有什麼意義呢?.
我們可以去到神那裡.
留意耶穌不是指示人.
怎樣去找到神.
耶穌是告訴你.
你想去認識神.
你來找我.
我就是神.
耶穌是道路.
這個對人的意義就是說.

$^{201}$在人生迷失的時候.
祂會給方向.
人生感到絕望的時候.
祂會有出路.
人生去到臨死.
好像我這樣臨死我不怕.
因為我知道我會去哪裡.
我主懷安息.
所以這個道路的意義.
是很有意思的.
各位.
你有沒有用GPS?.
GPS失靈有沒有試過?.
我認識一位英國的長者.
他是西人.
他從伯明翰要去到蘇格蘭.
GPS在半路故障.
他越走越迷失.
老人家好像我的動作很煩躁.
很煩躁.
越走越不知道應該怎樣.
結果他都想著.
我自己都教人開車.
都要問人.
但他不知道怎樣走.
問那些人.
他跟那些人說.
我不是不懂開車.
我只不過是不知道去哪裡.
各位人生.
你有沒有試過有些情況.
你不知道走向哪裡.
不知道走去哪裡好.
在這個人生.
你會遇到.
不知道怎樣走下去.
我四歲就信耶穌.
因為我四歲的時候患了病.
在患病的過程中.
我很怕死.

$^{241}$因為我前面的哥哥都死了.
剩下我.
當然輪到我第四個.
這個四又不好.
死死死.
所以很怕.
我那時候不是很懂.
但阿麻說這個世界有個神.
我就相信神.
你真的醫好我.
我相信你.
一生跟隨你.
雖然小時候.
都很有屬靈的思想.
結果我真的康復了.
好了.
耶穌說第二樣.
祂是真理.
憑信相信主.
相信主的話主的道.
你就會看到真理.
十七言.
耶穌是真理.
對人有什麼意義.
真理是什麼意思.
可信的.
真實的.
永遠不落空的.
總是不改變的.
各位.
你有沒有看到.
英國的政治是不停變的.
很多的制度不停變的.
而且具有永恆價值.
有意義的.
今天這個世代.
你會發現很多東西不可信的.
謬理可以變成真理.
真假的理.
也有很多人相信的.

$^{281}$很多東西不對現的.
不穩定的.
沒有價值.
沒有真理的.
耶穌是真理.
是要使迷失的人.
找到他生命的答案.
找到他應該走的路.
我們看看大神.
他是怎樣發現真理的.
第一.
他起先對耶穌是不夠信心的.
他對耶穌就算有信心.
是一個有限的信心.
他要求耶穌怎樣.
求你立刻和我下去救我的兒子.
醫好我的兒子.
耶穌沒有答應他的.
耶穌沒有答應他的目的.
不是不醫他的兒子.
而是要改變他.
不是你看著我.
走這個神跡你才信我.
你看不到神跡的時候.
你能不能夠信我呢.
今天基督徒我們都經常面對這個挑戰.
我們遇到困難的時候.
我們可能會對講主照我的時間.
照我的方法.
照我的計劃.
這個大神都一樣.
你當然要去到我家裡.
你才能醫好我的兒子.
醫生當然是這樣.
不會說不去你家裡.
這樣就遙遠遙距.
但是耶穌說.
他不會下去.
不過你的兒子已經康復了.
那麼這個大神有沒有聽.

$^{321}$他有沒有聽.
聖經說.
他聽到之後.
他就真的下去.
耶穌提升了他的信心.
他不是看到神跡信.
而是看不到耶穌怎樣走神跡.
但是他信什麼.
信耶穌說.
其實他不信都不行.
因為耶穌不下去.
難道你敢不敢捉他下去.
他可以他是大神.
他可以找些兵拉他下去.
但是他沒有這樣做.
他真的相信.
因為耶穌曾經提醒他.
很多人都是要求神跡之信.
耶穌會邁他提醒他.
你能不能夠看不到神跡.
你都能夠信我的說話呢.
他說這個要有很大的信心.
22.
耶穌說.
回去吧.
你的兒子活了.
那人就信耶穌所說的.
就真的回去.
但是他看不到的兒子真的好.
也都沒有現在的手機.
可以打個電話回家.
看看兒子是不是搞定.
沒有的.
他真的要憑信心回去.
結果怎麼樣呢.
憑信經歷看到神跡.
神跡發生.
這個大神在回家的途中.
他的僕人可能正在跑來通知他.
知道他會回來.

$^{361}$跑的時間.
那個僕人就說.
你的兒子已經活了.
活了除了是子醫好.
也都有可能他的兒子.
曾經病到死了.
但是突然之間活了.
這個不是巧合.
這個是神的大能.
神跡獨特之處.
22.
耶穌沒有到現場.
耶穌沒有按手在這個病人那裡.
耶穌沒有施行任何儀式.
就是說他遠離這個病人.
沒有接觸這個病人.
又不用用什麼法事來醫好這個人.
他用什麼來醫好他.
你的兒子活了這句話.
大家記不記得我曾經說過.
耶穌基督平靜.
大風大浪的時候.
他都說了一句話.
在馬可四章三十九節.
住了吧.
靜了吧.
風就住了.
不再吹了.
浪就靜了.
主的話是真實出現.
主的話的出現.
這句話本身也有意義.
耶穌說這個兒子活了.
對於這個大神是多麼大的.
可能他人生最大最寶貴的.
就是這個兒子.
因為這個兒子會承繼他.
繼續做大神.
所以是很重要的.
對於他.

$^{401}$十九世紀.
英國的孤兒院創辦人.
叫做喬治莫勒.
有一天他的孤兒院沒有食物.
數百名的孤兒.
沒有東西吃.
當時的莫勒.
他很相信神會照顧他.
於是他聚集所有的孤兒.
他說各位我們要來禱告.
我以為他的祈禱是求耶穌.
你給我們有食物.
不是.
他說什麼.
主啊.
他很相信主耶穌一定會給.
祈禱一言不擇.
神跡就出現.
有一輛牛奶車去到他們的門口.
突然之間故障.
不能動.
去不到別處.
正好他們需要牛奶.
另外有一輛麵包車.
竟然不知道有誰叫他.
他們有感動送些食物來.
這個是不是神跡.
神在工作.
各位我們做基督徒的.
你信神能不能信到這樣.
未有神跡給你看到.
你已經相信.
神是會幫你的.
好了.
第三.
耶穌說我是生命.
耶穌的生命說的不是.
好像我們地上這種生命.
耶穌說的生命.
是要似一個永恆的生命.

$^{441}$不是我們在世這段時間的生命.
還有耶穌說祂就是生命的意思.
就是祂是掌管生命.
和可以賜你生命的那一位主.
耶穌生命 the life.
人的生命.
人.
人有生理的生命.
也都有屬靈的生命.
你有沒有.
我有.
有限的生命.
你看到的.
我三十年之後.
我現在七十歲.
三十年之後.
你未必見到我的.
如果我還在這裡.
有限.
但我有屬靈的生命.
每個人都有的.
永恆的生命.
主耶穌給我們的生命.
第一祂會赦免人的罪.
令人和神和好.
第二會賜給人有平安的生命.
你在世其實最重要.
是不是平安各位.
最實際是平安.
第二祂應許給你永遠的生命.
意思是勝過死亡的生命.
耶穌是生命.
還有一個意思.
令失喪的人有活力.
有生命力.
有盼望.
信耶穌的人的生命.
永遠活在希望中.
誰是我們的希望.
基督徒說.

$^{481}$這麼小聲.
耶穌或神是我們的希望.
而且我們這個希望是美好的希望.
不過最好的神蹟.
最好的神蹟.
都不是這樣.
最好的神蹟不是這個人病好了.
耶穌醫好他.
甚至他死過復活.
最好的神蹟是.
這個大神.
他發現原來.
耶穌在他回到家之後.
他的僕人趕著來通知他好消息.
他的兒子沒事.
我看看時間.
還有.
他的兒子沒事.
他回到家.
就問家人.
你們是什麼時候.
我的兒子你們看到他復活了.
他痊癒了.
他說就在昨天下午一點鐘.
兒子痊癒了.
這個大神.
他們都有錶.
不過不是我們這些錶.
他計算一下時間.
這句話不是耶穌叫他.
你快點回去.
你的兒子康復了.
在那個時間裡面.
你會看到耶穌基督不單止有大能.
他甚至掌握時間.
神是掌握時間的神.
神掌握時間.
是不是都可以掌握你和我的生命.
他想我許開明你都幾乖.
一百歲給你.

$^{521}$或者許開明你一般般.
再留下來都沒什麼發揮.
好吧.
我現在七十歲.
我不想太早.
耶穌或者另外有時間.
我們一個很出名的宣教士.
英國瑪麗遜.
牧師Robert Morrison.
他在Scotland那邊去中國.
一百零七年他去的時候.
孤身一人.
沒教會沒童工沒警察會.
經濟的支持.
警察會對他說.
你去中國我們沒錢給你.
因為我們警察會都沒錢.
你只是靠上面.
我們可以為你禱告為你求錢.
但沒錢.
當他來到中國的時候.
有一個西人就問他.
中國人會不會因為你一個人.
所存的道就改變.
你是什麼水.
你又沒童工沒教會沒經濟支持.
但他說我不能.
但是上帝能.
各位二十多年之後.
中文聖經完成.
瑪麗遜打開了福音的大門.
數以千計的傳教士.
來到我們中國.
福音一路一路發展.
沒停的傳福音.
基督徒是數以千萬甚至.
過億在中國.
瑪麗遜沒看到這一切.
還有瑪麗遜帶來了外國.
第一是教育.

$^{561}$英華書院第一間的中學.
第一間的教育機構.
他幫很多在澳門的時候.
幫很多有牙周病和白內障.
我們中國人這兩種東西.
是搞得很痛苦.
白內障和牙周病.
他帶來了醫術.
後來的傳福音.
不單止說耶穌是神.
神的道.
還有整個社會都因為.
基督教的離道.
民主的精神.
自由的精神.
全部都是基督教.
由英國美國傳到我們中國.
帶來了中國的現代化.
貢獻是非常大的.
瑪麗遜看不到.
我們看到.
原來神的大能這麼厲害.
耶穌是道路.
二十八頁.
耶穌是真理.
耶穌是生命.
耶穌是道路.
就是說沒有路走的時候.
你找到出路.
在混亂的裡面.
你不知道靠誰.
可以靠耶穌.
耶穌是生命.
在絕望裡面.
讓你有盼望.
甚至起死回生.
你不信耶穌是不是很不著數.
獲得人生的三寶.
我們要知道怎麼樣呢.
第一.

$^{601}$唯有神擁有這三寶.
耶穌說祂是神.
祂就有這三個寶.
神為什麼要給我們這三個寶.
這三個寶其實是耶穌未到世界.
已經有猶太的拉比.
已經說出這三個寶是來自神的.
是基於神對人的愛.
沒有這三寶我們做人都會出現問題.
第二,第三.
人要和神有聯繫.
你信靠神,你信靠耶穌.
這三寶才能夠進到你生命的裡面.
我們再清晰一點說這三個寶.
耶穌要顯出猶太拉比傳統的智慧.
人生有三寶就活得好.
一寶就是指人擁有道路.
你做人有方向有目標.
遇到困難你有出路.
人懂得怎樣走出困局.
第二寶真理.
人擁有真理就懂得明辨.
什麼是有價值.
什麼是對的.
可以有分辨的能力.
做出有意義人生的抉擇.
有生命就有永活的生命.
這個生命是代表有生命力.
有活力.
有動力.
有能力.
人懂得依靠神.
他就能夠活出一個這麼美好的生命.
正不正啊.
就洗禮那幾位.
聽得耐會想打瞌睡.
因為周公很喜歡在這些時間找我們.
最近太古城母女先後跳樓.
這個悲劇我不知道有沒有人知道.
媽媽先死.

$^{641}$跟著就跟女兒.
七個小時之後女兒.
聽到很傷心.
我請教一些自殺的專家.
因為我在中大.
2016年我在中大有11個學生跳樓.
所以我就跟一些專家去挽救一些準備跳樓的.
我負責那間明華堂.
Ava好像去過.
明華堂聽聞有個學生想跳樓.
所以我要找一些專家去幫幫忙.
我又沒試過跳樓.
又不知道怎麼處理.
通常是怎樣呢.
第一個原因.
是這個人為什麼會跳樓.
因為他沒有出路.
沒有出路的時間.
魔鬼就會跳樓.
你死吧,沒用的了.
是解決不了的了.
如果你沒有信仰.
你在這段時間一定死的.
第二,沒有冷靜的去判斷.
可能有兩個聲音.
一個聲音說你跳下去吧.
你死吧.
一個就說還是不要住了.
有拉力的.
但到最後你搞到太辛苦了.
太痛苦了.
就因為太痛苦了.
不如死了.
死了就不痛苦了.
死了你可能不痛苦.
你周圍的人是不是痛苦.
可憐.
所以我猜丈夫很痛苦.
太太跟女兒一家人沒了兩個.
第三是生命.

$^{681}$人擁有生命力.
很多時候自殺的人.
是他連活著的能力都沒有.
媽媽死了.
怎麼辦呢.
是我害死媽媽的.
不是我.
我說的那個人.
他越想越不對.
不如我陪他吧.
那就跳下去吧.
魔鬼最喜歡的.
你們團聚.
所以人是沒有道路.
沒有真理.
沒有生命.
你不能做一個成功.
或者健康.
有活力的人.
各位.
我沒有了神.
我都沒有今天.
大神所信的各位是誰.
原先這個大神.
我相信他只是信耶穌是先知.
或者術士.
能夠行神蹟.
多謝.
但耶穌是怎樣.
耶穌不單止能夠行神蹟.
耶穌是顯示祂是誰.
祂是神.
不是術士先知做得到的.
第一.
耶穌自己說.
我是道路.
我是真理.
我是生命.
「我是」這個字.
原文就是「我是事」.

$^{721}$是神的自稱.
Ego-Amy.
我以前說過了.
不再詳細說.
神自己稱呼自己.
Ego-Amy.
祂就是神.
至於「成」.
我能夠給你道路.
真理,生命.
這個「至成」也是一個.
祂自己的承認.
或者宣告.
這三樣東西.
都是唯有神才擁有.
和掌握得好的.
神蹟.
耶穌基督說.
風浪靜止罷就靜止.
耶穌說.
你的兒子痊癒了.
祂的兒子就痊癒了.
以上是一個神學論證.
耶穌是誰.
是真神.
不是假的.
是真的.
大神不斷向前的信心.
祂又信靠耶穌.
可能初期信祂是先知.
後來知道是神的兒子.
祂又信靠.
變成信服.
耶穌叫祂回去.
祂有沒有回去.
祂不回去的話.
祂看不到這個神蹟.
可能.
祂肯.
由小信到深信.

$^{761}$原來.
這個大神.
祂信了耶穌.
就回去了.
那個信.
原文是一個淺信的字.
到後來.
祂看到祂的兒子.
真是復活了.
回到家裡.
看到真是復活了.
又聽到.
原來.
昨天.
兒子已經復活了.
已經醫好了.
祂一再這樣看到.
這個耶穌真是神.
祂的信.
是用確信的原文.
所以各位.
多讀聖經.
最好讀原文.
沒有原文都看聖經書.
幫你.
由小信到確信.
第三.
從客觀道理上的相信.
耶穌叫祂回去.
祂回去是不是好過不回去.
可能耶穌真是可以.
不過聖經說.
祂都是有信祂的.
但這個是客觀道理上.
耶穌這樣說.
那我都信祂了.
即是有人問你.
你信不信星球人.
我會答他.
我都信的.

$^{801}$既然地球有愛.
其他星球可能都有.
有些又說.
聖經又說.
有些沒有說.
但是我不是確信有外星人.
對不對.
不是確信的.
因為我都沒見過.
是不是.
所以大家要明.
祂由客觀真理上的相信.
去到主觀經歷上的相信.
弟兄姊妹.
我相信你信耶穌.
你一定有主觀經歷過.
耶穌真是一個神幫過你.
有的舉手.
我有了.
很多位.
由個人的信.
這個大神叫全家人信.
全家人.
我查過一些解釋.
不單止他家人.
還有他的兄弟姊妹.
這個大神.
還有他的僕人.
所有的僕人.
都是他的家人.
他一定不少.
還有甚至他部下.
這個大神都信了神.
他的部下.
他的軍隊都可能信了.
影響是不是很大.
生命的改變.
弟兄姊妹.
你信了主.
有沒有去帶你的家人.

$^{841}$或者你的朋友都信.
皇室大神.
因為信他靈命.
是一路來到去增長.
他相信耶穌之後.
他聽從耶穌.
接著大神全家都信.
他們信了之後.
他們當時沒有禮拜堂.
但是他們一定要去會堂.
早期基督教和猶太教是一起.
他們要去會堂.
敬拜主.
或者去祈禱的地點.
去敬拜神.
我的見證.
這裡不說這麼多.
我的見證.
三十四.
三十二還是三十三.
我的見證.
我今天不說了.
因為時間所限.
我相信耶穌.
我是聽從耶穌的.
我又親近耶穌的.
好了.
三十五.
耶穌給我有三寶.
叫我活得好.
我就不說了.
慈父.
最後一點.
說完.
希望給我.
從慈父的愛認識天父的愛.
三十七.
這個大神的故事.
他反映了一個慈父.
愛他這個兒子.

$^{881}$愛到怎樣.
為這個兒子奔走流淚.
去到耶穌的面前.
他是跪下的.
他是一品大觀.
向著耶穌下跪.
於尊降跪.
主耶穌也是.
祂是神.
祂來到這個世界.
為了愛我們.
要救我們.
讓我們生命活得好.
不是單單信祂.
而是你信了祂之後.
你有永恆的生命.
有美好的生命.
為了你.
祂親自來到.
好像這個慈父一樣.
這個慈父.
三十九頁.
祂不放棄.
不放棄.
不斷要去拯救人.
幫助人.
好像我們的神一樣.
祂從不放棄.
要救我們.
然後引領我們去認識耶穌.
認識了耶穌之後.
你就能夠認識到神了.
好像這個慈父.
帶領他家人信主.
好了.
我們最後.
有三個對大家的願望.
也是我今天說的重點.
一.
我們一起說.

$^{921}$一.
願我們憑信見神蹟.
從神蹟認識主耶穌是神.
大家留意.
大神在信心的功課上.
祂是不斷地求進步的.
今天很多基督徒.
你信了耶穌幾十年.
其實你有沒有進步.
你問一下你自己.
你有沒有像這個大神.
這麼有信心.
由小信到深信.
由客觀的信.
或者口頭的信.
而去到主觀.
真的從心裡面發出.
我信耶穌的.
你有沒有感受到.
你信了祂之後.
神怎樣愛你.
保守你.
由香港來到英國.
神怎樣去幫助你.
帶領你的生命.
各位今天第一次來的來賓.
我猜到不少.
我都希望你們能夠相信耶穌.
我們的耶穌是神.
祂愛我們.
為我們死.
祂來到世界受盡苦楚.
就是希望我們能夠認識祂.
得到一個美好的人生.
今天時間關係.
我不用你表示.
但如果你真的信神的話.
希望你繼續回我們的教會.
或者去你附近的華人教會.
認識神.

$^{961}$求主這樣幫助我們.
阿們.
\newpage



\section{士師記 11:30-31}
\label{sec:032355}
\textbf{聖經人聖經事:耶弗他以他女兒許願}~~~~黃克勤牧師博士
\newline
\newline
連結: \href{http://www.church.com.hk/acms/content.asp?site=cdc\&op=show\&type=product\&code=032355\&layout=sermon}{\texttt{032355}} ~~~~ 語音日期: 2026-06-27
\newline
\newline
\hyperref[sec:032354]{\small{< < < PREV SERMON < < <}}
~
\hyperref[sec:index_chronic]{\small{[返順時目]}}
~
\hyperref[sec:032356]{\small{> > > NEXT SERMON > > >}}
\newline
\newline
士師記 11:30-31
\newline
\begin{longtable}{cl}
\hline
\hline
章節 & 經文 (和合本修訂版)\\
\hline
11:30 & \begin{tabularx}{0.7\textwidth}{X} 耶弗他向耶和華許願，說：「你若真的將亞捫人交在我手中， \end{tabularx} \\ \\ \relax
11:31 & \begin{tabularx}{0.7\textwidth}{X} 我從亞捫人那裡平平安安回來的時候，無論誰先從我家門出來迎接我，就要歸給耶和華，我必將他獻上作為燔祭。」 \end{tabularx} \\ \\
[1ex]
\hline
\hline
\end{longtable}
$^{1}$各位電影姐妹大家好.
我們又一起讀聖經了.
又是《生命說書教》的網絡.
跟大家一起看聖經人.
看聖經的事件.
看看上帝在背後的心意.
那麼繼續講時詩記的耶福他時詩.
那麼這次呢.
哎呀.
真的很開心的.
打贏仗.
不過又一時口快.
又失誤.
不是失誤的.
我想不是耶福他特別想的.
不過就是一時很開心.
那麼就許了一個大願.
結果就令到他真的抱憾了.
好,看看耶福他以他女兒許願.
就是今天的題目.
其實耶福他向上帝許願.
就不是用他的女兒來許願的.
不過呢.
我是將整件事情用一句話來講.
是耶福他用了他的女兒來許願.
其實是耶福他許願.
不過結果就是關乎他的女兒.
他最疼愛的獨女.
好,看看.
時詩記十一章.
這次我就會這樣的.
就今天開始試一下.
就是改一改形式.
我會整個故事講完.
才去到最後有三點跟大家思想的.
就不會在中間開始思想了.
耶福他就向耶和華許願.
那麼你看到這個黃色的字.
就是螢光了.
就是最重要的地方了.

$^{41}$耶福他向耶和華許願就說了.
上個星期讀過這段聖經的了.
說.
「你若將阿們人交在我手中.
我從阿們人那裡平平安安回來的時候.
無論什麼人.
先從我家門出來迎接我.
就必歸你.
我也必將他獻上為梵祭」.
這個上一次就昨天已經講過了.
不過今天我再用它.
將今天的故事聯繫有關係.
耶福他向神許願.
起誓,立願.
「如果我贏了阿們人.
而且平安回來.
我是會回家的.
我不知道的.
總之家裡第一個跑出來迎接我的.
我就將他歸給你.
我會叫他獻上為梵祭」.
我們上次解釋過了.
將一個人獻上為梵祭.
剛剛跑出來是一隻羊也沒問題.
跑出來是一隻牛也沒問題.
跑出來是一個人來迎接我.
怎麼會做梵祭呢?.
耶福他也知道不會的.
是因為上帝早就已經非常反對用人為祭.
不過這是一個說法.
有時候我們也會說.
「主啊,我獻上我自己做活祭」.
真的用火燒死自己嗎?不是的.
梵祭的意思是.
「Ola」.
「Ola」是火祭或者升高.
將自己交給主.
遞上去的意思.
於是耶福他就往阿們人那裡去了.
與他們爭戰.

$^{81}$耶和華說將他們交在他手中.
就是耶福他打贏了.
他就大大殺敗他們.
從阿羅爾到米尼.
直到阿彼納.
就是基拉明.
這些名字在哪裡說的.
什麼意思不說了.
就攻取了二十座城.
也不是說不厲害.
而是神特別的手幫助他.
因為你要攻城.
其實是很難的.
而且很危險.
而且很多時候都會損傷很大.
攻城,不過呢.
攻了二十座城.
這樣阿們人就被以色列人制服.
本來是帶著大軍.
現在我們收回我們的地.
就是以前你們.
以色列人出埃及的時候.
經過我們以東摩押阿們.
拿了我們的地.
做了你們支派的分地.
現在我們拿回.
你弱,我強.
我就兵臨城下.
耶福他就問他.
你為什麼來打我們.
是因為我想拿回我的地.
如果你平平安安.
交回給我.
那沒事了.
大家各得其所.
不過.
如果你不肯我就打你.
結果呢.
就被耶福他打勝了.
耶福他就回米斯巴度了自己的家.

$^{121}$不料.
哎呀.
沒想到.
是他女兒拿著鼓跳舞出來迎接他.
哎呀.
是他獨生的.
此外無兒無女.
你知道只有一個女兒.
哎呀.
又大個又會跳舞.
又獨生女兒.
第一個走出來迎接他.
這樣真的很好.
《士師記》是薩姆爾寫的.
薩姆爾真的厲害.
耶福他回米斯巴度自己的家.
是不料的.
沒想到是誰出來迎接他.
是他的女兒來迎接他.
他的女兒就要獻給神了.
因為許了願要還願的.
那麼.
耶福他看見她.
就撕裂衣服說.
哀哉.
慘了我的女兒.
你使我甚是愁苦.
叫我作難了.
因為我已經向耶和華開口許願.
不能挽回了.
就是將你獻給神了.
你就要屬於他了.
這樣子.
那麼留意一下.
耶福他真的很慘.
哀哉.
慘了.
哀哉真的慘了.
我的女兒啊.
我真的好痛苦啊.

$^{161}$這個愁苦為止.
我真的好痛苦啊.
好困難啊.
為什麼呢?.
我現在左右做人難了.
我已經向耶和華開口許願.
是不能夠挽回的了.
黃色自然螢光.
他的女兒厲害.
他女兒回答說.
父啊.
你既向耶和華開口.
就當照你口中所說的.
向我行.
他的女兒說.
哎呀,爸爸,為什麼你這樣亂許願?.
你害死我了.
沒有啊.
他的女兒真的厲害.
不過你要知道.
在整件事情上.
聖經是沒有記載他的女兒的名字的.
這個真的是無名的.
很靈性的女孩子.
哇,真的厲害.
你照著口中所說的.
向我行就可以了.
為什麼呢?.
因為耶和華已經在受敵阿們人身上.
為你報仇了.
其實就是說.
你帶領我們以色列人向整個阿們報仇.
所以我現在獻給神.
很值得.
你而且開口許了願.
你照著做就可以了.
哇,他的女兒真的厲害.
不過這樣.
又對父親說.
不如這樣吧.

$^{201}$我也想求你一件事.
求你給我.
容我去兩個月.
與同伴在山上.
好哀哭.
我終為處女.
就是說.
其實獻給主.
主要就是不嫁人而已.
或者應該是事奉的.
做一些事奉的事情.
你知道以往比較少的獻祭等等各樣事奉.
或者是那些靈性上回望的事奉.
比較少是女性做的.
都是男性的.
如果女的去做.
你不會做女祭司的.
不可能的.
所以頂多都是做一些打點的事情.
也有一些帶領.
耶佛他女兒來幫忙.
當然是比較好一點.
不過他去了服侍神.
就好像今天.
女兒許願就是.
如果我這次打聖仗.
如果我這個病好.
我就想我女兒給你做傳道人.
很多人是這樣的.
很好笑的.
於是.
竟然痊癒了.
不過這裡有個特別.
就是爸爸我想去兩個月上山.
和我的朋友一起.
在那裡哭一哭.
我要做處女不可以出嫁了.
這麼有趣.
你說讓我去山上兩個月.
讓我開心一下.

$^{241}$Picnic一下.
然後游山玩水兩個月.
開心一點.
然後我就去服侍神.
這樣也對一點.
他說我哭兩個月.
其實.
就.
因為這種做法.
在《成本聖經》只有這一次.
這個上山去兩個月哭一哭.
這件事情.
沒有其他參考.
在山上有沒有儀式做.
怎樣做.
怎樣哭.
不知道.
還是爸爸讓我上山哭兩個月.
其實是上山去玩的.
我不知道.
因為下山之後就去事奉了.
所以先玩兩個月.
不知道.
這是我自己想的.
哭兩個月.
耶乎他就好了.
你去吧.
就容他去兩個月.
他便和同伴去了.
在山上為他終為處女哀哭.
這樣就很慘了.
兩月已滿.
他回到父親那裡.
父親便照所許的願向他行了.
即是行什麼呢?.
沒有什麼要做的.
沒有特別的儀式.
很有可能就是.
帶他去回眸那裡事奉.
就是這樣.

$^{281}$有一點像.
哈娜.
就是撒母耳的媽媽.
許願.
如果我有一個兒子.
我就叫他終身服侍耶和華.
所以他有一個兒子終身服侍耶和華.
撒母耳出世.
他便放在祭司爾尼手下.
祭司你和我帶他去終身服侍.
還願.
不過應該照計在回眸.
我推測.
我不知道是不是回眸.
那怎樣.
屬於耶和華就要終身不嫁.
我學識少.
我沒有這個印象要這樣做.
為什麼要這樣呢?.
終身沒有親近男子.
嗯.
好像不是這樣的.
有這樣的規例.
不過他大概是去了服侍.
那些去回眸服侍的都是沒有結婚的.
以後以色列終有一個規矩.
因為這件事.
每年以色列的女子.
都會為羈烈人耶佛他的女兒.
你看到這裡都沒有名字.
是哀哭四天.
那四天怎樣?整天哭嗎?.
即是說有四天是.
憂愁一點.
悲哀一點.
不要那麼多自助餐.
不要那麼多派對.
不要那麼多去玩.
不要那麼多笑.
不要那麼多歡樂.

$^{321}$是為了他的女兒.
來到有四天的安靜哀哭.
就是這樣.
這就是今天要看的.
好了.
這裡就開始和大家看看.
一是這個藍色字有幾點要講的.
第一.
耶佛他向耶和華許願.
我想大家參考傳道書第五章.
就是說你向神許願.
許願是什麼?你留意一下.
向神許願.
許了願.
是要償還不可辭賢.
傳道書說.
因為他不起願愚昧人.
所以你許的願應當償還.
你許願不還.
就不如不許了.
好.
講一下這裡.
既然耶佛他向耶和華許願.
其實.
是不是一定要許願呢?.
譬如耶佛他如果要打仗了.
他的不許願.
他只是禱告神.
求你幫助我.
能夠戰勝.
我就交託給你了.
在祈禱上不是許願.
將來你帶領我一家人.
如果我死了.
你就保守我家人.
之類各樣.
你可以在祈禱交託.
而不是許願.
主要是如果我能夠平安生還回來.
打勝仗.

$^{361}$我的女兒做傳道人.
其實這個願.
我個人認為耶佛他.
是起得有少少衝動了.
不管了.
總之第一個出來迎接我.
我就歸你了.
事奉你的.
你要知道.
這個願是你起的.
當然我作為耶佛他一家之主.
整個以色列的主.
我說什麼都可以.
但是.
其實在我們今天來說.
有些不太舒服的原因就是.
這個願不是我個人的事.
是和第一個走出來的人有關.
我是將他的生命.
做了一個安排.
而沒有問過他.
當然在當日是需要問的.
我叫我女兒嫁給誰.
嫁給誰.
不嫁不嫁.
誰離開香港.
我提親我不給.
都可以的.
我有權的.
是有權.
所有家庭.
因為其實耶佛他.
應該在當時他全家人當中.
他最大有權處理任何人.
怎樣生活法都可以.
所以他說.
神啊.
我去這個願.
誰第一個跑出來迎接我.
我就叫他歸你.

$^{401}$可以.
不過.
我覺得如果他不去這個願.
他應該會開心很多.
或者之後才再想.
現在回來了.
神就聽我祈禱.
女兒不如你服侍神.
或者怎樣.
或者我做什麼.
可能更好一點.
但他既然開口了.
女兒真厲害.
既然開口了.
許了願就要償還.
這個是以色列當中眾所周知.
而且還要快些還.
回到去.
女兒來的.
我將她都獻給神.
不過遲些遲些.
等她八十歲先.
不行.
傳道師說你許了願.
償還不可遲然.
不要假裝.
不要彈弓手.
等她遲些 長大些.
現在還小不行.
其實你這樣.
變成一個愚昧人.
你當上帝是誰.
當他傻.
你才傻.
其實他不不起悅愚昧人.
意思就是說.
那些遲然的人以為自己很厲害.
其實很蠢.
你許的願應該償還.
如果不是神是對你不客氣的.

$^{441}$如果你不想還.
不如不許.
就不要許了.
第六誌.
不可任你的口使肉體犯罪.
也不可在祭司面前說時.
哎呀 sorry.
錯許了.
為何使神因你的聲音發怒.
敗壞你所做的呢.
因此.
這裡是鼓勵大家記得.
今天我們都會的.
有時遇到疾病.
大件事.
兒子病得很嚴重.
我有個關過不了.
我有一件官司.
我女兒怎樣怎樣.
我兒子怎樣怎樣.
要讀哪間學校.
哎呀.
我想移民去哪裡.
很大件事.
如果我媽又入院了.
如果康復了.
我就怎樣怎樣.
我鼓勵你.
就祈禱好了.
不要隨便.
許願.
如果你還不了.
就不如不許.
又不能在祭司面前跟神說.
對不起.
一時衝動起錯了.
不可以的.
跟神說了.
是就是了.
這樣不跟著做.

$^{481}$是神會因為你所起的聲音發怒.
敗壞你所做的工作.
即是說.
祂不祝福你.
而且祂不成就你的工作.
祂還敗你的工作.
你更慘了.
不要任你的口使肉體犯罪.
肉體犯罪的意思不是犯什麼罪.
而是不要任你的口亂說話.
因為你這樣亂說話.
你的手所做的一定不祝福.
就是這樣.
你也不能夠說.
哎呀神啊對不起.
那天喝多了.
一時衝動.
那天怕得太多.
一時衝動.
不算不算.
就抹掉祂寫字.
粉筆字抹掉算了.
不行.
那你記得.
神是又真又活.
祂聽到.
祂是算數的.
聖經這件事屢試不爽.
當那些法利賽人.
當那些猶太人要殺死耶穌.
釘十字架給得多.
留無辜人的血.
那怎樣?歸誰?.
歸到我們.
亡我們子孫身上.
神說好.
收到聽到.
中以色列人一直到猶太人.
是非常慘的.
二次大戰已經.

$^{521}$只是這樣一次在德國已經死六百萬人.
其實不計其數.
歷代除了流放不同的國家.
周圍跑之外.
只是那一次都死那麼多人.
那你見到上帝.
德歸到你的身上.
你留些血給我看看.
第二點要說的.
真的很欣賞他的女兒.
真的想知道她什麼名字.
將來在天上應該見到.
他女兒回答說.
父啊,你既向耶和華開口.
就當照你口中所說的向我行.
嘩,這個女兒真是多好.
不知是否自小已經受很好的宗教教養.
靈性很好.
因為耶和華已經做了她的一部分.
為我們報了仇.
因此你照著做吧.
他的女兒反而很信服.
很棒.
希望我們都學到他的女兒這種心態.
向神許願要還的.
什麼跟我有關?.
既然父親是這樣.
我們信服他的女兒.
是一個信服的女兒.
是很蒙恩的.
雖然我們不知道她的名字.
第三,兩個月已滿.
她回到父親那裡.
照著她的御所行.
終身沒有親近男子.
於是以後有個規矩.
哭四個月,三哭四日.
大家要留意.
我想說的是.
這個女兒終身不嫁.

$^{561}$那很慘.
上山哭兩個月.
然後以後所有以色列婦女.
都每年為她哭四天.
為什麼呢?.
不過可能當日和今天.
真的有很多文化上的差異.
如果今天沒有結婚.
那就沒有了.
更加好了.
自由自在了.
有結婚好嘛.
是好的.
沒有結婚也不差.
今天就不要說.
我很慘.
因為我還沒有男朋友.
結不了婚.
很慘.
怎麼慘呢?.
不是很慘的.
有的那些我見到很多.
有男朋友結了婚之後.
可能比沒有更慘.
其實不知道的.
如果真是有心服侍神的.
沒有結婚.
就沒有了.
沒有男孩子.
沒有了,不要緊.
有也未必好.
千萬不要.
人有我有.
就一言連晚煎堆.
大家都要一個.
這樣不行的.
這樣其實不好的.
也不需要這麼哀哭.
不如將它轉化成為一個.
寬闊一點的胸襟.

$^{601}$神帶領我這條路.
竟然這樣安排.
這樣就好了.
我可以走一條更加寬闊.
適逢的路.
去開心地迎接這個新的轉變.
或者在這裡下去更加好一點.
亦無需我見到很多.
獨身的姊妹.
因為其實今天.
很多姊妹是很厲害的.
很厲害.
無論是讀書的背景.
教育背景.
工作能力.
領導才能.
口才.
各樣東西.
收入都比很多男士一點都不差.
因此其實她都會想選一個.
好一點.
配到一點的男士.
但是很多時候未必找到.
找了久了.
年紀又大了.
每個人都問.
「哎呀,你還未結婚呀?」.
很尷尬.
於是她就使出一套絕招.
就是說.
「神帶領我是獨身的」.
「獨什麼身?都沒有這樣的帶領」.
強行說自己有這樣的帶領.
於是到最後怎樣呢?.
就獨身久了之後.
又說自己有負擔獨身.
又說神有帶領我獨身.
又說我有獨身恩賜.
誰知道久了之後.
又見到一個很心儀的男生.

$^{641}$大家又情投意合.
「糟糕了,怎麼辦?」.
以前說自己有獨身恩賜.
又說有帶領.
所以其實不要隨便說.
不用這麼快跟自己下定論.
也無需許願.
「如果我有一個男生認識我怎麼辦?」.
不要隨便許願.
而且如果有獨身的.
不獨身的不要緊.
總之都是開心地事奉.
開心地生活.
這是最重要的.
叫神的心要滿足.
不一定要哀哭.
今天跟大家分享到這裡.
我們一起禱告.
多謝主我們學到功課.
耶復他很勇敢.
他拼命戰勝敵人.
也在神面前誠心許願.
最令我們感動的地方.
除了他倚靠你而成長之後.
我們認為他有一些衝動的許願.
但是他也沒有反悔.
他仍然尊重.
既然向神開口了.
他就走到尾.
雖然自己覺得痛苦不堪.
不想見到女兒.
或者覺得害了女兒不能結婚.
但是他也唯有走下去.
他的女兒更叫我們覺得是欣賞.
既然爸爸開口了.
我就信服了.
他也可能不開心一段時間.
哭了一兩個月.
他就很樂意地走上這條事奉的路.
願主親自保守祝福.

$^{681}$讓我們同樣有信服你的心.
知道你又真又活.
我們在你面前說了要算數.
又知道你又真又活.
你必定帶領合你心意的人.
多謝主將耶復他女兒成為我們的榜樣.
對爸爸,對你是信服的.
求主祝福我們面前的生活.
每一個時刻都抬頭仰望你.
又真又活的主.
知道你是會帶領的.
感恩祈禱.
奉耶穌名求.
阿們.
\newpage



\section{約書亞記 24:14-28}
\label{sec:032356}
\textbf{至於我和我家,我們必定事奉耶和華}~~~~梁成裕牧師
\newline
\newline
連結: \href{http://www.church.com.hk/acms/content.asp?site=cdc\&op=show\&type=product\&code=032356\&layout=sermon}{\texttt{032356}} ~~~~ 語音日期: 2026-06-21
\newline
\newline
\hyperref[sec:032355]{\small{< < < PREV SERMON < < <}}
~
\hyperref[sec:index_chronic]{\small{[返順時目]}}
~
\hyperref[sec:032357]{\small{> > > NEXT SERMON > > >}}
\newline
\newline
約書亞記 24:14-28
\newline
\begin{longtable}{cl}
\hline
\hline
章節 & 經文 (和合本修訂版)\\
\hline
24:14 & \begin{tabularx}{0.7\textwidth}{X} 「現在你們要敬畏耶和華，誠心誠意事奉他，除掉你們列祖在大河那邊和在埃及事奉的神明，事奉耶和華。 \end{tabularx} \\ \\ \relax
24:15 & \begin{tabularx}{0.7\textwidth}{X} 若你們認為事奉耶和華不好，今日就可以選擇所要事奉的：是你們列祖在大河那邊所事奉的神明，或是你們所住這地亞摩利人的神明呢？至於我和我家，我們必定事奉耶和華。」 \end{tabularx} \\ \\ \relax
24:16 & \begin{tabularx}{0.7\textwidth}{X} 百姓回答說：「我們絕不離棄耶和華去事奉別神。 \end{tabularx} \\ \\ \relax
24:17 & \begin{tabularx}{0.7\textwidth}{X} 因為耶和華－我們的神曾領我們和我們的祖宗從埃及地為奴之家出來，在我們眼前行了那些大神蹟，並在我們所行的一切路上，和所經過的各民族中保護了我們。 \end{tabularx} \\ \\ \relax
24:18 & \begin{tabularx}{0.7\textwidth}{X} 耶和華又把各民族和住此地的亞摩利人都從我們面前趕出去。所以，我們也必事奉耶和華，因為他是我們的神。」 \end{tabularx} \\ \\ \relax
24:19 & \begin{tabularx}{0.7\textwidth}{X} 約書亞對百姓說：「你們不能事奉耶和華，因為他是神聖的神，是忌邪的神，必不赦免你們的過犯罪惡。 \end{tabularx} \\ \\ \relax
24:20 & \begin{tabularx}{0.7\textwidth}{X} 你們若離棄耶和華去事奉外邦的神明，耶和華在降福之後，必轉而降禍給你們，把你們滅絕。」 \end{tabularx} \\ \\ \relax
24:21 & \begin{tabularx}{0.7\textwidth}{X} 百姓對約書亞說：「不，我們要事奉耶和華。」 \end{tabularx} \\ \\ \relax
24:22 & \begin{tabularx}{0.7\textwidth}{X} 約書亞對百姓說：「你們選擇耶和華，要事奉他，你們自己作證吧！」他們說：「我們願意作證。」 \end{tabularx} \\ \\ \relax
24:23 & \begin{tabularx}{0.7\textwidth}{X} 「現在，你們要除掉你們中間外邦的神明，專心歸向耶和華---以色列的神。」 \end{tabularx} \\ \\ \relax
24:24 & \begin{tabularx}{0.7\textwidth}{X} 百姓對約書亞說：「我們必事奉耶和華---我們的神，聽從他的話。」 \end{tabularx} \\ \\ \relax
24:25 & \begin{tabularx}{0.7\textwidth}{X} 那日，約書亞就與百姓立約，在示劍為他們制定律例典章。 \end{tabularx} \\ \\ \relax
24:26 & \begin{tabularx}{0.7\textwidth}{X} 約書亞把這些話寫在神的律法書上，又拿一塊大石頭立在橡樹下耶和華聖所的旁邊。 \end{tabularx} \\ \\ \relax
24:27 & \begin{tabularx}{0.7\textwidth}{X} 約書亞對眾百姓說：「看哪，這石頭可以向我們作見證，因為它聽見了耶和華所吩咐我們的一切話；這石頭將向你們作見證，免得你們背叛你們的神。」 \end{tabularx} \\ \\ \relax
24:28 & \begin{tabularx}{0.7\textwidth}{X} 於是約書亞解散百姓，各自回到自己的地業去了。 \end{tabularx} \\ \\
[1ex]
\hline
\hline
\end{longtable}
$^{1}$為什麼呢?因為有一次.
我聽過一個這樣的分享.
有一個老師問一個少年人.
究竟父親節和母親節有什麼分別呢?.
那個少年人就說.
沒什麼分別的.
不過父親節的禮物小一點.
所以其實我們真的很需要多謝我們的爸爸.
今晚呢.
有沒有人出去慶祝父親節的?.
有沒有?.
有的,應該有的.
我也很慶祝父親節.
因為我爸爸說想北上吃飯.
所以今晚下午完了青少年崇拜.
我就回北上跟我爸爸吃飯了.
希望我們每一個爸爸.
都有一個很開心的父親節.
其實爸爸是一個很重要的遺訓.
因為我們知道.
首先我們要多謝我們每一個的爸爸.
因為他為我們的家庭去辛勞.
他為我們家庭成為我們家庭的支柱.
有供應我們的需要.
但是今天原來作為爸爸.
你有另外一個很重要的遺訓.
就是原來你是帶領著你的家庭.
以及帶領著你家庭的方向.
所以今天我們特別要看一段經文.
當然這段經文不只是跟我們作為爸爸的去講.
但是今天.
對於我們每一個作為爸爸的.
我們又怎麼帶領我們的家庭.
而今天我們帶領我們的家庭.
去事奉哪一個呢?.
有時候可能作為爸爸已經很辛苦了.
要為家庭去努力.
要為生活去打拼.
但是原來今天爸爸很多時候每一天的選擇.
原來都在影響我們的下一代.

$^{41}$也在帶領著我們的家庭.
去走向一個方向.
在教會裡面.
我其實經常看到爸爸對我們下一代的影響.
今天連續兩天我也看到.
其實在敬拜上面.
有一些父子.
就是有一些父親和有一些兒子.
有時候我會看到.
有時候就是這個兒子打鼓很厲害的.
原來爸爸也會打鼓的.
有時候我又會看到.
爸爸又喜歡敬拜.
兒子又很喜歡敬拜.
同樣有時候.
我會跟一些初職聊天的時候.
還有跟一個初職聊天我就會問.
你不喜歡你會做什麼呢?.
他說我會跟隨爸爸.
去學做一位會計師.
我也在少年部看到有一些男孩子.
煮東西也很厲害的.
我就問.
為什麼你煮東西這麼厲害呢?.
他說因為我爸爸煮東西也很厲害的.
所以原來今天爸爸的生命.
也在影響著我們的下一代.
或者今天如果爸爸很重視成績.
我們的下一代可能就會用成績去定義自己.
今天如果爸爸很重視金錢.
可能我們的下一代就會用收入去衡量價值.
但是今天如果作為爸爸的重視信仰.
也會影響我們的下一代.
如何很真實地去跟隨神.
所以今天原來我們不只是為自己做決定.
因為原來我們的決定也在影響著我們的下一代.
今天作為爸爸真的是辛苦的.
要供樓.
要養家.
希望給最好的家.

$^{81}$這些都是很重要的.
但是今天我們進入一個更加深的問題.
就是爸爸不僅是養家.
原來也是在帶領著這個家的方向.
所以今天在聖經裡面.
約書亞比以色列人一個的挑戰.
就是今天你要選擇所事奉的是哪一個.
而今天約書亞首先在聖經裡面作出了一個宣告.
我們一起讀好不好?.
一二三.
至於我和我家.
我們必定事奉耶和華.
今天我們未必每一個都是爸爸的.
但是今天我們也在影響著我們下一代.
今天我們可能當中有做爸爸的.
我們當中有做媽媽的.
我們當中呢,剛剛Pandy說我們有屬靈的父母的.
但是可能你覺得這些字實在太老了.
叫我是屬靈的父母.
那麼今天你可能是屬靈的哥哥.
屬靈的姐姐.
同樣今天我們每一個都是一個基督的門徒.
我們作為門徒.
原來我們的生命也在影響著我們身邊的人.
影響著我們的下一代.
也在影響著我們教會的下一代.
所以今天原來很重要.
我們不僅是作為爸爸作為媽媽.
或者我們作為屬靈的父母.
或者屬靈的哥哥姐姐.
今天其實我們作為門徒.
原來我們的生命也在影響著我們身邊的人.
影響著我們的下一代.
甚至是教會的下一代.
今天我們又怎麼樣去影響身邊的人.
讓我們的生命一生事奉耶和華呢?.
很想今天去用一段我們可能也熟悉的聖經.
去學習三個很重要的功課.
今天我們第一個功課.
是我們要選擇認定我們所事奉的.

$^{121}$我們首先一起去讀二十四章十四至十八節.
好不好呢?.
一起讀.
一二三.
現在你們要敬畏耶和華.
誠心實意地事奉他.
將你們列祖在大河那邊.
和在埃及所事奉的神除掉.
去事奉耶和華.
若是你們以事奉耶和華為不好.
今天就可以選擇所要事奉的.
是你們列祖在大河那邊所事奉的神呢?.
是你們所住這地的阿摩利亞的神呢?.
至於我和我家.
我們必定事奉耶和華.
百姓回答說.
我們斷不敢離棄耶和華.
去事奉別神.
因耶和華我們的神.
曾將我們和我們列祖.
從埃及地的遺牢之間引出來.
在我們眼前行了那些大神蹟.
在我們所行的道上.
所經過的諸國都保護了我們.
耶和華又把著此地的阿摩利亞.
都從我們面前趕出去.
所以我們必事奉耶和華.
因為祂是我們的神.
當我們看完聖經.
我們從第十四節開始看.
原來第一到第十三節.
是在回顧神拯救以色列人的歷史.
而在這裡約書亞就呼籲以色列人.
你的生命要作出回應.
既然神這樣對待我們.
我們這一刻就要作出生命的回應.
所以在十四節一開始.
這裡說你們要敬畏耶和華.
敬畏就不是恐懼.
敬畏是尊崇.

$^{161}$並且知道神在我們生命中間的主權.
而誠心實意就是說.
我們不是只有一張嘴去說我們事奉神.
而是我們真的整個生命內外一致地去跟隨.
而以色列人要去除掉.
他們在大海那邊和埃及所事奉的神.
我們知道以色列人在他們的歷史裡面.
他們有拜偶像的歷史.
因為他們被身邊不同的群體影響.
但是在這裡.
今天他們要搞清楚.
究竟他們所事奉的是哪一個.
而到了第十五節.
我們看到約書亞用一個反面的問題問.
他說如果你覺得事奉耶和華是不好的話.
你今天就要選擇你所事奉的是哪一個.
如果你不選擇耶和華.
你會選擇什麼呢?.
你可能會選擇列祖在大河那邊所事奉的.
或者你會去選擇阿摩利人的神.
其實今天我們都面對這兩個選擇.
你可能說.
不是的,梁牧師.
我沒有拜過以色列列祖拜過的神.
是什麼來的?.
我又沒有拜過阿摩利的神.
因為其實阿摩利在哪裡?.
地圖在哪裡我也不是很清楚.
但其實這裡是在說兩件事.
原來我們可能在我們的原生家庭.
或者我們文化傳統裡面都拜過民間一些的神.
而那些神的價值觀可能正在影響我們.
同樣什麼叫做阿摩利人的神.
就是現在當今世界流行的價值觀.
今天我們也可能被現在世界流行的價值觀影響我們.
面對我們可能原生家庭.
或者傳統民間宗教的神.
面對現在世界的價值觀.
我們要選擇我們所事奉的是哪一個.
而約書亞選擇了.

$^{201}$他說至於我和我家.
我們必定事奉耶和華.
可能到了這裡你會覺得梁牧師我沒有問題.
是不是?.
我們這裡誰有事奉崗位的?.
你舉舉手.
你有事奉崗位的.
是的,謝謝你.
有事奉崗位的.
剛才你有沒有敬拜?.
幫Peggy跟我們說.
萬有都歸你.
你是永耀,你有沒有唱?.
有,我聽到你唱得很大聲.
今天你說不是.
梁牧師,我又有事奉崗位.
我又有敬拜.
那就是我選擇了.
但其實這裡說的事奉.
是更加深層的事奉.
當約書亞說至於我和我家.
我們必定事奉耶和華.
這個事奉不只是事奉崗位這麼簡單.
這個事奉不只是我們星期日回來敬拜這麼簡單.
這個事奉是你跟隨著哪一個神.
跟隨著哪一個神.
就是你的生命追求著什麼.
你生命所追求的目標是什麼.
你生命追求的方向是什麼.
今天我們可能有事奉.
但今天我追求的可能是名利.
今天我可能有事奉.
但我追求的是金錢.
今天我可以有事奉.
但我追求這個世界的地位.
今天我可能有事奉.
我追求這個世界的享樂.
所以原來事奉不只是事奉崗位這麼簡單.
事奉是說你的生命追求的目標.
你生命追求的方向.

$^{241}$你跟隨著什麼.
而在這裡約書亞選得很清楚.
他說至於我和我家.
我們必定追隨跟隨事奉耶和華.
所以今天我們也要選.
但是以色列人怎麼選呢?.
我們看看十六至十八節.
十六至十八節以色列人在這裡做了三個部分的回應.
第一個部分就是他們宣告他們的生命不離棄耶和華.
而第二.
他們說出他們去事奉耶和華的理由.
這裡有四個部分.
都是分別說過去耶和華在他們生命裡面的恩典.
他們說我們必定事奉耶和華.
第一.
因為耶和華就是領我們從埃及圍爐之地出來的神.
而第二.
就是在這個過程當我們離開埃及進入迦南地的時候.
我們經歷了神的大神蹟.
無論當時過紅海.
在曠野裡面的公營.
躍蕩河的乾地.
我們都經歷了神的大神蹟.
第三個不只是這樣.
他說在我們走的路上.
經過了不同的國家.
耶和華都保護了我們.
而最後當我們進入耶和華為我們預備地的時候.
他又將那個地方的阿摩利人.
從我們前面趕出去.
所以到了第三個部分.
他們再次宣告.
所以我們必事奉耶和華.
因為祂是我們的神.
今天我們也要選擇.
究竟我們的生命追隨的目標是什麼.
我們的生命追隨的方向是什麼.
今天我們究竟跟隨著哪一個神.
今天我們選好了沒有.
你問一問你旁邊的那個.

$^{281}$你選好了沒有.
大家一起笑而已.
選好了沒有.
選好了.
但是我們要小心.
因為我們可能也會被傳統的價值觀影響.
或者我們也會被世界的價值觀影響.
我曾經認識一位肢體.
他父親是典型上一代拼搏的香港人.
我想問哪一個是上一代典型拼搏的香港人.
有的.
辛苦你了.
我不是.
但是我父親是.
我父親是上一代典型拼搏的香港人.
但是有時候我們在這個典型拼搏香港人的環境成長.
我們也會被這個價值觀影響.
所以我認識這位肢體.
他從小到大.
他父親也是這樣教導他的.
要出人頭地.
所以要努力讀書.
努力讀書.
以前他說考第一.
每個人都會很尊崇你.
但是現在我做少年人.
我發現你考第一.
人們就會說你是學霸.
他說我們那個年代要很努力讀書.
然後最好努力賺錢.
努力賺錢之後要做什麼.
真的很好.
四堂大家都能回答我.
努力讀書.
努力賺錢.
然後買房子.
他說他在這個環境裡面成長.
他說原來他發現他信了主之後.
雖然他的頭腦知道要以神為首.
但是他說他的時間表.

$^{321}$他的價值觀仍然是成就至上.
工作至上.
所以有一次他說.
人生裡面有一個很大的掙扎.
就是他的老闆要他進入一些灰色地帶.
有一次老闆就走來跟他說.
我們現在要拿到這個大廈.
所以待會你做那個報告.
你美化一下數據.
每個人都是這麼做的.
你想不想升職?.
想升職就做吧.
在那一刻他說他裡面有一個很大的掙扎.
如果他知道他 say no.
一定沒得升的.
一定不是他.
就是他裡面有一個掙扎.
他說如果我跟老闆 say yes.
他說原來我事奉的不是耶和華.
我事奉的就變成了這個職場文化的神.
所以事奉不只是說一個崗位這麼簡單.
事奉是說我們生命的方向.
他在那一刻有一個掙扎.
我跟老闆 say yes.
我不再事奉耶和華了.
我事奉這個職場文化的神.
但是跟老闆 say no.
你不用自以為升職了.
他說在那一刻他有很大的掙扎.
但是他說最後他決定跟老闆 say no.
他說我要事奉耶和華.
當他跟老闆 say no 之後.
你猜升職的人是不是他?.
當然不是他了.
升職的人是其他人了.
但是他說那一刻他發現他要重整他的生命.
他說以前我是成功至上.
工作至上.
但是他說我重整我的時間表.
我要給回時間我的家.

$^{361}$給回時間我的教會事奉.
然後他說雖然我得不到升職帶來額外的收入.
但是我生命卻有額外的滿足和喜樂.
原來今天我們的決定.
或者我們正在做的事情.
是在影響我們的下一代的.
特別是今天父親節.
我們作為爸爸.
我們正在做的事情是在影響我們下一代的.
有一次我跟少年去唱K.
我經常覺得去唱K也輪不到我.
他認識的也不是我認識的.
突然間看到一個少年人唱.
為何偏偏喜歡你.
他就說哇.
可以唱啊.
拿個麥克風來.
大家如果知道這是誰唱的.
你跟我同一個年代了.
陳百強.
沒錯.
陳百強不是陳百祥.
然後.
我問他.
你為何會認識的.
他說因為我爸爸經常唱.
是啊.
我說你跟你爸爸一起唱.
他說不是啊.
我只是看他唱而已.
原來今天我們正在做的事情.
是在影響我們的下一代.
今天我們怎麼去選擇我們的生命.
都在影響我們的下一代.
所以今天很重要.
我們要問自己.
我們今天是不是選擇好.
我們跟隨的是誰.
因為我們今天所選擇的.
其實是在影響我們的下一代.

$^{401}$而當我們選擇了之後.
我們去到今天的第二點.
就是我們要去移除.
攔阻我們生命事奉的障礙.
我們一起從第十九讀到二十一節.
好嗎.
我們一起讀.
一二三.
若書下對百姓說.
你們不能事奉耶和華.
因為他是聖潔的神.
是忌邪的神.
必不赦免你們的過犯罪惡.
你們若離棄耶和華.
去事奉愛邦神.
耶和華在降福之後.
必轉移降福與你們.
把你們滅絕.
百姓回答.
若書下說.
不然我們定要事奉耶和華.
這段聖經其實有點奇怪.
如果我們記得十八節.
明明以色列人說得很清楚.
我們必定事奉耶和華.
但是為什麼這麼奇怪.
當我們說完之後.
就是十九節.
若書下說.
你們不能事奉耶和華.
今天如果我問你.
有沒有想過在敬拜堂事奉.
你考慮完之後.
你回應.
我會在敬拜堂事奉.
你繼續跟你說.
你不能在敬拜堂事奉.
你可能會說.
那你問完之後又不能.
為什麼呢?.

$^{441}$其實這裡不是這樣的意思.
在這裡若書下不是否定事奉本身.
但是若書下說的很重要的事情.
有些事情你要首先搞清楚.
因為如果你沒有搞清楚.
你就不能夠事奉耶和華了.
所以這裡不是說你不可以事奉耶和華.
但是他看到他們的生命.
他說有些事情你要先搞清楚.
如果這些事情你沒有搞清楚.
你就不能夠事奉耶和華了.
而在這裡.
其實是說你要移除.
你們能夠事奉耶和華的障礙.
其實若書下在這裡看到.
當時以色列人.
他們那種是一種膚淺的形象.
我們事奉耶和華.
他們以為靠自己的能力.
不需要神的幫助.
其實若書下看到他們當時.
是輕率含糊.
以及三心兩意.
所以在這裡.
若書下說你要首先有些事情搞清楚.
你才可以事奉耶和華.
你要除去那些障礙.
你才可以事奉耶和華.
除去什麼障礙呢?.
第一.
他說神是聖潔的神.
所以今天如果你說你要事奉耶和華.
你首先要去除去你裡面的罪.
另一方面他說因為神是忌邪的神.
就是說你不可以今天突然間又說事奉耶和華.
接著又去三心兩意去拜其他的神明.
因為神是忌邪的神.
神要我們是專一的愛.
神也要我們有一個專一的事奉.
就好像婚約裡面.

$^{481}$是不可以有第二個情人的.
在婚約裡面只可以有一位情人.
有沒有人覺得在婚約裡面.
可以有第二位情人的?.
你舉手.
接著一會兒我帶你去見驅魔.
是不可以的.
在婚約裡面只可以有一位的情人.
同樣今天我們敬拜神.
我們只可以有一位神.
所以在這裡講得很清楚.
書也說你不可以三心兩意.
如果你決定了事奉神.
你就不可以再去拜其他的神明.
因為你去拜其他的神明.
你就是犯了罪.
必不赦免.
特別他在這裡說.
當我們跟隨神的時候.
我們會經歷恩典的.
我們會神降福給我們的.
但是今天如果你說你跟隨耶和華.
你經歷了恩典.
但是你接著又去拜其他的神.
你就會經歷神的降禍.
所以在這裡離棄.
是一個背若的詞語.
去投去另外一個神明的意思.
其實如果你說事奉耶和華.
又去事奉其他神.
其實是犯了十誡的第一和第二誡.
特別有時候我們有一個錯覺.
我們覺得經歷了恩典.
我們就做什麼都可以了.
因為有恩典.
但是其實正正相反.
恩典是讓我們知道神的真實.
我們經歷了恩典.
知道神的真實.
我們更加知道我們不可以拜其他的神.

$^{521}$我們經歷了恩典.
經歷了神的真實.
我們就更加明白背若的嚴重性.
所以到了第二十一節.
百姓這樣回覆說.
不然我們定要事奉耶和華.
在這裡其實給我們三個很重要的教導.
今天首先當我們說去事奉耶和華的時候.
其實我們要搞清楚.
我們事奉神的本質是什麼.
剛剛不是有人舉手說.
今晚吃父親節飯的.
有誰說今晚你會吃父親節飯的?.
有沒有人是你請你的父親吃飯的?.
祝福你.
你請你爸爸吃飯的時候.
你不會請他吃他喜歡吃的東西的.
是不是?.
今晚你不會請你爸爸吃父親節飯.
你爸爸不吃辣的.
然後你說.
爸爸今晚我們一起蒸麻辣火鍋吃.
你不會的.
你請爸爸吃飯.
就是請爸爸喜歡吃的東西.
所以今晚是我爸爸選北上的.
今晚我們北上吃飯.
就是說今天同樣我們去事奉神.
我們要去認識神的本質.
祂是聖潔的神.
我們要除去我們的罪.
祂是忌邪的神.
我們要除去我們的三心兩意.
第二就是當我們經歷恩典.
是經歷了神的真實.
而我們經歷神的真實.
就更加要明白背若的嚴重性.
而到了第三.
我們就要再一次立志回應.
我們單單事奉耶和華.

$^{561}$我不知道大家記不記得.
你有沒有看過我們十週年的話劇.
有.
有啊.
裡面其實有一個預言.
裡面的預言就是.
阿根廷二十週年的時候.
教會就沒落了.
我真的祈禱希望不是這樣子.
因為我是十週年接任做主任目標.
我真的不想二十週年的時候.
教會就沒落了.
大家一起努力.
但是他說其中一個原因是為什麼.
你記不記得.
就是那些父母說一套做一套.
他們在教會就是一套.
但原來他們的生命.
有些的罪.
有些的爭執.
有些的張力未除去.
今天原來我們.
要首先除去我們生命的罪.
如果不是.
其實我們會影響我們的下一代.
有一次我在教會裡面.
和一個小朋友吃東西.
和他一起吃薯條.
是的.
那麼.
我老婆就跟他說.
要擦了手才吃薯條.
誰知道他不開心.
然後就叫我的太太做.
嘟婆.
是的.
那個嘟字是一個號碼來的.
是的.
我們很害怕.
一個小朋友.

$^{601}$你叫他先擦手.
行不行?.
是的.
叫他先擦手.
不用.
嘟婆.
這樣罵我老婆.
是不是?.
但是為什麼原來.
我發現原來他爸爸常常在家罵他媽媽.
所以他可能以為這個是對媽媽的稱呼.
我們要除去我們的罪.
如果你不知道那個號碼是什麼.
你問問你旁邊的老婆.
不是那個不是那個是雙數的那個字.
是的.
不是單數.
單數就更加害怕了.
我告訴你.
雙數.
那個字是雙數.
是的.
我看到大家開始在想.
究竟是單數還是雙數.
雙數.
是的.
但是也害怕的.
其實也害怕的.
但是原來因為他爸爸常常說.
所以我們要小心.
我們要除去我們的罪.
第二件事.
我們要除去我們的三心兩意.
我曾經認識一位同學.
我跟他聊天的時候.
我就問你是不是信耶穌的?.
他說是的,我信耶穌的.
我說為什麼你信耶穌?.
他說因為我爸爸回教會.
他就帶我回教會.

$^{641}$但是接著我開始看到他生命有些迷信的行為.
有一天我跟他聊天.
過年了.
過年有什麼做?.
他說過年我爸爸會給我一封利是.
然後那封利是放在枕頭底.
說帶來好運的.
有沒有曾經你家裡給你一封利是放在枕頭底?.
是的,也試過.
他這樣跟我講.
我說不是的.
你明明跟我說你爸爸信耶穌的.
他說是的,他信耶穌的.
不過他也給我一封利是放在枕頭底的.
然後有一次他爸爸做生意.
開店鋪的.
開張.
他跟我分享很開心.
哇!昨天我爸爸的店鋪開張.
拜神斬燒豬.
我說你爸爸不是信耶穌的嗎?.
為什麼拜神斬燒豬?.
是的.
他又會從神祈禱.
但是開張的時候拜神斬燒豬.
說燒豬真的很好吃.
無限量供應.
我說這麼開心.
但是我說你爸爸不是信耶穌的嗎?.
然後有一次他去日本旅行.
我去日本旅行.
他跟我回來.
說日本做什麼?.
我說我去了那些廟.
那些廟可以做什麼?.
在那裡祈福.
我說怎麼祈福?.
我說會有一些牌子的.
我不知道你有沒有去過.
你不知道有沒有做過.

$^{681}$寫了那些牌子之後.
就放在那裡.
然後.
他說得很繪形繪色.
我說你又祈福?.
不是我祈.
我爸爸祈.
你爸爸不是信耶穌的嗎?.
是的.
但是他去到那裡又祈福.
他寫什麼?.
寫我學業進步.
我就經常想.
其實今天究竟我們是.
在事奉耶穌.
還是在事奉其他神.
但是原來我們的下一代是看在眼內的.
要回到我認識另外一位的肢體.
他信耶穌的時候.
他剛剛經歷了生意失敗.
又經歷了三個兒女的反叛期.
他常常分享那時候是他人生最難的時候.
但是在他人生最難的時候.
他說他堅定相信神.
他說他一天那時候要打兩份工.
但是他堅持星期天回教會.
堅持星期天去事奉神.
但是接著當面對那些兒女的反叛期.
甚至有些認識了不好的朋友.
他不斷繼續去關心他們.
用基督的愛愛他們.
以及沒有停止鼓勵他們回教會.
他說那些兒女看到爸爸這樣.
他們慢慢回轉.
有一個回到去讀書.
讀完書回來在教會服侍青少年.
有另外一個其實之後讀不了書.
去了做工.
但是他說他翻shift的.
他一有時間就去照顧那些老人家在教會.

$^{721}$第三個看到哥哥姐姐這樣.
之後回到去讀書.
然後繼續在教會的青少年團體成長.
你會發覺今天我們都會有困難和挑戰.
但是今天原來我們的生命.
如果我們堅持跟隨神.
原來我們不單止影響我們自己.
我們會影響我們的家庭和下一代.
今天我們願不願意去除去跟隨神的障礙.
我們的生命.
單單去事奉我們的主呢.
然後到今天的第三點.
是至於我和我間.
我們必定事奉耶和華.
我們由最後22讀到28節.
我們一起讀.
1,2,3.
若孫夏對百姓說.
你們選定耶和華要事奉他.
你們自己作見證吧.
他們說我們願意作見證.
若孫夏說.
你們現在要除掉你們中間的愛邦神.
專心歸向耶和華以色列的神.
百姓回答若孫夏說.
我們必事奉耶和華.
我們的神聽從他的話.
當日若孫夏就與百姓立約.
在事件為他們立定律例典章.
若孫夏將這些話都寫在《神的律法書》上.
又將一塊大石頭.
立在橡樹下耶和華的聖所旁邊.
若孫夏對百姓說.
看啊!這石頭可以向我們作見證.
因為是聽見了耶和華所吩咐我們的一切話.
講話你們背棄你們的神.
這石頭就可以向你們作見證.
於是若孫夏打發百姓.
各歸自己的地業去了.
當那些以色列人決定說.

$^{761}$我們要跟隨耶和華.
我們要移除事奉的障礙.
跟隨耶和華.
在這裡若孫夏就說.
好!現在你們就要有行動了.
第一個行動是你們自己見證立約.
今天你選擇去事奉耶和華.
在古晉東立約.
他們的簽署也是他們的見證.
就是說當我們今天去見證的時候.
代表如果將來我們違約.
我們不可以說.
哎!我不知道啊.
或者說.
我沒有答應過.
是你自己簽署和見證的.
今天你自己見證.
你選擇了生命要去跟隨耶和華.
接著二十三節.
若孫夏就說.
現在你們要做事了.
現在你們就要除掉你們中間的愛邦臣.
專心歸向耶和華.
應該當時在以色列營裡面.
他們仍然有些是在拜其他的神.
可能是被埃及.
或者在曠野裡面受別人的影響.
或者是被附近迦南那些人的影響.
所以.
他在那裡說.
你今天就要除掉在你們中間的愛邦臣.
可能你說.
這個跟我沒有什麼關係.
是不是?.
現在我家裡不會.
你看到有其他的神像.
我又不會拜其他的神明.
但是今天其實我們的生命可能有其他的偶像.
我們的偶像可能是名譽.
我們的偶像可能是我們的事業.

$^{801}$今天我們的偶像可能是股票.
可能是什麼?.
當然是賺錢的股票.
可能是金錢.
可能是成就.
或者今天我們有另外一個偶像.
不過偶像有時候我也提醒自己.
就是手機.
我們經常黏著手機的.
最近我也經常被人拍到.
我經常想著看手機.
所以有時候手機成為了我們的偶像.
有時候可能是娛樂.
以前我經常覺得.
煲劇是年輕人才煲的.
現在我發現.
原來我們真的從大到小.
從小到大我們都要煲.
不是.
好的劇是很好的.
但是我們小心那個成為了我們的偶像.
今天我們要除去這些生命的偶像.
因為我們跟這些偶像的時候.
其實我們的下一代會看到的.
好了.
到了二十四節.
接著有另外一件事情他們要做了.
就是百姓回答約書案說.
我們必事奉耶和華.
我們的神聽從他的話.
原來今天當我們說我們要事奉耶和華的時候.
我們是要聽從耶和華的話的.
聽從是在說聽和從.
就是說我們不僅要聽.
我們的生命還要信服.
因為真正的聽.
是包括了信服的.
昨天我太太來了崇拜我.
我就已經向她認罪了.
因為很多時候我老婆叫我做事情的時候.

$^{841}$我都沒有做的.
我沒有做的時候她就會問我一條問題.
你聽不聽到我說話?.
你有沒有試過這樣問你的配偶?.
你聽不聽到我說話?.
那我就很自然地回答她.
我當然有聽到.
但是問題是我聽到沒有做.
你明白嗎?.
其實她問我聽不聽到.
不是問我是不是真的耳朵聽到.
是因為你沒有做.
所以很多時候.
昨天我太太回來第二課.
我向她認罪了.
因為很多時候她叫我做事情.
是不是?.
我就不做.
不做的時候她就問.
喂!你聽不聽到?.
我說我聽到.
但是其實不是真的聽到.
因為你聽到你沒有做.
就不是真的聽到.
同樣今天以色列人也不可以就這樣聽到.
他要聽到並且去跟從.
而到了二十五節.
約書亞就在那裡跟百姓立約了.
在這裡是立刻立約的.
這裡講當日.
不是說多好.
我說明年開始吧.
好不好?.
不是.
當日約書亞就跟百姓立約.
在士堅這個地方.
士堅是一個很重要的地方.
因為亞伯拉罕曾經在這裡築壇.
而亞國也曾經在這裡立約.
所以在這麼重要的地方.

$^{881}$約書亞跟他們立約.
但是立約有一個很重要的事情.
就是要立定律例典章.
就是說.
今天神的話語要成為我們生活的規章.
神的話語要成為我們生活的標準.
神的話語要成為我們生活行事為人的準則.
今天我們也要問自己.
當我們說事奉耶和華的時候.
究竟我們的生命有沒有神的律例典章呢?.
我們的生命有沒有神的話語呢?.
你問我很感恩.
因為接下來約書亞區.
第十一屆門女就要開始了.
我們一起一個月.
我們說我們每一天靈修.
然後我很開心.
因為我們分成六個小組.
這六個小組裡面.
每天我都看到弟兄姊妹一起靈修.
並且分享我們中間靈修的得著.
我們每一天都需要有神的話語.
因為神的話語要成為我們生命的律例典章.
而當他們這麼做的時候.
第二十六節到二十八節.
就是在說納若的過程.
第二十六節我們看到.
約書亞就將這些東西寫在律法書上.
又將一塊大石頭.
納在象樹下.
耶和華的聖所旁邊.
其實這個大石頭.
就好像一個紀念碑一樣.
紀念碑是什麼呢?.
就是提醒他們和他們的後代.
你們在這裡跟神納若.
一生事奉耶和華.
而二十七節有一個很特別的擬人法.
約書亞說.
石頭也可以跟我們作見證.

$^{921}$因為今天石頭也聽到耶和華吩咐我們的一切話.
所以以後你們離棄神.
這個石頭可以跟你作見證.
當然這是一個擬人法.
石頭就不是真的聽到我說話的.
雖然今天我想說.
其實我們的手機是聽到我們說話的.
我曾經試過.
就這樣放在手機上.
說我想買一個沙發.
接著我開始開那些社交媒體.
那些沙發的廣告就這樣彈出來.
所以雖然石頭聽不到你說話.
但是手機聽到你說話的.
但是那個石頭就不是聽到他說話的.
不過當時約書亞是用一個誇張的手法.
他說你今天納若.
石頭也聽到你說話的.
所以你要記住你今天跟神的納若.
接著第二十八節.
不要以為是最後.
是一段很重要的聖經.
在這裡他說.
約書亞就打發百姓.
各歸自己的地業.
為什麼呢?.
因為原來今天我們領受了神的話語.
我們領受了祂的律例.
典章不是留在這裡的.
我們要帶回我們的家.
所以今天我們作為爸爸的.
我們作為丈夫.
我們要將神的話語帶回我們家.
成為我們家的動力和祝福.
所以今天大家都不會留很久的.
現在12點44分.
在這裡一點鐘就結束了.
為什麼呢?.
因為我們都要回到我們自己的家.
將神的話語帶回我們家.

$^{961}$去成為我們的祝福.
至於我亡我家.
我們必定事奉耶和華.
在年頭的時候.
在約書亞區的區序裡面.
有一個很感動的見證.
就是我們邀請了一個家庭.
一家四口.
一起上來分享他們生命的見證.
這個家庭的爸爸.
在一次的特會裡面.
他立志要成為耶穌的門徒.
於是他不斷尋找一個合適的教會.
他很想尋找一個合適的教會.
是讓他和他一家都做耶穌的門徒.
所以他分享他找教會的時候.
不只是找自己很喜歡的.
他會看我的太太能不能夠適應.
我的女兒能不能夠適應.
他有兩個女兒.
所以他說在找教會的時候.
他不斷找神.
你給一個教會讓我和我一家可以做你的門徒.
神就將這一家帶了來雅庚堂.
而他的爸爸.
首先去參加門徒成長轉化旅程.
在過程裡面.
看到這個爸爸生命的改變.
其實有些改變我不是真的這麼看得到.
為什麼呢?.
因為我認識他的時候他已經改變了.
他分享說他以前是一個不說話的爸爸.
但是我認識他之後他很窩心.
他不只是對家庭窩心.
他對牧者也很窩心.
有一次他寫了一個很窩心的訊息給我.
他說我們一班弟兄看到你的狀況很擔憂.
很覺得你需要跑步.
很想邀請你和我們一起跑步.
我很少收到弟兄這麼窩心的訊息.

$^{1001}$因為他這麼窩心.
所以我現在一個月也說跑兩次.
兩次也不是很幫得上忙.
所以他的窩心我相信不只是對他家庭.
他真的對弟兄子母也是這麼窩心.
他參加了門女之後.
他反思自己的生命.
他更加花時間和家人相處.
而在這個過程裡面.
他特別和他女兒的關係更加盡心和融洽.
所以不只是這位爸爸去參加門女.
而之後家裡的媽媽和兩個女兒.
都參加了我們成人牧區不同的門女.
我真的看到原來一個爸爸生命的改變.
一個爸爸願意帶領整個家庭走一個方向.
我就真的看到這一段聖經.
至於我和我家.
我們必定是鳳爺和華.
而我真的親身在一個星期的下午.
看到一個很感動的場面.
在星期日的下午.
大部分我們星期日下午都不在教會.
但是我就看到這一家人.
星期日的下午一起在教會.
我就看到爸爸和媽媽.
在我們愛倫社的崇拜那裡服侍.
姐姐在我們星期日的少年崇拜那裡教英文.
也是我們星期日少年崇拜的導師.
去關心我們星期日的少年人.
而妹妹那天是在參加一個門女關顧的一個祈禱會.
接著他們一起等爸爸開車一起走.
我看到他們的時候我真的很感動.
因為我正正看到.
就是一個爸爸的生命.
怎樣帶領他的家庭.
去到至於我和我家.
我們必定是鳳爺和華.
所以今天當我們去到總結的時候.
我們請敬拜罪.
我發現原來這一段的經文.

$^{1041}$給我們一個很重要的訊息.
就是今天人是不會不事奉的.
那你說不是啊.
我沒有事奉光偉就不事奉了.
其實事奉永遠不只是在說這件事.
事奉是在說我們跟隨的方向.
敬拜的方向.
今天我們每一個都事奉的.
不過你未必事奉神而已.
今天我們可能不事奉神.
事奉什麼呢?.
就是事奉我剛才說的那些.
今天我們可能會事奉名譽.
事奉金錢.
事奉地位.
事奉人的讚賞.
事奉享樂.
事奉娛樂.
大家記得我曾經說過.
我以前認識一位肢體.
他經歷了很多.
神幫助了他很多之後.
他沒有再回來教會.
我說為什麼你經歷了神那麼多之後.
你不回來教會啊?.
接著他不回來他沒有回答我.
接著我問他男朋友.
他說他現在最喜歡就是.
回家看劇集.
是的.
原來我們的生命都是在事奉一些東西的.
我們是不會不事奉的.
只不過我們要選.
我們事奉的是什麼.
而第二件事很重要.
事奉不是口說的.
我們不是說今天我事奉神.
事奉是我們的生命堅持的根除.
而第三.
原來今天我們的事奉.

$^{1081}$是會影響我們的下一代的.
所以當約書亞說這句話.
我們再讀一次好不好?.
一二三.
至於我和我家.
我們必定事奉耶和華.
約書亞不是說你事奉耶和華.
你先去.
不是.
約書亞說我自己作為一個榜樣.
我自己和我的家.
首先作為一個榜樣.
就是在這個充滿偶像的世代裡面.
我們要做一個堅定的決定.
今天在父親節裡面.
我特別很想邀請我們當中.
不僅僅是做爸爸的.
我也很想邀請我們當中每一位做丈夫的.
因為聖經說得很清楚.
做爸爸.
做丈夫.
是家中的頭.
今天原來我們有一個位份.
原來我們是帶領著我們家庭的方向.
今天原來我們所重視的.
我們所追求的.
我們所事奉的.
是會影響我們的下一代.
今天可能我們的生命仍然有掙扎.
我們也可能掙扎在不同的價值觀當中.
但是今天神給我們一個機會.
我們重新立志.
而今天我們當中.
我們未必是做爸爸,做丈夫.
但是今天神都呼召我們每一個作主的門徒.
當我們說我們做門徒的時候.
其實同樣我們的生命.
和我們的家庭.
都是必定事奉耶和華.
當我們這樣做的時候.

$^{1121}$我們的生命就讓人看到.
我們生命裡面的亮光.
我們生命的意義.
就讓人的生命都願意去跟隨耶穌.
今天你願不願意.
你和你家都一生事奉耶和華.
\newpage



\section{約翰福音 1:37-51}
\label{sec:032357}
\textbf{蒙主呼召,更大的事}~~~~羅婉明傳道
\newline
\newline
連結: \href{http://www.church.com.hk/acms/content.asp?site=cdc\&op=show\&type=product\&code=032357\&layout=sermon}{\texttt{032357}} ~~~~ 語音日期: 2026-06-21
\newline
\newline
\hyperref[sec:032356]{\small{< < < PREV SERMON < < <}}
~
\hyperref[sec:index_chronic]{\small{[返順時目]}}
~
\hyperref[sec:032431]{\small{> > > NEXT SERMON > > >}}
\newline
\newline
約翰福音 1:37-51
\newline
\begin{longtable}{cl}
\hline
\hline
章節 & 經文 (和合本修訂版)\\
\hline
1:37 & \begin{tabularx}{0.7\textwidth}{X} 兩個門徒聽見他的話，就跟從了耶穌。 \end{tabularx} \\ \\ \relax
1:38 & \begin{tabularx}{0.7\textwidth}{X} 耶穌轉過身來，看見他們跟著，就問他們說：「你們要甚麼？」他們對他說：「拉比，你在哪裡住？」（「拉比」翻出來就是老師。） \end{tabularx} \\ \\ \relax
1:39 & \begin{tabularx}{0.7\textwidth}{X} 耶穌說：「你們來看。」他們就去看他在哪裡住。這一天他們就跟他同住；那時大約是下午四點鐘。 \end{tabularx} \\ \\ \relax
1:40 & \begin{tabularx}{0.7\textwidth}{X} 聽了約翰的話而跟從耶穌的那兩個人，其中一個是西門‧彼得的弟弟安得烈。 \end{tabularx} \\ \\ \relax
1:41 & \begin{tabularx}{0.7\textwidth}{X} 他先找到自己的哥哥西門，對他說：「我們遇見彌賽亞了。」（「彌賽亞」翻出來就是基督。） \end{tabularx} \\ \\ \relax
1:42 & \begin{tabularx}{0.7\textwidth}{X} 於是安得烈領西門去見耶穌。耶穌看著他，說：「你是約翰的兒子西門，你要稱為磯法。」（「磯法」翻出來就是彼得。） \end{tabularx} \\ \\ \relax
1:43 & \begin{tabularx}{0.7\textwidth}{X} 又過了一天，耶穌想要往加利利去。他找到腓力，就對他說：「來跟從我！」 \end{tabularx} \\ \\ \relax
1:44 & \begin{tabularx}{0.7\textwidth}{X} 這腓力是伯賽大人，是安得烈和彼得的同鄉。 \end{tabularx} \\ \\ \relax
1:45 & \begin{tabularx}{0.7\textwidth}{X} 腓力找到拿但業，對他說：「摩西在律法書上所寫的，和眾先知所記的那一位，我們遇見了，就是約瑟的兒子拿撒勒人耶穌。」 \end{tabularx} \\ \\ \relax
1:46 & \begin{tabularx}{0.7\textwidth}{X} 拿但業對他說：「拿撒勒還能出甚麼好的嗎？」腓力說：「你來看。」 \end{tabularx} \\ \\ \relax
1:47 & \begin{tabularx}{0.7\textwidth}{X} 耶穌看見拿但業向他走來，就論到他說：「看哪，這真是個以色列人！他心裡是沒有詭詐的。」 \end{tabularx} \\ \\ \relax
1:48 & \begin{tabularx}{0.7\textwidth}{X} 拿但業對耶穌說：「你從哪裡認識我的？」耶穌回答他說：「腓力還沒有呼喚你，你在無花果樹底下，我就看見你了。」 \end{tabularx} \\ \\ \relax
1:49 & \begin{tabularx}{0.7\textwidth}{X} 拿但業回答他：「拉比！你是神的兒子，你是以色列的王。」 \end{tabularx} \\ \\ \relax
1:50 & \begin{tabularx}{0.7\textwidth}{X} 耶穌回答他說：「因為我說在無花果樹底下看見你，你就信嗎？你將看見比這些更大的事呢！」 \end{tabularx} \\ \\ \relax
1:51 & \begin{tabularx}{0.7\textwidth}{X} 他又說：「我實實在在地告訴你們，你們將要看見天開了，神的使者在人子身上，上去下來。」 \end{tabularx} \\ \\
[1ex]
\hline
\hline
\end{longtable}
$^{1}$今天講到耶穌呼召門徒.
上星期聶牧師講到施洗若寒.
當你聽蒙主呼召.
即是信耶穌的經歷時.
我相信大家都有不同的信主經歷.
或者大家在信主的經歷中.
都有很多與上帝的相遇中所經歷的事.
我都講一下我自己的信主經歷.
我其實是在內地出生的.
我十幾歲來香港.
我來到香港的時候.
其實我對於耶穌是未聽過的.
我學校有一位老師真是很好.
他帶我回教會.
他帶我回教會那天其實是報道會.
我那天其實真是人生裡面.
第一次回教會 第一次去報道會.
亦都真是很正正式式.
第一次聽耶穌的名字.
其實我都覺得很奇妙.
那段日子我其實一直在尋索.
其實人生是為了什麼.
到底我人生是什麼人生來的.
所以當我去到報道會的時候.
你說報道會講什麼呢.
我猜大家相信耶穌.
應該都不太記得報道會講什麼.
但我只記得一句經文.
是上帝用一句經文打入我心裡面.
那句經文是說耶穌來是叫人得生命.
並且得得更加豐盛.
我聽完之後我覺得好像在解答我的事.
又好像在解答 又好像還未完全解答.
我就在想 真是豐盛的生命.
是否真是豐盛的生命呢.
我真的很想.
我真是很搞笑地想.
應該不會騙我的吧.
就算騙都沒有理由幾百人在騙.
他再講一些見證就歷勢歷勞.

$^{41}$沒有理由這麼多人在騙你.
騙你為什麼呢 你是哪位呢.
那一刻我就覺得 好 我相信.
我很想相信.
我很想要豐盛的生命.
我覺得回望這麼多年.
不要講這麼多 相信了多少年.
回望這麼多年.
我覺得上帝真的給我有豐盛的生命.
那個豐盛是在神裡面見到祂的榮耀.
見到祂的能力.
見到祂怎樣帶領我走過.
甚至是高高低低都好.
但卻是很微.
那種微是因為在上帝裡面的微.
我相信當我們一路去講.
耶穌呼召的時候.
可能你都會想起你昔日的順主的經歷.
或者你順主的歷程是很不同的.
但很重要的是耶穌遇見你.
耶穌拯救你.
我們一起祈禱.
主我懇求你.
當我們在你裡面去看你的話語的時候.
真你的聖靈打開我們的心.
打開我們屬靈的眼睛.
以致我們能夠真是被你的話語再一次更新我們.
最重要是上帝你親自的捉緊我們的心.
以致我們是緊緊的跟隨你.
求你與我們每一個同在親切的帶領我們.
幫助我們.
多謝主.
我們禱告奉耶穌基督明求.
阿們.
講蒙主呼召.
我們就講幾個門徒他怎樣遇見耶穌.
剛才你聽到由37節開始讀的.
但我就給35,36節.
因為都是有少少的背景.
37節一開始就講兩個門徒.

$^{81}$那兩個門徒其實是誰的門徒.
最開始.
約翰的.
施洗約翰的.
施洗約翰真的很特別.
他自己帶著門徒.
但他又給你看其他的師傅.
他見到耶穌走過來的時候.
我全世界都幾大聲.
哇 漢啊 叫你漢啊一定不會很小聲.
漢啊 寫字神的羔羊.
漢啊 這個是神的羔羊.
即是說除去我們世人罪孽.
約翰應該是第二次在他的門徒面前這樣講.
於是37節就講這兩個門徒.
很快反應.
我用紅色字給你看.
他有多快反應.
用了幾個動詞.
漢 聽見.
跟從.
很快.
我先回到正題.
我太快了.
當他們跟從的對面.
有一個叫安德烈.
看看他們有什麼反應.
他們跟從了耶穌.
然後耶穌轉身.
見到他們跟著.
耶穌問了他們一個問題.
這個問題實在太好了.
什麼時候問都可以.
耶穌就問你們要什麼.
問你們要什麼.
即是你們跟從我 你們要什麼.
然後他們就回答.
夫子在哪裡住.
這些問題很好笑.
如果你代入今天香港的社會.

$^{121}$你會覺得.
這樣問的.
然後你很快就會想.
騙子.
提防騙子.
因為我問你要什麼.
在哪裡住.
然後就說你們看吧.
看完之後.
他們很快就跟了.
就住在那裡.
我想這個跟他們的背景有關係.
我想昔日的使徒.
或者師父帶徒弟.
或者是夫子帶弟子.
其實我想昔日是會一起住的.
關係是很親密的.
跟我們今天都不同.
因為我們學校的老師都見得很長.
但你總是會回家的.
會自己住的.
如果是回來教會主日學的老師也好.
或者是阿灣娜老師也好.
其實就很短時間.
一個多小時.
見完就回家了.
缺席一兩次.
一個月見兩次.
但他們那時候是不同.
是同食 同住.
師父去哪裡.
他跟著去哪裡.
是一個很親密的關係.
所以他們很快就跟耶穌一起住了.
然後你看到安德烈的行動很快的.
看 聽.
聽完之後就跟從.
跟從完就去看.
然後就一起住了.
安德烈再做了一件事.

$^{161}$做了什麼呢.
好 給你看看.
去找誰呢.
西門 西門是他哥哥.
很快就回家找家人了.
是好事.
是好事 馬上去找.
很快就去找哥哥.
然後就跟他說.
我們遇見彌賽亞了.
這三個字是不得了的.
為什麼呢 對猶太人來說.
彌賽亞代表什麼呢.
救世主.
是他們等候了很久的.
日盼夜盼.
都是盼著救世主來的.
所以安德烈跟西門說.
為什麼安德烈馬上去找西門.
馬上去找西門.
因為他覺得太好了 太棒了.
因為找到救世主.
找到我們等了很久的那個.
他就馬上去找西門.
找阿波 找他回來.
就去見耶穌了.
你見到安德烈就帶人去信主.
其實有一些教會.
有些叫安德烈行動.
知不知道安德烈行動是什麼意思呢.
就是帶家人去信主.
就是這樣 帶家人去信主.
因為他真的很快.
他信了之後就很快帶家.
這個就是安德烈了.
下一個我看一看.
下一個就是西門.
他找阿波 找西門.
西門來到之後.
耶穌就看著西門.

$^{201}$看著他.
然後就跟西門說了一些話.
我們一起讀李氏的.
李氏約翰的兒子西門.
你要稱為基發.
基發翻出來就是彼得.
很特別的.
這裡西門有沒有說話的.
沒有說話的.
只是耶穌.
鏡頭好像是停不下來.
停在那裡就是他看著西門.
他看著西門.
然後他就給了他一個名字.
給了他一個名字.
這個名字就叫基發.
基發是亞蘭文來的.
為什麼是彼得.
都是同一個意思.
不過是希臘文.
彼得是希臘文.
大家知道亞蘭文是當日他們說的語言.
希臘文是寫出來的.
或者是官方的.
所以都是同一個意思.
是什麼意思呢.
就是石頭 磐石.
大家會知道了.
之後耶穌就跟彼得說.
我要將我的教會建在磐石上.
彼得是初期教會一個非常非常重要的領袖.
這裡很特別.
你會發現.
其實西門有沒有做什麼.
沒有.
只是站在那裡.
我們說是佈景.
但是是耶穌選擇他.
耶穌一選擇他就告訴他.
我給了你一個身份.

$^{241}$整個的選擇是更新了他整個身份.
名字是很重要的.
名字是代表著那個人.
代表著那個東西.
所以猶太人的名字是很重要的.
給了他就說.
其實你就是磐石.
我要將我的教會建在上面.
是一個很重要的.
弟兄姊妹.
其實你看到這一段的時候.
你會發覺.
我自己看.
我就覺得很感動.
感動是什麼呢.
彼得什麼都沒有做.
但是神說我就要選擇你.
其實就是這樣.
就是這樣.
這個就是恩典.
我記得六月頭.
六月七號的時候.
星期天我在兒童那裡說了一個主題.
那個主題就是.
我的被造.
我的被造很奇妙.
就是說.
神造我們的時候.
神造我們很獨一無二.
神造我們是有一個很奇妙的計劃.
每一個人都是.
你不需要羨慕你旁邊的人.
我跟孩子說.
其實你很有價值在神裡面.
你的價值不是因為你讀哪間學校.
你的價值不是因為我是名校.
所以我厲害一點.
不是.
你也不是因為你讀什麼班.
是精英班還是普通班都不關事.

$^{281}$也不是因為你的成績.
就是因為神體你就是這麼有價值.
神裡面神造你就是我受造奇妙可畏.
我還送了兩塊盾牌.
當別人笑你也好.
或者你覺得自己很沮喪也好.
你就用這兩塊盾牌來抵擋.
我受造奇妙可畏.
我是神的工作.
我覺得很重要.
因為我很想我們的孩子知道.
你不要被外面的東西.
彷彿很混亂.
我們的價值在哪裡.
有時候他們都不知道.
甚至不好意思我說一句.
有時候我們的家長都不知道.
但我很想親愛的家長.
我們要守護我們的孩子.
就是他很有價值.
他的價值不在乎他做什麼出來.
就正如耶穌選擇彼得.
你看這裡其實沒有說話.
就是因為我選擇你.
因為彼得我就是選擇你.
又不是誰重要一點.
安德烈也一樣這麼重要.
我們再看下一個.
下一個就是肥力.
肥力和耶穌的相遇.
不是人帶他去的.
這裡只是說耶穌要去加利利.
就遇見肥力.
然後就跟肥力說.
你跟從我吧.
是不是很簡單.
有時候我也覺得是不是太簡單.
就是這麼簡單.
就選擇了肥力.
肥力有跟從嗎.

$^{321}$有啦.
跟從完才去找他的朋友.
馬上去找朋友.
他信了之後.
跟從了耶穌.
就馬上去找他的朋友.
他的朋友就是拿旦葉.
你看他們的對話.
其實是很學術的.
很學術的.
我讀完解釋給你聽.
肥力就找拿旦葉.
我們在律法書上寫的那個.
和在眾先知書上所記的那個.
我遇到了.
這兩句說話.
這句說話其實很深.
為什麼很深.
他說經上.
經上就是舊約聖經.
因為新約的時候還沒有新約聖經.
舊約是什麼呢.
就是律法書.
先知書和著作.
著作就是百思針傳歌.
你想一下是不是全部都說了.
他們平時聊天聊什麼內容.
聊聖經.
還要聊這麼多篇幅.
你想一下多學術.
很厲害的.
很宗教性的.
所以他說.
我遇到了.
就是那個拿撒勒人耶穌.
拿旦葉有不少東西.
容許我這樣說.
拿旦葉有幾串的.
他說了一句什麼話.
拿撒勒什麼.

$^{361}$有什麼好東西.
意思是.
我看完聖經都沒有提過拿撒勒.
厲害的出處的.
我都知道了.
拿撒勒有什麼好東西.
那麼.
肥力就很有智慧.
我們傳福音都要找一下.
這些智慧和這些主意.
肥力有跟他爭辯嗎.
沒有爭辯.
他說你來看看吧.
回來教會看看吧.
Come and see.
如果很多人不停去挑戰你.
問你很多很多問題.
其實不懂回答的時候.
都不要緊的.
帶他回來教會.
相對會解答他的.
這個就是肥力.
肥力都是很有行動的.
遇到耶穌.
跟從.
然後就令人見耶穌.
跟誰很像.
跟安德烈都很像.
我們來看看這個.
那麼獨特的拿撒勒.
拿撒勒就這樣說.
然後耶穌就回答他.
耶穌就說.
這個真是一個真意識的人.
心裡面沒有鬼詐的.
沒有鬼詐.
為什麼.
有碗就有碗.
有碟就有碟.
想什麼就說什麼.

$^{401}$然後呢.
其實最奇妙的一個位置.
就是他講完這句話之後.
其實拿撒勒問了耶穌一句話.
他說你在哪裡知道我.
意思是.
為什麼你這麼認識我.
為什麼你這麼知道我.
然後耶穌就回答他.
他說其實肥力沒找你之前.
你在無花果樹底下的時候.
我就見到你了.
其實意思是什麼.
就是說其實我一早就見到你了.
我一早已經見到你了.
我一早已經知道你了.
弟兄姊妹感不感動.
我自己就感動.
神告訴你.
其實我一早就知道你了.
你沒串我之前.
我一早就知道你了.
意思就是.
其實耶穌不知道拿撒勒是一個怎樣的人.
他知道的.
但他也知道.
他一早就知道他是一個怎樣的人.
他很想告訴他.
我一早見到你.
我知道你是怎樣的人.
所以其實拿撒勒很特別.
他從耶穌的回應裡面.
認識到自己.
耶穌認識他.
他也認識到自己.
意思是.
他回應了耶穌.
這是拿撒勒的扭轉.
因為他知道耶穌認識他.
他知道耶穌見到他.

$^{441}$他知道耶穌明白他.
所以這是拿撒勒人生裡的其中一個扭轉.
那個扭轉就是他認順.
他那一刻從說有什麼好東西.
到說你真是神的兒子.
你是以色列的王.
這是一個認順.
認順就是.
這個真是神差你的受高者.
是真是我們期待的那個.
是真是神的兒子.
是差下來拯救我們那位.
是以色列王有權柄的.
耶穌就跟他說.
你只是覺得我認識你就已經信.
其實你將來要信更大的東西.
或者你將來要經歷更大的東西.
其實這是一個信心的開啟.
可能我們都在經歷不停不停信心的開啟.
當主打開我們.
當主跟我們說.
我們是誰.
或者跟我們一起經歷不同的東西的時候.
其實是讓我們看更大更大的事.
當我自己回望.
耶穌說我要叫你的生命得得更加豐盛的時候.
我真的確實覺得是越走越豐盛.
和越多神的恩典.
當然不是說一帆風順.
從來都不是.
但是他是在困難.
在困苦裡面.
他是見到神的恩典.
神的同在.
神的實在.
並且是神的榮耀.
那種豐富.
那種很多之材.
所以是很奇妙的一件事.
四個門徒你會見到.

$^{481}$還有讓你看一件事.
就是剛才你一直讀聖經的時候.
不知道你有沒有留意.
有一個字是常出現的.
就是漢字.
我不會問你出現了多少次.
因為太難了.
考小朋友的時候可能會的.
因為看著這本聖經.
出現了很多次.
十三次.
如果你看聖經的話.
常常重複代表什麼.
重要.
這個也是約翰很想表達的.
就是漢.
很多不同的漢.
來漢去漢.
有些是耶穌叫我們來看.
有些是我們去看.
這個漢其實說的是.
耶穌看見我們.
你見不見到耶穌看見我們.
耶穌看見我們之後.
你能不能夠同樣漢見到你自己.
那個漢真的很重要.
你看到什麼.
你眼睛.
你熟悉的眼睛裡面.
看到什麼.
那個影響著你的人生.
很重要.
所以在經文裡面.
說了很多耶穌的描述.
剛才我們在讀的時候.
漢是神的羔羊.
利塞亞是神的兒子.
以色列王和仁子.
說了很多耶穌的描述.
不單止這些.

$^{521}$但是我想問一句.
你看見這麼多描述.
我問你一個問題.
耶穌是誰.
你心想.
答完了.
不是啊.
弟兄姊妹.
我是問耶穌對你來說.
耶穌是誰.
耶穌在你生命裡面.
當你經歷困難的時候.
耶穌是誰.
當你經歷一些抉擇的時候.
耶穌是誰.
當你背負著很多事情的時候.
耶穌是誰.
耶穌是走在一邊.
還是真是這些.
當你面對一些沉溺.
或者你覺得很難.
去原諒自己的一些事情的時候.
耶穌是誰.
那個是神的羔羊.
救贖你.
將你的罪孽背在祂身上.
讓你能夠釋放.
讓你的綑綁被剪開.
讓你能夠享受這份自由.
祂是神的兒子.
祂有那個權柄.
祂有那個能力.
是不是.
當你很煩你的事情.
或者很煩你的子女的時候.
前路的時候.
你問一句.
耶穌是誰.
耶穌對你的子女是誰.
祂也是你的子女的主.

$^{561}$在我的生命裡.
耶穌是誰.
這是一個問題.
是常常都要問.
不只是聖經裡說耶穌是誰.
是今天你問你自己.
耶穌啊.
你是我誰.
耶穌是誰.
很重要的一個問題.
所以我們看回不同的門徒.
他們經歷不同的呼召.
剛才說了四個.
其實還有一個.
有沒有留意.
剛才不是說約翰有多少個門徒.
有兩個.
他只說了一個安德烈.
其實還有一個.
沒有名字.
有些人說是約翰.
因為經常有人封音.
經常喜歡不出自己的名字.
無論如何.
你看到當耶穌呼召門徒的時候.
不同的性格.
不同的際遇.
不同的經歷.
不同的位分.
神給他們不同的位置去經歷上帝.
有名也好.
沒有名也好.
在聖經裡沒有名字.
在上帝裡面一定有名.
因為神見到.
神看見.
我相信在你不同的人生裡.
都有很多不同的經歷.
在人生裡不同階段.
你都會經歷到神呼召你.

$^{601}$有不同的使命.
神召你去做不同的事.
我們昨天有講座.
我們昨天講座.
陳醫生來講.
我講一些他的見證.
其實他是CIA醫生.
他之前在政府醫院出來的時候.
就執業了一段時間.
後來他就祈禱.
他問神.
神你想我下一步做些什麼.
然後神就回答他.
叫他去一個地方.
十幾年前.
十幾年前有一個.
什麼悲情.
悲情城市.
在哪裡.
天水圍.
正值那段時間.
很多家庭的悲劇在那裡發生.
神叫他去那裡開診所.
那裡歪一點可以嗎.
屯門 元朗可以嗎.
不行.
就去那裡.
去到的時候開診所.
開診的時候.
不單是他.
還有一些在教會跟教會的連結.
有些牧者.
他們祈禱之後的領受.
其實他們的診所.
不單是醫人的身體.
他覺得心理 靈性都是很重要.
所以他的診所同時間.
都有輔導.
有些是有開藥.
同時間有些是輔導.

$^{641}$長話短說.
當他一路做一路做的時候.
其實一直擴充.
也有不同的輔導員.
他們就發現.
其實在那個城市裡.
在天水圍有很多的需要.
他們就決定.
有一次就決定.
很想用一些很漂亮的情緒病的藥.
比較漂亮的藥.
他們就去訂.
跟藥廠訂.
訂完之後沒有回音.
去到有一天.
然後就有一個藥廠的代表來.
一個很成熟的代表.
那位姑娘就說.
藥廠代表來了.
就是我們申請那個.
這麼成熟.
為什麼這樣問呢.
因為通常去賣藥的.
去診所的那些.
通常是年輕的.
成熟已經上去了.
他就坐下來聊天.
聊天之後發現.
原來真的是高層來的.
他就說.
為什麼他會來呢.
他說其實全香港沒什麼診所.
會訂我們這麼貴的藥.
我們覺得很奇怪.
所以我就要下來看看.
後面那句就是.
看你有什麼厲害.
訂這麼貴的藥.
然後他就跟代表.
分享他們的領袖.

$^{681}$他們看到整個社區.
聽完之後.
藥廠的代表就說.
如果你們真的找到那些有掃描員.
你會發覺他們真的很有需要.
家庭也是很有困難的.
我們藥廠免費資助.
他就小聲說.
藥廠很少做這件事.
你看到原來神呼召的時候.
神會開路.
很奇妙.
後來也有不同的藥廠.
提供了優惠給他們.
我想說的是.
當神呼召我們.
有不同的崗位的時候.
上帝一定開路.
只要你認準跟從的時候.
神會使用你成為很多人的祝福.
今天是父親節.
我也分享一下我們家裡的事.
我老公很喜歡唱歌.
不是當你唱歌.
他在家裡除了喜歡唱歌.
還會彈琴.
自己自彈自唱.
如果你喜歡唱歌的人.
你會知道.
總有一些時期.
你會很喜歡不停地唱一首歌.
大約兩年多前.
他很喜歡唱一首歌.
叫做《我知誰掌管明天》.
有些人有共鳴.
他經常唱.
有一天晚上.
我兒子洗澡.
洗完澡就衝出來很興奮地對我說.
「媽媽,剛才神回應了我」.

$^{721}$「什麼事?」.
「在洗澡的時候回應你」.
他說「我最近有些困惑」.
「有些事情我想不通」.
「想不到結果」.
「我祈禱但覺得神沒回答我」.
「但剛才洗澡神回答了我」.
「神讓我想起」.
「爸爸經常唱那首歌」.
「《我知誰掌管明天》」.
「死味」.
他很興奮.
其實我也很興奮.
我也很感動.
我感動的地方是.
我覺得是他和神之間的經歷.
這是很重要的.
他祈禱.
然後神回應他.
這是很重要的.
是他的信仰.
是父母的信仰.
當然我們在旁邊.
我們也在做我們要做的事.
那時候爸爸還沒回來.
他很興奮.
要等爸爸回來.
等爸爸回來.
一起唱這首歌.
一起敬拜.
我兒子現在在教會.
參與敬拜隊.
他自己也很喜歡在裡面敬拜神.
我想在這個位置.
很想表達一件事.
父親在家裡所做的.
其實你的孩子都看到了.
你不要想到光環.
我老公彈琴.
大家圍著.

$^{761}$其實不是的.
他有他自己.
但問題是.
原來一些很細微的事.
影響著我們的子女.
你記得李牧師經常說.
他和他的兒子做什麼.
看球賽.
看球賽是他們很喜歡的事.
我覺得是.
不同的父親和他們的孩子.
那個互動是不同的.
那個影響也是不同的.
所以我很想鼓勵.
也很想鼓勵我們的父親.
你在家裡所做的.
是有用的.
是神紀念的.
我都在耶穌裡面honor你們.
其實你們都做了很多.
做的事.
很辛勞.
很努力.
雖然我們經常說.
父親在家裡好像被遺忘的角色.
但我覺得上帝不會遺忘.
耶穌罕見.
一定罕見.
我想表達的就是.
很珍惜.
也很重要.
你們在家庭裡.
你想想舊約的雅各.
為他的兒子祝福.
祝福就是這樣.
神都是給權柄.
父親作為一個家庭.
其實祂都是覆蓋著他的兒女.
是一個帶頭的位分.
所以當我們一直說的時候.

$^{801}$不同的崗位.
不同的使命.
神呼召你的時候.
你就去做.
神會給你能力.
剛才我們說.
主耶穌很知道.
很看到我們.
然後呢.
但有一件事.
當神帶領耶穌呼喚我們的時候.
耶穌都問一件事.
你們要什麼.
這句話我覺得也是.
隨時都會帶來一個很大的迴響.
也很值得我們去深思.
耶穌都在問我們.
你要什麼.
對我來說.
我覺得.
信主那一刻那一句經文.
我真的抓得很緊.
我相信在耶穌裡面.
就有豐盛的生命.
是在祂裡面.
要在祂裡面才行.
在祂裡面才有.
所以那個對我來說.
是很重要的.
你要什麼.
而作為一個媽媽.
我都問我自己.
你要什麼.
我要的是.
我的子女一生跟隨著.
我和你分享一個.
我的WhatsApp.
其實是這樣的.
我以前.
教會.

$^{841}$之前教會服侍的同工.
發了一張相片給我.
因為我之前.
服侍那裡有十年了.
我的子女在那裡成長.
所以他們繼續都留在那裡.
一路去聚會.
一路成長.
我的同工就發了一張相片給我.
那天剛好是生日.
送個生日禮物給你.
為什麼呢.
前面這兩個.
剛好.
一個是我兒子.
一個是我女兒.
這是我兒子.
這是我女兒.
他們剛好在聚會裡.
彼此祈禱.
同工就拍了一張相片給我.
他說作為媽媽.
你莫過於最想的.
就是你的子女.
能夠一生跟隨著.
這是我.
在我裡面.
如果耶穌問.
你想你的子女要什麼.
我就一定會要.
就是一生跟隨著.
我可以為他.
籌算些什麼呢.
我以為我籌算到.
其實我根本籌算不到.
是不是.
我以為我給他.
或者我站在他旁邊.
好像很厲害.
其實.

$^{881}$有時我自己都不是很厲害.
但是如果他能夠.
一生跟隨著.
捉緊主耶穌.
走在他人生的路裡的時候.
你真的不用怕.
因為耶穌照著他.
還有一些路.
你都不能陪他走.
但是主陪他走.
最重要的一件事.
就是他一生跟隨著.
所以我很想.
去鼓勵我們在座的家長.
又或者可能.
我是公公 婆婆.
爺爺 媽媽.
都不要緊.
你的生命.
都在影響你的下一代.
你的後.
就是你的孫.
因為他們都在看著你.
很想.
對我來說.
作為媽媽.
這是我最想的.
作為兒童部的牧者.
我最想.
也是我們的孩子.
能夠在神的話語裡.
去成長.
站立得住.
去面對這個.
彎曲薄茂的世代.
以致他們很知道.
他們的價值.
不是外在人給他們的東西.
但是同樣都很想.
作為父母.

$^{921}$你們成為這個家.
很重要.
你們最多時間.
守護著你們的孩子.
神給你們一個權柄.
你可以保護你們的孩子.
遮蓋他們.
為他們禱告.
但是很重要的是.
你怎樣經歷神.
他們都看到.
但是我們一起努力.
我們一起努力.
我們家裡就是會.
做壞了的事.
我們都會分享.
然後告訴他們.
我們是這樣的.
在我們裡面.
我們覺得其實.
人來的.
我們就一起去走.
一起靠耶穌.
站立得穩.
所以我覺得.
我們裡面.
一生養我們的孩子.
跟隨主.
作為家長.
作為教會.
我們一起做.
我很想和家長一起.
我們一起做.
去抵抗這個世代.
那種衝擊.
不然他們真的站不穩.
所以對我來說.
神擺在我心裡.
就是我們很想.
祂站在神裡面.

$^{961}$站得穩.
這是我最想的.
這個就是邀請大家.
7月5日我們有一個.
無論升小一升中.
都是一個很大的挑戰.
但是我很想.
家長來的時候.
我們很想.
我們都能夠.
我們家長.
我們一起.
我們一起去面對.
我特意用一個字.
我們能夠接得住.
孩子的改變.
轉變和他的情緒.
以致我們和他們.
一起探索.
一起成長.
不是你為他.
籌算什麼.
而是他和你一起.
這是重要的.
更加重要的是.
有神在當中.
好,我們一起禱告.
天父我求你.
幫助我們每一個弟兄姊妹.
無論我們做父母的.
或者是我們在職場的.
在家庭的.
學生的.
或者是不同的崗位的.
但神我深信一件事.
就是你都給我們.
有不同的呼召.
在這個呼召裡面.
你也都會賜下.
你的恩典和能力.

$^{1001}$以致我們能夠.
靠著你去走下去.
但上帝被我們捉緊.
當你問我們.
要什麼的時候.
我們懂得回應.
我們要的是.
主緊緊跟隨你.
主無論如何.
都要緊緊跟隨你.
不單是我們.
我們的下一代.
我們屬靈的下一代.
或者我們屬靈的上一代.
我們都要這樣.
一起緊緊跟隨你.
去走下去.
多謝主.
我們這樣禱告.
奉耶穌基督名口.
阿們.
\newpage



\section{馬太福音 11:25-30}
\label{sec:032431}
\textbf{主裡得享安息}~~~~陳雅莉傳道
\newline
\newline
連結: \href{http://www.church.com.hk/acms/content.asp?site=cdc\&op=show\&type=product\&code=032431\&layout=sermon}{\texttt{032431}} ~~~~ 語音日期: 2026-07-12
\newline
\newline
\hyperref[sec:032357]{\small{< < < PREV SERMON < < <}}
~
\hyperref[sec:index_chronic]{\small{[返順時目]}}
~
\hyperref[sec:032432]{\small{> > > NEXT SERMON > > >}}
\newline
\newline
馬太福音 11:25-30
\newline
\begin{longtable}{cl}
\hline
\hline
章節 & 經文 (和合本修訂版)\\
\hline
11:25 & \begin{tabularx}{0.7\textwidth}{X} 那時，耶穌說：「父啊，天地的主，我感謝你！因為你把這些事向聰明智慧的人隱藏起來，而向嬰孩啟示出來。 \end{tabularx} \\ \\ \relax
11:26 & \begin{tabularx}{0.7\textwidth}{X} 父啊，是的，因為你的美意本是如此。 \end{tabularx} \\ \\ \relax
11:27 & \begin{tabularx}{0.7\textwidth}{X} 一切都是我父交給我的；除了父，沒有人知道子；除了子和子所願意啟示的人，沒有人知道父。 \end{tabularx} \\ \\ \relax
11:28 & \begin{tabularx}{0.7\textwidth}{X} 凡勞苦擔重擔的人都到我這裡來，我要使你們得安息。 \end{tabularx} \\ \\ \relax
11:29 & \begin{tabularx}{0.7\textwidth}{X} 我心裡柔和謙卑，你們當負我的軛，向我學習；這樣，你們的心靈就必得安息。 \end{tabularx} \\ \\ \relax
11:30 & \begin{tabularx}{0.7\textwidth}{X} 因為我的軛是容易的，我的擔子是輕省的。」 \end{tabularx} \\ \\
[1ex]
\hline
\hline
\end{longtable}
$^{1}$各位節目早安.
感恩.
可以一齊來聆聽神的話語了.
今天讀這段經文.
我想應該你們都有聽過.
從前在香港的時候.
這段經文經常會在大廈的外牆.
看到「凡奴婦擔重擔的人到我這裡來」.
是一段我們都很熟悉的經文.
我們都會很喜歡這段經文.
因為很安慰.
實在是否因為我們太多「奴婦重擔」呢?.
我們真是很「奴婦」.
昨晚看一場球賽.
都可能令你覺得很「奴婦」.
今早起床回來.
甚麼是「奴婦重擔」呢?.
是否真的單單照字面的意思.
就是「奴婦」的意思呢?.
我們今天一齊來再細讀這一段.
其實這一段是一個耶穌的禱告.
祂是禱告之後再發出邀請的.
如果你剛才有留意主席讀的時候.
讓我們更深刻認識這一段.
我們認為很熟悉的經文.
深信當中有更值得我們留心的地方.
我們一齊禱告.
父啊天地的主.
你知道我們有多少的「奴婦重擔」.
今日願你親自.
對每一個願意謙卑來到你的面前.
願意聆聽你話語的人說話.
懇求主.
讓我們能夠.
帶著我們的重擔來到你的面前.
你就使我們得安息.
求主幫助.
聽祈禱奉主名求.
Amen.
如果你手上有聖經.

$^{41}$我也想你翻開.
無論是紙本還是電子版的也沒問題.
不過你打開你的電子版就不要看WhatsApp.
不要滾到別處.
真的要打開你的聖經.
有時候我講道的時候會叫大家翻開聖經.
為什麼呢.
因為有時候我講的經文是要包括前文後理的.
前前後後這樣看.
你才能看到整段經文整體的意思.
所以有時候我會叫你們打開聖經.
你會更清楚看到整段經文的脈絡是怎樣的.
這段經文一開始在25節.
那時耶穌說.
「那時」是一個什麼時候呢?.
是一個什麼背景.
耶穌在一個什麼背景下講這個禱告呢?.
如果我們往上一點看.
你慢慢翻一下.
《馬太福音》11章25節.
如果我們往前一點看20節.
往前一點看.
前幾節而已.
11章20節.
裡面有講到.
「那時」.
耶穌在那些城裡行了很多義能.
因為城裡的人不肯悔改.
就責備那個城.
所以你看聖經你會看到.
21,22,23,24節就是耶穌責備那些不肯悔改的人.
接著就來到25節.
原來耶穌當時做這個禱告的時候.
祂在那些城裡行了很多義能.
但可惜很多人不肯悔改.
耶穌就是面對著一群很心硬.
很不信的人.
所以耶穌就責備他們.
看到那麼多神蹟.
看到那麼多義能.

$^{81}$你都不肯悔改.
耶穌在這個背景之下.
就去講那些責備的說話.
但是很神奇.
你有沒有覺得很奇怪.
作者這樣轉折.
明明正在責備那些不信的人.
為什麼突然間會變成.
即是他們不相信「哈利路亞」這樣?.
雖然在普遍的裡面.
很多人拒絕耶穌.
但似乎耶穌看到的不是這些.
耶穌看到的是.
父神的旨意仍然繼續成就著.
耶穌沒有因為看到不悔改的人而灰心.
相反.
祂立刻轉向父.
這些城明明是看到耶穌行了很多二能神蹟.
他們已經得到更多的光照.
本應應該快點悔改.
但竟然心硬不悔改.
照人的角度來看.
如果我們今天這樣出去傳福音.
如果我們出去傳福音.
我們報道.
全城的人都心硬.
全城的人都不悔改.
我們可能都會有點灰心.
有點氣餒.
有點不開心 覺得自己被拒絕.
彷彿最有機會明白的人竟然拒絕.
最有宗教背景的人.
反而看不見.
耶穌還發出頌讚.
很明顯.
耶穌沒有停留在那個灰心失望裡面.
耶穌著眼的.
不只是看到很多人悔改不信.
祂就覺得神的工作失敗.
不是.

$^{121}$耶穌很清楚.
福音被拒絕.
不代表神的工作失敗.
耶穌很清楚知道自己來地上是做什麼的.
祂就只管做好祂的工作.
只管繼續在地上.
要顯明祂自己.
要顯明上帝.
雖然有人拒絕福音.
但耶穌沒有因此氣餒.
甚至不僅不氣餒.
耶穌繼續用很溫柔.
很溫柔的心.
去邀請人.
你到我這裡來.
祂繼續不放棄地呼召人.
你來吧.
這是耶穌.
耶穌仍然畏著.
雖然很多人悔改.
但神的工作沒有停止.
接著耶穌說什麼呢?.
然後祂說「因為你將這些事.
向聰明智慧人隱藏起來.
向嬰孩啟示出來.
父啊是你的美意.
本是如此」.
耶穌說.
原來有一些啟示.
是會向一些人隱藏.
而向一些人就會開啟出來.
而且這是父上帝的美意.
祂在說什麼呢?.
誰是聰明智慧人呢?.
誰是嬰孩呢?.
我不知道你自認.
是否覺得自己很聰明.
很有智慧.
智商也很高.
我怎麼看都不像嬰兒.

$^{161}$聰明智慧人.
在聖經裡我們知道.
不單止是指頭腦很聰明.
很有智慧.
很懂得分辨是非這麼簡單.
祂是指一些很高學歷.
很高智商.
很高教育.
但卻是自以為很有學歷.
很有知識.
很有屬靈判斷.
還可以到處評斷別人的人.
相信耶穌說的聰明智慧人.
是指當時的法理才人.
文士.
祭司長.
縮篤律法 擁有很多知識.
但耶穌說.
祂的啟示是向這些人隱藏.
相反呢.
嬰孩是什麼人呢?.
我們知道嬰孩不是指年紀小.
不是指嬰兒.
嬰孩是指比喻那些.
自知無能為力.
很謙卑.
自知自己很需要耶穌.
所以從當時的預警來看.
很可能耶穌說的這些嬰孩.
是當時的門徒.
或者那些跟隨祂的人.
他們可能很多都是沒有學問.
在社會上沒有地位.
沒有資本.
但他們有的是什麼呢?.
他們有的只是一顆願意向耶穌打開的心.
這些人不是天然比較聰明.
也不是比較聰明.
比較做得很好.
很守規矩.

$^{201}$而是他們不像那些有很多知識.
知道很多東西的人.
那麼封閉.
嬰孩.
耶穌其實是在說一種心靈的狀態.
是一種沒有先入為主.
很願意受教.
很願意承認自己很需要.
是這種樣的人.
所以那些自以為.
什麼都知道 看得到的人.
其實他們是黑眼的.
那些謙卑像嬰孩.
以為自己什麼都不知道.
反而會得著耶穌的亮光.
耶穌說.
祂的啟示向嬰孩顯明.
問題不是在於有多高智力.
而是驕傲還是謙卑.
一個很高學問的人.
也可以在神面前很謙卑.
好像嬰孩一樣領受主的道.
一個文化不高的人.
一個文化很高的人.
如果他心裡驕傲.
他一樣得不到神的啟示.
神的事.
是向自以為有聰明的人隱藏.
卻是向謙卑的嬰孩啟示.
這裡耶穌指出.
原來要進入神的國.
不是靠聰明的.
不是靠你的智慧的.
要進入神的國.
是靠著你單單很渴慕耶穌.
不是靠你的努力.
是靠神願意向你啟示.
不過說到這裡.
我想如果我在座問大家.
你有沒有覺得自己很驕傲?.

$^{241}$或者你渴不渴慕做一個謙卑的.
如嬰孩一樣那麼謙卑?.
我相信我們每一個人都渴慕的.
我想我們沒有人會說.
我很想驕傲.
我很想自高自大.
我就是喜歡這樣.
我相信我們在座的.
都渴慕像嬰孩.
很想很謙卑地跟從神.
很渴慕.
但是你會發現.
即使很渴慕.
原來都不容易做到.
有時不是我們不知道要謙卑.
或者不是不懂得怎麼謙卑.
而是我們有些力不從心.
有時候不是我們故意要不像嬰孩.
我們不是故意.
而是我們習慣性.
我們習慣了很自然就流露那個.
很舊,很老的我.
那個舊我,老我.
我們很自然就讓它出來.
當我們放鬆它就出來了.
我們稍微不留意.
我們就習慣性讓它又出來.
回想我們初信主的年日.
我們當時可能比較單純.
比較單純地認識主,愛主.
聖經是什麼.
不知道,都沒看過.
很多不明白.
所以我們很願意被人修正.
很願意謙卑聽人意見.
主耶穌今日仍然向.
這些願意謙卑如嬰孩的人去呼召.
祂仍然願意向你啟示祂自己.
在今日這個世代.
雖然有很多自以為是的人拒絕耶穌.

$^{281}$但耶穌看到是仍然有嬰孩.
願意謙卑承認需要耶穌.
願意跟隨耶穌.
耶穌接著從向天賦予的一個頌讚禱告.
轉向一個很重要的宣告.
你會發現又轉折了.
轉折到一個宣告裡面.
祂宣告了這件很重要的事.
耶穌在第27節宣告.
一切都是我父交給我的.
除了父沒有人知道子.
除了子和子願意啟示的沒有人知道父.
這句看完也不知道祂在說什麼.
其實這是一句很深的基督論說話.
如果要在讀神學裡面.
單單讀基督論這科研究這句也要很長時間.
不過簡單一點.
耶穌其實是更盡心說了一件很重要的事.
祂說神的美意本來就是向嬰孩啟示.
那怎樣才可以做到嬰孩.
祂其實是帶領著我們.
既然我告訴你.
你要如謙卑如嬰孩.
你才聽到我的啟示.
那現在我告訴你.
你怎樣可以做到嬰孩.
祂說.
出路在祂那裡.
單單只有耶穌.
沒有人能夠藉著任何其他的方法認識父.
沒有.
單單要透過子.
單單透過耶穌基督這個人子.
唯有是祂.
你要來祂那裡.
你要來認識耶穌.
你要單單倚靠耶穌.
你才有機會認識父.
因為唯有只能夠啟示給人知道父是誰.
如果不跟從耶穌.

$^{321}$你是沒有辦法認識這位父.
所以耶穌其實是從一個頌讚禱告.
轉向了發出對眾人的邀請.
「我就是那位人子」.
「你就是能夠透過我認識我的父」.
「我是誰呢?」.
「我是與父如此親近的子」.
接下來.
說完這麼崇高的基督論.
說了祂是誰.
這麼有權柄.
這麼有威嚴.
不是應該叫人來膜拜祂.
敬拜祂.
尊祂為大嗎?.
但是耶穌再一次的轉向.
祂竟然再一次轉向說.
「凡勞苦擔重擔的人.
可以到我這裡來」.
每一次讀到這段經文.
由二十七節轉到二十八節的時候.
都有一種很深的感受.
可能我們平時習慣直讀.
就從二十八節開始.
很少讀祂前面那一句.
但是你再想深一層.
我們這位耶穌是一位怎樣的耶穌呢?.
祂是一位無所不知.
無處不在,無所不能的神.
祂是與父有非常異常獨特的關係的神.
祂擁有何等大的權柄.
以致祂可以呼喚任何人來膜拜祂.
敬拜祂,尊崇祂.
但是祂沒有這樣做.
祂竟然說.
「凡勞苦擔重擔的人.
可以到我這裡來」.
因為超乎萬有的上帝.
祂知道一切的上帝.
祂竟然告訴你.

$^{361}$祂看到你的.
祂不是只是高高在上的上帝.
祂是與人親近的神.
祂不是只是看到祂自己有多尊榮.
祂很高超.
沒有我啟示.
你不知道父.
祂不是這樣.
祂乃是說.
「我看到你」.
「我看到你勞苦擔重擔」.
祂不是遙不可及.
祂是主動邀請人來親近祂.
不僅如此.
祂願意與你同行.
耶穌在這裡說.
「凡勞苦擔重擔的人」.
哪一個是勞苦擔重擔的人呢?.
是不是今天的你和我呢?.
有時候我們常常都會很自然地覺得.
勞苦擔重擔的人.
就是指生活上很艱難.
貧窮.
生活很艱難.
工作很勞苦.
或者周身疾病.
當然都是的.
從廣義來看.
的確勞苦擔重擔都指著這一班人.
指著在罪裡面掙扎著.
未認識他們的人.
也指著面對著死亡.
疾病,痛苦.
生活壓力.
破碎等等.
這些疲乏的人.
當你感覺到自己在生活上已經撐不住了.
背負不了.
你不得安息.
耶穌叫你來.

$^{401}$祂不只是在邀請那些.
在社會最低階層的人.
祂是在邀請那些.
所有感覺到沉重的人.
不過.
這是一個廣義.
如果我們從狹窄一點的意思來看.
從狹義來看.
耶穌指著的勞苦擔重擔的人是甚麼人呢?.
剛才記得我.
我們看到有聰明智慧人.
和有嬰孩.
似乎耶穌這裡更加特別指著.
勞苦擔重擔的人.
是那些聰明智慧人.
聰明智慧人.
他們拒絕耶穌.
他們屬靈核眼.
如果我們再看下文.
我們剛才說十一章.
我們今天讀到三十節.
如果我們再看下去十二章.
你再看下去十二章.
你會發現.
我們大概能夠再揣測到.
甚麼是勞苦擔重擔的人.
去到十二章.
法利賽人控告耶穌說.
他在安息日做不合法的事.
他的門徒.
去到第十節.
就再說他要控告耶穌.
因為安息日治病不合法.
你會發現.
耶穌一直在說的勞苦擔重擔的人.
似乎不是真的生活上那種勞苦重擔.
他是在說.
那些死守律法.
緊緊守住安息日的人.
其實不得安息.

$^{441}$那些人外表很勁竭.
但其實心靈很沉重.
表面守足全律法.
安息日甚麼都不可以做.
但耶穌說他們勞苦重擔.
有些人自以為無可指責.
但其實不得安息.
所以似乎耶穌說的勞苦擔重擔的人.
不單只是生活上.
他不只是說那些繁瑣生活的壓力壓迫.
他更是在說繁瑣的規條.
道德的壓力.
是在說在自以不斷地追逐.
永遠覺得自己做得不夠好.
永遠被規條.
被道德壓得透不過氣的人.
那些很努力很努力想靠自己.
可以以為在神面前得稱為義的人.
結果卻越走越累.
耶穌說他們是勞苦擔重擔的人.
耶穌說勞苦擔重擔的人可以怎樣呢?.
「到我這裡來」.
耶穌沒有說凡勞苦擔重擔的人.
快點放過大家.
他沒有說勞苦擔重擔的人.
甚麼都不用做甚麼都不用管.
快點去安逸,享福,享受世界,睡覺.
不是.
耶穌說凡勞苦擔重擔的人.
「到我這裡來」.
這一份是主極為親近的恩典.
祂看到守律法守到很累的人.
祂叫他們「你可以到我這裡來」.
祂看到那些未信耶穌的人.
每一天在罪裡面的掙扎.
每一天要面對罪的捆綁 很勞苦.
「到我這裡來」.
有時候我們在信仰的路上.
都會覺得很勞苦.
很多重擔.

$^{481}$還記得幾個星期前.
黃嘉良博士來過我們中間.
他說過浪子的比喻.
他說不是浪子 是兩個兒子.
他其中一點說那個大兒子.
其實沒有在他的父神裡面享安息.
他只是和父等價交換.
他每一天很努力地做好他自己.
然後又期望父神你要獎賞我.
你要對我好.
然後他又很努力地用自己的努力.
要做好一點 再做好一點.
再做好一點 就要討神喜悅.
或者是.
有時候我們真的覺得.
再靠自己的努力做多一點.
我們會安心一點.
再靠著做很多的好行為.
讓人看到 我們會感覺良好一點.
甚至覺得自己優越感多一點.
就好像當時的法理賽人.
守律法守到足夠.
但他們卻是自以為是.
他們自以為比別人好一點.
甚至他們要求身邊的人.
和他一樣那麼勞苦重擔.
他自己手都不夠.
他要求身邊的所有人都要如他一樣.
今天我們是不是都是這個.
勞苦擔重擔的人呢?.
我們的主.
不是當我們做得很好.
當我們很多事奉撐得很漂亮,很美麗.
人人都看我很得體的時候.
祂看見我.
我們的主是.
當你勞苦重擔.
沒有人看見你撐得多辛苦的時候.
祂看見你.
主說不要再靠自己的努力了.

$^{521}$不要再靠了.
你最終只會不得安息.
祂開了一條出路給我們.
祂說.
來到我的面前.
你不用很完美.
你不用做好所有事情.
你才可以來跟隨我.
不用.
耶穌說你是可以帶著你的勞苦重擔.
來到我的面前的.
我們不需要再靠著自己.
再做好一點.
再完美一點.
再守規矩一點.
才能夠來到耶穌的面前.
不需要.
我們中間如果有人仍然覺得.
我不配跟隨耶穌.
我未做得夠好.
萬一跟隨了.
我又做不到那麼多事情怎麼辦.
不是.
耶穌是說.
你帶著你的勞苦重擔來.
祂說我要使你得安息.
什麼是安息呢?.
前幾天我有機會參加一個.
在職的青年人.
二十多歲.
剛剛大學畢業的一群弟兄姊妹的團體.
有一個年輕的弟兄問我.
他說.
明明耶穌說煩惱不負擔重擔的人.
可以到他那裡得安息.
但為什麼我到了耶穌那裡.
也不能得安息呢?.
他這樣問我.
我就叫他.
你今天聽我的說法吧.

$^{561}$他說因為每天他的工作都是很勞苦.
很沉重.
很累.
很多挑戰.
他是一個很年輕的醫生.
他在這邊做醫生.
但是很年輕.
很沉重.
每天面對著病人.
我們的耶穌.
為什麼好像叫我們來祂那裡.
就說會得安息.
但是我們卻不安息呢?.
我不知道你們有沒有這個想法.
你有沒有也覺得.
我好像到了.
但是也不安息呢?.
或者有時候我們對安息有些誤會.
我們可能會覺得.
耶穌說煩惱不負擔重擔的人.
到我這裡來.
就是使你無災無病.
無艱難無壓力無痛苦.
不用造嘆世界.
有些誤會了.
原來耶穌邀請我們進入安息.
不是進入一個不用負代價.
不需要信服的安逸.
什麼是安息?.
安息不是什麼都不做.
而是停止那些徒然的掙扎.
安息不是外面停工什麼都不用做.
而是哪怕很忙碌.
但是心靈可以很平靜.
這個是安息.
安息是你能夠真正從罪的牢獄裡面得釋放.
這個是安息.
安息就是我們植根於一個.
耶穌給我們的得救的安穩.
以致我們不再被罪控告.

$^{601}$我們良心得潔淨.
真正的安息.
是已經完成了耶穌的託付.
進入安穩裡面.
這個是安息.
所以在基督裡面的安息.
不是懶惰.
偷懶什麼都不用做.
而是不再靠你自己去證明你自己.
不需要了.
你不需要再靠自己去證明我有多有價值.
我多珍貴.
我多聰明.
我多能幫忙.
不需要了.
你到耶穌那裡.
這個安息是一份員工.
一份安穩.
一份和上帝和好的關係.
耶穌未講完的.
這個邀請.
耶穌邀請凡奴婦參眾膽人可以到我這裡來.
我就使你得安息.
耶穌沒有停在這裡.
耶穌接著再講了一些東西.
祂說:我心裡面柔和謙卑.
你們要負我的額.
向我學習.
那你們的心靈就必得安息了.
因為我的額是容易的.
我的擔子是輕散的.
原來耶穌不是說.
你凡奴婦擔重擔.
什麼都不用揹.
什麼都不用負.
什麼都不用背.
不是的.
耶穌說:你背我的額.
你擔我的擔子.
原來真正進入.

$^{641}$耶穌所賜的安息.
是要到耶穌那裡來.
放下你的奴婦重擔.
但同時我們是要背負主耶穌的額.
學耶穌的樣式.
額是什麼來的呢?.
耶穌的額.
這個額是什麼來的呢?.
額是一個工具.
能夠在古代的時候.
使兩隻動物一起向前走.
有同一個方向.
所以你有聽過「同父一額」.
他們兩個一起背負著那個額.
就能夠統一.
能夠分擔重量.
能夠提高效率.
所以額不是一個刑具.
其實它是什麼來的呢?.
額是一個受管束.
是受管束的.
是能夠受管束.
而有方向地工作.
是一個這樣的工具.
當然.
耶穌這裡所說的「額」.
並不是真的叫我們背負著這塊木頭.
祂說「你負我一額」.
並不是真的叫你背負著這塊木頭.
是一個象徵.
象徵你願意服在耶穌的筒子.
在祂的約束.
信服在祂的帶領.
信服在祂給你的方向.
然後去服侍.
這是「負耶穌的額」.
耶穌並不是真的叫我們.
背負著一件器具.
而是其實祂在這裡邀請人來.
你來做我的門徒.

$^{681}$你接受我帶領你的方式.
你接受我的教導.
你按我所教你的生活.
這就是我們福音最核心.
和最奇妙的地方.
世界的邏輯告訴我們什麼呢?.
世界的邏輯告訴我們.
只要有東西背著.
就不得安息.
所以要扔掉它.
「負額」是不得安息的.
但耶穌的邏輯不是這樣.
耶穌的邏輯是說.
你背負了錯的「額」.
就會使人沉重.
但如果你背負我的「額」.
你是會得安息的.
這是耶穌的邏輯.
這也是我們福音奇妙的地方.
不是全部拿走.
什麼都不用背.
而是你之前背錯了.
背錯了.
你現在背對了.
背耶穌那個.
你就得安息了.
今天.
這句話其實是.
不論對我們所有在座已經相信耶穌的.
還是你還未相信耶穌的說話.
耶穌是向每一個人親自去呼召.
因為耶穌看到.
如果你還未相信耶穌.
你每一天都在最終掙扎.
你背負了沉重.
人給你的眼光.
人給你的壓力.
你要努力用很多的好行為.
證明自己做得好一點.
這個是勞苦重擔.

$^{721}$信了主的我們都是.
我們仍然拿著那些律法.
拿著那些道德標準,宗教禮儀.
我們都是很勞苦.
仍然是未真正學會.
怎樣叫「倚靠主」.
我們只會越走越累.
所以耶穌不是說你逃出一個權力.
沒有自由.
沒有重擔.
沒有人管你.
你就自由自在.
很安息.
不是.
耶穌是說.
你換過一個額.
重綑綁你的額.
換成.
有一位柔和謙卑的主的額.
耶穌的額之所以容易.
不是因為不需要付代價.
耶穌的額之所以容易.
是因為我們這位主是柔和謙卑.
祂不是只懂得嚴厲管束的主.
祂是賜恩典給我們與祂同行的主.
所以祂的額容易是因為祂柔和謙卑.
我們的擔子之所以輕散.
不是因為沒有擔子任意放蕩.
所以輕散.
我們的擔子輕散.
是因為我們有這位柔和謙卑的主.
與我們同赴一額.
是祂與我們同行.
是祂與我們一起背這個擔子.
所以可以輕散.
所以跟從主從來都不是不需要付代價.
但我們的主不是將一堆命令.
一堆規條丟給我們.
然後我們就搜飽它.
不是.

$^{761}$我們的主是用祂自己柔和謙卑的榜樣.
做給你看.
你來學我.
你來跟從我學習.
祂沒有任由我們自生自滅.
祂只是將祂的管治.
你要的順服講給你聽.
你跟從祂.
你就很輕散.
所以安息不是放縱.
不是不需要悔改.
不是不需要跟從.
而是不需要再靠自己努力掙扎.
不需要再靠很多外表的宗教表演.
不需要再你自己一個人苦撐.
祂讓你忍到.
祂讓你安慰.
祂會承托你的生命.
這段經文我們讀的時候.
很多時候都會覺得很安慰.
是不是?好像我剛才所說.
是的,你這樣讀固然是沒有錯的.
「法勞負擔重等人,到我這裡來就得安息」.
心裡柔和謙卑.
真是好了.
很輕散.
很容易.
不是不好.
這段經文如果你敢讀.
不過你錯過了很多很深層的意義.
耶穌帶出的訊息.
是讓我們生命真的得轉化.
安息不是靠方法的.
不是靠很多安靜,靜坐,靈修,默想.
那些都可以幫助我們進入安息.
但真正的安息.
不是方法.
而是耶穌.
而是耶穌你要真正認識祂.
這段經文都提醒我們.

$^{801}$有時我們背負我們自己的厄.
背得很重的時候.
我們很容易都會令我們身邊的人.
都一樣背負這麼重的厄.
我們會不其然.
這樣要求自己的同時.
都這樣要求別人.
就好像當時的法利賽人一樣.
好像我們剛才看到.
他們自己要守安息日.
守到有病的不准醫.
死的都不准救.
但當耶穌致病.
當時耶穌在安息日救人的時候.
他們要求耶穌都不准這樣.
他自己死守不夠.
他要求所有人都不准這樣.
今日都幫助我們.
要避免這種如昔日的律法教師.
這樣專門捉人錯處.
挑剔別人.
這種都是勞苦重擔.
經常只見到別人的錯.
這種都是勞苦重擔.
這個世代抑鬱症.
焦慮症的數字每一天都增加.
我們不需要一個宗教.
令人更焦慮和更抑鬱.
我們要帶人真正進入耶穌的安息.
我曾經在一間幼稚園做過.
我在幼稚園做的時候.
那些小朋友很有趣.
在幼稚園生四歲.
有些再小一點.
小朋友之間總是有很多紛爭.
常常都會說.
他欺負我 他打我.
他搶了我的玩具.
我聽到在幼稚園裡最常聽到的一句.
你知道是什麼嗎.

$^{841}$就是說.
這樣不公平.
英文常常都聽到.
It's not fair.
很有趣.
歲幾兩歲已經懂得說It's not fair.
當然很可愛的小朋友.
就是從他自己的角度出發.
他看到的東西很少.
他的世界很小.
他從自己很簡單的角度出發.
就下定論了.
老師 不公平.
他就會跟我說.
我不明白為什麼那個小朋友受了傷.
就要其他人扶他.
不公平.
為什麼他受了傷.
還要其他人幫他拿書包.
還要其他人幫他拿袋子.
不公平.
哦 他扭傷了腳.
為什麼可以搭電梯.
不用走樓梯.
不公平.
有些小朋友比較膽小.
比較害羞.
常常都可以坐在我旁邊.
坐在老師旁邊.
他就會說.
為什麼他常常都可以坐在你旁邊.
不公平.
是的.
這個是小朋友.
他們按照很天然.
按照他們所認識的.
他們所知道的.
他就很直接去表達了.
很正常.
這個正常度不得了.

$^{881}$因為這個是最自在.
最自然.
沒有掩飾的他.
他直接就會來投訴了.
不過我覺得問題的關鍵.
不是在於他直接的投訴.
他怎樣看這件事.
他怎樣下判斷.
問題是在於.
如果他只是覺得他自己永遠是對的.
如果他只是單單從自己的角度出發.
不願意聽老師的教導.
不願意看老師再帶他看多一層的意義.
這個才是問題的關鍵.
我很感恩.
他們都很單純.
很可愛.
當我解釋給他們聽.
不是這樣的.
因為他痛.
其實不是不公平的.
解釋給他們聽.
因為他比較害羞.
你比較勇敢.
不是很公平吧.
你這麼勇敢.
你說話大聲一點.
他說話小聲一點.
不是很公平吧.
上帝讓你這麼大聲說話.
他慢慢會發現.
慢慢打開他的耳朵看到.
原來他看到的事物很單一.
我慢慢發現.
小朋友不單沒有覺得不公平.
他們開始願意聽.
願意學.
主動馬上去搶.
幫忙拿別人的書包.
我教我教.

$^{921}$我幫他拿.
會主動扶著那個人.
讓老師旁邊的位置.
給那個害羞的小朋友坐.
是.
從同學裡接納.
包容.
輔助裡面.
那些比較弱的.
比較沒那麼厲害的.
比較軟弱的.
加上訓練.
加上老師的教導.
慢慢.
可以那些有特殊需要的小朋友.
都能夠融入大家的當中.
耶穌的公平.
和我和你看的公平很不同.
我們的公平是.
要人人平等.
你比較弱嗎?.
我為什麼要遷就你?.
這個是我們的公平.
你表達能力沒那麼厲害.
說話比較得罪人嗎?.
為什麼我要體諒你?.
這樣對我不是很不公平嗎?.
為什麼.
為什麼你想事情慢一點.
我就要幫你做你的那一份?.
這樣不是很不公平嗎?.
這個是我們的公平.
聖經的公平是.
耶穌說在基督的身體裡面.
是要將加倍的體面.
給欠缺的肢體.
這個是耶穌的公平.
如果我們都學會.
體會耶穌的這份心信.
如果我們都學會.

$^{961}$背主的額.
我想我們不單止會自己輕散一點.
也會令我們身邊的人輕散不少.
勞苦擔重擔.
不是因為我們要刻意.
我知道我們沒有刻意.
我們都很想輕散.
但你要到耶穌那裡來.
你要放下你的成見.
放下你的偏見.
重新來學主的樣式.
背他的額.
重新信服.
信服在耶穌的教導底下.
你會發現你會得真正的安息.
勞苦重擔.
不是耶穌將你所有的問題拿走.
我們能夠得安息.
不是因為沒有問題.
所以很安息.
得安息是即使問題沒有離開.
但當我倚靠耶穌的時候.
耶穌使我心裡平靜.
耶穌賜我恩典.
賜我力量面對.
問題仍然在.
累不累呢?.
每天仍然都很累.
但心靈輕不輕散呢?.
很輕散.
我再邀請你們.
今天這段經文.
再一次邀請我們每一個.
如果這一刻你覺得你自己勞苦重擔.
再一次來到主耶穌的面前.
放下你的重擔.
不再靠你自己了.
你願意背起耶穌的教導.
背起耶穌的呃.
學習耶穌的養息.

$^{1001}$願我們每一個都真正得安息.
我們一起在座位上有安靜的時間.
如果這一刻.
你疲累了.
無論生命大小的事.
無論是宗教的規條.
無論是與罪的掙扎.
無論是面對生離死別.
每一次的失去.
每一次的患難.
工作的壓力.
將我們壓垮.
願你今天再一次來到主耶穌的面前.
祂要使你得安息.
壓傷的蘆葦.
祂不截斷.
傷殘的燈火.
祂不催滅.
耶穌看到你過去.
因為靠自己的努力.
想去勝過罪的捆綁.
祂看到你的掙扎.
耶穌也聽到.
因為這個世界沒辦法合乎你的標準的時候.
你所發出的投訴.
耶穌也聽到.
耶穌也知道你已經很努力.
但卻在很多事上仍然未能如願以上.
祂再一次對你說.
犯勞苦擔重擔的人可以到我這裡來.
我要使你們得安息.
這位心裡面柔和謙卑的主.
邀請你來負祂的額.
學祂的樣式.
跟從祂.
心靈就必得安息.
親愛的主耶穌.
向每一位願意如嬰孩般謙卑的人.
每一位自認已經擔緊勞苦重擔的人.
你向他們啟示你自己.

$^{1041}$多謝主.
多謝主你不是等我們做到完美.
你才接納我們.
而你是打開雙臂.
懷抱我們每一個勞苦重擔的人.
願我們來到這位.
心裡柔和謙卑的主的面前.
主讓我們和你說.
我們願意學你的樣式.
求主你叫我們輕散.
叫我們喜樂地跟隨你.
讓我們能夠在你裡面.
得真正的安息.
奉耶穌的名字祈求.
阿們.
\newpage



\section{俄巴底亞書 1:10}
\label{sec:032432}
\textbf{自作孽,滿招損}~~~~黃紹權牧師
\newline
\newline
連結: \href{http://www.church.com.hk/acms/content.asp?site=cdc\&op=show\&type=product\&code=032432\&layout=sermon}{\texttt{032432}} ~~~~ 語音日期: 2026-07-12
\newline
\newline
\hyperref[sec:032431]{\small{< < < PREV SERMON < < <}}
~
\hyperref[sec:index_chronic]{\small{[返順時目]}}
~
\hyperref[sec:032433]{\small{> > > NEXT SERMON > > >}}
\newline
\newline
俄巴底亞書 1:10
\newline
\begin{longtable}{cl}
\hline
\hline
章節 & 經文 (和合本修訂版)\\
\hline
1:10 & \begin{tabularx}{0.7\textwidth}{X} 因你向兄弟雅各施暴，你必蒙羞，永被剪除。 \end{tabularx} \\ \\
[1ex]
\hline
\hline
\end{longtable}
$^{1}$好,弟子姐妹,讓我們在聆聽主的說話之前.
我們再次同心低頭禱告仰望愛主愛的一神.
父神多謝你讓我們有生命有氣息.
也讓我們活在一個非常漂亮的國度裡面.
讓我們知道在運行的當中.
我們應當怎樣在上主裡面.
踏著重棧來活在其中.
但頑固在今天我們仍活在地上的時候.
上主你讓我們看見這個世界是充滿罪惡.
人性的醜惡,人性的詭詐.
甚至在親屬的之間也都能夠.
顯出人怎樣互相的攻姦彼此.
好在我們在這個世界.
見到一個這麼恐怖的世界的時候.
上帝你讓我們促成我們更深的思想.
究竟我們在上帝的倉庫裡面從來有沒有去感恩上主你的賦予呢?.
頑固今天的世界的每一個人.
為著自己的喜好,為著自己的利益而奔走.
上主求你幫助我們,我們一班屬你的子民.
讓我們每一個都能夠真真正正能夠活出有上主你的生命樣式.
讓我們體貼上帝你的心腸.
好想我們能夠在這個地上.
雖然按照你的意願來活.
並不是容易一件的事,甚至可能會受到遭害.
但是求主你讓我們明白到.
只要我們活出一個上主你喜悅的生命.
主你是最大的喜樂.
讓我們能夠在其中以上帝你的喜悅為我們中心.
走在這條屬靈的人生的路上.
讓今天的講道,無論聽的,講的.
生命也能夠得到更深化.
求你的帶領我們同心的禱告.
奉靠主耶穌基督得勝的名字而求.
阿們.
好了,今天我們會回來看俄巴底亞書.
我們今天只是看一節的經文,就是第十節的經文.
第十節進入到俄巴底亞書的第十八節的地步.
其實我們可以看到這一部分.
可以說是俄巴底亞對於以東的介紹是非常仔細的.
所以當我們去看這一部分的時候.

$^{41}$我們一定要專注看一下.
究竟俄巴底亞先知將耶和華神的說話記載下來.
上帝是怎樣看這個近親的國家,以東國.
首先大家要了解一個觀念.
國不一定是我們所想的一個國家的行政運作.
當時很多國家的運作是一個我們叫做tribal country.
就是一個族群的角度來的.
是每一個民族為核心的一個角度.
現在的國家就不是了.
尤其是你看一下加拿大,美國,英國.
現在你去英國不像去英國的.
你加拿大不知道加拿大是什麼.
去美國你更加不知道美國是什麼.
你只是靠一面旗.
所以大家有時候聽一下新聞.
那些政治人物就說我忠於我國家之旗.
我與我國家之旗站在一起.
你就會明白是什麼.
甚至中國也是一樣的.
現在中國搞了一樣是搞到這樣的情況.
中華民族搞了一個出來.
哇,不知多少族,五十多族,百多個民族.
我不知道了.
這種的觀念和以前是很不同的.
我不是說誰對誰錯.
我完全沒有任何的批判性在裡面.
不過當我們看回在古近東的世界.
特別是俄巴底亞書裡面所記載的.
上帝對於每一個民族,國度.
他們的記載都有一個很透徹的思想.
並且能夠準確地看到這幫人.
他們的問題以及他們的問題所衍生的結果.
而上帝就在一個這樣的歷史的世代環境之中去運作.
所以今天我會和大家一起去思想一個題目.
叫做「自作孽滿招損」.
這兩句話是很文雅的.
因為這兩句話是出自於兩篇不同的中國的古典.
都同樣在這個中國古典上書.
所以大家知不知道什麼叫上書?.
你們牧師以前是讀文學的.

$^{81}$所以對這些詩書禮樂這些東西很熟悉的.
書其實在中國文學裡面往往就是指道上書.
上書是非常難讀的東西.
當年考高級程度會考的時候.
如果出上書我老師說你還是不要答了.
是很難答的.
因為它的字眼很簡潔.
簡潔到你根本接近不明白.
什麼叫做「自作孽滿招損」呢?.
「自作孽」這句話是出自於上書的太甲書.
裡面有句這樣的說話說.
「天作孽猶可為,自作孽不可活」.
相信大家都聽過了.
何謂「自作孽」?.
「自作孽」是說一些壞事所做出來的結果.
其實種種都是因由你自己去主動去做一些惡事.
上天如果降災.
所謂「孽」這個字眼.
「孽」不一定是一些什麼惡事.
可能上帝是來懲罰,警戒.
用中國傳統的思想叫做「天」來去管理的時候.
降下天災,疾病.
降下人間.
人還可以有個理由用方法去逃避或者去減緩.
但是如果自己做出來的結果是一件壞事的話.
那你就不得逃過責任.
這個所謂「自作孽」.
至於「滿招損」是出自於上書另一卷叫做《大禹末》.
這篇記載裡面有句這樣的說話.
它說「滿招損謙受益,時乃天道」.
很文雅,很深奧,很難讀的,我明白的.
但是簡單一點來說.
「滿招損」的「滿」的意思就是一個很自大,自滿的人.
是很容易去為自己帶來損害.
而「謙受益」就是說一個人凡事不是吹噓的.
就算你是這個世界上最厲害的人.
PhD.
你都不會以你自己以為在這方面的專家而去驕傲起來.
一個謙卑的人.
他反而是很多時候受著他的謙卑和他的學識.

$^{121}$而帶來一些的逆轉.
所以我很感恩我們教會很多專家.
有藥劑的,有工程的,有電腦的,有財經的,有牙的.
這些全部都是專業的.
但是我們個個都沒有表現出一種「我很厲害的」.
不是,雖然你很厲害.
但是你就用你的專業和上帝給你的恩賜.
來在你的行業裡面去發揮.
我們反而可能在我們的環境裡面.
去不單止我們自己得到好處.
得到人欣賞.
並且上帝是用得著我們.
當一個人如果他是自己喜歡做壞事.
或者很自滿在這個世界上驕傲地去活著的時候.
很多時候就會造成很大的惡果.
而這些惡果往往是在其他人的身上.
好像他因為他的作惡而令自己得益對方受損.
但是從俄巴底亞書的第十節.
就告訴我們知道上帝有一個很重要的原則在這個世界上.
原來在這裡悄悄地運作.
雖然未必我們會即時見到那個果效或者結果.
但是上帝在他的運作裡面.
卻使得這個滿招損自作孽的人.
原來接受他自己所做出的惡果的結局.
而這個正正是上帝的運作在這個世界裡面.
我們不能排除的事實.
但是這個做法那個演繹是什麼呢?.
是不是上帝只是在這個世界上希望警惡懲奸是這樣呢?.
還是有另外更深層的意義呢?.
在前文如果我們上個禮拜去看完的.
其實從第一節到第九節裡面.
我們已經看到上帝即使宣告向以東去作出一連串的懲罰.
又說要將他在他最驕傲以為自己最強盛的時候.
將他拉倒下來.
甚至在眾國之間讓他能夠被藐視.
上帝這樣的做法不是警惡懲奸這麼簡單的一件事.
而是反映上帝對於這個國家或者這個民族國以東.
以及他的兄弟國,雅各或者我們稱之為猶大國.
其實都有一個很另一層的深層意思.
是讓他們同時從罪當中去反思走出來.

$^{161}$這個的動機或者上帝的宗旨.
是反而我們作為基督徒要去專注去看的.
而這個的題目在今天從第十節往後.
我們甚至到最後第二十一節的時候.
我們就會看到.
是不是二十一節我有沒有記錯?.
是呀二十一節.
我們就會看到上帝在整個過程裡面.
他的心路歷程是怎樣.
俄巴底亞書從我的《釋經》的全釋裡面.
我想看到的不是一般傳道人所看的.
你對懲罰以東人.
伸張公義為臣子民報仇.
這個絕對不是臣的看法.
臣的心思是要讓眾人得救.
他用盡各樣方法來將他們挽回.
其中懲罰,執行,訓導,指教.
就是以挽回為目的.
這個觀念和在《新約聖經》.
《迦拉太書》談論到我們作為臣子民在這個世界上.
要去作挽回的工作是無獨有偶地相同的.
只不過我們的程度和上帝的程度.
可能有某地方我們是適可而止.
上帝卻是做到盡給我們去看.
我們一起去看看.
今天我們看單純第十節.
從第十節裡面我們會看到.
上帝有兩樣東西對於以東是看得非常之準確的.
一會兒詳細去看哪兩方面.
當然講到已經分開兩部分.
我重複講一下少少前文所講的東西.
在前文如果大家還記得.
是從第五節至第九節裡面.
我們講到的就是.
當一個人有錯就要去改.
要去承認自己的錯.
透過認罪悔改.
來回歸上帝的期望.
很多人看認罪悔改.
其實只不過是一種難堪的感受.

$^{201}$但是請大家弟兄姊妹留意.
在聖經裡面認罪悔改.
尤其是悔改這個字眼.
在聖經新約裡面.
往往都和感覺是沒有關係的.
用英文所講就是.
Repentance is not just about emotion.
很多時候我們以為.
有人去信主了.
要使他在呼召的時候.
舉手認罪悔改.
讓他痛哭流涕.
發洩了他心裡面那種憂傷被罪的困鎖之後.
他就真的信耶穌了.
弟兄姊妹.
感性上我們可以流露.
上帝沒有否定過這件事.
但是如果你的悔改只停留在感性層面上.
是一件非常危險的事情.
因為感性只不過是一個發洩.
不是一個人生的取向的改變.
真正的悔改在聖經裡面所指的是.
整個人的思想以及靈性的一種轉移.
而這個轉移是180度的轉移.
叫那個人立定心意.
再不會重蹈覆轍.
這種就是真正的悔改.
不是只是在報道會裡面.
哭著說.
哎呀 我好慘呀 我曾經做過甚麼壞事.
得罪過哪個人.
哎呀 我傷害過哪個人.
我家人我怎樣傷害過他.
我好痛哭後悔.
求神讓我能夠在這個時候釋放.
沒錯 你可以釋放.
但是你的人生取向是否跟著下去去改變呢.
我相信如果讓你得罪過的人.
傷害過的人.
他更大的感恩是看見你的改變.

$^{241}$你的傷害在他曾經身上做過的.
你以後立志不再做.
那個人才是最大的喜樂.
所以換句話來說.
你要洗人家.
譬如我自己做過一些錯事.
以前我分享過給大家.
以前有一個人.
他整天都在取笑我家.
取笑我有一次在放學的時候.
向著這個阿伯說了一句粗口.
接著被我媽捉了一輪.
罵了一輪.
還要真心下去跟這個阿伯.
承認自己錯.
我好記得當時我媽媽.
在還沒帶我下去之前.
她說了一番話.
她說你知不知道你自己錯了甚麼.
我說錯了甚麼.
是不是我說了粗口罵阿伯.
我媽說不是說粗口罵阿伯.
是錯.
這個阿伯其實應該要用多些粗口罵.
真是壞人.
不過你錯的是錯在收家.
這個不是我教你的.
你是我家裡的兒子.
是代表著我家裡.
你的行事為人.
是一個家庭裡的榮耀.
同樣作為弟兄姊妹.
我們屬神的子民.
我們在這個地上活著.
我們在做上帝的榮耀的代表.
你的行事為人.
做得不好的時候.
是羞辱上帝.
當你做得好的時候.
別人看到的時候.

$^{281}$是榮耀上帝.
所以當我真正明白這一點的時候.
我媽帶我下去跟她道歉的時候.
自那一刻之後.
我沒有跟這個人再說過一句粗口.
你說牧師你之後有沒有說粗口.
我當然有說粗口.
我不會否定的.
但是我不會再用這樣的情況.
很衝動地去說粗口.
你說說粗口小事而已.
弟兄姊妹.
不是的.
我家門受羞辱.
這個是我的榮辱.
我家庭的榮辱.
我們同樣在上帝裡面.
有沒有認真地悔改.
不會再重蹈覆轍.
如果你真誠地改變.
那個被你罵過的人.
他就會覺得.
這個人也有得救.
所以當我們作為基督徒.
看到上帝怎樣懲罰一些惡人.
在這個世界的上面.
為我們伸冤.
因為這個人曾經怎樣惡待過我.
反過來我們要去問上帝.
上帝.
我的囚敵被你去作出一個警告.
甚至去作出管教.
在這個時候.
其實這個是不是同樣對我的一個提醒.
和一個令我靈性更新的機會呢.
我有沒有在這個地方的上面.
能夠更加仔細去明白你的心思.
而活在這個地上.
使得我成為上帝你在這個地上.
能夠見證出你的生命給人看.

$^{321}$的一個媒介呢.
這個是上文我們所看的.
好了來到今天.
來到這段聖經.
第十節.
第十節不是很長.
數起來好像.
中文都沒有多少個字.
我想大約十五個字.
十四個字左右.
十五.
三六一十二.
六一十二.
十三十四.
十五十六.
十六個字.
十六個字和合本.
說一堂道.
牧師你是不是很忙.
沒空.
不是.
這句話.
如果當我們再重新仔細看的時候.
它裡面所說的話.
其實透露了.
上帝一方面.
看出.
以東人或者以素.
他的本性上面的核心問題.
癥結在哪裡.
另外一樣東西.
就是他從時間看到這個性情.
以素或者以東.
他這個性情所招致的下場是怎樣.
當上帝看到這兩點的時候.
他就不想人走到這個結局.
即或以東以素人是.
曾經欺負過甚至陷害過猶大國.
上帝都想他走出來.
我們看看上帝這份心腸是怎樣.

$^{361}$第一件事我們看看第十節.
應該怎樣去翻譯這一句話.
和合本這樣說.
他說因你向兄弟雅各行強暴.
愧必遮蓋你.
你也必永遠斷絕.
這句話的翻譯.
應該是翻譯如下.
由於你們.
你們我加進去的.
加進去其實是指.
以東或者以素.
由於你們殺戮.
及殘暴虐待你的兄弟雅各.
於是乎羞辱就把你們.
隱藏或者包圍著.
你們將會自招被徹底摧毀的結果.
如果你想再清楚一點.
你可以再讀多一次.
如果你再不清楚.
你可以聽聽廣道之後的錄音.
和字幕上面有詳細介紹.
這句話應該這樣翻譯.
由於你們殺戮.
及殘暴虐待你的兄弟雅各.
於是乎羞辱把你們.
隱藏或者包圍著.
你們將會自招被徹底摧毀.
我們解拆一下這句話是說些什麼.
聖詹前文我們看見.
受敵受害.
似乎神好像為神的子民.
出一口氣.
甚至連敵人以東他的城堡提曼城.
亦都將會被摧毀的時候.
但是我們同時間聖經亦都記載到.
從第六節至第八節當中.
神用上了多次以素這個尊有名詞.
來取代以東人這個族.
又同時間又多次提及到.

$^{401}$以素和雅各兩兄弟這份原有關係.
從這兩者並排擺在這裡的時候.
我們就會發現到.
上帝好像在說為神子民出氣.
但是同時間.
上帝又好像很關心這兩兄弟.
以素,雅各或者來到這一段時間.
的以東人和猶大國的人.
他們本來那份親密的兄弟情.
所以在上文我們曾經說過.
上帝面對著以東要宣佈以東他的懲罰的時候.
上帝表現出一種耿焉的狀態.
為什麼你要走到這麼折墮的地步呢.
或者以東或者廣東話也可以改變成.
上帝說:以東為什麼你要走到這麼墮落的狀態呢.
所以廣東話很好.
將這些字混合起來也可以合理的.
好像希伯來文,希涅文那樣.
所以廣東話比傳統中文複雜很多.
是一個很成熟的語言.
從這段聖經裡面我們可以看見.
上帝有一種好像很矛盾的情況.
既要處理他的罪.
又同時間有一份情在當中.
這個就是我們在聖經我經常跟大家說.
神是既有異也有愛.
兩者並排融合其中.
今天教會最大的問題只是講愛.
不講神的異.
但是一講神的異就錯了.
就認為上帝要來處理這個世界.
警惡懲奸處理惡人.
你有沒有看到原來在神的異.
和神的愛兩者是在融合其中.
俄巴拉底亞書就將這兩者混合在一起.
所以為什麼在接下來這一節裡.
我們會看見要用的字眼是用雅國.
不是用猶大國.
在前文為什麼不只是用以東而要用以蘇.
因為上帝看到的不是只是兩個國度.

$^{441}$大家在爭端這麼簡單.
而是看著怎樣將這兩兄弟.
同屬原本都是亞伯拉罕的後裔.
將他挽回過來.
所以今天活到今天我們看這個世界.
應該有一個世界觀很不同.
為什麼上帝要放以色列人.
在一個充滿軍事威嚇的中東裡面.
面對著阿拉伯人.
面對著這些伊斯蘭教的世界.
為什麼上帝要這樣.
是不是上帝只是愛上以色列.
然後要取替其他的國度.
不是.
上帝想的是拯救普世的人.
所以刻意將雅國一名去取代了猶大國.
也用了以蘇這個名來取代了以東的描述.
讓我們看到這兩兄弟的情誼.
上帝怎樣將它挽回.
我們分開兩部分去看.
第一部分我們所看的就是.
上帝怎樣看見以東或者以蘇的問題.
或者同時間上帝怎樣看見.
我們屬靈生命的問題癥結.
在第十節的上半節那裡說了給我們知道.
十節的上半節這樣說.
因你向兄弟雅國行強暴.
上帝在這裡開宗明義地將以蘇.
他的問題他的性情他的靈性的損壞嚴重程度.
將它說了出來.
用了一個字眼來描述.
是充滿行強暴.
在原文聖經裡面翻譯這個行強暴的字眼.
其實不是只是一個字.
原來是分別有兩個希伯來文字.
去指道他是行殺戮和殘暴虐待.
所以為什麼我一開始要重新和大家翻譯這節的經文.
中文很簡單行強暴講完.
但原來實際上在希伯來文裡面.
是指道他行殺戮和殘暴虐待.

$^{481}$你說牧師沒有看見.
原來在我們的聖經裡面.
第十節是貼著上面和合本的第九節的最後那兩句.
我們讀和合本的俄巴底下書第九節.
他提問人.
你的勇士必驚惶.
甚至以素山的人都被殺戮剪除.
在和合本裡面好像只是說到.
以素山的人被殺戮剪除.
但原來有一小片語.
其實是提及到.
以素人或者提曼人本來他們是充滿殺戮.
這些沒辦法.
你一定要懂原文才可以解釋到這段聖經.
如果你沒有這句話.
你就不明白為什麼是這樣有記載.
原來他們是充滿殺戮和虐待殘暴.
換言之.
雅各面對著一個親屬的家族.
或者民族國.
是面對著一種兇殘.
什麼叫殘暴虐待或者殺戮.
殺戮就很清楚.
是將那個人殺死.
殘暴虐待這個字眼.
原來不是純粹是指到一個行動上的傷害.
按照《創世記》第十六章第五節.
在那節經文裡面.
我們掀一掀《創世記》第十六章第五節.
我讀給大家聽好不好.
《創世記》第十六章第五節.
撒萊對阿伯難說,我因你受屈,我將我的侍女放在你懷中,他見自己有了孕,就小看我,耶和華在你我中間判斷..
在這裡談論到他受屈辱這個陳述就是同一個字眼.
同一個講殘暴或者虐待是同一個字眼.
昔日在亞伯拉罕的年代.
原來阿伯難他的妻子撒萊因為她的丈夫和他的侍女下嫁.
懷孕已經使撒萊的正室受屈.
受屈是什麼意思?是在心靈層面上的一種折磨.
上帝在這裡講給異動人.
包括聽俄巴底亞書或者讀俄巴底亞書的猶太人.

$^{521}$猶大國的子民,被擄的子民.
讓他們明白到這班異數人,異動人.
上帝其實老早看見他們的行為是什麼.
他不只是想搶你的財物,他是在思想上折磨你.
大家好,為什麼異動人要這樣做呢?.
為什麼昔日的異數的人,包括他的後裔要這樣做呢?.
要在思想層面上折磨猶大國子民.
他瘋了嗎?你說這句話就對了,他真的瘋了.
當一個人與人之間的關係破裂之後.
他第一個行動未必是用刀劍槍來殺死你.
首先在你思想層面上折磨你.
這種關係在創世記從開始我們就已經看見了.
當我們去到創世記的第四章裡面.
看見當人進入墮落在罪的捆綁之下.
上帝給第一對兄弟的關係是誰呢?.
該人與阿伯,阿伯是哥還是弟弟呢?.
弟弟,記住了,不要聽見是阿伯就以為他老一點.
他是弟弟來的.
兩個為甚麼要搞到最後要殺人呢?就是為了敬拜上帝的獻祭事情上.
上帝沒有看中該人的祭物.
很多弟兄姊妹就說不公平,你只是看中了阿伯那些東西.
你是不是只喜歡吃肉不喜歡吃齋?.
上帝你究竟是想怎樣的?.
為甚麼我們不調轉問題去問.
上帝是用這個機會去發掘出.
該人他生命上的問題讓他看得到呢?.
當他看到自己的罪的時候.
他有沒有向上帝認罪悔改呢?沒有.
他相反在這件事上就認為是與兄弟關係的破裂開始的爭端.
其實開始的爭端不是他與兄弟的獻祭.
是他的生命上的問題.
誰看到人的生命上有問題?上帝.
所以上帝在這個時候不看中該人的獻祭.
看中了阿伯的獻祭.
是一個示憐.
暴露出人的問題是讓人可以悔改.
這種關係後來發展到成為兇殺.
由心靈上的糟蹋.
以至於行動上的謀殺.
當上帝問到該人說你的兄弟在哪裡的時候.

$^{561}$他講了一句什麼話?.
要我用廣東話翻譯一次.
我誰知道我盯著我的兄弟嗎?.
他的眼睛其實一直都盯著他的兄弟.
因為他恨他於入骨.
就是這個俄巴底亞書所講的殘暴虐待的狀態.
當世界上第一對兄弟的關係打破之後.
我們就會連續看見聖經裡很多兄弟的關係都發生這樣的問題.
第二對就正正就是以素和雅各的關係.
兩兄弟又是在爭.
由未出世爭到出世.
由出世出埋之後到死埋後裔都在爭.
因為他們的心裡面恨對方於入骨.
是你搶了我長子的名分.
是你搶了我長子的祝福.
於是乎我既然得不到的話.
我就要大家同歸於盡.
像不像很多小朋友在家打架.
我有份玩具.
我媽媽買給我的.
弟弟見到或者妹妹見到.
為什麼我沒有呢就走來跟哥哥或者姐姐說.
拿來給我然後就將它打爛.
一起沒得玩.
但是何曾有人明白過爸爸媽媽的心.
是買一件玩具讓你兩個一起玩呢.
為什麼父母的心腸子女捉不到呢.
問多條問題.
為什麼基督徒的心捉不到.
上帝的心是怎麼想呢.
明明上帝是想以掃瓦閣是修和的.
你們就將他替做成為.
上帝在這裡做成芥蒂.
現在來到不單止是兩個兄弟的關係.
是兩個兄弟的幫國的關係.
而這一場這樣的爭奪破裂的關係.
就不斷地重複上演.
現在就來到主前的602年左右的時間.
連猶大國都被滅了.
猶大國的皇室貴族被擄到新巴比倫國裡面.

$^{601}$我們有沒有看見.
上帝在這個時候讓以東人.
可以有一個看到自己真正那份殘暴和虐待的關係呢.
而當讀到俄巴底亞書.
當先知向著他的猶大國的子民去陳述這件事.
而上帝要去審判以東國的時候.
這班猶大國的子民又有沒有去體會到兄弟守望呢.
還是該你死了.
你都吃得和米多了.
你虐待我.
今天都輪到你了.
我們是不是落在一種幸災樂禍的關係裡面呢.
還是我們去看到上帝從創世記已經有的心願.
你和你的兄弟應該是彼此守望的呢.
丁子梅 神看見我們每個人的貞潔.
很多時候我們為了自己一己的慾望.
我們甚至連親情友情都可能會放棄.
並且以一種剝削掠奪的行動.
來使自己得益處.
兄弟都無情講了.
這個是不是基督教裡面的聖經的啟示呢.
原來基督教的神卻是讓我們看到.
當你或者我在災難裡面的時候.
其實更加體會上帝怎樣去讓我們.
從我們的惡行的生命裡面跳出來.
有一次我和湯泉道坐下來.
湯泉道你家的後院都弄得很漂亮.
我就說我隔壁的房子真是一般般.
亂七八糟種東西亂來.
然後湯泉道說你想怎樣.
我說都沒什麼心機搞.
他又亂糟糟 我又亂糟糟.
原來一起亂下去.
湯泉道他就講了一句話.
為什麼你不試想想你弄得它更漂亮.
我說為什麼要這樣做.
他說當你弄得漂亮的時候.
可能會導致到他會覺得.
嘩 隔壁這麼漂亮.
我都不要撿書了 我都要弄得漂亮了.

$^{641}$我說我又沒有這樣想過.
他說你想你的房子漂亮還是想大家一起漂亮.
弟兄姊妹.
上帝給我們每個屬他的子民.
其實都是想我們重新歸回上帝的樣式.
好像上帝那麼漂亮.
上帝呼籲了你和我成為他的子民.
我們會不會在這個地上活出上帝的樣式.
吸引更多上帝要選的人.
都同樣和我們箴言鬥禮.
做得更好的基督徒.
還是我們就會說.
想不做這麼認真的基督徒.
你看那個基督徒.
我看見他上天才去賭場巴士那裡.
不知道去哪裡.
你看那個基督徒他賭錢的.
甚至我認識一個人.
基督徒來的.
經常和我媽媽打麻將.
我媽媽說你是基督徒.
為什麼你經常認我來找你打麻將.
她第一句.
我是一個不好的基督徒.
我不知道不好的基督徒.
他的意思是什麼.
是不是可以不用負責任.
都上天國.
我想他們是這樣想的.
但是弟兄姊妹.
從這段聖經裡面.
上帝讓以東人知道上帝的行動.
也給猶大國或者雅國的後裔.
給他們看見.
上帝怎樣去處理罪的時候.
我們兩者一起看到.
上帝其實進行一個挽回的工作.
讓你們箴言鬥利.
都比互相廝殺好.
因為上帝的心是什麼.

$^{681}$我再講多一次.
是想萬人都得救.
但是得救的基礎是要悔改.
上帝看準了以東人這個問題.
講給他們聽.
現在講給雅國或者猶大國的人.
透過俄巴底亞書讀到了.
知道之後怎樣呢.
以東人這個性格上.
靈性上的問題.
上帝就要講這個靈性上的問題的結果給他們知道.
講這個結果給他們知道.
應當怎樣去回轉.
所以第二部份第十節的下半節.
上帝就講給他們知道.
如果你們不面對好你的下場.
現在我要講給你聽你的下場會怎樣.
你只是作繭自縛.
所以第二個部份十節的下半節.
我給一個附題.
我們以自身的問題向自己加害.
其實我們今天痛苦困在罪裡面.
其實是因為我們自己作孽.
將自己的下場.
困在自己所作的罪孽的結果上面.
我們看第十節的下半節講什麼.
修葵必遮著你,你也必永遠斷絕.
首先聖經讀到這些陳述.
不要以為上帝要滅絕以東人.
如果要滅絕就不用講這麼多.
一二三紅綠燈就搞定他.
所以說真的現在我們處於中東打仗.
有些人經常說打仗沒有實力的.
不是的,給足時間給足機會他改.
我不敢說現在美國總統是一個信上帝的人.
但至少他這個行動.
是給足人認罪悔改的機會.
他很厲害嗎?美國總統很厲害嗎?.
他不厲害他也要認罪悔改.
他做了多少壞事?大家知不知道?.

$^{721}$不用我講了,那些最討厭美國總統的人全部都知道.
他要悔改,他有什麼身份使其他人悔改?.
這段聖經說什麼我知道.
就是叫以東被懲罰的時候用以東來提醒猶大國的人.
你都不是好的了,猶大國.
你有沒有想一下?.
你不要只用手指指著以東壞才行.
你都不是好.
所以有俄巴大教書寫出來給他們看.
以素有什麼問題?以素在這裡給我們看到.
他其實加害害他的同胞兄弟.
這個我們之後在第十二節至第十四節裡面.
我們會詳細講的.
你心急可以去看一下.
今晚整個禮拜靈修一下.
但是上帝卻講了一個結局.
當你做了這些事情之後.
那個結果是將你厚厚地包住.
所以剛才如果大家還記得我的翻譯是什麼.
下半節是說.
於是乎羞辱把你們隱藏或者包圍著.
使你們將會在自招的這些罪的結果裡面徹底摧毀.
是一個警告來的.
然後我都用回我媽教我的那句話.
我媽媽在我很小的時候.
我很頑皮.
打完之後.
尤其是有一次.
講多一點你們喜歡聽牧師的童年事情.
我小時候我媽媽知道我讀書需要有人幫忙.
她不懂.
請了一個補習老師回來.
我的補習老師是禮拜六來的.
一個很好的補習老師.
女孩子來的.
讀完大學.
她通常就放了學.
在大學放了學.
上班之後她來補習.
我禮拜六上午上學.

$^{761}$長短週.
她通常在我吃完飯之後.
一兩點就來補習.
補習到三至四點左右.
兩個小時左右.
逢禮拜六我家裡.
我媽媽就會叫所有表弟妹.
全部一起去我阿公家裡玩.
全部去那裡吃飯.
玩.
每個禮拜六都是這樣.
因為我禮拜一要測驗.
補習老師就幫我運輸.
無論如何都有些事情不懂.
搞到六點多七點.
表弟妹通常都去了阿公家裡.
我就跟補習老師說.
今天不補了.
補習老師很溫柔地跟我說.
讀完再說吧.
你為什麼要這麼急.
讀完考好一點就好了.
接著我就忽發奇想.
我就跟她說了一句.
我不想運輸.
接著轉頭.
我家裡有塊黑板.
因為我家裡做生意.
很多東西有塊黑板寫在那裡.
提醒工人要做什麼.
我就在上面.
不知為什麼中文那麼厲害.
寫了四個字.
叫做抗議不溫.
幾有共產思想.
四字詞.
很精簡.
抗議你懂.
不溫你也懂.
不溫輸.

$^{801}$你明白嗎.
中國現在所有字眼都是.
很長的說話都變成兩個字.
抗議不溫.
那個補習老師沒有我扶.
他只是去我媽媽那裡算生.
他說王太.
我今天不補了.
跟你的兒子.
接著他哭.
我看到他哭.
他說他不溫.
那我還有什麼做法呢.
死了那晚.
那晚精彩.
不說給你聽.
不說.
大兄姐們.
你有這麼經歷.
大兄姐們.
當如果你做了這一件事之後.
受害的人不是那個補習老師.
不是我媽媽.
其實是我.
結果很淒涼.
多淒涼我不說了.
你還笑.
你可能真的很有經驗.
大兄姐們.
我想你也會有這樣的經驗.
當你自己做了一些事的時候.
那個結果.
惡果回到你身上.
我只是抗議不溫.
玩少日無所謂.
其實不是玩少日.
是玩少很多日.
未至於玩完.
但我媽媽就這樣說.
你讀不成書.

$^{841}$你日後做乞丐你就知死.
當然我知道我一定不會做乞丐.
但是大兄姐們.
誰先受害.
以東人.
以素人.
你作惡.
去害你的兄弟.
思想上行動上.
怎樣去折磨他.
你以為你自己一定.
就會得到益處嗎.
在希伯來文裡面.
最特別的地方就是.
那句說到.
必永遠斷絕這句說話.
他用了一個很特別的.
動詞形態.
希伯來文叫Nevalstem.
什麼叫Nevalstem.
你不需要知道.
但你知道它的作用是什麼.
就是一個自我反省的動詞.
希伯來文很特別.
和希臘文都很相似.
每一個動詞都有規律.
有些動詞是反思的.
有些動詞是勸導的.
有些動詞只是平鋪直敘的.
叫做記序的,中文都有的.
分門得很仔細.
其實在這裡告訴.
以素人,以東人.
你永遠被斷絕.
其實是讓你去反思.
在這樣的情況之下.
你是會有這樣的可能性的.
你是不是真的想這樣走呢.
怎麼走呢.
看回那句.

$^{881}$羞愧必遮蓋你.
和合本翻譯為.
羞愧必遮蓋你.
遮蓋.
記不記得我的翻譯是什麼.
是讓你被隱藏和包圍.
你被你自己的罪包住了.
你自己都不知道.
我用了一個文雅一點的說法叫作繭自縛.
你是自己被罪捆綁住都不知道.
因為你覺得在罪裡面很happy.
你以為你得到一些益處.
以東人,以素人.
你以為因為出賣你兄弟.
在之後的十二節至十五節裡面說的時候.
你以為出賣你得到好處.
結局.
你都沒有賣.
在這樣的情況之下.
讓我們看見.
其實以東人你那種的情況.
其實好像有一個井底之蛙.
只是看到一點點利益.
但是你卻不知道這個世界的凶險.
而這個凶險是.
罪會反噬你.
反噬到的情況.
連你自己都不知道.
包裹住你.
好像hold住你那樣.
但是原來你已經被窒息在裡面.
而這個是你作繭自縛的結局.
不是上帝使你受害.
你說是嗎.
你還記不記得第一節.
那班縱國的人組織了一個團隊.
大肆團隊.
一齊說興起上來.
一齊動你的.
記不記得.

$^{921}$上帝只需要容許人性的可怕.
當你和他們以為同一黨.
你就可以和他們一起同撈同煲的時候.
原來結局就是這班人在利益當前的時候.
第一個先犧牲的就是你.
在這裡我不作任何的演繹或者運用.
因為實在太過多層面可以去反映出來.
但是在這個時候.
我想和大家分享我對現在的.
加拿大的國家的痛心.
記不記得自去年開始.
加拿大很多超級市場.
在不斷地在我們偉大的安省省的省長.
和我們國家的新元首.
一路去鼓勵我們要買加拿大食物.
加拿大產品.
我們只買加拿大貨.
不買其他國家.
尤其是南邊美國的貨.
是不是這樣.
又說要關掉電掣又說要拔插頭.
不供電給你.
是不是這樣.
以為很聰明.
很多基督徒都這樣看.
我們愛國.
你看看我們多愛我們的國家.
不買美國貨.
不要用他們的東西.
最近你們有沒有去超級市場.
我都少了去超級市場.
不過最近我去了一次.
我就發現到.
很多東西原本標誌為加拿大貨.
現在改了些.
你有沒有留意改了些什麼.
叫做Buy Canadian Brand.
買加拿大品牌.
有什麼分別牧師.
有啊.

$^{961}$因為那些加拿大品牌.
很多是有美國的產品所做出來的.
銷售的貨物.
你留心看一下.
很多Canadian Brand寫到明的.
裡面寫到明.
貨物來源是美國.
說什麼.
你去搞這些.
以為愛國主義.
煽動國家的人.
來達至抗衡.
另一個國家.
尤其是友邦的.
應該是美國.
導致加拿大今時今日.
多少人多憎美國.
這些仇恨.
誰做出來的.
就是那些省督.
就是我們這個國家的總理.
我不是做政治宣傳.
我是想告訴大家.
如果做基督徒看不清這些.
隨波逐流跟著這些人去做.
我很憎美國的Donald Trump.
你這樣想.
You are too naive.
讓我借這句說話.
好像正式那句說話.
You are sometimes too naive.
Too simple sometimes too naive.
多謝你.
結局就是沒有生意做.
你不靠美國的產品.
來做你加拿大的產品.
你沒有生意做.
鬼鬼祟祟又說Canadian brand.
是不是反世.
還有來.

$^{1001}$陸續有來.
但我們基督徒是否在這個世界上.
沒有智慧地走下去.
我不是叫你們擁抱美國.
不是.
它對的東西你要支持.
不對的東西你要道出它的錯誤.
這個是做基督徒應有的.
屬神生命有的誤性.
不是被人煽動之後.
以為自己還好像很愛國.
包裹著自己一個很漂亮的外表.
原來結局反世的是你自己.
以東人最後自招徹底摧毀.
其實是因為他沒有正面面對自己的罪.
於是乎他在罪的內卷之下.
我現在很喜歡用內卷.
In position.
內卷自己谷爆自己.
然後死在自己的手上.
上帝就真正不想以東人.
轉眼又有沒有透過.
上帝在這個時候向以東這番話.
去想想你落得如此被老的局面.
又是誰的錯呢.
上帝的錯.
錯不在我.
上帝的錯.
每逢說錯不在我那些國家.
一定是沒有上帝的人.
因為沒有上帝的人就是這樣.
專是賴別人錯.
自己從來不想自己會錯.
所以剛才節目以東的結局.
會不會是猶大國的結局呢.
被自己的罪捆索.
被羞辱牢牢包固.
魔鬼撒旦就想.
魔鬼撒旦很想我們是這樣的下場.
神的子民.

$^{1041}$親愛的弟兄姊妹.
我們看不看到魔鬼撒旦這一套呢.
而學會怎樣去真正認罪悔改.
面對自己的罪.
而不會自己捆綁在罪當中.
讓我們一起和上帝祈禱.
父神多謝你讓我們在這段聖經裡面.
看到摸到你的心腸.
相信你有這份的愛心.
有這份的公義.
湊合出一份給我們看到.
怎樣你對我們真正的挽救.
去到在新約聖經當中.
你甚至透過浪子的爸爸.
來去剖析.
給我們看到.
上帝你何等關懷.
那個流浪在外方的浪子.
與及那個身在身邊.
原來沒有體會上帝心腸的大兒子.
求你讓我們這班熟悉你的子民.
真真正正能夠活出有你的樣式.
明白你的心思心腸.
並且以一樣的心思心腸.
活在這個地上.
縱使世人未必會喜悅.
甚至恥笑.
乘機傷害我們的時候.
讓我們仍然堅守在其中.
因為上帝你的喜悅.
才是我們最大的期望.
求你給我們這份勇氣.
屬靈上有這份勇氣.
站在真理裡面.
來活出一個.
上帝你要把舊恩普世挽回.
你所選的人身上.
禱告奉靠主耶穌基督德性的名字而求.
阿們.
\newpage



\section{申命記 2:2-25}
\label{sec:032433}
\textbf{人生的冤枉路}~~~~禤凱榮牧師
\newline
\newline
連結: \href{http://www.church.com.hk/acms/content.asp?site=cdc\&op=show\&type=product\&code=032433\&layout=sermon}{\texttt{032433}} ~~~~ 語音日期: 2026-07-12
\newline
\newline
\hyperref[sec:032432]{\small{< < < PREV SERMON < < <}}
~
\hyperref[sec:index_chronic]{\small{[返順時目]}}
~
\hyperref[sec:032434]{\small{> > > NEXT SERMON > > >}}
\newline
\newline
申命記 2:2-25
\newline
\begin{longtable}{cl}
\hline
\hline
章節 & 經文 (和合本修訂版)\\
\hline
2:2 & \begin{tabularx}{0.7\textwidth}{X} 耶和華對我說： \end{tabularx} \\ \\ \relax
2:3 & \begin{tabularx}{0.7\textwidth}{X} 『你們繞行這山已經夠久了，要轉向北方。 \end{tabularx} \\ \\ \relax
2:4 & \begin{tabularx}{0.7\textwidth}{X} 你要吩咐百姓說：你們弟兄以掃的子孫住在西珥，你們要經過他們的邊界。他們必懼怕你們，但你們要分外謹慎。 \end{tabularx} \\ \\ \relax
2:5 & \begin{tabularx}{0.7\textwidth}{X} 不可向他們挑戰；他們的地，連腳掌可踏之處，我都不給你們，因我已將西珥山賜給以掃為業。 \end{tabularx} \\ \\ \relax
2:6 & \begin{tabularx}{0.7\textwidth}{X} 你們要用錢向他們買糧吃，也要用錢向他們買水喝。 \end{tabularx} \\ \\ \relax
2:7 & \begin{tabularx}{0.7\textwidth}{X} 因為耶和華---你的神在你手裡所做的一切事上已賜福給你。你走這大曠野，他都知道。這四十年，耶和華---你的神與你同在，因此你一無所缺。』 \end{tabularx} \\ \\ \relax
2:8 & \begin{tabularx}{0.7\textwidth}{X} 「於是，我們經過我們弟兄以掃子孫所住的西珥，從亞拉巴的路，經過以拉他、以旬‧迦別，轉向摩押曠野的路去。 \end{tabularx} \\ \\ \relax
2:9 & \begin{tabularx}{0.7\textwidth}{X} 耶和華對我說：『不可侵犯摩押，也不可向他們挑戰。他們的地，我不賜給你為業，因我已將亞珥賜給羅得的子孫為業。』 \end{tabularx} \\ \\ \relax
2:10 & \begin{tabularx}{0.7\textwidth}{X} 先前，以米人住在那裡，百姓又大又多，像亞衲人一樣高大。 \end{tabularx} \\ \\ \relax
2:11 & \begin{tabularx}{0.7\textwidth}{X} 他們跟亞衲人一樣，也算是利乏音人，但摩押人卻稱他們為以米人。 \end{tabularx} \\ \\ \relax
2:12 & \begin{tabularx}{0.7\textwidth}{X} 從前，何利人也住在西珥，但以掃的子孫把他們除滅，佔領了他們的地，接續他們在那裡居住，如同以色列在耶和華賜給他們為業之地所做的一樣。 \end{tabularx} \\ \\ \relax
2:13 & \begin{tabularx}{0.7\textwidth}{X} 『現在，起來，過撒烈溪！』於是我們過了撒烈溪。 \end{tabularx} \\ \\ \relax
2:14 & \begin{tabularx}{0.7\textwidth}{X} 從離開加低斯‧巴尼亞到渡過撒烈溪，這段時期共三十八年，直到這一代的戰士都從營中滅盡，正如耶和華向他們所起的誓。 \end{tabularx} \\ \\ \relax
2:15 & \begin{tabularx}{0.7\textwidth}{X} 耶和華的手也攻擊他們，將他們從營中除滅，直到滅盡。 \end{tabularx} \\ \\ \relax
2:16 & \begin{tabularx}{0.7\textwidth}{X} 「百姓中所有的戰士滅盡死亡以後， \end{tabularx} \\ \\ \relax
2:17 & \begin{tabularx}{0.7\textwidth}{X} 耶和華吩咐我說： \end{tabularx} \\ \\ \relax
2:18 & \begin{tabularx}{0.7\textwidth}{X} 『你今日要經過摩押的邊界亞珥， \end{tabularx} \\ \\ \relax
2:19 & \begin{tabularx}{0.7\textwidth}{X} 走到亞捫人之地。不可侵犯他們，也不可向他們挑戰。亞捫人的地，我不賜給你們為業，因我已將那地賜給羅得的子孫為業。』 \end{tabularx} \\ \\ \relax
2:20 & \begin{tabularx}{0.7\textwidth}{X} 那地也算是利乏音人之地，因為先前利乏音人住在那裡，亞捫人稱他們為散送冥人。 \end{tabularx} \\ \\ \relax
2:21 & \begin{tabularx}{0.7\textwidth}{X} 那裡的百姓又大又多，像亞衲人一樣高大，但耶和華從亞捫人面前除滅他們，亞捫人就佔領他們的地，接續他們在那裡居住。 \end{tabularx} \\ \\ \relax
2:22 & \begin{tabularx}{0.7\textwidth}{X} 這正如耶和華從前為住在西珥的以掃子孫，將何利人從他們面前除滅，使他們得了何利人的地，接續他們在那裡居住，直到今日一樣。 \end{tabularx} \\ \\ \relax
2:23 & \begin{tabularx}{0.7\textwidth}{X} 亞衛人先前住在鄉村直到迦薩；從迦斐託出來的迦斐託人將亞衛人除滅，接續他們在那裡居住。 \end{tabularx} \\ \\ \relax
2:24 & \begin{tabularx}{0.7\textwidth}{X} 你們起來往前去，過亞嫩谷。看，我已將亞摩利人希實本王西宏和他的地交在你手中，你要開始去得他的地為業，向他挑戰。 \end{tabularx} \\ \\ \relax
2:25 & \begin{tabularx}{0.7\textwidth}{X} 從今日起，我要讓天下萬民驚慌懼怕你，聽見你的名聲，就因你發顫傷慟。」 \end{tabularx} \\ \\
[1ex]
\hline
\hline
\end{longtable}
$^{1}$不知大家有無試過走冤枉路呢.
如果大家開車的話,有時我相信大家都試過.
明明開緊車,調了個導航,導航失靈.
明明十分鐘ge 車程,最後行jor 兩個鐘.
有陣時我地係職場上面拼搏好多年.
不過到升職ge 時候,發覺原來隔離gwo 個升職.
有陣時你會發覺,無論係家庭事業愛情.
你花jor 好多ge 心血,好多ge 心機.
卻似乎一無所得.
甚至有時一個意外,一個衝擊.
你會發覺多年ge 努力.
成為泡影,花為烏有.
正所謂一夜回到解放前.
當呢D 咁ge 事情發生.
當你覺得人生好似停滯不前ge 時候.
唔知你ge 感覺係點樣ge 呢.
或者好多人都會咁樣祈禱.
神呀,點解我仲係呢份工作gwo 度架.
神呀,點解我ge 家庭ge 景況.
仲係咁樣毫無改善架.
神呀,點解我個病仲未好返架.
神呀,點解你唔聽我個禱告架.
好多好多ge 點解.
我地大家都經驗過,頭先所講ge 情況.
就好似一個.
妙無人煙,相當荒蕪ge 曠野一樣.
我地叫呢個係曠野ge 一個經驗.
無任何ge 改變,無論我地如何ge 努力.
我地就會問,神呀,係咪我走錯路呢.
神呀,係咪我走ge 一D 冤枉路呢.
點解其他人都係一路ge 向前.
我係硬係原地踏步.
點解呢.
今日ge 經文好長,好多ge 地名,好多ge 人名.
但係話俾我地真係好簡單ge 一個情況.
就係當時ge 以色列民.
佢地就好似我地頭先所講ge 景況一樣.
佢地從埃及去到迦南所謂ge 應許之地.
正常黎講係十零日ge je .
不過佢地竟然走jor 四十年.

$^{41}$係人睇黎呢,gwo 條冤枉路啦.
十幾日行jor 四十年仲唔係走冤枉路.
不過係神ge 眼中睇黎呢.
gwo 段路並唔係迷路黎架喎.
而係一條管教ge 路.
一條重建ge 路.
同埋一條修正ge 路.
讓呢班原本係叛逆ge 子民能夠重回神ge 道路上面.
鼎姐妹.
你有冇試過走冤枉路呀.
如果你曾經係迷失冤枉ge 道路上面.
你點樣走返出黎架.
又或者如果今日你仍然係呢一個ge 迷失或者冤枉ge 道路上面ge 時候.
我希望大家.
我地可以通過今日剛才所讀ge 經文.
我地一齊學習.
如果我地係迷路ge 當中.
我地能夠返返到去神ge 身邊.
我地能夠係倚靠佢ge 時候.
我地係呢個所謂冤枉ge 道路上面.
每一步都可以成為一個有意義ge 路程.
用我一條黎到學習.
我地上個星期就提到呢.
由於當時ge 以色列人唔信福神.
最後引黎神ge 憤怒.
神ge 憤怒正應話比我地聽.
神ge 審判.
審判jor gwo 個不信ge 世代.
gwo 班仲埃及出黎.
二十歲以上ge 藍丁.
全部都要死係曠野.
冇得入神所應許ge 美地.
好嚴厲ge 一個管教.
原本係幾日之內就可以入迦南.
最後行jor 三十八年.
然而大家留意.
你三十八年.
雖然係懲罰.
但我地睇到ge 神.
冇因為懲罰ge 緣故.

$^{81}$黎到離棄佢ge 百姓.
係《申命記》二章七節.
gwo 度黎到話比我地聽.
因為耶和華你ge 神.
係你手裡面所辦ge 一切事.
已經賜福比你.
雖然係四十年走在曠野兜兜轉轉.
但耶和華你ge 神.
常與你同在.
故此.
你一無所缺.
何等寶貴ge 一句說話.
而刑罰之中有公認.
管教當中有神ge 同在.
大家諗o下喎.
無錯.
神係gwo 幾十年係曠野ge 路程上面.
神係管教緊佢地.
但係佢地曠野冇餓死到.
冇因為冇水而渴死.
每日有嗎哪供應比佢地食.
唔知行邊度ge 時候.
白天有雲柱.
晚上有火柱.
衣服同鞋都冇著爛.
好大個奇蹟係咪呀.
我想問問大家.
你著件衣服.
最耐著多幾多年可以.
我有兩件衫我數過.
兩件衫我差唔多著過三十年.
三十年.
嘩你咁犀利ge .
你兩件衫可以著三十年ge .
因為我唔係成日著架嘛.
我間唔中著下je .
但我想諗下大家.
當日ge 以色列人.
走在曠野上面.
我唔相信佢有幾十套衣服.

$^{121}$俾佢日日咁樣換架係咪.
咁鞋呢.
你最耐ge 鞋著jor 幾耐呀.
我諗下計一計.
我冇一對鞋應該超過五年.
應該冇.
你呢.
但係呢個聖經話比我地聽.
當日走在曠野上面ge 以色列民.
四十年之久.
佢ge 鞋同埋衣服冇著爛到.
冇著爛到.
同埋大家去諗一諗.
佢地日日著住架喎.
好明顯.
係呢幾十年ge 曠野ge 日子裡面.
神與他門同在.
神冇因為呢一個刑罰.
而停止jor 佢對子民ge 供應.
弟兄姊妹.
有陣時我地接受神ge 管教.
好得時我地留意ge 就係神.
lor 去jor 我地D 咩ye ge ye .
但係我地好少留意.
係呢個管教ge 過程當中.
神仍然ge 保留jor 幾多ge 恩典.
係我地ge 生命裡面.
不過雖然係咁樣.
我地千祈唔好輕看.
神對我地管教ge 嚴肅性.
聖經呢.
係度呢.
就有咁ge 記載ge .
自從呢一個佢地ge 不信服開始.
一共有三十八年.
呢三十八年點呀.
神話等那世代ge 兵兵都從刑中滅盡.
正如耶和華向他們所起的誓.
好明顯我地呢幾次簡單ge 經文話比我地聽.
我地所信ge 神係一個公義ge 神.

$^{161}$我地人必須要為自己ge 不信.
同埋叛逆承擔後果.
gwo 一代ge 世人ge 結局呢.
其實係提醒緊我地信仰.
唔係玩泥沙.
唔係兒戲ge .
神ge 應許.
你唔好當係理所必然ge .
我地係需要用信心黎到作為回應ge .
不過冇錯啦.
呢個走在曠野ge 幾十年係一個懲罰.
不過想大家留意.
如果我地單係將呢個曠野ge 日子.
單單睇為神對百姓ge 刑罰ge 時候.
我估我地錯失jor 神更深遠ge 心意.
38年ge 曠野漂流.
大家諗一諗.
對新一代ge 以色列黎講.
係實戰ge 環境當中黎到被塑造.
係被預備ge 一個過程.
因為點解呀.
佢地要進入呢個迦蘭美地ge 時候.
要面對好多ge 爭戰.
神就係呢幾十年裡面預備一班新ge 一代.
通過呢38年ge 預備同埋訓練.
我地睇到新一代ge 以色列人.
無論係身體同埋靈性上.
不斷ge 壯大.
不斷ge 黎到被裝備.
正如呢個希伯來書所講.
如果我睇呢個希伯來書ge 十二章十到十一節.
佢話我地ge 地上ge 父親.
隨住自己ge 意思黎到教養兒女ge .
唯有萬靈ge 父.
我地天上ge 父親.
管教我地.
是要我地得益處.
使我地係佢ge 聖潔上有份.
大家都知道.
比爹地媽咪.

$^{201}$管教ge 時候.
我地被管教ge 時候.
一D 都唔感到快樂.
反而覺得受苦.
不過當我地經歷jor 管教.
然後望返轉頭ge 時候.
gwo 一個ge 管教.
反而係我地經歷練而結出屬靈ge 果子.
得著屬靈ge 福氣.
所以弟兄姊妹.
或者你正處係人生ge 所謂曠野ge 階段.
呢個曠野可能係事業ge 停滯.
可能係關係ge 破裂.
可能係健康ge 危機.
可能係靈性ge 干擾.
你覺得自己無論做咩都冇用ge .
改善唔到個情況.
好似走緊冤枉路一樣.
覺得一切都係白費ge .
不過呢段經文.
神話比我地聽.
曠野絕對唔係終點黎ge .
而係一個必經之路.
只要我地願意迴轉.
只要我地有受教ge 心.
管教絕對唔係.
神對佢ge 百姓ge 拒絕或者放棄.
反而係塑造.
讓我地可以經歷呢一個ge 管教之後.
得著更大ge 屬靈好處.
你話神如何.
係一個代表住管教ge 曠野裡面.
無助佢ge 百姓呢.
如果我地睇第十五節ge 時候.
呢度講到.
耶和華ge 首也攻擊他們.
將他們從營中除滅.
直到滅盡.
好嚴厲ge 一個說話.
我想大家留意呢一句說話當中ge 攻擊.

$^{241}$呢一個字.
原文就係同神.
與敵人爭戰時候.
所用ge 攻擊.
係同一個字黎ge .
即係話比我地聽.
gwo D 唔信服神ge 百姓.
最終係以神ge 敵人ge 身份.
黎到死去.
好可怕ge 一個提醒.
所以弟兄姊妹唔好輕看神ge 命令.
當面對神ge 帶領ge 時候.
千祈唔好話.
不過如果.
但係或者不如果.
唔好用呢D ge 說話黎到去推搪.
係關鍵ge 時候.
唔好懷疑唔好憂慮唔好恐懼.
我地要信服.
我地要倚靠.
我地要信靠.
以色列ge 歷史.
俾我地一個好好ge 尷尬.
從呢一個ge 曠野地.
以色列人ge 經驗.
我覺得我地可以學到兩個好重要ge 功課.
第一個.
不信服.
係帶有嚴重ge 後果.
大家記唔記得最嚴重ge 後果.
大家可以想像到係咩.
記唔記得人類ge 始祖華當.
佢唔信服神.
神講到明.
成個園裡面ge 果子你都可以食.
除jor 呢一個生命樹.
不過呢一個ge 亞當卻係冇理會.
卻係唔信服.
照食.
最後帶黎ge 後果係不能想像ge 後果.

$^{281}$唔單止佢受到咒詛.
係整個ge 創造.
整個ge 世界受到創造.
因為呢一個神ge 咒詛.
罪進入人世間.
自此.
人要面對死亡.
自此.
人.
活係罪ge 綑綁裡面.
係呢一個ge .
不信服ge 世代.
被完全ge 滅盡之後.
好啦.
新ge 一代將要進入.
應許之地.
神對新ge 一代ge 以色列人講.
當你地要進入呢個應許之地ge 時候.
你會經過其他ge 民族.
包括以東人.
摩訶人.
阿滿人.
你地會經過或者走近呢一D ge 土地.
不過你地千祈唔好.
妖害佢地.
唔好同佢地爭戰.
點解呀.
係咪因為呢班人好多人驚呀.
打唔贏呀.
唔係.
神講得好清楚.
唔好搞佢地.
因為點解.
因為神已將那地賜給他們.
我識難D ge 句話.
好緊要.
話俾我地聽下一樣好重要ge 事情.
我地需要知道.
神唔單係以色列ge 神.
神係列國ge 主.

$^{321}$連以東摩訶阿滿ge 土地.
神都有所安排.
而更加值得注意ge .
就係經文話比我地聽.
呢D 民族.
亦都曾經擊敗.
原本係gwo 度住ge 巨人族.
得到佢地ge 地圖.
如果我地睇呢一個二章十到十二節.
當時呢幾個民族.
入住呢一個迦南地之前.
gwo 度原本係有好多ge 人居住jor .
民眾民數眾多.
身體高大.
就好似阿納人一樣.
但係神讓佢地進入gwo 個地圖.
得著gwo 個ge 地圖.
最後gwo 句.
就如以色列在耶和華.
其他遺孽之地所行的一樣.
咩意思.
意思就即係話.
神亦都係幫助以東摩訶阿摩人.
黎到去驅趕.
當時係gwo 度居住.
好人數眾多ge 阿納人.
我想大家諗一諗.
大家覺唔覺得奇怪.
我地可以理解.
神點解會幫助以色列民.
黎到入住呢個迦南美地.
點解.
因為以色列民係神ge 選民.
但係大家要留意.
呢班ge 以東摩訶阿摩人.
佢地唔係神ge 選民.
亦都唔相信.
亦都唔認識.
亦都唔跟隨呢個神ge .
既然係咁ge 時候.

$^{361}$神點解要幫助佢地.
點解.
因為神曾經承諾過.
呢幾個民族ge 祖先.
神所作出ge 應許與承諾.
神必履行 神必英文保守.
縱然佢地ge 下一代並不認識佢.
我覺得我地可以明白.
神一個好重要ge 屬性.
係列國ge 裡面.
神有主權.
無論你信與唔信.
無論你認唔認識呢個神.
正如係呢個史唐寺有一段咁ge 經文.
神從一本做出萬族.
萬族住在地上.
神預先定準佢地ge 年限和所住ge 疆界.
由於係咁樣.
每一個人 每一個國家.
我地ge 生活動作存留都在乎佢.
係呢個但以理書亦都有類似ge 經文.
gwo 度講到.
但以理同神講.
神呀 你ge 名係應當稱重.
因為從耿古到永遠.
智慧與能力都係屬乎呢個神.
有語句.
他改變時候.
日其 廢亡 立亡.
將智慧與.
賜比智慧人.
將知識賜予聰明人.
神係度話比以色人聽.
我唔單止賜地比你地.
我亦都賜地比以素ge 子孫.
比羅德ge 子孫.
列國ge 疆界都係我ge 手當中.
不在乎呢一班人認識我與否.
唔在乎呢一班人相唔相信我.
呢個神.

$^{401}$都係整個宇宙ge 創造者.
呢個神都係列國萬國ge 主宰.
既是如此.
我話大家聽.
世界ge 歷史並不偶然.
世界ge 歷史完全係掌管係神ge 手裡面.
每一個人 每一個國家.
都係神ge 主權ge 裡面.
唔單止神.
叫當時ge 以色列人.
你唔好去搞擾gwo D ge .
你所經過或者走近ge 民族.
神要更加ge 話比你聽.
你要用錢向佢地買糧食.
你要用錢向佢地買水喝.
我想大家思考.
覺唔覺得奇怪.
幾十年係抗疫.
神一直用嗎哪供應佢地.
神讓水從盤石流出黎供應佢地.
從來唔缺食物.
從來唔缺水.
點解到而家呢一刻.
神只係叫佢地.
你去gwo 度ge 時候.
你用錢買糧食.
用錢黎到買水.
點解唔繼續供應.
我相信除jor 一個好客觀ge 原因.
因為佢地離開抗疫.
佢地現在要進入人所聚居ge 地圖.
呢個固然係一個客觀ge 原因.
但我相信有另外一個原因.
而呢個原因.
亦都係我地第二個.
可以從呢個故事學到ge .
熟齡功課.
呢個係咩功課.
尊重神係別人身上ge 主權.
神其實咁樣講.

$^{441}$gwo D ge 地圖唔係你地ge .
唔好去搶唔好去爭.
今日我地活係世上.
其中一個最大ge 試探係咩.
探.
英文有個greener pasture next door.
鄰居ge 土地梗係綠色過我地gwo 個.
有冇咁ge 感覺.
點解隔離gwo 個ge 份工好D .
個老婆靚D .
個家庭ge 仔女就乖D .
佢ge 恩典就多D .
又成功多D .
佢間教會都大D ge .
個牧師講ye 叻D ge .
有冇咁ge 感覺.
神要我地學習.
神比別人的.
唔好妒忌.
神比自己的.
我地需要珍惜同埋運用.
以色列人雖然係神ge 選民.
但佢地唔能夠因為神選民ge 身份.
就任意奪取別人ge 產業.
神比邊一個就係gwo 個.
即使gwo 個民族唔係神ge 選民.
神要我地學習一樣重要ge 功課.
佢既係掌管萬有ge 主.
佢比每一個人ge 份都唔同.
所以我地應該要感恩.
領受自己ge 份.
同一個時候.
尊重神係其他人身上ge 安排.
我地再睇落去.
跟住神對摩西講.
佢話你地繞行這山的日子夠了.
要轉向北去.
夠了.
想大家諗一諗.
呢兩個字係何等ge 寶貴.

$^{481}$夠了.
咩意思呀.
即係話比我地聽.
縱然我地係人生裡面.
走在曠野ge 日子.
但是曠野不是永遠的.
曠野不是永遠的.
神ge 管教係有定時ge .
當時候滿足ge 時候.
神就會同我地講.
夠啦.
而家轉向.
繼續向前走.
接落黎.
神俾以色列人一個好新ge 使命.
一個任務.
一個工作.
就話你而家要去同阿摩利王西弘爭戰.
得佢ge 地為業.
面對呢一個命令.
記唔記得.
進入去得地為業ge 時候.
係佢地ge 父輩.
嚇到佢地恐懼到不得了.
所以以色列人.
到新ge 一代ge 時候會唔會都係懼怕呢.
始終呢一個新ge 一代.
係曠野地長大ge 呢一班人.
都未有經過ge 咩ge 爭戰.
冇咩ge 戰績.
但係神對佢地話.
從今日起.
我要使天下萬民.
聽見你D 明星.
都驚恐懼怕.
且因你發戰相動.
我想大家留意一下.
經文提jor 幾個ge 民族.
神叫以色列人對每一個民族ge 處理ge 方式.
都有所不同ge .

$^{521}$但對以東人黎講.
神話唔好同佢地爭戰.
你要和平經過.
甚至用錢黎到買佢地ge 糧食同埋水.
對摩訶人就係.
不可擾攘.
不可爭戰.
對阿摩人都係一樣.
不可擾害.
不可爭戰.
但對阿摩利人.
就要.
與他爭戰.
得他的地為業.
點解唔同ge 民族有唔同ge 對待方式ge 呢.
因為神話.
你要爭戰就爭戰.
神話唔好爭戰就唔好爭戰.
呢個話比我地聽個好重要我地需要理解.
神有絕對ge 主權.
亦都話比我地聽一樣ye .
係屬靈ge 道路上面.
唔係樣樣ye 都要逞強ge .
唔係樣樣ye 都係勇往直前硬碰硬.
用強權用勇力黎到去絕對取ge .
好多時我地有個錯誤ge 感覺.
神走在我地前面為我地爭戰.
所以我地每一次ge 向前行.
都係諗住爭戰爭戰.
搞錯jor .
有一D ge 時候.
有一D ge 關係.
神要我地和平相處.
有一D ge 試探.
神要我地勇敢ge 爭戰.
有一D ge 機會.
神要我地憑信心黎到去爭取.
有一D ge 東西.
神要我地學識放手.
關鍵係在於咩呀.

$^{561}$關鍵在於唔係我地要用強.
關鍵在於我地要聆聽.
分辨神ge 聲音.
跟隨神ge 帶領.
神對以色列人講.
從今日起.
我要使天下萬民聽見你的名聲.
都驚恐懼怕.
大家有冇睇到呢個心態上面.
完全ge 逆轉jor .
記唔記得當日.
當呢一個以色列人.
聽到gwo D 巨人係個.
加拿大ge 時候.
怕到不得了.
而家神話呢個列國會懼怕以色列人.
呢個懼怕完全ge 逆轉jor .
呢一個逆轉點解會發生呢.
並唔係因為以色列人.
突然之間勁jor 猛jor 大隻jor 高大jor .
唔係咁ge 原因.
而係因為係呢幾十年ge 磨練裡面.
呢班以色列人終於學識jor .
咩叫信服.
終於學識jor 咩叫跟隨.
當我地唔再憑自己ge 眼光.
黎到睇環境ge 時候.
當我地係用信心.
睇神ge 同在同埋帶領ge 時候.
我地原本應該有個懼怕.
就會成為別人ge 懼怕.
並唔係因為我地有D 咩得人驚ge 地方.
而係因為我地有一個.
當受敬畏ge 神與我地同在.
所以大家我地都要明白.
我地每一個人都會走冤枉ge 路.
無論你點咁醒目.
無論你點咁幸運.
人生總係有一D ge 冤枉ge 經歷.
當我地睇到.

$^{601}$一D 浪費ge 歲月.
失敗ge 經歷.
破碎ge 關係.
呢D 都係我地走冤枉路ge 時候所經歷到.
但係當我地找緊神ge 應許.
用信心ge 眼光.
重新審視我地所謂ge 冤枉路ge 時候.
其實我話比大家聽.
每一次曠野ge 歷程.
都可以充滿神ge 恩典ge 痕跡.
每一次ge 管教.
我地願意ge 話.
都可以成為我地更加堅強.
可以成長ge 機會.
當管教過jor 之後.
當神好清楚話比我地聽.
夠啦,而家要再向前走.
咁樣ge 時候.
我可以話比大家聽.
就係一個時候話比我地聽.
我地擾行ge 日子夠啦.
我地自怨自艾ge 日子夠啦.
我地活係過去ge 日子夠啦.
gwo 個時候.
係時候我地離開我地ge 曠野.
進入神所應許ge 時候.
唔好怕,唔好退縮.
我地亦都或者面對過一個需要靠信心跨越ge 關口.
我地都好多時候我地需要一個好艱難ge 決定.
我地需要靠勇氣ge 見證.
我地需要一個付代價ge 信服.
神對我地話.
起身啦.
時間到了.
我已將那地交在你手中.
大家都同大家傾講過.
爭戰得勝ge 秘訣.
唔在於我地ge 力量.
而係在於神是否走係我地ge 面前.
作為總結.

$^{641}$再重申我地三個今日學到ge 功課.
一個三個不變ge 真理.
曠野係神管教我地ge 地方.
當我地做錯ye .
當我地唔信服ge 時候.
我地就被神管教.
聖經講得好清楚.
神所愛的他必管教.
所以當我地受管教ge 時候.
也許我地經歷好多痛苦.
覺得係一個好艱難ge 日子.
不過係呢個整個ge 管教與艱難ge 當中.
神仍然供應我地一切ge 需要.
我地並不因神ge 管教.
神而離棄我地.
第二.
曠野係學習之地.
讓我地係呢個歷程被管教ge 當中.
學習咩叫信服.
學習點樣分辨善與惡對與錯.
神要我地學習信服.
學習尊重別人身上ge 主權.
學習唔好貪戀唔屬於自己ge ye .
第三.
曠野係必經之路.
我地要進入應許.
我地要進入神所為我地預備ge 一切ge 時候.
好多時候我地必須要經過曠野.
被磨練.
被操練.
被無做.
但係呢一個曠野.
好似我頭先所講.
絕對唔係終點.
當時候滿足ge 時候.
神會親自同我地講夠啦.
唔好仍然活在過去當中.
而要挺起胸膛.
向前邁步.
因為神走在我地ge 前面.

$^{681}$為我地爭戰.
所以願意我地唔好似第一代ge 以色列人咁樣.
因為不信而倒閉.
係呢一個曠野gwo 度.
我地要好似當時ge 約書亞.
或者係迦勒一樣.
專心跟隨主.
最後能夠進入神所應許ge 豐盛之地.
好似詩篇所講.
詩篇話.
神使曠野變為水潭.
叫罕地變為水泉.
呢一個能夠將無變有.
能夠將曠野變為水源ge 神.
今日.
仍然係我地ge 生命當中.
做神蹟奇事.
我地ge 曠野.
亦都可以因為神ge 原因.
可以成為恩典ge 傳言.
我地兜兜轉轉ge 人生.
亦都可以成為.
我地一次.
得勝向前行ge 機會.
所以願神.
賜福我地每一個人.
走在人生路上ge 時候.
得著更大ge 恩典.
請我地同低頭祈禱.
親愛主耶穌基督.
我地ge 天父上帝.
我地到你面前ge 時候.
因為你就係gwo 個以色列ge 神.
亦都係我地ge 神.
我地人生ge 曠野.
也許我地會經歷好多ge 歷練.
但係神呀.
你從來冇放棄我地.
你ge 同在同埋供應.
係我地走在人生路上.

$^{721}$必須要有ge 力量傳言.
神呀求你赦免我地.
假如我地當中有不信.
有不信服ge 時候.
神呀求你賜我地.
好似新一代以色列人ge 信心.
我地信服一樣.
幫助我地係冤枉ge 人生路上.
學習到功課.
讓我地係當中睇到你ge 恩典.
睇到你ge 帶領.
以至我地人生得著更豐盛ge 恩典.
多謝你比我地一切祈禱.
奉耶穌基督得勝的名字而求.
阿們.
\newpage



\section{士師記 16:19}
\label{sec:032434}
\textbf{聖經人聖經事:耶和華已經離開了}~~~~黃克勤牧師博士
\newline
\newline
連結: \href{http://www.church.com.hk/acms/content.asp?site=cdc\&op=show\&type=product\&code=032434\&layout=sermon}{\texttt{032434}} ~~~~ 語音日期: 2026-07-09
\newline
\newline
\hyperref[sec:032433]{\small{< < < PREV SERMON < < <}}
~
\hyperref[sec:index_chronic]{\small{[返順時目]}}
~
\hyperref[sec:032435]{\small{> > > NEXT SERMON > > >}}
\newline
\newline
士師記 16:19
\newline
\begin{longtable}{cl}
\hline
\hline
章節 & 經文 (和合本修訂版)\\
\hline
16:19 & \begin{tabularx}{0.7\textwidth}{X} 大利拉哄參孫睡在她的膝上，叫一個人來剃掉參孫頭上的七條髮綹。於是大利拉開始制伏參孫，他的力氣就離開他了。 \end{tabularx} \\ \\
[1ex]
\hline
\hline
\end{longtable}
$^{1}$看完節目又做一日.
眨眼生命樹又同大家讀聖經.
聖經人,聖經事.
看完聖經背後的上帝工作.
今日的題目:耶和華已經離開了.
這是一個很大的嘆息.
耶和華離開了.
離開了,神不在了,不和他一齊了.
士師記十六章:大利拉見他把心中所藏的刀告訴了他.
就打發人到非利士人的首領那裡對他們說.
他已經把心中所藏的刀告訴我了.
請你們再上來一次.
如果我應該加一個括號.
記得帶錢.
再上來一次.
於是非利士的首領就手裡拿著銀紙.
上到婦人那裡.
沒有說多少個首領.
沒有說帶多少人.
沒有說帶多少錢.
終於這次成功了.
說了.
其實每次都是這樣.
今天會不會都是半信半疑呢?.
也不奇怪.
大利拉使參孫浸著他的膝睡覺.
並且叫了一個人來剃除他頭上的七條髮柳.
於是大利拉克制他.
他的力氣就離開他了.
今天你看到用黃色螢光的字.
其實參孫他想著我講真實的原因給你聽.
我都不怕.
這些頭髮長在我頭上.
你怎會拿到?怎會剃到?.
你當我傻的嗎?.
我綁著你就剃到.
綁我的繩子只有一根線.
我拍一拍就會斷了.
所以你怎會剃到我的頭?.
所以我告訴你我都不怕.

$^{41}$頭髮幸好我沒有能力了.
參孫想著你剃不到我的.
你剃眼眉嗎?.
你做不到.
剃我的頭髮我醒了你做不到.
綁著我十個人扎我都推開你.
於是大利拉這裡聰明了.
問他問得很差.
這裡聰明了.
使參孫浸著他的室睡覺.
你知道參孫去找這些女士.
都可能是比較上有他們做的事.
於是你比較累都不一定.
這裡很有可能有一件事.
你怎可以叫人剃到他頭上的髮瘤.
你一剃他就醒了.
你剃剩一條都不行.
我推測當時是否很累.
睡覺吧睡覺吧.
可能一直掃著參孫的頭髮.
掃啊掃啊長的.
睡覺吧睡覺吧休息休息.
因為你一直動不動.
剪了都不知道.
剃了都不知道.
所以一直梳理他的頭髮.
令他很舒服地睡覺.
在這時候.
不要弄我的頭髮.
我的頭髮不能動.
他會醒的.
但他令他浸著他的膝蓋睡了.
很安全很舒服.
可能還在弄他的頭髮.
所以他很舒服地睡了.
剃了力氣就離開他.
這樣開始失手了.
這就是他的原因.
太過舒服毫無警覺.
蜀靈上第一點要講的是.

$^{81}$他在蜀靈上也睡了覺.
人睡了覺也不要緊.
但最重要蜀靈上沒有睡.
他在蜀靈上睡了覺.
所以力氣就失去了.
大尼拉就說參孫.
非利士人來捉你了.
他跟她講過一次.
參孫你起來吧.
非利士人來捉你了.
他有時逆正逆邪.
起來吧非利士人來捉你了.
參孫從睡中醒來.
叫醒非利士人來捉你了.
他就想我要像前幾次出去活動身體.
讓我去給他幾招.
卻不知道第三點.
當你在蜀靈上睡覺的時候.
你的敵人就會來了.
頭髮就剪了.
敵人就來了.
耶和華就離開了.
所以這三部曲就很慘了.
第一部曲是什麼?.
就是參孫睡著了.
頭髮被剃了.
所以他的力氣.
蜀靈上的強壯力氣沒有了.
所以要警醒.
免得仇敵來迷惑你們.
就是這樣講的.
要警醒蜀靈上睡著了.
沒有防備.
第二.
你一睡著敵人就來捉你了.
而且還是你身邊的人帶他來.
非利士人來捉你了.
捉我?不是吧.
還很自負.
我出去動他才行.

$^{121}$我出去打他.
出去活動身體.
你知道我多厲害.
誰不知不知道.
第一,睡覺了.
力量失去了.
然後敵人就來了.
卻不知道耶和華離開了.
什麼都假.
耶和華離開了.
就不行了.
就被捉了.
輸了.
非利士人將他拿著.
這裡是碗了他的眼睛,讀"挖".
我查過,可以讀"碗",也可以讀"彎".
無論是哪個星,"彎"或者"碗".
都是指"挖".
而且原來有這個一棟一堤的邊.
這是表示用刀.
好像有匕首的刀.
用刀挖掛了他的眼睛.
用刀這樣綁住他.
很害怕,眼睛刺進去.
然後挖斷他的眼睛.
很可怕,很痛,很害怕.
將他拿住,然後弄盲了他.
其實摩西很早已.
不是摩西,經常說摩西,對不起.
其實參孫一祖屬靈的眼已經盲了.
現在是盲了肉身的眼.
這樣都沒有看到.
還要經常去非利士人那裡.
還要在他的膝蓋那裡睡了.
於是將他帶回去加薩.
加薩就是非利士人的總部.
總部城市.
而且用銅鏈驅索他.
為什麼用銅鏈呢?.
銅鏈很貴,而且很難找.

$^{161}$銅鏈其實沒有鐵鏈那麼好.
但銅鏈有什麼好處呢?.
粗一點的銅鏈不生鏽.
要困他多久呢?.
雖說沒有想過很快放的.
也沒有想過很快弄死他的.
所以用銅鏈鎖住他.
扔在監裡面,做什麼呢?.
推磨,大力,推磨.
推磨是什麼呢?.
是磨菊,磨麥粉.
或者很大塊石,中間凹了的.
將橄欖扔進去.
然後有一個站立的輪子.
輪子是有一點斜的.
然後中間穿一條木.
綁在你的頭上,綁在你的肩膀上.
或者你推著走.
不停地轉,不停地轉.
那塊石就磨爛了中間的橄欖.
通常是驢仔或者牛做的.
在監裡面推磨橄欖,磨麥粉之類.
磨石很重.
他的頭髮被剃,但又漸漸長起來.
你見到參孫經歷了這件事.
他的眼睛瞎了,被挖出來.
很害怕,而且被困住.
失去了自由.
他已經被拘禁了.
要看的是,耶和華已經離開了.
今天講得長,今天講短一點.
大家要留意.
當參孫失去了他的能力的來源之後.
你見到「我要出去活動一下」.
不過很可惜,離開了.
要知道最近那幾件事,幾個「了」是什麼呢?.
頭髮被剃了,即是能力失去了.
敵人馬上就「來了」.
制服你之後,你就失去了自由.
然後你會見到,為什麼有人制服你呢?.

$^{201}$最重要是因為耶和華離開了.
這個耶和華離開了,慘了.
所以弟兄姊妹,你不要失望.
如果我們的神沒有離開,在那裡就可以了.
你仍然可以依靠祂,祂的祝福隱秘仍然存在.
一間教會,有時我見那些教會是亂得不得了.
大家爭權奪利各樣東西,亂得不得了.
不過如果你堅信耶穌基督沒有離開.
只要耶穌基督沒有離開,沒問題,放心.
但是如果神離開了,一個人一個地方的教會就沒有了.
受敵已經來了,你的能力已經失去了.
最慘的是耶和華已經離開了參宣.
他是被自己的自負與自我而害死的.
就是這兩點要跟大家講的,讓我們一起禱告吧.
天父,我們在這裡覺得很嘆息.
我們一邊看參宣經歷的時候.
心裡也覺得心裡是砰砰跳的.
為什麼這麼有力量的人竟然到最後是這樣的收場.
而收場最可怕的地方固然是眼睛被刮去,被刀刮去.
眼睛很痛,看不到了,從此失明了,不再見到.
最重要的是,最開始是他屬靈的眼睛已經失明了.
幫助我們,保守我們的心眼,讓我們看得明亮,看到你的心意.
幫助我們,保守我們的能力是從你而來.
神啊,不要離開我們,讓我們心裡知道你一路的同在.
我們讚美你就是,其實你一直都應許不會離開我們.
你會一路做我們的牧者,如果會失去,是我們自己走失了.
而你是要來尋找我們,多謝你不要放棄我們.
讚美主,主我們繼續要看下去,參孫怎樣頭髮漸漸長起來.
祂的心智重新回來,祂的力量重新回來,都是從你而來的.
幫助我們,警醒,不要屬靈上,睡著了.
祈禱奉耶穌名求,阿們.
\newpage



\section{羅馬書 8:1-11}
\label{sec:032435}
\textbf{有靠山,活得帥!}~~~~林玉英牧師
\newline
\newline
連結: \href{http://www.church.com.hk/acms/content.asp?site=cdc\&op=show\&type=product\&code=032435\&layout=sermon}{\texttt{032435}} ~~~~ 語音日期: 2026-07-12
\newline
\newline
\hyperref[sec:032434]{\small{< < < PREV SERMON < < <}}
~
\hyperref[sec:index_chronic]{\small{[返順時目]}}
~
\hyperref[sec:032436]{\small{> > > NEXT SERMON > > >}}
\newline
\newline
羅馬書 8:1-11
\newline
\begin{longtable}{cl}
\hline
\hline
章節 & 經文 (和合本修訂版)\\
\hline
8:1 & \begin{tabularx}{0.7\textwidth}{X} 如今，那些在基督耶穌裡的人就不被定罪了。 \end{tabularx} \\ \\ \relax
8:2 & \begin{tabularx}{0.7\textwidth}{X} 因為賜生命的聖靈的律，在基督耶穌裡從罪和死的律中把你釋放出來。 \end{tabularx} \\ \\ \relax
8:3 & \begin{tabularx}{0.7\textwidth}{X} 律法既因肉體軟弱而無能為力，神就差遣自己的兒子成為罪身的樣子，為了對付罪，在肉體中定了罪， \end{tabularx} \\ \\ \relax
8:4 & \begin{tabularx}{0.7\textwidth}{X} 為要使律法要求的義，實現在我們這不隨從肉體、只隨從聖靈去行的人身上。 \end{tabularx} \\ \\ \relax
8:5 & \begin{tabularx}{0.7\textwidth}{X} 因為，隨從肉體的人體貼肉體的事；隨從聖靈的人體貼聖靈的事。 \end{tabularx} \\ \\ \relax
8:6 & \begin{tabularx}{0.7\textwidth}{X} 體貼肉體就是死；體貼聖靈就是生命和平安。 \end{tabularx} \\ \\ \relax
8:7 & \begin{tabularx}{0.7\textwidth}{X} 因為體貼肉體就是與神為敵，對神的律法不順服，事實上也無法順服。 \end{tabularx} \\ \\ \relax
8:8 & \begin{tabularx}{0.7\textwidth}{X} 屬肉體的人無法使神喜悅。 \end{tabularx} \\ \\ \relax
8:9 & \begin{tabularx}{0.7\textwidth}{X} 如果神的靈住在你們裡面，你們就不屬肉體，而是屬聖靈了。人若沒有基督的靈，就不是屬基督的。 \end{tabularx} \\ \\ \relax
8:10 & \begin{tabularx}{0.7\textwidth}{X} 基督若在你們裡面，身體就因罪而死，靈卻因義而活。 \end{tabularx} \\ \\ \relax
8:11 & \begin{tabularx}{0.7\textwidth}{X} 然而，使耶穌從死人中復活的神的靈若住在你們裡面，那使基督從死人中復活的，也必藉著住在你們裡面的聖靈使你們必死的身體又活過來。 \end{tabularx} \\ \\
[1ex]
\hline
\hline
\end{longtable}
$^{1}$各位弟子們平安.
首先要請大家低頭祈禱.
因為無論世界是否已經動蕩.
你的恩典總是不會離開我們.
你是我們懇求聖靈親自開啟我們的心.
讓我們聽見你對我們的生命.
說出安慰,鼓勵和提醒的說話.
天火取我們的重擔.
賜下你的自由.
我們這樣禱告直主耶穌.
基督的名而求.
阿們.
我請大家繼續閉上眼睛.
跟大家一起想像一個畫面.
可能這個畫面在香港這一刻.
很多家庭每天都會面對.
想像一下在我們客廳裡.
有一位中年的爸爸.
他坐在梳化上.
現在是星期一的上午.
全家人都出門了.
上班的上班.
上學的上學.
只剩下他一個人.
他看著手上那封上星期才收到的財源信.
他看著手機上的銀行存款.
又再看看KOL不斷提醒AI的技術.
如何在一夜之間取代了他奮鬥了二十年的專業.
他的心有很大的焦慮.
不敢出門.
不過也不敢留在家裡.
他覺得自己成為了家庭的負累.
他覺得自己是這個社會的失敗者.
我們再將視線移到廚房.
媽媽在廚房裡煮晚餐.
她心裡記掛著在房裡的孩子.
她的心情很複雜很矛盾.
既有恐懼又有擔心又有無力.
不過同時也很生氣.
為什麼別人的兒子考得這麼好.

$^{41}$入了名校.
我的兒子卻連學都不想上.
為什麼我這麼努力去經營我的家.
換來的是冷冰冰的牆.
我們又將我們的視線移到房裡.
門關得很緊.
我們聽到這位年輕人雙手劈哩啪啦地敲打著鍵盤.
戴著耳機將遊戲的聲音開得很大.
其實他真的不太喜歡玩遊戲機.
他真的很害怕.
他看著手機上的資訊.
知道自己這輩子都不可能買到房子.
也覺得自己迎合不了這個社會的標準.
他只是躲在網絡的世界後面躺平.
因為只有在螢幕的背後.
只有在虛擬的世界裡.
他才不會被他父母的叫聲嘆氣.
不會被他父母那種失望的眼神所定罪.
第二節我們可以張開眼睛.
我相信這不是某一個家庭的寫照.
這是現在香港社會的縮影.
兩千年前當史奴保羅去寫羅馬書的時候.
羅馬帝國統治之下的社會.
就好像現在的香港.
都是籠罩在低壓槽底下.
這種低壓槽每一個人.
都是被社會的標準日日去審判.
日日去定罪.
中年失業.
在這個世界的標準就是廢物.
年輕人不願意參與殘酷的內卷.
就被嘲笑是廢青.
而做家長因為自己的子女不出眾.
自己就變成了失職的罪人.
每一個人都想用自己的方法去衝出重圍.
結果全家人在同一個屋簷底下各自掙扎.
又彼此去傷害.
今天保羅在羅馬書的八章一至十一節.
要對我們這個充滿了創傷和焦慮的社會.
發出一個不一樣的宣告.

$^{81}$他要告訴我們.
在上帝的系統裡面.
我們不再需要去參與這場瘋狂的遊戲.
今天因為我們有上帝這一載的大靠山.
我們可以不再用自己的方法去死頂.
我們可以生活在一種.
在逆境裡面依然可以挺起我們的胸膛.
做一個充滿了神義和諸賢的一個人.
我們首先來看看羅馬書的八章一至三節.
去找一下我們的生命為何會乾涸.
真正的源頭是甚麼.
我們一起讀一至三節.
預備起.
如今那些在基督耶穌裡的人就不被定罪了.
因為似生命的聖靈的律.
在基督耶穌裡從罪和死的律把你釋放出來.
律法既因肉體軟弱而無能為力.
上帝就差衍自己的兒子成為罪身的樣子.
為了對付罪在肉體中定了罪.
請大家特別留意第三節這一句話.
律法既因肉體軟弱而無能為力.
我們必須明白甚麼是保羅所講的律法.
在聖經裡面律法是上帝所頒佈.
是神聖是公義是良善的標準.
對應到我們現代的生活.
律法可能就是我們口裡經常所講的人生大道理.
家長指南或者是成功必勝法.
這些大道理對不對.
完全對.
做人要有贊感.
我們要為我們的前途著想.
我們要振作.
這些說話本來是出於一份的愛一份的關懷.
我們做家長看到在房間裡面日打夜打的子女.
我們的心裡面很焦急.
想敲門進去想敲醒我們的子女的頭.
你知不知道以前父母是多麼辛苦.
當年我們在你這個年紀的時候就已經擔起整個家.
我們沒可能第一時間就找到我們喜歡的做.
先做著騎牛搵馬.

$^{121}$你現在這樣我們真的很擔心.
不過保羅指出了一個很殘酷的現實.
這些道理這些標準雖然是好.
卻是無能為力.
為什麼.
因為聽大道理的人.
他的肉體即是他的意志力.
他的精神能力他的情緒資本.
早已經在長期的焦慮和高壓底下.
徹底的已經燒盡了.
他不是不知道這些大道理.
他是知道他根本做不到.
因為當標準很高的時候.
而人的軟弱是極度.
而人的肉體是極度軟弱的時候.
原本出於良善的大道理.
在軟弱的人聽起來.
就變成了最殘酷的審判和定罪.
有一位媽媽跟我說.
她得到自己的兒子日打夜打.
在家裡躺平.
她實在不知道怎麼辦.
她心裡一方面怪責自己.
是否沒有教過兒子人生的大道理.
於是她鼓起勇氣衝入房間.
將這些人生大道理全部搬出來.
以為可以用她的道理去叫醒兒子.
結果兒子轉過頭來.
他的眼神沒有憤怒.
只有一種讓人發冷.
令人害怕的空洞.
兒子對她說.
我又沒有叫你生我出來.
你既然這麼辛苦.
為什麼當初你要生我出來呢.
這句話就好像一把刀.
直接插入媽媽的心裡.
當場流血.
這就是律法的無能為力.
大道理不是不好.

$^{161}$不過當一個人的心裡已經乾到.
連明天都不想面對的時候.
你給他再多的正確標準.
對他來說都只是廢話.
不單止不能對症下藥.
反而變成了最令人傷心的控訴.
讓我們的孩子的心鎖得更加緊.
讓在失業的人越來越自卑.
我們這個年代迷信什麼呢.
迷信只要你加倍努力.
你就一定能夠解決問題.
不過我們大家都知道.
我們正面對一個很巨變的時代.
AI海嘯 經濟轉型.
越是我們用肉體去死頂.
只會讓我們更加陷入絕望當中.
保羅說 律法因為人的軟弱做不到的.
上帝親自下場做了.
上帝處理我們的家庭.
處理我們人生的危機.
不是坐在天上冷眼旁觀.
然後扔一本完美的家庭手冊.
成為人生的法則去定我們的罪.
而是他拆印他的親生兒子耶穌.
穿上人的肉體.
走入我們的恐懼.
和我們的失敗無力當中.
他用自己的身體承受了世間所有.
給我們的罪 定我們的罪.
然後宣判說在基督裡面.
我不再定你的罪.
究竟什麼叫榮耀上帝呢.
榮耀上帝的第一步就是.
停止用世俗的成功律法.
去審判和定我們自己.
和定我們身邊的人的罪.
唯有當家人將這些窒息的.
令人窒息的成功指標放下.
我們全家人才會呼吸到.
聖靈第一口的自由空氣.

$^{201}$當我們被恩典釋放.
去定罪之後 第二步.
本來要我們檢視一下.
我們自己的思維系統.
我們的腦平日.
想得最多的是什麼呢.
我們來看看第五至七節.
我們一起讀出 預備起.
因為隨從肉體的人.
體貼肉體的事.
隨從聖靈的人.
體貼聖靈的事.
體貼肉體就是死.
聖靈就是生命和平安.
因為體貼肉體.
就是和上帝為敵.
對上帝的律法不信服.
事實上也無法信服.
提醒我們什麼呢.
要我們將我們的注意力.
去清楚一下.
我們的注意力放在哪裡.
因為我們這個世界.
正處在注意力的.
競爭戰當中 是什麼呢.
我相信大家都聽過搶four.
整個世界都在搶我們的focus.
我們的焦點.
我們的注意力.
你每一天給什麼資訊.
進入你的大腦.
在你的思想裡面.
安排些什麼呢.
你安排些什麼.
會直接決定你的生活.
是在天堂還是地獄.
什麼叫隨從肉體.
隨從肉體的事.
體貼肉體的事.
對於生活在恐懼的家長.

$^{241}$或者失業的中年人來說.
體貼肉體就是.
你的大腦24小時.
都被現實的匱乏感.
綁架了.
你每天睜開眼睛.
或者晚上睡不著的時候.
跳出來的念頭是什麼呢.
會不會是.
我的積蓄還剩下多少呢.
我的安揭怎麼辦呢.
我的求積信.
為什麼還石沉大海.
一封都沒有回音呢.
為什麼人家的兒子這麼出眾呢.
這種想要靠自己去掌控一切.
靠自己的意志去死定的心態.
保羅說.
這就是靈性上.
從上帝唱反調的死胡同.
對於.
關上房門的年輕人來說.
什麼叫體貼肉體呢.
體貼肉體就是.
將你所有的注意力.
完全逃避.
和沉浸在虛擬世界的刺激當中.
因為外面實在太殘酷了.
完全去自我沉迷和封閉.
保羅說這個不是自由.
這個叫精神孤獨.
每天睡在床上.
看著自己的生命去流逝.
到頭來.
只是換來更深的無力感.
抑鬱.
整個人都灰了.
沒有動力.
保羅說.
這是精神的死亡.

$^{281}$即使我是牧師.
你不要以為我很厲害.
一不留神.
我的思想的插頭.
就時時從上帝那裡掉下來.
或者被這個世界的思想所霸佔.
上個星期.
當我想起自己快要退休了.
其實很多弟兄姊妹的美言.
你不像66歲.
其實早就應該退下來的了.
當我看著銀行的儲蓄.
又記起有KOL說.
你沒有幾千萬.
你沒有資格退休.
當你看完之後.
你的心是不是很寒.
當我一體貼肉體.
我的平安立即就消失了.
就好像原本走在水面上.
平平安安的彼得.
走得好好的.
一看著身邊的風浪.
就立即沉下去了.
就在那個時候.
上帝的微聲就提醒我.
你的插頭鬆了.
你可以用你世界的標準.
來審判你自己.
我即時就跟神說.
上帝我承認我的.
我無能為力.
所以我將我的主權.
再一次交給上帝.
當我重新去接駁著.
上帝這個電源的時候.
上帝應許我們.
這個世界奪不走的平安.
又重新回到我的心裡.
最奇妙的是什麼呢.

$^{321}$第二天的靈修.
我正在用一個靈修的app.
所以每一天彈出來的經文.
不是我自己選的.
第二天的經文.
上帝藉著這個經文提醒我.
他說狐狸有洞.
天空的飛鳥有窩.
人子沒有枕首的地方.
是上帝揀選呼召我的.
我只管繼續去遵命而行.
有什麼好擔心呢.
今早的靈修又是另一段經文.
起初書的第五十五章提醒我.
萬物都是祂的恩賜.
上帝掌管一切.
各位弟兄姊妹.
什麼叫做體貼聖靈.
什麼叫順從聖靈.
不是叫你時運高時運高.
什麼都看不到.
不是叫你欺騙自己.
那些單不用交.
順從聖靈的意思是.
拿回你的注意力的主導權.
當全香港的人都因為很怕.
瘋狂內捲.
焦慮不安的時候.
當你來到上帝的面前.
你把你的注意力.
從我不見了什麼.
轉移到上帝今日依然賜我生命.
上帝的恩典救我勇.
你就會有神賜給你的生命和平安.
當然我再講一下.
這個需要操練的.
不是你馬上就聽到聖靈的吩咐.
又或者馬上可以插上插頭.
如果你不是每天都已經建立起.
你靈修的生活節奏的時候.

$^{361}$這個插頭可能已經封了塵很久.
當你接回上帝的時候.
我不是叫你每天苦讀幾十張聖經.
而是當你每天醒的時候.
當你發現恐懼排山倒海.
來侵襲你的時候.
你記得.
你願意花時間.
來向聖靈說.
聖靈啊,我無力去掌管我的人生.
也無力去掌管今日.
我把我的注意力交回給你.
請你賜給我平安.
這個就是重新接駁了.
我記得有一次.
很久很久之前.
我要用三個小時的祈禱.
安靜在神的面前.
禱詞只有兩個字.
上帝,請你給我平安.
跌下來的心.
慢慢覺得實了.
可以面對眼前一切的挑戰.
所以當身邊的同事.
在公事裡面因為很害怕.
在明爭暗鬥,講是講非的時候.
你卻因為接通了.
聖靈給你的平安.
而能夠很誠實.
很腳踏實地去做事.
甚至能夠去安慰身邊.
快要崩潰的同事.
當身邊的家長.
日日因為焦慮.
而對兒女進行情緒勒索的時候.
你卻願意很平靜地.
每天為你的兒女祈禱.
你心裡真的很想告訴你的兒女.
就算你今天什麼都做不到.
你依然是上帝.

$^{401}$你依然是我寶貴的兒女.
今早我在家長群組.
收到一位爸爸的禱文.
這篇禱文很有幫助.
親愛的聖靈.
求你賜下安定和平靜的心.
去幫助我的孩子.
去跨越文化和環境的適應期.
在預備跳躍的日子裡面.
拿開他的焦慮.
深刻看到你為他們量身打造的恩典.
願你的同在成為他們最大的底氣.
勇敢且踏實地走上.
上帝你為他們預備的恩典之路.
奉著耶穌基督的名去祝福我們的孩子.
是有我們需要彼此的支持.
各位弟兄姊妹.
這就是榮耀上帝.
榮耀上帝絕對不是讓我們升職加薪.
讓我們的子女考上名校.
我們才能夠站出來做見證.
不是.
在你最難的環境裡面.
在全香港的人都在焦慮和恐懼裡面沉溺.
你卻能夠靠著聖靈去活出一種不同的人生.
在逆境裡面.
你依然有著跌倒再起身.
和神一起嘆息的堅忍.
這種的活法就叫做「睡」.
就叫做「榮耀神」.
可能你會問.
牧師你講得真是容易.
我現在連呼吸都覺得很累.
我很陰毛.
情緒低落.
毫無動力.
怎麼辦.
這就是為什麼保羅在第九至十一章.
給了我們最堅實的應許.
我們一起來讀十一節.

$^{441}$預備.
使耶穌從死人中復活的上帝的靈.
若住在你們心.
那使基督從死人中復活的也必藉著.
住在你們裡面的聖靈.
使你們必死的身體又活過來.
讓這句話刻在我們的心裡面.
各位聽眾請你認清你的身份.
你不可能因為不能工作.
你學校成績包尾.
你就對自己有懷疑.
保羅告訴你.
上帝的靈是住在你裡面.
你的心不是廢墟.
你的心是天國的大本營.
你背後的大靠山是創造宇宙的上帝.
你根本不需要配合這個世界.
來看低你自己.
你根本不需要配合這個世界.
來看低你自己.
保羅在這裡用了一個很震撼的詞.
叫做「復大能」.
「復大能」的意思不是我們自己裡面.
產生出來的意志力.
而是聖靈住在我們裡面.
祂會重塑我們一份很超然的力量.
大家想一想.
上帝連死亡都可以擊敗.
打爆.
讓耶穌活生生從墳墓裡面走出來.
難道祂搞不定我們眼前的中年失業.
難道祂搞不定一個.
坐在房間裡受傷打機的年輕人.
AI技術沒錯是很聰明.
不過不要忘記.
它只不過是冷冰冰的數據.
經濟環境再差.
它只不過是歷史週期裡的一個低谷.
沒錯世界可以奪走你的職位.
可以拿走你的財富.

$^{481}$不過永遠壓殺不了的.
就是上帝在你體內已經建立起這個復活的能量.
羅馬書第八章最後的經文告訴我們.
整個受造之物雖然一同嘆息勞苦.
不過聖靈會親自用說不出來的嘆息.
去替我們禱告.
所以我們的遂氣不是假裝.
很痛苦去強顏歡笑.
而是知道不論這個世界這個現實多麼完美.
多麼令人疲累.
不過這個裡面復活的力量.
隨時可以讓我們復活.
讓我們在逆境裡仍然繼續可以昂首挺胸.
去活得遂氣.
我們之所以能夠這樣.
是因為上帝是我們的大靠山.
今天這段經文不是一個廉價的心靈雞湯.
它給我們是一個很重要的應許.
上帝要將這股復活的能力.
去注入我們充滿傷痕累累的家庭.
記住昨天我們跟婦女團契說.
我們不需要一步登天.
青年會次下給我們對她更渴望的心.
以致我們渴望要跟她連接.
以及給我們一些微小的力量.
讓我們從今天開始回到家之後.
去做一件很微小的事.
不過是最反叛這個世界的小事.
所以我們面對失業和經濟焦慮.
的父母和任何我們當中的弟兄姊妹.
要講這番說話.
記住你要靠著聖靈.
內在你裡面的力量.
當你回到家的時候.
首先放下那些讓你越看越焦慮的財經新聞.
或者是求職的網站.
你首先去到鏡子的面前.
向著自己說.
我的價值不由市場經濟去決定.
是由洗耶穌復活的上帝所決定.

$^{521}$當你願意在逆境之中去仰望上帝.
重新振作挺起你的胸膛的時候.
你這一種不屈不撓和平安.
就是上帝在全香港人的面前.
最大的榮耀.
在此我也願意給.
對孩子焦慮心痛的家長講以下的話.
靠著聖靈內在你裡面的力量.
今天將你的掌控權和標準放下.
我們不要再試圖去做我們孩子的上帝.
所以當你回到家之後.
當你經過你孩子的房門.
不要去敲門.
不要去訓話.
你站在房門外靜靜的來祈禱.
上主我承認我真的無能為力.
我將我的孩子的前路交給你.
試試讓聖靈賜給你的忍耐.
不要去責備不要給他面色去看.
安安靜靜的為家人煮一餐飯.
當你願意放手去信靠上帝.
這就是你做家長最帥氣的時刻.
接下來我想給在房裡不知所措.
甚至抗拒這個世界的年青人講以下的話.
上帝很愛你.
他一早就已經赦免你了.
他完全知道你多麼的累.
多麼的害怕.
多麼的憂慮.
上帝今天對你說的話就是.
你天天躺在床上看著自己生命的流逝.
其實是太便宜了這個瘋狂的世界.
今天請你靠著聖靈給你一點微小的力量.
試著做一件反悶這個世界的事情.
今晚吃飯的時候.
關掉你的遊戲機或者電話.
十分鐘.
(咳嗽).
出去和你爸爸媽媽一起吃飯.
或者你不想說話也不要緊.

$^{561}$請你們一起吃飯就行了.
當聖靈引導你而你有願意的話.
吃完飯之後去擁抱你的位置.
在外面已經受盡風浪已經快要崩潰的爸爸媽媽.
想像一下.
(咳嗽).
就算只是坐在那裡.
這也是上帝傅大能的彰顯.
(咳嗽).
各位靜姐妹想像一下.
當我失業的爸爸不再用成就去定義自己.
當我焦慮的媽媽學會放手去信靠上帝.
停下閘壁的口和手.
一個本來要躺平的孩子.
願意走出房門坐在飯桌的面前.
當聖靈的愛賜力量給我們.
可以大家一起同行的時候.
(咳嗽).
當我們放下一枚靠自己死頂的重擔.
一個平凡傷痕累累的家庭.
就是神跡發生的現場.
這裡就是上帝得到最高榮耀的時刻.
我們不需要世界幫我們撐腰.
因為我們有上帝做我們的靠山.
因為耶穌.
今天我們一起要聯手去活得遂氣.
榮耀上帝.
我們同心祈禱.
(咳嗽).
聖愛天父.
有時我們感到無能為力.
甚至覺得自己是失敗者.
賜下你的獨生正子耶穌基督.
耶穌來這個世界不是要頂我們的罪.
而是幫我們脫離世界對我們的評價.
脫離綁架我們的焦慮恐懼和不安.
幫助我們將焦點放在你的身上.
嫌棄世界對成功的定義.
這樣我們就能夠有真正的平安.
我們求聖靈賜給我們復活的大能.

$^{601}$讓我們不隨從世界的主流.
即使在嘆息和勞苦當中.
仍能夠活出你給我們的愛和正直.
和身邊同樣支離破碎的人同行.
這樣我們就能真正榮耀你.
我們祈禱在耶穌基督的名義下.
阿們.
\newpage



\section{創世記 25:19-34}
\label{sec:032436}
\textbf{我都就死,這於我有甚麼益處?}~~~~李均熊牧師博士
\newline
\newline
連結: \href{http://www.church.com.hk/acms/content.asp?site=cdc\&op=show\&type=product\&code=032436\&layout=sermon}{\texttt{032436}} ~~~~ 語音日期: 2026-07-12
\newline
\newline
\hyperref[sec:032435]{\small{< < < PREV SERMON < < <}}
~
\hyperref[sec:index_chronic]{\small{[返順時目]}}
~
\hyperref[sec:032437]{\small{> > > NEXT SERMON > > >}}
\newline
\newline
創世記 25:19-34
\newline
\begin{longtable}{cl}
\hline
\hline
章節 & 經文 (和合本修訂版)\\
\hline
25:19 & \begin{tabularx}{0.7\textwidth}{X} 這是亞伯拉罕的兒子以撒的後代。亞伯拉罕生以撒。 \end{tabularx} \\ \\ \relax
25:20 & \begin{tabularx}{0.7\textwidth}{X} 以撒四十歲時娶利百加為妻。利百加是巴旦‧亞蘭地的亞蘭人彼土利的女兒，是亞蘭人拉班的妹妹。 \end{tabularx} \\ \\ \relax
25:21 & \begin{tabularx}{0.7\textwidth}{X} 以撒因他妻子不生育，就為她祈求耶和華。耶和華應允他的祈求，他的妻子利百加就懷了孕。 \end{tabularx} \\ \\ \relax
25:22 & \begin{tabularx}{0.7\textwidth}{X} 胎兒們在她腹中彼此相爭，她就說：「若是如此，我為甚麼會這樣呢？」她就去求問耶和華。 \end{tabularx} \\ \\ \relax
25:23 & \begin{tabularx}{0.7\textwidth}{X} 耶和華對她說：兩國在你腹中；兩族要從你身上分立。這族必強於那族；將來大的要服侍小的。 \end{tabularx} \\ \\ \relax
25:24 & \begin{tabularx}{0.7\textwidth}{X} 到了生產的日期，看哪，腹中是對雙胞胎。 \end{tabularx} \\ \\ \relax
25:25 & \begin{tabularx}{0.7\textwidth}{X} 先出生的身體帶紅，渾身有毛，好像皮衣；他們就給他起名叫以掃。 \end{tabularx} \\ \\ \relax
25:26 & \begin{tabularx}{0.7\textwidth}{X} 隨後，以掃的弟弟也出生，他的手抓住以掃的腳跟，因此給他起名叫雅各。兩個兒子出生時，以撒六十歲。 \end{tabularx} \\ \\ \relax
25:27 & \begin{tabularx}{0.7\textwidth}{X} 兩個孩子漸漸長大，以掃善於打獵，常在田野；雅各為人安靜，常住在帳棚裡。 \end{tabularx} \\ \\ \relax
25:28 & \begin{tabularx}{0.7\textwidth}{X} 以撒愛以掃，因為常吃他的野味；利百加卻愛雅各。 \end{tabularx} \\ \\ \relax
25:29 & \begin{tabularx}{0.7\textwidth}{X} 有一天，雅各熬了湯，以掃從田野回來，疲憊不堪。 \end{tabularx} \\ \\ \relax
25:30 & \begin{tabularx}{0.7\textwidth}{X} 以掃對雅各說：「我累死了，請你讓我吃這紅的，這紅的湯吧！」因此以掃又叫以東。 \end{tabularx} \\ \\ \relax
25:31 & \begin{tabularx}{0.7\textwidth}{X} 雅各說：「你今日把長子的名分賣給我吧。」 \end{tabularx} \\ \\ \relax
25:32 & \begin{tabularx}{0.7\textwidth}{X} 以掃說：「看哪，我快要死了，這長子的名分對我有甚麼用呢？」 \end{tabularx} \\ \\ \relax
25:33 & \begin{tabularx}{0.7\textwidth}{X} 雅各說：「你今日對我起誓吧。」以掃就向他起誓，把長子的名分賣給了雅各。 \end{tabularx} \\ \\ \relax
25:34 & \begin{tabularx}{0.7\textwidth}{X} 於是雅各把餅和豆湯給了以掃，以掃吃喝以後，起來走了。這樣，以掃輕看他長子的名分。 \end{tabularx} \\ \\
[1ex]
\hline
\hline
\end{longtable}
$^{1}$弟兄姊妹平安.
真的要適應一下.
我今早下了床再走上來.
我想深水清的弟兄姊妹都會留意到.
這段時間的舊約經火都是選自創世記.
由亞伯拉罕到今日的雅各爾蘇.
接下來會繼續背讀創世記的經文.
在今日的舊約經火當中.
雅各蘇說了一句說話.
「我都就死了,這於我有什麼益處?」.
這句說話聽上去像上一輩的朋友經常說.
「信耶穌難道有飯吃嗎?」.
大家聽過這句說話嗎?.
可能你都聽過你爸爸祖父輩說過這句說話.
這句說話表面很世俗很功利.
甚至你會覺得很無賴.
但有時正正是這種說話.
說到人心底裡最深切的關注.
信仰 沒錯是說應許 說將來 說上主的帶領.
但人活在當下 活在眼前.
會肚餓 會累 會害怕 會覺得撐不住.
當眼前的需要太近.
將來的價值就變得很遠.
所以如果我們太快否定這句說話.
「我都快死了,這個長子名分於我有什麼益處?」.
很快將它讀成一個反面教材的話.
「華爾素他貪食短視 輕看屬靈的事」.
其實傳統都會這樣理解.
你見到今天的經課選讀.
羅馬書正正提醒我們.
不要掛著肉體 要掛著聖靈.
但如果我們太快否定.
以素所說的話.
或經驗他的想法的話.
可能這段經文對我們的意義就不會太大.
創世記二十五章的選讀.
大家可以打開來慢慢看.
它開始是說以撒亞伯拉罕的兒子.
然後說到後面.
你會發覺以素還未出世 鬥爭已經開始.

$^{41}$兩個兄弟在武福中相爭.
有些光所以我看不到稿.
以素雅各未出世已經在武福中相爭.
出世的時候一個先出來.
另一個捉住哥哥的腳跟.
於是人們就將他改成雅各.
雅各的意思就是「猥實」.
而這樣看來這場不是普通的兄弟相爭.
不是普通的兄弟爭執.
如果我們一直看下去的話.
我們會發覺整個故事是說.
一段關於身份 位置 承繼權的家庭糾紛.
經文也說了.
兩個國家在利伯加的肚子裡.
學者提醒我們.
雅各和以素從來不是兩個人的故事.
以素後來是以東民族的起源.
與以東民族的起源重疊.
以東是後來在新約時期的底土馬.
也是大希律的民族.
除此之外.
雅各更是以色列十二支派的祖先的爸爸.
他是祖先的祖先.
所以這不是純粹兩個人家庭的糾紛.
而是兩個大民族之間有次兄弟有次敵人的關係.
如果你留意一下歷史記載的話.
以東和以色列基本上是一個世仇.
但後來他們又彼此合作.
他們的關係又是次敵人又是兄弟.
創世記一直記載.
這幾段家庭鬧劇.
從母腹中的掙扎.
出生時的抓緊.
長子名分的交易.
祝福的爭奪.
都在預示著一段關係是怎樣拉長進入漫長的歷史.
所以二素不是一個輕易放在反面教材的人.
從經文中我們可以看見.
他是有體毛有顏色.
是一個屬於田野的人.

$^{81}$他從外面回來累到一個地步.
為了眼前的紅豆湯.
就比任何長遠的身份都覺得更加真實.
其實二素甚至可以說是直率.
不會想太多太遠.
但雅各也不是一個單純的英雄.
他懂得安靜地等待.
懂得看時機.
懂得在有限的資源裡算計.
這種生存方式.
我想香港人都很明白.
有時我們會稱這些人為香港仔.
懂得抓緊時機.
懂得運籌為額.
經文不是叫我們立刻對比誰比較熟稔.
誰比較好誰比較壞.
正正就是將兩種不同的生活模式.
放在同一個家庭裡.
兩個不同的人在家庭裡.
遲早都會出事.
遲早都會有衝突發生.
正所謂我們常說.
沒有比較沒有傷害.
兩個不同的人放在家裡.
遲早都會有衝突發生.
所以當二素說那句說話.
我都快死了.
這個長子名分對我有什麼益處呢?.
我聽到的不是一句低俗無賴又愚蠢的說話.
而是一個人被疲乏壓倒.
剩下眼前的需要.
長子名分不是不重要.
應許不是不重要.
但那一刻他不能容易滿足當時他即時的需要.
不容易回答他過度的疲累.
我們要想想.
究竟長子名分為什麼這麼重要呢?.
在古代世界裡.
這不是一個小問題.
今天我們很多人關心個人的抉擇.

$^{121}$個人前途.
職場的選擇.
我們為DSE的同學禱告.
決定讀哪間大學.
讀什麼科目.
全部都是個人的決定.
我們甚至用這種個人取捨的觀念.
來理解我們的人生.
理解我們前途.
但在古代氏族社會.
一個長子名分.
不是分多一點家產這麼簡單.
而是承繼整個家族的名號.
記憶.
甚至是家族的遺傳.
所謂我們常說的legacy.
現在我們有些基金強調.
將你的legacy傳下去.
所以長子名分在當時氏族社會裡.
不是這麼簡單.
一個名分這麼簡單.
不是因為家產分多一點.
而是你承繼家族的legacy.
家族的名號.
令家族在一個古代艱難的世界裡.
繼續存在.
令家族這個姓氏.
不會在艱難的世界裡消失.
所以其實長子名分的重要性.
在古代世界.
在一個氏族社會裡是很重要的.
我不知道大家有沒有從遺傳那裡出來.
或者回到鄉下.
他們的氏族.
誰是長子.
誰是兒子.
兒子才可以繼承家族.
因為他才會生.
兒子才會跟隨姓氏.
所以這些都是很重要.

$^{161}$我不知道現代家庭.
我們在一個核心家庭長大.
特別在現代社會強調個人的自決.
這種對於家族的承擔.
對於姓氏.
對於所謂宗族的一些想法.
究竟我們了解有多少.
認識有多少.
看不看到那種深刻.
或者那種長子名分的不簡單.
最近我煲劇.
煲權力遊戲.
我不知道大家知不知道.
很兒童不宜的.
十八歲以上才好看.
權力遊戲.
因為最近他又有另一套.
推出叫做龍之家族.
我不知道大家知不知道我在說什麼.
權力遊戲.
或者龍之家族.
這樣的電視劇.
就是說.
他借用中世紀的.
所謂封建社會的制度.
一個封地.
一個領主.
他怎樣繼承.
他有沒有下一代去繼承.
所謂婚嫁.
生不生子.
是女還是兒子.
這些是整個故事裡面.
很重要的一件事.
或者可能我們想想.
舊世紀那些大閘門.
電視劇經常都說這些.
離我們好像很遠.
但在古代社會.
在耳塞瓦割的故事裡面.

$^{201}$卻是一個很重要的元素.
如果你不掌握這個長子名分.
對於整個家族的重要性.
你不明白.
為什麼雅各要這樣欺騙.
不是欺騙.
應該是成人之危.
想拿走這個長子名分.
而以素興翰長子名分.
又是怎樣讓後世覺得.
被唾棄不正確的一件事.
所以長子名分不是你私人喜好.
而是整個家庭故事的核心.
甚至你看回創世記的話.
亞伯拉罕,耳塞瓦割的故事裡.
這個都不是人間的制度這麼簡單.
因為這個承接著上帝的應許.
上帝對亞伯拉罕.
說賜他後裔土地.
令他祝福萬民.
而不是放在天上的一句話.
而是實際在一個被呼召家族當中.
承載和延續的.
所以長子名分.
不單止關乎一個個人.
它是超越了整個個人層面.
關乎一個家.
甚至乎是關乎上主的應許.
能否延續下去.
對於亞伯拉罕的家族來說.
這個長子名分.
還是承載著上主的應許.
上主的祝福.
所以用一碗紅豆湯.
將長子名分交換.
不是我們想像中這麼簡單.
是比我們所想像的更加複雜.
一邊是馬上入口的食物.
一邊是用很多年.
用歷史.

$^{241}$在將來.
在後世.
你才會見到它的果效.
一個人承受長子名分.
不是分多一點.
而是承載整個家族的名號.
責任,記憶.
整個家族的遺傳,遺產.
將上一代和下一代連接起來.
牽涉一個家族.
在古代艱難的世界裡.
如何繼續存在.
所以耳熟聞那句.
究竟有什麼益處.
不是因為長子名分沒有益處.
它的益處其實很大.
但它的益處是太長.
太遠.
太依賴將來.
對他當時面對眼前的需要.
很迫切的需要來說.
這些對他來說都好像太遙遠了.
重要的東西仍然都重要.
但重要的東西.
不等於我們每個人.
每一刻都感受到它是這麼重要的.
信仰也一樣.
我們講呼召.
講基督徒的身份.
講上主的同在.
這些都是無比重要的.
只是當我們在疲累裡.
當我們在壓力裡.
當我們受生活的催逼.
就未必會感受到它有多真實.
也未必會感受到它有什麼即時的益處.
所以上一代經常說.
信耶穌難道有飯吃嗎?.
其實是反映他們的生活實在太艱苦.
太艱難.

$^{281}$太貧困.
基督的信仰未必能夠即時帶來滿足.
所以我們不要隨便說這句話.
拿來開玩笑.
信耶穌有飯吃嗎?.
到了這裡.
基督徒聽到自然會問.
這個故事教訓我們什麼?.
究竟哪個對哪個錯?.
就好像我們小時候看電視劇.
很想知道哪個是壞人哪個是好人.
我繼續看下去才能搞清楚他們的關係.
好人我支持他.
壞人我就推翻他.
但是今天的戲就不是這麼簡單.
我很喜歡看電視劇.
很喜歡看劇集.
現在的電視劇.
你不知道哪個是中哪個是奸.
到最後中的會變奸.
奸的會變中的.
其實這個對於我們現代人來說.
看聖經的故事是很正確的.
因為聖經不是告訴我們哪個是中哪個是奸.
它不是一個道德的榜樣這麼簡單.
所以當耶穌說.
我都快死了.
這個長子的名分與我有什麼處呢?.
我們不要以為是一句低俗無聊愚蠢的話.
一個人被疲累壓倒.
只剩下眼前的需要.
什麼長子的名分都不重要.
應許都不重要.
在那一刻.
什麼滿足到我們即時的需要呢?.
才是最重要的.
創世記的記載沒有告訴我們.
哪個是對哪個是錯.
其實.
經文沒有說將雅各寫成.

$^{321}$因為異行所以被揀選.
也沒有因為耳掃.
因為什麼惡事而被拒絕.
我們要小心.
不要以為被揀選就是得勝的證明.
就是他是好人所以上帝揀選他.
失去就變成了懲罰.
經文一直沒有這麼簡單的告訴我們.
哪個是對哪個是錯.
雅各的行動為了抓住.
乘虛而入.
乘人之危.
耳掃的行動.
帶著輕率.
受困於當下的真實.
但其實我們也不能說他是直率的人.
經文到最後就說.
耳掃輕看了長子的名分.
但他沒有因此說雅各的手段是對的.
沒有將雅各的手段洗白.
很多決定都是這樣做出來的.
就是眼前的感受.
往往比我們認為重要的東西.
更加重要更加即時.
當我們頂不住撐不住壓力太大的時候.
人就只會處理眼前的問題.
重要的東西沒有消失.
不過暫時就放在後面.
遲些再說.
長子的名分在這裡.
或甚至上帝應許在這裡.
都是一樣的.
對那一刻來說.
耳掃一碗紅豆湯.
是更加有益的.
是不是.
這不是一個計較的問題.
是我們在人生某些情況下.
都可能會問的一個問題.
所以耳掃不是一個遠古的人物.

$^{361}$我們今天看的不是一個遠古的家庭鬧劇.
而是其實是說今天我們的人生.
今天我們的人性.
我們未必會見到一個衝動世俗.
傾向屬靈福分的失敗哥哥.
反而我們見到的是人的境況.
什麼東西重要.
什麼東西可以放棄.
信仰是重要.
關係是重要.
承諾是重要.
但是當我們被疲累.
壓力 焦慮 壓垮的時候.
太靠近的時候.
那些覺得重要的東西.
信仰告訴我們重要的東西.
好像突然間變得很遠.
它不是消失了.
但是那一刻我們承受不了.
故事並沒有停在紅豆湯裡.
甚至最後那句話.
耳掃興翰了自己的長子名分.
並不是完結在這裡.
如果你讀過《創世記》.
或者我剛才所說.
你繼續追我們的舊約經課的話.
你會發覺.
事情並不是停在那裡.
《創世記》繼續向前走.
如果你讀過《創世記》之後.
你會知道雅各後來是怎樣的.
他逃亡.
他被欺騙 被算計.
被他和他差不多.
都認識成人之危的.
拉班舅父騙了14年.
但是如果你繼續看下去的話.
耳掃和雅各兩個兄弟.
他們會再相遇.
一起埋葬他們的父親.

$^{401}$《創世記》給我們看到的.
應許永遠不是一條直線.
今天選擇對了就一直都對.
人的選擇是有後果的.
使數交易不是無關痛癢的.
雅各的所謂賄賂.
都是要付出代價的.
上主沒有因為他們這樣去做.
而停下他的工作.
我不知道說到這裡.
你會不會想起自己的家.
我們會不會都覺得.
家裡有些衝突.
其實不是那麼劇烈的.
不會一開始見到面就打架.
很少是這樣的.
但是你可以想一下.
有些事情就是不願意提.
有些事情從小覺得你偏.
爸爸媽媽偏心你 偏心他.
為什麼爸爸媽媽有事的時候.
要我給多點錢.
要我照顧他多點.
是不是.
什麼是重要 什麼是不重要的呢.
我付出有沒有益處呢.
爸爸媽媽偏心.
這些事情可能你一直在心裡.
沒有說出來.
每次見面都忍住那句話.
到爸爸媽媽有事的時候.
你就會問.
我為家庭付出對我有什麼益處呢.
這個不是創世記的鬧劇.
這個基本上我們在社會.
在我們家庭會見到的一些事情.
不能夠不否認.
創世記沒有寫到這個家庭.
所有問題解決了.
上主才將他應許工作繼續完成.

$^{441}$上主是在一個有問題的家庭裡.
實踐祂的祝福和應許.
所以講到最後.
我想和大家回到我們今天講題的那句話.
或者開始所講的那句話.
我都快要死了.
這些對我有什麼益處呢.
信耶穌真的可以開飯嗎.
因為我們有時候會用不同的方法.
問類似的問題.
祈禱是不是可以解決我現在的情況呢.
有什麼用呢.
我讀到聖經.
上班是不是會容易一點呢.
對我來說可能是.
因為我是教神學院的.
但對你來說未必是.
上教會是不是就可以改善家庭關係呢.
每個星期回來有什麼意思呢.
做人這麼辛苦究竟上帝去了哪裡呢.
當眼前的壓力太大.
問題太迫切的時候.
我們可能都會用耳熟的說話.
但是形式未必一樣.
但我們其實在問同一個問題.
信耶穌有什麼益處呢.
這些是上帝應許的同在.
對我來說這一刻有什麼幫助呢.
這不是沒有信心的表現.
這是人之常情.
甚至我很建議大家有空問這個問題.
搞清楚為什麼要堅持在信仰裡.
堅持跟從上帝.
堅持在基督做他的門徒.
所以我們不要太快.
去為一個關係.
一句說話去下一個判語.
因為故事未完.
上帝一直沒有離開過.
今天我們未必看得明白.

$^{481}$未必感受得到.
但上主的應許仍然在故事裡.
在這個家庭裡.
在亞伯拉罕 以撒雅各的家庭裡.
雅各還要面對欺騙帶來的逃亡.
面對舅父拉班的另一種算計和出賣.
甚至到最後.
如果你熟悉的話.
他要跟上帝.
跟自己.
跟耳塞的影去摔跤.
因為他要渡過那個雅各渡口.
去見到耳塞.
去清算他們兄弟的恩怨.
創世記是很精彩的.
有空有機會和大家一起看.
但告訴我們.
不是今天我們做對了.
明天就會清楚.
也不會告訴我們.
今天我們做錯了.
明天就會出局.
永世都不能翻身.
應許不是一個整整齊齊的家庭裡發生的.
應許是在人間很混亂的世俗裡.
上帝在當中工作去成就.
人的選擇是會帶來後果的.
但這一點不能取消上帝的工作.
耳塞交易不是無關痛癢的.
雅各的死活.
也不是沒有代價的.
但應許告訴我們.
眾言人有種種的特質.
上帝沒有離開我們.
起碼他沒有離開創世記裡.
亞伯拉罕整個家庭.
到後來出埃及記.
到後來被擄回歸的時候.
他們都說這是亞伯拉罕耳塞雅各的神.
為什麼我們可以一直承傳下去呢.

$^{521}$我們家庭未必有劇烈的衝突.
但隔一段時間我們的關係可能會很隔離.
無論如何我們要明白.
不要太快為家人的關係下判斷.
就像創世記耳塞雅各.
不要太快判斷誰是對誰是錯.
因為這樣是沒有好處的.
反而我們停留在這裡.
慢慢再看上主的工作的時候.
我們反而會更加明白.
上帝在混亂的人間.
在我們的張力當中.
來成就他一切的偉大.
耳塞和雅各其實離我們不是太遠.
我們家庭可能未必會完全反目成仇.
但也未必有真正的和好.
但關係在這裡.
創世記沒有將這個家庭.
最後寫到大家和和氣氣.
齊歡暢和慶賀.
將問題全部解決.
然後上主才繼續工作.
相反上主正在這個仍然有距離.
仍然有很多舊的傷痕.
仍然未完全清楚.
關係如何復和和解.
而繼續前行.
這個不是叫我們將發生的事情.
一筆勾銷.
然後說神有計劃.
也不是將我們人生的傷痛取消.
相反是提醒我們.
人的家庭的故事.
人生的故事是很混亂的.
今天我們感受不到上主祝福的益處.
不代表我們可以取消上帝的恩典和幫助.
我們的膚照.
我們的身份.
上主的同在.
不會因為我們即時的需要.

$^{561}$滿足不了.
就變得沒有價值.
所以到最後.
我就想你們都問一問這個問題.
如果有一天.
你覺得我都快死了.
這些一切對我來說有什麼益處呢?.
就像我剛才所說的.
這麼辛苦祈禱.
是不是就可以幫到我呢?.
是不是讀多點聖經.
上班就沒那麼辛苦呢?.
是不是上教會.
我就和家人,朋友,同事的關係就會改善呢?.
未必.
我甚至告訴你.
通常都不會.
不要這樣期望.
但是雅各和耳俗的故事.
其實告訴我.
這些就是我們家人的故事.
就是我們生活的故事.
很多的恩怨情仇還沒說完.
但是上主沒有離席.
上主一直都在看著.
就算今天我們未必看到.
未必感受得到.
但是上主的應許仍然在我們生命裡工作.
你又相不相信呢?.
我們一起祈禱吧.
主人我們有時會覺得.
聖經就是道德的故事.
誰對誰錯.
我們應該跟誰.
應該學誰.
但是人生其實.
比經文所記載的更加複雜.
在我們家庭裡.
在我們人生裡.
在我們的抉擇裡.

$^{601}$有時會做錯.
在關係裡有時會緊張.
我們有時會憎恨某些人.
有時也會因為偏心.
會喜歡某些人.
不喜歡某些人.
但是讓我們看見.
在這些人生的混亂裡.
你沒有離棄我們.
正如荒書所說的.
羅馬書所說的.
我們在你裡面不被定罪.
求主幫助我們繼續向前行.
有勇氣.
有勇敢.
更加相信.
就是你的應許.
你的祝福.
是有益處的.
是有價值的.
今天我們看不到.
今天我們體會不到.
但是讓我們能夠.
繼續堅持去跟從你.
去見證你的偉大和榮美.
我們這樣禱告.
交託是奉主耶穌基督的名字.
阿們.
\newpage



\section{詩篇 119:49-56}
\label{sec:032437}
\textbf{靈命成長筆記 (2) 回歸永活的聖言}~~~~陸超明牧師博士
\newline
\newline
連結: \href{http://www.church.com.hk/acms/content.asp?site=cdc\&op=show\&type=product\&code=032437\&layout=sermon}{\texttt{032437}} ~~~~ 語音日期: 2026-07-12
\newline
\newline
\hyperref[sec:032436]{\small{< < < PREV SERMON < < <}}
~
\hyperref[sec:index_chronic]{\small{[返順時目]}}
~
\hyperref[sec:032438]{\small{> > > NEXT SERMON > > >}}
\newline
\newline
詩篇 119:49-56
\newline
\begin{longtable}{cl}
\hline
\hline
章節 & 經文 (和合本修訂版)\\
\hline
119:49 & \begin{tabularx}{0.7\textwidth}{X} 求你記念你向僕人所說的話，這話使我有盼望。 \end{tabularx} \\ \\ \relax
119:50 & \begin{tabularx}{0.7\textwidth}{X} 你的話將我救活了；這是我在患難中的安慰。 \end{tabularx} \\ \\ \relax
119:51 & \begin{tabularx}{0.7\textwidth}{X} 驕傲的人極度地侮慢我，我卻未曾偏離你的律法。 \end{tabularx} \\ \\ \relax
119:52 & \begin{tabularx}{0.7\textwidth}{X} 耶和華啊，我記念你從古以來的典章，就得了安慰。 \end{tabularx} \\ \\ \relax
119:53 & \begin{tabularx}{0.7\textwidth}{X} 我因惡人離棄你的律法，怒火中燒。 \end{tabularx} \\ \\ \relax
119:54 & \begin{tabularx}{0.7\textwidth}{X} 我在世寄居，以你的律例為詩歌。 \end{tabularx} \\ \\ \relax
119:55 & \begin{tabularx}{0.7\textwidth}{X} 耶和華啊，我夜間記念你的名，我也要遵守你的律法。 \end{tabularx} \\ \\ \relax
119:56 & \begin{tabularx}{0.7\textwidth}{X} 這臨到我，是因我謹守你的訓詞。 \end{tabularx} \\ \\
[1ex]
\hline
\hline
\end{longtable}
$^{1}$各位記者你們好.
我從來沒有試過連續兩個主日在一個教會裡面講道.
希望你們不會覺得悶 又見到我.
今天實在是很開心跟大家分享我自己的一些看法.
其實上個禮拜的講道.
還有今天跟大家分享的訊息.
其實你可以很濃縮地看到.
我的靈命成長的看法和路向.
很少有人在這麼短的時間.
將我要講的東西都可以講給大家知道.
希望大家都能夠在當中.
真的能夠得到上帝的祝福.
上個禮拜我講的題目你還記不記得.
通常都不會記得 一個小時之後已經忘記了.
何況是一個禮拜.
叫做定格永恆的生命.
上個禮拜我的主題是強調了生命.
神聖的生命.
今天我跟大家講的題目叫做.
回歸永活的聖言.
我的焦點放在聖言 即是聖經的說話.
我很肯定沒有一個基督徒會否定聖經的重要性.
今天我的側重點.
將一個好像很冰冷的一個題目.
或者好像是側重了頭腦知識的一個課題.
我試圖換一個新的面貌.
呈現在大家的面前.
讓大家看得到一件事.
原來聖經是一個對我們活生生的道理.
絕對不是一本書.
絕對不是一些很枯燥的事情.
而是上帝給信徒.
我們每一個人極之寶貴的一個禮物.
如果你能夠明白到這一點的話.
將你整個人生改變過.
值得我們一生去追求.
我們先來做一個祈禱.
你將你的話.
藉著聖靈的引導寫在這本聖經.
將你的啟示直接告訴我們.

$^{41}$求主叫我們不單止要讀聖經.
不單止看為是一種責任.
而是更加要珍惜這本聖經.
開了我們的心橋.
讓我們成為我們手中握著聖靈的寶劍.
為我們的主打個美好的著.
我們的禱告是奉耶穌基督的命求.
感恩.
海岸有個地方叫做Bell Rock.
Bell Rock原來是一種礁石的小島.
在潮漲的時候是看不到的.
因為水位很高.
但潮退的時候.
這些礁石的小島就會出現.
在這樣的情況下其實對於航海是很危險的.
因為你一中了礁之後你就會沉船.
大概是18世紀末期的時候.
有人計算過每一年在Bell Rock的地方.
都發生起碼六次沉船的意外.
當時的政府就決定在那裡建一個燈塔.
在19世紀初期建一個燈塔.
來警告周圍的船隻要小心.
不要靠近它.
並且有指示方向讓那些船隻應該怎樣走.
直到大概20世紀初期.
自從第一次世界大戰之後.
1914至1918年.
第一次世界大戰之後.
這個地方只發生過一次有沉船的意外.
原因是因為打仗所以完全不能有燈.
所以就發生意外.
各位 其實我們的人生往往都是在一個.
好像波濤洶湧的大海裡面.
一隻很小的船.
其實那個經歷是非常危險的.
我們看前路的時候好像很絕望.
所以我們希望能夠找到有一種希望在前面.
更加最重要的就是能夠邁向一個出路.
在基督徒來說.
我們的出路是甚麼呢.

$^{81}$我們的出路就是神的說話.
神的話就是我們的明燈.
神的話就是我們的燈塔.
指向一條道路.
叫我們脫離永遠的死亡.
而且得著一個救贖的機會.
在詩篇119篇的81節他是這麼說的.
我渴想你的救恩.
仰望你所應許的華語.
清清楚楚告訴我們知道.
上帝的華語是多麼寶貴.
請大家去設計聖經最重要最重要的功用在哪裡.
不是讓我們的頭腦裡面多一些的學識在打滾.
不是,他的作用在羅馬書15章第4節.
保羅講得很清楚.
他說從前所寫的聖經都是為教訓我們而寫的.
叫我們因為聖經所生的忍耐和安慰.
可以得著盼望.
這裡有幾個「虹」字是相當之重要.
我們要特別留意.
我們的神口出的說話.
在聖經裡面講得出來的時候.
就必定會成就.
他說從前所寫的聖經在保羅的時代.
從前所寫的聖經就是紙巾舊約聖經.
到今天我們新約時代的信徒.
包括了新約聖經在內.
都是這一半的聖經.
都是為教訓我們而寫的.
為什麼要教訓呢?.
就是因為我們不知道上帝的心意.
所以要教訓我們而寫.
得到的後果是什麼呢?.
是因聖經所生的忍耐和安慰.
我們會問為什麼要得到忍耐和安慰呢?.
原因很簡單.
因為世界上實在太多非常艱難的處境.
還要說人生不如意的事.
十常八九.
特別當我們的弟兄姊妹稍為年長一些.

$^{121}$或者是中年以上的弟兄姊妹.
特別都會有很多的經歷.
就提醒我們一句.
不如意的事是十常八九.
一點都不奇怪.
無論你是未信耶穌也好.
你是信了耶穌的信徒也好.
都是有這樣的問題.
所以情況就是如何去面對.
你又有問題.
我又有問題.
問題是什麼呢?.
就是你如何去面對.
差別就在這裡.
在《賽亞書》43章第二節的上半節.
他說:「你從水中經過,我必與你同在.」.
這個「我」是指神.
「我」是指耶和華神.
在43章的第二節.
在下面有說到.
「水不能夠慢過你,火不能夠燒掉你.」.
詩篇66章的12節.
有提到一個相類似的說話.
他說:「我們經過水火,.
你卻使我們到豐富之地.」.
經過水火就是經過困難.
到達一個豐富之地.
就是去到一個出乎意外的經歷.
弟兄姊妹.
當我們信了耶穌之後.
你有沒有試過有這樣的經歷呢?.
經過水火是很慘的.
但是經過了之後.
回頭一望.
原來是已經去到豐富之地.
是誰帶領我們有這樣的經歷呢?.
乃是神自己.
因為神就是掌管一切的那一位.
聖經的說話.
是給我們最大的安慰.

$^{161}$我們今天親眼看不到耶穌.
但是我們拿著聖經.
我們就可以好像面對面.
與神有一個直接的溝通.
剛才我們也唱過這首詩歌.
叫做「我心靈得安寧」.
我相信很多人都唱過.
但是你唱這首詩歌的時候.
如果你略為知道這首詩歌的背景.
你會更加明白.
為什麼這個作者有這麼深刻的領受.
將天上降臨的安慰的訊息.
帶到人間就可以得享一個平安的確具.
這首詩歌的主角.
就是 Horatio Spetford 先生.
他是一個十九世紀的人物.
史柏福先生.
他經歷了人生很大的苦難.
他寫成這首詩歌.
其實他經過的是事業的失敗.
最慘的就是他連家人差不多死光了.
在他悲痛之餘.
他就藉著這首詩歌寫出.
在主裡面原來帶給他真正的平安.
而且強調主在來的時候.
得到最大的盼望.
史柏福先生在十九世紀.
在美國芝加哥進行地產生意的一位律師.
他的遭遇真是可謂禍不單行.
不知道你有沒有經歷過.
或者你聽到有人禍不單行.
往往都是這樣的.
不是一次.
不是一次半次.
而是一次又一次又一次又一次.
可能你是其中一個經歷過這樣的狀態.
他的兒子首先四歲的時候離世.
1871年的時候.
他所經營的生意在芝加哥發生巨大的火災.
燒了他很多很多的財產.

$^{201}$兩年之後1873年.
他因為美國經濟很蕭條.
他的業務是非常艱難之際.
他就計劃不如這樣.
他和他家人一起去歐洲拜訪.
當時在英國的時候.
有一位出名的報道家叫做慕迪先生.
希望能夠得到一些指引.
再繼續計劃以後的前路.
但很奇怪.
因為有一些緊急的業務.
他不能夠離開.
所以他就叫他太太帶著他四個女兒上船.
先去歐洲然後下一班船.
他才去到那裡和他們會合.
當他太太和他四個女兒搭的船離開了之後.
就在北太大西洋與另外一艘船衝撞.
那艘船很快就沉沒在大海當中.
四個女兒全部撞上大海.
獨是他的太太能夠保存了他的生命.
史伯福先生收到他太太的電報.
只有兩個英文字.
Saved Alone.
僅僅我能夠生存.
他萬分悲痛之餘就立刻訂另外一班船的船票.
希望能夠和他太太會面.
當那艘船慢慢駛到他四個女兒沉下去的海域的時候.
在心如刀割的時刻.
他仰望著那些海水.
離開,越去他女兒的海面.
他沒有憤怒也沒有絕望.
心中卻是湧現了非常之大的平安.
就寫了這首「我心靈得安寧」這首詩歌.
在一百多年來.
這首詩歌傳到世界各地.
張史伯福先生很誠實去面對困難.
卻是對上帝有堅定的信仰.
這個見證充分地表達了出來.
每當有人身處於痛苦,絕望,失望之餘.
這首詩歌總會觸動,總會安慰,總會醫治人心.

$^{241}$史伯福先生很珍貴的領袖不是憑空而來的.
如果你稍微留意他詩歌的歌詞.
原來都是有聖經的根據.
都是從聖經的教訓裡面演繹出來.
對於我們的心靈是帶動很透徹的共鳴.
那個共鳴是甚麼呢?.
是無論在任何一個環境也好.
主賜下的平安就是最為寶貴.
主救贖實在是最大的恩情.
幫助了我們這群受害,受苦的人能夠得到釋放.
更加期待他們快點回來.
以致我們全權的可以得贖.
因此我們學習的重點在史伯福先生這首詩歌裡面.
就是要緊靠上帝的信仰來渡過我們人生每一種困局.
無論大困局,小困局.
都是要依靠神去渡過.
身處痛苦的時候,失意的時候.
仍然得到安慰和醫治.
他的經歷其實今日同樣我們都可以有這樣的經歷.
今日我要和大家分享的那個聖經的經文.
出自詩篇119篇.
詩篇119篇有一些很特別的特色.
不說你們未必會知道.
它是全本聖經最長的一章.
它也有唯一單一的焦點就是放在上帝的話語裡面.
整篇119都是在說上帝的話語.
那個訊息是很特別的.
它有幾個特色.
可能簡單介紹一下.
它的排列是很工整的.
工整到什麼地步我們根本都不可以想像.
它總共有22個段落.
每一個段落都是很工整.
都是一篇有8節的聖經.
有8句的說話.
而每一句的說話裡面的頭一個字的字母.
都是希伯來文字裡面的字母.
按著次序排出來的.
舉個例子第一個段落裡面的第一個字就是XXX然後大X.
希伯來字是從右向左的.

$^{281}$不同英文的英文是從左向右的.
希伯來字就是右向左.
它第一個說話的第一個字就是那個紅色的交叉.
那個是希伯來字母裡面的第一個字母.
到第二個段落又有同一個字XX又有一個紅色的X.
第二個紅色的字母就是希伯來字母的第二個.
如此類推下去.
所以是很工整的.
不是隨便來的.
給我們有一個很大別的很特別的一個感受就是.
上帝要特別向祂的人指紋.
要啟示他一些說話.
特別書寫出來.
詩篇119篇只有一個核心.
就是用八種的方法.
用八種的方法來表達上帝的啟示.
它用的字眼可以是律法.
可以是法道.
可以是判語.
可以是訓詞.
可以是律令.
可以是命令.
可以用華這個字.
用了兩個華這個字華一或者華二.
我特別要講講華一和華二是什麼來的.
很特別的.
華那個字第一個的解法就是Imra.
它的解法是什麼?.
是神口中.
口.
口裡面講的說話.
發出來的說話.
是上帝的言語和上帝的應許.
好像詩篇119篇的十一節說.
我將你的華藏在我心裡.
就是用這個Imra這個字.
但是第二個字.
希伯來的字眼讀華的字有第二個意思.
就是Taba這個字眼.
它的意思就是上帝的說話的總稱.

$^{321}$一個總稱.
人行動的一個指標.
好像是八十九節裡面說.
你的華安定在天.
那個就用Taba這個字眼.
所以如果我們略為看到原本的文字的描述.
我們明白詩篇119是更加精彩的.
要很努力的.
你就會明白更加更加多.
兩大主題.
神的話和神自己.
今天早上我們選了第七個段落裡面的四十九至五十六節.
為什麼會選這一段呢?.
有原因的.
因為這一段特別講到人遇到困難的時候.
怎樣靠我們這位神.
怎樣靠神的話來面對.
所以如果你看這一段的說話.
其實都可以分為三個小的段落.
頭兩節四十九和五十節.
是講到回顧詩篇這個作者.
這個詩篇的作者很可能是大衛寫的.
是大衛困難的時候.
大衛流放的時候.
他就定睛在上帝的應許底下來面對.
因此得到平安.
所以第一個段落講到他回顧.
回顧就是對神的一種肯定.
第二個段落五十一節到五十二節.
是講到他的掙扎.
他不是超人.
他一樣好像和我們一樣.
是有掙扎.
掙扎是說你要怎樣處理所遇到的艱難.
或者困境.
最後那四節就是第三個段落.
就是再一次肯定.
再一次肯定表現出這個詩人可能是大衛.
他堅定的信仰究竟是怎樣.
往往我們遇到困難的時候.

$^{361}$我們就可以根據這個圖案.
回顧上帝是誰.
清楚明白去了解我們當中的掙扎.
告訴神知道.
而第三再一次肯定.
神會帶領我們走出這個困境.
我們先看第一個段落.
第一個段落是四十九和五十節.
這裡同一節聖經.
話那一節出現了兩次.
第一個就是你僕人所應許的"話".
那個"話"就是上帝話語的總稱.
相似於新約聖經的logos.
道的意思.
上帝的話叫這個詩人.
有產生盼望.
第五十節講到.
他的"話"將他救活了.
這個就是"Inra".
是另一個翻譯.
這個"Inra"講到上帝口中真實講出來.
那種對他的回應.
讓他在患難當中得到安慰.
所以這裡帶給我們看三個重點.
有盼望被救活和得安慰.
以致他有希望可以繼續生存下去.
可以被救活和平復了他心中的鬱結.
其實這三樣東西.
其實都發生在人的困惑當中.
人有困惑沒什麼出奇.
但在困惑當中你要認清楚.
和要回顧.
究竟哪個人可以帶你離開這個困境.
我明白到"Who God is".
究竟這位神是誰來的.
我們往往信了耶穌之後.
我們都會take it for granted.
理所當然覺得神是會的.
但沒有清晰了解這位神究竟是一位什麼的神.
這個詩篇的作者講得很清楚.

$^{401}$這位神就是上帝施恩的上帝.
是施恩典的上帝.
而我就是一個得益的人.
基於他的說話.
我們得到上帝的應許.
就好像這裡再列出詩篇121篇.
第一節到第二節.
很多人都很熟悉這篇詩篇.
他說我要向山舉目.
我的幫助從何而來呢.
第二節他說過.
我的幫助是從造天地的耶和華而來.
如果你能夠在這個困難的時候.
你懂得背誦121的第一節到第二節.
你心裡面就有一個力量跑出來.
這個都是上帝的話語帶出來的力量.
弟兄姊妹.
我不是在你的教會裡面.
所以我不清楚你現在遇到什麼困難.
或者什麼處境.
但通過聖經的話.
很奇妙的.
他就能夠給我們.
通過上帝的啟示.
上帝要我們知道.
無論你的困難是什麼也好.
無論你的困難有多大也好.
只有上帝是唯一一個出路.
再說多一次.
無論你的困難是什麼.
無論你的難處是哪種.
都是一樣.
只有一條出路.
或者更加可以說得清楚.
只有唯一一個生路.
就是回到上帝那裡.
因為上帝是完全準備好.
他完全都在我們左右隨時隨地可以幫助我們.
問題就是我們願不願意去找他.
我們願不願意去找他.

$^{441}$如果我們願意去找他的話.
我們就可以解決很多的困難.
其實在詩篇裡面.
詩人其實都是spot guide.
在詩篇二十篇的第七節.
他講過.
你可能會聽過這段詩篇.
他說有人靠車.
有人靠車,有人靠馬.
但是我卻替耶和華.
是什麼意思呢?.
就是說耶和華是真正的倚靠.
如果你要落入一個地步.
已經不知道出路的時候.
你唯一一個倚靠.
就是耶和華自己.
沒有人能夠幫到你.
甚至你最親的人都幫不到你的時候.
唯有上帝在你旁邊可以幫助你.
我們陸續在這個經文下去.
你就會知道是什麼回事.
第二個段落.
二十五十一節和二十五十二節.
當我們小心去看的時候.
這裡提到有一堆人是驕傲的人.
很明顯是敵黨神的那些人.
是驕傲的人.
很狂妄.
傲慢我就是一個狂妄的人.
你有沒有留意作者詩人的回應是什麼呢?.
「我卻未曾偏離你的律法」.
講到他的決心.
各位聽眾.
一個有信仰的人.
其實都不是很艱難.
一個有信仰堅定的人才是最艱難.
無論什麼情況下.
他都堅定相信這位神.
是可以帶領他去解決這個問題.
因此在二十五十二節裡面講到.

$^{481}$祂要紀念你從古以來的典章.
以致我可以得到安慰.
典章是什麼意思呢?.
就是上帝的判決的言語.
判於.
就好像法官要判人.
他講出一段話.
上帝的判於.
上帝是用祂的公義作出判斷.
因此我就不需要怕.
我心裡面就得到安慰.
我們很多時候都會提到.
我們有問題的時候.
我們要彼此代討.
我們有困難的時候.
我們找我們的牧師.
牧師了解你的情況.
為你一切祈禱.
但是最深切的安慰是誰呢?.
就是上帝自己.
當我們了解到這件事.
我們就不需要怕.
往往我們面對困難的時候.
我們要面對的是即時的困擾.
那個困擾是很厲害的.
有時是旁邊的人困擾我們.
有時是整個世界裡面對我們嘲諷.
覺得基督徒是傻的.
向空氣說話.
哪有可能的.
我靠這件事.
我靠那件事.
我靠我的關係.
我靠我的錢財.
我靠我的勢力.
我才可以解決那個問題.
怎麼會靠祈禱呢?.
開玩笑的.
世界的人對於我們來說是嘲諷.
慢慢這些情況就會侵蝕我們對神的信心.

$^{521}$一點一滴的信心就會有問題.
其實解決的方法就只是這樣.
第一件事.
就是找尋從古以來的典章.
找尋從古以來上帝的話語.
原來是有能力的.
作為一個基礎.
而且心中一直堅持所聽到的信息.
這種堅持.
這種信心去持守.
原則就是來自這裡.
紅色的字.
一定要找到一個叫做定海神針.
定海神針就是好像武俠片那樣的情況.
很安定的.
拿到那個準則的.
來回應前面的掙扎.
聖經從來沒有說否定你的掙扎.
你要很誠實去面對.
好像剛才我們說的士伯福先生.
他真的很誠實面對現在.
老婆雖然救回來了.
但是四個女兒都死了.
這是很現實的一個問題.
掙扎是有的.
直接要面對它就不要迴避它.
你可以跟神說.
我現在真的為了這件事.
我真的很傷.
我真的很不明白.
這樣的情況會出現.
弟兄姊妹.
我們有時候面對一種困難是怎樣呢.
當我信了主不久的時候.
你會發覺一件事.
你求神是很成功的.
而且很快就成功.
但是當你信得久了之後.
信了十年之後.
二十年之後.

$^{561}$你會發覺.
神聽你的禱告好像慢了很多.
祂要你學習這個情況.
有時候人家告訴我們.
我記得我初信的時候.
有人給我一本小書.
《如何尋求神的旨意》.
這本小書你一定看過.
那時候半懂半不懂.
看完就算了.
有人告訴我們.
如何尋求神的旨意.
如果是順順利利.
就是神的旨意.
如果不順利.
就不是神的旨意.
那時候我信得很厲害.
覺得那條路就是這樣.
很通達.
那就是神的旨意.
但是要告訴大家.
當你慢慢成熟之後.
當你慢慢在主裡面開始有更加扎根的時候.
上帝會給你一個額外的功課.
在環境看來絕對是順利的.
但是卻不是祂的心意.
卻是所謂無緣無故.
無緣無故就會出現一些很奇怪的事情.
那你怎麼辦呢?.
那你說為什麼要推翻我以前祈禱求神那件事.
我就不再信祂了.
你耍我嗎?你有一種這樣的心態.
你耍我嗎?.
那你就信心越來越軟弱.
弟兄姊妹,這些情況將會非常普遍.
你遲早都會遇到這樣的情況.
因為我們的神會跟著我們的靈命.
來給你有不同的挑戰.
以致你更加成熟.
無論是順境,無論是逆境也好.

$^{601}$原來神都要帶領我們經過水火.
去達到豐富之地.
這個境界真是很高.
你都可以學習到.
我都未算學習到.
但是你總是試過一兩次.
真的摸不著頭腦,為什麼會這樣.
就是覺得很奇怪.
但是幸好都沒有碰釘到一個地步.
就是跨台.
真的很真實.
你不要以為牧師不會跨台.
你們牧師可能不用.
我這些可能會.
牧師都是人,都可以跨台.
如果有一天我們對神的信心是沒有.
或者欠缺的話.
一樣,好像你們一樣,會跨台.
所以在以賽亞書41章第10節.
這段說話真是很精彩.
祂說你們不用怕.
你們不需要驚慌.
因為祂是神.
祂一定會兼顧你.
祂會幫助你.
祂會扶持你.
因為你是祂的子民.
我們有時候說笑.
誰會照顧你.
意思就是說誰能夠捉住你.
保護你.
因為祂有權有勢.
祂可以保住你.
我們這位神就是保住我們那位神.
其實這裡告訴我們一個很大的真理.
在耶利敏書24章第7節下半節.
有兩個說話是非常厲害的.
第一句說話是.
神說他們要做我的子民.
他們這群人,這群以色列人.

$^{641}$是要做我的子民.
第二句說話.
我要做他們的神.
是什麼意思呢.
他們要做我的子民的意思就是說.
這位神是我們的山寨.
這位神是我們的避難所.
這位神就是那位充滿萬佑.
隨時都可以幫助我們那位神.
因此我們要得連率蒙恩惠.
作隨時的幫助.
就是那位.
在希伯來書裡面.
第4章16節下半節.
就是講這個說話.
當祂說我要做祂們的子.
我要做祂們的神.
是什麼意思呢.
就講到神跟人之間的關係在哪裡.
我們的信仰不是在講一些條文.
我們的信仰乃是在講.
我們跟上帝之間的關係在哪裡.
中國人有一句說話很精彩.
其實都有相類似的意思.
有關係就沒有關係.
有沒有聽過.
有關係就沒有關係.
意思就是說.
我跟對方有些來往.
有些關係的話.
什麼事情都沒有所謂.
什麼事情都沒有問題.
這樣的意思.
換句話來說.
如果我們跟神有關係的時候.
我們跟神是天上的父親的時候.
我們就不需要害怕.
沒什麼事情會離開祂的保護.
對不對.
盼望我們能夠得到這樣的保護.

$^{681}$其實我們跟神之間的關係.
是打風也打不掉的.
你記住.
在這樣的狀態下.
我們就可以去看看.
第三個段落究竟是講什麼.
第三個段落是講到53節到56節.
這裡講到有兩類型的情況.
一類型的情況就是神祂自己.
另外一類型的情況就是那些惡人.
當然那些惡人就是敵黨神的人.
跟神對立的.
他說我見到惡人離棄神理的律法.
就怒氣發作猶如火燒.
他真的很生氣.
他發了一些異怒.
他就很生氣.
我們會不會這樣呢.
我們可能會當作看不到.
這是他死他的事情.
發異怒.
基於有神的生命.
我們發異怒.
54節,55節和56節.
講到更加奇妙的情況.
他說因為你的律例作為我的詩歌.
是什麼意思呢.
原來詩歌就好像一個興奮劑一樣.
是一種勵志的作用.
你有沒有經歷過一個階段.
就是你很低落的.
有些詩歌真的可以提升你的情緒.
將你重新再起身.
弟兄姊妹.
我特別想跟弟兄講一下話.
弟兄們.
不要以為流淚是一種軟弱.
你記住.
不要以為流淚是一種軟弱.
流淚是表達了我們聽到一首詩歌.

$^{721}$我們心受感動之後.
很自然我們就眼濕濕.
很自然就受到一個安慰.
很自然就受到一種激勵.
這個很重要的.
所以下一次如果你唱詩歌的時候.
你流淚.
其實這個我可以恭喜你.
你真的進入到跟上帝之間有關係.
其實這個特性是什麼.
遵守他的律例.
因此最後我堅守你的份詞.
堅守你的份詞就是說.
他具體的要求指示.
他一一去做.
一一去做.
但這裡說到一件事.
這些人有一個對比.
那些惡人就是嫉惡如仇的狀態.
但有些人他是堅定.
走在上帝的路.
我沒什麼興趣知道你信主有多少年日.
因為信主年日不代表一切.
這個不是一個最好的指示.
我反而重視的是.
你信主之後.
有沒有一個重生的表現.
以致你能夠很清楚將上帝的生命活出來.
有些人信主不是很多年.
他的生命非常之高.
有些人信了主50年60年70年.
卻變成了什麼.
對不起說一句.
如果是說中你的.
我不是特意針對你的.
就變成老油條.
老油條的信徒太極極之是高手.
亦都是對神沒什麼感覺的.
對不起.
我真的對不起說以下的說話.

$^{761}$神如果要拯救這些人.
可能會用一些特效藥.
用一些開劑厲害的東西給你癱瘓.
往往是你的健康.
你信不信呢.
你周圍看看旁邊那個.
不是指你現在坐的旁邊那個.
我的意思就是說.
你真的看一下周圍.
有些信徒已經沉睡了.
信你仍然愛他用什麼方法.
怎樣提醒他.
就是大的苦難.
我們當然不會希望有苦難.
只是拿來搞.
但是上帝的愛在背後就出現了.
一會兒我們還會說這個題目.
在這個情況還要提到一句很精彩.
54節那裡還有很精彩.
我特別留意到.
我在世記據.
這是一個關鍵的觀念.
就是說我們一班信耶穌的人.
其實怎樣看我們生活在這個世代.
在這個世界裡面的時刻.
有人長一些有人短一些.
這個完全都沒有問題的.
有人長一些有人短一些.
完全都沒有問題的.
我們要明白到我們在世的時候.
只不過是暫時而已.
弟兄姊妹.
什麼叫做暫時.
就好像你去旅行一樣.
你背著行李.
暫時而已.
你去一間酒店覺得不太好.
你暫時過一兩晚就走人了.
OK.
只不過是暫時而已.

$^{801}$但是你買一間房子就不同了.
你可能是永遠的.
暫時和長久就有這樣的分別.
記居的意思就是說.
我們渴望的是永遠的家鄉.
弟兄姊妹.
當我們信耶穌信到一個地步.
我們很認同這個說法的時候.
我就要恭喜你了.
你真的拿到屬靈長進的若時的關鍵.
心裡面有永恆.
以致我可以如常去遵守上帝的命令.
無論環境是怎樣都好.
再一次肯定上帝的真實.
如果我們還沒看到上帝的真實.
那時候求祂讓我們看到祂的真實.
以致我們可以心底裡面.
有一個最大的一個改變.
其實無論你的環境是多麼艱難都好.
如果你跟著這個方向去做.
我相信你會靈明.
一定會長進.
好像保羅在哥林多後書.
這段聖經是很厲害的.
尤其是十七至十八次.
真的是很厲害.
十八次他就說.
保羅看到了.
原來當我們的視野被擴闊了之後.
擴闊了之後.
我們的動力就不一樣了.
他看到有兩堆人.
一堆人是平常的人.
平常的人他就看到什麼呢.
看到一些肉眼可以看得見的東西.
看到一些只不過是暫時的東西.
這些是平常人的一半的情況.
但是他又說到.
原來有另外一班人是不同的.
是屬神的人.

$^{841}$他的境界被提升了.
他是看到一些你看不到的東西.
他是看到永遠的東西.
而追求看不見的.
和追求永遠的.
在這樣的情況下.
你的靈明就已經是升了一個層次.
或者是兩個層次.
好不好啊.
是不是.
這些是好東西.
鼓勵你真的嚮往這件事.
在詩篇一百一十九篇.
我們會再重複一下.
要留意的是什麼.
左邊有說到正面的提醒.
右邊說到沒有做一些負面的東西.
兩部分其實都很重要的.
正面的提醒就是說.
這些困境都是很真的.
完全都是真真確確來到我們的身上.
所以最重要是回到神那裡.
來得力量去處理它.
要處理它才行.
它不會自己走了.
要處理它.
就好像哈巴局書裡面.
第三章十七至十九節.
那種最後的讚美.
就是一個最典型的一個鼓勵.
它提到有一個結構的.
就是雖然……然而…….
如果你讀過哈巴局書第三章.
這一段話的時候.
你心裡面有很大的安慰.
雖然無花果樹不發黃.
葡萄樹不結果.
橄欖樹也不效力 田地不出糧食.
圈中絕了羊 棚內也沒有牛.
這個就是十七節裡面說.

$^{881}$這些都是在外在的環境裡.
是很真實的困難.
在農業的社會來說.
是很慘的現狀.
但是哈巴局書的第三章十八至十九節.
說到另外一個境界.
他就說 然而 我要因耶和華而歡呼.
因救我的神而喜樂.
至於耶和華是我的力量.
它使得我的腳快如無碌的蹄.
它又使得我穩行在高處.
當你很低落的時候.
當你頭頭碰著黑的時候.
請你讀一讀哈巴局書第三章.
十七至十九節.
你就知道你如何.
應該的態度是怎樣.
雖然是這麼差.
雖然外在的環境是這麼差.
但是然而我有一天有神.
我就有完全不一樣的態度.
我們看回右手邊.
右手邊有一些很好的提醒.
叫我們不要做一些負面的事.
負面的事是什麼呢.
不要埋怨當中.
將你的身完全埋葬在一個埋怨當中.
怎樣埋怨呢.
常常都會問.
為什麼.
why me.
哇 現代人很喜歡這個詞.
why me.
你知不知道why me這句話.
隱含了有些什麼意思在裡面呢.
就是對神有一個很大的不信任.
這個是罪來的.
這個罪不單止是道德問題的罪.
對神不信任也是一種罪.
不信任就是不認為萬事神都知道.

$^{921}$萬事神都容許.
到最後卻是帶領我們去到有益處的狀態.
就是說你今天遇到這種困難.
我不知道你是什麼類型的困難.
具體的我不知道.
這個困難其實都是神容許它發生.
你以為神不知道嗎.
當然是知道的.
問題就是祂要熬煉你.
磨煉你這個艱難的時刻.
怎樣去依靠祂.
以至到你可以走出這個困境.
當你走出這個困境的時候.
是因為沒有埋怨這個困境.
而達到一個豐盛的生命境界.
這個就是那個目標.
因此整體來說.
詩篇119篇的這一段話.
是要我們幫助我們建立.
以一個神為主導的人生的體會.
總結有兩件事要做.
第一件事就是回歸到神那裡.
你回歸到神那裡.
將神視為你是唯一一個盼望.
第二回歸祂的話.
就是聖經.
所以其實不是很難.
具體來說有以下三點請你特別注意.
以下三點告訴我們知道.
第一 聖經是有上帝的應許.
就是給人有指引.
祂指引你.
祂未必會很清楚告訴你.
你應該選擇A這份工作.
或者B這份工作.
祂沒有這樣說.
但是祂給你一些指引.
這裡有一句話就是Richard Shipps.
就是十六世紀,十七世紀的時候.
英國清教徒路線一個主流的人物之一.

$^{961}$他說了一句話.
他說當我們在神的話語當中.
聽到有任何應許的時候.
就將祂變為一個禱告.
這個就是你的回應.
這個就是你的行動.
因為神的應許就是祂的承諾.
祂很願意我們跟祂角力.
是什麼意思呢?.
我舉一個例子你就明白.
譬如你有子女.
或者你甚至有孫.
那你曾經跟祂說過的話.
你說如果他考試90分以上就有獎.
你曾經說過的.
那你的孫或者你的子女也好.
他就聽了你的話.
他就偷偷地去祈禱.
向神祈禱.
神啊.
麻煩你叫我爸爸或者我爺爺.
就是提醒他.
記得提醒他.
買一份禮物給我.
因為我90分以上的.
他這樣祈禱.
你千萬不要忘記.
因為祂這個祈禱.
帶給一個力量就是這樣.
我告訴你肯定的後果.
你肯定的後果就是.
第一.
你甘心樂意去買那份禮物.
第二.
多麼貴也好.
都去買.
為什麼呢?.
這個是你的承諾.
你不可以在祂的面前失信於祂.
對不對?.

$^{1001}$如果你有子女.
或者你有孫的.
你多多少少都遇過這樣的經歷.
自己銀河包.
不捨得吃.
不捨得去買某某東西.
但是買給孫的那些.
怎樣也要.
我真的心裡面太多這樣的例子.
因為我太太經常都是.
非常慷慨買東西給我的孫.
我有時候說有沒有搞錯.
因為你不是我的孫.
所以你沒有份.
就不簡單.
你承諾了就必定要做.
拜託你.
我們的神承諾了一些東西.
他的應許.
聖經說的應許.
他必定會實行.
第二點.
聖經是匹於我們有堅定的信仰人生.
你有堅定的信仰.
你必須要聖賢為基礎.
這裡列了一句話.
以賽亞書第四十五章十七字.
他說唯有以色列必蒙耶和華的拯救.
得永遠的救恩.
你們必不蒙羞.
也不抱愧.
直到永世無盡.
兩次出現了一個字眼.
必.
你看聖經的時候不要這麼快.
這麼快一跳就看不到.
必的字眼.
必定的意思.
很肯定的說話.
一節的聖經提到有.

$^{1041}$兩次永遠的觀念.
我們的神要做事.
不是隨隨便便.
馬馬虎虎.
說過不算數的那些.
他說得出就做得到.
說得出就做得到.
所以是相當之有份量.
看到永遠的東西.
我有以下的一個感受.
我不知道你有沒有以下的同感.
現在這個世代.
這個時代.
有哪些東西是永遠的呢.
商品會變成壞的.
關係會變成冷淡.
承諾會失去效力.
多得很.
但唯有上帝給人的救恩.
是永遠不會失落.
你大可以放心.
因為我們的神對於救恩的給予我們.
是很堅定的.
他的承諾是非常堅定的.
他甚至連他自己的獨生子.
犧牲在十字架上都做得出.
因此他對救恩這個事情.
實在是很認真的.
不是隨隨便便做的.
所以我們大可以放心.
堅持上帝的救恩.
在這節經文裡面帶出有兩個訊息.
特別是第二點對我們有關.
第一點就是說.
這段經文說給以色列民知道.
上帝有主權和永遠不會落空的應許.
叫他們得拯救.
和離開被擄之災.
但對於我們來說.
就是引申到所有信徒們.

$^{1081}$上帝是應許救贖和永恆的生命.
給哪些人呢.
信靠他的人.
就是你和我.
如果你和我都學習信靠他的話.
你就同樣得到上帝的應許.
就同樣有很大的得福.
那你可能會問.
你怎樣會知道是這樣的情況.
或者你是怎樣那麼肯定.
為什麼不肯定呢.
答案很簡單.
因為聖經清清楚楚寫了出來.
我們俗語有句話.
口講到無憑.
清清楚楚寫出來是白字黑字寫出來的.
那沒得爭辯了.
所以你可以和神爭辯.
好像不對辦.
這裡的經文是這麼說的.
我覺得你一定能爭辯的.
一定能贏的.
因為我們的上帝是牙齒當今世.
不會說了不算數的.
所以你聖經一定要熟.
一定要堅定去支持.
你可能會覺得剛才的第一點和第二點比較抽象.
第三點就更加實際一點.
是很多實際的應用在裡面.
其實是歷世歷代的信徒.
已經是有很多的見證.
告訴我們知道.
聖經能夠幫助軟弱的我.
可以重新再起.
所以你一定要承認一件事.
無論你的地位是多麼慷恭.
無論你聖經多麼熟.
無論你是怎樣的情況下.
在屬靈的事情上.
你和我都相當之軟弱.

$^{1121}$但是軟弱不怕.
因為軟弱得來的時候.
我們可以有機會重新再起身.
為什麼?.
不是因為我們有什麼特異功能.
而是因為上帝是一個.
非常樂意幫助我們的一位神.
祂願意給我們重新再起的機會.
因為祂給我們很大的憐憫.
有三個細點.
第一點就是.
聖經內中有生命的應許.
這個真的很重要.
你一定要抓住.
這個就是我們上個星期說的.
神聖的生命那段話.
約翰一書的五章.
十一至十三節.
這裡說到.
這個見證就是.
神賜給我們這群人有永生.
這個永生就在.
耶穌基督的兒子裡面.
然後他說.
人走了.
上帝的兒子就有生命.
沒有神的兒子就沒有生命.
很簡單的道理.
有耶穌基督的生命.
有耶穌的兒子就有生命.
沒有耶穌的兒子就沒有這種的生命.
沒有這種的福分.
我將這些說.
社給你們信奉神的兒子的人.
要叫你們知道.
自己.
有永生.
已經有了永生.
你信了耶穌之後.
已經有永生.

$^{1161}$是從天而降.
進到你的內心裡面.
無法抵賴的.
無法抵賴.
你已經有純粹.
你已經有這樣的強烈的信任.
已經存在.
你根本跑不掉.
在這樣的狀態下.
聖經就可以將這種神聖的生命.
可以更加固化了它.
更加落實了它.
那你就很棒了.
就很棒了.
是不是這樣.
一定是這樣的.
這個道理就是這樣.
第二點.
神的話在客觀上帶來它的道.
其實帶了來信道的後果.
在羅馬書十章十七次.
講到有一段說話.
原來信道是怎樣來的呢.
就是從聽道來的.
聽道是怎樣來的呢.
從基督的話來的.
就是從聖經的話來的.
聖經的話就令到有人聽了道之後.
就會接受這個真道.
所以你聽是很重要的.
今天如果你是信心軟弱的.
如果某某弟兄或者某某姊妹.
信心軟弱的.
最重要最重要鼓勵他.
能夠再繼續來聽道.
因為上帝會使用他的道.
上帝會使用他的聖經的話語.
藉著傳道人也好.
或者藉著普通的弟兄姊妹也好.
能夠甦醒人心.

$^{1201}$那個能力不是那個人有什麼厲害.
那個能力就在上帝的話那裡有厲害.
以致他可以相信.
第三.
他的話有不變的承諾.
不變的承諾.
重現對上帝的信靠.
提摩大後書二章十三次.
這段非常好的信號.
希望你背了它.
真的太好了.
他說我們縱然失信.
我們縱然間失信也好.
他就是神.
他仍是可信的.
因為這位神他不能背負自己.
說明了他不可以背負他自己.
Guarantee.
所以我們縱然間很壞也好.
但是神仍然都是可信.
仍然都是可信.
弟兄姊妹.
如果是這樣的話.
你還有什麼可以害怕的.
你仍然都可以重新站起來.
絕對是可以重新站起來.
我跳過這一章.
我們總的領會就是.
上帝原來是堅定不移的天父.
即是說我們已經偏離了他.
或者我們已經忘記了他.
我們仍然都可以尋找他.
我們可以尋找他.
其實真正尋找不是我們去尋找神.
原來是耶穌去尋找我們.
詩篇23篇有說.
你的脹你的肝都安慰我.
耶穌很想尋找我們.
我們軟弱我們遺棄了他.
都尋找我們.

$^{1241}$因為他還有一樣很重要.
他會保守我們作為聖徒.
能夠堅忍到底.
不是堅忍一次兩次.
是堅忍到底.
到最後那一天.
為何他這樣說呢.
原來我們隨時都可以倒下.
隨時都可以倒下.
我們縱然軟弱也懷疑他.
但上帝仍然是那種力量.
可以幫助我們.
你知不知道跑馬拉松要跑多遠.
馬拉松.
有人跑過馬拉松嗎.
全馬42公里.
42公里.
大概英里是26公里.
聽說2026年的世界紀錄.
跑全馬男子的紀錄.
有一個肯亞的人.
兩個小時內已經跑了42公里.
真的很厲害.
你知不知道跑馬拉松.
我沒機會可以跑全馬.
半馬都不行.
可能一馬差不多.
你知不知道跑馬拉松.
最難是哪個時段.
最難不是最初要怎樣跑出來.
帶領那些人.
最難也不是最後怎樣衝刺.
而最艱難是甚麼.
是中間的時候.
即是跑到大概一個小時後.
即是跑到大概一半.
你怎樣堅持下去.
堅持下去.
我最近做跑步這件事.
要鍛鍊好身體.

$^{1281}$最初的時候不敢跑那麼遠.
跑到最後氣來氣喘的時候.
怎樣堅持下去.
那個真的很艱難.
祂知道我們的軟弱.
祂明白我們艱難.
祂會為我們打氣.
加油.
所以這裡的茅有甚麼提醒呢.
我們有一個指望.
如同靈魂的茅.
祂說那艘船要安定不要移來移去.
它拋了那個茅之後.
那艘船就穩定下來.
而且進入萬內.
進入到上帝好的關係裡面.
我們希望.
我們真的可以.
回歸到聖賢的路.
就是以下一個重點.
以下一個重點.
還記不記得我上一次和大家說過.
怎樣用顯微鏡.
怎樣用望遠鏡.
怎樣用一個停診器.
醫生的青筒.
不單是這樣.
還看到聖經裡面.
有一個真理的架構.
一個框架.
更加闊的框架.
究竟是甚麼.
上帝更大的應許.
我們有一個永遠的救恩.
你和我實在太重要了.
在神的眼中太重要.
是一個永遠的救恩.
而不單止在那裡停頓.
更加好像在醫生裡面聽筒.
聽我們心裡面有甚麼困難.

$^{1321}$進入到內在裡面的艱難的地方.
從日常的生活裡面.
開始建立到.
上帝的恩典裡面.
你欣賞上帝的恩典.
保羅有一句說話.
很珍貴的.
他說在臨前十五章十節.
他說我成為何等的人.
都是蒙神的恩才成的.
我們真的要去到.
這樣的境界.
就會有很好的幫助.
所以在這個時刻.
請大家都閉上眼睛.
我和大家說幾句心裡的話.
讓我們都很認真去面對.
聖經在我們心目中的地位.
是甚麼.
或者你覺得.
很可能它是一本可有可無的書.
還是它是一個寶藏.
讓我們發掘到.
上頭而來的真智慧.
又或者你已經信主.
很多年了.
卻沒有真正明白.
聖經的可貴.
我們去讀聖經.
是不是只是循例去作的呢.
又或者你已經是一個.
很認真的信徒.
很仔細去研讀.
聖經的每一句.
每一個字.
你從來沒有注意到.
聖經指的是一個媒介.
幫助我們追求生命的主.
並且和神.
要發展不斷的.

$^{1361}$建立一個關係.
這個目標.
你是不是都有這個問題呢.
請你記住.
我們只是要重視聖經.
不是崇拜聖經.
真正值得追求的.
是聖經指示的.
主耶穌基督.
這個時候.
可不可以你也和上帝.
做一些表示呢.
第一.
今天我們講到的訊息.
你開始明白到.
聖經真正的用途.
乃是和.
生命的主連接在一起.
回歸聖賢的追求.
我才明白到.
人生很多困難的時候.
真正的出路在哪裡.
對於神給你的恩典.
讓你重新認識聖經.
有沒有人願意.
向神表達.
如果有的話.
請你輕輕舉一舉你的手.
讓神知道.
多謝多謝.
讓神知道你是真正的.
明白聖經原來是這麼寶貴.
我再講有沒有人願意.
如果有請你舉手.
多謝多謝.
第二點.
既然聖經對我們屬靈的生命.
是這麼重要.
從今天開始.
我要立志好好地.

$^{1401}$看聖經.
追求上帝的道.
使我們的信仰根基.
可以牢牢的紮根.
有沒有弟兄姊妹.
願意這樣的立志.
現在這一刻開始.
好好地讀聖經.
追求這位生命的主宰.
如果有人可以舉一舉手.
多謝多謝.
不要放棄這個立志的機會.
因為是一個很嚴肅的時刻.
有沒有這樣的人.
再等多一兩秒我們就完了.
這個時候.
我就為當前.
這兩樣的事情.
舉手的弟兄姊妹.
求主見定.
這些弟兄姊妹的信心.
直覺永活的聖賢.
得以回歸到你那裡.
你的兒女都能夠.
成為堅強的信徒.
帶領他們.
一步一步走向.
榮耀的永恆.
我們的祈禱乃是奉耶穌基督的命求.
阿們.
願主真的祝福大家.
真的在聖經裡面的追求.
乃是一生之救.
我們的經歷.
往往未必是.
完全都是很順暢.
但是原來背後.
是得到上帝很大的祝福.
願上帝祝福大家.
\newpage



\section{啟示錄 6:1-8}
\label{sec:032438}
\textbf{審判的七印——逼近的災難}~~~~譚綺敏傳道
\newline
\newline
連結: \href{http://www.church.com.hk/acms/content.asp?site=cdc\&op=show\&type=product\&code=032438\&layout=sermon}{\texttt{032438}} ~~~~ 語音日期: 2026-07-12
\newline
\newline
\hyperref[sec:032437]{\small{< < < PREV SERMON < < <}}
~
\hyperref[sec:index_chronic]{\small{[返順時目]}}
~
\hyperref[sec:032439]{\small{> > > NEXT SERMON > > >}}
\newline
\newline
啟示錄 6:1-8
\newline
\begin{longtable}{cl}
\hline
\hline
章節 & 經文 (和合本修訂版)\\
\hline
6:1 & \begin{tabularx}{0.7\textwidth}{X} 我看見羔羊揭開七個印中第一個印的時候，聽見四活物中的一個活物，聲音如雷，說：「你來！」 \end{tabularx} \\ \\ \relax
6:2 & \begin{tabularx}{0.7\textwidth}{X} 我就觀看，看見一匹白馬，騎在馬上的拿著弓，並有冠冕賜給他。他出來征服，勝而又勝。 \end{tabularx} \\ \\ \relax
6:3 & \begin{tabularx}{0.7\textwidth}{X} 羔羊揭開第二個印的時候，我聽見第二個活物說：「你來！」 \end{tabularx} \\ \\ \relax
6:4 & \begin{tabularx}{0.7\textwidth}{X} 就另有一匹馬出來，是紅色的；有權柄賜給了那騎馬的，要從地上奪去太平，使人彼此相殺；他又接受了一把大刀。 \end{tabularx} \\ \\ \relax
6:5 & \begin{tabularx}{0.7\textwidth}{X} 羔羊揭開第三個印的時候，我聽見第三個活物說：「你來！」我就觀看，看見一匹黑馬；騎在馬上的，手裡拿著天平。 \end{tabularx} \\ \\ \relax
6:6 & \begin{tabularx}{0.7\textwidth}{X} 我聽見在四個活物中似乎有聲音說：「一個銀幣買一升麥子，一個銀幣買三升大麥；油和酒不可糟蹋。」 \end{tabularx} \\ \\ \relax
6:7 & \begin{tabularx}{0.7\textwidth}{X} 羔羊揭開第四個印的時候，我聽見第四個活物說：「你來！」 \end{tabularx} \\ \\ \relax
6:8 & \begin{tabularx}{0.7\textwidth}{X} 我就觀看，看見一匹灰色馬；騎在馬上的，名字叫作「死」，陰間也隨著他；有權柄賜給他們，可以用刀劍、饑荒、瘟疫、野獸，殺害地上四分之一的人。 \end{tabularx} \\ \\
[1ex]
\hline
\hline
\end{longtable}
$^{1}$今日的講題就是「審判的七任逼近的災害」.
我首先講一講我自己的一個故事.
有一年我去宣導總會的退休營.
當我開車回家的時候.
我聽到「啪啪」的聲音.
我不知道是從哪裡發出來.
是很持續和有規律的.
我走著走著 繼續聽到.
我的反應是「好像沒什麼事」.
「車開著都挺好 沒什麼事的」.
於是就不理會它了.
想著還有一個多小時才回家.
到時候回家就看看它.
怎知道還不到一個小時.
我就聞到車上有少許煙味.
又不是車頭冒煙.
當時我已經走上了起碼80公里的高速公路.
不是高速 是大的公路.
不能停下來 因為沒有小街讓我停.
唯有不要開得那麼快.
一直祈禱「神啊 你讓我平安」.
不知道發生什麼事.
大家知不知道發生什麼事?.
離開公路我就駛向一個小的街道.
我下車看看什麼事.
發覺車輪完全平了.
怪不得有點不舒服.
因為用輪胎在刮路.
一邊走 我都看見有些黑色飛鳥飛過來.
我不知道從哪裡飛過來.
我就只好去找NRMA 即是路邊服務.
我不懂得換輪胎.
後來收車的師傅跟我說.
「哇 你的輪胎爆成這樣 連鋼鐵都外露」.
「它的駕駛就很不穩定 而且很危險」.
我還開到80,90公里的高速公路.
所以問題不是車輪突然爆胎.
問題是我忽略了很早出現的記號.
就是車輪胎在路上被釘插進去.
然後就發出「砰砰」的聲音.

$^{41}$一直開車都有「砰砰」「砰砰」的聲音.
很小的聲音 但我聽到.
我當時不知道發生什麼事.
我又沒有預早停下來看什麼事.
所以將自己放在一個危險的環境裡面.
所以以後大家留意一下.
同樣的 我們今天活在一個危機四伏的時代.
當我們打開新聞 聽一些新聞.
你發覺不論是國際性的地緣政治.
或者是地區上一些的戰爭,煙火.
或者全球性的 現在都有通貨膨脹.
我們都很難以預測 也都有病毒.
或者是氣候變化帶來一些災難.
就是地震,海嘯,熱浪,熱死人都不少.
歐洲也都熱死人的.
所以其實這些東西我們都很難去預測.
你看見世界的人的道德觀念都要越來越下降.
人與人之間也都很緊張.
就算工作上,家庭裡面都有一些不平安.
所以我們經常都好像有些危機感.
或者有些人都會問 神在哪裡啊?.
祂為什麼不管?.
這麼多災難 神為什麼會讓這些災難去發生?.
其實啟示錄告訴我們一個我們很意想不到的答案.
就是神不單止在管.
神是向人已經發出了警號.
所以今天我們講到 剛才讀的經文.
講到揭開的印 其實是告訴我們一些災難正在逼近我們.
其實這些就是神要發出的警號.
讓我們看見世界是會有這樣的走向.
我們要怎樣去回應神呢?.
我們一起先來祈禱 好嗎?.
神愛天父 我們看見世界上的混亂.
我們看見你來的日子很近.
你要審判這個大地.
當我們看見這些混亂 不知道怎麼辦的時候.
如何求你去教導我們.
怎樣去數算我們的日子.
好叫我們得著智慧的心.
讓我們謙卑下來 願意去回到你的根前.

$^{81}$去聽從你的話語.
以致我們心裡面有平安.
因為神你說過 在你裡面有平安.
雖然世界上是有苦難.
願主你帶領我們.
多謝你聽我們祈禱.
奉主耶穌基督得勝的命祈求 阿們.
上星期Nelson傳道講到約翰福音.
約翰在《啟示錄》裡面被提到天上.
在《啟示錄》第四,第五章.
他看見神坐在寶座上.
那時候神拿著一張書卷.
在第五章的一節開始.
講到這張書卷是用七個印去填封的.
沒有人是配得去展開這張書卷.
書卷用七個印去填封.
是很實在的 很重要的.
當時猶太人是有些印去填封.
好像圍竹那些.
不准死之前都不可以去掀開的.
是很重要的書卷.
所以約翰看見神手裡面拿著這張書卷的時候.
他就大哭了.
為什麼他要哭呢?.
因為約翰知道如果神的計劃不可以實現的時候.
他知道這不是普通的書卷.
是一個好像別人納完圍竹之後很重要的書卷.
所以他知道是關於神要向普天下.
講的一些很重要的信息.
可能是關於一些審判的計劃.
更加是關於一些救贖的計劃.
因為神是很愛我們.
所以他就很緊張.
他很怕這些書卷不能夠打開.
你想想在第一世紀教會的處境.
其實他們正是受到逼迫的.
當時是羅馬王去管治他們.
很多基督徒他們都會被殺的.
如果你不敬拜羅馬王的時候.
所以有股惡勢力好像不斷去得勝.

$^{121}$不斷去擴大.
所以約翰也是受逼迫去流放一個荒島.
拔摩海島.
所以他一定很期待.
如果神的國度來臨.
公義就會彰顯.
苦難就可以結束.
所以如果沒有人去打開這個書卷.
代表神的審判或者神的救贖不可以開始.
所以惡人可能一直在當道.
教會也一直受苦.
公義也得不到彰顯.
所以約翰就去痛哭.
很大很大的哭泣.
很大很大的痛哭.
啟示錄五章五節有一個好消息.
說到耶穌基督已經復活了.
得勝了死亡.
所以耶穌基督可以去掀開這個書卷.
掀開這個嚴封的書卷.
即是說耶穌基督開始去執行神的計劃.
所以今天我們集中去看六章一到八節.
前面那四個印.
就說到有四匹馬.
剛才我們讀到.
是什麼意思呢?.
我們再看六章一到二節.
說到高揚揭開七印中.
第一印的時候.
有一匹白馬.
騎在馬上的拿著弓.
並有冠冕賜給他.
他出來勝了又要勝.
有些人認為這個白馬是贏的.
會不會是主耶穌基督來呢?.
主耶穌基督這匹白馬是代表.
好像啟示錄一直看下去.
十九章也說到耶穌基督騎著白馬來.
但其實這個解釋就很牽強.
因為你看見這個白馬是拿著弓的.

$^{161}$代表他會攻擊.
主耶穌基督從來都不會用弓箭去攻擊人.
所以從上下文去看.
這匹馬是和災難有關的.
他不是主耶穌基督.
那麼是誰呢?.
這個白馬要模仿基督.
所以是代表一個假基督的勢力.
他拿著弓代表他要去征戰.
而且是攻擊遠方.
不是近距離.
攻擊遠的東西.
是一個軍事的勢力.
可能是攻打遠方的地圖.
他要去有一個權力的鬥爭.
他要去擴張自己的勢力.
而當時第一世紀的人.
如果看到這個情景的時候.
他們會想到.
羅馬人最崇尚的就是騎著白馬.
所以你會發覺白馬是很好的.
我們說白馬王子.
所以當時很多人崇尚白馬.
因為白馬代表軍隊.
海船歸來都是騎著白馬.
代表勝利.
而且有冠冕賜給他.
所以想到是勝利的場面.
而羅馬當時管治的人民來說.
他一直宣稱.
我們羅馬多麼強大.
我們可以帶來羅馬的和平.
給他的子民.
所以當時他不斷擴充他的勢力.
去打別人的國家.
所以猶太人也是被他們打敗了.
以至他們是在羅馬統治的一個民族.
這個騎著馬的拿著弓.
但是很奇怪地他又沒有箭.
你看他拉著弓他都沒有箭.

$^{201}$代表他是虛假.
不是真的靠軍事實力去征服.
征服這個世界.
征服一些東西.
所以他是靠政治的謊言.
或者經濟操控.
或者一些宣傳.
或者一些偽裝的宣傳和平.
去叫別人去臣服他.
要他去迷惑別人的心.
結果帶來更多的災難.
所以第一匹馬告訴我們.
是一個虛假的和平和欺騙.
你看看今天的世界都是這樣.
我們覺得很和平.
但是如果你突然間接個電話.
然後又信他講什麼.
又將你的錢傳給他.
其實是什麼意思?.
就是被人騙了.
但是有很多人上吊.
澳洲好像是其中一個最容易被人騙的地方.
澳洲是被人騙了每人平均有一至兩元.
如果你沒被騙過其實你很厲害.
你想想二千萬人口.
每人騙了一至兩元就是四千萬.
就是不少人騙了很多錢給別人.
所以另外今天的世界就是這麼虛假.
我們覺得以為很和平.
但是突然間有人用一些虛假的東西去欺騙.
去迷惑你.
你發覺大國與大國之間都是.
弄個和平協議出來.
我們不再打仗了.
有沒有成功呢?.
幾天可能都不知道有沒有.
接著人道主義的口號.
就是我們要和平的.
我們要有人道的.
要講人權的.

$^{241}$很多這樣的口號.
但是其實背後都是騙人的.
大有大的騙.
小有小的騙.
小的騙就像你剛才說的.
電話理財服務其實騙了你的錢.
街頭祈福黨去詐騙又騙你的錢.
所以我們其實是很希望和平的.
我們都以為條約可以去製造和平.
我們以為宣傳和平可以製造和平.
但是你發覺人心比萬物都詭詐.
所以結果是帶來更多的衝突.
一個衝突可以毀掉一個國家.
一個仇恨可以毀滅人的一生.
或者一個謊言可以毀掉一個家庭.
所以我們發覺不容易要和平.
我想唯有我們回到天家才有真正的和平.
所以第一個白馬被我們看到的.
就是虛假的和平和欺騙.
會淋到這個世界.
真是越來越多.
越有知識就越知道怎樣去騙人.
這個真是今天我們面對的世界.
我們再來看第二個印第三第四節.
揭開第二印的時候.
有一匹馬又出來了 是紅的.
有權兵給了騎馬的.
可以從地上奪去太平.
使人彼此相殺.
又有一把大刀賜給他.
紅色我們都想到了.
代表暴力或者流血的顏色.
所以當太平被奪去.
可以奪去太平.
戰爭就沒有了太平.
所以當太平被奪去的時候.
表面上羅馬帝國被人覺得很有太平.
但其實國內的軍隊都是軍閥割據.
不是很服對方.
因為這是一個虛假的征服能力.

$^{281}$所以第一世紀的信徒.
他看見紅馬有把大刀的時候.
他知道.
說到羅馬帝國.
其實他們在公元69年.
剛剛寫啟示錄的前幾十年.
就有一個很悲慘的內戰.
一年內換了四個皇帝.
軍閥割據,血流成河.
其實他們都在經歷這些征戰.
所以紅馬是一個戰爭和人性殘殺的記號.
一個警號.
所以當虛假被揭開的時候.
接下來就是征戰.
人的野心和慾望沒有停止過.
所以不斷都有征戰.
經文講到使人彼此去相殺.
我覺得很恐怖.
人可以彼此相殺.
我們看見今天的征戰,今天的戰爭.
突然間一天可以扔個飛彈下來.
你就沒有了和平.
接著戰爭就開始.
昔日已經差不多四年了.
俄烏戰爭就是這樣.
接著現在的美伊衝突和中東的戰爭也是.
接著還有很多的內戰.
蘇丹,緬甸等內戰都沒有停止過.
所以我看見人最大的問題不是我們的資源不夠.
資源太少.
或者我們住的房子太小.
不是這些問題.
是人的心永遠都不能滿足.
你看見做到首富是不是可以滿足呢?.
你的投資幾代都吃不完.
你是不是可以滿足呢?.
你會發覺不是.
而是人的心不斷去不滿足.
夢想要多一點.
正如在《雅各書》也講到.

$^{321}$它講到你們中間的爭鬥.
按錯了 回頭看.
你們中間的爭鬥是從哪裡來的呢?.
不是從你們百體之中的私慾來的嗎?.
所以其實是我們的戰爭,我們的爭鬥.
不是外面的環境.
是我們的心.
所以我們的戰場其實是我們心裡面.
你想想家庭裡面的戰場.
有多少是在我們自己的心裡面.
就是夫妻之間的衝突.
所以如果人心裡面只有自己.
我這樣的 為什麼你不遷就我呢?.
我要怎樣你就應該遷就我.
如果是自我看得這麼大.
而心裡面沒有了神,沒有了愛.
也沒有了道德觀念.
其實罪惡就會不斷去產生.
而導致更加嚴重的衝突和爭戰.
前兩天我看到一個香港的新聞.
我覺得挺唏噓的.
那個標題大家可能都聽過.
「小巴站兩女打鬥 釀成命案」.
小巴站而已.
因為這個標題這麼說.
我還以為是女生 不是男的.
通常這些我都覺得男的會比較容易.
所以我就進去看看詳細的情形是怎樣.
他說有兩名五十多歲的女子.
也不是二十多歲 血氣方剛那些.
五十多歲.
晚上十點多乘坐小巴到一個屋村的時候.
下車的時候 因為一些肢體碰撞.
有人踩到有人的腳.
然後就引起衝突 踩到腳而已.
被踩腳的女人就將這個婦人倒在地上.
然後用雙腳跪在肚子裡面.
用兩隻手插著她的頸.
襲擊了幾分鐘之後就走了.
這個警方就接報到場.

$^{361}$發現死者昏迷 頸部臉部有明顯的傷痕.
然後就去送醫院搶救.
之後就證實不治.
你會發覺 嘩 為什麼這麼緊要 什麼事啊?.
所以這個新聞令我覺得人性很可怕.
小小的肢體碰撞 叫兩句 突然間就可以大打出手.
家庭也是這樣的 夫妻之間有愛也是這樣的.
所以你會發覺 可以是這麼大的報復行動.
這個可能和我們心理上要征服人的心態有關.
就是你踩到我 我要還顏色給你看看.
讓人知道你不是弱者.
然後就要控制一下 我就要贏.
看你知不知道我的厲害.
其實這個就是熊馬的縮影.
我們為什麼要戰爭呢?.
我們為什麼要爭鬥呢?.
就是因為我們的心不可以滿足.
我們的心要贏.
其實人性的暴戾也挺可怕的.
我們再來看第五節 第六節 六章的第五節.
第三章揭開的時候 黑馬就出現了.
黑馬手裡拿著天秤.
然後有一個聲音說.
「一錢銀紙買一升墨紙 遊和走不可糟蹋」.
一錢墨 先說黑馬.
黑馬其實是代表饑荒.
就是沒得吃 到很瘦.
所以你看到人都很瘦.
貧窮和經濟崩潰的現象.
所以天秤就是象徵著食物都要斤斤計較過.
才可以配給你.
就是天秤的意思 象徵的意思.
所以戰亂帶來的就是饑荒或者經濟的危機.
有可能就是我們接下來世界的走向.
所以第三匹馬告訴我們.
就是饑荒和經濟的崩潰.
當時一錢銀紙是什麼意思呢?.
一錢銀紙就是一個工人一千的工資.
平時一錢銀紙可以買八到十二公升的墨紙.
但是這裡說到現在只可以買一升的墨紙.

$^{401}$原來一升的墨紙當然夠一個人一天的食用.
就是說你辛辛苦苦一天的工作.
只是可以養到你自己.
一家人都養不到 都不夠吃.
所以這個代表極度嚴重的通貨膨脹.
但諷刺是「油和酒不可糟蹋」.
就是說有些人可以享受油和酒.
還說你不要糟蹋這些東西.
就是他在糟蹋 你叫他不要糟蹋.
就是這個意思.
所以就是說有人可以很富有.
有人可以很貧窮.
有人可以沒有飯吃.
有人喝酒有油都可以糟蹋得很嚴重.
所以基本上說到糧食也都被壟斷了.
窮人就餓死 富人就繼續享樂.
這個也是公元92年的一個寫照.
就是《啟示錄》的一個時代的寫照.
那個時候小亞細亞 就是七個教會所在的地方.
也都爆發了嚴重的糧食短缺.
為什麼會這樣呢?.
就是因為皇帝 那時候是羅馬皇帝.
為了保護意大利的本土葡萄酒的經濟.
就是有錢人的經濟.
竟然下令整個小亞細亞的地方不可以種蜜.
停止去種小蜜 改要種葡萄 就是葡萄.
所以令到當地的小蜜就沒有出產了.
也都導致糧食的價格突然間暴漲.
窮人連一星的蜜都買不到.
一銀子買三星大蜜 大蜜是粗糙一點的蜜.
那時候如果你吃一些比較差的蜜.
很難吃的那些 你就可以買三星.
所以其實小蜜大蜜都沒有什麼栽種.
以致人去飢餓 窮人去飢餓.
所以你發覺很不公平.
糧食很不公平地去分配.
上高的黃華點 你就要怎麼做.
不用理民間的情況.
也都講一下今天的數據給你聽一下.
聽完你都會有些驚訝的可能.

$^{441}$我們現在的世界又是怎樣呢?.
我們目前的世界有7至10億人.
是處於飢餓的狀態 即是長期營養不良.
每天工資低於一至兩美元.
你可能會問 是不是全球的生產不夠呢?.
即是糧食不夠呢?.
全球食物的供應不夠呢?.
你猜是不是?.
你都猜到了 不是 不是.
而是全世界如果是食物供應量.
是有150\%的食物.
即是說這些食物可以供應全世界的人.
可以有150\%的人吃.
即是說人還不夠多.
食物多到可以有150\%的人吃都吃不完.
即是說我們根本是夠食物.
我們不是不夠食物.
全球都是夠食物的.
我們住在澳洲.
你知不知道澳洲的糧食出產夠不夠我們吃呢?.
很夠吧 雖然有些東西很貴.
現在很貴 都是貴的.
完全是足夠的 而且是綽綽有餘的.
知不知道澳洲的食物生產是多少呢?.
原來澳洲的食物生產可以養活.
我們自己澳洲人兩至三倍.
多少呢?即是說你吃到肚子都吃不完.
即是說你吃不到那麼多.
所以澳洲最大的出口就是食物.
兩至三倍 不要浪費.
所以我們有時候都買貴的東西.
因為有些出口多了.
所以為什麼會有十億左右的人.
在飢餓的狀態裡長期沒得吃呢?.
很慘呢?.
真正的問題不是糧食不夠.
真正的問題就是不公平地去分配.
亦都是嚴重的食物浪費.
和地區衝突所引致.
或者是一些資訊不正當地去傳播.

$^{481}$所以今天我們真的要不要浪費我們的食物.
所以我看見糧食是這樣.
但世界的經濟體系又是怎樣呢?.
其實都很脆弱.
我們經歷過金融風暴.
一天之內你看見那些人拿著那些股票.
突然間變成廢紙.
現在都是物價高漲.
因為銀紙有些國家可以不斷印.
不斷印銀紙就全球物價高漲.
不只是他的國家是全球都會.
所以十年前的一百萬.
我們說一百萬富翁.
無論你拿著港幣澳幣都叫百萬富翁.
今天一百萬真是平民來的.
屋都不止一百萬都買不到.
如果我們的安全感建立在存款,投資,房屋,房地產.
這個黑馬就告訴我們.
我們的安全感隨時會動搖.
神提醒我們.
今天如果我們的心放在財富裡面.
當你的保障.
你發覺你的心就很不平安.
因為財富控制了你的心.
財富可以高高低低.
唯有神的話語才是我們的保障.
我們不會落空.
我們有安全感.
因為神的應許是不會落空的.
我們再來看第七,第八節.
第四印揭開的時候.
第四印揭開一匹灰馬就出來了.
在馬上的名字就叫做「死」.
跟著陰間也跟隨著它.
也有權柄給它.
可以用刀劍,饑荒,瘟疫,野獸.
殺害地上四分之一的人.
所以第四匹灰色的馬是一個死亡的陰影.
我們人類最害怕的是死亡.
所以第四匹馬就是告訴你死亡期會臨到.

$^{521}$還是用不同的方法去臨到.
不同的災難去臨到.
有刀劍,饑荒,瘟疫,野獸一起出現.
造成有四分之一人的死亡.
有些譯本就說是管轄四分之一地方.
以致那四分之一人可以死亡.
不管是怎樣.
譯本是怎樣.
但總之都會有四分之一的死亡.
會發生.
我想想2019年的病毒.
新冠的病毒大家是不是覺得很厲害.
那時候我們全球都恐慌.
不只是一個地方.
不只是西班牙發現病毒.
而是全球都有病毒.
每天的新聞就是看著全球有多少人死.
每天都有多少人被傳染.
大家就不敢出街.
如果有病毒就要隔離.
隔離也很慘.
到你死的時候你的親人都不在你旁邊.
小朋友隔離突然就死了.
連家人,父母的陪伴都沒有.
或者老人家都是.
死亡的時候連親人都不在旁邊.
所以其實這些病毒都很恐怖.
帶動我們的心情都起伏很沉重.
抗疫了三年就暫告一段落.
我們都經歷過.
所以一提起病毒瘟疫.
我們都提高了警覺.
這個警號我們都開始害怕.
雖然新冠病毒我們覺得這麼厲害.
大家有沒有去查過資料.
其實死了人口有多少.
我又查一查.
我講起這篇報道查很多資料.
那時候大約有多少人死亡呢.
原來大約死了全球的.

$^{561}$有些因為很難統計.
所以是0.1\%到0.4\%.
即是全球人口的半個百分點都沒有.
所以如果說到有些災禍是有四分之一人死亡.
其實這些災禍是很大的.
亦都很恐怖.
所以這四匹馬提醒我們.
死亡的陰影其實每個人都害怕.
亦都會淋到.
所以當我們看完這四匹馬.
你有什麼感受.
真是害怕.
但這四匹馬其實是末世反復出現的現象.
因為我們人性的罪惡帶來欺騙.
或者想征服別人.
更加想征服別人就帶來一些戰爭.
戰爭亦都帶來一些饑荒.
饑荒亦都帶來死亡.
所以我們都會感到恐懼.
亦都覺得世界是否失控呢.
但啟示佬告訴我們.
或者表達給我們的訊息其實是剛剛相反.
當我們感到恐懼或者驚慌的時候.
其實啟示佬剛才所說的四匹馬.
不是要我們驚慌.
而是告訴我們這個世界真的會有災難.
因為災難是因人的罪惡而起.
不是代表神失控.
更加災難是神要介入告訴我們.
我要審判這個世界.
警號已經來了.
你要去留意神的掌管.
神要介入 要審判.
祂讓我們看到災難逼近.
是祂的審判來到了.
世界上的警號已經響了.
所以神容許這些災難逼近.
容許地上的秩序好像在動搖.
其實祂是要打破我們一個.
對物質世界一個虛假的安全感.

$^{601}$今天我們以為我們住得好就炫耀了.
以為小孩很聰明就炫耀了.
其實這些事情都會過去的.
一場山火 一個地震.
你發覺最近都是.
完全房屋倒塌了下來.
所以神已經發出了警號.
催促人們要悔改 回轉歸向神.
所以第四匹馬其實殺害四分之一人.
你發覺反而是一個神的憐憫.
不是大部份人都死了沒有機會.
而是一個警號告訴你神的憐憫.
神帶著恩慈去發出這個管教和憐憫.
所以災難的意義是什麼呢?.
是呼喚我們要去悔改.
你見人的罪惡這麼滿盈的時候.
不是你要跟隨他們一份.
你怎樣對我 我就怎樣對你.
好像小巴那兩個女子那樣.
你要踩我一腳 我就踩你一腳.
不是這樣的.
告訴你災難的意義就是要我們去悔改.
要我們懂得去反省.
而且災難你發覺是有終點 有盡頭的.
不只是不斷的災難.
盡頭就是引向我們回到神那裡.
回到天家才有盼望.
是一個新天新地的方向.
如果大家有留意剛才讀的經文.
每個印都是誰揭開的呢?.
是羔羊 即是耶穌基督.
祂有資格去揭開審判的印.
所以重點是告訴我們.
不是災難有多可怕.
而是即使災難臨到.
是耶穌自己去掌管這個世界.
去揭開這個印.
所以就算有邪惡也好.
世界上的邪惡是由神許可產生的.
所以我們知道今天我們有邪惡.

$^{641}$但是神是掌管的.
所以經文出現.
你發覺有冠冕賜給他.
即是暫時讓他贏.
或者有權柄賜給騎馬.
意思就是我給你權.
不是你有權.
或者有一把大刀賜給他.
意思就是戰爭上你可以暫時讓他贏.
這些都是說到很凶殘的災禍.
但是主權不在他的手裡.
而是在神的手裡.
所以今天我們看到.
啟示錄最大的安慰.
不是預告災難是什麼時間發生.
而是預告我們.
告訴我們.
神仍然在掌權.
神仍然作爵為王.
所以只要我們深信神仍然掌權的時候.
當我們看到這些災禍的時候.
我們不是要驚慌.
而是為這些災禍祈禱.
求神的憐憫臨到我們.
以致我們不要失控.
知道神是愛我們.
為我們捨己的神.
以致祂一定去看顧祂的百姓.
所以今天我們看到這些災禍的時候.
神不是要我們驚慌.
而是要我們知道.
誰才是主宰呢?.
亦都讓第一世紀的教會.
他們面對逼迫的時候.
其實約翰寫這封信告訴他們.
羅馬王不是歷史的主宰.
撒旦亦都不是靈界的主宰.
真正掌權的是主耶穌基督.
所以這是帶給教會一個很大的安慰.
即使我們面對逼迫,面對死亡.

$^{681}$面對一些災難.
我們知道世界的走向.
可能世界會滅亡.
但是最終的走向.
是神的角度會臨到.
耶穌基督應許我們新天新地.
我們信祂的人.
祂為我們預備了一個榮耀的居所.
就是新天新地.
之後看啟示錄我們會再留意到.
所以當世人因為罪惡帶來的驚慌.
帶來的恐慌,帶來的經濟崩潰.
當世界的人有這些事情發生的時候.
造成災難的時候.
我們看見不是一個混亂的情況.
我們看見神是坐著為王.
神是在寶座上面.
所以一開始約翰看見這些.
以致我們第四,第五章啟示錄.
人是被約翰一看見這些.
以致他的心定下來.
才被你看見災難.
所以這是啟示錄一個很大的訊息.
至於今天我們要怎樣去回應這些警號呢?.
以致我們不要驚惶失措.
第一樣就是我們不要再沉迷在這個世界.
去追尋這個世界.
今天我們沉迷世界杯.
今天打邊場可能大家都知道.
世界會過去.
但這個世界物質的事情都會過去.
唯有神的國和神的愛是不會過去的.
祂給我們有盼望的.
所以我們不要失去盼望.
我們盼望不是世界上的事.
我們盼望的是神仍然去掌權.
我們將來就算我們死的時候.
我們知道我們回到的是天家.
是神為我們預備一個不會動搖的角度.
所以我們要對神有信心.

$^{721}$我們不要怕 只要信.
第三就是趁著今天還有機會悔改.
我們應當去歸向神 悔改歸向神.
所以今天我都很含重地去和大家說.
你們有誰還沒有信主耶穌基督?.
我知道在座當中都有人還沒有信主.
我都很盼望.
很盼望你今天去做一個決定.
去相信愛我們的這位神.
你要去趁著還有機會歸向祂.
你去相信祂.
祂是那位愛你的神.
也都趁著有機會.
也是把握時間 我們更加要傳福音.
我們的家人需要福音.
我們的朋友需要福音.
我們不知道明天的世界會怎樣.
我們要活在現在的時候.
我們要活出愛.
做一個和平的使者.
做一個珍惜你家人朋友的信徒.
以致我們願意將福音帶給我們未信的家人.
帶給我們需要愛的家人.
所以審判雖然是逼近.
但是救恩的門是很大大地開著.
災難雖然是增多.
但是我們記得神仍在掌權.
最後我想分享我自己很喜歡的經文.
在希伯來書12章1到2節.
說到我們既有許多的見證人.
如同雲彩圍著我們.
即是說我們有教會.
我們有很多周圍很美麗的弟兄姊妹.
好像雲彩圍著我們.
其實是加力給我們的.
要我們放下各樣的重擔.
斷去容易纏繞我們的罪.
我們很容易放棄.
我們有很多重擔圍著我們.
以致我們經常都很累.

$^{761}$但是神叫我們要脫離這些.
然後全心忍耐.
奔那擺在我們前頭的路程.
仰望為我們順心創始成中的耶穌.
所以當我們去忍耐的時候.
你會發覺神一定會開出路.
耶穌是怎樣的一個愛我們的神呢?.
祂因擺在前面的喜樂.
就興旱羞辱.
忍受了十字架的苦難.
便坐在神寶座的右邊.
所以耶穌也試過很大的死亡.
死的時候多辛苦在十字架上.
流血流死的.
幾個小時流血.
祂愛我們愛到一個地步.
就是要我們去有個救恩.
我特別喜歡的就是我提到的那句.
我希望大家我們一起來讀.
「全心忍耐」開始.
我們一起來讀兩次.
看看可不可以記得背了它.
全心忍耐.
奔那擺在我們前頭的路程.
仰望為我們順心創始成中的耶穌.
我們再來一次.
全心忍耐.
奔那擺在我們前頭的路程.
仰望為我們順心創始成中的耶穌.
所以我很喜歡這個奔字.
我們不是去走.
不是去赤赤裸裸地這樣走.
我們是去奔.
我們去盡力去走.
神要為我們預備的路程.
秘訣就是為何我們可以有力可以奔走呢?.
秘訣就是仰望為我們順心創始的.
創始成中的耶穌.
祂是多麼愛我,為我去捨棄.
以致我們今天可以有能力去愛人.

$^{801}$也可以活出愛.
願我們一起去奔走這條道路.
我們一起去祈禱.
恩愛的阿爸父.
我們知道你很想我們去奔走這條天國的道路.
你很想我們去仰望你.
但當我們看見災難的時候.
我們很容易去懼怕.
很容易去失落.
很容易害怕.
求主去兼顧我們的信心.
無論什麼境況.
讓我們都能夠仰望.
為我們順心創始成中的耶穌.
求主去給我們先嘗主恩的滋味.
我們如果沒有嘗過主恩的滋味的時候.
我們就不能體驗到這份愛是多麼的大.
求主去進到我們的心裡.
讓我們看見神的恩典.
讓我們嘗到神的恩典.
讓我們看見世界上雖然有很多的災難.
但你的愛從來沒有離開我們.
讓我們懂得去投靠你.
讓我們知道你是美善的.
如你在聖經裡說過.
我們要先嘗嘗主恩的滋味.
便知道主你是美善的.
投靠你的人是有福的.
當我們活在這個混亂的世代的時候.
我們最有福的事情是投靠你.
很多事情在我們的掌控之下.
雖然我們很想去掌管.
但我們不知道明天是怎樣.
唯有神你知道.
讓我們都可以上傳忍耐的心.
願意去盡力奔走.
擺在我們前面的路程.
積極的向我們家人去傳福音.
以示我們可以活出愛.
去吸引更多人去跪向你.

$^{841}$求主來帶領我們.
多謝您聽我們祈禱.
奉主耶穌的獨特性的名祈求.
大家好.
\newpage



\section{列王記上 11:1-43}
\label{sec:032439}
\textbf{愚昧的智者:心偏離神一寸,人生偏離神一里}~~~~袁惠鈞牧師博士
\newline
\newline
連結: \href{http://www.church.com.hk/acms/content.asp?site=cdc\&op=show\&type=product\&code=032439\&layout=sermon}{\texttt{032439}} ~~~~ 語音日期: 2026-07-12
\newline
\newline
\hyperref[sec:032438]{\small{< < < PREV SERMON < < <}}
~
\hyperref[sec:index_chronic]{\small{[返順時目]}}
~
\hyperref[sec:032440]{\small{> > > NEXT SERMON > > >}}
\newline
\newline
列王記上 11:1-43
\newline
\begin{longtable}{cl}
\hline
\hline
章節 & 經文 (和合本修訂版)\\
\hline
11:1 & \begin{tabularx}{0.7\textwidth}{X} 所羅門王在法老的女兒之外，又寵愛許多外邦女子，就是摩押女子、亞捫女子、以東女子、西頓女子、赫人女子。 \end{tabularx} \\ \\ \relax
11:2 & \begin{tabularx}{0.7\textwidth}{X} 論到這些國的人，耶和華曾吩咐以色列人說：「你們不可跟他們通婚，他們也不可跟你們在一起，因為他們一定會誘惑你們的心去隨從他們的神明。」所羅門卻為了愛，緊緊跟從他們。 \end{tabularx} \\ \\ \relax
11:3 & \begin{tabularx}{0.7\textwidth}{X} 所羅門娶七百個公主，三百個妃嬪。這些妻妾誘惑他的心。 \end{tabularx} \\ \\ \relax
11:4 & \begin{tabularx}{0.7\textwidth}{X} 所羅門年老的時候，他的妻妾誘惑他的心去隨從別神，不像他父親大衛以純正的心順服耶和華－他的神。 \end{tabularx} \\ \\ \relax
11:5 & \begin{tabularx}{0.7\textwidth}{X} 所羅門隨從西頓人的女神亞斯她錄和亞捫人可憎的米勒公。 \end{tabularx} \\ \\ \relax
11:6 & \begin{tabularx}{0.7\textwidth}{X} 所羅門行耶和華眼中看為惡的事，不像他父親大衛專心順從耶和華。 \end{tabularx} \\ \\ \relax
11:7 & \begin{tabularx}{0.7\textwidth}{X} 那時，所羅門為摩押可憎的基抹和亞捫人可憎的摩洛，在耶路撒冷對面的山上建造丘壇。 \end{tabularx} \\ \\ \relax
11:8 & \begin{tabularx}{0.7\textwidth}{X} 他為所有的妻妾，就是那些向自己神明燒香獻祭的外邦女子，也是這樣做。 \end{tabularx} \\ \\ \relax
11:9 & \begin{tabularx}{0.7\textwidth}{X} 耶和華向所羅門發怒，因為他的心偏離了向他顯現兩次的耶和華－以色列的神。 \end{tabularx} \\ \\ \relax
11:10 & \begin{tabularx}{0.7\textwidth}{X} 耶和華曾吩咐他這件事，不可隨從別神，他卻沒有遵守耶和華所吩咐的。 \end{tabularx} \\ \\ \relax
11:11 & \begin{tabularx}{0.7\textwidth}{X} 所以耶和華對他說：「你既然是這樣，不遵守我所吩咐你守的約和律例，我必定把國度撕裂離開你，將它賜給你的大臣。 \end{tabularx} \\ \\ \relax
11:12 & \begin{tabularx}{0.7\textwidth}{X} 然而，因你父親大衛的緣故，我不在你的日子行這事，而要從你兒子的手中撕裂這國。 \end{tabularx} \\ \\ \relax
11:13 & \begin{tabularx}{0.7\textwidth}{X} 只是我不撕裂全國，卻要因我僕人大衛和我所選擇的耶路撒冷，保留一個支派給你的兒子。」 \end{tabularx} \\ \\ \relax
11:14 & \begin{tabularx}{0.7\textwidth}{X} 耶和華使以東人哈達興起，作所羅門的敵人；他是以東王的後裔。 \end{tabularx} \\ \\ \relax
11:15 & \begin{tabularx}{0.7\textwidth}{X} 大衛在以東的時候，約押元帥上去埋葬陣亡的人，殺了以東所有的男丁。 \end{tabularx} \\ \\ \relax
11:16 & \begin{tabularx}{0.7\textwidth}{X} 約押和以色列眾人在以東住了六個月，直到把以東的男丁盡都剪除。 \end{tabularx} \\ \\ \relax
11:17 & \begin{tabularx}{0.7\textwidth}{X} 那時哈達還是幼童；他和他父親的臣僕，以及幾個以東人逃往埃及。 \end{tabularx} \\ \\ \relax
11:18 & \begin{tabularx}{0.7\textwidth}{X} 他們從米甸起行，到了巴蘭，再從巴蘭帶著幾個人來到埃及，到埃及王法老那裡。法老給他房屋，吩咐給他糧食，又把地賜給他。 \end{tabularx} \\ \\ \relax
11:19 & \begin{tabularx}{0.7\textwidth}{X} 哈達在法老眼前大蒙恩寵，法老就把王后答比匿的妹妹嫁給他。 \end{tabularx} \\ \\ \relax
11:20 & \begin{tabularx}{0.7\textwidth}{X} 答比匿的妹妹給哈達生了一個兒子，叫基努拔。答比匿使基努拔在法老的宮裡斷奶，基努拔就與法老的眾子一同住在法老的宮裡。 \end{tabularx} \\ \\ \relax
11:21 & \begin{tabularx}{0.7\textwidth}{X} 哈達在埃及聽見大衛與他祖先同睡，約押元帥也死了，就對法老說：「請你讓我走，我要回本國去。」 \end{tabularx} \\ \\ \relax
11:22 & \begin{tabularx}{0.7\textwidth}{X} 法老對他說：「你在我這裡有甚麼缺乏？看哪，你竟想要回你本國去！」他說：「我沒有缺乏甚麼，只是懇求王准我回去。」 \end{tabularx} \\ \\ \relax
11:23 & \begin{tabularx}{0.7\textwidth}{X} 神又使以利亞大的兒子利遜興起，作所羅門的敵人。他曾逃避主人瑣巴王哈大底謝。 \end{tabularx} \\ \\ \relax
11:24 & \begin{tabularx}{0.7\textwidth}{X} 大衛擊殺瑣巴人的時候，利遜召集了一群人，自己作他們的領袖。他們往大馬士革，住在那裡，在大馬士革建立王國。 \end{tabularx} \\ \\ \relax
11:25 & \begin{tabularx}{0.7\textwidth}{X} 所羅門活著的時候，除了哈達為患之外，利遜也作以色列的敵人。他憎恨以色列，作了亞蘭人的王。 \end{tabularx} \\ \\ \relax
11:26 & \begin{tabularx}{0.7\textwidth}{X} 尼八的兒子耶羅波安也舉起手來攻擊王。他是所羅門的臣僕，以法蓮支派的洗利達人；他母親是個寡婦，名叫洗魯阿。 \end{tabularx} \\ \\ \relax
11:27 & \begin{tabularx}{0.7\textwidth}{X} 他舉手攻擊王是因先前所羅門建造米羅，修補他父親大衛城缺口的這件事。 \end{tabularx} \\ \\ \relax
11:28 & \begin{tabularx}{0.7\textwidth}{X} 耶羅波安是個大有才能的人。所羅門見這青年殷勤，就派他監管約瑟家所有服勞役的工人。 \end{tabularx} \\ \\ \relax
11:29 & \begin{tabularx}{0.7\textwidth}{X} 那時，耶羅波安出了耶路撒冷，示羅人亞希雅先知在路上遇見他；亞希雅身上穿著一件新衣。田野中只有他們二人，沒有其他的人。 \end{tabularx} \\ \\ \relax
11:30 & \begin{tabularx}{0.7\textwidth}{X} 亞希雅拿起穿在自己身上的新衣，把它撕成十二片， \end{tabularx} \\ \\ \relax
11:31 & \begin{tabularx}{0.7\textwidth}{X} 對耶羅波安說：「你可以拿十片。耶和華－以色列的神如此說：『看哪，我必從所羅門手裡撕裂這國，把十個支派賜給你。 \end{tabularx} \\ \\ \relax
11:32 & \begin{tabularx}{0.7\textwidth}{X} 我因我僕人大衛和我在以色列眾支派中所選擇的耶路撒冷城的緣故，仍為所羅門留一個支派。 \end{tabularx} \\ \\ \relax
11:33 & \begin{tabularx}{0.7\textwidth}{X} 因為他們離棄我，敬拜西頓人的女神亞斯她錄、摩押的神明基抹和亞捫人的神明米勒公，沒有像他父親大衛一樣遵從我的道，行我眼中看為正的事，守我的律例典章。 \end{tabularx} \\ \\ \relax
11:34 & \begin{tabularx}{0.7\textwidth}{X} 但我不從他手裡奪走整個國家，卻使他在活著的日子作君王，是因我所揀選的僕人大衛遵守我的誡命律例。 \end{tabularx} \\ \\ \relax
11:35 & \begin{tabularx}{0.7\textwidth}{X} 我必從他兒子手裡將王國奪走，賜給你十個支派， \end{tabularx} \\ \\ \relax
11:36 & \begin{tabularx}{0.7\textwidth}{X} 只留一個支派給他的兒子，使我僕人大衛在我所選擇立我名的耶路撒冷城那裡，在我面前常有燈光。 \end{tabularx} \\ \\ \relax
11:37 & \begin{tabularx}{0.7\textwidth}{X} 我選你，使你照你心裡一切所願的作王，成為以色列的王。 \end{tabularx} \\ \\ \relax
11:38 & \begin{tabularx}{0.7\textwidth}{X} 你若聽從我一切所吩咐你的，遵行我的道，行我眼中看為正的事，謹守我的律例誡命，像我僕人大衛所行的，我就與你同在，為你立堅固的家，像我為大衛所立的一樣，將以色列賜給你。 \end{tabularx} \\ \\ \relax
11:39 & \begin{tabularx}{0.7\textwidth}{X} 我必因這事使大衛的後裔遭受患難，但不是永遠的。』」 \end{tabularx} \\ \\ \relax
11:40 & \begin{tabularx}{0.7\textwidth}{X} 所羅門想要殺耶羅波安，耶羅波安起身逃往埃及。他到了埃及王示撒那裡，就住在埃及，直到所羅門死了。 \end{tabularx} \\ \\ \relax
11:41 & \begin{tabularx}{0.7\textwidth}{X} 所羅門其餘的事，凡他所做的和他的智慧，不都寫在《所羅門記》上嗎？ \end{tabularx} \\ \\ \relax
11:42 & \begin{tabularx}{0.7\textwidth}{X} 所羅門在耶路撒冷作全以色列的王四十年。 \end{tabularx} \\ \\ \relax
11:43 & \begin{tabularx}{0.7\textwidth}{X} 所羅門與他祖先同睡，葬在他父親大衛的城裡，他兒子羅波安接續他作王。 \end{tabularx} \\ \\
[1ex]
\hline
\hline
\end{longtable}
$^{1}$《列王記》,《合一王朝》系列的第十講.
我每五講或者十講就會提醒大家.
我講這麼多道,這麼多的系列做什麼?.
就是希望鼓勵大家一起來讀經.
如果你是一個星期聽我五篇道.
你可以是兩年半就讀完整部聖經.
而且你是會對聖經有很深入的認識.
為什麼?因為我是用聖經來解釋聖經.
我是用字面來解經.
所以如果大家一起讀經.
是會明白怎樣來閱讀聖經.
能夠有很深入的了解.
而我在每一個系列第一篇道.
我都是會講一個介紹.
給大家很多的資料,很多的背景.
以致是能夠更加深入明白聖經.
我今天要講的是《列王記》上的第十一章.
要講所羅門最後的故事.
他人生最後的故事.
我給這個題目是愚昧的智者.
誰是愚昧的智者?.
就是所羅門.
神給他智慧成為世上最有智慧的人.
但是他就沒有用到這些智慧.
他是用了他自己的手段.
很多時作出很愚昧的行為.
所以我們都要小心.
今天的副標題就是.
心偏離神一寸.
人生偏離神一里.
當我們在心裡面.
慢慢慢慢離棄神的時候.
我們可能不覺得.
但是當我們的心偏離一寸.
我們所做出來的事.
我們的人生就會偏離神很遠.
所以我們要很小心.
愚蠢的智者就是所羅門.
聖經絕對不會忽視.
在聖經裡面英雄人物的缺點.

$^{41}$有人說聖經就好像一幅大特寫.
將人的皺紋.
將人的雀斑全部顯出來.
我有一天在家裡.
突然間找到一個鏡子.
我是無意中找到的.
這個鏡子是凹了的.
凹了進去中間那裡.
原來我發覺是一個放大鏡.
我可以見到自己的臉上.
那些皺紋和雀斑全部放大了.
我說見鬼了.
為甚麼會有這樣的鏡子.
後來我才發覺.
原來是那些女士要來做甚麼.
發現自己有皺紋和雀斑.
知道用甚麼面相.
用甚麼辦法來護理.
聖經就是這樣.
你不要以為那些英雄人物.
聖經就不會說出他們的弱點.
所羅門也是一樣.
所羅門是神給他最大的智慧.
而且有難以置信的財富.
但他始終有一段時間離棄了神.
作出了愚蠢的決定.
以致使到以色列國分裂.
在《箴言》的十九章三節.
是所羅門自己的寫作.
他為人的愚昧傾拜了他的道.
一個人不是不知道神的話語說甚麼.
但很多時會怎樣.
很多時在自己的身上疏忽了.
甚至像所羅門一樣.
他懂得教人.
但他自己就作出了愚蠢的行為.
我們很容易在所羅門的身上.
就能夠發覺他離棄神的原因是甚麼.
這就是我大綱第一和第二點.
第一點是甚麼.

$^{81}$所羅門違背了神的話語.
列王記上的十一章一至八節.
第二點是他沒有聽到神的勸告.
甚至勸告他他都不聽.
九至二十五節.
於是神就興起了他的敵人.
其中一個就是以法連人耶羅波安.
在二十六至四十節.
還有所羅門最後是與列祖同睡.
四十一至四十三節.
我剛才說我是鼓勵大家一齊讀經.
如果你今天聽我這篇道.
就要讀列王記上的十一章.
但不要忘記讀歷代志下的第九章.
因為所講的都是一樣的內容.
第一點我要和大家講的.
就是所羅門違背神的話語.
神的話語怎麼說.
不可回到埃及去.
那所羅門何時回到埃及.
他第一點就是娶了一個埃及的公主.
在列王記上的第三章一至.
然後他去到埃及頻頻賣馬匹.
他這兩個行為.
讓我們見到所羅門不是真的信靠神.
為什麼他要娶埃及的公主.
他希望和法老王聯盟.
一起來對抗外敵.
他為什麼要買馬匹.
他相信自己的實力多過相信神.
他根本不相信耶和華可以保護這片土地.
他想著自己一定要做些事.
所以他和埃及王聯婚.
他是買馬匹來鞏固自己的軍隊.
他忘記了他自己爸爸的話.
在《說詩篇》的二十篇第七節.
大衛怎麼說.
有人靠車有人靠馬.
但我們要提到耶和華我們神的名.
神才是最重要.

$^{121}$你是和埃及王來納文約.
沒有用的.
等一下我們就會見到.
買了這麼多馬匹都是沒有用.
所羅門娶外邦妻子.
他造成了很壞的榜樣.
直至到四世紀以後.
在《以斯拉記》和《尼希米記》裡面.
都有提到.
以色列人都和他一樣.
是和外邦的女子和外邦人來通婚.
從聖經的角度來看.
出埃及是什麼意思.
在信徒的身上我們出埃及.
就是我們脫離罪的綑綁.
就是接受耶穌基督之後.
我們得救.
被神稱為義人.
我和大家說過.
基督徒的人生有三個步驟.
清義,成聖和得榮耀.
所以當我們說我們出埃及.
就是我們接受了神脫離了罪的綑綁.
我們已經是被神稱為義人.
那你說出埃及去哪裡?.
去曠野.
以色列人去曠野做什麼?.
鍛鍊他們.
好像我們基督徒一樣.
在我們每天的路上.
我們都能夠成聖.
慢慢越來越像耶穌基督.
所以除了清義第二步就是成聖.
就好像以色列人在曠野.
最後進入迦南地.
就是當我們真真正正.
走了我們人生的路程.
我們能夠得到榮耀的身體.
我們能夠與神同在.
這個就是最後一步.

$^{161}$得榮耀.
所以說一個基督徒.
回到埃及不是一件好事.
即是說這個人仍然被俗世纏繞.
不能夠真真正正信靠神.
亦都好像馬太福音十三章.
撒種的比喻.
其中一種的泥土是什麼?.
是有荊棘的泥土.
種子進入泥土裡面.
很快發芽.
而且很快飆高.
但是慢慢慢慢就被俗世.
仍然纏繞著他.
以致不能夠結出果子.
不能夠結出果子.
就不是一個真正的信徒.
所以這個我們要很小心.
所以說基督徒.
信了之後仍然回到埃及.
絕對不是一件好事.
還有一個可能性.
我們可以說一個基督徒.
仍然在曠野徘徊.
代表什麼?.
就是一個基督徒不成熟.
不是真正服從神.
沒有使命感.
在屬靈的生命當中.
停留在一個地方不前進.
不相信神的應許.
毫無目的地在這裡生活.
這個都是非常非常之危險.
我們都要很注意.
在《禮王紀》上的十一章.
三節,四節和九節.
我們剛剛讀過經文.
這個作者是引用了.
《申命記》的七章四節.
《申命記》的七章四節說什麼?.

$^{201}$那些愛邦的女子.
是會使你的心轉離.
不跟從你的主.
去事奉別神.
而這一句話.
用在所羅門的身上.
是非常之的適合.
所羅門經文說.
他寵愛或者戀愛這些女子.
在經文希伯來文的原文.
就是dabak這個字.
《申命記》裡面.
叫以色列人要做什麼?.
要專靠神.
《申命記》的十章,十一章.
十三章和三十章.
專靠神都是dabak這個字.
那你說為什麼?.
去戀愛那些愛邦女子.
又是dabak.
專靠神又是dabak.
是呀,看你向哪一邊.
你去專靠那些愛邦女子.
專愛她們.
那你就dabak她們.
如果你專靠神.
那你就dabak神.
所以都是同一個字.
而在《約書亞記》裡面.
二十三章八節和十二節.
約書亞都用了很相似的言語.
他跟以色列人說.
你們要專靠耶和華.
你們的神dabak.
不要稍微轉去.
與你們中間所剩下的那些國民聯絡.
聯絡這個字又是dabak.
聯絡了就會怎樣?.
彼此結親互相往來.
有沒有看到?.

$^{241}$又是兩面倒.
你可以去專靠這些愛邦人.
跟他們結盟約,結婚.
但你也可以專靠神.
聖經裡面給我們有很大的警告.
所羅門很可惜.
他專靠這些愛邦的女子.
而沒有專靠到耶和華.
而在今日的經文.
列王記上的十一章.
四次又提到nata這個字.
希伯來文是什麼?.
被引誘.
三次和合本翻譯成.
這些愛邦女子引誘所羅門去拜偶像.
一次被翻譯為.
帶領他偏離離開神.
在十一章的九節.
讀者讀到所羅門獻聖殿的時候.
他怎樣說?.
要以色列的人民.
專向神.
遵行神的道.
這個人有口說.
以色列人你們要專向神.
他自己偏離神.
被愛邦女子引誘他離棄神.
耶和華還兩次提醒他.
第一次是欺騙.
就是給所羅門有智慧的時候.
第二次所羅門獻了聖殿.
神都來提醒他.
但是最後所羅門都沒有遵守到.
神的律例殿章.
非常之可惜.
所羅門他要人民.
以誠實的心來向著神.
在列王記上的八章六十一節.
誠實的心就是純正的心.
不可以分割的心.

$^{281}$他要以色列人完全馴服神.
而他自己呢?.
他沒有完全偏離神.
你不要以為他不去信神.
他只不過是把耶和華.
變成他其中敬拜的偶像之一.
也就是說他一面拜偶像.
一面敬拜耶和華.
他是個兩頭蛇.
這個就是聖經裡面所講.
心懷異議的人.
雅各書的一章八節.
心懷異議的人.
在他一切所行的路上.
都沒有定見.
而當所羅門真的這樣做的時候.
他就違背了十誡第一和第二誡.
是什麼?.
除我以外不可有其他的神.
還有第二誡呢?.
不可為自己雕刻偶像.
他就是離棄了神.
他就是在自己的生命當中有偶像.
而這個也正是.
以塞瓦書的四十二章第八節.
神說我是耶和華.
這是我的名.
我必不將我的榮耀歸給假神.
什麼意思?.
除我以外別無他神.
也不將我的稱讚歸給雕刻的偶像.
就是第二誡.
不可為自己雕刻偶像.
而和以塞瓦書的四十六章九節.
都是很相似.
神說我是神.
並無別神.
除我以外別無他神.
然後說我是神.
再沒有能給我的.

$^{321}$這些偶像只不過是石頭.
只不過是銅銅鐵鐵造.
怎能夠給耶和華.
所以這兩次的經文其實就是提醒我們.
十誡當中的第一誡和第二誡.
我們要知道.
所羅門離棄神.
不是突然睡醒.
我今天開始要敬拜偶像.
我不敬拜耶和華.
不是的.
他是慢慢一步一步離棄神.
開始就是在列王記上的三章一節.
他首先娶了埃及的公主.
愛邦女子.
然後他是隨從愛邦的妻妾的神.
在列王記上的第十一章的四節和第五節.
然後他在山上建一些幽壇.
來敬拜這些假神.
就是他妻妾的這些假神.
列王記上的十一章七節.
然後最初在宮外.
在山上來敬拜這些假神.
然後他就將這些假神搬了進去皇宮.
准許他的妻妾在皇宮裡面敬拜.
十一章的八節.
他以前只是向耶和華來獻祭燒香.
現在他居然是在宮殿裡面.
來敬拜他妻妾的假神.
當以色列的王自己都敬拜偶像的時候.
他能不能夠成為外方人的光?.
他不能夠.
所羅門敬拜的是甚麼偶像?.
聖經有講得很清楚.
有摩洛神.
就是阿們人所敬拜的神.
他們敬拜摩洛神的時候是很殘忍.
甚至那些人會將自己的子女.
那些嬰兒殺死.
來獻給摩洛神.

$^{361}$這樣的行為是神所憎惡的行為.
還有就是姬末.
姬末就是摩洛人的主神.
無論是阿們人還是摩洛人.
都是亞伯拉罕的侄子.
羅德的後人.
在創世記的十九章告訴我們.
所羅門還敬拜甚麼偶像?.
阿斯他祿的女神.
這個神是推羅和西頓人所敬拜.
是一個新玉的女神.
所以當他們敬拜阿斯他祿的時候.
無論男女在廟裡.
都會和廟祭來行人.
因為這是一個新玉的女神.
這樣的行為是神所憎惡的行為.
這個阿斯他祿.
也是巴比倫人所敬拜的女神.
我講完了第一點.
就是所羅門沒有服從道神的話語.
第二點就是他沒有聽到神的勸告.
其實所羅門剛生出來的時候.
神很愛他.
透過那丹仙芝給了他一個很特別的名字.
就是耶底底亞.
甚麼意思?.
就是神所愛.
但是現在經文告訴我們.
神向所羅門發怒.
因為他的心轉離了神.
所羅門違背了神賜給他的智慧.
忽視了神給他的亮光.
其實神給所羅門一個爸爸大衛.
雖然大衛不是完美的人.
但是大衛是一個忠於神的人.
這一點所羅門是能夠看得到.
大衛也曾經為所羅門禱告.
鼓勵他遵行神的旨意.
鼓勵他來起聖典.
神也兩次向所羅門顯現.

$^{401}$提醒他神和大衛所納的約.
所羅門絕對知道.
神和以色列人立約.
在《說明記》的28章至30章.
是講到這個文約的條本.
也就是說如果以色列人守約的話.
神就會祝福.
但是如果以色列人違約的話.
神就會審判,會詛咒.
所以所羅門絕對不能夠推搪.
說我不知道這些事.
為什麼?.
因為他自己在獻聖典的時候.
他已經提過這些文約.
他講到如果以色列人能夠守約.
神會祝福.
如果違約神就會詛咒.
而他自己就是沒有遵守到神的約.
我們剛才說.
神向所羅門來發怒.
因為所羅門的心轉離了神.
神是看內心的.
在撒姆以上的16章7節也是這樣說.
神是監察內心的.
而所羅門自己在箴言的4章23節.
所羅門也勸人要做什麼?.
要保守你的心.
因為生命的果效是由心發出.
有沒有看到?.
所羅門知道要怎麼做.
但是他自己就做不到.
他的心是遠離了神.
其實我們的心真的很重要.
我們現在看到Annie Wong姊妹的心出現問題.
會影響整個人.
心就是將血液泵到我們全身的細胞.
全身的器官.
來供給氧氣,營養.
所以心是很重要.
這是實體.

$^{441}$但是在靈裡面心也是很重要.
神召我們要怎樣?.
申命記的6章5節.
要盡心盡性盡力來愛耶和華我們的神.
神甚至將祂的話語刻在我們的心板上.
箴言的7章1至3節.
神一早就要以色列人行心的覺禮.
不單止是身體上的覺禮.
身體上只不過是外在.
但是神是看內心.
所以是心的覺禮.
當我們的心出現了問題的時候.
我們整個生命都會出現問題.
我們不知道神是什麼時候.
向所羅門發出這些警告.
我相信可能是透過一個先知.
神警告所羅門.
因為你不遵守我的約和我的律例.
我一定會將這個國從你奪回.
給你的其中一個臣子.
其中一個僕人.
這個遲些我們就知道.
就是耶羅伯安.
他就是所羅門的一個臣僕.
但是因為神和大衛所納的約.
神說我仍然會留一個支派給大衛的後人.
問題就來了.
有人說南國起碼有兩個支派.
有猶大和汴雅文.
為什麼一個支派?.
如果你認識約書亞.
在約書亞記裡面分地的話.
你就知道最初西面支派.
在猶大支派的境內.
你說加上西面猶大汴雅文.
三個支派怎麼會是一個支派?.
但是你要知道在南國.
所有忠以敬拜偶像不歸向神的人.
已經不遷去了以色列北國.
他們在那裡拜偶像.

$^{481}$南國是去耶和華的聖殿.
在耶路撒冷來敬拜.
剩下來那些真的歸向神的人.
他們同化了猶大支派.
所以經文說剩下一個支派給大衛的後人.
絕對是正確.
而的確在南國最大的支派是誰?.
就是猶大支派.
在列王記上的十一章.
十一節,十二節和十三節.
神說我會奪回你的王國.
奪回這個字在原文的希伯來文.
就是撕裂的意思.
將它撕裂了.
的確在列王記上的十一章,二十九節.
神差派一個先知叫做阿希雅.
用行動表現出這個寓言.
這個阿希雅在耶羅波安面前做什麼?.
將自己的新衣服撕成十二塊.
你能不能夠想像.
穿一件很漂亮的衣服外出.
突然間神給你一個啟示.
將新衣服撕裂成十二塊.
所以在聖經裡面.
很多先知他們做出來的行為.
都是很特別.
你不要以為他們瘋了.
他們不是.
他們是得到神的啟示.
然後就將十二塊撕裂的衣服.
其中十塊給了耶羅波安.
所以我們見到.
這個王國的分裂.
是不是人的行為?.
絕對不是.
是神的作為.
將以色列國分開南北國.
其實在歷史裡面.
以前亦都出現過這些事.
將衣服撕裂了.

$^{521}$記不記得.
蘇羅在《撒母耳》上的十五章.
他拉著撒母耳的衣服不肯放.
撒母耳就不理他,就走.
他一拉撒母耳的衣服.
將撒母耳的外袍撕裂了.
其實那裡都有一個預表.
就是神將會從蘇羅手上.
拿走王國.
然後給了大衛.
我們都講過.
在所羅門的時候.
我們沒有見到先知來勸所羅門.
大衛的時候.
頻頻你會見到先知來勸大衛.
似乎先知好像無庸無之地.
但是這個時間.
就快去到猶大國的末期.
先知的職位會再一次成為重要.
先知時代就會透過.
以利亞和以利沙來到這個世上.
而神就會再次透過先知.
來跟君王和祭司講說話.
當這些君王祭司.
背叛神的話語的時候.
這些先知就會來勸他們.
如果他們不聽的話.
甚至先知是會講出神的審判.
神的懲罰.
神給所羅門的懲罰.
神給所羅門的審判.
其實是延遲了.
神是沒有在所羅門在世的時候.
將他的國拿走.
將他的國分為南北兩個國.
而是等他死了先.
你說為什麼?.
當然這個是神的恩典.
是神的憐憫.
但是還有一個原因.

$^{561}$就是他是等待所羅門.
來回轉歸回.
我曾經講過回轉歸回.
或者悔改希伯來文是哪個字?.
是sub這個字.
就是180度的轉回.
神是希望所羅門會改過.
所以他沒有立刻執行這個審判.
在聶王紀上的十三章三十三節.
我們也是看得到.
也是發生在耶羅波安的身上.
就是那個拿走了.
十個支派的所羅門的僕人.
經文講什麼?.
聶王紀上的十三章三十三節.
這是以後耶羅波安仍不離開他的惡道.
離開這個字就是轉離悔改.
就是sub這個字.
所以你看到無論耶羅波安.
所羅門也是.
神是不斷給他們機會來改過.
這一點我們在耶利米書的十八章七至八節.
我們也看到.
神和耶利米說.
你去告訴我的國民聽.
如果我宣佈有一個民族或者王國.
我將會使他連根拔起.
我會拆毀他,我會毀滅他.
但是如果這一幫這一國.
願意轉移離開他們的惡.
我就必後悔.
不將我想要施行的災禍.
臨於他們.
有沒有看到我們這位的神.
就是一位憐憫有恩典的神.
他是不斷地在等人來悔改.
今時今日都是一樣.
為什麼耶穌基督還沒有回來.
神有祂的時候.
大神就願萬人得救.

$^{601}$不願一人沉淪.
所以我們等了這麼久.
我們說這個系列的時候.
我們也說過.
神其實和大衛納的約是無條件.
雖然神會從所羅門的兒子羅波安的手裡.
拿回這個王國.
是會給了耶羅波安.
但是神仍然會留一個支派給大衛的後人.
因為神和大衛立約.
這個是無條件的約.
神是會保存大衛的皇朝.
這個是神的恩典.
而的確耶穌基督是會透過猶大支派來到世上.
你有沒有看到神的作為.
真的非常巧妙.
雖然神和大衛納的約是無條件.
但是神和大衛後人納的約是有條件.
包括所羅門.
你看看聶王紀上的三章十四子.
他和所羅門說.
如果你服從我的話.
我必賜給你長壽.
等一下我們會回來看這一點.
看所羅門究竟他的壽命.
究竟是長還是短.
神的確賜給所羅門有智慧.
有辨別的能力.
能夠辨別是非.
這一點在所羅門審判兩個妓女的案子.
我們都看到.
的確是很有智慧.
但是這個智慧用在別人的身上.
他就能夠分別是非.
用在自己的身上.
污糟糟.
完全盲目.
這個就是所羅門.
其實神的管教.
所羅門絕對不會陌生.

$^{641}$在他自己爸爸大衛的身上.
所羅門就見到神的管教.
而且當神在機編第一次和所羅門見面.
在聶王紀上的第三章的時候.
神已經提醒了所羅門.
如果你不服從我的話.
我的管教會臨到你的身上.
而當所羅門奉獻聖殿.
神又再一次重覆這一句話.
而在說詩篇八十九篇.
亦都是有同樣的話.
但是很可惜.
所羅門是沒有回轉.
是沒有歸回.
沒有悔改.
沒有失.
那你說所羅門最後有沒有接受神.
有沒有歸向神.
這一點我們遲些也會回來看.
所羅門以為他自己用自己的手段.
和鄰國和外邦來通婚.
就可以和這些國家和平相處.
他自己還口出狂言.
誇口說.
他說現在耶和華我的神使我四周平安.
沒有仇敵沒有災禍.
在聶王紀上的五章四節.
他遲些很快就會見到.
神會差遣他的敵人來攻擊他.
他這個四周平安只不過是暫時.
神就會興起三個敵人來管教所羅門.
來攻擊所羅門.
其中第一個就是以東人哈達.
這個人是以東人.
我們知道剛才讀經文的時候.
所羅門所娶的外邦妻子是有以東人.
她當然希望和以東來立約.
大家和平相處.
但是她不能夠阻止.
以東王哈達來攻擊自己.

$^{681}$有沒有看到.
她這些手段完全是白費心機.
看著她爸爸的朝代.
大衛和約納曾經戰勝以東人.
滅絕了所有的南丁.
但是這個哈達就逃脫了.
有一班領袖帶領他.
去到埃及法老王那裡.
來尋求庇護.
法老王不單止收留了他.
給他食物,給他居住.
還將皇后的妹妹嫁了給他.
和他生了一個兒子.
你說這個埃及王有沒有搞錯.
不是吃了所羅門茶禮嗎?.
他不是嫁了女兒給所羅門嗎?.
現在又和以東王來交涉.
這裡就讓我們看到.
所羅門白費心機.
靠自己的方法.
埃及王的確將女兒嫁給所羅門.
但是這個收留哈達的埃及王.
是一個新王.
他不念舊情.
可能他根本不知道.
忘記了原來前面的埃及王.
曾經將女兒嫁給所羅門.
所以有沒有看到.
其實所羅門應該倚靠的.
是耶和華,是可靠的.
但是他就靠了自己的手段.
希望和這些國家立約.
結果如何?得不償失.
這個新的法老王.
遲些我們就會見到.
他是侍殺.
不是前面我們所說的.
埃及王西亞蒙.
現在大衛和約阿都已經死了.
哈達是能夠回到以東.

$^{721}$不再害怕大衛.
他回到以東又如何?.
增強了他自己的國家.
開始向以色列進行一連串的攻擊.
以東是不夠勢力擊敗所羅門.
因為所羅門的王國比以東大.
但是神就用哈達來教訓所羅門.
管教所羅門.
希望所羅門會得回心轉意.
除了這個以東王哈達.
還有第二個敵人.
是大馬士革的尼遜.
在《樹與薩姆》以下的第八章.
大衛曾經在所巴這個地方.
擊敗敘利亞人.
其中有一個士兵.
叫做尼遜.
他帶著士兵逃往大馬士革.
自立為王.
似乎大衛是承認了這個國王.
尼遜也是一個非常有能力的人.
敘利亞在他的領導之下不斷增強.
而且尼遜還和以東納文約.
一起來攻打以色列國.
敘利亞在尼遜的帶領之下越來越強.
記不記得在以賽亞先知的時代.
有亞蘭國.
亞蘭國就是敘利亞.
記不記得在以賽亞書的第七章.
曾經尼遜和以色列王一起聯手.
來攻打猶大國.
記不記得.
猶大王當時很害怕.
以賽亞先知就來和猶大王說.
你不用害怕.
你信靠耶和華就可以了.
這兩個只不過是無煙的火把.
所以我們見到.
你真的信靠耶和華.
而不是靠自己的手段.

$^{761}$來希望鞏固你自己的王國.
還有第三個敵人.
就是神興起以法連仁耶羅波安.
哈達是以東.
所以在南面來攻擊以色列.
敘利亞在北面.
所以在北面來攻打以色列.
你說多慘.
南北夾攻.
神還怎樣?.
比以色列還內斂.
耶羅波安就是以色列人民.
是所羅門自己的手下.
亦都是造反.
你想想.
所羅門是有多大的災難.
我們說耶羅波安是所羅門的臣子之一.
他是一個非常有財能.
管理財能的人.
所羅門一早就很賞識這個人.
而因為這個人是以法連支派的人.
所羅門就差派他來管理約瑟家.
約瑟家是什麼?.
就是以法連支派和嗎哪西支派.
記不記得這兩個就是約瑟的兩個兒子.
其實當時整個以色列國已經很厭倦.
所羅門不斷建築的工程.
亦都很厭倦所羅門時時徵召猶太人來做苦工.
所以當耶羅波安發起叛亂的時候.
北國十個支派就讓他來帶領.
仍然保留猶太支派.
而耶穌基督就是透過猶太支派來世上.
我們說神和大衛立的約是無條件.
在列王記上的十一章二十六節.
說還留一個支派給他的兒子.
使我獨人大衛在我所選擇.
立我名的耶路撒冷城旅.
在我面前長有燈光.
這個是無條件的約.
大衛的後人會直到永遠.

$^{801}$直至到耶穌基督來到世上.
所以這個燈會繼續來燃點.
但是雖然神和大衛立的約是無條件.
神和耶羅波安立的約是有條件.
就算耶羅波安能夠拿到十個支派.
都是神給他的.
神和耶羅波安立的約是怎樣?.
是有條件的.
你若聽從神一切所吩咐的 遵行他的道 行他眼中看為正的事 謹守他的律例誡命.
神說我就與你同在 堅固你的家.
有沒有見到?.
你是要守我的律例誡命.
這個就是有條件的約.
耶羅波安的王朝不能夠代替大衛的王朝.
因為神和大衛立了約.
還有一個經文我覺得很有意義.
就是列王記上的十一章三十九節.
神使大衛的後裔受苦難.
但不至於永遠.
有人問不是永遠的話.
什麼時候會沒有苦難?.
大衛的後裔.
有很多學者說將來去到天堂.
那就沒有苦難.
大衛的後裔是不是每個人都信主?.
不是的.
不是每個人都信主.
如果這個應許只是給信主的人.
所有不信主大衛的後裔.
這個應許在他們身上就落空了.
神的應許就落空了.
這個不是去到將來天上.
這個是地上有一個千年的國度.
以致大衛的王朝會再一次在這個世上.
大衛的後裔不會再有患難.
而這個王朝也會一直直至到永遠.
當彌賽亞回來的時候.
以色列分裂的國就會得到醫治.
由彌賽亞由耶穌基督親自來帶領.
他會親自統治整個以色列國.

$^{841}$神最後給所羅門勸告.
這個就是說如果你不聽話的話.
你的王國將會分裂.
我會將你大部分的王國拿走.
賜給你的神父.
這個真是最後的一個勸告.
還有什麼是大過這個?.
除了拿走所羅門的命.
他自己的王國將會被分裂.
如果這件事都不能勸醒所羅門的話.
我相信沒有什麼事是能夠勸醒他.
但是最後所羅門還是沒有醒過來.
他應該撫服在神的面前.
來認罪悔改.
但是他沒有這樣做.
與其是這樣.
他反而走去攻擊耶羅波安.
是想殺死他.
而他也應該要知道.
耶羅波安是得到神的允許.
才能夠拿走他國裡面十個支派.
但是他就想殺死耶羅波安.
耶羅波安又逃到哪裡去?.
又逃到埃及.
你真的看到.
所羅門和埃及的公主聯婚.
就是希望埃及和他聯盟.
結果有沒有用?.
沒有.
以東王哈達又是躲在埃及.
現在耶羅波安又是走到埃及.
來要求庇護.
這個法老王就叫做弒殺.
他根本不買大衛的情.
也不念所羅門和以前的公主曾經立約.
他容許耶羅波安在埃及居住.
直至到所羅門去世為止.
這個就帶領我們去到最後一點.
所羅門與列祖同睡.
在列王記上的十一章四十一節.

$^{881}$所羅門其餘的是凡他所行的.
和他的智慧都寫在所羅門記上.
關於他一切的事.
不單止記錄在列王記.
不單止記錄在歷代志.
也有好幾部失傳了的書.
在歷代志下的九章二十九節說.
包括所羅門記那丹先知書.
還有阿希雅就是剛才撕裂自己衣服的仙姿.
他也有一部預言書.
還有亦多的密士書.
這些全部都已經是失傳.
如果有一天我們能夠找回這些書.
都是好的.
但我們要注意.
就算找回這些書是沒有神的.
密士是沒有被列入聖經的正典.
只不過是一些歷史的記載.
是很好的歷史的資料.
但絕對不是神的密士.
經文說所羅門在耶路撒冷.
以色列眾人作王四十年.
我們剛才說聖經以外都有很多資料.
記不記得我和大家曾經說過.
在埃及王西亞夢.
他在宮殿裡有一幅雕刻的畫.
上高雕刻了他和斐理士人真正的圖面.
記不記得.
這件事就讓我們看到.
西亞夢他和斐理士人來到真賤.
應該是公元前的九百七十至九百六十年間.
而我們亦都知道.
這個正是所羅門登基日子很接近.
因為他娶了西亞夢的女兒.
而西亞夢亦都是將姬息賜給他.
不單止將自己的女兒給了他.
所以所羅門登基的日子.
就是很接近公元前的九百七十年.
這是第一個的日期.
而後來西亞夢視殺這個新的埃及王.

$^{921}$在聖經有記載他攻打耶路撒冷.
是在羅波安在位的第五年.
羅波安就是所羅門的兒子.
即是說所羅門死了五年之後.
這個義煞就攻打耶路撒冷.
這件事應該是公元前的.
九百三十年至九百二十五年間.
如果你是九百二十五年.
你多加五年羅波安在位.
那就是九百三十年.
完全和聖經所講的是一致.
剛才我說所羅門在位多少年?.
四十年.
公元前的九百七十年到九百三十年.
剛剛就是四十年.
而且還有其他俗世歷史的見證.
只是在亞述王薩曼以色三世.
就有兩個記錄.
可以給我們確定所羅門統治的時間.
第一個就是在一個石碑上.
他是記載了在公元前的八百五十三年.
有一個以色列王叫做阿哈.
阿哈這個王應該對我們很熟悉.
是一個非常壞的皇帝.
他參加了夸夸之戰.
就在公元前的八百五十三年.
另外在薩曼以色三世在位第十八年的年史.
記載了耶護在公元前的八百四十一年進貢給他.
這個耶護就是阿哈之後.
散了阿哈王朝.
這個耶護就是自立為王.
所以阿哈之後他有兩個兒子.
在位一共十二年.
然後就到這個耶護.
我們剛才說耶護是多少年.
進貢給薩曼以色三世.
是公元前的八百四十一年.
是當他登基的時候.
他進貢給亞述王.
那就是八百四十一年.

$^{961}$我剛才說阿哈兩個兒子在位多少年?.
十二年.
八百四十一年加十二年是多少?.
八百五十三年.
剛才說阿哈去參加了夸夸之戰.
所以我們知道阿哈王參加了夸夸之戰.
不久就去世.
這兩個兒子成為非常之重要.
是計算聖經自織的歷史.
如果大家記得的話.
我在說列王記的介紹的時候.
我在這個PowerPoint裡面有兩個表.
一個就是記載北國王.
所有他們登基的時候.
做了王多少年.
一個就是南國王登基和做了王多少年.
很多這些歷史都是和續世的歷史是吻合.
所以我們見到聖經非常非常之真實.
你說所羅門死了.
那他留下些什麼?.
留下了他所見的聖殿和皇宮.
還有他著作了《箴言傳道書》和《雅歌》.
他一生追求財富成就.
很可惜為以色列人帶來很沉重的負擔.
他徵稅,勞動人民.
在所羅門在位的時候.
這個國家外表看起來很團結.
但是內裡已經充滿了分歧不滿.
所以去到所羅門去世.
這個國家就分裂了.
所羅門他自己拜偶像就好了.
很可惜他也影響了所有跟著他的王.
他們見到所羅門大衛的兒子都拜偶像了.
我們也可以.
有壞的見證,有壞的榜樣.
而最後累到以色列人被擄.
離開了迦南地應許之地.
就好像回到埃及一樣.
仍然被捆綁.
而列王記就是要告訴以色列人.

$^{1001}$為什麼神對以色列人的審判這麼大.
源頭在哪裡?.
可以說是所羅門.
列王記也是要給正在流亡的以色列人聽.
神沒有忘記你們.
神跟大衛立的約是去到永遠的.
所以你們不用怕.
只要你們願意認罪悔改.
轉回歸向神的話.
神會讓你們回到以色列地.
所以列王記是給正在流亡的以色列人.
一個非常大的鼓勵.
我剛才也說了.
神給所羅門是有他的應許.
精文說什麼?.
我們看看列王記上的三章十四節.
「你若效法你的父親大衛遵行我的道.
遵守我的律例誡命.
我就必使你長壽」.
我們知道所羅門統治了以色列四十年.
你說他什麼時候在位?.
有些蛛絲馬跡.
在列王記上的三章七節.
所羅門叫自己做什麼?.
其實很多學者說不是的.
他只不過是自己謙卑.
又注意到在列王記上的二章九節.
已經叫他做聰明人.
不是一個弟弟.
所以有很多學者說.
所羅門登基的日期可能是他二十歲的時候.
他在位四十年是什麼時候去世?.
六十歲去世.
很年輕.
我想我們這裡大部分人都年長過六十歲.
有幾個就不是.
相對來說大衛是怎樣?.
在列王記上的第一章.
大衛臨死的時候年紀很大.
甚至不能夠保暖.

$^{1041}$可能古代沒有暖氣.
要用人造的暖氣是誰?.
就是亞比薩.
找一個童女回來.
讓大衛抱著睡覺.
這樣就可以保暖.
所以見到他的身體真的很差.
所以大衛的年紀很大.
所羅門很年輕.
有可能是因為他沒有遵守神的律例典章.
所以他不是那麼長命.
在結束的時候.
和大家一起來看看.
智慧等不等於成功?.
神給所羅門智慧不代表什麼.
我們要懂得運用這些智慧才行.
不是要用我們自己的想法.
很多時候基督徒都是一樣.
我們不是不知道什麼是對的.
聖靈也是會提醒.
但很多時候就是質疑自己想做什麼就去做什麼.
這個就是問題.
當我們讀傳道書的時候.
我們見到所羅門真的有一段時間.
他是離棄了神.
他是奮起疾逐.
他甚至是絕望.
懷疑自己的人生是不是有意思的呢?.
所以他不斷地說虛空.
然後他追求生活裡面的享樂.
他去追求建築為了什麼?.
他就希望在這些物質上的東西.
能夠使自己快樂.
但最後他快不快樂?.
一點都不快樂.
在傳道書裡面超過38次.
他說虛空的虛空.
我們也看到.
當我們屬靈的生命.
是被物質的東西代替了之後.

$^{1081}$我們的心就會慢慢慢慢地離棄神.
這個就是我們今天的主題.
我們的心離開神一寸.
我們的生命離開神一里.
所以我們真的要很小心.
我們要保守我們的心懷意志.
在聖經裡面也會有警告我們.
如果我們離開神一寸.
首先就會與世俗為友.
雅各書的四章四節.
然後沾染世俗.
雅各書的一章二十七節.
然後就會愛這個世界.
約翰一書的二章十五至十七節.
最後會效法這個世界.
羅馬書的十二章二節.
很不幸的.
如果我們這樣做的時候.
可能會與世人一起被證罪.
或失去一切.
我不是說失去我們的救贖.
一個真正相信的人.
不會失去救贖.
這個只有我們自己之神也知道.
我們是否真的走在神的旨意.
我們是否真的繼續追求神.
認識神.
如果是的話.
你是一個真正得救的人.
但如果你說你每個星期日.
你回來聽到.
你根本自己不會追求的話.
你未必是一個真正得救的人.
但就算是一個真正得救的人.
也可能像羅德一樣失去一切.
記不記得在創世記的十九章.
當神審判所多瑪俄摩拉的時候.
他是與所有那些不信的人.
一起來被審判.
你說羅德.

$^{1121}$這個人是不得救的.
可能是一個不得救的人.
不是,聖經叫他做二人.
在比特後書的二章七節.
說羅德是一個二人.
你說所羅門是否得救.
學者都爭論不休.
很多不同的意見.
但我相信聖經裡面.
既然有三部書都是所羅門所寫.
《雅歌》,《傳道書》和《箴言》.
我相信所羅門是一個真正得救的人.
而在《傳道書》的十二章.
十三至十四節.
就是《傳道書》結尾的時候.
所羅門怎樣說.
他勸人要敬畏耶和華.
他追求了俗世這麼多的東西.
樣樣都虛空.
然後他最後怎樣說.
要敬畏神.
人所做的一切.
就算是隱瞞的事,隱藏的事.
無論善惡神必審問.
這裡就讓我見到.
所羅門真的有個知罪的心.
所以我相信他最後.
是有回到神那裡.
最後是有接受神.
但他一生就好像羅德一樣.
是一個追求俗世的基督徒.
你不要以為很過癮.
我都要像羅德,所羅門那樣玩盡.
可能你不是得救那個.
你不要說你真的盡量放縱你自己.
這裡也要做一個預告.
我很快就會講哥林多前書.
哥林多前書的十章.
就是講到不要做一個放縱的基督徒.
我們下一年會講哥林多前書的第十章.

$^{1161}$願意我們都能夠立定心志.
繼續來效忠耶穌基督.
耶穌基督就是比所羅門更大的那一位.
耶穌基督就是為我們死.
為我們復活.
也是為我們將來回來迎接我們.
帶領我們一起去到新天新地.
願神祝福他自己的話語.
讓我們低頭禱告.
讓我們見到所羅門最後的生平.
我們見到他慢慢地離棄了神.
他的心離開神一寸.
他的一生就離開神一里.
以至最後他使到以色列國分裂.
雖然這件事沒有在他在生的時候發生.
這個全靠是你的恩典.
願意我們都不要濫用你的恩典.
願意我們都能夠有這個心.
繼續持守我們的信仰.
繼續效忠於耶穌基督.
因為你就是比所羅門更大的那一位.
我們知道你為我們死,為我們復活.
為我們將來會回來.
帶領我們是德性進入你的國度.
願意在你回來的時候.
我們都能夠被你稱為一個忠心良善的僕人.
將那榮耀還賺得貴給你.
我們這樣禱告.
奉主耶穌基督的名而求.
阿們.
\newpage



\section{士師記 16:26}
\label{sec:032440}
\textbf{聖經人聖經事:我要靠一靠}~~~~黃克勤牧師博士
\newline
\newline
連結: \href{http://www.church.com.hk/acms/content.asp?site=cdc\&op=show\&type=product\&code=032440\&layout=sermon}{\texttt{032440}} ~~~~ 語音日期: 2026-07-10
\newline
\newline
\hyperref[sec:032439]{\small{< < < PREV SERMON < < <}}
~
\hyperref[sec:index_chronic]{\small{[返順時目]}}
~
\hyperref[sec:032441]{\small{> > > NEXT SERMON > > >}}
\newline
\newline
士師記 16:26
\newline
\begin{longtable}{cl}
\hline
\hline
章節 & 經文 (和合本修訂版)\\
\hline
16:26 & \begin{tabularx}{0.7\textwidth}{X} 參孫對牽他手的僮僕說：「讓我摸摸支撐這廟宇的柱子，我要靠一靠。」 \end{tabularx} \\ \\
[1ex]
\hline
\hline
\end{longtable}
$^{1}$大家好,今天題目是「我要靠一靠」.
靠一靠不是靠人,是靠一根柱子.
要靠一靠,為何要靠一靠?.
因為累,辛苦,索氣.
有個好笑的事,很多非洲人喜歡這樣靠一靠.
經常靠一靠,好像很沒神氣,為何呢?.
原因是有一段時間,非洲有很多地方是饑荒的.
很艱難的,人們瘦得像棒骨,為何經常靠一靠?.
無力,食不足,無力,所以經常靠一靠.
我有一段時間身體不好,都靠一靠.
我發覺自己站起來要扶一扶.
上樓梯要扶手,扶手覺得有點骯髒.
一見到有椅子,如果沒有人爭,我都先坐一坐.
雖然坐一坐又起身,但坐一坐都好,為何?.
無力,這個無力會否是參酸呢?.
參酸不會,恐無有力,不過這是參酸的真實說話.
我要靠一靠,我要挨一挨.
《生命說》和大家繼續看聖經人,聖經事.
看看上主的工作.
《禪師經》疑言十六章.
非利士人的首領聚集,經常這樣說,薩姆爾寫出來.
有幾個首領沒有說,其實不緊要,薩姆爾都知道這件事.
非利士人的首領聚集,第一件要留意的是.
要給他們的神大菌獻大祭,留意讀字.
要給他們的神,即是非利士的神.
非利士的神叫大菌,讀菌,病菌的菌.
大菌,不要讀大菌,大菌是什麼呢?.
是一個男性的愚昧人生,男性.
很大隻,恐無有力,大菌.
因此推測非利士是由海上上岸登陸的民族.
要獻不是普通的祭,要獻個大祭,因為很開心.
因為我們的神大菌將我們的仇敵參孫交在我們手中了.
這麼大力的參孫都被我們贏了.
其實他們當時以往的人都非常spiritualized.
即是將他們日常生活很多東西都非常迷信地歸因於靈界的工作.
我們能夠打贏是因為神幫助我們.
我們打輸是因為神不喜歡我們.
我們打輸是因為他的神比我們的神厲害.
他們經常想,今天不精神,因為我的神沒有給我精神.
因為今天不精神,是因為我的神不喜歡我.

$^{41}$他們要找一個辦法去討好神.
我在印度見到很多,他們什麼都是神今天對我好或不好.
這個神作弄我,我要找一個大力的神去壓這個神.
很開心,他們要獻祭,歡樂,讚美他們的神.
原來這個參孫的被捉失敗是令到外邦的神被榮耀.
這真是很不值.
本來應該是我們的神得到榮耀.
但現在外邦的假神反而要接受大祭禮慶祝.
這真是很不好.
眾人看見參孫,就讚美他們的神.
你看,我們的神將毀壞我們的地殺害我們許多人的仇.
你搞在我們手中,見到參孫就讚美他們的神.
他們獻祭,獻大禮,參孫反而令到大菌神被讚美.
真是很諷刺,也很可悲.
他們正燕落的時候,就說不如叫參孫出來在我們面前戲耍戲耍.
戲耍戲耍即是甚麼?我們捉他出來看看他,玩他,整他,當猴子般搞他,玩他.
不如捉參孫出來.
於是就將參孫從監裡提出來,在眾人面前戲耍他.
湯姆斯特站在兩處中間,就戲耍他.
然後他瞎了,其實也挺慘,他瞎了,你想想當時有沒有眼藥,有沒有眼科.
醫生專科沒有,有沒有手術刀,有沒有毛菌技術來拿他的眼球呢?.
沒有,是用刀撩他的眼出來,又不流血,流血了.
有沒有藥,即是跟你做手術,連針,然後包住他毛菌,沒有.
所以我想參孫為他的眼痛,已經痛了一段很長的時間.
說不定他現在不是淚流滿面,淚腺可能沒有了.
而是兩隻眼凹陷了,眼腔凹陷了,裡面沒有眼球.
於是有些血流出來,流了又有膿,甚至有蒼蠅嬲.
很痛苦,很多東西流出來,發炎中,很嚴重.
又血又水又膿,又痛,不得了,又內心更加痛苦.
這就是參孫失敗之後的經歷,現在還要推磨.
推磨我看不到,不用看,你這樣向前推,不行我就鞭你.
向前推,向前推,磨呀磨呀磨,又痛,但繼續要做.
現在我大力不發,他沒有力氣了,變成一個普通人,於是很辛苦.
我們想像參孫的時候,內心的痛苦是怎樣,你可以想像到.
現在還被人抓出來,參孫,一巴掌打過去.
誰打你?告訴我,你的神不是很厲害,告訴我.
被人戲弄,或者找人潑他,或者找狗咬他,哎呀哎呀,死了.
你見到都很害怕,戲弄他,於是將他擺在兩柱中間.
這裡厲害了,兩柱中間是甚麼呢?.
就是在整個大菌的主苗裡面,有兩條柱頂住.

$^{81}$最重要的那兩條柱,那兩條柱通常是中心點.
那兩條柱中心點,擺在那裡,中心點是甚麼意思?.
個個都看到,哈哈哈哈,參孫你還不死這次.
那就活該了,跟我打個兩巴掌,鞭他找石釘.
這樣就很慘了,掃他,哎呀,很慘.
參孫向拉他手的童子說,求你讓我摸著托房子的柱子.
我要靠一靠,哎呀,你讓我靠一靠柱子.
不過我覺得上帝仍然憐憫參孫,是拉他手的童子.
我看到這裡心裡有少少感動與安慰是甚麼呢?.
有一個童子,不是少年人,可能是少年人,十多歲最大.
他做甚麼呢?他就是被派去看守著參孫.
他推木,推那條大木,然後中間有個輪來搖爛橄欖.
那條木在哪裡?這裡有人搭他手上去,搭他手上去.
這裡推吧,腳就損了,有個橄欖有條銅鏈.
有條龍鏈鎖住,手損了,腳損了,眼損了,痛得不得了.
渾身被人打到瘀,在推.
其實這個童子也很仁,每天看著他,有些童子很壞的.
會不會先整你,最慘的可能是這個童子.
一腳兜他飯,他沒得吃,給他污糟水喝.
明晚跟他做壞事,他也沒得動,踢他也可以.
但這個童子挺好,拉他手的童子.
如果你的童子會不會拉他的手,嫌骯髒.
就找條棍,這邊這邊,好像牛一樣趕,或者拉他條鏈還好一點.
他是拉他手的童子,即是說這個童子對他有憐憫的心.
也是他的朋友,可能大家聊天時是這個,在監獄裡.
他就跟這個童子說,求你讓我摸著托房子的柱子.
為什麼呢?我想靠一靠,很辛苦,我想靠一靠.
第三點我要跟大家說,我要靠一靠.
這個正是你永遠想不到這麼大力的參選會說,我要靠一靠.
他能夠將希伯倫,不是,參選能夠將加薩的城門.
連框大門,還塞住他,整個門這樣搖了下來,托在單手.
這扇門,城門很大,可能是二三十尺的一扇門,尺.
托住他,托住去到這麼遠的地圖,加薩走到希伯倫.
然後扔希伯倫前面一個山的山頂,還不是平地.
要走上山,扔在山頂,讓你知道我有多大力.
單手面色不轉,手都沒有轉,膊都沒有轉.
馬上去到山頂,加薩希伯倫,都很怕,都走一天路.
不說笑,就是這樣,要不要靠一靠?.
「哎呀,哥哥仔,帶我去一條柱子,一條大一點的,頂住這間屋的,我想靠一靠」.
這個小朋友不如有責,靠一靠,為甚麼要找最大的一條?.

$^{121}$他沒有問他,這條其實已經站在兩條柱子中間,就那兩條柱子.
他就說「我要靠一靠,其實如果我們一早認識,神呀,讓我靠一靠」.
「讓我靠一靠到你那裡,哎,這樣就好了」.
可惜他靠在大利拉的膝蓋那裡,而沒有靠在自己頭髮的能力那裡.
他沒有靠在神那裡,這個時間我要靠一靠.
「啊,懂了,靠在這裡,熟悉靠在神那裡,最好了」.
那時房裡充滿男女,非利士人的眾首領都在那裡.
哇,這個原來是大劇院,下面有一條柱頂住,上面坐了很多人.
可能大家可以想像,譬如我們坐在這裡.
你們看這裡是一個寬闊的露天的牆.
我們這裡是一個平頂,高了一層,有柱頂住,我們坐在上面.
亦在這裡開圓遊會,很開心.
那就可以看下去,哈哈哈哈,你站在那裡,參算站在下面.
我站在上面看,參算死了沒有?.
他走了過來我們下面,童子說「我要靠一靠,你帶我過去」.
帶了他到我們下面,我們有三千男女在上面玩耍.
這個時候「我要靠一靠」,就是參算講的一句重要的說話.
非常重要,我要靠一靠,很辛苦.
如果你要靠一靠,記得去到神面前,不是找一條柱來靠.
不是靠自己強而有力的頭髮與膀胱那裡,而是靠神那裡.
我要靠一靠,放下以前的自負,放下以前的自我.
甚至請一個看來不夠我這麼厲害,差很遠的童子.
你帶我去,我看不到,你帶我去柱那裡,我要靠一靠.
靠一靠在神的身上,比靠在任何地方都重要,祈禱.
主啊,我們今天到你面前,都需要找一個地方去叫一叫.
發覺我們已經軟弱無力,不做到什麼了.
這個時間眼睛看得不清楚,很多問題發生在自己身上.
很多難題去到很極端,解決不了.
已經是經歷了這些困難,痛苦,重擔,疾病很長時間.
主啊,我已經很累了,不受得了.
主啊,這個時間你叫我們,你要靠一靠.
幫助我們記起,靠在柱那裡找一個有錢人靠一靠,沒什麼用.
找一個名醫靠一靠,沒什麼用.
找一條柱靠一靠,沒什麼用,找一些錢靠一靠,都沒什麼用.
最重要是在神面前,主啊,讓我在你那裡靠一靠.
叫一叫,叫我們重新有力.
求主祝福,好過奉耶穌命求,阿們..
\newpage



\section{以弗所書 2:14}
\label{sec:032441}
\textbf{拆毀與更新}~~~~陳淳牧師
\newline
\newline
連結: \href{http://www.church.com.hk/acms/content.asp?site=cdc\&op=show\&type=product\&code=032441\&layout=sermon}{\texttt{032441}} ~~~~ 語音日期: 2026-07-12
\newline
\newline
\hyperref[sec:032440]{\small{< < < PREV SERMON < < <}}
~
\hyperref[sec:index_chronic]{\small{[返順時目]}}
~
\hyperref[sec:032442]{\small{> > > NEXT SERMON > > >}}
\newline
\newline
以弗所書 2:14
\newline
\begin{longtable}{cl}
\hline
\hline
章節 & 經文 (和合本修訂版)\\
\hline
2:14 & \begin{tabularx}{0.7\textwidth}{X} 因為他自己是我們的和平，使雙方合而為一，拆毀了中間隔絕的牆，而且以自己的身體終止了冤仇， \end{tabularx} \\ \\
[1ex]
\hline
\hline
\end{longtable}
$^{1}$這段經文的十一節到二十二節.
其實是處理猶太人和外邦人之間.
長期敵對,隔閡的原因.
我選擇的第十四節.
我沒可能十一節到二十二節說.
我猜大家應該要留很久了.
所以我只是按照我的習慣.
只選擇一節的經文.
這一節的經文.
其實是這段經文裡面的靈魂中樞.
如果我們說這所書的核心主題.
是萬物在基督裡面同歸於一.
那麼二章第十四節.
其實它就不只是一個歷史記錄.
它是一個很關鍵的切入點.
所以我將它的題目定在.
拆毀與重塑.
這一段的經文.
它刻畫,勾畫出一個偉大的藍圖.
在基督的國度裡面.
人與人之間不再需要彼此防備.
不同民族之間無需劃分敵我.
教會的意義.
就是讓這份合一真正的實現.
福音的影響力.
不只是打破了,消離了神與人之間的鴻溝.
讓人個人得救.
救恩的延伸.
是能夠延伸人際關係的重建.
讓我們這群本來被牆所隔開的人.
能夠在基督裡面.
因著愛,因著基督的大能.
能夠有真正的和好.
真正的復和.
真正的平安.
真正的和平.
在講到之前.
讓我們一起同心同意.
親愛的阿爸父.
我們來到你面前.

$^{41}$求你以聖靈親自的運行.
拆毀我們心中自以為是的高牆.
粉碎我們為了防衛自己而築起的警戒線.
願你的話語在我們裡面動工.
重塑我們的生命.
優先次序.
禱告乃奉我主耶穌基督的聖名而求.
阿們.
保羅說.
因基督使我們和睦.
將兩下合而為一.
拆毀了中間隔斷的牆.
提起牆.
我想我們大家熟悉的.
我們現在就坐在四面牆的中間.
牆幫助了我們.
和外面炎熱的天氣隔斷了.
牆保護了我們很多的家居.
甚至我們現在很熟悉的.
網絡的防火牆.
它可以過濾惡意的侵害.
心理學上的界線.
其實是能夠守護我們的情緒健康.
防止過度的消耗.
或者被情緒勒索.
在教會裡面.
我們會築起一個一個的信仰的邊界.
我們用教義.
用傳統.
用約定俗成.
劃出一塊又一塊的安全區域.
原意就是保護我們實踐信仰的純正.
讓我們在混亂分裂的世界裡面.
我們能夠在世而不屬世.
合乎中道的.
持守聖經真理.
持守那份聖潔的真愛.
然而.
這一切所有為了保護而設的屏障.
如果被過度的濫用.

$^{81}$它會靜悄悄的變質.
當我們習慣依賴牆的格阻.
來帶來安全感.
原本用來保護的界線.
牆就會很容易變成一個區分敵我的鐵籠.
因此.
保羅在這裡.
這段經文.
他要挑戰這種變質的囚禁.
變質的網絡.
他說.
人基督已經使我們和睦.
將兩下合而為一.
拆毀了中間間斷的牆.
中間隔斷的牆.
這個隔斷的牆.
原文是在說隔斷和尼巴.
中文四個字.
但是原文只有一個字.
而這個字.
不是比喻.
不是形容.
保羅寫這個字的時候.
他很可能在想耶路撒冷聖殿裡面.
外邦人院和聖殿內院.
中間那裡實際存在的實體圍牆.
石欄.
那裡是一道區分內外.
嚴禁外邦人越界的冰冷石欄.
這個石欄.
這個牆.
猶太人認為是守護他們的聖潔.
區分上帝子民和外邦異類的分隔牆.
所以在這幅牆上.
每隔一段距離.
就有一個牌寫著.
嚴禁外邦人進入聖殿周圍的尼巴和圍欄之內.
萬被捉拿者.
其後果自負.
死罪難逃.

$^{121}$你見過可以讀.
當然他不會用中文寫.
但是現在有些博物館就要收藏這些牌.
這道牆背後.
其實隱藏著猶太人一段極為沉重的創傷.
如果你看《歷代志》.
看《聶旺記》.
看《先知書》.
你就見到.
他們曾經因為信仰上的妥協和混雜.
他們因此而失去家園.
失去國家.
最後.
在慘痛經驗之後.
在他們回歸重建聖殿的時候.
他們就築起一座防線.
試圖嚴格遵守對神的忠貞.
很可惜.
這份原本出於敬畏的舉動.
最終卻讓他們將這個.
本來應該是萬國禱告的殿的入口.
就變成了一個嚴苛的排他防線.
弟兄姊妹 當你聽到這些教導的時候.
或者這一個聖經的記錄的時候.
你心裡面有什麼感覺呢.
其實這不是我們內心的寫照嗎.
這豈不是我們內心常常經歷的事情.
當我們每次忐忑不安.
當我們見到和我們不同.
和我們有差異的時候.
我們很自然.
很不自覺.
我們在面上就會掛起了這座石碑.
對於那些挑戰我們立場.
挑戰我們文化.
挑戰我們生活習慣的人.
我們很容易就在面上寫著.
不可靠近.
否則後果自負.
保羅透過這篇經文想告訴我們.

$^{161}$那個曾經死罪難逃的律法規條.
那座隔斷了的圍欄.
已經在基督十字架下.
渠然崩塌.
永遠抹去.
請留意.
這個不是一場拆毀舊建築物.
或者舊有結構的行動.
它是一場拆毀重塑的工程.
當然當我們這樣說的時候.
我們雖然是需要小心謹慎.
因為這件不是別人的事來的.
這不是其他人要做的事.
這是你和我都要共同面對.
互相勉勵的路來的.
因為當我們這樣看經文的時候.
我們很容易就看見自己的影子.
其實你和我.
我們很多人.
都是一個又一個的築牆人.
所以當我們看這篇經文的時候.
今天我宣講這篇聖經的話語的時候.
我們不是用來自責.
或者不是用來指責人的鞭子.
它是成為一面鏡.
照出我們的本相.
看見我們的軟弱.
看見我們的掙扎.
然後我們藉著保羅寫在聖經裡面的思路.
彼此警醒.
我們看看耶穌基督.
祂怎樣親自拆毀那種人與人之間的隔膜.
帶領我們進入祂所賜的真自由.
保羅說.
因基督使我們和睦.
將兩下合而為一.
拆毀了中間隔斷的牆.
拆毀與重塑.
是從焦點的跨越開始.
從自我的防禦.

$^{201}$去到轉眼仰望耶穌基督.
中文的聖經或者原文的聖經.
開首那兩個字.
因他.
因他這兩個字很容易你就會看到是因果關係.
是不是?.
因為是印.
但事實上原文不是這樣的.
希臘文原文是.
他自己就是.
保羅用語非常有震撼力.
他不是說因為耶穌基督帶來了和平.
他是說耶穌基督就是我們的和平.
這種和平.
這個用詞.
不是一般的紫色紛爭.
不是一般說平息糾紛.
不是靜靜雞.
他是承載希伯來文語系裡面.
平安.
Salon.
那個用詞.
那個豐厚的意義.
我曾經有機會在這裡用一堂的時間宣講.
所以我沒有辦法在這裡再解釋這個字.
否則的話我想我們就兩點鐘都走不到.
因此.
如果你有興趣你可以回到城鎮港島的資料庫裡面.
我再一次我們要說.
這種Salon.
是在說.
完好無缺.
整體的生命秩序.
這種平安.
是極具顛簸的宣告.
耶穌基督.
就是那份完全整全的和平.
整全的平安.
所以.
一切沒有和耶穌基督連結的人.

$^{241}$我們所追求.
所保護的所有的安穩.
都只是脆弱的假象.
只是用尼巴堆積起來.
薄弱的隔離區.
保羅宣告.
耶穌基督.
就是我們的和平.
其實他是直接.
否定了我們過去透過用梳理.
用分隔.
用防禦.
各種方法保護自己的方式.
所以.
我們要學會.
將焦點跨越.
從自我防禦的視角.
轉向仰望基督的和平.
基督的平安.
不過當你這樣說的時候.
陳牧師你說得很口響.
我知道.
但怎麼做呀.
或者說.
其實.
你難道.
拆了自己的牆.
你沒有心中不安的嗎.
是的 我們需要承認.
保護牆.
令我們覺得.
好像有點安全感.
牆背後.
總是有著你和我種種的不安.
種種忐忑.
種種對一些未知處境.
而帶來的徬徨.
但是福音的力量就在這裡.
基督親自成為我們的和平.
我們的平安.

$^{281}$他承擔了我們的焦慮.
所以.
我們應該可以靠著.
基督的大能.
基督的大愛.
學會.
不是靠一面牆.
不是靠一些分隔.
不是靠一些保護.
來保護令我們覺得有安全.
因為我們每一次的保護.
每一次的劃分.
每一次找一面牆困著我們.
我們其實都活在孤單和困擾當中.
我們總是覺得自己孤軍作戰.
所以耶穌基督要我們出來.
釋放了我們.
因為基督是用祂的寶血.
用祂的生命.
接納了最真實的你.
最真實的我.
最真實的祂.
唯有我們確認.
我們每一個人.
都是被基督無條件的接納.
我們才不再需要.
用保護牆.
用分隔線.
來保護自己.
來證明自己.
我們裡面那份忐忑不安.
因為基督的大愛.
因為基督的大能.
我們就能夠慢慢安靜下來.
我們不需要再用我們的躁動.
去作出反應.
應急.
應激.
保羅說.
因祂使我們和睦.

$^{321}$將兩下合而為一.
拆毀了中間隔斷的牆.
所以第二個.
其實跨越.
是轉化的跨越.
從剛剛我們所說的.
恐懼隔離.
轉化為.
甘願被拆毀.
第十四節.
裡面最主要的動詞.
就是拆毀.
不過你知道.
這個拆毀.
它在用什麼事態嗎.
它在用的是過去分詞.
所以意味著.
這場拆毀的行動.
並不是一個建議.
也不是一個今天你要做的事.
也不是明天你要立下心.
準備拆了.
過去分詞.
說的是.
它在基督十字架下.
已經被基督徹底完成.
圓滿的結構.
這個字.
不是一個簡單的動作.
不是推倒一段牆.
它是徹底的.
打破框架的能力.
它極具破壞力.
它不單止推倒這堵牆.
它更加重要的.
是連根拔起.
我們裡面那份平安.
安全的習慣.
我們時常見到有一些不妥.
見到一些不安.

$^{361}$見到一些不同.
見到一些差異.
見到一些步伐.
和我們不一致的時候.
我們就疏離.
就退開.
就隔斷它.
就劃清界線.
就走開.
我們這些壞習慣.
這些舊習慣.
其實是冷漠對立的思維造成.
耶穌基督的大愛.
耶穌基督的大能.
就是要徹底粉碎它.
我們回到剛才說.
耶路撒冷那幅中間隔斷的牆.
其實它不是一個建築物那麼簡單.
它其實是一個神學上的傲慢.
他告訴外邦人.
猶太人告訴外邦人.
你不配.
你不配去到上帝面前.
這幅牆.
其實令兩個群體陷入雙重的災難.
他拒絕了外邦人.
排斥了外邦人.
令外邦人失去了認識真神.
悔改歸向真神的機會.
同樣.
他又令築牆者,建牆的人.
猶太人陷入了僵化,恣意,自高自大.
失去了他們當日領受上帝使命的要求.
頂智慧.
那幅牆.
耶穌基督已經拆了.
不過在二千年來.
這幅牆存在.
在你和我的心中.
我們很多時候.

$^{401}$因為觀點的差異.
因為步伐的不一致.
因為方法的不同.
我們選擇疏離.
選擇推開.
不是他推開它.
就是自己走.
其實我們無意之間.
又捉起那座防衛疏離的牆.
或者.
你不會像猶太人那樣對外邦人說.
你不配.
你不敢這樣說的.
我知道.
我們比較謙虛.
不過我們用腳投票.
我們用心靈疏離投票.
我們離開那些影響我們的人.
離開和我們所有不同的人.
其實不也是猶太人推開外邦人一樣.
所以保羅在這裡提醒我們.
十字架大能的重點.
就是在於這幅牆已經摧毀.
已經拆毀.
我們不應該選擇疏離.
我們不應該因為不同.
因為差異.
因為步伐不一致.
因為大家觀點和角度的不同.
所以.
我們選擇起牆.
劃清界線.
你做你的 我做我的.
弟兄姊妹.
因為基督的愛.
是呼籲我們要跨越.
以至進入彼此.
跨越那道無形卻實質的牆.
如果我們要經驗耶穌基督的拆毀.
將人與人之間 文化與文化之間.

$^{441}$觀點與觀點之間.
那種拆毀的大能.
我需要讓我們弟兄姊妹清楚一件事.
我們其實是需要有意志的轉化.
意志的轉化.
當然並不是否定心理學上健康的心理界線.
我們再說.
適度的界線其實是守護的必需.
適度的界線是上帝賦予的禮物.
保羅在這裡要挑戰的.
是那種因為恐懼.
因為高傲.
因為疏離.
因為怯懦.
而產生的防衛機制.
弟兄姊妹.
讓我們一起學習.
放慢我們的腳步.
當我們審查內心.
有一種想退縮.
真的麻煩了.
想走開.
覺得都不是我的想法.
想隔斷它.
想捉長.
在這個時候.
每次你的防衛線想起的時候.
你的防衛感想起的時候.
請你請我.
在主的面前停下來.
停下來溫柔的問一下自己.
主啊.
這是你要守的分別衛性的界線.
抑或這是為了證明自己.
守護自己而起的疏離防線呢.
讓我們學會擁抱這種脆弱.
擁抱這種忐忑.
而不是選擇防衛.
仰望耶穌基督的接納.
將那份不安.

$^{481}$將那份徬徨.
將那份未知.
交給上帝.
因為拆毀這個字.
原文就是廢除.
原文就是取消約束.
在基督裡面.
在基督的大愛裡面.
他廢除了那套以律法排他的邏輯.
他用他的寶血.
為那些常常定義.
誰是聖潔 誰是罪人的標準.
徹底的抹去.
你和我在基督的寶血面前.
不是誰犯罪大 誰犯罪小.
我們都欠了上帝的愛.
我們都欠了基督的救贖.
所以在基督裡面.
我們都是蒙恩的.
我請你填充題行不行.
你是蒙恩的.
我知道你一定會寫罪人.
我想和你說.
我今天想提醒你.
你靠著耶穌基督的寶血.
已經成為一個蒙恩的新人.
我們都被基督所救贖.
我們都被基督所拯救.
成為新造的人.
舊事已過.
下一句悉背嗎.
這個是對我們教會的生活.
對我們信仰的生活一個極大的挑戰.
幫助我們跳出我們的安書區.
幫助我們跳出我們的習慣.
我們不止只是和我們志同道合的人聊天.
我們不止只是和我們.
那些想法相同一致的人共同工作.
我們在團契裡面.
是要學習更寬闊更寬廣的團契.

$^{521}$去聆聽那些不同的觀點.
就算他和我們不盡相同.
天南地北.
大能且不邁.
甚至他會挑戰你和我.
一些行事為人的準則.
但是這個豈不是罪在主的家中.
我們應有的反應嗎.
所以我們每一次.
每一步跨出去的時候.
其實這個就是我們見證著.
基督拆毀的大能.
坦白說.
不容易的.
因為你和我很清楚.
每次遇到那些.
我們很重視.
但是又和我們很不同的人.
那些人不是很遠的.
每個家裡都會有一個或兩個.
不是夫妻.
就是父子子女.
每一個人.
不用在教會找的.
你家裡就有些人.
柴米油鹽醬醋茶吵一頓.
就為了一些雞毛蒜皮的事.
不知道為什麼.
有說或者不說.
本來無所謂的.
但就是因為你不說出事.
這一切的所有的不同.
讓我們經歷很多的孤單.
讓我們經歷了很多的圍牆.
經歷很多的違憲.
但是請記得.
你和我不是單打獨鬥.
你和我是在那位已經拆毀.
猶太人和外邦人石牆的耶穌.
在填補了神與人之間來鴻溝的基督.

$^{561}$我們在祂的帶領下.
在祂的幫助下.
親自引領我們成為新造的人.
邁向真正的合一.
這份轉化對我們每一個人來說.
每一個人都是一種社稷.
當我們願意參與社稷.
背十字架.
和耶穌基督一起拆牆.
你將會感受到前所未有的.
不是堅強.
是脆弱.
因為你赤裸躺開你的自我.
你沒有防禦的赤裸躺開.
不過正正是這份脆弱.
你才能夠真正體驗.
基督是我們的和平.
基督是我們的平安.
這份平安不是來自你將其他人推開.
不是你將其他人擋在牆外.
不是你自己走開.
這份平安是來自你這樣基督的愛.
流進你的內心.
流進你所愛的人.
流進你覺得他麻煩的人.
這份愛貫穿於我們每一個不同的個體.
讓我們每一個人的生命裡面.
都產生一個內在的革命.
我們是慣常的.
為了維護關係.
所以你好我好.
大家相見搖手不作長談.
為了保持關係.
不要說那麼多.
不要見那麼多.
其實這是我們的藉口.
是掩飾我們內心深處的排他性.
高傲性.
怯懦性.
或者習以為常的孤單.

$^{601}$弟兄姊妹.
保羅告訴我們.
如果你想成為基督的見證.
你就不要再捉牆.
你就不要再和人之間劃很多保護的界線.
我再說健康的界線是需要存在的.
但是在基督裡面.
在我們生命當中.
在基督的大愛底下.
我們曾經堅韌強硬的偏見.
成見.
曾經冷漠的隔閡.
都因為基督的大能化為烏有.
我們要學會.
要看見.
這就是福音的震撼力.
如果耶穌基督在兩千年前.
祂能夠拆掉猶太人和外邦人.
那道路不可破的鴻溝.
祂怎會拆掉不了.
你和我心中那種對弟兄姊妹的隔閡冷漠呢.
不過請你記得.
這是福音底下.
今天當不信耶穌基督的人.
那兩族到現在還在打仗.
是不是.
唯有在基督裡面.
我們可以抹去這些仇恨.
在基督裡面.
我們才知道.
拆毀已成事實.
所以你和我回到家.
你和我看看旁邊左右.
那個既親近.
但同樣又有吉下你的人.
讓我們記得.
你和我中間的圍牆已經拆下.
因為基督的怨顧.
你和我可惜彼此接納.
雖有忐忑不安.

$^{641}$因為你和我真是很不同.
步伐真是很不一致.
但因為基督的愛.
我們可以學會接納對方.
共同交融.
保羅說.
因基督使我們和睦.
將兩下合而為一.
拆毀了中間隔斷的牆.
拆毀與重塑.
最後的跨越.
是使命的跨越.
是從各自孤立.
去到主內的合一.
經文說將兩者合而為一.
這個經文.
在一個原文裡面是一個動詞的短語.
但它並不是簡單的一加一等於二的加法.
保羅的用字不是這樣.
他本來可以選擇用合併.
合併就是一加一等於二.
他的動詞用的是製造.
什麼意思.
他在說的是一種本質上的徹底融合.
這個其實是一個極具神學張力的用詞.
他在宣告一種不會逆轉的事實.
兩群本來在血緣,文化,宗教習慣上.
勢不兩立的宗族.
在基督裡面不再是兩群人.
而變成一群新造的人.
所以沒有基督,仗繼續打.
有基督,我可以活在主的家當中.
這種合而為一.
並不是妥協式的共存.
這種合而為一.
不是為了和平而犧牲真理的苟且.
這種合而為一.
讓我們成為新造的人.
不用很遠.
你打開聖經,看回十四節下面.

$^{681}$第十五節.
保羅就是這樣說.
造成一個新人.
這是基督裡面救贖工作的最高使命.
也是傑出的果子.
他不只是來修補兩道牆之間的裂縫.
他要徹底將兩堆頑固的石頭.
將兩個心中頑固的人徹底用化.
重新鑄造成為一個全新前所未有的新人.
舊事已過,都變成新的了.
真正的合一是以基督為中心的合一.
基督是我們唯一的標準.
不過我們需要承認.
我們其實有很多人不喜歡合一.
或者我們很多人想的合一是你來合我的一.
我們總是害怕失去自我.
我們總是害怕欠缺自己的獨特性.
我看到你看不到.
我做到你做不到.
我有這些經歷你又沒有.
我們這種各種各樣的獨特性.
害死了我們.
但願讓我們記得.
唯有在基督裡面.
我們才敢於放棄各種區分敵我的邏輯.
我們的差異是事實.
我們的不同是真實的.
你有經歷過我沒有經歷過.
這個是事實.
但這一切無需區分敵我.
我們的種種不同是融合成一個更大的群體.
而不是要分裂成各種不同的獨特的敵我分明的群體.
所以我們彼此的使命不是平行線.
不是各有各做各有各行.
而是成為交織線交織點.
你中有我我中有你.
這個才是新的家.
這個才是教會.
我們很多時候常常以為.
保持禮貌互不干擾就是合一.

$^{721}$但聖經向我們揭示一個更豐富的可能.
合一是進入彼此的生命.
所以坦白說要進入彼此的生命有多難.
很難.
難到我們每一天都經歷磨合.
我們很多夫婦關係.
都在磨合當中產生劇烈的碰撞劇烈的討論.
因為我們都很不同.
坦白說如果兩個完全相同有人結不到婚.
通常都是一個快一個慢.
一個邏輯思維一個感性.
這種的不同要融為一體.
二人成為一體.
我們時常見證的.
我們用一生去見證這種磨合.
但為什麼我們仍然可以因愛如初呢.
因為有愛.
弟兄姊妹我們將這套東西搬回教會就可以了.
弟兄姊妹你和我的不同.
和家裡的不同不一樣.
只有基督的愛.
我們就可以融為一體.
甚至就算我們感覺到被挑戰.
甚至我們的步伐不一致.
甚至我們的方法不同.
但因為有愛.
我們可以經歷好像夫婦關係那種的甜蜜.
成為上帝給我們最好的禮物.
因為他不同.
所以我看得更闊的世界.
因為他和我不同.
所以我們可以做更大的事.
讓我們經驗這種因基督的愛而有的融合.
使我們成為基督裡面更完整的個體.
所以我們時常都這樣跟年輕人說.
你結了婚就知道了.
因為他就學會妥協.
他就學會一切共同的合作.
弟兄姊妹,這是基督裡面經歷生命豐盛的重要.
這不是剝奪自我,剝奪個性的方式.

$^{761}$這是我們生命能夠得到擴張的時候.
所以使命的跨越是各有立場的堅持.
去到共掃成為基督的新人.
從前我們都是各自誇口.
我有獨特性.
我有自己獨特的意見的群體.
但因為基督.
我們所有人都能夠跨越.
成為主內合一的生命共同體.
讓我們學會走近那些曾經我們避之則吉.
覺得他麻麻煩煩.
覺得他和我的看法不同.
我沒那麼多時間和你浪費生命去糾纏,去辯駁的人.
當我們學會拆下那幅無形卻有實的牆.
當我們學會和每一個不同意見的人.
不同觀點的人,不同步伐的人.
共同站在事奉的前線.
這時候我們就看見福音的震撼.
原來一個推崇自我.
自我狂傲,自我撕裂的世代.
真的可以因為耶穌基督.
在我們中間活出完整,完全的復和和好.
這不是一個團體活動.
這不是團體遊戲.
這是一個福音真實的見證.
我們不是要透過說服.
或者我沒有試圖說服你要彼此合一.
我們是需要一起仰望.
那個已經拆牆,已經鑄造新人的耶穌基督.
我們唯有在祂的大愛底下.
我們才找到那個連結點.
我們唯有在祂的大能的幫助底下.
我們才能夠跨出我們的怯懦.
所以如果我們要在教會裡面.
活出主內合一的生命.
讓世人看見隔閡不單止被廢除.
更加能夠轉化為一場超越血緣,超越立場.
超越一切人間標籤,神聖生命的聯合.
我們必須活在主的大愛當中.
我們必須認識主的教導.

$^{801}$我們唯有在上帝的教導底下.
在神的大能的幫助當中.
我們才能夠跨出這一切.
來到總結的時候.
我繼續用我的老習慣.
用四句屬靈的提醒.
成為今天拆毀重塑的提醒.
焦點跨越拆心網.
轉化跨越毀迷茫.
使命跨越重寬廣.
拆牆合一素仰望.
我們每一個使命門徒的拆毀與重塑.
是從拆毀我們心靈防禦的虛妄開始.
經歷基督的釋放.
我們擁有不再迷茫的自由.
擁有那份超越立場的寬廣.
超越個人看法感受的寬闊.
最終能夠在彼此接納當中.
一同仰望那位已經完全重塑了我們生命.
讓我們去領受這份新人.
新造的人.
生命的人.
那位基督.
弟兄姊妹.
我再說一次.
這個不是建議.
這個不是遊說.
這是一場徹底屬靈生命重塑合一的革命.
你願不願意參與呢.
如果你願意.
請你下一次遇到過去你曾經要迴避.
曾經要捉牆.
覺得他麻麻煩煩的人.
請你不要順從肉體的防禦.
讓我們記得.
基督已經拆毀了那面牆.
我們試試用溫柔的眼神.
用真誠的問候.
用謙卑的聆聽.
去回應那個你想避之則吉的麻煩人.

$^{841}$可能那個麻煩人在你家中生活.
可能那個麻煩人在你團契成長班當中.
可能那個麻煩人在你工作的場景下.
可能那個麻煩人就是你自己.
讓我們記得.
這些麻煩人.
都是神眼中同樣尊貴的人.
都是被塑造.
被大愛所塑造,新造的人.
這個就是福音.
這個就是拆毀,隔閡,隔閡之後.
神即搬在基督裡面.
我們竟然能夠合一.
靜音姐妹,讓我們一切同心同意.
我們不願意築起一幕牆.
活在孤單寂寞當中.
我們願意放下我們所有的靈角.
領受你的裝備和塗造.
求你將我們這些原本疏離的生命.
因為基督而合而為一.
讓世人看見你就是我們的和平.
看見你藉著我們所彰顯那份整全的榮耀.
我們能夠如此祈禱.
乃奉我主耶穌基督的聖名而祈禱.
阿們.
\newpage



\section{腓立比書 4:1-9}
\label{sec:032442}
\textbf{靠主與同心}~~~~陸幸泉牧師
\newline
\newline
連結: \href{http://www.church.com.hk/acms/content.asp?site=cdc\&op=show\&type=product\&code=032442\&layout=sermon}{\texttt{032442}} ~~~~ 語音日期: 2026-07-12
\newline
\newline
\hyperref[sec:032441]{\small{< < < PREV SERMON < < <}}
~
\hyperref[sec:index_chronic]{\small{[返順時目]}}
~
\hyperref[sec:032443]{\small{> > > NEXT SERMON > > >}}
\newline
\newline
腓立比書 4:1-9
\newline
\begin{longtable}{cl}
\hline
\hline
章節 & 經文 (和合本修訂版)\\
\hline
4:1 & \begin{tabularx}{0.7\textwidth}{X} 我所親愛、所想念的弟兄們，你們就是我的喜樂，我的冠冕。我親愛的，你們應當靠主站立得穩。 \end{tabularx} \\ \\ \relax
4:2 & \begin{tabularx}{0.7\textwidth}{X} 我勸友阿蝶和循都基要在主裡同心。 \end{tabularx} \\ \\ \relax
4:3 & \begin{tabularx}{0.7\textwidth}{X} 我也求你這真實同負一軛的，要幫助這兩個女人，因為她們在福音上曾與我、革利免和我其餘的同工一同勞苦，他們的名字都在生命冊上。 \end{tabularx} \\ \\ \relax
4:4 & \begin{tabularx}{0.7\textwidth}{X} 你們要靠主常常喜樂。我再說，你們要喜樂。 \end{tabularx} \\ \\ \relax
4:5 & \begin{tabularx}{0.7\textwidth}{X} 要讓眾人知道你們謙讓的心。主已經近了。 \end{tabularx} \\ \\ \relax
4:6 & \begin{tabularx}{0.7\textwidth}{X} 應當一無掛慮，只要凡事藉著禱告、祈求和感謝，將你們所要的告訴神。 \end{tabularx} \\ \\ \relax
4:7 & \begin{tabularx}{0.7\textwidth}{X} 神所賜那超越人所能了解的平安，必在基督耶穌裡，保守你們的心懷意念。 \end{tabularx} \\ \\ \relax
4:8 & \begin{tabularx}{0.7\textwidth}{X} 末了，弟兄們，凡是真實的、凡是可敬的、凡是公義的、凡是清潔的、凡是可愛的、凡是有美名的，若有甚麼德行，若有甚麼稱讚，你們都要留意。 \end{tabularx} \\ \\ \relax
4:9 & \begin{tabularx}{0.7\textwidth}{X} 你們從我所學習的，所領受的，所聽見的，所看見的事，你們都要繼續去做，賜平安的神就必與你們同在。 \end{tabularx} \\ \\
[1ex]
\hline
\hline
\end{longtable}
$^{1}$主席.
何俊仁議員.
主席.
何俊仁議員.
主席.
何志光議員.
主席.
何志光議員.
主席.
何志光議員.
主席.
何志光議員.
主席.
何志光議員.
主席.
何志光議員.
主席.
何志光議員.
主席.
何志光議員.
主席.
何志光議員.
主席.
何志光議員.
主席.
何志光議員.
主席.
何志光議員.
主席.
何志光議員.
主席.
何志光議員.
主席.
何志光議員.
主席.
何志光議員.
主席.
何志光議員.
明镜需要您的支持 欢迎收看訂阅明镜.
\newpage



\section{瑪拉基書 4:4-6}
\label{sec:032443}
\textbf{主角出場,君王降臨}~~~~石偉漢牧師
\newline
\newline
連結: \href{http://www.church.com.hk/acms/content.asp?site=cdc\&op=show\&type=product\&code=032443\&layout=sermon}{\texttt{032443}} ~~~~ 語音日期: 2026-07-12
\newline
\newline
\hyperref[sec:032442]{\small{< < < PREV SERMON < < <}}
~
\hyperref[sec:index_chronic]{\small{[返順時目]}}
~
\hyperref[sec:032444]{\small{> > > NEXT SERMON > > >}}
\newline
\newline
瑪拉基書 4:4-6
\newline
\begin{longtable}{cl}
\hline
\hline
章節 & 經文 (和合本修訂版)\\
\hline
4:4 & \begin{tabularx}{0.7\textwidth}{X} 「你們當記念我僕人摩西的律法，就是我在何烈山為以色列眾人所吩咐他的律例典章。 \end{tabularx} \\ \\ \relax
4:5 & \begin{tabularx}{0.7\textwidth}{X} 「看哪，耶和華大而可畏之日未到以前，我要差遣以利亞先知到你們那裡去。 \end{tabularx} \\ \\ \relax
4:6 & \begin{tabularx}{0.7\textwidth}{X} 他必使父親的心轉向兒女，兒女的心轉向父親，免得我來詛咒這地。」 \end{tabularx} \\ \\
[1ex]
\hline
\hline
\end{longtable}
$^{1}$我們一起先有個祈禱.
好 天父我們多謝你.
因為你是愛我們的主.
你是生命的主.
求主你幫助我們.
讓我們真的知道.
在我們這個人生裡面.
我們知道哪一個才是主角.
哪一個是配角.
也祝我們能夠找到我們的角色.
找到我們的定位.
讓主你自己真的顯出最大的光芒.
幫助我們學習思考你的說話.
祈禱奉耶穌名字而求.
阿們.
今天我們已經說到.
《馬拉基書》最後的部分.
《馬拉基書》其實只有四章.
四章的書卷.
所以在我們今天.
《馬拉基書》最後只有三節的部分裡面.
其實就是整個《馬拉基書》的一個總結.
一個呼籲.
所以我們藉著今天去看一個總結的時候.
我們先有一點重溫.
這個亦都是在過去七次的講道裡面.
都提及的一些思考和一些重點.
首先我們要重溫一下歷史.
我們知道《馬拉基書》裡面.
其實是隨著以色列國國度的一些轉變.
一些狀況而衍生了不同的一些階段.
事實上你見到在左邊中間的位置.
就是我們一路經常說的以色列國.
猶大國的地方耶路撒冷等等.
事實上在歷史裡面.
都是出現了不同的世代.
不同的朝代.
不同的國度.
所以如果你留意一下.
看一下這幅圖的右邊中間的上面.

$^{41}$就提到亞述帝國.
大家記不記得亞述帝國.
什麼時候講過.
上個星期約拿書就已經提及到.
亞述帝國有一個叫尼尼尾城.
這個地方裡面.
其實你見到是主前八世紀的時間.
主前八世紀有一個國度的興起.
但原來將時間流逝的時候.
又有新的國度興起.
接著你見到下一點就是巴比倫帝國.
大家記不記得哪一本書是講巴比倫帝國.
就是彈耳曆書.
就是Jenny講的彈耳曆書.
就是講巴比倫的帝國.
事實上亞述是主前八世紀.
巴比倫大約是主前七世紀的時間.
而去到後來這裡告訴你.
原來巴比倫的尼泊爾加密沙.
威脅了一段時間.
最後都威脅不完.
到他後代之後.
這個國都要滅亡.
而新的帝國又會興起.
就是波斯帝國.
你見到是主前六世紀的時間.
開始慢慢強大.
就是一個波斯的帝國.
其實就是隨著歷史的更替.
以色列國特別是猶大國.
隨著巴比倫的進攻.
三次的進城.
以致他們將猶大擄走.
以致基本上猶大亡國.
而亡國之後就隔了一些時間.
大約536年左右的時間.
他們才有機會回歸.
516年.
他們亡國是586年的時間.
這裡就說當他們回歸的時候.

$^{81}$為什麼能夠回歸呢.
就是因為波斯帝國的興起.
他們用了一個比較懷柔的政策.
讓他們能夠回歸.
也都隨著時間的轉變.
他們就開始回到耶路撒冷的地方.
你這裡見到就是.
回到耶路撒冷裡面.
主要有兩件很重大的事件.
包括了就是由索羅巴伯.
和大祭司約書亞.
就帶領他們重建了聖殿.
這裡就告訴你.
大約538年至515年的期間.
就是他們慢慢回歸的時間.
接著稍後的時間.
也都有文士以斯拉和尼希米.
去帶領他們重建城牆的時間.
就是這樣他們一批又一批.
回到耶路撒冷.
也都在這個環境之下.
我們說整個的馬拉基書.
就是他們回到耶路撒冷.
已經經歷了重建聖殿.
經歷了重建城牆.
應該整個生活都經歷重建的時候.
他們卻發現.
他們的生活.
他們的狀況.
並不是他們所期待的.
那個國家也都沒有因為.
他們重建了之後.
以致到能夠重新再一次.
以色列復國.
在眾多的失望之下.
他們就開始.
和上帝有一個對話的時間.
馬拉基書裡面就記載了.
他們和上帝的對話.
而這個對話的一開始.

$^{121}$就表明了他們.
神和百姓之間的那種狀況.
我想在馬拉基書的第一章二節裡面.
很清楚的說.
一個定位就是什麼呢.
神就說我曾愛你們.
神就很清楚說.
祂很愛以色列這班百姓.
這班人民.
所以整個定調就是說出.
神的愛.
去到馬拉基書裡面.
都是承接著一直在舊約裡面的訊息.
就是神顯明了祂愛.
祂所揀選的百姓.
祂所揀選的民族.
所以在這個基礎之下.
我們就看了七個不同的話題.
包括誰曾愛我.
A課的敬拜.
說他們的敬拜.
獻一些不合上帝心意的.
賜課的祭物.
然後也有與你有約.
說上帝提醒他們.
與上帝有一份的約.
是上帝一早就應許賜給這班百姓.
然而在這個過程裡面.
他們生活裡面.
行了很多的詭詐.
特別是在一些婚姻.
在一些關係上.
在一些偶像上.
都很多的詭詐的行為.
講到後面就說.
你也有今天了.
神要允許聽神說話的客人蒙福.
後面更加說.
我們要甘心樂意去奉獻的話.
我們會經歷到福浸眼眉的驚喜.

$^{161}$我們要樂於去事奉的時候.
有上帝給我們的賞哲.
讓我們經歷那種醫治的大樓.
就是這樣經歷了一連串的對話.
而在對話的過程裡面.
其實我們不斷聽到的以色列人.
有些什麼的回應呢.
他們一連串的回應.
都是一些些的抗辯.
一些的反駁.
對上帝很多的愛的表達.
幫助教導等等.
都是提出一些.
和上帝的方向不一樣的.
這裡說包括了他們問.
你為何事上愛我們呢.
他說煮的公菜是不潔淨的.
公菜上的果子和食物是可藐視的.
他說看看這些事何等繁瑣.
我們怎樣要轉向你呢.
我們何時奪取了你的物件呢.
我們在哪些事上.
我們藐視過你的名呢.
我們在何事上令到你不潔淨呢.
我們在何事上令到你厭煩呢.
似乎對上帝有很多不滿.
犯恨惡的也會說看為善.
並且喜愛他.
公平的上帝在哪裡呢.
我講了甚麼說話去頂撞你呢.
他們說事奉上帝是枉然的.
對於這班以色列的百姓.
已經有很多很多的投訴.
甚至很清楚有他們的教導和提醒給他們.
但這班人由始至終.
由去到第一港去到第七港.
都顯明了甚麼呢.
他們都自我中心.
自我中心自以為那位的主角.
我用四個字去描述.

$^{201}$死不悔改口是心非.
知而不幸自私自利.
博嘴博舌見利忘義.
以權謀私認賊作父.
顛倒是非顛三倒四.
巧言令色一言九鼎.
這個政府反映出的狀況就是.
以色列百姓只是在看自己的好處.
自己的期望.
自己想要的生活.
在這樣的情況下.
究竟他們怎樣去面對往後的日子呢.
特別在往後的日子裡面.
將要面對上帝的沉默無聲.
因為馬拉基書就是在舊約裡面的最後一卷書信.
最後一卷先知的書卷.
這個先知的書卷之後.
就有四百年的時間.
就再沒有任何先知的記錄.
這些書卷留下來.
面對這四百年的沉默.
他們怎樣能夠去面對呢.
在最後第四至六節.
四章四至六節.
上帝就留下了一些提醒.
一些盼望給他們.
所以今天我們就看到這個經文裡面.
究竟在當中我們看到了一些什麼.
他們要知道或者他們在總結的時候.
要看得到的事情.
第一個和大家看的分題就是.
有摩西頒佈律法.
但這裡我會說.
他是一個藍配角.
因為當那班以色列人期待自己是主角.
世界意義圍繞著他們.
中心運轉的時候.
就算摩西這個大人物.
都是在這個.
劇本裡面.

$^{241}$只是一個配角.
這個配角有多重要呢.
我們來看看第四章四節.
這裡一開始就講得很清楚.
就是上帝讓摩西在荷列山.
亦都是我們經常講到的西乃山.
在那裡神就讓摩西去頒佈律法.
而這些律法亦都成為了.
往後的一些典章律例.
摩西頒佈的律法.
我們最記得的就是那十屆.
但其實那十屆以外.
其實到後來.
按著摩西律法裡面所頒佈的每一條.
以色列人都將它重新歸納整理.
去化出613條的誡命.
這些誡命的本意.
就是讓這班以色列的百姓.
能夠和上帝有一個立約的基礎.
以致讓這班百姓.
能夠按著上帝的吩咐要求.
能夠過一個聖潔.
一個合乎上帝的.
一個吩咐的.
一個生活的標準.
一個參考.
上帝就留下了這些律法.
這裡第一個提醒.
就是摩西所頒佈的律法.
要我們怎樣做呢.
這班在回歸國土裡面的百姓.
要怎樣去對待摩西的律法呢.
在最後的總結裡面.
他提出一個很重要的觀念.
就是要去紀念.
要紀念這個僕人摩西的律法.
而紀念這個概念.
並不是只是回想.
有空想想.
而是要活出和進行.

$^{281}$馬拉基在這裡提醒我們.
要常常去紀念.
想起.
活出和進行這種律法.
就好像在申命記裡面.
很清楚.
教導以色列人.
怎樣去常常去記得.
活出這種律法.
申命記說.
六章四至九節.
以色列人你要盡心.
以色列人你要聽.
耶和華我們的上帝.
是獨一的主.
你要盡心盡性盡力.
愛耶和華你的上帝.
我今日吩咐你的這些話.
都要記在心上.
也要殷勤教導你的兒女.
無論你坐在家裡.
走在路上.
躺下起來都要吟誦.
要繫在手上作記號.
戴在額上作經集.
又要寫在你房屋的門框.
和你的城門上.
是的.
以色列人就是.
按著這個吩咐.
其實他們都很謹慎.
常常都去.
戴著這些經文.
在身上.
記得甚至背誦.
這些.
聖經裡面的說話.
神就是這樣吩咐我們要紀念.
摩西所頒佈的律法.
所以這個是.

$^{321}$生活裡面我們所要知道的.
因為這個律法能夠幫助我們.
能夠明白.
上帝的聖潔和公義的標準.
當我們每一次.
看這些律法的時候.
就好像一塊鏡.
你試過每天出街.
你通常都會怎樣?每天照一下鏡.
OK嗎?整不整齊.
才會出街.
而這個律法就讓我們每天都去.
掛在身邊的時候.
我們常常去紀念的時候.
讓我們知道自己是一個什麼狀況.
和神對照的標準.
之下他的狀況是怎樣.
使我們更加能夠明白.
我們自己在一個什麼的.
狀態裡面去面對上帝.
不過.
很重要的就是.
原來摩西所頒佈的.
這些律法確實是要讓我們.
常常去紀念.
不過摩西所頒佈的.
律法並不是.
上帝計劃裡面的主角.
並不是那個主角.
也不是.
上帝計劃裡面的主角.
律法有什麼用呢?.
在新約裡面解釋得很清楚.
到了新約時代.
就告訴我們其實律法.
最重要的是什麼呢?.
羅馬書三章二十節.
所以凡血肉之軀.
沒有一個能因.
律法的行為.

$^{361}$而在上帝面前稱賢.
因為律法本是要人.
認識罪.
這樣我們要怎麼說呢?.
律法是罪嗎?絕對不是.
但是若不藉著律法.
我就不知何謂罪.
若不是律法說.
不可貪心,我就不知何謂貪心.
這裡說出一個很重要的概念.
就是確實.
今天留下的以色列人.
他們要遵守.
要紀念摩西所頒佈的律法.
但這個不是他們.
最後的結局.
而這個留下的律法.
正正好能夠幫助他們.
能夠審查自己的罪.
羅馬書很清楚說.
律法本是要人認識罪.
要人知道罪.
知道自己在上帝裡面.
是不是一個罪的狀況.
弟兄姊妹.
今天我們要知道我們熟讀聖經.
我們熟讀神的說話.
並不是我們就藉著.
這些的手,這些的行為.
能夠讓我們生命能夠得救.
不過確實.
勤讀聖經讓我們明白.
上帝的心意,確實讓我們能夠.
知道究竟我們.
是一個什麼的狀況.
我們距離上帝的心意有多遠.
但當我們一都不認識上帝的.
律法的話,那個吩咐要求.
的時候,我們很多時候.
不知不覺就衰了.

$^{401}$我們就不知道.
或者我們不知道自己原來.
是衰到一個怎樣的地步.
我們又說世界盃.
世界盃裡面.
今年其中一個.
話題就是.
多了一樣的.
多了一些不同的.
安排,包括.
有VAR,VAR裡面.
整個入了球都可以.
不算數,這是其中一個.
昨晚都有看到.
但其中一個都挺有趣的.
就是.
今年的世界盃.
在球場裡面.
的要求是什麼呢?.
不准掩嘴.
師母剛剛就掩了嘴.
紅牌.
出場,離開這個教會.
今年就有一個.
規矩就是說.
當你每一個球員.
不准掩著嘴.
說話的時候.
當你掩著嘴的時候.
讓球證知道.
就紅牌了,所以你看到.
這幅圖就是.
第一個人掩著嘴.
對另一個球員說話.
接著那個球員就去.
跟球證說,舉報.
這個人掩著嘴,於是又怎樣?.
紅牌就出了場.
你覺得.
你不要理這個規矩.

$^{441}$怎樣都好.
但球員沒有.
知道這條規矩.
而做了這件事.
那又怎樣?就衰了.
出局了.
同樣就是.
如果上帝有很多的律法.
有很多的吩咐,有很多愛的教導.
給了我們,我們沒有看.
可能我們.
不知不覺就做了.
其實可能令我們.
被撒旦引誘我們.
離開了屬上帝.
這個場地.
所以我們為什麼.
要去明白.
這個律法.
我們確實需要明白律法.
但同時間我們要知道.
這個律法不是我們的終極.
不是我們的終極.
加勒太書裡面怎樣說.
加勒太書三章.
二十三至二十五節.
但這信還未來.
以前,我們被看守.
在律法之下,將被.
圈住,直到那將來.
的信顯明出來.
這樣,律法是.
我們的啟蒙教師.
直到基督來了,好使我們.
因信稱義.
但這信.
既然來到,我們從此.
就不再啟蒙教師的.
手下了.
原來.

$^{481}$摩西所頒布的律法.
不是終極.
原來有一個更大的終極.
就是.
我們可以因信稱義.
因信稱義.
律法只是讓我們去.
啟蒙我們,讓我們知道.
我們跟著律法.
去反思到自己,原來是.
不能夠靠律法能夠稱義.
但當我們能夠認識.
基督來到的時候.
我們能夠因信稱義.
以致我們就不再.
被這個律法所限制.
所以今天我們第一個.
我們要思考的重點就是.
我們要明白,摩西頒布.
律法為我們留下了.
一個很重要的.
我們可以按著.
神的吩咐去生活的.
一些參考,一些準則.
但是.
我們並不是因為這樣.
而讓我們能夠有一個真正的.
生命和生活.
所以今天第一個的思考.
就是我們今天要怎樣去紀念.
神的話呢?我們有沒有紀念.
神的話,常常去回顧.
實踐神的話的時候,以致能夠.
活出那種聖潔的生活呢?.
但是當我們.
一直很想去守著律法.
我們跟著律法的要求的時候.
我們又會不會不知不覺.
以我們守律法.
而驕傲自滿,就好像.

$^{521}$耶穌時代的文士和法理財人那樣.
守了613條的.
律例,守得到.
就覺得驕傲自滿.
以致到那個本質.
就是和上帝那份關係.
找回上帝那份的愛.
這麼重要的事情.
忘記了呢?.
所以今天我們要認真想想.
究竟我們的信仰生活.
我們往往是因行為.
稱義,還是因信.
稱義呢?.
求著幫助我們能夠明白.
摩西所留下的.
律法是幫我們的.
但同時間.
它並不是最重要.
我們終極要依靠的一些事情.
第二個我們要去看的.
藍配角是誰呢?.
第二個藍配角就是.
以利亞,以利亞就是一位.
先知,他要宣傳.
一個悔改的訊息.
這段經文裡面提到.
先知以利亞.
是一個什麼樣的人物呢?.
事實上他是一個.
很重要的舊約裡面的先知.
他能夠行神蹟.
斥責偶像,也能夠呼召.
王去悔改.
他是眾先知的.
一個代表人物.
這麼重要的先知,甚至乎聖經裡面.
說到他好像.
不見死亡一樣.
但這個先知裡面.

$^{561}$他留下一個重要的.
角色,吩咐是什麼呢?.
在《瑪拉基書》裡面總結的時候.
就說,原來我們要去看一看.
以利亞的時候.
就要看一看他,原來他帶來.
一個悔改,就是讓人轉向.
這裡提及就是.
讓父親的心.
要轉向兒女.
兒女的心要轉向父親.
轉向就是.
之前我們說的,好像180度的.
迴轉那樣,以致我們.
由一個破裂的關係.
去到一個復和的關係.
我們由一個向著滅亡的.
狀況,我們走回頭路.
我們向著一個永活的.
向導,所以以利亞.
就是一個這麼重要的先知.
帶來一種悔改,宣傳.
悔改轉向的訊息.
以致讓人不再活在.
一個咒座裡面.
不再活在咒座裡面.
而是邁向永生.
邁向神自己裡面.
人就要迎接.
真的要來的那一位.
事實上當你讀舊約.
你會明白,舊約裡面有很多不同的先知.
這裡說很多我們熟悉的.
以利亞,有以利亞的徒弟.
以利沙,也有我們.
所謂的大先知.
14位的大先知,或者.
所謂12位的小先知.
其實如果你用心留意一下.
不同的先知書裡面.

$^{601}$無論他講多少.
審判的訊息.
因為很多時候先知講的訊息.
其中一個很重要的範疇.
就是審判.
但同時間如果你熟讀聖經的話.
你會發現很多先知講完審判的時候.
很多先知書的結尾.
是講什麼?.
是講迴轉,是有出路的.
就是當你在.
上帝面前能夠悔改的時候.
就能夠有出路.
所以其實不只是以利亞.
宣傳一個悔改的先知.
事實上.
眾多舊約的先知.
都是傾向藉著他們的宣告.
以至讓帶領人.
面向悔改.
所以.
上個禮拜你聽了.
約拿書.
約拿書裡面都是講一連串的悔改.
包括.
船裡面的船長水手.
他們經歷悔改.
呼求神.
尼尼美這個大成.
第三章上個禮拜講什麼?.
他們尼尼美大成連王都一起.
去呼求悔改.
而約拿書最後.
就講到約拿都要思考.
他是不是要悔改.
這裡講的就是.
原來一連串的先知都是帶來一個悔改的訊息.
所以對於.
馬拉基書的百姓.
讀者他們都要知道.

$^{641}$他們的生命裡面.
要能夠捉緊和.
上帝的愛.
第一個就是要遵守這些律法.
但同時間.
要回轉.
轉向歸向神.
事實上當我們.
想想我們在講不同的先知.
帶來一些悔改的訊息的時候.
其實.
去到新約都有一位先知.
一樣的人物.
其實.
亦都是一位以利亞的人物.
就是哪一位呢?.
當我們在說.
神說要差遣一位.
以利亞先知來到.
原來去到耶穌年代.
正正好就有一個以利亞.
另外一個的.
配角以利亞出現.
而這個以利亞其實就是.
施浸約翰或者施洗約翰.
另外一個.
那個以利亞是怎樣呢?.
路加弗林就是這樣描述.
他說他將有.
以利亞的精神和能力.
走在主的前面.
叫父親的心轉向兒女.
叫搏逆的人.
轉向二人的智慧.
又為主預備.
迎接他的百姓.
這裡講得很清楚.
就是原來.
大人轉向的除了.
本身舊約那位的先知.

$^{681}$以利亞.
而去到新約的時候.
施浸約翰也承接著這個角色.
承接著.
馬拉基書裡面說.
我們要等待那一位的以利亞.
而這位約翰就似乎.
跟著要來那位的以利亞.
而他也是.
帶來一個轉向.
讓人去重回.
一份關係 知道悔改.
耶穌就這樣說.
如果你們願意接受的話.
就是要來的那位的以利亞.
就是那位的以利亞.
我想問.
如果你是以利亞.
成為萬眾等待的.
你會有什麼心情呢?.
可能你會興奮 你會雀躍.
開始自我膨脹.
今天我們的生命.
都是這樣.
如果摩西覺得.
全部人都要聽他所頒佈的律法.
而高興自滿.
自我膨脹的話.
他就離開了他本身有的角色.
以利亞來到.
如果他也同樣.
因為很受歡迎.
因為知道當時.
施禎約翰是很受歡迎的.
很多人都爭著跟著他.
如果施禎約翰.
也是這麼驕傲自滿的話.
他也離開了他本身.
帶著以利亞那份的角色.
但是以利亞怎麼說呢?.

$^{721}$我們看看約翰福音怎麼說.
約翰福音三章二十七.
至三十節說.
約翰回答說.
施禎約翰很清楚說.
他只是新郎身邊的朋友.
去見證新郎來到.
而他就說.
他必興旺.
我必衰名.
施禎約翰知道自己的角色.
就是他讓接下來的.
要興旺.
而他要退後一線.
退後一線.
他只需要做好他的本份.
他的角色.
他的角色是什麼?.
宣傳一個悔改的訊息.
要帶人去面對悔改.
而不是他令人悔改.
不是他令人悔改.
能夠找到出路的那一位.
而是他只是那個渠道.
去宣傳悔改.
所以施禎約翰.
就一直做好他自己的角色.
他自己的崗位.
以致他願意勇敢站出來.
對希律王說話.
他直接走出來.
指出希律王的不是.
希律王的問題.
第二點.
我們今天要想一下.
今天的教會的角色.
我們能不能夠.
以利亞式的聲音.
發出.
讓我們能夠像以利亞一樣.

$^{761}$發出一種讓人悔改的訊息.
這裡說我們今天有沒有.
以利亞的心智.
大膽宣傳悔改.
讓人與神修復關係.
轉向神.
我們要讓人轉向.
反過來.
看回頭.
看回神那裡.
以利亞就在做這個角色.
而且他很清楚.
他必興旺,我必衰微.
有時候我們事奉神的人.
我們是不是很介意別人的看法.
介意自己的面子.
介意自己是否興旺.
還是我們太介意自己的高低.
以致我們阻礙了福音興旺.
我們要想一想.
我們要如何忠誠地擔當.
傳福音的責任.
傳福音面對很多拒絕.
\newpage



\section{路加福音 2:1-7}
\label{sec:032444}
\textbf{耶穌的出生與一生}~~~~羅杰才牧師
\newline
\newline
連結: \href{http://www.church.com.hk/acms/content.asp?site=cdc\&op=show\&type=product\&code=032444\&layout=sermon}{\texttt{032444}} ~~~~ 語音日期: 2025-12-25
\newline
\newline
\hyperref[sec:032443]{\small{< < < PREV SERMON < < <}}
~
\hyperref[sec:index_chronic]{\small{[返順時目]}}
~
\hyperref[sec:032445]{\small{> > > NEXT SERMON > > >}}
\newline
\newline
路加福音 2:1-7
\newline
\begin{longtable}{cl}
\hline
\hline
章節 & 經文 (和合本修訂版)\\
\hline
2:1 & \begin{tabularx}{0.7\textwidth}{X} 在那些日子，凱撒奧古斯都降旨，叫全國人民都登記戶籍。 \end{tabularx} \\ \\ \relax
2:2 & \begin{tabularx}{0.7\textwidth}{X} 這第一次登記戶籍是在居里扭作敘利亞總督的時候行的。 \end{tabularx} \\ \\ \relax
2:3 & \begin{tabularx}{0.7\textwidth}{X} 眾人各歸各城，辦理登記。 \end{tabularx} \\ \\ \relax
2:4 & \begin{tabularx}{0.7\textwidth}{X} 約瑟也從加利利的拿撒勒城上猶太去，到了大衛的城名叫伯利恆，因為他是大衛家族的人， \end{tabularx} \\ \\ \relax
2:5 & \begin{tabularx}{0.7\textwidth}{X} 要和他所聘之妻馬利亞一同登記戶籍。那時馬利亞已經懷孕。 \end{tabularx} \\ \\ \relax
2:6 & \begin{tabularx}{0.7\textwidth}{X} 他們在那裡的時候，馬利亞的產期到了， \end{tabularx} \\ \\ \relax
2:7 & \begin{tabularx}{0.7\textwidth}{X} 就生了頭胎的兒子，用布包起來，放在馬槽裡，因為客店裡沒有地方。 \end{tabularx} \\ \\
[1ex]
\hline
\hline
\end{longtable}
$^{1}$各位弟兄姊妹,聖誕平安!.
很感恩,在大家當中,和大家一起過聖誕.
剛才施蕃獻了一首詩.
我很留意,很留心地聽和看那些歌詞.
我發覺它和我今天的講道是出奇地吻合.
真是同感一零.
待會我講完的時候,你再看那首詩歌.
你會覺得,真是非常吻合.
其實有很多聖誕歌,他選擇了這一首.
我今天的題目是耶穌的出生和祂的一生.
放大一點來說.
我們先有一個禱告.
親愛主,在你誕降,紀念你誕降的日子.
我們一起聚集,一起思想這個奇妙的救恩.
更加思想你的降生和你很特別的一生.
這和我們每一個人,特別是每一個基督徒.
都是非常息息相掛.
求主今日藉著聖經幫助我們.
讓我們更加明白.
你降生,你的救恩,你的一生.
我們這樣祈禱,奉耶穌基督的名.
誠心所願.
我曾經講過耶穌的一生有兩個定格.
一個是十字架,大家也知道.
基督教的記號,icon就是十字架.
去到哪裡都見到.
另外一個就是包著布,臥在馬槽裡.
一出生的時候,很清楚,這就是記號了.
這就是認出耶穌的記號.
很樸素,很簡單,甚至很卑微.
包著布,擺在馬槽裡.
其實一個嬰孩是不應該擺在馬槽裡.
這說明了耶穌出生的地方的處境是一個怎樣的處境.
一個是講到耶穌的死,一個是講到耶穌的生.
其實耶穌的一生是貫穿馬槽和十字架.
祂就過了祂的一生.
聖經是把耶穌基督.
告訴我們,上帝把耶穌基督賜給我們.
但不是一份禮物給了我們.
我們就收了祂,然後我們就享受這份禮物.

$^{41}$不是這樣.
祂把耶穌基督賜給我們.
祂就是我們的救主,也是我們生命的老師.
一個帶領者.
祂叫我們可以跟隨祂.
又可以無成好像祂那樣的生命.
有一個這樣的過程.
這才是救恩.
救恩不是我們拿了一份禮物,放下祂,已經很開心了.
拿了通行證,就可以上天堂了.
不是.
而是我們認識了耶穌.
我們就學習祂,跟隨祂.
像祂那樣來過一生.
上帝目標耶穌成為一個很被羨慕的對象.
我們有時會羨慕很多人.
這個人真的好,生得好.
有些人不信耶穌,有些人很會投胎.
生得好.
一個很羨慕的人,羨慕他的福氣.
羨慕他的幸運.
上帝沒有叫耶穌成為一個被羨慕的對象.
但是祂要耶穌成為一個被冒犯的.
或者學習的一個對象.
我們要學他.
耶穌的一生是怎樣的呢?.
一般人的一生,都有很多所謂的身不由己.
是被決定的.
但是如果我們想,祂是上帝的兒子.
那上帝應該會有參與當中.
就算我們是某一個人的孩子,某一個家庭的孩子.
我們很多事情是被決定的.
但是我們的父母是會幫我們做一些好的決定的.
是會將最好的東西給我們.
現在看我們的金科羅,很厲害的.
真的.
我們那一代,父母的生活條件不同.
我們那一代長大的.
和現在那一代長大的,真的很不同.
父母雖然是被決定.

$^{81}$但是父母是會幫他做一些對他最好的決定.
那主耶穌,祂來到這個世界.
天父參與祂來到這個世界.
那天父都要負責去幫祂做一些決定的.
幫祂做一些安排的.
這個是很合理,很順理成章的.
但是關於耶穌的出生.
一開始的時候.
正如我們讀過一段聖經.
就是居你柳作對你阿秦府的時候.
頭一次報名上冊,叫做戶籍登記.
就是我們現在的戶口登記是距離這裡.
以前的戶口登記就是.
你要回到原居地那裡做這個登記.
才查到你那些的資料.
有些資料顯示.
當時的戶口登記都很大的印象.
是十四年才做一次.
十四年一遇.
竟然遇到耶穌出生.
是很特別的.
這裡也有很多值得我們去想的.
就是祂一出生的時候.
就是成為了一個義務.
可以這麼說.
作為一個人的義務祂要去做登記.
那麼上帝是不是沒得選呢?.
還是祂是故意選的呢?.
你要想一想這個問題.
安排這個時候叫耶穌出生.
如果不是的話.
祂就不會回到這個伯利恆那裡出生.
怎麼不行呢?.
上帝告訴祂.
你回到伯利恆生就行了.
很簡單.
不需要弄得這麼複雜.
要安排一個報名上冊祂才可以回去.
周居勞頓全部的人都這樣回去.
是多麼的貧破呢?.

$^{121}$所以肯定不是沒得選.
而是故意選的.
我們可以深入思想一下.
上帝為什麼故意選.
我近來很喜歡進入聖經的時空裡.
嘗試去看看那些人物,關係,情節.
這對我思想聖經多了很多的發現和樂趣.
是故意選的.
聖經說祂及之時候滿足神.
就是猜遍祂的獨生子的降世.
是及之時候滿足.
不是及之時候滿意.
不是這個時候好了.
我們選對了這個時候是最好的.
我們現在的人生要選醫院.
又要選地方.
又要選國家.
要選很多很多東西.
上帝是時候滿足不是時候滿意.
這個時候是上帝選的.
為什麼呢?.
其實是為了和我選的.
我們值得去想一想.
這個差不多可以說是一個最差的時間.
是最差的時間.
是最不合適的時間.
我們經常選這個時間.
其實告訴我們.
哪怕是最壞的時間.
神都能夠將它變成美好.
神亦都不會去迴避.
整個耶穌的一生,耶穌降生.
是一路這樣重複說一件事.
神願意進入到一個最壞的處境裡面.
告訴我們.
那些在最壞處境裡面的人.
這些神都沒有迴避到這些最壞的處境.
這個對我們來說.
是一個很重要的信仰上的認知.
讓我們有信心,有力量去經過.

$^{161}$時間是這樣.
人呢?.
神是揀選了藥室和瑪利亞做耶穌的父母.
這兩個人.
聖經說他是一個木匠.
住在拿撒勒.
是普通到不能再普通的一個人.
是不是?.
是絕對的普通人.
他不是有名,有利,有身份,有知識.
另外一個跟耶穌差不多是同期.
那麼出名的人.
就叫做施洗約翰.
施洗約翰帶耶穌六個月.
是出生在祭司的家庭.
在耶路撒冷.
施洗約翰的出生就好很多了.
非常之好.
這個做耶穌的先鋒的人.
這個幫耶穌提鞋都不配的人.
竟然生在一個這麼好的家庭.
而耶穌呢?.
生在哪裡?.
在這個約瑟圖馬利亞的家庭那裡.
我們知道.
一個家庭對一個人的影響是很大的.
我們現在很多家庭都有傭工.
我們請一個傭人照顧小朋友都不行.
這個好不好?.
選一個來照顧我的小朋友.
我們請一個補習老師我們都不行.
現在上帝將祂自己最寶貴的兒子.
交給人去照顧祂.
祂選了約瑟圖馬利亞.
祂選撒迦利亞和伊利沙伯都好些.
有知識,環境又好.
一定是栽培得好些.
竟然將祂擺在蘭薩勒.
擺在約瑟圖馬利亞的手上.
字都不認識多少個.

$^{201}$可能真的不認識.
這樣的環境裡面.
這樣的人裡面來照顧祂.
神想表達些什麼呢?.
神揀選些卑微的人.
因此你和我才有份.
有時我都覺得有些搞笑的.
我們是神的僕人.
要配嗎?.
你有資格嗎?.
你讀到這裡才夠膽說.
應該都可以吧.
因為約瑟圖馬利亞都可以.
神是揀選那些很卑微的人.
所以在神面前你不要驕傲.
都不要自卑.
神又揀選了一個地方.
來讓耶穌成長.
差不多是一生.
不是一段時間.
那個地方叫蘭薩勒.
出生是一時.
多不運氣好.
過了就過了.
馬槽就是馬槽.
都是湊一會兒.
一會兒就出來.
但是成長是一生的.
祂揀選了約瑟圖馬利亞.
已經是一個沒什麼知識.
又什麼都很普通的一個人.
又揀選了一個地方.
約瑟圖馬利亞在耶路撒冷做平民.
那麼耶穌都可以出去旁邊.
找那些有知識的人討論聊天.
走遍整個蘭薩勒.
可能都沒有多少人識字.
就在那裡.
十二歲的時候.
耶穌做了一個行動.

$^{241}$祂就留在耶路撒冷.
不肯回去.
約瑟圖馬利亞就來找祂.
如果你細心一點進入這個時空.
每年約瑟圖馬利亞都帶耶穌去耶路撒冷.
似乎是一個很大很不尋常.
很不容易的一個行動.
去到耶路撒冷.
我相信.
祂應該會去找施洗約翰.
他們的家庭.
因為馬利亞懷孕的時候.
祂曾經去探望以利沙白.
以利沙白看到祂就說.
哎呀我主的母.
你懷胎是有福的.
所以他們是很認識到.
這個將要來到的一個小孩子.
那個很特別的身份.
對他們一定是非常好.
而且馬利亞在那裡住了三個月.
如果我在你家住三個月.
交情都不錯.
不是坐一坐.
是非常之好親密的一個關係.
當耶穌出世的時候.
祂每年去耶路撒冷.
會不會過門不入呢?.
不是很合理.
不是很可能.
所以一定會去找.
撒迦利亞和馬利亞的家庭.
即是說.
祂會找到施洗約翰.
耶穌和施洗約翰.
是相差了半年.
同期的.
他們兩個一定是非常之好的.
那些小朋友帶他去旅行.
去哪裡都好.

$^{281}$帶他去那些親戚那裡玩.
那些小朋友就會.
每年見一見.
夢都想著明年什麼時候再見.
是不是?.
我的孫女見到那些.
從外國回來的玩完幾天.
就已經不肯分開了.
他們是這樣長大的.
耶穌應該是很羨慕.
我是這樣說.
很羨慕施洗約翰.
他這樣的生活.
這樣的學習.
我真是很羨慕.
到他十二歲的時候.
他覺得他自己都長大了.
我要留在耶路撒冷.
我有一個這麼好的哥哥.
跟我一起學.
我又相信.
這個家庭都會接待.
都會歡迎.
如果他真的留在耶路撒冷.
是一件多好的事呢?.
十二歲不是小朋友.
十二歲出來工作.
可以很獨立.
所以耶穌很決定地告訴父母.
我是應該留在那裡的.
他不肯回去.
他不是撒脾氣.
不是什麼.
其實他是向著自己的理想去追尋.
他身邊有一個.
他模範的對象是十二歲的哥哥.
每天在聖殿的範圍裡跟拉比聊天.
學習這麼多.
自己從拿撒勒回去對著那群樹木.
你想想.

$^{321}$他的父母不知道為什麼.
你一定要回耶路撒冷.
一定要回拿撒勒.
有什麼理由?.
拿撒勒還有什麼好的嗎?.
沒有.
你只是糟蹋了一個天才.
如果他在耶路撒冷長大.
跟拉比聊天一起長大.
他那種成長.
他那種知識.
會弘大很多也不奇怪.
你要回去.
我又另一方面相信.
瑪利亞和約瑟.
他們是很清楚很堅定地相信.
耶穌是不可以留在耶路撒冷長大的.
如果他沒有這麼強的一種信念,肯定.
如果我們不妨讓上帝告訴他.
這個孩子你要帶在身邊.
你要讓他在拿撒勒這個地方長大.
我們知道耶穌在很小的時候.
已經顯示了他對上帝一個很信服的心.
瑪利亞和約瑟跟他這麼說.
這是上帝在他生命上的安排.
不是我們的意思.
不是我們父母捨不得你.
不是想你做裙腳仔.
耶穌就會被這個說服.
而且十二歲他回去之後.
他十三歲沒有再重複.
要一次,十四歲一次,十五歲一次.
不斷爭取,直到他成功爭取.
沒有.
他就這樣留在拿撒勒.
你知不知道聖經叫耶穌做什麼?.
拿撒勒人耶穌.
這就是他的身份.
他是拿撒勒人耶穌.
彼得說:我奉拿撒勒人耶穌.

$^{361}$成為了一個標記.
拿撒勒是什麼?.
是一個窗下.
一個不毛之地.
上帝竟然將耶穌.
擺在一個這麼貧瘠的地方裡面.
基本上是過了他的一生.
去到三十歲.
他的教育是自己摸的.
他可以去找人.
他不同施洗約翰.
施洗約翰去到成年的時候.
他就出來.
是另外一個進路.
他一生都可以說是.
困在拿撒勒.
所以我可以這樣說.
耶穌是沒有獲得過好的待遇.
更加沒有好的享受.
而最終的時候.
他都被釘十字架.
到他臨終的時候.
連衣服都被人拿走了.
只有一塊布.
就好像降生的時候一樣.
只有一塊布.
放在木頭上面.
當耶穌十二歲的時候.
他的父母叫他回去拿撒勒.
耶穌說他就順從了他們.
順服了他們.
這就是一句.
這就是耶穌的一生.
一句我願意.
那你就一世了.
順從一世.
耶穌順從地過了.
一種這樣的生活.
或者困在這裡.
為了什麼呢?.

$^{401}$不是他沒有能力.
不是他懦弱.
不是他沒有機會.
他又為了誰呢?.
我看主耶穌.
他的一生是向我們展示.
怎樣過你的一生.
或者可以怎樣過你的一生.
聖經說主耶穌來到降世是為要拯救罪.
其實我們是被每一個人.
無論你有錢,沒錢,有知識,沒知識.
都被很多東西困住.
有些比較明顯.
譬如你窮,你沒知識,你病.
是比較明顯.
有些是沒那麼明顯.
你覺得我又有工人,又照顧什麼.
但我們的生命都是被困住.
在不同的困境當中.
逃不掉.
這就是我們整個人的處境.
主耶穌這樣為了什麼呢?.
是為了讓祂向我們展示.
即使在一個這麼極端的困境中.
都困不住耶穌.
耶穌那種才是出來的.
祂好像什麼都沒有.
但祂又在一切上面.
過著祂自己.
做著祂自己要做的事.
這真是你想像不到的一種.
很奇特的一種結果.
我有幫助一個自閉症的機構.
所以有時候都會接觸到一些自閉症的家長.
我有一次聽到一個這樣的.
不可以說是故事,是真實的.
有一個爸爸.
他兒子是自閉症的.
所以很難跟他講道理.
很執著的自閉症的小朋友.

$^{441}$他緊張的時候就會尿尿.
幾歲的時候緊張的時候會尿尿.
他尿尿就站著不動.
那很難處理的.
他幾歲的時候就不動了.
那你身邊又沒有東西.
他糾纏得很厲害.
爸爸沒辦法.
後來就想到一個方法.
怎麼做呢.
那些小朋友尿尿.
他就用水淋濕自己的褲襠.
就好像尿尿一樣.
就告訴兒子,爸爸也尿尿了.
爸爸也沒事了.
小朋友尿尿已經有點異常了.
一個男人是很異常的.
是異常到非常.
在街上,在公眾場所.
我們沒事了.
一起去玩吧,去吃東西吧.
是真的,是人生的事.
這個爸爸真的把小孩教得很厲害.
他不是打他,不是罵他.
就帶著小朋友走出那個處境.
再慢慢教他.
聽完之後真的很有感.
我們的上帝何嘗不是這樣.
你在自己的處境裡.
祂走到你的處境裡.
祂甘願貧窮,甘願受壓抑.
祂告訴你,我也是這樣.
然後祂帶我們出來.
所以耶穌的救恩不只是在十字架上幫我們死一次.
你以為一個人死是很偉大.
不過不是最偉大.
你以為一個人活才難.
你猜不到吧.
你以為你死也沒問題.
要死的時候,我立刻拿出去.

$^{481}$我站在那裡幫你擋子彈.
子彈不會來的.
很難才會來的.
但是我活呢?活是很難的.
不單止是為我們死.
更是由瑪曹開始.
他的一無所有的一生.
他一生的坎坷.
一生的受壓抑.
一生過著一種很平凡.
被人藐視,歧視,受苦.
以賽卓這樣形容.
他無階型美容使我們注視他.
也沒有美貌使我們羨慕.
他被藐視,被人嫌棄.
多受痛苦,常經憂患.
為了什麼?.
為了你?.
為了我?.
救恩不只是把我們拉出去.
他還先走下來.
在我們的處境當中和我們一樣.
他為了那些在坎坷裡面.
在壓抑裡面.
在平凡,受歧視,受苦的人.
這樣過他的生命.
這樣過他的生活.
耶穌降生他的一生是告訴我們.
你不用害怕.
你不用感覺到羞恥.
你不用因為坎坷受壓抑.
很平凡受歧視.
你就覺得很自卑.
不需要.
耶穌就是這樣.
他幾生都是這樣.
他一生都是這樣.
他可以活出那種從容.
那種自信.
那種智慧.

$^{521}$那種仁愛.
完全無礙.
你不要以為你讀那麼多書.
你那麼有慧智.
那麼有身份.
你屬靈.
你敬虔.
很應該的.
你那麼好栽培.
不一定需要.
每個人都可以.
耶穌就演活了.
活出了.
但願這個降卑與我們同行的主耶穌.
他帶領我們成長.
帶領我們突破.
這個就是救恩.
這個就是福音.
我們祈禱.
親愛的主.
在這個紀念你誕降的日子.
我們思想你的降生.
你的一生.
求主的恩.
在我們心中.
激勵我們.
好像你那樣來生活.
也活出你的生命.
我們這樣祈禱.
奉主耶穌基督的名.
深受人欣賞.
\newpage



\section{腓立比書 2:1-11}
\label{sec:032445}
\textbf{與基督同一心思心志的合一表現}~~~~余大峰傳道
\newline
\newline
連結: \href{http://www.church.com.hk/acms/content.asp?site=cdc\&op=show\&type=product\&code=032445\&layout=sermon}{\texttt{032445}} ~~~~ 語音日期: 2026-07-12
\newline
\newline
\hyperref[sec:032444]{\small{< < < PREV SERMON < < <}}
~
\hyperref[sec:index_chronic]{\small{[返順時目]}}
~
\hyperref[sec:032446]{\small{> > > NEXT SERMON > > >}}
\newline
\newline
腓立比書 2:1-11
\newline
\begin{longtable}{cl}
\hline
\hline
章節 & 經文 (和合本修訂版)\\
\hline
2:1 & \begin{tabularx}{0.7\textwidth}{X} 所以，在基督裡若有任何勸勉，若有任何愛心的安慰，若有任何聖靈的團契，若有任何慈悲憐憫， \end{tabularx} \\ \\ \relax
2:2 & \begin{tabularx}{0.7\textwidth}{X} 你們就要意志相同，愛心相同，有一致的心思，一致的想法，使我的喜樂得以滿足。 \end{tabularx} \\ \\ \relax
2:3 & \begin{tabularx}{0.7\textwidth}{X} 凡事不可自私自利，不可貪圖虛榮；只要心存謙卑，各人看別人比自己強。 \end{tabularx} \\ \\ \relax
2:4 & \begin{tabularx}{0.7\textwidth}{X} 各人不要單顧自己的事，也要顧別人的事。 \end{tabularx} \\ \\ \relax
2:5 & \begin{tabularx}{0.7\textwidth}{X} 你們當以基督耶穌的心為心： \end{tabularx} \\ \\ \relax
2:6 & \begin{tabularx}{0.7\textwidth}{X} 他本有神的形像，卻不堅持自己與神同等； \end{tabularx} \\ \\ \relax
2:7 & \begin{tabularx}{0.7\textwidth}{X} 反倒虛己，取了奴僕的形像，成為人的樣式；既有人的樣子， \end{tabularx} \\ \\ \relax
2:8 & \begin{tabularx}{0.7\textwidth}{X} 就謙卑自己，存心順服，以至於死，且死在十字架上。 \end{tabularx} \\ \\ \relax
2:9 & \begin{tabularx}{0.7\textwidth}{X} 所以神把他升為至高，又賜給他超乎萬名之上的名， \end{tabularx} \\ \\ \relax
2:10 & \begin{tabularx}{0.7\textwidth}{X} 使一切在天上的、地上的和地底下的，因耶穌的名，眾膝都要跪下， \end{tabularx} \\ \\ \relax
2:11 & \begin{tabularx}{0.7\textwidth}{X} 眾口都要宣認：耶穌基督是主，歸榮耀給父神。 \end{tabularx} \\ \\
[1ex]
\hline
\hline
\end{longtable}
$^{1}$讓我們同心祈禱.
親愛的阿爸天父.
親愛的主.
求你差派聖靈.
讓我們內注的聖靈.
讓我們的心思.
意念.
價值觀.
能夠像主耶穌基督一樣.
幫我們.
看到我們內心的自我.
老我.
自嬌自我.
輕看別人.
求主憐憫幫助.
將現在的時間全部擺上.
求主你拖大.
祈禱交到奉主名而求.
阿們.
活著就像基督一樣.
好像耶穌基督在他身上.
好像見到他一樣.
有沒有.
反過來看到很多人都沒有反應.
更加危險.
我看到反面.
在他身上.
我完全看不到一些基督徒.
或者主義的樣式.
我看到的是相反.
我很難和他同行.
祈禱也好.
都想找一個人祈禱.
因為我很難和他同行.
更加重要的問題.
我們.
我們在神面前.
我們是一個基督徒.
在神面前.
神怎樣給分我們.

$^{41}$當我們這樣想的時候.
亦是一個很靠近的問題.
我經常問自己.
亦難以避免.
我在台上經常說我和我太太.
因為在我心目中.
在我心目中.
怎樣能夠最顯現.
究竟我是多愛主.
或者我是多改變.
就在家裡.
而莫過是我和我太太的關係.
我經常說.
愈愛就可能會愈恨.
我太太令我最憤怒.
所以我的臥上我太太見好了.
倒轉亦然.
但是弟兄姊妹.
我們都經歷過.
人生當中.
大家相信同一個主.
大家應該有很好的關係.
但為什麼關係有時會不太好.
有時不能夠同心合意.
夫婦如是.
家庭如是.
父母子女如是.
兄弟姊妹如是.
大家同一個主.
回到教會小組可能都如是.
多難有一個完美的小組.
5G竟然有些小組說.
我們小組近乎完美.
很難得.
他未必想.
一個近乎完美的小組.
但他不想拆組.
是否真的那麼完美.
可能是我們問的問題.
教會今天面對的是518.

$^{81}$但其實518可能是一個過程.
透過518.
可能反映.
究竟我們怎樣行神那條路.
是否真的像主耶穌基督.
我看這個當然是重點.
因為518總會過去.
但我們人生這條路.
我們是一個基督徒.
我們一直都要想這個問題.
今天是透過肥牙比書.
這個系列.
已經開頭了.
今天我要講一個很重要的經文.
二章一至十一節.
可以說是整個肥牙比書的膽.
整個肥牙比書.
就是透過基督頌.
大家都很熟悉.
熟悉到差不多會唸.
所以我特別想將.
這個基督頌.
放在整個肥牙比書的背景.
為什麼在肥牙比書.
面對什麼問題.
要處理什麼事情.
肥牙比書其實說什麼.
四章的書卷應該很容易明白.
是不是說合一.
是不是說受苦.
是不是說受苦當中喜樂.
是不是說謙卑.
其實是說什麼呢.
肥牙比書.
可能當我們了解背景.
就比較清楚.
肥牙比的地方.
你現在看到.
大家很熟悉的馬其頓.
即是現在在希臘附近.

$^{121}$在歐洲那邊.
隔離那邊是小亞細亞.
就是亞洲或者土耳其.
肥牙比是什麼城市呢.
肥牙比獨特之處.
是當時羅馬帝國.
少有的羅馬殖民地.
你說羅馬殖民地.
很少有.
你就自然想起.
因為我們很多人在英國殖民地出生.
英國殖民地.
不知你還記不記得.
英國殖民地.
以前在殖民地時代.
你入什麼行業都好.
公務員都好.
甚至在早期.
你看看歷史.
那些大班.
全部差不多都要英國人扣頭.
早期香港沒有法律學院.
你要去英國讀.
全部都英國文化.
英國牌才有用.
其他什麼專業都好.
其他外國牌都不行.
一定要英國牌.
當時你去英國浸鹹水.
你個人都特別聰明.
我是殖民地的地方.
去英國浸鹹水.
那個是英國牛津大學.
我們就是在這個地方長大.
所以我和肥牙比很有默契.
什麼默契呢.
作為一個羅馬的.
我們叫小友的.
殖民地.
他和馬其頓的首都.

$^{161}$當時馬其頓的首都是什麼呢.
是帖撒羅尼迦.
帖撒羅尼迦人口起碼比他.
最起碼是大十倍.
很多學者覺得.
人口可能去到二十萬.
十五至二十萬.
肥牙比多少人.
一萬至一萬五.
肥牙比是很小的地方.
因為.
當時羅馬帝國.
在二十年三十年打完仗之後.
他特別是將很多.
退伍軍人.
軍官.
一去哪裡退就走到那裡退.
你試想像在那個城市.
羅馬殖民地.
羅馬人是特別低一點.
西里羅馬人當中.
還有不同的官階.
大家不知道做過紀律部隊沒有.
我以前在香港和大家分享過.
我在做輔警.
你做輔警.
在街上見到警察.
在警署見到警察.
你一定要見他.
你看他的肩膀多一顆花.
多一顆花.
完全不同說法.
階級觀念是非常重的.
你看他如是.
他看不到你.
他直覺看不到你.
因為階級重到我不知道要看你.
你要給我面子.
我不用給你面子.
這個很重要.

$^{201}$整個菲律賓.
這個地方就是這樣的背景.
所以大家都會.
很吹噓自己的身份.
將一切榮耀.
你不給我也高舉.
你摸一摸肩膀.
你摸一摸身份.
我是羅馬公民.
所以大家都記得的時候.
當時保羅在第二次宣教旅程的時候.
透過聖靈的感動.
很出名的.
馬其頓意象.
去到馬其頓.
第一個城市是菲律賓.
所以聖經都說.
他雖然不是馬其頓的首輔.
但可能是頭一個城.
頭一個城的意思.
就是頭一個傳福音的地方.
當時說回保羅.
發生了什麼事.
傳福音的時候.
有一個侍女.
被鬼附.
結果.
他就將她釋放了.
驅鬼.
結果他反而被人捉去坐牢.
但被人捉去坐牢.
他就在監獄裡.
仍然是唱詩敬拜神.
他怕不怕.
他受很苦.
不怕.
又地震.
然後他的綑綁都鬆了.
結果嚇到過獄.
他就想自殺.

$^{241}$不要自殺.
我還在.
全家都相信主.
這個經歷很重要.
在保羅心目中.
坐牢不要緊嗎.
不要緊.
因為他活著就成為基督.
今天我們在座每一個人.
若然你要為耶穌基督傳福音的緣故.
你坐牢.
你心裡會想什麼.
還會敬拜讚美神.
我們所受的苦.
其實很少.
為主受的苦更少.
可能我們受的苦.
是你自己應該受的.
自己拿來壞.
但也不是為二受苦.
像保羅那樣為二受苦.
真的少或難得.
我自問我未必做到這個版.
一定要求神的幫助.
但當一天.
他在坐牢的時候.
這樣的事發生.
官長後來聽見.
原來你是羅馬公民.
馬上害怕.
羅馬公民.
特別有特權.
不要這樣.
我坐牢你沒死過.
你離開吧.
這個就是城市的背景.
肥納比教會.
保羅傳了福音之後.
保羅離開了.
肥納比教會是什麼教會.

$^{281}$在保羅心目中.
肥納比教會是好的教會還是不好的教會.
根據肥納比書.
保羅自己所說.
肥納比教會是我的喜樂.
我的冠冕.
傳福音是有份的.
所以就算他在監獄的時候.
他也很有愛心關顧.
肥納比教會.
因為也是從保羅開始.
和提摩泰.
所以他們也派了伊巴忽提.
去照顧保羅.
當時保羅相信.
家居被囚.
也可能在不同的地方.
但無論如何.
他們也派了伊巴忽提去照顧.
但伊巴忽提當時照顧他.
是患重病.
所以其中一個原因.
為什麼要寫肥納比書.
但也跟他們說.
伊巴忽提大病.
但現在康復了.
所以他們不用太擔心.
所以肥納比書.
是正面的偏向.
那為什麼要寫肥納比書.
自己在監獄.
那麼多地方都可以寫肥納比書.
為什麼寫給肥納比教會.
不用擔心之外.
第三點更重要.
就是第四章.
原來肥納比教會有兩個姊妹.
一個叫友雅蝶.
一個叫頓秦.
道姊友雅蝶.

$^{321}$和秦道姊.
要在主裡面有同樣的思想.
即使他們兩個有沒有同樣的思想.
是沒有的.
兩個是有爭拗的.
但這兩個女人.
曾經跟他在福音上.
是一同牢符的.
也保羅認定他們的名字.
在生命冊上.
保羅認定他們是信耶穌的.
但信耶穌也好.
全福音也好.
是不合的.
所以大家要分開.
而這裡用的是.
把整個肥納比書前後.
灌輸的一個希臘文.
叫做Forneo.
這個字可以解釋為思想.
但思想是解釋不到的.
因為它是灌輸整個人的.
無論你的心思.
你的意念.
所以今天我的題目是.
同一心思心智.
是一個心但是一個思念.
它不只是思想.
它是愛顧你的.
它有意志去做這件事.
三樣東西並齊.
這種.
是一種人生的態度.
我們怎樣對弟兄姊妹.
怎樣對你家的太太.
丈夫就是這個態度.
他就用這個字去串連.
若然主要是串這兩個人.
叫他們同一個思想.
若然這樣做.

$^{361}$一開始就直接講第四章.
為什麼要講頭三章呢.
特別要問一個問題.
第一個什麼問題呢.
就是上個星期所講.
狂牧師所講第一章.
為什麼要講第一章.
因為第一章是拉開了.
在開講第一章之前.
今天我要講第二章.
你發覺.
剛才所說的這個字.
這個所謂心思意念.
是整個全面書.
它是一個主線.
或者很重的字眼.
重複地用.
特別是在自己.
一開頭去跟.
菲律賓教會講的時候.
我對你們大家有同樣的想法.
我很掛念你們.
我經常都掛念你們.
想著你們.
在我心中.
之前我講過.
如果你在一個小組.
然後你去放假.
不是我.
你去放假.
然後你回來的時候.
弟兄姊妹.
真心.
不是假意.
幾個星期沒有掛念你.
沒有見到你.
很掛念你.
你老公老婆放了工之後.
幾個小時沒有見到你.
真的很掛念你.

$^{401}$你的心馬上甜了.
本來不同的思想.
本來就是這樣.
對菲律賓教會.
今日的經文.
中間點.
特別是用同一個字.
有同樣的思想.
思想一致.
同一個字出現了兩次.
這段經文主要出現了兩次.
後面這樣講.
各人不要只顧自己的事.
也要顧及別人的事.
不只是思想的問題.
思想去到一個地步.
不只是顧及.
去到一個地步.
是影響你做任何事的動機.
那個動機.
就是整份肥牛秘書所講.
一個很重要的問題.
究竟我們回教會.
我們有什麼動機.
我們在教會事奉.
有什麼動機.
在家結婚.
其實結婚有什麼動機.
有沒有人想過結婚有什麼動機.
或者生子女有什麼動機.
有些人生子女的動機是什麼.
說真一句.
為什麼生子女.
可能等我年紀大了.
之後子女去照顧我.
是不是.
希望子女孝順我.
希望子女聽話.
希望我去管他.
希望他做到我做不到的事.

$^{441}$之後讓我有榮耀.
很多不同的動機.
肥牛秘書.
再深入講.
要同一個意念.
同一個動機.
同一個價值觀.
說回第一章.
為什麼第一章這麼重要.
其中一個問題.
第一章.
保羅是透過自己的自白.
去引伸整份肥牛秘書所講的重點.
肥牛秘書的重點.
是要以基督的心為心.
而他第一章.
他就將這件事.
活現出來.
他將自己.
以基督的心為心.
去看肥牛秘教會.
再講一次.
他心裡有肥牛秘教會.
而肥牛秘教會的心裡有他.
因為他們是差派.
而他坐牢的時候.
都很掛念他.
知道他們有兩個姐妹吵架.
他專寫信給他們.
因為他掛念他.
他感恩掛念.
也祈求.
希望他能暫時得穩.
是不是很好的表達.
這樣就引伸了一個問題.
很好表達.
為什麼講到第一章.
一個很重要的.
上個禮拜大家也聽過.
是什麼問題呢.

$^{481}$問題在於.
弟們,我願你們知道.
我所遭遇的反而使福音.
更加興旺.
以致整個宮裡.
的衛隊和其餘一切的人.
都知道我是為了基督.
才受困所的.
而且大多數主內的弟兄.
因我所受的困所.
就篤信不疑,毫無畏懼.
更勇敢地傳講神的道.
他寫信給肥牛秘.
大家肥牛秘教會.
都知道他在坐牢.
他說,我在坐牢,但你不用擔心我.
雖然我在坐牢.
但問題是福音仍然廣傳.
所以他的著眼點是他坐牢還是福音.
他的價值觀.
最重要的是福音.
今日我們教會一樣.
我們是否最重要的是福音.
還是最重要的是自己的感受.
我猜很多人很重要是回教會,而我有什麼感受.
而不重要的是福音是否廣傳.
頂至沒有被詛咒.
但保羅完全不是這樣看.
而當他講到這一點的時候.
大多數主內的弟兄就跟著講.
但當中有些人.
出於妒忌.
出於自私.
動機不純正.
另外有些人出於好意.
出於愛心,上星期也講了.
我想問你一個問題.
這班人,現在講的是傳福音的人.
這班人.
在菲臘比傳福音.

$^{521}$還是在被囚的城市傳福音.
你覺得呢.
被囚的城市.
如果被囚在羅馬.
即是羅馬的人.
在菲臘比的地方.
在羅馬被囚的時候.
羅馬的人怎樣做.
跟菲臘比教會的人有什麼關係.
這個就是重點.
重點就是說.
雖然這樣,但也傳了出去.
都是好事.
所以這個很早的時候,我們教會也用.
是誤用了.
傳福音當然是好事.
但如果傳福音只是好事.
這樣的話,就講完了.
不用說出於什麼.
因為出於什麼不重要.
最重要是傳福音.
在教會的時候這樣講.
尤其是我們華人教會,最重要是做到事.
部門總之做到事.
開會就要講做到事.
所以做到事最重要.
518最重要,談了之後.
入了火,入了火就最重要.
過程不是很重要.
如果這樣講.
用這個講法是對的,那就不用寫這封信.
但他要專心寫這封信.
雖然福音傳了出去.
是好的,因為他最重要是福音傳出去.
但他們背後的動機.
就有分別.
所以在講這個的時候.
保羅有一個祈禱.
就是希望你們的愛心.
祈禱.

$^{561}$希望你們的愛心.
在知識和色見當中.
要增長.
即是你要分辨.
你要分辨到弟兄姊妹.
所傳福音背後的動機.
是什麼動機.
什麼心態,什麼價值觀.
他為什麼傳福音.
他為什麼在教會事奉.
他為什麼做執事,他為什麼做傳道人.
不要只看他傳福音.
看他為什麼傳福音.
這樣就不同了.
不要看他為什麼去小組.
不要看他是不是在小組.
而是他為什麼有什麼動機去小組.
這個很深沉.
因為這個就是影響心緣.
因為這個就不是基督的心為心.
而特別這句說話.
有些人在羅馬傳福音.
是出於自私.
或者另一個字叫結黨.
原文另一個解釋叫結黨.
他的動機不純正.
他只想加重我在捆綁中.
捆鎖中的煩惱.
你們想一想,這些人在外面傳福音.
保羅被人坐監.
為什麼這些人在外面傳福音.
會加重了.
保羅在監獄中的煩惱.
大家想得通嗎.
有沒有想過.
大家有沒有想過.
想到可以在座位輕微舉手.
一個傳福音的人.
他在這裡傳福音.
但他是出於歪心的.

$^{601}$動機不良.
但見到保羅一個傳福音很有效的人.
他在監獄中.
他在外面傳福音.
他覺得傳福音就令到你痛苦.
加添你痛苦.
這樣怎麼想呢.
大家想不想通.
希望大家不要想通.
因為這是小人之心.
因為他們是小人.
什麼叫小人呢.
他們根本是在爭榮耀.
在爭地盤.
這是他們的動機.
所以如果你進去坐監.
我在外面傳福音.
我便會有很多人跟隨我.
我便升職.
你便降職,我便高興.
我高興,你便很慘.
因為他們將心比己.
以小人的心.
去度君子之福.
就是他們.
今天教會有沒有人.
有人說有.
是什麼人會這樣想呢.
他有事奉.
人有事奉.
但見到人家恩賜比他高.
人家恩賜比他高.
最近他生病了.
心裡微微笑.
今次輪不到我.
但他會不會跟你說.
有誰在教會事奉呢.
他會跟你說.
我是小人,我是這樣反駁你的.
我出來是想威風你.

$^{641}$我講話是大聲上來聽.
有沒有.
但他會用什麼.
來講合一呢.
如果是領袖的時候.
他用合一的方法很容易講.
我是領袖,我有這個方案.
我們弟兄姊妹合一.
跟著什麼.
你跟我的方案就合一.
不過你不跟就不順服.
又不合一.
我想問你,這個講法.
是不是以基督的心為心.
他想人家跟他.
因為在他階級觀來講.
他才重要.
你不是這麼重要.
講到底.
他用任何基督教的口號.
去蓋過.
無論合一.
或者順服,或者愛心.
但講到底動機是什麼呢.
是想你跟他所做的想法.
但他沒這樣講出來.
就是這樣做,而被你看穿了.
結果是合一還是分裂呢.
就算被你跟他也好.
你們不是走基督的路.
所以這個背景很重要.
所以保羅一開始的時候.
他用了第一章所發生.
用了羅馬所發生的事.
他先以分辨出來.
他是以基督為主.
他要盼望的是什麼呢.
他要盼望的就是這件事.
就是在基督來到的時候.
他無可指責.

$^{681}$為什麼呢.
因為他所想要的.
在心中最重要的.
就是能夠讓耶穌基督.
讓到神.
得到榮耀.
原來這是他的心思.
他的心思,他整個人生的目的.
就是讓耶穌基督得到榮耀.
活著就好像耶穌一樣.
而他的祈求是什麼呢.
他的祈求就是說.
他們在一起的時候.
他們能夠站立得穩.
所以活著的時候為耶穌.
他自己講.
今天弟兄姊妹.
我想問大家一個很誠懇的問題.
如果你問你太太.
你問你丈夫.
你問你家人,你問你周圍身邊的弟兄姊妹.
究竟.
你是不是一個誠實的基督徒.
你只要問他一個心.
其實我們做事.
行事為人.
究竟有沒有考慮過他.
行事為人這個字在原文.
很奇怪,是什麼意思呢.
是作公民.
為什麼保羅要在這裡.
和在三七二十節.
用行事為人和作公民這個字.
來做一個串連呢.
因為在保羅來說.
最重要是一個天國公民的身份.
不是羅馬公民的身份.
而天國公民的身份.
就要耶穌基督的心態去思考.
而耶穌的心態去思考是什麼呢.

$^{721}$就是謙卑虛己.
而謙卑虛己最重要.
謙卑的相反是什麼呢.
驕傲.
但是驕傲怎樣表現出來呢.
原來驕傲的表現.
就是只顧自己.
而我經常說.
看別人比自己強.
我想了很久什麼叫.
看別人比自己強呢.
明明我唱歌比他好聽.
我不會走音,他走音.
那我怎樣看他比我強呢.
你明白我的意思嗎.
我明明神比我愛心大過他.
愛心爆棚,愛心那麼少.
我怎樣看別人比我強呢.
在經文裡的看他比他強呢.
又是finale.
但他有兩個finale.
你要和耶穌有同一個心智.
再加上要同一個愛心.
但再加上.
你要有這個finale.
什麼意思呢.
你是要有同一個心事.
而這個心智是有愛心.
之餘你是要mind他們.
你是要顧及他們.
這個叫同一個心事.
mind他們.
你要重視他們.
看別人比自己強.
是要重視他們的地步.
好像你重視你最愛的人.
你重視他的感受.
昨天我帶茶經的時候.
說起一件事很有趣.
我最近在帶一些.

$^{761}$每個人性格的問題.
我經常都沒有記性.
有時候在家裡工作我都會沒有記性.
我太太看到沒有記性就會很不高興.
因為我太太覺得.
我的沒有記性是什麼呢.
我說我沒有記性.
那太太會說什麼.
聽著.
可能你們這樣說.
其他東西又不見你沒有記性.
你自己要的.
你喜歡的你緊的東西.
你就有記性.
這句話是說什麼.
你不緊他.
mind.
others.
即是你緊他.
同一個心思.
看別人比自己強.
即是你緊他.
在教會的合一.
我想問大家.
是不是少數服從多數.
就是合一的表現.
如果用聖經來看.
這個不是合一的表現.
少數服從多數不是合一的表現.
多數服從少數是不是合一的表現.
都不是合一的表現.
領袖.
你去馴服領袖的決定.
是不是合一的表現.
領袖只聽你的表現.
你的意見最好.
在經文來說.
兩者都不是合一的表現.
因為兩者.
只不過是將其他人的意見.

$^{801}$去拉到是你的意見.
而在過程當中.
你是沒有.
重視其他人.
你沒有緊其他人.
我舉例.
我們很多移民.
或者是靠子女方法不同.
往往意見是不同的.
但對方.
最難受的是.
你沒有聽到我裡面.
所講的東西.
你沒有重視我.
我不是這樣跟你說了嗎.
為什麼你不記得.
是因為你不緊我.
你不重視我.
用我的舉例.
我太太跟我說.
其實余大峰是一個很驕傲的人.
余大峰是一個很自私的人.
余大峰你完全不謙卑.
因為你心裡面.
只想你自己.
你是沒有顧及我.
你不知道耶穌基督.
耶穌基督有什麼不同.
今日經文最主要講的是.
虛己.
虛己是.
袋淑.
道成肉身.
所以虛己.
顯示他是謙卑.
也顯示他不只是為自己.
他不是為自己.
他為誰呢.
他就為我們.
但今日經文.

$^{841}$重點.
是兩個字.
為什麼耶穌基督.
這麼謙卑.
這麼虛偽.
這麼倒空.
玩遊戲.
馴服.
馴服至死.
他馴服誰.
馴服我們.
他馴服父神.
耶穌基督的心.
他怎麼想.
他馴服父神.
當然其他經文補充.
他愛父神.
他的愛顯然在他馴服.
而他馴服.
而他的馴服.
是因為神想用這個方法救贖.
因為這個方法.
符合你的公義.
符合你的愛.
所以我因為馴服你.
我就會寫記.
今日我們問同一個問題.
我們待會回教會.
我們待會在事奉.
我們和家人吵架的時候.
意見不同的時候.
我們很想將別人的意見拉到你的意見.
因為這個就是你的階級觀念.
你覺得你這樣才有Faith.
你這樣才給我面子.
你當我是你老公的時候.
你聽我說.
不是,你當我是你老婆的時候,你聽我說.
如果你當我是牧師,你聽我說.
如果你當我是執事,你聽我說.

$^{881}$如果你當我是中文堂的,英文堂聽我說.
如果是中文堂的話.
如果你是甚麼,你聽我說.
個個都是想拉攏不同的人.
聽他所講的不同意見.
如果你這樣做的話.
是完完全全不是合一.
因為合一.
是倒轉問題.
你有沒有考慮別人的意見.
而考慮當中,你有沒有重視別人.
而這件事,為甚麼這樣做.
因為當你能夠寫記.
而重視別人的時候.
別人就在你生命當中.
見到基督的樣式.
你活著.
就是基督,你活著.
就是為了基督,因為別人見到.
你是卸下了自己的利益.
所以這堂課讓我有很大的反思.
因為我要問自己的動機.
我要問.
甚麼是合一的重心.
甚麼是謙卑.
原來每當我在家.
在小組.
在團契,在同工.
在執事會,在不同的教會.
若然我只是看自己的向度.
我沒有去關心別人.
沒有去著緊別人.
沒有去重視別人.
我只是想別人認同我.
因為我很feel good.
很有榮耀.
我就無意間.
是榮耀歸給自己.
而原來.
after all,很多基督徒.

$^{921}$都是很重視.
自己的榮耀.
而很容易在教會.
事奉.
最終.
可能只是血氣之爭.
這時候.
在神來看.
在耶穌來看.
我們是否一個人敗的基督徒.
是很值得去反思.
大家一起祈禱.
當我們去到你的面前的時候.
因為是你的捨己.
你每一步.
有你的道成肉身.
有你的.
教導.
有你怎樣去被人唾罵.
但你仍然是.
向著這些罪人.
去伸手.
哪怕你明知一些你醫好了.
患麻風病的人.
他不會回應你.
你仍然是著緊他們.
你仍然是去醫治他們.
因為你說.
醫生來到.
不是醫這些健康的人,是醫病的人.
你心目中是著緊每一個病的人.
但今天.
可能我們的意識形態.
不和我們.
剛剛談的弟兄姊妹.
就不和他們談.
不和我們剛剛談的聖駕弟兄姊妹.
千萬不要放在我一組.
在教會的討論.
開很多的會當中.

$^{961}$可能是血氣之爭.
我們用不同的觀點.
不同的角度.
去證明我的看法.
我更多的論據.
更多的證據.
但整個討論當中.
竟然我們沒有考慮.
真真正正著緊對方.
而也表現了出來.
結果.
會議開了.
決定做了.
但我們不是更加合一.
我們只是沒有說出來.
內心很難受.
可能我們都是成為一個.
以小人之心度君子之腹的人.
我們也是這樣度保羅.
但保羅反轉和我們說.
那又何妨.
他被人這樣冤枉.
又何妨.
他被人這樣奚落又何妨.
因為他單以耶穌基督.
為中心.
而他亦都是單以.
耶穌基督的想法為想法.
耶穌基督是中心.
耶穌基督.
都是能夠達到這個中心的途徑.
今日經文更加說到.
如何能夠做到這件事.
若然.
我們在.
基督裡有這個.
鼓勵.
若然我們有這個愛心.
若然我們有.
聖靈傳契的時候.

$^{1001}$三個都是若然.
但誰不知.
經文所說的.
其實是忌然.
因為在希臘文當中.
其實是一個忌然的表達.
換句話說.
我們怎能夠有合一,怎能夠謙卑.
原來經文提醒我們.
既然我們經歷了.
三位一體神.
的恩典.
既然.
我們就要更加去活出.
人與人之間那種.
關懷,憐憫.
所以反過來說.
為何今日我們不夠謙卑.
我們今日只顧自己的事.
我們自私.
我們結黨.
原來可能.
我們花了很少時間跟你結連.
我們真的沒有好好.
經歷主理恩典.
原來我們的信仰.
是那麼膚淺.
所以保羅.
為肥納比書的祈禱.
今天求主都幫助我們.
我們的愛心.
在知識和識見上.
去增長.
我們的信心.
愈發增多.
從而我們在全福音.
能夠站立得穩.
從而我們能夠一條心.
像耶穌基督的心.
好使我們見主面的時候.

$^{1041}$問心無愧.
無可指責.
人們見到我們的行為.
表現,然後將榮耀.
不是歸給我們.
而是歸給為主.
但願今日我們面對五一八的挑戰.
最終我們在考慮.
一切的事情.
更加考慮別人.
而最終將一切榮耀.
頌讚歸給主你.
感恩站美祈求,奉我主耶穌基督的聖名.
阿們.
\newpage



\section{士師記 16:30}
\label{sec:032446}
\textbf{聖經人聖經事:死時殺的人比活著時更多}~~~~黃克勤牧師博士
\newline
\newline
連結: \href{http://www.church.com.hk/acms/content.asp?site=cdc\&op=show\&type=product\&code=032446\&layout=sermon}{\texttt{032446}} ~~~~ 語音日期: 2026-07-11
\newline
\newline
\hyperref[sec:032445]{\small{< < < PREV SERMON < < <}}
~
\hyperref[sec:index_chronic]{\small{[返順時目]}}
~
\hyperref[sec:032447]{\small{> > > NEXT SERMON > > >}}
\newline
\newline
士師記 16:30
\newline
\begin{longtable}{cl}
\hline
\hline
章節 & 經文 (和合本修訂版)\\
\hline
16:30 & \begin{tabularx}{0.7\textwidth}{X} 然後他說：「讓我與非利士人一起死吧！」他盡力彎腰，廟宇就倒塌了，壓住領袖和廟宇內的眾人。這樣，參孫死的時候所殺的人比活著所殺的還多。 \end{tabularx} \\ \\
[1ex]
\hline
\hline
\end{longtable}
$^{1}$弟兄姊妹,我們又讀聖經了.
《生命說書》網絡聖經人聖經士.
可以看看神做了什麼工作.
「死時殺的人比活著時更多」.
「死時殺的更多」.
這是我今天想和大家看的題目.
這句說話也是聖經有的.
不過不完全是.
我將它簡化濃縮了一點.
「參孫在死時殺的人比活著時更多」.
這是聖經說的,這就是參孫在他的事先生涯裡的評價.
16章第28字.
「參孫在這個時候叩高夜禍」.
話說他叫小孩子帶他到柱子裡.
「我要唉唉,頂不住」.
他知道這柱子頂著整個建築物的大柱.
上面有三千人.
我不知道參孫知不知道.
只不過聽到上面的聲音.
很多人在看,很多人在笑他.
他只聽到.
可能他也會問「哥哥仔,上面有幾多人」.
「聽說有三千人」.
「這個可能是我最後的機會了」.
他就叩高夜禍,懇切祈禱.
第一點:「主耶和華,求你眷念我」.
這個藍色字兼畫了.
其實聖經裡面很多「主耶和華阿」.
「神阿」,「耶和華阿」.
那些其實是沒有「阿」字的.
中文的感嘆語氣是寫不出來的.
所以要加個字.
「耶和華阿」.
即是要加個感嘆.
那種懇切在那裡表現.
其實聖經是「主耶和華」.
祂的禱告是「主耶和華,求你眷念我」.
祂的稱呼沒有特別「阿」的.
所以希伯來文是沒有這個字的.
「主耶和華,求你眷念我」.

$^{41}$眷念,這個禱告我們經常都要做的.
可以不妨在你困難的時候.
在你痛苦的時候.
「主耶和華,求你眷念我」.
如果你覺得你去到人生最低谷.
簡直好像參宣一樣.
你更加要做這個禱告.
很簡單的禱告就夠了.
不用長篇大論,像品臣那樣.
怎麼行呢?.
「主耶和華,求你眷念我」.
禱告啦.
「眷念」這個字我特別要提的.
「眷念」這個希伯來文.
「Sakka」是紀念,回憶,Remember.
其實是「你記得我」.
「你記翻起我」或者「你記得我」.
有時候我去到不同的教會.
很多人就出來問我.
「你記不記得我?」.
其實我最怕這個問題.
「你記不記得我?」.
哎呀,怎麼回答呢?.
特別現在更加難.
因為這幾年有疫情.
香港人在世界很多地方.
特別是香港人必定要戴口罩.
不戴是犯法的.
戴著口罩過來,又戴著帽子.
接著問「你記不記得我?」.
當然不認得了.
怎麼回答呢?.
不戴口罩我都不認得.
我又真實,我沒有騙他.
我就說「樣子我就面熟」.
「不過名字我叫不出」.
這個答案很厲害吧,大家都可以學.
「樣子就熟,不過時間叫不出名字」.
「不好意思,你講個名字給我聽吧」.
「陳大文啊,陳大文」.

$^{81}$「是喎,陳大文我真的記得」.
「不過哪間教堂呢?」.
「啊,不記得」.
我又講一下,讓他自己講.
他說「求你紀念回憶,你記不記得我?」.
「你記住我啊,求你記住我」.
「主耶和華,你記返起我啊」.
意思是這樣,你記返起我.
記返起什麼?.
記返起我就是你給多一次力量我.
記返起我是什麼呢?.
我就是你叫我做拿西爾人.
要我做拿西爾人.
我聽我爸爸媽媽講.
我爸爸媽媽,我媽媽是不能生育不懷孕.
現在生了我出來.
是指明要做拿西爾人,屬於你的.
「主耶和華,你紀念我,我是屬於你,我是拿西爾」.
「嗯?現在才講?」.
「遲了少少喎,那時候你不講的」.
「大隻那時候你不講的,強那時候你不講的」.
「又不服侍我,就走去找這個找那個」.
「又特地去非利士人那裡招搖」.
「現在,抵啦,我不會聽你的啦」.
「哎呀,我這個神不是喎」.
他的禱告就是,他就是怕神是這樣.
所以他很誠懇地.
「主耶和華,求你回憶,是你叫我做拿西爾」.
「我雖然得罪你,不過我來到這一個位已經沒有了」.
他就說「上帝呀,耶路撒冷,神呀,求你賜我這一次的力量」.
「我只是要一次力量而已」.
「那些力量在頭髮那裡,頭髮都說漸漸長回來了」.
「是不是,給我一次力量吧」.
很有趣,參宣來到這裡,雖然求神你紀念我.
但是他都還挺自我,脫不了.
「求你給我一次力量,使我在非利士人身上布那碗我雙眼的仇」.
原來他到現在還記住「你把我的眼睛挖了,這個仇我一定要報」.
他就說「神呀,我都完了,走到這裡都知道絕路了」.
「求你給我可以再有一次力量報這個仇吧」.
「記得,我對於他下面我就一般般贊同」.

$^{121}$「求神記得他原來是給我繼續有力量報我被挖眼的仇」.
不過神都紀念他.
你要記得,主耶和華求你紀念我是一個非常好的祈禱.
於是參宣抱著托房的那兩根柱子.
左手抱一根,右手就抱一根.
這兩根柱子離開不遠,即是說房子的結構集中在這兩根柱子.
不過也有其他柱子幫忙,最大的這兩根柱子出來了.
這叫抱,抱就是我們想著全部抱過來.
但其實不是,抱可能是這樣,兩根柱子都能推,他盡力推.
左手一根,右手一根.
「我情願與非利士人同死,就盡力屈身」.
他開始有力量了,原來他頭髮漸漸長一點已經有力量了.
他被剃了,完全沒有,長一點已經重新有力量.
「盡力屈身,防止倒塌,壓住首領和房內的眾人」.
即是說他在下面兩根柱子,房子頂著上面,有三千人在上面.
上面應該還有一個頂的,所以他推倒底下這兩根柱子.
其他柱子支持不住了,開始倒下來,一倒就連上面的頂都倒了.
於是就壓住首領和房內的眾人,三千人,第二點.
「參孫死時所殺的人,比活著所殺的還多」.
這個就是他最後關於參孫的貢獻.
他死的時候,殺的人比活著所殺的還多.
那又怎樣呢?.
「參孫的弟兄和他父的全家都下去取他的屍首」.
「抬上來葬在索拉和以實圖中間,就是他父嗎哪亞的墳墓裡」.
好,總結,第三點.
「參孫著以色列的事司二十年」.
這個就是他的總結,這樣做了事司二十年.
這二十年做了什麼呢?.
做了死時所殺的比活著所殺的還多.
這個就是參孫人生事司生涯的總結,二十年.
其實你作為一個事司,要做什麼呢?.
不只是殺敵,事司要做法官記.
如果事司記就是記錄所有法官.
他們要做的是調解,是要判斷這件事這樣不好.
你要賠給他,你要賠給他,你要賠什麼,你要賠什麼.
是做判斷,做法官.
另一方面,他要帶領他的軍隊去跟他的外邦人打仗.
另一件事,他是要讓以色列人可以懂得敬畏耶和華.
跟隨神的法則.
亦有神父呼籲他做那西亞,是做那西爾人的.

$^{161}$哎呀,不過他沒有做到.
其實他殺的人,死的時候殺的人比生前還多.
做事司二十年,有沒有貢獻呢?.
他是歷史上殺菲爾士人最多的人,可能真的不少.
不過不算是非常有屬靈的貢獻.
反而他留給我們屬靈的警惕多些.
殺敵很多,但是不算是榜樣.
所以在參宣來說,他的墓志銘寫什麼呢?.
借過人死的時候,殺的人比活著的時候更多.
就是他的墓碑所寫的.
借過人生前,帶領很多人跟隨耶和華.
又帶領他們戰勝仇敵,那多好.
這個墓志銘寫著死者約參宣.
死的時候殺的人比在世的人殺的更多.
做了事司二十年.
並不是一個最光彩的墓志銘.
他幫助我們記錄自己一生回望的時候.
不是看到我們多麼強壯,多少人被我們打倒.
多少人被我們踩在下面,多少人跟我搵食.
那巴不巴閉?巴閉,不過死了.
他生前都做過很多事,不是的,是盲的.
屬靈上盲的,那這一生人浪費了.
記得在需要的時候,起碼要做一個倚靠上帝的人.
靠一靠,靠頭髮長,屬靈的能力.
倚靠的話,不然他離開了,仇敵就來了.
然後力氣消失了,那就糟糕了.
比我們在參宣一生裡面有更多的學習.
參宣是死時比活著時殺的人更多,一起祈禱.
我們看完參宣一生,如果以個人去戰勝敵人.
他一個人比一隊軍隊還強,而且他很聰明.
都很聰明,都很醒覺,不過太自我,不過太任性.
不過太自負驕傲,仰賴自己的肌肉,手乖起展.
很大力,目中無人,結果他是失敗的.
是輸的,甚至失去性命.
如果他懂得做一個好的拿西爾人.
帶領著做榜樣,帶領以色列人歸向你.
傳言屬於你好像拿西爾人一樣.
那整個以色列的蒙福就更加大,求主幫助我們.
主啊我們時都稱呼你做主啊,讓我們這個人屬於你.
雖然我們不是做拿西爾人,不過是做一個屬於主的人.

$^{201}$是為主而生活的人,求主教導我們.
隨天祈禱奉耶穌基督名禱.
阿們.
\newpage



\section{約翰一書 1:1-4}
\label{sec:032447}
\textbf{生命之道,團契生活}~~~~劉燕飛教師
\newline
\newline
連結: \href{http://www.church.com.hk/acms/content.asp?site=cdc\&op=show\&type=product\&code=032447\&layout=sermon}{\texttt{032447}} ~~~~ 語音日期: 2026-07-12
\newline
\newline
\hyperref[sec:032446]{\small{< < < PREV SERMON < < <}}
~
\hyperref[sec:index_chronic]{\small{[返順時目]}}
~
\hyperref[sec:032448]{\small{> > > NEXT SERMON > > >}}
\newline
\newline
約翰一書 1:1-4
\newline
\begin{longtable}{cl}
\hline
\hline
章節 & 經文 (和合本修訂版)\\
\hline
1:1 & \begin{tabularx}{0.7\textwidth}{X} 論到從起初原有的生命之道，就是我們所聽見、所看見、親眼看過、親手摸過的－ \end{tabularx} \\ \\ \relax
1:2 & \begin{tabularx}{0.7\textwidth}{X} 這生命已經顯現出來，我們看見了，現在又作見證，把原與父同在，並且向我們顯現過的那永遠的生命傳揚給你們－ \end{tabularx} \\ \\ \relax
1:3 & \begin{tabularx}{0.7\textwidth}{X} 我們把所看見、所聽見的傳揚給你們，為要使你們也與我們有團契，而我們的團契是與父和他兒子耶穌基督所共有的。 \end{tabularx} \\ \\ \relax
1:4 & \begin{tabularx}{0.7\textwidth}{X} 我們把這些事寫給你們，使我們的喜樂得以滿足。 \end{tabularx} \\ \\
[1ex]
\hline
\hline
\end{longtable}
$^{1}$謝主仔細睇下呢首歌ge 時候.
發覺佢同今日ge 經文ge 訊息係非常之配合.
我相信示範都花jor 好多ge 心思去選曲.
詩歌都話當我地願意去順服主導ge 時候.
我地係可以經歷得到團契交易ge 甜蜜.
今日我就比個題目佢叫做.
生命之道 團契生活.
D 兩組ge 四字詞相信大家都唔陌生.
我地都經常會講.
首先講下生命之道.
可能比起你去讀聖經睇到呢四個字.
你更加多ge 機會就係從教會gwo 度聽到.
呢一個ge 查經課程.
因為原來生命之道係.
中辰2007年ge 時候就已經開始ge 一個信徒讀經運動.
然後呢我地教會係2016年ge 時候就開始去參與呢個課程.
咁所以呢我去一路參加ge 時候.
我數下自己D 筆記呢我就列出jor .
原來已經不經不覺讀jor 咁多個書卷.
咁我都唔知有冇漏到.
原來呢我地都好熟悉架喇喎.
原來生命之道呢一樣ye .
咁賣個廣告呢.
咁如果你平時ge 星期四晚係有時間.
都想同我地一齊去讀聖經.
聽講解做功課ge 話呢.
咁我地鼓勵你可以報名參加啦.
咁無論你之前有冇參加過.
我地下一期呢就係九月開始呀.
咁就係讀貼薩羅利加前後書ge .
咁跟住講到團契生活呢.
我地都唔陌生啦.
因為呢講起團契.
咁一般黎講你都會諗起係我地教會裡面有一D 單位.
佢地就叫做團契.
咁唔知你數得出教會裡面幾多個團契ge 名呢.
除jor 你自己返緊gwo 個啦.
咁你仲可以數得出幾多個團契ge 名呢.
你知唔知教會裡面.
我地教會其實有幾多個團契ge 呢.

$^{41}$咁你仲記得係個零月之前呀.
即係五月三十一號呢.
咁教會係主日黎ge .
就下晝呢舉辦jor 一個活動喎.
咁大家有冇人記得gwo 日裡面我地做jor D 乜ye 呀.
食ye 跟住呢.
行樓梯打乒乓球.
係.
gwo 日ge 下晝呢教會就好似一個家庭聚會咁樣.
好似一個開放日咁樣.
咁當日呢有導賞團啦.
仲有D 導賞員呢就帶我地去行下教會唔同ge 樓層.
可能平時好少機會去ge .
咁我地就睇下原來gwo 個樓層.
ge gwo 樣ye 或者個設計呢係有D 故事係裡面ge .
咁我地一齊去聽下啦.
咁我地又會有乒乓球比賽啦.
咁好深刻gwo 個波亞斯團.
波亞斯團就gwo 隊叫做咩滴滴波呀.
無論你贏定輸佢地都會嗌滴滴波.
一定勝.
無論贏定輸佢地都會嗌ge .
咁呀非常之好氣氛呀.
咁呀所以呢.
係呀教會呢唔同ge 團契.
咁呀gwo 日呢就.
唔知大家呢有冇印象呢lor 過呢一份ye 喎.
咁樣.
有冇lor 過呀.
咁如果你未lor 過呢.
出面應該仲有頭先見到呢仲有少量你可以lor ge .
或者上網下載都得架喎.
咁呢份係乜ye 黎ge 呢原來.
教會裡面呢真係唔同年齡層呢都有唔同ge 團契架.
咁呢係呢一日之前呢就邀請jor 每一個團契.
你諗一樣食物去形容下你gwo 個團契.
就好似一D 咩食物咁呢.
咁呢我就再captureD 圖係呢度好啦.
咁呀有人話個團契呢似菠蘿油呀.
似蘋果.

$^{81}$似咖啡.
呃.
仲有咩呀.
仲有盤菜.
咁呀.
例如呢呢個帕提全味豆咁樣.
咁呢我都唔識ge .
原來係一D 唔同味道ge 糖.
咁呀仲有一D 呢.
可能你睇到佢揀乜ye 呢.
你都估到佢係咩年紀ge 團契.
有人呢就揀苦瓜啦.
清水啦.
白飯啦咁樣.
咁年青D gwo D 可能就係gwo D 糖呀盤菜呀gwo D 咁.
都幾有特色ge 喎.
咁我數一數呢原來裡面呢就講緊我地教會其中11個團契.
4個小組.
佢地係.
ge 一D 介紹喎.
咁呢講到團契生活啦.
咁呀唔知你又會唔會諗到呢一本都耐耐地.
但係呢都幾出名ge 書.
佢個書名呢就係叫做團契生活啦.
咁原來呢本書呢就係德國ge 神學家潘霍華所寫ge .
話說當年呀納粹主義高漲ge 時候.
當時係社會好多意識形態啦.
好多諗法出現jor 啦.
其中包括有反猶太主義啦.
反基督教ge 等等ge 極端思想.
咁當時呢.
德國裡面ge 基督教會.
都會受到影響啦.
咁係.
1933年ge 時候呢當時係教會呢就有一班牧者.
佢地呢就組織jor 一個叫做緊急牧者同盟呀.
就希望呢係呢個咁樣ge 時代呢去力挽狂瀾.
希望呢去維持到教會同埋呢挽回一班年青人呀.
咁當時ge 潘霍華呢.
廿幾歲架喎原來.

$^{121}$佢臨危受命.
當時呢就開辦jor 一間小小ge 地下神學院.
咁呢本書呢原來點解會寫呢.
就係為jor 係神學院裡面.
係咁樣ge 時代背景底下仲有一班年青人.
佢地有一個願意獻上自己生命ge 心.
佢地願意去係呢度.
去受栽培.
一齊去生活去操練希望呢一個群體將來成為一個有力ge 宣教團體.
咁所以呢本書就係講緊當一班信徒佢地聚埋一齊ge 時候.
究竟有D 乜ye 元素應該要係佢地中間呢.
究竟佢地如果要過個敬虔生活.
係要可以有D 乜ye ge 操練呢.
所以呢本書對於今日ge 信徒群體今日ge 團體生活呢依然有好多ge 啟發.
咁呢今年我地提摩太團我地都用周會一系列ge 周會呢一齊去睇呢本書.
咁我地暫時做jor 四次周會啦睇jor 兩個chapter啦.
咁就真係唔好睇佢本書個書名咁容易喎.
我地都搞jor 一輪之後發覺嘩真係其實係神學家寫出黎ge ye 呢.
都係其實都唔係咁容易去明白.
但係呢我地都好希望我地一齊去學習一齊去實踐.
咁今日呢個經文呢我就嘗試用呢兩個大家都熟悉ge 四字詞語.
就係生命之道團契生活去睇一睇今日我地要睇ge 經文.
今日呢我地會睇約翰一書一章一至四節.
用我地先一齊讀一次呢四節ge 經文啦.
論到從起初原有的生命之道.
就是我們所聽見所看見親眼看過親手摸過的.
這生命已經顯現出來.
我們看見了現在又作見證.
把原與附同在.
並且向我們顯現過的那永遠的生命傳揚給你們.
我們把所看見所聽見的傳揚給你們.
為要使你們也與我們有團契.
而我們的團契是與父和他兒子耶穌基督所共有的.
我們把這些事寫給你們使我們的喜樂得以滿足.
約翰一書ge 作者一般都認為就係耶穌十二門徒之一ge 約翰.
佢係西比太ge 兒子.
係聖經裡面約翰除jor 寫jor 呢一個約翰一書之外呢.
仲寫jor 約翰福音.
約翰二三書.
同埋仲有一卷就係啟示錄.

$^{161}$好多人都知道呀.
大家唔知仲記唔記得約翰係一個點樣ge 人啦.
佢同耶穌親唔親近ge 呢.
有人up 頭喇.
係喎係福音書裡面呢.
我地知道呢個約翰.
佢原本呢就係跟隨施洗約翰ge .
咁但係去到有一個時刻.
當施洗約翰見到耶穌.
然後就同佢D 門徒講.
你睇下gwo 個就係神ge 羔羊啦.
於是呢一個約翰呢就開始跟從jor 耶穌喇.
而係十二門徒之中呢.
其中有三個呢就同耶穌就比較親近架喎.
有陣時呢唔會帶住一大隊十二個人一齊出門ge 時候.
或者去做一D ye ge 時候呢.
就會帶呢三個去架喇.
佢地就包括係彼得.
雅各.
約翰非常之好.
冇錯喇就係呢一個約翰喇.
咁所以呢約翰同耶穌真係關係好好好近架喎.
年青時候ge 約翰呢有個綽號叫做.
狠口.
年青時候呢gwo D 人嗌佢做雷子呀.
點解呢即係佢性格好火爆架.
即係喂我做緊一D arm ge ye .
如果你唔同意我就要去懲罰你咁樣.
所以呢即係有一個雷子ge 綽號.
但係呢個約翰隨住佢年紀增長.
佢經歷jor 好多唔同ge 事之後.
呢個成熟ge 約翰唔單止唔再火爆喇.
係佢所寫ge 約翰一二三書裡面呢.
我地更加睇到呢佢被稱為一個愛ge 使徒呀.
因為呢佢好重視作為耶穌基督ge 門徒呀.
有愛係當中呀.
係寫約翰一書ge 時候呢就大概係主後九十年左右.
gwo 陣時呢一個使徒約翰呢已經年老喇.
因為距離耶穌升天已經有幾十年ge 時間喇.
gwo 一D 當時後生ge 時候同佢一齊跟從耶穌.

$^{201}$gwo 個時代一齊親眼見過耶穌ge gwo D 人呢.
已經陸陸續續呢都唔係世界上面喇.
就係咁樣ge 時刻呢個年老ge 約翰.
佢就寫jor 呢幾封ge 書信.
而且呢當時係一個咩ge 時候呢.
當時呢ge 社會上面呢.
亦都充斥住流傳住一D 異端ge 思想呀.
例如呢有人就會宣傳就認為.
耶穌呢唔係一開始就存在架.
佢係某一刻先出現架.
又或者呢有一個異端呢就叫做幻影說呀.
就話呢靈魂呢就係好ge .
肉身呢就係邪惡架喎.
咁所以既然神係好ge .
咁所以呢神係唔會變成一個人.
唔會lor 呢個邪惡ge 肉身架.
所以你其實睇到gwo 個呢唔係真ge .
佢唔會真係道成肉身架咁樣.
咁其實呢D 呢係錯架.
即係呢D 係異端ge 講法喎.
咁但係當呢個社會充斥住呢D .
唔同ge 諗法ge 時候呢.
其實呢信徒係好容易無所適從架.
咁所以我地睇到呢唔同ge 使徒呢.
佢地即係保羅呀或者其他人呀咁樣.
佢地都會就住當時個社會流傳jor 一D 錯ge ye .
佢地就可能會寫信啦.
或者會用佢ge 方法呢去教導返D 門徒.
正確ge 信仰應該係點樣.
咁所以呢呢本書我地睇到呢.
都係針對住當時呢一D ge 錯ge 流傳ge 講法啦.
但係呢本書呢都唔係淨係同你講理論概念架喎.
而係呢好真實咁樣講緊.
基督徒當你實踐信仰ge 時候.
當真係信仰真係係你生命裡面ge 時候.
其實係一個點樣ge 樣黎ge 呢.
咁未來呢今日會開始第一次啦.
咁未來呢ge 主日崇拜我地ge 同工呢.
都會即係一路去陸續同大家去講約翰一二三書ge .
咁我地期待當中ge 寶藏啦.

$^{241}$我地先睇頭兩節啦.
佢話論到從起初原有ge 生命之道.
就係我地所聽見所睇見.
親眼睇過親手摸過.
呢個生命已經顯現出黎啦.
我地睇到啦而家我地又作見證.
將gwo 個原本與父母同在.
並且向我地顯現過ge gwo 個永遠ge 生命傳揚比你地.
論到從起初原有ge 生命之道.
起初 道.
呢D 字眼呢好熟悉喎.
好容易呢就會令到我地去諗起另一卷書就係.
約翰福音一章一節.
無錯啦.
gwo 度一開始都係咁講架.
佢話太初有道.
道與神同在.
道就是神.
如果你一開始接觸基督教呢.
一開始同你睇約翰福音呢.
可能都暈低jor 啦.
好似咁難咁ge 咁樣.
咁呢度D 字眼好接近喎.
譬如佢呢度有太初.
係《楊一書》gwo 度就係講緊起初.
係原文裡面呢其實兩個字係一樣ge .
只不過呢.
係兩個字講緊呢又有少少唔同ge 意思.
係太初ge 時候呢.
呢個字係講緊.
一開始未有呢個世界ge 時候.
就已經有呢個道ge 存在啦.
英文呢就係in the beginning.
in.
咁呢度講緊ge 起初已經有ge 道呢.
起初就係一個世界開始ge 時候.
即係一開始ge 時候就有呢個道架啦.
就係from the beginning.
從一開始ge 時候.
咁呢個道呢.

$^{281}$係說話ge 意思.
英文就係word啦.
說話生命ge 道呢.
就係gwo 個裡面含有生命ge gwo D 說話.
嘩好複雜好抽象係咪呀.
以人ge 腦袋.
以我地肉身能夠感受到ge ye .
我地點能夠明白呢一個含有生命ge 說話.
即係點樣呀.
真係唔明啦.
所以呢.
係約翰福音一章十四節.
神就接住講啦.
呢個道為jor 讓我地能夠明白.
佢成為jor 一個人呀.
成為jor 一個.
咦我地真係聽到佢同我地講緊.
講緊我地ge 說話喎.
我地真係摸到佢真係同我一樣係一個人黎架喎.
呢個道成為jor 一個人.
去展現去表達出黎.
就住係人ge 中間.
有形有體咁樣ge 表達.
所以同接住黎約翰一書講緊ge .
就係gwo 個生命已經顯現出黎.
就係我地所聽得到ge .
佢用人ge 說話.
人話同我講架.
所睇得到ge .
就係佢作為一個人.
佢唔係一個靈體.
佢係一個真係睇得到ge 人.
約翰佢就親眼睇過.
親手摸過ge .
用佢ge 肉眼佢睇得到.
用佢ge 肉手佢可以摸得到.
ge 一個活生生ge 人.
係當時黎講.
所以呢.
就係要推翻.

$^{321}$唔係幻影黎架.
係真架.
呢個生命已經展現jor 出黎.
就係耶穌.
我地睇見.
而家亦都去作見證.
咁呢個生命之道.
其實係指緊乜ye 呢.
成為jor 肉身ge 人.
呢個生命ge 道.
呢個人一定係講緊.
耶穌.
呢個生命之道.
一定係講緊耶穌.
咁不過呢.
當人去仔細D 去睇聖經ge 時候呢.
就會留意到一樣ye 喎.
就係有一個代名詞.
當我地講緊一個人.
或者一樣ye .
咁呢.
再講ge 時候就會有一個代名詞.
例如我地ge 詩琴.
佢係姊妹黎ge .
咁所以呢.
佢ge 代名詞.
下次再講佢ge 時候呢.
就係一個他.
係女字邊gwo 個他啦.
係咪.
咁呢.
另外如果執事.
佢係男ge .
咁我地下一次再講佢ge 時候呢.
咁我地就用男性ge 他啦嘛.
係咪.
咁如果呢個生命之道.
係淨係講耶穌呢個人呢.
再講佢ge 時候.
佢個代名詞就應該係男性ge 他.

$^{361}$係咪.
咁就.
咁就好簡單啦.
係咪.
咁就講緊.
啊個生命之道咪就係他囉.
就係.
gwo 個耶穌囉.
咁不過呢.
啊.
仔細D 睇ge 時候呢.
即係.
人呢又睇得出.
咦.
係用jor 中性ge 他喎.
點解會咁ge .
點解唔係淨係講耶穌.
直接講耶穌ge .
咁原來呢.
啊佢地.
即係研究完一輪啦.
咁我淨係講個結論啦.
就係哦.
佢地相信呢.
呢個唔係淨係講耶穌呢個人呀.
呢個道唔係淨係gwo 個人.
呢個道係包括jor 佢個人.
同埋佢所做ge 一切.
成件事.
佢所做過ge 一切.
咁所以.
呢一個耶穌.
啊.
約翰gwo 個親手.
摸得到.
睇得到ge 耶穌.
其實係二千年前.
淨係係二千幾年前ge .
gwo 一個地區ge 人.
先至有機會去睇得到.

$^{401}$去摸得到.
係咪.
咁但係呢.
耶穌佢ge 存在.
同埋佢ge 工作.
冇停係二千幾年前ge gwo 個地區.
而係一直去到而家.
去到唔同ge 地方.
耶穌ge 存在.
同埋佢所做ge ye .
都依然一路發生緊.
耶穌復活升天之後.
聖靈降臨之後ge .
每一代ge 信徒.
就係咁樣去經歷.
啊當我地信jor 耶穌之後.
呢個生命ge 改變呢.
好奇妙架.
有D ye 唔係人.
可以解釋得到.
唔係話我幫你做一D 出黎.
就可以做得到ge 事.
呢D 係每個信徒.
能夠經歷得到.
活生生ge 見證.
相信我地每一個人呢.
我地多多少少都聽過.
唔同ge 基督徒ge 見證.
基督徒就會分享下.
啊佢係點樣開始認識耶穌ge 呢.
佢點解決定信耶穌呢.
信jor 耶穌之後.
生命經歷jor D 乜ye 呢.
其實呢有冇諗過呢.
都好特別架喎.
雖然我地好自然.
成日都聽到呢D .
但當你諗諗下ge 時候.
其實一般人.
唔係咁容易.

$^{441}$將自己gwo D ye .
攤出黎比人睇.
去講比人聽架.
係咪呀.
有D 可能你同佢同事.
係同學jor 一段唔短ge 時間.
間中都見下面.
你又唔可以話同佢唔friend.
但係呢.
硬係就.
其實你地係唔會講一D .
好深入ge ye .
你唔會將自己.
一D 好軟弱ge 地方.
去講出黎架嘛.
係咪呀.
咁呢個好正常.
又或者呢.
其實你有D 叻ye 呢.
咁除非gwo 個人真係.
好想得到別人ge 稱讚啦.
如果唔係.
你都唔會特登將一D .
自己叻ge ye .
係數比人聽架嘛.
係咪呀.
咁但係呢.
基督徒.
基督徒又好特別架喎.
基督徒好似唔介意.
好樂意.
有機會.
就會分享.
甚至呢.
同gwo D 完全唔識ge 人呢.
佢都可以話比你聽.
啊 我係一個點樣ge 人 我經歷jor D 咩事.
然後之後發生咩事.
因為基督徒想講ge .
唔係佢自己呀.

$^{481}$係咪呀.
我講比你聽我自己ge ye .
但係個主角唔係我呀.
個主角我要話比你聽.
耶穌係我身上做jor D 乜ye 呀.
你同我ge 經歷唔會一模一樣架.
如果我病.
啊 今次好返.
係咪代表所有聽到我講呢個見證ge 人.
做同一樣ge ye .
就一定會有同樣ge 結局呢.
唔係.
但係.
我ge 分享就係要話比你聽.
啊 你可以經歷同一位神架.
用你ge 方法.
按照你ge 狀態.
神就會比你經歷到佢.
咁所以呢個係.
基督徒點解會肯講見證.
會肯公開咁將自己gwo D 衰ye .
咁講晒比人知道ge .
咁呢唔知你有冇聽jor 幾多見證呢.
你有冇諗起呢.
其實呢 我曾經都聽過唔少ge 見證.
曾經呢都有D 見證呢.
好幫助到我.
可能係某一次ge 聚會.
或者某一個你認識ge 人.
你見到佢信耶穌.
你睇到佢所講ge ye .
佢所做ge ye .
你睇到佢ge 生命經歷.
吸引到你.
都想去認識o下呢個耶穌.
又或者呢.
你聽到呢一D 經歷.
你未必有一模一樣ge 經歷.
但係你會覺得.
原來耶穌係咁樣ge .

$^{521}$原來耶穌可以咁樣去做事係人ge 身上.
呢D 見證呢.
都建立緊你對耶穌ge 信心.
有冇諗起呀.
可以諗o下.
曾經邊一個人.
佢同我分享過佢ge 生命.
曾經邊一個我見證到佢ge 生命改變.
佢幫助到我.
去更加親近耶穌.
會唔會係你聽到罪人.
能夠悔改ge 經歷.
讓你去認識到.
耶穌ge 寬恕.
同埋佢個改變生命ge 能力.
會唔會係病得到醫治ge 見證.
讓你知道原來耶穌係一位最大ge 醫生呀.
會唔會係係充滿困難.
挑戰ge 生活裡面.
gwo 個難處都好大架啦.
但係gwo 個人點解都可以.
咁安定咁安心.
有咁大ge 平安.
佢仲笑得出.
佢甚至可以鼓勵佢身邊ge 人.
啊.
係呢度呢.
睇到耶穌佢真係超越D 環境.
佢可以似人地心底裡面最深gwo 份平安.
或者當你見到.
有D 人呢.
佢可以為jor 別人ge 好處.
寧願甘心樂意放低自己ge 權利.
或者放低自己ge 舒適.
去為別人去做一D ye .
你就見到佢信ge 呢個耶穌.
係一個願意去為別人去犧牲自己.
去愛人ge 神喎.
又或者當有D 人呢.
佢尋求人生方向ge 時候.

$^{561}$佢呢.
去問神架喎.
佢會去問.
求問佢所相信ge 耶穌基督.
然後問下啊神.
你想我.
我人生點樣用呢.
我應該行去乜ye ge 方向呢.
我黎緊應唔應該做呢一樣ye 啦.
咦 原來呢.
神係一個好願意人得著豐盛生命.
並且好想去指引人.
去行一條正路ge 神啊.
又當你聽到呢.
原來有一D 係比較困難ge 地方.
佢地要成為一個基督徒.
原來都要受到好多ge 風險危險.
可能有機會會比人拉.
會比人捉jor 去打佢.
或者甚至冇埋條命.
但係點解gwo D 人.
都仲可以.
喂 你講一句.
我唔係基督徒我唔識佢架咁樣.
好簡單就得架喇.
但係gwo D 人點解.
寧願受罰受苦.
佢地都繼續堅持ge 呢.
你就發現.
咦.
耶穌ge 名.
甚至可以比生命更加寶貴.
更加值得持守.
好多好多呢D 啊.
當我地一路去睇下去聽下ge 時候.
我地發覺.
嘩 好多真實ge 經歷啊.
佢地唔係同你講道理啊.
佢地係話比你聽.
咦.

$^{601}$係我生命裡面我真係咁樣經歷緊.
無論佢係乜ye 國籍 乜ye 年紀.
乜ye 家庭背景.
乜ye ge 學歷 乜ye ge 際遇.
原來呢都可以親身經歷耶穌架.
呢D 活生生ge 見證呢.
讓我地去認識神.
呢個生命之道.
係耶穌同埋佢ge 作為啊.
呢D 生命之道呢.
唔係抽象ge 概念.
唔係一定要讀書先至可以得到ge 東西.
而係gwo 個好真實ge 神.
就住係我地中間.
係我地日常生活裡面係可以經歷得到.
就係咁奇妙啦.
接住落黎第三第四節.
佢話我地將所睇到.
所聽到ge 傳比你地.
為要使你地同我地有團契.
而我地ge 團契係同父同埋佢.
兒子耶穌基督所共有ge .
我地將呢D 事寫比你地.
使我地ge 喜樂得到滿足.
約翰跟住就講啦.
點解我地要將我地ge 經歷.
我地ge 見證去同你講呢.
點解要傳開去呢.
無論口講或者係行為上面ge 傳遞.
原來帶住一個目的.
帶住一個期望一個原因.
就係想我地同你地之間可以有個團契.
而且呢個團契唔係淨係D 人聚集.
唔係淨係人同人之間ge 交流.
呢個團契係有父神同埋佢兒子耶穌.
一齊係裡面ge 團契.
跟住相信同一位神.
去分享神係我地生命裡面ge 工作.
所以分享緊ge .
係黎自同一位神ge 贖命經歷.

$^{641}$同一個使命.
同樣係黎自神ge 福氣.
咁樣ge 團契呢.
原來係可以帶黎滿足喜樂ge .
我相信呢個喜樂呢.
唔係話我返jor 團契之後呢.
我就見到gwo D 人真係好好人啦.
佢地對我好好好關心成日送ye 比我.
咁咪好開心囉係咪呀.
我相信呢都應該好大機會都唔係咁樣係咪呀.
好多時候呢我地返團契.
返團契就發覺.
我又有D 唔可愛我都知道.
你我又見到一D 唔可愛ge 地方.
就好似上個禮拜講完就同我地講啦.
如果我地要搵問題ge 話呢.
其實我地搵到無限咁多.
呢個團契D 人又咁樣ge .
你又叫團契職員但係你咁樣ge .
如果你搵問題呢有無限咁多.
所以點樣喜樂呀.
呢個喜樂呢唔係因為.
我係團契當中見到gwo D 人有幾好.
或者佢地對我有幾好所以我開心.
呢個喜樂呢係係我地係團契裡面見到神.
佢係當中點樣仍然愛我呢個唔可愛ge 人.
然後神點樣係我地當中去叫我地.
因著呢一個聚集去一齊去成長.
見到神ge 作為見到神ge 恩典.
所以喜樂gwo D 問題未必解決到.
大家都仲棘棘地不過.
如果你見到神係當中呢個就可以喜樂啦.
你有冇同父神同耶穌同信徒一齊經歷過.
一D 團契呢.
你有冇感受過團契當中ge 喜樂呢.
如果我呢度括住話可唔可以分享下.
咁我係呢度冇打算叫大家分享.
因為我恐怕大家陣間靜jor 五分鐘我就唔知點算.
咁如果真係問你你有冇經歷過團契呀.
你有冇經歷過當中ge 喜樂呀.

$^{681}$你可唔可以同佳麗分享下如果一陣而家唔使講.
你會唔會好我諗到呀我有呀我有經歷過.
然後呢我其實都好好想呀可以同人分享呀咁樣.
咁可能呢有唔同ge 情況呀有D 人呢就會話啦.
哎呀我冇返團契喎咁樣.
咁一方面呢就係我地頭先講呀嘛.
呢個團契ge 意思係乜ye 呀.
就係同主同人有屬靈上面ge 交流交往.
所以呢一種交往呢不限於有一種叫改jor 個名呀.
呢個群體我比個名佢啦叫嘟嘟嘟團.
不限於一個團契ge 一個模式.
所以其實當我地能夠係其他時候任何時候.
打電話又得whatsapp又得見面飲茶又得.
當我地能夠一齊去分享下係主裡面gwo 個經歷.
其實呢D 團呢D 咪就係團契囉.
當我地可能陣間完jor 送拜之後你落到去赴堂.
或者食飯ge 時候你就會諗起.
今日聽到ge 時候呢.
講完就問我有冇聽過一D 經歷呢好激勵我架.
有呀我記返起喇曾經某一個人.
佢曾經呢佢ge 生命經歷呢激勵jor 我幫助jor 我.
兼顧jor 我好多呀咁樣.
你可以同人分享下甚至呢.
你諗起佢你親自去同佢分享下.
講聲多謝當年你咁樣兼顧jor 我.
呢D 咁樣ge 分享交流其實就係團契啦.
又或者呢呀最近呢經歷jor .
一D 事但當我睇聖經ge 時候呢.
我發現jor 有一D ge 經文呢佢好似回應緊我.
佢比到我一D ge 安慰啦指引啦.
教導啦甚至責備緊我呢.
當我心裡面有呢D ge 感受ge 時候.
都可以同人分享架喎.
呢D 呢都係屬靈ge 交流.
又或者呢你面對緊正正面對緊一個處境.
大家都未諗到點解決亦都幫唔到你.
不過呢大家都一齊將佢擺係肚角裡面.
去為你祈禱.
其實呢呢個都係團契ge 喜樂.
咁所以呢你話咁呀唔使返團契.

$^{721}$我有呢D 咪得囉係咪呀.
咁但係呢其實我地仍然鼓勵呢每一個弟兄姊妹呢.
你去加入一個團契或者小組ge .
正如呢有人就會講啦.
其實呢我唔返學都可以學ye 架係咪.
我有興趣學gwo 樣ye .
我自己會上網睇架啦.
跟住呢我又好自律ge .
我到時到後我就會去做啦會去練習啦.
就算我冇一個團契呢.
我都會得閒約下我D friend.
就同佢分享下我ge 信仰經歷.
咁如果呢真係真係咁樣做得到呢.
真係非常之好我真係好欣賞你.
即係唔返學都學到ye ge .
咁係好好啦係咪呀.
如果係咁呢我更加希望呢邀請你去返一個團契.
因為呢你咁健康係你咁有呢D 團契ge 意識呢.
其實你都可以黎祝福團契裡面ge 其他人.
咁又有人會話啦.
我有返團契.
所以就唔好再問我啦.
我有返團契呀係咪呀.
咁不過呢我都曾經同人傾偈ge 時候呢.
有一D 對話啦.
咁gwo 個人呢佢有穩定咁樣返團契ge gwo 陣時.
甚至係團契職員jor 一段時間.
咁我就同佢傾偈gwo 陣時就問佢啦.
我就話你過去做jor 咁耐職員呢.
或者返jor 咁耐團契啦.
你覺得係團契裡面有D 咩成長呀.
同神ge 關係有D 咩得著呀咁樣.
咁對方呢好快就同我講.
無呀無呀.
咁我就o下.
唔會無ge 諗o下點會無呀.
o下點會無呀真係無.
你返jor 團契.
我諗你花ge 時間係咪呀.
你花ge 時間你做埋職員你ge 心機.

$^{761}$你可能要付出錢係咪呀.
喂.
好.
點會無呀好可惜喎如果真係無ge 話咁樣.
咁唔知呢你而家係一個咩ge 景況啦.
如果你係有返一個叫做團契ge 單位.
而你又真係經歷到當中ge 團契.
同神同人gwo 種ge 交往.
你都經歷到gwo 份喜悅呢.
如果你冇返緊一個團契.
但係你仍然有唔同ge 機會你去關心.
同弟兄姊妹有交流有團契.
咁樣都係非常之好.
如果可以ge 話呢.
我都鼓勵你加入一個團契.
去祝福其他人.
如果你有返團契.
但係呢其實你冇經歷得到團契gwo 種喜樂.
你冇真正去覺得我有D 信仰ge 經歷.
好想同人分享.
我亦都係裡面聽到別人ge 經歷.
如果你冇感受到呢份喜樂.
已經維持jor 一段時間啦咁樣.
咁呢我都鼓勵你去搵一個信任ge 團友.
或者導師同佢傾下計.
究竟係我自己有D 盲點.
去令到我限制jor 我好似.
我係有做緊gwo 樣ye .
但係好似得著唔到當中ge 喜樂.
定係其實有D ye 未必完全關你事架喎.
其實可能係個團契.
呢班人真係有D ye 要正視ge .
真係有D ye 要處理ge .
提出黎一齊去諗o下.
係咪有D ye 係團契做得唔好.
我地去犧牲jor .
即係失去jor 好多可以去相交ge 時光呢.
咁所以呢.
即鼓勵你如果係咁樣ge 話呢.
真係搵個人一齊傾下一齊去祈禱.

$^{801}$咁如果呢你話我兩樣都冇.
我既冇返一個團契.
我亦都冇同其他人有團契ge 生活.
咁呢.
我亦都鼓勵你去諗一諗.
先同神傾下計先.
因為可能你都唔知搵邊個傾.
你同神傾下計先.
啊神啊我而家呢又冇返團契.
不過呢講到團契咁好.
團契D 人又唔係咁可愛.
但係呢你都係裡面會經歷到喜樂架喎.
講到咁好.
咁但係我又冇返.
可能又有D 障礙原因.
或者有D 事發生過.
咁但係其實我應該點樣前進呢.
現階段ge 我可以點樣前進呢.
同神傾下計.
如果可以ge 話呢.
可以搵弟兄姊妹或者屬靈長輩去傾下計.
咁呢係信仰上面呢.
我地成日睇到呢個三角形.
呢個三角形呢.
就係有一邊上面啦就係神.
一邊就係我啦.
另外一邊就係別人.
係呢個三角形裡面呢.
好似代表緊一個關係ge 距離呀.
如果我同神近D .
如果對方同神又近D ge 時候呢.
咁會點呀.
我同別人ge 關係呢.
都會一齊拉近D .
因為其實呢我地向住同一個方向行呀.
我地有住同一個目標啦同一個意象同一個使命.
我地分享緊同一個福氣.
咁樣呢我地ge 關係都會近D .
上個禮拜四晚呢.
我地就開jor 一個工作小組ge 會議啦.

$^{841}$咁呢個小組呢叫輔唐小組呀.
咁唔知你知唔知道係教會有個咁樣ge 工作小組啦.
但唔緊要ge .
咁呢我就係呢度呢介紹一下.
咁起源係點解會有個咁ge 小組ge 呢.
輔唐就係呢度樓下.
即係一入門口gwo 個空間我地叫輔唐啦.
咁起源呢就係咁ge .
話說呢幾年前呢我地開始留意得到呢.
香港ge 各個教會就係多jor 關注.
教會同社區ge 關係.
包括點樣運用.
可以更加有效咁樣運用教會已有ge 場地去實踐使命.
教會呢.
唔同ge 教會多多少少你都有D 地方架係咪呀.
你都要坐係度你都要有一定ge 地方架嘛.
咁但係呢通常最忙ge 時候呢就係星期日.
朝早啦或者星期六啦或者星期日啦咁樣.
咁但係呢平日一至五呢.
好多時候教會都係拍蒼蠅ge 咁樣啦.
咁我曾經有一次呢.
我就企係個樓下個接待處gwo 度啦望住gwo 個.
CCTV啦gwo 個閉路電視啦.
跟住我發覺點解每一格佢影住唔同樓層.
佢就發覺點解每一格都冇人ge .
點解全個教會都冇人ge .
即係平日返工ge 時候.
咁我就發覺係呀.
即係教會ge 地方點樣可以更加好咁樣用黎實踐使命呢.
咁另外呢由開始呢聽到有一個.
詞語啦叫做第三空間.
The third place 呀咁樣.
咁佢呢就講緊呢.
啊呢個世界呢有唔同ge 空間.
第一個空間呢就係你ge 屋企.
即係你成日都係gwo 度住ge 地方啦.
訓ge 地方啦食飯啦訓覺啦.
第二個地方呢就係你ge .
學校或者你工作ge 地方啦.
係gwo 度呢你要打醒十二分精神.

$^{881}$gwo 度係一個薄殺ge 地方.
第三個空間呢.
就係一個呢.
此外除兩個之外呢一個可以比你放鬆.
去思考.
去同人交流ge 一個場所.
咁唔一定係講緊教會ge .
佢就係講緊呢可能係一D 茶樓呀.
甚至呢有D 民族呢gwo D 非法舖都係一個空間.
去gwo 度講o下事非.
傾o下計呀.
咁樣ge 一個放鬆交流ge 地方.
咁呢亦都.
即gwo 段時間亦都適逢疫情啦.
疫情呢帶黎jor 好多ge 隔離呀.
但係呢疫情亦都創造jor 好多新ge 機會.
讓人可以連結.
即係例如疫情ge 時候呢.
即係我都有試過呢去幫人買sung .
其實我自己都唔識買sung .
我都唔使買sung ge 係咪.
咁但係呢疫情ge 時候呢.
我又可以幫人買sung 喎.
平時真係唔使理啦係咪.
我買gwo D 更加好啦咁樣.
咁所以係疫情ge 時候.
又刺激o下我地諗o下.
其實有D 新ge 方式係可以同人連結.
咁所以講jor 咁耐.
就係點解呢.
跟住我地就教會ge 見證報道組呢.
就開始諗o下我地ge 輔堂.
可以點樣更好咁樣去運用.
於是呢我地就成立jor 呢個輔堂小組啦.
咁如果你返jor 教會至少有幾年時間呢.
你應該都會留意得到呀.
以前我地ge 輔堂呢.
係冇聚會ge 時候呢.
其實係黑蒙蒙ge .
凳都冇一張呀因為唔.

$^{921}$唔預設你會入去gwo 度坐或者活動ge .
咁但係呢.
而家已經唔痛哂啦.
咁我地話唔讓輔堂成為教會ge 客廳.
呢個客廳係點樣ge 呢.
我地中途呢就嘗試呀.
都係慢慢嘗試啦.
即係搬D 梳化呀擺入去呀.
不如整D 小擺設等佢好舒適咁樣啦.
不如呢擺D 雜誌係度俾人睇啦.
之前我地都嘗試過呢報下音樂啦.
但係發覺後來.
係算啦都係冇報啦咁樣.
咁所以我地做過唔同ge 嘗試亦都使用呢個空間呢去有唔同ge 活動.
譬如呢係gwo 度易剪啦幫人剪頭髮啦係咪呀.
或者呢有D 茶座呢.
歡迎出面ge 人街坊朋友黎聽下D 主題ge 講座啦.
之前呢我地都試過搞市集啦.
你唔係返開教會唔緊要呀.
黎我地教會睇下ye 啦咁樣.
成為一個呢.
我地自己ge 人就可以使用去相聚ge 空間.
成為一個呢.
外邊ge 人你ge 親友呢都可以入黎.
你一個接待ge 空間.
咁所以呢.
呢個地方我地好希望係裡面呢讓人.
無論你信jor 住未信住.
無論你係唔係我地教會ge 會友.
黎到呢個空間呢可以係度認識神.
可以感受到神未必一定直接同你講.
但係呢有神係裡面流動.
有人係裡面.
有人ge 愛係裡面流動.
咁呢講返呢個會啦.
係呢個會上面呢其中呢我地有一個時刻呢.
傾傾下ge 時候呢就有D 驚喜喎.
點解呢咁就係.
咁呢就係有姐妹呢執ye ge 時候呀.
係教會D 櫃gwo 度執ye ge 時候.

$^{961}$咁呢就執出jor 有D 呢.
二十年前ge ye 呀D 寶物呀咁樣.
咁呢好新淨呀一箱ye 就係呢個.
寶勒罕呀.
有冇人見過呢本ye 呢.
未出世gwo D 笑呀.
未出世呀二十年前.
咁呢2007年出ge 原來.
咁呀寶勒寶勒罕.
咁呢佢原來唔單止.
佢有好多本ge 原來.
係喇.
咁係裡面呢我地揭下揭下啦.
我地就發覺嘩.
即係越揭越開心呀就係呢.
我地揭到呢好多我地熟悉ge 面孔呢.
二十年前個樣喎咁樣.
咁咁就真係好開心啦.
咁其中呢其中呢本呀.
咁呢一揭我一揭呢就揭到.
在座ge 你同埋你啦咁樣.
就發覺.
gwo 陣時呢我當然唔識佢啦.
但係呢.
嘩當年呢個講緊呢個家庭.
生jor 個BB.
個BB呢要做手術.
嘩作為媽媽呢爸爸媽媽呢.
係咪呀幾咁掛心啦.
即係跟住就話唔知個BB將來.
呢個手術會點影響佢ge 成長呢.
咁樣今日健健康康咁樣係度啦.
咁哇呢D 見證呢好開心呀.
呢D 見證呢唔係睇緊D 道理神學書呀.
呢見證係話比你聽.
喂信耶穌呢日常生活經歷jor D 咁ge ye .
生命裡面經歷jor D 咁ge ye .
但係耶穌點樣帶住佢地經過.
我相信呢呢廿年係咪呀.
好多好多恩典呀仲有排講呀係咪呀.

$^{1001}$咁所以呢我地搵返呢D 寶物之後呢.
跟住我地就話喂擺呢D 出黎啦.
擺係度俾人睇啦係咪呀.
讓呢個空間有人黎到坐低ge 時候.
咪可以睇下究竟咩叫信耶穌.
信耶穌其實經歷緊D 乜ye .
咁所以呢度就唔係好多je 係咪呀.
所以呢就再諗下點用啦或者再整D 出黎啦.
咁但係呢我地好期待好期待呢個地方呢.
可以讓人去睇到生命之道呀.
睇下耶穌同埋耶穌係度做緊D 乜ye .
係呢度呢讓人有團契生活呀.
可以讓人呢係裡面分享.
神係佢身上做過D 乜ye .
無論透過係刊物啦.
透過人同人之間傾偈啦.
透過呢你未必直接同個人講ye .
但係你係個空間gwo 度呢.
你係度同人一個微笑啦一個互動.
其實身邊D 人都感受得到睇得到.
所以今日呢一個生命之道呀.
呢一個約翰好期待.
係信徒群體當中有ge 呢種團契生活.
就係我地一齊去分享下我地屬靈裡面ge 經歷.
一齊去分享一齊去領受使命.
一齊去分享神ge 福氣ge 呢個團契生活呀.
我地今日仍然係繼續咁樣經歷緊.
願我地都能夠去真係實踐睇我地ge 教會好唔好呀.
即係我地呢班信徒.
真係一齊去有耶穌係我地當中.
一齊去經歷團契ge 喜樂.
雖然呢問題依然會有.
困難仍然未消失.
但係呢我地因為見到主仍然係當中做緊ye .
仍然去愛我地.
仍然去帶住我地向前行.
我地為此而喜樂而感恩.
我地一齊祈禱.
秋夜天父上帝我地好多謝你.
今日呢透過約翰一書.

$^{1041}$使徒約翰佢所寫ge .
寫出佢心裡面一D 好想好想同信徒講ge ye .
讓佢呢去親身見證.
gwo 位佢曾經親手摸過親眼見過ge 耶穌.
佢講緊呢係信徒呀.
你地都可以繼續去經歷呢一位耶穌呀.
今日呢我地睇到呢好多好多ge 見證.
好似雲彩咁樣包圍住我地.
話緊比我地聽信耶穌唔係D 好深奧ge 道理呀.
唔係要讀幾多書.
信耶穌每一個人只要你願意將你ge 生命去交比耶穌.
你願意信服.
問下耶穌.
耶穌話今日我點樣可以去更加親近你.
去認識你將我生命去交比你.
由你帶住我去行啦.
只要我地願意ge 時候呢奇妙ge 事係會發生.
gwo D 我地生命ge 改變.
gwo 個經歷呢就係咁樣去發生.
主約願你都係我地ge 教會係我地ge 群體當中.
唔單止係使用輔堂使用我地整間教會.
亦都使用我地每一個.
因為當我地行出去ge 時候.
我地就係帶住呢個生命之道.
當我地行出去去分享耶穌ge 恩典ge 時候呢.
呢個就係團契生活.
讓我地呢將佢亦都帶進去我地ge 社區.
我地ge 生活當中.
讓神去得榮耀.
讓人去得著益處.
我地多謝你.
如果我地當中有一D ge 障礙.
有一D 盲點係我地未做到.
我地唔願意或者我地睇唔到ge .
求你都去提醒我地.
幫助我地.
讓我地呢都去一步一步咁樣去向前行.
更加去親近你.
經歷到係你裡面gwo 份團契ge 喜悅.
我地多謝你.

$^{1081}$聽我地祈禱.
奉主耶穌名求.
阿們.
\newpage



\section{創世記 4:1-7}
\label{sec:032448}
\textbf{討主喜悅}~~~~林羅麗斯牧師
\newline
\newline
連結: \href{http://www.church.com.hk/acms/content.asp?site=cdc\&op=show\&type=product\&code=032448\&layout=sermon}{\texttt{032448}} ~~~~ 語音日期: 2026-07-12
\newline
\newline
\hyperref[sec:032447]{\small{< < < PREV SERMON < < <}}
~
\hyperref[sec:index_chronic]{\small{[返順時目]}}
~
\hyperref[sec:032449]{\small{> > > NEXT SERMON > > >}}
\newline
\newline
創世記 4:1-7
\newline
\begin{longtable}{cl}
\hline
\hline
章節 & 經文 (和合本修訂版)\\
\hline
4:1 & \begin{tabularx}{0.7\textwidth}{X} 那人和他妻子夏娃同房，夏娃就懷孕，生了該隱 ，她說：「我靠耶和華得了一個男的。」 \end{tabularx} \\ \\ \relax
4:2 & \begin{tabularx}{0.7\textwidth}{X} 她又生了該隱的弟弟亞伯。亞伯是牧羊的；該隱是耕地的。 \end{tabularx} \\ \\ \relax
4:3 & \begin{tabularx}{0.7\textwidth}{X} 過了一些日子，該隱拿地裡的出產為供物獻給耶和華； \end{tabularx} \\ \\ \relax
4:4 & \begin{tabularx}{0.7\textwidth}{X} 亞伯也把他羊群中頭生的和羊的脂肪獻上。耶和華看中了亞伯和他的供物， \end{tabularx} \\ \\ \relax
4:5 & \begin{tabularx}{0.7\textwidth}{X} 卻看不中該隱和他的供物。該隱就非常生氣，沉下臉來。 \end{tabularx} \\ \\ \relax
4:6 & \begin{tabularx}{0.7\textwidth}{X} 耶和華對該隱說：「你為甚麼生氣呢？你為甚麼沉下臉來呢？ \end{tabularx} \\ \\ \relax
4:7 & \begin{tabularx}{0.7\textwidth}{X} 你若做得對，豈不仰起頭來嗎？你若做得不對，罪就伏在門前。它想要控制你，你卻要制伏它。」 \end{tabularx} \\ \\
[1ex]
\hline
\hline
\end{longtable}
$^{1}$很快就踏進了七月.
我們進入了這個年度的第三季.
所以這個季題就叫做「土主喜悅」.
經文其實是記載在羅馬書的十四章.
十七到十八節.
但是又記得梁家倫博士.
剛好在六月底也說了.
說到已經很詳細的將這一段.
我們的主題經文去釋經來講.
相信你們姐妹都應該記得清楚.
怎樣能夠是上帝喜歡的東西呢?.
其實我們能夠真的常常活在公義和平.
並聖靈當中的喜樂.
這些的服侍就已經是神所喜悅了.
因此我盼望在這一季我們真的努力學習.
當我們聚會 當我們相見的時間.
我們在一起的時間.
大家彼此勉勵去追求和睦的事.
和彼此去建立得恆的事.
好讓我們過上帝喜歡的生活.
所以今天我就調了另一個態度.
來看看這個主題.
因為昨天突然有機會在青崇實體講道.
因為同青年人他們發覺.
這些古老的字眼.
「土主喜悅」,「土人喜歡」.
這一類跟我們上一代的理解不是很相似.
我們很casual.
「做些事聽媽媽開心」.
「老闆喜歡」.
我們就好過一點.
我們就沒有什麼.
但是現在這一代對於這些.
好像有一點委屈自己.
來哄別人開心.
他們是很難的.
所以你看到他們畫的東西是很清楚的.
「土人喜歡」.
這些是撒旦的手段.
意思是我要扭曲自己.

$^{41}$然後迎合對方做的心意.
這些不是很好.
所以這次我不知道你突然閒起來.
這樣的心態.
看見這個主題的時候.
你會不會有些很負能量的想法.
又要我們做些什麼.
又要逼迫自己去做一些上帝喜歡做的事.
或者教會喜歡做的事.
所以我想在這個主日.
我們先打好那個基礎.
否則整個季度所講的東西.
我們就有很多抗拒.
但是無論你覺得.
是不是撒旦的手段都好.
不過問真一句.
你覺得你自己喜不喜歡別人喜歡你.
嘩.
這麼冷酷的你們.
沒有表情的.
在想是不是.
不知道怎麼回答是不是.
那你喜歡自己走去哪裡是討人厭的嗎.
是不是.
就是不會的.
很正常的.
你去到哪裡.
別人都歡迎你.
別人都喜歡.
你覺得你好的.
這個都是我們追求生活的態度.
所以是不是真的這麼差.
我就覺得大家先回到平衡的位置.
我們從小在教會.
都是用這些terms.
就是說你們要做土著希越.
特別很多時候.
真的.
我們想得別人的歡心.
是要花心思.

$^{81}$還有一點點委屈自己的.
因此木井就很流行這樣教我們的.
就是說你不要專注在別人的讚賞裡面.
最重要的是什麼.
就是上帝喜歡你.
上帝得到榮耀.
就可以了.
做了幾十年基督徒.
退休了.
你問我.
是不是真的可以的呢.
我們當中有很多老餅的基督徒.
我們都知道上帝要怎麼做.
應該怎麼做.
我們應該是一個怎麼生活.
怎麼上班.
怎麼去做人老婆.
很多的東西都學了.
但是不是真的說一句.
上帝喜歡我就可以了.
我不用管我老公喜不喜歡我.
我不用管我老闆喜不喜歡我.
是不是真的可以這麼豁達呢.
因此我想說.
我們不要說一些很負能量的話.
問心在心底的裡面.
就算我們生活在世間.
我們都想人喜歡自己.
同樣我們在上帝的面前.
將來我們還說要回天家.
去交帳.
我們每一個人都期望.
其實都可以得到主的讚賞.
因此我們都有一點點壓力.
其實我在人世間的生活.
我現在做一個基督徒.
在教會.
有很多東西又叫我們去釋放.
又叫我們傳福音.
是不是會有很大的所謂的壓力呢.

$^{121}$特別我們當中在疫情之後的新一代.
不知道為什麼個個都說自己是哀人的.
就是說全部都是啞的.
不太喜歡見人的.
不喜歡跟人打招呼的那些.
是不是真的這樣呢.
所以我教一下那些父母.
現在我經常口頭禪叫他們.
請你functionally.
我不用管你的天生是什麼都好.
請你現在的表現.
可以warm一點 主動一點.
是的.
不同類型的人有不同的性格.
我們說有些小朋友.
或者是就算大人也好.
其實是很怕見光的.
很怕被別人highlight.
所以你說台上的事奉 做組長.
總之有title 他們都不做的.
不給名字就無所謂的.
什麼都做的.
跟我們那一代不同的.
我們以前那一代.
當然是覺得名正言順的.
我沒有這個title.
我就不做這些事.
現在反過來.
他們做完所有事都可以.
但是就不要title.
背後是因為什麼問題呢.
昨天在清崇很好笑的.
原來自殘凝聖是他們不知道的一個說話.
我們當中這裡應該都知道的.
我們應該不是那麼年輕.
某程度來說.
有些人為什麼會自信心不足.
或者是很害羞.
或者很怕別人提他的名字.
或者highlight.

$^{161}$背後可能是責任感的問題.
他的要求很高.
留下title.
他覺得他要做得很完美.
所以就覺得自己不行.
我做不到的.
我沒有你那麼好.
或者我沒有以前的組長那麼厲害.
我不行的.
所以就很怕他take up一些東西.
或者有一些比較極端的.
可能他曾經是很有自信的.
他是很樂意出來的.
但是經驗不是很好.
於是就怕了.
不想再有任何隔出來的機會.
無論我們是什麼狀態.
或者是什麼心態都好.
面對著我後面這個所謂使命的人生.
今年對我們弟兄姊妹來說都挺多挑戰的.
因為在小組討論裡面都覺得.
大部分好像很迷茫的.
不知道為什麼.
我使命是什麼呢.
我應該做些什麼呢.
弟兄姊妹.
我覺得我們不要經常上了天就忘記了落地.
我們現在都是一個人.
跟隔壁的人說我是人.
不要想得那麼複雜.
每一個人出生.
他其實都會有他的人生的.
有使命的意思就是說.
你有方向有目標.
你就走得清楚一點.
向那個方向努力.
如果你說真的長期都沒有.
其實都是挺納悶的.
人家叫你做什麼你就做什麼.
做又不知道喜不喜歡做.

$^{201}$自己想不想做.
因此作為一個基督徒.
我們已經是跟隨耶穌的人.
所以耶穌基督是我們的師父.
師父做什麼.
徒弟是應該做什麼的.
你跟一個木匠.
你當然是學做一個木工.
你去跟一個做壽司的.
你都期望有機會拿一把刀割一塊肉的.
因此當我們說我們是基督徒的時候.
如果你沒有很特別的呼召.
沒有很特別的事情.
很簡單.
耶穌做什麼.
你就做什麼.
你知道耶穌做什麼嗎.
耶穌是PhD來的.
徒弟說你也是PhD.
你說我不是的.
我很討厭讀書的.
不是的.
耶穌是做什麼.
Preaching.
傳道.
耶穌是做什麼.
Healing.
祂會為病人去禱告.
醫治.
Deliverance.
祂會做趕鬼.
將很多克制的心靈得著釋放.
所以如果你沒有特別的呼召.
其實作為基督徒.
跟著耶穌做就可以了.
所以在第一季裡.
我們已經很清楚說了.
很多上帝在聖經裡寫明的旨意.
忘記了不知道是什麼.
回去自己再聽.

$^{241}$上一季裡.
為什麼會花那麼多時間.
整天說更新.
在小組討論.
討論到不知道說什麼的時候.
因為我覺得.
我們要認識神.
以致我們跟著耶穌走的時候.
我們知道要跟什麼.
但當我們知道要跟什麼的時候.
骨子裡.
很誠實地說.
很多事情是難的.
難是做不到的.
而是裡面很多事情拉扯.
很多事情的糾結.
要將聖經的價值觀.
成為自己的價值觀.
由真的很不習慣.
到要去學習習慣.
是一個很長的過程.
因此我們常常鼓勵弟子去學習.
在耶穌基督的救恩裡.
多點去默想和思想.
唯有在十字架的救恩裡.
我們再去看改變.
才能充滿盼望和能力.
因為靠自己去改變一些腦的性情.
要去靠自己化解人際關係之間的糾結.
是難的.
真的唯有從救恩的角度重新去認識神.
神是不是這麼嚴厲的呢?.
很多人不信主人會挑戰的.
這個神這麼惡.
不信祂又要下地獄.
對不起.
不是因為你不信祂下地獄.
你一向都下地獄.
只不過你不信祂.
你上不了救護車.

$^{281}$你明白嗎?.
很多時候我們不能夠重新在十字架的救恩裡.
去看我們的神.
去看這個世界.
我們有很多錯誤的期盼.
以為這個世界會越來越好.
信了主一定是飛黃騰達.
美滿幸福的.
遇到不如意的事情.
我們在上上很多的糾結.
因為我們看錯了這個世界.
曾經說得很清楚.
在罪的世界一定會成為過去的.
這個世界罪惡越來越滔天的時候.
人類所面對的苦難.
大部分其實是人禍.
真的多過所謂天然的災害.
因此當在這兩個基礎.
大家能夠固定了的時候.
其實說下去.
我們才會覺得容易聽一點.
只是容易聽一點.
因為當你明白了上帝的心意.
能夠學習去依靠十字架的救恩.
你的生命得到改變.
於是你在正確的道路上奔跑.
你真的會輕鬆一點.
問題是我們向哪裡奔跑呢.
不同的教會.
不同的機構.
不同的基督徒.
甚至我們小組.
不同的組員.
他們各自都有靈愁.
各自都有自己很想去事奉.
為主發夢的東西.
如果我是一個剛剛來的.
我也不太知道向哪個方向走.
因此當這個叫你這樣做.
那個又叫你這樣做.

$^{321}$最近又說招招待.
又叫你去做招待.
那段時間招的師伯又叫你去唱師伯.
於是我們真的一頭煙.
根本就不知道自己可以做些甚麼.
再加上如果你有之前所謂自殘營慰的心情.
我沒有他唱的歌那麼好聽.
我不夠這個組長.
他那麼有愛心.
只會頂著我.
如果給我我自己會殺死他.
我做不到組長的.
因此怎樣去面對.
真的希望上帝喜歡自己多一點.
我今天給你一個好消息.
我希望你甚麼都不記得.
記住這句說話.
就是無論如何.
上帝一直都那麼喜歡你.
跟嘉莉亞說上帝很喜歡你.
你要記住這個大前提.
你做事不做事.
上帝都一樣那麼愛你.
如果不是.
耶穌基督就不需要降世.
死在十字架上面.
拯救所有的人.
其實是真的.
如果你論拯救恩.
是人人都得救的.
現在的問題是人不肯跟從耶穌.
因此在所謂學習去事奉.
學習去親近主.
去走世間路的時候.
我們一定要記住.
上帝是喜悅我們.
如果不是.
祂不需要拯救我們.
祂知道我們會失信.
祂知道我們答應了的事.

$^{361}$我們又做不到.
好像剛才壽司哥一樣.
其實上帝一做就知道.
但是在上帝裡面.
世間觀念跟我們不一樣.
上帝其實已經看到.
你是天堂的樣子.
我們已經是得救的.
魔鬼無論在當中.
令我們有多少挑戰.
我們有多少跌跌碰碰都好.
不影響了我們的救恩.
我們已經被稱為上帝的兒女.
在上帝的國度裡面.
是與上帝同在.
耶穌自己親自說的.
祂說.
祂既然愛世間屬自己的人.
就會愛他們到底.
愛到底包括.
我們的軟弱.
包括我們說了不算數.
其實無論在什麼情況裡面.
只要我們願意跟從主的時候.
主一定不會放棄我們.
因此我們當說要面對.
真的過一個.
希望上帝喜歡的生活的時候.
我們要在今天.
好好地去.
建基在一個共同的基礎裡面.
就是這個十字架上的救恩.
但當然有很多東西.
是需要彼此去學習.
和彼此去提醒.
所以我選擇了一段經文.
其實好像與這個題目.
不知為何有關係.
記載在《創世記》的第四章.
是亞當夏娃犯罪之後.

$^{401}$第一個的故事.
這裡第一節就說有一天.
有一天是多久呢?.
不知道.
哪一天?.
不知道.
沒有說.
但起碼知道.
在第三章.
就是人類犯罪.
被上帝趕走.
離開了樂園的時間.
直至到何時發生這件事呢?.
不知道.
只是說有一天.
為何會這樣說呢?.
有沒有想過.
亞當夏娃這個夫婦.
曾經你是我骨中的骨.
肉中的肉.
多浪漫相愛的夫妻.
因為吃禁果的事.
他們竟然離開了樂園.
再沒有辦法地為他效力.
不用做有東西吃.
大家很恨的日子.
不用做有東西吃.
突然間消失了.
汗流浹背.
也沒有辦法謀生.
你想想.
他們兩個離開了.
上帝的樂園之後.
如何面對這個世界.
天氣不是他預期.
究竟如何令自己可以糊口.
完全不知道如何.
最重要.
這兩個人心裡一定有條刺.
是嗎?.

$^{441}$在很辛苦的時候.
那條刺一定會出現.
最壞的都是你.
不然你吃禁果就什麼的.
可能對方也說.
最壞的都是你.
誰說話不清不楚.
弟兄.
很多時候你老婆也這樣說你.
說話不清不楚.
說些不知情的.
問題是.
你們夫妻總之有一天.
他們應該和好.
然後才會同房.
才會生到一個兒子.
這個兒子叫做戈仁.
然後很快.
快極也差不多九個月了.
才生了第二個兒子.
就是叫做阿伯.
第三節又說.
有一天.
有一天是多久呢?.
你想想這兩個人.
能夠找到自己的職業.
一個是耕田的.
農夫.
一個弟弟是牧羊的.
一個牧羊人.
起碼也十多歲了.
不然也做不到這些事情.
在這段時間竟然有一天.
這兩個.
我想當他們是青年人.
竟然做了歷史上第一件事.
就是獻祭.
之前是沒有說過獻祭的.
什麼叫做獻祭不知道.
為什麼要獻祭.

$^{481}$沒有說.
總之就是看到.
戈仁和阿伯.
都是他自己老奴婦工作裡面的一些成果.
他竟然拿來.
去獻給耶和華.
當中.
究竟這十多年.
從出生.
看著父母吵架.
因為一定吵架.
這麼嚴重.
說回這些故以前成年的往事.
然後.
看著父母很辛苦的去養大他們.
他們小時候就要學習賺錢.
可能.
他們都知道有一個叫做上帝的人.
本來很有能力.
他們是富二代的.
誰知道父母不爭氣.
搞到自己.
差不多要蹲街了.
在這個時候.
這兩個人竟然會懂得.
買份禮物去探望爺爺.
差不多是這個想法.
原來你爺爺是.
你什麼的.
現在原來你扔了出街.
沒飯吃.
突然間.
回去.
哄一下爺爺.
可能.
有一點點好處也說不定.
人送禮.
心底裡.
是一定有一點點目的.
所以.

$^{521}$基督徒我們都是人.
我們親近主.
我們向主禱告.
我們怎樣也有一點點的想法.
是希望上帝祝福我們.
我們沒理由特意來祈壽.
上帝是救了我們.
但是.
阿伯.
就是拿他自己.
辛苦工作的東西.
該人又拿他自己辛苦工作的東西出來.
但是很奇怪.
聖經裡面的形容.
是有些不同的.
在第四節裡面.
他將阿伯所獻的羊群.
有兩個很特別的形容詞.
就是「投生」.
和「自由」.
這兩個的要求.
剛好在往後的摩西五經.
就說到這是獻給上帝.
最好的禮物.
小朋友就會說.
原來上帝喜歡腳扒.
不吃齋的.
為什麼結局是.
他喜歡吃羊扒.
但是又不吃齋的呢.
經文不是這樣說.
經文故意這樣去形容.
是在說.
阿伯在獻祭裡面.
他有他的心思.
他將他自己所做的事情裡面.
最好的拿出來.
來獻給耶和華.
他對上帝的尊崇和敬畏.
一定敬畏的.

$^{561}$因為很害怕.
阿伯做錯事就趕了出來.
一定有敬畏在.
於是他將最好的獻上.
結果上帝真的看中了.
阿伯和他的貢物.
但是竟然不閱納.
該人的獻祭.
如果我們留心經文所寫的.
你要留心前面的寫法.
他說他是看中了.
阿伯和他的祭物.
他是看不中.
該人和他的貢物.
我們中堂很多做父母的.
你就明白的.
其實子女送什麼給我們.
其實真的不是很重要的.
所有的錢都是你的.
很坦白.
但是.
正因為他是你的子女.
他拿一張紙.
畫一隻豬不像豬.
狗不像狗的來.
說媽媽我送給你.
這隻是牛來的.
是牛來的.
我們不會說任何東西.
很漂亮.
你真是牛.
很坦白我自己電腦也貼了一張.
不知是什麼在那裡.
因為小朋友.
他畫什麼給你.
你就會閱納.
你就會喜歡.
因為你喜歡你的子女.
這裡就很清楚.
是有些事情發生的.

$^{601}$上帝是看中了.
喜歡阿伯的.
因為可能他獻誠一心.
他所預備的東西.
是上帝喜悅的.
但是.
該人一定是有些事情.
所以上帝看不中他.
以致不閱納他的貢物.
結果我們看結果.
推翻上次.
該人的反應.
是大大發怒.
變了臉色.
中文寫得很好.
我們很明白.
一個的後輩.
一個的年輕.
一個的小朋友.
他可以發脾氣的.
因為不開心.
但是他能夠向你變臉.
你就明白.
是代表著一些事情.
這個就是.
該人和阿伯.
去看待上帝.
一個很不同的反應.
敬畏神.
是很重要的一個態度.
同樣.
我們敬畏我們的父母.
都是一個很重要的態度.
但是面對著他大逆不道.
在第六第七節.
上帝所說的話.
又不是我們這樣想的.
罵他諸如此類.
反而上帝和該人說.
為什麼你要生氣呢.

$^{641}$為什麼你要變臉呢.
他說如果你行得好.
豈不蒙越納.
正如我剛才所說.
其實只要你行得好.
你就一定蒙越納.
無論你欠的是什麼.
不是因為吃齋.
所以不越納你.
但是很重要.
你若行得不好.
罪就伏在門前.
上帝不是說越納與否的問題.
上帝關心的.
不是祭物的問題.
是祂的問題.
原來你行得不好的時候.
你要知道你活在一個有罪的世界.
你已經離開了上帝的任避.
你已經在這段月裡了.
當你行得不好的時候.
那惡者就伏在你門前.
隨時預備將你吞噬一樣.
所以上帝其實對於該人.
不單有愛心.
也給他很大的警告.
如果該人能夠尊敬神.
敬畏神的時候.
接著的殺人事件.
是可以不發生的.
他講得很清楚.
罪是什麼.
罪是亂無理.
亂無這個字是第二次出現.
第一次是在前一章.
就是當女人犯罪之後的咒詛.
我們會亂無你的丈夫.
亂無是一個強烈.
渴求並且要佔有與控制的情感.
所以大部分女人都是愛情大過天.

$^{681}$但是在這樣的狀況之下.
你就失去了你自己的理性.
你也沒有辦法去控制自己.
所以上帝為了救女人.
就跟丈夫說.
但是你的丈夫會管轄你.
我們最初以為怎麼搞的.
怎麼會要他管轄我.
但是你要明白.
原來上帝是用了丈夫.
來管轄他的太太.
讓他的太太有自控的能力.
在安全的範圍裡面.
做一個好好的女人.
所以該人應該要有法路的心.
預備會有下一件事出現.
上帝在這裡千鈞一髮的時間.
就跟他講.
你要制服他.
制服是要控制自己的情緒.
控制這件事.
在一個正確的範圍裡面.
以致不會越搞越和.
因此你明白.
在做基督徒的時候.
為什麼說.
信主有平安有喜樂.
回來後每個人都愛我.
每個人都對我好.
我任意發脾氣.
無所謂的.
女祖長領得順的.
他們都很有愛心的.
為什麼這麼上下的時候.
發覺不是的.
他們又要求我這樣.
又要求我那樣.
我不合適又要跟我談一大輪.
為什麼呢.
做基督徒很不自由.

$^{721}$像我媽媽那樣都要管.
弟兄姊妹.
人世間所有事都是有秩序的.
在秩序裡面一出現錯誤的時候.
就已經是世界末日了.
地球不跟著這樣走.
月亮不跟著這樣走.
我們已經出事了.
人世間都一樣.
父母君臣.
有夫妻.
有小朋友.
都有兄長.
有弟弟.
我們有不同的階層.
這些沒錯.
在現在這個現代.
不喜歡講規範.
不喜歡講架構的時候.
就覺得我要有自由.
但是我們是群居的生活.
大家不在自己的規範裡面.
就很出事了.
你看看外面那兩條馬路.
大家不跟著這條線這樣走.
一定是撞車的.
因此去控制自己.
讓自己走在正確的裡面.
一定要有一個很權力的上帝.
去教導和指引.
我們需要去學習.
去跟從.
好讓我們的生活.
真的活在這個恩典當中.
因此你明白.
在《摩西五經》裡面.
為什麼祭禮是一件常常出現的事.
以前就會覺得很煩.
讀那些利未記.
又殺那樣那樣.

$^{761}$怎樣殺怎樣洗.
都講得這麼清楚.
很煩.
做了基督徒其實都一樣.
現在又要回來崇拜.
回來崇拜又不准遲到.
然後又說穿衣服又不可以暴露.
一二到死.
怎麼搞.
我們很多事情.
都會覺得很麻煩.
你明不明白.
問題是.
原來很多的規範.
背後是有一個目的.
它的好意.
我們要去明白.
當中很多做父母的.
我都想說.
小朋友為什麼有這麼多.
不肯跟我們說話.
是因為我們沒有將.
規範背後的好意.
告訴小朋友.
他應該知道.
他跟從這樣去走.
他有什麼好處.
他不跟從這樣走.
他會有什麼結果.
這樣他就學會為自己負責任.
沒有人喜歡自己慘的.
或者自己不好的.
他會有兩次.
他就會乖乖的.
去跟從這些所謂的律例典章.
因為律例典章背後.
其實都是讓我們得著好處的.
因此其實.
你看到在摩西五經裡面.
很多的制例.

$^{801}$是已經訂明了出來.
這個是古代的時間.
舊約的時間.
但是你知道.
「飯祭」和「素祭」.
特意想說的原因是.
又吃肉的.
「飯祭」是獻牛羊的.
肉來的.
不要說吃素.
上帝都吃素的.
「素祭」.
上帝都是喜悅的.
這兩個祭是天天要做的.
感恩上帝的公益.
最重要是宣告自己.
是屬於上帝.
以你一神為神.
我經常都見到媽媽一早起床.
不用說話.
你不知道她信什麼都好.
媽媽擦完牙洗完臉.
第一件事就是一炸香.
就是點了它.
在祖先那裡搖兩下.
之後請熄它之後.
接著就到處插.
到處都可以插的.
你明不明白.
她做的這個動作.
是代表什麼.
我這間屋子.
這些是神來的.
祂們會保佑我的.
今天.
我們做基督徒都是一樣的.
為什麼說要靈修.
為什麼要讀經.
為什麼要親近主.
讀經是讓你認識神的.

$^{841}$如果你和上帝很友好的.
你不看本聖經.
我都不怪你的.
你明不明白.
但是你說我又不知道上帝想怎樣.
又不知道上帝怎麼跟我說話.
你就好好地讀本聖經吧.
不然怎樣和上帝溝通呢.
當我們每一天都能和上帝溝通.
親近的時候.
你的生活.
就能夠跟從上帝的指引.
你也都活在上帝的祝福裡面.
何樂而不為呢.
當然有不同的祭禮.
是在以色列人裡面是可以教的.
教到他們如果有感恩的事項.
就是要來感恩.
然後他們很大件事.
大件事的時候.
就是因為有一個贖罪日.
這個贖罪是在講冤罪的.
是人類犯罪這麼大件事的.
所以一年大祭司會去聖殿.
為全國的百姓獻祭.
我們就感恩了.
我們不需要處理了.
因為耶穌基督已經為我們一次獻上.
我們對於罪的權勢.
魔鬼的所謂控制.
我們不需要擔心.
在耶穌基督十字架的救恩裡面.
我們已經是得勝的.
但是也有一件事.
我們需要在自己的親近裡.
特別是你覺得你做了一些事.
虧虧了別人.
對別人不好的.
說錯了兩句話.
得罪了別人.

$^{881}$傷害了別人.
我們也應該要認罪.
因為事實上.
這個是破壞了人與人之間.
一個很重要的關係.
很多時候.
人際關係不是做了什麼殺人放火的事.
都是一個無言之失.
一個嘴苗苗.
那就可能另一個很好的朋友.
大家都沒有了.
因此不同的祭禮.
在舊約那裡.
就是做了一件事.
就是叫做儀式化.
我們叫做什麼.
就是有那個儀式感.
為什麼求婚現在要鑽石戒指呢.
我們那個年代是不用的.
現在是一定要的.
還要在Facebook給別人看到的.
可以嗎.
這個叫做儀式感.
現在的男士就不幸了.
用心賺錢吧.
但是.
每一樣東西.
都是為了什麼.
讓我們知道我們在那個祭禮的背後.
學習的那個好意.
因此基督徒.
我們又覺得教會可能有很多的安排.
有很多的聚會.
但是你要記住.
一切都跟你的得救是無關的.
我們不是靠行為得救.
只不過.
一個教會.
我們幫助大家去學習的時候.
總會有一些的儀式感.

$^{921}$當在那些儀式感裡面提醒我們.
我們可能已經完成一個很重要的目的.
我們為此而興奮一下.
所以你看到基督生平.
那一百多個弟兄姊妹.
開心到瘋了.
其實吃一頓飯而已.
就是一個儀式感.
令到我們對於我們附上了那六年.
可能有些人讀了很久的.
那個努力的學習.
有一個慶祝.
所以當你從這樣的態度去看的時候.
不要再因為彼此之間的那個好意.
而成為了一個無形的壓力.
我們大家知道.
我們要活著的.
就是基督.
這個在主餐裡面是一個很好的提醒.
保羅更加說.
我們要將身體獻上當作活祭.
是聖潔的.
最重要的是什麼?.
講得很清楚.
這個是神所喜悅的.
弟兄姊妹.
究竟你是不是一定要在教會有個崗位.
才是神所喜悅的呢?.
絕對不是.
就算我們教會很小的時候.
我們都很鼓勵弟兄姊妹的.
這裡不是很多的崗位.
也都個人的恩賜未必用得到.
上帝的角度是遼闊的.
甚至你去幫其他教會做.
我們覺得都沒有問題的.
因為我們能夠去服侍上帝.
在神的角度裡面有一份是我們的福氣.
問題是.
當我們在身邊有很多不同的組員.

$^{961}$我們有不同的好意.
我們要懂得自己去消化.
不要成為壓力.
今天就招待有壓力去做.
那天唱詩班 組長又叫你去做.
今天讀靈教又叫你去讀.
那天讀摩西又叫你去讀.
我們整天都被人逼的.
覺不覺得.
跟家裡人說.
放鬆一點.
真的放鬆一點.
教會都要讀的.
讀一輩子都慢慢讀沒問題.
我們又不是趕著畢業的.
但是.
為什麼要去讀聖經.
為什麼要去學習.
因為我們需要人幫助我們去認識.
你說不是.
外面有神學院可以讀.
好 你可以去日神學院讀.
真的沒問題.
我們只不過是幫助一些.
他們未有心智去到神學院讀書的弟兄姊妹.
都可以去認識聖經多一點.
一直到我們在這個基礎.
大家明白的時候.
我們再說我們去討主喜悅.
我們去工作也好.
我們去讀書也好.
我們去服事也好.
不要受很多框框的限制.
我們同屬天國的子民.
凡是天國的事情.
你是上帝的兒女.
你都可以去服事.
你都可以去幫手.
問題是.
這裡說的是作活祭.

$^{1001}$活祭不是死的祭.
意思是我們不是單單說做什麼的事.
不是羊才蒙月臘.
就是獻個菜.
獻個番茄.
神不月臘.
最重要的是你的心態.
究竟面對著給你的呼召和負擔的時候.
你用一個什麼心態去回應它.
很多時候教會用了很多東西去幫.
但是也很難所有盡解.
因為四大本分.
大家很努力地去做.
這是好的.
因為不說一二三四.
一個新的基督徒都不知道從哪裡開始.
因為每個人都專家地跟他說.
他真的不知道在搞什麼.
定了一個規矩出來.
讓別人可以去跟隨.
是不是那個規矩就已經等於上帝的呢.
跟別人說.
上帝是大於一切的規矩.
都很坦白.
如果你說你有感動做什麼.
你問到我們做傳道牧師的.
我們都不敢反你.
除非我們真的很清楚.
覺得上帝是KO了你這個想法.
我們都是彼此尊重.
因為每一個的創造是上帝自己獨特的.
上帝給每一個弟兄姊妹.
在不同的環境裡.
只要他獻上活祭.
他就能夠去榮耀神.
我們在台上說話的人.
我們未必和你一樣.
可以去你家.
或者你去探老人院的地方.
去接觸到那些人.

$^{1041}$上帝是要用每一個他的子民.
在這個地上成為其他人的祝福.
因為上帝的心意是想每一個人都有機會.
知道有一位愛他們的神.
他們不需要活在罪裡面.
只要他們願意來跟從主.
但最可惜的是.
很多時候我們弟兄姊妹.
不記得跟從主.
進行一切律例典章背後的好意.
我們只覺得為什麼要回來.
為什麼我要做什麼.
我要做什麼.
當我們明白十字架.
九因這份愛的時候.
其實一切都是你和上帝愛的關係.
在這個世界裡面.
如果我當一個天秤去看.
很多時候為什麼弟兄姊妹.
特別是回教會的時候.
有牧者有祖宗去幫你的時候.
好像總是叫你在那裡做很多事情.
你要知道原來在另一邊的位置.
是因為我們活在自己的肉體裡面.
肉體的喜歡.
我們要吃東西.
我們喜歡別人喜歡.
我們可以任意妄為.
這些所謂的自由.
一直是拉著我們.
這個力量是很大的.
背後是罪.
為什麼?.
因為我們的肉體是我們的重點.
是我們最軟弱的地方.
兩夫妻無事無干.
一天到晚吵架.
吵什麼架呢?.
其實我覺得大家吵架的位置.
都不是很嚴重的位置.

$^{1081}$但是可以破壞一個婚姻.
身體累.
多事做.
煩.
一個不控制自己的.
就亂說廿四.
一次說笑話.
兩次說笑話.
第三次就說不出笑話.
不知道彼此那個.
忍讓的能力有多高.
於是為什麼會鬧事.
在家裡吵架呢?.
其實解釋不了.
所以做輔導其實沒什麼用.
因為輔導是攤出來的.
攤出來沒什麼用.
但是就是大家不控制自己.
就一下子在一個無端端的嘈吵裡面.
人的心是很脆弱的.
你想想我們的肉體.
我們的身體也有病.
我們怎樣吃都是病的.
現在差在分開吃.
可能慢點死的意思.
每樣都有問題.
但是上帝為什麼要教導我們.
要一二三四五六七呢?.
就是讓我們有一個平衡力.
我們如果不竭力去愛主.
我們不竭力去跟從上帝.
我們一定是被自己的肉體拖垮的.
今早你起床不努力.
你都是回到床上的.
頂智慧我們每天都經歷.
這個所謂肉體的軟弱和掙扎.
因此我鼓勵在一個這樣的題目.
接下來這一季.
大家可能會教我們怎樣做.
哄上帝開心.

$^{1121}$諸如此類.
我希望第一講.
是讓大家知道.
我們用一個什麼態度.
來回應這些所謂的要求.
土主的喜悅.
是因為我們喜歡上帝.
你有沒有想過.
你的小朋友說那隻豬不像豬.
狗不像狗的走過來.
他都給你的.
為什麼?.
他喜歡你.
他不知道自己畫得漂不漂亮.
他喜歡你就給你.
有沒有見過小朋友喜歡你.
見到你給你吃的時候.
他吃過的雪糕都是給你吃的.
他真的沒有想過.
現在手足口病.
你這樣送給我吃.
沒有啊 我喜歡你.
我就這樣給你.
當我們去明白.
愛這個吸引力的時候.
你就知道上帝喜悅的是我們.
無論我們怎樣來服侍他.
上帝都會悅覽我們.
我們盡我們的能力.
我們不需要和隔壁那個比較.
窮寡婦那兩個小錢.
是上帝大大的誇獎.
所以當我們在這個基礎之下.
我們大家用自己的方法.
去向這位愛我們的神.
去表達自己的愛意.
我知道現在年輕一代.
很多事情做得有點困難.
很多時候他們都找了一些.
叫做Shift的工.

$^{1161}$我們現在的小組.
因為經常都Shift.
要看監表才知道他能不能回家.
變成了我們一個小組的經歷.
就是每個月都看監表祈禱.
希望他可以回來.
但是真的不容易.
星期天上崇拜.
他們都有小組的.
有在耳目的事奉.
接著我們又有小組時間.
一個不幸運的.
不是不幸運 是幸運的.
高明謙的講座又來了.
就很困難.
你調監怎麼可以調到.
就算每天給錢買監.
你都不可能一個星期走開幾次.
你肯回家很早 回家很晚.
盡量在當中調.
已經是一件很難的事.
偏監不是自己可以預計的.
但是我覺得.
當我們有這個心想愛主的時候.
上帝就起緣.
我們這幾個月.
每個小組經常都在過這些日子.
究竟能不能回到小組.
調到到監.
上帝能不能安排到.
是令到其他不關事的組員.
都在這個過程中.
好好看到上帝的恩典.
上帝一定是為他所愛的.
去預備.
丁姐妹 上帝愛你.
上帝也愛我.
連耶穌基督都願意放在十字架上.
其實不要再計較自己做到什麼.
做不到什麼.

$^{1201}$反而你日常的生活.
你起床 你見到人.
你走在這個城市當中.
究竟你能夠真的活出.
上帝的愛呢.
人會不會見到你不一樣呢.
這個才是最大的是逢.
我們年紀已經差不多了.
教會開始.
老人家比較多.
我們很多時候年輕都已經有病痛.
病痛是沒有辦法避免的.
因為人怎樣都要死.
但問題是.
當我們遊走在各大醫院.
在排隊.
我們可能都在做治療.
大家都是這麼痛的時候.
人家見到在裡面有那份平安.
有那份穩妥.
這個就是吸引他.
為什麼你和我是不同的.
弟兄姊妹.
不需要一定要懂得說懷傷.
很難說的.
也不需要斑馬.
整個幸福小組做這麼多事情.
你可能只是坐在醫院等覆診.
但是你那份信心.
你那份平安.
你對身邊人一個溫馨的微笑.
見到身邊的辛苦.
一句很簡單的慰問.
這個就是服侍.
這個就是活祭.
不需要讀書.
也不需要受訓練.
讓上帝的愛在你的生命當中永留.
你就白白的得來.
就白白的去捨去.

$^{1241}$我選了這段金句是因為.
我很想我們記得.
上帝喜歡什麼.
我們一起讀好不好.
耶和華喜悅泛濟和平安濟.
喜如喜悅人聽從他的話.
聽命勝於憲制.
順從勝於公羊的自由.
是的.
上帝喜歡我們聽命.
上帝喜歡我們順從.
但是不是一個刻薄的老闆.
或者是一個很控制的父母.
為他們的好處.
要我們聽命.
要我們順從.
乃是上帝知道.
在這個在世的路裡面.
實在有太大的扯鈴.
他指引我們的路.
我們不走的時候.
真是可憐的.
我們就活在很多困苦的裡面.
我們上班的努力做好自己.
我們真的努力犧牲奉獻.
因為實在太多建政了.
當我們真的犧牲奉獻.
說知道老闆被解僱了.
但每個頂智慧都看到.
他們生活上是公認的.
夫妻間.
上帝要太太的順服丈夫.
然後要丈夫愛妻子如同自己.
這個聽清楚.
愛妻子如同愛你自己.
這麼簡單.
其實任何的夫妻關係.
在主角裡面.
因為跟從他的.
就可以得到修補.

$^{1281}$做父母的.
不要惹兒女的氣.
按著真理去教他.
做子女的.
孝敬父母.
就會得福在世長壽.
所有東西其實是很清晰的.
問題是我們願不願意去走.
你想想我們剛才說的那些.
關上帝什麼事.
對上帝一點好處都沒有.
你夫妻吵架就吵架.
你兒女吵架就吵架.
不關上帝的事.
很多上帝教導的律例典章.
叫我們走的路.
是讓我們無福的.
所以希望在第一講究裡面.
我們真的去明白.
上帝對我們的接納.
和我們的愛.
上帝是沒有要求的.
祂要求的.
就是我們能夠以祂為我們的神.
祂要求的是我們去跟從祂.
背後的好意.
不是為了上帝的好處.
是為了我們得到.
祂應許的一切的福樂.
願意我們一起努力去學習.
在前面的日子裡.
我們再一次跟從我們的主.
讓我們無論在什麼環境裡.
在什麼階層.
在什麼退休 生病.
什麼狀況都好.
我們每一個弟兄姊妹.
都能夠成為活祭.
去榮耀我們主的名.
我們一起讚揚好嗎?.

$^{1321}$我們鼓勵你開心地為自己去禱告.
在這一季的學習中.
幫助我們真的能夠知道.
如何能夠禱告上帝的喜悅.
最重要的是我們知道.
當我們跟從上帝的時候.
我們一定走在無福的路當中.
\newpage



\section{哥林多前書 10:23-33}
\label{sec:032449}
\textbf{得自由,失自由?}~~~~梁永善牧師
\newline
\newline
連結: \href{http://www.church.com.hk/acms/content.asp?site=cdc\&op=show\&type=product\&code=032449\&layout=sermon}{\texttt{032449}} ~~~~ 語音日期: 2026-07-12
\newline
\newline
\hyperref[sec:032448]{\small{< < < PREV SERMON < < <}}
~
\hyperref[sec:index_chronic]{\small{[返順時目]}}
~
\hyperref[sec:032450]{\small{> > > NEXT SERMON > > >}}
\newline
\newline
哥林多前書 10:23-33
\newline
\begin{longtable}{cl}
\hline
\hline
章節 & 經文 (和合本修訂版)\\
\hline
10:23 & \begin{tabularx}{0.7\textwidth}{X} 「凡事都可行」，但不都有益處；「凡事都可行」，但不都造就人。 \end{tabularx} \\ \\ \relax
10:24 & \begin{tabularx}{0.7\textwidth}{X} 無論甚麼人，不要求自己的益處，而要求別人的益處。 \end{tabularx} \\ \\ \relax
10:25 & \begin{tabularx}{0.7\textwidth}{X} 凡市場上所賣的，你們只管吃，不要為良心的緣故問甚麼， \end{tabularx} \\ \\ \relax
10:26 & \begin{tabularx}{0.7\textwidth}{X} 「因為地和其中所充滿的都屬於主」。 \end{tabularx} \\ \\ \relax
10:27 & \begin{tabularx}{0.7\textwidth}{X} 倘若有一個不信的人請你們吃飯，而你們也願意去，凡擺在你們面前的，只管吃，不要為良心的緣故問甚麼。 \end{tabularx} \\ \\ \relax
10:28 & \begin{tabularx}{0.7\textwidth}{X} 若有人對你們說：「這是獻過祭的物」，那麼為了那告訴你們的人，並為了良心的緣故就不吃。 \end{tabularx} \\ \\ \relax
10:29 & \begin{tabularx}{0.7\textwidth}{X} 我說的良心不是你自己的，而是他的。我的自由為甚麼被別人的良心評斷呢？ \end{tabularx} \\ \\ \relax
10:30 & \begin{tabularx}{0.7\textwidth}{X} 我若謝恩而吃，為甚麼因我謝恩的物被人毀謗呢？ \end{tabularx} \\ \\ \relax
10:31 & \begin{tabularx}{0.7\textwidth}{X} 所以，你們或吃或喝，無論做甚麼，都要為榮耀神而做。 \end{tabularx} \\ \\ \relax
10:32 & \begin{tabularx}{0.7\textwidth}{X} 你們不要使猶太人、希臘人，或神教會中的人跌倒； \end{tabularx} \\ \\ \relax
10:33 & \begin{tabularx}{0.7\textwidth}{X} 但要像我一樣，凡事都使眾人喜歡，不求自己的益處，只求眾人的益處，使他們得救。 \end{tabularx} \\ \\
[1ex]
\hline
\hline
\end{longtable}
$^{1}$我不知大家去傳福音的時候.
有沒有這些人的反應就是.
有沒有叫我信.
信了耶穌.
這件事又不准做.
那件事又不准做.
我有一個很深的印象.
就是我剛剛做傳道人.
第一個要探訪的家庭.
因為他的媽媽和三個兒女都回了我的教會.
很好啊.
他們現在確有所成.
兒子已經做了校長.
女兒也是做了主任.
即是學校的主任.
我記得有一次我上去探望他.
他正在吃晚飯.
用現在來說.
那個人是做物流的.
但那個時候的意思是做搬運的.
我就說.
你家人都回了教會.
不如你也回去信耶穌.
神經病.
做什麼.
信耶穌.
這件事又不准做.
那件事又不准做.
什麼自由都沒有.
那.
要信一下.
我說是嗎.
什麼自由都沒有嗎.
我說哪些東西沒有自由.
他就在耍我.
擔著口煙.
你是基督徒可不可以吸煙呢.
我說真的不能吸煙.
我說是啊.
你也不能吸煙.

$^{41}$什麼我不能吸煙.
我說你是噴煙而已.
你吸了也不是噴.
是這樣吸的.
你有吸過煙嗎.
我就說.
做基督徒其實也可以吸煙.
聖經沒有說不可以.
吸煙不影響你的德格.
但會影響你早點見耶穌.
那我就問.
那你不吸煙會怎樣.
可能你們沒聽過這個術語.
是那個年代的.
那我就會吊煙龜.
吸煙吸得多叫煙鏟.
就是.
吸煙很辛苦叫吊煙龜.
什麼意思.
就說我會很辛苦.
哦.
那我就說就是了.
我就經常跟他說.
我是有自由去吸煙的.
不過我是運用了我的自由權.
放棄不吸煙而已.
你不吸煙很辛苦.
但我不吸煙.
我一點也不辛苦.
他當然不忿氣了.
然後他又駁嘴了.
那你能不能喝酒.
因為在吃飯.
那些很傳統的.
斬一大塊叉燒.
喝一杯肉冰燒那些.
那聖經又說.
是可以喝酒的.
身體不好可以喝一點點酒.
但聖經也提到我們.

$^{81}$酒能使人放蕩.
所以不可醉酒.
那我就不懂得喝酒.
那我就問他.
那你不喝酒怎麼辦.
他不喝酒哪有力氣做事.
那我就說到原來.
我不是沒有自由.
是你沒有自由.
我是有不犯罪的自由.
但你是沒有.
好了.
原來今天.
每一個在罪裡面被黑榜的人.
都是失去了不犯罪的自由.
主耶穌怎麼說呢.
我們一起來讀《約翰福音》.
第八章第34節.
我們一起來讀出.
耶穌回答說.
我實實在在的告訴你們.
所有犯罪的就是罪的奴僕.
你們必曉得真理.
真理必叫你們得以自由.
但如果我們有一個成熟的生命與神的關係好.
我們抗罪的能力.
是強很多的.
分別就在這裡其中一點.
就是我們能夠有能力抗拒罪惡.
原來信主的時候我們很自由.
但是有規範的.
是真理裡面可以得到自由.
就是在神的真理裡面.
你做什麼都可以.
很簡單舉個例子.
魚在水裡面是可以不斷游水的.
游去哪裡都可以.
不過是不可以游蛙式的.
魚不懂得游蛙式的.
但是如果那條魚說我不喜歡在水裡面游.

$^{121}$我要在陸地裡面游.
一下子跳上岸的時候.
牠都會是魚.
變成什麼魚呢?.
死魚.
或者鹹魚.
就是說牠在水裡面.
牠是可以很自由的.
但是水就是牠的規範.
我們作為基督徒我們是有自由的.
但是在真理裡面有自由.
但是不可以因為信主而輕率犯罪.
不要以為信主之後我犯罪神就會赦免.
因為神會有刑罰的.
好了.
我們更加留意到.
原來聖經說到我們以前被罪合制住.
但是主耶穌拯救了我們.
我們就轉了主人.
我們不是再被罪合制.
就不是罪的奴僕了.
我們要成為義的奴僕.
義的奴隸.
某個程度.
我又好像沒什麼自由.
我們一起來讀羅馬書第六章的.
十五到第十八節.
我們一起來讀出.
我們從前雖然作罪的奴僕.
現今卻從心裡信服了所傳給你們的道理的無法.
你們既從罪裡得了釋放.
就作了義的奴僕.
這裡很清楚說我們以前被罪合制住.
我們是罪的奴隸.
我們沒有不犯罪的自由.
好了.
耶穌拯救了我們的時候.
原來我們轉了.
我們是屬於主的.
但是我們還是義的奴隸.

$^{161}$那不就好像沒有了自由.
好像很複雜一樣.
但是保羅就提醒我們.
他說到原來我們信了主.
不只是為了上天堂.
我們有一個使命.
這個使命是什麼呢?.
哥林多前書第六章的第二十節.
我們一起來讀出.
因為你們是重價買來的.
所以要在你們的身子上榮耀神.
那個奴隸服侍的那個是主人.
但是後來有一個人說.
我都挺喜歡你這個奴隸.
可不可以轉過來我那裡.
行,沒問題.
我已經給了錢買斷他的.
如果你真的要他的時候.
行,你給錢.
而且一定要怎麼樣?.
那筆價錢是會大過他以前買的.
於是有一個人.
他給了一筆更加大的價錢.
給這個主人.
就可以將這個奴隸怎樣?.
轉了屬於自己.
你們是重價買來的.
這個就是一個買屬的意思.
主耶穌為了我們的緣故.
被釘十字架完成了救恩.
我們肯接受相信祂的時候.
我們就不再被罪核制.
我們就會成為主的兒女.
不只是.
但原來我們有責任.
去做異議的奴僕.
而我們有一個使命是什麼呢?.
不是為了上天堂這麼簡單.
這已經是一個事實.
主耶穌已經不在世上.

$^{201}$原來別人是透過我們的生命.
我們的生活.
去認識這個神.
原來我們有一個使命和責任.
就是我們要在身上榮耀神.
那怎麼榮耀神呢?.
保羅就用了一個例證.
就是應不應該吃拜過菩薩的食物.
在多倫多前書第十章前面.
已經講了一些信息.
接著保羅就做一個小結.
「凡事都可行,但不都有益處」.
「凡事都可行,但不都造就人」.
就是有很多事都可以做.
其實是沒有問題的.
但是不是每一次做.
都是對你有益處呢?.
同樣.
如果你做這些事.
會不會幫到人呢?.
接著就引申說.
可不可以吃拜過菩薩.
拜過偶像的食物.
到最後又再一個結論.
這個結論都是很相似的.
我們再讀第十章第三十三節.
「凡事都可行,但不都有益處」.
「不求自己的益處,只求眾人的益處」.
「叫他們得救」.
保羅就說到.
我其實做每一件事都想一想.
會不會影響到人.
我希望可以幫到人.
那究竟可不可以吃過拜菩薩的食物呢?.
我不知道在我們當中.
應該很多人都是第一代的基督徒.
第一代的基督徒.
因為你父母都未信主.
有哪些是第一代的基督徒?.
最後說很多,差不多一半以上.

$^{241}$我是第一代的基督徒.
如果熟悉我已經經常告訴你.
我家是很傳統的家庭.
是很迷信的.
我未信耶穌之前.
我是很喜歡他們拜菩薩的.
為什麼呢?.
我們六十年代是沒有什麼東西吃的.
有時候真的連菜都沒有.
到什麼地步呢?.
沒有菜吃就去買酸梅醬拌飯吃.
你們未經過這些日子.
但是一有什麼拜神.
他們的那些拜菩薩.
那就不同了.
有燒肉.
燒肉一定是整條豬橫舌的.
一條豬舌.
我家會買雞的.
那時候是自己宰雞的.
不單整隻雞.
還有什麼叫扶持.
知不知道什麼叫扶持?.
有些人知道,有些人不知道.
就是雞腸,雞卵,雞雜.
那些叫扶持放在那裡.
還有一些紅包.
晚上還要很多菜吃.
所以我很喜歡.
但是為什麼呢?.
拜菩薩有什麼?.
有好菜吃.
我住摩羅街.
當然想去吃摩羅.
由摩羅下街上去文武廟五分鐘.
我阿麻去廟.
我跟著她去.
除了有東西吃之外.
我可以到處玩.
我家的菩薩很漂亮的.

$^{281}$無論生辰死忌都要拜.
因為拜完之後.
那些燒肉沒有了很多.
燒肉皮脆的都沒有了.
最特別.
那些雞卵又不見了.
而最奇妙的是什麼呢?.
那些紅包仍然在那裡.
但是內涵已經沒有了.
因為為什麼呢?.
當我阿麻去化寶.
就是燒完寶拿出去燒的時候.
我就立刻拿了一包.
從底部撕開.
用手指挖掉那些蓮蓉.
很美味的.
然後放進去之後蓋上.
所以前面的紅包是很完整的.
但是裡面沒有了內涵.
被我家人罵到我死.
但是那時候我真的喜歡吃東西.
不如現在也是.
後來我十三歲開始信主了.
當我在家裡拜菩薩的時候.
用雞翅膀吃.
媽媽就說.
喂.
你信耶穌的.
你拜過菩薩.
拜過祖先的東西.
你吃?.
你怕不怕耶穌罵你?.
我說那就不吃了.
我是為了主的緣故.
我放棄.
我就不會再吃這些.
不過是在他們拜之前.
我才偷了來吃.
就是沒有問題.
我沒有拜過.

$^{321}$我只是把我的步驟提早了去做.
所以我們要靈活的.
所以我們仍然可以吃.
好了.
但是哥林多教會.
保羅就說到.
絕對不可以去吃.
在廟宇裡面的那些筵席.
他們拜過菩薩的東西.
絕對不可以去吃.
作為基督徒不可以去.
我們一起來讀哥林多前書第十章.
二十一到二十二節.
我們一起來讀出.
你們不能喝主的杯 又喝鬼的杯.
不能吃主的筵席 又吃鬼的筵席.
我們可以吃主的憤恨嗎?.
我們比他還有能力嗎?.
保羅就斬釘截鐵地說到.
你們作為基督徒是不應該.
去廟宇裡面來吃這些筵席.
為什麼呢?.
在神學上面.
就是因為他們仍然相信.
其實某個程度的信念.
和我們很相似.
我們基督徒都會領聖餐.
我們領聖餐的時候.
當然我們不是說.
耶穌降臨在餅裡面.
但是有一個象徵意義.
我們吃這個餅的時候.
就是進入我們生命裡面.
我們喝這個葡萄汁.
都進入我們生命當中.
就是我們與主結連.
我們都相信.
原來當我們來祭祀那些祭物.
那些神靈就會附在那些祭物上面.
你和他們一起在這些廟宇裡面吃的時候.

$^{361}$就好像認同了這個信仰.
人們會覺得.
原來相信耶穌也可以這樣.
就影響人的信仰.
其實這個不單單是一個神學的問題.
還有就是當時的社會.
大家知道原來哥林多.
是一個非常淫亂的城市.
那麼他們的道德生活是很敗壞.
如果你看哥林多的前書.
保羅責備他們.
他當中竟然有人是基督徒.
與自己的繼母發生關係.
但是教會的人都沒有責備他.
覺得沒有問題.
所以你想一下.
原來他們是很受外邦人感染的.
原來他們去拜這些假神的時候.
是有一些女祭司.
但是這些女祭司其實是什麼呢?.
就是妓女.
是妙妓.
就好像他舊約裡面.
我們拜巴黎和阿斯塔諾.
他都會很淫亂地發生關係.
甚至這些女祭司.
在鞋子上是有號碼的.
你可以選擇.
這個我只知道是弟兄.
這個漂亮.
她神聖一點.
我會跟她有好的關係.
於是可能就會跟那個女祭司.
就是妙妓發生關係.
是可以這樣.
你們的鞋子有沒有號碼的?.
有的.
我十號半.
為什麼你們沒有的?.
你們的鞋子沒有號碼的?.

$^{401}$我那些有的.
每次都有的.
一是十號一是十號半.
但是不關事的.
你明白嗎?.
好了.
大約就是這樣.
他說到.
原來你去這些妙語裡面.
你要吃這些筵席.
就是影響別人對你的信仰的看法.
甚至影響你的道德生活.
就是跟我們那些拜菩薩是不同的.
但是那時候.
一個很繁榮的城市.
很繁榮的城市.
它幾乎要做生意的.
大家知道以前在國內做生意.
簽約就不是早上在辦公室簽的.
是早幾天就去卡拉OK.
去夜總會.
喝酒.
然後接著做很多事情.
搞定了.
然後才簽約.
明白這個道理嗎?.
所以他們也為了做生意.
有時候就會參與這些.
但是保羅就說你們要放棄.
所以絕對不可以吃.
在妙語裡面那些筵席.
不過在家裡.
又不同.
可以有一些微信的人請你去家裡吃東西的時候.
你不要問那麼多.
別人請你吃什麼.
你就吃吧.
好,我們一起來看看去讀.
就是《哥林多前書》第十章二十五到二十六節.
凡事上所賣的.

$^{441}$你們只管吃.
不要為良心的緣故問什麼話.
因為地和其中所充滿的.
都屬符主.
他說到你去.
別人請你吃飯.
你不要問那些東西是從哪裡來的.
為什麼呢?.
原來有很多人拿一些祭生去祭祀.
祭祀完之後.
未必一定馬上吃完.
於是會運回市場賣.
甚至市場會特別來講.
這一隻是祭過神的.
有神的祝福.
賣貴一點的.
那些人可能買回去弄來吃.
你就不要問那些東西是從哪裡來的.
又或者.
他們自己家裡也有他們的家神.
他們會有一些好的食物.
先給祖先吃.
我家裡就是.
我家裡有生辰死祭.
紀念我爺爺.
這些東西.
就會拜神.
一定是很好的食物.
甚至我們請人回來吃飯.
如果那些非常好的菜.
原來我們事先都要做些什麼呢?.
是放在祖先的神位上面.
先供奉祂.
跟現在其實是很相似的.
現在你想想你們出去吃飯的時候.
都要先擺一擺的.
是給誰吃的?.
是給手機吃的.
就算做爸爸那些.
出街說吃飯了.

$^{481}$然後.
你一上去夾的時候.
可能兒子都會說.
喂!.
相機還沒拍,吃什麼?.
你會被人罵的.
雖然是你給錢.
原來現在很重要的.
年輕人是手機先吃的.
手機成為另類的祖先.
我們的年代就是祖先要吃這些東西.
又或者特別很多的神.
那時候我們家裡有很多神拜的.
我們有Balcony的.
有露台的.
那時候.
就是叫騎樓.
你們知道騎樓有什麼神嗎?.
大家知不知道?.
天官賜福.
叫當天.
門口呢?.
有土地的.
看到我說嗎?.
那些土地過不了來英國的.
因為不懂英文.
好了.
還有廚房有什麼?.
造君.
很多這些神拜的.
所以呢.
多倫多教會裡面.
他們那時候的社會.
都很多這些神.
可能他們吃的東西.
是用來祭祀神的.
你不要問.
總之別人給你吃.
你就吃吧.
但是如果你沒有問.

$^{521}$但是那個屋主主動告訴你.
那些是拜過神的.
那你又怎樣呢?.
我們一起來讀.
哥林多前書第十章的第二十八節.
我們一起來讀.
若有人對你們說.
這是獻過祭的物.
就要為那告訴你們的人.
並為良心的緣故不吃.
如果那個人無端端告訴你.
這些是拜過菩薩的.
又可以拜過我的祖先的.
或者這些食物.
原來是在廟宇運過出來的.
但是我買回來的.
他就為了這個緣故.
你馬上不要吃.
甚至離開.
但是用人的角度來說.
那就是很不給面子的屋主.
可能令到關係之間.
有點惡劣的.
但是保羅就是為了良心的緣故.
特別你留意這個良心.
不是自己的良心.
你留意保羅說.
我吃東西有什麼問題.
為什麼要為了他們.
我要來但是他放棄呢?.
短短經文出現了五次關於良心的事.
是為了對方.
免得他誤會了我們的信仰.
其實萬物都是從神而來的.
我們的神是食物的賜予者.
但是為了他們的緣故.
所以我們會放棄某些東西.
那就好像沒什麼自由.
那你就很無奈.
很痛苦.

$^{561}$但是原來我們放棄自由.
不是因為我們的無奈.
有一個很大的動力.
是什麼呢?.
我們一起來讀哥林多後書第五章十四十五兩節.
我們一起來讀出.
原來基督的愛激勵我們.
因我們上一人既替眾人死.
眾人就都死了.
並且他替眾人死.
是叫那些活著的人不再為自己活.
乃為替他們死而復活的主活.
我們是為神而活.
我們希望別人透過我們.
可以見到神的榮美.
那為神而活是很抽象.
某個程度是為人而活.
但是那個動力在哪裡呢?.
不是無奈.
是因為愛.
其實我很相信.
今天這幾年來香港的弟兄姊妹.
你已經明白這個什麼道理了.
你為什麼會過來?.
大部分很多人都是為了自己兒女的幸福.
為了他們的前途.
為了他們的學業.
不想受一些虛假的教導.
但是可能你在香港.
有一份好的職業.
有很好的高薪厚職.
但是你來到這裡不是這樣的.
但是因為愛你的子女.
所以你放棄了你的權利.
不是被迫.
是因為愛.
愛你的家人.
很簡單.
我上個星期天就經歷到.
這個真理.

$^{601}$我想問大家.
有誰上個星期凌晨兩點鐘.
有看《英格蘭對墨西哥》的?.
請舉手.
有幾個.
我有看.
那天以為一點.
原來改了兩點.
但是我怎麼看呢?.
韓國有一套電影叫《無聲的吶喊》.
你知道什麼叫《無聲的吶喊》嗎?.
我知道.
因為我的女兒早上七點鐘上班.
六點鐘起床.
六點半出車.
如果吵醒她.
是對她很辛苦.
我太太又睡了.
所以我下廳看.
還要關上兩邊的門.
其實都沒什麼作用.
聲音都會傳出來.
我看電視的時候.
看直播.
是沒有聲音的.
我只是看畫面.
但是當我一進去.
(吶喊聲).
這叫《無聲的吶喊》.
我很興奮.
大聲叫.
但是沒有聲音出來.
最後半小時是不是很刺激?.
因為被人攻到傻了.
(吶喊聲).
最後吹長雞.
我興奮到在客廳走了幾個圈.
但是還是沒有聲音.
為什麼?.
如果我理性一點.

$^{641}$講道理我贏了.
房子我出錢買的.
電視機我買的.
電費我付的.
我星期天已經講了三堂了.
現在我看一場球賽.
有什麼問題?.
講理我是贏的.
但是不是.
你是因為愛.
你就放下自己.
但是沒有人逼的.
我女兒沒叫我這樣做.
我太太沒叫我這樣做.
是我自願.
這就是愛的激勵.
所以你為了別人的緣故.
你願不願意這樣擺上呢?.
有時候我很擔心我們作為基督徒.
我們只喜歡抽一節經文來講.
沒有看上下文.
有時候失去了那種精要.
我們最厲害的一段時間就是講.
耶穌怎麼做.
我們就要怎麼做.
當以基督耶穌的心為心.
這個是絕對正確的.
但是我們有沒有看前面那一句是什麼?.
我們一起來讀這段經文.
腓立比書第二章第四到第八節.
我們一起來讀出.
我們常常都說以基督耶穌的心為心.
接著下面就講耶穌怎麼放棄自己的權利.
為了我們.
大家有沒有留意第四節是講什麼?.
各人不要單顧自己的事.
但是要顧別人的事.
你要去關心人.
你要去幫人.
然後還用耶穌作一個榜樣.

$^{681}$就是我們放棄自己的權利.
主耶穌基督本身就是神.
受著天使天君的敬拜服侍.
但是為了我們的緣故.
反倒虛己.
成為人.
我們常常都說耶穌愛的高峰就是釘十字架.
是.
但是我告訴大家.
耶穌愛的表現就是道成肉身.
因為釘十字架.
祂只是痛幾個小時就已經死了.
但是祂一生在世的時候.
是有身體的局限的.
祂本來是零的.
不需要受著身體的局限.
我相信今天在我們當中.
很多人已經五加一.
已經申請了.
曾經有一個財經報告.
包括施永青怎麼說.
香港未來的樓市會復甦.
因為有很多在英國五加一的人.
就會回流回去.
所以樓市就會升了.
有一集永善蘭說了.
他們錯了.
因為他們不了解.
等了五年沒理由不等到那一年.
要等到那一加一我才走.
怎麼會今年馬上回去呢.
但是是不是真的回去呢.
可能會有.
但是不是全部.
因為我們如果為了兒女.
他已經融入了這個社會.
你們也已經融入了.
已經六年了.
或者五年了.
回去?.

$^{721}$不是這麼容易的了.
華富村你見到很恐怖.
有個女人打人.
弄死了人.
大家知不知道這件事.
如果你留意到.
在法庭裡面那個案件的名字.
你就知道原來那個是新香港人.
是國內人.
就是國內的拼音.
所以可以這樣來打死人.
好了.
還有就是.
耶穌他鍍成肉身.
被身體限制住.
你有沒有這個經歷呢.
我們過了來這裡.
可能我們住的地方不是超豪宅.
但是也可能有兩層.
甚至可能我們是住排屋.
但是也有個什麼呢.
花園仔是不是.
可能只是三四百萬而已.
你去到中部或者再北上.
三百五十萬.
已經有四房五房了.
但是你三百四十萬.
三百五十萬回去香港買房子.
你在市區買什麼房子.
是非常安全的.
那些叫安全屋.
因為很小.
你洗澡的時候.
遊客都不用怕.
跌倒不倒地上的.
因為太窄了.
是完全吸不到地氣的.
所以是很安全的.
那些叫安全屋.
是不是.

$^{761}$你住慣了較為闊落的地方.
要你住一些狹隘的地方.
你容易接受嗎.
你不容易的.
但是耶穌卻是成為人.
有肉體的限制.
不單單成為人.
他不是人中人.
不是在皇宮裡面出生.
是在一個貧苦的家庭裡面出生.
而他的身份不是萬人之上.
是做勞瀑.
甚至為門徒洗腳去服侍人.
一個盲眼的黑衣耶穌.
很有禮貌地問他.
我可以為你做什麼.
是一個瀑人的身份.
降卑.
這個還未貼地.
接著還要壞到什麼地步.
是被釘在十字架上.
是一個極度的重犯.
被人輕視.
耶穌原來為了我們.
一直放棄了自己應有的權利.
一步一步地向下滑.
但當然這不是一個完結.
他從死裡面復活的時候.
接著被升為至高.
接著超乎萬明者之名.
我們信主的人.
不斷去歌頌讚美我們的神.
萬物都稱他為主.
去將榮耀歸恩給神.
所以留意到原來他的放棄.
並不是一個損失.
原來是一個祝福.
原來仍然是我們個人.
不要單顧自己的事.
也要顧別人的事.

$^{801}$我們活在這個世界.
我們不可以只是自己做忠心.
不去理人.
作為基督徒.
我們願意為了主.
為了別人的緣故.
放棄自己的權利.
所以某個程度.
我們千萬不要做絆腳石.
就是說我們做的事影響了別人.
我們一起來讀.
或者說第十八章.
第六到第八節.
我們一起來讀出.
「凡是者信我的.
一個小子跌倒的.
倒不如把大磨石散在.
這人的頸項上.
沉在深海裡.
這世界有禍了.
因為將人叛倒.
叛倒人的是.
是免不了的.
但那叛倒人的.
就禍了.
倘若你一隻手.
或是一隻腳.
叫你跌倒.
就撼下來丟掉.
你缺一隻手.
或是一隻腳.
進入永生.
強如有兩手兩腳.
被丟在永火裡」.
主耶穌說到.
在這個世界上有很多的事.
可能令到人對信仰.
產生了疑惑.
這個免不了一定會有的.
不過要提醒跟隨的人.

$^{841}$你要小心.
但你不要做絆腳石.
不要做叛倒人的事.
如果真的有這樣的問題的時候.
有什麼在影響你.
或者影響別人的時候.
如果是你一隻手.
就斬了那隻手下來.
那就比你全身落地獄好.
如果是一隻腳就斬一隻腳下來.
當然我們不可以照字面解釋.
如果是這樣的時候.
在天國裡面.
全部都是傷殘人士.
要麼沒有手.
要麼沒有腳.
要麼盲.
要麼啞的.
意思就是一些很大的問題.
我們一定要好地對付.
免得成為別人的影響.
有時候有些人說自己是基督徒.
我真的很害怕.
不要了.
你這樣說還害死我們.
正如耶穌那時候.
有些污鬼婦.
你們永生神的兒子.
耶穌就喝醉不讓他說話.
為什麼?.
在過去這幾年裡面.
有很多人都說.
我們基督教很反感.
他們做這些事.
完全是違反了原則.
就好像剛才所說.
哥林多教的基督徒.
是不應該去偶像的廟宇裡面.
我們要小心.
會不會有時候我們做的事.

$^{881}$是影響別人對我們信仰的看法呢?.
今天我們記住.
有時候是我們的生活.
都影響了別人.
這張照片是我兩三年前.
我坐郵輪去挪威.
好像七日那些.
嘩.
有一群香港人.
我覺得很慘.
我不敢認我是香港人.
他們一定要走.
在上面吃自助餐的地方.
公然地玩撲克牌.
喝酒.
大聲說話.
旁若無人.
甚至說髒話.
我記得我在挪威那裡.
有一天真的要坐地下巴士上山.
他們又是在那裡坐巴士.
坐在車尾.
又是高談闊論.
全部都是粗口爛事.
我唯一有一點覺得.
這樣都沒那麼慘.
因為整個車都是西人.
應該不懂得聽那些是粗口.
但是我覺得.
你知道這樣會怎樣嗎?.
影響別人對我們的看法.
這個不是你個人的問題.
是你個人的問題.
影響別人對你甚至對你的民族的看法.
世界盃.
日本隊盡力.
輸了給巴西.
但是很多人對日本的球隊或者球迷.
是急切讚賞.
為什麼呢?.

$^{921}$原來遠自當年法國.
就是舉辦世界盃的時候.
他們已經做這件事了.
很多年之前.
就是那些球迷去到那個地方裡面.
他們完了之後.
他們是會收回乾淨的地方.
撿走垃圾.
他說我們難得來到這個地方.
我們要保持那個場地的清潔.
雖然有人說他們裝模作樣.
在家裡不做事.
甚至日本隊輸了.
但是當他們離開那個更衣室之後.
那些人都拍到那個更衣室.
是撿得乾乾淨淨的.
輸了已經很不開心.
但是他們都要撿回這些東西.
你應該去過日本的很多人.
你沒來過日本.
街上沒有垃圾桶的.
很煩的,你拿著垃圾桶撿去哪裡好?.
原來他們是習慣在家裡分類來來去去.
雖然別人對他們會有這個好感.
但是相反地.
有些人去到不同的地方.
都破壞別人的東西.
租一些B and B.
全部是骯髒的.
廁所都塞滿了.
如果你知道住在B and B你應該要收拾乾淨.
但是沒有這個東西.
甚至近期最新的那個是怎樣?.
啞巴那個人.
不是在歌劇院嗎?.
拍照拍了很久.
阿姐不要拍那麼久.
接著他就罵人.
原來啞巴挺會罵人的.
五星旗你不要欺負我們中國人.

$^{961}$我們是強大的民族.
但是這個是不是真的這樣呢?.
所以我們要很小心.
原來你自己個人做的事是影響了別人.
我們要很小心就是不要做絆腳石.
這張圖也可以做絆腳石的.
你要檢視一下.
我舉一個例子.
我幾個星期前去一間教會講道.
我都很嚴厲去責備他們.
最多你下次不要請我.
十點半崇拜.
我十點二去到.
只有十多個人.
我記得很久之前我來的時候.
都有八,九十人.
這間教會為甚麼少這麼多人呢?.
後來我坐在後面.
教會介紹我.
認識一對夫婦.
跟我坐在一旁.
老師很好.
他們這次是第三次來.
他們沒有來幾個星期了.
我跟他說.
你講道挺好聽.
所以他們肯來.
我多謝他.
但我有點尷尬.
為甚麼?.
到我講道的時候.
都差不多坐滿了人.
有八,九十人.
即是有五,六十人是遲到的.
有甚麼問題?.
這是我個人的問題.
我肯回來算了.
昨天看完英格蘭.
再看阿根廷.
你以為我不累嗎?.

$^{1001}$遲一點沒甚麼問題.
這可能是你的問題.
但你知不知道.
有些人遲到.
如果有未信的人.
怎麼看你這個信仰?.
這群基督徒.
這樣去崇拜嗎?.
這個神很低級.
你想想你坐飛機.
都要早兩小時去辦理登機.
辦理登機後沒甚麼事做.
很悶的.
你都不可以不早點去.
你有機會面試.
你會不會遲?.
你約了NHS的預約.
你可不可以遲?.
遲十五分鐘.
你的預約就會取消.
如果我們的神是超越這一切.
但我們會不會是這樣?.
所以我們做信徒.
恩事會怎麼辦?.
原來我們做了一些事.
我們不知道.
我們影響了別人對信仰的看法.
一個不信的人來到.
可能見到.
有些基督徒在看手機.
但人家不是看經文.
是看Youtube.
是呀.
牧師們就轉台了.
但你知道這件事影響到別人嗎?.
還有就是.
最小心是我們的言語.
我早幾天.
我很感恩我的屬靈孩子.
被我養大的那些.

$^{1041}$其中這一位.
已經奉獻做傳道人了.
他過來英國旅行.
住在姐姐那裡.
姐姐就在Southampton.
他就過來Redding找我.
我就問他喜歡吃甚麼.
我真的沒想過.
我喝茶.
你在香港還要喝茶?.
你早幾天還拍給我看一張照片.
在金漢宮.
在Southampton喝茶.
你來這裡又喝茶?.
我跟他說你知不知道.
第一.
我們Redding沒甚麼茶喝.
我們的飲茶都要來到Southampton.
不過最近開了一間叫YY.
在中超舖.
擺一個位.
但生意很好.
做飲茶.
但我說那裡只有一間酒樓.
所謂的酒樓.
但我不會去.
因為很擠.
食物不好吃.
八九張桌子.
背對背.
我可以聽到隔兩桌的人說話.
我說話這麼大聲.
隔三桌的人都聽到.
那時候我不想去這些地方.
免得被人聽到我們說話.
所以我給他改了吃泰國菜.
那天下午.
只有三桌人.
有兩桌都是外國人.
那我就可以暢所欲言.

$^{1081}$會不會有時我們都犯了這個錯.
我們不自知.
我們華人圈子很小.
MK有多少間酒樓.
大班.
我上星期.
我不是說假的.
我上星期都來MK廣道.
十多人約我吃飯.
就跟我談信仰問題.
去了那間.
喝珍珠奶茶那間叫甚麼.
如善坊.
我遇了三桌人.
其中一桌坐在我旁邊.
就是嘉樂山的弟兄姊妹.
牧師就在那裡.
我就認不出是哪間教會.
嘉樂山.
我就去那裡見你.
我都不敢這樣說話.
你明不明白.
原來就是隔壁那桌已經是.
有時我們在吃飯的時候.
說這個不好.
說那個不好.
說這個牧師不行.
那個牧師不行.
例如想有人聽到.
他還敢不敢回教會.
還不如我回去之後給你說.
這個就是我們.
那時做信徒的.
在我們的生活裡面.
我們做了半腳石.
所以我們為了人.
為了神的緣故.
我們要去禁制自己.
生活的問題就是這樣.
我見過.

$^{1121}$那些人去Costco買東西.
那些一盒芒果.
可以切開逐個逐個選擇.
我懷疑他以前在香港做QC.
是逐個逐個拿出來驗的.
有些蛋是叫Mix的.
他可以將幾盒蛋打開.
其實大的和小的差不多很遠.
他就真的Mix.
他覺得他Mix得不好.
他就自己Mix.
甚至我上一次去Costco.
被一個西人問我.
我就沒有做這件事.
我在看西瓜.
他問我.
為什麼你們這麼喜歡拍西瓜.
他說是什麼原因.
我說刺激它.
那些水份會多一些.
所以是要拍拍的.
他們西人買就這樣舔一舔就走.
我真的就這樣舔一舔就走.
我沒有拍.
但都已經很甜.
但有時候我們會做這些事.
但這次我們是沒有想過.
原來影響別人對我們的看法.
對我們信仰的看法.
所以我們要很小心.
作為牧者.
我們更加要小心.
因為稍後會說.
牧者其實和我們一樣.
都是神的兒女.
但我們是要更加小心.
不只是說一趟道.
很簡單.
牧者的家庭已經很重要.
很不幸地.

$^{1161}$我見到有牧者鬧離婚.
離了婚.
我認識的.
不是因為有第三者.
只是所謂性格不合.
但如果你和學生說.
我們要接納.
要寬恕.
但原來.
你不可以寬恕你的配偶.
一些生活的習慣.
甚至要離婚.
下面的大英姐妹怎麼會服你.
是不會接受.
所以我們作為牧者.
我們更加要小心.
我舉一個例子.
我在41年前.
算回41年前.
應該是43年前.
我就在信義會事奉.
那間叫馬士紀念堂.
在太子道.
現在還在.
看電影站那裡.
當時在教會裡住的人.
都是附近來的.
那些婆婆住在上海街或砵蘭街.
你知道我很疼愛婆婆.
但我不敢上去探望他們.
他們在唐樓四樓或五樓.
我不是因為怕走樓梯.
我只會跟婆婆說.
約定你.
哪天在你對面那裡.
有一間餐廳.
我請你吃東西.
不要破壞.
你上來.
我煲湯給你喝.

$^{1201}$我說不,我不能上你那裡.
為什麼?.
總之不能上.
為什麼?.
因為上面很多樓梯.
寫著小紅,小花,小梅.
在二樓,三樓.
如果我上去四樓.
被人看到.
他質問我都沒問題.
可以解釋.
但他不會質問我.
就在那裡唱.
千萬不要回靈教會.
我看見梁永善.
上去一樓的房間.
怎麼搞的?.
那我就死了.
舉著牌.
我上去探訪.
那我就更死了.
你即是此地無銀三百兩無私顯見私.
我是事奉神的.
我可以用理性說.
神是會清楚知道的.
可以神為我來解釋.
甚至身冤在主.
但可能我做這些事.
就成為半腳石.
所以要小心.
另一個例子.
我不知道有沒有人住大埔.
大埔見到一間叫星輝戲院.
當時大埔有很多間戲院.
一間叫銀星戲院.
一間叫虹基戲院.
不過全部倒閉了.
這個星輝戲院.
在我家對面.
我要學過這間戲院.

$^{1241}$就更快過由這個港台.
出去你們門口.
我一下樓.
走到對面就是了.
有一晚九點回到家.
大家跟我說.
對面做黑回歸.
對黑回歸.
不知道有沒有看過這套戲.
梁家斌轉做邪國.
好看.
我們當時沒有看過.
現在又再做.
好啊.
我說不要去看.
為什麼.
沒問題.
這套戲是沒問題.
不過這間戲院有問題.
你很熟悉.
住在大埔.
原來那間戲院.
經常都是做三級戲.
來不來一兩天換一套正常戲.
而且那晚是最後一晚.
最後一場.
看完那些半戲.
已經換了上面的海報.
我和我弟弟出來.
很搞笑.
回家的時候.
突然間見到.
丁良和黃小慕坐在那裡.
在這間戲院出來.
做什麼.
好色老公淫賤妻.
不是吧.
那我就死給你看.
我仍然說.
他質問我.

$^{1281}$我可以回答.
他質問我的.
原來這樣就會成為.
別人的絆腳石.
所以我們做牧者.
是要承擔更大的代價.
所以求神要幫助我們.
但最後記住.
看聖經不是只看一面.
我們不要做絆腳石.
但是聖經裡面很清楚說.
我們是為了別人.
願意放下自己.
這本經文我們不讀了.
為什麼耶穌.
他們有這個世界.
後來他出來工作的時候.
我要他洗禮.
施洗約翰說.
我怎麼可以和你洗禮.
我不配的.
你是聖傑.
但耶穌說.
你暫且許我.
因為我們理當.
盡諸般的義.
耶穌受洗.
都是怎麼樣.
是為了別人的緣故.
我要你去做這些事.
保羅為什麼要和提摩太行割禮.
保羅一直都說恩順清義.
但為了提摩太.
為了他要傳道.
要面對那些猶太人.
保羅和他行割禮.
所以這個就是.
甘心放下自己.
好了.
我們不可以做絆腳石.

$^{1321}$但是記住.
特別做牧者的.
我們一起來讀.
哥林多號書第六章.
第三第四兩節.
我們一起來讀.
我們凡事都不叫人有妨礙.
免得這積分被人毀謗.
反倒在太陽的事上.
表明自己是神的用人.
就如在許多的忍耐.
患難 窮乏 困苦.
保羅就說.
我們作為神的僕人.
我們盡量做一些事.
不要妨礙別人.
所以我們要想更多的事.
好了.
我們不可以做絆腳石.
但是不要忘記.
聖經提醒我們.
眾人以為美的事.
你都要留心去做.
不是每個人認為對的事.
就是真的對的事.
可能會錯的.
耶穌是盆石 是不是?.
但是不要忘記.
耶穌都是絆腳石.
不是我說的.
是聖經說的.
我們一起來讀.
彼得前書第二章.
第七第八兩節.
又說.
很清楚說到耶穌就是盆石.
我們要紮根在盆石上 是不是?.
但是聖經都說.
耶穌基督是絆腳石.
大家有沒有想過這一點?.

$^{1361}$所以看真理.
我們要看全面的.
耶穌在世的時候.
一群死守自己定的律法的.
民事和法利賽人.
是耶穌離經半道的人.
安息日.
主要就是要你休息.
去敬拜神.
但他們定了很多規矩.
耶穌在安息日.
治病趕鬼.
他們不是看耶穌做了.
這些偉大的事.
而是破壞了規矩.
耶穌所說的.
所做的.
往往就成為這群.
死守人所定的規矩的.
法利賽民事的絆腳石.
還有.
最明顯.
當耶穌去到聖殿.
見到外邦人敬拜神的地方.
竟然被人拿去.
收受賄賂.
帶去販賣牛羊.
兌換銀錢.
耶穌怎樣?.
不如這樣吧.
你再退一點.
給外邦人一些位置.
去敬拜.
不要衝突.
大家有商有量吧.
耶穌不是這樣.
耶穌把桌椅推倒.
用一整片鑿鞭子.
把牛羊趕出去.
但祂沒有打人.

$^{1401}$然後怒斥這群人.
為何把神的殿變成了.
一個賊窩.
但你想一想.
在那一刻.
人們看到耶穌這樣的行為.
怎樣看待耶穌?.
又說自己心裡柔和謙卑.
現在很柔和嗎?.
人們可能對耶穌這個行動.
會有反感.
某程度上.
耶穌所做的事.
會成為這些人的絆腳石.
但是.
耶穌不會因為這樣的緣故.
而默然不語.
今天.
很簡單舉個例子.
同性戀.
LGBT.
我仍然堅持這是一個罪.
我會接納這些人.
但我不是接納這些人的行為.
但如果你跟你的兒女這樣說.
或者你跟其他人這樣說.
人家說你不懂得包容大愛.
你這麼狹隘.
你這麼迂腐.
可能我們說這些的時候.
正如剛才保羅提醒我們.
那些人跟你說.
那些食物原來是祭過神的.
或者在街上買回來.
都是拜過偶像的.
你馬上不好吃.
要走.
關係馬上有惡劣.
但是你會不會為了討好人.
而放棄了自己的原則.

$^{1441}$所以我們要分辨.
我們怎樣做絆腳石.
剛才我們已經說.
看聖經要看全段.
剛才耶穌基督是絆腳石的時候.
是什麼呢.
就是在《彼得前書》第二章.
第七第八兩節.
第九節是我們獨立抽出來說的.
我們一起去讀第九節.
其實是跟上文有關的.
原來剛剛說完所謂的絆腳石.
接著就說我們的身份.
我們不單單是神的兒女的身份.
你留意到幾個都是一個collective land.
我們是被揀選的族類.
聖潔的國度.
屬神的子民.
是一個群體.
但是君尊的祭司就是個人.
因為你和我都是君尊的祭司.
我們要宣揚.
照我們出黑暗入奇妙光明者的美德.
我們要透過我們的生命.
去分享神的美麗.
但是不等於.
我們什麼都認同.
我有我們的立場.
可能我們做的一些事.
會被人說的.
正如我經常這樣說.
我講一篇道.
一百個人可能有九十個人會認同你.
但可能有十個人不認同.
又講那些好像政治的.
其實那些不是政治的.
我講的是事實.
但是可能就會覺得.
這個叛倒我了.
我會檢討.

$^{1481}$但我覺得我是對的.
就繼續下去.
所以求神幫助我們.
讓我們有智慧.
從神而來的智慧知道.
我什麼時候要做絕對.
什麼時候我可以妥協.
我們為了信仰.
為了別人的緣故.
我甘心的.
放下我應有的一些權利.
免得我做絆腳石.
但我們都要有智慧.
有我的立場.
我願意為真理來站得硬.
求神來幫助我們.
原來我們一個結論是什麼呢.
就很簡單.
我們最後.
《道德論》的前書第十章第三十一節.
所以你們或吃或喝.
無論做什麼.
都要為榮耀神而行.
就是說我們的生活.
是要見證我們的神.
求神幫助我們.
特別是我們生活上所謂的細節.
剛才我所說.
去超市買東西.
排隊.
說話.
就是這些.
會影響別人對我們的看法.
這個是求神幫助我們.
我低頭為禱告.
其他事項.
\newpage



\section{以西結書 47:9}
\label{sec:032450}
\textbf{得著腳掌所踏之地 (華)}~~~~劉彤牧師
\newline
\newline
連結: \href{http://www.church.com.hk/acms/content.asp?site=cdc\&op=show\&type=product\&code=032450\&layout=sermon}{\texttt{032450}} ~~~~ 語音日期: 2026-07-12
\newline
\newline
\hyperref[sec:032449]{\small{< < < PREV SERMON < < <}}
~
\hyperref[sec:index_chronic]{\small{[返順時目]}}
~
\hyperref[sec:032451]{\small{> > > NEXT SERMON > > >}}
\newline
\newline
以西結書 47:9
\newline
\begin{longtable}{cl}
\hline
\hline
章節 & 經文 (和合本修訂版)\\
\hline
47:9 & \begin{tabularx}{0.7\textwidth}{X} 這兩條河所到之處，凡滋生的動物都必存活；這水流到那裡，使那裡的水變淡，因此裡面有極多的魚。這河水所到之處，百物都必存活。 \end{tabularx} \\ \\
[1ex]
\hline
\hline
\end{longtable}
$^{1}$我們讀一節聖經,在以西結書47章第9節,.
只有一節的聖經,我們一同來念,來..
這河水所到之處,凡滋生的動物都必生活..
並且因這流來的水,必有極多的魚,海水也變甜了..
這河水所到之處,百物都必生活..
昨天晚上我們有非常圓滿的一個結束..
有多少人參加這幾天晚上的聚會的?.
我們鼓掌一下好不好?.
有多少人七個晚上都來的,舉一下手好不好?.
我們更熱烈了,謝謝他們!.
我們敬拜禱告,一同呼喊,也一同宣告,.
仇敵魔鬼一切的兼顧引領都要倒塌,Amen!.
神的復興必要降臨在我們中間,.
而今天下午我們更是要實際的走進到我們的社區,.
每一個角落做禱告行徑,.
歡迎耶穌在我們的社區城市當中來掌權,Amen!.
Hallelujah!.
你知道這次耶利哥行動對我們教會是非常有意義,.
因為這是第一次我們的聚會不是以我們自己的需要為訴求,.
而是單單為了神的國而代求,.
這幾天晚上當我們聚集在一起,.
我們為影響世界的每一個領域我們禱告,.
也為著我們所居住的城市跟社區來禱告的時候,.
忽然間在我的靈裡,.
我真的看見一支屬神的大軍隊正在我們中間興起啊!.
而昨天晚上當我們真的是跪在神的面前,.
我們真的是為美國,為我們的社區求復興的時候,.
我就想起在去年我們慶祝三十週年的時候,.
神很清楚對我說神要設立生命河在矽谷成為一個復興的基地,.
生命河水要從這裡流到美國,.
甚至從這裡流到萬國萬民,Hallelujah!.
在以細節說四十七章裡面我們看見,.
當時神帶領以細節看見一個意象,.
看見生命河水要從聖殿湧流出來,.
而且越流越遠,越流越深,.
在剛才所讀的聖經第九節裡面講到這河水所到之處,.
凡資深的動物都必生活,.
並且因這流來的水有極多的魚,海水變甜,.
這河水所到之處百物都必生活..
哇!我希望你能夠看見這個圖畫,.

$^{41}$這麼清楚的圖畫擺在我們面前,.
你是不是能夠看見這個生命河水所湧流到的地方,.
一切都要改變,Amen!.
給你媽說生命河水流到這地方都要帶來改變,跟她說,.
這是為什麼今天主日崇拜之後,我們要進入社區做禱告行經,.
因為我們深深的相信這是神的應許,.
生命河水無論湧流到哪裡,.
我們看見都要帶來改變,.
我們看見死亡的氣息要活過來,重新活過來,.
那些沒有生命的地方重新恢復生命,Hallelujah!.
那你說,為什麼我們要出去做禱告行經呢?.
難道我們在教會裡面,在家裡面禱告不可以嗎?.
當然是可以,.
我們是可以家在家裡或在教會裡面為我們所居住的土地祝福,.
為我們的社區求平安,.
但是今天如果我們想要為神得地為業的話,.
我們若是想我們所居住的土地能夠迴轉來歸向神,.
那我告訴你,我們就必須在這塊土地上來做禱告行經,.
我就想起在我們教會這麼多年,.
我們在做跨文化宣教職堂的工作的時候,.
其實我告訴你,每一次我們在每一個國家開始,.
都從禱告行經開始的,.
在1998年,當我們進入蒙古的時候,.
那個時候我們對蒙古真的一無所知,.
我還記得當我們在蒙古的城市裡面做禱告行經的時候,.
我們就開始看見蒙古的貧窮,.
男人沒有工作,酗酒,家暴,.
在路上充滿了無家可歸的孩子,.
他們叫Street Kids,街童,.
當你走在蒙古烏蘭巴特的街道,.
你看不見任何的盼望,.
我們當時就站在烏蘭巴特那個獨立紀念三岡上,.
我們站在那裡看著烏蘭巴特,.
我們就來宣告,讓生命和來永留,.
帶來復興,帶來盼望,.
我們看見如今的蒙古經濟起飛,教會也成長,.
如今我們在蒙古已經超過55間分堂了,.
Hallelujah,讚美我們的神..
在公元2000年的時候,.
我們也第一次我們的腳踏上非洲的土地,.

$^{81}$當時非洲對我們來講更是一個完全陌生的環境,.
我們既不懂他們的語言,也不了解他們的文化,.
當然他們看見我們這些華人也覺得很奇怪,.
到這裡幹什麼?.
在當時在非洲不是太多中國人在那裡,.
這些人到底要幹什麼?.
其實當時我們也不知道我們能做些什麼,.
我們只是進到他們的土地,為土地來禱告,.
一邊禱告,一邊祝福他們的社區,.
一邊禱告,一邊看有什麼神的機會為我們預備..
所以在2001年的時候,.
在二月份我們在多哥的首都Lome,.
我們就開始了非洲第一間分堂..
如今我們從一個國家已經發展到18個國家,.
從一間分堂現在發展到超過130間分堂,.
Hallelujah!.
我們真是經歷到耶穌要求一張三節繩的應許,.
神說凡你們腳掌所踏之地怎麼樣,.
我都要賜給你們..
那你說我們在禱告行徑當中,.
我們當做些什麼事情呢?.
在路加福音第十章裡面很清楚的教導我們,.
禱告行徑,出去禱告應當做些什麼?.
首先第一個,什麼是禱告行徑?.
禱告行徑就是要去站立在神工作的地方,.
跟你媽說要站立在神工作的地方..
路加福音十章第一節,.
這裡特別是你看見主耶穌拆遣了門徒就是去禱告行徑,.
這事以後主又設立七十個人,.
拆遣他們兩個兩個在他前面,.
往自己所要到的各城各地方去..
很有意思,主在這裡並沒有叫門徒只是待在家裡禱告,.
或在教會會堂裡禱告,祝福城市而已..
主是叫他們出去,為什麼要出去?.
兩個原因,第一個要了解需要,.
跟你媽說要了解需要,.
因為只有當我們站立在這塊土地上面,.
我們才能夠了解這塊土地上真實的需要..
事實上這是為什麼主耶穌倒臣肉身來到世界上的原因..
其實如果神只是讓我們知道祂很愛我們,.

$^{121}$祂只要在天上掛個布條就好了,.
嘿!世人啊!我很愛你們啊!.
但是如果神想要拯救世人,.
那我告訴你光是掛個布條是不夠的,.
祂就必須倒臣肉身來到這個世界上,.
祂才能明白我們在肉身上所面對的試探軟弱,.
祂才能夠真正擔當我們去的憂患跟痛苦嘛!.
Amen!.
所以有人說,我們必須要讓我們的腳沾上塵土啊!.
否則我們不可能為這塊土地帶來影響力的..
如果你想帶來影響力,你的腳一定要沾上塵土..
今天我們很容易在教會裡面,.
我們可以發出高音部的讚美啊!多麼榮耀!.
但我們卻一點聽不見城市裡面那個低音部所發出來那個吶喊啊!.
只有當我們的腳踏出去,走在城市,社區當中,.
我們才可能體會主耶穌憐憫的心腸,Amen!.
這第一個,了解需要..
第二個,要領受全病..
只有當我們站在這塊土地上,神才會把這個土地的全病給我們..
所以在我們剛才讀的約蘇雅經一章第三節,神對約蘇雅說,.
腳掌所踏之地,我要賜給你..
其實這句話真正的意思是,神對我們說,只要你憑著信心踏出去,.
我就會把你腳掌所踏這個地方,那個全病賞賜給你..
只有當我們擁有這個土地,神所給我們這個全病的時候,.
我們才可能得地為業,而且為這塊土地帶來真正的影響力,Amen!.
那你說我們應當往哪裡去呢?.
主很清楚在路加福音十章第一節講到,所以主採檢他們是往自己要去的地方..
主要去哪裡?主要去的地方就是他要工作的地方..
所以主在這打發門徒,在他去以前,就去到他要工作的地方..
換句話說,今天我們若是期待神要在哪裡工作,.
我們就應當先到那個地方,在那裡做禱告行經..
今天我們若期待神在我們的城市社區來工作,.
我們就必須要在這個影響城市的領域,.
站立在那裡為這個城市做祝福禱告..
已故的青年使命團創辦人Lauren Conahan來到我們教會很多次,.
她就曾提出影響城市,影響世界的七大領域..
那就是政府,商業,教育,家庭,宗教,媒體以及藝術文化..
我們就要到每一個城市的市政中心,商業中心,教育中心,.
科技文化中心,宗教中心去做禱告行經..
在我們這次發的耶利哥行動禱告手冊,.

$^{161}$如果你還沒有的話,可以在外面,好像還有幾本你可以去領取..
上面就很清楚的說明,在每一個領域上你怎麼來禱告,.
能夠拆毀堅固的淫壘,歡迎神來掌權,Amen?.
我記得Mark Geppert牧師也曾經在我們中心很多次,.
他講到有時候我們就是得故意要去到一些你想神工作的地方..
他說你要故意去到你希望,期待神要工作的地方,.
在那裡做禱告行經..
所以今天下午就是我們生命盒的故意時刻,Amen?.
Hallelujah!我們要故意去到我們期待神要工作的地方,.
要在灣區每一個角落宣告,讓聖靈的江河來湧流,Hallelujah!.
這是第一個我們要站立在神工作的地方..
第二個,要呼求神打發工人,跟你們說要呼求神打發工人..
路加福音十章第二節,主就對他們說了,.
要收的莊稼多,做工的人少..
所以你們當求莊稼的主打發工人出去收他的莊稼..
你知道今天當我們要走進社區做禱告行經的時候,.
我們一定會看見在社區裡面有很多的不同的問題,.
有道德上的問題,有家庭的問題,有貧窮的問題,教育的問題,怎麼辦?.
有的時候我們光看我們自己的生活環境,.
我們也常常看見在我們自己生活當中也充滿了無解的問題,怎麼辦?.
主在這裡教導我們做一個很重要的禱告,.
就是主啊!求你打發工人啊!.
為什麼要做這個禱告?.
因為我告訴你很重要一個觀念,神是要來工作,.
但很多時候神是藉著人來工作,Amen?.
我們看見今天當我們面對這一切問題的時候,.
我們承認我們沒有解答,.
但我們也相信主擁有這一切問題的答案..
祂不但擁有這一切問題的答案,.
而且主也知道有哪些人能夠把這些答案帶到我們中間..
所以我們向神求說,主啊!求你打發工人吧!.
因為你知道答案,.
你也知道哪些人可以把這些答案帶來,.
求你打發工人把你的解答帶到我們中間,.
帶到我們的社區來,Amen?.
主啊!求你興起更多偉大的政治家來解決社會的問題,.
求你興起更多基督徒的老師來教育我們的下一代,.
求你興起更多愛你的丈夫妻子來建立基督化的家庭,.
求你打發更多的工人傳揚你的福音,.
讓更多人心歸向真神,Amen?.

$^{201}$我們要向神來呼求啊!.
在沒有解答的地方主求你打發工人帶來答案,Amen?.
當然當我們做著禱告的時候,.
很可能神所差遣的就是你啊!.
那我們要學習謙卑下來聆聽神的聲音,.
尋求神的心意,知道我們自己當做的事..
我還記得有一位牧師曾經做見證,.
他年輕的時候,他就按照神的話常常向神禱告,.
主啊!求你打發工人收取莊稼吧!.
有一次主就回答他說,我已經聽見你的禱告了,.
而要差遣的就是你!.
不是我!不是我!.
我不是為我自己求啊!主啊!.
我是為別人求啊!求你打發別人,我不是求自己咧!.
主就對他說,如果你都不肯為我去,.
有誰肯為我而去?.
他就順從神的心意走上全職事奉的道路..
我也知道在我們中間有弟兄,.
他是參加學校的家長會,.
看見學校充滿各式各樣的困境難處,.
他就回應神的呼召,開始參與了公職的選舉,.
現在成為學區委員,Hallelujah!.
當然我的意思不是說每一次當你這樣祈求打發工人,.
神猜的就是你了,不見得啦!.
所以不要緊張,都不敢做這個禱告了..
但是我們的確需要做這個禱告,.
因為知道神擁有答案嘛!.
我們要求主打發工人啊!.
在我們所有擁有一切無解的問題上,.
求神帶來答案,Amen!.
我告訴你,不要以為我們今天不能做什麼,.
很多時候,你知道我們華人生活在美國,.
我們常常以為我們是Minority,少數族裔,.
我們怎麼能夠去影響這個主流文化的人群,.
我們能夠做些什麼呢?.
但是讓我提醒你啊,.
這幾天晚上如果你有來參加我們的耶利哥行動,.
我們每天晚上介紹不同的城市,.
你就會發現我們所有在灣區的城市,.
我告訴你,亞裔的人口都是人口比例最高的!.

$^{241}$白人是Minority,.
而且非常少數的Minority,我告訴你!.
換句話說,今天我們不再是Minority了!.
我們不再是少數族裔了!.
我們就是主流!.
我要再一次說,你可能沒聽清楚,.
我們就是主流!.
This is our land!.
這就是我們的土地!.
我們真的要醒過來啊!.
這麼多年我們覺得我們是Minority,.
we cannot do anything,.
卻不知道今天我們是主流!.
我們有責任!.
有使命!.
有託付!.
來使這塊土地變得更為美好!.
Amen!.
我們就要宣告,.
生命河所流過的地方要帶來改變!.
生命河所流過的城市都可以成為神的城!.
Hallelujah!.
我深深的相信,.
生命河就是神在矽谷所揀選的復興基地!.
我在去年領受這個復興基地的時候,.
我還很難相信,.
這一個禮拜,.
當我們在耶利哥信眾每天晚上在禱告的時候,.
我越來越看見這是可能的!.
復興是可以從矽谷開始!.
生命河是可以成為一個復興的基地!.
而你我,.
就是神手中的啟明!.
Amen!.
所以我們求神,.
打發工人!.
而這個工人,.
很可能就是你跟我!.
Amen!.
第三,.

$^{281}$要去祝福社區!.
路加福音十章第五節,.
你看見主耶穌教導他們說,.
無論進哪一家,.
先要說什麼?.
願這家平安啊!.
聽說今天我們出去,.
不是尋找這個社區的問題啊!.
也不是要定罪這個黑暗的世代,.
都不是!.
因為只有神能夠定罪!.
我們的責任,.
只是把耶穌基督的愛,.
能夠顯明給這個世代看!.
所以無論我們走到哪裡,.
我們都應當為這個社區,.
裡面所有的人,.
我們來祝福禱告!.
那怎麼樣能夠具體為他們祝福禱告呢?.
我們就是要求神的國降臨嘛!.
神的國降臨在他們身上!.
什麼是神的國?.
羅馬書第十四章十七節,.
講得很清楚!.
羅馬書十四章十七節,.
我們一同來念,來!.
因為神的國不在乎吃喝,.
只在乎公義和平,.
並聖靈中的喜樂!.
所以你到每一個地方,.
不要再找小吃攤啦!.
神的國不在乎吃喝,.
在乎什麼?.
公義和平,.
聖靈中的喜樂!.
所以每到一個地方,.
我們要求三件事情,.
講得很清楚!.
第一個,我們要求神的公義降臨!.
求神的公義降臨!.

$^{321}$求神讓這些人認識耶穌基督嘛!.
然後按照神的心意,.
能夠遵行神的話!.
因為什麼是神的意啊?.
什麼是神的公義啊?.
當我們按照神所吩咐我們去做,.
那就是神的意嘛!.
我們遵行神的旨意,.
那就是神的意嘛!.
我們求神來光照,.
讓神的公義降臨,.
讓這些人能夠認識耶穌基督,.
按照神所吩咐的去做!.
Amen!.
第二,要求神賜下和平啊!.
求神使他們與神和好,.
與人和好,.
建立正確的人際關係,.
因為只有當一個人生活在,.
一個和神心意的人際關係,.
他才能夠得到真正的祝福嘛!.
第三,要求神賜下聖靈中的喜樂!.
這個世界上人都在追求快樂,.
但我告訴你,.
只有聖靈能夠帶給人心中,.
真實滿足的喜樂!.
求神賜下信心來,.
使他們能夠得到因著信靠神,.
而得到的喜樂!.
Amen!.
求神的國降臨在他們中間,.
祝福社區!.
第四,要期待神帶來改變!.
跟你們說,要期待神帶來改變!.
路加福音十章十七節,.
當那些七十個人回來,.
歡歡喜喜地回來,.
他們就說,主啊!.
因你的名,就是什麼?.
鬼也服了我們!.

$^{361}$Hallelujah!.
弟兄姊,你可知道,.
今天當我們在地上,.
我們是可以帶來改變的嗎?.
我們可以不再被動地,.
被這個黑暗世界所影響,.
反而在耶穌基督裡面,.
我們有權柄,.
可以改變我們生活當中,.
一切屬靈的環境!.
這裡說,因主你的名,.
連鬼也服了我們!.
在原文裡面,.
這個服這個意思,.
就好像電燈的開關一樣,.
一個電器的開關!.
換句話說,.
魔鬼可以把它打開,.
可是今天我們有權柄,.
也可以把它關閉!.
換句話說,魔鬼今天雖然可以.
在我們生活環境當中,.
啟動一些罪惡黑暗的事情,.
但弟兄姊妹,.
今天在基督裡面,.
我們有權柄,.
可以關閉魔鬼撒旦一切的作為,.
關閉在社區裡面一切的作為!.
Amen!.
你說有可能嗎?.
我們這麼軟弱!.
保羅在格雷多後書第四章第七節,.
很清楚地教導我們說,.
我們有這寶貝放在瓦器裡,.
要顯明這莫大的能力是怎麼樣?.
出於神,.
不是出於我們!.
保羅講得很清楚,.
是的我們很軟弱!.
像瓦器一樣的軟弱!.

$^{401}$瓦器是很平常,.
非常脆弱!.
雖然我們這麼瓦器一樣的軟弱,.
但保羅提醒我們,.
嘿!.
在我們裡面有個寶貝呀!.
這個寶貝是什麼?.
這個寶貝是從一種出於神的能力!.
什麼能力呀?.
如果你仔細讀格雷多後書第四章,.
你明白這個能力不是指別的能力,.
乃是指的叫耶穌基督從死裡復活的大能!.
Hallelujah!.
換句話說,.
今天我們做一個基督徒,.
我們不但有主與我們同在而已,.
更是復活的主,.
帶著復活的大能,.
住在我們這裡面!.
Hallelujah!.
這個大能是勝過一切死亡權勢的大能!.
這個大能是勝過一切黑暗權勢的大能!.
這個大能是勝過一切仇敵攻擊的大能!.
Amen!.
Hallelujah!.
高揚戰士的作者Mark Gapon牧師,.
也來過我們中間很多次,.
他在他的書裡面特別提到很有意思的事情,.
他說今天主差遣我們進入世界,.
如同羊進入狼群一樣!.
哇!.
你想想看,.
一個羊進入狼群,.
多可怕的事情!.
不是一進去就被吞噬了嗎?.
但要注意,.
主耶穌是說,.
我們好像羊進入狼群啊!.
Mark Gapon提醒我們說,.
我們雖然好像是羊,.

$^{441}$但是不是隻普通的羊!.
跟你媽說,.
你不是隻普通的羊!.
什麼時候我們一讚美,.
一宣告,.
才發現原來有一個偉大的獅子在裡面!.
主與我們同在!.
那個偉大的獅子住在我們裡面!.
Hallelujah!.
魔鬼以為我們是隻羊,.
想吞噬我們,.
一抓我們才發現,.
哇!羊皮裡面是隻獅子!.
牠就逃跑了!.
Amen!.
正好像神的話說,.
那住在我們裡面的是比那在世界上的更大!.
Hallelujah!.
這為什麼今天我們可以奉主名,.
我們拆毀一切仇敵的堅固銀雷!.
個人的後書第十章,.
第四到第五節,.
我們一同來這兩節的聖經很重要,.
我們征戰的兵器,.
本不是屬血器,.
乃是在神面前有能力,.
可以攻破堅固的銀雷!.
將各樣的計謀,.
各樣攬阻人認識神那些至高之事,.
一概攻破,.
又將人所有的心意奪回,.
使他都順服嫉妒!.
這不就是今天我們在我們生活環境裡面,.
看見仇敵許多的堅固銀雷,.
挾滯人的心,.
叫人不能夠來親近神嗎?.
保羅在這裡告訴我們,.
今天在神面前有能力呀!.
我們是可以拆毀這一切的堅固銀雷!.
Amen!.

$^{481}$什麼是堅固銀雷?.
堅固銀雷是在每一個地方,.
都有不同的攬阻神的工作嘛!.
就好像戴伊底的時候,.
他禱告尋求神,.
只等到21天天使才讓他顯現,.
天使就告訴他說,.
他禱告其實第一天神就聽見了,.
跟你們說神第一天就聽見了,.
但他為什麼等了21天呢?.
因為天使有波斯的魔君攬阻,.
這個波斯的魔君就是一種堅固銀雷呀!.
當然堅固銀雷不只是邪靈的工作,.
很可能是在我們思想中理性的束縛,.
也可能在一些地區有些偶像的敬拜,.
異教的敬拜的束縛,.
或者是一些社會潮流,文化潮流的影響,.
每個地方都有他黑暗的一個罪惡的一面,.
但是今天神給我們一個權柄,.
藉著禱告我們能夠推開黑暗的權勢,.
讓人的心靈得到自由跟釋放,Amen!.
所以今天當我們出去禱告行進的時候,.
除了要看見每個地方有不同的需要之外,.
更重要的,歡迎主耶穌到這個地方來掌權呀!.
因為我告訴你,感觸黑暗最好的方法,.
不是斥責黑暗,.
你光是斥責黑暗,黑暗不會離開,.
要感觸黑暗最好的方法,.
是歡迎光照進來,不是嗎?.
光一照進來,黑暗就自然而然地消失了,Amen!.
所以當我們在每個地方做禱告行進的時候,.
我們就歡迎聖靈降臨,.
歡迎主耶穌在這裡來掌權,Amen!.
Hallelujah!.
第五點,今天下午,你就要走出去進入社區,.
做禱告行進,.
在不同的領域裡面帶來祝福,帶來恩典,.
你興奮嗎?.
當你出去的時候,你要記得,.
你就是這條生命河,Amen!.

$^{521}$跟你們說,你就是生命河!.
就像我們一開始讀的聖經裡面,.
以細節書47章講到,.
生命河水從聖殿開始流出來,.
越流越遠,越流越深,.
但無論流到哪裡,.
今天無論你走到哪裡,.
你就要來宣告,.
生命河要帶來改變,Hallelujah!.
死亡的氣息要重新活過來,Amen!.
枯乾的樹要重新結果子,Amen!.
苦澀的水要變成甘甜,Amen!.
甚至連死海都變甜,重新充滿生命,Amen!.
無論你走到哪裡,.
你就要宣告,生命河來永流,Hallelujah!.
就算你說你不會禱告,.
就說這句話就可以了,.
生命河來永流,That's it!.
這是最偉大的宣告,Amen!.
當我們走到哪裡,.
做這個簡單的宣告的時候,.
我們就要看見神就要為我們興起,.
在我們的社區,城市帶來改變,Amen!.
我們一起禱告!.
謝謝主!.
你是在全地上掌權的神,.
你也在我們的社區,.
在我們的生命,.
在我們的家庭,.
在這個國家掌權的神,.
今天當我們走出去的時候,.
我們就是要宣告,.
讓生命的江河來永流,.
從我們的生命開始永流,.
流進我們的社區,.
流向這個國家,.
我們要宣告,.
這是你給我們的土地,.
這個土地是屬於你的,.
願你的榮光照耀下來,.

$^{561}$生命河水永流的地方都要帶來改變,.
謝謝我們的主,.
讓我們生命河真的成為一個祝福的管道,.
讓我們真的走在社區裡面,走在人群當中,.
我們就成為你施恩典的一個channel,.
主要願你的聖靈大大充滿,.
澆灌我們,使用我們,.
真的讓生命河果然成為一個復興的基地,.
讓你的恩典祝福從這裡流向社區,.
美國甚至萬國萬民,.
謝謝主,.
將一切榮耀歸給你,.
Amen!.
\newpage



\section{哥林多前書 11:2-16}
\label{sec:032451}
\textbf{再思女性的事奉 (2)}~~~~陳崇基牧師
\newline
\newline
連結: \href{http://www.church.com.hk/acms/content.asp?site=cdc\&op=show\&type=product\&code=032451\&layout=sermon}{\texttt{032451}} ~~~~ 語音日期: 2026-07-12
\newline
\newline
\hyperref[sec:032450]{\small{< < < PREV SERMON < < <}}
~
\hyperref[sec:index_chronic]{\small{[返順時目]}}
~
\hyperref[sec:032452]{\small{> > > NEXT SERMON > > >}}
\newline
\newline
哥林多前書 11:2-16
\newline
\begin{longtable}{cl}
\hline
\hline
章節 & 經文 (和合本修訂版)\\
\hline
11:2 & \begin{tabularx}{0.7\textwidth}{X} 我稱讚你們，因為你們凡事記得我，又堅守我所傳授給你們的。 \end{tabularx} \\ \\ \relax
11:3 & \begin{tabularx}{0.7\textwidth}{X} 但是我要你們知道：基督是男人的頭；男人是女人的頭；神是基督的頭。 \end{tabularx} \\ \\ \relax
11:4 & \begin{tabularx}{0.7\textwidth}{X} 凡男人禱告或講道，若蒙著頭，就是羞辱自己的頭。 \end{tabularx} \\ \\ \relax
11:5 & \begin{tabularx}{0.7\textwidth}{X} 凡女人禱告或講道，若不蒙著頭，就是羞辱自己的頭，因為這就如同剃了頭髮一樣。 \end{tabularx} \\ \\ \relax
11:6 & \begin{tabularx}{0.7\textwidth}{X} 女人若不蒙著頭，就該剪了頭髮；女人若以剪髮剃髮為羞愧，就該蒙著頭。 \end{tabularx} \\ \\ \relax
11:7 & \begin{tabularx}{0.7\textwidth}{X} 男人本不該蒙著頭，因為他是神的形像和榮耀；但女人是男人的榮耀。 \end{tabularx} \\ \\ \relax
11:8 & \begin{tabularx}{0.7\textwidth}{X} 起初，男人不是由女人而出，女人卻是由男人而出。 \end{tabularx} \\ \\ \relax
11:9 & \begin{tabularx}{0.7\textwidth}{X} 而且男人不是為女人造的，女人卻是為男人造的。 \end{tabularx} \\ \\ \relax
11:10 & \begin{tabularx}{0.7\textwidth}{X} 因此，女人為天使的緣故應當在頭上有服權柄的記號。 \end{tabularx} \\ \\ \relax
11:11 & \begin{tabularx}{0.7\textwidth}{X} 然而，照主的安排，女人不可沒有男人，男人也不可沒有女人。 \end{tabularx} \\ \\ \relax
11:12 & \begin{tabularx}{0.7\textwidth}{X} 因為女人原是由男人而出，男人是藉著女人而生；但萬有都是出於神。 \end{tabularx} \\ \\ \relax
11:13 & \begin{tabularx}{0.7\textwidth}{X} 你們自己要判斷，女人禱告神，不蒙著頭合宜嗎？ \end{tabularx} \\ \\ \relax
11:14 & \begin{tabularx}{0.7\textwidth}{X} 你們的本性不也教導你們，男人若留長頭髮是他的羞辱嗎？ \end{tabularx} \\ \\ \relax
11:15 & \begin{tabularx}{0.7\textwidth}{X} 但女人留長頭髮是她的榮耀，因為這頭髮是給她蓋頭的。 \end{tabularx} \\ \\ \relax
11:16 & \begin{tabularx}{0.7\textwidth}{X} 若有人想要辯駁，我們卻沒有這樣的規矩，神的眾教會也沒有。 \end{tabularx} \\ \\
[1ex]
\hline
\hline
\end{longtable}
$^{1}$所以今天我們繼續講一下處理另外一段經文.
這些經文都是比較難的.
首先跟大家解釋清楚.
雖然比較難.
或者你剛剛來到教會.
或者你可能還沒有接觸過聖經.
你發現這些經文裡面講的東西都不知道它講什麼.
這個是很正常的.
首先大家不要被它嚇怕.
因為講到最後的時候大概25分鐘左右.
那個應用一定是跟大家有關的.
就算你聽不明白.
最後的時候一定會有大家應用的地方.
不過雖然打開這些經文的時候.
就算你讀了聖經很多年.
都可能不明白.
有沒有同感.
聖經引使是很難明白的.
原因不是因為它特意寫得這麼難.
原因就是我們跟聖經的受眾.
那個時代和文化是有不同的地方.
正如有不同的地方.
所以我們需要下一點功夫才能夠明白.
當時的人他們怎樣了解經文的意思.
然後我們才今天可以說.
有些什麼原則裡面是藉著當時的文化.
當時的語言來表達的.
那些原則可以應用在我們現在這個時代.
所以大家用一用腦筋.
大家看完世界盃看得這麼精彩.
應該腦筋轉數很快.
OK,那我們來吧.
我選了這段經文是哥林多前書第十一章.
第十一章第二至十六節.
比較長一點.
這麼多字.
不過不要緊,我帶大家一起看.
十一章第二節.
這裡哥林多前書是寫給哥林多的教會.
或者哥林多的宗教會.

$^{41}$當時的一些家庭教會.
這封書信裡面說了很多其他東西.
但到十一章的時候.
他就討論到蒙頭的問題.
女人要蒙頭.
有veil over her head的問題.
十一章第二節.
我稱讚你們.
因為你們事事紀念我.
又堅守我傳給你們的傳統.
但我希望你們知道.
基督是每個男人的頭.
男人是女人的頭.
神是基督的頭.
每個男人禱告或宣講神的信息.
如果蒙著頭就是羞辱自己的頭.
每個女人禱告或者宣講神的信息.
如果不蒙頭的話.
就是羞辱自己的頭.
因為這如同剃了頭一樣.
女人如果不蒙著頭.
就應該把頭髮剪掉.
女人如果以剪髮剃頭為可恥.
就應該蒙著頭.
男人不應該蒙著頭.
因為他是神的形象和榮耀.
女人卻是男人的榮耀.
因為男人不是從女人而出.
女人卻是從男人而出.
並且男人不是為女人而做.
女人卻是為男人而做.
因此為天使的緣故.
女人頭上應該有權柄的記號.
但是在主裡面.
女人不可沒有男人.
男人也不可沒有女人.
這個我贊成.
因為正如女人是從男人而出.
照樣男人也是藉著女人生出來的.
但萬有都是出於神.

$^{81}$你們自己判斷.
女人向神禱告.
不蒙著頭恰當嗎?.
難道你們的本性沒有教導你們男人?.
若留長髮是他的羞辱.
女人若留長髮卻是他的榮耀嗎?.
因為刺激他的頭髮.
就跟蒙頭的外衣寬領一樣.
如果有人想爭辯.
我們可沒有這種習俗.
神的宗教會也都沒有.
這段經文好聽嗎?.
你不知道他在說什麼.
但先不要被他嚇怕.
因為我們不知道他在說什麼.
不代表他沒有意思.
我們不知道他在說什麼.
或是爭辯的時候.
原因可能是我們跟他的文化差距不同.
或是我們不知道他們當時的處境.
有時候看聖經.
特別是看新約的書信.
就好像聽人講電話.
你只聽到一邊.
聽不到另外一邊.
你只聽到另一邊.
他就是這樣.
是的.
你說得對.
那邊到底在說什麼呢?.
你不是很清楚.
所以我們不需要被他嚇怕.
首先我們想想.
到底他當時的背景.
或是他那邊的電話在說什麼呢?.
以致他這邊有這樣的回應.
而我們聽到他的回應.
保羅就是在回應.
哥林多教會裡面的一些問題.
哥林多教會裡面有很多的問題.

$^{121}$當然.
在哥林多的書信裡面.
你也知道這個教會是問題多多的教會.
所以叫哥林多.
所以我們理解一下他的前文後理的時候.
我們就可能會理解多一點.
剛才我們講那段經文.
我讀出來的時候.
大概你get到的意思.
女人需不需要蒙頭?.
是的.
大概get到的意思是這樣的.
正因為這樣的緣故.
所以有部份教會.
起碼我聽過的一些小部份的教會.
他們真的會這樣做.
他們就會說.
女人講道.
女人上台講道.
或者去做一些神的話語的教導的時候.
她就需要頭上有一條頭巾.
就包一條頭巾上去.
其實是表面的意思去應用經文的意思.
就是這麼簡單.
所以表面上我們就覺得.
應該是這樣處理的.
所以有些教會就會照辦主管.
經文講什麼我們就照樣去apply.
照樣這樣講.
所以這段經文我們一般的理解就是這樣的.
一般理解.
一般理解就是為什麼他們蒙頭呢?.
就是因為應該是女人的問題.
女人有什麼問題呢?.
女人的問題可能就是.
你可以想像一下當時的背景.
當時以男為上的文化.
以父權的社會.
像我們今天比較男女平等的文化.
你可以理解.

$^{161}$當時有些哥林多的女性.
她們信了主.
然後她們從那些父權的文化裡面.
以男為上的文化裡面.
她們聽到這個福音.
然後她們發現了保羅所講的福音.
耶穌基督的福音原來給我們一個新的身份.
在基督裡面我們大家都是同一的.
在基督裡面我們大家男女是平等的.
所以當時女人可能就是公開場合的時候.
她們就需要蒙頭.
然後她們去到教會的時候.
她們就覺得因為我們有這樣的身份.
在基督裡面我們就有一個新的身份.
所以我們索性就除了我們的頭巾.
然後就披頭散髮.
因為我們已經有了新的身份.
新的身份.
所以可能有這樣的問題.
所以保羅就說不如你們蒙頭.
不要破壞大家公開的見證.
因為敬拜,聚會是一些公開的東西.
所以如果有些人進來的時候.
在社會裡面進來的時候.
他見到那些女人是這樣的時候.
這樣就很忌諱的.
這樣就很麻煩的.
正因為這樣的緣故.
所以通常我們的解說就是說.
一定是女人的問題.
是女人有問題.
所以就需要要求她去到蒙頭.
如果你這樣想的時候.
其實也不是那麼麻煩的.
我覺得.
不是那麼難明白的.
你可以想像一下.
如果你是一些重男輕女的環境裡面出生的話.
女人經常被歧視的.
男人經常看低女人的.

$^{201}$其實我們華人的文化裡面.
我們傳統都有這樣的思維的.
女人無才便是德.
對不對.
如果你拿了幾個PhD你是女人的.
為什麼要拿這麼多PhD.
男人就很自然的.
我們學識上面是高一點等等.
女人就應該怎樣.
無才便是德.
你少認識一些東西.
或者你只是躲在廚房裡面就好了.
等等等等.
或者女人一定要跟尾的.
如果你經常聽這些東西.
或者這種環境裡面出生的時候.
其實都很大壓力的.
有沒有同感.
有吼.
如果你是女孩子.
從小到大你要做某一件事.
或者你在一個box裡面一個框架裡面.
你才會受人接納.
或者受人尊重的.
你一聽到這個福音的時候.
你就會覺得一個解放.
一個liberation.
原來不需要這樣的框架.
因為這些框架不一定是代表真理的.
耶穌基督的福音就告訴我們.
男女平等.
教會是一個新的社會.
一個基督裡面新的人類.
無分貴賤的.
無分男女的.
是很吸引的.
對不對.
你知道早年的時候.
大概200年前.
西方的宣教士去到東方的時候.

$^{241}$譬如進到中國的時候.
他們第一.
其中一件事他們想廢除的是什麼.
就是女人紮腳.
是不是.
因為西教士覺得這些太殘忍了.
太殘忍了.
而且這樣是控制女性的.
所以第一件事他們希望能夠廢除這些東西.
以及建立一些女子中學.
女子的小學,中學等等.
因為女孩子是不能讀書的.
是不是.
所以很多西教士去到東方的時候.
或者去到這樣的國家的時候.
他們就想改變那個社會的制度.
是什麼緣故呢.
是因為福音的緣故.
是因為福音給他們的精神.
所以他們覺得男女是要平等的.
所以我們看這段經文的時候.
我們就可能會想.
一般的解說就是女人有問題.
女人應該要蒙頭.
因為當時的髮色是這樣的.
這些就是我們西方的.
西化了的社會.
我們不是很明白的.
因為我們通常.
衣著,衣服,頭髮,長短.
這些我們就不太理會的.
或者我們比較開放一點.
但是比較保守的一些地區.
或者東方的社會就會比較緊張一點.
譬如去中東.
去中東.
你進去一些特別的聖堂的地方.
你就不可以露這麼多肉的.
是不是.
皮膚不可以露這麼多.

$^{281}$上一次我們去土耳其的時候.
我就經常吩咐一些弟兄姊妹.
我說我們今天會去聖堂.
特別吩咐一些姐妹.
就會說你們穿衣服就要保守一點.
雖然土耳其是一個很現代化的地方.
一個很現代化的地方.
但是你去一些聖堂.
譬如清真寺.
你就需要包得保守一點.
因為那裡有它的規矩.
當然你要去respect人家.
respect人家的時候.
你才可以進去.
誰知道我忘記了.
我就穿了短褲.
我忘記了,一不給別人吩咐.
我就因為熱.
我就很自然就穿了短褲.
去到門口的時候.
那個入口的看監.
他就說拿一塊布給我.
你圍住他.
這次我有一點點差池.
因為我通常都記得.
這麼多年我都記得.
結果我就要圍一條這樣的布.
我就發現很多人和我拍照.
你看到很多人和我一起拍照.
我很開心.
不是因為他們覺得我上鏡.
而是因為我穿了他們的裙子之後.
我就覺得特別上鏡.
所以東方的社會.
有時候會比較保守一點.
好不好看?.
挺好看吧?.
好看嗎?.
穿多一次是吧?.
我應該上台的時候.

$^{321}$穿這條這樣的東西.
尊重一下大家.
所以我想明白.
衣著或者外表.
在社會裡面確實是比較緊張一點.
所以你可以明白.
一般理解亂經民的時候.
就是女人有問題.
她們索性非投散法.
或者在聚會裡面.
保羅就覺得你過了火.
你會有損教會在社會裡面的聲譽.
所以保羅最緊張的是什麼?.
保羅最緊張的是.
基督徒的敬拜.
我們的敬拜.
是一個公開的場合.
我們就應當要有一個公開的.
好的見證.
不要被人誤會福音.
所以他就要求女人是蒙頭的.
這個是一般的解說.
但是如果你想清楚.
裡面有一點點這段經文.
裡面有一些前言不對後語.
保羅到底是要蒙頭還是不要蒙頭呢?.
我們其中一個方法.
我們就可以找一下.
當時的社會裡面.
到底女人是不是真的個個都要蒙頭呢?.
今天你可以去博物館那裡.
看一下羅馬的雕像.
那些頭的雕像.
有男有女的.
是不是個個女人都是蒙頭呢?.
又不是.
是不是男人個個都不蒙頭呢?.
又不是.
有時候男人又蒙頭.
所以就有一點點混合的證據.

$^{361}$到底當時的習俗是怎樣的呢?.
就比較難去找到一個一致的背景的看法.
或者可能我們看聖經的時候.
看得不通透.
會不會其實不應該這樣看呢?.
我自己一向.
以前這麼多年.
我都是看這段經文.
我都覺得應該是那個解說.
就是那麼多的女性是有問題.
就是覺得被解放.
然後她們就很自由很開放.
保羅就在那個情況裡面去做一些改正.
但是呢.
會不會有另外一個提問我們可以發出呢?.
就是保羅這裡會不會其實不是說女人有問題.
其實可能是男人有問題呢?.
會不會是這段經文裡面.
保羅是在引用哥林多的一些說法.
或者哥林多男人的一些說法.
然後他在引用當中就去到.
加插他自己的反駁在那裡.
這個早年的時候.
我就看了一本書.
就拜讀一個女色經家.
是英國人來的.
Lucy Papiat.
她就認為.
不應該這樣看這段經文的.
那為什麼要看一個女色經家寫的書呢?.
就很簡單.
因為連色經界裡面.
都很多時候是男人寫書多的.
可能你的印象都是.
但是近年裡面已經越來越多好的.
非常頂尖的女色經家出來了.
女色經家問的問題就和我們男人問的問題.
很多時候自然就不同的.
因為如果我看哥林多詞書十一章.
我們覺得女人有問題的.

$^{401}$我作為一個男人.
我當然覺得很舒服的.
當然女人有問題.
我們經常都是這樣.
我們覺得只有女人有問題.
但是如果女色經家講的時候.
可能她的角度就會很不同.
她就會問.
可不可以是另外一邊呢?.
這些絕對不是.
勉強去將字句曲解的.
而是用另一個角度去看一看.
會不會有這樣的可能性呢?.
是保羅去引用哥林多男人講的話.
然後再加以反駁.
其實哥林多前書或者其他保羅的書信.
都會引用一些所謂對手講的話.
譬如羅馬書等等.
在哥林多前書都有的.
譬如哥林多前書六章十二節.
他說「凡事都可行,但不是凡事都有益處」.
「凡事都可行,但無論哪一件我都不受他克制」.
你見到那個引號加了上去的.
那個原文是沒有的.
以前我年輕的時候.
看中文聖經或者甚至英文聖經.
都沒有那個引號的.
沒有引號的時候我們就會覺得.
當然是保羅出口了.
因為這封信是他寫的.
保羅一定是六章十二節的時候.
他說「凡事我都可行,但不是都有益處」.
凡事我都可行.
在基督裡面凡事都有自由.
凡事都可行.
但無論哪一件我不應該被他克制.
這句應該是保羅出口.
但現在慢慢我們開始有一個共識.
原來保羅有一些修辭上的東西.
是在他的書上用的.

$^{441}$凡事都可行.
其實應該是哥林多教會的口號.
因為哥林多教會男人特別.
他們就會覺得我們信了主.
我們就是anything we can do.
everything is permissible.
所有事情我們都可以做.
有什麼所謂.
我們基督裡面有自由.
但現在我們慢慢有共識.
其實保羅是引用哥林多教會裡面所說的話.
然後他就反駁.
你們說凡事都可行.
但不是凡事都有益處.
你想清楚.
不是每件事都可以做.
你最重要記住.
那件事不要被他克制.
世界盃可以看.
但不可以被世界盃控制你.
你的情緒也好等等.
所以他是在引用哥林多的說話.
但用這個角度看的時候.
你再看哥林多的前書.
可能就不同了.
我在這裡再讀這段聖經.
橙色的地方.
可能是哥林多的男人說的.
第二節.
我稱讚你們.
因為你們時時都紀念我.
凡事都堅守我給你的傳統.
但我希望你們知道.
基督是每個男人的頭.
男人是女人的頭.
神是基督的頭.
你們就說.
每個男人禱告.
或者宣講神的信息.
如果蒙著頭就是羞辱自己的頭.

$^{481}$如果女人講道.
或者禱告宣講神的信息.
如果不蒙著頭.
就是羞辱自己的頭.
就如同剃自己的頭一樣.
保羅就說.
保羅就在這裡說.
這樣的話.
如果你們說了這些話.
這樣的話.
那就第六節了.
女人如果不蒙著頭.
就應該將頭髮剪掉.
女人如果以剪髮剃頭為可恥.
就應該蒙著頭.
如果你們這樣去說四五節的話.
第六節你們就自然有這樣的結論.
第七節.
你們多不多人這樣說.
男人不應該蒙著頭.
因為他就是神的形象和榮耀.
女人卻是男人的榮耀.
因為男人不是從女人而出的.
女人卻是從男人而出的.
並且男人不是為女人而做.
女人卻是為男人而做.
這些你們說的.
這些你們多不多人說的.
因此為天使的緣故.
女人的頭上就應該有權柄的記號.
這些全部橙色的都是你們說的.
保羅就去跟他們說.
如果你們說這些話.
結論就會是這樣.
我們繼續看下去.
現在我就要矯正你.
十一節.
保羅就說什麼呢.
他說但是在主裡面.
女人不可沒有男人.

$^{521}$男人不可沒有女人.
我們每個人都有母親.
我們是需要女人的.
所以男女是平等的.
十二節.
因為正如女人是從男人而出.
照樣男人也是藉著女人生出來.
但當然萬有就是出於上帝.
你們自己判斷一下.
這裡就直接保羅跟哥林多教會說.
你們自己的判斷.
女人向神禱告.
不蒙著頭恰當嗎.
當然是恰當.
不蒙著頭當然是恰當.
為什麼.
難道你們的本性沒有教你們.
男人如果留長頭髮就是他的羞辱.
女人如果留長頭髮卻是他的榮耀嗎.
因為刺逼她的頭髮.
就是跟她蒙頭的外衣寬領一樣.
因為她已經有長頭髮了.
所以她已經好像她的外衣蒙頭一樣.
所以就不需要再蒙頭.
第十六節.
這裡就最清楚.
如果你們想再爭辯的話.
我們使徒們可沒有這樣的習俗.
神的宗教會.
Churches.
其他地方的Churches也沒有.
OK.
所以希望梳理了一點.
這些問題.
到底問題在哪裡.
如果你這樣看的時候.
你就發現.
他去到引用哥林多的對手說的話.
那些男人說的話.
然後他予以去反駁.

$^{561}$就說你們這樣想是不對的.
你們這樣想.
最後你的結論都是自相矛盾的.
所以他就說.
我們根本就沒有這樣的習俗.
宗教會也沒有這樣的習俗.
誰有這樣的習俗.
哥林多你們教會有這樣的習俗.
我們這樣看的時候.
這個看法有什麼可以梳理呢.
起碼首先保羅不會前言不對後語.
好像又蒙頭又不蒙頭.
原來他就是在反駁.
這樣的看法的時候.
保羅就是和哥林多的大男人神學.
去做爭辯.
哥林多教會有大男人.
他就用了.
曾經保羅說過.
男人是女人的頭.
我們就想去work out這個神學.
其實就是壞鬼神學.
就是大男人神學.
就是我是男人.
我就不用蒙頭.
但是如果女人要講道.
或者要禱告的話.
就請他蒙頭.
因為女人頭上一定有一個記號.
就是她要可以信服的.
對不對.
所以在反駁的當中.
保羅就用他的本性來帶出.
女人的長髮就自然是榮耀.
自然是honor.
他就不需要再加上一塊根來蒙頭.
也解釋到最後那句.
最後那句是什麼呢.
就是神的宗教會也沒有這樣的習俗.
是不是.

$^{601}$如果你這樣看的時候.
到底是男人有問題.
還是女人有問題.
如果你說女人有問題的話.
那你就很難.
就有些東西卡住了.
前言不對後語.
還有最重要是十六節.
如果有人和我去爭辯的話.
我們沒有這樣的習俗.
神的宗教會也沒有這樣的習俗.
是不是.
那會不會應該轉一轉那個角度呢.
其實可能是男人有問題.
是不是.
是男人有問題.
就是多不多有一班大男人.
就弄了一個大頭神學出來.
他們就要女人蒙頭.
這些就是我們開頭的時候說.
某些教會今天看這段聖經的時候.
他們就照主管去應用.
如果這樣應用的時候.
除了表面去看聖經.
也不是很明白那個上下文.
或者那個背景.
是不是.
那就會重複了.
哥林多教會的大男人蒙頭神學.
所以保羅寫這段話其實是在矯正這件事.
十六節也講得通.
我們沒有這樣的習俗.
神的教會也沒有.
哥林多為什麼會有呢.
所以整段經文其實針對的就是大男人.
哥林多的大男人.
如果你這樣看的時候.
就配合哥林多的前書整本書.
因為整本書裡面都針對很多領袖.
或者愛權的人.

$^{641}$他們想做的事.
和他們令到教會混亂的一些事項.
十一章就是處理女人蒙頭的問題.
如果這樣看的時候.
這個是男人的問題.
保羅就是斥責他們.
然後就講到聖餐的問題.
我們平時守聖餐的時候.
不是會讀一段聖經.
在十一章後面出來.
那裡就有一些混亂的情況.
因為他們就分開貧富.
守聖餐有些富有的人早到了.
他們就吃飽東西.
然後貧窮的人他們下班的時候.
他們來到聚會的時候.
他們就沒有東西吃.
他就去改正他們聖餐所有的混亂.
十二章就講到才幹和恩賜的使用.
就是保羅就去罵哥林多的人.
他說你們有一些的恩賜.
譬如謊言的恩賜就看低其他的人.
他就去予以改正.
十三章就是愛篇.
愛是恆久忍耐又有恩賜那篇.
很出名的那篇.
經常婚禮會用的那篇.
為什麼他會講這些.
因為上文下就是講你們都不懂愛.
我現在告訴你什麼是愛.
愛就是恆久忍耐又有恩賜等等.
是為其他人緣故的.
不是爭權奪利的.
然後十四章十五章就去處理其他的問題.
你這樣看整封信其實保羅就好像一個家長.
那些教會開始長大了.
他就不只是告訴他們do's and don'ts.
他就是告訴他們.
有一個新的人類在基督裡面誕生了.
這個就是教會.

$^{681}$在基督裡面有的弟兄或者姐妹.
我們都是平等的.
我們是按恩賜來事奉的.
因為上帝要在我們中間有一個新的創造.
但是那麼多的教會.
可能就是很多大男人.
愛權的領袖等等.
他們就做了很多這樣的東西.
在制度上或者習慣上.
就會令到其他人覺得他們是二等公民.
保羅就在這本書信裡面來做改正.
所以如果你這樣看的時候.
似乎就不是要女人蒙頭.
而是保羅相反.
就是不應蒙頭.
不應該蒙頭.
而那些大男人為什麼搞這樣的東西出來.
宗教會都沒有這樣的習俗.
為什麼你要搞這樣的東西出來.
因為你們亂去應用聖經.
這個值得我們去思考一下.
怎樣去看聖經.
有幾點我們可以點出來.
第一就是結論的時候.
第一就是在基督裡面.
男與女.
我們有的恩賜 有的才幹.
是應當可以拿來服侍的.
如果有一個女性 有一個姐妹.
她受上帝的呼召.
去信仰群體.
去確認她有這樣的恩賜的時候.
我是何許人去攔阻那位姐妹.
如果她有教導恩賜.
如果她有識經的恩賜.
如果她有講道的恩賜.
她有蒙上帝去呼召她.
又有一個好的見證.
其他人對她的群體上對她有引證的時候.
我是誰去阻止她.

$^{721}$Am I God enough.
Of course not.
在基督裡面這個新的群體.
在基督裡面我們能夠平起平坐.
是按照我們的恩賜.
不是按照我們的性別.
來去事奉.
這是基督有的意思.
第二件事我們想到的.
在讀經裡面.
我們看聖經裡面.
或者應用聖經的時候.
我們真的要去下一點功夫.
明白一下上文下來.
不要只是抽一些東西.
然後我們胡亂去應用.
這樣說不是將聖經困在當時的文化.
而是我們理解聖經是藉著當時的文化.
透過當時的背景.
來轉達上帝永恆的信息.
所以我們一定不能跳過那一步.
我們不理當時的背景.
我曾經說過.
我覺得每一次看聖經的時候.
都好像去旅遊.
去一個你陌生的地方去旅遊.
一個從來沒有去過的地方去旅遊.
你就聽一下別人說什麼.
聽一下別人的習俗.
聽一下別人的文化背景是怎樣.
聽一下那裡的信息是怎樣.
別人是怎樣接收.
看聖經其實都是一樣.
你去哪裡旅遊.
去一個聖經的世界.
去那裡聽一下別人說話.
明白別人的背景.
然後知道上帝所轉達的永恆信息是怎樣.
這是第二點.
第三點就是我們需要反問我們自己.

$^{761}$有什麼藉著聖經可以反問我們自己.
到底我們有沒有濫權.
我們的實踐上.
甚至我們的制度上.
我們有沒有濫權的問題.
如果我們濫權的時候.
或者我們胡亂用我們的權柄的時候.
或者限制其他人的時候.
那這個跟基督的福音是不是有違反呢.
我們自己生活.
我們有沒有像基督那樣謙卑呢.
有沒有濫用我們自己的權柄或者地位.
第四點.
因為我們講到女性.
我們不妨思想一下.
在我們人生裡面.
對你生命有影響的女性.
你想到有誰嗎.
你想一下她的樣子.
想一下她的名字.
想一下她做過什麼.
可能是牧師.
可能是傳道人.
可能是導師.
或者朋友的姐妹.
在你生命裡面有很多很多這樣的人.
其實是祝福過你的.
這些都是上帝所用的生命.
去感激他們.
對不對.
特別可以跟大男人說話.
在我們生命裡面有沒有這樣的人.
我想起我媽咪.
想起在我教會以前特別年輕的時候.
曾經傾注過她的生命.
在我的成長裡面.
這些可能是女傳道人.
可能是一些導師.
可能是身邊的一些姐妹.
對我生命裡面有影響的.

$^{801}$謝謝上帝.
我們低頭祈禱的時候.
我們不妨思想一下這個問題.
給大家一點時間.
有沒有一些姐妹.
一些女性.
在生命裡面有一個影響.
如果有的話.
上帝多謝你.
因為你做男做女.
在基督的救贖裡面.
在基督的恩典當中.
我們無分地位.
你去呼召男去服侍你.
也呼召姐妹.
在歷世歷代裡面有很多女性.
是你自己親自呼召出來的.
每一個都是被你所使用.
在生命當中.
你藉著這幾天.
去祝福我們生命.
去造就我們.
我們獻上感恩讚美.
雖然我們只是說一些.
女性事奉的經文.
但都是對我們每一個人.
有形容的.
你幫助我們.
要我們用我們的恩賜.
去改變我們.
矯正我們心裡面.
或者在我們背景裡面.
有一些不對的事情.
思想.
主我們求你.
求你讓我們在基督裡面.
我們真的享有這個自由.
在基督裡面.
我們能夠在我們角色功能上.
在我們服侍上.

$^{841}$我們沒有分男與女.
讓我們一起在神的國度裡面.
我們去服侍你.
有如基督去服侍教會一樣.
為了我們生命裡面.
有一些影響過我們的女性.
我們獻上感恩讚美.
多謝他們謙卑的服侍.
多謝他們所說的話.
多謝他們所做的事.
甚至默默無言地在我們身邊.
在我們生命裡面.
去傾倒他們的生命.
主我們獻上感恩.
獻上讚美.
求你讓我們這個群體.
在基督裡面形成一個教會.
我們是一個高舉.
每一個人.
在基督裡面.
我們去榮耀神利.
我們去體重每一個人.
有我們的事奉.
在基督裡面得目稱讚.
我們這樣祈禱.
奉主名禱.
\newpage



\section{士師記 17:6}
\label{sec:032452}
\textbf{聖經人聖經事:沒有王,各人任意而行 (1)}~~~~黃克勤牧師博士
\newline
\newline
連結: \href{http://www.church.com.hk/acms/content.asp?site=cdc\&op=show\&type=product\&code=032452\&layout=sermon}{\texttt{032452}} ~~~~ 語音日期: 2026-07-12
\newline
\newline
\hyperref[sec:032451]{\small{< < < PREV SERMON < < <}}
~
\hyperref[sec:index_chronic]{\small{[返順時目]}}
~
\hyperref[sec:032453]{\small{> > > NEXT SERMON > > >}}
\newline
\newline
士師記 17:6
\newline
\begin{longtable}{cl}
\hline
\hline
章節 & 經文 (和合本修訂版)\\
\hline
17:6 & \begin{tabularx}{0.7\textwidth}{X} 那時，以色列中沒有王，各人照自己眼中看為對的去做。 \end{tabularx} \\ \\
[1ex]
\hline
\hline
\end{longtable}
$^{1}$各位姐妹大家好.
祝大家平安.
生命樹宣教網絡.
再和大家一起讀聖經.
看聖經人.
聖經事.
聖經裡面的神.
今天的題目是.
「那時以色列中沒有王,各人任意而行」.
我不夠位置.
所以是「以色列中沒有王,各人任意而行」的第一集.
有好幾集.
因為這裡是《士師記》來到最後一段.
第十七章至第二十一章.
總共有五章聖經.
都可以說是一個很長篇的故事.
但又可以分作每一段.
因此我就分開去說.
但是.
「那時以色列中沒有王,各人任意而行」這句話.
是出現了很多次.
來表達「士師記」就是「那時以色列中沒有王,各人任意而行」.
特別是到尾段才這樣說.
用五章來表達混亂.
你會看到.
《士師記》開頭都挺好的.
後來就開始很混亂.
怎麼混亂呢?.
由今天就看到了.
「各人任意而行」到今天這樣.
第一集.
我們來看看.
上面這兩行字紅色黑色呢?.
紅色藍色呢?.
上面這兩行字紅色藍色呢?.
是我的題目.
「那時以色列中沒有王,各人任意而行」.
「任意而行」的第一項是什麼呢?.
竟然是「各設神堂」.
其實這裡只說了一個神堂.

$^{41}$不過.
這樣都自己開一個神堂了.
我猜其他地方都會的.
所以我就叫它「各設神堂」.
「各設神像」什麼神像呢?.
耶和華的像.
竟然是「自製」.
沒人理會.
「自立祭司」.
你知道祭司.
在會幕裡面敬拜祭司.
只有那一群.
太遠不去我這邊.
我自己開一個神堂.
我自己一個神像.
我自己一個祭司.
這樣都可以.
那是不是「任意而行」呢?.
我們來看看是怎樣的.
十七章第一節.
這次先看十七章.
「以法連山地有一個人名叫米迦」.
紅色字黑色字是經文和合本裡面用的字.
綠色字是我的注釋.
「以法連山地有一個人名叫米迦」.
「以法連」大家都聽過我說很多次了.
「以法連」.
「Ephraim」是「兩堆灰塵」.
「雙倍果子和祝福」的意思.
「兩堆灰塵」.
「雙倍果子和祝福」.
那是什麼來的?.
其實「以法連」就是說有雙重的祝福.
「以法連山地有一個人名叫米迦」.
「米迦」中文只翻譯了兩個音.
「米迦」其實是「米迦耶乎」.
「米迦耶乎」是四個音.
「米迦耶乎」.
意思就是「有誰將臣」.
「有誰將臣」.

$^{81}$「耶乎」像不像「耶和華」.
「米迦」就是「有誰」.
「誰」的意思.
或者「誰」是誰.
「米迦耶乎」.
如果是「米迦耶乎」變成「耶乎」.
就是「神」.
「耶乎」就是「耶和華」的意思.
「耶乎」也可以是「神」的意思.
「米迦」就是「誰」.
「誰將臣」這個字也翻譯為「米迦」.
但其實聖經出現過.
就是「米該亞」.
其實也是「米迦耶乎」.
大家知道了,這個「W」是「O」.
「Who」.
「米迦耶乎」.
「米迦」.
這個「米迦」讀的經文.
你們留意一下也很離譜.
任意而行.
這個「米迦」跟他媽媽說.
大家留意一下.
「媽媽,你拿一千一百赦赫勒銀紙被人拿去」.
「一千一百赦赫勒」是很多的.
「一赦赫勒」是11.3克左右.
這裡是說20多公斤的銀.
這個也是很有錢人.
「你那一千一百赦赫勒銀紙被人拿去不見了」.
「你因此咒詛」.
你生氣了,誰拿了我的銀紙.
願意他死了.
「你因此咒詛,並且告訴了我」.
「將這件不見了銀紙的事告訴我」.
「吵冤巴閉,周圍咒詛」.
然後她說「媽媽,看看」.
真的想不到.
「借銀紙在我借雷」.
「你的錢在這裡」.
「其實你不見了那二十多公斤銀紙在那裡」.

$^{121}$「是我拿去了」.
米加這個兒子跟媽媽說.
「媽媽,你前一陣子不見了一千一百赦赫勒銀紙」.
「二十多公斤,很多錢在那裡」.
「當時你吵冤巴閉,周圍咒詛」.
「告訴我,其實媽媽,我想告訴你」.
「錢是我拿了,我拿了」.
「沒有說我偷了,是我拿去了」.
「銀紙在我那裡」.
他母親就說了什麼,當時的人離譜得很厲害.
他母親說.
「我兒啊,怨耶和華,似福與禮」.
兒子偷了母親的錢,而且這麼大一筆真金白銀.
「兒子,怨耶和華,祝福你」.
「不是吧,巴升到去」.
「你為什麼拿我的錢,不問我」.
「不問自取,是為賊也」.
「我要教兒子」.
「怨耶和華,似福與禮」.
米加就把這一千一百赦赫勒銀紙還給他母親.
他母親竟然這樣說.
她母親說.
「我會分出一隻銀紙來,是為你獻給耶和華」.
「兒子,你這樣,我真的很感恩」.
「我會把這一千一百赦赫勒的錢拿出來」.
「獻給神,是為你而做的,在你名下的」.
然後做什麼呢,留意一下,他媽媽也算是.
很雕刻一個像,然後鑄成一個像.
做什麼呢?.
「你把這一千一百赦赫勒的銀紙還給我」.
「我會拿一些錢出來,鑄成一個像」.
然後做一個模,再鑄成一個像.
鑄成一個像是什麼呢?.
例如先用木,先用石鑄成一個像.
用泥吸住它,兩邊張開,用金屬鑄成一個模.
鑄成一個像,多數是用銀做的,很貴重.
「現在我還是交給你」.
即是那一千一百赦赫勒,你這麼好給媽媽.
「願耶和華祝福你,我還會拿一部分出來鑄成一個像」.
「是為你而做,是獻給神的,獻給耶和華」.

$^{161}$獻給耶和華,什麼來的?.
如果你是以色列人,一早就知道.
可不可以雕刻偶像來敬拜,不可雕刻,不可跪拜.
「現在我會為你做一個像獻給耶和華」.
那媽媽就麻煩了,不過她媽媽說.
「我會為你這樣做的,你真乖,願耶和華祝福你」.
「而且這一千一百赦赫勒銀紙,我還交給你,你給我管」.
這個媽媽真是慈母多敗兒.
米家將銀紙還他母親,他母親說「我給你管」.
他母親將二百赦赫勒銀紙交給銀將,雕刻一個像,鑄成一個像.
安置在米家的屋內.
即是米家有自己的家,而且看來都不小.
即是說米家是一個很有錢的人.
他媽媽銀紙也有一千一百赦赫勒,被人偷.
「給兒子你管,還拿二百赦赫勒出來做銀像」.
而且米家有自己的獨立屋.
在當日來說真的不得了,有錢.
第五節.
於是米家有了一個神堂.
自己有一個神堂.
然後他還做了一個以弗德.
以弗德即是大祭司穿的那件外袍的一部分.
以弗德,然後胸牌掛在以弗德上面.
做了一件以弗德的衣服在那裡,表示甚麼呢?.
即是說這裡象徵有祭司,有大祭司在那裡.
做了一件以弗德和家中的神像,還有其他神像.
家中的神像主要就是剛才做出來的二百赦赫勒銀那個.
另外,這個米家真厲害.
自立祭司分派他一個兒子做祭司,這樣也行.
即是說有錢,那筆錢本來以為不見了.
原來那個壞蛋拿了,現在拿出來,這個兒子真乖.
我會為你做一個像,獻給銀河華家,放在你家裡.
不如你乖乖地開一個神堂.
又做一件以弗德象徵,大祭司在那裡.
然後為這個神堂派一個兒子做祭司.
那時以色列終末有王,各人任意而行.
這樣,自己開一個神堂,又弄一個神像.
又自己做一件以弗德,然後叫兒子做祭司.
離譜至極,所以撒母耳記載到這裡是司祭的時候.
那時以色列終末有王,各人任意而行,就是這個題目.

$^{201}$嘆息不堪.
第七節,話鋒一轉,猶大伯利恆有一個少年人.
他在猶大伯利恆,即是說他本身是猶大族.
猶大族而且還是利未人.
他在那裡居住,在猶大伯利恆.
這個人離開猶大伯利恆,不想在那裡居住.
什麼原因?沒有說.
要找一個可居住的地方,即是說他覺得那裡不好居住.
不適合他等等.
他走路的時候,到了以法連山地,出來找地方.
誰知走到以法連山地,去到米家的家.
到處找地方居住,然後到以法連山地.
見到一家大戶人家,就去到米家的家.
很多人出入,門不關.
說得難聽一點,中國人說「以香火頂盛」.
看看是什麼來的,去到那裡.
碰到主人米家,米家就問他.
「你年青人,你在哪裡來的?見你都是外地人,在哪裡來的?」.
他就說「我是猶大的伯利恆那裡來的」.
「我本身是屬於利美族」.
這個都是他們一般人當時見面的稱呼.
的聊天.
「這位兄台,你在外地來的話,看知你外地來」.
風塵撲撲,他自我介紹.
「我是來自猶大的伯利恆的」.
「我本身的支派是屬於利美族」.
「來這裡是想找一個地方居住而已」.
這樣介紹.
例如你見到新界都是這樣介紹.
「我是變雅敏支派的人」.
介紹我們是什麼支派.
當時是這樣.
米加就說.
「這樣就好了」.
「你不如住在我這裡」.
「我會以你為父,為祭師」.
為父即是我尊敬的人.
而且做我的祭師.
「我會以你為父,為祭師」.
「我每年會給你十四克銀子做人工」.

$^{241}$「一套衣服和度日的食物」.
「在我這裡包吃包住」.
「包一年一套衣服」.
「和十四克銀子」.
在當時來說.
十四克銀子都是一個好的數字.
於是利未人就說.
「這樣就好了」.
「我就試試住在你這裡」.
利未人就情願與那人同住.
這是伯利恆來的利未人.
沒有名字.
就說「好啊」.
「住在這裡很好」.
「而且那個少年人」.
「那人」即是說米家看待這個少年人.
好像自己的兒子一樣.
很疼愛他.
因為他是他的祭師.
他真的說口齒.
他就會尊敬他.
做為父為祭師.
幫他打理神堂.
利未人是屬於大領敬拜的.
話說回來.
做祭師的.
如果從舊約來看.
只有利未人的亞倫族人才能做祭師.
不是每個利未人都做祭師.
很多利未人會管理會幕.
聖殿裡面的工作.
有些負責做湯.
牛羊.
有些負責唱歌.
有些負責將牛羊的灰倒掉.
很多事情.
不是利未人都做祭師.
是亞倫後裔才做祭師.
現在隨便一個利未人來.
找你做祭師.

$^{281}$而且你多疼愛他.
其實他都不能做祭師.
就是這樣發生.
所以.
你說他們是不是不敬威爺和華呢.
又不是.
因為整個神像都是威爺和華.
但是他又完全不跟舊約聖經的辯論.
當時摩西傳下來的律法是這樣做的.
治國立神堂.
治祭神像.
自己做祭師.
立了祭師.
治成一國.
其實當時的會幕在哪裡.
在士劍.
大家記得嗎.
在士劍有會幕.
是有大祭師在那裡.
有祭師在那裡.
是去那裡獻祭的.
不是治成一國.
遠了我不去了.
我自己開個神堂.
不行的.
他們是這樣做的.
所以米家分派少年的利未人做祭師.
住在米家的家裡.
現在我知道威爺和華必定賜福給我.
因為我還有一個利未人做祭師.
米家就很傻.
他以為現在我一定有更多的祝福.
為什麼.
我不單止有神堂.
我還有神像.
我還有自製的伊忽德.
然後我本來沒有兒子.
找我兒子做祭師.
現在真的有個利未人做祭師.
因為這樣.

$^{321}$我相信威爺和華必定更賜福給我.
因為我有一個利未人做祭師.
那些敬拜的觀念全部錯了.
離譜亂來.
沒有人管的嗎.
沒有人管.
那時以色列中沒有王國人任意而行.
這是第一講的關於祭師記.
沒有王國人任意而行.
在這故事裡看到一個例子.
這件事裡看到一個例子.
就是沒有王國人任意而行.
到自己在家裡開個神誕.
這樣都可以.
而且神堂的神像是雕刻鑄造出來.
是威爺和華向祂跪拜.
剛好十屆第一屆第二屆都不行.
不過他們沒有理會.
這也很光榮.
現在神一定會祝福我.
因為我還有一個利未人做祭師.
我一定會蒙福.
祂是否真誠呢 祂是真誠的.
不過道理不對.
大家記不記得.
主耶穌和撒瑪利亞婦人談道.
撒瑪利亞婦人說.
猶太人敬拜在耶路撒冷.
我們被迫在基里森山.
哪裡才是.
主耶穌的答案是.
其實在哪裡都不是最緊要.
最緊要因為耶和華神是靈.
所以拜祂的要用心靈和誠實.
記得嗎.
心靈就是去到敬拜神.
一定要用我的心靈去到敬拜主.
誠實的意思是我真真誠誠的.
真誠有誠則靈.
這個是中國人的想法.

$^{361}$有誠不一定靈的.
誠心啊 誠心就大了.
不是的 有誠不一定靈.
不一定通.
因為神是靈.
所以拜祂的要用心靈.
和誠實 誠實這個字.
大家要留意 要留意 要留意.
心靈和誠實.
這個誠實的意思是甚麼.
可以直接翻譯做真理.
是按著我對神真誠地敬拜的心.
和按著敬拜的真理而行.
不能夠隨便做敬拜.
這個是有誠.
不過它沒有真理.
真理是錯了.
不管了 我喜歡 我覺得好.
就任意而行.
你看到他媽媽也很有趣.
媽媽的心態也敗壞了.
不見得像兒子說.
願意祝福你.
她不罵的.
然後我還要為你拿一部分出來.
做一個像是獻給神.
很明顯她覺得這個兒子很乖.
這個兒子很乖.
亂了龍 只能夠說.
任意而行的第一講就是亂了龍.
不要這樣 以心靈敬拜主.
又要跟隨主是按著.
神的心意與原則與道理的.
真理 祈禱.
求主幫助我們.
讓我們讀這些經文的時候.
覺得為什麼這麼離譜.
我們不在這裡去責怪去挑戰.
只不過讀完之後.
心裡面去警醒提醒自己.

$^{401}$真的要很小心.
因為神是興邁不得.
讓我們很小心.
免得我們任意而行.
以為我做得落力.
神就一定欣賞.
以為我這樣做.
是帶著誠意.
神一定歡喜.
其實不是的.
求主真理而行.
求主教導我們.
求真理的聖靈教導我們.
管理我們 幫助我們.
甚至限制我們.
順聽祈禱.
奉耶穌基督的名字.
阿們.
\newpage



\section{羅馬書 10:4-21}
\label{sec:032453}
\textbf{宣認耶穌是主}~~~~周曉暉牧師
\newline
\newline
連結: \href{http://www.church.com.hk/acms/content.asp?site=cdc\&op=show\&type=product\&code=032453\&layout=sermon}{\texttt{032453}} ~~~~ 語音日期: 2025-12-14
\newline
\newline
\hyperref[sec:032452]{\small{< < < PREV SERMON < < <}}
~
\hyperref[sec:index_chronic]{\small{[返順時目]}}
~
\hyperref[sec:032454]{\small{> > > NEXT SERMON > > >}}
\newline
\newline
羅馬書 10:4-21
\newline
\begin{longtable}{cl}
\hline
\hline
章節 & 經文 (和合本修訂版)\\
\hline
10:4 & \begin{tabularx}{0.7\textwidth}{X} 律法的總結就是基督，使所有信他的人都得著義。 \end{tabularx} \\ \\ \relax
10:5 & \begin{tabularx}{0.7\textwidth}{X} 論到出於律法的義，摩西寫著：「行這些事的人，就必因此得生。」 \end{tabularx} \\ \\ \relax
10:6 & \begin{tabularx}{0.7\textwidth}{X} 但出於信的義卻如此說：「你不要心裡說：誰要升到天上去呢？（就是說，把基督領下來。） \end{tabularx} \\ \\ \relax
10:7 & \begin{tabularx}{0.7\textwidth}{X} 或說：誰要下到陰間去呢？（就是說，把基督從死人中領上來。）」 \end{tabularx} \\ \\ \relax
10:8 & \begin{tabularx}{0.7\textwidth}{X} 他到底怎麼說呢？「這話語就離你近，就在你口中，在你心裡，」（就是說，我們傳揚所信的話語。） \end{tabularx} \\ \\ \relax
10:9 & \begin{tabularx}{0.7\textwidth}{X} 你若口裡宣認耶穌為主，心裡信神叫他從死人中復活，就必得救。 \end{tabularx} \\ \\ \relax
10:10 & \begin{tabularx}{0.7\textwidth}{X} 因為，人心裡信就可以稱義，口裡宣認就可以得救。 \end{tabularx} \\ \\ \relax
10:11 & \begin{tabularx}{0.7\textwidth}{X} 經上說：「凡信靠他的人必不蒙羞。」 \end{tabularx} \\ \\ \relax
10:12 & \begin{tabularx}{0.7\textwidth}{X} 猶太人和希臘人並沒有分別，因為人人都有同一位主，他也厚待求告他的每一個人。 \end{tabularx} \\ \\ \relax
10:13 & \begin{tabularx}{0.7\textwidth}{X} 因為「凡求告主名的就必得救」。 \end{tabularx} \\ \\ \relax
10:14 & \begin{tabularx}{0.7\textwidth}{X} 然而，人未曾信他，怎能求告他呢？未曾聽見他，怎能信他呢？沒有傳道的，怎能聽見呢？ \end{tabularx} \\ \\ \relax
10:15 & \begin{tabularx}{0.7\textwidth}{X} 若沒有奉差遣，怎能傳道呢？如經上所記：「報福音、傳喜信的人，他們的腳蹤何等佳美！」 \end{tabularx} \\ \\ \relax
10:16 & \begin{tabularx}{0.7\textwidth}{X} 但不是每一個人都聽從福音，因為以賽亞說：「主啊，我們所傳的有誰信呢？」 \end{tabularx} \\ \\ \relax
10:17 & \begin{tabularx}{0.7\textwidth}{X} 可見，信道是從聽道來的，聽道是從基督的話來的。 \end{tabularx} \\ \\ \relax
10:18 & \begin{tabularx}{0.7\textwidth}{X} 但我要問，人沒有聽見嗎？當然聽見了。「他們的聲音傳遍全地；他們的言語傳到地極。」 \end{tabularx} \\ \\ \relax
10:19 & \begin{tabularx}{0.7\textwidth}{X} 我再問，以色列人不知道嗎？先有摩西說：「我要以不成國的激起你們嫉妒；我要以愚頑的國惹起你們發怒。」 \end{tabularx} \\ \\ \relax
10:20 & \begin{tabularx}{0.7\textwidth}{X} 又有以賽亞放膽說：「沒有尋找我的，我要讓他們尋見；沒有求問我的，我要向他們顯現。」 \end{tabularx} \\ \\ \relax
10:21 & \begin{tabularx}{0.7\textwidth}{X} 關於以色列人，他說：「我整天向那悖逆頂嘴的百姓招手。」 \end{tabularx} \\ \\
[1ex]
\hline
\hline
\end{longtable}
$^{1}$錄音/錄影的版權是屬於基督教宣導會不各堂所有.
我們的屬靈傳統之一,是以真理教導為本.
歡迎您將此錄音/錄影的連結傳給親友.
懇請注意,未經授權,請勿轉錄或剪輯本錄音/錄影.
甚願好像《使徒行傳》20章32節下所說.
「這道能建立你們,叫你們和一切誠誠的人同得基業」.
主耶穌基督在紀念您降臨的日子.
也預備迎接您再來的日子.
我們心裡真的要去審查自己.
將我們裡面有忍耳未言的罪,願你都赦免.
我們所身處的社群,這個世界充滿著罪惡.
主啊,我們都為了這一切,求你赦免.
引領我們今天要讀你的話語,求聖靈指教我們每一個.
讓我們聽進去,聽得明白.
並且生命因著你的話語有改變.
多謝你,我們的祈禱是奉耶穌基督,得勝靈指祈求,阿們.
這幾節熟悉的經文經常出現,就是我們今天所讀的第十章.
但是除了這幾節經文之外,這一大段說什麼呢?.
這一大段由九章開始,今天是第十章.
之後我們會再看第十一章.
這三章經文在過去一段比較長的日子裡.
學者認為它似乎是一個附錄.
明明第一章到第八章講恩訊清義,講福音.
講救主怎樣幫助我們得到這個身份.
講罪的問題,講怎樣相信.
講我們自己無助的時候怎樣禱告.
聖靈也說了,去到八章最後的結論是.
沒有任何東西可以令我們與神的愛隔絕.
完美,在這裡收筆.
然後去到十二章開始的應用篇.
怎樣有感恩的心回應神的愛.
將自己的身體獻上當作活祭.
去進入社群和人相處,這樣不就好了嗎?.
所以學者認為,為什麼在九章到十一章.
又講猶太人,以色列人得救的問題.
所以那些學者就說.
是不是加插進去的附錄呢?.
但近年卻有研究羅馬書的指出.
我們應該尊重作者寫的次序.
並且他當時寫,不是有些空位寫上去.

$^{41}$要寫一卷書,那些羊皮是很貴的.
每一個空位都好好應用.
就好像我們比較年長的知道.
以前打電報,每一顆字都是錢.
所以不是應該說的就不說了.
是省下了的.
所以保羅在這裡特別要和羅馬教會說.
用以色列救恩歷史的事情.
其實是在針對.
現在在他們羅馬教會裡面.
有事情你們要處理不對.
因為在紹行傳十八章我們記得.
先前羅馬皇帝格魯德就是曾經下命令.
在他的首都羅馬城所有猶太人走.
離開我們這個城市.
然後在教會裡面的猶太基督徒.
到壽千年都要離開.
於是突然間一下子.
本來有兩班人的教會.
有猶太人,有外邦人.
即不是猶太人的人.
特別是希臘人,羅馬人的信徒.
最早成的教會.
一下子沒有了一大半的猶太基督徒.
慢慢地他們又繼續在那裡成長.
聚會,傳福音.
有新的信仰人加入教會.
都是外邦人.
結果這個教會就越來越少猶太的元素.
文化,崇拜,很多禮儀等等.
越來越淡化.
那麼這個教會會變成怎樣呢?.
也很好.
但是到這個皇帝死了之後.
第二個凱撒就讓他們回來了.
讓猶太人回到羅馬.
有些以前在這個教會成長的猶太基督徒.
就回到羅馬城.
回到這個他成長的教會.
來到教會.

$^{81}$他就發現.
為什麼現在會這樣呢?.
很多事情都不同了.
臨近聖誕.
外地有很多頂子妹妹回來香港.
有些是我們自己的老會友.
早幾天有一個弟兄告訴我.
他六十年前的少年團友回來找他.
那個團友當時還在馬保道的年代.
一下來到康宅這個地方.
地方已經不同了.
人也很不同.
以前認識的人大部分都升天了.
那時候敬拜的方式講道.
是用國語,普通話的.
藤木詩.
今天主要是講廣東話.
當然我們也有英文堂友.
有普通話崇拜.
但是主體他所經驗的.
和當年很不同.
就連唱詩.
以前我們唱生命聖詩.
是用國語唱的.
我們的三一頌.
都是國語普通話.
現在全部是廣東話了.
那些文化習慣都改變了.
有一種親切又陌生的感覺.
對於這些猶太基督徒.
他們回到自己的教會.
他們感受到那種文化很不同.
所以會產生這種矛盾.
甚至開始爭拗.
很想將以前的一些做法帶回來.
新的信徒.
新的愛邦人信徒.
根本沒有經驗過.
他們甚至不認識這些猶太基督徒.
就好像我和你說過.

$^{121}$六十年前的會友.
你也不認識他.
大家的意見可能有衝突,有矛盾.
保羅在這裡要特別說.
以色列人救恩的歷史.
是想去處理一群信徒的心結.
到底我們這群猶太人和愛邦人.
我們怎樣可以在一起?.
愛邦人心裡可能想.
神不是已經叫保羅將福音傳給愛邦.
那個心意好像要愛邦人更多去信主.
猶太人還這麼抗拒.
好像不要理會他們.
有這種矛盾.
而保羅整本羅馬書.
還有一個底層的思想.
他寫這本書.
是很想這群人能夠明白福音.
並且同心.
以致力量一致地繼續宣教.
在當時.
保羅在羅馬書最後一章說.
他很想去西班牙宣教.
西班牙根據歷史.
是當時羅馬帝國新近佔領的地區.
現在羅馬鋪了路去那裡.
可以去到.
所以保羅有機會去宣教.
就是這樣.
羅馬這個首都是一個戰略的位置.
可以將福音帶到更遠.
如果這群人整天吵吵閉閉.
不和睦.
根本不能夠做到宣教的工作.
所以保羅在應用之前.
繼續要說.
我們要明白上帝怎樣想.
這群原本神的選民.
以色列人的得救問題.
稍為停一停.

$^{161}$這個背景都令我們想到.
我們今天反覆說福音的原因.
這一季我們已經一路說福音的核心.
整本羅馬書一路說這件事.
每個星期都說.
其實不說羅馬書.
每個星期都提.
甚至我想我們前面.
剛才打手鈴的小朋友.
經常都聽主一學.
耶穌基督為我們受死.
為我們復活.
經常都說.
祈禱有期.
重重複複說.
是為了什麼?.
大概不是怕我們不記得.
而是怕我們不常常將這件事擺為先.
保羅想福音是這群人擺為先.
對於我們今天來說.
我們常常聽的福音.
是要你再重溫.
重溫之後.
不是就停下來.
你要去跟別人說.
要跟別人分享.
不難的.
你可以做到的.
你覺得很困難.
起碼你都邀請人來聽福音.
這個聖誕節.
我相信在我們香港這個時刻.
上帝一定有祂的作為.
讓人得著福音的好處.
將那種患難轉化成為福氣.
弟兄姊妹.
不要放輕福音.
並且持續地.
去完成教會和睦之後要做的.
我們要宣教.

$^{201}$要傳揚福音.
保羅接續去處理九到十一章的問題.
以色列人為什麼還是不相信?.
如果你細心留意.
一到八章一直形容猶太人.
希臘人.
包括今天我們讀的.
都有形容這兩個類別.
但他特別提以色列得救的問題.
為什麼?.
他想這群人再回到立約的名字.
神同雅各立約的時候.
改名你從今以後叫以色列.
所以保羅特意用這個名字的轉換.
讓他們想起.
到底神起初的立約這個盟約是為了什麼?.
而這個盟約裡面.
為什麼這個原先立約的選民.
他們今天不相信.
失落了得救不了?.
上一章你聽過.
是說神的揀選主權.
並且對著硬心的人怎樣去處理.
而且不是肉身亞伯拉罕的後裔.
才是神的選民.
我們這群人都是外邦人.
我們卻成為神的兒女.
也同樣是亞伯拉罕的後裔.
也同樣成為真以色列人.
這是神的主權裡的揀選擺放.
上一章第九章你聽過.
第九章最後問一個問題.
如果神都揀選了.
不關我們事.
他不相信.
那些人不相信.
神令他們這樣.
不是.
第九章已經開始說.
他們的不信.

$^{241}$他們是有份的.
因為他們選擇.
他們進行一大堆律法.
都不願意相信耶穌.
做一大堆事都不願意承認耶穌.
這是第九章最後說到的原因.
於是接駁我們今天讀的第十章.
其實第十章我們應該從第一節開始.
看到第四節.
已經概略地告訴你.
它後面要說什麼.
一至四節說了三件事.
第一,這群人不按照知識的熱心.
令他們得不到救.
他們不明白.
不跟隨神的義.
即耶穌基督.
令我們稱義這個教導.
他們要行律法.
不要信基督.
即是怎樣呢?.
在十章的三節.
最後的三行.
我請大家幫我讀.
因為神的義在此.
告訴他們他們不明白.
說了之後又不跟隨.
只是想用自己的方法去做.
為什麼他們一直覺得.
我要做完這些行為.
我就得到神的喜愛.
我就得救.
為什麼他們是這樣呢?.
因為摩西律法頒佈了這麼久.
他們就覺得是這樣.
在九章尾.
其實給了我們一個關鍵的詞語.
我打了藍色給你們看.
第四行.
他們倒以為.

$^{281}$以為這個律法.
就是令我們得到神稱義身份的方式.
一個途徑.
一個工具.
一個方法.
所以他們不接受耶穌.
以為.
其實以為是很嚴重的.
你們有沒有試過在家裡.
跟太太或丈夫.
以為一些事情搞錯了.
就吵架.
在工作上.
老闆派了一些事情給你做.
你以為是這樣.
但搞錯了.
你的工作也沒有了.
在真理上.
以為是生和死的問題.
不要隨便以為.
所以這群人以為.
上帝跟他們立約.
給了他們這些條文.
他們以為這就是能夠給神稱讚.
給神喜愛的方法.
以為.
很嚴重.
所以他們失落了.
他們追求的不是神的義.
而是自己納自己的義.
我要堅持做這件事.
我是對的.
而且他們很熱心地去做.
所以在第二節.
十章二節裡.
保羅說他們很有熱心.
我可以為他造證.
其實他自己也是一份子.
保羅未信主前叫什麼名?.
素羅.

$^{321}$素羅是做什麼的?.
四處抓基督徒.
對的.
他覺得他是正義的代表.
代表上帝去抓那些基督徒.
弄死他們.
對的.
很熱心地去做.
耶穌光照他們.
讓他們看見.
你錯了.
你以為錯了.
他們有熱心.
但以為錯了.
就很大件事.
用自己的方法.
證明自己是神的子民.
還很驕傲地去做這一切.
是不可以的.
所以保羅在第四節.
將人的焦點聚在耶穌基督身上.
第四節請大家讀.
基督.
終結在原本的文字代表.
目的地.
終極的目標.
耶穌基督就是.
律法的終極目標.
不是不要律法.
不是以前律法廢了.
而是表達耶穌去到那個位置.
祂會成全這一切律法所表達的.
最重要的意義.
而律法這一堆.
他們做不到的事.
不可能守每一條.
這些是要逼人去想.
我搞不定.
原來我是要依靠耶穌.
我要相信耶穌.

$^{361}$我要依靠神的恩典.
相信神的恩典.
恩信,稱義.
以色列人.
依然沒有選擇信心的道路.
他們選擇行為的道路.
是不是他們無知呢?.
不是的.
第二節你看清楚.
說他們有熱心.
第三項說.
他們並非按照知識.
不是說他們無知.
不是說他們不知道.
而是他們知道了.
要照著做.
知道要這樣走才行.
不跟就此.
有一個人.
很想從柴灣.
去到長洲.
打Google Map告訴他.
只要你這輛車去到中環.
轉小輪就可以去到.
於是他滿心一腔熱血.
很想去長洲.
於是下樓開自己的車.
一路駕駛經阿公岩道.
然後上去東區走廊.
東區走廊來到我們附近的中環繞道.
然後去到中環繞道.
一出差不多交易廣場.
做什麼?.
轉入碼頭.
轉到碼頭他不下車.
不停車,不泊車.
為什麼我不可以繼續開這輛車去長洲?.
於是他就繼續走.
去到西營盤.
沒有地方出沒可能去到長洲.

$^{401}$去到博富林都沒可能.
去到南區越來越遠.
他不肯下車.
轉唯一方式.
乘小輪去長洲.
他無論如何都去不到.
他開車去.
有律法錯嗎?.
不是.
但去到一個位置他遇到基督.
他要接受基督.
在那個時刻.
耶穌已經來了.
這群以色列人依然不信.
還將他釘死.
他們選擇.
繼續走自己的意義.
無論如何都要自駕遊.
為什麼不可以開車去長洲?.
沒可能的.
保羅於是.
繼續和他們討論第五到第八節.
進一步對比.
律法和信心是什麼.
他知道他們一定會爭辯.
摩西說的.
進行這些事.
進行啊保羅.
走出來就因此而活.
走完這些律法.
保羅問.
摩西這樣寫.
難道你認為.
你們這樣走.
可以完成這個救恩嗎?.
信心的意義.
是告訴他們.
這一切的事.
打通這個人去到神這條路.
你的行為.

$^{441}$行義是做不到的.
難道你想天上.
使耶穌降生嗎?.
難道你進入.
陰間使耶穌在死人中.
復活嗎?.
通通都不關你人做的事.
不關的.
由頭到尾都是神選擇.
降生成為人.
在這裡上十字架死亡.
並且神讓.
耶穌基督復活.
這一切.
不好意思和你自駕遊無關.
全部都是信心.
你只能夠相信接受.
你沒有什麼可以做.
如果這樣的話.
這些猶太人.
最初決定要做神的子民.
跟隨神.
起初是沒有問題的.
就像我們有時候.
平時做決定,做選擇.
我們按當時最好的.
去選擇最對的.
過了一段很長的日子.
新的資訊,新的發展.
不同了.
我們原本的決定已經不合用.
你要改的.
你要再按當時真實.
更好的去轉變.
不是說.
以前的決定錯了.
不過到了新的階段.
你要選擇的事情.
你就要做對.
以色列人在這個位置.

$^{481}$他們不是以前決定錯.
而是這一刻有耶穌基督.
他們依然決定錯.
一定要走自己.
要走的路.
所以兩個選擇.
放在他們面前.
怎麼選擇呢?.
走律法?.
信基督?.
走律法我習慣了.
都是這樣.
信基督.
信基督很陌生.
不接受.
保羅繼續問.
這個信基督的方式很簡單.
我請大家讀第八節.
「這就是我們所說的關於信心的話」.
關於信心的話.
這個話是什麼話?.
就是你經常背的金句.
「口裡承認耶穌是主,心裡相信神使祂從死人中復活,就必得救」.
就是在說我們福音最核心.
耶穌基督所做的一切.
而祂現在對你來說是哪一位?.
這個核心你也懂得說.
耶穌基督是我們的救主.
是我們生命的主宰.
你這樣宣認出來.
說出來.
並且你心裡要相信.
祂的死亡與復活.
是真的.
但有些人就問.
「心裡相信」是什麼意思?.
可能對一個初信人來說.
這個概念有點模糊.
相信是不是我接受耶穌復活的那一位?.
是合理的.

$^{521}$他知道這是事實.
相信這是事實.
但是在雅各書也說過.
魔鬼也相信.
魔鬼也相信耶穌是復活了.
他知道這是事實.
但不會救到他.
如果這樣對我們人來說.
我們知道耶穌復活了.
也救不了你.
這個不是心裡相信.
因為相信有這個事實.
不等於真正的信靠.
我相信小輪可以去到長洲.
但我一直都開車.
我一直都不上小輪.
和我沒有關係.
我不是說小輪去不到長洲.
但我一直都要試開車.
相信並且你要接受這個事實.
你要應用在你的人生中.
在你的人生中.
這件事影響你往後的人生.
這個叫做心裡相信.
神使耶穌從死人中復活.
祂的復活影響你抉擇.
因為我們都會像耶穌一樣有復活的生命.
今天我們到了死所選擇的都不同了.
我們不再怕死亡.
我們不會因為被威嚇而退縮改變.
這個心裡相信實現出來了.
在十八世紀一百五百年.
有一個宣教士叫John Gibson Patton.
當時他被派去南太平洋一個小島去宣教.
那個小島今天的名字叫做阿魯雅圖.
之前叫你移民去那裡.
很漂亮的一個太平洋小島.
不過對於當時一百五百年的人來說.
是很驚嚇的地方.
因為那個島上住的島民是食人族.

$^{561}$宣教士去.
帶上太太和初生的孩子.
很快他的太太和初生孩子沒有被吃了.
不過染了疫病死了.
Patton繼續留在那裡.
三十多年去宣教傳福音.
並且翻譯聖經.
他在翻譯聖經的過程中.
他發現了一個難題.
就是如何用當地的語言去翻譯「信心」這個字.
你知道「信心」有時候對我們來說是很哲學性的.
我們問小朋友什麼是信心.
可能他也未必能夠回答你.
對於那些島民.
如何用他們的語言去講述「信心」或翻譯這件事.
就成為了Patton一路去拗頭的地方.
有一天他在他的宣教士的茅屋工作.
這真是一個阿魯亞圖的茅屋.
剛才第二課有一個姐妹告訴我.
她有張這樣的照片.
她的女兒去世了.
我馬上加回去PowerPoint.
宣教士的茅屋裡有一個島民.
就來找她聊天.
宣教士那時候還在思考翻譯的事.
當時她坐在一張椅子上.
好像大家一樣兩隻腳踩著地上.
坐在一張椅子上.
他突然間想.
先問一問她.
他就問這個島民.
我現在這樣的姿勢.
是應該如何表達在你的語言.
這個島民就說.
Kokai ana missi.
宣教士先生.
你坐在這裡.
接著他就換一個姿勢.
整個人靠在椅子的椅背後.
靠後很癱瘓的樣子.

$^{601}$兩隻腳吊起.
不再踩著地上了.
他又問.
那我這個姿勢呢?.
她就說 Fakaron rango.
宣教士先生.
你是整個人的重量.
壓在這個地方.
他突然間發出一聲.
This is the term.
我要找的自由.
這個就是相信,信靠,信心.
他往後翻譯.
翻譯《約翰福音》三章十六節.
當地的信徒就會聽到.
神愛世人.
甚至將他獨生子賜給他們.
叫一切將重量放在耶穌基督身上的人.
不致滅亡.
凡得永生.
將你整個人壓在耶穌身上.
投放在耶穌身上.
你一生擺在耶穌身上.
這個叫做相信.
Fakaron rango.
當他晚年的時候.
他回顧.
他見證這個島的島民.
不再是食人族.
建立了教會,學校,孤兒院.
這個島的島民.
相信耶穌之後的改變.
他回顧說.
我為這個島上的人去宣告.
耶穌基督是主.
不是頭腦知道這個事實.
他是影響了這群人的生活,生命改變.
整個人,全個重量.
我們生命的重擔,重心.
全放在耶穌身上.

$^{641}$這個叫做心裡相信.
其實我們中文已經很好了.
完全依靠.
依靠著.
全靠著這位主.
做什麼決定.
怎樣想事.
我們應該怎樣看人生.
怎樣看死亡.
怎樣看災難.
怎樣看我們和人的相處.
怎樣看社會.
怎樣看世界.
全部都是靠在耶穌身上.
從那裡去看.
這個叫做相信.
所以保羅說.
如果這樣相信的人.
必不羞愧.
他一定會去到上帝預備給他榮耀的地方.
不會再被羞辱.
一個對將基督作為救主和主宰的人.
持續信靠法的Ron Ron Go.
全個人的重量,重擔都在耶穌身上.
祂帶著你.
會怎樣呢?.
一位弟兄失業.
原本在他的公司都做得很高職位.
突然被裁員.
面對經濟的壓力.
身份的危機.
前途茫茫.
在這個時候他會求告主明.
他會怎樣祈禱.
當然他一定會祈禱.
主啊,我需要一份工作.
求你讓我能夠找到工作.
同時他也會祈禱.
耶穌基督你是我人生的主宰.
你掌管一切.

$^{681}$無論你讓我遇到怎樣的事.
我願意信靠.
接受你的安排.
求你幫助我.
繼續跟實你.
我一生都在你手上.
求你引領我.
這就是信仰.
不只是求你給我工作,救我.
還是我們人生的方向.
都在耶穌手上.
祂帶著我們走.
所以保羅說明了.
這個相信的事情.
他在十四到二十一節.
回到不信以色列人.
以色列人不信的問題.
他就這樣表達.
「凡求告主明的就必得救」.
是的,但是人未相信怎樣求呢?.
未聽見怎樣信呢?.
沒有人去傳,怎樣聽得見呢?.
沒有人奉差遣,怎樣傳呢?.
他用倒敘法.
五四三二一去問.
我幫你排在最底.
要有人差遣人.
他繼續用宣教思維.
我們全部人都是小宣教士.
從宣教人生的角度去看.
沒有人奉差遣,怎樣去傳揚?.
沒有人傳揚,那些人怎樣可以聽到?.
聽不到,他怎可能信呢?.
他不信的話,他不會去求告的.
這個邏輯去問.
那這班以色列人.
有沒有這五樣東西?.
有沒有人奉差遣?.
有沒有人傳揚?.
有沒有人讓他們聽見?.

$^{721}$有沒有人叫他們相信?.
以致他們能夠求告?.
有的,全部都有.
但是聽到的人的結果是什麼呢?.
不是所有人都聽從福音.
不是他們沒有聽見.
是他們聽了不跟從.
聽了不照著做.
下車了,轉船了.
這樣才能去到長洲.
不聽,不跟.
這樣的狀態.
是不是他們選擇的責任呢?.
當然是.
而且他們繼續去問.
難道他們真的聽不到嗎?.
我請大家讀十八節.
「可視」.
保羅羅斯篇讓他們很熟悉的話語.
這些已經傳了出去.
你應該聽到的.
以色列人是聽到的.
為什麼呢?.
耶穌基督釘十字架這件事.
在各地方都知道.
從耶路撒冷,猶太傳到撒馬利亞.
一直傳到羅馬.
都有人在說耶穌的事情.
你們都應該知道.
猶太人,以色列人都知道.
所以迫害基督徒的人繼續.
聽見.
但是他們不選擇這個信心.
只需要口裡宣認,心裡確實相信.
不選擇這條路.
責任真的在他們身上.
但是即使如此.
有一節最後這一句.
那裡說到.
神對著這群以色列人.

$^{761}$神依然去表達.
「我終日向那背逆,頂嘴的子民伸出雙手」.
不只是在羅馬書寫成的年代.
在往後的二千年.
上帝依然在等待他的以色列子民.
回到他身邊.
去歸順他.
到今天依然是這樣.
兩年多前.
以哈戰爭開始.
其中在軍隊裡有一個軍官.
他是基督徒.
是以色列人.
他大膽地向他的將軍申請.
「將軍,我看見我們的新兵」.
「個個都很疼」.
以色列每一個少男少女.
讀完書就要先當兵.
男和女都要當兵.
強制的.
所以當時打仗很多是新兵.
「這班新兵都退縮了」.
「我可不可以申請」.
「你讓我派聖經給他們?」.
派一些小小的聖經.
讓他可以戴著軍隊的行李背囊.
將軍說「好」.
於是這個軍官就去找人印了.
希伯來文舊約的聖經.
當然是用現在的現代希伯來文.
並且印上翻譯好.
現代希伯來文的新約.
整本聖經派給這些兵.
很重要.
為什麼?.
這些新一代的以色列人.
其實他們不是很猶太教的.
不會去聽拉比講道.
未必去敬拜.
他們的生活好像一般世人一樣.

$^{801}$但是到了戰爭.
隨時會死的.
回到以前猶太教的那些人.
「糟了!我這一代都沒有做過」.
「沒有行過義」.
「沒有按照條文做足」.
「又不守安息日」.
「怎麼辦?死了」.
「死後去哪裡?」.
但是當他們打開這本聖經.
有了新約.
他們看到有耶穌基督.
只要口裡承認.
心裡相信.
「我這一條命」.
「即使我沒有命」.
「我都知道我去哪裡」.
上帝的手依然攤開對他的子民.
同樣是對我們每一個人.
這段經文也告訴我們.
無論是猶太人還是希臘人.
神都是他們的主.
每一個都需要服膺.
頂著這個兩千年前的真理.
對我們今天的信徒依然重要.
因為我們的生活.
你記得.
不要再以為那麼多了.
尤其我們在今天.
對一個發展了的教會來說.
我們會以為很多事情.
比耶穌更重要.
比福音更重要.
我們以為我們所做的活動.
我們所擺出來的東西.
我們所說的.
我們所做的事情.
很重要很重要.
到一個地步.
我們忙於做這些.

$^{841}$忘記了福音的核心.
是和耶穌有真實的.
相信的關係.
包括今天我們對我們的孩子.
他們都知道.
耶穌基督釘十字架復活的事實.
但是否和祂的生命有關.
祂是否真的.
你正在教導祂.
將祂一生整個人的重量.
全部放在主身上.
以至讓耶穌帶著祂做每一個選擇.
還是世界教你怎樣選擇.
跟著選擇吧.
這樣會好.
這樣會成功.
弟子們我們不要有太多的疑惑.
讓他和主有真實的相信連結.
讓他能夠求告主名.
宣認耶穌是主.
不單知道祂是拯救那位.
還要祂的人生跟著這位主宰走.
最後就是我們一生.
我們一切.
今天明天.
整個人的重量都交給耶穌.
法格隆.
隆高.
純孝.
好像這幅圖一樣.
不是隻羊在走.
耶穌背著你走.
弟兄姊妹神愛世人.
甚至將祂的獨生子賜給他們.
叫一切將全人重量都交給耶穌的人.
不致滅亡.
反得永生.
弟兄姊妹這個真理.
今天我們再說.
尤其重要.

$^{881}$我們這幾個星期所經歷的.
很慘痛.
很多人在等著聽到這個得救的道理.
不再靠他們自己的方式.
好像我剛才和你說的.
以色列的新兵.
他知道只需要相信耶穌.
他即使戰死.
他都知道他在哪裡.
弟兄姊妹.
人生無常.
生死的事情.
完全不是我們人能夠作主.
但是我們可以做一個選擇.
今天放下自己的異位.
接受耶穌.
跟實耶穌.
一生都這樣做.
這是福音的核心.
我們一起祈禱.
多謝您我們的祈禱.
在奉耶穌基督的聖名之祈求.
阿們.
\newpage



\section{約翰福音 8:48-59}
\label{sec:032454}
\textbf{榮耀君王:我是救主}~~~~林鍵牧師
\newline
\newline
連結: \href{http://www.church.com.hk/acms/content.asp?site=cdc\&op=show\&type=product\&code=032454\&layout=sermon}{\texttt{032454}} ~~~~ 語音日期: 2026-07-11
\newline
\newline
\hyperref[sec:032453]{\small{< < < PREV SERMON < < <}}
~
\hyperref[sec:index_chronic]{\small{[返順時目]}}
~
\hyperref[sec:032455]{\small{> > > NEXT SERMON > > >}}
\newline
\newline
約翰福音 8:48-59
\newline
\begin{longtable}{cl}
\hline
\hline
章節 & 經文 (和合本修訂版)\\
\hline
8:48 & \begin{tabularx}{0.7\textwidth}{X} 猶太人回答他：「我們說你是撒瑪利亞人，並且是被鬼附的，這話不是很對嗎？」 \end{tabularx} \\ \\ \relax
8:49 & \begin{tabularx}{0.7\textwidth}{X} 耶穌回答：「我沒有被鬼附；我尊敬我的父，你們卻不尊敬我。 \end{tabularx} \\ \\ \relax
8:50 & \begin{tabularx}{0.7\textwidth}{X} 我不尋求自己的榮耀，但有一位為我尋求榮耀，判斷是非。 \end{tabularx} \\ \\ \relax
8:51 & \begin{tabularx}{0.7\textwidth}{X} 我實實在在地告訴你們，人若遵守我的道，就永遠不經歷死亡。」 \end{tabularx} \\ \\ \relax
8:52 & \begin{tabularx}{0.7\textwidth}{X} 於是猶太人對他說：「現在我們知道你是被鬼附了。亞伯拉罕死了，眾先知也死了，你還說：『人若遵守我的道，就永遠不經歷死亡。』 \end{tabularx} \\ \\ \relax
8:53 & \begin{tabularx}{0.7\textwidth}{X} 難道你比我們的祖宗亞伯拉罕還大嗎？他死了，眾先知也死了，你把自己當作甚麼人呢？」 \end{tabularx} \\ \\ \relax
8:54 & \begin{tabularx}{0.7\textwidth}{X} 耶穌回答：「我若榮耀自己，我的榮耀就算不了甚麼；榮耀我的是我的父，就是你們所說的你們的神。 \end{tabularx} \\ \\ \relax
8:55 & \begin{tabularx}{0.7\textwidth}{X} 你們不認識他，我卻認識他。我若說不認識他，我就是說謊的，像你們一樣；但我認識他，也遵守他的道。 \end{tabularx} \\ \\ \relax
8:56 & \begin{tabularx}{0.7\textwidth}{X} 你們的祖宗亞伯拉罕歡歡喜喜地仰望我的日子，他看見了，就快樂。」 \end{tabularx} \\ \\ \relax
8:57 & \begin{tabularx}{0.7\textwidth}{X} 猶太人就對他說：「你還沒有五十歲，難道見過亞伯拉罕嗎？」 \end{tabularx} \\ \\ \relax
8:58 & \begin{tabularx}{0.7\textwidth}{X} 耶穌對他們說：「我實實在在地告訴你們，還沒有亞伯拉罕我就存在了。」 \end{tabularx} \\ \\ \relax
8:59 & \begin{tabularx}{0.7\textwidth}{X} 於是他們拿石頭要打他，耶穌卻躲開，走出了聖殿。 \end{tabularx} \\ \\
[1ex]
\hline
\hline
\end{longtable}
$^{1}$歡迎大家黎到我地ge 周六崇拜.
上年七月我地開始研讀新約聖經當中ge 約翰福音.
我地講到系列ge 題目係.
耶穌榮耀君王.
我地花jor 半年ge 時間去宣講約翰福音第一章至第七章.
今年一月我地開始研讀約翰福音第五章.
兩個月前我地黎到約翰福音第七章.
一個月前我地進入約翰福音第八章.
約翰福音第七章同第八章ge 靈魂同核心.
係耶穌係鑄棚節最後一日大聲宣講ge 二合一自我啟示.
耶穌說:我就是,我就是係耶和華上帝ge 名字.
我就是上帝,我是生命ge 活水.
我是世界ge 光.
不過係呢度作者約翰比jor 我地一個溫馨提示.
第七章第三十九節咁.
佢話:耶穌者話:是子信他的人要受聖靈說的.
不過那時候還沒有賜下聖靈.
因為耶穌還沒有得到榮耀.
因為耶穌仲未上十字架.
所有一切ge 預言耶穌是世界ge 光.
是生命ge 活水必須要等到耶穌上十字架.
讓聖靈降臨ge 時候先至會成全.
因此真正福音ge 大能係黎自耶穌基督ge 榮耀十架.
同埋聖靈ge 洗禮.
係十字架真光同埋聖靈活水ge 應許底下.
呢個淨係講緊一個應許.
因為係將來ge 事呢件事仲未完全發生.
但係呢件事必然好快仲有六個月就會發生.
係十字架ge 真光聖靈活水ge 應許之下.
耶穌係耶路撒冷聖殿駐彭捷ge 最後一日.
開jor 一場萬人報道會.
報道會ge 訊息一共有四部分.
每一部分都有負面ge 同埋正面ge .
呢四部分分別係死在最終人子拯救.
罪的奴隸使你自由.
魔鬼謊言剩若兒女.
跟住係流亡火狐永遠不死.
上個禮拜我地講jor 呢篇好長好長ge 報道會講章ge 第三部分.
魔鬼謊言剩若兒女.
今日我地要講過ge 係整個報道會ge 最後高潮.

$^{41}$流亡火狐永遠不死.
上個周末報道會第三部分完結之後.
我收到一位朋友寫比我一段非常之寶貴ge 反思.
我得到gwo 個作者ge 允許.
今日係報道會終結篇開始之前.
我想先同大家分享.
亦都祈求上帝通過我地去睇呢一段ge 反思ge 時候.
係預備我地ge 心.
呢個係我收返黎ge 非常寶貴ge 反思.
佢話背起十字架跟隨主.
和變賣一切跟從主不是天國大學高材生ge 考卷.
呢個係我上個禮拜講ge .
而是天國幼兒園ge 入門題目.
我從前也跟很多人一樣.
認為這是保羅ge 考卷.
我只是一個小信的人.
讓我慢慢學習慢慢成長.
其實哪個人根本不是想要成長.
只是得過且過地上教會.
盡一切ge 義務.
就覺得自己是天國ge 兒女.
話自己小信.
話自己需要慢慢成長.
原來係得過且過ge 一個藉口.
這是極大的謊言.
極大的欺騙.
我已經看清楚呢個騙局了.
我要做真正ge 兒女.
不是自以為是的兒女.
求主憐憫 求主幫助.
是的 我過去一直都在漠視耶穌ge 十字架.
確實是罪加一等.
我這次真的看見十字架ge 寶貴.
親自嘗了它的口音.
我立志一生要回轉.
將我的餘生交給上帝.
最後 佢話 不管你是上了教會多久.
已經受洗多少年.
在教會服事多年.
甚至已經成為長執領袖.

$^{81}$只要你開始聽明白真ge 福音.
開始對自己ge 信心有懷疑.
你就真係需要停下來.
進入黑房.
單獨面對上帝.
深入反思.
願上帝弘大ge 厚愛.
喚醒我地每一個沉睡ge 心靈.
起來回應耶穌基督如此榮耀ge 福音.
阿們.
帶我地返返去耶穌報道會ge 大圖畫.
如果你有跟住聽我地ge 講道系列.
你會知道從約翰福音第8章33節開始.
之後ge 26節裡面.
亞伯拉罕一共出現過11次.
係今日ge 呢段聖經裡面.
有12節ge 聖經.
只有12節ge 聖經.
而亞伯拉罕出現過5次.
因此.
係我地未正式進入約翰福音今日ge 經文之前.
我想帶你返返去舊約聖經.
先再睇一次上帝對亞伯拉罕ge 呼召.
同埋神人之間.
gwo 個所納ge 恩典請求ge 蒙約.
我想帶你返返去創世記第12章第1至第3節.
要我話對亞伯拉罕說.
你要離開本地本族父家.
往我所指示你的地方去.
我必使你成為大國.
如果你會離開你會前往.
我必使你成為大國.
我必賜福給你.
我必使你的名為大.
你要使別人得福.
為你祝福的.
我必賜福給他.
詛咒你的.
我必詛咒他.
地上萬族都必因你而得福.

$^{121}$我想你知道呢個唔係一個普通尋常的移民呼召.
唔係話你而家所住ge 地方.
gwo D 氣候唔係太好.
於是上帝話不如我帶領你去一個更好的地方.
一個更美ge 家鄉.
如果你上個禮拜係度你有聽我講到ge 時候.
你會記得我當時引用jor 約書亞記第24章第2節.
黎說明亞伯拉罕此行ge 重要意義.
約書亞記第24章第2節咁講.
古時你們的列祖就是亞伯拉罕.
他住在大河那邊.
做乜ye 呀.
服侍必臣.
他在大河那邊服侍必臣.
亞伯拉罕必須要離開那地.
因為那地是偶像之家.
是詛咒之地.
如果亞伯拉罕唔離開佢就必死無疑.
係呢個時候.
你ge 腦袋裡面應該要呈現一個逃難ge 圖畫.
可能你會諗起.
即係年零前逮捕宏福縣ge 大火.
於是亞伯拉罕需要逃跑.
需要逃難.
需要逃生.
呢個呼召.
唔單止係亞伯拉罕ge 經歷.
亦都係所有亞伯拉罕兒女ge 經歷.
係基督徒人生ge 寫照.
你明唔明.
我地必須要逃離我地ge 老我老家.
我地需要逃離巴別巨塔.
我地需要逃離偶像假神探.
我地需要逃離聖經教ge 唔係我講ge .
呢個世界ge 引誘.
約翰一書又係約翰寫ge .
第二章第十五至十七字咁講.
不要愛世界和世上ge 東西.
若有人愛世界.
愛父的心不在他裡面.

$^{161}$因為凡世上的東西.
好比肉體的情慾.
眼目的情慾.
今生的驕傲.
都不是從父來的.
而是從世界來的.
這世界和世上的情慾都要消逝.
都要過去.
都要被上帝所審判.
唯獨那遵行上帝旨意的人.
永遠尚存.
我地要逃離.
我地要逃跑.
但係我地唔單止要逃離.
我地更要前往前往上天.
所指示我地要去ge 地方.
係一個更美ge 家鄉.
係人子ge 國度.
係天上ge 國度.
係永恆ge 國度.
我想你記住呢個圖畫.
因為呢個圖畫.
同我地今日呢段聖經.
有好多ge 關連.
因為今日呢段聖經裡面.
亞伯拉罕每隔一節就會出現一次.
但我地返返去約翰福音.
今日ge 經文.
第八章.
共有十二節可以分為三部分.
經文ge 主題係.
耶穌是誰.
答案有三.
第一.
耶穌係獨一ge 愛子.
傅榮耀.
第二.
耶穌係太初ge 真道.
永不死.
第三.

$^{201}$耶穌係亞伯拉罕應許ge 承傳.
早仰望早死.
列祖亞伯拉罕仰望耶穌.
我地先睇第一部分.
約翰福音第八章第四十八節.
猶太人回答耶穌說.
你是撒瑪利亞人.
並且是被鬼附的.
這話不是很對嗎.
一開波ge 時候.
你發現呢班猶太人怒罵耶穌.
稱佢為撒瑪利亞人.
稱佢為被鬼附.
聽起來.
佢地好毒.
但其實原因係因為.
耶穌先罵佢地.
係耶穌出招.
罵佢地.
罵佢地出於魔鬼.
係魔鬼ge 兒女.
第八章第四十四節.
你地是出於你地ge 父魔鬼.
因為耶穌咁罵佢地.
出於魔鬼.
係魔鬼ge 兒女.
唔係上帝ge 兒女.
唔係亞伯拉罕ge 兒女.
於是.
佢地頂耶穌唔信.
就反罵耶穌.
你先係被鬼附ge .
你係從淫亂而生ge .
因為佢地話佢.
耶穌話佢地唔係亞伯拉罕兒女.
你先唔係.
你係淫亂而生ge .
甚至而家加多一句.
你係撒瑪利亞人.
你係一個愛邦人.

$^{241}$都唔知你個老爸係邊個黎ge .
我知道你阿媽叫做瑪利亞.
我鬼知唔知你個老爸係邊個.
可能係一個愛邦人.
係一個愛邦人同一個.
猶太人所生ge .
雜種.
你冇再另外一個名.
比呢個名更加污糟邋遢.
就係你係撒瑪利亞人.
換句話說.
佢地話耶穌.
你先唔係上帝ge 兒女.
你先唔係亞伯拉罕ge 兒女.
於是.
耶穌回答.
佢地說第八章.
第四十九節.
我沒有被鬼附.
我尊敬我的父.
你們卻不尊敬我.
跟住第五十節.
我不尋求自己的榮耀.
但有一位.
為我尋求榮耀.
他是判斷.
是非的.
耶穌話.
我冇被鬼附.
我所做ge 一切其實係出於.
我ge 父親.
我ge 父親出於.
我對父ge 尊敬.
如果你睇個英文版係honour.
係因為I honour my father.
其實係出於.
舊約聖經第五誡.
當孝敬父母.
耶穌話我所做ge 一切.
係我對我父親ge 絕對ge .

$^{281}$信從因為我係.
第五誡ge 兒子.
我所做ge 一切係基於.
我對天父ge .
孝順.
跟住耶穌意味住.
我所講ge 一切.
頭先我同你講你ge 身份.
你係魔鬼ge 兒女.
你係罪ge 奴隸.
你必死在罪中等等.
你覺得好難聽ge ye .
其實全部都出於我ge 父.
係我父ge 指示.
係我父ge .
啟示.
耶穌話我尊重我ge 父.
但係你地卻不.
尊重我.
你係不尊重我.
你唔尊重我.
咩意思呀.
你話耶穌.
發嬲囉.
你唔尊重我.
我而家唔會尊重你.
你係咪咁諗呀.
唔係架喎.
有冇諗過信不信由你.
耶穌係呢度唔係抱怨.
係嘆惜.
點解嘆惜.
佢地唔尊敬佢.
咱返去創世記第十二章.
第三節.
為你祝福的.
我似普吉他.
OK 講緊ge 係邊個呀.
亞伯拉罕.
其實係講緊亞伯拉罕ge .

$^{321}$其中一個獨一無二ge 後裔.
叫耶穌基督呀.
亞伯拉罕土裡面佢話.
如果有一日你會祝福亞伯拉罕.
土裡面gwo 個.
獨一無二ge 後裔ge 時候.
你就會得到福氣.
所以.
詛咒你的.
其實你睇下.
其實呢個中文翻譯.
我唔知點解佢會咁翻.
其實佢唔應該翻詛咒你.
應該翻不尊敬你的.
Dishonor.
翻不尊敬你的.
我今日朝又早查返去.
希伯來文ge 原文.
希伯來文ge 原文ge 意思係.
唔係詛咒你的就被詛咒.
你唔使.
大聲鬧耶穌.
粗口ge 你知唔知呀.
原文ge 意思係小看耶穌.
小看耶穌.
冇將你.
睇到你應該有ge .
大細尺碼.
咁你就會得到詛咒.
我想問下你耶穌有幾大.
耶穌好大.
好大好大.
你將耶穌.
縮細jor .
你就得詛咒你明唔明呀.
你唔使講粗口鬧佢.
你只要將耶穌.
縮細jor .
我係教會裡面我記得.
早兩個月我講.

$^{361}$小小神服侍人.
大大神差遣民.
耶穌係大大神.
你竟然.
將佢好似黃大仙咁拜.
呢個叫咩呀?詛咒耶穌.
縮細jor 耶穌.
你記得個原文其實.
英文翻譯就好似係.
你不尊敬佢.
希伯來文更有趣.
你少睇jor 佢.
你將佢縮細jor .
你將上帝縮細jor .
上帝係大大神差遣民.
你將佢變jor 做小小神.
服侍你ge 時候.
你死硬啦.
你一定被上帝所詛咒.
所以係呢度.
耶穌話你門卻不尊敬我.
佢係悲嘆緊.
你唔好咁做喎.
你好蝕底.
我唔係嬲你地對我.
我係話比你知.
如果你唔尊敬我.
你大鑊ye .
因為我不尋求自己的榮耀.
但係有一位.
就係我個父.
我個父會幫我尋求榮耀.
最後呢句.
如果你睇佢英文.
He is the judge.
佢係審判官.
唔知點解中文.
判斷是非.
講得好忠誠.
唔係呀.

$^{401}$佢係審判官.
佢話我父會為我尋求榮耀.
有人話.
父會幫我討回公道.
父會為我伸張正義.
會為我報仇.
係唔係.
係不過又唔係.
因為係楊寶欽裡面.
當父要尋求榮耀ge 時候.
成全舊恩.
父要通過上帝.
佢自己獨身的兒子.
為罪人成全舊恩.
呢個叫榮耀.
於是你開始明白.
耶穌其實同佢傳福音.
耶穌唔係話好嬲嬲死我.
你豈有此理.
你不尊敬你唔比respect我.
我當你小強咁打死你.
唔係.
佢係祈求緊佢地.
大佬你有讀過舊約聖經.
你唔可以張誠不尊敬上帝.
因為你不尊敬亞伯拉罕.
不尊敬亞伯拉罕.
gwo 個認許之子.
你知道詛咒會臨到你.
上帝會顯榮耀.
上帝會審判.
然而.
你可以因著十字架而得救.
榮耀.
你去到楊寶欽第十七章.
第十一節.
耶穌要上十字架前ge .
gwo 一刻.
向王父話事後到了.
將你榮耀.

$^{441}$願你榮耀你ge 兒子.
使你兒子也榮耀你.
個父會榮耀佢ge 仔.
點樣榮耀.
當你唔肯悔改ge 時候.
父當然會審判你.
呢個係上帝ge 榮耀.
但當你悔改ge 時候.
上帝會通過耶穌基督ge 十字架.
去拯救你.
呢個係上帝ge 榮耀.
十字架ge 時候事後到了.
楊寶欽ge 尊友評辭.
十字架ge 時候事後係十字架ge 時候.
榮耀係十字架ge 榮耀.
天父要施行拯救.
令我都諗起.
整本楊寶欽ge .
核心經文.
上帝愛世人甚至將佢ge 獨生子賜給他們.
叫一切相信的.
不知滅亡反得永生.
因為上帝差獻佢ge 兒子.
來到世上.
不是要定世人ge 罪.
是要令世人因耶穌.
而得救.
最後有個警告.
因為上帝係終極ge 審判官.
信他的不被定罪.
不信的人已經被定罪了.
因為他們不信上帝.
獨一而至ge 名.
當然我地要問咩叫信.
唔緊要.
下面陸續有來.
不斷會同你講解乜ye 係信.
於是我地繼續去睇.
第八章第五十一節.
耶穌說.

$^{481}$我實實在在地告訴你們.
人若遵守我的道.
就永遠.
不經歷死亡.
我實實在在.
告訴你人若遵守.
我的道.
我的道係logos.
你知道在楊福音裡面.
logos就等於耶穌自己.
從某個角度來講.
你要遵守我的道.
你要遵守我.
好personal好個人.
聖經學者告訴我們.
遵守是說.
黐得好埋.
不是一行一善.
今日做一件.
耶穌覺得開心的事.
一行一善.
是黐得好埋.
耶穌係楊福音裡面經常會講.
你常在我裡面.
我也常在你裡面.
你要常在我的道裡面.
你要常在我裡面.
廣東話做一個比喻.
耶穌講緊.
你要遵守我的道.
你要黐住我.
你要黐實我.
細路你要黐實我.
前面.
有好危險的道路.
不過你黐實我就無事.
但你黐不實我.
我無法保你.
但如果你黐實我.
你就永遠不經歷死亡.

$^{521}$你要super glue yourself to me.
yourself to me.
超級萬能膠.
我唔知你點諗.
你要將耶穌同你的人生.
幫我做個手術.
黐埋手術.
我通常做個手術就分開手術.
不是我而家要做個手術.
將我同耶穌黐埋.
因為我要同他黐埋.
我就永遠不經歷死亡.
如果你熟.
楊福音的時候.
你會知道這節聖經第八章五十一節.
其實有另外一節.
姐妹經文.
三章詞我們已經讀了.
我實實在在告訴你.
第五章第二十四節.
那聽我話.
聽我道的.
又順差我來的那為的.
就有永生.
不至於被定罪.
而且已經出死.
入生.
你開始對應所有的東西.
什麼叫遵守我的道.
就是聽我的道.
然後相信.
我開始明白相信原來.
相信是代表黐住耶穌.
我們一直在說.
什麼叫相信.
我說在神學家的眼中.
相信有三個層次.
三合一的層次.
相信是清楚理解福音的內容.
清楚.

$^{561}$相信是.
從心裡面確信.
確認福音內容的真實.
所以是notitia.
全部三個拉丁文.
心裡面很明白.
人從你個世人都犯了罪.
虧缺了上帝的榮耀.
罪之公近乃是死死後且有審判.
你開始明白那個內容.
第二你相信那個內容跟你有關.
世人都犯了罪當然.
豬阿狗很多人犯罪.
不過不關你的事.
不是.
除了notitia之後.
從心裡面認順.
這個福音跟我有關.
林牧師你上星期說的.
所有事都是關於我.
通常人們說林牧師說得好.
因為都是關於其他人.
我想起很多千千萬萬人值得要聽.
那你呢.
那你只有notitia.
你沒有ascensus.
但是你有notitia.
你有清楚明白.
你有ascensus 心裡確信.
你還要有fiducia.
最後是信靠鬼身.
鬼身.
剛才我們說過.
亞伯拉罕的故事.
亞伯拉罕就是曾幾何時.
上帝的話.
臨到亞伯拉罕.
你有沒有想過.
這個亞伯拉罕會在.
約翰福音第八章不斷出現.

$^{601}$亞伯拉罕的故事就是.
上帝的話臨到亞伯拉罕.
臨到一個信心之父.
亞伯拉罕怎麼回應上帝.
亞伯拉罕馴服.
亞伯拉罕寫記.
寫了他舊世界裡.
所有的東西 他要離開.
亞伯拉罕跟從.
你明白嗎.
所以這個信心是在說一個.
逃跑的故事.
逃離老我老家.
逃離巴別巨塔.
逃離偶像假神.
逃離這個世界.
如果你跟著亞伯拉罕的蹤跡.
前往上帝.
知情更美的家鄉.
你就有永生.
就有永生.
那你就.
不被定罪.
那你就上面那裡.
永遠不經歷死亡.
英文是更加清楚.
他說will never see death.
永遠不見死亡.
永遠.
不見死亡.
在這裡我們要問這個問題.
你說的是什麼死亡.
就是不見死亡.
所有信耶穌的人.
全部都死了.
保羅死了.
彼得死了 約翰也死了.
約翰沒有說自己不死.
有人說彼得好像沒有死.
不是 約翰說他也死了.

$^{641}$全部都死了.
那這裡說的死亡.
約信靠耶穌 黏著耶穌.
就不看見死亡.
死亡說的是什麼.
原來說的不是肉體上的死亡.
而是第二次的死亡.
上帝的憤怒.
地獄的烈火.
永遠的刑罰.
我帶你去看.
我們必須要看的一段聖經.
啟示錄.
又是誰寫的.
又是約翰寫的.
所以這兩本書有很多.
平衡的東西.
啟示錄第21章.
差不多去到整本聖經的.
終結了.
一開始就成了.
他在抄自己.
耶穌說成了.
現在又成了.
我要把生命的泉水.
百百次給那口渴的人.
很面熟的.
生命活水.
第21章第6節.
接著進入我們今天要看的東西.
德勝的.
要承受這些遺孽.
新天新地.
德勝的.
跟之前的那些連繫.
什麼叫德勝.
遵守我的話的人.
黏著我的人.
緊緊跟隨我的人.
聽到我話語.

$^{681}$相信的人.
那種人.
就是德勝的人.
他就會承受.
新天新地.
沒有沒有.
至於膽怯的不信的.
他們將在燒著流亡的火.
湖裡有份.
這是第二次的死.
跟隨耶穌.
黏著耶穌的.
可以逃過流亡的火.
第二次的死.
就是耶穌所說的.
黏著我的.
永遠不見死亡.
你明白嗎.
shall never see death.
但是如果你不跟著我.
你一定和第二次的死亡.
二度死亡.
有關係.
希伯來書第九章.
二十七節這麼說.
我們很熟悉的.
按照定命人人都有一死.
死後且有審判.
人人都有一死.
除非你在耶穌.
回來那一刻之前.
你還沒死耶穌就回來了.
但普遍來說.
人人都有一死.
但是死後且有審判.
有些人會經歷二度的死亡.
那些人.
所謂的二度死亡.
是一個永遠的刑罰.
那個二度的死亡.

$^{721}$是地獄的烈火.
是上帝的憤怒.
我不知道.
怎樣跟你說.
你不用怕死.
但你需要怕死.
你不需要怕第一個死.
但你需要怕第二個死.
第一個死.
是人人都會經歷的.
但你不需要怕.
不是因為人人都會經歷.
你不需要死.
是因為死後會翻生.
但是第二個死就麻煩了.
第二個死.
是你最怕最怕的東西.
我不知道大家記不記得.
我想是一年多前.
在一個聖誕節的聚會裡.
我說過.
我以前有個神學老師.
跟我們講解.
什麼叫第二次的死.
今天我們在這個世界裡.
其實我們這個世界.
充滿著上帝給我們的恩典.
我們叫做普遍的恩典.
這個世界其實在墮落之後.
應該就已經在第二次的死裡.
大家明不明白.
我們應該已經在第二次的死裡.
不停地笑.
不停地狂笑就對了.
但是上帝在他的恩典言文裡.
第一.
容許這個世界繼續存在.
第二 在大家第一次死之前.
給我們有機會聽到福音.
在這個世界裡.

$^{761}$仍然有次序.
仍然有福音.
耶穌來承傳.
在這個世界裡.
我們不會一落街就被人斬.
幾十次.
我們還有可能上班上學.
上寫字樓 上教會.
但是什麼叫第二次死.
就是這個世界裡的普遍恩典.
上帝會停止.
有時我們很不明白.
我們見到一個人.
或者我們自己的親人在醫院裡.
很痛苦.
然後我們說.
終於他.
安靜了下來.
安靜了下來即是死了.
有沒有想過.
這邊安靜那邊不安靜.
你這裡下去那邊上來.
在另外一個世界裡面.
再上來.
如果他不認識上帝的時候.
那個世界就會令人害怕.
有時我們想很辛苦.
我見到一個朋友很辛苦.
我見到一個家人很辛苦.
加點麻煩好不好.
所謂的普遍恩典.
就是上帝在這個世界裡.
已經給了一個很巨大的麻煩.
給大家.
我們在這個世界經歷的最大的痛苦.
都是在上帝.
普遍恩典以下.
是減重了.
止痛了的東西.
但是有一天在第二個審判.

$^{801}$第二個死亡裡面.
那些麻煩會全部再沒有.
上帝再不會有普遍恩典.
這個世界裡面.
最大的痛苦.
在那個世界裡面是最小的痛苦.
你明不明白.
我們沒辦法想像是怎樣的一件事.
但是我必須要.
我們不是經常會講啟示錄第21章.
我們不是經常會講.
第二度的死亡.
我必須要警告你們.
我必須要告訴你們.
你現在面對的是.
第二度的死亡.
第一度的死亡很快.
很快.
沒有了就沒有了.
有沒有想過這裡沒有了.
第二邊又上來.
隔離屋有第二個世界.
那裡有.
二度的死亡.
那個.
沒得救.
那個沒得停.
沒得停.
你明不明白.
這是永恆一件很可怕的事情.
我想告訴你們.
我們不需要經歷第二度死亡.
這是一個選擇.
對於你聽了耶穌的人.
是一個選擇.
這是一個不必要的死亡.
但是你必須要.
遵守耶穌的道.
你要.
癡實耶穌.

$^{841}$Super glue yourself to 耶穌.
你就永遠不見.
死亡.
這個是約翰的訊息.
是啟示錄的訊息.
是約翰福音的訊息.
是同一個人的訊息.
兩本書是同一個人寫.
是同一個上帝的訊息.
因為無論是約翰福音還是啟示錄.
都是上帝所寫的訊息.
啟示錄第一章第二節.
又會和我們說上帝的道.
上帝的道太重要了.
上帝的道既是上帝的話語.
所以為什麼你一定要讀聖經.
上帝的道也是耶穌基督.
為什麼你讀聖經的時候會讀到耶穌.
為什麼所有的一切都是朝向耶穌.
因為.
約翰就將上帝的道.
這個道的原文是Logos.
和耶穌基督的見證.
這個和字有幾個翻譯.
就是耶穌基督的見證.
這個道就是耶穌的自己.
關於耶穌的事.
凡自己所見的都見證出來.
見證出來給你做什麼.
給你抱住.
給你黏住.
中讀者書上的寓言.
聽見又遵守.
所以聽見遵守.
你現在一想起聽見遵守.
就想起SuperGuru.
跟他黏在一起.
永不分離.
將你聽見他的道.
遵守他的道.

$^{881}$所記載的.
那你就有福了.
什麼叫有福.
就是你不需要經歷.
第二度的死亡.
你明白嗎.
因為那個詛咒就是.
你必須要經歷第二度死亡.
因為.
時候緊要.
聖經告訴我們.
我們之所以可以逃出烈火.
進入永生.
是因為上帝的道.
為我們道成了肉身.
不是單單說你黏住耶穌就搞定.
為什麼我們黏住耶穌.
最後可以.
離開逃離烈火.
上到上帝永恆之界.
是因為道成了肉身.
住在我們當中.
約翰福音第一章第十四節.
道成了肉身住在我們當中.
充充滿滿的.
有慈愛有純善.
我改了它.
恩典有真理.
我說更好的翻譯.
其實是翻譯舊約聖經裡面的.
hazard和image.
舊約聖經裡面的慈愛信實.
就是當上帝與亞伯拉罕立約的時候.
這個約.
其實可以拯救世界.
這個約是上帝慈愛的彰顯.
當上帝承傳了這個約的時候.
就是上帝.
信實的彰顯.
一個慈愛一個信實.

$^{921}$慈愛信實就是上帝.
與盟約的發生.
在最後.
道成了肉身住在我們當中.
充充滿滿有慈愛有純善.
接著第十七節.
律法是.
藉著摩西所頒布的.
舊約聖經只是一個影子.
新約聖經才是一個實體.
慈愛與信實現在有個實體的人.
去彰顯.
那個人是.
耶穌基督.
大家知道.
我們由第一章開始.
已經講到現在.
最終上帝對墮落罪人.
所有的慈愛信實.
所有的信實.
榮耀彰顯在耶穌基督的十字架上.
十字架.
是上帝的榮耀.
十字架是慈愛.
十字架是信實.
十字架是.
我們可以蒙恩得救的原因.
我們抱著十字架.
我們捨己.
跟隨耶穌.
就可以不看見.
第二次的死亡.
我想你這樣想.
如果道成肉身.
道成了肉身.
由天降下.
是耶穌的離開.
祂的本地本族父家.
榮耀的十字架就是天父上帝.
指示祂要去的地方.

$^{961}$祂離開本地本族父家.
道成了肉身.
但還要去一個地方.
去十字架.
耶穌最後的人生的終點就是十字架.
是因為那是上帝指示.
祂去的地方.
於是你開始.
更加.
深刻的明白.
整個福音我們繼續.
第八章第五十二節.
於是猶太人對他說.
現在我們知道.
你是被鬼附的.
亞伯拉罕死了.
眾先知也死了.
你還說.
人若遵守我的道.
就永遠.
不見死亡.
猶太人聽不明白耶穌所說的.
當然他們聽得明白.
也不會接受耶穌的警告.
耶穌的警告是永世.
第二次的死亡.
接著耶穌也讓他們認祂.
認祂是永生.
不單止不用經歷第二次的死亡.
還有永遠的生命.
我今天也要和你說同一樣東西.
我不知道你.
會不會接受耶穌的警告.
你先想想警告.
你再想想.
耶穌的應許.
耶穌的警告 耶穌的應許.
全部在他的話語裡面.
但是.
當時的猶太人是聽不明白的.

$^{1001}$他們以為耶穌說的是.
第一次的死亡.
他們以為耶穌說的是肉身的死亡.
於是他說.
阿拉罕死了.
眾仙姿死了.
眾仙姿包括.
摩西和大衛.
所有以色列人.
歷史裡面最偉大的人物.
每個人都死了.
你還說.
人若遵守你的道.
就永遠不經歷死亡.
所以這是他們的問題.
因為他們想的是第一次死亡.
於是他們繼續.
難道你比我們的祖宗.
阿拉罕還大嗎.
祂死了 再次重複.
眾仙姿也死了.
你把你自己當作.
什麼人.
背後的邏輯是這樣的.
如果耶穌可以.
真的可以使人不死.
第一次死亡 他不是說第二次死亡.
他想不到那麼遠.
又或者想到第二次死亡跟他們沒有分別.
他說如果耶穌.
可以使我們不死.
他自己必須要有不死之身.
他說難道.
你自己有不死之身.
這個大概.
是他們的邏輯.
於是他們在數.
阿拉罕最大的也死了.
所有的仙姿都死了.
你當你是誰.

$^{1041}$你怎麼可能逃過死亡呢.
這裡我自己加了一句.
這句.
可能是他當時有說.
也可能是他心裡的說話.
他就說耶穌.
榮耀你自己.
為什麼我要加這句說話.
因為這句說話是對應之後耶穌的回答.
他們一定是在想.
耶穌在榮耀他自己.
所以之後那句說話.
耶穌是說榮耀的.
耶穌回答說.
我若榮耀自己.
我的榮耀就算不了什麼.
榮耀我的是我的父.
就是你們所說的.
你們的上帝.
在這裡.
雞同鴨講的情況.
繼續出現.
猶太人以為耶穌.
是高抬自己榮耀他自己.
就好像他們很喜歡.
跟別人比較.
他把自己說得那麼大.
你現在大過阿拉罕.
你大過其他仙姿.
你大過摩西大過大衛.
他們以為耶穌是像他們那樣.
喜歡跟別人比較.
接著抬高自己.
接著去建立自己地上的榮耀.
是從父那裡來的榮耀.
是從天上降下的榮耀.
是救贖的榮耀.
是十字架的榮耀.
你明不明.
雞同鴨講的.

$^{1081}$他們將他們自己想的榮耀.
代入耶穌.
其實耶穌從來沒有追求地上的榮耀.
耶穌告訴他們.
我的父會榮耀我.
他們以為耶穌跟他們鬥.
其實耶穌想告訴他們.
我的父會將我送上.
十字架.
他們誤以為.
耶穌在炫耀自己.
貶低他們.
其實耶穌一直在呼喚他們.
迴轉,悲改.
有一天你會仰望.
那個十字架現在還沒到.
但過了半年就會出來.
有一天你會見到榮耀的十字架.
你要仰望這個榮耀的十字架.
你就可以得到救贖.
你就可以成為真正天賦的兒女.
你現在不是.
你現在是魔鬼的兒女.
但有一天你仰望十字架.
黏著它,抱著它.
你就可以成為天賦的兒女.
耶穌繼續.
是報道會.
耶穌說你們不認識他,我卻認識他.
我若說我不認識他.
我就是說方的像你們一樣.
但我認識他.
我也遵守他的道.
大家知道.
認識不代表認識那麼簡單.
認識代表愛慕.
耶穌說.
你們不愛慕上帝.
但我卻心愛我的天賦上帝.
我若說我不愛慕上帝.

$^{1121}$我就在說方.
很多人.
很難這樣說.
有人跟我說林牧師.
如果我說我不愛上帝.
我真的在說謊.
你就是在說謊.
但耶穌這樣說就不一樣.
只有他才可以.
我若說我不愛上帝.
我就在說方.
但我告訴你,我真的愛他.
我怎樣愛他?我遵守他的道.
耶穌說我們愛他.
我們遵守他的道.
現在耶穌愛上帝.
耶穌遵守上帝的道.
他心愛上帝,他馴服上帝.
他怎樣馴服上帝的道?.
他甘心樂意,道成了肉身.
住在我們的當中.
他甘心樂意降生成人.
他馴服至死.
被舉十字架.
為要.
遵行上帝給他的道.
要成全上帝.
差遣他來世上要成全.
的救贖.
耶穌還沒說完.
他繼續.
你們的祖宗亞伯拉罕.
他.
歡歡喜喜地.
仰望我的日子.
他看見了.
他就快樂.
原來當時有.
很多的猶太拉比.
都相信亞伯拉罕有.

$^{1161}$千里眼信奉儀.
他們都相信.
亞伯拉罕睡覺的時候.
上帝會給他一些意象.
亞伯拉罕可以見到.
先知般的遠見.
亞伯拉罕夷夢也好,意象也好.
他可以遙遠.
看見將來.
上帝給他的.
恩典盟約的承傳.
那個未塞亞.
時代的開始.
上帝給他一個很宏大的.
認許,他死的時候還沒見到.
是什麼?.
他說,你必成為大國.
你看看我兒子.
變不變成大國.
你必成為大國.
你必得福.
你的名必為大.
你會使其他人得福.
地上的萬族.
都會因他得福.
這些東西.
要遙遠張望.
才可以見到.
沒有承傳在他身上.
沒有承傳在他自己的耳塞身上.
在承傳他耳塞不知幾多代之後.
有一個獨一無二.
亞伯拉罕的後裔.
耶穌基督.
會承傳在耶穌基督的身上.
所以我想告訴你.
耶穌講這句說話.
其實當時的猶太人可以明白.
亞伯拉罕歡歡喜喜.
仰望的日子.

$^{1201}$與其同時.
如果你有機會讀舊約的先知書.
你會發現舊約的先知書裡面.
經常講耶和華的日子.
耶和華的日子.
這個也是眾所周知.
猶太人所期待耶和華的日子.
所以.
耶穌在講的是.
尼采亞的日子.
耶和華的日子.
但是這裡耶穌做了一個小動作.
大改造.
是什麼呢?.
他將耶和華的日子.
尼采亞的時代.
改變了成為.
我的日子.
這件事.
是不是應該很生氣呢?.
你將.
耶和華的日子.
變成了你的日子.
我們在想.
亞伯拉罕應該見到尼采亞的尼郎.
然後你說又是你?.
令我想起耶穌很喜歡做這件事.
在《約翰福音》第四章.
他與聖瑪利亞婦人的對話裡面.
我讀一讀給你聽.
第四章第25節.
婦人對耶穌說.
我知道尼采亞.
就是那稱為基督的.
祂要來.
祂來了會把這一切事.
告訴我們.
然後耶穌下一句.
我就是.
和你說話的.

$^{1241}$就是.
耶穌很喜歡將事情改了.
別人跟祂說尼采亞.
祂說就是我.
別人跟祂說耶和華的日子.
就是我.
我就是.
你們期待已久的.
尼采亞.
我就是耶和華.
你們的上帝.
祂仰望的是.
耶和華的日子.
祂仰望的是我的日子.
如果你聽得明白的時候.
你會覺得大逆不道.
褻瀆上帝.
但是.
幸好當時猶太人.
是聽不明白.
猶太人說.
他這一刻應該已經拿石頭.
扔他了.
但他想不通.
猶太人說.
50歲難道你見過.
亞伯拉罕嗎.
亞伯拉罕是2000多年前的時候.
耶穌說我實實在在.
告訴你們.
還沒有亞伯拉罕的時候.
我就存在了.
我們先說第一句.
他們以為耶穌在癡人說夢話.
當然他們不太相信.
亞伯拉罕.
基本上是跟著那句.
耶穌說.
亞伯拉罕看見了我就快樂.
看見你快樂.

$^{1281}$你跟亞伯拉罕見過面.
你都不夠50歲.
30歲而已.
你這麼年輕.
亞伯拉罕2000年前就死了.
難道你有2000多歲.
你說笑什麼.
但他們其實誤解了.
耶穌所說的話.
耶穌將他自己代替了耶和華.
這次.
耶穌不容許他們繼續誤會.
於是.
他堅持要啟示.
他屬天屬上帝的身份.
於是有第58節.
耶穌這次不放過他們.
耶穌說我實實在在告訴你們.
還沒有亞伯拉罕.
我就存在了.
更好的翻譯可能是.
在亞伯拉罕生前.
我就是.
OK.
也就是說.
如果耶穌當時的說.
在亞伯拉罕生前.
我就已經出生.
英文就是.
Before Abraham was, I was.
如果他都是用過去式的話.
就是說.
亞伯拉罕.
是生在這個世上.
不過我比他生得更早.
他又說你騙人.
你說你還老過亞伯拉罕.
但耶穌不是這麼說.
耶穌說Before Abraham became.
曾幾何時.

$^{1321}$亞伯拉罕不屬於這個世界.
在亞伯拉罕未出生之前.
我永遠長存.
I am.
現在式.
他不是說過去式.
而I am.
如果你們知道我們一直說.
約翰福音是代表了上帝.
上帝的名.
約翰將我們帶回去.
約翰福音第一章第一字.
太初有道.
都與上帝同在.
都就是上帝.
而耶穌是永恆真道.
在亞伯拉罕未成形.
未出生之前.
我永遠長存.
於是.
耶穌說.
你要生命你要得福.
你要不死.
你就要抱著我.
其實帶我們回到這件事.
你抱著的不是.
某一個先知的說話.
你抱著的是上帝.
我實實在在告訴你.
人若遵守我的道.
人若在我裡面.
在我的道裡面.
不是日行一善.
不是牧師說.
這個禮拜人人每日都要.
一件事做.
有沒有試過這樣的事.
我小時候都被人騙得很透.
你回教會有時說.
這個禮拜試試.

$^{1361}$每日找一件事.
每日讀兩節聖經.
每日做一個應用.
下個禮拜回來.
一日一個應用.
七日就有七個應用.
你回來可以覺得感覺好.
神經病.
耶穌不是跟你說這樣的事.
耶穌跟你說.
你要黏著我的道.
你要常在我道裡面.
常在我裡面.
你黏著我.
你就永遠不會經歷.
第二次的死亡.
大家真的要想想.
第二次的死亡.
第一次死亡不用怕.
我知道大家都怕第一次死亡.
但你不應該怕第一次.
第一次死亡就這樣沒了.
就這樣沒了.
很快的.
第二次死亡就慘了.
你想停都停不到.
第二次死亡是沒有麻痺針.
所有的止痛劑都停了.
再沒有上帝普遍的恩怨.
上帝的烈火濃濃的.
去衝向你自己.
但又不死.
永遠的不死.
是一個我們不可以想像.
但卻又最適合你.
最公道.
最公平.
對於不迴轉罪人.
最公平聖潔.
的審判.

$^{1401}$但你不需要經歷第二度的死亡.
你要去聽耶穌的道.
你要去信耶穌的道.
你要清楚的明白.
你要從心裡確信.
你要信靠.
鬼神.
你要記住你的祖宗亞伯拉罕.
如何離開本地.
本族富家.
去往上帝指示他去的地方.
你要和亞伯拉罕一樣.
離開.
你的老家.
你的老我.
離開你的芭比巨塔.
離開你的偶像假神.
離開這個世界.
不要愛這個世界.
如果你跟從.
這個腳蹤.
你就有永生.
你就不被定罪.
你就永遠看不見第二度的死亡.
我們又要.
好像我們的祖宗亞伯拉罕.
歡歡喜喜.
仰望耶穌的日子.
看見他就.
歡喜快樂.
當你聽到.
福音的時候.
你就好像亞伯拉罕一樣.
仰望耶穌的日子.
你看見耶穌.
你就歡喜快樂.
你想更加親近的去跟隨他.
你想和他黏著.
抱著.
永遠不分離.

$^{1441}$我們要.
看見耶穌.
給我們.
榮耀的福音.
是他離開了.
本地本族.
父家.
往上帝指示他的地方.
那裡去.
我們的離開.
是地上的離開.
從下面上去.
耶穌的離開是從上面下來.
耶穌從上面下來.
再上去.
嘩.
真是一條魚.
兩個菜.
三次聖經.
有兩個宣講.
我們怎樣可以得救.
第一我們需要耶穌承傳.
這個恩憐的盟約.
他離開本地本族.
父家去上帝指示他的地方.
接著在十字架上.
耶穌呼召我們.
我那邊沒有整領.
我們千千萬萬個小我們.
我們由地上.
走去.
天上.
跟隨上帝.
耶穌來為我們承傳了.
救恩的盟約.
承傳了你和我的.
救贖.
讓地上的萬族萬民.
可以因信耶穌.
什麼叫信耶穌 再想想.

$^{1481}$有沒有信耶穌 有啊 什麼時候.
做過 不知道多少年前.
做過的祈禱 不是.
有沒有黏著耶穌 有沒有超級黏著耶穌.
有沒有他去到.
那裡 我永不離開他.
他要我去那裡.
我就去那裡 他去那裡.
我就要去那裡 這個叫做信.
我們要跟著他.
直到永恆.
但我們回到今天講到的.
大圖畫 耶穌是誰.
耶穌到底是誰.
答案有三.
獨一愛子 父榮耀.
太初真道 永不死.
應許承傳.
早仰望.
我們應該怎樣.
回應耶穌.
報道會的訊息.
當耶穌現在報道會.
我們現在要結束了.
我們應該回轉悔舉.
我們應該全人全心.
信靠.
跟隨他.
我們真的要相信他.
但那個信不是.
我們說跟他一起做個祈禱.
我們的信是.
一生的.
貼著他 一生的貼實他.
與耶穌一起.
永遠不見.
第二次的死亡.
這個是唯一.
正確合理的回應.
聽完這個報道會之後.

$^{1521}$可惜的.
保護的是猶太人.
不願意相信.
不願意降服耶穌.
你看猶太人怎樣回應耶穌.
這次終於聽得明白了.
搞了這麼久.
最後耶穌那句.
我事終於搞定了.
於是他們拿石頭要打耶穌.
耶穌卻躲開.
走出了.
聖殿.
這次猶太人終於聽明白了.
耶穌的說話.
他們心硬不相信.
叛逆不跟隨.
他們拿石頭打耶穌.
因為他知道耶穌.
你講這句說話.
等於你將自己變成上帝.
你攝逐神靈.
於是他們拿石頭要打耶穌.
耶穌卻.
避開了.
走出聖殿.
奧古斯丁有句名言.
奧古斯丁看到這一節.
八章五十九節.
心裡慨嘆.
他說.
耶穌逃過了石頭.
但那些.
心如鐵石的人.
有禍了.
因為上帝離開了他們.
大件事了.
耶穌離開了聖殿.
上帝離開了聖殿.
上帝的榮耀離開了聖殿.

$^{1561}$聖殿可以拆毀.
因為是賊窩.
奧古斯丁講得好.
耶穌逃過了石頭.
耶穌也逃過了.
這些心如鐵石的人.
耶穌任由他們.
從此.
他們再無希望.
我們又如何.
如果我們今天不回應耶穌.
耶穌也離開我們.
只剩下你和我.
心如鐵石的人.
上帝離開了我們.
我們就在等.
等不是.
第一次的死亡.
等第二次的死亡.
第一次死亡無所謂.
第二次死亡.
是你必須要懼怕的.
今天仍有時間.
讓我們重寫.
改寫我們的人生.
這是我自己作出來的.
這是我們應該回應.
於是我們走進.
十甲的真光.
Walk in the light.
跟著那個真光而走.
我們領受聖靈.
活水的施洗.
我們重生.
誕誕為主而活.
求主幫助我們.
讓我們回應.
耶穌報道會上.
福音的訊息.
我們去祈禱.

$^{1601}$多謝你.
終於讀完.
楊福音第八章.
終於讀完耶穌.
那篇.
法人心醒.
去到盡.
的報道會.
福音訊息.
很難想像.
一個人.
頂著.
上千上萬的猶太人.
耶穌自己都開了.
四張底牌.
你們想要殺我.
去到最後.
拿起石頭.
要打死耶穌.
但真的死的不是耶穌.
是罪人.
因為上帝的榮耀.
離開了他的聖殿.
如果你熟悉舊約聖經.
你知道最得人驚.
的地方就是當上帝離開了.
他的聖殿.
因為起初上帝建立聖殿的時候.
他的榮光會.
進入聖殿.
但現在上帝離開了.
因為這個賊窩.
將要.
接受上帝.
終極的審判.
主耶穌幫助我們.
因為我們都在一個.
非常危險的狀態.
裡面.
除非我們領受耶穌.

$^{1641}$跟隨耶穌.
癡實耶穌.
如果不是我們的下場.
跟他們一樣.
幫助我們思想.
耶穌所說的.
每一句說話.
耶穌說.
我們需要.
信祂.
因為我們是罪的奴役.
因為我們會死在罪中.
因為我們是撒旦魔鬼.
謊言的兒女.
因為有一天.
除非我們相信祂.
我們必在.
流亡的火湖裡面.
永遠燃燒.
但耶穌不願意我們.
走向.
那個永恆的滅亡.
在四段.
的故事裡面.
耶穌不斷呼喚我們.
人子將要為我們.
被舉起.
罪的奴役.
奴隸可以得釋放.
撒旦的兒女.
可以成為聖若宗的兒女.
最後.
流亡與火的人.
可以.
永遠不死.
我們天父.
求主你感動我們的心.
光照我們的心.
讓我們回應你.
信靠你.

$^{1681}$一生跟隨你.
我們這樣祈禱 奉主耶穌基督名求.
阿們.
\newpage



\section{加拉太書 5:13-15}
\label{sec:032455}
\textbf{自由的真諦 (華)}~~~~葛國光牧師
\newline
\newline
連結: \href{http://www.church.com.hk/acms/content.asp?site=cdc\&op=show\&type=product\&code=032455\&layout=sermon}{\texttt{032455}} ~~~~ 語音日期: 2026-07-12
\newline
\newline
\hyperref[sec:032454]{\small{< < < PREV SERMON < < <}}
~
\hyperref[sec:index_chronic]{\small{[返順時目]}}
~
\hyperref[sec:032456]{\small{> > > NEXT SERMON > > >}}
\newline
\newline
加拉太書 5:13-15
\newline
\begin{longtable}{cl}
\hline
\hline
章節 & 經文 (和合本修訂版)\\
\hline
5:13 & \begin{tabularx}{0.7\textwidth}{X} 弟兄們，你們蒙召是要得自由；只是不可把這自由當作放縱情慾的機會，總要用愛心互相服侍。 \end{tabularx} \\ \\ \relax
5:14 & \begin{tabularx}{0.7\textwidth}{X} 因為全部律法都包括在「愛鄰如己」這一句話之內了。 \end{tabularx} \\ \\ \relax
5:15 & \begin{tabularx}{0.7\textwidth}{X} 你們要謹慎，你們若相咬相吞，恐怕要彼此消滅了。 \end{tabularx} \\ \\
[1ex]
\hline
\hline
\end{longtable}
$^{1}$加勒賽書第五章.
我們紀念是十三節到十五節.
在前面裡面.
保羅在這邊講說受割禮不受割禮.
他說什麼.
這都不是最重要的.
重要的是使人深發仁愛的信心才有功效.
對不對.
第五章在這個地方講到第六節.
然後他說你們向來跑得好.
有誰攔住你們呢.
然後就提了說.
你們的一點面教人死全團都發起來.
記得嗎.
因為當時候有一些從耶路撒冷來的這些母會的.
有一些的這些猶太教的背景的基督徒.
他們說你們信耶穌啊.
你們外邊的信耶穌還不夠.
你們還必須要受割禮.
所以保羅才寫了加勒賽書這本書.
嚴厲的告訴他們.
保羅的語氣裡面看整個加勒賽書是非常嚴厲的.
你很少看到保羅在其他的書信有這麼嚴厲的在講.
因為他說你們傳的是別的福音.
所以整個加勒賽書裡面講到這個律法主義.
也就是靠律法得救與靠恩典.
靠信心得救的一個差別.
他一直講講講從第一章一直講到第五章.
他還在講這件事情.
所以因為神的話這麼看重.
所以不是我要反覆的講這個東西.
是因為神的話一再的講.
接著你大概有這樣子的背景之後.
我們來看一下.
第三節他就提到說弟兄們你們蒙昭是要得自由.
這句話我想保羅一開始在這邊提到.
你們蒙昭是要得自由.
自由 兄弟妹 好不好.
想不想自由.
其實今天很多生活在中國大陸的人.

$^{41}$他的生活坦白說比你在這邊過得都好.
懂嗎 你們國內都知道我在講什麼.
現在國內真是方便到一個地步.
你坐動車 你覺得現在肚子餓了.
然後你到了下一站.
你跟那個 我不知道他叫 我沒有叫過.
你只要你APP上面叫一下.
你說我要吃牛肉麵.
你說我要吃什麼.
隨便你點 蔥油餅.
你到了站的時候馬上就送上來了.
鐵桶水 生活真方便.
但是為什麼還是有這麼多人要run出來.
為什麼.
因為還是有相當程度的不夠自由.
懂這意思吧 是不是.
那在美國實在太自由了 是不是.
自由不自由 也不自由.
那你到天國去最自由了.
在1990年的時候.
台灣有一個歌星.
這個叫李明依 聽過嗎.
沒聽過 那就知道你哪個年代的人了.
你大概是五六年級的人.
你就知道我在講什麼.
那麼李明依她當時候這個.
她在一個廣告 是廣告那個蜂蜜檸檬水.
叫做奇萌子.
她就講了一句話 她說.
只要我喜歡有什麼不可以.
那麼這句話雖然是一個廣告詞.
後來也編了歌.
但是這句話在當時候.
在台灣那種比較戒嚴時期.
剛剛開始開放.
那麼其實是相當具有一個震撼性的.
為什麼呢.
甚至到後來教育部.
後來的一些的媒體.
起了一些的論戰.

$^{81}$就認為說李明依教壞了我們的孩子.
因為她在挑戰當時候的一種威權.
你懂意思嗎.
她認為她在鼓動了這些年輕的人.
讓他們 可是事實上後來你會發現.
這個東西到現在的台灣.
甚至如果今天我們在美國.
你更會發現這句話根本是稀中平常.
你的孩子今天你管他你怎麼管.
你沒有辦法像以前那樣管.
所以那這個自由.
像這樣子的自由好不好.
只要我喜歡有什麼不可以.
保羅也說保羅說你們蒙昭就是要得自由.
對不對.
你們都沒有Amen.
但是保羅在這邊講到蒙昭要得自由.
你要了解當時候的背景環境.
從猶太教出來.
在猶太教裡面你必須要受隔離.
男子生下來第八天受隔離.
你必須要這個那個就業裡面的.
有很多的立法的規定條條框框告訴你.
所以你知道嗎.
到後來在新約時代.
耶穌責備最厲害的就是這些法利賽的.
對不對.
為什麼呢.
因為他們利率加利率規定了很多的繁文縮節.
把這些該做什麼.
禮拜天.
這不叫禮拜天.
它是安息日不可以做什麼.
不可以做什麼.
洗手要怎麼洗髮.
通通都規規矩矩.
我不要講說猶太教這樣規定.
我請問你.
今天我們教會規定.
你上廁所的時候.

$^{121}$你只能夠用一張紙擦手.
你會不會覺得被騙子.
不會.
可能你不會.
我們下一代一定會覺得.
還有這種規定啊.
一張紙怎麼能夠擦手.
我不要做調查.
我基本上面一張紙來擦手.
親愛的同學.
可是如果它變成是一個規定.
教會裡面如果有很多這一類的規定.
你走進禮拜堂.
你不可以說話.
其實我是鼓勵你.
你進入禮拜堂.
你真的不要喧喧嘩嘩的像菜市場一樣.
但是我要講的意思就是說.
親愛的同學.
律法.
這個舊約的很多的律法裡面.
它已經是超過了聖經裡面所說的.
所以保羅說你們蒙昭是要得自由.
他在強調的就是.
過去你們活在律法的底下.
你們活在舊約的這個律法的要求之下.
你必須要靠著行律法.
你才能夠得救.
你才能夠進入神的國.
對不對.
但是今天我告訴你們.
弟兄們你們蒙昭是要得自由.
馬丁路德在1517年改教.
當時候改教其中一個最主要的一個引爆點是什麼.
就是贖罪戒.
天主教發送這個贖罪戒.
他們叫做大赦.
我們知道美國總統有大赦.
但是天主教也發明了這個大赦這個用詞.
這是在1095年的時候教宗叫做烏爾班二世.

$^{161}$他發動第一次的十字軍東征.
那麼那時候開始.
當時候他宣佈凡是參加十字軍東征的人.
可以獲得教會減免罪罰的恩賜.
怎麼.
這個對我們今天基督教來講.
我們就有一個問題了.
什麼.
還有減免罪罰.
罪罰是什麼.
原來當時候教宗這樣子的規定.
那麼後來再加上商業的推展.
煉獄說還有功德庫這個理論.
他們的神學理論的發展.
教會慢慢的就把這個大赦發一個證書.
發一個證明書.
那麼就是因為這樣子的關係.
到了最高點就是教宗當時候為了籌措.
這個聖殿的這個.
你們到那個.
到那個哪裡到那個羅馬去看.
那個聖保羅的那個聖殿.
那個禮拜堂他要蓋的時候需要大筆的資金.
所以他就派著神職員到各歐洲各國.
當時候因為是神職人物.
歐洲都是信天主教的.
去宣傳告訴他們說要奉獻.
那麼這個奉獻到一個地步.
甚至是靈前叮噹響.
歷時靈魂上天堂.
到達這種地步.
那個宣傳人員.
所以在這種情況之下.
馬丁路德當時候就開始了這個改教.
直到今天雖然經過了改教之後.
天主教也做過一次的天特會議裡面.
做了一些的修正.
可是這種所謂的煉獄,功德庫.
還有這種所謂的這種大赦的觀念.
這個神學教義依然在天主教的裡面.

$^{201}$只不過他們廢止了.
用這種所謂的金錢上的交易來換取贖罪.
現在他們比較強調心靈的洗滌等等這些.
但是我要問一件事情.
親愛的女子們.
任何想要不管是有沒有牽扯到金錢.
任何想要藉著靠著自己的力量.
行好事,做好事來獲取恩寵的.
上帝的赦免的.
這是什麼.
這不是聖經的教導.
是不是.
不是.
因為德就是本乎恩.
因者性.
保羅說得清楚.
不是出於自己.
對不對.
也不是靠著我們的行為.
是因者恩.
就是恩典.
什麼叫恩典.
你如果付出了任何一文錢.
你如果做了任何一件事情.
這就不叫恩典.
懂這意思嗎.
所以保羅一直在強調這件事情.
馬丁路德改教改教的.
這麼樣子的付上生命的危險.
就是為了要強調一件事情.
因性成異.
好了.
那聽我準備.
講到這裡我們都沒問題了.
講到這裡我們真是感覺到.
我們要為著這些歷史歷代.
這些神的僕人使禮.
這些使徒先知.
還有馬丁路德加爾文.
這些改教的領袖們.

$^{241}$他們付上甚至可能有這種生命危險的.
他們所帶來的結果.
使得我們今天可以在這邊.
自由的來敬拜聖母.
但是保羅接下來提醒我們.
他說只是不可將你們的自由.
當作放縱情慾的機會.
我們中文的翻譯放縱情慾.
很容易你就會想到什麼.
想到什麼.
這個是我們也趁著機會.
大家知道怎麼讀經.
放縱情慾你會怎麼想.
你會怎麼想.
很簡單.
我喜歡.
直接了當.
是不是.
我們一看到放縱情慾.
我們就會想到.
特別是姐妹們可能會覺得.
這跟我沒關係.
這句話.
聖經這句話跟我沒關係.
我不會放縱情慾.
對不對.
但是原文裡面.
其實他這個放縱情慾的機會.
機會就是一個藉口.
就是一個occasion.
就是一個opportunity.
所以他說不是.
那放縱情慾呢.
其實這句話放縱情慾的機會.
原文的意思就是.
不要給肉體留機會.
不要給你的肉體.
an opportunity for the flesh.
請問.
弟兄你有沒有肉體.

$^{281}$姐妹你有沒有肉體.
是不是只有弟兄會有肉體.
只有弟兄的肉體.
會有機會犯罪呢.
不會.
姐妹會不會有肉體.
一樣有.
什麼叫肉體.
肉體聖經裡面.
保羅強調的.
長著這個肉體.
就是我們裡面的這個救人.
就是我們裡面這個老我.
有沒有.
夫妻之間的問題.
最大的是哪裡.
你知道嗎.
是這個救人.
是這個肉體.
是老我.
所以保羅在這裡說.
不要把你們的自由.
當作放縱情慾的機會.
意思就是說.
這個自由不要成為你們的藉口.
不要成為你們的一個.
肉體要去不遵行神話語的一個藉口.
一個機會.
一個opportunity.
不過.
剛才講到這個問題.
我們還是先從這個.
這個性方面這個問題.
來談這個問題.
根據統計.
2026年5月的統計.
美國人對LGBTQ.
方面的這一個權益的支持度呢.
我們來看看.
右邊這個圖表.

$^{321}$有三個你看了很明顯的數據.
有沒有.
左邊的65\%.
65\%.
就是他們認為同性婚姻合法.
同性婚姻合法.
他們支持.
全美國65\%的人支持同性婚姻.
第二個62\%是指同性戀的關係.
聽懂我意思嗎.
結婚了不用講.
那已經是真正事實.
同性戀的關係.
就算他沒有結婚.
但是他們有同性之間的關係.
是可以接受的.
這是一般美國人的.
那37\%是什麼呢.
是指的他們接受所謂的變性.
他們對於性別的問題.
變性啦跨性別啦.
他們接受.
美國有37\%的人.
OK好.
那我們看看基督徒的比例.
基督徒的比例.
支持同婚的這個支持率.
在天主教還有主流的新教.
什麼叫主流新教呢.
就是像長老會 信義會.
還有呢 衛禮公會 聖公會.
路德會跟信義會是差不多的.
基本上這些所謂的.
Main Street的這個主流的基督教的教派.
現在都淪陷了.
除了美蘭基信會.
這些大教派都淪陷.
意思是什麼.
他們全都去支持同婚.
因此造成了他們本身教會內部的分裂.

$^{361}$因為不贊成的人就離開了.
離開了教會了.
那麼所謂的福音派的基督徒呢.
有35\%的人也接受同婚 支持同婚.
福音派是基督教為了區別新派 自由派.
所以叫做Evangelical.
我們教會是Evangelical Church.
正道神學院叫做Evangelical Seminary.
華神叫做China Evangelical Seminary 懂嗎.
為什麼要加個Evangelical.
就表示說福音 聖經.
聖經是我們所相信的.
福音是我們今天教會基督徒存在的最重要的使命.
因為新派神學告訴你說.
福音 普究論全部的人都可以得救.
那還要傳福音幹嘛.
是不是.
甚至這些新派神學裡面把所謂的愛.
把所謂的做好事去救濟窮人.
當作是他們最高的使命.
而把福音丟掉了.
OK好 那麼.
連福音派的美國基督徒.
都已經有35\%的人支持同性婚姻.
下一代更慘.
下一代18到29歲的年輕人.
而且是福音派的基督徒.
有53\%的將近一半的比例.
他們是支持同婚的.
那今天你的孩子.
你看著他.
你不要以為你需要跟他早點談這些問題.
你不談 世界上的人都在影響他.
很多的孩子.
甚至不要等他到大學.
他的思想觀念已經被整個.
自由派 左派的思想所影響了.
好 那第五準備.
這個是做什麼.
什麼叫自由.

$^{401}$我們蒙昭德自由.
但是親愛的兄弟妹妹.
這個自由不是放縱.
這個自由是在基督裡.
我們在前面講了很多次自由對不對.
什麼叫真自由.
基督說 主耶穌說.
神的兒子叫你們自由.
你們就真自由了 是不是.
又說什麼 真理必叫你們什麼呀.
得以自由 是不是.
今天我們的宣昭經文.
主的靈在哪裡.
哪裡就得以自由.
所以你看見沒有.
自由是有一個範圍的.
是在真理的裡面才有自由的.
不在真理裡面.
那不叫自由 那叫放縱.
所以我們基督徒.
當然我相信可能我們多數的弟兄姊妹.
我們不會去認同這裡所說的.
支持同婚 支持跨性別.
但是親愛的兄弟妹妹.
我們可能有其他的方面.
在這一個所謂自由派的思想.
文化的影響之下.
我們的也漸漸開始會被影響.
好 那麼求主幫助我們.
那我們真的是.
不要讓任何的一點點的破口.
我們不要留任何一點點的破口.
嘗試習慣性的去操練.
當有一些思想觀念來的時候.
第一個你回去問說.
聖經怎麼說.
來 跟你旁邊說.
聖經怎麼說.
大部分的這些問題.
聖經都會告訴你.

$^{441}$可能有一些細節的地方.
聖經也許不一定馬上有直接了當的.
但是它的原則總是可以給我們代理.
親愛的兄弟妹妹.
如果基督徒對神的話沒有一個堅持.
你注定要被這個世界吞吃.
所以保羅在這裡提醒我們.
我們蒙昭是得自由的.
我們不需要像喇嘛教.
跪在那邊爬.
在那邊膜拜.
無底投地.
一天不知道幾千次.
在那個地方這樣子膜拜.
來表示他要去.
他的虔誠.
表示他要去贖他的罪.
我們也不需要像佛教徒.
我們要去做好事.
我們要捐錢.
所以拜登為我順便提一下弟兄姊妹.
不信主的人錢捐給教會可不可以.
可不可以.
可以啊.
多多益善.
什麼叫聖潔.
聖潔其中一個是叫做分別為聖.
聖的跟不聖的.
聖與俗不可以在一起.
我這是上次在長職會跟他們提醒這個原則.
如果不信的人奉獻的錢.
應當是要作為.
取之於社會用之於社會.
這是可以的.
了解意思嗎.
譬如今天可能有一個大規模的一個風災或怎麼樣.
那麼教會也成為一個來幫助這些的.
那這些錢也是去幫助這個社會的.
這一個是可以的.
但是如果教會的施工.

$^{481}$不要接受不信主的人奉獻.
教會如果不能夠有這種信心.
相信神會供應我們一切.
教會至少遲早會被世俗化.
所以保羅在這裡提醒我們.
我們蒙昭是自由的.
可是不要讓自由成為我們放縱的機會.
接下來他就提到了.
我是把他負面的消極的.
我先講.
總要用愛心互相服侍.
因為權力法都包在愛人如己這句話了.
你們要謹慎若相咬相吞.
只怕要彼此欺騙.
他就講到說如果你們濫用自由的一個結果是什麼呢.
他說就會相咬相吞.
相咬相吞 彼此消滅.
相咬相吞.
好 同志們.
教會如果.
好 先等一下.
你先看看這個圖.
你有沒有注意這個圖.
這是不同的動物在彼此爭鬥嗎.
有沒有注意.
看不清楚.
它是同一群的.
這是獅子還是豹.
獅子啊.
好 貓啊.
好 不管它.
反正它是同一群動物.
所以這叫做相咬相吞.
保祿在這裡提醒.
你們都信了耶穌.
但是你今天因為有這些人來.
他告訴你說你們必須要受割禮才能得救.
好了 在你們中間就起了紛爭了.
有的人支持割禮派.
有人支持信心派.

$^{521}$教會裡面就分成兩派了.
一個是割禮派.
一個是信心派.
懂這個意思嗎.
教會就相咬相吞了.
當然今天可能我們沒有這個問題.
但是我們會有很多其他類似的東西.
會讓我們相咬相吞的 會不會.
會讓我們彼此殘殺的 會不會.
也許你不要講殘殺.
殘殺好像太嚴重了.
會不會彼此相爭 請告訴我.
所以我一再強調.
幾幅碩書告訴我們一件事情.
保羅在講到教會這個藍圖的時候.
他所賜的有先知 有使徒.
先知傳扶於牧師和教師.
為了成全聖徒 各盡其職.
然後怎麼樣呢.
此等到我們認識神的兒子.
滿有基督長成的聲量.
然後接著他一直講下去.
講到最後 然後他講了什麼.
直到我們在愛中建立自己.
那個建立自己就是建立教會.
在哪裡建立 在愛中.
今天我們在北美.
我們來自五湖四海.
雖然我們都會講中文 是不是.
但是呢 請我山東人跟上海人.
兩個人兩個思維不一樣.
不一樣 是不是.
何況我們不僅僅還有國語堂.
我們還有什麼堂.
還有粵語堂.
講廣東話呢.
不懂聽不懂講 是不是.
所以你看語言的不同.
想法的不同.
思維也不同 是不是.

$^{561}$再加上個英語堂.
所以真的是variety 是不是.
所以在這樣子的不同當中.
我們更是需要竭力保守.
聖靈所賜 合而為一的心.
容易不容易 不容易.
但是如果你不覺得它重要.
你很容易就會放棄 是不是.
今天你看見不僅僅是華人教會.
韓國教會一樣.
在北美的韓國教會很多有分裂的問題.
特別移民教會.
因為移民教會大家都來自不同的地方.
那到了本土又受到本土的影響.
而且要複雜加複雜 你懂意思嗎.
我記得有一次我到一個地方.
到一個國家 我不要講哪個國家.
我到一個國家去講道.
我才到了parking.
他們接我到那邊.
才到了parking.
結果接我的弟兄就跟我介紹.
這個是我們教會的長老.
那個長老聽到我是新來的講員.
寒暄一下 打打招呼就好了 是吧.
沒有 怪怪的.
他在那邊抓著我.
跟我在那邊足足講了大概10分鐘.
我沒有仔細算時間.
至少10分鐘 講什麼.
講他教會主任路司的怎麼樣怎麼樣不好.
請問這樣的人合適做長老嗎.
對不起 不是說你.
你懂意思嗎.
怎麼會這樣子呢.
我一個外人 我一個剛來的講員.
你跟我講這些東西幹嘛.
了解意思嗎.
可是求主憐憫我們.
我相信他一定會覺得有很多受不了這個主任路司的.

$^{601}$我是沒有機會去問這個主任路司.
可能這個主任路司也很受不了他.
對不對.
unfortunately.
這個主任路司在那個教會裡面很多年了.
unfortunately.
這個長老是他們教會唯一的長老.
為什麼.
保羅說你們如果這樣子相咬相吞.
就會怎麼樣.
彼此消滅了.
跟你旁邊說不要相咬相吞.
要相親相愛.
家庭也是如此.
根據統計.
在歐洲 日本.
他所做的一些相關的統計.
大概大約有三分之一的夫妻.
是分床睡或者分房睡.
那這個原因很多.
有人說我老公晚上會打呼.
我沒有辦法睡.
他那個打呼真的是太大聲了.
有人說我太太晚上睡覺的時候會磨牙.
我告訴你.
我大學的寢室裡面.
上面是打呼下面是磨牙.
怎麼辦.
我們那個年代說什麼.
抗議 跟學校抗議.
換室友.
沒得事.
怎麼辦.
忍耐 跟你旁邊說忍耐.
我到現在都還記得我大學同學.
那個打呼磨牙的樣子.
因著那樣子的經驗.
我能夠去忍受.
現在有什麼.
在倫敦有一對夫妻.

$^{641}$他們決定.
為了讓兩個人有比較好的睡眠品質.
所以他們就決定.
他們因為丈夫是夜貓子.
而且會打呼.
那麼妻子需要早起.
而且她睡眠比較淺.
所以他們就為了自己的睡眠的自由.
他們就把公寓裝潢成兩間主臥室.
什麼意思.
個人住各一間主臥室.
現在都準備.
聖經怎麼說我們來看看.
哥林多前書第七章第五節我們起來唸.
我請弟兄來唸這一節.
夫妻不可彼此窺父.
除非兩相情願.
暫時分房.
未要專心禱告方可.
以後仍要同房.
免得撒旦趁著你們情不自禁.
引誘你們.
請姐妹唸彼得前書三章七節.
來 預備 請.
好 這裡說.
第一個哥林多前書第七章第五節說.
他說夫妻不可彼此窺父.
這是指在性關係上面.
他說除非兩相情願.
暫時分房.
原因是什麼.
因為打呼.
沒有.
他說為什麼.
未要專心禱告.
哇 非常屬靈的理由.
如果你未要專心禱告.
保羅說可以.
但是是暫時.
你不會說我一年到頭都在專心禱告.

$^{681}$那你不應該結婚.
你說我最近在禁食14天禱告.
或者我最近在三天禁食禱告.
我要專心禱告.
我到禱告山去了.
暫時分房.
但是保羅說以後要怎麼樣.
人要.
Come on.
我講的都是神的話.
弟兄姊妹你要對神的話的回應.
是一個信心的有把握的.
我是在操練你.
你懂意思嗎.
如果你不是這種態度.
你讀聖經對你沒有幫助.
你讀的只是知識.
你要相信神.
神這樣說.
他說免得怎麼樣.
撒旦趁著你們情不自禁.
引誘你們.
所以今天很多的大豔家庭.
夫妻分兩地.
這是非常不好的事情.
在幕會的中間.
從過去到現在.
我們一直提醒這些弟兄姊妹.
沒有任何原因.
不要這樣子分開.
很多分分分分的結果.
就是永遠分了.
可是這個時代.
整個大環境.
好像壓迫著讓很多的家庭走向這條路.
即使同住在國內也分.
那現在剛才講的三分之一這個比例.
那更高了.
更高了.
而且理由更有道理了.

$^{721}$我們沒有分.
我們沒有離婚.
我們還是住在一起.
只不過他睡他的床.
我睡我的床.
他有一間.
我有一間.
我們房間很多.
有時候覺得房間少有一個好處.
夫妻吵架沒得跑.
房間多了.
好.
換一間睡.
但是親愛的弟兄姊妹.
如果是這樣子的心態.
你們的關係已經出現了紅燈.
而且當你這樣子分的時候.
你在給下一代做了一個不好的榜樣.
下一代其實他們最在乎的是.
爸爸媽媽你們是不是彼此相愛.
我永遠記得.
我去探訪高長老的太太的爸爸媽媽.
所以我知道緣績.
因為講緣績可能很多人不知道.
緣績在嗎.
我記得那時候.
緣績姓什麼.
李伯伯李媽媽.
我印象很深刻.
當然我們聽到她講他們的戀愛史等等這些.
但我印象最深刻的就是他們分享.
他們每天晚上.
他們都是握著手一起睡覺的.
不單單不分房.
不單單不分床.
還每天握著手一起.
從結婚初一直到最後一天.
連並床上面都是這樣子.
我就跟旁邊說你的婚姻要這樣子.
我在努力當中.

$^{761}$所以那是我一直所羨慕的.
當然我想我們不是要每一個人都是標準模式.
可是我們在這裡提醒的就是.
我覺得很喜歡這有一個家庭的一個機構.
他們有一個算是他們的一個標語.
我們起來念好不好.
婚前不同房.
婚後不分房.
他這一篇文章就是專門在講到這件事情.
我覺得蠻經典的.
婚前不同房.
婚後不分房.
這兩個都是今天這個社會非常普遍的現象.
求主幫助我們.
不僅僅是如此.
不僅僅是分房睡或者分居.
或者是這些問題.
如果我們今天每個人都是想要捍衛自己的權利.
自己的自由.
你會發現問題都從這裡開始.
是不是.
只要我喜歡有什麼不可以.
為什麼你要管我.
我花錢.
你幹嘛管我.
夫妻之間常常吵架的一個因素是什麼.
金錢.
所以呢.
那既然如此.
結婚之前我們就講好了.
你的錢歸你的.
我的錢歸我的.
那家裡面的共同花費.
我們每個月一個人放多少錢在裡面.
那個是共同花費.
你幹嘛要結婚.
為什麼要結婚.
問你旁邊為什麼結婚.
如果你要自由.
你最笨的一件事情就是你結了婚.

$^{801}$你懂我意思嗎.
婚姻當中.
如果你沒有看見上帝設計婚姻的那個榮耀的目標.
藉著基督與教會的關係.
藉著妻子與丈夫的關係.
做丈夫的要為妻子捨己.
做妻子的要順服丈夫.
那都是要捨己的.
有姐妹說.
哎呀.
穆斯林不要老是講順服.
你講講弟兄要愛吧.
你以為愛比較容易.
還是順服比較容易.
我跟你講.
這兩個都有個共同點.
都是要怎麼樣.
捨己.
所以這兩個都一樣的.
親愛的弟兄姊妹.
只有當我們願意捨己.
這也是主呼召我們的.
主呼召我們.
剛剛講蒙昭德之友.
是不是.
但是主耶穌呼召我們的時候怎麼說.
祂說若有人跟從我.
就當怎麼樣.
捨己天天背起他的十字架.
來跟從我.
捨己怎麼捨呢.
背十字架.
十字架你喜歡背嗎.
沒有人喜歡背的.
十字架代表的是一個痛苦.
刑法.
但是親愛的弟兄姊妹.
十字架卻成為了榮耀.
為什麼.
因為耶穌基督為我們捨命.

$^{841}$所以第一件事情.
保羅提醒我們說.
我們不要彼此紛爭.
不要彼此爭.
不要彼此為了自我的權利.
在那邊爭這個爭那.
我做多少你做多少.
你洗了多少碗我洗多少.
乾脆好了.
每天你炒一半飯我炒一半飯.
是不是.
如果夫妻之間分到這個地步.
那沒有什麼.
這個婚姻已經不存在價值了.
是嗎.
每一個人都應當說.
我怎麼樣可以多為這個家.
多為你付出一點.
這就是和服神的心意了.
最後.
我們一起來讀這些聖經好不好.
來.
總要用愛心互相服侍.
因為神都包在愛人如己.
這一句話此類了.
這句話裡面說.
總要用愛心互相服侍.
首先我們來看看.
他這裡講的互相服侍.
原文的意思就是.
成為彼此的奴隸.
所以這個服侍是奴隸的意思.
什麼是奴隸.
你要了解在羅馬時代.
奴隸做什麼.
奴隸是一天二十四小時.
沒有自己的權利.
完全聽憑主人的.
主人要他做什麼.
他就做什麼的.

$^{881}$甚至有不好的主人.
打這個奴隸或者怎麼樣子的.
當然我們不是在講那個.
但是保羅也好.
耶穌也好.
都用這一個比例來告訴我們一件事情.
今天我們都是彼此怎麼樣.
我們做奴隸的.
我是高長老的奴隸.
Amen.
高長老是我的奴隸.
如果我們夫妻之間.
是帶著這樣子的心.
我是你的奴隸.
Yes.
老公.
有什麼需要的請吩咐.
你大概只能夠找到.
現在只能找到AI這樣子.
那現在都是奴隸.
可是聖經這樣教導我們的.
我們彼此之間的關係.
應該是一個服侍對方.
而不是想著說.
我要從你那邊得著什麼.
而是說主啊.
你幫助我.
我要怎麼樣子可以更好地來服侍.
服侍我的先生.
服侍我的太太.
我做她的奴仆.
我完完全全地聽她所說的.
那問題來了.
所以才會發明這些.
這些不合聖經的一些說法出來了.
我們家沒有.
我們家什麼.
我們家大事老公決定.
小事我決定.
但是大事小事呢.

$^{921}$不是不是.
大事小事我決定.
這些有沒有都聽過吧.
聖經怎麼說.
聖經有沒有說大事老公決定.
小事老婆決定.
有嗎.
沒有.
聖經怎麼說.
丈夫是妻子的頭所以呢.
但是實際上面.
今天很多的基督徒家庭.
可能丈夫不是頭.
是腳.
妻子很多時候變成是家中的頭了.
是不是.
好那親愛的女子們.
這些都是很實際的.
所以當我們說我們成為彼此的服侍.
馬丁路德他在解釋哥林多前書第九章十九節.
我們來看看.
哥林多前書九章十九節.
來我們起來讀.
我雖是自由的無人挾管.
然而我甘心做了眾人的僕人.
未要多的人.
保羅一方面說我是自由的.
沒人挾管.
哇.
這多好.
哇.
前面這句我太喜歡了.
但是他說我甘心做了眾人的僕人.
未要什麼.
多的人.
所以馬丁路德在解釋.
在他的這本書基督徒的自由裡面.
他解釋這句話.
他用了兩個命題.
第一個命題說.

$^{961}$他說基督徒是全然自由的眾人之主.
不受任何人管轄.
第二他說基督徒又是最敬重的僕人.
為眾人之僕.
受每一個人管轄.
哇.
這是一個自由的弔詭.
這兩個看起來是相互矛盾的.
但是保羅表明一件事情.
真正的自由是什麼.
真正的自由.
是我從我的內心的裡面.
我因著基督.
我因著主的緣故.
我是一個為主活的人.
我是基督的僕人.
所以我甘願做眾人的僕人來服侍人.
這是真正的自由.
接著最後 保羅說.
剛才講說.
剛才那句話我們再來看一下.
他說總要用愛心互相服侍.
因為什麼.
全律法都包在愛人如己這一句話之內.
愛人如己.
跟你旁邊說 愛人如己.
這句話是從哪裡來的呢.
在舊約裡面 利未記裡面這樣說.
他不可報仇 也不可向本國人懷恨.
卻要愛人如己.
如果你看了這節聖經.
你就知道今天以色列所做的.
是不是符合聖經.
不要說新約 因為他們沒有新約.
舊約都這樣說.
不可報仇 也不可向本國人懷恨.
卻要愛人如己.
所以你看了沒有.
今天以色列這些的執政長權者.
他只看到以牙還牙 以眼還眼.

$^{1001}$他沒有看到這樣進門.
他聖經沒有讀熟.
新約的時候 耶穌講更清楚了.
我們一起來讀二十二章這兩節.
耶穌對他說.
你要盡心盡性 盡一愛主你的神.
這是誡命中第一 且是最大的.
其次也要相仿 就是愛人如己.
這兩條誡命是律法和先知.
一切道理的總綱.
聖奧古斯丁在解釋.
解釋到這一個愛的時候.
他講到.
基督徒的自由是愛神.
然後要做你愛做的事.
這句話你一定要仔細的了解.
否則你很容易誤解.
其實他原文裡面是沒有.
沒有God這個字的.
他原文裡面 他的拉丁文裡面.
原文就是Love and do what you want.
但是他後面比較完整的.
把這句話做了一個解釋.
我念給大家聽.
他說愛 然後做你愛做的事.
如果你保持沉默.
出於愛而沉默.
如果你大聲呼喊.
出於愛而呼喊.
如果你糾正別人.
出於愛而糾正.
如果你寬恕別人.
出於愛而寬恕.
只要內心深處紮根著愛的根基.
這根基就絕不會長出.
良善以外的果子了.
奧古斯丁要表達的意思就是說.
你一切所做的.
你都是根基於這個愛.
這個愛是從神來的.

$^{1041}$所以他在解釋這個的時候.
就是從剛才我們說的.
你是一個愛神的人.
你一定也會愛人.
一個人他說他愛神.
他卻不愛人.
他是說謊.
約翰告訴我們的.
如果我們的心是愛神的.
我們一定會去愛人.
如己的.
所以從哪裡看見你是愛神的呢.
一個很重要的標記就是.
你是不是愛人如己.
這是一體的兩面.
夫妻關係如此.
弟兄姊妹之間的關係.
同宮之間的關係.
甚至你跟你靈社的關係.
我們跟家人 跟父母親.
跟小孩子的關係都是如此.
雖然我可能.
我都是舉別人的例子.
但是我舉一個我們自己教會的例子.
我不是要去凸顯.
也不是要去宣揚哪一個人.
但是我們大家都看得見.
Wanda照顧她自己的先生 George.
多少年了.
從進來我們教會到現在我看到.
九年的時間.
我不知道你有沒有機會看過.
Wanda抱怨.
覺得真倒楣.
怎麼會嫁錯老公呢.
今天這個世界不要等到像George這樣情況.
稍微有點怎麼樣子.
你自己身體要顧好.
將來沒有人會照顧你的.
這真是可悲.

$^{1081}$今天的世界太現實了.
所以我提他們的例子.
是因為我們教會有很多.
我看到這些長輩們.
我上次提過左波波.
我想找老師.
登泉.
我們最近看到.
你看見當真是在患難中的時候.
需要另外一半的時候.
他們是怎麼樣子的來服侍.
這就像牧師在主持婚禮的時候.
還記得嗎.
怎麼說.
牧師都會要另外一位新人跟著講一句話.
他說怎麼樣.
不論是健康或者疾病.
富有或是貧窮.
別人都可以離棄他.
但是你呢總不可離棄他.
你說我不是在家常結婚.
我沒有說過這個誓言.
沒關係.
但是你信了耶穌.
主要我們所走的路就是這樣子的路.
求主幫助我們.
讓我們能夠活出主的愛.
在我們的家中.
在我們的教會.
在我們這個社會活出主宰.
\newpage



\section{撒母耳記下 2:1-32}
\label{sec:032456}
\textbf{毋讓玉帛變干戈}~~~~李思聰牧師
\newline
\newline
連結: \href{http://www.church.com.hk/acms/content.asp?site=cdc\&op=show\&type=product\&code=032456\&layout=sermon}{\texttt{032456}} ~~~~ 語音日期: 2026-07-12
\newline
\newline
\hyperref[sec:032455]{\small{< < < PREV SERMON < < <}}
~
\hyperref[sec:index_chronic]{\small{[返順時目]}}
~
\hyperref[sec:032457]{\small{> > > NEXT SERMON > > >}}
\newline
\newline
撒母耳記下 2:1-32
\newline
\begin{longtable}{cl}
\hline
\hline
章節 & 經文 (和合本修訂版)\\
\hline
2:1 & \begin{tabularx}{0.7\textwidth}{X} 此後，大衛求問耶和華說：「我可以上猶大的一個城去嗎？」耶和華對他說：「可以上去。」大衛說：「我上哪一個城去呢？」耶和華說：「希伯崙。」 \end{tabularx} \\ \\ \relax
2:2 & \begin{tabularx}{0.7\textwidth}{X} 於是大衛和他的兩個妻子，一個是耶斯列人亞希暖，一個是作過迦密人拿八妻子的亞比該，都上那裡去了。 \end{tabularx} \\ \\ \relax
2:3 & \begin{tabularx}{0.7\textwidth}{X} 大衛也把跟隨他的人和他們各人的眷屬一同帶上去，住在希伯崙的城鎮中。 \end{tabularx} \\ \\ \relax
2:4 & \begin{tabularx}{0.7\textwidth}{X} 猶大人來，在那裡膏大衛作猶大家的王。有人告訴大衛說：「埋葬掃羅的是基列的雅比人。」 \end{tabularx} \\ \\ \relax
2:5 & \begin{tabularx}{0.7\textwidth}{X} 大衛就派使者到基列的雅比人那裡，對他們說：「願耶和華賜福給你們！因為你們忠心對待你們的主掃羅，埋葬了他。 \end{tabularx} \\ \\ \relax
2:6 & \begin{tabularx}{0.7\textwidth}{X} 你們既做了這事，願耶和華以慈愛和信實待你們，我也要為此厚待你們。 \end{tabularx} \\ \\ \relax
2:7 & \begin{tabularx}{0.7\textwidth}{X} 現在，你們的主掃羅死了，猶大家也已經膏我作他們的王，你們的手要堅強，要作英勇的人。」 \end{tabularx} \\ \\ \relax
2:8 & \begin{tabularx}{0.7\textwidth}{X} 掃羅軍隊的元帥，尼珥的兒子押尼珥，曾將掃羅的兒子伊施‧波設帶過河，到瑪哈念， \end{tabularx} \\ \\ \relax
2:9 & \begin{tabularx}{0.7\textwidth}{X} 立他作王，治理基列、亞書利、耶斯列、以法蓮、便雅憫和以色列眾人。 \end{tabularx} \\ \\ \relax
2:10 & \begin{tabularx}{0.7\textwidth}{X} 掃羅的兒子伊施‧波設登基的時候年四十歲，作以色列王二年，但是猶大家卻隨從大衛。 \end{tabularx} \\ \\ \relax
2:11 & \begin{tabularx}{0.7\textwidth}{X} 大衛在希伯崙作猶大家的王，共七年六個月。 \end{tabularx} \\ \\ \relax
2:12 & \begin{tabularx}{0.7\textwidth}{X} 尼珥的兒子押尼珥和掃羅的兒子伊施‧波設的僕人從瑪哈念出來，往基遍去。 \end{tabularx} \\ \\ \relax
2:13 & \begin{tabularx}{0.7\textwidth}{X} 洗魯雅的兒子約押和大衛的僕人也出來，在基遍池旁與他們相遇；一隊坐在池的這邊，一隊坐在池的那邊。 \end{tabularx} \\ \\ \relax
2:14 & \begin{tabularx}{0.7\textwidth}{X} 押尼珥對約押說：「讓年輕人起來，在我們面前較量一下吧！」約押說：「讓他們起來吧。」 \end{tabularx} \\ \\ \relax
2:15 & \begin{tabularx}{0.7\textwidth}{X} 他們就起來，點了人數過來：屬掃羅兒子伊施‧波設的有便雅憫人十二名，大衛的僕人也有十二名。 \end{tabularx} \\ \\ \relax
2:16 & \begin{tabularx}{0.7\textwidth}{X} 每人抓住對方的頭，用刀刺對方的肋旁，一同仆倒。所以，那地叫做希利甲‧哈素林，就在基遍。 \end{tabularx} \\ \\ \relax
2:17 & \begin{tabularx}{0.7\textwidth}{X} 那日戰況激烈，押尼珥和以色列人敗在大衛的僕人面前。 \end{tabularx} \\ \\ \relax
2:18 & \begin{tabularx}{0.7\textwidth}{X} 在那裡有洗魯雅的三個兒子：約押、亞比篩、亞撒黑。亞撒黑的腳快如野地裡的羚羊； \end{tabularx} \\ \\ \relax
2:19 & \begin{tabularx}{0.7\textwidth}{X} 亞撒黑追趕押尼珥，直追趕他不偏左右。 \end{tabularx} \\ \\ \relax
2:20 & \begin{tabularx}{0.7\textwidth}{X} 押尼珥回頭說：「亞撒黑，是你嗎？」他說：「是我。」 \end{tabularx} \\ \\ \relax
2:21 & \begin{tabularx}{0.7\textwidth}{X} 押尼珥對他說：「你轉左或轉右，去抓一個年輕人，剝去他的戰衣吧。」亞撒黑卻不肯轉開而不追趕他。 \end{tabularx} \\ \\ \relax
2:22 & \begin{tabularx}{0.7\textwidth}{X} 押尼珥又對亞撒黑說：「轉開，不要再追我了！我何必把你擊殺在地上呢？我若殺了你，怎麼有臉見你哥哥約押呢？」 \end{tabularx} \\ \\ \relax
2:23 & \begin{tabularx}{0.7\textwidth}{X} 亞撒黑仍不肯轉開，押尼珥就用回馬槍刺入他的肚腹，甚至槍從背後穿出，亞撒黑就仆倒在那裡，當場死了。眾人趕到亞撒黑仆倒而死的地方，就都站住。 \end{tabularx} \\ \\ \relax
2:24 & \begin{tabularx}{0.7\textwidth}{X} 約押和亞比篩追趕押尼珥。日落的時候，他們到了通往基遍曠野的路旁，基亞對面的亞瑪山。 \end{tabularx} \\ \\ \relax
2:25 & \begin{tabularx}{0.7\textwidth}{X} 便雅憫人聚集在押尼珥後面，成為一隊，站在一座山頂上。 \end{tabularx} \\ \\ \relax
2:26 & \begin{tabularx}{0.7\textwidth}{X} 押尼珥呼叫約押說：「刀劍豈可永遠吞噬呢？你豈不知，結局必是痛苦的嗎？你要等到何時才叫百姓回去，不追趕他們的弟兄呢？」 \end{tabularx} \\ \\ \relax
2:27 & \begin{tabularx}{0.7\textwidth}{X} 約押說：「我指著永生的神起誓：你若沒有這麼說，百姓就必繼續追趕弟兄，直到早晨。」 \end{tabularx} \\ \\ \relax
2:28 & \begin{tabularx}{0.7\textwidth}{X} 於是約押吹角，眾百姓就站住，不再追趕以色列人，也不再打仗了。 \end{tabularx} \\ \\ \relax
2:29 & \begin{tabularx}{0.7\textwidth}{X} 押尼珥和他的人整夜行過亞拉巴。他們過了約旦河，走過畢倫，到了瑪哈念。 \end{tabularx} \\ \\ \relax
2:30 & \begin{tabularx}{0.7\textwidth}{X} 約押追趕押尼珥回來，聚集眾百姓，大衛的僕人中缺少了十九個人和亞撒黑。 \end{tabularx} \\ \\ \relax
2:31 & \begin{tabularx}{0.7\textwidth}{X} 但大衛的僕人殺了押尼珥的人，便雅憫人三百六十名。 \end{tabularx} \\ \\ \relax
2:32 & \begin{tabularx}{0.7\textwidth}{X} 他們把亞撒黑送到伯利恆，葬在他父親的墳墓裡。約押和他的人走了一整夜，天亮的時候他們才到希伯崙。 \end{tabularx} \\ \\
[1ex]
\hline
\hline
\end{longtable}
$^{1}$上周跟大家講的是撒米爾記下的第一章.
講到關乎大衛知道蘇羅已經死了.
但是發掘了一個方畫.
顯出了他的生命的正直.
特別是來表揚這種尊榮神的品格.
今天來到第二章 撒米爾記第二章記載了一件事蹟.
這件事蹟也很值得我們思想學習領受.
今天我給這個題目是 大家認不認識.
無釀玉白成干戈 知不知道是什麼意思.
有沒有聽過一句話叫做化干戈為玉白.
是什麼意思 化干戈為玉白.
什麼叫做玉白.
玉白其實就是我們的保育.
很多時候都是代表兩國之間.
大家和平友好 彼此送禮.
以保密為徵 結盟 立和約.
所以玉白通常表明和平 盟約 大家和好 善意.
干戈呢 干戈就是打仗.
所以在中國成語 化干戈為玉白是什麼意思.
本來大家在爭執 鬥爭.
但是大事化小 小事化無.
大家能夠和平相處.
所以化干戈為玉白 很好.
本來打仗 現在成為大家彼此和好 結盟.
但是在這張聖經裡面 我們看見就糟糕了.
所以給我們一個反面的學習.
他們本來和好 但是後來就變成爭戰.
所以讓我們得到尷尬.
所以這個題目叫做「毋讓玉白反而成為干戈」.
這幅圖是有意思的.
這幅圖是考古的地圖.
大家看到中文字有解釋.
就是雞片的一個詞者.
雞片我們讀舊約聖經都是很熟悉的.
特別約書亞記出現的時候.
就是記載這些雞片人來欺騙約書亞.
因為他們害怕伊斯蘭進去不單止滅.
其他迦南人也會滅了他們.
所以這個雞片其實是一個迦南的古代的大城.
如果大家記得約書亞記記載的艾城.

$^{41}$它是很近艾城的.
在西北大約九公里之近.
但是雞片出名的地方附近有一些橄欖和地下的泉水.
其實雞片這個字的意思是屬於山地.
山地不過能夠有泉水.
所以是一個很特別的地方.
有一個很優美同時也有水源的地方.
但是這個考古發現.
他們大概在鐵器早期時代發掘出來的.
就是有一個大的坑.
裡面有11公尺深.
但是你看到有樓梯.
樓梯下去下面有一條水道.
這條水道有12尺通道這麼長.
是一個儲水庫.
所以他們將這些水池儲在裡面.
考古後來發現這個就是在雞片地方的一個池.
一個井.
他們挖出來的一個井.
是有水可以儲水的.
所以這個叫做雞片池.
就是我們稍後會讀到的地方.
所以這個雞片池發生了什麼事呢.
今天我們一起來看看.
這個題目叫做.
「化乾戈為玉帛」.
就是不要讓玉帛成為乾戈.
其實也是一個很好的意思.
讓我們看到本來他們是很美好的.
不過就是因為這些爭執.
搞成戰爭.
搞成人命的傷亡.
好像形容我們今天來說.
我們本來很好的.
大家的關係.
很朋友的.
但是因為這些爭執小事.
後來就越來越有爭執.
變成傷害.
變成仇恨.

$^{81}$變成血腥.
所以這個就是很難過的事.
無論是我們的家庭也好.
我們朋友之間的關係也好.
我們工作的地方也好.
本來是很和好的地方.
教會更加如此.
教會應該不用說了.
應該是一個和平和好的地方.
但是人性也會有鬥爭爭執或者摩擦.
不要讓玉帛成為什麼.
不要成為乾戈.
這就是我們今天從這段經文裡面的很好的提醒.
先看看背景.
就是從第一節開始.
這裡聖經不讀也不行.
因為是一個很重要的背景.
所以聽我讀.
我讀完十一節給大家聽.
中間有些解釋.
《撒末爾記》第二章第一節.
「之後大衛問耶和華說」.
所以這裡我們看到.
大衛這個靈性美好的開始.
「大衛問耶和華說.
『我想由大的一個城去,可以嗎?』.
耶和華就回答他.
『耶和華說,可以』.
大衛說,『我想哪一個城去呢?』.
耶和華說,『想希伯倫去』」.
你看到這裡.
大衛與神對話.
就好像我們今天祈禱一樣.
聖靈跟我們說話.
我們向神禱告.
所以這是一個很好的親密的關係.
這真是很美好的事.
大衛很值得欣賞的就是這些地方.
第二節.
「於是大衛和他兩個妻.

$^{121}$一個是耶斯列人.
阿希倫一個是作為加密人.
那八妻的阿比戈.
都上那裡去了」.
「大衛也將跟隨他的人.
和他們各人的眷屬一同帶上去.
住在希伯倫的城入中」.
所以你看到很多人一起.
去了希伯倫.
神指示他的地方.
所以我們看到.
第四節開始.
「猶大人就來到希伯倫了.
在那裡告大衛作猶大家的王」.
有人告訴大衛說.
「作埋葬素羅的基列雅比人」.
我們看到大衛從這裡開始作王.
所以他先做猶大家的王.
後來做了七年之後.
其他的耶斯列家.
其他的族人.
就一起再擁立他作為全國的王.
但他現在這個階段.
就是先做他們自己的支派.
大衛屬於哪個支派?.
猶大.
所以他是猶大支派.
所以先作猶大支派的猶大家的王.
但是有人告訴他.
剛剛上一章記載.
素羅死了.
葬在哪裡?.
葬在基列雅比人那裡.
基列雅比人是什麼人?.
其實就是素羅自己的原族.
所以他葬在自己的故鄉.
自己的原族那裡.
就是基列雅比.
第五節.
大衛差人去見基列雅比人.

$^{161}$對他們說.
你們厚待你們的主素羅.
將他埋葬.
願耶和華賜福於你們.
大家記得上一次素羅死的時候慘死.
還被人割了手甲.
掛在了愛邦之地.
結果這些同族的人.
很有憐憫之心.
因為大家是同族.
所以就將素羅的屍首帶回來.
將他安葬.
所以大衛就稱讚他們.
你厚待你們的主素羅將他埋葬.
第六節.
你們既行了這事.
願耶和華以恩慈誠實待你們.
我也要為此厚待你們.
現在你們的主素羅死了.
尤大家已經告我作他們的王.
所以你們要剛強奮勇.
所以這是很好的事.
大家都是同是以色列人.
雖然不同族.
現在我已經被尤大家封立為王.
不過不要緊.
你們有你們的生活.
我們互相尊重.
我也不要你們一定歸屬王.
你們就好好剛強奮勇.
好不好.
是不是好像大家一種很好的關係.
彼此互勉.
彼此祝福.
他說願耶和華以慈愛誠實待你們.
願耶和華賜福給你們.
好不好.
真是很好.
所以能夠對人.
對我們的弟兄姊妹彼此祝福.

$^{201}$以恩慈誠實彼此相待.
真是很美好的事.
而且他鼓勵他們.
你們應當剛強奮勇.
所以彼此鼓勵彼此勉勵.
這些真是很美好的事.
第八節開始有些不同.
我們看見第八節是素羅的元帥.
即是素羅家那邊.
尼爾的兒子叫甚麼名.
叫做阿尼爾.
他曾將素羅的兒子伊斯波切帶過河.
到馬哈廉納他作王.
智利,基列,阿蘇利,耶穌列,.
以法廉變阿敏.
這些全部都是屬於後來北角的地方.
即是素羅家那邊的地方.
以及以色列的眾人.
不是猶太家那邊.
所以素羅有他們自己的主將.
素羅有自己的兒子.
就是伊斯波切.
所以第十節.
素羅的兒子伊斯波切登基的時候四十歲.
作以色列王二年.
唯獨猶太家歸從大衛.
大家很清楚了.
所以猶太家歸誰?.
大衛.
其他那些以色列人就歸.
伊斯波切.
伊斯波切作王兩年.
所以這裡就分開了.
猶太家和伊斯波切.
即是素羅家.
大衛在希伯倫作猶太家的王.
總共多久?.
共七年零六個月.
所以這個就是他在希伯倫.
做猶太家的王.

$^{241}$開始做王了.
總共有七年多.
好了.
接著十二節下去.
我們用啟英方式讀.
我讀一節.
請弟兄姊妹讀一節.
我們一起看看這個法字.
十二節我先讀.
尼爾的兒子.
阿尼爾和素羅的兒子.
伊斯波切的僕人.
從哈馬哈廉出來.
往基片去.
即是去猶太家那邊.
十三節請讀.
(十三節).
所以這個就是基片馳立.
所以他們往基片去.
就是剛才我們看考古的地圖.
那幅圖的地方.
這個叫基片馳.
(十三節).
可以啊.
這個就開始棄灑.
怎樣棄灑呢?.
十五節請你們讀.
(十五節).
十六節.
彼此抽頭用刀刺勒.
一同撲倒.
即是怎樣?.
即是死了.
玩到死.
所以那地叫做希利夾哈蘇林.
就在基片這個地方.
十七節請你們讀.
(十七節).
所以本來十二個出來.
大家玩下玩下.

$^{281}$越玩越大問題.
接著就死了.
接著就變成戰士兇猛.
慘了.
見過誰打勝仗?.
大衛博人勝仗.
所以我們見到這個就開始.
阿尼爾那些人走了.
十八節請聽我讀.
(十八節).
十九節請讀.
(十九節).
阿尼爾正在走.
阿尼爾是誰?.
阿尼爾是塑羅家的主將.
他現在打輸了.
走了.
走放過他嗎?.
不放過.
誰追他?.
阿薩克追他.
所以這個阿薩克就是約翰的弟弟.
約翰就是大衛的元帥.
所以我們見到阿薩克就追他.
二十節請聽我讀.
阿尼爾回頭說.
追緊的時候回頭說.
你是阿薩克嗎?.
回答說.
是.
請讀.
(十八節).
所以阿尼爾叫他不要追我.
隨便找一個做替死鬼.
剝了他的戰衣就算了.
我們就此結束.
二十二節.
阿尼爾又對阿薩克說.
你轉開不追趕我爬.
我何必殺你呢?.

$^{321}$若殺你有什麼念見你哥哥若蘭呢?.
請讀.
(十八節).
所以後來阿尼爾是高手.
是以色列的元帥.
是塑羅家的.
而一槍一個槍轉.
就刺入阿薩克的肚腹.
就死了.
二十四節.
約阿和阿比西.
就是他們兩個哥哥.
追趕阿尼爾.
日落的時候到了.
通基片曠野的路旁.
基亞對面的阿瑪山.
所以這兩兄弟見到.
他的弟弟被阿尼爾殺了.
繼續追趕他.
由此你報仇.
二十五節請讀.
(十八節).
變阿曼人.
就是塑羅家的人.
塑羅家的人聚集.
跟隨阿尼爾站在一個山頂上.
阿尼爾呼叫約阿薩克山頂.
和對面的大衛加約阿這些人說.
刀劍豈可永遠殺人.
你豈不知終究必有苦楚.
你要等何時才叫百姓回去.
不追趕弟兄呢.
二十七節請讀.
(十八節).
他們本來是稱兄道弟的.
本來同為一家的.
真是好可惜.
二十八節.
(二十八節).
於是約阿摧國.

$^{361}$中文就暫住.
不再追趕以色列人.
也不再打仗了.
請讀.
(十八節).
約阿追趕阿尼爾回來.
聚集民眾.
見大衛的僕人.
終缺少了十九個人.
和阿薩克.
(十八節).
最後一次一起讀.
請讀.
(十八節).
所以叫各回自己的地方.
一個回到馬哈廉.
一個回到希伯倫.
我們見到阿薩克.
他本來是猶太家的人.
送回他故鄉的伯利恆安葬.
事情就到這裡.
後來如果大家回去讀.
還有的.
大衛知道的時候就為他們悲哀.
所以大衛就很傷心.
見到阿薩克被刺死.
所以見到死傷很多.
三十節和三十一節.
所以我們見到這件事.
因此這個地方.
我們看見十六節就這樣說.
(十六節).
我曾經和大家說過.
在三維記裡面有幾個地方.
很特別的.
不是普通的名字.
他們把音譯譯出來.
而且很長.
其中一個就是.
希利夾哈蘇林.

$^{401}$所以他的希伯倫文就是.
Kathasrem.
這個地方因此改名叫做.
希利夾哈蘇林.
本來叫做基片祠.
基片祠我們再說.
是一個山丘.
是一個有水有池的地方.
非常優美.
但是因為這件事就改成這個名.
這個名是甚麼意思.
這個名的意思就是.
刀劍之地.
優美小山處水池而變成刀劍之地.
這個就是玉帛成為甚麼.
成為干戈.
真是玉帛成為干戈.
不過求主幫助我們.
在這個故事當中我們得到尷尬.
不要讓玉帛成為干戈.
好了 刀劍之地.
由於甚麼出現.
請問為何會變成刀劍之地.
本來是一個和諧的地方.
一個和諧的家.
一個和諧的團體.
為何變成相咬相吞.
刀劍之地.
第一 我們既然刀劍之地出於甚麼.
出於氣喪.
本來只是玩玩而已.
結果弄出人命.
禍愈闖愈大.
在這裡給我們真是好提醒.
氣喪的問題.
第一 我們見到這個氣喪的行為.
第十四節這樣說.
阿尼爾對玉帛說.
讓少年人起來.
在我們面前氣喪吧.

$^{441}$玉帛說 好啊.
那就開始來了.
大家都是懂武功的.
大家鬥一下武功.
可能先出拳 後來出腳 後來出刀.
越打越盛.
所以這個是氣喪的行為.
我們中國人說得好.
不知道大家有沒有聽過.
勤就有功.
氣無益.
當然不是說玩遊戲.
誰知道乒乒那些遊戲.
打乒乓球 打籃球.
氣無益即是說.
彼此的氣灑.
玩著玩著就玩出火.
這種氣灑是很危險的.
所以這種通常都是出於好勝.
一定要贏.
還有放縱.
這些行為在氣灑裡面成為危險.
有沒有 今天這個世界多不多.
多不多這些.
本來大家慶祝.
突然間越整越搞.
又說玩一下遊戲.
你們喝一杯 我喝一杯.
越整越 接著出拳 出刀.
接著出什麼 出槍.
美國就是這樣.
你們現在看不看世界盃.
不是每個人都看.
現在哪隊就哪隊.
很多時候這些本來是慶祝的.
哪隊贏了就大肆慶祝.
前陣子NBA就是這樣.
NBA知不知道是什麼.
NBA就是打籃球的.
這個東岸贏了冠軍.

$^{481}$結果第一場贏了已經暴動了.
本來慶祝.
全街就搞到什麼.
又放火 又推警察車.
又打爛玻璃 又搶東西.
贏了總冠軍不得了.
結果要全城警察出動.
很多時候就是這些問題.
玩開始 玩著玩著就放縱.
今年年頭不知道大家有沒有聽過.
在阿爾卑斯山發生一件慘劇.
今年年頭1月1日凌晨發生.
記得嗎.
他們年輕人在那裡慶祝過年.
結果因為樓頂設計不太好.
他們騎駁馬 香檳又是放火的.
香檳放火 誰不知一個不小心燙著.
燙著東西就燒起來.
這些就是我們明天開始的VBS的佈置.
為什麼這裡有這麼多蜘蛛網.
因為明天是兒童書記聖經班的佈置.
他們不是燒著了屋頂.
結果死了47個人.
就是因為這樣燒.
115人受傷.
所以差不多200個年輕人.
在滑雪聖地 瑞士慘遭罹難.
屍體也無法辨識 消毒.
所以很多人害怕.
正當的娛樂 輕鬆有益的運動活動是好的.
我們要鼓勵弟兄姊妹.
但如果氣灑太多 整天只顧著玩.
人們就這樣 打麻將 通宵打麻將.
結果就很大很大的傷害身體.
一開始只是玩玩 後來賭一下.
越整越大 開始只是口角之爭.
後來就出手出刀.
所以我們發覺真的要留意.
為什麼保羅經常勸我們要謹慎自守 謹慎禱告.
真的要留意.

$^{521}$所以今天我們教會回到頂尖會.
教導我們年輕人真的很重要.
所以從小開始回教會.
我們就開始從小教導他們.
當然不可以避免所有.
不過我們都知道這是很重要的.
從小開始培養和指導.
基督徒可不可以玩.
近代有一個叫休閒神學.
休閒神學就試探討究竟我們基督徒如何休閒.
有什麼聖經的教導.
其實有一個指標或考慮是不錯的.
如果大家打麻將打得很興奮.
我們郭猶團是很好的.
我們打球之前祈禱.
所以我們看看.
如果你能夠在當中讓神參與.
那就有些樂譜了.
如果我們發覺自己有自己的生活.
教會有一套 出去又另一套.
最近我和我女兒聊天.
很傷心 見到有些他們的同輩.
教會很事奉 很熱心.
出去派對的時候和普通人一樣.
他們覺得為何基督徒要將這些分教.
教會是一個價值觀.
去到別處有一種表現.
那就差不多了.
如果你在這裡看通宵聖靈感動你睡覺.
看劇好收檔.
真的有些問題.
大家明白嗎.
基督徒不是不可以休閒 不可以氣灑.
不過去到這些地方的時候就過界了.
除了氣灑的行為 氣灑的言語.
若說我指著永生的上帝氣勢.
也若不說氣灑的那句話.
今天早晨百姓就會去不追趕弟兄.
這個後來約自己說.
最壞都是你 說這些話.

$^{561}$你一說這句話就讓我們有爭執開始.
所以氣灑的說話.
由衝突引起的原因.
現在這個世代也一樣 不單是行為.
說話亂來.
喜歡說話.
說得你也知道你的感受如何.
他不會想一下再說.
有時候也會說笑而已.
說就說笑 不過很容易就過火了.
就過界了.
所以當然有些粗俗的言語.
不好的說話更加要避免.
甚至那些嬉笑的話.
即是用來取笑別人.
那些嬉笑的話.
挑逗別人的血氣.
挖苦你 貶低你.
這些都是很容易在這個時代會有的.
在《爾弗所書》提醒我們.
他說任詞 妄語 嬉笑的話都不相宜.
所以不要想到說就說.
然後說笑一下.
說笑和取笑一線之間.
你可以說一句話 笑一下自己.
笑一下環境 有些幽默感.
基督徒應該也是這樣的.
有些幽默感.
做基督徒不會很呆板.
不過不要變成取笑.
即是傷害別人 貶低別人.
取笑的話更加挑動血氣.
就會危險.
這很多時候發生在哪裡呢.
即是第一次見面會不會說這些話.
好笑 好笑.
所以多數發生在哪裡.
很熟悉的那些人.
熟悉要小心.
唉呀 此熟邁熟.

$^{601}$熟悉就不拘禮了.
因為不拘禮就容易得罪別人.
所以代表熟悉要更加小心.
誰最熟悉呢.
家人最熟悉.
為甚麼很多時候會在家裡發生.
因為你不留意家人最熟悉.
所以一下子就出來.
就收不回.
最慘是這樣 說了出來收不回.
所以真的代表我們要留意.
不要太喜歡責人.
不要太喜歡取笑人.
小心.
所以原來這些言語上引起這個問題.
所以我們第一看見刀劍之地出於甚麼.
出於氣灑.
留意我們要提防.
不要將欺騙詞.
本來是一個很好的彼此祝福的地方.
變成一個刀劍之地.
真的要小心.
氣灑的行為.
氣灑的言語.
讓我們謹慎我們的言語.
謹慎我們的行為.
我再說謹慎不是代表你.
做基督徒好像節日不少.
又不能說笑 又沒甚麼輕鬆.
基督徒不是這樣的 我們有喜樂.
但喜樂裡面有愛.
在愛裡面要小心.
在愛裡面用祝福的說話.
所以這個很重要.
求神幫助我們.
第二我們看見刀劍之地出於甚麼.
刀劍之地出於憎勁.
我們發覺這個就是最大背後的原因.
本來同出一家 雅各家.
大家都是雅各的子孫.

$^{641}$但是發覺以色列和猶太家.
大家就因為這些爭鬥.
利益,地盤,權力開始.
本來大衛的心胸很廣.
但下面的人不是.
所以他們發現下面的人開始鬥爭.
來呀來呀 看誰比較厲害.
憎勁.
中文字挺好的.
競爭和憎勁有些不同.
競爭可以壞也可以好.
良性的競爭.
比賽,遊戲競爭.
希望大家都進步.
良性的不會傷害別人.
也出於要改進自己.
不是要貶低別人.
這些競爭大家努力.
像現在打足球 落場當然要贏.
不過落場出矛招就不是了.
一出矛招 心地不好.
就變成憎勁.
憎勁就是血氣.
憎勁就是要抬高自己貶低別人.
要比別人厲害.
那種心態不是為了一種正義的取向.
得勝利的標準.
良性的益處.
不是為了這些.
是為了一些血氣的競爭.
一些嫉妒,紛爭.
一些好勝.
所以我們看到很明顯.
這次整件事大家讀了這麼多經節.
你會留意到.
這是出於一種憎勁.
是肉體的一種罪性.
不能輸 不能輸.
所以心裡面其實是一種驕傲.
所以這種驕傲就是背後的原因.

$^{681}$所以聖經裡面提醒我們.
特別在《箴言書》有很多關乎憎勁的.
記載的原因和一些教導.
今天我們溫習有幾樣東西.
在聖經的《箴言》裡面.
第一 我們發覺很多問題出於憎勁.
出於什麼?.
剛才我們說憎勁出於什麼?.
出於驕傲.
《箴言》第十三章第十節.
一起讀好嗎?請讀.
驕傲旨愧憎勁.
聽勸賢的卻有自慰.
這真是自慰.
所以驕傲就帶來憎勁.
因為他覺得自己最厲害.
自己很驕傲.
我賺多少錢.
我多厲害.
我有什麼擁有.
他開始驕傲的時候就和別人比較.
和別人憎勁.
發覺有個人比我厲害.
不行.
我一定要比他厲害.
開始來了.
這個就是憎勁開始.
所以驕傲旨愧憎勁.
但是這個卻帶來一個.
彼此爭鬥.
彼此破壞的結果.
所以聽勸賢就有自慰.
不要留意驕傲的心態.
這正是我們的肉體.
新玉就知道了.
這是出於肉體的.
需要對付的.
需要靠聖靈致死它.
《箴言》十七章第十九節.
同樣請讀好嗎?.

$^{721}$喜愛憎勁的是喜愛過分.
高立家門的乃自取敗壞.
什麼叫高立家門?.
就是覺得自己比別人厲害.
自視過高.
這就是驕傲.
自視過高.
自己比別人厲害.
義氣憎勁.
所以這些人就憎勁了.
憎勁其實就是喜愛過分.
專搞事搞憎勁.
搞鬥爭.
這些就是自取敗壞.
我們見到聖經提醒我們.
小心這些驕傲.
在夫婦之間.
在親友之間.
大家有一種驕傲.
彼此鬥爭的心.
就會產生不好的後果.
產生憎勁.
第二我們見到.
這個憎勁在《箴言》提醒我們.
就是放任.
在《箴言》十七章第十四節.
我們再一起讀好嗎?.
請讀.
「紛爭的起頭如水放開.
所以在爭鬧之先.
必當止息憎勁」.
所以原來紛爭的起頭如水放開.
大家有沒有開過水喉?.
愈開愈猛.
那怎樣才解決呢?.
他說在開頭的時候快點關掉.
這個我們一會兒再看.
這種是指什麼呢?.
指放任.
因為他開頭已經覺得有點不妥.

$^{761}$不過他沒有聽勸告.
也沒有留心,耐心的感覺.
更加基督徒沒有留意聖靈的感動.
他就放縱.
照自己的意思而行.
所以憎勁就愈來愈大.
他沒有留意.
這是屬於肉體的.
這是破壞他自己.
也破壞他和別人的關係.
還有我們看到這種逞強.
什麼叫逞強?.
就是逞英雄.
《箴言二十五章》第八節.
請讀.
「不要務失出去與人爭競.
免得至終被他羞辱.
你就不知道怎樣行了」.
拍拍胸口走出去.
就是逞英雄.
他說不要務失出去.
想過這是否需要的.
如果你挺身而出.
為正義而行的.
付代價的.
沒問題.
但是如果你務失出去.
與人爭競.
這樣至終會產生羞辱.
根本你都不是別人的那一批人.
更加是逞強,逞英雄.
這樣就自找麻煩.
所以這種逞英雄心態.
今天很多人都有.
接著還有留言.
《箴言二十六章》二十到二十一節.
我們一起讀.
「火缺了柴就必熄滅.
無人傳洩爭競便自息.
好爭競的人善或爭端.

$^{801}$就如如火加炭火上加柴一樣」.
所以這種是甚麼?.
撥火.
所以原來爭競.
很多時候出於甚麼?.
出於傳洩.
傳洩即是甚麼?.
又講人行嗎?.
本來火都熄滅了.
接著給你一撥,又熄滅.
如如火加炭火上加柴.
所以這些人叫好爭競.
好爭競當然喜歡這樣做.
很大問題.
喜歡生事.
無端生事.
搞到你一塌糊塗就開心.
這些出於非常壞的心態.
但不知不覺就做了這種事.
所以真是好爭競的人要留心.
如果你發現自己很喜歡.
靈異認叻,怕心口.
說得出就說.
好爭競的時候.
小心會變成自己都不知道.
成為這種煽風點火的人.
流言.
指於甚麼?.
流言指於自己.
所以爭競就是知識.
沒有人傳洩.
接著就爭辯.
在提摩太后書2章23節.
唯有那愚拙無學問的辯論.
總要棄拙.
因為知道者等事事起爭競地.
保羅提醒提摩太后留心教會裡面.
即是他提醒教會.
這個伊弗梭教會.
保羅心目讓的.

$^{841}$提摩太后心目讓的.
他說這些愚拙無學問的辯論.
總要棄拙.
因為這些事很多時候最後都是.
大家都咬到臉紅耳熱.
最後又怎樣?.
最後就產生爭競.
所以這些是.
教會少談政治.
不要說我是藍你是紅.
就爭辯一輪.
所以是無謂的.
沒有絕對.
更加不是真理.
為甚麼要這麼守住自己的立場.
要爭辯到差不多要出拳出腳.
又不要因為大家的小小的看法不同.
又說咬到臉紅耳熱.
還有這個其實講到有些愚拙無學問的辯論.
有些世上的小學.
甚至有些神學辯論.
其實那些無關痛癢的也不是絕對真理.
聖經裡面有支持這邊有支持那邊.
你覺得自己這邊對.
你覺得我分享我的看法我的領受是這樣.
不過也尊重其他人的看法.
除非他是異端.
除非他是違反基要真理.
所以這些也是無謂辯論.
所以這個也是會引起爭辯.
爭辯開始.
那怎樣避免?.
怎樣避免增進之道?.
這裡也告訴我.
避免增進之道.
剛才我讀過.
第一就是甚麼?.
蓄止.
紛爭的起頭如水開放.
所以爭鬧之先.

$^{881}$必當甚麼?.
必當止息增進.
避免血氣之爭.
避免小事化大.
一發覺開始有些火藥味.
那又怎樣?.
我們談到這裡就好了.
就這樣.
不要再爭辯下去.
無謂爭辯.
我們可以agree to be disagree.
所以無謂傷感情.
所以說話不要太大聲.
不要太權威的.
因為留點餘地給別人.
所以我認為會不會是這樣呢?.
這樣好一點.
不要太絕對.
所以避免大家無謂的爭執.
事實上很多事情是.
有些事情不需要太絕對.
你的領袖未必是對的.
所以除非很清楚有聖經根據.
這個權威不是來自你.
而是聖經的權威.
這樣就不同.
所以我們看到這是很重要的.
一種止息增進之道.
第二就是最好和解.
所以箴言提醒我們.
十八章十八字.
請讀.
祭籤能止息增進.
也能解散強性的人.
所以你說你要跟我鬥.
一起喝.
你喝一杯我喝一杯.
你說不要用這些方法解決.
我們不如猜情沉算了.
或者讓他.

$^{921}$或者找其他方法抽籤.
所以不需要用這些血氣爭鬥的事.
用和平解決最好.
和平解決不要用鬥爭力鬥解決.
通常力鬥都是出於血氣.
所以好爭勁好出頭.
我們要提防.
免得落在血氣鬥爭之上.
小爭執本來出來玩玩.
不過化成大.
小糾紛不平息.
就變成嚴重的結果.
死了多少人.
死了很多人.
按著定數所計.
贖蘇羅兒子伊斯波帖的辯鴉民過去12名.
大衛的僕人過去有12名.
大家一起12名出去.
死了多少個.
死了24個.
約瓦追趕阿里回來.
聚集群眾見大衛的僕人.
少了多少個.
大衛那邊少了19個.
加上阿薩克20個.
總共32個.
大衛加上有32個.
但大衛的僕人殺了辯鴉民.
跟隨阿里的人有多少個.
360加12.
372個.
真是不值.
無端端死了400人做什麼呢.
四傷三重.
我們再講本來同國同家.
大家親同手足.
稱兄道弟.
今年我們教會的主題是什麼.
叫全家事奉主.
所以全家不代表我們肉身家庭.

$^{961}$全部是兄姊妹神的家.
神的家最好一起做什麼.
一起事奉主就最好.
Amen.
不要搞到全家.
彼此鬥爭就慘了.
家庭如是 教會如是 團體如是.
不要讓肉百變成幹局.
最後我們看見.
順便一提一下.
今天我們受媒體影響很大.
大家知不知道.
你經常看YouTube.
經常看電影.
很多時候都宣導這些.
你知不知道.
很多時候宣導這些鬥爭.
很厲害.
打打殺殺.
對人很mean.
有一套電影叫做Mean Girls.
你們有沒有聽過.
Mean怎麼翻譯.
很難翻譯.
你對我很mean.
很壞.
所以今天很多這些意識.
弟兄姊妹你們在教會裡.
受教導聖經.
每天靈修十分鐘.
每個星期加起來多少分鐘.
自己有沒有讀聖經.
有沒有思想用神的話.
但是你看那些媒體.
看那些電影 看那些劇集.
宣導你影響你的思想.
有多少分鐘.
懂不懂得分辨.
懂不懂得留意.
懂不懂得拒絕.

$^{1001}$第三我們看見刀劍之地.
因為出於氣灑.
出於憎勁.
我們跳一跳.
這節要講一下.
耶穌說他不憎勁.
不宣揚 街上有沒有聽見他的聲音.
這是應用主耶穌基督.
所以我們看到耶穌的榜樣就是這樣.
他不憎勁 不宣揚.
這就是主的榜樣.
求主幫助我們學習這個榜樣.
我們看看刀劍之地.
出於什麼.
出於慢神.
什麼叫慢神.
慢就是怠慢.
輕忽 輕視.
或者不敬.
輕慢.
神本是輕慢不得.
你發覺.
整個問題.
氣灑也好 憎勁也好.
背後就是因為他們.
把神放得太低.
因為他們.
太過慢神.
把自己抬得太高.
自己的血氣.
過於神的真理.
所以這張起頭.
本來我們看見很好.
最初大衛求問耶和華.
真的很好.
開頭開得.
第五節.
大衛就祝福他們.
願耶和華的慈愛.
恩代你們.

$^{1041}$所以大家留意.
大衛全部都是以耶和華為中心.
求問耶和華 耶和華和他們對話.
他就以耶和華的真理.
用他的愛對人.
以耶和華為.
他心中的大靈.
所以他就願耶和華賜福給你們.
願耶和華的.
慈愛誠實待你們.
所以我們看見.
這個就是以神為中心.
的言語行為.
祝福別人 恩代別人.
鼓勵別人不去憎勁.
我們看見這種就是靈命.
相反我們看見.
在整件事.
剛才讀了這麼長的第二章經文.
這些鬥爭裡面.
一句都沒有提及上帝.
就是人和人之間.
的仇恨 人和人之間的鬥爭.
人和人之間的血氣.
沒有提及上帝在我們中間.
為什麼會這樣呢?.
為什麼教會.
有時候都會有這些問題呢?.
因為沒有將神放在中間.
這裡本來是神聖之地.
怎麼會變成.
刀劍之地呢?.
因為他們.
萬神.
在唯一一句.
唯一一句.
約亞講.
我指著永生的上帝起誓.
你若不說.
氣散的話 今日早晨.

$^{1081}$百姓就會去.
不趕追趕弟兄了.
他只講了一句是提到神.
那麼多在哪裡呢?.
其實提到神的意思都是說.
我指著永生的上帝起誓.
借上帝的名字來講.
都是你壞.
我現在講的這句是真話.
可以在神面前講真話.
但其實在整個過程裡.
我們看到.
他們只不過是.
用自己的意思.
事實上約亞.
這句話特別的是約亞講.
約亞就出了名.
是一個血氣之輩.
約亞出了名就是非常自我自大.
約亞出了名.
就是不當神.
是很重要的.
由他講出來.
其實是口提上帝.
他口提上帝.
其實約亞本人.
就不是真是.
我們見到是一個.
真是敬畏神的人.
曾敬.
這些嫉妒.
紛爭.
是因為人大過神.
所以我們分析.
為何人大過神.
因為人的自我.
因為人的自私.
因為人的血氣.
因為人的自大.
驕傲狂妄.

$^{1121}$因為人的萬神.
當神無道.
所以我們見到.
這就是問題.
求神幫助我們.
讓我們了解人性就是這些.
人性就是自我.
自私自大.
但是靈性就是將神.
擺在我們心中.
管理我們.
敬畏上帝.
神聖之地應該是和平的地方.
不要一不喜歡自己的意思.
就隨便說隨便做.
破壞和諧.
因為其實問題在於.
心裡沒有尊敬上帝.
所以求主幫助我們.
讓我們能夠了解到這些道理.
明白這些真理.
實踐這些真理.
讓我們在我們的人際關係.
我們的家人,我們的朋友.
弟兄姊妹,最熟的那些.
更加小心.
在教會裡面.
不要讓這些.
本來.
是欺騙子.
變成.
刀劍之地.
求主幫助我們.
求主讓我們帶著恩慈,帶著鼓勵.
帶著美好的見證.
拒絕這些欺騙的言語.
尊敬的心態.
尊敬的行為.
更加的敬畏上帝.
被神管理.

$^{1161}$我們低頭祈禱.
求主幫助我們.
能夠活出耶穌基督的生命.
因為你不尊敬,不宣揚.
你說.
我心裡柔和謙卑.
你們當負我的額.
學我的養成.
這樣我們就能夠得向.
安息.
讓我們真的不要讓.
這些無謂的.
血氣的內心.
驕傲的心態.
自我的表現.
傷害了自己和你的關係.
更加不要讓我們.
透過這些.
傷害和別人的關係.
讓我們反而應該.
化乾戈成玉帛.
活出.
基督的養成.
求主聽我們禱告.
幫助我們在你面前.
藉著聖靈的更生.
再一次活出.
美好英傑的果子.
聽我們祈禱.
奉耶穌基督的名.
阿們.
\newpage



\section{馬可福音 3:13-19}
\label{sec:032457}
\textbf{耶穌為何呼召你}~~~~冼嘉琪牧師
\newline
\newline
連結: \href{http://www.church.com.hk/acms/content.asp?site=cdc\&op=show\&type=product\&code=032457\&layout=sermon}{\texttt{032457}} ~~~~ 語音日期: 2026-07-12
\newline
\newline
\hyperref[sec:032456]{\small{< < < PREV SERMON < < <}}
~
\hyperref[sec:index_chronic]{\small{[返順時目]}}
~
\hyperref[sec:032458]{\small{> > > NEXT SERMON > > >}}
\newline
\newline
馬可福音 3:13-19
\newline
\begin{longtable}{cl}
\hline
\hline
章節 & 經文 (和合本修訂版)\\
\hline
3:13 & \begin{tabularx}{0.7\textwidth}{X} 耶穌上了山，把自己所要的人召來，他們就來到他那裡。 \end{tabularx} \\ \\ \relax
3:14 & \begin{tabularx}{0.7\textwidth}{X} 於是他設立十二個人，又稱他們為使徒，要他們常和自己同在，也要差他們去傳道， \end{tabularx} \\ \\ \relax
3:15 & \begin{tabularx}{0.7\textwidth}{X} 並給他們權柄趕鬼。 \end{tabularx} \\ \\ \relax
3:16 & \begin{tabularx}{0.7\textwidth}{X} 他設立的十二個人有西門－耶穌又給他起名叫彼得， \end{tabularx} \\ \\ \relax
3:17 & \begin{tabularx}{0.7\textwidth}{X} 還有西庇太的兒子雅各和雅各的弟弟約翰—耶穌又給他們起名叫半尼其，就是雷的兒子— \end{tabularx} \\ \\ \relax
3:18 & \begin{tabularx}{0.7\textwidth}{X} 又有安得烈、腓力、巴多羅買、馬太、多馬、亞勒腓的兒子雅各、達太和激進黨的西門， \end{tabularx} \\ \\ \relax
3:19 & \begin{tabularx}{0.7\textwidth}{X} 還有出賣耶穌的加略人猶大。 \end{tabularx} \\ \\
[1ex]
\hline
\hline
\end{longtable}
$^{1}$弟兄姊妹,我想問一下,我們當中有沒有來過40年?.
40年,有啊,40年,暴露了大家的年齡..
40年前,有一個轟動世界的新聞,當時1986年的時候,.
NASA有一個火箭叫做挑戰者號,當時它升空了73秒之後就爆炸了..
事故的發現是為什麼會發生呢?.
因為工程師其實就發現在火箭裡面有一個零件叫做O-ring..
O-ring是一個如果很低溫的時候,有機會影響到這個機件是detect不到..
當時的工程師其實有提議過不如延後一天才發射,.
但是因為種種的因素,NASA最終決定預期要發射..
最終發射上去的時候,火箭升了上去73秒就爆炸,然後裡面有7個太空人全部罹難..
事後調查為什麼這件事發生,政府調查,就發現不只是O-ring這個零件上的問題,.
發現原來NASA裡面有一些決策上的問題..
當時這個決策是因為他們將要趕進度,外界的壓力,.
要趕著,因為如果延後的話,其實可能是虧蝕錢,.
或者推到其他的項目完了等等不同的原因,凌駕在安全之上..
所以自從那一次之後,NASA就要決定,所有東西不可以凌駕在安全之上..
弟兄姊妹,不是所有重要的東西都可以同時放在第一位,.
進度重要,完成項目重要,成功很重要,但是不可以全部都放在第一位..
今天我們看《馬可福音》一段經文,.
當耶穌呼召門徒的時候,有一些次序去要求門徒不可以倒轉..
同一樣東西都是這麼重要的時候,我們必須要認清次序..
耶穌說這個次序不可以顛倒,一定要知道..
我們今天來看看《馬可福音》裡面,.
耶穌說到一個門徒有什麼是最重要最重要,不可以倒轉次序..
我們一起來看看《馬可福音》..
我們一起來讀今天的經文,三章十三到十九節..
「耶穌上了山,把自己所要的人召來,他們就來到他那裡,.
於是他設立十二個人,又稱他們為使徒,.
要他們常和自己同在,也要差他們去傳道,給他們權柄廣貴..
他設立的十二個人有一個西門,耶穌又給他起名叫彼得,.
還有西比太的兒子雅各和雅各的弟弟約翰,.
耶穌又給他們起名叫伴尼奇,就是雷的兒子,.
又有安德烈,肥力,巴多羅邁,馬太,多瑪,阿勒肥的兒子雅各,達泰和激進黨的西門,.
還有出賣耶穌的加勒人猶大.」.
我們一起多一次祈禱..
親愛的主,我們一起安靜在你的面前,.
求你的話語光照我們,讓我們對準你,.
因為我們知道你的話語是寶貴..
天父,我們將時間交給你,為同心肚告奉主名求,Amen..
當耶穌呼召門徒去做使徒的時候,第一件事,.

$^{41}$我們看看他今天經文講什麼,.
他說第一件事,一定要和耶穌自己一起,.
第一件事,一定要常與耶穌同在..
耶穌開始的時候,他沒有第一件事分配說,.
十二使徒,你先去看了這本書,你先去教導人們,.
你先去接受一連串的訓練,去建立你的使宮..
耶穌呼召十二使徒,第一件事,.
他說你來,常和我一起,常和我同在..
在聖經裡面,其實一直都在講這件事,.
如果有人問聖經裡面主題是什麼的話,.
可能我們會想,救恩,十字架,耶穌基督,救贖,.
當然都是很重要的主題,很重要的神學主題,.
但裡面更加有一個主題其實貫穿了,.
由創世到啟示錄,這個主題都沒有變過,.
一直都在,那個主題是什麼呢?.
原來神說,他渴望和我們一起..
所以原來由創世記開始已經是這樣講,.
當創世的時候,神做阿當夏娃,.
阿當夏娃犯罪,離開了伊甸園的時候,.
原來她不是沒有了花園,更重要的是她沒有了什麼?.
她沒有了神同在,沒有了伊甸園,沒有了神同在,.
但後來神的心意是什麼呢?.
神在舊約的時候,他要設立會幕,為什麼?.
因為祂要和人一起,祂說要住在人間,.
進入新約的時候,祂差派耶穌基督來到這個世界上,.
所以馬太福音告訴我們耶穌是誰?.
耶穌是以馬來尼,代表著神同在,.
一直去到啟示錄的時候,啟示錄告訴我們,.
祂說,看啊,神的帳幕在人間,神要與人同住,.
一直去到啟示錄的時候,都是告訴我們神渴望和人一起,.
從創世記一直到啟示錄,神的心意沒有變過,.
神不是要一個人去完成祂派給我們的工作,.
神是想人和祂一起,神沒有說我們一定要完成那些工作,.
才做祂的使徒,神說要和我們一起,.
所以當耶穌呼召十二門徒的時候,第一句不是說你去,.
我們想起使徒,門徒的時候都是去,.
但耶穌說不是去,你是來,來,來到我當中,.
經文說,耶穌召他們來,來到耶穌當中,.
先要來才可以去,這個次序不能倒轉,.
你不來,你不能去,耶穌是這個意思,.

$^{81}$但當時的門徒不是完全明白什麼是跟隨耶穌,.
因為門徒可能只是想,耶穌是喊話跟隨他,要做師父,.
他們只是好像當時猶太的傳統,跟隨師父是怎樣?.
就是跟隨他生活,學他,好像走進少林寺一樣,.
學他和他生活,然後聽他的教導,以為跟隨師父是這麼多,.
他們不知道原來每一日經歷的,和耶穌每一餐飯說的每一句話,.
原來他們是經歷一些從舊約一直期待,歷世歷代都期待的,.
是神親自處在人間,神親自處在人間,這個很寶貴的經歷,.
有時候我們會想,不知道跟耶穌吃飯的時候,.
耶穌吃東西的樣子是怎樣的呢?.
不知道當耶穌安慰人的時候,他的眼神是怎樣的呢?.
不知道耶穌會不會說笑呢?.
有時候會在想,他吃飯的時候會怎樣呢?.
但當時門徒其實每一日經歷,每一日跟耶穌生活的時候,.
他望見耶穌,他不知道自己原來正在經歷最寶貴的事,.
他正在經歷一個最大的恩典,神親自處在人間..
弟兄姊妹,有時候我們也是,我們有時候正在經歷最大最大的恩典,.
但我們自己不覺,我們以為很平常,神親自處在人間..
今天雖然我們不能像以前的十二門徒,十二使徒那樣,.
跟耶穌一起走加利利的路,望見耶穌平靜風浪,.
雖然不可以親眼見到他怎樣醫病降鬼,.
但耶穌說,他呼召人跟他一起,這個心意一直都沒有變過..
耶穌現在怎樣跟我們一起呢?.
他說,透過聖靈,我們可以跟他一起..
聖靈內住在我們裡面的時候,我們就可以經常跟耶穌一起..
所以弟兄姊妹,神同在不只是知道神存在這麼簡單,.
神同在也不只是我們聚會有時候很感動,一下子的感動這麼簡單,.
神同在是一種生命,神同在是在家裡,在工作裡,.
在教會,在開車,在我們情緒爆發的時候,神都同在..
我自己有時經常都反省,因為我有個踏入teenager的女兒,.
大家聽到teenager就,很近是吧?.
也不是的,她很好,但我自己是我的忍耐程度越來越低,我發覺..
因為十幾歲,我就開始覺得十幾歲你應該要懂什麼,.
你應該要做什麼,慢慢我就發覺自己越來越沒有耐性,.
所以很多時候跟她說話,從很溫柔,就很快變成火就來了..
有一次火火地的時候,神突然提醒我,.
祂說,你這句話如果耶穌在你旁邊的話,你會不會說呢?.
如果這句話,耶穌在你旁邊的話,我會不會說呢?.
我那一刻突然停下來,我知道其實要改的,.
不是只是我的女兒,或者我們下一代,而是我都要改,.

$^{121}$我們一起要在神裡面改,看一下神在哪裡,看一下神在哪裡..
各位姐妹,我們有沒有活在神的裡面呢?.
很多時候我們不是不愛神,很多時候我們知道自己很愛神,.
但是因為很忙,因為有些事情,說這樣說那樣,.
慢慢我們好像沒怎麼靜下來聽一下,神跟我們說什麼..
但是神說,第一件事,我們要跟祂一起,跟祂一起..
有一套基督教的電視劇,我最近很喜歡,叫Encounter,.
我喜歡的是電視劇,因為它有電影,電視劇那個是正好,.
還是精彩一點,它每一集都有不同的故事,.
不同的故事其實是說我們人生很多,真的很生活化的一些情境,.
但是每一次這些情節的時候,突然間會有一個很智慧的陌生人出現,.
那個智慧的陌生人說了幾句話,就會扭轉整個的心情,.
後來他就會慢慢說,那個就是耶穌..
我初初看的時候,我覺得神學好像有點問題,.
不可以說自己是耶穌,一個平凡的人,.
但是其實這套劇集它想說什麼呢?.
它是想說其實耶穌在我們的生活裡面,其實耶穌在我們的身邊,.
但是每一個人就是讓耶穌在我們旁邊,但是擦身而過,沒有理會祂,.
耶穌說話的時候沒有聽祂說話,這套電視劇很有意思,.
我們發現是我們沒有留意,沒有留意耶穌跟我們說話,.
神是無所不在,但是不是人人都活在祂的同在裡面,.
神是無所不在,但是不是人人都活在神的同在裡面,.
第一次因為耶穌呼召你,不是要你完成工作,.
耶穌呼召你,是要先得著你這個人,先得著你這個人..
我們看《馬可福音》頭幾章的時候,發現一件很有趣的事,.
就是當耶穌面對人群越來越多的時候,耶穌會做什麼呢?.
耶穌好像越多人的時候就越退後,第一章說西門和那群徒弟去找耶穌,.
他跟耶穌說周圍很多人在找你耶穌,眾人都在找你,.
接著38次怎麼說呢?耶穌跟他們說,讓我們去另一個地方,我們去隔離那條村,避開..
接著又怎麼說呢?第三章,當很多人去擠擁耶穌的時候,.
耶穌就跟那些門徒去到海邊,退走,還要開一隻小船,走遠一點..
為什麼耶穌會這樣的呢?為什麼耶穌會這樣的呢?.
很有趣的一個情節,一個鋪排,三章十節說到耶穌要避開那群人群,不想別人擁著他..
但是今天經文多少節?三章十四節,十三節,就說到耶穌不是要那群人群,.
祂要的是一群使徒,祂不是要一大群人,祂要一世群人,為什麼呢?.
因為耶穌不是找人去圍觀祂,祂不是要人去擁戴祂,祂要人跟祂一起,就是要得著那個人..
祂說:來,與我同在,常和我同在..
弟兄姊妹,如果今天神走到我們當中,相信祂第一句問的說話,祂不是問你最近的服事怎樣,你的影片搞得好不好,你的歌唱得好不好..
弟兄姊妹,祂不是問我們事公上做得怎樣,相信祂會問我們,你最近有沒有跟我一起,有沒有跟我一起..
很多事公開始的時候,我們會很懇切祈禱,禁食祈禱,懇切到..

$^{161}$慢慢開始一段時間,因為太忙,因為種種原因,我們就開始會忘記了要祈禱,忘記了要跟神一起..
為什麼耶穌第一件事就要跟門徒說,你要跟我一起,你要跟我一起..
因為神不是要我們完成工作,神要我們這個人,神不是要一個工人,神要兒女回到祂身邊..
神呼召的門徒是要不斷,持續,經常的常跟祂一起,常跟祂一起..
當我們經常跟神一起的時候,我們經常在神的裡面的時候,有些事是我們會忍不住自然流露出來的..
弟兄姊妹,我們想像一下,如果耶穌呼召門徒,只是要他們日日陪著自己的話,那些門徒會怎樣呢?.
我很爽,好像看電影一樣,天天看神蹟,聽教導看神蹟,但這個不是神的心意,不是神的心意..
原來很奇妙的事情會發生,當我們常和耶穌一起的時候,我們自然會很想很想被神給我們做一些事,被神猜我們做一些事..
我們看看之後怎樣說,耶穌說:你常和我同在,然後耶穌猜他們去傳道,你都被他們權柄去講鬼..
留意原來這兩節經文有兩個要,要和耶穌常在一起,要猜他們去傳道,原來同在和猜險是不會分割的..
我們常和耶穌一起的時候,很自然就想祂叫我們做一些事,被祂猜險,這兩件事不會分割,但次序也不會擾亂,不可以擾亂..
今天很多基督徒有時會分割這兩件事,不小心,有些人很重視靈修,很重視敬拜,很重視退修,因為很棒和神一起,那種安靜是很棒..
但慢慢整個信仰就好像只剩下祂和神兩個,有些人很喜歡事奉,喜歡查經,喜歡不停去招待,喜歡和人說話,很喜歡照顧人,.
但慢慢在事奉當中,沒有了神同在的喜樂,自己不知道.耶穌說這兩件事不可分割,要同在,也要被猜險,這兩件事不可分割,要來才可以去..
所以我們教會的call value要賣廣告,教會的call value是什麼?是什麼?Being before doing,我們是要同在,然後才被神猜險,真正的使命是要有神同在才開始,.
耶穌說祂猜他們去傳道,並給權柄他們講鬼,但這件事我不知道讀到這裡大家會否很好奇,為什麼給權柄他們講鬼,而不說給能力他們教書,給能力他們講道,給能力他們醫病,為什麼要說給權柄他們講鬼呢?.
我第一個反應會問,當時很多鬼要講,是不是呢?我們留意一下,《馬學灰音》又是很有趣的,它真的經常講到耶穌趕走很多污鬼,我們見到一章,三章,五章,七章,九章都有提到耶穌怎樣趕走污鬼,.
耶穌有很多污鬼,富鬼的人來找耶穌幫忙,我們見到《馬學灰音》是鋪了一個脈絡,它要告訴人聽一些東西,是不是耶穌會講鬼這麼簡單呢?當然不是,是不是耶穌來到特別多鬼呢?當然不是,為什麼呢?.
黑暗一直都在,但是當耶穌這個真光來到這個世界上的時候,光一照落黑暗,什麼都出來了,就好像這個黑漆漆的街,一個黑漆漆的後巷,原本黑漆漆沒有人理的,但是當有光照下去的時候,所有蛇蟲鼠蟻會怎樣呢?.
馬上走來走去,找洞鑽,就是這個情況,當耶穌這個真正的君王,真正的真光來到,照亮的時候,炸彈,他就害怕,因為他知道這個才是真正的君王..
所以當《馬學灰音》記載,《灰音書》也是記載,為什麼很多人未知道耶穌是誰的時候,但是撒旦已經會講話,烏鬼已經知道,我知道你是誰,你是神的聖者,你是至高神的兒子,真正掌權的那個來到,所以照到他們害怕..
所以當耶穌猜門徒,說你拿著權柄出去講鬼的時候,其實不只是叫他們出去講鬼這個功能,他要提門徒,你拿著這個權柄,要出去宣告,神的國已經勝過黑暗的國度,神的國勝過黑暗的國度..
記不記得在《馬學灰音》裡面,其實經常提到耶穌是很有權柄的,特意記載到他是很有權柄,當然不只是特意記載,而是在文學上記錄,特意要這樣彰顯出來..
第一章講到什麼呢?他說耶穌在會堂裡面交代的時候,他是很有權柄的,人們看到他不像一個文士,然後又講到耶穌是有赦罪的權柄,然後耶穌又說他是安息日的主,即是他是神,為什麼要這樣鋪下去呢?.
原來就是在講一件事,馬學一直在講一件事,耶穌是憑什麼去做這些事情呢?耶穌憑什麼去醫治,去講道呢?耶穌是憑著他屬神的權柄,屬神的權柄,弟兄姊妹..
原來這個這麼厲害的權柄,這麼厲害的權柄,神說交給我們,交給他的使徒,交給他揀選的門徒,神授權給我們,代表他去宣告,代表他去宣告..
我們看到第一到三章其實一直在重複講這件事,當耶穌要猜他們這個權柄出去的時候,其實門徒同時就是要宣告這個是什麼神,這個是什麼神.弟兄姊妹,我們有online banking,你們有沒有online banking?有呀,大家會不會講你online banking的密碼給人的?.
不會吧?會的話是什麼意思呢?就是說我們授權給那個人去用你的account,是不是?耶穌就是這樣授權給他揀選的人去用他的account..
我們有些什麼配得給耶穌這樣信任我們呢?有些什麼配得他這樣信得過我們呢?.
有時做基督徒一段時間之後開始事奉,事奉一段時間就會發現慢慢忘記了權柄是來自神的,有些事工領袖就很容易評價一些事情,因為領袖嘛,當然要評估啦..
慢慢慢慢就形成很多的評價,很多的不同的事情發生,但是很感恩,很感恩就是我們Living Grace是沒有的,又不是賣花讚花香..
我們教會真是,弟兄姊妹是,每一個侍工的領袖,每一個弟兄姊妹,真是很愛神,很愛教會..
我們見到每一個leader,他們不是去judge弟兄姊妹行不行,每一個都是很用心去幫助大家,一起serve..
因為我們大家都看到,這個是屬神的權柄幫助我們去完成,不是我們的權柄完成侍工..
有一個牧者講了一番話,他說當所有人都認為你有權去決策,去領導,去帶頭做一些事的時候,.
但是自己能夠在神的裡面謙卑下來,承認自己沒有這個權柄,這件事是有難度,但是要學習,他的意思是要學習..
弟兄姊妹,屬神的權柄,但是我們要謙卑,現在我們每一個都很謙卑,我們一起繼續努力,繼續一起努力,在神裡面繼續謙卑..
我們知道完成侍工永遠不是靠著人的能力,我們要靠著基督的權柄,靠著基督的權柄..
我們那次去短孫的時候更加體會到,完成侍工一定不是靠人的能力,因為當我們那次去仙台短孫,去之前有一個姐妹病得很嚴重,.
病得臨飛前一天,她就想,還要不要去呢?是否真的要去呢?如果去的話要被人照顧,是否適合去呢?.

$^{201}$她就祈禱,求神,如果讓我下午一點前舒服一點,讓我知道我應該去還是不去..
誰知神真的讓她在下午一點鐘,因為祂設定了時間,祂覺得需要跟團隊交代,祂說真的讓她在下午一點鐘好很多,然後第二天照樣上飛機..
我們二十個小時的旅程,飛到去的時候更加奇妙的是,她病情沒有惡化,也不需要躺在床上,神奇妙地醫治,讓她一天可以完全恢復她的體力..
這是神彰顯她的權能,不是人可以做得到.短孫去之前,我們知道可能有風暴,看天氣預告,日本有風暴,我們能不能飛呢?.
我們祈禱,去到我們的旅程,一滴雨都沒有,神沒有被風暴影響我們,有幾滴雨,無細雨,在晚上,我們都沒有出去..
我們知道完成事工是神自己的權能,當我們再看看使徒的人名表的時候,我們更加會知道,完成神的事工不是靠人的能力,不是靠人的能力..
我們看看這個人名表,如果你是耶穌,你會不會選這十二個人呢?我們看看,彼得是什麼人?彼得是很衝動的人..
然後,雅各和約翰是什麼人?是兩兄弟,為什麼會有這樣的名字?叫雷子,因為他們脾氣很火爆,所以叫雷子,幫他們改名叫雷子.然後,馬太是什麼人?馬太是收稅的,稅吏..
馬太在這個時候,不知道辭職了沒有,但她是稅吏,幫政府做事.然後還有激進黨的西門,我們看和合本的時候會知道,粉瑞黨的西門,是什麼呢?什麼叫粉瑞黨呢?.
即是說他是反政府的,很憎恨羅馬政府.一個是幫政府做事,一個是反政府的,在這個團隊裡面,可能第一個會議已經鬧大架了..
如果我們是耶穌,會不會選擇他呢?十九節,重大件事,如何形容呢?這是什麼人呢?出賣耶穌的加略人猶大,結果我們都知道,耶穌呼召他們的時候,除了知道他們本身已經是這樣,已經不容易,.
頂得住.更加,你說知不知道他們之後是怎樣呢?你說知不知道他們之後是怎樣呢?他們之後,我們都知道,給了三次不認署,耶穌有事的時候,都走了,猶大出賣他..
你猜耶穌知不知道猶大出賣他的時候,說的話,那一個眼神,那一些背叛,你猜耶穌知不知道呢?耶穌每一天都知道,每一天都知道..
但是耶穌仍然呼召,仍然讓這12個人和他相處,仍然開了一個帳戶給他們用,授權給他們..
為什麼呢?這就是恩典.弟兄姊妹,我們都是領受同樣的恩典,我們都是領受同樣的恩典..
有些弟兄姊妹會想,我相信耶穌,不要叫我做門徒,門徒可能要做這樣,做那樣,很辛苦,要有些犧牲,很多責任..
但是原來耶穌今天和我們說,我們看經文的時候,耶穌今天和我們說,做門徒,第一要緊的事,你是要和我一起,你來,回到我身邊先,你來,回到我身邊先說..
耶穌從來沒有弄一個清單說,你要接受多少門徒訓練,你要看多少次聖經,看多少次福音書,你才可以做我的門徒..
耶穌第一件事就是說,來,與我同在,然後自然你就會很享受,很喜歡,出去,開著我的帳號任用.耶穌從來都是說,來,才去..
丁姊妹,如果我們今天要回應耶穌對我們這個呼召的時候,我們會怎樣回應呢?原來祂說做門徒,就是來,與耶穌同在,才去..
我們可不可以嘗試每一天問耶穌一個問題?主啊,今天我要怎樣和你一起走?今天你想我拿著你的信用卡怎麼用?開著你的戶口怎麼用?.
好嗎?丁姊妹,我們一起低頭祈禱.親愛主,多謝你,多謝你讓我們可以有你在我們當中..
多謝你和我們說,原來最重要最重要,有些次序不可以倒轉,就是先要和你一起,做你門徒就是要在你的同在裡面.主啊,這是很無比的福分,這是很大很大的恩典..
父啊,多謝你,你沒有嫌棄過我們曾經怎樣背叛你,你沒有嫌棄過我們有時害怕,有時遠離你,你沒有生氣過我們有背叛你的時候,你仍然呼召我們..
你說,來,與我同在.當你呼召我們做門徒的時候,你同時跟我們說,兒子女兒,回來,坐在我身邊..
父啊,多謝你,多謝你愛,多謝你的愛,不離不棄,以致我們可以有力,以致我們可以有你的愛,被你拆出去..
天父,多謝你,我們同心禱告奉主耶穌基督的名,Amen..
\newpage



\section{士師記 18:9}
\label{sec:032458}
\textbf{聖經人聖經事:沒有王,各人任意而行 (2) 但人以神之名謀劃殺人奪地}~~~~黃克勤牧師博士
\newline
\newline
連結: \href{http://www.church.com.hk/acms/content.asp?site=cdc\&op=show\&type=product\&code=032458\&layout=sermon}{\texttt{032458}} ~~~~ 語音日期: 2026-07-13
\newline
\newline
\hyperref[sec:032457]{\small{< < < PREV SERMON < < <}}
~
\hyperref[sec:index_chronic]{\small{[返順時目]}}
~
\hyperref[sec:032459]{\small{> > > NEXT SERMON > > >}}
\newline
\newline
士師記 18:9
\newline
\begin{longtable}{cl}
\hline
\hline
章節 & 經文 (和合本修訂版)\\
\hline
18:9 & \begin{tabularx}{0.7\textwidth}{X} 他們說：「起來，我們上去攻打他們吧！我們已經窺探了那地，看哪，那地非常好。你們還要待在這裡嗎？不要再遲延了，立刻出發去得那地為業吧！ \end{tabularx} \\ \\
[1ex]
\hline
\hline
\end{longtable}
$^{1}$粉絲們大家平安.
又和大家讀聖經了.
生命樹預備了聖經人,聖經時.
最重要看到聖經後面.
上帝做些什麼.
那時以色列沒有王.
各人任意而行.
減少了以色列的字.
因為不夠位置.
這個沒有王,沒有王管.
沒有王,沒有王管.
各人任意而行.
這個任意而行的第二輯.
哇 做什麼呢.
就是彈人.
彈這個支派.
彈支派的人.
不是全部 部分.
以神之名.
以神之名即是大殺.
對殺,arm 殺.
站在道德的高位.
去謀劃殺人奪地.
上帝派我來的.
我應該要拿你的地方.
殺死你.
很可怕.
就是沒有王管.
各人任意而行.
而出現的事.
好 看看經文.
未看經文先看地圖.
猶大支派.
彈支派就在這附近.
然後只是字遮了.
以法連.
猶大支派,以法連支派.
馬來西在約旦河的河西.
有馬來西.
約旦河的河東有馬來西.

$^{41}$然後有迦德留邊.
即是馬來西,迦德留邊.
三個支派.
都在約旦河的東邊.
西邊有猶大,以法連,彈馬來西.
以撒加,西布倫,阿徹,拿夫塔利等等.
在約旦河的西岸.
這個海岸線就是地中海.
地中海在北邊.
今天屬於黎巴嫩.
有西頓和推羅.
推羅,西頓.
兩個都是靠地中海邊的城市.
推羅後來發展成一個很大的海港商業城市.
西頓沒有它那麼厲害.
就算在主耶穌的時代.
推羅,西頓這些地方.
都屬於外邦人的地方.
今天比較屬於黎巴嫩那邊.
大家看看.
這裡有兩個綠色字.
就是索拉和以實圖.
是兩個很小的附近小城市.
索拉,以實圖有很多彈人住在這裡.
這裡一直都是屬於非利士地.
索拉和以實圖裡面彈人的一部分.
由於自己沒有自己的土地.
於是他們竟然周圍派探子去找地.
找到很北的.
還高過加利利.
去到拉翼.
拉翼這個小城市.
發覺那裡很太平.
沒有防守,很平靜.
又沒有朋友,沒有人來幫他.
於是他們想.
不如我們攻佔了那裡來自己用.
由金南沿海去到幾乎屬於南面.
然後走到金北搶人的地方.
搶的.

$^{81}$中間還發生了一些事.
我們看看彈人是多麼差.
但是很差的.
那時以色列中沒有王.
各人任意而行.
然後第二輯.
如何任意而行.
就是以神之名謀劃殺人奪地.
第二輯.
十八章第一節.
那時以色列中沒有王館.
沒有王.
但支派的人仍是尋地居住.
什麼意思呢.
就是由於當約書亞帶他們進入迦南的時候.
大家還記得.
我都提過了.
講解過了.
就是將十二支派分了地.
然後叫他們自己去拿這個地.
不是約書亞帶齊十二支派去打了整個迦南地.
然後分的.
是看著地圖分了.
然後你們自己去拿那些城市.
由大應該拿了非利士.
不過他又沒有.
當時神說你們進去的時候.
要將那些人.
起碼如果可以的.
就將他們殺死.
那麼殘忍.
其實是神用以色列人來作為皮鞭.
懲罰這些很壞的迦南人.
很多種族的.
如果不是你就要趕他們走.
和要拆爛他們的祭壇.
不要和他們結婚.
不要要他們的東西.
如果不是.
他就會成為你將來的網絡.

$^{121}$和肋骨中間的荊棘.
撬著你 刺著你.
千萬不要.
他好像網似的看著你.
千萬不要結婚.
不要拜他們的神.
拆爛他們.
將他們趕走.
或者甚至將他們殺死.
但是很多人都沒有做.
後來見到他們都住在那裡.
聊天都挺好的.
於是想一想.
如果我打得贏他.
不如我抓了他的人.
和我煮飯洗衣服做奴隸更好.
有些就說好啊.
你不要殺我.
我就幫你做事.
所以就很多人在裡面.
接著全部很多以色列人.
都開始鑽研了他們的假神偶像.
大部分的迦南的神.
不同種族的.
有個通用的字.
在他們的信仰裡面.
就叫巴力.
就是主.
每個人都拜巴力.
不過不同的種族裡面.
有不同的巴力.
他們就在那裡這樣.
哎呀.
但支派仍是沉地居住.
就是他們還沒有落實到.
這麼大的支派.
那些人住在哪裡.
那麼很多已經落實了.
不過部分就沒有.
沒有那怎麼辦.

$^{161}$那麼就.
因為到那日子.
他們還沒有.
在以色列支派中得地為業.
那麼就沒有派給我.
不是.
給了.
劃分了.
不過他沒有去得著.
但是他又不去得著.
不去得著可能是.
姑息.
或者是.
怕著不夠人打.
種種原因.
總之他.
還沒有得到地為業.
不是全部彈.
所有彈的支派.
彈支派裡面的.
少數人.
好了.
第二次彈的人呢.
彈.
那麼就代表著支派了.
從索拉和以實圖.
打發本族中的.
五個勇士.
去仔細窺探那地.
那地在哪裡.
就是整個迦南地.
還看哪裡有沒有洞.
沒有什麼人的地方.
我們去那裡是合用的.
我們就用那塊地.
那麼我們剩下彈的人就住在那裡.
那麼就叫做窺探那地.
如果照這樣看.
就應該他們走遍整個迦南了.
看看地圖.

$^{201}$他們由南邊一直向上走.
為什麼不去下面呢.
下面那些就是曠野了.
南地.
不適合人類居住的.
由底到最南的.
他們真的走遍整個迦南地.
我們看看有沒有什麼地可以住.
就去到最北.
沒有什麼.
其實都沒有什麼以色列人在那裡住了.
那麼他就去到這個拉翼這裡.
好,看下去.
那麼你們去窺探那地.
吩咐他們說.
你們去窺探那地.
然後呢.
就回來再說.
他們來到以法連山地.
就進了米迦的住宅.
就在那裡住宿.
咦,米迦記得嗎?.
就是昨天說的那個米迦.
他就招待了一個猶太支派.
上來以法連山地.
又要翻回地圖.
不好意思.
猶太支派這裡.
所拉爾什圖.
以法連山地就是這個位置.
以法連山地.
那麼猶太支派有一個利未人.
就找地方住.
就去到以法連山地.
那麼就去到一個人家裡.
就叫做這個米迦.
昨天前天都說了.
米迦是一個很壞的人.
偷了他媽媽1100赤霞蘭銀子.
然後有一天就跟媽媽說.

$^{241}$媽媽,你不是被人偷了1100赤霞蘭銀子.
20多公斤的銀子.
那麼你當時就罵了所有人.
就坐了.
其實媽媽的錢在我這裡.
是我偷了你的.
現在還給你.
那麼他媽媽說.
哎呀,真是願耶和華賜福給你.
那麼我就拿這些錢出來.
做一個銀的神像.
然後就給你用.
敬拜耶和華.
另外那些錢都交給你.
反正都已出之物,不見了.
全給了米迦.
米迦自己有間很大的豪宅.
那麼當時算是很多錢.
那麼就見到這個.
猶大的利美族人.
就說不如你在我這裡住.
我以你做我的傅.
做我的老師.
做我的祭司.
那麼我都聽你的了.
我家裡會開一個神堂.
這樣子.
然後就做了一件.
大祭司的爾伏特自己做了一件.
唉,這樣也有的.
那麼這個米迦真的不是好人.
然後呢.
又做一個銀的神像.
還有各樣神像放在這裡.
就開一個神壇.
或者叫神堂.
就給人去拜了.
那麼猶大利美族人走上來.
就做了祭司了.
啊?.

$^{281}$那麼利美族做祭司是否正確呢?.
當然不正確了.
因為祭司是要阿倫的後裔才可以做的.
不是利美派每個人都可以做的.
只是阿倫.
及阿倫的後裔.
沒有你.
米迦自己我是金主.
有什麼不行的?.
何必去到士劍這麼遠的地方.
回望敬拜.
我這裡就可以了.
我上去就可以了.
就這樣吧.
那麼那五個勇士去窺探的.
就進了米迦的住宅.
就在那裡住.
好好地在那裡住.
那麼臨近米迦的住宅.
就去到這個臨近的去到.
就聽到這個少年人.
這裡的祭司是利美口音.
是猶大利美口音.
那麼就問他了.
就是誰領你到這裡.
你在這裡做什麼.
你在這裡得到什麼.
探子嘛.
盤人先.
盤了他.
喂.
你哪裡來的.
可能去到米迦住宅門口.
就利未人出來迎接他.
以為他當然是來光顧.
都不一定了.
這五個外人.
那麼就聽他說話.
先生.
你們想要.

$^{321}$你們要做什麼.
咦?.
利美口音.
就是猶大的利美口音.
雖然是以色列人.
但是口音不同.
那麼就盤他了.
問他什麼呢.
誰帶你來的.
你在這裡做什麼.
在這裡賺了多少錢.
得到什麼.
三個問題.
都很重要的.
第一點要和大家說.
他們去仔細窺探那地.
去找.
就是這裡這樣出的.
那麼為什麼還要去找呢.
就是他們開頭的時候不賣力.
變成自己不夠地方住.
那麼這一群彈的人.
就要去找地了.
那麼如果他們在附近.
出力找都應該找到的.
不過不想就擠迫.
各種原因.
派五個勇士去找.
那麼那個勇士就發生這件事了.
就問你在哪裡來的.
你為什麼來這裡的.
賺了多少錢.
就去到這個豪宅那裡.
他就回答他了.
就是米家這個豪宅的主人.
待我如此如此.
還請我做祭司.
所以我做了祭司了.
那你做了祭司了.
那你是不是阿倫後人了.

$^{361}$不是啊.
這樣都可以.
好啊.
不要緊.
你做了就做了.
那麼他們去窺探.
然後他們就說了.
第一他們去窺探.
當然剛才我們也問了幾個問題.
就是我們的人生都是這樣的.
誰領你到這裡來.
你今天在做什麼.
你得到些什麼.
我們經常都要問自己的.
那麼辛苦.
我得到些什麼.
就是我的生命這樣用法.
對還是不對.
好還是不好.
我有沒有投資.
我的人生有沒有問題.
那麼主是怎樣帶領我.
好了.
來到第二點了.
他就說你既然是祭司.
那不如你和我們去求問耶和華.
就是讓我們知道.
我們所行的道路通不通達.
就是這次我們去信不信你.
你就下馬問前程了.
就是你和我們占卜一下前面.
求問耶和華.
請你求問神.
使我們知道.
我們這次去做窺探.
信不信.
有沒有危險.
問這件事.
都可以.
當你有很多人生問題的時候.

$^{401}$你要做的就是求問神.
求問神.
神是會給答案你的.
誰知道祭司對他們說.
你們可以平平安安地去.
所走的路都在耶和華面前.
那就是說什麼.
他問都沒有問.
馬上就回答.
你們可以平安去.
這一程平安.
那你就不怕了.
你走吧.
儘管去.
耶和華在你面前.
祂一直帶領你在你面前.
帶領你保護你.
安排給你.
祂照著你的.
你不妨前行.
平安的平安的.
這樣.
這些你看到舊約裡面說的假先知.
經常說平安嘛平安平安.
儘管去儘管去.
神與你同在祝福你.
說這些話.
這樣就中聽.
這個猶大的一個走出來的利未人.
又不是亞倫之派.
不過他做了祭司之後.
就滑頭了.
你看到他.
平不平安啊.
跟我求問.
不用求問了.
我知道了.
平安平安.
去吧.
就耶和華在你面前照著你.

$^{441}$這樣.
那他們五人就走了.
就來到拉益.
就是說從下面走到相當遠了.
走遍全個迦南地.
見那裡的民安居無慮.
如同西頓人安居一樣.
而且那地是沒有人掌權擾亂他們.
很有趣的.
掌權的人會擾亂人嗎.
是啊.
兇死嘛.
全部趴生嘛.
這樣就沒有人掌權.
沒有人擾亂.
離西頓人都遠.
與別人沒有來往.
再看一次.
悶悶的.
不過還是要知道的.
由索拉和以實圖兩個地方.
一路走到整個迦南地.
走到北.
就去到拉益.
就見到這個地方.
與西頓隔離很遠.
與推羅隔很遠.
這些大城市.
附近都沒有什麼大事.
那就是.
又沒有人騷擾.
又沒有人保護.
有小朋友.
小朋友不是小朋友.
是沒有朋友.
朋友很少.
而且在這裡.
又沒有人話事.
一班人住在那裡.
他們就想.

$^{481}$咦,這裡挺好.
你再看.
有很多河流.
水源充足.
好地方.
不過北部沒什麼人住.
於是他們覺得好.
他們去到.
挺漂亮.
這個地方叫拉益.
五個人就回到夏索和以實圖.
由北邊又走回去.
見他們的兄弟.
弟兄就問他.
有沒有什麼要告訴我們.
你們有什麼.
有什麼結果.
他們回答說.
起來,來吧.
第三點要說的.
那五個人說什麼.
我們想去攻擊他們.
嘩,很aggressive.
這樣很兇.
我們去攻擊他們.
攻擊哪裡呢.
攻擊拉益.
我們已經窺探過那地.
見那地甚好.
喂,為什麼你們為何靜坐不動.
還坐在那裡幹什麼.
急速起來.
慈望去得那地為業.
不可遲言.
馬上去.
不要遲言.
要得那地為業.
我們上去幹什麼呢.
攻擊他們.
不要坐在那裡不動.

$^{521}$馬上去.
急速去得那地.
不要遲言.
快.
為什麼.
心急.
除了心急.
其實有沒有怕別人去搶他呢.
應該沒有.
不過是心急.
嘩.
雄心萬丈.
去搶別人的地.
我們上去.
攻擊他們.
不是和平.
不是傾.
我們都很多人想在那裡住.
行不行呢.
那我們一起住吧.
我們很和平的.
不是.
我們去攻擊他.
攻擊.
即是說.
動手.
這樣很差的.
我們去攻擊他們.
為什麼還要坐在那裡.
馬上去.
攻擊他.
他們是.
神帶我們去.
還問我們去.
平不平安.
求問神.
然後.
去了平安之後說.
神帶領我們的.
所以.

$^{561}$但這班人是以神之名去籌劃殺人.
搶他們的地.
覺得很可憎是什麼呢.
他們任意而行.
自己說.
就是神之名帶領我們去的.
問了一個這樣的.
艾亞祭司.
祭司就說平安平安.
是的.
神帶領我們了.
帶我們去看到.
即是神帶領我們去得那個地.
我們快點去.
神給我們的.
將他們攻擊殺了.
我們拿到那塊地了.
馬上去.
嘩.
很差.
是沒有王任意而行的第二輯.
這樣就看到沒有王.
那個人任意而行.
以神之名去謀劃殺人奪地.
很差的.
這樣的人也有的.
有.
舊的也有.
新的也有.
今天我們的社會也有.
甚至今天我們的教會.
也有類似情況.
不過沒有殺人.
搶奪權位.
在教會很多.
排斥其他人殺死他.
殺死他的意思是.
等待別人離開教會.
甚至離開主也很多.
我們一起祈禱吧.

$^{601}$天護上帝我們看到.
事司祭的結束是很混亂的.
因為沒有王.
只有事司是每一個民族自己的法官.
帶領他們抗拒外族的武力入侵.
有時是跟我們判斷生活上的爭執糾紛.
沒有什麼聽到有先知.
沒有什麼聽到有祭司出來說話.
都很沉寂.
以至他們只記得耶和華.
於是自己是將很多心裡面想要的東西.
投射到耶和華.
神這個名稱裡面.
就是神叫我們做的.
這個很多時候是我們一個很大的弱點.
這個都是叫網稱耶和華的名.
主說其實我們全部都認識這些十戒.
全部都認識這些不可殺人.
全部都認識這些東西.
不過他們就網稱網用你的名.
就說急速啊.
我們很趕快上去德勒地攻擊他們.
是神給我們的.
主求你憐憫.
給我們小心.
我們更加以愛相繫.
給我們更加和平地共處.
這是聖經一直給我們的教訓.
到了新約更加明顯.
求主幫助我們用愛來維繫.
垂聽禱告.
奉耶穌基督名求.
阿們.
\newpage



\section{馬可福音 1:1-8}
\label{sec:032459}
\textbf{為福音,好好預備自己}~~~~梁柏堅弟兄
\newline
\newline
連結: \href{http://www.church.com.hk/acms/content.asp?site=cdc\&op=show\&type=product\&code=032459\&layout=sermon}{\texttt{032459}} ~~~~ 語音日期: 2026-07-11
\newline
\newline
\hyperref[sec:032458]{\small{< < < PREV SERMON < < <}}
~
\hyperref[sec:index_chronic]{\small{[返順時目]}}
~
\hyperref[sec:032460]{\small{> > > NEXT SERMON > > >}}
\newline
\newline
馬可福音 1:1-8
\newline
\begin{longtable}{cl}
\hline
\hline
章節 & 經文 (和合本修訂版)\\
\hline
1:1 & \begin{tabularx}{0.7\textwidth}{X} 神的兒子，耶穌基督福音的起頭。 \end{tabularx} \\ \\ \relax
1:2 & \begin{tabularx}{0.7\textwidth}{X} 正如以賽亞先知書上記著：「看哪，我要差遣我的使者在你面前，他要為你預備道路。 \end{tabularx} \\ \\ \relax
1:3 & \begin{tabularx}{0.7\textwidth}{X} 在曠野有聲音呼喊著：預備主的道，修直他的路。」 \end{tabularx} \\ \\ \relax
1:4 & \begin{tabularx}{0.7\textwidth}{X} 照這話，施洗約翰來到曠野，宣講悔改的洗禮，使罪得赦。 \end{tabularx} \\ \\ \relax
1:5 & \begin{tabularx}{0.7\textwidth}{X} 猶太全地和全耶路撒冷的人都出去，到約翰那裡，承認他們的罪，在約旦河裡受他的洗。 \end{tabularx} \\ \\ \relax
1:6 & \begin{tabularx}{0.7\textwidth}{X} 約翰穿駱駝毛的衣服，腰束皮帶，吃的是蝗蟲和野蜜。 \end{tabularx} \\ \\ \relax
1:7 & \begin{tabularx}{0.7\textwidth}{X} 他宣講，說：「有一位在我以後來的，能力比我更大，我就是彎腰給他解鞋帶也不配。 \end{tabularx} \\ \\ \relax
1:8 & \begin{tabularx}{0.7\textwidth}{X} 我用水給你們施洗，他卻要用聖靈給你們施洗。」 \end{tabularx} \\ \\
[1ex]
\hline
\hline
\end{longtable}
$^{1}$經文是這樣的,是《馬可福音》的一章一至八節.
如果你們可以看著我聽,聽我讀出.
「神的兒子耶穌基督福音的起頭.
正如先知耳措書上記著說.
看那我要差遣我的使者.
在你前面預備道路.
在曠野有人聲喊著說.
預備主的道修正他的路.
召者說:約翰來了.
在曠野施洗.
全悔改的洗禮.
使罪得赦.
猶太傳地和耶路撒冷的人.
都出去去到約翰那裡.
承認他們的罪.
再約旦河裡受他的洗.
約翰穿落駝毛的衣服.
腰束皮帶吃的是旺蟲野蜜.
他傳道說:有一位在我耳後來的能力比我更大.
我就是彎腰和他解鞋帶也是不配的.
我是用水給你們施洗.
他卻要用聖靈給你們施洗」.
這裡我會用三個的向度.
來跟大家說這個預備是會怎樣發生.
第一個預備我會看第一到第五節.
第一節的時候一開始就是這樣的.
「神的兒子耶穌基督福音的起頭」.
這句好像很平平無奇的說話.
它是一個怎樣的.
對我們好像平平無奇.
但其實一點也不平平無奇.
福音書的編排.
我們讀本聖經的時候.
就是馬太,馬可,路加,約翰.
但聖經出現的經卷時序.
福音書的順序.
一般學者會理解.
馬可福音的.
是眾多福音書裡的第一卷出現的福音書.
而換言之這句話一章一節.

$^{41}$這句福音書第一次有福音書出現.
然後第一句話.
就是在說「神的兒子耶穌基督福音的起頭」.
福音這個字我們可能都很.
大家都傳福音.
有些人聽過福音.
信耶穌又要用福音三章十六節.
大家可能聽過很多這樣的說話.
但福音在約翰這個時代.
司祭約翰的這個時代.
其實它不是一個教會信仰的字眼.
事實上這個字眼的意思是誰用的呢?.
其實是凱撒用的.
凱撒在那個時代被認為是福音.
他為世界帶來福音.
凱撒為當時的世界帶來福音.
羅馬帝國為當時的世界帶來福音.
因為他平定戰亂.
他將所有的戰亂鎮壓了.
壓低了,剷平了敵對勢力.
所以當時世界由大混亂開始變成平定.
所以當時的人就認為.
凱撒為世界帶來福音,好的消息.
然後這個好的消息.
以至他甚至被當時的人.
形容他是神的兒子.
當時的錢幣的後面.
會寫神的兒子是凱撒.
凱撒的樣子是印在上面.
神的兒子這樣描述.
換言之第一章一節這個故事開始的時候.
他就說神的兒子不是凱撒.
是耶穌基督,祂的福音在這裡開始.
是一個宣告.
宣告一個新的好消息.
這個宣告原來不是突然間平地一聲雷.
他在一章二至三節裡面描述.
這個宣告其實是有一個源頭.
源頭是來自以塞俠書.
以塞俠先知書上記著說.

$^{81}$這句說話是「看那我要差遣那使者在你前面預備道路」.
「在曠野人聲喊著說,預備主那道修節他的路」.
是甚麼來的呢?.
原來他是來自以塞俠書的四十章.
四十章又是一件甚麼事呢?.
以塞俠書頭三十九章是講很多審判.
他在描述一件事.
描述當時的以色列那些猶太人.
是不跟從上帝的說話.
接著他們就去搞那些政治結盟.
經濟的力量各類各類的事情.
然後他不跟上帝的道去行.
那些孤兒寡婦,城門口給人公義全部不要來.
以致他們被擄亡國.
一個這樣的結局.
這個就是審判的結局.
在以塞俠書頭三十九章就是不斷描寫這件事.
去到第四十章這句說話開始的時候.
意思是甚麼呢?.
就是世界有一個新的一頁要打開.
你們的懲罰夠了.
上帝對你們的懲罰夠了.
這個世界有一個新的一頁要打開.
只要你悔改.
你就可以連接和上帝的關係.
我們重新去建設這個拯救的道路.
一個這樣的說話,以塞俠書第二三節.
這個故事就是耶穌基督福音的起頭的一個連接.
這個連接有甚麼出奇呢?.
特別的地方就是這句說話.
你要想像一下它是何時出現的.
你見到旁邊有一卷像爛布的東西嗎?.
角落的螢幕上有一個像爛布的東西.
那卷是甚麼來的呢?.
那個是我們叫死海古卷的東西.
其實有一個叫昆蘭社團.
它在死海北面的地方.
他們有一個聚集地方.
他們會抄寫聖經.
然後他們會有一個等待上帝國度的聚集.

$^{121}$他們在那裡抄寫聖經的時候.
整個舊約的經卷他們都抄了.
有趣的地方就是他們被發現的時候.
裡面有一個經卷很多的.
就是以塞俠書.
以塞俠書在他們的copies裡面.
有二十多個copies是和以塞俠書抄寫出來的.
有趣的地方就是原來他們重視以塞俠書.
那個時代的人重視以塞俠書.
為甚麼呢?.
因為你要想像抄寫一卷聖經.
不是像今天這樣印出來的.
是抄出來的.
據說最近到八月尾.
以色列博物館在展示一些以塞俠書的經卷.
如果你有興趣可以去看一下.
裡面在描述它和今天的舊約是一樣的.
抄寫的時候你很仔細.
不可以錯一個這樣的狀態來抄寫.
然後在這樣的力量之下抄寫了二十多個copies.
意思是甚麼呢?.
意思是有很多人在讀書.
所以才要抄寫這麼多.
如果沒有人讀書我不會抄寫這麼多.
不會花這麼多人力物力去做這件事.
花這麼多人力物力.
意思是有很多人在讀書.
不單止在這裡有很多人在讀書.
你記得耶穌在會堂裡面讀甚麼嗎?.
會堂裡面都是在讀以塞俠書.
神的靈在我身上用高高我.
叫我參與報告被擄得釋放.
乞訝得漢奸.
受壓制得自由.
報告神月立人的欺凌.
是不是?.
都是以塞俠書.
換言之以塞俠書這些經卷這些內容.
一直都在當時非常流行.
特別在約翰的時代.

$^{161}$好了.
約翰的時代距離這句話出現相隔多遠的時間呢?.
新約舊約一般來說我們都說是四百年.
而舊約這段以塞俠書的經卷這個時間.
去到這個位置的時間.
學者會認為大概是五百年.
換言之就是五百年來.
這句話在社群裡面一直流傳流通.
甚至流通到約翰的時代.
很多人讀了.
問題是這句話有多少人用他的生命走出來呢?.
有多少人用他的生命將這句話.
預備著的道收直他的腦.
這句話說了很多次了.
有多少人將它讀入心.
然後將它變成自己的生命的訊息.
然後他將自己的生命連上一章一節.
上帝的故事.
耶穌基督福音的起頭.
這五百年來有多少人做這件事?.
或者我們可以特別留意第四節的頭三個字.
召者說.
有時候我們會覺得召者說好像只是一個預言.
一個預言 召者說這句話就發生了.
好像是一個連接著一個預言的意思.
但是召者說你再讀下去的時候.
有一種味道的就是.
誰召者說?.
上帝有沒有呼喚你 呼喚我在這裡召者說.
這句話有沒有入到你的心?.
上帝的說話有沒有入到你的心?.
我們可能在返教會有一段日子.
我們在這裡每個星期都會聽到很多的說話.
重點不是在於這些說話解釋得有多精妙.
重點是就算解釋得不精妙.
就算我今晚解釋得很差.
解釋到你每個人都不想聽.
但是讀完經文之後 你看完經文.
那段經文感動你.
你召者說這句經文就在你心中開始發動力量.

$^{201}$我在上年做了一件奇怪的事.
就是我做書籍編輯 入特朗普的時候.
接著我就幫了一個牧師 叫做伍維文牧師.
他就幫我編了一本書.
牧師八十多歲.
他是崇基學院 中文大學崇基學院的校務.
他在崇基學院做了十多年的校務.
是崇基學院裡面最長的校務.
然後他將崇基學院裡面的崇基禮拜堂.
那個教會的崇拜.
本來只是在上下學期的時候有聚會.
接著他將它變成全年都有聚會.
以致慢慢形成了一個信仰群體.
伍維文牧師就在崇基學院裡面做了一件這樣的事.
有一天有一個朋友inbox我.
問我可不可以幫伍維文牧師編一本書.
其實那時候我已經很久沒有編過書了.
所以就以這件事.
但是一想到伍維文牧師是因為他曾經幫我的神學院老師.
招攬了他去信義中神學院.
我在信義中神學院讀神學.
以致我認識了我當時一個對我很重要的恩師.
叫做鄧兆光老師.
因為鄧老師影響了我很深.
所以我就幫他做這本書.
這本書做的時候就編完了.
然後牧師就問我.
這本書你覺得最深刻的是哪句說話呢.
有沒有一些片段是最深刻的呢.
我就選了這一段給他看.
這一段其實是來自啞國.
他和上帝摔跤.
然後就將很多不情願.
很多尋求神的祝福.
很多恐懼各種的狀況.
然後他仍然死在上帝.
求主去幫他.
的一個過程裡面.
他和上帝的故事連接.
以致十二支派從他而出.

$^{241}$一個這樣的故事的延展.
牧師就說.
其實我們對上帝國度的認識.
對神的認識.
是多過啞國很多的.
因為啞國是舊約時代的人物.
啞國是一個舊約時代的人物.
他不知道是誰.
他不知道耶穌.
他不知道上帝會有道成肉身.
叫耶穌來到中間.
他不知道有十字架.
為我們的罪去代贖.
他不知道十字架之後.
有三天的復活.
他不知道上帝會差派聖靈來.
啞國也不知道聖靈來到.
會建立一個信仰群體叫教會.
他不知道以後發生的事.
但是啞國憑著他對上帝的.
一輪絆肘的認識.
他將自己的生命投入去.
然後他走經他自己的時代 他的時間.
然後他將生命照著他走下去的時候.
留下了上帝國度的痕跡.
接著我說牧師你留下最後那三個字.
真是很硬心的.
也很難的.
他說 我們呢?.
啞國就是這樣走經時間.
留下痕跡 我們呢?.
到底我們知道了更多東西.
我們有更多學術的知識.
我們有更多人去演繹聖經.
我們有很多 甚至今天.
我們有時候看聖經.
你問AI也會答得很頭頭是道.
是不是?.
我們也有很多的書.
很多的資料.

$^{281}$也有很多人被我們問.
你問一下牧師 問一下傳道.
很多人都會很樂意地回答你.
我們知道更多.
但是我們更多的反思.
有多少會變成照這話.
成為我們生命的力量呢?.
求主幫我們.
可能這樣的事情是不容易的.
因為要照這話 有時候.
上帝你的話甚難.
我們有時候就是這麼軟弱.
這麼不容易.
但是只要你開口求主幫你.
上帝的說話會一點一滴.
在你的生命裡面實現.
然後你會在上帝的國度裡面.
留下痕跡.
這是第一個預備.
來到的人 他要尋找的.
是尋找一些會照這話.
將他感動到聖經的說話.
感動他的東西.
是會在生命裡面演繹出來的.
來的人要尋找這樣的人.
這是第一個預備.
第二個預備.
第二個預備是這樣的.
第二個預備在第六節.
若翰穿駱駝毛的衣服.
腰束皮帶吃的是蝗蟲野物.
我會先說它的外表.
他穿的是駱駝毛的衣服.
腰束皮帶.
這是一個什麼形象呢?.
我就叫AI幫我拍一張照片.
他拍了一張野人的照片.
穿駱駝毛的衣服.
腰束皮帶抗日的樣子.
很有趣.

$^{321}$對以色列人來說.
這個樣子是一個什麼樣子呢?.
其實是一個相當熟悉的樣子.
他熟悉的樣子是什麼原因呢?.
因為這個樣子.
是和以前另一個歷史人物很相似的樣子.
在列王記下一章八節的時候.
有一個王叫阿哈謝.
有一次那些蝗蟲和王說.
有一個人來求見你.
這個人是什麼人呢?.
他身穿毛衣.
腰束皮帶.
王就說這個人一定是.
提斯比人以利亞.
原來這個形象.
其實在以色列人的腦海裡面.
是有一個印象.
就是以利亞.
身穿毛衣腰束皮帶.
就像鱷駝毛衣腰束皮帶一樣.
所以這個樣子太厲害太震撼了.
所以印象很深刻.
而實際上很有趣的地方就是.
以色列人裡面.
他對以利亞的理解是什麼呢?.
他有沒有死呢?.
他是沒有死的.
為什麼呢?.
因為在以色列人的傳統裡面.
以利亞是被旋風接了升天的.
所以他沒有在世上被記錄了他的死亡.
換言之就是.
在那裡的人的理解就是他沒有死.
他去了上帝的國度.
被旋風接了升天.
所以在後來的一些先知書的描述裡面.
就是當上帝大而可畏的日子來的時候.
上帝拯救審判的日子來到的時候.
以利亞會再回來.

$^{361}$再出現.
所以這個形象一出現的時候.
以色列人會怎樣理解呢?.
上帝的拯救日子.
大而可畏的日子來了.
所以有些人就問.
施洗約翰.
就問他.
你是不是那位要來的?.
其實你是不是就是以利亞?.
在問他是不是這件事.
那以色列人就說.
我不是.
我不是這樣的人.
你可能會說.
那這個預備是不是.
預備我們回去馬上找一下有沒有命?.
找一件毛衣出來.
腰出皮底.
下個禮拜cosplay回來.
是不是這樣呢?.
我這裡就很想.
你看看剛才畫面裡面的另一半.
powerpoint裡面的另一半的畫面.
旁邊那裡你會看到有另一個畫面.
那畫面是什麼呢?.
那畫面呢.
是有一些人穿著一些比較華麗的衣服.
那是什麼地方來的呢?.
你猜猜.
是聖殿.
是聖殿的地方.
在聖殿裡面.
我也是叫AI畫的.
描述一下聖殿當時是什麼樣子.
大概是這樣.
下面那些祭司.
中間有一個人好像抱著口水肩在中間.
那個是什麼來的呢?.
那個人就是大祭司.

$^{401}$大祭司中間有一塊布.
那塊布上面是鑲了一些寶石.
有12粒.
為什麼呢?.
因為它代表著十二支派的獻祭.
它代表著十二支派獻祭的一些寶石鑲在身上.
穿著華麗的衣服.
然後就在聖殿獻祭.
你就會看到有一個很大的對比.
召者說.
如果約翰要召者說的時候.
為什麼不在那裡而是跑去曠野.
你明白嗎?.
召者說如果他要去事奉神.
為什麼你不在聖殿事奉神呢?.
為什麼你會跑去曠野事奉神呢?.
你召者說嘛.
你召者說跟著上帝嘛.
為什麼你去了那裡.
為什麼不在這裡.
或者可能你會說.
其實理由也很簡單.
因為進聖殿做事.
除了你是有社會地位.
因為你知道他們對祭司很有必公必敬.
因為有社會地位.
有專業人士.
除了這個之外.
要做得勝殿祭司.
其實還有另一個要求.
那個要求是什麼?.
祂有血統.
因為你做祭司不是任何人做的.
祭司是上帝揀選的族類.
才可以去做祭祀獻祭的工作.
是不是?.
你還沒有人.
換言之就是.
因為祂有沒有血統呢?.
你想做嗎?.

$^{441}$不一定到你做的.
你有血統的.
但是有趣的地方就是.
就是施洗約翰其實有血統的.
因為他爸爸是誰?.
他爸爸是撒迦利亞.
撒迦利亞有個兒子.
天使說這個兒子影響了很多人.
是不是這樣?.
所以撒迦利亞本來是祭司一員.
換言之就是其實.
施洗約翰.
他一來有一個服侍的心.
跟隨主的心.
跟從天國的心.
去照著祂的信仰連結著自己生命的心.
另一方面他也有一個血統.
可以在聖殿裡面事奉.
問題是.
他為什麼不在那裡呢?.
或者可能.
原因是在於當時的聖殿.
其實慢慢遠離了上帝的國度.
當時的聖殿.
它對聖潔律的要求.
是非常要求高的.
以致到.
例如聖潔律.
來到的時候.
安息日.
你能不能夠治病呢?.
不能治病.
因為你安息日治病.
就違反了聖潔律.
違反了上帝的律法.
好了.
但是.
安息日.
是為生命而設的.
上帝創造.

$^{481}$第七日.
竊了一切的功.
要他們欣賞上帝的創造.
但是安息日不能治病.
所以他病.
就病了.
他只能夠病.
呃這樣.
那你就會想.
似乎他們對律法的期望和想法.
慢慢取消了.
取代了上帝對生命的憐憫的心意.
聖殿慢慢偏離了上帝的國度.
再加上.
當時的聖殿.
是倒塌了再建的.
他怎樣再建?.
誰給錢再建?.
希律和羅嗎哪資源.
用各種資源來建.
當時的聖殿.
是有羅馬的首肯.
羅馬願意讓你去建.
然後有個叫希律.
連結著羅馬的勢力.
來一起去.
讓裡面找到一個空間.
羅馬當時也有要求.
不同類型的祭祀.
譬如君王的崇拜.
有時在裡面可能要放一隻金鷹.
在裡面作為他們的標誌.
代表著羅馬的帝國.
嘩.
當時施洗約翰.
他看著一個這樣的畫面的時候.
他就會問.
如果我要跟隨上帝的國度.
這個是不是我可以去的呢?.
然後他就選擇了不在那裡事奉.

$^{521}$然後他選擇了.
上帝的說話.
自由.
只要每一個願意悔改的.
上帝都會將他攬住.
所以他就去了曠野去招聚人.
換言之他放下了什麼呢?.
他放下了一個能夠.
逐步攀升的支力架構.
因為他本來是可以慢慢在聖殿裡攀升的.
他放下了這個支力架構.
他放下了那個安全感.
然後本來一個能夠.
給他一個這樣的攀升的階梯.
他選擇去了一個曠野.
沒有人招聚的地方.
這是第一個的預備.
第二個預備的第一個要點.
第二個預備的第二個要點在下一頁.
他吃了什麼呢?.
他要粗皮蛋.
一家的蝗蟲野蜜.
為了要去驗證一下.
蝗蟲是不是可以吃呢?.
我自己就驗證不到.
親口吃過.
但我就上網找.
找視頻.
找視頻看看蝗蟲是不是真的可以吃.
還有怎麼吃的.
終於就在我們.
什麼都會吃的祖國裡面.
我們就找到一條視頻.
這條視頻就是.
我不會播放的.
請大家放心.
因為很恐怖.
整個反嘔出來.
原來蝗蟲不是一定可以吃的.
因為蝗蟲是.

$^{561}$當牠變成黑色的時候是不可以吃的.
當牠聚集變成黃災黃禍的時候.
那時候變成黑色.
就不可以吃的.
所以這裡看到的畫面裡面全部都是綠色的.
牠在分散的時候就是綠色的蝗蟲.
這個時候牠的身體是沒有毒的.
當牠聚集黑色的時候就是有毒的.
所以有毒的時候你不可以吃.
綠色的蝗蟲他們就要用水洗.
你不可以拍片或用刀來處理.
在生的時候用水洗.
在生的時候就用尖刺穿著牠.
然後就燒牠.
我就發現了一個飲食的真理.
就是只要加了孜然粉就能吃了.
很震撼.
這個飲食的真理.
然後我就很驚訝.
在片子裡面那個燒的人真的吃給你看.
然後就這樣.
原來聖經描述這個的時候.
原來吃這件事.
大家看聖經要保持著好奇心.
到底是不是真的這麼說呢.
追查每一樣東西你會發現很多奇特的地方.
吃蝗蟲.
然後吃野蜜.
我小時候以為野蜜是花蜜.
但實際上野蜜就不是花蜜.
野蜜是圖上面的另一種果實.
那種果實叫做椰棗樹的果實.
椰棗.
椰棗是什麼來的呢.
在中東地方的沙漠偏在新疆的地方.
如果你在新疆的店鋪你會買到椰棗的乾.
椰棗其實像什麼呢.
像蜜棗.
是一隻長身大顆的蜜棗.
特別的地方就是它的糖分含量非常高.

$^{601}$超過一半的含量都是糖分.
換言之這兩樣東西是什麼來的呢.
就是糖分,熱量加上蛋白質.
健身房的朋友留意一下.
蛋白質.
這段時間有些年紀就常常被人提醒.
要小心肌少症,趁有機會盡量鍛鍊.
所以要盡量吸收.
所以我這段時間經常在早餐裡面.
只吃雜蛋和雞胸肉.
真的是非常努力的日子.
蛋白質.
哇.
加孜然粉就可以吃了.
這裡真的要再做一個比較.
你要記得.
施洗藥項它本來可以在哪裡呢.
聖殿.
聖殿吃什麼的.
吃獻祭的肉.
獻祭的東西.
獻祭要獻什麼.
有錢的人要獻什麼.
牛.
還有什麼.
羊.
有錢的人要獻羊.
沒錢的人要獻什麼.
百家.
還有.
斑鳩.
就是甲子和斑鳩.
知道是什麼嗎.
斑鳩就是大家在維多利亞公園旁邊.
大家應該很容易認得出的.
就是那些像甲子的東西.
但頸上一點點的就是斑鳩.
甲子和斑鳩.
加起來.
就有一個共通名詞.

$^{641}$叫做咕咕咕.
OK.
換句話說.
有錢的人.
沒錢的人.
他們在聖殿獻祭.
獻什麼祭.
獻烤肉.
他本來是吃烤肉的.
如果你有資源份你要撒到哪裡.
當然是撒那些.
我有資源份.
還是吃黃蟲鮮.
不會的.
你在一個選擇裡面.
你會見到.
約翰他選擇了離開聖殿之後.
他不單止沒有了一個上升的階梯.
安全感沒有了.
一個保障沒有了.
社會地位沒有了.
他並且飲食上他要靠野地的供應.
然後那些看到很好吃的東西沒有了.
他就要吃那些黃蟲野物.
其實這是一個很大的選擇.
是一個價值的選擇.
當他選擇了走這條路的時候.
蕉姐說的重量.
其實就在於這幾個選擇裡面.
跟你說一下一些事奉的路.
事奉的路有時候最難的地方就是.
當你身邊的人.
走另一條路的時候.
開始飛黃騰達.
很多青年人在他年輕的時候.
他想走一條事奉的路.
但是那些事奉的路通常是比較窮的.
我自己第一份工作其實就不是突破的.
我在另一個機構.
都是機構.

$^{681}$我第一份工作是在一個機構裡做.
是一個推動宣教的機構.
我還記得那天見工的時候是這樣的.
那個人跟我說.
我們會出這個價錢.
這個價錢是一個什麼價錢呢.
就是當時一個大學畢業生.
入息教書的時候的一半工資.
50\%.
現在可能大學畢業生2萬多元.
入職一些比較政府崗位.
大概就是一晚頭的今天.
我那時候再咳嗽.
接著我看著他.
我說 哦 這樣.
接著他看到我猶豫的眼神.
接著他下一步就跟我說.
哦 其實我們也有包飯的.
OK.
黃松野說.
包飯.
接著我看著他.
我真的不知道怎樣回答他.
我就說.
哦.
就是.
有上班就有飯吃.
他說 不是的.
接著後來我都.
後來輾轉反側.
接著我都覺得.
我還是做吧.
我未至於去到一個召者這麼大的感動.
不過都覺得.
年輕人為什麼不試一試呢.
那時候我就這樣去試一試.
接著下來的挑戰就是.
另一個挑戰就是.
接著呢.
就是畢業之後.

$^{721}$第一個同學聚會是最難的.
在那個同學聚會裡面.
大家還會問你.
誒 找到工作了嗎.
接著你要回答.
找到了.
接著呢.
接著下一個問題是什麼.
基層工.
接著我為了這個.
有機會要回答這個問題的時候.
我心裡掙扎到一個地步.
我沒有出現.
其實代價有時候就是這樣.
事奉一路的代價有時候就是.
會去到.
當人做多比較.
心裡面是很難不去比較的.
比較是一個我們很.
很怎麼說呢.
很本質的.
很本性的.
但是呢.
那些幽暗的使者就是.
在裡面拿你這些不安.
你要怎樣在上帝的說話裡面.
找到那個安定呢.
是很多生命裡面的體會.
耶穌裡面是怎樣提醒我們這些體會呢.
耶穌在《馬太福音》裡面提到.
你看看.
耶利的花.
你看看天上的飛鳥.
你看看他們.
不種也不收.
上帝尚且養活他.
他在說一件事.
他不是在說.
上帝尚且養他們飛黃騰達.
不是這樣 是不是.

$^{761}$上帝尚且養活他.
你走這段路.
我有時聽到這句說話的時候.
就想起我那時見工的時候.
看到那個人的樣子.
我們有包飯.
但是感恩的地方就是.
是的 我由那時開始的事奉.
接著轉去突破的事奉.
到今天.
真是一樣也不缺.
不是沒有掙扎.
不是沒有犧牲.
不是沒有心裡有苦毒.
完全不是 也有.
但是上帝在每一次這樣的事情出現的時候.
上帝都會用他很不同的方式來安慰.
鼓勵.
繼續走一條路.
弟子們.
來到的人.
他想見到些什麼呢.
來到的人他想見到的是.
見到將上帝的說話落在自己生命裡面的人.
他見到的是第二種.
要見到的是什麼呢.
就見到有人願意付出代價的人.
為這個信仰付出代價.
以致他當你見到.
有個人他為這個信仰付出代價.
為福音付出代價的時候.
他就知道福音的重量.
是比那個代價更重.
以致他願意付出這個代價來換福音.
他要見到福音有多重.
就看我們放在多重的代價裡面.
你說 很難啊.
每一個人可以付出的代價都不同的.
其實也不需要每一個人都立刻拋頭露骰而血.
我想說的就是.

$^{801}$包括後面的嬰兒.
帶回來容易嗎.
一點都不容易.
特別是我的小孩很大.
我兒子.
剛剛今早跟他說電話.
在台灣生活非常艱難.
現在回香港了.
他已經24歲了.
我女兒26歲了.
所以他們BB的時候.
要把他帶回來很難的.
很多預備的.
但是弟兄姊妹一直幫忙.
彼此補位.
以致我們跟小孩能走下去.
這些付出照顧彼此的支援.
其實就是我們付出的代價.
你願意帶小孩回來.
你願意彼此扶著.
幫弟兄姊妹帶小孩回來.
你願意為每一個流眼淚的人承載著他的眼淚.
這個就是你的代價.
求主幫我們.
願意我們每一個人都多走一步.
彼此承托.
以致人見到福音.
在這裡是有人願意付代價去換取的.
這是第二個我們描述的.
最後一個.
最後兩節.
他傳道說.
有一位在我以後來的能力比我更大.
就是我彎腰和他解鞋帶是不配的.
我用水來和你們施洗.
他們要用聖靈.
他們要用聖靈和你們施洗.
這兩句說話在.
在橫教會裡面是一個很容易被人用來描屬靈音運動的.
因為他在描述的時候好像有一個比較.

$^{841}$在這個比較之下.
我只是用水和你們施洗.
但這個人會用聖靈施洗.
他見到我彎腰和他解鞋帶是不配的.
因為在這個比較之下.
你就會覺得.
聖靈施洗更厲害.
所以當時後來有一些演繹的教導就是.
我們每一個人都要追求聖靈的洗.
我們每一個人追求這個特別的經歷.
但聖經不是一個為今天而寫的.
聖經在他那個時代有他要面對的處境和角度.
以至於這句說話其實出現的時候.
其實他不是真的在說.
我們今天在說的聖靈的洗.
那些靈音運動的現象.
不是這樣的.
他當時在說他當時的處境.
當時的處境是怎樣呢.
當時好像約翰那樣在曠野裡焦聚人的人.
不只是約翰一個.
當時也有一些人.
他是在焦聚革命軍.
在焦聚革命軍在曠野裡聚集.
那些人通常會說什麼呢.
有個人他就這樣說.
他說我是能夠令到約旦河分開.
我有異能能夠令到約旦河分開.
就好像摩西分開紅海一樣.
所以你來跟我吧.
有個人他就會這樣說.
我們可以圍著耶路撒冷城.
圍著祂.
祈禱圍著祂.
圍到祂塌下.
以至到耶路撒冷城一塌下的時候.
我們就攻進去.
就好像以前圍耶利哥城那樣.
當時有這樣的政治領袖 軍事領袖.
一些宣稱自己有異能的人.

$^{881}$就在那個時代裡焦聚著不同的人.
然後去做一些特別的事情.
但是 但是.
約翰在這裡說這個宣告.
你真的會覺得神經病.
如果我真的想做焦聚人的工作.
我就會強調自己有多厲害.
我有很多異能 很多東西.
但是約翰在這裡說什麼.
我約翰是用水來為你施洗.
我知道自己不是那個跟著要來的那個.
我知道自己不是.
不過我約翰是會用水來為你施洗.
水是什麼.
水就是你和我都可以用的東西.
水是你和我每一個普通人都可以使用的東西.
然後上帝使用的時候.
每一樣普通的事情.
都會在人的生命裡留下恩典的記號.
在這裡他留下了悔改的記號.
我約翰為你們施洗一個悔改的洗禮.
這個悔改的洗禮透過水的施洗.
他記住他生命要有一個轉向.
本來是這樣走.
本來是離棄上帝 離開上帝的道.
然後我轉向.
我要走回上帝的道路.
我約翰用的只是水.
你和我都可以用的水.
換言之就是其實每一個普通人.
每一個普通人.
好像你和我那樣的普通人.
我們不一定要有很多特別的能力.
才可以為主所用.
我自己在機構認識了很多很厲害的人.
你知道我們機構有歷史悠久.
有很多偉大的人.
偉大的名字.
很多這樣的.
你見到很多恩賜的人.

$^{921}$你見到他們一出來的時候就壓場.
我們小時候看著全部都.
嘩 真的仰望他們.
眨眨眼睛.
那時候去到突破路做的時候.
我基本上是覺得.
嘩 每一次見到他們的時候.
就好像見到有很多很厲害的.
很屬靈的人在你面前出現.
然後你就覺得嘩 聖光照耀.
我有時候都會見到這樣的場面畫面.
但是.
這麼多年的日子裡面.
給我最深刻的記憶和感受.
給我最深刻覺得上帝的祝福.
其實很多時候都是來自一些很.
叫做什麼呢.
前線甚至乎是很老人家.
好像不會站上台的人.
他甚至乎一站上台.
他就不好不好不好.
然後就走到一邊.
然後你叫他.
譬如見到很多年了.
大家都很喜歡你.
不如上來吧.
沒有沒有沒有.
他們通常會做些什麼呢.
很多時候他會無緣無故在你的位置.
留下一張紙條.
留下一塊餅.
有時候他們會是在.
我們一些很不容易的地方的時候.
需要有人幫忙的時候.
他第一個衝出來.
有什麼要幫忙我來.
有時候我看著他們的時候.
其實他們是很忠心.
他們的心很純.
他們的心只是在想.

$^{961}$有什麼主上帝你要我做.
我就馬上去做.
我是每一次見到他們的時候.
我就覺得.
神真是使用每一個人的.
神真是使用每一個人.
不單止在機構裡面是這樣.
我在自己教會也是這樣.
有一段時間我是非常艱難的日子.
回到教會的時候.
我是不跟別人打招呼.
最好你見不到我.
最好就走去一些燈很暗的地方的角落.
你最好見不到我.
我來到的時候就會坐在那裡.
閉上眼睛.
我閉上眼睛你就覺得.
也沒什麼大不了他.
做一個勢出來.
你就不會走過來.
你明白那種感覺.
我不知道你有沒有試過.
我不是說任何一個人.
我不認識你們的.
我說我自己.
我是這樣的.
我是這樣去回教會.
有一段日子很艱難的時候.
就是不想有人提起任何事情.
不想有人關心.
什麼都不想.
不想有人問起我的近況.
國內都不想.
然後那段日子.
其實我大概都感受到自己.
可能有些depress的狀態.
然後有一天就是這樣.
就是有一位.
那時候我突然感覺到有人拍我的肩膀.
就是跟我名字一樣.

$^{1001}$拍肩.
拍我的肩膀.
OK.
突然拍一拍.
那些黑漆漆的地方.
突然有人拍肩膀.
用手語.
然後那一刻的感受就是.
他見到我.
其實誰見不到你.
他見到我.
他並且不單見到我.
並且沒有因為我的氣場.
他越過了我的氣場.
然後來給一個鼓勵和安慰我.
有人拍你.
那你當然要看看是什麼人.
我看了一下.
然後.
他是教會裡面的一個叔叔.
一個大我十年左右的叔叔.
那時候已經是叔叔.
然後他是一個.
叫做矮矮細細的一個男人.
不善說話.
沒怎麼說話的.
就是教會一有什麼.
譬如擔椅子.
搬桌子.
教會搬辦公室 收拾東西.
他就會出現.
是一個這樣的叔叔.
掃地 洗廁所 他都會出現.
然後他每一次這樣.
是一個漫漫無聲的一個人.
然後他上台說話的時候.
很固執.
明知其實他不習慣面對很多人的叔叔.
然後他拍一拍我的時候.
他也沒有跟我說話.

$^{1041}$他只是拍一拍我.
拍一拍 按一按 然後就離開了.
然後我那一刻的感覺就是.
這真是一個最溫柔的安慰.
上帝使用了他做給我最溫柔的安慰.
他鼓勵了我當時在很幽暗的日子裡.
我以為這個世界不會有光.
我以為這個世界就是會連在一起.
我以為這個世界不會有未來.
但是上帝使用他這個拍一拍.
其實拍一拍.
你跟我.
不用特別的恩賜 你都做到的.
其實拍一拍.
一句慰問.
不用特別的恩賜 你跟我都做到.
就好像約翰使用的水一樣.
是你跟我都可以用到去祝福人.
你慰問到那一刻 你先去祈禱紀念.
然後你再做.
那件事就算多普通.
上帝都會在適合的時間.
去使用成為人的祝福記號.
弟姐妹 我有時候都會想起.
去到不同的教會去講道的時候.
其實最記掛的是什麼呢?.
我最記掛的就是去到每一個地方.
都希望祝福那個教會.
能夠彼此承托.
讓生命來到中間的時候.
無論他這次是第一次來到.
還是已經來過幾次.
又或者在尋求一些東西.
不知道在尋求什麼.
但是在信仰上想要有些經歷的人.
他來到中間的時候.
我們其實很希望每一個地方.
每一間教會都能夠彼此承托.
要人來到中間的時候.
他不是好像在筲箕之中將水流走.

$^{1081}$而是每一個來到中間的時候.
他在裡面見到.
好像圖金一樣.
見到晶金在中間.
他見到生命是被承載.
在裡面每一個破碎的心.
他在那裡被上帝觸摸.
上帝的觸摸.
原來是可以透過我們每一個人的付出.
然後將那些位置去填滿去補充.
求主幫我們.
讓他的聖靈在我們中間運行.
有時候我們說聖靈在我們中間.
什麼叫在我們中間呢?.
有時候我們會覺得.
聖靈就在我們中間的時候.
好像在我們心裡面的中間.
但其實可能是.
如果按著潘福華.
因為我上學的時候.
最後的畢業論文之一.
就是寫了潘福華.
他講團體生活的時候.
我就引用了一些經文.
引用了一些說話.
引用了一些他另一些的著作.
裡面提到基督在我們中間的時候.
他用的中間的字是哪個字呢?.
他不是中心的中間.
他是在我們之間.
在我們之間有上帝填滿.
有基督填滿.
以至於來到的生命.
他不會在我們中間溜走.
而是在裡面中間經歷神的幫助,拯救.
求主幫我們.
讓聖靈在我們中間.
我們求主幫助我們一起祈禱好嗎?.
上帝求你幫我們.
主願意上帝你的說話.

$^{1121}$上帝你透過.
施洗若寒的經歷.
你鼓勵我們.
鼓勵我們的心在你的裡面.
抓緊你.
抓住你.
讓上帝的說話.
在我們的生命留下痕跡.
以至於我們也都在.
我們所經歷過的每一個地方.
每一個群體都留下痕跡.
在這個時代裡面留下痕跡.
讓上帝的恩典留下痕跡.
求主你幫助我們.
願意這些說話去激動我們的時候.
我們不只是聽完之後坐在這裡.
是聽完之後行動.
願意你的說話來到我們中間的時候.
我們有一個重量可以承載.
我們願意付出代價去將它換取福音.
求主你都幫助我們.
不要看我們自己太小.
覺得我們自己手上只有水.
只有很普通的東西.
只有很多沒什麼恩不恩切的東西.
但主讓每一件事情.
在你的手上都被你使用.
祝福我們的生命.
祝福我們的生命在你裡面.
讓你幫助我們.
讓我們每一樣細微的事物.
都有上帝祝福的恩典.
以致我們彼此連結.
彼此在這裡中間的時候.
基督在我們中間.
承托著每一個來到的生命.
讓每一個生命裡面.
找到上帝你恩典的意好.
找到你的愛.
幫助帶領我們.

$^{1161}$我們求你在我們同在.
讓我們得到祝福.
我們這樣好好交託.
奉主耶穌的生命祈求.
阿們.
\newpage



\section{啟示錄 2:1-7}
\label{sec:032460}
\textbf{教會如何在黑暗中作光}~~~~劉沙崙牧師
\newline
\newline
連結: \href{http://www.church.com.hk/acms/content.asp?site=cdc\&op=show\&type=product\&code=032460\&layout=sermon}{\texttt{032460}} ~~~~ 語音日期: 2026-07-12
\newline
\newline
\hyperref[sec:032459]{\small{< < < PREV SERMON < < <}}
~
\hyperref[sec:index_chronic]{\small{[返順時目]}}
~
\hyperref[sec:032461]{\small{> > > NEXT SERMON > > >}}
\newline
\newline
啟示錄 2:1-7
\newline
\begin{longtable}{cl}
\hline
\hline
章節 & 經文 (和合本修訂版)\\
\hline
2:1 & \begin{tabularx}{0.7\textwidth}{X} 「你要寫信給以弗所教會的使者，說：『那右手拿著七顆星，在七個金燈臺中間行走的這樣說： \end{tabularx} \\ \\ \relax
2:2 & \begin{tabularx}{0.7\textwidth}{X} 我知道你的行為、勞碌、忍耐，也知道你不容忍惡人。你也曾察驗那自稱為使徒卻不是使徒的，看出他們是假的。 \end{tabularx} \\ \\ \relax
2:3 & \begin{tabularx}{0.7\textwidth}{X} 你能忍耐，曾為我的名勞苦而不困倦。 \end{tabularx} \\ \\ \relax
2:4 & \begin{tabularx}{0.7\textwidth}{X} 然而，有一件事我要責備你，就是你把起初的愛心拋棄了。 \end{tabularx} \\ \\ \relax
2:5 & \begin{tabularx}{0.7\textwidth}{X} 所以你要回想你是從哪裡墜落的，並且要悔改，做起初所做的工作。你若不悔改，我要到你那裡去，把你的燈臺從原處挪去。 \end{tabularx} \\ \\ \relax
2:6 & \begin{tabularx}{0.7\textwidth}{X} 然而你還有一件可取的事，就是你恨惡尼哥拉派的行為，這種行為也是我所恨惡的。 \end{tabularx} \\ \\ \relax
2:7 & \begin{tabularx}{0.7\textwidth}{X} 凡有耳朵的都應當聽聖靈向眾教會所說的話。得勝的，我必將神樂園中生命樹的果子賜給他吃。』」 \end{tabularx} \\ \\
[1ex]
\hline
\hline
\end{longtable}
$^{1}$當我來到禱告有什麼訊息和你分享的時候.
我特別和大家分享一個題目,一個訊息.
就是教會怎樣能夠在黑暗裡面成為光.
你看到有個燈台.
其實聖經裡面耶穌說我們是這個世界的燈.
是這個世界的光.
要照亮黑暗的地方.
我們知道為什麼我們教會存在的目的.
除了我們很開心我們可以一齊回到去敬拜.
去有團契,有交通的時間.
但其實教會最重要的責任或者我們的功能是什麼呢.
就是要在這個世界上成為光.
大家同不同意.
雖然我們看這個世界似乎越來越黑暗.
似乎越來越多問題.
但是聖經神卻提醒我們.
教會要在黑暗裡面作光.
我相信根據聖經的預言.
雖然耶穌回來之前會越來越多問題,越來越多艱難.
但是同時間教會都會越來越復興.
都會越來越多人在教會裡面真的找到神的愛.
找到真理,找到生命.
所以今天我也會特別從一個書信來分享.
大家都知道《新約聖經》除了四福音,《使徒行傳》等等.
很多都是書信來的.
其中誰寫得最多書信?.
就是保羅.
除了保羅之外.
當然我們都知道有雅各,約翰,彼得,猶大等等.
他們都有寫書信.
但是原來其實耶穌都有寫信的.
那耶穌的信是記載在哪一本書呢?.
有些如果有看過的就是在《啟示錄》.
不知道誰有讀過《啟示錄》的舉手看看.
誰明白《啟示錄》說什麼的舉手看看.
沒有人敢舉手了.
馬上明白了 厲害啊.
其實《啟示錄》好像很難明白.
但是其實又不是真的那麼難明白.
簡單地說我們要知道《啟示錄》寫的對象是什麼.

$^{41}$是寫給教會的.
為什麼呢?.
一開始《啟示錄》已經說明了要給教會.
特別是一開始第二,第三章.
已經有七封信給當時的七間教會.
其實《啟示錄》從一個宏觀的角度看.
不是那麼難明白.
簡單地說就是耶穌得勝.
魔鬼被打敗.
說完了.
是不是?.
就是這樣.
然後耶穌會回來.
然後我們會與基督一同作王.
是不是?.
就會有一個完全不同的世界.
一個新天新地.
簡單地說就是這樣.
當然細節.
如果我們要研究.
究竟敵基督是誰.
666是什麼.
災前災後.
災中被提.
這些就要爭辯很久了.
但是總之我們知道大前提很簡單.
就是耶穌得勝.
Amen.
魔鬼完全被打敗.
教會要與神一同作王.
是的.
當你從這個大前提去看《啟示錄》.
就不會覺得懶惰了.
就不會覺得災難很恐怖了.
只不過是過程的一個部分.
耶穌會得勝的.
所以今天我會特別從耶穌寫給以弗所的教會.
這一封信.
來看看神如何提醒我們.
教會如何在黑暗裡作光.

$^{81}$在講這封信之前.
也講一下以弗所的一些歷史背景.
因為其實以弗所這個城市.
是有點像香港的.
你看看.
以弗所當時是一個商業中心.
又靠近海.
又是一個很重要的河出口的地方.
所以它是很繁榮的.
也被稱為亞細亞的大都會.
香港也被稱為亞洲的大都會.
大家知不知道亞細亞其實是哪裡.
就是今天的土耳其.
所以我們看到七個教會.
都是位於今天的土耳其.
就是那個地區.
而以弗所很特別的.
就是它政治上也是特區.
有點像香港.
我們也是特區.
因為當時雖然是羅馬帝國統治.
但是它是有自治權的.
無論政治還是經濟上.
都是相當自由和繁榮的.
我們可以看看這個地圖.
我們可以看到.
耶穌對七個教會來說話.
七個教會的地理位置.
當時約翰在摩海島被放逐.
我們知道其實.
啟示錄是寫於很後的日子.
就是大約主後90年左右.
約翰也是唯一一個使徒.
他是真的可以老死的.
其他都是在很早年的時候.
就已經為主去殉道.
他當時在拔摩海島雖然被放逐.
但是耶穌向他啟示.
就寫了啟示錄.
當中七個教會.

$^{121}$我們看到以弗所其實是最近拔摩海島的.
以弗所.
接著有善美娜.
別加摩.
推亞推拉.
撒迪.
菲拉鐵菲.
和魯迪加.
而在這七個教會裡面.
其實五間教會的信.
耶穌都是有讚也有責備.
但是有兩間是沒有被責備的.
就是善美娜和菲拉鐵菲.
為什麼呢.
因為這兩間教會都是大大受逼迫的.
接著我們繼續看以弗所的地方.
當時我們要知道這個城市雖然很繁榮.
經濟很繁榮.
但是他們整個的城市.
其實是敬拜一個叫做阿底米的女神.
她的廟宇是很大的.
到今時今日都是被稱為世界七大奇景之一.
這個阿底米的女神.
她是很醜的.
讓大家看看她的樣子.
是這樣的.
很多的乳房.
這個就是她的女神的廟宇.
當時是很大的.
裡面是容納了一千多個所謂的女祭師.
但是其實是做淫亂的事情.
他們敬拜的時候都會做淫亂的.
所以當時的以弗所其實是充滿著淫亂的.
而整個城市的經濟發展.
其實都是圍繞著這個廟宇的.
為什麼這樣說呢.
當時有些人談生意.
是一定會在廟宇那裡談的.
也都會有淫亂的事情發生.
大家明白我說什麼嗎.

$^{161}$有些類似現代的夜總會.
在那裡又談生意.
不過又做很多不應該做的事情.
當時的以弗所就是這樣的.
很淫亂.
而且當時以弗所的很多人.
都是靠做偶像來發財.
因為那些人是很迷信的.
家家戶戶都會擺很多偶像.
所以他們是靠這件事來發財的.
接著我們看到.
雖然在這樣的環境裡面.
但是神就差派了保羅.
去到以弗所那裡來建立.
當時以弗所.
保羅是第三次宣教的時候.
當時和阿居拉和伯基拉.
是兩夫婦來的.
就一起來到以弗所這個地方.
來傳福音.
建立教會.
以弗所教會特別幸福的.
為什麼呢.
因為保羅停留在那裡有三年.
其他的教會都可能是停留比較短的時間.
建立了.
接著幾個月後興起人.
他就去下一個地方了.
但是以弗所他停留了三年.
而且聖經裡面有一本書叫以弗所書.
裡面很精彩的.
講教會是怎樣去彰顯神的榮耀.
又講到應該怎樣彼此相愛.
講到夫妻.
又講到屬靈爭戰.
一個很傳備的教導.
所以我們知道以弗所當時是這樣.
而當時也因為傳福音的時候.
是滿有神蹟騎士.
我們看回《賜徒行傳》第19章.

$^{201}$記載當時大家有沒有讀過那段經文.
就是保羅因為太多人想他為他祈禱了.
沒有那麼多時間.
所以他就祈禱在那些毛巾上面.
接著就派給那些人.
那些人都得醫治的.
就是發生在以弗所.
也都有一些本來行巫術的人.
都得到釋放.
所以當時是很多的神蹟騎士發生.
以致於本來敬拜亞底米女神的人.
都轉向敬拜耶穌.
但是也都因為這樣.
其實保羅當時也是面對逼迫的.
我們看回《賜徒行傳》第19章.
就說當時就是因為教會的影響力大到.
搞到整個城市都騷動起來.
特別是那些賣偶像的人的生意受影響.
沒有人再賣偶像了.
所以就想藉著騷動來趕走這些基督徒.
但是結果成不成功.
當然是不成功.
這個教會繼續成為當時以弗所的一個光.
看到是非常復興的一間教會.
接著保羅離開了以弗所之後.
就是誰去牧養的呢.
就是提摩泰了.
哇 他的屬靈仔.
所以這間教會是不是很厲害.
由保羅親自建立.
提摩泰牧養.
接著後來也都使徒約翰.
也都在晚年的時候去牧養這間教會.
有些傳統的說法都說.
耶穌的媽媽都是回這間教會的.
為什麼呢.
因為耶穌當祂釘十字架的時候.
祂就叫約翰要照顧祂媽媽的.
大家記不記得.
所以很多人就說.

$^{241}$祂其實一直跟著約翰.
特別是去到晚年的時候.
都是在以弗所那裡回這間教會.
所以想一想這間教會.
有耶穌的媽媽在那裡.
聖母瑪利亞.
有保羅.
有提摩泰.
接著又有約翰.
嘩 這間教會是不是很棒.
是一間滿有神的同在的教會.
而且當時候這間教會.
也是在主後70年之後.
成為了最重要的教會.
之前耶路撒冷也是很重要的.
但是主後70年我們知道.
當時候耶路撒冷的聖殿被毀.
那些基督徒都被趕散.
來到不同的地方.
所以以弗所可以說.
取代了成為當時候.
這麼多間教會的大佬.
一哥的位置.
所以我們看以弗所教會.
我們明白當時候.
耶穌寫這封信給以弗所教會.
是一個非常復興.
反轉整個城市.
也有真理.
有恩賜.
有教導的一間教會.
嘩 應該讚到飛起.
但是我們看下去.
原來耶穌對這間教會也有提醒.
好 我們一起來讀.
《啟示錄》二章一到七節.
一二三 「你要寫信給教會的使者說.
那右手拿著七星.
在七個金燈台中間行走的說.
我知道你的行為勞碌 忍耐.

$^{281}$也知道你不能容隱惡人.
你也曾試驗那自稱為使徒.
卻不是使徒的.
看出他們是假的來.
你也能忍耐.
曾為我的名勞苦.
並不罰權.
然而有一件事我要責備你.
就是你把起初的愛心離棄了.
所以應當回想.
你是從哪裡墜落的.
並要悔改.
行起初所行的事.
你若不悔改.
我就臨到你那裡.
把你的燈台從原處挪去.
然而你還有一件可取的事.
就是你恨惡尼哥拉一黨人的行為.
這也是我所恨惡的.
聖靈向宗教會所說的話.
凡有耳的就應當聽.
得勝的.
我必將神樂園中生命樹的果子賜給他人」.
Amen.
雖然這封信是耶穌寫給二佛所.
但這裡告訴我們.
這封信也是寫給我們.
因為他說.
聖靈向宗教會所說的話.
凡有耳的就應當聽.
大家有沒有耳朵?.
我們有耳朵.
我們就要聽神對他的教會所說的話.
我相信.
其實神可能對每一間教會都有一封信.
不知道祂對同福堂的信是怎樣呢?.
但今日願意我們從二佛所教會的信裡.
我們去學習.
教會怎樣在黑暗裡作光.
首先在這封信裡.

$^{321}$我們看到.
耶穌告訴教會.
其實祂是掌管和恩高教會的那一位.
為什麼教會能夠在黑暗裡作光呢?.
不是因為我們每個人都很厲害.
每個人都博士畢業.
每個人都很有成就.
或者我們每個人都完美.
不是.
一定要知道是因為耶穌恩高祂的教會.
其實當我們有時候去看教會的時候.
可能都會覺得教會都不是那麼完美.
每個教會都有它的問題.
如果一個人說我要找一間完美的教會.
只要你去到那間完美的教會.
那間教會已經不完美了.
明白我說什麼嗎?.
所以教會始終都會有問題.
因為有人的地方就會有不同的問題.
但是你看到.
耶穌祂知道人是有問題.
但是祂的選擇要透過教會.
去讓世人看見祂的榮耀.
祂稱呼我們為祂的身體.
祂是頭.
我們是祂的身體.
祂選擇使用我們一群不完美的人.
但是我們在不完美當中.
我們學習去愛神和愛人.
也學習怎樣去彰顯神的榮耀.
所以我們看第一節.
二章第一節.
這樣說.
右手拿著七星.
在七個金燈台中間行走的說.
啟示錄七個教會.
這些信是很特別的.
耶穌一開始第一句.
就是用一句句子去形容祂自己.
每一封信都有不同的形容.

$^{361}$而對以弗所教會.
祂就形容祂自己是什麼呢.
就是你想想祂的右手.
拿著七顆星.
有七顆星在祂的手上.
然後就走在七個金燈台中間.
就是這幅圖畫.
所以啟示錄有時候覺得很難明白.
因為有很多這樣的預表.
但是其實所有預表都有解釋的.
所以又不是真的那麼難明白.
這裡的解釋是什麼呢.
原來第一章已經解釋了.
因為當時約翰看到耶穌這樣的形象.
然後耶穌已經說.
你所看見在我右手中的七星.
和七個金燈台的奧秘是什麼呢.
七個星就是代表七個教會的使者.
七個燈台就是代表七間教會.
所以其實當我們明白預表.
又不是那麼難明白.
首先說一下右手是什麼意思.
其實聖經裡面.
往往右手都是代表神的保護.
神的能力.
祝福神的喜悅.
我們看到很多時候.
爸爸去祝福兒子.
例如亞伯拉罕祝福以撒.
去到雅各祝福他的兒子的時候.
都是用右手來按手的.
而我兒子的名字叫做Benjamin.
即是辯雅文.
即是雅各最小的兒子.
為什麼會改這個名呢.
因為我當時聽了一篇道.
那個牧師就說.
原來辯雅文的意思就是右手之子.
右手就是代表祝福神的保護等等.
所以我就說好名.

$^{401}$我就改我的兒子叫做Benjamin.
就是這樣的.
有沒有人叫阿Ben的.
是好名來的.
即是代表你是神的右手之子.
蒙福的一個人來的.
神所保護的 所保守的.
接著七星.
就是這裡很清楚的說.
就是代表教會的使者.
教會的使者.
這個翻譯.
你可以翻譯做天使.
或者翻譯做messenger.
即是傳遞訊息的.
首先我都相信.
這個教會.
每一間教會.
都真的有天使和我們一起的.
去保護我們.
去為我們征戰.
雖然我們見不到他們.
都不需要見到他們.
但是當我們來到去禱告敬拜的時候.
其實祂在這裡為我們征戰.
當我們剛才一起敬拜的時候.
其實我們看不到天空發生什麼事.
但是祂在這裡為我們征戰.
來捆綁受敵的作為.
以致神的聖靈.
能夠在我們當中運行.
所以我們相信有天使在我們當中.
另外亦都說到.
這個使者就是傳遞訊息的.
其實所說到的.
就是每一個願意被神使用的人.
包括教會的領袖.
但領袖不只是牧師.
領袖是每一個去服侍的人.
無論我們的小組長.

$^{441}$代討者.
或者我們有很多不同的服侍.
其實每一個願意去服侍別人的人.
都是神所恩告的.
都是在神的手中.
所以我們知道教會的盼望.
是在於什麼呢.
是有神的同在.
還有這裡說.
耶穌祂是走在七個金燈台中間.
就是走在教會的當中.
你知道今天祂就在我們當中.
聖經很清楚說.
只要兩三個人聚集.
祂就在我們當中.
祂就與我們同在.
所以我們知道原來神就是這麼愛我們.
藉著祂的聖靈與我們同在.
所以我們要相信.
我們每一個其實都是被神恩告的.
接著在二弗所書.
這個就是保羅寫給教會的信.
不過也是被聖靈感動的.
裡面保羅其實有一個禱告.
這個禱告都是我們很需要去學習的.
也相信這個禱告是很重要的.
其實二弗所書前後一共有兩個禱告.
一個就是第一章 第二個就是第三章.
第三章主要的禱告就是.
為了我們能夠一同明白基督的愛.
是何等的強乎高深來禱告.
而第一章就是禱告.
我們要明白神如何與我們同在.
那個能力有多大.
我們一起來讀二弗所書一章.
十七到十九節.
一二三.
我求我們主耶穌基督的上帝.
的父把持智慧和啟示的靈.
給你們.

$^{481}$好使你們更深地認識他.
我也求上帝照亮你們心中的眼睛.
使你們知道他的呼召.
給你們帶來了何等的盼望.
他應許給眾聖徒的產業.
有何等豐富的榮耀.
並且他在我們這些信的人身上.
所運行的能力是何等浩大.
教會是基督的身體.
蘊含著那無所不在.
充滿萬有者的豐盛.
最後這一句你記住.
教會是基督的身體.
裡面是有什麼.
無所不在.
充滿著萬有者的豐盛.
原來神是怎麼看教會的呢.
怎麼看我們每一個的呢.
願意我們都去求主.
開啟我們的眼睛.
看到我們除了神很想我們.
一起去更深地認識他之外.
他也要照亮我們的眼睛.
讓我們知道.
神是呼召和揀選我們每一個的.
Amen.
你知不知道神是揀選了你的.
是啊.
剛才說到我爸爸.
其實應該是.
不是51年.
是57年.
在香港.
他1969年來香港的.
已經差不多57年了.
為什麼他會來香港呢.
為什麼來做神要他做的.
就是他聽到神的呼召.
在他16歲信主的時候.
神已經呼召他.

$^{521}$要去中國傳福音.
其實我爸爸如果認識他.
會知道他是一個極度害羞的人.
極端內向和害羞的人.
但是神可以使用一個這樣的人.
所有人都能使用的.
我爸爸也不是完美.
不過他有一樣東西很好.
就是他是一個禱告的人.
也是這麼多年我學習.
因為他知道自己不行.
所以他每一天都花很多時間.
到今時今日.
我知道他每天都花3個小時.
來祈禱.
所以其實很多的事情.
我們見到能夠透過他的服事.
來成就.
都是完全將榮耀歸給神的.
但是神也呼召我們.
每個人我們最低限度.
可以做的都是祈禱.
你都可以向我們身邊的人傳福音.
我自己也是15歲的時候.
神呼召我.
很簡單.
他呼召我做的事情就是.
要成為中國人的祝福.
八個字.
這個就是神給我的呼召.
所以我現在.
總之有中國人的地方我就去.
現在的服事.
我早期都是多些教會的牧羊.
但是現在我和我先生.
都是多些會去到國內.
或者去到不同的地方.
本來我也是想香港大陸兩邊.
誰知道中國人全世界都去的.
所以我都很感恩.

$^{561}$有機會可以去到不同的地方.
去服事不同的華人.
這個就是神給我的呼召.
那你們呢.
你也有的.
最低限度.
這些不需要很大的啟示.
最低限度.
你都要向你的家人傳福音.
你所認識的人傳福音.
去愛你身邊的人.
好好地忠心去做.
實要你做的事.
我相信有一天.
我們很想遇到主的時候.
都是主祂說.
我們是祂又良善又忠心的僕人.
神給我們有多少的恩賜.
我們好好地去用.
而且這裡也說到.
神在我們這些信的人身上.
所運行的能力.
是何等的浩大.
你知不知道.
不是因為我有能力.
但是聖靈在我裡面.
是大有能力的.
我們的禱告是大有能力的.
所以剛才聽到.
你們有為宣教去禱告.
其實真的很需要更加多的禱告.
因為禱告是大有能力的.
二人祈禱所發的力量.
是大有功效的.
我們為人禱告.
我們傳福音的時候.
大家都很渴望見到神蹟奇事.
神蹟奇事什麼時候會發生.
就是當我們去傳福音.
為人禱告的時候.

$^{601}$我們就會見到神應運我們的禱告.
所以首先我們要知道.
為什麼教會能夠在黑暗裡做光.
是因為神與我們同在.
因為聖靈在我們裡面.
聖靈住在哪裡.
住在哪裡.
在我們裡面.
我們的身體就是聖靈的電.
祂是與我們同在的.
好像剛才的詩歌.
無論我們去到哪裡.
祂的靈都是與我們同在的.
所以願意我們常常提醒我們自己.
我們不是驕傲.
也不是自卑.
其實我們要謙卑.
謙卑是什麼.
就是少一點看自己.
少一點看自己為之謙卑.
多一點去看神的作為.
而是願意被神去使用.
因為我們知道.
當我們願意去的時候.
神的榮耀會在我們生命裡面彰顯.
為什麼教會能夠在黑暗裡成為光.
耶穌說是因為他們的好行為.
因為我們的好行為.
是能夠讓世人看見我們的不一樣.
所以耶穌在對爾伐索教會的信裡面說.
我知道你的行為牢固忍耐.
也不會容隱惡人.
也能夠去分辨哪些是假的使徒.
也能夠忍耐為神的明路虎.
他特別提到有一件你做得好的事.
就是你恨惡尼加拉一黨人的行為.
這個也是我所恨惡的.
所以我們看到爾伐索教會.
當時也是一間.
相信是很多士工.

$^{641}$這裡說得很勤力.
也很努力去為主去作供.
也能夠去分辨哪些是神喜悅的.
尼加拉黨是什麼呢.
其實有不同的說法.
但是我個人比較傾向相信的說法.
就是來自於使徒約翰的弟子.
這個人叫做愛淫瘤.
他是在第二世紀的人.
當時他寫了一本叫做《反二端》的書.
裡面他很清楚地說.
尼加拉一黨就是使徒行傳裡面所記載.
七個執事之一的尼加拉的跟隨者.
他們是做什麼的呢.
就是當時他們過著毫無節制的生活.
教導人可以行淫.
和吃拜偶像之物.
認為這些事都是無所謂的.
原來當時的尼加拉.
本來是七個執事.
大家有沒有看過《啟示錄》.
有七個執事的.
其中有兩個很出名的.
一個叫做司提反.
就是第一個為主殉道的.
另一個叫做肥力.
說一個很厲害的.
主要將他從一個地方提到第二個地方.
和安提亞伯太監傳福音那個.
但是原來其中有一個執事.
他後來出了問題.
他後來教導的是什麼呢.
恩順清義.
你只要恩順清義是對的.
我們都信恩順清義.
但是他恩順清義就說.
因為恩順清義.
所以你就是自由的了.
因為你的靈魂無論你做什麼都好.
都不會影響你的救恩的.

$^{681}$你無論怎樣犯罪.
怎樣淫亂都好.
都沒有問題的.
接著他就開始這樣去教導.
其實這個正正.
就是無論是保羅所說.
或者在希伯來書裡面說到.
就是用基督的恩典.
當作放縱情慾的藉口.
所以當時教會原來都出現這個問題.
特別對於以弗所當時的信徒來說.
其實這種教導很吸引的.
因為剛才我不是說.
他們要談生意.
他們很多時候工作或者各樣.
都要去廟裡面的.
還有大部分的食物都是拜過偶像的.
拜過偶像.
你去街市買菜 很多東西都買不到的.
所以其實對於當時的信徒來說.
這個教導太好了.
我做什麼都行 無所謂.
總之我得救就行.
總之我信耶穌就行.
所以其實是一個很大的破壞.
但我們看到當時的以弗所教會.
他們是怎樣的.
他們恨惡尼哥拉一黨的人的教訓.
其實這個都是提醒我們今時今日的教會.
說真的.
有沒有尼哥拉一黨的人的教訓在我們當中呢.
可能我相信同福堂沒有的.
但是我會看 特別一些西方的教會.
很多的妥協.
有些很大的宗派.
現在都在教會裡面掛這個彩虹旗.
我都很震驚的.
我去年去到英國.
去了一間聖公會的教會.
香港的聖公會很好的.

$^{721}$不要誤會.
是很好的.
但是英國的聖公會竟然是掛彩虹旗.
我看到都很震驚.
接著原來有某些的宗派.
包括英國的聖公會.
現在無論你是變性人也好.
你是同性戀也好.
你是又離婚又再結婚又離婚又再結婚.
那種牧師也好.
都是可以繼續做牧師的.
大家明白我說什麼嗎.
完全是和這個世界妥協的.
沒錯.
我們相信.
我們要傳福音給所有的人.
包括所謂的這些性小眾.
我們都要愛他們.
傳福音給他們.
但是每一個人都要悔改的.
罪就是罪.
不能夠來妥協.
但是我們在今時今日.
特別是歐美的教會.
其實是很多的問題就是這樣.
特別我知道有些最大的教會.
就是常常都會說.
總之最重要是你成功.
你回教會就是怎樣蒙福.
怎樣恩典.
接著就是不需要認罪的.
因為甚至有一些教導.
我聽到都很震驚.
就是說當你信耶穌那一刻認罪悔改.
就已經夠了.
之後如果你再認罪.
就是你不夠信心.
耶穌已經赦免了你所有的罪.
因為他說你信耶穌那一刻.
信耶穌那一刻悔改.

$^{761}$就已經是過去.
現在將來的罪都赦免了.
所以從此以後.
不需要再認罪悔改.
你聽到這樣的說法.
不知道你會怎樣看呢.
其實我們大家都知道.
如果你信了耶穌.
代不代表你就不再犯罪.
不會.
不過我們最大的分別.
你怎樣知道你是重新得救呢.
就是你犯罪之後.
你會很辛苦的.
你會知道要悔改的.
你要回到神的裡面.
去呼求祂的恩典.
呼求耶穌的寶血遮蓋.
一個真的有聖靈在裡面的人.
是不能夠不停犯罪而沒有感覺的.
你明白我說什麼嗎.
所以當時這個教導.
就是要破壞教會的聖潔.
因為我們知道其實耶穌說回來.
祂是迎接一個聖潔的新婦.
所以我們都真是.
耶穌去提醒.
我們是世上的光.
教會要和這個世界不一樣.
我知道根據美國的統計.
基督徒的離婚率.
和非基督徒是沒有分別的.
這個其實很可悲的.
很多人可能回教會.
但是有沒有重新得救.
有沒有經歷神.
這個都不知道.
所以我們真正信耶穌的人.
我們的生命會改變的.
沒錯.

$^{801}$可能我們都有破碎的經歷.
可能都曾經有破碎的婚姻.
但是在神的裡面.
神是能夠醫治的.
我們會呼求神的醫治.
我們會追求聖潔的.
接著.
講完有兩個很重要的.
就是教會怎樣來做光.
就是因為耶穌恩高教會.
與教會同在.
也都是教會的好行為.
但是對這所謂的教會.
耶穌都有一個提醒.
祂說如果失去了起初的愛.
教會就不能夠在黑暗裡面去做光.
所以原來很多東西都很重要.
我們的神的同在 恩高 福士.
很重要.
有好行為很重要.
聖潔很重要.
不過有一件事是不可少的.
就是起初的愛.
所以這節經文我們再讀.
二章四到五節.
一二三.
然而有一件事我要責備你.
就是你把起初的愛心離棄了.
所以應當回想.
你是從那裡墜落的.
並要悔改.
行起初所行的事.
你若不悔改.
我就臨到你那裡.
打你的燈台.
從原處羅去.
這個起初.
也都是說.
首先第一最重要的.
就好像英文那個字first.

$^{841}$可以有很多的解釋一樣.
所以這個神是強調.
愛是最重要的.
不可以離棄這個愛.
如果離棄了.
要快點去回想去悔改.
如果不是.
耶穌的警告是什麼呢.
會將你的燈台從原處羅去.
如果教會不再是燈台.
其實已經沒有意思了.
當一個黑暗的地方沒有了光.
那就繼續黑暗.
但這裡耶穌提醒原來.
起初的愛是很重要的.
而大家都很熟悉.
耶穌給我們的誡命.
一點都不複雜.
就是要盡心盡性盡意盡力.
愛主我們的神.
其次就是愛人如己.
其實基督徒.
我們在地上.
一生人要做什麼呢.
其實兩件事.
很簡單的.
兩件事.
跟你旁邊那個說.
很簡單的.
兩件事.
我一生人就是做這兩件事.
沒有了.
兩件事.
愛神.
和愛人如己.
愛神愛人.
是不是很簡單.
其實也不是很簡單.
愛神比較容易.
因為耶穌那麼可愛.

$^{881}$當然是愛他.
他那麼愛我.
愛神不是太難.
不過愛人呢.
你看看你旁邊那個.
你覺得愛他難不難.
有時也OK.
有時也不知道.
因為我們人的愛.
真的很有限.
我們必須承認.
人很容易愛那些可愛的.
對我好的.
疼我的.
不過有些本來疼我.
後來又好像不太疼我的.
也未必那麼容易愛.
是不是.
可能我們也很容易.
會經歷一些不同的情況.
以致我們都受了傷害.
在愛裡是有懼怕.
而不是沒有懼怕.
但是耶穌說.
祂賜給我們一個新的命令.
要我們彼此相愛.
而耶穌的標準是什麼.
祂說我怎樣愛你們.
你們就要怎樣相愛.
這個標準也很高.
耶穌怎樣愛我們.
愛到為我們死.
愛到我們一次又一次.
令祂失望.
祂也繼續愛我們.
大家有沒有試過.
回去祈禱.
主啊對不起.
我答應你們的事.
我又做不到.

$^{921}$求你繼續憐憫我.
幫助我.
然後祂又繼續幫助我們.
耶穌就是這樣.
我們無論什麼時候.
就算離開多遠也好.
只要願意悔改.
回到耶穌那裡.
祂也繼續愛我們.
但是耶穌要我們這樣去愛人.
不容易的.
我們很容易在愛的過程裡.
可能開始的時候.
起初的愛.
我想起我初初服事的時候.
也是很單純的.
我已經服事了主33年.
是不是看不出來.
33年了.
我一出生就全職了.
不是說笑.
33年.
開頭的時候.
我也記得.
當時服事很單純.
主說要成為中國人的祝福.
那就愛.
很快又嫁給中國人的老公.
然後又繼續去愛.
但是原來在過程裡不容易.
開頭我也記得跟主說.
主啊 你要我做什麼都可以.
為你死我都願意.
甚至我曾經做過這樣的祈禱.
現在想起來也很厲害.
我就說.
主啊 反正每個人都有一死.
我最想為你殉道而死.
這樣死法是最好的.
是不是很瘋狂.

$^{961}$現在回想起來.
可能我未必敢做這樣的祈禱.
但是當時就是這樣.
很單純.
去服事.
在中國.
在教會.
去服事.
但是過了一段時間之後.
原來服事是會受傷的.
原來你以為你很愛人.
原來別人未必去領情.
或者我愛的方法未必對.
而且在服事的過程裡.
我也有很多軟弱的.
可能我也會有比較.
我也會有懼怕.
別人說一句話.
我又很敏感.
諸如此類.
原來人真的很軟弱.
甚至有一段時間.
我是經歷過所謂的burn out.
就是去到一個地步.
就是覺得.
唉 我沒力再服事下去了.
我記得那段日子.
我也繼續在教會裡面.
繼續講道.
大效敬拜.
諸如此類.
但是我裡面開始跟人之間.
我發覺自己有一道很厚的牆.
因為經歷了一些事情之後.
我對人失去了信任.
我覺得原來人是很恐怖的.
原來人會在背後說你.
原來人會這樣說你閒話.
原來這樣這樣.
可能其實是幾個人傷害我.

$^{1001}$不過我覺得全世界都是這樣.
就是這樣.
然後就很多懼怕.
服事其實很明顯自己知道是沒力的.
很明顯知道是沒愛的.
講話講完都等於廢話.
自己的感覺就是這樣.
但是那時候我就禱告.
知道這樣下去不行.
有一次在禱告裡面.
神給我看到一個異象.
在異象裡面.
我看到自己站在一個很高的高台上.
下面全部都是羊咩咩.
咩咩咩咩.
很多羊咩咩.
然後我就看到.
自己在高台上.
把食物扔到下面給羊吃.
吃啦吃啦.
扔給羊吃.
然後有個聲音跟我說.
餵羊不是這樣餵的.
你要下去下面.
去看羊的需要.
因為有些是受傷的羊.
有些羊是很飢餓的.
你這樣餵牠是吃不到的.
你要去愛牠們.
你要去逐隻去顧念牠們的需要.
我馬上的反應是.
不行的羊很恐怖的.
會咬人的.
我下去一定會被牠們咬死的.
當時就是這個異象.
然後就完結了.
我很明顯知道.
是聖靈跟我說話.
我們聽聖靈的聲音.
其中一件事就是這件事.

$^{1041}$祂讓我看到.
反映著我內心的恐懼.
我覺得羊很恐怖.
然後我就知道.
神要我開始禱告.
是不是羊真的這麼恐怖呢.
沒錯羊也有牠的問題.
羊有時候也很蠢.
但是我也蠢.
我也是主的羊.
也有牠的問題.
但是神的眼光是怎麼看的呢.
神是愛到願意為每一隻羊寫命的.
甚至走迷了路的.
牠也會去找回牠的.
我就開始呼求主.
主求你幫助我.
能夠回復這個起初單純的愛.
拿走我內心的恐懼.
然後神開始感動我.
特別有一些我覺得是傷害我的人.
我來找回他們.
我發現原來有很多是誤會.
同時我覺得牠有問題.
但其實我也有傷害牠的地方.
然後神的愛開始教我.
如何修補這些關係.
也開始我的禱告就是.
主求你幫助我.
不要再用自己的眼光去看人.
因為我看人很多都會錯.
我覺得這個好那個不好.
諸如此類.
但是神才能看到人的內心.
求主讓我用你的眼光去看人.
這樣禱告之後.
我發現身邊的人真的可愛很多.
特別是我的老公.
馬上就可愛了很多.
原來神看牠是很可愛的.

$^{1081}$所以回復起初的愛.
保羅也說.
就算我能夠說萬人的謊言.
天使的話語.
但如果沒有愛.
也成了明的樂響的筆.
就算我有各樣的知識.
沒有愛也是沒有益處的.
就算我有信能夠而生.
也不算什麼.
我們知道.
可能我們也很想見到能力恩賜.
但最重要的是愛.
我們需要追求愛.
耶穌教導我們.
如何回復起初的愛.
有三個步驟.
祂說.
要我們去回想.
在哪裡墜落了.
接著要悔改.
接著要行起初所行的.
所以今天.
我相信聖靈也在提醒我們.
我們有沒有失去起初的愛呢?.
可能是對神和神的關係.
也可能是跟人的關係.
究竟哪裡墜落了呢?.
我就知道.
是因為有些傷害.
令我對人產生不信任.
產生懼怕.
不知道你們有沒有呢?.
悔改.
悔改的意思.
就是要去到神的面前.
求神用祂的愛.
除去攔阻我們去愛神.
和愛人的事情.
你呼求祂.

$^{1121}$祂真的會幫助你的.
神很好的.
接著就是行動.
就是花時間.
如果有關係要修補.
就求神使我們成為一個.
使人和好的人.
接著跟人的關係.
重新去恢復.
在《約翰一書》.
這也是約翰所寫的.
他提醒我們.
我們相愛.
不要只是在言語和舌頭上.
不要只有講字.
而是要在行為和誠實上.
這樣做的時候.
我們就是屬真理的.
我們的心就可以安穩了.
就可以在神的面前.
沒有懼怕了.
而且他特別提到.
我們一切所求的.
就會從祂得著.
大家想不想.
你的禱告目應運呢?.
原來很重要的.
就是我們要彼此相愛.
彼此相愛的時候.
我們的禱告.
神就會應運.
而且我們所做的事情.
就是祂所喜悅的.
最後.
在這段經文裡.
耶穌的應許.
給那些聽從聖靈聲音的人.
祂說會成為得勝者.
其實每一個書信.
最後祂都說.

$^{1161}$聖靈向宗教會所說.
凡有耳聽的.
就應當聽.
得勝的.
就會有一個應許.
而對於以弗所教會的應許.
是什麼呢?.
祂說得勝的.
就會賜神樂園中.
生命樹的果子給他吃.
對於以弗所教會來說.
得勝者是什麼呢?.
就是能夠勝過.
能夠回復起初愛的人.
就是得勝者.
本來失去了.
可能種種原因.
可能在服事裡有爭執.
可能大家都很努力.
不過你又不欣賞我.
我又不欣賞你.
或者有嫉妒.
有紛爭.
但是當他們.
靠著聖靈去得勝之後.
這裡說.
他們能夠吃生命樹的果子.
在創世記裡面.
有兩棵很出名的樹.
一棵叫生命樹.
一棵叫分辨是/或是樹.
我們都知道.
亞當夏娃.
因為吃了分辨是/或是樹的果子.
所以就帶來死亡.
聚集到世界.
但是今天我們作為.
亞當夏娃的子孫.
都求主幫助我們.
不要吃分辨是/或是樹的果子.

$^{1201}$分辨是/或是樹是什麼呢?.
就是我用我自己.
去分辨什麼是好.
什麼是不好.
我們人很容易的.
批評論斷.
這個好人.
這個壞人.
這個人有得救.
這個人愛我.
這個人不愛我.
我們很容易.
用自己的眼光去看.
但是我們要回到.
生命樹那裡.
生命樹是什麼呢?.
就是回到.
和神的關係.
我們知道.
生命是從神而來的.
愛也是從神而來的.
我們要更多靠著聖靈.
就是聖靈求你幫助我.
在這個世界裡面.
能夠去愛.
好像你愛我那樣.
我們更多去經歷.
神的愛.
神的醫治.
我們就能夠去愛.
我自己都很感恩.
其實我都不完美.
我都曾經.
經歷過抑鬱症.
我都曾經.
對人完全失去信心.
但是呢.
我今天能夠繼續去服事.
不是因為我厲害.
我只是有一樣東西厲害.

$^{1241}$就是我懂得呼求主.
就是我知道.
我需要主.
我沒有主就不行.
到今時今日.
我都繼續求主.
繼續給我愛的力量.
因為我都.
繼續可能在以後的日子.
都有機會再受傷害的.
但是都求主幫助我.
就算受傷害.
都要繼續去愛.
就好像耶穌愛我們一樣.
所以願意我們.
都一齊能夠去吃.
生命樹的果子.
因為二人所結的果子.
就是生命樹.
有智慧的.
必能得人.
我相信神要祝福.
同福堂的.
每一個弟兄姊妹.
都能夠成為.
得人如得魚的人.
藉著愛.
去彰顯.
神自己的榮耀.
\newpage



\section{希伯來書 2:5-18}
\label{sec:032461}
\textbf{希伯來書系列 (3) 救贖的元帥}~~~~張健庭牧師博士
\newline
\newline
連結: \href{http://www.church.com.hk/acms/content.asp?site=cdc\&op=show\&type=product\&code=032461\&layout=sermon}{\texttt{032461}} ~~~~ 語音日期: 2026-07-12
\newline
\newline
\hyperref[sec:032460]{\small{< < < PREV SERMON < < <}}
~
\hyperref[sec:index_chronic]{\small{[返順時目]}}
~
\hyperref[sec:032462]{\small{> > > NEXT SERMON > > >}}
\newline
\newline
希伯來書 2:5-18
\newline
\begin{longtable}{cl}
\hline
\hline
章節 & 經文 (和合本修訂版)\\
\hline
2:5 & \begin{tabularx}{0.7\textwidth}{X} 我們所說將來的世界，神沒有交給天使管轄。 \end{tabularx} \\ \\ \relax
2:6 & \begin{tabularx}{0.7\textwidth}{X} 但有人在某處證明說：「人算甚麼，你竟顧念他；世人算甚麼，你竟眷顧他。 \end{tabularx} \\ \\ \relax
2:7 & \begin{tabularx}{0.7\textwidth}{X} 你使他暫時比天使微小，賜他榮耀尊貴為冠冕，你派他管理你手所造的， \end{tabularx} \\ \\ \relax
2:8 & \begin{tabularx}{0.7\textwidth}{X} 使萬物都服在他的腳下。」既然使萬物都服他，就沒有剩下一樣不服他的了。只是如今我們還不見萬物都服他； \end{tabularx} \\ \\ \relax
2:9 & \begin{tabularx}{0.7\textwidth}{X} 惟獨見那成為暫時比天使微小的耶穌，因為受了死的痛苦，得了尊貴榮耀為冠冕，好使他因著神的恩，為人人經歷了死亡。 \end{tabularx} \\ \\ \relax
2:10 & \begin{tabularx}{0.7\textwidth}{X} 原來那為萬物所屬、為萬物所本的，為要領許多兒子進入榮耀，使救他們的元帥因受苦難而得以完全，本是合宜的。 \end{tabularx} \\ \\ \relax
2:11 & \begin{tabularx}{0.7\textwidth}{X} 因那使人成聖的，和那些得以成聖的，都是出於一。為這緣故，他稱他們為弟兄也不以為恥， \end{tabularx} \\ \\ \relax
2:12 & \begin{tabularx}{0.7\textwidth}{X} 說：「我要將你的名傳給我的弟兄，在會眾中我要頌揚你。」 \end{tabularx} \\ \\ \relax
2:13 & \begin{tabularx}{0.7\textwidth}{X} 他又說：「我要依賴他。」他又說：「看哪！我與神所給我的兒女都在這裡。」 \end{tabularx} \\ \\ \relax
2:14 & \begin{tabularx}{0.7\textwidth}{X} 既然兒女同有血肉之軀，他也照樣親自成了血肉之軀，為能藉著死敗壞那掌管死權的，就是魔鬼， \end{tabularx} \\ \\ \relax
2:15 & \begin{tabularx}{0.7\textwidth}{X} 並要釋放那些一生因怕死而作奴隸的人。 \end{tabularx} \\ \\ \relax
2:16 & \begin{tabularx}{0.7\textwidth}{X} 誠然，他並沒有幫助天使，而是幫助了亞伯拉罕的後裔。 \end{tabularx} \\ \\ \relax
2:17 & \begin{tabularx}{0.7\textwidth}{X} 所以，他凡事應當與他的弟兄相同，為要在神的事上成為慈悲忠信的大祭司，為百姓的罪獻上贖罪祭。 \end{tabularx} \\ \\ \relax
2:18 & \begin{tabularx}{0.7\textwidth}{X} 既然他自己被試探而受苦，他能幫助被試探的人。 \end{tabularx} \\ \\
[1ex]
\hline
\hline
\end{longtable}
$^{1}$主席.
今次我們繼續看希伯來書.
今天就看到第二章第五章第十八節.
上一次就說到希伯來書.
就比對了主耶穌和天使的比較.
天使在猶太人來說已經是一個非常超越的一個受造物.
因為天使本身有很大的能力.
都是在猶太人的傳統當中.
我們相信甚至乎在西乃山頒佈律法的時候.
都是經過天使的手來做中寶來賜予給人.
所以猶太人對於天使都是覺得非常值得尊崇.
但是上一次就說到.
作者就比較天使都是神計劃當中的獨極.
但是主耶穌就是神的兒子.
兩個是不能比較的.
而主耶穌一定是比天使更加超越的.
今天就繼續來到去講的時候.
似乎又將話題轉了一點.
一開始第五節.
我們所說將來的世界神如沒有交給天使管轄.
但有人在經上某處證明說.
人算甚麼你竟顧念他.
世人算甚麼你竟眷顧他.
你叫他比天使微小一點.
賜他榮耀尊貴為冠冕.
並將你手所做的派他管理.
叫萬物都服在他的腳下.
既叫萬物服他都服他.
就沒有剩下一樣不服他的.
只是如今我們還不見萬物都服他.
接著作者就將那個焦點現在就轉一轉.
他就從天使當中轉到人那裡.
在某個程度來說.
天使當然是比人的能力更加大.
你在舊約的故事當中.
你看到譬如兩個天使.
就可以將羅德救了出來.
在索多瑪俄摩拉的天房墳城的時候.
或者是神差一個天使.
就將在希西加的時代.

$^{41}$將亞述的軍隊18萬的整個軍隊打敗了.
所以天使是很厲害的.
但是作者就說到.
其實天使是要完成神的計劃.
但是在詩篇第八篇.
作者就引用詩篇第八篇.
就說到其實在神的計劃當中.
有另外一位.
或者有另外一個角色.
是比天使更加重要的.
你知不知道是誰?.
就是人.
在神的整個宇宙萬物的計劃當中.
人是扮演著一個非常重要的角色.
人和天使在能力上是完全無法較量的.
但是人在神的救贖計劃當中.
他是比天使重要很多.
因為人是神的計劃當中.
其中一個很重要的位置.
你叫他比天使微小一點.
賜他榮耀尊貴為冠冕.
並將你所做的都派他管理.
叫萬物都服在他的腳下.
作者引用詩篇第八篇就說出一點.
人之所以在神的救贖計劃當中這麼重要.
原來因為人是神的代表.
神創造了宇宙萬物的時候.
將人好像放在中間.
叫人代表神來管理宇宙萬物.
所以說你派他管理.
叫萬物都服在他的腳下.
人之所以為人有尊貴.
不是因為人的能力很大.
如果你說人的能力很大.
那麼天使就比人大.
人之所以為人的尊貴.
是因為神給了他一個位置.
是代表神來管理.
神所創造的宇宙萬物.
所以弟兄姊妹.

$^{81}$你看到這裡是一個非常特別的位置.
作者引用詩篇第八篇的時候.
他就想說出這一點.
原來在宇宙萬物當中.
人是很特別的.
人是代表神.
所以在創世記一開始的時候就已經說.
人是按著神的形象而被造的.
所有其他的動物.
各樣的東西.
甚至天使本身.
都沒有說是按著神的形象而被造的.
就是說是神的.
如果用圖像或者比喻來說.
我們是神的帽子吸出來的.
如果你做月餅的那個月餅帽子.
上面有什麼字.
中秋快樂.
一吸就吸了中秋.
做出來的月餅就寫著中秋快樂.
而人的形象.
當然就不是說神有眼耳口鼻.
有手有腳的那個形象.
是不是.
而是說人代表神.
人在整個宇宙萬物當中.
神的計劃當中.
有一個很獨特的位置.
就是我們是代表神.
所以宇宙萬物的東西.
神就派給人來管理.
以前你在中國的朝代.
那個人拿著聖旨.
就等於皇帝.
所以那個人拿著聖旨.
即使是太監.
他去到那裡.
拿聖旨出來.
就算是一品兩品的官員.
都要下跪接旨.

$^{121}$同樣的原來.
我們是有一個很尊貴的身份.
人之所以在宇宙萬物之中獨特.
不單是我們懂得思想.
我們懂得創造.
我們懂得自我超越.
而是人是有神的形象.
在宇宙萬物當中.
我們就是代表神.
但是這裡就講得很特別.
既叫萬物都服他就沒有剩下.
一樣不服他.
只是如今我們還不見萬物都服他.
這裡就將籠統的人.
是收窄了範圍.
是講為一個人.
這個他不是他們.
他說叫萬物都服他們.
不是所有人.
是他.
這裡究竟想講什麼呢.
原來這裡所講的不是普通的人.
當然原先所有的人都是很尊貴.
但是人因為犯罪墮落的時候.
已經失落了人之所以為人那種尊貴.
而人之所以為人那種尊貴.
最後如何得到恢復呢.
就只是藉著一個人.
一個很特別的人.
這個人就是耶穌基督自己.
當講到這裡的時候.
所以耶穌基督就成為了人類的代表.
所以他說唯獨見那成為比天使少一點的耶穌.
在這個角度來說.
耶穌降世為人的時候.
事實上是比天使更加微小.
上一堂都講過天使都挺有權威的.
司祭項目的爸爸薩加利亞 駁了一句嘴.
然後天使就啞到司祭項目出世才可以講話.
你看到天使很有權威.

$^{161}$但是耶穌基督的時候經常都被人駁嘴.
耶穌傳道或者耶穌教導門徒的時候.
經常都被人頂嘴.
好像也沒那麼厲害.
所以你看到耶穌作為人的代表.
後面會講到他完完全全成為人的一份子.
而作為一個不折不扣100\%完完全全的人.
耶穌站的位置所處的能力比天使還要少.
雖然他比天使還要少.
但是他為人恢復了人的尊嚴.
因為受死的苦就得了尊貴榮耀為冠冕.
叫他因著神的恩為人賞了死命.
這部份關於救贖等一下會再講.
但是我想停下來給大家思想一個問題.
你有沒有想過究竟信耶穌是什麼意思呢?.
為什麼你要信耶穌呢?.
或者再調一調問題再問一問.
就是如果信耶穌就是代表接受神的救恩.
那麼這個救恩是要達到什麼目的呢?.
有沒有想過呢?.
神的救恩要達到什麼目的呢?.
如果你要知道.
譬如你要去一個地方上車.
你要知道這個車的目的地是哪裡你才上的.
你不是亂上的.
同樣都是一樣.
你要做基督徒你要信耶穌.
你知不知道這個救恩會帶你去哪裡達到什麼目的呢?.
有沒有想過?.
我相信很多人第一個答案給出來的答案是什麼呢?.
就是我要上天堂不下地獄.
是不是?.
很多人最直接的答案就是這個.
是不是?.
那我就再問一個問題.
除了上天堂下地獄還有什麼呢?.
還有什麼呢?.
那似乎很多人又不懂得回答了.
那我告訴你.
如果你只是想上天堂下地獄.

$^{201}$那這個就是你將救恩的豐富.
是完全的無平了.
就等於去參觀海洋公園或者去參觀水族館.
問你去完水族館之後有什麼看呢?.
看了些什麼呢?.
那些小朋友可能只是喜歡八爪魚.
他就只是這樣.
有八爪魚看.
但水族館不只是有八爪魚.
有很多其他的東西.
他就只是看到八爪魚.
其實是錯失了很多東西.
很多時候.
救恩是什麼你明白嗎?.
救恩其中一方面就是.
人之所以為人.
是在神的救贖計劃裡面.
是很尊貴.
可以代表神.
可以跟神去相交.
這是神創造人原先的目的.
神創造.
在舊約創世記一開始的時候.
神創造亞當夏娃的時候.
不是創造他預備天堂地獄給他的.
根本風馬牛不相及.
創世記一創造亞當就預備他上天堂嗎?.
根本不是.
完全不是這樣的一回事.
神創造人的時候.
原先就是在神的心裡面.
一個好像掌上明珠.
將人放在心上.
一個掌上明珠的位置.
就好像將人可以成為他的代表.
好像分享著神自己的尊貴.
人之所以為人的尊嚴.
是因為人就是代表神.
另外人也可以在神的愛當中.
跟神相交,交流.

$^{241}$在伊甸園裡面.
你完全看到這些.
神在圓中行走.
跟亞當夏娃完完全全面對面交流.
一個成年人.
跟一個小朋友的交流是不對等的.
因為很多東西你是交流不到的.
因為他的智力還沒去到你的水平.
或者他的感受還沒去到你的水平.
你跟他說什麼失戀很痛苦.
是都不知道是什麼來的.
除了智力的問題.
他的感受,經歷那些不好.
現在再低一層.
如果他不是一個小朋友.
是一隻小狗.
又再低一點.
再低一點可能是一條金魚.
又低一點.
你跟金魚的交流.
每天放工跟他講講今天有什麼工作.
怎麼受氣那些浪費氣.
金魚當他聽不聽到你的一回事.
當他懂得聽.
他都不能夠明白你.
再低一點就更加不行.
但是.
就是因為交流本身是需要有某程度的容量.
有那個容量才可以交流.
大家沒有對等的容量是交流不到的.
而人的獨特的位置.
在創造宇宙萬物當中.
在創世記可以看到.
人是有一個容量.
可以跟神交流的.
正正就是因為人是按著神的形象被造.
所以人跟神是可以彼此相交.
有那個Fellowship.
這個也是伊甸園.
稱為樂園的重點所在.

$^{281}$所以當在救贖當中.
主耶穌要恢復的.
他代表著人的位置要恢復的.
就是恢復人的尊嚴.
以致人可以跟神再一次交流.
就好像去到啟示錄最後的時候.
就是神的帳幕在人間.
好像伊甸園再一次恢復.
而且是更加榮耀的恢復.
以前是一棵生命樹.
現在變成了十二棵生命樹.
就是神的榮耀計劃.
不單止是恢復.
而且是更加榮耀.
更加豐盛地去彰顯出來.
而這個就是救贖的第一部分.
原來救贖本身不是說.
我死了之後上天堂落地獄的決定.
當然那是有相關的.
但這裡不是說這件事的最低限度.
如果你說死了之後上天堂.
你說我還沒死.
所以現在我還是謀生.
謀生比其他的重要.
但是如果你說.
救贖本身根本就不是.
重點就是恢復你和神的關係.
而去恢復你作為人的尊嚴的時候.
就和你今時今刻已經非常有關係.
為什麼那麼多時候.
人好像活著找不到那種尊嚴.
好像那種尊嚴是需要世界給的.
有功有尊嚴.
有錢有尊嚴.
有榮譽有尊嚴.
或者諸如此類.
但是你那些東西又會失去.
所以你很痛苦.
但是原來不是.
人之所以最高的尊嚴原來.

$^{321}$就是在神的計劃當中.
可以和神分享神自己的榮耀.
而這個人犯罪和神脫離的關係就已經沒有了.
所以人用其他的東西替代就很空虛.
這個時候.
現在耶穌就要恢復了.
而耶穌的恢復的時候.
等下會再說為人人嘗死味.
祂成為人.
為人死在十字架上是有重點.
等下會再說.
但是在第十節那裡.
祂說出了一點.
原來那為萬物所屬.
為萬物所本.
要領許多兒子進榮耀裡去.
使救他們的元帥因受苦難.
得以完全本事合宜的.
這裡說出了一個很有意思的觀念.
這個觀念本身.
現代人是很少去這樣想的.
第一就是.
這裡所說的就是要進入榮耀裡面.
除非你是看那些.
以前暫時那些.
譬如羅馬時代那些.
國鬥士那些片.
很有榮耀.
打贏他們有榮耀.
或者你是中世紀那些武士.
要和人生死決鬥.
可能有人羞辱了你.
你要和人生死搏鬥.
拿回你的榮耀.
這些我們覺得很過時.
看電影覺得很過癮.
但是真實你不會覺得很重要.
但是這個觀念本身.
其實在聖經那裡是很重要的一部分.
當然聖經的榮耀.

$^{361}$和剛才所說的那些get data.
那些元卓武士那種榮耀是不同的.
但是聖經那裡是有一個進入榮耀的觀念.
是什麼意思呢.
就是說人是有那個屬神的.
因為剛才說人被神創造的時候.
是有了神的榮耀.
而這個榮耀是失去了.
當人犯罪的時候.
這個榮耀是失去了.
現在救贖是要恢復這個榮耀.
是要恢復這個榮耀.
你說這個好像很遙遠.
我個人還是比較低調一點.
榮榮耀都沒有什麼重要的.
最重要是生活安恕就行了.
或者沒事沒事 平安無事.
平平安安又到一天就最好了.
沒有什麼榮耀.
但是你不要這樣想.
因為這個榮耀本身.
和現代人可能沒有用榮耀這個字.
但是從另外一個角度來說.
就是有一種人之所以為人屬於你的東西.
你失去了.
如果我用現代人心理學的說話來說.
就是人失去了做人的目的和意義.
本來神造我們.
我們有這個榮耀的時候.
我們很清楚自己的目的和意義.
這個是我們去代表神的榮耀.
但是當你沒有神.
你活在世上的時候.
你也找不到終極的目的和意義的時候.
其實就是因為你失去了神造你的榮耀.
這個就是問題所在.
罪所帶來的.
用聖經的術語來說.
除了人和神脫離的關係之外.
人也失去了榮耀.

$^{401}$但是用今天的說話來說.
不用聖經的術語來說.
罪也令人失去了.
人自己的存活的目的和意義.
人看不到自己有終極.
可以不靠著這個世界所給你的東西.
而有任何的榮耀本身.
你明白我的意思嗎?.
本來神造你.
我舉例來說.
你應該是自己懂得發光的.
因為神造你.
那個榮耀神的形象在你裡面.
不是在外面來的.
那你就自己懂得發光.
但是你黑暗了.
沒有了光.
那你就要在外面的東西.
要靠外面的電筒照下來.
那你才發光.
就好像太陽.
本來自己懂得發光的.
現在太陽沒了光.
就變成像月亮一樣.
要靠太陽的光照.
要靠外面的東西照.
那外面的東西是什麼呢?.
剛才你可以想像到.
代入到你在世間上.
你所追求到的東西.
令你覺得自己有一種榮耀感.
有些人追求金錢.
有些人追求名譽.
有些人追求.
很被人欣賞.
有自我認定的價值觀.
諸如此類的各樣東西.
就好像要靠外面的電筒照過來.
那你就發光.
但是那些電筒有時候是沒電的.

$^{441}$人家照你的時候.
有時候是不照你的.
所以有時候你就會很失落.
因為你本身.
你已經失去了進入榮耀的路.
你自己失去了.
你也不知道怎樣走回到神的榮耀裡面.
而這裡就說到主耶穌就是那個.
要領許多兒子進榮耀裡去.
就說主耶穌.
他怎樣領許多兒子進榮耀裡去呢?.
他說:他就成了使救他們的元帥.
這裡就說到.
主耶穌就是成為了那個元帥.
帶許多人進榮耀裡面去.
元帥這個字是什麼意思呢?.
有些英文的翻譯就是Captain.
就是那個.
Captain就是什麼呢?.
如果是船就船長.
就是帶頭走在前面的.
中文元帥就是可能打仗的時候.
那個長官就是那個元帥.
總之就是帶領在前面衝鋒的那個.
現在聖經就用了一個這樣的觀念.
人心失去了人之所以被造的榮耀.
剛才所引用的《詩篇第八章》所說的.
人之所以為人的榮耀尊貴.
人失去了.
現在有一個特別的人.
就是耶穌基督自己.
做這個元帥.
做這個我們的Captain.
再一次帶領我們進入神的榮耀裡面.
這個是救贖.
而他為什麼可以有資格做救贖的元帥.
帶領我們進去呢?.
他說那為萬物所屬.
為萬物所本的.
他這樣去形容.

$^{481}$就是說這個.
本來宇宙萬物.
你讀希伯來書前面已經說過.
第一章第一至第三節已經說過.
宇宙萬物都是靠耶穌基督而做的.
宇宙萬物都是靠耶穌基督而維持的.
但是現在這個創造宇宙萬物的.
創造主維持者.
現在進入人間.
成為我們的元帥.
成為人的一份子.
成為我們的元帥.
再一次好像衝鋒陷陣的走在前面.
現在要帶領那些屬他的人.
再一次進入神的榮耀.
這裡就是在說這個事.
救贖.
救恩.
不只是死了之後上天堂.
是現在神要將屬他的人.
在此時此刻.
讓他成為你的元首.
你跟從他.
他已經衝在前面.
你一直跟從他帶領的時候.
再一次進入神原先創造人.
應該有的那種榮耀.
當然了.
這個救贖計劃.
主耶穌承受了極大的苦難.
使救他們的元帥因受苦難.
得以完全本事合二的.
他自己成為這個元帥本身.
他不是用神的身份.
用神的能力.
去做這個元帥.
他是自己成為一個人.
去做這個元帥.
他成為一個人去做這個元帥的時候.
他就經歷了人間的苦難.

$^{521}$就是這樣的時候.
所以他說.
主耶穌都是因為他要做人的元帥.
在苦難當中可以這麼說.
使他成為完美的元帥.
這就是救贖的故事.
在這個救贖故事的背後.
他說.
因使人誠誠的和那.
捨得以誠誠的都市出於一.
所以他稱他們為弟兄.
不以為恥.
說我要將你的名傳遠我的弟兄.
在會中我要重揚你.
又說我要依賴他.
又說看啊.
我與神所給我的兒女.
這裡就說出了.
主耶穌為什麼要道成肉身.
之前我也說過.
救贖的故事根本.
與那些神仙下凡.
希臘神話那些.
中國神話的神仙下凡.
希臘神話那些.
那些什麼什麼神.
變了做人.
然後在人間出現.
做一些很厲害的事.
完全不同.
因為主耶穌是要進入去人類當中.
成為一個不折不扣的人.
所以他經歷過人的每一個階段.
由嬰兒出世.
其實一開始由十月懷胎.
到嬰兒出世.
到成長.
聖經有說主耶穌做小孩子的時候.
十二歲的時候就上聖殿.
然後又經過成年期.

$^{561}$他的爸爸.
若失應該是.
如果研究聖經的學者.
推斷.
應該是一早已經死了.
耶穌出來傳道的時候已經沒有了若失.
他經歷過提早喪父.
也都經歷過他要做一家之主.
要去照顧母親.
最後當然也都經歷過人的排擠.
門徒的出賣.
人類的誤解.
甚至乎將他釘在十字架.
然後受苦.
然後死亡.
然後才復活.
所以他絕對有資格成為我們救贖的元帥.
因為他經歷過人的每一個階段.
經歷過人的每一個苦難.
所以他說使人成聖的.
他本來是神.
使人成聖.
但是他說和那些得以成聖的.
就是我們.
可以因為依靠他們而得以成聖.
到時出於一.
到時出於一的意思就是說.
都是同一個肉體.
同樣的生命.
活在世上.
而他也都不羞恥地稱我們為弟兄姊妹.
這個就是說道成肉身.
神之間.
神與人之間的認同.
如果墮落是人因為背叛了神.
失去了人之所以為人的尊貴榮耀的時候.
救贖就是神自己成為人.
去將人的尊貴和榮耀再次恢復.
以致人與神再次有關係.
而要這樣做.

$^{601}$不是創造天地萬物的主.
在外面踏踏隻手指.
或者.
你想想.
神創造宇宙萬物是怎樣創造的呢?.
怎樣創造?.
是不是.
踏踏隻手指這樣做?.
變魔術這樣.
一扭開帽子.
就飛了隻白鴿出來.
飛了個地球出來.
當然不是.
根本都不用這樣.
這樣更容易.
神說有就有.
根本手指都不用踏.
但是對於救贖來說.
就完全另一個檔次的難度.
你可以這樣說.
你想想.
神創造宇宙萬物只是說.
但是神要救贖的時候.
祂就親自成為人.
去進入人.
來我們中間.
這樣去完成救贖.
所以.
是一個非常奇妙.
也是一個非常偉大的救贖.
如果你.
不體會到神自己所付的代價的時候.
你對於神的救贖.
可能你都不懂得.
怎樣去欣賞和珍惜.
如果你也不體會到.
原來救贖要帶你去哪裡的時候.
你對於恢不恢復這種關係.
也不覺得有什麼需要這麼重視.
是不是.

$^{641}$而這個正正就是.
神自己的作為.
神自己救贖的作為.
這個.
也讓我想到一點.
其實.
我之前都沒有這樣想過.
就是.
神.
最配得過.
稱為我們的神.
你知不知道.
知不知道是什麼.
有什麼元素神最配得過稱為.
是我們的神.
是不是神的能力.
是不是神的全知.
是不是神的智慧.
這些全部都很重要.
但是其實.
最主要.
是因為神.
能夠體會我們和我們的認同.
為什麼你這樣說.
因為最近.
說一點.
又是題外話又不是題外話.
因為這是一個挺特別的例子.
最近其實有一本科幻小說.
也不是最近的.
也出了好幾年.
我拿了很多獎的一本科幻小說.
科幻小說.
其實當中.
當然有一些小說的情節.
就是說將來人的AI很厲害.
人創造了AI出來.
AI很厲害.
AI基本上.
差不多.

$^{681}$如果用人的角度來說.
差不多好像有神那麼大的能力.
又無所不知.
又無所不能.
當然這些是小說.
是一個比喻.
又無所不知又無所不能.
要做什麼都做得到.
他要控制.
人全部都控制得到.
很厲害.
差不多升級到神的檔次.
但是.
人呢.
當不當AI是神呢.
不能當AI是神.
為什麼呢.
因為.
AI不明白人.
所以有時候.
AI的決定.
在人來說很殘忍.
很邪惡也有.
因為.
他要懲罰人.
AI定了一些很公義的規矩.
但是其實.
對於那些人來說.
那些人懲罰得很殘忍.
諸如此類.
那個故事.
那個科幻小說的故事.
就說到.
最後.
其實他問一個哲學問題.
究竟神怎樣才稱為神呢.
是不是有能力.
才能稱為神呢.
是不是有無所不知.
才能稱為神呢.

$^{721}$是不是很高的智慧才能稱為神呢.
那些AI完全都符合了.
但是都不能稱為神.
因為甚至在人的眼中.
這是很邪惡的.
是一個很恐懼怕的.
當然.
只是一個小說.
說到回到時光機.
可以回到.
回到古代的時候.
回到沒有AI之前的時候.
他告訴他重新做一次.
人要做一個甚麼AI出來.
是可以.
有這麼有能力.
差不多去到神的當次的時候.
人都才可以.
不懼怕.
可以稱為神的呢.
那你知不知道.
那個小說就說.
他回到以前.
時光機回到以前的時候.
當他做那個AI的時候.
是要有那個.
empathy.
是要有那個跟人的認同.
到那個AI.
慢慢慢慢越來越進步的時候.
又有能力又有智慧.
又有各方面的.
其他那些.
很厲害的.
發展出來的時候.
他都不會去傷害人.
因為他會認同.
人的苦難.
認同人的感受.
其實那個小說背後.

$^{761}$是講一個哲學的問題.
甚麼才稱為神.
而主耶穌.
唯一.
我們基督教的信仰.
就是這樣.
我們的神.
是100\%完全跟我們認同.
這個世界.
有很多其他的神很厲害的.
很有智慧的.
當然那些真是假.
我就不在這裡講.
即使當是有的.
但即使是有這樣的神.
都未必一定是神.
是不是.
魔鬼都很厲害的.
魔鬼都很聰明的.
魔鬼都是殺人命如草芥.
但是可能都是.
可能要達到這個.
世界人口太多.
就殺緊一點.
主耶穌.
唯有我們的神.
他100\%與我們認同.
他才可以做我們.
救贖的元帥.
而.
更加.
聖經說.
更加講到一件事.
去到.
人最終極的敵人.
就是面對死亡.
而.
我們.
好像被死亡的權勢.
操控著.

$^{801}$沒辦法有人可以擺脫.
但主耶穌.
他親自的.
成為血肉之體.
以致他可以經歷死亡.
而在死亡中復活.
去勝過死亡.
以致可以解放.
我們這一生.
我們未必真的.
好像說每天都很害怕.
害怕受怕.
每天都害怕死的那種.
但是.
無可否認我們.
死亡是將人控制著.
而這種控制.
好像將人成為了.
五指山走不出.
好像成為了.
死亡的奴役的時候.
主耶穌的復活.
幫我們去釋放.
而.
去到這裡的時候.
作者就.
再將這個論點.
再推前一步.
在神的事上.
成為慈悲忠信的大祭司.
為百姓的罪.
獻上作挽回祭.
作者就再在前面.
再推進一步.
就是說當主耶穌成為人.
成為救贖的元帥的時候.
其實他也都成為.
可以成為了我們.
一個慈悲忠信的.
大祭司.

$^{841}$大祭司是.
當時的人他們很明白.
有兩個要求.
第一個要求.
是大祭司要.
體恤到那些.
去求.
求問的人的一種.
連述.
就是.
如果你去.
求神的憐憫的時候.
你不知道神會不會應運.
大祭司就好像代表神.
來宣告.
那個神的憐憫.
所以大祭司要體恤人的軟弱.
是必要的條件.
第二個條件.
就是忠信.
大祭司是忠於.
他所託的託付.
而完成這個職任.
而主耶穌.
他成為人的時候.
他的救贖.
因為他也都能夠完全成為.
大祭司的身份.
而當然.
他成為大祭司有所獻的時候.
他獻的就是他自己.
這一部份.
在希伯來書後面.
會講到大祭司.
獻祭那些.
第七第八.
第九章那些會再講.
現在這裡.
只是帶出一個前提.
整體來說你看到.

$^{881}$今天我們就明白到.
原來.
作者要我們明白救恩本身.
是沒有第二個人.
可以為我們完成的.
是只有耶穌基督.
自己.
創造宇宙萬物的主.
成為人.
那為萬物所屬,為萬物所本.
來到我們中間.
成為那個.
救贖的Captain.
元帥.
帶領我們進入.
神的榮耀.
恢復人之所以為人.
的榮耀.
求神幫助我們.
當我們明白這些的時候.
我們就更加明白救恩是多麼的偉大.
我們多麼的蒙恩.
可以領受救恩.
\newpage



\section{箴言 12:27}
\label{sec:032462}
\textbf{懶到出汁}~~~~曹子光牧師
\newline
\newline
連結: \href{http://www.church.com.hk/acms/content.asp?site=cdc\&op=show\&type=product\&code=032462\&layout=sermon}{\texttt{032462}} ~~~~ 語音日期: 2026-07-12
\newline
\newline
\hyperref[sec:032461]{\small{< < < PREV SERMON < < <}}
~
\hyperref[sec:index_chronic]{\small{[返順時目]}}
~
\hyperref[sec:032463]{\small{> > > NEXT SERMON > > >}}
\newline
\newline
箴言 12:27
\newline
\begin{longtable}{cl}
\hline
\hline
章節 & 經文 (和合本修訂版)\\
\hline
12:27 & \begin{tabularx}{0.7\textwidth}{X} 懶惰的人不烤獵物；殷勤的人卻得寶貴的財物。 \end{tabularx} \\ \\
[1ex]
\hline
\hline
\end{longtable}
$^{1}$好 今日想同大家讀一節經文.
這節經文是來自《尊嚴書》第十二章 第二十七節.
《懶惰出執》.
我們一齊讀一讀這節經文 請讀.
懶惰的人不孝打獵所得的.
殷勤的人卻得寶貴的財物.
我們讀多一次.
懶惰的人不孝打獵所得的.
殷勤的人卻得寶貴的財物.
大家都知道.
蜜蜂基本上是象徵著勤力的代表.
大家知不知道要製造一公斤的蜂蜜.
整個蜂巢的蜂蜂總共要飛多遠呢.
總共要飛十四萬公里.
概念就是足足圍繞著地球三圈半.
他們要探訪至少四百萬朵花.
才可以做到這樣的效果.
而一隻夏季的蜂蜂.
其實牠的生命力只有六個星期.
嘩 原來非常短.
夏季的蜂蜂只有六個星期的壽命.
而蜂蜂雖然是男生.
但是牠們全部都是化育不健全.
所以牠們不能夠生小孩.
所以牠們的老大賀爾蒙.
退化到只剩下一樣的衝動.
就是瘋狂地做.
不斷地做 所謂的蜂蜂.
而牠們會隨著年齡來升職.
例如牠會做清潔.
由清潔會做保姆 照顧.
然後就起蜂巢.
最後就會去到最危險的前線.
就是外勤的財物.
蜂蜂的老大基本上是沒有放假.
也沒有退休這個概念.
所以牠們的人生.
不是 蜂蜂的生命裡面只有兩個字.
就是儲糧.
儲糧 然後都是儲糧.

$^{41}$當蜂蜂飛到翅膀斷的時候.
再飛不起的時候.
牠們就會選擇寧願死在蜂巢外面.
慢慢爬出去.
希望不要白佔了蜂巢的空間.
所以牠們的勤奮一生.
都是為別人 不是為自己.
都是用牠的生命來換整個族裔的賢族.
這一種奉獻.
勤奮 殷勤的相對就是懶惰.
不知道大家覺得自己是殷勤還是懶惰呢.
你的小朋友是殷勤還是懶惰呢.
心裡有個數 相信大家.
大家知道 前幾天 也不是前幾天.
今天也是熱到出汁.
大家看到你潑到這樣就知道熱到出汁.
你就算什麼都不用做.
只是坐著 你也會開始慢慢流汗.
真的熱到出汁.
那麼懶惰到出汁又怎麼樣呢.
懶惰到出汁就明明是懶惰到什麼都不做.
但是都居然出汁.
懶惰到出汁 是不是.
其實這個意思就是說.
人懶惰的那個點.
是連汁都逼出來.
帶出一種人去到盡.
這種懶惰 懶惰到盡.
懶惰到極致.
那種逼真感.
大家知不知道這個世界衛生組織裡面有一個調查.
原來每年有200萬人死於懶惰出汁病.
其中有六成的人有心臟病 中風 糖尿病 腎病 癌症等等的疾病.
專家更指出 懶的病因是什麼呢.
懶去廁所.
有沒有試過懶到不去廁所.
忍 頂住.
試過有多少小時沒有去過廁所.
大家想想那個數字.
有一次我女兒和我搭飛機.

$^{81}$她13個小時沒有去過廁所.
真的嚇到我.
懶去廁所.
懶得喝水.
懶得煮飯.
還有懶做運動.
這些都是致命的點.
但是今天我們要講的是懶到出汁.
其實這個是一個很悲的畫面.
為什麼呢.
因為根據經文告訴我們.
一隻本來新鮮獲得的獵物.
因為主人懶得去處理它.
最後變成腐爛 發臭 出汁.
徹底地報銷.
今天就和大家看一節的真面.
十二章的第27節.
看看怎麼樣我們能夠拆解.
拒絕成為一個懶到出汁的人.
第一件事我們要別再爛尾收場.
我們先看看上半句.
懶惰人就不孝打獵所得.
這句說話如果放在古代的世界裡.
其實是一個很諷刺的語境.
孝 不是英文.
是火字邊的孝.
孝這個字在整本舊約聖經裡是很罕見的.
是指透過火的溫度來燒烤.
而所謂孝是什麼呢.
就是這個短句的最後一個字.
就是獵物.
到底是誰孝呢.
是懶惰的人.
其實懶惰給人的感覺就是.
什麼都不動.
像不像你家的貓呢.
整天都不動.
躺在那裡 躲起來.
整天躺平.
但其實他不是一個天天躺平的人.

$^{121}$他不是一個連手指都不動的人.
這幅圖畫裡面的主角.
相反根據這個經文告訴我們.
我們有多肯定.
他居然是願意拿起打獵的器具.
弓箭 矛 槍 不知道.
總之他花很大的體力時間.
走出去打獵.
在古代世界裡面.
你知不知道打獵是一個很高的運動量.
而且很危險.
你知不知道.
你要追蹤那個獵物.
還要埋伏等.
還要小心樹林的樹枝刮傷你.
還有樹林裡面有很多猛獸.
還有很多獵戶.
他們設下的陷阱.
那些捕獸器.
所以你知道這個所謂的懶惰人.
其實他竟然可以做到.
他甚至可以成功捕捉到那隻獵物.
他不是一個我們所想像的懶惰出汁的人.
但最荒謬的事件就是.
當他把獵物拿回家之後.
他竟然懶得去洗.
懶得去拔毛.
懶得去煮.
當時你們知道沒有冰箱.
你試一下.
你買一排肉回來.
你就放在外面.
放在外面.
我想不用一天.
他馬上會變.
當時天氣多麼炎熱.
是不是新鮮的肉放在屋裡面.
幾個小時.
真的當時都可能分分鐘會發臭.
會腐爛.

$^{161}$變得就沒有用.
你有沒有試過在某些超市買了一些肉回來.
然後沒有處理就會發霉.
或者有些臭味在裡面.
我真的試過一個C字頭.
我不說名字了.
大家經常去的.
在C字頭這間很大間的所謂的超市.
我就買了牛尾回來.
說起牛尾你都知道哪一間.
通常都是那間大店才可以買.
就很大一排.
我看著它的日期.
那個日期寫明是那天.
你知道我的人比較懶.
我的人通常到期前一天.
我才去處理問題.
結果我就在到期那天.
嘗試去打開一剪.
簡直臭味紛圈.
當你如果不理那件事的時候.
假設到了期的話.
分分鐘有這樣的後果.
而在中世紀的猶太人裡面.
有一句很古老的.
他們認為這是一個古老的諺語.
其實他覺得這一句的真言.
其實是在說他們有能力去獲取資源.
但是因為缺乏了最後的一步堅持.
最終都沒有辦法享受自己所得來的效果.
這個正正就是中國人所說的半途而廢.
我們真的要學習.
不要浪費你已經擁有的.
更加不要只是做那個令你興奮的前半段.
而放棄了可能你認為很沉悶.
很關鍵的後半段.
不過細心想一想.
這裡其實還有一個更值得我們去深思的一個疑問.
一個這麼願意頂著烈日當空.
冒著風險.

$^{201}$還有相當的技能.
才能夠捕獲獵物的人.
他怎麼會捕獲獵物之後.
懶得去孝贖他呢.
如果他連孝贖他這個相對較容易的步驟都去懶理的話.
其實他怎麼捱得過打獵那個步驟呢.
其實大概.
我相信答案不在於他是懶.
而是在於有另外兩種的.
跟我們生活相關的東西.
第一種就是他分心.
曾經讀過一篇的研究.
不知道大家認不認同.
這篇研究很有可能是香港的.
我懷疑他是香港的.
他說原來辦公室.
每十一分鐘就會有一個指令.
有一樣東西令到你停下來.
你在做的事情.
是會令到你干擾了你的.
例如老闆跟你說.
可不可以幫幫忙做一樣東西.
又或者你的同事跟你說.
昨天的蛋糕很好吃.
又或者接著打電話跟你說.
你有沒有這個產品.
總而言之.
每十一分鐘.
你就會被干擾一次.
另外.
有時除了是這些外在的干擾之外.
也有另一個調查顯示.
原來人在每十四秒.
是十四秒.
他的想法就會變一次.
不知道大家同不同意.
十四秒.
你的想法就會變一次.
而另外有些人.
聲稱自己可能多功能運作的人.

$^{241}$其實這些多功能的人.
他們以為自己的能力.
可以做很多不同的事情.
在同一個時段.
其實專家說.
這個根本是騙局.
這個只是反而令到我們.
每天在做的事情.
會損失28\%的時間.
情況就好像你的手機.
你打開你的手提電腦.
同一時間.
你開著越來越多的工作的話.
例如APP.
例如網頁.
例如背景的程式等等.
你的電力.
是會耗損得更多.
回到這一節的真言.
大概這位仁兄.
他打獵回來之後.
很有可能.
他有機會被其他事情吸引了他.
搶走他的注意力.
例如.
坐下吃一件蛋糕.
或者他先聽個音樂.
或者他休息幾分鐘.
休息十分鐘.
又或者洗個澡.
小休一會兒再回來的時候.
原來已經忘記了.
這一件事都未定.
你知道天氣一熱的時候.
那些食物就會變壞.
又或者朋友突然間.
call他出來玩一下.
又或者突然之間有個娛樂.
令到他分了心.
又或者有另一件更重要的事.

$^{281}$他覺得先處理好一點.
有一些東西吸引了他.
令到他分了心.
做了那一樣東西.
以致到忘記了那塊肉本身.
這個正正就是這個世代.
可能我們面對的寫照.
我們不是沒有能力.
去完成那件事.
而是我們的專注力.
太容易被下一個通知.
下一個訊息.
彈出來的那個video.
拿了我們的.
那個注意力.
你有沒有試過讀書讀到一半的時候.
拿起手機.
想著告訴自己.
我看五分鐘.
然後其實是.
你放下手機的時候.
可能已經隔了半個鐘.
或者分分鐘更長的時間.
而且.
更加特別的地方就是.
你要知道有時候.
你買一些食材正如你剛才所說.
你放在雪櫃的時候.
然後你打算煮了.
怎知隔了幾天過期了.
又浪費了那個東西.
因為你在這幾天裡面.
你分心做其他東西.
以致到那些獵物腐爛了.
大學生可能.
有這個感覺也不一定.
你有沒有發現.
廚房的升盤.
是特別多碗碟.
我去過很多.

$^{321}$跟大學生聊天吃飯.
我常常都見到廚房裡面.
碗是多到不得了.
然後我就問他為什麼你們不用碗.
不是我的那些碗.
那裡起碼.
我相信不可能是每個學生.
那些碗來的.
一定是有其他人再加碗上去.
不止一個單位.
很多人就懶得去做那樣東西.
很有可能他們想著吃完飯.
洗個澡再弄.
弄著弄著別人已經開始吃完.
第二些人已經吃完了.
碗又放在裡面.
如此類推.
升盤裡面永遠都是碗碟.
最後大家忍.
誰忍不了就誰洗.
第二個可能性除了分心之外.
結果.
對於打獵的人來說.
可能他只是享受過程.
例如那種追蹤.
那種刺激.
就好像我們玩爬山.
行山 健身.
玩樂器 釣魚.
你很享受過程.
至於捕獲之後的獵物.
你根本都不在乎.
你不在乎.
如果是這樣的話.
這不是我們所想的難.
而是一種價值觀的錯置的問題.
他很享受過程.
但他傾看那個結果.
這種的解讀.
就剛剛和這一節的經文下半句.

$^{361}$就可以說是一個非常好的伏筆.
下半句說什麼.
因勤的人,卻得寶貴的財物.
這裡所說的寶貴.
不單止是說那件事本身.
那個物件本身值錢.
而是說.
更加是一個重要的態度.
因勤的人.
他珍惜自己努力所有的.
他期待和看重那個成果.
所以他才很用心去保存它.
去善用它.
原來決定一件事有沒有價值.
不是在於你給了多少.
而是你有多在乎的結果.
無論是分心也好.
還是不在乎的結果也好.
這兩種可能性.
都比純粹是懶惰.
來得更貼我們.
每一個人的真實處境.
因為大多數人.
都不是完全不動的人.
我看大家的樣子都不是完全躺平那類人.
你們都是肯動的人.
我們的問題往往都是.
心思散漫.
或者根本沒有去完成.
或者珍惜那些東西.
所以我們都是開局的大師.
但卻是爛尾的.
重症患者.
現在很流行一種所謂的.
假裝很努力的速食文化.
就是熱情來得快.
去得更快.
這個就是.
這位行動很積極的獵人.
他很輕易.

$^{401}$或者他很努力之下.
打倒獵物.
但從來對那個結果都不覺得有什麼責任.
或者要負責那個結果.
我突然之間有沒有這個共鳴呢.
例如你在社交平台.
買了一本書.
你熱血沸騰到一個地步.
馬上下單.
把整套電子書買回來.
這些叫做打獵所得.
但買了之後.
你連點開第一頁.
你的耐性都沒有.
這種就叫做不孝的獵物.
又例如在新的學期裡.
你下載了一個最新的.
時間app.
做完一份很溫馨.
很完美的計劃進度表.
沒錯.
你做到一個捕捉獵物的意思.
在裡面.
但幾天之後.
你開始覺得這些內容很枯燥.
也跟進漏了一天.
你就開始熱情大減.
然後你就報廢了這個時間表.
這些叫做不孝的獵物.
又或者你很辛苦.
買了一件樂器.
練幾下學幾堂.
玩一段時間.
就再沒有玩過了.
很多.
但完成的是很少.
我對音樂.
也有丁點的興趣.
沒有天分.
興趣.

$^{441}$有一次我第二個兒子.
拉Violin.
有空沒事做.
就拿起Violin出來拉兩下.
原來真的好像宰雞一樣.
真的有這樣的感覺在裡面.
但我告訴自己.
這個世界沒有事做不到.
我結果.
嘗試拉奇異音調.
讓我拉到幾個音.
真的拉到.
我只用了十分鐘.
就拉到第一行.
我心裡想.
以我這樣的速度.
一日之內我一定可以拉到這首歌.
然後我就放下了這個Violin.
想著喝杯飲料聊聊天.
然後我再沒有下氣.
所以到今時今日.
我只會拉到頭一行.
再不可以再拉下去.
是開始到你覺得很興奮.
但你沒有這樣的.
完成下去的心態.
有沒有一件事.
兄姐妹朋友.
是你已經打獵所得.
但你一直丟在一個角落裡.
沒有去敲的東西.
可能是一本書.
可能是你報了名一直缺席的興趣班.
我年輕的時候.
有讀過德文.
有讀過法文.
但每次都完成不到.
那個課程.
完全沒有一次.
完成得到.

$^{481}$又或者你很想學的東西.
永遠只能夠學一半.
沒有學下去.
甚至信仰上.
我們都是這樣.
很堅持想著每日去祈禱.
每日去讀聖經.
每次去小組.
每個禮拜去教會.
但我們到最後.
都沒有學到.
不過感恩.
我們在這個年度.
第一次成人主義學.
我們只有四堂.
我親眼目測到.
至少有八成或以上的人.
完成了四堂.
以前你知不知道.
去教會教成人主義學.
通常都是一季.
有十三次.
通常第一次是最多人去.
最後一堂.
永遠都是最少人.
因為我們永遠都是這種心態.
開始.
但完成不到.
可能我們要先問自己一個問題.
令我們卻步的是分心.
還是不在乎.
如果是分心.
解決方法很簡單.
將你的手機放在一邊.
或者放在抽屜裡.
花一個特別的時間.
做完那件事才拿出來.
如果是不在乎.
問題就大一些.
你需要重複問自己一個問題.

$^{521}$這件事對我來說.
值不值得去珍惜.
值不值得去重視.
例如.
我去學.
單車游泳.
我以前為何要學會這兩樣東西.
我心裡有一天想.
如果有一天我去到別處.
去牧會也好,宣家也好.
這兩樣都是求生的技能.
你似乎要把山橫.
就學會了.
如果你沒有動力.
你不知道重要性.
你是很難有動力去完成那件事.
學一件事.
到底你要想想.
你珍不珍惜那件事.
你重不重視那件事.
值不值得.
這正正是荷里活巨星史太龍給我們一個功課.
他是荷里活一個很出名的巨星.
但當年他很年輕的時候.
他也是一個很窮的人.
他帶著自己寫的一個劇本.
逐間荷里活的製片商.
電影公司.
去敲門.
全部都被拒絕.
甚至有些人一見到這個劇本.
就立刻可以扔垃圾桶.
他試了一輪.
兩輪,三輪.
都沒有人理他.
直至到第三輪之後.
去到第350間公司的時候.
第三輪的第350間公司的時候.
有人說.
不如你可以留下你的劇本.

$^{561}$這個就是Walkie.
很出名的一個劇本.
這個就是一個烈火裡面.
拿去敲.
直到它去發光發熱.
要成功甚至要做到.
從來都不是一次到位.
而是屢戰屢敗.
直至那件事情做完.
這個對於我們每一個現在摸索自己方向的人來說.
是一個很大的釋放.
你不需要第一次.
就去做到完美.
而是你需要繼續走下去.
做那件事.
就一定能夠完成.
還有第二件事值得我們去學習的就是.
不要靠捷徑和運氣.
我們很容易忽略了一個很重要的事.
就是原來懶惰的人.
這個希伯來文.
其實在舊約裡面.
同時有個意思叫取巧.
或者是奸詐.
或者叫鬼詐.
換句話說.
人的問題可能不純粹是.
一個懈怠或者是懶惰.
更可能是想用最少的力氣.
去換取更大的成果.
這裡其實是.
藏著兩個可能性.
如果這個字的意思是指.
懈怠或者是懶惰的意思的話.
這個故事的邏輯就是.
這個人的能力完成.
但是他因為分心.
或者是不在乎.
最後都沒有去敲到那些獵物.
但如果這個字是指.

$^{601}$鬼詐的意思的話.
故事就會變成這樣.
這個人用了一些不義的.
走捷徑的方法.
不合法的渠道或者是詭計.
的手段去拿到那些東西.
而不去敲.
不是他的原意.
他都不是想的.
但是最終他都沒有機會享用.
因為不義的東西會流走得很快.
最後都保不到.
無論是哪個解釋都好.
結論都是指向一個方向.
靠懶惰浪費的.
固然是白費.
靠取巧而得到的.
更加是留不到.
真正能夠持守到底的.
真正能夠享用成果的.
就只有腳踏實地.
誠實的路徑.
留意執言的下半句.
他說恩勤的人.
卻得寶貴的財物.
按照漁民的次序.
最後那個字.
可以理解為恩勤.
但同樣可以理解為正金.
如果是正金的話.
這句說話就會變成這樣.
高尚的人的財富.
乃是正金.
正金.
是所有金石裡面.
最需要經過火煉.
亦都是經不起偷工減料的一種金屬.
你沒有辦法用捷徑.
去偽造正金.
亦都沒有辦法靠運氣.

$^{641}$去將它來到去挖到現成的正金回來.
你只能夠透過反覆的高溫.
不斷來到去提煉.
將這些雜質一層一層地消走.
才能夠拿到九九九純度的真金出來.
這個和捷徑.
和運氣.
可以說是完全相反.
賭博贏回來的錢.
留不到很久.
很快你就會輸回給它.
抄功課換來的後果.
沒錯你分數可能很高.
但是一考試.
你就會落淚.
走後門的機會.
沒錯你可能拿到這個機會.
但是你沒有真材實料的支撐.
你遲早都是寸步.
這個政策回應了.
我們經常聽到的一種叫懶人哲學.
大家知不知道Microsoft的老闆比爾蓋茨.
他曾經說過.
我會讓一個懶惰的人.
去處理困難的工作.
因為懶人總能夠找最簡單的方法.
完成任務.
這種說法似乎.
因為懶出門.
所以就弄了這些社交平台出來.
不然怎會有Zoom.
怎會有Google meet這些東西.
因為懶不用出門.
就做到這些東西.
因為那些人就會發明這些東西.
而懶自己寫報告.
懶得去查字典.
懶得找人來查grammar.
所以我們很歡迎AI.
我們非常歡迎.

$^{681}$因為一秒就可以幫你搞定所有的東西.
很多人.
因此就會覺得.
走捷徑才是王道.
但千萬不要誤會.
這些是懶人的哲學.
那些所謂的懶人哲學.
發明AI的天才.
表面上在幫人建立了省力.
但他們其實背後的思維.
是建立了一個很驚人的高效的殷勤.
絕對不是聖經所說的懶惰.
更加不是靠捷徑.
運氣等等的東西.
甚至他們背後其實是一個非常強大的經濟體系.
例如那些科技巨頭.
他們是一些賺大錢的商道.
要做這些東西.
所有AI的模型.
例如check GPT.
靠emphobic.
他們都需要極高階的圖形處理器.
即是所謂的GPU.
都需要這些東西來去運算.
目前這個板塊裡面幾乎全部都被NVIDIA壟斷了.
NVIDIA單季的營業額.
已經高達幾百億的美元.
市值更加是超過五兆美元的大關.
甚至一早已經超越了蘋果和Google.
而那些微軟,Google,Amazon.
這些科技巨頭.
他們就建立了龐大的雲端基礎的設施.
任何想用AI的公司.
他們必須向他們租用伺服器.
而這些大型的科技公司巨頭.
為了維持AI的運作.
他們都必須用電力,核能,綠色能源來配合.
所以這些大行.
他們就不斷發財.
不斷來建立他們的生意.

$^{721}$這些才是他們真正的商道的世界.
沒錯.
這才是AI背後的真實.
最殘酷.
但也是最龐大的交易真相.
世界上沒有憑空出來的捷徑.
我們眼前的便利.
全部都是因為這些科技巨頭.
在背後默默耕耘.
用極致的努力和資本.
去鋪設這樣的道路出來.
AI的大爆發.
名義上是幫人省了力.
實際上是講這些憑硬實力.
精準去收割的商機.
你可以用智慧去尋找高效的方法.
但千萬不要誤會這些科技的便利.
將它當成不需要努力的運氣.
你可能有智慧去找最省力的方法.
但你不能夠缺乏面對責任.
腳踏實地的承擔力.
所以你就不斷去複製複製複製.
走捷徑.
甚至用鬼債的方式.
最後你都沒有辦法轉變成真正的養分料.
的確這個世界最大的誘惑.
不是懶得去做.
而是你懶得去想.
當你利用AI的時候.
利用科技為自己省下來的時間.
突然之間我想問.
那你省下來的時間你拿去做什麼.
是讓你看多幾套片.
還是玩遊戲機.
還是睡覺.
舉一個例子.
可能你要做一份project.
大家都用AI生成一份的內容.
很快已經做成了.
如果你直接交了這件事.

$^{761}$沒錯AI可以幫你做到這件事.
但是如果你願意在AI幫你做多這些事的時候.
你多出來的時間.
再check下它.
再做一些功夫下去的時候.
你就會發現.
可能你在裡面學到更多東西.
能夠做得更好的東西.
你要修正一些你前所未有可能修正的東西.
你交出來的東西.
就不是一份打獵所得的獵物.
而是一份真正敲熟了的成果.
我可以告訴大家.
我講道都用AI.
我都是用大家AI那些東西.
來check下那些東西.
很快去問一問那些字眼為什麼.
但是我從來都不會完全照抄.
我絕對不會完全照抄.
甚至我做完一輪之後.
AI都說.
你都想不到你用這些東西.
它都會跟你對話.
它都會講一輪這樣的東西出來.
所以你要見到.
你在小事上面如何對待.
AI幫你節省的那些時間.
將來就會決定到.
你將來如何面對更大機會的時間.
這個是一個精心的一個學習.
而不是只是複製貼上.
靠僥倖.
靠這些捷徑.
最後第三樣東西.
就是不要錯過那個火熱的心.
不知道你有沒有留意.
回到這個真言本身.
我們經常都忽略了一個細節位.
就是火.
火當然不能玩.

$^{801}$我以為你現在熱到好像火滾一樣.
但是留意這隻經文的動作.
敲這個字.
敲是要用火的.
如果下半句那句說話.
是有精金的含義的話.
就是說精金要煉出來.
都是需要提煉的火.
將這些雜質燒走.
才能夠令到它有黃金的效應出來.
所以這節經文.
如果你將這個放為精金的意思的話.
那解釋上面就是徹頭徹尾圍繞著火這個意象.
懶惰的人不肯生火去處理他的獵物.
而殷勤的人他所得的.
他是經過火煉.
經得起考驗的精金.
獵物需要用火.
煉精金都需要用火.
第一種火是烤肉的火.
我們看見經文表面上是說打獵和處理獵物.
實際上其實是說.
你願不願意讓火進入你手上的東西.
懶惰人避開了第一種火.
他不願意生火去處理獵物.
結果用很辛苦拿回來的成果.
白白浪費了腐壞了.
他不是沒有獵物.
而是基於不同的原因.
將獵物可以變成好的食物的這個機會浪費了.
但另一種人呢.
叫做所謂殷勤更有質素的人.
或者叫高尚的人.
他願意讓第二種火進入他人生的生命裡面.
就是提煉的火.
你要知道火在這裡不是說懲罰.
而是說分辨.
火不是會創造精金原本沒有的東西.
而是將原本藏在裡面的東西.
將它的雜質燒走.

$^{841}$讓真正的價值浮出來.
這個正正是煉金的邏輯.
他不是要加添甚麼.
而是除去甚麼.
這種火比第一種火更慢.
更磨人.
因為他燒的不是外表的東西.
而是裡面的東西.
裡面的雜質.
他要處理的不是有沒有.
而是做得夠不夠徹底.
夠不夠誠懇.
夠不夠誠實.
精金之所以是貴重.
正正是因為.
正正是因為他經過了這種火.
而沒有去逃避.
雜質一層一層.
結果就燒走了.
反覆地燒的時候.
就提煉出九九九金.
兩種的火其實對應著兩種人.
不同人生的態度.
第一種火就是考驗我們願不願意動手.
第二種火就是考驗我們願不願意去被提煉.
針面似乎在提醒我們.
懶惰的所失去的.
從來都不是只是一餐的獵物.
他失去的.
是一次讓我們被煉成的機會.
而殷勤所得到的.
從來都不是眼前的精金.
他得到的是一個願意被火試煉.
而變得更真實.
更有份量的生命.
如果我們放大來看針面第十二章的整體框架的話.
你就會發現智慧人做了大量鮮明的對比.
善人和詭詐的人對比.
義人和惡人對比.
智慧人和愚蠢人對比.

$^{881}$在這個分類裡.
懶惰的人竟然被編排在惡人的裡面.
我們可能會覺得.
懶惰最多是性格的一個缺點.
沒有壞心眼.
只是一個懶惰.
但是在聖經裡.
就打破了這個興旱的觀念.
在上帝的眼中.
不負責任的懶惰就是惡.
因為懶惰本質.
在本質上是對上帝所賜的生命.
時間.
恩賜的一種傲慢和浪費.
當然如果懶.
是因為行詭詐就更加顯出.
惡的本質出來.
弟兄姊妹.
我們每一分擁有的時間.
每一個我們遇上的機會.
每一個在我們身上的恩賜.
都是上帝所賜的.
當我們懶惰去浪費這些東西的時候.
我們其實就是浪費神親手.
交給我們的禮物.
這個不是一個效率的問題.
而是一個對神態度的問題.
所以讓我們老實問自己一個問題.
我們可以很勤奮.
但是我們到底為哪一個來去殷勤.
為哪一個來去勤奮呢.
我們可以通宵打網絡遊戲機.
通宵和隊友去想一些攻略的東西.
我們不達到目標就不罷休.
我們也見過一些人為了MIRROR的演唱會.
EASON的演唱會.
軒仔的演唱會.
日夜去撲飛.
這些叫做極致的投入.
同意的.

$^{921}$但當我們翻開聖經讀不夠五分鐘.
我們就覺得沉悶枯燥.
我們可以為了一場考試.
為了社交媒體上面經營的自己一些東西.
連續幾天很細緻地去做這些東西.
經常留意有沒有人看你,like你.
有沒有人看了你的view數.
甚至焦慮到失眠.
我們這些叫做高尚的勤奮奮鬥.
但一講到教會的聚會.
一講到身邊有需要的朋友.
一講到每天下來讀經.
或者怎樣建立自己的靈命.
我們會馬上就會幻想.
會變成那個不孝的獵物.
這樣一個擺爛的mode出來.
這個是不是我們真實的寫照呢.
這個世界的事情上面.
我們是一個殷勤的策劃高手.
那在信仰上面呢.
我們連火都不願意.
上帝有很多機會擺在我們面前.
教會的事工,身邊認識的人.
很多人可能他剛剛來到這裡.
可能他覺得很孤獨,很迷茫.
在社交的媒體上面.
來可能甚至會求救.
甚至是光道部的移民.
可能有不同的意外發生.
又或者有個病,或者心靈軟弱.
社區你接觸了很多人.
其實這些都是上帝親自.
為我們打獵所得的獵物.
擺在我們面前的都是契機.
但是我們面對這些機會的.
曬度又是怎樣呢.
例如我們放長假期的時候.
我們將整副的心思就來拍卡.
拍照,旅行.
但當別人邀請我們去參與一個事工.

$^{961}$短暫的,我們會馬上搖頭.
太累了,沒有時間了.
還要給錢的,還要放長假.
不好了吧.
上帝將機會給了我們.
給了我們很多獵物.
但是我們懶,還是不懶呢.
可不可以花時間去敲熟它呢.
請我們記得,第一.
不可以再做爛尾收場了.
第二,不要靠運氣,不要靠捷徑.
第三,不要錯過那個火熱的心.
我們一起祈禱吧.
來讀這一節這麼寶貴的真言.
我們求你將這句真言放在我們的心裡.
讓我們學會不要再懶惰.
因為上帝在我們人生,在我們教會裡.
其實你給了我們很多獵物.
問題是我們有沒有去打.
有沒有去敲熟它.
當我們能夠敲的時候.
當我們不分心.
當我們在乎那個結果的時候.
我們就能夠享受那個成果.
求你閱立我們的禱告.
奉耶穌基督恩主的命求,Amen.
\newpage



\section{利未記 23:39-43}
\label{sec:032463}
\textbf{住棚}~~~~陳崇基牧師
\newline
\newline
連結: \href{http://www.church.com.hk/acms/content.asp?site=cdc\&op=show\&type=product\&code=032463\&layout=sermon}{\texttt{032463}} ~~~~ 語音日期: 2026-06-28
\newline
\newline
\hyperref[sec:032462]{\small{< < < PREV SERMON < < <}}
~
\hyperref[sec:index_chronic]{\small{[返順時目]}}
~
\hyperref[sec:032464]{\small{> > > NEXT SERMON > > >}}
\newline
\newline
利未記 23:39-43
\newline
\begin{longtable}{cl}
\hline
\hline
章節 & 經文 (和合本修訂版)\\
\hline
23:39 & \begin{tabularx}{0.7\textwidth}{X} 「但是，從七月十五日起，你們收藏了地的出產之後，要守耶和華的節期七日。第一日為要完全安息，第八日也要完全安息。 \end{tabularx} \\ \\ \relax
23:40 & \begin{tabularx}{0.7\textwidth}{X} 第一日，你們要拿美好樹上的果子、棕樹枝、樹葉茂密的枝條和河邊的柳枝，在耶和華－你們的神面前歡樂七日。 \end{tabularx} \\ \\ \relax
23:41 & \begin{tabularx}{0.7\textwidth}{X} 每年你們要向耶和華守這節七日。你們在七月裡所守的節，要成為世世代代永遠的定例。 \end{tabularx} \\ \\ \relax
23:42 & \begin{tabularx}{0.7\textwidth}{X} 你們要住在棚裡七日；凡以色列家出生的人都要住在棚裡， \end{tabularx} \\ \\ \relax
23:43 & \begin{tabularx}{0.7\textwidth}{X} 好叫你們世世代代知道，我領以色列人出埃及地的時候，曾使他們住在棚裡。我是耶和華－你們的神。」 \end{tabularx} \\ \\
[1ex]
\hline
\hline
\end{longtable}
$^{1}$我們一直都在討論這個系列.
這次我們應該是第五講.
我們來討論鑄棚這件事.
鑄棚是什麼一回事呢?.
我們先看聖經.
然後大家可以思想一下這段聖經對我們今天有什麼意思.
還有講到最後的時候.
我會告訴大家有些講堂上的更新.
利美的二十三章這裡說什麼呢?.
他說你們對以色列人說.
神的子民.
他說你收藏了地的出產.
從七月十五日起.
要守耶和華的節有七天.
第一天你要休息.
第八天也要休息.
中間做什麼呢?.
第一天.
首先你要拿美好樹上的枝子,果子.
棕樹枝.
密茂樹的枝子.
還有河邊的楊柳等等.
就是拿一些草,樹枝.
在耶和華的神面前.
來去歡喜七天這麼多.
你要守節是七天的.
作為耶和華的節.
這個就是你們世世代代永遠的定例.
要在這七月.
守這個節日.
你們要住在棚裡面七天.
以色列人中凡本地的人都要住在棚內.
好讓你們世世代代知道.
我領以色列人出埃及的時候.
曾使他們住在棚內.
我是耶和華.
你的神.
我們說的這個大系列.
就是從創世記開始直到聖經最後.
我們說聖經好像一個長篇的劇集一樣.

$^{41}$有很多集的.
你當它每一卷書是一集.
就有66集.
我們就說整個的聖經.
就是告訴我們.
這個劇集裡面.
最主要的主角是誰呢?.
是上帝自己.
上帝在人間,在人類,在自己的子民裡面做了很多事情.
然後祂就藉著這些故事,這些事蹟.
來告訴我們啟示.
介紹這位上帝是一個怎樣的上帝.
縱使這個世界是祂所創造的.
是祂安排得井井有條的.
但是因為人.
我們叛逆.
以致搞亂了這個世界.
但是上帝從來沒有放棄過這個世界.
一直到推進祂自己救贖的工作.
我們知道在創世記裡面.
祂就揀選了阿巴拉罕等等的族長.
由他們生出一個大國.
這個大國叫做舊約的以色列.
這個大國就領受了上帝的律法.
出了埃及.
然後成為一個新的子民.
然後他們能夠成為上帝的工具.
以致上帝能夠去祝福,傳遞所有裡面的人.
今天這段經文就是摩西五經的第三本書.
叫做什麼?.
利未記.
利未記那裡講了一些我們經常覺得是一種的規條.
就是跟我們現代人沒什麼相關的.
因為裡面講到會幕裡面辦理的祭祀.
或者那些儀式,規矩等等.
祭物要怎麼處理等等.
然後又講了一大堆的規條.
是講到時間怎麼使用的.
譬如節日.
什麼時候會有不同的節日等等.

$^{81}$這些都是利未記我們的感覺.
但是利未記其實是寫給每一個神的子民的.
因為利未記第一章第一節講到說.
上帝要曉如告訴摩西.
你要對以色列人說.
不只是以色列國裡面.
以色列人民裡面的祭司.
那些所謂神職人員.
去講這些所謂祭祀裡面的手冊或者規條等等.
原來每一個以色列人都需要明白這些東西.
就是說如果你拿著利未記看這一集的時候.
這個劇集.
或者這個書卷的時候.
你就會覺得我都好像大祭司一樣.
我要怎麼處理我自己.
怎麼去確保我自己能夠去到上帝的面前.
能夠讓上帝住在我們中間呢.
所以在這裡講了一些規條.
其中今天的經文就講到一些首節的規條.
當然有大大小小的節日.
如果你認識一些猶太人的朋友.
或者同事的話.
他們經常放假的.
久不久就放假.
他們每個月都有假放的.
就不像我們美國人.
我們美國人是最少假放的.
他們經常都放假的.
每個月都有假放的.
這裡其中講到一個月.
就是他們的農曆.
7月左右的時候.
7月15號.
就是今天我們的西曆大概9月尾.
10月的時候.
就講到一些節期的.
在摩西五經裡面.
告訴以色列人聽神的指紋.
他們時間是由上帝的原則來定的.
不是說你喜歡怎麼樣就怎麼樣的.

$^{121}$時間是有一種規律的.
而且那種規律是一種屬靈的.
spiritual(靈性)的規律.
就是上帝告訴大家.
你每逢七天裡面有一天.
是需要安息去敬拜我的.
你不能夠七天都做事的.
六天你要努力去做事.
但是你第七天的時候就是安息日了.
就是所謂今天新約裡面.
基督已經承傳了所有的東西.
所以很多時候是第七天的第一天.
我們去敬拜.
所以我們信了主的時候.
我們的時間就會改變的.
是不是?.
我們時間的習慣.
我們時間的運作.
我們都會有所改變的.
你有沒有想過我們的信仰.
跟我們的schedule是有關的.
跟我們的calendar是有關的.
yes 是有關的.
在舊約裡面.
神的子民領受這些律法的時候.
就告訴他們.
大大小小的節日裡面有三個最大的.
第一個就是.
預節passover.
就是西曆大概四月初我們復節的時候.
第二個就是七七四十九日之後.
第五十日.
五旬節Pentecost.
就他們收了他們的大麥小麥.
有收割節的時候.
那個大概是西曆的五月.
然後駐彭節就是今天我們所講的那段經文.
就是到他們.
大概今時今日的西曆九月尾.
十月的時候.

$^{161}$他們就需要去做祝鵬.
來到所謂的慶節,首節.
講到這些的節日.
我們每一個人.
都會想起的.
每逢我們有節日的時候.
我們就會想起那些典故.
或者那些由來.
最近是什麼節日?.
最近是端午是不是?.
是不是端午?.
是端午我沒有教錯.
因為我收了很多突然之間給我的愛心的種.
解決了很多頓飯.
我們會想起那些源頭.
為什麼有這樣的故事等等.
我們基督教最重要的節日是哪幾個?.
復活節和聖誕節.
可以這麼說.
我們比較西方的西化的基督教.
就是有這樣的重要節日.
在舊約裡面就是這三個最重要的節日.
今天這段經文是講到.
駐彭節的吩咐.
有幾件事值得留意的.
第一件事.
就是這些的駐彭節.
或者這些節日.
都是叫我們停下來.
來紀念一下.
想清楚上帝在我們心靈裡面有的恩典.
這裡很搞笑.
他說7月15日就開始守節.
什麼時候呢?.
首七天.
他說第一天你要怎麼樣.
然後第八天你要怎麼樣.
但是有趣的地方就是.
他說你要在出產之後.
你收藏了出產.

$^{201}$到7月15日起.
就要守節了.
什麼意思呢?.
就是你出糧之後.
出了糧.
你收藏了你的土產.
這是秋季的時候.
你有秋收.
然後你就要停下來.
去守節.
這個不是很有趣.
因為農作物入倉.
你收到這些大麥,小麥.
或者這些農作物的時候.
最安逸的時候就到了.
是不是?.
我問你你出糧那天.
你去哪裡呢?.
出糧那天你去哪裡呢?.
我還記得那時候我在香港.
就是中學畢業.
然後我找了一份工作.
那份工作當時是OK的.
那個工資.
就是你肯做就行了.
在銀行做的.
我還記得我第一張的支票.
出糧的支票.
出了的時候.
我去了哪裡呢?.
很有趣.
那時候銀行是出現現金的.
那時候我是拿一倍,拿一袋.
出現現金.
整疊現金在手上.
我就去了哪裡呢?.
我就去了酒吧.
我就和朋友一起去酒吧.
我請大家吃一餐.
就是你出了糧.

$^{241}$當然要去開心.
去安逸的時候.
最開心的時候.
去慶典的時候.
吃大餐的時候.
應該是這樣.
但是這裡很奇怪.
他說你豐收之後.
你就要做什麼呢?.
要搭棚.
然後在那裡住七天.
這個不只是以色列人的問題.
其實是說我們人性的問題.
我們最好景的時候.
我們最容易忘記上帝.
我們的心很容易傾向於物質的東西.
朱鵬哲偏偏安排在豐收之後.
神好像告訴我們.
你最危險的時候是什麼時候呢?.
你和上帝之間的關係.
如果你說你相信耶穌的話.
你是一個基督徒的生命的時候.
你最危險的那一刻不是你缺的時候.
是你有的時候.
是你豐收的時候.
就是當你拿到上帝的祝福.
手裡面有金有銀的時候.
那時候才是最大的危機.
《申命記》第八章.
提到以色列人進入迦南的時候.
進入豐收之地.
進入安逸之地的時候.
他們有一個警告.
這個警告是什麼呢?.
以色列人走過曠野四十年.
《申命記》五經第五本.
就告訴他們.
你進入迦南地有豐收有安逸的時候.
要記得什麼呢?.
你要謹守耶和華的律法.

$^{281}$以至你不會忘記耶和華.
因為什麼呢?.
因為在說人性的一件事.
就是你建造美好的房屋的時候.
有牛有羊的時候.
金銀增添的時候.
你心裡面就會驕傲.
驕傲不一定是.
被人家.
黑黑霸霸或者很威風的態度.
驕傲可以藏在心裡面.
驕傲就是我們有一種自足.
當我們有自足有自滿的時候.
我們就容易忘記耶和華.
忘記我們的神.
所以有人說得很對.
他說錢可以跟我們說很多東西.
並且每句說話幾乎是上帝可以跟我們說的.
因為如果有錢的話.
我們很富足的時候.
錢就好像上帝一樣可以跟我們說.
你很安全.
是吧?.
當你查一下你的銀行帳戶.
或者你的 portfolio的時候.
錢好像在跟你說話.
不用怕.
你有保障.
你是有和平的.
你有安全.
你是有平安的.
等等等等.
這些不一定是很錯的事.
錢可以跟我們說很多東西.
是上帝可以跟我們說的.
是吧?.
是取代上帝所說的話.
你未來一定是無憂的等等.
但是聖經裡面反反覆覆的告訴我們.
只有上帝才會賜給我們真正的平安.

$^{321}$所以在律法書裡面.
是有一個智慧的.
他將這個鑄盤子嵌入在一個時間.
就是你最豐收的時候.
拿完所有東西的時候.
人最富足的時候.
你就需要停下來.
你要想一下.
上帝怎樣恩待過你.
如果你不記得的話.
在忙碌裡面.
或者你在善忘的當中.
你很快上帝就會淡忘了.
我們教會去到第十年.
每逢教會去到唐慶.
或者新的一年的時候.
或者有新的發展的時候.
很多時候我們就會很自然的.
我們會數一數我們的成績等等.
我們做過多少事.
當然是有需要的.
有需要.
我最近在反反覆覆的在想.
我們教會從開始的時候.
我們是從零開始的.
從零開始.
我們一開始是沒有糧出的.
很緊的.
如果不是上帝的恩典.
真的不會成為一個教會的.
如果不是因為大英姐妹.
如果不是因為大家在禱告.
在信心裡面擺上.
是根本不會走到今天的.
對不對.
因為大部分.
我講北美.
大部分新開的教會.
是不過三年的.
三年就死的.

$^{361}$要不就是同工.
內部的領袖.
他們有瓦解.
或者談不攏.
或者有不同的外在.
或者內在的因素等等.
開教會的時候的同工.
大部分都還在.
大部分還在.
普中心留在這裡.
那些是中流砥柱.
但我想講一件事就是.
我們在堂面的時候.
我們可以數一下業績.
可以數一下成功的地方.
但是一件事.
千萬不要忘記.
想起零開始的時候.
原來是上帝的恩典.
這是朱鵬的智慧.
豐收之後.
你就需要停下來.
因為我們要學習一下.
去紀念一下上帝的恩典.
上帝在生命裡面.
有沒有祝福.
有沒有恩典.
你當然會說有.
但重點就是.
你有沒有停下來去數一數.
停下來紀念一下.
寫下來也好.
自己去到銘記也好.
以致你跟上帝的關係.
不會因為忙碌.
豐收.
而忘記了.
第二方面.
朱鵬哲的智慧是什麼.
我要面對人生的.

$^{401}$軟弱和脆弱.
這裡說朱鵬.
第四十節.
他說第一天.
你要拿樹上的果子.
樹枝等等.
在耶和華你神的面前.
你要歡喜快樂七天.
你要住在棚裡面四十二節.
幾天.
七天.
有一個星期.
有一個棚子.
希伯來文就是Sukkah.
或者縱數是Sukkot.
如果你認識猶太人的朋友.
他們首節的時候.
或者看一下那些的閱歷.
猶太人的閱歷.
你要知道有Sukkot這件事.
Sukkot就是.
朱鵬哲的時候.
就是今天西曆大概九月尾的時候.
棚子算什麼.
這裡很有趣.
後期的拉比.
去到解釋這些的.
律法的時候.
他們認為整個Sukkah.
是應該怎麼整的.
拉比文學裡面有一堆東西.
叫做Halaqah.
Halaqah.
Halaqah是什麼意思.
Halaqah就是Halaq這個字.
Halaq就是走出來的意思.
就是實踐出來的意思.
所以有一大批的.
律法的解釋條文.
是歸納了成為所謂Halaq.

$^{441}$Halaqah.
裡面解釋一個棚子.
應該是怎麼整的.
聽著有幾樣東西值得留意.
原來每逢朱鵬哲來到的時候.
你就要預備材料.
然後你就在後yard.
或者你有的地方.
你的所謂Patio.
你就會整一個棚子.
這個棚子的高度.
是不可以高過20爪的.
就是不可以高過.
大概今天我們算是33尺.
因為拉比認為.
如果你整得太高的話.
太過宏偉的話.
就會怎樣.
就會變成永久.
你就不記得了.
你就住在那裡.
你就覺得.
我永久就住在這個地方.
太舒服的時候.
你就忘記了.
你所住的.
只不過是一個很脆弱的棚子.
所以一定不可以太穩陣的.
第二樣東西.
就是要用天然的材料.
樹枝等等那些.
枯枝.
會損壞的.
會枯萎的.
還有是不太防水的.
不要用一些很好的材料.
就是你不要Home Depot.
你找外面隨便收拾一些東西好了.
是不是特別是打完風我們大拍外面.
斬一點下來.

$^{481}$然後你就搭一個棚子.
並且一個棚子是需要通天的.
棚的帳幕的周圍.
就是要簡陋的.
因為你住七天而已.
你住七天.
你就不會覺得.
真的是住一個家.
你不會真的搬去那個地方.
只不過讓你.
去到明白一件事.
就是其實.
一個屬靈上面的一個反思.
你住在那個棚子七天.
你就可以明白到.
人生真的很脆弱.
就像棚子一樣.
是不是.
我們現代人生活.
我們花很多錢.
我們去買保障或者.
買一些安全感回來.
買一些穩定回來.
住棚.
就是跟這個的.
概念是完全相反的.
你住七天一個會塌的棚子.
會塌的棚子.
並且規定你.
要弄一個會塌的棚子.
不可以弄得太實際.
因為人生就是脆弱的.
人生是脆弱的.
上個星期.
上兩個星期.
我們.
那個大風.
真的很可怕.
那你們有沒有停電.
有誰有停電.

$^{521}$我停了二十多電.
通常天氣就是這樣.
你最需要冷氣的時候.
它就停電.
你最需要那些東西的時候.
它就剛好會壞.
接下來這幾天又會有熱浪.
我告訴你.
你有冷氣就會壞.
很有趣.
停了一天幾電.
有美國以外的一些朋友.
跟我有些通訊的.
他說不是嗎.
你美國大國居然停電.
美國大國就會停電.
因為我們給很多錢給軍.
打仗.
軍備我們很厲害.
但是電力我們就不夠.
甚至NBC Chicago有一個報道.
問一個問題.
他說龍捲風的通道.
和龍捲風的大道.
因為有些地方特別.
譬如Kansas.
下面那些州.
就經常會有龍捲風.
這個所謂的Alley.
Tornado的Alley.
是不是會Shift.
因為氣候變化的緣故.
想都沒有想過.
我們都有份.
不要說.
別人會打龍捲風.
我們不會的.
可能我們都會.
鑄鵬的鵬子是什麼.
鑄鵬Tabernacle.

$^{561}$或者Squat.
就是說我們是脆弱的.
我們是頂不住的.
天氣氣候所變化的.
你怎麼可能有控制.
人生就是這樣.
人生就是有脆弱.
我們需要去承認的.
真正保護你的.
真的不是鵬子.
真正保護你的.
是在曠野裡面跟你同行的上帝.
祂的回眸住在你們中間.
祂去給你最終的保障.
祂給你最終的Security.
是住在你們中間的上帝.
不是你所踏的棚.
生命裡面我們可以踏多少棚都好.
都是會脆弱.
大風.
雨水臨到的時候.
我們就會倒下.
最重要的是什麼.
我們不是說不要見那棚.
我們要做好的管家.
但我們要明白我們本質是脆弱的.
教會也是一樣.
教會也可以是脆弱的.
因為我們都是不堪一擊.
我們都是有缺陷.
不是大英姐妹我們夾手夾腳的話.
我們早已經玩完的.
對嗎.
所以在申命記.
那裡說什麼呢.
他說.
申命記29章這裡沒有打出來的.
他說回顧40年裡面.
神怎樣帶領以色列人的時候.
他說什麼呢.

$^{601}$他說我領你們走過曠野40年.
你們身上的衣衫沒有穿破.
你們的鞋也沒有穿壞.
鞋沒有爛到衣服沒有穿破.
當然他不是在說神蹟.
他在說你所基本要有的東西.
所有必須的東西.
這40年裡面我都供應了給你們.
在40年的曠野生命裡面.
不是因為你碰到.
其實你很脆弱.
但是上帝一直去保守著.
去帶領著你.
除了去到面對我們人生的脆弱之外.
第三件事就是去到學會.
怎樣去承傳上帝的使命.
怎樣去傳承它.
這段鑄盤節的吩咐裡面.
其中一段最重要的說話.
就是好讓你們世世代代都知道.
我曾經領以色列人出埃及的時候.
曾經使他們在曠野裡面.
住在棚裡面.
我是耶和華.
你不是的神.
這可以算是最重要的一句宣告.
是一個神學的宣告.
因為在說一個鑄盤的日子.
這個之日.
並非是一個文化.
純粹文化的之日.
是一個活生生的.
一個重演的.
一個神學的宣告.
去告訴人聽屬靈上面.
哪一個是我們的上帝.
你知道猶太人的家庭.
他們到這個鑄盤節的時候.
他們會出去一個很脆弱的棚.
他們會住七天.

$^{641}$或者在那裡吃七晚餐.
他們有一些猶太人的.
叫做律法的移民.
那些ritual.
即是說父母會和小孩下一代.
來到去有問有答的.
好像那些問答的冊子那樣.
你打開那些冊子.
你就會問.
小孩會問一個問題.
然後父母就會回答.
其中一些問題是問什麼呢.
小孩會問.
為什麼我們今天會搬出來.
為什麼我們這七天要住在這個地方.
為什麼我們放棄.
我們有磚有瓦的.
很舒服的屋.
然後我們搬出來這裡呢.
在這裡紮營呢.
父母的回答就是什麼呢.
因為我們曾經有家.
曾經.
sorry.
因為我們曾經是沒有家.
但是上帝成為我們的家.
We did not have a home.
But God is our home.
為什麼我們會出來這裡住.
因為God is our home.
第二個問題就是.
為什麼屋頂會有洞呢.
那些sugar.
和那些棚紙.
為什麼會是通洞會漏水呢.
父母就回答.
不是因為是壞的.
是故意的.
是故意弄到這樣是通風的.
因為真正遮蓋我們的不是屋頂.

$^{681}$是上帝.
屋頂不是我們的shelter.
God is our shelter.
上帝才是我們的shelter.
在住棚節裡面有很多這樣的神學的反省.
雖然我們未必會走去搭棚.
去守住棚節.
但是.
我們需要明白.
聖經裡面所講的.
告訴我們這些屬靈的.
所謂wisdom.
這些的智慧.
並且後期.
住棚節.
猶太人的住棚節.
他們會在猶太人的會堂裡面.
去供讀某些聖經的書卷.
聖經裡面舊約.
希伯來的聖經.
有三十九卷.
你覺得在住棚節裡面.
猶太人的會堂.
很多時候供讀的是什麼書卷.
哪一卷.
創出利文森 舒斯德 習王 代代拉尼斯等等.
三十九卷.
哪一卷.
是傳道書.
傳道書.
為什麼.
因為傳道書說.
我們人生.
是虛幻.
好像一口氣.
我不講太多.
因為Joyce傳道已經在講傳道書的.
系列.
那裡講什麼.
傳道書最強調的就是.

$^{721}$既然人生在日光之下.
是脆弱的.
你自己應該怎樣去調教你人生的方向.
最值得你追求的東西.
是什麼.
傳道書最後的結論就是.
你追求上帝.
以上帝的事業為念.
以致你人生的方向.
是不會走錯.
你可以忙於你每一日.
由早到晚.
24/7.
裡面所要經營的東西.
但是不要搞錯你最終的方向.
是應該去到敬畏上帝.
曾經有人問我.
就是今年.
他說如果時光倒流的話.
回到十年前.
當初開教會的時候那麼艱難.
你都知道我開教會前.
我開過Uber.
有時候.
有時候那些人會跟我開玩笑.
他說多好玩.
我說好玩.
當你兩三點要去酒吧載人的時候.
當你戶口的時候.
你覺得好玩嗎.
但是有人問我.
如果你時光倒流十年前.
你會再來過嗎.
我幾乎不用思考.
我都會答.
I will.
I will do it again.
你知道我們人生裡面.
特別是基督徒的人生裡面.
有幾何我們能夠去建立教會.

$^{761}$有幾何我們有這樣的福份.
有這樣的機會.
我們能夠去參與建立教會的工作.
這個不是要自豪.
不是要去.
不是要去說我做了很偉大的事.
No.
而是我們有這樣的福份.
有這樣的機會.
有這樣的恩典.
可以參與在上帝的工作.
參與在上帝的工作.
最近大家經常很留意的是甚麼.
世界盃是嗎.
世界盃.
如果你可以和美斯踢一場球.
你願意嗎.
願意嗎.
即使你未有機會碰到那個球.
起碼你可以和他同場.
對嗎.
和他同場和他拍張照.
對嗎.
給他衣服和簽個名.
這是甚麼.
不是因為你聰明.
球你沒有碰.
對嗎.
是因為你和他同場.
你可以和他參與那場賽事.
只是這樣.
出場三分鐘都可以.
就算出場三分鐘被人踢到腳.
踢倒進醫院.
你都願意.
真的不願意嗎.
你以為每個人都可以嗎.
你願意嗎.
就是因為你可以參與.
就有這樣的機會可以參與.

$^{801}$人生好像三分鐘一樣.
如果我可以參與和球王同場的話.
那是多麼的榮幸.
上帝給我們的生命都是一樣.
能夠在這裡建立上帝的家.
能夠在我三分鐘的人生裡面.
我能夠去參與.
在上帝的遊戲之中.
世界球賽.
多大的福分.
棚子很快就會離開.
我們應該怎樣去把握.
待會我們按照我們的系統.
我們每一次都會給大家一些紀念品.
這個瓶子.
如果你今天是新來的朋友.
或者最近才來的.
你不知道在搞什麼.
你可以拿回去.
因為每一次我們系列裡面.
說到這個大故事.
上帝是一個願意祝福世界.
願意使用我們.
給我們機會來祝福世界的.
這個訊息.
我們就放在這個瓶子裡.
每一件紀念品我們放在這個瓶子裡.
放在你的桌上.
你見到我就放在我的電腦桌上.
讓我記得.
是多大的福分.
三分鐘可以參與這場賽事.
但是拿這件紀念品的時候.
請你記得.
想一想.
我最近怎樣經歷過上帝的恩典.
怎樣經歷過上帝的恩典.
這次拿這個紀念品是一支竹.
一支很短的竹.
代表搭棚.

$^{841}$豬棚.
這支竹提醒我們.
好像以色列人過去四十年.
他們走過曠野的時候.
上帝是一個很脆弱的棚.
仍然堅持著他們.
去保守著他們.
這是第一個問題.
第二個問題.
我們應當怎樣去回應.
上帝給我們的祝福和保守.
有沒有停下來想一想.
身體力行我們應該怎樣去回應.
第三我應該怎樣去傳承神的功用.
特別是下一代.
下一代.
他會知道神仍然在的.
好像小孩住在棚內.
他問我們.
為甚麼我們會在這裡.
為甚麼我們來參與這樣的工作.
為甚麼我們不睡晚一點.
為甚麼我們要做這樣做那樣.
我們怎樣回答.
我們回答就是上帝仍然在.
上帝是很真實.
我自己經歷過.
我都希望你能夠經歷.
一些未信主的朋友.
一些我們身邊的一些鄰舍.
我們要告訴他們.
上帝是真實.
上帝是仍然在的.
今天請你拿這個紀念品回去.
如果你沒有那個瓶子的話.
你整個瓶子可以拿走.
我今天預備了一些瓶子.
有這些紀念品在這裡.
作為一個提醒.
我們同心低下祈禱.

$^{881}$上主,我們多謝你.
因為我們能夠和你同場.
有這樣的福分.
有這樣的榮幸.
可以參與你自己的工作.
是一些有永恆價值的事物.
神啊,我們多謝你.
在我們脆弱的人生.
三分鐘的人生裡面.
我們可以有機會.
領受你的祝福,領受你的恩典.
能夠和你一齊去服侍.
一齊去事奉,建立.
你自己的神國的子民.
你自己的工作.
主啊,我們很軟弱.
我們很脆弱的.
但我們求你來使用我們.
多謝你給我們這個權柄,這個機會.
讓主在我們生命當中.
在我們家庭裡面.
流露出上帝你自己將來的國度.
以自己的國度給人的感覺.
是更加真實的.
你是真實的.
你的感覺是更加真實的.
祈禱奉主,彼得.
\newpage



\section{士師記 18:17}
\label{sec:032464}
\textbf{聖經人聖經事:沒有王,各人任意而行 (3) 六百但人搶劫神像,挾持祭司}~~~~黃克勤牧師博士
\newline
\newline
連結: \href{http://www.church.com.hk/acms/content.asp?site=cdc\&op=show\&type=product\&code=032464\&layout=sermon}{\texttt{032464}} ~~~~ 語音日期: 2026-07-14
\newline
\newline
\hyperref[sec:032463]{\small{< < < PREV SERMON < < <}}
~
\hyperref[sec:index_chronic]{\small{[返順時目]}}
~
\hyperref[sec:032465]{\small{> > > NEXT SERMON > > >}}
\newline
\newline
士師記 18:17
\newline
\begin{longtable}{cl}
\hline
\hline
章節 & 經文 (和合本修訂版)\\
\hline
18:17 & \begin{tabularx}{0.7\textwidth}{X} 那窺探這地的五個人上前去，進入裡面，拿走雕刻的像、以弗得、家中的神像，以及鑄成的像。祭司和帶兵器的六百人一同站在門口。 \end{tabularx} \\ \\
[1ex]
\hline
\hline
\end{longtable}
$^{1}$一日過得很快,弟兄姊妹,大家平安.
一日之後我們又讀聖經.
生命樹為大家預備的聖經人,聖經事.
今天來到以色列會有王國人任意而行的第三講.
第三就是六百個彈人.
彈支派的人,六百個.
去搶劫了神像,挾持了祭司.
嘩,這樣也行?.
搶神像.
其實這神像是耶和華的神像.
當然耶和華哪有神像都說不可以雕刻跪拜.
不過他照做.
就是米加做的,跟他媽媽做的.
然後去到那個叫做阿雅祭司.
但是這六百個彈的人出發.
那五個探子帶著他出發.
就來搶劫了神像,挾持了祭司.
這樣也有的嗎?.
真的在聖經裡面沒聽過.
以為這些事在電影裡面才有.
今天是沒有的.
其實如果《士司記》的十七章.
到二十一章如果拍成電影真的很好看.
不過這套電影很血腥的.
很可怕的.
是兒童不宜的.
好了.
又看回這裡.
所拉爾什圖他們就在這裡出發.
嘩,六百人浩浩蕩蕩.
就去拉益攻打人.
搶他的地.
他們經過以法連山地.
又去經過米家那間屋.
又見到那個阿雅祭司.
那個阿雅祭司在聖經裡面出現《士司記》.
但是到現在都不知道他什麼名字.
我不知道他將來在天上是否會見到這個.
阿雅猶大之派的利未人祭司.
好了.

$^{41}$第三就是來到這裡了.
就是第十八章第十一節.
於是彈軸裡面的六百人.
有六百人.
那六百人是什麼呢.
其實是軍隊.
勇士,即是可以打的.
但是否只有這六百人呢?.
不是,不過呢.
他們那裡可能不夠住等等.
我們要找地,於是就去了.
各帶兵器.
刀劍,劍不應該沒有的.
刀多過矛那些.
脫石的機緣,盾牌那些.
六百人彈,很兇啊這班人.
即是叫兇,你看下去就知道了.
各帶兵器.
從索拉和以實圖前往.
這六百人彈的民族就前往了.
上到猶大的基列耶林.
在基列耶林後面安營.
安營之後那個地方叫做.
瑪哈尼丹.
改了名字.
不要緊,這些不解釋了.
上去安營.
然後再上去以法連山地.
由南至北路走,走到上去.
以法連山地.
走回舊路,因為那五個人帶.
然後去到哪裡呢?.
又去到米加的豪宅.
住宅.
就是那裡了,走回舊路.
為什麼呢?這五個探子都算壞了.
真是,很奸的.
從前窺探拉翼地的五個人.
對他們的弟兄說.
各位兄弟,六百位兄弟,大家留意.

$^{81}$我們來過這裡.
這裡的豪宅,就是宅子裡有以弗德.
大祭司那件袍,外袍.
外面那件背心型的以弗德.
上面掛著一個胸牌.
如果是大祭司的話.
有個胸牌是金的.
有以色列十二支牌.
代表著十二塊寶石.
不過有以弗德是象徵祭司.
大祭司.
只有一件,就是以弗德.
我自己做了一件.
真厲害,我自己做.
還有家中的神像.
這裡有很多神像.
這一間很厲害.
有以弗德和神像在裡面.
雕刻的像也有.
鑄成的像也有.
雕刻就是說,譬如用金銀寶石.
或者什麼,木等等.
雕刻,雕出來的.
就成了形.
鑄成的就是雕了一個形之後.
就用泥土吸著它做一個模.
然後打開.
就變成兩個一半的模.
然後就將一些金屬倒進去.
然後吸回去.
或者它吸住,就留一個洞.
將一些溶了的銀流進去.
於是就會變成一個鑄造出來的像.
你們知不知道.
你們知不知道這間屋裡面.
這間屋裡面有以弗德和很多神像.
有些是直接雕刻.
有些是用銀鑄造.
很貴重的,你們知不知道.
想一想,應該怎麼做.

$^{121}$這五個探子很厲害.
其實他們應該說.
弟兄,兄弟,不如我們去搶它,拿它.
這裡有這麼多好東西.
大家想一想,怎麼做好.
「識做」嘛.
現在呢.
就如果用今天的術語.
香港廣東話就是.
你們「識做」嘛.
上一檔走一走.
你們「識做」嘛.
我們「識做」啦.
「識唔識做」呀.
搶它咯,這樣的意思.
五個人就進入米家的住宅.
五個人進去了.
就去到少年利未人的房內問他好.
「識晒路」,記住了.
不愧為探子.
直接就進了他家之後.
通常就是以為有人來到.
獻祭或者什麼給香油求問.
但是可能那邊很靜.
利未人就在自己的房間.
沒有人,他們直接進去到利未人的家.
因為他在那裡住過.
去到利未人的房間裡面.
「有沒有人呀?有,在這裡」.
「哎呀,好嘛,好久沒有見你」.
「哎呀,你們五位呀」.
「你們都是彈支派的嘛」.
「是呀,是呀」.
「又是你回來」.
「是呀,是呀,我們回來找你」.
「上次不是請你和我們求問嗎」.
「我們前程是怎樣的?」.
「你說平安嘛」.
「真是平安」.
「好ye ,好準」.

$^{161}$進去問他好,看他在不在.
那六百人沒有進去.
那六百人就站在門口.
各帶兵器站在門口.
六百人站在門口.
拿著兵器就等他們出來.
窺探地的五個人走進去.
進去了,進去做什麼呢?.
「你好呀,二濟師,多好嘛,平安嘛,平安,好呀好呀」.
「回來呀?」.
「是呀,是呀」.
跟著他五個人二話不說.
就去前面神堂那裡.
將雕刻的像.
以忽德加宗的神像.
雕刻的像.
以忽德加宗的神像.
鑄成的像.
這個鑄成出來的可能是用銀做的.
代表耶和華的像.
是怎樣的呢?沒有人知道.
不過他做了一個像.
象徵耶和華.
都拿了去.
第一點就是.
那幫人去搶劫.
搶劫神像.
搶劫.
搶劫了這些神像.
以忽德全部拿了,搶劫.
「濟師和帶兵器的六百人一同站在門口」.
做什麼呢?.
搶著就拿出門口走.
跟著那六百兄弟站在那裡.
那濟師看見他們拿東西就說「喂喂喂,你們做什麼呀?」.
拿了所有東西的.
那濟師就出來跟他們一起站在門口.
就是這一句.
當然表達得不算很清楚.
「濟師和帶兵器的六百人都在門口」.

$^{201}$其實濟師跟著他們.
「喂喂喂,為什麼拿了神像,你們做什麼呀?」.
「怎麼可以這樣做?」.
就全部出來門口了.
那五個人就進了米家的住宅.
拿了雕刻的像.
以忽德加宗的神像.
鑄成的像.
拿了出來.
這裡重新再講.
於是這個利未人.
就是做了濟師.
這個埃雅濟師.
就問他們.
「你們做什麼呀?」.
「拿了神像出來,你們做什麼呀?」.
「你們偷東西嗎?」.
「為什麼拿出來呢?」.
「搞什麼呀?」.
「你們做什麼呀?」.
他們就回答說.
他們回答他.
說難聽一點.
「這是少年人來的,這個是什麼呢?」.
「都拿少年利未人的房來問他」.
「少年人來的,年輕人」.
他們就說「喂,年輕人」.
現在就不是「濟師」.
「以你為父」.
現在就不是「濟師」.
「喂,年輕人」.
「如果今天我們香港廣東話是什麼呢?」.
「喂,年輕人」.
那個神高神大.
六百個士兵在那裡.
跟那個年輕人利未濟師說.
「不要作聲」.
「收聲」.
「不要說話」.
「用手摀口」.

$^{241}$就是掩著嘴巴.
不要說話.
掩著嘴巴.
這樣嚇他.
「跟我們一起走吧」.
「我們以你為父為濟師」.
「跟我們一起走」.
「做我們的濟師」.
「米家不要做他的濟師」.
「不要做這二發連山地的濟師」.
「跟我們走」.
「那你做什麼呢?」.
「那去哪裡?」.
「不要管」.
「收口」.
「掩著嘴巴」.
「不要再說話」.
很兇的.
「跟我一起去」.
「做我們的濟師」.
「你想想,小孩」.
「你做一家人」.
「一家的濟師」.
「就是米家家裡」.
「自己的濟師」.
「一家私人濟師」.
「還是第二」.
「做以色列一族」.
「一個支派的濟師好呢?」.
「現在你跟我們」.
「不要說話」.
「跟著我們回去」.
「你做我們彈支派的濟師」.
「如果你今天留在這裡」.
「小孩,你在做什麼?」.
「你在做米家一家人家裡的濟師」.
「是有點香火」.
「不過不夠的」.
「如果你跟我們」.
「你是做我們一族」.

$^{281}$「整個彈支派的濟師」.
「誰好些?誰勁些?」.
「我?」.
「這樣的埃雅濟師」.
「真的不要得」.
「濟師心裡喜悅」.
「好啊,這件事」.
「做彈支派的濟師」.
「好過做米家家裡的濟師」.
「有什麼好呢?」.
「不夠大」.
「現在我又上一層樓」.
「由一家的濟師」.
「變成一族的濟師」.
「是一個支派的濟師」.
「更加厲害」.
「這個濟師心裡喜悅」.
「於是他拿鞭拿雀」.
「你拿了東西出來做什麼?」.
「拿著的」.
「這樣啊」.
「去你那裡做濟師」.
「好啊,升級了」.
「升很大級」.
「這個宣傳很好」.
「前景很好」.
「於是他馬上拔起那些爾忽德」.
「拿那些神雀自己拿著」.
「好啊好啊」.
「那你一齊去吧」.
「我做你們的濟師」.
「這些是道具來的」.
「記得全部拿回來」.
「那些人看見他肯跟著他」.
「就讓他拿回去」.
「這些東西都是濟師拿的」.
「你拿吧」.
「他就馬上拿著那些道具」.
「爾忽德家中神像」.
「雕刻的像」.

$^{321}$「進入他們中間」.
「和他們一起」.
「他們就轉身離開那裡」.
「好啊,我們走吧」.
「轉身走了」.
「這個哎呀濟師真的不行」.
「他沒有說等等」.
「我叫我太太和子女先出來」.
「他沒有」.
「拿回那些爾忽德的道具」.
「我們走吧」.
「他沒有想他太太和子女的嗎?」.
「他們轉身走離開那裡」.
「原來妻子兒女」.
「生畜財物都在前頭」.
「就帶出來了」.
「原來究竟是他帶」.
「還是那些彈人把他們帶出來了」.
「不知道」.
「總之妻子兒女」.
「生畜財物全部拿出來」.
「那就開始了」.
「起行了」.
「好,走」.
600人加上這個哩個濟師.
哩個哩個濟師一家就走了.
即是任意而行.
這個濟師又任意而行.
這600彈支派的人搶劫了那些神像.
然後嚇那個濟師.
挾持他一起走.
「你究竟還想不想生存?」.
「你要做一家的濟師」.
「還是做我們一族一個支派的濟師?」.
「族是比支派小的」.
「支派有很多族,然後有很多家」.
「那你想怎樣?」.
「小朋友,少年人」.
「好啊,那當然跟你啦」.
「拿些道具拿在手上然後跟他啦」.

$^{361}$「妻子兒女生後都在前頭」.
「都在前頭令我想起」.
「因為這些不清楚」.
「那時候他們是否已經在那裡呢?」.
「已經帶了出來呢?」.
「這樣很厲害」.
那他就在那裡走了.
離米家的住宅遠離開了.
在米家的近鄰聽到.
「嘩,600人在門口」.
「你猜你怕不怕?」.
「什麼事呢?」.
「看一看,裝一裝」.
「嘩,帶著濟詩歌哥走」.
「這個哥哥仔」.
「嘩,太太,子女全部一齊」.
「還串通了伊忽德全部的神像拿了」.
「怎麼辦呢?快點通知米家啦」.
米家就立刻聚集起來追趕他們.
留意這一句很厲害.
一路追一路叫.
呼叫彈人.
「前面彈的兄弟」.
「我們都是以色列人」.
「彈的兄弟留意留意」.
「等等等等」.
「彈的那600人回頭問米家」.
「你聚集這麼多人來幹什麼?」.
「你幹什麼跟著來?」.
「這麼多人,人馬這麼大,來幹什麼?」.
「為什麼追著我們?」.
真是厲害.
搶劫了神像,挾持了濟師走.
又嚇他,你不肯跟我們?.
你命都沒有,跟我們好一點.
又升了級,多好.
大了,變了名濟師,整個彈都在你掌管.
令他心魔怒放.
米家追出來的時候.
「你帶這麼多人來追我們幹什麼?」.

$^{401}$「這樣都可以,為什麼來追我?」.
米家說.
「你把我所做的神像和濟師都帶走了」.
「我還有所剩的嗎?」.
「你什麼都拿了,我現在沒有剩了」.
「你竟然問我幹什麼?」.
「你還問我幹什麼?」.
「我現在要追回那些東西」.
「人你又挾持了神像,搶劫了濟師,你又挾持了」.
「我當然要追你」.
「還問我惡人先告狀,有沒有搞錯?」.
「我沒有這麼多人,怎樣跟你鬥?」.
「快點站住,不然我不客氣」.
誰知道這班蛋人,你猜他有沒有害怕?.
蛋人對米家說.
第三點.
「你不要使我們聽見你的聲音」.
老老實實告訴你.
閉嘴,馬上走.
不要追著來.
不要使我們聽見你的聲音.
恐怕有些人,性暴的人.
就是一個很暴躁的人.
聽到你的聲音追上來,他會反過來攻擊你.
「如果我的兄弟不高興,他反過來攻擊你」.
「他們是很殘暴的,不要讓我們再聽到你的聲音了」.
「快點走,閉嘴,馬上走,不要追上來」.
「如果不是我們惡的兄弟聽到你的聲音」.
「他反過來攻擊你,以致你全家盡都喪命」.
我怕他反過來殺光你們,我的兄弟.
惡到這樣,厲害嗎?.
六百個蛋人搶劫神像,挾持祭司,還說不要再追來.
「不要讓我們聽到你的聲音」.
「聽到你的聲音,我們有些惡的兄弟,他真的會生氣」.
「他反過來殺光你全家」.
蛋人繼續走.
米家看到.
他們的勢力被自己強盛.
就轉身回家去了.
米家看到他們這麼厲害這麼惡.

$^{441}$看來又這麼好打,又這麼惡.
「不要追來,是不是我殺光你們」.
他唯有頭低低轉身回家去了.
沒有了,什麼都沒有了,被他們全拿了.
一件這樣的事發生了.
所以以色列中沒有王館.
各位電影姐妹,大家看到.
這張結尾告訴我們六百個蛋人搶劫神像.
挾持祭司,細胸甲狼.
不要再追,再追我一些很憤怒的兄弟會殺光你們.
米家就頭低低回去了.
這樣都可以,任意而行.
我們祈禱.
基本上我們心裡面是覺得悲傷的.
看到蛋屬於你的子民.
其實他都認識耶和華的.
他們是用你的名,是你賜給我們地方.
更加回來搶了那些神像.
搶了伊忽德,搶了你的神像.
搶了象徵你的神像,搶了這個祭司走.
還細胸甲狼.
警告人不要再追,再追你會沒命.
我們看到今天很多人非常窮凶極惡.
要得到手中的利益.
只求你自己憐憫這個多苦的世代.
我們看到這個米家,雖然他本身都不是好人.
但是他在這次裡面失去他所有的.
其實有很多人在這些戰亂當中.
在一個不好的環境當中.
在一個沒有王的時代裡面.
那個時代裡面是有很多人受苦的.
所以今天我們這個世界都看到很多人受苦.
是被欺凌,是被壓迫,是不講道理.
是被強壯的人搶奪.
無論是給你有錢,他就搶奪你的金錢.
用他的資本搶你的資本.
有些軍備很強的就用他的武器.
來搶你自己的,搶那些人.
另外搶一些人的生活.
搶他們的財富,搶他們的土地.

$^{481}$主要我們求你真是憐憫.
憐憫這個多苦的世界.
讓我們做基督徒跟隨主的.
是更多高舉愛心.
更多去服侍需要的人.
也更加要堅持深信神你是會審判的.
你是又真又惡的主你看到的,你知道的.
主要是今天來到這個世界的末後的時候.
求我們的主你親自憐憫.
警醒我們更多傳揚福音.
更多高舉愛心,更多服侍需要的人.
求主幫助,垂聽祈禱,奉耶穌名求,阿們.
\newpage



\section{創世記 30:25-31:16}
\label{sec:032465}
\textbf{參與神的作為}~~~~林天行牧師
\newline
\newline
連結: \href{http://www.church.com.hk/acms/content.asp?site=cdc\&op=show\&type=product\&code=032465\&layout=sermon}{\texttt{032465}} ~~~~ 語音日期: 2026-07-12
\newline
\newline
\hyperref[sec:032464]{\small{< < < PREV SERMON < < <}}
~
\hyperref[sec:index_chronic]{\small{[返順時目]}}
~
\hyperref[sec:032466]{\small{> > > NEXT SERMON > > >}}
\newline
\newline
創世記 30:25-31:16
\newline
\begin{longtable}{cl}
\hline
\hline
章節 & 經文 (和合本修訂版)\\
\hline
30:25 & \begin{tabularx}{0.7\textwidth}{X} 拉結生約瑟之後，雅各對拉班說：「請讓我走，回到我的本鄉本土去。 \end{tabularx} \\ \\ \relax
30:26 & \begin{tabularx}{0.7\textwidth}{X} 請你把我服事你所得的妻子和孩子給我，讓我走吧！我怎樣服事你，你都知道。」 \end{tabularx} \\ \\ \relax
30:27 & \begin{tabularx}{0.7\textwidth}{X} 拉班對他說：「願你看得起我，因我占卜得知，耶和華賜福給我是因你的緣故。」 \end{tabularx} \\ \\ \relax
30:28 & \begin{tabularx}{0.7\textwidth}{X} 又說：「請為我定你的工資，我就給你。」 \end{tabularx} \\ \\ \relax
30:29 & \begin{tabularx}{0.7\textwidth}{X} 雅各對他說：「我怎樣服事你，你的牲畜在我這裡變得怎樣，你都知道。 \end{tabularx} \\ \\ \relax
30:30 & \begin{tabularx}{0.7\textwidth}{X} 我未來以前，你擁有的很少，現在卻已大量增加，因為耶和華隨著我的腳步賜福給你。現在，我到甚麼時候才可以成家立業呢？」 \end{tabularx} \\ \\ \relax
30:31 & \begin{tabularx}{0.7\textwidth}{X} 拉班說：「我該給你甚麼呢？」雅各說：「你甚麼也不必給我，只要你為我做這件事，我就繼續牧放你的羊群。 \end{tabularx} \\ \\ \relax
30:32 & \begin{tabularx}{0.7\textwidth}{X} 今天我要走遍你的羊群，把綿羊中凡有點的、有斑的，和小綿羊中凡是黑色的羊；以及山羊中凡有斑的、有點的，都從那裡挑出來，作為我的工資。 \end{tabularx} \\ \\ \relax
30:33 & \begin{tabularx}{0.7\textwidth}{X} 以後你來當面查看我的工資，任何我這裡的山羊不是有點有斑的，小綿羊不是黑色的，就算是我偷的。這就可以證明我是正直的。」 \end{tabularx} \\ \\ \relax
30:34 & \begin{tabularx}{0.7\textwidth}{X} 拉班說：「看哪，就照你所說的做吧。」 \end{tabularx} \\ \\ \relax
30:35 & \begin{tabularx}{0.7\textwidth}{X} 當日，拉班把有紋的、有斑的公山羊，一切有點的、有斑的、有少許白色的母山羊，以及小綿羊中所有黑色的，都挑出來，交在他兒子們的手裡， \end{tabularx} \\ \\ \relax
30:36 & \begin{tabularx}{0.7\textwidth}{X} 又使自己和雅各相隔三天的路程。雅各就牧放拉班其餘的羊。 \end{tabularx} \\ \\ \relax
30:37 & \begin{tabularx}{0.7\textwidth}{X} 雅各拿楊樹、杏樹、楓樹的嫩枝，把皮剝出白色的條紋，使枝子露出白色來。 \end{tabularx} \\ \\ \relax
30:38 & \begin{tabularx}{0.7\textwidth}{X} 他把剝了皮的枝子對著羊群，插在羊喝水的水溝和水槽裡。羊來喝水的時候，牠們彼此交配。 \end{tabularx} \\ \\ \relax
30:39 & \begin{tabularx}{0.7\textwidth}{X} 羊對著枝子交配，就生下有紋的、有點的、有斑的來。 \end{tabularx} \\ \\ \relax
30:40 & \begin{tabularx}{0.7\textwidth}{X} 雅各把小綿羊分出來，讓羊對著拉班羊群中有紋的和所有黑色的。於是他把自己的羊群分開，不叫牠們和拉班的羊混在一起。 \end{tabularx} \\ \\ \relax
30:41 & \begin{tabularx}{0.7\textwidth}{X} 當肥壯的羊交配的時候，雅各就把枝子插在水溝裡，使羊對著枝子交配。 \end{tabularx} \\ \\ \relax
30:42 & \begin{tabularx}{0.7\textwidth}{X} 可是當瘦弱的羊交配的時候，他就不插枝子。這樣，瘦弱的就歸拉班，肥壯的就歸雅各。 \end{tabularx} \\ \\ \relax
30:43 & \begin{tabularx}{0.7\textwidth}{X} 於是這人極其發達，擁有許多的羊群、奴僕、婢女、駱駝和驢。 \end{tabularx} \\ \\ \relax
31:1 & \begin{tabularx}{0.7\textwidth}{X} 雅各聽見拉班兒子們的話，說：「雅各把我們父親所有的都奪去了！他從我們父親所擁有的獲得這一切的財富。」 \end{tabularx} \\ \\ \relax
31:2 & \begin{tabularx}{0.7\textwidth}{X} 雅各見拉班的臉色，看哪，待他不如從前了。 \end{tabularx} \\ \\ \relax
31:3 & \begin{tabularx}{0.7\textwidth}{X} 耶和華對雅各說：「你要回你祖先之地，到你本族那裡去，我必與你同在。」 \end{tabularx} \\ \\ \relax
31:4 & \begin{tabularx}{0.7\textwidth}{X} 雅各就派人叫拉結和利亞到田野他的羊群那裡去， \end{tabularx} \\ \\ \relax
31:5 & \begin{tabularx}{0.7\textwidth}{X} 對她們說：「我看你們父親待我的臉色不如從前了，但我父親的神向來與我同在。 \end{tabularx} \\ \\ \relax
31:6 & \begin{tabularx}{0.7\textwidth}{X} 你們也知道，我盡了全力服事你們的父親。 \end{tabularx} \\ \\ \relax
31:7 & \begin{tabularx}{0.7\textwidth}{X} 可是你們的父親欺騙我，十次更改我的工資，但神不容許他害我。 \end{tabularx} \\ \\ \relax
31:8 & \begin{tabularx}{0.7\textwidth}{X} 他若說：『有點的歸給你作工資』，羊群所生的都是有點的；他若說：『有紋的歸給你作工資』，羊群所生的都是有紋的。 \end{tabularx} \\ \\ \relax
31:9 & \begin{tabularx}{0.7\textwidth}{X} 這樣，神把你們父親的牲畜拿來賜給我了。 \end{tabularx} \\ \\ \relax
31:10 & \begin{tabularx}{0.7\textwidth}{X} 「羊群交配的時候，我在夢中舉目一看，看哪，跳母羊的公羊都是有紋的、有點的、有花斑的。 \end{tabularx} \\ \\ \relax
31:11 & \begin{tabularx}{0.7\textwidth}{X} 神的使者在夢中呼叫我說：『雅各。』我說：『我在這裡。』 \end{tabularx} \\ \\ \relax
31:12 & \begin{tabularx}{0.7\textwidth}{X} 他說：『你舉目觀看，跳母羊的公羊都是有紋的、有點的、有花斑的。拉班向你所做的一切，我都看見了。 \end{tabularx} \\ \\ \relax
31:13 & \begin{tabularx}{0.7\textwidth}{X} 我是伯特利的神；你曾在那裡用油膏過柱子，向我許過願。現在你起來，離開這地，回你本族之地去吧！』」 \end{tabularx} \\ \\ \relax
31:14 & \begin{tabularx}{0.7\textwidth}{X} 拉結和利亞回答雅各說：「在我們父親家裡還有我們可分得的產業嗎？ \end{tabularx} \\ \\ \relax
31:15 & \begin{tabularx}{0.7\textwidth}{X} 我們不是被他看作外人嗎？因為他賣了我們，還吞吃了我們的銀錢。 \end{tabularx} \\ \\ \relax
31:16 & \begin{tabularx}{0.7\textwidth}{X} 神從我們父親所拿走的一切財物，都是我們和我們孩子的。現在，凡神所吩咐你的，你只管去做吧！」 \end{tabularx} \\ \\
[1ex]
\hline
\hline
\end{longtable}
$^{1}$感激主 今天可以有機會和大家一起.
敬拜 一起聽聖後話語.
在此祝賀窩福堂慶祝五十週年.
代表洛恩堂將我們的祝福獻上.
多謝窩福堂多年來對不同的子堂.
都有不同的關心.
我上一次在此講道的時候.
是四年前.
洛恩堂慶祝二十週年.
我們尋根返回母堂聚會.
當時窩福堂借出一個地方.
讓我們可以在此聚會.
所以很感恩.
今天在此知道過去五十年.
大家都經歷了很多不同神的帶領.
在此展望未來的日子.
相信神會繼續去拓展.
求主繼續使用當中的每一位.
在此講道之前.
也想邀請大家做一個動作.
如果你有聖經的話.
我們可以一起擁抱神的話語去禱告.
據聞現在大家比較多用電子版.
所以你可以抱著你的電話.
雖然手感有少許不同.
或者程序表內也印了部分的聖經.
雖然都是較薄.
但都是神話話語.
讓我們一起擁抱神話語去禱告.
將你的話語賜給我們.
讓我們可以在你的話語裡面.
更多地去認識你.
讓我們成為基督的身體的一部分.
讓我們可以一起去敬拜.
一起去彼此造就.
我求你繼續帶領倭福堂.
這個教會繼續走面前的路的時候.
讓他們能夠豐豐滿滿地.
去經歷你的恩典和福氣.
看見你的大能在他們當中.

$^{41}$也幫助我們每一個.
願意躺開我們的心去聽你的話語.
願意被聖靈觸動去改變我們.
讓我們都能夠更常基督.
我將以下時間恭敬的交上.
求你親自向我們說話.
觸動我們的心靈.
改變我們的生命.
聽我們的禱告.
奉靠耶穌基督 得勝的名字祈求.
阿們.
如果要放手.
不去用自己的方法.
拿我想要的東西.
是要有相當大的決心才能做到的.
我的兒子現在十歲了.
快要十一歲了.
正在讀小五.
剛剛放暑假.
大家可以想像.
在香港的孩子讀書壓力很大.
每天追功課.
墨書測驗都不簡單.
所以做功課通常都要講求效率.
要快點做完 又能夠學到.
這才是最重要的.
他做功課其中一樣比較浪費時間.
即是費時間 需要多時間的功課.
就是作文.
中英作文都是比較長時間的.
有一次他要用新的作文方法去作.
就是要寫作一些好像新聞報道的文章.
用英文來寫.
自然花多一點時間.
他做了很久.
他做完之後.
我一看.
我的下巴掉到地上.
不是他的內容.
因為我還未看他的內容.

$^{81}$我的下巴已經掉下來了.
他整篇文章都用經典的新聞字體來寫.
即是S字應該這樣寫.
他就這樣寫.
每個字都慢慢逗出一個字體.
那怎麼搞.
兒子你做什麼.
這份是英文功課.
不是美術功課.
他說這樣寫好看一點.
像寫新聞字一樣.
重點是你的內容寫成怎樣.
不是你的手寫的字寫成怎樣.
這樣寫老師會讚的.
我真的想著他的頭.
你應該是希望你的內容寫得好.
老師會讚.
不是你這樣逗些東西出來.
逗得好看的樣子.
老師會讚的.
說著說著.
他都不被說服.
都要用自己的方法.
得到老師的稱讚.
我估計他心裡知道.
他靠的內容是老師不會讚他的.
所以他就要靠這樣寫才行.
但我估計他這個表現.
都反映了我們當中.
或多或少都有這樣的心態.
就是我們會一意孤行.
用我的方法去得到我想要的東西.
他就覺得我這樣寫.
我會得到老師的稱讚.
我就去做.
我相信在我們個人.
或者在群體的層面.
我們有時候都會想.
我想要些什麼.
我怎樣去得到.

$^{121}$而往往或多或少都忽略了去問.
神想要我們什麼.
神怎樣給我們.
我們很多時候都會用自己的方法.
而不是求問什麼是用神的方法.
舉一個比較極端的例子.
我讀神學的時候在美國讀.
所以都有很多不同的留學生.
在不同的國家去.
都有很多不同的交流.
其中有一位尼日利亞尼的同學.
他說在非洲.
有些牧者去找一些巫師去學巫術.
因為覺得教會沒有什麼人.
如果他能夠行異能的話.
就會吸引到更多人去教會.
所以就去找巫師學巫術.
聽起來都頗可怕.
當然我們不會去到這個狀態.
但是會不會有些時候.
我們都想著.
都是想為教會好.
都是想做些事.
不過會不會是用錯了方法.
而不是用神的方法呢?.
這是我們要有的危機的意識.
免得我們落入一個這樣的陷阱.
我相信可能我們訂立一些個人的目標.
或者對教會的期望.
有些時候我們都會有些偏差.
都會只是想我想要些什麼.
而不是先去求問神想要些什麼.
我們怎樣可以做到呢?.
今天想藉著大家去看《創世記雅各》那段經文.
我們思想一下.
我們怎樣去參與神的作為.
雅各的故事相信很多人都認識.
今天這段是少少生悶一點.
又少少冷門一點.
但是都希望藉著這段經文.

$^{161}$我們一起去更多的思想.
究竟作者想給我們一個怎樣的訊息.
進入今天經文之前.
我們先看一下作者怎樣去寫雅各的故事.
一開始的時候是寫雅各怎樣去欺騙耳掃.
他半騙半買.
用一碗紅豆漿換了這個長子的名分回來.
都有少少詭詐.
不是全身為欺騙的.
但是都是看準一個時機.
趁哥哥很口渴很餓的時候.
就這樣買了回來.
接著他明明地去欺騙爸爸.
欺騙了長子的祝福回來.
當時他哥哥很生氣.
說要殺他.
所以他就離開了.
離家出走的時候.
他就遇見神和神的使者.
大家比較熟悉的就是遇見天梯.
神和他講說話.
這個都很重要.
之後我會再回來看.
然後他就開始漂泊的日子.
他就在這裡流浪.
最後他就去到找到他的舅父拉班.
即是他媽媽的哥哥.
拉班他的舅父.
逐漸就成為他的外父.
接著又成為他的老闆.
但是拉班是一個吃人不吐的吸血鬼.
他對著拉班都遇上不少苦頭.
接著他就想離開.
但他離開的時候就遇見神和神的使者.
他就和神出交過一幕.
大家都記得了.
最後他就見到耳掃.
和哥哥和好.
有一個大團圓結局.
那裡就很感人.

$^{201}$所以作者去寫作雅各故事的時候.
其實是很工整地ABCBA地寫.
而在過程中.
讓我們看見初期的雅各.
是用很多詭詐用自己的方法.
拿他想要的東西.
他未出生的時候.
就拿著哥哥的腳想爭做大.
想用自己的方法去拿.
但這個也很可悲.
因為是全家都沒有理會神所說的話.
在雙子未出生的時候.
在媽媽的肚子裡爭著的時候.
神已經說將來大哥要服侍小哥哥.
但是出生之後.
個個都是爭著做大哥哥.
他去騙長子的名分回來.
都是用自己的方法.
看著機會去跟哥哥做買賣.
他知道爸爸要祝福哥哥的時候.
他就明明地騙爸爸.
假扮哥哥去騙了長子的祝福回來.
一直以來.
他都是用自己的計謀去拿他想要的東西.
那雅各什麼時候作出改變呢?.
在今天這段經文裡面.
他改變了沒有呢?.
讓我們去看看.
他裡面所發生的事.
再從而思想一下.
作者給我們一個怎樣的訊息.
我們先看看有一個不合理的條件.
他開了一個不合理的條件.
跟拉班說他要離開.
拉傑生了弱勢之後.
雅各就跟拉班說.
請讓我走.
回到我的本鄉本土去.
請你把我服侍你所得的妻子和孩子給我.
讓我走吧.

$^{241}$我怎樣服侍你.
你都知道.
拉班對他說.
願你看得起我.
因我占卜特質.
若說似乎給我是因你的緣故.
若說請為我定你的工資.
我就給你.
雅各說要走了.
拉班說好啊.
你開個價.
開個價我就給你.
好不好啊?.
就這樣看好像挺好的.
但是.
如果你知道發生什麼事的話.
你就知道這個桶千萬不要踩.
他現在跟他說.
請為我定你的工資.
我就給你.
他上次說工資是什麼時候?.
就是當年他要娶老婆的時候.
拉班跟雅各說.
雖然你是我的親戚.
但是怎麼可以讓你在這裡白白服侍我呢?.
告訴我.
你要什麼作工資.
他就很清楚地說.
我要你的小女兒拉傑做我的妻子.
好啊.
我答應你.
當然給你.
比你好給外人.
最後睡醒後發現.
為什麼是姐姐呢?.
他說要他定工資的時候.
就是他晚上看不到的時候.
大小交換了.
這個正正就是他這樣騙他爸爸.
是在哄他爸爸年紀老邁.

$^{281}$眼睛看不到的時候.
大小交換了.
他騙了祝福.
他面對著拉班.
他強中自有強中手.
他常用絞計.
比一個更加奸詐的人用絞計去贏了他.
現在再問你要不要工資的時候.
他怎麼對他呢?.
還有沒有更多的計謀呢?.
只用自己的方法?.
只用神的方法呢?.
雅各對他說.
我怎樣服侍你.
你的生畜在我這裡變成怎樣.
你都知道.
我未來的時候 你有的又很少.
現在大量增加.
因為耶和華隨著我的腳步賜福給你.
現在我到什麼時候才可以成家立業呢?.
拉班對他說.
你想我做些什麼?.
我說你什麼也不必給我.
只要你為我做這件事.
我就繼續募放你的羊群.
他這句話.
連什麼都不用給我.
有幾分他爺爺亞伯拉罕的氣勢在這裡.
當年索多瑪王和亞伯蘭.
當年未改名的亞伯拉罕.
說「你將人給我.
財物你自己拿去吧」.
亞伯蘭對索多瑪王說.
「我只著耶和華至高神天地的註起誓.
凡是你的東西.
就算一根線,一根鞋帶.
我都不拿」.
免得你說「是我使亞伯蘭付足」.
亞伯拉罕當時說.
「你什麼我都不要」.

$^{321}$現在亞國也說.
「你的東西我什麼都不要.
你什麼都不用給我」.
似乎有少少像他爺爺那種.
「我要依靠神,不依靠你讓我付足」.
那種氣勢在這裡.
作者究竟怎樣說呢?.
我們仔細看下去.
他在說什麼要求呢?.
「今日我要走遍你的羊群.
將綿羊中有頂的,有斑的.
和小綿羊中凡是黑色的羊.
以及山羊中凡是有斑的,有頂的.
都可以從那裡挑出來.
作為我的公知.
以後你當面去查看我的公知.
任何我這裡的山羊.
不是有斑,有頂的.
小綿羊不是黑色的.
就算是我偷的.
這樣就可以證明我是正直的」.
先不要說他開了什麼盤.
他開了一個要求.
他說要這個東西去證明我是正直的.
但作者沒有告訴我們.
究竟他是否正直.
但是我們對於經文的詮釋.
是很重要的.
就是要知道雅各是否正直.
他開了一個什麼盤呢?.
就是在羊群裡面你選擇.
選擇有斑,有頂的才是我的.
不是有斑,有頂的話.
如果我拿走的話.
這就算是我偷的.
拉班說「看吧 照你所做」.
當日拉班就將有斑,有頂的山羊.
一切有斑,有頂的.
有少許白色的沒有山羊.
以及小綿羊,中細的,黑色的.

$^{361}$都挑出來.
交在他兒子的手裡.
又指自己和雅各的羊群.
隔三天的路程.
雅各就放拉班其餘的羊.
將他們分到很遠.
究竟發生什麼事呢?.
羊群一般來講.
通常是白色的,啡色的,黑色的.
淨色的比較多.
但也有些有斑,有頂的.
相對上是比較少的.
現在雅各就說.
在羊群裡面.
我養你的羊.
生出來有斑,有頂的就是我的.
如果是淨色的就是你的.
所以拉班說「好」.
當他說「好」的時候.
他就將有斑,有頂的.
全部給他的兒子.
離開到很遠.
相隔三天的路程.
剩下雅各的就是淨色的羊.
你不用讀遺傳學.
也知道淨色的羊.
即是白色和白色的.
生出來通常是白色的.
要生出來有斑,有頂的.
是很少見的.
而拉班還絕情地.
將所有斑,有頂的分到很遠.
我看你拿了什麼羊出來.
如果白色裡面.
生了有斑,有頂的.
就是雅各了.
如果雅各是拿著白色的走.
一看就知道他是偷的.
這個就清清楚楚地分得出來.
我讀中學的時候.

$^{401}$有一個同學.
我們列了一個黑名單.
是不會和他吃麥當勞的.
我不知道現在的中學生是怎樣.
但是我的年代的中學生男生.
麥當勞薯條是神聖不可侵犯的.
我們吃的時候.
通常我一邊吃麵包.
一邊吃薯條.
或者吃完麵包才吃薯條.
但是這個同學.
他就是先吃薯條.
當他的薯條吃到七七八八的時候.
他就將大家的薯條混在一起.
一起吃.
我說 嘩.
這麼無恥的事情都做得出來.
但是你全部倒出來.
都分不出誰是誰.
這樣就很難搞.
所以羊如果這樣混在一起.
你就混不出來.
但是如果說他在吃的是.
只有他在吃紐紋薯條.
我們每個人都是普通少年.
你混在一起都還分得出來.
那些是你的.
你不要拿走我們那些.
現在雅各分得很清楚.
有班有點的才是我的.
不是有班有點的就不是我的.
分得很清楚.
這樣就表明我們是正直的.
他開了一個盤.
都算是下下風.
很不合理的條件.
有沒有人這麼傻呢?.
為什麼要這樣做呢?.
他怎樣回應呢?.
雅各就將羊樹,杏樹和楓樹的枝子.

$^{441}$將皮去剝出白色的紋條.
使枝子露出白色來.
他就將剝了皮的枝子.
插在對著羊群.
插在羊羊喝水的地方.
羊羊喝水的時候.
他們彼此交配.
他們對著枝子交配.
就長出來是有班有紋有點的.
雅各就將小綿羊分出來.
讓羊對著拉班羊群裡面.
所有帶著黑色的羊群分開.
不和拉班羊混在一起.
當肥狀交配.
他就插枝子.
長出來肥狀的就是有班有點.
如果瘦弱的就不插枝子.
長出來的就是歸拉班的.
很奇怪.
他在搞什麼?.
好像很不拉家似的.
羊樹的皮是布滿班紋的.
杏樹的皮是條狀的紋.
楓樹的皮是很多班點的.
他就將這些有班有點的樹.
割開一層層皮.
露出的底是白色.
將這些露出有白色.
本身是有紋有點的樹葉.
露出有白色的枝子插著對著羊群.
強壯的羊交配.
看著這些露出白色的枝子.
就長出強壯的羊.
但就成為有班有點.
瘦弱的羊他就不插.
任由他們長出瘦弱的.
瘦弱的就是正式的.
那些就歸拉班了.
他所做的事是很不科學的.
拗了頭.

$^{481}$為什麼會這樣做?.
一般人都覺得.
白貓和白貓生出來的就是白貓.
很少白貓和白貓生出來的班點.
他現在在做什麼呢?.
我們試試看原文.
在原文裡面.
「拉班」這個名「拉班」就是白色的意思.
所以白色都是「拉班」.
「枝子」這個字.
原文裡面是「拉班」.
所以和「拉班」的名字都有些相似.
現在他就是說.
他將這些「枝子」「拉班」.
露出白色「拉班」.
就將「拉班」的「羊」由正式變成「有班」,「有點」.
「雅閣」原文是「accode」.
「有文」的是「accode」.
「有點」的是「nacode」.
換句話說.
這些「枝子」白色的位置.
「拉班」的東西.
就變成「有班」,「有點」.
就成為「雅閣」.
他在這裡是在玩一個字義遊戲.
一個workplay.
或者廣東話是在「食」字.
換句話說.
他將「拉班」,「拉班」白色的.
變成「雅閣」.
「有班」,「有點」.
「accode」,「nacode」.
換句話說,他在這裡是一個非常不科學.
你可以說是一個white magic.
是將白色的東西變成「有班」,「有點」的東西.
是用「拉班」的東西.
因為「拉班」是白色.
我用你的東西.
用你的白色.
用你的「枝子」.

$^{521}$將你的東西轉化.
成為「有班」,「有點」.
「雅閣」的東西.
這個方法,如果「拉班」事先聽到.
他絕對不會擔心.
因為這個非常不科學.
非常不合常理.
但是雅閣用這個做法.
他就是用「拉班」的東西.
將「拉班」的東西轉成「雅閣」.
如果用廣東話來「食」字.
你可以想像.
如果「拉班」是姓「白」的「白先生」.
如果「雅閣」是姓「黃」的「黃先生」.
他就是用「白色」的東西.
將「白先生」的東西變成「黃色」.
就成為「黃先生」的東西.
類似是這樣.
或者英文可能是「Mr.White」和「Mr.Brown」.
透過這個方法.
他自己就很富裕.
經文裡面說.
「這人極其發達.
擁有許多羊群奴僕.
婢女樂陀和奴」.
《創世記》30章形容雅閣的富裕字眼.
和之前形容亞伯拉罕.
和形容以撒的字眼很相近.
彷彿作者是在告訴我們.
如何令到亞伯拉罕和以撒富足.
他同樣的賜福給雅閣讓他富足.
雅閣願意開出這個不合理的條件.
是因為他背後.
他相信有一個不能擋的大銀.
是神的大銀.
以致他不需要懼怕.
「雅閣聽見拉班的兒子們的話.
說雅閣把我們父親所有的都奪去了.
他從我們父親所擁有的.
獲得這一切的財富.

$^{561}$雅閣見拉班的臉色.
看呀 待他不如從前了」.
這是一個強中自有強中手.
食人不吐的吸血鬼.
他亦很絕情.
現在他要離開拉班的手.
一點也不容易.
而他可以在拉班的手下.
在這麼低的成功率的情況下.
可以大大的富足.
拉班很不高興.
他的兒子也不高興.
要離開一點也不簡單.
你看電影也知道.
你進了黑社會.
你要離開那麼容易讓你走嗎?.
現在他要走的時候.
他也很害怕.
但是耶和華對雅閣說.
「你要回你祖先知地.
到你本族那裡去.
我必與你同在」.
這裡作者讓我們知道.
神要他走.
所以他之前跟拉班說要離開的時候.
是對的.
是合乎神的心意.
作者在第三十章沒有評論.
只是描述事情的發生.
到了第三十一章的時候.
開始讓我們知道.
背後是神在做事.
雅閣派人去叫他兩個老婆.
對他們說.
「我看你們父親待我的面色不如從前了.
但我父親的神向來與我同在.
你們也知道.
我盡了全力服侍你們的父親.
可視你們的父親卻欺騙我.
十次跟趕我的公知.

$^{601}$但神不容他害我」.
他說.
「有點的歸你作公知.
羊群所生也是有點的.
他若說有文的作你的公知.
羊群所生也是有文的.
這樣神把你們父親的生畜拿來賜給我了」.
接著說羊交配的時候.
「我在夢中一見看.
母羊的公羊是有文的.
有點的 有花紋的.
神在夢中對我說.
雅閣我說我在這裡.
跳母羊的公羊也是有花紋的.
拉班所對你做的一切.
我都看見了.
我是伯格利的神.
你曾在那裡用油.
囂觀的柱子向我許願.
現在你起來.
離開這地.
回你本族之地去吧」.
這裡神說.
「昔日你離家.
徬徨之際.
我怎樣向你顯現.
在天梯裡讓你見到.
你怎樣去回應.
曉油在柱子上.
我怎樣答應你.
我會帶你回來.
我會和你同在.
你現在回家吧」.
雅閣記得.
昔日神在二十八章的時候說.
「祂做夢看見一個柱子在地上.
梯子的端伸到天上.
看那神的使者在梯子上.
上下往來.
耶和華在梯子上說.

$^{641}$我是耶和華.
你祖父亞伯拉罕的神.
以撒的神」.
這是神第一次親自和雅閣說話.
當年祂對他們說.
「我是你爺爺.
是你爸爸的神」.
他們很自然就會想.
「這個是不是也是我的神呢?.
我帶我兒子找學校.
說這是爺爺的學校.
也是爸爸的學校」.
他們自然會問.
「這個會不會是我的學校?」.
神當時應許雅閣.
「你看見的地方.
你躺在那裡.
都會是你的.
你的後裔也會很多.
萬眾會因你得福.
看啊.
我必與你同在.
無論你往哪裡去.
我必保佑你.
領你歸回這地.
我還不離棄你.
直到我實現了對你所說的話」.
當年雅閣張信張義.
不太肯定.
經歷了那麼多之後.
他記得神應許的事.
他現在回應.
「我要回去」.
神當年在我離開的時候.
叫我要回去.
我現在回去.
神當年答應我.
祂必保佑我.
與我同在.
領我歸回.

$^{681}$我現在要憑頌心.
去讓這個應許銳現.
拉潔和尼雅回應雅閣.
「在我們父親的家裡.
我們還有什麼份可以給他產業.
我們不是被他當作外人嗎.
他賣了我們.
還吞吃了我們的銀錢.
神從我們父親所拿走的一切財物.
都是我們和我們孩子的.
現在凡神所吩咐你的.
你只管去做」.
在第三十一章的時候.
有三次說到是神所做的.
雅閣自己說.
「神將你們父親的身畜.
拿來賜給我.
神自己說.
拉班向你所做的.
我都看見了」.
他兩個老婆說.
「凡神所吩咐你的.
你只管去做」.
三十章賣了關子.
我們不知道.
三十章很清楚讓我們知道.
是神的作為.
是神的旨意.
雅閣不要說.
「我終於贏了拉班.
我終於騙了他」.
他不要在這裡表示自己的勝利.
因為他承認不是他贏的.
是神的能力讓他得勝.
似乎讓我看見雅閣學了功課.
他知道不是用自己的方法.
去拿他自己想要的東西.
他作出了改變.
我沒有機會再說下去.
所以容許我劇透一下.

$^{721}$可能我要到六十周年才回來跟大家說.
所以不差讓我講完.
之後他和神率交.
神說他贏了.
但他明明輸了.
他贏了誰?.
他贏了自己.
他終於學到了.
不是用自己的方法.
他率交時輸了.
神說要走.
他捉住他.
他不讓他走.
「你不祝福我.
我就不讓你走」.
他抓住神.
他終於知道.
福氣是從神而來.
不是用自己的方法.
而是要抓住上帝.
他見回耳掃時.
大團圓結局.
他送很多禮物過去.
在原文裡.
禮物這個字跟福氣是同一個字.
換句話說.
當年我騙了福氣.
我現在還給你.
他騙了祝福是.
你的兄弟要向你下拜.
他見到他哥哥時.
他向哥哥下拜.
騙了的東西全還給他.
雅各是一個奸狡的人開始.
但他改變了.
我們見到他這個改變開始時.
是因為他抓住神.
他相信神的大能.
他相信要從神得到福氣.
是用神的方法.

$^{761}$不是用自己的方法.
弟兄姊妹.
我們看《亂經文》時.
我們不是單單去模仿跟從.
我們看作者說了什麼.
作者讓我們知道.
他開出了不合理的條件.
我們不是每次都開出不合理的條件.
但這個不合理的條件.
是背後他相信不能阻擋的大能.
是神的能力.
所以他倚靠神.
所以當我們今天去想怎樣應用時.
我們要學習.
不要輕看我們的參與.
我們要參與神的作為.
不是靠自己的方法.
不是用自己的能力.
不是拿到自己想要的東西.
而是倚靠神.
藉著參與神的作為.
我們見證神的大能.
你不要小看你的參與.
不是你做什麼.
是神做什麼.
我們不是要有什麼雄才偉略.
做什麼大事.
因為背後我們知道.
不是我們可以成就什麼.
是神自己親自成就他的工作.
我們要做的是.
我們參與其中.
讓我們見證到神的作為.
我被邀請交代一下諾恩堂的近況.
諾恩在疫情之前.
大概有二百人參加成人崇拜.
但這幾年.
我們都像香港普遍的教會一樣.
面對很多困難,很多風風雨雨.
在一九年的社會事件.

$^{801}$我們大概有十個家庭離開.
接著在疫情期間的移民潮湧現.
我們有大概二十個家庭離開.
所以我們的成人崇拜.
由二百人到一百人.
差不多沒有了一半人.
離開的很多都是有家庭有孩子的.
都是為孩子而離開.
所以我們的兒童聚會.
由本來八,九十人到現在二,三十人.
離開的很多都是領袖.
所以事奉人員都有一個很大的打擊.
來得開可以移民走的.
通常都是經濟會好一點.
所以走了一半人.
風險就不止跌一半.
你見到我們有很多赤字.
青少年都失去了一代人.
因為有一班那時候是讀DSE.
沒有怎麼回來.
接著考完試.
或者期間我們在網上聚會.
接著復常了.
回來後發覺青少年全都去了外國.
有一個資深導師又移民了.
有另一個資深導師又轉了教會.
回到後面目全非.
我們都沒有回來.
整個領袖都沒有了.
我們重新再來過.
遇到很多不同的艱難.
在面對這麼多的困境裡面.
就好像昔日雅各都面對很高的困境.
對著一個很難勝過的仇敵拉班的時候.
可以怎樣做呢?.
我都沒有什麼很特別的方法.
我只能夠跟執事說.
我們就專心做教會應該要做的.
做我們應該要做的事.
不要去計算做了什麼會有什麼回報.

$^{841}$就衷心去做當作的事.
三年前我們重新啟動我們的社區服事.
之前都有做過.
但是做得不是很好.
我們同工落水.
大力地參與.
做很多的預備.
有些不同的培訓.
鼓勵跟不同機構合作.
終於在開放的時候.
我們就開了兩條線.
一條就是做婦女的.
一條就是做長者的.
我們就透過新福士宮接觸到一班婦女.
大部分都是新來港的.
有些還是持有雙程證.
有些是單親媽媽.
我們就開設了婦女的侍工.
我就沒有怎樣參與.
他們聚會的時候我都會走開.
通常都是他們搬家的時候我才幫忙.
燒烤的時候我就燒烤給他們吃.
但是慢慢地就有十幾二十.
大概十至十五個婦女在穩定地聚會.
當中他們都有很多不同的需要.
你都可以想像.
新來港的或者單親的怎麼教呢?.
即是其中有一個單親媽媽.
帶著兩個S型的孩子.
每天都很難走過.
當中建立了很多的信任.
都有人信主.
跟那些女童工.
尤其是我太太.
就比較跟得貼一點.
會是半夜打給她.
我會說.
「我見到我兒子寫了一封遺書」.
「很害怕 怎麼辦?」.
會打給她說.

$^{881}$「老公兩天沒有回家」.
「你可不可以幫我找他?」.
吵架,打架.
扔爛了家裡的東西.
很多的困難.
但卻跟他們同行.
建立了很多的信任.
到現在 過了三年之後.
我們婦女的聚會.
由十五人到零.
我們結束.
因為當中有些慢慢找到工作.
找到工作就回不來了.
弄來弄去地走.
會不會很可惜呢?.
沒有問題.
我常常跟教會童工說.
對的事不怕做得失敗.
最怕的是錯的事做得很成功.
你做了事.
建立了關係.
不會失去的.
你建立了人.
不會失去的.
失去了聚會就不再聚會了.
其中一個婦女.
有一次哭著跟我說.
「多謝教會」.
其實我沒有做什麼.
我只是間接幫忙.
但她說「多謝教會」.
一路無私地付出.
她說.
她在內地.
從來不會有人不求回報地幫她.
她體驗到神的愛.
什麼都夠了.
不是要講什麼人多.
不是要怎樣奉獻多.
只是做當做的事.

$^{921}$改變了生命.
這樣就夠了.
我們服侍長者.
有大概二十五至三十個長者.
穩定地回來.
一開始我們都是跟他們玩的多.
教他們怎樣防跌.
防騙.
怎樣用手機.
跳舞.
玩得開心.
慢慢我們就引入多些福音元素.
正正當當地跟他們查經.
他們差不多有一半人.
從來沒有聽過福音.
我成為第一個跟他們講耶穌的人.
多榮耀啊.
當然他們來得很遠.
所以不是那麼容易信主.
但這幾年來.
都有兩三個信了主.
都很開心.
他們可以很深入.
很坦白地跟我們講.
他們不信主的原因.
這樣也好.
這樣就進一步對話了.
很難得.
我們資源有限.
都是一季帶他們請他們一次喝茶.
他們都不是自在式的.
就是一群人開心.
但你可以想像.
三十幾個長者走進茶樓.
都挺有氣勢的.
試過一次進去的時候.
他們問.
嘩 你什麼事這麼多人來.
其中一個老後記就說.
是啊 我教會請我喝茶.

$^{961}$然後隔壁台的長者說.
我教會不見我教會請我喝茶.
我們建立了一個群體.
很多互相的尊重.
其中有些街坊會跟我說.
他們自己吃飯的時候.
會聊天會說.
為什麼教會會對我們這麼好.
他們為什麼會那麼真誠地關心我們.
弟兄姊妹.
我們在做的事.
未必有什麼回報.
我和執事說.
我們做社區的服務.
未必會令到洛恩壯大.
但我們做教會應該要做的事.
就讓我們進入人群.
見證基督.
不是我們有什麼能力.
但在過程中.
我們見到神的作為.
我們見到神的工作.
這樣就夠了.
今天我是一個外人.
我復航面對一些什麼挑戰.
我不知道.
怎麼走.
我也不知道.
我只能說.
相信大家面前都有不同的挑戰.
有很多的未知.
但願意鼓勵大家.
努力地.
積極地.
去做神要你做的事.
衷心地去回應上帝.
參與在神的工作.
無論是教會的發展方向.
或者你個人的回應.
都在自己的綱外上尊重.

$^{1001}$不是我們有什麼能力.
甚至正正就是我們的無能.
是我們的軟弱.
是我們的破碎.
我們做不到什麼.
我們只能夠參與在神的作為.
但藉著這個參與.
我們可以見證神的大能.
雅歌的故事讓我們見到.
即使面對多大的難關.
祂都可以做.
彷彿勝算很低的事情.
不需要計較.
我們的付出和收入是否情情比.
要想的就是.
什麼是神所喜悅的.
什麼是神的方法.
就讓我們衷心去做.
神想我們做的事.
在結束的時候.
讓我舉多一個例子.
話說兩三年前.
有一晚我帶兒子去游泳.
游泳後在更衣室換衣服後.
我拿起電話.
見到很多錯誤的電話.
我知道大件事了.
我一看.
是有一個我們教會的會友.
他去英國的.
打了很多電話給我.
我馬上回覆.
他爸爸病危.
他請我去跟爸爸傳福音.
他方寸大亂.
但是他離開香港之前.
我曾經和他探望過他爸爸.
相互認識.
我馬上放下兒子回家.
馬上趕去醫院.

$^{1041}$我去到醫院的時候.
我再透過視像給這位弟兄.
了解發生什麼事.
去到世伯已經無話可說.
沒有回應.
家人就算未姓朱.
都知道我來了.
所以我簡單地向他招呼.
我就跟他說.
「世伯 認不認識我」.
「之前探望過你」.
「你的兒子想我來探望你」.
「也想再和你談談」.
「我們所信的什麼」.
我就再跟他說一次福音.
正如我所說.
他給不到任何回應.
他各方面的衛生指數.
都是勉勉強強調的.
所以在這樣的情況.
我都請家人逐一去跟他道別.
最後有什麼話想說的.
都說了.
讓大家沒有遺憾的離開.
無論是在現場的.
或是視像的.
都說了.
說完之後.
他仍然繼續在那裡.
通常我經歷這些情況.
說完了斷斷心事.
他就會很快地離開.
他還在那裡.
然後我就問家人.
他什麼時候開始無話可說.
他說大概中午的時候開始無話可說.
我計算一下.
都六七個小時無話可說.
他可能有話想說.
我就走過去跟他說.

$^{1081}$「世伯 你是不是有話想說?」.
我就按我對世伯的認識.
我就說.
「你是否想告訴家人」.
「你很知道大家都很愛你」.
「你也很愛大家」.
大家明白的.
都收到了.
「你是否想告訴大家」.
「你也很不捨得大家」.
大家知道的.
不過都放手讓你走.
「你是否想告訴大家」.
「叫大家好好的去照顧伯母」.
一定的.
大家都會好好地去照顧她.
你不需要擔心.
我就在這裡按我對她家庭的認識.
我幫她說我顧她想說的說話.
我一說完.
家人就說「心跳停了」.
我陪家人到遺體去到殮房的時候.
我就離開.
我回到家.
我兒子第一句就問我.
「爹爹怎樣?有沒有好消息?」.
你想像他也想像.
經歷到我要趕著去醫院的情況.
我說這是一個好問題.
有沒有好消息?.
很典型的.
我們期望的好消息就是.
我臨終的時候傳福音.
他順了主.
最好他還做了水禮.
那就確信了.
就是這樣.
但往往的結果不是如我們所預期.
去到他表達不到什麼.
但我說有好消息.

$^{1121}$就是我覺得我見證到神的工作.
有沒有這麼巧?.
是我說完.
他就馬上心跳停止.
似乎是他聽到我所說的話.
所以之前我跟他說了福音.
他也聽到.
不過表達不到.
所以之前家人跟他道別最後的說話.
他也聽到.
不過表達不到.
所以當我說了心底的說話的時候.
他釋懷.
他就離開了.
他有沒有順主我不確定.
但這不是重點.
因為不是我做了什麼.
重點是我經歷了神的工作.
我參與了神的作為.
那就夠了.
弟兄姊妹.
讓我們明白.
不是我要成就什麼大事.
而是我衷心做神給我的託付.
讓我參與神的作為.
是基於我相信從神而來的最好.
是基於我相信神的方法比我的方法好.
是基於我相信神的大銀勝過一切.
所以我不用自己的方法.
我不去拿我要的東西.
我是衷心去會願上帝給我的託付.
積極地參與在神的工作.
這就是我和你應該要做的事情.
楊凱彤弟彤.
弟兄姊妹.
我想請你有少許神面前反思的時間.
你有多相信神給你的是最好的.
會不會有些時候.
我們都會用自己的方法.
去拿自己想要的東西.

$^{1161}$無論是一些個人的目標.
或者教會的期望.
你願不願意一一放下.
單單要從神而來的福氣.
用神的方法.
有些什麼你知道是神想你做的.
你會不會願意跟神說.
我會衷心積極毀身.
去完成神給我的託付.
讓我去參與神的作為.
不需要計較有什麼後果.
因為那個回報不是我們預期的.
那個回報是我們能夠經歷到神的同在.
經歷到神的大能.
天主上帝我們再一次在你面前.
去承認你的王權.
讓我們再一次謙卑在你面前.
去承認我們的軟弱.
我們有懼怕 我們有小信.
我們有忘記你的時候.
但求你幫助我們.
開我們的眼睛.
讓我們看見你的大能.
看見你的信實.
看見福氣是從你而來.
讓我們緊緊的去抓住.
讓我們有這個心.
相信你的大能勝過一切.
好叫我們能夠衷心去回應.
祈求你帶領臥福堂.
在未來的日子.
讓我繼續行在你的心意裡面.
讓每一位弟兄姊妹.
都能夠衷心去完成自己的託付.
好讓我們看見.
你是我們當中的.
聽我們禱告奉耶穌名求.
阿們.
\newpage



\section{路加福音 21:25-38}
\label{sec:032466}
\textbf{AI,UAP與末世:僥倖或挺身昂首}~~~~陳傳華牧師博士
\newline
\newline
連結: \href{http://www.church.com.hk/acms/content.asp?site=cdc\&op=show\&type=product\&code=032466\&layout=sermon}{\texttt{032466}} ~~~~ 語音日期: 2026-07-12
\newline
\newline
\hyperref[sec:032465]{\small{< < < PREV SERMON < < <}}
~
\hyperref[sec:index_chronic]{\small{[返順時目]}}
~
\hyperref[sec:032467]{\small{> > > NEXT SERMON > > >}}
\newline
\newline
路加福音 21:25-38
\newline
\begin{longtable}{cl}
\hline
\hline
章節 & 經文 (和合本修訂版)\\
\hline
21:25 & \begin{tabularx}{0.7\textwidth}{X} 「日月星辰要顯出預兆，地上的邦國也有困苦，因海中波浪的響聲而惶惶不安。 \end{tabularx} \\ \\ \relax
21:26 & \begin{tabularx}{0.7\textwidth}{X} 人想到那要臨到世界的事，就都嚇得魂不附體，因為天上的萬象都要震動。 \end{tabularx} \\ \\ \relax
21:27 & \begin{tabularx}{0.7\textwidth}{X} 那時，他們要看見人子帶著能力和大榮耀駕雲來臨。 \end{tabularx} \\ \\ \relax
21:28 & \begin{tabularx}{0.7\textwidth}{X} 一有這些事，你們就當挺身昂首，因為你們得救贖的日子近了。」 \end{tabularx} \\ \\ \relax
21:29 & \begin{tabularx}{0.7\textwidth}{X} 耶穌對他們講了一個比喻說：「你們看無花果樹和各樣的樹， \end{tabularx} \\ \\ \relax
21:30 & \begin{tabularx}{0.7\textwidth}{X} 樹葉一長出來，你們看了自然就知道夏天近了。 \end{tabularx} \\ \\ \relax
21:31 & \begin{tabularx}{0.7\textwidth}{X} 同樣，當你們看見這些事發生，就知道神的國近了。 \end{tabularx} \\ \\ \relax
21:32 & \begin{tabularx}{0.7\textwidth}{X} 我實在告訴你們，這世代還沒有過去，一切都要發生。 \end{tabularx} \\ \\ \relax
21:33 & \begin{tabularx}{0.7\textwidth}{X} 天地要廢去，我的話卻絕不廢去。」 \end{tabularx} \\ \\ \relax
21:34 & \begin{tabularx}{0.7\textwidth}{X} 「你們要謹慎，免得被貪食、醉酒和今生的憂慮壓住你們的心，那日子就忽然臨到你們， \end{tabularx} \\ \\ \relax
21:35 & \begin{tabularx}{0.7\textwidth}{X} 如同羅網一樣，因為那日子要臨到所有居住在地面上的人。 \end{tabularx} \\ \\ \relax
21:36 & \begin{tabularx}{0.7\textwidth}{X} 你們要時時警醒，常常祈求，使你們能逃避這一切要來的事，得以站立在人子面前。」 \end{tabularx} \\ \\ \relax
21:37 & \begin{tabularx}{0.7\textwidth}{X} 耶穌每日在聖殿裡教導人，每夜出城到橄欖山住宿。 \end{tabularx} \\ \\ \relax
21:38 & \begin{tabularx}{0.7\textwidth}{X} 眾百姓清早上聖殿，到耶穌那裡聽他講道。 \end{tabularx} \\ \\
[1ex]
\hline
\hline
\end{longtable}
$^{1}$今天的講題是比較特別一點的.
也是一個進到21章後節.
講關於末世的事情.
我記得上個星期.
當梁牧師講到前半段.
講到當日他們見到耶路撒冷.
聖殿等等的災禍.
今天就擴闊到全球.
前兩天我有幾位好的同工.
就跟我說.
喂 陳牧師.
有東西看啊.
我說有什麼事呢.
他說去到七十一.
會有大搶購.
他說有大件事.
我說什麼大件事呢.
是不是疫情要搶東西呢.
是不是打仗呢.
結果就去了.
嘩 真是很多人在那裡.
還有幾位女同事在搶購.
看到嗎.
背後你認不出他們是哪幾個.
而右邊真是他們幾位姐妹.
就買了這兩箱東西.
當他們進去的時候.
陳牧師你看到多厲害.
接著他們就說.
你先回去吧.
因為要很長時間才能買單.
很長時間要排隊.
究竟是什麼事呢.
是不是疫情來要搶購呢.
不是.
是因為什麼.
七十一.
七七節.
有十天左右.
七十一.

$^{41}$昨天是七月十一日.
哦 原來如此.
因為這個緣故很多人買東西.
但是還有很多東西可以買.
入定貨.
我們是知道的.
如果真的疫情有大件事.
到時候的樣子應該是這樣的.
大家記不記得.
疫情2020年.
買不到東西的.
沒有東西買的.
是這樣的狀況.
真的會這樣.
我們知道.
今天在香港或者全球.
很多新奇的事情正在發生.
例如我們知道.
疫情來了幾年.
又例如現在氣候變化.
大家知道昨天香港多少度嗎.
三十五.
有些地方三十六點多.
大家知不知道上個月.
六月中美歐洲多少度.
四十六度.
西班牙.
知不知道法國和德國多少度.
四十到四十四點多.
知不知道六月最後兩個星期.
死多少人.
熱死多少人.
八千人.
知不知道.
其實全球所有事情都亂了.
我們都知道.
氣候亂了.
因為氣候的亂.
會有暴雨.
有十號風球.

$^{81}$有海嘯等等.
我想大家都知道.
我們曾經在香港試過.
一天有一次還是多少次的黑雨.
試過.
兩次.
知不知道.
當然我們都知道現在全球都混亂了.
或者有很多打仗.
有很多很怪的現象.
以往的世紀從來不發生.
或者核武.
或者AI的工具等等.
但有些我不明白.
什麼不明白呢.
我們知道的.
凡是暑假.
香港都會有大暴雨.
會有八號風球.
九號風球.
十號風球.
都會有停車場是水浸的.
知不知道我不明白什麼.
我不明白為什麼車還會停在那些必定水浸的停車場.
這個我真的不明白.
如果你不明白就回答我.
我每堂都問的.
為什麼呢.
例如大家知道.
在柴灣有個叫做環翠村停車場.
差不多一定會浸的.
但每次浸我都明白的.
都有車停在那裡.
為什麼呢.
我想了很久.
是不是想著可以僥倖.
應該浸不到我的.
又或者覺得.
真的很懶找.
另外來這一場.

$^{121}$可能距離家裡遠一點.
又或者.
應該輪不到我的.
可能有特別的魔法.
即是浸到隔壁的車.
去到那裡突然間封鎖.
接著魔勢出現.
紅海都要分開.
我不明白的.
或者孤陋寡聞.
沒那麼巧吧.
又或者我剛剛拍的沒事吧.
我都是在說.
我是不明白的.
每一次其實我們知道.
有這些災難我們要準備好的.
是不是.
明白嗎.
剛才我說的.
正如現在天氣又會反常.
或者戰爭也都越來越澎湃.
好像幾個國家.
大背後有幾十個國家在打.
又或者這樣.
我不知道大家知不知道.
兩天前又有件事.
除了七七節之外.
再公佈了第四批UAP.
知不知道UAP.
UAP即是UFO.
即是Unidentified Flying Object的新版.
叫UAP.
是什麼呢.
Unidentified.
A是什麼.
Anomalous.
即是.
你是不知道一個異象.
Phenomena.
不單止飛.

$^{161}$在不同的現象.
出現第四批.
很多人問.
陳牧師那些是什麼.
那些聖經是怎麼演繹的.
是不是真的有外星人.
我也會試.
有些方向去.
不單止這樣.
還有AI.
例如AI教父.
一個諾貝爾得獎者.
叫Jeffrey Hinton.
他說到AI已經發展到一個地步.
有百分之十到二十的機會.
人類被AI殲滅.
我不知道大家有沒有看過.
一套叫Terminator.
我覺得好看.
每次看我都幻想自己.
不是諾斯華辛尼加.
那種肌肉.
就是另一個叫John.
被人追殺的人類.
因為我都叫John.
在旁邊說.
原來陳牧師叫John.
但是原來不是一個科幻片.
這個諾貝爾得獎者.
其實三年前已經辭了.
Google.
本來是首席幫助他發展AI.
他公認是AI教父.
還有牛津的數學教授.
叫John Lennart.
是一個基督徒.
他也是非常出名的數學教授.
他在屬靈場覺得.
魔鬼很可能在末世利用AI毀滅世.
這代表什麼呢.

$^{201}$剛才我講到上個星期.
梁牧師去到前一半.
講到耶路撒冷當日被毀.
接著慢慢耶穌講未來的時代.
牽涉到全宇宙終極的事件.
災難的細節.
今天我們要進行這一篇講題.
而這一篇講題是非一般的講題.
所以跟我們旁邊的人說一說.
準備好啊.
每一堂講完之後.
大家的樣子都像你們現在這麼兇.
我們看看究竟有什麼講.
今天我們有三個主點我們會看.
第一個主點.
我馬上進去25節到28節.
我們一起說好不好.
123 日月星辰要顯出異兆.
一個場景.
場景涉及什麼呢.
我先用藍色.
藍色第一個是日月星.
接著還有地上海中還有天勢.
即是包括什麼.
包括所有東西.
異兆都包括了.
無可避免.
無人可以豁免.
他說日月星顯出異兆.
列國有地上的邦國有困苦.
海中的波浪.
可能是說當時將會面臨極端的事情.
地震海嘯風暴.
或者日月星顯出異兆.
即是同一時間.
不簡單.
是涉及什麼呢.
並且可能是宇宙界有關的東西.
當時他描述的狀況.
當時不明白.

$^{241}$但今天我們很明白.
因為我們已經去了地球很多地方.
天勢什麼地方都會有.
而看完藍色看黃色.
當時人的反應會害怕.
方方不定.
見到震動.
而人會魂不附體.
什麼叫魂不附體呢.
中國人有句說話很有趣.
三魂不見了甚麼.
七魄.
你旁邊那個魄.
不是那個魄.
那個是道教思想.
好像什麼都沒有.
不知如何是好.
很怕無助.
什麼都做了.
都沒有用.
末世是這樣.
我們現在很多事情.
在疫情困難.
我們做了有用.
可以嗎.
而到時是沒有用.
做完都沒有了.
明白了.
我當然知道.
除了在這裡描述.
還有啟示錄.
馬太坊等等.
但以理書.
甚至他在其他書卷描述到.
當日發生的事.
可以無助到.
無論物價.
經濟崩潰.
人與人之間的爭鬥.
等等.

$^{281}$不過綠色就出現了.
他沒有停在那裡.
綠色就在那一刻.
人會很害怕.
我們回想以前.
疫情等等初時都很害怕.
但是.
那時不單是這樣.
疫情我們怕都買到東西.
遲一個星期.
這些醋買這麼多沒用.
你不喝的.
你口罩買幾箱.
最多都是種一下.
到時是買不到的.
完了.
沒得補給.
你明白嗎.
你想想如果這樣不斷發生.
人會害怕的.
不過.
耶穌說.
我告訴你的預兆.
就是告訴你不用怕.
你們在摧毀地球.
你們在打仗.
你們可以令到.
溫度上升.
碳排放不理.
你摧毀的時刻.
我要回來.
而回來.
你就不用怕.
因為越去到這樣的狀況.
就等於.
耶穌會臨近.
我們要從驚慌失措到挺身昂首.
我們不單止到未來.
希望我們的世代不發生.
不過.

$^{321}$聖經描述的這個狀況.
真的這幾十年才有.
過去二十年都沒有發生.
這幾十年才有可能.
因為科技的發展.
特別是這二十年.
就有這樣的條件.
所以這裡說一件事.
在那個狀況裡面.
我們更加要知道.
主的日子臨近.
今天我們給你的就是.
我們清晰地用眼光去看看.
不是正正在那個環境的困難.
震動.
宇宙的動盪.
是上帝迎而掌權.
我告訴你將會有什麼事.
災難可以很大.
但是救主越來越臨近.
這個要非常重要.
我們現在不是說科幻電影.
而是那些事情正在發生.
因為我會說多一點.
但是我又會用馬太福音來多個角度.
看看究竟又會是什麼.
我們先看看.
馬太福音二十四章二十九到三十字.
一起讀.
哪些日子的災難一過去.
入口就變黑了.
月亮也不放光.
眾星要從以上墜落.
一切都要震動.
哪時人人子的笑頭要延在天上.
地上的蠻族都要愛護.
他們要看見人子有能力.
有大榮耀.
踏出天上的黃昏.
他要猜疑使者.

$^{361}$用武統的大聲.
將他的算盲.
從四方.
從天這邊.
到天那邊.
都照墜落.
他說入口變黑.
月亮不放光.
我相信不是單單是日蝕.
是同一時間日日夜夜發生.
並且眾星掉落.
什麼狀況會是這樣.
大災難.
打仗.
核爆.
核廢料.
遮蓋了.
不知道.
我不知道.
是今天才可能會發生.
我不是叫大家害怕.
我不是叫大家聽完之後.
去買米放在家裡.
我不是那種.
我描述的狀況.
是以往沒有這個科技.
不會發生的事.
還有他說到.
蠻族都要哭號.
現在還未到.
現在有很多地方受苦.
有很多還在享受.
歐洲已經熱死.
因為不賣冷氣.
歐洲不賣冷氣.
賣冷氣要申請.
他們的目標是要碳排放平衡.
不過有些地方就管你.
你就受苦.
我就繼續碳.

$^{401}$行不行.
是這樣.
不過歐洲佔了全球幾個百分比.
他繼續節省.
是幫不了九十幾個百分比.
就是說還未到蠻族.
不過就快會有些地方受苦了.
他跟我們說.
不單止我們有眼光.
要知道什麼是仰望主.
還有我們要聚集.
他在說天那邊.
天這邊.
我們聚集起來.
我們順道不要怕.
聚集起來.
當日其實疫情也是.
如果大家記得.
如果你們回雅庚.
當時我們知道.
有人繼續有工作.
有糧出.
正.
但是有人沒有工作.
沒有糧出.
我們當時就提議.
一件事就是.
如果你有工作.
又不用上班.
又有糧出.
你先不要這麼開心.
請你拿一個月半的人工.
給你周邊的親友.
沒有工作.
沒有糧出.
你幫他.
行不行.
記不記得.
我們不可以散.
我們要互相守望.

$^{441}$行不行.
更加珍惜.
究竟要怎樣走下去.
有些很重要的東西要知道.
好了.
裡面描述.
看見人之有靈.
大榮耀.
他未降臨之前.
會有一段日子.
我們要度過.
好了.
再看下去.
我們繼續讀.
第六第八節.
一二三.
你們也要聽見.
打仗和打仗的風聲.
他們繼續說.
你說打仗.
以往也有.
全世界大戰.
是現在才有.
這一百年.
第一次世界大戰.
第二次世界大戰.
不單止國打國.
民打民.
一個國家裡面.
現在差不多每個地區都在打.
裡面.
不同的政黨.
示威.
打人殺人.
單說英國.
又競選了.
他的首相.
十年第七個.
很多事情在發生.
以前不是這麼混亂.

$^{481}$這麼恐怖.
並且同一個民族.
不同的看見.
都可以打到裡面.
可以流人血.
饑荒.
有些人很多東西吃.
有些人沒東西吃.
有東西吃的不理沒東西吃.
地震.
瘟疫.
我們知道地球繼續暖化.
那些菌越來越出現.
他說這是災難.
是甚麼.
是起頭.
起頭.
不是完結.
這樣說是甚麼意思.
這樣說是耶穌說了給你聽.
不用怕.
可以嗎.
不用怕.
一定會來.
不過不用怕.
我們知道一個充滿不安的世代.
戰爭衝突.
今天已經多到我們開始冷感.
打了幾年.
死了很多人.
一邊有東西一邊沒東西不理.
以前不是這樣的.
可以嗎.
打仗.
死人.
一會兒我會說到.
人會開始走向那個的狀態.
不單止是這樣.
我想說會說AI整合一下.
AI吹哨者.

$^{521}$一個Google的前高層.
現在不做了.
他說到大家要醒了.
你們上網都找到的.
他說AGI出現了.
AGI和AI又不同.
AGI是Artificial General Intelligence.
AI就是它執行指令.
再生成你要它生成的東西.
OK嗎.
現在很多設計師不用做了.
你給他一說.
搞定了.
開車都不用了.
因為已經有微模.
這個小東西.
香港的AI都很差.
外面很多.
因為不給你用.
什麼叫做AGI.
它已經和你一樣.
有智慧.
有想法.
它自己開始創造.
而他在說.
明年將來30\%的工作沒有了.
以往工業革命.
就是取締勞動力.
可以嗎.
AI就是取締你的努力.
很簡單.
我女兒剛從美國回來.
她男朋友在Meta.
做最厲害的電腦公司.
現在炒到.
人炒到很厲害.
最厲害的大學進去了.
以前是鐵飯碗.
你知不知道現在是怎樣.
她和她男朋友.

$^{561}$男朋友在裡面工作.
他們聊天.
我女兒昨天聽完第一堂.
我都告訴她.
老爸你說話還不夠厲害.
現在更厲害.
她男朋友和她聊天.
等等.
用幾秒.
和14個AI人員.
14個.
指令他工作.
他寫文章.
本來是14個同事.
現在不用了.
一個人.
15個變1個.
誰懂得用AI就取締了.
還有恐怖就是.
現在他要請人去訓練AI.
訓練一年之後.
訓練的人就沒有工作了.
你做不做.
不做一年都沒有了.
很厲害.
他說你醒了.
這個出口者.
不是這麼簡單.
不只是方便.
接下來.
不單止這樣.
剛才我說的.
人工智能教父.
諾貝爾得獎者.
他三年前又離開了Google.
為什麼.
他說很不妥.
我要告訴別人.
很不妥.
再發展下去.

$^{601}$發展下去.
已經去到一個地步.
不單止會毀滅.
我們人與人之間.
裡面會不停開始.
有些我們估計不到的傷害.
他有很多訪問.
Fortune有訪問.
很多人訪問.
我看了很多訪問.
其中一件事.
我做研究.
現在最新的.
接著那個訪問人問他.
如果我問你.
你下一代.
你要去讀什麼.
做什麼.
才不會被AI逮捕.
你知道他怎麼回答嗎.
他說我就77歲了.
我的影響不大.
不過你問我.
下一代要做什麼.
我回答不了.
回答不了.
於是我看看Elon Musk.
Elon Musk就真的有下一代了.
77歲的下一代.
可能會有四五十.
Elon Musk你也知道.
他越大就越生.
現在最小的也是幾歲.
我估計他還有很長時間生.
有人訪問Elon Musk.
下一代有什麼讀有什麼做.
停了.
回答不了.
我就讀他喜歡讀的東西.
沒什麼用.

$^{641}$如果真的這樣發生.
我希望不是.
也不想在我的年代.
希望以後也不會.
如果真的是.
就是牛津教授John Lennon.
一個很出名的護教學.
也是很厲害的教授.
他說末世撒旦將會用AI摧毀世界.
我又說一下第二個東西.
UAP.
剛才說的不明物體.
很多人問.
末世代表什麼.
聖經代表什麼等等.
我估計有幾個可能性.
第一個可能性就是.
其實這些現象.
都是某國家的武器.
他不承認就最好了.
第二個可能性.
就是有些天文現象.
我們還沒認識.
測度到就最好了.
最好了.
當然有可能.
也是一些外星的智慧生物.
接著我說的這兩分鐘.
你未必明白.
你明白不明白就算了.
明白你不明白也沒所謂.
明白你就明白吧.
他是同一維度.
同一宇宙.
OK吧.
第三個可能.
第四個可能.
就是不同維度.
不同宇宙.
不同宇宙就是multiverse.

$^{681}$英文叫universal.
同一維度同一宇宙.
就是說有些不同維度.
進入了我們的維度.
什麼叫不同維度呢.
我講小小的.
就是說我們試過家裡.
留了一些飯.
忘記收拾.
幾天後那些螞蟻.
幾十隻螞蟻圍著東西吃.
試過沒有.
見過這場面嗎.
見過吧.
如果你見過很多螞蟻.
一群.
然後你用手.
壓過去.
那些螞蟻會怎樣.
散啊.
很害怕吧.
為什麼.
因為牠看不到你的維度.
Dimension.
螞蟻通常的維度.
是直線.
平面.
橫向.
牠不是看不到.
牠一直這樣走.
牆都是這樣.
但你上面壓下去.
牠不知道你壓下來的.
所以你壓完牠散了.
你拿回手.
又不見到你了.
你明白嗎.
好了.
維度是一定有.
如果你不相信有維度.

$^{721}$你就不要相信耶穌了.
因為耶穌不是.
我們物質世界三維度內出現的.
耶穌升天.
保羅見到大馬士革.
耶穌出現.
那個出現.
是從另一個維度出現.
保羅說三重天.
記得吧.
所以我們經歷神.
不是這三個維度這麼簡單.
祂進入我們的維度.
我給個例子再說.
我以前說過.
我兒子出生.
2000年.
早一個月.
出生時我很害怕.
如果大家記得.
我還沒準備走老袋.
走到一隻生仔袋.
接著還沒有智能電話.
我又無法通知家人.
我進去了.
凌晨一,兩點.
然後第二天早上四,五點出生.
第二天發現一件事.
就是Pastor Pam.
神叫醒他.
兩點起床.
為我的兒子出世祈禱.
你記不記得.
沒有人通知他.
誰通知他.
神啊.
神是從哪裡來的.
祂的維度進入了物質世界.
跟他說話.
就是這麼簡單.

$^{761}$然後祂祈禱四點祈完.
第二天上班.
他問為什麼呢.
John有很長時間還沒出世.
不是.
我兒子有很長時間還沒出世.
還有四個多星期.
我兒子出世很輕.
五磅多.
但是神提醒他.
所以進入了一個維度.
說完祂祈禱四點幾.
祈完禱.
就是四點幾.
兒子出世就祈完.
OK吧.
好了.
就是說有可能.
不同維度.
不同的宇宙進入了.
又或者.
靈界在不同的維度.
進入了我們的維度.
OK吧.
不過無論怎樣都好.
我們不只是問那些是什麼.
是在問主啊.
你想提醒什麼.
很密集的.
這二三十年發生的.
無論網絡.
電話等等等等.
沒有的.
我們小時候想都想不到.
沒有可能的.
突然間.
哇.
Bluetooth.
Siri出現.
是不是.

$^{801}$沒有的.
現在突然間發現.
還有啊.
以前你覺得祈禱.
什麼我們在天下.
天生愛天父.
祂就會聽你講.
現在都很容易明白了.
你都不用碰電話了.
Siri這樣.
是不是.
什麼什麼捅蛇.
什麼什麼.
蛇蟻捅蛇.
是啊.
其實.
現在人都發明到.
這種的.
這樣的通話的方法.
你明白我說什麼嗎.
好了.
那是什麼呢.
突然間發生了.
現在啟示說.
我們買東西到時候.
都不用金錢.
用密碼而已.
以前不知道什麼叫密碼嘛.
現在有了.
還怎麼拿錢啊.
是不是.
你嘟而已.
以後不用嘟了.
你走過去認了樣子.
為什麼會發生呢.
為什麼會這麼密集.
就是這幾十年.
我們一定問這個問題.
否則我們錯過了.
耶穌藉著這件事.

$^{841}$想跟我們講的.
而很有趣的.
第27次.
剛才我們讀了.
那種驚慌.
那種混亂.
不用怕.
轉折就是.
我們要看見.
人只有能力.
有大雲.
要駕雲降臨的時候.
就是說.
世界可能未來會問.
如果真的.
科學家發生.
未來五到十年.
是很大件事的.
如果是.
希望不是.
我不是在發預言.
我也不知道.
我在說現在看到的東西.
就他們說.
完了.
世界末日.
但是對我們來說.
正啊.
耶穌回來了.
夠不夠.
救主回來了.
仰望主.
這個是.
耶穌很想告訴我們的事.
而在這個狀況.
祂說.
挺身昂首.
我們更加要知道.
人只有能力.
有大雲.

$^{881}$要當我們搞到.
世界大混亂.
打仗大混亂.
什麼都混亂的時候.
得贖的日子近.
這個不是毀滅.
是一個更新.
新天新地將開始.
但是之前.
我們會經歷.
陣痛.
OK.
好了.
這個時刻第一點我們要反思.
就是面對危機.
檢視自己的反應.
是不是嚇得魂不附體.
就算疫情也是.
是不是魂不附體.
還是我們知道上帝掌管.
挺身昂首.
當新聞讓你焦慮的時候.
我們知道.
其實自己這次會過去.
但是我們主是不過去.
永恆的眼光.
解讀今天的困難.
人經常覺得自己可以解決所有的挑戰.
That is wrong.
一定不可能.
還有.
信心的根基.
基督徒的優勢不是知道所有答案.
而是知道誰掌管答案.
這個就重要了.
記住.
我們和日本主義.
和這個世界最不同的地方.
他們以為人可以解答所有問題.
我記得十幾年前.

$^{921}$和我兒子去迪士尼.
在奧蘭多美國.
其中去到Tomorrowland.
一進去已經說.
人類的科技有多偉大.
我們有什麼不可以解答和解決.
所以我們要繼續去研究.
因為人就可以這樣帶進未來.
我聽著聽著.
就知道背後是很恐怖的.
好了.
我們去到第二點.
第一點的結尾.
我們一起讀這一句.
1 2 3.
萬有都是靠他做的.
所以跟隔壁說.
不用怕.
萬有都是他做的.
第二點.
我們看下去.
第29到33節.
我們一起讀.
1 2 3.
耶穌由切比喻對他們說.
你們看無花果樹和各樣的樹.
下百年的時候.
你們一看見.
只要有得下天的恩典.
這樣你們看見.
這些事千千要承受.
也看見得神的國恩典.
我實在無數你們.
這些世代還沒有過去.
這些事都要承受.
天地要廢去.
我的話卻不能廢去.
他說天地都要廢去.
我的話卻什麼.
不能廢去.

$^{961}$他有些事要說.
然後他說.
無花果樹和各樣的樹.
無花果樹是什麼呢.
曾經在舊約有代表以色列.
所以有解說.
當以色列再發芽.
以色列已經滅了很多年.
1948年再復國.
有人說這樣很近.
有其中一個可能性.
第二就是還有其他樹發芽.
如果你看見這種季節出現.
你就知道發生了什麼.
如果你看見.
耶穌說末世的情景.
那些狀況開始出現.
你就知道末世快到了.
所以第二樣我們要學.
我們不需要推算日子.
不過要認出季節.
推算不到.
耶穌說他也不知道.
我們也不用推算.
不過耶穌要告訴我們.
我們要準備.
所以從第一個.
我們眼睛要仰望主.
第二我們要辨別出.
這些季節出現.
推算何時發生到.
辨認什麼正發生.
作為我們.
我們不是要精確計算那個日子.
是需要有熟齡的敏銳度.
知道這些時候等於什麼.
我信了主三十八.
七十九了.
以往信主真的沒有這些思想.
沒有這些涉及到.

$^{1001}$末世的什麼越來越像樣.
當然我不知道等於這些.
不過上帝創造的秩序開始大扭轉.
剛才說天,地,海,星月大扭轉.
在神學上也在表示一種大混亂.
什麼叫做敏銳季節呢.
我講一下你就明白.
我們知道每四年有個波就會出現.
那個叫做世界什麼.
背.
季節出現.
大家今早有沒有看.
英格蘭對什麼.
挪威.
我有啊.
昨天早上睡.
四點多先看完.
行不行啊.
我想著他不需要加時.
可以再睡半個小時.
誰知道加時.
看完.
好看.
因為我是捧英國的.
多謝.
挪威也很厲害.
其實五百多萬人打到去.
八強真的超乎想像.
行不行啊.
你知道啊.
那個時候就會出現.
我剛剛三月尾去了日本看什麼花.
櫻花.
你跟別人說.
你真是厲害啊.
為什麼啊.
櫻花什麼時候出啊.
三月尾.
通常出幾天啊.
十幾天而已.

$^{1041}$是吧.
你知道了.
那個時候出.
就有櫻花.
你就做櫻花的事.
是吧.
知道了.
多說一個比喻.
上兩個禮拜.
有個姐妹.
真的很疼我.
秋仙.
季節來了.
要訂食了.
是什麼.
榴槤.
真的.
你知不知道兩個禮拜前的榴槤.
多做到便宜到怎樣.
是平時的六折.
六折啊朋友.
本來兩千元六個.
現在是一千二百元六個.
你不按著季節.
要做季節的事.
是吧.
你就錯過了.
你就要吃貴榴槤.
行不行.
耶穌跟我們說.
你開始知道.
那些事發生.
就要做那些事發生的事.
OK.
所以用這種的季節.
作為比喻.
讓我們知道.
什麼時候來.
不過很慘.
就是兩個禮拜前沒有訂.

$^{1081}$因為剛剛那禮拜.
我家聚集不到那麼多人.
其實買六個.
我一個人都可以照顧你們兩個.
不過我現在減肥.
你明白嗎.
真是很掙扎.
這些掙扎的事還是不說了.
繼續講道理好一點.
他在說.
要我們改變主啊.
什麼時候來改為主啊.
今天你要我在這個季節.
做什麼.
他有說.
剛才我說了.
聚集.
眼光要仰望主.
接下來要分辨.
很多事我們應該去回應.
否則我們就錯過了.
耶穌告訴我們.
要準備的內容.
還有當然我們要珍惜.
我們身邊的人.
告訴你.
時日無多了.
你更加要愛他.
更加要同行.
更加要將福音告訴他.
因為他們不明白.
我們知道.
還有我們要默想反覆.
宣告主的話.
比天地更堅固.
剛才那段經文講得很清楚.
天地都要什麼.
廢去.
我的話卻什麼.
不廢去.

$^{1121}$所以站在神的話語上.
不是站穩在環境上.
非常重要.
第二件事我們要學習的.
所以到今天最後一點.
第三點.
亦都推到越來越緊密.
我們要學習的事.
是三十四到三十六節.
我們一起讀.
預備起.
你們要謹慎.
恐怕因貪食.
遵守被禁食的.
思慮淺盡的心.
那日子就如同.
罔顧不見淫蹈靈魂.
因為那日子.
要這淫蹈傳遞上.
一切居住的人.
那個思慮.
內心還是.
好像跟平時日復日一樣.
做著平時的事.
要去到警醒祈求.
什麼東西會累著我們的心呢.
繼續去肉體的享樂.
繼續沉醉著.
禁觀的麻木.
繼續玩著平時的事.
做著平時的事.
繼續生活的憂慮.
我怎樣賺多些.
怎樣才可以更多更大更靚更好.
錯了.
他說很大件事.
他完全失去了.
上帝提醒我們的內容.
繼續起床滾下手機.
吃飽資訊就去吃飽早餐.

$^{1161}$問題不是滾手機.
問題是我們沒有察覺.
那個時間越來越緊密.
以往我做投資銀行.
認識了很多富裕的人.
近來我常常跟富裕的人聊天.
剛剛有一個很富裕的.
真的相當富裕.
富裕到.
跟他聊天的時候.
我曾經說過.
他很憂心.
差不多depression.
我問他為什麼這麼憂心.
他心想那些錢都幾世都吃不完.
AI來都影響不到你.
因為你不用做了.
我跟他說你為什麼這麼憂心.
你給我一個零頭.
我都夠用了.
零頭馬上改話題.
當然不給你了.
不過他說為什麼這麼憂心.
我常常想怎樣可以賺多些.
繼續安全.
繼續安穩.
他找我牧者當然有原因.
你miss the point.
人生不是繼續要多些靚些大些好些.
你要怎樣用那些工具.
成為你生命裡面.
是有意義的人生.
你搞錯了.
作為信徒亂了路.
耶穌就是說這件事.
你搞什麼呀大家.
清楚一點那些是工具.
那些叫我們找到主.
找到意義人生.
叫我們可以有著那種生命力.

$^{1201}$那種方向力.
不要搞錯事事竟成.
就是很多僥倖覺得.
災難不會輪到自己.
以為還有時間可以放縱.
網羅臨到你都不知道.
剛才有沒有記得讀其中一樣的經文說.
那日子就如同網羅.
突然臨到你知道為什麼嗎.
網羅那個網.
即是這樣.
而網羅即是說你想想.
有些野獸在森林.
那些紀錄片我看過很多.
不過不知道是真是假.
走一走突然有個陷阱.
突然有個網就包住它.
就衝上一棵樹.
就捉了它.
大家好像很欣賞這個.
突然間來到網住了.
它說你不準備.
你突然間就成為網住了.
它警告我們.
之後.
它告訴我們.
要我們可以警醒.
可以渡過.
要從那個.
思慮累住.
到時事警醒.
上上其教.
我們必須分辨和世界的想法.
我們在世界.
我們都活在世界.
我們都會食榴槤.
但我們怎樣分辨到它們.
跌到那個陷阱.
那個重要.
很重要.

$^{1241}$有個例子.
現在人們開始冷感.
你死你事打仗.
你沒東西食你事.
我繼續有得食.
繼續安全.
很恐怖.
剛才講的狀態.
近來我見到一個.
一個人.
他的文章叫林旭智.
挺好看的文章.
巴比倫時代的文章描述.
巴比倫我解釋一下.
歷史上主要是.
五六百年.
當時是一個最大的國家.
歷史是先亞述.
很大.
接著被巴比倫打倒.
接著被誰打倒呢.
波斯.
接著被誰呢.
希臘.
中間還有些短時間不提.
我提一大些.
羅馬.
很大巴比倫.
有個描述我覺得很正.
巴比倫當時有個尼布格利薩王.
歷史又出現了.
描述有一個狀況.
有一次他做了一個夢.
在《但爾里書》第二章.
夢見但爾里解夢.
有一個像.
頭是金的.
身是銀的.
腳是泥的.
代表不同的年代.

$^{1281}$不同的國家.
但爾里解釋.
你是金的巴比倫.
不過你過去有個銀的.
結果被人打敗了.
當時的寓言是波斯.
歷史是神控制著的.
大家不用怕.
我在說有另一個維度.
好了.
他講完之後.
第二章尼布格利薩.
哇這個神勁.
不過第三章尼布格利薩王.
他做了一個像是全金的.
你說我的頭金.
身被人取締.
我就告訴你我要全金.
我不會被人取締.
他就在說.
我不要被取締.
巴比倫王.
我的國度不會過去.
我決定歷史.
我要決定誰是真正的王.
我要和誰下令.
誰就要跪拜我.
誰音樂一響.
就要向我金像跪拜.
不跪拜就跌入火湖.
你記得這段記載.
音樂一響你就要拜.
不拜捉到你.
扔你入火湖.
火炎.
就燒死你了.
當時很嚴厲.
繼續描述.
這個文章寫得又短又精彩.
然後他說.

$^{1321}$事情到了這裡.
問題已經不是拜偶像崇拜.
因為尼布格尼薩.
不單止想做世界王.
他想.
掌控人的敬拜.
他想所有人.
忠誠於他都歸向自己.
如果人要保住性命.
很簡單.
跪下就行了.
官員跪下百姓跪下.
整個巴比倫帝國跪下.
他管轄的地方跪下.
但當時有三個人不跪.
但是你的朋友.
當時他不跪.
他當然要面對火搖.
不過最恐怖的不是火搖.
最恐怖的是當時如果你這樣.
你就會面對孤單.
壓力被眾人注視.
這件事比.
火搖更加離得早.
如果我們今天再去修讀.
不同世界做那一套.
你將會跟他做的事有點不同.
你就要.
堅持.
所以為什麼要聚集.
很重要.
值得注意的是.
巴比倫的官員會主動.
跑去告發他們.
監察他們等待機會.
就在王面前告他們.
所以王一下令官員就監察.
同僚觀察.
捉把柄就捉了他們.
當時.

$^{1361}$他們這三個人.
跟王回答.
我所事奉的神可以拯救我.
即或不言.
最出名.
就算我會燒死.
也不拜金像.
在那一刻.
他不知道會發生什麼事.
不過我覺得很棒.
文章最後說.
有一個王想奪取.
神的榮耀的時代裡.
仍然有人相信.
真正的掌權是神.
最後被證明.
站立不住的不是那三位朋友.
是彌補格拉薩王的驕傲.
金像再高.
始終都是人手所足.
巴比倫最終.
崩潰了.
只有神的國.
永遠不過去.
巴比倫的國家沒有了.
即是今天的伊拉克.
一帶.
當日很繁榮.
所以我們要知道.
在這個世代.
我們要警醒祈求.
活出一個不一樣的生命.
我們要問.
我們跟不信的人在末世.
活的分別是什麼.
什麼東西.
還累著我們的心.
什麼工作.
可能令我們憂慮.
什麼物質.

$^{1401}$我們追求到不能放手.
究竟什麼關係.
令我們扭曲不健康.
我們要刻意安排.
注意.
長長期求.
主啊我要把這些都交給你.
真的.
當我們憂慮的時刻.
我們回到主那.
如果用我剛才的例子.
當我兒子出世的時候.
神會用不同的維度.
進入我們的三道.
其實這裡是三道空間.
他會在這裡顯明.
不用怕.
神是真實.
總結更有趣.
當耶穌說了這麼重量的說話.
耶穌告訴我們出路.
我們要聚集.
我們要有眼光分辨.
我們要辨跡.
並且我們要競爭.
還有.
耶穌做了一件很奇怪的事.
我們一起讀.
三十七,三十八字.
一二三.
耶穌每日在電裡教訓人.
每日在電裡教訓人.
\newpage



\section{民數記 15:1-41}
\label{sec:032467}
\textbf{審判,恩典,敬拜}~~~~林誠信牧師
\newline
\newline
連結: \href{http://www.church.com.hk/acms/content.asp?site=cdc\&op=show\&type=product\&code=032467\&layout=sermon}{\texttt{032467}} ~~~~ 語音日期: 2026-07-12
\newline
\newline
\hyperref[sec:032466]{\small{< < < PREV SERMON < < <}}
~
\hyperref[sec:index_chronic]{\small{[返順時目]}}
~
\hyperref[sec:032468]{\small{> > > NEXT SERMON > > >}}
\newline
\newline
民數記 15:1-41
\newline
\begin{longtable}{cl}
\hline
\hline
章節 & 經文 (和合本修訂版)\\
\hline
15:1 & \begin{tabularx}{0.7\textwidth}{X} 耶和華吩咐摩西說： \end{tabularx} \\ \\ \relax
15:2 & \begin{tabularx}{0.7\textwidth}{X} 「你要吩咐以色列人，對他們說：你們到了我所賜給你們居住的地， \end{tabularx} \\ \\ \relax
15:3 & \begin{tabularx}{0.7\textwidth}{X} 你們要從牛群羊群中取牲畜獻給耶和華為火祭，無論是燔祭或祭物，為要還所許特別的願或甘心祭，或節期的祭，作為獻給耶和華的馨香之祭， \end{tabularx} \\ \\ \relax
15:4 & \begin{tabularx}{0.7\textwidth}{X} 那獻供物的要將十分之一伊法細麵和四分之一欣油調和作素祭，獻給耶和華。 \end{tabularx} \\ \\ \relax
15:5 & \begin{tabularx}{0.7\textwidth}{X} 無論是燔祭或祭物，要為每隻小綿羊預備四分之一欣酒作澆酒祭。 \end{tabularx} \\ \\ \relax
15:6 & \begin{tabularx}{0.7\textwidth}{X} 要為每隻公綿羊預備十分之二伊法細麵，和三分之一欣油調和作素祭， \end{tabularx} \\ \\ \relax
15:7 & \begin{tabularx}{0.7\textwidth}{X} 又用三分之一欣酒作澆酒祭，獻給耶和華為馨香之祭。 \end{tabularx} \\ \\ \relax
15:8 & \begin{tabularx}{0.7\textwidth}{X} 你預備公牛獻給耶和華作燔祭或祭物，為要還所許特別的願，或平安祭， \end{tabularx} \\ \\ \relax
15:9 & \begin{tabularx}{0.7\textwidth}{X} 就要把十分之三伊法細麵和半欣油調和作素祭，和公牛一同獻上， \end{tabularx} \\ \\ \relax
15:10 & \begin{tabularx}{0.7\textwidth}{X} 又用半欣酒作澆酒祭，獻給耶和華為馨香的火祭。 \end{tabularx} \\ \\ \relax
15:11 & \begin{tabularx}{0.7\textwidth}{X} 「獻公牛、或公綿羊、或小綿羊、或小山羊，每隻都要這樣處理； \end{tabularx} \\ \\ \relax
15:12 & \begin{tabularx}{0.7\textwidth}{X} 無論你們所獻的數目多少，照著數目每隻都要這樣處理。 \end{tabularx} \\ \\ \relax
15:13 & \begin{tabularx}{0.7\textwidth}{X} 凡本地人將馨香的火祭獻給耶和華，都要照樣處理。 \end{tabularx} \\ \\ \relax
15:14 & \begin{tabularx}{0.7\textwidth}{X} 若有外人寄居在你們那裡，或有人世世代代住在你們中間，願意將馨香的火祭獻給耶和華，你們怎樣處理，他也要照樣處理。 \end{tabularx} \\ \\ \relax
15:15 & \begin{tabularx}{0.7\textwidth}{X} 至於會眾，無論是你們或寄居的外人都要遵守同一條例；這是你們世世代代永遠的定例。在耶和華面前，你們怎樣，寄居的也要怎樣。 \end{tabularx} \\ \\ \relax
15:16 & \begin{tabularx}{0.7\textwidth}{X} 你們和寄居在你們那裡的外人要遵守同一律法，同一典章。」 \end{tabularx} \\ \\ \relax
15:17 & \begin{tabularx}{0.7\textwidth}{X} 耶和華吩咐摩西說： \end{tabularx} \\ \\ \relax
15:18 & \begin{tabularx}{0.7\textwidth}{X} 「你要吩咐以色列人，對他們說：你們到了我領你們進去的那地， \end{tabularx} \\ \\ \relax
15:19 & \begin{tabularx}{0.7\textwidth}{X} 吃那地的糧食時，要把舉祭獻給耶和華。 \end{tabularx} \\ \\ \relax
15:20 & \begin{tabularx}{0.7\textwidth}{X} 你們要用初熟的麥子磨麵，做成餅當舉祭獻上。你們要舉上，如同舉禾場的舉祭。 \end{tabularx} \\ \\ \relax
15:21 & \begin{tabularx}{0.7\textwidth}{X} 你們世世代代要用初熟的麥子磨麵，當舉祭獻給耶和華。 \end{tabularx} \\ \\ \relax
15:22 & \begin{tabularx}{0.7\textwidth}{X} 「你們若犯了錯，不遵守耶和華所吩咐摩西的這一切命令， \end{tabularx} \\ \\ \relax
15:23 & \begin{tabularx}{0.7\textwidth}{X} 就是耶和華藉摩西一切所命令你們的，從耶和華命令的那日直到你們的世世代代， \end{tabularx} \\ \\ \relax
15:24 & \begin{tabularx}{0.7\textwidth}{X} 會眾因沒有察覺而犯了無心之過，全會眾就要將一頭公牛犢作燔祭，遵照典章把素祭和澆酒祭一同獻給耶和華為馨香的祭，又要獻一隻公山羊作贖罪祭。 \end{tabularx} \\ \\ \relax
15:25 & \begin{tabularx}{0.7\textwidth}{X} 祭司要為以色列全會眾贖罪，他們就必蒙赦免，因為這是無心之過。他們要因自己的無心之過，把供物，就是向耶和華當獻的火祭和贖罪祭，帶到耶和華面前。 \end{tabularx} \\ \\ \relax
15:26 & \begin{tabularx}{0.7\textwidth}{X} 以色列全會眾和寄居在他們中間的外人就必蒙赦免，因為這是眾百姓的無心之過。 \end{tabularx} \\ \\ \relax
15:27 & \begin{tabularx}{0.7\textwidth}{X} 「若有一個人無意中犯了罪，他就要獻一隻一歲的母山羊作贖罪祭。 \end{tabularx} \\ \\ \relax
15:28 & \begin{tabularx}{0.7\textwidth}{X} 這誤犯罪的人因無意中犯了罪，祭司要在耶和華面前為他贖罪，他就必蒙赦免。 \end{tabularx} \\ \\ \relax
15:29 & \begin{tabularx}{0.7\textwidth}{X} 以色列中的本地人和寄居在他們中間的外人，若無意中犯了罪，都要遵守同一律法。 \end{tabularx} \\ \\ \relax
15:30 & \begin{tabularx}{0.7\textwidth}{X} 但那故意犯罪的人，無論是本地人是寄居的，褻瀆了耶和華，這人必從百姓中剪除。 \end{tabularx} \\ \\ \relax
15:31 & \begin{tabularx}{0.7\textwidth}{X} 因為他藐視耶和華的話，違背耶和華的命令，這人一定要剪除；他的罪孽要歸到自己身上。」 \end{tabularx} \\ \\ \relax
15:32 & \begin{tabularx}{0.7\textwidth}{X} 以色列人還在曠野的時候，發現有一個人在安息日撿柴。 \end{tabularx} \\ \\ \relax
15:33 & \begin{tabularx}{0.7\textwidth}{X} 發現他撿柴的人把他帶到摩西、亞倫以及全會眾那裡。 \end{tabularx} \\ \\ \relax
15:34 & \begin{tabularx}{0.7\textwidth}{X} 他們把他收押在監裡，因為還不知道要怎樣懲罰他。 \end{tabularx} \\ \\ \relax
15:35 & \begin{tabularx}{0.7\textwidth}{X} 耶和華吩咐摩西說：「這人應當處死；全會眾要在營外用石頭打死他。」 \end{tabularx} \\ \\ \relax
15:36 & \begin{tabularx}{0.7\textwidth}{X} 於是全會眾把他帶到營外，用石頭打死他，是照耶和華所吩咐摩西的。 \end{tabularx} \\ \\ \relax
15:37 & \begin{tabularx}{0.7\textwidth}{X} 耶和華對摩西說： \end{tabularx} \\ \\ \relax
15:38 & \begin{tabularx}{0.7\textwidth}{X} 「你吩咐以色列人，對他們說，他們世世代代要在衣服邊上縫繸子，並在邊上的繸子釘一條藍色帶子。 \end{tabularx} \\ \\ \relax
15:39 & \begin{tabularx}{0.7\textwidth}{X} 你們要佩帶這繸子，好叫你們看見它就記起耶和華一切的命令，並且遵行，不隨從自己內心和眼目的情慾而跟著行淫。 \end{tabularx} \\ \\ \relax
15:40 & \begin{tabularx}{0.7\textwidth}{X} 這樣，你們就必記得並遵行我一切的命令，成為聖，歸你們的神。 \end{tabularx} \\ \\ \relax
15:41 & \begin{tabularx}{0.7\textwidth}{X} 「我是耶和華－你們的神，曾把你們從埃及地領出來，要作你們的神。我是耶和華－你們的神。」 \end{tabularx} \\ \\
[1ex]
\hline
\hline
\end{longtable}
$^{1}$好,我們一齊來看一看我們今天這一段的經文,好不好呀?.
這一段經文是在十五章.
我們看十五章之前,其實是應該先看一看.
上一次和大家分享的.
在十三,十四章所講的那些事情.
十三,十四章是講到.
他們差遣十二個探子去看這個應許之地.
但回報的時候是不盡不實.
是扭曲了上帝的真心意.
因為這些不信的探子.
他們只是看眼前各樣的困難.
他永遠是看不到神的信實,神的豐盛.
神的應許就告訴我們.
今天是沒有人可以僭越上帝他自己.
也沒有人可以改變神他自己的作為.
所以讓我們今天在他自己的面前.
都有這個敬畏上主的心.
神宣佈當時是對成年的一代人.
那個是清楚的審判.
他們沒有一個人可以入應他們.
為什麼?因為他們不信.
他們多處埋怨神.
他們是頂撞神.
但是只有迦勒和約書亞可以進入應許地.
而不信神的應許嚴重的後果就是死亡.
也是他們錯過了上帝這個恩典.
我盼望我們人生不要有遺憾.
如果你的遺憾是因為你不遵從上帝的心意.
你去抵擋上帝在你身上要成就的美意的時候.
這個是最大的一個遺憾.
而文訴記去到十五章的時候.
是再一次讓我們看見.
上帝在審判的背後是有恩典的.
是要給年輕的一代他們有盼望.
他們知道他們不是絕望.
不是無望無助無奈地去進入應許地.
但是神要求他們要對神絕對的忠誠.
我們要分辨誤犯和顧犯的這兩類人.
這個什麼叫誤犯顧犯呢.
即是你是不小心.

$^{41}$你是錯誤地.
因為無知或者忽略了神的命令.
所以犯錯.
神是給你仍有恩典可以來贖罪的.
但是故意頂撞神去違反上帝的命令.
真的當神是誤導的.
這一種的人.
神說要在他們當中剪除.
我們今天來看這段經文.
讓我們我試一下用三個大的段落給大家看.
先我們來看審判之後.
其實仍有恩典和盼望的.
我們一起來讀頭那三節.
耶和華吩咐摩西說.
你要吩咐以色列人對他們說.
你們到了我所賜給你們居住的地.
你們要從羊群中取生畜.
獻給耶和華為火祭.
無論是繁祭或祭物.
為要還所許特別的怨或甘心祭.
或節期的祭.
作為獻給耶和華的興香之祭.
頭這三節開宗明義說得很清楚.
但是有一個關鍵點.
要我們在第二節裡面看見.
是你們到了我賜給你們居住的地方.
就是說他們進入英雄之地之後.
神吩咐他們要如此來獻祭.
這個獻祭的要求就是為了新一代的神的子民.
讓他們不要忘記昔日神在他們身上的恩典.
他們可能他們的父母.
他們的長輩.
他們在曠野路裡面.
他們偏行幾路.
那個日子他們過去.
他們沒有機會入英他們.
但是你們這一班新的一代.
他們要好進行上帝的心意.
上帝不因為人的軟弱或者失敗.
他的計劃就不能夠成就.

$^{81}$上帝永遠都是那個超越歷史.
和一切人為錯誤的終極的掌權者.
我們在第四節到十二節的下文裡面.
我們就更加會清楚看見.
這個要求是包含是很完善的.
你看見在這裡分了三個column.
由四節去到第十節.
首先在最左邊的濟生.
這裡特別講到那些獻煩祭.
要獻上的那些是動物.
而右邊那些是掃祭.
是包括那些細面和那些油.
那些油呢.
那些橄欖樹壓榨成橄欖油.
那些細面就是那些麥穗子.
將它變成那些麵粉.
麥的麵粉然後製成這些麵.
我們今天都有蕎麥麵.
這些健康的食品.
昔日他們是去到這個應許地.
是一個流奶與麥之地.
他們在那裡就看見神的豐富的供應.
上帝就說他們進了應許地後.
他們需要不單止獻煩祭.
再加上這些掃祭.
再加上最右邊那些電祭.
電祭是包括那些酒.
酒是由葡萄釀製的.
葡萄都是來自植物的.
所以這個給我們看見.
是一個很完整的.
一個獻祭裡面所包含的東西.
就是要我們對上帝的感恩.
特別在掃祭這裡提到.
這些由麥而製成的麵.
象徵著耶穌的無瑕疵.
柔和,完美的人性和品格.
而這些柔是象徵著聖靈.
基督是被聖靈充滿.
聖靈感恩基督而生.

$^{121}$也是囂冠耶穌基督.
在地上生活的.
祂作成傳道的工作.
完成上帝天父的心意.
為人類寫明十佳料贖罪.
如果你由上而下.
由第四節一直看到第十節.
這些祭身是越來越貴重的.
由綿羊羔到弓羊.
所以你在掃祭或者電祭.
那些分量裡面是增加的.
是相應增加的.
你看那個分量是多少.
十分一依法這個細面.
這個依法是那些舊約.
他們量度那些乾糧那些食物.
一個成年人他們每一天.
那個口糧的分量.
是十分一的依法.
十分一的依法.
所以你說一個成年人的細面.
那個分量.
而這個一因是說四公升的.
四分一因就是多少公升.
就是一公升.
一公升有多少.
都很多.
你現在買一瓶的那些礦泉水.
有一般來說都是三百或者五百.
再大瓶一些就八百.
如果是一公升的礦泉水都不少.
所以這個一公升的油.
是淋在那些細面.
是一起獻上.
這個是掃祭.
再加上就是四分一因的酒.
也是一公升的酒.
很和味.
你那個就是獻給上帝.
表明為什麼有電製.

$^{161}$就是毫無保留的.
傾倒的.
在舊約到新約.
聖經是沒有禁止我們去喝酒.
吩咐我們不要醉酒.
有羅亞醉酒就闖大禍.
新約裡面.
保羅多次提醒我們.
不要醉酒.
乃要被聖靈充滿.
酒.
但是在耶穌迦蘭的婚宴裡面.
都將水變為酒.
酒是在一個筵席裡面.
是讓主人家和客人.
一起去享受那個筵席的這份喜樂.
這是預表著我們在屬靈裡頭那種喜樂.
由這個酒.
是象徵神和人之間.
我們在敬拜這個祭祀裡面.
我們甘心樂意的獻上喜樂地.
敬拜神.
所以我們今天敬拜.
將剛才God Alone這首歌唱出來.
我覺得心裡頭很寬恩.
很感受到這個真理.
正正是我們今天這個.
是彎曲薄茂的世代.
我們要很鄭重去抓緊.
而且你看到這件事的combination.
就是和梵祭一起.
與掃祭以及電祭獻上.
是表示獻祭者與做物主.
這一種完美的喜樂的敬拜.
一種心靈我們是親近神.
是甘心樂意的獻上.
而帶出來的.
這個祭祀裡面.
亦都帶給我們一個要求.
對獻祭者的要求在13到16節.

$^{201}$無論是本地人或外人.
本地人就是說.
以色列人入住了神應許之地.
那些寄居在當中的外族人.
都要遵守統一的律法.
這個Torah就是舊約摩西五經裡面.
所強調的神的律例典章.
他們一代傳一代.
都要傳頌下去.
都要來遵守.
統一的典章這個respect這個字.
是表示那個訴訟層面的公義.
是表明我們的上帝.
為何那麼配受我們的尊榮呢.
因為上帝不會以有罪為無罪.
上帝亦都不會偏袒.
以色列人你是我的選民.
所以你們犯了罪.
我都姑且望開一面.
我當看不到.
上帝不會.
無論是以色列人或外地人.
他們都是要遵守統一的律法.
統一的典章.
這個是世世代代永遠的一個定例.
這個是表示上帝.
由舊約至新約一個的心意.
上帝的恩典.
是臨到萬民一齊可以共享的.
上帝在舊約很願意.
透過以色列人成為外邦人的光.
去祝福和將舊恩帶給列國.
可惜.
以色列人並沒有.
很完美地去履行上帝這份的心意.
上帝在新約.
給我們很清楚看到.
主耶穌基督是延續舊約.
這一種的精神.
是將上帝這個完全公義聖潔和慈愛.

$^{241}$成就了十加的救贖.
福音臨到萬民萬族.
耶穌能夠作成的這件事.
是沒有其他任何一個人可以造成.
因為只有耶穌.
for him and him alone.
他能夠是沒有帶著任何一個罪.
死在十字架上.
這個為什麼我們說是好消息呢?.
這個要告訴我們今天地上的人知道.
是人人都可以.
是平等地我們去接受這個救恩.
在上帝的救贖的計劃裡面.
是沒有二等三等的公民.
神願意就是說.
無論你是以色列人與否.
你是神的選民與否.
你都能夠照著上帝的說話.
馴服在祂面前.
我們在敬拜裡面.
就能夠與神建立一個真正的關係.
這個是寶貴的.
弟兄姊妹.
你看到這個段落的一個總結.
我們可以說是獻梵祭或者平安祭.
這些都是好像一個禮物獻過給神.
是配上掃祭和奠祭.
就是一個很完美的.
一個是整全的敬拜的圖畫.
為什麼說這個是一個完美的組合呢.
因為他們進入了應許之地.
他們不只是獻祭生.
他們獻上在那個土地出產的.
這些是掃祭和奠祭.
都是從植物那裡而取得的.
這些的修成.
他們不要忘記上帝的恩典.
接著在第十七到二十一節.
在這一段的經文.
更加是接續前文講的.

$^{281}$他們要獻初熟的墨紙.
作為一種舉祭.
什麼為之舉祭呢.
我們來看看.
在十七到十八節我們一起讀.
耶和華吩咐摩西說.
你要吩咐以色列人對他們說.
你們到了那領你們進去的那地.
進到神領去他們那個地方.
就是應許之地.
在這個應許之地.
他們要進行的.
就說作這個舉祭.
是包含了十九到二十一節.
而十九到二十一節是說.
他們將修成那些墨碎紙.
是製成一些細面.
然後也可以將這些墨粉.
來做成來敲製那些墨餅.
好了 做成這些麵和餅之後.
他們成為舉祭.
舉祭的意思是將這些日常他們飲食的.
這個是他們進入榮美之地之後.
他們的平日.
他們的食物都是包括這些麵和墨紙製成的餅.
他們先獻給神.
雙手恭敬地.
所以叫舉祭.
舉起雙手獻給神.
也是獻給祭司.
讓祭司可以有良.
供給他們.
讓他們能夠專心在聖所事奉.
繼續維持神的子民與神之間.
透過敬拜.
這種有的.
屬靈的關係.
而當然最後他們自己和家人.
也可以一起享用.
那麼你看一看這段經文的佈局的時候.

$^{321}$它是一個平衡的結構.
是十九節和這個二十節的下半節.
就是AA派.
都是說到舉祭.
然後就說將獻的這些磨成的麵和這些餅.
就是作為舉祭獻給神.
四次出現的這個舉祭.
是表示你們進入了英雄之地.
不要忘記你們的修成.
這些的初熟的物質.
就是在地裡頭最首先出產的.
你不要就當這些是我自己努力的修成.
弟兄姊妹.
你可不可以和你隔壁的位置說一聲.
你今天所擁有的一切.
其實全是神的恩典.
不是因為你多了不起.
你可以擁有你今天所有.
好嗎 和隔壁的位置說一句.
我今天擁有的全是神的恩典.
我們彼此提醒.
全是神的恩典.
好.
你們很踴躍.
如果坐在你旁邊的是你的配偶.
你也要說一聲.
全是恩典.
能夠和你的配偶一起去敬拜神.
這個是恩典.
我已經很多年沒有機會和師母坐在一起去敬拜神.
接下來退休再職就開始可以了.
很感恩.
我覺得我們可以一起來去敬拜上帝.
而我們總結一下.
去到二十一節.
這整個段落裡面.
是神審判之後帶出的.
這個恩典.
而這個恩典是透過我們學習.
來藉著這些的獻祭.

$^{361}$你不要忘恩負義.
你要討神喜悅.
獻這些繁祭之餘.
也要將你日常飲食的每一樣東西.
你都不要忘記.
是出於上帝的供應.
所以整個整傳的敬拜圖畫.
就是說.
不只是那些濟生的獻祭.
是包括你平日工作.
你耕種.
甚至你糊口那些食物.
都是來自神的.
我們是一個整傳的敬拜圖畫的話.
如果你套用我們今天的場景.
就不只是禮拜六禮拜日.
你回來講福堂.
在這個禮拜堂裡面才叫做敬拜神.
你平時週間每一天都是敬拜.
你起身的時候就不要忘記上帝.
你呼第一口氣.
享受新的一天的生活.
讓我仍有恩典.
還有精神,有能力.
我可以去工作.
這一切都是要我們向神來頌讚.
所以不只是在宗教的場所裡面.
我們才去敬拜神.
而是我們生活中每一餐飯.
我們不只是在所謂謝飯祈禱的時候.
你在這裡吃真一點.
你吃下午那餐飯.
你嘗試吃真一點.
慢一點吃好不好.
不要在那裡咬青.
原來每一口飯.
你不知道原來飯都有味的.
有滋味的.
每一口的那些菜.
我做牧者已經三十多年了.

$^{401}$我自己很體會.
你今天可以吃到東西.
可以喝到水.
很大恩典.
有時我們去醫院探望一些長期的病患者.
或者一些是癌症的末期者.
連吞一口水都有困難.
更加不要說他吃東西.
所以我們今天擁有從上帝而來的.
這一切給我們的.
我們不要忘恩.
不要負義.
所以接著第二個段落.
就要提醒我們.
既然神在審判之後給我們的恩典.
要透過這些獻祭.
神就提醒我們.
你仍然在恩典裡面.
你要分別為聖.
你要分別誤犯和顧犯的要處理.
如果你是無意.
是因為無知.
你忽略.
你違背了神的命令.
你是有贖罪的.
這些祭祀可以來挽救.
但如果你是明知顧犯.
很大一件事.
我等一下會跟大家來看一看.
這個段落.
由二十二到三十六節.
是說這兩件事情.
誤犯者的贖罪的處理.
如果是匯眾一個群體來說.
就獻一頭公牛毒.
這個繁祭.
再加上掃祭.
奠祭.
再加一頭公山羊.
這個公山羊就是贖罪祭.

$^{441}$公山羊是傳所有.
是獻情.
是這個公山羊.
是代替神的百姓所犯的罪.
他流的血.
遮蓋.
這個也都預表著.
耶穌就是那個完美的贖罪祭.
耶穌是毫無保留.
在十字架上為我們寫明.
而個人的.
只是需要獻一歲的母山羊.
一個山羊羔.
這個當然在價值上來說.
是少於匯眾的那個.
這個就是贖罪祭.
我們必須都是以一個.
祭神的寫明.
來預表著耶穌.
為我們寫明去贖罪.
好了.
這個剛才說的是誤犯的.
我應該怎樣.
回到這裡.
在故意犯罪的那些人.
這件大事.
在三十到三十六節.
在中文的譯本裡面.
我們和修館.
可以譯作故意犯罪.
但有些譯本是說.
明目張膽.
或者是膽大妄為.
膽大妄為.
就是說這一班人.
他那個的意識裡面.
他是神的.
說是神的子民.
但竟然他褻瀆神.
即是他來埋怨上帝到一個地步.

$^{481}$是來褻瀆.
是不當上帝為一個敬畏的對象.
而且藐視耶和華.
還有第三樣.
他違背耶和華的命令.
如果你看這三樣的表現.
這個人.
你覺得他怎樣.
真是該死.
怎樣該死.
結果就是他要剪除.
在刑外用石頭打死.
好了.
如果你就這樣看這段經文.
接著他在三十二到三十六節.
就用一個例證去說.
例如.
是什麼為之明知故犯.
他就說在曠野裡面.
有人在安息日去撿柴.
撿那些柴回家.
被人看到.
被人舉報.
他在安息日做了這些工作.
他就說把他帶出去刑外.
用石頭打死.
有些人看那段經文就很詫異.
馬上大罵.
什麼舊約的神這麼殘忍.
撿些柴而已.
為什麼這樣就違規呢.
什麼事呢.
但你要記住.
那個講到在曠野.
在安息日.
你要留意這個場景.
在曠野.
安息日.
他們有沒有東西吃.
他們在曠野.

$^{521}$他們每天.
得到上帝的供應.
是吃些什麼.
嗎哪嘛.
嗎哪是什麼時候給他們的.
每天清晨.
他們出去撿回來.
就給自己和家人吃.
是很平均的.
在營裡.
這麼多人.
都不會缺乏的.
神是信實的.
但安息日.
上帝有沒有過嗎哪.
那怎麼辦.
星期六安息日.
星期五那天.
他們就要收傷份.
傷份.
這個意義是什麼.
就是他們不要來作供.
他們結下一切.
為了他們要來全家.
一起去敬拜神.
是這樣的意思.
是神和他們納的藥.
而這個藥.
是一個盟約.
這個盟約.
對於我們今天.
所以你說他撿柴.
不只是撿柴回家.
他把柴儲在家裡.
做什麼.
他撿柴回去做什麼.
準備開火.
點火.
生火.
然後BBQ.

$^{561}$是這樣的意思.
所以是不容許的.
絕對是違規的.
而且這個是明明的.
漠視上帝的命令.
上帝說不可以在安息日作供.
但你竟然這樣做.
所以這件事對他們來說.
是很嚴重的.
很嚴重.
好了.
我們今天去到新約的時代.
請問我們.
是不是守禮拜六的安息日.
不是.
那你守禮拜幾.
我們說主日.
我們分別為聖.
因為我們說的是.
紀念耶穌復活的主日.
這個是初月教會.
一直至今都是這樣.
但是仍然我們要處理一個問題.
在星期日也好.
星期六也好.
我們有四堂崇拜.
你說禮拜六.
禮拜日.
我們都有一些頂子妹.
是因為輪班的緣故.
她回不來崇拜.
有沒有這樣的可能.
絕對有.
有些可能是消防員.
有些可能是醫生.
有些可能是護士.
我們說.
他們違規了.
是不是.
新約裡面說的安息.

$^{601}$是一個更美的安息.
是因為耶穌為我們死而復活.
上帝仍要我們.
是守住這個安息的原則.
就不是只是驅離在哪一天.
那個西書都有這樣跟我們說.
是因為耶穌基督.
叫我們生命裡得著豐盛.
而且得著自由.
這一種得著釋放的.
不再被任何的這些儀節.
去驅離於一些宗教的律例.
儀節裡面.
我們今天都不需要帶牛羊來.
但是你的心你要來.
你今天奉獻的金錢.
你的時間.
你的恩賜的服侍.
所有這些背後的原則.
是甘心樂意.
你完全的擺上.
你不會說你今天奉獻.
哎呀.
我擺錯了.
我應該是擺100元的.
我擺了500元進去.
不行.
我起一個奉獻袋.
我找回另外400元上來.
行不行.
你不可能這樣.
你奉獻就是你甘心樂意.
你全心全意.
同樣你對神的服侍都一樣.
你不是在這裡斤斤計較.
所以我們今天說.
處理這件事的時候.
你問我們的牧師傳道.
Ronica.
你放星期幾.

$^{641}$你的例假.
星期一.
牧師你的例假.
星期二.
我的例假.
星期三.
原來我們每個人的例假不同.
為什麼.
因為我們七天都開.
所以我們每一個傳道同工.
每一個星期有一天的例假.
就是你那天.
你沒有什麼其他事.
你要你負責.
所以你在那天.
你盡量給自己.
有你休息的時候.
讓你可以享受在神的面前.
重新得力.
去休潔.
不過作為我們牧師傳道.
我們知道.
不一定你可能.
你的例假.
你就全然休息.
有些時候人家安息禮拜.
正正是你的例假.
那你怎樣.
對不起.
你去找別的日子.
我要休息.
不可以.
我們作為牧者傳道人就是僕人.
所以我們寧願.
就是我們自己就一些時間.
所以我們當我們願意這樣做.
仍然有原則.
你一定要爭取那個時間.
你要休息就休息.
休息是為了讓我們走更長的路.

$^{681}$好好在你的事奉裡面得力.
你不要這麼快耗盡了自己.
所以這些讓我們一起有這樣的學習.
我們再來看最後這個段落.
既然在恩典裡面.
我們都要分辨.
是誤犯還是固犯.
你在恩典裡面.
你要學習.
真的每一天.
你都去敬拜神.
去敬拜神.
這裡有幾個提醒.
我們生活要做到與神親密的結連.
首先三十七,三十八節.
跟我們講一個.
在衣服上我們怎樣提醒自己.
我們一起讀.
耶和華對摩西說.
你吩咐以色列人.
對他們說.
你們世世代代要在衣服邊上.
縫錘子.
並在邊上的錘子釘一條藍色帶子.
這個藍色帶子.
在原文的這個字.
其實也都是.
有些中文的譯本都會翻譯做紫色.
或者是深藍色.
深藍去到好像紫色一樣.
這個是代表著皇族那種尊貴的顏色.
也都在屬靈上預表著.
藍色.
我們講藍色的天空.
是一個藍色是天上而來的.
一個君尊祭司的身份.
所以祭司袍裡面他們的錘子.
是有這些帶紫色.
藍色的這些帶子.
是表明他們的身份.

$^{721}$是尊貴.
他們是屬於神所命定的.
去事奉他的這些祭司.
但是在舊約裡面.
有三次出現這一類的經文.
在民數記裡面.
是將他變成是針對神的指紋.
我們每個人都有份.
這條藍色帶子.
我這裡用現代的例子來講.
我們不是穿著那些祭司袍.
但是好像那些牧師袍.
他戴著一個褲子.
如果這個牧師穿著那件祭司袍.
他今天走去蘭桂坊.
他在那裡.
跟別人去喝酒.
跟著去醉酒.
跟著在那裡所謂抱.
我們說什麼.
你自己有沒有搞錯.
你自己都不尊重你自己的身份.
這個是一個提醒.
隔壁那個.
有些人穿校服.
穿著校服他都不敢亂來.
因為他出了學校.
他都知道.
他如果被人看到.
譬如你是喇沙仔.
你是佩貞.
或者你是某某一些名校.
你自己更加要檢點你自己.
原來我們的服飾.
是代表著你的一個身份.
不過今天你不要只是這個服飾.
最重要你是心裡頭.
你知道你是神的兒女.
跟著三十九到四十節.
是提醒我們.

$^{761}$我們要在這個睡紙裡面.
我們要一看到.
就要記得神的一切的命令.
不只是記得.
還要遵行.
不要隨從自己的內心和眼目的情慾.
去跟著行人.
這樣你們就必記得.
定遵行我一切的命令.
成為聖歸你們的神.
我們要歸神為聖.
就是我們每天都要提醒我們.
你每天.
你怎樣提醒你自己.
你是神的兒女.
請問大家.
你每一天.
你起身的時候.
你一定會做的一些習慣.
包括什麼.
祈禱.
還有呢.
擦牙洗面.
這些是.
不然你不要出來見人.
好了.
如果你這麼緊張.
你見人之前都要擦牙洗面.
你的儀容.
請問.
你今天回來教會之前.
你有沒有向神請安.
你有沒有向神祈禱.
神有吩咐你什麼.
你讀了神的什麼說話.
有沒有.
有.
是不是真的有.
你不要敷衍我.
你真的有的時候.

$^{801}$你會發覺這個很不同.
我們現在每一天.
屬靈的旅程.
推行這個門徒的.
所謂門女.
門徒的旅程.
是一個旅程.
是我們一起去經歷去體會.
我們每一天.
以前自己靈修就自己看自己的祈禱.
現在靈修看完.
還要和我小組的弟兄姊妹去分享.
分享.
今早我看的經文就是.
約百記二十三章.
然而祂知道我所行的路.
祂試煉我之後.
我必如精今.
原來我們每一天.
我們要走的路.
或者在你現在人生裡面.
你不知道你下一個階段.
你應該怎樣走.
神在這裡給我一個很大的提醒.
和安慰.
是神知道我所行的路.
你不需要擔心.
但你需要求問神.
你需要緊緊的跟隨神.
神給我們叫行的路.
是一個恩典的路.
是一個蒙福的路.
不過你不去將神當為你最親密.
每天要和祂.
不只是請安.
是你二十四小時你的生活.
包括你睡覺.
原來我們都是與神一個親近.
神給我們可以安眠.
都是恩典.

$^{841}$讚美神.
原來神在我們人生不同的階段.
給我們有不同的試煉.
但考驗之後.
我們的信心要變成精今一樣.
你的品格.
你的為人.
你更加的成熟.
更加的成熟.
讓我人生在今年.
其實是我第三十七年.
我出來默會已經三十七年.
所以很多人都不相信我.
要退休.
他說牧師你這麼年輕.
為什麼要退休.
我每次都要把袋子拿出來.
這個是我的身份證.
這個是我的老人卡.
我是要承傳.
是要給下一代.
是要更年輕的一代.
來去接棒.
要鼓勵他們繼續在神的國.
我是要退任.
是這個必須.
是要讓教會走更遠的路.
這個是要全個教會都一起去認同.
一起去擁抱.
所以我們要息息不斷.
門徒繼續不斷的成長.
增長.
和我們一起彼此扶持.
今天你每天有沒有看手機.
有.
答得這麼快.
因為機不離手.
有些人出了門.
發覺忘記帶手機.
馬上第一件事.

$^{881}$他說沒有手機.
他出不了門.
因為他所有錢包起了手機.
所以我們今天.
今天這個時代.
我們看見手機真是.
是一個工具.
但是我將一些最近看的一些研究.
給大家知道.
今年.
有一些國家是倡議.
十四歲以下的小朋友.
不可以用手機.
有些國家也都通過一些法案.
要求社交媒體公司.
禁止十六歲以下的兒童.
和青少年.
使用社交媒體平台.
為什麼.
因為.
他說有些青少年一機在手.
無限上網.
變成網絡成癮.
要好像戒毒那樣.
才戒得掉.
不過我很想補充一句.
這個文章的作者沒有提.
我不是要為青少年去平反.
我只不過.
我認為.
今天這個手機.
不只是青少年的問題.
是包括你和我.
是中年人.
長輩.
現在很多奇英很厲害.
看手機.
真是很懂.
不斷滾動.
滾動到平台.

$^{921}$我都不懂.
他們說牧師.
你不見了手機.
有一隻這樣的平台.
幫你找手機.
原來.
幫我說的那位.
都接近九十歲了.
所以你說我們長者是不是奇英.
是智慧型的.
所以今天我們說.
這個是工具.
你今天可以透過手機去看聖經.
不過.
你不要走入一些.
不應該去的.
一些平台.
一些網站去看.
包括一些暴力.
包括一些色情的地方.
也有學者研究顯示.
手機裡面.
看文章.
遠不及你實體書.
那個專注力那麼強.
所以.
你儘管機不離手.
請你要克制自己.
每天不要超過.
多少個小時.
你最好立志在神面前.
不要超過.
四小時.
或者三小時.
或者一兩小時.
給自己一個極限.
你不需要每個小時.
在那裡查你的WhatsApp.
行不行.
你現在應該.

$^{961}$專心去敬拜.
你剛才一時疏忽.
你無知地.
故犯.
不斷地響.
請你調回靜音.
是不是.
我們都有恩典的.
我們不會馬上.
扔你出去.
找石頭扔死你.
我們需要在這裡.
用新約聖經的原則.
在耶穌基督.
十字架.
恩典下.
有憐憫.
有憐憫之後.
提醒了之後.
我們記住.
要很好的.
我們怎樣.
真的來分別為聖.
最後.
這一句也是這段經文.
最重要的一個.
是結尾的經節.
在四十一節.
「我是耶和華你們的神,曾擺你們從埃及地領出來,要作你們的神.」.
「我是耶和華你們的神.」.
這句說話重覆了兩次.
「我是耶和華你們的神.」.
就是神和我們立約的提醒.
我們永遠都是神的子民.
今天.
你真的要將自己.
好好分別出來.
在恩典中.
你要有這種敬拜.
真正的自由釋放.

$^{1001}$不再被罪嚇制你.
不再被任何的重擔.
壓住你.
不再因為你和某個弟兄或者姐妹之間.
或者你某個家人.
你們的關係破裂.
你在耶穌基督.
那個敬拜裡面.
你要真正得著.
復和.
和醫治.
得著真正全人的平安.
這個Salom.
就是聖經說的恩典的福音.
而我們今天要被提醒.
就是我們歸神為聖.
神將我們分別.
在萬民中分別出來.
作祂的子民.
就是要叫我們生活聖潔.
我們遵行祂自己所吩咐一切的命令.
故此.
我們作為神的子民.
你不要濫用神的恩典.
如果你今天是誤犯神的命令.
你記住認罪悔改.
你就可以重獲真正的自由.
如果你今天你是明知故犯.
請你記住.
你趕快去認罪悔改.
你不要死撐.
你不要當你自己就是皇帝.
你自己就是你.
無論公司也好.
或者你自己的家庭也好.
包括你自己的人生.
其實我們沒有一個人可以做自己人生的主人.
死定了.
我們的主導權在神那裡.
我們要聽命於祂.

$^{1041}$我們要遵從祂的吩咐.
祂的心意.
求主幫助我們.
施恩於我們.
我們一起祈禱.
天父求你幫助我們廣福家.
每一個弟兄姊妹.
每一個家庭.
我們都來遵從你的心意.
我們知道你的審判背後有恩典.
在恩典當中.
你要叫我們有真正的自由和釋放.
我們懂得去悔改去認罪.
讓我們的敬拜是全人的一種敬拜.
是全人的一種自由釋放地.
向神發出頌讚.
求主施恩於我們.
讓我們每天的生活.
都是帶著這種生命裡的自由和釋放.
我們去服侍主服侍人.
跟隨你直至到有一天.
你呼喚我們要歸回你的天家的時候.
我們能夠無悔無憾地.
我們走過人生.
每一里路.
每一個階段.
都是能夠榮神益人.
同心合意祈禱.
奉主明求.
阿們.
\newpage



\section{路加福音 19:11-27}
\label{sec:032468}
\textbf{天國人的使命}~~~~俞大鳴牧師
\newline
\newline
連結: \href{http://www.church.com.hk/acms/content.asp?site=cdc\&op=show\&type=product\&code=032468\&layout=sermon}{\texttt{032468}} ~~~~ 語音日期: 2026-07-05
\newline
\newline
\hyperref[sec:032467]{\small{< < < PREV SERMON < < <}}
~
\hyperref[sec:index_chronic]{\small{[返順時目]}}
~
\hyperref[sec:032469]{\small{> > > NEXT SERMON > > >}}
\newline
\newline
路加福音 19:11-27
\newline
\begin{longtable}{cl}
\hline
\hline
章節 & 經文 (和合本修訂版)\\
\hline
19:11 & \begin{tabularx}{0.7\textwidth}{X} 眾人正聽見這些話的時候，耶穌因為將近耶路撒冷，又因他們以為神的國快要顯現，就接著講了一個比喻， \end{tabularx} \\ \\ \relax
19:12 & \begin{tabularx}{0.7\textwidth}{X} 說：「有一個貴族往遠方去，為要取得王位，然後回來。 \end{tabularx} \\ \\ \relax
19:13 & \begin{tabularx}{0.7\textwidth}{X} 他叫了自己的十個僕人來，交給他們十錠銀子，說：『你們去做生意，直到我回來。』 \end{tabularx} \\ \\ \relax
19:14 & \begin{tabularx}{0.7\textwidth}{X} 他本國的百姓卻恨他，打發使者隨後去，說：『我們不願意這個人作我們的王。』 \end{tabularx} \\ \\ \relax
19:15 & \begin{tabularx}{0.7\textwidth}{X} 他得了王位回來，就吩咐叫那領了銀子的僕人來，要知道他們做生意賺了多少。 \end{tabularx} \\ \\ \relax
19:16 & \begin{tabularx}{0.7\textwidth}{X} 頭一個上來，說：『主啊，你的一錠銀子已經賺了十錠。』 \end{tabularx} \\ \\ \relax
19:17 & \begin{tabularx}{0.7\textwidth}{X} 主人對他說：『好，我善良的僕人，你既在最小的事上忠心，你有權柄管十座城。』 \end{tabularx} \\ \\ \relax
19:18 & \begin{tabularx}{0.7\textwidth}{X} 第二個來，說：『主啊，你的一錠銀子已經賺了五錠。』 \end{tabularx} \\ \\ \relax
19:19 & \begin{tabularx}{0.7\textwidth}{X} 主人也對這個說：『你管五座城。』 \end{tabularx} \\ \\ \relax
19:20 & \begin{tabularx}{0.7\textwidth}{X} 又有一個來說：『主啊！看哪，你的一錠銀子在這裡，我把它包在手巾裡存著。 \end{tabularx} \\ \\ \relax
19:21 & \begin{tabularx}{0.7\textwidth}{X} 我向來怕你，因為你是嚴厲的人：沒有放的，也要去拿；沒有種的，也要去收。』 \end{tabularx} \\ \\ \relax
19:22 & \begin{tabularx}{0.7\textwidth}{X} 主人對他說：『你這惡僕，我要憑你的話定你的罪。你既知道我是嚴厲的人，沒有放的也去拿，沒有種的也去收， \end{tabularx} \\ \\ \relax
19:23 & \begin{tabularx}{0.7\textwidth}{X} 為甚麼不把我的銀子存在銀行，等我來的時候，連本帶利都取回來呢？』 \end{tabularx} \\ \\ \relax
19:24 & \begin{tabularx}{0.7\textwidth}{X} 於是他對那些站在旁邊的人說：『把他這一錠奪過來，給那有十錠的。』 \end{tabularx} \\ \\ \relax
19:25 & \begin{tabularx}{0.7\textwidth}{X} 他們對他說：『主啊，他已經有十錠了。』 \end{tabularx} \\ \\ \relax
19:26 & \begin{tabularx}{0.7\textwidth}{X} 主人說：『我告訴你們，凡有的，還要給他；沒有的，連他所有的也要奪過來。 \end{tabularx} \\ \\ \relax
19:27 & \begin{tabularx}{0.7\textwidth}{X} 至於我那些仇敵，不要我作他們王的，把他們拉來，在我面前殺了！』」 \end{tabularx} \\ \\
[1ex]
\hline
\hline
\end{longtable}
$^{1}$領主們平安,我們一起祈禱.
我還在耶利哥那裡講,因為耶穌是將近耶路撒冷,亦都是我們這個與我的王同行一路的系列最後一講,下個星期就是主任牧師會來講一個系列,一連三個禮拜,.
他的那個叫做broken hearted,亦叫傷心的時候怎樣全心全意來愛主,今年的主題是hold hearted,全心全意,但是人都會遇到一些時候不能夠全心全意愛主的,就是受傷的時候,傷心的時候,怎樣呢?主任牧師就想講這個系列,.
他來講三個禮拜就講到,我就順便放三個禮拜假,休息一下,清感一點的年假,上個星期就提到,當主耶穌開始這一次的耶路撒冷之旅的時候,其實他是知道這一次是他在地上最後一次去耶路撒冷,.
因此他就要快快加速的將天國的道理傳開,一開始在第十章,我們就看到他猜派七十二位門徒去不同的城鎮傳福音,在頭三本福音書裡面,在這裡記載派七十二位門徒,其他福音書有記載派十二位門徒,十二位使徒,.
這裡就特別提他們派了七十二位,怎麼派這麼多人?多些人就可以做多些工作,去多些城鎮傳福音,所以我們說是很緊張,在這一次的旅途裡面,我們又看到路加醫生記載了主耶穌是藉著教導,醫病和其他的神蹟,.
來顯示出他就是要來建立天國的那位微才,正如我曾經講過,天國的來臨是分兩個階段的,頭一個階段就是要建立一群天國子民的階段,這個階段的天國是在人的心裡面,.
所以耶穌說上帝的國就在你們的心裡,在路加方案十七章二十一節裡面有記載,第二個階段就是主耶穌再來審判世人,那什麼時候來呢?這個就沒有人知道,因為耶穌在什麼時候回來其實不是有時間限制,.
不是有時間限制,而是受上帝的計劃完成了未來決定,我們人就受時間限制,搭火車又要看時間表,搭飛機就一定要按時間,並且還要早些去到飛機場辦登機手續,.
但耶穌回來的時間就不是固定了某個時間,而是等到他在第一階段所定下來得救的人數達標,那他就可以回來了,那我們要問,當我們等候耶穌再回來的時候,上帝對我們有什麼要求呢?.
我相信有幾個重要的要求,第一就是要我們完全信服上帝,主耶穌在路加福音十八章十七次講得很清楚,很絕對,他說「凡要接受上帝國的若不像小孩子,絕不能進去.」.
像小孩子的意思不是叫我們幼稚無知,而是要信服上帝的主權,改變我們的思想,改變我們的價值觀,甚至我們的情緒的反應,就好像細蚊仔一樣,會聽父母和大人講的話,.
第二方面我們要行徹禱告,在路加福音十八章一到十四節,主耶穌講了兩個比喻講明禱告的重要性,第三方面就是今天的經文要告訴我們,我們要執行上帝給我們的使命,活出一個有使命感的人生,我們一起來思想今天的經文,.
十一節告訴我們耶穌已經快到耶路撒冷了,從耶利哥到耶路撒冷,大約有三十公里的路,如果步行的話大概六到七個小時就到了,.
我們看一幅地圖,我們了解距離耶利哥來耶路撒冷的確不是很遠,這個時候跟著耶穌的門徒他們有一個想法,以為上帝的國快要顯現,他們以為天國很快就要降臨,.
我們大概還記得耶穌曾經三次告訴門徒,他這次去耶路撒冷是要受害的,不過路加福音十八章三十四節告訴我們,其實那幫門徒是夢察察的,耶穌講了三次,他們仍然對這些事門徒一點也不明白,.
這話的意思是隱藏的,他們不知道所說的是什麼,耶穌講了三次,那些門徒都不知道他在說什麼,這個就是今天這段經文,耶穌講這個比喻的原因,是要幫助門徒明白天國的降臨是有兩個階段的,.
比喻一開始的時候在十二節,耶穌就說有一個貴族往遠方去,說要取得王位然後回來,那位貴族就是代表主耶穌,這裡很清楚幫助門徒和讀者明白耶穌的國度的降臨是有兩個階段,.
第一個階段就是耶穌要往遠方去,要離開門徒一段長的時間,之後才會回來在地上做黃絲行審判,這個比喻其實是有當時的歷史背景的,猶大不是一個擁有自己主權的獨立國家,.
他當時是受羅馬人的統治,如果我們還記得主耶穌出生的時候,是有一位希律王在位的話,其實那位希律王是被羅馬的凱撒大帝冊封為猶大地區的王,.
他是猶大地區的王,只是很小的地區,他所管轄的範圍只限於猶大地區,希律王是人來的,人都會死的,他死了之後會怎樣呢?.
他就想把王位傳給他的兒子阿基魯,不過猶大地區的人不喜歡阿基魯做王,他們就寫信向凱撒大帝抗議,在耶穌的比喻當中都提到有一批人是不喜歡那位貴族做王的,.
就說有些人是接受主耶穌的王權,有些人就不肯接受主耶穌的王權,那位貴族在出門之前就叫了十位僕人來,十位只是一個象徵的數字,象徵所有的僕人都叫來了,.
每人就給了一個金幣,吩咐他們要去做生意賺錢,僕人就代表耶穌基督的門徒,也可以說是所有的基督徒,銀紙就代表上帝給我們的生命,.
生命包括了我們生存的時間,我們的健康智慧學識和才幹,聖經中有一個名詞就叫做恩慈,意思就是說上帝給我們一切的才智和能力,.
用現代社會的例子來說,這些僕人可以比喻做投資經理,現代人都很熟悉基金這個名詞,我知道很多基金都是請一些專業的投資經理負責去做投資買賣賺錢的,.
基金的錢不是屬於投資經理的,那些錢是屬於投資人的,投資經理只是受委託去做投資賺錢,所賺到的錢就不屬於那位投資經理的,仍然都是屬於投資人的,基金經理就有收一份薪水,.
我們每一個基督徒都好像是上帝的投資經理,上帝給了我們恩賜和財富,但是我們只是受託人,我們千萬不要懼為己有,在耶穌還沒回來之前我們應該把握機會去運用我們一切的資源,為上帝來賺取利潤,.
甚至說他得了王位回來就吩咐叫領銀子的僕人來,要知道他們做生意賺了多少,被冊封回來之後就是代表主耶穌第二次回來,也是天國的第二個階段的實現,.
我們留意到經文講到,他回來就吩咐僕人來,吩咐僕人要來,要來做什麼呢?就是做審判,做交場,每一個僕人都要向他匯報,向他交場,16至14節就給我們看到那位王對僕人的審判,.
僕人當中有兩類,第一類就是依照主人所吩咐的去做生意,賺了錢的,頭兩個僕人都賺了錢,賺了錢的僕人主人稱讚他們都是好僕人,並且委以重任,賺了十個金幣的給他十座城去管,賺了五個金幣的就給他五座城去管,.
這個做法其實可以說是合情合理的,如果有一位好的投資經理,很有智慧,很會投資,董事局當然會升他的職位,給他多幾個基金去管理,從比喻的角度來看,其實用一個忠心僕人就夠了,用一個就夠了,不用用兩個,.
但耶穌就用了兩個忠心的僕人,耶穌是要強調主人是必定會獎賞那些忠心的僕人,無論他賺多賺少都是會獎賞,並且那個獎賞是公平的..
我們是不是想得到主耶穌的獎賞呢?如果我們渴望得到獎賞的話,我們就要善用主耶穌所給我們的機會,使用上帝給我們的生命,因此去為上帝服務..
我從前的教會有位會友是醫生,醫生都每次向病人傳福音,不做醫生的人其實都可以一樣傳福音,你在工作場所裡面可以向同事傳福音..
我女兒以前在公司工作的時候就和幾位同事在公司辦了一次小型報道會,請同事來吃飯,然後請一位特別嘉賓來分享福音..
我從前的教會也有位講國語的姐妹,叫Lydia,她也沒有工作,但是她有一個房子租了幾間房子給學生住..
我時時請他們吃飯,就向他們傳福音,有時還請我去傳福音.所以其實我們不用做牧師,我們也可以傳福音.不用做醫生也可以傳福音,你就在你的環境裡面想想有沒有什麼適合的方法來傳福音..
有些人可能會說我沒有傳福音,那怎麼辦呢?你也可以傳福音,你在教會裡面有不同的崗位一起去事奉,一起去努力傳福音.你控制電腦,幫忙煮飯,清理場地,都可以的..
我們英文堂有很多義工教英文來傳福音,我們可以和其他的弟兄姊妹一起努力去建立天國子民的工作,做這些工作..

$^{41}$無論一個人的崗位是什麼,恩賜是什麼,忠心事奉上帝的人,一樣會得到同樣的獎賞..
第二類人就是不願意使用我們生命和恩賜去事奉上帝的人,他們有什麼下場呢?有的,我們看二十四節到二十六節,這個僕人就是不用恩賜事奉主的人的代表..
他不肯事奉,也有一個原因的,他說我向來怕你,因為你試過嚴厲的人,言下之意他就是怕蝕本,而會被主人責罵..
他的說話其實還有一個可能的意思,就是說我不去做生意,主人是由於你平時太嚴厲了,所以害得我很害怕..
你不能怪我,你要怪就怪你自己,你平時那麼嚴厲幹什麼?你害得我很害怕..
其實如果他真的有這個想法,就是典型的魯亞當的表現,當亞當犯罪之後,他怎樣呢?.
給上帝心,他對上帝說,上帝啊,你給我作伴侶的,你給我女人的,給我那些我就吃了,不只是我的錯啊上帝,你做了女人給我就害了我,是你的錯啊上帝,將責任推給上帝..
在比喻中的主人其實不理會那部人的狡辯,因為他一方面知道主人是嚴厲的,另一方面他連最穩陣最保險的投資,.
怎樣呢?千萬不要買股票了,就把錢放在銀行收利息,他有沒有做呢?不做啊,連放下錢放在銀行收利息都不做啊,把錢埋在地上,收起來就算了..
嚴厲的主人當然就要懲罰他,懲罰的方法就是拿走他的金幣,交給有賺到十個金幣的人..
有些部人就覺得,有沒有搞錯啊主人,不是很公平啊,那個人已經有十個金幣了,還給他?還給他更多?他們覺得那個懲罰對吳忠德部人不太公平,.
就把那個金幣給那個已經有很多的更加不公平,主人就把他的想法原則再講出來,他說我告訴你,凡有的還要給他,沒有的連他所有的也要奪下來..
事實上這個原則是很公平的,跟我們今天的商業社會也沒什麼大分別,任何一個老闆請工人都會棄從那位忠心能幹的工人,.
懶惰又不負責的工人,最後當然是被老闆解僱了,辭退也有份的,這個其實也是主耶穌的原則..
好了,看完這個比喻了,對我們今天已經信主的弟兄姊妹來說,讓我們從這個比喻裡面學習三件事..
第一,作為天國子民,在等候主耶穌再回來的時候,我們應該祈禱等候,也都要使用上帝所給我們的生命和才幹去事奉上帝..
當我們願意運用恩賜事奉主的時候,我們會得到更多的,我們的恩賜會越用越好,不是越用越虧損,不是越用越少..
另一方面,當主再來的時候,祂會獎賞忠心的僕人,可能有些人只是想得到今生的榮耀,但我鼓勵大家要將眼光放遠一點,看長遠一點..
人生幾十年相對於永恆來說,是微不足道,是非常短暫的,眨眼就過了.我們不要被情與慾,名與利所吸引..
在世上的風光,無論你多風光,都是很短暫的.我想起我們長大的時候的那些名歌星,羅文,梅艷芳,鄧麗君,現在全部都成為古人了..
只是留一些歌給我唱,卡拉OK而已.我們這代人過去之後,可能沒有人會繼續唱了.我鼓勵大家要爭取的是永恆的獎賞,得到永恆的光..
但《義理書》十二章三節說:「智慧人要發光,如同天上的光,那令許多人歸於夷的,必發光而升,直到永永遠遠.」.
大人信主是很重要的,傳福音是很重要的,這是上帝給所有忠心的人的獎賞,永遠的光..
第二點,許多基督徒有時會記得主耶穌是慈愛的,而忘記了主耶穌是上帝,上帝也是嚴厲的..
很多人信主之後就不怕上帝,忘記了上帝是嚴厲的,只記得耶穌是很慈愛的,有罪的,會殺死人的.於是他們不怕上帝,那怎麼辦呢?.
偏行紀律,不遵循上帝的教導,聖經的教誨,聖經教導我們要有正確的行為,但今天很多年輕人對婚前性行為覺得不是犯罪,跟著世俗潮流,無所謂,想做就做..
聖經說上帝恨惡離婚,很多人,包括基督徒,運用性格不合離婚,不理會上帝的感受,自己想做就去做..
不要以為耶穌還沒有回來才不會審判我們,我們可以隨意犯罪.隨意犯罪的基督徒其實就是藐視上帝.希伯來書有一個警告,在10章30-31節說:主要審判他的百姓,落在永生上帝的手裡,真是可怕..
在希伯來書第12章,作者提醒我們,上帝一定會管教他的兒女,我們不是靠好行為成為上帝的兒女,但做了上帝的兒女之後,就要學習聽上帝的話去行事為人..
耶穌說要像小孩子一樣,要聽上帝的話去行事為人,否則就難免被上帝管教,懲罰,審判,就好像被父母管教一樣..
第三點,每一個人都要為自己的行為交仗,當主耶穌第二次再來的時候,他要執行審判,我們每一個人都要為自己的行為負責..
使命和責任是不能推卸,盡責的會得獎賞,不盡責的就不可以賴牧師母告訴我..
我今天已經講得很清楚了,每個人都要為自己的選擇和行為負責任,不能推卸,我們一起祈禱..
又讓我們的生命活出耶穌的光輝,照亮這個黑暗的世界,奉耶穌基督的名祈禱,阿們..
\newpage



\section{申命記 2:26-3:11}
\label{sec:032469}
\textbf{戰勝巨人三步曲}~~~~禤凱榮牧師
\newline
\newline
連結: \href{http://www.church.com.hk/acms/content.asp?site=cdc\&op=show\&type=product\&code=032469\&layout=sermon}{\texttt{032469}} ~~~~ 語音日期: 2026-07-19
\newline
\newline
\hyperref[sec:032468]{\small{< < < PREV SERMON < < <}}
~
\hyperref[sec:index_chronic]{\small{[返順時目]}}
~
\hyperref[sec:032470]{\small{> > > NEXT SERMON > > >}}
\newline
\newline
申命記 2:26-3:11
\newline
\begin{longtable}{cl}
\hline
\hline
章節 & 經文 (和合本修訂版)\\
\hline
2:26 & \begin{tabularx}{0.7\textwidth}{X} 「我從基底莫的曠野派遣使者到希實本王西宏那裡，用和平的話說： \end{tabularx} \\ \\ \relax
2:27 & \begin{tabularx}{0.7\textwidth}{X} 『求你讓我穿越你的地，我走路的時候，只走大路，不偏左右。 \end{tabularx} \\ \\ \relax
2:28 & \begin{tabularx}{0.7\textwidth}{X} 你可以賣糧給我吃，賣水給我喝；只要讓我步行過去， \end{tabularx} \\ \\ \relax
2:29 & \begin{tabularx}{0.7\textwidth}{X} 就如住在西珥的以掃子孫和住在亞珥的摩押人待我一樣，等我過了約旦河，進入耶和華－我們神所賜給我們的地。』 \end{tabularx} \\ \\ \relax
2:30 & \begin{tabularx}{0.7\textwidth}{X} 但希實本王西宏不肯讓我們從他那裡經過，因為耶和華－你的神使他性情頑梗，內心剛硬，為要把他交在你手中，像今日一樣。 \end{tabularx} \\ \\ \relax
2:31 & \begin{tabularx}{0.7\textwidth}{X} 耶和華對我說：『看，我已開始把西宏和他的地交給你了，你要開始得他的地為業。』 \end{tabularx} \\ \\ \relax
2:32 & \begin{tabularx}{0.7\textwidth}{X} 「西宏和他的眾百姓出來迎擊我們，在雅雜與我們交戰。 \end{tabularx} \\ \\ \relax
2:33 & \begin{tabularx}{0.7\textwidth}{X} 耶和華－我們的神把他交給我們，我們就殺了他和他的眾兒子，以及他所有的百姓。 \end{tabularx} \\ \\ \relax
2:34 & \begin{tabularx}{0.7\textwidth}{X} 那時，我們奪了他一切的城鎮，毀滅各城的男人、女人、孩子，沒有留下一個倖存者。 \end{tabularx} \\ \\ \relax
2:35 & \begin{tabularx}{0.7\textwidth}{X} 只有牲畜和所奪各城的財物，我們都取為自己的掠物。 \end{tabularx} \\ \\ \relax
2:36 & \begin{tabularx}{0.7\textwidth}{X} 從亞嫩谷旁的亞羅珥和谷中的城，直到基列，沒有一座城是高得我們不能攻取的；耶和華－我們的神把它們全都交給我們了。 \end{tabularx} \\ \\ \relax
2:37 & \begin{tabularx}{0.7\textwidth}{X} 只有亞捫人之地，雅博河沿岸，以及山區的城鎮，你沒有挨近，這全是耶和華－我們神所吩咐的。」 \end{tabularx} \\ \\ \relax
3:1 & \begin{tabularx}{0.7\textwidth}{X} 「我們又轉回，朝巴珊的路上去。巴珊王噩和他的眾百姓出來迎擊我們，在以得來與我們交戰。 \end{tabularx} \\ \\ \relax
3:2 & \begin{tabularx}{0.7\textwidth}{X} 耶和華對我說：『不要怕他！因我已把他和他的眾百姓，以及他的地，都交在你手中；你要待他像從前待住在希實本的亞摩利王西宏一樣。』 \end{tabularx} \\ \\ \relax
3:3 & \begin{tabularx}{0.7\textwidth}{X} 於是耶和華－我們的神也把巴珊王噩和他的眾百姓都交在我們手中；我們殺了他，沒有給他留下一個倖存者。 \end{tabularx} \\ \\ \relax
3:4 & \begin{tabularx}{0.7\textwidth}{X} 那時，我們奪了他一切的城鎮，共六十座，沒有一座城不被我們所奪，這是亞珥歌伯的全境，巴珊王噩的國度。 \end{tabularx} \\ \\ \relax
3:5 & \begin{tabularx}{0.7\textwidth}{X} 這些堅固的城都有高的城牆，有門有閂，此外，還有許多無城牆的鄉村。 \end{tabularx} \\ \\ \relax
3:6 & \begin{tabularx}{0.7\textwidth}{X} 我們把這些都毀滅了，像從前待希實本王西宏一樣，毀滅各城的男人、女人、孩子； \end{tabularx} \\ \\ \relax
3:7 & \begin{tabularx}{0.7\textwidth}{X} 只有一切牲畜和城中的財物，我們取為自己的掠物。 \end{tabularx} \\ \\ \relax
3:8 & \begin{tabularx}{0.7\textwidth}{X} 那時，我們從兩個亞摩利王的手裡把約旦河東邊的地奪過來，從亞嫩谷直到黑門山， \end{tabularx} \\ \\ \relax
3:9 & \begin{tabularx}{0.7\textwidth}{X} 這黑門山，西頓人稱為西連，亞摩利人稱為示尼珥。 \end{tabularx} \\ \\ \relax
3:10 & \begin{tabularx}{0.7\textwidth}{X} 我們奪了平原的各城、基列全地、巴珊全地，直到撒迦和以得來，都是巴珊王噩國內的城鎮。 \end{tabularx} \\ \\ \relax
3:11 & \begin{tabularx}{0.7\textwidth}{X} 利乏音人所剩下的只有巴珊王噩。看哪，他的床是鐵床，按照人肘的度量，長九肘，寬四肘，現今不是在亞捫人的拉巴嗎？」 \end{tabularx} \\ \\
[1ex]
\hline
\hline
\end{longtable}
$^{1}$我想問大家 如果有人問你.
在你的心目中最理想的人生是怎樣的人生呢?.
如果你可以選擇的話 你會怎樣選擇?.
最理想的人生.
我想 大概很多人都會說.
哇 四字詞 馬上就出來了.
一帆風順 無病無痛 心想事成 出入平安.
龍馬精神 生意興隆 夫妻恩愛 父慈子厚.
可能隨時加多一個年生貴子.
不過這個可能大家都不太想.
總之一句話 所有新年的賀詞.
最好全部實現.
這就是一個理想的人生.
這是不是你所追求的理想人生呢?.
如果是 當然是.
對不起 如果是的話 我猜大概你會很失望.
因為你和我住了人生某一段時間之後.
我相信我們都理解到.
剛才所說的那些.
都不會在現實上完全發生.
如果你真的這麼想的時候.
我想大概你未完全明白神的心意.
如果我們用一艘船來代表人的一生的時候.
我相信神從來不是做好一艘船.
是把船放在碼頭 海港那裡來美一下.
神做好一艘船 必定是要讓凡出海.
乘風破浪 顯出這艘船被造的價值.
這才是船應該有的 是不是?.
今天的經文告訴我們.
縱然神將這個應許之地給過以色列人.
不過神不是抬著花轎送他們進去的.
或者今天的用詞就是.
神不會開Limousine.
讓你們這些以色列人威風進去的.
聖經告訴我們.
他們要進入應許之地.
得這個迦南美地是需要征戰的.
是會經歷很多的戰爭的.
而且他們要和這些人的征戰.
不是說那些幼弱的民眾.

$^{41}$這裡所說的都是一些巨人.
就是說他們要得這個神的應許之地.
他們就要在戰爭裡面擊敗這些巨人.
今天你和我如果認真面對我們的人生的時候.
你會發覺生命上有很多巨人.
當然我所說的巨人不是一些高十尺.
又帶著長矛盾牌的巨人.
我在這裡所說.
我們生命上總是遇到一些困難.
也許是一個醫生跟你說的報告.
聽完之後嚇得你半死.
又或者是一個經濟的危機.
看到這個經濟危機的話.
就是完全沒有對前面的希望.
又或者是一個破裂的關係.
面對這個破裂的關係.
你知道沒有辦法修補.
又或者是一個潛藏在你生命裡面的一個恐懼.
沒有辦法超越.
我想問你.
面對這些生命上沒有辦法逃避的挑戰.
或者我們叫做巨人的時候.
怎麼搞.
我們怎麼面對.
讓我們通過剛才我們所讀的經文.
一同來看看.
聖經其實在這個故事告訴我們.
三個我們需要知道,明白,實踐的真理.
以致我們可以在人生的路上.
無論遇到什麼巨人.
我們都能夠克服,能夠戰勝.
好,我們看第一個.
第一個.
今天的經文開始的時候.
就講到一個故事.
摩西派人去到這個希實本王叫做西弘.
那幫人就用一些很和睦.
就是好聲好氣.
講什麼.
他說.

$^{81}$王啊王啊.
求你容我們通過你的地.
我們只是走大道而已.
不會亂走的.
我們經過的時候.
你都可以向我們賣糧賣水.
這樣挺好的.
就是說.
他們不單止沒有損失.
這一幫人也可以賺錢.
給他們賺錢的機會.
以色列又可以通過.
是一個雙贏局面.
你贏我又贏.
不過很有趣.
聖經告訴我們.
這個希實本王西弘.
卻是拒絕了.
我想大家想一想.
不知道你有沒有一個印象.
就是說.
我們信了耶穌之後.
哇,真是棒.
每天天色常藍.
為什麼?.
神又有能力又慈愛又愛我們.
神必定會賜福給我們人生的每一個細節.
我曾經聽過這樣的傳福音.
你信耶穌吧.
你信了耶穌之後.
你的問題神會幫你解決.
你有病嗎?.
神會幫你醫好的.
你關係破裂嗎?.
神會幫你修改的.
我甚至聽到這些人越講到眉飛色舞的時候.
誰誰都是這樣的.
信了耶穌之後馬上找到工作.
信了耶穌之後.
我的病馬上康復.

$^{121}$那個毒瘤就馬上.
不知道為什麼醫生也解釋不了.
消失了.
我相信神絕對有這樣的能力.
行這些神蹟.
不過我想大家知道.
如果我們用這些東西來作為我們信耶穌的賣點的時候.
我不覺得這個是傳福音.
我覺得這個像是推銷.
還是用一些很美化的包裝來作為促銷.
快點賣掉的感覺.
如果我們看聖經.
看回早期教會的歷史.
或者看我們周圍一些我們尊敬的基督徒的時候.
你會發覺一樣事實.
這個事實是什麼.
沒有一個基督徒.
他的人生.
是每一天都天色常藍.
每一天都風和日麗.
沒有這樣東西.
而當時的以色列人.
也都是一樣.
雖然神已經賜了一個應許之地給他們.
是過去式賜了的.
但是不等於他們沒有任何的障礙.
而這個西弘.
這個王.
就是以色列人進入迦南地的第一個障礙.
留意聖經怎麼說.
這個西弘王不讓他們經過他的地圖.
聖經說是什麼原因呢.
聖經說.
因為耶和說你的神使他心中剛硬.
神使他心中剛硬.
神使到他.
是不是意思就是說.
神使到這個西弘王拒絕以色列人經過呢.
如果這樣的時候.
那就是.

$^{161}$怎麼解釋呢.
這段經文我們必須要解釋.
必須要領悟的一樣東西.
這段經文不是告訴我們.
神是強行將一些觀念.
或者決定.
硬生生執入西弘的腦海裡面.
不是這個意思.
神不是待人好像木偶那樣.
我喜歡戳圓就圓.
喜歡按扁就剪.
神不是這樣待人的.
神絕對不會將一個善良的人.
變成一個邪惡的人.
我們應該怎麼理解這段經文呢.
這段經文就是告訴我們.
神是用已經在西弘心裡面的頑梗.
驕傲.
使用他來彰顯.
在這個過程裡面.
神要彰顯的能力.
和他要成就的心意.
很多時候在現實上你會發覺.
無論我們怎樣的努力.
追求和平.
追求和睦.
我們做事就是有規有矩.
不要有衝突.
不過我們仍然很多時候.
都會遇到一些頑固的反對.
就在我們面前明明是機會的門.
一下子就被人關了.
接著衝突升級.
在這個時候我們都會質疑自己.
我們究竟是否走在神的心意上面.
為什麼會有這麼多的挫敗呢.
我就想大家留意這段經文.
神的回答是怎麼樣的.
神的回答是.
從此起手.

$^{201}$我要將西弘和他的地.
交給以色列人.
你要得西弘王的地為業.
很明顯.
這個西弘的硬心.
他的倔強.
並不是告訴我們.
告訴以色列人.
神已經離棄了他的子民.
你們自生自滅.
自己搞定.
並不是這個意思.
相反來說.
這個王的硬心.
反而成為了上帝最終勝利的鋪墊.
預備了最終的勝利.
神容讓西弘的心.
硬起來.
以至神可以將他交在以色列人的手中.
所以弟兄姊妹.
我相信我們人生路上總是有這些風浪.
總是有這些巨人.
不過今天的經文告訴我們.
第一樣東西.
還有一個很重要的東西.
我們需要明白.
在人生路上必有困難.
在人生路上甚至有苦難.
所以不要當有苦難來的時候.
我們就怨天尤人.
神為什麼會有這麼多東西.
你是不是有能力.
你是不是愛我的.
為什麼會有這麼多苦難.
第一樣我們需要知道.
人必定會遇到風浪.
不要讓困難挑戰巨人.
就想著神是離棄我們.
反而是一個機會.
讓我們為神爭戰.

$^{241}$是一個機會為神得著勝利.
讓我們更加的熟齡,長進,成熟.
靠近我們的主,我們的神.
以至神可以得著榮耀.
所以有人說得好.
不要在神面前過分祈求一個.
順風順水,每天天色常藍的人生.
反而我們要求神賜給我們一雙有力的腳.
可以行走在神為我們預備的人生路上.
神從來沒有承諾過我們一個無風無浪的人生.
神並沒有.
好了,我們明白,亦理解.
我們需要用信心來回應神的應許和承諾.
但大家同一個時間亦需要信心.
並沒有移除了行動.
信心並沒有移除了行動.
所以大家可以想像.
當這個西弘王帶著他的軍隊.
來與以色列人爭戰.
以色列人基本上可以說是沒有選擇.
只能夠與出來的西弘王的民眾爭戰.
而正如神之前亦很清楚地說.
我已經將西弘王的地給了你們.
所以很明顯他們打勝仗擊敗了這一班人.
我想大家留意一個很平鋪直敘的描述.
這裡說:耶和華,我們的神將他交給我們.
我們就把他和他的兒女.
並他的眾民都擊殺了.
這是一段很簡單的經文.
但其實我們認真看的時候.
我們發覺其實是一個次序在當中.
次序的第一是什麼?.
次序的第一就是神有所行動.
第二,人,百姓,信服.
接著,人,回應,用行動來回應.
所以神沒有說過你們儘管進去.
你們盡你們的力打,拿到多少就多少.
神並不是這樣說.
神說:去得你,我已經賜給你的地.
以色列亦是很認真的遵守神的命令和吩咐.

$^{281}$所以你看回整個故事說的時候.
這一班人沒有走偏了路來攻打其他村莊.
沒有,他們單是按著神的指示來攻打.
這個勝利可以說是很徹底的.
勢宏王被消滅.
他的軍隊被消滅.
他的地就落在以色列人的手中.
我想大家留意一下.
摩西怎樣來形容當時的戰績.
我們看這個24到36節.
他形容戰績是這樣的.
我們奪了他的一切成就和其中的財物.
耶和說:我們的神都教給我們了.
沒有一座城高得使我們不能共取的.
我想大家留意最後一句.
沒有一座城高得使我們不能共取的.
大家覺不覺得這個說話很面熟.
幾時聽過.
記不記得他們準備進去的時候.
他們派了12個探子走進去探路.
拿情報.
他們回來的情報是甚麼.
不行不行.
這些人民又高大.
他們的城牆又高又堅固.
我們這次一定死定了.
就是因為這個回報.
所以他們怎樣.
全個以色列民陷在恐慌當中.
他們的聖經仍說他們的心都消化了.
最後怎樣.
最後不肯上去.
每個人都非常害怕.
我想問同一個城和同一個城牆.
為何有這麼大分別.
一開始說打不到.
現在就說沒有一座城不能被共取的.
為何有這麼大的分別.
我想大家想一想.
一個小朋友晚上走路.

$^{321}$非常害怕自己一個人.
突然一走的時候.
聽到他們的聲音.
馬上回頭看一看.
嘩 牆上有一個這樣的影子.
這是甚麼來的 老虎.
那個小朋友嚇得不得了.
不過再看清楚一點.
不是老虎來的.
是一隻貓來的.
頂曉梅 我告訴你.
恐懼是不必要地.
亦都可以無限地放大某些事情.
甚至歪曲了某些事情.
而信心是將我們的恐懼放在一旁.
讓我們的焦點仍然集中在我們所相信的天賦上帝那裡.
而這個神仙行動.
人民跟隨的這個模式.
亦都重複在他們下一個戰役.
和這個巴山王岳的戰役那裡.
同一個模式亦都有發生.
不過這一次有些不同.
巴山王聖經說是一個更加恐怖的人物.
如何恐怖呢.
聖經說他有六十座城.
這些城他們的高牆又高又堅固.
有門有山.
所以在十一節更加告訴我們.
這個巴山王岳是一個利弗欽人.
利弗欽人是甚麼人.
是一個超級的巨人.
很高大的.
我想大家想像一下.
一個高十呎的巨人.
帶著一大班他的士兵出城.
去對陣以色列人.
我們可以想像出來的時候呼聲震天.
接著鼓聲雷動.
很大的威勢.
大致上這班人出來的時候.

$^{361}$都是頭戴這個盔甲.
拿著長矛拿著這個盾牌.
做足了這個有了準備.
很大的威勢.
但你猜以色列人是甚麼樣子的.
以色列人是不過在之前的幾十年.
在曠野地兜來兜去的一班人民.
我想大概他們那個戰爭的裝備.
遠不及這一班人.
所以我們可以想像.
在一個這樣的懸殊的情況下.
以色列人感到恐懼.
必然的事可以理解的.
但是神這樣跟摩西說.
神說不要怕他.
因我已將他和他的眾民.
並他的地都交在你手中.
不用怕.
留意這是一個過去式的.
我已經將他們的地給了你.
就是說在戰爭未開始之前.
勝負已經神定了.
這班以色列人進去.
就是拿神已經給過他們的地的意思.
也許今天.
我們所面對的巨人.
不是一個十呎高.
拿著盾牌.
拿著長矛的敵人.
但是大家要留意.
舊約裡面其實差不多每一個的戰爭.
每一個以色列人所打的戰爭.
其實都是指向耶穌基督的.
當我們這樣理解的時候.
我們就明白一個很重要的真理.
其實好像以色列人一樣.
我們今天所面對的敵人.
都是我們戰勝不了的敵人.
不過當然我們面對的敵人.
不是西弘.

$^{401}$不是我.
我們最大的敵人是什麼.
就是我們的罪.
死亡.
和永恆的奏狀.
不過聖經講得很清楚.
在十字架上面.
主耶穌為人的緣故.
面對一個比我.
比西弘更加強大的敵人.
但是他仍然的戰勝.
他的戰勝讓他可以為人承擔所有的罪惡.
他為人忍受了神的懲罰.
他為人進入死亡.
他為人的緣故.
第三天來都是從死裡得到榮耀的復活.
這個復活.
其實就是神向整個世界宣稱.
我已經勝過你們的敵人.
所以不用怕他.
就正如保羅在哥林多前書的第十五章這樣講.
藉著我們的主耶穌基督得勝.
所以弟兄姊妹.
一個已經發生了的事實.
就讓耶穌在十字架上成就的救恩.
讓我們有一個確信.
有一個能力.
來過一個得勝的生活.
並不因我們的能力.
而是因為耶穌基督.
在十字架上所成就的救恩.
當日擊敗西弘和岳的神.
沒有改變.
當日讓耶穌基督從死裡復活的神.
都沒有改變.
他就是今天與我們同行人生.
當我們面對挑戰困難的時候.
可以依靠的那個神.
所以弟兄姊妹讓我們.
不要按我們對自己的信心.

$^{441}$自己的能力 自己的意志.
來面對我們每一個人生風浪.
而是依靠走在我們前面.
為我們爭戰的神.
連這個神從來沒有違背祂的承諾.
所以正如今天的經文不斷提醒我們.
不要怕祂.
無論我們生命裡什麼巨人都不需要怕祂.
因耶和華我們的神.
已將祂和祂的眾民.
並祂的地都交在我們手中.
所以大家要明白.
謹記一個很重要的事實.
這個世界上沒有一個巨人.
是大得過神能夠擊倒的.
也沒有一個山高得過神不能攀爬的.
也沒有一個難關.
是困難得過神沒辦法去克服的.
所以當我們理解這個真理的時候.
當我們面對人生每一個困境.
縱然是一個最大最難.
沒可能跨越.
沒可能克服的困難來到的時候.
不要被它打垮.
反而讓那個危機.
那個困難成為一個我們成長.
要榮耀祂.
彰顯祂能力的機會.
所以願意這個應許.
鼓勵我們每一個人在人生的每一個的爭戰上.
我們都依靠祂.
直到祂帶我們進入應許之地.
你聽到這裡的時候.
大家會不會覺得.
哇 牧師你真的講得很精彩.
不過講是容易的.
做是很難的.
你不要說爭戰勝利.
你在困難.
在這個危難.

$^{481}$在風暴裡面.
你過一天.
你都很難的.
牧師你不要講得這麼誇張.
這個是不是你的想法.
首先我想大家去明白一件事.
在這裡我想跟大家分享兩個真理.
這兩個真理很多時候我們會忽略.
甚至是誤解的.
我講第一個的時候.
我想首先問大家一個問題.
什麼為之在屬靈的爭戰裡面得著勝利.
什麼才算是勝利.
我用疾病來做一個例子.
當醫生跟你講.
哎呀 不好意思.
你患了末期的不治之症.
你會怎麼做.
啊 聽到這個消息的時候.
我個個都馬上.
馬上去找靈醫醫治.
馬上來祈禱.
叫上頂尖幫我們一起祈禱.
求神醫治.
這些做法當然是對的.
不過我想大家想一想.
如果神沒有醫治呢.
如果你的治療是無效呢.
如果神好像沒有聽你的禱告.
沒有讓你的病得醫治.
甚至你的病況越來越差呢.
你會怎麼樣.
你會不會因此而感到極度的不滿.
極度的埋怨 投訴.
神啊 為什麼會這樣.
神啊 你是不是比我英文.
你會走在我前面為我爭戰的嗎.
神啊 你是不是跟我說過.
當我們倚靠你的時候.
我們就能夠贏每一個的屬靈爭戰的嗎.

$^{521}$這是不是大家的想法.
你怎樣理解屬靈爭戰的勝利.
如果你理解為.
如果你是這樣理解.
病得醫治才叫做爭戰成功的.
如果你是這樣理解的時候.
對不起 你搞錯了.
病得醫治是一個勝利.
這是屬世的看法.
這是世界的觀念.
但是對於基督徒來說 在一個屬靈的世界裡面.
絕對不是這樣.
在一個屬靈的世界裡面 告訴我們.
爭戰的勝利不一定等於病得醫治.
大家想想 如果神要醫治每一個病.
當我們祈禱的時候 神就醫治的時候.
基督徒基本上就不用死了.
每一個病都得醫治.
基督徒可以長命百歲 不需要死.
既然基督徒仍然需要面對肉身的死亡.
病未必一定得醫治的時候.
我們怎樣理解 甚麼是屬靈爭戰的勝利呢?.
贏得屬靈爭戰.
不是要得到你想要的東西.
而是要忠於耶穌基督.
使祂的旨意能夠在你和我的生命當中得以成就.
讓人 讓神得著榮耀.
這就是屬靈爭戰的勝利.
我想大家明白.
當一個人知道他離世的日子不遠.
他能夠用一個正面積極的態度.
面對他餘下的日子.
心裡面滿有平安喜樂.
盡量使用他餘下的日子.
見證神 事奉神.
我告訴你 榮耀神強過病得醫治.
當人面對死亡而無懼於色.
仍然滿有感恩 平安與喜樂.
這些是裝不出來的.
但當你倚靠神的時候.

$^{561}$神就容讓你過一個得勝的生活.
包括面對死亡而不懼怕.
我們再看第二.
大家有沒有留意.
在最後第十一節的時候.
摩西好像無端端說了一些離題的話.
在第十一節摩西說到.
這個巴山王國的床是鐵造的.
長九爪 闊四爪.
好像跟故事沒有什麼關係.
所以大家千萬不要誤會.
這個補充的資料.
不是讓我們預備問答遊戲的時候.
多一個問答文.
不是這個意思.
這個補充的資料.
是讓整個以色列民甚至比我們更加的一個記錄.
為什麼要記錄呢.
一張很大的鐵床.
用今天的長度是十三尺乘六尺的大小.
是一個很具體.
你摸得到 看得到的一個提醒.
告訴我們昔日以色列人.
怎樣靠神的恩典來打敗這個巨人.
有多大.
就是這張床那麼大.
看得到 摸得到.
甚至他們擺在某一個地方.
好像紀念館那樣.
你不夠信心你去那裡自己看一看.
一個記號.
這個記號我告訴大家.
很重要.
因為神能夠戰勝一個人體為巨大到不得了的戰爭.
大家可以想像一下.
N年之後.
當這個以色列人仍然感到很大的挑戰.
有每一個挑戰來到的時候.
我們怎麼搞.
搞不定.

$^{601}$不要緊.
你記不記得這個巴山王岳.
記不記得神怎樣來戰勝他.
他有多大.
就是這張床那麼大.
如果神能夠戰敗這個巨人岳的時候.
他就能夠幫我們去戰勝下一個的人生挑戰.
而事實上.
這個當日擊敗這個西弘和岳.
在以色列人的歷史裡面.
是一個很重大的勝利.
如果大家熟悉詩篇的話.
你看這個一三五和一三六篇.
兩篇的詩篇.
其實都是紀念這個神替以色列人戰勝的一個感恩.
所以頂主妹.
你和我都需要這個鐵床.
在我們的生命裡面提醒我們.
神在我們的生命的歷程裡面.
曾經為你打勝過仗.
既然昔日神已經為我們打勝仗的時候.
我們就不要那麼懼怕明天的征戰.
大家記不記得這個大衛王的故事.
大衛王是當時剛剛出來的時候.
是個小朋友.
是個少年人.
是為了打歌.
當時的蘇勒王說:你行不行啊小朋友.
當時的大衛王就說.
他沒有跟蘇勒王說.
我打算怎樣怎樣打.
這是我的策略.
我這樣打就會贏了.
你會發覺大衛王沒有這樣說.
大衛王是怎麼說的.
大衛王只是很簡單地在這個三十七節裡面說.
夜禍說救我脫離獅子和紅的爪.
也必救我脫離這非利士人的手.
大家看不看得到.
贏取屬靈征戰的秘訣.

$^{641}$有什麼秘訣.
除了一個事實.
就是信心是神所賜予之外.
第二樣東西我們都需要理解.
我們也可以通過我們跟神的經驗.
累積建立我們對神的信心.
神每一次在我們人生中讓我們得到的勝利.
這個勝利就是將來我們戰勝下一個挑戰的信心.
裡面被建立起來.
神很小心一次過將所有人生的困難全部丟下來.
怕我們不死.
神從來不會這樣做.
神是按部就班一步一步培育我們.
讓我們經年之後信心就得到被建立.
當我們面對每一次新的巨人的時候.
這個故事告訴我們.
我們都可以回頭看一看.
神在我們人生軌跡上的信實.
神昔日帶領每一個祂的子民通過曠野.
這個神也是為我們戰勝我們過去的巨人.
也是今天與我們同在危難的時候與我們同在的神.
所以在我們過往的人生裡面.
大家有沒有經歷過獅子與熊的爪.
神有沒有曾經將你從獅子與熊的爪裡面拿出來呢?.
神有沒有為你戰勝你生命中的西王或是巴山王岳呢?.
我想大家知道.
如果我們平常的人生從來都沒有依靠過神.
來過得勝的生活的時候.
我們怎麼能夠期望在一個重大的人生風暴裡面.
我們有足夠的信心面對呢?.
大家去過學校.
其實就好像上學一樣.
當你上學的時候.
久不久就會有功課要做.
久不久也會有一個小小的測試.
讓你去跟著課程.
如果每一個的軍閘.
每一個的小測試.
你都是懶懶散散.
隨便吧,沒所謂的.

$^{681}$你猜你到大考的時候可不可以過關呢?.
弟兄姊妹.
我們依靠神是需要我們每天累積我們的人生經驗.
來到今天經文的訊息的結尾.
我想重新再提醒大家這三個很重要的功課.
第一.
神從來沒有應許過天色常藍.
神從來都沒有說人生是無風無浪.
所以當我們遇到困難的時候.
首先我們不要怨天尤人.
不要因為風浪的來臨.
我們就否定了這個神.
這是第一步我們需要確認的.
第二件事.
當風暴來臨的時候.
我們必定要抓緊神.
深信神是與我們同在.
祂應許走在我們前面為我們爭戰.
但是不是按我們的意思.
是按神的意思.
所以我們不需要懼怕.
而是依靠祂來過一個祂已經贏了的戰爭.
第三.
無論大事還是小事.
我們都需要學習怎樣來依靠祂過一個得勝的生活.
每一個過去的勝利就是建立我們將來面對將來挑戰的信心.
當我們有了這幾樣的三部曲.
當我們理解了的時候.
我告訴大家.
我們就能夠靠著神的恩典.
來戰勝我們生命的巨人.
請我停低頭祈禱.
耶穌基督我們承認.
我們很多時候都用恐懼來遮蓋我們的信心.
很多時候我們所追求的.
未必是我們所認真的想法.
但很多時候我們都無意識地在追求世界的成功和勝利.
而忘記了真實屬靈爭戰的勝利.
是在風暴之中.
或者風暴過後.

$^{721}$我們仍然站在神的道路上.
而你在整個過程中得著榮耀.
我們在整個過程中得著成長.
所以神下求你教導我們.
幫助我們每一天依靠你過著一個得勝的生活.
以致將來現在我們能夠面對人生每一個風浪.
讓你得著當得的榮耀.
多謝你讓我們一起祈禱奉主命強.
阿們.
\newpage



\section{羅馬書 1:20-23}
\label{sec:032470}
\textbf{誰是我「王」}~~~~梁錫強醫生
\newline
\newline
連結: \href{http://www.church.com.hk/acms/content.asp?site=cdc\&op=show\&type=product\&code=032470\&layout=sermon}{\texttt{032470}} ~~~~ 語音日期: 2026-07-19
\newline
\newline
\hyperref[sec:032469]{\small{< < < PREV SERMON < < <}}
~
\hyperref[sec:index_chronic]{\small{[返順時目]}}
~
\hyperref[sec:032471]{\small{> > > NEXT SERMON > > >}}
\newline
\newline
羅馬書 1:20-23
\newline
\begin{longtable}{cl}
\hline
\hline
章節 & 經文 (和合本修訂版)\\
\hline
1:20 & \begin{tabularx}{0.7\textwidth}{X} 自從造天地以來，神的永能和神性是明明可知的，雖然眼不能見，但藉著所造之物就可以了解看見，叫人無可推諉。 \end{tabularx} \\ \\ \relax
1:21 & \begin{tabularx}{0.7\textwidth}{X} 因為，他們雖然知道神，卻不把他當作神榮耀他，也不感謝他。他們的思想變為虛妄，無知的心昏暗了。 \end{tabularx} \\ \\ \relax
1:22 & \begin{tabularx}{0.7\textwidth}{X} 他們自以為聰明，反成了愚昧， \end{tabularx} \\ \\ \relax
1:23 & \begin{tabularx}{0.7\textwidth}{X} 將不能朽壞之神的榮耀變為偶像，仿照必朽壞的人、飛禽、走獸、爬蟲的形像。 \end{tabularx} \\ \\
[1ex]
\hline
\hline
\end{longtable}
$^{1}$弟兄姊妹 主內平安.
很高興再一次在台上跟大家有機會分享上帝的話語.
大家如果有留意英國的政局.
明天早上就會有一件很大的事情.
就是我們會更換我們的國家領導人.
卸任的首相一早就會去白金漢宮.
為什麼去白金漢宮呢?.
就是要見我們的國王查理三世.
要求辭職.
不是說辭就辭 要皇帝批准他才可以辭職.
理由就是他不可以再有國會的支持.
新上任的大家都知道是誰.
這個Andy 我都不知道他的中文怎麼翻譯.
有不同的翻譯方法.
接著他就會.
一個出來 跟著第二個進去.
好像貨源輪轉那樣 他就會進去.
他就要求國王安立他做新任的首相.
很奇怪 究竟我們的領導人是誰呢?.
首相還是皇帝呢?.
其實英國是一個很特別的地方.
其實跟日本差不多.
我們叫做君主立憲.
就是保留有這個君王的制度.
但是就透過立憲就是將人民的權力擺在第一.
就限制了皇帝的權力.
所以皇帝又不是廢的.
有人說他廢君 其實不是 是很重要的.
譬如他可以不批准的 不批准你下台.
也有很多的法律最後要他簽名 他可以不簽.
不過如果不簽又不批准就會有很多政制危機.
幾百年來都沒有試過.
希望明天不會.
所以其實皇帝也是最高的權位在英國.
雖然這個Andy說他自己也是皇帝.
不過有人封他 他說King of the North.
有沒有聽過?.
因為他之前在曼城做了市長.
做得有聲有色的 人們說他.
所以人們叫他北方之王 King of the North.

$^{41}$現在空降來了.
統治我們全英國了 看看會怎樣.
有些人拭目以待 有些人很期望.
有些人就很怕.
他說要去回復英國70年代.
上世紀的輝煌.
上世紀70年代我來英國讀書.
我就覺得一點都不輝煌.
英國那時候很窮.
所以我們拭目以待.
皇帝這件事給我們的印象中是最厲害的.
最強的.
最能說服人的.
尤其是大家看中國歷史的皇帝.
他可以定你生死的.
他說你要死就死 說你生就生.
你不能反駁 不能上訴.
而且就算皇帝上了任.
他的名字你都不可以跟著他.
所有人都要改名.
因為這樣你是會一記遺的.
不可以亂.
你的名字不可以和皇帝的名字一樣.
所以權力很大.
也讓我們知道最強的就是皇帝.
所以踢球打球有球王.
打拳有拳王.
下棋有棋王.
在自然界有萬獸之王.
通常都是說獅子是萬獸之王.
當然在樂壇裡面也有四大天王.
當然這四大天王都比較過氣了.
所以給我們的印象中.
皇帝是很厲害的.
而我們其實心底裡面.
我們每一個人都想做王.
有誰不想考試考第一呢?.
有誰不想在公司越爬越高.
做到最高的CEO呢?.
我們想做到最高 做到最好.

$^{81}$給人景仰的.
我不是說這個是錯.
但這個其實是人性最深的渴望.
我們就算做不到也好.
我想我兒子出行投地.
我做不到王 我兒子做也可以.
我們很多時候就是有這樣的心態.
就算基督徒也好.
很多時候我們有這樣的心態.
在《撒慕爾記》第八章.
回看《撒慕爾記》上.
很多人就會想想.
究竟《撒慕爾記》上是說什麼.
很多人都看過了.
其實主角不是撒慕爾.
撒慕爾不是他.
主角也不是大衛.
大家熟悉的大衛.
當然真正的主角當然是上帝.
但是人的主角來說.
其實是說掃羅多.
掃羅是花了很多篇幅.
是說掃羅.
所以大家回去看回《撒慕爾記》上.
看回掃羅的生平.
看完《撒慕爾記》上就可以了.
我們現在開始《撒慕爾記》第八章.
第八章到第八章的時候.
我看到第四節.
長老聚集在拉瑪地方.
來到撒慕爾那裡說話.
有什麼說呢?.
他說:漢娜,你年紀老了.
你的兒子又不行你的道.
現在請你為我立一個王.
治你我門將列國一樣.
為什麼會這樣說呢?.
因為回看第八章開頭的三節.
他形容這兩個兒子.
撒慕爾其實有兩個兒子.

$^{121}$就是很不正當.
《聖經》用了幾個字去形容他們.
說他貪圖財利.
收取賄賂.
屈枉正直.
其實就是做什麼都不好的事.
濫用他們的權力.
他們也有一些權力.
也是濫用了他父親撒慕爾.
所以就做一些人家個個都可憎的事.
所以如果你要他們這兩個兒子.
擔任以色列人的領袖.
他們就很不想.
所以他說:既然你年紀老了.
你兩個兒子又不爭氣了.
不如你為我們立一個王治你我門.
因為他知道.
真的要有一個王.
就不是他們推舉出來的.
是要撒慕爾按照上帝的意思.
才可以立.
重點是最後那幾個字.
將列國一樣.
撒慕爾說人家個個國家周圍.
都有皇帝的.
為什麼我們沒有.
今天我們的心態就是.
有些事人家個個都是這樣做的.
在新的時代.
我們基督徒不做不跟風.
我們就out.
如果我們out.
我們就怎樣去傳福音.
怎樣去埋他們堆呢.
很多時候我們將這件事.
做了我們的藉口.
不是說人家其他國家有王.
我們就要有王.
不要忘記以色列.
他們一個國家是很不同的.

$^{161}$是神權統治的.
他們真正的王是上帝.
這麼多年來都是的.
他們的王不是約書.
也不是摩西.
也不是宗善師.
但是現在他們想要找個王.
你可以說找個王都可以的.
他也under上帝的.
但是人的心目中.
王就是能說話的.
最大權力的.
其實他們的意思.
就是他們不想再聽上帝的說話.
他們就自作主張.
撒母耳聽了這句話.
就很不開心.
第六節.
他說立一個王.
他就跟耶和華禱告說.
但是耶和華怎樣回答他呢.
你只管聽從百姓向你說的一切話.
因為他們不是厭棄你.
而是厭棄我.
不要我作他們的王.
自從我領他們從埃及出來到現在.
他們都離棄我.
事奉別的神.
正像他們從前所做的一切事.
現在他們也照樣做了.
現在你只管告訴他們.
不過你要嚴厲警告.
告訴他們將來王會用什麼方式管轄你們.
這個就是很重要.
這番話看到我們上帝的性情.
或者是屬性.
上帝是不開心的.
不喜歡的.
但是上帝是知道為什麼.
說你立一個王.

$^{201}$不是想著你的制度會好一點.
你以後就會統治得更加好的國家.
其實他們的心底.
上帝已經知道了.
他們想離開上帝.
不再聽上帝的說話.
上帝是知道的.
但是上帝就任憑他們去做.
但是有一個警告.
就是警告了他們.
你今天這樣做.
你就看到將來的王會怎樣對待你們.
結果後來歷史是應驗了上帝的說話.
當然是應驗了.
上帝是自有永有的.
今天上帝和我們的關係都是一樣.
很多時候我們任意去做.
隨從自己的私慾去做事的時候.
上帝都是准許我們去做的.
他不阻止 不攔阻.
但是他有警告.
他藉著聖經的說話.
弟兄姊妹的提醒.
牧者的教導.
他會警告我們.
提醒我們這樣不行.
你這樣會承擔後果.
上帝是這樣的.
他給我們某一個自由.
所以不要說上帝和耶穌沒有自由.
不是 上帝是有自由.
但是上帝很關心我們.
他不會說任由你去死.
他不會.
他會給一個警告.
當然 以色列人是聽不出這個警告的.
今天我們其實為什麼要立一個皇帝呢?.
自己做不了皇帝.
也想立一個皇帝出來.
其實很簡單的.

$^{241}$就是我們想自己的事自己去管.
上帝是自有永有的.
上帝你是無所不在的.
上帝你看不到 摸不到的.
我不知道你是誰.
和你的關係好像很遠似的.
我找回一個人.
我信服的.
我可以去管我們的.
那就更好.
我們很多時候想自己管自己的人生.
我們真的是.
我們在我們日常生活當中.
很多時候就要在我的掌控之內.
你不想失控的.
最好每樣東西你是控制得到的.
但是上帝很飄忽的.
有時候不知道他在想什麼.
我怎麼掌控呢?.
所以我們就是這樣.
就算我們信了主很久.
到時候我們很多時候用人的意思.
我們不是再聽上帝的說話了.
我回來教會.
我們就不如做一個平順圖.
就是我們坐在後面.
我們不要搞這麼多東西了.
因為變成我們不是聽上帝的說話.
我自己心中就當我自己是王.
所謂平順圖.
其實我們平到真的躺平那種.
上帝不是要我們這樣.
當我們與世人為伍的時候.
我們就很擔心別人怎麼看我們.
我們不想做到自己太過凶 是不是?.
忘記了上帝要我們.
就是正正要我們分辨為性.
但是我們就擔心.
被人小看 被人歧視.
特別是在這種世俗的社會當中.

$^{281}$他們會覺得我們怪.
星期天天氣這麼好.
我們不去郊遊 回來教會.
很多時候我們就決定了.
我們的人生.
我們再不信服上帝了.
我們就順從自己.
順從自己就很危險.
順從自己一定是順從了.
自己的邪情私慾去走的.
一定是的.
不會不是的.
因為罪的性在我們當中.
所謂的肉體的情慾.
眼目的情慾.
今生的驕傲就會出來.
我們自己決定自己做的事.
那你問不問上帝?.
都會問.
你說我會祈禱.
但是你不是祈禱求問上帝.
你是祈禱決定了之後.
你知會上帝.
求上帝開路.
開你那條路.
所以結果是什麼呢?.
結果是人作王的時候.
神就退位了.
好像明天首相更換那樣.
上帝你退位了.
讓我上吧.
或者讓我喜歡的人上去吧.
自己管自己吧.
以色列人就是這樣的心態.
大家剛才看那一段聖經.
就看到羅馬書一章那裡.
有幾句話就是提醒了我們.
第一就是.
他們知道是有神.
但是不當作神的榮耀.

$^{321}$今天我們也是.
我們知道有神的.
我們知道上帝的.
大家信了主的大部分人.
我們都洗了禮.
歸入教會的.
但是我們知道有神.
但我們不是做每一樣事以神做榮耀.
他們的思念變為虛妄.
無知的心昏暗.
自作聰明.
反成了愚蠢.
保羅經常說這些話.
以為自己很聰明.
其實是很愚笨.
最慘就是將不能扭壞之神的榮耀變為偶像.
彷彿扭壞的人和飛禽走獸昆蟲的樣式.
什麼意思呢?.
很多時候今天.
我們仍然是記得上帝的.
星期天我們仍然回來敬拜的.
但是這個多半是一個偶像式的敬拜.
當上帝是一個偶像.
我們來敬拜你.
你保我平安一個星期.
有事找你就會聽我說.
我們是這樣.
想保個平安.
跟你去求神問卜.
去拜觀音菩薩分別不大.
但我們的神不是這樣.
敬拜的原因不是這樣.
我們不要當上帝是一個偶像去拜.
當我們當自己是王的時候.
其他的就沒那麼重要.
當上帝就淪落成為一個偶像.
你回來教我們敬拜.
你很高傲很多東西.
但是你的心在哪裡呢?.
你是否將榮耀歸給上帝呢?.

$^{361}$是否將主權仍然歸給上帝呢?.
還是你要唱得好聽的歌給人聽呢?.
是榮耀自己呢?.
這是其中一個.
我們要留意的.
我們是很愚昧的.
因為我們自作聰明.
回到《撒母耳記》那裡.
我們看到蘇羅.
蘇羅是一個很典型的例子.
蘇羅是一個什麼人呢?.
上次也說過.
蘇羅其實是一個外表很出眾的人.
很有才能.
其實不用計算.
看他的樣子就知道了.
聖經怎麼形容他呢?.
他比眾民高過一個頭.
你離遠已經看到蘇羅了.
天生出來就是領袖的樣子.
很方便的.
你離遠看到他了.
那就做領袖.
我記得小時候參加一些旅行團.
以前很流行的.
美加旅行團.
那個領隊是怎樣的呢?.
那個領隊很矮的.
那怎樣呢?.
他靠一支旗.
有一支很高的旗.
看到那支旗就知道領隊不會走失.
如果認他就慘了.
他那麼矮.
眾人之中都找不到他.
但蘇羅不是.
蘇羅是高大帥哥.
有氣勢的.
走出來大家都說.
哇!.

$^{401}$那種的.
所以他在先天來的優點是很多的.
上帝是揀選他的.
但是.
聖經一路看下去.
大家知道.
其實他沒有讓上帝真正作他的王.
上帝是揀選他的.
蘇羅不是自己霸佔來做的.
但是他的問題是什麼呢?.
他的問題就是從來沒有讓上帝進入他心裡面.
做他真正的王.
聖經提醒我們.
《迦太書》一章十節那裡說.
我現在要得人的心還是得上帝的心呢?.
難道要討人喜歡嗎?.
我若仍舊想討人喜歡.
我就不是基督的僕人.
這句話提醒我們什麼呢?.
就是蘇羅.
他怕人多過怕上帝.
要討人的開心.
而不是討上帝的開心.
今天我們都是一樣.
很多時候我們很緊張別人怎麼看我們.
但是我們有沒有緊張過上帝怎麼看我們呢?.
沒有的.
很多時候都沒有的.
我們比較緊張別人怎麼看我們.
我們做的事情要別人開心.
但是有沒有想過要討上帝開心呢?.
為什麼我們要討人開心呢?.
因為我們要榮耀自己.
我們要把自己弄得這麼厲害.
別人開心.
別人說你哪個好.
我們就當自己是王.
不是上帝的王.
所以加泰書保羅就很清楚說.
如果你仍舊想討人開心的話.

$^{441}$你就不是基督的僕人.
在《撒末記》上13章.
上次講到也說過.
就是蘇羅有一件事.
蘇羅是做錯的.
當時非理士人.
就是以色列人的死對頭.
就打到來了.
押警了.
去到他們的周邊.
那些百姓就害怕得很慘.
害怕得很慘.
他們打仗之前.
他們通常都會做一個祈禱.
今天我們崇拜之前做一個祈禱.
他們就是做一個獻祭的.
撒母耳講得很清楚.
等我來.
做了獻祭之後出去打仗.
將一切交託給上帝.
但是那些人在吵在怕.
蘇羅為了要討他們開心.
蘇羅就看到撒母耳遲了來.
他就等我來做.
這是違反了上帝的教訓.
最重要不是他獻不獻祭.
最重要是他沒有聽上帝的說話.
他沒有等.
上帝有他的時間.
他沒有等.
所以撒母耳來到.
第一句話就是.
蘇羅.
你做了糊塗的事.
你這次失敗了.
很嚴重的錯犯.
為什麼?.
就因為蘇羅.
他沒有把上帝當心中的王.
他當自己是.

$^{481}$他要討人開心.
讓人敬仰他.
結果他做了一些上帝不喜悅的事.
所以加泰遜很重要的.
講得很清楚.
如果你想討人喜歡的話.
你不是基督的僕人.
另外一道我們看到就是.
《約翰福音》.
《約翰福音》之前是講什麼呢?.
之前我們也不太重要.
就是講一堆異賽亞書.
接著講到當時有一些官長.
有一些文士.
有一些法利賽人.
這裡是講耶穌當時的教導完之後.
第42節就說.
雖然如此.
他是官長.
他是猶太人的官長.
做官的.
有些名目的.
有很多是信他的.
信誰呢? 耶穌.
但是他不敢說出來.
為什麼呢?.
因為法利賽人的緣故不敢承認.
恐怕被趕出會堂.
如果你公然說他們信耶穌是彌賽亞.
或者聽耶穌的話.
法利賽人就會很生氣.
法利賽人就是那些律法師.
就是那些很熟讀舊約的知識份子.
會趕他們走.
那麼耶穌就說.
因為約翰寫出來就是.
因他們愛人給的尊榮.
過於上帝給的尊榮.
他們看重人怎麼看他.
看重人給什麼尊榮他.

$^{521}$不是上帝給什麼尊榮他.
就是他聽人說多於聽上帝.
他介意人怎麼說.
多於介意上帝怎麼看他們.
為什麼呢?.
因為他想有個尊榮在當中.
他當自己是王.
他要別人景仰.
他怕自己壞了.
不敢承認基督.
不敢承認耶穌.
他們是信的.
他們明白的.
但是他們不敢說.
今天我們是怎樣呢?.
我們介意別人怎麼看我們嗎?.
還是我們介意上帝怎麼看我們?.
我們信上帝之後.
我們願不願意出去傳福音呢?.
就是說我們很怕.
怕沒有面子.
怕別人覺得我們怪.
在AI時代你還相信這些東西.
蘇羅就是這樣.
他怕.
他怕別人怎麼看他.
他保住自己的王位.
他當自己是真正的王.
就忘記了上帝.
第二點就是.
看到很多時候.
我們是用宗教.
我們的信仰來包裝我們自己很漂亮.
蘇羅就是這樣.
有一件事是發生於撒母耳記上十五章.
當時他們去打亞馬利人.
蘇羅帶軍隊打亞馬利人.
上帝已經提醒他們.
就好像以前在打耶利哥.
在《約書下記》裡說.

$^{561}$就是你要滅絕.
滅盡亞馬利人.
滅盡的意思就是.
他們所有神族人.
還有牲畜.
牛,羊.
都要滅絕.
這是上帝的命令.
當撒母耳來到蘇羅的時候.
打完仗之後.
聽到有牛在叫.
羊在叫.
他說你沒有殺光他們嗎?.
留下一點嗎?.
他真的留下一點.
如果你在蘇羅.
為什麼被人揭穿了?.
說對不起.
下次不敢了.
有點貪心的.
看到羊很肥.
牛又肥.
今晚吃火鍋也有用.
你可能會這樣想.
會這樣說.
但不是.
蘇羅還在狡辯.
他怎麼說呢?.
他還把上帝放在桌子上.
他很厲害.
他說這些是亞馬利人得來的.
其實我們殺光很多了.
剩下一些好的東西.
是留下來獻給耶和華的.
他用這個信仰.
這個宗教包裝得自己很漂亮.
放了上帝出來.
沒有人敢反駁他.
是不是?.
是真的.

$^{601}$上帝會不會要外族人的生畜獻祭給他呢?.
是不會的.
其實他是貪心的.
也可能是因為百姓要求.
不要浪費.
他本來聽了百姓的說話.
沒有聽上帝的說話.
所以我們真的要看到這裡的時候.
我們想一下.
今天很多時候我們都是一樣.
我們包裝得自己很好.
每一次都打響招牌.
我們基督徒.
把信仰,基督教,道德的光環放在頭上.
沒有人敢說你.
特別是你說.
經過祈禱反省.
聖靈的引導.
我這樣這樣.
沒有人敢說你.
我不是說你不是.
但當你放這句出來.
好像蘇羅那樣說.
是我拿來獻祭給上帝的.
其實沒有人敢出聲的.
很多時候我們會不會犯這個錯誤呢?.
宗教的語言可以說得很漂亮.
但其實心是為了自己的.
你打響道德的光環.
去批評弟兄姊妹的時候.
你目的是幫弟兄姊妹.
還是想抬高自己的道德地位呢?.
所以保羅說得很清楚.
你所說的話是求別人的益處.
就是求自己的益處.
我們需要想一想.
這件事你是榮耀了上帝.
就是榮耀了你自己.
保羅很直的跟哥林多前書說.
早四個星期也說了.

$^{641}$拜偶像的食物就是這樣說.
你是求自己的榮耀.
就是求上帝的榮耀.
你是求自己的益處.
就是求別人的益處.
很多時候我們就是.
甚至我們把上帝放在桌上.
甚至用蜀靈的膠紙來包著自己.
包得很漂亮給別人看.
讓別人不敢罵你.
蘇羅就是這樣.
他恃著自己高大衛猛.
有型有款.
恃著自己是撒母耳的案立.
他做上帝的子兒.
他做皇帝的.
但是他心不在上帝.
他心在自己.
他要保住自己的江山.
他要別人去讚賞他.
當犯錯的時候撞狡辯.
把蜀靈包裝得自己是多麼漂亮.
撒母耳的一句就拆穿了他.
這句是在十五章那裡說.
他說.
大家很熟悉的.
其實那句就很長.
不過其實記住這幾個字就行了.
信服勝於獻祭.
耶和華是愛繁祭的祭物.
但是更加愛什麼?.
聽從祂說話的人.
什麼繁祭什麼最好的肥羊肥牛.
沒有用的.
你不聽上帝的說話.
你奉獻擺上什麼都沒有用.
今天我們來到教會.
我們做多少的事奉.
我們參加多少的活動.
從早到晚都在這裡.

$^{681}$不代表你是聽從上帝的說話.
這個是屬靈外表的包裝.
耶和華是喜歡這些東西.
不是不喜歡.
不過更加喜歡.
如果你聽祂說話.
信服勝於獻祭.
撒慕爾很清楚告訴保羅.
很多時候我們借了上帝的光環.
道德的光環.
擺在自己的頭上.
當作自己的光環.
在《馬太福音》.
這句也是很直接的.
耶穌說的.
「不是每一個稱呼我主呀主呀的人.
都能進天國.
唯有進行我天父旨意的人才能進去」.
不是你的口經常說.
「My God, My God, My God, My God, My God, My God, Help me God」.
什麼什麼的.
不是你滿口都是那些屬靈的漂亮說話.
就是你一個可以進入天國的人.
《聖經》說唯有進行上帝旨意的人才可以.
我們不知道.
我不知道你是不是.
你也不知道我是不是.
但上帝是知道你是否進行祂旨意的人.
不是口裡經常把《聖經》幾句金句.
牧者幾句說話放一個隨口說.
就可以說你是屬靈的人.
你出發點是對的.
可能是 但未必一定是.
《聖經》說得很清楚.
耶穌說的那些說話.
你口裡掛著主呀主呀不代表什麼.
行動緊要.
口說是沒什麼用.
口說我們每個人都可以說.
說是可以說得很漂亮.

$^{721}$說是可以說得很好.
訴羅做的事是很漂亮.
被人看到.
他們要他獻祭他就獻祭.
要他留一些生畜.
他又留一些生畜.
他很聽話.
聽人民話.
但他不需要聽人民話.
其實他不是民選出來的.
我們套用今天的詞語就是.
他是上帝安立他的.
他一定要聽上帝的話.
到最後就很可怕.
訴羅的路是越走越窄.
什麼叫越走越窄.
大家有沒有走過窄路.
我不是說屬靈的窄路.
真的是一條窄的路.
那些路真的很窄.
很難的.
以前香港新界很多路是沒什麼行人路的.
跌跌撞撞就算有也只能一個人走.
對面有人來你就要避開他.
打側讓他過.
不然就走出馬路.
這是窄路.
窄路是很難走.
窄路也是越走越迷路.
其實是不好的.
我們喜歡康莊大道.
窄路是難走的.
但訴羅的生命就是越走越窄.
為什麼呢.
初時他都是犯幾個錯.
每個人都會.
但他就想了一遍.
他沒有來到上帝面前悔改.
他沒有求上帝饒恕他.
他繼續去做自己想做的事.

$^{761}$隨從他自己個人的私慾繼續去做.
結果那條路是越走越窄.
越走越迷路.
越走越窄.
大家如果看《撒母耳記》上.
從18章到28章.
分了10章的篇幅去形容.
訴羅後期的生活.
他是怎樣的.
他是遭遇.
當你是偏行幾路.
你隨從自己的心態去做事的時候.
第一你內部的罪性就出來了.
嫉妒.
嫉妒是一個很大的罪性.
他妒忌大衛.
人民開始擁戴大衛.
他不喜歡.
打完仗回來.
明明他是主帥.
他帶頭的.
大衛跟著他.
人們就在喊什麼.
那些婦女.
訴羅就說這些無知婦女.
有沒有搞錯.
你們在喊什麼.
你說訴羅殺死千千.
大衛殺死萬萬.
不當我是王.
當大衛是王.
皇帝當然是萬歲.
不會千歲那麼少.
他妒忌大衛.
他追殺大衛.
到最後.
他甚至是去求問鬼神.
即是他徹底離開了上帝.
他去求問鬼神.
本來他是一個上帝揀選的人.

$^{801}$但後來他變成這樣的樣子.
我相信連他自己都認得自己的樣子.
變成這樣.
樣子可能沒有變.
蒼老了一點.
打了那麼多仗.
但內心是完全變了.
他求問鬼神.
今天我們又怎樣呢.
很多時候.
路越走越窄的時候.
我們所謂想一遍.
走一遍想一遍.
最怕就是想一遍.
人們有情緒病.
我們最怕就是想一遍.
路越走越窄.
開始封閉自己.
不再接觸外界.
不再交朋友.
不再回教會.
不再見人.
整整下一年就不去看醫生.
又不去吃藥.
路會越走越窄.
接著會做很多反常的事.
人們都認不出你是基督徒.
素羅去求問鬼神.
要幫他.
借助那些鬼神的力量去幫他.
忘記了上帝的能力.
這就是素羅的收場.
變成這樣.
路越走越窄是很可怕.
想一遍是很可怕.
最後只是走去滅亡的路.
唯有回到上帝那裡悔改.
但素羅沒有這樣做.
路越窄的時候.
就越來越黑暗.

$^{841}$越來越黑暗的時候.
就越來越孤單.
人生的路就是這樣.
你封閉了自己.
你靠自己.
你就不再靠上帝.
開始連教會都不想回來了.
初時就在家裡聽道.
上網崇拜.
整整又沒空.
又很多事要做.
很多隱憂.
很多事影響到你.
慢慢就離開了.
基督徒朋友想找你.
初時也會見一下.
後來每次找我都叫我回教會.
不好意思.
不再當沒空.
你就離開了這個群體.
去交那些損友.
交那些世俗的朋友.
開心很多.
一起去飲飲食食.
我喜歡說什麼就說什麼.
沒有人說我了.
沒有人叫我回教會.
我們開始走那條路.
我見過的.
我有朋友他從小到大回教會.
後來他就離開了上帝.
初時就是因為忙.
世俗的隱憂太大.
他很多工作.
星期天都要開工.
有很多應酬.
要和這個去吃東西.
和那個去巴結.
搞好一些網絡.
搞好一些關係.

$^{881}$整整一下就根本沒時間給上帝了.
自己是王.
那些朋友是最重要的.
那些朋友捧到他很高.
結果到最後.
他是越走越孤單.
離開了.
連他的家庭也散了.
這個就是很多人害怕的一件事.
另外認識一個姊妹.
她也是一個基督徒的姊妹.
她很多年前她患病.
重病.
初時對上帝都有信心.
但是醫來醫去都醫不好.
她沒有交托.
她反而埋怨.
基督徒朋友來找她.
就是叫她祈禱.
她只有這樣.
還有其他事情嗎?.
開始對他們沒有信心.
開始遠離基督徒朋友.
然後自己走到一邊.
路越窄.
到最後什麼鬼神.
什麼的事情.
她都走去試.
走去問.
有殺過沒有放過.
她說她希望可以醫好.
結果路就走到黑暗滅亡那裡.
在《雅各書》一章十至十五節那裡.
她說每一個人被誘惑.
是因自己的私慾牽引而被誘惑.
私慾既懷了胎就生出死.
這三樣東西是一環扣一環的.
就是私慾,罪和死.
我們犯罪的原因就是有私慾這件事.
之前也說過了.

$^{921}$那些今生的驕傲.
肉體的情慾.
眼目的情慾.
這個在我們的DNA裡面.
很多時候我們就是一個不經意就走出來的.
每個人都會有的.
如果我們不是靠上帝.
不是靠耶穌基督的寶血.
不是上帝的教導提醒.
我們就一定會跟著這個世界去犯罪.
因為這個世界就是魔鬼管轄.
這個世界屬世的世界是屬肉體.
屬肉體的是有罪的.
是惡者管轄.
而到最後《羅馬書》說得很清楚.
罪的功架就是死.
這個就是我們一個定律.
是走不出來的.
除了有基督耶穌的寶血.
照著自己的私慾去做.
當自己是王.
我們介意別人怎麼看.
不是上帝怎麼看.
我們將榮耀歸給自己.
不是歸給上帝.
我們求自己益處.
不求上帝益處的時候.
這條路就是這樣走的.
就好像訴羅那樣.
到最後的結局是什麼呢?.
《撒慕爾記》上最後一章的時候.
在基利波那裡.
訴羅倒下在戰場上死了.
這個不是罪可恃.
一個元帥帶領人去打仗.
就是自生死在道外的.
你隨時知道沒命的.
最大的悲劇就是.
屬靈的悲劇.
就是他離開了上帝.

$^{961}$他沒有再依靠上帝.
他堅持自己做王帝的時候.
生命就失控.
以為自己可以控制自己的人生.
結果是失控.
搞不定.
最徬徨的時候.
他又要逃跑.
又要求神問鬼.
連基本的方向都沒有了.
當我們離開上帝.
我們連人生的方向都沒有了.
我們這條路只會失控.
最後的結局就是死亡.
《羅馬書八章》.
五至八節那裡說.
隨從肉體的人.
體貼肉體的事.
隨從聖靈和體貼聖靈的事.
肉體和聖靈.
我們經常說.
是一個mutually exclusive.
在數學當中.
兩樣東西是不會重疊的.
你不可以隨從一點點肉體.
隨從一點點上帝,聖靈.
沒有.
要麼一樣要麼那樣.
沒有中間的.
Sorry 沒有的.
體貼肉體的就是死.
體貼聖靈的就是生命和平安.
因為體貼肉體就是與神為敵.
因為體貼肉體就是體貼這個世界.
體貼那惡者魔鬼的掌權.
魔鬼是與神為敵.
對上帝的律法不信從.
事實上也無法信服.
屬肉體的人無法使上帝喜悅.
所以保羅是沒有辦法使上帝喜悅.

$^{1001}$因為他一生之中就是封自己做王.
不是以上帝做王.
結果他就徹底的失敗.
最後就是賠償生命.
這個不只是肉體的死.
是屬靈上的死.
它離開上帝.
今天我們以為自己可以.
控制到自己的人生.
結果就是我們走了很多冤枉路.
我們給了很多壓力我們自己.
你封自己做王.
不容易做的.
做皇帝其實.
你做一個公司的高層.
CEO.
其實很多事情要搞的.
老套說「是」.
「是」你就能說「是」.
「是」你就很威風.
很多人都怕你.
「是」「是」「真」.
但其實帶來的是有很多很多的壓力.
壓力令到我們是透不過氣.
令到我們走歪路.
走歪一邊.
想正面.
走偏門.
我們很多時候就是這樣開始.
罪就是這樣開始來.
大家有沒有看球賽?.
有沒有?.
有嗎?.
所以睜開眼睛睡就是這個意思.
昨晚那場球賽不是很晚.
10點開場.
我很乖.
因為我知道要早堂講道.
我沒有看的.
就是英格蘭這次打的唯一一場我沒有看的.

$^{1041}$結果是最好看的.
一場球賽本來是沒有人想看.
沒有人想打的.
連球員都不想打的.
結果就是打到不想停.
看到不想停.
結果就是英格蘭贏.
法國6比4.
姑勿論你覺得這個結果是怎樣.
你開不開心都好.
為什麼突然間打得這麼好.
兩隊都好.
沒有了壓力.
毫無壓力了.
只差季軍.
贏了世界盃捧不到盃.
打得好.
是不是?.
是很簡單.
人生就是這樣的.
就是這樣.
其實不錯的.
雖然今早我們的主席.
他不在現在.
就板著臉.
因為英國輸了.
但是我提醒他.
原來英格蘭是1966年以來.
這次打得最好.
季軍.
以前沒有試過的.
打軍都沒有試過.
而且不在自己主場.
這次是最好的.
因為上次66年贏了.
我們在自己主場.
所以如果這樣想之下.
又不錯.
今天我們如果當自己是王的時候.
我們給自己很多壓力.

$^{1081}$人生就好像一場球賽一樣越打越差.
因為太大壓力了.
越打越差.
上帝說.
他提醒我們.
為什麼你不將你的勞苦重擔.
帶到主耶穌基督面前呢?.
其實很簡單的.
凡勞苦重擔帶到主耶穌基督面前.
你就可以安息.
你可以將東西射下去.
讓耶穌基督幫你揹著.
射一下這個包袱.
這條路就更加輕散.
長波就越打越漂亮.
應該是的.
做人都是一樣.
你執著這個王位做什麼.
沒什麼好處的.
沒什麼好處.
我們的首相打算多做兩年.
知道嗎?.
沒想到這麼快.
起碼多做兩年.
任期還沒到期.
其實是辛苦多兩年.
被人罵也是兩年.
一定是.
是不是這樣是最好呢?.
今天我們執著世上的東西.
執著很多個人的榮辱.
我們的路只會越走越難.
儘管我們用屬靈的包裝.
包裝自己很漂亮也好.
儘量我們已經很尊重人意.
很討人開心.
但是我們沒有討上帝的開心.
所以我們要提醒我們自己.
當我們來到上帝面前的時候.
我們謙卑下來.

$^{1121}$真正的王是上帝.
因為祂是創始成中.
祂是自有永有.
祂是無所不在.
無所不知.
無所不能的.
為什麼你不去依靠祂.
你生命的路越走越窄.
越走越難走.
你很灰的時候.
是不是要回到上帝那裡.
這是唯一的辦法.
真是唯一的辦法.
藥物不會幫到你.
朋友不會幫到你.
唯一的辦法一定是回到上帝那裡.
你將上帝封為王.
去掌管你一生.
不是說以後就一帆風順.
生命不會有起跌.
一樣有.
你仍然會遇到很多的困難.
仍然有生老病死.
仍然有很多的事情要面對.
但是有上帝和你一起去面對.
這個已經差很遠.
大家低頭我們有個反思一下.
今天我們是不是好像掃羅那樣呢?.
我們是不是怕人多過怕上帝?.
我們是不是介意別人怎麼說.
多過上帝怎麼看我們?.
我們是不是很多時候用信仰.
用屬靈.
用道德的光環來包裝自己.
以為這樣就是給人一個印象.
我們是最好的.
還是我們已經內心的邪情私慾.
我們的妒忌.
我們的仇恨已經走了出來.
我們不知道.

$^{1161}$我們的生命的路會不會越走越窄.
上帝不是要奪我們的權.
而是上帝要救我們脫離恐懼.
脫離壓力.
世俗的比較仇恨和嫉妒.
在《斐納比書》2:9-11說.
所以上帝把他升為至高.
又賜給他超乎萬名之上的名.
使一切在天上的,地上的和地底下的.
因耶穌的名眾室都要跪下.
眾口都要承認.
耶穌基督是主.
願榮耀給付上帝.
耶穌基督是我們真正的王.
萬室都要跪拜他.
上帝不是用他的權能來統治我們.
而是藉著耶穌基督在十字架的救恩那裡拯救我們.
所以今天我們可以成為他的兒女.
可以帶著勞苦沖擔來到他面前.
今天上帝問我們願不願意.
將我們的生命交給他.
以上帝為王.
以後我們的壓力變為一個倚靠.
我們的恐懼變為一個信心.
嫉妒變為一個安息.
我們的生命不會越走越窄.
而是會越走越闊.
天父上帝今天我們安靜地來到你的面前.
我們承認很多時候我們想自己作王.
我們想控制我們自己的人生.
我們很怕別人怎樣看我們.
我們討好其他人.
求主你憐憫.
除去我們心裡面的恐懼.
比較,固執,嫉妒,仇恨.
讓我們可以重新學會怎樣去信服.
怎樣謙卑來到你的面前.
去敬拜,尊你為主.
求主繼續帶領我們.
成為我們生命的主.

$^{1201}$生命的王.
帶領我們走在你的路上.
讓我們有勇氣,有膽量.
每一天都有智慧,平安和喜樂在當中.
願你的國降臨.
願你的旨意成就在我們今天每一節的事情當中.
我們奉主耶穌基督的名祈求.
阿們.
\newpage



\section{士師記 18:31}
\label{sec:032471}
\textbf{聖經人聖經事:沒有王,各人任意而行 (4) 六百但人屠城,奪地,私設神堂}~~~~黃克勤牧師博士
\newline
\newline
連結: \href{http://www.church.com.hk/acms/content.asp?site=cdc\&op=show\&type=product\&code=032471\&layout=sermon}{\texttt{032471}} ~~~~ 語音日期: 2026-07-15
\newline
\newline
\hyperref[sec:032470]{\small{< < < PREV SERMON < < <}}
~
\hyperref[sec:index_chronic]{\small{[返順時目]}}
~
\hyperref[sec:032472]{\small{> > > NEXT SERMON > > >}}
\newline
\newline
士師記 18:31
\newline
\begin{longtable}{cl}
\hline
\hline
章節 & 經文 (和合本修訂版)\\
\hline
18:31 & \begin{tabularx}{0.7\textwidth}{X} 神的家在示羅多少日子，但人為自己設立米迦所雕刻的像也在但多少日子。 \end{tabularx} \\ \\
[1ex]
\hline
\hline
\end{longtable}
$^{1}$大家好,一起讀聖經.
生命樹為大家預備聖經人,聖經事.
今天是講士師記來到一個很差的故事.
士師記最後那五章,大家知道我講過了.
是講了很多當時國內非常醜惡的事情.
主要是因為沒有王國人任意而行.
比較多是講怎樣任意而行的.
來到第四集,任意而行的第四集是六百個彈支派的人.
去屠城,然後搶別人的地,名字也改了.
然後自己私設一個神堂在那裡.
自己私設神堂這件事就是說.
我自己也有神的敬拜地方.
我不用去使劍,不用去會幕那裡.
這樣也可以,自己製造祭司.
又有人肯做,這樣就可以了.
是不是?好了,真是氣死.
好,下去看看.
那時色烈中會有王國人任意而行.
六百個彈的人在那裡屠城奪地私設神堂.
二十七節,士師記.
彈的人是將米家所造的神像和他的祭司.
都帶到拉翼這個地方.
大家記得嗎?他們由南邊一直上到最北.
幾乎去到黎巴嫩.
本來是在耶路撒冷以南的.
在猶大的支派的地上.
就是以法連的山地那裡.
一直上到最北.
幾乎去到黎巴嫩那麼遠.
將米家所造的神像.
和祭司一個猶大支派的利尾人帶進去.
帶到那個地方就是他們去五個探子.
去窺探過叫拉翼這個地方.
見到居民是安居無慮的民.
於是他們這六百個彈的人怎樣做呢?.
屠城.
用刀,第一在他屠城的地方,很差的地方.
第一就是用刀殺了那民.
又放火燒了那城.
所以就屠城.

$^{41}$用刀殺了那些人.
然後放火燒了那個城.
第二十八支.
並無人答救,拉翼沒有人去救他.
因為離開西頓很遠.
他又與別人沒有來往,在一個孤城.
自己一個世外桃源小城市生活.
城又在平原,平原靠近百里峽.
但人又在那裡收城居住.
就在那裡無人答救.
照著他們始祖以色列之子.
蛋的名字.
就吸那城.
起名叫蛋.
原先那個城的名字叫拉翼.
大家知道啦,只不過是作者薩姆爾將它再講一次.
這個蛋的祖先是誰呢?.
他們的始祖,始祖亞當夏娃.
不是,他們的始祖是以色列之子蛋的名字.
以色列就是亞國.
亞國有十個兒子.
十個兒子裡面有一個叫蛋.
他們是蛋的子孫.
所以他們說始祖以色列就是指亞國的兒子蛋.
就用我們這一族最早的人.
就是蛋,亞國的兒子的名字.
來起這個城的名.
所以這個蛋是城的名.
舊名叫拉翼.
他就將它改了名叫蛋.
就在那裡建了個城,剷了,然後就在那裡住了.
這個蛋跟大家解釋一下.
講到這裡才跟大家解釋.
蛋呢,Darlin.
蛋這個希伯文叫Darlin.
只不過通常英文就讀了.
darn,darn,darn.
其實是Darlin.
黎是在後面的,Dar,D-A.
Darlin不過就讀了一個音,Darn,Darlin.

$^{81}$Darlin的意思是審判.
蛋的意思是審判.
自己的名字都叫審判,又不怕審判.
原先那個城的名叫拉翼.
拉翼呢,Liege.
Liege是獅子的意思.
可能以往那裡有獅子.
所以就叫做獅子,Liege.
本來人家的名叫拉翼.
他改了做蛋,裡面的人全部殺死了.
殺死了之後,這個蛋的人.
第三個就是施設神堂.
這件事就犯了很大的罪.
神應該很不高興.
是為自己設立那個雕刻的像.
為自己設立那個雕刻的像是什麼意思呢?.
就是他們去到米家家裡.
把米家家裡那個銀的耶和華.
象徵耶和華的像是什麼樣子呢?不知道.
把他搶了.
然後其他神像也拿走了.
有雕刻,有鑄造的,全部拿走了.
然後連那個祭司那個猶大族的利未人也拘捕了.
猶大支派利美族的人也拘捕了.
脅持他,帶著子女和老婆一起走.
財物那樣去.
由下南走到北,去到那裡.
而且當時還威脅脅持他.
喂,你究竟想做一家米家家的祭司.
還是做我們整個蛋族的祭司呢?.
你不如來我們那裡做吧.
識趣一點,識時務者為俊傑.
不如來我們那裡做吧.
那個猶大支派利美族的祭司.
沒有名字,我到現在都沒有記錄他的名字.
他覺得好啊,其實他不是祭司.
只不過米家叫他做祭司.
他就做起祭司上來了.
其實只不過是一個利美族的人而已.
他就做起祭司了.

$^{121}$做起祭司,我們整個蛋族都找你做祭司.
跟著去.
但是其實去到他們拿了拉益.
然後改了整個名字.
開始又自設神堂的時候.
就自己設立了雕刻的像.
摩西的子孫,摩西有幾個兒子.
其中一個兒子叫甲順.
甲順的兒子若拿丹.
和他的子孫.
就做蛋支派的祭司.
於是去到自己又有地方了.
想著想著那個猶大支派裡面.
利美支派那個少年人.
年輕人做祭司.
都是不行的.
不如我們找些正統一點的.
勁一點的.
找摩西的兒子甲順.
摩西的孫甲順的兒子若拿丹.
和他的子孫做我們的祭司.
即是找了摩西的室若拿丹.
做他們的祭司.
和子孫都是繼續做下去.
即是說那個猶大支派利美年輕人.
就被廢了.
所以你見到跟這些惡人都沒什麼好處.
以為我們以你為父.
以你為祭司.
跟你,你和我們一齊去做.
怎知銀杖搶了.
祭司沒得做.
結果是怎樣呢?沒講到.
直到那地方遭擄掠的日子.
後來這個地方都被聖國擄掠了.
就沒得做了.
再看下去.
神的殿在士羅紀內.
丹人為自己設立米家所雕刻的像.
也在丹多少日子.

$^{161}$所以丹是一個很頑皮的.
以色列十月支派裡面其中一個.
其實當時以色列的首都.
是叫士羅,首府叫士羅.
而神的殿其實是繪墨在那裡.
繪墨在那裡.
神所在的地方.
所以稱為神的殿.
繪墨在那裡.
藥櫃在那裡.
這個時候.
藥櫃在神的殿,在繪墨裡面.
地點就是士羅.
神的殿在那裡有多久.
直到所羅門建好第一聖殿之前.
其實繪墨都不是在耶路撒冷.
都在士羅為主.
神的藥櫃,神的殿,神的繪墨只有一個.
但是丹人為自己施設一個.
我們設立了一個神堂.
擺了一個雕刻的像.
是誰呢?.
就是他們在米家家搶來的那個.
在那裡敬拜耶和華.
他們不是拜第二個神,是敬拜耶和華.
這樣都可以?.
那就是說找了另一個像出來代替耶和華.
這件事是很糟糕的.
也在丹多久?.
神的殿和藥櫃在士羅多久?.
他施設的神堂都在丹有多久?.
這件事是很糟糕的,是神很不歡喜的.
今天跟大家看到.
六百個丹人去到那裡屠城奪地施設神堂.
其實是神很不歡喜的一件事.
沒有王館.
人最怕沒有王館.
求主幫助我們,讓我們更加謙卑.
當這班丹的百姓,那六百人.
或者整個丹的支派,他們覺得自己強.

$^{201}$覺得自己恐怖有力,好打,有計謀.
又惡,他們就開始到處去搶奪.
先搶奪了米家家的寶貴財物.
搶了祭司.
然後去到北邊搶別人的城市.
殺了所有人,然後在裡面改名.
再在那裡去發展.
再建造一個他們自己的城市.
然後還在那裡建一個神堂.
於是他們不用再去事件敬拜了.
我們自己有私家祭司,私家神堂.
自己變得越來越大.
我相信神是不歡喜,是會有工作的.
你自己在丹,到後來這個支派都是消失了.
在北邊的支派多數都是.
一路下去都是行耶和華.
眼中行為惡的事.
最後是被擄掠,被消滅.
求主自己憐憫,幫助我們.
不會自己去用自己私心.
在教會裡面,在整個角度裡面.
為自己設立些什麼,想要些什麼.
想要些權力,想要些什麼.
到最尾,其實都是遭遇到上帝的難阻與毀滅.
求主幫助我們敬畏你.
誰聽祈禱,奉耶穌基督的名字.
阿們.
\newpage



\section{以賽亞書 44:6-8}
\label{sec:032472}
\textbf{真實與虛無}~~~~林家棋宣教師
\newline
\newline
連結: \href{http://www.church.com.hk/acms/content.asp?site=cdc\&op=show\&type=product\&code=032472\&layout=sermon}{\texttt{032472}} ~~~~ 語音日期: 2026-07-19
\newline
\newline
\hyperref[sec:032471]{\small{< < < PREV SERMON < < <}}
~
\hyperref[sec:index_chronic]{\small{[返順時目]}}
~
\hyperref[sec:032473]{\small{> > > NEXT SERMON > > >}}
\newline
\newline
以賽亞書 44:6-8
\newline
\begin{longtable}{cl}
\hline
\hline
章節 & 經文 (和合本修訂版)\\
\hline
44:6 & \begin{tabularx}{0.7\textwidth}{X} 耶和華－以色列的君王，以色列的救贖主－萬軍之耶和華如此說：「我是首先的，也是末後的；除我以外再沒有神。 \end{tabularx} \\ \\ \relax
44:7 & \begin{tabularx}{0.7\textwidth}{X} 自從古時我設立了人，誰能像我宣告，指明，又為自己陳說呢？讓他指明未來的事和必成的事吧！ \end{tabularx} \\ \\ \relax
44:8 & \begin{tabularx}{0.7\textwidth}{X} 你們不要恐懼，也不要害怕。我豈不是從上古就告訴並指示你們了嗎？你們是我的見證人！除我以外，豈有神呢？誠然沒有磐石，就我所知，一個也沒有！」 \end{tabularx} \\ \\
[1ex]
\hline
\hline
\end{longtable}
$^{1}$如果我們要去形容我們三位一體獨一的真誠的時候.
大家會怎樣形容我們對於這個三位一體的神的屬性.
我們會怎樣形容他呢?.
給大家一些提示.
如果大家剛才祈禱的時候有留意到少兒姐妹.
她一開始就提了兩個上帝的屬性.
是哪兩個呢?.
慈愛和公義.
除了慈愛和公義之外.
我們還會怎樣去形容我們的上主呢?.
聖潔,全知,全能,無所不在,自有永有,永恆不變等等.
在今天的經文裡面提到.
我是首先的也是末後的.
除我以外再沒有神.
我是首先的我是末後的.
其實這個說話連同今天的經文.
在二賽亞書裡面總共大概出現了三次.
分別就是41章的第4節.
以及48章的第12節.
這句說話其實是上主對於自己的自我表述.
而上主這一番的自我表述和我們今天的猜拳主人有什麼關係呢?.
我先邀請大家去打開今天的經文.
二賽亞書44章第6至第8節.
二賽亞書44章第6至第8節的經文.
大家先保留著這一頁的聖經.
因為我們今天除了看這一段經文外.
也會看其他的經文.
打開了這個二賽亞書44章第6至第8節的經文.
今天的經文其實都已經在之前講二賽亞書.
或者我們說到第二二賽亞的主題裡面.
我們經常都會提及背景是什麼呢?.
就是以色列文已經是被擄巴比倫已經有一段日子.
當以色列文經歷了上主審判之後.
亦即是二賽亞書的第一章至第39章的經文內容.
但是到了四十至五十五章的經文內容由原本審判的主題轉為救贖的主題.
特別我們看一看大家保留著44章第6至第8節的經文.
我們再揭一前少少四十章一到二節.
大家揭一揭四十章一到二節的經文.
二賽亞書四十章一到二節的經文.
如果揭開的我邀請大家一起去讀.

$^{41}$預備一二三.
若果二賽亞書我們俗稱第二二賽亞.
即是四十章至到五十五章的經文裡面.
是關於安慰以及救贖這個主題.
但是我們揭回去四十四章六到八節今天的經文.
好像又沒有什麼直接關係和這個主題.
在四十四章六到八節的經文更似.
剛才我提及的就是上主的自述多過救贖的主題.
弟兄姊妹我們要細心去閱讀的時候.
我們會發現其實四十四章六到八節的經文.
很強烈表達一個主題.
就是上主是獨一的真神.
這一種的一神觀其實和救贖的主題息息相關.
對於當時的以色列民而言.
他們已經身處在巴比倫已經有一段日子.
巴比倫當時充滿敬拜偶像的風氣.
在四十四章九至到二十節.
其實反映當時的人是會為假神去製造偶像.
並且會向這些的神像去叩拜.
向他們祈禱去求他們施行拯救.
我們一起看一看四十四章十七節.
它裡面這樣提到.
然後又用剩下的一半做了一個神明.
就是雕刻的偶像.
向這偶像服服叩拜.
向他禱告說求你拯救我.
因你是我的神明.
我們看見其實當時有一班人.
為了假神去製造偶像.
亦都去拜這些偶像.
甚至乎在四十四章十八節.
他們的經文才知道去形容他們.
其實是無知又不思想.
這班人是一班心硬眼狠耳朵發沉的群體.
以為他們所製造出來的神像.
可以幫到他們.
可以拯救他們.
但是在四十四章二十節提到.
這班人的舉動.
其實沒有辦法得到真正的拯救.

$^{81}$基於這種背景情況之下.
先知福術上主的自述之中.
刻意去表達上主他是真實.
以及他是獨一的神.
目的其實就是要對抗.
當時人民向偶像追求.
敬拜偶像的環境和風氣.
我們一起再一次去看一看.
第六至第七節.
關於上主的自述.
第六至第七節裡說到.
我是首先也是末後.
除我以外再沒有神.
自從古時我設立人.
誰能將我宣告陳明.
又為自己陳述呢.
讓他指明未來的事和必成的事吧.
我是首先也是末後.
這一番的自述.
表達上主其實是歷史的主.
歷史的開始歷史的終結.
都是屬於上主.
上主是一切的開始.
也是一切的結束.
他才是歷史上的主宰.
沒有任何神明可以和他媲美.
所以去到第七節上主提到.
誰能將我宣告.
指明又為自己陳述呢.
讓他指明未來的事和必成的事吧.
其實這句說話.
如果更貼馬索拉文本的翻譯.
其實我們可以參考新譯本.
新譯本裡面這樣翻譯第七節.
誰將我呢讓他宣告.
讓他述說又在我面前陳明吧.
自從我設立古時的子民.
未來的事和將要發生的事.
讓他們把他們述說出來吧.
宣告陳明述說.

$^{121}$三個其實都是同義詞.
而在新譯本的翻譯.
其實很全神的表達.
上主其實是向這班假神去發出挑戰.
就是哪一個神明.
能夠指明以及確實宣告將來能夠發生的事.
邀請大家揭一揭.
以賽亞書四十一章二十一到二十四節.
裡面講明一個神明.
怎樣才能夠證明他是真實.
這段經文有提到.
以賽亞書四十一章二十一到二十四節.
他裡面這樣講到.
耶和華說你們要呈上你們的案件.
雅各的君王說你們要提出你們的理由.
讓他們近前來告訴我們將來要發生什麼事.
你們要說明先前發生的事.
好讓我們思索.
或者告訴我們將來的事.
使我們得知事情的結局.
你要指明未來的事.
使我們知道你們是神明.
你們挖礦禍挖礦福.
看啊!你們屬乎虛無.
你們的作為也屬虛空.
那選擇你們的是可真悟的.
如果我們用這段經文去理解的時候.
假的神要證明他自己是真實.
就必須清楚指出過去所發生的事的原因.
以及對未來有一個準確的預測.
這裡講的預測.
不是一般性的預測未來.
而是關於以色列民被擄他的原因.
以及他們將來的發展.
但這些假神卻是沒有辦法做到.
唯有上主先至清楚知道.
以色列民他們被擄之後能夠回歸耶路撒冷.
甚至是上主親自去策劃這件事.
他用一個令人意想不到的方法.
去讓到以色列民回歸耶路撒冷.

$^{161}$這個方法就是興起波斯帝國.
假神在這件事上沒有辦法參與.
甚至預測不到以色列用這種方法來回歸耶路撒冷.
上主那番宣告其實是向以色列民表達.
他們身處巴比倫.
縱然充滿著偶像的地方.
看似這些偶像可以為人民帶來風調雨順的生活.
但卻只不過是虛無.
就好像44章第9和第12節所說.
這些事這些偶像對他們無益.
沒有辦法真正為他們帶來拯救.
唯有那位首先幕後掌管歷史的上主.
才是以色列民真正的救贖主.
這位掌管歷史的主.
怎樣成為以色列百姓的救贖呢.
各位姐妹我們觀察一下第6節.
先知在引述上主的自述之前.
他提到這位上主是以色列的君王.
亦都是以色列的救贖主.
救贖我們很常聽.
但救贖這個觀念是來自利未記25章25節.
他裡面根據律法.
如果有一個人因為貧困.
他去將自己所屬的土地去抵押給人.
他所抵押的土地.
必須由他最近的親屬去幫他贖回.
如果賣地的人沒有辦法.
沒有能力去贖回.
大家知不知道有什麼方法可以拿回.
將我的土地抵押給人.
只有一個方法是可以拿回這塊土地.
50年的欺凌.
所抵押的土地.
會無條件由賣主去還給抵押的主人.
但問題是如今被擄的一班以色列民.
他們要回去應許之地的時候.
大家都知道他不可以隨便選一個地方去居住.
而是回去他們應該從列祖的時代.
被分派的土地裡面去居住.
問題就是.

$^{201}$我已經有大班人.
大班的以色列民.
我已經被擄.
誰幫我拿回這塊土地.
誰幫我贖回這塊土地.
沒有.
幾乎是沒有.
先知提到.
上主是以色列的救贖者.
即是上主才有能力和資格.
去幫這班被擄的以色列民.
去贖回他們自己原本的地方.
上主贖回土地這個概念.
其實也是出自利未記.
二十五章二十八節希年的觀念.
上主如今興起波斯.
去讓到以色列民.
能夠回到自己所屬的土地.
而且以色列民.
不需要附上任何的抵押品.
去贖回自己的土地.
這個情況就好像希年的觀念一樣.
所以當以色列民歸回應許之地的時候.
我們可以理解.
那一年其實就是上主的希年.
先知在第六節提到.
上主的身份是以色列的救贖主.
其實就是要指出希年這個觀念.
上神沒有辦法幫他們贖回土地.
唯有那位能夠掌管歷史興起波斯.
甚至用這個意想不到的方法.
的上主才可以贖回他們原本的應許之地.
這個上主才是真神.
聽姊們,其實這個世界上的人.
很多都和當時被擄巴比倫的以色列民很相似.
這個世界的人.
其實仍然被很多的偶像去迷惑.
這些偶像不是單單指向一些神明.
更加可以是一些觀念,物質.
甚至是自己.

$^{241}$以致他們看不到上主的真實.
他們仍然繼續沉淪在偶像的虛無之中.
他們不將虛無的偶像視為真實.
將真實的上主視為虛謬和虛假.
甚至有些否定上主的存在.
上星期我和黃秀英姊妹招呼一班文藝學院的中一同學.
我們當中有分享關於這個世界的來源.
很多同學都相信這個世界是由大爆炸而來.
但他們選擇否定這個世界是上主創造.
當中的原因是覺得上主創造很不科學.
大爆炸的說法相對更加可信.
他們認為上主他們根本看不見.
沒有辦法證實.
不過有些同學在當中再分享多一點的時候.
他們相信這個世界有鬼.
當他們聽到這一刻的時候.
我心裡有一句說話就是.
鬼都看不見.
你都沒有辦法用科學去驗證.
但為什麼他們會寧願相信這個世界有鬼.
但否定上主呢.
其實這班同學可以說是這個世界的縮影.
正正顯出人心裡那種對超自然的安全感和渴望.
以及他們理性主義之間的矛盾和掙扎.
其實對他們來說.
上主是明明可知的.
但他們選擇否認和拒絕去承認.
寧願他們將信耶穌當為一種迷信.
他們只信科學.
但心裡一層歸根究底.
其實他們只是信自己.
他們以自己為中心.
他們選擇自己認為哪一樣才是真實.
哪一樣是虛假.
哪一樣可以幫到自己.
哪一樣對自己實際沒有任何幫助.
在我身邊的朋友之中.
其實他們都是不信耶穌的人.
而他們都是一些專業的職業.
從事一些專業的職業.

$^{281}$但他們當中的人.
不少人都相信一些民間的傳統.
甚至是相信風水.
背後的原因是因為他們很想知道自己的前途.
以及前途.
當然他們亦都為了自己去求平安.
於是他們會留意風水去擺風水陣.
甚至乎他們逢新年都會去廟裡走一圈.
讓自己轉一轉運.
其實他們不相信耶穌是獨一的神.
是真實的時候.
但是他們轉為去求一些虛妄的事情.
背後的真相其實.
他們很想企圖去透過一些等價交換的關係.
來掌握掌控自己的命運.
所謂求神明.
本質上其實都是以自己為中心.
他們以為自己付出所付出的金錢製品.
神明就會保佑自己.
好像和我們的信仰.
又或者我們的救贖觀.
贖回觀.
恰恰相反.
無關乎在羅馬書保羅他這樣說.
因為他們雖然知道上帝.
卻不把他當作上帝榮耀他.
他們的思想變為虛妄.
無知的心昏暗.
他們自以為聰明.
反成為愚昧.
將不能撈壞的上帝的榮耀變為偶像.
仿照必撈壞的人.
飛禽走獸爬蟲的形象.
人不信耶穌.
不信上主是獨一真神.
除了是來自人的自我中心之外.
其實亦都受制於文化歷史社會價值觀等等的因素.
好像日本.
我們ViuTV都去了日本那裡.
有一年的短宣.

$^{321}$當地其實都分享過幾次.
當地的基督教信徒的比例其實不到1\%.
是屬於福音眼圖.
背後除了他們固有的神道教.
以及一系列的歷史包袱.
這些歷史包袱包括.
奧母真理教以及統一教的事件.
這些都影響他們對於宗教信仰的看法.
文化亦都是一個很大的集體枷鎖.
特別日本人很講禮節.
他們不輕易去和其他人成為親密的朋友.
昨天我和一些年輕人說.
IG上的綠色朋友.
很難.
不知道什麼叫綠色朋友.
可以問一下年輕人.
日本人很難和他們做綠色朋友.
做一些很親密的關係的朋友.
其實他們是屬於一個封閉式的社交.
他們只會相信自己圈內的人.
如果要和日本人傳福音.
你必須要打進這個社交圈子.
你並且要得到他的信任和認同.
他覺得你值得信得過.
他才會聽你說話.
日本人就是有這種文化.
另外日本人也有另一種文化.
就是OT視日常.
即是加班.
加班他們不加班不行.
如果有一個人很想和公司的老闆.
或者和上司說.
我這個禮拜不能進去我要回教會.
我要回團體.
我不加班可以.
對於我們香港人來說很輕易.
和老闆說一句話.
不回就不回.
有多要命呢.
但對於日本人來說.

$^{361}$其實是一件很難開口的事.
因為背後其實他們覺得.
這種行為很不合群.
我不加班.
我只是和我的同事很疏離.
我背叛了整班同事.
甚至會為公司和同事帶來不便.
這種為人帶來不便.
對於他們來說.
是一種極大的羞恥.
極大的羞恥.
其實在宣教學裡面.
這種他們認為為人帶來不便.
所帶來的羞恥.
是屬於羞恥和尊榮的文化.
其實每一個國家.
都會存在著不同的文化.
而在宣教學裡面.
其中會提及一個主題.
叫做三維福音.
三維福音就是將人.
分為三種類型的文化.
除了剛才提到的羞恥.
尊榮的文化.
另外兩種就是.
罪疚,清白以及懼怕能力.
第一個的文化.
罪疚,清白就是指人.
他企圖去靠一些善行.
修行去贖罪.
去填補他那種的罪疚感.
但我們都清楚知道.
其實人沒有辦法用種種的行為.
和功德去救贖自己.
去贖回自己.
到頭來所做的事其實很虛無.
我們都知道唯有耶穌基督.
才是那位生命裡面真正的救贖主.
另一種文化是關於.
懼怕又或者我叫恐懼能力的文化.

$^{401}$恐懼是指人很怕一些鬼神看不見的東西.
又或者很怕未來那種不可預期而帶來的不安.
對於我們華人文化其實都很常見.
所以以至於有些人為何要去看風水.
去拜黃大仙等等這些虛無的偶像.
但唯獨是我們那位掌管歷史.
是首先也是幕後的上主.
才能夠為我們帶來那種真正的平安.
世界上的人其實仍然沉淪於種種的虛無之中.
對於他們來說他們不是不需要神.
而是他們被自身的文化和社會結構等等去蒙蔽.
去用虛無去取代真實.
這個其實就是我們仍然需要去推動.
大家要去往普天下去的原因.
如果上主只是香港人的上主.
我們只是關心自己人就夠了.
但上主說除我以外再無神.
世上的人仍然沉淪於種種虛無之中.
仍未找到那個真實的依靠.
所以在第八節 四十四章第八節經文裡面提到.
你們不要恐懼也不要害怕.
我豈不是從上古就告訴你們.
並指示你們不了嗎.
你們是我的見證人.
除我以外再沒有神.
誠然沒有盆石.
就我所知一個也沒有.
首先要解釋一下第八節下半節裡面說到.
仍然沒有盆石.
就我所知一個也沒有.
這番說話其實所要表達的意思就是.
世上還有哪一個可以自稱為盆石.
我實在不知道.
這番說話是上主的自問自答.
他問題的核心就是作為整個宇宙以及歷史的開始和終結的上主.
以他的全知全能.
他都不知道除了他自己之外.
還有哪一個是盆石.
盆石的意思是可以依靠和給予保護的意思.
連這個獨一主宰他都不知道的說話.

$^{441}$世界上還有哪一個神明可以說有另一個神明.
可以比上主更能被人依靠和提供保護呢.
作為上主的自問除了要認識上主是歷史的主宰.
真正的救贖主之外.
同時經歷上主救贖之後.
我們卻是要成為上主的見證人.
成為見證人不是因為我們有什麼特別的恩賜和才能.
而是來自上主的獨特揀選.
是上主給我們的恩典.
一切都是來自上主的主權.
所以當我們經歷上主的拯救之恩.
我們被賦予成為他的見證人.
然而我們不是只向香港人做見證.
而是向世界未得之民做見證.
上主不只是屬於以色列的主.
也不是單單屬於香港人的主.
上主是全世界的主.
這個世上所有人都是由他創造而來.
都是他所愛.
所以我們要見證上主不只是向香港人做見證.
而是向普天下的人做見證.
上主並不只是向以色列人傳得救的訊息.
也不只是優待我們香港人.
只是讓香港人得救.
世上還有很多人仍然被罪捆綁.
仍然未被上主贖回.
在《士徒人傳》17章30至保羅這麼說.
世人愚昧無知的時候.
上帝並不追求.
如今卻吩咐各處的人都要悔改.
弟兄姊妹.
今天《義創下書》44章8節.
上主對我們說.
你們不要恐懼不要害怕.
你們是我的見證人.
世上仍然還有很多福音眼土.
但是上主提醒我們.
他是那位首先是幕後的歷史之主.
每一個地方國家的文化枷鎖.
其實他都有能力去打破.

$^{481}$然而他是邀請我們去參與其中.
就是要我們成為他的見證人.
差傳不只是宣教士被賦予的召命.
而是每一個被上主贖回的人的使命.
鼓勵大家可以先從禱告開始去紀念世界未聞福音的國家地方去禱告.
同樣大家不知道要從哪個國家去入手.
哪個地方去入手.
鼓勵弟兄姊妹這個星期.
可以嘗試在網上.
按進各大差會的網頁去瀏覽.
裡面當中會簡介到不同的地方國家的需要.
嘗試按進去看看.
開一下自己的眼界.
開一下這個世界其實不是只有香港.
這個世界仍然有很多人未得聞福音.
面對困境.
嘗試去為這些國家地方人民去祈禱.
嘗試去以禱告作為關心差傳的開始.
踏出見證世上唯一的盤石.
讓我們痛心低頭禱告.
上主我們都求你.
你是那位守仙門後的主.
你是那位歷史的主.
同時你也是這個全宇宙的主.
求你叫我們去不單單只是關心我們所居住的地方.
世界上仍然有很多人都未得聞福音.
他們都是屬於你所愛.
我們不單單只是關心我們自己的地方.
求主你開我們的眼界.
讓我們有你的視覺去騙測整個地方.
整個地球.
我們不單單只顧自己的事.
也要顧世界.
顧很多仍然未得聞福音的國家地方人民.
他們仍然需要被上主贖回.
回到真正的愛裡.
求你讓我們開這個眼界.
去關心這個世界的需要.
去關心這個世界很多未聞福音的地方和國家.
求你幫助我們.

$^{521}$我們要向祂禱告,奉主耶穌基督,聖名祈求,阿們..
\newpage



\section{那鴻書 1:2-3}
\label{sec:032473}
\textbf{榮耀的釋放}~~~~蕭壽華牧師
\newline
\newline
連結: \href{http://www.church.com.hk/acms/content.asp?site=cdc\&op=show\&type=product\&code=032473\&layout=sermon}{\texttt{032473}} ~~~~ 語音日期: 2025-12-21
\newline
\newline
\hyperref[sec:032472]{\small{< < < PREV SERMON < < <}}
~
\hyperref[sec:index_chronic]{\small{[返順時目]}}
~
\hyperref[sec:032474]{\small{> > > NEXT SERMON > > >}}
\newline
\newline
那鴻書 1:2-3
\newline
\begin{longtable}{cl}
\hline
\hline
章節 & 經文 (和合本修訂版)\\
\hline
1:2 & \begin{tabularx}{0.7\textwidth}{X} 耶和華是忌邪、報應的神。耶和華施報應，大有憤怒；耶和華向他的敵人報應，向他的仇敵懷怒。 \end{tabularx} \\ \\ \relax
1:3 & \begin{tabularx}{0.7\textwidth}{X} 耶和華不輕易發怒，大有能力，但耶和華萬不以有罪的為無罪。他的道路在旋風和暴風之中，雲彩為他腳下的塵土。 \end{tabularx} \\ \\
[1ex]
\hline
\hline
\end{longtable}
$^{1}$此錄音或錄影的版權是屬於基督教宣導會不覺堂所有.
我們的屬靈傳統之一是以真理教導為本.
歡迎你將此錄音或錄影的連結傳給親友.
懇請注意未經授權請勿轉錄或剪輯本錄音或錄影.
甚願好像《使徒行傳》二十章三十二節下所講.
「這道能建立你們,叫你們和一切聖聖的人同得基業」.
各位弟兄姊妹平安.
也思念主耶穌他帶給我們的救恩.
我請大家和我一起祈禱我們同心仰望主.
滿有恩惠滿有憐憫的上帝.
雖然我們見到人世間有很多叫我們難過灰心的事情.
但是主啊你帶給我們盼望.
在黑暗當中你讓光照進我們的生命中.
因此雖然在困難裡面我們仍然可以快樂.
仍然可以感恩.
思想你的話.
我們祈求天上的主.
藉著你的話.
藉著你寶貴的聖靈.
在我們心中工作.
讓我們更加深刻的知道.
你帶給我們的是何等大的盼望.
我們就更加知道.
每一天歡喜快樂地跟隨你.
聽我們祈禱奉主耶穌成名祈求.
今天跟大家思想的是兩段經文.
一段是聖經舊約的拿紅書.
第二段就是盧家甫音撒迦利亞的一個頌詞.
當我們這段時間留意周圍的世界.
社會的情況的時候.
你一定都會體會到.
這個世界越來越多各種的災難.
有一些是自然的災難.
有一些是人帶來的禍患.
在大埔的火災之後.
很多人都說.
他說沒錯天災很恐怖.
但是人禍比天災更加可怕.
很多的災難背後是因為人的罪.
帶來的很多很多的傷害.

$^{41}$今天是聖誕主日.
我們仍然可以因為聖誕.
我們可以快樂可以歡慶可以感恩.
因為當我們思想到人類的問題.
很多是因為罪的緣故.
就體會到當主耶穌降生的時候.
帶給我們榮耀的釋放.
主耶穌的降生.
叫我們從很多的捆綁很多的情慾很多不同的私慾裡面.
釋放出來.
叫我們能夠在神的面前脫離罪惡.
進入永遠的和平永遠的平安.
因此我們知道在一切的情況之下.
我們仍然可以感恩.
因為知道上帝給我們最徹底最完全的釋放.
我想回看一段新聞.
關於黃福遠的情況.
一位姓張的先生.
他講述他自己當時在自己的家裡.
但是四面被火海包圍著.
整個環境是很窒息.
因為周圍濃煙密佈.
他不可能出去離開自己的家居.
只能夠躲在自己的家裡.
一種很無助的情況.
自己心裡也想定是輸了.
就在那個時候.
他說突然間有消防員是劈門而入.
入了他的家.
然後就見到他.
就將他背起.
又帶他出去.
在一個很漆黑的環境裡面.
帶他走下樓梯.
過程裡面他說我曾經跌倒了.
我想自己起來.
嘗試摸一下牆自己走.
但是一碰到牆.
是完全不可以碰的.
因為牆是火熱的牆.

$^{81}$牆壁是很燙手的.
周圍有些水.
但是那些水是滾水.
這個消防員告訴他.
你不要走.
你上我的背.
因此在那個時候.
他再次爬上這位消防員的背上.
消防員就向前走.
在他頭上有一盞燈照著前面的路.
就帶他一級一級的下去.
下去得到幫助.
當我看到這個畫面的時候.
我心裡就想到.
其實我們每一個.
每一個在地上生活的人都體會到.
我們其實都被各種的事情捆綁住.
特別被罪惡的權勢捆綁住.
自己都覺得沒有盼望.
但是突然間我們看到.
神給我們一個很奇妙的釋放.
祂來到我們當中.
將我們從死亡當中拯救過來.
讓我們能夠領受祂給我們的平安.
各位電影姐妹.
是你熟悉的訊息.
但是在聖誕期間.
你繼續去回想.
是一件何等大的恩典.
何等大的釋放.
截至今天.
你帶著平安帶著盼望.
繼續走你人生的路.
我們剛才提及.
我們今天會讀《拿紅書》.
《拿紅書》很特別.
是一本很重要的小先知書.
《拿紅書》裡面特別提及一件事情.
就是神藉著先知拿紅.
向當時的亞述帝國.

$^{121}$也是首都尼尼美.
宣佈一個重要的審訊.
就是神要向他們宣佈.
他們要快被滅絕.
當時的猶太人.
他們知道這件事情的時候.
他們很高興.
因為猶太他們很清楚知道.
亞述帝國是他們歷代的強敵.
或者說勢手.
他們不能夠面對這麼大的勢力.
因為尼尼美是一個很巨大的城市.
不單止是當時地方的人口來說.
而英國的國家.
其實是一個擁有極強大軍隊的國家.
這個國家也是兇殘.
兇暴殘酷著稱.
他們面對著他們的敵人的時候.
就會不留情.
不斷殘殺很多的敵人.
但是他們因為太強大.
所以沒有人能夠抵禦他們.
所以他們稱為無法抵禦的軍隊.
但是那個時候他們聽到一個訊息.
就是神要滅絕這麼強大的國家.
或者說神好像沒有憐憫他們.
為什麼會滅絕他們呢?.
如果你留心.
聖經裡面給我們留心看到的就是約拿書.
你知道約拿書.
若拿是向哪裡發出警告呢?.
尼尼美城.
當時尼尼美面對著神的審判.
但是岑念敏他們差派若拿.
去向他們宣告悔改的訊息.
他們迴轉歸向上帝就得蒙赦罪.
尼尼美當時他們真的願意悔改.
後來神就赦免他們的罪.
但是過了大概一百年之後.
這個城市又重新再陷在罪惡當中.

$^{161}$有各種邪惡的事情.
這個時候神透過拿紅再次發出他的審判.
宣告要滅絕這個城市這個國家.
在聖經拿紅書一開始的時候.
聖經就說.
耶和華是忌邪報應的神.
耶和華施報應大有憤怒.
耶和華向他的敵人報應.
向他的仇敵懷怒.
然後聖經說.
耶和華不輕易發怒大有能力.
但耶和華萬不以有罪的為無罪.
聖經特別說到一位上帝是忌邪的上帝.
我們不要誤會這裡所說的忌邪.
不是說人那種因為自私而妒忌別人比自己好.
很想傷害別人否定別人那種的妒忌.
聖經說的忌邪.
是說到一位上帝.
他是喜愛公義恨惡罪惡.
上帝的本性是恨惡罪惡.
他不會容許邪惡代替良善.
正如我們不會希望我們身體會出現癌細胞.
而癌細胞吞噬其他健康的細胞.
因為上帝他忌邪.
也意味著他不願意邪惡的事情去傷害良善.
正如一個家庭裡面夫妻有盟約的關係.
但是妻子背棄了丈夫寧結新歡.
丈夫不會接受到.
因為他愛他的太太.
聖經說到一位上帝他是忌邪的上帝.
當你讀到舊約聖經的時候你會發覺.
忌邪之後接著說上帝是大有憤怒.
他是滿有憤怒.
英文所說the wrath of God.
神的憤怒淋到那些在罪惡當中.
不願意悔改歸向上帝的人.
聖經提到神的憤怒.
也顯明了上帝是完全公義的上帝.
因為上帝他不會去胡混黑白是非.
他不會去顛倒是非.

$^{201}$因此聖經說他不會以有罪的為無罪.
今次你只要想一想.
在很多不同地區的法庭.
本來應該審訊很多犯罪的人.
但是因為混亂了是非黑白.
帶來多少傷害多少問題.
因此你知道當神發怒的時候.
正是顯出他公義的時刻.
同時神的怒氣也顯明他是愛.
他不會容許不義來傷害義.
他不會容許不義來傷害生命.
因此你看到當聖經說到.
因為上帝大有憤怒的時候.
正正是我們要感恩的時候.
就是他要潔淨罪.
他要除滅罪.
但當我們剛才讀到這段經文第三節的時候.
聖經就立刻說.
耶和華不輕易發怒.
沒錯神是會憤怒的.
但是他不輕易發怒.
意味著他很有寬容.
人犯罪之後.
神會給很多機會呼籲人悔改.
好像他呼籲若那去勸勉你你你悔改一樣.
但是我們不要以為這個寬容.
是一個自己可以繼續犯罪的機會.
不要誤會神的寬容是表示他很軟弱.
很多人描寫這位上帝.
是一位白髮奴翁.
坐在沙發上滿臉慈祥地告訴你.
沒事沒事沒問題任你做.
神的寬容不代表神的軟弱.
聖經就說他是大有能力的.
他能勝過一切的邪惡.
勝過一切強的勢力.
包括了亞述大帝國.
神是大有能力的.
因此聖經就說你不要輕率去看你的罪.
以為神會不理不知.

$^{241}$因此下帝經文就說.
他發憤怒的時候誰能立得住呢?.
他發烈怒的時候誰能當得起呢?.
他的憤怒如火傾瀉而出.
盤石因他崩裂.
聖經是自然界一些很強的災難.
來描述神自己的憤怒.
聖經說他必以將來的洪水淹沒其地方.
說來來來的地方.
或者說惡人的地方.
又驅逐他們進入黑暗.
聖經描述這些事情的時候.
讓我們看見.
這位上帝是大有能力的.
他不會面對惡勢力的時候.
會退縮.
他有能力去勝過一切的罪惡.
但是聖經在說的時候.
第六第八節之間.
聖經仍然不斷強調一件事情.
耶和華本為善.
因為他是良善的.
他樂意去保障人.
保護人.
他樂意給我們平安.
但是他期望.
尋求他的人.
先會得到平安.
所以跟著他說.
並且認識那些.
什麼人?.
投靠他的人.
投靠他的人是神所認識的.
是神樂意向他們施慈愛的.
去保護,去看顧他們.
去幫助他們.
我想在這段時間.
我們常常都會想起.
大埔火災的事情.
很多人都提起聖經.

$^{281}$《路加福音》第十三章的經文.
這段經文大家可能會記得.
當時主耶穌提到.
就是當時的一座樓.
叫做西諾亞樓.
這座樓倒塌了.
其中有十八個人死了.
當時的人就說.
為什麼會有這樣的災難呢?.
是不是這些人.
他們犯了很多的罪.
所以上天不容.
容許他們在塔樓當中死去.
主耶穌當時告訴他們.
你們不要以為.
他們比其他人更有罪.
我告訴你們不是的.
你們若不悔改.
都照樣要滅亡.
要滅亡.
不是因為.
他住的那座樓.
不是很安全.
這是一個原因.
但是更重要的是.
人在罪當中.
每一個人.
如果不願意承認他的罪.
悔罪歸向上帝.
都要滅亡.
然後聖經就用了一個比喻.
主耶穌就說.
有一個園丁.
他跟他主人說.
主人告訴他.
我三年來想找一些果子.
聖經說的是一些悔罪的果子.
但是都找不到.
好了,要將它砍掉了.
但是園丁說.

$^{321}$他說不要止.
在他周圍挖一些.
挖鬆一些土.
加多一些肥田料.
它可以更好地生長.
但是最後一句.
你看到.
以後若結果子.
變罷.
就不要砍掉它.
但是不然再把它.
陷了.
主耶穌很想讓我們知道.
神是滿有憐憫.
祂是良善的.
但是祂仍然是公義的上帝.
祂不會姑息罪.
在足夠的機會.
之後仍然不願意悔改.
最終.
是要面對審判.
是被陷掉的.
因此當我們今天來到聖誕節的時候.
我們想回主耶穌的救恩的時候.
我們首先去.
清楚地知道一件事情.
神的寬容,神的拯救.
不是一個.
讓我們可以輕率的.
一個時刻.
或者說我們會以為.
祂會遷就我們的罪.
但是不是.
神是公義的主.
我們再看.
再看下面一段經文的時候.
我們看到一件事情.
大家都會知道.
在聖經裡面.
當撒迦利亞說出這個寓言的時候.

$^{361}$當時是主耶穌.
將要降生的時候.
但是當時的人.
很清楚知道.
在新約開始之前.
當時的社會.
有一個叫做四百年的.
黑暗期.
或者稱為一個沉靜期.
在那段時間裡面.
雖然說很沉靜.
但其實不是在說平安.
好安穩.
是在說當時.
整個世界.
在一個很動盪不安的日子.
各種的災難.
各種人為的災禍.
在當時出現.
當時的社會.
當時的政治是極之混亂.
同時也不再看到.
有神透過先知.
讓他們有默示.
讓他們有提醒.
一切都好像進入一個很絕望的處境.
就在這個時候.
一位撒迦利亞.
竟然遇見天使.
寓言讓他知道他的太太.
年邁的老邁的太太.
可以生孩子.
當時撒迦利亞.
不相信.
因為他的不信.
他不能再說話.
他突然間失聲.
直至孩子出生.
事情實現的時候.
他開聲.

$^{401}$他一開聲的時候.
他就說這句話.
主以色列的神.
是應當稱頌的.
一開口的時候.
他心裡面湧不出那種讚美.
含忍不住那種感恩的心.
就讚美上帝.
神是值得稱頌的.
為什麼?.
因祂眷顧祂的百姓.
因為這位上帝.
眷顧百姓.
雖然人在罪當中.
帶來了給自己的傷害.
對別人的傷害.
帶來了很多的災難.
但是神仍然眷顧.
這個眷顧原來的意思.
是神來到的意思.
神來到的意思.
因為上帝來到.
他們等候了四百年.
在那種困難.
黑暗當中.
好像沒有盼望.
但是神來到.
拯救他們.
有點像.
在火場當中.
被火圍困.
心裡面是絕望的.
突然間消防員來到.
聖誕節的時候.
上上上到.
神來到.
你的生命當中.
帶給你盼望.
帶給你奇妙的恩惠.
因祂眷顧祂的百姓.

$^{441}$為他們施行拯救.
在他僕人大衛的家中.
為我們興起了拯救的國.
大衛的家中.
是指著彌賽亞.
彌賽亞要來到了.
他成為了拯救的國.
當時的人.
當說到國的時候.
是在表示能力.
好像野牛的國.
大有能力.
去抵擋很多他的敵人.
很多的野獸.
師兄說到.
神的拯救是大有能力的.
是勝過很多的罪惡的.
他能夠成為了.
我們生命中真正的幫助.
師兄裡面提到.
他拯救了我們.
在當時來說.
聖經在七十一節繼續說.
他拯救我們脫離一切.
是我們之人的手.
很明顯他們就會想到.
因為上帝要救我們脫離.
很多政治的.
很多的捆綁.
壓力.
特別是當時的社會.
大家知道當時猶太國.
他們是被羅馬政府.
來到拯救.
被羅馬政府來到統治.
他們在羅馬政權的.
壓迫之下來到生活.
他們可能想到.
因為利塞雅來到.
就會幫助我們脫離了.

$^{481}$這些政治的權勢.
但是再看的時候.
你會發覺原來不是.
政治的權勢.
而是罪惡的權勢.
在七十四節.
這段經文我請大家讀一讀.
好不好.
無所懼怕去事奉神.
不再是事奉哪個政權.
不再是服侍世界.
服侍自我.
而是事奉那位永生的上帝.
請提到.
是無所懼怕.
用聖經.
來祈禱.
去事奉神.
請提到.
無所懼怕.
去事奉神.
是無所懼怕.
無所懼怕.
用聖潔去事奉祂.
這裡正是意味著.
夢神拯救的人.
他們已經是屬於神.
因為聖經說.
人非聖潔不能見主.
這些人能夠.
可以以聖潔.
來跟隨上帝.
意味著他已經屬於上帝.
在神的裡面.
他們是用.
公義事奉上帝.
顯示出他們已經成為神的子民.
被神潔淨.
他們可以行在光明當中.
所以他們可以用.

$^{521}$聖潔公義來事奉.
一位天地的主.
在過程當中.
無所懼怕正正是表達出.
那種生命的關係.
面對上帝.
不再需要面對.
神的憤怒.
你記得我們剛才說到.
神的憤怒.
面對著.
上帝不再需要.
去見到.
神的憤怒.
在當中安然地與.
神同行.
我.
記得有一次看到一段.
分享一位宣教士.
說到他自己在一個.
國家的經歷.
宣教士提及到.
他自己在當時.
他去到一個地方.
學了當地的.
考了當地的牌.
走一條路的時候.
就開車駛入去.
但是他.
沒有留意到.
路口寫了一個牌子.
不准駛入.
但是他很奇怪.
因為那條路平時是雙程路.
但是那天後來.
剛好改為單程路.
單線行車.
所以他駛入去之後.
糟了,錯了.
但是竟然發覺.

$^{561}$已經有警察.
捕捉在那裡.
好像知道有人進來.
馬上就抄牌.
罰了三分,罰了款.
當地的.
一些規矩.
如果罰六分就要停牌.
他就因此回家.
他說以後.
一出街駕車.
很尷尬.
很不安.
因為知道隨時會被抄牌.
後來他就心想.
為什麼會這麼不安.
他說因為他有了.
扣分的記錄.
這個記錄不斷告訴他.
你犯了法.
你做錯了.
但當他想到那裡的時候.
他突然間心裡有一個.
很大的感恩.
我過去所扣的分.
主已經完全.
赦免.
抹圖.
將那些扣分記錄.
全部抹開.
我不需要面對神的憤怒.
我在神的面前.
被祂完全接納.
各位靈姐妹.
當你想到聖經裡面所說.
你可以無懼地.
用聖潔公義.
事奉那位生命的主.
背後是何等大的恩惠.
然後你再讀的時候.

$^{601}$你會更加體會到.
聖經裡面說到.
在過程裡面.
神將我們從死刃的地方.
帶到平安的地方.
撒迦利亞繼續說的時候.
他提到他自己的孩子.
大家知道施洗約翰.
出生之後.
這件事情.
他是帶領很多的百姓.
來悔改.
歸向上帝.
為基督預備道路.
什麼叫做預備道路呢.
當人願意悔改的時候.
就是鋪平了一條路.
讓基督.
能夠來到他的生命中.
當他願意悔改的時候.
就鋪平了基督.
前來拯救他的路.
不是約翰.
能夠去拯救.
約翰是呼籲人悔改.
但當人悔改的時候.
基督就來到.
拯救人.
因此叫我們認識.
舊恩.
我們因罪得到赦免.
而認識舊恩.
經歷舊恩.
我們剛才.
看到.
約翰來到的時候.
他預備.
這一條路.
因此你記得當他來到的時候.
他很清楚和周圍的人說.

$^{641}$看到主耶穌的時候就說.
看吶 神的羔羊.
然後呢.
除去世人.
罪孽.
我們沒有其他的盼望.
是唯有主耶穌基督.
他能夠除去我們的罪孽.
叫我們領受盼望.
他幫助我們.
除去我們的過犯.
我們的罪疚.
我們的罪污.
讓我們心靈領受那份平安.
我很歡喜一節聖經.
就是聖經二兆書記中的.
四十三章二十五節.
聖經裡面說.
唯有我為自己的緣故.
你留意到.
不是因為你好.
不是因為你究竟乾.
而是因為神為他自己的憐憫.
他就.
除去你的過犯.
不再.
記得你的罪惡.
我曾經好多年前.
跟弟兄姊妹說過一個.
自己好深刻的經歷.
什麼經歷呢.
就是在我年幼的時候.
我和我的哥哥.
都幾多時候有爭執.
有些時候.
甚至打起架來.
在打架過程.
小時候.
都不知道自己有些蠻力.
有一次.

$^{681}$哥哥有些什麼事情.
令到我很不開心.
我就一手推一推他.
我想不到.
竟然他整個人掉在地上.
而且是正面面對著地上.
啪一聲.
門下崩了一半出來.
我想不到自己.
竟然可以這樣傷害自己的大哥.
他站起來.
沒有流血.
不過就崩了一邊.
但是因為家庭.
當時是基層的家庭.
父母都沒有什麼.
餘錢可以給他去補.
這一次我哥哥有幾年.
他都崩了一邊.
但是對我來說.
每次當他笑的時候.
就是我.
羞愧的時候.
我每次看到他笑.
那隻崩了出來的牙.
我就好內疚.
為什麼我當時這樣做.
我明白.
什麼是內疚.
什麼是罪疚.
直到後來當我哥哥自己出來工作.
他賺到錢.
他去補牙.
然後對著我.
笑一笑.
沒事了.
幾大一種的內心釋放.
各位聽眾你明白的.
在我們人生裡面.
如果我們要數我們過去人生.

$^{721}$曾經所經歷的事情.
那些曾經.
自己都不敢再提出的.
惡習.
那些曾經傷害別人.
那些曾經怎樣詭詐.
去得到一些利益.
每當你想起來的時候.
你自己有時候都不會原諒自己.
但你看到聖經給你的應許.
在主耶穌的救恩裡面.
不再記得你的罪惡.
祂不再記得你的罪.
也都呼籲你.
放心在祂的恩典裡面.
你忘記這些一切過去的事.
當聖經裡面提及這件事情的時候.
你知道不單止神赦免你的罪.
是跟著經文再說.
因我們神憐憫的心腸.
叫清神的日光.
從高天臨到我們.
照亮坐在黑暗中.
死淫裡面的人.
沒錯我們過去.
真的好像.
坐在死淫裡面.
我們沒有盼望.
我們在黑暗當中.
神的光竟然照進.
我們的生命當中.
叫我們能夠.
去到平安的路上.
但這一切你知道.
都是因為一個原因.
這個原因就是.
神憐憫的心腸.
你記得我們開始的時候.
講到神的眷顧嗎.
神的眷顧 神來到了.

$^{761}$祂拯救我們.
叫我們脫離神給我們.
應該有的公義的審判.
但你想想.
神現在.
祂讓高天的日光.
意味著一種.
從天而來的.
榮耀的光.
這些一切.
是因為祂的憐憫.
我自己.
很喜歡這首聖詩.
這首聖詩是我很多時候.
遇到一些困難的時候.
我心裡就會唱出來.
這首詩歌.
是主凡事引導.
詩歌第一節說.
救世主凡事引導我.
我何須別有所求.
主愛憐我.
何用多疑.
我一生夢主願有.
我通常自己會用英文.
來唱.
但每當我唱到.
Can I doubt His tender mercy.
這就變成一個很大的提醒.
我遇到一個問題.
難道神因此不愛你.
你忘了神對你一生的.
看顧是多麼大.
因此我勉勵自己.
不要懷疑.
不要懷疑神的憐憫.
Who through life.
has been my guide.
一生一世.
祂一路眷顧我.

$^{801}$保守我.
今天面對新的問題.
祂何嘗不會.
像過往一樣.
一路來施恩.
保守我.
Tender mercy.
這句詞.
神的很柔軟的憐憫.
上上給我很深刻的鼓勵.
亦因為是神的憐憫.
你會知道.
就是叫清晨的光.
從高天臨到你.
將你的腳.
引到平安的路上.
這個平安不是說.
沒有困難這麼簡單.
而是一個完全的.
生命的尾線.
的一種狀況.
神將平安賜給你.
叫你能夠可以.
一生走在一條無恩的路上.
各位靈靈姐妹.
沒錯我們這段日子.
我們都會有很多.
見到災難的事情.
甚至包括自己有的經歷.
但我們仍然.
心裡面可以很感恩.
因為神已經將.
榮耀的釋放賜了給我們.
我們不再需要.
面對著神的憤怒.
我們在祂的憐憫當中.
我們進入了平安.
因此你.
繼續放膽.
面對你人生的問題.

$^{841}$帶著那份平安.
去每一點.
去經歷神的恩典.
因為每一個困難.
都是你盡心.
去經歷神恩典的時刻.
我們一齊來祈禱.
我們的主.
我們的神.
上人將你.
致勝致榮.
可重可畏.
你的憐憫.
大過我們一切的鴛鴦.
你的慈愛.
遮蓋我們一切的罪.
主我們在罪當中.
曾經掙扎.
我們曾經.
深刻的體會到.
那種無助.
自己那種的困惑.
但是主啊.
你讓你的救恩.
來到我們生命中.
照耀我們的生命.
將我們從黑暗當中救拔出來.
讓我們進入了.
光明的路.
主啊我們為了你奇妙的救恩.
求你教導我們.
在前面日子.
我們繼續放膽跟隨你.
靠著你的恩惠.
讓我們無論有什麼困難.
我們仍然相信.
你的憐憫.
永遠是我們的把握.
是我們可以.
寬起跟隨.

$^{881}$的一種基礎.
聽我們祈禱.
奉耶穌基督.
成名祈求.
阿們.
\newpage



\section{使徒行傳 16:6-10}
\label{sec:032474}
\textbf{直搗羅馬}~~~~鄧瑞強牧師博士
\newline
\newline
連結: \href{http://www.church.com.hk/acms/content.asp?site=cdc\&op=show\&type=product\&code=032474\&layout=sermon}{\texttt{032474}} ~~~~ 語音日期: 2026-07-19
\newline
\newline
\hyperref[sec:032473]{\small{< < < PREV SERMON < < <}}
~
\hyperref[sec:index_chronic]{\small{[返順時目]}}
~
\hyperref[sec:031601]{\small{> > > NEXT SERMON > > >}}
\newline
\newline
使徒行傳 16:6-10
\newline
\begin{longtable}{cl}
\hline
\hline
章節 & 經文 (和合本修訂版)\\
\hline
16:6 & \begin{tabularx}{0.7\textwidth}{X} 因為聖靈禁止他們在亞細亞講道，他們就經過弗呂家、加拉太一帶地方。 \end{tabularx} \\ \\ \relax
16:7 & \begin{tabularx}{0.7\textwidth}{X} 到了每西亞的邊界，他們想要往庇推尼去，耶穌的靈卻不許。 \end{tabularx} \\ \\ \relax
16:8 & \begin{tabularx}{0.7\textwidth}{X} 他們就越過每西亞，下特羅亞去。 \end{tabularx} \\ \\ \relax
16:9 & \begin{tabularx}{0.7\textwidth}{X} 夜間，有異象向保羅顯現。有一個馬其頓人站著求他說：「請你過來，到馬其頓來幫助我們！」 \end{tabularx} \\ \\ \relax
16:10 & \begin{tabularx}{0.7\textwidth}{X} 保羅既看見這異象，我們就立即設法往馬其頓去，認為神呼召我們傳福音給那裡的人。 \end{tabularx} \\ \\
[1ex]
\hline
\hline
\end{longtable}
$^{1}$弟兄姊妹早晨.
願神的恩典常與你們同在.
講完四福音.
今天就講《仕途行傳》.
我們看看耶穌升天之後.
教會如何發展.
教會的發展並非一帆風順.
很多波折.
但在神的介入下.
亦成為了很多機遇.
今天我們看看教會發展的一個關鍵時刻.
我們看今天講到的經文.
是在《仕途行傳》第十六章六至第十節.
聖靈既然禁止他們在亞細亞講道.
他們就經過佛雷加加拉太一帶的地方.
到了美西亞的邊界.
他們想要往比推尼去.
耶穌的靈卻不許.
耶穌的靈不許他們.
他們就越過美西亞.
下到特羅亞去.
在夜間有異象現與保羅.
有一個馬其頓人站著求他說.
求你過來馬其頓來幫助我們.
保羅既看見這異象.
我們隨即想要往馬其頓去.
以為上帝召我們傳福音給那裡的人聽.
一般來說.
經文的故事被稱為「馬其頓異象」.
很多機構和仕工都以此為名.
這個故事相信大家都已熟能詳.
今天我講一下這個故事.
背後涉及的一些歷史典故.
以及它的神學意義.
首先要講發生異象的地方.
叫做特羅亞.
一個關鍵的問題是.
特羅亞是一個甚麼地方呢?.
從字源來說.
特羅亞其實是指特洛伊人的地區.

$^{41}$這個地區最重要的城市.
當然就是特洛伊.
在文學上.
很多時特羅亞和特洛伊是同義互用的兩個詞.
互相通用的.
在我的講道裡.
特羅亞和特洛伊這兩個地名都是互用的.
特羅亞有甚麼重要性呢?.
它的重要性就是它的地理位置.
現在很清楚地看得到.
特羅亞或者特洛伊.
特羅亞是指一個地區.
特洛伊是指一個城市.
它位於現在的土耳其.
在歐洲的希臘本土出發.
要通往亞洲.
而最直接最簡單的方法.
就是從水路越過愛琴海.
就會去到亞洲的特羅亞.
簡單來說.
這是歐亞貿易的主要通道.
也因為它是一個最主要的歐亞通道.
它也成為了兵家必爭之地.
有點像現在的霍爾木茲海峽.
它握著能源貿易的咽喉.
它就成為了兵家必爭之地.
特洛伊其實是一個著名的古戰場.
亞洲的帝國想征服歐洲.
或者歐洲的帝國想征服亞洲.
就會在特洛伊一蹶雌雄.
當史諾恆傳提到特洛亞的時候.
當地人一聽.
就會直接聯想到特洛伊這個古戰場.
每個文化都有一些很出名的古戰場.
一提到這些戰場的名字.
其實就不用多說了.
那個古戰場相關的歷史.
就自然會在文化圈裡面的人的腦海裡.
就會自然浮現出來.
譬如在中國的文化圈裡面.

$^{81}$我們一提到赤壁.
不用再多說了.
你就會直接聯想到孫權和劉備.
與曹操大戰.
你會想到很多故事和典故.
孔明草船借箭.
借東風火燒連環船.
你會聯想到大戰之後.
曹操大敗就要拜就華容道.
在那裡遇到關公.
義氣大過天的關公就放了他走.
當然你也會想到蘇東坡的赤壁懷古.
一提到這個地名.
一個古戰場.
不用多說了.
所有東西就會浮現在文化圈裡面的人的腦海裡.
所以經文一提到「特羅亞」.
其實那些人一看到這裡.
他們馬上就會想到「特洛伊」.
而他們又會聯想到「特洛伊戰爭」.
「特洛伊戰爭」是非常有名的.
就是我們所熟悉的「木馬屠城記」的故事.
「木馬屠城記」這場大戰的故事.
其實是一場神與人共同參與的戰爭.
有點像中國「封神榜」的故事.
「封神榜」講到.
武王佛袖其實是神與人共同參戰的.
很多神明在這裡.
在「木馬屠城記」這個故事裡面.
有一個神明的人物不得不說.
他叫做「阿基里斯」.
他是希臘的第一勇士.
刀槍不入.
但是有一個死門.
就是他的腳zoung .
後人稱腳zoung 上的根結.
那個結叫做「阿基里斯結」.
早前我已經斷了這個「阿基里斯結」.
故此我一生都會記得他.
好了.

$^{121}$回到主題.
「特洛伊戰爭」的本質.
當然不是為了爭奪一個美女Helen.
而是歐洲的希臘和亞洲的特洛伊.
為了爭奪貿易的黃金航道.
而爆發的一場地緣戰爭.
結果歐洲贏了.
作為古戰場.
特羅亞或特洛伊令人想起的.
其實就是戰爭.
一提到這個名字.
或者我們說在中國的文化裡面.
一提到「赤壁」.
我們就知道這是王者之戰.
打贏的就能統一歐亞.
一統當時的天下.
成為王者中的王者.
馬其頓二將的核心地點.
就是特羅亞.
或者他的同義詞「特洛伊」.
這段經文背後說的.
其實是王者之戰.
這個故事出現了一個.
有時我們忽略他的人物.
經文說有一個馬其頓人.
這個人是誰呢?.
在講希臘話的文化圈裡面.
一聽到有一個馬其頓人.
其實在他們的文化圈裡面.
他們又會馬上聯想到一個人物.
就是亞歷山大大帝.
馬其頓人亞歷山大大帝.
正如我們在我們的文化圈裡面.
一聽到烏干.
不用再講夏文.
你們就會想到西楚霸王.
一聽到曲阜.
你就會聯想到孔子.
事實上經文提到的地點特羅亞.
其實有一個全名.

$^{161}$聖經沒有寫出來.
他的全名是亞歷山大特羅亞.
這個名字背後原來有亞歷山大這個名字.
其實是特羅亞地區最重要的一個港口城市.
這個港口城市離開特洛伊古城只有25公里.
古代的帝國除了埃及.
全部都是亞洲的帝國.
你看看聖經出現的舊約的帝國.
除了埃及.
不算中國.
因為聖經的地理不包括中國.
全部都是亞洲的帝國.
亞述,巴比倫,波斯.
西方的帝國入侵亞洲的帝國.
首開先河的就是亞歷山大.
他滅了波斯.
一直殺到印度的邊界.
正正就是這個馬其頓人.
一統歐亞.
一統當時的天下.
來到亞洲的時候.
亞歷山大準備進軍波斯.
他先去哪裡呢?.
他就是先去特洛伊.
亞歷山大故意去特洛伊古城.
在那裡獻祭.
他聲稱自己就是希臘那位刀槍不入的勇士.
阿基里士的後裔.
然後他把這場歐洲的文明.
進攻亞洲傳統的帝國波斯.
把這場大戰宣傳為新的特洛伊戰爭.
這種說法大大激勵了當時的希臘士兵士氣.
最後歐洲的力量真的打敗了亞洲的霸主波斯.
講起亞歷山大入侵亞洲.
在第一世紀.
猶太歷史學者約瑟夫的《猶太古史》.
這本書裡有記載.
並且記載了亞歷山大大帝的一件日事.
話說亞歷山大大帝準備進軍亞洲的時候.
順道經過耶路撒冷.

$^{201}$大軍殺入耶路撒冷.
每個人都以為將會屠城.
誰知入城後.
亞歷山大大帝見到耶路撒冷的大祭司Jedua.
竟然向大祭司致敬.
並且讚頌大祭司的冠冕上.
寫著神的名字「耶和華」.
亞歷山大大帝讚頌耶和華.
記載在歷史的古籍裡.
當時的人都很驚訝.
亞歷山大大帝為何會這樣做呢?.
原來話說亞歷山大大帝還在馬其頓的時候.
他很猶疑是否真的要進軍亞洲呢?.
因為這場是世紀決戰.
在夢裡面.
有一個人向他顯現.
叫他勇敢一些.
渡過愛琴海.
這個人在夢中還承諾.
會帶領他的軍隊奪取波斯.
而這個夢中人所穿的衣服.
正正就是他現在眼見的.
耶路撒冷大祭司的那件衣服.
亞歷山大大一眼就認得出.
我做夢的那個就是他了.
他認為正正就是這個大祭司所敬拜的神.
耶和華幫助他得勝.
在這個歷史的日史裡面.
亞歷山大這個馬其頓人.
認為是耶和華幫助他.
戰勝了當時的霸主波斯.
將這些故事放回聖經做一些對比.
這位馬其頓人.
在夢裡面看到一個人指示他.
進軍亞洲.
一個未知的領域.
然後那個人承諾會幫助他.
一統當時的天下.
現在當書寫《仕途行傳》的時候.
那個情節就是保羅在夢中.

$^{241}$見到一個馬其頓人.
這個馬其頓人求他去歐洲的馬其頓.
我們相信這個技術.
這個經文的潛台詞就是.
保羅的福音將會由亞洲橫跨到歐洲.
並且在耶和華的說領下.
將會一統天下.
關於這段經文有兩點值得一提.
第一就是在《仕途行傳》裡面.
一般來說.
說到「聖靈」的時候.
就會直接用「聖靈」這個字.
大家看《仕途行傳》.
你經常看到「聖靈」「聖靈」.
但是你看看這段經文.
是否很有趣呢?.
刻意用了「耶穌的靈」來稱呼「聖靈」.
我們不知道為什麼.
不過我們有一個推論.
也很合理的.
或者這正是因為要對應亞歷山大的夢境.
亞歷山大的夢境.
是認到那個人的.
那個人是耶和華的僕人.
有面目.
亞歷山大是認到他的.
如今來到這段經文.
或者經文很想對比那個故事.
就說這個「聖靈」是會被人認得出的.
你要給他一個認得出的面貌.
故此很特別的.
在《仕途行傳》這裡.
就說這個是「耶穌的靈」.
有一個面貌.
於是可以認得到.
是他了.
是耶穌.
然後經文的第二個地方值得一提的就是.
這個馬其頓人是求保羅去幫助他們.
如果這個馬其頓人是暗指亞歷山大大帝的話.

$^{281}$這個說法就是一反亞歷山大大帝征服其他人的時候的那種霸氣.
亞歷山大一統歐亞是靠武力 如今一統亞歐是靠福音.
真正能夠讓天下歸一的是神恩典的福音.
沒有任何的霸力.
只有互相的幫助.
只有愛的共融.
接著我想說的是特洛伊.
其實和羅馬有一個很複雜的淵源.
話說特洛伊戰爭被毀了之後.
皇族的後人.
艾利亞斯·阿利亞斯.
他從特洛伊一直往西走.
走了去歐洲.
然後他成為了後來一對很出名的.
孿生兄弟的祖先.
這對孿生兄弟很出名的.
他是被狼養大的.
這對孿生兄弟後來建立了羅馬城.
其實在羅馬歷史的敘事中.
羅馬其實就是特洛伊的賢族.
你可能覺得很奇怪.
羅馬當時那麼偉大的一個帝國.
怎麼會認祖歸宗於特洛伊戰爭的失敗者呢?.
特羅亞是輸了的.
一個當時的霸主.
怎麼會認祖歸宗於特羅亞呢?.
這當然有很多的分析.
其中有可能的就是.
羅馬人要宣稱他們是源於一個不同於希臘文明.
卻又與希臘文明平起平坐的血統.
大家要明白.
羅馬興起的時候是希臘文明.
如何對抗它呢?.
要有一個敘述.
就是我的祖先源自一個文明.
是與希臘文明能夠同時遠古平起平坐的.
而當後來羅馬要征服希臘的時候.
就被宣傳成為特洛伊的後人.
要對希臘人進行一個報復的大戰.
這樣就給了羅馬對希臘的統治有一種合法性.

$^{321}$說了那麼久.
簡單來說.
談論特洛伊或者談論特羅亞.
終極來說.
其實就是談論羅馬.
從特洛伊或者從特羅亞往西行.
會越過馬其頓.
然後最後就會直指羅馬.
說起羅馬.
又不得不提凱撒大帝.
話說當年凱撒大帝很厲害.
打遍整個高盧的地區.
現在大概是法國的地區.
他的兵權極盛.
羅馬的元老院很怕那麼厲害的凱撒.
威脅到自己的利益.
於是就下令要解除他的兵權.
主前第49年.
凱撒連同他的大軍.
來到一條河.
叫做盧比孔河.
Rubicon River.
來到這條河邊.
這條河為什麼很重要呢?.
這條河是一條界線.
一條界線.
任何的軍隊未經授權.
越過這條界線.
就視為背叛羅馬.
放下兵權.
回到羅馬城.
凱撒預計到元老院會處死他.
但是不放下兵權.
而越過這條界線.
其實意味著與羅馬帝國主力正面交鋒.
何去何從呢?.
就在盧比孔河的河邊.
凱撒心靈交戰.
那一晚.
過還是不過?.

$^{361}$交兵權還是不交兵權?.
凱撒到底在想什麼呢?.
這成為了很多文學家發揮想像的題材.
這些是心靈交戰的重大時刻.
有一個傳統的說法.
當凱撒猶疑之際.
他忽見異象.
神明指示他.
帶著大軍過河.
原來夜晚凡是異象發夢.
就會牽涉到所謂大戰.
這些是古時的想法.
大戰前夕沒有神明的指引.
千萬不要動.
凱撒在夜晚.
見到神明顯示.
指示他帶著大軍過河.
於是凱撒就說出那句.
命垂千古的話.
The die is cast.
即是命運的骰子已經摘下了.
一切就交由命運決定.
於是他就帶著大軍.
越過界線.
直至羅馬.
說到這裡我們回到《朝道行傳》.
其實整段經文.
一段你們很熟悉的經文.
其實他到底想說甚麼呢?.
其實在文學上.
特羅亞即是特洛伊.
一聽到特洛伊.
大家都知道這是一個古戰場.
是王者爭雄的地方.
在這個地方有很多祭壇.
任何大戰作戰之前.
任何一方都要求問神明.
要求神指引,求神幫助.
在這裡出現了很多重要的戰役.
其中一個出名的戰役.

$^{401}$就是特洛伊戰爭.
西方的馬其頓戰勝了亞洲的特洛伊.
然後馬其頓的亞歷山大大帝.
曾經也在這裡獻祭.
求神明指引他.
在整個歷史的日史記載裡面.
真正在夢裡面指引亞歷山大大帝的.
就是他在馬其頓時候夢見的耶和華的僕人.
這個僕人預告.
他的軍隊必定能夠戰勝亞洲的波斯.
在這麼簡單的故事裡面.
已經預示了真正掌管歷史的.
是耶和華.
馬其頓人亞歷山大以武力征服天下.
但是真正統一天下的.
《使徒行傳》講得很清楚.
是福音.
是耶穌的福音.
馬其頓人真正需要的.
不是武力而是和平的福音.
最後特洛伊繼續講下去.
原來特洛伊就是羅馬的根源.
談到特洛伊就不能不談羅馬.
而我們剛才也說了.
凱撒大帝在夢裡面.
得到神明的指引.
越過他不應該超越的界線.
直指羅馬.
今日《使徒行傳》.
保羅得到耶穌的靈的指引.
他越過一個界線.
這個界線是當時的保羅他們認為.
是不應該越過的.
我們留在亞洲繼續傳福音就可以了.
就不用去歐洲了.
但是得到耶穌的靈的指引.
他越過了亞歐的界線.
到達了福音前所未到的地區.
而從保羅的位置.
特羅亞或者特洛伊.

$^{441}$穿過馬其頓.
再往西行的話.
就直指羅馬了.
在保羅的時代.
羅馬是一統地中海天下的帝國.
他以武力主宰著歷史的軌跡.
《使徒行傳》卻是說.
在萬有之上.
在羅馬帝國之上.
其實是耶穌在掌權.
而耶穌才是真正的萬王之王.
這個王者之戰.
最後能勝利的.
只有萬王之王的耶穌.
是祂指引著他的僕人.
突破人間的界限.
以人愛的福音.
直取羅馬.
真正能夠令天下一家的.
不是靠武力.
只能靠福音.
到《使徒行傳》結尾的時候.
所記載的事就是.
保羅就身處羅馬.
全港天國.
真正的國度.
真正一統天下的帝國.
是人間真正美善的天國.
上帝的國度.
而這就是我們能夠在歷史裡所盼望的.
願神賜福給大家.
也願榮耀歸給聖父,聖子,聖靈的名.
阿們.
\newpage

\allsectionsfont{\centering}

\setlength\parindent{0pt}
\setlength{\columnsep}{1.25em}
\setlength{\parfillskip}{0pt}
\setlength{\tabcolsep}{1em}
\raggedbottom

\pagenumbering{gobble}


\newfontfamily\leftfont[Path=../fonts/fell_french_canon/, Ligatures=TeX, ItalicFont=IMFeFCit29C.otf, BoldFont=AveriaLibre-Bold.ttf]{IMFeFCrm29C.otf}
\newfontfamily\leftcitationfont[Path=../fonts/frankruehl/]{FrankRuehlCLM-Medium.ttf}
\newfontfamily\centerfont[Path=../fonts/garamond/, Ligatures=TeX, ItalicFont=EBGaramond-SemiBoldItalic.ttf]{EBGaramond-SemiBold.ttf}
\newfontfamily\rightfont[Path=../fonts/averia/, Ligatures=TeX, ItalicFont=AveriaLibre-RegularItalic.ttf, BoldFont=AveriaLibre-Bold.ttf, BoldItalicFont=AveriaLibre-BoldItalic.ttf]{AveriaLibre-Light.ttf}
\newfontfamily\rightcitationfont[Path=../fonts/rashi/]{Mekorot-Rashi.ttf}
\definecolor{hcolor}{HTML}{D3230C}
\definecolor{rcolor}{HTML}{D36F0C}
\newcommand{\chfont}[1]{\centerfont{\huge\textcolor{hcolor}{#1}}}
\newcommand{\leftcitation}[1]{\leftcitationfont{\Large\textcolor{hcolor}{#1}}}
\newcommand{\rightcitation}[1]{\rightcitationfont{\normalsize\textcolor{rcolor}{#1}}}
\newfontfamily\flowerfont[Path=../fonts/fell_flowers/]{IMFeFlow2.otf}

\begin{sloppypar}

\chapter*{\chfont{編按結語}}

\columnratio{0.5,0.5}\begin{paracol}{2}

\fontsize{11}{13}\leftfont \Large \leftcitation{א} \leftfont 余少好文．宏志博覽群書而不忘．善存藏經籍文獻備後時之用。\leftcitation{ב} \leftfont 歸主年時．受友所薦．聞道網海．\switchcolumn\fontsize{11}{13}\rightfont \Large \leftcitation{ח} \rightfont 有見粵道之危．國之封講道千言亦將就至．急之何則為？\leftcitation{ט} \rightfont 嘗聞猶太者之傳承．在其力守口述之

\end{paracol}


\columnratio{0.32,0.32,0.32}\begin{paracol}{3}

\fontsize{11}{13}\leftfont \Large 尤以吳約翰遜者 \switchcolumn[2]\fontsize{11}{13}\rightfont \Large 統．以煉千載不

\end{paracol}

\columnratio{0.32,0.32,0.32}
\begin{paracol}{3}\fontsize{11}{13}\leftfont \Large 為重．其載上之粵語講道緩緩入耳．收之藏其音頻．善妥整存．反復而嚼．受益無窮。\leftcitation{ג} \leftfont 我城我國既限．歷一四一九之不測．肺疫延年．信徒靈長屢受圍創．神州燈臺數盡指日可待．粵道之求與日俱增。\leftcitation{ד} \leftfont 觀乎社、經、法、媒、言、信、網之地．愈趨受鋤．自翔不果．授受壓力．粵道聖言亦愈漸艱難。\leftcitation{ה} \leftfont 況崇基例乎．學苑講道屢逆權勢者．其言末強受壓．舊章盡刪以存其身。講道釋數失傳．徒嘆奈何。

\switchcolumn

\fontsize{11}{13}\centerfont 
\begin{tikzpicture}
    \node (0,0) [xshift=-0.10cm, yshift=-1.0cm, opacity=0.10]{\includegraphics[width=0.30\textwidth]{../ot_frontcover.png}} ;
    \node (0,0) [xshift=+0.20cm, yshift=+2.0cm, opacity=0.10]{\includegraphics[width=0.20\textwidth]{../christ_on_cross.png}} ;
\end{tikzpicture}
\Large 

\leftcitation{ס} \centerfont 詩百又廿七載：
\leftcitation{ע} \centerfont 非耶和華建屋宇.則匠人之經營徒.
\leftcitation{פ} \centerfont 非耶和華衛城邑.則守者之儆醒徒.
\leftcitation{צ} \centerfont 余獻是卷予華人社區．願為福音流通之器．願獻斯微材為祭榮耀上帝．
\leftcitation{ק} \centerfont 阿門

\switchcolumn

\fontsize{11}{13}\rightfont \Large 滅．時越次聖殿期及當今。\leftcitation{י} \rightfont 猶太者力廣納之．筆錄以卷軸．便以傳、閱、頌、攜、守、鎖、抄、譯、釋、編，得書塔木德、密示拿等經傳．家喻戶曉．傳流若芳。\leftcitation{כ} \rightfont 猶太者文以載道．傳其口述．今我輩粵道之傳應當作如是．遂力行粵音識辨之法．載言載道．以盡忠傳粵道以待興。\leftcitation{ל} \rightfont 蒙下賜恩惠．無畏海量字音文書．既馭上帝之道．今廣及粵語講道．重駛編程之技．匯導粵音遂字稿．重塑講道現場．以傚猶太卷軸之舉便以傳流。\leftcitation{מ} \rightfont 是卷乃粵音口述傳之屬．莫通華文白話之語．

\end{paracol}

\columnratio{0.5,0.5}
\begin{paracol}{2}\fontsize{11}{13}\leftfont \Large \leftcitation{ו} \leftfont 斯殺一違儆百逆．既禁壓之．我輩聞風無奈．在所難免。\leftcitation{ז} \leftfont 另有異人例乎．以版權之名．脅網絡頻道之舉．同授礙予粵道之存流。

\switchcolumn

\fontsize{11}{13}\rightfont \Large 惟待後繼來者之傚．以譯釋傳之於神州華文地。\leftcitation{נ} \rightfont 今能排程驅馭圖靈以編彙文檔，其碼長共數千千亦無逢大礙．全蒙上帝保守。

\end{paracol}



\columnratio{1}\begin{paracol}{1}

\fontsize{11}{13}\rightfont \Large
~~~~~~~~~~~~~~~~~~~~~~~~~~~~~~~~~~~~~~~~~~~~~~~~~~~~~~~~~~~~~~~~~~~~~~~~~~~~~~~\leftcitation{ר} \rightfont 二零二三年二月一日

~~~~~~~~~~~~~~~~~~~~~~~~~~~~~~~~~~~~~~~~~~~~~~~~~~~~~~~~~~~~~~~~~~~~~~~~~~~~~~~\leftcitation{ש} \rightfont 米迦勒

~~~~~~~~~~~~~~~~~~~~~~~~~~~~~~~~~~~~~~~~~~~~~~~~~~~~~~~~~~~~~~~~~~~~~~~~~~~~~~~\leftcitation{ת} \rightfont 書於香港

\end{paracol}

\end{sloppypar}
\end{document}
